

\maketitle

\tableofcontents

%% ---- begin inlined file: body ----
\clearpage
\setcounter{section}{-1}
\refstepcounter{section}
\section*{Chapter~\thesection\quad Conventions, numbering and glossary of symbols\addcontentsline{toc}{section}{\protect\numberline{\thesection}Conventions, numbering and glossary of symbols}\label{ch:0}}
%% ---- begin inlined file: chapters/c00 ----
The book is in three parts. Its chapters are numbered consecutively from $0$ and run straight
through the three parts. Within each chapter a single counter is shared by theorems, propositions,
lemmas, corollaries, definitions and remarks. A statement is therefore numbered by its chapter and
its place in that count, whatever its kind. Equations are numbered within the chapter in the same
two-part form. The clauses of the main theorem stated in Part~I are lettered (0)--(vi),
with sub-letters for their parts, and are printed in the order (0), (i), (iii), (ii), (iv), (v),
(vi), clause (iii) before (ii) because the proof of (ii) uses (iii).
The hypotheses of the theorem are (H1)--(H7), and the hypothesis of clause (vi) is the far-sphere
flux hypothesis $(\mathrm{IR}\text{-}\mathrm{fl})_{\omega_0}$.

Cross-references within the book are by number only, never by page. Ba{\l}aban's papers and the
other published works are cited by their bibliography keys, with the number of the equation,
theorem or page meant.

A numbered statement of Part~II appears in one of the following forms, indicated after its
name in its heading.
\begin{itemize}
\item No indication: the statement is proved in full and the proof follows; a statement proved
from other named statements lists them, by number or by name, as its explicit hypotheses.
\item A citation with its locus: a theorem of a published paper, quoted by its statement, whose
proof is that of the cited paper and is not reproduced.
\item The words ``printed without proof in'' with the citation and locus, followed by ``cited as
printed'': a sentence that a published paper states without proof. The sentence is quoted, set
off as a displayed quotation, and used as printed.
\item ``Stated; proof incomplete at the step marked $(\ast)$'': the proof is carried as far as it
goes, and the step where it stops is marked.
\item The words ``Proof outstanding.'' followed, in place of a proof, by a paragraph headed
``Route.'' that describes in mathematical terms the route a proof would take.
\item ``Model'', for a few statements: the first sentence says what the model is, and the statement
is proved in that model, not for Ba{\l}aban's renormalization step.
\end{itemize}
Wherever a proof appeals to a statement not proved in full, it says so at the point of use.

In this book a statement whose proof is printed in full carries no status sentence. Part~III
gives the standing of a statement by the first four phrases below and, for a sentence of
Ba{\l}aban's papers used as printed, by the fifth. These phrases have the
following meanings.
\begin{itemize}
\item ``proved'': the proof is complete and is printed with the statement in this book, or is that
of the reference cited with it; only a statement with this status is counted as proved.
\item ``proved as an implication from $\langle$named statements$\rangle$'': a complete proof that
the named statements imply the statement is given here or in the reference cited; the phrase is
used only when some named statement is not itself counted as proved, and each such statement is
named with its own status.
\item ``stated; proof incomplete'': the statement is formulated precisely and a proof is given in
part, but a named piece of that proof, not isolated as a hypothesis, is missing, and the statement
is not counted as proved.
\item ``not proved'': no complete proof is given here or in a reference cited with the statement,
and none is asserted to exist.
\item ``printed without proof; cited as printed'': a statement printed without proof in one of
Ba{\l}aban's papers and used here as printed, with a citation of the place where it is stated;
it is not counted as proved.
\end{itemize}
Besides the forms and phrases above, a statement of Part~II may also appear in either of two
further forms, in neither of which it is counted as proved. The first is a remark that carries
the number under which a lemma, proposition, theorem or corollary is cited and reads ``The
statement of this lemma is not written out here; parts of its proof are printed below.'' or
``The statement and proof of this lemma are not written out here.'', with ``proposition'',
``theorem'' or ``corollary'' in place of ``lemma'' as the case may be; the printed parts of the
proof then follow the remark, as lemmas and as proofs whose headings name that statement. The
second is a proposition headed ``Proof outstanding'' whose text reads ``The statement of this
proposition is not written out here.'' and which is followed by its route. Part~III records as
``proof outstanding'' a statement in the fifth form of the first list above, as ``model; proved
in its model'' one in its last form, and as ``cited as printed at its use site'' a sentence of
Ba{\l}aban's papers quoted where it is used, and lists a statement without a number by its
content.

All constants are symbolic. There is one existential constant $L_0 = L_0(N)$, depending only on
$N$, where $\mathrm{SU}(N)$ is the gauge group. It is bound once and never given a value. Every
condition on the blocking factor $L$ is a lower bound, and none is numerical. Every hypothesis is
one-sided.

The glossary of symbols is the list of Definition~\ref{0.2:not-glossary}, which follows. The last
chapter, in Part~III, indexes every statement of Part~II whose proof is outstanding, incomplete or
printed without proof.

\begin{definition}[Glossary]\label{0.2:not-glossary}
The following list records the symbols used in the displayed statements of the book. Each entry
gives the symbol and its meaning where this list determines it, and otherwise the kind of object it
denotes and the part of the theory in which it occurs; where the symbol is fixed in a particular
chapter, the entry names that chapter. The entries are grouped by subject.

\medskip
\noindent Lattice and torus.
\begin{description}
\item[$d=4$] the dimension of the lattice.
\item[$N$] the integer $N$ of the gauge group $\mathrm{SU}(N)$.
\item[$\mathrm{SU}(N)$] the gauge group.
\item[$\operatorname{tr}$] the trace on $N\times N$ matrices, normalised by
$\operatorname{tr}\mathbf{1}=1$.
\item[$L$] the blocking factor, an odd integer (Definition~\ref{1.1:def-lattice}).
\item[$L_0=L_0(N)$] the one existential constant of the book, the constant of hypothesis (H2) of
Definition~\ref{2.2:not-hypotheses}. It is bound once, by ``there exists'',
and is never given a value.
\item[$m$] the index of the torus (Definition~\ref{1.1:def-lattice}).
\item[$\mathbb{T}_m$] the torus of side $2L^m$ (Definition~\ref{1.1:def-lattice}).
\item[$K$] the number of renormalisation steps; the bare lattice spacing is $L^{-K}$
(Definition~\ref{1.1:def-lattice}).
\item[$K(g_0,g)$] the first-passage number of steps, as a function of the bare coupling $g_0$ and the
reference coupling $g$ (Definition~\ref{1.2:def-flow-constants}).
\item[$k$] the level, $0\le k\le K$.
\item[$s$] the depth below the cutoff; depth $s$ is level $K-s$.
\item[$p$] a plaquette (Definition~\ref{1.1:def-lattice}).
\item[$\partial p$] the boundary of the plaquette $p$.
\item[$x(p)$] the barycentre of the plaquette $p$.
\end{description}

\medskip
\noindent Action and couplings.
\begin{description}
\item[$U$] the bond variables (Definition~\ref{1.1:def-wilson-action}).
\item[$U(\partial p)$] the ordered product of the bond variables around the boundary $\partial p$
of the plaquette $p$ (Definition~\ref{1.1:def-wilson-action}).
\item[$A(U)$] the action (Definition~\ref{1.1:def-wilson-action}),
\begin{equation*}
A(U)=\sum_p\bigl[1-\operatorname{Re}\operatorname{tr}U(\partial p)\bigr].
\end{equation*}
\item[$g_0$] the bare coupling.
\item[$g$] the reference coupling.
\item[$g_k$] the coupling at level $k$ (Definition~\ref{1.2:def-couplings}).
\item[$x_k$] the inverse square of the coupling at level $k$, $x_k=g_k^{-2}$
(Definition~\ref{1.2:def-couplings}).
\item[$\beta_{k+1}$] the coupling increment at step $k+1$ (Definition~\ref{1.2:def-couplings}).
\item[$\lambda^{\sharp}$, $\underline{\lambda}$] constants fixed in
Definition~\ref{1.2:def-flow-constants} and related by
\begin{equation*}
\lambda^{\sharp}\ge\underline{\lambda}=c_0/4 .
\end{equation*}
\item[$B^{\sharp}$, $\bar{B}$] constants fixed in Definition~\ref{1.2:def-flow-constants} and
related by
\begin{equation*}
B^{\sharp}\le\bar{B}=C_\beta+1 .
\end{equation*}
\item[$c_5$] the constant $c_5=2c_2$, fixed with the constant $c_2$ in
Definition~\ref{1.2:def-flow-constants}.
\item[$\lambda_U$] a constant of the coupling flow.
\item[$g_-(g)$] the function of the reference coupling $g$
(Definition~\ref{1.2:def-flow-constants}) given by
\begin{equation*}
g_-(g)=\bigl(g^{-2}+B^{\sharp}\bigr)^{-1/2}.
\end{equation*}
\item[$m_{\mathbb{T}}(g')$] the torus exponent (Definition~\ref{1.2:def-flow-constants})
\begin{equation*}
m_{\mathbb{T}}(g')=11+m'+\log_L R(g'),
\end{equation*}
where $m'$ is the exponent in $M=L^{m'}$ among the auxiliary parameters $\mathfrak{p}$ below, and
$R(g')$ is the function of the coupling fixed with $m_{\mathbb{T}}$ in that definition; this $R$
is distinct from the ball radius $R$ below.
\item[$\mathfrak{p}$] Ba{\l}aban's auxiliary parameters $\kappa$, $\kappa_1$, $M_1$, $M_2=L^2M_1$,
$M=L^{m'}$, $A_0$, $A_1$, $C_0$, $C_1$, the radii $\alpha$, $C_\beta$, $\varepsilon^{I}_0$,
$\varepsilon^{I}_1$, $p_0$, $q$, and the further parameters of his construction; they are chosen
in one fixed order, after $L$ (Definition~\ref{2.1:not-prefix}).
\item[$\gamma$] a parameter.
\item[$\gamma_H$] a parameter, distinct from $\gamma$.
\item[$\gamma_J$] a parameter, distinct from $\gamma$.
\item[$\gamma_W$] a parameter, distinct from $\gamma$.
\item[$\gamma_U$] a parameter, distinct from $\gamma$.
\item[$\varrho$] a parameter.
\item[$\varrho_2$] half the parameter $\varrho$, $\varrho_2=\varrho/2$.
\item[$\vartheta$] the second parameter of the class $\mathfrak{T}(\mathsf{d};\vartheta)$ of test
functions listed under the heading ``The witness'' below (Chapter~\ref{ch:8}).
\end{description}

\medskip
\noindent Ba{\l}aban's step.
\begin{description}
\item[$U_k$, $V_k$] the block fields at level $k$, produced by the block-averaging map
(Definition~\ref{1.5:def-block-average}).
\item[$\rho_k$] the effective density at level $k$ (Definition~\ref{1.5:def-densities}).
\item[$\chi_k$] a function at level $k$ of Ba{\l}aban's step.
\item[\normalfont\itshape small-field and large-field regions] the small-field and large-field
regions at level $k$ of Definition~\ref{1.5:def-densities}; the large-field operation acts on the
latter. The words used of large-field regions are listed at the end of this glossary.
\item[$\mathbf{R}$, $\mathbf{B}$, $\mathbf{T}$, $\mathbf{T}'_k(X)$] the large-field operations, as
Ba{\l}aban defines them (Definition~\ref{1.5:def-densities}).
\item[$X$, $Z$] large-field regions.
\item[$d_j(X)$] a quantity attached to a large-field region $X$ and a step $j$; here the letter $d$
does not denote the dimension.
\item[$d_{j,Z}(X)$] a quantity attached to a large-field region $X$, a step $j$ and a second
large-field region $Z$.
\item[$d^{\theta}(X)$] a quantity attached to a large-field region $X$ and a parameter $\theta$.
\item[$E_+$, $E_-$] two quantities of the analysis of Ba{\l}aban's step.
\item[$N_k$] an object at level $k$ of Ba{\l}aban's step; distinct from the integer $N$ of
the gauge group.
\item[$h$] a history.
\item[$F^{(h)}$] an object indexed by the history $h$.
\item[$\tau_h$] the term indexed by the history $h$ in the representation of the large-field terms.
\end{description}

\medskip
\noindent Observables.
\begin{description}
\item[$\mathcal{O}(f)$] the action density smeared with the test function $f$
(Definition~\ref{1.3:def-curvature-field}),
\begin{equation*}
\mathcal{O}(f)=\sum_p f\bigl(x(p)\bigr)\bigl[1-\operatorname{Re}\operatorname{tr}U(\partial p)\bigr].
\end{equation*}
\item[$\|f\|_{\sharp}$] the norm (Definition~\ref{1.3:def-curvature-field})
\begin{equation*}
\|f\|_{\sharp}=\max\bigl\{\|f\|_\infty,\;40M\operatorname{Lip}f\bigr\}.
\end{equation*}
\item[$\mathsf{d}$] the separation of the supports of the test functions
(Definition~\ref{1.3:def-curvature-field}).
\item[$R$] the radius of a ball.
\item[$Z(z)$] a function of the variable $z$; distinct from the large-field regions $Z$ above.
\item[$S^T_{n;K,m}$] the connected Schwinger functions, $n\ge2$, with $K$ renormalisation steps on
the torus $\mathbb{T}_m$.
\item[$S_q$] functions indexed by $q$.
\item[$C_n(\mathsf{d})$] a quantity indexed by $n$ that depends on the separation $\mathsf{d}$.
\item[$C^{\infty}_n(\mathsf{d};R)$] a quantity indexed by $n$ that depends on the separation
$\mathsf{d}$ and on the ball radius $R$.
\item[$\ell(n,R,g)$] a quantity depending on $n$, on the ball radius $R$ and on the reference
coupling $g$.
\item[$\mathsf{Q}$] a quantity attached to the observables.
\item[$w$] the complex source
\begin{equation*}
w=\sum_i z_if_i ,
\end{equation*}
with complex numbers $z_i$ and test functions $f_i$.
\item[$\zeta$, $\zeta_k$] quantities attached to the complex sources; $\zeta_k$ is indexed by the
level $k$.
\item[$\zeta'_k$] the running shift (Chapter~\ref{ch:5}).
\item[$b_{k+1}$] a quantity attached to the complex sources at step $k+1$.
\end{description}

\medskip
\noindent Limits.
\begin{description}
\item[$S^T_{n;m}$, $S^T_{n;\infty}$] the limit connected Schwinger functions, on the torus of
index $m$ and in infinite volume respectively; the latter are fixed with the infinite-volume
theorem in Chapter~\ref{ch:7}.
\item[$S^{(\hat{w})}$] the limit Schwinger functions labelled by the datum $\hat{w}$ listed under
the heading ``Uniqueness'' below.
\item[$W_4$, $W_3$] the hyperoctahedral groups in four and three dimensions.
\item[$\mathrm{O}(4)$] the orthogonal group in four dimensions.
\item[$\mathcal{H}^{\mathrm{curv}}$] the Hilbert space of the curvature observables.
\item[$\Omega$] the vacuum vector.
\item[$H$] the Hamiltonian.
\item[$e^{-tH}$] the semigroup generated by $H$.
\end{description}

\medskip
\noindent The witness.
\begin{description}
\item[$b_0$] the constant of the two-point witness construction given by
\begin{equation*}
b_0=\frac{N^2-1}{2}.
\end{equation*}
\item[$\mathcal{A}$] the free lattice field.
\item[$Q$] the linearised block-averaging map.
\item[$F_B$] a quantity of the two-point witness construction, fixed in Chapter~\ref{ch:8}.
\item[$\mathcal{K}_s(f_1,f_2)$] a kernel at depth $s$ evaluated on the test functions $f_1,f_2$,
fixed in Chapter~\ref{ch:8}.
\item[$\mathfrak{T}(\mathsf{d};\vartheta)$] a class of test functions depending on $\mathsf{d}$ and
$\vartheta$, fixed in Chapter~\ref{ch:8}.
\item[$s(\mathsf{d})$, $s_W(g)$, $s_*$] depths of the two-point witness construction, fixed in
Chapter~\ref{ch:8}.
\item[$r(s)$] the extraction depth of clause (iv) of Theorem~\ref{3.1:thm-main}, as a function of
the depth $s$, fixed in Chapter~\ref{ch:8}.
\item[$k_W$] a level of the two-point witness construction, fixed in Chapter~\ref{ch:8}.
\item[$C_W(s)$] a constant of the two-point witness construction depending on the depth $s$,
fixed in Chapter~\ref{ch:8}.
\item[$c_{\mathcal{K}}$, $C_{\mathcal{K}}$] two constants attached to the kernel
$\mathcal{K}_s$, fixed in Chapter~\ref{ch:8}.
\item[$R_K$] a quantity of the two-point witness construction depending on the number $K$ of
renormalisation steps, fixed in Chapter~\ref{ch:8}.
\item[$\bar{x}_s$, $\underline{x}_s$] two quantities at depth $s$ of the two-point witness
construction, fixed in Chapter~\ref{ch:8}.
\item[$u_\gamma(g)$] a function of the reference coupling $g$ depending on the parameter $\gamma$,
fixed in Chapter~\ref{ch:8}.
\end{description}

\medskip
\noindent Uniqueness.
\begin{description}
\item[$X$] the inverse square of the reference coupling, $X=g^{-2}$; this use of the letter is
distinct from the large-field regions $X$ above.
\item[$\hat{x}$] the canonical comparison point (Chapter~\ref{ch:9}).
\item[$t$] the displacement of $X$ from the canonical comparison point, $t=X-\hat{x}$; this use of
the letter is distinct from the parameter $t$ of the semigroup $e^{-tH}$.
\item[$\rho$] the contraction constant of clause (v-d) of Theorem~\ref{3.1:thm-main}, fixed in
Chapter~\ref{ch:9}; distinct from the effective densities $\rho_k$ above.
\item[$a_n$] a sequence of the uniqueness argument, indexed by $n$ (Chapter~\ref{ch:9}).
\item[$\hat{w}$] the datum labelling the limit Schwinger functions $S^{(\hat{w})}$, fixed with the
uniqueness theorem in Chapter~\ref{ch:9}.
\item[$\upsilon_w$] a quantity of the uniqueness argument, indexed by $w$ (Chapter~\ref{ch:9}).
\item[$\mathfrak{n}_n$] a quantity of the uniqueness argument, indexed by $n$
(Chapter~\ref{ch:9}).
\item[$\Upsilon_n$] a quantity of the uniqueness argument, indexed by $n$ (Chapter~\ref{ch:9}).
\item[$N_n$] a quantity of the uniqueness argument, indexed by $n$ (Chapter~\ref{ch:9}); distinct
from the integer $N$ of the gauge group and from $N_k$ above.
\item[$\hat{\mathsf{N}}_n$] a quantity of the uniqueness argument, indexed by $n$
(Chapter~\ref{ch:9}).
\item[$i''$] the re-filing level (Chapter~\ref{ch:9}).
\end{description}

\medskip
\noindent $\mathrm{O}(4)$.
\begin{description}
\item[$\Theta_{\nu\sigma}$] the stress-tensor density (Chapter~\ref{ch:10}).
\item[$O_1$] the one hypercubic scalar of dimension six (Chapter~\ref{ch:10}),
\begin{equation*}
O_1=\sum\operatorname{tr}\bigl(D_\mu F_{\mu\nu}\,D_\mu F_{\mu\nu}\bigr),
\end{equation*}
where $F_{\mu\nu}$ denotes the curvature and $D_\mu$ the covariant derivative; the range of the sum
is fixed in Chapter~\ref{ch:10}.
\item[$\mathcal{R}$] an object of the $\mathrm{O}(4)$ argument, fixed in Chapter~\ref{ch:10}.
\item[\normalfont\itshape shell functional] a functional of the $\mathrm{O}(4)$ argument, fixed in
Chapter~\ref{ch:10}.
\item[$\omega_0$] the parameter carried by the far-sphere flux hypothesis
$(\mathrm{IR\text{-}fl})_{\omega_0}$, the hypothesis under which clause (vi) of
Theorem~\ref{3.1:thm-main} is stated, fixed in Chapter~\ref{ch:10}: the flux of the lattice
rotation current through far spheres vanishes along one sequence of radii. A sufficient condition,
displayed beside the hypothesis in Theorem~\ref{3.1:thm-main} and in Chapter~\ref{ch:10}, is that
the connected
correlation of one far-away stress-tensor insertion with the observables decays faster than the
inverse fourth power of the distance of the insertion from the observables, uniformly in the
cutoff.
\end{description}

\medskip
\noindent Two decay exponents, the degree function, and the word \emph{non-trivial}.
\begin{description}
\item[$\varkappa$] the decay rate of the two-point witness theorem, defined in
Chapter~\ref{ch:8} by $\varkappa=(\kappa-c_{\mathrm{an}})/(2c_dM)$, with constants
$c_{\mathrm{an}}$ and $c_d$ fixed there; it is distinct from $\kappa'$ below.
\item[$\kappa'$] the decay exponent of the kernel estimate fixed at Stage~0 of the
order of choice of all constants (Definition~\ref{2.3:not-stages-first}); the letter denotes only
this exponent. The inequality it
satisfies opens the reference window of Definition~\ref{1.2:def-flow-constants};
that inequality is not a hypothesis of Theorem~\ref{3.1:thm-main}. The letter $\kappa'$ is never
used for the decay rate $\varkappa$.
\item[$p_0(\cdot)$] Ba{\l}aban's degree function of the coupling,
\begin{equation*}
p_0(g')=A_0\bigl(\log g'^{-2}\bigr)^{p_0},
\end{equation*}
the same letter denoting the function and its fixed exponent $p_0$, one of the auxiliary parameters
$\mathfrak{p}$, subject to the one-sided condition $p_0\ge5r$ imposed by Ba{\l}aban. When the
function is meant, we say ``the degree function $p_0(\cdot)$''.
\item[\normalfont\itshape non-trivial] said of a continuum limit. The observable is the action
density $\operatorname{tr}F\cdot F$ smeared with a test function, in the normalisation in which the
coupling multiplies the action and not the field. The connected two-point function of two such
observables supported in balls of radius $d$ (here $d$ denotes a length, not the dimension) at
mutual distance of order $d$, normalised by the sup
norms of the test functions, is the fourth power of an effective coupling at scale $d$ times a
dilation-invariant kernel. A limit is non-trivial if this quantity is not identically zero and tends
to zero, as $d\to0$, more slowly than every power of $d$. This excludes that all connected
correlations of the action density vanish, that the limit is a free Maxwell field at some charge,
and that it is scale invariant. It is a statement about a two-point function only; see
Definition~\ref{3.1:def-non-trivial}.
\end{description}

\medskip
\noindent Large-field and age vocabulary. These words are defined at their first use in
Chapter~\ref{ch:8}; the canonical comparison point, first used there, is fixed in
Chapter~\ref{ch:9}.
\begin{description}
\item[\normalfont\itshape large-field region created at step $j$] a large-field region that the
large-field operation introduces at step $j$ of the iteration.
\item[\normalfont\itshape age at level $k$] of a region created at step $j$: the number $k-j$, the
number of steps since the step that created it.
\item[\normalfont\itshape young, old] a region is young if its age is below the age threshold
displayed in Chapter~\ref{ch:8}, old if it is above it.
\item[\normalfont\itshape live] still carrying its large-field factor at the level read.
\item[\normalfont\itshape healed] with its large-field factor absorbed back into the small-field
format by a later step, and the region re-filed into the regular expansion.
\item[\normalfont\itshape never-healed old-age tail] the contribution of regions created more than
the prescribed number of steps before the witness depth and not healed by then.
\item[\normalfont\itshape depth floor] a lower bound on the depth, fixed at its first use in
Chapter~\ref{ch:8}.
\item[\normalfont\itshape separation window] the two-sided condition in the definition of the class
$\mathfrak{T}(\mathsf{d};\vartheta)$; it is a condition on test functions, not a hypothesis of
Theorem~\ref{3.1:thm-main}.
\item[\normalfont\itshape canonical comparison point $\hat{x}$] the point $\hat{x}$ listed under
the heading ``Uniqueness'' above.
\item[\normalfont\itshape pocket] a region whose fluctuation variables were integrated out at a
given level, as defined at the first use of the word in Chapter~\ref{ch:8};
Chapter~\ref{ch:4} defines the word separately, for a large-field region together with a level
attached to it by the construction.
\end{description}
\end{definition}

%% ---- end inlined file: chapters/c00 ----

\part{The statement of Theorem~0}
\clearpage
\setcounter{section}{0}
\refstepcounter{section}
\section*{Chapter~\thesection\quad The lattice, the action and the smeared curvature observables\addcontentsline{toc}{section}{\protect\numberline{\thesection}The lattice, the action and the smeared curvature observables}\label{ch:1}}
%% ---- begin inlined file: chapters/c01 ----
\begin{definition}[The torus, the bare lattice, bonds and plaquettes; physical and lattice units]
\label{1.1:def-lattice}
Fix the blocking factor $L$, an odd integer; an integer $m\ge 1$; and an integer $K\ge 0$, the
number of renormalization steps. Ba{\l}aban assumes that $L$ is an odd integer greater than $11$
and that $m$ is a positive integer \cite{B-CMP109}.

\smallskip
\noindent\emph{The torus.} The \emph{torus of side $2L^m$} in physical units is
\begin{equation}\label{1.1:eq-lattice-1}
\mathbb{T}_m:=\mathbb{R}^4/(2L^m\mathbb{Z}^4) .
\end{equation}

\smallskip
\noindent\emph{The bare lattice.} At cutoff $K$ the \emph{bare lattice spacing} is $L^{-K}$. The
\emph{bare sites} are the points of
$L^{-K}(\mathbb{Z}+\tfrac12)^4$ taken modulo $2L^m\mathbb{Z}^4$, that is, regarded as
points of $\mathbb{T}_m$ \cite[(0.1)]{B-CMP109}. Since the side $2L^m$ of the torus is
$2L^{m+K}$ times the spacing $L^{-K}$, there are $2L^{m+K}$ bare sites per side and
$(2L^{m+K})^4$ bare sites in all.

\smallskip
\noindent\emph{Bonds.} A \emph{bond} is an ordered pair $b=\langle x,y\rangle$ of bare sites at
distance $L^{-K}$. The \emph{reversed bond} is $-b:=b^{-1}:=\langle y,x\rangle$; we write
$-b$ throughout.

\smallskip
\noindent\emph{Plaquettes.} A \emph{plaquette} $p$ is an elementary square of the bare lattice.
Its boundary $\partial p$ is a closed path of four bonds, and $x(p)\in\mathbb{T}_m$ denotes its
\emph{barycentre}, the centre of the square.

\smallskip
\noindent\emph{Units.} Two systems of units are used.
\begin{itemize}
\item In \emph{physical units} the torus $\mathbb{T}_m$ has side $2L^m$ and the bare spacing is
$L^{-K}$.
\item After $k$ renormalization steps, $0\le k\le K$, the level-$k$ block lattice, defined
together with Ba{\l}aban's block-averaging map, has spacing
$L^{k-K}$. In \emph{level-$k$ lattice units} this spacing is $1$, so that the torus has side
$2L^{m+K-k}$. \emph{Bare lattice units} are the case $k=0$.
\end{itemize}
Clauses (0) and (i) of Theorem~0 at scale $k$ are stated in level-$k$ lattice units; clauses
(ii)--(vi) are stated in physical lengths.
\end{definition}

\begin{definition}[Gauge fields, Haar measure, the Wilson action and its weighted form, the bare density and the expectation at cutoff $K$]\label{1.1:def-wilson-action}
Let $\mathbb{T}_m$, the bare lattice of spacing $L^{-K}$ on it, its bonds, its plaquettes $p$ with boundaries $\partial p$ and their barycentres $x(p)$ be as in Definition~\ref{1.1:def-lattice}.
\begin{enumerate}
\item A \emph{gauge field} is a map $U$ from the oriented bonds to $\mathrm{SU}(N)$, $N\ge2$, with
$U(-b)=U(b)^{-1}$.
Choosing one orientation of each bond identifies the configuration space with a finite product
of copies of $\mathrm{SU}(N)$. We write $dU$ for the product of the normalised Haar measures.
\item If $\partial p$ is traversed as $b_1b_2b_3b_4$, the \emph{plaquette holonomy} is
$U(\partial p):=U(b_1)U(b_2)U(b_3)U(b_4)$. The trace is normalised, $\operatorname{tr}\mathbf{1}=1$.
For the unoriented plaquette $p$ the quantity
\begin{equation}\label{1.1:eq-wilson-action-1}
1-\operatorname{Re}\operatorname{tr}U(\partial p)\in[0,2]
\end{equation}
is well defined.
\item The \emph{Wilson action} and its \emph{weighted form}, for a real- or complex-valued function
$c$ on $\mathbb{T}_m$ with the weight attached at the barycentre, are
\begin{equation}\label{1.1:eq-wilson-action-2}
A(U):=\sum_p\bigl[1-\operatorname{Re}\operatorname{tr}U(\partial p)\bigr],\qquad
A(c,U):=\sum_p c\bigl(x(p)\bigr)\bigl[1-\operatorname{Re}\operatorname{tr}U(\partial p)\bigr],
\end{equation}
the sums running over all unoriented plaquettes $p$ of the bare lattice on $\mathbb{T}_m$
(Definition~\ref{1.1:def-lattice}).
The map $c\mapsto A(c,U)$ is linear, and $A(1,U)=A(U)$. For a real Lipschitz function $f$ on
$\mathbb{T}_m$ set $\mathcal{O}(f):=A(f,U)$; this is the smeared curvature field of
Definition~\ref{1.3:def-curvature-field}. In particular, for a bare coupling $g_0>0$, complex numbers
$z_1,\dots,z_n$ and real Lipschitz functions $f_1,\dots,f_n$ on $\mathbb{T}_m$,
\begin{equation}\label{1.1:eq-wilson-action-3}
-g_0^{-2}A(U)+\sum_{i=1}^n z_i\,\mathcal{O}(f_i)
=-A\Bigl(g_0^{-2}-\sum_{i=1}^n z_if_i,\;U\Bigr).
\end{equation}
\item The \emph{bare density} is
\begin{equation}\label{1.1:eq-wilson-action-4}
\rho_0(U):=\exp\bigl[-g_0^{-2}A(U)\bigr],
\end{equation}
and it satisfies $0<\int\rho_0\,dU<\infty$.
\item The \emph{Wilson measure} is the probability measure
$\mu_{K,m}:=\rho_0\,dU\big/\!\int\rho_0\,dU$. For bounded Borel functions $F$ of the bare field,
the \emph{expectation at cutoff $K$} is
\begin{equation}\label{1.1:eq-wilson-action-5}
\langle F\rangle_K=\langle F\rangle_{K,m}:=\int F\,d\mu_{K,m}
=\frac{\int F\rho_0\,dU}{\int\rho_0\,dU}.
\end{equation}
\end{enumerate}
\end{definition}

\begin{proof}
The definition contains three assertions: the well-definedness and range in
\eqref{1.1:eq-wilson-action-1}, the identity \eqref{1.1:eq-wilson-action-3}, and the bounds
$0<\int\rho_0\,dU<\infty$ on the normalisation of the bare density \eqref{1.1:eq-wilson-action-4}.

Changing the starting point of the traversal of $\partial p$ replaces $U(b_1)U(b_2)U(b_3)U(b_4)$
by a cyclic permutation of the product, which is a conjugate of $U(\partial p)$; reversing the
orientation replaces it by $U(-b_4)U(-b_3)U(-b_2)U(-b_1)=U(\partial p)^{-1}$. The trace is
invariant under conjugation, and for a unitary matrix $V$ one has
$\operatorname{tr}V^{-1}=\operatorname{tr}V^{*}=\overline{\operatorname{tr}V}$, which has the same
real part as $\operatorname{tr}V$. Hence $\operatorname{Re}\operatorname{tr}U(\partial p)$ depends
only on the unoriented plaquette $p$. Since $U(\partial p)$ is unitary, its eigenvalues have
modulus one, and the normalised trace is their average; thus
$|\operatorname{tr}U(\partial p)|\le1$, so $\operatorname{Re}\operatorname{tr}U(\partial p)\in[-1,1]$
and $1-\operatorname{Re}\operatorname{tr}U(\partial p)\in[0,2]$.

For \eqref{1.1:eq-wilson-action-3}, linearity of $c\mapsto A(c,U)$ is immediate from
\eqref{1.1:eq-wilson-action-2}, since the sum over $p$ is finite and each term is linear in the
value $c(x(p))$; and $A(1,U)=A(U)$ by comparing the two sums in \eqref{1.1:eq-wilson-action-2}.
Writing $A(U)=A(1,U)$ and $\mathcal{O}(f_i)=A(f_i,U)$ and using linearity with the complex
coefficients $-g_0^{-2}$ and $z_1,\dots,z_n$,
\[
-g_0^{-2}A(1,U)+\sum_{i=1}^n z_iA(f_i,U)
=A\Bigl(-g_0^{-2}+\sum_{i=1}^n z_if_i,\;U\Bigr)
=-A\Bigl(g_0^{-2}-\sum_{i=1}^n z_if_i,\;U\Bigr).
\]

Finally, the configuration space is a finite product of copies of the compact group
$\mathrm{SU}(N)$, hence compact, and $dU$ is a probability measure on it. The function $A$ is a
finite sum of continuous functions of $U$, so $\rho_0$ is continuous, and it is strictly positive
because it is an exponential of a real number. A continuous strictly positive function on a compact
space is bounded above and bounded below by a positive constant, so $0<\int\rho_0\,dU<\infty$.
Consequently $\mu_{K,m}$ is a well-defined probability measure, and the integral in
\eqref{1.1:eq-wilson-action-5} exists for every bounded Borel function $F$.
\end{proof}

The bare lattice of Definition~\ref{1.1:def-lattice} has spacing $L^{-K}$.
A \emph{gauge transformation} is a map $\lambda$ from bare sites to $\mathrm{SU}(N)$. It acts on gauge
fields by
\begin{equation}\label{1.1:eq-invariance-1}
(U^\lambda)(\langle x,y\rangle):=\lambda(x)\,U(\langle x,y\rangle)\,\lambda(y)^{-1}.
\end{equation}
Let $\sigma$ be an isometry of $\mathbb{T}_m$ that maps bare sites to bare sites. It acts by
\begin{equation}\label{1.1:eq-invariance-2}
(\sigma U)(\langle x,y\rangle):=U(\langle\sigma^{-1}x,\sigma^{-1}y\rangle).
\end{equation}
Such isometries include:
\begin{itemize}
\item the translations by $L^{-K}\mathbb{Z}^4$;
\item the hyperoctahedral group $W_4$ of coordinate permutations and sign changes, which preserves
$L^{-K}(\mathbb{Z}+\frac12)^4$ and $2L^m\mathbb{Z}^4$;
\item the reflections $r_{\nu,c}\colon x_\nu\mapsto 2c-x_\nu$ with $2c\in L^{-K}\mathbb{Z}$, each a
sign change followed by a translation by $2c$ in the $\nu$-th coordinate direction.
\end{itemize}
On $\mathbb{T}_m$ the reflection $r_{\nu,c}$ fixes the two hyperplanes $\{x_\nu=c\}$ and
$\{x_\nu=c+L^m\}$. It exchanges the two components
\begin{equation}\label{1.1:eq-invariance-3}
H^+_{\nu,c}:=\{c<x_\nu<c+L^m\},\qquad H^-_{\nu,c}:=\{c-L^m<x_\nu<c\}
\end{equation}
of their complement. The hyperplanes contain no bare site if $c\in L^{-K}\mathbb{Z}$, and they
contain bare sites if $c\in L^{-K}(\mathbb{Z}+\frac12)$. We write
\begin{equation}\label{1.1:eq-invariance-4}
(\Theta_{\nu,c}F)(U):=\overline{F(r_{\nu,c}U)}
\end{equation}
for the reflection composed with complex conjugation. Reflection positivity of the Wilson measure
with respect to the operations $\Theta_{\nu,c}$ is proved in Chapter~\ref{ch:7}.

\begin{lemma}[Gauge transformations, lattice isometries, the hypercubic group and reflections;
invariance of the Wilson measure]\label{1.1:lem-invariance}
For every $g_0>0$, every gauge transformation $\lambda$, every isometry $\sigma$ of $\mathbb{T}_m$
mapping bare sites to bare sites, and every bounded Borel function $F$,
\begin{equation}\label{1.1:eq-invariance-5}
\langle F(U^\lambda)\rangle_K=\langle F(\sigma U)\rangle_K=\langle F\rangle_K .
\end{equation}
\end{lemma}

\begin{proof}
We first show that the action $A(U)$ of Definition~\ref{1.1:def-wilson-action} is invariant.
Let $\partial p$ run through the sites $x_1,x_2,x_3,x_4$ and back to $x_1$, so that
\[
U(\partial p)=U(\langle x_1,x_2\rangle)\,U(\langle x_2,x_3\rangle)\,U(\langle x_3,x_4\rangle)\,
U(\langle x_4,x_1\rangle).
\]
Under $U\mapsto U^\lambda$, by \eqref{1.1:eq-invariance-1}, the
intermediate factors $\lambda(x_i)^{-1}\lambda(x_i)$ cancel, and
$U^\lambda(\partial p)=\lambda(x_1)\,U(\partial p)\,\lambda(x_1)^{-1}$: each $U(\partial p)$ is
conjugated by $\lambda$ at the starting site of $\partial p$. Under $U\mapsto\sigma U$, by
\eqref{1.1:eq-invariance-2}, plaquettes are permuted, possibly with a change of orientation or of
starting point. A change of starting point replaces $U(\partial p)$ by a cyclic permutation of its
four factors, which is a conjugate of $U(\partial p)$; a change of orientation replaces
$U(\partial p)$ by $U(\partial p)^{-1}=U(\partial p)^{*}$. Since
$\operatorname{tr}$ is invariant under conjugation and $\operatorname{Re}\operatorname{tr}W^{*}
=\operatorname{Re}\overline{\operatorname{tr}W}=\operatorname{Re}\operatorname{tr}W$ for
$W\in \mathrm{SU}(N)$, in both cases each term $1-\operatorname{Re}\operatorname{tr}U(\partial p)$ is
either left as it is or moved to another plaquette. Summing over $p$ gives
$A(U^\lambda)=A(\sigma U)=A(U)$.

Next we show that $dU$ is invariant. In the coordinates given by the positively oriented bonds,
$U\mapsto U^\lambda$ multiplies each coordinate on the left and on the right by fixed elements of
$\mathrm{SU}(N)$. The map $U\mapsto\sigma U$ permutes the coordinates and inverts those whose
orientation is reversed, the variable of a reversed bond being the inverse of the variable of the
bond. Haar measure on the compact group $\mathrm{SU}(N)$ is invariant under left and right
translations and under inversion. Hence the product measure $dU$ is invariant, and, together with
the invariance of the action shown above, so is $\rho_0\,dU$, where $\rho_0$ is the bare density of
Definition~\ref{1.1:def-wilson-action}.

Finally, $\langle\cdot\rangle_K$ is the expectation with respect to the normalised measure
$\rho_0\,dU/\int\rho_0\,dU$, and $F$ is bounded and Borel, so all integrals below are finite.
Write $U'=U^\lambda$, respectively $U'=\sigma U$. The substitution $U\mapsto U'$ is a bijection of
the space of gauge fields which preserves $\rho_0\,dU$, so
$\int F(U')\,\rho_0(U)\,dU=\int F(U)\,\rho_0(U)\,dU$. Dividing by
$\int\rho_0\,dU$ yields \eqref{1.1:eq-invariance-5}.
\end{proof}

\begin{definition}[The running inverse couplings and the coupling increment]\label{1.2:def-couplings}
Let $g_0>0$ be the bare coupling of Definition~\ref{1.1:def-wilson-action}. The \emph{running
inverse couplings} $x_k$ and the \emph{running couplings} $g_k$ are defined by
\begin{equation}\label{1.2:eq-couplings-1}
x_0:=g_0^{-2},\qquad x_{k+1}:=x_k-\beta_{k+1},\qquad g_k:=x_k^{-1/2}\quad(x_k>0),
\end{equation}
the recursion being run forward from $g_0$. Thus $x_k=g_k^{-2}$ whenever $x_k>0$, and
\eqref{1.2:eq-couplings-1} is the orientation
\begin{equation*}
g_k^{-2}=g_{k+1}^{-2}+\beta_{k+1}
\end{equation*}
of Ba{\l}aban's coupling flow, \cite[(0.18), (0.20)]{B-CMP109}: the inverse coupling at level
$k+1$ is the inverse coupling at level $k$ diminished by $\beta_{k+1}$.

The number $\beta_{k+1}$, the \emph{coupling increment} of step $k+1$, is Ba{\l}aban's coupling
renormalization: it is defined in \cite[(1.22)]{B-CMP109}, and Ba{\l}aban identifies it with
\cite[(3.61)]{B-CMP119}. Its construction within Ba{\l}aban's step, its independence of the torus
and of $K$, and the bound $|\beta_{k+1}|\le B_{\sharp}$, with $B_{\sharp}$ the flow constant of
Definition~\ref{1.2:def-flow-constants}, belong to Chapter~\ref{ch:4}.

The positivity $x_k>0$ for $0\le k\le K$, which is what makes $g_k$ defined at every level of
the recursion, is not assumed in this definition; it is part of clause~(0) of
the main theorem.
\end{definition}

\begin{definition}[Flow constants, the reference window, first passage and the volume function]
\label{1.2:def-flow-constants}
The one-loop constant is
\begin{equation}\label{1.2:eq-flow-constants-1}
c_0:=\frac{11N^2}{12\pi^2}\log L .
\end{equation}
The other flow constants are as follows.
\begin{itemize}
\item $c_2>0$ is the constant in the finite part of the one-loop law of the coupling flow; it
depends only on $L$ and $N$, and $c_5:=2c_2$.
\item $\lambda_\sharp$ and $B_\sharp>0$ are the two slopes of clause (0) of Theorem~0. They
depend only on $L$, the auxiliary parameters $\mathfrak{p}$, $N$ and $\gamma$.
\end{itemize}
Let $C_\beta$ denote the member of the tuple $\mathfrak{p}$ of Ba{\l}aban's auxiliary parameters so
written. With the rest of $\mathfrak{p}$ it is chosen after $L$ and before $\gamma$, and its choice is
the one made in Chapter~\ref{ch:4}.
The flow constants satisfy
\begin{equation}\label{1.2:eq-flow-constants-2}
\lambda_\sharp\ge\underline{\lambda}:=c_0/4,\qquad B_\sharp\le\bar B:=C_\beta+1 ,
\end{equation}
so that $\bar B\ge B_\sharp$ is a bound independent of $\gamma$.
The constants $c_2$, $C_\beta$, $\lambda_\sharp$, $B_\sharp$ are fixed in Chapter~\ref{ch:4} (the coupling
flow). Here they are named with their dependences exactly as stated above, and nothing more is
claimed about them.

Every constant fixed before $\gamma$, and every ceiling on $\gamma$, depends on the flow constants
only through the triple $(c_0/4,c_5,\bar B)$, which is fixed before $\gamma$. None of them reads
$\lambda_\sharp$ or $B_\sharp$. Quantities fixed after $\gamma$ may read $\lambda_\sharp$ and
$B_\sharp$. Among them are $\gamma_H$ (see the thresholds chosen last and the clause-internal
quantifiers, Remark~\ref{2.3:rem-thresholds}), $g_-(g)$, $K(g_0,g)$,
$m_{\mathbb{T}}(g_-(g))$ and the lower bound $\varrho\ge6(B_\sharp+1)$ of clause (v).

Let $x_k=g_k^{-2}$ be the running inverse couplings of Definition~\ref{1.2:def-couplings}. Further,
\begin{equation}\label{1.2:eq-flow-constants-3}
\begin{split}
g_-(g)&:=(g^{-2}+B_\sharp)^{-1/2},\\
K(g_0,g)&:=\min\{k\ge0:\ x_k\le g^{-2}+B_\sharp\},\\
R(g')&:=\min\{L^i:\ i\in\mathbb{Z}_{\ge0},\ L^i\ge(\log g'^{-2})^r\},\\
m_{\mathbb{T}}(g')&:=11+m'+\log_L R(g'),\qquad M=L^{m'} .
\end{split}
\end{equation}
Here $r$ is one of the integer exponents fixed before $L$, in the first stage of the order of
choice, which reads only $(d,N)$; it is the exponent entering hypothesis (H5) among the hypotheses
(H1)--(H7) of Notation~\ref{2.2:not-hypotheses}. The
integer $m'$ is the exponent of Ba{\l}aban's auxiliary parameter $M=L^{m'}$ in the tuple
$\mathfrak{p}$.

The couplings $g'$ with $g_-(g)\le g'\le g$, equivalently $g^{-2}\le g'^{-2}\le g^{-2}+B_\sharp$,
form the \emph{reference window} at $g$; since $B_\sharp>0$, one has $g_-(g)<g$, and $g_-(g)$ is the
lower end of the window. In the quantifier prefix of Theorem~0 the bare coupling ranges over
$0<g_0\le g_-(g)$. The quantity $K(g_0,g)$ is called the
\emph{first-passage cutoff}, and the bare lattice spacing is $L^{-K}$ with $K:=K(g_0,g)$. That
$K(g_0,g)$ is a well-defined finite integer under the hypotheses is clause (0) of Theorem~0; it is
not asserted here. The function $m_{\mathbb{T}}(\cdot)$ is the volume function: in the quantifier
prefix of Theorem~0 the torus exponent ranges over $m\ge m_{\mathbb{T}}(g_-(g))$.
\end{definition}

\begin{definition}[The smeared curvature field (action density), the test-function norm, the generating function, the connected Schwinger functions on separated supports and the centred moments]
\label{1.3:def-curvature-field}
We use the torus $\mathbb{T}_m$, the bare lattice of spacing $L^{-K}$ and its plaquettes $p$, with
boundary $\partial p$ and barycentre $x(p)$, of Definition~\ref{1.1:def-lattice}. We also use the gauge
fields $U$, the Haar measure $dU$, the convention $\operatorname{tr}\mathbf{1}=1$ and the Wilson action
$A(U)$ of Definition~\ref{1.1:def-wilson-action}.

\begin{enumerate}
\item \emph{The curvature field.} Let $f$ be a real Lipschitz function on $\mathbb{T}_m$. Put
\begin{equation}\label{1.3:eq-curvature-field-1}
\mathcal{O}(f)=\sum_p f(x(p))\,\bigl[1-\operatorname{Re}\operatorname{tr}U(\partial p)\bigr].
\end{equation}
The sum runs over all bare plaquettes $p$ of all six orientations.

Thus $\mathcal{O}(f)$ is the smeared gauge-invariant curvature density, that is, the Wilson action
density with plaquette weights $f$. We call it the \emph{smeared action density}, or the curvature
field smeared with $f$; it is not the tensor $\operatorname{tr}F_{\mu\nu}F_{\rho\sigma}$, whose
components are considered nowhere in this book.

A function $f\in C^1_c(\mathbb{R}^4)$ is transported to $\mathbb{T}_m$, and
\eqref{1.3:eq-curvature-field-1} is then read on $\mathbb{T}_m$, once twice the radius of a ball
containing its support is less than $L^m$.

\item \emph{The test-function norm.} For such $f$ put
\begin{equation}\label{1.3:eq-curvature-field-2}
\|f\|_\sharp=\max\{\|f\|_\infty,\,40M\operatorname{Lip}f\}.
\end{equation}
Here $M=L^{m'}$ is the cube size among the auxiliary parameters of Ba{\l}aban's construction, and
$\operatorname{Lip}f$ is the Lipschitz constant of $f$.

\item \emph{The generating function.} For test functions $f_1,\dots,f_n$ as above and
$z=(z_1,\dots,z_n)$ put
\begin{equation}\label{1.3:eq-curvature-field-3}
Z(z)=\int \exp\Bigl[-g_0^{-2}A(U)+\sum_{i=1}^n z_i\,\mathcal{O}(f_i)\Bigr]\,dU .
\end{equation}

\item \emph{The connected Schwinger functions.} Let $n\ge2$, and let $f_1,\dots,f_n$ be test functions
whose supports are pairwise separated: there is $\mathsf{d}>0$ with
\begin{equation}\label{1.3:eq-curvature-field-4}
\operatorname{dist}(\operatorname{supp}f_i,\operatorname{supp}f_j)\ge\mathsf{d}\qquad(i\neq j).
\end{equation}
For such $f_1,\dots,f_n$ we define
\begin{equation}\label{1.3:eq-curvature-field-5}
S^T_{n;K,m}(f_1,\dots,f_n)=\partial_{z_1}\cdots\partial_{z_n}\log Z(0)
=\langle\mathcal{O}(f_1);\dots;\mathcal{O}(f_n)\rangle^{\mathrm{conn}}_{K,m}.
\end{equation}
The second equality defines the connected expectation
$\langle\,\cdot\,;\dots;\,\cdot\,\rangle^{\mathrm{conn}}_{K,m}$ at cutoff $K$ on $\mathbb{T}_m$.

These are the Schwinger functions of the main theorem. They are taken only for $n\ge2$ and only for
test functions whose supports are pairwise at distance at least $\mathsf{d}>0$, as in
\eqref{1.3:eq-curvature-field-4}. For what lies outside this class, see the remarks following this
definition. When $m$ is fixed we write $S^T_{n;K}$.

\item \emph{The centred moments.} The centred Schwinger functions $S_q$ are the centred moments
determined by the connected ones. For $f_1,\dots,f_q$ as in \eqref{1.3:eq-curvature-field-4}, each is
a sum over partitions of $\{1,\dots,q\}$ into blocks of size at least two:
\begin{equation}\label{1.3:eq-curvature-field-6}
S_q(f_1,\dots,f_q)=\sum_{\pi}\ \prod_{B\in\pi}S^T_{|B|;K,m}\bigl((f_i)_{i\in B}\bigr).
\end{equation}
Here $\pi$ runs over the partitions of $\{1,\dots,q\}$ all of whose blocks $B$ have $|B|\ge2$. Every
subfamily $(f_i)_{i\in B}$ again satisfies \eqref{1.3:eq-curvature-field-4}, so each factor is
defined by \eqref{1.3:eq-curvature-field-5}.

\item \emph{Further notation.}
\begin{itemize}
\item $W_4$ is the hypercubic group of Lemma~\ref{1.1:lem-invariance}.
\item $W_3$ is the subgroup of $W_4$ that fixes a distinguished coordinate axis, the time axis.
\item $(\mathbb{R}^4)^n_{\neq}$ and $(\mathbb{T}_m)^n_{\neq}$ are the sets of $n$-tuples of pairwise
distinct points.
\end{itemize}
\end{enumerate}
\end{definition}

The Schwinger functions of Theorem~0 are the connected functions
$S^T_{n;K,m}(f_1,\dots,f_n)$, $n\ge 2$, of Definition~\ref{1.3:def-curvature-field}, taken
with test functions whose supports are pairwise separated by a distance at least
$\mathsf{d}>0$, and their limits; with them go the centred moments $S_q$ they determine.
The following lie outside this class, and nothing is stated about any of them anywhere in
the book:
\begin{itemize}
\item the case $n=1$;
\item coinciding or overlapping supports (the action density $\operatorname{tr} F\cdot F$ has
dimension four, and $|x-y|^{-8}$ is not locally integrable);
\item block-averaged observables;
\item Wilson loops of fixed physical size;
\item the state on Ba{\l}aban's block algebra.
\end{itemize}
Block variables and the block algebra appear in this book only as the vehicle of the proofs,
through the block-averaging map and the block fields of Definition~\ref{1.5:def-block-average},
and inside Remark~\ref{3.2:rem-block-laws} on the block-law statements, which are not proved in
this book and which no clause of Theorem~0 uses.

This section recalls Ba{\l}aban's step only to the extent that the clauses of Theorem~0 need words
for it; it asserts nothing. Throughout, $G:=\mathrm{SU}(N)$, and the torus $\mathbb{T}_m$, the bare
lattice, its plaquettes $p$ and the transports $U(\partial p)$ are those of
Definitions~\ref{1.1:def-lattice} and~\ref{1.1:def-wilson-action}.

\begin{definition}[Block lattices, Ba{\l}aban's block-averaging map extended to all
configurations, block fields and depth, the block algebra (a vehicle), the exact
push-forward]\label{1.5:def-block-average}
\emph{Block lattices.} Let $\mathbb{T}^{(0)}$ be the bare lattice of
Definition~\ref{1.1:def-lattice}. For $0\le k\le K$ let $\mathbb{T}^{(k)}$ be the lattice
$L^{k-K}(\mathbb{Z}+\frac12)^4$ taken modulo $2L^m\mathbb{Z}^4$. Its spacing is $L^{k-K}$ in
physical units, that is, $L^k$ in bare lattice units.

Let $1\le k\le K$. Since $L$ is odd, $\frac{L-1}{2}$ is an integer, and
$L(n+\frac12)=Ln+\frac{L-1}{2}+\frac12\in\mathbb{Z}+\frac12$ for every $n\in\mathbb{Z}$. Hence
$L^{k-K}(\mathbb{Z}+\frac12)\subset L^{k-1-K}(\mathbb{Z}+\frac12)$, and every site of
$\mathbb{T}^{(k)}$ is a site of $\mathbb{T}^{(k-1)}$. The \emph{blocks} of $\mathbb{T}^{(k-1)}$ are
the sets $B(y)$, $y\in\mathbb{T}^{(k)}$. They are given by \eqref{1.5:eq-block-average-2} below
with $k-1$ in place of $k$, and each contains $L^4$ sites. They partition $\mathbb{T}^{(k-1)}$.

For $0\le k\le K$ let $\mathcal{V}_k:=G^{\mathrm{bonds}(\mathbb{T}^{(k)})}$ be the configurations
on positively oriented bonds, extended by $V(-b)=V(b)^{-1}$. Let $N_k:=|\mathbb{T}^{(k)}|$ be the
number of sites. For $0\le k<K$ and a bond $c$ of $\mathbb{T}^{(k+1)}$, the \emph{double block} of
$c$ is the union of the two blocks of $\mathbb{T}^{(k)}$ centred at the endpoints $c_-$, $c_+$ of
$c$. We write $\mathcal{P}_c$ for the set of plaquettes of $\mathbb{T}^{(k)}$ contained in the
double block of $c$.

\emph{The block average.} Fix a number $\omega_M(L)>0$, depending only on $N$ and $L$ ($d=4$ and
$N$ being fixed first; the subscript refers to the map $\mathcal{M}$ below, not to the parameter
$M$), such that the formulas of \cite[(0.10)--(0.12)]{B-CMP109} are defined on
configurations whose plaquette variables lie within $\omega_M(L)$ of $\mathbf{1}$.
Ba{\l}aban requires only that the averaged configurations have sufficiently small diameter.

Let $0\le k<K$. For a bond $c$ of $\mathbb{T}^{(k+1)}$ put
\begin{equation}\label{1.5:eq-block-average-1}
\mathcal{S}_c:=\bigl\{U\in\mathcal{V}_k:\ |U(\partial p)-\mathbf{1}|<\omega_M(L)\ \text{for all }
p\in\mathcal{P}_c\bigr\}.
\end{equation}
Here $|Y|:=(\operatorname{tr}Y^*Y)^{1/2}$ with $\operatorname{tr}\mathbf{1}=1$. The condition does
not depend on the corner and orientation chosen for $\partial p$, since
$|uV^{\pm1}u^{-1}-\mathbf{1}|=|V-\mathbf{1}|$ for unitary $u$, $V$.

The map $Y\mapsto e^{iY}$ is injective on the set of traceless hermitian $N\times N$ matrices $Y$
whose largest and smallest eigenvalues differ by less than $2\pi$. The \emph{intrinsic logarithm}
of $G$ is its inverse, defined on the image of that set, whose complement in $G$ is closed and
null for Haar measure.

The map $\mathcal{M}:\mathcal{V}_k\to\mathcal{V}_{k+1}$ is defined bond by bond, as follows; in (a)
and (b) the bond $c$ of $\mathbb{T}^{(k+1)}$ is positively oriented.
\begin{itemize}
\item[(a)] For $U\in\mathcal{S}_c$, $\mathcal{M}(U)(c)$ is Ba{\l}aban's symmetrised average
\cite[(0.11)--(0.12)]{B-CMP109}, built on the inner group mean \cite[(0.10)]{B-CMP109}. Every
logarithm in it is the intrinsic logarithm of $G$.
\item[(b)] For $U\notin\mathcal{S}_c$, $\mathcal{M}(U)(c)$ is the double-contour average of
intrinsic logarithms \cite[(0.4)]{B-CMP109}, with the contours and weights of that display. This
applies at every $U$ at which all logarithms entering the display are defined. On the remaining
set, which is null for Haar measure, $\mathcal{M}(U)(c)$ is the parallel transport of $U$ along the
straight path joining the endpoints of $c$.
\item[(c)] For the reversed bond, $\mathcal{M}(U)(-c):=\mathcal{M}(U)(c)^{-1}$.
\end{itemize}
Item (b) extends the map of \cite[(0.11)--(0.12)]{B-CMP109} off the sets $\mathcal{S}_c$, so that
$\mathcal{M}$ is defined on all of $\mathcal{V}_k$; this is the extended map referred to below.

The paths, contour families, inner mean and weights entering (a) and (b) are those of
\cite[(0.3), (0.4), (0.10)--(0.12)]{B-CMP109}, with the logarithm named in (a) and (b); we record
only the blocks. Put
$\eta:=L^{k-K}$, the spacing of $\mathbb{T}^{(k)}$, and $h_L:=(L-1)/2$. For a site $y$ of
$\mathbb{T}^{(k+1)}$, the block centred at $y$ is
\begin{equation}\label{1.5:eq-block-average-2}
B(y):=\Bigl\{y+\eta\sum_{\nu}n_\nu e_\nu:\ n\in\mathbb{Z}^4,\ |n_\nu|\le h_L\Bigr\},
\end{equation}
where $e_1,\dots,e_4$ are the coordinate unit vectors of $\mathbb{R}^4$.
Since $m\ge1$ and $k<K$, a double block lies in a coordinate box of side $(2L-1)\eta$, and so does
every path entering the formulas cited in (a) and (b). This side is less than $L^m$, half the
period, so the covering map $\mathbb{R}^4\to\mathbb{T}_m$ is injective on such a box. All paths
are read in such a box.

\emph{Block fields and depth.} For a bare configuration $U$ on $\mathbb{T}^{(0)}$ and
$0\le k\le K$ put $V_k:=\mathcal{M}^kU\in\mathcal{V}_k$. Here $\mathcal{M}^k$ is the composition
of the maps $\mathcal{M}$ of levels $0,\dots,k-1$, and $V_0=U$. The \emph{depth} of $V_k$ is
$s:=K-k$, and we write $V^{(s)}:=\mathcal{M}^{K-s}U=V_{K-s}$. A block field of depth $s$ lives on
a lattice of physical spacing $L^{-s}$, whatever $K\ge s$.
A parenthesised superscript thus denotes the level on a lattice, as in $\mathbb{T}^{(k)}$, and the
depth on a block field: $V^{(s)}$ is a configuration on $\mathbb{T}^{(K-s)}$, that is, an element
of $\mathcal{V}_{K-s}$.

\emph{The block algebra.} The algebra $\mathfrak{A}$ consists of the bounded, gauge-invariant,
continuous functions of finitely many block variables $V^{(s)}(b)$, of any depths. For
$F\in\mathfrak{A}$ of largest depth $s$, the function $F\circ\mathcal{M}^{K-s}$ of the bare field
is defined for every $K\ge s$.

\emph{The exact push-forward.} Let $dU$ be the Haar probability measure on $\mathcal{V}_k$, and let
$\delta$ be the Dirac distribution on $\mathcal{V}_{k+1}$ normalised against Haar measure. For a
finite positive Borel measure $\rho(U)\,dU$ on $\mathcal{V}_k$ put
\begin{equation}\label{1.5:eq-block-average-3}
(\mathbf{T}\rho)(V):=\int_{\mathcal{V}_k}dU\,\delta\bigl(\mathcal{M}(U)V^{-1}\bigr)\,\rho(U),
\qquad V\in\mathcal{V}_{k+1}.
\end{equation}
This is the image measure under $\mathcal{M}$, taken without any restriction on the integration
variables \cite[(0.1)]{B-CMP119}.
\end{definition}

The algebra $\mathfrak{A}$ is introduced only as the domain of the bounded functions $F$ of block
fields of which the exact-insertion identity of clause~(i) speaks. No clause of Theorem~0 is a
statement about it.

Let $\rho_k$ be the effective density of Definition~\ref{1.5:def-densities}. Inside
$\mathbf{T}\rho_k$ Ba{\l}aban inserts two devices \cite[\S3]{B-CMP119}. The first consists of
exact decompositions of unity into small-field and large-field parts; the second is the gauge
fixing \cite[(1.5)]{B-CMP119}, which is that of \cite[(0.16)]{B-CMP109}. These devices organise
$\mathbf{T}\rho_k$ into terms.

The properties of the extended map $\mathcal{M}$ which the proofs use are stated and established
where they are used in Part~II (Chapter~\ref{ch:4} and Chapter~\ref{ch:7}).

\begin{definition}[The large-field operation in outline, the effective densities and their decomposition
over histories]\label{1.5:def-densities}
\emph{The large-field operation.} Ba{\l}aban's operation $\mathbf{R}$ \cite{B-CMP122a,B-CMP122b}
acts on a connected component of the large-field region of a term of $\mathbf{T}\rho_k$, where
$\mathbf{T}$ is the exact push-forward of Definition~\ref{1.5:def-block-average}. It acts at the
first step $k$ at which certain conditions on the component hold, two of them being Ba{\l}aban's
\cite{B-CMP122a}. These conditions, and the precise set whose meeting by a component counts as the
creation of a new large-field region, are specified in Chapter~\ref{ch:4}. That set is a union of
cubes meeting the large-field sets of \cite[(3.2)--(3.3)]{B-CMP119} and certain further sets. At
such a step $\mathbf{R}$ does two things.
\begin{itemize}
\item[(1)] It integrates the term over the block field $V_k$ inside the component, regards the
result as independent of $V_k$ there \cite{B-CMP122a}, and multiplies it by the normalized reference
laws of \cite[(1.100)--(1.102)]{B-CMP122a}, the display \cite[(1.100)]{B-CMP122a} being the one
that states their normalization.
\item[(2)] On components of Ba{\l}aban's first class, in his classification of large-field
components into two classes \cite{B-CMP122a,B-CMP122b}, it removes the gauge-fixing
$\delta$-functions. This is done by the gauge-fixing identity that introduced them, read in reverse
\cite{B-CMP122b}.
\end{itemize}
Components at which these conditions do not hold at step $k$ are not acted on by $\mathbf{R}$ at
that step, and a term of $\mathbf{T}\rho_k$ none of whose components qualifies is left unchanged
by $\mathbf{R}$.
Ba{\l}aban states the normalization property \cite[(0.4)]{B-CMP122a}:
\begin{equation}\label{1.5:eq-densities-1}
\int dV\,(\mathbf{R}\rho)(V)=\int dV\,\rho(V).
\end{equation}

\emph{Effective densities.} With $L$ and the auxiliary parameters $\mathfrak{p}$ fixed, put
\begin{equation}\label{1.5:eq-densities-2}
\rho_{k+1}:=\mathbf{R}\mathbf{T}\rho_k\quad(0\le k<K),\qquad\text{hence}\qquad
\rho_k=(\mathbf{R}\mathbf{T})^k\rho_0\quad(0\le k\le K),
\end{equation}
where $\rho_0$ is the bare density of Definition~\ref{1.1:def-wilson-action}. Following
Ba{\l}aban we call $\rho_k$ a density; the object defined, once $\mathbf{R}$ is well defined at each
step (the positivity of its denominators is treated in Chapter~\ref{ch:4}), is the finite positive
Borel measure $\rho_k\,dV_k$ on the space $\mathcal{V}_k$ of block
fields $V_k$ of level $k$.

\emph{Histories.} Ba{\l}aban writes $\rho_k\,dV_k$ as a sum over histories $h$ of positive measures
$\tau_h$ \cite[(2.18)]{B-CMP119},
\begin{equation}\label{1.5:eq-densities-3}
\rho_k\,dV_k=\sum_h\tau_h .
\end{equation}
A history records the outcomes of the decompositions of unity at the steps
up to $k$, that is, an admissible sequence of large-field domains. Once $\mathbf{R}$ has acted,
components of the second class keep in this representation the gauge-fixing $\delta$-functions,
which $\mathbf{R}$ does not remove on them. Accordingly, bounds on $\rho_k$ are bounds on measures,
not pointwise bounds on a function.

\emph{Tree domains and the tree length.} Fix one level $\ell$; write $\mathcal{D}_\ell$ for the
domains of level $\ell$ and $d_\ell(\cdot)$ for their length, as in
Definition~\ref{1.5:def-domains}, $\pi_\ell$ for the partition into the cubes of which the domains
of level $\ell$ are unions (cubes are closed), and $N(A)$ for the number of
cubes of $\pi_\ell$ in a union $A$ of such cubes. Let $\mathsf{I}$ be a finite family of pairwise
non-touching boxes, the \emph{islands}, and $W:=\bigcup\mathsf{I}$. Put
\[
\mathcal{D}^{\mathsf{I}}:=\{X\in\mathcal{D}_\ell:\ X\cap W\ \text{is a union of islands},\ X\setminus
W\neq\emptyset\}.
\]
For $X\in\mathcal{D}^{\mathsf{I}}$ let $d_{\ell,W}(X)$ be the length of $X$ relative to $W$, in the
notation $d_{j,Z}(X)$ of Definition~\ref{1.5:def-domains}, and fix a tree $\mathcal{T}_X$ attaining
the minimum in the definition of $d_{\ell,W}(X)$, chosen by a rule reading $(\pi_\ell,\mathsf{I},X)$
only (any fixed tie-break). Put
$\mathsf{c}_d:=2\cdot3^d$, $\lambda_\theta:=(1+3^d)\mathsf{c}_d\log2$ and
$\mathsf{k}_d:=1+\mathsf{c}_d(2d\log3+(1+3^d)\log2)$; then
$\mathsf{k}_d-1=2d\,\mathsf{c}_d\log3+\lambda_\theta\ge\lambda_\theta$. For a box $I$ with sides
$\sigma_1,\dots,\sigma_d$ (in cubes) and largest side $\sigma$, let $I^{(1)}$ be the box with sides
$\sigma_i+2$ containing $I$ centrally, and put $\mathsf{b}(I):=2d(\sigma+2)^{d-1}$. Then
\begin{align*}
N(I^{(1)}\setminus I)&=\prod_{i}(\sigma_i+2)-\prod_i \sigma_i\\
&=\sum_{j=1}^d2\prod_{i<j}(\sigma_i+2)\prod_{i>j}\sigma_i\le2d(\sigma+2)^{d-1}=\mathsf{b}(I),
\end{align*}
the middle identity by telescoping, replacing $\sigma_j$ by $\sigma_j+2$ one index at a time, and
the inequality because each of the $d$ summands is at most $2(\sigma+2)^{d-1}$. For
$A\in\mathcal{D}_\ell$ let $\mathsf{I}(A)$ be the set of islands that meet or touch $A$. The
\emph{tree domain} of $X\in\mathcal{D}^{\mathsf{I}}$ is the point set
\[
X^\theta:=\bigcup\{\Delta\in\pi_\ell:\ \Delta\cap\mathcal{T}_X\neq\emptyset\},
\]
the union of the cubes of $\pi_\ell$ met by the tree; for $X\in\mathcal{D}_\ell$ meeting no island,
$X^\theta:=X$. The \emph{tree length} is $d^\theta_\ell(X):=d_\ell(X^\theta)$. For
$X\in\mathcal{D}^{\mathsf{I}}$ these quantities satisfy
\[
d_\ell(X^\theta)\le d_{\ell,W}(X)\le d_\ell(X),\qquad
N(X^\theta)\le\mathsf{c}_d\bigl(d^\theta_\ell(X)+1\bigr),\qquad
\#\mathsf{I}(X^\theta)\le3^dN(X^\theta).
\]
\end{definition}

\begin{definition}[Position-dependent couplings and the background configuration, the degree
function, localization domains and their lengths, analyticity domains read
hereditarily]\label{1.5:def-domains}
\emph{Position-dependent couplings.} Near large-field regions the coupling is position
dependent \cite[(2.24)]{B-CMP119}. For a function $c$ on the torus, $A(c,U)$
is the weighted Wilson action of Definition~\ref{1.1:def-wilson-action}, with $x(p)$ the
barycentre of the plaquette $p$. Let $\Lambda_j$ be Ba{\l}aban's small-field domain at
level $j$, and let $\Lambda_j^{\sim-1}$ be the domain $\Lambda_j$ with one layer of cubes
removed. Let $\varphi_j\in C_0^\infty(\Lambda_j)$ with $\varphi_j=1$ on
$\Lambda_j^{\sim-1}$. The position-dependent inverse couplings $x_j(\cdot)$ satisfy
\begin{equation}\label{1.5:eq-domains-1}
x_j(x)=x_{j-1}(x)-\beta_j\,\varphi_j(x).
\end{equation}
With $x_0(\cdot)\equiv x_0$, the recursion \eqref{1.5:eq-domains-1} gives
\begin{equation}\label{1.5:eq-domains-2}
x_j(\cdot)-x_j=\sum_{1\le i\le j}\beta_i\,(1-\varphi_i).
\end{equation}
Ba{\l}aban states, for the nested small-field domains of an admissible sequence
\cite[\S2]{B-CMP119}, that $x_j(x)=x_j$ on $\Lambda_j^{\sim-1}$ \cite[p.~259]{B-CMP119};
he writes this as $g_j^2(x)=g_j^2$ on the last domain.

\emph{The background configuration.} The effective action at level $k$ contains
$-A(x_k(\cdot),U_k)$ \cite[(2.23)]{B-CMP119}. Here $U_k=U_k(V_k)$ is Ba{\l}aban's
background configuration, the minimal configuration determined by the block field $V_k$ of
Definition~\ref{1.5:def-block-average}.

\emph{The degree function.} The degree function is
\begin{equation}\label{1.5:eq-domains-3}
p_0(g):=A_0\,(\log g^{-2})^{p_0}.
\end{equation}

\emph{Localization domains.} $\mathcal{D}_k$ is the set of finite connected unions of cubes
of Ba{\l}aban's partitions at level $k$ \cite{B-CMP109}, \cite[\S\S1--2]{B-CMP116}. In
\cite{B-CMP116} only unions of connected families of $M_1$-cubes are required. For
$X\in\mathcal{D}_k$, $d_k(X)$ is Ba{\l}aban's size of $X$ at level $k$ \cite{B-CMP109}, and
$d_{k,Z}(X)$ is his relative length \cite[(1.67)]{B-CMP122b}. For $X\in\mathcal{D}_j$, the
length $d^\theta(X)$ appearing in clause (i-a) of Theorem~0 is a tree length defined
in Chapter~\ref{ch:4}, not exceeding $d_{j,Z}(X)$.

\emph{Analyticity domains.} For $X\in\mathcal{D}_j$ and radii $\alpha_0,\alpha_1>0$,
$U^c_j(X,\alpha_0,\alpha_1)$ \cite[(1.11)--(1.16)]{B-CMP109} is the set of complex
configurations $\mathbf{U}=U'U$, in the notation there, with a current, that satisfy the
following conditions.
\begin{itemize}
\item[(R1)--(R3)] Pointwise and local gauge bounds on $\mathbf{U}$ and on the current. The
radius for the current is a third argument, equal to $\alpha_0$ unless indicated otherwise.
\item[(R4)] Bounds, on $X$ less its collar, for the constrained minimizers
\cite{B-CMP102b} built from block averages of $\mathbf{U}$ along the maximal tower of
subdomains of $X$ \cite[(1.3)--(1.6)]{B-CMP99a}. The collar is the two layers of level-$j$
cubes of $X$ adjacent to its boundary.
\end{itemize}
In clause (i-a) of Theorem~0, $U^c_j(X,\alpha_0,\alpha_1)$ is read hereditarily: it
denotes the set of configurations of the set just defined that also satisfy (R4), with
Ba{\l}aban's constants, at every subregion of every tower of subdomains of $X$ formed under
either of Ba{\l}aban's two rules for forming these subdomains \cite{B-CMP109},
\cite[(2.13)]{B-CMP119}, the collar of each subregion being the two layers of its own cubes
adjacent to its boundary.
\end{definition}

%% ---- end inlined file: chapters/c01 ----

\clearpage
\setcounter{section}{1}
\refstepcounter{section}
\section*{Chapter~\thesection\quad The order of quantifiers and the hypotheses\addcontentsline{toc}{section}{\protect\numberline{\thesection}The order of quantifiers and the hypotheses}\label{ch:2}}
%% ---- begin inlined file: chapters/c02 ----
\begin{definition}[The quantifier prefix of Theorem~\ref{3.1:thm-main} and its modifications inside
clauses (iii), (iv), (v), (vi)]\label{2.1:not-prefix}
We write $\mathfrak{p}$ for the tuple of Ba{\l}aban's auxiliary parameters listed in hypothesis (H3)
(Definition~\ref{2.2:not-hypotheses}). A \emph{Stage-0 datum} is a finite collection of integer
exponents and of numbers that depend on $(d,N)$ only. Theorem~\ref{3.1:thm-main} has the prefix
\begin{multline}\label{2.1:eq-prefix-1}
\exists\,(\text{Stage-0 datum})\ \ \exists\,L_0=L_0(N)\ \ \forall\,L\ge L_0\ \text{odd}\\
\exists\,\mathfrak{p}\ (\text{in one fixed order})\ \ \exists\,\gamma\ \ \exists\,\gamma_H\\
\forall\,g\in(0,\gamma_H]\ \ \forall\,m\ge m_{\mathbb{T}}(g_-(g))\\
\forall\,g_0\in(0,g_-(g)]:\quad K:=K(g_0,g),
\end{multline}
where $K(g_0,g)$ is the first-passage level of Definition~\ref{1.2:def-flow-constants},
\begin{equation}\label{2.1:eq-prefix-2}
K(g_0,g)=\min\{\,k\ge0:\ x_k\le g^{-2}+B_\sharp\,\},
\end{equation}
with $x_k$ the running inverse couplings of Definition~\ref{1.2:def-couplings}, and the bare lattice
spacing is $L^{-K}$. The constant $L_0=L_0(N)$ is bound existentially once in
\eqref{2.1:eq-prefix-1} and is never given a value; every condition on $L$ is a lower bound.

The prefix \eqref{2.1:eq-prefix-1} is modified inside clauses (ii)--(vi) and inside the paragraph of
part (i-d) on sourced densities, and only there. Each modification is an upper bound on a threshold
chosen last or a one-sided floor. The modifications for clauses (iii), (iv), (v) and (vi) are listed
below; clause (ii) and the paragraph of part (i-d) on sourced densities take the modification of
clause (iii), except that part (ii-$\mathbb{T}$) takes that of clause (v), the first-passage
quantifier returning in part (ii-$\mathbb{T}$-b) as in part (v-c)(2).
\begin{itemize}
\item Inside clauses (iii) and (vi), the segment $\forall\,g\in(0,\gamma_H]$ of
\eqref{2.1:eq-prefix-1} is replaced by
\[
\forall\,\varrho>0\ \ \exists\,\gamma_J\in(0,\gamma_H]\ \ \forall\,g\in(0,\gamma_J],
\]
followed by the same quantifiers over $m$ and $g_0$.
\item Inside clause (iv), which is stated in the regime of (iii), the same segment is replaced by
\begin{multline*}
\forall\,\varrho>0\ \ \exists\,\gamma_J\in(0,\gamma_H]\\
\forall\,\vartheta\in(0,1]\ \ \exists\,\gamma_W\in(0,\gamma_J]\ \ \forall\,g\in(0,\gamma_W],
\end{multline*}
followed by the quantifiers over $m$ and $g_0$; here $\vartheta$ is a shape constant of the test
functions.
\item Inside clause (v), the same segment is replaced by
\[
\forall\,\varrho>0\ \ \exists\,\gamma_U\in(0,\gamma_H]\ \ \forall\,g\in(0,\gamma_U].
\]
Then come an integer $m_0\ge0$, the torus exponent of clause (v), (H5) being assumed for that torus
as for every clause, and the depth $n$, the bare inverse coupling being tuned to
$g_0^{-2}=a_n-\hat w$ instead of being fixed by first passage. The first-passage quantifier over
$g_0\in(0,g_-(g)]$ with $K=K(g_0,g)$ returns only in part (v-c)(2), the identification of
subsequential limits, and in part (ii-$\mathbb{T}$-b); in part (v-c)(2) it requires in addition
$\varrho\ge6(B_\sharp+1)$ and $g$ below a further threshold, which depends on $\varrho$ and is at
most $\gamma_U$.
\end{itemize}
The thresholds $\gamma$, $\gamma_H$, $\gamma_J$, $\gamma_W$, $\gamma_U$ and the quantities that depend
on $\varrho$ and on $\vartheta$ are those of Remark~\ref{2.3:rem-thresholds}.

Fixing the Stage-0 numbers before $L$, rather than after $L$ together with the parameters of (H3), is
the stronger statement: an $\exists\exists\forall$ prefix implies the corresponding
$\exists\forall\exists$ prefix. Indeed, suppose a Stage-0 datum and a constant $L_0$ are such that for
every odd $L\ge L_0$ there are $\mathfrak{p}$, $\gamma$, $\gamma_H$ for which the statement formed by
the remaining quantifiers of \eqref{2.1:eq-prefix-1}, followed by the clause they govern, holds. Then
the prefix in which the datum is chosen after $L$, namely
\[
\exists\,L_0\ \ \forall\,L\ge L_0\ \text{odd}\ \ \exists\,(\text{Stage-0 datum})\ \
\exists\,\mathfrak{p}\ \ \exists\,\gamma\ \ \exists\,\gamma_H\ \cdots,
\]
holds with the same $L_0$: for each odd $L\ge L_0$ one
takes the same datum, which does not depend on $L$, and the $\mathfrak{p}$, $\gamma$, $\gamma_H$
supplied for that $L$.

The existence of a Stage-0 datum and of finitely many finite lower bounds on $L$, which define
$L_0=L_0(N)$, is Lemma~\ref{2.4:lem-stage-zero-data}; that the parameters of (H3), then $\gamma$,
$\gamma_H$ and, for every $\varrho$, a $\gamma_J$ can be chosen in the stated order for every odd
$L\ge L_0$ is Proposition~\ref{2.4:prop-order-consistent}.
\end{definition}

\begin{definition}[The hypotheses \textup{(H1)}--\textup{(H7)}]\label{2.2:not-hypotheses}
Theorem~\ref{3.1:thm-main} is stated under exactly the following hypotheses. They are read in the
order of choice fixed by the quantifier prefix of Theorem~\ref{3.1:thm-main}, displayed in
Definition~\ref{2.1:not-prefix}. Nothing among them is two-sided.
\begin{enumerate}
\item[\textup{(H1)}] The model is pure $\mathrm{SU}(N)$ Yang--Mills theory, $N\ge 2$, in dimension
$d=4$, with Wilson action $A(U)=\sum_p\bigl[1-\operatorname{Re}\operatorname{tr}U(\partial p)\bigr]$,
on the periodic torus $\mathbb{T}_m$ of side $2L^m$ with lattice spacing $L^{-K}$.
Clauses \textup{(0)} and \textup{(i)} of Theorem~\ref{3.1:thm-main} are asserted for every
$N\ge 2$, and clauses \textup{(ii)}--\textup{(vi)} at $N=2$. The passage of clause \textup{(i)} to
$\mathrm{SU}(N)$, $N\ge 3$, rests on two statements stated in Chapter~\ref{ch:4}. That of clauses
\textup{(iii)}--\textup{(iv)} rests on one stated in Chapter~\ref{ch:5}, and that of clause
\textup{(v)} on one stated in Chapter~\ref{ch:9}. None of them is proved in this book.
\item[\textup{(H2)}] $L$ is odd and $L\ge L_0(N)$. Here $L_0=L_0(N)$ is the one existential constant
of this book; it is bound once in the quantifier prefix and never given a value. Every condition on
$L$ is a lower bound and reads only the data fixed before $L$, which are finitely many numbers
depending only on $d$ and $N$.
\item[\textup{(H3)}] After $L$, Ba{\l}aban's auxiliary parameters $\mathfrak{p}$ are chosen in one
fixed order, the order displayed in Definitions~\ref{2.3:not-stages-first},
\ref{2.3:not-stages-middle} and~\ref{2.3:not-stages-last}. They are
\[
\kappa,\ \kappa_1,\ \dots,\ M_1,\ M_2=L^2M_1,\ M=L^{m'},\ \dots,\ A_1,\ A_0,\ C_0,\ C_1,\
\text{the radii }\alpha,\ C_\beta,\ \varepsilon^{I}_0,\ \varepsilon^{I}_1,\ \dots .
\]
Each parameter falls under one of three cases:
\begin{itemize}
\item it is determined by earlier ones;
\item it is bounded on one side by a positive function of earlier ones;
\item for finitely many of them only, it is confined to a window that the conditions on earlier
choices make non-empty (for instance $M_1<M_2=L^2M_1<M$, non-empty because $M$ is chosen after
$M_1$ and is subject to lower bounds only).
\end{itemize}
Then $\gamma$ is chosen, and then $\gamma_H$, subject to
\begin{equation}\label{2.2:eq-hypotheses-1}
\gamma_H\le\bigl(\gamma^{-2}+c_5+1\bigr)^{-1/2},
\end{equation}
with $c_5=2c_2$ as in Definition~\ref{1.2:def-flow-constants}. No condition bounds $L$, $M$,
$M_1$, $\kappa$, $\kappa_1$, $q$, $p_0$, $K$, $m$, an age, a depth or a number of steps from above.
No condition bounds any coupling from below. Write the cube sizes as powers of $L$:
$M_1=L^{m'_1}$, $M_2=L^{m'_2}$, $M=L^{m'}$. The integers $m'_1$, $m'_2-m'_1$ and $m'-m'_2$ do not
depend on $L$ and are fixed before $L$. This is a convention on the order of choice, not a further
hypothesis.
\item[\textup{(H4)}] $0<g\le\gamma_H$, with $\gamma_H$ as in \eqref{2.2:eq-hypotheses-1}.
\item[\textup{(H5)}] $m\ge m_{\mathbb{T}}(g_-(g))$. Here $g_-(g)=(g^{-2}+B^\sharp)^{-1/2}$, and
\[
m_{\mathbb{T}}(g')=11+m'+\log_L R(g'),
\]
with $m'$ the exponent of $M$ fixed in \textup{(H3)} and $R(g')$ as in
Definition~\ref{1.2:def-flow-constants}. Ba{\l}aban's coarsest cubes then fit
at every scale. This is a lower bound on the volume only.
\item[\textup{(H6)}] $0<g_0\le g_-(g)$, with no lower bound on $g_0$. The number of steps is the
first passage $K=K(g_0,g)$ of Definition~\ref{1.2:def-flow-constants}.
\item[\textup{(H7)}] For clauses \textup{(iii)}--\textup{(vi)}, a radius $\varrho>0$ is fixed after
$\gamma_H$, and then
\[
g\le\gamma_J\le\gamma_H ,
\]
where the threshold $\gamma_J$ depends on $L$, the parameters $\mathfrak{p}$, $\gamma$ and $\varrho$.
Inside clause \textup{(iv)}, further $g\le\gamma_W\le\gamma_J$. The threshold $\gamma_W$ depends on
the quantities fixed before it and on the parameter $\vartheta$ chosen inside that clause. Inside
clause \textup{(v)}, $g\le\gamma_U\le\gamma_H$.
\end{enumerate}
In clause \textup{(iii)}, $\varrho$ is the radius of the disc of complex sources $z$ defined by
$\sum_i|z_i|\,\|f_i\|_\sharp<\varrho$. Constant shifts of $g_0^{-2}$ are the case $f\equiv 1$.
Since $\gamma$ is fixed with the parameters, before $\varrho$, the threshold $\gamma_J$ may depend on
$\gamma$.

One dynamical hypothesis is made, in clause \textup{(vi)} alone: the far-sphere flux hypothesis
$\textup{(IR-fl)}_{\omega_0}$ displayed there, which states that the flux of the lattice rotation
current through far spheres vanishes along one sequence of radii. Displayed beside it there is a
sufficient condition for it: the connected correlation of one far-away stress-tensor insertion with
the observables decays faster than the inverse fourth power of its distance, uniformly in the cutoff.
Beyond it, no dynamical hypothesis is made: nothing is assumed on the sign of the coupling increments
$\beta_{k+1}$, and no convergence, mass gap or clustering property is assumed.

That the order of choice has no cycle is Lemma~\ref{2.4:lem-no-cycle}. That the conditions can be
met is Proposition~\ref{2.4:prop-order-consistent}, together with
Proposition~\ref{2.4:prop-m-freeness}, which is stated in this chapter, with proof incomplete at
the step marked $(\ast)$, and is treated in Chapter~\ref{ch:5}.
\end{definition}

\begin{definition}[Order tables: letters and their kinds, the letters a condition reads, attachment, admissible shape, openers]\label{2.3:def-order-table}
An \emph{order table} consists of finitely many \emph{letters} $\mathsf{a}_1,\dots,\mathsf{a}_\nu$,
listed in their order of choice, and finitely many conditions. Each letter is of one of four
kinds, and the range $R_j$ of the letter $\mathsf{a}_j$ may depend on the earlier letters
$\mathsf{a}_1,\dots,\mathsf{a}_{j-1}$.
\begin{itemize}
\item \emph{Free upward.} The range $R_j\subset(0,\infty)$ is unbounded above. Examples are
$\mathbb{N}$, the odd integers $\ge 3$, and the integer powers of $L$.
\item \emph{Free downward.} The range $R_j$ is an interval $(0,c^0_j]$ with $c^0_j>0$.
\item \emph{Window.} The range $R_j$ is an arbitrary set.
\item \emph{Determined.} The letter is given by $\mathsf{a}_j=\varphi_j(\mathsf{a}_1,\dots,
\mathsf{a}_{j-1})$, where the map $\varphi_j$ is defined whenever the earlier letters lie in their
ranges. The range of a determined letter is its set of values.
\end{itemize}
A \emph{condition} is a relation among the letters. Replace every determined letter in it,
recursively, by its defining expression. The free and window letters that then occur are the
letters the condition \emph{reads}; we require that there be at least one. The condition is
\emph{attached} to the latest letter it reads.

Let $\mathsf{a}_j$ be a free or window letter. An assignment of the free and window letters of
index $<j$ is \emph{admissible} if each $\mathsf{a}_i$ among them lies in its range $R_i$,
computed from the earlier letters, and every condition attached to one of them holds. For such an
assignment, $\mathcal{A}_j\subset R_j$ denotes the set of values of $\mathsf{a}_j$ at which every
condition attached to $\mathsf{a}_j$ holds. For a window letter $\mathsf{a}_j$ the set
$\mathcal{A}_j$ is also called its \emph{window}.

The table has \emph{admissible shape} if the following three properties hold for every
admissible assignment of the earlier free and window letters.
\begin{itemize}
\item Every condition attached to a free upward letter $\mathsf{a}_j$ is a \emph{floor}: there is
$t<\infty$ such that the condition holds for all $\mathsf{a}_j\in R_j\cap[t,\infty)$.
\item Every condition attached to a free downward letter $\mathsf{a}_j$ is a \emph{ceiling}: there
is $c>0$ such that the condition holds for all $\mathsf{a}_j\in R_j\cap(0,c]$.
\item For every window letter $\mathsf{a}_j$ the set $\mathcal{A}_j$ is non-empty.
\end{itemize}
A condition on an earlier letter from which the last property of a window is proved is called
the \emph{opener} of that window.
\end{definition}

\begin{definition}[Caps]\label{2.3:def-cap}
A \emph{cap} is a condition that bounds one of $L$, $M$, $M_1$, $\kappa$, $\kappa_1$, $q$,
$p_0$ or $m$ from above; or that bounds $K$ from above uniformly in $g_0$; or that bounds a
depth, the age of a term (the number of renormalization steps since the term was created) or a
number of renormalization steps from above by a quantity that does not grow with $K$; or that
bounds $g$ or $g_0$ from below. The exclusion of a fixed number of scales at either end of the
range of scales is not a cap.
\end{definition}

A cap would shrink the range over which Theorem~\ref{3.1:thm-main} is asserted, or make its
hypotheses (H1)--(H7) of Definition~\ref{2.2:not-hypotheses} depend on constants chosen later.

\begin{lemma}[Choice by shapes: in a table of admissible shape the letters can be chosen one after
another; adjoining a block]\label{2.3:lem-choice-shapes}
Let an order table, in the sense of Definition~\ref{2.3:def-order-table}, have admissible shape.
For every free or window letter $a_j$ and every admissible assignment of the earlier free and
window letters, the set $\mathcal{A}_j$ is non-empty. If $a_j$ is free upward, then
$\mathcal{A}_j$ contains $R_j\cap[t,\infty)$ for some finite $t$, where $R_j$ is the range of
$a_j$ fixed in Definition~\ref{2.3:def-order-table}. If $a_j$ is free downward, then
$\mathcal{A}_j$ contains $R_j\cap(0,c]$ for some $c>0$.

Hence the letters can be chosen one after another in the order of the table, each from a non-empty
set. Every assignment of all letters satisfying all conditions arises in this way.

Suppose that the table is enlarged by adjoining finitely many letters and conditions, the new
letters being placed together with the original ones in a single order, and suppose that every
condition of the original table reads only letters of the original table. Then every assignment of
the letters of the enlarged table satisfying all of its conditions restricts to an assignment of
the original letters that is admissible for the original table. If in addition the enlarged table
has admissible shape, that is, the new conditions on upward letters are floors, those on downward
letters are ceilings, and at each window letter the window stays non-empty, then the conclusions of
the first two paragraphs hold for the enlarged table.
\end{lemma}

\begin{proof}
Let $a_j$ be free upward, and let $\Gamma_i$, $i\in I$, be the conditions attached to it, where
$I$ is a finite index set. Fix
an admissible assignment of the earlier free and window letters. Together with the determined
letters computed from it through the maps $\varphi_j$ of Definition~\ref{2.3:def-order-table},
this assignment fixes every letter other than $a_j$ that a condition attached to $a_j$ reads.
Since a condition is attached to the latest letter it reads, each $\Gamma_i$ thereby becomes a
condition on $a_j$ alone. Because the table has admissible shape, each $\Gamma_i$ is a floor:
there are finite numbers $t_i$, $i\in I$, such that $\Gamma_i$ holds on $R_j\cap[t_i,\infty)$ for
every $i\in I$. If $I$ is non-empty, then on $R_j\cap[\max_{i\in I} t_i,\infty)$ all of the
conditions $\Gamma_i$, $i\in I$, hold at once, and hence
\begin{equation}\label{2.3:eq-choice-shapes-1}
R_j\cap\bigl[\max_{i\in I} t_i,\infty\bigr)\subset\mathcal{A}_j .
\end{equation}
The left-hand side is non-empty because $R_j$ is unbounded above. This gives the assertion with
$t=\max_{i\in I} t_i$. If $I$ is empty, no condition restricts $a_j$, and $\mathcal{A}_j=R_j$.

Let $a_j$ be free downward, and let $\Gamma_i$, $i\in I$, again be the finitely many conditions
attached to
it, read, as above, as conditions on $a_j$ alone once the earlier free and window letters and the
determined letters computed from them are fixed. By admissible shape they are ceilings; they give
numbers $c_i>0$, $i\in I$, such that $\Gamma_i$ holds for $0<a_j\le c_i$. Hence, with
$R_j=(0,c^0_j]$ the range of $a_j$ fixed in Definition~\ref{2.3:def-order-table},
\begin{equation}\label{2.3:eq-choice-shapes-2}
\bigl(0,\min\bigl(\{c^0_j\}\cup\{c_i: i\in I\}\bigr)\bigr]\subset\mathcal{A}_j ,
\end{equation}
and the left-hand side is non-empty, since the minimum of finitely many positive numbers is
positive. If $a_j$ is a window letter, then $\mathcal{A}_j\ne\emptyset$ by hypothesis, since
admissible shape requires the window at each window letter to be non-empty.

Choosing any element of $\mathcal{A}_j$ extends the given assignment to an admissible assignment
of all free and window letters of index $\le j$, by the definition of $\mathcal{A}_j$ in
Definition~\ref{2.3:def-order-table}. Starting from the first free or window letter and applying
this at each index in turn, induction on $j$ produces a full admissible assignment of the free and
window letters, each value being chosen from a set which is non-empty by the two preceding
paragraphs; the determined letters are then read off through the maps $\varphi_j$. Conversely, the
restriction of an admissible full assignment to the letters of index $\le j$ is admissible. Hence,
at every index $j$, the value that such an assignment gives to $a_j$ lies in the set
$\mathcal{A}_j$ formed with its own earlier values, so every such assignment arises by successive
choice.

Let an assignment of the letters of the enlarged table satisfy all of its conditions. Each
condition of the original table is among these conditions and reads only original letters. It is
therefore evaluated on the same values whether it is read in the enlarged table or in the original
one, and so it holds for the restriction of the assignment to the original letters. Hence that
restriction is admissible for the original table; no property of the shape of the adjoined
conditions is used here. If the enlarged table has admissible shape, the first two paragraphs of
the lemma, proved above for an arbitrary order table of admissible shape, apply to it.
\end{proof}

\begin{definition}[The order of choice of all constants, Stages 0--2: the numbers fixed before $L$,
the blocking factor $L$ and its floors, the numbers determined by $L$]
\label{2.3:not-stages-first}
The words ``free upward'', ``window'', ``opener'' and ``determined'' are used in the sense of
Definition~\ref{2.3:def-order-table}; a \emph{floor} (resp.\ \emph{ceiling}) on a letter is a
lower (resp.\ upper) bound on it by a positive function of earlier letters. In the list of
Stages~0--9 (this list
and Notations~\ref{2.3:not-stages-middle} and~\ref{2.3:not-stages-last}), $\beta_{\mathrm{s}}$
denotes Ba{\l}aban's enlargement constant, which \cite{B-CMP119} and \cite{B-CMP109} call $\beta$.
It is distinct from the coupling increment $\beta_{k+1}$ and from the H\"older exponent $\beta_H$.
The constant $B_3=B_3(d,L)$ is that of \cite{B-CMP102b}. Unless stated otherwise, $\delta_1$
denotes the decay rate of \cite{B-CMP116}; the different rate used in \cite{B-CMP109} is written
out explicitly where it occurs.

\begin{description}
\item[Stage 0 (before $L$; numbers depending on $d$, $N$ and one another).]
The letters of this stage are chosen once, before $L$, under conditions that read $d$, $N$ and
one another only; the tuple so chosen is called the \emph{Stage-0 datum}.
\begin{itemize}
\item $r\in\mathbb{N}$ with $r\ge2$, the exponent of the length $R(g)$, as in \cite{B-CMP119}.
\item $r_0\ge r$, the exponent of \cite{B-CMP122b}.
\item $p_0\in\mathbb{N}$, free upward, with $p_0\ge5r$ \cite{B-CMP119} and the floor
$p_0\ge(d+5)r_0+8r+3$, among others.
\item $p_1$, a window, under three kinds of condition:
\begin{itemize}
\item ``We assume that $2p_1-(d+5)r_0>p_0$'' \cite{B-CMP122b};
\item $p_1<p_0$ \cite{B-CMP122a};
\item finitely many lower bounds on $p_0-p_1$ by numbers at most $4r+1$.
\end{itemize}
The floor $p_0\ge(d+5)r_0+8r+3$ is the opener. With $p_1:=p_0-4r-1$,
\begin{equation}\label{2.3:eq-stages-first-1}
2p_1-(d+5)r_0-p_0=p_0-8r-2-(d+5)r_0\ge1 ,
\end{equation}
and $p_0-p_1=4r+1$. Raising $p_0$ only widens the window.
\item $q:=q_0=q_1$, free upward, with floors such as $q\ge p_0+7r+1$ and $q\ge3p_0+3r+1$. In
\cite{B-CMP119}, $q_0,q_1$ are bounded only from below: ``$q_0,q_1$ are integers greater than 1,
$C_0,C_1$ are sufficiently large positive numbers''.
\item $\kappa_0\ge17$, the exponent of \cite[(2.31)]{B-CMP119}. For it, \cite{B-CMP119} states
``$\kappa_0$ can be chosen arbitrarily large, similarly as $\kappa$''.
\item $v\ge2p_0+dr_0$, where $v$ is the exponent that \cite[p.~363]{B-CMP122b} calls $\nu$:
``we need $\nu\geqq2p_0+dr_0$''.
\item Fractions and numbers:
\begin{itemize}
\item $\beta_{\mathrm{s}}=1/4$ (``e.g., we can take $\beta=1/4$'' \cite{B-CMP119});
\item the fraction $\beta_0\in(0,\frac12]$ of \cite[(2.6)]{B-CMP119}, free in
\cite{B-CMP119} and fixed here at $\beta_0=1/7$; its admissibility is the ceiling
$2c_2\gamma^2\le1-(1+\beta_0)^{-2}$ of Stage~8 (Notation~\ref{2.3:not-stages-last});
\item the H\"older exponent $\beta_H\in(0,1)$, fixed before every constant that depends on it;
\item the number $\alpha_6$ of \cite{B-CMP116}, depending on $d$ only.
\end{itemize}
\item Further windows.
\begin{itemize}
\item The window of one further exponent, lying between functions of $(r,r_0,p_0,q)$, is opened
by a floor on $q$.
\item The window of one further integer exponent, written $\nu_I$, is the non-empty finite set
$\{r+1,\dots,8r\}$.
\item Several numbers range in fixed non-empty intervals between numbers, for example
$[1/2,8/9]$, $(8/7,3/2)$ and $(3/5,1)$.
\end{itemize}
Two further windows need a comment:
\begin{itemize}
\item one number carries only ceilings at Stage~0, so its admissible set there is an interval
$(0,c]$ with $c>0$; the lower bound that would complete its window is carried instead by a floor
on $L$ in Stage~1 whose right member involves this number;
\item one interval between numbers is non-empty provided that an exponent $\kappa'$, given by a
decay estimate outside the table, satisfies $\kappa'>5/12$. Only this inequality, and no
numerical value of $\kappa'$, opens the window. The inequality $\kappa'>5/12$ is an input of
Proposition~\ref{2.4:prop-order-consistent}, where it is displayed as an antecedent of standing
``stated; proof incomplete''.
\end{itemize}
\end{itemize}
\item[Stage 1 ($L$ odd, free upward; every floor reads Stage 0 only).]
\begin{itemize}
\item $L>11$ (``where $L$ is an odd, positive integer $>11$'' \cite{B-CMP109}).
\item $(1+7\beta_{\mathrm{s}})L^{-2}\le1$ \cite{B-CMP109}.
\item $L^2\ge2(1+\beta_{\mathrm{s}})(1+\beta_0)^2/(1-\beta_{\mathrm{s}}/8)$, together with
finitely many further floors by Stage-0 numbers arising in the arguments.
\item $L\ge4r_\star$, where $r_\star$ is a positive Stage-0 number. Under this floor the width
$\delta(L):=(1-2r_\star/L)/10$ lies in $[1/20,1/10)$.
\end{itemize}
Let $L_0'$ be the least odd integer above all floors on $L$ of the table of Stages~0--9; these
floors read the Stage-0 datum only, so $L_0'$ reads that datum only, and once the datum is fixed
it is a number $L_0'=L_0'(N)$. The constant $L_0=L_0(N)$ of hypothesis (H2) of
Notation~\ref{2.2:not-hypotheses} satisfies $L_0(N)\ge L_0'$ and also covers the floors on $L$ of
the conditions proper to clause (v) (Lemma~\ref{2.4:lem-stage-zero-data}). Neither threshold is
evaluated, and no numerical value is asserted for either.
\item[Stage 2 (determined by Stage 0 and $L$).] The following numbers are determined:
\begin{itemize}
\item the width $\delta=\delta(L)$;
\item the flow numbers $c_0=(11N^2/12\pi^2)\log L$, $c_2$, $c_5=2c_2$ and $(8c_2+4)^{-1/2}$ of
the running-coupling bounds;
\item the numbers $M_1^{\mathrm{thr}},\delta_0,a_0,B_0$ of \cite[Theorem~3.1]{B-CMP99b};
\item the decay rate $\delta_1$ of \cite{B-CMP116}, a number depending on $(d,L)$;
\item counting numbers.
\end{itemize}
No floor on $L$, $\kappa$, $\kappa_1$ or $M_1$ depends on $c_2$ or $c_5$.
\end{description}
\end{definition}

\begin{proof}
We verify those assertions of the list that are explicit: the window of $p_1$, the finite
windows, the width $\delta(L)$, and the last sentence of Stage~2.

\emph{The window of $p_1$.} Put $p_1:=p_0-4r-1$. Then
\begin{equation*}
2p_1-(d+5)r_0-p_0=2p_0-8r-2-(d+5)r_0-p_0=p_0-8r-2-(d+5)r_0 ,
\end{equation*}
and the floor
$p_0\ge(d+5)r_0+8r+3$ gives that this quantity is at least $1$; this is
\eqref{2.3:eq-stages-first-1}, so $2p_1-(d+5)r_0>p_0$. Next $p_0-p_1=4r+1>0$, so $p_1<p_0$.
Finally each of the finitely many lower bounds on $p_0-p_1$ is by a number at most $4r+1$, and
$p_0-p_1=4r+1$, so all of them hold. Thus the window of $p_1$ is non-empty once the opener holds.
For fixed $(d,r,r_0)$ the window of $p_1$ is the set of $p_1$ with
$p_1>\tfrac12\bigl(p_0+(d+5)r_0\bigr)$ and $p_1\le p_0-b$ for each of the finitely many numbers
$b$ bounding $p_0-p_1$ from below, together with $p_1<p_0$. When $p_0$ increases by $t>0$, the
lower end increases by $t/2$ and every upper end by $t$; hence raising $p_0$ only widens the
window, and the floor $p_0\ge(d+5)r_0+8r+3$ stays satisfied.

\emph{The finite windows.} Since $r\ge2$, we have $r+1\le8r$, so $\{r+1,\dots,8r\}$ is non-empty and
finite. The intervals $[1/2,8/9]$, $(8/7,3/2)$ and $(3/5,1)$ are non-empty because
$1/2<8/9$, $8/7<3/2$ and $3/5<1$. The choice $\beta_0=1/7$ lies in $(0,\frac12]$.

\emph{The width.} The floor $L\ge4r_\star$ gives $2r_\star/L\le1/2$, hence
$1-2r_\star/L\ge1/2$ and $\delta(L)\ge1/20$; since $r_\star>0$ and $L>0$, we have
$1-2r_\star/L<1$ and $\delta(L)<1/10$.

\emph{The last sentence of Stage~2.} Consider first the floors on $L$. Every floor on $L$ of the
table is listed at Stage~1 and reads Stage-0 numbers only, and the floors on $L$ proper to
clause~(v) read Stage-0 numbers only by Lemma~\ref{2.4:lem-stage-zero-data}; since $c_2$ and
$c_5=2c_2$ are determined only at Stage~2, by Stage~0 and $L$, no floor on $L$ depends on $c_2$ or
$c_5$. For $\kappa$, $\kappa_1$ and $M_1$ the assertion is
a property of their floors, which are listed in Notation~\ref{2.3:not-stages-middle}: although
$c_2$ and $c_5$ are determined at Stage~2, before these letters are chosen, no floor on $\kappa$,
$\kappa_1$ or $M_1$ in that list has $c_2$ or $c_5$ in its right member.
\end{proof}

\begin{definition}[The order of choice, Stages 3--6: $\kappa$, $\kappa_1$, $M_1$ with the
constants fixed after it, and $M$]
\label{2.3:not-stages-middle}
We continue the order of choice begun with Stages~0--2 in
Definition~\ref{2.3:not-stages-first}, in the terms of Definition~\ref{2.3:def-order-table}:
each stage is headed by the letter it fixes first and the side on which that letter is free,
and the conditions listed under a stage are the ones attached to the letters fixed there.
\begin{itemize}
\item Stage~3 ($\kappa$, free upward). Let $\kappa_{\mathrm{thr}}$ be the maximum of the
finitely many floors of this stage. Each of them reads only $d$, Stage-0 numbers and $L$, the
latter directly (as in $\kappa\ge10\log L+1$) or through $\delta=\delta(L)$ (as in
$\kappa\ge C_\delta(d)/\delta$, which also reads the Stage-0 number $\alpha_6$); none reads
$M_1$, $M$ or any later letter. The number $\kappa_{\mathrm{thr}}$ is the threshold that
\cite[Theorem~3]{B-CMP109} denotes $\kappa_0$, an object different from the exponent
$\kappa_0$ of Stage~0, and $\kappa\ge\kappa_{\mathrm{thr}}$. Among these floors are:
\begin{itemize}
\item $\kappa\ge10\log L+1$. For $n>j$ this gives
$\tfrac12\kappa(n-j)\ge(5\log L+\tfrac12)(n-j)$, hence
\[
\exp\bigl(-\tfrac12\kappa(n-j)\bigr)\le e^{-(n-j)/2}L^{-5(n-j)}<(L^jL^{-n})^5 ,
\]
the last step because $e^{-(n-j)/2}<1$; this is the inequality used in \cite{B-CMP119}.
\item $\kappa\ge C_\delta(d)/\delta$ (``for $\delta\kappa$ sufficiently large''
\cite{B-CMP116}).
\item Floors depending on $d$ and Stage-0 fractions alone.
\end{itemize}
\item Stage~4 ($\kappa_1$, free upward). Among the conditions are
$\frac12(\kappa_1-1)\ge2L\kappa$ (``Of course, we assume that
$\frac12(\kappa_1-1)\geqq2L\kappa$'' \cite{B-CMP116}) and
$\kappa_1\ge\max\{16\kappa+1,4L\kappa+1\}$. The others, finitely many, are floors on
$\kappa_1$ by $d$, $L$, $\kappa$ and counting numbers. These are always inequalities;
$\kappa_1$ is never defined by an equation.
\item Stage~5 ($M_1\in L^{\mathbb{N}}$, free upward, after $\kappa_1$).
\begin{itemize}
\item $M_1\ge M_1^{\mathrm{thr}}$ (``$M_1$ is the size of large cubes for which all the
theorems of the previous papers \dots\ are valid'' \cite{B-CMP119}).
\item $\delta_1M_1\ge c\,\kappa_1$, with $c\ge1$ depending on $d$. This is the requirement
``$\delta_1M\geqq\kappa_1$'' of \cite{B-CMP116}, applied on cubes of side $M_1$.
\item Then a further size parameter $M_1^\sharp$, free upward, after $L$ and $M_1$: the size
of the cubes on which the statement on the two lowest tower levels, stated in
\cite[(3.35), (3.37)]{B-CMP99b} without proof, is used. Its floor is
$M_1^\sharp\ge5L\,M_1$, relaxed to $M_1^\sharp\ge3L\,M_1$ on one sub-family of those cubes,
singled out where that statement is applied. This floor reads only $L$ and $M_1$.
\item Then $M_2:=L^2M_1$ (``we can take $M_2=L^2M_1$'' \cite{B-CMP119}); the determined
constants $B_1$, $B_3$, $B_5$, $a_1$ of \cite{B-CMP102b}; and the length parameter $R_1$ of
that paper, a power of $L$ chosen upward.
\item Then the auxiliary kernel and random-walk constants of the fluctuation-operator
expansion, all determined by $(d,L,M_1)$. Among them are the constants $(B_W,\delta_W)$ and
$\mu_0$ that appear in Proposition~\ref{2.4:prop-m-freeness}, whose proof is incomplete at
one step.
\item Then two width parameters $\Theta_1\ge1$ and $\varpi_T\in\mathbb{N}$, free upward, at
or above finitely many floors that read those kernel constants.
\item Then an auxiliary length $r'$, chosen upward, and finitely many quantities determined
by it.
\item Then a second auxiliary length $\ell$, a power of $L$ chosen upward.
\item Then one constant $\hat c$, chosen downward below a ceiling by earlier letters. In
Stage~7 of Definition~\ref{2.3:not-stages-last}, $\hat c$ determines
$\gamma_2:=\hat c\,(LM)^{-4}$.
\end{itemize}
The kernel constants placed here are the constants of the kernel bounds for backgrounds in
the global regularity class, hence functions of $(d,L,M_1)$ that do not depend on $M$: for
$(B_W,\delta_W)$ and $\mu_0$ by the kernel bound in the global regularity class recalled
after Proposition~\ref{2.4:prop-m-freeness}; for the other kernel
constants, among them the rate read by the ceiling on $\hat c$, by hypothesis~(A1) of
Proposition~\ref{2.4:prop-order-consistent}. So these constants,
and the letters $\Theta_1$, $\varpi_T$ and $\hat c$ whose conditions read them, stand before
$M$ by the form of their conditions, and the three floors on $M$ of Stage~6 that read $\mu_0$
and $\delta_W$ are floors by their form. Proposition~\ref{2.4:prop-m-freeness}, whose proof
is incomplete at one step, is the statement that the bounds hold with these same constants
at the local backgrounds at which they are applied; it is used for no placement and for no
side.
\item Stage~6 ($M=L^{m'}$, free upward).
\begin{itemize}
\item $M\ge L^3M_1$. With $M_2:=L^2M_1$ this reads $M\ge LM_2$ (``$M$ will be chosen much
larger than $M_2$'' \cite{B-CMP119}).
\item $M\ge M(\kappa)$ \cite[Theorem~3]{B-CMP109}.
\item $\delta_0M\ge\kappa$ and $\delta_1M\ge\kappa_1$ \cite{B-CMP116}, with $\delta_1$ the
rate of that paper, and $e^{-\delta_0M/3}e^{16\kappa_1}<1$.
\item $M\ge\max\{\varpi_T,3\}$.
\item With the $M$-free constants $\mu_0$, $\delta_W$ of Stage~5 (those of the bounds at
local data by Proposition~\ref{2.4:prop-m-freeness}, whose proof is incomplete at one step):
\begin{gather*}
\mu_0M\ge2\kappa_1,\qquad \delta_WM\ge2\kappa_1,\\
\delta_WM/4\ge32\kappa_1+2\log(LM)+C(d,L).
\end{gather*}
\end{itemize}
No floor on $M$ is imposed at the rate
$\delta_1^{\mathrm{alt}}:=\frac12\min\{\delta_0,\kappa M^{-1}\}$ that \cite{B-CMP109} uses
in its bound \cite[(5.10)]{B-CMP109}. At that rate
$\delta_1^{\mathrm{alt}}M\le\frac12\kappa M^{-1}M=\kappa/2$, and $\kappa/2<\kappa_1$ because
$\kappa_1\ge16\kappa+1$ by Stage~4; so such a floor would be empty.

The only inequality of the form $M\le F$, with $F$ free of $M$, that enters the construction
is the inequality $M\le R_1M_1a_1/\varepsilon^{\mathrm{run}}$ of \cite{B-CMP102b}. There
$\varepsilon^{\mathrm{run}}$ is the actual size of the small fields at the running coupling,
an object different from the parameter $\varepsilon^{\mathrm I}_1$ of Stage~7 of
Definition~\ref{2.3:not-stages-last}. That
inequality is not a condition of the table in this form. Writing
$\varepsilon^{\mathrm{run}}(x)$ for that size at inverse coupling $x$, the table imposes
instead the stronger ceiling
\begin{equation}\label{2.3:eq-stages-middle-1}
\sup_{x\ge\gamma^{-2}}O(1)\,LM\,\varepsilon^{\mathrm{run}}(x)\le R_1M_1a_1 ,
\end{equation}
with $L$, $M$, $M_1$, $R_1$, $a_1$ fixed earlier and the constant $O(1)$ taken at least $1$
(enlarging it only strengthens the ceiling). Since then $O(1)L\ge1$, \eqref{2.3:eq-stages-middle-1}
gives $M\varepsilon^{\mathrm{run}}(x)\le O(1)LM\varepsilon^{\mathrm{run}}(x)\le R_1M_1a_1$,
so the ceiling implies that inequality at every level. Since
$\varepsilon^{\mathrm{run}}(x)\to0$ as $x\to\infty$, this is a ceiling on $\gamma$ of the
form of Stage~8 of Definition~\ref{2.3:not-stages-last} and is placed there. Every other
inequality in which $M$ stands on the smaller side, such as $(LM)^4\alpha_i\le1$,
$O(1)M\alpha_1\le\beta_{\mathrm{s}}$ or the condition displayed under $C_1$ in Stage~7 of
Definition~\ref{2.3:not-stages-last}, is
attached to a letter chosen after $M$ and bounds that letter on its free side.
\end{itemize}
\end{definition}

\begin{definition}[The order of choice, Stages 7--9: the remaining parameters, the thresholds
$\gamma$ and $\gamma_H$ with their ceilings, and the variables $g$, $m$, $g_0$, $K$]
\label{2.3:not-stages-last}
The order of choice whose Stages~0--6 are fixed in Definitions~\ref{2.3:not-stages-first} and~\ref{2.3:not-stages-middle} is completed by the following three stages. The words
``free upward'', ``free downward'', ``determined'' and ``attached'' are used as in
Definition~\ref{2.3:def-order-table}.

\smallskip\noindent
Stage~7 (after $M$, in this order).
\begin{itemize}
\item $A_1>0$, then $A_0$ upward, with
\[
A_0/A_1\ge\max\{c(d)(1+\beta_0)B_3L^{d+2},\,c'(d,L)M^2\}
\]
(``we assume that the constant $A_0$ is much larger than $A_1$'' \cite{B-CMP119}).
Here $B_3$ is the constant of \cite[p.~277]{B-CMP109} entering the comparison of radii
quoted below; since this floor is attached to $A_0$, $B_3$ is a letter fixed before $A_0$.
Further, $\beta_0>0$ is the constant of \cite{B-CMP119} whose smallness is used in Stage~8
below, and $c(d)$, $c'(d,L)$ are numbers depending only on the arguments shown.
\item $C_0$ upward, with floors by multiples of $A_0$ and $A_1$. Then $C_1$ upward, with
\begin{equation}\label{2.3:eq-stages-last-1}
C_1>2B_3^2O(1)M\,C_0 .
\end{equation}
Here $O(1)$ is the unspecified constant of \cite{B-CMP109}. The inequality comes from
``Assume now that $B_3^2O(1)M\alpha_0<\frac12\alpha_1$'' \cite[p.~277]{B-CMP109}. At
$q_0=q_1=q$ the running radii $\alpha_{i,n}=g_nC_i(\log g_n^{-2})^q$ of
\cite[(2.28)]{B-CMP119}, where $g_n$ is the coupling at level $n$, satisfy
$B_3^2O(1)M\alpha_{0,n}<\frac12\alpha_{1,n}$ exactly when \eqref{2.3:eq-stages-last-1}
holds: both sides carry the common factor $g_n(\log g_n^{-2})^q$, which is positive
for $0<g_n<1$, and
dividing by it and multiplying by $2$ turns the comparison into
\eqref{2.3:eq-stages-last-1}. No condition on $C_0$ depends on $C_1$.
\item The radii, all free downward. They are chosen in the order $\alpha'_1,\alpha'_0$,
then $\alpha_1$, $\alpha_0$, $\alpha_2$. For the primed radii, \cite{B-CMP116} states:
``with constants $\alpha'_0,\alpha'_1$ much bigger than $\alpha_0,\alpha_1$, therefore we
can restrict them''. The representative conditions are:
\begin{itemize}
\item $(LM)^4\alpha_i\le1$ (``$(LM)^4\alpha_0$, $(LM)^4\alpha_1$, $(LM)^4\alpha_4$,
$(LM)^4\gamma_2$ are bounded by a constant independent of $M$, for example by $1$''
\cite{B-CMP116});
\item $O(1)M\alpha_1\le\beta_{\mathrm{s}}$ \cite[p.~278]{B-CMP109}, where
$\beta_{\mathrm{s}}$ is the constant bounding $O(1)M\alpha_1$ in the condition at
\cite[p.~278]{B-CMP109}, a letter fixed before the radii;
\item $B_3^2O(1)M\alpha_0<\frac12\alpha_1$ \cite[p.~277]{B-CMP109};
\item ``We assume that $(4B_3^2O(1)M)^2\alpha_0\leqq\beta$'' \cite[p.~278]{B-CMP109};
\item $\alpha_2\le\alpha_1/8$ and $\alpha_2\le\alpha_0/C_{\mathrm{IV}}$, with
$C_{\mathrm{IV}}=C_{\mathrm{IV}}(d,L)$ a number depending only on $d$ and $L$.
\end{itemize}
\item The absolute constant $E_0$ of \cite{B-CMP116}. Then $\alpha_4$ downward, with
$(LM)^4\alpha_4\le1$. Then the determined constants, among them the constant $C_\beta$ of
the bound on the coupling increment. Its rate is ``determined by $\delta_0$, $\kappa$, and
$M$ (e.g., $\delta_1=\frac12\min\{\delta_0,\kappa M^{-1}\}$)'' \cite{B-CMP109}, so
$C_\beta$ comes after $M$. This fixes the triple $(c_0/4,\,c_5,\,C_\beta+1)$.
\item $\varepsilon^{\mathrm I}_0$, then $\varepsilon^{\mathrm I}_1$, downward, under the
restrictions listed in \cite{B-CMP116} (``Lemma 3. Under all the above restrictions on the
constants $M,\kappa,\kappa_1,\alpha_0,\alpha_1,\alpha_4,\alpha_6,\gamma_2,\gamma,
\varepsilon_1$ \dots'').
\end{itemize}

\smallskip\noindent
Stage~8 ($\gamma$, free downward, last; then $\gamma_H$). By
\cite[Theorem~3]{B-CMP109}, ``The constant $\gamma$ depends on all other constants''.
Every condition attached to $\gamma$ is a ceiling of one of two forms: $\gamma\le c$ with
$c>0$, or
\[
\sup_{x\ge\gamma^{-2}}h(x)\le c ,
\]
where $c>0$ and $h$, with $h(x)\to0$ as $x\to\infty$, are built from letters of
Stages~0--7. Examples:
\begin{itemize}
\item $\gamma\le(8c_2+4)^{-1/2}$.
\item $\bar\varsigma(\gamma)\le\min\{c_0/4,1\}$. Here
$\bar\varsigma(\gamma)=C'\gamma^{\kappa_0-8}$ is the bound on the large-field part of the
coupling increment, Chapter~\ref{ch:4}, valid for couplings at most $\gamma$, where $C'$ is a
constant of Stage~7 and $\kappa_0$ is the exponent of Stage~0 fixed in
Definition~\ref{2.3:not-stages-first}. Since
$\lambda_\sharp=c_0/2-\bar\varsigma$ and
$B_\sharp=C_\beta+\bar\varsigma$, this gives $c_0/4\le\lambda_\sharp$ and
$B_\sharp\le C_\beta+1$.
\item $(C_\beta+1)\gamma^2\le1$ and $2c_5\gamma^2\le1$.
\item $2c_2\gamma^2\le1-(1+\beta_0)^{-2}$. With $y=c_5\gamma^2=2c_2\gamma^2$ it gives
$1+y\le(1-y)^{-1}\le(1+\beta_0)^2$: the first inequality holds because
$(1+y)(1-y)=1-y^2\le1$ and $0\le y<1$, the second because the condition reads
$1-y\ge(1+\beta_0)^{-2}$. This is the form in which the statement ``$\beta_0>0$ can be
chosen arbitrarily small, if $g$ is sufficiently small'' of \cite{B-CMP119} is used.
\item $\sup_{x\ge\gamma^{-2}}C_ix^{-1/2}(\log x)^q\le\alpha_i$ for $i=0,1$, so that the
running radii never exceed the absolute ones; indeed, at $x=g_n^{-2}$ the function under the
supremum is $\alpha_{i,n}$.
\item The ceiling from \cite{B-CMP102b} displayed after Stage~6 in
Definition~\ref{2.3:not-stages-middle}.
\end{itemize}
Every ceiling of Stage~7 that bounds $\alpha_0$ or $\alpha_1$ alone by earlier letters is
also imposed in this way, with the running radii $\alpha_{0,n}$, $\alpha_{1,n}$ in place of
$\alpha_0$, $\alpha_1$; this includes the ceilings whose right member contains the primed
radii $\alpha'_0,\alpha'_1$. Only the three conditions of Stage~7 that compare two of the
unprimed radii are not imposed at running radii as ceilings on $\gamma$:
\begin{itemize}
\item The comparison $B_3^2O(1)M\alpha_0<\frac12\alpha_1$: at running radii it is
\eqref{2.3:eq-stages-last-1} (Stage~7), free of the coupling.
\item The comparisons $\alpha_2\le\alpha_1/8$ and $\alpha_2\le\alpha_0/C_{\mathrm{IV}}$ are
conditions on $\alpha_2$ at the absolute radii only. The radius $\alpha_2$ has no running
version, and with $\alpha_{1,n}$ or $\alpha_{0,n}$ in place of $\alpha_1$ or $\alpha_0$ each
of them would be a lower bound on the coupling.
\end{itemize}
Every ceiling that involves the flow constants $(\lambda_\sharp,c_5,B_\sharp)$ of
Definition~\ref{1.2:def-flow-constants} is imposed at the triple $(c_0/4,c_5,C_\beta+1)$,
which is fixed before $\gamma$. Then $\gamma_H$ is chosen subject to the terminal ceiling
\[
\gamma_H\le(\gamma^{-2}+c_5+1)^{-1/2} .
\]
Inside clause (iii), $\gamma_J\in(0,\gamma_H]$ is then chosen downward under finitely many
ceilings.

\smallskip\noindent
Stage~9 (the couplings, the volume and the number of steps). The remaining choices are:
\begin{itemize}
\item $g\in(0,\gamma_H]$, and $g\le\gamma_J$ inside (iii);
\item $m\ge m_{\mathbb{T}}(g_-(g))$;
\item $g_0\in(0,g_-(g)]$;
\item $K$, the first passage $K(g_0,g)$ of Definition~\ref{1.2:def-flow-constants}.
\end{itemize}
Quantities indexed by the level $k$ (auxiliary lengths, counts and schedules of the
construction) are determined at each level from $x_0,\dots,x_k$, evaluated at
$\check x_k:=\min_{i\le k}x_i$. Some of them lie in two-sided ranges, which are non-empty at
every level $k$ with $\check x_k>\gamma^{-2}+c_5+1$. If $g_k<g$ for all $k\le K$, as
clause (0) asserts, then
\[
x_k=g_k^{-2}>g^{-2}\ge\gamma_H^{-2}\ge\gamma^{-2}+c_5+1\qquad\text{for every }k\le K,
\]
the middle inequality by $g\le\gamma_H$ and the last by the terminal ceiling; hence
$\check x_k>\gamma^{-2}+c_5+1$, as a minimum of finitely many numbers each exceeding this
bound, and every condition on $x$ placed in Stage~8 holds at every level, since each such
condition bounds a function on the half-line $x\ge\gamma^{-2}$, which contains $x_k$ and
$\check x_k$. These quantities impose no condition on a parameter: the only conditions they
require are the floors on $x$ of Stage~8, and no letter or condition of Stages~0--8 reads
them. The ceilings on $\gamma$ read only letters of Stages~0--7, and those on $\gamma_H$
only letters of Stages~0--7 and $\gamma$.
\end{definition}

\begin{remark}[The thresholds chosen last and the clause-internal quantifiers]\label{2.3:rem-thresholds}
In the quantifier prefix of Theorem~0 (Notation~\ref{2.1:not-prefix}), Ba{\l}aban's auxiliary
parameters $\mathfrak{p}$ are chosen after $L$, in one fixed order. Then $\gamma$ is chosen, and then
$\gamma_H$. Among the ceilings on $\gamma$ is
\begin{equation}\label{2.3:eq-thresholds-1}
\gamma\le(8c_2+4)^{-1/2}\le\tfrac12 .
\end{equation}
The threshold $\gamma_H$ is chosen after $\gamma$ and satisfies
\begin{equation}\label{2.3:eq-thresholds-2}
\gamma_H\le(\gamma^{-2}+c_5+1)^{-1/2}<\gamma ,
\end{equation}
with $c_2$ and $c_5=2c_2$ as in Definition~\ref{1.2:def-flow-constants}. The prefix then continues
with the quantifiers on $g$, $m$ and $g_0$ of Notation~\ref{2.1:not-prefix}.

Inside the clauses, and only there, further quantifiers occur; each is an upper bound on a threshold
chosen last or a one-sided floor:
\begin{itemize}
\item in (iii) and (vi):
\begin{equation*}
\forall\varrho>0\ \exists\gamma_J\in(0,\gamma_H]\ \forall g\le\gamma_J ;
\end{equation*}
\item in (iv):
\begin{equation*}
\forall\varrho\ \exists\gamma_J\ \forall\vartheta\in(0,1]\ \exists\gamma_W\in(0,\gamma_J]\
\forall g\le\gamma_W ;
\end{equation*}
\item in (v):
\begin{equation*}
\forall\varrho\ \exists\gamma_U\in(0,\gamma_H]\ \forall g\le\gamma_U ,
\end{equation*}
then a torus exponent, which we denote $m_0$, with $m_0\ge0$, and the depth $n$, with the bare
inverse coupling tuned, $g_0^{-2}=a_n-\hat{w}$,
instead of first passage; the first-passage quantifier returns only in (v-c)(2) and (ii-$\mathbb{T}$-b).
\end{itemize}
In the form of hypothesis (H7) of Notation~\ref{2.2:not-hypotheses}: for clauses (iii)--(vi) the
radius $\varrho>0$ is fixed after $\gamma_H$; then, in clauses (iii), (iv) and (vi),
$g\le\gamma_J(L,\mathfrak{p},\gamma;\varrho)\le\gamma_H$;
inside (iv) further $g\le\gamma_W\le\gamma_J$, where $\gamma_W$ depends on the data on which $\gamma_J$
depends and on $\vartheta$; inside (v),
$g\le\gamma_U\le\gamma_H$.

The radius $\varrho$ belongs neither to the parameter tuple nor to the Stage-0 datum. No condition
attached to a parameter chosen before $\gamma$ (Definition~\ref{2.3:not-stages-last} and the stages
preceding it), to $\gamma$ or to $\gamma_H$ depends on $\varrho$. The conditions
that depend on it are the ceilings on $\gamma_J$; the ceiling on $g$ of clause (iv), which lies below
$\gamma_J$; the ceilings on $\gamma_U$; and, for the identification statement of clause (v) only, the
floor $\varrho\ge6(B_\sharp+1)$ together with a further ceiling on $g$ that depends on $\varrho$ and is
at most $\gamma_U$.

Clause (iv) rests on clause (iii) and adds conditions of its own, all one-sided and all chosen after
the thresholds of clause (iii): for a fixed shape of the test functions, a ceiling on $g$ fixed inside
the clause, below the ceilings of clause (iii); an upper threshold on the radius $\mathsf{d}$ of the
balls containing the supports, depending on $g$ only through $\log g^{-2}$ (equivalently, a floor on
the depth of the scale of $\mathsf{d}$, which bounds no depth from above); the exclusion of finitely
many scales nearest the cutoff, their number
bounded below by a displayed function of the coupling and $L$; and the floor on the volume of clause
(iii), the volume being otherwise free. These conditions change no condition of the order of choice
and no hypothesis of clauses (0), (i) and (iii).
\end{remark}

\begin{proof}
The second inequality in \eqref{2.3:eq-thresholds-1} is equivalent to $8c_2+4\ge4$, that is to
$c_2\ge0$, which holds by Definition~\ref{1.2:def-flow-constants}. The strict inequality in
\eqref{2.3:eq-thresholds-2} is equivalent to $\gamma^{-2}+c_5+1>\gamma^{-2}$, that is to
$c_5+1=2c_2+1>0$, which again follows from $c_2\ge0$; hence $\gamma_H<\gamma\le\frac12$.

Every clause-internal ceiling is nested below $\gamma_H$:
\begin{gather*}
g\le\gamma_J\le\gamma_H\quad\text{in (iii) and (vi)},\\
g\le\gamma_W\le\gamma_J\le\gamma_H\quad\text{in (iv)},\\
g\le\gamma_U\le\gamma_H\quad\text{in (v)}.
\end{gather*}
Hence every coupling admitted inside a clause satisfies $0<g\le\gamma_H$,
which is hypothesis (H4), and the quantities $g_-(g)$, $m_{\mathbb{T}}(g_-(g))$ and, where first passage
is used, $K(g_0,g)$ of Definition~\ref{1.2:def-flow-constants} are defined for it exactly as in the
prefix.

Finally, in each of clauses (iii)--(vi) the quantifier on $\varrho$ stands after $\gamma_H$, and
$\gamma_H$ stands
after $\gamma$, after $\mathfrak{p}$, after $L$ and after the Stage-0 datum. A condition fixed at a
place of the order preceding the introduction of $\varrho$ is a condition on the parameters already
chosen there, and therefore cannot depend on $\varrho$. This applies to every condition attached to a
parameter chosen before $\gamma$, to $\gamma$ and to $\gamma_H$. A condition depending on $\varrho$ can
therefore only be
one of those placed after it; among these, the ones that do depend on $\varrho$ are the ceilings on
$\gamma_J$, the ceiling of clause (iv) below $\gamma_J$, the ceilings on $\gamma_U$, and, in the
identification statement of clause (v), the floor $\varrho\ge6(B_\sharp+1)$ with its ceiling on $g$
at most $\gamma_U$. The conditions of clause (iv) are placed after the thresholds of clause (iii), and
they are of three kinds: a further ceiling on $g$, below the ceilings of clause (iii), which only
shrinks the range of the universally quantified coupling and leaves every existing condition as it
stands; the upper threshold on $\mathsf{d}$ and the exclusion of finitely many scales nearest the
cutoff, which restrict only the test functions admitted in the clause and the scales at which the
clause is read; and the volume floor of
clause (iii) itself, unchanged. Hence they change no condition of the order of choice and no
hypothesis of clauses (0), (i) and (iii).
\end{proof}

\begin{lemma}[Stage-0 data exist; the lower bounds on $L$ are finitely many and read Stage-0
numbers only; the constant $L_0(N)$]\label{2.4:lem-stage-zero-data}
In this lemma the \emph{conditions considered} are the conditions of the order of choice of
Definitions~\ref{2.3:not-stages-first}, \ref{2.3:not-stages-middle}
and~\ref{2.3:not-stages-last} that are used by the clauses of Theorem~0. Conditions introduced
only for the block-law comparison statement of Remark~\ref{3.2:rem-block-laws}, which no clause
of Theorem~0 uses, are not among them. A \emph{Stage-0 datum} is a
choice of the finitely many integer exponents and numbers of Stage~0. These depend only on $d$,
on $N$ and on one another. Let $\kappa'$ be the Stage-0 decay exponent of the kernel estimate.
\begin{itemize}
\item[(a)] For every integer $r\ge2$ there is a Stage-0 datum with this $r$ that meets every
Stage-0 condition considered; no condition on $\kappa'$ is used in (a).
\item[(b)] The displayed order carries one further Stage-0 window, an interval between numbers
whose opener is the inequality $\kappa'>5/12$. It serves only the block-law comparison
statement of Remark~\ref{3.2:rem-block-laws} and is not among the conditions considered. If
$\kappa'>5/12$, that window is non-empty.
\item[(c)] Fix a Stage-0 datum as in (a). The conditions considered that are attached to $L$,
in the sense of Definition~\ref{2.3:def-order-table}, are finitely many. Each is a lower bound
on $L$ by a finite number that reads Stage-0 numbers only, and none bounds $L$ from above. We
write $L_0'=L_0'(N,\text{datum})$ for the least odd integer above all of them; every odd
$L\ge L_0'$ then meets them.
\item[(d)] We fix once and for all one Stage-0 datum as in (a). If the window of (b) is
non-empty, the number it constrains is taken in that window. The constant $L_0=L_0(N)\ge L_0'$
of hypothesis~(H2) of Definition~\ref{2.2:not-hypotheses} is the least odd integer above all
lower bounds on $L$ at this datum, those
of the block adjoined for clause~(v) in Chapter~\ref{ch:9} included; by the statement of their
shapes and dependences in Chapter~\ref{ch:9}, the latter are finitely many lower bounds on $L$
by finite functions of Stage-0 numbers. No value is assigned to $L_0(N)$.
\end{itemize}
\end{lemma}

\begin{proof}
We use the letters of Definition~\ref{2.3:not-stages-first}, except for the one exponent renamed
in the first item below. We go through its Stage~0 and
Stage~1 in the displayed order. For each condition we name what it constrains and which earlier
coordinates it reads. The words \emph{free upward}, \emph{window}, \emph{opener} and
\emph{determined} are those of Definition~\ref{2.3:def-order-table}.

\emph{Stage~0, part (a).} Let $r\ge2$ be given; $r$ is the exponent of the length $R(g)$, as
in \cite{B-CMP119}.
\begin{itemize}
\item We write $r_1$ here for the exponent of \cite{B-CMP122b} that is chosen next, denoted
$r_0$ in Definition~\ref{2.3:not-stages-first}, since in clause~(iv) of Theorem~0 the letter
$r_0$ denotes the depth offset of the Gaussian extraction. The exponent $r_1$ is constrained by
$r_1\ge r$, a floor that reads $r$.
\item The integer $p_0$ is free upward. It is constrained by finitely many floors, each by
earlier exponents, among them $p_0\ge5r$ \cite{B-CMP119} and $p_0\ge(d+5)r_1+8r+3$. These read
$r$ and $r_1$. We take $p_0$ above all of them.
\item The exponent $p_1$ lies in a window. Three kinds of condition constrain it:
$2p_1-(d+5)r_1>p_0$ \cite{B-CMP122b}; $p_1<p_0$ \cite{B-CMP122a}; and finitely many lower
bounds on $p_0-p_1$ by numbers at most
$4r+1$. All three read $p_0$, $r$ and $r_1$ only. The opener is the floor
$p_0\ge(d+5)r_1+8r+3$. Put $p_1:=p_0-4r-1$. Then
\begin{equation*}
2p_1-(d+5)r_1-p_0=p_0-8r-2-(d+5)r_1\ge1 ,
\end{equation*}
so the first condition holds. Moreover $p_0-p_1=4r+1>0$, so $p_1<p_0$, and $p_0-p_1$ exceeds
or equals each of the lower bounds by numbers at most $4r+1$. Raising $p_0$ further only widens
this window, so later floors on $p_0$ do not close it.
\item The exponent $q:=q_0=q_1$ is free upward. Its conditions are floors by earlier
exponents, for example $q\ge p_0+7r+1$, $q\ge3p_0+3r+1$ and $q\ge21p_0+1$, the last of which is
used by clause~(iv); each of these reads $p_0$ and $r$ only. The opener of the window of the one
further exponent is also a floor on $q$; that window lies between functions of
$(r,r_1,p_0,q)$. We take $q$ above all these floors, and that window is then non-empty. We take
the further exponent in it; it reads $(r,r_1,p_0,q)$ only.
\item The exponent $\kappa_0$ of \cite[(2.31)]{B-CMP119} is constrained by $\kappa_0\ge17$,
which reads nothing.
\item The exponent $v$, called $\nu$ in \cite[p.~363]{B-CMP122b}, is constrained by
$v\ge2p_0+dr_1$. This floor reads $p_0$ and $r_1$.
\item The fraction $\beta_{\mathrm{s}}=1/4$ \cite{B-CMP119} is determined. The fraction
$\beta_0$ of \cite[(2.6)]{B-CMP119}, free there in $(0,\frac12]$, is fixed at $\beta_0=1/7$.
The condition on $\beta_0$ is the ceiling $2c_2\gamma^2\le1-(1+\beta_0)^{-2}$
on $\gamma$, attached to $\gamma$ in Definition~\ref{2.3:not-stages-last}. Its right member is
positive since $\beta_0>0$. The H\"older exponent $\beta_H$ is any number in $(0,1)$. The
number $\alpha_6$ of \cite{B-CMP116} is determined by $d$.
\item The exponent $\nu_I$ ranges in $\{r+1,\dots,8r\}$. This set is finite and non-empty,
because $r+1\le8r$.
\item Several numbers range in fixed intervals between numbers, for example $[1/2,8/9]$,
$(8/7,3/2)$ and $(3/5,1)$. Each of these is non-empty.
\item One number carries only ceilings at Stage~0. Its admissible set there is an interval
$(0,c]$ with $c>0$, which is non-empty. The lower bound that would complete its window is a
floor on $L$ at Stage~1, treated below.
\end{itemize}
Each choice reads only coordinates chosen before it. None of these conditions involves
$\kappa'$. The one Stage-0 window of the displayed order not yet treated is the interval
between numbers whose opener is $\kappa'>5/12$; it serves only the block-law comparison
statement of Remark~\ref{3.2:rem-block-laws}, so it is not among the conditions considered, and
(a) requires nothing of it. This proves (a).

\emph{Part (b).} The opener of that window is the inequality $\kappa'>5/12$, and only this
inequality, not any value of $\kappa'$, is used. So the window is non-empty when the inequality
holds. Nothing here proves or evaluates $\kappa'>5/12$; that inequality is not a condition of
Theorem~0.

\emph{Stage~1, part (c).} The blocking factor $L$ is odd and free upward. The conditions
attached to it are the following.
\begin{itemize}
\item $L>11$ \cite{B-CMP109}, which reads nothing.
\item $(1+7\beta_{\mathrm{s}})L^{-2}\le1$ \cite{B-CMP109}, that is
$L\ge(1+7\beta_{\mathrm{s}})^{1/2}$, which reads $\beta_{\mathrm{s}}$.
\item The floor displayed below, which reads $\beta_{\mathrm{s}}$ and $\beta_0$; its right
member is finite because $1-\beta_{\mathrm{s}}/8=31/32>0$:
\begin{equation*}
L^2\ge\frac{2(1+\beta_{\mathrm{s}})(1+\beta_0)^2}{1-\beta_{\mathrm{s}}/8}.
\end{equation*}
\item Finitely many further floors by Stage-0 numbers. Among them is the floor whose right
member involves the number with admissible set $(0,c]$.
\item $L\ge4r_\star$, where $r_\star$ is a Stage-0 number. Under this floor the width
$\delta(L):=(1-2r_\star/L)/10$ lies in $[1/20,1/10)$.
\end{itemize}
Each of these is a lower bound on $L$ by a finite function of Stage-0 numbers, and none bounds
$L$ from above. Each condition is met by every odd integer above the lower bound on $L$ it
expresses (for the floor on $L^2$, above its square root). So every odd integer above the
largest of these finitely many bounds meets all of them. The least odd integer above them is
$L_0'(N,\text{datum})$, and every odd $L\ge L_0'$ meets every condition attached to $L$. This
proves (c).

\emph{Part (d).} The datum is fixed once. The lower bounds on $L$ of the block adjoined for
clause~(v) are stated, with their shapes and dependences, in Chapter~\ref{ch:9}; by that
statement they are finitely many, and each is a lower bound on $L$ by a finite function of
Stage-0 numbers. Adding them to the conditions of (c) replaces $L_0'$ by the least odd integer
$L_0(N)\ge L_0'$ above all of them, and every odd $L\ge L_0(N)$ meets every condition on $L$.
With $d=4$ the datum reads $N$ only, so $L_0(N)$ depends on $N$ alone.

The three depth floors of clause~(iv), $s_W(g)$, $s_{\mathrm{vol}}^\star(g)$ and
$s_T(g;\vartheta,\varrho)$, the last polylogarithmic in $g^{-2}$, whose maximum is the lower
bound that clause imposes on the depth $s(\mathsf{d})$, together with the constant $C_T(g)$ and
the threshold $\gamma_W(L,\mathfrak{p},\gamma,\vartheta;\varrho)\le\gamma_J$, are chosen after
$\gamma_H$, $\varrho$ and $\vartheta$: the floors read $g$ (and, for $s_T$, $\vartheta$ and
$\varrho$) only; the constant reads $g$; the threshold reads $L$, $\mathfrak{p}$, $\gamma$,
$\vartheta$ and $\varrho$. They impose no condition on the Stage-0 datum or on $L$.

\emph{Consistency with the prefix and the hypotheses.} The existence of $L_0(N)$ rests on the
fixed datum. Together, (a), (c) and (d) supply the first two existential quantifiers of the
prefix of Theorem~0 (Definition~\ref{2.1:not-prefix}), in this order: a Stage-0 datum, then
$L_0(N)$. This is in agreement with hypothesis~(H2), which requires that every condition on $L$
be a lower bound reading Stage-0 data only. It is also in agreement with the part of
hypothesis~(H3) stating that no condition bounds $L$ from above.
\end{proof}

\begin{lemma}[The order of choice has no cycle and is one-sided (no caps)]\label{2.4:lem-no-cycle}
Consider the conditions used by the clauses of Theorem~0 other than those of the block proper to
clause (v); of that block only the threshold $\gamma_U$, with the ceiling $\gamma_U\le\gamma_H$ of the
quantifier prefix, enters here, and its further conditions, its other letters and the one-sided
shapes of their conditions are stated in Chapter~\ref{ch:9}. Place them in the order of
choice of Notation~\ref{2.3:not-stages-first}, Notation~\ref{2.3:not-stages-middle} and
Notation~\ref{2.3:not-stages-last}, together with the thresholds chosen inside the clauses under
hypothesis (H7) of Notation~\ref{2.2:not-hypotheses}. In this order the letters come as follows.
\begin{itemize}
\item First the Stage-0 numbers, which read $d$, $N$ and one another, and then $L$. One Stage-0
window is opened by the inequality $\kappa'>5/12$ on the decay exponent $\kappa'$. This inequality
is not a condition on a letter of the order; it is carried as an antecedent of
Proposition~\ref{2.4:prop-order-consistent}.
\item Then Ba{\l}aban's auxiliary parameters $\mathfrak{p}$ of (H3) (Notation~\ref{2.1:not-prefix}),
in one fixed order. Each of them is determined by earlier letters, or bounded on one side by a
positive function of earlier letters, or, for finitely many of them, confined to a window that
earlier conditions make non-empty, for example $M_1<M_2=L^2M_1<M$.
\item Then $\gamma$, and then $\gamma_H$, subject to
\begin{equation}\label{2.4:eq-no-cycle-1}
\gamma_H\le(\gamma^{-2}+c_5+1)^{-1/2}.
\end{equation}
\item Then, inside the clauses, $\varrho$ and the thresholds chosen after it.
\end{itemize}
Then:
\begin{enumerate}
\item[(a)] among these letters the relation ``the condition on $a$ reads $b$'', $b\neq a$, has no
cycle, and no condition on a letter placed in Stages~0--9, on $\gamma_J$, on
$\gamma_W$ or on a letter of clause (iv) reads a letter of the block proper to
clause (v) other than $\gamma_U$;
\item[(b)] no condition bounds $L$, $M$, $M_1$, $\kappa$, $\kappa_1$, $q$, $p_0$, $K$, $m$, an age, a
depth or a number of renormalisation steps from above, and no condition bounds any coupling from
below; in particular the order contains no cap in the sense of Definition~\ref{2.3:def-cap}.
\end{enumerate}
\end{lemma}

\begin{proof}
By Definition~\ref{2.3:def-order-table}, a condition is attached to the latest letter it reads. We
verify, stage by stage, the two properties (P1) and (P2) below, and note for each condition the
side on which it bounds its letter; (b) is read off these notes at the end of the proof.
\begin{enumerate}
\item[(P1)] Each condition imposed in the order at the place of a letter $a$ reads, after its
determined letters are unfolded, the letter $a$ and otherwise only letters placed before $a$; thus
$a$ is the letter to which the condition is attached.
\item[(P2)] Each determined letter, and each constant taken from outside the order, reads only
letters placed before it.
\end{enumerate}
Given (P1) and (P2), suppose $y_1,\dots,y_n,y_{n+1}=y_1$ were a cycle, with the condition on $y_i$
reading $y_{i+1}\neq y_i$. Then $y_{i+1}$ would be placed strictly before $y_i$ for every $i$, so
$y_1$ would be placed strictly before itself, which is impossible. This proves the first part of
(a).

At Stage~0 (Notation~\ref{2.3:not-stages-first}) the letters are taken in the following order; we
record for each what its conditions read and on which side they bound it. That the windows of this
stage are non-empty is part of Lemma~\ref{2.4:lem-stage-zero-data}; it is not used here.
\begin{itemize}
\item $r\ge2$, and then the exponent $r_1\ge r$ of \cite{B-CMP122b}, which reads $r$ (the letter
is that of Lemma~\ref{2.4:lem-stage-zero-data}).
\item $p_0$, under floors reading $r$ and $r_1$, among them $p_0\ge5r$ and
$p_0\ge(d+5)r_1+8r+3$.
\item $p_1$, confined to a window. The conditions of this window are $2p_1-(d+5)r_1>p_0$,
$p_1<p_0$, and finitely many lower bounds on $p_0-p_1$ by numbers at most $4r+1$; they read $r$,
$r_1$ and $p_0$. All of them are attached to $p_1$, so they bound $p_1$ and not $p_0$.
\item $q$, under floors such as $q\ge p_0+7r+1$ and $q\ge3p_0+3r+1$.
\item $\kappa_0\ge17$, and $v\ge2p_0+dr_1$.
\item The fixed numbers $\beta_{\mathrm{s}}=1/4$, $\beta_0=1/7$, $\beta_H\in(0,1)$ and
$\alpha_6=\alpha_6(d)$.
\item Further windows. One exponent lies between functions of $(r,r_1,p_0,q)$, in a window opened
by a floor on $q$. The exponent $\nu_I$ ranges in $\{r+1,\dots,8r\}$, which reads $r$. Several
numbers range in fixed intervals between numbers and read no letter. One number carries only
ceilings at Stage~0; the lower bound that would complete its window is a floor on $L$ at Stage~1.
\item One interval between fixed numbers is non-empty provided the decay exponent $\kappa'$, given
by an estimate outside the order, satisfies $\kappa'>5/12$. Only this inequality enters, and no
condition reads a value of $\kappa'$.
\end{itemize}
No Stage-0 condition bounds $p_0$ or $q$ from above.

At Stage~1 every condition on $L$ is one of the floors listed at Stage~1 of
Notation~\ref{2.3:not-stages-first}, among them $L>11$ and the floor whose right member involves
the number carrying only ceilings at Stage~0. By Lemma~\ref{2.4:lem-stage-zero-data} these floors
are finitely many and finite, and they read Stage-0 numbers only. Every one of them is a lower
bound on $L$.

Stage~2 consists of determined letters that read Stage~0, $d$, $N$ and $L$. These are:
\begin{itemize}
\item $\delta(L)$;
\item $c_0=(11N^2/12\pi^2)\log L$, $c_2$, $c_5=2c_2$ and $(8c_2+4)^{-1/2}$;
\item $M_1^{\mathrm{thr}},\delta_0,a_0,B_0$ of \cite[Theorem~3.1]{B-CMP99b};
\item the rate $\delta_1$ of \cite{B-CMP116};
\item counting numbers.
\end{itemize}
No floor on $L$, $\kappa$, $\kappa_1$ or $M_1$ reads $c_2$ or $c_5$.

At Stages~3 and~4 (Notation~\ref{2.3:not-stages-middle}) every condition on $\kappa$ is a floor,
such as $\kappa\ge10\log L+1$ or $\kappa\ge C_\delta(d)/\delta$. Each of these reads only $d$,
Stage-0 numbers and $L$, the last either directly or through $\delta(L)$. Every condition on
$\kappa_1$ is a floor reading $d$, $L$, $\kappa$ and counting numbers, such as
$\tfrac12(\kappa_1-1)\ge2L\kappa$ and $\kappa_1\ge\max\{16\kappa+1,4L\kappa+1\}$.

At Stage~5 the letters and conditions are as follows.
\begin{itemize}
\item $M_1\in L^{\mathbb{N}}$ is bounded below by $M_1\ge M_1^{\mathrm{thr}}$ and by
$\delta_1M_1\ge \bar c(d)\,\kappa_1$, where $\bar c(d)\ge1$ depends on $d$ only.
\item Then the cube size $M_1^\sharp$ on which the two lowest levels \cite[(3.35), (3.37)]{B-CMP99b}
are applied, under the floor $M_1^\sharp\ge5LM_1$ ($3LM_1$ suffices on a sub-class of these cubes);
both are lower bounds reading $L$ and $M_1$ only.
\item Then come the determined letters $M_2:=L^2M_1$ and $B_1,B_3,B_5,a_1$ of \cite{B-CMP102b}.
\item Then $R_1$, chosen upward.
\item Then the kernel and random-walk constants, fixed before $M$ and reading at most $d$, $L$,
$M_1$, $M_1^\sharp$ and $\kappa_1$, all of which are placed earlier. Among them are the
constants $(B_W,\delta_W)$ and $\mu_0$ of Proposition~\ref{2.4:prop-m-freeness}. That proposition
asserts that the bounds hold with these constants at the
local backgrounds at which they are applied. It is used for no placement and for no side.
\item Then $\Theta_1$ and $\varpi_T$, under floors reading those constants.
\item Then $r'$ and $\ell$, chosen upward, together with quantities determined by $r'$.
\item Then $\hat c$, under a ceiling by earlier letters.
\end{itemize}

At Stage~6 every condition attached to $M=L^{m'}$ is a floor. The conditions are
\begin{gather*}
M\ge L^3M_1,\qquad M\ge M(\kappa),\qquad \delta_0M\ge\kappa,\\
\delta_1M\ge\kappa_1,\qquad e^{-\delta_0M/3}e^{16\kappa_1}<1,\qquad M\ge\max\{\varpi_T,3\},
\end{gather*}
together with the three floors
\begin{equation*}
\mu_0M\ge2\kappa_1,\qquad \delta_WM\ge2\kappa_1,\qquad
\delta_WM/4\ge32\kappa_1+2\log(LM)+C(d,L).
\end{equation*}
All of them read letters of Stages~0--5. Since $L>11$, we have $M_2=L^2M_1>M_1$ and
$M\ge LM_2>M_2$, so the window $M_1<M_2<M$ is non-empty.

No floor on $M$ is imposed at the rate $\tfrac12\min\{\delta_0,\kappa M^{-1}\}$ of
\cite[(5.10)]{B-CMP109}: at that rate
\begin{equation*}
\tfrac12\min\{\delta_0,\kappa M^{-1}\}\,M\le\kappa/2<\kappa_1 ,
\end{equation*}
so such a floor would be empty. The one inequality $M\le R_1M_1a_1/\varepsilon^{\mathrm{run}}$ of
\cite{B-CMP102b}, in which $\varepsilon^{\mathrm{run}}$ is the size of the small fields at the
running coupling, is not a condition of the order in this form. Writing
$\varepsilon^{\mathrm{run}}(x)$ for that size at inverse coupling $x$, the order imposes instead the
ceiling
\begin{equation}\label{2.4:eq-no-cycle-2}
\sup_{x\ge\gamma^{-2}}O(1)\,LM\,\varepsilon^{\mathrm{run}}(x)\le R_1M_1a_1 ,
\end{equation}
where $O(1)\ge1$. Since $O(1)L\ge1$, \eqref{2.4:eq-no-cycle-2} implies that inequality at every
level. Since $\varepsilon^{\mathrm{run}}(x)\to0$ as $x\to\infty$, \eqref{2.4:eq-no-cycle-2} is a
ceiling on $\gamma$, and it is placed at Stage~8. Every other inequality with $M$ on the smaller
side is attached to a letter chosen after $M$ and bounds that letter on its free side.

At Stage~7 (Notation~\ref{2.3:not-stages-last}) the letters are taken in the following order.
\begin{itemize}
\item $A_1>0$.
\item $A_0$, under the floor
\begin{equation*}
A_0/A_1\ge\max\{c(d)(1+\beta_0)B_3L^{d+2},\,c'(d,L)M^2\}.
\end{equation*}
\item $C_0$, under floors by multiples of $A_0$ and $A_1$. No condition on $C_0$ reads $C_1$.
\item $C_1$, under the floor $C_1>2B_3^2O(1)MC_0$, with $O(1)$ the unspecified constant of
\cite[p.~277]{B-CMP109}, not the constant of \eqref{2.4:eq-no-cycle-2}.
\item The radii $\alpha'_1,\alpha'_0,\alpha_1,\alpha_0,\alpha_2$, in this order, each under ceilings
by letters placed before it. Representative ceilings are
\begin{gather*}
(LM)^4\alpha_i\le1,\qquad O(1)M\alpha_1\le\beta_{\mathrm{s}},\\
B_3^2O(1)M\alpha_0<\tfrac12\alpha_1,\qquad (4B_3^2O(1)M)^2\alpha_0\le\beta_{\mathrm{s}},
\end{gather*}
together with $\alpha_2\le\alpha_1/8$ and $\alpha_2\le\alpha_0/C_{\mathrm{IV}}(d,L)$.
\item $E_0$, and then $\alpha_4$, under the ceiling $(LM)^4\alpha_4\le1$.
\item The determined constants. Among them is $C_\beta$, which reads $\delta_0$, $\kappa$ and $M$
and therefore comes after $M$. This fixes the triple $(c_0/4,c_5,C_\beta+1)$.
\item $\varepsilon^{\mathrm{I}}_0$ and then $\varepsilon^{\mathrm{I}}_1$, chosen downward under the
restrictions listed in \cite[Lemma~3]{B-CMP116}, which are stated there as restrictions on the
constants $M,\kappa,\kappa_1,\alpha_0,\alpha_1,\alpha_4,\alpha_6,\gamma_2,\gamma,\varepsilon_1$. Of
these, $M$, $\kappa$, $\kappa_1$, $\alpha_0$,
$\alpha_1$, $\alpha_4$ and $\alpha_6$ are placed earlier, and so is $\gamma_2:=\hat c\,(LM)^{-4}$,
which is determined by $\hat c$, $L$ and $M$; the restriction on $\varepsilon_1$ is the one imposed
here on $\varepsilon^{\mathrm{I}}_1$. A member of these restrictions that involves
$\gamma$ reads a letter placed after $\varepsilon^{\mathrm{I}}_1$; by
Definition~\ref{2.3:def-order-table} it is attached to $\gamma$, and it is among the conditions on
$\gamma$ of Stage~8.
\end{itemize}
The listed order $\kappa,\kappa_1,M_1,M_2,M,A_1,A_0,C_0,C_1$, the radii, $C_\beta$,
$\varepsilon^{\mathrm{I}}_0,\varepsilon^{\mathrm{I}}_1$ is the order displayed in (H3).

At Stage~8 every condition on $\gamma$ is a ceiling. It has either the form $\gamma\le c$ or the
form $\sup_{x\ge\gamma^{-2}}h(x)\le c$, where $c>0$ and $h(x)\to0$ as $x\to\infty$. Here $c$ and $h$
are built from letters of Stages~0--7. Examples are:
\begin{itemize}
\item $\gamma\le(8c_2+4)^{-1/2}$;
\item $(C_\beta+1)\gamma^2\le1$;
\item $2c_5\gamma^2\le1$;
\item $2c_2\gamma^2\le1-(1+\beta_0)^{-2}$;
\item $\sup_{x\ge\gamma^{-2}}C_ix^{-1/2}(\log x)^q\le\alpha_i$ for $i=0,1$;
\item \eqref{2.4:eq-no-cycle-2};
\item the requirement $C'\gamma^{\kappa_0-8}\le\min\{c_0/4,1\}$ on the bound of the large-field
part of the coupling increment. This is a ceiling because $\kappa_0\ge17$, and here $C'$ is a
constant of Stage~7. Write $\bar\varsigma(\gamma):=C'\gamma^{\kappa_0-8}$. Since
$\lambda_\sharp=c_0/2-\bar\varsigma(\gamma)$ and $B_\sharp=C_\beta+\bar\varsigma(\gamma)$, this
ceiling gives $c_0/4\le\lambda_\sharp$ and $B_\sharp\le C_\beta+1$.
\end{itemize}
Every ceiling involving $(\lambda_\sharp,c_5,B_\sharp)$ is imposed at the triple
$(c_0/4,c_5,C_\beta+1)$, which is fixed before $\gamma$.

The Stage-7 ceilings on $\alpha_0$ or $\alpha_1$ alone are imposed again at the running radii
$\alpha_{i,n}=g_nC_i(\log g_n^{-2})^q$ as ceilings on $\gamma$. Three comparisons are not imposed
in this way. The comparison $B_3^2O(1)M\alpha_{0,n}<\frac12\alpha_{1,n}$ is equivalent,
after division by $g_n(\log g_n^{-2})^q>0$, to the coupling-free condition
$C_1>2B_3^2O(1)MC_0$. The comparisons of $\alpha_2$ with $\alpha_1/8$ and with
$\alpha_0/C_{\mathrm{IV}}$ are imposed at the absolute radii only, since at running radii each of
them would bound the coupling from below.

Then $\gamma_H$ is chosen under \eqref{2.4:eq-no-cycle-1}, which reads $\gamma$ and $c_5$. The
ceilings on $\gamma$ read only letters of Stages~0--7, and those on $\gamma_H$ read only these
letters and $\gamma$.

Inside the clauses, $\varrho>0$ is fixed after $\gamma_H$. The thresholds are then chosen as
follows.
\begin{itemize}
\item $\gamma_J(L,\mathfrak{p},\gamma;\varrho)\in(0,\gamma_H]$ is chosen under finitely many
ceilings.
\item $\gamma_U\in(0,\gamma_H]$, the threshold of clause (v), is chosen after $\varrho$. The
condition $\gamma_U\le\gamma_H$ is a ceiling on $\gamma_U$ reading $\gamma_H$, which is placed
before it; by the quantifier prefix of Notation~\ref{2.1:not-prefix}, every condition on
$\gamma_U$ is an upper bound on it. The further ceilings on $\gamma_U$ belong to the block proper
to clause (v) and are stated, with the letters they read, in Chapter~\ref{ch:9}.
\item Inside clause (iv), $\vartheta\in(0,1]$ is taken first, and then the ceiling
$\gamma_W(L,\mathfrak{p},\gamma,\vartheta;\varrho)\le\gamma_J$ is imposed.
\end{itemize}
Also inside clause (iv) there are the constants $k_W(L)$, $c_{\mathcal{K}}$, $C_{\mathcal{K}}$, and,
after $g$, the three depth floors $s_W(g)$, $s_{\mathrm{vol}}^{\star}(g)$ and
$s_T(g;\vartheta,\varrho)$ of that clause, the second being the value at the reference torus
exponent $m_{\mathbb{T}}(g_-(g))$ of the volume-dependent depth floor $s_{\mathrm{vol}}$ of that
clause, $s_{\mathrm{vol}}^{\star}(g):=s_{\mathrm{vol}}(g,m_{\mathbb{T}}(g_-(g)))$, together with
the constant $C_T(g)$. Here $s_T$ is polylogarithmic in $g^{-2}$. None of these letters reads $K$,
$g_0$, $m$, $\mathsf{d}$ or a position. Positively, $k_W$ reads $L$; $s_W$,
$s_{\mathrm{vol}}^{\star}$ and $C_T$ are functions of $g$, and $s_T$ of $g$, $\vartheta$ and
$\varrho$, besides letters fixed before $g$; and $c_{\mathcal{K}}$, $C_{\mathcal{K}}$ are the lower
and upper constants of the explicit kernel of clause (iv). Thus each of $k_W$, the three depth
floors and $C_T$ reads only letters placed before it. Clause (iv) imposes on the separation
$\mathsf{d}$ of the test pair the two conditions
\begin{gather*}
\max\{s_W(g),s_{\mathrm{vol}}^{\star}(g),s_T(g)\}\le s(\mathsf{d}),\\
K-s(\mathsf{d})\ge k_W .
\end{gather*}
Like the separation window in the definition of the class $\mathfrak{T}(\mathsf{d};\vartheta)$,
they are conditions on the test pair, imposed
after $g$, $g_0$, $m$ and $K$; they bound no letter of the order.
The conditions proper to clause (v) form an adjoined block with its own letters. These letters are
chosen after $\gamma_H$, or at their own places on their free side, and are stated in
Chapter~\ref{ch:9}. Apart from $\gamma_U$, no letter of Stages~0--9 listed here, no threshold
listed above and no letter of clause (iv) belongs to that block, and the conditions on these
letters, as listed, read no letter of that block other than $\gamma_U$. Hence no condition on a
letter placed in Stages~0--9, on $\gamma_J$, on $\gamma_W$ or on a letter of
clause (iv) reads a letter of the block proper to clause (v) other than $\gamma_U$.

At Stage~9 the letters are chosen in the order $g,m,g_0$, under the following conditions.
\begin{itemize}
\item $g\le\gamma_H$, and inside the clauses also $g\le\gamma_J$, $g\le\gamma_W$ or
$g\le\gamma_U$. These are ceilings.
\item $m\ge m_{\mathbb{T}}(g_-(g))$, a floor on $m$. It reads $g$, which is chosen earlier.
\item $g_0\le g_-(g)$, a ceiling on $g_0$.
\end{itemize}
The number $K=K(g_0,g)$ is determined by first passage, and no condition is attached to it.

The level-indexed quantities are determined at each level from $x_0,\dots,x_k$. They impose no
condition on any letter: the only conditions they require are the floors on $x$ of Stage~8, and no
letter or condition of Stages~0--8 reads them. This establishes (P1) and (P2), and hence (a).

It remains to read (b) off the notes on sides. Every condition attached to $L$, $p_0$, $q$,
$\kappa$, $\kappa_1$, $M_1$ or $M$ is a floor. A condition that reads one of these letters but is
attached to a later letter bounds that later letter, on its free side. The letter $m$ is bounded
only from below. No condition is attached to $K$.

Ages and numbers of steps are indices of the construction, not letters of the order. No condition
of the order reads them. The depth $s(\mathsf{d})$ is read only by the two conditions of clause (iv)
on the separation $\mathsf{d}$ displayed above, which are conditions on the test pair and bound no
letter of the order.

Every condition on $\gamma$, $\gamma_H$, $\gamma_J$, $\gamma_W$, $\gamma_U$, $g$ or
$g_0$ is a ceiling. Every condition on a running inverse coupling is a floor on it, which acts as
a ceiling on $\gamma$. Hence no coupling is bounded from below.
\end{proof}

\begin{proposition}[$M$-freeness of the random-walk constants. Stated; proof incomplete at the step marked $(\ast)$ in Chapter~\ref{ch:5}]\label{2.4:prop-m-freeness}
Consider backgrounds given only as local data on the cells of the one family of cells, at the
decomposition of the density at which the following bounds are applied. Let $M$ range over its
values above its floors. For such backgrounds, the kernel $G_W$ of the conditional
fluctuation operator satisfies
\begin{equation}\label{2.4:eq-m-freeness-1}
|G_W(x,y)|\le B_W\,e^{-\delta_W|x-y|}.
\end{equation}
The companion kernel of the random-walk expansion of that operator satisfies its exponential bound
with decay rate $\mu_0$. Here $(B_W,\delta_W)$ and $\mu_0$ are the constants fixed at Stage~5 of
the order of choice (Notation~\ref{2.3:not-stages-middle}), that is, the $M$-free constants of
the same two bounds for backgrounds in the global regularity class.
\end{proposition}

The proposition is stated again in Chapter~\ref{ch:5}, next to Ba{\l}aban's bounds for old terms
beside a new region, together with what is proved toward it and with the step at which its proof
is incomplete marked.

The proposition is needed for the following reason. The floors of $\Theta_1$ and $\varpi_T$, the
width parameters of the order of choice (Notation~\ref{2.3:not-stages-middle}), and the three
floors on $M$ of Stage~6 are written in the constants of Stage~5. The bounds they serve, however,
are applied at local data. The words \emph{floor} and \emph{ceiling} are used here in the sense of
Notation~\ref{2.3:not-stages-first}.

Suppose the constants of the bounds at local data depended on $M$. Then the width parameters
$\Theta_1$ and $\varpi_T$ could be moved after $M$. The floor $M\ge\varpi_T$ would then be replaced
by a ceiling of sup form on $\gamma$. The three floors on $M$, however, would then compare $M$ with
an unspecified function of $M$, and it would not be determined on which side they bound $M$. Write
$\delta_W(M)$ for the decay rate in \eqref{2.4:eq-m-freeness-1} when it is allowed to depend on
$M$. If $\delta_W(M)\,M$ stayed bounded, these three conditions would even be empty. For them
there is no alternative placement.

\begin{proposition}[The conditions of Theorem~0 can be met]\label{2.4:prop-order-consistent}
Let $N\ge2$. In this proposition the conditions considered are the conditions used by the clauses
of Theorem~0, and no others. They are ordered as in the order of choice of
Notation~\ref{2.3:not-stages-first}, Notation~\ref{2.3:not-stages-middle} and
Notation~\ref{2.3:not-stages-last}. The condition of Notation~\ref{2.3:not-stages-first} confining
one Stage-0 number to an interval whose non-emptiness rests on an inequality for the exponent
$\kappa'$ is not among them, and no condition considered here reads that number or $\kappa'$. The
conditions proper to clause~(v) form a separate block with its own letters; statement (A1) below
does not concern them, and they are adjoined in part~(c) under the hypothesis stated there.
The thresholds chosen last and the clause-internal quantifiers
are those of Remark~\ref{2.3:rem-thresholds}. Assume the following two statements.
\begin{itemize}
\item[(A1)] Consistency of the displayed order with hypothesis (H3) of
Notation~\ref{2.2:not-hypotheses} (stated; proof incomplete).
\begin{itemize}
\item At each free or window letter of Stages~2--7, given any admissible choice of the earlier
letters, the conditions attached to it are as follows:
\begin{itemize}
\item at a letter free upward, finitely many floors, each with a finite positive right member in
the earlier letters;
\item at a letter free downward, finitely many ceilings, each with a finite positive right member
in the earlier letters;
\item at a window letter, the two ends of a window that conditions on earlier letters make
non-empty (an opener in the sense of Definition~\ref{2.3:def-order-table}).
\end{itemize}
Every determined letter is read off from earlier ones.
\item At $\gamma$, the conditions are finitely many, and each is a ceiling $\gamma\le c$ with
$c>0$, or has the form
\begin{equation}\label{2.4:eq-order-consistent-1}
\sup_{x\ge\gamma^{-2}}h(x)\le c ,
\end{equation}
where $c>0$ and $h$, with $h(x)\to0$ as $x\to\infty$, are built from letters of Stages~0--7.
\item At $\gamma_H$ the conditions are finitely many ceilings with positive right members. The
ceiling
\begin{equation}\label{2.4:eq-order-consistent-2}
\gamma_H\le(\gamma^{-2}+c_5+1)^{-1/2}
\end{equation}
is among them.
\item After $\gamma_H$ the earlier choices and $\varrho>0$ are held fixed. At these fixed values
the conditions on $\gamma_J$ are finitely many and read only $L$, $\mathfrak{p}$, $\gamma$ and
$\varrho$, besides the ceiling $\gamma_J\le\gamma_H$. Inside clause~(iv), after $\gamma_J$, also
$\vartheta\in(0,1]$ is held fixed; at these fixed values the conditions on $\gamma_W$ are finitely
many and read only $L$, $\mathfrak{p}$, $\gamma$, $\vartheta$ and $\varrho$, besides the ceiling
$\gamma_W\le\gamma_J$. Each condition on $\gamma_J$ or on $\gamma_W$ has a left member that tends
to its admissible side as the coupling tends to $0$.
\item Inside clause~(iv), after $\vartheta$, the further letters of that clause are: three floors
on the depth $s(\mathsf{d})$, among them $s_W(g)$, each a function of $g$ alone at fixed
$\vartheta$ and $\varrho$, and one of them polylogarithmic in $g^{-2}$; the integer $k_W$,
depending on $L$; the constants $c_{\mathcal{K}}$ and $C_{\mathcal{K}}$; and the constant of the
bound comparing the connected two-point functions on any two tori of the family whose exponents
are at least $m_{\mathbb{T}}(g_-(g))$, a function of $g$. Each of
them is determined or bounded on one side only, and none of them reads $K$, $g_0$, $m$,
$\mathsf{d}$ or the position of the supports.
\end{itemize}
The proof of (A1) is incomplete at one point: the enumeration of the order of choice gives these
shapes at every letter of Stages~2--8 and at the thresholds $\gamma_J$ and $\gamma_W$, except at
the following constants:
\begin{itemize}
\item the constants $(B_W,\delta_W)$ of a bound on the kernel of the conditional fluctuation
operator;
\item the decay rate $\mu_0$ of a second kernel estimate;
\item the constants fixed with them.
\end{itemize}
These constants are fixed after $M_1$ and before $M$, by formulas in $L$ and $M_1$ only. The
condition imposed at these constants, on the cells of the one family of cells over which the
effective densities are decomposed, is~(A2). The shapes of the letters proper to
clause~(iv), listed in the last item above, are those displayed in the statement of the two-point
witness theorem in Chapter~\ref{ch:8}.
\item[(A2)] $M$-independence of the random-walk constants (stated; proof incomplete). The
two kernel estimates hold at local data on the one family of cells, with the $M$-free constants
$(B_W,\delta_W)$ and $\mu_0$, for every cube size $M$. This is
Proposition~\ref{2.4:prop-m-freeness}. Its proof is incomplete: the $M$-independence is proved for
the kernel estimates on the whole lattice; what is missing is the local version, on the one
family of cells, which is the version used.
\end{itemize}
Then there are admissible Stage-0 data and a constant $L_0=L_0(N)$ such that, for every odd
$L\ge L_0$, the following choices can be made.
\begin{itemize}
\item[(a)] Ba{\l}aban's auxiliary parameters $\mathfrak{p}$ can be chosen in the one fixed order of
(H3). Each of them is determined by earlier ones, or bounded on one side by a positive function
of earlier ones, or, for finitely many of them, confined to a window that earlier conditions
make non-empty, for instance
\begin{equation*}
M_1<M_2=L^2M_1<M .
\end{equation*}
\item[(b)] Then there is $\gamma>0$ satisfying every condition attached to $\gamma$, and then there
is $\gamma_H>0$ satisfying every condition attached to $\gamma_H$, among them
\eqref{2.4:eq-order-consistent-2}.
\item[(c)] After $\gamma_H$, for every $\varrho>0$, there is a threshold $\gamma_J\in(0,\gamma_H]$,
depending on $(L,\mathfrak{p},\gamma;\varrho)$ only.
Inside clause~(iv), for every $\vartheta\in(0,1]$, there is a threshold
$\gamma_W\in(0,\gamma_J]$, depending on $(L,\mathfrak{p},\gamma,\vartheta;\varrho)$ only; the
further letters of that clause can then be chosen, each with a non-empty admissible set, the
condition on the depth being that $s(\mathsf{d})$ is at least the largest of the three depth
floors; and these letters add no choice to (H3). Inside clause~(v), assume that each condition
proper to that clause has the shape and dependence recorded in the statement on the conditions
proper to clause~(v), stated in Chapter~\ref{ch:9}: it is attached to the
latest letter it reads, it reads only letters chosen before it, and it is a floor on an upward
letter, a ceiling on a downward letter, or one end of a window on a letter proper to clause~(v)
made non-empty by numbers fixed earlier; and assume that the conditions of that clause attached
to $\gamma_U$ are, at fixed $\varrho$, finitely many ceilings with positive right
members. Then, for every $\varrho>0$, there is a threshold
$\gamma_U\in(0,\gamma_H]$.
\end{itemize}
In particular, the quantifier prefix of Theorem~0 (Notation~\ref{2.1:not-prefix}) is non-vacuous.
\end{proposition}

\begin{proof}
The proof runs through the order of choice, one group of letters at a time.

\emph{The Stage-0 data and $L$.} We use Lemma~\ref{2.4:lem-stage-zero-data}. It gives admissible
Stage-0 data for every $r\ge2$. It also shows that
the lower bounds on $L$ are finitely many, are
finite, and read Stage-0 numbers only. We fix one $r\ge2$ and one admissible datum once and for
all. The datum depends on $N$ only, since $d=4$ is fixed. Let $L_0(N)$ be the least odd integer
above all floors on $L$ of the conditions considered. This includes the floors on $L$ among the
conditions proper to clause~(v), which are lower bounds on $L$ by Lemma~\ref{2.4:lem-no-cycle}; by
Lemma~\ref{2.4:lem-stage-zero-data} they too are finitely many, finite and read Stage-0 numbers
only. Then $L_0(N)$ is
the constant of (H2), and every odd $L\ge L_0(N)$ satisfies all conditions on $L$. No value is
assigned to $L_0(N)$.

\emph{Admissible shape.} We use Lemma~\ref{2.4:lem-no-cycle}. It shows that every condition reads,
besides the letter to which it is attached, only earlier letters. It also shows that no condition
is a cap in the sense of Definition~\ref{2.3:def-cap}. Now consider the free letters of
Stages~2--7. By (A1), every condition attached to such a letter is a floor if the letter is free
upward and a ceiling if it is free downward, with a finite positive right member in earlier
letters. By (A1), every window letter has a window that conditions on earlier letters make
non-empty, hence a non-empty admissible set. Hence the part of the table of the conditions
considered formed by Stages~2--7, without the block of conditions proper to clause~(v), has
admissible shape in the sense of Definition~\ref{2.3:def-order-table}.

At the constants $(B_W,\delta_W)$ and $\mu_0$, placed after $M_1$ and before $M$, the step is the
following. These constants are given by formulas in $L$ and $M_1$ only. Hence the floors on $M$
that read $\mu_0$ and $\delta_W$ are floors by their form. The condition at these constants is
attached to no letter and bounds none. That it is met, namely that the two kernel estimates hold
with these same constants at local data on the one family of cells, for every $M$ chosen above its
floors, is (A2), that is, Proposition~\ref{2.4:prop-m-freeness} (stated; proof incomplete). This
is the one point at which (A2) enters.

\emph{The auxiliary parameters.} By Lemma~\ref{2.3:lem-choice-shapes}, applied to that part of
the table with $L$ fixed, its letters can be chosen one after another. For every odd
$L\ge L_0(N)$ this gives $\mathfrak{p}$ in the order of (H3), which is~(a).

\emph{The thresholds $\gamma$ and $\gamma_H$.} A ceiling $\gamma\le c$ with $c>0$ holds on
$(0,c]$. For a condition of the form \eqref{2.4:eq-order-consistent-1}, the function
$Y\mapsto\sup_{x\ge Y}h(x)$ is non-increasing, since the supremum is taken over a smaller set as
$Y$ grows, and it tends to $0$ as $Y\to\infty$, since $h(x)\to0$. So there is $Y_0>0$ with
$\sup_{x\ge Y}h(x)\le c$ for every $Y\ge Y_0$, and with $Y=\gamma^{-2}$ the condition holds for
every $\gamma\in(0,Y_0^{-1/2}]$. By (A1) the conditions at $\gamma$ are finitely many, so the
intervals on which they hold meet in $(0,c']$, where $c'>0$ is the least of their right ends,
and any $\gamma$ in this intersection is admissible. By (A1), $\gamma_H$ is then free downward
under finitely many ceilings with positive right members, among them
\eqref{2.4:eq-order-consistent-2}; each holds on an interval $(0,c]$ with $c>0$, and these
intervals meet in the same way. This gives~(b).

\emph{The thresholds $\gamma_J$ and $\gamma_W$.} Fix $\varrho>0$. By (A1), each condition on
$\gamma_J$ has a left member that tends to its admissible side as the coupling tends to $0$, so it
holds on some interval $(0,c]$ with $c>0$; as in the case of $\gamma$, these finitely many
intervals meet $(0,\gamma_H]$ in a non-empty interval, and $\gamma_J$ is chosen in it. By (A1) the
conditions on $\gamma_J$ read only $L$, $\mathfrak{p}$, $\gamma$, $\varrho$ and, through the ceiling
$\gamma_J\le\gamma_H$, the threshold $\gamma_H$, which is chosen from $L$, $\mathfrak{p}$ and
$\gamma$; so $\gamma_J$ depends on $(L,\mathfrak{p},\gamma;\varrho)$ only. Inside clause~(iv),
fix $\vartheta\in(0,1]$ after $\gamma_J$, as in Remark~\ref{2.3:rem-thresholds}. Likewise, with
$\gamma_J$ in place of $\gamma_H$ and $\vartheta$ among the letters read, there is
$\gamma_W\in(0,\gamma_J]$, and it depends on $(L,\mathfrak{p},\gamma,\vartheta;\varrho)$ only. The
further letters of clause~(iv) are then chosen
inside that clause, after $\vartheta$. By (A1) each of them is determined or bounded on one side
only, so its admissible set is non-empty. In particular each of the three depth floors bounds
$s(\mathsf{d})$ from below, so the admissible depths, those with $s(\mathsf{d})$ at least the
largest of the three floors, form a half-line. These letters are chosen inside clause~(iv), after
$\vartheta$ and after all choices of (H1)--(H7), so they add no choice to (H3); by (A1) none of
them reads $K$, $g_0$, $m$, $\mathsf{d}$ or the position of the supports.

\emph{The threshold $\gamma_U$.} Assume the hypothesis of the clause-(v) part of (c), and fix
$\varrho>0$. Where a condition proper to clause~(v) bears on a letter of the table, it bounds that
letter on its free side, since, by that hypothesis, it is a floor on an upward letter or a ceiling
on a downward one. The
choice at that letter is then made above the larger floor or below the smaller ceiling. The
enlargement part of Lemma~\ref{2.3:lem-choice-shapes} therefore applies to the table enlarged by
this block. Its hypothesis, that every condition of the original table reads only letters of that
table, is part of Lemma~\ref{2.4:lem-no-cycle}. By the same hypothesis of (c), at the fixed
$\varrho$ the conditions of the block attached to $\gamma_U$ are finitely many ceilings with
positive right members. Each holds on an interval $(0,c]$ with $c>0$, and, as in the case of
$\gamma$, these intervals meet $(0,\gamma_H]$ in a non-empty interval. This gives
$\gamma_U$ and completes~(c).

\emph{Non-vacuity of the prefix.} With these choices, by the one-sided form of (H4)--(H6), the
remaining quantifiers of the prefix of Notation~\ref{2.1:not-prefix} range over non-empty sets:
$g$ over $(0,\gamma_H]$ (over $(0,\gamma_J]$ or $(0,\gamma_W]$ inside the clauses that use these
thresholds), $m$ over the half-line of integers $m\ge m_{\mathbb{T}}(g_-(g))$, and $g_0$ over
$(0,g_-(g)]$; and $K=K(g_0,g)$ is determined by first passage as in
Definition~\ref{1.2:def-flow-constants}. Inside clause~(v), $g$ ranges over $(0,\gamma_U]$, and in
place of first passage the bare coupling is tuned, with the torus exponent $m_0\ge0$ and the depth
$n$ quantified as in Notation~\ref{2.1:not-prefix}.
\end{proof}

%% ---- end inlined file: chapters/c02 ----

\clearpage
\setcounter{section}{2}
\refstepcounter{section}
\section*{Chapter~\thesection\quad Statement of Theorem~0\addcontentsline{toc}{section}{\protect\numberline{\thesection}Statement of Theorem~0}\label{ch:3}}
%% ---- begin inlined file: chapters/c03 ----
\begin{definition}[Non-trivial]\label{3.1:def-non-trivial}
Throughout this book the word \emph{non-trivial}, said of a continuum limit, has the following
precise meaning.

The observable is the Yang--Mills action density $\operatorname{tr} F\cdot F$ smeared with a
test function, in the normalisation in which the coupling constant multiplies the action and not
the field. In this normalisation the connected two-point function of two such observables,
supported in balls of radius $\mathsf{d}$ at mutual distance of order $\mathsf{d}$ and normalised
by the product $\|f_1\|_\infty\|f_2\|_\infty$ of the sup norms of the two test functions $f_1$,
$f_2$, is the
fourth power of an effective coupling at scale $\mathsf{d}$ times a kernel invariant under
dilations.

A continuum limit is called \emph{non-trivial} if this normalised two-point quantity is not
identically zero and tends to zero, as $\mathsf{d}\to 0$, more slowly than every power of
$\mathsf{d}$.

This excludes the three degenerate outcomes of a limit taken at vanishing bare coupling:
\begin{itemize}
\item[(a)] that all connected correlations of the action density vanish;
\item[(b)] that the limit is a free Maxwell field at some charge, for which the quantity is a
non-zero constant;
\item[(c)] that the limit is scale invariant.
\end{itemize}

Non-triviality in this sense is a statement about a two-point function. It does not assert that
the limit theory is, in the sense of axiomatic field theory, anything other than a theory of
generalised free type, that is, a theory all of whose correlation functions are determined by its
two-point function through Wick's rule; no statement about a two-point function can assert this,
since every admissible two-point function is also the two-point function of a Gaussian field
with prescribed covariance. Nor would it be the right criterion that the law of the smeared
action density fails to be Gaussian, since the square of a free field already has a law that
fails to be Gaussian. Nothing
is asserted about correlations at three or more points, about Wilson loops, about rotation
invariance beyond the hypercubic group, or about a mass gap.
\end{definition}

\begin{theoremzero}[SU(2) lattice Yang--Mills theory in four dimensions along Ba{\l}aban's
renormalization group, for the curvature observables; $N\ge2$ in clauses (0)--(i)]
\label{3.1:thm-main}
Let $N\ge2$ and let $L_0=L_0(N)$ be as in hypothesis (H2). Let $L\ge L_0$ be odd, and let
Ba{\l}aban's auxiliary parameters $\mathfrak{p}$ and the thresholds $\gamma$, $\gamma_H$ be chosen
after $L$ in the order of hypothesis (H3), as in Notation~\ref{2.1:not-prefix} and
Remark~\ref{2.3:rem-thresholds}, the members of $\mathfrak{p}$ being subject also to the
one-sided lower bounds displayed in clauses (iv) and (v-e) and to one further lower bound on $M$,
mentioned after the theorem. After $\gamma_H$, the source radius $\varrho>0$ and the
threshold $\gamma_J$ are chosen as in Remark~\ref{2.3:rem-thresholds}; clause (v) quantifies
$\varrho$ afresh. Assume (H1)--(H6) of Notation~\ref{2.2:not-hypotheses}, and
(H7) where a clause says so. Unless a clause says otherwise, $g\in(0,\gamma_H]$,
$m\ge m_{\mathbb{T}}(g_-(g))$, $g_0\in(0,g_-(g)]$, and $K=K(g_0,g)$ is the first-passage cutoff of
Definition~\ref{1.2:def-flow-constants}, so that the bare spacing is $L^{-K}$. Clauses (0) and
(i) hold for every $N\ge2$; clauses (iii), (ii), (iv), (v) and (vi) are stated for $N=2$. The
constants of (0) and (i) depend only on $(L,\mathfrak{p},\gamma,N)$, and on the level $k$ only
through $g_k$; they depend on none of $K$, $m$, $g_0$ and the position, with three exceptions,
the quantities displayed as such with the proof of clause (i) in Chapter~\ref{ch:4}, none of
which enters any of clauses (ii)--(vi).

\begin{itemize}
\item[(0)] \emph{Coupling flow.} $K$ is a well-defined non-negative integer, and
\begin{equation*}
\frac{g_0^{-2}-g^{-2}}{B_\sharp}-1\;\le\;K\;\le\;1+\frac{g_0^{-2}-g^{-2}-B_\sharp+c_5}{\lambda_\sharp},
\end{equation*}
so that $K\to\infty$ as $g_0\downarrow0$. Moreover $g_k\in(0,g)\subset(0,\gamma_H]$ for $k\le K$,
$g_K\in[g_-(g),g)$, and
\begin{equation*}
\lambda_\sharp(j-i)-c_5\;\le\;x_i-x_j\;\le\;B_\sharp(j-i),\qquad 0\le i\le j\le K.
\end{equation*}
The $K$ renormalization steps on $\mathbb{T}_m$ produce exactly these couplings.
\emph{Fixed-depth windows.} Let $g_0^{(i)}\downarrow0$ and $K_i:=K(g_0^{(i)},g)$. For every
$s\ge0$ and every $i$ with $K_i\ge s$,
\begin{equation*}
x_{K_i-s}\in\bigl(g^{-2}+\max\{0,\lambda_\sharp s-c_5\},\;g^{-2}+B_\sharp(s+1)\bigr];
\end{equation*}
along a subsequence, for every $s$, $g_{K_i-s}\to g_{\infty,s}$, with
$g_{\infty,s}^{-2}\in[g^{-2}+\lambda_\sharp s-c_5,\;g^{-2}+B_\sharp(s+1)]$.
Not asserted: a sign of $x_k-x_{k+1}$ at a single step; the one-loop value of the individual
increments (for the cumulative law in the regime $g\le\gamma_J$ see (iii-d)); monotonicity of a
tuned bare coupling;
convergence of $g_{K-s}$ along first passage without passage to subsequences; regularity of
$\beta_j$ in the couplings.

\item[(i)] \emph{Ultraviolet stability.} Clause (i) is \cite[Theorem~1]{B-CMP122b}, its
hypothesis supplied by clause (0), used in the following form. No statement about the laws of the
block variables is a clause of
the present theorem (compare Remark~\ref{3.2:rem-block-laws}, which contains the block-law
statements; their proofs are outstanding and no clause below uses them). For every
$k\le K$:
\begin{itemize}
\item[(i-a)] $\rho_k$ has the representation \cite[(2.18)]{B-CMP119}: a sum over admissible
histories of large-field operations applied to characteristic functions, times the exponential of
the sum of $-A(g_k^{-2}(\cdot),U_k)$ and localised terms of every level $j\le k$; each localised
term is gauge
invariant, analytic on $U^c_j(X,\alpha_0,\alpha_1)$ read hereditarily
(Definition~\ref{1.5:def-domains}), and bounded by a constant times $e^{-\kappa d_j(X)}$. A boundary
term whose domain contains a large-field region $Z$ of the second class obeys the bound of
\cite[(1.99)]{B-CMP122b}
in the relative length $d_{j,Z}(X)$ when it is created beside no older such region, and otherwise
a bound decaying in the tree length $d^\theta(X)$, where $d^\theta(X)\le d_{j,Z}(X)$; no lower
bound of $d^\theta(X)$ by $d_{j,Z}(X)$ is asserted.
\item[(i-b)] On real fields,
\begin{equation*}
\chi_k\,e^{-g_k^{-2}A(U_k)-E_-N_k}\,dV_k\;\le\;\rho_k\,dV_k,
\qquad \int\rho_k\,dV_k\;\le\;e^{E_+N_k},
\end{equation*}
with $E_-=E_-(L,\mathfrak{p},N)$ and $E_+\le C(L,\mathfrak{p},\gamma,N)(1+\log g_k^{-2})$;
\cite[(0.1)]{B-CMP122b} states a pointwise upper bound, and the present clause uses the upper
bound in the integrated form displayed.
\item[(i-c)] $\int\rho_k\,dV_k=\int\rho_0\,dU$.
\item[(i-d)] \emph{Exact insertion.} For $0\le s\le K$ and every bounded gauge-invariant
continuous function $F$ of the block variables at level $K-s$,
\begin{gather*}
\langle F\rangle_K=\frac{\sum_h\int F^{(h)}\,d\tau_h}{\sum_h\int d\tau_h},\\
F^{(h)}=\mathbb{E}\bigl[F(\mathcal{M}^{K-s}U)\,\big|\,h,V_{K-s}\bigr],\qquad
|F^{(h)}|\le\sup|F|,
\end{gather*}
with $\mathcal{M}$ the one-step block-averaging map of Definition~\ref{1.5:def-block-average},
so that $V_{K-s}=\mathcal{M}^{K-s}U$.
For complex curvature sources $w=\sum_iz_if_i$ with $(z_i)$ in the source disc of the correlation
theorem (Chapter~\ref{ch:5}), the sourced densities $\rho^w_k$ keep the same inductive format,
with the running shift of (iii-c), and $\int\rho^w_k\,dV_k=\int\rho^w_0\,dU$. This is an
identity; it asserts nothing about the values of the laws of the block variables.
\end{itemize}

\item[(iii)] \emph{Cutoff-uniform correlation bounds} ($N=2$; (H7)).
\begin{itemize}
\item[(iii-a)] \emph{The torus bound.} For $n\ge2$ and real Lipschitz $f_1,\dots,f_n$ whose
supports are pairwise at distance at least $\mathsf{d}>0$,
\begin{gather*}
|S^T_{n;K,m}(\vec f)|\le C_n(\mathsf{d})\prod_{i=1}^n\|f_i\|_\sharp,\\
C_n(\mathsf{d})=\bigl[C_{UV}(\mathsf{d})+\log2\cdot(2\mathsf{Q}+2)^n\bigr](n/\varrho)^n,
\end{gather*}
where $\mathsf{Q}=e^{(E_+^\sharp+C_{\mathrm{low}})N_T^4}$ is the volume factor and
$C_{UV}(\mathsf{d})$, $E_+^\sharp$, $C_{\mathrm{low}}$, $N_T$ are fixed with the proof in
Chapter~\ref{ch:5}; $C_n(\mathsf{d})$ depends on $(L,\mathfrak{p},g,m,\varrho,n,\mathsf{d})$ only,
not on $K$, $g_0$ or the position. By itself this bound gives nothing in infinite volume.
\item[(iii-b)] \emph{The per-ball volume-free bound.} For $n\ge2$, $R>0$ and $\mathsf{d}>0$ there
is $C^\infty_n(\mathsf{d};R)$, independent of $K$, $g_0$, $m$ and the position and non-increasing
in $\mathsf{d}$ (no decay in $\mathsf{d}$ is asserted), such that
\begin{equation*}
|S^T_{n;K,m}(\vec f)|\le C^\infty_n(\mathsf{d};R)\prod_i\|f_i\|_\sharp
\end{equation*}
for every $m$ with
$L^m\ge\ell(n,R,g)$, the least power of $L$ that is at least
$\max\{8n(R+1),L^{m_{\mathbb{T}}(g_-(g))}\}$, every $g_0\in(0,g_-(g)]$, and all real Lipschitz
$f_i$, each supported in a ball of radius $R$, with supports pairwise at distance at least
$\mathsf{d}$.
\item[(iii-c)] \emph{The running shift.} Let $\zeta$ be a constant shift of $g_0^{-2}$, all
auxiliary choices being frozen at $\zeta=0$, and let $b_{k+1}(\zeta)$ be the coupling increment of
the $(k+1)$-st step of the shifted run. Put $\zeta_0=\zeta$ and
$\zeta_{k+1}=\zeta_k+b_{k+1}(\zeta)-b_{k+1}(0)$. Then $\zeta_k$ is analytic on $|\zeta|<2\varrho$
and real on reals, $|\zeta_k-\zeta|\le|\zeta|/3$,
$\sum_{k<K}\operatorname{Lip}_{\bar D(0,\varrho)}b_{k+1}\le\tfrac23$, and hence
$\zeta'_k:=\partial_\zeta\zeta_k(0)\in[\tfrac23,\tfrac43]$.
\item[(iii-d)] \emph{Exact one-loop running of the construction's coupling.} For every
$g\le\gamma_J$, every $g_0\in(0,g_-(g)]$ with $K=K(g_0,g)$, every $m\ge m_{\mathbb{T}}(g_-(g))$
and all
$0\le i<j\le K$,
\begin{equation*}
|x_i-x_j-(j-i)c_0|\le2c_2+\sum_{k=i}^{j-1}\bigl[C_\epsilon\,\epsilon(g_k)+v_k\bigr];
\end{equation*}
in particular there is a function $\Delta$ of $s$, with constants depending on
$(L,N,\mathfrak{p})$ only and not on $K$, $g$, $g_0$, $m$ or the position,
with $\Delta(s)/s\to0$, such that for $0\le s\le K$
\begin{equation*}
|x_{K-s}-x_K-s\,c_0|\le2c_2+\sum_{u=1}^{s}\bigl[C_\epsilon\,\epsilon(g_{K-u})+v_{K-u}\bigr]
\le\Delta(s).
\end{equation*}
Here $c_0=(11N^2/12\pi^2)\log L$ and
$c_2=c_2(L,N)$ are the one-loop constants of Definition~\ref{1.2:def-flow-constants}, $N=2$ as
throughout clause (iii);
$\epsilon(\cdot)$
is a function fixed in the correlation theorem of Chapter~\ref{ch:5}, tending to zero with the
coupling, $C_\epsilon$ is a constant fixed in that theorem (distinct from the auxiliary parameter
$C_\beta$ of Definition~\ref{1.2:def-flow-constants}), and $(v_k)$ is
the summable sequence of remnant bounds fixed in that theorem. The coupling is normalised so that
its
inverse square multiplies the Wilson action with the trace normalised to one,
$g^{-2}=2N/g_{\mathrm{std}}^2$, $g_{\mathrm{std}}$ the coupling in standard units; in standard
units the inverse square of the coupling grows by
$11N/(24\pi^2)=\beta_0/(8\pi^2)$ per e-fold of the scale, $\beta_0=11N/3$ the standard first
coefficient (not a member of the sequence $\beta_{k+1}$), and the factor $2N$
between the two coefficients is this normalisation.
\end{itemize}
Not asserted: $n=1$; coinciding points; decay in $\mathsf{d}$; $N\ge3$.

\item[(ii)] \emph{Limits of the curvature Schwinger functions; translation and hypercubic
invariance; reflection positivity; Osterwalder--Schrader reconstruction from the curvature
fields} ($N=2$; (H7); for (ii-$\mathbb{T}$-a) also $g\le\gamma_U$). For a space $Y$ we write
$Y^n_{\neq}$ for the set of $n$-tuples of pairwise distinct points of $Y$; $\ell_g$ is the
reference length. In (ii-$\mathbb{T}$-a) and in clause (v) the letter $n$ denotes a depth below
the reference scale and $\mathsf{p}$ the number of observables; elsewhere $n$ is the number of
observables.
\begin{itemize}
\item[(ii-$\mathbb{T}$-a)] \emph{Torus, the tuned family.} Fix $m_0\ge0$ and the torus
$\mathbb{T}$ of side $2L^{m_0}\ell_g$, hypothesis (H5) being assumed for that torus, and let
$S^n(x;\vec f)$ be the
connected $\mathsf{p}$-point function of $\mathcal{O}(f_1),\dots,\mathcal{O}(f_{\mathsf p})$ on the
bare lattice of $\mathbb{T}$ lying $n$ levels below the reference scale, at bare inverse coupling
$x$. For every $\mathsf{p}\ge2$, all real Lipschitz $f_i$ on $\mathbb{T}$ with supports pairwise
at distance at least $\mathsf{d}>0$, and every $\hat w$ with $|\operatorname{Re}\hat w|\le\varrho_2/2$
($\varrho_2:=\varrho/2$) and $|\operatorname{Im}\hat w|<\upsilon_w$, where $\upsilon_w>0$ depends
on the torus, $S^n(a_n-\hat w;\vec f)\to S^{(\hat w)}(\vec f)$ as $n\to\infty$, without passage to
subsequences, uniformly, at the rate $n^{-1/2}(\log n)^{p_0+1}$ in the bare labelling
(exponentially in the renormalised labelling of (v-d)); $a_n$ is the pin of (v-a). For real
$\hat w$ the $S^{(\hat w)}$ are multilinear, symmetric, invariant under all translations of
$\mathbb{T}$ (on $C^1$ test functions) and under $W_4$, reflection positive at every hyperplane
$x_\mu=c$ on centred separated tensors, and obey the bound of (iii-a). The first-passage family
converges only along subsequences, (ii-$\mathbb{T}$-b), and its reference offsets do not converge.
\item[(ii-$\mathbb{T}$-b)] \emph{Torus, first passage.} Fix $m\ge m_{\mathbb{T}}(g_-(g))$ and
$g_0^{(i)}\in(0,g_-(g)]$ with $K_i\to\infty$. Along a subsequence, for every $n\ge2$ and every
$n$-tuple of real $C^1$ functions on $\mathbb{T}_m$ with pairwise disjoint supports simultaneously,
$S^T_{n;K_i,m}(\vec f)\to S^T_{n;m}(\vec f)$; $S^T_{n;m}$ extends to a unique distribution of
locally finite order on $(\mathbb{T}_m)^n_{\neq}$, symmetric, invariant under all translations and
under $W_4$, and reflection positive at every hyperplane $x_\mu=c$. Under the proviso of (v-c)(2)
every such subsequential limit is some $S^{(\hat w)}$ with $|\hat w|\le\varrho_2/2$.
\item[(ii-$\mathbb{R}^4$)] \emph{Infinite volume along growing volumes.} With
$(g_0^{(j)},K_j)$ as in (ii-$\mathbb{T}$-b) and $m_j\to\infty$, along one joint diagonal
subsequence: for every $n\ge2$ and all real $f_i\in C^1_c(\mathbb{R}^4)$ with pairwise disjoint
supports, $S^T_{n;K_j,m_j}(\vec f)\to S^T_{n;\infty}(\vec f)$. The $S^T_{n;\infty}$ are
multilinear, symmetric, real, invariant under $\mathbb{R}^4\rtimes W_4$ exactly, obey the bound of
(iii-b), and are distributions of locally finite order on $(\mathbb{R}^4)^n_{\neq}$; the centred
Schwinger functions are reflection positive in the sense of \cite{OS73,OS75} at the hyperplane of
vanishing time coordinate, on separated time-ordered test functions in the open half-space.
Osterwalder--Schrader reconstruction from the linear span of centred separated-support monomials
of the curvature fields gives a separable Hilbert space $\mathcal{H}^{\mathrm{curv}}$ with a unit
vector $\Omega$, a strongly continuous semigroup $e^{-tH}$ of positive self-adjoint contractions,
$H\ge0$ with $H\Omega=0$, a strongly continuous unitary representation of $\mathbb{R}^3$ fixing
$\Omega$ whose generators commute strongly with $H$, with joint spectrum in
$[0,\infty)\times\mathbb{R}^3$, and a unitary action of $W_3$.
\end{itemize}
Not asserted: convergence on $\mathbb{R}^4$ without passage to subsequences in the volume;
uniqueness in the cutoff on $\mathbb{R}^4$; $n=1$; coinciding points; temperedness, field
operators, Wightman functions; the spectrum condition (joint energy--momentum spectrum in the
closed forward cone) beyond $H\ge0$; uniqueness of $\Omega$ and clustering; a mass gap;
$O(4)$ invariance (that is clause (vi)); regularisation independence; limit tuples of Lipschitz
test functions on $\mathbb{R}^4$; identification of $\mathcal{H}^{\mathrm{curv}}$ with spaces built
from the block variables; anything on the block algebra of Definition~\ref{1.5:def-block-average}.
The inequality $\dim\mathcal{H}^{\mathrm{curv}}\ge2$ in every limit is a consequence stated in
clause (iv).

\item[(iv)] \emph{Non-triviality: the curvature two-point witness, uniformly in the volume}
($N=2$; in the regime
of (iii)). The free integer exponent $q$ of \cite[(2.28)]{B-CMP119}, one of the auxiliary
parameters $\mathfrak{p}$, is chosen subject to the one-sided lower bound
$q\ge21\mathfrak{p}_0+1$, where $\mathfrak{p}_0$ is a second free integer exponent of
Ba{\l}aban's construction, also among $\mathfrak{p}$. For every $\vartheta\in(0,1]$ there exist
$\gamma_W^{\mathbb{T}}=\gamma_W^{\mathbb{T}}(L,\mathfrak{p},\gamma,\vartheta;\varrho)\le\gamma_J$;
depth floors $s_W(g)$, $s_{\mathrm{vol}}^\star(g)$, the volume depth floor of Chapter~\ref{ch:8}
evaluated at the reference torus exponent $m_{\mathbb{T}}(g_-(g))$,
and $s_T(g)=s_T(g;\vartheta,\varrho)$, each, for the given $\vartheta$ (and $\varrho$ for $s_T$),
depending on the couplings only through $g$, each a lower bound on the depth,
$s_W$ non-decreasing and depending on $g$ only through $\log_L g^{-2}$,
$s_T$ polylogarithmic in $g^{-2}$ of degree at most $4r_{\mathrm{pr}}+1$, with $r_{\mathrm{pr}}$ an
integer fixed in Chapter~\ref{ch:8}; an integer
$k_W=k_W(L,\mathfrak{p})\ge1$; constants $c_{\mathcal{K}},C_{\mathcal{K}}>0$; a constant $C_T(g)$; a
non-decreasing function $C_W(s)\ge1$ of $\bar x_s:=g^{-2}+B_\sharp(s+1)$, polylogarithmic in it;
and an integer $r_0$, the constant depth offset of the Gaussian extraction
(Chapter~\ref{ch:8}). None of these depends on $K$,
$g_0$, $m$, $\mathsf{d}$ or the position; $C_W$ and $c_{\mathcal{K}}$ depend on $\vartheta$.

\emph{Notation.} $b_0:=(N^2-1)/2$; $\zeta'_k$ is as in (iii-c). $\mathcal{A}$ is the real free
lattice gauge field at unit coupling on the bare lattice of $\mathbb{T}_m$, modulo gauge and
harmonic modes; $Q$ is the linearised block-averaging map to level $K-s$; $F_B$ is the lattice
curvature of $\mathbb{E}[\mathcal{A}\,|\,Q\mathcal{A}]$; and
\begin{equation*}
\mathcal{K}_s(f_1,f_2):=\sum_{p,p'}f_1(x(p))\,f_2(x(p'))\,\bigl(\mathbb{E}[F_B(p)F_B(p')]\bigr)^2,
\end{equation*}
whose continuum counterpart, multiplied by $b_0$, is
$3(N^2-1)\pi^{-4}\iint f_1(x)f_2(y)|x-y|^{-8}\,dx\,dy$; $\mathcal{K}_s^\star$ denotes
$\mathcal{K}_s$ read on the reference torus $\mathbb{T}_{m_{\mathbb{T}}(g_-(g))}$. Further
$s(\mathsf{d}):=\min\{s\ge0:C_W(s)L^{-s}\le\mathsf{d}\}$, which is independent of $K$ and tends to
$\infty$ as $\mathsf{d}\downarrow0$;
\begin{equation*}
r(s):=\lceil\log_L(20LC_W(s))\rceil+\lceil(\mathfrak{p}_0/2+1)\log_L\log\bar x_s\rceil+r_0;
\end{equation*}
$\underline{x}_s:=g^{-2}+\lambda_\sharp s-c_5$; $s_*$ is the depth, fixed in
Chapter~\ref{ch:8} and depending only on $L$ and the parameters fixed before it, from which on the
large-field and boundary contributions to the plaquette average at depth $s$ obey the smallness
conditions used there; and $s_W(g)$ is the least $s$ such that
\begin{equation*}
s-r(s)\ge s_*,\qquad s\ge m_{\mathbb{T}}(g),\qquad 2B_\sharp(r(s)+1)\le\underline{x}_s.
\end{equation*}

\emph{Test functions.} $\mathfrak{T}(\mathsf{d};\vartheta)$ is the set of pairs of non-negative
Lipschitz functions $f_1,f_2\not\equiv0$ on $\mathbb{T}_m$ with $\operatorname{supp}f_i$ in a ball
of radius $\mathsf{d}$, $\int f_i\ge\vartheta\|f_i\|_\infty\mathsf{d}^4$,
$\operatorname{Lip}f_i\le\|f_i\|_\infty/(\vartheta\mathsf{d})$, and
$\mathsf{d}\le\operatorname{dist}(\operatorname{supp}f_1,\operatorname{supp}f_2)\le6\mathsf{d}$;
this two-sided separation window is a condition on the test functions.

\emph{Statement.} For every $g\le\gamma_W^{\mathbb{T}}$, every $g_0\in(0,g_-(g)]$ with
$K=K(g_0,g)$, every $m\ge m_{\mathbb{T}}(g_-(g))$ (hypothesis (H5), and nothing more on $m$), every
$\mathsf{d}$ with
\begin{gather}
\max\{s_W(g),\,s_{\mathrm{vol}}^\star(g),\,s_T(g)\}\le s(\mathsf{d}),\label{3.1:eq-main-1}\\
K-s(\mathsf{d})\ge k_W,\label{3.1:eq-main-2}
\end{gather}
$s:=s(\mathsf{d})$, and every $(f_1,f_2)\in\mathfrak{T}(\mathsf{d};\vartheta)$,
\begin{equation}\label{3.1:eq-main-3}
S^T_{2;K,m}(f_1,f_2)=\zeta'^{\,2}_{K-s}\,b_0\,g_{K-s}^4\,\mathcal{K}_s^\star(f_1,f_2)
\,\bigl(1+R_{K,m}\bigr),\qquad |R_{K,m}|\le\tfrac12,
\end{equation}
with $\zeta'_{K-s}\in[\tfrac23,\tfrac43]$ and
\begin{equation*}
c_{\mathcal{K}}\|f_1\|_\infty\|f_2\|_\infty\le\mathcal{K}_s^\star(f_1,f_2)
\le C_{\mathcal{K}}\|f_1\|_\infty\|f_2\|_\infty.
\end{equation*}
Hence, uniformly in $K$, $g_0$, $m$ and the position,
\begin{align*}
&\tfrac29\,b_0\,g_{K-s}^4\,\mathcal{K}_s^\star\le S^T_{2;K,m}(f_1,f_2)
\le\tfrac83\,b_0\,g_{K-s}^4\,\mathcal{K}_s^\star,\\
&\tfrac29\,b_0c_{\mathcal{K}}\|f_1\|_\infty\|f_2\|_\infty\bigl(g^{-2}+B_\sharp(s+1)\bigr)^{-2}
\le S^T_{2;K,m}(f_1,f_2)\\
&\qquad\le\tfrac83\,b_0C_{\mathcal{K}}\|f_1\|_\infty\|f_2\|_\infty
\bigl(g^{-2}+\lambda_\sharp s-c_5\bigr)^{-2}.
\end{align*}
Moreover (\emph{volume transfer}), for all $m\ge\tilde m\ge m_{\mathbb{T}}(g_-(g))$ and all real
Lipschitz $f_1,f_2$ of scale $\mathsf{d}$ --- both supported in one ball of radius
$5\mathsf{d}$, with $\operatorname{Lip}f_i\le\|f_i\|_\infty/(\vartheta\mathsf{d})$ and
$L^{-K}\le\mathsf{d}\le1/40$, and with no condition on the sign of $f_i$, on $\int f_i$, or on
the distance between the supports ---
\begin{equation*}
\bigl|S^T_{2;K,m}(f_1,f_2)-S^T_{2;K,\tilde m}(f_1,f_2)\bigr|\le C_T(g)\,\mathsf{d}^2
\|f_1\|_\infty\|f_2\|_\infty.
\end{equation*}
\emph{Consequences}, in every limit of clause (ii), those of (ii-$\mathbb{R}^4$) included: the
smeared curvature has strictly positive variance, so $\dim\mathcal{H}^{\mathrm{curv}}\ge2$; the
limit is not scale invariant under any assignment of a scaling dimension to the curvature; it is
reproduced by no free Maxwell field at any charge; and the coupling read off the two-point function
at separation $\mathsf{d}$ is strictly positive and tends to zero logarithmically as
$\mathsf{d}\downarrow0$, at the two-sided rate of clause (0). Every limit of clause (ii) is thus
non-trivial in the sense of Definition~\ref{3.1:def-non-trivial}. By \eqref{3.1:eq-main-3}
together with (iii-d), the coupling $g(\mathsf{d}):=g_{K-s(\mathsf{d})}$ obeys
\begin{equation*}
1/g(\mathsf{d})^2=(c_0/\log L)\ln(1/\mathsf{d})\bigl(1+o(1)\bigr)\qquad(\mathsf{d}\to0),
\end{equation*}
uniformly in the
cutoff; the charge read off the two-point function, whose fourth power differs from
$g(\mathsf{d})^4$ by a factor in $[\tfrac29,\tfrac83]$, therefore runs with an inverse-square slope
between $(3/8)^{1/2}$ and $(9/2)^{1/2}$ times the one-loop Yang--Mills value of (iii-d), at $N=2$.
Not asserted: that the limit fails to be Gaussian; $O(4)$ invariance; a mass gap; anything on
block observables,
Wilson loops or the block algebra; one source, overlapping or signed sources, separation larger
than $6\mathsf{d}$, $\vartheta\to0$, depths outside the range \eqref{3.1:eq-main-1},
\eqref{3.1:eq-main-2}; $\zeta'_{K-s}\to1$, or any bound on $|\zeta'_{K-s}-1|$ beyond the displayed
range; $N\ge3$; a volume-uniform bound on the
large-field regions that are never healed (a large-field region is healed when a later step
absorbs its large-field factor back into the small-field format and re-files the region into the
regular expansion, Chapter~\ref{ch:8}); volume-uniform rarity of
large-field regions;
convergence or uniqueness in $m$; the volume transfer for $n\ge4$ observables; the power
$\mathsf{d}^4$ in place of $\mathsf{d}^2$ in the transfer bound.

\item[(v)] \emph{Uniqueness up to one real parameter, the renormalised coupling at the reference
length} ($N=2$). For every $\varrho>0$ there is $\gamma_U$ such that the following holds for every
$g\le\gamma_U$. Put $X:=g^{-2}$ and $\lambda_U:=c_0/2$, with $c_0$ the one-loop coefficient of
(iii-d) in the normalisation stated there (Definition~\ref{1.2:def-flow-constants}). The pure
numbers entering
this clause are among the numbers fixed at Stage~0 (Notation~\ref{2.3:not-stages-first}).
\begin{itemize}
\item[(v-a)] \emph{Exact pin.} For every depth $n\ge0$ there is exactly one real $a_n$, depending
on $(L,\mathfrak{p},g,n)$ and on no torus, such that the $n$-step recursion from $x_0=a_n$ (with
infinite-lattice coefficients and the coupling-dependent choices fixed in terms of $X$) is
admissible and ends at $x_n=X$ exactly; moreover, with a constant $c'$,
\begin{equation*}
\lambda_Un-2c_2\le a_n-X\le3\lambda_Un+2c_2,\qquad
\lambda_U(n-l)-c'\le a_n-a_l\le3\lambda_U(n-l)+c'.
\end{equation*}
No per-step monotonicity is asserted.
\item[(v-b)] \emph{The one-parameter limit family.} On each torus $\mathbb{T}$ of side
$2L^{m_0}\ell_g$, for $\mathsf{p}\ge2$ and separated Lipschitz $\vec f$, the limits
$S^{(\hat w)}(\vec f)$ of (ii-$\mathbb{T}$-a) exist for $|\operatorname{Re}\hat w|\le\varrho_2/2$;
$\hat w\mapsto S^{(\hat w)}(\vec f)$ is real-analytic on $(-\varrho_2/2,\varrho_2/2)$, real on
reals, and holomorphic on the strip $|\operatorname{Im}\hat w|<\upsilon_w$ (holomorphy on a larger
domain is not asserted). The depth-to-depth coupling map $\mathfrak{n}_n$, defined by equality of the
reference-scale couplings of $(n+1,a_{n+1}-\mathfrak{n}_n(\hat w))$ and $(n,a_n-\hat w)$, is
holomorphic on $D(0,\varrho)$, with $\mathfrak{n}_n(0)=0$ and
$|\mathfrak{n}_n(\hat w)-\hat w|\le\Upsilon_n|\hat w|$, where
$\Upsilon_n\le C\,n^{-3/2}(\log n)^{p_0+1}$, so $\sum_n\Upsilon_n<\infty$. The value $N_n$ of
$\zeta'$ at the tuned point lies in $[\tfrac23,\tfrac43]$, $N_{n+1}\mathfrak{n}_n'(0)=N_n$, and
$(N_n)$ converges. The $a_n$, $\mathfrak{n}_n$, $N_n$, $\Upsilon_n$ and all constants are
independent of the torus, which enters only through $\upsilon_w$, the rate constant and a floor
$n_0$ on $n$.
\item[(v-c)] \emph{Exhaustion.} (1) Every real sequence $(a'_n)$ with $|a'_n-a_n|\le\varrho_2/2$
has $S^n(a'_n;\vec f)-S^{(a_n-a'_n)}(\vec f)\to0$; hence two admissible bare families whose
offsets converge to the same $\hat w$ have equal limits $S^{(\hat w)}$, and an admissible family
converges (without passage to subsequences) if and only if $S^{(a_n-a'_n)}$ does, in particular
whenever its offset converges. (2) If moreover $\varrho\ge6(B_\sharp+1)$ and $g\le\gamma_{U'}$
for a further threshold $\gamma_{U'}\le\gamma_U$, every bare inverse coupling $x$ whose
first-passage run stops at depth $K$ has $|x-a_K|\le\varrho_2/2$, and every subsequential limit of
(ii-$\mathbb{T}$-b) is some $S^{(\hat w)}$ with $|\hat w|\le\varrho_2/2$. (3) The set of continuum
limits reachable on $\mathbb{T}$ from bare data $x$ with $|x-a_n|\le\varrho_2/2$ is exactly
$\{S^{(\hat w)}:|\hat w|\le\varrho_2/2\}$. Not asserted: calibration by a physical quantity (for
the injectivity of $\hat w\mapsto S^{(\hat w)}$ on each torus and on $\mathbb{R}^4$ see (v-f)).
\item[(v-d)] \emph{The marginal coordinate is the renormalised coupling.}
(i) Re-file each healed large-field remnant (healed in the sense recalled in clause (iv)) created
at step $i$ into the regular family at the
re-filing level $i''=i+\log_L(M\check R_i)$, at which its partition cube is a single cell
($M=L^{m'}$ the auxiliary parameter of Ba{\l}aban's large-field geometry among $\mathfrak{p}$,
$\check R_i$ the size parameter of the remnant); label the family by
$t:=X-\hat x$, where $\hat x$, the canonical comparison point, is the reference-scale inverse
coupling of the re-filed run (by ($\gamma$) below it is the limit $\hat x_0$ of the reference-scale
inverse couplings of the tuned runs), and normalise the
$\mathsf{p}$-point functions by $\widehat{\mathsf{N}}_n(t)^{\mathsf{p}}$. Then the $t$-labelled
normalised family converges at an exponential rate $n^c\rho^{n/2}$ for some constant $c$, with
$\rho<1$ a contraction constant bounded below in terms of the numbers fixed at Stage~0
(Notation~\ref{2.3:not-stages-first}) and the flow constants of
Definition~\ref{1.2:def-flow-constants}
only, and the run differences of the
re-filed expansion are at most $C_E\rho^{\,n-r}$, with $r$ the depth index of these differences
fixed in Chapter~\ref{ch:9}, and with no polynomial
floor. For the
family pinned on the $x_K$ of clause (0), for the bare operator, and for the family labelled by a
bare shift $\hat w\ne0$, no exponential rate in $n$ is asserted. The bare offset $\hat w$ and the
renormalised coordinate $t$ parametrise the same family, at each depth real-analytically with slope
in $[\tfrac23,\tfrac43]$ near $0$, by (iii-c).
(ii) Two admissible families on $\mathbb{T}$ whose reference-scale inverse couplings converge to
the same value have the same limit; this follows from (v-c)(1) and the fact that the map
$\hat w\mapsto x_{\mathrm{ref},n}(\hat w)$, the reference-scale inverse coupling of the run from
$a_n-\hat w$, is bi-Lipschitz on $|\hat w|\le\varrho_2/2$ with lower constant at least
$\tfrac13$ and upper constant at most $\tfrac53$, uniformly in $n$, and converges uniformly as
$n\to\infty$.
(iii) \emph{Up to scale.} The data of the construction for a limit are $(m_0,\hat x)$,
equivalently $(m_0,g,\hat w)$. For $s\ge1$ put $X_s:=X+c_0s$, with $c_0=2\lambda_U$ the constant
of (iii-d), $\tilde g_s:=X_s^{-1/2}<g$ (a reference coupling, not one of the running couplings
$g_k$),
$D_s(g):=2c_2+\sum_{u<s}\iota_u(X)$ with $\iota_u$ the per-level drift constants of the
correlation theorem, and $j_U:=1+\lceil(\varrho_2+c')/\lambda_U\rceil$. Write
$S^T_{\mathsf{p};n,m_0}[x](\vec f)$ for the connected $\mathsf{p}$-point function on the bare
lattice ($2L^{m_0+n}$ sites a side) of the torus $\mathbb{T}_{m_0}$ of side $2L^{m_0}\ell_g$ at
bare inverse coupling $x$, and $S^{(\hat w)}_{\mathbb{T}_{m_0}}[g]$ for its limit of (v-b).
\begin{itemize}
\item[($\alpha$)] For $0\le s\le n$, every bare $x_0$ and every tuple,
\begin{equation*}
S^T_{\mathsf{p};n,m_0}[x_0](\vec f)=S^T_{\mathsf{p};n-s,m_0+s}[x_0](\vec f(L^{-s}\,\cdot\,)):
\end{equation*}
one lattice integral, the plaquette barycentres being named with the reference length moved $s$
levels toward the cutoff. The tuned bare datum carries the two labels exactly:
$a_n(g)-\hat w=a_{n-s}(\tilde g_s)-(\hat w+\delta_{n,s}(g))$, with $\delta_{n,s}$ real,
$|\delta_{n,s}|\le\tfrac32D_s(g)$, and $\delta_{n,s}(g)\to\delta_s(g)$ summably in $n$.
\item[($\beta$)] Under the one-sided condition $\varrho\ge48c_2$ and
$\sum_{u<s}\iota_u(g^{-2})\le\varrho/24$ (a floor on $\varrho$ and, for each $s$, a ceiling on
$g$; equivalently, for each $g$, a range $s\le s_{\mathrm{dil}}(g;\varrho)$ with
$s_{\mathrm{dil}}\uparrow\infty$ as $g\downarrow0$), for every real $\hat w$ with
$|\hat w|<\varrho_2/4$ and all separated real Lipschitz $\vec f$ on $\mathbb{T}_{m_0}$,
\begin{equation*}
S^{(\hat w)}_{\mathbb{T}_{m_0}}[g](\vec f)
=S^{(\hat w+\delta_s(g))}_{\mathbb{T}_{m_0+s}}[\tilde g_s](\vec f(L^{-s}\,\cdot\,)),
\end{equation*}
both members being full limits, without subsequences, of the families of (v-b) at $(m_0,g)$ and
at $(m_0+s,\tilde g_s)$: the continuum limit of the member $(m_0,\hat x)$ is the $s$-fold change
of units
of the continuum limit of the member $(m_0+s,\hat x_s)$, with $\hat x_s$ as in ($\gamma$) and
$\hat x_s-\hat x\in[\lambda_\sharp s-c_5,B_\sharp s]$; the identifications compose in $s$.
\item[($\beta'$)] Under the floor $\varrho\ge6(B_\sharp+1)$ and the ceiling $g\le\gamma_{U'}$ of
(v-c)(2): for every $X'$ with $g':=X'^{-1/2}\le\gamma_{U'}$ (on either side of $X$) and every real
$|\hat w|\le\varrho_2/2$ there are integers $\varsigma$, $s$ and a real $\hat w'$ with
$|\hat w'|\le\varrho_2/2$, depending on $(L,\mathfrak{p},g,g',\hat w)$ and neither on the torus
nor on the test functions, with $|\varsigma-s|\le j_U-1$, $s\,(X'-X)\ge0$ (so $s$ and $X'-X$ are
not of opposite signs), and
\begin{equation*}
\frac{|X'-X|-B_\sharp}{B_\sharp}<|s|<\frac{|X'-X|+B_\sharp+c_5}{\lambda_\sharp},
\end{equation*}
such that for every $m_0$ above a floor depending on $g,g'$ only and all separated real Lipschitz
$\vec f$,
\begin{equation*}
S^{(\hat w)}_{\mathbb{T}_{m_0}}[g](\vec f)
=S^{(\hat w')}_{\mathbb{T}_{m_0+\varsigma}}[g'](\vec f(L^{-\varsigma}\,\cdot\,)):
\end{equation*}
the member
$(m_0;X,\hat w)$ is the exact $L^\varsigma$-dilate of the member $(m_0+\varsigma;X',\hat w')$, the
full limit of its own tuned family; one run of the common bare data with its own thresholds
first-passes at $X$ and, $|s|$ levels toward the cutoff, at $X'$. Moreover, for every
$s\ge s^\flat(g)$, a floor depending on $g$ only, there is such an $X'>X$ with
$\varsigma\in[s,s+\lfloor c_5/\lambda_\sharp\rfloor]$ ($\varsigma=s$ when $c_5<\lambda_\sharp$).
\item[($\gamma$)] For every fixed $s\ge0$ the inverse coupling $s$ levels above the reference of
the tuned runs from $(a_n(g)-\hat w)_n$ converges as $n\to\infty$, uniformly in
$|\hat w|\le\varrho_2/2$, to a real-analytic $\hat x_s(\hat w)$; $\hat x_0=\hat x$ and
$\hat x_s-\hat x\in[\lambda_\sharp s-c_5,B_\sharp s]$. Hence, along the tuned family, the limits
$g_{\infty,s}$ of the fixed-depth windows of clause (0) exist without passage to subsequences, and
$g_{\infty,s}=\hat x_s^{-1/2}$.
\item[($\delta$)] Nothing relates distinct $\hat x$ at one torus and one reference length. Not
asserted: that the tuned bare data $a_n(g)$ and $a_{n-s}(g')$ coincide for any $g'$; a prescribed
exponent across an arbitrary reference coupling beyond the $\lfloor c_5/\lambda_\sharp\rfloor$ of
($\beta'$); ($\beta$) outside its window (it follows from clause (v) read at $\tilde g_s$ with the
radius $\varrho_s:=48c_2+24\sum_{u<s}\iota_u(g^{-2})+8|\hat w|$ whenever
$\tilde g_s\le\gamma_U(\varrho_s)$; displayed, not adopted); convergence of the parameter $\hat w'$
of ($\beta'$); the
chart maps between the two labellings, $(X,\hat w)\mapsto(X',\hat w')$ and
$\hat x_s\mapsto\hat x'$; an exact continuous dilation (for the short-distance germ see (v-f));
dimensional transmutation of whole torus theories or of the sets of limits on $\mathbb{R}^4$.
\end{itemize}
\item[(v-e)] \emph{The one-lattice Lipschitz step.} On one lattice and with one frozen sequence of
the coupling-dependent auxiliary choices, for two input
lists in the class of the correlation theorem (complex sources allowed), the outputs of
Ba{\l}aban's complete step, the large-field operations $\mathbf{R}$, $\mathbf{B}$, $\mathbf{T}$
included, differ by at most cutoff-free multiples of the input difference, measured in the weighted
unit of the correlation theorem (Chapter~\ref{ch:5}): energy and small-field entries with weights
proportional to a function of $g_k$, the pure marginal slot at second order, with weight the
square of that function times the size of the slot, boundary and
large-field entries by sums weighted by $e^{-p_0(g_k)}$, with $p_0(\cdot)$ the degree function of
Definition~\ref{1.5:def-domains}; no constant
depends on $K$, $k$, the torus or the position. Old large-field regions are forgotten
geometrically from their first step: a fixed
factor strictly less than one per step
of life, paid from Ba{\l}aban's entropy constant under the one-sided condition $p_0\ge6r+1$ on
the exponent of the degree function, with
$r$ the auxiliary parameter of Ba{\l}aban that enters his condition $p_0\ge5r$. This
is the step that (ii-$\mathbb{T}$-a) and (v-b) iterate.
\item[(v-f)] \emph{Covariance and injectivity of the renormalised coupling.} Under the one-sided
floor $\varrho\ge\varrho_\Lambda$ on the radius of (H7), the ceiling $g\le\gamma_\Lambda$ and the
floor $m\ge m_\natural(g)$ on the torus exponent, where $\varrho_\Lambda$, $\gamma_\Lambda$ and the
function $m_\natural$ of $g$ alone are fixed in Chapter~\ref{ch:9}, the following holds for every
odd $L\ge L_0$. The tuned bare inverse couplings have convergent increments,
\begin{equation*}
a_{n+1}-a_n\to\alpha_\infty\in[\lambda_U,3\lambda_U]\qquad(n\to\infty),
\end{equation*}
and
\begin{equation*}
S^{(\hat w)}_{\mathbb{T}_m}(\vec f)=S^{(\hat w-\alpha_\infty)}_{\mathbb{T}_{m+1}}
\bigl(\vec f(\,\cdot\,/L)\bigr)
\qquad\text{whenever }\hat w,\ \hat w-\alpha_\infty\in(-\varrho_2/2,\varrho_2/2);
\end{equation*}
the slope $b:=\alpha_\infty/\log L$ depends neither on $L$ nor on the reference coupling; for any
two coprime odd $L,L'\ge L_0$ the ratio $\ell_L(x_0)/\ell_{L'}(x_0)$ of the physical lengths of
the two frames, $\ell_L(x_0)$ being the physical length of the frame of blocking factor $L$ at
bare inverse coupling $x_0$, converges as $x_0\to\infty$; and for all
$u_1,u_2\in(-\varrho_2/2,\varrho_2/2)$, all $m,m'\ge m_\natural(g)$, $\mathsf{p}\in\{2,3\}$,
every $\mathsf{d}\le\mathsf{d}_\Lambda$, a length fixed in Chapter~\ref{ch:9}, and every tuple
$\vec f\in\mathfrak{P}_{\mathsf{p}}(\mathsf{d};\vartheta)$ of $C^2$ functions of scale
$\mathsf{d}$, the class fixed in Chapter~\ref{ch:9}, the two- and three-point functions satisfy,
with a constant $C$,
\begin{equation*}
\Bigl|S^{(u_1)}_{\mathbb{T}_m}(\vec f)
-S^{(u_2)}_{\mathbb{T}_{m'}}\bigl(\vec f(e^{-(u_1-u_2)/b}\,\cdot\,)\bigr)\Bigr|
\le C\,\mathsf{d}^{4-\mathsf{p}}\prod_i\|f_i\|_\infty.
\end{equation*}
Consequently $\hat w\mapsto S^{(\hat w)}_{\mathbb{T}_m}$ is injective on
$[-\varrho_2/2,\varrho_2/2]$ for every $m\ge m_\natural(g)$, and the sets of limit points on
$\mathbb{R}^4$ at two distinct renormalised couplings are disjoint. Not asserted: dilation
covariance of whole torus theories or of the sets of limit points on $\mathbb{R}^4$; a rate; the
value of $b$; anything infrared.
\end{itemize}
Not asserted: $\mathsf{p}=1$; $\log Z$; coinciding points; Wilson loops; $N\ge3$; uniformity in the
torus size; uniqueness of the state on the block algebra (the block-law statements of
Remark~\ref{3.2:rem-block-laws}, whose proofs are outstanding);
uniqueness in the volume; convergence of the effective actions; exponential rates in any bare
labelling. No clause of the theorem uses \cite[Theorem~2, (0.31)]{B-CMP109}.

\item[(vi)] \emph{Full rotation invariance under the far-sphere flux hypothesis} ($N=2$; in the
regime of (iii); on the limits of (ii-$\mathbb{R}^4$) only). The hypothesis of the clause is the
following infrared flux hypothesis; the clause is conditional on it, and the theorem asserts
nothing about its truth.

\emph{Hypothesis} (IR-fl)$_{\omega_0}$ (the far-sphere flux hypothesis). Let $\omega_0\ne0$ be
one antisymmetric $4\times4$ matrix, the generator of the rotations whose current is considered;
for $\mathcal{R}>0$ let $\chi_{\mathcal{R}}$ be a radial cutoff function at radius $\mathcal{R}$,
and let $\mathfrak{B}^{\mathrm{sh}}_{\chi_{\mathcal{R}}}$ be the flux of the lattice rotation
current with generator $\omega_0$ through the spheres about the origin, averaged over a shell at
radius $\mathcal{R}$ by means of $\chi_{\mathcal{R}}$ as Chapter~\ref{ch:10} fixes it;
the current is that of the rotational Ward identity of
Chapter~\ref{ch:10}, whose boundary term is this flux, and these objects are defined there. The
hypothesis is: for every $n\ge2$ and every $n$-tuple $\vec f$ of test functions with pairwise
disjoint supports in the class of the rotational Ward identity of Chapter~\ref{ch:10},
\begin{equation*}
\liminf_{\mathcal{R}\to\infty}\,\limsup_{j\to\infty}\,
\bigl|\langle\mathfrak{B}^{\mathrm{sh}}_{\chi_{\mathcal{R}}};\mathcal{O}(f_1);\dots;
\mathcal{O}(f_n)\rangle^{\mathrm{conn}}_{K_j,m_j}\bigr|=0,
\end{equation*}
the upper limit in $j$ being taken first and the lower limit in $\mathcal{R}$ second; here
$(K_j,m_j)$ is the joint diagonal subsequence of (ii-$\mathbb{R}^4$), and
$\langle\cdot\rangle^{\mathrm{conn}}_{K_j,m_j}$ is the connected correlation at cutoff $K_j$ on
$\mathbb{T}_{m_j}$. In words: the flux of the lattice rotation current through far spheres
vanishes along one sequence of radii.

\emph{Sufficient conditions.} Let $\Theta_{\nu\sigma}$ be the smeared, dimension-four,
$W_4$-tensor lattice stress-tensor density of the Wilson action, taken traceless at tree level,
$x_0$-weighted, where $x_0=g_0^{-2}$ is the bare inverse coupling of
Definition~\ref{1.2:def-couplings}, with its $x_0$-free admixtures. The hypothesis holds if the
connected correlation of $\Theta$ at
distance $\mathcal{R}$ from the supports with $\mathcal{O}(f_1),\dots,\mathcal{O}(f_n)$ is
$o(\mathcal{R}^{-4})$ as $\mathcal{R}\to\infty$,
uniformly in the cutoff $K$ (and in $m$ along the joint diagonal), in the shell-averaged form with
the upper limit in $K$,
\begin{equation*}
\limsup_K\int_{\mathcal{R}\le|y|\le2\mathcal{R}}
\bigl|\langle\Theta(y);\mathcal{O}(f_1);\dots;\mathcal{O}(f_n)\rangle^{\mathrm{conn}}_K\bigr|\,dy
\;\to\;0
\end{equation*}
along one sequence $\mathcal{R}_j\to\infty$. This holds in turn under the following stress-tensor
decay form: for some $\delta'>0$, every closed
unit cube $\square$ with $\operatorname{dist}(\square,\bigcup_i\operatorname{supp}f_i)\ge1$ and
every Lipschitz $\psi$ supported in $\square$,
\begin{align*}
&\sup_K\bigl|\langle\Theta_{\nu\sigma}(\psi);\mathcal{O}(f_1);\dots;\mathcal{O}(f_n)
\rangle^{\mathrm{conn}}_K\bigr|\\
&\qquad\le C(\vec f)\,\|\psi\|_\sharp\,
\operatorname{dist}\Bigl(\square,\textstyle\bigcup_i\operatorname{supp}f_i\Bigr)^{-4-\delta'};
\end{align*}
that is, decay strictly faster than $|x|^{-4}$, uniformly in the cutoff; a uniform decay slower
than $|x|^{-4}$ does not suffice (at tree level the decay is $|x|^{-8}$, and a mass gap would give
exponential decay).

\emph{Statement.} Assume (IR-fl)$_{\omega_0}$. Then for every limit $S^T_{n;\infty}$ of
(ii-$\mathbb{R}^4$),
every $n\ge2$, every $\Lambda\in O(4)$ and all real $f_i\in C^1_c(\mathbb{R}^4)$ with pairwise
disjoint supports,
\begin{equation*}
S^T_{n;\infty}(f_1\circ\Lambda^{-1},\dots,f_n\circ\Lambda^{-1})=S^T_{n;\infty}(f_1,\dots,f_n);
\end{equation*}
hence the $S^T_{n;\infty}$ are invariant under $\mathbb{R}^4\rtimes O(4)$, and $O(3)$ acts on
$\mathcal{H}^{\mathrm{curv}}$ strongly continuously by unitaries commuting with $H$ and extending
the action of $W_3$. The ultraviolet half is carried inside the proof by two statements: the
rotational Ward identity, an exact lattice identity whose $O(4)$-breaking term is $a^2$, $a$ the
lattice spacing, times insertions of dimension-six operators in the antisymmetric rank-two
$W_4$-channel, plus a boundary flux, the one $O(4)$-breaking hypercubic scalar modulo total
derivatives being
\begin{equation*}
O_1=\sum_{\mu,\nu}\operatorname{tr}\bigl(D_\mu F_{\mu\nu}\,D_\mu F_{\mu\nu}\bigr)
\end{equation*}
(two as densities; no cubic term in $F$ breaks $O(4)$; nothing of dimension at most five occurs in
that channel); and the dimension-six coefficient bound, by which the $a^2$-insertion of $O_1$,
carried through Ba{\l}aban's flow for connected correlations with $n\ge2$ separated sources, is
bounded by a constant times $a^{2-\eta}$ for a fixed $\eta<2$ in the small-field part; in the
large-field part either the large-field operations return the insertion in locality-preserving
form with its accumulated contraction, giving the same bound, or the insertion merely tends to
zero as $a\to0$, no power of $a$ being asserted, which suffices equally.
Chapter~\ref{ch:10} establishes the ultraviolet half from clauses (0), (i), (ii-$\mathbb{R}^4$)
and (iii), so that (IR-fl)$_{\omega_0}$ is the only hypothesis of the clause.
Not asserted: tensor multiplets; rotation of block observables or of Wilson loops; anything on a
torus beyond the finite-volume covariance under rotations for which the rotated family of lattices
is commensurate with the original one; that
(IR-fl)$_{\omega_0}$ holds;
clause (vi) along the tuned family of (v) on $\mathbb{R}^4$.
\end{itemize}
\end{theoremzero}

Part~II proves the theorem clause by clause: (0) and (i) in Chapter~\ref{ch:4}; the last sentence
of (0), that the $K$ steps on every torus $\mathbb{T}_m$ produce these couplings, rests on the
torus-independence lemma of that chapter, by which the couplings of clause (0) do not depend on the
torus exponent $m$; that lemma is itself proved as an implication, as described below under (iv).
(iii-a) and (iii-c) are proved in
Chapter~\ref{ch:5}; (iii-b) in Chapter~\ref{ch:6}; (iii-d) in Chapter~\ref{ch:5},
as an implication from the flow statement of the correlation theorem (its bounds on the
coupling increments) and the
one-loop bound on the coupling increments whose constants $c_0$, $c_2$ are those of
Definition~\ref{1.2:def-flow-constants}, under the
identification of the correlation theorem's coupling increment with the increment $\beta_{k+1}$ of
Definition~\ref{1.2:def-couplings}, both being Ba{\l}aban's coefficient of the complete regular
part; (ii) in Chapter~\ref{ch:7}. Clause (iv) is proved in Chapter~\ref{ch:8} as an implication
from the statements from which (iii) is proved and, through the normalising factor in the passage
to all volumes, from the torus-independence lemma of Chapter~\ref{ch:4}. That lemma is itself
proved as an implication from earlier lemmas of this account and Ba{\l}aban's definitions
\cite[p.~394; p.~410, Corollary~3.8]{B-CMP99b}, the agreement of the localised objects across tori
of different size on their common domain being part of those definitions; it uses one further
one-sided lower bound on the auxiliary parameter $M$, displayed with that lemma in
Chapter~\ref{ch:4}, and the convergence and
analytic extension of the expansions on the joined localisation domains as stated in
\cite[p.~192, (1.70)]{B-CMP122a} and \cite[p.~377, (1.68)--(1.69)]{B-CMP122b}. The second-order
remainder estimate of the Gaussian extraction, on which (iv) rests, is proved as an implication
from the correlation theorem, one cluster-expansion lemma and Ba{\l}aban's regularity theorem
\cite[Theorem~1]{B-CMP102b}, under the lower bound $q\ge21\mathfrak{p}_0+1$ displayed in clause
(iv); it is proved in that form in Chapter~\ref{ch:8}. A volume-uniform bound on never-healed
large-field regions is not proved, and no clause of
the theorem uses it. Clause (v) is proved in Chapter~\ref{ch:9}, as an implication from the
statements listed at the end of that
chapter, some of which are stated there with proof incomplete or with proof outstanding;
(vi), as a deduction from the far-sphere flux hypothesis
(IR-fl)$_{\omega_0}$,
in Chapter~\ref{ch:10}; Chapter~\ref{ch:11} assembles the clauses; and Chapter~\ref{ch:13} indexes
every statement of Part~II not proved in full.
The short-distance uniqueness of the limits on $\mathbb{R}^4$ at a given renormalised coupling is
a corollary of clauses (iii), (iv) and (v), stated and proved after clause (v) in
Chapter~\ref{ch:9}.

For the local gauge-invariant field of Theorem~\ref{3.1:thm-main}, the smeared curvature density
$\mathcal{O}(f)$ (the action density $\operatorname{tr}F\cdot F$ smeared with a real test function;
for the non-abelian group the field strength itself is not gauge invariant), the applicable axiom
system is that of Osterwalder and Schrader \cite{OS73}, \cite{OS75} (see also \cite{GJ87}); the
theorem makes no assertion about Wilson loops. OS0--OS4 below follow the numbering E0--E4 of
\cite{OS73}. For $N=2$ the axioms stand as follows.

\begin{description}
\item[Regularity and growth (OS0).] This is clause (iii), treated in Chapters~\ref{ch:5}
and~\ref{ch:6}. For $n\ge 2$ and real Lipschitz $f_1,\dots,f_n$ with supports pairwise at distance
at least $\mathsf{d}>0$, the connected Schwinger functions are bounded by a constant times
$\prod_i\|f_i\|_\sharp$. This constant is free of the cutoff, the bare coupling and the position,
but not of the volume (the torus bound). For supports in balls of radius $R$, bare coupling in
$]0,g_-(g)]$ and tori above a volume floor depending on $n$, $R$ and $g$, it is free also of the
volume and non-increasing in $\mathsf{d}$ (the per-ball volume-free bound). With clause (ii), the
limits on $\mathbb{R}^4$ are therefore distributions of locally finite order on the configurations
of non-coinciding points that obey the per-ball bound. Nothing is asserted for $n=1$, at
coinciding points, or about decay in $\mathsf{d}$. Neither the growth condition in $n$ of
\cite{OS75} nor field operators are asserted. OS0 is proved as an implication. The torus bound is
the conclusion of the correlation theorem. The correlation theorem is proved as an implication
from inputs listed in Chapter~\ref{ch:5}. These include the weight lemma and the sixth-family
lemma (each stated; proof incomplete), one statement that is not proved, and sentences of
Ba{\l}aban (printed without proof; cited as printed). The per-ball bound is proved as an
implication from the correlation theorem and three further inputs: reflection positivity of the
lattice measure, a bound of Ba{\l}aban cited by statement, and uniformity of the last-scale
constants in the torus exponent (stated; proof incomplete).
\item[Translations and hypercubic covariance.] This is clause (ii), treated in Chapters~\ref{ch:6}
and~\ref{ch:7}. The torus limits are invariant under all translations (for $C^1$ test functions)
and under $W_4$, and the limits on $\mathbb{R}^4$ exactly under $\mathbb{R}^4\rtimes W_4$. This is
proved as an implication. For the first-passage torus limits, the antecedents are clause (iii) and
reflection positivity of the lattice measure. For the tuned family (with $g\le\gamma_U$), they are
the statements on which clause (v) depends. On $\mathbb{R}^4$, they are the per-ball volume-free
bound, reflection positivity of the lattice measure, clause (0) and the correlation theorem.
\item[Full Euclidean covariance (OS1).] This is clause (vi), treated in Chapter~\ref{ch:10}. It is
conditional on the infrared hypothesis displayed in clause (vi), the far-sphere flux hypothesis
$(\mathrm{IR}\text{-}\mathrm{fl})_{\omega_0}$: the flux of the lattice rotation current through
far spheres vanishes along one sequence of radii. This is not the clustering of OS4. A sufficient
condition is displayed beside it. It requires the connected correlation of one far insertion of
the lattice stress-tensor density with $n\ge2$ observables $\mathcal{O}(f_i)$ of disjoint supports
to be $o(\mathcal{R}^{-4})$ in the distance $\mathcal{R}$. This decay is taken in shell-averaged
form along one sequence of distances, uniformly in the cutoff (and in the volume along the joint
subsequence of clause (ii)). Under $(\mathrm{IR}\text{-}\mathrm{fl})_{\omega_0}$ every limit on
$\mathbb{R}^4$ is invariant under $\mathbb{R}^4\rtimes\mathrm{O}(4)$, and $\mathrm{O}(3)$ acts on
$\mathcal{H}^{\mathrm{curv}}$ by unitaries commuting with $H$. Clause (vi) is a complete deduction
from $(\mathrm{IR}\text{-}\mathrm{fl})_{\omega_0}$ and clauses (0), (i) and (iii) and the part of
clause (ii) on $\mathbb{R}^4$; the theorem does not assert that this hypothesis holds, and the
hypothesis is not proved.
\item[Reflection positivity (OS2).] Reflection positivity of the lattice measure is proved in
Chapter~\ref{ch:6}. For the limits it is part of clause (ii) and is proved as an implication from
the statements listed under translations. It holds for the centred field, in the form of
\cite{OS73}, on the linear span of products with pairwise separated supports in the open
half-space on one side of the hyperplane. On each torus it holds at every hyperplane $x_\mu=c$. On
$\mathbb{R}^4$ it holds at $x_0=0$ for time-ordered products with compact supports, and hence, by
the invariance under $\mathbb{R}^4\rtimes W_4$, at every coordinate hyperplane.
\item[Symmetry (OS3).] The limits are multilinear and symmetric. This is proved as an implication,
as part of clause (ii), from the same statements.
\item[Reconstruction.] This is clause (ii) on $\mathbb{R}^4$, treated in Chapter~\ref{ch:7}. It is
proved as an implication from the statements listed under translations for $\mathbb{R}^4$. The
construction starts from the linear span of centred monomials of the curvature fields with
separated supports; no state on the polynomial algebra is asserted. It gives the following
objects:
\begin{itemize}
\item a separable Hilbert space $\mathcal{H}^{\mathrm{curv}}$;
\item a unit vector $\Omega$, whose uniqueness is not asserted;
\item a positive self-adjoint operator $H$ with $H\Omega=0$, which generates a strongly continuous
contraction semigroup $e^{-tH}$;
\item strongly continuous unitary spatial translations fixing $\Omega$, whose generators commute
strongly with $H$, with joint spectrum in $[0,\infty[\times\mathbb{R}^3$;
\item a unitary action of $W_3$.
\end{itemize}
\item[Clustering and uniqueness of the vacuum (OS4).] Not asserted. No bound decays with the
separation, and no mass gap is asserted.
\end{description}

On each torus the limits exist without subsequences along the tuned family. For the first-passage
family on $\mathbb{T}_m$ with $m\ge m_{\mathbb{T}}(g_-(g))$, they exist along subsequences of
cutoffs. On $\mathbb{R}^4$ they exist along one joint subsequence of cutoffs and volumes, and
every statement above about $\mathbb{R}^4$ holds for each limit so obtained. Convergence on
$\mathbb{R}^4$ without subsequences is not asserted, and neither is uniqueness in the cutoff there.

Clause (iv) is stated for $N=2$ inside the regime of clause (iii) and is treated in
Chapter~\ref{ch:8}. Among its conditions, the two-sided separation window defining the class
$\mathfrak{T}(\mathsf{d};\vartheta)$ of test-function pairs is a condition on the test functions,
not a hypothesis. Under the conditions of clause (iv), the connected two-point function of two
smeared action densities is given by a product of four factors. These are $(N^2-1)/2$, an explicit
positive free-field kernel, the fourth power of the running coupling at scale $\mathsf{d}$, and the
square of a factor in $[2/3,4/3]$. The relative error is at most one half, uniformly in the
cutoff, the bare coupling, the volume and the position. Hence every limit of clause (ii), on a
torus or on $\mathbb{R}^4$, is non-trivial in the sense of Definition~\ref{3.1:def-non-trivial},
and $\dim\mathcal{H}^{\mathrm{curv}}\ge 2$ for the limits on $\mathbb{R}^4$. This is proved as an
implication. The deduction is complete at the standing of the correlation bounds of clause (iii).
It is also complete at the standing of the second-order extraction estimate of the Gaussian
extraction and of the continuum kernel (each stated; proof incomplete). Finally, through the
normaliser statement of the passage to all volumes, it is complete at the standing of the lemma of
clause (0) on independence from the choice of torus. That lemma is proved as an implication from
earlier lemmas of this book and Ba{\l}aban's definitions. The convergence of the expansions on the
joined localisation domains is cited from Ba{\l}aban as printed. Convergence or uniqueness of the
limits in the volume is not asserted.

\begin{corollary}[Short-distance uniqueness of the infinite-volume theory at given renormalised coupling]
\label{3.6:cor-short-distance-uniqueness}
Let $g$ be a renormalised coupling in the window of clause~(v) of Theorem~\ref{3.1:thm-main} with
\begin{equation*}
g\le\min\{\gamma_\natural,\gamma'_U\},
\end{equation*}
where $\gamma_\natural$ and $\gamma'_U$ are thresholds on the coupling. Let the source radius
$\varrho$ satisfy the one-sided floor $\varrho\ge 6(B_\sharp+1)$. Let $n\in\{2,3\}$, let the separation
$\mathsf{d}$ satisfy
\begin{equation*}
\mathsf{d}\le L^{\underline{s}'(g)}/40,
\end{equation*}
and let $(f_1,\dots,f_n)$ be an $n$-tuple of test functions admissible at scale $\mathsf{d}$ in the
sense of clause~(iii) of Theorem~\ref{3.1:thm-main}. Take any two continuum theories
of the construction at the coupling $g$, each being either the theory
on a torus $\mathbb{T}_m$ with
\begin{gather*}
m\ge m_\natural(g,n),\\
m_\natural(g,n)=j_U+\max\bigl\{m_{\mathbb{T}}(g_-(g)),\
\max\{m_\star,\lceil\log_L\ell_n\rceil\}+1+\bar s(g)\bigr\},
\end{gather*}
or an infinite-volume limit point along an arbitrary sequence of cutoffs and volumes,
the two limits taken in either order. Then the connected $n$-point functions of the smeared action
density of the two theories at
$(f_1,\dots,f_n)$ differ by at most
\begin{equation}\label{3.6:eq-short-distance-uniqueness-1}
2\,C^{\star}_{T,n}\,L^{-(4-n)\underline{s}'(g)}\,\mathsf{d}^{\,4-n}\prod_{i=1}^{n}\|f_i\|_\infty .
\end{equation}
Here $C^{\star}_{T,n}$ depends on nothing but $(L,\mathfrak{p},\gamma,\gamma_H,\varrho,\vartheta)$;
$m_\natural(g,n)$ is the floor on the torus exponent displayed above, in which $m_{\mathbb{T}}$ and
$g_-$ are as in the glossary, Definition~\ref{0.2:not-glossary}, $\ell_n$ is the least
power of $L$ with $\ell_n\ge 8n(R+1)$, $m_\star$ and $\bar s(g)$ are determined by $g$, and $j_U$ is
an integer depending on nothing, the same integer as in \eqref{3.6:eq-short-distance-uniqueness-2}
below; and the exponent
$\underline{s}'(g)$ satisfies the one-sided bound
\begin{equation}\label{3.6:eq-short-distance-uniqueness-2}
\underline{s}'(g)\ \ge\ \frac{g^{-2}-\gamma_\star^{-2}}{B_\sharp}-(j_U+1),
\end{equation}
with $j_U$ as above and
\begin{equation*}
\gamma_\star=\min\{\gamma_W^{\mathbb{T}},\gamma_U\},
\end{equation*}
where $\gamma_U$ is the threshold on the coupling of the uniqueness clause~(v) of
Theorem~\ref{3.1:thm-main} and $\gamma_W^{\mathbb{T}}$ is a threshold on the coupling; the threshold
$\gamma'_U$ of the hypothesis on $g$ is a distinct letter. In particular, all
infinite-volume limit points at the coupling $g$ agree, to the accuracy
\eqref{3.6:eq-short-distance-uniqueness-1}, with the
unique continuum theory on the reference torus $\mathbb{T}_{m_\natural(g,n)}$ (the torus whose
exponent is the floor displayed above); this theory depends real-analytically on the coupling.
On the pairs of test functions of clause~(iv) of
Theorem~\ref{3.1:thm-main}, at the reference scale, that is at separations $\mathsf{d}\le 1/40$ in
reference units, the relative discrepancy of the two connected two-point functions is at most
\begin{equation*}
A\,g^{-4}\exp\Bigl[-\frac{2\ln L}{B_\sharp}\,g^{-2}\Bigr],
\end{equation*}
with $A$ a constant depending on nothing.

In the first-passage labelling of clause~(ii-$\mathbb{R}^4$) of Theorem~\ref{3.1:thm-main} the same
holds with
\begin{equation*}
\underline{s}(g)=\Bigl\lfloor \frac{g^{-2}-\gamma_\star^{-2}-B_\sharp}{B_\sharp}\Bigr\rfloor
\end{equation*}
in place of $\underline{s}'(g)$, each limit point lying within
\begin{equation*}
C^{\star}_{T,n}\,L^{-(4-n)\underline{s}(g)}\,\mathsf{d}^{\,4-n}\prod_{i=1}^{n}\|f_i\|_\infty
\end{equation*}
of the reference-torus theory at some coupling of the window.
\end{corollary}

Proved as an implication --- complete deduction at the standing of the correlation bounds (iii) and of
the uniqueness clause (v) of Theorem~\ref{3.1:thm-main}.

In words, writing $\ell_\star=L^{\underline{s}'(g)}$ for the length, in reference units, below one
fortieth of which the separations of the corollary lie, so that the bound
\eqref{3.6:eq-short-distance-uniqueness-1} equals
\begin{equation*}
2\,C^{\star}_{T,n}\,(\mathsf{d}/\ell_\star)^{4-n}\prod_{i=1}^{n}\|f_i\|_\infty ,
\end{equation*}
all infinite-volume limit points at a given
renormalised coupling --- along any sequences of cutoffs and volumes, in any order --- have the same
connected two- and three-point functions of the smeared action density at separations
$\mathsf{d}\le\ell_\star/40$, to accuracy a constant depending only on
$(L,\mathfrak{p},\gamma,\gamma_H,\varrho,\vartheta)$ times $(\mathsf{d}/\ell_\star)^2$
(respectively $\mathsf{d}/\ell_\star$) times the
test-function norms, and they are equal, to that accuracy, to those of the unique continuum theory on one
fixed reference torus, which depends real-analytically on the coupling; in reference units the
discrepancy is exponentially small in the inverse coupling squared --- smaller than every power of the
running coupling. The passage to infinite volume
uses, of the large volume, only reflection positivity of the bare Wilson measure and the lower stability
bound.

The corollary concerns $n=2$ and $n=3$ only; no statement is made for $n\ge 4$.
The following are not asserted:
\begin{itemize}
\item uniqueness of the infinite-volume limit;
\item uniqueness of the vacuum;
\item clustering of the infinite-volume limit points;
\item continuity, analyticity or injectivity of the map that assigns to a coupling the corresponding
theory on $\mathbb{R}^4$.
\end{itemize}
None of these is implied by the statements proved in this book. What is asserted is the common
short-distance structure stated above.

\begin{remark}[The block-law statements]\label{3.2:rem-block-laws}
The two statements below concern Ba{\l}aban's block laws and the limit states on the block
algebra. They are not clauses of Theorem~\ref{3.1:thm-main}, and no clause of
Theorem~\ref{3.1:thm-main} uses them. Clause~(0) of Theorem~\ref{3.1:thm-main} enters
the deduction of~(b) and the consequences drawn from~(a) below. Conversely, no clause of
Theorem~\ref{3.1:thm-main} uses any statement other than clause~(0) from which (a) or~(b) is
deduced below. Both statements are made under the hypotheses (H1)--(H6) and in the notation of
Theorem~\ref{3.1:thm-main}, on a fixed torus $\mathbb{T}_m$ and at a fixed reference coupling
$g$. Statement~(b) is made moreover for $N=2$ and under (H7).

The \emph{block law at depth $s$} is the positive measure
$\nu^B_{K,s}=\rho_{K-s}\,dV_{K-s}/\int\rho_0\,dU$ on the block fields at level $K-s$. In words, it
is the measure on the block fields at level $K-s$ whose density with respect to $dV_{K-s}$ is the
effective density $\rho_{K-s}$ of Definition~\ref{1.5:def-densities}, divided by the total mass of
the bare density. Further, $\mathfrak{A}$ denotes the block algebra of
Definition~\ref{1.5:def-block-average}, $\langle\cdot\rangle_K$ is the expectation in the lattice
measure of Definition~\ref{1.1:def-wilson-action} on $\mathbb{T}_m$ at bare spacing $L^{-K}$,
and $p_0(\cdot)$ is the degree function of Definition~\ref{1.5:def-domains}.
\begin{itemize}
\item[(a)] \emph{Block-law comparison.} There are constants $C,c>0$, not depending on $K$, $s$ or
the position (independence of $m$ is not asserted), and an integer $k_*\ge1$, depending only on
$L$ and the auxiliary parameters. For
every $K$, every depth $s$ with $s_*\le s\le K-k_*$, and every $C^2$ gauge-invariant block-local
function $F$ of depth $s$, they satisfy
\[
\bigl|\langle F\rangle_K-\nu^B_{K,s}(F)\bigr|\le C\,e^{-c\,p_0(g_{K-s})}
\bigl(\|DF\|+\|D^2F\|\bigr).
\]
Here $\|DF\|$ is the sum, over the block bonds on whose variables $F$ depends, of the suprema of
the first
derivatives of $F$ in the block variables, and $\|D^2F\|$ is the corresponding sum over pairs of
such bonds of the second derivatives. Nothing is claimed at the $k_*$ depths nearest the cutoff.
This inequality would be combined with clause~(0) and with bounds on the block laws themselves,
which are not part of Theorem~\ref{3.1:thm-main}. The combination would give two-sided bounds on
the expectation of a depth-$s$ plaquette variable, that is, of
$1-\operatorname{Re}\operatorname{tr}V_{K-s}(\partial p)$ for a plaquette $p$ of the block lattice
at depth $s$. These bounds would in turn show the following for every limit state on
$\mathfrak{A}$, that is, every weak-$*$ limit point, on the fixed torus and at fixed $g$, of the
restrictions of the states $\langle\cdot\rangle_K$ to $\mathfrak{A}$ along $g_0\to0$ with
$K=K(g_0,g)$: for every depth $s$ not smaller than $s_*$ and than a floor depending on $g$, the
law of a depth-$s$ plaquette variable under such a state is not the point mass at zero, so that
the state is not that of the trivial connection. This
property concerns states on the block algebra; it is not the non-triviality of
Definition~\ref{3.1:def-non-trivial}.

Statement~(a) is not proved. Its display is the sentence at one depth of a statement called
here the comparison lemma, which compares $\langle\cdot\rangle_K$ with Ba{\l}aban's block laws
at all depths $s_*\le s\le K-k_*$ at once; none of the constituent statements of that lemma is
proved in this book.

\item[(b)] \emph{Block-law uniqueness in the cutoff.} Let $g_0^{(i)}\downarrow0$, with
$K_i=K(g_0^{(i)},g)$, be a calibrated family. By this we mean a sequence along which, for every
depth $s$, the running couplings $g_{K_i-s}$ converge to a limit $g_{\infty,s}$ and the block laws
$\nu^B_{K_i,s}$ converge weakly to a limit $\nu^\infty_s$. Then the states
$\langle\cdot\rangle_{K_i}$ converge on the whole of $\mathfrak{A}$, without extraction of
subsequences, to a state $\omega$. This state satisfies $\omega(F)=\lim_{s}\nu^\infty_s(F)$ for
every $F\in\mathfrak{A}$, and
\[
\bigl|\omega(F)-\nu^\infty_s(F)\bigr|\le C\,e^{-c\,p_0(g_{\infty,s})}
\bigl(\|DF\|+\|D^2F\|\bigr)
\]
for every depth $s$ and every $C^2$ gauge-invariant block-local function $F$ of depth $s$, with
$C$ and $c$ as in~(a). Moreover, for every fixed $C^2$ gauge-invariant block-local function $H$
there is a constant $C_H$ with $|\omega(H)-\nu^\infty_s(H)|\le C_H\,e^{-c\,p_0(g_{\infty,s})}$
at every depth $s$ not smaller than that of $H$.

Statement~(b) itself is not proved: it is proved as an implication from the following three
statements, the first two of which are not proved:
\begin{itemize}
\item the comparison lemma of~(a) read at every depth finer than that of $F$;
\item one further statement on the convergence of Ba{\l}aban's block laws at each fixed depth;
\item clause~(0).
\end{itemize}
No converse is asserted.
\end{itemize}

From the block side, the proofs of Theorem~\ref{3.1:thm-main} use in particular the
exact-insertion identity of clause~(i), the depth $s_*$ of clause~(iv), and the block fields of
Definition~\ref{1.5:def-block-average} as their vehicle; they use no statement about the values
of the block laws.
\end{remark}

In addition to what each clause lists as not asserted, Theorem~\ref{3.1:thm-main} does not
assert any of the following.
\begin{itemize}
\item Invariance under $\mathrm{O}(4)$ beyond the hypercubic group $W_4$, except in clause (vi).
That clause deduces it from its hypothesis $(\mathrm{IR}\text{-}\mathrm{fl})_{\omega_0}$, the
far-sphere flux hypothesis: the flux of the lattice rotation current through far spheres
vanishes along one sequence of radii. Clause (vi) displays a sufficient condition for it: the
connected correlation of one far insertion of the lattice stress-tensor density with $n \ge 2$
curvature observables of pairwise disjoint supports decays faster than $|x|^{-4}$, uniformly in
the cutoff. Nothing is asserted about the truth of $(\mathrm{IR}\text{-}\mathrm{fl})_{\omega_0}$,
and it is assumed in no other clause.
\item Convergence on $\mathbb{R}^4$ without passage to subsequences in the volume. Uniqueness of
the infinite-volume limit. Uniqueness of the vacuum $\Omega$. Clustering.
\item A mass gap. Temperedness. Field operators or Wightman functions. The spectrum condition
beyond the non-negativity $H \ge 0$ of the Hamiltonian.
\item Anything about clauses (ii)--(vi) for $\mathrm{SU}(N)$ with $N \ge 3$. Clauses (0) and (i)
are stated for every $N \ge 2$.
\item That the parameter $\hat w$ of clause (v) is determined by the limit, that is, injectivity
of $\hat w \mapsto S^{(\hat w)}$. Uniformity of clause (v) in the torus size. Uniqueness in the
cutoff of the limits on $\mathbb{R}^4$.
\item Anything about one-point functions, coinciding supports, or components of the curvature
other than the action density of Definition~\ref{1.3:def-curvature-field}. Anything about
block-averaged observables, Wilson loops, or the state on the algebra of block-averaged
observables. The block-law statements are the subject of Remark~\ref{3.2:rem-block-laws}; they
are not proved, and no clause of Theorem~\ref{3.1:thm-main} uses them.
\item Non-triviality in any sense other than that of Definition~\ref{3.1:def-non-trivial} and
clause (iv), which concern the connected two-point function of the smeared action density only.
For the limits on $\mathbb{R}^4$, Theorem~\ref{3.1:thm-main} asserts existence along
subsequences of volumes, invariance under translations and under $W_4$, reflection positivity,
and the reconstruction of a Hilbert space $\mathcal{H}^{\mathrm{curv}}$ with vacuum $\Omega$ and
Hamiltonian $H \ge 0$. By clause (iv) it asserts also, for reference couplings $g$ below the
ceiling of that clause, that these limits are non-trivial and that $\mathcal{H}^{\mathrm{curv}}$
has dimension at least two; clause (iv) is proved as an implication from the statements used in
the proof of clause (iii) and, through the passage from one torus to all tori, from the
torus-independence lemma of clause (0). That lemma is proved as an implication from earlier
lemmas of this book and Ba{\l}aban's definitions, together with the convergence of the expansions
on the joined localization domains, which is printed without proof in \cite[(1.70)]{B-CMP122a}
and \cite[(1.68)--(1.69)]{B-CMP122b} and cited as printed.
\item A sign of the coupling increment $x_k - x_{k+1}$ at a single step. The statements of
\cite[Theorem~2, (0.31)]{B-CMP109}, which no clause uses.
\end{itemize}

Part~II proves Theorem~\ref{3.1:thm-main} clause by clause, in the order of the clauses (0), (i),
(iii), (ii), (iv), (v), (vi). Chapter~\ref{ch:4} treats the coupling flow and ultraviolet stability,
clauses (0) and (i), by one induction on the level; the one hypothesis of clause (i) is supplied by
clause (0). Chapter~\ref{ch:5} treats the correlation theorem, which carries clauses (iii-a) and
(iii-c), and Chapter~\ref{ch:6} the per-ball volume-free bound, clause (iii-b).
Chapter~\ref{ch:7} treats clause (ii), covering limits, invariance, reflection positivity and
reconstruction. Clause (ii) consumes clause (iii), which is why (iii) precedes (ii) in the theorem.
Chapter~\ref{ch:8} treats clause (iv), the two-point witness, which is stated inside the regime of
clause (iii) and uniformly in the volume. Also stated there in place, with its proof outstanding,
is an estimate that no clause's proof uses. It bounds the contribution of large-field regions
created more than the prescribed number of renormalization steps before the depth at which the
two-point function is read and whose large-field factor has not by then been absorbed by a later
step back into the small-field format of the representation of clause (i-a); this contribution is
the never-healed old-age tail.
Chapter~\ref{ch:9} treats clause (v), uniqueness up to the renormalized coupling.
Chapter~\ref{ch:10} treats clause (vi), $\mathrm{O}(4)$ invariance, as a deduction from its
displayed hypothesis $(\mathrm{IR}\text{-}\mathrm{fl})_{\omega_0}$, the far-sphere flux
hypothesis: the flux of the lattice rotation current through far spheres vanishes along one
sequence of radii; beside it is displayed a sufficient infrared decay condition: one far-away
stress-tensor insertion decorrelates from the curvature observables faster than the inverse
fourth power of the distance, uniformly in the cutoff. Chapter~\ref{ch:11} assembles
the clauses; the dependency among them is drawn once there and shown to have no cycle. Part~III
consists of Chapter~\ref{ch:13}, the index of every statement of Part~II whose proof is
outstanding or incomplete, or which is printed without proof.

Each numbered statement of Part~II declares its standing where it is stated, in one of the forms
fixed in Chapter~\ref{ch:0}: proved in full; proved as an implication from named statements
listed among its hypotheses; a published theorem, cited by statement with its
exact reference; a sentence Ba{\l}aban printed without proof, quoted at its use with its
reference and marked ``printed without proof; cited as printed''; stated, with its proof
incomplete at a step marked $(\ast)$; stated, with its proof outstanding and the route of a proof
set out in its place; or a model statement, proved in the model that its first sentence
describes. Wherever an argument appeals to a statement that is not proved in full, the appeal
says so at the place where it is made.

Two reading conventions hold throughout. Constants are symbolic, and the quantities on which each
constant may depend are stated at each statement; no existentially bound constant is given a
numerical value; in particular the blocking factor $L$ is subject to one existential threshold
$L \ge L_0 = L_0(N)$, bound once and never evaluated, and every further condition on $L$ is a
lower bound. Ba{\l}aban's displays are used as stated
in the cited papers and read through his definitions; where a reading of a display has to be
declared, it is declared at the place where the display is used, and the few misprints that are
named all occur in displays used in Chapter~\ref{ch:4}.

For $\mathrm{SU}(N)$ with $N \ge 3$, the statements on which the passage of clauses (i), (iii),
(iv) and (v) to that group rests are stated in place in Chapters~\ref{ch:4}, \ref{ch:5}
and~\ref{ch:9}, each with its standing declared where it is stated; some of them have proofs that
are outstanding or incomplete.

%% ---- end inlined file: chapters/c03 ----

\part{The proof, clause by clause}
\clearpage
\setcounter{section}{3}
\refstepcounter{section}
\section*{Chapter~\thesection\quad The coupling flow and ultraviolet stability: clauses (0) and (i)\addcontentsline{toc}{section}{\protect\numberline{\thesection}The coupling flow and ultraviolet stability: clauses (0) and (i)}\label{ch:4}}
%% ---- begin inlined file: chapters/c04 ----
\begin{theorem}[Clause (0) of Theorem~0: the coupling flow]\label{4.1:thm-clause-zero}
In the setting and under the hypotheses of Theorem~0, let $K=K(g_0,g)$ be the first-passage
index of Definition~\ref{1.2:def-flow-constants}. Let $x_0=g_0^{-2}$, let $x_k$ be the running
inverse couplings of Definition~\ref{1.2:def-couplings}, whose increments $\beta_{k+1}$ are the
coupling increments constructed in Definition~\ref{4.2:def-increment}, and put
$g_k=x_k^{-1/2}$. Then the following hold.
\begin{enumerate}
\item $K$ is a well-defined non-negative integer, and
\begin{equation}\label{4.1:eq-clause-zero-1}
\frac{g_0^{-2}-g^{-2}}{B_\sharp}-1\;\le\;K\;\le\;
1+\frac{g_0^{-2}-g^{-2}-B_\sharp+c_5}{\lambda_\sharp}.
\end{equation}
In particular $K\to\infty$ as $g_0\downarrow 0$.
\item For every $k\le K$ one has $g_k\in(0,g)\subset(0,\gamma_H]$.
\item $g_K\in[g_-(g),g)$.
\item The running inverse couplings satisfy the two-sided flow bound
\begin{equation}\label{4.1:eq-clause-zero-2}
\lambda_\sharp(j-i)-c_5\;\le\;x_i-x_j\;\le\;B_\sharp(j-i),
\qquad 0\le i\le j\le K.
\end{equation}
\item The $K$ renormalisation steps on $\mathbb{T}_m$ produce exactly these couplings
$g_0,\dots,g_K$.
\item (Corollary: fixed-depth windows.) Let $(g_0^{(i)})_{i\ge 1}$ be a sequence of bare
couplings with $g_0^{(i)}\downarrow 0$, put $K_i:=K(g_0^{(i)},g)$, and for each $i$ let $x_k$
and $g_k$ denote the running couplings started at $x_0=(g_0^{(i)})^{-2}$ (the dependence on $i$
is suppressed in the notation). Then for every integer $s\ge 0$ and every
$i$ with $K_i\ge s$,
\begin{equation}\label{4.1:eq-clause-zero-3}
x_{K_i-s}\in\bigl(g^{-2}+\max\{0,\lambda_\sharp s-c_5\},\;g^{-2}+B_\sharp(s+1)\bigr];
\end{equation}
and, along a subsequence, $g_{K_i-s}\to g_{\infty,s}$, where the limit satisfies
\begin{equation}\label{4.1:eq-clause-zero-4}
g_{\infty,s}^{-2}\in\bigl[g^{-2}+\lambda_\sharp s-c_5,\;g^{-2}+B_\sharp(s+1)\bigr].
\end{equation}
\end{enumerate}
Not asserted: a sign of $x_k-x_{k+1}$ at each single step; monotonicity of a tuned $g_0$;
convergence of $g_{K-s}$ along first passage; regularity of $\beta_j$ as a function of the
couplings.
\end{theorem}

This is clause (0) of Theorem~0, restated as the object of this chapter. The sections preceding
Proposition~\ref{4.19:prop-clause-zero} supply the inputs of a single induction on the level $k$,
shared with clause (i) of Theorem~0, and Proposition~\ref{4.19:prop-clause-zero}, together with
Lemma~\ref{4.19:lem-flow-bounds} and Corollary~\ref{4.19:cor-windows}, concludes clause (0).

Clauses (0) and (i) of Theorem~0 are restated in this chapter as Theorem~\ref{4.1:thm-clause-zero} and
as Theorems~\ref{4.20:thm-representation} and~\ref{4.33:thm-stability}. They are not proved one after
the other but are the two halves of a single induction on the level $k$. At stage $k$ the increments
$\beta_l$, $l\le k$, each a function of $g_0,\dots,g_{l-1}$, are known with their per-step bounds, as
are the two-sided flow bounds
\[
\lambda_\sharp(j-i)-c_5\le x_i-x_j\le B_\sharp(j-i),\qquad 0\le i\le j\le k,
\]
where $x_j=g_j^{-2}$, and the representation of the effective densities $\rho_j$, $j\le k$, with the
bounds of clause (i) at every level $j\le k$.

Four statements are specific to clause (0): the independence of the localised terms from the
small-field region, the operator bounds at free current, the representation and bound of the
large-field terms, and the large-field part of the coupling increment. Clause (0) depends in
addition on further inputs, some of them taken from Ba{\l}aban's papers. One step $k\to k+1$
proceeds in the following order.
\begin{enumerate}
\item The fluctuation integral of step $k+1$, taken at coupling $g_k\in(0,\gamma]$, produces the
small-field effective action on the whole lattice. The number $\beta_{k+1}$ is its $\mathbb{Z}^4$
coefficient (Definition~\ref{4.2:def-increment}), read off the Gaussian normalizer and the term of the
fluctuation integral. Both terms contain only $g_k$, the cutoff function and the propagators of step
$k$, and the effective action at level $k$, into which the large-field remainders of earlier steps
have been merged. Neither contains an object of index $k+1$, the number of steps $K$, or the torus.
Hence $\beta_{k+1}$ is a function of $(g_0,\dots,g_k)$ alone, and the induction is well founded.
\item The per-step bounds at level $k+1$ are the finite part of the one-loop law, the bounds on the
unmarked part of the increment (Lemmas~\ref{4.2:lem-one-loop} and~\ref{4.2:lem-unmarked}), and the
bound on the large-field part of the coupling increment, that is, on the contribution to
$\beta_{k+1}$ of the terms made by the operation $\mathbf{R}$. These use clause (i) and the flow
bounds at levels at most $k$, and the containment $x_0,\dots,x_k\ge\gamma^{-2}$.
\item The recursion sets $x_{k+1}=x_k-\beta_{k+1}$ and, when $x_{k+1}>0$, $g_{k+1}=x_{k+1}^{-1/2}$.
The two-sided flow bounds then hold at indices at most $k+1$. Then the cube size of step $k+1$ is
fixed, and then the thresholds of the decompositions of unity with which step $k+1$ opens.
\item Suppose moreover that $x_{k+1}\ge\gamma^{-2}$, so that $g_{k+1}\in(0,\gamma]$. Then the
operations $\mathbf{T}$ and $\mathbf{R}$ of step $k+1$ are performed. The representation and bounds
of $\rho_{k+1}$ that they yield are clause (i) at level $k+1$.
\end{enumerate}

Along the first-passage trajectory one has
\[
x_k>g^{-2}\ge\gamma_H^{-2}\ge\gamma^{-2},\qquad 0\le k\le K.
\]
So the containment required in the first, second and last items above holds at every step; it is
the one hypothesis on the coupling in Ba{\l}aban's one-step renormalization theorem.
No upper bound on $K$ and no continuity of the increments in
$g_0$ is used. The constants of the induction depend neither on the torus exponent $m$ nor on $K$.
The torus $\mathbb{T}_m$ enters only through a lower bound on $m$ and through extensive quantities
such as the site counts $N_k$. The sequence $(g_j)_j$ is therefore defined before any torus is
declared, and the first passage $K(g_0,g)$ afterwards.

The next section gives the first-passage arithmetic and defines the increment. The section after it
cites Ba{\l}aban's theorems by statement. The inputs of clause (0) then follow, each in its own
section and in the form it has: proved, cited from a published paper, printed without proof and
cited as printed, or with its proof outstanding. A section on the two-sided flow bounds then
concludes clause (0). Clause (i) comes next. Its first part is the inductive representation of the
effective densities, whose proof is an induction occupying several sections. Its second part is the
integrated stability bound and the equality of integrals. The chapter closes with exact insertion,
gauge fixing, and the statements needed for $\mathrm{SU}(N)$ with $N\ge3$.

\begin{lemma}[First-passage arithmetic]\label{4.2:lem-first-passage}
Let $\lambda_\sharp>0$, $B_\sharp>0$, $c_5\ge0$, $g>0$ and $t\ge g^{-2}+B_\sharp$. Let
$(x_k)_{0\le k<n_{\max}}$, with $n_{\max}\in\{1,2,\dots\}\cup\{\infty\}$, be a real sequence with
$x_0=t$ and the following property. For every integer $n\ge1$ such that $x_0,\dots,x_{n-1}$ are
defined and all exceed $g^{-2}$, the number $x_n$ is defined and
\begin{equation}\label{4.2:eq-first-passage-1}
\lambda_\sharp(j-i)-c_5\;\le\;x_i-x_j\;\le\;B_\sharp(j-i),\qquad 0\le i\le j\le n .
\end{equation}
Then
\begin{equation*}
K:=\min\{k:\ x_k\text{ is defined and }x_k\le g^{-2}+B_\sharp\}
\end{equation*}
is a well-defined non-negative integer, and the following hold.
\begin{itemize}
\item[(a)] $x_k>g^{-2}+B_\sharp$ for $k<K$, and $g^{-2}<x_K\le g^{-2}+B_\sharp$. Equivalently, with
$g_k:=x_k^{-1/2}$ and $g_-(g):=(g^{-2}+B_\sharp)^{-1/2}$ (for the constant $B_\sharp$ of
Definition~\ref{1.2:def-flow-constants} this is the function $g_-$ defined there), one has
$g_k\in(0,g_-(g))$ for $k<K$ and $g_K\in[g_-(g),g)$. In particular $g_k\in(0,g)$ for all $k\le K$.
\item[(b)] \eqref{4.2:eq-first-passage-1} holds for $0\le i\le j\le K$.
\item[(c)] The number of steps satisfies
\begin{equation}\label{4.2:eq-first-passage-2}
\frac{t-g^{-2}}{B_\sharp}-1\;\le\;K\;\le\;1+\frac{t-g^{-2}-B_\sharp+c_5}{\lambda_\sharp}.
\end{equation}
In particular, for any family of such sequences indexed by $t$, $K\to\infty$ as $t\to\infty$.
\end{itemize}
\end{lemma}

\begin{proof}
We first show that $K$ exists. Since $\lambda_\sharp>0$, there is an integer $n_0\ge1$ with
\begin{equation*}
n_0>\frac{t-g^{-2}-B_\sharp+c_5}{\lambda_\sharp}.
\end{equation*}
Suppose that no defined $x_k$ with $k\le n_0$ satisfies $x_k\le g^{-2}+B_\sharp$.

We show by induction on $n\le n_0$ that $x_0,\dots,x_n$ are defined and exceed
$g^{-2}+B_\sharp$. For $n=0$, the number $x_0=t$ is defined, and it exceeds $g^{-2}+B_\sharp$ by the
assumption just made. Let $1\le n\le n_0$, and suppose that $x_0,\dots,x_{n-1}$ are defined and exceed
$g^{-2}+B_\sharp$. Since $B_\sharp>0$, they exceed $g^{-2}$, so the hypothesis of the lemma, applied
with this $n$, makes $x_n$ defined; as $n\le n_0$, the assumption gives $x_n>g^{-2}+B_\sharp$.

In particular $x_0,\dots,x_{n_0-1}$ are defined and exceed $g^{-2}$, so
\eqref{4.2:eq-first-passage-1} holds with $n=n_0$. Its left inequality with $(i,j)=(0,n_0)$ reads
$\lambda_\sharp n_0-c_5\le t-x_{n_0}$, and by the choice of $n_0$ we have
$\lambda_\sharp n_0>t-g^{-2}-B_\sharp+c_5$; hence
\begin{equation*}
x_{n_0}\le t-\lambda_\sharp n_0+c_5<g^{-2}+B_\sharp ,
\end{equation*}
which contradicts $x_{n_0}>g^{-2}+B_\sharp$. Hence the minimum defining $K$ is attained and
$K\le n_0$. Since the set of indices at which the sequence is defined is the initial segment
$\{k:0\le k<n_{\max}\}$ and contains $K$, the numbers $x_0,\dots,x_K$ are all defined.

For (a), minimality of $K$ gives $x_k>g^{-2}+B_\sharp$ for $k<K$. For $x_K$ we distinguish two cases.
\begin{itemize}
\item If $K=0$, then $g^{-2}+B_\sharp\le t=x_0\le g^{-2}+B_\sharp$, by the hypothesis on $t$ and
the definition of $K$; so $x_K=g^{-2}+B_\sharp>g^{-2}$.
\item If $K\ge1$, then $x_0,\dots,x_{K-1}$ exceed $g^{-2}+B_\sharp>g^{-2}$, so
\eqref{4.2:eq-first-passage-1} holds with $n=K$; its right inequality with $(i,j)=(K-1,K)$ gives
$x_K\ge x_{K-1}-B_\sharp>g^{-2}$.
\end{itemize}
So all $x_k$ with $k\le K$ exceed $g^{-2}>0$, and $g_k=x_k^{-1/2}$ is a well-defined positive number
for $k\le K$. The map $x\mapsto x^{-1/2}$ is strictly decreasing on $(0,\infty)$ and sends
$g^{-2}+B_\sharp$ to $g_-(g)$ and $g^{-2}$ to $g$. Therefore $x_k>g^{-2}+B_\sharp$ is the statement
$g_k<g_-(g)$, and $g^{-2}<x_K\le g^{-2}+B_\sharp$ is the statement $g_K\in[g_-(g),g)$. Since
$g_-(g)<g$, it follows that $g_k\in(0,g)$ for all $k\le K$.

For (b), if $K\ge1$, the numbers $x_0,\dots,x_{K-1}$ exceed $g^{-2}$ by (a), and the assertion is
\eqref{4.2:eq-first-passage-1} with $n=K$. If $K=0$, the only pair is $i=j=0$, and the assertion
reads $-c_5\le0\le0$, which holds because $c_5\ge0$.

For (c), the lower bound follows from the right inequality of (b) with $(i,j)=(0,K)$ together
with $x_K\le g^{-2}+B_\sharp$:
\begin{equation*}
t-g^{-2}-B_\sharp\le t-x_K\le B_\sharp K;
\end{equation*}
dividing by $B_\sharp>0$ gives the left inequality of \eqref{4.2:eq-first-passage-2}. For the upper
bound, suppose first that $K\ge1$. Then (a) gives $x_{K-1}>g^{-2}+B_\sharp$, and the left inequality
of (b) with $(i,j)=(0,K-1)$ gives $\lambda_\sharp(K-1)-c_5\le t-x_{K-1}$; together,
\begin{equation*}
g^{-2}+B_\sharp<x_{K-1}\le t-\lambda_\sharp(K-1)+c_5 .
\end{equation*}
This rearranges to $\lambda_\sharp(K-1)<t-g^{-2}-B_\sharp+c_5$, and dividing by $\lambda_\sharp>0$
gives the right inequality of \eqref{4.2:eq-first-passage-2}. If $K=0$, the right side of
\eqref{4.2:eq-first-passage-2} is at least $1$, because $t-g^{-2}-B_\sharp\ge0$, $c_5\ge0$ and
$\lambda_\sharp>0$; so it is at least $K$. Finally, the lower bound in
\eqref{4.2:eq-first-passage-2} tends to $\infty$ as $t\to\infty$, and hence so does $K$.
\end{proof}

Fix $j\ge 1$ and work in scale-$j$ units; we first fix the lattices, domains and tori at this scale.
The block field lives on the bonds of the unit lattice
$\mathbb{Z}^4$, and the minimal configurations live on the fine lattice $L^{-j}\mathbb{Z}^4$. We write
$\mathcal{D}_j^\infty$ for Ba{\l}aban's localisation domains of the infinite lattice, that is,
connected unions of $M$-cubes of the partition \cite{B-CMP109}.

For an even multiple $\varsigma$ of $M$ we write $\mathbb{T}^{(j)}_\varsigma$ for the torus of side
$\varsigma$ (the letter $s$ being reserved for the depth below the cutoff), with covering projection
$p_\varsigma$ and localisation domains $\mathcal{D}_j(\mathbb{T}^{(j)}_\varsigma)$. For the torus
$\mathbb{T}_m$ at cutoff $K$ we have $\varsigma=2L^{m+K-j}$.

A domain $X\in\mathcal{D}_j(\mathbb{T}^{(j)}_\varsigma)$ is \emph{non-wrapping} if it lies in a
closed cube of side at most $\varsigma/4$; otherwise $X$ \emph{wraps}. Suppose $X$ lies in such a
cube $Q$. Then $p_\varsigma$ maps each cube of $\mathbb{R}^4$ lying over $Q$ injectively onto $Q$.
Moreover, $p_\varsigma^{-1}(X)$ is the disjoint union of the translates by $\varsigma\mathbb{Z}^4$ of
one set $\hat X$, a \emph{lift} of $X$, on which $p_\varsigma$ is injective. We write $d_j(X)$ for
Ba{\l}aban's tree length of $X$ in units of $M$ \cite{B-CMP109}.

We use the following properties of $d_j$ and $p_\varsigma$; here $|\cdot|$ denotes the sup norm.
\begin{itemize}
\item[(g1)] There are constants $\kappa(d)$ and $K_0$ (Ba{\l}aban's constants in
\cite[(0.26)]{B-CMP109}) such that
\[
\sum_{X\ni z}e^{-\kappa(d)\,d_j(X)}\le K_0 ,
\]
both on $\mathbb{Z}^4$ and on every $\mathbb{T}^{(j)}_\varsigma$.
\item[(g2)] Every $\hat X\in\mathcal{D}_j^\infty$ lies in a cube of side $M(d_j(\hat X)+2)$. This
follows from the definition of $d_j$ in \cite{B-CMP109} as a tree length in units of $M$; the two
additional cube sides account for the cubes at the two ends of a path in the tree. On every
$\mathbb{T}^{(j)}_\varsigma$, every $X$ lies in the image under $p_\varsigma$ of a cube of side
$M(d_j(X)+2)$, so its diameter in the torus distance is at most $M(d_j(X)+2)$.
\item[(g3)] A wrapping $X$ has $d_j(X)>\varsigma/(4M)-2$. Indeed, by (g2) the domain $X$ lies in
the image under $p_\varsigma$ of a cube of side $M(d_j(X)+2)$; if this side were at most
$\varsigma/4$, that image would be a closed cube of side at most $\varsigma/4$ containing $X$, and
$X$ would be non-wrapping. A non-wrapping $X\ni z$ has exactly one lift $\hat X\ni\hat z$, for any
fixed lift $\hat z$ of $z$, and $d_j(\hat X)=d_j(X)$. Conversely, suppose $\hat X\ni\hat z$ and
$d_j(\hat X)<\varsigma/(4M)-2$. Then $M(d_j(\hat X)+2)<\varsigma/4$, so by (g2) the set $\hat X$
lies in a box of side less than $\varsigma/4$; hence $p_\varsigma$ is injective on $\hat X$ and
$p_\varsigma\hat X$ is non-wrapping.
\end{itemize}

\begin{definition}[The coupling increment as the whole-lattice second-variation coefficient of the
complete regular part of one step]\label{4.2:def-increment}
Suppose steps $1,\dots,j-1$ have been performed, so that $g_0,\dots,g_{j-1}$ are defined
(Definition~\ref{1.2:def-couplings}). The \emph{complete regular part} produced by step $j$ is the
sum of two pieces.
\begin{itemize}
\item[(a)] The Gaussian normaliser
\[
\log Z^{(j-1)}(U_j)-\log Z^{(j-1)}(\mathbf{1}).
\]
Here $Z^{(j-1)}$ is the Gaussian integral that normalises the fluctuation measure of step $j$, as in
\cite[(3.28)]{B-CMP119} and \cite{B-CMP116}. It is evaluated at the minimal configuration $U_j$,
the function of the block field written $U_k(V_k)$ at level $k$ in
Definition~\ref{1.5:def-domains}, here at level $j$, and at the trivial configuration $\mathbf{1}$.
\item[(b)] The logarithm of the fluctuation integral, that is, the last logarithm in
\cite[(3.28)]{B-CMP119}. It is built from the effective action $E_{j-1}$. As in \cite{B-CMP119},
the aged large-field remainders $R^{(i)}$ produced at the earlier steps $i\le k_1(j-1)$ are
included in $E_{j-1}$; their marking is fixed in the section on the large-field part of the
coupling increment.
\end{itemize}
In its localised representation this regular part is a sum of terms $E^{(j)}(X;(U,J),z)$. Here $X$
is a localisation domain, $z\in X$ is a base point, and $(U,J)$ is a configuration together with a
current. On $\mathbb{T}^{(j)}_\varsigma$ these terms are written $E^{(j)}_\varsigma$.

Consider the map $B\mapsto (U_j,J_j)(e^{iB})$. It sends a real Lie-algebra-valued field $B$ on
unit-lattice bonds to the minimal configuration $U_j(e^{iB})$ determined by the block field
$e^{iB}$ and to its current $J_j(e^{iB})$ \cite[(1.8)]{B-CMP109},
\cite{B-CMP102b}. Let $\mathfrak{h}_\infty$ be the \emph{free-boundary response} on $\mathbb{Z}^4$:
$\mathfrak{h}_\infty\delta_x$ is the sum, over all walks of the infinite lattice, of the random walk
expansion \cite[Theorems~3.10, 3.12]{B-CMP99b} of the linearisation of this map at $B=0$,
$U=\mathbf{1}$ (Ba{\l}aban's operator with free boundary conditions \cite{B-CMP109}). This sum
converges absolutely, and $\mathfrak{h}_\infty$ is the bondwise limit of the torus linearisations
$\mathfrak{h}_\varsigma$, the linearisations at $B=0$ of the same map on
$\mathbb{T}^{(j)}_\varsigma$. The bondwise limit is established in the proof of the lemma on the
existence of the increment. That lemma is stated under two hypotheses, bounds on the whole-lattice
effective action of step $j$ and the torus-independence of the terms localised in non-wrapping
domains. From that lemma the following enter here: the bondwise limit just stated, established in
its proof; two of its conclusions, namely the convergence of the sum over base points below and its
dependence on a pair of bonds only through their difference; and the whole-lattice family of
localised terms supplied by its torus-independence hypothesis, which enters
\eqref{4.2:eq-increment-1}.

For unit-lattice bonds $x,y$ (colour indices suppressed), let $\delta_x$ denote the field supported
on $x$, and let $\hat z\in\mathbb{Z}^4$. Set
\begin{equation}\label{4.2:eq-increment-1}
\Pi^{(j)}_\infty(x,y,\hat z):=
\sum_{\substack{\hat X\in\mathcal{D}_j^\infty\\ \hat X\ni\hat z}}
\bigl\langle \partial^2E^{(j)}(\hat X;(\mathbf{1},0),\hat z);\
\mathfrak{h}_\infty\delta_x,\ \mathfrak{h}_\infty\delta_y\bigr\rangle ,
\end{equation}
where $\langle\partial^2E^{(j)}(\hat X;(\mathbf{1},0),\hat z);v_1,v_2\rangle$ denotes the second
derivative at $(\mathbf{1},0)$ of $(U,J)\mapsto E^{(j)}(\hat X;(U,J),\hat z)$ in the directions
$v_1,v_2$, and $E^{(j)}(\hat X;\cdot)$ is the family of whole-lattice localised terms with which the
torus terms on non-wrapping domains coincide (the torus-independence hypothesis of the lemma on the
existence of the increment). Set also
\begin{equation}\label{4.2:eq-increment-2}
\Pi^{(j)}_\infty(x,y):=\sum_{\hat z}\Pi^{(j)}_\infty(x,y,\hat z).
\end{equation}
By the lemma on the existence of the increment, this sum converges and depends
on $(x,y)$ only through $x-y$. We write it $\Pi^{(j)}_\infty(x-y)$, with components
$\Pi^{(j)}_{\infty,\mu\nu}(u)$, $u\in\mathbb{Z}^4$, for bonds in the directions $\mu,\nu$. The
\emph{coupling increment} is
\begin{equation}\label{4.2:eq-increment-3}
\beta_j:=\sum_{u\in\mathbb{Z}^4}\Pi^{(j)}_{\infty,\mu\nu}(u)\,u_\mu u_\nu\qquad(\mu\neq\nu),
\end{equation}
which is \cite[(1.22)]{B-CMP109}, identified with \cite[(3.61)]{B-CMP119}.

The terms split into an unmarked and a marked (large-field-remainder) sub-family,
\[
E^{(j)}=E^{(j),\mathfrak{A}}+E^{(j),\mathfrak{R}},
\]
the marking being the one fixed in the section on the large-field part of the coupling increment.
The same formulas \eqref{4.2:eq-increment-1}--\eqref{4.2:eq-increment-3} applied to each sub-family
define $\beta^{\mathfrak{A}}_j$ and $\beta^{\mathfrak{R}}_j$. The superscripts $\mathfrak{A}$ and
$\mathfrak{R}$ label the two sub-families and carry no other meaning.

The increment $\beta_j$ is also distinct from two other letters:
\begin{itemize}
\item Ba{\l}aban's enlargement constant $\beta_{\mathrm{s}}$ (printed as $\beta$;
\cite{B-CMP122a} assumes $\beta_{\mathrm{s}}\le 1/4$);
\item the H\"older exponent $\beta_H$.
\end{itemize}
\end{definition}

\begin{remark}[How the torus enters the increment]
Given the torus-independence of the terms localised in non-wrapping domains, which is a hypothesis
of the lemma on the existence of the increment and rests, among other inputs, on the coincidence of
the localised terms with those of the whole lattice \cite[p.~278]{B-CMP119}, the torus enters the
second variation in only two ways.
\begin{itemize}
\item Through the terms localised in wrapping domains.
\item Through the argument: \cite[(1.20)]{B-CMP109} inserts the minimal configuration
$U_j(e^{iB})$ of the torus $\mathbb{T}^{(j)}_\varsigma$ into the terms from outside.
\end{itemize}
This is why \cite[(1.22)]{B-CMP109} is defined only after a limit in which the torus grows to the
whole lattice. No rate for that limit is needed here, and $\beta_j$ is not computed on any
finite torus.
\end{remark}

Let $Z^{(l-1)}$ denote the normalising factor of the Gaussian fluctuation integral of step $l$,
regarded as a function of the background configuration, and let $Z^{(l-1)}(\mathbf{1})$ be its value
at the trivial background. Put
\[
F_{l-1}:=\log Z^{(l-1)}-\log Z^{(l-1)}(\mathbf{1}).
\]
For a $\mathfrak{g}$-valued field $B$ on the block lattice, let $U_l(e^{iB})$ be the minimal
configuration determined by the block field $e^{iB}$. We define
$\beta^0_l$ by applying the construction of the coefficient in Definition~\ref{4.2:def-increment} to
the functional $B\mapsto F_{l-1}(U_l(e^{iB}))$ alone. In other words, we apply it only to the piece of
the functional of that definition that is formed by the Gaussian normaliser.

This functional involves no coupling and no earlier increment. By part~(a) of
Lemma~\ref{4.2:lem-one-loop} below, its localised family satisfies hypotheses (A) and (B) of the
existence lemma for the coefficient of Definition~\ref{4.2:def-increment} with $j=l$.
Hence, by the existence lemma, $\beta^0_l$ is a well-defined
real number depending only on $l$, $L$, $N$ and the auxiliary parameters $\mathfrak{p}$. In
particular it does not depend on the couplings $g_0,\dots,g_{l-1}$, on $K$ or on the torus.

Realise the Lie algebra $\mathfrak{g}$ of $\mathrm{SU}(N)$ as the traceless Hermitian $N\times N$
matrices, with the inner product $\langle Y_1,Y_2\rangle=\operatorname{tr}(Y_1Y_2)$, where
$\operatorname{tr}\mathbf{1}=1$. Let $\{t_a\}$ be an orthonormal basis of $\mathfrak{g}$ for this
inner product, and define the Casimir constant $C_{\mathfrak g}$ by
\[
\sum_a[t_a,[t_a,Y]]=C_{\mathfrak g}Y\qquad(Y\in\mathfrak{g}).
\]
Write $\operatorname{Tr}=N\operatorname{tr}$ for the unnormalised trace. Let $\{T_a\}$ be a basis of
$\mathfrak{g}$ with $\operatorname{Tr}(T_aT_b)=\tfrac12\delta_{ab}$. Then
\[
\operatorname{tr}(t_at_b)=\tfrac1N\operatorname{Tr}(t_at_b)=\delta_{ab}
\quad\text{for}\quad t_a=\sqrt{2N}\,T_a .
\]
Moreover $\sum_a[T_a,[T_a,Y]]=NY$, so $C_{\mathfrak g}=2N\cdot N=2N^2$. Hence the constant $c_0$ of
Definition~\ref{1.2:def-flow-constants} is
\begin{equation}\label{4.2:eq-one-loop-1}
c_0=\frac{11\,C_{\mathfrak g}}{24\pi^2}\log L=\frac{11N^2}{12\pi^2}\log L>0\qquad(N\ge2).
\end{equation}

\begin{lemma}[{the one-loop coefficient and the finite part of the one-loop law;
\cite[(1.22), Theorem~3]{B-CMP109}, \cite[(2.41)]{B-CMP116}}]
\label{4.2:lem-one-loop}
With $F_{l-1}$, $\beta^0_l$ and $c_0$ as above, the following hold.
\begin{itemize}
\item[(a)] Let $l\ge1$. The localised family of the functional $B\mapsto F_{l-1}(U_l(e^{iB}))$
satisfies hypotheses (A) and (B) of the existence lemma for the coefficient of
Definition~\ref{4.2:def-increment} with $j=l$, read for
this family alone. In hypothesis (A) the terms of the family are analytic on polydiscs $P_a(X)$
and bounded there by $E^\natural e^{-\kappa d_l(X)}$; the bound $E^\natural$ and the decay rate
$\kappa$ depend only on $L$, $N$ and $\mathfrak{p}$. The analyticity radius $a=a_l>0$ depends in
addition on $l$, through the scale
factor of the potentials; in the norms of \cite[(1.13)]{B-CMP109} the radius is free of $l$. Only
$a_l>0$ is used.
\item[(b)] There is a constant $c_2\ge0$, depending only on $L$ and $N$, such that for every
$k\ge0$
\begin{equation}\label{4.2:eq-one-loop-2}
\Bigl|\sum_{l=1}^{k}\beta^0_l-k\,c_0\Bigr|\le c_2 .
\end{equation}
Consequently, for all $0\le i<n$, with $c_0$ as in \eqref{4.2:eq-one-loop-1},
\begin{equation}\label{4.2:eq-one-loop-3}
\sum_{l=i+1}^{n}\beta^0_l\ \ge\ (n-i)\,c_0-2c_2 .
\end{equation}
The bound \eqref{4.2:eq-one-loop-3} has the same form for every $N\ge2$.
\end{itemize}
\end{lemma}

The lemma is used here as stated in \cite[(1.22), Theorem~3]{B-CMP109} and
\cite[(2.41)]{B-CMP116}.

For part~(a), the localised family of $B\mapsto F_{l-1}(U_l(e^{iB}))$ is obtained from
the random walk expansion of the fluctuation covariance \cite[Theorems~3.10, 3.12]{B-CMP99b}. Three
ingredients enter.
\begin{itemize}
\item The terms of the family are analytic on the polydisc $P_{a_l}(X)$ and decay exponentially
in $d_l(X)$. This gives the bounds of hypothesis (A).
\item Walks whose footprints lie in Dirichlet cubes inside a non-wrapping domain coincide on the
torus and on $\mathbb{Z}^4$, and their terms are equal there. This gives hypothesis (B).
\item The covariance of the totals under reflections uses the three statements on the covariance
of the totals that accompany the existence lemma for the coefficient of
Definition~\ref{4.2:def-increment}.
\end{itemize}
For (b), the telescoping identity and the two comparison bounds with $k$-independent constants are
used as stated in \cite[Theorem~3]{B-CMP109} and \cite[(2.41)]{B-CMP116}; the derivation of (b)
from them is as follows.

\emph{Outline of (b).} Put $\eta:=L^{-k}$. The bound \eqref{4.2:eq-one-loop-2} rests on two facts.

\emph{The composite functional.} At zero coupling the fluctuation integral of step $l+1$,
$0\le l<k$, is the Gaussian integral normalised by $Z^{(l)}$. The Gaussian integrals of steps
$1,\dots,k$ compose, because the block-averaging maps compose to $\mathcal{M}^k$; here $\mathcal{M}$
denotes the block-averaging map of Definition~\ref{1.5:def-block-average}. For a block field $V$ at
scale $k$, read each $F_l$ with $l<k$ at the
scale-$(l+1)$ block configuration derived from the minimal configuration $U_k(V)$. Let $B_k$ be the
coefficient of Definition~\ref{4.2:def-increment} of the composite functional
$V\mapsto\sum_{l<k}F_l$ read in this way.

The coefficient of Definition~\ref{4.2:def-increment}, regarded as a function of the functional to
which it is applied, has three properties. It is linear in the functional. It is the coefficient of
order two in the momentum of a translation-invariant, gauge-invariant quadratic form. It is not
changed when a functional of a finer block field is read through the minimal configuration
determined by a coarser one. These three properties give the telescoping identity
\[
\sum_{l=1}^{k}\beta^0_l=B_k .
\]

\emph{The one-loop value.} Let $D^2_{\mathrm{W}}$ be the Wilson Hessian in background Feynman gauge
and $\Delta_0$ the Faddeev--Popov operator, both on the same fine lattice. Let
$B^{\mathrm{ref}}_k$ be the coefficient of Definition~\ref{4.2:def-increment} of the reference
functional
\[
-\tfrac12\log\det(D^2_{\mathrm{W}}+\eta^2)+\log\det(\Delta_0+\eta^2),
\]
minus its value at the trivial background. There are constants $C_{\mathrm{det}}(N)$ and
$C_{\mathrm{cmp}}(L,N)$, independent of $k$, such that
\begin{gather*}
|B_k-B^{\mathrm{ref}}_k|\le\tfrac12C_{\mathrm{cmp}}(L,N),\\
\Bigl|B^{\mathrm{ref}}_k-\frac{11C_{\mathfrak g}}{24\pi^2}\log\frac1\eta\Bigr|\le
C_{\mathrm{det}}(N).
\end{gather*}

The first bound is the comparison of the composite Gaussian normaliser with the reference
determinants at $\eta=L^{-k}$. The second is the lattice one-loop computation. There the gluon
determinant contributes $5C_{\mathfrak g}/(12\pi^2)$ per unit of $\log(1/\eta)$, and the ghost
determinant contributes $C_{\mathfrak g}/(24\pi^2)$. Their sum is $11C_{\mathfrak g}/(24\pi^2)$, since
$\tfrac{5}{12}+\tfrac{1}{24}=\tfrac{11}{24}$.

Since $\log(1/\eta)=k\log L$, \eqref{4.2:eq-one-loop-1} gives
\[
\frac{11C_{\mathfrak g}}{24\pi^2}\log\frac1\eta=k\,c_0 .
\]
Put
\[
c_2:=C_{\mathrm{det}}(N)+\tfrac12C_{\mathrm{cmp}}(L,N);
\]
this is the constant $c_2$ of Definition~\ref{1.2:def-flow-constants}.
The telescoping identity, the triangle inequality and the two bounds above then give
\[
\Bigl|\sum_{l=1}^{k}\beta^0_l-k\,c_0\Bigr|
\le|B_k-B^{\mathrm{ref}}_k|+|B^{\mathrm{ref}}_k-k\,c_0|\le c_2 ,
\]
which is \eqref{4.2:eq-one-loop-2}.

For \eqref{4.2:eq-one-loop-3}, split the sum as in the first line of the display below. Apply
\eqref{4.2:eq-one-loop-2} with $k=n$ to bound the first sum on the right from below, and with $k=i$
to bound the second sum from above. This gives
\begin{align*}
\sum_{l=i+1}^{n}\beta^0_l&=\sum_{l=1}^{n}\beta^0_l-\sum_{l=1}^{i}\beta^0_l\\
&\ge(nc_0-c_2)-(ic_0+c_2)=(n-i)\,c_0-2c_2 .
\end{align*}

No sign of an individual $\beta^0_l$ is asserted: \eqref{4.2:eq-one-loop-3} with $n=i+1$ gives only
$\beta^0_l\ge c_0-2c_2$.

The order of choice of the constants (Notation~\ref{2.3:not-stages-first},
Notation~\ref{2.3:not-stages-middle} and Notation~\ref{2.3:not-stages-last}) imposes the upper bound
$\varepsilon_1^{\mathrm{I}}\le c_0/(2C_{\mathrm{rem}})$ on Ba{\l}aban's small-field parameter
$\varepsilon_1^{\mathrm{I}}$, where $C_{\mathrm{rem}}$ is the constant of
Lemma~\ref{4.2:lem-unmarked}. Since $C_{\mathrm{rem}}$ is fixed before $\varepsilon_1^{\mathrm{I}}$,
this bound is one-sided in the order of hypothesis (H3) of Notation~\ref{2.2:not-hypotheses}.

Assume this bound, and assume the hypotheses of Lemma~\ref{4.2:lem-unmarked} for $l\le n$.
By Lemma~\ref{4.2:lem-unmarked}, whose proof is incomplete at one step, the unmarked part
$\beta^{\mathfrak A}_l$ of the increment satisfies
$|\beta^{\mathfrak A}_l-\beta^0_l|\le C_{\mathrm{rem}}\varepsilon_1^{\mathrm{I}}$. Summing this bound
over $i<l\le n$ and combining it with \eqref{4.2:eq-one-loop-3} gives
\[
\sum_{l=i+1}^{n}\beta^{\mathfrak{A}}_l\ge\frac{(n-i)\,c_0}{2}-2c_2 .
\]

\begin{lemma}[Bounds on the unmarked part of the increment. Stated; proof incomplete at the
step marked $(\ast)$]\label{4.2:lem-unmarked}
There are constants $C_{\mathrm{rem}}$ and $C_\beta$ with the following property. The constant
$C_{\mathrm{rem}}$ depends only on $L$, $N$ and the parameters of hypothesis \textup{(H3)}
(Notation~\ref{2.2:not-hypotheses}) chosen, in the order of choice
(Notations~\ref{2.3:not-stages-first} to~\ref{2.3:not-stages-last}), before
$\varepsilon_1^{\mathrm{I}}$, so not on $\varepsilon_1^{\mathrm{I}}$. The constant $C_\beta$ depends
only on $L$, $N$ and the parameters of hypothesis \textup{(H3)} fixed before $\gamma$.

Let $l\ge1$. Write $E^{(l),\mathfrak{A}}$ for the unmarked part, in the sense of the definition of
the coupling increment (Definition~\ref{4.2:def-increment}), of the whole-lattice small-field
effective action produced by the fluctuation integral of step $l$ at coupling $g_{l-1}$. By that
definition, $E^{(l),\mathfrak{A}}$ is the sum of the Gaussian normaliser and a second summand,
which we call the remainder part of $E^{(l),\mathfrak{A}}$.
Assume the following.
\begin{itemize}
\item $g_0,\dots,g_{l-1}\in(0,\gamma]$, and steps $1,\dots,l-1$ have been performed along these
couplings.
\item The inductive bounds in the statement of Theorem~\ref{4.20:thm-representation} (the
representation of the effective densities) hold at scales below $l$.
\item $E^{(l),\mathfrak{A}}$ satisfies, at step $l$, the bounds of Ba{\l}aban's one-step
renormalization theorem.
\item The remainder part of $E^{(l),\mathfrak{A}}$ obeys Ba{\l}aban's bound
\cite[(2.41)]{B-CMP116}, with its factor $\varepsilon_1^{\mathrm{I}}$. For the whole-lattice
small-field effective action produced at step $l$ by the proposition of this chapter that
constructs the family of effective actions, this hypothesis holds by the second part of that
proposition, under the two upper bounds on $\gamma$ stated there.
\end{itemize}
Then
\begin{equation}\label{4.2:eq-unmarked-1}
\bigl|\beta^{\mathfrak{A}}_l-\beta^0_l\bigr|\ \le\ C_{\mathrm{rem}}\,\varepsilon_1^{\mathrm{I}},
\qquad
\bigl|\beta^{\mathfrak{A}}_l\bigr|\ \le\ C_\beta .
\end{equation}
Here $\beta^0_l$ denotes the coefficient that defines the coupling increment in
Definition~\ref{4.2:def-increment}, evaluated on the Gaussian normaliser alone, and
$\varepsilon_1^{\mathrm{I}}$ is Ba{\l}aban's small-field parameter $\varepsilon_1$ of
\cite{B-CMP109,B-CMP116}. It carries a superscript to distinguish it from the running thresholds
$\varepsilon_k$.
\end{lemma}

Stated; the proof below is incomplete at the step marked $(\ast)$. The estimates of
\cite{B-CMP116} and \cite{B-CMP109} that the proof combines are quoted, not
re-derived.

\begin{proof}[Outline of the proof]
By the definition of the coupling increment (Definition~\ref{4.2:def-increment}),
$E^{(l),\mathfrak{A}}$ is the sum of the Gaussian normaliser and its remainder part, and
$\beta^0_l$ is the coefficient that defines the increment there, evaluated on the Gaussian
normaliser alone. That coefficient is linear in the family to which it is applied.
Hence $\beta^{\mathfrak{A}}_l-\beta^0_l$ is that coefficient evaluated on the remainder part
of $E^{(l),\mathfrak{A}}$.

$(\ast)$ The argument is Ba{\l}aban's
for \cite[Lemma~3]{B-CMP116}, carried out with the small-field cutoff of \cite{B-CMP119}, at
the radius $\delta_{l-1}/g_{l-1}$, where $\delta_{l-1}$ is the small-field radius parameter of
step $l-1$ in \cite{B-CMP119}. The bound
\cite[(2.41)]{B-CMP116} holds on Ba{\l}aban's analyticity domain
$U^c_l(X,\alpha_0,\alpha_1)$; here $X$ is the localisation domain of \cite{B-CMP116}, not a
large-field region. The radii $\alpha_0$, $\alpha_1$ of this domain and the decay constant of
that bound do not depend on the domain or on the
scale index, because Ba{\l}aban's constants are absolute, independent of $X$ and of the scale
\cite{B-CMP109}. In the order of hypothesis (H3) they are chosen before
$\varepsilon_1^{\mathrm{I}}$. Below, the existence lemma is the lemma of this chapter that
establishes the existence of the coefficient of Definition~\ref{4.2:def-increment} for families
obeying, with index $j=l$, the bounds assumed in the third hypothesis above. The scale-dependent
constant $C_{\mathfrak{h}}$ that enters the
proof of the existence lemma is not used here.

Measure the response in the norms that define $U^c_l(X,\alpha_0,\alpha_1)$, in which a potential
carries the scale factor of \cite[(1.13)]{B-CMP109}. In these norms the random-walk expansion of
\cite[Theorems~3.10, 3.12]{B-CMP99b} gives, for every bond $b$,
\begin{equation*}
|(\mathfrak{h}_\infty\delta_x)(b)|\le C\,e^{-\delta_0\operatorname{dist}(x,b)},
\end{equation*}
with $C$ independent of $l$, where $\mathfrak{h}_\infty$ is the free-boundary response operator
of Definition~\ref{4.2:def-increment}, $\delta_x$ is the unit source at $x$, and $\delta_0$ is
Ba{\l}aban's decay constant, the one that enters the rate of \cite[(5.10)]{B-CMP109} below.
Therefore, in the Cauchy
estimate of the proof of the existence lemma, the scale factor of the response and the scale
factor of the analyticity radius cancel.

Write $\Pi_{\mu\nu}$ for the tensor $\Pi^{(l)}_{\infty,\mu\nu}$ of
Definition~\ref{4.2:def-increment}, formed from $E^{(l),\mathfrak{A}}$ in place of the complete
regular part of step $l$, so that $\beta^{\mathfrak{A}}_l=\sum_u\Pi_{\mu\nu}(u)u_\mu u_\nu$.
Write $\Pi^{\mathrm{r}}_{\mu\nu}$ for the same tensor formed from the remainder part of
$E^{(l),\mathfrak{A}}$; by the first paragraph of this proof,
\begin{equation*}
\beta^{\mathfrak{A}}_l-\beta^0_l=\sum_u\Pi^{\mathrm{r}}_{\mu\nu}(u)\,u_\mu u_\nu .
\end{equation*}
By the cancellation of scale factors just described, the pointed tensor of the remainder part of
$E^{(l),\mathfrak{A}}$ (the three-point tensor whose sum over its base point is
$\Pi^{\mathrm{r}}_{\mu\nu}$) is bounded by $\varepsilon_1^{\mathrm{I}}$ times an exponential decay
in the distances to the base point, whose constants involve neither $l$ nor
$\varepsilon_1^{\mathrm{I}}$. Summing over the base point gives a bound of the same kind for
$\Pi^{\mathrm{r}}_{\mu\nu}(u)$, and summing $|u|^2$ against that decay gives the first bound in
\eqref{4.2:eq-unmarked-1}.

The second bound follows by estimating $\sum_u|u|^2|\Pi_{\mu\nu}(u)|$ with the decay
\cite[(5.10)]{B-CMP109}. That decay has rate $\tfrac12\min\{\delta_0,\kappa/M\}$. Its prefactor
depends on $d$, $L$, $\kappa$, $M$, the analyticity radii $\alpha_0$, $\alpha_1$ and the amplitude
$E_0$ of the bound \cite[(1.18)]{B-CMP109}. All of these are parameters fixed before $\gamma$, and
none of them depends on the step $l$. So $C_\beta$ involves only parameters fixed before $\gamma$.
\end{proof}

\begin{theorem}[Ba{\l}aban's one-step small-field renormalization theorem, cited by statement, in the family form in which this book uses it; {\cite[Theorem~3]{B-CMP109}}]
\label{4.3:thm-small-field-step}
Assume the operator bounds at free current (Lemma~\ref{4.5:lem-free-current-operators}), the
torus-independence of the terms localized in
non-wrapping domains, and the statement on the one-power hereditary class. Assume that $\kappa$
and $M$ satisfy
Ba{\l}aban's lower bounds in \cite[Theorem~3]{B-CMP109}. Assume that $\gamma$ is at most the
threshold of \cite[Theorem~3]{B-CMP109}, which depends on all the constants chosen before it
(among them $\kappa$, $M$, the radii $\varepsilon^{\mathrm{I}}_0$, $\varepsilon^{\mathrm{I}}_1$,
$\alpha_0$, $\alpha_1$). Assume also that $\gamma$ is at most the threshold of
\cite[(2.31)]{B-CMP116}, which depends on $M$ and $\kappa$. The constant $\gamma$ is chosen after
all of them; these are among the conditions of hypothesis (H3) in
Notation~\ref{2.2:not-hypotheses}.

Let $k\ge0$, and let $\{(E^{(j)},\beta_j)\}_{j\le k}$ be any family of effective actions
$E^{(j)}$ and numbers $\beta_j$ with Ba{\l}aban's inductive properties, namely:
\begin{itemize}
\item the properties \cite[(1.7), (1.18), (1.19)]{B-CMP109}, the analyticity in
\cite[(1.18)]{B-CMP109} being read on the one-power hereditary class, for which the passage to
subregions built from finer cubes is an identity;
\item the Euclidean clause among Ba{\l}aban's inductive assumptions
\cite[(1.1)--(1.22)]{B-CMP109};
\item torus-independence of the terms of $E^{(j)}$ localized in non-wrapping domains;
\item $\beta_j$ is the coupling increment of Definition~\ref{4.2:def-increment}, the
$\mathbb{Z}^4$ coefficient \cite[(1.22)]{B-CMP109} of $E^{(j)}$.
\end{itemize}
Let $g\in(0,\gamma]$ be arbitrary; in this statement only, $g$ denotes this coupling. Then the
fluctuation integral \cite[(2.12)--(2.13)]{B-CMP109}
at coupling $g$ produces an effective action $E^{(k+1)}$ with the same properties and the same
constants. The recursion that follows the integral only relabels.

Let $p_1(g'):=(\log g'^{-2})^{p_1}$, where $p_1$ is the auxiliary parameter of that name in
$\mathfrak{p}$ (Notation~\ref{2.1:not-prefix}), the exponent in the radius $C_1p_1(g)$ of the
fluctuation domain \cite[(3.43)]{B-CMP119}. Suppose moreover that:
\begin{itemize}
\item $C_1\,g'\,p_1(g')\le\varepsilon^{\mathrm{I}}_1$ for every $g'\in(0,\gamma]$;
\item the cutoff radius $(A_1/A_0)\,p_0(g)$, where $p_0(\cdot)$ is the degree function, exceeds
the constant of \cite[(2.31)]{B-CMP116}, which depends on $M$ and $\kappa$. When $g=g_k$, this
radius is the quantity $\delta_k/g_k$ of \cite{B-CMP119}.
\end{itemize}
Then the same conclusion holds with the following two changes. First, the fluctuation domain is
that of \cite[(3.43)]{B-CMP119}, on which the modulus of the fluctuation field is strictly less
than $C_1p_1(g)$. Second, the cutoff radius is the one just described. In addition, the
part of $E^{(k+1)}$ other than its Gaussian part obeys \cite[(2.41)]{B-CMP116} with the factor
$\varepsilon^{\mathrm{I}}_1$.
\end{theorem}

\cite[Theorem~3]{B-CMP109} states the following, where we write $\gamma_0$ for the
constant that it calls $\gamma$, so that the constant $\gamma$ of the theorem above is at most
$\gamma_0$: there exist positive constants
$\kappa_0$, $M(\kappa)$, $\gamma_0$, $\varepsilon^{\mathrm{I}}_0$, $\varepsilon^{\mathrm{I}}_1$,
$\alpha_0$, $\alpha_1$ such that, if $\kappa\ge\kappa_0$, $M\ge M(\kappa)$ and $0<g_k\le\gamma_0$
for $k=0,1,\dots,K$, then the sequence of actions defined inductively by the small-field
renormalization transformations \cite[(0.17)--(0.20)]{B-CMP109} satisfies all the inductive
assumptions \cite[(1.1)--(1.22)]{B-CMP109}; the constants $\varepsilon^{\mathrm{I}}_0$,
$\varepsilon^{\mathrm{I}}_1$, $\alpha_0$, $\alpha_1$ depend on $M$, and $\gamma_0$ depends on all
the other constants. Its general form is \cite[Theorem~1]{B-CMP109}. The first part of
Theorem~\ref{4.3:thm-small-field-step} is \cite[Theorem~3]{B-CMP109}
read for an arbitrary
family with the inductive properties listed above. Its second part is the argument of
\cite{B-CMP116}, including its bound \cite[(2.41)]{B-CMP116}, read with the fluctuation domain of
\cite[(3.43)]{B-CMP119} under the two conditions stated above. For the proofs we refer to
\cite[\S\S2--5]{B-CMP109} and \cite{B-CMP116}.

In Ba{\l}aban's argument the couplings enter only through the containment $0<g_k\le\gamma$ and,
in the second part of the theorem above (the fluctuation domain of \cite[(3.43)]{B-CMP119} with
the cutoff radius there), through the two conditions stated there; the remark following
\cite[Theorem~1]{B-CMP109} states that no special asymptotic behaviour of the couplings is
assumed. Of Ba{\l}aban's hypotheses on the couplings only containment is invoked in this book.
\cite[Theorem~2]{B-CMP109}, together with \cite[(0.31)]{B-CMP109}, would remove it by
second-order perturbation theory; \cite{B-CMP109} refers its proof to a separate paper, and
\cite{B-CMP122b} refers to that proof as unpublished; it is not used here.

\begin{theorem}[Ba{\l}aban's ultraviolet stability theorem {\cite[Theorem~1]{B-CMP122b}}]
\label{4.3:thm-uv-stability}
Let $\gamma$ be the constant of Ba{\l}aban's one-step renormalization theorem, and assume that $\gamma$ is at
most the smallness threshold of \cite[Theorem~1]{B-CMP122b}. If the effective
coupling constants $g_0,\dots,g_K$ all lie in the interval $(0,\gamma]$,
then for every $k \le K$ the effective density $\rho_k$ has the form, and satisfies all the
conditions and bounds, described in \cite[\S2]{B-CMP119}. In particular, $\rho_k$ has the
inductive representation \cite[(2.18)]{B-CMP119}, and this representation satisfies the
inductive bounds of \cite[\S2]{B-CMP119}.
\end{theorem}

This is \cite[Theorem~1]{B-CMP122b}. The large-field renormalization operation that enters
the representation, and the bounds on it, are constructed in \cite{B-CMP122a} and
\cite{B-CMP122b}.

The representation \cite[(2.18)]{B-CMP119} of $\rho_k$ is a sum over admissible sequences of
large-field domains. Each term is a composition of large-field operations applied to a product
of characteristic functions with the exponential of an effective action.

The theorem has one hypothesis on the effective couplings, namely that they all lie in the
interval $(0,\gamma]$; the smallness of $\gamma$ assumed in the statement is a condition on the
constant $\gamma$, not on the couplings.
Ba{\l}aban remarks \cite[pp.~355--356]{B-CMP122b} that \cite[Theorem~2]{B-CMP109} would allow
this hypothesis to be removed and that the proof of that theorem has not been published, and he
states the theorem with the hypothesis. In this book that hypothesis is supplied by clause~(0) of
Theorem~0
(Theorem~\ref{4.1:thm-clause-zero}),
through the containment statement (Proposition~\ref{4.19:prop-containment}). That
proposition gives $g_k \in (0,\gamma]$ for every $k \le K$.

Clause~(i) of Theorem~0 is the theorem above, with two qualifications. First, the displays of
\cite{B-CMP119}, \cite{B-CMP109} and \cite{B-CMP122b} that its bounds use are read with every
quantity in them taken as Ba{\l}aban defines it. Second, for
boundary terms beside
live regions (large-field regions that still carry their large-field factor at the level read)
the bounds read are those stated in the theorem on the inductive representation of the
effective densities (Theorem~\ref{4.20:thm-representation}), in place of the corresponding
bounds of \cite[\S2]{B-CMP119}.

\begin{remark}[Ba{\l}aban's independence assertion]\label{4.4:rem-independence-assertion}
Let $\Lambda_{k+1}$ be the small-field region of the $(k+1)$-st step. The following two sentences are
stated in \cite[p.~278]{B-CMP119} without proof:
\begin{quote}
The important remark is that the terms $\mathbf{E}_0^{(k+1)}(\Lambda_{k+1},X,z)$ of this
representation, for $X \subset \Lambda_{k+1}$, do not depend on $\Lambda_{k+1}$, and coincide with
the corresponding terms arising from the expressions defined on the whole lattice, i.e.\ in the
framework of [I].
\end{quote}
\begin{quote}
The terms with localization domains contained in $\Lambda_{k+1}$ have the same important property
as the corresponding $\mathbf{E}^{(k+1)}$-terms, namely they do not depend on $\Lambda_{k+1}$.
\end{quote}
In \cite{B-CMP119} the reference [I] is \cite{B-CMP109}.

We write $\mathcal{S}$ for the setting of \cite{B-CMP119}. In $\mathcal{S}$ the $(k+1)$-st
$\mathbf{T}$-step is carried out with the small-field region $\Lambda_{k+1}$. We write
$\mathcal{S}_0$ for the whole-lattice setting of \cite{B-CMP109}, that is, the framework [I] of the
first quotation, in which the same expressions are defined on the whole lattice.

Let $\mathcal{P}$ be one of the operators, listed in this chapter, through which non-locality enters
the $(k+1)$-st $\mathbf{T}$-step of \cite{B-CMP119}. That list contains in particular the operators
\begin{equation*}
C^{(k)}(\Lambda_{k+1}),\quad \bigl(C^{(k)}(\Lambda_{k+1})\bigr)^{1/2},\quad
(xI+Q_{\Lambda\Lambda})^{-1},\quad (Q_{\Lambda\Lambda}+\lambda)^{-1},\quad C^{(k)}(Z_0),
\end{equation*}
that is, the conditioned covariance $C^{(k)}(\Lambda_{k+1})$ and its square root, the two
inverses formed from $Q_{\Lambda\Lambda}$, and the covariance $C^{(k)}(Z_0)$ on the
region $Z_0$. These operators are written in the notation of \cite{B-CMP119} and
\cite{B-CMP99b}; in the display above the letters $x$, $\lambda$, $I$ (the identity), $Q$, and the
letter $\Lambda$ in the subscript of $Q_{\Lambda\Lambda}$, have the meaning they have there, not the
meaning given to them in the glossary.
Fix once and for all, independently of the setting, the walk-expansion rule constructed later in
this chapter, which writes $\mathcal{P}$ as a sum over walks,
\begin{equation}\label{4.4:eq-independence-assertion-1}
\mathcal{P} = \sum_{\omega} \mathcal{P}_\omega ,
\end{equation}
in which each walk $\omega$ carries a footprint $|\omega|$, in the sense fixed with that rule. We
make the two printed sentences above precise as the following statement.

\medskip
\noindent\emph{For every operator $\mathcal{P}$ of that list, expanded by the fixed rule as in
\eqref{4.4:eq-independence-assertion-1}, and for every walk $\omega$ with
$|\omega| \subset \Lambda_{k+1}$, the walk $\omega$
occurs in the expansion \eqref{4.4:eq-independence-assertion-1} in the setting $\mathcal{S}$ if and
only if it occurs there in the setting $\mathcal{S}_0$, and}
\begin{equation}\label{4.4:eq-independence-assertion-2}
\mathcal{P}_\omega[\mathcal{S}] = \mathcal{P}_\omega[\mathcal{S}_0]
\quad\text{as functions of } (\mathbf{U},\mathbf{J})|_{|\omega|} .
\end{equation}

\medskip
Here $\mathcal{P}_\omega[\mathcal{S}]$ and $\mathcal{P}_\omega[\mathcal{S}_0]$ denote the term of
index $\omega$ in \eqref{4.4:eq-independence-assertion-1} in the two settings. The restriction
$(\mathbf{U},\mathbf{J})|_{|\omega|}$ is the restriction of the complexified bond configuration and
of the current variable to the footprint of $\omega$.

This statement concerns the walk factors in the proof of the independence of the localised terms
from the small-field region. Together with a lemma of this chapter, it is exactly what the walk
factors require in the step of that proof in which the ingredients of the two settings are shown to
agree on $\Lambda_{k+1}$. The statement is itself proved as a theorem later in this chapter.
\end{remark}

\begin{definition}[Units, distances and bond sets]\label{4.4:not-units-distances}
Throughout this section the level $k$ is fixed, and lengths are measured in units of level $k$:
$T^{(k)}$ denotes the unit lattice, and we write $\eta = L^{-k}$.

\emph{Distances.} The symbol $|\cdot|$ denotes the $\ell^1$-distance used in \cite{B-CMP96}.
Tree lengths and the diameters of sets used in what follows are measured in this distance. The
words \emph{neighbourhood} and \emph{cube}, on the other hand, refer to the sup-distance.

\emph{Partitions and cut-offs.} The partitions used are the standard nested ones, in the sense of
this book's convention on partitions and of its statement on nested partitions. They are
compatible in the sense of \cite{B-CMP119}. By this book's convention on cut-offs, every cut-off
function is a fixed function of the
sup-distances from its argument to the faces of its cube, as in \cite[(1.118)]{B-CMP95},
\cite[(2.36)]{B-CMP96}, \cite{B-CMP109} and \cite[(3.40)]{B-CMP119}.

\emph{Cube sizes.} The cube sizes $M_1$, $R_1 M_1$ and $M$ of \cite{B-CMP116} satisfy
\begin{equation*}
M_1 \le R_1 M_1 \le M,
\end{equation*}
where $R_1 \ge 1$ is the factor of \cite{B-CMP116} for which $R_1 M_1$ is the intermediate cube
size: the cubes of the partition $\pi_k(1)$ of that paper are of size $M_1$ rather than $M$, and
$R_1 M_1 M^{-1} \le 1$.

\emph{Bond sets.} For a region $\Omega$, the set $\mathbb{B}(\Omega)$ consists of the bonds of
$T^{(k)}$ in $\Omega$, with the eliminated bonds $b_0(c)$, $c \in T^{(k+1)}$, removed
\cite{B-CMP109}, \cite{B-CMP99b}; here $T^{(k+1)}$ is the lattice of the next level, of spacing
$L$ in the present units, and $b_0(c)$ denotes, as in \cite{B-CMP109}, the eliminated bonds
attached to $c$. Thus
\begin{equation*}
\mathbb{B}(\Omega)
= \{\, b : b \text{ a bond of } T^{(k)} \text{ in } \Omega \,\}
\setminus \{\, b : b \text{ among the eliminated bonds } b_0(c),\ c \in T^{(k+1)} \,\}.
\end{equation*}
We write
\begin{equation*}
\mathbb{B}^{+} = \mathbb{B}(\Omega_{k+1}),
\end{equation*}
where $\Omega_{k+1}$ is the member of level $k+1$ of the nest of regions $\Omega_j$ in
Ba{\l}aban's construction.
\end{definition}

\begin{definition}[Settings, data domains and the analyticity class]\label{4.4:def-settings-data-domains}
Let $k$ be fixed, with the units and bond sets of Definition~\ref{4.4:not-units-distances}, and
write $T$ for the unit lattice $T^{(k)}$ fixed there.
A \emph{setting} is a quadruple
\begin{equation*}
  \mathfrak{s} = \bigl(\boldsymbol{\Omega}, Y;\, D, (\mathbf{U},\mathbf{J})\bigr)
\end{equation*}
whose entries are as follows.
\begin{enumerate}
  \item $\boldsymbol{\Omega}$ is a decreasing sequence of regions
  \begin{equation*}
    \boldsymbol{\Omega} = \bigl(\Omega_0 \supset \Omega_1 \supset \cdots \supset \Omega_k \supset \Omega_{k+1}\bigr),
  \end{equation*}
  in which the sequence $(\Omega_j)_{j \le k}$ is admissible in the sense of
  \cite[(2.1)--(2.2)]{B-CMP96}, with the amendment made in \cite{B-CMP99b}.
  The last member $\Omega_{k+1}$ is the region that carries the next block structure:
  the $L$-blocks, the block quantities $b_0(\cdot)$ and the operators written $C$, $h$,
  $\tilde{D}$ in \cite{B-CMP109}; these letters are used here in the sense of that paper,
  not in the sense of the glossary's $b_0$ and $h$.
  \item $Y \subset \mathbb{B}(\Omega_{k+1})$ is the set of \emph{integration bonds}.
  The pair $(\boldsymbol{\Omega}, Y)$ is called the \emph{geometry} of $\mathfrak{s}$ and is
  denoted $\mathfrak{g}$ (not to be confused with the complexified Lie algebra $\mathfrak{g}^c$).
  \item $D \subset T$ is a closed set, the \emph{data domain} of $\mathfrak{s}$.
  If $D = T$, the setting is called \emph{global}.
  If $D \neq T$, the setting is said to carry \emph{local data}; for such a setting the only
  quantities ever formed are the localised quantities $K^W$, $W \subset D$ (the superscript $W$
  is an index, not a power), which are defined later in this chapter.
  \item $(\mathbf{U}, \mathbf{J})$ is a datum on the bonds of $D$, belonging to a class of
  configurations on which the operators and the quadratic forms and covariances of the
  construction are used; the principal such class, $\mathbf{P}'$, is defined next.
\end{enumerate}

The class $\mathbf{P}'$ consists of the configurations $(\mathbf{U}, \mathbf{J})$ on the bonds of
$D$ that satisfy the regularity conditions (i)--(iii) of \cite[(1.11)--(1.14)]{B-CMP109} (cited
as I.(i)--(iii) in \cite{B-CMP116}) on $D$, with the constants $\alpha_0, \alpha_1$ replaced
by the constants $\alpha'_0, \alpha'_1$ of \cite{B-CMP116}, which are much bigger than
$\alpha_0, \alpha_1$; the current radius $\gamma_0$ of \cite[(1.14)]{B-CMP109} is replaced by a
radius of its own, written $\gamma'_0$.
In \cite{B-CMP116} it is stated that the quadratic forms and covariances in the expression
written $H(Z)$ in that paper are analytic functions on the space of configurations
$(\mathbf{U},\mathbf{J})$ that satisfy I.(i)--(iii) on a domain, there written $Z$, with the
constants $\alpha'_0, \alpha'_1$.

In $\mathbf{P}'$ the second variable $\mathbf{J}$ is free. The only condition imposed on it is
that its modulus lie below its own radius:
\begin{equation}\label{4.4:eq-settings-data-domains-1}
  |\mathbf{J}| < \gamma'_0 \quad \text{on the bonds of } D.
\end{equation}
This is the primed form of the current condition $|J| < \gamma_0$ of
\cite[(1.14)]{B-CMP109}. The letters $\alpha_0,\alpha_1$ are not used in the same places in
\cite{B-CMP109} and in \cite{B-CMP116}; here $\alpha_0,\alpha_1$ are those of
\cite[(1.11)--(1.14)]{B-CMP109}, and $\gamma'_0$ is the enlarged current radius.

A multi-scale version of these conditions is given in \cite[(2.34)--(2.39)]{B-CMP119}, as a more
detailed and precise form of the regularity conditions for background fields. The condition on
the current among them, \cite[(2.36)]{B-CMP119}, carries the radii $\alpha_{0,n}$.

No second-order regularity condition, in the sense of \cite[(3.36)]{B-CMP99b}, is part of
$\mathbf{P}'$. As explained in \cite{B-CMP109}, such conditions are not formulated there;
instead, the current expression $J_j$ of \cite{B-CMP109} is replaced by the new, independent
variable $\mathbf{J}$, and
\eqref{4.4:eq-settings-data-domains-1} is the only requirement placed on that variable.

The quantities formed from a setting, in particular the localised quantities $K^W$, are regarded
as functions of $(\mathbf{U},\mathbf{J})$; wherever such a quantity is said below to be analytic,
analyticity on $\mathbf{P}'$ is meant. The weight bound and the one-family statement of this
chapter are likewise stated on $\mathbf{P}'$. They are also read on a larger domain, defined
later in this chapter, which contains $\mathbf{P}'$ together with a second class of data; the
corollary of the weight bound holds on that larger domain. At the full-torus geometry, in which
$\Omega_j = T$ for all $j$ and $Y = \mathbb{B}(T)$, every representative of a datum of
$\mathbf{P}'$ lies in that larger domain, provided $\alpha'_0$, $\alpha'_1$, $\gamma'_0$ satisfy
three upper bounds imposed after $M$ is chosen; this is a corollary stated later in this
chapter.
\end{definition}

\begin{definition}[Local data classes and the gauge condition]\label{4.4:def-local-data-classes}
Let $\mathsf{o}$ be the number written $O(1)$ in the cube side $O(1)LM$ of
\cite[(1.12)]{B-CMP109}. It is a fixed absolute number, and the only property of
$\mathsf{o}$ used below is
\begin{equation*}
\mathsf{o}L\ge 1 .
\end{equation*}
Let $B^{109}$ be the constant of \cite[(1.12)]{B-CMP109}. Its value is fixed in the
gluing construction of this chapter, as the larger of $4C_{\mathrm{gl}}B$ and a floor on
$B^{109}$ fixed at the same place; it is chosen before $M$ in the order of choice of
Definitions~\ref{2.3:not-stages-first}--\ref{2.3:not-stages-last}. The bare letter $B$
keeps its meaning as the absolute
constant of
\cite[(2.38)]{B-CMP119}. With $\alpha_0'$ the primed radius of the class $\mathbf{P}'$ of
Definition~\ref{4.4:def-settings-data-domains}, put
\begin{equation*}
\mathsf{a}_U:=\mathsf{o}LMB^{109}\alpha_0' .
\end{equation*}

\emph{Interpretation of the cube size.} In \cite[(1.12)]{B-CMP109} the phrase ``for each
cube $\square\subset X$ of a size $O(1)LM$'' is read as ``for each cube
$\square\subset X$ of a size at most $\mathsf{o}LM$''. At level $j$ this is the condition
on every cube $\square\subset X$ of side at most $\mathsf{o}LM\cdot L^{j}\xi$. One bound,
$\mathsf{o}LMB^{109}\alpha_0$, is imposed on all these cubes.

\emph{The local class.} Let $D\subset T$ be closed, as in
Definition~\ref{4.4:def-settings-data-domains}, and let $G=SU(N)$. Then
$\mathbf{P}'_D$ is the set of pairs $(\mathbf{U},\mathbf{J})$ on the bonds of $D$ with
the following three properties.
\begin{itemize}
\item $\mathbf{U}=U'U$, where $U$ is $G$-valued and $U'=\exp(i\xi A')$, and $A'$ obeys
\cite[(1.13)]{B-CMP109} on $D$ at the primed radii.
\item $\mathbf{J}$ obeys the current condition of \cite[(1.14)]{B-CMP109} at the primed
radius, namely $|\mathbf{J}|<\gamma_0'$ on $D$.
\item The local form of the gauge condition of \cite[(1.12)]{B-CMP109} holds. That is,
every cube $\square\subset D$ of side at most $M/12$ carries a $G$-valued $u$ with
\begin{equation}\label{4.4:eq-local-data-classes-1}
U^{u}=\exp(i\xi A),\qquad |A|,\ |\nabla^{\xi}A|<\mathsf{a}_U\quad\text{on }\square .
\end{equation}
\end{itemize}
Elements of $\mathbf{P}'_D$ are the \emph{local data on $D$} carried by the settings of
Definition~\ref{4.4:def-settings-data-domains}.

\emph{Units.} Every condition in the definition of $\mathbf{P}'_D$, and every condition in
the further classes of data of this chapter, each at its definition, is
measured in the units of level $k$. In these units the local unit of length equals one on
$\Omega_k$. This applies to the potentials, to their difference quotients, to
$\mathbf{J}$, and to the cube sides.

The plaquette condition \cite[(1.11)]{B-CMP109} and the first condition of
\cite[(1.14)]{B-CMP109} are not imposed. A datum that satisfies them in addition lies a
fortiori in $\mathbf{P}'_D$.

The following hold.
\begin{enumerate}
\item[(a)] The conditions of \cite{B-CMP109} for the operators of step $k$ are written at
level $k+1$, in the units of level $k+1$. They imply the corresponding conditions in the
units of level $k$. So the classes just defined contain the corresponding classes of
\cite{B-CMP109}.
\item[(b)] If $D'\subset D$ is closed, the restriction to $D'$ of local data on $D$ is
local data on $D'$.
\item[(c)] Let $X\supset D$ be a domain. Under the interpretation of the cube size fixed
above, the restriction to $D$ of a point of $\mathbf{P}'$ on $X$ lies in $\mathbf{P}'_D$.
So does the restriction of every representative of the space of configurations of
\cite[(1.11)--(1.14)]{B-CMP109} on $X$.
\item[(d)] Assume the ceilings on the primed radii $\alpha_0',\alpha_1',\gamma_0'$ imposed
after $M$ in the order of choice of
Definitions~\ref{2.3:not-stages-first}--\ref{2.3:not-stages-last}. Then $\mathbf{P}'_D$
is contained in the class
$\mathcal{U}^{\flat}_D(a_1^{\flat},a_0^{\flat})$ of local data on which the inclusion lemma
is stated. It is also contained in the class
$\mathcal{U}^{\flat\nabla}_D(a_1^{\sharp},a_0^{\sharp})$ of local data of its gradient
version. Hence every term whose footprint lies in $D$ is defined on $\mathbf{P}'_D$.
\end{enumerate}
\end{definition}

\begin{proof}
(a) Let a potential, its covariant difference quotient and a current have values $A_t$,
$\nabla^{\xi}A_t$ and $J_t$ in the units of level $k+1$. In the units of level $k$ their
values are
\begin{equation*}
A_k=A_t/L,\qquad \nabla A_k=L^{-2}\nabla^{\xi}A_t,\qquad J_k=L^{-3}J_t .
\end{equation*}
Since $L>1$, none of these moduli increases in passing to the units of level $k$. A bound
imposed at level $k+1$ therefore holds, with the same radius, in the units of level $k$.
This proves the inclusion stated in (a).

(b) The bonds of $D'$ are bonds of $D$. The bound on $A'$ of \cite[(1.13)]{B-CMP109} and
the bound $|\mathbf{J}|<\gamma_0'$ hold on $D'$, because they hold on $D$. A cube
$\square\subset D'$ of side at most $M/12$ is a cube contained in $D$ of side at most
$M/12$. So it carries a gauge $u$ with \eqref{4.4:eq-local-data-classes-1}.

(c) Let $(\mathbf{U},\mathbf{J})$ be a point of $\mathbf{P}'$ on $X$, or a
representative of the space of configurations of \cite[(1.11)--(1.14)]{B-CMP109} on $X$.

First take the conditions on $A'$ and on $\mathbf{J}$. They hold on $D$ because they hold
on $X\supset D$. For the space of configurations of \cite[(1.11)--(1.14)]{B-CMP109}, these
conditions are written in the
units of level $k+1$ and are carried to the units of level $k$ by (a). Its radii are at
most the primed ones; \cite{B-CMP116} takes $\alpha_0',\alpha_1'$ much bigger than
$\alpha_0,\alpha_1$.

Now take the gauge condition. Let $\square\subset D$ be a cube of side at most $M/12$ in
the units of level $k$. Since $\mathsf{o}L\ge 1$,
\begin{equation*}
\frac{M}{12L}\le\frac{M}{12}\le M\le\mathsf{o}LM .
\end{equation*}
Hence $\square$ is itself a cube contained in $X$ of side at most $\mathsf{o}LM$, in the
units of level $k$ and in those of level $k+1$. Under the interpretation of the cube size
fixed above, \cite[(1.12)]{B-CMP109} therefore provides a $G$-valued $u$ on $\square$ with
$U^{u}=\exp(i\xi A)$.
\begin{itemize}
\item For a point of $\mathbf{P}'$ it gives $|A|,|\nabla^{\xi}A|<\mathsf{o}LMB^{109}\alpha_0'$,
which is $\mathsf{a}_U$.
\item For the space of configurations of \cite[(1.11)--(1.14)]{B-CMP109} it gives
$|A|,|\nabla^{\xi}A|<\mathsf{o}LMB^{109}\alpha_0$
in the units of level $k+1$. By (a) this yields a bound by $\mathsf{o}LMB^{109}\alpha_0$ in
the units of level $k$, and $\alpha_0\le\alpha_0'$.
\end{itemize}
In both cases \eqref{4.4:eq-local-data-classes-1} holds on $\square$.

(d) The proof of the inclusion lemma and the proof of its gradient version both proceed
cube by cube. Those proofs use families $\mathcal{Q}_{\mathsf{m}_i}$ of cubes, indexed by
the scales $\mathsf{m}_i$ of those proofs, the hat $\hat q$ attached to each cube $q$, and
the hulls $H_i$. Each proof fixes one hat $\hat q$, its gauge $u$ and its potential $A$,
and computes on $\hat q$. Consequently, membership of a datum on $D$ in
$\mathcal{U}^{\flat}_D(a_1^{\flat},a_0^{\flat})$ or in
$\mathcal{U}^{\flat\nabla}_D(a_1^{\sharp},a_0^{\sharp})$ asks for gauges only on the hats
$\hat q$ of the cubes $q\in\mathcal{Q}_{\mathsf{m}_i}$ and on the hulls $H_i$, all lying
in $D$.
These sets have side at most $3\mathsf{m}_4$. The number $M$ is bounded below by the floor
$M_{\flat}$ on $M$ of the order of choice of
Definitions~\ref{2.3:not-stages-first}--\ref{2.3:not-stages-last}, and
$M\ge M_{\flat}\ge 16L^2\mathsf{m}_4$. Hence
\begin{equation*}
3\mathsf{m}_4\le\frac{16L^2\mathsf{m}_4}{12}\le\frac{M_{\flat}}{12}\le\frac{M}{12}.
\end{equation*}
The first inequality uses that the blocking factor $L$ is an odd integer larger than one,
so that $16L^2/12\ge 12>3$.

Each set on which those proofs require a gauge is therefore a cube contained in $D$ of
side at most $M/12$, and the local gauge condition \eqref{4.4:eq-local-data-classes-1}
supplies one. The remaining requirements are bounds on $A'$ and $\mathbf{J}$ on $D$. Under
the ceilings on the primed radii, the argument of those proofs therefore applies
to every point of $\mathbf{P}'_D$. This gives both inclusions.

The terms whose footprint lies in $D$ are defined on these classes of local data. They are
therefore defined on $\mathbf{P}'_D$.
\end{proof}

Clause (i) of the one-power hereditary class $U^{\Sigma^1}$ is \cite[(1.12)]{B-CMP109}
under the interpretation of the cube size fixed above.
By part (i) of the theorem on the one-power hereditary class, the gauge condition passes
from a domain $Y$ at level $i$ to every $X\subset Y$ at level $i'\le i$. Write $a_0$ for
the first radius of the class $U^{\Sigma^1}_j(X;a_0,a_1,\gamma_0)$, and $B$ for the
constant of its clause (i), as that theorem writes it. The passage is the identity
\begin{quote}
A cube $\square\subset X$ of side at most $\mathsf{o}LM\cdot L^{i'}$ lattice steps is a
cube contained in $Y$ of side at most $\mathsf{o}LM\cdot L^{i}$ lattice steps; by the
interpretation of the cube size, the datum on $Y$ carries a $G$-valued $u$ on $\square$
with $U^{u}=e^{i\xi_iA}$ and $|A|,|\nabla^{\xi_i}A|<\mathsf{o}LMBa_0$; rewritten at the
spacing of level $i'$, $U^{u}=e^{i\xi_{i'}A^{\flat}}$ with
$|A^{\flat}|,|\nabla^{\xi_{i'}}A^{\flat}|<\mathsf{o}LMB(\mu a_0)$, where
$\mu=L^{-(i-i')}$.
\end{quote}
This passage holds at thin $X$ as at every $X$.

Under the interpretation in which the cubes of \cite[(1.12)]{B-CMP109} have side exactly
$\mathsf{o}LM$, no such passage exists at a thin $Y$, for instance a domain one cube wide.
This is why the interpretation ``at most'', under which the gauge condition is hereditary
on $U^{\Sigma^1}$, is adopted both for the class $U^{\Sigma^1}$ and
for $\mathbf{P}'_D$.

Under the exact-size interpretation, statement (c) still holds whenever $\square$ lies in
a cube of the prescribed size contained in $X$. But then \cite[(1.12)]{B-CMP109} is void
on a domain $X$ too thin to contain such a cube. The space of configurations of
\cite[(1.11)--(1.14)]{B-CMP109} then carries
no first-order gauge condition there, and no theorem of \cite{B-CMP99b} concerns it. The
interpretation ``at most'' is the one under which the words ``on the domain $Z$'' of
\cite{B-CMP116} have content.

\begin{definition}[Real regular backgrounds and polydiscs around them]\label{4.4:def-regular-backgrounds}
Let $G=\mathrm{SU}(N)$, let $T$ be the lattice on which the configurations of the current level
are defined, and let $\xi$ be its spacing.
\begin{enumerate}
\item The space $\mathcal{U}_{\mathbb{R}}$ of \emph{real regular backgrounds} consists of the
$G$-valued configurations $U$ on $T$ with the following property: every cube of
side at most $6M$ carries a $G$-valued gauge transformation $u$ such that the transform $U^{u}$
of $U$ by $u$ satisfies $U^{u}=\exp i\xi A$ on that cube, with
\begin{equation*}
|A|<2\mathsf{a}_U,\qquad |\nabla^{\xi}A|<2\mathsf{a}_U .
\end{equation*}
Here $A$ takes values in the Lie algebra of $G$, $\nabla^{\xi}$ is the lattice difference at
spacing $\xi$, and $\mathsf{a}_U>0$ is the regularity constant of this chapter, namely the
constant at which the regularity condition \cite[(1.12)]{B-CMP109} of the primed class
$\mathbf{P}'$ of Definition~\ref{4.4:def-local-data-classes} is taken in (P3) below.
We also say that such a configuration $U$ is first-order regular.
\item For $\sigma>0$, the set $\mathbf{P}_{\sigma}$ consists of the pairs
$(\mathbf{U},\mathbf{J})$ lying within $\sigma\cdot(\alpha_0,\alpha_1)$ of a $G$-valued base
point $(U,0)$ with $U\in\mathcal{U}_{\mathbb{R}}$. Explicitly, the following two conditions
hold.
\begin{itemize}
\item $\mathbf{U}=e^{i\xi A'}U$, where $A'$ obeys the condition \cite[(1.13)]{B-CMP109}
(written \cite[(2.39)]{B-CMP119} in the lettering of that paper) with $\sigma\alpha_1$ in place
of $\alpha_1$;
\item $\mathbf{J}$ obeys the condition \cite[(1.14)]{B-CMP109} (written \cite[(2.36)]{B-CMP119}
in the lettering of that paper) with $\sigma$ times the un-primed radius of the
$\mathbf{J}$-variable in place of that radius.
\end{itemize}
No condition on the plaquette variables $\partial\mathbf{U}$ is part of the definition of
$\mathbf{P}_{\sigma}$.
\item For a data domain $D$, the set $\mathbf{P}_{\sigma,D}$ is defined in the same way on $D$.
The regularity of $U$ is then required in the localised form of \cite[(1.12)]{B-CMP109} on $D$.
\item Put
\begin{equation}\label{4.4:eq-regular-backgrounds-1}
\rho:=\frac{\mathcal{R}_{\mathbf{P}}\,\gamma_0^X}{16\,\mathsf{B}},\qquad
\mathcal{R}_{\mathbf{P}}:=\min\Bigl(\frac{\alpha_0'}{\alpha_0},\,\frac{\alpha_1'}{\alpha_1}\Bigr).
\end{equation}
Here $\mathsf{B}$ and $\gamma_0^X$ are the positive constants so denoted in the order of choice
of the parameters (Definitions~\ref{2.3:not-stages-first}, \ref{2.3:not-stages-middle}
and~\ref{2.3:not-stages-last}); the superscript in $\gamma_0^X$ is part of the name of the
constant and does not denote a large-field region. In this chapter $\rho$ denotes the scale factor
\eqref{4.4:eq-regular-backgrounds-1} of the polydiscs $\mathbf{P}_{\sigma}$, and
$\mathcal{R}_{\mathbf{P}}$ the ratio of radii displayed there; $\rho$ is distinct from the
effective densities $\rho_k$ of Definition~\ref{1.5:def-densities} and from the
contraction letter of clause (v-d).
In the lettering of \cite{B-CMP116} and \cite{B-CMP119}, the ratio $\alpha_0'/\alpha_0$ is
the ratio of the primed to the un-primed radius of the $\mathbf{J}$-variable. In the lettering
of \cite{B-CMP109}, the same ratio is written $\gamma_0'/\gamma_0$ and is read in its place.
\end{enumerate}
The following properties hold.
\begin{itemize}
\item[(P1)] $\mathcal{U}_{\mathbb{R}}$ is invariant under $G$-valued gauge transformations.
\item[(P2)] $\mathcal{U}_{\mathbb{R}}$ contains the identity configuration $1$.
\item[(P3)] $\mathcal{U}_{\mathbb{R}}$ contains the $G$-valued factor $U$ of every point
$\mathbf{U}=e^{i\xi A'}U$ of the primed class $\mathbf{P}'$ of
Definition~\ref{4.4:def-local-data-classes}, the regularity condition \cite[(1.12)]{B-CMP109}
of $\mathbf{P}'$ being taken at the constant $\mathsf{a}_U$.
\item[(P4)] Under the floor
\begin{equation}\label{4.4:eq-regular-backgrounds-2}
\mathcal{R}_{\mathbf{P}}\ \ge\ \frac{32\,\mathsf{B}}{\gamma_0^X},
\end{equation}
imposed in the order of choice of the parameters (Definitions~\ref{2.3:not-stages-first},
\ref{2.3:not-stages-middle} and~\ref{2.3:not-stages-last}), one has $\rho\ge2$. Consequently
$\mathbf{P}_{\rho}$ contains, around every base point, the
polydisc of radii $(\alpha_0,\alpha_1)$ on which Ba{\l}aban's analysis is carried out.
\end{itemize}
The space $\mathcal{U}_{\mathbb{R}}$ coincides with the space of configurations satisfying the
two regularity conditions of the corollary to the operator theorem at free current, taken with
the cube side $6M$ and the bound $2\mathsf{a}_U$ of the present definition. The first of these
conditions, at the constant $a_{\mathrm{reg}}^{\oplus}$ of that corollary, holds by restriction
to the cubes on which that corollary imposes it, under the auxiliary condition stated there. The
second condition is the condition on the cubes of side $6M$ above.
Of $\mathcal{U}_{\mathbb{R}}$, the arguments of this chapter
use only the following:
\begin{itemize}
\item the first-order bound \cite[(2.38)]{B-CMP119};
\item property (P1);
\item property (P2).
\end{itemize}
The sets defined here are used as follows in the rest of this chapter:
\begin{itemize}
\item the basic analyticity estimate is proved on $\mathbf{P}_{2\rho}$;
\item items (ii) and (iii) of the tower lemma, clauses (c2) and (c5) of the response
proposition, and clause (d) of the principal theorem of this chapter are read on
$\mathbf{P}_{\rho}$;
\item item (v) of the tower lemma and the lemma on localised data are read on
$\mathbf{P}_{\rho,D}$.
\end{itemize}
The set $\mathbf{P}_{\rho,D}$ is the polydisc on which the localisation lemma of the
cluster-expansion step is applied.
\end{definition}

\begin{proof}[Proof of properties (P1)--(P4)]
(P1). Let $U\in\mathcal{U}_{\mathbb{R}}$ and let $v$ be a $G$-valued gauge transformation. Fix
a cube of side at most $6M$. By hypothesis it carries a $G$-valued $u$ with $U^u=\exp i\xi A$
and $|A|,|\nabla^{\xi}A|<2\mathsf{a}_U$. Let $v'$ be the gauge transformation on the cube
obtained by transforming first by the inverse of $v$ and then by $u$; it is $G$-valued because
$v$ and $u$ are, and transforming first by $v$ and then by $v'$ is the same as transforming
by $u$. Hence $(U^{v})^{v'}=U^{u}=\exp i\xi A$, with the same $A$. Thus the cube carries a
transformation as required for $U^{v}$. The cube was arbitrary, so $U^{v}\in\mathcal{U}_{\mathbb{R}}$.

(P2). For $U=1$, take $u=1$ and $A=0$ on every cube. Then both bounds hold because
$\mathsf{a}_U>0$.

(P3). Let $\mathbf{U}=e^{i\xi A'}U$ be a point of $\mathbf{P}'$. Its $G$-valued factor $U$
satisfies the regularity condition of \cite[(1.12)]{B-CMP109} at the constant $\mathsf{a}_U$.
That condition is imposed on every class cube of Definition~\ref{4.4:def-local-data-classes},
the class cubes being the cubes of side $\mathsf{o}L^2M$ in units of the level-$k$ lattice.
Taken at the constant $\mathsf{a}_U$, it provides on each class cube a $G$-valued gauge
transformation $u$ and a field $A$ with values in the Lie algebra of $G$ such that
$U^u=\exp i\xi A$ and $|A|,|\nabla^{\xi}A|<2\mathsf{a}_U$ on that class cube.
Since $\mathsf{o}L^2M\ge 6M$, each cube of side at most $6M$ is contained in one of them.
Restricting to it the gauge transformation $u$ and the field $A$ of that class cube
gives a $G$-valued $u$ with $U^u=\exp i\xi A$ and both bounds on the smaller cube, since
pointwise bounds persist on restriction. Hence $U\in\mathcal{U}_{\mathbb{R}}$.

(P4). By \eqref{4.4:eq-regular-backgrounds-1} and \eqref{4.4:eq-regular-backgrounds-2},
\begin{equation*}
\rho=\frac{\mathcal{R}_{\mathbf{P}}\,\gamma_0^X}{16\,\mathsf{B}}
\ \ge\ \frac{32\,\mathsf{B}}{\gamma_0^X}\cdot\frac{\gamma_0^X}{16\,\mathsf{B}}=2 .
\end{equation*}
The conditions defining $\mathbf{P}_{\sigma}$ around a fixed base point bound $A'$ by a multiple
of $\sigma\alpha_1$ and $\mathbf{J}$ by a multiple of $\sigma$ times the un-primed
$\mathbf{J}$-radius. They therefore weaken as $\sigma$ increases, so $\sigma\le\sigma'$ implies
$\mathbf{P}_{\sigma}\subset\mathbf{P}_{\sigma'}$. Since $\rho\ge2>1$, we get
$\mathbf{P}_1\subset\mathbf{P}_{\rho}$. Around every base point, $\mathbf{P}_1$ is the polydisc
of radii $(\alpha_0,\alpha_1)$.
\end{proof}

\begin{definition}[The quadratic form and the compressed operators]\label{4.4:def-quadratic-form}
Fix the level $k$ and a complexified configuration $(\mathbf{U},\mathbf{J})$ as in
Definition~\ref{4.4:def-settings-data-domains}. Throughout, the operators named below without
further comment carry Ba{\l}aban's notation and meaning from \cite[\S\S A--D]{B-CMP99b}, in the
form of \cite[(3.156)]{B-CMP99b} ``with the $\delta$-function gauge fixing term replaced by the
exponential one'', and from \cite{B-CMP109} and \cite[(1.2)--(1.6)]{B-CMP116}. In particular
$D$, $R$, $P$, $Q'$, $G'$ and $a$ in item~(b) are the operators of \cite[\S\S A--D]{B-CMP99b},
and $T$ in item~(b) is the operator defined there; they are not the local data $D$ or the
letter $T$ of Definition~\ref{4.4:def-settings-data-domains}, nor the ball radius $R$. Likewise
$H$, $H_0$ and $H_1$ are the operators of \cite[\S\S A--D]{B-CMP99b} and
\cite[(1.2)--(1.4)]{B-CMP116}, not the Hamiltonian $H$. In
Ba{\l}aban's words \cite{B-CMP119}, when the geometry changes ``only the symbols represent
different operators, determined by the sequence $\{\Omega_j\}$''.

\begin{enumerate}
\item[(a)] \emph{Two geometries.} The geometry $\mathcal{S}$ is Ba{\l}aban's setting following
\cite[(3.20)--(3.25)]{B-CMP119}:
\begin{equation*}
\mathcal{S} = (\boldsymbol{\Omega}, Y), \qquad
\boldsymbol{\Omega} = \{\Omega_j\}_{j \le k+1}, \qquad Y = \mathbb{B}(\Lambda_{k+1}).
\end{equation*}
Here $Y$ is the set of integration bonds of ``the conditional integral with respect to $A_k$
restricted to $\Lambda_{k+1}$'' \cite{B-CMP119}. Let $T_\eta$ be the lattice torus of
\cite{B-CMP119}. The
geometry $\mathcal{S}_0$ is
\begin{equation*}
\mathcal{S}_0 = \bigl((T_\eta, \dots, T_\eta),\ \mathbb{B}(T_\eta)\bigr),
\end{equation*}
that is, $\Omega_j = \Lambda_j = T_\eta$ for $j = 1, \dots, k+1$, which is the case treated in
\cite{B-CMP119} ``assuming $\Omega_j = \Lambda_j = T_\eta$''. In $\mathcal{S}_0$ the small-field
thresholds are those of \cite{B-CMP119}; this is how the phrase ``in the framework of [I]''
of \cite{B-CMP119} (there [I] is \cite{B-CMP109}) is read throughout this chapter. Both
$\mathcal{S}$ and $\mathcal{S}_0$ are geometries in the sense
of Definition~\ref{4.4:def-settings-data-domains}. No datum is part of them: a datum, global or
local, is attached only when a quantity is evaluated.

\item[(b)] \emph{Sequence operators.} The list $\mathbb{L}_1$ below and the list
$\mathbb{L}_2$ of item~(d) are parts of the list $\mathbb{L}$ of
Definition~\ref{4.4:def-settings-data-domains}. The list $\mathbb{L}_1$ consists of operators
determined
by $(\Omega_j)_{j \le k}$, by the block structure and by $(\mathbf{U},\mathbf{J})$, and not by
$Y$. Its members are
\begin{gather*}
G', \quad T := (Q' G'^{2} Q'^{*})^{-1}, \quad P, \quad R, \quad
G_0 := (\Delta + D R D^{*} + Q^{*} a Q)^{-1}, \\
G, \quad G_1, \quad \Delta^{(2)}[\mathfrak{j}_0], \quad
(Q G_0 Q^{*})^{-1}, \quad (Q G Q^{*})^{-1}, \quad (Q G_1 Q^{*})^{-1}, \quad H, \quad H_1,
\end{gather*}
where $\Delta^{(2)}[\mathfrak{j}_0]$ denotes the operator $\Delta^{(2)}$ evaluated at
$\mathfrak{j}_0$, together with the following:
\begin{itemize}
\item the operators $H_0$ and $\tilde G$ of \cite[(1.2)--(1.4)]{B-CMP116}, compare
\cite[(129)--(131), (143)]{B-CMP102b};
\item the operator $\Delta^{(k)}$;
\item the form $\mathsf{A}$ of item~(c).
\end{itemize}
The list $\mathbb{L}_1$ also contains the analogues of $H$, $H_0$, $H_1$ and $\tilde G$ at
every level $j \le k$, that is, the same operators built from the truncated sequence
$(\Omega_i)_{i \le j}$. These are the functions
\begin{equation*}
\mathbf{H}_j, \qquad \frac{\delta}{\delta B}\mathbf{H}_j
\end{equation*}
of \cite[(3.6)--(3.9)]{B-CMP109}, where $B$ is the variable of $\mathbf{H}_j$ there with
respect to which the derivative is taken, with the decoupling parameters of the cluster
expansion of \cite{B-CMP116} inserted. The operators $\mathfrak{P}$ and
$\mathfrak{G}$ of \cite[(3.147), (3.153)]{B-CMP99b} are not members of $\mathbb{L}_1$;
$\mathfrak{P}$ contains the composite $G_1 D R D^{T}$ of \cite[(3.152)]{B-CMP99b}, in which the
transpose $D^{T}$ occurs bare, and neither operator is used in what follows.

\item[(c)] \emph{The quadratic form.} Let $C$ be the operator of the block structure carried by
$\Omega_{k+1}$ (Definition~\ref{4.4:def-settings-data-domains}). Set
\begin{equation}\label{4.4:eq-quadratic-form-1}
\mathsf{A} := C^{*} \Delta^{(k)} C
\qquad \text{on } \ell^2\bigl(\mathbb{B}(\Omega_{k+1}); \mathfrak{g}^c\bigr).
\end{equation}
In Ba{\l}aban's exponential block gauge \cite{B-CMP109}, $\mathsf{A}$ is the quadratic form at
the free source field, the current being treated as a free variable and $\mathbf{J}$ substituted
for it.

\item[(d)] \emph{Compressed operators.} The admissible bond sets are
\begin{equation*}
Y' \in \{Y\} \cup \bigl\{\mathbb{B}(Z_0) : Z_0 \in \mathbf{D}_k,\ Z_0 \subset
\text{the } Y\text{-region}\bigr\},
\end{equation*}
where the $Y$-region is taken in the sense of Definition~\ref{4.4:def-settings-data-domains}
(in the geometry $\mathcal{S}$, where $Y = \mathbb{B}(\Lambda_{k+1})$, it is $\Lambda_{k+1}$).
For such $Y'$, let $\chi_{Y'}$ be multiplication by the indicator function of $Y'$, and set
\begin{equation*}
\mathsf{A}^{Y'} := \chi_{Y'} \mathsf{A} \chi_{Y'},
\end{equation*}
regarded as an operator on $\ell^2(Y';\mathfrak{g}^c)$.
The list $\mathbb{L}_2$ consists of the operators
\begin{equation}\label{4.4:eq-quadratic-form-2}
\mathsf{R}^{Y'}(x) := (x + \mathsf{A}^{Y'})^{-1} \quad (x \ge 0), \qquad
C^{(k)}(Y') = \mathsf{R}^{Y'}(0),
\end{equation}
together with
\begin{equation}\label{4.4:eq-quadratic-form-3}
\mathsf{M}^{Y} := (\mathsf{A}^{Y})^{-1/2}
= \frac{1}{\pi} \int_0^\infty x^{-1/2}\, \mathsf{R}^{Y}(x)\, dx .
\end{equation}
In \eqref{4.4:eq-quadratic-form-2}, $C^{(k)}(Y')$ is Ba{\l}aban's propagator of
\cite[(2.5)]{B-CMP116}, compressed to $Y'$; item~(e) lists the cases. On the reduced variables
it is
\begin{equation*}
C^{(k)}(Y') = (C^{*} \Delta_k C)^{-1},
\end{equation*}
with $\Delta_k$ the operator so denoted in \cite[(3.158)]{B-CMP99b}, where this identity is
stated; compare \cite{B-CMP109}.
In \eqref{4.4:eq-quadratic-form-3}, $\pi$ denotes the number $\pi$ and not the projection onto
$\mathfrak{g}^c$. The root in \eqref{4.4:eq-quadratic-form-3} is given by the integral on its
right-hand side, as in \cite[(2.5)--(2.7)]{B-CMP116}. When $\mathsf{A}^{Y}$ is self-adjoint and
positive definite, this integral coincides with the spectral root $(\mathsf{A}^{Y})^{-1/2}$.
This follows from the spectral theorem and the scalar identity
\begin{equation*}
\int_0^\infty x^{-1/2} (x + a)^{-1}\, dx = \pi a^{-1/2} \qquad (a > 0),
\end{equation*}
which follows from the substitution $x = a t^2$ together with
$\int_0^\infty 2 (1 + t^2)^{-1}\, dt = \pi$.

\item[(e)] \emph{Ba{\l}aban's operators in this notation.} For a region $Z$ write
$C^{(k)}(Z)$ for $C^{(k)}(\mathbb{B}(Z))$, and write $\mathsf{A}_{\Lambda\Lambda}$ for
$\mathsf{A}^{Y}$ in the geometry $\mathcal{S}$. In particular the following five operators of
Ba{\l}aban's construction are members of $\mathbb{L}_2$, all with $Y' = Y$ in the geometry
$\mathcal{S}$ except $C^{(k)}(Z_0)$, for which $Y' = \mathbb{B}(Z_0)$:
\begin{itemize}
\item $C^{(k)}(\Lambda_{k+1})$ and its square root $\mathsf{M}^{Y}$. These are the operators
of \cite[(2.5)--(2.7)]{B-CMP116}, taken ``restricted to $\Lambda_{k+1}$'' as in
\cite{B-CMP119}.
\item $C^{(k)}(Z_0)$ of \cite[(2.5)]{B-CMP116}.
\item $(x + \mathsf{A}_{\Lambda\Lambda})^{-1}$ of \cite[(2.7)]{B-CMP116}.
\item $(\mathsf{A}_{\Lambda\Lambda} + \lambda)^{-1}$ of \cite[(3.36)--(3.37)]{B-CMP119}, with
$\lambda$ the parameter so denoted in that display; for $\lambda \ge 0$ it is
$\mathsf{R}^{Y}(\lambda)$.
\end{itemize}
\end{enumerate}
\end{definition}

\begin{definition}[Walk systems, footprints and the weighted norm]\label{4.4:def-walk-systems}
Let $K$ be an operator. A \emph{walk system} for $K$ is a family $(K_\omega)_{\omega\in\mathcal{W}}$
of operators indexed by a set $\mathcal{W}$ of walks, such that
\begin{equation*}
K=\sum_{\omega\in\mathcal{W}}K_\omega .
\end{equation*}
Here a \emph{walk} $\omega$ is a finite arrangement of \emph{steps}, linear or tree-like as in
\cite{B-CMP99b}. Each step carries a closed set, its \emph{domain}, and the domains of consecutive
steps intersect. The \emph{footprint} $|\omega|$ of $\omega$ is the union of the domains of its
steps.

Let $K$ be an operator on the unit lattice with a walk system $(K_\omega)_{\omega\in\mathcal{W}}$,
and write $K_\omega(b,b')$ for the kernel of $K_\omega$ on bonds $b,b'$.
Lengths are measured in the conventions of Notation~\ref{4.4:not-units-distances}. We put
\begin{itemize}
\item $d(\omega;b,b')$: the length of a shortest tree that contains the initial points $b_-$ and
$b'_-$ of the bonds $b$ and $b'$ and meets every domain of $\omega$;
\item $d(\omega)$: the length of a shortest tree that meets every domain of $\omega$.
\end{itemize}
For $\theta\ge 0$, the \emph{weighted norm} of the walk system is
\begin{equation}\label{4.4:eq-walk-systems-1}
\begin{split}
\lvert\!\lvert\!\lvert K\rvert\!\rvert\!\rvert_{\theta}
:=\max\Big\{&\sup_{b}\sum_{\omega}\sum_{b'}|K_\omega(b,b')|\,e^{\theta d(\omega;b,b')},\\
&\sup_{b'}\sum_{\omega}\sum_{b}|K_\omega(b,b')|\,e^{\theta d(\omega;b,b')}\Big\}.
\end{split}
\end{equation}
The quantity \eqref{4.4:eq-walk-systems-1} is the row-wise norm of this chapter, taken over the
pairs $(\omega,b')$ at fixed $b$ and over the pairs $(\omega,b)$ at fixed $b'$.
These definitions, and the lemma that follows, apply without change to any family of kernels indexed
by walks, whether or not the family sums to an operator.
\end{definition}

\begin{lemma}
Let $\theta\ge 0$, and let $(K_\omega)_{\omega\in\mathcal{W}}$ and $(K'_{\omega'})_{\omega'\in\mathcal{W}'}$
be walk systems for unit-lattice operators $K$ and $K'$.
\begin{enumerate}
\item[(a)] For a pair $(\omega,\omega')$, let $\omega\omega'$ be the arrangement of the steps of
$\omega$ followed by those of $\omega'$. The family $(K_\omega K'_{\omega'})_{(\omega,\omega')}$ sums
to $KK'$, and, in the sense of the last sentence of Definition~\ref{4.4:def-walk-systems}, its
weighted norm satisfies
\begin{equation*}
\lvert\!\lvert\!\lvert KK'\rvert\!\rvert\!\rvert_{\theta}
\le\lvert\!\lvert\!\lvert K\rvert\!\rvert\!\rvert_{\theta}\,
\lvert\!\lvert\!\lvert K'\rvert\!\rvert\!\rvert_{\theta}.
\end{equation*}
\item[(b)] $\|K\|_{\ell^2\to\ell^2}\le\lvert\!\lvert\!\lvert K\rvert\!\rvert\!\rvert_{0}$.
\item[(c)] For all $\omega$, $b$ and $b'$,
\begin{equation*}
|K_\omega(b,b')|\le\lvert\!\lvert\!\lvert K\rvert\!\rvert\!\rvert_{\theta}\,e^{-\theta d(\omega;b,b')},
\qquad d(\omega;b,b')\ge\max\{|b-b'|,\,d(\omega)\}.
\end{equation*}
\end{enumerate}
\end{lemma}

\begin{proof}
(a) The kernel of $K_\omega K'_{\omega'}$ is $\sum_{b''}K_\omega(b,b'')K'_{\omega'}(b'',b')$.

First we compare the tree lengths. Take a shortest tree for $d(\omega;b,b'')$ and a shortest tree for
$d(\omega';b'',b')$. Both contain $b''_-$, so we may join them at $b''_-$. The result contains $b_-$
and $b'_-$. It meets every domain of $\omega$ and every domain of $\omega'$, hence every domain of
$\omega\omega'$. Its length is at most the sum of the two lengths. Therefore
\begin{equation}\label{4.4:eq-walk-systems-2}
d(\omega\omega';b,b')\le d(\omega;b,b'')+d(\omega';b'',b').
\end{equation}

Now fix $b$. Since $\theta\ge 0$, the inequality \eqref{4.4:eq-walk-systems-2} gives
\begin{equation*}
\begin{split}
\sum_{\omega,\omega'}\sum_{b'}\big|(K_\omega K'_{\omega'})(b,b')\big|\,e^{\theta d(\omega\omega';b,b')}
\le\sum_{\omega,b''}&|K_\omega(b,b'')|\,e^{\theta d(\omega;b,b'')}\\
&\times\sum_{\omega',b'}|K'_{\omega'}(b'',b')|\,e^{\theta d(\omega';b'',b')}.
\end{split}
\end{equation*}
We sum first over $(\omega',b')$. The inner sum is at most
$\lvert\!\lvert\!\lvert K'\rvert\!\rvert\!\rvert_{\theta}$, whatever $b''$ is. We then sum over
$(\omega,b'')$, and the remaining sum is at most $\lvert\!\lvert\!\lvert K\rvert\!\rvert\!\rvert_{\theta}$.

For fixed $b'$ the argument is the same with the roles of the two factors exchanged. We sum first over
$(\omega,b)$, which gives at most $\lvert\!\lvert\!\lvert K\rvert\!\rvert\!\rvert_{\theta}$. We then sum
over $(\omega',b'')$, which gives at most $\lvert\!\lvert\!\lvert K'\rvert\!\rvert\!\rvert_{\theta}$.

(b) By the triangle inequality, $|K(b,b')|\le\sum_\omega|K_\omega(b,b')|$. Hence every row sum
$\sum_{b'}|K(b,b')|$ and every column sum $\sum_{b}|K(b,b')|$ is at most
$\lvert\!\lvert\!\lvert K\rvert\!\rvert\!\rvert_{0}$.

We now apply the Schur test. For $f\in\ell^2$, the Cauchy--Schwarz inequality gives
\begin{equation*}
\begin{split}
\Big|\sum_{b'}K(b,b')f(b')\Big|^2
&\le\Big(\sum_{b'}|K(b,b')|\Big)\Big(\sum_{b'}|K(b,b')|\,|f(b')|^2\Big)\\
&\le\lvert\!\lvert\!\lvert K\rvert\!\rvert\!\rvert_{0}\sum_{b'}|K(b,b')|\,|f(b')|^2 .
\end{split}
\end{equation*}
We sum over $b$ and use the column bound:
\begin{equation*}
\|Kf\|^2\le\lvert\!\lvert\!\lvert K\rvert\!\rvert\!\rvert_{0}\sum_{b'}|f(b')|^2\sum_b|K(b,b')|
\le\lvert\!\lvert\!\lvert K\rvert\!\rvert\!\rvert_{0}^2\,\|f\|^2 .
\end{equation*}

(c) Fix $\omega$, $b$ and $b'$. The term $|K_\omega(b,b')|\,e^{\theta d(\omega;b,b')}$ is one of the
terms of the row sum in \eqref{4.4:eq-walk-systems-1} at this $b$. All terms of that sum are
non-negative, so this term is at most $\lvert\!\lvert\!\lvert K\rvert\!\rvert\!\rvert_{\theta}$, which is
the first claim.

For the second claim, consider a tree admissible for $d(\omega;b,b')$. It contains a path from $b_-$
to $b'_-$. The length of that path is at least the distance between its endpoints, which is $|b-b'|$.
The same tree also meets every domain of $\omega$, so it is admissible for $d(\omega)$, and its length
is at least $d(\omega)$.
\end{proof}

\begin{definition}[Locality of a walk-expansion rule]\label{4.4:def-local-rules}
A \emph{rule} $K$ assigns to every geometry $\mathfrak{g}=(\boldsymbol{\Omega},Y)$ at level $k$,
consisting of a nest $\boldsymbol{\Omega}$ and a set $Y$ of bonds, the following data:
\begin{itemize}
\item an index set $\mathcal{W}_K(\mathfrak{g})$;
\item for every $\omega\in\mathcal{W}_K(\mathfrak{g})$, its steps and its footprint $|\omega|$, in the
sense of Definition~\ref{4.4:def-walk-systems};
\item for every such $\omega$, a function $\mathsf{u}\mapsto K_\omega[\mathfrak{g}](\mathsf{u})$ on the
local class $\mathbf{P}'_{|\omega|}$ of Definition~\ref{4.4:def-local-data-classes}; for a rule of the
class $\mathbb{L}_2$, whose terms at global data are described in the remark following this
definition, the domain is instead the local polydisc $\mathbf{P}_{\rho,|\omega|}$ at radius $\rho$,
one of the polydiscs of Definition~\ref{4.4:def-regular-backgrounds}.
\end{itemize}
For a setting $\mathfrak{s}=(\mathfrak{g};D,\mathsf{u})$, that is, a geometry $\mathfrak{g}$ together
with a domain $D$ and a datum $\mathsf{u}$ on $D$, and a footprint
$|\omega|\subset D$ we put
\begin{equation*}
K_\omega[\mathfrak{s}]:=K_\omega[\mathfrak{g}]\bigl(\mathsf{u}|_{|\omega|}\bigr).
\end{equation*}

Let $\mathfrak{g}=(\boldsymbol{\Omega},Y)$ and $\mathfrak{g}'=(\boldsymbol{\Omega}',Y')$ be two geometries.
They may lie on one torus. They may also lie on two tori, in which case $W$ below is read on a
common lift of the two tori. A closed union $W$
of $M_1$-cubes is \emph{uniform for} $(\mathfrak{g},\mathfrak{g}')$ if the following two conditions
hold.
\begin{itemize}
\item $W\subset\Omega_{k+1}\cap\Omega'_{k+1}$, where $\Omega_{k+1}$ and $\Omega'_{k+1}$ are the
level-$(k+1)$ members of the nests $\boldsymbol{\Omega}$ and $\boldsymbol{\Omega}'$.
\item $Y\cap W=Y'\cap W$. This equality means that $Y$ and $Y'$ contain the same bonds among those
bonds that lie in $W$ at distance $\ge 1$ from $\partial W$. By condition \textup{(iii)} below, the
support property of $K_\omega$, no other bond enters.
\end{itemize}
Let the settings be $\mathfrak{s}=(\mathfrak{g};D,\mathsf{u})$ and
$\mathfrak{s}'=(\mathfrak{g}';D',\mathsf{u}')$, and write their data as
$\mathsf{u}=(\mathbf{U},\mathbf{J})$ and $\mathsf{u}'=(\mathbf{U}',\mathbf{J}')$. The set $W$ is
\emph{uniform for} $(\mathfrak{s},\mathfrak{s}')$ if it is uniform for $(\mathfrak{g},\mathfrak{g}')$,
if $W\subset D\cap D'$, and if
\begin{equation*}
(\mathbf{U},\mathbf{J})|_W=(\mathbf{U}',\mathbf{J}')|_W .
\end{equation*}
In that case we say that $\mathfrak{s}$ and $\mathfrak{s}'$ have the same \emph{trace on $W$}.

A rule $K$ is \emph{local} if the following three conditions hold for all geometries
$\mathfrak{g},\mathfrak{g}'$ and every $W$ uniform for $(\mathfrak{g},\mathfrak{g}')$.
\begin{itemize}
\item[\textup{(i)}] The two index sets agree on walks inside $W$:
\begin{equation*}
\{\omega\in\mathcal{W}_K(\mathfrak{g}):|\omega|\subset W\}
=\{\omega\in\mathcal{W}_K(\mathfrak{g}'):|\omega|\subset W\}.
\end{equation*}
\item[\textup{(ii)}] For each $\omega$ in the set of \textup{(i)}, the functions $K_\omega[\mathfrak{g}]$
and $K_\omega[\mathfrak{g}']$ are equal as functions on the local class of $|\omega|$.
\item[\textup{(iii)}] The support property of $K_\omega$: for every $\omega$, the kernel of $K_\omega$
vanishes unless both of its arguments lie
in $|\omega|$ at distance $\ge 1$ from $\partial|\omega|$.
\end{itemize}
This is the locality condition for walk systems, with one change: the
traces of $\boldsymbol{\Omega}$ and of $Y$ on $W$ are added to the data.
\end{definition}

For a local rule $K$, a setting $\mathfrak{s}=(\mathfrak{g};D,\mathsf{u})$, and a set $W\subset D$, put
\begin{equation}\label{4.4:eq-local-rules-1}
K^W[\mathfrak{s}]:=\sum_{\omega\in\mathcal{W}_K(\mathfrak{g}),\ |\omega|\subset W}K_\omega[\mathfrak{s}].
\end{equation}
Then $K^W[\mathfrak{s}]$ depends on $\mathfrak{s}$ only through its trace on $W$. That is, if $W$ is
a closed union of $M_1$-cubes uniform for $(\mathfrak{s},\mathfrak{s}')$, then
$K^W[\mathfrak{s}]=K^W[\mathfrak{s}']$, and the two sums \eqref{4.4:eq-local-rules-1} agree term by term.

\begin{proof}[Proof that {$K^W[\mathfrak{s}]$} depends on {$\mathfrak{s}$} only through its trace on {$W$}]
Let $W$ be uniform for $(\mathfrak{s},\mathfrak{s}')$. Then $W$ is in particular uniform for
$(\mathfrak{g},\mathfrak{g}')$, so condition \textup{(i)} of the definition applies. It shows that the
sums defining $K^W[\mathfrak{s}]$
and $K^W[\mathfrak{s}']$ run over the same set of indices $\omega$.

Fix such an $\omega$. Since $W\subset D\cap D'$, we have $|\omega|\subset D$ and $|\omega|\subset D'$, so
both $K_\omega[\mathfrak{s}]$ and $K_\omega[\mathfrak{s}']$ are defined. Since
$(\mathbf{U},\mathbf{J})|_W=(\mathbf{U}',\mathbf{J}')|_W$ and $|\omega|\subset W$, the two data restrict
to the same element:
\begin{equation*}
\mathsf{u}|_{|\omega|}=\mathsf{u}'|_{|\omega|}.
\end{equation*}
Unwinding the definition of $K_\omega[\,\cdot\,]$ at a setting, and then applying condition
\textup{(ii)}, we obtain
\begin{align*}
K_\omega[\mathfrak{s}]
&=K_\omega[\mathfrak{g}]\bigl(\mathsf{u}|_{|\omega|}\bigr)
=K_\omega[\mathfrak{g}']\bigl(\mathsf{u}|_{|\omega|}\bigr)\\
&=K_\omega[\mathfrak{g}']\bigl(\mathsf{u}'|_{|\omega|}\bigr)
=K_\omega[\mathfrak{s}'].
\end{align*}
Summing over the common index set gives $K^W[\mathfrak{s}]=K^W[\mathfrak{s}']$, term by term.

It remains to justify the reading of $Y\cap W=Y'\cap W$ used in the definition. Let $|\omega|\subset W$,
and let $b$ be an argument at which the kernel of $K_\omega$ does not vanish. By the support property
\textup{(iii)}, $b$ lies in
$|\omega|$, hence in $W$, at distance $\ge 1$ from $\partial|\omega|$. Suppose that some point
$x\in\partial W$ lies at distance $<1$ from $b$, and let $y$ be a point of the bond $b$ with
$|y-x|<1$. If $x\in|\omega|$, then $x\in\partial|\omega|$, because
$|\omega|\subset W$ and a point interior to $|\omega|$ is interior to $W$; this contradicts the
support property \textup{(iii)}.
If $x\notin|\omega|$, then the straight segment from $y$ to $x$, of length $<1$, joins a point of
$|\omega|$ to a point outside $|\omega|$ and therefore meets $\partial|\omega|$ at a point at distance
$<1$ from $y$, hence at distance $<1$ from $b$; this again contradicts the support property
\textup{(iii)}. Hence $b$ lies in $W$
at distance $\ge 1$ from $\partial W$.
\end{proof}

\begin{remark}
The definition of uniformity for $(\mathfrak{s},\mathfrak{s}')$ compares the data of $\mathfrak{s}$ and
$\mathfrak{s}'$ on $W$ only. It does not ask for an extension of a local datum to the whole lattice, and
the definition is stated in this form for the following reason. A set $W$ as in the definition, lying
in the small-field region $\Lambda_{k+1}$, may be
multiply connected, its holes being large-field regions. Suppose the restriction of $U$ to $W$ is flat,
with holonomy about a hole lying in the centre of $\mathrm{SU}(N)$. Then an extension across the hole
needs flux
at least of order one at curvature below $\alpha_0'\xi^2$, where $\xi$ is the lattice spacing at the
level read and $\alpha_0'\xi^2$ is the curvature bound that the extended datum would have to obey. This
is a lower bound on the product of $\alpha_0'$ with
the area of the hole, whereas $\alpha_0'$ is subject to upper bounds only. Such an extension need
therefore not exist, and in the form above nothing is extended.

For a rule of the class $\mathbb{L}_2$ and a global datum $\mathsf{u}\in\mathbf{P}_\rho$, the term at
$\mathsf{u}$ is the one that the lemma constructing the terms of these rules assigns to a global datum.
In general
$\mathbf{P}_\rho|_W\not\subset\mathbf{P}_{\rho,W}$. The reason is that the space $\mathcal{U}_{\mathbb{R}}$
of Definition~\ref{4.4:def-regular-backgrounds} is defined with the bound $2\mathsf{a}_U$, while the
local gauge condition of Definition~\ref{4.4:def-local-data-classes} is imposed with the bound
$\mathsf{a}_U$. Some of the data to which these terms are applied are obtained by extension; they are
global points of this kind, lying in
$\mathbf{P}_{2\rho}$. Where both the global term and the local function are defined, they agree, by
the same lemma. Every later use of these terms takes its base configuration in the class
at the bound $\mathsf{a}_U$.
\end{remark}

\begin{definition}[Interpolation of walk expansions over cubes]\label{4.4:def-interpolation-rule}
Let $T$ be the region on which the walks live, and let $\sigma_0$ be a family of closed cubes
$\Delta\subset T$. Let
\begin{equation*}
K=\sum_{\omega}K_{\omega}
\end{equation*}
be a walk expansion. The sum runs over its family of walks $\omega$, and each walk has a
footprint $|\omega|\subset T$. We assume throughout that every footprint is the closure of an
open subset of $T$. For the families of walks constructed in this chapter, the statement that
constructs them makes the footprint of a walk the closure of its own interior, and that
interior is connected; so the assumption holds for them.

A cube $\Delta\in\sigma_0$ \emph{meets} $|\omega|$ if
$\operatorname{int}\Delta\cap|\omega|\neq\emptyset$. Any fixed convention would serve equally
well; this one, which is the convention adopted in this chapter for the cluster expansion of
\cite{B-CMP116}, is used in every setting of the chapter. For an assignment $s$ of a number
$s(\Delta)$ to each cube $\Delta\in\sigma_0$, put
\begin{equation}\label{4.4:eq-interpolation-rule-1}
K(s):=\sum_{\omega}K_{\omega}\prod_{\Delta\in\sigma_0:\ \Delta\text{ meets }|\omega|}s(\Delta),
\end{equation}
where the empty product is $1$.

For a closed region $W\subset T$ write $K^{W}$ for the sum of the terms $K_{\omega}$ over the
walks with $|\omega|\subset W$. For $\sigma'\subset\sigma_0$ let $1_{\sigma'}$ be the
indicator function of $\sigma'$ on $\sigma_0$, and put
\begin{equation*}
W(\sigma'):=\overline{T\setminus\textstyle\bigcup_{\Delta\in\sigma_0\setminus\sigma'}\Delta}.
\end{equation*}
Then
\begin{equation}\label{4.4:eq-interpolation-rule-2}
K(1_{\sigma'})=K^{W(\sigma')}.
\end{equation}
\end{definition}

\begin{proof}
Fix a walk $\omega$. By hypothesis $|\omega|=\overline{O}$ for some open set $O\subset T$. Fix
also a cube $\Delta\in\sigma_0$; it is closed, and it is the closure of its interior.

We first show that $\Delta$ meets $|\omega|$ if and only if $O\cap\Delta\neq\emptyset$. The set
$\operatorname{int}\Delta$ is open. An open set meets the closure of $O$ exactly when it meets
$O$, so $\operatorname{int}\Delta\cap|\omega|\neq\emptyset$ if and only if
$\operatorname{int}\Delta\cap O\neq\emptyset$.

Next, $O$ misses $\operatorname{int}\Delta$ if and only if it misses $\Delta$. One direction
holds because $\operatorname{int}\Delta\subset\Delta$. For the other, let
$x\in O\cap\Delta$. Since $x$ lies in the closure of $\operatorname{int}\Delta$, its open
neighbourhood $O$ meets $\operatorname{int}\Delta$. Together the two statements give the
claim.

Put $F:=\bigcup_{\Delta\in\sigma_0\setminus\sigma'}\Delta$, so that
$W(\sigma')=\overline{T\setminus F}$. We claim that no cube of $\sigma_0\setminus\sigma'$
meets $|\omega|$ if and only if $|\omega|\subset W(\sigma')$.

Suppose first that no cube of $\sigma_0\setminus\sigma'$ meets $|\omega|$. By the preceding
step $O$ misses every such cube, so $O\subset T\setminus F$. Taking closures gives
$|\omega|\subset W(\sigma')$.

Conversely, suppose $|\omega|\subset W(\sigma')$, and suppose some
$\Delta\in\sigma_0\setminus\sigma'$ meets $|\omega|$. By the first step there is a point
$x\in O\cap\operatorname{int}\Delta$. The set $O\cap\operatorname{int}\Delta$ is an open
neighbourhood of $x$ contained in $F$, so it does not meet $T\setminus F$. Hence
$x\notin W(\sigma')$. But $x\in O\subset|\omega|\subset W(\sigma')$, a contradiction.

Now insert $s=1_{\sigma'}$ into \eqref{4.4:eq-interpolation-rule-1}. Each factor
$1_{\sigma'}(\Delta)$ is $1$ or $0$. So the product attached to $\omega$ equals $1$ if every
cube of $\sigma_0$ meeting $|\omega|$ lies in $\sigma'$, and equals $0$ otherwise. When no cube
meets $|\omega|$ the product is empty and equals $1$, in agreement with this description.

Equivalently, the product is $1$ exactly when no cube of $\sigma_0\setminus\sigma'$ meets
$|\omega|$. By the claim this happens exactly when $|\omega|\subset W(\sigma')$. Therefore
\begin{equation*}
K(1_{\sigma'})=\sum_{\omega:\ |\omega|\subset W(\sigma')}K_{\omega}=K^{W(\sigma')},
\end{equation*}
which is \eqref{4.4:eq-interpolation-rule-2}.
\end{proof}

The rule \eqref{4.4:eq-interpolation-rule-1} is the interpolation applied in \cite{B-CMP116}
to the terms of \cite[(1.6)]{B-CMP116}; it reads:
\begin{quote}
\dots for a random walk $\omega$ localized in
$\tilde X_0^5\cup\tilde X_1^5\cup\dots\cup\tilde X_n^5$ we take the
$\{\Delta_1,\dots,\Delta_m\}$ of all cubes from $\sigma_0$ which intersect this localization
domain, and we multiply the term in (1.6) corresponding to $\omega$ by
$s(\Delta_1)\cdots s(\Delta_m)$.
\end{quote}
By the preceding sentence on the same page of \cite{B-CMP116}, each factor in
\cite[(1.6)]{B-CMP116} depends on the external gauge field configuration restricted to the
corresponding domain $\tilde X_j^5$; here the $X_j$ are the localization domains of that
formula and $\tilde X_j^5$ is the domain attached to $X_j$ there. The kernel of each such
factor vanishes in a sufficiently thick neighbourhood of $\partial X_j$, for instance in one
of width $\tfrac13 M_1$. Thus in Ba{\l}aban's construction the domain entering the rule for a
step is the domain $\tilde X_j^5$ on which the step depends, and not the support of its
kernel. For the families of walks of this chapter the domain of a step is the cube listed
with it, and the step reads the configuration at radius $0$ from that cube, so the listed cube
is the domain on which the step depends; for these families too the domain entering
\eqref{4.4:eq-interpolation-rule-1} is the dependence domain, and
\eqref{4.4:eq-interpolation-rule-1} is Ba{\l}aban's rule read for them.

\begin{lemma}[The vertex formula for interpolated functionals]\label{4.4:lem-vertex-formula}
Let $\sigma_0$ be a finite set, and let $f$ be a function on $[0,1]^{\sigma_0}$, with points written
$s=(s(\Delta))_{\Delta\in\sigma_0}$, which is $C^1$ in each variable. For $\sigma\subset\sigma_0$ write
$\sigma^c=\sigma_0\setminus\sigma$, write $(s(\sigma),s(\sigma^c)=0)$ for the point whose coordinates
are $s(\Delta)$ for $\Delta\in\sigma$ and $0$ for $\Delta\in\sigma^c$, and let $1_{\sigma'}$ be the
indicator of $\sigma'\subset\sigma_0$, with $1=1_{\sigma_0}$. Then
\begin{equation}\label{4.4:eq-vertex-formula-1}
f(1)=\sum_{\sigma\subset\sigma_0}f_\sigma ,
\end{equation}
where
\begin{equation}\label{4.4:eq-vertex-formula-2}
f_\sigma:=\prod_{\Delta\in\sigma}\int_0^1 ds(\Delta)\,\frac{\partial}{\partial s(\Delta)}\,
f\bigl(s(\sigma),s(\sigma^c)=0\bigr)
=\sum_{\sigma'\subset\sigma}(-1)^{|\sigma\setminus\sigma'|}f(1_{\sigma'}) .
\end{equation}
Here the product over $\Delta\in\sigma$ is the composition of the operations
$\int_0^1 ds(\Delta)\,\partial/\partial s(\Delta)$, each acting on the variable $s(\Delta)$ with the
others held fixed; for $\sigma=\emptyset$ it is the identity, so that $f_\emptyset=f(0)$.
\end{lemma}

\begin{proof}
For $\Delta\in\sigma_0$ let $E^1_\Delta$ and $E^0_\Delta$ denote evaluation at $s(\Delta)=1$ and at
$s(\Delta)=0$ respectively. Applied to a function that is $C^1$ in each variable, each of them yields
a function of the remaining variables that is again $C^1$ in each of them, so all compositions below
are defined. By the fundamental theorem of calculus, applied in the variable $s(\Delta)$ with the other
variables fixed,
\begin{equation*}
\int_0^1 ds(\Delta)\,\frac{\partial}{\partial s(\Delta)}=E^1_\Delta-E^0_\Delta .
\end{equation*}
Evaluations in distinct variables commute, and hence so do the operations $E^1_\Delta-E^0_\Delta$ for
distinct $\Delta$; in particular the product in \eqref{4.4:eq-vertex-formula-2} does not depend on the
order of its factors. Since $f(s(\sigma),s(\sigma^c)=0)=\prod_{\Delta\notin\sigma}E^0_\Delta f$, we
obtain
\begin{equation*}
f_\sigma=\prod_{\Delta\in\sigma}\bigl(E^1_\Delta-E^0_\Delta\bigr)\prod_{\Delta\notin\sigma}E^0_\Delta f .
\end{equation*}
Expanding the first product by distributivity, a term is determined by the set $\sigma'\subset\sigma$ of
those $\Delta$ at which the factor $E^1_\Delta$ is chosen; it carries the sign
$(-1)^{|\sigma\setminus\sigma'|}$ and equals
\begin{equation*}
\prod_{\Delta\in\sigma'}E^1_\Delta\prod_{\Delta\in\sigma_0\setminus\sigma'}E^0_\Delta f=f(1_{\sigma'}).
\end{equation*}
This is the right member of \eqref{4.4:eq-vertex-formula-2}. Finally, expanding the product over all
$\Delta\in\sigma_0$ in the same way, now with a term for each $\sigma\subset\sigma_0$ (the set at which
$E^1_\Delta-E^0_\Delta$ is chosen),
\begin{equation*}
\sum_{\sigma\subset\sigma_0}f_\sigma
=\prod_{\Delta\in\sigma_0}\Bigl(\bigl(E^1_\Delta-E^0_\Delta\bigr)+E^0_\Delta\Bigr)f
=\prod_{\Delta\in\sigma_0}E^1_\Delta f=f(1),
\end{equation*}
which is \eqref{4.4:eq-vertex-formula-1}.
\end{proof}

The expansions \cite[(1.9)--(1.10)]{B-CMP116} and \cite[(2.8)]{B-CMP116} are instances of
\eqref{4.4:eq-vertex-formula-1}--\eqref{4.4:eq-vertex-formula-2}. The non-linear functionals
denoted $D$, $\mathbf{A}_0$, $\mathbf{H}_k$ in \cite{B-CMP116} are functionals of the operators
interpolated there \cite[(1.5)]{B-CMP116}, which solve separate equations in components
\cite[\S1]{B-CMP116}, and the
function $G(Z_0,X,\cdot,\cdot,\cdot)$ of \cite[(2.6)]{B-CMP116} is an explicit function of its three
operator arguments. Hence, since by the interpolation rule of
Definition~\ref{4.4:def-interpolation-rule} every interpolated operator $K$ evaluated at $1_{\sigma'}$
equals $K^{W(\sigma')}$, the second equality in \eqref{4.4:eq-vertex-formula-2} shows that the term of
\cite[(1.10)]{B-CMP116} indexed by $Y_0$, and the term $H(Z,Z_0)$ of
\cite[(2.8)--(2.9)]{B-CMP116}, are finite alternating sums of one and the same functional evaluated with
every operator $K$ replaced by $K^{W'}$, where $W'$ runs over unions of cubes with $W'\subset Y_0$,
respectively $W'\subset Z$. The Cauchy circles $|s(\Delta)|=e^{\kappa_1}$ enter only the bounds, not
these identities. Consequently, in the first step of the proof in which these terms are estimated, the
ingredient concerning the interpolation parameters on the Cauchy circles is not needed, and the
operator ingredient of that step consists of the operators $K^{W'}$ with $W'\subset W$, where
$W=Y_0$, respectively $W=Z$.

\begin{remark}\label{4.4:lem-walk-expansion-locality}
The statement of this lemma is not written out here; parts of its proof are printed below.
\end{remark}

\begin{proof}[Proof of Lemma~\ref{4.4:lem-walk-expansion-locality}, continued: steps (e)--(g)]
\label{4.4:prf-expansion-words}
Throughout, $\mathsf{K}$ denotes one of the operators of
Lemma~\ref{4.4:lem-walk-expansion-locality}. The settings $\mathfrak{s}$, $\mathfrak{s}'$ and the region
$W$ are those of its statement, and $\mathfrak{s}_\square$ is the setting attached to a cube $\square$
in step (b). The families of cubes $\mathcal{D}$, $\mathcal{D}'(\square)$, $\mathcal{D}''(\square)$ are the
covers of \cite{B-CMP99b}. For a set $X$ and a cube $\square$, $\tilde X^5$ and $\tilde\square^3$ are the
enlargements of \cite{B-CMP99b}.

\medskip\noindent\emph{Step (e): the expansions are sums of words in the letters.}
The operators are taken in the order of the following table. This order has no cycle: each row
uses only operators of earlier rows. Among the operators of the fourth to seventh rows the order is
the one in which the two lemmas and the proposition of this chapter on the inverses
$(QG_0Q^*)^{-1}$ and $(QGQ^*)^{-1}$ write out their expansions: $G$; then
\cite[(3.132)]{B-CMP99b} for $G$; then \cite[(3.133)]{B-CMP99b}; then $H^*\mathbf{J}$; then
\cite[(3.137)]{B-CMP99b} and \cite[(3.138)]{B-CMP99b}; and finally \cite[(3.132)]{B-CMP99b} for
$G_1$, which belongs to the seventh row.

\medskip
\begin{center}
\begin{tabular}{p{0.17\textwidth}p{0.19\textwidth}p{0.24\textwidth}p{0.18\textwidth}}
operator & expansion & letters & earlier operators replaced by their expansions\\
\hline
$G'$ & \cite[(3.87)--(3.90)]{B-CMP99b} & $h_\square$, $K(h_\square)$, $G'_\square$ & none\\
$T$; $P$, $R$ & \cite[(3.95)--(3.98)]{B-CMP99b} & $\chi_\square$, $Q'$, $Q'^*$, $C_\square$,
$G'_\square$, $G'_{\square_0}$, $h'$ & $G'$ (first row), on $\mathcal{D}'(\square)$\\
$G_0$ & \cite[(3.100)--(3.107)]{B-CMP99b} & $D$, $D^*$, $Q$, $Q^*$, $a$, $K(h_\square)$,
$S(\partial h)$, $\zeta_\square$, $P_\square$, $P_{\square,1}(\partial h_\square)$, $G_\square$ &
$P$ (first two rows), on $\mathcal{D}'(\square)$, $\mathcal{D}''(\square)$\\
$G$ & \cite[(3.120), (3.130)]{B-CMP99b} & $\mathbf{J}$-vertices with legs $I$, $G'RD^*$,
$DG'RD^*$ & $G_0$, $G'$, $R$ (first three rows)\\
$(QG_0Q^*)^{-1}$, $(QGQ^*)^{-1}$ & written out in this chapter & $h_\square$, $\chi_\square$, $Q$,
$Q^*$, $C^Q_\square$, $G_\square$ & $G_0$, $\Delta'_\pi$ (third and fourth rows)\\
$H$, $H_0$, $\tilde G$ & \cite[(3.126)]{B-CMP99b}; \cite[(129)--(131)]{B-CMP102b} & $Q$, $Q^*$ &
third to fifth rows\\
$\Delta^{(2)}$, $\Delta^{(2)}_\pi$, $G_1$, $(QG_1Q^*)^{-1}$, $H_1$ &
\cite[(3.134)--(3.138), (3.129)]{B-CMP99b}; tree-like walks & $C^{(2)}$, $\mathbf{J}$ &
$H$, $G'$, $R$, $G_0$ (first six rows)\\
$\Delta^{(k)}$, $\mathsf{A}$ & \cite[(3.156)]{B-CMP99b}; \cite[(2.5), (2.11)]{B-CMP109} &
see below & see below
\end{tabular}
\end{center}
\medskip

The expansions of the first four rows read as follows. For the first row,
\begin{equation*}
G' = \sum_n G'_0 (R')^n,\qquad G'_0 = \sum_\square h_\square G'_\square h_\square,\qquad
R' = \sum_\square K(h_\square) G'_\square h_\square .
\end{equation*}
In the second row $P = G'Q'^*TQ'G'$ and $R = I - P$. The expansion of $T$ is
\begin{equation*}
T = \sum_n C_0\mathcal{R}^n,\qquad C_0 = \sum_\square h_\square C_\square h_\square,
\end{equation*}
where
\begin{equation}\label{4.4:eq-expansion-words-1}
\begin{split}
-\mathcal{R} ={}& \sum_\square (1-\chi_\square)\, Q'G'^2Q'^*\, h_\square C_\square h_\square
+ \sum_\square \chi_\square\, Q'(G'^2 - G'^2_\square)Q'^*\, h_\square C_\square h_\square\\
&+ \sum_\square \bigl[\chi_\square Q'G'^2_\square Q'^*,\, h_\square\bigr] C_\square h_\square .
\end{split}
\end{equation}
For the third row,
\begin{equation*}
G_0 = \sum_n G_{00}\mathcal{R}_0^{\,n},\qquad G_{00} = \sum_\square h_\square G_\square h_\square .
\end{equation*}
Here $\mathcal{R}_0$ denotes the operator written $\mathcal{R}$ in \cite[(3.105)]{B-CMP99b}. It is a
different operator from the one in \eqref{4.4:eq-expansion-words-1}. For the fourth row,
\begin{equation*}
G = \sum_n G_0(\Delta'_\pi G_0)^n .
\end{equation*}

For the fifth row, \cite{B-CMP99b} states only that the two inverses are treated ``in the same way
as the operator $(Q'G'^2Q'^*)^{-1}$. We will not repeat these considerations'', and it defines no
terms. The terms are those written out in the two lemmas on the walk expansions of
$(QG_0Q^*)^{-1}$ and $(QGQ^*)^{-1}$ in this chapter and in the proposition that accompanies them.

The sixth row consists of
\begin{equation*}
H = GQ^*(QGQ^*)^{-1},\qquad H_0 = G_0Q^*(QG_0Q^*)^{-1},\qquad \tilde G = G_0 - H_0QG_0 .
\end{equation*}

In the last row $\mathsf{A} = C^*\Delta^{(k)}C$. The letters of \cite{B-CMP99b}, namely $a$,
$\tilde D^{(2)}$, $\mathbf{J}$, $\mathfrak{g}$ and $C$, together with $(QG_1Q^*)^{-1}$ and $H_1$
replaced by their expansions from the seventh row, describe this operator. They do not fix its
decomposition into terms. The decomposition used for $\Delta^{(k)}$ and $\mathsf{A}$ is the one
constructed in the decomposition theorem for $\mathsf{A}$ in this chapter. Its terms are not walks of
the type of \cite[(3.107)]{B-CMP99b}. Whether the walks of \cite{B-CMP99b} satisfy the walk bound for
$\mathsf{A}$ is neither shown nor used. For this row the rest of step (e) and steps (f), (g) are
replaced by the locality corollary of that theorem (see the end of step (g)).

It remains to check that no letter of a finished expansion is a global object of $\mathfrak{s}$.
Consider the identities \cite[(3.95), (3.105), (3.120), (3.130), (3.134)--(3.138), (3.156)]{B-CMP99b}.
The only global objects in them are operators from earlier rows of the table. Ba{\l}aban replaces each
of these by its own expansion before he speaks of a walk expansion. For $T$ he writes:
\begin{quotation}
This expansion is unsatisfactory yet, because it is expressed in terms of the unlocalized operator
$G'$. To get an expansion having the desired properties we replace the operator $G'$ everywhere in
the series (3.96) by an expansion similar to (3.90)
\end{quotation}
For $G_0$: ``It does not satisfy yet the properties we need, because of the nonlocality of the
operator $P$ \dots\ we replace the operator $P$ by its random walk expansion''. For $G$: ``we replace
each operator in (3.130) by its random walk expansion''. For the seventh row: ``Replacing the
operators in (3.138) by their random walk expansions''. Also: ``$\Delta^{(2)}$ \dots\ involves the
operator $H$. To have a localization in $U$ we have to expand this operator also''. All these
quotations are from \cite{B-CMP99b}.

The expansions of the second and of the fifth row (the latter made in the same way) rest on one
identity. For each cube $\square$ we have
\begin{equation}\label{4.4:eq-expansion-words-2}
\begin{split}
Q'G'^2Q'^*h_\square C_\square h_\square ={}& (1-\chi_\square)\,Q'G'^2Q'^*h_\square C_\square h_\square
+ \chi_\square Q'(G'^2-G'^2_\square)Q'^*h_\square C_\square h_\square\\
&+ \bigl[\chi_\square Q'G'^2_\square Q'^*,\, h_\square\bigr]C_\square h_\square
+ h_\square\bigl(\chi_\square Q'G'^2_\square Q'^*C_\square\bigr)h_\square .
\end{split}
\end{equation}
This is obtained in three moves. First write $1 = (1-\chi_\square)+\chi_\square$ on the left. Then
write $G'^2 = (G'^2-G'^2_\square) + G'^2_\square$ in the $\chi_\square$ part. Finally move
$h_\square$ to the left across $\chi_\square Q'G'^2_\square Q'^*$, which produces the commutator.

The last term of \eqref{4.4:eq-expansion-words-2} equals $h_\square^2$, for three reasons:
\begin{itemize}
\item $\chi_\square h_\square = h_\square$;
\item $C_\square$ inverts $Q'G'^2_\square Q'^*$ on the coarse set of $\mathfrak{s}_\square$;
\item the operator $Q'$ of $\mathfrak{s}$ and that of $\mathfrak{s}_\square$ coincide on
$\tilde\square^3$.
\end{itemize}
Summing \eqref{4.4:eq-expansion-words-2} over $\square$, and using $\sum_\square h_\square^2 = 1$ and
the definition \eqref{4.4:eq-expansion-words-1} of $\mathcal{R}$, gives
$Q'G'^2Q'^*C_0 = I - \mathcal{R}$, which is the identity behind the expansion
$T = \sum_n C_0\mathcal{R}^n$ of the second row.

We now induct along the first seven rows of the table. Each global operator is replaced by its
expansion from an earlier row, and the letters of the resulting words are the letters introduced in
steps (b)--(d). Hence
\begin{equation*}
\mathsf{K} = \sum_{\mathfrak{w}} c_{\mathfrak{w}}\operatorname{val}(\mathfrak{w}).
\end{equation*}
The sum runs over words $\mathfrak{w}$ in the letters of steps (b)--(d), and each coefficient
$c_{\mathfrak{w}}$ is $\pm 1$ or a fixed rational number. Moreover:
\begin{itemize}
\item the set of words is described by incidences among cubes of $\mathcal{D}$,
$\mathcal{D}'(\square)$ and $\mathcal{D}''(\square)$, and by nothing else;
\item $\operatorname{val}(\mathfrak{w})$ is the composition of the values of the letters of
$\mathfrak{w}$.
\end{itemize}

\medskip\noindent\emph{Step (f): from words to walks.}
Ba{\l}aban groups the words into factors $R_\alpha(X)$ by rules that are syntactic: they read the word,
not the setting. The rules are of three kinds. The first is the attachment rule:
\begin{quotation}
We attach each operator $h'_{\square_1}G'_{\square_1}h'_{\square_1}$ to a closest operator
$K(h'_{\square_2})G'_{\square_2}h'_{\square_2}$ on the right, and the operator
$h_\square C_\square h_\square$ to a closest operator of that type on the left
\end{quotation}
The second is the rewriting by commutators together with the pairing \cite[(3.97)]{B-CMP99b} of two
words with the same cubes. The third is the cancellation used in the expansion of $G_0$ in
\cite{B-CMP99b}. Of the result, \cite{B-CMP99b} says: ``Each factor is a product of several
operators localized correspondingly in cubes $\square_1, \square_2, \dots$ A sum of these cubes
determines a localization domain $X$''. It adds that such a factor ``is localized in $X$, i.e.\ its
kernel has a support in $X\times X$, it depends on $U$ restricted to $\tilde X^5$''.

Ba{\l}aban's list of the types of factors ends with ``etc.'' Nothing depends on completing it. Whatever
the grouping, a factor is a finite sum of words. The letters of these words are attached to cubes of
the family that makes up $X$, and by step (d) all their dependence sets lie in $\tilde X^5$.

\medskip\noindent\emph{Step (g): conclusion.}
Let $\omega$ be a term of the expansion, with localisation domains $X_i$, put
$|\omega| = \bigcup_i \tilde X_i^5$, and suppose that $|\omega|\subset W$. The three requirements of
locality in Lemma~\ref{4.4:lem-walk-expansion-locality} are checked as follows.
\begin{itemize}
\item By steps (a) and (b), the cubes and partitions that can occur in $\omega$ are the same for
$\mathfrak{s}$ and for $\mathfrak{s}'$. This is the requirement on the index set.
\item By steps (b)--(d), every letter of $\omega$ has the same value for $\mathfrak{s}$ and for
$\mathfrak{s}'$. This is the requirement on the values.
\item The kernels are supported in $X_i\times X_i$, at distance at least $5M_1$ from
$\partial\tilde X_i^5$. Let $c(d)$ be a constant depending only on the dimension; the letters of
finite range have range $c(d)L$, and this range is absorbed because $M_1\ge c(d)L$. This is the
requirement on the support.
\end{itemize}

For $\Delta^{(k)}$ and $\mathsf{A}$ the three requirements are the content of the locality corollary
of the decomposition theorem for $\mathsf{A}$. Write $\mathsf{A}_\omega$ for the term of
$\mathsf{A}$ indexed by $\omega$ and $[\omega]$ for its footprint. That theorem gives three facts:
\begin{itemize}
\item $\mathsf{A}_\omega$ reads the data of the setting and the geometry on $[\omega]$ itself, that
is, with reading radius zero;
\item $\mathsf{A}_\omega(b,b') = 0$ unless the bonds $b$, $b'$ lie in $[\omega]$ at distance at least
one from $\partial[\omega]$;
\item a comparison clause of that theorem applies to two settings that agree on $[\omega]$.
\end{itemize}
\end{proof}

\medskip\noindent\emph{Relation to \cite{B-CMP99b}.} Steps (a)--(g) transcribe the step on which the
proof of \cite[Theorem~3.14]{B-CMP99b} rests. There Ba{\l}aban writes: ``We take random walk
expansions for both operators, and in the difference all terms for walks with localizations contained
in $\Omega$ are cancelled''.

\medskip\noindent\emph{No margin is spent.} The footprint of a term is defined through $\tilde X_i^5$,
respectively through $[\omega]$. Hence the inclusion $|\omega|\subset W$ already contains the
neighbourhoods that enter the definitions.

\medskip\noindent\emph{Levels $j\le k$.} The proof holds without change when $k$ is replaced by $j$,
for the operators of the truncated sequence $(\Omega_i)_{i\le j}$. Indeed, $W\subset\Omega_{k+1}$ and
$\Omega_{k+1}\subset\Omega_j$, so $W$ lies in the top layer of the truncated sequence. The cubes of
the cover that lie in $W$ are the standard cubes of side $2M_1L^{j-k}$, and the local sequences are
concentric cubes.

\medskip\noindent\emph{Bounds at free current.} Convergence and bounds of these expansions at free
current $\mathbf{J}$ are not theorems of \cite{B-CMP99b}. The statement \cite[Theorem~3.12]{B-CMP99b}
is made under \cite[(3.35)--(3.36)]{B-CMP99b}, that is with $J = J(U)$, and it concerns the series,
not the individual terms. The bounds \cite[(1.6)--(1.7)]{B-CMP116} and the proposition accompanying
the expansions of $(QG_0Q^*)^{-1}$ and $(QGQ^*)^{-1}$ stand under the same hypotheses.

For $\mathsf{A}$ the bound is the walk bound for $\mathsf{A}$. It is the restriction of the bound of
the decomposition theorem, which supplies it, to the set $\mathcal{S}_0$ at which the hypothesis of
the lemma of this chapter that uses these bounds is stated (see the end of this paragraph). The
remaining operators at free $\mathbf{J}$ are $G$, $G_1$, $H$, $H_0$, $H_1$, $\tilde G$ and the
inverses they contain; the statements of this chapter that use them read them as infinite series in
the norms of \cite{B-CMP116}. For these operators the bound is supplied by the one-family bound,
restricted to the operators other than $\mathsf{A}$. That bound enters as a hypothesis, and it
concerns its own family of expansions, not the words of the table above. The hypothesis of the lemma
of this chapter that uses these bounds is therefore the conjunction of the walk bound for
$\mathsf{A}$ and the one-family bound, both at $\mathcal{S}_0$. The walk bound for $\mathsf{A}$ is
used in the one form that follows.

\begin{definition}[Hypothesis on the quadratic form's walk expansion]\label{4.4:def-form-walk-hypothesis}
Let $\rho_0$, $B_{\mathsf{Q}}$ and $\delta_{\mathsf{Q}} = 1/\rho_0$ be the numbers of the walk-expansion
theorem for the quadratic form. These numbers depend only on $(d,L)$, on the structure group $G$ and
on the function $M(\cdot)$ of that theorem, which is an object distinct from the cube side $M$. They are
fixed before $M$, and they involve none of $k$, $K$, $M$, the nest, the volume or $g$. Put
\begin{equation}\label{4.4:eq-form-walk-hypothesis-1}
\rho_w := \sqrt{d}\,\rho_0, \qquad
\delta_{\mathsf{A}} := \frac{1}{2\sqrt{d}\,\rho_0}, \qquad
\mathsf{B} := d\,(\rho_0+2)^d\,B_{\mathsf{Q}} .
\end{equation}

We say that the operator $\mathsf{A}$ of the quadratic form satisfies the \emph{hypothesis on the
quadratic form's walk expansion} if the following hold with the constants
\eqref{4.4:eq-form-walk-hypothesis-1}.
\begin{enumerate}
\item $\mathsf{A} = \sum_\omega \mathsf{A}_\omega$ is a walk expansion in the sense of
Definition~\ref{4.4:def-walk-systems}, and its weighted norm satisfies
$\lvert\!\lvert\!\lvert\mathsf{A}\rvert\!\rvert\!\rvert_{\delta_{\mathsf{A}}} \le \mathsf{B}$.
\item Every domain of every walk $\omega$ of $\mathsf{A}$ has $\ell^1$-diameter at most $\rho_w$.
\item Every $\mathsf{A}_\omega$ is analytic on the domain $\mathbf{P}'$. It is also analytic on the
analyticity domain $\mathcal{U}(1)$ of the walk-expansion theorem, a domain which, under the closeness
condition on the backgrounds in Definition~\ref{4.4:def-regular-backgrounds}, contains
$\mathbf{P}' \cup \mathbf{P}_{2\rho}$; here $\mathbf{P}_{2\rho}$ is the polydisc $\mathbf{P}_\sigma$ of
Definition~\ref{4.4:def-regular-backgrounds} at $\sigma = 2\rho$. The bound of the first clause holds
at every point of $\mathcal{U}(1)$, in particular at points of $\mathbf{P}_{2\rho}$ outside $\mathbf{P}'$.
\item Every $\mathsf{A}_\omega$ depends on $(\mathbf{U},\mathbf{J})$ only through their restriction to the
defining neighbourhood $[\omega]$ of the walk $\omega$. It depends on the geometry in the same way,
through $[\omega]$ alone.
\item Every $\mathsf{A}_\omega$ is defined and analytic, with its bound, on local data given on
$[\omega]$ only. That is, $\mathsf{A}_\omega$ is a function on the local class
$\mathbf{P}'_{[\omega]}$ of data on $[\omega]$. Moreover, the majorants $a_\omega$ in the bound of the
walk-expansion theorem are numbers attached to $\omega$ and do not depend on the datum. Consequently
every row sum and every column sum entering the bound
$\lvert\!\lvert\!\lvert\mathsf{A}\rvert\!\rvert\!\rvert_{\delta_{\mathsf{A}}} \le \mathsf{B}$ is likewise a
number free of the datum.
\item The rule $\mathfrak{g} \mapsto (\mathsf{A}_\omega[\mathfrak{g}])_\omega$ is local in the sense of
Definition~\ref{4.4:def-local-rules}. It has, in addition, the four further properties of such rules
listed in the proposition on walk-expansion rules.
\item The clauses above hold uniformly in $k$, $K$, $M$, the volume and $g$, among admissible nests,
and uniformly on $\mathbf{P}'$. Two facts hold at the trivial nest. First, the admissibility condition
is void there. Second, every representative of $\mathbf{P}'$ lies in $\mathcal{U}(1)$ under the three
ceilings imposed in the corollary for the trivial nest.
\item The clauses above hold at every level $k \ge 0$, with the one set of numbers
\eqref{4.4:eq-form-walk-hypothesis-1} for all $k \ge 0$.
\end{enumerate}
\end{definition}

The clauses of Definition~\ref{4.4:def-form-walk-hypothesis} are the conclusions of the walk-expansion
theorem for the quadratic form and of its two corollaries, specialised to the trivial nest. The
walk-expansion theorem covers every level $k \ge 1$.

At $k = 0$, which is Ba{\l}aban's first step, the same clauses hold by part (iii) of the level-zero lemma
for the quadratic form. This is the step used, with $j = 1$, in the proof of independence of the
localised terms from the small-field region.

The single set of numbers in the last clause of Definition~\ref{4.4:def-form-walk-hypothesis} is
obtained from the level-zero values and the values for $k \ge 1$ as follows. $B_{\mathsf{Q}}$ is the
larger of the two values. The constant $a_1^{\flat}$ appearing in the walk-expansion theorem and in the
level-zero lemma is the smaller of its two values. Since $\rho_0$ belongs to this single set, the
constants \eqref{4.4:eq-form-walk-hypothesis-1} are the same at every $k \ge 0$.

\begin{definition}[One walk family for the other operators]\label{4.4:def-one-family-hypothesis}
Let $k\ge 1$. Let $\mathbb{L}_1$ be the list of operators of the step treated in this chapter, and let
$\mathsf{A}$ be the operator of the quadratic form of Definition~\ref{4.4:def-form-walk-hypothesis}.
The remaining members of the list are
\begin{gather*}
G',\quad T',\quad P,\quad R;\qquad G_0,\quad T^{-1}:=(QG_0Q^*)^{-1},\quad H_0,\quad \tilde G;\\
G_\pi\ (=G),\quad T_\pi^{-1},\quad H;\qquad \Delta^{(2)}[\mathfrak{j}_0];\qquad
G_1,\quad T_1^{-1},\quad H_1;\qquad \Delta^{(k)}.
\end{gather*}
These operators fill the first eleven rows of the table of the walk-expansion theorem for the
operators of the list, $\Delta^{(k)}$ occupying the eleventh row; the operator of its twelfth row,
which is not among the operators displayed, is called here the \emph{twelfth-row operator} and
enters below only through the definition of trivial traces. Here $\Delta^{(2)}[\mathfrak{j}_0]$ is the
member $\Delta^{(2)}$ with the argument $\mathfrak{j}_0$ of that theorem.
For an operator $K$ of $\mathbb{L}_1\setminus\{\mathsf{A}\}$, a \emph{walk-expansion rule} assigns to
every setting $\mathfrak{g}$ of the geometry an index set $\mathcal{W}_K(\mathfrak{g})$ of walks and, to
each walk $\omega$, a set $\operatorname{dom}\omega$ of domains, two end sets $X_-(\omega)$,
$X_+(\omega)$, a weight $\kappa_\omega$ and an operator $K_\omega$, with $K=\sum_\omega K_\omega$.
The \emph{footprint} $[\omega]$ of $\omega$ is the closure of the union of the open cubes of
$\operatorname{dom}\omega$; it is the set written $|\omega|$ in
Definitions~\ref{4.4:def-local-rules} and~\ref{4.4:def-interpolation-rule}. We write
$\mathcal{S}_0$ for the trivial setting and $\mathcal{S}$ for a general admissible one.

A walk $\omega$ has \emph{trivial traces} if either the setting is $\mathcal{S}_0$, or, in a general
admissible geometry, $X_n\cap[\omega]=\emptyset$ for every $n<k_\ast$ and,
when $K$ is $\Delta^{(k)}$ or the twelfth-row operator, moreover
$[\omega]\subset\Omega_t$. Here $k_\ast$ is the step index of the walk-expansion theorem for the
operators of the list, and the members $X_n$, $\Omega_n$ of the nest are indexed as in that theorem;
this is the indexing in which the admissibility-margin statement used below is written. Further,
$t$ is the index of the member $\Omega_t$ of the nest that the
walk-expansion theorem for the operators of the list uses for these two operators, in its indexing;
the admissibility-margin statement, which places the range of this chapter in that indexing, also
carries a second member of the same index, written $\Omega'_t$.

The \emph{one-family hypothesis} is the following statement. For each
$K\in\mathbb{L}_1\setminus\{\mathsf{A}\}$ there is a walk-expansion rule
\begin{equation*}
\mathfrak{g}\longmapsto\bigl(\mathcal{W}_K,\ \operatorname{dom}\omega,\ X_\pm(\omega),\
\kappa_\omega,\ K_\omega\bigr),
\end{equation*}
belonging to the class of walk-expansion rules at every admissible geometry, with
$K=\sum_\omega K_\omega$. The following four properties hold together, for one and the same terms
$K_\omega$ and one and the same set of weights $\kappa_\omega$. The class conditions hold at every
setting. The bounds (f1), (f3) and (f4) are read at $\mathcal{S}_0$ on the set $\mathbf{P}'$ of
Definition~\ref{4.4:def-form-walk-hypothesis}, on which the terms of $\mathsf{A}$ are analytic, and
(f4) is moreover read at trivial traces as stated there.
\begin{itemize}
\item[(f1)] \emph{Bound, with its local half.} By the zeroth and third conditions of the class, the
walk sum exists and converges absolutely at free
$\mathbf{J}$, in its two-ended form. Its constants depend only on $d$, $L$ and the gauge group, and
they are fixed before $M$. Hence, for the members of the list that act on the unit lattice,
$\lvert\lvert\lvert K\rvert\rvert\rvert_{\theta}\le B_K$ at a rate $\delta_K$, where
$\lvert\lvert\lvert\cdot\rvert\rvert\rvert_{\theta}$ is the triple norm of
Definition~\ref{4.4:def-form-walk-hypothesis}, here carrying the index $\theta$ (a local letter, not
the schedule fraction), and $\delta_K$ is the rate of the bound. Moreover, by the second and third
conditions of the class, every
$K_\omega$ is defined and analytic, together with its bound $\kappa_\omega$, on local data given on
$[\omega]$ only.
\item[(f2)] \emph{Locality}, for walks with trivial traces and for each operator of the list:
clauses (b1), (b2) and (b4) of the walk-expansion theorem for the quadratic form, together with the
fourth and fifth conditions of the class of walk-expansion rules. Explicitly:
\begin{itemize}
\item the domains are closed cubes, and the union of the open cubes of $\operatorname{dom}\omega$ is
connected, so that $[\omega]$ is the closure of its connected interior (clause (b1));
\item $K_\omega=1_{X_-(\omega)^\circ}\,K_\omega\,1_{X_+(\omega)^\circ}$, where $X^\circ$ denotes the
core of a set $X$, and the cores of the end
sets lie at distance at least $4L$ inside $X_\pm(\omega)$ (clause (b4));
\item $K_\omega$ is a function of the local datum on $[\omega]$, read at radius $0$ (clause (b2));
\item the fourth condition of the class of walk-expansion rules holds, the class being that of the
definition which accompanies the walk-expansion theorem for the quadratic form;
\item whether $\omega$ is an index, its domains, its weight and the function $K_\omega$ are determined
by the traces of the nest on $[\omega]$ (the fifth condition, at radius $0$).
\end{itemize}
These items yield the rules (L-ind) and (L-val) of Definition~\ref{4.4:def-local-rules} and the
support property of $K_\omega$ stated there,
the vertex rule $K(1_{\sigma'})=K^{W(\sigma')}$ of Definition~\ref{4.4:def-interpolation-rule}, and
clause (c4) of the proposition of this chapter on the walk-expansion family of the operators of
$\mathbb{L}_1$. The margin of (b4) (the
end-cube margin)
and the norms of (f4) are stated for each operator of the list at trivial traces only, and (f2) is
assumed in that scope. In the admissible geometry this scope is the hypothesis of the first
corollary (C1) below, for a closed set $W$: $W\cap X_n=\emptyset$ for every $n<k_\ast$, so that
$W\subset\Omega_{k_\ast}\subset\Omega_0$, and, for $\Delta^{(k)}$ and the twelfth-row operator,
moreover $W\subset\Omega_t$.
\item[(f3)] \emph{Covariance.} By the sixth condition of the class, lattice covariance holds
termwise, with
$\kappa_{r\omega}=\kappa_\omega$, together with the domain clause of clause (c) of the walk-expansion
theorem for the quadratic form. Here $r$ ranges over the translations by $M\mathbb{Z}^4$ and the
reflections and elements of $W_4$ that are used in the third corollary of the theorem on a uniform
region $W$ and in the mean-value lemma.
\item[(f4)] \emph{Norms of \cite{B-CMP116}, on the same terms with the same $\kappa_\omega$.}
Let $\beta^\sharp\in(0,1)$ be the H\"older exponent of the walk-expansion theorem for the operators
of the list. For a function $u$ set
\begin{equation*}
N_{\beta^\sharp}(u):=\sup|u|+\sup|\nabla_{\mathbf{U}}u|
+\bigl[\nabla_{\mathbf{U}}u\bigr]_{\beta^\sharp},
\end{equation*}
where $[\,\cdot\,]_{\beta^\sharp}$ is the H\"older norm of exponent $\beta^\sharp$, as in
\cite[(3.39)--(3.40)]{B-CMP99b}. Below, $c_{\mathrm{vol}}$ is the volume constant of the
walk-expansion theorem for the operators of the list, $|f|_Y$ denotes the supremum norm of $f$ over
a set $Y$, in the notation of \cite[(3.39)--(3.40)]{B-CMP99b}, and the coarse operators are those
members of the list whose kernels $K_\omega(y,y')$ are functions of two points $y,y'$ of the coarse
lattice of the step. The following bounds hold, for functions $f$ on which the operators of the
first line act and fields $\eta$ on which those of the second line act, for every walk $\omega$ with
trivial traces and every datum of the local analyticity space of data on $[\omega]$ with radii
$a_1^\sharp,a_0^\sharp$, written $\mathcal{U}^{\flat\nabla}_{[\omega]}(a_1^\sharp,a_0^\sharp)$:
\begin{align*}
N_{\beta^\sharp}(K_\omega f)&\le c_{\mathrm{vol}}\,\kappa_\omega\,|f|_{X_+(\omega)}
&&\text{for }K\in\{G_0,G_\pi,G_1,\tilde G,G',P\},\\
N_{\beta^\sharp}(K_\omega \eta)&\le c_{\mathrm{vol}}\,\kappa_\omega\,|\eta|_{X_+(\omega)}
&&\text{for }K\in\{H_0,H,H_1\},\\
|K_\omega(y,y')|&\le\kappa_\omega&&\text{for the coarse operators }K.
\end{align*}
In \cite{B-CMP116} Ba{\l}aban writes that ``we can replace the propagators by arbitrary operators
having the same regularity properties and satisfying the same bounds. Only these bounds were
important in the analysis of Sect. C, E [15]''. He also writes that ``This bound holds for all
operator norms in formulations of theorems in [13], e.g. in Theorem 3.1''. (The bracketed numbers
refer to the reference list of \cite{B-CMP116}; there [15] is \cite{B-CMP102b}.) The bounds meant in
(f4) are therefore those read in \cite[Sects.~C,~E]{B-CMP102b}:
\begin{itemize}
\item the first two bounds of \cite[(3.42)]{B-CMP99b};
\item the first bound of \cite[(3.43)]{B-CMP99b};
\item the bounds \cite[(3.132), (3.133)]{B-CMP99b};
\item the bound for the site operator on H\"older data.
\end{itemize}
They are not the bound for $G\nabla^*$ on sup data, nor \cite[(3.44)--(3.45)]{B-CMP99b} at the
full $\nabla^*$, nor \cite[(3.49)]{B-CMP99b}, nor any bound for a Laplacian.
\end{itemize}

The hypothesis is read together with the following three corollaries of (f1)--(f4), which are quoted
here by statement from the walk-expansion theorem for the operators of the list. In
these corollaries $K^W$ is the restricted operator of Definition~\ref{4.4:def-interpolation-rule},
and $B^{\natural}$ is the constant of that theorem which enters the bounds of (C1) and (C2).
\begin{itemize}
\item[(C1)] Let $W$ be closed and such that all its walks have trivial traces. In an admissible
geometry this hypothesis is taken in the form: $W\cap X_n=\emptyset$ for every $n<k_\ast$, so that
$W\subset\Omega_{k_\ast}\subset\Omega_0$, and, when $K$ is $\Delta^{(k)}$ or the twelfth-row
operator, moreover
$W\subset\Omega_t$. It holds for any $W$
at $\mathcal{S}_0$, and for $W\subset\Lambda_{k+1}$ in $\mathcal{S}$. Then $K^W$ converges absolutely
in the norms of (f4), and
\begin{equation*}
N_{\beta^\sharp}(K^W f)\le 7\,c_{\mathrm{vol}}\,B^{\natural}\,|f|_W .
\end{equation*}
Moreover $K^W$ is analytic and reads the datum on $W$ only. It satisfies
$K^W[\mathcal{S}]=K^W[\mathcal{S}_0]$ term by term, and it is $r$-covariant on the common domain.
\item[(C2)] This is the form of \cite[(1.7), (1.11)]{B-CMP116}, in which the printed factor
$e^{16\kappa_1}$ is replaced by $e^{5^d\kappa_1}$. Suppose that $|s(\Delta)|\le e^{\kappa_1}$, and
that the cubes of $\sigma_0$ have side $M'$ with $M'\ge\max\{\rho_0,\,4\cdot 5^d\kappa_1\rho_0\}$,
where $\rho_0$ is the constant of the walk-expansion theorem for the operators of the list;
this is a lower bound only. Then
\begin{equation*}
N_{\beta^\sharp}(K(s)f)\le 7\,c_{\mathrm{vol}}\,B^{\natural}\,e^{5^d\kappa_1}\,|f|.
\end{equation*}
\item[(C3)] The same holds at every level $1\le j\le k$ on the truncated trivial nest. At $j=0$ one
has $H=H_0=H_1=1$ and $\tilde G=0$.
\end{itemize}

\emph{Range of $k$.} The walk-expansion theorem for the operators of the list speaks at $k\ge 1$.
At $k=0$ the operators read in \cite[(1.2)--(1.4)]{B-CMP116} are $1$, $1$ and $0$. The others are
$0$, $1$, finite words in local operators, or $G'$ and $G_0$ obtained by inversion (clauses (i) and
(iv) of the level-zero lemma). Nothing is required there.

\emph{Use in this chapter.} Locality (f2) is used in clause (a) of the theorem on a uniform region
$W$, for the operators of $\mathbb{L}_1\setminus\{\mathsf{A}\}$. The bound (f1) is used in clause
(d) of that theorem. Both (f1) and (f2), with the local half of (f1), are used in the lemma
accompanying that theorem. The pair (f1), (f4) is used in the three corollaries of the theorem on a
uniform region $W$. It is also
used for the unique small solutions, written $D$, $\mathbf{A}_0$, $\mathbf{H}_k$ where they are
introduced, with the
operators $H^{W'}$, $\tilde G^{W'}$, $H_0^{W'}$. Covariance (f3) is used in item (h10) of the
hypothesis list of this chapter, in the third of those corollaries
and in the mean-value lemma.

\emph{Scope of use.} In this chapter (f2) is read only at $\mathcal{S}_0$, or for walks $\omega$ with
$[\omega]\subset W$, where $W$ is a closed union of $M_1$-cubes with $W\subset\Omega_{k+1}$. This is
the range of the theorem on a uniform region $W$ and of its lemma and corollary. It is also the range
of the vertex rule and of clause (c4) of the proposition on the walk-expansion family of the
operators of $\mathbb{L}_1$, inside a region contained in $\Omega_{k+1}$, and of the local
symmetry statement, whose cubes lie in a region $W'\subset\Omega_{k+1}$. Every such walk has trivial
traces.
\end{definition}

\begin{proof}
We prove the last assertion of the definition, that every walk in the scope of use has trivial
traces.

Suppose first that the setting is $\mathcal{S}_0$. By the definition of trivial traces, every walk at
$\mathcal{S}_0$ has trivial traces.

Suppose next that the setting is a general admissible one, $\mathcal{S}$. Let $W$ be a closed union
of $M_1$-cubes with $W\subset\Omega_{k+1}$. In the indexing of the same nest used by the
walk-expansion theorem for the operators of the list, the admissibility-margin statement gives, from
$W\subset\Omega_{k+1}$, that $W\subset\Omega_t\cap\Omega'_t$, and the same statement supplies the
margin
\begin{equation*}
\operatorname{dist}_\infty\bigl(\Omega_t,\Omega_{k_\ast}^{c}\bigr)\ge 6M
\end{equation*}
in that indexing. By this margin, $W$ satisfies both parts of the hypothesis of (C1) in the form
taken in an admissible geometry: first,
$W\cap X_n=\emptyset$ for every $n<k_\ast$, so that $W\subset\Omega_{k_\ast}\subset\Omega_0$;
second, $W\subset\Omega_t$, which is the additional condition for $\Delta^{(k)}$ and the twelfth-row
operator.

Now let $\omega$ be a walk with $[\omega]\subset W$. For every $n<k_\ast$,
\begin{equation*}
X_n\cap[\omega]\subset X_n\cap W=\emptyset .
\end{equation*}
When $K$ is $\Delta^{(k)}$ or the twelfth-row operator, we also have
$[\omega]\subset W\subset\Omega_t$. These are exactly the two conditions in the definition of trivial
traces in an admissible geometry. Hence $\omega$ has trivial traces.
\end{proof}

\begin{remark}
Nothing of (f2) or (f4) is read in this chapter at walks whose traces are not trivial. The statement
on the independence of the localised terms from the small-field region reads the operators of this
family only at the value $\nu=0$ of its index $\nu$. Two things lie
outside the one-family hypothesis.

The first concerns seam walks, by which we mean the walks of this family whose traces are not
trivial. The bounds of (f4) at non-trivial traces, that is, the norms of
\cite{B-CMP116} on the seam walks of this same family, are not
proved. The walk-expansion theorem for the operators of the list makes no claim for general
nests beyond the conditions of its class. Now $K(s)=\sum_\omega K_\omega\prod s(\Delta)$ is one
decomposition of $K[\mathcal{S}]$. A straddling large-field region $X$, that is, one met by walks of
the family whose footprints meet some $X_n$, therefore collects
seam walks of this family. There the printed argument no longer applies verbatim to two results:
\begin{itemize}
\item the straddling half of the transfer of the analysis of \cite[\S3]{B-CMP109} to the cover,
which is printed without proof in \cite[p.~276]{B-CMP119} and cited as printed;
\item the convergence in the representation and bound of the large-field terms.
\end{itemize}
This missing bound falls within what that transfer leaves unproved.

The second concerns the covariant-Laplacian bound. For $G_0$, $H_0$ and $\tilde G$ this bound is
read in \cite[Sect.~E]{B-CMP102b}, and it is supplied by (f4) for no operator. For these three
operators it is the content of the covariant-Laplacian corollary of the walk-expansion theorem for
the operators of the list. For $G_\pi$ and $G_1$
Ba{\l}aban excludes it: ``It does not hold for $G$, $G_1$'' \cite{B-CMP99b}. For $H$ and $H_1$ it is
not proved at free $\mathbf{J}$; it belongs to the operator bounds at free current.
\end{remark}

\begin{definition}[Ceilings on the primed radii]\label{4.4:def-radius-ceilings}
Let $B^{109}$ and $\mathsf{o}$ be as in Definition~\ref{4.4:def-local-data-classes}; thus
$\mathsf{o}LM$ is the cube side of \cite[(1.12)]{B-CMP109}, and the constant $B$ of
\cite[(1.12)]{B-CMP109} is written $B^{109}$. Let $c_{\mathfrak{g}}$ denote
the norm of the bracket $[\cdot,\cdot]$ on $\mathfrak{g}^c$, and let $\mathsf{a}_U$ be the constant of
the space $\mathcal{U}_{\mathbb{R}}$ of real regular backgrounds
(Definition~\ref{4.4:def-regular-backgrounds}).

The corollary on representatives in the analyticity domain $\mathcal{U}(1)$ asserts the
following. Assume the floor $M \ge M_{\flat}$. Then
$\mathcal{U}(1)$ contains every representative
$(\mathbf{U},J)=(U'U,J)$ that obeys conditions I.(i)--(iii) of \cite{B-CMP109}, that is
\cite[(1.11)--(1.14)]{B-CMP109}, at $j=k+1$, with constants $(\alpha_0',\alpha_1',\gamma_0')$ such that
\begin{equation*}
\mathsf{o}LMB^{109}\alpha_0' \le a_{\mathrm{reg}}^{\oplus},
\qquad \alpha_1' \le a_1^{\oplus},
\qquad \gamma_0' \le a_0^{\oplus}.
\end{equation*}
These are three ceilings on the primed radii of \cite{B-CMP116}. They are imposed after $M$ is
fixed, and they make precise the requirement of \cite{B-CMP109} that these radii be sufficiently
small.

We say that the primed radii and the constant $\mathsf{a}_U$ satisfy the \emph{halved ceilings} if
\begin{equation}\label{4.4:eq-radius-ceilings-1}
\begin{aligned}
&2\mathsf{o}LMB^{109}\alpha_0' \le a_{\mathrm{reg}}^{\oplus},
\qquad 2\alpha_1' \le a_1^{\oplus},\\
&2\gamma_0' \le a_0^{\oplus},
\qquad 2c_{\mathfrak{g}}\mathsf{a}_U \le 1,
\end{aligned}
\end{equation}
and that they satisfy the \emph{regularity ceiling} if
\begin{equation}\label{4.4:eq-radius-ceilings-2}
36\,\gamma_+^{\mathsf{Q}}\,\mathsf{a}_U \le a_1^{\oplus}\gamma_{\oplus},
\qquad \text{where } \gamma_+^{\mathsf{Q}} := B_{\mathsf{Q}}\,S(\delta_{\mathsf{Q}}).
\end{equation}
Here $\gamma_+^{\mathsf{Q}}$ is the constant of the operator bounds for $\mathsf{Q}$, formed there
from the bound $B_{\mathsf{Q}}$, the rate $\delta_{\mathsf{Q}}$ and the function $S(\cdot)$ as
displayed, and $\gamma_{\oplus}$ is the constant that accompanies $a_1^{\oplus}$.
In \eqref{4.4:eq-radius-ceilings-1} and \eqref{4.4:eq-radius-ceilings-2} the constants
$a_{\mathrm{reg}}^{\oplus}$, $a_1^{\oplus}$, $a_0^{\oplus}$ are formed by the recipe of the
corollary on representatives, but from radii $a_1^{\sharp}$ and $a_0^{\sharp}$ satisfying
\begin{equation*}
a_1^{\sharp} \le \tfrac12 a_1^{\flat}, \qquad a_0^{\sharp} \le a_0^{\flat},
\end{equation*}
in place of the pair $(a_1^{\flat},a_0^{\flat})$ of radii from which the corollary forms them.
With this choice the localised current increment lies in its target class
$\mathcal{U}^{\flat\nabla}$.
\end{definition}

Both \eqref{4.4:eq-radius-ceilings-1} and \eqref{4.4:eq-radius-ceilings-2} are displayed
because the hypothesis on the quadratic
form's walk expansion (Definition~\ref{4.4:def-form-walk-hypothesis}) and the hypothesis of one
walk family for the other operators (Definition~\ref{4.4:def-one-family-hypothesis}) are read on
the polydisc $\mathbf{P}'$ in the setting of the corollary on representatives, and locally on the
polydisc $\mathbf{P}'_D$, and because the former hypothesis is also read on
\begin{equation*}
\mathcal{U}(1) \supseteq \mathbf{P}' \cup \mathbf{P}_{2\rho}
\end{equation*}
at the extended points $\hat{\mathsf{u}}$; here $\mathbf{P}_{2\rho}$ is the polydisc
$\mathbf{P}_{\sigma}$ of Definition~\ref{4.4:def-regular-backgrounds} at $\sigma = 2\rho$.

The proof of the cell bound uses only the first and the last of the four inequalities in
\eqref{4.4:eq-radius-ceilings-1}. The second and the third are not used there; the factor $2$ in
them is spare margin.

The regularity ceiling \eqref{4.4:eq-radius-ceilings-2} is the ceiling required by the positivity
statement, read at the constant $\mathsf{a}_U$ of $\mathcal{U}_{\mathbb{R}}$.

All of these conditions are one-sided: each bounds a primed radius, or $\mathsf{a}_U$, from
above. Their relation to the lower bounds on the primed radii is as follows. A separate
condition, the ratio condition on the polydiscs, involves two ratios. The first is the ratio of
the radii in the current slot. In the notation of \cite{B-CMP109} it is $\gamma_0'/\gamma_0$, with
$\gamma_0$ the radius of \cite[(1.14)]{B-CMP109}; in the notation of \cite{B-CMP116} it is
$\alpha_0'/\alpha_0$. The second ratio is $\alpha_1'/\alpha_1$. The ratio condition bounds these
two ratios from below. Together with that floor, the pair $(\gamma_0',\alpha_1')$, where
$\gamma_0'$ is the radius written $\alpha_0'$ in \cite{B-CMP116}, ranges over a window.

This window is opened by ceilings, imposed after $M$, on the unprimed radii. These are
restrictions of the same kind as those of \cite{B-CMP109}, where one reads: ``The constants
$\varepsilon_0, \varepsilon_1, \alpha_0, \alpha_1$ depend on $M$ and satisfy numerous
restrictions''.

The first inequality of \eqref{4.4:eq-radius-ceilings-1},
$2\mathsf{o}LMB^{109}\alpha_0' \le a_{\mathrm{reg}}^{\oplus}$, is a plain ceiling on the primed
radius of the bond slot. It is paired with no floor other than the trivial one,
$\alpha_0' \ge \alpha_0$, and no upper bound is imposed on the ratio $\alpha_0'/\alpha_0$.

\begin{lemma}[Coercivity of the form on the polydisc]\label{4.4:lem-form-coercivity}
Let $\mathsf{A}=\sum_{\omega}\mathsf{A}_{\omega}$ be a family of forms on
$\ell^2(\mathbb{B}(\Omega_{k+1});\mathfrak{g}^c)$, indexed by $(\mathbf{U},\mathbf{J})$ and
given with its walk expansion, and assume the following.
\begin{enumerate}
\item[(a)] The walk hypothesis of Definition~\ref{4.4:def-form-walk-hypothesis} holds on the
domain $\mathcal{U}(1)$. In particular every $\mathsf{A}_{\omega}$ is analytic on $\mathcal{U}(1)$, and
$|||\mathsf{A}|||_{\delta_{\mathsf{A}}}\le\mathsf{B}$ at every point of $\mathcal{U}(1)$.
\item[(b)] The positivity lemma at real regular backgrounds holds, for the form on any region
$\Lambda$ of the class treated in that lemma, in the following form. Let $\mathcal{U}$ be a space of
$G$-valued configurations with four properties:
\begin{itemize}
\item it is invariant under $G$-valued gauge transformations;
\item it contains $1$;
\item $(U,0)\in\mathcal{U}(1)$ for every $U\in\mathcal{U}$;
\item it satisfies the \emph{gauge clause}: on every cube $Q$ of side $6M$ there is a $G$-valued $u$
with $U^u=e^{i\xi A}$, and a radius $r_{\mathrm{reg}}\ge 12\gamma_+^{\mathsf{Q}}/\gamma_{\oplus}$,
such that the path $\mathsf{u}_z:=(e^{iz\xi\eta A}\cdot 1,0)$ lies in $\mathcal{U}(1)$ for
$|z|\le r_{\mathrm{reg}}$.
\end{itemize}
Here $\eta$ is the cut-off of that lemma, with values in $[0,1]$ and Lipschitz constant $2/M$, and
$\gamma_+^{\mathsf{Q}}$ is the constant of that lemma entering its radius condition. The remaining
hypotheses of that lemma hold: the walk hypothesis; the lower bound
$\mathsf{A}(1,0)\ge\gamma_{\oplus}>0$ at the pair $(1,0)$ of every nest; and its floors on $M$. Under
all of these, for every $U\in\mathcal{U}$ and $\mathbf{J}=0$ the operator $\mathsf{A}(U,0)$ is real
symmetric, and
\begin{equation*}
\mathsf{A}(U,0)\ge\gamma_0^X,\qquad \gamma_0^X:=\gamma_{\oplus}/2 .
\end{equation*}
\item[(c)] The ceiling on the regularity radius of Definition~\ref{4.4:def-radius-ceilings} holds at
the constant $\mathsf{a}_U$ of $\mathcal{U}_{\mathbb{R}}$, that is,
\begin{equation*}
r_{\mathrm{reg}}:=\frac{a_1^{\oplus}}{3\mathsf{a}_U}\ \ge\ \frac{12\gamma_+^{\mathsf{Q}}}{\gamma_{\oplus}} .
\end{equation*}
\item[(d)] The ratio $\mathcal{R}_{\mathbf{P}}$ of Definition~\ref{4.4:def-regular-backgrounds}
obeys the floor $\mathcal{R}_{\mathbf{P}}\ge 32\,\mathsf{B}/\gamma_0^X$.
\end{enumerate}
Put $\rho:=\mathcal{R}_{\mathbf{P}}\gamma_0^X/(16\,\mathsf{B})$. Then $\rho\ge 2$. Moreover, for every
$Y\subset\mathbb{B}(\Omega_{k+1})$ and at every point of $\mathbf{P}_{2\rho}$ ($\supset\mathbf{P}_{\rho}$),
\begin{equation}\label{4.4:eq-form-coercivity-1}
\operatorname{Re}\langle v,\mathsf{A}v\rangle\ \ge\ \gamma_*\,|v|^2
\qquad\text{for all } v\in\ell^2(Y;\mathfrak{g}^c),
\end{equation}
with $\gamma_*:=\tfrac34\gamma_0^X$. This is a number $\gamma_*=\gamma_*(d,L,G)>0$, and
$\gamma_*\le\mathsf{B}$.
\end{lemma}

\begin{proof}
From (d) and the definition of $\rho$ we get
\begin{equation*}
\rho=\frac{\mathcal{R}_{\mathbf{P}}\gamma_0^X}{16\,\mathsf{B}}\ge\frac{32}{16}=2 .
\end{equation*}

\emph{The base point.} We apply (b) with $\mathcal{U}=\mathcal{U}_{\mathbb{R}}$ of
Definition~\ref{4.4:def-regular-backgrounds}.

First, the region. The region $\Lambda_{k+1}$ is a union of cubes of side $LMR_{k+1}$
\cite{B-CMP119}, and
\begin{equation*}
\operatorname{dist}(\Lambda_{k+1},\Omega_{k+1}^c)\ge 2LMR_{k+1}
\end{equation*}
by \cite[(3.20)]{B-CMP119}. It is therefore a region of the class treated in the positivity lemma, and
so is the whole torus.

The positivity lemma asks four things of the space: invariance under $G$-valued gauge maps,
$1\in\mathcal{U}_{\mathbb{R}}$, $(U,0)\in\mathcal{U}(1)$, and the gauge clause. The first two hold by
Definition~\ref{4.4:def-regular-backgrounds}, and the third by the inclusion recorded in
Definition~\ref{4.4:def-radius-ceilings}. It remains to check the gauge clause for
$\mathcal{U}_{\mathbb{R}}$, as follows.

Let $U\in\mathcal{U}_{\mathbb{R}}$, and let $Q$ be a cube of side $6M$. Let $u$ be the $G$-valued gauge
transformation on $Q$ with $U^u=e^{i\xi A}$ and $|A|,|\nabla^{\xi}A|<2\mathsf{a}_U$. Since the cut-off
$\eta$ takes values in $[0,1]$ by (b), we have $|\eta A|<2\mathsf{a}_U$. Since $\eta$ has Lipschitz
constant $2/M$, the product rule for differences gives
\begin{equation*}
|\nabla^{\xi}(\eta A)|\le 2\mathsf{a}_U+\frac{2}{M}\,2\mathsf{a}_U
=2\mathsf{a}_U\Bigl(1+\frac{2}{M}\Bigr)\le 3\mathsf{a}_U ,
\end{equation*}
where the last inequality is $2/M\le\frac12$, which holds under the floors on $M$ assumed in (b).
Hence $\|(z\eta A,0)\|_{\star}\le 3|z|\mathsf{a}_U/a_1^{\oplus}$, where $\|\cdot\|_{\star}$ denotes the
norm on the fibre data of $\mathcal{U}(1)$ used in the positivity lemma, in which the slot of
$z\eta A$ is measured in units of $a_1^{\oplus}$. This is at most one for
$|z|\le r_{\mathrm{reg}}:=a_1^{\oplus}/(3\mathsf{a}_U)$, so the path $\mathsf{u}_z$ lies in
$\mathcal{U}(1)$ for these $z$, and $r_{\mathrm{reg}}$ serves as the radius of the gauge clause. The
requirement $r_{\mathrm{reg}}\ge12\gamma_+^{\mathsf{Q}}/\gamma_{\oplus}$ is exactly the ceiling (c).

The lower bound at the pair $(1,0)$ assumed in (b) is, for Ba{\l}aban's operators, the content of
\cite[(2.153), (2.157)]{B-CMP96}. For operators at $U=1$, \cite{B-CMP99b} refers for it to
\cite[Lemma~2.4]{B-CMP96}. In this book it is supplied by the lower bound at the trivial background.

Hence (b) applies. At every $U\in\mathcal{U}_{\mathbb{R}}$ the operator $\mathsf{A}(U,0)$ is real
symmetric, with $\mathsf{A}(U,0)\ge\gamma_0^X$ on real vectors. For complex $v=a+ib$, with $a,b$ real,
symmetry gives
\begin{equation*}
\langle v,\mathsf{A}(U,0)v\rangle=\langle a,\mathsf{A}(U,0)a\rangle+\langle b,\mathsf{A}(U,0)b\rangle
\ge\gamma_0^X|v|^2 .
\end{equation*}
So the bound holds on $\ell^2(\mathbb{B}(\Omega_{k+1});\mathfrak{g}^c)$.

A compression to a subset $Y$ inherits the bound. Indeed, extend $v\in\ell^2(Y;\mathfrak{g}^c)$ by zero.
This leaves both $\langle v,\mathsf{A}v\rangle$ and $|v|$ unchanged.

\emph{Two norm facts.} For a kernel $K=\sum_\omega K_\omega$, write
$\|K\|_0:=|||K|||_0$, the weighted norm of Definition~\ref{4.4:def-walk-systems} at weight zero.

First, $\|K\|_0$ bounds the operator norm. By the Cauchy--Schwarz inequality,
\begin{equation*}
|\langle u,Kv\rangle|\le\sum_{b,b'}|u(b)|\Bigl(\sum_\omega|K_\omega(b,b')|\Bigr)|v(b')|
\le|||K|||_0\,|u|\,|v| .
\end{equation*}
This is the Schur test, with the row and column sums both bounded by $|||K|||_0$.

Second, $|||K|||_0\le|||K|||_{\delta_{\mathsf{A}}}$, since $e^{\theta d(\omega;b,b')}\ge1$ for
$\theta\ge 0$. By (a), therefore, $\|\mathsf{A}\|_0\le\mathsf{B}$ at every point of $\mathcal{U}(1)$.

\emph{From $(U,0)$ to $\mathbf{P}_{2\rho}$.} Fix $(\mathbf{U},\mathbf{J})\in\mathbf{P}_{2\rho}$ over
the base $(U,0)$, with $\mathbf{U}=e^{i\xi A'}U$. Consider the complex line
\begin{equation*}
z\longmapsto\bigl(e^{iz\xi A'}U,\ z\mathbf{J}\bigr).
\end{equation*}
At $z=1$ it passes through $(\mathbf{U},\mathbf{J})$, and at $z=0$ through $(U,0)$.

For $|z|<\mathcal{R}_{\mathbf{P}}/(2\rho)$ the point $z$ of the line has $A'$-slot bounded by
\begin{equation*}
|z|\cdot 2\rho\alpha_1<\alpha_1'\le a_1^{\oplus},
\end{equation*}
and current slot less than $\gamma_0'\le a_0^{\oplus}$; here $\alpha_1'$, $\gamma_0'$, $a_0^{\oplus}$
and $a_1^{\oplus}$ are the constants of Definition~\ref{4.4:def-radius-ceilings}. The fibres of
$\mathcal{U}(1)$ over real backgrounds are convex and circled, by the fibre-convexity statement for
$\mathcal{U}(1)$. Hence this segment of the line lies in the fibre of $\mathcal{U}(1)$ over $U$ (the
condition on $|\partial\mathbf{U}-1|$ in the definition of $\mathbf{P}'$ is not preserved along the
line, which is why the bound of (a) is used on $\mathcal{U}(1)$).

Let $f(z)$ be $\mathsf{A}$ evaluated on the line. Then $f$ is analytic on the disc
$|z|<\mathcal{R}_{\mathbf{P}}/(2\rho)$, and $\|f(z)\|_0\le\mathsf{B}$ there by (a). Write
$f(z)=\sum_{m\ge0}a_mz^m$. By Cauchy's formula on circles of radius
$r<\mathcal{R}_{\mathbf{P}}/(2\rho)$, followed by the limit $r\to\mathcal{R}_{\mathbf{P}}/(2\rho)$,
\begin{equation*}
\|a_m\|_0\le\mathsf{B}\,(2\rho/\mathcal{R}_{\mathbf{P}})^m .
\end{equation*}
Set
\begin{equation*}
t:=\frac{2\rho}{\mathcal{R}_{\mathbf{P}}}=\frac{\gamma_0^X}{8\mathsf{B}} .
\end{equation*}
By the base point, $\mathsf{A}(U,0)\ge\gamma_0^X$, so its operator norm is at least $\gamma_0^X$. By the
first norm fact and (a), that operator norm is at most $\mathsf{B}$. Hence
$\gamma_0^X\le\|\mathsf{A}(U,0)\|\le\mathsf{B}$, and so $t\le\frac18$. In particular $z=1$ lies in the
disc, and $f(1)-f(0)=\sum_{m\ge1}a_m$. Therefore
\begin{equation}\label{4.4:eq-form-coercivity-2}
\|\mathsf{A}(\mathbf{U},\mathbf{J})-\mathsf{A}(U,0)\|_0\le\sum_{m\ge1}\mathsf{B}t^m
=\frac{\mathsf{B}t}{1-t}\le\frac{\gamma_0^X}{8}\cdot\frac87
=\frac{\gamma_0^X}{7}\le\frac{\gamma_0^X}{4}.
\end{equation}

Now take $v\in\ell^2(Y;\mathfrak{g}^c)$. Combine the base bound (compressed to $Y$) with
\eqref{4.4:eq-form-coercivity-2} and the first norm fact:
\begin{equation*}
\operatorname{Re}\langle v,\mathsf{A}(\mathbf{U},\mathbf{J})v\rangle
\ge\langle v,\mathsf{A}(U,0)v\rangle
-\bigl|\langle v,(\mathsf{A}(\mathbf{U},\mathbf{J})-\mathsf{A}(U,0))v\rangle\bigr|
\ge\Bigl(1-\frac17\Bigr)\gamma_0^X|v|^2 .
\end{equation*}
The right-hand side is $\frac67\gamma_0^X|v|^2\ge\frac34\gamma_0^X|v|^2$, which is
\eqref{4.4:eq-form-coercivity-1} with $\gamma_*=\frac34\gamma_0^X$.

Since $\gamma_0^X=\gamma_{\oplus}/2$, the constant $\gamma_*$ depends on $(d,L,G)$ only, and
\begin{equation*}
\gamma_*=\tfrac34\gamma_0^X\le\gamma_0^X\le\mathsf{B}.
\end{equation*}
\end{proof}

\begin{remark}
The floor (d) is a floor on the ratio which \cite{B-CMP116} takes to be much larger than one. The
fourth clause of the lemma on compressions, resolvents and square roots, and the three clauses of the
localisation theorem on $\Lambda$, are algebraic consequences of \eqref{4.4:eq-form-coercivity-1}.
They therefore hold at every point of $\mathbf{P}_{2\rho}$.

The bound at the base point corresponds to two statements in the literature. A lower bound
$\gamma_0>0$, independent of $k$ and $U$, is stated in \cite{B-CMP99b}. The symmetry of the operators,
and the positivity of the measure, at the pair $(U,0)$ are stated in \cite{B-CMP116}. The argument
above does not rely on them: at a global base $(U,0)$ with $U\in\mathcal{U}_{\mathbb{R}}$, the lower
bound and the real symmetry of the total form $\mathsf{A}(U,0)$ are the conclusion of the positivity
lemma.
\end{remark}

\begin{definition}[Auxiliary parameters and the near--far split]\label{4.4:def-auxiliary-parameters}
Let $\mathsf{A}=\sum_{\omega}\mathsf{A}_{\omega}$ satisfy the walk-expansion hypothesis of
Definition~\ref{4.4:def-form-walk-hypothesis}, with norm bound $\mathsf{B}$, decay rate
$\delta_{\mathsf{A}}$ and walk-domain radius $\rho_w$, and let $\gamma_*$ be the constant of
the coercivity hypothesis of Lemma~\ref{4.4:lem-form-coercivity}. The four numbers
$\mathsf{B},\delta_{\mathsf{A}},\gamma_*,\rho_w$ depend only on $d$, $L$ and the gauge group, and
are fixed before $M$. The numbers defined below are the same in every setting in which they are
used. They are chosen after $\mathsf{B},\delta_{\mathsf{A}},\gamma_*,\rho_w$, in the order
(k1), (k2), (k2$'$), (k3), and $M$ is chosen last.

For $r'>0$ put
\begin{equation}\label{4.4:eq-auxiliary-parameters-1}
\begin{gathered}
r:=r'+2\rho_w,\qquad
\varepsilon(r'):=\mathsf{B}\,e^{-\delta_{\mathsf{A}}r'/2},\\
\theta_0:=\min\Bigl\{\frac{\gamma_*}{4e\,\mathsf{B}\,r},\ \frac{\delta_{\mathsf{A}}}{2}\Bigr\},
\qquad \theta:=\frac{\theta_0}{4},\\
S:=\sup_{b}\sum_{b'\in\mathbb{B}(T)}e^{-(\theta_0/2)|b-b'|}\ \le\ d\,(1+4/\theta_0)^{d}.
\end{gathered}
\end{equation}
Here $T$, its set of bonds $\mathbb{B}(T)$, the distance $|b-b'|$ between bonds and the extent
$d(\omega)$ of a walk $\omega$ are those of Definition~\ref{4.4:def-form-walk-hypothesis}.
The letter $r$ here is distinct from the exponent $r$ of \cite{B-CMP119}.
\begin{enumerate}
\item[(k1)] Fix $r'$ with $8S\varepsilon(r')\le\gamma_*$.
\item[(k2)] Fix $\ell\in L^{\mathbb{N}}$ with $\ell\ge 2r+4$ and
\begin{equation*}
\varrho:=\frac{2\mathsf{B}}{\gamma_*}\,S\,e^{-\theta_0\ell/8}\le\frac12 .
\end{equation*}
\item[(k2$'$)] In addition, $\ell$ satisfies the second lower bound
\begin{equation}\label{4.4:eq-auxiliary-parameters-2}
2^{d+1}c_h^{2}\,r^{2}\,\mathsf{B}\ \le\ \tfrac18\,\gamma_*\,\ell^{2},
\qquad c_h:=\sup|h'|,
\end{equation}
where $h$ is the profile fixed in the localised coercivity lemma of this chapter. The profile
$h$ is an absolute function, chosen before all the numbers of this definition. The factor
$2^{d+1}$ is twice the factor $2^{d}$ used in the localised coercivity lemma; the surplus
leaves a margin $\tfrac1{16}$ in that lemma.
\item[(k3)] $M\ge 48\ell$.
\end{enumerate}

\emph{Cells.} Let $\mathcal{Q}$ be the standard partition of $T$ into cubes $q$ of side $\ell$.
For $q\in\mathcal{Q}$ let $q^*$ be the concentric cube of side $3\ell$, and let $q^{\sharp}$ be
the concentric cube of side $3\ell+2r+4$; by (k2) the side of $q^{\sharp}$ is at most $4\ell$.
For a bond $b$ with endpoints $b_-,b_+$, ``$b\subset q^*$'' means that both endpoints of $b$
lie in $q^*$. Put
\begin{equation}\label{4.4:eq-auxiliary-parameters-3}
\kappa_q(b):=\tfrac12\bigl(\mathbf{1}[b_-\in q]+\mathbf{1}[b_+\in q]\bigr),
\end{equation}
and let $\kappa_q$ also denote the diagonal operator
$(\kappa_q v)(b)=\kappa_q(b)\,v(b)$ on $\ell^2(\mathbb{B}(T);\mathfrak{g}^c)$. These operators
satisfy $\sum_{q\in\mathcal{Q}}\kappa_q=I$.

\emph{Near and far.} Put
\begin{equation}\label{4.4:eq-auxiliary-parameters-4}
\mathsf{A}_{\mathrm{loc}}:=\sum_{d(\omega)\le r'}\mathsf{A}_{\omega},\qquad
\mathsf{A}_{\mathrm{far}}:=\mathsf{A}-\mathsf{A}_{\mathrm{loc}} .
\end{equation}
\end{definition}

\begin{proof}
We check that the choices can be made in the stated order and verify the assertions about
$q^{\sharp}$ and $\kappa_q$.

\emph{(k1).} By the definition of $\theta_0$ in \eqref{4.4:eq-auxiliary-parameters-1},
\begin{equation*}
\frac{4}{\theta_0}
=\max\Bigl\{\frac{16e\,\mathsf{B}\,r}{\gamma_*},\ \frac{8}{\delta_{\mathsf{A}}}\Bigr\}
\le\frac{16e\,\mathsf{B}\,(r'+2\rho_w)}{\gamma_*}+\frac{8}{\delta_{\mathsf{A}}} .
\end{equation*}
The bound on $S$ recorded in \eqref{4.4:eq-auxiliary-parameters-1} then gives
\begin{equation*}
S\le d\Bigl(1+\frac{16e\,\mathsf{B}\,(r'+2\rho_w)}{\gamma_*}+\frac{8}{\delta_{\mathsf{A}}}\Bigr)^{d}.
\end{equation*}
The right-hand side is a polynomial of degree $d$ in $r'$. Its coefficients depend only on
$d,\mathsf{B},\delta_{\mathsf{A}},\gamma_*,\rho_w$. On the other hand,
$\varepsilon(r')=\mathsf{B}e^{-\delta_{\mathsf{A}}r'/2}$ with $\delta_{\mathsf{A}}>0$. Hence
$8S\varepsilon(r')\to0$ as $r'\to\infty$, and some $r'$ satisfies $8S\varepsilon(r')\le\gamma_*$.

\emph{(k2) and (k2$'$).} Once $r'$ is fixed, the numbers $r$, $\theta_0>0$ and $S$ are fixed.
Let $\ell\to\infty$ through $L^{\mathbb{N}}$. Then $\ell\ge2r+4$ eventually holds. Moreover,
$\varrho=(2\mathsf{B}/\gamma_*)Se^{-\theta_0\ell/8}$ tends to $0$. The left-hand side of
\eqref{4.4:eq-auxiliary-parameters-2} does not depend on $\ell$, since $c_h$ is fixed with $h$.
Its right-hand side grows like $\ell^2$. So all three conditions hold for every sufficiently
large $\ell\in L^{\mathbb{N}}$, and we fix one such $\ell$.

\emph{(k3).} This is a lower bound on $M$ alone. Since $M$ is chosen after $\ell$, it can be met.

\emph{Side of $q^{\sharp}$.} By (k2), $2r+4\le\ell$. Hence $3\ell+2r+4\le4\ell$, so
$q^{\sharp}$ has side at most $4\ell$.

\emph{Partition of the identity.} Since $\mathcal{Q}$ is a partition of $T$, each endpoint of a
bond $b$ lies in exactly one cube $q\in\mathcal{Q}$. Therefore
$\sum_{q}\mathbf{1}[b_-\in q]=\sum_{q}\mathbf{1}[b_+\in q]=1$, and by
\eqref{4.4:eq-auxiliary-parameters-3}, $\sum_{q}\kappa_q(b)=1$ for every bond $b$. Each
$\kappa_q$ acts by multiplication by $\kappa_q(b)$ at $b$. Hence
$\sum_{q\in\mathcal{Q}}\kappa_q=I$.

\emph{Covariance.} Let $\sigma$ be a lattice symmetry with $\sigma\mathcal{Q}=\mathcal{Q}$. It
maps a bond $b$ to the bond $\sigma b$ whose endpoints are $\sigma b_-$ and $\sigma b_+$,
possibly in exchanged order. The right-hand side of \eqref{4.4:eq-auxiliary-parameters-3} is
symmetric in the two endpoints of $b$, and $\mathbf{1}[\sigma x\in\sigma q]=\mathbf{1}[x\in q]$.
Hence $\kappa_{\sigma q}(\sigma b)=\kappa_q(b)$.

Consider the sharp assignment, which assigns each bond $b$ to the cube of $\mathcal{Q}$
containing $b_-$. Let $\sigma$ be a reflection, preserving $\mathcal{Q}$, in a
hyperplane orthogonal to $b$ that contains neither endpoint of $b$. Then $\sigma$ exchanges the
roles of the two endpoints. If $b_-$ and $b_+$ lie in different cubes of $\mathcal{Q}$, the cube
assigned to $\sigma b$ is the image of the cube containing $b_+$, not of the cube containing
$b_-$. So the sharp assignment is not covariant.
\end{proof}

Because $r=r'+2\rho_w$, the values of $\theta_0$, $\theta$, $S$, $r'$ and $\ell$, and of every
lower bound on $M$ that involves them, depend on the radius $\rho_w$. The form of no bound
depends on it.

\begin{remark}
The weights $\kappa_q$ are used in place of the sharp assignment of each bond to the cube
containing $b_-$. That assignment is not covariant under reflections, which exchange $b_-$ and
$b_+$. By contrast, $\kappa_q$ is covariant under every lattice symmetry preserving
$\mathcal{Q}$. This covariance enters the termwise covariance of the remainder of the
localisation later in this chapter.
\end{remark}

\begin{lemma}[Cell-resolvent expansion of the compressed resolvents]\label{4.4:lem-cell-resolvent}
Let $\mathsf{A}=\sum_\omega\mathsf{A}_\omega$ satisfy hypothesis $(\mathrm{W}')$ of
Definition~\ref{4.4:def-form-walk-hypothesis}. Let $Y$ be a set of bonds of $\Omega_{k+1}$ as in
Lemma~\ref{4.4:lem-form-coercivity}, satisfying the coercivity condition $(\gamma)$ of that
lemma. Assume the conditions (k1) and (k2) on the
auxiliary parameters (Definition~\ref{4.4:def-auxiliary-parameters}).

The following objects are those of Definitions~\ref{4.4:def-walk-systems}
and~\ref{4.4:def-auxiliary-parameters}:
\begin{itemize}
\item walks, their domains and footprints $|\cdot|$;
\item the cells $q$ with the cubes $q^*$ and $q^\sharp$, and the operators $\kappa_q$ attached to
the cells;
\item the near--far split $\mathsf{A}=\mathsf{A}_{\mathrm{loc}}+\mathsf{A}_{\mathrm{far}}$;
\item the lengths $d(\omega)$ and $d(\omega;b,b')$ of walks, the radii $r'$ and
$r=r'+2\rho_w$, the small quantity $\varepsilon(r')$, and the decay parameters $\theta_0$
and $\theta$;
\item the constant $S$ and the norms ${|\!|\!|}\cdot{|\!|\!|}_\theta$.
\end{itemize}
For a set of bonds $X$, $\chi_X$ denotes multiplication by its indicator. We write
$\mathsf{A}^Y:=\chi_Y\mathsf{A}\chi_Y$ for the compression of $\mathsf{A}$ to
$\ell^2(Y;\mathfrak{g}^c)$. Then, for every $x\ge 0$, the following hold.
\begin{enumerate}
\item[(i)] The operator $\mathsf{R}^Y(x)=(x+\mathsf{A}^Y)^{-1}$ exists on the polydisc
$\mathbf{P}_\rho$ of Lemma~\ref{4.4:lem-form-coercivity}, and
$\mathsf{R}^Y(x)=\sum_\varpi\mathsf{R}^Y_\varpi(x)$. The sum runs over walks $\varpi$ each of
whose steps is either a cell $q$, with domain $q^\sharp$, or a far walk $\omega$ of $\mathsf{A}$,
with the domains of $\omega$. Explicitly, put
\begin{equation*}
c^*_q:=\{b\in Y: b\subset q^*\},\qquad T_q:=\chi_{c^*_q}\mathsf{A}_{\mathrm{loc}}\chi_{c^*_q},
\end{equation*}
and restrict $\kappa_q$ to $Y$ (its support lies in $c^*_q$). Set
\begin{equation}\label{4.4:eq-cell-resolvent-1}
P_q(x):=\chi_{c^*_q}(x+T_q)^{-1}\kappa_q,\qquad
B_q(x):=\chi_{Y\setminus c^*_q}\mathsf{A}_{\mathrm{loc}}\chi_{c^*_q}(x+T_q)^{-1}\kappa_q .
\end{equation}
Then
\begin{equation}\label{4.4:eq-cell-resolvent-2}
\begin{aligned}
\mathsf{R}^Y(x)&=\sum_{m\ge0}\mathsf{N}(x)\bigl[-\chi_Y\mathsf{A}_{\mathrm{far}}\chi_Y\mathsf{N}(x)\bigr]^m,
\\
\mathsf{N}(x)&:=(x+\chi_Y\mathsf{A}_{\mathrm{loc}}\chi_Y)^{-1}
=\sum_{n\ge0}\sum_{q_0,\dots,q_n}P_{q_0}(-B_{q_1})\cdots(-B_{q_n}).
\end{aligned}
\end{equation}
\item[(ii)] The bound
\begin{equation*}
{|\!|\!|}\mathsf{R}^Y(x){|\!|\!|}_\theta\le 4S\bigl(x+\tfrac12\gamma_*\bigr)^{-1}
\end{equation*}
holds uniformly in $Y$, in $x$, in the setting label $\mathfrak{s}$, and on $\mathbf{P}_\rho$.
Every $\mathsf{R}^Y_\varpi(x)$ is analytic on $\mathbf{P}_\rho$.
\item[(iii)] Put
\begin{equation*}
\mathsf{M}^Y:=\pi^{-1}\int_0^\infty x^{-1/2}\,\mathsf{R}^Y(x)\,dx,\qquad
\mathsf{M}^Y_\varpi:=\pi^{-1}\int_0^\infty x^{-1/2}\,\mathsf{R}^Y_\varpi(x)\,dx .
\end{equation*}
Then $\mathsf{M}^Y=\sum_\varpi\mathsf{M}^Y_\varpi$; moreover
${|\!|\!|}\mathsf{M}^Y{|\!|\!|}_\theta\le 4S(2/\gamma_*)^{1/2}$ and
$(\mathsf{M}^Y)^2=(\mathsf{A}^Y)^{-1}$.
\item[(iv)] (Intrinsic terms.) The operators $\mathsf{R}^Y_\varpi(x)$ and $\mathsf{M}^Y_\varpi$ are
functions of $x$, of the matrices $\{\mathsf{A}_\omega:|\omega|\subset|\varpi|\}$ and of the set
$Y\cap|\varpi|$ only. Their kernels vanish unless both arguments lie in $|\varpi|$ at distance
at least $1$ from $\partial|\varpi|$.
\item[(v)] (Local form.) Assume in addition the positivity hypothesis, the inclusion condition
on the analyticity polydiscs and the regularity hypothesis fixed in the lemma establishing the
local form of the cell-resolvent expansion, below. Let $W'\subset\Omega_{k+1}$ be a closed union of
$M_1$-cubes (the set $W'$ is unrelated to hypothesis $(\mathrm{W}')$), and let $\mathsf{u}$ be
a local datum in the polydisc $\mathbf{P}_{\rho,W'}$ of that lemma. Then
the terms $\mathsf{R}^Y_\varpi(x)$ and $\mathsf{M}^Y_\varpi$ with $|\varpi|\subset W'$ are defined
and analytic in $\mathsf{u}$. Moreover the families
\begin{equation*}
\mathsf{R}^{Y,W'}(x):=\bigl(\mathsf{R}^Y_\varpi(x)\bigr)_{|\varpi|\subset W'},\qquad
\mathsf{M}^{Y,W'}:=\bigl(\mathsf{M}^Y_\varpi\bigr)_{|\varpi|\subset W'}
\end{equation*}
obey
\begin{equation*}
{|\!|\!|}\mathsf{R}^{Y,W'}(x){|\!|\!|}_\theta\le 4S\bigl(x+\tfrac12\gamma_*\bigr)^{-1},\qquad
{|\!|\!|}\mathsf{M}^{Y,W'}{|\!|\!|}_\theta\le 4S(2/\gamma_*)^{1/2}.
\end{equation*}
\end{enumerate}
\end{lemma}

\begin{proof}[Proof of Lemma~\ref{4.4:lem-cell-resolvent}, first three steps]
All operators are finite matrices, since the torus is finite. Kernels take values in
$\operatorname{End}(\mathfrak{g}^c)$. On kernel entries $|\cdot|$ denotes the operator norm of
$\operatorname{End}(\mathfrak{g}^c)$; on vectors of $\ell^2$ it denotes the $\ell^2$-norm; and
$\|\cdot\|$ is the operator norm on $\ell^2$. For bonds $b,b'$, $|b-b'|$ is their
$\ell^1$-distance.

\emph{First step (near and far walks).} Let $\omega$ be a walk with $d(\omega)\le r'$. By
$(\mathrm{W}')$ each domain of $\omega$ has $\ell^1$-diameter at most $\rho_w$; since
$d(\omega)\le r'$, all of them meet a tree of length at most $r'$. Take two points of
$|\omega|$, lying in domains $D$ and $D'$, and join them
through points of $D$ and $D'$ on the tree. Their distance is then at most
$\rho_w+r'+\rho_w$. Hence
\begin{equation*}
\operatorname{diam}|\omega|\le r'+2\rho_w=r .
\end{equation*}
Since $\mathsf{A}_{\mathrm{loc}}$ is the sum of the $\mathsf{A}_\omega$ with $d(\omega)\le r'$, it
follows that $\mathsf{A}_{\mathrm{loc}}(b,b')=0$ for $|b-b'|>r$. Suppose moreover that $\omega$ has a
row in $q^*$. Then $|\omega|$ meets $q^*$, and since $\operatorname{diam}|\omega|\le r$,
$|\omega|$ lies in the $r$-neighbourhood $(q^*)^{+r}$ of $q^*$. By the construction of
$q^\sharp$ (Definitions~\ref{4.4:def-walk-systems} and~\ref{4.4:def-auxiliary-parameters}),
this neighbourhood is contained in $q^\sharp$ with a margin of width $2$. The same holds for a
column. Thus
\begin{equation}\label{4.4:eq-cell-resolvent-3}
d(\omega)\le r'\ \text{and}\ \omega\ \text{has a row or a column in}\ q^*
\quad\Longrightarrow\quad |\omega|\subset q^\sharp .
\end{equation}

On far walks, (n3), one of the conditions fixed with the walk systems and the near--far split
in Definitions~\ref{4.4:def-walk-systems} and~\ref{4.4:def-auxiliary-parameters}, gives
$d(\omega;b,b')\ge d(\omega)>r'$. Hence, on far walks,
\begin{equation*}
e^{\frac12\delta_{\mathsf{A}}d(\omega;b,b')}\le e^{-\frac12\delta_{\mathsf{A}}r'}
e^{\delta_{\mathsf{A}}d(\omega;b,b')}.
\end{equation*}
The sums defining ${|\!|\!|}\mathsf{A}_{\mathrm{far}}{|\!|\!|}_{\delta_{\mathsf{A}}/2}$ run over a
subset of the walks of $\mathsf{A}$. Therefore
\begin{equation*}
{|\!|\!|}\mathsf{A}_{\mathrm{far}}{|\!|\!|}_{\delta_{\mathsf{A}}/2}
\le e^{-\delta_{\mathsf{A}}r'/2}\,{|\!|\!|}\mathsf{A}{|\!|\!|}_{\delta_{\mathsf{A}}}\le\varepsilon(r').
\end{equation*}
In the same way $(\mathrm{W}')$ gives
${|\!|\!|}\mathsf{A}_{\mathrm{loc}}{|\!|\!|}_{\delta_{\mathsf{A}}}\le\mathsf{B}$.
Multiplying by $\chi_Y$ on either side only sets kernel entries to zero. It therefore does not
increase these norms.

\emph{Second step (coercivity).} Let $v\in\ell^2(Y;\mathfrak{g}^c)$, and write
$\mathsf{A}_{\mathrm{loc}}=\mathsf{A}-\mathsf{A}_{\mathrm{far}}$. By $(\gamma)$,
$\operatorname{Re}\langle v,\mathsf{A}v\rangle\ge\gamma_*|v|^2$. Since $v=\chi_Yv$, we have
$\langle v,\mathsf{A}_{\mathrm{far}}v\rangle=\langle v,\chi_Y\mathsf{A}_{\mathrm{far}}\chi_Yv\rangle$;
by the bound on $\mathsf{A}_{\mathrm{far}}$ of the first
step, the row and column sums of $|\chi_Y\mathsf{A}_{\mathrm{far}}\chi_Y(b,b')|$ are at most
$\varepsilon(r')$ (the weights being at least $1$), so the Schur test gives
$|\langle v,\mathsf{A}_{\mathrm{far}}v\rangle|\le\varepsilon(r')|v|^2$. By the conditions (n2),
fixed together with (n3), and (k1), with $S\ge1$,
$\varepsilon(r')\le\frac18\gamma_*$. Hence
\begin{equation*}
\operatorname{Re}\langle v,\mathsf{A}_{\mathrm{loc}}v\rangle\ge(\gamma_*-\varepsilon(r'))|v|^2
\ge\tfrac78\gamma_*|v|^2 .
\end{equation*}
For $X\subset Y$ put $T_X:=\chi_X\mathsf{A}_{\mathrm{loc}}\chi_X$. For $v\in\ell^2(X)$ we have
$\langle v,T_Xv\rangle=\langle v,\mathsf{A}_{\mathrm{loc}}v\rangle$, because $v=\chi_Xv$ and
$\ell^2(X)\subset\ell^2(Y)$. Hence
$\operatorname{Re}\langle v,T_Xv\rangle\ge\frac78\gamma_*|v|^2$ on $\ell^2(X)$ for every
$X\subset Y$; in short, $\operatorname{Re}T_X\ge\frac78\gamma_*$.

\emph{Third step (Combes--Thomas bound, uniform in $x$ and $X$; cf.\ \cite{CT73}).} Let
$\theta_0$ be the decay parameter fixed among the auxiliary parameters
(Definition~\ref{4.4:def-auxiliary-parameters}). Let
$X\subset Y\cap q^*$. Here $q^*$ is a cube of side $3\ell$, $\ell$ the cell side of
Definition~\ref{4.4:def-auxiliary-parameters}, on which the coordinates are
single-valued, so $\vartheta\cdot b$ is defined for bonds $b\in X$ and $\vartheta\in\mathbb{R}^d$.
Let $|\vartheta|_\infty\le\theta_0$, and let $e^{\pm\vartheta\cdot b}$ act by multiplication.
Then
\begin{equation*}
e^{\vartheta\cdot b}\,T_X\,e^{-\vartheta\cdot b}=T_X+F_\vartheta,\qquad
F_\vartheta(b,b')=\bigl(e^{\vartheta\cdot(b-b')}-1\bigr)\,T_X(b,b').
\end{equation*}
By the first step, $T_X(b,b')$ vanishes unless $|b-b'|\le r$, so only this range matters.

By the choice of $\theta_0$ there, $e\,\theta_0r\,\mathsf{B}\le\gamma_*/4$; this enters the
bound on $F_\vartheta$ below. We also use $\gamma_*\le\mathsf{B}$. If $X=\emptyset$ there is
nothing to prove, so let $X\neq\emptyset$, hence $Y\neq\emptyset$. At any point of
$\mathbf{P}_\rho$, $(\gamma)$ and the Schur bound
$\|\chi_Y\mathsf{A}\chi_Y\|\le{|\!|\!|}\mathsf{A}{|\!|\!|}_{\delta_{\mathsf{A}}}\le\mathsf{B}$,
applied to a non-zero $v\in\ell^2(Y;\mathfrak{g}^c)$, give
$\gamma_*|v|^2\le\mathsf{B}|v|^2$, that is, $\gamma_*\le\mathsf{B}$. These two facts give
$\theta_0r\le1$. Hence, on the range $|b-b'|\le r$,
\begin{equation*}
|\vartheta\cdot(b-b')|\le|\vartheta|_\infty|b-b'|\le\theta_0r\le1 .
\end{equation*}
For $|z|\le1$ one has $|e^z-1|\le|z|e^{|z|}\le e|z|$, so $|e^{\vartheta\cdot(b-b')}-1|\le
e\theta_0r$ there.

The weights in ${|\!|\!|}\cdot{|\!|\!|}_{\delta_{\mathsf{A}}}$ are at least $1$. Hence the row and
column sums of $|T_X(b,b')|$ are bounded by
${|\!|\!|}\mathsf{A}_{\mathrm{loc}}{|\!|\!|}_{\delta_{\mathsf{A}}}\le\mathsf{B}$.
The Schur test then gives
\begin{equation*}
\|F_\vartheta\|\le e\,\theta_0r\,\mathsf{B}\le\tfrac14\gamma_* .
\end{equation*}
For any $F$, $\operatorname{Re}\langle v,Fv\rangle\ge-\|F\|\,|v|^2$. Since $x\ge0$ is real, the
second step gives, on $\ell^2(X)$,
\begin{equation*}
\operatorname{Re}(x+T_X+F_\vartheta)\ge x+\tfrac78\gamma_*-\tfrac14\gamma_*=x+\tfrac58\gamma_* .
\end{equation*}
If $\operatorname{Re}\langle v,\Xi v\rangle\ge c|v|^2$ for an operator $\Xi$ and some $c>0$, then
$c|v|^2\le|\langle v,\Xi v\rangle|\le|v|\,|\Xi v|$. So $\Xi$ is injective, hence invertible on
the finite-dimensional space, and $\|\Xi^{-1}\|\le1/c$. Applied with $c=x+\frac58\gamma_*$, and
also at $\vartheta=0$, this shows that $x+T_X$ and $x+T_X+F_\vartheta$ are invertible and that
\begin{equation*}
\|(x+T_X+F_\vartheta)^{-1}\|\le\bigl(x+\tfrac58\gamma_*\bigr)^{-1}
\le\bigl(x+\tfrac12\gamma_*\bigr)^{-1}.
\end{equation*}

Inverting the conjugation identity gives
\begin{equation*}
e^{\vartheta\cdot b}(x+T_X)^{-1}e^{-\vartheta\cdot b}=(x+T_X+F_\vartheta)^{-1}.
\end{equation*}
In kernels this reads
\begin{equation*}
(x+T_X)^{-1}(b,b')=e^{-\vartheta\cdot(b-b')}(x+T_X+F_\vartheta)^{-1}(b,b').
\end{equation*}
For given $b,b'$ take $\vartheta_\mu=\theta_0\operatorname{sgn}(b-b')_\mu$. Then
$|\vartheta|_\infty\le\theta_0$ and $\vartheta\cdot(b-b')=\theta_0|b-b'|$. A kernel entry is
bounded by the operator norm. Hence
\begin{equation}\label{4.4:eq-cell-resolvent-4}
|(x+T_X)^{-1}(b,b')|\le\bigl(x+\tfrac12\gamma_*\bigr)^{-1}e^{-\theta_0|b-b'|}.
\end{equation}

Parts (i)--(iv) are deduced from the inclusion \eqref{4.4:eq-cell-resolvent-3}, the bounds on
$\mathsf{A}_{\mathrm{loc}}$ and $\mathsf{A}_{\mathrm{far}}$ of the first step, the coercivity of
the second step and the estimate \eqref{4.4:eq-cell-resolvent-4} in the lemma completing the
proof of the cell-resolvent expansion, below; part (v) is proved in the lemma establishing the
local form of the cell-resolvent expansion, below.
\end{proof}

\begin{proof}[Proof of Lemma~\ref{4.4:lem-cell-resolvent}, continued]
\label{4.4:prf-parametrix-root}
We keep the notation and the hypotheses (W$'$), ($\gamma$), (k1), (k2) of the lemma and of the
earlier part of its proof. In particular, for every cell $q$ of the partition $\mathcal{Q}$ the cell
resolvent $(x+T_q)^{-1}$ on $\ell^2(c^*_q)$ exists and obeys the Combes--Thomas bound (CT)
established there, uniformly in $q$. Throughout, $x\ge 0$; complex $x$ is discussed
in the remark following the proof.

\emph{The parametrix.} Let $\eta\in\ell^2(c^*_q)$. Since $c^*_q\subset Y$, we have
$\chi_Y=\chi_{c^*_q}+\chi_{Y\setminus c^*_q}$ and $\chi_{c^*_q}\eta=\eta$. Hence, on $\ell^2(Y)$,
\begin{equation}\label{4.4:eq-parametrix-root-1}
\bigl(x+\chi_Y\mathsf{A}_{\mathrm{loc}}\chi_Y\bigr)\chi_{c^*_q}\eta
=(x+T_q)\eta+\chi_{Y\setminus c^*_q}\mathsf{A}_{\mathrm{loc}}\chi_{c^*_q}\eta ,
\end{equation}
because $\chi_{c^*_q}\mathsf{A}_{\mathrm{loc}}\chi_{c^*_q}=T_q$.

A bond touching $q$ has both endpoints within distance $1$ of $q$, so it lies in $q^*$. Hence
$\operatorname{supp}\kappa_q\subset c^*_q$, and for $v\in\ell^2(Y)$ the vector
$\eta:=(x+T_q)^{-1}\kappa_q v$ is defined and lies in $\ell^2(c^*_q)$. With this choice,
$\chi_{c^*_q}\eta=P_q(x)v$, the first term on the right of \eqref{4.4:eq-parametrix-root-1} is
$\kappa_q v$, and the second is $B_q(x)v$. Summing over $q$ and using
$\sum_q\kappa_q=I_Y$, we obtain
\begin{equation}\label{4.4:eq-parametrix-root-2}
\bigl(x+\chi_Y\mathsf{A}_{\mathrm{loc}}\chi_Y\bigr)P=I_Y+B,
\qquad P:=\sum_q P_q,\quad B:=\sum_q B_q .
\end{equation}

We regard $\{P_q\}$ and $\{B_q\}$ as walk systems in which each term has the single domain
$q^\sharp$. The rows and columns of $P_q$ and $B_q$ lie in $q^\sharp$, so
$d(q;b,b')\le|b-b'|$ for every matrix element. It therefore suffices to bound the weighted row
and column sums with the weight $e^{\theta|b-b'|}$.

\emph{Norm of $P$.} By definition,
$P_q(b,b')=\chi_{c^*_q}(b)\,(x+T_q)^{-1}(b,b')\,\kappa_q(b')$. We use (CT), which is uniform
in $q$, together with $\sum_q\kappa_q(b')=1$. For a fixed row $b$ this gives
\begin{align*}
\sum_q\sum_{b'}|P_q(b,b')|\,e^{\theta|b-b'|}
&\le\sum_{b'\in Y}\Bigl[\sum_q\kappa_q(b')\Bigr]\bigl(x+\tfrac12\gamma_*\bigr)^{-1}
e^{-(\theta_0-\theta)|b-b'|}\\
&\le S\bigl(x+\tfrac12\gamma_*\bigr)^{-1},
\end{align*}
the lattice sum over $b'$ being at most $S$. For a fixed column $b'$, we first sum over $b$ and
then write the sum over $q$ as $\sum_q\kappa_q(b')$. The same bound results. Hence
\begin{equation*}
\lvert\lvert\lvert P\rvert\rvert\rvert_\theta\le S\bigl(x+\tfrac12\gamma_*\bigr)^{-1}.
\end{equation*}

\emph{Norm of $B$.} By definition,
\begin{equation*}
B_q(b,b')=\sum_{b''\in c^*_q}\mathsf{A}_{\mathrm{loc}}(b,b'')\,(x+T_q)^{-1}(b'',b')\,\kappa_q(b'),
\end{equation*}
where $b\not\subset q^*$, $b'$ touches $q$, and $|b-b''|\le r$.

Since $b\not\subset q^*$, one endpoint of $b$ lies at sup-distance greater than $\ell$ from $q$.
Hence the initial point $b_-$ lies at sup-distance greater than $\ell-1$ from $q$. On the other
hand, $b'_-$ lies within $1$ of $q$. Therefore
\begin{equation*}
|b-b'|\ge|b-b'|_\infty>\ell-2 .
\end{equation*}
By the triangle inequality and (k2),
\begin{equation*}
|b''-b'|\ge\ell-2-r\ge\ell/2 .
\end{equation*}
Using also $|b-b'|\le|b-b''|+|b''-b'|$, we obtain
\begin{equation}\label{4.4:eq-parametrix-root-3}
e^{\theta|b-b'|}e^{-\theta_0|b''-b'|}
\le e^{\theta|b-b''|}\,e^{-(\theta_0/2)|b''-b'|}\,e^{-\theta_0\ell/8}.
\end{equation}

For the row sums we use three further facts. First, $\sum_q\kappa_q(b')=1$. Second,
$(x+\tfrac12\gamma_*)^{-1}\le 2/\gamma_*$ for $x\ge0$. Third,
\begin{equation*}
\sum_{b''}|\mathsf{A}_{\mathrm{loc}}(b,b'')|\,e^{\theta|b-b''|}
\le\lvert\lvert\lvert\mathsf{A}\rvert\rvert\rvert_{\delta_{\mathsf{A}}}\le\mathsf{B},
\end{equation*}
which holds because $|b-b''|\le d(\omega;b,b'')$ for every term $\mathsf{A}_\omega$ and
$\theta\le\delta_{\mathsf{A}}$. We insert (CT) and \eqref{4.4:eq-parametrix-root-3}, extend the
sum over $b''\in c^*_q$ to all $b''$, and bound the lattice sum over $b'$ of
$e^{-(\theta_0/2)|b''-b'|}$ by $S$. For a fixed row $b$ this gives
\begin{align*}
\sum_q\sum_{b'}|B_q(b,b')|\,e^{\theta|b-b'|}
&\le\Bigl[\sum_{b''}|\mathsf{A}_{\mathrm{loc}}(b,b'')|\,e^{\theta|b-b''|}\Bigr]
\bigl(x+\tfrac12\gamma_*\bigr)^{-1}e^{-\theta_0\ell/8}\,S\\
&\le\mathsf{B}\,(2/\gamma_*)\,S\,e^{-\theta_0\ell/8}.
\end{align*}

For a fixed column $b'$, we proceed in three stages. We first sum over $b$, using the column
norm of $\mathsf{A}_{\mathrm{loc}}$. We then sum over $b''$. Finally we sum over $q$ with the
weight $\kappa_q(b')$. The same bound results.

The quantity $\mathsf{B}(2/\gamma_*)Se^{-\theta_0\ell/8}$ is at most $\tfrac12$ by (k2).
Hence
\begin{equation*}
\lvert\lvert\lvert B\rvert\rvert\rvert_\theta\le\tfrac12 .
\end{equation*}
No count of the cells meeting a given bond enters this bound.

The operator $x+\chi_Y\mathsf{A}_{\mathrm{loc}}\chi_Y$ is invertible, as shown earlier in the
proof, and $\|B\|\le\tfrac12$ by (n2). Hence $I+B$ is invertible, with
$(I+B)^{-1}=\sum_{n\ge0}(-B)^n$. Multiplying \eqref{4.4:eq-parametrix-root-2} on the right by
$(I+B)^{-1}$ and then on the left by $(x+\chi_Y\mathsf{A}_{\mathrm{loc}}\chi_Y)^{-1}$, we obtain
\begin{equation*}
\mathsf{N}(x)=P(I+B)^{-1}=P\sum_{n\ge0}(-B)^n .
\end{equation*}
By (n1), the series converges in $\lvert\lvert\lvert\cdot\rvert\rvert\rvert_\theta$, and
\begin{equation*}
\lvert\lvert\lvert\mathsf{N}(x)\rvert\rvert\rvert_\theta
\le\lvert\lvert\lvert P\rvert\rvert\rvert_\theta\sum_{n\ge0}2^{-n}
\le 2S\bigl(x+\tfrac12\gamma_*\bigr)^{-1}.
\end{equation*}

\emph{The far part.} Since $\mathsf{A}=\mathsf{A}_{\mathrm{loc}}+\mathsf{A}_{\mathrm{far}}$, we have
\begin{equation*}
x+\mathsf{A}^Y=\bigl(x+\chi_Y\mathsf{A}_{\mathrm{loc}}\chi_Y\bigr)+\chi_Y\mathsf{A}_{\mathrm{far}}\chi_Y .
\end{equation*}
The bound on $\mathsf{N}$ at $x\ge0$ gives
$\lvert\lvert\lvert\mathsf{N}(x)\rvert\rvert\rvert_\theta\le 4S/\gamma_*$. Here
$\theta\le\delta_{\mathsf{A}}/2$, and for such $\theta$ the far part of $\mathsf{A}$ is bounded
by $\varepsilon(r')$ in the weighted norm. By (n1) and (k1),
\begin{equation*}
\lvert\lvert\lvert\chi_Y\mathsf{A}_{\mathrm{far}}\chi_Y\mathsf{N}\rvert\rvert\rvert_\theta
\le\varepsilon(r')\cdot 4S/\gamma_*\le\tfrac12 .
\end{equation*}
Moreover,
\begin{equation*}
x+\mathsf{A}^Y=\bigl(I+\chi_Y\mathsf{A}_{\mathrm{far}}\chi_Y\mathsf{N}\bigr)
\bigl(x+\chi_Y\mathsf{A}_{\mathrm{loc}}\chi_Y\bigr).
\end{equation*}
Both factors are invertible, the first by the Neumann series. Hence
\begin{equation}\label{4.4:eq-parametrix-root-4}
\mathsf{R}^Y(x)=\mathsf{N}\sum_{m\ge0}\bigl(-\chi_Y\mathsf{A}_{\mathrm{far}}\chi_Y\mathsf{N}\bigr)^m,
\qquad
\lvert\lvert\lvert\mathsf{R}^Y(x)\rvert\rvert\rvert_\theta\le 4S\bigl(x+\tfrac12\gamma_*\bigr)^{-1}.
\end{equation}

We now expand $\mathsf{A}_{\mathrm{far}}=\sum_{d(\omega)>r'}\mathsf{A}_\omega$, and $\mathsf{N}$ as
above, in \eqref{4.4:eq-parametrix-root-4}. This gives clauses (i) and (ii) of
Lemma~\ref{4.4:lem-cell-resolvent}. In a nonvanishing term, consecutive steps have a bond in
common, since otherwise the matrix product vanishes. Hence the footprints of the walks are
connected. They also have connected interior, by the strengthened connectivity of the domains
$[\omega]$ of the terms of $\mathsf{A}$; this property of the footprints is needed further on.

The terms of the expansion are analytic, for two reasons. First, $x+T_q$ is an analytic matrix
function with $\operatorname{Re}(x+T_q)\ge x+\tfrac78\gamma_*$, so its inverse exists and is
analytic. Second, the $\mathsf{A}_\omega$ are analytic.

\emph{The square root.} The spectrum of $\mathsf{A}^Y$ lies in
$\{z:\operatorname{Re}z\ge\gamma_*\}$. For $z\notin\,]-\infty,0]$,
\begin{equation*}
z^{-1/2}=\pi^{-1}\int_0^\infty x^{-1/2}(x+z)^{-1}\,dx .
\end{equation*}
The Dunford calculus is multiplicative. Hence the operator
\begin{equation*}
\mathsf{M}^Y:=\pi^{-1}\int_0^\infty x^{-1/2}\,\mathsf{R}^Y(x)\,dx
\end{equation*}
satisfies $(\mathsf{M}^Y)^2=(\mathsf{A}^Y)^{-1}$. At real backgrounds it is the positive square
root of $(\mathsf{A}^Y)^{-1}$.

By clause (ii), the walk series of $\mathsf{R}^Y(x)$ converges in
$\lvert\lvert\lvert\cdot\rvert\rvert\rvert_\theta$ uniformly in $x$ after multiplication by
$(x+\tfrac12\gamma_*)$. It may therefore be integrated term by term against
$\pi^{-1}x^{-1/2}\,dx$. The resulting factor is
\begin{equation*}
\pi^{-1}\int_0^\infty x^{-1/2}\bigl(x+\tfrac12\gamma_*\bigr)^{-1}dx=(2/\gamma_*)^{1/2}.
\end{equation*}
This is clause (iii).

\emph{Intrinsic character of the terms.} The operators $P_q$ and $B_q$ are built from the
following data:
\begin{itemize}
\item $\kappa_q$ restricted to $Y$;
\item the set $c^*_q$;
\item the bonds of $Y\setminus c^*_q$ within distance $r$ of $q^*$;
\item the matrices $\chi\mathsf{A}_\omega\chi'$ with $d(\omega)\le r'$ having a row or a column
in $q^*$.
\end{itemize}
The margin of $q^\sharp$ over $q^*$ is $r+2$, so all the bonds in the first three items are bonds
of $Y$ at distance at least $1$ from $\partial q^\sharp$. For the matrices in the last item,
$|\omega|\subset q^\sharp$ by ($\dagger$).

A far step is $\chi_Y\mathsf{A}_\omega\chi_Y$. By the support condition in hypothesis (W$'$),
its kernel lives at distance at least $1$ inside $|\omega|$. The split into near and far terms
reads only $d(\omega)$, which is a function of $\omega$. The parameters $r'$ and $\ell$, the
partition $\mathcal{Q}$ and the point $x$ are not data of the setting. This is clause (iv).

Clauses (i)--(iv) are thus proved. Clause (v) of Lemma~\ref{4.4:lem-cell-resolvent} is
established separately in what follows.
\end{proof}

\begin{remark}
We record four observations on the construction in the proof above.
\begin{itemize}
\item The domain assigned to the cell terms is the footprint $q^\sharp$, which is strictly larger
than $q^*$. An inverse $(x+T)^{-1}$ built from the operator $\mathsf{A}$ compressed to $c^*_q$
without its expansion into terms would contain terms $\mathsf{A}_\omega$ with
$|\omega|\not\subset c^*_q$. Such an inverse would not be intrinsic to $c^*_q$. Here, instead,
$T_q$ is built from the expanded operator truncated at $d(\omega)\le r'$, and the remaining terms
are distributed over $B_q$ and $\mathsf{A}_{\mathrm{far}}$. By ($\dagger$) the inverse is then
intrinsic to the declared domain $q^\sharp$.
\item The weights $\kappa_q$ are symmetric, in place of sharp indicators of cells. This is so
that the construction is covariant under reflections.
\item The invertibility of the local operator, the Combes--Thomas bounds for the cell
resolvents, and the parametrix and far-part arguments above hold verbatim for complex $x$ with
$\operatorname{Re}x\ge0$.
\item The level $k$ enters nowhere in the construction or in the bounds.
\end{itemize}
\end{remark}

\begin{lemma}[Near-part coercivity on one cell]\label{4.4:lem-cell-coercivity}
Assume the following:
\begin{itemize}
\item hypothesis (W$'$) of Lemma~\ref{4.4:lem-form-coercivity}, together with its local half;
\item the positivity statement for the base setting $\mathcal{S}_0$;
\item hypothesis (k1), under which the cell-resolvent expansion (Lemma~\ref{4.4:lem-cell-resolvent}) is
stated;
\item the halved ceilings on the primed radii and the primed regularity condition of
Definition~\ref{4.4:def-radius-ceilings}.
\end{itemize}
Let $q\in\mathcal{Q}$ be a cell. Let $\mathsf{u}=(\mathbf{U},\mathbf{J})\in\mathbf{P}_{\rho,q^\sharp}$ be a
local datum on $q^\sharp$ in the sense of Definition~\ref{4.4:def-local-data-classes}. For a set $X$ of bonds
of $\mathbb{B}(T)$ lying in $q^*$ put
\begin{equation*}
T_X(\mathsf{u}):=\sum\bigl\{\chi_X\,\mathsf{A}_\omega\bigl(\mathsf{u}|_{[\omega]}\bigr)\,\chi_X\;:\;
d(\omega)\le r',\ [\omega]\subset q^\sharp\bigr\}.
\end{equation*}
Then
\begin{equation}\label{4.4:eq-cell-coercivity-1}
\operatorname{Re}\langle v,T_X(\mathsf{u})v\rangle\ \ge\ \tfrac78\,\gamma_*\,|v|^2
\qquad\text{for all } v\in\ell^2(X;\mathfrak{g}^c).
\end{equation}
\end{lemma}

\begin{proof}
Throughout, $\xi$ is the lattice spacing of the level read, and $\nabla$ denotes the difference quotient at spacing
$\xi$ between neighbouring parallel bonds. For a set $E$ of points, $E^{+r}$ denotes the set of points within
distance $r$ of $E$. We write $\mathbf{U}=U'U$ with $U$ group-valued and $U'=e^{i\xi A'}$. For a
$G$-valued $W$, $R(W)$ is the adjoint action.

\emph{Step (1): what $T_X$ reads.}
Let $\omega$ be a walk with $d(\omega)\le r'$ and $\chi_X\mathsf{A}_\omega\chi_X\neq0$. Then
$\mathsf{A}_\omega$ has a row and a column in $X\subset q^*$. By the support bracket for walks with
$d(\omega)\le r'$, therefore, $[\omega]\subset(q^*)^{+r}$. Recall that $q^\sharp=(q^*)^{+(r+2)}$, so that
$(q^*)^{+r}\subset q^\sharp$. By the locality property of the walk operators, each $\mathsf{A}_\omega$
depends on the datum only through its restriction to $[\omega]$. Hence $T_X$ is a function of
$\mathsf{u}|_{(q^*)^{+r}}$. The walks concerned are walks of the base setting $\mathcal{S}_0$. Their traces
are trivial, by the trace statement for the walk operators applied at radius $0$.

\emph{Step (2): one gauge.}
The set $q^\sharp$ is a cube of side at most $4\ell\le M/12$. The gauge condition of
Definition~\ref{4.4:def-local-data-classes}, in its local form, therefore gives a $G$-valued gauge
transformation $u$ on $q^\sharp$ such that $U^u=e^{i\xi A}$ there, with
\begin{equation*}
|A|<\mathsf{a}_U,\qquad |\nabla A|<\mathsf{a}_U .
\end{equation*}
By the gauge covariance of the walk operators, applied term by term,
\begin{equation*}
T_X(\mathsf{u}^u)=R(u)\,T_X(\mathsf{u})\,R(u)^{-1}.
\end{equation*}
Here $R(u)$ is unitary bond by bond. The transformed datum is
\begin{equation*}
\mathsf{u}^u=\bigl(e^{i\xi R(u)A'}e^{i\xi A},\,R(u)\mathbf{J}\bigr).
\end{equation*}
Its entries satisfy the following identities:
\begin{equation*}
|R(u)A'|=|A'|,\qquad \nabla_{U^u}R(u)A'=R(u)\nabla_U A',\qquad |R(u)\mathbf{J}|=|\mathbf{J}|.
\end{equation*}
So $\mathsf{u}^u$ obeys the same bounds as $\mathsf{u}$.

The operator $R(u)$ maps $\ell^2(X;\mathfrak{g}^c)$ onto itself isometrically. For every
$v\in\ell^2(X;\mathfrak{g}^c)$ we have
\begin{equation*}
\operatorname{Re}\langle v,T_X(\mathsf{u}^u)v\rangle
=\operatorname{Re}\langle R(u)^{-1}v,\,T_X(\mathsf{u})R(u)^{-1}v\rangle,
\end{equation*}
and $|R(u)^{-1}v|=|v|$. Hence \eqref{4.4:eq-cell-coercivity-1} for $\mathsf{u}$ is equivalent to
\eqref{4.4:eq-cell-coercivity-1} for $\mathsf{u}^u$, and we may assume $U=e^{i\xi A}$ on $q^\sharp$.

\emph{Step (3): a cut-off inside the margin.}
Put
\begin{equation*}
\eta_q(x):=\min\bigl\{1,\ \tfrac12\operatorname{dist}_\infty(x,T\setminus q^\sharp)\bigr\},
\end{equation*}
with distances in the units of level $k$. Then $\eta_q=1$ on $(q^*)^{+r}$ and $\eta_q=0$ off $q^\sharp$.
Moreover $\eta_q$ is $\tfrac12$-Lipschitz. On the bonds of $q^\sharp$ put
\begin{equation*}
\hat A:=\eta_qA,\qquad \hat A':=\eta_qA',\qquad \hat J:=\eta_q\mathbf{J},
\end{equation*}
with $\eta_q$ evaluated at the midpoint of the bond. On all other bonds of $T$ put
$\hat A:=\hat A':=\hat J:=0$. Set
\begin{equation*}
\hat U:=e^{i\xi\hat A},\qquad \hat{\mathsf{u}}:=\bigl(e^{i\xi\hat A'}\hat U,\ \hat J\bigr).
\end{equation*}
This is a global datum.

We first bound $\hat A$. Let $b,b'$ be neighbouring parallel bonds of $q^\sharp$. Then
\begin{equation*}
\hat A(b')-\hat A(b)=\eta_q(b')\bigl(A(b')-A(b)\bigr)+\bigl(\eta_q(b')-\eta_q(b)\bigr)A(b),
\end{equation*}
and $|\eta_q(b')-\eta_q(b)|\le\tfrac12\xi$. Hence $|\nabla\hat A|<\mathsf{a}_U+\tfrac12\mathsf{a}_U$.

Now let $b$ be a bond of $q^\sharp$ and $b'$ a bond outside $q^\sharp$. Then $\hat A(b')=0$. The midpoint of
$b$ lies within $3\xi/2$ of $T\setminus q^\sharp$, so $\eta_q(b)\le\tfrac34\xi$. The difference quotient is
then at most $\tfrac34|A(b)|$. At a pair of bonds both outside $q^\sharp$ it vanishes. Therefore, everywhere
on $T$,
\begin{equation*}
|\hat A|<\mathsf{a}_U,\qquad |\nabla\hat A|<\mathsf{a}_U+\tfrac34\mathsf{a}_U<2\mathsf{a}_U .
\end{equation*}
Taking the trivial gauge on every cube, these bounds give $\hat U\in\mathcal{U}_{\mathbb{R}}$.

In the same way, $|\hat A'|<\rho\alpha_1$. For the covariant difference of $\hat A'$, let $b,b'$ be
neighbouring parallel bonds of $q^\sharp$ and let $c$ be the bond that joins them; $c$ is a bond of
$q^\sharp$. Then
\begin{equation*}
\xi\nabla_{\hat U}\hat A'
=\eta_q(b')\bigl[R(\hat U(c))A'(b')-A'(b)\bigr]+\bigl(\eta_q(b')-\eta_q(b)\bigr)A'(b),
\end{equation*}
and
\begin{equation*}
R(\hat U(c))A'(b')-A'(b)=\xi\nabla_UA'+\bigl(R(\hat U(c))-R(U(c))\bigr)A'(b').
\end{equation*}

Since $A$ is Hermitian, $\operatorname{ad}A$ is Hermitian for the trace inner product on $\mathfrak{g}^c$,
and $R(e^{i\xi A})$ is an isometry. With $A=A(c)$ and $\eta_q=\eta_q(c)$,
\begin{equation*}
R(e^{i\xi\eta_qA})-R(e^{i\xi A})
=R(e^{i\xi A})\bigl(e^{-i\xi(1-\eta_q)\operatorname{ad}A}-1\bigr).
\end{equation*}
The norm of the left side therefore equals the norm of $e^{-i\xi(1-\eta_q)\operatorname{ad}A}-1$; as
$|e^{-is\lambda}-1|\le s|\lambda|$ for real $\lambda$ and $s\ge0$, that norm is at most
$\xi(1-\eta_q)\|\operatorname{ad}A\|$, and $0\le1-\eta_q\le1$. Hence
\begin{equation*}
\bigl\|R(e^{i\xi\eta_qA})-R(e^{i\xi A})\bigr\|\le c_{\mathfrak{g}}\,\xi\,|A|,
\end{equation*}
where $c_{\mathfrak{g}}$ is the Lie-algebra constant with $\|\operatorname{ad}A\|\le c_{\mathfrak{g}}|A|$.

The local datum satisfies $|\nabla_UA'|<\rho\alpha_1$ and $|A'|<\rho\alpha_1$, and
$|\eta_q(b')-\eta_q(b)|\le\tfrac12\xi$. Combining the last three displays therefore gives
\begin{equation*}
|\nabla_{\hat U}\hat A'|<\rho\alpha_1\bigl(1+c_{\mathfrak{g}}\mathsf{a}_U+\tfrac12\bigr)\le2\rho\alpha_1,
\end{equation*}
where the last step uses the condition $2c_{\mathfrak{g}}\mathsf{a}_U\le1$ on the constant $\mathsf{a}_U$.

At a pair with $b$ in $q^\sharp$ and $b'$ outside it, $\hat A'(b')=0$, and the difference quotient is at most
$\tfrac34|A'(b)|$, by the bound $\eta_q(b)\le\tfrac34\xi$ obtained above. Finally, $|\hat J|\le|\mathbf{J}|$.
Hence $\hat{\mathsf{u}}\in\mathbf{P}_{2\rho}$.

\emph{Step (4): conclusion.}
Since $\eta_q=1$ on $(q^*)^{+r}$, we have $\hat{\mathsf{u}}=\mathsf{u}$ on $(q^*)^{+r}$. Let
$\mathsf{A}_{\mathrm{loc}}(\hat{\mathsf{u}})$ be the near part, in the near--far split of
Definition~\ref{4.4:def-auxiliary-parameters}, of the operator of the global setting
$(\mathcal{S}_0;T,\hat{\mathsf{u}})$.

By Step (1), every near walk $\omega$ with $\chi_X\mathsf{A}_\omega\chi_X\ne0$ has
$[\omega]\subset(q^*)^{+r}$. On this set $\hat{\mathsf{u}}$ and $\mathsf{u}$ coincide, and it is contained in
$q^\sharp$. Consequently
\begin{equation*}
T_X(\mathsf{u})=\chi_X\,\mathsf{A}_{\mathrm{loc}}(\hat{\mathsf{u}})\,\chi_X .
\end{equation*}

Hypothesis $(\gamma)$ of Lemma~\ref{4.4:lem-form-coercivity} is formulated on $\mathbf{P}_{2\rho}$. For the
base setting $\mathcal{S}_0$ on $T$ with $Y:=\mathbb{B}(T)$ it is the positivity statement for the base
setting, which is assumed above, and $\hat{\mathsf{u}}\in\mathbf{P}_{2\rho}$ by Step (3); hypothesis
(W$'$) of Lemma~\ref{4.4:lem-form-coercivity} is assumed as well. Apply the second step of the proof
of Lemma~\ref{4.4:lem-form-coercivity} at the point $\hat{\mathsf{u}}$, with $Y:=\mathbb{B}(T)$. It gives
\begin{equation*}
\operatorname{Re}\langle v,\mathsf{A}_{\mathrm{loc}}(\hat{\mathsf{u}})v\rangle\ge\tfrac78\gamma_*|v|^2
\qquad\text{for all } v\in\ell^2(\mathbb{B}(T)).
\end{equation*}
For $v\in\ell^2(X;\mathfrak{g}^c)$ we have $\chi_Xv=v$, so
\begin{equation*}
\langle v,\chi_X\mathsf{A}_{\mathrm{loc}}(\hat{\mathsf{u}})\chi_Xv\rangle
=\langle v,\mathsf{A}_{\mathrm{loc}}(\hat{\mathsf{u}})v\rangle.
\end{equation*}
Compressing to $X$ in this way yields \eqref{4.4:eq-cell-coercivity-1}.
\end{proof}

\begin{remark}
\begin{enumerate}
\item In the collar $q^\sharp\setminus(q^*)^{+r}$ the curvature of $\hat U$ is of size $\mathsf{a}_U$, not
$\alpha_0'$. This is why the classes $\mathcal{U}_{\mathbb{R}}$ and $\mathbf{P}_\sigma$, for every radius
$\sigma$, are defined by gauge conditions only. Neither the domain $\mathcal{U}(1)$ of hypothesis (W$'$)
of Lemma~\ref{4.4:lem-form-coercivity} nor the proof of the positivity statement for the base setting
reads the curvature on a plaquette.
\item The cube $q^\sharp$ has side at most $4\ell$. This is less than the side of the torus $T$, by the
standing lower bound on that side.
\item Only the datum on a cube is ever extended. The topology of the ambient region therefore plays no part.
\end{enumerate}
\end{remark}

\begin{proof}[Proof of clause (v) of Lemma~\ref{4.4:lem-cell-resolvent}]
\label{4.4:prf-local-resolvent}
Let $\mathsf{u}\in\mathbf{P}_{\rho,W'}$. We define every object of the cell-resolvent expansion
from the datum $\mathsf{u}$ alone, using only the cells $q$ with $q^\sharp\subset W'$. For such a
cell put
\begin{equation*}
T_q:=T_{c^*_q}\bigl(\mathsf{u}{\restriction}q^\sharp\bigr),
\end{equation*}
where $T_X$ is the operator of near-part coercivity on one cell
(Lemma~\ref{4.4:lem-cell-coercivity}) with $X=c^*_q=\{b\in Y: b\subset q^*\}$. We define
$P_q$ and $B_q$ by the formulas of clause (i) of Lemma~\ref{4.4:lem-cell-resolvent},
\begin{align*}
P_q(x)&:=\chi_{c^*_q}(x+T_q)^{-1}\kappa_q,\\
B_q(x)&:=\chi_{Y\setminus c^*_q}\,\mathsf{A}_{\mathrm{loc}}\,\chi_{c^*_q}(x+T_q)^{-1}\kappa_q,
\end{align*}
in which $\mathsf{A}_{\mathrm{loc}}$ now stands for the sum of the kernels
$\mathsf{A}_\omega(\mathsf{u}{\restriction}[\omega])$ over those $\omega$ with $d(\omega)\le r'$
and a column in $q^*$. For each such $\omega$ the near-part inclusion of
Lemma~\ref{4.4:lem-cell-resolvent} gives
$[\omega]\subset q^\sharp$, and $q^\sharp\subset W'$, so the restriction
$\mathsf{u}{\restriction}[\omega]$ is defined. The far steps are the kernels
\begin{equation*}
\chi_Y\,\mathsf{A}_\omega\bigl(\mathsf{u}{\restriction}[\omega]\bigr)\,\chi_Y,
\qquad d(\omega)>r',\quad [\omega]\subset W',
\end{equation*}
and the term of a walk $\varpi$ is the product of its steps, taken as in clause (i).

The support $|\varpi|$ of a walk is the union of the domains of its steps: $q^\sharp$ for a
step that is a cell $q$, and $[\omega]$ for a far step $\omega$. Hence $|\varpi|\subset W'$ if
and only if every step of $\varpi$ belongs to the list just defined. The term of every walk
$\varpi$ with $|\varpi|\subset W'$ is therefore defined, and it is the function of clause (iv)
of Lemma~\ref{4.4:lem-cell-resolvent} evaluated at the family
$\{\mathsf{A}_\omega:|\omega|\subset|\varpi|\}$ and at the set $Y\cap|\varpi|$. The family
$\mathsf{R}^{Y,W'}(x)$ of clause (v) is thus the family of these terms, $\varpi$ ranging over
the walks with $|\varpi|\subset W'$.

\emph{Existence.} Let $q$ be a cell with $q^\sharp\subset W'$. Applied to the restriction
$\mathsf{u}{\restriction}q^\sharp$, Lemma~\ref{4.4:lem-cell-coercivity} gives, for every set
$X\subset Y\cap q^*$,
\begin{equation*}
\operatorname{Re}\langle v,T_X v\rangle\ \ge\ \tfrac78\gamma_*|v|^2
\qquad\text{for all } v\in\ell^2(X;\mathfrak{g}^c).
\end{equation*}
This bound takes the place of the second step of the proof of the preceding clauses of
Lemma~\ref{4.4:lem-cell-resolvent}. In particular, for $X=c^*_q$ and $x\ge0$,
\begin{equation*}
|v|\,\bigl|(x+T_q)v\bigr|\ \ge\ \operatorname{Re}\langle v,(x+T_q)v\rangle
\ \ge\ \bigl(x+\tfrac78\gamma_*\bigr)|v|^2 ,
\end{equation*}
so $x+T_q$ is injective on $\ell^2(c^*_q;\mathfrak{g}^c)$ and has closed range. Since
$\operatorname{Re}\langle v,(x+T_q)^*v\rangle=\operatorname{Re}\langle v,(x+T_q)v\rangle$, the
same lower bound holds for the adjoint, which is therefore injective, so the range of $x+T_q$
is dense. Hence $x+T_q$ is invertible, and its inverse has norm at most
$(x+\frac78\gamma_*)^{-1}$. Thus $P_q(x)$ and $B_q(x)$ are defined for every $x\ge0$.

\emph{The conclusion of the third step.} The third step of that proof holds verbatim. It reads
only three things: the range $r$ of the kernel of $T_X$; Schur's test, applied with the row and
column sums of $\sum_\omega|\mathsf{A}_\omega(b,b')|$ bounded by $\mathsf{B}$; and the lower bound
$\operatorname{Re}T_X\ge\frac78\gamma_*$.
The bounds by $\mathsf{B}$ are the datum-free majorants furnished by the local half of
hypothesis $(\mathrm{W}')$ of Lemma~\ref{4.4:lem-cell-resolvent}. The lower bound was
established in the preceding paragraph. Hence
the conclusion of the third step, the bound established at the end of that step, holds for the
operators $T_q$ defined from $\mathsf{u}$.

\emph{Norms.} The two estimates of the fourth step of that proof are estimates of row sums and
of column sums, whose terms are non-negative. Summed over the cells $q$ with
$q^\sharp\subset W'$ only, they hold with $\sum_q\kappa_q(b')\le1$ in place of
$\sum_q\kappa_q(b')=1$, and they give
\begin{equation*}
|\!|\!|(P_q)|\!|\!|_\theta\ \le\ S\bigl(x+\tfrac12\gamma_*\bigr)^{-1},
\qquad
|\!|\!|(B_q)|\!|\!|_\theta\ \le\ \varrho\ \le\ \tfrac12 ,
\end{equation*}
the norms being taken over the families indexed by these cells. The first step of that proof
gives, for the family of far steps,
\begin{equation*}
|\!|\!|(\mathsf{A}_\omega)_{d(\omega)>r',\,[\omega]\subset W'}|\!|\!|_{\delta_{\mathsf{A}}/2}
\ \le\ \varepsilon(r').
\end{equation*}
We now apply the product inequality for families of kernels used in the proof of the preceding
clauses of Lemma~\ref{4.4:lem-cell-resolvent}, which is a statement about families of kernels
and not about operators. By it, the family of products $P_{q_0}(-B_{q_1})\cdots(-B_{q_n})$ has
norm
\begin{equation*}
\le\ S\bigl(x+\tfrac12\gamma_*\bigr)^{-1}\sum_{n\ge0}\varrho^n
\ \le\ 2S\bigl(x+\tfrac12\gamma_*\bigr)^{-1},
\end{equation*}
the last inequality because $\varrho\le\frac12$. Since $x\ge0$, this bound is at most
$2S(\frac12\gamma_*)^{-1}=4S/\gamma_*$. Clause (i) expands the resolvent as
\begin{equation*}
\mathsf{R}^Y(x)=\sum_{m\ge0}\mathsf{N}(x)\bigl[-\chi_Y\mathsf{A}_{\mathrm{far}}\chi_Y
\mathsf{N}(x)\bigr]^m .
\end{equation*}
In the family $\mathsf{R}^{Y,W'}(x)$ the place of each occurrence of $\mathsf{N}(x)$ in the
expansion just displayed is taken by the family of products
$P_{q_0}(-B_{q_1})\cdots(-B_{q_n})$, and the place of $\chi_Y\mathsf{A}_{\mathrm{far}}\chi_Y$ by
the family of far steps; no operator $\mathsf{N}(x)$ is formed. Bounding the first family of
products by $2S(x+\frac12\gamma_*)^{-1}$, each of the $m$ further ones by $4S/\gamma_*$, and
each of the $m$ families of far steps by $\varepsilon(r')$, the product inequality gives, in
the norm it furnishes,
\begin{equation*}
|\!|\!|\mathsf{R}^{Y,W'}(x)|\!|\!|
\ \le\ 2S\bigl(x+\tfrac12\gamma_*\bigr)^{-1}
\sum_{m\ge0}\Bigl(\varepsilon(r')\,\frac{4S}{\gamma_*}\Bigr)^m
\ \le\ 4S\bigl(x+\tfrac12\gamma_*\bigr)^{-1},
\end{equation*}
where hypothesis $(\mathrm{k}1)$ of Lemma~\ref{4.4:lem-cell-resolvent} gives the last inequality.

The operator identities of the fourth and fifth steps of that proof,
\begin{equation*}
(x+\chi_Y\mathsf{A}_{\mathrm{loc}}\chi_Y)P=I_Y+B,
\qquad
\mathsf{N}=P(I+B)^{-1},
\end{equation*}
are not used here: there is no operator to invert. The object $\mathsf{R}^{Y,W'}(x)$ is an
absolutely summable family of kernels, and absolute summability of this family, term by term,
is all that its later application in this chapter requires. Where that application needs a
Gaussian whose covariance has an inverse with positive real part (the vertex Gaussians of
\cite[(2.8)--(2.9)]{B-CMP116}), the local coercivity lemma proved later in this chapter, which
is distinct from Lemma~\ref{4.4:lem-cell-coercivity}, supplies it:
for the choice of $Y$ and $W'$ made there, the sum of the family is the resolvent
$(x+\mathsf{T}_{W'})^{-1}$ of the operator $\mathsf{T}_{W'}$ of that lemma. Where the Gaussian
needs a real symmetric base (for the cancellation of
\cite[(2.23)--(2.26)]{B-CMP116}), the local symmetry lemma and its companion proposition in a
later section of this chapter supply it.

\emph{Analyticity} follows by the argument of the last sentence of the fifth step of that
proof. \emph{The family $\mathsf{M}^{Y,W'}$} is obtained by the termwise integration of the
sixth step of that proof, applied to the absolutely summable family $\mathsf{R}^{Y,W'}(x)$.
\end{proof}

\begin{lemma}[Coercivity of the localised quadratic form]\label{4.4:lem-local-coercivity}
Assume the hypotheses of Lemma~\ref{4.4:lem-cell-coercivity}. Fix once and for all a function
$h\in C^\infty(\mathbb{R})$ with $0\le h\le 1$ and $\operatorname{supp}h\subset(-1,1)$ such that
\begin{equation}\label{4.4:eq-local-coercivity-1}
\sum_{n\in\mathbb{Z}}h(t-n)^2=1\qquad\text{for all } t\in\mathbb{R},
\end{equation}
and put $c_h:=\sup|h'|$. Assume moreover that the cell side $\ell$ satisfies the lower bound
\begin{equation}\label{4.4:eq-local-coercivity-2}
2^dc_h^2r^2\mathsf{B}\,\ell^{-2}\le\tfrac1{16}\gamma_* .
\end{equation}
Let $W'$ be a closed union of $M_1$-cubes with $W'\subset\Omega_{k+1}$,
let $\mathsf{u}\in\mathbf{P}_{\rho,W'}$ be a local datum, and let $Y'\subset\mathbb{B}(T)$ be a set
of bonds such that $q^\sharp\subset W'$ for every cell $q\in\mathcal{Q}$ for which $q^*$ contains
a bond of $Y'$. Put
\begin{equation*}
\mathsf{T}_{W'}:=\chi_{Y'}\Big(\sum_{\omega:\ [\omega]\subset W'}
\mathsf{A}_\omega(\mathsf{u}\restriction[\omega])\Big)\chi_{Y'},
\end{equation*}
a finite matrix on $\ell^2(Y';\mathfrak{g}^c)$. Then
\begin{equation*}
\operatorname{Re}\langle v,\mathsf{T}_{W'}v\rangle\ge\tfrac58\,\gamma_*|v|^2
\qquad\text{for all } v\in\ell^2(Y';\mathfrak{g}^c).
\end{equation*}
\end{lemma}

\begin{proof}
\emph{The profile $h$.} A function $h$ as required exists. Take $\eta\in C^\infty_0((-1,1))$ with
$\eta\ge0$ and $\eta>0$ on $[-\tfrac12,\tfrac12]$. The function
$\Sigma_\eta:=\sum_{n\in\mathbb{Z}}\eta(\cdot-n)^2$ is
smooth (locally the sum has at most two non-zero terms), $1$-periodic, and strictly positive,
because every $t\in\mathbb{R}$ lies within $\tfrac12$ of an integer. Put $h:=\eta/\Sigma_\eta^{1/2}$.
Then $h$ is
smooth, $\operatorname{supp}h=\operatorname{supp}\eta\subset(-1,1)$, $0\le h\le1$ because
$\eta^2\le\Sigma_\eta$,
and by the periodicity of $\Sigma_\eta$,
\begin{equation*}
\sum_n h(t-n)^2=\Sigma_\eta(t)^{-1}\sum_n \eta(t-n)^2=1 .
\end{equation*}
This is a partition of unity of the kind used in \cite[(1.118)]{B-CMP95}; only the properties
listed in the statement are used below.

\emph{Splitting.} Write $\mathsf{T}_{W'}=\mathsf{T}_{\mathrm{near}}+\mathsf{T}_{\mathrm{far}}$,
where $\mathsf{T}_{\mathrm{near}}$ collects the terms with $d(\omega)\le r'$ and
$\mathsf{T}_{\mathrm{far}}$ those with $d(\omega)>r'$, according to the near--far split of
Definition~\ref{4.4:def-auxiliary-parameters}. Consider the family
$\{\mathsf{A}_\omega : d(\omega)>r',\ [\omega]\subset W'\}$. By the local half of the standing
hypothesis on the walk operators, its members have majorants that do not depend on the datum,
exactly as in part~\textup{(v)} of Lemma~\ref{4.4:lem-cell-resolvent}. The first step of the proof
of the
walk-operator bounds, applied to this family, therefore gives
$|\!|\!|\mathsf{T}_{\mathrm{far}}|\!|\!|_0\le\varepsilon(r')$. Since the norm
$|\!|\!|\cdot|\!|\!|_0$ dominates the
operator norm on $\ell^2$ (Schur bound), and since the condition on $r'$ among the hypotheses of
Lemma~\ref{4.4:lem-cell-coercivity} yields $\varepsilon(r')\le\tfrac18\gamma_*$, because the
constant $S$ occurring in that condition satisfies $S\ge1$, we obtain
\begin{equation}\label{4.4:eq-local-coercivity-3}
|\langle v,\mathsf{T}_{\mathrm{far}}v\rangle|\le\varepsilon(r')|v|^2\le\tfrac18\gamma_*|v|^2 .
\end{equation}
The first step of the proof of the walk-operator bounds also shows that
$\mathsf{T}_{\mathrm{near}}(b,b')=0$ unless $b,b'\in[\omega]$ for some
walk $\omega$ with $\operatorname{diam}[\omega]\le r$; in particular
$\mathsf{T}_{\mathrm{near}}(b,b')=0$ unless $|b_--b'_-|_1\le r$, where $b_-$ denotes the initial
point of the bond $b$. Moreover $|\!|\!|\mathsf{T}_{\mathrm{near}}|\!|\!|_0\le\mathsf{B}$.

\emph{Weights.} For a cell $q\in\mathcal{Q}$ with centre $c(q)$ and a bond $b$ put
\begin{equation*}
h_q(b):=\prod_{\mu=1}^{d}h\big((b_--c(q))_\mu/\ell\big),
\end{equation*}
the coordinate differences being read on the torus modulo its side. By the lower bound on the
torus side and by the second part of the condition relating the torus side to the cell side, the
side of the torus is a multiple of $\ell$ and at least $2\ell$; hence at most two translates of the
window of $q$ meet a given point, and the sum over cells below is the periodised sum. Since $\ell$
divides $M$ and $M$
divides the side of the torus, the centres of the cells form a coset $c(q_0)+\ell\mathbb{Z}^d$,
where
$q_0$ is one fixed cell. Consequently, by \eqref{4.4:eq-local-coercivity-1} applied in each
coordinate,
\begin{equation}\label{4.4:eq-local-coercivity-4}
\sum_{q}h_q(b)^2=\prod_{\mu=1}^{d}\ \sum_{n\in\mathbb{Z}}
h\big((b_--c(q_0))_\mu/\ell-n\big)^2=1 .
\end{equation}
If $h_q(b)\ne0$, then $|(b_--c(q))_\mu|<\ell$ for every $\mu$, because
$\operatorname{supp}h\subset(-1,1)$; hence both end points of $b$ lie in the cube concentric with $q$
of side $2\ell+1\le3\ell$, that is, $b\subset q^*$. For each $t$ at most two values of $n$ give
$h(t-n)\ne0$; hence for a given bond $b$ at most $2^d$ cells $q$ have $h_q(b)\ne0$. Finally, for
numbers $a_\mu,a'_\mu\in[0,1]$ one has
$|\prod_\mu a_\mu-\prod_\mu a'_\mu|\le\sum_\mu|a_\mu-a'_\mu|$, so that, with $c_h=\sup|h'|$,
\begin{equation}\label{4.4:eq-local-coercivity-5}
\begin{aligned}
|h_q(b)-h_q(b')|
&\le\sum_{\mu}\big|h\big((b_--c(q))_\mu/\ell\big)-h\big((b'_--c(q))_\mu/\ell\big)\big|\\
&\le c_h\,|b_--b'_-|_1/\ell .
\end{aligned}
\end{equation}

\emph{The localisation identity.} Let $\mathsf{K}$ be any matrix on $\ell^2(Y';\mathfrak{g}^c)$ and
let
$(h_q)_q$ be real weights with $\sum_q h_q^2=1$ pointwise. Expanding the square and using
$\sum_q h_q(b)^2=\sum_q h_q(b')^2=1$ gives, entrywise,
\begin{equation*}
\sum_q h_q(b)h_q(b')=1-\tfrac12\sum_q\big(h_q(b)-h_q(b')\big)^2 .
\end{equation*}
Writing $\langle v,\mathsf{K}v\rangle=\sum_{b,b'}\overline{v(b)}\,\mathsf{K}(b,b')v(b')$, multiplying
the last identity
by $\overline{v(b)}\,\mathsf{K}(b,b')v(b')$ and summing over $b,b'$, we obtain for every complex
$v\in\ell^2(Y';\mathfrak{g}^c)$, with $h_qv$ the pointwise product,
\begin{equation}\label{4.4:eq-local-coercivity-6}
\langle v,\mathsf{K}v\rangle=\sum_q\langle h_qv,\mathsf{K}h_qv\rangle
+\frac12\sum_q\sum_{b,b'}\big(h_q(b)-h_q(b')\big)^2\,
\overline{v(b)}\,\mathsf{K}(b,b')v(b') .
\end{equation}

\emph{The near part, cell by cell.} Apply \eqref{4.4:eq-local-coercivity-6} with
$\mathsf{K}=\mathsf{T}_{\mathrm{near}}$ and the weights $h_q$, which satisfy
\eqref{4.4:eq-local-coercivity-4}. Fix a cell $q$. If a term $\mathsf{A}_\omega$ of
$\mathsf{T}_{\mathrm{near}}$ has $h_q(b)\mathsf{A}_\omega(b,b')h_q(b')\ne0$, then it has a row $b$
with $h_q(b)\ne0$, hence $b\subset q^*$ by the weights step, and therefore $[\omega]\subset q^\sharp$
by the localisation property of near walks established in the first step of the proof of the
walk-operator bounds. Conversely, suppose that $h_qv\neq0$, so that $q^*$ contains a bond $b$ of
$Y'$. Then $q^\sharp\subset W'$ by the hypothesis on $Y'$, and so every walk $\omega$ with
$d(\omega)\le r'$ and $[\omega]\subset q^\sharp$ occurs in the sum defining $\mathsf{T}_{W'}$.
Since restrictions compose,
$\mathsf{A}_\omega(\mathsf{u}\restriction[\omega])
=\mathsf{A}_\omega((\mathsf{u}\restriction q^\sharp)\restriction[\omega])$ for such $\omega$.
Let $X$ be the set of bonds of $Y'$ lying in $q^*$; then $h_qv\in\ell^2(X;\mathfrak{g}^c)$, and
\begin{equation*}
\langle h_qv,\mathsf{T}_{\mathrm{near}}h_qv\rangle
=\langle h_qv,T_X(\mathsf{u}\restriction q^\sharp)h_qv\rangle ,
\end{equation*}
where $T_X$ is the operator of Lemma~\ref{4.4:lem-cell-coercivity}, built, as there, from near
walks only. (If $h_qv=0$ both sides vanish.) The restriction of a local datum is a local datum, so
$\mathsf{u}\restriction q^\sharp\in\mathbf{P}_{\rho,q^\sharp}$ is a local datum on $q^\sharp$, and
Lemma~\ref{4.4:lem-cell-coercivity} applies:
\begin{equation*}
\operatorname{Re}\langle h_qv,\mathsf{T}_{\mathrm{near}}h_qv\rangle\ge\tfrac78\gamma_*|h_qv|^2 .
\end{equation*}
By \eqref{4.4:eq-local-coercivity-4}, $\sum_q|h_qv|^2=|v|^2$; summing over $q$,
\begin{equation}\label{4.4:eq-local-coercivity-7}
\operatorname{Re}\sum_q\langle h_qv,\mathsf{T}_{\mathrm{near}}h_qv\rangle\ge\tfrac78\gamma_*|v|^2 .
\end{equation}

\emph{The double sum.} On the support of $\mathsf{T}_{\mathrm{near}}(b,b')$ we have
$|b_--b'_-|_1\le r$ by the splitting step, hence
$(h_q(b)-h_q(b'))^2\le(c_hr/\ell)^2$ by \eqref{4.4:eq-local-coercivity-5}. A term of the sum over
$q$ is non-zero only if $h_q(b)\ne0$ or $h_q(b')\ne0$, which by the weights step happens for at most
$2\cdot2^d$ cells. Hence
\begin{equation*}
\sum_q\big(h_q(b)-h_q(b')\big)^2\le2^{d+1}(c_hr/\ell)^2
\end{equation*}
on the support of $\mathsf{T}_{\mathrm{near}}$, and the modulus of the double sum in
\eqref{4.4:eq-local-coercivity-6} (with $\mathsf{K}=\mathsf{T}_{\mathrm{near}}$) is bounded as
follows, using
the Schur bound for $|\!|\!|\cdot|\!|\!|_0$, the bound
$|\!|\!|\mathsf{T}_{\mathrm{near}}|\!|\!|_0\le\mathsf{B}$,
and, in the last but one step, the lower bound on $\ell$ assumed in the statement, in the form
\eqref{4.4:eq-local-coercivity-2}:
\begin{equation}\label{4.4:eq-local-coercivity-8}
\begin{aligned}
&\frac12\cdot2^{d+1}\Big(\frac{c_hr}{\ell}\Big)^2
\sum_{b,b'}|v(b)|\,|\mathsf{T}_{\mathrm{near}}(b,b')|\,|v(b')|
\le2^d\Big(\frac{c_hr}{\ell}\Big)^2|\!|\!|\mathsf{T}_{\mathrm{near}}|\!|\!|_0\,|v|^2\\
&\qquad\le2^dc_h^2r^2\mathsf{B}\,\ell^{-2}|v|^2
\le\tfrac1{16}\gamma_*|v|^2\le\tfrac18\gamma_*|v|^2 .
\end{aligned}
\end{equation}

\emph{Conclusion.} By \eqref{4.4:eq-local-coercivity-6}, \eqref{4.4:eq-local-coercivity-7} and
\eqref{4.4:eq-local-coercivity-8},
$\operatorname{Re}\langle v,\mathsf{T}_{\mathrm{near}}v\rangle\ge(\tfrac78-\tfrac18)\gamma_*|v|^2$;
together with \eqref{4.4:eq-local-coercivity-3},
\begin{equation*}
\operatorname{Re}\langle v,\mathsf{T}_{W'}v\rangle
\ge\big(\tfrac78-\tfrac18-\tfrac18\big)\gamma_*|v|^2=\tfrac58\gamma_*|v|^2 .
\end{equation*}
The hypothesis on $\ell$ is a lower bound only. The weights $h_q$, which use the initial points
$b_-$, occur only inside this proof of a lower bound and define no term of the expansion, so the
covariance property of the expansion is untouched; nothing is extended, the
topology of $W'$ plays no part, and the level $k$ enters nowhere.
\end{proof}

For every cell $q$ let $q^*$, $q^\sharp$ and $c^*_q$ be its enlargements, $\ell$ the cell side and $\kappa_q$ the partition operators. The remaining objects are those of the local form of the cell-resolvent expansion (\ref{4.4:prf-local-resolvent}):

\begin{itemize}
\item $\mathbb{B}(Z_0)$ is the bond set on which the covariance $C^{(k)}(Z_0)$ acts, and $Z'_0$ is the enlargement of
$Z_0$ in the construction of the localisation domains. For a set $Y$ of bonds, $\chi_{Y}$ is the projection onto
$\ell^2(Y)$.
\item The localised operator is written $\mathsf{A}_{\mathrm{loc}}$; its terms are $\mathsf{A}_\omega$, with support
$[\omega]$ and size $d(\omega)$, and $r'$ is the range. The operator localised in $W'$ is $\mathsf{T}_{W'}$.
\item $\mathsf{R}^{\mathbb{B}(Z_0)}_{\varpi}(x)$ is the contribution of a walk $\varpi$ with trace $|\varpi|$, and
$C^{(k)}(Z_0)^{W'}$ is the localised covariance.
\item $\mathbf{P}_{\rho,W'}$ is the set of local data on $W'$ with parameter $\rho$, and $\gamma_*$ is the coercivity
constant, both as fixed in the setting of Lemma~\ref{4.4:lem-local-coercivity}. The kernel norms are written
$|\!|\!|\cdot|\!|\!|_0$ and $|\!|\!|\cdot|\!|\!|_{\delta_{\mathsf{A}}}$. The constant $\mathsf{B}$ is the one in the
bound $|\!|\!|\mathsf{A}|\!|\!|_{\delta_{\mathsf{A}}} \le \mathsf{B}$ fixed in that setting. The region
$\mathcal{S}_0$ is the one whose family of terms $\{\mathsf{A}_\omega : [\omega]\subset\mathcal{S}_0\}$ carries this
bound. The constants $\varepsilon(r')$, $S$ and $\varrho$ are those of clause~(v) of the local form of the
cell-resolvent expansion.
\end{itemize}

For a set $Z$ and $r \ge 0$, $Z^{+r}$ denotes the set of points within sup-distance $r$ of $Z$. For an operator $T$
on $\ell^2(Y')$ we write $\operatorname{Re}T := \tfrac12(T+T^*)$. The inequality $\operatorname{Re}T \ge c$ means
that $\operatorname{Re}\langle Tu,u\rangle \ge c|u|^2$ for all $u$.

\begin{corollary}[The localised covariance inverts the localised form]
\label{4.4:cor-localised-covariance}
Work in the setting of Lemma~\ref{4.4:lem-local-coercivity} and of the local form of the cell-resolvent expansion
(\ref{4.4:prf-local-resolvent}). Let $Z_0 \in \mathbf{D}_k$ and $Y' := \mathbb{B}(Z_0)$. Let $W'$ be a closed union of
$M_1$-cubes with
\begin{equation*}
Z'_0 \subset W' \subset \Omega_{k+1}.
\end{equation*}
Then the following holds at every local datum $\mathsf{u} \in \mathbf{P}_{\rho,W'}$, complex data included. It also
holds for every value of the parameter $x$ admitted in the local form of the cell-resolvent expansion
(\ref{4.4:prf-local-resolvent}). On the finite-dimensional space $\ell^2(Y')$ we have
\begin{equation}\label{4.4:eq-localised-covariance-1}
\sum_{|\varpi| \subset W'} \mathsf{R}^{\mathbb{B}(Z_0)}_{\varpi}(x) = (x+\mathsf{T}_{W'})^{-1}.
\end{equation}
In particular
\begin{gather}
C^{(k)}(Z_0)^{W'} = \mathsf{T}_{W'}^{-1}, \qquad
\operatorname{Re}\bigl[\bigl(C^{(k)}(Z_0)^{W'}\bigr)^{-1}\bigr] = \operatorname{Re}\mathsf{T}_{W'}
\ge \tfrac58\gamma_*, \label{4.4:eq-localised-covariance-2}\\
\bigl\|C^{(k)}(Z_0)^{W'}\bigr\| \le \frac{8}{5\gamma_*}, \qquad
\operatorname{Re} C^{(k)}(Z_0)^{W'} \ge \frac{5\gamma_*}{8\mathsf{B}^2}. \label{4.4:eq-localised-covariance-3}
\end{gather}
No base point and no operator on the whole lattice enter.
\end{corollary}

\begin{proof}
\emph{The hypothesis of Lemma~\ref{4.4:lem-local-coercivity}.} That lemma requires, of $W'$, that
$q^\sharp \subset W'$ for every cell $q$ that touches a bond of $Y'$ or whose enlargement $q^*$ contains one. We
verify this.

By the construction of the localisation domains \cite{B-CMP116}, $Z'_0$ contains the set
$\{\operatorname{dist}(\cdot,Z_0)\le M\}$ of points at distance at most $M$ from $Z_0$. Moreover $Z'_0 \subset W'$.

Let $q$ be a cell that touches a bond of $Y'$, or whose enlargement $q^*$ contains one. Such a cell lies within
sup-distance $\ell+1$ of $Z_0$. The concentric cell $q^\sharp$ has side at most $4\ell$, and $4\ell \le M/12$.
Therefore
\begin{equation*}
q^\sharp \subset Z_0^{+4\ell} \subset \{\operatorname{dist}(\cdot,Z_0)\le M\} \subset Z'_0 \subset W'.
\end{equation*}
Thus the hypothesis of Lemma~\ref{4.4:lem-local-coercivity} holds.

A cell $q$ with $\kappa_q\chi_{Y'} \neq 0$ is among the cells considered in the preceding paragraph, so
$q^\sharp \subset W'$. Consequently the matrices $T_q$, $P_q(x)$, $B_q(x)$ of clause~(v) of the local form of the
cell-resolvent expansion (\ref{4.4:prf-local-resolvent}) are defined from the restriction $\mathsf{u}|_{W'}$. The
cells with $\kappa_q\chi_{Y'} = 0$ contribute $P_q = B_q = 0$ in clause~(i) of that statement and are absent from
clause~(v). So the two index sets of cells agree.

\emph{(i) The parametrix.} Set
\begin{equation*}
\mathsf{A}^{W'}_{\mathrm{loc}} := \sum\bigl\{\mathsf{A}_\omega(\mathsf{u}|_{[\omega]}) :
d(\omega) \le r',\ [\omega] \subset W'\bigr\},
\end{equation*}
and put $P := \sum_q P_q$ and $B := \sum_q B_q$.

By the support property of the cell construction, a term $\mathsf{A}_\omega$ of $\mathsf{A}_{\mathrm{loc}}$ with a
column in $c^*_q$ has $[\omega] \subset q^\sharp$. By the preceding paragraph $q^\sharp \subset W'$, so every such
term is a term of $\mathsf{A}^{W'}_{\mathrm{loc}}$. Hence
\begin{equation*}
\mathsf{A}^{W'}_{\mathrm{loc}}\chi_{c^*_q} = \mathsf{A}_{\mathrm{loc}}\chi_{c^*_q},
\qquad
T_q = \chi_{c^*_q}\mathsf{A}^{W'}_{\mathrm{loc}}\chi_{c^*_q}.
\end{equation*}
These two identities, together with $\sum_q \kappa_q = I_{Y'}$, are the inputs of the parametrix algebra of the
step on the parametrix, far part and square root (\ref{4.4:prf-parametrix-root}). That algebra therefore applies
verbatim and gives
\begin{equation}\label{4.4:eq-localised-covariance-4}
\bigl(x + \chi_{Y'}\mathsf{A}^{W'}_{\mathrm{loc}}\chi_{Y'}\bigr)P = I_{Y'} + B.
\end{equation}
By clause~(v) of the local form of the cell-resolvent expansion (\ref{4.4:prf-local-resolvent}),
$\|B\| \le \varrho \le \tfrac12$. Hence the Neumann series $\sum_{n\ge0}(-B)^n$ converges, and
$(I_{Y'}+B)\sum_{n\ge0}(-B)^n = I_{Y'}$. Put
\begin{equation*}
\mathsf{N}^{W'}(x) := P\sum_{n\ge0}(-B)^n.
\end{equation*}
By \eqref{4.4:eq-localised-covariance-4}, $\mathsf{N}^{W'}(x)$ is a right inverse of the square matrix
$x + \chi_{Y'}\mathsf{A}^{W'}_{\mathrm{loc}}\chi_{Y'}$. Since $Y'$ is finite, $\ell^2(Y')$ is finite-dimensional,
and a right inverse of a square matrix is its inverse.

\emph{(ii) The far part.} Let $\mathsf{T}_{\mathrm{far}}$ be the far part in clause~(v) of the local form of the
cell-resolvent expansion (\ref{4.4:prf-local-resolvent}). It is the part of $\mathsf{T}_{W'}$ not contained in
$\chi_{Y'}\mathsf{A}^{W'}_{\mathrm{loc}}\chi_{Y'}$, so that
\begin{equation*}
x + \mathsf{T}_{W'} = \bigl(x + \chi_{Y'}\mathsf{A}^{W'}_{\mathrm{loc}}\chi_{Y'}\bigr) + \mathsf{T}_{\mathrm{far}}.
\end{equation*}
By that clause~(v) and the standing smallness condition on $\varepsilon(r')$ of this section,
\begin{equation*}
\bigl\|\mathsf{T}_{\mathrm{far}}\mathsf{N}^{W'}\bigr\| \le \varepsilon(r')\cdot\frac{4S}{\gamma_*} \le \frac12 .
\end{equation*}
So the series $\sum_{m\ge0}(-\mathsf{T}_{\mathrm{far}}\mathsf{N}^{W'})^m$ converges. By (i),
\begin{align*}
(x+\mathsf{T}_{W'})\,\mathsf{N}^{W'}\sum_{m\ge0}\bigl(-\mathsf{T}_{\mathrm{far}}\mathsf{N}^{W'}\bigr)^m
&= \bigl(I_{Y'} + \mathsf{T}_{\mathrm{far}}\mathsf{N}^{W'}\bigr)
\sum_{m\ge0}\bigl(-\mathsf{T}_{\mathrm{far}}\mathsf{N}^{W'}\bigr)^m \\
&= I_{Y'}.
\end{align*}
Again by finite dimension, $x+\mathsf{T}_{W'}$ is invertible, and its inverse is
\begin{equation*}
\mathsf{N}^{W'}\sum_{m\ge0}\bigl(-\mathsf{T}_{\mathrm{far}}\mathsf{N}^{W'}\bigr)^m .
\end{equation*}
Expand $\mathsf{N}^{W'}$ and $\mathsf{T}_{\mathrm{far}}$ into their terms. This produces exactly the walks $\varpi$ of
clause~(i) of the local form of the cell-resolvent expansion (\ref{4.4:prf-local-resolvent}) all of whose steps
lie on the list of clause~(v) for $W'$. These are the walks with $|\varpi| \subset W'$. This proves
\eqref{4.4:eq-localised-covariance-1}.

\emph{The bounds.} In particular, \eqref{4.4:eq-localised-covariance-1} gives
$C^{(k)}(Z_0)^{W'} = \mathsf{T}_{W'}^{-1}$. Hence $\bigl(C^{(k)}(Z_0)^{W'}\bigr)^{-1} = \mathsf{T}_{W'}$.
Lemma~\ref{4.4:lem-local-coercivity}, whose hypothesis was shown to hold above, gives
$\operatorname{Re}\mathsf{T}_{W'} \ge \tfrac58\gamma_*$. This gives \eqref{4.4:eq-localised-covariance-2}.

Let $T$ be invertible on $\ell^2(Y')$ with $\operatorname{Re}T \ge c > 0$. For every $u$,
\begin{equation*}
c|u|^2 \le \operatorname{Re}\langle Tu,u\rangle \le |Tu|\,|u|,
\end{equation*}
so $|Tu| \ge c|u|$ and $\|T^{-1}\| \le 1/c$. Now let $w \in \ell^2(Y')$ and $u := T^{-1}w$. Since
$|w| = |Tu| \le \|T\|\,|u|$,
\begin{equation*}
\operatorname{Re}\langle w, T^{-1}w\rangle = \operatorname{Re}\langle Tu,u\rangle
\ge c|u|^2 \ge \frac{c\,|w|^2}{\|T\|^2}.
\end{equation*}
Apply these two facts with $T = \mathsf{T}_{W'}$ and $c = \tfrac58\gamma_*$. The first gives
$\|C^{(k)}(Z_0)^{W'}\| \le 8/(5\gamma_*)$.

For the second, we need $\|\mathsf{T}_{W'}\| \le \mathsf{B}$. This follows from
\begin{equation*}
\|\mathsf{T}_{W'}\| \le |\!|\!|\mathsf{T}_{W'}|\!|\!|_0 \le \mathsf{B}.
\end{equation*}
The first inequality is the comparison of the operator norm with the kernel norm $|\!|\!|\cdot|\!|\!|_0$. For the
second inequality, note that the family $\{\mathsf{A}_\omega : [\omega]\subset W'\}$ is a sub-family of the index set
attached to $\mathcal{S}_0$. Its terms therefore carry the datum-free majorants that yield
$|\!|\!|\mathsf{A}|\!|\!|_{\delta_{\mathsf{A}}} \le \mathsf{B}$. The same majorants bound the terms of
$\mathsf{T}_{W'}$, and therefore $|\!|\!|\mathsf{T}_{W'}|\!|\!|_0 \le \mathsf{B}$. The second fact then gives
$\operatorname{Re}C^{(k)}(Z_0)^{W'} \ge 5\gamma_*/(8\mathsf{B}^2)$, which is
\eqref{4.4:eq-localised-covariance-3}.

None of these steps uses a base point or an operator on the whole lattice. Each holds at every
$\mathsf{u} \in \mathbf{P}_{\rho,W'}$, complex data included.
\end{proof}

The bounds \eqref{4.4:eq-localised-covariance-2}--\eqref{4.4:eq-localised-covariance-3} are valid at every local
datum in $\mathbf{P}_{\rho,W'}$.

\begin{remark}[The vertex integral needs a symmetric base]\label{4.4:rem-symmetric-base-need}
The vertex value $f(1_{\sigma'})$ of \cite[(2.8)]{B-CMP116} is, in the reading of
\cite[p.~13]{B-CMP116} adopted in this chapter, the iterated integral
\begin{equation}\label{4.4:eq-symmetric-base-need-1}
f(1_{\sigma'})
=\int d\mu_0(X)\,e^{-\frac12\langle \Gamma X,\,C\,\Gamma X\rangle}
\int d\mu_C(B)\,e^{-\langle B,\,\Gamma X\rangle}\,F ,
\end{equation}
where $X$ and $B$ are the integration variables of \cite[(2.8)]{B-CMP116}, $d\mu_0$ is the measure
of $X$, $\Gamma$ is the operator through which $X$ enters the $B$-integral, $F$ is the remaining
integrand, and $d\mu_C$ is the normalised Gaussian measure whose density is proportional to
$e^{-\frac12\langle B,C^{-1}B\rangle}$. Here $\sigma$ denotes the parameters of the cluster
expansion of \cite[\S2]{B-CMP116} on which the operators of that section depend, and $\sigma'$,
$1_{\sigma'}$ (above) and $\sigma(Z)$ (below) are the parameter values so denoted there; the
$\sigma=0$ base is formed by those operators at $\sigma=0$. Two conditions make
\eqref{4.4:eq-symmetric-base-need-1} a number:
\begin{itemize}
\item[(G$^+$)] $\operatorname{Re}\langle B,C^{-1}B\rangle>0$ at every local datum; this is what
Corollary~\ref{4.4:cor-localised-covariance} supplies;
\item[(X-cv)] the $d\mu_0(X)$-integral converges.
\end{itemize}
Let $T:=C^{-1}$ be real, with symmetric part $T_s$ and antisymmetric part $T_a$, $T=T_s+T_a$, and
suppose $T_s>0$. Put $y:=\Gamma X$. Then the integral of $e^{-\langle B,y\rangle}$ against
$d\mu_C$ equals $e^{\frac12\langle y,T_s^{-1}y\rangle}$, while the prefactor equals
$e^{-\frac12\langle y,(T^{-1})_s y\rangle}$, where
\begin{equation}\label{4.4:eq-symmetric-base-need-2}
(T^{-1})_s=T^{-1}T_sT^{-\mathrm{T}}=\bigl(T_s+T_a^{\mathrm{T}}T_s^{-1}T_a\bigr)^{-1}.
\end{equation}
The net exponent is therefore $+\frac12\langle y,\mathsf{D}y\rangle$, with
\begin{equation}\label{4.4:eq-symmetric-base-need-3}
\mathsf{D}:=T_s^{-1}-\bigl(T_s+T_a^{\mathrm{T}}T_s^{-1}T_a\bigr)^{-1}\ \ge 0,
\qquad \mathsf{D}=0\iff T_a=0,
\qquad \|\mathsf{D}\|\le\|T_s^{-1}\|^3\,\|T_a\|^2 .
\end{equation}
Hence the cancellation that \cite[\S2]{B-CMP116} uses (``the zeroth order term cancels the first
quadratic form'') is precisely the symmetry of the $\sigma=0$ base covariance. The readings of the
vertex analysis of \cite[\S2]{B-CMP116} in this chapter take this symmetry as an input, in two
forms. The first is the statement that the base operator $\mathsf{A}_{\mathrm{base}}$ of the
$\sigma=0$ vertex analysis, taken at the datum $(\mathbf{U},\mathbf{J})=(U,0)$ and $\sigma=0$, is
real symmetric and bounded below by the constant $\gamma_0$ of that reading; here
$\mathsf{A}_{\mathrm{base}}$ is the operator written $Z_0C^*\Delta_kCZ_0$ in the reading of
\cite[\S2]{B-CMP116}, with $Z_0$, $C$, $C^*$, $\Delta_k$ the operators so denoted there. The
second consists of the statements that the operator $\Gamma$ of
\eqref{4.4:eq-symmetric-base-need-1}, taken at $\sigma=0$ and at the datum
$(\mathbf{U},\mathbf{J})=(U,0)$, namely $\Gamma_0:=\Gamma(0;U,0)$, is real, and that the quantity
$E_2$ of the reading of the vertex analysis is defined as the real part of the quantity $E_1$ of
that reading, $E_2:=\operatorname{Re}E_1$, over the real inverse
$(C_0^{\mathrm{cv}})^{-1}$ of the base covariance $C_0^{\mathrm{cv}}:=\mathsf{A}_{\mathrm{base}}^{-1}$.
Ba{\l}aban's base, ``the corresponding operators with $\sigma(Z)=0$, $\mathbf{U}=U$,
$\mathbf{J}=0$'' \cite[\S2]{B-CMP116}, is a partial sum: it collects the walks inside
$W(\emptyset)=Z'_0$ of \cite[(2.5)]{B-CMP116}, at a global datum as at a local one. The base
operator $\mathsf{A}_{\mathrm{base}}$ of the readings is the same partial sum: it is the sum of the
terms $\mathsf{A}_\omega$ of the total form below over the walks $\omega$ inside $W(\emptyset)$, at a
global datum as at a local one, so that the first input above is a statement about a partial sum of
the total form. The total form
$\mathsf{A}(U,0)=\sum_\omega\mathsf{A}_\omega$ of Lemma~\ref{4.4:lem-form-coercivity}, the sum over
all walks $\omega$, is real symmetric, by the positivity of the form at
$\sigma=0$ together with clause (a) of the symmetry theorem for the quadratic forms of this chapter,
but its partial sums are not closed under transposition term by term. The cellwise coercivity bound
and the localised coercivity bound, the latter in the form used in
Corollary~\ref{4.4:cor-localised-covariance}, bound a Hermitian part and say nothing of $T_a$ or of
$\operatorname{Im}T$. The lemma that follows supplies both, up to tails, from the total.
\end{remark}

\begin{proof}
Since $T$ is real and $T_a^{\mathrm{T}}=-T_a$, one has $\langle B,TB\rangle=\langle B,T_sB\rangle$
for real $B$, so $d\mu_C$ is the centred Gaussian measure of covariance $T_s^{-1}$, which exists
because $T_s>0$. With $\langle\cdot,\cdot\rangle$ the bilinear pairing, its Laplace transform is
\begin{equation*}
\int d\mu_C(B)\,e^{-\langle B,y\rangle}=e^{\frac12\langle y,T_s^{-1}y\rangle}
\end{equation*}
for every vector $y$, real or complex: for real $y$ this is the Gaussian formula, and both sides
are entire functions of $y$, so the identity extends by analytic continuation.
The matrix $T$ is invertible, since $\langle v,Tv\rangle=\langle v,T_sv\rangle>0$ for real $v\ne0$.
Since the pairing is bilinear, the antisymmetric part of $C=T^{-1}$ drops out of $\langle y,Cy\rangle$,
so $\langle y,Cy\rangle=\langle y,(T^{-1})_sy\rangle$, which gives the prefactor.

For \eqref{4.4:eq-symmetric-base-need-2}: $(T^{-1})_s=\frac12(T^{-1}+T^{-\mathrm{T}})$, and
$T^{-1}(T+T^{\mathrm{T}})T^{-\mathrm{T}}=T^{-\mathrm{T}}+T^{-1}$, so
$(T^{-1})_s=T^{-1}T_sT^{-\mathrm{T}}$. Inverting, and using $T^{\mathrm{T}}=T_s-T_a$,
\begin{align*}
\bigl((T^{-1})_s\bigr)^{-1}
=T^{\mathrm{T}}T_s^{-1}T
&=(T_s-T_a)\,T_s^{-1}\,(T_s+T_a)\\
&=T_s+T_a-T_a-T_aT_s^{-1}T_a
=T_s+T_a^{\mathrm{T}}T_s^{-1}T_a .
\end{align*}
The product of the two Gaussian factors is then $e^{\frac12\langle y,\mathsf{D}y\rangle}$ with
$\mathsf{D}$ as in \eqref{4.4:eq-symmetric-base-need-3}.

Put $\mathsf{P}:=T_a^{\mathrm{T}}T_s^{-1}T_a$ and $\mathsf{S}:=T_s+\mathsf{P}$. Since
$\langle v,\mathsf{P}v\rangle=\langle T_av,T_s^{-1}T_av\rangle\ge0$, we have $\mathsf{S}\ge T_s>0$.
Then $T_s^{-1/2}\mathsf{S}T_s^{-1/2}\ge1$, hence its inverse $T_s^{1/2}\mathsf{S}^{-1}T_s^{1/2}\le1$,
that is $\mathsf{S}^{-1}\le T_s^{-1}$; so $\mathsf{D}=T_s^{-1}-\mathsf{S}^{-1}\ge0$. Next,
$\mathsf{D}=0$ if and only if $\mathsf{S}=T_s$, if and only if $\mathsf{P}=0$; by the displayed
expression for $\langle v,\mathsf{P}v\rangle$ and $T_s^{-1}>0$, this holds if and only if $T_av=0$
for every $v$, that is $T_a=0$. Finally,
\begin{equation*}
\mathsf{D}=T_s^{-1}(\mathsf{S}-T_s)\mathsf{S}^{-1}=T_s^{-1}\mathsf{P}\mathsf{S}^{-1},
\end{equation*}
and $0<\mathsf{S}^{-1}\le T_s^{-1}$ gives $\|\mathsf{S}^{-1}\|\le\|T_s^{-1}\|$, while
$\|\mathsf{P}\|\le\|T_a\|^2\|T_s^{-1}\|$. Therefore
\begin{equation*}
\|\mathsf{D}\|\le\|T_s^{-1}\|\,\|\mathsf{P}\|\,\|\mathsf{S}^{-1}\|\le\|T_s^{-1}\|^3\,\|T_a\|^2,
\end{equation*}
which is the bound in \eqref{4.4:eq-symmetric-base-need-3}.

By \eqref{4.4:eq-symmetric-base-need-3}, the two quadratic forms in the exponent cancel exactly when
$T_a=0$, that is when the base covariance is symmetric; otherwise $\mathsf{D}\ne0$, so the quadratic
form $\frac12\langle y,\mathsf{D}y\rangle$ is not identically zero in $y$; for real $y=\Gamma X$ the
factor $e^{\frac12\langle y,\mathsf{D}y\rangle}$ is $\ge1$, and condition (G$^+$) of the remark,
namely $\operatorname{Re}\langle B,C^{-1}B\rangle>0$, which concerns only
$\operatorname{Re}\langle B,TB\rangle$, does not control it. Since the base is a
partial sum of the real symmetric total $\mathsf{A}(U,0)$ and partial sums are not
transposition-closed term by term, the symmetry of the base does not follow term by term from that
of the total; this is the gap that the lemma that follows closes up to tails.
\end{proof}

\begin{definition}[Symmetrisation of kernels and the three families]
\label{4.4:not-kernel-symmetrisation}
In this section $\mathfrak{g}\subset\mathfrak{g}^c$ is a real form. We write $(\cdot,\cdot)$ for the
complex-bilinear extension to $\mathfrak{g}^c$ of the invariant inner product of $\mathfrak{g}$,
$v\mapsto\bar v$ for complex conjugation of $\mathfrak{g}^c$ with respect to $\mathfrak{g}$, and
$\langle v,v'\rangle:=(\bar v,v')$ for the associated Hermitian inner product, with norm $|v|$.
Throughout this section $K$ denotes a kernel or one of the three families of (d), $H$ a cube of
sites and $g$ an element of $G$; none of these letters stands here for the number of
renormalisation steps, the Hamiltonian or the reference coupling.

\begin{enumerate}
\item[(a)] For $X\in\operatorname{End}(\mathfrak{g}^c)$, its \emph{conjugate} $\bar X$ is defined by
$\bar X v:=\overline{X\bar v}$. We set $\operatorname{Re}X:=\tfrac12(X+\bar X)$ and
$\operatorname{Im}X:=(X-\bar X)/2i$. The \emph{transpose} $X^{\mathsf T}$ is defined by
$(Xv,v')=(v,X^{\mathsf T}v')$, and $X^*$ is the adjoint for $\langle\cdot,\cdot\rangle$. Finally,
$|X|$ is the operator norm.

\item[(b)] Let $K$ be a kernel on $Y'\times Y'$ with values in $\operatorname{End}(\mathfrak{g}^c)$,
where $Y'$ is the compression set of (d) below.
Then $\bar K$, $\operatorname{Re}K$ and $\operatorname{Im}K$ are taken entrywise, and
$K^{\mathsf T}(b,b'):=K(b',b)^{\mathsf T}$. The kernel $K$ is \emph{real} if $\bar K=K$ and
\emph{symmetric} if $K^{\mathsf T}=K$. We put $\operatorname{sym}K:=\tfrac12(K+K^{\mathsf T})$,
$K_a:=\tfrac12(K-K^{\mathsf T})$, and
$\langle v,Kv\rangle:=\sum_{b,b'\in Y'}\langle v(b),K(b,b')v(b')\rangle$.

\item[(c)] The gauge action $R(u)$, as fixed earlier in this chapter, is diagonal in the bond $b$.
Each diagonal entry $R_b(u)$ is
$\operatorname{Ad}$ of a value of the gauge transformation $u$ at an end of $b$; this is
Ba{\l}aban's orthogonal action \cite[(2.17)--(2.18)]{B-CMP109}.

\item[(d)] \emph{The three families.} The letter $K$ ranges over $\mathsf{A}$, $\mathsf{R}(x)$ with
$x\ge0$ real, and $\mathsf{M}$, at a compression set $Y'\subset\mathbb{B}(T)$, where $\mathbb{B}(T)$ is
the bond set attached to $T$. Let $W'\subset\Omega_{k+1}$
be a closed union of $M_1$-cubes, and fix a local datum on $W'$. Here and below the constants
$\mathsf{B}$, $\delta_{\mathsf{A}}$, $\gamma_*$, the radius $\rho_w$ and the parameters $\ell$, $M$
are those of Definition~\ref{4.4:def-auxiliary-parameters}. We set
\begin{equation}\label{4.4:eq-kernel-symmetrisation-1}
K^{W'}:=
\begin{cases}
\chi_{Y'}\sum\{\mathsf{A}_\omega(\,\cdot\,|_{[\omega]}):[\omega]\subset W'\}\,\chi_{Y'}
=\mathsf{T}_{W'}, & K=\mathsf{A},\\[2pt]
\sum_{|\varpi|\subset W'}\mathsf{R}^{Y'}_{\varpi}(x), & K=\mathsf{R}(x),\\[2pt]
\sum_{|\varpi|\subset W'}\mathsf{M}^{Y'}_{\varpi}, & K=\mathsf{M}.
\end{cases}
\end{equation}
In the first line $\omega$ runs over steps with site set $[\omega]$, $\mathsf{A}_\omega$ is evaluated at
the restriction of the datum to $[\omega]$, and $\chi_{Y'}$ is multiplication by the indicator of $Y'$;
the compressed sum on the right, $\chi_{Y'}\sum\{\dots\}\chi_{Y'}$, is the chapter's kernel
$\mathsf{T}_{W'}$ on $W'$.
The last two lines are the two families, summed, supplied by clause (v) of the sixth-family lemma,
which constructs $\mathsf{R}^{Y'}_\varpi(x)$ and $\mathsf{M}^{Y'}_\varpi$; $\varpi$ runs over the
indices of those families and $|\varpi|$ is the set of sites of $\varpi$.

\item[(e)] \emph{Norm and rate.} We set
\begin{equation*}
(B_K,\delta_K):=
\begin{cases}
(\mathsf{B},\delta_{\mathsf{A}}), & K=\mathsf{A},\\
\bigl(4S(x+\tfrac12\gamma_*)^{-1},\,\theta\bigr), & K=\mathsf{R}(x),\\
\bigl(4S(2/\gamma_*)^{1/2},\,\theta\bigr), & K=\mathsf{M}.
\end{cases}
\end{equation*}
Here $B_K$ is the family's norm bound and $\delta_K$ its rate, $S$ being the norm constant fixed with
the families $\mathsf{R}$ and $\mathsf{M}$ earlier in this chapter (not the side of a lattice box).
The rate $\theta_K:=\theta_0/4=\theta$, with $\theta_0$ the base rate fixed earlier in this chapter
(a decay rate, not one of the functions $\theta_i$ of the partition of unity), is the
smaller one, $\theta\le\delta_{\mathsf{A}}/8$.

\item[(f)] \emph{Domain diameter.} $\varrho_K$ is the bound on the sup-diameter of a step's domain
(a constant attached to the family $K$, not the constant $\varrho$ with $\varrho_2=\varrho/2$). We
set $\varrho_{\mathsf{A}}:=\rho_w$ and $\varrho_K:=4\ell$ for $K\in\{\mathsf{R}(x),\mathsf{M}\}$.

\item[(g)] \emph{Cubes around a bond.} For a bond $b$ with initial site $b_-$ and $\mathsf{r}>0$, let
\begin{equation}\label{4.4:eq-kernel-symmetrisation-2}
H_b(\mathsf{r}):=\{y\in T:|y-b_-|_\infty\le\mathsf{r}+\varrho_K+3\}.
\end{equation}
This is a cube of side $2(\mathsf{r}+\varrho_K+3)$. For such an $H$, $H^{-2}$ is the set of sites of
$H$ at sup-distance at least $2$ from $T\setminus H$. We write ``$[\omega]\subset H^{-2}$'' to mean
that every site of $[\omega]$ lies in $H^{-2}$.
\end{enumerate}

With this notation the following hold.
\begin{enumerate}
\item[(A)] $\bar X$ is complex-linear, $X^*=\bar X^{\mathsf T}$, $|\bar X|=|X^{\mathsf T}|=|X|$, and
$|\operatorname{Re}X|,|\operatorname{Im}X|\le|X|$.

\item[(B)] For $g\in G$:
\begin{equation*}
\overline{\operatorname{Ad}(g)}=\operatorname{Ad}(g),\qquad
\operatorname{Ad}(g)^{\mathsf T}=\operatorname{Ad}(g)^{-1}=\operatorname{Ad}(g)^*.
\end{equation*}

\item[(C)] $K$ is real if and only if every entry of $K$ maps $\mathfrak{g}$ into $\mathfrak{g}$.

\item[(D)] For real $\mathfrak{g}$-valued $v$:
\begin{align*}
\langle v,\bar Kv\rangle&=\overline{\langle v,Kv\rangle},\\
\langle v,K^{\mathsf T}v\rangle&=\langle v,Kv\rangle,\\
\langle v,(\operatorname{sym}\operatorname{Re}K)v\rangle&=\operatorname{Re}\langle v,Kv\rangle.
\end{align*}

\item[(E)] Whatever the site convention in (c), $\bar R_b=R_b$ and $R_b^{\mathsf T}=R_b^{-1}$.
Hence, for every $X\in\operatorname{End}(\mathfrak{g}^c)$,
\begin{align*}
\operatorname{Im}(R_bXR_{b'}^{-1})&=R_b(\operatorname{Im}X)R_{b'}^{-1},\\
(R_{b'}XR_b^{-1})^{\mathsf T}&=R_bX^{\mathsf T}R_{b'}^{-1},\\
|R_bXR_{b'}^{-1}|&=|X|.
\end{align*}

\item[(F)] The value $\varrho_{\mathsf{A}}$ is a bound on the sup-diameter of the domains of the steps
$\omega$. The value $4\ell$ is such a bound for $\mathsf{R}(x)$ and $\mathsf{M}$.

\item[(G)] The bounds $\lvert\!\lvert\!\lvert\cdot\rvert\!\rvert\!\rvert_{\delta_K}\le B_K$, for the
chapter's kernel norm $\lvert\!\lvert\!\lvert\cdot\rvert\!\rvert\!\rvert_{\delta}$ with decay rate
$\delta$, hold in two
cases. First, for the local $W'$-family \eqref{4.4:eq-kernel-symmetrisation-1}, at every datum of
the local domain of data $\mathbf{P}_{\rho,W'}$. Second, for the global family in the global setting
$\mathcal{S}_0$, at every point of the domain of data $\mathbf{P}_\rho$. Here $\rho$ is the parameter
of these domains of data, not the effective density $\rho_k$ nor the contraction letter $\rho$ of
clause (v-d) of Theorem~\ref{3.1:thm-main}.
\end{enumerate}
\end{definition}

\begin{proof}
Write $v=v_1+iv_2$ and $v'=v_1'+iv_2'$ with $v_1,v_2,v_1',v_2'\in\mathfrak{g}$. Expanding,
\begin{equation*}
(v,v')=(v_1,v_1')-(v_2,v_2')+i\bigl((v_1,v_2')+(v_2,v_1')\bigr).
\end{equation*}
Since the inner product is real on $\mathfrak{g}$, this gives
$\overline{(v,v')}=(\bar v,\bar v')$. Applying the expansion to $(\bar v,v)$ with
$\bar v=v_1-iv_2$ gives $\langle v,v\rangle=|v_1|^2+|v_2|^2$, which
is unchanged under $v\mapsto\bar v$. Hence conjugation is an antilinear isometry.

(A) For $\lambda\in\mathbb{C}$ we have
\begin{equation*}
\bar X(\lambda v)=\overline{X(\bar\lambda\bar v)}=\lambda\,\overline{X\bar v}=\lambda\bar Xv,
\end{equation*}
so $\bar X$ is complex-linear. By definition $\bar X\bar v=\overline{Xv}$. Therefore
\begin{equation*}
\langle Xv,v'\rangle=(\overline{Xv},v')=(\bar X\bar v,v')=(\bar v,\bar X^{\mathsf T}v')
=\langle v,\bar X^{\mathsf T}v'\rangle,
\end{equation*}
that is, $X^*=\bar X^{\mathsf T}$. Because conjugation is an isometry,
$|\bar Xv|=|X\bar v|\le|X|\,|v|$, so $|\bar X|\le|X|$. Since $\bar{\bar X}=X$, the reverse inequality
also holds. Replacing $X$ by $\bar X$ in $X^*=\bar X^{\mathsf T}$ gives $X^{\mathsf T}=\bar X^{*}$,
hence
\begin{equation*}
|X^{\mathsf T}|=|\bar X^{*}|=|\bar X|=|X|.
\end{equation*}
Finally, $|\operatorname{Re}X|$ and $|\operatorname{Im}X|$ are at most $\tfrac12(|X|+|\bar X|)=|X|$.

(B) The map $\operatorname{Ad}(g)$ preserves $\mathfrak{g}$, so by the argument given for (C) below it
commutes with conjugation. This is
the identity $\overline{\operatorname{Ad}(g)}=\operatorname{Ad}(g)$. It also preserves the inner
product of $\mathfrak{g}$, and hence, by bilinearity, the form $(\cdot,\cdot)$. So
$\operatorname{Ad}(g)^{\mathsf T}\operatorname{Ad}(g)=1$, which gives
$\operatorname{Ad}(g)^{\mathsf T}=\operatorname{Ad}(g)^{-1}$. By (A) and the first identity,
$\operatorname{Ad}(g)^*=\overline{\operatorname{Ad}(g)}^{\mathsf T}=\operatorname{Ad}(g)^{\mathsf T}$.

(C) Suppose $X$ maps $\mathfrak{g}$ into $\mathfrak{g}$. Then
$\overline{X(v_1+iv_2)}=Xv_1-iXv_2=X\overline{(v_1+iv_2)}$, so $\bar X=X$. Conversely, if $\bar X=X$
and $v_1\in\mathfrak{g}$, then $\overline{Xv_1}=\bar Xv_1=Xv_1$, so $Xv_1\in\mathfrak{g}$. Apply this
to every entry of $K$.

(D) Let $v$ be real and fix $b,b'$. Then
\begin{equation*}
\langle v(b),\bar K(b,b')v(b')\rangle=(v(b),\overline{K(b,b')v(b')})
=\overline{(v(b),K(b,b')v(b'))}.
\end{equation*}
Summing over $b,b'$ gives the first identity. For the second, use the definition of the transpose,
the symmetry of $(\cdot,\cdot)$, and exchange the names of $b$ and $b'$:
\begin{align*}
\langle v,K^{\mathsf T}v\rangle
&=\sum_{b,b'}(v(b),K(b',b)^{\mathsf T}v(b'))
=\sum_{b,b'}(K(b',b)v(b),v(b'))\\
&=\sum_{b,b'}(v(b'),K(b',b)v(b))=\langle v,Kv\rangle.
\end{align*}
For the third, write $\operatorname{sym}\operatorname{Re}K$ as
$\tfrac14\bigl(K+\bar K+K^{\mathsf T}+(\bar K)^{\mathsf T}\bigr)$. Applying the first two identities,
with $\bar K$ in place of $K$ for the last term, gives
\begin{equation*}
\langle v,(\operatorname{sym}\operatorname{Re}K)v\rangle
=\tfrac14\bigl(2\langle v,Kv\rangle+2\,\overline{\langle v,Kv\rangle}\bigr)
=\operatorname{Re}\langle v,Kv\rangle.
\end{equation*}

(E) Each $R_b$ is $\operatorname{Ad}(g)$ for some $g\in G$. So (B) gives $\bar R_b=R_b$ and
$R_b^{\mathsf T}=R_b^{-1}=R_b^*$, and this uses nothing about which end of $b$ carries $g$.
Conjugation is multiplicative, $\overline{XY}=\bar X\bar Y$, since
$\overline{XY\bar v}=\bar X\,\overline{Y\bar v}$. Hence
$\overline{R_bXR_{b'}^{-1}}=R_b\bar XR_{b'}^{-1}$, which gives the identity for $\operatorname{Im}$.
Next, $(R_{b'}XR_b^{-1})^{\mathsf T}=(R_b^{-1})^{\mathsf T}X^{\mathsf T}R_{b'}^{\mathsf T}$, and this
equals $R_bX^{\mathsf T}R_{b'}^{-1}$. Finally, $R_b$ and $R_{b'}$ are unitary, so the operator norm
is unchanged.

(F) By the standing bound on the terms $\mathsf{A}_\omega$ stated earlier in this chapter, every domain
of a step $\omega$ has $\ell^1$-diameter at most $\rho_w$; and sup-diameter is
bounded by $\ell^1$-diameter. For $\mathsf{R}(x)$ and $\mathsf{M}$: the enlarged cell $q^\sharp$ of a
cell has side at most $4\ell$, and each domain of a far step $\omega$ has diameter $\rho_w<\ell/4$ by
the near--far condition of Definition~\ref{4.4:def-auxiliary-parameters}; in either case the
sup-diameter is at most $4\ell$, the value of $\varrho_K$ set in (f).

(G) For $K=\mathsf{A}$, the local bound is the local half of that standing bound on the terms
$\mathsf{A}_\omega$, and the global bound is that standing bound itself. For $K=\mathsf{R}(x)$ and
$K=\mathsf{M}$, the local bound is clause (v) of the sixth-family lemma, and the
global bound is its clauses (ii)--(iii).
\end{proof}

\begin{remark}\label{4.4:lem-reality-symmetry}
The statement of this lemma is not written out here; parts of its proof are printed below.
\end{remark}

\begin{proof}[Proof of Lemma~\ref{4.4:lem-reality-symmetry}, concluded]
\label{4.4:prf-tail-estimates}
We keep the notation of the preceding steps of the proof. Fix $b,b'\in Y'$ with
$H_b(\mathsf{r})\subset W'$. For a walk $\varpi$ of the walk expansion of the family, $K_\varpi(b,b')$
denotes its contribution to the entry $(b,b')$ and $d(\varpi;b,b')$ the quantity attached to $\varpi$
in the preceding step; $\varpi$ is
\emph{far} if $d(\varpi;b,b')>\mathsf{r}$. The kernels $F$ and $\hat F$ collect the far walks of
the local family on $W'$ and of the global family, respectively. The matrices $R_b$ are real and
orthogonal; $\hat K(b,b')$ is real and the kernel $\hat K$ is symmetric (see the preceding step).
The preceding step gives the decomposition
\begin{equation}\label{4.4:eq-tail-estimates-1}
R_b K^{W'}(b,b') R_{b'}^{-1} = \hat K(b,b') - \hat F(b,b') + R_b F(b,b') R_{b'}^{-1}.
\end{equation}
For a kernel $X$ on $Y'\times Y'$ we write $X^{\mathsf{T}}$ for the kernel
$(b,b')\mapsto X(b',b)^{\mathsf{T}}$. Recall that ${|\!|\!|}X{|\!|\!|}_0$ is the larger of the row-sum norm
$\sup_b\sum_{b'}|X(b,b')|$ and the column-sum norm $\sup_{b'}\sum_b|X(b,b')|$. The norm $|\cdot|$ on
entries is unchanged under transposition and under multiplication on either side by the orthogonal
matrices $R_b$.

\emph{Tails.} Let $\varpi$ be a far walk, so $d(\varpi;b,b')>\mathsf{r}$. By the decay clause of the
hypotheses on the family (the clause that also gives the entry bound $B_Ke^{-\delta_K|b-b'|}$ used
below), $d(\varpi;b,b')\ge|b-b'|$. Hence
\begin{equation*}
d(\varpi;b,b')\ge \mathsf{m}(b,b'):=\max\{\mathsf{r},|b-b'|\}.
\end{equation*}
It follows that
\begin{align*}
|F(b,b')| &\le \sum_{\varpi:\,d(\varpi;b,b')>\mathsf{r}} |K_\varpi(b,b')| \\
&\le e^{-\delta_K \mathsf{m}(b,b')} \sum_{\varpi} |K_\varpi(b,b')|\, e^{\delta_K d(\varpi;b,b')}
 \le B_K\, e^{-\delta_K \mathsf{m}(b,b')} .
\end{align*}
In the last step the sum over $\varpi$ is bounded by the row sum at $b$ of the weighted majorant
that defines the norm ${|\!|\!|}\cdot{|\!|\!|}_{\delta_K}$. By the norm bound in the definition of the
family this norm is at most $B_K$ for the local family on $W'$ at the datum $(U,0)\in\mathbf{P}_{\rho,W'}$.
For the family $\mathsf{A}$ (the one of the three families whose majorants do not depend on the
datum) the datum-free majorants are used instead.

Since $d(\varpi;b,b')>\mathsf{r}$ on every far walk, the same argument gives the row and column
sums:
\begin{align*}
\sum_{b'}|F(b,b')| &\le e^{-\delta_K\mathsf{r}} \sum_{b'}\sum_{\varpi}|K_\varpi(b,b')|\,
 e^{\delta_K d(\varpi;b,b')} \le e^{-\delta_K\mathsf{r}} B_K,\\
\sum_{b}|F(b,b')| &\le e^{-\delta_K\mathsf{r}} B_K .
\end{align*}
The same three bounds hold for $\hat F$. There one uses the norm of the global family at the
extended point $(\hat U,0)\in\mathbf{P}_\rho$, which the norm bounds in the definition of the global
family and the hypotheses of Lemma~\ref{4.4:lem-reality-symmetry} on the global family bound
by $B_K$. The extended point $\hat U$ depends on the row $b$ through the cube $H_b$. This is
harmless: each row, and each column, is estimated with its own extended point, and the bounds are
uniform on $\mathbf{P}_\rho$.

\emph{Proof of (a).} The matrices $R_b$, $R_{b'}$ are real. Hence
\begin{equation*}
\operatorname{Im}\bigl(R_bK^{W'}(b,b')R_{b'}^{-1}\bigr) = R_b\bigl(\operatorname{Im}K^{W'}(b,b')\bigr)R_{b'}^{-1}.
\end{equation*}
Since $\hat K(b,b')$ is real, \eqref{4.4:eq-tail-estimates-1} gives
\begin{equation*}
R_b\bigl(\operatorname{Im}K^{W'}(b,b')\bigr)R_{b'}^{-1}
= \operatorname{Im}\bigl(-\hat F(b,b') + R_bF(b,b')R_{b'}^{-1}\bigr).
\end{equation*}
The norm of an imaginary part does not exceed the norm of the matrix, and conjugation by the
orthogonal matrices $R_b$, $R_{b'}$ preserves the norm. By the tail bounds,
\begin{equation*}
|\operatorname{Im}K^{W'}(b,b')| \le |\hat F(b,b')| + |F(b,b')|
\le 2B_K e^{-\delta_K\mathsf{m}(b,b')}.
\end{equation*}
In the case $\mathsf{n}(b,b')=0$ of the preceding step, the term $\hat F(b,b')$ does not enter, and
the same argument gives
\begin{equation*}
|\operatorname{Im}K^{W'}(b,b')| \le |F(b,b')| \le B_K e^{-\delta_K\mathsf{m}(b,b')} .
\end{equation*}

\emph{Proof of (b).} For the entry $(b',b)$ we use the same cube $H=H_b(\mathsf{r})$. The tree
condition defining $d(\varpi;b,b')$ is symmetric in $b$ and $b'$, so $d(\varpi;b',b)=d(\varpi;b,b')$.
Hence the near walks for $(b',b)$ are those for $(b,b')$, all contained in $H^{-2}$. The
decomposition \eqref{4.4:eq-tail-estimates-1} at $(b',b)$ reads
\begin{equation*}
R_{b'}K^{W'}(b',b)R_b^{-1} = \hat K(b',b) - \hat F(b',b) + R_{b'}F(b',b)R_b^{-1}.
\end{equation*}
We transpose this identity. Since $R_b$ and $R_{b'}$ are orthogonal,
$(R_{b'}XR_b^{-1})^{\mathsf{T}} = R_bX^{\mathsf{T}}R_{b'}^{-1}$ for every matrix $X$. Since
$\hat K$ is symmetric, $\hat K(b',b)^{\mathsf{T}}=\hat K(b,b')$. Therefore
\begin{equation*}
R_bK^{W'}(b',b)^{\mathsf{T}}R_{b'}^{-1}
= \hat K(b,b') - \hat F(b',b)^{\mathsf{T}} + R_bF(b',b)^{\mathsf{T}}R_{b'}^{-1}.
\end{equation*}
Subtracting this from \eqref{4.4:eq-tail-estimates-1}, the terms $\hat K(b,b')$ cancel and
\begin{equation}\label{4.4:eq-tail-estimates-2}
\begin{split}
R_b\bigl[K^{W'}(b,b') - K^{W'}(b',b)^{\mathsf{T}}\bigr]R_{b'}^{-1}
&= -\hat F(b,b') + \hat F(b',b)^{\mathsf{T}} \\
&\quad + R_b\bigl[F(b,b') - F(b',b)^{\mathsf{T}}\bigr]R_{b'}^{-1}.
\end{split}
\end{equation}
The right side consists of four tails. We have $\mathsf{m}(b',b)=\mathsf{m}(b,b')$,
$|X^{\mathsf{T}}|=|X|$, and the norm is invariant under the orthogonal conjugation. So each tail
has norm at most $B_Ke^{-\delta_K\mathsf{m}(b,b')}$ by the tail bounds. Hence
\begin{equation*}
|K^{W'}(b,b') - K^{W'}(b',b)^{\mathsf{T}}| \le 4B_K e^{-\delta_K\mathsf{m}(b,b')},
\end{equation*}
which is (b). If $\mathsf{n}(b,b')=0=\mathsf{n}(b',b)$, only two tails occur.

\emph{The Schur bounds and the kernel $A_K$.} Assume now $H_b(\mathsf{r})\subset W'$ for every
$b\in Y'$. Summing the bound of (a) over $b'$ and using the row sums of the two tails at row $b$,
\begin{equation*}
\sum_{b'}|\operatorname{Im}K^{W'}(b,b')|
\le \sum_{b'}\bigl(|\hat F(b,b')|+|F(b,b')|\bigr)
\le 2B_K e^{-\delta_K\mathsf{r}}.
\end{equation*}
The column sums obey the same bound. To see this, run the argument given for (a) with the cube
$H_{b'}(\mathsf{r})$ in place of $H_b(\mathsf{r})$.

By \eqref{4.4:eq-tail-estimates-2}, the sum over $b'$ of $|K^{W'}(b,b')-K^{W'}(b',b)^{\mathsf{T}}|$
is bounded by two groups of terms: the row sums of $\hat F(b,\cdot)$ and $F(b,\cdot)$, and the
column sums of $\hat F(\cdot,b)$ and $F(\cdot,b)$. Hence
\begin{equation*}
\sum_{b'}|K^{W'}(b,b') - K^{W'}(b',b)^{\mathsf{T}}| \le 4B_K e^{-\delta_K\mathsf{r}},
\end{equation*}
and the column sums likewise.

Put
\begin{equation*}
A_K := \tfrac12\bigl[\operatorname{Re}K^{W'} + (\operatorname{Re}K^{W'})^{\mathsf{T}}\bigr],
\qquad
(\operatorname{Re}K^{W'})_a := \tfrac12\bigl[\operatorname{Re}K^{W'} - (\operatorname{Re}K^{W'})^{\mathsf{T}}\bigr].
\end{equation*}
By construction $A_K$ is real and symmetric. Writing
$K^{W'}=\operatorname{Re}K^{W'}+i\operatorname{Im}K^{W'}$ and subtracting $A_K$ gives
\begin{equation*}
K^{W'} - A_K = i\operatorname{Im}K^{W'} + (\operatorname{Re}K^{W'})_a .
\end{equation*}
Real parts commute with transposition, and the norm of a real part does not exceed the norm of the
matrix. Hence, by (b),
\begin{equation*}
|(\operatorname{Re}K^{W'})_a(b,b')|
= \tfrac12\bigl|\operatorname{Re}\bigl[K^{W'}(b,b')-K^{W'}(b',b)^{\mathsf{T}}\bigr]\bigr|
\le 2B_K e^{-\delta_K\mathsf{m}(b,b')}.
\end{equation*}
With (a) this yields $|K^{W'}-A_K|(b,b')\le 4B_Ke^{-\delta_K\mathsf{m}(b,b')}$. The row and column
sums of $\operatorname{Im}K^{W'}$ were bounded above by $2B_Ke^{-\delta_K\mathsf{r}}$. Those of
$(\operatorname{Re}K^{W'})_a$ are at most one half of the sums bounded by
$4B_Ke^{-\delta_K\mathsf{r}}$. Hence
\begin{equation*}
{|\!|\!|}K^{W'} - A_K{|\!|\!|}_0 \le 2B_K e^{-\delta_K\mathsf{r}} + 2B_K e^{-\delta_K\mathsf{r}} .
\end{equation*}
Finally, by the triangle inequality,
\begin{equation*}
|A_K(b,b')| \le \tfrac12\bigl(|K^{W'}(b,b')|+|K^{W'}(b',b)|\bigr)
\le \max\{|K^{W'}(b,b')|,|K^{W'}(b',b)|\}.
\end{equation*}
By the decay clause of the hypotheses on the family, and since $|b'-b|=|b-b'|$, this is at most
$B_Ke^{-\delta_K|b-b'|}$. Moreover
\begin{equation*}
{|\!|\!|}A_K{|\!|\!|}_0 \le {|\!|\!|}K^{W'}{|\!|\!|}_0 \le B_K .
\end{equation*}
Indeed, taking real parts does not increase any entry's norm. Transposition exchanges row sums and
column sums, so it preserves their maximum. Averaging two kernels does not increase that maximum,
by the triangle inequality. The last identity asserted in Lemma~\ref{4.4:lem-reality-symmetry}
holds by the definitions fixed in the notation preceding that lemma.
\end{proof}

\emph{Remarks.}
We record three observations on the hypotheses and the use of
Lemma~\ref{4.4:lem-reality-symmetry}.
\begin{itemize}
\item[(i)] The hypothesis $2(\mathsf{r}+\varrho_K+3)\le M/12$ and the inclusions
$H_b(\mathsf{r})\subset W'$ are one-sided. Where the lemma is applied they are met by the choice
$\mathsf{r}:=\mathsf{r}_s(M)$, growing with $M$. The required smallness
$e^{-\delta_K\mathsf{r}_s}\le e^{-\theta\mathsf{r}_s}$, with $\theta$ the constant of the floor
condition on $M$ imposed where $\mathsf{r}_s$ is chosen, then becomes the lower bound on $M$ stated
there. No upper bound on $M$ is imposed and the level $k$ enters nowhere. Nothing but the datum of
one cube is extended (see clause (d) of the remark following the cell lemma of this chapter), and
the topology of $W'$ plays no part.
\item[(ii)] The proof uses only the following hypotheses on the families:
\begin{itemize}
\item the clause furnishing the matrices $R_b$, with the matrices $R_b$ orthogonal;
\item two further clauses of the hypotheses on the families in
Lemma~\ref{4.4:lem-reality-symmetry}, the second of them only at radius zero;
\item the row sums and the column sums of the decay clause, at both ends;
\item the clause of the positivity hypothesis of Lemma~\ref{4.4:lem-reality-symmetry} stating
that the total kernel at $(\hat U,0)$ is real and symmetric.
\end{itemize}
All of these are already among the inputs of the theorem on independence of the localised terms
from the small-field region. In particular, no hypothesis is made that the underlying operators map
$\mathfrak{g}$-valued fields to $\mathfrak{g}$-valued fields.
\item[(iii)] Consider a global base $(U,0)$, where $U$ is the $G$-valued factor of a point of
the domain $\mathbf{P}'$ of the setting, so that the restriction of $U$ to $W'$ obeys the local
regularity condition of
Lemma~\ref{4.4:lem-reality-symmetry}. Wherever the partial sums of the setting are concerned, the
lemma holds at such a base with $W'=T$. It is also needed there for the partial sum of the setting
with index $\sigma=0$ inside $Z'_0$. Every base to which the lemma is applied is such a factor, and
global data are treated by restriction of the base to $W'$.
\end{itemize}

\begin{corollary}[The symmetric base form on a polymer]\label{4.4:cor-polymer-base-form}
Assume the hypotheses of Lemma~\ref{4.4:lem-reality-symmetry} and of
Lemma~\ref{4.4:lem-local-coercivity}, together with the standing condition $M\ge 48\ell$ on $M$
among the conditions on the auxiliary parameters.
For a set $X$ of lattice points and $r\ge 0$, write $X^{+r}$ for the set of points at sup-distance
at most $r$ from $X$, and for a bond $b$ let $b_-$ denote the endpoint of $b$ distinguished in
Lemma~\ref{4.4:lem-reality-symmetry}. Let
$Z_0\in\mathbf{D}_k$, and let $Z'_0$ be the $LM$-hull
enlargement of $Z_0$ attached to the class $\mathbf{D}_k$; it satisfies
\begin{equation*}
Z'_0\supseteq\{x:\operatorname{dist}_\infty(x,Z_0)\le M\}=Z_0^{+M}.
\end{equation*}
Let $W'\subset\Omega_{k+1}$ be a closed union of $M_1$-cubes with $Z'_0\subset W'$, and let $(U,0)$
be a real local datum on $W'$ in the sense of Lemma~\ref{4.4:lem-reality-symmetry}. Take
$Y':=\mathbb{B}(Z_0)$ for the family $\mathsf{T}$ and for the localised covariance $C$ of
Corollary~\ref{4.4:cor-localised-covariance}, and
$Y':=Y$ for the family $\mathsf{M}:=\mathsf{M}^Y$ and for the factor of the family $\mathsf{A}$
that enters the expression $\Gamma$ of this section,
where $\mathbb{B}(Z'_0)\subseteq Y\subseteq\mathbb{B}(T)$ (and $Y=\mathbb{B}(T)$ in the setting
$\mathcal{S}_0$).
Let $\mathsf{R}(x)$, indexed by a point $x$, be the further family of this section whose
$Z'_0$-family is treated in (iv) together with that of $\mathsf{M}$, and let $S$ be the constant
attached to the family $\mathsf{M}$ in this section, in terms of which its constant
$B_{\mathsf{M}}$ is given in (iv).
Let $\varepsilon_0^{P'}$ be the threshold fixed for the symmetric base form among the conditions
on the auxiliary parameters,
and let it satisfy $\varepsilon_0^{P'}\le\min\{3,\mathsf{B}\}$; put
\begin{equation}\label{4.4:eq-polymer-base-form-1}
\mathsf{r}^*:=\frac{2}{\theta_K}\log\frac{12\mathsf{B}}{\varepsilon_0^{P'}},
\end{equation}
and assume the floor stated among the same conditions, one-sided and imposed on $M$ after $\ell$ and
after $\theta_K$, $\mathsf{B}$ and $\varepsilon_0^{P'}$ have been fixed,
\begin{equation}\label{4.4:eq-polymer-base-form-2}
M\ge 24\,(4\ell+3+\mathsf{r}^*).
\end{equation}
Set
\begin{equation}\label{4.4:eq-polymer-base-form-3}
\mathsf{r}_s:=\frac{M}{24}-4\ell-3,\qquad
\eta_{\mathrm{sym}}:=4\mathsf{B}\,e^{-\theta_K\mathsf{r}_s/2},
\end{equation}
where $\theta_K=\theta\le\delta_{\mathsf{A}}/8$ is the decay rate of
Lemma~\ref{4.4:lem-reality-symmetry}, the same for each of its families $K$.
Then the following hold.
\begin{enumerate}
\item[(i)] $\mathsf{r}_s\ge\mathsf{r}^*>0$ and $\eta_{\mathrm{sym}}\le\varepsilon_0^{P'}/3$.
\item[(ii)] For each of the three families, clauses (a) and (b) of
Lemma~\ref{4.4:lem-reality-symmetry} hold with $\mathsf{r}=\mathsf{r}_s$ at every row index
$b\in\mathbb{B}(Z_0)$ of the kernel and every $b'\in Y'$.
\item[(iii)] Let $T_0:=\mathsf{T}_{Z'_0}(U|_{Z'_0},0)$ and let $A_0:=\operatorname{sym}\operatorname{Re}T_0$
be the symmetric part of its real part, that is, the operator with kernel
$A_0(b,b'):=\tfrac12\bigl(\operatorname{Re}T_0(b,b')+\operatorname{Re}T_0(b',b)\bigr)$. Then $A_0$ is
real symmetric and, with $T_0$ and $A_0$ regarded as kernels on $Y'=\mathbb{B}(Z_0)$, for all
$b,b'\in\mathbb{B}(Z_0)$,
\begin{align}
&\tfrac58\gamma_*\le A_0\le\mathsf{B},\qquad
|A_0(b,b')|\le\mathsf{B}\,e^{-\delta_{\mathsf{A}}|b-b'|},\label{4.4:eq-polymer-base-form-4}\\
&|T_0(b,b')-A_0(b,b')|\le 4\mathsf{B}\,e^{-\delta_{\mathsf{A}}\max\{\mathsf{r}_s,|b-b'|\}}
\le\eta_{\mathrm{sym}}\,e^{-(\delta_{\mathsf{A}}/2)|b-b'|},\label{4.4:eq-polymer-base-form-5}\\
&|\!|\!|T_0-A_0|\!|\!|_0\le 4\mathsf{B}\,e^{-\delta_{\mathsf{A}}\mathsf{r}_s}\le\eta_{\mathrm{sym}}.
\label{4.4:eq-polymer-base-form-6}
\end{align}
\item[(iv)] For the $Z'_0$-family of $\mathsf{M}$, and for that of $\mathsf{R}(x)$, clauses (a) and
(b) of Lemma~\ref{4.4:lem-reality-symmetry} hold with $\mathsf{r}=\mathsf{r}_s$ at every row $b_1$
with $\operatorname{dist}_\infty(b_{1,-},Z_0)\le 5M/6$, with the constants
\begin{equation*}
(B_{\mathsf{M}},\theta_K):=\bigl(4S(2/\gamma_*)^{1/2},\,\theta\bigr).
\end{equation*}
\end{enumerate}
\end{corollary}

\begin{proof}
\emph{The radius $\mathsf{r}_s$.} For each of the three families one has $\varrho_K\le 4\ell$, hence
\begin{equation*}
\mathsf{r}_s+\varrho_K+3=\frac{M}{24}-4\ell+\varrho_K\le\frac{M}{24},
\end{equation*}
that is, $2(\mathsf{r}_s+\varrho_K+3)\le M/12$, which is the condition on $\mathsf{r}$ in
Lemma~\ref{4.4:lem-reality-symmetry}; equivalently, $H_b(\mathsf{r}_s)$ has side at most $M/12$.
The positivity of $\mathsf{r}_s$ does not follow from $M\ge 48\ell$, which gives only
$M/24\ge 2\ell$; it follows from the floor \eqref{4.4:eq-polymer-base-form-2}, which is equivalent
to $M/24-4\ell-3\ge\mathsf{r}^*$, that is, to $\mathsf{r}_s\ge\mathsf{r}^*$. Moreover
$\mathsf{r}^*>0$: from $0<\varepsilon_0^{P'}\le\min\{3,\mathsf{B}\}$ we get $\mathsf{B}>0$ and
$\varepsilon_0^{P'}\le\mathsf{B}<12\mathsf{B}$, so the logarithm in
\eqref{4.4:eq-polymer-base-form-1} is positive. The restriction
$\varepsilon_0^{P'}\le\min\{3,\mathsf{B}\}$ costs nothing, since $\varepsilon_0^{P'}$ is a threshold
whose existence is asserted, and such a threshold may be decreased. The bound by $\mathsf{B}$ is
the one used here; the bound by $3$ is imposed at the same time for the sake of another floor in
the list of conditions on the auxiliary parameters and plays no role below. Finally,
$\mathsf{r}_s\ge\mathsf{r}^*$ gives
\begin{equation*}
e^{-\theta_K\mathsf{r}_s/2}\le e^{-\theta_K\mathsf{r}^*/2}=\frac{\varepsilon_0^{P'}}{12\mathsf{B}},
\qquad\text{hence}\qquad
\eta_{\mathrm{sym}}\le 4\mathsf{B}\cdot\frac{\varepsilon_0^{P'}}{12\mathsf{B}}
=\frac{\varepsilon_0^{P'}}{3}.
\end{equation*}
This proves (i).

\emph{The row condition on $\mathbb{B}(Z_0)$.} Each choice of $Y'$ is contained in
$\mathbb{B}(T)$, and each contains $\mathbb{B}(Z_0)$, since
$\mathbb{B}(Z_0)\subset\mathbb{B}(Z'_0)\subset Y\subset\mathbb{B}(T)$. Let $b\in\mathbb{B}(Z_0)$. Then
$b_-\in Z_0^{+1}$, so $H_b(\mathsf{r}_s)\subset Z_0^{+(M/24+1)}$. The standing condition on $M$
provides $M/8\ge 1$, so $M/24+1\le M/24+M/8\le M$, and therefore
\begin{equation*}
H_b(\mathsf{r}_s)\subset Z_0^{+(M/24+1)}\subset Z_0^{+M}\subset Z'_0\subset W'.
\end{equation*}
Together with the first paragraph, all hypotheses of Lemma~\ref{4.4:lem-reality-symmetry} are met
for each of the three families with $\mathsf{r}=\mathsf{r}_s$, at every row $b\in\mathbb{B}(Z_0)$;
this is (ii).

\emph{The operator $A_0$.} By its definition, $A_0$ has real entries and
$A_0(b,b')=A_0(b',b)$. For a real vector $v$,
\begin{equation*}
\langle v,A_0v\rangle=\operatorname{Re}\langle v,T_0v\rangle\ge\tfrac58\gamma_*|v|^2
\end{equation*}
by Lemma~\ref{4.4:lem-local-coercivity}. For a complex vector $v=v_1+iv_2$ with $v_1$, $v_2$
real, the cross terms $i\langle v_1,A_0v_2\rangle-i\langle v_2,A_0v_1\rangle$ cancel because $A_0$
is real symmetric, so
\begin{equation*}
\langle v,A_0v\rangle=\langle v_1,A_0v_1\rangle+\langle v_2,A_0v_2\rangle
\ge\tfrac58\gamma_*\bigl(|v_1|^2+|v_2|^2\bigr)=\tfrac58\gamma_*|v|^2.
\end{equation*}
This is the lower bound in \eqref{4.4:eq-polymer-base-form-4}. For the upper bound and the kernel
bound, the standing bounds on the localised family $\mathsf{T}$, with constants $\mathsf{B}$ and
$\delta_{\mathsf{A}}$, bound $|T_0(b,b')|$ by $\mathsf{B}e^{-\delta_{\mathsf{A}}|b-b'|}$; since by
definition $|A_0(b,b')|\le\tfrac12\bigl(|T_0(b,b')|+|T_0(b',b)|\bigr)$ and $|b-b'|=|b'-b|$, this
gives $|A_0(b,b')|\le\mathsf{B}e^{-\delta_{\mathsf{A}}|b-b'|}$. The same bounds give
$|\!|\!|A_0|\!|\!|_0\le\mathsf{B}$, and $\|A_0\|\le|\!|\!|A_0|\!|\!|_0$ by the comparison of the
operator norm with the norm $|\!|\!|\cdot|\!|\!|_0$.

\emph{The distance from $T_0$ to $A_0$.} By the definition of $A_0$,
\begin{equation*}
T_0-A_0=i\operatorname{Im}T_0+\tfrac12\bigl(\operatorname{Re}T_0-(\operatorname{Re}T_0)^{\top}\bigr).
\end{equation*}
Clauses (a) and (b) of
Lemma~\ref{4.4:lem-reality-symmetry}, applied by (ii) to the family $\mathsf{T}$, whose constants
there are $\mathsf{B}$ and $\delta_{\mathsf{A}}$, with $\mathsf{r}=\mathsf{r}_s$, bound the first
and the second term respectively, at every row $b\in\mathbb{B}(Z_0)$ and in the norm
$|\!|\!|\cdot|\!|\!|_0$ of kernels on $Y'=\mathbb{B}(Z_0)$, and together give the first
inequality of \eqref{4.4:eq-polymer-base-form-5} and the first
inequality of \eqref{4.4:eq-polymer-base-form-6}. For the second inequality of
\eqref{4.4:eq-polymer-base-form-5}, use
$\max\{\mathsf{r}_s,|b-b'|\}\ge\tfrac12\mathsf{r}_s+\tfrac12|b-b'|$ and
$\theta_K\le\delta_{\mathsf{A}}/8\le\delta_{\mathsf{A}}$:
\begin{equation*}
4\mathsf{B}\,e^{-\delta_{\mathsf{A}}\max\{\mathsf{r}_s,|b-b'|\}}
\le 4\mathsf{B}\,e^{-\delta_{\mathsf{A}}\mathsf{r}_s/2}\,e^{-(\delta_{\mathsf{A}}/2)|b-b'|}
\le\eta_{\mathrm{sym}}\,e^{-(\delta_{\mathsf{A}}/2)|b-b'|}.
\end{equation*}
For the second inequality of \eqref{4.4:eq-polymer-base-form-6}, $\delta_{\mathsf{A}}\mathsf{r}_s
\ge\theta_K\mathsf{r}_s/2$, so $4\mathsf{B}e^{-\delta_{\mathsf{A}}\mathsf{r}_s}\le\eta_{\mathrm{sym}}$.
This completes (iii).

\emph{Rows near $Z_0$ for $\mathsf{M}$ and $\mathsf{R}(x)$.} Let $b_1$ be a row with
$\operatorname{dist}_\infty(b_{1,-},Z_0)\le 5M/6$. As above, $H_{b_1}(\mathsf{r}_s)$ has side at most
$M/12$ and $H_{b_1}(\mathsf{r}_s)\subset Z_0^{+(5M/6+M/24+1)}$. Since
$\tfrac56+\tfrac1{24}=\tfrac78$ and $M/8\ge 1$, we have $\tfrac78M+1\le M$, whence
\begin{equation*}
H_{b_1}(\mathsf{r}_s)\subset Z_0^{+(5M/6+M/24+1)}\subset Z_0^{+M}\subset Z'_0.
\end{equation*}
Hence the row condition of Lemma~\ref{4.4:lem-reality-symmetry} holds at $b_1$ for the
$Z'_0$-families of $\mathsf{M}$ and of $\mathsf{R}(x)$, and clauses (a) and (b) hold there with
$(B_{\mathsf{M}},\theta_K)=(4S(2/\gamma_*)^{1/2},\theta)$. This is (iv).
\end{proof}

\begin{definition}[Floors and threshold of the base comparison]\label{4.4:def-base-floors}
Let $\mathsf{B}$, $\delta_{\mathsf{A}}$, $\gamma_*$ be the constants fixed before $M$ that are recalled in
Definition~\ref{4.4:def-auxiliary-parameters}, and let $S$ and $\theta_K>0$ be the further positive
quantities of this section that are fixed before $M$, $\theta_K$ being the decay rate of this section
that enters the exponentials below; the constant $S$ is distinct from the constant $\mathsf{S}$ that
appears further on. Let $\ell$ be the parameter chosen before $M$, and let $\beta$ be Ba{\l}aban's
schedule constant ($\beta=\tfrac14$; only $\beta\ge0$ is used below). Let $\gamma_0^X>0$ be the
constant of this section, and let
\begin{equation*}
\mathcal{R}_{\mathbf{P}}:=\min\Bigl\{\frac{\alpha_0'}{\alpha_0},\frac{\alpha_1'}{\alpha_1}\Bigr\}
\end{equation*}
be the ratio of the analyticity radii that Ba{\l}aban takes large in \cite{B-CMP116}. The letters
$\varsigma$ and $\rho$ are those of the proposition of this section in which the floors below are
used; in the floors they are taken at
\begin{equation*}
\varsigma=1+\beta,\qquad
\rho=\frac{\mathcal{R}_{\mathbf{P}}\,\gamma_0^X}{16\,\mathsf{B}}.
\end{equation*}

\emph{Threshold.} The threshold $\varepsilon_0^{P'}=\varepsilon_0^{P'}(\mathsf{B},\gamma_*,S,\theta_K,\delta_{\mathsf{A}})>0$
is a positive number, depending on the quantities displayed and on nothing else, below which two
arguments of this section close: the lemma of this chapter that is stated for every value of its
smallness parameter up to a threshold, that threshold being taken here to be $\varepsilon_0^{P'}$; and
the estimates \cite[(2.23)--(2.26)]{B-CMP116}, together with the lemma of this chapter that derives
\cite[(2.16)--(2.17)]{B-CMP116} under Ba{\l}aban's condition that $\alpha_5$ is sufficiently small.
The constants $C_P$, $C_s$, $\mathsf{S}$, $B_{\mathsf{M}}$ of this section are absorbed into
$\varepsilon_0^{P'}$. Since a threshold whose existence is asserted may be
decreased, we assume without loss of generality
\begin{equation}\label{4.4:eq-base-floors-1}
\varepsilon_0^{P'}\le\min\{3,\mathsf{B}\}.
\end{equation}
No value is assigned to $\varepsilon_0^{P'}$; it is a function of quantities fixed before $M$, and it is
not the constant $\varepsilon_0$ of \cite[Theorem~3]{B-CMP109}.

\emph{Floors.} Write $\bar\eta=e^{-\theta_K M/64}$, and put
\begin{equation*}
\mathsf{r}_s:=\frac{M}{24}-4\ell-3,\qquad
\eta_{\mathrm{sym}}:=4\mathsf{B}\,e^{-\theta_K\mathsf{r}_s/2},\qquad
\mathsf{r}^*:=\frac{2}{\theta_K}\log\frac{12\,\mathsf{B}}{\varepsilon_0^{P'}}.
\end{equation*}
The three floors of the base comparison are the conditions
\begin{align}
&\bar\eta\le\tfrac13\,\varepsilon_0^{P'}, \label{4.4:eq-base-floors-2}\\
&\mathcal{R}_{\mathbf{P}}\ge
\frac{16\,\mathsf{B}\,(1+\beta)\bigl(1+3/\varepsilon_0^{P'}\bigr)}{\gamma_0^X},
\label{4.4:eq-base-floors-3}\\
&M\ge 24\,\bigl(4\ell+3+\mathsf{r}^*\bigr). \label{4.4:eq-base-floors-4}
\end{align}
Condition \eqref{4.4:eq-base-floors-2} is a lower bound on $M$, imposed after $\theta_K$;
condition \eqref{4.4:eq-base-floors-3} is a lower bound on the ratio $\mathcal{R}_{\mathbf{P}}$;
condition \eqref{4.4:eq-base-floors-4} is a lower bound on $M$, imposed after $\ell$. Each floor is
one-sided, involves only $\mathcal{R}_{\mathbf{P}}$ itself or quantities fixed before $M$, imposes no
upper bound, and does not involve the level $k$. With \eqref{4.4:eq-base-floors-3} the lower endpoint
of the corresponding window in the order-of-choice table of this chapter becomes
\begin{equation*}
\frac{a^{\oplus}\,\gamma_0^X}{32\,\mathsf{B}\,(1+\beta)\bigl(1+3/\varepsilon_0^{P'}\bigr)},
\end{equation*}
where $a^{\oplus}$ is the constant of that table entering this endpoint.
The depth $\mathsf{r}_s$ and the quantity $\eta_{\mathrm{sym}}$ are those of the symmetric base form on a
polymer (Corollary~\ref{4.4:cor-polymer-base-form}).

\emph{Consequences.} The floors have the following consequences.
\begin{itemize}
\item[(a)] Condition \eqref{4.4:eq-base-floors-3} is equivalent to $\rho>\varsigma$ together with
\begin{equation*}
\frac{\varsigma}{\rho-\varsigma}\le\frac{\varepsilon_0^{P'}}{3}\qquad\text{at }\varsigma=1+\beta,
\end{equation*}
and it implies $\mathcal{R}_{\mathbf{P}}\ge 32\mathsf{B}/\gamma_0^X$ and $\rho\ge2\varsigma$.
\item[(b)] Condition \eqref{4.4:eq-base-floors-2} is equivalent to
\begin{equation*}
M\ge\frac{64}{\theta_K}\log\frac{3}{\varepsilon_0^{P'}},
\end{equation*}
whose right side is non-negative.
\item[(c)] Under \eqref{4.4:eq-base-floors-4}, $\mathsf{r}_s\ge\mathsf{r}^*>0$ and
$\eta_{\mathrm{sym}}\le\varepsilon_0^{P'}/3$.
\end{itemize}
\end{definition}

\begin{proof}
We prove the consequences (a)--(c) stated in Definition~\ref{4.4:def-base-floors}.

\emph{The floor \eqref{4.4:eq-base-floors-3}.} Since $\rho=\mathcal{R}_{\mathbf{P}}\gamma_0^X/(16\mathsf{B})$
with $\gamma_0^X,\mathsf{B}>0$, multiplying \eqref{4.4:eq-base-floors-3} by $\gamma_0^X/(16\mathsf{B})$
shows that it is equivalent to
\begin{equation}\label{4.4:eq-base-floors-5}
\rho\ge\varsigma\bigl(1+3/\varepsilon_0^{P'}\bigr),\qquad \varsigma=1+\beta .
\end{equation}
Under \eqref{4.4:eq-base-floors-5} we have $\rho-\varsigma\ge 3\varsigma/\varepsilon_0^{P'}>0$, hence
$\varsigma/(\rho-\varsigma)\le\varepsilon_0^{P'}/3$. Conversely, if $\rho>\varsigma$ and
$\varsigma/(\rho-\varsigma)\le\varepsilon_0^{P'}/3$, then $3\varsigma\le\varepsilon_0^{P'}(\rho-\varsigma)$,
which is \eqref{4.4:eq-base-floors-5}. Thus \eqref{4.4:eq-base-floors-3} is the condition
\begin{equation*}
\frac{\varsigma}{\rho-\varsigma}\le\frac{\varepsilon_0^{P'}}{3}\qquad\text{at }\varsigma=1+\beta .
\end{equation*}
By \eqref{4.4:eq-base-floors-1}, $\varepsilon_0^{P'}\le 3$, so $1+3/\varepsilon_0^{P'}\ge 2$ and
\eqref{4.4:eq-base-floors-5} gives $\rho\ge 2\varsigma$. Likewise \eqref{4.4:eq-base-floors-3} gives
\begin{equation*}
\mathcal{R}_{\mathbf{P}}\ge\frac{32\,\mathsf{B}\,(1+\beta)}{\gamma_0^X}\ge\frac{32\,\mathsf{B}}{\gamma_0^X},
\end{equation*}
the second inequality because $\beta\ge0$. So \eqref{4.4:eq-base-floors-3} contains the standing
conditions $\mathcal{R}_{\mathbf{P}}\ge 32\mathsf{B}/\gamma_0^X$ and $\rho\ge2\varsigma$.

\emph{The floor \eqref{4.4:eq-base-floors-2}.} Taking logarithms, $\bar\eta\le\varepsilon_0^{P'}/3$ is
equivalent to $M\ge(64/\theta_K)\log(3/\varepsilon_0^{P'})$, a lower bound on $M$ in terms of $\theta_K$ and
$\varepsilon_0^{P'}$, whose right side is non-negative by \eqref{4.4:eq-base-floors-1}.

\emph{The floor \eqref{4.4:eq-base-floors-4}.} Dividing \eqref{4.4:eq-base-floors-4} by $24$ and
subtracting $4\ell+3$ gives $\mathsf{r}_s\ge\mathsf{r}^*$. By \eqref{4.4:eq-base-floors-1},
$\varepsilon_0^{P'}\le\mathsf{B}<12\mathsf{B}$, so $12\mathsf{B}/\varepsilon_0^{P'}>1$, its logarithm is
positive, and, as $\theta_K>0$, $\mathsf{r}^*>0$. Hence $\mathsf{r}_s\ge\mathsf{r}^*>0$. Finally, since
$\theta_K>0$ and $\mathsf{r}_s\ge\mathsf{r}^*$,
\begin{equation*}
\eta_{\mathrm{sym}}=4\mathsf{B}\,e^{-\theta_K\mathsf{r}_s/2}
\le 4\mathsf{B}\,e^{-\theta_K\mathsf{r}^*/2}
=4\mathsf{B}\cdot\frac{\varepsilon_0^{P'}}{12\,\mathsf{B}}
=\frac{\varepsilon_0^{P'}}{3},
\end{equation*}
where the equality uses $\theta_K\mathsf{r}^*/2=\log(12\mathsf{B}/\varepsilon_0^{P'})$.
\end{proof}

\begin{proposition}[The local real base of a polymer. Proof outstanding]\label{4.4:prop-local-real-base}
Assume the following.
\begin{enumerate}
\item[(a)] The hypotheses of Lemma~\ref{4.4:lem-local-coercivity} hold, together with those of the
lemma on the symmetric part of the localised form, whose application is
Corollary~\ref{4.4:cor-polymer-base-form}. Among them are (W$'$) with its local half, (k1), (k2),
(k2$'$) and (k3).
\item[(b)] The two inequalities
\begin{equation*}
\theta_K L M \ge 648\,\kappa_1, \qquad e^{243\kappa_1} \le e^{\theta_K M/64}
\end{equation*}
hold. Put $\bar\eta := e^{-\theta_K M/64}$.
\item[(c)] The walk-counting bound holds. Consider cubes of side $LM$ and walks whose step domains have
diameter at most $LM/4$; here the diameter is at most $4\ell \le M/12$. A walk $\varpi$ from $b$
to $b'$ that meets $m$ of these cubes satisfies
\begin{equation*}
m \le 81\Bigl(\frac{2\,d(\varpi;b,b')}{LM} + 1\Bigr),
\end{equation*}
where $d(\varpi;b,b')$ is the length of $\varpi$ between $b$ and $b'$ in the sense of that bound.
\item[(d)] The three floors of Definition~\ref{4.4:def-base-floors} hold.
\end{enumerate}

We work in the setting of \cite[(2.14)]{B-CMP116}, which we now describe.
\begin{itemize}
\item $Z_0 \in \mathbf{D}_k$.
\item $Z \supset Z_0'$, where $Z_0'$ is the enlargement of $Z_0$ fixed in
Corollary~\ref{4.4:cor-polymer-base-form}, and $Z$ is a closed union of $M_1$-cubes with
$Z \subset \Omega_{k+1}$. In the setting $\mathcal{S}_0$ any $Z$ is admitted.
\item $\sigma_0$ is the family of cubes of side $LM$ of Definition~\ref{4.4:def-interpolation-rule},
and
\begin{equation*}
\sigma := (\sigma(\Delta))_{\Delta\in\sigma_0,\ \Delta\subset Z\setminus\operatorname{int}Z_0'}
\end{equation*}
is a family of complex numbers with $|\sigma(\Delta)| \le e^{\kappa_1}$.
\end{itemize}
Let $K$ be one of the three families
\begin{equation*}
\mathsf{A}, \qquad C := \mathsf{R}^{\mathbb{B}(Z_0)}(0), \qquad \mathsf{M} := \mathsf{M}^{Y}.
\end{equation*}
Here $Y$ is the region carrying the family $\mathsf{M}$, as in
Corollary~\ref{4.4:cor-polymer-base-form}. Each of the three families has the walk expansion
$K=\sum_\varpi K_\varpi$ of Lemma~\ref{4.4:lem-cell-resolvent}(v). For a walk $\varpi$ of it,
let $\sigma_0(\varpi)$ be the set of cubes of $\sigma_0$ that meet its
footprint $|\varpi|$, and put $m(\varpi) := \#\sigma_0(\varpi)$. For a local datum $\mathsf{u}$ on
$Z$, set, according to the interpolation rule of Definition~\ref{4.4:def-interpolation-rule},
\begin{equation*}
K(\sigma;\mathsf{u}) := \sum_{|\varpi|\subset Z} K_\varpi\bigl(\mathsf{u}|_{|\varpi|}\bigr)
\prod_{\Delta\in\sigma_0(\varpi)} \sigma(\Delta).
\end{equation*}
Since $W(\emptyset) = Z_0'$, we have $K(0;\mathsf{u}) = K^{Z_0'}(\mathsf{u})$. Put
$C(Z_0,\sigma;\mathsf{u}) := C(\sigma;\mathsf{u})$. By Corollary~\ref{4.4:cor-localised-covariance},
\begin{equation*}
C(Z_0,0;\mathsf{u})^{-1} = \mathsf{T}_{Z_0'}(\mathsf{u}).
\end{equation*}
Let $\chi_{Z_0}$ and $\chi_{Z_0^c}$ denote multiplication by the indicators of $Z_0$ and of its
complement, and set
\begin{equation*}
\Gamma(\sigma;\mathsf{u}) := \chi_{Z_0}\,\mathsf{A}(\sigma;\mathsf{u})\,\chi_{Z_0^c}\,
\mathsf{M}(\sigma;\mathsf{u}).
\end{equation*}
This is the product of the two separately weighted families: the interpolation rule is applied to
each operator on its own.

Let $(U|_Z,0)$ be the real local base of the polydiscs $\mathbf{P}_{\cdot,Z}$ of
Definition~\ref{4.4:def-regular-backgrounds}. Let
$\mathsf{u} = (e^{i\xi A'}U,\mathbf{J}) \in \mathbf{P}_{\varsigma,Z}$ with
$0 < \varsigma \le 1+\beta$ and $\rho \ge 2\varsigma$. Define the
\emph{local real base} by
\begin{equation}\label{4.4:eq-local-real-base-1}
\begin{gathered}
T_0 := \mathsf{T}_{Z_0'}(U|_{Z_0'},0), \qquad A_0 := \operatorname{sym}\operatorname{Re} T_0, \qquad
C_0^{\mathrm{cv}} := A_0^{-1},\\
\Gamma_0 := \Gamma(0;U,0), \qquad \Gamma_0^{\mathbb{R}} := \operatorname{Re}\Gamma_0 .
\end{gathered}
\end{equation}
Here $\operatorname{sym}$ denotes the symmetric part:
\begin{equation*}
\operatorname{sym}\operatorname{Re}T_0 = \tfrac12\bigl(\operatorname{Re}T_0 + (\operatorname{Re}T_0)^{\top}\bigr).
\end{equation*}
In \textup{(2s)}, \textup{(2$\Gamma$)} and \textup{(4)} below, $\mathsf{r}_s$ and
$\eta_{\mathrm{sym}} = 4\mathsf{B}e^{-\theta_K\mathsf{r}_s/2}$ are those of
Corollary~\ref{4.4:cor-polymer-base-form}. Then the following hold.

\smallskip
\noindent\textup{(1)} \emph{Base.} The operators $A_0$ and $C_0^{\mathrm{cv}}$ are real symmetric,
and
\begin{equation}\label{4.4:eq-local-real-base-2}
\tfrac58\gamma_* \le A_0 \le \mathsf{B}, \qquad
\mathsf{B}^{-1} \le C_0^{\mathrm{cv}} \le \frac{8}{5\gamma_*}, \qquad
|A_0(b,b')| \le \mathsf{B}e^{-\delta_{\mathsf{A}}|b-b'|}.
\end{equation}
The operator $\Gamma_0^{\mathbb{R}}$ is real, and
\begin{equation*}
|\Gamma_0^{\mathbb{R}}(b,b')| \le |\Gamma_0(b,b')|
\le \mathsf{B}B_{\mathsf{M}}\mathsf{S}\,e^{-(\theta/2)|b-b'|}.
\end{equation*}
Here
\begin{equation*}
\mathsf{S} := \sup_b \sum_{b'\in\mathbb{B}(T)} e^{-(\theta/4)|b-b'|}
\end{equation*}
is a number that depends only on $d$, $L$ and the gauge group, and is fixed before $M$.

\smallskip
\noindent\textup{(2)} \emph{Displacement of the true operators from the base, kernelwise with
decay.} The following hold at every $\mathsf{u}\in\mathbf{P}_{\varsigma,Z}$ and every $\sigma$ with
$|\sigma(\Delta)|\le e^{\kappa_1}$.

\smallskip
\noindent\textup{(2$\sigma$)} This clause holds at every $\mathsf{u}\in\mathbf{P}_{\rho,Z}$, the
full radius; it is in this form that \textup{(2U)} and \textup{(3)} use it. Let
$K\in\{\mathsf{A},C,\mathsf{M}\}$. Let $b$ be a row with
$\operatorname{dist}_\infty(b_-,Z_0)\le M/4$, where $b_-$ is the initial point of the bond $b$;
for $K = C$, let $b$ be any row and $b'$ any column.
Then
\begin{equation}\label{4.4:eq-local-real-base-3}
\sum_{\varpi:\,m(\varpi)\ge1} |K_\varpi(\mathsf{u})(b,b')|\,e^{\kappa_1 m(\varpi)}
\le B_K\,\eta_1\,e^{-(3\theta_K/8)|b-b'|},
\end{equation}
where
\begin{equation*}
\eta_1 := e^{81\kappa_1-\theta_K M/8} \le e^{-23\theta_K M/192} \le \bar\eta .
\end{equation*}
Hence, at such rows,
\begin{equation*}
\bigl|[K(\sigma;\mathsf{u})-K^{Z_0'}(\mathsf{u})](b,b')\bigr|
\le B_K\,\eta_1\,e^{-(3\theta_K/8)|b-b'|},
\end{equation*}
and
\begin{equation*}
\lvert\!\lvert\!\lvert C(\sigma;\mathsf{u})-\mathsf{T}_{Z_0'}(\mathsf{u})^{-1}
\rvert\!\rvert\!\rvert_0
\le \frac{8S}{\gamma_*}\,e^{81\kappa_1-\theta_K M/4}
\le \frac{8S}{\gamma_*}\,e^{-47\theta_K M/192}.
\end{equation*}

\smallskip
\noindent\textup{(2U)} For $K\in\{\mathsf{A},C,\mathsf{M}\}$,
\begin{align*}
\bigl|[K^{Z_0'}(\mathsf{u})-K^{Z_0'}(U,0)](b,b')\bigr|
&\le B_K\,e^{-\theta_K|b-b'|}\,\frac{\varsigma}{\rho-\varsigma},\\
\lvert\!\lvert\!\lvert K^{Z_0'}(\mathsf{u})-K^{Z_0'}(U,0)\rvert\!\rvert\!\rvert_0
&\le B_K\,\frac{\varsigma}{\rho-\varsigma}.
\end{align*}
The same two bounds hold for the $\sigma$-families, that is, for
$K(\sigma;\mathsf{u}) - K(\sigma;U,0)$, after a replacement of constants. At all rows,
$(B_K,\theta_K)$ is replaced by $(2B_Ke^{81\kappa_1},\,3\theta_K/8)$. At rows within sup-distance
$M/4$ of $Z_0$, it is replaced by $(2B_K,\,3\theta_K/8)$. These two classes of rows are those
of the bounds on the interpolated families, read with the count of \textup{(2$\sigma$)}.

\smallskip
\noindent\textup{(2s)} The symmetric part is close to the localised form:
\begin{equation*}
|T_0(b,b')-A_0(b,b')|
\le 4\mathsf{B}\,e^{-\delta_{\mathsf{A}}\max\{\mathsf{r}_s,|b-b'|\}}
\le \eta_{\mathrm{sym}}\,e^{-(\delta_{\mathsf{A}}/2)|b-b'|},
\end{equation*}
and $\lvert\!\lvert\!\lvert T_0-A_0\rvert\!\rvert\!\rvert_0\le\eta_{\mathrm{sym}}$. Moreover,
\begin{equation*}
\bigl|[T_0^{-1}-C_0^{\mathrm{cv}}](b,b')\bigr|
\le C_s\,\eta_{\mathrm{sym}}\,e^{-\theta_s|b-b'|},
\end{equation*}
with $C_s,\theta_s>0$ depending only on $(\mathsf{B},\gamma_*,S,\theta,\delta_{\mathsf{A}})$.

\smallskip
\noindent\textup{(2$\Gamma$)} For $b\in Z_0$,
\begin{align*}
|\operatorname{Im}\Gamma_0(b,b')|
&\le \mathsf{B}B_{\mathsf{M}}\mathsf{S}\,\bigl[4e^{-\theta\mathsf{r}_s/2}
+e^{-\theta M/4}\bigr]\,e^{-(\theta/4)|b-b'|}\\
&\le B_{\mathsf{M}}\mathsf{S}\,(\eta_{\mathrm{sym}}+\mathsf{B}\bar\eta)\,
e^{-(\theta/4)|b-b'|}.
\end{align*}

\smallskip
\noindent\textup{(3)} \emph{Coercivity on the whole $\sigma$-polydisc.} Let
$\mathsf{u}\in\mathbf{P}_{\rho,Z}$ (full radius) and let $|\sigma(\Delta)|\le e^{\kappa_1}$. Then
$C(Z_0,\sigma;\mathsf{u})$ is invertible and
\begin{equation}\label{4.4:eq-local-real-base-4}
\operatorname{Re} C(Z_0,\sigma;\mathsf{u})^{-1} \ge \tfrac38\gamma_* .
\end{equation}
Of the hypotheses, this clause uses only the two inequalities of (b) together with
(k2) and (k3), and none of the three floors of (d).

\smallskip
\noindent\textup{(4)} \emph{Applications.} Put
\begin{equation*}
\hat\varepsilon := \bar\eta + \frac{\varsigma}{\rho-\varsigma} + \eta_{\mathrm{sym}}.
\end{equation*}
Under the three floors of Definition~\ref{4.4:def-base-floors}, each of the three summands is at
most $\varepsilon_0^{P'}/3$. The quantity
$\varepsilon_0^{P'} = \varepsilon_0^{P'}(\mathsf{B},\gamma_*,S,\theta_K,\delta_{\mathsf{A}})$ is
the threshold, given by an existence statement, below which two results apply. The first is the
perturbation lemma for the Gaussian integral over the base, which holds for perturbation size
$\varepsilon\le\varepsilon_0^{P'}$. The second is the estimate of
\cite[(2.23)--(2.26)]{B-CMP116}, where the threshold enters as the smallness of $\alpha_5$. The
threshold $\varepsilon_0^{P'}$ is a function of quantities fixed before $M$. Without loss of
generality it is shrunk to
$\varepsilon_0^{P'}\le\min\{3,\mathsf{B}\}$ (Corollary~\ref{4.4:cor-polymer-base-form},
Definition~\ref{4.4:def-base-floors}). Then the following hold.

\smallskip
\noindent\textup{(4a)} The bounds on the $\sigma$-interpolated families and on their product
hold on the three local families $\mathsf{A}(\sigma;\mathsf{u})$, $C(\sigma;\mathsf{u})$ and
$\mathsf{M}(\sigma;\mathsf{u})$. Their proofs use only the count in \textup{(2$\sigma$)} and the
bound in \textup{(2U)}.

The perturbation lemma for the Gaussian integral, parts (i)--(iv), holds verbatim with the base
$(\Gamma_0^{\mathbb{R}},C_0^{\mathrm{cv}})$ in place of $(\Gamma(0;U,0),C(Z_0,0;U,0))$, with the
following changes.
\begin{itemize}
\item The perturbation size $\varepsilon$ is replaced by $C_P\hat\varepsilon$, where
$C_P = C_P(\mathsf{B},\gamma_*,S,\theta,\delta_{\mathsf{A}})$.
\item The lower bound for the base covariance is replaced by $\tfrac58\gamma_*/\mathsf{B}^2$.
Indeed, by \eqref{4.4:eq-local-real-base-2},
\begin{equation*}
C_0^{\mathrm{cv}} \ge \mathsf{B}^{-1} \ge \tfrac58\gamma_*/\mathsf{B}^2 .
\end{equation*}
\item The lower bound for the reciprocal of the norm of the base covariance is replaced by
$\tfrac58\gamma_*$. This satisfies $\tfrac58\gamma_*\le 1/\|C_0^{\mathrm{cv}}\|$.
\item The decay bound for the inverse base covariance is replaced by
\begin{equation*}
|(C_0^{\mathrm{cv}})^{-1}(b,b')| = |A_0(b,b')| \le \mathsf{B}e^{-\delta_{\mathsf{A}}|b-b'|},
\end{equation*}
which holds by \eqref{4.4:eq-local-real-base-2}. The Combes--Thomas step of part (i)
is therefore not needed.
\item The perturbation is $E' := C(Z_0,\sigma;\mathsf{u})-C_0^{\mathrm{cv}}$, bounded by
\textup{(2$\sigma$)}, \textup{(2U)} and \textup{(2s)}.
\end{itemize}

\smallskip
\noindent\textup{(4b)} Two estimates apply: the lemma that establishes
\cite[(2.16)--(2.17)]{B-CMP116}, and
the estimate of \cite[(2.23)--(2.26)]{B-CMP116}. Their input is that
$A_0 := \operatorname{sym}\operatorname{Re}\mathsf{T}_{Z_0'}(U,0)$ is real symmetric with
$A_0\ge\tfrac58\gamma_*$. This input is not assumed: at local data it is proved in \textup{(1)}.
In \cite[(2.24)]{B-CMP116} the base covariance $C_0^{\mathrm{cv}} = A_0^{-1}$ is symmetric.
Hence the zeroth-order term cancels exactly, and the integral over $d\mu_0(X)$ in
\cite[(2.25)]{B-CMP116} converges.

\smallskip
\noindent\textup{(4c)} Consider the functional $f(\sigma)$ of \cite[(2.8)--(2.9)]{B-CMP116}, formed
by the interpolation rule of Definition~\ref{4.4:def-interpolation-rule} with every $K$ replaced by
$K(\sigma;\mathsf{u})$. It is defined and analytic on the closed polydisc
\begin{equation*}
|\sigma(\Delta)| \le e^{\kappa_1}, \qquad \Delta\in\sigma_0,
\end{equation*}
which contains $[0,1]^{\sigma_0}$. This holds at every $\mathsf{u}\in\mathbf{P}_{\varsigma,Z}$,
$\varsigma\le1+\beta$, that represents a point of the first of the two classes
\begin{align*}
&U^{\Sigma^1}_{k+1}(Z;(1+\beta)\alpha_0,(1+\beta)\alpha_1,\alpha_0),\\
&U^{\Sigma^1}_{k+1}(Y';(1+\beta)\alpha_0,(1+\beta)\alpha_1,\alpha_0),
\end{align*}
with $Z\in\mathbf{D}_k\supset\mathbf{D}_{k+1}$. For such $\mathsf{u}$, the restriction
$\mathsf{u}|_{Y'}$ represents a point of the second class for every
$Y'\in\mathbf{D}_k$ with $Y'\subset Z$. This follows from the restriction theorem for the
one-power hereditary class at equal levels, which rests on
$\mathfrak{T}_{k+1}(Y')\subset\mathfrak{T}_{k+1}(Z)$. These restrictions are the arguments of the
potentials $V_k(Y',\cdot)$, $Y'\subset Z_0$, of $F(Z_0,B)$, in the space of
\cite[(1.34)]{B-CMP116} on $Y'$. All of this holds whether or not $\mathsf{u}$ extends to a point
of $\mathbf{P}_\varsigma$. In this form the hypothesis of the vertex formula
(Lemma~\ref{4.4:lem-vertex-formula}) is met, and in this form the later uses of $f$ depend on it.
\end{proposition}

Proof outstanding.

\smallskip
\noindent\emph{Route.} A proof has to establish the clauses from the hypotheses (a)--(d) as
follows. Clauses \textup{(1)} and \textup{(2s)} are to be obtained from the bound on the symmetric
part of the localised form applied in Corollary~\ref{4.4:cor-polymer-base-form}, which compares
$T_0$ with the total operator $\mathsf{A}(U,0)$, real and symmetric, up to contributions of range
at least $\mathsf{r}_s$, together with the coercivity of
Lemma~\ref{4.4:lem-local-coercivity}; the bound on $\Gamma_0$ in \textup{(1)} and clause
\textup{(2$\Gamma$)} are to be obtained from the kernel bounds of the two factors of $\Gamma_0$ and
from \textup{(2s)}. Clause \textup{(2$\sigma$)} is to be obtained at the full radius $\rho$ from the
walk-counting bound of (c), the weight $e^{\kappa_1 m(\varpi)}$ of a walk being absorbed by part of
the decay of $K_\varpi$ under the inequalities of (b). Clause \textup{(2U)} is to be obtained by a
Cauchy estimate in the radius of the polydisc from the bounds at the full radius $\rho$, read on
the count of \textup{(2$\sigma$)} for the $\sigma$-families. Clause \textup{(3)} is to be obtained
from \textup{(2$\sigma$)} at the full radius, the inequalities of (b), and (k2), (k3). Clause
\textup{(4)} is to be obtained from \textup{(1)}--\textup{(3)} and the three floors of (d).

\begin{proof}[Proof of Proposition~\ref{4.4:prop-local-real-base}, continued: the base covariance
and the imaginary part of the vertex kernel]
\label{4.4:prf-imaginary-part}
Throughout this part of the proof, kernels are indexed by bonds, and $|b-b'|$ is the $\ell^1$
distance on the torus between bonds. We write $\theta=\theta_K$ for the common decay rate of the
local families, $\delta_{\mathsf{A}}$ for the decay rate of $\mathsf{A}$, and $C_0^{\mathrm{cv}}=A_0^{-1}$
for the base covariance. For a rate $c>0$ we put
\begin{equation*}
\mathcal{L}(c):=\sup_b\sum_{b_1}e^{-c|b-b_1|}.
\end{equation*}

\emph{The base covariance.} The first of the two assertions concerning the base covariance is the
symmetric base form on a polymer, Corollary~\ref{4.4:cor-polymer-base-form}, applied as it stands.
For the second we start from the identity
\begin{equation*}
T_0^{-1}-C_0^{\mathrm{cv}}=T_0^{-1}-A_0^{-1}=T_0^{-1}(A_0-T_0)A_0^{-1},
\end{equation*}
which holds because $T_0^{-1}A_0A_0^{-1}-T_0^{-1}T_0A_0^{-1}=T_0^{-1}-A_0^{-1}$. We bound the kernels
of the three factors in turn.

The inverse $T_0^{-1}$ is the $Z'_0$-family of the covariance $C^{(k)}(Z_0)$ at the datum $(U,0)$.
By item~(v) of the cell-resolvent expansion, Lemma~\ref{4.4:lem-cell-resolvent}, at $x=0$,
together with the exponential kernel bounds of the local families,
\begin{equation*}
|T_0^{-1}(b,b')|=\bigl|C^{(k)}(Z_0)^{Z'_0}(U,0)(b,b')\bigr|\le\frac{8S}{\gamma_*}\,e^{-\theta|b-b'|}.
\end{equation*}
The kernel of the middle factor obeys
\begin{equation*}
|(A_0-T_0)(b,b')|\le\eta_{\mathrm{sym}}\,e^{-(\delta_{\mathsf{A}}/2)|b-b'|}
\le\eta_{\mathrm{sym}}\,e^{-\theta|b-b'|}.
\end{equation*}

For the third factor we use the Combes--Thomas device \cite{CT73}. By the first part of the proof,
$A_0$ is real symmetric with $A_0\ge\frac58\gamma_*$ and
$|A_0(b,b')|\le\mathsf{B}\,e^{-\delta_{\mathsf{A}}|b-b'|}$. Fix $b'$, let
$\delta_{\mathrm{CT}}>0$ be a rate to be chosen, and put $\Phi(b):=\delta_{\mathrm{CT}}|b-b'|$. Since
$\Phi$ is built from the $\ell^1$ torus distance alone, no single-valued coordinate on the torus is
needed, and by the triangle inequality
$|\Phi(b)-\Phi(b'')|\le\delta_{\mathrm{CT}}|b-b''|$. Regarding $e^{\pm\Phi}$ as multiplication
operators, define $F$ by $e^{\Phi}A_0e^{-\Phi}=A_0+F$, so that
$F(b,b'')=\bigl(e^{\Phi(b)-\Phi(b'')}-1\bigr)A_0(b,b'')$. From $|e^t-1|\le|t|e^{|t|}$ we obtain
\begin{equation*}
|F(b,b'')|\le\bigl(e^{\delta_{\mathrm{CT}}|b-b''|}-1\bigr)|A_0(b,b'')|
\le\delta_{\mathrm{CT}}|b-b''|\,e^{\delta_{\mathrm{CT}}|b-b''|}\,\mathsf{B}\,
e^{-\delta_{\mathsf{A}}|b-b''|}.
\end{equation*}
Let
\begin{equation*}
S'':=\sup_b\sum_{b''}|b-b''|\,e^{-(\delta_{\mathsf{A}}/2)|b-b''|}.
\end{equation*}
If $\delta_{\mathrm{CT}}\le\delta_{\mathsf{A}}/2$, the last bound is at most
$\delta_{\mathrm{CT}}\mathsf{B}|b-b''|e^{-(\delta_{\mathsf{A}}/2)|b-b''|}$, which is symmetric in
$(b,b'')$; so its row sums and its column sums are both bounded by $\mathsf{B}\delta_{\mathrm{CT}}S''$,
and the Schur test gives $\|F\|\le\mathsf{B}\delta_{\mathrm{CT}}S''$. We choose
\begin{equation*}
\delta_{\mathrm{CT}}:=\min\Bigl\{\frac{\delta_{\mathsf{A}}}{2},\,
\frac{5\gamma_*}{16\,\mathsf{B}S''}\Bigr\},
\qquad\text{so that}\qquad \|F\|\le\frac{5\gamma_*}{16}.
\end{equation*}
Then for every complex vector $v$ we have
\begin{equation*}
\operatorname{Re}\langle v,(A_0+F)v\rangle\ge\tfrac58\gamma_*\|v\|^2-\|F\|\,\|v\|^2
\ge\tfrac{5}{16}\gamma_*\|v\|^2 .
\end{equation*}
Hence $\|(A_0+F)v\|\ge\frac{5}{16}\gamma_*\|v\|$, so $A_0+F$ is invertible with
$\|(A_0+F)^{-1}\|\le16/(5\gamma_*)$. Since $A_0^{-1}=e^{-\Phi}(A_0+F)^{-1}e^{\Phi}$, we have
$A_0^{-1}(b,b')=e^{-\Phi(b)+\Phi(b')}(A_0+F)^{-1}(b,b')$. Here $\Phi(b')=0$, and every entry of a
matrix is bounded by its operator norm, so
\begin{equation*}
|A_0^{-1}(b,b')|\le\frac{16}{5\gamma_*}\,e^{-\delta_{\mathrm{CT}}|b-b'|}.
\end{equation*}

It remains to multiply the three decaying kernels, halving the rate at each Schur product. Put
$a:=\min\{\theta,\delta_{\mathrm{CT}}\}$. For nonnegative $x,y$ with $x+y\ge|b-b_2|$ we have
$ax+ay\ge\frac a2|b-b_2|+\frac a2x$. Using this with $x=|b-b_1|$, $y=|b_1-b_2|$, and the bound
$\theta\ge a$, we get
\begin{equation*}
\bigl|[T_0^{-1}(A_0-T_0)](b,b_2)\bigr|
\le\frac{8S}{\gamma_*}\,\eta_{\mathrm{sym}}\sum_{b_1}e^{-a|b-b_1|}e^{-a|b_1-b_2|}
\le\frac{8S}{\gamma_*}\,\eta_{\mathrm{sym}}\,\mathcal{L}(a/2)\,e^{-(a/2)|b-b_2|}.
\end{equation*}
The same inequality at rate $a/2$ applies to the product with $A_0^{-1}$, because
$\delta_{\mathrm{CT}}\ge a\ge a/2$. We conclude that
\begin{equation}\label{4.4:eq-imaginary-part-1}
\bigl|[T_0^{-1}-C_0^{\mathrm{cv}}](b,b')\bigr|\le C_s\,\eta_{\mathrm{sym}}\,e^{-\theta_s|b-b'|},
\qquad \theta_s:=\tfrac14\min\{\theta,\delta_{\mathrm{CT}}\},
\end{equation}
with
\begin{equation*}
C_s:=\frac{8S}{\gamma_*}\cdot\frac{16}{5\gamma_*}\,\mathcal{L}(a/2)\,\mathcal{L}(a/4).
\end{equation*}
Since $\delta_{\mathrm{CT}}$ is determined by $\mathsf{B}$, $\gamma_*$ and $\delta_{\mathsf{A}}$, the
constant $C_s$ depends only on $(\mathsf{B},\gamma_*,S,\theta,\delta_{\mathsf{A}})$.

\emph{The imaginary part of the vertex kernel.} Let $b\in Z_0$, and let $b_-$ be its lower endpoint,
so that $b_-\in Z_0^{+1}$, the set of points at sup-distance at most one from $Z_0$. At the datum
$(U,0)$ we have
\begin{equation*}
\Gamma_0(b,b')=\sum_{b_1}\mathsf{A}^{Z'_0}(b,b_1)\,\mathsf{M}^{Z'_0}(b_1,b').
\end{equation*}
For complex numbers, $\operatorname{Im}(XY)=(\operatorname{Im}X)(\operatorname{Re}Y)
+(\operatorname{Re}X)(\operatorname{Im}Y)$. We split the intermediate bonds $b_1$ into two
classes. A bond $b_1$ is \emph{near} if $\operatorname{dist}_\infty(b_{1-},Z_0)\le5M/6$, and
\emph{far} otherwise. For a far bond, since $\operatorname{dist}_\infty(b_-,Z_0)\le1$ and $M\ge3$,
\begin{equation*}
|b-b_1|\ge|b_--b_{1-}|_\infty\ge\frac{5M}{6}-1\ge\frac M2 .
\end{equation*}

By Corollary~\ref{4.4:cor-polymer-base-form}, since the row $b$ lies in the set
$\mathbb{B}(Z_0)$ of that corollary,
\begin{equation*}
|\operatorname{Im}\mathsf{A}^{Z'_0}(b,b_1)|\le2\mathsf{B}\,
e^{-\delta_{\mathsf{A}}\max\{\mathsf{r}_s,|b-b_1|\}},
\end{equation*}
and at near bonds $b_1$
\begin{equation*}
|\operatorname{Im}\mathsf{M}^{Z'_0}(b_1,b')|\le2B_{\mathsf{M}}\,
e^{-\theta\max\{\mathsf{r}_s,|b_1-b'|\}}.
\end{equation*}
At all bonds the exponential kernel bounds of the local families give
\begin{equation*}
|\mathsf{A}^{Z'_0}(b,b_1)|\le\mathsf{B}\,e^{-\delta_{\mathsf{A}}|b-b_1|},\qquad
|\mathsf{M}^{Z'_0}(b_1,b')|\le B_{\mathsf{M}}\,e^{-\theta|b_1-b'|}.
\end{equation*}
We also use $|\operatorname{Re}X|\le|X|$ and $|\operatorname{Im}Y|\le|Y|$. At near bonds we bound
both products by the estimates just listed. At far bonds we bound the first product in the same
way, and the second product by $|\mathsf{A}^{Z'_0}(b,b_1)|\,|\mathsf{M}^{Z'_0}(b_1,b')|$. With
$\delta_{\mathsf{A}}\ge\theta$ this gives
\begin{align*}
|\operatorname{Im}\Gamma_0(b,b')|
&\le\sum_{b_1}2\mathsf{B}\,e^{-\theta\max\{\mathsf{r}_s,|b-b_1|\}}\,B_{\mathsf{M}}\,
e^{-\theta|b_1-b'|}\\
&\quad+\sum_{b_1\ \mathrm{near}}\mathsf{B}\,e^{-\theta|b-b_1|}\,2B_{\mathsf{M}}\,
e^{-\theta\max\{\mathsf{r}_s,|b_1-b'|\}}\\
&\quad+\sum_{b_1\ \mathrm{far}}\mathsf{B}\,e^{-\theta|b-b_1|}\,B_{\mathsf{M}}\,
e^{-\theta|b_1-b'|}.
\end{align*}

We treat the three sums with two elementary inequalities. First, since $\max\{r,u\}\ge(r+u)/2$,
\begin{equation*}
e^{-\theta\max\{r,u\}}\le e^{-\theta r/2}e^{-\theta u/2}.
\end{equation*}
Second, for nonnegative $x,y$ with $x+y\ge|b-b'|$,
\begin{equation*}
e^{-(\theta/2)x}e^{-(\theta/2)y}\le e^{-(\theta/4)|b-b'|}e^{-(\theta/4)x}.
\end{equation*}
In the first two sums we apply the first inequality, weaken the remaining factor $e^{-\theta(\cdot)}$
to $e^{-(\theta/2)(\cdot)}$, extend the second sum to all $b_1$, and apply the second inequality with
$x=|b-b_1|$, $y=|b_1-b'|$. Each of the two sums is then at most
\begin{equation*}
2\mathsf{B}B_{\mathsf{M}}e^{-\theta\mathsf{r}_s/2}e^{-(\theta/4)|b-b'|}
\sum_{b_1}e^{-(\theta/4)|b-b_1|}
\le2\mathsf{B}B_{\mathsf{M}}\mathsf{S}\,e^{-\theta\mathsf{r}_s/2}e^{-(\theta/4)|b-b'|}.
\end{equation*}
In the third sum $|b-b_1|\ge M/2$, so that
$e^{-\theta|b-b_1|}\le e^{-\theta M/4}e^{-(\theta/2)|b-b_1|}$. After the same two steps, the third
sum is at most
\begin{equation*}
\mathsf{B}B_{\mathsf{M}}\mathsf{S}\,e^{-\theta M/4}e^{-(\theta/4)|b-b'|}.
\end{equation*}
In total,
\begin{equation*}
|\operatorname{Im}\Gamma_0(b,b')|\le\mathsf{B}B_{\mathsf{M}}\mathsf{S}
\bigl[4e^{-\theta\mathsf{r}_s/2}+e^{-\theta M/4}\bigr]e^{-(\theta/4)|b-b'|}.
\end{equation*}
Here $4\mathsf{B}e^{-\theta\mathsf{r}_s/2}=\eta_{\mathrm{sym}}$ (recall $\theta_K=\theta$), and
$e^{-\theta M/4}\le e^{-\theta M/64}=\bar\eta$. Therefore, for $b\in Z_0$,
\begin{equation}\label{4.4:eq-imaginary-part-2}
|\operatorname{Im}\Gamma_0(b,b')|\le B_{\mathsf{M}}\mathsf{S}\,
\bigl(\eta_{\mathrm{sym}}+\mathsf{B}\bar\eta\bigr)\,e^{-(\theta/4)|b-b'|}.
\end{equation}
\end{proof}

\begin{proof}[Proof of Proposition~\ref{4.4:prop-local-real-base}, parts \textup{(3)} and \textup{(4)}]
\label{4.4:prf-polydisc-positivity}
The argument uses the hypotheses of the proposition and, in addition, the following statements as hypotheses. The last of them enters only in part~(4c).
\begin{itemize}
\item The walk-count lemma. A walk $\varpi$ whose domains meet $m(\varpi)$ cubes of $\sigma_0$ satisfies
$m(\varpi)\le 81\bigl(2d(\varpi;b,b')/(LM)+1\bigr)$.
\item The kernel lemma for walk families, with its two classes of rows.
\item The product lemma for walk kernels. It multiplies at most three such kernels and estimates their difference in the datum $(U,\mathbf{J})$.
\item Parts (i)--(iv) of the perturbation lemma for the localised vertex covariance. These hold for $\varepsilon\le\varepsilon_0^{P'}$.
\item The positivity lemma for the Gaussian vertex integral, together with its bound on \cite[(2.23)--(2.26)]{B-CMP116}.
\item The restriction theorem for $\Sigma$-representatives at equal levels.
\end{itemize}
From the hypotheses of the proposition we use the two conditions $\theta_K LM\ge 648\kappa_1$ and
$e^{243\kappa_1}\le e^{\theta_K M/64}$, and we write $\bar\eta:=e^{-\theta_K M/64}$.
Parts (1), (2$\sigma$), (2U) and (2s) (base bounds and displacement of the operators) and
part~(2$\Gamma$) (imaginary part of the vertex kernel) are the preceding parts of the proof of
Proposition~\ref{4.4:prop-local-real-base}; below they are cited by these part numbers.

\smallskip\noindent\emph{Part} (3).
Let $\mathsf{u}\in\mathbf{P}_{\rho,Z}$ and $|\sigma(\Delta)|\le e^{\kappa_1}$, and put $\mathsf{T}_{\mathsf{u}}:=\mathsf{T}_{Z_0'}(\mathsf{u})$. By Corollary~\ref{4.4:cor-localised-covariance},
$C(Z_0,0;\mathsf{u})=\mathsf{T}_{\mathsf{u}}^{-1}$, with $\operatorname{Re}\mathsf{T}_{\mathsf{u}}\ge\frac58\gamma_*$ and $\|\mathsf{T}_{\mathsf{u}}\|\le\mathsf{B}$. Set $\mathsf{E}_\sigma:=C(Z_0,\sigma;\mathsf{u})-\mathsf{T}_{\mathsf{u}}^{-1}$. When $Z_0$ and $\mathsf{u}$ are fixed we abbreviate $C(\sigma)=C(\sigma;\mathsf{u})=C(Z_0,\sigma;\mathsf{u})$.
The operator norm is bounded by the Schur norm, so the Schur bound of part~(2$\sigma$) gives
\begin{equation*}
\|\mathsf{E}_\sigma\|\le\lvert\!\lvert\!\lvert\mathsf{E}_\sigma\rvert\!\rvert\!\rvert_0
\le(8S/\gamma_*)e^{81\kappa_1-\theta_K M/4}.
\end{equation*}

Put $\mathsf{t}:=\gamma_*/(4\mathsf{B}S)$; then $\mathsf{t}\le\frac14$, since $\gamma_*\le\mathsf{B}$ and $S\ge1$.
The hypothesis of the proposition that $(2\mathsf{B}/\gamma_*)Se^{-\theta_0\ell/8}\le\frac12$ says exactly
that $e^{-\theta_0\ell/8}\le\mathsf{t}$. The hypothesis $M\ge48\ell$ of the proposition, with
$\theta_K=\theta_0/4$, gives
\begin{equation*}
\theta_KM\ge(\theta_0/4)\cdot48\ell=12\theta_0\ell .
\end{equation*}
Hence
$e^{-\theta_KM}\le(e^{-\theta_0\ell/8})^{96}\le\mathsf{t}^{96}$.

With the last bound of part~(2$\sigma$), $e^{81\kappa_1-\theta_KM/4}\le e^{-47\theta_KM/192}$, and with
$\mathsf{B}(8S/\gamma_*)=2/\mathsf{t}$ we get
\begin{equation*}
\mathsf{B}\|\mathsf{E}_\sigma\|\le\frac{2}{\mathsf{t}}\,\mathsf{t}^{96\cdot47/192}
=2\mathsf{t}^{45/2}\le\tfrac12.
\end{equation*}
Since $\mathsf{B}=\gamma_*/(4S\mathsf{t})$, also
\begin{equation*}
\mathsf{B}^2\|\mathsf{E}_\sigma\|\le\frac{\gamma_*}{4S\mathsf{t}}\cdot2\mathsf{t}^{45/2}
=\frac{\gamma_*\mathsf{t}^{43/2}}{2S}\le\frac{\gamma_*}{16}.
\end{equation*}

Now $C(Z_0,\sigma;\mathsf{u})=\mathsf{T}_{\mathsf{u}}^{-1}+\mathsf{E}_\sigma=\mathsf{T}_{\mathsf{u}}^{-1}(I+\mathsf{T}_{\mathsf{u}}\mathsf{E}_\sigma)$, and $\|\mathsf{T}_{\mathsf{u}}\mathsf{E}_\sigma\|\le\mathsf{B}\|\mathsf{E}_\sigma\|\le\frac12$. So $C(Z_0,\sigma;\mathsf{u})$ is invertible, and
\begin{equation*}
C(Z_0,\sigma;\mathsf{u})^{-1}=(I+\mathsf{T}_{\mathsf{u}}\mathsf{E}_\sigma)^{-1}\mathsf{T}_{\mathsf{u}}
=\mathsf{T}_{\mathsf{u}}-\mathsf{T}_{\mathsf{u}}\mathsf{E}_\sigma(I+\mathsf{T}_{\mathsf{u}}\mathsf{E}_\sigma)^{-1}\mathsf{T}_{\mathsf{u}}.
\end{equation*}
Consequently
\begin{equation*}
\|C(Z_0,\sigma;\mathsf{u})^{-1}-\mathsf{T}_{\mathsf{u}}\|
\le\frac{\mathsf{B}^2\|\mathsf{E}_\sigma\|}{1-\mathsf{B}\|\mathsf{E}_\sigma\|}
\le2\mathsf{B}^2\|\mathsf{E}_\sigma\|\le\frac{\gamma_*}{8},
\end{equation*}
and therefore
\begin{equation*}
\operatorname{Re}C(Z_0,\sigma;\mathsf{u})^{-1}\ge\tfrac58\gamma_*-\tfrac18\gamma_*=\tfrac12\gamma_*\ge\tfrac38\gamma_* .
\end{equation*}

Of the parameter conditions of the proposition, only the two conditions recalled at the beginning
of the proof, the bound $(2\mathsf{B}/\gamma_*)Se^{-\theta_0\ell/8}\le\frac12$ and the bound
$M\ge48\ell$ were used for the value bound, and no Schur constant entered it.

\smallskip\noindent\emph{Part} (4a).
The kernel lemma for walk families is the count of part~(2$\sigma$) combined with the pointwise
decay bound of the families assumed in the proposition. The distance it requires is at least $M/3$ at rows within $M/4$ of $Z_0$. Its class of all rows is the walk count without the distance to $D$.

The product lemma for walk kernels multiplies at most three such kernels seen from $Z_0$, in its two regions. It takes the difference in $(U,\mathbf{J})$ by the Cauchy estimate, which is now part~(2U), inside $\mathbf{P}_{\rho,Z}$. Both lemmas therefore hold for the local families at every $\mathsf{u}\in\mathbf{P}_{\varsigma,Z}$, with
\begin{equation*}
(B',\delta')=\bigl(\max\{\mathsf{B},\,8S/\gamma_*,\,4S(2/\gamma_*)^{1/2}\},\ \theta\bigr),
\qquad\varepsilon:=\bar\eta+\frac{\varsigma}{\rho-\varsigma}.
\end{equation*}
Here $B'=\max_KB_K$ and $\delta'=\theta_K$. Since $8S/\gamma_*=4S(2/\gamma_*)$, this is the standing
substitution $(\mathsf{B},\delta_{\mathsf{A}})\mapsto(B',\delta')$ in the constants of the kernel and
product lemmas, under which the two lemmas are invoked where the proposition is applied.

\emph{Part (i) of the perturbation lemma.} Its inputs are a base $C_0^{\mathrm{cv}}$ with positive real part and a decaying kernel of $(C_0^{\mathrm{cv}})^{-1}$, together with a small, decaying $E':=C(Z_0,\sigma)-C_0^{\mathrm{cv}}$. Here
\begin{equation*}
C_0^{\mathrm{cv}}=A_0^{-1}\ge\mathsf{B}^{-1}\ge\tfrac58\gamma_*/\mathsf{B}^2,
\end{equation*}
the last inequality since $\gamma_*\le\mathsf{B}$. This is the lower bound required of the base. Next, $|(C_0^{\mathrm{cv}})^{-1}(b,b')|=|A_0(b,b')|\le\mathsf{B}e^{-\delta_{\mathsf{A}}|b-b'|}$ directly, by the
pointwise decay bound of the proposition, so the lemma's Combes--Thomas step is not needed. Also $1/\|C_0^{\mathrm{cv}}\|=\min\operatorname{spec}A_0\ge\frac58\gamma_*$, which is at least the required $1/(2B'S)$ because $B'\ge8S/\gamma_*$ and $S\ge1$. Finally,
\begin{equation*}
E'=[C(\sigma;\mathsf{u})-\mathsf{T}_{Z_0'}(\mathsf{u})^{-1}]+[\mathsf{T}_{Z_0'}(\mathsf{u})^{-1}-T_0^{-1}]
+[T_0^{-1}-C_0^{\mathrm{cv}}].
\end{equation*}
The first bracket is $\mathsf{E}_\sigma$ of part~(3); its kernel is controlled by the kernel bound of
part~(2$\sigma$), with $\eta_1:=e^{81\kappa_1-\theta_KM/8}$ as there. The second is bounded by
part~(2U) applied to $K=C$, whose $Z_0'$-family is $\mathsf{T}_{Z_0'}(\cdot)^{-1}$. The third is bounded
by the base comparison of part~(2s), with its constants $C_s$ and $\theta_s$ and with the quantity
$\eta_{\mathrm{sym}}$ of Definition~\ref{4.4:def-base-floors}. With $\eta_1\le\bar\eta$
(part~(2$\sigma$)), the kernel of $E'$ is at most
\begin{equation*}
\frac{8S}{\gamma_*}\eta_1e^{-(3\theta/8)|\cdot|}+\frac{8S}{\gamma_*}e^{-\theta|\cdot|}\frac{\varsigma}{\rho-\varsigma}
+C_s\eta_{\mathrm{sym}}e^{-\theta_s|\cdot|}\le C_P\,\hat\varepsilon\,e^{-\theta_s|\cdot|},
\end{equation*}
where $C_P$ is a constant of this proof depending only on $\mathsf{B}$, $\gamma_*$, $S$, $\theta$ and
$\delta_{\mathsf{A}}$, and $\hat\varepsilon:=\varepsilon+\eta_{\mathrm{sym}}$ with $\varepsilon$ as above.
Each of the three terms $\bar\eta$, $\varsigma/(\rho-\varsigma)$, $\eta_{\mathrm{sym}}$ of
$\hat\varepsilon$ is at most $\varepsilon_0^{P'}/3$ by the three floors of
Definition~\ref{4.4:def-base-floors}, so that $\hat\varepsilon\le\varepsilon_0^{P'}$.
The Neumann series of part (i) and its bound on the remainder $R_2$ there apply without change.
Its conclusion $\operatorname{Re}C(Z_0,\sigma)^{-1}>0$ is also given directly by part~(3).

\emph{Part (ii) of the perturbation lemma.} Write
\begin{equation*}
\Gamma(\sigma;\mathsf{u})-\Gamma_0^{\mathbb{R}}=[\Gamma(\sigma;\mathsf{u})-\Gamma_0]+i\operatorname{Im}\Gamma_0 .
\end{equation*}
By the product lemma (two factors, seen from $Z_0$), with its constant $C_2$, the first term has
kernel at most $C_2\varepsilon e^{-(\theta/32)|\cdot|}$. By the bound of part~(2$\Gamma$), the second
has kernel at most $B_{\mathsf{M}}\mathsf{S}(\eta_{\mathrm{sym}}+\mathsf{B}\bar\eta)e^{-(\theta/4)|\cdot|}$.
Together they are at most $C_P\hat\varepsilon e^{-(\theta/32)|\cdot|}$. Moreover
$R_3:=\operatorname{Re}\Gamma(\sigma;\mathsf{u})-\Gamma_0^{\mathbb{R}}=\operatorname{Re}[\Gamma(\sigma;\mathsf{u})-\Gamma_0]$ needs no reality of $\Gamma_0$, since $|\operatorname{Re}z-\operatorname{Re}z_0|\le|z-z_0|$. This is where the base $\operatorname{Re}\Gamma_0$ takes the place of a real $\Gamma_0$.

\emph{Part (iii) of the perturbation lemma.} It telescopes
$\Gamma^{\mathrm T}C\Gamma-(\Gamma_0^{\mathbb{R}})^{\mathrm T}C_0^{\mathrm{cv}}\Gamma_0^{\mathbb{R}}$. The transpose, not the adjoint, is what enters the analytic prefactor. The two differences are the ones just bounded, and the result is read as a form in real $X$. It applies without change.

\emph{Part (iv) of the perturbation lemma.} It writes
\begin{align*}
E_1&:=C(Z_0,\sigma;\mathsf{u})^{-1}-(C_0^{\mathrm{cv}})^{-1}=C(Z_0,\sigma;\mathsf{u})^{-1}-A_0,\\
E_2&:=\operatorname{Re}E_1,\qquad\det\operatorname{Re}C(\sigma)^{-1}=\det(A_0+E_2).
\end{align*}
It uses that $(C_0^{\mathrm{cv}})^{-1}=A_0$ is real, which part~(1) gives. Moreover $\|C_0^{\mathrm{cv}}E_i\|_0\le(8/(5\gamma_*))\,C_P\hat\varepsilon$ times a lattice sum, and this is at most $\frac12$ under the floors of Definition~\ref{4.4:def-base-floors}. So part (iv) applies without change.

\smallskip\noindent\emph{Part} (4b).
The positivity lemma for the Gaussian vertex integral asks, at the base $\sigma=0$, for a real symmetric $A_0$ bounded below by a positive constant. Part~(1) supplies this with the constant $\frac58\gamma_*$. Its parts (ii)--(iv) ask for a small $E$ with decaying kernel. Parts (2$\sigma$), (2U) and (2s) supply this, and part (iii) there already admits $E'\in\{\operatorname{Re}E,\operatorname{sym}(\operatorname{Re}E)\}$. Only Schur norms of size $O(\hat\varepsilon)$ enter \cite[(2.21)--(2.23)]{B-CMP116}, so,
in Ba{\l}aban's notation of \cite[p.~20]{B-CMP116}, where $\alpha_4$ and $\gamma_2$ are parameters
left free there,
$\alpha_5=O(1)\hat\varepsilon+O(1)\alpha_4+\gamma_2$.

In \cite[(2.23)--(2.24)]{B-CMP116} the base covariance is $C_0^{\mathrm{cv}}=A_0^{-1}$. Thus
$T:=(C_0^{\mathrm{cv}})^{-1}=A_0$ is symmetric and its antisymmetric part vanishes. In the
computation of Remark~\ref{4.4:rem-symmetric-base-need}, with the vector $y=\Gamma X$ there taken at
the base, $y=\Gamma_0^{\mathbb{R}}X$, this gives $\mathsf{D}=0$: the exponent
$\frac12\langle y,A_0^{-1}y\rangle$ of the $B$-integral cancels the exponent
$-\frac12\langle y,C_0^{\mathrm{cv}}y\rangle$ of the prefactor identically. This is the cancellation
of the zeroth-order term in \cite[(2.23)--(2.24)]{B-CMP116}, noted in \cite[p.~17]{B-CMP116}.
What remains of the $X$-exponent is Ba{\l}aban's remainder of order $O(\alpha_5)\|\cdot\|^2$. Hence
the $d\mu_0(X)$-integral of \cite[(2.25)]{B-CMP116} converges absolutely once $\alpha_5$ is small.
This is the smallness of $\alpha_5$ that Ba{\l}aban assumes, for example in the form
$\alpha_5\|C^{(k)}(Z_0,0)\|<\frac12$ \cite[p.~17]{B-CMP116}; here
$\|C_0^{\mathrm{cv}}\|\le8/(5\gamma_*)$. For the share of $\hat\varepsilon$ in $\alpha_5$, the
assumption is delivered by $\hat\varepsilon\le\varepsilon_0^{P'}$, that is, by the three floors of
Definition~\ref{4.4:def-base-floors}. For the shares of $\alpha_4$ and $\gamma_2$, it is delivered by
Ba{\l}aban's own ceilings on these free parameters in the expansion of \cite[p.~20]{B-CMP116}, in
which $\alpha_0$, $\alpha_1$ and $\delta_0$ are likewise quantities of that paper:
\begin{equation*}
\alpha_5=O(1)e^{-\frac13\delta_0M}+O(\alpha_0+\alpha_1)+O(1)\alpha_4+\gamma_2 .
\end{equation*}

The remainder is estimated as follows; in \cite[p.~17]{B-CMP116} it is bounded by
$\frac12O(\alpha_5)\|ZX\|^2$, where $\|ZX\|$ is the norm appearing in
\cite[(2.24)--(2.25)]{B-CMP116}; the letter $Z$ there is Ba{\l}aban's and is not the region $Z$ of
the datum $\mathsf{u}$. Let $\alpha_5\|C_0^{\mathrm{cv}}\|<\frac12$. Then
\begin{equation*}
(A_0-\alpha_5)^{-1}-A_0^{-1}=A_0^{-1}\bigl(A_0-(A_0-\alpha_5)\bigr)(A_0-\alpha_5)^{-1}
=\alpha_5A_0^{-1}(A_0-\alpha_5)^{-1},
\end{equation*}
with norm at most
\begin{equation*}
\alpha_5\|C_0^{\mathrm{cv}}\|\cdot\|C_0^{\mathrm{cv}}\|/(1-\alpha_5\|C_0^{\mathrm{cv}}\|)
\le2\alpha_5\|C_0^{\mathrm{cv}}\|^2 .
\end{equation*}
After the cancellation, the $X$-exponent consists of $\frac12\langle\Gamma_0^{\mathbb{R}}X,[(A_0-\alpha_5)^{-1}-A_0^{-1}]\Gamma_0^{\mathbb{R}}X\rangle$ and Ba{\l}aban's remainder from the first exponential. It is therefore at most
\begin{equation*}
\tfrac12\alpha_5\bigl(1+2\|C_0^{\mathrm{cv}}\|^2\|\Gamma_0^{\mathbb{R}}\|^2\bigr)\|ZX\|^2 .
\end{equation*}
The $d\mu_0(X)$-integral of \cite[(2.25)]{B-CMP116}, bounded in this way, converges absolutely if and only if the coefficient of $\frac12\|ZX\|^2$ is less than one, that is, once
$\alpha_5(1+2\|C_0^{\mathrm{cv}}\|^2\|\Gamma_0^{\mathbb{R}}\|^2)<1$. By part~(1), $\|C_0^{\mathrm{cv}}\|\le8/(5\gamma_*)$. Schur's bound applied to the kernel bound of part~(1) gives $\|\Gamma_0^{\mathbb{R}}\|\le\mathsf{B}B_{\mathsf{M}}\mathsf{S}^2$. So this is a smallness condition on $\alpha_5$ in terms of constants fixed before $M$, and the dependence of $\varepsilon_0^{P'}$ stated in Definition~\ref{4.4:def-base-floors} is unchanged.

\smallskip\noindent\emph{Part} (4c).
At a vertex $\sigma=1_{\sigma'}$, Definition~\ref{4.4:def-interpolation-rule} gives
$K(1_{\sigma'};\mathsf{u})=K^{W(\sigma')}(\mathsf{u})$ for every family. Here
$W(\sigma')\supseteq Z_0'$ and $|1_{\sigma'}(\Delta)|\le1\le e^{\kappa_1}$. The positivity
$\operatorname{Re}\langle B,C^{-1}B\rangle>0$ at the vertex is
Corollary~\ref{4.4:cor-localised-covariance} with $W'=W(\sigma')$. This choice is admissible because
$W(\sigma')\subset Z\subset\Omega_{k+1}$; by the hypothesis of the proposition on $\sigma_0$, the
$\sigma_0$-cubes lie off $\operatorname{int}Z_0'$. The convergence of the $X$-integral is
part~(4b) at $\sigma=1_{\sigma'}$.

Now consider the closed polydisc $|\sigma(\Delta)|\le e^{\kappa_1}$, which contains $[0,1]^{\sigma_0}$. On it the families converge absolutely and uniformly, with majorants $\sum_\varpi|K_\varpi|e^{\kappa_1m(\varpi)}<\infty$ by the count of part~(2$\sigma$). The positivity and the convergence hold at every $\sigma$ there, by parts~(3) and~(4b). The function $f(\sigma)$ is the double integral $\int d\mu_0(X)\int d\mu_{C(\sigma)}(B)[\dots]$ of \cite[(2.8)--(2.9)]{B-CMP116}. Its integrand is analytic in $\sigma$ term by term. It is dominated, uniformly on the closed polydisc, by the absolute majorants of \cite[(2.15)--(2.25)]{B-CMP116}. By dominated convergence and Morera's theorem, $f$ is analytic on the polydisc. Hence $f$ is $C^1$ on the cube, which is the hypothesis of Lemma~\ref{4.4:lem-vertex-formula}.

The datum enters only through $\mathsf{u}|_Z$. No extension of $\mathsf{u}$ to the whole torus is
required, and multiply connected $Z$ are included. The potentials in $F(Z_0,B)$ require
$\mathsf{u}|_{Y'}$, with $Y'\subset Z_0$, to be a $\Sigma$-representative at $(k+1,Y';\cdot)$, that
is, to satisfy \cite[(1.34)]{B-CMP116} on $Y'$. The restriction theorem for $\Sigma$-representatives
at equal levels gives this from a $\Sigma$-representative $\mathsf{u}$ at
$(k+1,Z;(1+\beta)\alpha_0,(1+\beta)\alpha_1,\alpha_0)$. For this reason part~(4c) is stated for such
representatives; they form the domain of the corollary in which the present proposition is applied.
\end{proof}

\begin{remark}[The local base in the cluster-expansion lemmas]\label{4.4:rem-cluster-lemmas-base}
In \cite[\S2]{B-CMP116} the cluster expansion is built on a base taken at the value $\sigma = 0$
of the polydisc variable $\sigma$ (the variable of the polydisc on which
Proposition~\ref{4.4:prop-local-real-base} is stated), and at the configuration $U$ with zero
current, $(\mathbf{U},\mathbf{J}) = (U,0)$. In the notation of \cite[\S2]{B-CMP116} this base is
the pair
\begin{equation*}
\bigl(\Gamma(0;U,0),\, C(Z_0,0;U,0)\bigr).
\end{equation*}
At local data this base is replaced by the pair
\begin{equation*}
\bigl(\Gamma_0^{\mathbb{R}},\, C_0^{\mathrm{cv}}\bigr),
\end{equation*}
where $\Gamma_0^{\mathbb{R}}$ is the first member of the local real base of
Proposition~\ref{4.4:prop-local-real-base}, and $C_0^{\mathrm{cv}}$ is the inverse of its second
member, a real symmetric operator. Both members of the pair are real, and $C_0^{\mathrm{cv}}$ is
symmetric.

With this replacement, each of the following is read word for word, with the small parameter
$\varepsilon$ of these statements replaced by $C_P\hat{\varepsilon}$, where $\hat{\varepsilon}$ is
the small parameter of this chapter at local data and $C_P$ is a constant fixed once in this
chapter:
\begin{itemize}
\item the two cluster-expansion lemmas and the cluster-expansion proposition built on
\cite[\S2]{B-CMP116};
\item the further proposition built on \cite{B-CMP116}.
\end{itemize}

Let $Z$ be a region at level $k+1$, and consider every local datum that is a representative of the
one-power hereditary class
\begin{equation*}
U^{\Sigma^1}_{k+1}\bigl(Z;(1+\beta)\alpha_0,(1+\beta)\alpha_1,\alpha_0\bigr).
\end{equation*}
At every such datum the statements concerning the following hold:
\begin{itemize}
\item the vertex functional;
\item the function $f$ of the auxiliary lemma;
\item the exponentiated series of the corollary on exponentiation;
\item the persistence of \cite[Lemma~3]{B-CMP116} and of the proposition of this chapter built on
\cite{B-CMP119} under the operation that this chapter applies to them.
\end{itemize}
By the heredity theorem for the one-power hereditary class, applied with equal levels, every such
datum is then also a representative of that class at $(k+1, Y')$, with the parameters supplied by
that theorem, for every $Y' \in \mathbf{D}_k$ with $Y' \subset Z$; and the statements concerning
the four items above hold at $(k+1, Y')$ as well.
\end{remark}

\begin{proof}
The pair $(\Gamma_0^{\mathbb{R}}, C_0^{\mathrm{cv}})$ is built from the local real base of
Proposition~\ref{4.4:prop-local-real-base}. That base is real, and its second member is a real
symmetric operator. Hence $\Gamma_0^{\mathbb{R}}$ and $C_0^{\mathrm{cv}}$ are real, and
$C_0^{\mathrm{cv}}$, the inverse of a real symmetric operator, is symmetric.

The pair $(\Gamma_0^{\mathbb{R}}, C_0^{\mathrm{cv}})$ stands in the
place of $(\Gamma(0;U,0), C(Z_0,0;U,0))$ in the two cluster-expansion lemmas, the
cluster-expansion proposition and the further proposition listed in the remark; with the small
parameter $\varepsilon$ replaced by $C_P\hat{\varepsilon}$, these four statements are then read word
for word at local data.

Now let a local datum be a representative of the class
\begin{equation*}
U^{\Sigma^1}_{k+1}\bigl(Z;(1+\beta)\alpha_0,(1+\beta)\alpha_1,\alpha_0\bigr).
\end{equation*}
At such a datum the vertex functional is the one constructed in
Proposition~\ref{4.4:prf-polydisc-positivity} (polydisc positivity and the vertex functional).
The function $f$, the exponentiated series, and
the persistence of \cite[Lemma~3]{B-CMP116} and of the proposition built on \cite{B-CMP119} under
the operation of this chapter hold by the auxiliary lemma, by the corollary on exponentiation and by
the statement of this chapter that establishes that persistence, respectively, each read at such a
datum.

Now let $Y' \in \mathbf{D}_k$ with $Y' \subset Z$. The passage from $Z$ to $Y'$ is given by the
heredity theorem for the one-power hereditary class, applied with equal levels, that is, with
$i = i'$ and scale ratio $\mu = L^{-(i-i')} = 1$ in the notation of that theorem. By it, such a
datum is also a representative of the class at $(k+1, Y')$, with the parameters that theorem
supplies, and the statements concerning these four items hold at $(k+1, Y')$.
\end{proof}

\begin{definition}[The fixed walk-expansion rule]\label{4.4:def-expansion-rule}
The fixed walk-expansion rule (in this chapter simply \emph{the rule}) assigns a walk expansion to
each operator $K$ of the list $\mathbb{L}$ of operators fixed earlier in this chapter, with its
sublists $\mathbb{L}_1$ and $\mathbb{L}_2$; $\mathsf{A}$ is the operator of the quadratic form.
The rule is fixed once and does not depend on the setting in which it is applied. It consists of
the following four clauses.
\begin{itemize}
\item[\textup{(R1)}] Let $K \in \mathbb{L}_1 \setminus \{\mathsf{A}\}$. The expansion of $K$ is
the one given by the explicit construction of the one-family theorem. That construction assigns to
a datum $\mathfrak{g}$ the following objects:
\begin{itemize}
\item the set $\mathcal{W}_K$ of walks;
\item for each walk $\omega \in \mathcal{W}_K$, its list of domains $\operatorname{dom}\omega$;
\item its two end sets $X_\pm(\omega)$;
\item the factor $\kappa_\omega$ that the construction attaches to the walk;
\item the term $K_\omega$.
\end{itemize}
The family is built from the cells $\mathcal{Q}_{\mathsf{m}}$ of the walk expansion of the
quadratic form, from the words of the inversion rule, and from the composition rules for local
letters, products, sums and branchings. Three further choices enter the construction.
\begin{itemize}
\item[(F1)] The smooth product profile
\begin{equation*}
  \eta_q(x) := \prod_{\mu} \psi\Bigl(\frac{x_\mu - c_\mu(q)}{\mathsf{m}\,\ell_n}\Bigr)
\end{equation*}
attached to a cell $q$, where $c_\mu(q)$ are the coordinates attached to the cell $q$ by the
one-family construction.
\item[(F2)] The lower bounds on $\mathsf{m}_1$ and $\mathsf{m}_3$, and the upper bounds on the
radii, that go with this profile.
\item[(F3)] The decompositions of the operators $W_\pi$ and $W_1$ of the one-family construction,
\begin{equation*}
  W_\pi = \mathsf{W}^0 + D\,\mathsf{W}^1, \qquad W_1 = \mathsf{W}_1^0 + D\,\mathsf{W}_1^1,
\end{equation*}
together with the operators
\begin{equation*}
  H_\bullet := G_\bullet\, Q^{T}\, T_\bullet^{-1} ,
\end{equation*}
which are the same operators as those of the one-family construction.
\end{itemize}
The domain of a step is the cube listed for it. The footprint of a walk $\omega$ is
\begin{equation*}
  [\omega] := \bigcup \operatorname{dom}\omega .
\end{equation*}
At a level $j \le k$, the family is the one given at that level by the corresponding corollary of
the one-family theorem. That this family has the properties (f1)--(f4) is the one-family
hypothesis of Definition~\ref{4.4:def-one-family-hypothesis}; it is assumed wherever the rule is
used for $K \in \mathbb{L}_1 \setminus \{\mathsf{A}\}$.

\item[\textup{(R1$'$)}] Let $K = \mathsf{A}$. The expansion of $\mathsf{A}$ is the family of walks
of the walk-expansion theorem for the quadratic form. This is the decomposition
\begin{equation*}
  \mathsf{A} = \sum_\omega \mathsf{A}_\omega
\end{equation*}
of Definition~\ref{4.4:def-form-walk-hypothesis}. Wherever two expansions of $\mathsf{A}$ meet in
a single object that is compared term by term, the terms are those of this family. The domain of a
step is its cube from $\operatorname{dom}\omega$, and the footprint is $[\omega]$.

\item[\textup{(R2)}] Let $K \in \mathbb{L}_2$. The expansion of $K$ is the cell-resolvent
expansion of Lemma~\ref{4.4:lem-cell-resolvent}. It coincides term by term with the expansion
produced by the inversion rule; this is the one-rule statement. The domain of a step that is a
cell $q$ is $q^\sharp$.

\item[\textup{(R3)}] The interpolation parameters $s(\Delta)$, $\Delta \in \sigma_0$, are attached
to every walk of every $K \in \mathbb{L}$ by the interpolation rule of
Definition~\ref{4.4:def-interpolation-rule}.
\end{itemize}
Ba{\l}aban's expansions on $M_1$-cubes in \cite{B-CMP116} are not part of the rule. In those
expansions all domains are small, and the domain of a step $(\alpha, X)$ is $\tilde{X}^{5}$. Their
locality is proved by the locality lemma for those expansions, and Ba{\l}aban's theorems bound them
under the condition on the configuration fixed for these expansions, with current $J = J(U)$. They
constitute a second expansion, distinct from the one fixed here.
\end{definition}

\begin{proposition}[Properties of the walk-expansion rule]\label{4.4:prop-rule-properties}
Let the walk-expansion rule of Definition~\ref{4.4:def-expansion-rule} be in force, with the
auxiliary parameters and the near--far split of Definition~\ref{4.4:def-auxiliary-parameters}.
The rule is used for the expansions of \cite[\S\S1--2]{B-CMP116} and of \cite[\S3]{B-CMP119}.

Throughout, $\mathsf{A}$ is the quadratic form, and $\mathbb{L}_1$ is the list of operators whose walk
families the rule takes from the hypotheses below: $\mathsf{A}$ from the hypothesis on the walk
expansion of the quadratic form, the others from the one-family hypothesis. $\mathbb{L}_2$ is the
list of operators whose families are produced
from that of $\mathsf{A}$ by the cell-resolvent expansion (Lemma~\ref{4.4:lem-cell-resolvent}).
For an operator $\mathsf{K}$ of either list:
\begin{itemize}
\item $\mathsf{K}_\omega$ is its term indexed by the walk $\omega$, and $\mathsf{K}_\omega[\mathfrak{s}]$ is
that term when the setting $\mathfrak{s}$ (with its nest of small-field regions) is displayed;
\item $\mathsf{K}(s)$ is the operator obtained when the walks carry the interpolation parameters $s$
of \cite[\S2]{B-CMP116} ($s$ is Ba{\l}aban's letter here, not the depth below the cutoff).
\end{itemize}
The sets $\mathbf{P}_\rho$, $\mathbf{P}_{2\rho}$, $\mathbf{P}'$ and $\mathcal{U}(1)$ are the analyticity
domains of the pairs $(\mathbf{U},\mathbf{J})$ used with the rule, and $\mathcal{S}_0$ is the setting
to which the hypothesis on the walk expansion of the quadratic form refers.

Assume:
\begin{enumerate}
\item[(i)] the hypothesis on the walk expansion of the quadratic form $\mathsf{A}$,
Definition~\ref{4.4:def-form-walk-hypothesis};
\item[(ii)] clauses (f1)--(f3) of the one-family hypothesis of
Definition~\ref{4.4:def-one-family-hypothesis}, for the operators of $\mathbb{L}_1$ other than
$\mathsf{A}$;
\item[(iii)] the reflection-covariance statement for the walk expansion of $\mathsf{A}$ (the
covariance clause of the theorem that supplies the walk expansion of $\mathsf{A}$), which is as
follows.

Let $r$ be an isometry of the torus that preserves the partition into level-$k$ cubes of side
$\mathsf{m}_4$ and that carries the nest of a setting $\mathfrak{s}$ to an admissible nest $r\mathfrak{s}$.
In the setting $\mathcal{S}_0$ every isometry preserving that partition carries the nest in this way;
in general, every isometry preserving the level-$k$ $LM$-cubes does.

For every such $r$ there is a bijection $\omega\mapsto r\omega$ of the walks of $\mathsf{A}$ in the
setting $\mathfrak{s}$ onto those in $r\mathfrak{s}$. It satisfies $\operatorname{dom}(r\omega)=r(\operatorname{dom}\omega)$ and
$X_\pm(r\omega)=rX_\pm(\omega)$, where $X_\pm(\omega)$ are the end sets of the walk. It satisfies
$a_{r\omega}=a_\omega$, where $a_\omega$ are the coefficients attached to the walks. Write
$\mathsf{u}=(\mathbf{U},\mathbf{J})$ for a point of the analyticity domain and $r\cdot\mathsf{u}$ for
its image under $r$. Finally, the bijection
satisfies
\begin{equation*}
\mathsf{A}_{r\omega}[r\mathfrak{s}](r\cdot\mathsf{u})
=R_r(\mathsf{u})\,\mathsf{A}_\omega[\mathfrak{s}](\mathsf{u})\,R_r(\mathsf{u})^{-1},
\end{equation*}
where $R_r$ is the orthogonal action on the fluctuation field of
\cite[(2.17)--(2.18)]{B-CMP109}.

This identity is one of analytic functions. It holds wherever both members are defined, in
particular at every point $\mathsf{u}$ with $\ell^3|\mathbf{J}|<a_0^{\oplus}/e$ and at every real point;
here $a_0^{\oplus}$ is the radius of the domain $\mathbf{P}'$ in the current variable $\mathbf{J}$, in
the normalisation $\ell^3|\mathbf{J}|$.
\end{enumerate}

Then the following hold.
\begin{enumerate}
\item[(c1)] At $s\equiv1$ every $\mathsf{K}(s)$ is the exact operator $\mathsf{K}$. Consequently the following
coincide with Ba{\l}aban's:
\begin{itemize}
\item the integral \cite[(3.25)]{B-CMP119} and its logarithm;
\item the term $\mathbf{E}^{(k+1)}_0(\Lambda_{k+1},b)$;
\item every total;
\item the quantities denoted $\Pi$ and $\beta$ in \cite[\S3]{B-CMP119}.
\end{itemize}
The rule fixes only the split of a total into terms localised in domains $X$.
\item[(c2)] The three families used in \cite[(2.6)--(2.8)]{B-CMP116}, namely the family of $\mathsf{A}$,
the family of the operator $\mathsf{M}^Y$ attached to the region $Y$ in
\cite[(2.6)--(2.8)]{B-CMP116}, and the family of the localised covariance $C^{(k)}(Z_0)$ of
Corollary~\ref{4.4:cor-localised-covariance}, satisfy
${|\kern-0.25ex|\kern-0.25ex|}\cdot{|\kern-0.25ex|\kern-0.25ex|}_\theta\le B_\theta$ on
$\mathbf{P}_\rho$, where ${|\kern-0.25ex|\kern-0.25ex|}\cdot{|\kern-0.25ex|\kern-0.25ex|}_\theta$ is the
walk norm of the hypothesis on the walk expansion of the quadratic form, at decay rate $\theta$.
Here $B_\theta$ is a number depending only on $d$, $L$ and the gauge group.
\item[(c3)] A walk that carries an interpolation factor of \cite[\S2]{B-CMP116} has a domain at
distance greater than $\tfrac23M$ from $Z_0$.
\item[(c4)] Let $Z$ be a union of cubes and let $s=0$ off $Z$. The kernels of $\mathsf{K}(s)$ vanish
unless both arguments lie in one connected component of $Z$; whence \cite[(2.10)]{B-CMP116}.
\item[(c5)] Each term is analytic on $\mathbf{P}_\rho$; the terms of the families of the hypotheses are
analytic on $\mathcal{U}(1)\supseteq\mathbf{P}'\cup\mathbf{P}_{2\rho}$. Each term depends on
$(\mathbf{U},\mathbf{J})$ only in its footprint.
\item[(c6)] For every gauge transformation $v$, with $\mathbf{U}^v$ the transformed configuration and
$R(v)$ the adjoint action of $v$ (not to be confused with the action $R_r$ of an isometry $r$ in
hypothesis (iii)),
\begin{equation}\label{4.4:eq-rule-properties-1}
\mathsf{K}_\omega\bigl(\mathbf{U}^v,R(v)\mathbf{J}\bigr)=R(v)\,\mathsf{K}_\omega(\mathbf{U},\mathbf{J})\,R(v)^{-1}.
\end{equation}
\item[(c7)] Let $r$ be a lattice symmetry that preserves the partition into $M$-cubes and the level-$k$
partition into $\mathsf{m}_4$-cubes, and that carries the nest to an admissible nest. Then
\begin{equation}\label{4.4:eq-rule-properties-2}
\mathsf{K}_{r\omega}[r\mathfrak{s}]=r\,\mathsf{K}_\omega[\mathfrak{s}]\,r^{-1}.
\end{equation}
In the notation of hypothesis (iii) this reads
\begin{equation*}
\mathsf{K}_{r\omega}[r\mathfrak{s}](r\cdot\mathsf{u})
=R_r(\mathsf{u})\,\mathsf{K}_\omega[\mathfrak{s}](\mathsf{u})\,R_r(\mathsf{u})^{-1}.
\end{equation*}
This holds on $\mathbf{P}'\cap r^{-1}\mathbf{P}'$, which contains all real points and all points
$(\mathbf{U},\mathbf{J})$ of the analyticity space whose current satisfies
$\ell^3|\mathbf{J}|<a_0^{\oplus}/e$, and it holds on $\mathbf{P}_{2\rho}$.
\end{enumerate}
\end{proposition}

\begin{proof}
The items are proved in order. Lemma~\ref{4.4:lem-cell-resolvent} is applied under its hypotheses:
hypothesis (i), and the remaining conditions of that lemma, $(\gamma)$, (k1) and (k2), which are
in force with the rule and the auxiliary parameters of
Definitions~\ref{4.4:def-expansion-rule} and~\ref{4.4:def-auxiliary-parameters}.

\emph{(c1).} Each family sums to its operator:
\begin{itemize}
\item for $\mathsf{A}$ this is hypothesis (i);
\item for the other operators of $\mathbb{L}_1$ it is clause (f1) of hypothesis (ii);
\item for the operators of $\mathbb{L}_2$ it is parts (i) and (iii) of
Lemma~\ref{4.4:lem-cell-resolvent}.
\end{itemize}
At $s\equiv1$ every interpolation factor equals one. So $\mathsf{K}(s)$ is the sum of the whole family
of $\mathsf{K}$, that is, the exact operator $\mathsf{K}$.

The integral \cite[(3.25)]{B-CMP119} and all the quantities built from it are defined through the
exact operators. These are its logarithm, $\mathbf{E}^{(k+1)}_0(\Lambda_{k+1},b)$, every total, and
the quantities $\Pi$ and $\beta$. They are therefore Ba{\l}aban's. Only the way a total is written as
a sum of terms localised in domains $X$ depends on the choice of the families.

\emph{(c2) and (c5).} The norm bounds, the analyticity and the dependence on $(\mathbf{U},\mathbf{J})$
through the footprint come from three sources:
\begin{itemize}
\item for the family of $\mathsf{A}$, from hypothesis (i), whose terms are analytic on
$\mathcal{U}(1)\supseteq\mathbf{P}'\cup\mathbf{P}_{2\rho}$;
\item for the families taken through the one-family hypothesis, from its clause (f1);
\item for the families produced by the cell-resolvent expansion, from parts (ii)--(iii) of
Lemma~\ref{4.4:lem-cell-resolvent}.
\end{itemize}
The constant $B_\theta$ is determined by the constants of these statements; through them it depends
on the radius $\rho_0$ among the constants fixed before the auxiliary parameters of
Definition~\ref{4.4:def-auxiliary-parameters}, which is itself a number depending only on $d$, $L$
and the gauge group. It is therefore a number depending only on $d$, $L$ and the gauge group.

\emph{(c3).} In \cite[\S2]{B-CMP116} the interpolation parameters are attached to the cubes of the
family $\sigma_0$. Let $\tilde Z_0$ be the domain of \cite[\S2]{B-CMP116} chosen around $Z_0$,
which contains $\{\operatorname{dist}(\cdot,Z_0)\le M\}$, and let $Z_0'$ be its $LM$-cube hull, the
smallest union of $LM$-cubes containing it; this is the reading of Ba{\l}aban's set $Z_0'$ adopted
in this chapter. The cubes of $\sigma_0$ are $LM$-cubes outside $Z_0'$, and
\begin{equation*}
Z_0'\supseteq\tilde Z_0\supseteq\{\operatorname{dist}(\cdot,Z_0)\le M\}.
\end{equation*}
Hence every point of such a cube lies at distance at least $M$ from $Z_0$.

A domain of a walk that meets one of these cubes has diameter at most $\max(\rho_w,4\ell)$. The
bound $\rho_w$ is that of hypothesis (i) for the domains of walks of $\mathsf{A}$; the bound $4\ell$
is that for the cell domains $q^\sharp$ of Lemma~\ref{4.4:lem-cell-resolvent}.

Write $r_{\mathrm{c}}$ and $r'_{\mathrm{c}}$ for the two range parameters of conditions (k2)--(k3) of
Definition~\ref{4.4:def-auxiliary-parameters}; the letter $r$ is reserved in this proposition for
isometries. Those conditions give
\begin{equation*}
\ell\ \ge\ 2r_{\mathrm{c}}+4\ =\ 2r'_{\mathrm{c}}+4\rho_w+4\ >\ 4\rho_w,
\qquad M\ \ge\ 48\ell .
\end{equation*}
Hence $\max(\rho_w,4\ell)=4\ell\le M/12$.

In particular
\begin{equation*}
M\ >\ 192\rho_w\ \ge\ 12\rho_w .
\end{equation*}
The bound $M>12\rho_w$ is thus implied by (k2)--(k3). It is a one-sided floor on $M$, of the same
kind as Ba{\l}aban's requirement that $M$ be chosen much larger than $M_2$ \cite{B-CMP119}. It is not
an additional condition and places no upper bound on $M$.

Consider now a domain that contains a point at distance at least $M$ from $Z_0$ and has diameter
at most $M/12$. Each of its points lies at distance at least
\begin{equation*}
M-\frac{M}{12}=\frac{11}{12}M>\frac23M
\end{equation*}
from $Z_0$. This proves (c3).

The argument replaces Ba{\l}aban's special class of domains, introduced in \cite[p.~13]{B-CMP116}
by the words \emph{``We choose $\tilde Z_0$ as one of the domains, more exactly we take a domain
$X_0$ such, that $\tilde X_0^5=\tilde Z_0$ \dots\ $\operatorname{dist}(X,Z_0)>\frac23M$''}. This replacement is
permitted by the reading of the interpolation domains of \cite[\S2]{B-CMP116} adopted in this
chapter.

\emph{(c4).} Two distinct components of $Z$ share no cube and no wall \cite[\S2]{B-CMP116}.

First, $\operatorname{int}Z$ is the disjoint union of the interiors of the components. Let $y$ be an interior
point of $Z$. Since $y$ is interior to $Z$ and the cubes have pairwise disjoint interiors, every cube
containing $y$ lies in $Z$. These cubes form a wall-connected family, and two components share no
wall; so they lie in one component. Their union is a neighbourhood of $y$, so
$y$ is an interior point of that component.

Second, each domain lies in a single component. The domains in question are the cells $q^\sharp$,
the domains $[\omega]$ of the walks of $\mathsf{A}$, the footprints $[\omega]$ of the families of the
one-family hypothesis, and Ba{\l}aban's domains $\tilde X^5$. The cells and Ba{\l}aban's domains are
closures of connected open sets by construction; for the walk domains of $\mathsf{A}$ and the
footprints this is part of hypotheses (i) and (ii). Let such a domain
be contained in $Z$. The connected open set of which it is the closure is contained in
$\operatorname{int}Z$, hence, by the first step, in the interior of one component. Its closure, the domain,
then lies in that component, which is closed.

Third, a whole walk lies in one component. With $s=0$ off $Z$, a walk one of whose domains is not
contained in $Z$ carries a vanishing interpolation factor, so only walks all of whose domains lie
in $Z$ contribute to the kernel. Two consecutive steps of such a walk share a bond at distance at
least one from the boundary of each of the two domains. The points of that bond are interior points
of $Z$, so the two
domains lie in the same component. By induction along the walk, all its domains lie in one component
of $Z$; in particular the two arguments of the kernel do. This gives (c4), and with it
\cite[(2.10)]{B-CMP116}.

\emph{(c6).} The identity holds for the families of $\mathbb{L}_1$ as follows:
\begin{itemize}
\item for the family of $\mathsf{A}$, it is the gauge covariance carried by hypothesis (i);
\item for the other families of $\mathbb{L}_1$, it is the gauge-covariance clause of the one-family
hypothesis, which belongs to its clause (f2).
\end{itemize}
The same identity holds for the objects of Ba{\l}aban's expansion
\cite[(3.28)--(3.34)]{B-CMP99b} and \cite[Lemma~2, last clause]{B-CMP116}.

For the families of $\mathbb{L}_2$, note that $R(v)$ is diagonal in the bond $b$. Hence it commutes
with every characteristic projection $\chi_X$. The operators $T_q$ and $(x+T_q)^{-1}$ of
Lemma~\ref{4.4:lem-cell-resolvent} are built from $\mathsf{A}$ and such projections. By the
identity for $\mathsf{A}$ they transform by conjugation with $R(v)$. This uses only algebra; the
unitarity of $R(v)$ is not used.

\emph{(c7).} \emph{The class of $r$ and the domain.} The isometries considered are those of (c7): they
preserve the $M$-partition and the level-$k$ $\mathsf{m}_4$-partition, with $\mathsf{m}_4$ dividing
$M$, and they carry the nest to an admissible nest. In the setting $\mathcal{S}_0$ every isometry
preserving these partitions is of this kind.

By hypothesis (iii), the identity for $\mathsf{A}$ holds wherever both members are defined, in
particular at every real point and wherever $\ell^3|\mathbf{J}|<a_0^{\oplus}/e$. On
$\mathbf{P}_{2\rho}$, the radius in the variable $\mathbf{J}$ is a fraction at most $\frac18<\frac1e$
of the corresponding radius of $\mathbf{P}'$. Hence $\mathbf{P}_{2\rho}$ lies in that domain.

\emph{Translations.} Consider translations by vectors of $M\mathbb{Z}^4$. Together with lifts, these are
the only symmetries that the lemma on the transport of the localised terms under lifts and
translations uses. Here the identity
follows by transport of structure, because each of the following is
defined by a recipe that commutes with translations:
\begin{itemize}
\item the family of $\mathsf{A}$ and the families of the one-family hypothesis (hypothesis (iii) and
clause (f3) of hypothesis (ii) cover every isometry preserving the level-$k$ $\mathsf{m}_4$-partition);
\item the cell data of Lemma~\ref{4.4:lem-cell-resolvent};
\item every object of Ba{\l}aban's expansion \cite[(3.28)--(3.34)]{B-CMP99b}, in his notation (these
letters are his and are distinct from the coefficients $a_\omega$ and the projections $\chi_X$
above): his cut-off and partition-of-unity functions $h_\square$, $\chi_\square$, $\zeta_\square$,
his averaging operators $Q$, $Q'$, his operator $a$, the operators $K(h)$, $S(\partial h)$, and
the Dirichlet operators of the concentric settings $\mathfrak{s}_\square$.
\end{itemize}
In addition, translations preserve initial sites of bonds.

\emph{Reflections.} Reflections, and more generally the elements of the hyperoctahedral group $W_4$,
are read term by term. This happens in the covariance of the localised terms that accompanies the
theorem on independence of the localised terms from the small-field region. That covariance is the
reflection covariance of the localised terms of \cite{B-CMP119}, read term by term. For
reflections, therefore, the assignment of terms in each family must
itself be reflection covariant:
\begin{itemize}
\item for $\mathsf{A}$ this is hypothesis (iii);
\item for the other operators of $\mathbb{L}_1$ it is clause (f3) of hypothesis (ii).
\end{itemize}

For Ba{\l}aban's own expansion this fails term by term. Ba{\l}aban's cut-offs multiply bond fields
at the initial site of the bond, \cite[(3.100), (3.102), (3.103)]{B-CMP99b}. Here, in Ba{\l}aban's
notation, $h$ is the cut-off function, $\eta$ the lattice spacing, and $b_-$, $c_-$ are the initial
sites of the bonds $b$, $c$:
\begin{align*}
(D_\mu hA_\nu)(x)&=h(x)(D_\mu A_\nu)(x)
+(\partial_\mu h)(x)R\bigl(U(x,x+\eta e_\mu)\bigr)A_\nu(x+\eta e_\mu),\\
(Q_jhA)(c)&=h(c_-)(Q_jA)(c)+\dots,\\
(Q^*aQhA)(b)&=h(b_-)(Q^*aQA)(b)-\dots .
\end{align*}
A reflection that reverses $e_\nu$ carries the initial site $b_-$ of a bond in direction $\nu$ to the
final site of the image bond. Ba{\l}aban's walks therefore do not carry reflection symmetry term by
term; this is one reason why the families of the rule are not his.

For the operators of $\mathbb{L}_2$, the cell partition is the standard partition into cubes of side
$\ell$, and $\ell$ divides $M$. A symmetry preserving the $M$-partition therefore preserves, for
reflections as well:
\begin{itemize}
\item the cell partition;
\item the cubes $q^*$ and $q^\sharp$;
\item the relation $b\subset q^*$;
\item the weights $\kappa_q$.
\end{itemize}
The quantity attached to a walk $\omega$ that decides the near--far split of
Definition~\ref{4.4:def-auxiliary-parameters} depends on the domains of $\omega$ only, and hence so
does that split. The terms of
Lemma~\ref{4.4:lem-cell-resolvent} thus inherit (c7) from those of $\mathsf{A}$.

Hence the term-by-term reflection covariance used for the localised terms follows from hypothesis
(iii) and from clause (f3) of hypothesis (ii). The distance $|b-b'|$ between bonds enters bounds only,
never a definition.
\end{proof}

\begin{remark}[The cluster-expansion lemmas under the fixed rule at local data]
\label{4.4:rem-cluster-lemmas-local}
Let $Z_0\in\mathbf{D}_k$, let $Z'_0$ be the enlarged domain of
Corollary~\ref{4.4:cor-polymer-base-form}, on which the localised operator
$\mathsf{T}_{Z'_0}$ is formed, so that $Z'_0$ contains every point within sup-distance $M$
of $Z_0$, and let $Z\supseteq Z'_0$. A \emph{local datum} is a point $\mathsf{u}$ of a polydisc
$\mathbf{P}_{\varsigma,Z}$ of data on $Z$ with radius index $\varsigma$. We write $\sigma$ for the
complex parameter of the polydisc on which Proposition~\ref{4.4:prop-local-real-base}
establishes positivity. Assume the hypotheses of Proposition~\ref{4.4:prop-local-real-base},
in particular its lower bound (R9) on $M$, read with the pair $(\mathsf{B},\delta_{\mathsf{A}})$
of (A) below replaced by
\begin{equation*}
\Bigl(\max\bigl\{\mathsf{B},\,4S(2/\gamma_*),\,4S(2/\gamma_*)^{1/2}\bigr\},\ \theta\Bigr),
\end{equation*}
where $\theta$ is the decay rate of the norm of item (c2) of
Proposition~\ref{4.4:prop-rule-properties}, $S$ is the constant of the base-covariance lemma in
(A) below and $\gamma_*$ is the coercivity constant of Lemma~\ref{4.4:lem-local-coercivity};
assume its three further lower bounds on $M$, and its
hypotheses (k2) and (k3), and assume the following statements.
\begin{enumerate}
\item[(A)] The two walk-expansion lemmas and the base-covariance lemma by which the bounds
\cite[(2.16)--(2.17)]{B-CMP116} are obtained for the walk-expansion rule. The two walk-expansion
lemmas apply to families of kernels that satisfy their weighted summability hypothesis. The
base-covariance lemma has items (i)--(iv) and is formulated with constants $\varepsilon$,
$\tfrac12 c_1$, $1/(2B'S)$, with the decay bound $(4/c_1)e^{-\theta_1|b-b'|}$ on the inverse
base covariance, and with a constant $E'$; here $B'$, $S$, $c_1$, $\theta_1$ are constants of
that lemma. We write $\mathsf{B}$ and $\delta_{\mathsf{A}}$ for the size constant and the decay
rate of the base operator in the hypotheses of the lemmas of (A) and (B).
\item[(B)] The lemma that completes the proof of the bounds \cite[(2.16)--(2.17)]{B-CMP116}
under the walk-expansion rule, one of whose inputs is the inverse base covariance (the symmetrised
base operator of Proposition~\ref{4.4:prop-local-real-base}), and the accompanying bound
for \cite[(2.23)--(2.26)]{B-CMP116} stated with it.
\item[(C)] The local real base of a polymer, Proposition~\ref{4.4:prop-local-real-base}.
\item[(D)] The restriction theorem for $\Sigma$-representatives at equal levels $i=i'$. A
$\Sigma$-representative at a level, a domain and three radii is understood in the sense of that
theorem.
\end{enumerate}
Then \cite[Lemma~3]{B-CMP116} and the proposition of this book that is its analogue for the
expansions of \cite[\S3]{B-CMP119} hold under the walk-expansion rule of
Proposition~\ref{4.4:prop-rule-properties}, with unchanged form, at every local datum
$\mathsf{u}\in\mathbf{P}_{\varsigma,Z}$, $\varsigma\le 1+\beta$, which is a
$\Sigma$-representative at
\begin{equation*}
\bigl(k+1,Z;(1+\beta)\alpha_0,(1+\beta)\alpha_1,\alpha_0\bigr).
\end{equation*}
For every such $\mathsf{u}$ and every $Y'\in\mathbf{D}_k$ with $Y'\subset Z$, the restriction
$\mathsf{u}\upharpoonright Y'$ is a $\Sigma$-representative at level $k+1$ on $Y'$, with the radii
assigned by the restriction theorem of (D). Moreover, items (i)--(iv) of the base-covariance
lemma of (A) hold as stated at the local real base of (C), with $C_P\hat\varepsilon$,
$\tfrac58\gamma_*/\mathsf{B}^2$, $\tfrac58\gamma_*$ in place of $\varepsilon$, $\tfrac12c_1$,
$1/(2B'S)$, with the decay bound $\mathsf{B}e^{-\delta_{\mathsf{A}}|b-b'|}$ in place of
$(4/c_1)e^{-\theta_1|b-b'|}$, with the constant $E'$ replaced by the bound supplied by the two
walk-expansion lemmas of (A) increased by the share that comes from the symmetrisation, and
under the three further lower bounds on $M$ among the hypotheses of
Proposition~\ref{4.4:prop-local-real-base}; here $C_P\hat\varepsilon$ is the smallness parameter
of the lemma of (B) and $\gamma_*$ is the coercivity constant of
Lemma~\ref{4.4:lem-local-coercivity}.
\end{remark}

\begin{proof}
We first consider the families. The three families named in items (c2) and (c3) of the
properties of the walk-expansion rule (Proposition~\ref{4.4:prop-rule-properties}), measured in
the norm $\lvert\!\lvert\!\lvert\cdot\rvert\!\rvert\!\rvert_{\theta}$ of item (c2), where
$\theta$ is the decay rate of that norm, satisfy, by the bound
$\lvert\!\lvert\!\lvert\cdot\rvert\!\rvert\!\rvert_{\theta}\le B$ of item (c2) together with
item (c3), the weighted summability hypothesis of the two walk-expansion lemmas in (A).
Consequently the lemmas listed in (A) and (B) may be applied to these three families on the
polydisc $\mathbf{P}_{\rho,Z}$, where $\rho$ is the radius index of the polydisc
$\mathbf{P}_{\rho}$ of item (c2) of Proposition~\ref{4.4:prop-rule-properties}, that is, at
local data on $Z\supseteq Z'_0$, so that the words ``on the domain $Z$'' in \cite{B-CMP116}
apply literally. In their standing lower bound on $M$, which is hypothesis (R9)
of Proposition~\ref{4.4:prop-local-real-base}, the pair $(\mathsf{B},\delta_{\mathsf{A}})$ is
replaced by
\begin{equation*}
(\mathsf{B},\delta_{\mathsf{A}})\longmapsto
\Bigl(\max\bigl\{\mathsf{B},\,4S(2/\gamma_*),\,4S(2/\gamma_*)^{1/2}\bigr\},\ \theta\Bigr),
\end{equation*}
where $\gamma_*$ is the coercivity constant of Lemma~\ref{4.4:lem-local-coercivity}.

Next we identify the base. At local data, the base at $\sigma=0$ and at the real datum $(U,0)$ of
Proposition~\ref{4.4:prop-local-real-base}, from which
\cite[(2.15)--(2.17)]{B-CMP116} and \cite[(2.23)--(2.26)]{B-CMP116} are measured, is the local
real base of Proposition~\ref{4.4:prop-local-real-base}. It is the pair
\begin{equation*}
\bigl(\Gamma_0^{\mathbb{R}},\,\mathsf{C}^{\mathrm{cv}}_0\bigr)
:=\bigl(\operatorname{Re}\Gamma(0;U\upharpoonright Z,0),\ \mathsf{G}_0^{-1}\bigr),
\qquad
\mathsf{G}_0:=\operatorname{sym}\operatorname{Re}\mathsf{T}_{Z'_0}(U\upharpoonright Z'_0,0),
\end{equation*}
where $\Gamma$ denotes the vertex functional and $\operatorname{sym}$ denotes the symmetric part,
as in Proposition~\ref{4.4:prop-local-real-base}.
This base is real. The coercivity of the localised quadratic form
(Lemma~\ref{4.4:lem-local-coercivity}) and the symmetrisation lemma, in the form of the
symmetric base form on a polymer (Corollary~\ref{4.4:cor-polymer-base-form}), show that
$\mathsf{G}_0$ is real symmetric and satisfies
\begin{equation*}
\tfrac58\gamma_*\le\mathsf{G}_0\le\mathsf{B},
\qquad
|\mathsf{G}_0(b,b')|\le\mathsf{B}\,e^{-\delta_{\mathsf{A}}|b-b'|}.
\end{equation*}
The operator $\mathsf{G}_0$ is the inverse base covariance that the lemma of (B) takes as input.
It is also the inverse
base covariance $(\mathsf{C}^{\mathrm{cv}}_0)^{-1}$ of the base-covariance lemma in (A). At
local data these properties are therefore proved, and are not assumed.

We turn to the items of the base-covariance lemma. Items (i)--(iv) of that lemma hold as they
are stated, at this base, after the following substitutions:
\begin{equation*}
\varepsilon\longmapsto C_P\hat\varepsilon,\qquad
\tfrac12 c_1\longmapsto \tfrac58\gamma_*/\mathsf{B}^2,\qquad
1/(2B'S)\longmapsto\tfrac58\gamma_*,
\end{equation*}
where $C_P\hat\varepsilon$ is the smallness parameter of the lemma of (B).
The decay bound $(4/c_1)e^{-\theta_1|b-b'|}$ is replaced by
\begin{equation*}
|(\mathsf{C}^{\mathrm{cv}}_0)^{-1}(b,b')|\le\mathsf{B}\,e^{-\delta_{\mathsf{A}}|b-b'|},
\end{equation*}
which is supplied by the decay estimate for the symmetrised base operator; it is the decay bound
for $\mathsf{G}_0$ displayed above. The
constant $E'$ is the bound supplied by the two walk-expansion
lemmas, increased by the share that comes from the symmetrisation. All of this holds under the
three further lower bounds on $M$ among the hypotheses of
Proposition~\ref{4.4:prop-local-real-base}: the one on the small factor $\bar\eta$ of that
proposition, the one for the walk-expansion rule, and the one for the symmetrisation.

Positivity on the whole $\sigma$-polydisc is already available. Part~(3) of
Proposition~\ref{4.4:prop-local-real-base} gives
\begin{equation*}
\operatorname{Re}C(Z_0,\sigma;\mathsf{u})^{-1}\ge\tfrac38\gamma_*,
\end{equation*}
where $C(Z_0,\sigma;\mathsf{u})$ is the covariance on $Z_0$ at the parameter $\sigma$ and the
datum $\mathsf{u}$ in that proposition, and it does so from hypothesis (R9), the lower bound on
$M$, together with hypotheses (k2) and (k3) of Proposition~\ref{4.4:prop-local-real-base}.

By the restriction theorem at $i=i'$ in (D), in the form of part~(4c) of
Proposition~\ref{4.4:prop-local-real-base}, for every local datum $\mathsf{u}$ of
$\mathbf{P}_{\varsigma,Z}$, $\varsigma\le 1+\beta$, which is a $\Sigma$-representative at
$(k+1,Z;(1+\beta)\alpha_0,(1+\beta)\alpha_1,\alpha_0)$, and for every $Y'\in\mathbf{D}_k$ with
$Y'\subset Z$, the restriction $\mathsf{u}\upharpoonright Y'$ is a $\Sigma$-representative at
level $k+1$ on $Y'$, with the radii assigned by that theorem. The inputs established in the three
preceding paragraphs --- the summability of the three families, the bounds at the local real
base and the positivity on the $\sigma$-polydisc --- are available at every such local datum
$\mathsf{u}$. Therefore \cite[Lemma~3]{B-CMP116} and the proposition of this book that is its
analogue for the expansions of \cite[\S3]{B-CMP119} hold under the walk-expansion rule with
unchanged form, at every such local datum $\mathsf{u}$.
\end{proof}

\begin{theorem}[Locality of the rule and agreement on uniform regions]
\label{4.4:thm-uniform-agreement}
Let $\mathbb{L}$ be the list of operators through which non-locality enters the $(k+1)$-st
$\mathbf{T}$-step of \cite{B-CMP119}. It consists of the operators of $\mathbb{L}_1$, among
them the quadratic form $\mathsf{A}$, and of the operators of $\mathbb{L}_2$, which are
expanded by the cell-resolvent expansion of Lemma~\ref{4.4:lem-cell-resolvent}. The fixed
walk-expansion rule of Definition~\ref{4.4:def-expansion-rule} writes every
$\mathsf{K}\in\mathbb{L}$ as a sum $\mathsf{K}=\sum_\omega\mathsf{K}_\omega$ over walks
$\omega$, each carrying a footprint $|\omega|$. Let $\mathcal{S}$ and $\mathcal{S}_0$ be the
two settings of the $(k+1)$-st step. Assume:
\begin{enumerate}
\item the hypothesis on the walk expansion of the quadratic form
(Definition~\ref{4.4:def-form-walk-hypothesis});
\item the one-walk-family hypothesis (Definition~\ref{4.4:def-one-family-hypothesis}) for
the operators of $\mathbb{L}_1\setminus\{\mathsf{A}\}$;
\item the positivity hypothesis at the setting $\mathcal{S}_0$;
\item the walk-covering hypothesis;
\end{enumerate}
all four read with the reading convention that accompanies them in this chapter. For the
assertions of part~(d) concerning the setting $\mathcal{S}$, assume moreover the
admissibility lemma for the nests made by $\mathbf{R}$. Then the following hold.
\begin{enumerate}
\item[(a)] The rule of Definition~\ref{4.4:def-expansion-rule} is local, in the sense of
Definition~\ref{4.4:def-local-rules}, for every $\mathsf{K}\in\mathbb{L}$.
\item[(b)] If a region $W$ is uniform for a pair of settings
$(\mathfrak{s},\mathfrak{s}')$, in the sense in which Definition~\ref{4.4:def-local-rules}
uses the term, then $\mathsf{K}^{W'}[\mathfrak{s}]
=\mathsf{K}^{W'}[\mathfrak{s}']$, term by term, for every $\mathsf{K}\in\mathbb{L}$ and
every closed union $W'\subset W$ of $M_1$-cubes; here $\mathsf{K}^{W'}$ is the part of the
walk expansion of $\mathsf{K}$ carried by the walks $\omega$ with $|\omega|\subset W'$.
\item[(c)] Every closed union $W\subset\Lambda_{k+1}$ of $M_1$-cubes is uniform for the pair
of settings $(\mathcal{S},\mathcal{S}_0)$. Hence the locality statement for the
$(k+1)$-st step holds: for every $\mathsf{K}\in\mathbb{L}$ and every walk $\omega$ with
$|\omega|\subset\Lambda_{k+1}$, $\omega$ is a walk of $\mathsf{K}$ in $\mathcal{S}$ if and
only if it is a walk of $\mathsf{K}$ in $\mathcal{S}_0$, and
\begin{equation*}
\mathsf{K}_\omega[\mathcal{S}]=\mathsf{K}_\omega[\mathcal{S}_0]
\end{equation*}
as functions on the local class of $|\omega|$; the terms whose large-field region $X$ touches
the rim are included, by the rim-touching clause of the standing hypotheses of this chapter.
\item[(d)] In the setting $\mathcal{S}_0$, the hypothesis on the walk expansion of the
quadratic form, clause~(f1) of the one-walk-family hypothesis, and clauses~(ii) and~(iii) of
Lemma~\ref{4.4:lem-cell-resolvent} hold with constants depending only on $d$, $L$ and the
gauge group $G=SU(N)$.
In the setting $\mathcal{S}$, the hypothesis on the walk expansion of the quadratic form
holds among admissible nests, by the walk-expansion theorem for the quadratic form and, at
the nests made by $\mathbf{R}$, by the admissibility lemma; Lemma~\ref{4.4:lem-cell-resolvent}
holds among admissible nests likewise. The one-walk-family hypothesis restricted to
$\mathcal{S}$ is neither assumed nor used. In $\mathcal{S}$, of the two properties of the
walk family that Definition~\ref{4.4:def-expansion-rule} assigns to the operators of
$\mathbb{L}_1\setminus\{\mathsf{A}\}$, the first holds in $\ell^2$ on every admissible
nest, and the second holds at trivial traces only. Not asserted: the norms of
\cite{B-CMP116} on seam walks, which are the seam-norm bounds resting on the transfer of
the analysis of \cite[\S3]{B-CMP109} to the cover, printed without proof in
\cite[p.~276]{B-CMP119}; and the operators at free current $\mathbf{J}$ in $\mathcal{S}$,
which are those of the operator bounds at free current restricted to $\mathcal{S}$.
\end{enumerate}
\end{theorem}

\begin{proof}
(a) Let $\mathfrak{s}$, $\mathfrak{s}'$ be two settings and $W$ a region uniform for them;
we verify the index condition, the value condition and the support condition of
Definition~\ref{4.4:def-local-rules} for each $\mathsf{K}\in\mathbb{L}$.

For $\mathsf{K}\in\mathbb{L}_1\setminus\{\mathsf{A}\}$ the three conditions are clause~(f2)
of the one-walk-family hypothesis (Definition~\ref{4.4:def-one-family-hypothesis}).

For $\mathsf{K}=\mathsf{A}$ the three conditions are part of hypothesis~(1): they are among
the statements of the corollary through which the hypothesis on the walk expansion of the
quadratic form (Definition~\ref{4.4:def-form-walk-hypothesis}) is given.

Let $\mathsf{K}\in\mathbb{L}_2$. By Lemma~\ref{4.4:lem-cell-resolvent}, a walk $\varpi$ of
$\mathsf{K}$ is a sequence whose steps are cells $q$ of the cell cover $\mathcal{Q}$, with
domains $q^{\sharp}$, and far walks of $\mathsf{A}$. Suppose $|\varpi|\subset W$. The cells
$q$ with $q^{\sharp}\subset W$ are the same in both settings, and the far walks $\omega$ of
$\mathsf{A}$ with $|\omega|\subset W$ are the same in both settings by the index condition
for $\mathsf{A}$, established in the preceding paragraph. Hence the walks $\varpi$ of
$\mathsf{K}$ with $|\varpi|\subset W$ are the same in both settings; this is the index
condition for $\mathsf{K}$.

By clause~(iv) of Lemma~\ref{4.4:lem-cell-resolvent}, the term $\mathsf{K}_\varpi$ is a
function of the family $\{\mathsf{A}_\omega : |\omega|\subset|\varpi|\}$ and of the bonds of
the bond set $Y'$ to which $\mathsf{K}$ is compressed that lie at distance $\geq 1$ from
$\partial|\varpi|$. The family is the same in both settings by the value condition for
$\mathsf{A}$. The bonds are the same in both settings because $W$ is uniform for
$(\mathfrak{s},\mathfrak{s}')$; when $Y'$ is the bond set $\mathbb{B}(Z_0)$ of a region
$Z_0$, the set $Z_0$ belongs to the index
of the quantity considered and not to the setting. Hence
$\mathsf{K}_\varpi[\mathfrak{s}]=\mathsf{K}_\varpi[\mathfrak{s}']$ as functions on the
polydisc $\mathbf{P}_{\rho,|\varpi|}$, where the term exists by clause~(v) of
Lemma~\ref{4.4:lem-cell-resolvent}; this is the value condition for $\mathsf{K}$. The
support condition for $\mathsf{K}$ is clause~(iv) of Lemma~\ref{4.4:lem-cell-resolvent}.

(b) This is the definition of $\mathsf{K}^{W'}$, whose terms are the terms
$\mathsf{K}_\omega$ with $|\omega|\subset W'\subset W$; by part~(a) their index sets and
their values agree in the two settings.

(c) Let $W\subset\Lambda_{k+1}$ be a closed union of $M_1$-cubes. By
\cite[(3.20)]{B-CMP119}, $\Lambda_{k+1}\subset\Omega_{k+1}$, so that
$W\subset\Lambda_{k+1}\subset\Omega_{k+1}$; moreover $W\subset T$, the region of the setting
$\mathcal{S}_0$. A bond inside $W$ at distance $\geq 1$ from $\partial W$ has both endpoints
in the interior of $\Lambda_{k+1}$. It therefore belongs to $\mathbb{B}(\Lambda_{k+1})$ under
any convention for the bonds of a domain \cite{B-CMP96}, and it belongs to $\mathbb{B}(T)$.
The bond $b_{\ast}(c)$ attached to each cube $c$ and the $L$-blocks are the standard ones in
both settings. Hence $W$ is uniform for $(\mathcal{S},\mathcal{S}_0)$. By part~(a), applied
to these regions $W$, the walks with footprint in such a region are the same in the two
settings and their terms satisfy $\mathsf{K}_\omega[\mathcal{S}]=\mathsf{K}_\omega[\mathcal{S}_0]$
as functions on the local class of $|\omega|$; hence the locality statement holds, the terms
of large-field regions $X$ touching the rim included by the rim-touching clause.

(d) The hypothesis on the walk expansion of the quadratic form is uniform along the
sequence by its statement (Definition~\ref{4.4:def-form-walk-hypothesis}). The bounds of the
one-walk-family hypothesis are stated at $\mathcal{S}_0$
(Definition~\ref{4.4:def-one-family-hypothesis}). The coercivity condition $(\gamma)$ of
Lemma~\ref{4.4:lem-form-coercivity} holds for $\Lambda_{k+1}$ and for $T$. This gives the
assertions at $\mathcal{S}_0$ with constants depending only on $d$, $L$ and $G$; in the
setting $\mathcal{S}$ the assertions are those supplied, among admissible nests, by the
walk-expansion theorem for the quadratic form and by the admissibility lemma at the nests
made by $\mathbf{R}$, which is assumed.
\end{proof}

\begin{lemma}[Interior sufficiency of the whole-lattice setting]\label{4.4:lem-interior-sufficiency}
Assume the hypothesis on the quadratic form's walk expansion
(Definition~\ref{4.4:def-form-walk-hypothesis}) and items (f1) and (f2) of the one-family
hypothesis (Definition~\ref{4.4:def-one-family-hypothesis}), each together with its local form.
Assume also the positivity hypothesis at the whole-lattice geometry $\mathcal{S}_0$ and the
ceilings on the primed radii (Definition~\ref{4.4:def-radius-ceilings}).
Let $\mathfrak{s}$ be any geometry, and let $W$ be uniform for the pair
$(\mathfrak{s},\mathcal{S}_0)$. For instance, one may take $\mathfrak{s}=\mathcal{S}$ and $W$ a
closed union of $M_1$-cubes in $\Lambda_{k+1}$; this $W$ is uniform for $(\mathcal{S},\mathcal{S}_0)$
by part~(c) of Theorem~\ref{4.4:thm-uniform-agreement}. Let
$\mathsf{u}$ be a local datum on $W$
(Definition~\ref{4.4:def-local-data-classes}), with $\mathsf{u}\in\mathbf{P}'_W$ for the
constituents $K$ of the rule in the class $\mathbb{L}_1$ and $\mathsf{u}\in\mathbf{P}_{\rho,W}$ for
those in the class $\mathbb{L}_2$. Then for every constituent
$K\in\mathbb{L}$ of the rule and every $W'\subset W$ the series $K^{W'}[\mathfrak{s}](\mathsf{u})$
converges absolutely and coincides, term by term, with $K^{W'}[\mathcal{S}_0](\mathsf{u})$. No bound
in the geometry $\mathcal{S}$ enters, hence no admissibility condition on $\mathcal{S}$; and no
point of $\mathcal{S}_0$ extending $\mathsf{u}$ is used.
\end{lemma}

\begin{proof}
For a walk $\omega$ we write $[\omega]$ for its domain,
$|\omega|$ for its support, $d(\omega)$ for its size, and $\mathsf{u}|_{[\omega]}$ for the
restriction of the local datum $\mathsf{u}$ to $[\omega]$.

We first compare the index sets and the terms of the two series. On $W$ the traces of the nests
of both geometries are trivial: $W$ is contained in $\Omega_{k+1}$, which lies in the interior of
$\Omega_k$, and
\begin{equation*}
\Omega_k\subset\Omega_{k-1}\subset\cdots\subset\Omega_0 .
\end{equation*}
Of the geometry $\mathfrak{s}$ nothing else is used than this nestedness and one margin. When
$W\subset\Lambda_{k+1}$, that margin is the one of \cite[(3.20)]{B-CMP119}.

We call the theorem that supplies the hypothesis on the quadratic form's walk expansion
(Definition~\ref{4.4:def-form-walk-hypothesis}) the walk-expansion theorem. By its locality item,
the part of the rule attached to $W$ is determined by the traces of the nest on
$W$ alone. It is therefore defined for any geometry whose traces on $W$ are trivial, whether or not
the remainder of its nest satisfies the admissibility conditions.

Consider first the members of $\mathbb{L}_1$ other than the quadratic form $\mathsf{A}$.
By item (f2) of the one-family hypothesis, the walks $\omega$ with
$[\omega]\subset W'$ are the same for $\mathfrak{s}$ and for $\mathcal{S}_0$. They carry the same
domains and the same weights, and the term $K_\omega$ is, in both geometries, the same function of
the datum on $[\omega]$. For $\mathsf{A}$ the same conclusions follow from the shell item of the
walk-expansion theorem at radius~$0$ and from the first of its two kernel bounds for $\mathsf{A}$.

The term $K_\omega$ is defined at $\mathsf{u}|_{[\omega]}$ by the locality of the increments.
Consequently
\begin{equation}\label{4.4:eq-interior-sufficiency-1}
K^{W'}[\mathfrak{s}](\mathsf{u})
=\sum_{\omega:\ [\omega]\subset W'}K_\omega\bigl(\mathsf{u}|_{[\omega]}\bigr)
=K^{W'}[\mathcal{S}_0](\mathsf{u}).
\end{equation}
These equalities hold as series with one index set and identical terms, and they do so before any
question of convergence arises.

Next we prove convergence. Let $\omega$ be a walk with $[\omega]\subset W'$. The term bound of the
walk-expansion theorem holds on the local domain $\mathcal{U}^{\flat}_{[\omega]}$ by the local form
of item (f1). It bounds the norm of $K_\omega(\mathsf{u}|_{[\omega]})$ by
$\kappa_\omega\,\ell(X_-(\omega))^{\alpha}$, a number attached to $\omega$; here $\kappa_\omega$,
$\ell$, $X_-(\omega)$ and $\alpha$ are the quantities of that term bound. For the quadratic form,
the two kernel bounds of the walk-expansion theorem bound $|\mathsf{A}_\omega(b,b')|$, for every
pair of bonds $(b,b')$, by $a_\omega e^{-\delta_{\mathsf{Q}}\tau}$: an amplitude $a_\omega$, which
depends on $\omega$ only, times an exponential decay factor, with $\delta_{\mathsf{Q}}$ and $\tau$
the rate and the quantity of those kernel bounds. Thus the majorants are numbers
attached to the walks, independent of the datum at which the terms are evaluated.

Their sums over the walks with $[\omega]\subset W'$ are sub-sums of sums over the index set of
$\mathcal{S}_0$:
\begin{itemize}
\item for $K\neq\mathsf{A}$, the sums of these majorants are sub-sums of the two-ended sums that
accompany the term bound of the walk-expansion theorem;
\item for $\mathsf{A}$, the sums of these majorants are sub-sums of the sums behind the bound
$\lvert\lvert\lvert\mathsf{A}\rvert\rvert\rvert_{\delta_{\mathsf{A}}}\le\mathsf{B}$ of
Definition~\ref{4.4:def-form-walk-hypothesis}.
\end{itemize}
Hence the series \eqref{4.4:eq-interior-sufficiency-1} converges absolutely and in $\ell^2$.
Since all traces on $W$ are trivial, it also converges in the norms of \cite{B-CMP116}. This follows
from item (f4) of the one-family hypothesis together with the first corollary to the construction of
the walk family that supplies it. These conclusions hold whatever the local datum is. What is taken
from $\mathcal{S}_0$ is its index set and the weights attached to it, not a point of
$\mathcal{S}_0$.

It remains to treat the members of $\mathbb{L}_2$, which involve the cell matrices $T_q$ and the
resolvents $(x+T_q)^{-1}$ that enter the local form of the cell-resolvent expansion
(Lemma~\ref{4.4:prf-local-resolvent}), where $T_q$ and the parameter $x$ are fixed. Here $q$ runs
over the set $\mathcal{Q}$ of cells, with the enlargements $q^{*}$ and $q^{\sharp}$ of
Lemma~\ref{4.4:lem-cell-coercivity}. Item (v) of
the cell-localisation lemma states the following: if a walk $\omega$ with $d(\omega)\le r'$ has a
row or a column in $q^{*}$, then
\begin{equation*}
|\omega|\subset q^{\sharp}.
\end{equation*}
Let $q^{\sharp}\subset W'$. The matrix $T_q$ is built from the terms $\mathsf{A}_\omega$ of the
walks with $d(\omega)\le r'$ that have a row or a column in $q^{*}$, all of which lie in
$q^{\sharp}$ by this inclusion, evaluated at the restriction $\mathsf{u}|_{q^{\sharp}}$, a local
datum in $\mathbf{P}_{\rho,q^{\sharp}}$. The near-part coercivity on one cell
(Lemma~\ref{4.4:lem-cell-coercivity}) bounds from below by $\tfrac78\gamma_*|v|^2$ the real part of
the quadratic forms built from the terms $\mathsf{A}_\omega$ with $d(\omega)\le r'$ and
$[\omega]\subset q^{\sharp}$ at such a local datum; by that lemma, the resolvent $(x+T_q)^{-1}$
exists.
\end{proof}

\begin{remark}
We note which inputs the two halves of the independence of the localised terms from the small-field
region use.

Part (d) of Theorem~\ref{4.4:thm-uniform-agreement}, read in the geometry $\mathcal{S}$, is used by
no statement that depends on that independence theorem. Those statements compare interiors, or two
whole-lattice settings. That the terms of $\mathcal{S}$ attached to the large-field regions $X$
sum to the integral is the half of item (i-a) that concerns the
operation $\mathbf{T}$.

The part of the independence theorem used for clause (0) of Theorem~0 consists of parts (a)--(c) of
Theorem~\ref{4.4:thm-uniform-agreement} and the first four corollaries. It uses no admissibility of
$\mathcal{S}$. It depends on:
\begin{itemize}
\item the walk-expansion hypothesis, the one-family hypothesis and the positivity hypothesis at
$\mathcal{S}_0$;
\item a further hypothesis on the walk expansion, in restricted form;
\item the local form, at thin regions $X$, of the condition \cite[(1.12)]{B-CMP109}, together
with the convention that fixes which domain class is used there.
\end{itemize}
The remaining half of the independence theorem is the one resting on a lemma of the large-field
section of this chapter; we call it the large-field half.
The one-power hereditary class enters in two places. It enters through the large-field half. It
also enters through the vertex-functional theorem at equal levels, for the
second and third parts of the first corollary; these concern the potentials of the vertex
functionals on the subregions $Y'\subset X$.

The large-field half at the levels $n\le k$ and item (i-a) depend on further inputs:
\begin{itemize}
\item the admissibility of $\mathcal{S}$;
\item the operator bounds at free current, in particular their operator part in $\mathcal{S}$;
\item a corollary and a further statement, both of the large-field section of this chapter.
\end{itemize}
\end{remark}

\begin{corollary}[Independence of the localised terms from the small-field region]
\label{4.4:cor-localised-independence}
Let $\mathcal{S}$ denote the setting in which the $(k+1)$-st renormalisation step is performed
with small-field region
$\Lambda_{k+1}$, let $\mathcal{S}_0$ denote the whole-lattice setting, and let
$\mathcal{R}$ be the walk-expansion rule, with $\mathbb{L}$ the family of operators it expands and
$K^{W'}$ the $W'$-indexed term of $K\in\mathbb{L}$. In this corollary and its proof the letters
$\mathcal{R}$, $\rho$, $s$ and $t$ have the following meanings: $\mathcal{R}$ is the rule just
named, $\rho$ is the index of the local class $\mathbf{P}_{\rho,X}$ of
Definition~\ref{4.4:def-local-data-classes}, $s$ denotes the interpolation parameters of
$\mathcal{R}$, and $t$ is the remaining argument of the terms of \cite[(3.47)]{B-CMP119}.
Assume the following.
\begin{enumerate}
\item[(i)] For every closed union $W$ of $M_1$-cubes, every $W$-indexed quantity of
\cite[\S\S1--2]{B-CMP116} and of
\cite[(3.44)--(3.47)]{B-CMP119} is a finite sum of products of walk factors all of whose steps
have their domains (the grains of the walk) contained in $W$,
local densities on $W$, characteristic functions of bonds of $W$, interpolation integrals, and
universal coefficients.
\item[(ii)] For $W\subset\Lambda_{k+1}$ the local densities, the characteristic functions and the
interpolation integrals are the same in $\mathcal{S}$ and in $\mathcal{S}_0$.
\item[(iii)] The terms $\mathbf{E}_0^{(k+1)}(\Lambda_{k+1},X,b)$ are finite sums of the products
in (i), and a product all of whose ingredients agree in $\mathcal{S}$ and $\mathcal{S}_0$ agrees in
$\mathcal{S}$ and $\mathcal{S}_0$.
\item[(iv)] The conclusions of Theorem~\ref{4.4:thm-uniform-agreement} (locality of the rule and
agreement on uniform regions), of Lemma~\ref{4.4:lem-interior-sufficiency} (interior sufficiency
of the whole-lattice setting) and of Lemma~\ref{4.4:lem-vertex-formula} (the vertex formula)
hold, and conditions (f1) and (f4) of the one-family hypothesis, as it appears among the
hypotheses of Lemma~\ref{4.4:lem-interior-sufficiency}, hold for the operator families expanded
by $\mathcal{R}$.
\item[(v)] For a representative of the one-power hereditary class at $(k+1,X;\cdot)$ and every
$Y'\in\mathbf{D}_k$ with $Y'\subset X$, the restriction to $Y'$ is a representative at
$(k+1,Y';\cdot)$ (the restriction statement for the one-power hereditary class, in the case
where its two levels coincide).
\end{enumerate}
Then for every $X\in\mathbf{D}_{k+1}$ with $X\subset\Lambda_{k+1}$ and every bond $b$ of $X$
(respectively every point $z$ of $X$, for the terms of \cite[(3.47)]{B-CMP119} indexed by points),
the term $\mathbf{E}_0^{(k+1)}(\Lambda_{k+1},X,b)$ (respectively
$\mathbf{E}_0^{(k+1)}(\Lambda_{k+1},X,z)$) of \cite[(3.47)]{B-CMP119}, for
each of the three parts of \cite[(3.37)]{B-CMP119}, computed in $\mathcal{S}$ under
$\mathcal{R}$, equals the term with the same $(X,b)$ (respectively $(X,z)$) computed in
$\mathcal{S}_0$, as a function of $((\mathbf{U},\mathbf{J})|_X;t)$, with $t$ the remaining
argument of the terms in \cite[(3.47)]{B-CMP119}. The equality holds
\begin{itemize}
\item for the first part, on the local class $\mathbf{P}_{\rho,X}$
(Definition~\ref{4.4:def-local-data-classes});
\item for the second and third parts, on the representatives of the one-power hereditary class
$U^{\Sigma^1}_{k+1}(X;(1+\beta)\alpha_0,(1+\beta)\alpha_1,\alpha_0)$, the
class on which these parts are evaluated.
\end{itemize}
These representatives lie in the local class $\mathbf{P}_{1+\beta,X}$
(Definition~\ref{4.4:def-local-data-classes}), where the vertex functionals formed from these
terms are
defined (Lemma~\ref{4.4:lem-local-coercivity}, Proposition~\ref{4.4:prop-local-real-base} and a
standing hypothesis of this chapter).
\end{corollary}

\cite[p.~278]{B-CMP119} states that the terms $\mathbf{E}_0^{(k+1)}(\Lambda_{k+1},X,z)$ with
$X\subset\Lambda_{k+1}$ do not depend on $\Lambda_{k+1}$ and coincide with the corresponding terms
arising from the expressions defined on the whole lattice; the corollary gives this statement for
the three parts, on the domains displayed above.

Hypotheses (i), (ii) and (iii) are, respectively, the structure step, the agreement step for the
local ingredients, and the assembly step together with its assembly proposition, in the proof of
the representation lemma for the localised terms.

\begin{proof}
By (i) each $W$-indexed quantity is a finite sum of products of five kinds of ingredient. By (ii)
the local densities, the characteristic functions and the interpolation integrals agree in
$\mathcal{S}$ and $\mathcal{S}_0$ when $W\subset\Lambda_{k+1}$. The coefficients are universal and
are therefore the same in both settings. By (iii) it therefore suffices to show that the walk
factors agree.

\emph{The walk factors.} Every product in (i) is built from operators of $\mathbb{L}$ through the
rule $\mathcal{R}$. By the vertex formula (Lemma~\ref{4.4:lem-vertex-formula}), applied to the
interpolation parameters $s$ of $\mathcal{R}$, the value at full interpolation is the sum over
$\sigma$ of the terms $f_\sigma$ of Lemma~\ref{4.4:lem-vertex-formula}. Each $f_\sigma$ is a
finite alternating
sum of values at vertices $1_{\sigma'}$, $\sigma'\subset\sigma$. By the same lemma, a $W$-indexed
quantity contains the operators of $\mathbb{L}$ only through the terms $K^{W'}$ with
$W'\subset W$.

Now let $W\subset\Lambda_{k+1}$. By part (c) of Theorem~\ref{4.4:thm-uniform-agreement},
$K^{W'}[\mathcal{S}]=K^{W'}[\mathcal{S}_0]$ term by term for every $K\in\mathbb{L}$ and every
$W'\subset W$ as above. Interior
sufficiency (Lemma~\ref{4.4:lem-interior-sufficiency}) makes this equality one of functions of the
local datum alone. So the walk factors agree, and by (iii) every $W$-indexed quantity with
$W\subset\Lambda_{k+1}$ agrees in $\mathcal{S}$ and $\mathcal{S}_0$. We now apply this to the
three parts of \cite[(3.37)]{B-CMP119}, computed under $\mathcal{R}$.

\emph{First part.} For $a>0$ and $\lambda_0>0$,
\begin{equation*}
\frac{1}{\lambda+\lambda_0}-\frac{1}{\lambda+a}
=\frac{a-\lambda_0}{(\lambda+\lambda_0)(\lambda+a)} .
\end{equation*}
Integrating over $\lambda\in[0,\infty)$ gives
\begin{equation}\label{4.4:eq-localised-independence-1}
\log a-\log\lambda_0
=\int_0^\infty (a-\lambda_0)(\lambda+\lambda_0)^{-1}(\lambda+a)^{-1}\,d\lambda .
\end{equation}
We apply \eqref{4.4:eq-localised-independence-1} with $\mathsf{A}^Y$ in place of $a$; here
$\mathsf{A}^Y$ is the operator $\mathsf{A}$ of the first part compressed to the set of bonds $Y$
over which the determinant in the first part of \cite[(3.37)]{B-CMP119} is taken,
as in the cell-resolvent expansion (Lemma~\ref{4.4:lem-cell-resolvent}), and $\lambda_0>0$ is
fixed. We take the determinant term of the first part of \cite[(3.37)]{B-CMP119} with the sign
fixed in the determinant-normalisation statement, under which it is the term
$-\tfrac12\log\det\mathsf{A}^Y$, and write
$\mathsf{R}^Y(\lambda)=(\lambda+\mathsf{A}^Y)^{-1}$. This gives
\begin{equation}\label{4.4:eq-localised-independence-2}
-\tfrac12\log\det\mathsf{A}^Y
=\mathrm{const}-\tfrac12\int_0^\infty d\lambda\,(\lambda+\lambda_0)^{-1}
\operatorname{Tr}\bigl[(\mathsf{A}^Y-\lambda_0)\mathsf{R}^Y(\lambda)\bigr],
\end{equation}
with a constant independent of the configuration.

Write $\mathsf{A}=\sum_\omega\mathsf{A}_\omega$ for the walk expansion of $\mathsf{A}$ under
$\mathcal{R}$. By the cell-resolvent expansion (Lemma~\ref{4.4:lem-cell-resolvent}(i), with
$x=\lambda\ge 0$) we also have
$\mathsf{R}^Y(\lambda)=\sum_\varpi\mathsf{R}^Y_\varpi(\lambda)$. The $(b,b)$-entry of the
integrand in \eqref{4.4:eq-localised-independence-2} is therefore the sum of two terms:
\begin{itemize}
\item the sum over pairs $(\omega,\varpi)$ of
$\operatorname{tr}[\mathsf{A}_\omega\mathsf{R}^Y_\varpi(\lambda)](b,b)$;
\item the sum over walks $\varpi$ of $-\lambda_0\operatorname{tr}\mathsf{R}^Y_\varpi(\lambda)(b,b)$.
\end{itemize}
The $X$-term collects the pairs with
$\operatorname{hull}_{k+1}(|\omega|\cup|\varpi|)=X$, and likewise the walks with
$\operatorname{hull}_{k+1}(|\varpi|)=X$. Here $\operatorname{hull}_{k+1}$ assigns to a set of bonds
its localisation domain at level $k+1$.

The $\lambda$-integral converges absolutely in the walk norm $|\!|\!|\cdot|\!|\!|_\theta$ of
Proposition~\ref{4.4:prop-rule-properties}, because
\begin{equation*}
\int_0^\infty(\lambda+\lambda_0)^{-1}\bigl(\lambda+\tfrac12\gamma_*\bigr)^{-1}d\lambda<\infty ,
\end{equation*}
where $\gamma_*>0$ is the constant in the coercivity hypothesis of
Lemma~\ref{4.4:lem-cell-resolvent}.
For $X\subset\Lambda_{k+1}$, every factor $\mathsf{A}_\omega$ and $\mathsf{R}^Y_\varpi(\lambda)$
that enters the $X$-term is indexed by a set contained in $X\subset\Lambda_{k+1}$. By
part (c) of Theorem~\ref{4.4:thm-uniform-agreement} each such factor is therefore the same in
$\mathcal{S}$ and $\mathcal{S}_0$. Hence the $X$-term of the first part agrees, on
$\mathbf{P}_{\rho,X}$.

\emph{Second and third parts.} These are given by \cite[(2.1)--(2.13)]{B-CMP116} and
\cite[(3.44)--(3.47)]{B-CMP119}, with operator arguments $\mathsf{A}(s)$, $\mathsf{M}^Y(s)$,
$C^{(k)}(Z_0,s)$: the three operator families of Proposition~\ref{4.4:prop-rule-properties}(c2),
with $\mathsf{M}^Y$ compressed to $Y$ and $C^{(k)}(Z_0)$ the family attached to the region $Z_0$,
at the interpolation parameters $s$ of $\mathcal{R}$. Inside the potential $V_k(Y',\cdot)$ of
\cite[(1.34)]{B-CMP116} they also involve the operators of the sub-family $\mathbb{L}_1$ of
$\mathbb{L}$ that enter
this potential. Conditions (f1) and (f4) of the one-family hypothesis, as it appears among the
hypotheses of Lemma~\ref{4.4:lem-interior-sufficiency}, provide the terms $K^{W'}$
of these operators at free $\mathbf{J}$ in the norms of \cite{B-CMP116}. By the
vertex formula, every $X$-indexed quantity of these parts contains all of these operators only
through terms $K^{W'}$ with $W'\subset X\subset\Lambda_{k+1}$. By
part (c) of Theorem~\ref{4.4:thm-uniform-agreement}, these terms agree in $\mathcal{S}$ and
$\mathcal{S}_0$.

The quantities are evaluated on the representatives of the one-power hereditary class named in
the statement. These lie in $\mathbf{P}_{1+\beta,X}$ by the domain inclusion for this class, and
there the vertex functionals are defined: the polydisc positivity they require holds on
$\mathbf{P}_{1+\beta,X}$ by Lemma~\ref{4.4:lem-local-coercivity}, and the property of
Proposition~\ref{4.4:prop-local-real-base} holds there by that proposition together with a
standing hypothesis of this chapter. By (v), the restriction of each such
representative to every
$Y'\in\mathbf{D}_k$ with $Y'\subset X$ is a representative at
$(k+1,Y';\cdot)$, that is, an admissible argument of the potential $V_k(Y',\cdot)$ of
\cite[(1.34)]{B-CMP116} on $Y'$. Hence the second and third
parts agree on this class, which completes the proof.
\end{proof}

The terms of Ba{\l}aban's expansion that are localised inside the small-field region $\Lambda_{k+1}$ are compared in two constructions. These are fixed as follows.

We write $\mathcal{S}_0$ for the construction of \cite{B-CMP119} on the whole lattice. In it the analyticity radii are those of \cite[(2.28)]{B-CMP119}, proportional to $g$. A second function of the coupling, the function $p_1(\cdot)$ of \cite{B-CMP119}, distinct from the degree function $p_0(\cdot)$, enters evaluated at the running coupling, as $p_1(g_k)$. The large-field operations $\mathbf{R}^{(j)}$ with $j\le k_1$ are merged as in \cite{B-CMP119}; here $k_1$ is the level, so denoted in \cite{B-CMP119}, up to which these operations are merged.

We write $\mathcal{S}$ for the construction of \cite{B-CMP119} in which the assertion quoted in part (b) of the corollary below is made, with small-field region $\Lambda_{k+1}$.

We write $\mathbf{D}_k$ for the set of localisation domains at level $k$, $\mathbf{U}$ for the complexified bond configuration, $\mathbf{J}$ for the independent current variable, and $A$ for the fluctuation field of the step, as in \cite{B-CMP119}; here $A$ is not the action $A(U)$. We write $V_k(Y',\cdot)$ for the term with localisation domain $Y'$ of the expansion \cite[(1.41)]{B-CMP116}; this is not the block field $V_k$ of Ba{\l}aban's step.

\begin{corollary}[Agreement of the localised potentials inside the small-field region]
\label{4.4:cor-potential-agreement}
\leavevmode
\begin{itemize}
\item[(a)] Let $Y'\in\mathbf{D}_k$ with $Y'\subset\Lambda_{k+1}$. Consider the term $V_k(Y',\cdot)$ of the expansion \cite[(6.41)]{B-CMP109}, which is the expansion \cite[(1.41)]{B-CMP116}. Computed in $\mathcal{S}$, it is the same function of the restriction $(\mathbf{U},\mathbf{J},A)|_{Y'}$ as the term computed in $\mathcal{S}_0$.
\item[(b)] Consequently, consider the assertion of \cite{B-CMP119}, in which $[\mathrm{I}]$ denotes \cite{B-CMP109} and ``those lemmas'' are \cite[Lemmas~1 and~2]{B-CMP116}, that ``the terms of the expansion (I.6.41), with localization domains $Y$ contained in $\Lambda_{k+1}$, satisfy exactly all the conditions of those lemmas''. In $\mathcal{S}$ this assertion holds as soon as it holds in $\mathcal{S}_0$. For these terms, therefore, the assertion is used here as an assertion about the whole-lattice construction $\mathcal{S}_0$ of \cite{B-CMP119}, and not as one about the construction of \cite{B-CMP109}. The remaining passage from $\mathcal{S}_0$ to the hypotheses of \cite[Lemmas~1 and~2]{B-CMP116} is provided by the dictionary between the whole-lattice construction $\mathcal{S}_0$ and those hypotheses. Hence the assertion holds in $\mathcal{S}$ for all terms with localisation domains contained in $\Lambda_{k+1}$.
\end{itemize}
\end{corollary}

\begin{proof}
\emph{Part (a).} Fix $Y'\in\mathbf{D}_k$ with $Y'\subset\Lambda_{k+1}$. The computation of $V_k(Y',\cdot)$ in $\mathcal{S}$ uses exactly three inputs:
\begin{itemize}
\item the first locality statement of this section;
\item the averaged-configuration lemma;
\item clause $(\beta)$ of the second step of the interior construction.
\end{itemize}
Each of these three inputs is stated for the data $(\mathbf{U},\mathbf{J},A)$ restricted to a localisation domain $Y'$ contained in $\Lambda_{k+1}$, and in that form it is the same statement in $\mathcal{S}$ and in $\mathcal{S}_0$. Each therefore produces, in both constructions, the same output from the same restricted data $(\mathbf{U},\mathbf{J},A)|_{Y'}$.
The series that occur in these three inputs are evaluated at free current $\mathbf{J}$. They converge there by the first and fourth clauses of the one-family convergence statement.

No other ingredient of $\mathcal{S}$ enters. Hence the function of $(\mathbf{U},\mathbf{J},A)|_{Y'}$ that these inputs produce in $\mathcal{S}$ is the function they produce in $\mathcal{S}_0$. This is assertion (a).

\emph{Part (b).} Let $Y\in\mathbf{D}_k$ with $Y\subset\Lambda_{k+1}$. By part (a), applied with $Y'=Y$, the term $V_k(Y,\cdot)$ of $\mathcal{S}$ and the term $V_k(Y,\cdot)$ of $\mathcal{S}_0$ are one and the same function of the same variables $(\mathbf{U},\mathbf{J},A)|_{Y}$.

A condition of \cite[Lemmas~1 and~2]{B-CMP116} imposed on such a term is a condition on this function, measured in the norms of \cite{B-CMP116}. Therefore, once these norms are fixed, the term satisfies it in $\mathcal{S}$ if and only if it satisfies it in $\mathcal{S}_0$. Applied to every such $Y$, this shows the following. The quoted assertion of \cite{B-CMP119}, restricted to localisation domains contained in $\Lambda_{k+1}$, holds in $\mathcal{S}$ as soon as it holds in $\mathcal{S}_0$. Thus, for these terms, $\mathcal{S}$ is reduced to $\mathcal{S}_0$, the whole-lattice construction of \cite{B-CMP119}. In $\mathcal{S}_0$ the hypotheses of \cite[Lemmas~1 and~2]{B-CMP116} are reached through the dictionary between $\mathcal{S}_0$ and those hypotheses. Hence the assertion holds in $\mathcal{S}$ for every term whose localisation domain is contained in $\Lambda_{k+1}$.

It remains to say how the norms of \cite{B-CMP116} are read in $\mathcal{S}_0$. The constants $K^{W'}$ and $K(s)$ of the norms of \cite{B-CMP116} are, in $\mathcal{S}_0$, the constants supplied by the two corollaries on the constants of the weighted norms. The factor $e^{5^d\kappa_1}$ stands in place of the factor $e^{16\kappa_1}$ of \cite{B-CMP116}.

Accordingly, \cite[Lemmas~1 and~2]{B-CMP116} are read on the family of backgrounds prescribed by the first reading rule for these lemmas. For the terms whose localisation domains meet both $\Lambda_{k+1}$ and its complement, the constants are those of the transfer of the analysis of \cite[\S3]{B-CMP109} to the cover, that is, of the sentence
\begin{quote}
We repeat the analysis of Sect.~3 [I] using the above cover $\{\square\}$, and the partition of unity $\{\zeta_\square\}$.
\end{quote}
printed without proof in \cite[p.~276]{B-CMP119}; cited as printed. In this sentence $[\mathrm{I}]$ denotes \cite{B-CMP109}.

Finally, these straddling terms are defined on the spaces \cite[(3.43)]{B-CMP119}. They are not treated in this corollary and fall under the transfer quoted above.
\end{proof}

\begin{definition}[Ba{\l}aban's layered analyticity spaces and their gauge condition
{\cite[p.~261, (2.34)--(2.39)]{B-CMP119}}]
\label{4.4:def-layered-spaces}
The symbols $\xi$, $G$, $\Omega_n$, $\mathbf{D}_j$, $\mathbf{B}_p(X)$, $\Gamma_i$,
$M^{\cdot}(\cdot)$, $\alpha_{0,n}$, $\alpha_{1,n}$, $M$ and $C$ below, the configurations
$U(\cdot,\cdot)$ and $\mathbf{J}_{p,X}(\cdot)$ determined by averaged data, the index $p$ of
the sets $\mathbf{B}_p(X)$ (which is not a plaquette), and the parameters $\tilde\alpha_0$,
$\tilde\alpha_1$ have the meaning they have in \cite[\S2]{B-CMP119}, and $\beta$ is the
constant fixed after the definition. Let $X\in\mathbf{D}_j$. The space
$\tilde{U}^c_j(X,\tilde\alpha_0,\tilde\alpha_1)$ is the set of configurations
$(\mathbf{U},\mathbf{J})$ defined on $X$ that satisfy the following three conditions.
\begin{enumerate}
\item[(i)] $\mathbf{U}=U'U$, where $U$ takes values in the group $G$. Put
\begin{equation*}
U_{p,X}(M^{\cdot}(\mathbf{U}))
=U\bigl(\mathbf{B}_p(X)\cup\{\Gamma_i\}_{i<p},\,M^{\cdot}(\mathbf{U})\bigr).
\end{equation*}
The configurations $U$, $\mathbf{U}$, $U_{p,X}(M^{\cdot}(\mathbf{U}))$, $\mathbf{J}$ and
$\mathbf{J}_{p,X}(M^{\cdot}(\mathbf{U}))$ satisfy the bounds
\begin{align}
|\partial U-1|,\ |\partial\mathbf{U}-1|
&<\bigl(1-\beta(1-2^{-(j-n)})\bigr)\alpha_{0,n}\,\xi^2(L^n\xi)^{-2},
\label{4.4:eq-layered-spaces-a}\\
|\partial U_{p,X}(M^{\cdot}(\mathbf{U}))-1|
&<\bigl(1-\beta(1-2^{-(j-n)})\bigr)\alpha_{0,n}\,L^{-2p}(L^nL^{-p})^{-2},
\label{4.4:eq-layered-spaces-b}\\
|\mathbf{J}|
&<\bigl(1-\beta(1-2^{-(j-n)})\bigr)\alpha_{0,n}\,(L^n\xi)^{-3},
\label{4.4:eq-layered-spaces-c}\\
|\mathbf{J}_{p,X}(M^{\cdot}(\mathbf{U}))|
&<\bigl(1-\beta(1-2^{-(j-n)})\bigr)\alpha_{0,n}\,L^{n-\min(p,n)}(L^nL^{-p})^{-3},
\label{4.4:eq-layered-spaces-d}
\end{align}
on $X\cap(\Omega_n\setminus\Omega_{n+1})$ for $n=1,\dots,j-1$, and on $X\cap\Omega_j$ for
$n=j$.
\item[(ii)] Consider, for $n=1,\dots,j-1$, each cube $\square$ of size $CML^n\xi$ with
$\square\subset\Omega_n\setminus\Omega_{n+2}$ and $\square\cap\Omega^c_{n+1}\neq\emptyset$.
Consider also, for $n=j$, each cube $\square\subset\Omega_j$ of size $CM$. For every such cube
there is a gauge transformation $u$, defined on $\square\cap X$, such that $U^u=\exp i\xi A$ and
\begin{equation}
\label{4.4:eq-layered-spaces-e}
L^n\xi\,|A|,\ (L^n\xi)^2\,|\nabla^{\xi}A| < BCM\alpha_{0,n}
\quad\text{on } \square\cap X,
\end{equation}
where $B$ is an absolute constant.
\item[(iii)] $U'=\exp i\xi A'$, where $A'$ takes values in the algebra $\mathfrak{g}^c$, and
\begin{equation}
\label{4.4:eq-layered-spaces-f}
L^n\xi\,|A'|,\ (L^n\xi)^2\,|\nabla^{\xi}_U A'|
<\bigl(1-\beta(1-2^{-(j-n)})\bigr)\alpha_{1,n}
\end{equation}
on $X\cap(\Omega_n\setminus\Omega_{n+1})$ for $n=1,\dots,j-1$, and on $X\cap\Omega_j$ for
$n=j$.
\end{enumerate}
\end{definition}

\begin{remark}
Throughout, $\beta=1/4$. In \cite{B-CMP119} the number $\beta$ in the factors
$1-\beta(1-2^{-(j-n)})$ is introduced with the words ``The number $\beta$ is a small positive
constant, but not too small, e.g., we can take $\beta = 1/4$''. The same constant $\beta$ is
meant in the corresponding factors of the spaces of \cite{B-CMP109}.

The definition is used in the following precise form.

\begin{enumerate}
\item[(a)] The radius in each bound of condition (i) is $\alpha_{0,n}$. This includes the two
bounds on the current, \eqref{4.4:eq-layered-spaces-c} and \eqref{4.4:eq-layered-spaces-d}.
The radius $\alpha_{1,n}$ enters only through \eqref{4.4:eq-layered-spaces-f}, that is, only
through the factor $U'$.
\item[(b)] Condition (ii) constrains only the $G$-valued factor $U$. The bound
\eqref{4.4:eq-layered-spaces-e} has the constant $BCM\alpha_{0,n}$ and carries no factor
$1-\beta(1-2^{-(j-n)})$.
\item[(c)] Condition (i) bounds the configurations $U_{p,X}(M^{\cdot}(\mathbf{U}))$ and
$\mathbf{J}_{p,X}(M^{\cdot}(\mathbf{U}))$, which are determined by $M^{\cdot}(\mathbf{U})$.
It contains no bound for $U_{p,X}(M^{\cdot}(U))$ or for $\mathbf{J}_{p,X}(M^{\cdot}(U))$.
By contrast, the corresponding condition \cite[(1.16)]{B-CMP109} ends with the clause that
the same bounds hold for $U$ instead of $\mathbf{U}$.
\item[(d)] The cubes of condition (ii) have size $CML^n\xi$ for $n=1,\dots,j-1$, and size
$CM$ inside $\Omega_j$. The corresponding condition \cite[(1.12)]{B-CMP109} is imposed instead
on each cube $\square\subset X$ of a size $O(1)LM$.
\item[(e)] Condition (ii) is required for every cube of the lattice that has the stated size
and satisfies the stated position conditions. This holds whether or not the cube is aligned
with a partition. In that condition $u$ is $G$-valued and $A$ takes values in $\mathfrak{g}$,
as in \cite{B-CMP122a}. The constant $C$ is the one of the same condition in
\cite[(1.69)]{B-CMP122a}, where it is described as
\begin{quote}
a positive integer. It is usually small, because of geometric constraints on the
cubes, for example we can assume that $C \leqq 2$.
\end{quote}
We take $C\in\{1,2\}$; in particular $C=1$ is admitted.
\end{enumerate}
\end{remark}

\begin{definition}[The post-step space, the fluctuation polydisc and the twist size]
\label{4.4:def-post-step-space}
Let $\beta$ be the schedule constant of the layered analyticity spaces
(Definition~\ref{4.4:def-layered-spaces}), let $\xi$ be the lattice spacing of that definition,
and put, for integers $n\le j$,
\begin{equation}\label{4.4:eq-post-step-space-1}
\ell_n := L^n\xi, \qquad f_{j,n} := 1-\beta\bigl(1-2^{-(j-n)}\bigr).
\end{equation}
Since $\beta>0$ and $2^{-(j-n)}$ decreases as $j$ increases, $f_{j,n}$ is decreasing in $j$.
Moreover $f_{j,j}=1$. The number $f_{j,n}$ is the factor that multiplies the radius
$\alpha_{0,n}$ or $\alpha_{1,n}$ in the bounds of Definition~\ref{4.4:def-layered-spaces}
at level $j$ on layer $n$; see, for instance, \eqref{4.4:eq-layered-spaces-a} and
\eqref{4.4:eq-layered-spaces-f}.

\smallskip
\emph{The post-step space.} Fix a number $\theta_S\in\,]0,1[$. For a region $Y$, the
\emph{post-step space} $\mathfrak{U}^{\mathrm{post}}_{k+1}(Y)$ is the space of
Definition~\ref{4.4:def-layered-spaces} at level $k+1$ on $Y$, with one change. In the bounds
\eqref{4.4:eq-layered-spaces-a}--\eqref{4.4:eq-layered-spaces-d} and
\eqref{4.4:eq-layered-spaces-f}, the schedule factor $f_{k+1,n}$ in front of the radius is
replaced by
\begin{equation}\label{4.4:eq-post-step-space-2}
F_{\mathrm{top}} := 1+\beta \quad\text{on } Y\cap\Omega_{k+1},
\end{equation}
and, for $n=1,\dots,k$, by
\begin{equation}\label{4.4:eq-post-step-space-3}
F_n := f_{k+1,n}+\theta_S\,\beta\,2^{-(k+1-n)}
\quad\text{on } Y\cap(\Omega_n\setminus\Omega_{n+1}).
\end{equation}
The gauge condition of Definition~\ref{4.4:def-layered-spaces}, which carries no schedule
factor, is unchanged.
In this space, the bounds \eqref{4.4:eq-layered-spaces-b} and \eqref{4.4:eq-layered-spaces-d}
are read with $X=Y$. They are read for the configurations
\begin{equation}\label{4.4:eq-post-step-space-a}
U\bigl(\mathbf{B}_p(X'')\cup\{\Gamma_i\}_{i<p},\,M^{\cdot}(\mathbf{U})\bigr),
\qquad X''\subset Y \text{ admissible},
\end{equation}
and they are required for every admissible $X''\subset Y$, not only for $X''=Y$.
Here a subset $X''\subset Y$ is called admissible if it belongs to the class of subsets for
which Ba{\l}aban's construction forms the determining sets $\mathbf{B}_p(X'')$.

\smallskip
\emph{The fluctuation polydisc.} Let $C_1^{\mathrm{fl}}$ be a positive constant, the polydisc
constant. The \emph{fluctuation polydisc} is the set of fluctuation data
$(A,A_k,\mathbf{H}^{(k)})$ given by
\begin{equation}\label{4.4:eq-post-step-space-4}
\begin{split}
\mathfrak{D} := {}&\bigl\{(A,A_k) : |A|,\ |A_k| < C_1^{\mathrm{fl}}\,p_1(g_k)\bigr\}\\
&\times\bigl\{\mathbf{H}^{(k)} : |\mathbf{H}^{(k)}| < C_1^{\mathrm{fl}}\,\varepsilon_k\,e^{-R_k}\bigr\}.
\end{split}
\end{equation}
Here $R_k$ is the partition radius at level $k$, $p_1(\cdot)$ is Ba{\l}aban's degree function
$p_1$ of the coupling, and $\varepsilon_k$ is Ba{\l}aban's small-field parameter at level $k$.

\smallskip
\emph{The twist size.} The \emph{twist size} at level $k$ is
\begin{equation}\label{4.4:eq-post-step-space-5}
w_k := 2\,\|C\|\,C_1^{\mathrm{fl}}\,A_1\,g_k\,\bigl(\log g_k^{-2}\bigr)^{p_0},
\end{equation}
where $A_1$ and the exponent $p_0$ are among the auxiliary parameters $\mathfrak{p}$ of
Ba{\l}aban's construction, and $\|C\|$ is a constant of that construction. On
$\mathfrak{D}$, the data twist $\mathsf{b}$ of the level-$k$ data satisfies
$|\mathsf{b}|\le w_k$.
\end{definition}

\begin{definition}[Induced bounds for the group-valued factor]\label{4.4:def-group-factor-bounds}
Let $Y$ be a domain at level $k+1$. Let $\mathfrak{U}^{\mathrm{post}}_{k+1}(Y)$ be the post-step space
of Definition~\ref{4.4:def-post-step-space}, with the factors $F_{\mathrm{top}}=1+\beta$ on
$Y\cap\Omega_{k+1}$ and $F_n$ on layer $n\le k$. In that space the bounds
\eqref{4.4:eq-layered-spaces-b} and \eqref{4.4:eq-layered-spaces-d} of
Definition~\ref{4.4:def-layered-spaces} are read as in \eqref{4.4:eq-post-step-space-a}, that is, for
$U(\mathbf{B}_p(X'')\cup\{\Gamma_i\}_{i<p},M^{\cdot}(\mathbf{U}))$ for every admissible $X''\subset Y$.
Let $(\mathbf{U},\mathbf{J})\in\mathfrak{U}^{\mathrm{post}}_{k+1}(Y)$ and write $\mathbf{U}=U'U$ with
$U'=e^{i\xi A'}$, where $U$ takes values in the group $G$.

We say that $(\mathbf{U},\mathbf{J})$ satisfies the \emph{induced bounds for the group-valued factor}
if these bounds also hold with $U$ in place of $\mathbf{U}$:
\begin{equation}\label{4.4:eq-group-factor-bounds-a}
\begin{gathered}
\text{\eqref{4.4:eq-layered-spaces-b} and \eqref{4.4:eq-layered-spaces-d}, read as in
\eqref{4.4:eq-post-step-space-a},}\\
\text{hold on $Y$ with the factors $F_{\mathrm{top}}$, $F_n$,}\\
\text{for } U\bigl(\mathbf{B}_p(X'')\cup\{\Gamma_i\}_{i<p},M^{\cdot}(U)\bigr)
\text{ and for } \mathbf{J}_{p,X''}(M^{\cdot}(U)),
\end{gathered}
\end{equation}
for every admissible $X''\subset Y$ and every index $p$ at which \eqref{4.4:eq-post-step-space-a} reads
them. Here $\mathbf{J}_{p,X''}(M^{\cdot}(U))$ is the current built on
$U(\mathbf{B}_p(X'')\cup\{\Gamma_i\}_{i<p},M^{\cdot}(U))$. We set
\begin{equation*}
\mathfrak{U}^{\mathrm{post},U}_{k+1}(Y):=\bigl\{(\mathbf{U},\mathbf{J})\in
\mathfrak{U}^{\mathrm{post}}_{k+1}(Y):\ \eqref{4.4:eq-group-factor-bounds-a}\text{ holds}\bigr\}.
\end{equation*}
\end{definition}

\begin{remark}
Condition \eqref{4.4:eq-group-factor-bounds-a} is used in exactly two later statements of this section:
the closure of the post-step space under the induction through $\mathbf{T}$, and the conclusion on the
new terms $\mathbf{B}^{(k+1)}(Y)$. No other statement uses it. Its standing depends on the position of
$Y$.
\begin{enumerate}
\item[(a)] \emph{Real points and fibres.} At a real point $\mathbf{U}=U$ the condition coincides with
the bounds \eqref{4.4:eq-layered-spaces-b} and \eqref{4.4:eq-layered-spaces-d}, read as in
\eqref{4.4:eq-post-step-space-a}. Now consider a fibre
$\{(e^{i\xi A'}U,\mathbf{J})\}$ over a fixed $G$-valued $U$. These are the polydiscs on which Cauchy
estimates vary $A'$, $\mathbf{J}$ and the fluctuation $A$, and never $U$. On such a fibre the condition
involves the base $U$ alone. It holds on the whole fibre as soon as the real point over $U$ lies in
$\mathfrak{U}^{\mathrm{post}}_{k+1}(Y)$.
\item[(b)] \emph{Interior domains $Y\subset\Lambda_{k+1}$.} Here the space of
\cite[(1.34)]{B-CMP116} is
\begin{equation*}
U^c_{k+1}\bigl(Y,(1+\beta)\alpha_0,(1+\beta)\alpha_1,\alpha_0\bigr)\times
\bigl\{B:\ |B|<\varepsilon_1 g_k^{-1}\text{ on }Y\bigr\},
\end{equation*}
where $B$ and $\varepsilon_1$ are the field variable and the constant so denoted in
\cite[(1.34)]{B-CMP116}.
Its first factor is the space of \cite{B-CMP109}. Condition (iv) of that space, the last clause of
\cite[(1.16)]{B-CMP109}, contains the requirement by definition. This requirement is read on the
one-power hereditary class $U^{\Sigma^1}_j(X;a_0,a_1,\gamma_0)$. There the restriction theorem for
that class carries it, for $\mathbf{U}$ and for $U$ alike, to the tower of every localisation domain of
every level $p\le j$ inside every sub-domain at equal or lower level. The corresponding conclusion at
interior domains is stated later in this section. It is derived from that restriction theorem and from
the statement that the pair produced by the operator bounds at free current lies in the one-power
hereditary class.
\item[(c)] \emph{Straddling domains}, those $Y$ that are not interior in the sense of item (b),
treated in \cite{B-CMP119}. The layered space
$\tilde{U}^c_j(X,\tilde\alpha_0,\tilde\alpha_1)$ of Definition~\ref{4.4:def-layered-spaces} contains
no such clause. There \eqref{4.4:eq-group-factor-bounds-a} is a condition imposed on the space on
which the new terms $\mathbf{B}^{(k+1)}(Y)$ are asserted to be analytic. It is preserved by the
induction through $\mathbf{T}$, and by item (a) it holds automatically at real points. Along the chain
of large-field operations of \cite{B-CMP122a} and \cite{B-CMP122b} the composite operations leave the
group-valued factor unchanged, and the group-valued factor that enters there is that of a real
configuration, so item (a) applies; a re-blocking inside the large-field region $Z''_k$ treats that
factor in the same way. Both facts are taken from the statement on the group-valued factor along the
chain of large-field operations, which supplies the condition there. At straddling
domains only, it enters as a hypothesis of the theorem on the independence of the localised terms
from the small-field region.
\end{enumerate}

The condition is not derived from membership in $\mathfrak{U}^{\mathrm{post}}_{k+1}(Y)$, for two
reasons.
\begin{itemize}
\item The route of \cite{B-CMP109}, through \cite[Proposition~9]{B-CMP102b}, gives bounds
$2B_3\alpha'_0\xi^2$ and $2B_3\alpha'_0(L^n\xi)^2$, with $\alpha'_0$ the constant so denoted in
\cite{B-CMP109}. On layer $j$ these carry a factor $2B_3$, and no
factor $L^{-(n-j)}$ is available to absorb it. Absorbing it would require an inequality of the form
$4B_3\le L^2$, which is not assumed.
\item The perturbative comparison of $U$ with $\mathbf{U}$ uses the response bound at the data twist
$\mathsf{b}_{\mathrm{tot}}$. It costs $cC_H\alpha_{1,n}$, with $c$ a constant, to be paid from a margin
$(1-\theta)\beta 2^{-(k+1-n)}\alpha_{0,n}$. This is not uniform in $k$.
\end{itemize}

Nor can the condition be dropped. In \cite[\S3]{B-CMP109} the pair \cite[(1.15)]{B-CMP109} is
substituted into $\mathbf{E}^{(j)}$. That substitution requires \cite[(1.11)]{B-CMP109} for its
group-valued factor $U_j(M^{\cdot}(U))$. In the words of \cite{B-CMP109}:
\begin{quote}
The last condition (iv) is connected with the fact that in some important constructions in the next
sections we have to substitute the functions (1.15) in place of the variables $\mathbf{U}$,
$\mathbf{J}$. We have to make sure, then, that they have required properties
\end{quote}
At interior domains this requirement is met as described in item (b). At straddling domains it is the
imposed condition of item (c).
\end{remark}

\begin{proof}[Proof of item (a)]
At a real point the decomposition $\mathbf{U}=U'U$ has $U'=1$, so $\mathbf{U}=U$. Then
$M^{\cdot}(\mathbf{U})=M^{\cdot}(U)$, and the configuration
$U(\mathbf{B}_p(X'')\cup\{\Gamma_i\}_{i<p},M^{\cdot}(U))$ and the current
$\mathbf{J}_{p,X''}(M^{\cdot}(U))$
are exactly those at which \eqref{4.4:eq-post-step-space-a} reads \eqref{4.4:eq-layered-spaces-b} and
\eqref{4.4:eq-layered-spaces-d}, with the same factors. Hence \eqref{4.4:eq-group-factor-bounds-a}
coincides with those bounds.

On a fibre $\{(e^{i\xi A'}U,\mathbf{J})\}$ the $G$-valued factor is the fixed $U$ at every point. The
objects named in \eqref{4.4:eq-group-factor-bounds-a} are the same functions of $U$ at every point of
the fibre. So the condition does not depend on $A'$ or on $\mathbf{J}$. At the real point over $U$ it is
the bounds \eqref{4.4:eq-layered-spaces-b} and \eqref{4.4:eq-layered-spaces-d}, read as in
\eqref{4.4:eq-post-step-space-a}, which hold if that point lies in
$\mathfrak{U}^{\mathrm{post}}_{k+1}(Y)$. Therefore it holds on the whole fibre.
\end{proof}

\begin{definition}[The twisted argument of an old term]\label{4.4:def-twisted-argument}
Let $w$ denote the twist parameter of the step, on which the response and the current increment
below depend. Let $(\mathbf{U},\mathbf{J})$ be a point of Ba{\l}aban's analyticity spaces
\eqref{4.4:eq-layered-spaces-f}, $\mathbf{U}$ a complexified
bond configuration and $\mathbf{J}$ the independent current variable. An old term is a function of
$(\mathbf{U},\mathbf{J})$ only. Put
\begin{equation*}
\boldsymbol{\mathcal{H}} := \boldsymbol{\mathcal{H}}_w(\mathbf{U},\mathbf{J}),
\end{equation*}
the solution of the variational equations \cite[(179)--(180)]{B-CMP102b} with $(\mathbf{U},\mathbf{J})$
inserted in the operators. The \emph{twisted argument} of the old term is the pair
$(\mathbf{U}^{\flat},\mathbf{J}^{\flat})$ given by
\begin{equation}\label{4.4:eq-twisted-argument-1}
\mathbf{U}^{\flat} := e^{i\xi\boldsymbol{\mathcal{H}}}\,\mathbf{U},
\qquad
\mathbf{J}^{\flat} := \mathbf{J} + \mathfrak{J}^{(\mathbf{U},\mathbf{J})}(w).
\end{equation}
Here $\mathfrak{J}^{(\mathbf{U},\mathbf{J})}(w)$ is the current increment as a kernel object: the
kernel lemma constructs it as the analytic extension of
$J(e^{i\xi\boldsymbol{\mathcal{H}}}\mathbf{U})-J(\mathbf{U})$, and shows that at true pairs, in the
sense of that lemma, $\mathfrak{J}^{(\mathbf{U},\mathbf{J})}(w)=J(\mathbf{U}^{\flat})-J(\mathbf{U})$.
Consequently, if moreover $\mathbf{J}=J(\mathbf{U})$, then $\mathbf{J}^{\flat}=J(\mathbf{U}^{\flat})$:
at such pairs the twisted argument is the pair formed by $\mathbf{U}^{\flat}$ and its own current,
and the extension restricts to the old term evaluated there.

Choice of decomposition: for $\mathbf{U}=U'U$, with $U$ group-valued and $U'=\exp i\xi A'$, put
\begin{equation}\label{4.4:eq-twisted-argument-2}
U^{\flat} := U,
\qquad
U'^{\flat} := e^{i\xi\boldsymbol{\mathcal{H}}}\,U' =: \exp i\xi A'^{\flat}.
\end{equation}
\end{definition}

This is the argument at which old terms are evaluated in \cite{B-CMP119}, where it is introduced by
reference to \cite[(3.13)]{B-CMP109}. The pair \eqref{4.4:eq-twisted-argument-1} is formed directly
from $(\mathbf{U},\mathbf{J})$; it is not the minimiser of the variational problem with twisted data.

The current component of the pair carries one object, the kernel increment
$\mathfrak{J}^{(\mathbf{U},\mathbf{J})}(w)$. It is not split into members. The reason is the one
Ba{\l}aban gives in \cite[\S3]{B-CMP109}, where the argument passes from the operator $D^{*}D$ to
$\Delta-P_1$ plus lower-order terms: bounds of the required kind do not hold for the operator
$D^{*}D$. For this reason he does not formulate second-order regularity conditions, and instead
replaces the expression $J_j$ for the current by the new variable $\mathbf{J}$. Accordingly, the
increment in the second component is bounded through the kernel and not member by member.

The decomposition \eqref{4.4:eq-twisted-argument-2} is a decomposition of $\mathbf{U}^{\flat}$:
\begin{equation*}
e^{i\xi\boldsymbol{\mathcal{H}}}(U'U)=(e^{i\xi\boldsymbol{\mathcal{H}}}U')\,U,
\qquad\text{so}\qquad
\mathbf{U}^{\flat}=U'^{\flat}\,U^{\flat}.
\end{equation*}
The whole twist, real or complex, goes into the factor $U'$. This is possible because
$\mathfrak{g}\subset\mathfrak{g}^c$, and $A'^{\flat}$ is in general $\mathfrak{g}^c$-valued. The
layered analyticity spaces \eqref{4.4:eq-layered-spaces-f} are defined through some decomposition
$\mathbf{U}=U'U$ with $U$ group-valued, not through a canonical one. Since an old
term is a function of $(\mathbf{U},\mathbf{J})$ only, its value at the twisted argument does not
depend on this choice.

\begin{definition}[The response bound at free current and the layer decay margins;
{\cite[Proposition~9, (190), (172)]{B-CMP102b}}]
\label{4.4:not-response-bound}
Let $d(\cdot,\cdot)$ be the weighted distance of \cite[(2.46)]{B-CMP96}. Let $\delta_0$ be the
constant of \cite[Proposition~9]{B-CMP102b}, in whose bound \cite[(190)]{B-CMP102b} the kernel
decays at rate $\delta_0/8$. Let $c_*$, with $0<c_*\le 1$ and depending only on $L$ and on the
dimension, which is four, be the constant of the operator bounds at free current, stated later in
this chapter, in which the kernel of the current increment decays at rate $c_*\delta_0$; summed by
\cite[Lemma~2.1]{B-CMP96}, this kernel gives the rate $c_*\delta_0/2$. Put
\begin{equation}\label{4.4:eq-response-bound-1}
\delta_H:=\min\Bigl\{\frac{\delta_0}{16},\ \frac{c_*\delta_0}{2}\Bigr\}.
\end{equation}

The \emph{response bound at free current} is the following statement. Admissibility of the sequence
of regions is its standing hypothesis. Let the determining set be of the form $\{\Gamma_i\}$ or of
the form $\mathbf{B}_p(X'')\cup\{\Gamma_i\}_{i<p}$, for a region $X''$ and an index $p$ as in the
definition of the determining sets. Let the background lie in the regularity class of
that determining set, and let $\mathsf{b}$ be a data twist. Then, at every point $y$ of layer $n$,
the response $\boldsymbol{\mathcal{H}}$ to the data twist $\mathsf{b}$ and the corresponding
current increment $\mathfrak{J}$ (Definition~\ref{4.4:def-twisted-argument}) satisfy
\begin{equation}\label{4.4:eq-response-bound-a}
\begin{gathered}
\ell_n\,|\boldsymbol{\mathcal{H}}(y)|,\quad \ell_n^2\,|\nabla\boldsymbol{\mathcal{H}}(y)|,\quad
\ell_n^3\,|\mathfrak{J}(y)|
\;\le\; C_H\,\sup_{y'}\Bigl[\,|\mathsf{b}(y')|\,e^{-\delta_H d(y,y')}\Bigr],\\
C_H\ge 1 .
\end{gathered}
\end{equation}
Here the constant $C_H$ depends only on $L$ and on the dimension.
The statement includes a second case. Let $U_0$ be a real minimiser base, and let the complex data $B$
lie in the polydisc \cite[(172)]{B-CMP102b}. In this case all five members of
\cite[(190)]{B-CMP102b}, each multiplied by the power of $\ell_n$ matching its order as in
\cite[(190)]{B-CMP102b}, are bounded by the right member of \eqref{4.4:eq-response-bound-a} with
$B$ in place of $\mathsf{b}$. This
applies in particular to the second-order members of $\mathcal{H}(B)$ in \cite[(190)]{B-CMP102b}
containing the operators written there $D^*D$ and $P_1$.

The bound \eqref{4.4:eq-response-bound-a} is taken by statement. At real minimiser bases with
complex data in the polydisc \cite[(172)]{B-CMP102b}, the five members are
\cite[Proposition~9, (190)]{B-CMP102b}, integrated in $B$ and summed over $y'$ by
\cite[Lemma~2.1]{B-CMP96}; this summation halves the rate $\delta_0/8$ of \cite[(190)]{B-CMP102b}
and gives the rate $\delta_0/16$. For complexified backgrounds $(\mathbf{U},\mathbf{J})$ with
independent current variable $\mathbf{J}$, the two members of
\eqref{4.4:eq-response-bound-a} in $\boldsymbol{\mathcal{H}}$ and the member in $\mathfrak{J}$ are
the operator bounds at free current, which are stated later in this chapter.

The \emph{layer decay margin} of layer $n$ is
\begin{equation}\label{4.4:eq-response-bound-2}
D_n:=\delta_H\, d\bigl(\text{layer } n,\ \Omega_{k+1}\cup S_{k+1}\bigr).
\end{equation}
Let $\mathsf{w}_*$ denote the least width of a layer, measured in that layer's own unit. Then
\begin{equation}\label{4.4:eq-response-bound-3}
D_k=0,\qquad D_n\ \ge\ \delta_H\,\mathsf{w}_*\,(k-n).
\end{equation}
Let $R$ denote the constant of the
admissibility condition \cite[(2.1)--(2.2)]{B-CMP96}. A layer made by the operation $\mathbf{T}$
at level $j$ has width at least $2MR_j$ in its own unit, where $R_j$ is the partition radius at
level $j$, by \cite[(3.5)]{B-CMP119}. Every layer of an admissible
sequence, the thin shells made by the operation $\mathbf{R}$ included, has width greater than $RM$
in its own unit, by \cite[(2.2)]{B-CMP96}; hence $\mathsf{w}_*>RM$.
\end{definition}

\begin{lemma}[Ratio of radii along the coupling flow]\label{4.4:lem-radii-ratio}
Write $\alpha_{\cdot,i}$ for either of the radii $\alpha_{0,i},\alpha_{1,i}$ of \cite[(2.28)]{B-CMP119}
at level $i$, and $q_\cdot$, $C_\cdot$ for the corresponding exponent ($q_0$ or $q_1$) and constant
($C_0$ or $C_1'$). Let $\beta_0>0$ be given, and let $\gamma>0$, the ceiling on the couplings, satisfy
\begin{equation}\label{4.4:eq-radii-ratio-1}
\begin{gathered}
\gamma^2\le e^{-1},\qquad B^{\sharp}\gamma^2\le 1,\\
q_\cdot\, c_5\gamma^2\le \log(1+\beta_0)\,\log\gamma^{-2}\quad(q_\cdot=q_0,q_1).
\end{gathered}
\end{equation}
Let $g_i\le\gamma$ be the couplings at the levels $i$ considered, and put $x_i:=g_i^{-2}$.
Fix a pair of levels $n\le k'\le k+1$ and assume as hypothesis the two-sided slope bound of clause~(0)
of Theorem~0 at the pair $(n,k')$:
\begin{equation}\label{4.4:eq-radii-ratio-2}
-c_5\;\le\; x_n-x_{k'}\;\le\; B^{\sharp}(k'-n).
\end{equation}
Then the first line of the following display holds; if moreover $q_0,q_1\ge p_0$, $n\le k$, and
\eqref{4.4:eq-radii-ratio-2} holds at the pair $(n,k)$, the second line holds as well:
\begin{equation}\label{4.4:eq-radii-ratio-a}
\begin{gathered}
\frac{\alpha_{\cdot,k'}}{\alpha_{\cdot,n}}\le(1+\beta_0)(1+k'-n)^{1/2},\\
\frac{w_k}{\alpha_{\cdot,n}}\le\frac{2\|C\|C_1A_1}{C_\cdot}\,(1+\beta_0)(1+k-n)^{1/2},
\end{gathered}
\end{equation}
where $w_k=2\|C\|C_1A_1\,g_k(\log g_k^{-2})^{p_0}$ is the twist size of
Definition~\ref{4.4:def-post-step-space}.

The lemma is applied with $k'=k$ for the layers $n\le k$ against data at level $k$ and for $w_k$; and
with $n=j\le k$ and $k'\in\{j+1,\dots,k+1\}$, for the upper layers of the post-step space of
Definition~\ref{4.4:def-post-step-space} and in verifying the condition \cite[(2.38)]{B-CMP119}. At
$k'=k+1$ the hypothesis \eqref{4.4:eq-radii-ratio-2} is the slope bound at the pair $(n,k+1)$, which
is available as soon as $g_{k+1}$ is defined by the coupling flow, before step $k+1$ is carried out.
Its upper side at the pair $(j,j+1)$ is the inequality $x_j\le x_{j+1}+B^{\sharp}$.
\end{lemma}

\begin{proof}
By the upper side of \eqref{4.4:eq-radii-ratio-2}, and since $x_{k'}^{-1}=g_{k'}^2\le\gamma^2$,
\begin{equation*}
\frac{g_{k'}^2}{g_n^2}=\frac{x_n}{x_{k'}}
\le 1+\frac{B^{\sharp}(k'-n)}{x_{k'}}
\le 1+B^{\sharp}\gamma^2(k'-n)\le 1+k'-n,
\end{equation*}
the last step by the second condition in \eqref{4.4:eq-radii-ratio-1}. Hence
$g_{k'}/g_n\le(1+k'-n)^{1/2}$.

By the lower side of \eqref{4.4:eq-radii-ratio-2}, $x_{k'}\le x_n+c_5$, so, using
$\log(1+u)\le u$ and $x_n^{-1}=g_n^2\le\gamma^2$,
\begin{equation*}
\log x_{k'}\le \log x_n+\log\Bigl(1+\frac{c_5}{x_n}\Bigr)\le \log x_n+c_5\gamma^2 .
\end{equation*}
Both $\log x_{k'}$ and $\log x_n$ are at least $\log\gamma^{-2}>0$, because $x_i\ge\gamma^{-2}$. Dividing
by $\log x_n$ and using $1+u\le e^u$,
\begin{equation*}
\Bigl(\frac{\log x_{k'}}{\log x_n}\Bigr)^{q_\cdot}
\le\Bigl(1+\frac{c_5\gamma^2}{\log\gamma^{-2}}\Bigr)^{q_\cdot}
\le\exp\Bigl(\frac{q_\cdot c_5\gamma^2}{\log\gamma^{-2}}\Bigr)\le 1+\beta_0,
\end{equation*}
the last step by the third condition in \eqref{4.4:eq-radii-ratio-1}. In \cite{B-CMP119} the
exponents $q_0,q_1$ are integers greater than $1$ and the constants $C_0$ and $C_1'$ (the latter
written $C_1$ in \cite{B-CMP119}) are sufficiently large positive numbers; no upper bound on $q_0,q_1$
is imposed, and the third condition in \eqref{4.4:eq-radii-ratio-1} is a condition on $\gamma$,
imposed after $q_0,q_1$ are fixed.

By \cite[(2.28)]{B-CMP119} the radius $\alpha_{\cdot,i}$ is $C_\cdot\,g_i(\log g_i^{-2})^{q_\cdot}$, so
\begin{equation*}
\frac{\alpha_{\cdot,k'}}{\alpha_{\cdot,n}}
=\frac{g_{k'}}{g_n}\Bigl(\frac{\log x_{k'}}{\log x_n}\Bigr)^{q_\cdot},
\end{equation*}
and the two bounds just proved give the first line of \eqref{4.4:eq-radii-ratio-a}. For the second
line,
\begin{equation*}
\frac{w_k}{\alpha_{\cdot,n}}
=\frac{2\|C\|C_1A_1}{C_\cdot}\,\frac{g_k}{g_n}\,
\frac{(\log x_k)^{p_0}}{(\log x_n)^{q_\cdot}} .
\end{equation*}
Since $\log x_k\ge\log\gamma^{-2}\ge 1$, by the first condition in \eqref{4.4:eq-radii-ratio-1}, and
$q_\cdot\ge p_0$, we have $(\log x_k)^{p_0}\le(\log x_k)^{q_\cdot}$. The two bounds above at the pair
$(n,k)$ then give the second line of \eqref{4.4:eq-radii-ratio-a}.
\end{proof}

Throughout, $\xi>0$ is the lattice spacing of the bonds considered. For a bond $b$, $b_-$ and $b_+$ are its initial and final points and $\bar b$ is the reversed bond, and $e_\nu$ is the $\nu$-th unit vector. $G=SU(N)$ is the gauge group and $\mathfrak{g}^c$ is the complexification of its Lie algebra.

On $N\times N$ matrices we put $|X|^2:=\operatorname{tr}X^*X$, with $\operatorname{tr}1=1$, and write $\|X\|$ for the operator norm, so that $|X|\le\|X\|\le\sqrt N\,|X|$. We set $R(W)X:=WXW^{-1}$ and $\operatorname{Im}W:=(W-W^{-1})/2i$, and $\pi$ denotes the projection of the $N\times N$ matrices onto $\mathfrak{g}^c$. For a configuration $V$ on the bonds,
\[
J(V):=D^{\xi*}_V\,\xi^{-2}\pi\operatorname{Im}\partial V
\]
as in \cite[(1.8)]{B-CMP109}. Here $D^{\xi}_V$ is the covariant derivative of a bond field \cite[(3.4)]{B-CMP99b}, $D^{\xi*}_V$ is its transpose, and $\partial V$ is the plaquette holonomy. The projection $\pi$ commutes with $R(W)$.

In \cite[p.~272, (3.11)]{B-CMP109} the identity
\begin{multline*}
D^{\xi*}_{\exp(i\xi\mathbf{A})U}\,\xi^{-2}\pi\operatorname{Im}\partial\bigl(\exp(i\xi\mathbf{A})U\bigr)\\
=D^{\xi*}_U\,\xi^{-2}\pi\operatorname{Im}\partial U+D^{\xi*}_UD^{\xi}_U\mathbf{A}+\mathbf{F}(U,\mathbf{A})
\end{multline*}
is written for any regular configuration $U$. There $\mathbf{F}$ is a local operator depending on $U$, $\partial U$, $\mathbf{A}$, $\nabla^{\xi}_U\mathbf{A}$ only. We take this identity as the definition of $\mathbf{F}$. For a $G$-valued configuration $U_0$ and a $\mathfrak{g}^c$-valued field $A$ it reads
\begin{equation}\label{4.4:eq-current-remainder-1}
F(U_0,A):=\mathbf{F}(U_0,A)
=J\bigl(e^{i\xi A}U_0\bigr)-J(U_0)-D^{\xi*}_{U_0}D^{\xi}_{U_0}A .
\end{equation}

On a reversed bond we put $A(\bar b'):=-R(U_0(b'))^{-1}A(b')$. With this convention $V:=e^{i\xi A}U_0$ satisfies $V(c)=e^{i\xi A(c)}U_0(c)$ on every oriented bond $c$. Indeed,
\[
V(\bar b')=U_0(b')^{-1}e^{-i\xi A(b')}=e^{i\xi A(\bar b')}U_0(b')^{-1}.
\]

For a bond $b$ and a plaquette $p\ni b$, let $\partial_bp$ be the boundary contour of $p$ that starts at $b_-$ and runs through $b$ first. We write $V(\partial_bp)$ for the product of the values of $V$ along $\partial_bp$, in the order of traversal. Then
\begin{equation}\label{4.4:eq-current-remainder-a}
J(V)(b)=\xi^{-3}\sum_{p\ni b}\pi\operatorname{Im}V(\partial_bp).
\end{equation}

To see \eqref{4.4:eq-current-remainder-a}, take a plaquette $p=\langle x,y,z,w\rangle$ with $b_1=\langle x,y\rangle$, $b_2=\langle y,z\rangle$, $b_3=\langle w,z\rangle$ and $b_4=\langle x,w\rangle$. By \cite[(3.4)]{B-CMP99b},
\[
(D^{\xi}_VA)(p)=\xi^{-1}\bigl(A(b_1)+R(V(b_1))A(b_2)-R(V(b_4))A(b_3)-A(b_4)\bigr).
\]
Its transpose carries a plaquette field $f$ to the following contributions from $p$:
\begin{itemize}
\item $\xi^{-1}f(p)$ at $b_1$;
\item $\xi^{-1}R(V(b_1))^{-1}f(p)$ at $b_2$;
\item $-\xi^{-1}R(V(b_4))^{-1}f(p)$ at $b_3$;
\item $-\xi^{-1}f(p)$ at $b_4$.
\end{itemize}

Apply this to $f=\xi^{-2}\pi\operatorname{Im}V(\partial p)$, with $V(\partial p)=V(b_1)V(b_2)V(b_3)^{-1}V(b_4)^{-1}$. We use three facts: $\pi$ commutes with $R(W)$; $R(W)^{-1}\operatorname{Im}V(\partial p)=\operatorname{Im}(W^{-1}V(\partial p)W)$; and $-\operatorname{Im}V(\partial p)=\operatorname{Im}V(\partial p)^{-1}$. With these, in each of the four cases the contribution at $b_i$ is $\xi^{-3}\pi\operatorname{Im}$ of $V$ along the contour of $p$ re-based at $(b_i)_-$ and oriented along $b_i$, that is, along $\partial_{b_i}p$. For instance, at $b_3$,
\[
\bigl(V(b_4)^{-1}V(\partial p)V(b_4)\bigr)^{-1}=V(b_3)V(b_2)^{-1}V(b_1)^{-1}V(b_4).
\]

\begin{lemma}[The quadratic remainder in the current and its Lipschitz bound]
\label{4.4:lem-current-remainder}
Let $\xi\le\ell$ and let $b$ be a bond. Let $U_0$ be $G$-valued on the bonds of the $2(d-1)$ plaquettes containing $b$, with
\[
|U_0(\partial p)-1|\le\mathsf{c}\,\xi^2\ell^{-2},\qquad \mathsf{c}\le1,
\]
on these plaquettes. For $\mathfrak{g}^c$-valued $A$ on these bonds put
\[
\|A\|:=\max\bigl\{\ell|A|,\ \ell^2|\nabla_{U_0}A|\bigr\}.
\]
Here $|A|$ is the supremum of $|A(b')|$ over the bonds $b'$ of this set. The quantity $|\nabla_{U_0}A|$ is the supremum of the moduli of the covariant difference quotients between neighbouring parallel bonds of this set. For parallel bonds $b'$ and $b'+\xi e_\nu$, with $e$ the bond from $b'_-$ to $b'_-+\xi e_\nu$, this quotient, the covariant difference of \cite[(3.4)]{B-CMP99b}, is
\[
\xi^{-1}\bigl(R(U_0(e))A(b'+\xi e_\nu)-A(b')\bigr).
\]
There is $c_F=c_F(d,N)$ with the following two properties.
\begin{enumerate}
\item[(i)] For $\|A\|\le1$,
\[
\ell^3|F(U_0,A)(b)|\le c_F(\mathsf{c}+\|A\|)\|A\|.
\]
\item[(ii)] For $\|A_1\|,\|A_2\|\le\mathsf{a}$ and $2\mathsf{a}+\mathsf{c}\le1$,
\begin{equation}\label{4.4:eq-current-remainder-b}
\ell^3\bigl|F(U_0,A_2)(b)-F(U_0,A_1)(b)\bigr|
\le4c_F(\mathsf{c}+\mathsf{a})\,\|A_2-A_1\| .
\end{equation}
\end{enumerate}
\end{lemma}

\begin{proof}
(i) \emph{The plaquette factor.} Let $b$ point in the direction $e_\mu$. Let $p\ni b$ lie in the plane of $e_\mu,e_\nu$, on the side $\sigma\in\{+,-\}$ of $b$. Then $\partial_bp=(b,c_2,\bar b_3,\bar c_4)$, where:
\begin{itemize}
\item $c_2$ is the oriented bond from $b_+$ to $b_++\sigma\xi e_\nu$;
\item $b_3:=b+\sigma\xi e_\nu$;
\item $c_4$ is the oriented bond from $b_-$ to $b_-+\sigma\xi e_\nu$.
\end{itemize}

Thus $V(\partial_bp)=V(b)V(c_2)V(b_3)^{-1}V(c_4)^{-1}$. Insert $V(c)=e^{i\xi A(c)}U_0(c)$ and move the factors $U_0$ to the right, using $We^{X}=e^{R(W)X}W$ and $U_0(b)U_0(c_2)U_0(b_3)^{-1}=U_0(\partial_bp)U_0(c_4)$. This gives
\[
V(\partial_bp)=E_p\,U_0(\partial_bp),\qquad
E_p:=e^{i\xi a_1}e^{i\xi a_2}e^{-i\xi a_3}e^{-i\xi a_4},
\]
where
\begin{align*}
a_1&:=A(b), & a_2&:=R(U_0(b))A(c_2),\\
a_3&:=R(U_0(\partial_bp))R(U_0(c_4))A(b_3), & a_4&:=R(U_0(\partial_bp))A(c_4).
\end{align*}

\emph{Isometries and the bound on $\delta_i$.} Since $U_0$ is unitary, each $R(U_0(\cdot))$ is an isometry for $|\cdot|$. For unitary $W$ we have $R(W)X-X=(W-1)XW^{-1}+X(W^{-1}-1)$ and $\|W^{-1}-1\|=\|W-1\|$. Hence $\|R(W)-1\|\le2\|W-1\|$, in the sense that $|R(W)X-X|\le2\|W-1\|\,|X|$. In particular $|a_i|\le|A|$.

The differences $R(U_0(c_4))A(b_3)-A(b)$ and $R(U_0(b))A(c_2)-A(c_4)$ are, up to such isometries, $\xi$ times covariant difference quotients of $A$ between neighbouring parallel bonds. When $\sigma=-$, re-basing the reversed bonds adds terms which, up to such isometries, are of the form $(R(U_0(\partial p))^{\pm1}-1)A$. Moreover
\begin{align*}
a_3-a_1&=R(U_0(\partial_bp))\bigl[R(U_0(c_4))A(b_3)-A(b)\bigr]
+\bigl(R(U_0(\partial_bp))-1\bigr)A(b),\\
a_4-a_2&=-\bigl[R(U_0(b))A(c_2)-A(c_4)\bigr]+\bigl(R(U_0(\partial_bp))-1\bigr)A(c_4).
\end{align*}
Every term of the form $(R(W)-1)A(c)$ occurring here has $W$ a conjugate of $U_0(\partial p)^{\pm1}$. Hence it is bounded by
\[
2\|W-1\|\,|A|\le2\sqrt N\,\mathsf{c}\,\xi^2\ell^{-2}|A|,
\]
and at most four such terms occur in each of $a_3-a_1$ and $a_4-a_2$. Put $\delta_1:=\xi^{-1}(a_3-a_1)$ and $\delta_2:=\xi^{-1}(a_4-a_2)$. Using $\xi\le\ell$, $\mathsf{c}\le1$, $\ell|A|\le\|A\|$ and $\ell^2|\nabla_{U_0}A|\le\|A\|$, we obtain
\begin{equation}\label{4.4:eq-current-remainder-2}
|\delta_i|\le|\nabla_{U_0}A|+8\sqrt N\,\mathsf{c}\,\xi\ell^{-2}|A|\le9\sqrt N\,\ell^{-2}\|A\| .
\end{equation}
We write $|a|:=\max_i|a_i|\le|A|$ and $|\delta|:=\max_i|\delta_i|$. When the dependence on $p$ matters we write $a_2^p$, $\delta_i^p$.

\emph{Expansion of $E_p$.} Put $x_i:=\xi a_i$ and $y_i:=\xi^2\delta_i$ for $i=1,2$. Then
\[
E_p=e^{ix_1}e^{ix_2}e^{-i(x_1+y_1)}e^{-i(x_2+y_2)},
\]
a convergent power series in the four letters $x_1,x_2,y_1,y_2$. Its terms of first order in $x$ cancel. Those of first order in $y$ are $-i(y_1+y_2)$. Those of second order in $x$ are those of the group commutator, $[ix_1,ix_2]=-[x_1,x_2]$. Every other monomial has degree at least three in $x$, or at least one in $x$ and at least one in $y$, or at least two in $y$.

The same holds for
\[
E_p^{-1}=e^{i(x_2+y_2)}e^{i(x_1+y_1)}e^{-ix_2}e^{-ix_1}=1+i(y_1+y_2)+[x_1,x_2]+\cdots .
\]
By $\xi\le\ell$, $\|A\|\le1$ and \eqref{4.4:eq-current-remainder-2} we have $\|x_i\|\le\sqrt N$ and $\|y_i\|\le9N$. Hence, with a constant $c=c(N)$ (here and below the letter $c$ denotes a constant depending on $d$ and $N$ only, whose value may change from line to line; it is distinct from the constant $\mathsf{c}$ of the hypothesis),
\begin{equation}\label{4.4:eq-current-remainder-3}
\begin{gathered}
\frac{E_p-E_p^{-1}}{2i}=-\xi^2(\delta_1+\delta_2)+i\xi^2[a_1,a_2]+\rho_p,\\
|\rho_p|\le c\bigl(\xi^3|a|^3+\xi^3|a||\delta|+\xi^4|\delta|^2\bigr),\qquad
|E_p^{\pm1}-1|\le c\,\xi^2\bigl(|\delta|+|a|^2\bigr).
\end{gathered}
\end{equation}

\emph{The current difference.} Let $W:=U_0(\partial_bp)$. Then
\[
\operatorname{Im}(E_pW)-\operatorname{Im}W=\frac{E_p-E_p^{-1}}{2i}
+\frac1{2i}\Bigl[(E_p-1)(W-1)-(W^{-1}-1)(E_p^{-1}-1)\Bigr].
\]
Indeed, expanding the bracket, the right side equals
\[
\frac1{2i}\bigl[E_pW-W-W^{-1}E_p^{-1}+W^{-1}\bigr],
\]
which is $\operatorname{Im}(E_pW)-\operatorname{Im}W$ because $(E_pW)^{-1}=W^{-1}E_p^{-1}$.
In the second term, $|W^{\pm1}-1|\le\mathsf{c}\,\xi^2\ell^{-2}$ and $|E_p^{\pm1}-1|$ is bounded by \eqref{4.4:eq-current-remainder-3}; also $|XY|\le\|X\|\,|Y|$. We apply \eqref{4.4:eq-current-remainder-a} to $V$ and to $U_0$, use $\pi$ as a fixed linear map on a space of dimension depending on $N$, and absorb the number $2(d-1)$ of plaquettes into the constants. This gives
\begin{equation}\label{4.4:eq-current-remainder-4}
\begin{split}
J(V)(b)-J(U_0)(b)={}&\sum_{p\ni b}\pi\Bigl\{-\xi^{-1}(\delta_1^p+\delta_2^p)+i\xi^{-1}[a_1,a_2^p]\Bigr\}\\
&+\sum_{p\ni b}O\Bigl(|a|^3+|a||\delta|+\xi|\delta|^2+\xi\mathsf{c}\ell^{-2}\bigl(|\delta|+|a|^2\bigr)\Bigr).
\end{split}
\end{equation}
By \eqref{4.4:eq-current-remainder-1}, $F(U_0,A)(b)$ is the right side of \eqref{4.4:eq-current-remainder-4} minus $(D^{\xi*}_{U_0}D^{\xi}_{U_0}A)(b)$. We estimate the three resulting contributions in turn.

$(\alpha)$ \emph{The linear terms.} We have $-\xi^{-1}(\delta_1^p+\delta_2^p)=\xi^{-2}(a_1+a_2-a_3-a_4)$. By the transpose rule used in the proof of \eqref{4.4:eq-current-remainder-a}, $(D^{\xi*}_{U_0}D^{\xi}_{U_0}A)(b)$ is $\xi^{-2}\sum_{p\ni b}$ of the same four fields, but transported as in \cite[(3.4)]{B-CMP99b}. These $\mathfrak{g}^c$-valued fields are left unchanged by $\pi$. The transports of \cite[(3.4)]{B-CMP99b} differ from those in $a_1,\dots,a_4$ by factors $R(U_0(\partial p))^{\pm1}$ on two of the four fields when $\sigma=+$, and on three when $\sigma=-$. By the bound on $R(W)-1$ above, the difference between the linear terms of \eqref{4.4:eq-current-remainder-4} and $(D^{\xi*}_{U_0}D^{\xi}_{U_0}A)(b)$ is at most
\[
2(d-1)\cdot\xi^{-2}\cdot2\sqrt N\,\mathsf{c}\,\xi^2\ell^{-2}\cdot3|A|\le c\,\mathsf{c}\,\ell^{-2}|A|,
\]
with $c=c(d,N)$.

$(\beta)$ \emph{The commutator terms.} Take the two plaquettes, $\sigma=\pm$, of one plane. Put $y:=b_+$ and $e:=\langle y-\xi e_\nu,y\rangle$, and write $A_\nu(y')$ for the value of $A$ on the bond from $y'$ to $y'+\xi e_\nu$. Then $c_2^+$ runs from $y$ to $y+\xi e_\nu$ and $c_2^-=\bar e$. By the reversed-bond convention and the definition of the covariant difference quotient,
\[
A(c_2^+)+A(c_2^-)=A_\nu(y)-R(U_0(e))^{-1}A_\nu(y-\xi e_\nu)
=\xi\,R(U_0(e))^{-1}(\nabla_\nu A_\nu)(y-\xi e_\nu).
\]
Since $a_2^{\nu\sigma}=R(U_0(b))A(c_2^\sigma)$ and $R(U_0(\cdot))$ is an isometry,
\[
\Bigl|\xi^{-1}\sum_{\sigma=\pm}[a_1,a_2^{\nu\sigma}]\Bigr|\le c_{\mathfrak g}\,|A|\,|\nabla_{U_0}A| .
\]
Here $c_{\mathfrak g}=c_{\mathfrak g}(N)$ is a constant with $|[X,Y]|\le c_{\mathfrak g}|X||Y|$. Hence the commutator terms of \eqref{4.4:eq-current-remainder-4}, summed over the $d-1$ planes, are bounded by $c\,\ell^{-3}\|A\|^2$.

$(\gamma)$ \emph{The remainder terms.} By \eqref{4.4:eq-current-remainder-2}, $\|A\|\le1$ and $|a|\le\ell^{-1}\|A\|$, the $O$-terms of \eqref{4.4:eq-current-remainder-4} are bounded by
\[
c\,\ell^{-3}\Bigl(\|A\|^3+\|A\|^2+\frac{\xi}{\ell}\|A\|^2+\frac{\xi}{\ell}\,\mathsf{c}\,\|A\|\Bigr).
\]

\emph{Collecting.} $F(U_0,A)(b)$ is the sum of the difference in $(\alpha)$, the commutator terms in $(\beta)$ and the $O$-terms in $(\gamma)$. The three bounds, together with $\ell|A|\le\|A\|\le1$ and $\xi\le\ell$, give
\[
\ell^3|F(U_0,A)(b)|\le c_F(\mathsf{c}+\|A\|)\|A\|
\]
with $c_F=c_F(d,N)$. In particular, the computation exhibits $F(U_0,A)(b)$ as a function of $U_0$, $U_0(\partial p)$, $A$ and $\nabla_{U_0}A$ on the bonds of the plaquettes containing $b$ only.

(ii) \emph{Holomorphy.} The map $A\mapsto F(U_0,A)(b)$ is entire in the finitely many variables $A(b')$. Indeed, by \eqref{4.4:eq-current-remainder-1} and \eqref{4.4:eq-current-remainder-a} it is built from products of exponentials of linear functions of $A$, their inverses (again such products), and the holomorphic maps $\pi$ and $\operatorname{Im}$, together with the linear term $D^{\xi*}_{U_0}D^{\xi}_{U_0}A$.

\emph{The derivative bound.} If $\mathsf{a}+\mathsf{c}=0$, then $A_1=A_2=0$ and there is nothing to prove. Otherwise let $\|A\|\le\mathsf{a}$, let $\delta\ne0$ be a $\mathfrak{g}^c$-valued field on the same bonds, and put $R:=(\mathsf{a}+\mathsf{c})/\|\delta\|$. For $|\zeta|=R$,
\[
\|A+\zeta\delta\|\le\mathsf{a}+R\|\delta\|=2\mathsf{a}+\mathsf{c}\le1 ,
\]
so (i) applies to $A+\zeta\delta$. By Cauchy's formula,
\[
dF(U_0,A)[\delta]=\frac1{2\pi i}\oint_{|\zeta|=R}\zeta^{-2}F(U_0,A+\zeta\delta)(b)\,d\zeta .
\]
Together with (i) this yields
\[
\ell^3\bigl|dF(U_0,A)[\delta]\bigr|\le R^{-1}c_F(2\mathsf{c}+2\mathsf{a})(2\mathsf{a}+\mathsf{c})
\le4c_F(\mathsf{a}+\mathsf{c})\|\delta\| .
\]

\emph{Integration.} Integrate this bound along the segment from $A_1$ to $A_2$, on which $\|\cdot\|\le\mathsf{a}$ by convexity:
\[
F(U_0,A_2)(b)-F(U_0,A_1)(b)=\int_0^1dF\bigl(U_0,A_1+t(A_2-A_1)\bigr)[A_2-A_1]\,dt .
\]
This gives \eqref{4.4:eq-current-remainder-b}.
\end{proof}

\begin{lemma}[Stability of the layered spaces under the twist]\label{4.4:lem-twist-stability}
Let $\alpha_{0,n},\alpha_{1,n}$ be the radii of \cite[(2.28)]{B-CMP119} at level $n$. Where an
assertion concerns either of them, we write $\alpha_{\cdot,n}$, and $C_{\cdot}$ for the
corresponding constant $C_0$ or $C_1'$ of \cite[(2.28)]{B-CMP119}. Let $C_1$, $\|C\|$, $A_1$,
$\theta\in\,]0,1[$ and the twist size $w_k$ be as in Definition~\ref{4.4:def-post-step-space},
$C_1$ being there the constant of the polydisc $\mathfrak{D}$,
and let $C_H\ge 1$, $\delta_H$, the decay margins $D_n$ and the least layer width
$\mathsf{w}_*$ be as in Definition~\ref{4.4:not-response-bound}. Put
\begin{equation*}
h_n:=C_H\,w_k\,e^{-D_n},\qquad c:=8c_F(1+B_3),
\end{equation*}
where $c_F=c_F(d,N)$ is the constant of Lemma~\ref{4.4:lem-current-remainder} and $B_3$ is the
constant in the bound of \cite[Theorem~1, Proposition~9]{B-CMP102b} for the plaquette holonomies
of a real minimiser, in the form in which \cite{B-CMP109} uses it. Let $G=\mathrm{SU}(N)$.
Assume:
\begin{itemize}
\item the response bound at free current, Definition~\ref{4.4:not-response-bound}, that is
\eqref{4.4:eq-response-bound-a} together with all five members of \cite[(190)]{B-CMP102b} at
real minimiser bases with complex data in the polydisc \cite[(172)]{B-CMP102b};
\item condition \eqref{4.4:eq-post-step-space-a} of Definition~\ref{4.4:def-post-step-space};
\item $q_0,q_1\ge p_0$, the conclusion \eqref{4.4:eq-radii-ratio-a} of
Lemma~\ref{4.4:lem-radii-ratio} for every pair of levels $n\le k'\le k+1$, and the ceilings on
$\gamma$ under which that lemma derives it;
\item the conditions
\begin{equation}\label{4.4:eq-twist-stability-a}
\begin{aligned}
&(\mathrm{s}0)\quad \gamma\le\gamma(C_0,C_1',C_H),\\
&(\mathrm{s}1)\quad 2\sqrt{2}\,e^{-\delta_H\mathsf{w}_*}\le 1,\\
&(\mathrm{s}2)\quad L\ge 2\sqrt{2}\,(1+\beta)(1+\beta_0),\\
&(\mathrm{s}3)\quad C_0,\ C_1'\ \ge\ \frac{8c_{\flat}(1+\beta_0)\,\|C\|\,C_1A_1C_H^2}{(1-\theta)\beta},\\
&(\mathrm{s}3')\quad C_0>2C_{IV}(d,L)\,C_H(1+\beta_0)\|C\|C_1A_1 .
\end{aligned}
\end{equation}
\end{itemize}
Here (s0) is a ceiling on $\gamma$, depending on $C_0,C_1',C_H$ and the constants fixed
before them, of the kind made in \cite{B-CMP109}, where $\gamma$ is chosen depending on all
other constants; it is the ceiling under which, for all levels $n,n'$,
\begin{equation*}
h_n\le C_H^2w_k\le\alpha_{\cdot,k}\le 1,\qquad
c\,\bigl(\alpha_{0,n}+C_H(\alpha_{1,n'}+h_{n'})\bigr)\le 1 .
\end{equation*}
Since $\mathsf{w}_*$ exceeds $M$ times the radius parameter of the admissibility condition
\cite[(2.2)]{B-CMP96} for every admissible sequence, condition (s1) follows as soon as
$\delta_H$ times this product is at least $\log(2\sqrt{2})$, a floor on
$M$. Condition (s2) is a floor on $L$ by an absolute number: since $\beta=\tfrac14$ and
$\beta_0\le 1$, its right side is at most $5\sqrt{2}<8$, so (s2) is implied by the standing
condition of \cite{B-CMP109} that $L$ is an odd integer greater than $11$. In (s3),
$c_{\flat}:=\max(4,c_1')$, with $c_1'$ the
absolute constant of \eqref{4.4:eq-twist-stability-2} below. In (s3'), $C_{IV}(d,L)$ is the
constant of the uniform-increment lemma among the operator bounds at free current. All
constants appearing in (s3) and (s3') are fixed before $C_0$, and (s3') is assumed only for
conclusion $(2^{\mathrm{int}})$. Condition (s3) contains the inequality
$C_1'>16\,C_H(1+\beta_0)\|C\|C_1A_1$, since $c_\flat\ge4$, $C_H\ge1$ and $(1-\theta)\beta<1$.

Then for all $1\le j\le k<\infty$, all regions $Y$ on which the space
$\mathfrak{U}^{\mathrm{post}}_{k+1}(Y)$ of Definition~\ref{4.4:def-post-step-space} is
considered, all $X\in\mathbf{D}_j$ with $X\subset Y$ and all
fluctuation data $(A,A_k,\mathbf{H}^{(k)})\in\mathfrak{D}$, the twisted argument
$(\mathbf{U}^{\flat},\mathbf{J}^{\flat})$ of Definition~\ref{4.4:def-twisted-argument} has the
following properties, each under the additional hypotheses stated with it:
\begin{equation}\label{4.4:eq-twist-stability-c}
\begin{aligned}
&\text{(1) if $(\mathbf{U},\mathbf{J})\in\mathfrak{U}^{\mathrm{post}}_{k+1}(Y)$:}\\
&\qquad(\mathbf{U}^{\flat},\mathbf{J}^{\flat})|_X\in
\tilde{U}^c_j(X,\tilde\alpha_0,\tilde\alpha_1);\\
&\text{(1$'$) if moreover \eqref{4.4:eq-group-factor-bounds-a} holds on $Y$:}\\
&\qquad\text{\eqref{4.4:eq-layered-spaces-b}, \eqref{4.4:eq-layered-spaces-d} hold at level
$j$ on $X$ with $U$ in place of $\mathbf{U}^{\flat}$;}\\
&\text{(2) if moreover $X\subset\Lambda_j$, \eqref{4.4:eq-group-factor-bounds-a} holds on $Y$
and \eqref{4.4:eq-twist-stability-b} holds:}\\
&\qquad(\mathbf{U}^{\flat},\mathbf{J}^{\flat})|_X\in U^c_j(X,\alpha_{0,j},\alpha_{1,j});\\
&\text{($2^{\mathrm{int}}$) if instead $Y\subset\Lambda_{k+1}$ and $(\mathbf{U},\mathbf{J})$
is a}\\
&\qquad\text{representative of the class in item ($2^{\mathrm{int}}$) below:}\\
&\qquad(\mathbf{U}^{\flat},\mathbf{J}^{\flat})|_X\in
U^{\Sigma^1}_j(X;\alpha_{0,j},\alpha_{1,j})\subset U^c_j(X,\alpha_{0,j},\alpha_{1,j}).
\end{aligned}
\end{equation}
The four conclusions are meant in the following sense.
\begin{enumerate}
\item[(1)] In (1) the membership means that all of
\eqref{4.4:eq-layered-spaces-a}--\eqref{4.4:eq-layered-spaces-f} hold for
$(\mathbf{U}^{\flat},\mathbf{J}^{\flat})$ on $X$ at level $j$.
\item[(1$'$)] If moreover \eqref{4.4:eq-group-factor-bounds-a} holds on $Y$, then
\eqref{4.4:eq-layered-spaces-b} and \eqref{4.4:eq-layered-spaces-d}, read for the sets of
\eqref{4.4:eq-post-step-space-a}, hold at level $j$ on $X$ also with $U^{\flat}=U$ in place of
$\mathbf{U}^{\flat}$.
\item[(2)] If $X\subset\Lambda_j$, if \eqref{4.4:eq-group-factor-bounds-a} holds on $Y$, and if
the further ceiling on $\gamma$ (imposed after $L$, $M$, $B$, $C_0$, $q_0$)
\begin{equation}\label{4.4:eq-twist-stability-b}
6(\mathsf{o}LM)^2B\,\alpha_{0,j}\le r_{\mathrm{gl}}(d,G)\qquad\text{for all }j
\end{equation}
holds, then conclusion (2) holds, where $U^c_j(X,\alpha_{0,j},\alpha_{1,j})$ is the space of
\cite[(2.27)(ii)]{B-CMP119}, that is, the space of \cite[(1.11)--(1.16)]{B-CMP109} at level $j$,
with the constant of \cite[(1.12)]{B-CMP109} equal to
\begin{equation}\label{4.4:eq-twist-stability-e}
B^{109}:=\max\{4C_{\mathrm{gl}}(d,G)\,B,\ B_{\mathrm{P}}\}.
\end{equation}
Here $B$ is the absolute constant of \eqref{4.4:eq-layered-spaces-e}; $\mathsf{o}LM$ is the
cube side in \cite[(1.12)]{B-CMP109}, $\mathsf{o}$ the absolute number there;
$C_{\mathrm{gl}}$ and $r_{\mathrm{gl}}$ are the constants of the gluing lemma for boxes and its
corollary for cubes, which use \eqref{4.4:eq-twist-stability-e} only through the inequality
$B^{109}\ge 4C_{\mathrm{gl}}B$; and $B_{\mathrm{P}}$ is the floor on the constant of
\cite[(1.12)]{B-CMP109} under which the pairs of \cite[(3.38)]{B-CMP109} lie in the one-power
hereditary class.
\item[$(2^{\mathrm{int}})$] Suppose $Y\subset\Lambda_{k+1}$ and, in place of
$(\mathbf{U},\mathbf{J})\in\mathfrak{U}^{\mathrm{post}}_{k+1}(Y)$ (no condition
\eqref{4.4:eq-layered-spaces-e} being required), $(\mathbf{U},\mathbf{J})$ is a representative
of the one-power hereditary class
\begin{equation*}
U^{\Sigma^1}_{k+1}\bigl(Y;(1+\beta)\alpha_{0,k+1},(1+\beta)\alpha_{1,k+1},\alpha_{0,k+1}\bigr),
\end{equation*}
that is, it satisfies \cite[(1.11)--(1.16)]{B-CMP109} at level $k+1$ on $Y$, with
\cite[(1.12)]{B-CMP109} in the reading fixed for that class and with \cite[(1.16)]{B-CMP109}
imposed for the tower of every $(p,Z)\in\mathfrak{T}_{k+1}(Y)$ at one power of the scale ratio
(this class is contained in the space $U^c_{k+1}(Y,\cdot)$ of \cite[(1.34)]{B-CMP116}, read in
the same way). Assume (s3'). Take as stated the following two statements. The first is the
restriction theorem for the one-power hereditary class, by which a representative of the class
at level $i$ on a localization domain restricts to a localization domain of level $i'\le i$
contained in it as a representative at level $i'$, with the same group-valued factor and with
the first two parameters multiplied by $L^{-(i-i')}$. The second is the statement that the pair
produced by the operator bounds at free current from a representative and a field whose size
and covariant differences are less than $\min\{\alpha_{1,j}/8,\alpha_{0,j}/C_{IV}\}$ in
level-$j$ units lies in the one-power hereditary class at level $j$; it is applied with the
response $\boldsymbol{\mathcal{H}}$ as that field, and that (s3) and (s3') secure its size
condition is shown in the proof of conclusion $(2^{\mathrm{int}})$. Then for every
$X\in\mathbf{D}_j$ with $X\subset Y$ conclusion $(2^{\mathrm{int}})$ holds; neither
\eqref{4.4:eq-group-factor-bounds-a} nor \eqref{4.4:eq-twist-stability-b} nor the gluing lemma
is used.
\end{enumerate}
\end{lemma}

\begin{proof}
We prove conclusion (1). Throughout, $\Xi:=\Omega_{k+1}\cup S_{k+1}$, and $d(\cdot,\cdot)$ is
the weighted distance of Definition~\ref{4.4:not-response-bound}.

\emph{The response on layer $n$.} On $\mathfrak{D}$ the data twist of the level-$k$ data is
supported in $\Xi$ and is bounded there by $w_k$ (Definition~\ref{4.4:def-post-step-space}).
For a point $y$ of layer $n$ the right member of \eqref{4.4:eq-response-bound-a} is therefore at
most $C_Hw_ke^{-\delta_Hd(y,\Xi)}$, and $\delta_Hd(y,\Xi)\ge D_n$, because $D_n$ is $\delta_H$
times the distance from layer $n$ to $\Xi$. Hence
\begin{equation}\label{4.4:eq-twist-stability-1}
\ell_n|\boldsymbol{\mathcal{H}}|,\quad \ell_n^2|\nabla_{\mathbf{U}}\boldsymbol{\mathcal{H}}|,\quad
\ell_n^3|\mathfrak{J}|\ \le\ h_n\qquad\text{on layer }n .
\end{equation}
Moreover $h_n\le C_H^2w_k$, because $C_H\ge1$ and $D_n\ge0$, and $C_H^2w_k\le\alpha_{\cdot,k}\le1$
by (s3) and (s0).

\emph{The increments.} Write $\mathbf{U}=U'U$ with $U$ group-valued and $U'=\exp i\xi A'$. By
Definition~\ref{4.4:def-twisted-argument}, $\mathbf{U}^{\flat}=U'^{\flat}U^{\flat}$ with
$U^{\flat}=U$ and $U'^{\flat}=e^{i\xi\boldsymbol{\mathcal{H}}}U'=:\exp i\xi A'^{\flat}$. For each
quantity entering \eqref{4.4:eq-layered-spaces-a}--\eqref{4.4:eq-layered-spaces-f} (each such
quantity is called a slot below) we bound its
increment under $(\mathbf{U},\mathbf{J})\mapsto(\mathbf{U}^{\flat},\mathbf{J}^{\flat})$ at a
point of layer $n$, measured in the slot's own unit: a quantity of weight $\varpi$ is measured
with the $\varpi$-th power of the layer unit $\ell_n$, as in those displays. The weight is given
in parentheses.

\emph{The first member of \eqref{4.4:eq-layered-spaces-a}: $\partial U$ $(2)$.} The increment is
$0$, since $U^{\flat}=U$.

\emph{The second member of \eqref{4.4:eq-layered-spaces-a}: $\partial\mathbf{U}$ $(2)$.} The
increment is at most $4h_n$. Indeed, transporting the four exponential factors
$e^{i\xi\boldsymbol{\mathcal{H}}}$ on the bonds of the boundary of a plaquette $p$ to the base
point of $p$ and combining them by the Baker--Campbell--Hausdorff formula gives
\begin{equation*}
\partial(e^{i\xi\boldsymbol{\mathcal{H}}}\mathbf{U})(p)
=\exp\bigl[i\xi^2(\operatorname{curl}_{\mathbf{U}}\boldsymbol{\mathcal{H}})(p)
+O(\xi^2|\boldsymbol{\mathcal{H}}|^2)\bigr]\,\partial\mathbf{U}(p),
\end{equation*}
where $\operatorname{curl}_{\mathbf{U}}\boldsymbol{\mathcal{H}}$ is the covariant curl of
$\boldsymbol{\mathcal{H}}$; moreover
$|\operatorname{curl}_{\mathbf{U}}\boldsymbol{\mathcal{H}}|\le
2|\nabla_{\mathbf{U}}\boldsymbol{\mathcal{H}}|$, and $(\ell_n|\boldsymbol{\mathcal{H}}|)^2\le
h_n^2$ by \eqref{4.4:eq-twist-stability-1}. The first two members of
\eqref{4.4:eq-twist-stability-1}, those of the response at free current, rest on clauses (K1),
(K3), (K8) of the existence theorem for the response chart
(Theorem~\ref{4.5:thm-chart-existence}), as they enter the operator bounds at free current.

\emph{\eqref{4.4:eq-layered-spaces-f}: $A'$ $(1)$ and $\nabla_UA'$ $(2)$.} The increment is at
most $4h_n$. By the Baker--Campbell--Hausdorff formula applied to
$\exp i\xi A'^{\flat}=e^{i\xi\boldsymbol{\mathcal{H}}}e^{i\xi A'}$,
\begin{equation*}
A'^{\flat}=A'+\boldsymbol{\mathcal{H}}+O(\xi|\boldsymbol{\mathcal{H}}||A'|),\qquad
\xi|A'|\le\alpha_{1,n}\le1;
\end{equation*}
and $\nabla_U\boldsymbol{\mathcal{H}}=\nabla_{\mathbf{U}}\boldsymbol{\mathcal{H}}
+O(|A'||\boldsymbol{\mathcal{H}}|)$ with
$\ell_n^2|A'||\boldsymbol{\mathcal{H}}|\le\alpha_{1,n}h_n$. The members of
\eqref{4.4:eq-twist-stability-1} used are again the first two, resting on the same clauses.

\emph{\eqref{4.4:eq-layered-spaces-c}: $\mathbf{J}$ $(3)$.} The increment is at most $h_n$: by
the third member of \eqref{4.4:eq-twist-stability-1},
\begin{equation*}
\ell_n^3|\mathbf{J}^{\flat}-\mathbf{J}|=\ell_n^3|\mathfrak{J}|\le C_Hw_ke^{-D_n}=h_n .
\end{equation*}
This is the $\mathfrak{J}$-member of the response bound, which rests on clauses (K2), (K3), (K8)
of Theorem~\ref{4.5:thm-chart-existence}, as they enter the operator bounds at free current.
The increment $\mathfrak{J}$ contains the quadratic
remainder of the current, so no separate term for it arises here, and no ceiling on $\gamma$ is
used in this slot.

\emph{\eqref{4.4:eq-layered-spaces-b}, \eqref{4.4:eq-layered-spaces-d}: $\partial U_{p,X}$
$(2)$ and $\mathbf{J}_{p,X}$ $(3)$.} We show that the increment is at most
\begin{equation}\label{4.4:eq-twist-stability-2}
c_1'\,C_H^2\,w_k\,e^{-D_n},\qquad c_1'\ \text{an absolute constant}.
\end{equation}
Consider the data twist
$\mathsf{b}:=\frac1i\log[M^{\cdot}(\mathbf{U}^{\flat})M^{\cdot}(\mathbf{U})^{-1}]$. For
$y'\in\Gamma_{n'}$, by the estimate used at \cite[(3.3)]{B-CMP109} and by
\eqref{4.4:eq-response-bound-a},
\begin{equation*}
|\mathsf{b}(y')|\le 2\ell_{n'}\sup|\boldsymbol{\mathcal{H}}|
\le 2C_Hw_ke^{-\delta_Hd(y',\Xi)},
\end{equation*}
and for all points $y,y'$ the triangle inequality gives $d(y,y')+d(y',\Xi)\ge d(y,\Xi)$.

The quantities in these two slots are functions of $\mathbf{U}^{\flat}$ alone, through
$M^{\cdot}(\mathbf{U}^{\flat})$. With $U$ the $G$-valued factor (so that $U^{\flat}=U$) put
$\mathsf{b}_{\mathrm{tot}}(\cdot):=\frac1i\log[M^{\cdot}(\cdot)M^{\cdot}(U)^{-1}]$. Then
$|\mathsf{b}_{\mathrm{tot}}|\le2(\alpha_{1,n'}+h_{n'})$ on $\Gamma_{n'}$, which lies in the
polydisc \cite[(172)]{B-CMP102b}. Let $U_0:=U_{p,X}(M^{\cdot}(U))$. This is a real minimiser,
which exists and is regular by \cite[Theorem~1]{B-CMP102b}, its data satisfying
\cite[(7)]{B-CMP102b} by the first member of \eqref{4.4:eq-layered-spaces-a} and
\cite[Proposition~2]{B-CMP98}. For a data twist $\mathsf{b}'$ let
$\mathcal{H}_p(\mathsf{b}')$ be the response of $U_0$ to $\mathsf{b}'$. Then
$U_{p,X}(M^{\cdot}(\cdot))$ equals, up to a gauge transformation,
$e^{i\xi\mathcal{H}_p(\mathsf{b}_{\mathrm{tot}})}U_0$, and there
\begin{equation*}
\mathbf{J}_{p,X}=J(U_0)+D^*D\mathcal{H}_p+F(U_0,\mathcal{H}_p),
\end{equation*}
where $D$ is the covariant difference operator at the base $U_0$ and $D^*$ its adjoint, in the
notation of \cite{B-CMP102b}.
The second part of the response bound (Definition~\ref{4.4:not-response-bound}) gives all five
members of \cite[(190)]{B-CMP102b} at complex data in the polydisc \cite[(172)]{B-CMP102b}, at
the real base $U_0$. Integrate along the segment from $\mathsf{b}_{\mathrm{tot}}(\mathbf{U})$ to
$\mathsf{b}_{\mathrm{tot}}(\mathbf{U}^{\flat})$; writing
$\Delta\mathsf{b}:=\mathsf{b}_{\mathrm{tot}}(\mathbf{U}^{\flat})
-\mathsf{b}_{\mathrm{tot}}(\mathbf{U})$, one has on it
$|\Delta\mathsf{b}(y')|\le2C_Hw_ke^{-\delta_Hd(y',\Xi)}$. This gives
\eqref{4.4:eq-twist-stability-2}, with the quadratic remainder treated as follows.

The requirement that $\mathsf{b}_{\mathrm{tot}}$ lie in the polydisc
\cite[(172)]{B-CMP102b} is a ceiling on $\alpha_{1,\cdot}$ and on $h_{\cdot}$, imposed after
$C_H$; it is of the kind made in \cite{B-CMP109}, where $\alpha_1$ depends on $M$, and
$h_{\cdot}\le\alpha_{\cdot,k}$ by (s0) and (s3). The remainder $F(U_0,\mathcal{H}_p)$ is
explicit at the real base, and what enters the bound is its increment $\Delta F$ along the
segment. To bound it, note first that $U_0$ is real and that, in the slot's unit,
$|U_0(\partial p)-1|\le\mathsf{c}\,\xi^2\ell_n^{-2}$ with $\mathsf{c}\le2B_3\alpha_{0,n}$, by
\cite[Theorem~1, Proposition~9]{B-CMP102b}, in the form in which \cite{B-CMP109} uses it (the
bound $2B_3\alpha_0'\xi^2$ for $|\partial U_n(M^{\cdot}(\mathbf{U}))-1|$). Put
$\mathsf{b}_1:=\mathsf{b}_{\mathrm{tot}}(\mathbf{U})$ and
$\mathsf{b}_2:=\mathsf{b}_{\mathrm{tot}}(\mathbf{U}^{\flat})$. Integrating the first two members
of \cite[(190)]{B-CMP102b} from $\mathcal{H}_p(0)=0$, and writing $\|\cdot\|$ for the norm of
Lemma~\ref{4.4:lem-current-remainder} with $\ell=\ell_n$,
\begin{equation*}
\|\mathcal{H}_p(\mathsf{b}_i)\|\le C_H\sup|\mathsf{b}_i|\le2C_H(\alpha_{1,n'}+h_{n'}),
\qquad i=1,2.
\end{equation*}
Apply \eqref{4.4:eq-current-remainder-b}, part (ii) of
Lemma~\ref{4.4:lem-current-remainder}, with the two fields of part (ii) taken to be
$\mathcal{H}_p(\mathsf{b}_1)$ and $\mathcal{H}_p(\mathsf{b}_2)$ and with
$\mathsf{a}:=2C_H\sup_{n'}(\alpha_{1,n'}+h_{n'})$; the ceiling (s0) also secures
$2\mathsf{a}+\mathsf{c}\le1$. Since $4c_F(\mathsf{c}+\mathsf{a})\le
8c_F(1+B_3)\bigl(\alpha_{0,n}+C_H\sup_{n'}(\alpha_{1,n'}+h_{n'})\bigr)$, this gives, with
$c=8c_F(1+B_3)$, a number depending on $d,L,N$,
\begin{equation*}
\ell_n^3|\Delta F|\le c\bigl(\alpha_{0,n}+C_H(\alpha_{1,n'}+h_{n'})\bigr)
\max\{\ell_n|\Delta\mathcal{H}_p|,\ \ell_n^2|\nabla\Delta\mathcal{H}_p|\},
\end{equation*}
the supremum over the layers $n'$ being understood in the bracket, and a fortiori
\begin{equation*}
\ell_n^3|\Delta F|\le c\bigl(\alpha_{0,n}+C_H(\alpha_{1,n'}+h_{n'})\bigr)
\bigl(\ell_n|\Delta\mathcal{H}_p|+\ell_n^2|\nabla\Delta\mathcal{H}_p|\bigr).
\end{equation*}
Here $\Delta\mathcal{H}_p:=\mathcal{H}_p(\mathsf{b}_2)-\mathcal{H}_p(\mathsf{b}_1)$, and by the
first two members of \cite[(190)]{B-CMP102b}, the bound on $\Delta\mathsf{b}$ and the triangle
inequality above,
\begin{equation*}
\max\{\ell_n|\Delta\mathcal{H}_p|,\ \ell_n^2|\nabla\Delta\mathcal{H}_p|\}
\le C_H\sup_{y'}\bigl[|\Delta\mathsf{b}(y')|e^{-\delta_Hd(y,y')}\bigr]
\le2C_H^2w_ke^{-D_n}.
\end{equation*}
Hence
\begin{equation*}
\ell_n^3|\Delta F|\le c\bigl(\alpha_{0,n}+C_H(\alpha_{1,n'}+h_{n'})\bigr)\cdot
2C_H^2w_ke^{-D_n}.
\end{equation*}
Under
(s0) the coefficient $c(\alpha_{0,n}+C_H(\alpha_{1,n'}+h_{n'}))$ is at most $1$, so $c_1'$ in
\eqref{4.4:eq-twist-stability-2} is an absolute constant, and hence so is $c_\flat$ below. When
the fluctuation data vanish,
$\mathsf{b}_2=\mathsf{b}_1$ and $\Delta F=0$, while $F(U_0,\mathcal{H}_p(\mathsf{b}_1))$ itself
does not vanish at complex $\mathbf{U}$; it belongs to the untwisted
$\mathbf{J}_{p,X}(M^{\cdot}(\mathbf{U}))$, which \eqref{4.4:eq-layered-spaces-d} at level $k+1$
bounds as a whole.

The regularity of the background of the sets of \eqref{4.4:eq-post-step-space-a} is the
untwisted bound given by $(\mathbf{U},\mathbf{J})\in\mathfrak{U}^{\mathrm{post}}_{k+1}(Y)$.
These two slots thus rest on the second part of the response bound and on the first member of
\eqref{4.4:eq-twist-stability-1}, which enters the bound on $\mathsf{b}$ and rests on the
operator bounds at free current.

\emph{\eqref{4.4:eq-layered-spaces-e}: the gauge condition of $U$ on cubes.} There is no
increment, since $U^{\flat}=U$; the comparison on cubes is carried out after this proof.

\emph{The schedule.} Every increment above is therefore at most $\pi_n\alpha_{\cdot,n}$, where
\begin{equation*}
\pi_n:=c_\flat C_H^2\,\frac{w_k}{\alpha_{\cdot,n}}\,e^{-D_n},\qquad
c_\flat=\max(4,c_1')\ \text{absolute};
\end{equation*}
for the quantities whose increment was bounded by $4h_n$ or by $h_n$ this uses $C_H\ge1$. By
Lemma~\ref{4.4:lem-radii-ratio} with $k'=k$,
\begin{equation*}
\frac{w_k}{\alpha_{\cdot,n}}\le\frac{2\|C\|C_1A_1}{C_\cdot}\,(1+\beta_0)(1+k-n)^{1/2};
\end{equation*}
moreover
$D_n\ge\delta_H\mathsf{w}_*(k-n)$ and $(1+k-n)^{1/2}\le(\sqrt{2})^{k-n}$. Hence
\begin{equation*}
\begin{split}
2^{k+1-n}\pi_n
&\le\frac{4c_\flat(1+\beta_0)\|C\|C_1A_1C_H^2}{C_\cdot}
\bigl(2\sqrt{2}\,e^{-\delta_H\mathsf{w}_*}\bigr)^{k-n}
\le\tfrac12(1-\theta)\beta,
\end{split}
\end{equation*}
where the first inequality follows from these three facts and the second from (s1) and (s3);
that is,
\begin{equation}\label{4.4:eq-twist-stability-d}
\pi_n\le(1-\theta)\beta\,2^{-(k+1-n)} .
\end{equation}

\emph{Layers $n<j$.} By \eqref{4.4:eq-twist-stability-d} and the definition of $F_n$ and
$f_{j,n}$ in Definition~\ref{4.4:def-post-step-space},
\begin{equation*}
F_n+\pi_n\le f_{k+1,n}+\beta2^{-(k+1-n)}=f_{k,n}\le f_{j,n},
\end{equation*}
the equality because both sides equal $1-\beta+\beta2^{-(k-n)}$, and the last inequality
because $f_{j,n}$ decreases in $j$ and $j\le k$.

\emph{Layer $n=j$.} It is part of $X\cap\Omega_j$, where the factor at level $j$ is $1$, and
$F_j+\pi_j\le f_{k,j}\le1$ by the same computation.

\emph{Layers $j<n\le k$ and the top layer $n=k+1$.} These points lie in $X\cap\Omega_j$, where at
level $j$ the bounds are in the unit of scale $j$ with factor $1$. A bound of weight $\varpi$ at
scale $n$ converts to scale $j$ with the factor $L^{-\varpi(n-j)}$, and $\varpi\ge1$; for
\eqref{4.4:eq-layered-spaces-d} the additional factor $L^{n-\min(p,n)}$ that the slot carries at
scale $n$ is at most $L^{n-j}$ times the one it carries at scale $j$, and there $\varpi=3$.
Since $F_n\le1+\beta$ and $F_{\mathrm{top}}=1+\beta$, the untwisted part contributes a factor at
most
\begin{equation*}
(1+\beta)\frac{\alpha_{\cdot,n}}{\alpha_{\cdot,j}}L^{-(n-j)}
\le\frac{(1+\beta)(1+\beta_0)\sqrt{2}}{L}\le\frac12,
\end{equation*}
by Lemma~\ref{4.4:lem-radii-ratio} for the pair $(j,n)$, $n\le k+1$, together with
$\sup_{i\ge1}(1+i)^{1/2}L^{-i}=\sqrt{2}/L$, and by (s2). Put $\bar n:=\min(n,k)$; on
$\Omega_{k+1}$ the twist is smooth at scale $k$, and $D_k=0$. In the unit of scale $j$ the twist
adds at most
\begin{equation*}
\begin{split}
c_\flat C_H^2\,\frac{w_ke^{-D_{\bar n}}}{\alpha_{\cdot,j}}\,L^{-(\bar n-j)}
&\le c_\flat C_H^2\,\frac{2\|C\|C_1A_1}{C_\cdot}\,(1+\beta_0)\\
&\qquad\cdot(1+k-\bar n)^{1/2}e^{-\delta_H\mathsf{w}_*(k-\bar n)}
\cdot(1+\bar n-j)^{1/2}L^{-(\bar n-j)}\ \le\ \tfrac14 .
\end{split}
\end{equation*}
The first inequality is Lemma~\ref{4.4:lem-radii-ratio} for the pair $(j,k)$, which gives
\begin{equation*}
\frac{w_k}{\alpha_{\cdot,j}}\le\frac{2\|C\|C_1A_1}{C_\cdot}\,(1+\beta_0)(1+k-j)^{1/2},
\end{equation*}
together with
$D_{\bar n}\ge\delta_H\mathsf{w}_*(k-\bar n)$ and
$(1+k-j)^{1/2}\le(1+k-\bar n)^{1/2}(1+\bar n-j)^{1/2}$, which holds because
$(1+a)(1+b)\ge1+a+b$ for $a,b\ge0$. For the second,
\begin{equation*}
(1+k-\bar n)^{1/2}e^{-\delta_H\mathsf{w}_*(k-\bar n)}
\le\bigl(\sqrt{2}\,e^{-\delta_H\mathsf{w}_*}\bigr)^{k-\bar n}\le1
\end{equation*}
by (s1); $(1+\bar n-j)^{1/2}L^{-(\bar n-j)}\le1$; and
\begin{equation*}
c_\flat C_H^2\,\frac{2\|C\|C_1A_1}{C_\cdot}\,(1+\beta_0)\le\tfrac14(1-\theta)\beta\le\tfrac14
\end{equation*}
by (s3).
The bound does not depend on $k$, and no restriction on the depth $k-j$ is needed. The total is
at most $\frac12+\frac14<1$.

This proves conclusion (1) in all slots except the comparison on cubes for
\eqref{4.4:eq-layered-spaces-e}; that comparison and the proofs of conclusions (1$'$), (2) and
$(2^{\mathrm{int}})$ are given in the statements that follow.
\end{proof}

\begin{proof}[Proof of conclusions (1) and (1$'$) of Lemma~\ref{4.4:lem-twist-stability} at the gauge condition: the gauge condition by extension of cubes]\label{4.4:prf-cube-extension}
\emph{Hypotheses.} This part of the proof of the stability of the layered spaces under the twist (Lemma~\ref{4.4:lem-twist-stability}) assumes the following:
\begin{itemize}
\item the hypotheses of Lemma~\ref{4.4:lem-twist-stability}, including the ceiling (s0) on $\gamma$ and the floors (s1)--(s3) of \eqref{4.4:eq-twist-stability-a};
\item the induced bounds for the group-valued factor, \eqref{4.4:eq-group-factor-bounds-a}.
\end{itemize}
The notation is that of Lemma~\ref{4.4:lem-twist-stability}. In particular $\Omega_1\supset\Omega_2\supset\dots\supset\Omega_{k+1}$ are the members of the admissible sequence at level $k+1$ on $Y$. The old term has level $j$ and admissible domain $X\subset Y$. We write $\mathbf{U}=U'U$ as in Ba{\l}aban's layered analyticity spaces (Definition~\ref{4.4:def-layered-spaces}).

By the \emph{level-$(k+1)$ class of index $n$} we mean the cubes over which \eqref{4.4:eq-layered-spaces-e} at level $k+1$ quantifies for the index $n$: the lattice cubes $\square$ with the following properties:
\begin{itemize}
\item $\square$ has side $CML^n\xi$, where $C\in\{1,2\}$ is the positive integer of \eqref{4.4:eq-layered-spaces-e};
\item $\square\subset\Omega_n$;
\item $\square\cap\Omega^c_{n+1}\neq\emptyset$.
\end{itemize}
For such a cube, $\square\cap\Omega_{n+2}=\emptyset$ by admissibility of the sequence.
Such a cube may or may not be aligned with any partition. The gauge condition \eqref{4.4:eq-layered-spaces-e} asks, for each such cube, for a $G$-valued gauge transformation $u$ on $\square$ and a $\mathfrak{g}$-valued field $A$ obeying the bounds displayed there for $L^n\xi|A|$ and $(L^n\xi)^2|\nabla A|$. The constant $B$ is the absolute constant of \eqref{4.4:eq-layered-spaces-e}.

\emph{Reduction.} The twist does not act on the group-valued factor: $U^{\flat}=U$. Hence it suffices to prove the following implication: if the gauge condition \eqref{4.4:eq-layered-spaces-e} holds at level $k+1$ on $Y$ with the relaxed factor $(1+\beta)$ of the post-step space \eqref{4.4:eq-post-step-space-a}, then it holds at level $j$ on $X$ with factor $1$.

For indices $n\le j-1$ the cubes, their sizes and the bounds required at level $j$ on $X$ are literally the same as those at level $k+1$ on $Y$. Moreover $\square\cap X\subset\square\cap Y$ for every cube $\square$. So for these indices there is nothing to prove.

\emph{The index $j$.} Let $\square\subset\Omega_j$ have side $CML^j\xi$, $C\in\{1,2\}$. Put
\begin{equation*}
n_0:=\min\{n:\ \square\cap(\Omega_n\setminus\Omega_{n+1})\neq\emptyset\},
\end{equation*}
and $n_0:=k+1$ if $\square\subset\Omega_{k+1}$. Since $\Omega_n\setminus\Omega_{n+1}$ is disjoint from $\Omega_j\supset\square$ for $n<j$, we have $n_0\ge j$.

We first show that $\square\subset\Omega_{n_0}$. We have $\square\subset\Omega_j$. For $j\le n<n_0$ the set $\square\cap(\Omega_n\setminus\Omega_{n+1})$ is empty, by the minimality of $n_0$. Hence $\square\subset\Omega_n$ implies $\square\subset\Omega_{n+1}$ for each such $n$. Induction on $n$ from $n=j$ gives $\square\subset\Omega_{n_0}$.

\emph{The case $n_0=j$.} Then $\square$ is itself a cube of the level-$(k+1)$ class of index $j$:
\begin{itemize}
\item $\square\subset\Omega_j$;
\item $\square\cap\Omega^c_{j+1}\neq\emptyset$;
\item $\square\cap\Omega_{j+2}=\emptyset$, by admissibility of the sequence.
\end{itemize}
The same cube therefore carries the same bound $BCM\alpha_{0,j}$, and the condition at level $j$ holds for $\square$.

\emph{The case $n_0>j$.} The same condition is stated in \cite[below (1.69)]{B-CMP122a}, where $C$ is a positive integer which, because of geometric constraints on the cubes, may be assumed to satisfy $C\le 2$. In particular the value $1$ is admitted, and we read the level-$(k+1)$ condition at index $n_0$ for cubes with integer $1$.

Put $S':=ML^{n_0}\xi$. This is the mesh of the standard $M$-partition at scale $n_0$. Every member $\Omega_{n_0}$ of an admissible sequence is built in one of two ways:
\begin{itemize}
\item by the operation $\mathbf{T}$, as a union of $MR_{n_0}$-cubes \cite{B-CMP119};
\item by the operation $\mathbf{R}$, as a union of $M$-cubes.
\end{itemize}
By the compatibility of partitions, in either case $\Omega_{n_0}$ is a union of cubes of a partition of mesh $S''\ge S'$ compatible with the standard one.

\emph{Claim.} There is a cube $\square'\supset\square$ of side $S'$ in the level-$(k+1)$ class of index $n_0$.

\emph{Proof of the claim.} Fix a coordinate $\mu$ and let $I_\mu$ be the interval of $\square$ in that coordinate. Its length satisfies
\begin{equation*}
CML^j\xi\le 2ML^j\xi\le ML^{j+1}\xi\le S',
\end{equation*}
because $L\ge2$ and $n_0\ge j+1$. Since $S''\ge S'$, the interval $I_\mu$ meets the interiors of one or two intervals of the partition of mesh $S''$.
\begin{itemize}
\item If it meets the interior of one, extend $I_\mu$ to an interval $I'_\mu\supset I_\mu$ of length $S'$ inside that partition interval. This is possible because $S'\le S''$.
\item If it meets the interiors of two, extend $I_\mu$ to an interval $I'_\mu\supset I_\mu$ of length $S'$ inside their union, which has length $2S''\ge S'$, in such a way that both interiors are still met.
\end{itemize}
Put $\square':=\prod_\mu I'_\mu$. In each coordinate, $I'_\mu$ lies in the union of the partition intervals whose interiors $I_\mu$ meets, and meets the interiors of all of them. Hence $\square'$ meets the interiors of the same partition cubes as $\square$, and lies in their union.

Now $\Omega_{n_0}$ is a union of cubes of this partition and contains $\square$. So every partition cube whose interior meets $\square$ is one of the cubes making up $\Omega_{n_0}$, and therefore $\square'\subset\Omega_{n_0}$. Next,
\begin{equation*}
\square'\cap\Omega^c_{n_0+1}\ \supset\ \square\cap\Omega^c_{n_0+1}\ \neq\ \emptyset ,
\end{equation*}
and, when $n_0$ is given by the minimum, the right side is non-empty because $\square\cap(\Omega_{n_0}\setminus\Omega_{n_0+1})\neq\emptyset$.

Finally, $\square'\cap\Omega_{n_0+2}=\emptyset$ by admissibility. Indeed, by \cite[(2.2)]{B-CMP96},
\begin{equation*}
\operatorname{dist}(\Omega^c_{n_0+1},\Omega_{n_0+2})>R^{\mathrm{adm}}ML^{n_0+1}\xi ,
\end{equation*}
where $R^{\mathrm{adm}}$ denotes the constant of that display. This distance exceeds $S'$. Since $\square'$ meets $\Omega^c_{n_0+1}$ and its diameter in the maximum metric is $S'$, which is smaller than that distance, $\square'$ cannot meet $\Omega_{n_0+2}$.

The proof of the claim uses only $S''\ge S'$: in the first case $S'\le S''$, in the second $S'\le 2S''$. This ends the proof of the claim.

\emph{The bound for $\square$.} Let $u$ be the gauge transformation of $\square'$ and $A$ the associated $\mathfrak{g}$-valued field given by the level-$(k+1)$ condition at index $n_0$, read for the cube $\square'$, whose integer is $1$. In the post-step space \eqref{4.4:eq-post-step-space-a} this condition is the bound
\begin{equation*}
L^{n_0}\xi|A|<(1+\beta)BM\alpha_{0,n_0}.
\end{equation*}
Restrict $u$ and $A$ to $\square$. Then
\begin{align*}
L^j\xi|A| &= L^{-(n_0-j)}L^{n_0}\xi|A|
 < L^{-(n_0-j)}(1+\beta)BM\alpha_{0,n_0}\\
&\le BM\alpha_{0,j}\ \le\ BCM\alpha_{0,j}.
\end{align*}
The second inequality follows from the ratio of radii along the coupling flow (Lemma~\ref{4.4:lem-radii-ratio}), applied to the pair of levels $(j,n_0)$, which is admissible there because $n_0\le k+1$, together with the floor (s2) of \eqref{4.4:eq-twist-stability-a}. The last inequality holds because $C\ge1$.

The bound on $(L^j\xi)^2|\nabla A|$ required by \eqref{4.4:eq-layered-spaces-e} follows from the same chain. One writes $(L^j\xi)^2=L^{-2(n_0-j)}(L^{n_0}\xi)^2$ and uses the factor $L^{-2(n_0-j)}$, which is at most $L^{-(n_0-j)}$, in place of $L^{-(n_0-j)}$.

We remark that for the real backgrounds of Ba{\l}aban's construction the conclusion holds in any case. On $\Omega_j\setminus\Omega_{j+1}$ and inside it, the configuration $U_{k+1}$ satisfies at least the regularity bounds required of the level-$j$ configuration $U_j$ \cite{B-CMP119}, so the gauge condition at level $j$ holds for it, and $U^{\flat}=U$ preserves this.

Together with the bounds established in the preceding part of the proof of Lemma~\ref{4.4:lem-twist-stability}, this proves conclusion (1) of \eqref{4.4:eq-twist-stability-c}.

\emph{Conclusion (1$'$).} Here $U^{\flat}=U$. Let $X''\subset X\subset Y$ be admissible and let $p$ range over the indices of the determining sets. Then the quantities
\begin{equation*}
U\bigl(\mathbf{B}_p(X'')\cup\{\Gamma_i\}_{i<p},\,M^{\cdot}U\bigr),\qquad
\mathbf{J}_{p,X''}(M^{\cdot}U)
\end{equation*}
do not move. The normalised increments therefore vanish: $\pi_n=0$ for every $n$.

The three layer estimates in the proof of conclusion (1) then read, without the twist, as follows:
\begin{itemize}
\item for the layers $n<j$,
\begin{equation*}
F_n\le f_{k,n}\le f_{j,n};
\end{equation*}
\item for the layer $n=j$, $F_j\le 1$;
\item for the layers $j<n\le k+1$, the new factor is at most $\tfrac12$.
\end{itemize}
This proves conclusion (1$'$).
\end{proof}

\begin{definition}[Gauges on lattice boxes and estimates near the identity]
\label{4.4:not-box-conventions}
In what follows gauge transformations are constructed on lattice cubes of several sides and
compared near the identity of the structure group; we fix the conventions used for this.
\begin{enumerate}
\item Lengths are measured in the unit in which $L^j\xi = 1$, where $\xi \le 1$ is the
lattice spacing and $j$ is the level under consideration.
\item $|\cdot|$ and $\|\cdot\|$ denote two norms on $N\times N$ matrices, of which only the
following properties are used: the norm
$|\cdot|$ is unitarily invariant, and $|XYZ| \le \|X\|\,|Y|\,\|Z\|$ for all $N\times N$
matrices $X$, $Y$, $Z$; the norm $\|\cdot\|$ is submultiplicative with $\|1\| = 1$; and
$\|X\| \le \sqrt{N}\,|X|$ for every $N\times N$ matrix $X$.
\item Let $S > 0$ with $S/\xi$ a positive integer. A \emph{lattice cube of side $S$} is a
product of $d$ runs, each consisting of $S/\xi$ consecutive lattice sites. The bonds of a
lattice cube are the bonds with both end points in it.
\item Write $e_\mu$ for the unit vector in direction $\mu$. The symbol $\nabla$ denotes the
plain, non-covariant difference quotient: $\nabla_\mu f(x) = \xi^{-1}(f(x+\xi e_\mu) - f(x))$
between neighbouring sites. For a function on bonds it is the analogous quotient between a
bond and its neighbouring parallel bond.
\end{enumerate}
The following hold.
\begin{enumerate}
\item[(e1)] For Hermitian matrices $X$, $Y$ and unitary matrices $a$, $b$, $a'$, $b'$, the
three inequalities in the first two rows of \eqref{4.4:eq-box-conventions-a} hold.
\item[(e2)] Let $G$ be a closed subgroup of $\mathrm{U}(N)$, for instance $\mathrm{SU}(N)$,
and write
$\mathfrak{g}$ for its Lie algebra, realised as the space of Hermitian matrices $X$ with
$e^{itX} \in G$ for all real $t$. There is $r_L = r_L(G) \in (0, 1/(2\sqrt{N})]$ with the
following property. Every
$g \in G$ with $|g - 1| \le r_L$ is of the form $g = e^{iX}$ with
$X := \frac{1}{i}\log g \in \mathfrak{g}$, where $\log$ is the principal branch. Moreover, the
third row of \eqref{4.4:eq-box-conventions-a} holds for such $g$, and the fourth row holds
for any two such $g$, $g'$.
\end{enumerate}
The estimates of (e1) and (e2) read, in order,
\begin{equation}\label{4.4:eq-box-conventions-a}
\begin{gathered}
|e^{iX} - 1| \le |X|, \qquad |e^{iX} - e^{iY}| \le |X - Y|, \\
|ab - a'b'| \le |a - a'| + |b - b'|, \\
|X| \le \tfrac{\pi}{2}\,|g - 1|, \\
|\log g - \log g'| \le 2\,|g - g'| .
\end{gathered}
\end{equation}
\end{definition}

\begin{proof}
\emph{(e1).} By the spectral theorem, $X = \Phi\operatorname{diag}(\eta_1,\dots,\eta_N)\Phi^*$
with $\Phi$ unitary and $\eta_i$ real. Hence
\begin{equation*}
e^{iX} - 1 = \Phi\operatorname{diag}(e^{i\eta_1} - 1, \dots, e^{i\eta_N} - 1)\Phi^* .
\end{equation*}
Put $c_i := (e^{i\eta_i} - 1)/\eta_i$ if $\eta_i \ne 0$ and $c_i := 0$ if $\eta_i = 0$; then
$e^{i\eta_i} - 1 = c_i \eta_i$ and $|c_i| \le 1$, since
$|e^{i\eta} - 1| = 2|\sin(\eta/2)| \le |\eta|$ for real $\eta$.
Every complex number $c$ with $|c| \le 1$ is the mean of two numbers of modulus one: writing
$c = \cos \tau\, e^{i\varphi}$ with $\tau \in [0, \pi/2]$, one has
$c = \frac12\bigl(e^{i(\varphi + \tau)} + e^{i(\varphi - \tau)}\bigr)$. Hence
$\operatorname{diag}(c_1, \dots, c_N) = \frac12(\Xi_1 + \Xi_2)$ with $\Xi_1$, $\Xi_2$ diagonal
unitary, and
\begin{equation*}
e^{iX} - 1 = \tfrac12\,(\Phi \Xi_1 \Phi^*)\,X + \tfrac12\,(\Phi \Xi_2 \Phi^*)\,X .
\end{equation*}
The matrices $\Phi \Xi_1 \Phi^*$ and $\Phi \Xi_2 \Phi^*$ are unitary, so unitary invariance
and the triangle inequality
give $|e^{iX} - 1| \le \frac12|X| + \frac12|X| = |X|$.

For the second inequality, the derivative of $s \mapsto e^{isX}e^{i(1-s)Y}$ is
$e^{isX}(iX - iY)e^{i(1-s)Y}$. Integrating over $[0,1]$ gives
\begin{equation*}
e^{iX} - e^{iY} = i\int_0^1 e^{isX}(X - Y)\,e^{i(1-s)Y}\,ds .
\end{equation*}
The factors $e^{isX}$ and $e^{i(1-s)Y}$ are unitary. By unitary invariance the integrand has
norm $|X - Y|$ for every $s$, and the triangle inequality for the integral gives
$|e^{iX} - e^{iY}| \le |X - Y|$.

For unitaries, write $ab - a'b' = (a - a')b + a'(b - b')$. Unitary invariance gives
$|(a - a')b| = |a - a'|$ and $|a'(b - b')| = |b - b'|$, whence the third inequality.

\emph{(e2).} A closed subgroup $G$ of $\mathrm{U}(N)$ is an embedded Lie subgroup. Hence
$X \mapsto e^{iX}$ maps a neighbourhood of $0$ in $\mathfrak{g}$ onto a neighbourhood of
$1$ in $G$. We take that neighbourhood of $0$ inside the set of $X$ whose
eigenvalues lie in $(-\pi, \pi)$. On such $X$ the principal logarithm inverts the exponential:
$\log e^{iX} = iX$, computed on the eigenvalues. Choose $r_L \le 1/(2\sqrt{N})$ so small that
$\{g \in G : |g - 1| \le r_L\}$ lies in the image neighbourhood. Then every such $g$
equals $e^{iX}$ with $X = \frac{1}{i}\log g \in \mathfrak{g}$.

The eigenvalues $\eta$ of $X$ satisfy $|\eta| \le \pi$, and
$|e^{i\eta} - 1| = 2|\sin(\eta/2)|$. Jordan's inequality $\sin u \ge 2u/\pi$ for
$0 \le u \le \pi/2$ therefore gives $|\eta| \le \frac{\pi}{2}|e^{i\eta} - 1|$ for each
eigenvalue. Write, with $\Phi$ unitary,
\begin{equation*}
\begin{gathered}
X = \Phi\operatorname{diag}(\eta_1, \dots, \eta_N)\Phi^*, \\
g - 1 = \Phi\operatorname{diag}(e^{i\eta_1} - 1, \dots, e^{i\eta_N} - 1)\Phi^* .
\end{gathered}
\end{equation*}
Since $e^{i\eta_i} = 1$ forces $\eta_i = 0$ on $(-\pi, \pi)$, we may write
$\eta_i = \frac{\pi}{2}\,c_i'\,(e^{i\eta_i} - 1)$ with $|c_i'| \le 1$ (and $c_i' := 0$ if
$\eta_i = 0$). Decomposing $\operatorname{diag}(c_1', \dots, c_N')$ as the mean
$\frac12(\Xi_1' + \Xi_2')$ of two diagonal unitaries $\Xi_1'$, $\Xi_2'$, by the device used
in (e1), we obtain
\begin{equation*}
X = \tfrac{\pi}{2}\Bigl(\tfrac12\,(\Phi \Xi_1' \Phi^*)(g - 1)
+ \tfrac12\,(\Phi \Xi_2' \Phi^*)(g - 1)\Bigr),
\end{equation*}
and unitary invariance with the triangle inequality yields
$|X| \le \frac{\pi}{2}|g - 1|$.

For the Lipschitz bound, put $z = g - 1$ and $z' = g' - 1$. Since
$|z|, |z'| \le r_L$ and $r_L \le 1/(2\sqrt{N})$, the comparison of the two norms in item 2
gives
\begin{equation*}
\|z\| \le \sqrt{N}\,|z| \le \sqrt{N}\,r_L \le \tfrac12 ,
\end{equation*}
and likewise $\|z'\| \le \frac12$. On this range the principal logarithm is given by the power
series, so
\begin{equation*}
\log(1 + z) - \log(1 + z')
= \sum_{n \ge 1} \frac{(-1)^{n-1}}{n}\,(z^n - z'^n) .
\end{equation*}
Now $z^n - z'^n$ is the sum of the $n$ products $z^{n_1}(z - z')z'^{\,n_2}$ with
$n_1, n_2 \ge 0$ and $n_1 + n_2 = n - 1$. Apply
$|XYZ| \le \|X\|\,|Y|\,\|Z\|$ to each with $\|z^{n_1}\| \le \|z\|^{n_1}$ and likewise for
$z'$ (submultiplicativity of
$\|\cdot\|$, and $\|1\| = 1$ for a zero exponent). This gives
$|z^n - z'^n| \le n\,2^{-(n-1)}|z - z'|$. Summing,
\begin{equation*}
|\log g - \log g'| \le \sum_{n \ge 1} 2^{-(n-1)}\,|z - z'| = 2\,|g - g'| .
\qedhere
\end{equation*}
\end{proof}

\begin{proposition}[An equivariant smooth retraction onto the gauge group; classical, cited by statement]
\label{4.4:prop-equivariant-retraction}
Let $G\subset M_N(\mathbb{C})$ denote the gauge group. With the norm $|\cdot|$ fixed in the
conventions~\ref{4.4:not-box-conventions} on gauges on lattice boxes and estimates near the
identity, $G$ is a compact $C^\infty$ submanifold of the
Euclidean space
\begin{equation*}
(M_N(\mathbb{C}),|\cdot|)\cong\mathbb{R}^{2N^2}.
\end{equation*}
There are a radius $r_P=r_P(G)\in(0,r_L]$, where $r_L$ is the radius of the logarithm near the
identity, a constant $\Lambda=\Lambda(G)\ge 1$, and a $C^\infty$ map
\begin{equation*}
P\colon\{Z\in M_N(\mathbb{C}) : \operatorname{dist}(Z,G)<2r_P\}\longrightarrow G
\end{equation*}
onto $G$ with the following three properties: for all $g,h,h'\in G$, all $Z$ in the domain of $P$,
and all $H,H'\in M_N(\mathbb{C})$, where in the two derivative bounds $Z$ ranges over the ball
$\{Z : |Z-1|\le r_P\}$,
\begin{equation}
\label{4.4:eq-equivariant-retraction-a}
\begin{gathered}
P(g)=g,\qquad P(hZh')=h\,P(Z)\,h',\\
|DP(Z)[H]|\le\Lambda|H|,\\
|D^2P(Z)[H,H']|\le\Lambda|H|\,|H'|.
\end{gathered}
\end{equation}
\end{proposition}

\begin{remark}
The statement is the classical theorem on the existence and smoothness of the nearest-point
retraction of a tubular neighbourhood of a compact submanifold of a Euclidean space. The map $P$ may be
taken to be the nearest-point map: for a compact smooth submanifold of a Euclidean space, the point
of the submanifold nearest to $Z$ is unique and depends smoothly on $Z$ in a neighbourhood of the
submanifold. For $h,h'\in G$ the map $Z\mapsto hZh'$ is an isometry of
$(M_N(\mathbb{C}),|\cdot|)$ that preserves $G$; it therefore carries nearest points to nearest
points, and this is the equivariance in \eqref{4.4:eq-equivariant-retraction-a}. The bounds on the
first and second derivatives hold by compactness. For $G=U(N)$ one may take
\begin{equation*}
P(Z)=Z\,(Z^*Z)^{-1/2},
\end{equation*}
and for $G=SU(N)$ this map divided by the $N$-th root near $1$ of its determinant; either choice
serves equally well, since only the three properties displayed in
\eqref{4.4:eq-equivariant-retraction-a} are used in what follows.
\end{remark}

\begin{definition}[Admissible covers of a lattice box]\label{4.4:def-admissible-covers}
Lengths are measured as in the conventions of~\ref{4.4:not-box-conventions}: the unit of length is
$L^j\xi=1$, the lattice spacing is $\xi\le 1$, and the length $|J|$ of a run $J$ of consecutive
sites is $\xi$ times the number of its sites. A run of sites is written
$[p,p+l)$, meaning the sites $t$ with $p\le t<p+l$; it has length $l$.

Let $S\ge 80\xi$ and $S\ge 1$, and let $\square^*=\prod_\mu J_\mu$ be a lattice box, that is, a
product over the $d$ directions $\mu$ of runs $J_\mu$ of consecutive sites. Put
\begin{equation*}
S^*:=\max\Bigl\{S,\ \max_\mu |J_\mu|\Bigr\}.
\end{equation*}
An \emph{$S$-admissible cover} of $\square^*$ is, for each direction $\mu$, a family of lattice runs
$I^{(\mu)}_i=[p_i,p_i+l_i)\subset J_\mu$, $0\le i\le n_\mu$, such that
\begin{equation}\label{4.4:eq-admissible-covers-a}
\begin{aligned}
&\text{(ac1)}\quad p_0=\min J_\mu,\qquad p_{n_\mu}+l_{n_\mu}=\max J_\mu+\xi,\\
&\phantom{\text{(ac1)}}\quad\ \text{and both } i\mapsto p_i \text{ and } i\mapsto p_i+l_i
  \text{ are strictly increasing};\\
&\text{(ac2)}\quad l_i\le S;\\
&\text{(ac3)}\quad p_i+l_i-p_{i+1}\ge S/5\qquad (0\le i<n_\mu);\\
&\text{(ac4)}\quad I^{(\mu)}_i\cap I^{(\mu)}_{i'}=\emptyset\qquad (|i-i'|\ge 3);\\
&\text{(ac5)}\quad (n_\mu+1)\,S\le 4S^*.
\end{aligned}
\end{equation}
Thus by (ac3) consecutive runs overlap in length at least $S/5$. The boxes of the cover are
\begin{equation*}
c_m:=\prod_\mu I^{(\mu)}_{m_\mu},\qquad m\in\prod_\mu\{0,\dots,n_\mu\}.
\end{equation*}
\end{definition}

An $S$-admissible cover has the following properties.
\begin{enumerate}
\item[(a)] (Betweenness.) If $t\in I^{(\mu)}_i\cap I^{(\mu)}_{i'}$ and $i''$ lies between $i$
and $i'$, then $t\in I^{(\mu)}_{i''}$.
\item[(b)] If $n_\mu\ge 1$, then $l_i\ge S/5$ for $0\le i\le n_\mu$.
\item[(c)] For every index $m$, the at most $2^d$ boxes $c_{m'}$ with
$m'_\mu\in\{m_\mu,m_\mu+1\}\cap\{0,\dots,n_\mu\}$ for every $\mu$ have a common site; this is
the common site of the $2^d$ boxes round a nerve plaquette, in the language of the gluing
construction of this chapter.
\item[(d)] If $|J_\mu|<S$, the single run $I^{(\mu)}_0=J_\mu$, with $n_\mu=0$, satisfies
(ac1)--(ac5) in the direction $\mu$; a box thinner than $S$ in a direction is thus admitted,
with $n_\mu=0$ there.
\end{enumerate}

\begin{proof}
Fix $\mu$ and write $I_i$ for $I^{(\mu)}_i$.

(a) Let $i<i''<i'$ (the case $i'<i''<i$ is symmetric). By (ac1), $p_{i''}<p_{i'}\le t$ and
$t<p_i+l_i<p_{i''}+l_{i''}$, so $p_{i''}\le t<p_{i''}+l_{i''}$, that is, $t\in I_{i''}$.

(b) For $0\le i<n_\mu$, (ac1) gives $p_{i+1}>p_i$, hence by (ac3)
\begin{equation*}
l_i=(p_i+l_i)-p_i> (p_i+l_i)-p_{i+1}\ge S/5 .
\end{equation*}
For $i=n_\mu\ge 1$, (ac1) gives $p_{n_\mu}+l_{n_\mu}>p_{n_\mu-1}+l_{n_\mu-1}$, hence by (ac3)
\begin{equation*}
l_{n_\mu}>(p_{n_\mu-1}+l_{n_\mu-1})-p_{n_\mu}\ge S/5 .
\end{equation*}

(c) For $0\le i<n_\mu$, by (ac1) the intersection $I_i\cap I_{i+1}$ is the run
$[p_{i+1},p_i+l_i)$, whose length is at least $S/5\ge 16\xi$ by (ac3); so it contains a site.
For each $\mu$, choose a site $t_\mu$ in $I_{m_\mu}\cap I_{m_\mu+1}$ if $m_\mu<n_\mu$, and in
$I_{m_\mu}$ if $m_\mu=n_\mu$. Then $t_\mu\in I_{m'_\mu}$ for every admissible choice of
$m'_\mu$, and the site $(t_\mu)_\mu$ lies in every box $c_{m'}$ of (c).

(d) With $n_\mu=0$, $p_0=\min J_\mu$ and $l_0=|J_\mu|$ one has
$p_0+l_0=\max J_\mu+\xi$, and the monotonicity in (ac1) is vacuous; (ac2) is $|J_\mu|\le S$;
(ac3) and (ac4) are vacuous; (ac5) reads $S\le 4S^*$, which holds since $S\le S^*$.
\end{proof}

\begin{remark}\label{4.4:lem-gluing-boxes}
The statement of this lemma is not written out here; parts of its proof are printed below.
\end{remark}

\begin{proof}[Proof of Lemma~\ref{4.4:lem-gluing-boxes}, concluded]\label{4.4:prf-glued-gauge}
We keep the notation introduced so far in this proof. Thus $\mathsf{k} = S^*/S \ge 1$, and
$\{\chi_m\}$ is the partition of unity built from the cover: $\sum_m \chi_m = 1$ on $\square^*$, at
most $3^d$ of the $\chi_m$ are non-zero at any site, and $|\nabla\chi_m| \le c_\psi/S$,
$|\nabla\nabla\chi_m| \le c_\psi/S^2$. The support property in the corresponding display says
that if $\chi_m(x) \ne 0$, every site of $\square^*$ within sup-distance $4\xi$ of $x$ lies in $c_m$.
The gauges $\tilde u_m = G_m u_m$ and potentials $\tilde A_m = \operatorname{Ad}_{G_m} A_m$ are
rotated by the constants $G_m$ of the axial gauge on the nerve of the cover; thus
$U^{\tilde u_m} = e^{i\xi \tilde A_m}$ and $|\tilde A_m|, |\nabla \tilde A_m| \le \mathsf{a}$ on $c_m$.
Their transitions $\tilde g_{mm'} = \tilde u_m \tilde u_{m'}^{-1}$ satisfy the transition identity in
the corresponding display, which expresses $\tilde g_{mm'}(b_+)$ as
$e^{-i\xi \tilde A_m(b)}\tilde g_{mm'}(b_-)e^{i\xi \tilde A_{m'}(b)}$ for every bond $b$ of
$c_m \cap c_{m'}$, and the bound in the corresponding display, namely
$|\tilde g_{mm'}(x) - 1| \le \bar\varphi$ at every site $x$ of $c_m \cap c_{m'}$, where
$\bar\varphi = 64 d^3 S^* \mathsf{a}$.
For a bond $b$ from $b_-$ to $b_+$ we have $U^u(b) = u(b_-)U(b)u(b_+)^{-1}$.

Throughout, $S^*\mathsf{a} \le r_{\mathrm{gl}}$, with $r_{\mathrm{gl}}$ the constant fixed in
\eqref{4.4:eq-glued-gauge-7} below. Since $S^* \ge S \ge 1$, this gives
$\mathsf{a} \le S^*\mathsf{a} \le 1$ and $\mathsf{k}\mathsf{a} \le S^*\mathsf{a} \le 1$. Since
$r_{\mathrm{gl}} \le r_P/(64d^3)$, it also gives $\bar\varphi \le r_P$.

\emph{The glued gauge.} For $x \in \square^*$ put
\begin{equation}\label{4.4:eq-glued-gauge-3}
u(x) := P\Bigl(\sum_m \chi_m(x)\,\tilde u_m(x)\Bigr).
\end{equation}
This is well defined. If $\chi_m(x) \ne 0$, then $x \in c_m$ by the support property, so
$\tilde u_m(x)$ is defined. For two such indices $m, m'$, unitary invariance of $|\cdot|$ and
the corresponding display give
\begin{equation*}
|\tilde u_m(x) - \tilde u_{m'}(x)| = |\tilde g_{mm'}(x) - 1| \le \bar\varphi .
\end{equation*}
Fix one such $m'$. Since $\sum_m \chi_m(x) = 1$, the convex combination
$\sum_m \chi_m(x)\tilde u_m(x)$ lies within $\bar\varphi \le r_P$ of $\tilde u_{m'}(x) \in G$. Hence it
lies in the domain of the retraction $P$ of Proposition~\ref{4.4:prop-equivariant-retraction}, and
$u(x) \in G$.

Fix $x_0 \in \square^*$ and an index $m_0$ with $\chi_{m_0}(x_0) \ne 0$. Let $\mathcal{N}$ be the set
of sites of $\square^*$ within sup-distance $2\xi$ of $x_0$.

First, $\mathcal{N} \subset c_{m_0}$, by the support property at $x_0$. Next, suppose $\chi_m$ does not
vanish identically on $\mathcal{N}$, say $\chi_m(y) \ne 0$ with $y \in \mathcal{N}$. Every site of
$\mathcal{N}$ is within sup-distance $4\xi$ of $y$, so $\mathcal{N} \subset c_m$ by the support
property at $y$.

Sums over $m$ on $\mathcal{N}$ are taken over such indices only; the others contribute zero. On
$\mathcal{N}$ put
\begin{align*}
\tilde g_m &:= \tilde g_{m m_0}, \qquad
\mathsf{Z} := \sum_m \chi_m \tilde g_m = 1 + \sum_m \chi_m(\tilde g_m - 1), \\
p &:= P(\mathsf{Z}).
\end{align*}
The second expression for $\mathsf{Z}$ uses $\sum_m \chi_m = 1$. By the corresponding display,
$|\mathsf{Z} - 1| \le \bar\varphi \le r_P$, so $p$ is defined and $G$-valued.

Since $\tilde u_m = \tilde g_m \tilde u_{m_0}$ on $\mathcal{N}$, we have
$\sum_m \chi_m\tilde u_m = \mathsf{Z}\,\tilde u_{m_0}$ there. The equivariance
$P(Zh') = P(Z)h'$ for $h' \in G$ then gives $u = p\,\tilde u_{m_0}$ on $\mathcal{N}$.

Let $b$ be a bond of $\mathcal{N}$. Then
\begin{equation*}
U^u(b) = p(b_-)\,\tilde u_{m_0}(b_-)U(b)\tilde u_{m_0}(b_+)^{-1}\,p(b_+)^{-1}.
\end{equation*}
Since $U^{\tilde u_{m_0}} = e^{i\xi\tilde A_{m_0}}$, this reads
\begin{equation}\label{4.4:eq-glued-gauge-a}
U^u(b) = p(b_-)\, e^{i\xi \tilde A_{m_0}(b)}\, p(b_+)^{-1}.
\end{equation}

\emph{Bounds.} We first bound the transitions on $\mathcal{N}$.

By the corresponding display, $|\tilde g_m - 1| \le \bar\varphi$.

For the first difference quotients, let $\tilde A_{m,\mu}(x)$ denote the value of $\tilde A_m$ on the
bond from $x$ to $x + \xi e_\mu$. For $X, Y$ Hermitian and $g$ a matrix put
\begin{equation*}
\Gamma(X,Y,g) := \xi^{-1}\bigl(e^{-i\xi X} g\, e^{i\xi Y} - g\bigr).
\end{equation*}
By the transition identity in the corresponding display,
\begin{equation*}
(\nabla_\mu \tilde g_m)(x)
= \Gamma\bigl(\tilde A_{m,\mu}(x), \tilde A_{m_0,\mu}(x), \tilde g_m(x)\bigr).
\end{equation*}
The map $\Gamma$ is linear in $g$, and it has two properties.

The first property is
\begin{equation*}
|\Gamma(X,Y,h)| \le (\|X\| + \|Y\|)\,|h|.
\end{equation*}
To see this, write
\begin{equation*}
\xi\,\Gamma(X,Y,h) = (e^{-i\xi X} - 1)\,h\, e^{i\xi Y} + h\,(e^{i\xi Y} - 1),
\end{equation*}
and use $|XYZ| \le \|X\|\,|Y|\,\|Z\|$ together with the estimates
\eqref{4.4:eq-box-conventions-a} near the identity.

The second property holds for unitary $g$:
\begin{equation*}
|\Gamma(X',Y',g) - \Gamma(X,Y,g)| \le |X' - X| + |Y' - Y| .
\end{equation*}
Indeed, $\xi$ times the left side is
$|e^{-i\xi X'} g e^{i\xi Y'} - e^{-i\xi X} g e^{i\xi Y}|$. By the estimates
\eqref{4.4:eq-box-conventions-a} for products of unitaries and for differences of exponentials, this is
at most $\xi|X'-X| + \xi|Y'-Y|$.

Since $\Gamma(0,0,g) = 0$, the second property gives $|\Gamma(X,Y,g)| \le |X| + |Y|$ for unitary
$g$. As $\tilde g_m$ is unitary and $|\tilde A_m|, |\tilde A_{m_0}| \le \mathsf{a}$, this gives
$|\nabla \tilde g_m| \le 2\mathsf{a}$; the first property is used only for the second quotient.
For the second quotient, write $x' = x + \xi e_\nu$ and
split
\begin{align*}
\xi\,(\nabla_\nu\nabla_\mu \tilde g_m)(x)
&= \Gamma\bigl(\tilde A_{m,\mu}(x'), \tilde A_{m_0,\mu}(x'), \tilde g_m(x') - \tilde g_m(x)\bigr)\\
&\quad + \Gamma\bigl(\tilde A_{m,\mu}(x'), \tilde A_{m_0,\mu}(x'), \tilde g_m(x)\bigr)
- \Gamma\bigl(\tilde A_{m,\mu}(x), \tilde A_{m_0,\mu}(x), \tilde g_m(x)\bigr).
\end{align*}
The first term is at most $2\sqrt{N}\mathsf{a}\cdot 2\xi\mathsf{a}$ by the first property. The
difference is at most $2\xi\mathsf{a}$ by the second property and $|\nabla\tilde A_m| \le \mathsf{a}$.
Hence, on $\mathcal{N}$,
\begin{equation}\label{4.4:eq-glued-gauge-4}
|\tilde g_m - 1| \le \bar\varphi, \qquad |\nabla \tilde g_m| \le 2\mathsf{a}, \qquad
|\nabla\nabla \tilde g_m| \le 2\mathsf{a} + 4\sqrt{N}\mathsf{a}^2 .
\end{equation}

We now apply Leibniz' rule for difference quotients,
\begin{equation*}
\nabla_\mu(fh)(x) = (\nabla_\mu f)(x)\,h(x+\xi e_\mu) + f(x)\,(\nabla_\mu h)(x),
\end{equation*}
to $\mathsf{Z} - 1 = \sum_m \chi_m(\tilde g_m - 1)$. A first quotient involves two sites, so at most
$2\cdot 3^d$ indices $m$ contribute to it. The terms with $\nabla\chi_m$ are bounded by
$(c_\psi/S)\bar\varphi$. The terms $\sum_m \chi_m(x)\nabla\tilde g_m(x)$ are bounded by $2\mathsf{a}$,
because $\sum_m \chi_m = 1$.

A second quotient involves four sites, so at most $4\cdot 3^d$ indices contribute to it. Applying the
rule twice produces four kinds of term:
\begin{itemize}
\item $\nabla\nabla\chi_m$ times a value of $\tilde g_m - 1$, bounded by $c_\psi\bar\varphi/S^2$;
\item two terms $\nabla\chi_m$ times $\nabla\tilde g_m$, together bounded by $4c_\psi\mathsf{a}/S$;
\item $\sum_m\chi_m\nabla\nabla\tilde g_m$, bounded by $2\mathsf{a} + 4\sqrt{N}\mathsf{a}^2$.
\end{itemize}
Let $z_1 := \sup_{\mathcal{N}}|\nabla\mathsf{Z}|$ and $z_2 := \sup_{\mathcal{N}}|\nabla\nabla\mathsf{Z}|$,
the suprema over all first, respectively second, difference quotients of $\mathsf{Z}$ with sites in
$\mathcal{N}$. By \eqref{4.4:eq-glued-gauge-4},
\begin{align*}
z_1 &\le 2\cdot 3^d c_\psi\,\bar\varphi/S + 2\mathsf{a}, \\
z_2 &\le 4\cdot 3^d c_\psi\bigl(\bar\varphi/S^2 + 4\mathsf{a}/S\bigr) + 2\mathsf{a}
+ 4\sqrt{N}\mathsf{a}^2 .
\end{align*}
Now $\bar\varphi/S = 64d^3\mathsf{k}\mathsf{a}$, $S \ge 1$, $\mathsf{a} \le 1$ and $\mathsf{k} \ge 1$.
Therefore
\begin{equation}\label{4.4:eq-glued-gauge-5}
z_1,\ z_2 \le C'\mathsf{k}\mathsf{a}, \qquad C' = C'(d,N);
\end{equation}
for instance, since $c_\psi = c_\psi(d)$, one may take
\begin{equation*}
C' := 4\cdot 3^d c_\psi(64 d^3 + 4) + 2 + 4\sqrt{N}.
\end{equation*}

Next we bound $p$. The values of $\mathsf{Z}$ lie in the convex ball $\{|Z - 1| \le r_P\}$ of
Proposition~\ref{4.4:prop-equivariant-retraction}, on which
$|DP(Z)[H]| \le \Lambda|H|$ and $|D^2P(Z)[H,H']| \le \Lambda|H|\,|H'|$.

For neighbouring sites $x, x'$ of $\mathcal{N}$,
\begin{equation*}
p(x') - p(x) = \int_0^1 DP\bigl(\mathsf{Z}(x) + s(\mathsf{Z}(x') - \mathsf{Z}(x))\bigr)
[\mathsf{Z}(x') - \mathsf{Z}(x)]\,ds ,
\end{equation*}
and the segment lies in the ball by convexity. Hence $|\nabla p| \le \Lambda z_1$.

For the second quotient at $x$ in directions $\mu, \nu$, put
\begin{equation*}
H_1 := \nabla_\mu\mathsf{Z}(x), \qquad H_2 := \nabla_\nu\mathsf{Z}(x), \qquad
H_{12} := \nabla_\nu\nabla_\mu\mathsf{Z}(x),
\end{equation*}
and
\begin{equation*}
\Phi(s,t) := P\bigl(\mathsf{Z}(x) + s\xi H_1 + t\xi H_2 + st\xi^2 H_{12}\bigr),
\qquad s,t \in [0,1].
\end{equation*}
The argument of $P$ equals
\begin{equation*}
(1-s)(1-t)\mathsf{Z}(x) + s(1-t)\mathsf{Z}(x+\xi e_\mu) + (1-s)t\,\mathsf{Z}(x+\xi e_\nu)
+ st\,\mathsf{Z}(x+\xi e_\mu+\xi e_\nu).
\end{equation*}
It is therefore a convex combination of values of $\mathsf{Z}$ and lies in the ball. At the four
corners, $\Phi$ takes the values of $p$ at $x$, $x + \xi e_\mu$, $x + \xi e_\nu$ and
$x + \xi e_\mu + \xi e_\nu$. So the mixed second difference of $p$ is
\begin{equation*}
\xi^2(\nabla_\nu\nabla_\mu p)(x) = \int_0^1\!\!\int_0^1 \partial_s\partial_t\Phi\,ds\,dt,
\end{equation*}
where
\begin{equation*}
\partial_s\partial_t\Phi = D^2P\bigl[\xi H_1 + t\xi^2 H_{12},\ \xi H_2 + s\xi^2 H_{12}\bigr]
+ DP[\xi^2 H_{12}].
\end{equation*}
Since $H_1 + t\xi H_{12} = (1-t)\nabla_\mu\mathsf{Z}(x) + t\nabla_\mu\mathsf{Z}(x+\xi e_\nu)$ and
$H_2 + s\xi H_{12} = (1-s)\nabla_\nu\mathsf{Z}(x) + s\nabla_\nu\mathsf{Z}(x+\xi e_\mu)$, both are
convex combinations of first quotients of $\mathsf{Z}$ and have norm at most $z_1$. Hence
$|\nabla\nabla p| \le \Lambda(z_2 + z_1^2)$.

We now bound $A$. By \eqref{4.4:eq-glued-gauge-a},
\begin{equation*}
U^u(b) = e^{i\xi\operatorname{Ad}_{p(b_-)}\tilde A_{m_0}(b)}\; p(b_-)p(b_+)^{-1}.
\end{equation*}
By the estimates \eqref{4.4:eq-box-conventions-a} near the identity, unitary invariance and
$|\tilde A_{m_0}| \le \mathsf{a}$,
\begin{equation*}
|U^u(b) - 1| \le \xi\mathsf{a} + |p(b_-) - p(b_+)| \le \xi(\mathsf{a} + \Lambda z_1).
\end{equation*}
By \eqref{4.4:eq-glued-gauge-5}, $\xi \le 1$ and $\mathsf{k}\mathsf{a} \le S^*\mathsf{a}$,
\begin{equation*}
\xi(\mathsf{a} + \Lambda z_1) \le (1 + \Lambda C')S^*\mathsf{a} \le r_L,
\end{equation*}
because $r_{\mathrm{gl}} \le r_L/(1+\Lambda C')$. By the logarithm estimate near the identity in
\eqref{4.4:eq-box-conventions-a}, $U^u(b) = e^{i\xi A(b)}$ with $A(b) \in \mathfrak{g}$ and
\begin{equation}\label{4.4:eq-glued-gauge-6}
|A| \le \tfrac{\pi}{2}(\mathsf{a} + \Lambda z_1).
\end{equation}

Now let $b' = b + \xi e_\nu$ be a bond of $\mathcal{N}$ parallel and neighbouring to $b$. Let $b$
point in direction $\mu$, write $\tilde A = \tilde A_{m_0}$, and put $x := b_-$, $x' := b'_-$.

For the exponential factors, use
\begin{align*}
&p(x')\tilde A(b) p(x')^{-1} - p(x)\tilde A(b) p(x)^{-1}\\
&\qquad= (p(x') - p(x))\tilde A(b)p(x')^{-1} + p(x)\tilde A(b)\bigl(p(x')^{-1} - p(x)^{-1}\bigr)
\end{align*}
together with $|p(x')^{-1} - p(x)^{-1}| = |p(x) - p(x')|$. With the estimates
\eqref{4.4:eq-box-conventions-a} and $\|\tilde A\| \le \sqrt{N}\mathsf{a}$ this gives
\begin{align*}
\bigl|e^{i\xi\operatorname{Ad}_{p(x')}\tilde A(b')} - e^{i\xi\operatorname{Ad}_{p(x)}\tilde A(b)}\bigr|
&\le \xi\bigl(|\tilde A(b') - \tilde A(b)| + 2\|\tilde A\|\,|p(x') - p(x)|\bigr)\\
&\le \xi^2\bigl(\mathsf{a} + 2\sqrt{N}\Lambda z_1\mathsf{a}\bigr).
\end{align*}

For the remaining factors,
\begin{equation*}
p(b_-)p(b_+)^{-1} = 1 - \xi(\nabla_\mu p)(x)\,p(x+\xi e_\mu)^{-1}.
\end{equation*}
The difference of $(\nabla_\mu p)(y)\,p(y+\xi e_\mu)^{-1}$ between $y = x'$ and $y = x$ is at most
$\xi|\nabla\nabla p| + \|\nabla p\|\,\xi|\nabla p|$. Hence this factor moves, from $b$ to $b'$, by at
most
\begin{equation*}
\xi^2\bigl(|\nabla\nabla p| + \sqrt{N}|\nabla p|^2\bigr)
\le \xi^2\bigl(\Lambda(z_2 + z_1^2) + \sqrt{N}\Lambda^2 z_1^2\bigr).
\end{equation*}

Both factors of $U^u$ are unitary, so the two movements add. The logarithm estimate in
\eqref{4.4:eq-box-conventions-a} then gives
\begin{equation*}
|\nabla A| \le 2\bigl[\mathsf{a} + 2\sqrt{N}\Lambda z_1\mathsf{a} + \Lambda(z_2 + z_1^2)
+ \sqrt{N}\Lambda^2 z_1^2\bigr].
\end{equation*}

Finally, by \eqref{4.4:eq-glued-gauge-5},
$z_1^2 \le C'^2(\mathsf{k}\mathsf{a})\cdot\mathsf{k}\mathsf{a} \le C'^2\mathsf{k}\mathsf{a}$, since
$\mathsf{k}\mathsf{a} \le S^*\mathsf{a} \le 1$. Also $z_1\mathsf{a} \le C'\mathsf{k}\mathsf{a}$ and
$\mathsf{a} \le \mathsf{k}\mathsf{a}$. Inserting these into the last bound gives
$|\nabla A| \le C_{\mathrm{gl}}\mathsf{k}\mathsf{a}$. Inserting them into \eqref{4.4:eq-glued-gauge-6}
gives
\begin{equation*}
|A| \le \tfrac{\pi}{2}(1+\Lambda C')\mathsf{k}\mathsf{a} \le C_{\mathrm{gl}}\mathsf{k}\mathsf{a}.
\end{equation*}

The site $x_0$ was arbitrary. Every bond $b$ of $\square^*$ is a bond of $\mathcal{N}$ for
$x_0 = b_-$, and so is every bond of $\square^*$ parallel and neighbouring to it. Moreover $A$ is
determined by $u$ alone. Hence the lemma holds with
\begin{equation}\label{4.4:eq-glued-gauge-7}
\begin{gathered}
C_{\mathrm{gl}} := 2\bigl[1 + (1+2\sqrt{N})\Lambda C' + (1+\sqrt{N}\Lambda)\Lambda C'^2\bigr], \\
r_{\mathrm{gl}} := \min\bigl\{r_P/(64d^3),\ r_L/(1+\Lambda C'),\ 1\bigr\}.
\end{gathered}
\end{equation}
The choice $r_{\mathrm{gl}} \le r_P/(64 d^3)$ is what gives $\bar\varphi \le r_P$ in the construction
of $u$.
\end{proof}

\begin{corollary}[Gluing on cubes covered by sub-cubes of fixed side]\label{4.4:cor-gluing-cubes}
Let $C_{\mathrm{gl}}\ge 1$ and $r_{\mathrm{gl}}\in(0,1]$ be the constants of
Lemma~\ref{4.4:lem-gluing-boxes}, and measure lengths as in
Definition~\ref{4.4:def-admissible-covers}. Let $S\ge 80\xi$, $S\ge 1$, let $\square^*$ be a
lattice cube of side $S^*\ge S$, let $U$ be $G$-valued on the bonds of $\square^*$, and let
$\mathsf{a}>0$. Suppose that every lattice cube $c\subset\square^*$ of side $S$ carries a
$G$-valued $u_c$ on the sites of $c$ with $U^{u_c}=e^{i\xi A_c}$, where $A_c$ is
$\mathfrak{g}$-valued and $|A_c|,|\nabla A_c|\le\mathsf{a}$ on $c$. If
$S^*\mathsf{a}\le r_{\mathrm{gl}}$, then there is a $G$-valued $u$ on $\square^*$ with
$U^u=e^{i\xi A}$, $A$ $\mathfrak{g}$-valued, and
\begin{equation*}
|A|,\ |\nabla A|\ \le\ C_{\mathrm{gl}}\,(S^*/S)\,\mathsf{a}\qquad\text{on }\square^*.
\end{equation*}
\end{corollary}

\begin{proof}
Put $\mathsf{k}:=S^*/S\ge 1$. The length of a run of sites is $\xi$ times its number of sites,
so $S$, $S^*$ and $S^*-S$ are integer multiples of $\xi$. We use coordinates relative to the
corner of $\square^*$; then in every direction $\mu$ the run $J_\mu$ consists of the sites
$t\in\xi\mathbb{Z}$ with $0\le t<S^*$, so that $\min J_\mu=0$ and $\max J_\mu+\xi=S^*$.

\emph{The points $p_i$.} If $\mathsf{k}=1$, put $n:=0$ and $p_0:=0$ (a single run in every
direction). If $\mathsf{k}>1$, then $S^*-S>0$; let $n\ge 1$ be the least integer with
$\tfrac34 nS\ge S^*-S$, and put
\begin{equation*}
p_i:=\xi\Big\lfloor \frac{i(S^*-S)}{n\xi}\Big\rfloor ,\qquad 0\le i\le n .
\end{equation*}
The $p_i$ are multiples of $\xi$ and non-decreasing in $i$; $p_0=0$, and $p_n=S^*-S$ because
$(S^*-S)/\xi$ is an integer. Writing $p_i=i(S^*-S)/n-\varepsilon_i$ with
$0\le\varepsilon_i<\xi$, the steps $s_i:=p_{i+1}-p_i$, $0\le i\le n-1$, satisfy
\begin{equation*}
\Big|s_i-\frac{S^*-S}{n}\Big|=|\varepsilon_i-\varepsilon_{i+1}|<\xi .
\end{equation*}
By the choice of $n$, $(S^*-S)/n\le\tfrac34 S$. If $n\ge 2$, minimality of $n$ gives
$\tfrac34(n-1)S<S^*-S$, hence, since $(n-1)/n\ge\tfrac12$,
\begin{equation*}
\frac{S^*-S}{n}>\frac34\,S\,\frac{n-1}{n}\ge\frac38\,S .
\end{equation*}
Since $\xi\le S/80$, it follows that
\begin{equation}\label{4.4:eq-gluing-cubes-1}
s_i<\Big(\frac34+\frac1{80}\Big)S\le\frac45\,S\quad\text{for all } i,\qquad
s_i>\Big(\frac38-\frac1{80}\Big)S\ge\frac13\,S\quad\text{if } n\ge 2 .
\end{equation}
Moreover $n$ is the least integer not smaller than $\tfrac43(\mathsf{k}-1)$, so
$n<\tfrac43(\mathsf{k}-1)+1$, and therefore
\begin{equation}\label{4.4:eq-gluing-cubes-2}
n+1<\frac43\,\mathsf{k}+\frac23\le 2\mathsf{k},
\end{equation}
the last inequality because $\mathsf{k}\ge1$. For $\mathsf{k}=1$ one has $n+1=1\le2\mathsf{k}$
directly.

\emph{The runs.} In every direction $\mu$ put $n_\mu:=n$ and, for $0\le i\le n$,
$I^{(\mu)}_i:=[p_i,p_i+S)$, the set of sites $t$ with $p_i\le t<p_i+S$, and $l_i:=S$. Each
$I^{(\mu)}_i$ has $S/\xi$ sites, hence length $S$, and lies in $J_\mu$ because $p_i\ge0$ and
$p_i+S\le p_n+S=S^*$.

We show that this is an $S$-admissible cover of $\square^*$ in the sense of
Definition~\ref{4.4:def-admissible-covers}, i.e.\ that conditions (ac1)--(ac5) of
\eqref{4.4:eq-admissible-covers-a} hold. Since $|J_\mu|=S^*\ge S$ for every $\mu$, the
quantity $\max\{S,\max_\mu|J_\mu|\}$ entering (ac5) is the side $S^*$ of the cube.

(ac1) We have $p_0=0=\min J_\mu$ and $p_n+l_n=S^*=\max J_\mu+\xi$. For $n=0$ there is nothing
more to check. For $n=1$ the single step is $s_0=S^*-S>0$; for $n\ge2$ every step satisfies
$s_i\ge\tfrac13 S>0$ by \eqref{4.4:eq-gluing-cubes-1}. Hence $p_i$ is strictly increasing,
and so is $p_i+l_i=p_i+S$.

(ac2) $l_i=S$.

(ac3) Consecutive runs overlap in length $p_i+l_i-p_{i+1}=S-s_i\ge S/5$, since
$s_i\le\tfrac45 S$ by \eqref{4.4:eq-gluing-cubes-1}.

(ac4) If $n\le1$ there are at most two runs and the condition is vacuous. Let $n\ge2$ and
$i+3\le i'\le n$. By monotonicity and \eqref{4.4:eq-gluing-cubes-1},
\begin{equation*}
p_{i'}\ge p_{i+3}=p_i+s_i+s_{i+1}+s_{i+2}\ge p_i+S ,
\end{equation*}
so every site of $I_{i'}$ is $\ge p_i+S$, while every site of $I_i$ is $<p_i+S$; thus
$I_i\cap I_{i'}=\emptyset$. The condition is symmetric in $i,i'$.

(ac5) By \eqref{4.4:eq-gluing-cubes-2}, $(n+1)S\le 2\mathsf{k}S=2S^*\le 4S^*$.

\emph{Conclusion.} The boxes of this cover are
$c_m=\prod_\mu I^{(\mu)}_{m_\mu}=\prod_\mu[p_{m_\mu},p_{m_\mu}+S)$,
$m\in\prod_\mu\{0,\dots,n\}$. Each is a lattice cube of side $S$ contained in $\square^*$, so
by hypothesis it carries $u_m:=u_{c_m}$ with $U^{u_m}=e^{i\xi A_m}$, $A_m:=A_{c_m}$
$\mathfrak{g}$-valued and $|A_m|,|\nabla A_m|\le\mathsf{a}$ on $c_m$. Together with
$S^*\mathsf{a}\le r_{\mathrm{gl}}$, all hypotheses of the gluing lemma for boxes
(Lemma~\ref{4.4:lem-gluing-boxes}) hold for $\square^*$ with the cover $\{c_m\}$, and that
lemma yields a $G$-valued $u$ on $\square^*$ with $U^u=e^{i\xi A}$, $A$ $\mathfrak{g}$-valued,
and $|A|,|\nabla A|\le C_{\mathrm{gl}}(S^*/S)\mathsf{a}$ on $\square^*$.
\end{proof}

\begin{remark}[Remarks on gluing: sharpness, axial gauges and real backgrounds]
\label{4.4:rem-gluing-remarks}
Let $\square^*$, $S^*$, $S$ and $\mathsf{a}$ be as in the lemma on gluing small gauges on a covered
box (Lemma~\ref{4.4:lem-gluing-boxes}), and let $\alpha_0$ be the radius of the curvature bound
\cite[(1.11)]{B-CMP109}. Let $\mathsf{o}$ denote the factor of order one in the cube side
$\mathsf{o}LM$ of the gauge condition \cite[(1.12)]{B-CMP109}.
\begin{enumerate}
\item The factor $S^*/S$ in the conclusion of Lemma~\ref{4.4:lem-gluing-boxes} is structural. The
gauge condition \cite[(1.12)]{B-CMP109} is not a restriction of the gauge condition
\eqref{4.4:eq-layered-spaces-e} of the layered analyticity spaces read at level $j$, and its bound,
proportional to the side of the cube, is of the best possible form.
\item An axial gauge on $\square^*$ bounds $\sup|A|$ by a multiple of $S^*\alpha_0$, but it does
not bound $\nabla A$ under first-order hypotheses. Gluing, as in Lemma~\ref{4.4:lem-gluing-boxes},
keeps the bound on $\nabla A$.
\item At real backgrounds, a gauge on cubes with a bound proportional to the side is proved in
\cite[Proposition~6, p.~99]{B-CMP99a}, under a second-order hypothesis on the background. That
result does not apply to a first-order $G$-valued factor, which is what the spaces
$\mathfrak{U}^{\mathrm{post}}_{k+1}(Y)$ of this chapter carry. Where both apply, it has the form of
the conclusion of Lemma~\ref{4.4:lem-gluing-boxes}.
\end{enumerate}
\end{remark}

\begin{proof}
(1) Consider a configuration whose curvature saturates the bound \cite[(1.11)]{B-CMP109}. A square
loop of side $S^*$ in $\square^*$ then encloses a flux of order $(S^*)^2\alpha_0$. The loop has
length $4S^*$, and in any gauge $U^u=e^{i\xi A}$ on $\square^*$ this flux is carried by $A$ along the
four sides of the loop. Therefore every gauge on $\square^*$ satisfies
\begin{equation*}
\sup_{\square^*}|A| \gtrsim \frac{(S^*)^2\alpha_0}{4S^*} = \frac{S^*\alpha_0}{4}.
\end{equation*}
Thus no gauge on a cube can have $|A|$ bounded independently of the side of the cube under
\cite[(1.11)]{B-CMP109} alone.

The condition \cite[(1.12)]{B-CMP109} asks for a gauge on cubes of side $\mathsf{o}LM$. Its bound is
proportional to that side, so it is not a restriction of \eqref{4.4:eq-layered-spaces-e} read at
level $j$, and by the lower bound above its form is the best possible.

For the same reason the conclusion of Lemma~\ref{4.4:lem-gluing-boxes} cannot be improved to a bound
of the size $\mathsf{a}$ of the local gauges. The local bound $\mathsf{a}$ holds on boxes of the cover,
of side at most of order $S$, and the glued gauge lives on $\square^*$, of side $S^*$. A factor
proportional to $S^*/S$ is therefore unavoidable.

(2) Fix an axial gauge on $\square^*$: the bond variables are set to the identity along a fixed
family of coordinate paths, and $A$ at every other bond is determined by the plaquette variables
along those paths. From the plaquette bound alone this gauge gives
\begin{equation*}
\sup_{\square^*}|A| \le c(d)\,S^*\alpha_0 ,
\end{equation*}
with a constant $c(d)$ depending only on the dimension; this is the bound of the lemma on the axial
gauge of a cube. That lemma bounds $A$ bond by bond and states no bound on its derivatives.

Indeed, in the axial gauge $A$ is obtained by summing the curvature along the axial paths. Its
finite difference $\nabla A$ is therefore obtained by summing the differences of the curvature along
those paths. The first-order conditions bound the curvature itself and not its differences, so none
of them bounds $\nabla A$ in this gauge.

In the gluing construction the situation is different. The transition functions
$\tilde g_{mm'}$ between the local gauges satisfy the corresponding estimates. In
those estimates the second differences $\nabla\nabla\tilde g_{mm'}$ are controlled by the first
differences $\nabla\tilde A_m$ of the local fields. These are bounded by $\mathsf{a}$ by hypothesis.
Hence the glued field keeps a bound on $\nabla A$, namely
\begin{equation*}
|\nabla A| \le C_{\mathrm{gl}}(S^*/S)\,\mathsf{a}
\quad\text{on } \square^* ,
\end{equation*}
as stated in Lemma~\ref{4.4:lem-gluing-boxes}.

(3) In the notation of \cite{B-CMP99a}, Ba{\l}aban proves the following.
\begin{quote}
Proposition 6. Let $U_0, U_0', \square, \tilde\square$ be as described above, and let
$7dL^2M\alpha_0 \le c_1$. There exists a gauge transformation $u$ defined on $\tilde\square$ and
such, that $U_0'^{\,u^{-1}} = U_1 = e^{i\eta A}$ on $\tilde\square$,
\begin{multline*}
L^j\eta|A|,\quad (L^j\eta)^2|\nabla^\eta A|,\quad (L^j\eta)^3|\partial^{\eta *}\partial^\eta A|,
\quad (L^j\eta)^3|\Delta^\eta A| \\
< 7dL^2B_1M\alpha_0 \quad \text{on } \square_j ,
\end{multline*}
\end{quote}
The two displays quoted here are \cite[(1.135)--(1.136)]{B-CMP99a}.

In that statement $\square \subset \Omega_j$ is a cube of ``a size \dots\ equal to $ML^j\eta$,
where $M$ is a multiple of $R_1M_1$'' \cite{B-CMP99a}. The proposition therefore applies also with
$M$ replaced by $\mathsf{o}LM$, whenever $\mathsf{o}LM$ is such a multiple.

The background is assumed to satisfy $U_0 \in \mathfrak{U}_k(\{\Omega_j\},\alpha_0)$. This space is
defined by the conditions \cite[(1.7)--(1.9)]{B-CMP99a}, and these include the second-order condition
\cite[(1.9)]{B-CMP99a}. A first-order $G$-valued factor, which is what the spaces
$\mathfrak{U}^{\mathrm{post}}_{k+1}(Y)$ carry, need not satisfy that condition. The proposition
therefore does not serve such a factor.

Where both results apply, the proposition has the form of the conclusion of
Lemma~\ref{4.4:lem-gluing-boxes}, in two respects:
\begin{itemize}
\item Measured in the unit $L^j\eta$, the bound on $|A|$ and $|\nabla^\eta A|$ is proportional to
$M$, hence to the side $ML^j\eta$ of the cube.
\item The hypothesis $7dL^2M\alpha_0 \le c_1$ is a ceiling on the product of the side and
$\alpha_0$, as the hypothesis $S^*\mathsf{a} \le r_{\mathrm{gl}}$ is in
Lemma~\ref{4.4:lem-gluing-boxes}.
\end{itemize}
\end{proof}

\begin{proof}[Proof of conclusion (2) of Lemma~\ref{4.4:lem-twist-stability}]
\label{4.4:prf-small-field-landing}
We prove conclusion (2) under the hypotheses of Lemma~\ref{4.4:lem-twist-stability} together with
the following:
\begin{itemize}
\item the bounds \eqref{4.4:eq-group-factor-bounds-a} for the group-valued factor, on $Y$;
\item the ceiling \eqref{4.4:eq-twist-stability-b};
\item the gauge condition \cite[(1.12)]{B-CMP109} in the form admitting cubes of side at most
$\mathsf{o}LM$, where $\mathsf{o}$ is the bounded factor in the cube side of that condition;
\item the convention on the class of cubes in \eqref{4.4:eq-layered-spaces-e} fixed with
Definition~\ref{4.4:def-layered-spaces}, under which each lattice cube of side $M$ contained in
$\Omega_j$ belongs to the class at level $j$.
\end{itemize}
The constant $B^{109}$ is the one of \eqref{4.4:eq-twist-stability-e}. From these we derive that
$(\mathbf{U}^{\flat},\mathbf{J}^{\flat})|_X\in U^c_j(X,\alpha_{0,j},\alpha_{1,j})$. This is the
space of \cite[(2.27)(ii)]{B-CMP119}, that is, the conditions \cite[(1.11)--(1.16)]{B-CMP109} at
level $j$, with $B^{109}$ as the constant of \cite[(1.12)]{B-CMP109}.

Let $X\subset\Lambda_j$, so that $X\subset\Omega_j$, and take $L^j\xi$ as the unit of length,
$L^j\xi=1$. By conclusion (1) of Lemma~\ref{4.4:lem-twist-stability}
(see \eqref{4.4:eq-twist-stability-c}), the restriction $(\mathbf{U}^{\flat},\mathbf{J}^{\flat})|_X$
satisfies \eqref{4.4:eq-layered-spaces-a}--\eqref{4.4:eq-layered-spaces-f} at level $j$. The
comparisons of layer factors made in the proof of conclusion (1) are read here at scale $j$. Since
all of $X$ lies in $\Omega_j$, the factor of the layer $n=j$ is $1-\beta(1-2^{0})=1$. The
conditions of \cite{B-CMP109} are then found among those of
Definition~\ref{4.4:def-layered-spaces} as follows.
\begin{itemize}
\item Condition \cite[(1.11)]{B-CMP109} is the first inequality of
\eqref{4.4:eq-layered-spaces-a}.
\item Condition \cite[(1.13)]{B-CMP109} is \eqref{4.4:eq-layered-spaces-f}.
\item Condition \cite[(1.14)]{B-CMP109} is the second inequality of
\eqref{4.4:eq-layered-spaces-a} together with \eqref{4.4:eq-layered-spaces-c}. This holds at the
choice $\gamma_0=\alpha_0$, which is the choice made in \cite{B-CMP109}, and
\eqref{4.4:eq-layered-spaces-c} carries $\alpha_{0,n}$.
\item Item (iv) of the definition of the space in \cite{B-CMP109}, namely
\cite[(1.15)--(1.16)]{B-CMP109}, is \eqref{4.4:eq-layered-spaces-b} and
\eqref{4.4:eq-layered-spaces-d} under \eqref{4.4:eq-post-step-space-a}.
\end{itemize}
For the last of these, the normalisations at $n=j$, with $L^j\xi=1$, give the right members of
\cite[(1.16)]{B-CMP109}:
\begin{equation}\label{4.4:eq-small-field-landing-1}
L^{-2p}\bigl(L^jL^{-p}\bigr)^{-2}=L^{-2j}=\xi^2,\qquad
L^{j-p}\bigl(L^jL^{-p}\bigr)^{-3}=L^{-2j+2p}=(L^p\xi)^2 .
\end{equation}
The last clause of item (iv) of \cite{B-CMP109} has no counterpart among
\eqref{4.4:eq-layered-spaces-a}--\eqref{4.4:eq-layered-spaces-f} \cite[p.~261]{B-CMP119}. It is
obtained from the hypothesis \eqref{4.4:eq-group-factor-bounds-a} on $Y$. That hypothesis says
that \eqref{4.4:eq-layered-spaces-b} and \eqref{4.4:eq-layered-spaces-d} under
\eqref{4.4:eq-post-step-space-a} hold also with $U$ in place of $\mathbf{U}$. Since
$U^{\flat}=U$, the normalised increments vanish, $\pi_n=0$, and the same three comparisons of
layer factors carry these bounds from $Y$ at level $k+1$ to $X$ at level $j$. This is conclusion
(1$'$) of Lemma~\ref{4.4:lem-twist-stability}, whose hypothesis is exactly
\eqref{4.4:eq-group-factor-bounds-a} on $Y$.

It remains to verify the gauge condition \cite[(1.12)]{B-CMP109}, in the form admitting cubes of
side at most $\mathsf{o}LM$. Let $\square^*\subset X$ be a lattice cube of side $S^*\le\mathsf{o}LM$.
For a $G$-valued gauge $u$ on a cube we write $A$ for its field and $\nabla A$ for its lattice
gradient, the quantities bounded in \cite[(1.12)]{B-CMP109}.

Case $S^*\le M$. The set $\Omega_j$ is a union of cubes of a partition of mesh at least $M$.
Prolong the coordinate intervals of $\square^*$ exactly as in the claim in the proof of the gauge
condition by extension of cubes, within the proof of conclusion (1) of
Lemma~\ref{4.4:lem-twist-stability}. This yields
a cube $\square\subset\Omega_j$ of side $M$, with $\square\supset\square^*$, belonging to the class
of \eqref{4.4:eq-layered-spaces-e} at level $j$. By conclusion (1), this cube carries a gauge on
$\square\cap X\supset\square^*$, and its restriction to $\square^*$ satisfies
\begin{equation}\label{4.4:eq-small-field-landing-2}
|A|,\ |\nabla A|<BM\alpha_{0,j}\le \mathsf{o}LMB^{109}\alpha_{0,j}.
\end{equation}

Case $S^*>M$. Every lattice cube $c\subset\square^*$ of side $M$ lies in $X\subset\Omega_j$.
By the convention on the class stated above, $c$ therefore belongs to the class of
\eqref{4.4:eq-layered-spaces-e} at level $j$. By conclusion (1), $c$ carries a $G$-valued gauge
$u_c$ on $c\cap X=c$ with
\begin{equation}\label{4.4:eq-small-field-landing-3}
|A_c|,\ |\nabla A_c|<\mathsf{a}:=BM\alpha_{0,j}.
\end{equation}
We apply gluing on cubes covered by sub-cubes of fixed side
(Corollary~\ref{4.4:cor-gluing-cubes}) with $S:=M$. Its hypotheses are checked as follows.
\begin{itemize}
\item The floor on $M$ fixed in the first part of this chapter, with $\ell$ the constant entering
that floor, together with $\xi=L^{-j}\le1$, gives
\begin{equation*}
M\ge48\ell\ge80\ge80\xi,
\end{equation*}
so $S\ge80\xi$ and $S\ge1$.
\item The side of $\square^*$ satisfies $S^*\in(M,\mathsf{o}LM]$, so $S^*\ge S$.
\item Every lattice cube $c\subset\square^*$ of side $S$ carries $u_c$ with
\eqref{4.4:eq-small-field-landing-3}.
\item The smallness condition holds, by \eqref{4.4:eq-twist-stability-b} and
$\mathsf{o}L\ge1$:
\begin{equation}\label{4.4:eq-small-field-landing-4}
\begin{aligned}
S^*\mathsf{a}&\le\mathsf{o}LM\cdot BM\alpha_{0,j}\le(\mathsf{o}LM)^2B\alpha_{0,j}\\
&\le6(\mathsf{o}LM)^2B\alpha_{0,j}\le r_{\mathrm{gl}}.
\end{aligned}
\end{equation}
\end{itemize}
The corollary yields one $G$-valued gauge $u$ on $\square^*$. By the choice
\eqref{4.4:eq-twist-stability-e}, which gives $B^{109}\ge4C_{\mathrm{gl}}B$, it satisfies
\begin{equation}\label{4.4:eq-small-field-landing-5}
\begin{aligned}
|A|,\ |\nabla A|&\le C_{\mathrm{gl}}\frac{S^*}{M}\,BM\alpha_{0,j}
=C_{\mathrm{gl}}S^*B\alpha_{0,j}\le C_{\mathrm{gl}}\mathsf{o}LMB\alpha_{0,j}\\
&<4C_{\mathrm{gl}}\mathsf{o}LMB\alpha_{0,j}\le\mathsf{o}LMB^{109}\alpha_{0,j}.
\end{aligned}
\end{equation}
In both cases $\square^*$ carries a $G$-valued gauge with the bound $\mathsf{o}LMB^{109}\alpha_{0,j}$
required in \cite[(1.12)]{B-CMP109}. This completes the verification of
$(\mathbf{U}^{\flat},\mathbf{J}^{\flat})|_X\in U^c_j(X,\alpha_{0,j},\alpha_{1,j})$.

The ceiling \eqref{4.4:eq-twist-stability-b} is a ceiling on $\gamma$. By
\cite[(2.28)]{B-CMP119},
\begin{equation*}
\alpha_{0,j}=g_jC_0\bigl(\log g_j^{-2}\bigr)^{q_0}.
\end{equation*}
The derivative of $t\mapsto t(\log t^{-2})^{q_0}$ on $0<t<1$ is
\begin{equation*}
\bigl(\log t^{-2}\bigr)^{q_0-1}\bigl(\log t^{-2}-2q_0\bigr).
\end{equation*}
This is non-negative for $t\le e^{-q_0}$, so the function increases on $(0,e^{-q_0}]$. Hence,
whenever $g_j\le\gamma\le e^{-q_0}$,
\begin{equation}\label{4.4:eq-small-field-landing-6}
\alpha_{0,j}\le\gamma C_0\bigl(\log\gamma^{-2}\bigr)^{q_0}.
\end{equation}
Thus \eqref{4.4:eq-twist-stability-b} holds for all $j$ once $\gamma\le e^{-q_0}$ and
$6(\mathsf{o}LM)^2B\gamma C_0(\log\gamma^{-2})^{q_0}\le r_{\mathrm{gl}}$. This is a ceiling on
$\gamma$, imposed after $L$, $M$, $B$, $C_0$ and $q_0$, independent of $k$ and of $j$, and with no
lower bound on the coupling. The factor $6$ is an absolute number.

The constant in \eqref{4.4:eq-twist-stability-e} is the maximum of $4C_{\mathrm{gl}}B$ and the
lower bound on the constant of \cite[(1.12)]{B-CMP109} required by the one-power hereditary class.
The gluing argument above uses only $B^{109}\ge4C_{\mathrm{gl}}B$; the maximum serves both
requirements, which concern the same constant. The constant $B$ of \cite[(1.12)]{B-CMP109} is
required there only to be sufficiently large and is determined later, so it may be taken as large
as required. The choices that depend on $B^{109}$ are three ceilings on a small-field radius,
imposed after $M$ and $B^{109}$ (the first ceiling of the small-field corollary and two further
ceilings, named where they are imposed), and the ceilings on $\gamma$ in the operator bounds at
free current. All of these are made after $B^{109}$, and $B^{109}$ itself is determined from
$C_{\mathrm{gl}}$, $B$ and the lower bound of the one-power hereditary class; hence $B^{109}$ is
fixed before $M$ in the order of choice of the constants.
\end{proof}

\begin{proof}[Proof of conclusion $(2^{\mathrm{int}})$ of Lemma~\ref{4.4:lem-twist-stability}:
interior sources on the one-power hereditary class]
\label{4.4:prf-interior-sources}
Conclusion $(2^{\mathrm{int}})$ in \eqref{4.4:eq-twist-stability-c} is derived here from two
statements that enter as hypotheses and are used by statement only. We also use the response
bound at free current \eqref{4.4:eq-response-bound-a} and the
ratio of radii along the coupling flow, Lemma~\ref{4.4:lem-radii-ratio}, in the form
\eqref{4.4:eq-radii-ratio-a}. The two hypotheses are the following.
\begin{itemize}
\item[(R)] The restriction theorem for the one-power hereditary class. Let
$1 \le i' \le i$ and $\mu := L^{-(i-i')}$. Let $Y \in \mathbf{D}_{i-1}$ and
$X \in \mathbf{D}_{i'-1}$ with
$X \subset Y$, and let $(\mathbf{U},\mathbf{J})$ be a representative of
$U^{\Sigma^1}_{i}(Y;a_0,a_1,\gamma_0)$. Then $(\mathbf{U},\mu^3\mathbf{J})$ restricted to $X$
is a representative of $U^{\Sigma^1}_{i'}(X;\mu a_0,\mu a_1,\mu^3\gamma_0)$, with the same
$G$-valued factor $U$ and the restricted gauges. Its clause (iv)${}^{\Sigma}$ comes from the
inclusion $\mathfrak{T}_{i'}(X) \subset \mathfrak{T}_{i}(Y)$ together with a change of units.
Its clause (i) contains the cube reading of \cite[(1.12)]{B-CMP109}, in which that bound is
imposed on every cube of side at most $\mathsf{o}LM\cdot L^{i}\xi$ contained in the domain.
No condition is imposed
on $M$, on the position of $X$ in $Y$, on the shape or thickness of $Y$, or on $i-i'$.
\item[(P)] The statement that the pair produced by the operator bounds at free current lies in
the one-power hereditary class. Let $j \ge 1$ and $X \in \mathbf{D}_j$, and let
$(\mathbf{U},\mathbf{J})$ be a representative of a class $U^{\Sigma^1}_j(X;a_0,a_1,\gamma_0)$
to which the operator bounds at free current are applied under the ceilings below, with
response $\boldsymbol{\mathcal{H}}$. For every tower
$(p,Z) \in \mathfrak{T}_j(X)$ the statement instantiates the uniform-increment lemma of the
operator bounds at free current, whose constant is $C_{IV} = C_{IV}(d,L)$, at $(p,Z)$ with
parameters $(M/L, LB)$ and radii
\begin{equation*}
\bigl(L^{-(j-p)}\alpha_{0,j},\ L^{-(j-p)}\alpha_{1,j},\ L^{-(j-p)}\alpha_2\bigr);
\end{equation*}
here $L^{-(j-p)}$ is the scale ratio between the levels $j$ and $p$,
$\alpha_{0,j}$ and $\alpha_{1,j}$ are the first two radii of the target
class, and $\alpha_2$ is the number in the hypothesis on the increment $A$ below. It holds
under the floors
\begin{equation*}
c_{12}L \ge 2,\qquad L \ge 4d,\qquad \delta_0 M \ge 32L,\qquad 2LR_1M_1 \le M,
\end{equation*}
where $c_{12}$ and $R_1$ are constants of the uniform-increment lemma,
and under the ceilings of the uniform-increment lemma at
$(M,B^{109};\alpha_{0,j},\alpha_{1,j})$; these ceilings are imposed among the ceilings on
$\gamma$. Suppose the increment $A := \boldsymbol{\mathcal{H}}_w$, written in level-$j$ units,
satisfies
\begin{equation*}
|A|,\ |\nabla_{\mathbf{U}}A| \le \alpha_2 < \min\{\alpha_{1,j}/8,\ \alpha_{0,j}/C_{IV}\}.
\end{equation*}
Then clause (iv)${}^{\Sigma}$ of $U^{\Sigma^1}_j(X;\alpha_{0,j},\alpha_{1,j})$, that is,
\cite[(1.16)]{B-CMP109} for the tower of every $(p,Z) \in \mathfrak{T}_j(X)$, holds for the
tower built on the twisted configuration $e^{i\xi\boldsymbol{\mathcal{H}}}\mathbf{U}$.
\end{itemize}

Throughout, $1 \le j \le k$, $Y \subset \Lambda_{k+1}$, $X \in \mathbf{D}_j$, $X \subset Y$ and
$w \in \mathfrak{D}$. The pair $(\mathbf{U},\mathbf{J})$ is a representative of
\begin{equation*}
U^{\Sigma^1}_{k+1}\bigl(Y;(1+\beta)\alpha_{0,k+1},(1+\beta)\alpha_{1,k+1},\alpha_{0,k+1}\bigr),
\end{equation*}
a class defined for $Y \in \mathbf{D}_k$. In place of
$(\mathbf{U},\mathbf{J}) \in \mathfrak{U}^{\mathrm{post}}_{k+1}(Y)$ only membership in this
class is assumed; in particular the bound \cite[(2.38)]{B-CMP119} of the layered analyticity
spaces, the one carrying the absolute constant $B$, is not asked. The twisted pair is
\begin{equation*}
(\mathbf{U}^{\flat},\mathbf{J}^{\flat})
= \bigl(e^{i\xi\boldsymbol{\mathcal{H}}}\mathbf{U},\ \mathbf{J}+\mathfrak{J}\bigr),
\qquad \mathfrak{J} = \mathfrak{J}^{(\mathbf{U},\mathbf{J})}(w).
\end{equation*}
We must show that its restriction to $X$ lies in
$U^{\Sigma^1}_j(X;\alpha_{0,j},\alpha_{1,j})$.

We first determine the layer structure. Since
$Y \subset \Lambda_{k+1} \subset \Omega_{k+1} \subset \Omega_j$, the domain $Y$ meets only one
layer, the top one. Thus $n = k+1$, $\bar n = \min(n,k) = k$, and $D_{\bar n} = D_k = 0$.

Next we treat the pointwise bounds. At level $k+1$, consider the bounds
\cite[(1.11), (1.13), (1.14), (1.15)--(1.16)]{B-CMP109}. In the unit $L^{k+1}\xi$ and with
factors at most
$1+\beta = F_{\mathrm{top}}$, they are respectively:
\begin{itemize}
\item the first member of \eqref{4.4:eq-layered-spaces-a};
\item the bound \eqref{4.4:eq-layered-spaces-f};
\item the second member of \eqref{4.4:eq-layered-spaces-a} together with
\eqref{4.4:eq-layered-spaces-c};
\item the bound \eqref{4.4:eq-layered-spaces-b} together with \eqref{4.4:eq-layered-spaces-d}.
\end{itemize}
These are the identifications made in the proof of conclusion (2) of
Lemma~\ref{4.4:lem-twist-stability}, read at level $k+1$. The third radius of the space of
\cite[(1.34)]{B-CMP116} is $\gamma_0 = \alpha_{0,k+1} \le (1+\beta)\alpha_{0,k+1}$.

The layer-by-layer comparison in the proof of conclusion (1) of
Lemma~\ref{4.4:lem-twist-stability} bounds each increment separately, one pointwise bound at a
time. Of the incoming pair
it uses only the bounds just listed. Its line for the layers $j < n \le k$ and the top layer,
taken at $n = k+1$, gives the following in the unit of scale $j$: the new factor is at most
$\tfrac12$ and the contribution of the twist is at most $\tfrac14$. The total is therefore
$< 1$.

We turn to clause (iv)${}^{\Sigma}$ for the incoming pair. We apply hypothesis (R) with
$(i,i') := (k+1,j)$ and $\mu := L^{-(k+1-j)}$. Its hypotheses hold for the following reasons.
First, $1 \le j \le k < k+1$. Second, $Y \in \mathbf{D}_k$. Third, $X \subset Y$. Fourth,
$X \in \mathbf{D}_j \subset \mathbf{D}_{j-1}$, by the nesting of the localization domains.
Hypothesis (R) then shows that
$(\mathbf{U},\mu^3\mathbf{J})$ restricted to $X$ is a representative of
\begin{equation*}
U^{\Sigma^1}_j\bigl(X;\mu(1+\beta)\alpha_{0,k+1},\ \mu(1+\beta)\alpha_{1,k+1},\
\mu^3\alpha_{0,k+1}\bigr).
\end{equation*}
In particular, since $\mathfrak{T}_j(X) \subset \mathfrak{T}_{k+1}(Y)$, clause
(iv)${}^{\Sigma}$ holds for the tower of every $(p,Z) \in \mathfrak{T}_j(X)$, both for
$\mathbf{U}$ and for $U$, at the constant
\begin{equation}\label{4.4:eq-interior-sources-1}
(1+\beta)\alpha_{0,k+1}L^{-(k+1-p)}
= \bigl[\mu(1+\beta)\alpha_{0,k+1}\bigr]L^{-(j-p)}
\end{equation}
in level-$p$ units. The identity \eqref{4.4:eq-interior-sources-1} is
$\mu L^{-(j-p)} = L^{-(k+1-j)}L^{-(j-p)} = L^{-(k+1-p)}$.

The right-hand side of \eqref{4.4:eq-interior-sources-1} is the constant that clause
(iv)${}^{\Sigma}$ of the target class imposes at $a_0 := \mu(1+\beta)\alpha_{0,k+1}$. This
constant satisfies
\begin{equation}\label{4.4:eq-interior-sources-2}
a_0 = (1+\beta)\frac{\alpha_{0,k+1}}{\alpha_{0,j}}\,L^{-(k+1-j)}\,\alpha_{0,j}
\le \tfrac12\,\alpha_{0,j}.
\end{equation}
The factor in front of $\alpha_{0,j}$ is the new factor of the top-layer line. It is at most
$\tfrac12$ by Lemma~\ref{4.4:lem-radii-ratio}, applied to the pair $(j,k+1)$, and the floor (s2)
of \eqref{4.4:eq-twist-stability-a}.

This argument imposes no condition on $M$, on the position of $X$ in $Y$, or on the thickness
of $Y$. It treats towers of level equal to $j$ and towers of lower level in the same way.
By the corollary of hypothesis (R) that concerns the large-field part of the coupling
increment, the same restriction argument serves, unchanged, for the towers met there, at the
radii $\alpha_{0,\cdot},\alpha_{1,\cdot}$ of \cite[(2.28)]{B-CMP119}.

Next we treat the product. The twisted pair $(\mathbf{U}^{\flat},\mathbf{J}^{\flat})$ must itself
lie in $U^{\Sigma^1}_j(X;\alpha_{0,j},\alpha_{1,j})$. Its clauses (i)--(iii) are the bounds of
the top-layer line described above. Its clause (iv)${}^{\Sigma}$ is \cite[(1.16)]{B-CMP109} for
the tower of every $(p,Z) \in \mathfrak{T}_j(X)$ built on $\mathbf{U}^{\flat}$. This is
hypothesis (P) at $A := \boldsymbol{\mathcal{H}}_w$ in level-$j$ units, applied to
$(\mathbf{U},\mu^3\mathbf{J})$ restricted to $X$, which is a representative of the one-power
hereditary class at level $j$ on $X$ by the restriction step above; we verify its
hypothesis on $A$.

The first two members of \eqref{4.4:eq-response-bound-a}, taken in the top layer $\bar n = k$
with $D_k = 0$, give
$\ell_k|\boldsymbol{\mathcal{H}}|,\ \ell_k^2|\nabla\boldsymbol{\mathcal{H}}| \le C_Hw_k$.
Since $\ell_j/\ell_k = L^{-(k-j)}$, in level-$j$ units this yields
\begin{align*}
|A| &= \ell_j|\boldsymbol{\mathcal{H}}| = L^{-(k-j)}\ell_k|\boldsymbol{\mathcal{H}}|
\le L^{-(k-j)}C_Hw_k,\\
|\nabla_{\mathbf{U}}A| &= \ell_j^2|\nabla\boldsymbol{\mathcal{H}}|
= L^{-2(k-j)}\ell_k^2|\nabla\boldsymbol{\mathcal{H}}| \le L^{-2(k-j)}C_Hw_k
\le L^{-(k-j)}C_Hw_k.
\end{align*}
Lemma~\ref{4.4:lem-radii-ratio}, applied to the pair $(j,k)$, gives
\begin{equation*}
\frac{w_k}{\alpha_{0,j}} \le \frac{2\|C\|C_1A_1}{C_0}(1+\beta_0)(1+k-j)^{1/2},
\qquad
\frac{w_k}{\alpha_{1,j}} \le \frac{2\|C\|C_1A_1}{C_1'}(1+\beta_0)(1+k-j)^{1/2}.
\end{equation*}
We also use $(1+m)^{1/2}L^{-m} \le 1$ for $m \ge 0$ and $L \ge 2$. At $m = 0$ both sides
equal $1$. For $m \ge 1$, $(1+m)^{1/2} \le 2^m \le L^m$. We apply it with $m = k-j$.

Hence, by the inequality $C_1' > 16C_H(1+\beta_0)\|C\|C_1A_1$ contained in (s3) of
\eqref{4.4:eq-twist-stability-a},
\begin{equation*}
\frac{L^{-(k-j)}C_Hw_k}{\alpha_{1,j}}
\le \frac{2C_H(1+\beta_0)\|C\|C_1A_1}{C_1'}\,(1+k-j)^{1/2}L^{-(k-j)}
< \frac18 .
\end{equation*}
Likewise, by $(\mathrm{s}3')$ of \eqref{4.4:eq-twist-stability-a}, namely
$C_0 > 2C_{IV}C_H(1+\beta_0)\|C\|C_1A_1$,
\begin{equation*}
\frac{L^{-(k-j)}C_Hw_k}{\alpha_{0,j}}
\le \frac{2C_H(1+\beta_0)\|C\|C_1A_1}{C_0}\,(1+k-j)^{1/2}L^{-(k-j)}
< \frac{1}{C_{IV}} .
\end{equation*}
Therefore
\begin{equation*}
|A|,\ |\nabla_{\mathbf{U}}A| \le L^{-(k-j)}C_Hw_k
< \min\{\alpha_{1,j}/8,\ \alpha_{0,j}/C_{IV}\},
\end{equation*}
so hypothesis (P) applies with $\alpha_2 := L^{-(k-j)}C_Hw_k$.

The two inequalities used involve only $C_H$, $\beta_0$, $\|C\|$, $C_1$, $A_1$ and
$C_{IV}(d,L)$; they do not depend on $k$ or $j$. The growth $(1+k-j)^{1/2}$ coming from
Lemma~\ref{4.4:lem-radii-ratio} is absorbed by the unit factor $L^{-(k-j)}$. So the bound holds
for all $1 \le j \le k < \infty$, with no cap on the depth $k-j$.

The remaining hypotheses of (P) are the floors $c_{12}L \ge 2$, $L \ge 4d$,
$\delta_0M \ge 32L$ and $2LR_1M_1 \le M$, and the ceilings of the uniform-increment lemma at
$(M,B^{109};\alpha_{0,j},\alpha_{1,j})$. The four floors are among the floors on $L$ and $M$
fixed with the constants of this chapter, and the ceilings are among the ceilings on $\gamma$
assumed in Lemma~\ref{4.4:lem-twist-stability}. Hypothesis (P) then gives clause
(iv)${}^{\Sigma}$ for
$(\mathbf{U}^{\flat},\mathbf{J}^{\flat})$ on $X$.

We now treat the $G$-valued factor. Clause (iv)${}^{\Sigma}$ of the one-power hereditary class is
imposed for $U$ in place of $\mathbf{U}$ as well, as in the last clause of
\cite[(1.16)]{B-CMP109}. Hypothesis (R) keeps
the $G$-valued factor $U$ unchanged under restriction. The twist does not act on it:
$U^{\flat} = U$. The only change is the change of units, and its factor is at most
$\tfrac12$ by \eqref{4.4:eq-interior-sources-2}.

Finally we treat \cite[(1.12)]{B-CMP109}, which is clause (i) of hypothesis (R) in the cube
reading. Let $\square^* \subset X$ be a cube of side at most $\mathsf{o}LM\cdot L^j\xi$. Then
$\square^* \subset Y$, and since $L^j\xi \le L^{k+1}\xi$ its side is at most
$\mathsf{o}LM\cdot L^{k+1}\xi$. In the cube reading of \cite[(1.12)]{B-CMP109} it is therefore
one of the cubes on which the hypothesis on $(\mathbf{U},\mathbf{J})$ at level $k+1$ is imposed.
Hence it carries a $G$-valued gauge transformation $u$ with $U^u = e^{i\xi A}$, where here $A$
denotes the Lie-algebra-valued field of \cite[(1.12)]{B-CMP109} on $\square^*$, not the
increment of the preceding paragraphs, and
\begin{equation*}
L^{k+1}\xi|A|,\ (L^{k+1}\xi)^2|\nabla A| < \mathsf{o}LMB^{109}(1+\beta)\alpha_{0,k+1},
\end{equation*}
with $B^{109}$ as in \eqref{4.4:eq-twist-stability-e}. Consequently
\begin{align*}
L^j\xi|A| &= L^{-(k+1-j)}\,L^{k+1}\xi|A|\\
&< \mathsf{o}LMB^{109}\alpha_{0,j}\cdot(1+\beta)\frac{\alpha_{0,k+1}}{\alpha_{0,j}}
L^{-(k+1-j)}\\
&\le \tfrac12\,\mathsf{o}LMB^{109}\alpha_{0,j},
\end{align*}
where the last step uses Lemma~\ref{4.4:lem-radii-ratio} for the pair $(j,k+1)$ and the floor
(s2) of \eqref{4.4:eq-twist-stability-a}, as in \eqref{4.4:eq-interior-sources-2}. The same
chain bounds $(L^j\xi)^2|\nabla A|$, with $L^{-2(k+1-j)} \le L^{-(k+1-j)}$ in place of
$L^{-(k+1-j)}$. Moreover $U^{\flat} = U$, so the same $u$ serves for the twisted pair.

No gluing of cubes enters this argument. The cubes of side at most $\mathsf{o}LM\cdot L^j\xi$
inside $X$ are among the cubes on which the hypothesis on $(\mathbf{U},\mathbf{J})$ is already
imposed.

Together, these steps establish clauses (i)--(iii), clause (iv)${}^{\Sigma}$ for $\mathbf{U}$ and
for $U$, and the cube bound \cite[(1.12)]{B-CMP109} for the twisted pair restricted to $X$. Hence
$(\mathbf{U}^{\flat},\mathbf{J}^{\flat})|_X \in U^{\Sigma^1}_j(X;\alpha_{0,j},\alpha_{1,j})$, which
is conclusion $(2^{\mathrm{int}})$.
\end{proof}

\begin{lemma}[The new small-field region avoids old large-field collars]
\label{4.4:lem-collar-avoidance}
Let $k\ge 0$ and write $x_i=g_i^{-2}$. For every level $i$ let $R_i$ be the partition radius of
\cite[(2.5)]{B-CMP119}: the least power of $L$ such that $R_i\ge(\log x_i)^{r^{(2.5)}}$. For
$j\le k$ let $Z_j=\Lambda_j^c$ be the large-field region \cite[(2.3)]{B-CMP119}, and let $Z'_j$
be the union of all $LMR_{j+1}$-cubes that meet $Z_j$. For a set $Z$ that is a union of cubes of
a partition, $Z^{\sim}$ denotes $Z$ enlarged by one layer of cubes of that partition, and
$Z^{\sim n}$ the $n$-fold enlargement. Assume the following.
\begin{itemize}
\item The couplings satisfy $g_i\le\gamma$, that is $x_i\ge\gamma^{-2}$, and the first of the
flow inequalities read in Lemma~\ref{4.4:lem-radii-ratio} holds at every pair of consecutive
levels, $x_i\le x_{i+1}+B^{\sharp}$.
\item The ceiling
\begin{equation}\label{4.4:eq-collar-avoidance-a}
r^{(2.5)}B^{\sharp}\gamma^{2}\le\log L\cdot\log\gamma^{-2}
\end{equation}
holds.
\end{itemize}
Then:
\begin{enumerate}
\item[(a)] $LR_{j+1}\ge R_j$ for every $j\le k$.
\item[(b)] Let $j\le k$, and let $\Omega_{j+1}$ be produced from the density of level $j$ by
the $\mathbf{T}$-step through the construction of \cite[(3.5)]{B-CMP119}. If the enlargement
$Z_j^{\sim}$ is formed with the cubes of the partition used in \cite[(2.41)]{B-CMP119}, either
the $MR_j$-cubes of which $Z_j$ is a union or the $LMR_{j+1}$-cubes, then
\begin{equation*}
Z_j^{\sim}\subset {Z'}_j^{\sim}\subset {Z'}_j^{\sim 8}\subset\Omega_{j+1}^c .
\end{equation*}
\item[(c)] If, for a given $j\le k$, $\Omega_{j+1}$ is produced as in (b), then
$\Lambda_{k+1}\cap Z_j^{\sim}=\emptyset$. The same conclusion holds when an $\mathbf{R}$-step
intervenes before the $\mathbf{T}$-step producing $\Omega_{j+1}$, under the following hypothesis,
which forms part of the statement on the representation and bound of the large-field terms: the
$\mathbf{T}$-step producing $\Omega_{j+1}$ runs the construction of \cite[(3.5)]{B-CMP119} from the new large-field
region $Z_j\cup\bigcup_i Y_i$, where the $Y_i$ are the regions produced by the
$\mathbf{R}$-operation (they are not to be confused with the domain $Y$ of (d)).
\item[(d)] Under the hypotheses of (c), an old term of level $j\le k$ localised in a domain
$X\subset\Lambda_{k+1}$ (in particular, in a domain $X\subset Y$ with $Y\subset\Lambda_{k+1}$,
as in the interior case of Lemma~\ref{4.4:lem-twist-stability}) is not a $\mathbf{B}^{(j)}$-term;
it is an $\mathbf{E}^{(j)}$-term or an $\mathbf{R}^{(j)}$-term.
\end{enumerate}
\end{lemma}

Condition \eqref{4.4:eq-collar-avoidance-a} is a one-sided ceiling on $\gamma$, imposed after
$r^{(2.5)}$, $B^{\sharp}$ and $L$ have been fixed; it does not involve $k$.

\begin{proof}
\emph{Part (a).} Recall the definition of the radii in \cite[(2.5)]{B-CMP119}: $R_j$ is the
smallest number of the form a power of $L$ such that $R_j\ge(\log g_j^{-2})^{r^{(2.5)}}$. The same
definition is made at every level, and the exponent satisfies $r^{(2.5)}\ge 2$
\cite{B-CMP119}.

Fix $j\le k$ and suppose, to reach a contradiction, that $LR_{j+1}<R_j$. Since $LR_{j+1}$ and
$R_j$ are both powers of $L$, the strict inequality forces $L\cdot LR_{j+1}\le R_j$, that is,
$LR_{j+1}\le R_j/L$. The number $R_j/L$ is a power of $L$ smaller than $R_j$, so by the
minimality of $R_j$ it fails the defining inequality: $R_j/L<(\log x_j)^{r^{(2.5)}}$. Together
with $R_{j+1}\ge(\log x_{j+1})^{r^{(2.5)}}$ this gives
\begin{equation}\label{4.4:eq-collar-avoidance-1}
L(\log x_{j+1})^{r^{(2.5)}}\le LR_{j+1}\le R_j/L<(\log x_j)^{r^{(2.5)}} .
\end{equation}

Now apply the first of the flow inequalities read in Lemma~\ref{4.4:lem-radii-ratio} at the pair
of levels $(j,j+1)$: $x_j\le x_{j+1}+B^{\sharp}$. Taking logarithms and using $\log(1+t)\le t$ with
$t=B^{\sharp}/x_{j+1}$, then $x_{j+1}\ge\gamma^{-2}$ and $\log x_{j+1}\ge\log\gamma^{-2}$, we get
\begin{equation*}
\log x_j\le\log x_{j+1}+\frac{B^{\sharp}}{x_{j+1}}
\le\log x_{j+1}\Bigl(1+\frac{B^{\sharp}\gamma^{2}}{\log\gamma^{-2}}\Bigr).
\end{equation*}
Raising to the power $r^{(2.5)}$ and using $1+t\le e^{t}$,
\begin{equation*}
\Bigl(\frac{\log x_j}{\log x_{j+1}}\Bigr)^{r^{(2.5)}}
\le\Bigl(1+\frac{B^{\sharp}\gamma^{2}}{\log\gamma^{-2}}\Bigr)^{r^{(2.5)}}
\le\exp\Bigl(\frac{r^{(2.5)}B^{\sharp}\gamma^{2}}{\log\gamma^{-2}}\Bigr)\le L ,
\end{equation*}
the last inequality being \eqref{4.4:eq-collar-avoidance-a} divided by $\log\gamma^{-2}$ and
exponentiated. Hence $(\log x_j)^{r^{(2.5)}}\le L(\log x_{j+1})^{r^{(2.5)}}$, which contradicts
\eqref{4.4:eq-collar-avoidance-1}. Therefore $LR_{j+1}\ge R_j$.

Only this inequality is used below. It is consistent with the remark in \cite{B-CMP122a} that
``in some steps the number $R_k$ decreases by the factor $L^{-1}$''.

\emph{Part (b).} We first recall the construction of the nest member in \cite[pp.~264--265]{B-CMP119}.
At the $\mathbf{T}$-step that produces $\Omega_{k+1}$ from the density of level $k$,
Ba{\l}aban writes:
\begin{quote}
The term has a large field region $Z_k=\Lambda_k^c$, and we denote by $Z'_k$ the union of all
$LMR_{k+1}$-cubes intersecting the region $Z_k$ \dots\ compatible with all the previous
partitions. We surround $Z'_k$ by four layers of the $LMR_{k+1}$-cubes, i.e., we take the
domain ${Z'}_k^{\sim 4}$
\end{quote}
and then
\begin{quote}
Denote by $P^1_{k+1}$ the subset of $T^{(k+1)}$ such that
$B^{k+1}(P^1_{k+1})=({P'}^{\sim}_{k+1})^c\cap({Z'}^{\sim 5}_k)^c$
\end{quote}
and
\begin{quote}
we surround the domain $Q'_{k+1}\cup((B^{k+1}(P^1_{k+1}))^c)^{\sim}$ by two layers of such
cubes. Denote
\begin{equation*}
\Omega_{k+1}=\bigl({Q'}^{\sim 2}_{k+1}\cup(B^{k+1}(P^1_{k+1})^c)^{\sim 3}\bigr)^c .
\end{equation*}
\end{quote}
The last display is \cite[(3.5)]{B-CMP119}. In it, $T^{(k+1)}$ is a set whose subsets $P$ are
carried by the map $B^{k+1}$ to domains $B^{k+1}(P)$, and ${P'}^{\sim}_{k+1}$ and $Q'_{k+1}$ are
domains defined on the cited page; these objects enter below only through the two complements
just displayed. We use the same construction at the step $j\le k$ in place of $k$.

Since $Z'_j$ is the union of the $LMR_{j+1}$-cubes that meet $Z_j$, and these cubes cover
$Z_j$, we have $Z_j\subset Z'_j$. By the displayed definition of $B^{j+1}(P^1_{j+1})$, its
complement is $({P'}^{\sim}_{j+1})\cup{Z'}^{\sim 5}_j$, which contains ${Z'}^{\sim 5}_j$.
Enlargement by layers of cubes of a fixed partition is monotone with respect to inclusion, and
three layers added to ${Z'}_j^{\sim 5}$ give ${Z'}_j^{\sim 8}$. Taking the complement in
\cite[(3.5)]{B-CMP119} at step $j$ therefore gives
\begin{equation}\label{4.4:eq-collar-avoidance-2}
\Omega_{j+1}^c={Q'}^{\sim 2}_{j+1}\cup(B^{j+1}(P^1_{j+1})^c)^{\sim 3}
\supset({Z'}^{\sim 5}_j)^{\sim 3}={Z'}^{\sim 8}_j\supset{Z'}^{\sim}_j .
\end{equation}

It remains to show $Z_j^{\sim}\subset{Z'}_j^{\sim}$. The enlargement in \cite[(2.41)]{B-CMP119}
adds one layer of cubes of a partition compatible with the ones defining the sets: in
\cite{B-CMP119} the large-field regions are ``unions of $MR_j$-cubes'', and a set is enlarged by
``adding one layer of cubes taken from this family of cubes, in terms of which the sets are
defined''. Whichever of the two compatible partitions is used, the $MR_j$-cubes of $Z_j$ or the
$LMR_{j+1}$-cubes of \cite[p.~264]{B-CMP119}, its cubes are not larger than the
$LMR_{j+1}$-cubes: for the first this is $MR_j\le LMR_{j+1}$, which is part (a), and for the
second it is trivial. As the partitions are compatible, each cube added in forming $Z_j^{\sim}$
touches $Z_j$ and lies in an $LMR_{j+1}$-cube; that $LMR_{j+1}$-cube touches $Z_j$, hence
touches $Z'_j\supset Z_j$, and so lies in ${Z'}_j^{\sim}$. Thus
$Z_j^{\sim}\subset{Z'}_j^{\sim}$, and with \eqref{4.4:eq-collar-avoidance-2} this proves (b).

\emph{Part (c).} By the nesting of the small-field regions and nest members
\cite[(2.1)]{B-CMP119},
\begin{equation*}
\Omega_1\supset\Lambda_1\supset\Omega_2\supset\dots\supset\Omega_k\supset\Lambda_k ,
\end{equation*}
together with its extension by one level \cite[(3.20)]{B-CMP119},
\begin{equation*}
\Lambda_k\supset\Omega_{k+1}\supset\Lambda_{k+1} ,
\end{equation*}
and by part (b), we obtain the following. For $j<k$,
\begin{equation*}
\Lambda_{k+1}\subset\Omega_{k+1}\subset\Omega_{j+1}\subset(Z_j^{\sim})^c .
\end{equation*}
For $j=k$,
\begin{equation*}
\Lambda_{k+1}\subset\Omega_{k+1}\subset({Z'}^{\sim 8}_k)^c\subset(Z_k^{\sim})^c .
\end{equation*}
In both cases $\Lambda_{k+1}\cap Z_j^{\sim}=\emptyset$.

Now let an $\mathbf{R}$-step intervene. By the hypothesis stated in (c), the $\mathbf{T}$-step
that produces $\Omega_{j+1}$ runs the construction of \cite[(3.5)]{B-CMP119} from the new
large-field region $\widehat Z_j:=Z_j\cup\bigcup_i Y_i$ (see \cite{B-CMP122b}). Let
$\widehat Z'_j$ be the union of the $LMR_{j+1}$-cubes meeting $\widehat Z_j$. Since
$\widehat Z_j\supset Z_j$, every $LMR_{j+1}$-cube meeting $Z_j$ meets $\widehat Z_j$, so
$\widehat Z'_j\supset Z'_j$. By the monotonicity of enlargement,
${\widehat Z}'^{\sim 8}_j\supset{Z'}^{\sim 8}_j$. The argument of
\eqref{4.4:eq-collar-avoidance-2} with $\widehat Z_j$ in place of $Z_j$ gives
$\Omega^c_{j+1}\supset{\widehat Z}'^{\sim 8}_j$. Hence, using part (b) for the inclusion
$Z_j^{\sim}\subset{Z'}_j^{\sim}\subset{Z'}_j^{\sim 8}$, we again get
$Z_j^{\sim}\subset\Omega^c_{j+1}$. The two chains displayed above then apply unchanged, and
$\Lambda_{k+1}\cap Z_j^{\sim}=\emptyset$.

\emph{Part (d).} By \cite[(2.41)]{B-CMP119} (p.~261), a term $\mathbf{B}^{(j)}(X)$ satisfies
\begin{equation*}
X\cap\Omega_j\neq\emptyset\quad\text{and}\quad X\cap Z_j^{\sim}\neq\emptyset .
\end{equation*}
If $X\subset\Lambda_{k+1}$ and $j\le k$, part (c) gives $X\cap Z_j^{\sim}=\emptyset$, so
the second condition fails and the term is not a $\mathbf{B}^{(j)}$-term. It is therefore an
$\mathbf{E}^{(j)}$-term or an $\mathbf{R}^{(j)}$-term. In particular, every application of
Lemma~\ref{4.4:lem-twist-stability} to an old term localised in a domain $X\subset Y$ with
$Y\subset\Lambda_{k+1}$ is an application to interior sources, the case of interior sources on
the one-power hereditary class treated in Lemma~\ref{4.4:lem-twist-stability}.
\end{proof}

\begin{corollary}[Interior cluster-expansion terms are unchanged by conditioning]
\label{4.4:cor-interior-cluster-terms}
Let $\mathcal{S}$ denote the conditioned setting and $\mathcal{S}_0$ the whole-lattice setting of
\cite{B-CMP119}, let $\Lambda_{k+1}$ be the level-$(k+1)$ domain of the construction of
\cite{B-CMP119} that appears in the sentence of that paper quoted below, and let
$\mathbf{D}_k$ be the class of localisation domains at level $k$. Let $Y'\in\mathbf{D}_k$ with
$Y'\subset\Lambda_{k+1}$. Then the term $V_k(Y',\cdot)$ of the expansion
\cite[(1.41)]{B-CMP116}, computed in
$\mathcal{S}$, is the same function of the restriction $(\mathbf{U},\mathbf{J},A)|_{Y'}$ as the
term $V_k(Y',\cdot)$ computed in $\mathcal{S}_0$. Here $\mathbf{U}$ and $\mathbf{J}$ are the
sequences of bond variables and of currents, and $A$ is the third argument of the terms of the
expansion \cite[(1.41)]{B-CMP116} (not the Wilson action $A(U)$).

Consequently the sentence of \cite{B-CMP119}
\begin{quote}
the terms of the expansion (I.6.41), with localization domains $Y$ contained in
$\Lambda_{k+1}$, satisfy exactly all the conditions of those lemmas
\end{quote}
holds for the terms of $\mathcal{S}$ with localisation domains contained in $\Lambda_{k+1}$, these
terms being reduced to the corresponding terms of $\mathcal{S}_0$. In this sentence
\textup{(I.6.41)} is the reference of \cite{B-CMP119} to the expansion \cite[(1.41)]{B-CMP116}, and
the lemmas referred to are \cite[Lemmas~1 and~2]{B-CMP116}.

The setting $\mathcal{S}_0$ is the whole-lattice construction of \cite{B-CMP119}, and not the
construction of \cite{B-CMP109} itself. It has the following features:
\begin{itemize}
\item the radii of \cite[(2.28)]{B-CMP119}, proportional to the coupling constant;
\item the degree $p_1(g_k)$ of \cite{B-CMP119};
\item the large-field operations $\mathbf{R}^{(j)}$ with $j\le k_1$ merged, $k_1$ as in
\cite{B-CMP119}.
\end{itemize}
The hypotheses of \cite[Lemmas~1 and~2]{B-CMP116} are reached from $\mathcal{S}_0$ through the
dictionary between the whole-lattice construction of \cite{B-CMP119} and the setting of
\cite{B-CMP116}, which the present reduction leaves unchanged.
\end{corollary}

\begin{proof}
Fix $Y'\in\mathbf{D}_k$ with $Y'\subset\Lambda_{k+1}$. That the term $V_k(Y',\cdot)$ computed in
$\mathcal{S}$ and the term $V_k(Y',\cdot)$ computed in $\mathcal{S}_0$ agree as functions of
$(\mathbf{U},\mathbf{J},A)|_{Y'}$ is obtained from three statements alone:
\begin{itemize}
\item the sequence-operator lemma, which supplies the operators $\mathbb{L}_1$;
\item the further lemma of the localisation argument used together with them;
\item part $(\beta)$ of the second step of the proof of the localisation theorem.
\end{itemize}
The series entering these statements converge at free current $\mathbf{J}$ by clauses (f1) and
(f4) of the one-family statement. This gives the first assertion: $V_k(Y',\cdot)$ is the same
function of $(\mathbf{U},\mathbf{J},A)|_{Y'}$ in the two settings.

For the second assertion, take a term of the expansion \cite[(1.41)]{B-CMP116} whose localisation
domain $Y$ is contained in $\Lambda_{k+1}$. By the first assertion, applied with $Y'=Y$, it
coincides as a function of the restricted data with the corresponding term of $\mathcal{S}_0$.
Therefore the conditions of \cite[Lemmas~1 and~2]{B-CMP116} for the terms named in the quoted
sentence are conditions on terms of the whole-lattice construction $\mathcal{S}_0$, and for these
terms the hypotheses of \cite[Lemmas~1 and~2]{B-CMP116} are reached through the dictionary
between the whole-lattice construction of \cite{B-CMP119} and the setting of \cite{B-CMP116}
named at the end of the statement.

It remains to fix the constants with which \cite[Lemmas~1 and~2]{B-CMP116} are read.

In $\mathcal{S}_0$ the constants written $K^{W'}$ and $K(s)$ in the norms of \cite{B-CMP116} are
those supplied by the two corollaries on the constants of the hereditary norm bounds. These
corollaries carry the factor
\begin{equation*}
e^{5^{d}\kappa_1}
\quad\text{in place of the factor}\quad
e^{16\kappa_1}
\end{equation*}
of \cite{B-CMP116}.

Accordingly, \cite[Lemmas~1 and~2]{B-CMP116} are applied on the operator family of the first of
the two routes. At seams the constants used are those of the transfer of the
analysis of \cite[\S3]{B-CMP109} to the cover, which is printed without proof in
\cite[p.~276]{B-CMP119} and cited as printed.

The terms whose localisation domains straddle the boundary of $\Lambda_{k+1}$, on the spaces
\cite[(3.43)]{B-CMP119}, are not covered by the corollary. For them this book uses the transfer
sentence just named, printed without proof in \cite[p.~276]{B-CMP119}; cited as printed where it
is used.
\end{proof}

\begin{corollary}[Interior large-field terms of the $s$-expectation are unchanged by conditioning]
\label{4.4:cor-interior-large-field}
Let $\mathcal{S}$ be the conditioned setting and $\mathcal{S}_0$ the whole-lattice setting of
\cite{B-CMP119}, and consider the $s$-expectation of \cite[(3.26)]{B-CMP119}, where $s$ is the
letter of that display and not the depth $s$ of the glossary, together with its
representation of the type \cite[(3.47)]{B-CMP119}, whose terms are indexed by regions $X$. Here
$\Omega_{k+1}$ and $\Lambda_{k+1}\subset\Omega_{k+1}$ are the domains of \cite{B-CMP119} at
level $k+1$. Assume:
\begin{itemize}
\item[(a)] the half of the representation and bound of the large-field terms that is stated in
terms of the operations $\mathbf{R}'$ and $\mathbf{B}$, together with the six lemmas on which that
half rests, in the form in which it
gives that the whole-lattice $\mathbf{R}'$-terms of \cite{B-CMP122a,B-CMP122b} are label-free,
in the sense in which that statement uses the word;
\item[(b)] the corollary of this section on the operator arguments of the measure
\cite[(3.27)]{B-CMP119};
\item[(c)] for part (2) only: the reflection part of the covariance statement and clause (f3) of
the one-family statement.
\end{itemize}
Then the following hold.
\begin{itemize}
\item[(1)] Under (a) and (b), every $X$-term of the representation of the type
\cite[(3.47)]{B-CMP119} of the $s$-expectation in \cite[(3.26)]{B-CMP119} whose region satisfies
$X\subset\Lambda_{k+1}$ is the same in $\mathcal{S}$ as in $\mathcal{S}_0$.
\item[(2)] Under (a), (b) and (c), these $X$-terms, $X\subset\Lambda_{k+1}$, are reflected termwise,
which is the input used in verifying \cite[(2.32)]{B-CMP119} in the case of the operation that
\cite{B-CMP119} denotes $R''$; the corresponding
statement for translations by elements of $M\mathbb{Z}^4$ follows by transport.
\end{itemize}
\end{corollary}

\begin{proof}
We prove (1). The $X$-terms in question are computed from the measure \cite[(3.27)]{B-CMP119} at
$t=0$, where $t$ is the parameter of that display and not the letter $t$ that the glossary
assigns to the uniqueness chapter. This
measure has the operator arguments treated in the corollary named in
hypothesis (b) and the extra vertices
\begin{equation*}
\mathbf{R}^{(j)}(X'),\qquad k_1<j\le k,\quad X'\subset X ,
\end{equation*}
where $k_1$ is the level fixed in \cite{B-CMP119}, the operations $\mathbf{R}^{(j)}$ with
$j\le k_1$ being merged into the whole-lattice construction. The operator arguments coincide in
$\mathcal{S}$ and $\mathcal{S}_0$ by hypothesis (b). The extra vertices $\mathbf{R}^{(j)}(X')$ are
treated by the step of the conditioning argument of this section that is also used in the proof
of Corollary~\ref{4.4:cor-interior-cluster-terms} and whose input is that the whole-lattice
$\mathbf{R}'$-terms of \cite{B-CMP122a,B-CMP122b} are label-free; by hypothesis (a) this input
holds, and that step gives that the extra vertices are the same in
$\mathcal{S}$ and in $\mathcal{S}_0$.

The external field $A_k$ of the conditioned setting $\mathcal{S}$ lives on
$\Omega_{k+1}\setminus\Lambda_{k+1}$. It enters the measure through the operators
\begin{equation*}
\chi_{Y}\,\mathsf{A}_{\omega}\,\chi_{\mathbb{B}(\Omega_{k+1})\setminus Y},
\end{equation*}
where $Y$ is as in \cite{B-CMP119}, $\chi$ denotes characteristic functions and
$\mathbb{B}(\Omega_{k+1})$ is the set of bonds of
$\Omega_{k+1}$, and through the localisations that contain bonds of $A_k$. In both cases $A_k$
enters only through terms whose localisation is not contained in $\Lambda_{k+1}$. Hence no
$X$-term with $X\subset\Lambda_{k+1}$ depends on $A_k$, and, its operator arguments and its extra
vertices being the same in $\mathcal{S}$ and $\mathcal{S}_0$ by the preceding paragraph, it is the
same in $\mathcal{S}$ as in $\mathcal{S}_0$. This proves (1).

We prove (2). The termwise reflections of the $X$-terms of (1), which are the input used in
verifying \cite[(2.32)]{B-CMP119} in the case of the operation that \cite{B-CMP119} denotes
$R''$, follow from the two statements of
hypothesis (c): the reflection part of the covariance statement together with clause (f3) of the
one-family statement. The corresponding statement for translations by elements of
$M\mathbb{Z}^4$ holds by transport.
\end{proof}

\begin{corollary}[Local pieces of the conditioned covariance, its root and resolvents]
\label{4.4:cor-covariance-local-pieces}
Let $\mathsf{A}=C^{*}\Delta^{(k)}C$, and let $\{\mathsf{A}_{\omega}\}$ be the local operator family
of the decomposition of $\mathsf{A}$ into local terms. Here $C$ and $\Delta^{(k)}$ are the operators
fixed in the statement that constructs the family $\{\mathsf{A}_{\omega}\}$, and $\mathsf{A}$ is the
operator decomposed there;
$C^{(k)}(Z_0)$ and $C^{(k)}(\Lambda)$ are the conditioned covariances at level $k$ on the regions
$Z_0$ and $\Lambda$; $Y'\in\mathbf{D}_k$ is a localisation domain; $q^{*}$ and $q^{\sharp}$ are the
enlarged cubes around the cube $q$; $T_q$ is the cube operator of the resolvent lemma for the
conditioned covariance, assembled from the $\mathsf{A}_{\omega}$ with $|\omega|\subset q^{\sharp}$;
and $x$, $\lambda$, $\omega$ and $|\omega|$ are as in that lemma and in
the construction of the family $\{\mathsf{A}_{\omega}\}$. The operators below have the local pieces
stated in items (3) and (4); these are items (3) and (4) of the list of local pieces of the
operators of the conditioned step.
\begin{enumerate}
\item[(3)] The operator $\Delta^{(k)}$ is treated through $\mathsf{A}=C^{*}\Delta^{(k)}C$, and it is
$\mathsf{A}$, not $\Delta^{(k)}$, that the family $\{\mathsf{A}_{\omega}\}$ decomposes. The
covariances $C^{(k)}(Z_0)$ and $C^{(k)}(\Lambda)=\mathsf{R}^{Y'}(0)$ are given by the resolvent
lemma for the conditioned covariance. Their local piece is $(x+T_q)^{-1}$ on $Y'\cap q^{*}$.
\item[(4)] The operators $(C^{(k)})^{1/2}$, $(x+\mathsf{A})^{-1}$ and $(\mathsf{A}+\lambda)^{-1}$
are given by clauses (i) and (iii) of the resolvent lemma for the conditioned covariance. Their local
pieces are the same as in item (3).
\end{enumerate}
With these two items, claim $(\alpha)$ of the expansion statement at its step $j$ holds for every
item of that list.
\end{corollary}

\begin{proof}
\emph{Items (3) and (4).} In item (3) the operator $\Delta^{(k)}$ is treated through
$\mathsf{A}=C^{*}\Delta^{(k)}C$, and the construction of the family $\{\mathsf{A}_{\omega}\}$ is a
decomposition of $\mathsf{A}$. Both items then follow from the resolvent lemma for the conditioned
covariance: it gives $C^{(k)}(Z_0)$ and $C^{(k)}(\Lambda)$, the latter as $\mathsf{R}^{Y'}(0)$, and
its clauses (i) and (iii) give $(C^{(k)})^{1/2}$, $(x+\mathsf{A})^{-1}$ and
$(\mathsf{A}+\lambda)^{-1}$. In each case the piece attached to a cube $q$ is $(x+T_q)^{-1}$ on
$Y'\cap q^{*}$.

\emph{Claim $(\alpha)$ for all items.} We apply part (b) of the localisation theorem for conditioned
settings to each item. We take $\mathfrak{s}$ and $\mathfrak{s}'$ to be the whole-lattice settings,
in the sense of that theorem, of two sites, and we take $W$ to be a lifted cube common to both.

The hypotheses of part (b) are met:
\begin{itemize}
\item the locality hypothesis of part (b) is the locality property of the family
$\{\mathsf{A}_{\omega}\}$;
\item the one-family bounds and the two further conditions of part (b) of that theorem are
volume-free, and so are available in both settings $\mathfrak{s}$ and $\mathfrak{s}'$.
\end{itemize}

Part (b) of that theorem then yields claim $(\alpha)$ for each item of that list. This includes items
(3) and (4), whose local pieces were identified above.
\end{proof}

\begin{remark}[Algebraic exactness and the positivity input of localisation]
\label{4.4:rem-exactness-positivity}
Consider the localisation lemma for the conditioned covariance. We note four facts about the
inputs of its proof.
\begin{itemize}
\item[(a)] The exactness of the localisation lemma is algebraic: it follows from clause~(iv) of the
lemma (called the representation lemma in what follows) that represents the conditioned covariance,
its square root and its resolvents through the local operator family $\mathsf{A}_{\omega}$, and on
which the local pieces of Corollary~\ref{4.4:cor-covariance-local-pieces} are built.
\item[(b)] The positivity input of the localisation lemma is the positivity statement at order zero.
\item[(c)] The proof of the localisation lemma uses nothing from \cite[\S E]{B-CMP99b}.
\item[(d)] The proof of the localisation lemma also invokes three further statements; these serve
only the convergence of the expansions in that proof.
\end{itemize}
\end{remark}

\begin{proof}
By clause~(iv) of the representation lemma, the exactness asserted by the localisation lemma holds
as an algebraic identity in the family $\mathsf{A}_{\omega}$. No estimate enters it, and no
operator or bound of \cite[\S E]{B-CMP99b} is used to obtain it.

Item~(c) is read off the appeals in the proof of the localisation lemma: they are clause~(iv) of
the representation lemma, the positivity statement at order zero and the three statements of
item~(d), and none of them is taken from \cite[\S E]{B-CMP99b}.
\end{proof}

\begin{remark}[Conditioning as a change of Ba{\l}aban's domain sequence. Stated; proof
incomplete at the step marked $(\ast)$]\label{4.4:rem-domain-sequence-change}
This remark concerns the conditioned covariances $C^{(k)}(\Lambda)$ and $C^{(k)}(Z_0)$ when they are
taken in the form constructed in \cite[\S E]{B-CMP99b}.

Let $\boldsymbol{\Omega}=(\Omega_0,\dots,\Omega_k)$ be the sequence of domains at level $k$. Let
$\Lambda$ be the conditioning set and $B^k(\Lambda)$ the set formed from $\Lambda$ in
\cite[\S E]{B-CMP99b}; it becomes the last member of the extended sequence below. We put
\begin{equation}\label{4.4:eq-domain-sequence-change-1}
  \widetilde{\boldsymbol{\Omega}}=(\Omega_0,\dots,\Omega_k,\Omega_{k+1}),
  \qquad \Omega_{k+1}=B^k(\Lambda).
\end{equation}
Let $\Lambda'$ be a second conditioning set. Let $\boldsymbol{\Omega}'=(\Omega'_0,\dots,\Omega'_k)$ be
its sequence of domains, and let $\widetilde{\boldsymbol{\Omega}}'$ be formed from $\boldsymbol{\Omega}'$ and
$\Omega'_{k+1}=B^k(\Lambda')$ as in \eqref{4.4:eq-domain-sequence-change-1}.
\begin{itemize}
\item[(a)] In \cite[\S E]{B-CMP99b} the conditioning set enters in only two ways.
  \begin{itemize}
  \item[(i)] Through $\Omega_{k+1}=B^k(\Lambda)$, the last member of $\widetilde{\boldsymbol{\Omega}}$.
  \item[(ii)] It enters through the restriction to $\Lambda$ of the block-local maps $Q$ and $\mu(\cdot)$
  of \cite[(3.169)]{B-CMP99b}, and of $\bar D$ and $\tilde D^{(2)}$.
  \end{itemize}
\item[(b)] Every operator of the expansion of $C^{(k)}(\Lambda)$ in \cite[\S E]{B-CMP99b} is a word
  each of whose letters is of one of three kinds: an operator constructed for
  $\widetilde{\boldsymbol{\Omega}}$ (among them $\tilde G_0$, $\tilde G'$, $\tilde T$, $\tilde Q$,
  $\tilde Q'$); an operator constructed for $\boldsymbol{\Omega}$ ($G'$, $T$, $Q'$, $H_1$); or a
  block-local map of (a).
\item[(c)] Let $\Gamma$ be a region with
  \begin{equation}\label{4.4:eq-domain-sequence-change-2}
    \Gamma\subset B^k(\Lambda)\cap B^k(\Lambda')\cap\Omega_k\cap\Omega'_k .
  \end{equation}
  Consider a term of the random-walk expansion of $C^{(k)}(\Lambda)$ all of whose cubes lie in
  $\Gamma$. This term equals the corresponding term of $C^{(k)}(\Lambda')$.

  This is the cancellation of \cite[Theorem~3.14]{B-CMP99b}, read for the extended sequences
  $\widetilde{\boldsymbol{\Omega}}$ and $\widetilde{\boldsymbol{\Omega}}'$, with the set $\Omega$ of that
  theorem taken to be
  \begin{equation*}
    \Omega=\Omega_{k+1}\cap\Omega'_{k+1}=B^k(\Lambda)\cap B^k(\Lambda').
  \end{equation*}
  In this sense, in \cite[\S E]{B-CMP99b} conditioning on $\Lambda$ is a change of the domain sequence
  from $\boldsymbol{\Omega}$ to $\widetilde{\boldsymbol{\Omega}}$.
\item[(d)] The argument of (a)--(c) does not give the root $(C^{(k)})^{1/2}$. Nor does it give the
  resolvents $(x+\mathsf{A})^{-1}$ of the operator $\mathsf{A}=C^{*}\Delta^{(k)}C$ of
  Corollary~\ref{4.4:cor-covariance-local-pieces} at values $x$ of the resolvent parameter in the
  interval $[0,\gamma_1]$, where $\gamma_1>0$ is the upper end of the interval of resolvent
  parameters at which these resolvents are needed in
  Corollary~\ref{4.4:cor-covariance-local-pieces}. The root and these resolvents are treated in
  Corollary~\ref{4.4:cor-covariance-local-pieces}.
\end{itemize}
\end{remark}

\begin{proof}
\emph{Clause (a).} The construction of \cite[\S E]{B-CMP99b} contains the two sentences:
\begin{quote}
We form a new sequence adding the set $\Omega_{k+1} = B^k(\Lambda)$. Let us denote operators
constructed for this sequence by a wavy line
\end{quote}
The new sequence is $\widetilde{\boldsymbol{\Omega}}$ of
\eqref{4.4:eq-domain-sequence-change-1}.

Apart from this extension, $\Lambda$ occurs in \cite[\S E]{B-CMP99b} only through the restrictions to
$\Lambda$ of the block-local maps $Q$ and $\mu(\cdot)$, introduced in \cite[(3.169)]{B-CMP99b}, and of
$\bar D$ and $\tilde D^{(2)}$. This gives (i) and (ii).

\emph{Clause (b).} We use the lemma giving the decomposition of $G_2^{-1}$ into nine types of terms.
Let $\mathcal{T}$ be the index set of the nine types. The decomposition reads
\begin{equation}\label{4.4:eq-domain-sequence-change-3}
  G_2^{-1}=\tilde G_0^{-1}+W,\qquad
  W=\sum_{\tau\in\mathcal{T}}c_\tau\,K_\tau^{\dagger}M_\tau K'_\tau .
\end{equation}
Each factor $K_\tau$, $M_\tau$, $K'_\tau$ is a word each of whose letters is of one of the three kinds
listed in (b): an operator of the sequence $\widetilde{\boldsymbol{\Omega}}$ (such as $\tilde G'$,
$\tilde T$, $\tilde Q'$); an operator of the sequence $\boldsymbol{\Omega}$ ($G'$, $T$, $Q'$, $H_1$); or a
block-local map of (a).

In particular, the operator $C'^{(k)}(\Lambda)$ of \cite[(3.184)]{B-CMP99b} is given by
\cite[(3.174)]{B-CMP99b} as
\begin{equation}\label{4.4:eq-domain-sequence-change-4}
  C'=Q'\,\tilde{\mathfrak{G}}\,Q'^{*},\qquad
  \tilde{\mathfrak{G}}=(\tilde G')^{2}-(\tilde G')^{2}\,\tilde Q'^{*}\,\tilde T\,\tilde Q'\,(\tilde G')^{2}.
\end{equation}
Here $\tilde{\mathfrak{G}}$ is a word in operators of $\widetilde{\boldsymbol{\Omega}}$.

The remaining steps of the expansion introduce no letter of any other kind. These steps are the
following.
\begin{itemize}
\item The Neumann expansion of $(\tilde Q G_2\tilde Q^{*})^{-1}$ around $(\tilde Q\tilde G_0\tilde Q^{*})^{-1}$
  (the lemma on the inverse of $\tilde Q G_2\tilde Q^{*}$).
\item The representations \cite[(3.186)]{B-CMP99b} and \cite[(3.185)]{B-CMP99b} (the lemma on these
  two representations).
\item The insertion of the elementary expansions (the lemma on their insertion).
\item The exponential gauge of \cite{B-CMP109}, as used in the proposition on the exponential gauge,
  which contributes one more block-local Gaussian variable.
\end{itemize}
Each of these produces only words each of whose letters is of one of the three kinds listed in (b).
This gives (b).

\emph{Clause (c).} $(\ast)$ The proof of the lemma on the sequence operators $\mathbb{L}_1$ applies
verbatim to $\widetilde{\boldsymbol{\Omega}}$ with top index $k+1$; the cover cubes of the new top
member $\Omega_{k+1}=B^k(\Lambda)$ have side $2LM_1$ and lie inside $B^k(\Lambda)$. In a region
$\Gamma$ satisfying \eqref{4.4:eq-domain-sequence-change-2},
three things then hold:
\begin{itemize}
\item both extended sequences $\widetilde{\boldsymbol{\Omega}}$ and $\widetilde{\boldsymbol{\Omega}}'$
  are uniform, in the sense used in the proof of the lemma on the sequence operators
  $\mathbb{L}_1$;
\item every local sequence of domains, in the sense of that proof, is a family of concentric cubes;
\item every letter of the three kinds listed in (b) is the same for the two extended sequences.
\end{itemize}

Now take a term of the random-walk expansion of $C^{(k)}(\Lambda)$ all of whose cubes lie in $\Gamma$.
By (b), it is a word whose letters are of the three kinds listed there. By the third point, all letters
of the word, including the block-local maps, agree for $\Lambda$ and $\Lambda'$.
The term therefore equals the corresponding term of $C^{(k)}(\Lambda')$.

This is the cancellation of \cite[Theorem~3.14]{B-CMP99b}, read for the extended sequences
$\widetilde{\boldsymbol{\Omega}}$ and $\widetilde{\boldsymbol{\Omega}}'$ with
$\Omega=\Omega_{k+1}\cap\Omega'_{k+1}$. Hence, in \cite[\S E]{B-CMP99b},
conditioning on $\Lambda$ is a change of the domain sequence.

\emph{Clause (d).} In the notation of \cite{B-CMP116}, the resolvent-type operator at parameter $x$
is written $\tilde G_3(x)$, and \cite{B-CMP116} states without proof that $\tilde G_3(x)$ has the
same properties as the operator written $\tilde G_2$ there; that sentence is therefore not used here
as a theorem.

Nor does the argument of (a)--(c) supply these properties for $x\in[0,\gamma_1]$: for such $x$ the
product of $x$ with the local form, that is, with the block-local part of the operator whose
resolvent is taken, is not small. For these two reasons the root and the resolvents
$(x+\mathsf{A})^{-1}$ at $x\in[0,\gamma_1]$ are not obtained from this remark. They are obtained by
the route of Corollary~\ref{4.4:cor-covariance-local-pieces}.
\end{proof}

\begin{definition}[Conventions: lattice geometry, block sizes and weighted norms]
\label{4.5:not-geometry-conventions}
Throughout this chapter the following conventions are in force.
\begin{enumerate}
\item \emph{Group.} The dimension is $d = 4$ and $G = \mathrm{SU}(N)$. We write $G^{c}$ for the
complexification of $G$ and $\mathfrak{g}^{c} = \mathfrak{g} \otimes \mathbb{C}$ for the
complexified Lie algebra.

\item \emph{Gauge action.} Ba{\l}aban's gauge action on pairs is
\begin{equation*}
(\mathbf{U},\mathbf{J})^{u} = (\mathbf{U}^{u}, R(u)\mathbf{J})
= (u_{-}\mathbf{U}u_{+}^{-1}, R(u_{-})\mathbf{J}),
\end{equation*}
see \cite[(1.10)]{B-CMP109}; here $R(X)Y = XYX^{-1}$ for an invertible matrix $X$, as in
\cite[(56)]{B-CMP98}. On a bond $b$ with initial point $b_{-}$ and final point
$b_{+}$ it reads
\begin{equation*}
U^{u}(b) = u(b_{-})\,U(b)\,u(b_{+})^{-1}.
\end{equation*}
For two gauge transformations $a$, $c$ put $(ca)(x) := c(x)a(x)$. Then
\begin{equation}\label{4.5:eq-geometry-conventions-1}
(U^{a})^{c} = U^{ca}.
\end{equation}

\item \emph{Scales.} At level $k$ the lattice spacing is $\eta = L^{-k}$, and
$\ell_j := L^{j}\eta$. The sets $\Lambda_j$ and their union
$\mathfrak{B} = \bigcup_j \Lambda_j$, and, for $y, y' \in \mathfrak{B}$, the sets
$\Delta(y)$, $\tilde\Delta(y)$ and the distance $d(y,y')$, are those of
\cite{B-CMP99b}. We write $\ell$ for the operator of multiplication by $\ell_j$ on
$\Lambda_j$ and, for an admissible tower $\{\Omega_j\}$ as below, on
$\Omega_j \setminus \Omega_{j+1}$.

\item \emph{The block size of the propagator papers.} $\mathfrak{M}$ denotes the block size
of the propagator papers, that is, the number written $M_1$ in \cite{B-CMP98},
\cite{B-CMP99a}, \cite{B-CMP99b}, \cite{B-CMP102b} and \cite{B-CMP116}: in
\cite{B-CMP99b} it is the lower bound on the cube size of that paper (written $M$ there)
above which the propagator
theorems hold, with constants depending on $d$ and $L$ only, and in
\cite[p.~4]{B-CMP116} it is the size of the cubes used there in place of $M$.
It is a number depending on $(d,L)$ only, fixed together with the theorems of those papers.
$R_1$ denotes the gap number of \cite[(1.4)]{B-CMP99a} for the block $\mathfrak{M}$:
a sufficiently large positive integer, a power of $L$, such that all theorems on
propagators of the papers cited there hold for $R_1\mathfrak{M}$. The condition is a
condition on the product $R_1\mathfrak{M}$.

\item \emph{The block size of the convergent expansions.} $M_1$ denotes the block size of
\cite{B-CMP119}, in which each domain $\Omega_n$ is a union of cubes of side
$L^{n}\xi M_1$, where $\xi$ is the lattice spacing of \cite{B-CMP119}, and
$\operatorname{dist}(\Omega_n, \Omega_{n-1}^{c}) \ge L^{n}\xi M_1$.
By the statement on the regularity of the window rim in this chapter, $M_1$ is a power of
$L$ with
\begin{equation}\label{4.5:eq-geometry-conventions-2}
M_1 = R'_{\star}\,\mathfrak{M}, \qquad R'_{\star} \ge \max\{R_1, 8\}.
\end{equation}
In particular $\mathfrak{M} \le M_1$. Consequently:
\begin{enumerate}
\item every lower bound on $M$ written with $M_1$, of the form $c\,M_1 \le M$ with
$c \ge 0$, implies the same lower bound with $\mathfrak{M}$ in place of $M_1$;
\item every hypothesis of Ba{\l}aban's papers that names their block is taken at
$\mathfrak{M}$; this applies to the condition in \cite[(3.35)]{B-CMP99b} that $O(1)M$,
in that paper's letter $M$, is a
size of $\square$, to \cite[(1.4)]{B-CMP99a}, and to \cite[(1)]{B-CMP102b} and
\cite[(144)]{B-CMP102b};
\item every upper bound on a collar depth or on a hull side written with $M_1$ holds
a fortiori, since the actual object is built with the block $\mathfrak{M}$.
\end{enumerate}
The single letter $M_1$ of the cited papers thus corresponds to two letters here,
$\mathfrak{M}$ and $M_1$; every other object of those papers keeps its printed letter.

\item \emph{The cube size.} $M = L^{m'}$ is the cube size of \cite{B-CMP109}, chosen
much larger than the previously fixed scales. It enters the statements of this chapter
only through lower bounds.

\item \emph{Towers.} A sequence $\{\Omega_j\}_{j \le k}$ is \emph{admissible} if it is as
in \cite[(1)]{B-CMP102b}, with the block of that display taken at $\mathfrak{M}$. That is,
each $\Omega_j$ is a union of blocks of side $\mathfrak{M}L^{j}\eta$, and
\begin{equation}\label{4.5:eq-geometry-conventions-3}
(L^{j}\eta)^{-1}\operatorname{dist}(\Omega_j^{c}, \Omega_{j+1}) > R\,\mathfrak{M},
\qquad R \ge R_1.
\end{equation}
The scalar $R$ here is the gap ratio of \cite[(1)]{B-CMP102b}; it is unrelated to the
adjoint action $R(\cdot)$ above and to the ball radius $R$ of the glossary.
With each $\Omega_j$ one associates the set of bonds having at least one end-point in
$\Omega_j$. Some or all of the $\Omega_j$ may be equal to the whole lattice $T_\eta$, as in
\cite{B-CMP99a}.

For a term of the effective action localised in a domain $X$ one takes $\Omega_0 := X$
with the maximal tower of condition~(iv) of \cite{B-CMP109}, and $R = R_1$. The lemma of
this chapter on localised terms applies, through its clause on the one-power hereditary
class, to every admissible tower of $X$ whose collar is at most $\tfrac12 M$.

The tower $\{\square_n\}$ built from the window $\square_0$ by \cite[(2.13)]{B-CMP119}
is used in the later sections of this chapter. It is admissible with block $\mathfrak{M}$,
and its gap satisfies
\begin{equation}\label{4.5:eq-geometry-conventions-4}
\ell_n^{-1}\operatorname{dist}(\square_n^{c}, \square_{n+1}) \ge L M_1
> R'_{\star}\mathfrak{M} \ge R_1\mathfrak{M}.
\end{equation}
Here one takes $R := L R'_{\star}$, a power of $L$ larger than $R_1$.

\item \emph{Norms.} For a function $f$ on the hierarchy and a real $\alpha$ we set
\begin{equation*}
|f|_{(\alpha)} := \sup_j \sup \ell_j^{-\alpha}|f|,
\end{equation*}
where the inner supremum is over the arguments of $f$ at level $j$. For a field $A$,
with $\nabla_U$ the covariant gradient in the background $U$, we set
\begin{equation*}
N(A) := \max\bigl\{|A|_{(-1)},\, |\nabla_U A|_{(-2)}\bigr\}.
\end{equation*}
Write $[\nabla_U A]_{(-2-\beta)}$ for the H\"older seminorm of $\nabla_U A$ in the sense
of \cite[(3.40)]{B-CMP99b}, taken with the weight $\ell_j^{2+\beta}$ at level $j$. Then
\begin{equation*}
N_{\beta}(A) := \max\bigl\{|A|_{(-1)},\, |\nabla_U A|_{(-2)},\,
[\nabla_U A]_{(-2-\beta)}\bigr\}.
\end{equation*}
Here $0 \le \beta \le \beta_0 < 1$, with $\beta_0$ fixed.

\item \emph{Pairings and transposes.} $\langle\cdot,\cdot\rangle$ denotes the bilinear
pairings built with $\operatorname{tr}$, with measure $\eta^{d}$ on the finest lattice
and $\ell_j^{d}$ on $\Lambda_j$. The bilinear transpose is ${}^{T}$. The operators
$D^{*}_{U}$ and $Q^{*}$ denote the analytic continuations of the adjoints; these coincide
with the bilinear transposes.

\item \emph{Weighted norms.} For $S \subset \mathfrak{B}$ and every $\kappa \ge 0$ put
$\rho(y) := e^{\kappa d(y,S)}$. Writing the inner supremum in the definition of
$|f|_{(\alpha)}$ as $\sup_{y \in \Lambda_j}\sup_{\Delta(y)}$, we set
\begin{equation*}
|f|_{(\alpha),\rho} := \sup_j \sup_{y \in \Lambda_j} \rho(y)\sup_{\Delta(y)}
\ell_j^{-\alpha}|f|,
\end{equation*}
so that $|f|_{(\alpha)}$ is the case $\rho \equiv 1$. The norm $N^{\rho}$ is defined like
$N$, except that each supremum over $\Delta(y)$ is multiplied by $\rho(y)$.

\item \emph{Units.} Every object is written in the units of its own level. The factor
$\mathfrak{s} = (L^{j-1}\eta)^{3}$ of \cite[(3.12)]{B-CMP109} satisfies
$\mathfrak{s} \le L^{-3}$ and is absorbed into the bound in which it occurs.

\item \emph{Homonyms.} The letter $\varepsilon_1$ is that of \cite{B-CMP109}: it
measures the size of the fluctuation data in \cite[(3.16)]{B-CMP109}, and it keeps this
meaning in the fluctuation bounds of this chapter.

The datum-curvature letter of \cite[(7)]{B-CMP102b} is written
$\varepsilon_1^{(102)}$, the superscript referring to \cite{B-CMP102b}. It is frozen at
the $(d,L)$-value
\begin{equation}\label{4.5:eq-geometry-conventions-a}
\begin{aligned}
\bar\varepsilon_1 &:= \min\Bigl\{a_1^{(102)},\ a_6,\ \tfrac16 R_1 a_1^{(102)},\
\frac{1}{4C_1},\ \frac{2\,\mathfrak{r}_g}{3dLC_Hc_1C_1}\Bigr\},\\
\mathfrak{r}_g &:= \min\Bigl\{1,\ \frac{1}{24\,d\,c_Q},\ L\,\alpha_1^{(98)}\Bigr\}.
\end{aligned}
\end{equation}
The constants in \eqref{4.5:eq-geometry-conventions-a} are as follows.
\begin{itemize}
\item $a_1^{(102)}$, $a_0$ and $B_3$ are the constants $a_1$, $a_0$ and $B_3$ of
\cite[Theorem~1]{B-CMP102b}.
\item $a_6$ is the threshold, and $C_H$ the constant, of the corollary of this chapter
bounding the background potential of the frame construction. The constant $C_H$ is valid
at every $\varepsilon_1^{(102)} \le a_6$. The ratio $M\alpha_0/\varepsilon_1^{(102)}$ entering
$C_H$ is $O(C_1B_3)$ by \cite[(28)]{B-CMP102b}, so $C_H$ depends neither on
$\varepsilon_1^{(102)}$ nor on $M$.
\item $C_1$ is the $(d,L)$-constant of the bound $O(C_1B_3)$ of \cite[(28)]{B-CMP102b}
just quoted.
\item $c_1$ is the constant of \cite[Lemma~2.1]{B-CMP96}, taken at the smallest fraction
at which that lemma is applied in this chapter.
\item $c_Q$ is the constant of the Baker--Campbell--Hausdorff step of this chapter.
\item $\alpha_1^{(98)}$ is the radius of \cite[Proposition~6]{B-CMP98}.
\end{itemize}
The constants are chosen in the order $C_1$, then $C_H$, then $\bar\varepsilon_1$, and
$c_1$, $c_Q$, $\alpha_1^{(98)}$ are fixed before $\bar\varepsilon_1$. Hence
\eqref{4.5:eq-geometry-conventions-a} involves no circularity.

The five entries of $\bar\varepsilon_1$ serve, in order, the following purposes:
\begin{itemize}
\item the uniqueness clause of \cite[Theorem~1]{B-CMP102b};
\item the validity of $C_H$;
\item the first admissibility condition on the initial configuration, in the later step of
this chapter that verifies it;
\item the convergence of the logarithm in the construction of the axial frame;
\item the radius at which a further lemma of this chapter is applied in the first clause of
the lemma on localised terms, and the smallness conditions for the
frame map (the gauge transformation written $u_{k+1}$ in \cite{B-CMP109}).
\end{itemize}
Where \cite[(7)]{B-CMP102b} is used at another admissible value, namely for the real block
data of the real step of the axial and fine axial frame constructions, the letter is
written $\varepsilon_1^{(102),(7)}$, and it is used only there.

The subscript $\beta \in [0,\beta_0]$ of $N_{\beta}$ is a H\"older index and is unrelated
to the enlargement parameter $\beta$ of \cite{B-CMP109} entering $\alpha'''$ below; the
former occurs only as a subscript of $N$.

\item \emph{Enlargement.} The small enlargement parameter $\beta$ of \cite{B-CMP109}, which
appears in $(1+\beta)\alpha_0$ and $(1+2\beta)\alpha_0$, is a fixed positive number
$\le 1$. We put $\alpha''' := (1+2\beta)\alpha$.
\end{enumerate}
\end{definition}

\begin{proof}
We verify the assertions contained in the conventions above.

\emph{The composition rule \eqref{4.5:eq-geometry-conventions-1}.} Let $b$ be a bond. By
the definition of the action, applied first with $a$ and then with $c$,
\begin{align*}
(U^{a})^{c}(b) &= c(b_{-})\,U^{a}(b)\,c(b_{+})^{-1}
= c(b_{-})a(b_{-})\,U(b)\,a(b_{+})^{-1}c(b_{+})^{-1}\\
&= (ca)(b_{-})\,U(b)\,\bigl((ca)(b_{+})\bigr)^{-1} = U^{ca}(b).
\end{align*}
In the last line we used $a(b_+)^{-1}c(b_+)^{-1} = (c(b_+)a(b_+))^{-1}$.

\emph{The inequality $\mathfrak{M} \le M_1$ and its consequences.} By
\eqref{4.5:eq-geometry-conventions-2} we have $R'_{\star} \ge 8 > 1$. Since
$\mathfrak{M} > 0$, it follows that $M_1 = R'_{\star}\mathfrak{M} \ge \mathfrak{M}$.

For the first of the three consequences, let $c \ge 0$ and $cM_1 \le M$. Then
$c\,\mathfrak{M} \le c\,M_1 \le M$.

The second consequence is a convention. The block of each paper quoted there is the number
those papers call $M_1$, and that number is $\mathfrak{M}$ by the convention fixing
$\mathfrak{M}$.

The third consequence holds because the object bounded is built with the block
$\mathfrak{M} \le M_1$, while the bound is written with $M_1$.

\emph{The gap of the window tower \eqref{4.5:eq-geometry-conventions-4}.} The first
inequality, $\ell_n^{-1}\operatorname{dist}(\square_n^{c}, \square_{n+1}) \ge LM_1$, is
the gap supplied by the statement on the regularity of the window rim. For the second,
$L > 1$ and $R'_{\star}\mathfrak{M} > 0$ give
\begin{equation*}
LM_1 = L R'_{\star}\mathfrak{M} > R'_{\star}\mathfrak{M}.
\end{equation*}
The third inequality, $R'_{\star}\mathfrak{M} \ge R_1\mathfrak{M}$, is
$R'_{\star} \ge R_1$ from \eqref{4.5:eq-geometry-conventions-2}.

The number $R = LR'_{\star}$ satisfies
\begin{equation*}
R = L R'_{\star} \ge L R_1 > R_1,
\end{equation*}
and $L M_1 = R\,\mathfrak{M}$. Thus \eqref{4.5:eq-geometry-conventions-4} gives the strict
inequality \eqref{4.5:eq-geometry-conventions-3} for the tower $\{\square_n\}$ with
$R'_{\star} \ge R_1$ in place of $R$, and with $R = LR'_{\star}$ it gives
\begin{equation*}
\ell_n^{-1}\operatorname{dist}(\square_n^{c}, \square_{n+1}) \ge R\,\mathfrak{M}.
\end{equation*}
The argument above verifies the gap condition in the admissibility of $\{\square_n\}$; the
other condition, that each $\square_n$ is a union of blocks of side $\mathfrak{M}L^{n}\eta$,
is part of the convention stated in the item on towers and is not derived here.

\emph{The factor $\mathfrak{s}$.} We have $\eta = L^{-k}$ and $j \le k$. Hence
\begin{equation*}
L^{j-1}\eta = L^{j-1-k} \le L^{-1},
\end{equation*}
and therefore $\mathfrak{s} = (L^{j-1}\eta)^{3} \le L^{-3}$.

\emph{The first entry of \eqref{4.5:eq-geometry-conventions-a}.} The list of constants of
\cite[Theorem~1]{B-CMP102b} contains the inequality $B_3a_1 \le a_0$. The uniqueness
clause of that theorem states that the orbit is a unique critical orbit in the space
\cite[(6)]{B-CMP102b} if $B_3\varepsilon_1 \le \varepsilon_0$ and
$\varepsilon_0 \le a_0$.

Let $\varepsilon_1^{(102)} \le a_1^{(102)}$ and put
$\varepsilon_0 := B_3\varepsilon_1^{(102)}$. The first condition then holds with equality.
The second follows from
\begin{equation*}
\varepsilon_0 = B_3\varepsilon_1^{(102)} \le B_3a_1^{(102)} \le a_0 .
\end{equation*}
So the uniqueness clause holds at every $\varepsilon_1^{(102)} \le a_1^{(102)}$, in
particular at $\bar\varepsilon_1$.

\emph{The second entry.} Since $\bar\varepsilon_1 \le a_6$, the constant $C_H$ is valid
at $\bar\varepsilon_1$.

The third and fourth entries enter no inequality verified here; they are the thresholds
used at the places listed above.

\emph{The fifth entry and $\mathfrak{r}_g$.} From
$\bar\varepsilon_1 \le 2\mathfrak{r}_g/(3dLC_Hc_1C_1)$ we obtain
\begin{equation}\label{4.5:eq-geometry-conventions-5}
3dLC_Hc_1\cdot\tfrac12 C_1\,\bar\varepsilon_1 \le \mathfrak{r}_g .
\end{equation}
Thus any quantity bounded by the left side of \eqref{4.5:eq-geometry-conventions-5} is
at most $\mathfrak{r}_g$. This is the form in which the bound on the background quantity
entering the frame map is used in this chapter, both for its supremum over a cube and for
its norm $N$.

For a quantity $r$ with $0 \le r \le \mathfrak{r}_g$, the three entries of
$\mathfrak{r}_g$ give the following three inequalities:
\begin{align*}
r &\le 1,\\
d\,c_Q\,r &\le \tfrac{1}{24},\\
r/L &\le \alpha_1^{(98)}.
\end{align*}
The first gives that the norm $N$ of that quantity is at most $\mathfrak{r}_g \le 1$. The
second is the smallness condition of the
Baker--Campbell--Hausdorff step. The third places the quantity within the radius of
\cite[Proposition~6]{B-CMP98}.
\end{proof}

\begin{definition}[The class of admissible background pairs and its orbit extension]
\label{4.5:def-admissible-pairs}
Notation is that of the conventions fixed in Definition~\ref{4.5:not-geometry-conventions}. In
particular
$\ell_j = L^{j}\eta$, $\ell$ is multiplication by $\ell_j$ on $\Lambda_j$, $\mathfrak{M}$ is the
block size of the propagator papers, $N_{U}(\cdot)$ and $|\cdot|_{(\alpha)}$ are the scaled norms
over the hierarchy, and $(\mathbf{U},\mathbf{J})^{u} = (\mathbf{U}^{u}, R(u)\mathbf{J})$ is the gauge
action, with $R(u) = \operatorname{Ad}_u$ the adjoint action. Let $a_0$, $a_1$, $a_\gamma$ be
positive radii depending on $d$ and $L$ only; they are fixed in the course of this chapter.

\smallskip
\noindent\emph{The class $\mathfrak{S}$.} It consists of the pairs $(U,J)$ with the following
properties.
\begin{itemize}
\item[(S1)] $U = e^{i\eta\mathfrak{a}}U_r$, where $\mathfrak{a}$ is a $\mathfrak{g}^{c}$-valued bond
function and $U_r$ is $G$-valued and satisfies
\cite[(3.35)]{B-CMP99b}. That is, for every $j$ and every cube $\square$ at level $j$ of the class on
which that condition is imposed, there is a $G$-valued gauge transformation $u$ on $\square$ with
the properties listed in \eqref{4.5:eq-admissible-pairs-a}. These cubes are those of size
$O(1)\mathfrak{M}$ in level-$j$ units. By \cite[p.~409]{B-CMP99b} the number $O(1)$ may be taken
equal to $12$, so that the enlarged cube $\square^{5}$ of \cite[p.~409]{B-CMP99b}, of side
$11\mathfrak{M}$ in level-$j$ units, is contained in one of the cubes on which the condition holds.
\item[(S2)] The twist $\mathfrak{a}$ satisfies \cite[(3.37)]{B-CMP99b} in the form displayed below.
\item[(S3)] $J$ is a $\mathfrak{g}^{c}$-valued bond function obeying the bound displayed below.
\end{itemize}
In formulas:
\begin{equation}\label{4.5:eq-admissible-pairs-a}
\begin{aligned}
&\text{(S1)}\quad U = e^{i\eta\mathfrak{a}}U_r,\qquad U_r^{u} = e^{i\eta A},\qquad
\ell_j|A| < a_0,\quad \ell_j^{2}|\nabla A| < a_0 \ \text{ on } \square,\\
&\text{(S2)}\quad N_{U_r}(\mathfrak{a}) < a_1,\\
&\text{(S3)}\quad |J|_{(-3)} < a_\gamma .
\end{aligned}
\end{equation}
No condition corresponding to \cite[(3.36)]{B-CMP99b} or to \cite[(3.38)]{B-CMP99b} is imposed, and
no relation between $J$ and $U$ is required. Conditions (S1) and (S2) give, at level $j$,
\begin{equation}\label{4.5:eq-admissible-pairs-1}
|\partial U - 1| \le 4(a_0 + a_1)\,\eta^{2}\ell_j^{-2}.
\end{equation}

\smallskip
\noindent\emph{The data $B$.} Let $c_9$ be the radius of the data $B$, a positive constant. A
datum $B$ is a $\mathfrak{g}^{c}$-valued function on $\mathfrak{B}$ with $|B| := \sup|B| < c_9$.
We put $\tilde B := \ell^{-1}B$.

\smallskip
\noindent\emph{The orbit extension.} Put $\mathfrak{K} := 4$ and
\begin{equation*}
\mathfrak{S}^{\mathrm{orb}} := \bigl\{ ((U,J)^{u}, B^{\mathrm{orb}}) : (U,J) \in \mathfrak{S},\
\|\operatorname{Ad}_u^{\pm 1}\| \le \mathfrak{K},\ |B^{\mathrm{orb}}| < c_9/\mathfrak{K} \bigr\},
\end{equation*}
where $B^{\mathrm{orb}}$ is a $\mathfrak{g}^{c}$-valued function on $\mathfrak{B}$ and
$|B^{\mathrm{orb}}| := \sup|B^{\mathrm{orb}}|$.
Here $u$ ranges over the $G^{c}$-valued gauge transformations, and
$\|\operatorname{Ad}_u^{\pm 1}\|$ is the supremum over sites of the operator norms of
$\operatorname{Ad}_{u(x)}^{\pm 1}$ with respect to the norm on $\mathfrak{g}^{c}$ in which
$|B| = \sup|B|$ is taken. At every point of $\mathfrak{S}^{\mathrm{orb}}$,
and for every representative, the datum pulled back to the representative satisfies
\begin{equation}\label{4.5:eq-admissible-pairs-2}
|R(\bar u^{-1})B^{\mathrm{orb}}| \le \mathfrak{K}\,|B^{\mathrm{orb}}| < c_9 .
\end{equation}
Here $\bar u$ denotes the gauge transformation on the hierarchy $\mathfrak{B}$ associated
with $u$.

\smallskip
\noindent\emph{Inclusion of the classes of \cite{B-CMP109}.} Let $X$ be a domain as in
\cite[(1.11)--(1.14)]{B-CMP109} (the set $\Omega_0$ of
Definition~\ref{4.5:not-geometry-conventions}), let $(\alpha_0,\alpha_1,\gamma_0)$ be
the parameters of conditions (i)--(iii) of \cite[(1.11)--(1.14)]{B-CMP109} on $X$, let
$c_{12}$ be the constant $O(1)$ of
\cite[(1.12)]{B-CMP109}, and let $B^{109}$ be the large constant written $B$ in
\cite[(1.12)]{B-CMP109}. Assume the ceilings
\begin{equation*}
c_{12}LMB^{109}\alpha_0 \le a_0, \qquad \alpha_1 \le a_1, \qquad \gamma_0 \le a_\gamma .
\end{equation*}
Assume also the floors
\begin{equation}\label{4.5:eq-admissible-pairs-3}
17M_1 \le LM, \qquad 17M_1 \le c_{12}LM, \qquad c_{12}L \ge 2 .
\end{equation}
Then every pair $(\mathbf{U},\mathbf{J})$ satisfying conditions (i)--(iii) of
\cite[(1.11)--(1.14)]{B-CMP109} on $X$, with parameters $(\alpha_0,\alpha_1,\gamma_0)$, satisfies
(S1)--(S3). Condition (iv) of \cite{B-CMP109} is not part of the definition of $\mathfrak{S}$ and is
not derived from (S1)--(S3); where it is used, it is supplied separately.
\end{definition}

\begin{proof}
\emph{The bound \eqref{4.5:eq-admissible-pairs-2}.} Let $((U,J)^{u}, B^{\mathrm{orb}})$ be a point of
$\mathfrak{S}^{\mathrm{orb}}$, presented through a representative $(U,J) \in \mathfrak{S}$ and a
$G^{c}$-valued $u$ with $\|\operatorname{Ad}_u^{\pm1}\| \le \mathfrak{K}$. The pulled-back datum is
$R(\bar u^{-1})B^{\mathrm{orb}}$. Since $\bar u$ takes at every element of $\mathfrak{B}$ a value of
$u$, and $R(\bar u^{-1})$ acts at each element by the adjoint action of the inverse of that value,
the bound $\|\operatorname{Ad}_u^{-1}\| \le \mathfrak{K}$ gives the norm bound $\mathfrak{K}$ for the
operator $R(\bar u^{-1})$ on $\mathfrak{g}^{c}$-valued functions on $\mathfrak{B}$ in the norm
$\sup|\cdot|$; this is the first inequality of \eqref{4.5:eq-admissible-pairs-2},
$|R(\bar u^{-1})B^{\mathrm{orb}}| \le \mathfrak{K}\,|B^{\mathrm{orb}}|$. Since
$|B^{\mathrm{orb}}| < c_9/\mathfrak{K}$,
\begin{equation*}
|R(\bar u^{-1})B^{\mathrm{orb}}| \le \mathfrak{K}\,|B^{\mathrm{orb}}|
< \mathfrak{K}\cdot c_9/\mathfrak{K} = c_9 .
\end{equation*}
So the pulled-back datum belongs to the class of data $B$ defined above.

\emph{The inclusion.} In the notation of \cite{B-CMP109} the lattice spacing is $\xi$; here it is
$\eta$. We use three of the conditions there.
\begin{itemize}
\item Condition (i) states that $\mathbf{U} = U'U_r$ with $U_r$ $G$-valued (the factor written $U$
in \cite{B-CMP109}) and
$|\partial U_r - 1| < \alpha_0\xi^{2}$ on $X$ \cite[(1.11)]{B-CMP109}. It further states that for
each cube $\square \subset X$ of size $c_{12}LM$ there is a $G$-valued gauge transformation $u$ on
$\square$ with $U_r^{u} = \exp i\xi A$ and with $|A|$, $|\nabla^{\xi}A|$ less than
$c_{12}LMB^{109}\alpha_0$ on $\square$ \cite[(1.12)]{B-CMP109}. Here $\nabla^{\xi}$ is the
plain lattice gradient.
\item Condition (ii) states that $U' = \exp i\xi A'$, where $A'$ is $\mathfrak{g}^{c}$-valued and
$|A'|, |\nabla^{\xi}_{U_r}A'| < \alpha_1$ on $X$ \cite[(1.13)]{B-CMP109}.
\item Condition (iii) states that $|\partial\mathbf{U} - 1| < \alpha_0\xi^{2}$ and
$|\mathbf{J}| < \gamma_0$ on $X$ \cite[(1.14)]{B-CMP109}, where $\mathbf{J}$ is the
$\mathfrak{g}^{c}$-valued bond function of the pair, as in \cite{B-CMP109}.
\end{itemize}
In the letters of the definition the configuration is $\mathbf{U}$. Set $\mathfrak{a} := A'$, so
that $\mathbf{U} = U'U_r = e^{i\eta\mathfrak{a}}U_r$, and set $J := \mathbf{J}$.

For (S1), recall the cubes on which the constructions of \cite{B-CMP99b} use the configuration.
The random-walk representations there use gauges defined on the enlarged cubes $\square^{5}$
\cite[p.~409]{B-CMP99b} and boundary conditions imposed outside the enlarged cubes $\square^{8}$
\cite[p.~418]{B-CMP99b}. These have sides
$11\mathfrak{M}$ and $17\mathfrak{M}$ in level-$j$ units. Since $\mathfrak{M} \le M_1$, both sides
are at most $17M_1$. The cubes of the class in (S1) have side $12\mathfrak{M} \le 17M_1$ in
level-$j$ units as well, so the following argument covers them together with the enlarged cubes.

Let $\square$ be a cube at level $j$ whose side is at most $17M_1$ in level-$j$ units, that is, a
cube of the class in (S1) or a cube $\square^{5}$ or $\square^{8}$. Its side is at most
$17M_1\ell_j$, and since $\ell_j \le 1$ this is at most
$17M_1$ in the units of the top lattice. By the second floor in \eqref{4.5:eq-admissible-pairs-3},
it therefore lies in one cube of side $c_{12}LM$ of \cite[(1.12)]{B-CMP109}. On that cube,
condition (i) provides a $G$-valued gauge transformation $u$ with $U_r^{u} = e^{i\eta A}$ and with
$|A|$, $|\nabla A|$ less than $c_{12}LMB^{109}\alpha_0$. Restrict it to $\square$. Since
$\ell_j \le 1$,
\begin{equation*}
\ell_j|A| \le |A|,
\qquad \ell_j^{2}|\nabla A| \le |\nabla A|,
\end{equation*}
and both right-hand sides are less than $c_{12}LMB^{109}\alpha_0$, which is at most $a_0$ by the
first ceiling.
This is (S1).

For (S2), we have $N_{U_r}(\mathfrak{a}) = \max\{|\mathfrak{a}|_{(-1)}, |\nabla_{U_r}\mathfrak{a}|_{(-2)}\}$.
The weights $\ell_j$ and $\ell_j^{2}$ in these norms do not exceed $1$. Condition (ii) and the
second ceiling therefore give $N_{U_r}(\mathfrak{a}) < \alpha_1 \le a_1$.

For (S3), $J$ is $\mathfrak{g}^{c}$-valued. Since $\ell_j^{3} \le 1$, condition (iii) and the third
ceiling give
\begin{equation*}
|J|_{(-3)} \le \sup|J| < \gamma_0 \le a_\gamma .
\end{equation*}
\end{proof}

\begin{definition}[Reading conventions for the cited hypotheses]
\label{4.5:not-reading-conventions}
The following six reading conventions, $(\rho1)$--$(\rho6)$, are in force wherever the cited
hypotheses and conclusions of Ba{\l}aban's papers are applied in this chapter.
\begin{itemize}
\item[$(\rho1)$] Condition \cite[(1.12)]{B-CMP109} is imposed there, for a localisation domain
$X\in\mathbf{D}_j$ of \cite{B-CMP109}, on the cubes $\square\subset X$
of a size $O(1)LM$. Write $c_{12}$ for the constant $O(1)$ of that condition. The condition is read
as imposed on every cube $\square\subset X$ of size at most $c_{12}LM$.
\item[$(\rho2)$] The tower $\{\square_n\}$ of the window $\square_0$ is the tower of
\cite[(2.13)]{B-CMP119}. The same tower is the one read in clause (iv), condition
\cite[(1.16)]{B-CMP109}, of the space
\begin{equation*}
\mathfrak{U}=U'^{c}_{k+1}\bigl(\square_0,(1+2\beta)\alpha_0,(1+2\beta)\alpha_1,\alpha_0\bigr)
\end{equation*}
of \cite{B-CMP109}, where $\beta$ is the parameter of \cite[\S3]{B-CMP109} (not the coupling
$\beta_{k+1}$). It is
admissible in the sense of \cite[(1)]{B-CMP102b}, with block $\mathfrak{M}$. The clause on the
one-power hereditary class in the inductive lemma for the small-field conditions covers this tower
wherever clause (iv) is produced for it, provided every collar $X\setminus\square_n$ has depth at
most $\tfrac12 M$, in level-$j$ units, in the domain $X$ on which clause (iv) is read; this is
checked at the place where clause (iv) is produced, and imposes no condition on the parameters
beyond those already in force. This single rule fixes the tower of the window wherever a tower
of $\square_0$ is read in this chapter.
\item[$(\rho3)$] In \cite[Theorem~3.10]{B-CMP99b} the configurations $U$ are required to satisfy
\cite[(3.95)]{B-CMP99b}. This requirement is read as condition \cite[(3.35)]{B-CMP99b}.
\item[$(\rho4)$] The set \cite[(3.16)]{B-CMP109} is read at representatives. In \cite{B-CMP109},
the space $U^{c}_{j}(X,\alpha_0,\alpha_1,\gamma_0)$ is a union of orbits
$[(\mathbf{U},\mathbf{J})]$ \cite[p.~261]{B-CMP109} under all $G^{c}$-valued gauge maps. The
inductive lemma for the small-field conditions, however, and the bound $|A|<\alpha_2$ of clause
(ii), on the $\mathfrak{g}^{c}$-valued field $A$ by which the background is shifted, are
statements about a single representative, and they are not invariant under $\mathrm{Ad}$
beyond $\|\mathrm{Ad}_u\|\le 3/2$. Accordingly, the far-term statement asks that one representative,
modulo $G$-valued gauge maps, satisfy both members of \cite[(3.16)]{B-CMP109}. It passes to the
orbit through the gauge invariance of the far functional under $G$-valued maps $u$ and under maps
$u=u'u_G$ with $u_G$ $G$-valued and $\|\mathrm{Ad}_{u'}^{\pm1}\|\le\mathfrak{K}$, and not beyond.
Ba{\l}aban's own condition \cite[(3.14)]{B-CMP109} is read in the same way.
\item[$(\rho5)$] The bound \cite[(172)]{B-CMP102b} is applied to a datum $e^{iB}V_0$, with $B$ a
$\mathfrak{g}^{c}$-valued field on the block-bond set $\mathfrak{B}$ and $V_0$ the configuration
being twisted, with the margin
\begin{equation*}
\sup|B|\le \tfrac12 C_{1}\bar\varepsilon_1 ,
\end{equation*}
where $C_1$ is the constant of \cite[(172)]{B-CMP102b} and $\bar\varepsilon_1$ is the frozen
datum radius.
\item[$(\rho6)$] The frame map $v$ of the axial gauge used in this chapter is the inverse of the map
denoted $v$ in Ba{\l}aban's axial-gauge construction, so that our $v$ equals Ba{\l}aban's
$v^{-1}$.
\end{itemize}
\end{definition}

\begin{remark}[Bounds and walk expansion of the constrained covariance inverse; printed without proof
in {\cite[\S3, before (3.132)]{B-CMP99b}}]
\label{4.5:rem-covariance-inverse-walks}
Let $U$ satisfy conditions (S1) and (S2) of Definition~\ref{4.5:def-admissible-pairs}.
Here $U=e^{i\eta\mathfrak{a}}U_r$ is in general $G^{c}$-valued. Put
\begin{equation*}
\mathsf{q}_0 := (QG_0Q^{*})^{-1},
\end{equation*}
where $Q$ is the averaging operator of \cite[\S3]{B-CMP99b} in the background $U$, $G_0$ is the
Green's operator of this chapter at vanishing current $J$, which is the operator $G$ of
\cite[\S3]{B-CMP99b} in the background $U$,
and $Q^{*}$ is the analytic
continuation of the adjoint of $Q$, that is, its bilinear transpose. Two properties of
$\mathsf{q}_0$, for every such $U$, are used.
\begin{enumerate}
\item The operator $\mathsf{q}_0$ obeys the bounds \cite[(3.132)]{B-CMP99b}.
\item The operator $\mathsf{q}_0$ has a generalised random-walk expansion of the type
\cite[(3.98)--(3.99)]{B-CMP99b}, obtained as in the construction of
\cite[Theorem~3.9]{B-CMP99b}. Each term of the expansion depends on $U$ only through the values of
$U$ on the hulls of the corresponding walk, and the terms obey bounds of the type
\cite[(3.108)]{B-CMP99b}.
\end{enumerate}
Both properties are stated in \cite[\S3, the sentence preceding (3.132)]{B-CMP99b} without proof,
and are cited as printed. The sentence reads:
\begin{quotation}
The operators $(QGQ^{*})^{-1}$, or $(QG_1Q^{*})^{-1}$, can be analyzed in the same way as the
operator $(Q'G'^{2}Q'^{*})^{-1}$. We will not repeat these considerations here, let us write
only bounds
\end{quotation}
after which the bounds \cite[(3.132)]{B-CMP99b} are displayed.
The same two properties are stated in \cite{B-CMP102b} for the operator written $H_0$ there.
\end{remark}

\begin{remark}[{Analyticity of the derivative bounds for the averaging maps; printed without proof
in \cite[\S E, p.~43]{B-CMP98}}]
\label{4.5:rem-averaging-analyticity}
Let $Q_k$ and $C_k$ be the maps of the averaging construction of \cite{B-CMP98}, in the notation
there. In \cite[Proposition~5]{B-CMP98} the derivatives $\delta Q_k/\delta A$ and
$\delta C_k/\delta A$ with respect to the field $A$ satisfy the bounds
\cite[(156)--(157)]{B-CMP98}.

The extended form used in this book is the following. Let $U'U_0$ be a
configuration of the product form considered in \cite[\S E, (158)--(163)]{B-CMP98}, in which
$U_0$ is a background configuration and the factor $U'$ is determined by the field $A'$. The
derivatives $\delta Q_k/\delta A$ and $\delta C_k/\delta A$ satisfy the bounds
\cite[(156)--(157)]{B-CMP98} at $U'U_0$. They are analytic in $A'$, and these bounds hold
uniformly in $A'$.

This extension is stated in \cite[\S E, after Proposition~6 and (164), p.~43]{B-CMP98}
without proof, and it is cited as printed:
\begin{quote}
Similarly, Proposition 5 may be extended to include analyticity and uniformity statements. The
formulations are obvious
\end{quote}
\end{remark}

\begin{remark}[Small curvature of block averages at near-unitary complex data;
printed without proof in {\cite[p.~263 and (3.26), p.~275]{B-CMP109}}]
\label{4.5:rem-block-curvature-complex}
In this remark $p$ denotes a level index of the tower of block averages, not a plaquette. Let
$\mathbf{U} = e^{i\xi A'}U$ be a complex configuration satisfying the three near-unitarity
conditions imposed on complex configurations in this chapter, with constants
$\alpha_0^{\bullet},\alpha_1^{\bullet}$, the bond field $A'$ and the parameter $\xi$ being those of
these conditions; the third of these conditions is a bound on the curvature of
$\mathbf{U}$ and is called the curvature condition below. Assume moreover
\begin{equation*}
\alpha_1^{\bullet} \le \alpha_1^{(98)},
\end{equation*}
where $\alpha_1^{(98)}$ denotes the smallness threshold on the parameter $\alpha_1$ of
\cite{B-CMP98}. Let $M^{\cdot}(\mathbf{U})$ be the
tower of block averages $\bar{\mathbf{U}}^{p}$ on the layers $\Lambda_p$. The bounds on block
averages of \cite[Propositions~1 and~2]{B-CMP98} and the axial-gauge bounds of
\cite[Section~F]{B-CMP102b} hold at such data in the following form.
\begin{itemize}
\item[(a)] The block plaquettes within each level are small: on $\Lambda_p$,
\begin{equation*}
|\partial\bar{\mathbf{U}}^{p}-1| \le C_{98}\,\alpha_0^{\bullet}(L^{p}\xi)^{2},
\end{equation*}
with an absolute constant $C_{98}$.
\item[(b)] For the generalised axial gauge of $M^{\cdot}(\mathbf{U})$ on the hierarchy
$\mathfrak{B}(\square_0)$ of the window $\square_0$, the axial representative $V$ of
$M^{\cdot}(\mathbf{U})$ satisfies
\begin{equation*}
|V(c)-1| \le C_{V}\,M\alpha_0^{\bullet}\qquad\text{for every } c\in\mathfrak{B}(\square_0).
\end{equation*}
\end{itemize}

Statement (a) is the form at complex data, with the constant $2$ replaced by $C_{98}$, of
Ba{\l}aban's sentence on the tower of a complex configuration $\mathbf{U}$ (boldface in the
notation of \cite{B-CMP109}), stated in \cite[p.~263]{B-CMP109} without proof; cited as printed:
\begin{quote}
We have $M^{\cdot}(\mathbf{U}) = \bar{\mathbf{U}}^{p}$ on $\Lambda_p$,
$|\partial\bar{\mathbf{U}}^{p}-1| < 2\alpha_0'(L^p\xi)^2$.
\end{quote}
Statement (b) is Ba{\l}aban's display read at the same complex configuration $\mathbf{U}$,
stated in \cite[(3.26), p.~275]{B-CMP109} without proof; cited as printed:
\begin{quote}
$M^{\cdot}(\mathbf{U}) = V^{v}$, $|V(b)-1| < O(1)M\alpha_0$.
\end{quote}
Here $v$ is the gauge transformation with $M^{\cdot}(\mathbf{U}) = V^{v}$, so that $V$ is the axial
representative of the tower, and $b$ runs over the bonds of the hierarchy; the letters
$\alpha_0'$ and $\alpha_0$ in the quotations are those of \cite{B-CMP109}.

The published proofs of these bounds cover two cases only. For fields with values in
$\mathrm{U}(N)$ the
bounds on block averages are \cite[Propositions~1 and~2]{B-CMP98}, and the constant $2$ in the
sentence quoted from \cite[p.~263]{B-CMP109} is the constant of \cite[(54)]{B-CMP98} at real
fields; at complex data it is replaced by the constant $C_{98}$ of (a). For the real minimiser the
axial-gauge bounds are those of \cite[Section~F]{B-CMP102b}: on the two top levels of the hierarchy
\cite[(160)]{B-CMP102b},
\begin{quote}
$|B(x,x')| < (8d^2L^2+4L^2|x-y|)\varepsilon_1$,
\end{quote}
and on the lower layers \cite[(145), (146), (151)]{B-CMP102b}, the last of which reads
\begin{quote}
$|V''-1| < 9dL^2M\varepsilon_0-\varepsilon_0$ on $\mathfrak{C}_k$,
\end{quote}
with the letters $B$, $V''$, $\varepsilon_0$, $\varepsilon_1$ of these two quotations being those of
\cite[Section~F]{B-CMP102b}, and
where, in the notation of \cite[(148)]{B-CMP102b}, $\Lambda'_j$ denotes the set of sites of the
hierarchy at level $j$, and the set of sites of the hierarchy is their union
\begin{equation*}
\mathfrak{C}_k = \bigcup_{j\le k}\Lambda'_j ,
\end{equation*}
the set of sites of the hierarchy whose block bonds form $\mathfrak{B}(\square_0)$; these in turn
rest on
\cite[(1.65), Lemma~1]{B-CMP99a}. The proof in \cite{B-CMP102b} uses its curvature
assumption \cite[(7)]{B-CMP102b}:
\begin{quote}
Now we use the assumption (7) for $V$, hence for $V'$
\end{quote}
(\cite[p.~303]{B-CMP102b}). Without a curvature hypothesis, \cite{B-CMP98} states that
a complex perturbation adds
\begin{quote}
$|Y_x| = O(L^2\alpha_0+L\alpha_1)$
\end{quote}
in the notation of that paper; the curvature condition is what removes the term $L\alpha_1$, and
this is why (a) and (b) carry $\alpha_0^{\bullet}$ alone.

\emph{Route.} A proof at near-unitary complex data would re-run \cite[(45)--(53)]{B-CMP98} and
\cite[(145)--(160)]{B-CMP102b} for $\mathrm{GL}(N,\mathbb{C})$-valued data subject to
\begin{equation*}
\|V(b)^{\pm1}\|\le 1+\delta
\end{equation*}
with $\delta$ small. The local axial gauge $v_0(x) = v_0(y)\,V(\Gamma_{y,x})$, with
$\Gamma_{y,x}$ the contour of the axial gauge from $y$ to $x$, needs only invertibility;
\cite[(47)--(50)]{B-CMP98} are holomorphic identities; logarithms are taken by the power series
about $1$. The transports entering the averages are to be bounded uniformly in $p$ by
\cite[(159), (161), (163)]{B-CMP98}: for every level index $q<p$ (here $q$ is a level index, not
the auxiliary parameter of that name),
\begin{equation*}
\|\bar{\mathbf{U}}^{q}(c)^{\pm1}\|\le e^{c_Q\alpha_1^{\bullet}\ell_q},
\end{equation*}
with $\ell_q$ the factor of level $q$ in these bounds, and the product of
these bounds over the levels $q<p$ is at most $e^{2c_Q\alpha_1^{\bullet}}$. Since the leading term of
\cite[(50)]{B-CMP98} appears conjugated by these transports, its coefficient at real fields is
multiplied by a factor at most $e^{2c_Q\alpha_1^{\bullet}}$; the requirement
\begin{equation*}
e^{2c_Q\alpha_1^{\bullet}}\le \tfrac32
\end{equation*}
is an upper bound on $\alpha_1^{\bullet}$ depending on $(d,L)$ only, with no condition on the number
of renormalisation steps, and along this route the constant in (a) is absolute and of the form
$C_{98} = 2+O(\alpha_1^{\bullet})$.

In this chapter statements (a) and (b) are used at one place only: in the lemma constructing the
axial gauge of the window background through the gauge transformation $v$ above, which brings the
tower to the axial gauge, and the map from the
Landau gauge to the axial gauge, (a) in its bound on the plaquettes of the tower and (b) in its
pointwise bound on the axial representative. In particular they are not
used by the lemma on the fine-lattice axial frame, which obtains bounds of the same form, depending
on $\alpha_0^{\bullet}$ alone and linear in $M$, from the curvature condition directly, and hence
not by the bound among the operator bounds at free current that is obtained through that lemma.
\end{remark}

\begin{remark}[The cited sentences under the near bound and the choice of frame]
\label{4.5:rem-near-bound-frame}
The near bound among the operator bounds at free current is the pair of bounds
\cite[(3.31), (3.32)]{B-CMP109}. These are bounds in a frame in which the window background
$U_{k+1}(\square_0,\cdot)$ has the form $e^{i\eta\mathcal{A}}$. Here $\mathcal{A}$ is small, and
its bound involves $\alpha_0$ only and is linear in $M$. Three sentences that Ba{\l}aban states
without proof stand under the near bound: those of
Remarks~\ref{4.5:rem-covariance-inverse-walks}, \ref{4.5:rem-averaging-analyticity}
and~\ref{4.5:rem-block-curvature-complex}, which are stated in \cite{B-CMP99b},
\cite{B-CMP98} and \cite[(3.26)]{B-CMP109}, respectively, without proof. The first two are used
essentially; the third is used only in the frame of \cite[(3.26)--(3.28)]{B-CMP109}.
Concerning this frame:
\begin{enumerate}
\item In Ba{\l}aban's literal frame \cite[(3.26)--(3.28)]{B-CMP109}, a bound on the axial data
$V$ of \cite[(3.26)]{B-CMP109} in terms of $\alpha_0$ alone is, at a configuration $\mathbf{U}$
of the space $\mathfrak{U}$, the statement that block averages of a near-unitary field of small
curvature have small curvature: the sentence of Remark~\ref{4.5:rem-block-curvature-complex},
which is stated in \cite[(3.26)]{B-CMP109} without proof.
\item If \cite[Proposition~6]{B-CMP98} alone is used in place of that sentence, so that only
$|V-1|\le CM(\alpha_0+\alpha_1)$ is available, the assumption \cite[p.~277]{B-CMP109} that
$B_3^2\,O(1)M\alpha_0<\tfrac12\alpha_1$ becomes an upper bound on $M/L$.
\item The same frame bounds, involving $\alpha_0$ only and linear in $M$, are obtained
without that statement, through an axial gauge on the fine lattice.
\end{enumerate}
\end{remark}

\begin{proof}
In the frame of \cite[(3.26)--(3.28)]{B-CMP109} the window background is reached through two
gauge transformations. The first is the gauge transformation $v$ that brings $M^{\cdot}(\mathbf{U})$
into the generalised axial gauge,
\begin{equation*}
M^{\cdot}(\mathbf{U})=V^{v},\qquad |V(b)-1|<O(1)M\alpha_0,\qquad |v-1|<O(1)M\alpha_1 .
\end{equation*}
The second is the map $u_{k+1}$ from axial to Landau gauge for $U_{k+1}(\square_0,V)$.
By the response corollary at the base point $1$, the bound $B_3\,O(1)M\alpha_0$ of
\cite[(3.27)]{B-CMP109} on $\mathbf{H}_{k+1}(\square_0,(1/i)\log V)$, its gradient and its
H\"older norm on $\tilde\square^{4}$ is the response of the background function
$\mathbf{H}_{k+1}$ to the datum $(1/i)\log V$.
The size of that datum is the one given by \cite[(3.26)]{B-CMP109}. Hence the form of
\cite[(3.27)]{B-CMP109} involving $\alpha_0$ only rests on the form of \cite[(3.26)]{B-CMP109}
involving $\alpha_0$ only.

At a configuration $\mathbf{U}$ of the space $\mathfrak{U}$, the form of \cite[(3.26)]{B-CMP109}
involving $\alpha_0$ only is exactly the
statement of Remark~\ref{4.5:rem-block-curvature-complex}, which is stated in
\cite[(3.26)]{B-CMP109} without proof. That remark says that block averages of a
near-unitary field of small curvature have small curvature. This proves the first item.

For the second item, suppose only \cite[Proposition~6]{B-CMP98} is used. It gives
\begin{equation*}
|V-1|\le C M(\alpha_0+\alpha_1)
\end{equation*}
with a constant $C$. Carry this bound through \cite[(3.27)]{B-CMP109} and
\cite[(3.32)]{B-CMP109}, and then into the assumption of \cite[p.~277]{B-CMP109}:
\begin{equation*}
B_3^2\,O(1)M\alpha_0<\tfrac12\alpha_1 .
\end{equation*}
The term proportional to $\alpha_1$ in the corresponding condition then reads
\begin{equation*}
B_3B_bM\alpha_1L^{-1}<\tfrac12\alpha_1 ,
\end{equation*}
where $B_b$ is the constant in the data-axial frame lemma's bound
$|\mathcal{A}|\le B_bM(1+2\beta)\alpha_0$ on the frame potential $\mathcal{A}$.
Dividing by $\alpha_1>0$ gives $B_3B_b\,M/L<\tfrac12$, which is an upper bound on $M/L$; the
computation of this term is the one in the remark on the upper bound on $M/L$ that accompanies
the data-axial frame lemma.
This proves the second item. It also shows that, in the literal frame, the statement of
Remark~\ref{4.5:rem-block-curvature-complex}, which is stated in \cite[(3.26)]{B-CMP109}
without proof, is used in an essential way.

That statement is nevertheless a choice of route and not a necessity. The fine axial frame
lemma reaches the same frame bounds, involving $\alpha_0$ only and linear in $M$, by a
different route. It starts from the curvature condition (iii) of \cite[p.~275]{B-CMP109} in the
definition of the space $\mathfrak{U}$. It fixes an axial gauge of the field $\mathbf{U}$ itself
on the fine lattice,
where the curvature bound is a hypothesis and not a proposition. It then averages at the
base point $1$, where \cite[Proposition~6]{B-CMP98} is printed for complex twists.
The near bound of the operator bounds at free current is obtained through the fine axial
frame lemma. This proves the third item.

The data-axial frame lemma reproduces Ba{\l}aban's own construction in the literal frame of
\cite[(3.26)--(3.28)]{B-CMP109}. The sentence of Remark~\ref{4.5:rem-block-curvature-complex},
which is stated in \cite[(3.26)]{B-CMP109} without proof, enters only at that construction, and
the near bound does not depend on it.
\end{proof}

\begin{lemma}[Algebra of exponentially decaying kernels]\label{4.5:lem-kernel-algebra}
Use the lattice geometry, block sizes and weighted norms of
Conventions~\ref{4.5:not-geometry-conventions}. Let
$a\in\mathbb{R}$, $C>0$, $\delta>0$. A linear operator $T$ belongs to $\mathcal{K}_a$ \emph{with constant
$C$ and rate $\delta$} if, for all $y\in\Lambda_j$, $y'\in\Lambda_{j'}$ and every $f$ with
$\operatorname{supp}f\subset\Delta(y')$,
\begin{equation*}
\sup_{\Delta(y)}|Tf|\le C\,\ell_j^{a}\,e^{-\delta d(y,y')}\,\sup|f| .
\end{equation*}
Let $\delta_0$ be the decay constant of \cite{B-CMP96}. For $\varpi>0$ put
$c_0(\varpi):=\sum_{z\in\mathbb{Z}}e^{-\varpi\delta_0|z|}$, as in \cite[\S2]{B-CMP96}. Assume
\begin{equation}\label{4.5:eq-kernel-algebra-a}
\begin{gathered}
\tfrac{\delta}{4}\,R_1\mathfrak{M}\ge (d+4)\log L,\\
\tfrac14\,\varpi\,\delta_0\,R_1\mathfrak{M}>2d\log c_0(\tfrac12\varpi)+1
\quad\text{for }\varpi\in\{\tfrac14,\tfrac18,\tfrac1{16},c_{*}\}.
\end{gathered}
\end{equation}
The second line is \cite[(2.59)]{B-CMP96}, read with $R_1\mathfrak{M}$ in place of Ba{\l}aban's
product of gap number and block size. It is read at the fraction $\tfrac14$ of $\delta_0$ used in the
present lemma and at the further fractions $\tfrac18$, $\tfrac1{16}$ and $c_{*}$ of $\delta_0$, where
$c_{*}>0$ denotes the decay-rate fraction. Every condition in \eqref{4.5:eq-kernel-algebra-a} is a
lower bound on the one number $R_1\mathfrak{M}$. Let $c_1$ be the constant of
\cite[Lemma~2.1]{B-CMP96} at the least fraction of $\delta_0$ listed in
\eqref{4.5:eq-kernel-algebra-a}, and assume moreover that $\delta$ is such that the sum bound
\cite[(2.61)]{B-CMP96} holds in the form
\begin{equation}\label{4.5:eq-kernel-algebra-1}
\sum_{y'\in\mathfrak{B}}e^{-(\delta/4)d(y,y')}\le c_1\qquad(y\in\mathfrak{B}).
\end{equation}
Then:
\begin{enumerate}
\item[(i)] if $T\in\mathcal{K}_a$ with constant $C$ and rate $\delta$, $|\alpha|\le d+4$,
$S\subset\mathfrak{B}$, and $\rho(y)=e^{\kappa d(y,S)}$ is the weight of the conventions with
$0\le\kappa\le\delta/2$, then
\begin{equation*}
|Tf|_{(\alpha+a),\rho}\le C L^{d+4}c_1\,|f|_{(\alpha),\rho};
\end{equation*}
\item[(ii)] $\mathcal{K}_a\mathcal{K}_b\subset\mathcal{K}_{a+b}$, the product having rate $\delta/2$;
\item[(iii)] if $T\in\mathcal{K}_0$ with constant $C$ and rate $\delta$, and $CL^{d+4}c_1\le\tfrac12$,
then the Neumann series $\sum_{n\ge0}T^n$ belongs to $\mathcal{K}_0$ with constant $2$ and rate
$\delta/2$;
\item[(iv)] if $T$ has a kernel satisfying
\begin{equation*}
|T(x,x')|\le C\,\ell_j^{a}\,\ell_{j'}^{-d}\,e^{-\delta d(y,y')}
\qquad(x\in\Delta(y),\ x'\in\Delta(y'),\ y\in\Lambda_j,\ y'\in\Lambda_{j'}),
\end{equation*}
then $T\in\mathcal{K}_a$ with rate $\delta$, and if moreover $|d+a|\le d+4$, then
$T^{T}\in\mathcal{K}_a$ with rate $\tfrac34\delta$.
\end{enumerate}
\end{lemma}

\begin{proof}
(i) Fix $y\in\Lambda_j$ and write $f=\sum_{y'}f_{y'}$, where $f_{y'}$ is $f$ on $\Delta(y')$ and zero
elsewhere. For $y'\in\Lambda_{j'}$ the definition of the weighted norm gives
\begin{equation*}
\sup|f_{y'}|\le\rho(y')^{-1}\ell_{j'}^{\alpha}\,|f|_{(\alpha),\rho}.
\end{equation*}
By linearity of $T$ and the definition of $\mathcal{K}_a$ applied to each $f_{y'}$,
\begin{equation}\label{4.5:eq-kernel-algebra-2}
\rho(y)\,\ell_j^{-\alpha-a}\sup_{\Delta(y)}|Tf|
\le C\sum_{y'}e^{-\delta d(y,y')}\,\frac{\rho(y)}{\rho(y')}\,
\ell_{j'}^{\alpha}\ell_j^{-\alpha}\;|f|_{(\alpha),\rho}.
\end{equation}
Three estimates control the right-hand side.

First, the weight. Since $d(y,S)\le d(y,y')+d(y',S)$, we have
$\rho(y)\le\rho(y')e^{\kappa d(y,y')}$. Because $\kappa\le\delta/2$, this gives
$\rho(y)/\rho(y')\le e^{(\delta/2)d(y,y')}$.

Second, the level factors. Write $(t)_+:=\max\{t,0\}$. By \cite[(2.60)]{B-CMP96},
\begin{equation*}
d(y,y')\ge R_1\mathfrak{M}\,(|j-j'|-1)_+ .
\end{equation*}
Moreover $|\alpha|\log L\le(d+4)\log L$, which is at most $\tfrac{\delta}{4}R_1\mathfrak{M}$ by the
first line of
\eqref{4.5:eq-kernel-algebra-a}. Since $\ell_{j'}/\ell_j=L^{j'-j}$, these two facts give
\begin{equation*}
\begin{aligned}
\ell_{j'}^{\alpha}\ell_j^{-\alpha}
&\le L^{|\alpha|\,|j-j'|}\le L^{|\alpha|}\,L^{|\alpha|(|j-j'|-1)_+}\\
&\le L^{|\alpha|}\,e^{(\delta/4)R_1\mathfrak{M}(|j-j'|-1)_+}
\le L^{|\alpha|}\,e^{(\delta/4)d(y,y')} .
\end{aligned}
\end{equation*}
This is the exchange of level factors against decay made in \cite[p.~397]{B-CMP99b}.

Third, the sum. Inserting the first two estimates into \eqref{4.5:eq-kernel-algebra-2}, the
exponents combine to $-\delta+\frac{\delta}{2}+\frac{\delta}{4}=-\frac{\delta}{4}$. By
\eqref{4.5:eq-kernel-algebra-1},
\begin{equation*}
\rho(y)\,\ell_j^{-\alpha-a}\sup_{\Delta(y)}|Tf|
\le C L^{|\alpha|}\sum_{y'}e^{-(\delta/4)d(y,y')}\,|f|_{(\alpha),\rho}
\le C L^{d+4}c_1\,|f|_{(\alpha),\rho},
\end{equation*}
where $L^{|\alpha|}\le L^{d+4}$ because $L>1$ and $|\alpha|\le d+4$. Taking the supremum over $j$ and
$y\in\Lambda_j$ gives (i).

(ii) The statement on products is obtained from the composition bound \cite[(2.63)]{B-CMP96},
applied to the two factors.

(iii) Let $y'\in\mathfrak{B}$ and $\operatorname{supp}f\subset\Delta(y')$. We apply (i) with
$a=\alpha=0$, $S=\{y'\}$ and $\kappa=\delta/2$, so that $\rho(y)=e^{(\delta/2)d(y,y')}$. In the
decomposition used in the proof of (i), $f$ has the single term $f_{y'}=f$, and $\rho(y')=1$;
hence $|f|_{(0),\rho}=\sup|f|$. Applying (i) $n$ times and using $CL^{d+4}c_1\le\tfrac12$,
\begin{equation*}
|T^nf|_{(0),\rho}\le(CL^{d+4}c_1)^n\,|f|_{(0),\rho}\le 2^{-n}\sup|f| .
\end{equation*}
By the definition of the weighted norm, for $y\in\Lambda_j$ this gives
$\sup_{\Delta(y)}|T^nf|\le 2^{-n}e^{-(\delta/2)d(y,y')}\sup|f|$. Summing over $n\ge0$, the series
converges uniformly on $\Delta(y)$, and
\begin{equation*}
\sup_{\Delta(y)}\Bigl|\sum_{n\ge0}T^nf\Bigr|\le 2\,e^{-(\delta/2)d(y,y')}\sup|f| ,
\end{equation*}
which is membership in $\mathcal{K}_0$ with constant $2$ and rate $\delta/2$.

(iv) For $T$ itself, let $\operatorname{supp}f\subset\Delta(y')$ and $x\in\Delta(y)$. Summing the
kernel bound over $x'\in\Delta(y')$ with the weight $\ell_{j'}^{d}$ of the pairing, the factor
$\ell_{j'}^{-d}$ is cancelled, and
\begin{equation*}
\sup_{\Delta(y)}|Tf|\le C'\,\ell_j^{a}\,e^{-\delta d(y,y')}\,\sup|f|,
\end{equation*}
where $C'$ is $C$ times the largest number of points of a block $\Delta(y')$. Hence
$T\in\mathcal{K}_a$ with constant $C'$ and rate $\delta$.

For the transpose, $T^{T}(x,x')=T(x',x)$. For $x\in\Delta(y)$, $x'\in\Delta(y')$ with
$y\in\Lambda_j$, $y'\in\Lambda_{j'}$, the kernel bound with the roles of the two points exchanged
gives
\begin{equation*}
|T^{T}(x,x')|\le C\,\ell_{j'}^{a}\ell_j^{-d}e^{-\delta d(y,y')}
=C\,\ell_j^{a}\ell_{j'}^{-d}\Bigl(\frac{\ell_{j'}}{\ell_j}\Bigr)^{d+a}e^{-\delta d(y,y')} .
\end{equation*}
The factor $(\ell_{j'}/\ell_j)^{d+a}$, that is the exchange of $\ell_{j'}^{-d}\ell_j^{d}$, is
traded against decay exactly as the level factors in (i). Since $|d+a|\le d+4$ by hypothesis,
$d+a$ takes the place of $\alpha$, and the first line of \eqref{4.5:eq-kernel-algebra-a} gives
\begin{equation*}
\Bigl(\frac{\ell_{j'}}{\ell_j}\Bigr)^{d+a}\le L^{d+4}e^{(\delta/4)d(y,y')} .
\end{equation*}
Hence $T^{T}$ has a kernel of the same form as $T$, with constant $CL^{d+4}$ and rate
$\tfrac34\delta$. By the first part of (iv), applied with this constant and this rate,
$T^{T}\in\mathcal{K}_a$ with rate $\tfrac34\delta$.
\end{proof}

\begin{definition}[Propagators and constrained Green's functions at an admissible background
{\cite[Theorems~3.3, 3.4, (3.42)--(3.49), (3.69), Corollary~3.5]{B-CMP99b}}]
\label{4.5:not-background-propagators}
Throughout, $(U,J)$ is a pair of the class $\mathfrak{S}$ of Definition~\ref{4.5:def-admissible-pairs}.
Thus $U = e^{i\eta\mathfrak{a}}U_r$ satisfies the conditions recalled in
\eqref{4.5:eq-admissible-pairs-a}, with the constant $a_0$ of the first condition and the constant
$a_1$ of the second condition there.
The classes $\mathcal{K}_n$ of kernels of order $n$ with exponential decay are those of the kernel
algebra (Lemma~\ref{4.5:lem-kernel-algebra}). The bilinear transpose is written ${}^{T}$; the operators
$D^{*}_{U}$, $\nabla^{*}_{U}$, $Q^{*}$ are the analytic continuations of the adjoints, equal to the
transposes. The operators below are those of \cite{B-CMP99b} and \cite{B-CMP102b}, read at $U$.

\begin{enumerate}
\item \emph{Extension.} In \cite[Theorem~3.4]{B-CMP99b} the operators $G'$,
$\bigl(Q'G'^{2}Q'^{*}\bigr)^{-1}$, $R$ and $G$, defined at the real configurations of
\cite[Theorems~3.1--3.3]{B-CMP99b}, extend to the complex configurations written $U'U$ there, as analytic
functions of the variable written $A'$ there, under a smallness condition on $A'$ with a positive
constant $a_1$; the extended operators satisfy the inequalities of
\cite[Theorems~3.1--3.3]{B-CMP99b}. Here the real configuration is $U_r$, which satisfies
\cite[(3.35)]{B-CMP99b} by the first condition in \eqref{4.5:eq-admissible-pairs-a}; the complex
configuration is $U = e^{i\eta\mathfrak{a}}U_r$, the analytic variable is
$\mathfrak{a}$, and the smallness condition is the second condition in
\eqref{4.5:eq-admissible-pairs-a}, with the same constant $a_1$. On the pairs of $\mathfrak{S}$ these
operators are therefore defined at $U$ and analytic in $\mathfrak{a}$, and every bound below holds there.

\item \emph{The operators $G'$, $R$, $P$, $P_1$.} Among the four bounds of \cite[(3.42)]{B-CMP99b},
those on $G'$, $\nabla_UG'$ and $G'\nabla^{*}_{U}$ give
$G' \in \mathcal{K}_2$ and $\nabla_U G',\ G'\nabla^{*}_{U} \in \mathcal{K}_1$. With $P := I - R$ and
$P_1 := D_U P D^{*}_{U}$, the pointwise kernel bounds of \cite[(3.49)]{B-CMP99b} on $P$, $D_UP$,
$PD^{*}_{U}$ and $D_UPD^{*}_{U}$, at $x \in \Delta(y)$, $y \in \Lambda_j$, and $x' \in \Delta(y')$,
$y' \in \Lambda_{j'}$, are $O(1)$ times the factors $1$, $(L^{j}\eta)^{-1}$ and
$(L^{j'}\eta)^{-1}$ respectively, and, for $D_UPD^{*}_{U}$, the fourth factor printed there, each
together with the remaining factors of the right member printed there.
Read through the pointwise-kernel clause of the kernel algebra
(Lemma~\ref{4.5:lem-kernel-algebra}), these bounds give
\begin{equation*}
R,\ P \in \mathcal{K}_0, \qquad D_UP,\ PD^{*}_{U} \in \mathcal{K}_{-1}, \qquad P_1 \in \mathcal{K}_{-2}.
\end{equation*}

\item \emph{The Green's function $G_0$.} Let $\ell$ denote the multiplication operator by the scale
of the hierarchy, equal to $\ell_j$ on the layer $\Lambda_j$, as in the kernel algebra
(Lemma~\ref{4.5:lem-kernel-algebra}). Put $a := \ell^{-2}$, define
\begin{equation*}
\Delta_a := \Delta + D_U R D^{*}_{U} + Q^{*}aQ, \qquad G_0 := \Delta_a^{-1},
\end{equation*}
and write $\mathsf{q}_0 := (QG_0Q^{*})^{-1}$.

By \cite[Theorem~3.3]{B-CMP99b}, the operator $G(U)$ (with $a = 1$ there) satisfies
\cite[(3.42)--(3.47)]{B-CMP99b} with $G'(U)$ replaced by $G(U)$. Read with $a = \ell^{-2}$ in
place of $a = 1$, these inequalities give
$G_0 \in \mathcal{K}_2$ and $\nabla_UG_0,\ G_0\nabla^{*}_{U} \in \mathcal{K}_1$.

The H\"older bounds \cite[(3.43)]{B-CMP99b} hold on sup-norm data for every H\"older exponent
$\beta \le \beta_0$, with $\beta_0 < 1$ as in \cite[(3.43)]{B-CMP99b} (in this item $\beta$ is a
H\"older exponent, not the schedule constant of Definition~\ref{4.4:def-post-step-space}).
The operators $\nabla_UG_0\nabla^{*}_{U}$ and $\nabla_UG'\nabla^{*}_{U}$ are of
order $0$ on H\"older data of exponent $\beta + \varepsilon$, where $0 < \varepsilon \le 1$ and
$\beta \le \beta_0$. Explicitly, for the localised data of \cite[(3.45)]{B-CMP99b}, with a cut-off
$\varphi$ supported in $\Delta(y)$, $y \in \Lambda_j$, and a datum $\psi$ supported in $\Delta(y')$,
in the convention of Lemma~\ref{4.5:lem-kernel-algebra},
\begin{equation*}
\begin{split}
\|\varphi\,\nabla_UG'(U)\nabla^{*}_{U}\psi\|_{\beta}
 &\le B_0'(\varepsilon,\beta_0)\,(L^{j}\eta)^{-\beta}\,\bigl(\|\varphi\|_{\beta} + |\varphi|\bigr)\\
 &\qquad\times e^{-\delta_0 d(y,y')}\,\bigl(\|\psi\|_{\beta+\varepsilon} + |\psi|\bigr);
\end{split}
\end{equation*}
here $\|\cdot\|_{\gamma}$ denotes the H\"older norm of exponent $\gamma$ on the cube carrying the
support, $|\cdot|$ the supremum norm, and $\delta_0$ the decay rate, all as in
\cite[(3.43)--(3.45)]{B-CMP99b}.
The sup-norm bound on $\nabla_UG\nabla^{*}_{U}\psi$ for $\varepsilon$-H\"older data $\psi$ is
\cite[(3.44)]{B-CMP99b}.

\item \emph{Local and block-local operators.} The operator $\Delta' := \Delta - D^{*}_{U}D_U$ is local,
and by \cite[(3.69)]{B-CMP99b}
\begin{equation*}
|\Delta' A|_{(-3)} \le O(1)\,(a_0 + a_1)\,|A|_{(-1)}.
\end{equation*}
The operators $Q$, $Q^{*}$ and $Q'$ are block-local.

\item \emph{The constrained covariance inverse.} By Remark~\ref{4.5:rem-covariance-inverse-walks},
printed without proof in \cite[\S3]{B-CMP99b} and cited as printed, one has
$\mathsf{q}_0 \in \mathcal{K}_{-2}$,
and $\mathsf{q}_0$ has a generalised random-walk expansion.

\item \emph{The operator $\Lambda$.} Put $\Lambda := G'RD^{*}_{U}$ (an operator; not to be confused
with the layers $\Lambda_j$ of the hierarchy). Since $R$ and $G'$ are symmetric for
the bilinear form, $\Lambda^{T} = D_URG'$. Moreover
\begin{equation*}
|\Lambda A|_{(0)} \le C\,|A|_{(-1)}, \qquad |D_U\Lambda A|_{(-1)} \le C\,N(A).
\end{equation*}
We write $C_R := \|\Lambda^{T}\|_{(-4)\to(-3)}$ for the norm of $\Lambda^{T}$ from
$|\cdot|_{(-4)}$ to $|\cdot|_{(-3)}$.

\item \emph{Minimiser map and constrained Green's function.} Following \cite[(129), (131)]{B-CMP102b}, put
\begin{equation*}
H_0 := G_0Q^{*}\mathsf{q}_0\,\ell^{-1}, \qquad P_0 := I - G_0Q^{*}\mathsf{q}_0Q,
\qquad \tilde{\mathfrak{G}}_0 := P_0G_0 .
\end{equation*}
By \cite[p.~300]{B-CMP102b}, $\tilde{\mathfrak{G}}_0 = G_0P_0^{T}$. There $\tilde{\mathfrak{G}}_0$ is
written as $G_0$ minus $G_0Q^{*}(QG_0Q^{*})^{-1}QG_0$.
\end{enumerate}

\noindent The definitions give the identities
\begin{equation}\label{4.5:eq-background-propagators-1}
\begin{gathered}
QH_0 = \ell^{-1}, \qquad QP_0 = 0, \qquad P_0^{T}Q^{*} = 0,\\
Q\tilde{\mathfrak{G}}_0 = 0, \qquad \Delta_a\tilde{\mathfrak{G}}_0 = P_0^{T}, \qquad
\Delta_aH_0B = Q^{*}\mathsf{q}_0\tilde B, \quad \tilde B := \ell^{-1}B .
\end{gathered}
\end{equation}
Each identity in \eqref{4.5:eq-background-propagators-1} follows in one line.
\begin{itemize}
\item Since $(QG_0Q^{*})\mathsf{q}_0 = I$, the product $QH_0 = (QG_0Q^{*})\mathsf{q}_0\ell^{-1}$
equals $\ell^{-1}$.
\item For the same reason, $QP_0 = Q - (QG_0Q^{*})\mathsf{q}_0Q$ equals $Q - Q = 0$.
\item Transposing $QP_0 = 0$, with $Q^{T} = Q^{*}$, gives $P_0^{T}Q^{*} = 0$.
\item Since $Q\tilde{\mathfrak{G}}_0 = (QP_0)G_0$ and $QP_0 = 0$, one has $Q\tilde{\mathfrak{G}}_0 = 0$.
\item Since $\Delta_aG_0 = I$ and $\tilde{\mathfrak{G}}_0 = G_0P_0^{T}$, one has
$\Delta_a\tilde{\mathfrak{G}}_0 = P_0^{T}$.
\item Since $\Delta_aG_0 = I$, the product $\Delta_aH_0B = \Delta_aG_0Q^{*}\mathsf{q}_0\ell^{-1}B$
equals $Q^{*}\mathsf{q}_0\ell^{-1}B$.
\end{itemize}

\noindent Finally, with constants $B_0$ and $B_0'$ (the letter $B_0'$, written here without
arguments, denotes the constant of the second bound below),
\begin{equation}\label{4.5:eq-background-propagators-2}
N_{\beta}(H_0B) \le B_0\,|B|, \qquad N_{\beta}(\tilde{\mathfrak{G}}_0f) \le B_0'\,|f|_{(-3)},
\qquad P_0^{T} \in \mathcal{K}_0 .
\end{equation}
The bounds collected in this definition are quoted by statement from \cite[\S3]{B-CMP99b}, combined,
where $\mathsf{q}_0$ enters, with Remark~\ref{4.5:rem-covariance-inverse-walks} (printed without proof
in \cite[\S3]{B-CMP99b}; cited as printed), at the operators of \cite[(129), (131)]{B-CMP102b}; their
proofs are Ba{\l}aban's and are not reproduced.
\end{definition}

\begin{lemma}[The operators with the current as a free variable]
\label{4.5:lem-free-current-operators}
Assume the bounds and the generalised random-walk expansion of
$\mathsf{q}_0=(QG_0Q^{*})^{-1}$ given in Remark~\ref{4.5:rem-covariance-inverse-walks},
which is stated in \cite{B-CMP99b} without proof, and the analytic and uniform form of the
derivative bounds for the averaging maps given in Remark~\ref{4.5:rem-averaging-analyticity},
which is stated in \cite{B-CMP98} without proof.

Let $c^{(1)}$, $c^{(2)}$ be the second-order terms of the expansion \cite[(3.116)]{B-CMP99b}
of the gauge transform of a bond variable by a gauge function $\lambda$,
\begin{equation*}
\frac{1}{i\eta}\log\bigl[e^{i\lambda(b_-)}e^{i\eta A(b)}e^{-iR_b\lambda(b_+)}\bigr]
= A - D_U\lambda + c^{(1)}(A,\lambda) + c^{(2)}(\lambda,\lambda) + \cdots,
\end{equation*}
where $R_b$ is the operator of \cite[(3.116)]{B-CMP99b} acting on the value $\lambda(b_+)$ at the
final point of $b$ (it is distinct from the operator $R$ entering $\Lambda$ below), and where,
by the Baker--Campbell--Hausdorff formula,
\begin{equation*}
c^{(1)}(A,\lambda)(b) = \frac{i}{2}\bigl[\lambda(b_-)+R_b\lambda(b_+),A(b)\bigr],
\qquad
c^{(2)}(\lambda,\lambda)(b) = \frac{1}{2i}\bigl[\lambda(b_-),(D_U\lambda)(b)\bigr];
\end{equation*}
no negative power of $\eta$ occurs in $c^{(1)}$ or $c^{(2)}$. Let $c^{(2)}(\lambda_1,\lambda_2)$ be
the polarisation of the quadratic expression $c^{(2)}(\lambda,\lambda)$, and, for a current $J$,
define $\mathfrak{j}_J$, $\mathfrak{k}_J$, $\mathfrak{m}_J$ by
\begin{equation*}
\langle\lambda,\mathfrak{j}_J(A)\rangle = \langle A,\mathfrak{k}_J(\lambda)\rangle
:= \langle c^{(1)}(A,\lambda),J\rangle,
\qquad
\langle\lambda_1,\mathfrak{m}_J(\lambda_2)\rangle := 2\langle c^{(2)}(\lambda_1,\lambda_2),J\rangle .
\end{equation*}
Then $\mathfrak{j}_J$ and $\mathfrak{k}_J$ are pointwise in $J$, with
\begin{equation*}
|\mathfrak{j}_J(A)|_{(-4)}\le 2d\,|A|_{(-1)}|J|_{(-3)},
\qquad
|\mathfrak{k}_J(\lambda)|_{(-3)}\le 2|\lambda|_{(0)}|J|_{(-3)},
\end{equation*}
and $\mathfrak{m}_J(\lambda)=m_3+D^{*}_Um_4$, where $m_3$ is a pointwise expression in the
commutators $[D_U\lambda,J]$ and $m_4$ a pointwise expression in the commutators $[\lambda,J]$.

With $\Lambda:=G'RD^{*}_U$, $\Pi:=I-D_U\Lambda$ and $C^{(2)}$ the second-order term entering
\cite[(3.134)--(3.136)]{B-CMP99b} (an object distinct from the term $c^{(2)}$ above), consider
the operators
\begin{equation}\label{4.5:eq-free-current-operators-a}
\begin{gathered}
\Theta'_J:=\mathfrak{k}_J\Lambda+\Lambda^{T}\mathfrak{j}_J+\Lambda^{T}\mathfrak{m}_J\Lambda,
\qquad
G^{J}_{\pi}:=\sum_{n\ge0}G_0(\Theta'_JG_0)^{n},
\qquad
\mathsf{q}_J:=(QG^{J}_{\pi}Q^{*})^{-1},
\\
H^{J}:=G^{J}_{\pi}Q^{*}\mathsf{q}_J\ell^{-1},
\qquad
\langle A,\Theta_JA\rangle:=2\langle H^{J}C^{(2)}(\ell A),J\rangle,
\qquad
\tilde G^{J}:=(I-\tilde{\mathfrak{G}}_0\Theta_J)^{-1}\tilde{\mathfrak{G}}_0,
\\
G^{J}_{1}:=\sum_{n\ge0}G_0\bigl((\Theta'_J+\Pi^{T}\Theta_J\Pi)G_0\bigr)^{n},
\qquad
H^{J}_{1}:=G^{J}_{1}Q^{*}(QG^{J}_{1}Q^{*})^{-1}\ell^{-1}.
\end{gathered}
\end{equation}
There are constants $a_\gamma$, $C_J$, $C_\Theta$, depending only on $d$ and $L$, such that, with
$\varepsilon:=\frac12(1-\beta_0)$ fixed, the following properties (J1)--(J8) of the operators
\eqref{4.5:eq-free-current-operators-a} hold on the class $\mathfrak{S}$ of admissible
pairs $(U,J)$ of \eqref{4.5:eq-admissible-pairs-a}. The constants satisfy
$C_Ja_\gamma\le\frac12$ and $C_\Theta a_\gamma B_0'\le\frac12$. Here $\beta_0<1$ is the ceiling
of the H\"older exponents, and $B_0$, $B_0'$ are the constants of the bounds on $H_0$ and on
$\tilde{\mathfrak{G}}_0$, all as fixed in Definition~\ref{4.5:not-background-propagators}.
\begin{enumerate}
\item[(J1)] The operator $\Theta'_J$ obeys
\begin{equation*}
N(G_0\Theta'_JA)\le C_J|J|_{(-3)}N(A),
\end{equation*}
also with $N_\beta$, for every $\beta\le\beta_0$, on the left, and in the weighted norms $N^{\rho}$.
At $J=J(U)$, the current of the expansion of the action about the background (the weight
appearing in \cite[(3.120)]{B-CMP99b}), $\Theta'_J$ is the operator $\Delta'_\pi$ of
\cite[(3.120)]{B-CMP99b}, by the Ward identities under nonlinear gauge variations
\eqref{4.5:eq-ward-identities-a}.
\item[(J2)] The series defining $G^{J}_{\pi}$ converges in the following norms of
\cite[(3.42)--(3.47)]{B-CMP99b}: the members with sup data on the right (the bounds on
$|Gf|$, $|\nabla Gf|$ and the localised $\beta$-H\"older member of $\nabla Gf$) and the members
with $\nabla^{*}$ on the right (the bounds on $G\nabla^{*}\lambda$ and $\nabla G\nabla^{*}\lambda$
on data measured in $\|\lambda\|_\varepsilon$). The Laplacian member is not asserted for
$G^{J}_{\pi}$; it is used for $G_0$ and $H_0$ only. Moreover
$(G^{J}_{\pi})^{-1}=\Delta_a-\Theta'_J$.
\item[(J3)] $\mathsf{q}_J\in\mathcal{K}_{-2}$; this is the analogue of \cite[(3.132)]{B-CMP99b}.
\item[(J4)] $H^{J}\in\mathcal{K}_{-1}$ and $N_\beta(H^{J}B)\le B_0|B|$ for fields $B$ on the
coarse lattice, the analogue of \cite[(3.133)]{B-CMP99b}; moreover
\begin{equation*}
\ell QH^{J}=1,
\qquad
(\Delta_a-\Theta'_J)H^{J}=Q^{*}\mathsf{q}_J\ell^{-1}.
\end{equation*}
The identity $RD^{*}_UH^{J}=0$ is neither asserted nor used below.
\item[(J5)] The field $H^{J,T}J$, which is the weight of the quadratic form defining $\Theta_J$,
is of weight $-4$: $|H^{J,T}J|_{(-4)}$ is at most a constant depending only on $d$ and $L$
times $|J|_{(-3)}$. The operator $\Theta_J$ has finite range, and
\begin{equation*}
|\Theta_JA|_{(-3)}\le C_\Theta|J|_{(-3)}|A|_{(-1)},
\end{equation*}
the analogue of \cite[(3.137)]{B-CMP99b}.
\item[(J6)] $N_\beta(\tilde G^{J}f)\le 2B_0'|f|_{(-3)}$ and $Q\tilde G^{J}=0$.
\item[(J7)] The assertions of (J2)--(J4) hold for $G^{J}_{1}$ and $H^{J}_{1}$ in place of
$G^{J}_{\pi}$ and $H^{J}$ (compare \cite[(3.135), (3.138)]{B-CMP99b}).
\item[(J8)] All these operators are analytic in $(\mathfrak{a},J)$. Each of
\begin{equation*}
G_0,\ G',\ R,\ P,\ (Q'G'^{2}Q'^{*})^{-1},\ \mathsf{q}_0,\ G^{J}_{\pi},\ \mathsf{q}_J,\
\tilde G^{J},\ G^{J}_{1},\ (QG^{J}_{1}Q^{*})^{-1}
\end{equation*}
has a generalised random-walk expansion with bounds of the type of
\cite[(3.108)]{B-CMP99b} and \cite[(1.7)]{B-CMP116} in the norms listed in (J2), each term
depending on $(U,J)$ only on its domains. For $\mathsf{q}_0$ this is the content of
Remark~\ref{4.5:rem-covariance-inverse-walks}, which is stated in \cite{B-CMP99b} without proof.
For $G^{J}_{\pi}$, $\mathsf{q}_J$, $\tilde G^{J}$,
$G^{J}_{1}$ the series defining $G^{J}_{\pi}$ and $G^{J}_{1}$ and the Neumann series for
$\mathsf{q}_J$ and $\tilde G^{J}$ are expanded into concatenated walks, not re-summed: pieces of
walks of $G_0$ joined by the local insertions $\Theta'_J$, $\Pi^{T}\Theta_J\Pi$ and
$Q(G^{J}_{\pi}-G_0)Q^{*}$, each insertion carrying its side branch (the walk of $H^{J,T}J$ inside
$\Theta_J$; the walks of $\Lambda$ and of $RG'$ inside $\Theta'_J$), and the side branches are
counted among the domains of the term.
\end{enumerate}
\end{lemma}

\begin{proof}
Throughout, $C$ denotes a constant depending only on $d$ and $L$, and the properties of
$G_0$, $G'$, $R$, $P$, $\Lambda$, $H_0$, $\tilde{\mathfrak{G}}_0$ used are those fixed
in Definition~\ref{4.5:not-background-propagators}, the statement on propagators and constrained
Green's functions at an admissible background; in particular $\Lambda^{T}=D_URG'$,
$|\Lambda A|_{(0)}\le C|A|_{(-1)}$ and $|D_U\Lambda A|_{(-1)}\le CN(A)$.

\emph{The decomposition of $\mathfrak{m}_J$.} By polarisation,
\begin{equation*}
c^{(2)}(\lambda_1,\lambda_2)(b)=\frac{1}{4i}\Bigl(\bigl[\lambda_1(b_-),(D_U\lambda_2)(b)\bigr]
+\bigl[\lambda_2(b_-),(D_U\lambda_1)(b)\bigr]\Bigr).
\end{equation*}
In the pairing with $J$, the invariance of the pairing under the adjoint action moves each
commutator onto $J(b)$. The first term then pairs $\lambda_1(b_-)$ with a commutator of
$(D_U\lambda_2)(b)$ and $J(b)$; summing over the bonds with a given initial point gives
$\langle\lambda_1,m_3\rangle$ with $m_3$ pointwise in $[D_U\lambda_2,J]$. The second term pairs
$(D_U\lambda_1)(b)$ with a commutator of $\lambda_2(b_-)$ and $J(b)$, that is, it equals
$\langle\lambda_1,D^{*}_Um_4\rangle$ with $m_4$ pointwise in $[\lambda_2,J]$.

\emph{(J1).} Fix $A$ and put $\theta:=\Theta'_JA$. Write $m_3$, $m_4$ for the two pieces of
$\mathfrak{m}_J(\Lambda A)$, so that $\mathfrak{m}_J(\Lambda A)=m_3+D^{*}_Um_4$. Since
$\Lambda^{T}=D_URG'$, the definition of $\Theta'_J$ gives
\begin{equation}\label{4.5:eq-free-current-operators-1}
\theta=\psi+D_U\varphi,
\qquad
\psi:=\mathfrak{k}_J(\Lambda A),
\qquad
\varphi:=RG'\bigl(\mathfrak{j}_J(A)+m_3\bigr)+RG'D^{*}_Um_4 .
\end{equation}
By the bounds on $\Lambda A$ and $D_U\Lambda A$ just recalled, the pointwise bounds on
$\mathfrak{j}_J$ and $\mathfrak{k}_J$, and the pointwise form of $m_3$ (in
$[D_U\Lambda A,J]$) and of $m_4$ (in $[\Lambda A,J]$),
\begin{equation*}
|\psi|_{(-3)},\quad |\mathfrak{j}_J(A)+m_3|_{(-4)},\quad |m_4|_{(-3)}
\ \le\ C|J|_{(-3)}N(A).
\end{equation*}
Hence $|\varphi|_{(-2)}\le C|J|_{(-3)}N(A)$, and every H\"older norm of $\varphi$ of exponent less
than one, at weight $-2$, obeys a bound of the same form. Indeed, write $R=I-P$ in
\eqref{4.5:eq-free-current-operators-1}. The piece $G'(\mathfrak{j}_J(A)+m_3)$ is Lipschitz by the
gradient member of \cite[(3.42)]{B-CMP99b}; the piece $G'D^{*}_Um_4$ has H\"older norms of every
exponent less than one by the second member of \cite[(3.43)]{B-CMP99b}, applied to the sup
datum $m_4$; and the parts coming from $P$ are Lipschitz by \cite[(3.49)]{B-CMP99b}.

Now, by \eqref{4.5:eq-free-current-operators-1},
\begin{equation*}
G_0\theta=G_0\psi+(G_0\nabla^{*})\varphi,
\qquad
\nabla G_0\theta=\nabla G_0\psi+(\nabla G_0\nabla^{*})\varphi .
\end{equation*}
Both terms of $G_0\theta$ are bounded in $|\cdot|_{(-1)}$ by \cite[(3.47)]{B-CMP99b}, which holds for
$G_0$ in place of $G'$. In $\nabla G_0\theta$, the first term is bounded by
\cite[(3.47)]{B-CMP99b}, and the second by \cite[(3.44)]{B-CMP99b}, applied with the
$\varepsilon$-H\"older norm of $\varphi$. For the $\beta$-H\"older member of $\nabla G_0\theta$, the
first term is bounded by \cite[(3.43)]{B-CMP99b}, and the second by \cite[(3.45)]{B-CMP99b} with the
$(\beta+\varepsilon)$-H\"older norm of $\varphi$, which is available because, for $\beta\le\beta_0$,
\begin{equation*}
\beta+\varepsilon\le\beta_0+\tfrac12(1-\beta_0)=\tfrac12(1+\beta_0)<1 .
\end{equation*}
Collecting, $N(G_0\Theta'_JA)\le C_J|J|_{(-3)}N(A)$, and the same bound with $N_\beta$ on the left
for every $\beta\le\beta_0$. This is the mechanism of the corresponding estimate for $\Delta'_\pi$
in \cite[\S3]{B-CMP99b}: one of the three derivatives is applied to the expression on the right
or on the left of the perturbation, and the constants depend only on the constants of
\cite[(3.42)--(3.47)]{B-CMP99b} at a suitably chosen $\varepsilon$. Here
$\varepsilon=\frac12(1-\beta_0)$, so that both \cite[(3.44)]{B-CMP99b} and
\cite[(3.45)]{B-CMP99b} apply. Only $|J|_{(-3)}$ enters: no H\"older norm of $J$ and no
$D^{*}_UJ$ is used, and $\varepsilon$ is never sent to zero.

For the weighted norms $N^{\rho}$, apply the algebra of exponentially decaying kernels
(Lemma~\ref{4.5:lem-kernel-algebra}) to each factor of the chain above in turn; the insertions
$\mathfrak{j}_J$, $\mathfrak{k}_J$, $m_3$, $m_4$ are pointwise, and each costs at most a factor
$e^{\kappa c}$, where $c$ bounds the range of these pointwise insertions.

\emph{(J2).} By (J1), $A\mapsto G_0\Theta'_JA$ has norm at most $C_J|J|_{(-3)}$ in $N$, in each
$N_\beta$ with $\beta\le\beta_0$, and in $N^{\rho}$. Since
$G_0(\Theta'_JG_0)^{n}=(G_0\Theta'_J)^{n}G_0$, the series defining $G^{J}_{\pi}$ is the geometric
series $\sum_n(G_0\Theta'_J)^{n}$ applied to $G_0f$, for a datum $f$ of the member in question,
and it converges because
$C_Ja_\gamma\le\frac12$. For the members with sup data on the right, $G_0f$ is bounded by
\cite[(3.42)--(3.47)]{B-CMP99b} for $G_0$, and the geometric series in $G_0\Theta'_J$ preserves
these bounds. This is the argument of \cite[\S3]{B-CMP99b} for the series built on $\Delta'_\pi$,
with the parameter $\alpha_0$ there replaced by $|J|_{(-3)}$: the bound \cite[(3.36)]{B-CMP99b}
enters that argument only as the size of the weight. For the members with $\nabla^{*}$ on the
right, the first iterate $A=G_0\nabla^{*}\lambda$ has $|\nabla A|$ bounded only on data measured
in $\|\lambda\|_\varepsilon$, by \cite[(3.44)]{B-CMP99b}; thereafter the splitting
\eqref{4.5:eq-free-current-operators-1} of (J1) applies to every further iterate. Finally, telescoping,
\begin{equation*}
G^{J}_{\pi}(\Delta_a-\Theta'_J)=\sum_{n\ge0}(G_0\Theta'_J)^{n}(I-G_0\Theta'_J)=I,
\end{equation*}
since $(G_0\Theta'_J)^{n}\to0$ by (J1); likewise $(\Delta_a-\Theta'_J)G^{J}_{\pi}$ telescopes to
$I-\lim_n\Theta'_J(G_0\Theta'_J)^{n}G_0=I$, because $(G_0\Theta'_J)^{n}G_0f\to0$ in $N$ by (J1),
and, by \eqref{4.5:eq-free-current-operators-1}, $\Theta'_JA=\psi+D_U\varphi$ with $|\psi|_{(-3)}$
and $|\varphi|_{(-2)}$ at most $C|J|_{(-3)}N(A)$ for every $A$, so that, for
$A=(G_0\Theta'_J)^{n}G_0f$, both pieces tend to zero in these norms.
Hence $(G^{J}_{\pi})^{-1}=\Delta_a-\Theta'_J$.

\emph{(J3).} Since $QG_0Q^{*}\mathsf{q}_0=I$,
\begin{equation*}
QG^{J}_{\pi}Q^{*}=QG_0Q^{*}\bigl(I+\mathsf{q}_0Q(G^{J}_{\pi}-G_0)Q^{*}\bigr).
\end{equation*}
By Remark~\ref{4.5:rem-covariance-inverse-walks}, which is stated in \cite{B-CMP99b} without
proof, $\mathsf{q}_0\in\mathcal{K}_{-2}$; by (J2), $G^{J}_{\pi}-G_0=\sum_{n\ge1}(G_0\Theta'_J)^{n}G_0$
is of size $O(a_\gamma)$. By Lemma~\ref{4.5:lem-kernel-algebra},
$\mathsf{q}_0Q(G^{J}_{\pi}-G_0)Q^{*}$ is a perturbation in $\mathcal{K}_0$ of size $O(a_\gamma)$, so
its Neumann series converges and
\begin{equation*}
\mathsf{q}_J=\bigl(I+\mathsf{q}_0Q(G^{J}_{\pi}-G_0)Q^{*}\bigr)^{-1}\mathsf{q}_0\in\mathcal{K}_{-2}.
\end{equation*}

\emph{(J4).} $H^{J}\in\mathcal{K}_{-1}$ and the analogue of \cite[(3.133)]{B-CMP99b} follow by
composition, using Lemma~\ref{4.5:lem-kernel-algebra}, (J2) and (J3), in the decomposition
\begin{equation*}
H^{J}=H_0+(G^{J}_{\pi}-G_0)Q^{*}\mathsf{q}_J\ell^{-1}+G_0Q^{*}(\mathsf{q}_J-\mathsf{q}_0)\ell^{-1};
\end{equation*}
the piece $H_0$ obeys $N_\beta(H_0B)\le B_0|B|$ by
Remark~\ref{4.5:rem-covariance-inverse-walks}, which is stated in \cite{B-CMP99b} without proof.
The identities follow from the formula: $QH^{J}=QG^{J}_{\pi}Q^{*}\mathsf{q}_J\ell^{-1}=\ell^{-1}$, and
$(\Delta_a-\Theta'_J)H^{J}=Q^{*}\mathsf{q}_J\ell^{-1}$ because $(\Delta_a-\Theta'_J)G^{J}_{\pi}=I$ by (J2).

\emph{(J5).} By definition
\begin{equation*}
\langle H^{J}C^{(2)}(\ell A),J\rangle=\langle C^{(2)}(\ell A),H^{J,T}J\rangle .
\end{equation*}
Since $H^{J}\in\mathcal{K}_{-1}$ by (J4), Lemma~\ref{4.5:lem-kernel-algebra} gives
$|H^{J,T}J|_{(-4)}\le C|J|_{(-3)}$. With this weight in place of the weight of
\cite[(3.134)--(3.136)]{B-CMP99b}, the derivation of \cite[(3.137)]{B-CMP99b} uses
\cite[(149)]{B-CMP98} at complex configurations $U$, which is available by
Remark~\ref{4.5:rem-averaging-analyticity}, stated in \cite{B-CMP98} without proof; it yields the
finite range of $\Theta_J$ and $|\Theta_JA|_{(-3)}\le C_\Theta|J|_{(-3)}|A|_{(-1)}$.

\emph{(J6).} Expand
\begin{equation*}
\tilde G^{J}=\sum_{n\ge0}(\tilde{\mathfrak{G}}_0\Theta_J)^{n}\tilde{\mathfrak{G}}_0 .
\end{equation*}
By (J5) and
$N_\beta(\tilde{\mathfrak{G}}_0f)\le B_0'|f|_{(-3)}$ for every datum $f$, each factor
$\tilde{\mathfrak{G}}_0\Theta_J$ has
norm at most $C_\Theta|J|_{(-3)}B_0'$, and $C_\Theta a_\gamma B_0'\le\frac12$; the Neumann series
therefore converges and $N_\beta(\tilde G^{J}f)\le 2B_0'|f|_{(-3)}$. Every term begins with
$\tilde{\mathfrak{G}}_0$ on the left and $Q\tilde{\mathfrak{G}}_0=0$, so $Q\tilde G^{J}=0$.

\emph{(J7).} Since $\Pi^{T}=I-\Lambda^{T}D^{*}_U=I-D_URG'D^{*}_U$,
\begin{equation*}
\Pi^{T}\Theta_J\Pi A=f'-D_U\varphi',
\qquad
f':=\Theta_J\Pi A,
\qquad
\varphi':=RG'D^{*}_Uf'.
\end{equation*}
Here $|\Pi A|_{(-1)}\le CN(A)$ by the bound on $D_U\Lambda A$, so $f'$ is of weight $-3$ with
$|f'|_{(-3)}\le C|J|_{(-3)}N(A)$ by (J5), and $\varphi'$ is of the form of the last piece of
$\varphi$ in \eqref{4.5:eq-free-current-operators-1}. This is the structure of (J1), and the
arguments of (J1)--(J4) apply verbatim to $G^{J}_{1}$, $(QG^{J}_{1}Q^{*})^{-1}$ and $H^{J}_{1}$.

\emph{(J8).} Consider one of the cubes on which, after a gauge transformation $u$, the
configuration takes the form $U^{u}=e^{i\eta A^{c}}$ with a complex potential $A^{c}$,
$N(A^{c})\le2(a_0+a_1)$. By \cite[Corollary~3.5]{B-CMP99b} such a configuration lies in the
domain \cite[(3.37)]{B-CMP99b} with background $U=1$; this corollary serves the local
propagators at complex configurations as \cite[Corollary~3.6]{B-CMP99b} serves them at real
$U$, and the constructions of \cite[Theorems~3.7, 3.9 and 3.10]{B-CMP99b} are built from these
local propagators alone. This gives the random-walk expansions of $G_0$, $G'$, $R$, $P$ and
$(Q'G'^{2}Q'^{*})^{-1}$ with bounds of the type of \cite[(3.108)]{B-CMP99b}, analytically in
$\mathfrak{a}$ by the statement on propagators at an admissible background
in Definition~\ref{4.5:not-background-propagators}. For the series in $J$, each operator in the
series is
replaced by its random-walk expansion, as is done for \cite[(3.130)]{B-CMP99b}. The current
$J$ enters polynomially, through the pointwise insertions $\mathfrak{j}_J$, $\mathfrak{k}_J$,
$m_3$, $m_4$ and through the walk-expanded insertion $H^{J,T}J$; hence each term is analytic
in $(\mathfrak{a},J)$ and depends on $(U,J)$ only on its domains, the side branches included.
The insertion series converge uniformly, since $C_Ja_\gamma\le\frac12$ and
$C_\Theta a_\gamma B_0'\le\frac12$, so their sums are analytic in $(\mathfrak{a},J)$ and obey the
bounds in the norms of (J2). The expansion of $\mathsf{q}_0$ is that of
Remark~\ref{4.5:rem-covariance-inverse-walks}, which is stated in \cite{B-CMP99b} without proof;
the expansions of $\mathsf{q}_J$ and of $(QG^{J}_{1}Q^{*})^{-1}$ are obtained by concatenating the
walks of the Neumann series of (J3) and of its analogue in (J7).
\end{proof}

\begin{lemma}[The nonlinear correction to the linearising map at free current]
\label{4.5:lem-nonlinear-correction}
Assume the analytic and uniform form of \cite[Proposition~5]{B-CMP98} formulated in
Remark~\ref{4.5:rem-averaging-analyticity}, which is stated in \cite{B-CMP98} without proof.
Let $H^{J}$ be the operator of Lemma~\ref{4.5:lem-free-current-operators}, and let $B_0$ be the
constant in its bound. The scale factors $\ell$, $\ell_j$ and the layers $\Lambda_j$ are those of
Lemma~\ref{4.5:lem-kernel-algebra}, and $Q_j(U;\cdot)$ is the averaging operator onto $\Lambda_j$
of \cite[(48)]{B-CMP102b}.
Let $C$ be the nonlinear map of \cite{B-CMP102b} whose fixed-point problem
is treated in \cite[(50)--(58)]{B-CMP102b}, the map whose functional derivative is bounded in
\cite[Proposition~5, (157)]{B-CMP98}, and let $C'$, $C''$ be its first and second functional
derivatives. Let $a_3$ and $C_3$ be the constants of the argument of
\cite[(50)--(58), (63)--(73)]{B-CMP102b}, and let $C_2$ be the constant supplied by
\cite[Propositions~4 and~7]{B-CMP98}. Let $\varepsilon_3 \le a_3$ and
$|A'|_{(-1)} < \varepsilon_3$. Then the following hold.
\begin{enumerate}
\item[(a)] The map $X \mapsto C(\ell A' - \ell H^{J} X)$ has a unique fixed point $D^{J}(A')$
in the set of $X$ with $X = 0$ on $\Lambda_0$ and $|X| < \varepsilon_3/B_0$. It is analytic
in $(A', \mathfrak{a}, J)$, and $|D^{J}(A')| \le 4 C_2 |A'|_{(-1)}^{2}$.
\item[(b)] The map $\mathfrak{L}^{J}(A') := A' - H^{J} D^{J}(A')$ satisfies
$N(\mathfrak{L}^{J} A') \le 2 N(A')$ and the analogue of \cite[(48)]{B-CMP102b}:
\begin{equation}\label{4.5:eq-nonlinear-correction-1}
Q_j\bigl(U; \eta\, \mathfrak{L}^{J}(A')\bigr) = \ell_j\, Q_j A' \qquad \text{on } \Lambda_j .
\end{equation}
\item[(c)] Put $\mathfrak{R}^{J} := C'(\ell \mathfrak{L}^{J} A')\, \ell H^{J}$, the counterpart,
with the linearising map of \cite{B-CMP102b} replaced by $H^{J}$, of the operator that enters
\cite[(63)--(73)]{B-CMP102b}. The derivative
$\mathfrak{D}^{J} := \delta D^{J}/\delta A'$ is
\begin{equation}\label{4.5:eq-nonlinear-correction-2}
\mathfrak{D}^{J} = (I + \mathfrak{R}^{J})^{-1}\, C'(\ell \mathfrak{L}^{J} A')\, \ell ,
\end{equation}
and $\mathfrak{D}^{J} \in \mathcal{K}_1$, $\mathfrak{D}^{J,T} \in \mathcal{K}_1$, with constant
$O(C_3 \varepsilon_3)$. The difference $\mathfrak{D}^{J}_{3} := \mathfrak{D}^{J} - \mathfrak{D}^{(2)}$,
where $\mathfrak{D}^{(2)}$ denotes the leading part of $\mathfrak{D}^{J}$, linear in $A'$, namely
the functional derivative of the quadratic approximation to $D^{J}$ in the argument of
\cite[(63)--(73)]{B-CMP102b} with the linearising map replaced by $H^{J}$, is
$O(\varepsilon_3^{2})$. Put $\mathfrak{T}^{J} := I - H^{J} \mathfrak{D}^{J}$.
\item[(d)] For a direction $\mathfrak{A}$ of variation of $A'$,
\begin{equation}\label{4.5:eq-nonlinear-correction-3}
\partial_{\mathfrak{A}} \mathfrak{D}^{J}
 = (I + \mathfrak{R}^{J})^{-1}\, C''(\ell \mathfrak{L}^{J} A')
 \bigl[\ell \mathfrak{T}^{J} \cdot\,,\ \ell \mathfrak{T}^{J} \mathfrak{A}\bigr],
\end{equation}
and, with a constant $C_{D2}$, for every $Y$ and every index $\alpha$ for which the norms
$|\cdot|_{(\alpha)}$ are defined,
\begin{equation}\label{4.5:eq-nonlinear-correction-4}
\bigl|(\partial_{\mathfrak{A}} \mathfrak{D}^{J})^{T} Y\bigr|_{(\alpha+1),\rho}
 \le C_{D2}\, N^{\rho}(\mathfrak{A})\, |Y|_{(\alpha)} .
\end{equation}
\end{enumerate}
\end{lemma}

\begin{proof}
The lemma is \cite[Proposition~3]{B-CMP102b} with the linearising map $h$ of that proposition
replaced by $H^{J}$; we carry over the argument of \cite[(50)--(58), (63)--(73)]{B-CMP102b}
without change. That argument uses two properties
of $h$ and no others: the identity $\ell Q h = 1$, which is what gives the linearisation
property, here \eqref{4.5:eq-nonlinear-correction-1}; and the bound $|hX| \le B_0 \ell^{-1}|X|$
together with the exponential decay of the kernel of $h$. Ba{\l}aban remarks in
\cite{B-CMP102b} that there are many linearising transformations, that the choice $h = H$ is a
convenient one but by no means the only one, and that $D(A')$ has the same domain of analyticity
as $H(U_0)$. For $H^{J}$, the clause of Lemma~\ref{4.5:lem-free-current-operators} on $H^{J}$,
displayed in \eqref{4.5:eq-free-current-operators-a}, gives $\ell Q H^{J} = 1$,
$H^{J} \in \mathcal{K}_{-1}$ and $N_{\beta}(H^{J} X) \le B_0 |X|$, which is the bound
$|hX| \le B_0 \ell^{-1}|X|$ in the norms $N_{\beta}$; the analyticity clause of the same lemma
gives analyticity of
$H^{J}$ in $(\mathfrak{a}, J)$. Hence the fixed-point argument of \cite{B-CMP102b} yields the
existence and uniqueness in (a), the analyticity of $D^{J}(A')$ in $(A', \mathfrak{a}, J)$, and
(b). The constant $C_2$ in the bound of (a) is supplied by \cite[Propositions~4 and~7]{B-CMP98}.
The step corresponding to \cite[(72)]{B-CMP102b} uses \cite[Proposition~5]{B-CMP98} in its
analytic and uniform form, which is the hypothesis of the lemma (Remark
\ref{4.5:rem-averaging-analyticity}, stated in \cite{B-CMP98} without proof). The step
corresponding to \cite[(71)]{B-CMP102b} is the algebra of exponentially decaying kernels
(Lemma~\ref{4.5:lem-kernel-algebra}). The bounds of (c) are then \cite[(73)]{B-CMP102b} in the
present setting.

For the formula \eqref{4.5:eq-nonlinear-correction-2}, differentiate the fixed-point identity
$D^{J}(A') = C(\ell A' - \ell H^{J} D^{J}(A'))$ with respect to $A'$. Since
$\ell A' - \ell H^{J} D^{J}(A') = \ell \mathfrak{L}^{J} A'$, the chain rule gives
\begin{equation*}
\mathfrak{D}^{J} = C'(\ell \mathfrak{L}^{J} A')\,\bigl(\ell - \ell H^{J} \mathfrak{D}^{J}\bigr),
\end{equation*}
that is, $(I + \mathfrak{R}^{J}) \mathfrak{D}^{J} = C'(\ell \mathfrak{L}^{J} A')\, \ell$, and
the invertibility of $I + \mathfrak{R}^{J}$ is part of the argument just carried over. The same
identity reads $\mathfrak{D}^{J} = C'(\ell \mathfrak{L}^{J} A')\, \ell\, \mathfrak{T}^{J}$, and
$\delta \mathfrak{L}^{J}/\delta A' = I - H^{J} \mathfrak{D}^{J} = \mathfrak{T}^{J}$.

For (d), differentiate this identity once more, in the direction $\mathfrak{A}$. The derivative
of the argument $\ell \mathfrak{L}^{J} A'$ is $\ell \mathfrak{T}^{J} \mathfrak{A}$, and the
derivative of $\mathfrak{T}^{J}$ is $-H^{J} \partial_{\mathfrak{A}} \mathfrak{D}^{J}$; hence
\begin{equation*}
\partial_{\mathfrak{A}} \mathfrak{D}^{J}
 = C''(\ell \mathfrak{L}^{J} A')\bigl[\ell \mathfrak{T}^{J} \cdot\,,\ \ell \mathfrak{T}^{J}
 \mathfrak{A}\bigr] - \mathfrak{R}^{J}\, \partial_{\mathfrak{A}} \mathfrak{D}^{J},
\end{equation*}
and applying $(I + \mathfrak{R}^{J})^{-1}$ gives \eqref{4.5:eq-nonlinear-correction-3}. For the
bound \eqref{4.5:eq-nonlinear-correction-4}, the second derivative $C''$ is block-local. Of the
two arguments of the bilinear map $C''(\ell \mathfrak{L}^{J} A')[\cdot,\cdot]$, the one that is a
functional derivative is bounded by \cite[(157)]{B-CMP98}, in the analytic and
uniform form of Remark~\ref{4.5:rem-averaging-analyticity}, which is stated in \cite{B-CMP98}
without proof; the other is bounded by a Cauchy estimate.
\end{proof}

\begin{lemma}[Locality and size of the curvature remainder in the current]\label{4.5:lem-curvature-remainder}
Let $U$ be a configuration satisfying conditions (S1) and (S2) of Definition~\ref{4.5:def-admissible-pairs}, with the constants $a_0$ and $a_1$ appearing there. Put
\begin{equation*}
\mathcal{J}(U) := D^{*}_{U}\,\eta^{-2}\,\pi\operatorname{Im}\partial U ,
\end{equation*}
with $\pi$ and $\operatorname{Im}$ as in \cite[(3.11)]{B-CMP109}, and define
\begin{equation}\label{4.5:eq-curvature-remainder-1}
F(U,A) := \mathcal{J}(e^{i\eta A}U) - \mathcal{J}(U) - D^{*}_{U}D_{U}A .
\end{equation}
Then $F(U,A)$ is local and analytic, and $F(U,0) = 0$. Moreover, there are constants $C_F$ and $C_F'$ such that, for $N(A) \le 1$ and every direction $f$,
\begin{align}
|F(U,A)|_{(-3)} &\le C_F\,\bigl(a_0 + a_1 + N(A)\bigr)\,N(A), \label{4.5:eq-curvature-remainder-2}\\
|\partial_A F(U,A)[f]|_{(-3)} &\le C_F'\,\bigl(a_0 + a_1 + N(A)\bigr)\,N(f). \label{4.5:eq-curvature-remainder-3}
\end{align}
\end{lemma}

\begin{proof}
The definition \eqref{4.5:eq-curvature-remainder-1} is \cite[(3.11)]{B-CMP109}. On \cite[p.~272]{B-CMP109} the remainder is introduced through the identity
\begin{equation*}
\begin{split}
D^{\xi *}_{\exp i\xi\mathbf{A}\,U}\,\xi^{-2}\pi\operatorname{Im}\partial \exp i\xi\mathbf{A}\,U
&= D^{\xi *}_{U}\,\xi^{-2}\pi\operatorname{Im}\partial U \\
&\quad + D^{\xi *}_{U}D^{\xi}_{U}\mathbf{A} + \mathbf{F}(U,\mathbf{A}),
\end{split}
\end{equation*}
in which $\xi$ is the lattice spacing, written $\eta$ here, and the boldface $\mathbf{A}$ and $\mathbf{F}(U,\mathbf{A})$ are the field $A$ and the remainder $F(U,A)$ of \eqref{4.5:eq-curvature-remainder-1}. Written in the present letters, the identity expresses the left-hand side as the sum of $\mathcal{J}(U)$, of $D^{*}_{U}D_{U}A$ and of $F(U,A)$, which is \eqref{4.5:eq-curvature-remainder-1}.

\emph{Locality.} It is stated in \cite[p.~272]{B-CMP109} that, for any regular configuration $U$, $\mathbf{F}$ is a local operator depending on $U$, $\partial U$, $\mathbf{A}$ and $\nabla^{\xi}_{U}\mathbf{A}$ only. In the present letters, $F(U,A)$ is a local function of $U$, $\partial U$, $A$ and $\nabla_U A$.

\emph{Vanishing at $A = 0$.} At $A = 0$ we have $e^{i\eta A}U = U$ and $D^{*}_{U}D_{U}A = 0$. Hence \eqref{4.5:eq-curvature-remainder-1} gives
\begin{equation*}
F(U,0) = \mathcal{J}(U) - \mathcal{J}(U) = 0 .
\end{equation*}

\emph{Structure of the summands.} Expand $F(U,A)$ in powers of $A$. By \cite[(3.6)]{B-CMP99b}, every summand of this expansion is a product of factors, each of which is $A$, $\nabla_U A$ or $\eta^{-2}(\partial U - 1)$. At most one factor of each product is differenced by $D^{*}$. Two cases of a differenced factor occur.
\begin{itemize}
\item If the differenced factor is a curvature factor $\eta^{-2}(\partial U - 1)$, the difference carries its explicit factor $\eta^{2}$.
\item If the differenced factor is a commutator $[A,A]$, the difference gives a term of the form $(\nabla_U A)\cdot A$, by \cite[(94)--(96)]{B-CMP102b}.
\end{itemize}
Each factor is measured in the norms of the hierarchy. By the definition of $N(A)$,
\begin{equation*}
|A|_{(-1)} \le N(A), \qquad |\nabla_U A|_{(-2)} \le N(A).
\end{equation*}
By conditions (S1) and (S2) of Definition~\ref{4.5:def-admissible-pairs}, each curvature factor $\eta^{-2}(\partial U - 1)$ is bounded in the norms of the hierarchy by a constant times $a_0 + a_1$. Since $F(U,0) = 0$, every summand has order at least one in $A$, and so contains at least one factor $A$ or $\nabla_U A$. Moreover, $\mathcal{J}(U)$ is the part of $\mathcal{J}(e^{i\eta A}U)$ of order zero in $A$, and $D^{*}_{U}D_{U}A$ is its part of order one in $A$ that contains no curvature factor. Since both are subtracted in \eqref{4.5:eq-curvature-remainder-1}, every summand of $F(U,A)$ contains, besides one factor $A$ or $\nabla_U A$, at least one further factor $A$, $\nabla_U A$ or $\eta^{-2}(\partial U - 1)$.

\emph{Convergence and the first bound.} The series in $A$ converges for $N(A) \le 1$. In particular, $F(U,A)$ is analytic in $A$ there. Bound each summand by the product of the bounds on its factors: the factor $A$ or $\nabla_U A$ contributes $N(A)$, and the further factor contributes at most a constant times $a_0 + a_1 + N(A)$. Summing these bounds over the convergent series gives the quadratic bound \eqref{4.5:eq-curvature-remainder-2}.

\emph{The derivative bound.} Since $F(U,A)$ is given for $N(A)\le 1$ by the expansion above, we differentiate it term by term in the direction $f$. The derivative of a summand is the sum of the products obtained from it by replacing one factor $A$ by $f$, or one factor $\nabla_U A$ by $\nabla_U f$, leaving all other factors unchanged. By the definition of $N(f)$, the new factor is bounded by $N(f)$ in the norm in which the replaced factor was measured. By the structure of the summands established above, each summand of $F(U,A)$ contains at least two factors, of which at least one is $A$ or $\nabla_U A$; hence, after one factor $A$ or $\nabla_U A$ has been replaced, at least one factor $A$, $\nabla_U A$ or $\eta^{-2}(\partial U - 1)$ remains, and it contributes at most a constant times $a_0 + a_1 + N(A)$. The remaining factors are bounded as in the proof of \eqref{4.5:eq-curvature-remainder-2}, using $N(A)\le 1$. Summing these bounds over the series gives \eqref{4.5:eq-curvature-remainder-3}.
\end{proof}

In the estimates that follow, three further objects and one domain are used; we fix them here.

First, let $V_0(A;U)$ denote the sum of the terms of order at least three in $A$ of the action $\mathsf{A}(e^{i\eta A}U)$. It is an entire function of $A$, and its coefficients are local in $\eta^{-2}(\partial U - 1)$.

Second, put
\begin{equation*}
\mathsf{k} := (\mathsf{q}_J - a)\,\ell^{-1} ,
\end{equation*}
where $\ell$ denotes multiplication by $\ell_j = L^{j}\eta$ on the $j$-th layer of the hierarchy, and $a$ is a fixed constant.

Third, put
\begin{equation*}
D^{J}_{3} := D^{J} - C^{(2)}(\ell\,\cdot) ,
\end{equation*}
where $C^{(2)}$ is the second-order map that, composed with multiplication by $\ell$, is removed from $D^{J}$.

Fourth, the functional $V^{J}$ introduced next is defined for fields $A'$ of the same kind as the field $A$ above, in the set
\begin{equation}\label{4.5:eq-curvature-remainder-4}
\{\, A' : N(A') < \varepsilon_3 \,\},
\end{equation}
where $\varepsilon_3$ is a fixed positive radius. No Landau gauge condition is imposed on $A'$ in \eqref{4.5:eq-curvature-remainder-4}.

\begin{definition}[The higher-order action functional at free current]
\label{4.5:def-higher-order-functional}
Let $\varepsilon_3$ be as in Lemma~\ref{4.5:lem-nonlinear-correction}. Take
$D^{J}$, $\mathfrak{L}^{J}$, $\mathfrak{D}^{J}$, $\mathfrak{D}^{J}_{3}$ and
$\mathfrak{T}^{J} = I - H^{J}\mathfrak{D}^{J}$ from that lemma, and the pointwise insertions
$\mathfrak{j}_J$, $\mathfrak{k}_J$ from Lemma~\ref{4.5:lem-free-current-operators}. Let
$V_0(\cdot\,;U)$, the operator $\mathsf{k} = (\mathsf{q}_J - a)\ell^{-1}$ and
$D^{J}_{3} = D^{J} - C^{(2)}(\ell\,\cdot)$ be as fixed following
Lemma~\ref{4.5:lem-curvature-remainder}; recall that $V_0$ is
entire. Let $\Lambda := G'RD^{*}_{U}$.

On the set $\{A' : N(A') < \varepsilon_3\}$ of \cite[(77)]{B-CMP102b}, on which no Landau
condition is imposed, define $V^{J,\mathrm{disp}}$, $W^{J}$, $V^{J}$ and $\mathbb{V}'^{J}$ by the
definitions (the relations $:=$) in
\begin{equation}\label{4.5:eq-higher-order-functional-a}
\begin{split}
V^{J,\mathrm{disp}}(A') &:= -\langle H^{J}D^{J}_{3},J\rangle
 - \langle QA',\mathsf{k}D^{J}\rangle_{\mathfrak{B}}\\
 &\qquad + \tfrac12\langle \ell^{-1}D^{J},\mathsf{k}D^{J}\rangle_{\mathfrak{B}}
 + V_0(\mathfrak{L}^{J}A';U),\\
W^{J}(A') &:= -\langle \Lambda A',\mathfrak{j}_J(H^{J}D^{J}(A'))\rangle,\\
V^{J} &:= V^{J,\mathrm{disp}} + W^{J},\\
\mathbb{V}'^{J} &:= \nabla V^{J}\\
 &= -\mathfrak{D}^{J,T}_{3}H^{J,T}J - Q^{T}\mathsf{k}D^{J}
 - \mathfrak{D}^{J,T}\mathsf{k}^{T}QA' + \mathfrak{D}^{J,T}\ell^{-1}\mathsf{k}D^{J}\\
 &\qquad + \mathfrak{T}^{J,T}\nabla V_0(\mathfrak{L}^{J}A')
 - \Lambda^{T}\mathfrak{j}_J(H^{J}D^{J})
 - \mathfrak{D}^{J,T}H^{J,T}\mathfrak{k}_J(\Lambda A').
\end{split}
\end{equation}
Here $\nabla$ is the gradient in $A'$ for the bilinear pairing $\langle\cdot,\cdot\rangle$, and
${}^{T}$ denotes the bilinear transpose.
\end{definition}

The last equality in \eqref{4.5:eq-higher-order-functional-a}, the seven-term formula for
$\mathbb{V}'^{J}$, is not part of the definition; it is proved below. In it the operator
$\ell^{-1}\mathsf{k}$ is symmetric for the bilinear pairing, because $\Theta'_J$ is.

The functional $V^{J,\mathrm{disp}}$ in \eqref{4.5:eq-higher-order-functional-a} is the
counterpart of
\cite[(80)]{B-CMP102b}. Its second and third terms are written in the form that follows
\cite[(87)]{B-CMP102b}, and the whole expression is taken here as a definition.

\begin{proof}[Proof of the seven-term formula in \eqref{4.5:eq-higher-order-functional-a}]
The map $A'\mapsto D^{J}(A')$ is analytic by Lemma~\ref{4.5:lem-nonlinear-correction}, and
$V_0$ is entire. Hence $V^{J}$ is differentiable. The operators $H^{J}$, $Q$, $\mathsf{k}$,
$\ell$ and $\Lambda$ do not depend on $A'$. We differentiate each term in a direction
$\mathfrak{A}$ and write each derivative as $\langle \mathfrak{A},\cdot\rangle$.

By Lemma~\ref{4.5:lem-nonlinear-correction}, the derivative of $D^{J}$ is $\mathfrak{D}^{J}$.
The derivative of $D^{J}_{3} = D^{J} - C^{(2)}(\ell\,\cdot)$ is therefore $\mathfrak{D}^{J}$
minus the derivative of $A'\mapsto C^{(2)}(\ell A')$; the latter derivative is the operator
$\mathfrak{D}^{(2)}$ with which Lemma~\ref{4.5:lem-nonlinear-correction} forms
$\mathfrak{D}^{J}_{3} = \mathfrak{D}^{J} - \mathfrak{D}^{(2)}$, so the derivative of $D^{J}_{3}$
is $\mathfrak{D}^{J}_{3}$. Since $\mathfrak{L}^{J} = A' - H^{J}D^{J}$, the
derivative of $\mathfrak{L}^{J}$ is $\mathfrak{T}^{J}$.

The first term has derivative
\begin{equation*}
-\langle H^{J}\mathfrak{D}^{J}_{3}\mathfrak{A},J\rangle
= -\langle \mathfrak{A},\mathfrak{D}^{J,T}_{3}H^{J,T}J\rangle .
\end{equation*}

The second term is bilinear in $A'$ and $D^{J}$. Its derivative is
\begin{equation*}
-\langle Q\mathfrak{A},\mathsf{k}D^{J}\rangle_{\mathfrak{B}}
-\langle QA',\mathsf{k}\mathfrak{D}^{J}\mathfrak{A}\rangle_{\mathfrak{B}}
= -\langle \mathfrak{A},Q^{T}\mathsf{k}D^{J}\rangle
 -\langle \mathfrak{A},\mathfrak{D}^{J,T}\mathsf{k}^{T}QA'\rangle .
\end{equation*}

For the third term, $\ell$ acts on each layer as multiplication by a number, so $\ell^{-1}$ is
symmetric. The operator $\ell^{-1}\mathsf{k}$ is symmetric as well. Therefore the two terms
produced by the product rule are equal:
\begin{equation*}
\langle \ell^{-1}\mathfrak{D}^{J}\mathfrak{A},\mathsf{k}D^{J}\rangle_{\mathfrak{B}}
= \langle \mathfrak{D}^{J}\mathfrak{A},\ell^{-1}\mathsf{k}D^{J}\rangle_{\mathfrak{B}}
= \langle \ell^{-1}D^{J},\mathsf{k}\mathfrak{D}^{J}\mathfrak{A}\rangle_{\mathfrak{B}} .
\end{equation*}
With the factor $\tfrac12$, the derivative of the third term is
$\langle \mathfrak{A},\mathfrak{D}^{J,T}\ell^{-1}\mathsf{k}D^{J}\rangle$.

By the chain rule, the fourth term has derivative
\begin{equation*}
\langle \nabla V_0(\mathfrak{L}^{J}A'),\mathfrak{T}^{J}\mathfrak{A}\rangle
= \langle \mathfrak{A},\mathfrak{T}^{J,T}\nabla V_0(\mathfrak{L}^{J}A')\rangle .
\end{equation*}

For $W^{J}$ we use Lemma~\ref{4.5:lem-free-current-operators}. By that lemma,
\begin{equation*}
\langle\lambda,\mathfrak{j}_J(A)\rangle = \langle A,\mathfrak{k}_J(\lambda)\rangle
= \langle c^{(1)}(A,\lambda),J\rangle ,
\end{equation*}
and this expression is bilinear in $(A,\lambda)$. The
product rule therefore gives
\begin{equation*}
-\langle \Lambda\mathfrak{A},\mathfrak{j}_J(H^{J}D^{J})\rangle
-\langle \Lambda A',\mathfrak{j}_J(H^{J}\mathfrak{D}^{J}\mathfrak{A})\rangle .
\end{equation*}
The first of these two terms equals
$-\langle \mathfrak{A},\Lambda^{T}\mathfrak{j}_J(H^{J}D^{J})\rangle$. The identity above, with
$\lambda = \Lambda A'$ and $A = H^{J}\mathfrak{D}^{J}\mathfrak{A}$, turns the second term into
\begin{equation*}
-\langle H^{J}\mathfrak{D}^{J}\mathfrak{A},\mathfrak{k}_J(\Lambda A')\rangle
= -\langle \mathfrak{A},\mathfrak{D}^{J,T}H^{J,T}\mathfrak{k}_J(\Lambda A')\rangle .
\end{equation*}

Adding the seven contributions gives the seven-term formula in
\eqref{4.5:eq-higher-order-functional-a}.
\end{proof}

\begin{lemma}[Quadratic bound on the gradient of the higher-order functional]
\label{4.5:lem-gradient-bound}
Let $\mathbb{V}'^{J}=\nabla V^{J}$ be the gradient of the higher-order action functional at free
current, given by the seven-term formula \eqref{4.5:eq-higher-order-functional-a}. Let
$\mathbb{V}''^{J}$ be its derivative in $A'$. Then the following hold.
\begin{enumerate}
\item[(i)] $\mathbb{V}'^{J}$ is analytic in $(A',\mathfrak{a},J)$ on the product of the set
$\{N(A')<\varepsilon_3\}$ with the class $\mathfrak{S}$ of admissible pairs of
\eqref{4.5:eq-admissible-pairs-a}.
\item[(ii)] $\mathbb{V}'^{J}(0)=0$.
\item[(iii)] $\mathbb{V}''^{J}$ is symmetric.
\item[(iv)] The gradient is bounded by
\begin{equation*}
|\mathbb{V}'^{J}(A')|_{(-3)}\le C^{J}_{4}\,N(A')^{2},
\qquad C^{J}_{4}:=C_{4}+C'a_{\gamma}.
\end{equation*}
Here $\varepsilon_3$ is the radius of Lemma~\ref{4.5:lem-nonlinear-correction}, $C_4$ is the
constant of \cite[(98)]{B-CMP102b}, $a_{\gamma}$ is the bound on $|J|_{(-3)}$ for
$(U,J)\in\mathfrak{S}$ in \eqref{4.5:eq-admissible-pairs-a}, and $C'$ is a constant depending only
on $d$, on the kernel constants $C_R$ of $\Lambda=G'RD^{*}_{U}$ and $B_0$ of $H^{J}$, and on the
constants $C_2$, $C_3$ of Lemma~\ref{4.5:lem-nonlinear-correction}.
\end{enumerate}
\end{lemma}

\begin{proof}
The gradient has seven terms, listed in \eqref{4.5:eq-higher-order-functional-a}. We treat the
first five together and the last two separately.

Throughout, $C$ denotes a constant depending only on $d$, on the kernel constants $C_R$ of
$\Lambda=G'RD^{*}_{U}$ (statement~\ref{4.5:not-background-propagators}) and $B_0$ of $H^{J}$
(Lemma~\ref{4.5:lem-free-current-operators}), and on the constants $C_2$, $C_3$ of
Lemma~\ref{4.5:lem-nonlinear-correction}.

\emph{The first five terms.} Start from the terms of Ba{\l}aban's gradient formula and its
estimates \cite[(85), (88), (89), (90)--(96)]{B-CMP102b}. In them, replace Ba{\l}aban's operators
$H$, $G$, $\mathfrak{D}$, $D$ and his current $J(U_0)$ by $H^{J}$, $G^{J}_{\pi}$,
$\mathfrak{D}^{J}$, $D^{J}$ and the free current $J$. The result is exactly the first five terms.

Ba{\l}aban's estimates of these terms, which lead to \cite[(98)]{B-CMP102b}, use only two kinds of
input. The first is the propagator estimate read in \cite{B-CMP102b} at this point and the
estimates \cite[(55), (73), (46)]{B-CMP102b},
which bound the propagators, the minimiser map, the derivative $\mathfrak{D}$ of the nonlinear
correction and the nonlinear correction itself. The second is bounds on the plaquette-local
coefficients of $V_0$ in terms of $|\partial U-1|$; Ba{\l}aban computes these coefficients without
any restriction on the variations \cite{B-CMP102b}.

For the operators at free current, the corresponding bounds are supplied by three statements. The
first is the statement~\ref{4.5:not-background-propagators} on propagators and constrained Green's
functions at an admissible background. The second is Lemma~\ref{4.5:lem-free-current-operators}
on the operators with the current as a free variable. The third is the nonlinear correction
to the linearising map at free current (Lemma~\ref{4.5:lem-nonlinear-correction}), whose parts (a)
and (c) bound $D^{J}$ and $\mathfrak{D}^{J}$.

The coefficients of $V_0$ are the same functions of $\eta^{-2}(\partial U-1)$ as in
\cite{B-CMP102b}, and $|\partial U-1|$ is controlled for $(U,J)\in\mathfrak{S}$ by
\eqref{4.5:eq-admissible-pairs-a}, so the bounds on them hold with the same constants. Since
$H^{J}$ and $G^{J}_{\pi}$ are analytic in $(\mathfrak{a},J)$
by the statement~\ref{4.5:not-background-propagators} and
Lemma~\ref{4.5:lem-free-current-operators}, and
$D^{J}$, $\mathfrak{D}^{J}$, $\mathfrak{L}^{J}$ are analytic in $(A',\mathfrak{a},J)$ by
Lemma~\ref{4.5:lem-nonlinear-correction}, the first five terms are analytic in
$(A',\mathfrak{a},J)$. Ba{\l}aban's estimates then apply verbatim and bound their
$|\cdot|_{(-3)}$-norm by $C_4N(A')^2$, as in \cite[(98)]{B-CMP102b}.

\emph{The sixth term.} The sixth term is
$-\Lambda^{T}\mathfrak{j}_J(H^{J}D^{J})$. It is a product of contributions from $\Lambda^{T}$, from
the insertion $\mathfrak{j}_J$, from $H^{J}$ and from $D^{J}$. Its $|\cdot|_{(-3)}$-norm is bounded
by
\begin{equation*}
C_R\cdot 2d\,a_{\gamma}\cdot B_0\cdot 4C_2\,N(A')^{2}.
\end{equation*}
The four factors bound, in this order, the four contributions: $C_R$ is the constant of
$\Lambda^{T}$; $2d\,a_{\gamma}$ comes from the bound on $\mathfrak{j}_J$ in
Lemma~\ref{4.5:lem-free-current-operators}, whose constant is $2d|J|_{(-3)}$, with $a_{\gamma}$ the
bound on $|J|_{(-3)}$ for $(U,J)\in\mathfrak{S}$; $B_0$ is the constant of $H^{J}$; and
$4C_2N(A')^{2}$ is the bound on $D^{J}$ of Lemma~\ref{4.5:lem-nonlinear-correction}(a).

\emph{The seventh term.} The seventh term is
$-\mathfrak{D}^{J,T}H^{J,T}\mathfrak{k}_J(\Lambda A')$. Its $|\cdot|_{(-3)}$-norm is bounded by
\begin{equation*}
O(C_3N(A'))\cdot C\cdot 2a_{\gamma}\cdot C\,N(A').
\end{equation*}
Here $O(C_3N(A'))$ is the $\mathcal{K}_1$-constant of $\mathfrak{D}^{J,T}$: by
Lemma~\ref{4.5:lem-nonlinear-correction}(c), applied with $\varepsilon_3$ replaced by any radius
$\varepsilon\le\varepsilon_3$ exceeding $N(A')$, this constant is $O(C_3\varepsilon)$, and letting
$\varepsilon$ decrease to $N(A')$ gives $O(C_3N(A'))$. The first $C$ is the constant of $H^{J,T}$,
$2a_{\gamma}$ comes from the bound on $\mathfrak{k}_J$ in Lemma~\ref{4.5:lem-free-current-operators},
whose constant is $2|J|_{(-3)}$, and $C\,N(A')$ bounds $\Lambda A'$.

\emph{Analyticity of the last two terms.} Both terms are analytic in $(A',\mathfrak{a},J)$ on
$\{N(A')<\varepsilon_3\}\times\mathfrak{S}$: the operators $\Lambda$ and $H^{J}$ are analytic
by the statement~\ref{4.5:not-background-propagators} and
Lemma~\ref{4.5:lem-free-current-operators}, $D^{J}$ and
$\mathfrak{D}^{J}$ are analytic by Lemma~\ref{4.5:lem-nonlinear-correction}(a) and (c), and
$\mathfrak{j}_J$, $\mathfrak{k}_J$ are linear in their argument and pointwise in $J$. Together with
the first five terms this gives (i).

\emph{Conclusion.} On $\{N(A')<\varepsilon_3\}$, the bounds on the sixth and seventh terms add up to
at most $C'a_{\gamma}N(A')^{2}$ for a constant $C'$. Adding the bound on the first five terms gives
the estimate in (iv), with $C^{J}_{4}=C_4+C'a_{\gamma}$.

The estimate in (iv) forces $\mathbb{V}'^{J}(0)=0$, which is (ii).

For (iii), note that $V^{J}$ is analytic on $\{N(A')<\varepsilon_3\}\times\mathfrak{S}$, being built
from $D^{J}$, $\mathfrak{L}^{J}$ (Lemma~\ref{4.5:lem-nonlinear-correction}), $H^{J}$, $\Lambda$ and
the entire functional $V_0$. Since $\mathbb{V}'^{J}=\nabla V^{J}$ is its gradient, the derivative
$\mathbb{V}''^{J}$ is the second derivative of $V^{J}$. This is a
symmetric bilinear form.
\end{proof}

\begin{definition}[The response chart and its representation identity]\label{4.5:def-response-chart}
Let $(U,J)\in\mathfrak{S}$ be an admissible pair. Let $G_0$, $H_0$, $P_0$, $Q$, $\mathsf{q}_0$, the
constant $a$, the operators $\Delta_a$ and $\Delta'$, and the block field $\tilde B$ associated with
a block field $B$ be those of Definition~\ref{4.5:not-background-propagators}. Let
$\bar\varepsilon_4>0$ be a radius depending on $(d,L)$ only, such that the map
\eqref{4.5:eq-response-chart-1} below has exactly one fixed point in the set of the first item. The
existence of such a radius is established in Theorem~\ref{4.5:thm-chart-existence}\,(K1).
For a block field $B$ the following objects are defined.
\begin{enumerate}
\item The field $\mathcal{A}^{J}_{0}(B)$ is the unique fixed point, in the set
$\{\mathcal{A} : N(\mathcal{A})\le\bar\varepsilon_4\}$, of the map
\begin{equation}\label{4.5:eq-response-chart-1}
\mathcal{A}\longmapsto\tilde{\mathfrak{G}}_0\bigl[\Theta_J(\mathcal{A}+H_0B)
-\mathbb{V}'^{J}(\mathcal{A}+H_0B)\bigr].
\end{equation}
Here $\tilde{\mathfrak{G}}_0=P_0G_0$ is the constrained Green's function, and $\mathbb{V}'^{J}$ is
the gradient of the higher-order action functional at free current
(Definition~\ref{4.5:def-higher-order-functional}).
\item The background potential, its nonlinear correction, the linearised background and the
associated field are
\begin{align*}
A_G&:=\mathcal{A}^{J}_{0}(B)+H_0B, & X_G&:=D^{J}(A_G),\\
\mathcal{H}^{J}(B)&:=\mathfrak{L}^{J}(A_G)=A_G-H^{J}X_G, & \Xi&:=\Theta_JA_G-\mathbb{V}'^{J}(A_G),
\end{align*}
where $D^{J}$ and $\mathfrak{L}^{J}=I-H^{J}D^{J}$ are the nonlinear correction and the linearising
map of Lemma~\ref{4.5:lem-nonlinear-correction}.
\item With $\mathfrak{j}_J$ the pointwise insertion of Lemma~\ref{4.5:lem-free-current-operators},
$\Lambda=G'RD^{*}_{U}$, and $F(U,\cdot)$ the curvature remainder of
Lemma~\ref{4.5:lem-curvature-remainder}, put
\begin{align*}
\Lambda^{J}(B)&:=P_0^{T}\Xi+Q^{*}(\mathsf{q}_0-a)\tilde B-Q^{*}\mathsf{k}X_G
-\Lambda^{T}\mathfrak{j}_J(H^{J}X_G),\\
\mathfrak{J}(B)&:=\Lambda^{J}(B)-\Delta'\mathcal{H}^{J}(B)+F\bigl(U,\mathcal{H}^{J}(B)\bigr).
\end{align*}
\item For a direction $v$ of the block field (here $v$ denotes only such a direction) and
$t_{\square}\in\mathbb{C}$, with $\delta\mathcal{H}^{J}(B)[v]$ and $\delta\Lambda^{J}(B)[v]$ the
derivatives of $B\mapsto\mathcal{H}^{J}(B)$ and of $B\mapsto\Lambda^{J}(B)$ in the direction $v$, put
\begin{align*}
\mathcal{H}^{J}(B,t_{\square})&:=\mathcal{H}^{J}(B)+t_{\square}\,\delta\mathcal{H}^{J}(B)[v],\\
\mathfrak{J}(B,t_{\square})&:=\Lambda^{J}(B)+t_{\square}\,\delta\Lambda^{J}(B)[v]
-\Delta'\mathcal{H}^{J}(B,t_{\square})+F\bigl(U,\mathcal{H}^{J}(B,t_{\square})\bigr).
\end{align*}
\item The datum produced by the chart is the orbit of the pair
$\bigl(e^{i\eta\mathcal{H}^{J}(B)}U,\;J+\mathfrak{J}(B)\bigr)$.
\end{enumerate}
The fixed-point equation of the first item is equivalent to the equation
\begin{equation*}
\mathcal{A}=\tilde G^{J}\bigl[\Theta_JH_0B-\mathbb{V}'^{J}(\mathcal{A}+H_0B)\bigr],
\end{equation*}
which is the current-dependent form of \cite[(180)]{B-CMP102b}, that is, of the equation
\begin{equation*}
A_0=-GJ+G\Delta^{(2)}(A_0+H_0B)-G\tfrac{\delta}{\delta A'}V(A_0+H_0B)
\end{equation*}
of \cite[(143)]{B-CMP102b}, in Ba{\l}aban's notation there, with the source term $-GJ$ removed. In
\cite[(1.4)]{B-CMP116} the same background $A_G$ is determined by the equation
$A_0+G\bigl(\tfrac{\delta}{\delta A'}V\bigr)(H_0B'+A_0)=0$, in which the block field is written $B'$,
with the terms split differently: the
operator $H_0$ there already contains $\Delta^{(2)}$, so that equation carries no $\Theta$-source;
the use made of it in \cite[(1.18)]{B-CMP116} is linear in the bounds on the propagators, so that the
difference between the two splittings is immaterial there. The field $\mathcal{H}^{J}(B)$ is the
current-dependent form of the field of \cite[(179)]{B-CMP102b} and is the field of
\cite[(1.2)]{B-CMP116}.

At every $(U,J)\in\mathfrak{S}$ the chart satisfies the \emph{representation identity}
\begin{equation}\label{4.5:eq-response-chart-a}
\Delta_aA_G=P_0^{T}\Xi+Q^{*}\mathsf{q}_0\tilde B,
\end{equation}
and consequently
\begin{equation}\label{4.5:eq-response-chart-2}
\begin{split}
\mathfrak{J}(B)&=(\Delta_a-Q^{*}aQ-\Delta')A_G-\mathcal{K}_1X_G+F\bigl(U,\mathcal{H}^{J}(B)\bigr),\\
\mathcal{K}_1X_G&:=Q^{*}\mathsf{k}X_G+\Lambda^{T}\mathfrak{j}_J(H^{J}X_G)-\Delta'H^{J}X_G.
\end{split}
\end{equation}
In \eqref{4.5:eq-response-chart-2}, $\mathcal{K}_1X_G$ denotes the displayed expression in $X_G$; the
letter does not refer here to the class of kernels $\mathcal{K}_1$.
\end{definition}

\begin{proof}
We first prove \eqref{4.5:eq-response-chart-a}. By the first item, $\mathcal{A}^{J}_{0}(B)$ is a
fixed point of the map \eqref{4.5:eq-response-chart-1}. Since
$\mathcal{A}^{J}_{0}(B)+H_0B=A_G$ by the second item, the fixed-point equation reads
\begin{equation*}
\mathcal{A}^{J}_{0}(B)=\tilde{\mathfrak{G}}_0\bigl[\Theta_JA_G-\mathbb{V}'^{J}(A_G)\bigr]
=\tilde{\mathfrak{G}}_0\,\Xi=P_0G_0\,\Xi,
\end{equation*}
by the definition of $\Xi$ and $\tilde{\mathfrak{G}}_0=P_0G_0$. Hence
\begin{equation*}
A_G=P_0G_0\,\Xi+H_0B.
\end{equation*}
We apply $\Delta_a$ to the two terms on the right. The formulas of
Definition~\ref{4.5:not-background-propagators}, which hold at every $(U,J)\in\mathfrak{S}$, give
\begin{equation*}
\Delta_a\bigl(P_0G_0\,\Xi\bigr)=P_0^{T}\Xi,\qquad \Delta_a\bigl(H_0B\bigr)=Q^{*}\mathsf{q}_0\tilde B.
\end{equation*}
Adding the two identities gives $\Delta_aA_G=P_0^{T}\Xi+Q^{*}\mathsf{q}_0\tilde B$, which is
\eqref{4.5:eq-response-chart-a}.

We now derive \eqref{4.5:eq-response-chart-2}. By \eqref{4.5:eq-response-chart-a},
\begin{equation*}
P_0^{T}\Xi+Q^{*}(\mathsf{q}_0-a)\tilde B=\Delta_aA_G-Q^{*}a\tilde B,
\end{equation*}
so the definition of $\Lambda^{J}(B)$ in the third item becomes
\begin{equation*}
\Lambda^{J}(B)=\Delta_aA_G-Q^{*}a\tilde B-Q^{*}\mathsf{k}X_G-\Lambda^{T}\mathfrak{j}_J(H^{J}X_G).
\end{equation*}
By the second item, $\Delta'\mathcal{H}^{J}(B)=\Delta'A_G-\Delta'H^{J}X_G$. Inserting both expressions
into the definition of $\mathfrak{J}(B)$ gives
\begin{align*}
\mathfrak{J}(B)&=\Delta_aA_G-Q^{*}a\tilde B-Q^{*}\mathsf{k}X_G-\Lambda^{T}\mathfrak{j}_J(H^{J}X_G)
-\Delta'A_G+\Delta'H^{J}X_G+F\bigl(U,\mathcal{H}^{J}(B)\bigr)\\
&=(\Delta_a-\Delta')A_G-Q^{*}a\tilde B-\mathcal{K}_1X_G+F\bigl(U,\mathcal{H}^{J}(B)\bigr),
\end{align*}
where the last line collects the three terms in $X_G$ into $\mathcal{K}_1X_G$ as defined in
\eqref{4.5:eq-response-chart-2}. It remains to write $Q^{*}a\tilde B$ as $Q^{*}aQA_G$. By the formulas
of Definition~\ref{4.5:not-background-propagators}, $QP_0=0$ and $QH_0B=\tilde B$; applied to
$A_G=P_0G_0\,\Xi+H_0B$ they give $QA_G=\tilde B$, hence $Q^{*}a\tilde B=Q^{*}aQA_G$. With this
substitution the identity above becomes the first line of \eqref{4.5:eq-response-chart-2}.
\end{proof}

\begin{theorem}[Existence, uniqueness and analyticity of the response chart]
\label{4.5:thm-chart-existence}
Assume the following two statements:
\begin{itemize}
\item the conclusions of Lemma~\ref{4.5:lem-free-current-operators} (the operators with the current as a free variable);
\item the conclusions of Lemma~\ref{4.5:lem-nonlinear-correction} (the nonlinear correction to the linearising map at free current).
\end{itemize}
Then there exist the following constants:
\begin{itemize}
\item constants
\[
a_0,\ a_1,\ a_\gamma,\ a_3,\ \bar\varepsilon_4,\ c_9,\ C_f,\ C^{K}_{9},\ C'_{9},\ C_{\mathfrak{J}},\ C'_{\mathfrak{J}},\
C^{\boxplus,K}_{2},\ C^{K}_{189},\ C^{K}_{190},
\]
depending only on $d$ and $L$;
\item a number $c_*\in(0,1)$, a fraction of a decay rate, fixed in clause (K10);
\item an integer $c_{\mathbb{K}}$ independent of all parameters;
\item the constants $B_0, B'_0$ entering the propagator bounds of
Definition~\ref{4.5:not-background-propagators}, the constant $\delta_0$ of \cite[(2.59)]{B-CMP96},
the constants $C_2, C_3$ of Lemma~\ref{4.5:lem-nonlinear-correction}, and the constant $C_4$ of
Lemma~\ref{4.5:lem-gradient-bound}.
\end{itemize}
All of them are uniform in $k$, in $K$, in the sequence of domains $\{\Omega_j\}$ of
\cite[(1)]{B-CMP102b}, in the domain $X$ of which this sequence is the admissible tower
(arbitrarily large $X$ included), and in $M_1$. All
statements below hold under \eqref{4.5:eq-kernel-algebra-a} and under \cite[(2.59)]{B-CMP96} imposed
at finitely many absolute fractions of $\delta_0$, in the sense of
Lemma~\ref{4.5:lem-kernel-algebra}. For every $(U,J)\in\mathfrak{S}$, with $\mathfrak{S}$ as in
\eqref{4.5:eq-admissible-pairs-a}, and every $B$ as in Definition~\ref{4.5:def-response-chart} with
$|B|<c_9$, the clauses (K1)--(K10) hold.
Clause (K1) is stated and proved below; each of the clauses (K2), \dots, (K10) is stated, with the
same constants, and proved in the statement of this section that bears its letter in its name.

\smallskip
\noindent (K1) \emph{Existence, uniqueness, joint analyticity.} The field $\mathcal{A}^{J}_{0}(B)$ of Definition~\ref{4.5:def-response-chart} exists and is unique in $\{N\le\bar\varepsilon_4\}$. It satisfies
\begin{equation*}
N_{\beta}\bigl(\mathcal{A}^{J}_{0}\bigr)\le C_f|B|.
\end{equation*}
The maps
\[
\mathcal{A}^{J}_{0},\quad A_G,\quad D^{J}(A_G),\quad \mathcal{H}^{J},\quad \Lambda^{J},\quad \mathfrak{J}
\]
(with $D^{J}$ as in Lemma~\ref{4.5:lem-nonlinear-correction} and the others as in
Definition~\ref{4.5:def-response-chart}) are jointly analytic in $(B,\mathfrak{a},J)$ on
$\{|B|<c_9\}\times\mathfrak{S}$.
\end{theorem}

\begin{proof}[Proof of \textup{(K1)}]
Fix $(U,J)\in\mathfrak{S}$ and $B$ as in Definition~\ref{4.5:def-response-chart} with $|B|<c_9$.
Define the map
\begin{equation}\label{4.5:eq-chart-existence-1}
\Phi(\mathcal{A}):=\tilde{\mathfrak{G}}_0\bigl[\Theta_J(\mathcal{A}+H_0B)
-\mathbb{V}'^{J}(\mathcal{A}+H_0B)\bigr],
\end{equation}
where $H_0$ is the operator of Definition~\ref{4.5:def-response-chart}.
By Definition~\ref{4.5:def-response-chart}, $\mathcal{A}^{J}_{0}(B)$ is defined as the fixed point
of $\Phi$ in the ball $\{N(\mathcal{A})\le\bar\varepsilon_4\}$; we now show that $\Phi$ has exactly
one fixed point there. Choose the constants so that
\begin{equation}\label{4.5:eq-chart-existence-2}
B_0c_9\le\bar\varepsilon_4\le a_3/4 .
\end{equation}

\emph{The ball is mapped into itself.} Let $N(\mathcal{A})\le\bar\varepsilon_4$. Since
$N(H_0B)\le B_0|B|$ by the propagator bounds of Definition~\ref{4.5:not-background-propagators},
and $|B|<c_9$, the first inequality in \eqref{4.5:eq-chart-existence-2} gives
$N(H_0B)\le\bar\varepsilon_4$; hence the argument of $\Theta_J$ and $\mathbb{V}'^{J}$ in
\eqref{4.5:eq-chart-existence-1} satisfies $N(\mathcal{A}+H_0B)\le 2\bar\varepsilon_4$. Three
bounds enter:
\begin{itemize}
\item the norm of the constrained Green's function $\tilde{\mathfrak{G}}_0$ is bounded by $B'_0$, by
the propagator bounds of Definition~\ref{4.5:not-background-propagators};
\item the operator $\Theta_J$ of Lemma~\ref{4.5:lem-free-current-operators}, whose conclusions are
assumed in the hypotheses, is bounded by $C_\Theta a_\gamma$, with the constant $C_\Theta$ of that
lemma;
\item by the quadratic bound on the gradient of the higher-order functional
(Lemma~\ref{4.5:lem-gradient-bound}), $\mathbb{V}'^{J}$ is bounded by $C^{J}_{4}N(\cdot)^2$, with
the constant $C^{J}_{4}$ of that lemma.
\end{itemize}
Together they give
\begin{equation*}
N\bigl(\Phi(\mathcal{A})\bigr)
\le B'_0\bigl[2C_\Theta a_\gamma+4C^{J}_{4}\bar\varepsilon_4\bigr]\bar\varepsilon_4
\le\bar\varepsilon_4 ,
\end{equation*}
where the last inequality holds once $a_\gamma$ and $\bar\varepsilon_4$ are small enough.

\emph{Contraction.} Let $\mathcal{A}_1,\mathcal{A}_2$ lie in the ball. Since $\Theta_J$ is linear,
the $\Theta_J$-term contributes at most $B'_0C_\Theta a_\gamma N(\mathcal{A}_1-\mathcal{A}_2)$ to
$N\bigl(\Phi(\mathcal{A}_1)-\Phi(\mathcal{A}_2)\bigr)$. For the $\mathbb{V}'^{J}$-term we use the
Cauchy estimate at radius $2\bar\varepsilon_4$, which is admissible because
$2\bar\varepsilon_4\le a_3/2$ by \eqref{4.5:eq-chart-existence-2}. It is applied to the analytic map
$\mathbb{V}'^{J}$ with the bound of Lemma~\ref{4.5:lem-gradient-bound}. Combining the two
contributions we obtain the analogue at free current of \cite[(120)]{B-CMP102b}:
\begin{equation*}
\begin{aligned}
N\bigl(\Phi(\mathcal{A}_1)-\Phi(\mathcal{A}_2)\bigr)
&\le B'_0\bigl[C_\Theta a_\gamma+8C^{J}_{4}\bar\varepsilon_4\bigr]N(\mathcal{A}_1-\mathcal{A}_2)\\
&\le\tfrac12\,N(\mathcal{A}_1-\mathcal{A}_2),
\end{aligned}
\end{equation*}
again for $a_\gamma$ and $\bar\varepsilon_4$ small enough.

Thus $\Phi$ is a contraction of the closed ball $\{N\le\bar\varepsilon_4\}$ into itself. It therefore has exactly one fixed point there, which is $\mathcal{A}^{J}_{0}(B)$. This proves existence and uniqueness.

\emph{The bound.} From $\mathcal{A}^{J}_{0}=\Phi(\mathcal{A}^{J}_{0})$ and the contraction estimate with $\mathcal{A}_2=0$,
\begin{equation*}
N(\mathcal{A}^{J}_{0})
\le N\bigl(\Phi(\mathcal{A}^{J}_{0})-\Phi(0)\bigr)+N\bigl(\Phi(0)\bigr)
\le\tfrac12N(\mathcal{A}^{J}_{0})+N\bigl(\Phi(0)\bigr).
\end{equation*}
Hence $N(\mathcal{A}^{J}_{0})\le 2N(\Phi(0))$. At $\mathcal{A}=0$ the argument of $\Theta_J$ and
$\mathbb{V}'^{J}$ in \eqref{4.5:eq-chart-existence-1} is $H_0B$, with $N(H_0B)\le B_0|B|$, so the
three bounds above give
\begin{equation*}
\begin{aligned}
N\bigl(\Phi(0)\bigr)
&\le B'_0\bigl[C_\Theta a_\gamma B_0|B|+C^{J}_{4}B_0^2|B|^2\bigr]\\
&\le B'_0B_0\bigl[C_\Theta a_\gamma+C^{J}_{4}B_0c_9\bigr]|B| .
\end{aligned}
\end{equation*}
Hence $N(\mathcal{A}^{J}_{0})\le C_f|B|$ with
$C_f:=2B'_0B_0\bigl[C_\Theta a_\gamma+C^{J}_{4}B_0c_9\bigr]$. The corresponding bound in the norm
$N_{\beta}$ follows from the propagator bounds of Definition~\ref{4.5:not-background-propagators}.

\emph{Analyticity of the fixed point.} The iterates $\Phi^{n}(0)$, $n\ge0$, are analytic in $(B,\mathfrak{a},J)$, for three reasons:
\begin{itemize}
\item $\tilde{\mathfrak{G}}_0$ is analytic in $\mathfrak{a}$ on $\mathfrak{S}$ by Definition~\ref{4.5:not-background-propagators};
\item $\Theta_J$ is analytic in $(\mathfrak{a},J)$ by Lemma~\ref{4.5:lem-free-current-operators};
\item $\mathbb{V}'^{J}$ is analytic in $(A',\mathfrak{a},J)$ by Lemma~\ref{4.5:lem-gradient-bound}.
\end{itemize}
Moreover $H_0B$ is linear in $B$. By the contraction estimate the iterates converge to $\mathcal{A}^{J}_{0}$ uniformly on $\{|B|<c_9\}\times\mathfrak{S}$. A uniform limit of analytic maps is analytic, so $\mathcal{A}^{J}_{0}$ is jointly analytic in $(B,\mathfrak{a},J)$.

This is the argument of \cite[proof of Proposition~6]{B-CMP102b}. That proof uses only three
ingredients: the norm of the propagator, the gradient bound \cite[(98)]{B-CMP102b}, and the
analyticity of $\delta V/\delta A'$ as a function of $A'$, from which the analyticity of the
approximations follows. After \cite[(143)]{B-CMP102b} it is observed that that equation has all the
properties of \cite[(111)]{B-CMP102b}, so that Proposition~6 holds for it as well. In
\cite[p.~2]{B-CMP116} it is observed that in these equations the propagators may be replaced by
arbitrary operators with the same regularity properties and satisfying the same bounds, these bounds
being all that matters. The equation for $\mathcal{A}^{J}_{0}(B)$ is \cite[(143)]{B-CMP102b} with
the source term removed, the propagator $G$ of \cite{B-CMP102b} replaced by
$\tilde{\mathfrak{G}}_0$, the operator $\Delta^{(2)}$ of \cite[(143)]{B-CMP102b} replaced by
$\Theta_J$, and $\delta V/\delta A'$ replaced by $\mathbb{V}'^{J}$, with the bounds listed above.

\emph{The remaining maps.} By Definition~\ref{4.5:def-response-chart}, the remaining maps are obtained from $\mathcal{A}^{J}_{0}$ by composition:
\begin{itemize}
\item $A_G=\mathcal{A}^{J}_{0}+H_0B$;
\item $D^{J}(A_G)$ and $\mathcal{H}^{J}(B)=\mathfrak{L}^{J}(A_G)$;
\item $\Lambda^{J}(B)$, built from $\Theta_JA_G-\mathbb{V}'^{J}(A_G)$, from $B$, from $D^{J}(A_G)$ and from the insertion $\mathfrak{j}_J$;
\item $\mathfrak{J}(B)=\Lambda^{J}(B)-\Delta'\mathcal{H}^{J}(B)+F(U,\mathcal{H}^{J}(B))$.
\end{itemize}
Since $N(A_G)\le2\bar\varepsilon_4\le a_3/2$ by \eqref{4.5:eq-chart-existence-2}, $A_G$ lies in the
domain of Lemma~\ref{4.5:lem-nonlinear-correction} and of Lemma~\ref{4.5:lem-gradient-bound}.
Analyticity of the building blocks is supplied as follows:
\begin{itemize}
\item $D^{J}$ and $\mathfrak{L}^{J}$ are analytic by Lemma~\ref{4.5:lem-nonlinear-correction} (by
hypothesis);
\item $\Theta_J$ and $\mathfrak{j}_J$ are analytic by Lemma~\ref{4.5:lem-free-current-operators} (by
hypothesis);
\item the operators $H_0$, $H^{J}$, $P_0^{T}$, $Q^{*}$, $\mathsf{q}_0$, $\mathsf{k}$ and
$\Lambda^{T}$ entering $\mathcal{H}^{J}$ and $\Lambda^{J}$ are analytic in $\mathfrak{a}$, and in $J$
where they depend on it, by Definition~\ref{4.5:not-background-propagators} and by
Lemma~\ref{4.5:lem-free-current-operators} (by hypothesis);
\item $F$ is local and analytic by Lemma~\ref{4.5:lem-curvature-remainder};
\item $\mathbb{V}'^{J}$ is analytic by Lemma~\ref{4.5:lem-gradient-bound}.
\end{itemize}
The composition of these analytic maps with the analytic map $(B,\mathfrak{a},J)\mapsto\mathcal{A}^{J}_{0}$ is jointly analytic on $\{|B|<c_9\}\times\mathfrak{S}$.
\end{proof}

We prove clause (K2) of Theorem~\ref{4.5:thm-chart-existence}.
\label{4.5:prf-chart-structure}
Let $(U,J)\in\mathfrak{S}$. At $B=0$,
\begin{equation*}
\mathcal{A}^{J}_{0}=A_{G}=X=\mathcal{H}^{J}=\Lambda^{J}=\mathfrak{J}=0 .
\end{equation*}
Here $X$ is the object of the chart listed in clause (K1) of
Theorem~\ref{4.5:thm-chart-existence} among the jointly analytic quantities; it is not a
large-field region.
For $|B|<c_{9}$ and for every $j$, the identity
\begin{equation}\label{4.5:eq-chart-structure-a}
Q_{j}\bigl(U;\eta\,\mathcal{H}^{J}(B)\bigr)=B \quad\text{on } \Lambda_{j}
\end{equation}
holds exactly; it is the counterpart at free current $J$ of \cite[(20)]{B-CMP102b}.
No claim is made that the counterpart of \cite[(21)]{B-CMP102b}, the relation
$R D^{*}_{U}\mathcal{H}^{J}=0$, holds off the diagonal.

\begin{proof}
Vanishing at $B=0$. The gradient of the higher-order functional at free current vanishes
at the zero field, $\mathbb{V}'^{J}(0)=0$. Hence at $B=0$ the zero field is a fixed point
of the equation that defines $\mathcal{A}^{J}_{0}$. Since $N(0)=0\le\bar\varepsilon_{4}$,
the uniqueness assertion of clause (K1) shows that the zero field is the fixed point, that is,
$\mathcal{A}^{J}_{0}(0)=0$. In addition the map $F$ of the construction of the chart
satisfies $F(U,0)=0$. For the remaining quantities $A_{G}$, $X$, $\mathcal{H}^{J}$,
$\Lambda^{J}$, $\mathfrak{J}$ of the first display, the argument at $B=0$ consists of
inserting these two facts, $\mathcal{A}^{J}_{0}(0)=0$ and $F(U,0)=0$, into their defining
formulas, which then give zero; hence all six quantities of the first display vanish at
$B=0$. Both facts hold at every $(U,J)\in\mathfrak{S}$.

The identity \eqref{4.5:eq-chart-structure-a}. Write the background potential as
$A_{G}=\tilde{\mathfrak{G}}_{0}\varpi+H_{0}B$. Here $\varpi$ is the vector through which the
constrained Green's function $\tilde{\mathfrak{G}}_{0}$ contributes the first summand of
$A_{G}$, and $H_{0}$ is the minimiser map without current, the counterpart of $H^{J}$.
The defining formulas of $\tilde{\mathfrak{G}}_{0}$ and of $H_{0}$ give, with $\ell$ as in
the construction of the chart,
\begin{equation*}
\ell Q A_{G}=\ell Q\tilde{\mathfrak{G}}_{0}\varpi+\ell Q H_{0}B=0+B .
\end{equation*}
Insert this into the free-current form of \cite[(48)]{B-CMP102b}. The result is
\eqref{4.5:eq-chart-structure-a} on each layer $\Lambda_{j}$.

Only these two formulas enter the argument; no estimate is used. Hence
\eqref{4.5:eq-chart-structure-a} holds as an identity of formulas wherever the chart is
defined, that is, at every $(U,J)\in\mathfrak{S}$ and every $|B|<c_{9}$.
\end{proof}

\begin{lemma}[Linear bounds on the background potential and the current correction]
\label{4.5:prf-chart-linear-bounds}
Let the constants be those of Theorem~\ref{4.5:thm-chart-existence}, let
$(U,J)\in\mathfrak{S}$ and $|B|<c_9$. Assume the bounds and the generalised random-walk
expansion of $\mathsf{q}_0=(QG_0Q^{*})^{-1}$ of
Remark~\ref{4.5:rem-covariance-inverse-walks}, which are printed without proof in
\cite[(3.132)]{B-CMP99b} and cited as printed. Then clause (K3) of
Theorem~\ref{4.5:thm-chart-existence} holds:
\begin{align}
N_{\beta}(\mathcal{H}^{J}(B)) &\le C^{K}_{9}|B|,
\label{4.5:eq-chart-linear-bounds-1}\\
|\Lambda^{J}(B)|_{(-3)},\ |\mathfrak{J}(B)|_{(-3)} &\le C_{\mathfrak{J}}|B|,
\label{4.5:eq-chart-linear-bounds-2}\\
|P_1\mathcal{H}^{J}(B)|_{(-3)} &\le C\,C^{K}_{9}|B|,
\label{4.5:eq-chart-linear-bounds-3}
\end{align}
where $P_1=D_UPD^{*}_U$ and $C$ is a constant depending on $d$ and $L$ only.
\end{lemma}

\begin{proof}
We write $\mathsf{X}$ for the higher-order term of the response chart, the argument of
$H^{J}$ in the representation identity \eqref{4.5:eq-response-chart-a}, and $\mathsf{w}$ for
the first source term there, so that $\Delta_aA_G=P_0^{T}\mathsf{w}+Q^{*}\mathsf{q}_0\tilde B$,
with $\tilde B$ the datum entering that identity. Put
$\mathfrak{b}:=(C_f+B_0)|B|$. Apart from the three appeals named separately below (the
hypothesis taken from Remark~\ref{4.5:rem-covariance-inverse-walks}, \cite[(3.69)]{B-CMP99b}
and Lemma~\ref{4.5:lem-curvature-remainder}), every estimate in this proof is one of the
estimates of \cite[(133)--(136)]{B-CMP102b}, read as bounds on kernels in the classes
$\mathcal{K}_0,\mathcal{K}_{-1},\dots$; the constants $C_{\Theta}$, $C^{J}_{4}$ and $C_R$
below are constants depending on $d$ and $L$ only, and $C'$ denotes a constant depending on
$d$ and $L$ only, not necessarily the same at each occurrence.

\emph{The background potential.} By the estimates just cited,
\begin{equation*}
N_{\beta}(A_G)\le\mathfrak{b},\qquad |\mathsf{X}|\le 4C_2\mathfrak{b}^2,\qquad
N_{\beta}(H^{J}\mathsf{X})\le B_0|\mathsf{X}|.
\end{equation*}
In the response chart (Definition~\ref{4.5:def-response-chart}), $\mathcal{H}^{J}(B)$ is the
sum of $A_G$ and $H^{J}\mathsf{X}$, the latter with the sign fixed there. The ceiling on
$c_9$ in Theorem~\ref{4.5:thm-chart-existence}, together with $|B|<c_9$, makes
$4C_2B_0\mathfrak{b}\le1$ and $2\mathfrak{b}\le1$. Hence, by the triangle inequality,
\begin{equation*}
N_{\beta}(\mathcal{H}^{J}(B))\le N_{\beta}(A_G)+N_{\beta}(H^{J}\mathsf{X})
\le\mathfrak{b}+4C_2B_0\mathfrak{b}^2\le2\mathfrak{b}=2(C_f+B_0)|B|,
\end{equation*}
which is \eqref{4.5:eq-chart-linear-bounds-1} once $C^{K}_{9}$ is chosen at least
$2(C_f+B_0)$.

\emph{The current correction.} The representation identity
\eqref{4.5:eq-response-chart-a} expresses $\mathfrak{J}(B)$, and with it $\Lambda^{J}(B)$,
through the terms $P_0^{T}\mathsf{w}$, $Q^{*}(\mathsf{q}_0-a)\tilde B$, $\Delta'A_G$,
$Q^{*}\mathsf{k}\mathsf{X}$, $\Lambda^{T}\mathfrak{j}_J(H^{J}\mathsf{X})$,
$\Delta'H^{J}\mathsf{X}$ and $F(U,\mathcal{H}^{J}(B))$; here $\Lambda^{J}(B)$ is the part of
this expression that carries neither $\Delta'$ nor $F$
(Definition~\ref{4.5:def-response-chart}), so that its bound is read off the same list.
We bound the terms in turn.
First,
\begin{equation*}
|\mathsf{w}|_{(-3)}\le C_{\Theta}a_{\gamma}\mathfrak{b}+C^{J}_{4}\mathfrak{b}^2,
\end{equation*}
and $P_0^{T}\in\mathcal{K}_0$, which bounds $P_0^{T}\mathsf{w}$. Next, by the hypothesis
taken from Remark~\ref{4.5:rem-covariance-inverse-walks},
\begin{equation*}
|Q^{*}(\mathsf{q}_0-a)\tilde B|_{(-3)}\le C'|\tilde B|_{(-1)}=C'|B|.
\end{equation*}
Further,
\begin{equation*}
|Q^{*}\mathsf{k}\mathsf{X}|_{(-3)}\le C'|\mathsf{X}|,\qquad
|\Lambda^{T}\mathfrak{j}_J(H^{J}\mathsf{X})|_{(-3)}\le C_R\,2d\,a_{\gamma}B_0|\mathsf{X}|,
\end{equation*}
and the two terms carrying $\Delta'$ are estimated by \cite[(3.69)]{B-CMP99b}. Finally,
since the norm $N$ is dominated by $N_{\beta}$, the bound obtained in the first step of this
proof gives
\begin{equation*}
N(\mathcal{H}^{J}(B))\le N_{\beta}(A_G)+N_{\beta}(H^{J}\mathsf{X})\le2\mathfrak{b}\le1,
\end{equation*}
the last two inequalities resting on the ceiling on $c_9$; hence the locality and
size lemma for the curvature remainder (Lemma~\ref{4.5:lem-curvature-remainder}) applies
with $A=\mathcal{H}^{J}(B)$ and gives
\begin{equation*}
|F(U,\mathcal{H}^{J}(B))|_{(-3)}\le C_F(a_0+a_1+2\mathfrak{b})\,2\mathfrak{b}.
\end{equation*}
Each of these bounds is at most a constant depending on $d$ and $L$ times $|B|$, because
$\mathfrak{b}$ is a fixed multiple of $|B|$ and $\mathfrak{b}\le\tfrac12$. This proves
\eqref{4.5:eq-chart-linear-bounds-2}.

\emph{The last bound.} By the propagator bounds at an admissible background of
Definition~\ref{4.5:not-background-propagators} (\cite[(3.49)]{B-CMP99b}),
$P_1\in\mathcal{K}_{-2}$; applied to \eqref{4.5:eq-chart-linear-bounds-1} this gives
\eqref{4.5:eq-chart-linear-bounds-3}.
\end{proof}

\begin{lemma}[First variation of the chart and its exponential decay]\label{4.5:prf-chart-variation-decay}
Let the constants be those of Theorem~\ref{4.5:thm-chart-existence} and let $(U,J)\in\mathfrak{S}$.
Write $A_G$, $\mathcal{H}^{J}=\mathcal{H}^{J}(B)$, $\Lambda^{J}=\Lambda^{J}(B)$ and
$\mathfrak{J}=\mathfrak{J}(B)$ for the pieces of the response chart of
Definition~\ref{4.5:def-response-chart}. Write $X:=D^{J}(A_G)$ for the nonlinear correction of
Lemma~\ref{4.5:lem-nonlinear-correction} at the background potential. For a set $S$ and
$\kappa\ge 0$, let $\rho(y)=e^{\kappa d(y,S)}$ be the weight of the norms $N^{\rho}$ and
$|\cdot|_{(\alpha),\rho}$. Then clauses (K4) and (K8) of Theorem~\ref{4.5:thm-chart-existence}
hold in the following form.

\smallskip
\noindent\textup{(K4)} Let $|B|\le c_9/2$ and let $v$ be a direction of the block variable $B$.
The first variations along $v$ are
\begin{equation*}
\delta A_G=\mathfrak{A}_0+H_0v,\qquad \delta X=\mathfrak{D}^{J}(A_G)\,\delta A_G ,
\end{equation*}
where $\mathfrak{A}_0$ and $\delta\mathcal{H}^{J}[v]$ are given by
\begin{equation}\label{4.5:eq-chart-variation-decay-a}
\begin{gathered}
\bigl(I+\tilde G^{J}\mathbb{V}''^{J}(A_G)\bigr)\mathfrak{A}_0
=\tilde G^{J}\bigl[\Theta_J-\mathbb{V}''^{J}(A_G)\bigr]H_0v,\\
\delta\mathcal{H}^{J}[v]=\mathfrak{T}^{J}(A_G)\,\delta A_G ,
\end{gathered}
\end{equation}
and the variations of $\Lambda^{J}$ and $\mathfrak{J}$ are
\begin{equation}\label{4.5:eq-chart-variation-decay-1}
\begin{split}
\delta\Lambda^{J}[v]={}&P_0^{T}\bigl[\Theta_J-\mathbb{V}''^{J}(A_G)\bigr]\delta A_G
+Q^{*}(\mathsf{q}_0-a)\ell^{-1}v\\
&-Q^{*}\mathsf{k}\,\delta X-\Lambda^{T}\mathfrak{j}_J(H^{J}\delta X),\\
\delta\mathfrak{J}[v]={}&\delta\Lambda^{J}[v]-\Delta'\delta\mathcal{H}^{J}[v]
+\partial_AF(U,\mathcal{H}^{J})\,\delta\mathcal{H}^{J}[v].
\end{split}
\end{equation}
They obey $N_{\beta}(\delta\mathcal{H}^{J}[v])\le C_9'|v|$ and
$|\delta\mathfrak{J}[v]|_{(-3)}\le C_{\mathfrak{J}}'|v|$. No second derivative of a response is
taken.

\smallskip
\noindent\textup{(K8)} Let $A'$ lie in the domain on which the quadratic bound on the gradient
(Lemma~\ref{4.5:lem-gradient-bound}) is stated. Let $S$ be arbitrary and
$\kappa\le c_{*}\delta_0$. Let $c\in\mathfrak{B}$ be a block bond at $y'\in\Lambda_{j'}$, let
$y\in\Lambda_j$ and let $x\in\Delta(y)$. Let $\zeta$ be a localising function, and let
$\mathfrak{h}_{\zeta}(x,c)$ denote the $\zeta$-localised H\"older member of the kernel of
$\delta\mathcal{H}^{J}/\delta B$, in the sense of the third member of \cite[(190)]{B-CMP102b}.
Write $\mathfrak{E}:=e^{-c_{*}\delta_0 d(y,y')}$. Then
\begin{equation}\label{4.5:eq-chart-variation-decay-b}
\begin{gathered}
|\mathbb{V}''^{J}(A')\mathfrak{A}|_{(-3),\rho}\le C^{K}_{189}\,\varepsilon_3\,N^{\rho}(\mathfrak{A}),\\
\Bigl|\frac{\delta\mathcal{H}^{J}(x)}{\delta B(c)}\Bigr|
\le C^{K}_{190}\,\ell_j^{-1}\ell_{j'}^{-d}\,\mathfrak{E},\qquad
\Bigl|\nabla_U\frac{\delta\mathcal{H}^{J}(x)}{\delta B(c)}\Bigr|
\le C^{K}_{190}\,\ell_j^{-2}\ell_{j'}^{-d}\,\mathfrak{E},\\
\mathfrak{h}_{\zeta}(x,c)\le C^{K}_{190}\bigl(\|\zeta\|_{\beta}+|\zeta|\bigr)
\ell_j^{-2-\beta}\ell_{j'}^{-d}\,\mathfrak{E},\qquad
\Bigl|\frac{\delta\mathfrak{J}(x)}{\delta B(c)}\Bigr|
\le C^{K}_{190}\,\ell_j^{-3}\ell_{j'}^{-d}\,\mathfrak{E}.
\end{gathered}
\end{equation}
The first line of \eqref{4.5:eq-chart-variation-decay-b} is the weighted Hessian bound. The
next three bounds are the first three members of \cite[(190)]{B-CMP102b}. The last bound, the
kernel of $\mathfrak{J}$, stands in place of the fourth and fifth members of
\cite[(190)]{B-CMP102b}, those carrying $D^{*}D$ and $\Delta$. Sums over $y'$ carry the measure
$\ell_{j'}^{d}$.
\end{lemma}

\begin{proof}
Throughout, $C$ denotes a constant of the kind admitted in
Theorem~\ref{4.5:thm-chart-existence}, not always the same. The letters $H$, $\mathfrak{D}$,
$\mathfrak{D}_3$, $\mathfrak{T}$, $\mathfrak{L}$ and $D$ stand for $H^{J}$, $\mathfrak{D}^{J}$,
$\mathfrak{D}^{J}_3$, $\mathfrak{T}^{J}$, $\mathfrak{L}^{J}$ and $D^{J}$. All of them are taken
at the current $J$, and the last five are evaluated at $A'$.

\emph{The second variation of $V^{J}$.} We differentiate the seven-term formula for
$\mathbb{V}'^{J}$ in \eqref{4.5:eq-higher-order-functional-a} along a direction $\mathfrak{A}$.
The dependence on $A'$ enters through four identities, each of which can be checked on the line.
First, $\partial_{\mathfrak{A}}D=\mathfrak{D}\mathfrak{A}$, since
$\mathfrak{D}=\delta D/\delta A'$. Second, $\partial_{\mathfrak{A}}(\mathfrak{L}A')=\mathfrak{A}-H\mathfrak{D}\mathfrak{A}=\mathfrak{T}\mathfrak{A}$,
since $\mathfrak{L}A'=A'-HD(A')$. Third,
$\partial_{\mathfrak{A}}\mathfrak{T}^{T}=-(\partial_{\mathfrak{A}}\mathfrak{D})^{T}H^{T}$,
since $\mathfrak{T}=I-H\mathfrak{D}$. Fourth, $\mathfrak{j}_J$ and $\mathfrak{k}_J$ are linear.

Apply the product rule to each of the seven terms:
\begin{itemize}
\item the first and second terms give one term each;
\item the third, fourth and fifth terms give two terms each;
\item the sixth term gives one term;
\item the seventh term gives two terms.
\end{itemize}
This yields the complete list of eleven terms:
\begin{equation}\label{4.5:eq-chart-variation-decay-c}
\begin{split}
\mathbb{V}''^{J}\mathfrak{A}={}&-(\partial_{\mathfrak{A}}\mathfrak{D}_3)^{T}H^{T}J
-Q^{T}\mathsf{k}\mathfrak{D}\mathfrak{A}
-(\partial_{\mathfrak{A}}\mathfrak{D})^{T}\mathsf{k}^{T}QA'
-\mathfrak{D}^{T}\mathsf{k}^{T}Q\mathfrak{A}\\
&+(\partial_{\mathfrak{A}}\mathfrak{D})^{T}\ell^{-1}\mathsf{k}D
+\mathfrak{D}^{T}\ell^{-1}\mathsf{k}\mathfrak{D}\mathfrak{A}
+\mathfrak{T}^{T}\nabla^{2}V_0(\mathfrak{L}A')[\mathfrak{T}\mathfrak{A}]\\
&-(\partial_{\mathfrak{A}}\mathfrak{D})^{T}H^{T}\nabla V_0(\mathfrak{L}A')
-\Lambda^{T}\mathfrak{j}_J(H\mathfrak{D}\mathfrak{A})\\
&-(\partial_{\mathfrak{A}}\mathfrak{D})^{T}H^{T}\mathfrak{k}_J(\Lambda A')
-\mathfrak{D}^{T}H^{T}\mathfrak{k}_J(\Lambda\mathfrak{A}).
\end{split}
\end{equation}
The first eight terms are $\nabla^{2}V^{J,\mathrm{disp}}\mathfrak{A}$. The last three are
$\nabla^{2}W^{J}\mathfrak{A}$.

\emph{The weighted bound.} Each term of \eqref{4.5:eq-chart-variation-decay-c} has exactly one
slot occupied by $\mathfrak{A}$. The weight $\rho$ is carried through that slot, and the
remaining slots are bounded without weight.

The weight is carried as follows:
\begin{itemize}
\item for the linear operators, by the algebra of exponentially decaying kernels
(Lemma~\ref{4.5:lem-kernel-algebra});
\item for $\partial_{\mathfrak{A}}\mathfrak{D}$, by part (d) of
Lemma~\ref{4.5:lem-nonlinear-correction}, which gives
$|(\partial_{\mathfrak{A}}\mathfrak{D})^{T}Y|_{(\alpha+1),\rho}\le C_{D2}N^{\rho}(\mathfrak{A})|Y|_{(\alpha)}$;
\item for $\nabla^{2}V_0$, $\mathfrak{j}_J$ and $\mathfrak{k}_J$, by locality, at the cost of a
factor $e^{\kappa c_{\mathrm{loc}}}$, where $c_{\mathrm{loc}}$ bounds the range of these local
operators.
\end{itemize}

The factors in the slots not occupied by $\mathfrak{A}$ obey the following unweighted bounds:
\begin{align*}
&|H^{T}J|_{(-4)}\le Ca_{\gamma},\qquad |\mathsf{k}^{T}QA'|_{(-4)}\le C\varepsilon_3,
\qquad |D|\le 4C_2\varepsilon_3^{2},\\
&|\nabla V_0|_{(-3)}\le C\varepsilon_3^{2},\qquad
|\mathfrak{k}_J(\Lambda A')|_{(-3)}\le Ca_{\gamma}\varepsilon_3 .
\end{align*}
In addition, $\mathfrak{D}$ and $\mathfrak{D}^{T}$ are $O(C_3\varepsilon_3)$, and
$\partial_{\mathfrak{A}}\mathfrak{D}_3=O(\varepsilon_3)$. Finally,
\begin{equation*}
|\nabla^{2}V_0(\mathcal{A})[f]|_{(-3),\rho}\le C\varepsilon_3\,N^{\rho}(f),
\end{equation*}
where $\mathcal{A}$ is the argument of $V_0$, here $\mathfrak{L}A'$. This follows by Cauchy's
estimate applied to \cite[(90)--(96)]{B-CMP102b}.

In the eleventh term, $\Lambda\mathfrak{A}$ is written as $G'\nabla^{*}\mathfrak{A}$ minus
$G'PD^{*}\mathfrak{A}$. The operator $G'\nabla^{*}$ is the third member of
\cite[(3.42)]{B-CMP99b}.

Summing the eleven bounds gives the first line of \eqref{4.5:eq-chart-variation-decay-b}. This
holds provided $c_{*}\delta_0$ is at most one half of the least spare rate, that is, the least
rate of decay left to the finitely many operators involved once the decay spent on the factors of
clause (K10) of Theorem~\ref{4.5:thm-chart-existence} has been paid. The fraction $c_{*}$ is
fixed in that clause.

\emph{The first variation.} The weighted Hessian bound yields
\begin{equation*}
\|\tilde G^{J}\mathbb{V}''^{J}(A_G)\|_{N^{\rho}\to N^{\rho}}\le \tfrac12 .
\end{equation*}
This is the analogue at free current of \cite[(187)]{B-CMP102b}. It imposes a ceiling on
$\varepsilon_3$, and hence on $c_9$.

It follows that the first equation of \eqref{4.5:eq-chart-variation-decay-a} is solved by a
Neumann series in $N^{\rho}$, as in \cite[(188)]{B-CMP102b}:
\begin{equation*}
\mathfrak{A}_0=\sum_{n\ge 0}\bigl(-\tilde G^{J}\mathbb{V}''^{J}(A_G)\bigr)^{n}\,
\tilde G^{J}\bigl[\Theta_J-\mathbb{V}''^{J}(A_G)\bigr]H_0v .
\end{equation*}
Moreover, $N^{\rho}(\mathfrak{A}_0)$ is at most twice
$N^{\rho}\bigl(\tilde G^{J}[\Theta_J-\mathbb{V}''^{J}(A_G)]H_0v\bigr)$. Every other operator in
\eqref{4.5:eq-chart-variation-decay-a} and \eqref{4.5:eq-chart-variation-decay-1} lies in some
class $\mathcal{K}_a$ or is local. Hence the weighted bounds pass to $\delta A_G$, $\delta X$,
$\delta\mathcal{H}^{J}[v]$, $\delta\Lambda^{J}[v]$ and $\delta\mathfrak{J}[v]$.

\emph{Reading off the kernels.} Take $S:=\{y'\}$ and $v:=\ell_{j'}^{-d}\delta_c$, where
$\delta_c$ is the direction supported on the single block bond $c$. Take $\kappa=c_{*}\delta_0$
and read off the weighted bounds at $x\in\Delta(y)$. Since $\rho(y)=e^{c_{*}\delta_0 d(y,y')}$,
this gives the remaining four bounds of \eqref{4.5:eq-chart-variation-decay-b}, with the powers of
$\ell_j$ carried by the respective norms.

\emph{The sizes.} Take $\kappa=0$, so that $\rho\equiv 1$. The same bounds then give
$N_{\beta}(\delta\mathcal{H}^{J}[v])\le C_9'|v|$ and
$|\delta\mathfrak{J}[v]|_{(-3)}\le C_{\mathfrak{J}}'|v|$, which are the sizes in (K4).
\end{proof}

In what follows, $R(X)$ is the adjoint action $R(X)Y=XYX^{-1}$ of an invertible matrix $X$ on
$\mathfrak{g}^{c}$, applied pointwise to fields, as in Definition~\ref{4.5:def-admissible-pairs};
it is not the ball radius $R$. The quantity $\|\operatorname{Ad}_u^{\pm1}\|$ is the one bounded
by $\mathfrak{K}$ in the definition of $\mathfrak{S}^{\mathrm{orb}}$ there. Gauge transformations
act on bond variables by $U^{u}_{b}=u(b_-)U_b u(b_+)^{-1}$, so that $(U^{u})^{u'}=U^{u'u}$.
The norms $N^{U}_{\beta}$ and $N^{U}$ are $N_{\beta}$ and $N$ computed with the transports of the
background $U$. The path $\Gamma_{x,x'}$ from $x$ to $x'$ is as in \cite[(3.40)]{B-CMP99b}, and
$U(\Gamma_{x,x'})$ is the transport along it. The $t_{\square}$ are the parameters of the
cut-offs $\zeta_{\square}$.

\begin{lemma}[Gauge covariance and reality of the chart on the orbit class]
\label{4.5:prf-chart-covariance}
This is clause (K5) of Theorem~\ref{4.5:thm-chart-existence}. The class $\mathfrak{S}$, the orbit
class $\mathfrak{S}^{\mathrm{orb}}$ with $\mathfrak{K}=4$, the radii and the constants
$C_f$, $\bar\varepsilon_4$, $c_9$ are those of Definition~\ref{4.5:def-admissible-pairs} and
Theorem~\ref{4.5:thm-chart-existence}. The maps $\mathcal{A}^{J}_{0}(B)$, $\mathcal{H}^{J}(B)$,
$\Lambda^{J}(B)$, $\mathfrak{J}(B)$ are those of the response chart constructed there.

For a $G^{c}$-valued gauge transformation $u$, let $\bar u := u\restriction(\text{block sites})$
be its restriction to the block sites, the end-points of the block bonds in $\mathfrak{B}$ on
which $B$ is defined. Consider the map
\begin{equation}\label{4.5:eq-chart-covariance-1}
  (U,J,B)\longmapsto \bigl(U^{u},\,R(u)J,\,R(\bar u)B\bigr).
\end{equation}
\begin{enumerate}
\item[\textup{(1)}] Under \eqref{4.5:eq-chart-covariance-1} the maps $\mathcal{H}^{J}$,
$\Lambda^{J}$, $\mathfrak{J}$ are covariant: their values at the transformed datum are $R(u)$
applied to their values at $(U,J,B)$. This holds for $G$-valued $u$. It also holds for every
$G^{c}$-valued $u$ with $\|\operatorname{Ad}_u^{\pm1}\|\le\mathfrak{K}'$, where
$\mathfrak{K}'\le\mathfrak{K}^{2}=16$, as an algebraic identity wherever both sides are defined.
\item[\textup{(2)}] Consequently the chart is defined, single-valued and analytic on
$\mathfrak{S}^{\mathrm{orb}}$. It satisfies clauses (K1)--(K4) and (K8) of
Theorem~\ref{4.5:thm-chart-existence} with their constants multiplied by a factor at most
$\mathfrak{K}^{2}=16$. This holds under the ceiling
$\mathfrak{K}^{2}\cdot 4C_f c_9\le\bar\varepsilon_4$ and with the orbit radius $c_9/\mathfrak{K}$
for $B$.
\item[\textup{(3)}] For $f$ on bonds,
$N^{U^{u}}_{\beta}(R(u)f)\le\|\operatorname{Ad}_u\|\,N^{U}_{\beta}(f)$. Indeed, the difference
entering the norm of \cite[(3.40)]{B-CMP99b} obeys
\begin{equation}\label{4.5:eq-chart-covariance-2}
\begin{split}
  &R\bigl(U^{u}(\Gamma_{x,x'})\bigr)(R(u)f)(x')-(R(u)f)(x)\\
  &\qquad = R(u(x))\bigl[R\bigl(U(\Gamma_{x,x'})\bigr)f(x')-f(x)\bigr].
\end{split}
\end{equation}
\item[\textup{(4)}] At $G$-valued $U$ and real $J$, $B$ and $t_{\square}$, everything is real.
\end{enumerate}
\end{lemma}

\begin{proof}
\emph{Covariance of every constituent.} Every constituent of the chart is built from transports,
traces and the averaging maps of \cite[(3.28)--(3.34)]{B-CMP99b}. The averaging maps are
covariant in the sense of \cite[p.~27]{B-CMP98}:
\begin{quotation}
if we apply a gauge transformation $v$ to a configuration $V$,
$V^{v}_{b}=v(b_-)V_b v^{-1}(b_+)$, $b\subset\Omega'$, then
$(\overline{V^{v}})_c=v(c_-)\overline{V}_c v^{-1}(c_+)$
\end{quotation}
These conjugation laws are algebraic identities. By \cite[(56)--(57)]{B-CMP98}, they remain valid
for invertible matrices wherever one side is defined: ``$R(X)f(Y)=f(R(X)Y)$ for analytic
functions $f$''. The operators $c^{(1)}$, $c^{(2)}$ entering the insertions $\mathfrak{j}_J$,
$\mathfrak{k}_J$, $\mathfrak{m}_J$ are pointwise commutators with transports. They therefore
transform by $R(u)$ under \eqref{4.5:eq-chart-covariance-1}. Finally, the fixed points defining
the chart are unique. This proves (1), for $G$-valued $u$ and for $G^{c}$-valued $u$ as stated.

\emph{The two representatives.} Let a point of $\mathfrak{S}^{\mathrm{orb}}$ have two
representatives $(U_1,J_1,B_1)$ and $(U_2,J_2,B_2)$, with $(U_i,J_i)\in\mathfrak{S}$. Thus the
point is $((U_i,J_i)^{u_i},B^{\mathrm{orb}})$ with
$\|\operatorname{Ad}_{u_i}^{\pm1}\|\le\mathfrak{K}$ and $B_i=R(\bar u_i^{-1})B^{\mathrm{orb}}$,
$i=1,2$. Put $u_{21}:=u_2^{-1}u_1$. Since $U_1^{u_1}=U_2^{u_2}$ and $(U^{u})^{u'}=U^{u'u}$, we
have $(U_2,J_2,B_2)=(U_1,J_1,B_1)^{u_{21}}$ in the sense of \eqref{4.5:eq-chart-covariance-1},
and
\[
  \|\operatorname{Ad}_{u_{21}}^{\pm1}\|
  \le\|\operatorname{Ad}_{u_2}^{\mp1}\|\,\|\operatorname{Ad}_{u_1}^{\pm1}\|\le\mathfrak{K}^{2}.
\]
By the orbit radius, $|B_i|\le\mathfrak{K}|B^{\mathrm{orb}}|<c_9$. Hence both charts exist, by
clause (K1) of Theorem~\ref{4.5:thm-chart-existence}. By the algebraic covariance of every
constituent of the chart, $R(u_{21})\mathcal{A}^{J_1}_{0}(B_1)$ solves the fixed-point equation
at $(U_2,J_2,B_2)$. The norm $N$ transforms as $N_{\beta}$ does in (3); its transport
differences obey \eqref{4.5:eq-chart-covariance-2}. We also use the bound
$N^{U_1}(\mathcal{A}^{J_1}_{0}(B_1))\le C_f|B_1|$ of Theorem~\ref{4.5:thm-chart-existence}.
Together these give
\[
  N^{U_2}\bigl(R(u_{21})\mathcal{A}^{J_1}_{0}(B_1)\bigr)
  \le\|\operatorname{Ad}_{u_{21}}\|\,N^{U_1}\bigl(\mathcal{A}^{J_1}_{0}(B_1)\bigr)
  \le\mathfrak{K}^{2}C_f|B_1|<\mathfrak{K}^{2}C_f c_9\le\bar\varepsilon_4,
\]
the last inequality by the ceiling $\mathfrak{K}^{2}\cdot4C_f c_9\le\bar\varepsilon_4$. The
fixed point at $(U_2,J_2,B_2)$ is unique in $\{N\le\bar\varepsilon_4\}$ by clause (K1). Hence
\[
  R(u_{21})\mathcal{A}^{J_1}_{0}(B_1)=\mathcal{A}^{J_2}_{0}(B_2).
\]
Every quantity built from the fixed point therefore agrees at the two representatives, after
application of $R(u_{21})$.

The chart is extended to $\mathfrak{S}^{\mathrm{orb}}$ by covariance: at
$((U,J)^{u},B^{\mathrm{orb}})$ its value is $R(u)$ applied to the chart at
$(U,J,R(\bar u^{-1})B^{\mathrm{orb}})$. By the preceding paragraph this extension is
single-valued. On each open patch $\{(U,J)^{u}:(U,J)\in\mathfrak{S}\}$ with fixed $u$, the
extension is $R(u)$ applied to the chart at the representative, which is holomorphic by clause
(K1). Hence it is holomorphic on each patch, and therefore on their union
$\mathfrak{S}^{\mathrm{orb}}$, since holomorphy is a local property. In
\cite[p.~307]{B-CMP102b} the corresponding extension is made by ``treating the equality as a
definition for the remaining $v$''. On $\mathfrak{S}^{\mathrm{orb}}$ the argument above proves
that this definition is consistent. This proves the single-valuedness and analyticity asserted
in (2).

\emph{Derivatives and differences.} The covariant derivatives satisfy
$\nabla_{U^{u}}R(u)=R(u)\nabla_{U}$; no derivative of $u$ enters. For
\eqref{4.5:eq-chart-covariance-2}, note that
$U^{u}(\Gamma_{x,x'})=u(x)U(\Gamma_{x,x'})u(x')^{-1}$ and that $R$ is multiplicative. Hence
\[
  R\bigl(U^{u}(\Gamma_{x,x'})\bigr)R(u(x'))f(x')=R(u(x))R\bigl(U(\Gamma_{x,x'})\bigr)f(x').
\]
Subtracting $(R(u)f)(x)=R(u(x))f(x)$ gives \eqref{4.5:eq-chart-covariance-2}. Taking norms, with
$\|R(u(x))\|\le\|\operatorname{Ad}_u\|$, gives the inequality of (3).

\emph{Reality.} At $G$-valued $U$ and real $J$, $B$, $t_{\square}$, the iterates of the
fixed-point map are real. Hence the fixed point is real, and so is every quantity built from it.
This proves (4).
\end{proof}

\begin{lemma}[Uniformity of the chart constants]\label{4.5:prf-chart-constants}
This is clause (K7) of Theorem~\ref{4.5:thm-chart-existence}. Let $\mathfrak{S}$ be the class of
admissible pairs $(U,J)$ of that theorem.
\begin{enumerate}
\item Every constant of Theorem~\ref{4.5:thm-chart-existence} is a constant depending only on $d$
and $L$, multiplied by an absolute combinatorial factor (see \cite{B-CMP99b} and
\cite{B-CMP102b}). Each such constant is attached to the objects of a single renormalisation
step, and no product over levels is formed. The level $k$, the number of steps $K$, the region
$X$ and the sequence of domains $\{\Omega_j\}$ of Theorem~\ref{4.5:thm-chart-existence} are read
nowhere.
\item The blocking factor $L$ enters only through fixed powers, by the lemma on the kernel
classes $\mathcal{K}_a$. The constants depend on $M_1$ and on the parameter $R_1$ of the
construction of the response chart only through the fixed lower bounds (floors) imposed on
these two parameters, not through their values.
\item The following are never formed at a general pair $(U,J)\in\mathfrak{S}$:
\begin{itemize}
\item the covariant Laplacian $\Delta_U$ and the operators $D^{*}_{U}D_U$ and $D_U D^{*}_{U}$
built from the covariant derivative $D_U$, applied to the image under $H^{J}$ of the response
variable of the chart (the quantity whose first variation is $\mathfrak{D}^{J}(A_G)\delta A_G$ in
clause (K4) of Theorem~\ref{4.5:thm-chart-existence}, proved in
Lemma~\ref{4.5:prf-chart-variation-decay}) or to $\mathcal{H}^{J}$, to their variations, or to
their $s$-pieces, that is, the pieces of these objects labelled by $s$ in the construction of
the response chart;
\item the adjoint $D^{*}$ acting on the fine index, that is, on the fine-lattice argument of the
kernels, of $\mathfrak{D}^{J}$, of the operator $C_j$ of the construction of the response chart,
or of $Q$;
\item $D^{*}_{U}J$, and any H\"older norm of $J$;
\item Ba{\l}aban's smoothed supremum-norm estimates \cite[(3.44), (3.45)]{B-CMP99b},
applied to unsmoothed supremum data, or taken with the smoothing parameter
$\varepsilon \to 0$ (this exclusion rests on the smoothing lemma of this chapter);
\item the identities \cite[(3.124), (3.152)]{B-CMP99b}, the identity $RD^{*}H = 0$, in which $R$
and $D^{*}$ are the operators occurring in $\Lambda = G'RD^{*}_{U}$ and $H$ is Ba{\l}aban's
minimiser map, and the identities \cite[(170), (171), (178)]{B-CMP102b};
\item the identity \cite[(87)]{B-CMP102b}, as a step of a derivation.
\end{itemize}
\item The objects and identities listed in the preceding item are formed and used only at the
diagonal pairs, at which they are the objects and identities of Ba{\l}aban's papers.
\end{enumerate}
\end{lemma}

\begin{proof}
Each of the four items is a statement about which objects and which constants occur in two
places: the construction of the response chart, and the estimates on which
Theorem~\ref{4.5:thm-chart-existence} rests. Among these estimates are the first-variation
formulae and the decay bounds of clauses (K4) and (K8) of
Theorem~\ref{4.5:thm-chart-existence}, proved in Lemma~\ref{4.5:prf-chart-variation-decay}.

The items are established by inspecting that construction and those estimates. For each
constant introduced there we list the quantities on which it depends. The dependence is only
on $d$ and $L$, through fixed powers of $L$ supplied by the lemma on the kernel classes
$\mathcal{K}_a$, and on $M_1$ and $R_1$ only through the floors imposed on them. No constant
depends on $k$, $K$, $X$ or $\{\Omega_j\}$, and none arises as a product over levels. This gives
the first two items.

For each operator, each estimate and each identity that appears we list the configurations at
which it is formed or applied. The objects and identities of the third item occur only on the
diagonal, where they are the objects and identities of Ba{\l}aban's papers. None of them occurs
at a general pair $(U,J) \in \mathfrak{S}$. This gives the last two items.
\end{proof}

\begin{lemma}[Lipschitz dependence on the background in complex directions]
\label{4.5:prf-complex-lipschitz}
This is clause (K9) of Theorem~\ref{4.5:thm-chart-existence}. Let $(U,J)$ satisfy the conditions
(S1)--(S3) of \eqref{4.5:eq-admissible-pairs-a} with the radii $a_0,a_1,a_\gamma$ replaced by
$\tfrac12 a_0,\tfrac12 a_1,\tfrac12 a_\gamma$. Let $\mathsf{a}$ and $\mathsf{j}$ be $\mathfrak{g}^{c}$-valued
bond functions (the directions of the background twist and of the current) with
$\max\{N(\mathsf{a}),|\mathsf{j}|_{(-3)}\}\le 1$, and put
\begin{equation*}
\alpha^{\mathrm{ch}}:=\tfrac14\min\{a_1,a_\gamma\}.
\end{equation*}
For $\sigma\in\mathbb{C}$ with $|\sigma|<\alpha^{\mathrm{ch}}$ set $U_\sigma:=e^{i\sigma\eta\mathsf{a}}U$ and
$J_\sigma:=J+\sigma\mathsf{j}$, and write $\mathcal{H}^{(\sigma)}(B)$ and $\mathfrak{J}^{(\sigma)}(B)$ for the
background potential $\mathcal{H}^{J_\sigma}(B)$ and the current correction $\mathfrak{J}(B)$ of the
response chart at the base $(U_\sigma,J_\sigma)$. Then $(U_\sigma,J_\sigma)\in\mathfrak{S}$, and for every
$\mathfrak{g}^{c}$-valued function $B$ on $\mathfrak{B}$ with $|B|=\sup|B|<c_9$
\begin{equation}\label{4.5:eq-complex-lipschitz-1}
\begin{split}
&\max\Bigl\{N_U\bigl(\mathcal{H}^{(\sigma)}(B)-\mathcal{H}^{(0)}(B)\bigr),\,
|\mathfrak{J}^{(\sigma)}(B)-\mathfrak{J}^{(0)}(B)|_{(-3)}\Bigr\}
\le C^{\boxplus,K}_{2}\,|B|\,|\sigma|,\\
&\qquad
C^{\boxplus,K}_{2}:=\frac{2(1+2\alpha^{\mathrm{ch}})\max\{C^{K}_{9},C_{\mathfrak{J}}\}}{\alpha^{\mathrm{ch}}}.
\end{split}
\end{equation}
\end{lemma}

\begin{proof}
\emph{Membership.} Since $J$ and $\mathsf{j}$ are $\mathfrak{g}^{c}$-valued bond functions, so is
$J_\sigma$. Since $|J|_{(-3)}<\tfrac12 a_\gamma$ by (S3) at half radius, $|\mathsf{j}|_{(-3)}\le1$ and
$|\sigma|<\alpha^{\mathrm{ch}}\le\tfrac14 a_\gamma$, the triangle inequality gives
\begin{equation*}
|J+\sigma\mathsf{j}|_{(-3)}<\tfrac12 a_\gamma+\tfrac14 a_\gamma<a_\gamma,
\end{equation*}
which is (S3) for $J_\sigma$. Write $U=e^{i\eta\mathfrak{a}}U_r$ as in (S1). By the
Baker--Campbell--Hausdorff formula the product $e^{i\sigma\eta\mathsf{a}}e^{i\eta\mathfrak{a}}$ is a single
exponential $e^{i\eta\mathfrak{a}_\sigma}$, whose twist $\mathfrak{a}_\sigma$ is $\mathfrak{a}+\sigma\mathsf{a}$
plus a commutator remainder. Hence $U_\sigma=e^{i\eta\mathfrak{a}_\sigma}U_r$ with the same $G$-valued part
$U_r$ as $U$; since the regularity required in (S1) is a condition on $U_r$ alone, (S1) holds for
$U_\sigma$, and it remains to bound $N_{U_r}(\mathfrak{a}_\sigma)$. The first two terms obey
$N_{U_r}(\mathfrak{a})<\tfrac12 a_1$, by (S2) at half radius, and, with $N(\mathsf{a})$ read in the frame
$N_{U_r}$ of (S2), $N_{U_r}(\sigma\mathsf{a})=|\sigma|N(\mathsf{a})<\alpha^{\mathrm{ch}}\le\tfrac14 a_1$, so that
\begin{equation*}
N_{U_r}(\mathfrak{a})+|\sigma|N(\mathsf{a})<\tfrac12 a_1+\tfrac14 a_1 ,
\end{equation*}
and the commutator remainder has $N_{U_r}$-norm less than the remaining $\tfrac14 a_1$ by the estimate
of the Baker--Campbell--Hausdorff remainder for products of exponentials near the identity, applied at
the radii $\tfrac12 a_1$ and $\tfrac14 a_1$. This gives
$N_{U_r}(\mathfrak{a}_\sigma)<a_1$, which is (S2) for $U_\sigma$. Hence $(U_\sigma,J_\sigma)\in\mathfrak{S}$.

\emph{Change of frame.} For a bond function $f$, using $|\sigma|<\alpha^{\mathrm{ch}}$,
$\ell|\mathsf{a}|\le N(\mathsf{a})\le1$ and $\ell|f|\le N(f)$, we have
\begin{equation}\label{4.5:eq-complex-lipschitz-2}
\ell^{2}\,|\nabla_{U_\sigma}f-\nabla_U f|
\le 2|\sigma|\,\ell|\mathsf{a}|\cdot\ell|f|
\le 2\alpha^{\mathrm{ch}}N(f).
\end{equation}
Thus a bound in the frame of $U_\sigma$ becomes, after multiplication by $1+2\alpha^{\mathrm{ch}}$, a bound
in the frame $N_U$ of the base $U$; this is the origin of that factor in
\eqref{4.5:eq-complex-lipschitz-1}.

\emph{Holomorphy, size and zeros.} Let $F(B,\sigma)$ be a scalar component of
$\mathcal{H}^{(\sigma)}(B)-\mathcal{H}^{(0)}(B)$, or of
$\ell^{3}\bigl(\mathfrak{J}^{(\sigma)}(B)-\mathfrak{J}^{(0)}(B)\bigr)$, or of
$\ell^{2}\nabla_U\bigl(\mathcal{H}^{(\sigma)}(B)-\mathcal{H}^{(0)}(B)\bigr)$, the last after the change
of frame \eqref{4.5:eq-complex-lipschitz-2}. By clause (K1) of Theorem~\ref{4.5:thm-chart-existence},
$F$ is holomorphic on the bidisc
\begin{equation*}
\{|B|<c_9\}\times\{|\sigma|<\alpha^{\mathrm{ch}}\}.
\end{equation*}
The linear bounds on the background potential and the current correction,
Lemma~\ref{4.5:prf-chart-linear-bounds}, applied at the bases
$(U_\sigma,J_\sigma)$ and $(U,J)$ of $\mathfrak{S}$, together with the factor $1+2\alpha^{\mathrm{ch}}$ of the
change of frame and the triangle inequality, give
\begin{equation*}
\sup|F|\le 2(1+2\alpha^{\mathrm{ch}})\max\{C^{K}_{9},C_{\mathfrak{J}}\}\,c_9
=C^{\boxplus,K}_{2}\,\alpha^{\mathrm{ch}}c_9 .
\end{equation*}
$F$ vanishes at $\sigma=0$, the two terms then being equal; and it vanishes at $B=0$ by the vanishing
at zero datum, Lemma~\ref{4.5:prf-chart-structure}, which holds at every base, in particular at
$(U_\sigma,J_\sigma)$ and at $(U,J)$.

\emph{Schwarz in $\sigma$.} For fixed $B$, $\sigma\mapsto F(B,\sigma)/\sigma$ is holomorphic on
$\{|\sigma|<\alpha^{\mathrm{ch}}\}$, since $F(B,0)=0$. On each circle $|\sigma|=r$, $0<r<\alpha^{\mathrm{ch}}$,
it is bounded by $C^{\boxplus,K}_{2}\alpha^{\mathrm{ch}}c_9/r$; by the maximum principle, letting
$r\uparrow\alpha^{\mathrm{ch}}$,
$|F(B,\sigma)/\sigma|\le C^{\boxplus,K}_{2}c_9$ on the whole disc.

\emph{Schwarz in $B$.} For fixed $\sigma\ne0$, the function $B\mapsto F(B,\sigma)/\sigma$ is holomorphic
on $\{|B|<c_9\}$, bounded by $C^{\boxplus,K}_{2}c_9$ and vanishes at $B=0$. Fix $B\ne0$ with $|B|<c_9$,
and consider along the ray through $B$ the function $\varphi(\cdot):=F(\cdot\,B,\sigma)/\sigma$ of one
complex variable on the disc of radius $c_9/|B|$ about the origin, which contains the point $1$. It is
holomorphic, $\varphi(0)=0$ and $|\varphi|\le C^{\boxplus,K}_{2}c_9$, so Schwarz's lemma on this disc gives
$|\varphi(1)|\le C^{\boxplus,K}_{2}c_9\cdot|B|/c_9$, that is,
\begin{equation*}
|F(B,\sigma)|\le C^{\boxplus,K}_{2}\,|B|\,|\sigma|.
\end{equation*}
The same holds trivially for $B=0$ or $\sigma=0$. Taking the supremum over the components gives
\eqref{4.5:eq-complex-lipschitz-1}.
\end{proof}

\begin{proposition}[Walk expansions of the chart and localisation in the decoupling parameters]
\label{4.5:prf-chart-walk-localisation}
This is clause (K10) of Theorem~\ref{4.5:thm-chart-existence}. The pairs $(U,J)\in\mathfrak{S}$,
the datum $B$ and the constants are those of that theorem.

\smallskip\noindent\emph{Hypothesis.} The constrained covariance inverse
$\mathsf{q}_0=(QG_0Q^{*})^{-1}$ has the bounds and the generalised random-walk expansion of
Remark~\ref{4.5:rem-covariance-inverse-walks}, printed without proof in
\cite[after Theorem~3.9]{B-CMP99b}; cited as printed.

\smallskip\noindent\emph{Conclusion.} Let $\sigma_0$ be the family of decoupling cubes of
\cite{B-CMP116}, and let $s=(s(\Delta))_{\Delta\in\sigma_0}$ be complex decoupling parameters.
Consider the primitive operators
\begin{equation}\label{4.5:eq-chart-walk-localisation-1}
G_0,\quad G',\quad R,\quad P,\quad (Q'G'^{2}Q'^{*})^{-1},\quad \mathsf{q}_0,\quad
G^{J}_{\pi},\quad \mathsf{q}_J,\quad \tilde G^{J},\quad G^{J}_{1},\quad (QG^{J}_{1}Q^{*})^{-1}.
\end{equation}
Each of them is expanded into its generalised random-walk expansion; for $G^{J}_{\pi}$,
$\mathsf{q}_J$, $\tilde G^{J}$, $G^{J}_{1}$ the defining series are expanded into concatenated
walks with their side branches, as in the expansion clause of
\eqref{4.5:eq-free-current-operators-a}. These
series are not re-summed at $s$, so no condition couples $a_\gamma$ to $\kappa_1$. Every walk
term is multiplied by the factor
\begin{equation}\label{4.5:eq-chart-walk-localisation-2}
\prod_{\Delta}s(\Delta),
\end{equation}
the product running over the cubes $\Delta\in\sigma_0$ met by any of the localisation domains
of the term, the domains of its side branches included. This is the rule by which the terms of
\cite[(1.6)]{B-CMP116} receive their factors, stated in \cite[after (1.7)]{B-CMP116} as follows:
\begin{quote}
for a random walk $\omega$ localized in
$\tilde X_0^{5}\cup\tilde X_1^{5}\cup\dots\cup\tilde X_n^{5}$ we take the
$\{\Delta_1,\dots,\Delta_m\}$ of all cubes from $\sigma_0$ which intersect this localization
domain, and we multiply the term in (1.6) corresponding to $\omega$ by
$s(\Delta_1)\cdots s(\Delta_m)$
\end{quote}
Here the localisation domain of a walk $\omega$ is the union of the domains $\tilde X_i^{5}$
displayed in the quotation, which \cite{B-CMP116} attaches to $\omega$; a cube of $\sigma_0$ is
met by a walk term if it intersects the localisation domain of one of its pieces or side
branches. The operators so obtained are the primitive operators at $s$.

At these operators define the composites $H_0$, $P_0^{T}$, $H^{J}$, $H^{J}_{1}$, $\Lambda$,
$\Lambda^{T}$, $(I+\mathfrak{R}^{J})^{-1}$, and then $D^{J}$, $\mathbb{V}'^{J}$,
$\mathcal{A}^{J}_{0}$, $\mathcal{H}^{J}$, $\Lambda^{J}$, $\mathfrak{J}$, by the formulas of
Lemma~\ref{4.5:lem-free-current-operators} and Definition~\ref{4.5:def-response-chart}. In
particular $\mathcal{A}^{J}_{0}(s)(B)$ is taken in the form
\begin{equation}\label{4.5:eq-chart-walk-localisation-3}
\mathcal{A}=\tilde G^{J}\bigl[\Theta_JH_0B-\mathbb{V}'^{J}(\mathcal{A}+H_0B)\bigr],
\end{equation}
whose only term linear in any variable is the source term $\tilde G^{J}\Theta_JH_0B$.
Suppose that
\begin{equation*}
|s(\Delta)|\le e^{\kappa_1}\quad(\Delta\in\sigma_0),\qquad \delta_1M\ge\kappa_1,
\end{equation*}
where $\delta_1$ depends on $\delta_0$ only and $M$ is the side of the cubes of $\sigma_0$ in
\cite{B-CMP116}. Then the following hold.
\begin{enumerate}
\item Each primitive operator \eqref{4.5:eq-chart-walk-localisation-1} at $s$ is analytic in
$s$ and is bounded by $B_0e^{16\kappa_1}$ in the norms of the list in
\eqref{4.5:eq-free-current-operators-a}. Each composite is bounded by $B_0e^{c_{\mathbb{K}}\kappa_1}$.
\item Clauses (K1), (K3)--(K5) and (K8) of Theorem~\ref{4.5:thm-chart-existence}, together with
$\mathcal{H}^{J}(s)(0)=\mathfrak{J}(s)(0)=0$, hold with $B_0$ replaced by
$B_0e^{c_{\mathbb{K}}\kappa_1}$, under the ceiling $e^{2c_{\mathbb{K}}\kappa_1}|B|\le c_9$. In
particular
\begin{equation*}
N_{\beta}\bigl(\mathcal{H}^{J}(s)(B)\bigr),\ |\mathfrak{J}(s)(B)|_{(-3)}\le
Ce^{c_{\mathbb{K}}\kappa_1}|B|.
\end{equation*}
\item Let $\sigma\subset\sigma_0$ and suppose $s(\Delta)=0$ for $\Delta\notin\sigma$. Then all
kernels vanish across the components of $Y(\sigma)$, and the equations separate.
\item At $s\equiv1$ all the objects above coincide with those of
Lemma~\ref{4.5:lem-free-current-operators} and Definition~\ref{4.5:def-response-chart}.
\end{enumerate}
The identity \eqref{4.5:eq-chart-structure-a} and the exactness statements
\eqref{4.5:eq-diagonal-exactness-a} are asserted at $s\equiv1$ only.
\end{proposition}

\begin{proof}
\emph{Concatenated walks.} A concatenated walk is written
\begin{equation*}
\varpi=(\omega_0;\iota_1,\omega_1;\dots;\iota_{\mathsf{n}},\omega_{\mathsf{n}}).
\end{equation*}
Each $\omega_i$ is a walk of $G_0$, of $G'$ or of $\mathsf{q}_0$; the $\mathsf{q}_0$-walks are those
supplied by the hypothesis, Remark~\ref{4.5:rem-covariance-inverse-walks}, printed without proof
in \cite[after Theorem~3.9]{B-CMP99b}; cited as printed. The walk $\omega_i$ has anchored
localisation cubes
$\Delta(y_i)$ and $\Delta(y_i')$, the cubes containing the points $y_i$ and $y_i'$. Each
$\iota_i$ is a local insertion of range at most a fixed constant $c_\iota$. The local parts of
the insertions, namely $\mathfrak{j}_J$, $\mathfrak{k}_J$ and two further local insertions, are
pointwise in $J$. The non-local parts $\Lambda$, $RG'$ and $H^{J,T}J$ of an
insertion are not counted as insertions: they are themselves walk-expanded pieces, either of the
chain or of a side branch. Consequently
$\operatorname{dist}(y_i',y_{i+1})\le c_\iota$ for each $i$.

\smallskip\noindent\emph{The two spare exponentials.} By \cite[(1.7)]{B-CMP116}, each piece
$\omega_i$ carries two kinds of factor. The first are factors, summable with a bound by $B_0$,
that control the sum over walks. The second are two spare exponentials
$e^{-\delta_0|\omega_i|_{y_i,y_i'}}$, which are used to produce exponential bounds for localised
terms. Here $|\omega|_{y,y'}$ denotes the anchored length of a walk $\omega$ with anchors $y$,
$y'$.

For the whole concatenated walk $\varpi$ with end anchors $y=y_0$ and $y'=y'_{\mathsf{n}}$, let
$|\varpi|_{y,y'}$ be the length of a shortest tree that passes through all cubes met by $\varpi$,
its side branches included, and is anchored at $y$ and $y'$. The anchored lengths chain:
\begin{equation}\label{4.5:eq-chart-walk-localisation-4}
\sum_{i=0}^{\mathsf{n}}|\omega_i|_{y_i,y_i'}+c_\iota\mathsf{n}\ \ge\ |\varpi|_{y,y'}.
\end{equation}
Side branches are handled in the same way: their anchored lengths, each anchored at its
insertion point, are added to the left-hand side.

\smallskip\noindent\emph{The walk bound.} Let $\mathsf{m}(\varpi)$ be the number of distinct cubes
of $\sigma_0$ met by $\varpi$, branches included. Since $|s(\Delta)|\le e^{\kappa_1}$, the
factor \eqref{4.5:eq-chart-walk-localisation-2} of the term of $\varpi$ has modulus at most
$e^{\kappa_1\mathsf{m}(\varpi)}$. Two cases arise.
\begin{itemize}
\item If $\mathsf{m}(\varpi)\le2^4$, then $e^{\kappa_1\mathsf{m}(\varpi)}\le e^{16\kappa_1}$.
\item If $\mathsf{m}(\varpi)>2^4$, then the counting argument following \cite[(1.11)]{B-CMP116}
applies
to $\varpi$: a tree meeting $\mathsf{m}(\varpi)>2^4$ distinct cubes of $\sigma_0$ has
$\delta_0|\varpi|_{y,y'}\ge\delta_1\mathsf{m}(\varpi)M$, and $\delta_1M\ge\kappa_1$ turns this into
$\delta_0|\varpi|_{y,y'}\ge\kappa_1\mathsf{m}(\varpi)$.
\end{itemize}
In the second case, combining with \eqref{4.5:eq-chart-walk-localisation-4},
\begin{equation*}
e^{\kappa_1\mathsf{m}(\varpi)}\le e^{\delta_0|\varpi|_{y,y'}}\le
\prod_{i=0}^{\mathsf{n}}e^{\delta_0|\omega_i|_{y_i,y_i'}}\cdot e^{\delta_0c_\iota\mathsf{n}}.
\end{equation*}
Since every factor on the right of the following display is at least one, both cases give
\begin{equation}\label{4.5:eq-chart-walk-localisation-5}
\Bigl|\prod_{\Delta}s(\Delta)\Bigr|\le e^{\kappa_1\mathsf{m}(\varpi)}\le e^{16\kappa_1}
\prod_{i=0}^{\mathsf{n}}e^{\delta_0|\omega_i|_{y_i,y_i'}}\cdot e^{\delta_0c_\iota\mathsf{n}},
\end{equation}
with the branch pieces included in the product.

The two spare exponentials of each piece are used as follows.
\begin{itemize}
\item One of them pays the factor $e^{\delta_0|\omega_i|_{y_i,y_i'}}$ in
\eqref{4.5:eq-chart-walk-localisation-5}.
\item The other keeps the decay of the composite at the rate from which $c_{*}\delta_0$ is taken.
\item The remaining factor $e^{\delta_0c_\iota}$, one per insertion, is absorbed into the constants
$C_J$ and $C_\Theta$ of the insertion series of Lemma~\ref{4.5:lem-free-current-operators}. These
constants depend on $d$ and $L$ only; they are independent of $\kappa_1$ and of $M$.
\end{itemize}
The insertion series therefore converge under
\begin{equation*}
C_Ja_\gamma\le\tfrac12,\qquad C_\Theta a_\gamma B_0'\le\tfrac12,
\end{equation*}
and these conditions do not involve $\kappa_1$. This is the mechanism of
\cite[(1.15)]{B-CMP116}: smallness comes from the ceiling on the field for the non-linear part,
and none is needed for the source term.

Each walk term is the product of its term at $s\equiv1$ with a monomial in the $s(\Delta)$. By
\eqref{4.5:eq-chart-walk-localisation-5} and the summability just described, the series converge
uniformly on the polydisc $|s(\Delta)|\le e^{\kappa_1}$. This gives analyticity in $s$ and the
bound $B_0e^{16\kappa_1}$ for each primitive operator \eqref{4.5:eq-chart-walk-localisation-1} in
the norms of the list in \eqref{4.5:eq-free-current-operators-a}, as in \cite[(1.11)]{B-CMP116}.
The composites are formed from the primitive operators at $s$ by the formulas of
Lemma~\ref{4.5:lem-free-current-operators} and Definition~\ref{4.5:def-response-chart}, and those
formulas read of the primitive operators only their bounds in the norms of the list in
\eqref{4.5:eq-free-current-operators-a} and their locality. By the bound just proved for the
primitive operators, each composite is therefore bounded by $B_0e^{c_{\mathbb{K}}\kappa_1}$, with
the absolute integer $c_{\mathbb{K}}$ of Theorem~\ref{4.5:thm-chart-existence}. This gives the
first item.

\smallskip\noindent\emph{The clauses at $s$.} The remaining assertions rest on the expansion
clause of \eqref{4.5:eq-free-current-operators-a}. The proofs of clauses (K1),
(K3)--(K5) and (K8), in Theorem~\ref{4.5:thm-chart-existence},
Proposition~\ref{4.5:prf-chart-linear-bounds}, Proposition~\ref{4.5:prf-chart-variation-decay}
and Proposition~\ref{4.5:prf-chart-covariance}, use two properties of the operators only: their
bounds in the norms of the list in \eqref{4.5:eq-free-current-operators-a}, and their locality.
By the
preceding paragraphs both properties hold at $s$, with $B_0$ replaced by
$B_0e^{c_{\mathbb{K}}\kappa_1}$. Those proofs therefore apply verbatim at $s$, under the ceiling
$e^{2c_{\mathbb{K}}\kappa_1}|B|\le c_9$. This gives the second item, including the bounds on
$N_{\beta}(\mathcal{H}^{J}(s)(B))$ and $|\mathfrak{J}(s)(B)|_{(-3)}$.

The vanishing $\mathcal{H}^{J}(s)(0)=\mathfrak{J}(s)(0)=0$ holds for the same reason as at
$s\equiv1$ in Proposition~\ref{4.5:prf-chart-structure}. In
\eqref{4.5:eq-chart-walk-localisation-3} the only term linear in any variable is the source term
$\tilde G^{J}\Theta_JH_0B$, and $\mathbb{V}'^{J}$ vanishes at zero. Hence $\mathcal{A}=0$ is the
fixed point at $B=0$ for every $s$. Then $A_G$, $X$, $\mathcal{H}^{J}(s)(0)$,
$\Lambda^{J}(s)(0)$ and $\mathfrak{J}(s)(0)$ vanish in turn by the formulas, exactly as in
Proposition~\ref{4.5:prf-chart-structure}.

The explicit formulas, as opposed to the bounds and the locality, enter only into
\eqref{4.5:eq-chart-structure-a}, \eqref{4.5:eq-response-chart-a} and
\eqref{4.5:eq-diagonal-exactness-a}. This is why the first and the last of these are asserted at
$s\equiv1$ only.

\smallskip\noindent\emph{Separation, and the value at $s\equiv1$.} If $s(\Delta)=0$ for
$\Delta\notin\sigma$, the vanishing of all kernels across the components of $Y(\sigma)$, and the
resulting separation of the equations, follow by the argument of
\cite[between (1.7) and (1.11)]{B-CMP116},
applied kernel
by kernel to the walk terms defined by \eqref{4.5:eq-chart-walk-localisation-2}. At $s\equiv1$
every factor \eqref{4.5:eq-chart-walk-localisation-2} equals one. The expanded series then sum to
the operators of Lemma~\ref{4.5:lem-free-current-operators} and
Definition~\ref{4.5:def-response-chart}, which gives the fourth item.

The argument of \cite[(1.12)--(1.18)]{B-CMP116} is the same argument for the triple
$(H,\tilde G,H_0)$. Here the list of operators is longer. The list of norms is shorter: no
Laplacian of any piece is estimated.
\end{proof}

\begin{lemma}[Ward identities under nonlinear gauge variations]
\label{4.5:lem-ward-identities}
Let $U$ be a background configuration. Write the lattice action as a functional $\mathsf{A}$ of
the fluctuation $A$ about $U$,
\begin{equation}\label{4.5:eq-ward-identities-1}
\mathsf{A}=\mathsf{A}(U)+\langle A,J(U)\rangle+\tfrac12\langle A,\Delta A\rangle+V_0(A),
\end{equation}
where $J(U)$ is the current, $\Delta$ is the operator of the quadratic part, taken symmetric for
the bilinear form $\langle\cdot,\cdot\rangle$, and $V_0$ is of order at least three in $A$. Let
$c^{(1)}$, $c^{(2)}$ and $\mathfrak{j}_J$, $\mathfrak{k}_J$, $\mathfrak{m}_J$ be the second-order
terms and the insertions of Lemma~\ref{4.5:lem-free-current-operators}, and let $\Theta'_J$ be the
operator of that lemma. Then
\begin{equation}\label{4.5:eq-ward-identities-a}
\begin{gathered}
D^{*}_{U}J(U)=0,\qquad
D^{*}_{U}\Delta A=\mathfrak{j}_{J(U)}(A),\qquad
\Delta D_U\lambda=\mathfrak{k}_{J(U)}(\lambda),\\
\langle D_U\lambda_1,\Delta D_U\lambda_2\rangle
=-2\langle c^{(2)}(\lambda_1,\lambda_2),J(U)\rangle .
\end{gathered}
\end{equation}
Moreover, with $\Delta_\pi$ and $\Lambda:=G'RD^{*}_{U}$ as in \cite[(3.119)]{B-CMP99b} and
$\Delta'_\pi:=\Delta-\Delta_\pi$, one has $\Delta'_\pi=\Theta'_{J(U)}$; that is, the formula
for $\Theta'_J$ in Lemma~\ref{4.5:lem-free-current-operators} is \cite[(3.120)]{B-CMP99b}, the
sign of the term $\Lambda^{T}\mathfrak{m}_{J}\Lambda$ included, with the weight $J(U)$ replaced by
a free current $J$.
\end{lemma}

\begin{proof}
By gauge invariance of the action \cite[pp.~422--423]{B-CMP99b}, $\mathsf{A}$ is invariant under
\begin{equation}\label{4.5:eq-ward-identities-2}
A\mapsto A^{\lambda}=A-D_U\lambda+c^{(1)}(A,\lambda)+c^{(2)}(\lambda,\lambda)+\dots,
\end{equation}
where the omitted terms are of higher order in $(A,\lambda)$. Insert $A^\lambda$ into
\eqref{4.5:eq-ward-identities-1} and compare the homogeneous parts of each degree in
$(A,\lambda)$ of $\mathsf{A}(A^\lambda)$ and $\mathsf{A}(A)$. Since $A^\lambda$ vanishes to first
order at $(A,\lambda)=(0,0)$ and $V_0$ is cubic, $V_0(A^\lambda)$ contributes only in degree at
least three, as do the omitted terms of \eqref{4.5:eq-ward-identities-2} inserted into the linear
and quadratic parts. Recall that $D^{*}_{U}$ is the transpose of $D_U$:
$\langle D_U\lambda,X\rangle=\langle\lambda,D^{*}_{U}X\rangle$.

Degree one, linear in $\lambda$: $\langle -D_U\lambda,J(U)\rangle=0$ for all $\lambda$, that is,
$\langle\lambda,D^{*}_{U}J(U)\rangle=0$, whence $D^{*}_{U}J(U)=0$.

Degree two, bilinear in $(A,\lambda)$: using the symmetry of $\Delta$,
\begin{equation*}
\langle c^{(1)}(A,\lambda),J(U)\rangle-\langle D_U\lambda,\Delta A\rangle=0 .
\end{equation*}
By the definitions of Lemma~\ref{4.5:lem-free-current-operators},
$\langle c^{(1)}(A,\lambda),J\rangle=\langle\lambda,\mathfrak{j}_J(A)\rangle
=\langle A,\mathfrak{k}_J(\lambda)\rangle$, while
$\langle D_U\lambda,\Delta A\rangle=\langle\lambda,D^{*}_{U}\Delta A\rangle
=\langle A,\Delta D_U\lambda\rangle$. As $A$ and $\lambda$ are arbitrary,
$D^{*}_{U}\Delta A=\mathfrak{j}_{J(U)}(A)$ and $\Delta D_U\lambda=\mathfrak{k}_{J(U)}(\lambda)$.

Degree two, quadratic in $\lambda$:
$\langle c^{(2)}(\lambda,\lambda),J(U)\rangle+\frac12\langle D_U\lambda,\Delta D_U\lambda\rangle=0$.
Both sides are quadratic forms in $\lambda$, the second symmetric because $\Delta$ is; polarising
gives the last identity of \eqref{4.5:eq-ward-identities-a}. With the polarised definition of
$\mathfrak{m}_J$ this reads $D^{*}_{U}\Delta D_U=-\mathfrak{m}_{J(U)}$.

By \cite[(3.119)]{B-CMP99b}, $\Delta_\pi=\Pi^{T}\Delta\Pi$ with $\Pi=I-D_U\Lambda$, so
$\Pi^{T}=I-\Lambda^{T}D^{*}_{U}$. Expanding the product,
\begin{align*}
\Delta'_\pi=\Delta-\Pi^{T}\Delta\Pi
&=\Lambda^{T}D^{*}_{U}\Delta+\Delta D_U\Lambda-\Lambda^{T}D^{*}_{U}\Delta D_U\Lambda\\
&=\Lambda^{T}\mathfrak{j}_{J(U)}+\mathfrak{k}_{J(U)}\Lambda
+\Lambda^{T}\mathfrak{m}_{J(U)}\Lambda=\Theta'_{J(U)},
\end{align*}
where the second line substitutes the three operator identities just obtained, the last one
turning $-\Lambda^{T}D^{*}_{U}\Delta D_U\Lambda$ into $+\Lambda^{T}\mathfrak{m}_{J(U)}\Lambda$.
This is the expression of \cite[(3.120)]{B-CMP99b} for $\Delta'_\pi$, with the same sign of
the last term.
\end{proof}

\begin{lemma}[Reduction to Ba{\l}aban's operators at a true pair]\label{4.5:lem-true-pair-reduction}
Let $U_0$ be a real configuration for which the quantities in \cite[(3.35), (3.36)]{B-CMP99b}
are small, and put $J := J(U_0)$, so that Ba{\l}aban's operators $G_{\pi}$ and $H$ of
\cite{B-CMP99b} exist and satisfy $R D^{*}_{U_0} H = 0$, \cite[(3.124)]{B-CMP99b}. Then the
following hold.
\begin{enumerate}
\item[(i)] The free-current operators at $J = J(U_0)$ are Ba{\l}aban's operators:
\begin{equation*}
\Theta'_J = \Delta'_{\pi}, \quad G^{J}_{\pi} = G_{\pi}, \quad H^{J} = H, \quad
\Theta_J = \Delta^{(2)}, \quad D^{J} = D, \quad \mathfrak{D}^{J} = \mathfrak{D}.
\end{equation*}
\item[(ii)] Let $\mathfrak{L}(A') = A' - HD(A')$ be Ba{\l}aban's linearising transformation.
For every $A'$ in the domain \cite[(77)]{B-CMP102b},
\begin{equation}\label{4.5:eq-true-pair-reduction-1}
\begin{split}
\mathsf{A}\bigl(e^{i\eta\mathfrak{L}(A')}U_0\bigr)
&= \mathsf{A}(U_0) + \langle A', J\rangle \\
&\quad + \tfrac12 \bigl\langle A', (\Delta - \Delta^{(2)})A'\bigr\rangle + V^{J}(A'),
\end{split}
\end{equation}
with $V^{J}$ the higher-order action functional at free current of
Definition~\ref{4.5:def-higher-order-functional}.
That is, $V^{J} = V^{\Delta}$, where $V^{\Delta}$ denotes the functional
\cite[(80)]{B-CMP102b} ``with the operator $\Delta$ instead of $\Delta_{\pi}$''.
\item[(iii)] The gradient of the higher-order functional that enters the equations
\cite[(127), (128), (131)--(133), (143), (180)]{B-CMP102b} is $\nabla V^{\Delta}$.
\end{enumerate}
\end{lemma}

\begin{proof}
Throughout, $J = J(U_0)$, and we abbreviate
$\mathfrak{j} := \mathfrak{j}_J$, $\mathfrak{k} := \mathfrak{k}_J$, $\mathfrak{m} := \mathfrak{m}_J$,
and $D := D(A')$ when the argument is clear.

(i) The operators $\Theta'_J$, $G^{J}_{\pi}$, $H^{J}$, $\Theta_J$ of
Lemma~\ref{4.5:lem-free-current-operators}, and $D^{J}$, $\mathfrak{D}^{J}$ of
Lemma~\ref{4.5:lem-nonlinear-correction}, are defined by Ba{\l}aban's formulas with the
current $J(U)$ replaced by a free current $J$. At $J = J(U_0)$ the definitions are therefore
those of $\Delta'_{\pi}$, $G_{\pi}$, $H$, $\Delta^{(2)}$, $D$, $\mathfrak{D}$. For the first
identity this is the computation in the lemma on Ward identities under nonlinear gauge
variations (Lemma~\ref{4.5:lem-ward-identities}): with $\Lambda = G' R D^{*}_{U_0}$,
\begin{equation}\label{4.5:eq-true-pair-reduction-2}
\Delta'_{\pi} = \Delta - \Delta_{\pi}
= \Lambda^{T}\mathfrak{j} + \mathfrak{k}\Lambda + \Lambda^{T}\mathfrak{m}\Lambda
= \Theta'_J .
\end{equation}

(ii) Since $\Lambda = G' R D^{*}_{U_0}$ and $R D^{*}_{U_0} H = 0$, we have $\Lambda H = 0$.
Hence, for every $X$, by \eqref{4.5:eq-true-pair-reduction-2} and the linearity of
$\mathfrak{k}$ and $\mathfrak{m}$ in their arguments,
\begin{equation*}
\Delta'_{\pi} H X
= \Lambda^{T}\mathfrak{j}(HX) + \mathfrak{k}(\Lambda H X) + \Lambda^{T}\mathfrak{m}(\Lambda H X)
= \Lambda^{T}\mathfrak{j}(HX),
\end{equation*}
which is \cite[(138)]{B-CMP102b}. Taking the pairing with $HX$ and with $A'$, and using the
definition of the bilinear transpose, we obtain
\begin{equation}\label{4.5:eq-true-pair-reduction-3}
\langle HX, \Delta'_{\pi} H X\rangle = \langle \Lambda H X, \mathfrak{j}(HX)\rangle = 0,
\qquad
\langle A', \Delta'_{\pi} H X\rangle = \langle \Lambda A', \mathfrak{j}(HX)\rangle .
\end{equation}
Let $V_{(80)}$ denote the functional \cite[(80)]{B-CMP102b} itself. The operator
$\Delta_{\pi}$ enters it through its second and third terms,
$-\langle A', \Delta_{\pi} H D\rangle$ and $\tfrac12 \langle HD, \Delta_{\pi} H D\rangle$.
Replacing $\Delta_{\pi}$ there by $\Delta = \Delta_{\pi} + \Delta'_{\pi}$ and applying
\eqref{4.5:eq-true-pair-reduction-3} with $X = D$, we get
\begin{equation*}
\begin{split}
V^{\Delta} - V_{(80)}
&= -\langle A', \Delta'_{\pi} H D\rangle + \tfrac12 \langle HD, \Delta'_{\pi} H D\rangle \\
&= -\langle \Lambda A', \mathfrak{j}(H D)\rangle = W^{J}(A'),
\end{split}
\end{equation*}
the last equality by the definition of $W^{J}$ in \eqref{4.5:eq-higher-order-functional-a}
together with $H^{J} = H$ and $D^{J} = D$ from (i).

Next, $V_{(80)} = V^{J,\mathrm{disp}}$. By its definition in
\eqref{4.5:eq-higher-order-functional-a}, the first and the last term of
$V^{J,\mathrm{disp}}$ are those of \cite[(80)]{B-CMP102b} with $H^{J}$, $D^{J}$ in place of
$H$, $D$, so they coincide with the corresponding terms of \cite[(80)]{B-CMP102b} by (i).
For the second term, $R D^{*}_{U_0} H = 0$ and the formulas for the operators at free current
collected in \eqref{4.5:eq-free-current-operators-a}, taken at $J = J(U_0)$,
give
\begin{equation*}
\langle A', \Delta_{\pi} H D\rangle
= \bigl\langle A', (\Delta_{\pi} + D_{U_0} R D^{*}_{U_0}) H D\bigr\rangle
= \langle Q A', \mathsf{k} D\rangle_{\mathfrak{B}} .
\end{equation*}
In the same way, by $R D^{*}_{U_0} H = 0$ and the same formulas, the third term satisfies
\begin{equation*}
\tfrac12\langle HD, \Delta_{\pi} H D\rangle
= \tfrac12\bigl\langle HD, (\Delta_{\pi} + D_{U_0} R D^{*}_{U_0}) H D\bigr\rangle
= \tfrac12 \langle \ell^{-1} D, \mathsf{k} D\rangle_{\mathfrak{B}} .
\end{equation*}
Thus $V_{(80)} = V^{J,\mathrm{disp}}$, and with the previous paragraph and the definition
$V^{J} = V^{J,\mathrm{disp}} + W^{J}$ we obtain $V^{\Delta} = V^{J}$.

It remains to obtain the expansion. Write the action about $U_0$ in the form used in
Lemma~\ref{4.5:lem-ward-identities},
\begin{equation*}
\mathsf{A}(e^{i\eta A}U_0) = \mathsf{A}(U_0) + \langle A, J\rangle
+ \tfrac12\langle A, \Delta A\rangle + V_0(A),
\end{equation*}
with $V_0$ cubic. Substitute
$A = \mathfrak{L}(A') = A' - HD(A')$ in this expansion and rearrange as in
\cite[(78)]{B-CMP102b}; the result is
\eqref{4.5:eq-true-pair-reduction-1} with $V^{\Delta}$ in place of $V^{J}$. Since
$V^{\Delta} = V^{J}$, this proves (ii).

(iii) In \cite{B-CMP102b} the linearising transformation is applied and the functional
\cite[(74)]{B-CMP102b} is obtained ``with the operator $\Delta$ instead of $\Delta_{\pi}$'';
this functional is the left-hand side of \eqref{4.5:eq-true-pair-reduction-1}. Equation
\cite[(127)]{B-CMP102b} expresses the criticality of this functional against variations in
$\ker Q$, and \cite[(128), (131)--(133), (143), (180)]{B-CMP102b} do not alter its
higher-order term. By (ii) the gradient carried through these equations is therefore
$\nabla V^{\Delta}$.

The term $W^{J}$ is needed in the definition of $V^{J}$ for the following reason. If the
gradient $\nabla V_{(80)}$ were used in the free-current version of \cite[(180)]{B-CMP102b},
the identity between gradients that this version requires would not follow: at the argument
$A'_1$ at which that identity is evaluated, $\nabla V_{(80)}$ and $\nabla V^{\Delta}$ differ by
$-\Lambda^{T}\mathfrak{j}(HD(A'_1))$.
\end{proof}

\begin{proposition}[Exactness of the chart on the diagonal]\label{4.5:prf-diagonal-exactness}
Let $V_0$ be real and satisfy \cite[(7)]{B-CMP102b} with curvature parameter
$\varepsilon_1\le\bar\varepsilon_1$, where $\varepsilon_1$ denotes the parameter so named in
\cite[(7)]{B-CMP102b}, and let $U_0=U_k(V_0)$. In this situation $U_0$ lies in
the set \cite[(8)]{B-CMP102b}, $U_0$ satisfies the regularity condition
\eqref{4.5:eq-admissible-pairs-a} on every cube of the class considered in \cite{B-CMP99b}, and
\begin{equation*}
|J(U_0)|_{(-3)}\le C_1B_3\bar\varepsilon_1<a_\gamma ,
\end{equation*}
where $C_1$ and $B_3$ are fixed constants, not depending on $V_0$.
Put $J:=J(U_0)$. Let $B$ lie in the complex polydisc \cite[(172)]{B-CMP102b} with $|B|<c_9$, let
$A'_1$ be the critical configuration of \cite[(127)]{B-CMP102b} with block datum $B$, and let $U_1$
and $\mathcal{H}(B)$ be the configuration and the background potential attached to $A'_1$ there; in
this situation $N(A'_1-H_0B)\le\bar\varepsilon_4$, where $\bar\varepsilon_4$ is the radius of
Definition~\ref{4.5:def-response-chart}: this is the choice of $\bar\varepsilon_4$, since the
statement fixing that radius bounds $N(A'_1-H_0B)$ for Ba{\l}aban's critical configuration at every
$B$ in the polydisc with $|B|<c_9$. Let $v$ be a block direction and let $A(t_\square)$ and
$\mathfrak{J}(B,t_\square)$ be as in \eqref{4.5:eq-response-chart-a}. Then
\begin{equation}\label{4.5:eq-diagonal-exactness-a}
\begin{aligned}
&\text{(a)}\quad A_G=A'_1,\qquad RD^{*}_{U_0}A_G=0,\\
&\phantom{\text{(a)}}\quad \mathcal{H}^{J}(B)=A'_1-HD(A'_1)=\frac{1}{i\eta}\log U_1=\mathcal{H}(B),\\
&\text{(b)}\quad \Lambda^{J}(B)=\Delta(U_0)\mathcal{H}^{J}(B),\\
&\phantom{\text{(b)}}\quad \mathfrak{J}(B)=D^{*}_{U_0}D_{U_0}\mathcal{H}^{J}(B)+F(U_0,\mathcal{H}^{J}(B))
=J\bigl(e^{i\eta\mathcal{H}^{J}(B)}U_0\bigr)-J(U_0),\\
&\text{(c)}\quad J+\mathfrak{J}(B,t_\square)=J\bigl(e^{i\eta A(t_\square)}U_0\bigr)
\quad\text{for all complex } t_\square ,\\
&\text{(d)}\quad \text{the datum produced by the chart, the orbit of }
\bigl(e^{i\eta\mathcal{H}^{J}(B)}U_0,\,J+\mathfrak{J}(B)\bigr),\\
&\phantom{\text{(d)}}\quad \text{is a true pair.}
\end{aligned}
\end{equation}
\end{proposition}

This is clause (K6) of the existence, uniqueness and analyticity of the response chart
(Theorem~\ref{4.5:thm-chart-existence}), with its hypotheses made explicit.

\begin{proof}
\emph{Regularity of $U_0$.} Ba{\l}aban states in \cite{B-CMP102b}: ``The additional regularity
condition (3.35) in (1.33) is satisfied for $U_0$ also, because of the result of Sect.~F''; his
condition \cite[(3.35)]{B-CMP99b} is the condition \eqref{4.5:eq-admissible-pairs-a}. The result
invoked is
\cite[Theorem~1, (9)]{B-CMP102b}, which holds ``for an arbitrary cube $\square$ in the class
described above, of a size $2ML^j\eta$, $M\le M(\varepsilon_1)$ \dots\ More exactly
$M(\varepsilon_1)=R_1M_1(a_1/\varepsilon_1)$'', where Ba{\l}aban's $M_1$ is $\mathfrak{M}$ and
$a_1$ is his constant of that theorem. The cubes
on which \eqref{4.5:eq-admissible-pairs-a} is required have size $12\mathfrak{M}L^j\eta$, that is
$M=6\mathfrak{M}$, and
\begin{equation*}
12\mathfrak{M}\le 2R_1\mathfrak{M}\,a_1/\bar\varepsilon_1 ,
\end{equation*}
a ceiling on $\bar\varepsilon_1$ depending on $(d,L)$ only, built into the definition of
$\bar\varepsilon_1$ in \eqref{4.5:eq-geometry-conventions-a}. Since
$\varepsilon_1\le\bar\varepsilon_1$, this gives
\begin{equation*}
6\mathfrak{M}\le R_1\mathfrak{M}\,a_1/\bar\varepsilon_1
\le R_1\mathfrak{M}\,a_1/\varepsilon_1=M(\varepsilon_1),
\end{equation*}
so \cite[Theorem~1]{B-CMP102b} applies on these cubes. In particular $U_0$ is real and regular,
so that the reduction to Ba{\l}aban's operators at a true pair
(Lemma~\ref{4.5:lem-true-pair-reduction}) applies to $U_0$: $RD^{*}_{U_0}H=0$, $H^{J}=H$, $D^{J}=D$,
$\Theta_J=\Delta^{(2)}$, and the expansion of Lemma~\ref{4.5:lem-true-pair-reduction}(ii) holds.
These facts, together with the bound on $|J|_{(-3)}$ in the statement, place $(U_0,J)$ in the
class $\mathfrak{S}$ of admissible pairs of \eqref{4.5:eq-admissible-pairs-a}, so that
Theorem~\ref{4.5:thm-chart-existence} applies at $(U_0,J)$ for $|B|<c_9$.

\emph{(a), real $B$.} Let $B$ be real. By \cite[(127)]{B-CMP102b}, $A'_1$ is critical for the
functional $A'\mapsto\mathsf{A}(e^{i\eta\mathfrak{L}(A')}U_0)$ on the affine set $\{\ell QA'=B\}$; in
Ba{\l}aban's words, ``we obtain the following equation on $A'_1$, (127), for all $\delta A'$
satisfying $Q\delta A'=0$''. By Lemma~\ref{4.5:lem-true-pair-reduction}(ii),
\begin{equation*}
\mathsf{A}(e^{i\eta\mathfrak{L}(A')}U_0)=\mathsf{A}(U_0)+\langle A',J\rangle
+\tfrac12\langle A',(\Delta-\Delta^{(2)})A'\rangle+V^{J}(A'),
\end{equation*}
and
$\Delta^{(2)}=\Theta_J$. Differentiating at $A'_1$ in a direction $\delta A'$ with $Q\delta A'=0$,
the criticality reads
\begin{equation}\label{4.5:eq-diagonal-exactness-1}
\bigl\langle\delta A',\,J+(\Delta-\Theta_J)A'_1+\mathbb{V}'^{J}(A'_1)\bigr\rangle=0
\qquad\text{for all }\delta A'\text{ with }Q\delta A'=0 .
\end{equation}
Ba{\l}aban notes that ``The configuration $A'_1$ satisfies $RD^{*}A'_1=0$'' \cite{B-CMP102b},
that is, $RD^{*}_{U_0}A'_1=0$.
The formulas defining $\Delta_a$, displayed together with the representation identity
\eqref{4.5:eq-response-chart-a}, together with $RD^{*}_{U_0}A'_1=0$ give
$(\Delta_a-Q^{*}aQ)A'_1=\Delta A'_1$; only the formulas defining $\Delta_a$ are used here, not the
identity satisfied by $A_G$, which has not yet been identified with $A'_1$. Moreover
$\langle\delta A',Q^{*}aQA'_1\rangle=\langle Q\delta A',aQA'_1\rangle=0$ when $Q\delta A'=0$;
hence $\Delta$ may be replaced by $\Delta_a$ in \eqref{4.5:eq-diagonal-exactness-1}. Write
\begin{equation*}
\Psi:=\Theta_JA'_1-\mathbb{V}'^{J}(A'_1).
\end{equation*}
Since $\ker Q=\operatorname{range}P_0$, the admissible directions are exactly
$\delta A'=P_0A''$ with $A''$ arbitrary, and \eqref{4.5:eq-diagonal-exactness-1} becomes
$\langle A'',P_0^{T}[J+\Delta_aA'_1-\Psi]\rangle=0$ for all $A''$, that is
\begin{equation}\label{4.5:eq-diagonal-exactness-2}
P_0^{T}\bigl[J+\Delta_aA'_1-\Psi\bigr]=0 .
\end{equation}
By \cite[(178)]{B-CMP102b}, ``$\langle\delta A',J\rangle=0$ for $\delta A'$:
$Q(U_k)\delta A'=0$'', so $J\in(\ker Q)^{\perp}=\operatorname{range}Q^{*}$; writing $J=Q^{*}y$ we
get $P_0^{T}J=(QP_0)^{T}y=0$. Further
$P_0^{T}\Delta_a=\Delta_a-Q^{*}\mathsf{q}_0Q=\Delta_aP_0$, and $P_0A'_1=A'_1-H_0B$ because
$\ell QA'_1=B$. Hence \eqref{4.5:eq-diagonal-exactness-2} gives
$\Delta_a(A'_1-H_0B)=P_0^{T}\Psi$, and applying $G_0$,
\begin{equation*}
A'_1-H_0B=G_0P_0^{T}\Psi=\tilde{\mathfrak{G}}_0\Psi .
\end{equation*}
Since $\Psi=\Theta_J(\mathcal{A}+H_0B)-\mathbb{V}'^{J}(\mathcal{A}+H_0B)$ with
$\mathcal{A}=A'_1-H_0B$, this is the fixed-point equation defining the response chart
(Definition~\ref{4.5:def-response-chart}), and the bound $N(A'_1-H_0B)\le\bar\varepsilon_4$,
which holds by the choice of the radius $\bar\varepsilon_4$ of
Definition~\ref{4.5:def-response-chart} recalled in the statement, places
$A'_1-H_0B$ in the ball in which that fixed point is unique by the existence, uniqueness and
analyticity of the response chart (Theorem~\ref{4.5:thm-chart-existence}). Therefore
$\mathcal{A}^{J}_0(B)=A'_1-H_0B$, i.e.\ $A_G=A'_1$, and $RD^{*}_{U_0}A_G=0$. Since
$X=D^{J}(A_G)=D(A'_1)$ and $H^{J}=H$, the definition of the chart gives
$\mathcal{H}^{J}(B)=A_G-H^{J}X=A'_1-HD(A'_1)$, which \cite{B-CMP102b} identifies with
$\frac{1}{i\eta}\log U_1=\mathcal{H}(B)$.

\emph{(a), complex $B$.} Both sides of each identity in (a) are analytic in $B$ on the polydisc,
the chart side by Theorem~\ref{4.5:thm-chart-existence}; they agree at real $B$, hence on the
whole polydisc.

\emph{(b).} First, let $\tilde B$ be the block datum entering the representation identity
\eqref{4.5:eq-response-chart-a}. The quantity $\Theta_JA_G-\mathbb{V}'^{J}(A_G)$ enters that
identity and equals $\Psi$ at $A_G=A'_1$ by (a); the identity therefore gives
\begin{equation*}
P_0^{T}\Psi+Q^{*}(\mathsf{q}_0-a)\tilde B=(\Delta_a-Q^{*}aQ)A'_1=\Delta A'_1,
\end{equation*}
the last equality because $RD^{*}_{U_0}A'_1=0$. Second,
\begin{equation*}
\Delta HX=\Delta_\pi HX+\Delta'_\pi HX=Q^{*}\mathsf{k}X+\Lambda^{T}\mathfrak{j}_J(HX),
\end{equation*}
by \cite[(137)]{B-CMP102b}, which needs $RD^{*}_{U_0}H=0$ and $QH=\ell^{-1}$, and by
\cite[(138)]{B-CMP102b}, which needs $\Lambda H=0$; the condition $RD^{*}_{U_0}H=0$ enters both,
\cite[(138)]{B-CMP102b} through $\Lambda H=G'RD^{*}_{U_0}H=0$. As
$\Delta-\Delta'=D^{*}_{U}D_{U}$ by definition, on the diagonal
\begin{equation*}
D^{*}_{U_0}D_{U_0}(HX)=Q^{*}\mathsf{k}X+\Lambda^{T}\mathfrak{j}_J(HX)-\Delta'HX
=\mathcal{K}_1X ,
\end{equation*}
where $\mathcal{K}_1$ is the operator so denoted in \eqref{4.5:eq-response-chart-a}.
By \eqref{4.5:eq-response-chart-a} the current correction has the two expressions
\begin{align*}
\mathfrak{J}&=(\Delta_a-Q^{*}aQ-\Delta')A_G-\mathcal{K}_1X+F(U_0,\mathcal{H}^{J}),\\
\mathfrak{J}&=\Lambda^{J}-\Delta'\mathcal{H}^{J}+F(U_0,\mathcal{H}^{J}),
\end{align*}
the second being the definition of $\mathfrak{J}(B,t_\square)$ at $t_\square=0$, with
$\mathcal{H}^{J}=A_G-HX$. Comparing them,
\begin{equation*}
\Lambda^{J}=(\Delta_a-Q^{*}aQ)A_G-\mathcal{K}_1X-\Delta'HX ,
\end{equation*}
and by the two facts just proved
\begin{equation*}
\Lambda^{J}(B)=\Delta A'_1-\Delta HX=\Delta(A'_1-HX)=\Delta(U_0)\mathcal{H}^{J}(B).
\end{equation*}
Consequently
\begin{equation*}
\mathfrak{J}(B)=(\Delta-\Delta')\mathcal{H}^{J}+F(U_0,\mathcal{H}^{J})
=D^{*}_{U_0}D_{U_0}\mathcal{H}^{J}+F(U_0,\mathcal{H}^{J}).
\end{equation*}
Third, \cite[(3.11)]{B-CMP109} reads
\begin{equation*}
\begin{split}
D^{\xi *}_{\exp i\xi\mathbf{A}U}\,\xi^{-2}\pi\operatorname{Im}\partial\exp i\xi\mathbf{A}U
&=D^{\xi *}_{U}\,\xi^{-2}\pi\operatorname{Im}\partial U\\
&\quad+D^{\xi *}_{U}D^{\xi}_{U}\mathbf{A}+\mathbf{F}(U,\mathbf{A}),
\end{split}
\end{equation*}
whose left side and first term on the right are the currents of $\exp(i\xi\mathbf{A})U$ and of
$U$; here $\mathbf{F}$ is the quadratic remainder in the current, written $F$ above
(Lemma~\ref{4.4:lem-current-remainder}). At $\xi=\eta$, $U=U_0$, $\mathbf{A}=\mathcal{H}^{J}(B)$ it
gives
\begin{equation*}
D^{*}_{U_0}D_{U_0}\mathcal{H}^{J}+F(U_0,\mathcal{H}^{J})
=J(e^{i\eta\mathcal{H}^{J}(B)}U_0)-J(U_0).
\end{equation*}

\emph{(c).} The identity $\Lambda^{J}=\Delta\mathcal{H}^{J}$ holds as an analytic identity in $B$,
so differentiating in the block direction $v$ gives
$\delta\Lambda^{J}[v]=\Delta\,\delta\mathcal{H}^{J}[v]$, and therefore
\begin{equation*}
\delta\Lambda^{J}[v]-\Delta'\delta\mathcal{H}^{J}[v]=D^{*}_{U_0}D_{U_0}\delta\mathcal{H}^{J}[v].
\end{equation*}
Inserting this and (b) in the definition of $\mathfrak{J}(B,t_\square)$ in
\eqref{4.5:eq-response-chart-a},
\begin{equation*}
\mathfrak{J}(B,t_\square)=D^{*}_{U_0}D_{U_0}A(t_\square)+F(U_0,A(t_\square)),
\end{equation*}
and \cite[(3.11)]{B-CMP109} at $\mathbf{A}=A(t_\square)$, both sides being entire in
$t_\square$, gives $J+\mathfrak{J}(B,t_\square)=J(e^{i\eta A(t_\square)}U_0)$ for all complex
$t_\square$.

\emph{(d).} By (b), $J+\mathfrak{J}(B)=J(e^{i\eta\mathcal{H}^{J}(B)}U_0)$: the current of the
produced datum is the current of its configuration.
\end{proof}

Exactness of this kind is required, and holds, at true pairs only: at the step's configuration
$U_{k+1}$, which equals $U_j(M^{j}(U_{k+1}))$ for each $j$, with $M^{j}(U_{k+1})$ the block
averages of $U_{k+1}$ at level $j$ from the tower $M^{\cdot}(U_{k+1})$, and equals
$U_{k+1}(\square_0,M^{\cdot}U_{k+1})$ on the window $\square_0$, by the lemma identifying the
step's configuration on the window; at images of the chart at real data; at
images of real fields under the tower map of \cite[(1.15)]{B-CMP109}; at restrictions of such
pairs to sub-regions; and at
minimisers with histories. Complex pairs $(U,J)$, points with decoupling parameters
$s(\Delta)\not\equiv1$ in the sense of \cite{B-CMP116}, and the complex discs of the parameters
$t_\square$, $s(\Delta)$ and $\sigma$, the last being the parameter of the sets $Y(\sigma)$ of
\cite{B-CMP116}, serve only as points of Cauchy estimates: analyticity, bounds,
covariance and reality are read there, exactness is not.

\begin{remark}[Condition (iv): the uniform fine-scale curvature bound]\label{4.5:rem-condition-iv}
Condition (iv) of \cite[p.~262]{B-CMP109}, a condition on the configurations over a localisation
domain $X$ (not to be confused with clause (iv) of Theorem~0), reads, in the words of
\cite[(1.15)--(1.16)]{B-CMP109}:
\begin{quote}
We consider $X$ as $\Omega_0$, and we construct the sequence $\{\Omega_n\}$, $n = 1,\dots,j$,
of maximal possible domains satisfying the conditions (1.3)--(1.6) [14] (with $j$, $\xi$
instead of $k$, $\eta$, $R = R_1$). For this sequence, or rather for its subsequences
$\{\Omega_0,\Omega_1,\dots,\Omega_n\}$, $n = 1,\dots,j$, we construct the functions $U_n(V)$ in
the axial gauges, for regular $G^{c}$-valued configurations $V$. We consider the pair
\begin{equation*}
\bigl(U_n(M^{\cdot}(\mathbf{U})),\,J_n(M^{\cdot}(\mathbf{U}))\bigr),\qquad n = 1,\dots,j,
\tag*{\cite[(1.15)]{B-CMP109}}
\end{equation*}
where $M^{\cdot}(\mathbf{U},b) = M^{p}(\mathbf{U},b) = \bar{\mathbf{U}}^{p}(b)$ for
$b \in \Lambda_p$ \dots, and $J_n(M^{\cdot}(\mathbf{U}))$ is defined by the formula (1.8) with
$U_n(M^{\cdot}(\mathbf{U}))$ instead of $U_j$, $L^{-n}$ instead of $\xi$. This pair satisfies
the bounds
\begin{equation*}
|\partial U_n(M^{\cdot}(\mathbf{U})) - 1| < \alpha_0\xi^2,\qquad
|J_n(M^{\cdot}(\mathbf{U}))| < \alpha_0(L^n\xi)^2 \quad\text{on } \tilde X^{-2},
\tag*{\cite[(1.16)]{B-CMP109}}
\end{equation*}
and the same bounds hold for $U$ instead of $\mathbf{U}$. The domain $\tilde X^{-2}$ is obtained
from $X$ by taking away two layers of cubes from $\pi_j$, which are closest to the boundary of
$X$.
\end{quote}
Throughout we write $X^{\sim 2}$ for the domain $\tilde X^{-2}$ of this passage. Here $\xi = L^{-j}$
is the spacing at the scale of $X$, so that $L^n\xi = L^{n-j} \le 1$ for $n \le j$, and
$\bar{\mathbf{U}}^{p}$ is the $p$-fold block average of $\mathbf{U}$ on the layer $\Lambda_p$ of
the hierarchy. The two members of the bound of condition (iv) are of different kinds. The first,
the curvature bound $\alpha_0\xi^2$, is uniform: it is the same at every $n = 1,\dots,j$. The
second, $\alpha_0(L^n\xi)^2$, depends on $n$.
For $n = 1,\dots,j$ we put
\begin{equation*}
\lambda := L^n\xi = L^{n-j} \le 1,\qquad \eta_n := \xi/\lambda,
\end{equation*}
the latter being the fine spacing at level $n$; these are the level-$n$ units used below.

In \cite[p.~263, (1.17)]{B-CMP109} the passage is followed by the remark:
\begin{quote}
We have $M^{\cdot}(\mathbf{U}) = \bar{\mathbf{U}}^{p}$ on $\Lambda_p$,
$|\partial\bar{\mathbf{U}}^{p} - 1| < 2\alpha_0'(L^p\xi)^2$. From Proposition 9 [15] we obtain
\begin{gather*}
|\partial U_n(M^{\cdot}(\mathbf{U})) - 1| < B_3 2\alpha'_0 L^{-2n}(L^n\xi)^2
= 2B_3\alpha'_0\xi^2,\\
|J_n(M^{\cdot}(\mathbf{U}))| < B_3 2\alpha'_0(L^n\xi)^2 \quad\text{on } \tilde X^{-2}.
\tag*{\cite[(1.17)]{B-CMP109}}
\end{gather*}
It is obvious that for $\alpha'_0$ sufficiently small the above estimates imply the condition
(iv).
\end{quote}
Here the references [14] and [15] are \cite{B-CMP99a} and \cite{B-CMP102b}, and $B_3$ is the
constant of \cite[Proposition~9]{B-CMP102b}. The equality in the first line of
\cite[(1.17)]{B-CMP109} is the identity $L^{-2n}(L^n\xi)^2 = \xi^2$.

Three ways of obtaining condition (iv) for the pairs produced by a renormalisation step are not
used in this book.
\begin{enumerate}
\item The estimate \cite[(1.17)]{B-CMP109} itself, taken as a route to (iv) for a produced pair,
loses the factor $2B_3$ against the factor $\tfrac12$ of the half radii with which (iv) enters as
a hypothesis (see the paragraph following this list): \cite[(1.17)]{B-CMP109} bounds the curvature
member by $2B_3\alpha'_0\xi^2$, with the constant $2B_3$ where the hypothesis carries $\tfrac12$.
\item Sizing the complex data by the potential, as in \cite[(172)]{B-CMP102b}, with the curvature
radius $\varepsilon_1$ of \cite[(7)]{B-CMP102b} set equal to the curvature of the data,
orders the radii as $\alpha_1 \lesssim c\,\alpha_0\lambda$, $c$ a constant; together with the
condition of \cite{B-CMP109} that relates these radii to $M$, this confines $M$ to a window.
\item At a boldface configuration $\mathbf{U}$, the estimate \cite[(1.17)]{B-CMP109} is obtained
from \cite[Proposition~9]{B-CMP102b}, which is stated for data sized by the potential, not by the
curvature; this book does not use it in that form.
Obtaining it instead by the device of the uniform-increment lemma would cost
$\tfrac12 C_{IV}\alpha'_1\lambda^2$, with $C_{IV}$ the constant of
\eqref{4.5:eq-uniform-increment-a}; this orders $\alpha'_1$ against $\alpha_0$.
\end{enumerate}

The uniform-increment lemma takes condition (iv) for $\mathbf{U}$ as a hypothesis. The first member
of \cite[(3.16)]{B-CMP109} carries it, through the space
\begin{equation*}
U^{c}_{j}\bigl(X,\tfrac12\alpha_0,\tfrac12\alpha_1,\alpha_0\bigr)\cap\{\dots\}
\end{equation*}
of that display, and the lemma differentiates only the perturbation
$\mathbf{U} \mapsto e^{i\xi A}\mathbf{U}$ of radius $\alpha_2$. It does not use
\cite[(1.17)]{B-CMP109} in any step of its proof.

A levelwise route is also not taken. Suppose the increments of the data under the perturbation are
bounded level by level, with weights of the type of \cite[(19)]{B-CMP102b}, an increment of the
datum at level $p$ being of order $\alpha_2(L^p\xi)$ times the weight, and that
\cite[Proposition~9]{B-CMP102b} is then applied. This gives bounds layer by layer on the layers of
$X$: the increment is at most $C\alpha_2L^{-2p}$ on the layer $\Lambda_p$, with a constant $C$.
The curvature member of \cite[(1.16)]{B-CMP109}, however, is the uniform bound
\begin{equation*}
\alpha_0\xi^2 = \alpha_0L^{-2j}\quad\text{on } X^{\sim 2}.
\end{equation*}
The two agree only on the core layer $p = j$. On a deeper layer $p < j$ a smallness condition
$\alpha_2 \le c\,\alpha_0$ would have to be replaced by
\begin{equation*}
\alpha_2 \le c\,\alpha_0L^{-2(j-p)},
\end{equation*}
a smallness that grows with the depth $j - p$ of the layer. In the level-$n$ units $\lambda$ and
$\eta_n$ introduced above, the levelwise bound reads
\begin{equation*}
\bigl|\partial U_n(M^{\cdot}(e^{i\xi A}\mathbf{U})) - \partial U_n(M^{\cdot}(\mathbf{U}))\bigr|
\lesssim C\alpha_2\eta_n^2 = C\alpha_2\xi^2\lambda^{-2}.
\end{equation*}
This is two powers of $\lambda$ short of the uniform bound $\alpha_0\xi^2$ on the deep layers, where
$n \ll j$ and $\lambda \ll 1$. On that route the loss is cured only by
\begin{equation*}
\alpha_2 \lesssim \alpha_0\lambda^2 = \alpha_0L^{-2(j-n)},
\end{equation*}
again a smallness growing with the depth. Such a condition is not admissible, because $\alpha_2$
must be, in the words of \cite[\S3]{B-CMP109}, ``a constant $\alpha_2$, depending on $\alpha_0$ and
on some absolute constants''.

In level-$n$ units both members of \cite[(1.16)]{B-CMP109} carry the factor $\lambda^2$. Put
\begin{equation}\label{4.5:eq-condition-iv-1}
\Phi_1 := \eta_n^{-2}\bigl(\partial U_n(M^{\cdot}(\mathbf{U})) - 1\bigr),\qquad
\Phi_2 := J_n(M^{\cdot}(\mathbf{U})).
\end{equation}
Since $\eta_n^{-2}\xi^2 = \lambda^2$ and $(L^n\xi)^2 = \lambda^2$, the bounds
\cite[(1.16)]{B-CMP109} are equivalent to
\begin{equation}\label{4.5:eq-condition-iv-2}
|\Phi_1| < \alpha_0\lambda^2,\qquad |\Phi_2| < \alpha_0\lambda^2 \quad\text{on } X^{\sim 2}.
\end{equation}
The uniform-increment lemma obtains the two powers $\lambda^2$ for the increments of $\Phi_1$ and
$\Phi_2$ from three ingredients, used together in place of the levelwise route. First, every
$x \in X^{\sim 2}$ lies in the top layer $\Omega_n$ of every sub-tower $\{\Omega_0,\dots,\Omega_n\}$.
The kernel read at the point assigned to $x$ therefore carries the top unit, and no power of the
level of that point is lost. Second, an exact recentring of the perturbation places the two powers
$\lambda^2$ on the near sources. Third, with $a$ the decay constant of the kernels, the decay
factor of the far sources obeys
\begin{equation*}
e^{-a/\lambda} \le e^{-a/2}\lambda,
\end{equation*}
and together with the floor $\delta_0M \ge 32$ this supplies the two powers for the far sources.
The result is the bound \eqref{4.5:eq-uniform-increment-a}: at every level
$n \le j$ and every $x \in X^{\sim 2}$ the increments of $\Phi_1$ and $\Phi_2$ are bounded by
$\tfrac12 C_{IV}\alpha_2\lambda^2$, with one constant $C_{IV}$ that does not depend on the depth
$j - n$.
\end{remark}

\begin{lemma}[Reproduction of condition (iv) under a twist with a depth-free constant. Stated; proof incomplete at the step marked $(\ast)$]
\label{4.5:lem-uniform-increment}
Let $X\in\mathbf{D}_j$, the class of localisation domains of \cite{B-CMP109} at scale $\xi=L^{-j}$.
Let $X^{\sim2}$ be the domain obtained from $X$ by removing the two layers of cubes closest to its
boundary, as in Remark~\ref{4.5:rem-condition-iv}. The following numbers depend on $(d,L)$ only:
\begin{itemize}
\item $\bar\varepsilon_1$ of \eqref{4.5:eq-geometry-conventions-a};
\item $c_Q$ and $C_1$ of \cite[Proposition~6]{B-CMP98} and \cite[(172)]{B-CMP102b};
\item $C_7$, the number written $O(1)$ in the curvature condition \cite[(7)]{B-CMP102b};
\item $c_{12}$, the number written $O(1)$ in \cite[(1.12)]{B-CMP109}, read as in (\textrm{$\rho$}1)
of the corresponding display;
\item $C_K$, the constant of the kernel estimate used in the proof of
\eqref{4.5:eq-uniform-increment-a};
\item $C_{14}:=9(2d+2)$;
\item $c_\sharp:=\min\{\frac14,\frac12\alpha_1^{(98)}\}$, where $\alpha_1^{(98)}$ denotes the
constant $\alpha_1$ of \cite[Proposition~6]{B-CMP98}.
\end{itemize}
The letter $B$ denotes the constant which, in the third ceiling of
\eqref{4.5:eq-uniform-increment-b}, multiplies the product of the cube side $c_{12}LM$ of reading
(\textrm{$\rho$}1) with the amplitude $\frac12\alpha_0$. The letter $\delta_0$ denotes the constant
of the first floor of \eqref{4.5:eq-uniform-increment-b}.

There is a constant $C_{IV}=C_{IV}(d,L)\ge 4C_{14}$ with the following property.

Let $\alpha_0,\alpha_1,\alpha_2>0$, $M$ and $L$ satisfy the ceilings and floors, and let the data of
\cite[(172)]{B-CMP102b} be placed with the margin, as follows:
\begin{equation}\label{4.5:eq-uniform-increment-b}
\begin{gathered}
C_7\alpha_0\le\bar\varepsilon_1,\qquad
c_Q(\alpha_1+\alpha_2)\le\tfrac12C_1\bar\varepsilon_1,\qquad
\tfrac12c_{12}LMB\,\alpha_0+\tfrac12\alpha_1+\alpha_2\le c_\sharp,\\
\alpha_1\le1,\qquad
\alpha_2\le\min\{\alpha_1/8,\ \alpha_0/C_{IV}\};\\
\delta_0M\ge32,\qquad 2R_1M_1\le M,\qquad c_{12}L\ge2,\qquad L\ge4d;\\
\sup|\mathcal{B}|\le\tfrac12C_1\bar\varepsilon_1\ \text{ for every datum } e^{i\mathcal{B}}V_0
\text{ placed in \cite[(172)]{B-CMP102b}}.
\end{gathered}
\end{equation}
Here $R_1$ and $M_1$ are the numbers of \cite[(1)]{B-CMP102b} fixed in the conventions
of~\ref{4.5:not-geometry-conventions}, and $R_1\ge L$, since in \cite[(1.4)]{B-CMP99a} $R$ is a
sufficiently large positive integer and a power of $L$.

Data are placed in \cite[(172)]{B-CMP102b} with the margin of reading (\textrm{$\rho$}5) in
the corresponding display. That is, a datum of \cite[(172)]{B-CMP102b}, written here
$e^{i\mathcal{B}}V_0$ with $\mathcal{B}$ its exponent, is placed with
$\sup|\mathcal{B}|\le\frac12C_1\bar\varepsilon_1$. Then both readings of \cite[(172)]{B-CMP102b},
$|V'-1|<C_1\varepsilon_1$ and $|\mathcal{B}'|<2C_1\varepsilon_1$, where $V'$ is the output of that
display, $\mathcal{B}'$ its exponent and $\varepsilon_1$ the parameter of that display, hold with
$\varepsilon_1$ set equal to $\bar\varepsilon_1$. Indeed $C_1\bar\varepsilon_1\le\frac14$ by the
choice of $\bar\varepsilon_1$ in \eqref{4.5:eq-geometry-conventions-a}, and hence
$e^{\frac12C_1\bar\varepsilon_1}-1<C_1\bar\varepsilon_1$. Moreover the segment $t\mathcal{B}$,
$t\in[0,1]$, stays inside the same bound.

Let $(\mathbf{U},\mathbf{J})$ have a representative modulo $G$-valued gauge maps, in the sense of
(\textrm{$\rho$}4) of the corresponding display, with the following properties.
\begin{itemize}
\item The representative has the form $\mathbf{U}=e^{i\xi A'}U$ with $U$ $G$-valued.
\item It satisfies conditions (i)--(iii) of \cite[(1.11)--(1.14)]{B-CMP109} on $X$ with constants
$(\frac12\alpha_0,\frac12\alpha_1)$.
\item It satisfies condition (iv) (Remark~\ref{4.5:rem-condition-iv}) with $\frac12\alpha_0$ in the
following form. Let $\mathbf{\Omega}=\{\Omega_0=X\supset\Omega_1\supset\dots\supset\Omega_j\}$ be a
tower of $X$, admissible in the sense of \cite[(1)]{B-CMP102b} with block $\mathfrak{M}$ and gap
parameter $R\ge R_1$ in the sense of the conventions of~\ref{4.5:not-geometry-conventions}, each
$\Omega_n$, regarded as a set of bonds, consisting of the bonds with at least one end-point in
$\Omega_n$ (\cite[(1)]{B-CMP102b}), and such that for every $n\le j$ the collar $X\setminus\Omega_n$
has depth in $X$ at most $\frac12M$, in level-$j$ units. It is assumed that the bounds
\cite[(1.16)]{B-CMP109} with $\frac12\alpha_0$ hold on $X^{\sim2}$ for every sub-tower
$\{\Omega_0,\dots,\Omega_n\}$, $n=1,\dots,j$, both for $\mathbf{U}$ and for $U$.
\end{itemize}
Nothing of the tower is used beyond its admissibility, the collar bound and the bounds
\cite[(1.16)]{B-CMP109} just assumed.

The towers covered include the maximal tower of $X$ with $R=R_1$ of condition (iv) in
\cite{B-CMP109}. Its collars have depth at most $R_1M_1$, a bound belonging to the step marked
$(\ast)$ in the proof below, hence at most $\frac12M$ under the floor $2R_1M_1\le M$. They also
include the tower of
\cite[(2.13)]{B-CMP119} of reading (\textrm{$\rho$}2), wherever it satisfies the collar bound and
condition (iv) is established for it.

Let $A$ be $\mathfrak{g}^c$-valued on the bonds of $X$ with $|A|<\alpha_2$ and
$|\nabla^\xi_{\mathbf{U}}A|<\alpha_2$, in the conventions of \cite{B-CMP109}.

For $n\in\{1,\dots,j\}$ use level-$n$ units: set $\lambda:=L^n\xi=L^{n-j}\le1$ and
$\eta_n:=\xi/\lambda$. For a configuration $V$ write
\begin{equation*}
\Phi_1(V):=\eta_n^{-2}\bigl(\partial U_n(M^{\cdot}(V))-1\bigr),\qquad
\Phi_2(V):=J_n(M^{\cdot}(V)).
\end{equation*}
Here $M^{\cdot}$ is the tower of block averages, and $U_n$, $J_n$ are the functions of
Remark~\ref{4.5:rem-condition-iv} built on the sub-tower $\{\Omega_0,\dots,\Omega_n\}$ of
$\mathbf{\Omega}$; the $\Phi_i$ are read pointwise.
The constant $C_{IV}$ is independent of $M$, of $n$, $j$, of the level $k$ of the renormalisation
step at which the lemma is applied, of $X$ and of the depth $j-n$.

Then $e^{i\xi A}\mathbf{U}=e^{i\xi A''}U$ satisfies on $X$ the following.
\begin{enumerate}
\item[(i)] The $G$-valued factor is the same $U$, so \cite[(1.11)--(1.12)]{B-CMP109} hold unchanged.
\item[(ii)] The potential satisfies $|A''|<\frac12\alpha_1+4\alpha_2\le\alpha_1$ and
$|\nabla^\xi_UA''|<\frac12\alpha_1+4\alpha_2\le\alpha_1$.
\item[(iii)$_1$] The first member of condition (iii), \cite[(1.14)]{B-CMP109}, holds in the form
\begin{equation*}
|\partial(e^{i\xi A}\mathbf{U})-1|<(\tfrac12\alpha_0+C_{14}\alpha_2)\xi^2\le\alpha_0\xi^2 .
\end{equation*}
\item[(iv)] Condition (iv) holds with $\alpha_0$, both halves, for the same tower $\mathbf{\Omega}$.
The half for $e^{i\xi A}\mathbf{U}$ holds in the following uniform form: for every $n=1,\dots,j$ and
every $x\in X^{\sim2}$,
\begin{equation}\label{4.5:eq-uniform-increment-a}
|\Phi_i(e^{i\xi A}\mathbf{U})-\Phi_i(\mathbf{U})|\le\tfrac12C_{IV}\alpha_2\lambda^2,\qquad i=1,2 .
\end{equation}
Consequently, under $\alpha_2\le\alpha_0/C_{IV}$, the hypothesis $|\Phi_i(\mathbf{U})|<\frac12\alpha_0\lambda^2$
gives $|\Phi_i(e^{i\xi A}\mathbf{U})|<\alpha_0\lambda^2$. In $\xi$-units this reads
\begin{equation*}
|\partial U_n(M^{\cdot}(e^{i\xi A}\mathbf{U}))-1|<\alpha_0\xi^2,\qquad
|J_n(M^{\cdot}(e^{i\xi A}\mathbf{U}))|<\alpha_0(L^n\xi)^2 .
\end{equation*}
These are the bounds \cite[(1.16)]{B-CMP109} at $\alpha_0$ for every $n\le j$. The first of them is
the uniform fine-scale member.
\end{enumerate}

The smallness required of $\alpha_2$ here is $\alpha_2\le\min\{\alpha_1/8,\alpha_0/C_{IV}\}$. It does
not depend on $n$, $j$, $k$, the depth $j-n$, $M$ or $X$. No inequality of the lemma compares
$\alpha_0$ with $\alpha_1$. The letter $M$ occurs only in the floors and, through
$\frac12c_{12}LMB\,\alpha_0$, in the ceiling on $\alpha_0$.

The constants $C_{IV}$, $C_K$, $C_{14}$, $c_\sharp$, $\bar\varepsilon_1$, $c_Q$, $C_1$, $C_7$ do not
depend on $M$. Hence every constant, ceiling and floor of the lemma is uniform in $M=L^{m'}$ over
all $M$ satisfying the floors in \eqref{4.5:eq-uniform-increment-b}. The lemma may therefore be
applied at another
admissible cube letter. An example is $M/L$, with $LB$ in place of $B$ and
$(\mu\alpha_0,\mu\alpha_1,\mu\alpha_2)$ in place of $(\alpha_0,\alpha_1,\alpha_2)$ for a number
$\mu$ fixed by that application. The floors of that application, read at the original $M$, are
$\delta_0M\ge32L$ and $2LR_1M_1\le M$.
\end{lemma}

\begin{proof}
We prove clause (ii), the first member of the curvature bound, clause (i), and the half of
(iv) concerning $U$. The step marked $(\ast)$ is the increment bound
\eqref{4.5:eq-uniform-increment-a}; from it we then derive the bounds at $\alpha_0$.

\emph{Clause (ii), the potential.} We have $e^{i\xi A''}=e^{i\xi A}e^{i\xi A'}$. By the
Baker--Campbell--Hausdorff formula,
\begin{equation*}
A''=A+A'+\tfrac{i\xi}{2}[A,A']+\cdots .
\end{equation*}
By hypothesis $|A|<\alpha_2$. Condition (ii) for $\mathbf{U}$ with $\frac12\alpha_1$ gives
$|A'|<\frac12\alpha_1$. Since $\xi\le1$, the third ceiling in \eqref{4.5:eq-uniform-increment-b}
gives $\xi(|A|+|A'|)<\alpha_2+\frac12\alpha_1\le c_\sharp\le\frac14$, so the series converges; its
terms beyond the first commutator are bounded by the lemma on the Baker--Campbell--Hausdorff
remainder of this chapter. Summing the series yields
\begin{equation*}
|A''|<\alpha_2+\tfrac12\alpha_1+\xi\alpha_1\alpha_2\le\tfrac12\alpha_1+2\alpha_2 .
\end{equation*}
The last inequality holds because $\xi\le1$ and $\alpha_1\le1$ give $\xi\alpha_1\alpha_2\le\alpha_2$.

\emph{Clause (ii), the gradient.} The hypothesis bounds $\nabla^\xi_{\mathbf{U}}A$, whereas (ii)
asks for the gradient taken with $U$. On a bond $b$, with $\mathbf{U}(b)=e^{i\xi A'(b)}U(b)$,
\begin{equation*}
\nabla^\xi_UA-\nabla^\xi_{\mathbf{U}}A
=\xi^{-1}\bigl(\operatorname{Ad}U(b)-\operatorname{Ad}\mathbf{U}(b)\bigr)A(b_+).
\end{equation*}
The right side has size at most $2|A'||A|$. By $|A'|<\frac12\alpha_1$ and $|A|<\alpha_2$ this is at
most $\alpha_1\alpha_2$. Hence $|\nabla^\xi_UA|<\alpha_2+\alpha_1\alpha_2$.

Condition (ii) for $\mathbf{U}$ gives $|\nabla^\xi_UA'|<\frac12\alpha_1$. Differentiating the series
for $A''$ term by term, with the same bound on the terms beyond the first commutator, gives
\begin{equation*}
\begin{aligned}
|\nabla^\xi_UA''|
&\le|\nabla^\xi_UA|+|\nabla^\xi_UA'|
+\xi\bigl(|\nabla A||A'|+|A||\nabla A'|\bigr)+\cdots\\
&<\alpha_2+\alpha_1\alpha_2+\tfrac12\alpha_1+2\alpha_1\alpha_2 .
\end{aligned}
\end{equation*}
The right side equals $\frac12\alpha_1+\alpha_2(1+3\alpha_1)$. Since $\alpha_1\le1$, it is at most
$\frac12\alpha_1+4\alpha_2$.

Finally, $\alpha_2\le\alpha_1/8$ gives $4\alpha_2\le\frac12\alpha_1$. Hence
$\frac12\alpha_1+4\alpha_2\le\alpha_1$, which proves clause (ii).

\emph{The first member of the curvature bound.} Let $p$ be a plaquette. Then
$\partial(e^{i\xi A}\mathbf{U})(p)$ is the plaquette holonomy of $\mathbf{U}$ dressed by four
conjugated factors of the form $e^{i\xi A}$. Expanding these factors gives
\begin{equation*}
|\partial(e^{i\xi A}\mathbf{U})(p)-\partial\mathbf{U}(p)|
\le\xi^2\bigl(2d\sup|\nabla^\xi_{\mathbf{U}}A|+2|A|^2\bigr)\bigl(1+4\xi(\alpha_1+\alpha_2)\bigr)^2 .
\end{equation*}
We bound the two factors on the right.
\begin{itemize}
\item By the third ceiling in \eqref{4.5:eq-uniform-increment-b}, $\frac12\alpha_1+\alpha_2\le c_\sharp$.
Therefore $\alpha_1+\alpha_2\le2c_\sharp\le\frac12$. Since $\xi\le1$, this gives
$4\xi(\alpha_1+\alpha_2)\le2$, so $(1+4\xi(\alpha_1+\alpha_2))^2\le9$. This holds at every level,
$j=0$ included.
\item By hypothesis $|\nabla^\xi_{\mathbf{U}}A|<\alpha_2$ and $|A|<\alpha_2$. Moreover
$\alpha_2\le\alpha_1/8\le1$. Hence $2d\alpha_2+2\alpha_2^2\le(2d+2)\alpha_2$.
\end{itemize}
Together,
\begin{equation*}
|\partial(e^{i\xi A}\mathbf{U})(p)-\partial\mathbf{U}(p)|\le9(2d+2)\alpha_2\xi^2=C_{14}\alpha_2\xi^2 .
\end{equation*}
The first member of \cite[(1.14)]{B-CMP109} for $\mathbf{U}$ with $\frac12\alpha_0$ gives
$|\partial\mathbf{U}(p)-1|<\frac12\alpha_0\xi^2$. Adding,
$|\partial(e^{i\xi A}\mathbf{U})(p)-1|<(\frac12\alpha_0+C_{14}\alpha_2)\xi^2$.

By $\alpha_2\le\alpha_0/C_{IV}$ and $C_{IV}\ge4C_{14}$ we have $\alpha_2\le\alpha_0/(4C_{14})$.
Hence $C_{14}\alpha_2\le\frac14\alpha_0$, and so $(\frac12\alpha_0+C_{14}\alpha_2)\xi^2\le\alpha_0\xi^2$.

\emph{Clause (i) and the half of (iv) for $U$.} The twist $e^{i\xi A}$ changes only the complex
factor: $e^{i\xi A}\mathbf{U}=e^{i\xi A''}U$ with the same $G$-valued $U$. Therefore
\cite[(1.11)--(1.12)]{B-CMP109}, which concern $U$, are unchanged.

For the same reason the bounds \cite[(1.16)]{B-CMP109} for $U$ are unchanged. For every sub-tower
$\{\Omega_0,\dots,\Omega_n\}$ of $\mathbf{\Omega}$ they hold by hypothesis with $\frac12\alpha_0$,
hence with $\alpha_0$.

$(\ast)$ \emph{The increment bound.} For every $n\in\{1,\dots,j\}$ and every $x\in X^{\sim2}$ the
increment bound \eqref{4.5:eq-uniform-increment-a} holds. A proof places $x$ in the top layer
$\Omega_n$ of the sub-tower $\{\Omega_0,\dots,\Omega_n\}$, recentres the twisted data exactly in
level-$n$ units so that the contributions of nearby points carry the factor $\lambda^2$, reads the
kernel estimate at that top-layer point, and controls the contributions of distant points by
exponential decay in level-$n$ units under the floor $\delta_0M\ge32$; this argument is not carried
out in the present proof. The collar bound $R_1M_1$ for the maximal tower, used in the statement,
belongs to the same step.

\emph{From \eqref{4.5:eq-uniform-increment-a} to the bounds at $\alpha_0$.} Fix $n\in\{1,\dots,j\}$
and $x\in X^{\sim2}$. Since $\xi/\eta_n=\lambda$ and $L^n\xi=\lambda$, the hypothesis bounds for
$\mathbf{U}$ become
\begin{equation*}
|\Phi_1(\mathbf{U})|=\eta_n^{-2}|\partial U_n(M^{\cdot}(\mathbf{U}))-1|<\eta_n^{-2}\,\tfrac12\alpha_0\xi^2
=\tfrac12\alpha_0\lambda^2,
\end{equation*}
and
\begin{equation*}
|\Phi_2(\mathbf{U})|=|J_n(M^{\cdot}(\mathbf{U}))|<\tfrac12\alpha_0(L^n\xi)^2=\tfrac12\alpha_0\lambda^2 .
\end{equation*}
By \eqref{4.5:eq-uniform-increment-a} and $\alpha_2\le\alpha_0/C_{IV}$, for $i=1,2$
\begin{equation*}
|\Phi_i(e^{i\xi A}\mathbf{U})|
\le|\Phi_i(\mathbf{U})|+\tfrac12C_{IV}\alpha_2\lambda^2
<\tfrac12\alpha_0\lambda^2+\tfrac12\alpha_0\lambda^2=\alpha_0\lambda^2 .
\end{equation*}
Multiplying the case $i=1$ by $\eta_n^2$ and using $\eta_n^2\lambda^2=\xi^2$ gives
$|\partial U_n(M^{\cdot}(e^{i\xi A}\mathbf{U}))-1|<\alpha_0\xi^2$. The case $i=2$ is
$|J_n(M^{\cdot}(e^{i\xi A}\mathbf{U}))|<\alpha_0(L^n\xi)^2$.

These are the bounds \cite[(1.16)]{B-CMP109} at $\alpha_0$ for $e^{i\xi A}\mathbf{U}$, for every
$n\le j$ and the same tower $\mathbf{\Omega}$. Neither step uses a condition on $\alpha_2$ that
involves $n$, $j$ or $j-n$.
\end{proof}

\noindent\textit{Units, tower geometry and local gauge.}\label{4.5:prf-increment-geometry}
We begin the proof of Lemma~\ref{4.5:lem-uniform-increment}. In this first part we fix the units in
which the estimates are carried out, collect the properties of the tower that are used, and reduce
to a configuration given by a single potential near the point under consideration.

As in \cite[p.~257]{B-CMP109}, the term of index $j$ lives on the continuous space scaled so that it
corresponds to the lattice of spacing $\xi=L^{-j}$. This space is decomposed into the lattice of closed
cubes of side $M$, a power of $L$. We call lengths measured in this way level-$j$ units. In them:
\begin{itemize}
\item $X$ is a union of cubes of side $M$;
\item the cubes of \cite[(1.12)]{B-CMP109} have side $c_{12}LM$;
\item for $n\le j$, the blocks of level $n$ have side $\lambda:=L^{n-j}=L^{n}\xi\le 1$.
\end{itemize}

Fix $n\le j$ and $x\in X^{\sim 2}$, the set obtained from $X$ by removing two boundary layers of
cubes of side $M$ (the set $\tilde X^{-2}$ of \cite[(1.16)]{B-CMP109}). We now measure lengths in
level-$n$ units, that is, we multiply all lengths by $\lambda^{-1}$. In these units:
\begin{itemize}
\item the fine spacing is $\eta_n:=\xi/\lambda=L^{-n}$, which we write $\eta$ in this proof;
\item for $p\le n$, the block of level $p$ has side $\ell_p:=L^{p-n}$;
\item the cubes of side $M$ have side $M/\lambda$.
\end{itemize}
Potentials scale with $\lambda$ and their covariant derivatives with $\lambda^{2}$. Put
$\hat A:=\lambda A$ and $\hat A':=\lambda A'$. Then
\begin{align*}
&e^{i\xi A}=e^{i\eta\hat A},\qquad |\hat A|<\alpha_2\lambda,\qquad
|\nabla^{\eta}_{\mathbf{U}}\hat A|=\lambda^{2}|\nabla^{\xi}_{\mathbf{U}}A|<\alpha_2\lambda^{2},\\
&|\hat A'|<\tfrac12\alpha_1\lambda,\qquad |\nabla^{\eta}\hat A'|<\tfrac12\alpha_1\lambda^{2},\qquad
|\partial\mathbf{U}-1|,\ |\partial U-1|<\tfrac12\alpha_0\lambda^{2}\eta^{2}.
\end{align*}
These relations are justified as follows.
\begin{itemize}
\item The first identity holds because $\xi A=\eta\lambda A=\eta\hat A$.
\item A difference quotient with spacing $\eta=\xi/\lambda$ equals $\lambda$ times the corresponding
quotient with spacing $\xi$. Hence $\nabla^{\eta}_{\mathbf{U}}\hat A=\lambda\nabla^{\xi}_{\mathbf{U}}(\lambda A)$,
which gives the equality in the first line.
\item The bounds on $\hat A$ are the hypotheses $|A|,|\nabla^{\xi}_{\mathbf{U}}A|<\alpha_2$ of the lemma.
\item The bounds on $\hat A'$ and on the plaquettes are the hypotheses of the lemma with constants
$(\tfrac12\alpha_0,\tfrac12\alpha_1)$. For the plaquettes this uses $\xi^{2}=\lambda^{2}\eta^{2}$.
\end{itemize}

Condition \cite[(1.16)]{B-CMP109} takes the following form in level-$n$ units. For a plaquette $p$
based at $x$ and a bond $b$ with initial point $b_-=x$, put
\begin{equation}\label{4.5:eq-increment-geometry-1}
\Phi_1(\mathbf{U}):=\eta^{-2}\bigl(\partial U_n(M^{\cdot}(\mathbf{U}))(p)-1\bigr),\qquad
\Phi_2(\mathbf{U}):=J_n(M^{\cdot}(\mathbf{U}))(b).
\end{equation}
Here $J_n$ is defined by the formula \cite[(1.8)]{B-CMP109} with
$U_n(M^{\cdot}(\mathbf{U}))$ in place of $U_j$ and $L^{-n}$ in place of $\xi$; it is therefore already
expressed in level-$n$ units. Condition \cite[(1.16)]{B-CMP109} at $\alpha_0$ asks that
$|\Phi_i|<\alpha_0\lambda^{2}$ for $i=1,2$. Every plaquette and every bond of $X^{\sim 2}$ is reached
in this way, as $x$ ranges over $X^{\sim 2}$. Because $p$ and $b$ are based at $x$, the gauge
invariance used in the increment estimate below is exact.

\emph{The tower.} We first recall the conditions imposed on the tower. We write $\mathfrak{M}$ for the
block size of \cite{B-CMP99a} (denoted $M_1$ there), so that $\mathfrak{M}\le M_1$. In the conventions of
\cite[(1.3)--(1.6)]{B-CMP99a}, read with $j$ levels, spacing $\xi$ and gap $R_1$, the tower satisfies
$\Omega_0=X$, each $\Omega_p$ is a union of cubes of side $\mathfrak{M}L^{p}\xi$, and
\begin{equation*}
(L^{p}\xi)^{-1}\operatorname{dist}(\Omega_p^{c},\Omega_{p+1})>R_1\mathfrak{M}.
\end{equation*}
For the maximal tower of $X$, $\Omega_{p+1}$ is the largest such union.

Next we translate these conditions into level-$n$ units, where they read as follows:
\begin{itemize}
\item $\Omega_p$ is a union of cubes of side $\mathfrak{M}\ell_p$;
\item $\Omega_{p+1}$ has distance greater than $R_1\mathfrak{M}\ell_p$ from $\Omega_p^{c}$.
\end{itemize}

For the maximal tower of $X$ we claim that $\Omega_{p+1}$ contains every point $z$ of $\Omega_p$ at
distance greater than
$(R_1+dL)M_1\ell_p$ from $\Omega_p^{c}$. Let $z$ be such a point, and take the cube of side
$\mathfrak{M}\ell_{p+1}=L\mathfrak{M}\ell_p$ of the level-$(p+1)$ decomposition that contains $z$.
Its diameter is at most $dL\mathfrak{M}\ell_p$. Since $\mathfrak{M}\le M_1$, its distance from
$\Omega_p^{c}$ therefore exceeds
\begin{equation*}
(R_1+dL)M_1\ell_p-dL\mathfrak{M}\ell_p\ge R_1M_1\ell_p\ge R_1\mathfrak{M}\ell_p .
\end{equation*}
By maximality this cube lies in $\Omega_{p+1}$, which proves the claim.

The depth of a point of $X$ is its distance from $X^{c}$, and the depth of a subset of $X$ is the
largest depth of its points. By induction on $p$ and the triangle inequality, every point of
$X\setminus\Omega_{p+1}$ lies within $\sum_{q\le p}(R_1+dL)M_1\ell_q$ of $X^{c}$. Indeed, a point of
$\Omega_p\setminus\Omega_{p+1}$ lies within $(R_1+dL)M_1\ell_p$ of $\Omega_p^{c}$, and, since
$\Omega_p\subset X$, every point of $\Omega_p^{c}$ lies in $X^{c}$ or in $X\setminus\Omega_p$, to which
the induction hypothesis applies. Hence the collar $X\setminus\Omega_n$
satisfies
\begin{equation}\label{4.5:eq-increment-geometry-2}
\operatorname{depth}(X\setminus\Omega_n)\le\sum_{p<n}(R_1+dL)M_1\ell_p
\le\frac{2(R_1+dL)M_1}{L}\le R_1M_1\le\tfrac12 M\le\frac{M}{2\lambda}.
\end{equation}
The steps of \eqref{4.5:eq-increment-geometry-2} are justified as follows.
\begin{itemize}
\item The second inequality holds because
\begin{equation*}
\sum_{p<n}\ell_p=\sum_{q=1}^{n}L^{-q}<\frac{1}{L-1}\le\frac{2}{L}.
\end{equation*}
\item The third inequality is equivalent to $2dL\le(L-2)R_1$. The floors of the lemma give
$L\ge 4d$, hence $L\ge 4$, and $R_1\ge L\ge 4d$. The latter holds because $R_1$ is a sufficiently
large power of $L$ in the construction of \cite{B-CMP99a}. Therefore
$(L-2)R_1\ge 4d(L-2)\ge 2dL$.
\item The fourth inequality is the floor $2R_1M_1\le M$.
\item The last inequality holds because $\lambda\le 1$.
\end{itemize}
This is the instance, for the maximal tower, of the collar hypothesis in condition (iv) of
Lemma~\ref{4.5:lem-uniform-increment}.

For a general tower $\mathbf{\Omega}$ admissible in the sense of \cite[(1)]{B-CMP102b}, as allowed
in condition (iv), the bound $\operatorname{depth}(X\setminus\Omega_n)\le\tfrac12M$ in level-$j$ units,
that is $M/(2\lambda)$ in level-$n$ units, is the hypothesis itself. Of the tower only two things are
used: its collar bound, in statement (a) below, and its admissibility, namely block $\mathfrak{M}$,
gap at least $R_1\mathfrak{M}$, and bond sets consisting of the bonds with at least one end-point in
the corresponding set. Admissibility is all that \cite[Theorem~1, Proposition~9]{B-CMP102b} and the
corollary on the kernels of the hierarchy, all of which are read for towers as in
\cite[(1)]{B-CMP102b}, use of $\mathbf{\Omega}$ in the next part of the proof.
Nothing else of the tower enters below. By contrast, $X^{\sim 2}$ lies at depth at least $2M$ in
level-$j$ units, that is, at depth at least $2M/\lambda\ge 2M$ in level-$n$ units.

Consider the truncated tower $\{\Omega_0,\dots,\Omega_n\}$. Its last layer $\Lambda_n$ is the whole of
the level-$n$ lattice of $\Omega_n$. Hence the $n$-fold blocking of \cite[(1.6)]{B-CMP99a}, which is
the identity at level $0$, carries $\Lambda_n$ onto $\Omega_n$; in particular $\Omega_n$ is the top layer
of the truncated tower.

Let $y_x$ be the base point of the block $\Delta(y_x)$ of level $n$ that contains $x$. Put
$x_0:=y_x$, so that $|x-x_0|_\infty\le 1$. Put $r:=M/\lambda$, and let $Q$ be the cube with centre
$x_0$ and half-side $r$; the letters $r$ and $Q$ have this meaning in this proof only. We work
under the floor $2R_1M_1\le M$ for the maximal tower, and under the collar bound of condition (iv)
for a general tower. We prove two statements.
\begin{enumerate}
\item[(a)] Every point $z$ with $|z-x_0|_\infty\le r/2+2$ lies in the top layer $\Omega_n$.

Such a point satisfies $|z-x|_\infty\le r/2+3$. Its depth is therefore at least
\begin{equation*}
\frac{2M}{\lambda}-\frac{M}{2\lambda}-3\ge\frac{M}{\lambda}>\frac{M}{2\lambda}
\ge\operatorname{depth}(X\setminus\Omega_n).
\end{equation*}
The first inequality uses $M\ge 6\lambda$, which holds for every tower by the standing floor
$2R_1M_1\le M$ of the lemma:
\begin{equation*}
M\ge 2R_1M_1\ge 8d>6\ge 6\lambda .
\end{equation*}
For the maximal tower, the comparison with the collar also follows from
\eqref{4.5:eq-increment-geometry-2} and
\begin{equation*}
\frac{M}{\lambda}\ge M\ge 2R_1M_1>R_1M_1 .
\end{equation*}
Hence $z$ lies in the
top layer $\Omega_n$. In particular $x$, the block $\Delta(y_x)$ and the neighbourhood of $x$ used
below lie in $\Omega_n$. The local length unit $\ell(y)$ of the hierarchy, equal to $\ell_p$ on the
layer of level $p$, is therefore the top one at $y_x$, $\ell(y_x)=\ell_n=1$, and the kernels of the
corollary on the kernels of the hierarchy, read at $y_x$, carry no power of $\ell(y_x)$.
\item[(b)] $Q\subset X$, and there is a $G$-valued gauge map $u$ on $Q$ with $U^{u}=e^{i\xi A_0}$
on $Q$ and $|\hat A_0|\le c_{\sharp}\lambda$, where $\hat A_0:=\lambda A_0$.

The depth of $x_0$ is at least $2M/\lambda-1>r$, so $Q\subset X$. The cube $Q$ has side $2M$ in
level-$j$ units, and $2M\le c_{12}LM$ by the floor $c_{12}L\ge 2$. By the reading convention for \cite[(1.12)]{B-CMP109}, that condition is read for every
cube contained in $X$ of side at most
$c_{12}LM$. Applied at $\tfrac12\alpha_0$, it gives a single $G$-valued gauge map $u$ on $Q$ with
$U^{u}=e^{i\xi A_0}$ on $Q$ and
\begin{equation*}
|A_0|,\ |\nabla^{\xi}A_0|<c_{12}LMB\cdot\tfrac12\alpha_0 ,
\end{equation*}
where $B$ is the constant of \cite[(1.12)]{B-CMP109}. In level-$n$ units this reads
\begin{equation*}
|\hat A_0|\le c_{12}LMB\cdot\tfrac12\alpha_0\cdot\lambda\le c_{\sharp}\lambda ,
\end{equation*}
by the third ceiling of the lemma, $c_{12}LMB\cdot\tfrac12\alpha_0+\tfrac12\alpha_1+\alpha_2\le c_{\sharp}$.
\end{enumerate}

Statement (a) is the point at which the floor $2R_1M_1\le M$, respectively the collar hypothesis,
enters the proof. What it secures is that the top length unit is kept, so that no power of
$\ell(y_x)$ is lost in the kernels read at $y_x$. The two powers $\lambda^{2}$ in the conclusion are
obtained in the increment estimate below. At this point a direct Cauchy bound gives only
$C'\alpha_2\lambda$, with a constant $C'$, which is one power short. The second power is supplied by
the exact recentring field $\nu$ introduced there, together with the decay of the kernels.

\emph{Local gauge.} Extend $u$ by $1$ off $Q$. We check that every hypothesis of the lemma, and both
$\Phi_1$ and $\Phi_2$, are isometric under $G$-valued gauge maps.
\begin{itemize}
\item Plaquette variables are conjugated.
\item The block averages are covariant, $\overline{U^{u}}=(\bar U)^{u}$, by \cite[(11)]{B-CMP98}.
Together with \cite[(181)]{B-CMP102b} at a $G$-valued map, this gives
$\Phi_1\mapsto\operatorname{Ad}_{\bar u(x)}\Phi_1$. Here $\bar u$ denotes the gauge map by which $u$
acts on the averages of level $n$.
\item The current satisfies $J(U^{u})=R(u)J(U)$, where $R(u)$ denotes the action of the gauge map $u$
on $\mathfrak{g}$-valued bond fields in \cite[(1.11)]{B-CMP99a}, by the covariance of $D^{\eta *}$
there. Hence $\Phi_2\mapsto R(\bar u(b_-))\Phi_2$.
\item For the potential of hypothesis (ii), $A'\mapsto R(u)A'$, and
$\nabla_{U^{u}}R(u)=R(u)\nabla_{U}$.
\end{itemize}
Since $\operatorname{Ad}$ and $R$ of a $G$-valued element are isometries, all the quantities
entering the hypotheses and the conclusion of the lemma keep their norms. We may therefore, and do,
assume that $U=e^{i\eta\hat A_0}$ on $Q$.

\begin{proof}[Proof of Lemma~\ref{4.5:lem-uniform-increment}, continued: real base, complex datum and
the derivative kernel]\label{4.5:prf-increment-kernel}
We keep the level $n\le j$, the site $x\in X^{\sim 2}$, the level-$n$ units and the notation of
the geometric part of this proof. The fine spacing is
$\eta=\eta_n=L^{-n}$, and the level-$q$ block has side $\ell_q=L^{q-n}\le 1$ for $q\le n$. We
write $\hat A=\lambda A$ and $\hat A'=\lambda A'$, with $|\hat A|<\alpha_2\lambda$ and
$|\hat A'|<\tfrac12\alpha_1\lambda$. The two quantities to be controlled are
\begin{equation*}
\Phi_1(\mathbf{U})=\eta^{-2}\bigl(\partial U_n(M^{\cdot}\mathbf{U})(p)-1\bigr),
\qquad
\Phi_2(\mathbf{U})=J_n(M^{\cdot}\mathbf{U})(b),
\end{equation*}
for a plaquette $p$ based at $x$ and a bond $b$ with $b_-=x$. In this part we prove the
representation \eqref{4.5:eq-increment-kernel-a} of their derivatives along the path from
$\mathbf{U}$ to $e^{i\xi A}\mathbf{U}$.

\smallskip
\noindent\textit{The real base.}
Put $V_0:=M^{\cdot}(U)$, the tower of block averages of the $G$-valued configuration $U$. Then
$V_0$ is real, that is, $G$-valued. In \cite{B-CMP102b} the top level, written $k$ there, is our
$n$, and we write $q$ for the running level; in \cite{B-CMP98} the number of averaging steps,
written $k$ there, is written $q$ here. We apply
\cite[Proposition~2]{B-CMP98} at the $G$-valued configuration $U$; that proposition is local,
and it says the following. If $U$ satisfies
\cite[(52)]{B-CMP98} with $\alpha_0\le c_2$, then for every $q\le n$ the $q$-fold average obeys
\begin{equation*}
|\bar U^{q}(\partial p')-1|<\alpha_0+2C_0\alpha_0^2<2\alpha_0 .
\end{equation*}
Here $c_2$ and $C_0$ are the constants of that proposition.

Consequently $V_0$ satisfies the curvature condition \cite[(7)]{B-CMP102b}, namely
$|(\partial V)(p')-1|<\varepsilon_1^{(102)}$ for the plaquettes $p'$ of $\mathfrak{B}_n$, where
$\varepsilon_1^{(102)}$ is our name for Ba{\l}aban's curvature parameter in that condition (the
superscript refers to \cite{B-CMP102b}). For
the plaquettes joining two layers of the hierarchy, that condition is supplied by the
observation in \cite{B-CMP102b} that, by \cite[Proposition~2]{B-CMP98}, the averaged
configuration satisfies that condition with $\varepsilon_1^{(102)}=O(\varepsilon_0)$.

The curvature of the datum on a plaquette $p'$ of layer $q$ obeys
\begin{equation}\label{4.5:eq-increment-kernel-1}
|(\partial V_0)(p')-1|\le C_7\cdot\tfrac12\alpha_0\lambda^2\ell_q^2\le C_7\alpha_0\le\bar\varepsilon_1 ,
\end{equation}
where $C_7$ is a constant depending on $d$ and $L$ only, and we used $\lambda\le 1$ and
$\ell_q\le 1$. Condition \cite[(7)]{B-CMP102b} is an upper
bound. It therefore holds with $\varepsilon_1^{(102)}:=\bar\varepsilon_1$, the constant depending
on $d$ and $L$ only that is fixed in \eqref{4.5:eq-geometry-conventions-a}. We stress that this
value is a fixed constant; it is not the curvature of the datum.

The hypothesis that the data form a tower of averages holds automatically, because $V_0$
consists of the averages of one field. The tower in use is the tower $\mathbf{\Omega}$ of the
lemma. It is admissible in the sense of \cite[(1)]{B-CMP102b}, and this admissibility is all
that \cite[Theorem~1]{B-CMP102b}, \cite[Proposition~9]{B-CMP102b} and the corollary on
Ba{\l}aban's background function $\mathcal{H}$ of \cite[Proposition~9]{B-CMP102b}, recalled at
\eqref{4.5:eq-increment-kernel-3} below, require of it.

\smallskip
\noindent\textit{The real minimiser.}
Put $U_0:=U_n(V_0)$. By \cite[Theorem~1]{B-CMP102b}, every configuration $V$ satisfying
condition \cite[(7)]{B-CMP102b} with $\varepsilon_1^{(102)}\le a_1^{(102)}$ has a minimal orbit.
This orbit is the unique critical orbit in the space \cite[(6)]{B-CMP102b} provided
$B_3\varepsilon_1^{(102)}\le\varepsilon_0$ and $\varepsilon_0\le a_0^{(102)}$; here
$a_0^{(102)},a_1^{(102)}$ are the constants $a_0,a_1$ of that theorem.

Both conditions hold here. First, $\bar\varepsilon_1\le a_1^{(102)}$, this being the first entry
of the minimum in \eqref{4.5:eq-geometry-conventions-a}. Second, the constants of that theorem
satisfy $B_3a_1^{(102)}\le a_0^{(102)}$, so we may take $\varepsilon_0:=B_3\bar\varepsilon_1$.
Hence $U_0$ is the real
minimiser of \cite[Theorem~1]{B-CMP102b}. By the uniqueness of the critical orbit, $U_0$ does
not depend on which admissible value of $\varepsilon_1^{(102)}$ is used to name it.

\smallskip
\noindent\textit{The complex datum.}
For $t\in[0,1]$ put $\mathbf{U}_t:=e^{it\eta\hat A}\mathbf{U}$. Thus $\mathbf{U}_0=\mathbf{U}$,
and $\mathbf{U}_1=e^{i\xi A}\mathbf{U}$ because $e^{i\xi A}=e^{i\eta\hat A}$.

Three facts from \cite{B-CMP98} govern the averages of such twisted configurations.
\begin{itemize}
\item By \cite[Proposition~6]{B-CMP98}, for every $q\le n$ the $q$-fold average of a
configuration twisted by
$e^{i\eta A'}$ is an analytic function of $A'$, a field with values in the complexified algebra
$\mathfrak{g}^{c}$, for $|A'|<\alpha_1^{(98)}$; here $\alpha_1^{(98)}$ denotes the radius of that
proposition. Moreover, the average of the twisted configuration multiplied by the inverse
of the average of the untwisted one differs from $1$ by strictly less than
$O(1)\alpha_1^{(98)}$ \cite[(164)]{B-CMP98}.
\item The sizes level by level are given by \cite[(159), (161)]{B-CMP98}: the logarithm of the
level-$q$ factor of the twisted average has modulus less than
$2\alpha_1^{(98)}L^q\eta=2\alpha_1^{(98)}\ell_q$.
\item The accumulated factor $v_n$ (written $v_k$ in \cite{B-CMP98}) satisfies
\cite[(163)]{B-CMP98}:
\begin{equation*}
|v_n(x)-1|<O(1)\alpha_1^{(98)}\sum_{q=0}^{n-1}L^{q+1}\eta\le O(1)\alpha_1^{(98)} .
\end{equation*}
\end{itemize}

Put $V_t:=M^{\cdot}(\mathbf{U}_t)$. With the bounds on $\hat A'$ and $\hat A$ recalled above,
these three facts give $V_t=e^{iB_t}V_0$, where on $\Lambda_q$ the exponent obeys
\begin{equation}\label{4.5:eq-increment-kernel-2}
|B_t|\le c_Q\bigl(\tfrac12\alpha_1+\alpha_2\bigr)\lambda\ell_q\le c_Q(\alpha_1+\alpha_2)
\le\tfrac12C_1\bar\varepsilon_1 ,
\end{equation}
where $c_Q$ is the constant depending on $d$ and $L$ only that enters $\mathfrak{r}_g$ in
\eqref{4.5:eq-geometry-conventions-a}, the last inequality is the second ceiling of
\eqref{4.5:eq-uniform-increment-b}, and
$C_1$ is the constant of the polydisc \cite[(172)]{B-CMP102b} recalled next.

Now recall the polydisc \cite[(172)]{B-CMP102b}: a configuration with
$|V'-1|<C_1\varepsilon_1^{(102)}$ has the form $V'=e^{iB'}$, so that the logarithms $B'$ range
over the polydisc $|B'|<2C_1\varepsilon_1^{(102)}$ on $\mathfrak{B}_n$. Read at
$\varepsilon_1^{(102)}=\bar\varepsilon_1$, \eqref{4.5:eq-increment-kernel-2} places $B_t$ in
this polydisc of logarithms with room to spare, since
$\tfrac12C_1\bar\varepsilon_1<2C_1\bar\varepsilon_1$.

The average $e^{iB_t}$ of the twist is kept as a complex datum over the real base $V_0$. No part
of $\hat A'$ is treated as a gauge transformation.

\smallskip
\noindent\textit{The minimiser at the complex datum and the gauge map.}
By \cite[Proposition~9]{B-CMP102b}, the minimal configuration $U_n(V)=U_n(V'V_0)$ in the axial
gauge extends to an analytic function of $G^{c}$-valued small configurations $V'$. Written
relative to $U_0=U_n(V_0)$ and transformed to the Landau gauge, this extension is, by
definition, $e^{i\eta\mathcal{H}(B)}U_0$, where $B=\tfrac1i\log V'$. Returning to the axial
gauge by the map $g$ described below and applying this at $V'=V_tV_0^{-1}$, so that $B=B_t$, we
obtain
\begin{equation}\label{4.5:eq-increment-kernel-3}
U_n(V_t)=\bigl(e^{i\eta\mathcal{H}(B_t)}U_0\bigr)^{g(B_t)} .
\end{equation}
Here $g$ is the map from the Landau gauge to the axial gauge. It is normalised by the averaging
conditions of \cite[(1.29)]{B-CMP99a}, which read, as in \cite[p.~281]{B-CMP102b},
\begin{equation*}
\overline{R_0u}^{\,q}=1\ \text{on }\Lambda_q,\quad q=0,1,\dots,n\qquad
(\text{thus }u=1\text{ on }\Lambda_0).
\end{equation*}
Such a gauge transformation is determined uniquely by the configuration and is given by an
explicit formula; this is stated in \cite{B-CMP102b}.

On the block $\Delta(y_x)$, $g(B_t)$ is an explicit local function of $\mathcal{H}(B_t)$ there.
It is built from the nested transports, down the contour of $x$, of the averages of
$e^{i\eta\mathcal{H}}$, multiplied by the top value, and from the inverse of the iterated block
mean \cite[(86)--(87)]{B-CMP98} of those transports.

Its size is controlled by the logarithm bound for the averaging gauge map at an arbitrary axial
base. The top block has
side $1$ in level-$n$ units, and that bound gives
\begin{equation}\label{4.5:eq-increment-kernel-4}
|\log g|\le 4dc_Q\Bigl(1+O\bigl(dc_Q\sup_{\Delta(y_x)}|\mathcal{H}|\bigr)\Bigr)
\sup_{\Delta(y_x)}|\mathcal{H}| .
\end{equation}

Both premises of that bound hold at $\mathcal{H}=\mathcal{H}(B_t)$. To see this, sum the value and
gradient bounds of the corollary on $\mathcal{H}$, with the constant $C_H$ of that corollary,
over the sources, using
\cite[Lemma~2.1]{B-CMP96} with its constant $c_1$ and the site multiplicity $3dL$ of the site
count of the axial frame. Together with \eqref{4.5:eq-increment-kernel-2} this gives
\begin{equation}\label{4.5:eq-increment-kernel-5}
\sup_{\Delta(y_x)}|\mathcal{H}(B_t)|\le 3dLC_Hc_1\sup|B_t|
\le 3dLC_Hc_1\cdot\tfrac12C_1\bar\varepsilon_1\le\mathfrak{r}_g .
\end{equation}
The last step uses the fifth entry $2\mathfrak{r}_g/(3dLC_Hc_1C_1)$ of the minimum defining
$\bar\varepsilon_1$ in \eqref{4.5:eq-geometry-conventions-a}, where
$\mathfrak{r}_g=\min\{1,1/(24dc_Q),L\alpha_1^{(98)}\}$. In particular the two premises follow:
\begin{itemize}
\item $dc_Q\sup|\mathcal{H}|\le 1/24$, which is what the Baker--Campbell--Hausdorff step of the
bound needs;
\item $\sup|\mathcal{H}|/L\le\alpha_1^{(98)}$, the radius of
\cite[Proposition~6]{B-CMP98}, which is what that proposition needs in the averaged legs.
\end{itemize}

Moreover $g(0)=1$. From \eqref{4.5:eq-increment-kernel-4} and \eqref{4.5:eq-increment-kernel-5},
\begin{equation}\label{4.5:eq-increment-kernel-6}
\|\operatorname{Ad}_g^{\pm1}\|\le e^{8dc_Q(1+O(1))\sup|\mathcal{H}|}=O(1),
\end{equation}
and below only the fact that this quantity is $O(1)$ is used.

\smallskip
\noindent\textit{The two quantities along the path.}
We express $\Phi_1$ and $\Phi_2$ through \eqref{4.5:eq-increment-kernel-3}.

For $\Phi_1$: the plaquette $p$ is based at $x$, so a gauge transformation $g$ conjugates its
holonomy by $g(x)$.

For $\Phi_2$: the current transforms covariantly, $J(U^u)=R(u)J(U)$, by \cite[(1.11)]{B-CMP99a}.
The remainder $F$ of \cite[(3.11)]{B-CMP109} (Lemma~\ref{4.5:lem-curvature-remainder}) is defined
by
\begin{equation*}
F(U,A)=J(e^{i\eta A}U)-J(U)-D^{*}_{U}D_{U}A ,
\end{equation*}
where $J$ is the current in level-$n$ units.

Applying these two facts to \eqref{4.5:eq-increment-kernel-3} gives
\begin{align}
\Phi_1(\mathbf{U}_t)&=\operatorname{Ad}_{g(x)}\Bigl[\eta^{-2}
\bigl(\partial\bigl(e^{i\eta\mathcal{H}(B_t)}U_0\bigr)(p)-1\bigr)\Bigr],
\label{4.5:eq-increment-kernel-7}\\
\Phi_2(\mathbf{U}_t)&=R\bigl(g(b_-)\bigr)\Bigl[J(U_0)+D^{*}_{U_0}D_{U_0}\mathcal{H}(B_t)
+F\bigl(U_0,\mathcal{H}(B_t)\bigr)\Bigr](b).
\label{4.5:eq-increment-kernel-8}
\end{align}

\smallskip
\noindent\textit{Differentiation in $t$.}
Write $\dot{\mathcal{H}}:=\frac{d}{dt}\mathcal{H}(B_t)=\delta\mathcal{H}(B_t)[\dot B_t]$. Here
$\dot B_t$ is the differential of the logarithm at $V_tV_0^{-1}$ applied to the velocity
\begin{equation*}
\beta_t:=\tfrac1i\,\dot V_tV_t^{-1}.
\end{equation*}
On the polydisc \cite[(172)]{B-CMP102b} this differential has norm at most $2$.

Differentiating \eqref{4.5:eq-increment-kernel-7} in $t$, the derivative
$\frac{d}{dt}\Phi_1(\mathbf{U}_t)$ involves two things.
\begin{itemize}
\item It involves $\dot{\mathcal{H}}$ and $\nabla\dot{\mathcal{H}}$ at $p$; these are controlled
by the value and gradient bounds of the corollary on $\mathcal{H}$.
\item It involves $\dot g$, which contributes one more kernel term. This term consists of the
contour sums of $\eta\dot{\mathcal{H}}$ and their block mean, and is
$O\bigl(\sup_{\Delta(y_x)}|\dot{\mathcal{H}}|\bigr)$.
\end{itemize}

Differentiating \eqref{4.5:eq-increment-kernel-8} in $t$, the derivative
$\frac{d}{dt}\Phi_2(\mathbf{U}_t)$ involves three things.
\begin{itemize}
\item The term $D^{*}D\dot{\mathcal{H}}$, controlled by the bound on $D^{*}D\mathcal{H}$ in the
corollary on $\mathcal{H}$.
\item The term $\partial_AF\cdot(\dot{\mathcal{H}},\nabla\dot{\mathcal{H}})$, controlled by the
value and gradient bounds of that corollary together with
Lemma~\ref{4.5:lem-curvature-remainder}.
\item The commutator $[\delta\log g,\Phi_2]$ produced by differentiating $R(g(b_-))$.
\end{itemize}

For the second item, Lemma~\ref{4.5:lem-curvature-remainder} is applicable because
\begin{equation}\label{4.5:eq-increment-kernel-9}
N\bigl(\mathcal{H}(B_t)\bigr)\le 3dLC_Hc_1\sup|B_t|
\le 3dLC_Hc_1\cdot\tfrac12C_1\bar\varepsilon_1\le\mathfrak{r}_g\le 1 .
\end{equation}
The first inequality is the summation used for \eqref{4.5:eq-increment-kernel-5}: the value and
gradient bounds of the corollary on $\mathcal{H}$, summed over the sources with
\cite[Lemma~2.1]{B-CMP96} and the site multiplicity $3dL$. The second is
\eqref{4.5:eq-increment-kernel-2}. The third inequality is the fifth
entry of $\bar\varepsilon_1$. The last is $\mathfrak{r}_g\le1$.

\cite[Proposition~9]{B-CMP102b} contains a bound of the same kind, namely that the extended
configuration satisfies \cite[(19)--(21)]{B-CMP102b} with
$\varepsilon_2=B_5\varepsilon_1^{(102)}$. It is not used here.

For the third item, the second factor of the commutator is
$O(B_3\bar\varepsilon_1+C_H\cdot C_1\bar\varepsilon_1)=O(1)$. This is a constant depending on $d$
and $L$ only, not a power of $\lambda$.

\smallskip
\noindent\textit{The kernel.}
By item (a) of the geometric part of the proof, the site $y_x$ lies in the top layer, so
$\ell(y_x)=1$. The kernels of the corollary on $\mathcal{H}$ read at $y_x$ therefore carry no
power of $\ell(y_x)$.

We now collect the terms. Apply the decay bound of the corollary on $\mathcal{H}$ at the real
base $U_0$ to each $\dot{\mathcal{H}}$-term above, express $\dot{\mathcal{H}}$ through $\beta_t$
by the chain rule and the bound $2$ on the differential of the logarithm, and sum with
\cite[Lemma~2.1]{B-CMP96}. This gives, for $i=1,2$,
\begin{equation}\label{4.5:eq-increment-kernel-a}
\begin{gathered}
\frac{d}{dt}\Phi_i(\mathbf{U}_t)=\sum_{b'\in\mathfrak{B}_n}\ell_{b'}^{d}
\operatorname{tr}K_i(V_t;b')\,\beta_t(b'),\\
|K_i(V_t;b')|\le C_K\,\ell_{b'}^{-d}\,e^{-\frac14\delta_0 d(y_x,y')} .
\end{gathered}
\end{equation}
Here the bond $b'$ lies at the site $y'\in\Lambda_{q(b')}$ of the layer $q(b')$ that contains
it, $\ell_{b'}=\ell_{q(b')}$, the decay constant $\delta_0$ and the hierarchy distance
$d(y,y')$ are those fixed in Notation~\ref{4.5:not-geometry-conventions}, and the constant
$C_K=C_K(d,L)$ depends on $d$ and $L$ only.
\end{proof}

\begin{proof}[Uniform increment: recentring and two powers of the layer scale]
\label{4.5:prf-increment-recentring}
We show that the derivatives $d\Phi_i/dt$, $i=1,2$, represented by the kernel formula
\eqref{4.5:eq-increment-kernel-a} obtained in the preceding part of the proof
(\ref{4.5:prf-increment-kernel}), are bounded by $\tfrac12 C_{IV}\alpha_2\lambda^2$. The
mechanism is an exact infinitesimal recentring of $\hat A$ at the base point $x_0$ of the
block. Throughout, $\Phi_i$ is regarded as a function of the configuration through
$U_n\circ M^{\cdot}$.

\emph{The recentring field.} Let $Q$ be a cube about $x_0$ with
$\{z:|z-x_0|_\infty\le r/2\}\subset Q\subset\{z:|z-x_0|_\infty\le r\}$. Let $\chi$ be a
cut-off function with
$\chi=1$ on $\{z:|z-x_0|_\infty\le r/2\}$, $\chi=0$ off $Q$, and $|\nabla\chi|\le 4/r$, and put
\begin{equation}\label{4.5:eq-increment-recentring-1}
\nu(z):=-\chi(z)\sum_{\mu}\hat A_\mu(x_0)\,(z-x_0)_\mu .
\end{equation}
Thus $\nu$ is $\mathfrak{g}^{c}$-valued and linear in $z$ where $\chi=1$; $\nu(x_0)=0$, and
$\nu=0$ at the base point of the block $\Delta(y_x)$. For $s\in\mathbb{C}$ let $w_s:=e^{is\nu}$
be the corresponding site map, and let $\bar w_s:=w_s|_{\text{block sites}}$. The
tower of block averages is covariant:
\begin{equation}\label{4.5:eq-increment-recentring-2}
M^{\cdot}(\mathbf{U}_t^{w_s})=M^{\cdot}(\mathbf{U}_t)^{\bar w_s}
\end{equation}
by \cite{B-CMP98}; the similarity argument given there, together with
\cite[(56)--(57)]{B-CMP98}, uses only the invertibility of the site map, so it applies to the
complex map $w_s$. For the minimiser, \cite[(181)]{B-CMP102b} states that
$U_k(V^{v})=U_k(V)^{\tilde v}$, where $\tilde v$ is constant on blocks and equal to $v$ at the
block sites. We write $\tilde v_s$ for the block-constant map equal to $\bar w_s$ at the block
sites. This covariance is used here only through its derivative at $s=0$ along the
one-parameter family $w_s$, which passes through the identity map at $s=0$; we now justify
it without any radius condition on $\nu$.

\emph{Justification by the identity theorem.} Recall from the preceding part of the proof
that $V_t:=M^{\cdot}(\mathbf{U}_t)=e^{iB_t}V_0$ over the real base $V_0$. Put
\begin{equation*}
F_L(s):=U_n\bigl(M^{\cdot}(\mathbf{U}_t^{w_s})\bigr)=U_n\bigl((e^{iB_t}V_0)^{\bar w_s}\bigr),
\qquad
F_R(s):=U_n\bigl(M^{\cdot}\mathbf{U}_t\bigr)^{\tilde v_s}.
\end{equation*}
Choose $s_0>0$ so small that $(e^{iB_t}V_0)^{\bar w_s}$ stays in the polydisc
\cite[(172)]{B-CMP102b} over $V_0$ for $|s|<s_0$. For $|s|<s_0$, both $F_L$ and $F_R$ are
holomorphic functions of $s$, of $t$, and of the components of $\hat A$, $\hat A'$ in a basis
of $\mathfrak{g}$, on a connected product domain; this is \cite[Proposition~9]{B-CMP102b} at the
real base $V_0$ together with \cite[Proposition~6]{B-CMP98}.

Consider now real $(s,t)$ and $\mathfrak{g}$-valued $(\hat A,\hat A')$. Then every field and
every map is $G$-valued, since the base configuration $e^{i\eta\hat A_0}$ is. The data are
real and satisfy the curvature condition \cite[(7)]{B-CMP102b}. Indeed, put
$V':=(e^{iB_t}V_0)^{\bar w_s}V_0^{-1}$. It is $G$-valued, and it lies in the polydisc
\cite[(172)]{B-CMP102b} by the margin in \eqref{4.5:eq-uniform-increment-b}, so that
$|V'-1|<C_1\bar\varepsilon_1$. By the perturbation lemma for the real base tower, for
$G$-valued $V'$ with $|V'-1|<C_1\varepsilon_1$ the configuration $V'V_0$ satisfies
\cite[(7)]{B-CMP102b} with $(1+4C_1(1+O(\varepsilon_1)))\varepsilon_1$, the seams included.
Applied with $\varepsilon_1=\bar\varepsilon_1$, this gives condition \cite[(7)]{B-CMP102b}
with radius
\begin{equation*}
\bigl(1+4C_1(1+O(\bar\varepsilon_1))\bigr)\bar\varepsilon_1\le(1+5C_1)\bar\varepsilon_1 .
\end{equation*}
The list of bounds that fix the frozen datum radius $\bar\varepsilon_1$ contains
$\bar\varepsilon_1\le a_6$, where $a_6$ is an entry of that list with
$a_6\le a_1/(1+5C_1)$, and $a_1$ denotes the threshold of
\cite[Theorem~1]{B-CMP102b}. Hence
\begin{equation*}
(1+5C_1)\bar\varepsilon_1\le(1+5C_1)\,a_6\le a_1 .
\end{equation*}

Consequently both sides are the axial minimiser. For $G$-valued $V'$ the analytic extension
coincides with the minimal configuration \cite{B-CMP102b}, and the minimiser is unique by
\cite[Theorem~1]{B-CMP102b} and \cite{B-CMP99a}. The equality $F_L=F_R$ at these points is the
real form of the covariance \cite[(181)]{B-CMP102b} for minimisers; its two-line proof is
given with the literal frame, at \eqref{4.5:eq-literal-frame-b}. The real points form a
non-empty open subset of the maximal totally real subspace of the connected domain of
holomorphy. Hence $F_L\equiv F_R$ by the identity theorem. In particular
\begin{equation*}
\frac{d}{ds}\Big|_{s=0}F_L=\frac{d}{ds}\Big|_{s=0}F_R
\quad\text{at the given complex } (t,\hat A,\hat A').
\end{equation*}

Equivalently: for $|s|$ small, $w_s=e^{is\nu}$ lies within any prescribed distance of the
identity. It therefore lies within the hypotheses of the remark on frame maps near the
identity, which concerns frame maps of the form $v=v'v_G$ with $v'$ in a prescribed
neighbourhood of $1$, taken at $v_G=1$. No smallness of $\nu$ is used, only the
derivative at $s=0$; the bound $|\nu|\le dM\alpha_2$ obtained below need not be small.

\emph{Exact invariance of $\Phi_1,\Phi_2$.} Since $\tilde v_s$ is constant on blocks and equal
to $e^{is\nu(x_0)}=1$ on the block of $x$, we have $\tilde v_s=1$ on $\Delta(y_x)$. The
plaquette $p$ is based at $x$, and $b_-=x\in\Delta(y_x)$. Therefore, for $|s|<s_0$,
\begin{align*}
\Phi_1(\mathbf{U}_t^{w_s})&=\operatorname{Ad}_{\tilde v_s(x)}\Phi_1(\mathbf{U}_t)
=\Phi_1(\mathbf{U}_t),\\
\Phi_2(\mathbf{U}_t^{w_s})&=R(\tilde v_s(b_-))\,\Phi_2(\mathbf{U}_t)=\Phi_2(\mathbf{U}_t),
\end{align*}
exactly, where the first equality in each line is $F_L\equiv F_R$ followed by the
transformation law of $\Phi_i$ under the block-constant gauge map $\tilde v_s$ (by
$\operatorname{Ad}_{\tilde v_s(x)}$ for $\Phi_1$, by $R(\tilde v_s(b_-))$ for $\Phi_2$).

\emph{Replacement of $\hat A$ by the recentred field.} The derivative at $s=0$ of
$\mathbf{U}_t^{w_s}$ is the velocity field $-D_{\mathbf{U}_t}\nu$, where
\begin{equation*}
\eta\,D_{\mathbf{U}_t}\nu(b)=\operatorname{Ad}_{\mathbf{U}_t(b)}\nu(b_+)-\nu(b_-).
\end{equation*}
The velocity of $\mathbf{U}_t$ in $t$ is $\hat A$. Both velocities enter $d\Phi_i$ through the
same linear map, namely the derivative of $\Phi_i\circ U_n\circ M^{\cdot}$ at $\mathbf{U}_t$.
By the exact invariance, this map annihilates $-D_{\mathbf{U}_t}\nu$. Adding
$D_{\mathbf{U}_t}\nu$ to $\hat A$ therefore does not change $d\Phi_i/dt$. Put
\begin{equation}\label{4.5:eq-increment-recentring-3}
\tilde A_t:=\hat A+D_{\mathbf{U}_t}\nu .
\end{equation}
Hence in \eqref{4.5:eq-increment-kernel-a} the velocity $\beta_t$ may be replaced by
$DM^{\cdot}(\mathbf{U}_t)[\tilde A_t]$ without changing the sum.
Moreover,
\begin{equation}\label{4.5:eq-increment-recentring-4}
\bigl|DM^{\cdot}(\mathbf{U}_t)[\tilde A](b')\bigr|
\le 2c_Q\,\ell_{b'}\sup_{B(b'_-)\cup B(b'_+)}|\tilde A| .
\end{equation}
This follows from \cite[Propositions~6 and~7]{B-CMP98} by the Cauchy estimate at the point
$t\hat A+\hat A'$, which lies inside half the radius by the third ceiling in
\eqref{4.5:eq-uniform-increment-b}. The restriction of the supremum to
$B(b'_-)\cup B(b'_+)$ is the locality of the block averages \cite{B-CMP98}: the $k$-fold
average over a block bond $c$ of the region on which the $k$-fold averages are taken depends
only on the bond variables in the two blocks at the end-points of $c$.

\emph{Sizes on the inner cube $\{\chi=1\}$.} There
\begin{equation*}
D_{\mathbf{U}_t}\nu(b_\mu)=-\hat A_\mu(x_0)
+\eta^{-1}\bigl[\operatorname{Ad}_{\mathbf{U}_t(b)}-1\bigr]\nu(b_+),
\qquad
\mathbf{U}_t(b)=e^{it\eta\hat A}e^{i\eta\hat A'}e^{i\eta\hat A_0},
\end{equation*}
and therefore
\begin{align*}
|\tilde A_t(z)|
&\le|\hat A(z)-\hat A(x_0)|+2\bigl(|\hat A|+|\hat A'|+|\hat A_0|\bigr)|\nu(z)|\\
&\le\alpha_2\lambda^2\bigl(|z-x_0|_1+1\bigr)
+2\bigl(\alpha_2+\tfrac12\alpha_1+c_{\sharp}\bigr)\lambda\cdot d\alpha_2\lambda|z-x_0|_1\\
&\le 2d\alpha_2\lambda^2\bigl(|z-x_0|_1+1\bigr).
\end{align*}
Here $\hat A(z)-\hat A(x_0)$ telescopes bond by bond, through $\nabla_{\mathbf{U}}\hat A$ and
conjugations of size $2(|\hat A|+|\hat A'|+|\hat A_0|)\eta$ per bond. No long non-unitary
transport occurs, and no factor $M\alpha_1$ arises. The radius $\alpha_1$ enters only through
the coefficient $2(\alpha_2+\tfrac12\alpha_1+c_{\sharp})$, which multiplies $\alpha_2\lambda^2$.

\emph{Sizes on the collar $Q\setminus\{\chi=1\}$ and outside $Q$.} On the collar,
$|\nu|\le d\alpha_2\lambda\cdot r=dM\alpha_2$. Using $|\nabla\chi|\le 4/r$,
\begin{align*}
|D_{\mathbf{U}_t}(\chi\nu)|
&\le|D\nu|+|\nabla\chi|\,|\nu|+2\bigl(|\hat A|+|\hat A'|+|\hat A_0|\bigr)|\nu|\\
&\le\alpha_2\lambda+4d\alpha_2\lambda
+2\bigl(\alpha_2+\tfrac12\alpha_1+c_{\sharp}\bigr)\lambda\cdot dM\alpha_2\\
&\le C(d)(1+M)\alpha_2\lambda .
\end{align*}
Off $Q$, $\nu=0$ and $\tilde A=\hat A$.

\emph{The sums in \eqref{4.5:eq-increment-kernel-a}.} Block diameters are at most $d$. Hence a
source $b'$ at $y'$ with $|y'-x_0|\le r/2-d$ sees only the inner cube $\{\chi=1\}$ in
\eqref{4.5:eq-increment-recentring-4}. The local units satisfy $\ell_{b'}\le1$, and this also
gives $d(y_x,y')\ge|y'-x_0|-c$, for a geometric constant $c$ of the tower, which is the
constant $c$ in the estimates below. Combining the kernel bound
$|K_i(b')|\le C_K\ell_{b'}^{-d}e^{-\frac14\delta_0 d(y_x,y')}$ with
\eqref{4.5:eq-increment-recentring-3} and \eqref{4.5:eq-increment-recentring-4} gives the
following two estimates.

The near part is at most
\begin{equation*}
\sum_{y'}C_Ke^{-\frac14\delta_0 d(y_x,y')}\cdot 2c_Q\cdot
2d\alpha_2\lambda^2\bigl(d(y_x,y')+c+1\bigr)\le C'\alpha_2\lambda^2,
\end{equation*}
because, by \cite[Lemma~2.1]{B-CMP96},
\begin{equation*}
\sum_{y'}\bigl(d(y_x,y')+c\bigr)e^{-\frac14\delta_0 d(y_x,y')}\le C .
\end{equation*}

For the far part, $|y'-x_0|\ge r/2-d$ implies $d(y_x,y')\ge M/(2\lambda)-c$. We split
\begin{equation*}
e^{-\frac14\delta_0 d(y_x,y')}=e^{-\frac18\delta_0 d(y_x,y')}\,e^{-\frac18\delta_0 d(y_x,y')}
\end{equation*}
and use the collar bound, which also dominates $|\hat A|$ off $Q$. The far part is then at most
\begin{align*}
&2c_QC_KC(d)(1+M)\alpha_2\lambda\,e^{-\frac18\delta_0(M/(2\lambda)-c)}
\sum_{y'}e^{-\frac18\delta_0 d(y_x,y')}\\
&\qquad\le C''(1+M)\alpha_2\lambda\,e^{-\delta_0M/(16\lambda)}
\le C'''\alpha_2\lambda^2 .
\end{align*}
The last step is as follows. Let $0<\lambda\le1$ and $a:=\delta_0M/16\ge2$. Then
$e^{-a/(2\lambda)}\le e^{-1/\lambda}$ because $a/2\ge1$; $e^{-a/(2\lambda)}\le e^{-a/2}$
because $\lambda\le1$; and $e^{-1/\lambda}\le\lambda$ because $e^{1/\lambda}\ge1/\lambda$.
Together,
\begin{equation*}
e^{-a/\lambda}\le e^{-a/(2\lambda)}e^{-1/\lambda}\le e^{-a/2}\lambda .
\end{equation*}
Furthermore, since $\sup_{\sigma>0}\sigma e^{-\delta_0\sigma/32}=32/(e\delta_0)$,
\begin{equation*}
(1+M)e^{-\delta_0M/32}\le 1+\frac{32}{e\delta_0},
\end{equation*}
which is a $(d,L)$-number. The condition $a\ge2$ is the floor $\delta_0M\ge32$ in
\eqref{4.5:eq-uniform-increment-b}; the second bound holds for every $M\ge0$.

\emph{Conclusion.} Adding the near and far parts, $|d\Phi_i/dt|\le(C'+C''')\alpha_2\lambda^2$.
Hence
\begin{equation}\label{4.5:eq-increment-recentring-5}
\Bigl|\frac{d\Phi_i}{dt}\Bigr|\le\tfrac12C_{IV}\alpha_2\lambda^2,
\qquad
C_{IV}:=\max\bigl\{2(C'+C'''),\,4C_{14}\bigr\}.
\end{equation}
The constant $C_{IV}$ is independent of $M$, of $n$, of $j$ and of the region $X$.
\end{proof}

\begin{proof}[Proof of Lemma~\ref{4.5:lem-uniform-increment}, conclusion]
\label{4.5:prf-increment-conclusion}
Fix $n\in\{1,\dots,j\}$ and $x\in X^{\sim 2}$, and work in level-$n$ units, with $\xi=L^{-j}$ and
$\lambda=L^{n-j}\le 1$. Integrating the bound \eqref{4.5:eq-increment-kernel-a} over $t\in[0,1]$ gives,
for $i=1,2$,
\begin{equation}\label{4.5:eq-increment-conclusion-1}
\bigl|\Phi_i(e^{i\xi A}\mathbf{U})-\Phi_i(\mathbf{U})\bigr|\le \tfrac12 C_{IV}\alpha_2\lambda^2 .
\end{equation}
This is the first assertion of \eqref{4.5:eq-uniform-increment-a}, and it holds for every $n\le j$ and
every $x\in X^{\sim 2}$. The ceiling $\alpha_2\le\alpha_0/C_{IV}$ of the lemma bounds the right-hand side
of \eqref{4.5:eq-increment-conclusion-1} by $\tfrac12\alpha_0\lambda^2$.

The hypothesis of the lemma gives $|\Phi_i(\mathbf{U})|<\tfrac12\alpha_0\lambda^2$ at the same $n$ and $x$.
By the triangle inequality and \eqref{4.5:eq-increment-conclusion-1},
\begin{equation*}
\bigl|\Phi_i(e^{i\xi A}\mathbf{U})\bigr|
\le \bigl|\Phi_i(\mathbf{U})\bigr|+\bigl|\Phi_i(e^{i\xi A}\mathbf{U})-\Phi_i(\mathbf{U})\bigr|
<\tfrac12\alpha_0\lambda^2+\tfrac12\alpha_0\lambda^2=\alpha_0\lambda^2 ,
\qquad i=1,2 .
\end{equation*}
This is \cite[(1.16)]{B-CMP109} at the constant $\alpha_0$, for every $n\le j$ and every $x\in X^{\sim 2}$.
Both members are obtained, the first member in the uniform form recorded in
\eqref{4.5:eq-uniform-increment-a}. Hence the half of condition~(iv) of \cite{B-CMP109} that concerns the
configuration $e^{i\xi A}\mathbf{U}$ holds with $\alpha_0$. It holds for the tower $\mathbf{\Omega}$ of the
hypothesis, that is, for the same admissible tower of $X$ on whose sub-towers
$\{\Omega_0,\dots,\Omega_n\}$ the hypothesis asks condition~(iv), whether or not $\mathbf{\Omega}$ is the
maximal tower of $X$.

Since $e^{i\xi A}\mathbf{U}=e^{i\xi A''}U$ has the same $G$-valued factor $U$ as $\mathbf{U}$, the half of
condition~(iv) that concerns $U$, and condition~(i), are the ones assumed and are unchanged.

It remains to record the dependence of the constants. Every constant used above depends on $(d,L)$ only,
and none of $k$, $n$, $j$ or $j-n$ enters any of them. The quantities $M$, $R_1M_1$ and $\delta_0M$ enter
the argument only through the following conditions:
\begin{itemize}
\item the floors $\delta_0M\ge 32$ and $2R_1M_1\le M$ (or, in place of the second floor, the collar
hypothesis on the tower);
\item the ceiling $c_{12}LMB\cdot\tfrac12\alpha_0\le c_{\sharp}$;
\item the tail bound displayed below, whose right-hand side does not involve $M$.
\end{itemize}
The tail bound reads
\begin{equation}\label{4.5:eq-increment-conclusion-2}
(1+M)\,e^{-\delta_0M/32}\le 1+\sup_{y\ge 0} y\,e^{-\delta_0 y/32}=1+\frac{32}{e\,\delta_0}.
\end{equation}
It holds because $e^{-\delta_0M/32}\le 1$ and because the function $y\mapsto y\,e^{-\delta_0y/32}$ attains
its maximum on $[0,\infty)$ at $y=32/\delta_0$. Consequently every constant of the argument is uniform in
$M$ over the half-line on which these floors hold.

No inequality between $\alpha_0$ and $\alpha_1$ is used. The amplitude $\alpha_2$ is constrained only by
ceilings imposed after the $(d,L)$-constants have been fixed. This agrees with \cite{B-CMP109}, where
$\alpha_2$ is a constant depending on $\alpha_0$ and on absolute constants.
\end{proof}

\begin{proposition}[The glued window background lies in the orbit class]
\label{4.5:prop-glued-window}
Fix the level $k$, and let $\beta$ be the schedule constant of the layered analyticity spaces
(Definition~\ref{4.4:def-post-step-space}). Let $a_0$, $a_1$, $a_{\gamma}$, $c_9$ be the constants of
Theorem~\ref{4.5:thm-chart-existence}, let $\mathfrak{S}$ and $\mathfrak{S}^{\mathrm{orb}}$ be the
class of admissible background pairs and its orbit extension (Definition~\ref{4.5:def-admissible-pairs},
conditions \eqref{4.5:eq-admissible-pairs-a}), with $\mathfrak{K}=4$, and let $N_U$ denote the norm of
condition (S2) at level $k$. We use the following window and assume the following theorem.
\begin{enumerate}
\item[(W)] \emph{The window.} Let $\square$ be a cover cube, that is, a union of $2^{d}$ cubes of the
partition $\pi_k$ of \cite{B-CMP109}, and let $\square_0:=\tilde\square^{5}$ be its enlargement in
\cite{B-CMP109}. Its side is $12M$ in units $\ell_k$, that is, $12M/L$ in units $\ell_{k+1}$; here
$\ell_j$ is the unit of length at level $j$, and $\ell_k<\ell_{k+1}$. The window is of the type in
which the small-field region $\Omega_{k+1}$ at level $k+1$ is the whole lattice $T_{\eta}$; thus
$\square_0\subset\Omega_{k+1}$, and $b:=k+1$. Let $\{\square_n\}_{n=0}^{k+1}$ be the tower of
\cite[(2.13)]{B-CMP119} built on $\square_0$, read as in $(\rho2)$ of
the corresponding display; thus it is admissible in the sense of \cite[(1)]{B-CMP102b}
with block $\mathfrak{M}$ and admissibility radius $LR'_{\star}>R_1$, and it is the tower of
condition (iv) of the space $\mathfrak{U}$ below. Its
layers are $\Lambda_q:=\square_q\setminus\square_{q+1}$, $q=0,\dots,b-1$; its top region is
$\square_{b}$; $P:=\Lambda_0$; and
$\mathfrak{B}(\square_0):=\bigcup_{n\le k+1}\Lambda_n(\square_0)$ is its block-bond set, in the sense
of \cite[(1.12)]{B-CMP99a}. The sites of $\square_0$ include the end-points of its bonds. For a
configuration $\mathbb{U}$ on the bonds of
$\square_0$, the axial representative is
$S[\mathbb{U}]:=U_{k+1}(\square_0,M^{\cdot}(\mathbb{U}))$, built on $\mathfrak{B}(\square_0)$ from the
tower $M^{\cdot}(\mathbb{U})$ of block averages of \cite[(1.15)]{B-CMP109}.
\item[(R)] \emph{The rim regularity theorem.} Let $L_{\mathrm{rim}}$ and $M_{1,\mathrm{rim}}$ denote
the lower bounds on $L$ and on $M_1$ in the hypotheses of the rim regularity theorem. Let $L$ be odd
with $L\ge L_{\mathrm{rim}}$, let $M_1\ge M_{1,\mathrm{rim}}$, and let $M>L^{2}M_1$. There are
$c_{\mathrm{rim}}=c_{\mathrm{rim}}(M)>0$ and $K_{\mathrm{rim}}=K_{\mathrm{rim}}(M)<\infty$, depending on
$(d,G,L,M_1,M)$ only, with
\begin{equation*}
K_{\mathrm{rim}}=10^{3}d\,\mathsf{M}\,(K_{\mathbf{1}}+1),\qquad \mathsf{M}=M/L,
\end{equation*}
where $K_{\mathbf{1}}<\infty$ is a constant that does not depend on $M$, fixed in the proof of the rim
regularity theorem, such that the following holds for every window
$(\square_0,b)$ as in (W), every real $\ell_*\ge\ell_b$, every $a\le c_{\mathrm{rim}}$ and every
configuration $\mathbb{U}=e^{i\eta A'}U$ on the bonds of $\square_0$ satisfying:
\begin{enumerate}
\item[(h1)] $U$ is $G$-valued, and there is a $G$-valued site map $u_0$ on $\square_0$ with
$U^{u_0}=e^{i\eta A_0}$ and
$\ell_*|A_0|\le a$, $\ell_*^{2}|\nabla^{\eta}A_0|\le a$;
\item[(h2)] $A'$ is $\mathfrak{g}^{c}$-valued, with
$\ell_*|A'|\le a$, $\ell_*^{2}|\nabla^{\eta}_{U}A'|\le a$.
\end{enumerate}
Then:
\begin{enumerate}
\item[(a)] $S=S[\mathbb{U}]$ is defined and holomorphic in $A'$;
\item[(b)] there are a site map $h$ on $\square_0$ and a bond function $\Xi$ on $\square_0$ with
$S^{h}=e^{i\eta\Xi}\mathbb{U}$, such that $h=1$ at every site of $P$ and $\Xi=0$ on every bond not
in $\square_1$;
\item[(c)] the bounds of (e) hold on all bonds of $\square_0$: on every layer $\Lambda_q$,
$q=0,\dots,b-1$, on the interfaces between layers, on $P$ and on the top region alike;
\item[(d)] $\|\operatorname{Ad}_h^{\pm1}\|\le 2$; and if $\mathbb{U}$ is $G$-valued, then $h$ is
$G$-valued and $\Xi$ is $\mathfrak{g}$-valued;
\item[(e)] $S^{h}=e^{i\eta A''}U$, with $A''=A'$ on the bonds not in $\square_1$, and
\begin{equation*}
\ell_*|A''|,\ \ell_*^{2}|\nabla^{\eta}_{U}A''|\le K_{\mathrm{rim}}\,a,
\qquad |S^{h}(\partial p)-1|\le K_{\mathrm{rim}}\,a\,\eta^{2}\ell_*^{-2}
\end{equation*}
for the plaquettes $p$ of $\square_0$.
\end{enumerate}
\end{enumerate}
Assume $L\ge L_{\mathrm{rim}}$, $M_1\ge M_{1,\mathrm{rim}}$, $M>L^{2}M_1$, and $L^{2}\ge 12/c_{12}$, where
$c_{12}$ is the cube-size constant of the reading $(\rho1)$ of the corresponding display.
Assume moreover that the region $P$ of (W) is $LM_1$ bonds of the finest lattice thick (the thickness
statement for the tower of (W)).
Let $\mathfrak{U}$ be the space $U'^{c}_{k+1}(\square_0,(1+2\beta)\alpha_0,(1+2\beta)\alpha_1,\alpha_0)$
of \cite{B-CMP109}: it is defined by the conditions (i)--(iv) of \cite{B-CMP109}, except that the
configurations $\mathbf{U}$, $\mathbf{J}$ are defined, and satisfy (i)--(iii), on the whole lattice
$T_{\eta}$. Let $(\mathbf{U},\mathbf{J})\in\mathfrak{U}$, with $\mathbf{U}=e^{i\eta A'}U$, $U$
$G$-valued. Define the glued background $\mathbf{U}^{\mathrm{gl}}$ by
$\mathbf{U}^{\mathrm{gl}}:=S=S[\mathbf{U}]$ on the bonds of $\square_0$ and
$\mathbf{U}^{\mathrm{gl}}:=\mathbf{U}$ on all other bonds, and put
\begin{equation*}
\alpha_0''':=(1+2\beta)\alpha_0,\qquad \alpha_1''':=(1+2\beta)\alpha_1,\qquad
\mathbf{a}:=\max\{c_{12}LMB^{109}\,\alpha_0''',\ \alpha_1'''\},
\end{equation*}
with $B^{109}$ the constant of \cite[(1.12)]{B-CMP109} (Definition~\ref{4.4:def-local-data-classes}).
Assume the ceiling
\begin{equation}\label{4.5:eq-glued-window-a}
\begin{gathered}
\mathbf{a}\le c_{\mathrm{rim}}(M),\qquad K_{\mathrm{rim}}(M)\,\mathbf{a}\le a_1,\\
c_{12}LMB^{109}\,\alpha_0'''\le a_0,\qquad 2\alpha_0\le a_{\gamma}.
\end{gathered}
\end{equation}
Then $\mathbf{U}$, restricted to the bonds of $\square_0$, satisfies the hypotheses (h1) and (h2) of
(R) with $\ell_*=\ell_{k+1}$ and $a=\mathbf{a}$. Let $h$ and $A''$ be the maps which (R) then provides
for $\mathbb{U}=\mathbf{U}$ on the bonds of $\square_0$, and extend $h$ by $1$ and $A''$ by $A'$
outside $\square_0$. Moreover:
\begin{enumerate}
\item[(i)] one has
\begin{equation*}
((\mathbf{U}^{\mathrm{gl}})^{h},\operatorname{Ad}_h\mathbf{J})
=(e^{i\eta A''}U,\operatorname{Ad}_h\mathbf{J})\in\mathfrak{S},
\end{equation*}
with the trivial sequence (Definition~\ref{4.5:def-admissible-pairs}), at level $k$, and with
$U_r:=U$;
\item[(ii)] $(\mathbf{U}^{\mathrm{gl}},\mathbf{J})\in\mathfrak{S}^{\mathrm{orb}}$; data $B'$ with
$|B'|<c_9/8$ lie within the orbit radius, since $c_9/8<c_9/\mathfrak{K}$;
\item[(iii)] the map $\mathbf{U}\mapsto\mathbf{U}^{\mathrm{gl}}$ is holomorphic and single-valued; it
has no seam;
\item[(iv)] on all bonds of $\square_0$, namely on the graded collar $\square_0\setminus\square_{k+1}$
(every layer $\Lambda_q$ of the tower, the interfaces, and the region $P$, on which $h=1$) and on the
top region alike,
\begin{equation}\label{4.5:eq-glued-window-1}
S^{h}=e^{i\eta A''}U,\qquad N_U(A'')\le K_{\mathrm{rim}}(M)\,\mathbf{a}<a_1,
\end{equation}
with one $G$-valued factor $U$, gauged on $\square_0$ by $u_0$, and one twist $A''$; across
$\partial\square_1$ one has $A''=A'$, and across $\partial\square_0$ one has $S=\mathbf{U}$.
\end{enumerate}
\end{proposition}

\begin{proof}
We apply the rim regularity theorem (R) to the window of (W), for which
$\square_0\subset\Omega_{k+1}=T_{\eta}$ and $b=k+1$, with $\ell_*:=\ell_b=\ell_{k+1}$,
$\mathbb{U}:=\mathbf{U}$ on the bonds of $\square_0$, and $a:=\mathbf{a}$. We first check its
hypotheses.

\emph{Hypothesis (h1).} The cube $\square_0$ has side $12M/L$ in units $\ell_{k+1}$. Since
$L^{2}\ge 12/c_{12}$, we have $12/L\le c_{12}L$, hence $12M/L\le c_{12}LM$. Therefore $\square_0$ is a
single cube of the class of cubes on which condition \cite[(1.12)]{B-CMP109}, read as in $(\rho1)$ of
the corresponding display, holds for the level-$(k+1)$ space $\mathfrak{U}$. Let $u_0$ be
the $G$-valued gauge which this condition provides on $\square_0$; then $U^{u_0}=e^{i\eta A_0}$ with
\begin{equation*}
\ell_b|A_0|,\ \ell_b^{2}|\nabla^{\eta}A_0|<(1+2\beta)c_{12}LMB^{109}\,\alpha_0
=c_{12}LMB^{109}\,\alpha_0'''\le\mathbf{a}.
\end{equation*}
The factor $U$ of $\mathbf{U}$ is $G$-valued. Thus (h1) holds with $\ell_*=\ell_b$ and $a=\mathbf{a}$.

\emph{Hypothesis (h2).} Condition \cite[(1.13)]{B-CMP109} at level $k+1$ gives
\begin{equation*}
\ell_b|A'|,\ \ell_b^{2}|\nabla^{\eta}_{U}A'|<\alpha_1'''\le\mathbf{a},
\end{equation*}
and $A'$ is $\mathfrak{g}^{c}$-valued. Thus (h2) holds. Finally $a=\mathbf{a}\le c_{\mathrm{rim}}$ by
the first inequality of \eqref{4.5:eq-glued-window-a}. Hence the conclusions (a)--(e) of (R) hold.

\emph{Gluing.} By (b) and (e), $S^{h}=e^{i\eta\Xi}\mathbf{U}=e^{i\eta A''}U$ on the bonds of
$\square_0$, with $h=1$ at every site of $P$, $\Xi=0$ and $A''=A'$ on every bond not in $\square_1$.
We extend $h$ by $1$ and $\Xi$ by $0$ outside $\square_0$, and $A''$ by $A'$. The region $P$ is the
lowest layer $\Lambda_0$ of the tower. For this layer $B^{0}(\Lambda_0)=\Lambda_0$ in the notation of
\cite{B-CMP99a}, and the constraint of \cite[(3)]{B-CMP102b} reads $\bar U^{j}=V$ on $\Lambda_j$; at
$j=0$ the constraint is the field itself. Hence $S=\mathbf{U}$ on every bond not in $\square_1$. By
the assumption on its thickness, $P$ is $LM_1$ bonds of the finest lattice thick. Therefore no
plaquette and no difference quotient
involves both a bond of $\square_1$ and a bond outside $\square_0$, and a difference quotient which
reads one bond of $P$ and one bond outside $\square_0$ reads $A'$ on both of them. Since $h=1$ at every
site of $P$, and outside $\square_0$, the gauge-transformed glued configuration is
$(\mathbf{U}^{\mathrm{gl}})^{h}=e^{i\eta A''}U$ on every bond of $T_{\eta}$: on the bonds of
$\square_0$ by (e), and on the other bonds because there $h=1$,
$\mathbf{U}^{\mathrm{gl}}=\mathbf{U}=e^{i\eta A'}U$ and $A''=A'$.

\emph{Condition (S1).} We take $U_r:=U$. The $G$-valued factor is not changed by the gluing, so the
conditions \cite[(1.11)--(1.12)]{B-CMP109} hold on cubes which straddle $\partial\square_0$ or
$\partial\square_1$ exactly as on any other cube, with amplitude $c_{12}LMB^{109}\,\alpha_0'''$, which
is $\le a_0$ by the third inequality of \eqref{4.5:eq-glued-window-a}. The whole displacement of $S$
from $\mathbf{U}$, real part included, is carried by the twist $A''$. This is admissible, because
condition (S2) is a ball of absolute radius $a_1$ in the twist, as in \cite{B-CMP99b}.

\emph{Condition (S2).} A bound at level $k+1$ is a fortiori a bound at level $k$, since
$\ell_k<\ell_{k+1}$. On the bonds of $\square_0$, clause (e) with $a=\mathbf{a}$ gives the bound
$K_{\mathrm{rim}}\mathbf{a}$ for $A''$; on the other bonds $A''=A'$ and \cite[(1.13)]{B-CMP109},
valid on the whole of $T_{\eta}$, gives the bound
$\alpha_1'''$; by the gluing step no difference quotient mixes the two regions except through $A'$.
Hence
\begin{equation*}
N_U(A'')\le\max\{K_{\mathrm{rim}}\mathbf{a},\ \alpha_1'''\}<a_1 .
\end{equation*}

\emph{Condition (S3).} In \cite{B-CMP109} the third parameter of the space is $\gamma_0=\alpha_0$:
the bound for $\mathbf{J}$ carries a different constant from the bounds for $\mathbf{U}$. By (d),
$\|\operatorname{Ad}_h\|\le 2$, so
\begin{equation*}
|\operatorname{Ad}_h\mathbf{J}|\le 2|\mathbf{J}|<2\alpha_0\le a_{\gamma},
\end{equation*}
the last inequality being the fourth of \eqref{4.5:eq-glued-window-a}. This proves (i), with the
trivial sequence at level $k$.

\emph{Orbit.} By (d), $\|\operatorname{Ad}_h^{\pm1}\|\le 2\le 4=\mathfrak{K}$. The pair
$(\mathbf{U}^{\mathrm{gl}},\mathbf{J})$ is obtained from the pair of (i), which lies in
$\mathfrak{S}$, by the gauge transformation $h^{-1}$, whose adjoint action and its inverse have norm
at most $2$. Hence $(\mathbf{U}^{\mathrm{gl}},\mathbf{J})\in\mathfrak{S}^{\mathrm{orb}}$ by
Definition~\ref{4.5:def-admissible-pairs}. A datum $B'$ with $|B'|<c_9/8$ satisfies
$|B'|<c_9/4=c_9/\mathfrak{K}$. This proves (ii).

\emph{Holomorphy and single-valuedness.} On the bonds of $\square_0$, $\mathbf{U}^{\mathrm{gl}}=S$,
which is holomorphic in $A'$ by (a); on the other bonds $\mathbf{U}^{\mathrm{gl}}=\mathbf{U}$. The map
is given by this direct formula, hence is single-valued, and it has no seam, because $S=\mathbf{U}$ on
every bond not in $\square_1$. Single-valuedness of the chart on $\mathfrak{S}^{\mathrm{orb}}$ is
supplied by the gauge covariance of the chart on the orbit class (see~\ref{4.5:prf-chart-covariance}).
Membership in
$\mathfrak{S}^{\mathrm{orb}}$ is needed only pointwise in $\mathbf{U}$; holomorphy of
$\mathbf{U}\mapsto h(\mathbf{U})$ is neither needed nor claimed. The analyticity of the chart at
$\mathbf{U}^{\mathrm{gl}}$ comes from its direct formulas at the point $S$, which is holomorphic by
(a); the gauge map $h$ enters only the bounds. This proves (iii).

\emph{The collar representation.} By (e) and (c), on all bonds of $\square_0$ (every layer
$\Lambda_q$, the interfaces, $P$ and the top region alike), $S^{h}=e^{i\eta A''}U$ with the one
$G$-valued factor $U$, gauged on $\square_0$ by $u_0$ as in (h1), and the one twist $A''$, and by
condition (S2) above $N_U(A'')\le K_{\mathrm{rim}}(M)\mathbf{a}<a_1$. Across $\partial\square_1$ one has
$A''=A'$ by (e), and across $\partial\square_0$ one has $S=\mathbf{U}$ by the gluing step. This proves
(iv).

\emph{Reality.} If $\mathbf{U}$ is $G$-valued, clause (d) gives, moreover, that $h$ is $G$-valued and
$\Xi$ is $\mathfrak{g}$-valued.
\end{proof}

Each of the four inequalities of \eqref{4.5:eq-glued-window-a} is an upper bound on $\alpha_0$ or on
$\alpha_1$ once $L$, $M_1$ and $M$ are fixed; none of them is a cap in the sense of
Definition~\ref{2.3:def-cap}. Since $\mathbf{a}\ge\alpha_1'''$, the second inequality contains
$K_{\mathrm{rim}}(M)(1+2\beta)\alpha_1\le a_1$, which, by the value of $K_{\mathrm{rim}}$ in (R), reads
\begin{equation*}
10^{3}d\,(K_{\mathbf{1}}+1)(1+2\beta)\cdot\frac{M}{L}\cdot\alpha_1\le a_1 .
\end{equation*}
This is a smallness condition on $M\alpha_1$, of the kind made in \cite[p.~275]{B-CMP109}, where
$\alpha_0$, $\alpha_1$ are chosen so that $M\alpha_0$ and $M\alpha_1$ are still small.

In the frames of $\mathbf{J}$ and $B'$ the gauge map $h$ is rough: only
$\|\operatorname{Ad}_h^{\pm1}\|\le 2$ is known. This is harmless, because the chart reads $J$ only
through pointwise commutators, pairings and kernels summed against $J$, never through $D^{*}_{U}J$ or
a H\"older norm of $J$ (clause (K7) of Theorem~\ref{4.5:thm-chart-existence}), and it reads $B'$
only through $\sup|B'|$.

\begin{remark}[Ba{\l}aban's gauge frame for the near term
{\cite[(3.25)--(3.32), pp.~275--277]{B-CMP109}}]\label{4.5:rem-near-term-frame}
We use the notation of \cite{B-CMP109}:
\begin{itemize}
\item $\eta$ is the lattice spacing of the lattice $T_\eta$ on which the configurations $\mathbf{U}$ and the currents $\mathbf{J}$ are defined;
\item $\square_0$ is the window attached to a cube $\square$ of the cover, with the enlarged cubes $\tilde\square$, $\tilde\square^{3}$, $\tilde\square^{4}$ of \cite{B-CMP109}, and $y_0$ is the centre of $\square$;
\item $Q$ is the block-averaging operation, $M^{\cdot}(\mathbf{U})$ is the tower of block averages $\bar{\mathbf{U}}^{p}$ of $\mathbf{U}$, and $U_j(\square_0,\cdot)$, $U_{k+1}(\square_0,\cdot)$ are the configurations of the indicated levels that their arguments determine on $\square_0$;
\item $T^{(k+1)}$ is the lattice of level $k+1$, $\Lambda_n$ are the bond layers of the hierarchy, and $\square_j^{(j)}$ is the level-$j$ cube of the tower of the window $\square_0$;
\item $\mathbf{H}_k(B')$ and $\mathbf{H}_{k+1}(\square_0,\cdot)$ are Ba{\l}aban's background functions, $B'$ being the argument at which \cite{B-CMP109} evaluates $\mathbf{H}_k$ in the function \eqref{4.5:eq-near-term-frame-1} below; $\mathbf{E}^{(j)}(X,\cdot)$ is the effective-action term attached to a localisation domain $X\in\mathbf{D}_j$, and $R(u)$ denotes the action of a gauge transformation $u$ on the current $\mathbf{J}$ and on the further variable $B$ on which, together with $\mathbf{J}$, the function $\mathbf{H}_k(B')$ depends;
\item $|\cdot|$, $\nabla^{\eta}$ and $\|\cdot\|_{1,\beta}$ are the sup norm, the lattice gradient at spacing $\eta$ and the H\"older-type norm of \cite{B-CMP109}; $\alpha_0,\alpha_1,\alpha_2,\beta_0,B_3,\varepsilon_1$ are the parameters and constants of \cite{B-CMP109}; $M$ is Ba{\l}aban's parameter.
The letter $\beta$ occurs in two roles, as in \cite{B-CMP109}: in the radii $(1+2\beta)\alpha_0$, $(1+2\beta)\alpha_1$ of the space $\mathfrak{U}$ below it is the fixed schedule constant of Definition~\ref{4.4:def-post-step-space}; as the subscript of the norm $\|\cdot\|_{1,\beta}$ it is a H\"older exponent, and the bounds below are asserted for every $0\le\beta\le\beta_0<1$.
\end{itemize}

\emph{The two cut-offs.} The function \eqref{4.5:eq-near-term-frame-1} below is built with two cut-off functions.
\begin{itemize}
\item The wide cut-off $\tilde\zeta_\square\in C_0^\infty(\square_0)$ satisfies $\tilde\zeta_\square=1$ on $\tilde\square^{3}$ and $\tilde\zeta_\square=0$ outside $\tilde\square^{4}$.
\item The partition-of-unity member is
\begin{equation*}
\zeta_\square(x)=\prod_\mu\zeta\bigl(M^{-1}(x_\mu-y_\mu)\bigr),
\end{equation*}
where $\zeta(s)=1$ for $|s|\le 1/3$, $\zeta(s)=0$ for $|s|\ge 2/3$, and $y=(y_\mu)$ is the reference point of the cube $\square$ in the partition of unity of \cite{B-CMP109}; its support lies in $\tilde\square$.
\end{itemize}
The parameter $t$ multiplies the wide cut-off $\tilde\zeta_\square$, and the parameter $t_\square$ multiplies $\zeta_\square$.

\emph{The function.} The object is
\begin{equation}\label{4.5:eq-near-term-frame-1}
\mathbf{E}^{(j)}\Bigl(X,\,U_j\Bigl(\square_0,\,M^{\cdot}\Bigl(\exp i\eta\bigl(t\tilde\zeta_\square
+t_\square\zeta_\square\bigr)\mathbf{H}_k(B')\,U_{k+1}\bigl(\square_0,M^{\cdot}(\mathbf{U})\bigr)\Bigr)\Bigr)\Bigr).
\end{equation}
It is a well-defined analytic function of $(\mathbf{U},\mathbf{J})$ in a sufficiently small domain. Ba{\l}aban proves that it is analytic on the space
\begin{equation*}
\mathfrak{U}=U'^{\,c}_{k+1}\bigl(\square_0,(1+2\beta)\alpha_0,(1+2\beta)\alpha_1,\alpha_0\bigr).
\end{equation*}
This space is defined in \cite{B-CMP109} by the conditions (i)--(iv) with which \cite{B-CMP109} defines the corresponding unprimed space $U^{c}_{k+1}$, except that $\mathbf{U}$ and $\mathbf{J}$ are defined on the whole lattice $T_\eta$ and satisfy conditions (i)--(iii) there. The construction below is made on $\mathfrak{U}$.

\emph{Axial frame.} Ba{\l}aban \cite[p.~275]{B-CMP109} repeats the constructions of \cite[Sect.~F]{B-CMP102b}, and introduces the same generalized axial gauge for $M^{\cdot}(\mathbf{U})$ by a gauge transformation $v$:
\begin{equation}\label{4.5:eq-near-term-frame-2}
M^{\cdot}(\mathbf{U})=V^{v},\qquad |V(b)-1|<O(1)M\alpha_0,\qquad |v-1|<O(1)M\alpha_1 .
\end{equation}
The variables $V$ are expressed through the contour variables introduced in \cite[(77)]{B-CMP98}. In the simplest case, for a bond $b\in\tilde\square^{4}\cap T^{(k+1)}$ with end-points $b_-,b_+$, one has
\begin{equation*}
V(b)=\bar{\mathbf{U}}^{k+1}\bigl(\Gamma_{y_0,b_-}\cup b\cup\Gamma_{b_+,y_0}\bigr),
\end{equation*}
with $\Gamma_{y_0,b_-}$ and $\Gamma_{b_+,y_0}$ the contours of the axial gauge. The parameters $\alpha_0,\alpha_1$ are chosen so that $M\alpha_0$ and $M\alpha_1$ are still small.

\emph{Landau frame.} A second gauge transformation $u_{k+1}$ changes the axial gauge into the Landau gauge for $U_{k+1}(\square_0,V)$:
\begin{equation}\label{4.5:eq-near-term-frame-3}
U_{k+1}(\square_0,V)=\Bigl(\exp iL^{-1}\eta\,\mathbf{H}_{k+1}\bigl(\square_0,\tfrac1i\log V\bigr)\Bigr)^{u_{k+1}} .
\end{equation}
Moreover, for $0\le\beta\le\beta_0<1$,
\begin{equation}\label{4.5:eq-near-term-frame-4}
\begin{gathered}
\bigl|\mathbf{H}_{k+1}(\square_0,\tfrac1i\log V)\bigr|,\quad
\bigl|\nabla^{L^{-1}\eta}\mathbf{H}_{k+1}(\square_0,\tfrac1i\log V)\bigr|,\\
\qquad\bigl\|\mathbf{H}_{k+1}(\square_0,\tfrac1i\log V)\bigr\|_{1,\beta}
<B_3O(1)M\alpha_0\quad \text{ on }\tilde\square^{4},\\
|u_{k+1}-1|<B_3O(1)M\alpha_0\quad \text{ on }\square_0 .
\end{gathered}
\end{equation}

\emph{The transformed function.} Performing these gauge transformations in \eqref{4.5:eq-near-term-frame-1} gives
\begin{equation}\label{4.5:eq-near-term-frame-5}
\begin{aligned}
&\mathbf{E}^{(j)}\Bigl(X,U_j\Bigl(\square_0,M^{\cdot}\Bigl(\exp i\eta(t\tilde\zeta_\square+t_\square\zeta_\square)
\mathbf{H}_k(B')\,\exp iL^{-1}\eta\,\mathbf{H}_{k+1}\bigl(\square_0,\tfrac1i\log V\bigr)\Bigr)\Bigr)\Bigr)\\
&\quad=\mathbf{E}^{(j)}\Bigl(X,U_j\Bigl(\square_0,\exp iQ\Bigl(\tfrac1i\log\Bigl[
\exp i\eta(t\tilde\zeta_\square+t_\square\zeta_\square)\mathbf{H}_k(B')\\
&\qquad\qquad\times\exp iL^{-1}\eta\,\mathbf{H}_{k+1}\bigl(\square_0,\tfrac1i\log V\bigr)\Bigr]\Bigr)\Bigr)\Bigr).
\end{aligned}
\end{equation}
This step uses the gauge invariance of \eqref{4.5:eq-near-term-frame-1}. It therefore requires the gauge transformations to be applied also to the variables $\mathbf{J},B$, and not only to the configurations in \eqref{4.5:eq-near-term-frame-2} and \eqref{4.5:eq-near-term-frame-3}. Accordingly, in $\mathbf{H}_k(B')$ the variables $\mathbf{J},B$ are replaced by
\begin{equation*}
R(u_{k+1}^{-1})R(v^{-1})\mathbf{J},\qquad R(u_{k+1}^{-1})R(v^{-1})B .
\end{equation*}
Two remarks on these expressions. First, the bounds on $v$ and $u_{k+1}$ in \eqref{4.5:eq-near-term-frame-2} and \eqref{4.5:eq-near-term-frame-4} imply that these expressions obey almost the same bounds as $\mathbf{J},B$. Second, $v$ and $u_{k+1}$ are explicitly given analytic functions of $\mathbf{U}$ that depend only on the restriction of $\mathbf{U}$ to $\square_0$.

\emph{Gauge covariance.} The functions \eqref{4.5:eq-near-term-frame-1}, \eqref{4.5:eq-near-term-frame-2} are gauge invariant under the simultaneous transformations
\begin{equation}\label{4.5:eq-near-term-frame-6}
\mathbf{U}\to\mathbf{U}^{u},\qquad \mathbf{J}\to R(u)\mathbf{J},\qquad B\to R(u)B .
\end{equation}
Here $u$ ranges over the $G^{c}$-valued transformations in a sufficiently small neighbourhood of the $G$-valued ones, chosen so that the transformed configurations again belong to the proper spaces. Under \eqref{4.5:eq-near-term-frame-6} the expressions in \eqref{4.5:eq-near-term-frame-5} transform by $R(u(y_0))$.

\emph{The logarithm and its regularity.} Replace $(t\tilde\zeta_\square+t_\square\zeta_\square)\mathbf{H}_k(B')$ in \eqref{4.5:eq-near-term-frame-5} by a $\mathfrak{g}^{c}$-valued variable $\mathbf{A}$, and write $\hat{\mathbf{A}}$ for the logarithm that \cite[(3.30)]{B-CMP109} denotes by a lightface $A$. Thus
\begin{equation}\label{4.5:eq-near-term-frame-7}
\hat{\mathbf{A}}=\frac{1}{i\eta}\log\Bigl(\exp i\eta\mathbf{A}\,
\exp iL^{-1}\eta\,\mathbf{H}_{k+1}\bigl(\square_0,\tfrac1i\log V\bigr)\Bigr),\qquad
B=Q(\eta\hat{\mathbf{A}})\ \text{ on }\square_0 .
\end{equation}
The function $B$ so defined lives on the set of bonds that determines $U_j(\square_0)$, namely
\begin{equation*}
\Bigl(\bigcup_{n=0}^{j-1}\Lambda_n\Bigr)\cup\square_j^{(j)} .
\end{equation*}
The class of admissible $\mathbf{A}$ consists of the regular $\mathbf{A}$ satisfying, for $0\le\beta\le\beta_0<1$,
\begin{equation}\label{4.5:eq-near-term-frame-8}
|\mathbf{A}|,\ |\nabla^{\eta}\mathbf{A}|,\ \|\mathbf{A}\|_{1,\beta}<\alpha_2\ \text{ on }\square_0 .
\end{equation}
The function $(t\tilde\zeta_\square+t_\square\zeta_\square)\mathbf{H}_k(B')$ satisfies \eqref{4.5:eq-near-term-frame-8} for $\varepsilon_1$ sufficiently small. By \eqref{4.5:eq-near-term-frame-4}, $\mathbf{H}_{k+1}(\square_0,\frac1i\log V)$ satisfies the analogous bounds with $\alpha_2$ replaced by $B_3O(1)M\alpha_0$ and $\square_0$ replaced by $\tilde\square^{4}$. Assumption \eqref{4.5:eq-near-term-frame-8} implies
\begin{equation}\label{4.5:eq-near-term-frame-9}
|\hat{\mathbf{A}}|,\ |\nabla^{\eta}\hat{\mathbf{A}}|,\ \|\hat{\mathbf{A}}\|_{1,\beta}
<2\bigl(\alpha_2+B_3O(1)M\alpha_0\bigr)\ \text{ on }\tilde\square^{4}.
\end{equation}
These regularity conditions are the basis of the subsequent analysis of the near term.

This is \cite[(3.25)--(3.32), pp.~275--277]{B-CMP109}.
\end{remark}

\begin{definition}[The data-axial frame on the window]\label{4.5:not-axial-frame}
Throughout this definition $(\mathbf{U},\mathbf{J})\in\mathfrak{U}$ with $\mathbf{U}=e^{i\eta A'}U$, where $U$ is $G$-valued.
The constants are those of level $k+1$, namely $\alpha_0''':=(1+2\beta)\alpha_0$ and $\alpha_1''':=(1+2\beta)\alpha_1$.
Lengths are measured in the $\eta$-units of level $k$, in which the top block of the window has side $L$.
The cube $\square$, the window $\square_0$, its tower $\{\square_n\}$, the block-bond set $\mathfrak{B}(\square_0)$, the glued window background $S=\mathbf{U}^{\mathrm{gl}}\restriction\square_0$ and the gauge map $h$ are those of Proposition~\ref{4.5:prop-glued-window}.
The tower $M^{\cdot}\mathbf{U}$ is considered on $\mathfrak{B}(\square_0)$: on $\Lambda_n(\square_0)$ it consists of the $n$-fold
block averages, and on the lowest layer $\Lambda_0(\square_0)=P$, the pinned region of
Proposition~\ref{4.5:prop-glued-window}, it is the field itself.
For $n\le k+1$ put
\begin{equation*}
\ell_n:=L^{\,n-k-1}\le 1 ,
\end{equation*}
the relative block size, so that a block of level $n$ has side $L\ell_n$ in $\eta$-units.
Let $y_0$ be the centre of $\square$.
$T_\eta$ denotes the lattice torus of spacing $\eta$.
A bond $c$ runs from $c_-$ to $c_+$.
For a gauge transformation $u$ we write $U^{u}(c)=u(c_-)U(c)u(c_+)^{-1}$.
$\bar\Gamma$ is the path $\Gamma$ run backwards.

\smallskip
\noindent\emph{Contours.} These are the contours of the generalised axial gauge of \cite[Sect.~F]{B-CMP102b}, \cite[(77)]{B-CMP98} and \cite[(1.15), Sect.~F]{B-CMP99a}.

Let $x$ be a block site in the top region $\square_{k+1}$, and let $x_{k+1}$ be its top base point.
Then $\Gamma_{y_0,x}$ is a coordinate-ordered path of the top level from $y_0$ to $x_{k+1}$, followed by the nested descent
\begin{equation*}
\Gamma_{x_j,x_{j-1}},\qquad x_{j-1}\in B(x_j),
\end{equation*}
down to $x$, carried out inside the global $(k+1)$-block.
Here $B(x_j)$ denotes the set of sites of level $j-1$ in the block with base point $x_j$.

Now let $x\in\Lambda_p(\square_0)$ with $p\le k$ be a block site of the collar.
Such an $x$ has no base point of level $k+1$ in the window.
It does, however, lie in a global $(k+1)$-block $B^{k+1}(x_{k+1})\subset\square_0$ of $T_\eta$, the
$(k+1)$-block with base point $x_{k+1}$.
Indeed, $\mathfrak{U}$ imposes conditions (i)--(iii) on the whole lattice $T_\eta$, and $\square_0$ is a union of $(k+1)$-blocks.
For such $x$, $\Gamma_{y_0,x}$ is given by the same recipe.
It consists of the top-level path from $y_0$ to $x_{k+1}$, followed inside that global block by the nested descent \cite[(1.15)]{B-CMP99a}:
\begin{equation*}
\Gamma_{x_{k+1},x_k},\ \dots,\ \Gamma_{x_{p+1},x_p},\qquad x_p=x .
\end{equation*}
Thus there is one top-level tree over the whole of $\square_0$, collar included, followed by the nested descent.

This is Ba{\l}aban's own tree.
In \cite[Sect.~F, p.~98]{B-CMP99a} the configuration $U_0'$ satisfies the axial gauge conditions
\cite[(1.15)]{B-CMP99a}.
Its average $\bar U_0'^{\,k}$ satisfies the global axial gauge conditions $\bar U_0'^{\,k}(\Gamma_{y,x})=1$ for $x\in\tilde\square^{(k)}$, where $y$ is a centre of $\tilde\square^{(k)}$.
In \cite[(146)]{B-CMP102b} the conditions are imposed for bonds $\langle x,x'\rangle\subset\tilde\square^{(j)}$, $j=0,1,\dots,k$.
The estimates \cite[(145)--(151)]{B-CMP102b} are stated for these contours.

\smallskip
\noindent\emph{Leg counts.} The top path has at most $6dM/L$ top bonds.
Each nested leg has at most $dL$ bonds of its level $n$, each of $\eta$-length $L\ell_n$.
These counts hold for every block site, top region and collar alike; no leg runs sideways inside a collar layer.

\smallskip
\noindent\emph{Convention off the block-bond set.} Let $c$ be a bond of level $q$ of a global $q$-block meeting $\square_0$, and suppose $c$ is not a bond of $\mathfrak{B}(\square_0)$.
These are the bonds of the top path and of the upper legs $q>p$ of the contour of a collar site.
On such $c$ we put
\begin{equation*}
(M^{\cdot}\mathbf{U})(c):=\bar{\mathbf{U}}^{q}(c),
\end{equation*}
the global $q$-average.
On the bonds of $\mathfrak{B}(\square_0)$, the global $q$-average is the bond variable of the window datum itself.
The pointwise bound on $v^{\pm1}$ derived in this section from the levelwise bounds on block averages
holds on these legs verbatim, with the same constant: those levelwise bounds are applied at each
level $q$, and they apply to the global $q$-averages on these legs as to the averages on
$\mathfrak{B}(\square_0)$.

\smallskip
\noindent\emph{The frame map and the axial datum.} For a block site $x\in\Lambda_p(\square_0)$, top region or collar, put
\begin{equation}\label{4.5:eq-axial-frame-1}
\begin{split}
v(x):=(M^{\cdot}\mathbf{U})(\Gamma_{y_0,x})
&=\bar{\mathbf{U}}^{k+1}(\Gamma_{y_0,x_{k+1}})\,\bar{\mathbf{U}}^{k}(\Gamma_{x_{k+1},x_k})\\
&\qquad\cdots\bar{\mathbf{U}}^{p}(\Gamma_{x_{p+1},x}).
\end{split}
\end{equation}
This is the ordered product, along $\Gamma_{y_0,x}$, of the global averages $\bar{\mathbf{U}}^{q}$, $q\ge p$, of $\mathbf{U}\restriction\square_0$.
On each leg of level $q$ the $q$-average is used; on bonds of the window it is the bond variable of the datum.
Inverses are taken on backward bonds.
Bonds of the datum enter only where $\Gamma_{y_0,x}$ runs in $\Lambda_q(\square_0)$.

On $\mathfrak{B}(\square_0)$ put
\begin{equation}\label{4.5:eq-axial-frame-2}
V:=(M^{\cdot}\mathbf{U})^{v},\qquad\text{that is,}\qquad V(c)=v(c_-)\,(M^{\cdot}\mathbf{U})(c)\,v(c_+)^{-1}.
\end{equation}
For $c\in\mathfrak{B}(\square_0)$ let $\mathrm{hol}(c)$ be the ordered product of the same global averages around the loop $\Gamma_{y_0,c_-}\cup c\cup\bar\Gamma_{y_0,c_+}$.

Ba{\l}aban writes $M^{\cdot}(\mathbf{U})=V^{v}$ in \cite[(3.26)]{B-CMP109}.
The transformation denoted $v$ there is therefore $v^{-1}$ in the present notation.
Likewise, for the induced transformation of the currents, written $R(\cdot)$ in
\cite[p.~276]{B-CMP109}, the operator $R(v^{-1})$ there is $R(v)$ here.
For this reason, every bound on $v$ in what follows is stated for $v$ and $v^{-1}$ alike.

\smallskip
\noindent\emph{Real data.} Let $v_G(x)$ be the same product formed for $U$:
\begin{equation*}
v_G(x)=\bar U^{k+1}(\Gamma_{y_0,x_{k+1}})\cdots\bar U^{p}(\Gamma_{x_{p+1},x})\in G,
\end{equation*}
and put $V_G:=(M^{\cdot}U)^{v_G}$.
At real data, $V_G$ is the window datum in Ba{\l}aban's generalised axial gauge, the configuration
denoted $V''$ in his construction of that gauge in \cite[Sect.~F]{B-CMP102b}.
In that construction, $V_G(c)=\bar U'^{\,p}(c)$ for $c\in\Lambda_p(\square_0)$, where $U':=U^{u'}$ and
$u'$ is the generalised axial gauge map of the construction; the map $u'$ coincides with $v_G$ at
block base sites.

Next, $g$ denotes the gauge transformation written $u_{k+1}$ in \cite[(3.27)]{B-CMP109}.
It takes the window configuration from Landau gauge to the axial gauge, with base value $1$.
That is,
\begin{equation*}
U_{k+1}(\square_0,V)=\bigl(e^{i\eta\mathcal{A}}\bigr)^{g},
\end{equation*}
where $\mathcal{A}:=L^{-1}\mathbf{H}_{k+1}(\square_0,\tfrac1i\log V)$ is the Landau-gauge potential of
\cite[(3.27)]{B-CMP109}.

\smallskip
\noindent\emph{The twist polydisc.} Define
\begin{equation}\label{4.5:eq-axial-frame-3}
\begin{split}
\mathcal{P}_U:=\Bigl\{&A'\ \mathfrak{g}^{c}\text{-valued on the bonds of }T_\eta:\\
&\qquad L\sup|A'|<\alpha_1''',\ \ L^{2}\sup|\nabla^{\eta}_{U}A'|<\alpha_1'''\Bigr\},
\end{split}
\end{equation}
in $\eta$-units.
In the units $\xi=L^{-1}\eta$ of level $k+1$ of \cite[(1.13)]{B-CMP109} at $j=k+1$, the two conditions read $|A'|,|\nabla^{\xi}_{U}A'|<\alpha_1'''$.
Thus $\mathcal{P}_U$ is the set of twists of \cite[(1.13)]{B-CMP109} at level $k+1$ over the fixed $G$-valued $U$.
Condition~(iii) of $\mathfrak{U}$ is not part of its definition.
For $A'\in\mathcal{P}_U$ we write $\mathbf{U}_{A'}:=e^{i\eta A'}U$.

These objects have the following properties.
\begin{enumerate}
\item[(a)] $v$ depends on $\mathbf{U}\restriction\square_0$ only, and $v(y_0)=1$.
\item[(b)] $V=1$ on the contour bonds lying in $\mathfrak{B}(\square_0)$, and $V(c)=\mathrm{hol}(c)$ for every other $c\in\mathfrak{B}(\square_0)$.
In particular, for a bond $b$ of the top region this is Ba{\l}aban's formula
\eqref{4.5:eq-axial-frame-4} for his ``simplest situation'' in \cite[p.~275]{B-CMP109}.
\item[(c)] $\mathcal{P}_U$ is an open convex subset of a complex vector space of finite dimension.
It contains $0$ and the given $A'$.
Its intersection with the real subspace of $\mathfrak{g}$-valued $A'$ is a non-empty open subset of that subspace.
\end{enumerate}
\end{definition}

\begin{proof}
(a) Every block site $x$ of $\square_0$, top region or collar, lies in a global $(k+1)$-block contained in $\square_0$, because $\square_0$ is a union of $(k+1)$-blocks.
The contour $\Gamma_{y_0,x}$ consists of the top path and the nested descent inside that block.
The coordinate-ordered top path joins two points of the cube $\square_0$ and therefore stays in
$\square_0$; each $(k+1)$-block it meets is contained in $\square_0$, again because $\square_0$ is a
union of $(k+1)$-blocks.
Hence every global average entering \eqref{4.5:eq-axial-frame-1} is an average of $\mathbf{U}\restriction\square_0$.
These averages are defined because $\mathfrak{U}$ imposes (i)--(iii) on all of $T_\eta$.
The contour $\Gamma_{y_0,y_0}$ is empty, so the product in \eqref{4.5:eq-axial-frame-1} is empty and $v(y_0)=1$.

(b) Let $c\in\mathfrak{B}(\square_0)$.
Since inverses are taken on backward bonds, the ordered product along $\bar\Gamma$ is the inverse of the ordered product along $\Gamma$.
Hence, by \eqref{4.5:eq-axial-frame-1},
\begin{equation*}
v(c_+)^{-1}=\bigl((M^{\cdot}\mathbf{U})(\Gamma_{y_0,c_+})\bigr)^{-1}=(M^{\cdot}\mathbf{U})(\bar\Gamma_{y_0,c_+}).
\end{equation*}
Inserting this and \eqref{4.5:eq-axial-frame-1} for $c_-$ into \eqref{4.5:eq-axial-frame-2}, we obtain
\begin{equation*}
\begin{split}
V(c)&=(M^{\cdot}\mathbf{U})(\Gamma_{y_0,c_-})\,(M^{\cdot}\mathbf{U})(c)\,
(M^{\cdot}\mathbf{U})(\bar\Gamma_{y_0,c_+})\\
&=(M^{\cdot}\mathbf{U})\bigl(\Gamma_{y_0,c_-}\cup c\cup\bar\Gamma_{y_0,c_+}\bigr)=\mathrm{hol}(c).
\end{split}
\end{equation*}

If $c$ is a contour bond, the contours form one tree rooted at $y_0$.
The contour of one end-point of $c$ is then the contour of the other, extended by $c$ in the direction of the tree.
Therefore the loop $\Gamma_{y_0,c_-}\cup c\cup\bar\Gamma_{y_0,c_+}$ retraces itself, the ordered product collapses, and $V(c)=1$.

For a bond $b$ of the top region, all global averages on the loop are of level $k+1$.
The loop formula then reads
\begin{equation}\label{4.5:eq-axial-frame-4}
V(b)=\bar{\mathbf{U}}^{k+1}(\Gamma_{y_0,b_-}\cup b\cup\Gamma_{b_+,y_0}),
\end{equation}
with $\Gamma_{b_+,y_0}=\bar\Gamma_{y_0,b_+}$.

Finally, \eqref{4.5:eq-axial-frame-2} gives $M^{\cdot}\mathbf{U}=V^{v^{-1}}$.
Compared with $M^{\cdot}(\mathbf{U})=V^{v}$ in \cite[(3.26)]{B-CMP109}, this shows that the transformation $v$ there is $v^{-1}$ here, as stated in the definition.

(c) For fixed $U$, both $A'\mapsto A'$ and $A'\mapsto\nabla^{\eta}_{U}A'$ are complex-linear.
Hence $A'\mapsto L\sup|A'|$ and $A'\mapsto L^{2}\sup|\nabla^{\eta}_{U}A'|$ are seminorms.
The strict sublevel set of a seminorm is open and convex.
Therefore $\mathcal{P}_U$, the intersection of two such sets, is open and convex.

The torus $T_\eta$ is finite and $\mathfrak{g}^{c}$ has finite dimension.
Hence $A'$ ranges over a complex vector space of finite dimension, that is, $\mathcal{P}_U$ involves finitely many complex coordinates.

Both seminorms vanish at $0$, and $\alpha_1'''>0$, so $0\in\mathcal{P}_U$.
Since $(\mathbf{U},\mathbf{J})\in\mathfrak{U}$ and $\mathbf{U}=e^{i\eta A'}U$, the twist $A'$ obeys the bounds of \cite[(1.13)]{B-CMP109} at $j=k+1$ with constant $\alpha_1'''$.
In the units $\xi=L^{-1}\eta$ these are $|A'|,|\nabla^{\xi}_{U}A'|<\alpha_1'''$.
Written in $\eta$-units, these are exactly the two inequalities in \eqref{4.5:eq-axial-frame-3}; hence $A'\in\mathcal{P}_U$.

The intersection of the open set $\mathcal{P}_U$ with the real subspace of $\mathfrak{g}$-valued $A'$ is open in the relative topology.
It contains $0$, so it is non-empty.
\end{proof}

\begin{lemma}[The literal frame: axial data, frame maps and base-one response]\label{4.5:lem-literal-frame}
Work in the setting of Notation~\ref{4.5:not-axial-frame}, with the conventions of Notation~\ref{4.5:not-geometry-conventions} and the radius $\bar\varepsilon_1$ of~\eqref{4.5:eq-geometry-conventions-a}. Thus:
\begin{itemize}
\item $(\mathbf{U},\mathbf{J})\in\mathfrak{U}$ and $\mathbf{U}=e^{i\eta A'}U$ with $U$ $G$-valued;
\item $\alpha_0'''=(1+2\beta)\alpha_0$ and $\alpha_1'''=(1+2\beta)\alpha_1$ are the level-$(k+1)$ radii;
\item $\square_0$ is the window, with tower $\{\square_n\}$, layers $\Lambda_n(\square_0)$, $\Lambda_0(\square_0)=P$, and hierarchy $\mathfrak{B}(\square_0)$;
\item $\ell_n=L^{n-k-1}\le1$ are the relative block sizes, and $y_0$ is the centre of $\square$;
\item $\Gamma_{y_0,x}$ are the contours, $v,v_G$ the frame maps, $V=(M^{\cdot}\mathbf{U})^{v}$ and $V_G=(M^{\cdot}U)^{v_G}$ the axial data, and $g$ the Landau-to-axial map of the window at base one;
\item $\mathcal{P}_U$ is the twist polydisc, $S=\mathbf{U}^{\mathrm{gl}}\restriction\square_0$, and $h$ is the gauge map of the setting.
\end{itemize}
Let $c_Q$ and $\alpha_1^{(98)}$ be the constant of the level sizes and the radius of \cite[Proposition~6]{B-CMP98}, $C_{98}$ the absolute constant of Remark~\ref{4.5:rem-block-curvature-complex}, $c_1$ the constant of \cite[Lemma~2.1]{B-CMP96}, $C_H$ the constant of the kernel bound in hypothesis (A) below, and $C_1$ the $(d,L)$-constant that enters, together with $\bar\varepsilon_1$, the ceilings (c7) and (c8) below. Put
\begin{gather*}
C_v:=8dc_Q,\qquad C_V:=\max\{2^8d^2L^2,\,12C_v\},\qquad B_b:=6dLc_1C_HC_V,\qquad C_{14}:=9(2d+2),
\end{gather*}
all depending on $(d,L)$ only. Let $C_7\ge2$ be the constant implicit in the statement of \cite{B-CMP102b} that, if the space \cite[(6)]{B-CMP102b} is non-empty and $\varepsilon_0$ is sufficiently small, then by \cite[Proposition~2]{B-CMP98} the configuration $V$ there satisfies \cite[(7)]{B-CMP102b} with $\varepsilon_1=O(\varepsilon_0)$; thus $V$ satisfies \cite[(7)]{B-CMP102b} with $\varepsilon_1=C_7\varepsilon_0$. The lower bound $2$ is the constant of \cite[(54)]{B-CMP98}. Write
\begin{equation*}
b:=B_bM\alpha_0''',\qquad \mathfrak{v}:=C_VM\alpha_0''',\qquad b^{\mathrm{poly}}:=B_bM(\alpha_0'''+\alpha_1''').
\end{equation*}
Since $\alpha_0'''=(1+2\beta)\alpha_0$, the quantity $b$ carries the factor $1+2\beta$; accordingly $B_b$ is placed after $\beta$ in the order of choice.
Assume the ceilings, all imposed after $M$ is fixed (the third member of (c7$'$) is a ceiling on $\alpha_0$ after $M/L$):
\begin{equation}\label{4.5:eq-literal-frame-a}
\begin{aligned}
&\text{(c7)}\quad 4C_VM(\alpha_0'''+\alpha_1''')\le C_1\bar\varepsilon_1;\\
&\text{(c7$'$)}\quad 2C_vM\alpha_1'''\le\tfrac16,\quad 4dLc_QB_bM\alpha_0'''\le\tfrac16,
\quad B_b\tfrac{M}{L}\alpha_0'''\le\alpha_1^{(98)};\\
&\text{(c8)}\quad c_Q\alpha_1'''\le\tfrac12C_1\bar\varepsilon_1,\quad
C_7(\alpha_0'''+C_{14}\alpha_1''')\le\bar\varepsilon_1,\quad
\alpha_1'''\le\min\{\alpha_1^{(98)},1\}.
\end{aligned}
\end{equation}
The third member of (c7$'$) says $b/L\le\alpha_1^{(98)}$, the radius of \cite[Proposition~6]{B-CMP98} for the averaged legs of $g$ at every layer of $\square_0$. Assume also $\bar\varepsilon_1\le1/(4C_1)$ and the lower bound
\begin{equation*}
L\ge2\bigl(1+O_{\mathrm{CL\text{-}U}}(1/24)\bigr)
\end{equation*}
on $L$, where $O_{\mathrm{CL\text{-}U}}(x)$ denotes the supremum, over $0<dc_Qb\le x$, of the relative Baker--Campbell--Hausdorff remainder in the logarithm bound for $g$ in the proof of (AX3) below; it depends neither on $c_Q$ nor on $L$. Assume further:
\begin{itemize}
\item[(A)] the corollary on derivative bounds of Ba{\l}aban's background function, applied for the tower $\{\square_n\}$ at base one;
\item[(B)] the logarithm bound for the Landau-to-axial map: for a chart function $\mathcal{H}$ with Landau-to-axial map $g$, a point $x$ and the block $B^p(x_p)$ of level $p$ containing $x$, $|\log g(x)|\le4dc_Q\ell_p\sup_{B^p(x_p)}|\mathcal{H}|$; at base one, the case used here, it is verified in the course of the proof of (AX3) below;
\item[(C)] for the clauses (AX1-$\mathbf{U}$) and (AX2-pt) only, clauses (a) and (b) of Remark~\ref{4.5:rem-block-curvature-complex} at $\mathbf{U}$, with constants $(\alpha_0''',\alpha_1''')$; clause (a) is stated in \cite[p.~263]{B-CMP109} and clause (b) in \cite[(3.26)]{B-CMP109}, each without proof.
\end{itemize}
Then the following hold.
\begin{itemize}
\item[(AX1)] \emph{Data.} On all of $\mathcal{P}_U$,
\begin{equation}\label{4.5:eq-literal-frame-b}
M^{\cdot}\mathbf{U}=e^{iB'}\,M^{\cdot}U\ \text{ on }\mathfrak{B}(\square_0),\qquad
|B'(c)|\le c_Q\alpha_1'''\ell_n\quad(c\in\Lambda_n(\square_0)).
\end{equation}
At $n=0$, on $P$, the factor $e^{iB'}$ is the fine-bond twist factor of $\mathbf{U}=e^{i\eta A'}U$ itself.

Let $c$ be a level-$q$ bond of a global $q$-block meeting $\square_0$ that is not a bond of $\mathfrak{B}(\square_0)$; these are the contour legs of the setting. On every such bond, likewise,
\begin{equation*}
\bar{\mathbf{U}}^{q}(c)=e^{iB'_q(c)}\bar U^{q}(c),\qquad |B'_q(c)|\le c_Q\alpha_1'''\ell_q .
\end{equation*}
The within-level block plaquettes $p'$ of $M^{\cdot}U$ obey
\begin{equation*}
|\partial(M^{\cdot}U)(p')-1|<2\alpha_0'''\ell_n^2\le2\alpha_0''' .
\end{equation*}
\item[(AX1-$\mathbf{U}$)] At the given $\mathbf{U}$, for which condition (iii) in the definition of $\mathfrak{U}$ holds, the within-level block plaquettes obey
\begin{equation*}
|\partial(M^{\cdot}\mathbf{U})(p')-1|\le C_{98}\alpha_0'''\ell_n^2 .
\end{equation*}
Condition (iii) is not part of the definition of $\mathcal{P}_U$. The seam clause of \cite[(7)]{B-CMP102b} for $M^{\cdot}\mathbf{U}$ is neither asserted nor used.
\item[(AX2)] \emph{Frame maps.} On all of $\mathcal{P}_U$,
\begin{equation*}
\|(vv_G^{-1})^{\pm1}-1\|\le e^{C_vM\alpha_1'''}-1,\qquad
\|\operatorname{Ad}_v^{\pm1}\|\le e^{2C_vM\alpha_1'''}.
\end{equation*}
Only the non-unitary part of $v$ is bounded; the factor $v_G\in G$ is not small. This is the form taken here by the bound on the gauge map in \cite[(3.26)]{B-CMP109}.
\item[(AX2-poly)] For every $A'\in\mathcal{P}_U$ and every $c\in\mathfrak{B}(\square_0)$,
\begin{equation*}
|V(c)-1|\le C_VM(\alpha_0'''+\alpha_1''').
\end{equation*}
\item[(AX2-pt)] At the given $\mathbf{U}$, $|V(c)-1|\le\mathfrak{v}=C_VM\alpha_0'''$ for every $c\in\mathfrak{B}(\square_0)$. Only $\alpha_0'''$ enters this bound.
\item[(AX3)] \emph{Response at base one.} The field $B_V:=\frac1i\log V$ is defined on all of $\mathcal{P}_U$, with $\sup|B_V|\le\frac12C_1\bar\varepsilon_1$. Hence $B_V$ lies, with a margin, in the polydisc \cite[(172)]{B-CMP102b} over the base $V_0:=1$ of the window tower, whose minimiser is $U_0:=U_{k+1}(\square_0,1)=1$. At the given $\mathbf{U}$, $\sup|B_V|\le2\mathfrak{v}$.

Let $\mathcal{H}^{\square_0}$ be the function of \cite[Proposition~9]{B-CMP102b} for the window, in its own units $\xi=L^{-1}\eta$. Then
\begin{equation*}
U_{k+1}(\square_0,V)=(e^{i\eta\mathcal{A}})^{g},\qquad \mathcal{A}:=L^{-1}\mathcal{H}^{\square_0}(B_V);
\end{equation*}
in \cite[(3.27)]{B-CMP109} this function is written $L^{-1}\mathbf{H}_{k+1}(\square_0,\frac1i\log V)$. At the given $\mathbf{U}$, on $\tilde\square^{4}$,
\begin{equation*}
|\mathcal{A}|,\ |\nabla^{\eta}\mathcal{A}|,\ \|\mathcal{A}\|_{1,\beta}\le b=B_bM\alpha_0'''.
\end{equation*}
On a block of a collar layer $\Lambda_q(\square_0)$ the same holds in the unit of that layer:
\begin{equation*}
\ell_q\sup|\mathcal{H}^{\square_0}(B_V)|\le3dLC_Hc_1\sup|B_V|\le3dLC_Hc_1\cdot2C_VM\alpha_0'''=b .
\end{equation*}
Moreover, on $\square_0$, at every layer alike,
\begin{equation*}
|\log g|\le2dLc_Qb,\qquad |g-1|\le e^{2dLc_Qb}-1,\qquad \|\operatorname{Ad}_g^{\pm1}\|\le e^{4dLc_Qb}.
\end{equation*}
On all of $\mathcal{P}_U$ the same bounds hold with $b$ replaced by $b^{\mathrm{poly}}$.
\item[(AX4)] \emph{The frame.} Let $y_x$ be the base site of the block containing $x$, and let $\tilde v(x):=v(y_x)$ be the extension of $v$ from the base sites that is constant on blocks, as in \cite[(181)]{B-CMP102b}.
\begin{itemize}
\item[(i)] The following equality holds:
\begin{equation*}
U_{k+1}(\square_0,V)=S^{\tilde v}.
\end{equation*}
This is an equality between two charts of \cite[Proposition~9]{B-CMP102b} at two bases. The left side is the chart at base $1$ at the datum $V=e^{iB_V}\cdot1$. The right side is the chart at the real base $M^{\cdot}U$ at the datum $M^{\cdot}\mathbf{U}=e^{iB'}M^{\cdot}U$, transformed by the block-constant map $\tilde v$. It is proved by the identity theorem on $\mathcal{P}_U$; the identity \cite[(181)]{B-CMP102b} for the block-constant extension is used at $G$-valued maps only.
\item[(ii)] Let $w:=g^{-1}\circ\tilde v$ be the site map $x\mapsto g(x)^{-1}\tilde v(x)$: first $\tilde v$, then $g^{-1}$. Since $(U^{a})^{c}=U^{ca}$, on the bonds of $\square_0$
\begin{equation*}
S^{w}=(S^{\tilde v})^{g^{-1}}=e^{i\eta\mathcal{A}}.
\end{equation*}
In the notation of \cite[p.~276]{B-CMP109}, where the gauge map written $v$ there is $v^{-1}$ here and $u_{k+1}$ is $g$, the operator $R(w)$ is $R(u_{k+1}^{-1})R(v^{-1})$. This is exactly the replacement made there.
\item[(iii)] At every site of $\square_0$,
\begin{equation*}
\|\operatorname{Ad}_w^{\pm1}\|\le e^{4dLc_Qb+2C_vM\alpha_1'''}\le e^{1/3}<2 .
\end{equation*}
\item[(iv)] $w$ is an explicit analytic function of $\mathbf{U}\restriction\square_0$.
\item[(v)] Let $\tilde v_G$ be any $G$-valued extension of the block-constant map of $v_G$ (for instance by $1$). Let $w^{\mathrm{ext}}$ equal $w$ on the sites of $\square_0$, end-points of its bonds included, and $\tilde v_G$ off $\square_0$, and let $\bar w^{\mathrm{ext}}$ be its restriction to block sites. Then $\|\operatorname{Ad}_{w^{\mathrm{ext}}}^{\pm1}\|\le2$ at every site, and $\|\operatorname{Ad}_{w^{\mathrm{ext}}h^{-1}}^{\pm1}\|\le4=\mathfrak{K}$.

Let $t$ multiply the cut-off $\tilde\zeta_\square$ and $t_\square$ the partition-of-unity function $\zeta_\square$ of \cite[(3.25)]{B-CMP109}. Let $H_k(B';\cdot,\cdot)$ be the response chart of Theorem~\ref{4.5:thm-chart-existence} at the datum $B'$. By clause (K5) of that theorem, the gauge covariance of the chart, on $\square_0$,
\begin{equation}\label{4.5:eq-literal-frame-1}
\begin{split}
\mathbf{A}&:=(t\tilde\zeta_\square+t_\square\zeta_\square)\,R(w)H_k(B';\mathbf{U}^{\mathrm{gl}},\mathbf{J})\\
&=(t\tilde\zeta_\square+t_\square\zeta_\square)\,
H_k\bigl(R(\bar w^{\mathrm{ext}})B';(\mathbf{U}^{\mathrm{gl}})^{w^{\mathrm{ext}}},R(w^{\mathrm{ext}})\mathbf{J}\bigr).
\end{split}
\end{equation}
The identity is algebraic. $\mathbf{J}$ and $B'$ enter only through their suprema and no derivative of $w^{\mathrm{ext}}$ is taken, so the roughness of $w^{\mathrm{ext}}$ at the boundary of $\square_0$ does no harm. This is the function $(t\tilde\zeta_\square+t_\square\zeta_\square)\mathbf{H}_k(B')$ of \cite[(3.28)]{B-CMP109}, in which $\mathbf{J}$ and $B$ are replaced by $R(u_{k+1}^{-1})R(v^{-1})\mathbf{J}$ and $R(u_{k+1}^{-1})R(v^{-1})B$. Moreover $\operatorname{supp}\mathbf{A}\subset\tilde\square^{4}$. Both charts exist: for $|B'|<c_9/8$,
\begin{equation*}
|R(\bar w^{\mathrm{ext}})B'|\le2|B'|<c_9/4<c_9 .
\end{equation*}
\end{itemize}
\item[(AX5)] Put $C_N:=600C_9^K$. For the chart at the value $s\equiv1$ of the localisation parameters $s$ of clause (K10) of Theorem~\ref{4.5:thm-chart-existence}, on $\square_0$,
\begin{equation*}
|\mathbf{A}|,\ |\nabla^{\eta}\mathbf{A}|,\ \|\mathbf{A}\|_{1,\beta}\le C_N(|t|+|t_\square|)\sup|B'| .
\end{equation*}
For the $s$-localised chart of clause (K10) of Theorem~\ref{4.5:thm-chart-existence}, the same holds with $C_N^{s}:=C_Ne^{c_{\mathbb{K}}\kappa_1}$.

Consequently, let $|t|\le1$ and let the datum have sup-radius $\varepsilon_1':=\sup|B'|$. Assume the ceiling $3C_N^{s}\varepsilon_1'\le\alpha_2$ on the fluctuation radius of \cite{B-CMP109}, \cite{B-CMP116}, imposed after $\alpha_2$ and $\kappa_1$. Assume further that
\begin{equation*}
|t_\square|\le r_\square(B'):=\frac{\alpha_2}{3C_N^{s}\sup|B'|}.
\end{equation*}
Then, on $\square_0$,
\begin{equation*}
|\mathbf{A}|,\ |\nabla^{\eta}\mathbf{A}|,\ \|\mathbf{A}\|_{1,\beta}
\le\tfrac13\alpha_2+\tfrac13\alpha_2<\alpha_2 .
\end{equation*}
This is the bound \cite[(3.31)]{B-CMP109}.
\item[(AX6)] Let $\hat{\mathbf{A}}:=\frac1{i\eta}\log(e^{i\eta\mathbf{A}}e^{i\eta\mathcal{A}})$; this is the field $A$ of \cite[(3.30)]{B-CMP109}. Assume $\alpha_2+b\le1/(4C_{\mathrm{BCH}})$. Then, on $\tilde\square^{4}$,
\begin{equation*}
|\hat{\mathbf{A}}|,\ |\nabla^{\eta}\hat{\mathbf{A}}|,\ \|\hat{\mathbf{A}}\|_{1,\beta}
<2(\alpha_2+b)=2(\alpha_2+B_bM\alpha_0''').
\end{equation*}
This is \cite[(3.32)]{B-CMP109}. The same holds for the $s$-localised chart.
\end{itemize}
\end{lemma}

The fractions $\frac13$ in clause (AX5) of Lemma~\ref{4.5:lem-literal-frame} belong to the ceilings of that lemma; the corresponding assumption in \cite{B-CMP116} uses $\frac14$ and $\frac12$. The radius $r_\square(B')$ depends on the field and is of the kind of \cite[(1.22)]{B-CMP116}, but smaller than the circle taken there at equal constants. If $r_\square(B')$ is used for that circle, the factor $8$ of \cite[(1.24)]{B-CMP116} becomes $12$.

The proof below establishes clauses (AX1)--(AX3). Clauses (AX4)--(AX6) are proved separately, after this proof.

\begin{proof}[Proof of clauses \textup{(AX1)}--\textup{(AX3)}]
\emph{Clause} (AX1). Let $A'\in\mathcal{P}_U$. Then $L\sup|A'|<\alpha_1'''$, so
\begin{equation*}
\sup|A'|<\alpha_1'''/L\le\alpha_1^{(98)}
\end{equation*}
by (c8). We apply \cite[Proposition~6, (164)]{B-CMP98} at each averaging level over the $G$-valued field $U$, writing $e^{i\eta A'}U$ as a twist of $U$. The hypothesis of that proposition reads only the supremum of a complexified-algebra-valued $A'$; it involves no derivative and no curvature. It therefore applies on all of $\mathcal{P}_U$.

The size of the level-$n$ logarithm is amplitude times block side. This follows from the level sizes \cite[(159), (161)]{B-CMP98}, which bound the logarithm of the $j$-th averaged twist by $2\alpha_1L^j\eta$ in the units of that paper. It also uses \cite[(163)]{B-CMP98}, which puts the conjugating factor of \cite[(159)]{B-CMP98} within $O(1)\alpha_1$ of $1$. Since the $\eta$-side of a level-$n$ block is $L\ell_n$, for $c\in\Lambda_n(\square_0)$
\begin{equation*}
|B'(c)|\le c_Q\cdot\frac{\alpha_1'''}{L}\cdot L\ell_n=c_Q\alpha_1'''\ell_n .
\end{equation*}
At $n=0$ the average is the field itself, so $e^{iB'}$ is the fine-bond twist factor.

The same argument applies verbatim to the global $q$-averages $\bar{\mathbf{U}}^{q}(c)=e^{iB'_q(c)}\bar U^{q}(c)$, $q\le k+1$, on the contour legs off $\mathfrak{B}(\square_0)$. The field and the levels are the same, so $|B'_q(c)|\le c_Q\alpha_1'''\ell_q$.

The plaquette bound for $M^{\cdot}U$ is \cite[Proposition~2]{B-CMP98} at the $G$-valued $U$. It is applied layer by layer from condition (i) of $\mathfrak{U}$ at $\alpha_0'''$. It bounds the deviation of a block plaquette from $1$ by $\alpha_0'''$ plus a term quadratic in $\alpha_0'''$, hence by $2\alpha_0'''$ in the normalisation of that paper. In ours this reads $|\partial(M^{\cdot}U)(p')-1|<2\alpha_0'''\ell_n^2\le2\alpha_0'''$, since $\ell_n\le1$.

Clause (AX1-$\mathbf{U}$) is clause (a) of Remark~\ref{4.5:rem-block-curvature-complex}, with $(\alpha_0''',\alpha_1''')$ in place of its constants. This is admissible since $\alpha_1'''\le\alpha_1^{(98)}$ by (c8). Clause (a) of that remark is stated in \cite[p.~263]{B-CMP109} without proof.

\emph{Clause} (AX2). By the setting, $v(x)$ is the ordered product along $\Gamma_{y_0,x}$ of the global averages of $\mathbf{U}$, with inverses on backward bonds. By (AX1), with $q(c)$ the level of the leg carrying $c$,
\begin{equation*}
v(x)=\prod_{c\in\Gamma_{y_0,x}}\bigl(e^{iB'_{q(c)}(c)}\bar U^{q(c)}(c)\bigr)^{\pm1}
\end{equation*}
in path order. We move every factor $e^{\pm iB'}$ to the left through the unitary $\bar U$'s, using $\bar Ue^{iY}=e^{i\operatorname{Ad}(\bar U)Y}\bar U$ for $Y\in\mathfrak{g}^{c}$. This gives
\begin{equation*}
v=\Bigl[\prod_c e^{\pm i\operatorname{Ad}(u_c)B'(c)}\Bigr]v_G
\end{equation*}
with unitary $u_c$. Since $\operatorname{Ad}$ of a unitary preserves $|\cdot|$, each factor has the form $1+E_c$ with $\|E_c\|\le e^{|B'(c)|}-1$, and the same holds for the inverses. Hence
\begin{equation}\label{4.5:eq-literal-frame-2}
\|(vv_G^{-1})^{\pm1}-1\|\le\prod_c(1+\|E_c\|)-1\le\exp\Bigl(\sum_{c\in\Gamma}|B'(c)|\Bigr)-1 ,
\end{equation}
and in particular $\|(vv_G^{-1})^{\pm1}\|\le e^{\sum_c|B'(c)|}$.

By the leg counts of the setting, the top path has at most $6dM/L$ bonds, each with $\ell_{k+1}=1$. Each nested leg of level $n\le k$ has at most $dL$ bonds. Therefore
\begin{align*}
\sum_{c\in\Gamma}|B'(c)|
&\le c_Q\alpha_1'''\Bigl[\frac{6dM}{L}+dL\sum_{n\le k}\ell_n\Bigr]
\le c_Q\alpha_1'''\Bigl[\frac{6dM}{L}+4d\Bigr]\\
&\le8dc_Q\frac{M}{L}\alpha_1'''\le C_vM\alpha_1''' .
\end{align*}
Here $\sum_{n\le k}\ell_n\le2/L$, and $4d\le2dM/L$ because $M>L^2M_1\ge2L$.

Finally, $v=(vv_G^{-1})v_G$ and $\operatorname{Ad}_{v_G}$ is isometric, while $\|\operatorname{Ad}_a\|\le\|a\|\,\|a^{-1}\|$. Hence
\begin{equation*}
\|\operatorname{Ad}_v^{\pm1}\|\le\|\operatorname{Ad}_{vv_G^{-1}}^{\pm1}\|\le e^{2C_vM\alpha_1'''} .
\end{equation*}

\emph{Clause} (AX2-poly). We first bound $V_G$. The field $V_G=(M^{\cdot}U)^{v_G}$ is the generalised axial gauge of the real data $M^{\cdot}U$ of the $G$-valued field $U$, which satisfies $|\partial U-1|<\alpha_0'''\xi^2$, on the window $\square_0$ with its tower. The tower is of the class of \cite[(144)]{B-CMP102b}.

In \cite[Sect.~F]{B-CMP102b} the following is shown at real data. There $U'_k=U^{u'}$, where $u'$ is the generalised axial gauge map of that paper, and $V''$ is the axial datum of that paper, equal to $\bar U'^{\,j}_k$ on the level-$j$ hierarchy bonds $\Lambda'_j$ by \cite[(149)]{B-CMP102b}. If a configuration lies in the space \cite[(2)]{B-CMP102b} with $\varepsilon_0$ small, then $U'_k$ satisfies the conditions \cite[(2)]{B-CMP102b} on $\tilde\square$. From this, \cite[Proposition~2]{B-CMP98} and the generalised axial gauge conditions, the bound \cite[(145)]{B-CMP102b} follows:
\begin{equation*}
|\bar U'^{k}_k(x,x')-1|<|x-y|\,2L^2\varepsilon_0\le4d(M+R_1M_1)L^2\varepsilon_0\le8dL^2M\varepsilon_0 ,
\end{equation*}
where $y$ is the centre of the cube there (the role of $y_0$ here).
The inequality \cite[(1.65)]{B-CMP99a} then gives \cite[(146)]{B-CMP102b} for $\langle x,x'\rangle\subset\tilde\square^{(j)}$, $j=0,\dots,k$. It also gives
\begin{equation*}
|V''-1|<9dL^2M\varepsilon_0-\varepsilon_0
\end{equation*}
on the set $\mathfrak{C}_k$ of \cite[(148)]{B-CMP102b}, the hierarchy's sites; this is \cite[(151)]{B-CMP102b}. For the two top levels, \cite[(160)]{B-CMP102b} gives $|B(x,x')|<(8d^2L^2+4L^2|x-y|)\varepsilon_1$.

These arguments read of the field only its membership in the space \cite[(2)]{B-CMP102b}, which is a curvature bound, and \cite[Proposition~2]{B-CMP98}, which is proved for $\mathrm{U}(N)$-valued fields. They therefore apply verbatim to $U$ on the tree of the setting. That tree is the one used in \cite{B-CMP102b}: one top-level tree followed by the nested descent \cite[(1.15)]{B-CMP99a}.

By covariance of averages \cite{B-CMP98}, for $c\in\Lambda_p(\square_0)$,
\begin{equation*}
V_G(c)=v_G(c_-)\bar U^{p}(c)v_G(c_+)^{-1}=\bar U'^{p}(c)=V''(c)
\end{equation*}
by \cite[(149)]{B-CMP102b}, since $u'$ coincides with $v_G$ at block base sites. So the bounds \cite[(145)--(146), (151)]{B-CMP102b} are read as they stand. The half-side of the region in \cite[(145)]{B-CMP102b}, written $M$ there, is here at most $6M/L+2M_1\le8M/L$ in level-$(k+1)$ units, and $\varepsilon_0$ there becomes $\alpha_0'''$. This gives, for every $c\in\mathfrak{B}(\square_0)$,
\begin{equation*}
|V_G(c)-1|\le9dL^2\cdot\frac{8M}{L}\cdot2\alpha_0'''\le\tfrac12C_VM\alpha_0''' ,
\end{equation*}
since $C_V\ge2^8d^2L^2\ge288dL$.

Next we compare $V$ with $V_G$. The value $V(c)$ is the ordered product of the global averages of $\mathbf{U}$ around the loop
\begin{equation*}
\Gamma_{y_0,c_-}\cup c\cup\bar\Gamma_{c_+,y_0},
\end{equation*}
and $V_G(c)$ is the same product for $U$. By (AX1), on $\mathfrak{B}(\square_0)$ and on the legs, $\bar{\mathbf{U}}^{q}=e^{iB'_q}\bar U^{q}$ bondwise on every bond of the loop. So the factors are
\begin{equation*}
(e^{iB'(c')}\bar U(c'))^{\pm1}=\bar U(c')^{\pm1}+E_{c'},\qquad \|E_{c'}\|\le e^{|B'(c')|}-1 .
\end{equation*}
For unitaries $U_i$ one has $\|\prod(U_i+E_i)-\prod U_i\|\le\prod(1+\|E_i\|)-1$. By the leg sums above, applied to the two contours, and the bound on $|B'(c)|$,
\begin{align*}
|V(c)-V_G(c)|&\le e^{\sum_{\mathrm{loop}}|B'|}-1\le e^{(2C_v(M/L)+c_Q)\alpha_1'''}-1\\
&\le e^{3C_vM\alpha_1'''}-1\le6C_vM\alpha_1'''\le\tfrac12C_VM\alpha_1''' .
\end{align*}
Here $3C_vM\alpha_1'''\le\frac14$ by (c7$'$), and $C_V\ge12C_v$. Hence, on all of $\mathcal{P}_U$,
\begin{equation*}
|V(c)-1|\le\tfrac12C_VM\alpha_0'''+\tfrac12C_VM\alpha_1'''\le C_VM(\alpha_0'''+\alpha_1''') .
\end{equation*}
Every input to this bound is a published result.

\emph{Clause} (AX2-pt) is clause (b) of Remark~\ref{4.5:rem-block-curvature-complex}, which is stated in \cite[(3.26)]{B-CMP109} without proof. Nothing else is asserted about $V$ at the point. As regards (AX2-pt), the computation given above for (AX2-poly) serves only to fix the constant $C_V$.

\emph{Clause} (AX3). On $\mathcal{P}_U$, by (AX2-poly), (c7) and $C_1\bar\varepsilon_1\le\frac14$,
\begin{equation*}
|V-1|\le C_VM(\alpha_0'''+\alpha_1''')\le\tfrac14C_1\bar\varepsilon_1\le\tfrac1{16}.
\end{equation*}
So $B_V:=\frac1i\log V$ is defined by the power series, with $|B_V|\le2|V-1|\le\frac12C_1\bar\varepsilon_1$. Hence $B_V$ lies in the polydisc \cite[(172)]{B-CMP102b} over the base $V_0=1$ of the window tower $\{\square_n\}$. That tower is admissible in the sense of \cite[(1)]{B-CMP102b}, as recorded in Notation~\ref{4.5:not-geometry-conventions}. Its minimiser is $U_0=1$: it has zero action, meets the constraints, is axial, and is the unique critical orbit of \cite[Theorem~1]{B-CMP102b}. Condition \cite[(7)]{B-CMP102b} holds trivially at base $1$. At the given point, $|B_V|\le2\mathfrak{v}$ by (AX2-pt) and the same logarithm bound.

By \cite[Proposition~9]{B-CMP102b} at base $1$, $U_{k+1}(\square_0,V)=U_{k+1}(\square_0,e^{iB_V}\cdot1)$ is the axial chart, and its Landau-gauge form is, by definition, the exponential of the function of that proposition. The window has fine spacing $\xi=L^{-1}\eta$, and its top block has side $1$ in its own units ($L$ in $\eta$-units). For the window this reads
\begin{equation*}
U_{k+1}(\square_0,V)^{\mathrm{Landau}}=\exp\bigl(i\xi\mathcal{H}^{\square_0}(B_V)\bigr)
=\exp\bigl(i\eta L^{-1}\mathcal{H}^{\square_0}(B_V)\bigr)=e^{i\eta\mathcal{A}} .
\end{equation*}
The axial representative is $(e^{i\eta\mathcal{A}})^{g}$, where $g$ is the Landau-to-axial map, that is, $u_{k+1}$ of \cite[(3.27)]{B-CMP109}.

Hypothesis (A) is applied at base $U_0=1$. There the covariant gradient $\nabla_{U_0}$ is the plain gradient $\nabla$, and the H\"older member is the plain one. It is integrated along the segment $\tau B_V$, $\tau\in[0,1]$; this segment lies in the polydisc by the margin, and at $\tau=0$ the datum is the base, where $\mathcal{H}^{\square_0}$ vanishes. The result is summed over the sources $c\in\mathfrak{B}(\square_0)$ with \cite[Lemma~2.1]{B-CMP96}. Each block site $y'$ carries at most $d$ block bonds of its own level and at most $2dL$ bonds at an interface. Hence, with $k_c$ the kernel of the derivative in the datum at $c$, $\delta_0$ the decay rate of the kernel bound in (A), $y$ the block site of $x$ and $\ell_{y'}$ the relative block size of the level of $y'$,
\begin{equation*}
\sum_c|k_c(x)|\,|B_V(c)|\,\ell_{y'}^d\le3dL\,C_H\sum_{y'}e^{-\delta_0d(y,y')/4}\sup|B_V|
\le3dLC_Hc_1\sup|B_V| .
\end{equation*}
At the top layer $\square_{k+1}\supset\tilde\square^{4}$ the local unit is $1$ in window units. So in window units $|\mathcal{H}^{\square_0}|$, $|\nabla^{\xi}\mathcal{H}^{\square_0}|$ and $\|\mathcal{H}^{\square_0}\|_{1,\beta}$ are all at most $3dLC_Hc_1\sup|B_V|$. On a block of a collar layer $\Lambda_q(\square_0)$, where the members in (A) carry the factor $\ell_q^{-1}$, the bound $\ell_q\sup|\mathcal{H}^{\square_0}|\le3dLC_Hc_1\sup|B_V|$ holds likewise.

Converting to $\eta$-units ($\mathcal{A}=L^{-1}\mathcal{H}^{\square_0}$, $\nabla^{\eta}=L^{-1}\nabla^{\xi}$, and the H\"older quotients alike) only improves these bounds, by powers of $L^{-1}\le1$. Hence at the point, on $\tilde\square^{4}$,
\begin{equation*}
|\mathcal{A}|,\ |\nabla^{\eta}\mathcal{A}|,\ \|\mathcal{A}\|_{1,\beta}\le3dLC_Hc_1\sup|B_V|
\le3dLC_Hc_1\cdot2C_VM\alpha_0'''=b
\end{equation*}
by the definition of $B_b$. On $\mathcal{P}_U$ the same quantities are at most $b^{\mathrm{poly}}$, by $\sup|B_V|\le2C_VM(\alpha_0'''+\alpha_1''')$.

We turn to $g$. It is normalised by the averaging conditions \cite[(1.29)]{B-CMP99a}: the $j$-fold averages of the gauge map equal $1$ on $\Lambda_j$, $j=0,1,\dots,k$. By \cite{B-CMP98}, such a gauge map is determined uniquely by its values at the points of the coarsest lattice, and these values are computed from additional conditions. By \cite{B-CMP102b}, exactly one gauge transformation satisfies these conditions, and it is given by the formula \cite[(86)--(87)]{B-CMP98}. Since the $0$-fold average is the identity, $g=1$ on $\Lambda_0(\square_0)=P$ only, and not at higher base sites.

Fix a layer $\Lambda_q(\square_0)$, $1\le q\le k+1$, and a block $B^q(x_q)$ of it; the argument is the same at every layer. There, at the point,
\begin{equation*}
\ell_q\sup_{B^q(x_q)}|\mathcal{H}^{\square_0}|\le3dLC_Hc_1\sup|B_V|\le3dLC_Hc_1\cdot2C_VM\alpha_0'''=b
\end{equation*}
in window units. Since $\mathcal{A}$ is $L^{-1}\mathcal{H}^{\square_0}$ in $\eta$-units, $\sup_{B^q(x_q)}|\mathcal{A}|\le b/(L\ell_q)$ at the point, and $\le b^{\mathrm{poly}}/(L\ell_q)$ on $\mathcal{P}_U$.

Below its top base-site value $g(x_q)$, the map $g$ is the nested-contour transport of the averages of $e^{i\eta\mathcal{A}}$. By covariance of averages \cite{B-CMP98} and \cite[(1.15)]{B-CMP99a},
\begin{equation*}
g(x_p)=g(x_{p+1})\cdot\overline{(e^{i\eta\mathcal{A}})}^{\,p}(\Gamma_{x_{p+1},x_p}),\qquad p<q .
\end{equation*}
By the formula \cite[(86)--(87)]{B-CMP98}, the top value $g(x_q)$ is the inverse of the iterated block mean over $B^q(x_q)$ of these same nested transports; the mean is taken at every level.

The leg $\Gamma_{x_{p+1},x_p}$ has at most $dL$ bonds of the level-$p$ block lattice. Each carries an averaged bond variable $e^{iB_p(c)}$ of $e^{i\eta\mathcal{A}}$ at base $1$, with
\begin{equation*}
|B_p(c)|\le c_Q\cdot\frac{b}{L\ell_q}\cdot L\ell_p=c_Qb\,\frac{\ell_p}{\ell_q}
\end{equation*}
by \cite[Proposition~6]{B-CMP98} at $U_0=1$ with \cite[(159), (161), (163)]{B-CMP98}: amplitude times $\eta$-length of a block bond. Here the spare factor $L^{-1}$ and the layer's unit are kept. The proposition applies because the legs average up to level $q-1$, whose block side is $L\ell_{q-1}=\ell_q$ in $\eta$-units. So the rescaled amplitude is $\ell_q\cdot b/(L\ell_q)=b/L\le\alpha_1^{(98)}$, by the third member of (c7$'$).

Since $\sum_{p<q}L^{p-q}\le2/L$, every nested transport is within
\begin{equation*}
dLc_Qb\sum_{p<q}\frac{\ell_p}{\ell_q}\le2dc_Qb
\end{equation*}
of $1$. This bound does not depend on $k$ or $q$, because the legs shrink geometrically inside the layer. The inverse of the iterated mean is likewise within $2dc_Qb$ of $1$, up to Baker--Campbell--Hausdorff terms of order $(dc_Qb)^2$. Here $dc_Qb\le1/(24L)\le1/24$ by the second member of (c7$'$). The bound \cite[(162)--(163)]{B-CMP98} on a block mean of contour logarithms is of the same kind.

Recall that $O_{\mathrm{CL\text{-}U}}(x)$ is the supremum, over $0<dc_Qb\le x$, of the relative Baker--Campbell--Hausdorff remainder in this bound. It is an absolute function: it depends neither on $c_Q$ nor on $L$. Combining the two factors gives
\begin{equation*}
|\log g(x)|\le4dc_Qb\bigl(1+O_{\mathrm{CL\text{-}U}}(1/24)\bigr)\le2dLc_Qb
\qquad\text{for }L\ge2\bigl(1+O_{\mathrm{CL\text{-}U}}(1/24)\bigr),
\end{equation*}
at every layer, that is, on $\square_0$. This is hypothesis (B) at base $U_0=1$, whose quantity $\ell_p\sup|\mathcal{H}|$ equals $L\ell_q\cdot b/(L\ell_q)=b$ here; the argument just given verifies (B) at base one up to the factor $1+O_{\mathrm{CL\text{-}U}}(1/24)$, which is absorbed by the lower bound on $L$. From $|\log g|\le2dLc_Qb$ follow $|g-1|\le e^{2dLc_Qb}-1$ and $\|\operatorname{Ad}_g^{\pm1}\|\le e^{4dLc_Qb}$. On $\mathcal{P}_U$ the same holds with $b^{\mathrm{poly}}$ in place of $b$.

The bound $|u_{k+1}-1|<B_3O(1)M\alpha_0$ on $\square_0$ of \cite[(3.27)]{B-CMP109}, and the bounds \cite[(1.70)--(1.74)]{B-CMP99a} invoked in \cite{B-CMP102b}, are of the same kind; neither is used above.
\end{proof}

\begin{lemma}[Change of base by the identity theorem]\label{4.5:prf-base-change}
Work in the data-axial frame on the window (Definition~\ref{4.5:not-axial-frame}), under the
ceilings (c7), (c7$'$) and (c8) of~\eqref{4.5:eq-literal-frame-a}, and let
$(\mathbf{U},\mathbf{J})\in\mathfrak{U}$, $\mathbf{U}=e^{i\eta A'}U$, be the pair fixed there. Fix the
real configuration $U$ together with its data of condition~(i) and the map $u_0$ attached to it in
that frame. Let $\mathcal{P}_U$ be the twist polydisc of the literal frame lemma
(Lemma~\ref{4.5:lem-literal-frame}), which contains $A'$. Let $v$, $V$, $g$ and $\mathcal{A}$ be the
frame map, the axial data, the Landau-to-axial map and the frame potential of that lemma. For a
real base $V_0$ write $U^{[V_0]}_{k+1}(\square_0,\cdot)$ for the axial chart around $V_0$, defined
on the domain \cite[(172)]{B-CMP102b} over $V_0$. Let
$S=\mathbf{U}^{\mathrm{gl}}\restriction\square_0$ be as in the data-axial frame. It is the axial
chart around the real base $M^{\cdot}(U)$, evaluated at the datum $M^{\cdot}(\mathbf{U})$
(Definition~\ref{4.5:not-axial-frame}). Put
$\tilde v(x):=v(y_x)$, where $y_x$ denotes the base site of the block containing the site $x$. This
is the block-constant extension of the restriction of $v$ to the block base sites. Then the
following hold.
\begin{enumerate}
\item[(a)] On the bonds of $\square_0$,
\begin{equation}\label{4.5:eq-base-change-1}
U^{[1]}_{k+1}(\square_0,V)=S^{\tilde v},
\end{equation}
where $U^{[1]}_{k+1}(\square_0,\cdot)$ is the axial chart around the base $1$.
\item[(b)] With the site map $w(x):=g(x)^{-1}\tilde v(x)$, one has $e^{i\eta\mathcal{A}}=S^{w}$ on the
bonds of $\square_0$.
\item[(c)] At every site $x$ of $\square_0$,
\begin{align*}
\|\mathrm{Ad}_{w}^{\pm1}\|&\le\|\mathrm{Ad}_{g^{\mp1}}\|\,\|\mathrm{Ad}_{\tilde v}^{\pm1}\|,\\
\|\mathrm{Ad}_{\tilde v(x)}^{\pm1}\|&=\|\mathrm{Ad}_{v(y_x)}^{\pm1}\|.
\end{align*}
\item[(d)] The map $w$ is an explicitly given analytic function of the restriction of
$\mathbf{U}$ to $\square_0$.
\item[(e)] The definition of the field $\mathbf{A}$ on $\square_0$ in item (AX4-v)
of~\eqref{4.5:eq-literal-frame-b}, and
every bound on it, uses $w$ at the sites of $\square_0$ only. The identity expressing $\mathbf{A}$
through the chart at the transformed background holds with the extension $w^{\mathrm{ext}}$ of
Lemma~\ref{4.5:lem-literal-frame}. Finally, $\mathbf{A}$ vanishes off $\tilde\square^{4}$.
\end{enumerate}
\end{lemma}

\begin{proof}
For a site map $u$ and a bond configuration $W$ we write $W^{u}$ for the gauge-transformed
configuration. The composition rule $(W^{a})^{c}=W^{ca}$ holds. For a bond, or a contour, $\Gamma$
from $x$ to $y$ it reads $W^{u}(\Gamma)=u(x)W(\Gamma)u(y)^{-1}$. For a site map $u$, $R(u)$ denotes
its adjoint action on Lie-algebra valued fields, as in Lemma~\ref{4.5:prf-chart-covariance}, and
$\bar u$ denotes the restriction of $u$ to the
block sites. Write $\mathbf{U}_{A''}:=e^{i\eta A''}U$ for $A''\in\mathcal{P}_U$; thus
$\mathbf{U}_{A'}=\mathbf{U}$. Identify $\mathcal{P}_U$
with a polydisc in $\mathbb{C}^{D_U}$, where $D_U$ is the number of real coordinates of a
$\mathfrak{g}$-valued field $A''$ on the bonds concerned. Then the points of
$\mathcal{P}_U\cap\mathbb{R}^{D_U}$ are the $\mathfrak{g}$-valued $A''$.

\emph{The two holomorphic functions.}
Put $V_0^U:=M^{\cdot}(U)$. Since $U$ is $G$-valued with $|\partial U(p)-1|<\alpha_0'''\xi^{2}$,
\cite[Proposition~2]{B-CMP98} applies to it. The consequence of that proposition drawn
in~\cite{B-CMP102b}, recalled in (c8), states that the block configuration satisfies
\cite[(7)]{B-CMP102b} with $\varepsilon_1=O(\varepsilon_0)$. Hence $V_0^U$ satisfies
\cite[(7)]{B-CMP102b} with $\varepsilon_1\le C_7\alpha_0'''$, where $C_7$ is the constant of (c8). By
(c8),
\begin{equation*}
C_7\alpha_0'''\le C_7(\alpha_0'''+C_{14}\alpha_1''')\le\bar\varepsilon_1 .
\end{equation*}
So $V_0^U$ is real data with \cite[(7)]{B-CMP102b} at $\varepsilon_1^{\mathrm{vp}}=\bar\varepsilon_1$;
here and below the superscript $\mathrm{vp}$ marks the parameter $\varepsilon_1$ of
\cite[(7)]{B-CMP102b} and the constants of \cite[Theorem~1]{B-CMP102b}. The
minimiser $U_{k+1}(\square_0,V_0^U)$ exists by \cite[Theorem~1]{B-CMP102b}.

Let $U^{[V_0^U]}_{k+1}(\square_0,\cdot)$ denote the axial chart around the real base $V_0^U$, and
define
\begin{equation*}
S[A'']:=U^{[V_0^U]}_{k+1}\bigl(\square_0,M^{\cdot}(\mathbf{U}_{A''})\bigr),
\qquad
F_1(A''):=S[A'']^{\tilde v[A'']} .
\end{equation*}
Here the datum is $M^{\cdot}(\mathbf{U}_{A''})=e^{iB'[A'']}V_0^U$. By item (AX1)
of~\eqref{4.5:eq-literal-frame-b},
$B'[A'']$ is holomorphic on $\mathcal{P}_U$ with $\sup|B'|\le c_Q\alpha_1'''$, and by (c8)
$c_Q\alpha_1'''\le\tfrac12C_1\bar\varepsilon_1$. Hence the datum lies inside the domain
\cite[(172)]{B-CMP102b} of the chart, with a margin. Next, $\tilde v[A''](x):=v[A''](y_x)$, where
$v[A'']$ is the product, defined in the data-axial frame, of the global averages of
$\mathbf{U}_{A''}$ along the contours $\Gamma_{y_0,\cdot}$. The map $v[A'']$ is holomorphic on
$\mathcal{P}_U$, and hence
so is $\tilde v[A'']$. The gauge action $S\mapsto S^{\tilde v}$ is polynomial in
$(S,\tilde v^{\pm1})$. Therefore $F_1$ is holomorphic on $\mathcal{P}_U$.

Let $U^{[1]}_{k+1}(\square_0,\cdot)$ be the axial chart around the base $1$, and define
\begin{equation*}
F_2(A''):=U^{[1]}_{k+1}\bigl(\square_0,V[A'']\bigr),
\qquad
V[A'']:=\bigl(M^{\cdot}(\mathbf{U}_{A''})\bigr)^{v[A'']}=e^{iB_V[A'']}\cdot1 .
\end{equation*}
By items (AX2-poly) and (AX3) of~\eqref{4.5:eq-literal-frame-b}, $B_V[A'']$ is holomorphic on
$\mathcal{P}_U$ with values in the
domain \cite[(172)]{B-CMP102b} over the base $1$. Hence $F_2$ is holomorphic on $\mathcal{P}_U$. At
the given point $A'$ we have $F_1(A')=S^{\tilde v}$ and $F_2(A')=U^{[1]}_{k+1}(\square_0,V)$. So
\eqref{4.5:eq-base-change-1} is the equality $F_1(A')=F_2(A')$.

\emph{Real points.}
Let $A''\in\mathcal{P}_U$ be $\mathfrak{g}$-valued. Then $\mathbf{U}_{A''}$ is $G$-valued.
Condition~(iii) in the definition of $\mathfrak{U}$ (Proposition~\ref{4.5:prop-glued-window}) is not
available for it. Instead, its curvature is bounded by the plaquette expansion in item~(ii) of
Lemma~\ref{4.5:lem-uniform-increment}. Writing $\partial W(p)$ for the plaquette
variable of $W$ at $p$ and $\nabla^{\xi}_U$ for the covariant lattice derivative,
\begin{equation*}
\begin{split}
|\partial\mathbf{U}_{A''}(p)-\partial U(p)|
&\le\xi^{2}\bigl(2d\sup|\nabla^{\xi}_UA''|+2\sup|A''|^{2}\bigr)(1+4\xi\alpha_1''')^{2}\\
&\le C_{14}\alpha_1'''\xi^{2},
\end{split}
\end{equation*}
so that $|\partial\mathbf{U}_{A''}-1|<(\alpha_0'''+C_{14}\alpha_1''')\xi^{2}$. The block data
$M^{\cdot}(\mathbf{U}_{A''})$ are then $G$-valued by \cite[Proposition~2]{B-CMP98}, which is stated
for configurations with values in $\mathrm{U}(N)$. By the consequence of that proposition in
\cite{B-CMP102b}, which is stated at real fields, they satisfy \cite[(7)]{B-CMP102b} in all its
clauses with
\begin{equation*}
\varepsilon_1^{\mathrm{vp},(7)}:=C_7(\alpha_0'''+C_{14}\alpha_1''')\le\bar\varepsilon_1,
\end{equation*}
the inequality being (c8). Thus $M^{\cdot}(\mathbf{U}_{A''})$ is real data satisfying
\cite[(7)]{B-CMP102b} at $\varepsilon_1^{\mathrm{vp}}=\bar\varepsilon_1\le a_1^{\mathrm{vp}}$. Here
$a_1^{\mathrm{vp}}$
is the bound on $\varepsilon_1$ in the list of constants of \cite[Theorem~1]{B-CMP102b}.
Furthermore, $v[A'']\in G$. The configuration $V[A'']$ is $G$-valued and has the same plaquette
variables as $M^{\cdot}(\mathbf{U}_{A''})$. Since condition \cite[(7)]{B-CMP102b} is invariant under
gauge maps on the block sites, $V[A'']$ satisfies it as well. Finally,
$|V[A'']-1|\le\tfrac14C_1\bar\varepsilon_1$.

\emph{Step $(\alpha)$.}
The value $S[A'']$ is the chart around the real base $V_0^U$, evaluated at the $G$-valued datum
$e^{iB'}V_0^U$ near $V_0^U$. By the text preceding \cite[(181)]{B-CMP102b}, at such data the chart
coincides with the minimal configuration constructed in the previous sections there. Hence
$S[A'']=U_{k+1}(\square_0,M^{\cdot}(\mathbf{U}_{A''}))$, the
axial minimiser of \cite[Theorem~1]{B-CMP102b} for the real data $M^{\cdot}(\mathbf{U}_{A''})$.

This minimiser is well defined for two reasons. First, \cite[Theorem~1]{B-CMP102b} gives a unique
minimal orbit. It is applied with $\varepsilon_0:=B_3^{\mathrm{vp}}\bar\varepsilon_1$, where
$B_3^{\mathrm{vp}}$ is a constant of that theorem and $a_0^{\mathrm{vp}}$ is its bound on
$\varepsilon_0$; by its list of constants,
$B_3^{\mathrm{vp}}\bar\varepsilon_1\le B_3^{\mathrm{vp}}a_1^{\mathrm{vp}}\le a_0^{\mathrm{vp}}$.
Second, the axial conditions
\cite[(1.15)]{B-CMP99a}, together with the condition \cite[(1.14)]{B-CMP99a} on gauge
transformations, determine a unique element in each orbit of the subgroup of gauge transformations
satisfying \cite[(1.14)]{B-CMP99a}.

\emph{Step $(\beta)$.}
In the same way, $F_2(A'')$ is the chart around the base $1$ at the $G$-valued datum $V[A'']$ near
$1$. Hence $F_2(A'')=U_{k+1}(\square_0,V[A''])$, the axial minimiser for the real data $V[A'']$.

\emph{Step $(\gamma)$: covariance of the real minimiser.}
Let $\mathsf{V}$ be $G$-valued data satisfying \cite[(7)]{B-CMP102b}. Let $\mathsf{v}$ be a $G$-valued
map on the block sites, with block-constant extension $\tilde{\mathsf{v}}$. We prove that
\begin{equation}\label{4.5:eq-base-change-2}
U_{k+1}(\square_0,\mathsf{V}^{\mathsf{v}})=U_{k+1}(\square_0,\mathsf{V})^{\tilde{\mathsf{v}}} .
\end{equation}
This is the covariance \cite[(181)]{B-CMP102b} at $G$-valued maps.

The action functional $\mathsf{A}$ of the variational problem is invariant under every gauge map on
the fine lattice. In particular, it is invariant under the block-constant map $\tilde{\mathsf{v}}$.

The constraints transform as follows. By~\cite{B-CMP98}, for a block-constant $\tilde{\mathsf{v}}$
equal to $\mathsf{v}$ at the block base sites,
\begin{equation*}
M^{\cdot}(W^{\tilde{\mathsf{v}}})=\bigl(M^{\cdot}(W)\bigr)^{\mathsf{v}} .
\end{equation*}
Indeed, the one-step averages \cite[(42)]{B-CMP98} are conjugated by the values of the map at the
base sites that are the end-points of the block bond. Iterating preserves this, because
$\tilde{\mathsf{v}}$ is constant on every sub-block.

Consequently $W\mapsto W^{\tilde{\mathsf{v}}}$ maps the space \cite[(6)]{B-CMP102b} for the data
$\mathsf{V}$ onto the space \cite[(6)]{B-CMP102b} for the data $\mathsf{V}^{\mathsf{v}}$. Its
inverse is $W\mapsto W^{\tilde{\mathsf{v}}^{-1}}$. The same map sends the curvature space
\cite[(2)]{B-CMP102b} onto itself. Since the action is invariant, it sends minima onto minima. By the
uniqueness in \cite[Theorem~1]{B-CMP102b}, the minimal orbit for $\mathsf{V}^{\mathsf{v}}$ is the
image of the minimal orbit for $\mathsf{V}$.

The axial conditions \cite[(1.15)]{B-CMP99a} are preserved by a block-constant map. Each nested
contour $\Gamma_{x_j,x_{j-1}}$ lies inside one block, on which $\tilde{\mathsf{v}}$ is constant.
Hence $W(\Gamma)=1$ implies
\begin{equation*}
W^{\tilde{\mathsf{v}}}(\Gamma)=\tilde{\mathsf{v}}\,W(\Gamma)\,\tilde{\mathsf{v}}^{-1}=1 .
\end{equation*}
So the image of the axial representative is the axial representative of the image orbit. This is
\eqref{4.5:eq-base-change-2}.

\emph{Application of $(\gamma)$.}
Take $\mathsf{V}:=V[A'']$ and $\mathsf{v}:=v[A'']^{-1}$. Then
$\tilde{\mathsf{v}}=(\tilde v[A''])^{-1}$.
By the composition rule,
\begin{equation*}
\mathsf{V}^{\mathsf{v}}
=\Bigl(\bigl(M^{\cdot}(\mathbf{U}_{A''})\bigr)^{v[A'']}\Bigr)^{v[A'']^{-1}}
=\bigl(M^{\cdot}(\mathbf{U}_{A''})\bigr)^{v[A'']^{-1}v[A'']}
=M^{\cdot}(\mathbf{U}_{A''}) .
\end{equation*}
Hence~\eqref{4.5:eq-base-change-2} gives
\begin{equation*}
U_{k+1}\bigl(\square_0,M^{\cdot}(\mathbf{U}_{A''})\bigr)
=U_{k+1}\bigl(\square_0,V[A'']\bigr)^{(\tilde v[A''])^{-1}} .
\end{equation*}
By $(\alpha)$ and $(\beta)$, this reads $S[A'']=F_2(A'')^{(\tilde v[A''])^{-1}}$. Applying
$\tilde v[A'']$ and the composition rule gives $F_2(A'')=S[A'']^{\tilde v[A'']}=F_1(A'')$. This
holds at every $\mathfrak{g}$-valued $A''\in\mathcal{P}_U$.

\emph{Identity theorem.}
Let $f$ be any matrix entry of any bond variable of $F_1-F_2$. It is a holomorphic function on
$\mathcal{P}_U$, which is connected and open in $\mathbb{C}^{D_U}$. It vanishes on the non-empty,
relatively open set $\mathcal{P}_U\cap\mathbb{R}^{D_U}$. At a real point of this set, all derivatives
of $f$ in real directions vanish. Since $f$ is holomorphic, its complex partial derivatives of every
order equal the corresponding real-direction derivatives. Hence all Taylor coefficients of $f$ at
that point vanish, and $f$ vanishes on a neighbourhood of the point. By connectedness of
$\mathcal{P}_U$, $f$ vanishes on $\mathcal{P}_U$. The given $A'$ lies in $\mathcal{P}_U$, so
$F_1(A')=F_2(A')$ there. This proves~(a).

No radius of the small neighbourhood of $G$-valued transformations in \cite[(181)]{B-CMP102b} enters.
The covariance is used at $G$-valued maps only. There it is the statement of
\cite[(181)]{B-CMP102b} for minimisers, and $(\gamma)$ proves it.

Two further points concern the constants. The sizes of item (AX2-pt)
of~\eqref{4.5:eq-literal-frame-b}, which involve $\alpha_0'''$ only, and the number $b$, are read at
the given point $A'$, where condition~(iii) holds. The
polydisc $\mathcal{P}_U$ serves holomorphy only. Its price is the ceiling (c7), which contains
$\alpha_1'''$ as well as $\alpha_0'''$. That ceiling bounds $\alpha_0'''$ and $\alpha_1'''$ after
$M$ and orders neither against the other, so it is not a cap in the sense of
Definition~\ref{2.3:def-cap}.

\emph{Proof of~(b).}
By item (AX3) of~\eqref{4.5:eq-literal-frame-b}, $(e^{i\eta\mathcal{A}})^{g}=U^{[1]}_{k+1}(\square_0,V)$,
and by (a) this equals $S^{\tilde v}$. Applying $g^{-1}$ and the composition rule twice,
\begin{equation*}
e^{i\eta\mathcal{A}}
=\bigl((e^{i\eta\mathcal{A}})^{g}\bigr)^{g^{-1}}
=(S^{\tilde v})^{g^{-1}}
=S^{g^{-1}\tilde v}
=S^{w} .
\end{equation*}

\emph{Proof of~(c).}
From $w=g^{-1}\tilde v$ we get $\mathrm{Ad}_w=\mathrm{Ad}_{g^{-1}}\mathrm{Ad}_{\tilde v}$ and
$\mathrm{Ad}_w^{-1}=\mathrm{Ad}_{\tilde v}^{-1}\mathrm{Ad}_{g}$. Submultiplicativity of the operator
norm gives the first claim. The second claim holds because $\tilde v(x)=v(y_x)$.

\emph{Proof of~(d).}
The map $v$ is a product of averages of the restriction of $\mathbf{U}$ to $\square_0$ and of fine
bond variables on $P$. It is therefore analytic by \cite[Proposition~6]{B-CMP98}. The map $g$ is
explicit in the averages of $e^{i\eta\mathcal{A}}$. Here $\mathcal{A}=L^{-1}\mathcal{H}^{\square_0}(B_V)$
is analytic in $B_V$ by the analyticity of the axial chart around the base $1$. In turn, $B_V$ is
analytic in $V$, and $V$ is analytic in the restriction of $\mathbf{U}$ to $\square_0$. Composing, $w$
is an explicit analytic function of that restriction.

\emph{Proof of~(e).}
On $\square_0$ the field $\mathbf{A}$ is defined, in item (AX4-v)
of~\eqref{4.5:eq-literal-frame-b}, by
\begin{equation*}
\mathbf{A}:=\bigl(t\tilde\zeta_{\square}+t_{\square}\zeta_{\square}\bigr)\,
R(w)\,H_k\bigl(B';\mathbf{U}^{\mathrm{gl}},\mathbf{J}\bigr).
\end{equation*}
This definition, and every bound derived from it, uses $w$ on the sites of $\square_0$ only.

The identity in item (AX4-v) of~\eqref{4.5:eq-literal-frame-b} interprets $\mathbf{A}$ as
Ba{\l}aban's background function in which the variable $B'$, the background $\mathbf{U}^{\mathrm{gl}}$
and the current $\mathbf{J}$ are replaced by $R(\bar w^{\mathrm{ext}})B'$,
$(\mathbf{U}^{\mathrm{gl}})^{w^{\mathrm{ext}}}$ and $R(w^{\mathrm{ext}})\mathbf{J}$, that is, as the
chart at the transformed background. That identity needs a map on all sites, and $w^{\mathrm{ext}}$
supplies one. We apply the gauge covariance of the chart
(Lemma~\ref{4.5:prf-chart-covariance}) with the conjugating map $w^{\mathrm{ext}}h^{-1}$, relative to
the pair $(S^{h},\mathrm{Ad}_h\mathbf{J})\in\mathfrak{S}$ of
Proposition~\ref{4.5:prop-glued-window}. For this map
$\|\mathrm{Ad}^{\pm1}\|\le2\cdot2=\mathfrak{K}$, since
$\|\mathrm{Ad}^{\pm1}_{w^{\mathrm{ext}}}\|\le2$ at every site by item (AX4-v)
of~\eqref{4.5:eq-literal-frame-b} and $\|\mathrm{Ad}^{\pm1}_{h}\|\le2$. The covariance is an
algebraic identity wherever both of its sides are defined, and here both are. On the one hand,
since $|B'|<c_9/8$, where $c_9$ is the $B$-radius of Lemma~\ref{4.5:prf-chart-covariance},
\begin{equation*}
|R(\bar w^{\mathrm{ext}})B'|\le2|B'|<c_9/4<c_9 .
\end{equation*}
On the other hand, the chart at the representative $(S^{h},\mathrm{Ad}_h\mathbf{J})\in\mathfrak{S}$
exists by the existence property (K1) of the chart, which holds on the orbit class
$\mathfrak{S}^{\mathrm{orb}}\supset\mathfrak{S}$ by Lemma~\ref{4.5:prf-chart-covariance}.

As to the support, $\zeta_{\square}$ is supported in $\tilde\square$ and $\tilde\zeta_{\square}$
vanishes off $\tilde\square^{4}$; see~\cite{B-CMP109}. Hence the prefactor in
this definition vanishes off $\tilde\square^{4}$, and so does $\mathbf{A}$.
\end{proof}

\begin{proof}[Proof of the plain regularity bounds of Lemma~\ref{4.5:lem-literal-frame}]
\label{4.5:prf-literal-frame-bounds}
We prove the two bounds in \eqref{4.5:eq-literal-frame-b} that concern the fluctuation field
$\mathbf{A}$ and the combined field $\hat{\mathbf{A}}$. The first is the bound on the sup, the
plain gradient and the plain $\beta$-H\"older member of $\mathbf{A}$, at $s\equiv1$ and for the
$s$-localised chart. The second is the bound on $\hat{\mathbf{A}}$ on $\tilde\square^{4}$.
Throughout, $|B'|$ stands for $\sup|B'|$.

We begin with the bound in the glued window background. By
Proposition~\ref{4.5:prop-glued-window}, the glued window background $\mathbf{U}^{\mathrm{gl}}$ lies
in the orbit class $\mathfrak{S}^{\mathrm{orb}}$, with axial representative $S$ and gauge map $h$
satisfying $\|\mathrm{Ad}_h^{\pm1}\|\le 2$.
We apply the gauge covariance of the chart on the orbit class
(see~\ref{4.5:prf-chart-covariance}). The conjugating map is $h^{-1}$, and the number $2$ takes
the place of $\mathfrak{K}$. We write
$\bar h$ for the restriction of $h$ to the block sites, $R(u)$ for the action of a gauge map $u$
on functions on bonds and on data, and $S^{u}$ for the gauge transform of $S$ by $u$, both in
the convention of the covariance statement. Covariance gives
\begin{equation}\label{4.5:eq-literal-frame-bounds-1}
\begin{aligned}
N_{\beta}^{\mathbf{U}^{\mathrm{gl}}}\bigl(H_k(B';\mathbf{U}^{\mathrm{gl}},\mathbf{J})\bigr)
&=N_{\beta}\bigl(R(h^{-1})\,H_k(R(\bar h)B';S^{h},\mathrm{Ad}_h\mathbf{J})\bigr)\\
&\le 2\cdot C_9^{K}\cdot 2|B'|\cdot 4\;\le\;16\,C_9^{K}|B'| .
\end{aligned}
\end{equation}
The inequality is the size bound $C_9^K|B'|$ on the background potential, taken from the linear
bounds on the background potential and the current correction
(see~\ref{4.5:prf-chart-linear-bounds}), where the constant $C_9^K$ is fixed, carried to the
orbit class by the covariance statement
with the number $2$ in place of $\mathfrak{K}$; the numerical factors so incurred multiply to
at most $16$.

For the $s$-localised chart we use the walk expansions of the chart and its localisation in the
decoupling parameters (see~\ref{4.5:prf-chart-walk-localisation}). The same estimate holds with
$C_9^K$ replaced by $C_9^K e^{c_{\mathbb{K}}\kappa_1}$. Indeed, the size bound and the covariance
statement hold for the $s$-localised chart with constants multiplied by
$e^{c_{\mathbb{K}}\kappa_1}$. Moreover, the frame maps $h$ and $w$ do not depend on $s$.
Finally, the family $\sigma_0$ of decoupling cubes consists of cubes of side $M$ disjoint from the
interior of $\tilde\square^{4}$ \cite{B-CMP116}. The $s$-localisation therefore lives off the
frame's support $\tilde\square^{4}$.

Next we transport to the frame. Let $f$ be a function on bonds and let $\Gamma_{x,x'}$ be the path
from $x$ to $x'$ that is used in the covariant gradient and H\"older members. We write
$S(\Gamma_{x,x'})$ for the parallel transport of $S$ along $\Gamma_{x,x'}$. The gauge
transformation by $w$ satisfies the identity
\begin{equation}\label{4.5:eq-literal-frame-bounds-2}
R\bigl(S^{w}(\Gamma_{x,x'})\bigr)\,(R(w)f)(x')-(R(w)f)(x)
=R(w(x))\bigl[R\bigl(S(\Gamma_{x,x'})\bigr)f(x')-f(x)\bigr].
\end{equation}
We apply \eqref{4.5:eq-literal-frame-bounds-2} bondwise and to pairs of points within
$\tilde\square^{4}$; only the values of $w$ at the sites of $\square_0$ enter. This gives, on
$\tilde\square^{4}$,
\begin{equation*}
N_{\beta}^{e^{i\eta\mathcal{A}}}\bigl(R(w)f\bigr)\le\|\mathrm{Ad}_w\|\,
N_{\beta}^{\mathbf{U}^{\mathrm{gl}}}(f).
\end{equation*}
We take $f=H_k(B';\mathbf{U}^{\mathrm{gl}},\mathbf{J})$ and use $\|\mathrm{Ad}_w\|\le 2$, which is
part of the bound $\|\mathrm{Ad}_w^{\pm1}\|\le2$ on the frame map in
Lemma~\ref{4.5:lem-literal-frame}, together
with \eqref{4.5:eq-literal-frame-bounds-1}. This gives the bound $32\,C_9^K|B'|$ on
$\tilde\square^{4}$.

In the norm $N_\beta$, take for the hierarchy the trivial sequence of domains. Its top layer is the
whole lattice and carries the scale factor one, so no scale factor intervenes. Hence each
covariant member of $R(w)H_k$, namely the sup, the covariant gradient and the covariant
$\beta$-H\"older quotient, is at most $32\,C_9^K|B'|$.

We now bring in the cut-offs. The field $\mathbf{A}$ is the transported function $R(w)H_k$
multiplied by $t\tilde\zeta_\square+t_\square\zeta_\square$. The cut-offs satisfy
\begin{equation*}
|\zeta_\square|,\ |\tilde\zeta_\square|\le1,\qquad
|\nabla^{\eta}\zeta_\square|,\ |\nabla^{\eta}\tilde\zeta_\square|\le O(1)M^{-1}\le 1,\qquad
\|\zeta_\square\|_{\beta},\ \|\tilde\zeta_\square\|_{\beta}\le 2 .
\end{equation*}
These bounds hold because the cut-offs $\zeta_\square,\tilde\zeta_\square$ of
\cite[\S3]{B-CMP109} have derivatives up to the second order bounded by an absolute constant, and
because $M$ lies above an absolute floor.

By the Leibniz rule for multiplication by a scalar function, multiplying by $\zeta_\square$ or by
$\tilde\zeta_\square$ changes the members as follows:
\begin{itemize}
\item the sup member is multiplied by at most $1$;
\item the covariant gradient member is at most $1+1=2\le3$ times the corresponding bound;
\item the covariant H\"older member is at most $2+1=3$ times the corresponding bound.
\end{itemize}
Hence the covariant members of $\mathbf{A}$ are at most
\begin{equation*}
3\bigl(|t|+|t_\square|\bigr)\cdot 32\,C_9^K|B'| .
\end{equation*}
The two cut-offs enter symmetrically, so no bound depends on which of them carries the tilde.

We now pass from covariant to plain members. Let $e$ be the direction of differentiation. For a
bond $b'$ write $b'_+$ for the translate of $b'$ by $\eta e$, and $(x,e)$ for the bond from the
initial point $x$ of $b'$ to $x+\eta e$; let $\nabla_{\mathcal{A}}$ be the covariant gradient in
the background $e^{i\eta\mathcal{A}}$.
Then
\begin{equation*}
\nabla^{\eta}\mathbf{A}(b')=\nabla_{\mathcal{A}}\mathbf{A}(b')
-\eta^{-1}\bigl(\mathrm{Ad}\,e^{i\eta\mathcal{A}(x,e)}-1\bigr)\mathbf{A}(b'_+),
\end{equation*}
and consequently
\begin{equation*}
|\nabla^{\eta}\mathbf{A}|\le|\nabla_{\mathcal{A}}\mathbf{A}|+2b\,|\mathbf{A}|\,e^{\eta b}.
\end{equation*}
The frame potential $\mathcal{A}$ is Lipschitz with $|\mathcal{A}|,|\nabla\mathcal{A}|\le b$.
Hence, for $|x-x'|\le1$, the parallel transports of $e^{i\eta\mathcal{A}}$ differ from $1$ by at
most $2b|x-x'|$. It follows that
\begin{equation}\label{4.5:eq-literal-frame-bounds-3}
\|\mathbf{A}\|^{\mathrm{plain}}_{1,\beta}\le(1+4b)\bigl(\|\mathbf{A}\|^{\mathrm{cov}}_{1,\beta}
+|\nabla_{\mathcal{A}}\mathbf{A}|+|\mathbf{A}|\bigr).
\end{equation}
This passage uses only $|\mathcal{A}|,|\nabla\mathcal{A}|\le b$ and
$\|\mathrm{Ad}_w^{\pm1}\|\le2$. No derivative of $w$ enters, because
\begin{equation*}
w(x)^{-1}w(x+\eta e)=\mathrm{Ad}_{w(x)^{-1}}\bigl(e^{-i\eta\mathcal{A}(x,e)}\bigr)\,S(x,e).
\end{equation*}

We now fix the constant. By the ceiling $4dLc_QB_bM\alpha_0'''\le\tfrac16$ in
\eqref{4.5:eq-literal-frame-a} we have
\begin{equation*}
b=B_bM\alpha_0'''\le\frac{1}{24dLc_Q}\le\frac16 ,
\end{equation*}
so that $1+4b\le2$. Inserting the bounds on the covariant members into
\eqref{4.5:eq-literal-frame-bounds-3}, and into the bound on $|\nabla^\eta\mathbf{A}|$, we see
that each plain member of $\mathbf{A}$ is at most
\begin{equation*}
2\cdot(3+3+1)\cdot32\,C_9^K\bigl(|t|+|t_\square|\bigr)|B'|
=448\,C_9^K\bigl(|t|+|t_\square|\bigr)|B'|
\le C_N\bigl(|t|+|t_\square|\bigr)|B'|,
\end{equation*}
where $C_N=600\,C_9^K$. For the $s$-localised chart the same argument gives the bound with
$C_N^s=C_Ne^{c_{\mathbb{K}}\kappa_1}$ in place of $C_N$.

The consequence stated in the lemma, which is the bound required in \cite[(3.31)]{B-CMP109}, is
arithmetic. Let $|t|\le1$ and let $\varepsilon_1'$ be the
sup-radius of the datum, with $3C_N^s\varepsilon_1'\le\alpha_2$. Let
$|t_\square|\le r_\square(B')=\alpha_2/(3C_N^s|B'|)$, the field-dependent Cauchy radius of the
lemma, of the kind of \cite[(1.22)]{B-CMP116}. Then
\begin{equation*}
C_N^s\cdot1\cdot\varepsilon_1'\le\tfrac13\alpha_2,
\qquad
C_N^s\cdot r_\square(B')\cdot|B'|=\tfrac13\alpha_2 ,
\end{equation*}
and therefore
\begin{equation*}
|\mathbf{A}|,\ |\nabla^{\eta}\mathbf{A}|,\ \|\mathbf{A}\|_{1,\beta}
\le\tfrac13\alpha_2+\tfrac13\alpha_2<\alpha_2 .
\end{equation*}

It remains to bound the combined field
\begin{equation*}
\hat{\mathbf{A}}=\frac{1}{i\eta}\log\bigl(e^{i\eta\mathbf{A}}e^{i\eta\mathcal{A}}\bigr).
\end{equation*}
By the Baker--Campbell--Hausdorff formula,
\begin{equation}\label{4.5:eq-literal-frame-bounds-4}
\hat{\mathbf{A}}=\mathbf{A}+\mathcal{A}+\frac{i\eta}{2}[\mathbf{A},\mathcal{A}]
+\eta^{2}\mathfrak{r}(\mathbf{A},\mathcal{A}).
\end{equation}
Here the remainder $\mathfrak{r}(\mathbf{A},\mathcal{A})$, defined by this identity, satisfies
\begin{equation*}
|\mathfrak{r}(\mathbf{A},\mathcal{A})|\le C_{\mathrm{BCH}}\bigl(|\mathbf{A}|+|\mathcal{A}|\bigr)^{3}
\qquad\text{whenever } \eta\bigl(|\mathbf{A}|+|\mathcal{A}|\bigr)\le1 .
\end{equation*}
This follows from the remainder estimate \cite[(38)]{B-CMP98}, whose constant is absolute.

On $\tilde\square^{4}$ the bound just proved gives $|\mathbf{A}|<\alpha_2$, the bound on the
frame potential gives $|\mathcal{A}|\le b$, and the ceiling $\alpha_2+b\le1/(4C_{\mathrm{BCH}})$
of the lemma gives $|\mathbf{A}|+|\mathcal{A}|\le\alpha_2+b\le1$. Hence
\begin{equation*}
|\hat{\mathbf{A}}|\le\alpha_2+b+\eta\alpha_2b+C_{\mathrm{BCH}}(\alpha_2+b)^{3}<2(\alpha_2+b).
\end{equation*}

We next apply $\nabla^\eta$ to each term of \eqref{4.5:eq-literal-frame-bounds-4}. With an
absolute constant $C$ coming from the remainder $\mathfrak{r}$, this gives
\begin{align*}
|\nabla^{\eta}\hat{\mathbf{A}}|
&\le|\nabla^{\eta}\mathbf{A}|+|\nabla^{\eta}\mathcal{A}|
+\eta\bigl(|\nabla^{\eta}\mathbf{A}||\mathcal{A}|+|\mathbf{A}||\nabla^{\eta}\mathcal{A}|\bigr)\\
&\qquad+C(\alpha_2+b)^{2}\bigl(|\nabla^{\eta}\mathbf{A}|+|\nabla^{\eta}\mathcal{A}|\bigr)\\
&<(\alpha_2+b)\bigl(1+C(\alpha_2+b)\bigr)<2(\alpha_2+b).
\end{align*}
The $\beta$-H\"older member is bounded in the same way. Each term of
\eqref{4.5:eq-literal-frame-bounds-4} is a product of H\"older functions whose sups are at most
$\alpha_2+b\le1$. For the $s$-localised chart the same computation applies, since it uses only
the bounds on $\mathbf{A}$ just proved, and these hold at $s$ as well.

Finally, $b=B_bM\alpha_0'''=B_b(1+2\beta)M\alpha_0$ is the quantity $B_3O(1)M\alpha_0$ of
\cite[(3.32)]{B-CMP109}; of the two radii $\alpha_0,\alpha_1$ only $\alpha_0$ enters it.
\end{proof}

\begin{remark}[Role of small block-averaged curvature in the literal frame]
\label{4.5:rem-averaged-curvature-role}
Let $\mathbf{U}$ be near-unitary complex data as in Lemma~\ref{4.5:lem-literal-frame}. Let $V$ be the
axial data on the hierarchy $\mathfrak{B}(\square_0)$ of the window and $v$ the data-axialiser, and
let $C_V$, $B_b$ and $b := B_b M\alpha_0'''$ be as in that lemma. Let $B_3$ be Ba{\l}aban's constant
of \cite[\S3]{B-CMP109}, the one appearing in his assumption quoted in the proof below.

Suppose the small curvature of block averages at near-unitary complex data
(Remark~\ref{4.5:rem-block-curvature-complex}, printed without proof in \cite[(3.26)]{B-CMP109};
cited as printed) is not used. Then the quantity $b$ entering \cite[(3.27), (3.32)]{B-CMP109} is
replaced by $b^{\mathrm{poly}} := B_b M(\alpha_0'''+\alpha_1''')$. Ba{\l}aban's assumption on
\cite[p.~277]{B-CMP109} then implies
\begin{equation}\label{4.5:eq-averaged-curvature-role-1}
B_3 B_b M < \tfrac12 L ,
\end{equation}
that is, a cap on $M/L$ in the sense of Definition~\ref{2.3:def-cap}.

Consequently the form of the frame potential of Lemma~\ref{4.5:lem-literal-frame} whose bound $b$
involves $\alpha_0'''$ alone and is linear in $M$ is needed by the argument of this chapter, in this
frame or in any other. In Ba{\l}aban's literal frame, at complex near-unitary data $\mathbf{U}$, this
form is supplied by Remark~\ref{4.5:rem-block-curvature-complex} (printed without proof in
\cite[(3.26)]{B-CMP109}; cited as printed). The use of that remark is a choice, not a necessity:
the same bounds, involving $\alpha_0$ alone and linear in $M$, are reached without it by an axial
gauge of $\mathbf{U}$ itself on the fine lattice, followed by block averaging.
\end{remark}

\begin{proof}
We take the steps in order.

\emph{Step 1: the bound on $V$ without the remark.} The pointwise bound on $V$ involving $\alpha_0'''$
alone, in \eqref{4.5:eq-literal-frame-b}, is obtained from clause (b) of
Remark~\ref{4.5:rem-block-curvature-complex}, printed without proof in \cite[(3.26)]{B-CMP109};
cited as printed. If
that remark is not used, only the polynomial form of the bound in \eqref{4.5:eq-literal-frame-b}
remains. It gives, at each point $c \in \mathfrak{B}(\square_0)$,
\begin{equation}\label{4.5:eq-averaged-curvature-role-2}
|V(c)-1| \le C_V M(\alpha_0'''+\alpha_1''') .
\end{equation}
Accordingly, the quantity entering \cite[(3.27), (3.32)]{B-CMP109} in place of $b$ is
$b^{\mathrm{poly}} = B_b M(\alpha_0'''+\alpha_1''')$ instead of $b = B_b M\alpha_0'''$.

\emph{Step 2: the assumption on page 277.} Ba{\l}aban assumes
\begin{equation*}
\text{\lq\lq Assume now that } B_3^2 O(1) M\alpha_0 < \tfrac12\alpha_1\text{\rq\rq}
\end{equation*}
in \cite[p.~277]{B-CMP109}. In the argument that follows this assumption, the background function
$\mathbf{H}_j(\square_0, \tau Q(L^{-1}\eta\mathbf{H}_{k+1}))$, in the notation of
\cite[\S3]{B-CMP109}, is bounded by $B_3\, b$ times the factor $L^{j-1}\eta$
\cite[p.~277]{B-CMP109}. With $b^{\mathrm{poly}}$ from Step~1 in place of $b$ the assumption reads
\begin{equation}\label{4.5:eq-averaged-curvature-role-3}
B_3 B_b M(\alpha_0+\alpha_1) L^{-1} < \tfrac12 \alpha_1 .
\end{equation}

\emph{Step 3: the cap on $M/L$.} The radius $\alpha_0$ is non-negative. Dropping the term
$B_3 B_b M \alpha_0 L^{-1}$ from the left side of \eqref{4.5:eq-averaged-curvature-role-3} therefore
gives
\begin{equation*}
B_3 B_b M \alpha_1 L^{-1} < \tfrac12 \alpha_1 .
\end{equation*}
Dividing by $\alpha_1 > 0$ and multiplying by $L$ yields \eqref{4.5:eq-averaged-curvature-role-1},
a cap on $M/L$ in the sense of Definition~\ref{2.3:def-cap}. This proves the first assertion.

\emph{Step 4: why the $\alpha_0$-only form is needed.} By Steps~1--3, a frame potential whose bound
contains $\alpha_1$ with the factor $M$ imposes \eqref{4.5:eq-averaged-curvature-role-1}, which is a
cap on $M/L$. This holds whatever frame the bound is derived in. The order of choice of the
constants admits no cap (Lemma~\ref{2.4:lem-no-cycle}), so such a bound cannot be fed into the step.
Hence the form involving $\alpha_0$ alone and linear in $M$ is needed.

\emph{Step 5: how the literal frame supplies it.} In Ba{\l}aban's literal frame
\cite[(3.26)--(3.28)]{B-CMP109}, at complex near-unitary data $\mathbf{U}$, this form is the content
of Remark~\ref{4.5:rem-block-curvature-complex}, printed without proof in \cite[(3.26)]{B-CMP109};
cited as printed. There the radius $\alpha_1$ sits entirely in $v$:
Ba{\l}aban has \lq\lq$|v-1| < O(1)M\alpha_1$\rq\rq{} \cite[(3.26)]{B-CMP109}. In
Lemma~\ref{4.5:lem-literal-frame} this costs only the ceilings
\begin{equation*}
2C_v M\alpha_1''' \le \tfrac16, \qquad 4dLc_Q B_b M\alpha_0''' \le \tfrac16, \qquad
B_b (M/L)\alpha_0''' \le \alpha_1^{(98)}
\end{equation*}
of that lemma. The complex data-axial gauge then removes $\alpha_1$ from $V$. That removal rests on a
single statement: block averaging preserves small curvature at near-unitary fields. This statement
is Remark~\ref{4.5:rem-block-curvature-complex}, printed without proof in \cite[(3.26)]{B-CMP109};
cited as printed.

\emph{Step 6: the remark is not needed.} An axial gauge of the complex data $\mathbf{U}$ itself on the
fine lattice, followed by block averaging, reaches the same frame bounds, involving $\alpha_0$ alone
and linear in $M$. It starts from the curvature bound assumed of the complex data $\mathbf{U}$ itself
and uses published inputs only. Hence the appeal to
Remark~\ref{4.5:rem-block-curvature-complex} (printed without proof in \cite[(3.26)]{B-CMP109};
cited as printed) is confined to the literal frame of
Lemma~\ref{4.5:lem-literal-frame}. This proves the last assertion.
\end{proof}

\begin{definition}[Axial gauge of the fine field: idea and units at level $k+1$]
\label{4.5:not-fine-frame-units}
The data-axial frame on the window (Definition~\ref{4.5:not-axial-frame}) brings the data
$M^{\cdot}(\mathbf{U})$ into axial gauge. At a complex configuration $\mathbf{U}$ the
residual holonomies of the data in this gauge are small of order $\alpha_0$ only by way of an
averaging theorem for the block-averaged data. In what follows we proceed differently.

\begin{itemize}
\item We first bring the fine field $\mathbf{U}$ itself into axial gauge on the window. Its
plaquettes are controlled directly by condition (iii) below, and no averaging theorem is
needed for this.
\item Only then do we average, at base $1$, that is, in the units fixed below, in which the
window's top block has side $1$. There \cite[Proposition~6]{B-CMP98}, which is proved there
for complex twists, applies.
\end{itemize}

The twist of size $\alpha_1'''$ of condition (ii) goes into the frame map. This is the place
of the bound $|v^{\pm1}-1| < O(1)M\alpha_1$ on the frame map in Ba{\l}aban's construction
\cite{B-CMP109}. What remains in the residual holonomies is the curvature of $\mathbf{U}$
summed over a family of fine plaquettes spanning the contour loop. This is the place of the
bound $|V(b)-1| < O(1)M\alpha_0$, $b$ a bond of the window, in that construction.

\emph{Level-$(k+1)$ units.} Lengths are measured in the units in which the block lattice
$T^{(k+1)}_{1}$ of level $k+1$ has unit spacing, so that the window's top block has side $1$
and the fine lattice, on whose bonds $U$ and $\mathbf{U}$ live, has spacing $\xi$. The fine
spacing $\xi$ and the length $\mathsf{M}$ are
\begin{equation*}
\xi := L^{-(k+1)} = L^{-1}\eta, \qquad \mathsf{M} := M/L ,
\end{equation*}
where $\eta = L^{-k}$ is the fine spacing in the units of the data-axial frame
(Definition~\ref{4.5:not-axial-frame}).
The window $\square_0 = \tilde\square^{5}$ has side $12\mathsf{M}$. Its centre is $y_0$, the
centre of $\square$.

We let $\square_0^{+}$ be the set of sites that are end-points of bonds of $\square_0$, in
accordance with the convention that the sites of $\square_0$ include the end-points of its
bonds. The set $\square_0^{+}$ lies inside the box of centre $y_0$ and half-side
$6\mathsf{M}+\xi$.

For $n \le k+1$ a block of level $n$ has side
\begin{equation*}
\ell_n = L^{n-k-1} \le 1 .
\end{equation*}

\emph{The conditions \textup{(i)}--\textup{(iii)} of $\mathfrak{U}$ on $U$ and $\mathbf{U}$,
rewritten in these units.} Recall from Proposition~\ref{4.5:prop-glued-window} that
$\alpha_0''' = (1+2\beta)\alpha_0$ and $\alpha_1''' = (1+2\beta)\alpha_1$, with $\beta$ the
schedule constant of Definition~\ref{4.4:def-post-step-space}. For $(\mathbf{U},\mathbf{J})$
in the space $\mathfrak{U}$ of that proposition these conditions read as follows; they hold on
the whole lattice \cite{B-CMP109}. Here $\partial U$ denotes the family of plaquette variables
$U(\partial p)$, and $\nabla^{\xi}$, $\nabla^{\xi}_{U}$ denote the lattice gradient and the
lattice gradient covariant with respect to $U$, both at spacing $\xi$.
\begin{enumerate}
\item[(i)] The configuration $U$ is $G$-valued on all bonds of the fine lattice (the whole
torus), and it satisfies
\begin{equation*}
|\partial U - 1| < \alpha_0'''\xi^{2} .
\end{equation*}
There are a cube $Q_0 \supseteq \square_0^{+}$ of the kind of \cite[(1.12)]{B-CMP109}, a gauge
transformation $u_0$ and a potential $A_0$ on the bonds of $Q_0$ such that
$U^{u_0} = e^{i\xi A_0}$ on $Q_0$. The side of $Q_0$ is $c_{12}LM \ge 12\mathsf{M}+1$, and
\begin{equation*}
|A_0|,\ |\nabla^{\xi}A_0| < c_{12}LMB\alpha_0''' .
\end{equation*}
Here $B$ is the constant of the condition \cite[(1.12)]{B-CMP109}; it is not to be confused
with $B_\sharp$ or $\bar B$.
\item[(ii)] The configuration $\mathbf{U}$ is a twist of $U$: we have
$\mathbf{U} = e^{i\xi A'}U$, where the twist $A'$ is written in level-$(k+1)$ units and
satisfies
\begin{equation*}
|A'|,\ |\nabla^{\xi}_{U}A'| < \alpha_1''' .
\end{equation*}
\item[(iii)] The configuration $\mathbf{U}$ satisfies
\begin{equation*}
|\partial\mathbf{U} - 1| < \alpha_0'''\xi^{2} .
\end{equation*}
\end{enumerate}

The twist polydisc $\mathcal{P}_U$ is the set introduced with the data-axial frame
(Definition~\ref{4.5:not-axial-frame}); it is the same set here. The frame potential
constructed in these units is, at the end, converted to the $\eta$-units of the data-axial
frame.
\end{definition}

\begin{lemma}[The fine axial frame and its bounds from fine curvature alone]
\label{4.5:lem-fine-axial-frame}
Let $(\mathbf{U},\mathbf{J})$ lie in the space $\mathfrak{U}$ of
Proposition~\ref{4.5:prop-glued-window}. Throughout we use the level-$(k+1)$ units, the window
$\square_0$ with centre $y_0$, the set $\square_0^{+}$ of end-points of bonds of $\square_0$, the
block sides $\ell_n=L^{n-k-1}\le 1$ and the hypotheses (i), (ii), (iii) of $\mathfrak{U}$ written in
these units, all as in the conventions of~\ref{4.5:not-fine-frame-units}. In particular
$\xi=L^{-(k+1)}=L^{-1}\eta$ is the fine spacing and $\mathsf{M}=M/L$.

Let $u_0$ be the gauge of hypothesis (i). Let $\mathcal{P}_U$, the layers $\Lambda_n(\square_0)$ and
the set $\mathfrak{B}(\square_0)$ of block bonds of the window hierarchy be as in
Lemma~\ref{4.5:lem-literal-frame}, and likewise the constants $c_Q$, $C_1$, $\bar\varepsilon_1$,
$C_{14}$. Let $\alpha_1^{(98)}$ denote the radius of \cite[Proposition~6]{B-CMP98}.

For a datum $V$ on $\mathfrak{B}(\square_0)$ and a base of the window tower $\{\square_n\}$, write
$U_{k+1}(\square_0,V)$ for the chart of \cite[Proposition~9]{B-CMP102b}. Write
$\mathcal{H}^{\square_0}$ for the function of that proposition for the window tower, in the window's
own units $\xi$. The constants $C_H$ and $c_1$ are those of the first of the two statements
assumed below.

Assume the following two statements:
\begin{itemize}
\item the corollary bounding the background function $\mathcal{H}^{\square_0}$ of the window, at
the base $1$;
\item the lemma bounding the Landau-to-axial gauge map of the chart; only its statement is used,
for the map $g'$ below, in the form in which it enters (AX3) of~\eqref{4.5:eq-literal-frame-b}.
\end{itemize}

Put
\begin{equation*}
\bar{\mathfrak{c}}:=2\bigl(\alpha_1'''+c_{12}LMB\alpha_0'''\bigr),\qquad
\delta_f:=e^{7d\mathsf{M}\bar{\mathfrak{c}}}-1 .
\end{equation*}
Let $\rho_{\mathrm{rim}}=\rho_{\mathrm{rim}}(d,G,L)$ be the constant such that the identity
\cite[(181)]{B-CMP102b} remains valid for $G^{c}$-valued maps of the block sites whose logarithms
are bounded by $\rho_{\mathrm{rim}}/8$. It is fixed in the order of choice together with the
constants of the theorem that furnishes the axial representative $S$ and the gauge map $h$ of the
glued window background.

Assume the ceilings
\begin{equation}\label{4.5:eq-fine-axial-frame-a}
\begin{gathered}
\text{(cf1)}\quad 7d\mathsf{M}\,\bar{\mathfrak{c}}\le\tfrac{1}{10},\qquad
\delta_f\le\min\bigl\{\tfrac19,\tfrac{\rho_{\mathrm{rim}}}{8}\bigr\},\\
\text{(cf2)}\quad 56d\mathsf{M}\bigl(\alpha_0'''+C_{14}\alpha_1'''\bigr)
\le\min\bigl\{\alpha_1^{(98)},\tfrac14\bigr\},\\
\mathfrak{v}_f^{\mathrm{poly}}:=448d^2c_QM\bigl(\alpha_0'''+C_{14}\alpha_1'''\bigr)
\le\tfrac12 C_1\bar\varepsilon_1,\\
\text{(cf3)}\quad 4dLc_Q b'\le\tfrac16,\qquad \frac{b'}{L}\le\alpha_1^{(98)},\\
b':=B_b^{f}M\alpha_0''',\qquad B_b^{f}:=3dLc_1C_H\cdot 448d^2c_Q ,
\end{gathered}
\end{equation}
together with the ceilings (c8) of~\eqref{4.5:eq-literal-frame-a}. Here $c_{12}$ and $B$ are the
constants of the cube and of the bound in hypothesis (i).

In (cf3), the condition $b'/L\le\alpha_1^{(98)}$ is the radius condition of
\cite[Proposition~6]{B-CMP98} for the block averages that enter $g'$ at every layer of
$\square_0$. The second member
$\rho_{\mathrm{rim}}/8$ of the minimum in (cf1) is idle in (AXF1)--(AXF6): no radius of
\cite[(181)]{B-CMP102b} at complexified maps is used, and (AXF4-i) is proved by the identity
theorem.

Each of these conditions is an upper bound on $\alpha_0$ or on $\alpha_1$ once $L$, $M$ and $B$ are
fixed. None of them orders $\alpha_0$ against $\alpha_1$, and none is a cap in the sense of
Definition~\ref{2.3:def-cap}. No further hypothesis is made. In contrast with
Lemma~\ref{4.5:lem-literal-frame}, no bound of \cite{B-CMP109} on averages of a complexified
configuration is assumed. Then the following hold.
\begin{enumerate}
\item[(AXF1)] \emph{Fine axial gauge of the field on $\square_0$.} For $x\in\square_0^{+}$ let
$\Gamma^{0}_{y_0,x}$ be the lexicographic path of the fine lattice from $y_0$ to $x$: all steps in
direction $1$ first, then all steps in direction $2$, and so on up to direction $d$, each stage
monotone.

The path stays inside the box spanned by $\square_0^{+}$, and it has
$N_x:=|x-y_0|_1/\xi\le 7d\mathsf{M}/\xi$ bonds. The set $\square_0^{+}$ is not itself a box: it is
$\square_0$ together with one fine layer of sites. Late stages of the path may run on bonds with
both end-points in $\square_0^{+}\setminus\square_0$. Below, $\mathbf{U}\restriction\square_0^{+}$
denotes $\mathbf{U}$ on every bond with both end-points in $\square_0^{+}$.

For a bond $b$ let $b_-$ and $b_+$ be its initial and final site, and put
\begin{equation*}
\begin{gathered}
\mathsf{P}_f(x):=\mathbf{U}\bigl(\Gamma^{0}_{y_0,x}\bigr),\qquad \mathsf{P}_f(y_0)=1,\\
W:=\mathbf{U}^{\mathsf{P}_f},\qquad
W(b)=\mathsf{P}_f(b_-)\,\mathbf{U}(b)\,\mathsf{P}_f(b_+)^{-1}.
\end{gathered}
\end{equation*}
Here $\mathbf{U}(\Gamma)$ is the ordered product of the bond variables of $\mathbf{U}$ along
$\Gamma$, with inverses on bonds traversed backwards. Thus $\mathsf{P}_f$ is a function of
$\mathbf{U}\restriction\square_0^{+}$ only, and the gauge $u_0$ does not enter its definition.
Moreover:
\begin{enumerate}
\item[(a)] Put $\mathsf{P}_{f,c}(x):=u_0(y_0)\,\mathsf{P}_f(x)\,u_0(x)^{-1}$, so that
$\mathsf{P}_f(x)=u_0(y_0)^{-1}\mathsf{P}_{f,c}(x)\,u_0(x)$. Then
\begin{equation*}
\bigl\|\mathsf{P}_{f,c}(x)^{\pm1}-1\bigr\|\le\delta_f ,
\end{equation*}
and $\delta_f\le\tfrac19$ by (cf1); hence
\begin{equation*}
\bigl\|\mathrm{Ad}_{\mathsf{P}_f(x)}^{\pm1}\bigr\|\le(1+\delta_f)^2\le\bigl(\tfrac{10}{9}\bigr)^2 .
\end{equation*}
\item[(b)] $W(b)=1$ for every bond $b$ of the tree
$\mathcal{T}:=\bigcup_{x\in\square_0^{+}}\Gamma^{0}_{y_0,x}$.
\item[(c)] For every bond $b=(x,x+\xi e_\mu)$ of $\square_0$ not in $\mathcal{T}$,
\begin{equation*}
\|W(b)-1\|\le 4N_x\alpha_0'''\xi^2\le 28d\mathsf{M}\alpha_0'''\xi .
\end{equation*}
Hence $W=e^{i\xi\mathcal{A}_f}$ on the bonds of $\square_0$, with
\begin{equation*}
\sup|\mathcal{A}_f|\le b_f:=56d\mathsf{M}\alpha_0''' .
\end{equation*}
This bound involves $\alpha_0'''$ only and one power of $\mathsf{M}$. No bound on the lattice
gradient of $\mathcal{A}_f$ is asserted, and none is used.
\item[(d)] The size $\alpha_1'''$ of hypothesis (ii) enters only through the non-unitary part
$\mathsf{P}_{f,c}$ of the frame map: $\delta_f$ is comparable to
$\mathsf{M}(\alpha_1'''+c_{12}LMB\alpha_0''')$. This has the form $|v-1|<O(1)M\alpha_1$ of the
bound on Ba{\l}aban's frame map in \cite[\S3]{B-CMP109}, up to the term of order $M^2B\alpha_0'''$
coming from the gauge $u_0$ of hypothesis (i); it is not that bound exactly.

The residue $\mathcal{A}_f$ is controlled by the curvature of $\mathbf{U}$, hypothesis (iii), summed
over at most $7d\mathsf{M}/\xi$ fine plaquettes. This gives a bound of the form
$|V(b)-1|<O(1)M\alpha_0$ of \cite[\S3]{B-CMP109} without any averaging theorem at a complex field.
\end{enumerate}

\item[(AXF2)] \emph{Window data at base $1$.} Let $\bar{\mathsf{P}}_f$ be the restriction of
$\mathsf{P}_f$ to the base sites of the blocks. By the covariance of block averages
\cite{B-CMP98},
\begin{equation*}
M^{\cdot}(W)=\bigl(M^{\cdot}(\mathbf{U})\bigr)^{\bar{\mathsf{P}}_f}\quad\text{on }
\mathfrak{B}(\square_0).
\end{equation*}
On $\Lambda_0(\square_0)$ the data are the fine field $W$ itself. By
\cite[Proposition~6]{B-CMP98} at $U_0:=1$, for $c\in\Lambda_n(\square_0)$ with $n\ge1$ one has
$M^{n}(W)(c)=e^{iB_f(c)}$, with $B_f(c)$ analytic and
\begin{equation*}
|B_f(c)|\le c_Qb_f\ell_n\le c_Qb_f .
\end{equation*}
On $\Lambda_0(\square_0)$, $|B_f|=\xi|\mathcal{A}_f|\le b_f\xi$.

The data of the hierarchy on block bonds that join or cross blocks (\cite[(1.31)]{B-CMP99a}, and
the products $V(x,y)\bar V(y,z)\cdots$ entering \cite[(7)]{B-CMP102b}) are explicit analytic
products of at most $2dL$ averaged or fine variables of $W$. Each factor is within $c_Qb_f$ of $1$,
so each such product is within $4dLc_Qb_f$ of $1$.

Hence the window data $e^{iB_f}:=M^{\cdot}(W)$ obey, using $L\mathsf{M}=M$,
\begin{equation*}
\sup_{\mathfrak{B}(\square_0)}|B_f|\le\mathfrak{v}_f:=8dLc_Qb_f=448d^2c_QM\alpha_0''' ,
\end{equation*}
which involves $\alpha_0'''$ only and the first power of $M$. On all of $\mathcal{P}_U$,
$\sup|B_f|\le\mathfrak{v}_f^{\mathrm{poly}}\le\tfrac12 C_1\bar\varepsilon_1$ by (cf2). Thus $B_f$
lies, with margin, in the polydisc \cite[(172)]{B-CMP102b} over the base $V_0=1$.

\item[(AXF3)] \emph{The analogue of \cite[(3.27)]{B-CMP109} at base $1$.} This is (AX3)
of~\eqref{4.5:eq-literal-frame-b} with $(V,B_V,\mathfrak{v})$ replaced by
$(e^{iB_f},B_f,\mathfrak{v}_f)$.

By \cite[Proposition~9]{B-CMP102b} at base $1$,
\begin{equation*}
U_{k+1}\bigl(\square_0,e^{iB_f}\bigr)=\bigl(e^{i\eta\mathcal{A}'}\bigr)^{g'},\qquad
\mathcal{A}':=L^{-1}\mathcal{H}^{\square_0}(B_f),
\end{equation*}
in the units $\eta$ of Lemma~\ref{4.5:lem-literal-frame}, where $g'$ is the Landau-to-axial gauge
map of the chart.

Apply the corollary bounding $\mathcal{H}^{\square_0}$ at base $1$ along the segment from $0$ to
$B_f$, together with \cite[Lemma~2.1]{B-CMP96}. This gives, on $\tilde\square^{4}$ and indeed on all
of $\square_{k+1}$,
\begin{equation*}
|\mathcal{A}'|,\ |\nabla^{\eta}\mathcal{A}'|,\ \|\mathcal{A}'\|_{1,\beta}\le
b'=3dLC_Hc_1\mathfrak{v}_f=B_b^{f}M\alpha_0''' .
\end{equation*}
Likewise, on a block of a collar layer $\Lambda_q(\square_0)$, in window units (the factor
$\ell_q^{-1}$ of that corollary),
\begin{equation*}
\ell_q\sup\bigl|\mathcal{H}^{\square_0}(B_f)\bigr|\le 3dLC_Hc_1\mathfrak{v}_f=b' .
\end{equation*}
By the lemma bounding the Landau-to-axial gauge map, as in (AX3), $|\log g'|\le 2dLc_Qb'$ on
$\square_0$ at every layer, and $\|\mathrm{Ad}_{g'}^{\pm1}\|\le e^{4dLc_Qb'}$. On $\mathcal{P}_U$ the
same bounds hold with $b'$ replaced by
$b'^{\,\mathrm{poly}}:=3dLC_Hc_1\mathfrak{v}_f^{\mathrm{poly}}$.

\item[(AXF4)] \emph{The frame; agreement of the charts.} For a site $x$ let $y_x$ be the base site
of the block containing $x$, and put
$\tilde{\bar{\mathsf{P}}}_f(x):=\bar{\mathsf{P}}_f(y_x)=\mathsf{P}_f(y_x)$.

This map is constant on blocks. It coincides with $\mathsf{P}_f$ at the sites of $\square_0^{+}$
that lie in the finest layer $\Lambda_0(\square_0)$, where the blocks are single sites, and at the
boundary sites of the window hierarchy.
\begin{enumerate}
\item[(AXF4-i)] Let $S$ be the axial representative of the glued window background, as in (AX4-i)
of~\eqref{4.5:eq-literal-frame-b}. Then
\begin{equation*}
U_{k+1}\bigl(\square_0,M^{\cdot}(W)\bigr)=S^{\tilde{\bar{\mathsf{P}}}_f}.
\end{equation*}
This is an equality between two charts of \cite[Proposition~9]{B-CMP102b}: on the left the chart
at base $1$ at the datum $e^{iB_f}$; on the right the chart at base $M^{\cdot}(U)$ at the datum
$M^{\cdot}(\mathbf{U})$, transformed by the block-constant map $\tilde{\bar{\mathsf{P}}}_f$. It is
proved by the identity theorem on $\mathcal{P}_U$.
\item[(AXF4-ii)] Let $w':=g'^{-1}\circ\tilde{\bar{\mathsf{P}}}_f$ be the site map
$x\mapsto g'(x)^{-1}\tilde{\bar{\mathsf{P}}}_f(x)$ (first $\tilde{\bar{\mathsf{P}}}_f$, then
$g'^{-1}$; we use the convention $(U^{a})^{c}=U^{ca}$). Then
\begin{equation*}
S^{w'}=\bigl(S^{\tilde{\bar{\mathsf{P}}}_f}\bigr)^{g'^{-1}}=e^{i\eta\mathcal{A}'}
\quad\text{on the bonds of }\square_0 .
\end{equation*}
Moreover, by (cf1) and (cf3),
\begin{equation*}
\bigl\|\mathrm{Ad}_{w'}^{\pm1}\bigr\|\le\bigl(\tfrac{10}{9}\bigr)^2e^{4dLc_Qb'}
\le\bigl(\tfrac{10}{9}\bigr)^2e^{1/6}<2 .
\end{equation*}
\item[(AXF4-iii)] $w'$ is an explicit analytic function of $\mathbf{U}\restriction\square_0^{+}$.
It is built from products of bond variables of $\mathbf{U}$ (through $\mathsf{P}_f$), from block
averages (\cite[Proposition~6]{B-CMP98}) and from the function $\mathcal{H}^{\square_0}$
(\cite[Proposition~9]{B-CMP102b}).

It is therefore of the kind required of the frame maps in \cite[\S3]{B-CMP109}: an explicitly given
analytic function of $\mathbf{U}$ depending only on the restriction of $\mathbf{U}$ to the window.
The gauge $u_0$ enters the bounds only, not the definition of $w'$.
\item[(AXF4-iv)] Let $w'^{\,\mathrm{ext}}$ be $w'$ on $\square_0^{+}$ and $1$ elsewhere. Then
$\|\mathrm{Ad}^{\pm1}_{w'^{\,\mathrm{ext}}}\|\le2$ at every site. With $h$ the gauge map of (AX4-v)
of~\eqref{4.5:eq-literal-frame-b}, the map $w'^{\,\mathrm{ext}}h^{-1}$ satisfies
$\|\mathrm{Ad}^{\pm1}\|\le\mathfrak{K}$.

By the gauge covariance of the chart,~\ref{4.5:prf-chart-covariance}, applied as in (AX4-v) and
for the data $B'$ admitted there, with $\bar w'^{\,\mathrm{ext}}$ the restriction of
$w'^{\,\mathrm{ext}}$ to block sites, one has on $\square_0$
\begin{equation*}
\begin{split}
\mathbf{A}'&:=\bigl(t\tilde\zeta_{\square}+t_{\square}\zeta_{\square}\bigr)\,
R(w')\,H_k\bigl(B';\mathbf{U}^{\mathrm{gl}},\mathbf{J}\bigr)\\
&=\bigl(t\tilde\zeta_{\square}+t_{\square}\zeta_{\square}\bigr)\,
H_k\bigl(R(\bar w'^{\,\mathrm{ext}})B';(\mathbf{U}^{\mathrm{gl}})^{w'^{\,\mathrm{ext}}},
R(w'^{\,\mathrm{ext}})\mathbf{J}\bigr),
\end{split}
\end{equation*}
and $\operatorname{supp}\mathbf{A}'\subset\tilde\square^{4}$.
\end{enumerate}

\item[(AXF5)] \emph{The bound \cite[(3.31)]{B-CMP109}, for the chart and for its $s$-localised
version.} (AX5) of~\eqref{4.5:eq-literal-frame-b} holds verbatim with $(w,\mathcal{A},b)$ replaced
by $(w',\mathcal{A}',b')$. That is, on $\square_0$,
\begin{equation*}
|\mathbf{A}'|,\ |\nabla^{\eta}\mathbf{A}'|,\ \|\mathbf{A}'\|_{1,\beta}\le
C_N\bigl(|t|+|t_{\square}|\bigr)\sup|B'| ,
\end{equation*}
with $C_N$ replaced by $C_N^{s}$ for the $s$-localised chart, the constants being those of (AX5).

Hence, for $|t|\le1$ and
\begin{equation*}
|t_{\square}|\le r_{\square}(B'):=\frac{\alpha_2}{3C_N^{s}\sup|B'|},
\end{equation*}
under the ceiling $3C_N^{s}\varepsilon_1'\le\alpha_2$ with $\varepsilon_1':=\sup|B'|$, we obtain
\begin{equation*}
|\mathbf{A}'|,\ |\nabla^{\eta}\mathbf{A}'|,\ \|\mathbf{A}'\|_{1,\beta}
\le\tfrac13\alpha_2+\tfrac13\alpha_2<\alpha_2\quad\text{on }\square_0 .
\end{equation*}

\item[(AXF6)] \emph{The analogue of \cite[(3.32)]{B-CMP109}.} (AX6)
of~\eqref{4.5:eq-literal-frame-b} holds verbatim. Under $\alpha_2+b'\le 1/(4C_{\mathrm{BCH}})$, the
field
\begin{equation*}
\hat{\mathbf{A}}':=\frac{1}{i\eta}\log\bigl(e^{i\eta\mathbf{A}'}e^{i\eta\mathcal{A}'}\bigr)
\end{equation*}
obeys, on $\tilde\square^{4}$,
\begin{equation*}
|\hat{\mathbf{A}}'|,\ |\nabla^{\eta}\hat{\mathbf{A}}'|,\ \|\hat{\mathbf{A}}'\|_{1,\beta}
<2(\alpha_2+b')=2\bigl(\alpha_2+B_b^{f}M\alpha_0'''\bigr),
\end{equation*}
and likewise for the $s$-localised chart.
\end{enumerate}
In summary, (AXF1)--(AXF6) give
\begin{equation}\label{4.5:eq-fine-axial-frame-b}
\begin{gathered}
\sup|\mathcal{A}_f|\le b_f=56d\mathsf{M}\alpha_0''',\qquad
\sup_{\mathfrak{B}(\square_0)}|B_f|\le\mathfrak{v}_f=448d^2c_QM\alpha_0''',\\
|\mathcal{A}'|,\ |\nabla^{\eta}\mathcal{A}'|,\ \|\mathcal{A}'\|_{1,\beta}\le
b'=B_b^{f}M\alpha_0''',\qquad \bigl\|\mathrm{Ad}_{w'}^{\pm1}\bigr\|<2,\\
\text{for }|t|\le1,\ |t_{\square}|\le r_{\square}(B')
\text{ and under }3C_N^{s}\varepsilon_1'\le\alpha_2:\\
|\mathbf{A}'|,\ |\nabla^{\eta}\mathbf{A}'|,\ \|\mathbf{A}'\|_{1,\beta}<\alpha_2
\quad\text{on }\square_0,\\
\text{under }\alpha_2+b'\le 1/(4C_{\mathrm{BCH}}):\\
|\hat{\mathbf{A}}'|,\ |\nabla^{\eta}\hat{\mathbf{A}}'|,\ \|\hat{\mathbf{A}}'\|_{1,\beta}
<2(\alpha_2+b')\quad\text{on }\tilde\square^{4}.
\end{gathered}
\end{equation}
\end{lemma}

The proof of Lemma~\ref{4.5:lem-fine-axial-frame} uses the following inputs and no others.
\begin{itemize}
\item Hypothesis (i) of $\mathfrak{U}$, through which the gauge $u_0$ enters bounds only.
\item Hypothesis (ii), through its supremum bound only.
\item Hypothesis (iii).
\item The telescoping identity for ordered products established in the proof.
\item From \cite{B-CMP98}:
  \begin{itemize}
  \item the locality of the averaging and the similarity argument;
  \item the covariance of averages;
  \item the relations \cite[(56)--(57)]{B-CMP98};
  \item \cite[Proposition~2]{B-CMP98} at $G$-valued fields only, and only inside the real step of
  (AXF4-i);
  \item \cite[Proposition~6]{B-CMP98} at $U_0=1$ and at $U_0=U$, with the level sizes
  \cite[(159), (161), (163), (164)]{B-CMP98};
  \item for the normalisation and the top value of $g'$, the relations \cite[(76)--(87)]{B-CMP98}
  (the averaging conditions \cite[(77), (81)]{B-CMP98} and the iterated block mean
  \cite[(86)--(87)]{B-CMP98}), with the closed formula \cite[(106)]{B-CMP98} as cited in
  \cite{B-CMP102b}.
  \end{itemize}
\item The relation \cite[(1.29)]{B-CMP99a}.
\item From \cite{B-CMP102b}:
  \begin{itemize}
  \item the passage of \cite{B-CMP102b} on the normalisation of the gauge map;
  \item \cite[Theorem~1]{B-CMP102b};
  \item \cite[Proposition~9]{B-CMP102b} at the real bases $1$ and $M^{\cdot}(U)$;
  \item the polydisc \cite[(172)]{B-CMP102b};
  \item the condition \cite[(7)]{B-CMP102b}, together with the sentence following it;
  \item the identity \cite[(181)]{B-CMP102b} at $G$-valued maps only, as established in
  (AX4-i)$(\gamma)$ of~\eqref{4.5:eq-literal-frame-b}.
  \end{itemize}
\item The corollary bounding $\mathcal{H}^{\square_0}$.
\item The lemma bounding the Landau-to-axial gauge map (only its statement), for $g'$.
\item \cite[Lemma~2.1]{B-CMP96}.
\end{itemize}

The following are not used:
\begin{itemize}
\item \cite[Propositions~1 and~2]{B-CMP98} at any complex field;
\item \cite[(145)--(160)]{B-CMP102b} at complex data;
\item the bounds of \cite{B-CMP109} on the block plaquettes and on the axial data of the averaged
configuration $M^{\cdot}(\mathbf{U})$ at complex $\mathbf{U}$. These are the bounds assumed in
(AX1-$\mathbf{U}$) and (AX2-pt) of~\eqref{4.5:eq-literal-frame-b}.
\end{itemize}

The proof of Lemma~\ref{4.5:lem-fine-axial-frame} is given below.

\begin{proof}[Proof of Lemma~\ref{4.5:lem-fine-axial-frame}]
\label{4.5:prf-lattice-stokes}
Throughout, $\mathbf{U}=e^{i\xi A'}U$ is the field of Lemma~\ref{4.5:lem-fine-axial-frame}, and $U$, $A'$,
$u_0$ are as in hypotheses (i)--(iii), under which that lemma is stated. For a
site map $v$ we write $V^{v}(b)=v(b_-)V(b)v(b_+)^{-1}$, so that $(V^{u})^{v}=V^{vu}$, and
$\operatorname{Ad}_uX=uXu^{-1}$. The letter $b$ denotes a bond until the last paragraph of the
proof.

\smallskip\noindent\textit{The non-unitary part of the frame.}
On the region $Q_0\supseteq\square_0^{+}$ of hypothesis (i) of
Lemma~\ref{4.5:lem-fine-axial-frame}, on which that hypothesis represents the base field, with the
potential $A_0$ given there, as
$U(b)=u_0(b_-)^{-1}e^{i\xi A_0(b)}u_0(b_+)$, we have, for every bond $b$ there,
\begin{align*}
\mathbf{U}(b)&=e^{i\xi A'(b)}U(b)=e^{i\xi A'(b)}u_0(b_-)^{-1}e^{i\xi A_0(b)}u_0(b_+)\\
&=u_0(b_-)^{-1}e^{i\xi\operatorname{Ad}_{u_0(b_-)}A'(b)}e^{i\xi A_0(b)}u_0(b_+)
=u_0(b_-)^{-1}e^{i\xi\mathfrak{c}(b)}u_0(b_+),
\end{align*}
where the third equality is $e^{i\xi A'}u_0^{-1}=u_0^{-1}e^{i\xi\operatorname{Ad}_{u_0}A'}$, and
where $\mathfrak{c}(b)$ is defined as a logarithm:
\begin{equation*}
\mathfrak{c}(b):=(i\xi)^{-1}\log\bigl(e^{i\xi\operatorname{Ad}_{u_0(b_-)}A'(b)}
e^{i\xi A_0(b)}\bigr).
\end{equation*}
Since $u_0(b_-)$ is unitary, $\operatorname{Ad}_{u_0(b_-)}$ is isometric; with
$\alpha_1'''$ and $c_{12}LMB\alpha_0'''$ the bounds on $|A'|$ and on $|A_0|$ provided by
hypotheses (ii) and (i), the bound on the Baker--Campbell--Hausdorff remainder
\cite[(38)]{B-CMP98} gives
\begin{equation*}
|\mathfrak{c}|\le\alpha_1'''+c_{12}LMB\alpha_0'''
+C_{\mathrm{BCH}}\,\xi\,(\alpha_1'''+c_{12}LMB\alpha_0''')^2\le\bar{\mathfrak{c}},
\end{equation*}
the factor $C_{\mathrm{BCH}}\xi(\alpha_1'''+c_{12}LMB\alpha_0''')$ being at most one, and
$\bar{\mathfrak{c}}=2(\alpha_1'''+c_{12}LMB\alpha_0''')$.

Multiply these representations along the path $\Gamma^{0}_{y_0,x}$; on a bond traversed
backwards, $\mathbf{U}(b)^{-1}=u_0(b_+)^{-1}e^{-i\xi\mathfrak{c}(b)}u_0(b_-)$. At every
intermediate site of the path the factor $u_0(\cdot)$ of one bond meets its inverse from the
next, and these inner factors cancel:
\begin{equation*}
\mathsf{P}_f(x)=u_0(y_0)^{-1}\Bigl[\prod_{\Gamma^{0}_{y_0,x}}e^{\pm i\xi\mathfrak{c}}\Bigr]u_0(x)
=:u_0(y_0)^{-1}\,\mathsf{P}_{f,c}(x)\,u_0(x),
\end{equation*}
the product in path order. For matrices $E_j$, expanding the product gives
$\|\prod_j(1+E_j)-1\|\le\prod_j(1+\|E_j\|)-1$, and $\|e^{X}-1\|\le e^{\|X\|}-1$. The path has
$N_x$ bonds, each factor of $\mathsf{P}_{f,c}(x)$ and of its inverse has exponent of size at
most $\xi\bar{\mathfrak{c}}$, and $\xi N_x\le 7d\mathsf{M}$; hence, by the ceiling
$7d\mathsf{M}\bar{\mathfrak{c}}\le\frac1{10}$ of \eqref{4.5:eq-fine-axial-frame-a},
\begin{equation*}
\|\mathsf{P}_{f,c}(x)^{\pm1}-1\|\le e^{\xi N_x\bar{\mathfrak{c}}}-1
\le e^{7d\mathsf{M}\bar{\mathfrak{c}}}-1=\delta_f\le e^{1/10}-1\le\tfrac19 .
\end{equation*}
The adjoint action of a unitary is isometric, and
$\|\operatorname{Ad}_{\mathsf{P}_{f,c}}\|\le\|\mathsf{P}_{f,c}\|\,\|\mathsf{P}_{f,c}^{-1}\|
\le(1+\delta_f)^2$, likewise for the inverse. Since
$\operatorname{Ad}_{\mathsf{P}_f(x)}=\operatorname{Ad}_{u_0(y_0)^{-1}}
\operatorname{Ad}_{\mathsf{P}_{f,c}(x)}\operatorname{Ad}_{u_0(x)}$ with $u_0$ unitary, this
gives $\|\operatorname{Ad}^{\pm1}_{\mathsf{P}_f(x)}\|\le(1+\delta_f)^2\le(10/9)^2$.

\smallskip\noindent\textit{Tree bonds.}
Let $b=(b_-,b_+)$ be a bond of the tree $\mathcal{T}$ with
$\Gamma^{0}_{y_0,b_+}=\Gamma^{0}_{y_0,b_-}\cup b$. Then
$\mathsf{P}_f(b_+)=\mathsf{P}_f(b_-)\mathbf{U}(b)$, so
\begin{equation*}
W(b)=\mathsf{P}_f(b_-)\mathbf{U}(b)\mathbf{U}(b)^{-1}\mathsf{P}_f(b_-)^{-1}=1 .
\end{equation*}
If instead $\Gamma^{0}_{y_0,b_-}=\Gamma^{0}_{y_0,b_+}\cup b$ (with $b$ traversed backwards),
then $\mathsf{P}_f(b_-)=\mathsf{P}_f(b_+)\mathbf{U}(b)^{-1}$, and the same computation with the
roles of the end-points exchanged gives $W(b)=1$.

\smallskip\noindent\textit{Lattice Stokes in an axial gauge.}
Let $b=(x,x+\xi e_\mu)$ be a bond of $\square_0$ not in $\mathcal{T}$, and put
$\mathsf{p}_0:=\Gamma^{0}_{y_0,x}$, $\mathsf{p}_1:=\Gamma^{0}_{y_0,x+\xi e_\mu}$. Since $x$ and
$x+\xi e_\mu$ have the same coordinates in every direction $\nu\ne\mu$, the two paths coincide
through the stages $1,\dots,\mu-1$ and through a common part of stage $\mu$. Let $z_0$ be the
end of the $\mu$-stage of $\mathsf{p}_0$, that is, the site whose coordinates in the directions
$\nu\le\mu$ are those of $x$ and whose coordinates in the directions $\nu>\mu$ are those of
$y_0$; write $\mathsf{r}_0:=(z_0,z_0+\xi e_\mu)$. If $x_\mu\ge y_{0,\mu}$, the $\mu$-stage of
$\mathsf{p}_1$ is the $\mu$-stage of $\mathsf{p}_0$ followed by $\mathsf{r}_0$. If
$x_\mu<y_{0,\mu}$, the $\mu$-stage of $\mathsf{p}_0$ is the $\mu$-stage of $\mathsf{p}_1$, which
ends at $z_0+\xi e_\mu$, followed by $\mathsf{r}_0$ reversed. In either case
$\mathsf{r}_0\in\mathcal{T}$, and $W(\mathsf{r}_0)=1$ by the preceding paragraph.

After stage $\mu$ both paths perform the same moves, in the directions $\nu>\mu$: the remainder
of $\mathsf{p}_0$ is a sequence of sites $(z_0,z_1,\dots,z_{\mathsf{n}}=x)$, and the remainder
of $\mathsf{p}_1$ is its translate $(z_0+\xi e_\mu,\dots,z_{\mathsf{n}}+\xi e_\mu)$, with
$\mathsf{n}\le N_x$; all bonds of both remainders are bonds of $\mathcal{T}$. Define the rungs
and the plaquettes
\begin{equation*}
\mathsf{r}_i:=(z_i,z_i+\xi e_\mu),\qquad
p_i:=(z_{i-1},\,z_i,\,z_i+\xi e_\mu,\,z_{i-1}+\xi e_\mu),\qquad 1\le i\le\mathsf{n},
\end{equation*}
so that $\mathsf{r}_{\mathsf{n}}=b$. The plaquette variable of $W$ on $p_i$, based at $z_{i-1}$,
is
\begin{align*}
W(\partial p_i)&=W(z_{i-1},z_i)\,W(\mathsf{r}_i)\,
W(z_i+\xi e_\mu,z_{i-1}+\xi e_\mu)\,W(\mathsf{r}_{i-1})^{-1}\\
&=W(\mathsf{r}_i)\,W(\mathsf{r}_{i-1})^{-1},
\end{align*}
because the first and the third factor are variables on bonds of $\mathcal{T}$ (the third
traversed backwards), which equal $1$. Hence
$W(\mathsf{r}_i)=W(\partial p_i)W(\mathsf{r}_{i-1})$ for $1\le i\le\mathsf{n}$, and
telescoping from $W(\mathsf{r}_0)=1$ gives the exact identity
\begin{equation*}
W(b)=W(\partial p_{\mathsf{n}})\,W(\partial p_{\mathsf{n}-1})\cdots W(\partial p_1).
\end{equation*}
Plaquette variables conjugate under gauge maps: for a plaquette $p$ based at $z$, the inner
factors $\mathsf{P}_f(\cdot)^{\mp1}$ cancel around the loop and
$W(\partial p)=\mathsf{P}_f(z)\,\mathbf{U}(\partial p)\,\mathsf{P}_f(z)^{-1}$. Hypothesis (iii),
$|\mathbf{U}(\partial p)-1|<\alpha_0'''\xi^2$, holds on the whole lattice, hence also for the
plaquettes having a vertex in $\square_0^{+}\setminus\square_0$; with the bound on
$\operatorname{Ad}_{\mathsf{P}_f}$ proved above,
\begin{equation*}
\|W(\partial p)-1\|\le\|\operatorname{Ad}_{\mathsf{P}_f(z)}\|\,|\mathbf{U}(\partial p)-1|
<(10/9)^2\alpha_0'''\xi^2<2\alpha_0'''\xi^2 .
\end{equation*}
Write $W(\partial p_i)=1+E_i$. Then $2\mathsf{n}\alpha_0'''\xi^2\le14d\mathsf{M}\alpha_0'''\xi
\le\frac12$, a bound contained in the ceiling
$56d\mathsf{M}(\alpha_0'''+C_{14}\alpha_1''')\le\frac14$ of
\eqref{4.5:eq-fine-axial-frame-a}; using $e^{t}-1\le 2t$ for $0\le t\le\frac12$,
\begin{align*}
\|W(b)-1\|&\le\prod_{i=1}^{\mathsf{n}}(1+\|E_i\|)-1\le(1+2\alpha_0'''\xi^2)^{\mathsf{n}}-1
\le e^{2\mathsf{n}\alpha_0'''\xi^2}-1\\
&\le4\mathsf{n}\alpha_0'''\xi^2\le 28d\mathsf{M}\alpha_0'''\xi ,
\end{align*}
where $\mathsf{n}\le N_x\le 7d\mathsf{M}/\xi$ was used in the last step. The power series of the
logarithm defines $\xi\mathcal{A}_f(b):=-i\log W(b)$ and gives
$|\xi\mathcal{A}_f(b)|\le2\|W(b)-1\|$, so that
\begin{equation*}
|\mathcal{A}_f|\le56d\mathsf{M}\alpha_0'''=b_f .
\end{equation*}
On bonds of $\mathcal{T}$, $\mathcal{A}_f=0$. Adjacent ladders of plaquettes carry independent
fluxes, so no smallness of the differences of $\mathcal{A}_f$ in transverse directions is
asserted; none is used below, since \cite[Proposition~6]{B-CMP98} reads of its argument only
the supremum of $|A'|$.

\smallskip\noindent\textit{Covariance of the window data at the site map $\mathsf{P}_f$.}
For $G$-valued $U$ and $\mathfrak{g}$-valued $A'$ the field $\mathbf{U}$ is $G$-valued, hence
$\mathsf{P}_f\in G$, and the identity $M^{\cdot}(W)=(M^{\cdot}\mathbf{U})^{\bar{\mathsf{P}}_f}$
on $\mathfrak{B}(\square_0)$, with $\bar{\mathsf{P}}_f$ the restriction of $\mathsf{P}_f$ to the
block base sites, is the covariance of averages of \cite{B-CMP98}: if a gauge transformation
$v$ acts on a configuration $V$ by $V^{v}(b)=v(b_-)V(b)v(b_+)^{-1}$, then the averaged
configuration obeys $\overline{V^{v}}(c)=v(c_-)\bar V(c)v(c_+)^{-1}$. Its proof there is that
the matrices entering the average are unitarily equivalent, hence have equal eigenvalues, hence
unitarily equivalent logarithms. For an invertible, non-unitary map one reads ``similar''
instead of ``unitarily equivalent'': by \cite[(56)--(57)]{B-CMP98}, for an arbitrary invertible
matrix $X$ the operator $R(X)Y=XYX^{-1}$ satisfies $R(X)f(Y)=f(R(X)Y)$ for analytic $f$, and
this is applied to $f$ the logarithm in \cite[(42)--(43)]{B-CMP98}, provided its arguments at
every level lie within the disc of the power series. This is the case at $\mathfrak{g}$-valued
$A'$ by \cite[Proposition~2]{B-CMP98}; for general $A'\in\mathcal{P}_U$ we argue instead by
holomorphy. Both sides of
$M^{\cdot}(\mathbf{U}^{\mathsf{P}_f})=(M^{\cdot}\mathbf{U})^{\bar{\mathsf{P}}_f}$ are
holomorphic on $\mathcal{P}_U$. The left side is so by \cite[Proposition~6]{B-CMP98} at base $1$
applied to $W=e^{i\xi\mathcal{A}_f}$: on $\mathcal{P}_U$ the lattice Stokes argument above, run
with hypothesis (iii) replaced by the bound $(\alpha_0'''+C_{14}\alpha_1''')\xi^2$ of the
plaquette expansion used in the change of base by the identity theorem in the proof of
Lemma~\ref{4.5:lem-literal-frame}, gives
\begin{equation*}
|\mathcal{A}_f|\le b_f^{\mathrm{poly}}:=56d\mathsf{M}(\alpha_0'''+C_{14}\alpha_1''')
\le\alpha_1^{(98)}
\end{equation*}
by the ceiling \eqref{4.5:eq-fine-axial-frame-a}, where $\alpha_1^{(98)}$ is the radius
$\alpha_1$ of \cite[Proposition~6]{B-CMP98}. The right side is holomorphic by
\cite[Proposition~6]{B-CMP98} at base $U$ together with the polynomial dependence of the gauge
action on $\bar{\mathsf{P}}_f^{\pm1}$. The identity theorem, as in the change of base in the
proof of Lemma~\ref{4.5:lem-literal-frame}, therefore transports the identity from the
$\mathfrak{g}$-valued points to all of $\mathcal{P}_U$. Either way,
$M^{\cdot}(W)=(M^{\cdot}\mathbf{U})^{\bar{\mathsf{P}}_f}$.

\smallskip\noindent\textit{Sizes of the window data.}
We use \cite[Proposition~6, (164)]{B-CMP98} at the base $U_0:=1$, in the following form, in
which the lattice spacing of \cite{B-CMP98} is the window's fine spacing $\xi$ and the number
of averaging steps is the level $n$: if $U_0$ satisfies \cite[(52)]{B-CMP98}, then the
$n$-fold average of $U'U_0$ is an analytic function of the logarithm $(i\xi)^{-1}\log U'$, for
logarithms with values in the complexified algebra and of supremum below the radius
$\alpha_1^{(98)}$, and it satisfies
$|\overline{U'U_0}^{\,n}(\bar U_0^{\,n})^{-1}-1|<O(1)\alpha_1^{(98)}$. This comes with
the level sizes \cite[(159), (161)]{B-CMP98}, by which the logarithm of the level-$j$ averaged
field has modulus less than $2\alpha_1^{(98)}L^{j}\xi$, and \cite[(163)]{B-CMP98},
\begin{equation*}
|v_n(x)-1|<O(1)\alpha_1^{(98)}\sum_{j=0}^{n-1}L^{j+1}\xi\le O(1)\alpha_1^{(98)} ,
\end{equation*}
for the conjugating factor $v_n$. In these bounds the radius enters as the bound on the
supremum of the logarithm; they are applied below with that supremum at most $b_f$. Apply this
with $U':=W=e^{i\xi\mathcal{A}_f}$, whose logarithm $(i\xi)^{-1}\log W$ is $\mathcal{A}_f$;
$U_0=1$ satisfies \cite[(52)]{B-CMP98} trivially and its averages are $1$. Let $\ell_n$ denote
the length of a level-$n$ block bond in units of the
top block. For $c\in\Lambda_n(\square_0)$, $n\ge1$, the level-$n$ averaged bond variable is
$e^{iB_f(c)}$ with
\begin{equation*}
|B_f(c)|\le c_Q\,b_f\,\ell_n\le c_Q\,b_f ,
\end{equation*}
the amplitude $b_f$ times the block-bond length. On the seam and crossing block bonds of the
hierarchy the bound is the one given in the statement, their data being the explicit analytic
products of at most $2dL$ averaged or fine variables of $W$ described there, each within
$c_Qb_f$ of $1$. Finally, since $L\mathsf{M}=M$,
\begin{equation*}
\mathfrak{v}_f=8dL\cdot c_Q\cdot56d\mathsf{M}\alpha_0'''=448d^2c_QL\mathsf{M}\alpha_0'''
=448d^2c_QM\alpha_0''' .
\end{equation*}
On $\mathcal{P}_U$ the same computation with $b_f^{\mathrm{poly}}$ in place of $b_f$ gives
\begin{equation*}
\sup|B_f|\le8dLc_Qb_f^{\mathrm{poly}}=448d^2c_QM(\alpha_0'''+C_{14}\alpha_1''')
=\mathfrak{v}_f^{\mathrm{poly}} .
\end{equation*}

\smallskip\noindent\textit{The chart at base one.}
The argument is the one given for the chart at base one in the proof of
Lemma~\ref{4.5:lem-literal-frame}, with $\sup|B_f|\le\mathfrak{v}_f$ in place of the bound on
$B_V$ (and $\sup|B_f|\le\mathfrak{v}_f^{\mathrm{poly}}\le\frac12C_1\bar\varepsilon_1$ on
$\mathcal{P}_U$, by \eqref{4.5:eq-fine-axial-frame-a}). At base $1$,
\cite[Proposition~9]{B-CMP102b} gives the axial chart
$U_{k+1}(\square_0,e^{iB_f})=(e^{i\eta\mathcal{A}'})^{g'}$ with
$\mathcal{A}'=L^{-1}\mathcal{H}^{\square_0}(B_f)$ and $g'$ the Landau-to-axial map. The decay
bound for the kernel of $\mathcal{H}^{\square_0}$ at base one, in which the covariant gradient is
the plain one, is integrated along $tB_f$, $t\in[0,1]$, and summed over the sources with
\cite[Lemma~2.1]{B-CMP96}, each block site carrying at most $3dL$ block bonds; the passage from
window units to $\eta$-units improves the bounds by a factor $L^{-1}\le1$. This yields
\begin{equation*}
|\mathcal{A}'|,\ |\nabla^{\eta}\mathcal{A}'|,\ \|\mathcal{A}'\|_{1,\beta}
\le3dLC_Hc_1\mathfrak{v}_f=b' ,
\end{equation*}
where $C_H$ is the constant of the kernel decay bound and $c_1$ the constant bounding the sum of
its exponential decay factors over block sites, the constants entering $B_b^{f}$ in
\eqref{4.5:eq-fine-axial-frame-a}.
The map $g'$ is treated as $g$ there, with the same normalisation and bound, layer by layer.
Fix a layer $\Lambda_q(\square_0)$ and a block $B^{q}(x_q)$ of it, with base site $x_q$. There
$\ell_q\sup|\mathcal{H}^{\square_0}(B_f)|\le b'$ in window units, so $|\mathcal{A}'|\le
b'/(L\ell_q)$ on the legs. Below its top base-site value, $g'$ is the nested-contour transport of
the averages of $e^{i\eta\mathcal{A}'}$, each leg consisting of at most $dL$ bonds of the
level-$p$ block lattice, $p<q$. By \cite[Proposition~6]{B-CMP98} at base $1$, whose radius
condition is the bound $b'/L\le\alpha_1^{(98)}$ on the rescaled amplitude, the second member of
the ceiling on $b'$ in \eqref{4.5:eq-fine-axial-frame-a}, each such averaged bond variable is
$e^{iB_p(c)}$ with $|B_p(c)|\le c_Qb'\ell_p/\ell_q$. Hence every nested transport is within
\begin{equation*}
\sum_{p<q}dL\,c_Qb'\,\frac{\ell_p}{\ell_q}\le2dc_Qb'
\end{equation*}
of $1$, and the top value likewise. The bound $dc_Qb'\le\frac1{24}$ follows from
$4dLc_Qb'\le\frac16$ in \eqref{4.5:eq-fine-axial-frame-a}, and as for $g$ it gives
\begin{equation*}
|\log g'|\le4dc_Qb'\bigl(1+O(dc_Qb')\bigr)\le2dLc_Qb' .
\end{equation*}
The bounds $|\log g'|\le2dLc_Qb'$ and
$\|\operatorname{Ad}^{\pm1}_{g'}\|\le e^{4dLc_Qb'}$ of the statement hold on $\square_0$.

\smallskip\noindent\textit{Agreement of the two charts.}
We argue as in the change of base by the identity theorem in the proof of
Lemma~\ref{4.5:lem-literal-frame}, now with
\begin{equation*}
F_1(A'):=S[A']^{\tilde{\bar{\mathsf{P}}}_f[A']},\qquad
F_2(A'):=U_{k+1}\bigl(\square_0,M^{\cdot}(W[A'])\bigr)
=U_{k+1}\bigl(\square_0,e^{iB_f[A']}\cdot1\bigr),
\end{equation*}
where $S[A']=U_{k+1}(\square_0,M^{\cdot}\mathbf{U}_{A'})$ and $\mathbf{U}_{A'}=e^{i\xi A'}U$.
Both are holomorphic on $\mathcal{P}_U$. For $F_1$, $S$ is holomorphic as in the change of base,
and $\mathsf{P}_f$ is a product of bond variables of $e^{i\xi A'}U$ and of their inverses, hence
entire in $A'$, with the gauge action polynomial in $\tilde{\bar{\mathsf{P}}}_f^{\pm1}$. For
$F_2$, $B_f[A']$ is holomorphic with values in the set \cite[(172)]{B-CMP102b} over the base $1$
by the analyticity and size bounds for the window data proved above and the ceiling
\eqref{4.5:eq-fine-axial-frame-a} on $\mathfrak{v}_f^{\mathrm{poly}}$.

Let $A'\in\mathcal{P}_U$ be $\mathfrak{g}$-valued. Then $\mathsf{P}_f\in G$, because
$\mathfrak{c}$ is $\mathfrak{g}$-valued and $u_0\in G$. Hence $W=\mathbf{U}^{\mathsf{P}_f}$ is
$G$-valued, its plaquette variables are those of $\mathbf{U}_{A'}$ conjugated, and
\begin{equation*}
|W(\partial p)-1|=|\mathbf{U}_{A'}(\partial p)-1|<(\alpha_0'''+C_{14}\alpha_1''')\xi^2 .
\end{equation*}
By \cite[Proposition~2]{B-CMP98} and the statement of \cite{B-CMP102b}, made for real fields,
that the block configuration of such a field satisfies \cite[(7)]{B-CMP102b} with
$\varepsilon_1=O(\varepsilon_0)$, the data $M^{\cdot}W$ and $M^{\cdot}\mathbf{U}_{A'}$ are real
and satisfy \cite[(7)]{B-CMP102b} with the constant $\varepsilon_1$ there equal to
$C_7(\alpha_0'''+C_{14}\alpha_1''')\le\bar\varepsilon_1$, by the ceiling
\eqref{4.5:eq-literal-frame-a}. Both charts are then the axial minimisers for their real data:
the chart coincides with the minimal configuration \cite{B-CMP102b}, the minimal orbit is
unique by \cite[Theorem~1]{B-CMP102b}, applicable because $\bar\varepsilon_1\le a_1$ and
$B_3a_1\le a_0$ in the constants of that theorem, as in the change of base, and the axial
conditions determine a unique element in each orbit \cite{B-CMP99a}. The real form of
\cite[(181)]{B-CMP102b} at $G$-valued block-site maps, proved in the change of base in the
proof of Lemma~\ref{4.5:lem-literal-frame} and displayed as \eqref{4.5:eq-literal-frame-b},
applied at the
$G$-valued map $\bar{\mathsf{P}}_f$, together with the covariance proved above, gives
\begin{align*}
U_{k+1}(\square_0,M^{\cdot}W)
&=U_{k+1}\bigl(\square_0,(M^{\cdot}\mathbf{U}_{A'})^{\bar{\mathsf{P}}_f}\bigr)\\
&=U_{k+1}(\square_0,M^{\cdot}\mathbf{U}_{A'})^{\tilde{\bar{\mathsf{P}}}_f}
=S[A']^{\tilde{\bar{\mathsf{P}}}_f[A']} .
\end{align*}
Thus $F_1=F_2$ at every $\mathfrak{g}$-valued point of $\mathcal{P}_U$. By the identity theorem,
exactly as in the change of base, $F_1=F_2$ on $\mathcal{P}_U$, in particular at the given $A'$.
No radius of the neighbourhood of $G$-valued transformations in \cite[(181)]{B-CMP102b} is used,
since that statement enters at $G$-valued maps only. Accordingly the second member of the
condition on $\delta_f$ in \eqref{4.5:eq-fine-axial-frame-a} is not used in this proof, and
$\delta_f\le\frac19$ serves only the bound on $\|\operatorname{Ad}_{\mathsf{P}_f}\|$. The
holomorphy on $\mathcal{P}_U$ is paid for by the ceiling on $\mathfrak{v}_f^{\mathrm{poly}}$ in
\eqref{4.5:eq-fine-axial-frame-a}.

\smallskip\noindent\textit{The frame map.}
By the chart at base one and the agreement of the two charts,
$(e^{i\eta\mathcal{A}'})^{g'}=U_{k+1}(\square_0,M^{\cdot}W)
=S^{\tilde{\bar{\mathsf{P}}}_f}$. Acting with $g'^{-1}$ and using
$(V^{u})^{v}=V^{vu}$,
\begin{equation*}
e^{i\eta\mathcal{A}'}=\bigl(S^{\tilde{\bar{\mathsf{P}}}_f}\bigr)^{g'^{-1}}
=S^{g'^{-1}\tilde{\bar{\mathsf{P}}}_f}=S^{w'} .
\end{equation*}
Since $\tilde{\bar{\mathsf{P}}}_f(x)=\mathsf{P}_f(y_x)$,
\begin{equation*}
\|\operatorname{Ad}^{\pm1}_{w'}\|\le\|\operatorname{Ad}^{\mp1}_{g'}\|\cdot
\|\operatorname{Ad}^{\pm1}_{\mathsf{P}_f(y_x)}\| ,
\end{equation*}
and the bounds on $\operatorname{Ad}_{\mathsf{P}_f}$ and $\operatorname{Ad}_{g'}$ proved above
give the bound of the statement. The remaining two assertions about the frame hold as stated.
By construction $w'$ is built from products of bond variables of $\mathbf{U}$ (through
$\mathsf{P}_f$), from averages, analytic by \cite[Proposition~6]{B-CMP98}, and from
$\mathcal{H}^{\square_0}$, analytic by \cite[Proposition~9]{B-CMP102b}; the map $u_0$ enters
bounds only. The extension $w'^{\mathrm{ext}}$ and the identity for $\mathbf{A}'$ are obtained
exactly as the extension $w^{\mathrm{ext}}$ and the identity for $\mathbf{A}$ in the proof of
Lemma~\ref{4.5:lem-literal-frame}, with $w'$ in place of $w$.

\smallskip\noindent\textit{Plain bounds.}
The proofs of the corresponding plain bounds for $\mathbf{A}$ and for its combination with
$\mathcal{A}$, among the plain regularity bounds in the proof of
Lemma~\ref{4.5:lem-literal-frame}, read of
the frame only the relation $S^{w}=e^{i\eta\mathcal{A}}$, the bounds
$|\mathcal{A}|,|\nabla\mathcal{A}|\le b\le1/(24dLc_Q)$, where $b=B_bM\alpha_0'''$ is the
constant of that lemma, and $\|\operatorname{Ad}^{\pm1}_{w}\|\le2$; no gradient of $w$ is used.
These three facts hold for $(w',\mathcal{A}',b')$: the first and the third by the paragraph on
the frame map, the second by the chart at base one, and $b'\le1/(24dLc_Q)$ by
$4dLc_Qb'\le\frac16$ in \eqref{4.5:eq-fine-axial-frame-a}. Hence those proofs apply verbatim
with $(w,\mathcal{A},b)$ replaced by $(w',\mathcal{A}',b')$.
\end{proof}

\begin{remark}[Scope and gauge covariance of the fine axial frame]\label{4.5:rem-fine-frame-scope}
The frame of Lemma~\ref{4.5:lem-fine-axial-frame} is axial for the \emph{field}: the fine field
$W=\mathbf{U}^{\mathsf{P}_f}$ equals $1$ on every bond of the tree $\mathcal{T}$ of lexicographic paths
issuing from $y_0$. The frame of \cite[(3.26)]{B-CMP109} is axial for the \emph{data}: there the block
data $V$ equal $1$ on the tree of the hierarchy. The near member of the operator bounds at free current,
and the arguments that use it (\cite[(3.33)--(3.43)]{B-CMP109}, the argument of this chapter that
carries out those equations, and \cite{B-CMP116}), use only the
following properties of the frame.
\begin{enumerate}
\item The background in the frame is $e^{i\eta\mathcal{A}}$ on $\square_0$. The bounds
$|\mathcal{A}|,|\nabla^{\eta}\mathcal{A}|,\|\mathcal{A}\|_{1,\beta}\le b$ hold on $\tilde\square^{4}$,
with $b$ depending on $\alpha_0$ only and on the first power of $M$.
\item The frame map is an explicit analytic function of $\mathbf{U}\restriction\square_0^{+}$ (in the
words of \cite{B-CMP109}, of $\mathbf{U}$ restricted to the cube $\square_0$), and its
adjoint action satisfies $\|\operatorname{Ad}^{\pm1}\|\le 2$.
\item The data of the cube $\square_0$ in the frame are $e^{iB}$ over the base $1$, with
$\sup|B|\le\tfrac12 C_1\bar\varepsilon_1$. Consequently the expansions of \cite[(3.33)]{B-CMP109}
at $U_j(\square_0,1)$ apply, since the constraint that the block averages $M^{\cdot}$ of the minimiser
equal the prescribed data holds for arbitrary data.
\end{enumerate}
Lemma~\ref{4.5:lem-fine-axial-frame} supplies all three. The literal $V$ of
\cite[(3.26)]{B-CMP109} lies one further block-site gauge map away, namely the map that makes the
data $e^{iB_f}$ axial. That map is small of order $\alpha_0$, and Ba{\l}aban's $V$ then satisfies
$|V-1|=O(M^{2}\alpha_0/L)$, which is a bound
of the same form with a worse constant. The single power of $M$ in Ba{\l}aban's bound for his $V$ is
the only gain that the averaging estimates at a complex field would bring there, and none of the
arguments listed above uses it.

Moreover, the frame is covariant under the gauge transformation \cite[(3.29)]{B-CMP109}. This
covariance is not used by the operator bounds at free current; it is recorded because it shows that
the sentence of \cite{B-CMP109} quoted below holds in the fine axial frame. Let $u$ be a
sitewise gauge map, and put $\bar u:=u\restriction(\text{block sites})$. Suppose that the covariance
\cite[(181)]{B-CMP102b} of the axial minimiser under block-site maps holds at $\bar u$. This is the
case if $u$ is $G$-valued, or if $u=u'u_G$ with $u_G$ $G$-valued and $u'$ a $G^{c}$-valued map in the
range in which that covariance is established for complex block-site maps (the complex block-site
maps within the radius $\rho_0/8$, with $\rho_0$ the radius entering the hypotheses of
Lemma~\ref{4.5:lem-fine-axial-frame}). Let
$\tilde u(x):=\bar u(y_x)$, where, as in Lemma~\ref{4.5:lem-fine-axial-frame}, $y_x$ denotes the block
base site of the block containing $x$; thus $\tilde u$ is the block-constant extension of $\bar u$.
Write $R(X)Y:=XYX^{-1}$ for Ba{\l}aban's conjugation operator \cite[(56)]{B-CMP98}. Under
\cite[(3.29)]{B-CMP109} with this $u$, the background function $\mathbf{A}'$ of
\eqref{4.5:eq-fine-axial-frame-b} transforms as
\begin{equation}\label{4.5:eq-fine-frame-scope-1}
\mathbf{A}'\longmapsto R(u(y_0))\,\mathbf{A}'.
\end{equation}
Thus the sentence of \cite{B-CMP109}, ``The expressions in (3.28) transform in a very simple way
under (3.29), namely by $R(u(y_0))$'', holds in the
fine axial frame.
\end{remark}

\begin{proof}
We first check that Lemma~\ref{4.5:lem-fine-axial-frame} supplies the three properties, with
$(\mathcal{A},b)=(\mathcal{A}',b')$, frame map $w'$, and $B=B_f$; all references are to the bounds
\eqref{4.5:eq-fine-axial-frame-b}.

For the first property, the chart agreement gives $S^{w'}=e^{i\eta\mathcal{A}'}$ on the bonds of
$\square_0$. The bounds on the background give
$|\mathcal{A}'|,|\nabla^{\eta}\mathcal{A}'|,\|\mathcal{A}'\|_{1,\beta}\le b'=B_b^{f}M\alpha_0'''$ on
$\tilde\square^{4}$; this involves $\alpha_0$ only and $M$ to the first power.

For the second property, $w'=g'^{-1}\tilde{\bar{\mathsf{P}}}_f$ is built from products of bond
variables of $\mathbf{U}$, block averages and $\mathcal{H}^{\square_0}$. It is therefore an explicit
analytic function of $\mathbf{U}$ on $\square_0^{+}$, and
$\|\operatorname{Ad}_{w'}^{\pm1}\|<2$.

For the third property, the data are $e^{iB_f}:=M^{\cdot}(W)$ over the base $1$, and on the twist
polydisc $\mathcal{P}_U$ they satisfy
\begin{equation*}
\sup|B_f|\le\mathfrak{v}_f^{\mathrm{poly}}\le\tfrac12 C_1\bar\varepsilon_1 .
\end{equation*}

We turn to the covariance. We write $\mathbf{U}^{u}(b)=u(b_-)\mathbf{U}(b)u(b_+)^{-1}$.

Along $\Gamma^{0}_{y_0,x}$ the inner factors of $u$ cancel in the ordered product, on backward bonds
as well. Hence $\mathbf{U}^{u}(\Gamma^{0}_{y_0,x})=u(y_0)\,\mathbf{U}(\Gamma^{0}_{y_0,x})\,u(x)^{-1}$,
that is, $\mathsf{P}_f\mapsto u(y_0)\mathsf{P}_fu^{-1}$.

The new fine field on a bond $b$ is
\begin{align*}
W^{\mathrm{new}}(b)
&=u(y_0)\mathsf{P}_f(b_-)u(b_-)^{-1}\,u(b_-)\mathbf{U}(b)u(b_+)^{-1}\,
u(b_+)\mathsf{P}_f(b_+)^{-1}u(y_0)^{-1}\\
&=\operatorname{Ad}_{u(y_0)}W(b).
\end{align*}
Apply the covariance of block averages under site maps, as in the proof of the window-data bound of
Lemma~\ref{4.5:lem-fine-axial-frame}, at the constant map $u(y_0)$;
since the logarithm is a power series, it commutes with conjugation. Together these give
$B_f\mapsto\operatorname{Ad}_{u(y_0)}B_f$.

By \eqref{4.5:eq-fine-axial-frame-b}, $\mathcal{A}'=L^{-1}\mathcal{H}^{\square_0}(B_f)$ and
$U_{k+1}(\square_0,e^{iB_f})=(e^{i\eta\mathcal{A}'})^{g'}$. The constant map $u(y_0)$ fixes the base
$1$, so the gauge covariance and reality of the chart on the orbit class, proved in
step~\ref{4.5:prf-chart-covariance} and applied at $u(y_0)$, gives
$\mathcal{A}'\mapsto\operatorname{Ad}_{u(y_0)}\mathcal{A}'$, and the Landau-to-axial map built from it
satisfies $g'\mapsto\operatorname{Ad}_{u(y_0)}g'$.

The data $M^{\cdot}\mathbf{U}$ become $(M^{\cdot}\mathbf{U})^{\bar u}$. By \cite[(181)]{B-CMP102b} at
$\bar u$, this gives $S\mapsto S^{\tilde u}$.

For the frame map, the block-constant map changes pointwise as
\begin{equation*}
\tilde{\bar{\mathsf{P}}}_f(x)\mapsto u(y_0)\tilde{\bar{\mathsf{P}}}_f(x)\tilde u(x)^{-1}.
\end{equation*}
Substituting this and the new $g'$ into $w'=g'^{-1}\tilde{\bar{\mathsf{P}}}_f$ gives
\begin{equation*}
w'^{\mathrm{new}}=u(y_0)\,g'^{-1}u(y_0)^{-1}\,u(y_0)\,\tilde{\bar{\mathsf{P}}}_f\,\tilde u^{-1}
=u(y_0)\,w'\,\tilde u^{-1}.
\end{equation*}

Finally, by \eqref{4.5:eq-fine-axial-frame-b},
$\mathbf{A}'=(t\tilde\zeta_{\square}+t_{\square}\zeta_{\square})\,R(w')
H_k(B';\mathbf{U}^{\mathrm{gl}},\mathbf{J})$. The scalar cut-off factor is unchanged. By the
gauge covariance of the chart proved in step~\ref{4.5:prf-chart-covariance},
\begin{equation}\label{4.5:eq-fine-frame-scope-2}
R(w'^{\mathrm{new}})\,H_k\bigl(R(\bar u)B';(\mathbf{U}^{\mathrm{gl}})^{\tilde u},
R(u)\mathbf{J}\bigr)
=R(u(y_0))\,R(w')\,H_k(B';\mathbf{U}^{\mathrm{gl}},\mathbf{J}).
\end{equation}
The transformed background function is $(t\tilde\zeta_{\square}+t_{\square}\zeta_{\square})$ times
the left member of \eqref{4.5:eq-fine-frame-scope-2}, so it equals $R(u(y_0))\mathbf{A}'$,
which is \eqref{4.5:eq-fine-frame-scope-1}.
\end{proof}

\begin{lemma}[The window chart reproduces the whole-lattice minimiser]\label{4.5:lem-window-minimiser}
Let $U_{k+1}$ be the minimiser on the whole lattice for real data satisfying
\cite[(7)]{B-CMP102b}, with $\varepsilon_1^{[15]}$ small, taken in the nested axial gauge
\cite[(1.15)]{B-CMP99a}. Let the block set $\mathfrak{B}(\square_0)$ of the window
$\square_0$ be built on the blocks and contours of the hierarchy of the whole lattice
(the reading $(\rho 2)$ of the corresponding display). Then
\begin{equation}\label{4.5:eq-window-minimiser-1}
U_{k+1}(\square_0, M^{\cdot}U_{k+1}) = U_{k+1}
\qquad\text{on the bonds of } \square_0 .
\end{equation}
Hence $\mathbf{U}^{\mathrm{gl}}[U_{k+1}] = U_{k+1}$; the pair
$(\mathbf{U}^{\mathrm{gl}}[U_{k+1}], J_{k+1})$, formed with the current $J_{k+1}$ of the
step, is a true pair, that is, it is the physical point of the chart, the real
whole-lattice minimiser taken with its own current; and, by the exactness of the chart on
the diagonal, \eqref{4.5:eq-diagonal-exactness-a}, the function $H_k(B')$ of the chart at
the datum $B'$, evaluated at this pair, is the function $\mathcal{H}(B')$ of
\cite{B-CMP102b}.
\end{lemma}

\begin{proof}
Throughout, the bonds and plaquettes of a domain $\Omega_j$ (the letter denotes here the
domains of \cite{B-CMP102b}) are those ``with at least one
end-point belonging to $\Omega_j$'' \cite{B-CMP102b}, \cite{B-CMP99a}. Here
$U_{k+1}(\square_0, M^{\cdot}U_{k+1})$ is the minimiser of the variational problem in the
window $\square_0$ whose constraints prescribe, on the layers $\Lambda_p(\square_0)$ of the
window tower $\{\square_n\}$, the block averages $\bar U_{k+1}^{\,p}$ of the tower
$M^{\cdot}U_{k+1}$; we call this the window problem, and we call the minimisation on the
whole lattice solved by $U_{k+1}$ the global problem.

First, the restriction $U_{k+1}{\upharpoonright}_{\square_0}$ meets the constraints of the
window problem, since these constraints are its own block averages.

Second, $U_{k+1}{\upharpoonright}_{\square_0}$ is critical for the window problem. By
\cite[(42)--(43)]{B-CMP98} the $(p+1)$-st block average is the one-step average of the
$p$-th, and
``$\bar U^k_c$, $c \subset \Omega^{(k)}$, depends only on the bond variables $U_b$ for
$b \subset B^k(c_-) \cup B^k(c_+)$'' \cite{B-CMP98}. We show by induction on $p$ that
every $p$-bond $c$ of the hierarchy of the whole lattice with both end-points outside
$\square_{p+1}$ has its average $\bar U^{\,p}_c$ fixed by the constraints of the window
problem, that is, left unchanged by every variation which preserves these constraints.
Let $c$ be such a bond, with end-points $c_-$, $c_+$, and write $B(c_-)$, $B(c_+)$ for the
blocks of these end-points. If an end-point of $c$ lies in $\Lambda_p(\square_0)$, then
$\bar U^{\,p}_c$ is itself one of the constraints. If both end-points of $c$ lie outside
$\square_p$, then $\bar U^{\,p}_c$ is the one-step average of $(p-1)$-bonds lying inside
$B(c_-)\cup B(c_+)$, and
\begin{equation*}
B(c_-)\cup B(c_+) \subset \square_p^{\,c};
\end{equation*}
these $(p-1)$-bonds have both end-points outside $\square_p$, so their averages are fixed
by the induction hypothesis, and therefore so is $\bar U^{\,p}_c$. It remains to consider
a $p$-bond $c$ with exactly one end-point in $\square_p$; by the convention on the bonds of
a domain recalled at the beginning of the proof, such a bond is a bond of the layer at
level $p$, so its average is pinned at level $p$ itself as one of the constraints. At
$p = k+1$ all the constraints of the global
problem are among the averages so fixed. Consequently a variation supported on bonds of
$\square_1$ which fixes $\bar U^{\,p}$ on $\Lambda_p(\square_0)$ for every $p$ fixes every
global average $\bar U^{\,k+1}_c$. Thus the variations of the window problem form a
subclass of the variations of the global problem, and $U_{k+1}$ is critical for the
latter; hence $U_{k+1}{\upharpoonright}_{\square_0}$ is critical for the window problem.

Third, $U_{k+1}{\upharpoonright}_{\square_0}$ lies in the space \cite[(6)]{B-CMP102b} of
the window problem, at
\begin{equation*}
\varepsilon_0 := O(1)\,B_3^2\,\varepsilon_1^{[15]} \le a_0 ,
\end{equation*}
where $a_0$ is the bound on $\varepsilon_0$ required in \cite{B-CMP102b} and $B_3$ is
Ba{\l}aban's constant entering the passage from $\varepsilon_1^{[15]}$ to $\varepsilon_0$.
Indeed, in \cite{B-CMP102b} ``by Proposition 2 [4] the configuration $V$ satisfies (7)
with $\varepsilon_1 = O(\varepsilon_0)$'', and the condition \cite[(2)]{B-CMP102b} on
$\Omega_j$ is weaker for smaller $j$. By the statement of \cite{B-CMP102b}, ``This orbit is
a unique critical orbit in the space (6)'', it follows from the first and second steps
that $U_{k+1}{\upharpoonright}_{\square_0}$ lies on the minimal orbit of the window problem
under the group of gauge transformations $u$ with $u = 1$ on $\mathfrak{B}(\square_0)$.

Fourth, we identify the representative on this orbit. The gauge condition
\cite[(1.15)]{B-CMP99a} is nested: in the notation of \cite{B-CMP99a}, where $B(x)$ is the
block of the site $x$ and $\Gamma_{x,y}$ the contour of the hierarchy from $x$ to $y$, it
reads, for the points $x_j, x_{j-1}, \dots, x_1, x$ of the hierarchy,
\begin{equation*}
\begin{gathered}
\bar U^{\,j-1}(\Gamma_{x_j,x_{j-1}}) = 1 \quad\text{for } x_{j-1}\in B(x_j), \qquad
\bar U^{\,j-2}(\Gamma_{x_{j-1},x_{j-2}}) = 1, \quad \dots, \\
U(\Gamma_{x_1,x}) = 1 \quad\text{for } x \in B(x_1)
\end{gathered}
\end{equation*}
\cite[(1.15), p.~78]{B-CMP99a}. Since $\mathfrak{B}(\square_0)$ is built on the blocks and
contours of the hierarchy of the whole lattice, the conditions of the window problem at a
point $x_p \in \Lambda_p(\square_0)$ are the members of levels $< p$ of the list of
conditions of the global problem; $U_{k+1}$ is in the nested axial gauge, so it satisfies
them. By \cite[after (1.15)]{B-CMP99a}, ``these equations together with (1.14) for gauge
transformations determine uniquely an element in each orbit''. Hence the representative
of the minimal orbit of the window problem selected by these conditions, which is
$U_{k+1}(\square_0, M^{\cdot}U_{k+1})$, coincides with $U_{k+1}$ on the bonds of
$\square_0$. This is \eqref{4.5:eq-window-minimiser-1}.

For the remaining assertions, recall from Proposition~\ref{4.5:prop-glued-window} that the
glued window background $\mathbf{U}^{\mathrm{gl}}[\mathbf{U}]$ equals
$U_{k+1}(\square_0, M^{\cdot}\mathbf{U})$ on the bonds of $\square_0$ and equals
$\mathbf{U}$ elsewhere. At $\mathbf{U} = U_{k+1}$, by \eqref{4.5:eq-window-minimiser-1} it
equals $U_{k+1}$ on the bonds of $\square_0$, and it equals $U_{k+1}$ elsewhere by
construction; hence $\mathbf{U}^{\mathrm{gl}}[U_{k+1}] = U_{k+1}$. The pair
$(\mathbf{U}^{\mathrm{gl}}[U_{k+1}], J_{k+1})$ is therefore the pair
$(U_{k+1}, J_{k+1})$, the real whole-lattice minimiser with its own current, which is a
true pair, and the exactness of the chart on the diagonal,
\eqref{4.5:eq-diagonal-exactness-a}, applied at this pair gives that $H_k(B')$ there
is the function $\mathcal{H}(B')$ of \cite{B-CMP102b}.
\end{proof}

\begin{remark}[Order of choice of the constants, cube size and amplitudes]
\label{4.5:rem-constants-order}
The letters of this chapter are fixed in the order below. At each stage a letter is chosen
only after every letter it reads. Here $\alpha''' := (1+2\beta)\alpha$ for an amplitude
$\alpha$, and
\begin{equation*}
b := B_b M\alpha_0''', \qquad b' := B_b^{f} M\alpha_0''',
\end{equation*}
where $\beta$, $B_b$, $B_b^{f}$ (see \eqref{4.5:eq-constants-order-1}), $M$ and $\alpha_0$ are
the letters fixed at the stages listed below.
The ceilings (c7), (c7$'$) and (c8) are those so labelled in the hypotheses of
Lemma~\ref{4.5:lem-literal-frame}, displayed in \eqref{4.5:eq-literal-frame-a}. The ceilings
(cf1), (cf4) and (cf5) are those so labelled in the hypotheses of
Lemma~\ref{4.5:lem-fine-axial-frame}, displayed in \eqref{4.5:eq-fine-axial-frame-a}.
\begin{enumerate}
\item The blocking factor $L$.
\item The constants that depend on $d$ and $L$ only and are taken from \cite{B-CMP96},
\cite{B-CMP98}, \cite{B-CMP99a}, \cite{B-CMP99b}, \cite{B-CMP102b} and \cite{B-CMP116},
or are introduced in this chapter, in the notation fixed in the
conventions~\ref{4.5:not-geometry-conventions} and in the lemmas of this section, together
with a last group of constants that depend on $(d,G,L)$. These are:
  \begin{itemize}
  \item the block size $\mathfrak{M}$ of the propagator papers and the gap number $R_1$;
  \item $\delta_0$, $\delta_1$, $B_0$, $B_0'$, $B_3$ and $C_0,\dots,C_4$, where $C_1$ is the
  constant that enters the definition of $\bar\varepsilon_1$ in
  \eqref{4.5:eq-geometry-conventions-a};
  \item $C_7$, $c_Q$, $C_{98}$, $C_{\mathrm{BCH}}$ and $\alpha_1^{(98)}$;
  \item $c_1$, the constant of \cite[Lemma~2.1]{B-CMP96} at the least fraction at which that
  lemma is read;
  \item the constant $C_H$ of the bound on the background field;
  \item $a_0^{[15]}$, $a_1^{[15]}$ and $a_3,\dots,a_6$;
  \item the numbers $\bar\varepsilon_1$ and $\mathfrak{r}_g$ of
  \eqref{4.5:eq-geometry-conventions-a}, where $C_1$ is chosen before $C_H$, and $C_H$ before
  $\bar\varepsilon_1$, and where $c_1$, $c_Q$ and $\alpha_1^{(98)}$ are chosen before
  $\bar\varepsilon_1$;
  \item the radii $a_0$, $a_1$, $a_\gamma$, $a_3$, $\bar\varepsilon_4$ and $c_9$;
  \item $C_{IV}$, $C_V$, $C_v$, $C_N$, $c_*$, $c_{\mathbb{K}}$, $C_F$, $C_F'$ and $C_{14}$;
  \item the constants $C_G$, $\rho_0$, $c_{\mathbf{1}}$ and $K_{\mathbf{1}}$ of the window
  regularity estimates. These depend on $(d,G,L)$ and do not read $M$. Of these, $\rho_0$
  enters only through the member $\rho_0/8$ of the bound $\delta_f\le\min\{1/9,\rho_0/8\}$ in
  (cf1), and that member is not used.
  \end{itemize}
\item The block size $M_1$ of \cite{B-CMP119}. It is a power of $L$ not smaller than
$R'_\star\mathfrak{M}$, where $R'_\star\ge\max\{R_1,8\}$ is the gap multiplier, a power of $L$,
of the window regularity estimates.
\item $\kappa$, then $\kappa_1$.
\item The enlargement parameter $\beta$ of \cite{B-CMP109}, a fixed positive number
$\beta\le 1$ through which the enlarged amplitudes $(1+\beta)\alpha$ and $(1+2\beta)\alpha$ are
formed, and the constant $B$ of \cite{B-CMP109}, which enters the ceiling
$c_{12}LMB\alpha_0'''\le a_0$ below.
\item The numbers
\begin{equation}\label{4.5:eq-constants-order-1}
B_b := 6dLc_1C_HC_V, \qquad B_b^{f} := 3dLc_1C_H\cdot 448d^2c_Q .
\end{equation}
Their products with $1+2\beta$ are read only after $\beta$.
\item The cube size $M = L^{m'}$, a power of $L$, subject to the floors
\begin{equation}\label{4.5:eq-constants-order-2}
\begin{gathered}
17M_1\le LM,\quad 17M_1\le c_{12}LM,\quad c_{12}L\ge 2,\quad \delta_0M\ge 32,\quad
2R_1M_1\le M,\\
\delta_1M\ge\kappa_1,\quad M>L^2M_1,\quad L\ge 27\ \text{odd},\quad L^2\ge 12/c_{12},\quad
L\ge 4d,\\
L\ge 2\bigl(1+O_{\mathrm{CL\text{-}U}}(1/24)\bigr).
\end{gathered}
\end{equation}
Here $O_{\mathrm{CL\text{-}U}}(x)$ denotes the supremum of the Baker--Campbell--Hausdorff
remainder over arguments of size $s$ with $dc_Qs\le x$; it is a function of $x$ alone and
reads no amplitude. The floor
$L\ge 2(1+O_{\mathrm{CL\text{-}U}}(1/24))$ serves the bound on the Landau-to-axial map $g$ in
Lemma~\ref{4.5:lem-literal-frame} and in Lemma~\ref{4.5:lem-fine-axial-frame}.
There are two further floors:
  \begin{itemize}
  \item the condition on $M_1$ imposed at the stage of $M_1$;
  \item the floor on $R_1\mathfrak{M}$ formed by \eqref{4.5:eq-kernel-algebra-a} together with
  \cite[(2.59)]{B-CMP96}, read at the fractions $\tfrac14$, $\tfrac18$, $\tfrac1{16}$ of
  $\delta_0$ and at the fraction $c_*$.
  \end{itemize}
\item $\alpha_1$, under the following ceilings:
  \begin{itemize}
  \item $\alpha_1'''\le a_1$;
  \item those of \eqref{4.5:eq-uniform-increment-b} and of \eqref{4.5:eq-glued-window-a};
  \item the first member of (c7$'$) and the ceiling (c8), both in
  \eqref{4.5:eq-literal-frame-a};
  \item the $\alpha_1$-halves of the sum ceilings (c7) of \eqref{4.5:eq-literal-frame-a} and
  (cf1), (cf5) of \eqref{4.5:eq-fine-axial-frame-a};
  \item the ceiling $O(1)M\alpha_1\le\beta$ of \cite{B-CMP109}.
  \end{itemize}
\item $\alpha_0$, under the following ceilings:
  \begin{itemize}
  \item the ceilings $c_{12}LMB\alpha_0'''\le a_0$ and $\alpha_0'''\le a_\gamma$ attached to
  the radii $a_0$ and $a_\gamma$, and the $\alpha_0$-ceilings of
  \eqref{4.5:eq-uniform-increment-b} and \eqref{4.5:eq-glued-window-a};
  \item (c7), which reads the already fixed $\alpha_1$;
  \item the second and third members of (c7$'$), and (c8);
  \item (cf1), (cf5) and (cf4), including the second member of (cf4);
  \item $\max\{b,b'\}<1/(4C_{\mathrm{BCH}})$.
  \end{itemize}
The third member of (c7$'$) and the second member of (cf4) read
$b/L\le\alpha_1^{(98)}$ and $b'/L\le\alpha_1^{(98)}$ respectively. In addition, $\alpha_0$
satisfies the two ceilings of \cite{B-CMP109}:
\begin{equation}\label{4.5:eq-constants-order-3}
\alpha_0<\frac{\alpha_1}{2B_3^2\,O(1)\,M},\qquad \bigl(4B_3^2\,O(1)\,M\bigr)^2\alpha_0\le\beta .
\end{equation}
\item $\alpha_2$, with
\begin{equation}\label{4.5:eq-constants-order-4}
\alpha_2\le\min\Bigl\{\frac{\alpha_1}{8},\ \frac{\alpha_0}{C_{IV}},\ \frac{1}{4C_F'},\
\frac{1}{4C_{\mathrm{BCH}}}-\max\{b,b'\}\Bigr\}.
\end{equation}
In \cite{B-CMP109} this radius is introduced as ``a constant $\alpha_2$, depending on
$\alpha_0$ and on some absolute constants''.
\item The fluctuation radius $\varepsilon_1$ of \cite{B-CMP109} and \cite{B-CMP116}. Its
ceilings are:
  \begin{itemize}
  \item the ceiling $\varepsilon_1\le\tfrac13\alpha_2$, from the bound on the fluctuation data;
  \item the ceiling on $\varepsilon_1$ imposed by the localisation estimate, in which the
  factor $e^{2c_{\mathbb{K}}\kappa_1}$ enters;
  \item the ceiling $3C_N^{s}\varepsilon_1'\le\alpha_2$ of
  Lemmas~\ref{4.5:lem-literal-frame} and~\ref{4.5:lem-fine-axial-frame}.
  \end{itemize}
\end{enumerate}
The Cauchy radius $r_{\square}(B')$ is a function of the datum, not a constant. Apart from the
first inequality of \eqref{4.5:eq-constants-order-3}, no condition compares $\alpha_0$ with
$\alpha_1$. The cube size $M$ is bounded from below only.
\end{remark}

\begin{proof}
We check that the order is consistent. This means that no letter reads a letter chosen after
it, and that at every stage the conditions imposed admit a strictly positive choice.

First, the definition of $\bar\varepsilon_1$ creates no cycle. By
\eqref{4.5:eq-geometry-conventions-a},
\begin{gather*}
\bar\varepsilon_1=\min\Bigl\{a_1^{[15]},\ a_6,\ \frac{R_1a_1^{[15]}}{6},\ \frac{1}{4C_1},\
\frac{2\mathfrak{r}_g}{3dLC_Hc_1C_1}\Bigr\},\\
\mathfrak{r}_g=\min\Bigl\{1,\ \frac{1}{24dc_Q},\ L\alpha_1^{(98)}\Bigr\}.
\end{gather*}
The number $\mathfrak{r}_g$ reads only $d$, $L$, $c_Q$ and $\alpha_1^{(98)}$, and all of these
are fixed at the stage of the constants depending on $(d,L)$, before $\bar\varepsilon_1$. The
number $\bar\varepsilon_1$ reads
$a_1^{[15]}$, $a_6$, $R_1$, $C_1$, $c_1$, $C_H$ and $\mathfrak{r}_g$. Of these, $C_H$ is chosen
after $C_1$. The bound that carries $C_H$ holds at every $\varepsilon_1^{[15]}\le a_6$. In that
bound the ratio $M\alpha_0/\varepsilon_1^{[15]}$ is $O(C_1B_3)$ by \cite[(28)]{B-CMP102b}, so
$C_H$ is free of $\varepsilon_1^{[15]}$ and of $M$. Hence $C_H$ reads neither
$\bar\varepsilon_1$, nor $M$, nor an
amplitude. The second entry of the minimum gives $\bar\varepsilon_1\le a_6$, so $C_H$ may be
read at $\varepsilon_1^{[15]}=\bar\varepsilon_1$. Every entry of both minima is strictly
positive. Therefore $\bar\varepsilon_1$ is a strictly positive number depending on $(d,L)$
only, fixed after every letter it reads.

The fifth entry has two roles. It is the radius used in the first step of the proof of
Lemma~\ref{4.5:lem-uniform-increment}. It also carries the two premises of the bound on the
Landau-to-axial map $g$ in Lemma~\ref{4.5:lem-literal-frame}.

Next, $B_b$ and $B_b^{f}$ in \eqref{4.5:eq-constants-order-1} are built from $d$, $L$, $c_1$,
$C_H$, $C_V$ and $c_Q$, all fixed at the stage of the constants depending on $(d,L)$. So they
depend on $(d,L)$ only. The letters
$b$ and $b'$ involve $1+2\beta$, and they are read only after $\beta$ is fixed.

The conditions of \eqref{4.5:eq-constants-order-2}, together with the two further floors
imposed at the stage of $M$, bound $M$, $R_1\mathfrak{M}$ or $L$ from below. They read only
letters fixed before $M$. None of them bounds any letter from above.

Each ceiling imposed at the stages of $\alpha_1$ and $\alpha_0$ either bounds the amplitude
being chosen from above by a strictly positive number built from letters fixed at earlier
stages, or bounds a sum from above, in the form
\begin{equation}\label{4.5:eq-constants-order-5}
c\,M(\alpha_0+c'\alpha_1)\le C,
\end{equation}
where $c,c'>0$ and $C>0$ are built from earlier letters. The ceilings (c7), (c8), (cf1) and
(cf5) are of the second form.

For a sum ceiling we proceed in two steps. At the stage of $\alpha_1$ we impose its
$\alpha_1$-half, $cc'M\alpha_1\le C/2$. At the stage of $\alpha_0$ we impose
$cM\alpha_0\le C/2$. Adding the two gives
\eqref{4.5:eq-constants-order-5}. For (c7), which reads $\alpha_1$, one may instead impose the
whole of \eqref{4.5:eq-constants-order-5} at the stage of $\alpha_0$ as a ceiling on
$\alpha_0$. This is possible because $\alpha_1$ is then already fixed.

We choose $\alpha_1$, and then $\alpha_0$, as one half of the minimum of the finitely many
strictly positive bounds imposed at its stage. Then every ceiling of that stage holds, the
strict ones strictly. In particular
the first inequality of \eqref{4.5:eq-constants-order-3} is one more upper bound on $\alpha_0$
after $\alpha_1$. It is therefore compatible with every other ceiling, all of which are upper
bounds. As stated in the remark, no condition besides this one compares $\alpha_0$ with
$\alpha_1$.

The condition $\max\{b,b'\}<1/(4C_{\mathrm{BCH}})$ is a ceiling on $\alpha_0$. Indeed,
$b$ and $b'$ equal $M\alpha_0'''$ times $B_b$ and $B_b^{f}$ respectively, and these factors are
fixed at the stage of $B_b$ and $B_b^{f}$. By the same reasoning, the conditions
$b/L\le\alpha_1^{(98)}$ and
$b'/L\le\alpha_1^{(98)}$ are ceilings on $\alpha_0$ read after $M/L$. Once $\alpha_0$ is fixed,
the fourth entry of \eqref{4.5:eq-constants-order-4} is therefore strictly positive. So are
the other three entries. Hence $\alpha_2$ may be chosen strictly positive, and it reads only
$\alpha_0$, $\alpha_1$ and earlier letters.

Finally, the ceilings imposed at the stage of $\varepsilon_1$ read $\alpha_2$,
$c_{\mathbb{K}}$, $\kappa_1$ and $C_N^{s}$,
all of which are fixed before. So $\varepsilon_1$ may be chosen strictly positive last.

The radius $r_{\square}(B')$ enters none of these choices. Every condition on $M$ is a floor,
so $M$ is bounded from below only.
\end{proof}

\begin{theorem}[Operator bounds at free current]\label{4.5:thm-operator-bounds}
Work in the setting of this section, with the reading conventions and with the constants chosen in the order of
Remark~\ref{4.5:rem-constants-order}. Assume the following statements:
\begin{itemize}
\item the bounds and the generalised random-walk expansion of the constrained covariance
inverse $\mathsf{q}_0=(QG_0Q^{*})^{-1}$ (Remark~\ref{4.5:rem-covariance-inverse-walks}),
which is stated in \cite[\S3]{B-CMP99b} without proof;
\item the analyticity and uniformity of the derivative bounds for the averaging maps
(Remark~\ref{4.5:rem-averaging-analyticity}), which is stated in \cite[p.~43]{B-CMP98}
without proof;
\item the corollary that furnishes the constant $C_H$ of
Remark~\ref{4.5:rem-constants-order}, together with the two auxiliary statements
accompanying it;
\item the regularity theorem for the glued window background $\mathbf{U}^{\mathrm{gl}}$;
\item the identity $U_j(\square_0,M^{\cdot}U_{k+1})=U_{k+1}$ at level $j$, the level-$j$
counterpart of Lemma~\ref{4.5:lem-window-minimiser};
\item the bound on the Baker--Campbell--Hausdorff remainder that controls the gauge maps
$g$ and $g'$ of the frames, whose statement also stands among the displays
\eqref{4.5:eq-literal-frame-b};
\end{itemize}
and, for the literal-frame variant of \textup{(N-b)} only, the small curvature of block
averages at near-unitary complex data (Remark~\ref{4.5:rem-block-curvature-complex}), which
is stated in \cite[(3.26)]{B-CMP109} without proof. Then the following hold.

\smallskip
\noindent\textup{(O)} \emph{Operators at free current.} Let $(\mathbf{U},\mathbf{J})$
satisfy the conditions (i)--(iii) of \cite{B-CMP109} at any constants
$\alpha_0',\alpha_1'$ under the ceilings for which the pairs with (i)--(iii) lie in the class
$\mathfrak{S}$ of \eqref{4.5:eq-admissible-pairs-a} (in \cite{B-CMP109} these are constants
that are sufficiently small but much larger than $\alpha_0,\alpha_1$). Then the operators
\begin{equation*}
G,\quad G_1,\quad H,\quad H_1,\quad H_0,\quad \tilde G,\quad (QGQ^{*})^{-1},\quad
(QG_1Q^{*})^{-1},\quad \Delta^{(2)},
\end{equation*}
with $J_{k+1}$ replaced by $\mathbf{J}$ in all propagators and operators, exist as the
operators at free current of Lemma~\ref{4.5:lem-free-current-operators}, are analytic,
obey the bounds of \cite[Theorem~3.3]{B-CMP99b} listed in
\eqref{4.5:eq-free-current-operators-a} (those with sup-norm data and those involving
$\nabla^{*}$, but not the bound involving the Laplacian, which is kept for $G_0$, $H_0$ only)
and the bounds \cite[(3.132), (3.133), (3.137)]{B-CMP99b}, have the generalised random-walk
expansions of \eqref{4.5:eq-free-current-operators-a}, and their versions localised in the
decoupling parameters $s$ as in clause (K10) of Theorem~\ref{4.5:thm-chart-existence} obey
\cite[(1.11)]{B-CMP116} in these norms.

\smallskip
\noindent\textup{(F)} \emph{Far terms.} Let $X\in\mathbf{D}_j$ meet the complement of
$\tilde\square^{2}$, and let $(\mathbf{U},\mathbf{J})\in\mathcal{U}_X$, where
$\mathcal{U}_X$ denotes the set \cite[(3.16)]{B-CMP109},
\begin{equation*}
\begin{split}
\mathcal{U}_X:=U^{c}_{j}\bigl(X,\tfrac12\alpha_0,\tfrac12\alpha_1,\alpha_0\bigr)
\cap\bigl\{(\mathbf{U},\mathbf{J}):\ &\mathbf{U},\mathbf{J}\text{ satisfy (i)--(iii) for }
j=k+1\\
&\text{with the constants }(1+\beta)\alpha_0,\ (1+\beta)\alpha_1,\ \alpha_0\bigr\}.
\end{split}
\end{equation*}
Membership in $\mathcal{U}_X$ is understood representative-wise, in the sense of
the corresponding display: one representative modulo $G$-valued gauge maps
satisfies both members of the intersection, and the orbit is recovered by \textup{(F4)}
below. Put $J:=\mathfrak{s}\mathbf{J}$ with $\mathfrak{s}=(L^{j-1}\eta)^{3}\le L^{-3}$.
Take the chart of Theorem~\ref{4.5:thm-chart-existence} at level $j$, in units of $\xi$ and
with the trivial sequence of domains, and let $B(t)$ and $v:=\partial_{t_\square}B(t)$ be
as in \cite[(3.6)--(3.8)]{B-CMP109}, taken as functions of $H_k:=\mathcal{H}^{J}$ at level
$k$. Replace the functional of \cite[(3.13)]{B-CMP109} by
\begin{equation*}
\Phi^{\mathbb{K}}:=\mathbf{E}^{(j)}\bigl(X,\ e^{i\xi A(t_\square)}\mathbf{U},\
J+\mathfrak{J}(B(t),t_\square)\bigr),
\end{equation*}
where $A(t_\square)$ is the field of \cite[(3.13)]{B-CMP109}, now formed from
$\mathcal{H}^{J}$, and $\mathfrak{J}(B(t),t_\square)$ is the current correction
$\mathfrak{J}$ of the chart formed in the same way;
replace the conditions of \cite[(3.14)]{B-CMP109} by
$|A|,|P_1A|,|\nabla_{\mathbf{U}}A|,|\mathfrak{J}|<\alpha_2$ on $X$, and define the radius
$r$ by
\begin{equation*}
r\cdot\max\bigl\{|\delta\mathcal{H}[v]|_X,\ |P_1\delta\mathcal{H}[v]|_X,\
|\nabla\delta\mathcal{H}[v]|_X,\ |\delta\mathfrak{J}[v]|_X\bigr\}=\tfrac13\alpha_2 ,
\end{equation*}
where $\delta\mathcal{H}[v]$, $\delta\mathfrak{J}[v]$ are the first variations of clause (K4)
of Theorem~\ref{4.5:thm-chart-existence}, displayed in
\eqref{4.5:eq-chart-variation-decay-a}.
If $\varepsilon_1$ is sufficiently small, then:
\begin{itemize}
\item[(F1)] $\Phi^{\mathbb{K}}$ is analytic on $\mathcal{U}_X\times\{|t_\square|\le r\}$,
and $|\Phi^{\mathbb{K}}|\le E_0e^{-\kappa d_j(X)}$.
\item[(F2)] (The counterpart of \cite[(3.9)]{B-CMP109}.) On $X$ each of the four quantities
in the definition of $r$ is bounded by
\begin{equation*}
C L^{j}\eta\, e^{-c_{*}\delta_0\operatorname{dist}^{(\xi)}(X,\operatorname{supp}\zeta_\square)}
\,|\zeta_\square H_k(B')| ,
\end{equation*}
whence \cite[(3.17)]{B-CMP109} holds at the rate $c_{*}\delta_0$.
\item[(F3)] At $(\mathbf{U},\mathbf{J})=(U_{k+1},J_{k+1})$ one has
$\Phi^{\mathbb{K}}=\mathbf{E}^{(j)}(X,e^{i\xi A(t_\square)}U_{k+1})$, which is the composite
function of \cite[(3.7)]{B-CMP109}.
\item[(F4)] (The invariance \cite[(1.19)]{B-CMP109}.) $\Phi^{\mathbb{K}}$ is invariant under
$(\mathbf{U},\mathbf{J},B)\mapsto(\mathbf{U}^{u},R(u)\mathbf{J},R(\bar u)B)$ for every
$G$-valued $u$ and for every $u=u'u_G$ with $u_G$ $G$-valued and
$\|\mathrm{Ad}_{u'}^{\pm1}\|\le\mathfrak{K}$, as in clause (K5) of
Theorem~\ref{4.5:thm-chart-existence}. This is the scope of the restriction statement for
complexified gauge maps near $G$-valued ones (maps $v=v'v_G$ with $|v'-1|$ small), and it is
the scope used in \cite[(3.29)]{B-CMP109}, namely $G^{c}$-valued transformations in a
sufficiently small neighbourhood of $G$-valued transformations. The invariance is used
only within this scope. The
statements that use \textup{(F)} --- \cite[(3.17)]{B-CMP109}, the inclusion of analyticity
spaces of \cite{B-CMP116}, and the real integrals of the step --- use representatives.
\item[(F5)] The same statements hold for the functions localised in the decoupling
parameters, with constants multiplied by $e^{c_{\mathbb{K}}\kappa_1}$.
\item[(F6)] The same statements hold for the two further far terms that \cite{B-CMP109}
treats in exactly the same way --- the second term of \cite[(3.21)]{B-CMP109} and the
bracket removed after \cite[(3.22)]{B-CMP109} --- for $X\subset\tilde\square^{2}$ with
$\operatorname{dist}^{(\xi)}(X,\operatorname{supp}(1-\tilde\zeta_\square))\ge
M(L^{j}\eta)^{-1}$; for these terms the chart is the window chart $H_j(\square_0,\cdot)$ on
the tower $\{\Omega_n\}:=\{\square_n\}_{n\le j}$.
\end{itemize}

\smallskip
\noindent\textup{(N)} \emph{Near terms.} Let $(\mathbf{U},\mathbf{J})\in\mathfrak{U}$, the
space of Proposition~\ref{4.5:prop-glued-window}.
\begin{itemize}
\item[(N-a)] $H_k(B')$ at the background $(\mathbf{U}^{\mathrm{gl}},\mathbf{J})$ is defined
and analytic in $(\mathbf{U},\mathbf{J},B')$ for $|B'|<c_9/8$, with
$N_\beta(H_k(B'))\le 16C^{K}_{9}|B'|$; \cite[(1.21)]{B-CMP116} holds for its version
localised in the decoupling parameters, with the factor $e^{c_{\mathbb{K}}\kappa_1}$; and it
depends on $(\mathbf{U},\mathbf{J})$ only through their restrictions to
$Y_0\cup\square_0$.
\item[(N-b)] There is a map $\mathfrak{w}=\mathfrak{w}[\mathbf{U}{\upharpoonright}\square_0^{+}]$
on the sites of $\square_0^{+}$, an explicitly given analytic function of the restriction of
$\mathbf{U}$ to $\square_0^{+}$, with $\|\mathrm{Ad}_{\mathfrak{w}}^{\pm1}\|\le2$, such that
on the bonds of $\square_0$ the background is $(\mathbf{U}^{\mathrm{gl}})^{\mathfrak{w}}
=e^{i\eta\mathcal{A}}$. Here $\mathfrak{w}:=w'$ is the frame map of
Lemma~\ref{4.5:lem-fine-axial-frame} and $\mathfrak{b}:=B_b^{f}(1+2\beta)$; in the
literal-frame variant, $\mathfrak{w}:=w=g^{-1}\circ\tilde v$ is the frame map of
Lemma~\ref{4.5:lem-literal-frame}, which is $u_{k+1}^{-1}v^{-1}$ in the notation of
\cite[(3.26)--(3.27)]{B-CMP109}, and $\mathfrak{b}:=B_b(1+2\beta)$. The data of the window
in this frame are $e^{iB}$ with $\sup|B|\le\frac12C_1\bar\varepsilon_1$ at base one. Put
\begin{equation*}
\mathbf{A}:=(t\tilde\zeta_\square+t_\square\zeta_\square)\cdot R(\mathfrak{w})H_k(B'),
\qquad
\hat{\mathbf{A}}:=\frac{1}{i\eta}\log\bigl(e^{i\eta\mathbf{A}}e^{i\eta\mathcal{A}}\bigr);
\end{equation*}
$\mathbf{A}$ is the field $(t\tilde\zeta_\square+t_\square\zeta_\square)\mathbf{H}_k(B')$
of \cite[(3.28)]{B-CMP109} in which the variables $\mathbf{J}$, $B$ are replaced by
$R(u_{k+1}^{-1})R(v^{-1})\mathbf{J}$, $R(u_{k+1}^{-1})R(v^{-1})B$, and $\hat{\mathbf{A}}$ is
the logarithm of \cite[(3.30)]{B-CMP109}. Then
\begin{equation}\label{4.5:eq-operator-bounds-a}
\begin{gathered}
|\mathcal{A}|,\ |\nabla^{\eta}\mathcal{A}|,\ \|\mathcal{A}\|_{1,\beta}\le
\mathfrak{b}M\alpha_0\quad\text{on }\tilde\square^{4},\\
|\mathbf{A}|,\ |\nabla^{\eta}\mathbf{A}|,\ \|\mathbf{A}\|_{1,\beta}<\alpha_2
\quad\text{on }\square_0,\\
|\hat{\mathbf{A}}|,\ |\nabla^{\eta}\hat{\mathbf{A}}|,\ \|\hat{\mathbf{A}}\|_{1,\beta}
<2(\alpha_2+\mathfrak{b}M\alpha_0)\quad\text{on }\tilde\square^{4}.
\end{gathered}
\end{equation}
The first line, the counterpart of \cite[(3.27)]{B-CMP109}, involves $\alpha_0$ alone (not
$\alpha_1$) and the first power of $M$. The second line, the counterpart of
\cite[(3.31)]{B-CMP109}, is stated with the plain lattice derivative $\nabla^{\eta}$ and holds
for $|t|\le1$, $|t_\square|\le r_\square(B')$ under $3C_N^{s}\varepsilon_1'\le\alpha_2$, both
at $s\equiv1$ and for the localised chart $H_k(s(Y_0),B')$ of \cite{B-CMP116}. The third
line is the counterpart of \cite[(3.32)]{B-CMP109}.
\item[(N$'$)] At $(\mathbf{U},\mathbf{J})=(U_{k+1},J_{k+1})$ one has
$\mathbf{U}^{\mathrm{gl}}=U_{k+1}$, and $H_k(B')$ is the function $\mathcal{H}(B')$ of
\cite{B-CMP102b}.
\end{itemize}
\end{theorem}

\begin{proof}
\emph{Clause} (O). The operators listed are those of
Lemma~\ref{4.5:lem-free-current-operators}, formed from Ba{\l}aban's operators by the
replacement $J_{k+1}\mapsto\mathbf{J}$ in all propagators and operators, as defined before
that lemma. That lemma gives their existence, their analyticity, the bounds of
\cite[Theorem~3.3]{B-CMP99b} and \cite[(3.132), (3.133), (3.137)]{B-CMP99b} in the form
listed in \eqref{4.5:eq-free-current-operators-a}, and their generalised random-walk
expansions. The bounds \cite[(1.11)]{B-CMP116} for their versions localised in the
decoupling parameters are clause (K10) of Theorem~\ref{4.5:thm-chart-existence}.

\emph{Clause} (F1). The second member of the intersection defining $\mathcal{U}_X$,
together with the ceilings under which pairs satisfying (i)--(iii) lie in the class
$\mathfrak{S}$ of \eqref{4.5:eq-admissible-pairs-a}, places the representative
$(\mathbf{U},J)$ in $\mathfrak{S}$ at the levels $j$ and $k$; for the current this uses
\begin{equation*}
|J|_{(-3)}=\mathfrak{s}|\mathbf{J}|\le L^{-3}\alpha_0<a_\gamma .
\end{equation*}
Hence the chart of Theorem~\ref{4.5:thm-chart-existence} is available at this pair. By the
linear bounds, clause (K3) of Theorem~\ref{4.5:thm-chart-existence}, for $\varepsilon_1$
sufficiently small the functions $\mathcal{H}(B(t))$ and $\mathfrak{J}(B(t))$ satisfy the
conditions that replace \cite[(3.14)]{B-CMP109} with $\frac13\alpha_2$ in place of
$\alpha_2$; this is the smallness of $\varepsilon_1$ under which, in \cite{B-CMP109},
$\mathbf{H}_j(B(t))$ satisfies \cite[(3.14)]{B-CMP109} with $\frac13\alpha_2$. By the
definition of $r$, on $|t_\square|\le r$ the first-order
term in $t_\square$ contributes at most $\frac13\alpha_2$ to each quantity, so that
$|A(t_\square)|,\ |\nabla A(t_\square)|\le\frac23\alpha_2$. For the current correction,
the difference $\mathfrak{J}(B,t_\square)-\mathfrak{J}(B)-t_\square\,\delta\mathfrak{J}[v]$
is the second-order Taylor remainder of the map $F$ of the chart (clause (K4) of
Theorem~\ref{4.5:thm-chart-existence}), bounded by $C_F'\alpha_2^{2}$, where $C_F'$ is the
second-derivative constant of $F$; therefore
\begin{equation*}
|\mathfrak{J}(B,t_\square)|\le\tfrac23\alpha_2+C_F'\alpha_2^{2}
\le\tfrac23\alpha_2+\tfrac14\alpha_2<\alpha_2
\end{equation*}
under the ceiling $\alpha_2\le1/(4C_F')$.

By the first member of the intersection, the representative $\mathbf{U}$ satisfies the
conditions (i)--(iv) at the constants $(\frac12\alpha_0,\frac12\alpha_1)$. Apply
Lemma~\ref{4.5:lem-uniform-increment} with $A:=A(t_\square)$, whose amplitude is
$\frac23\alpha_2\le\alpha_2$: it gives, for $e^{i\xi A(t_\square)}\mathbf{U}$, the conditions
(i), (ii), the first member of \cite[(1.14)]{B-CMP109}, and condition (iv) in the uniform form
\eqref{4.5:eq-uniform-increment-a}, at the constants $(\alpha_0,\alpha_1)$. For the current,
\begin{equation*}
|J+\mathfrak{J}|\le L^{-3}\alpha_0+\alpha_2<\alpha_0 ,
\end{equation*}
which is the second member of \cite[(1.14)]{B-CMP109}, whose constant there is $\alpha_0$.
Hence the orbit of the produced pair lies in $U^{c}_{j}(X,\alpha_0,\alpha_1)$. On this space
$\mathbf{E}^{(j)}$ is defined and analytic, the constants $\alpha_0,\alpha_1$ being absolute
(independent of $X$ and $j$), and it obeys the bound
\cite[(1.18)]{B-CMP109}, which gives $|\Phi^{\mathbb{K}}|\le E_0e^{-\kappa d_j(X)}$.
Composing with the joint analyticity of the chart, clause (K1) of
Theorem~\ref{4.5:thm-chart-existence}, gives the analyticity of $\Phi^{\mathbb{K}}$ on
$\mathcal{U}_X\times\{|t_\square|\le r\}$.

\emph{Clause} (F2). The direction $v=\partial_{t_\square}B(t)$ is supported within one
level-$j$ block of $\operatorname{supp}\zeta_\square$, and
$|v|\le c_QL^{j}\eta\,|\zeta_\square H_k(B')|$. The first variations
$\delta\mathcal{H}[v]$, $\delta\mathfrak{J}[v]$ are given by clause (K4) of
Theorem~\ref{4.5:thm-chart-existence}, displayed in \eqref{4.5:eq-chart-variation-decay-a},
and their exponential decay away from the support of $v$, at the rate $c_{*}\delta_0$ in
$\operatorname{dist}^{(\xi)}$, is clause (K8), displayed in
\eqref{4.5:eq-chart-variation-decay-b}. This bounds $|\delta\mathcal{H}[v]|_X$,
$|\nabla\delta\mathcal{H}[v]|_X$ and $|\delta\mathfrak{J}[v]|_X$ as asserted; since
$P_1\in\mathcal{K}_{-2}$, the same bound holds for $|P_1\delta\mathcal{H}[v]|_X$. Inserting
these bounds in the definition of $r$ gives \cite[(3.17)]{B-CMP109} at the rate
$c_{*}\delta_0$.

\emph{Clause} (F3) follows from the diagonal exactness clause (K6)(c) of
Theorem~\ref{4.5:thm-chart-existence}, displayed in \eqref{4.5:eq-diagonal-exactness-a},
combined with \cite[(1.9)]{B-CMP109}.

\emph{Clause} (F4) is the covariance clause (K5) of Theorem~\ref{4.5:thm-chart-existence}
applied with the maps stated.

\emph{Clause} (F5). By clause (K10) of Theorem~\ref{4.5:thm-chart-existence}, the assertion
of \cite{B-CMP109} that the new localised functions satisfy the bounds
\cite[(3.14), (3.9)]{B-CMP109} holds for
the list of conditions above, with constants multiplied by $e^{c_{\mathbb{K}}\kappa_1}$; the
proofs of (F1)--(F4) then apply to the localised functions.

\emph{Clause} (F6). Theorem~\ref{4.5:thm-chart-existence} is uniform in the sequence
$\{\Omega_j\}$. By the definition of $\mathfrak{S}$ in \eqref{4.5:eq-admissible-pairs-a},
$(\mathbf{U},\mathfrak{s}\mathbf{J}){\upharpoonright}\square_0\in\mathfrak{S}$ for the window
tower. Lemma~\ref{4.5:lem-uniform-increment} depends on the tower only through its clause for
the one-power hereditary class, which reads from a tower only its admissibility in the sense of
\cite[(1)]{B-CMP102b} and the collar bound $\le\frac12M$. Here, as in (F1), that lemma is
applied to the maximal tower of $X$, which is condition (iv) of the first member of
$\mathcal{U}_X$; the window tower $\{\square_n\}$ is covered by the same clause wherever
condition (iv) is established for it. The diagonal exactness (K6) holds at
$U_j(\square_0,M^{\cdot}U_{k+1})=U_{k+1}$: this is the identity at level $j$ assumed in the
hypotheses, whose proof is that of Lemma~\ref{4.5:lem-window-minimiser} at level $j$. With
these three facts the proofs of (F1)--(F5) apply word for word to the two further far terms.

\emph{Clause} (N-a). Proposition~\ref{4.5:prop-glued-window} places
$(\mathbf{U}^{\mathrm{gl}},\mathbf{J})$ in the orbit class $\mathfrak{S}^{\mathrm{orb}}$ of
\eqref{4.5:eq-admissible-pairs-a}, through the gauge map $h$ of
\eqref{4.5:eq-glued-window-a}. By (K1), (K3) and the covariance clause (K5), applied with
$\mathfrak{K}=2$ for $h$,
\begin{equation*}
H_k(B';\mathbf{U}^{\mathrm{gl}},\mathbf{J})
=R(h^{-1})\,H_k\bigl(R(\bar h)B';\,S^{h},\,\mathrm{Ad}_h\mathbf{J}\bigr),
\end{equation*}
and consequently
\begin{equation*}
N_\beta\bigl(H_k(B')\bigr)\le 2\cdot C^{K}_{9}\cdot2|B'|\cdot4=16C^{K}_{9}|B'| .
\end{equation*}
The bound
\cite[(1.21)]{B-CMP116} for the localised version is (K10). Each walk factor depends on
$\mathbf{U}^{\mathrm{gl}}$ only on its own domain, and $\mathbf{U}^{\mathrm{gl}}$ depends on
$\mathbf{U}$ only on $\square_0$; this gives the locality statement. Analyticity in
$\mathbf{U}$ follows from part (a) of the regularity theorem for the glued window background
and composition.

\emph{Clause} (N-b) is Lemma~\ref{4.5:lem-fine-axial-frame}, whose conclusions are displayed
in \eqref{4.5:eq-fine-axial-frame-b}. The literal-frame variant is
Lemma~\ref{4.5:lem-literal-frame}, whose conclusions are displayed in
\eqref{4.5:eq-literal-frame-b}, under the additional hypothesis of
Remark~\ref{4.5:rem-block-curvature-complex}, which is stated in \cite[(3.26)]{B-CMP109}
without proof.

\emph{Clause} (N$'$) follows from Lemma~\ref{4.5:lem-window-minimiser} and the diagonal
exactness clause (K6)(a), displayed in \eqref{4.5:eq-diagonal-exactness-a}.
\end{proof}

\begin{remark}[The one-sided size conditions and their joint satisfiability]
\label{4.5:rem-size-conditions}
The constants, the cube side $M$ and the amplitudes of this section are chosen in the order of
Remark~\ref{4.5:rem-constants-order}. We write $\alpha''' := (1+2\beta)\alpha$ for
$\alpha = \alpha_0, \alpha_1$, where $\beta$ is the parameter of \cite{B-CMP109}, the schedule
constant of Definition~\ref{4.4:def-post-step-space}. We write $B^{109}$ for Ba{\l}aban's constant
$B$ of \cite{B-CMP109}, as in Definition~\ref{4.4:def-radius-ceilings}. We write $a_1^{[15]}$ for the
constant $a_1$ of \cite[Theorem~1]{B-CMP102b} and $\varepsilon_1^{[15]}$ for the parameter
$\varepsilon_1$ of \cite[(7)]{B-CMP102b}. We also
write $b := B_bM\alpha_0'''$ and $b' := B_b^{f}M\alpha_0'''$. The conditions imposed in this section
are of two kinds, floors and ceilings, and there are no caps.

\emph{Floors.} These are lower bounds on $M$, on $R_1\mathfrak{M}$ and on $L$, each imposed after
the constants it reads:
\begin{enumerate}
\item $17M_1 \le LM$ and $17M_1 \le c_{12}LM$; both follow from $M > L^2M_1$, $L \ge 17$ and
$c_{12} \ge 1$;
\item $c_{12}L \ge 2$;
\item the kernel condition \eqref{4.5:eq-kernel-algebra-a} together with \cite[(2.59)]{B-CMP96},
imposed on $R_1\mathfrak{M}$ after $L$. It is imposed at those fractions $\tfrac14$, $\tfrac18$,
$\tfrac1{16}$ of $\delta_0$ that are used in this section and at the rate fraction $c_{*}$ of
Theorem~\ref{4.5:thm-chart-existence}, where $\delta_0$ is Ba{\l}aban's constant fixed in
Remark~\ref{4.5:rem-constants-order}; it counts as one floor;
\item $\delta_0M \ge 32$ and $2R_1M_1 \le M$;
\item $\delta_1M \ge \kappa_1$;
\item the floors of the rim-regularity theorem: $L \ge 27$ with $L$ odd, the floor on $M_1$
imposed there, and $M > L^2M_1$;
\item $L^2 \ge 12/c_{12}$;
\item $L \ge 4d$, together with $R_1 \ge L$;
\item $L \ge 2(1 + \mathsf{c}_g)$, where $\mathsf{c}_g$ is the supremum of the
Baker--Campbell--Hausdorff remainder over arguments with $dc_Qb \le 1/24$, an absolute number.
This floor is required for the bound on the frame map $g$ in clause (AX3) of
Lemma~\ref{4.5:lem-literal-frame} and in clause (AXF3) of Lemma~\ref{4.5:lem-fine-axial-frame};
\item $M$ lies above the absolute floor that makes
$|\nabla\zeta_{\square}| \le 1$ and $|\nabla\tilde\zeta_{\square}| \le 1$.
\end{enumerate}
By its definition in \cite{B-CMP116},
``The number $R_1$ is a power of $L$ satisfying the condition $R_1M_1M^{-1} \leqq 1$''.
Now $M$ is a power of $L$ with $M \ge 2R_1M_1$. Since $R_1$ is a power of $L$ by this definition
and $M_1$ is a power of $L$ by its choice in Remark~\ref{4.5:rem-constants-order}, the product
$R_1M_1$ is a power of $L$; so $M$, a power of $L$ at least $R_1M_1$, is a multiple of $R_1M_1$,
and hence of its divisor $R_1\mathfrak{M}$. This meets the requirement of \cite{B-CMP102b} that
``$M$ is a multiple of $R_1M_1$'' (in the notation of \cite{B-CMP102b}, whose $M_1$ is the block
written $\mathfrak{M}$ here).

The two floors $\delta_0M \ge 32$ and $2R_1M_1 \le M$ are those of
Lemma~\ref{4.5:lem-uniform-increment}, where they are stated for that lemma's cube side. Where the
lemma is applied with cube side $M/L$, they read $\delta_0M \ge 32L$ and $2LR_1M_1 \le M$. These are
floors of that application, and no new floors are added here.

\emph{Ceilings.} These are radii. Each bounds one amplitude, or a sum of the two amplitudes, from
above, and each is imposed after the constants it reads.
\begin{enumerate}
\item The radii $a_0, a_1, a_\gamma, a_3, \bar\varepsilon_4, c_9$ of
Theorem~\ref{4.5:thm-chart-existence} satisfy
\begin{equation*}
\mathfrak{K}^2\cdot 4C_fc_9 \le \bar\varepsilon_4, \qquad
B_0c_9 \le \bar\varepsilon_4 \le a_3/4 ,
\end{equation*}
and $e^{2c_{\mathbb{K}}\kappa_1}$ times the size of the data to which
Theorem~\ref{4.5:thm-chart-existence} is applied is at most $c_9$.
Here $\mathfrak{K} = 4$, $C_f$ is the constant of Theorem~\ref{4.5:thm-chart-existence}, and $B_0$
is Ba{\l}aban's constant fixed in Remark~\ref{4.5:rem-constants-order}; on the orbit class the
radius in the variable $B$ of Theorem~\ref{4.5:thm-chart-existence} is $c_9/\mathfrak{K}$.
\item The three ceilings against the radii $a_0$, $a_1$, $a_\gamma$ of
Theorem~\ref{4.5:thm-chart-existence}; the third radius $a_\gamma$ bounds a third amplitude
$\gamma_0$, besides $\alpha_0$ and $\alpha_1$, read by that theorem; here $\gamma_0$ is taken equal
to $\alpha_0$:
\begin{equation*}
c_{12}LMB^{109}\alpha_0 \le a_0, \qquad \alpha_1 \le a_1, \qquad \gamma_0 = \alpha_0 \le a_\gamma .
\end{equation*}
Each is imposed on the amplitudes in the form in which they occur at its point of use in
Theorem~\ref{4.5:thm-chart-existence}, that is, on the primed amplitudes of that theorem or on
$\alpha'''$.
\item The ceilings of Lemma~\ref{4.5:lem-uniform-increment}, as in
\eqref{4.5:eq-uniform-increment-b}:
\begin{gather*}
C_7\alpha_0 \le \bar\varepsilon_1, \qquad
c_Q(\alpha_1+\alpha_2) \le \tfrac12 C_1\bar\varepsilon_1, \qquad
c_{12}LMB^{109}\cdot\tfrac12\alpha_0 + \tfrac12\alpha_1 + \alpha_2 \le c_{\sharp},\\
\alpha_1 \le 1, \qquad
\alpha_2 \le \min\{\alpha_1/8,\ \alpha_0/C_{IV}\}.
\end{gather*}
\item The ceiling \eqref{4.5:eq-glued-window-a} of Proposition~\ref{4.5:prop-glued-window}.
\item The ceilings of Lemma~\ref{4.5:lem-literal-frame}, as in \eqref{4.5:eq-literal-frame-a}:
\begin{align*}
&\text{(c7)}\quad 4C_VM(\alpha_0'''+\alpha_1''') \le C_1\bar\varepsilon_1;\\
&\text{(c7$'$)}\quad 2C_vM\alpha_1''' \le \tfrac16,\quad
4dLc_QB_bM\alpha_0''' \le \tfrac16,\quad
B_b(M/L)\alpha_0''' \le \alpha_1^{(98)};\\
&\text{(c8)}\quad c_Q\alpha_1''' \le \tfrac12 C_1\bar\varepsilon_1,\quad
C_7(\alpha_0'''+C_{14}\alpha_1''') \le \bar\varepsilon_1,\quad
\alpha_1''' \le \min\{\alpha_1^{(98)},1\}.
\end{align*}
The third member of (c7$'$) is $b/L \le \alpha_1^{(98)}$.
\item The ceilings of Lemma~\ref{4.5:lem-fine-axial-frame}, as in
\eqref{4.5:eq-fine-axial-frame-a}:
\begin{align*}
&\text{(cf1)}\quad 14d(M/L)\bigl(\alpha_1'''+c_{12}LMB^{109}\alpha_0'''\bigr) \le \tfrac1{10},\qquad
\delta_f \le \min\{\tfrac19,\ \rho_0/8\};\\
&\text{(cf5)}\quad 56d(M/L)(\alpha_0'''+C_{14}\alpha_1''') \le \min\{\alpha_1^{(98)},\tfrac14\},\\
&\phantom{\text{(cf5)}\quad} 448d^2c_QM(\alpha_0'''+C_{14}\alpha_1''')
\le \tfrac12 C_1\bar\varepsilon_1;\\
&\text{(cf4)}\quad 4dLc_QB_b^{f}M\alpha_0''' \le \tfrac16,\qquad
B_b^{f}(M/L)\alpha_0''' \le \alpha_1^{(98)} .
\end{align*}
Here $\rho_0$ is the radius fixed in the rim-regularity theorem. The entry $\rho_0/8$ of (cf1) is
idle: no estimate reads it. The second member of (cf4) is $b'/L \le \alpha_1^{(98)}$.
\item The condition $\beta_0 + \varepsilon < 1$ of clause (J1), the first clause of the statement
in which the constant $\beta_0$ is introduced; $\beta_0$ is fixed before this condition is imposed,
and $\varepsilon := \tfrac12(1-\beta_0)$. Since $\beta_0+\varepsilon =
\tfrac12(1+\beta_0)$, it holds exactly when $\beta_0 < 1$.
\item The ceilings $\alpha_2 \le 1/(4C_F')$, with $C_F'$ the constant of
Remark~\ref{4.5:rem-constants-order}, and
$\alpha_2 + \max\{b,b'\} \le 1/(4C_{\mathrm{BCH}})$.
\item The ceilings on the fluctuation radius. The first, from clause (F1) of the fluctuation-radius
lemma, which is applied at the first step of the proof of Lemma~\ref{4.5:lem-uniform-increment},
bounds it by
$\tfrac13\alpha_2$. The second is $3C_N^{s}\varepsilon_1' \le \alpha_2$, from clause (AX5) of
Lemma~\ref{4.5:lem-literal-frame} and clause (AXF5) of Lemma~\ref{4.5:lem-fine-axial-frame}. The
third is the ceiling of clause (K10) of Theorem~\ref{4.5:thm-chart-existence}.
\item The datum radius and the radius of the frame map $g$ are as in
\eqref{4.5:eq-geometry-conventions-a}:
\begin{gather*}
\bar\varepsilon_1 := \min\Bigl\{a_1^{[15]},\ a_6,\ \frac{R_1a_1^{[15]}}{6},\ \frac{1}{4C_1},\
\frac{2\mathfrak{r}_g}{3dLC_Hc_1C_1}\Bigr\},\\
\mathfrak{r}_g := \min\bigl\{1,\ 1/(24dc_Q),\ L\alpha_1^{(98)}\bigr\}.
\end{gather*}
Both are $(d,L)$-numbers, and $\mathfrak{r}_g$ contains the radius of the twist condition on the
background $U_0$ in clause (K6) of Theorem~\ref{4.5:thm-chart-existence}. The fifth
entry of $\bar\varepsilon_1$ is
the radius of the fluctuation-radius lemma required at the first step of the proof of
Lemma~\ref{4.5:lem-uniform-increment}, together with the two premises on the frame map $g$.
\end{enumerate}
None of these ceilings orders $\alpha_0$ against $\alpha_1$.

Ba{\l}aban's own conditions in \cite{B-CMP109} are
``$B_3^2O(1)M\alpha_0 < \tfrac12\alpha_1$'',
``$O(1)M\alpha_1 \leqq \beta$'' and
``$(4B_3^2O(1)M)^2\alpha_0 \leqq \beta$''.
We display them only for the order they impose. The first is the only one among the conditions
listed here that orders $\alpha_0$ against $\alpha_1$.

\emph{Caps.} There are none: no condition bounds $L$, $M$, $K$, $k$ or $N$ from above, and none
bounds the coupling from below.

Once $L$ and the constants that precede $M$ in the order of Remark~\ref{4.5:rem-constants-order}
have been fixed subject to the conditions above that concern them (the floors on $L$ and on
$R_1\mathfrak{M}$, and the ceilings of the first item among the radii of
Theorem~\ref{4.5:thm-chart-existence}), all floors on $M$ and all ceilings on the amplitudes can be
met simultaneously. The admissible cube sides are exactly the powers of $L$ that are at least a power
$M_\flat$ of $L$ determined by the floors, that is, the powers of $L$ in the half-line
$[M_\flat,\infty)$.
\end{remark}

\begin{proof}
\emph{The floors.} We first derive the first floor from its stated premises. Since $M > L^2M_1$ and
$L \ge 17$,
\begin{equation*}
LM > L^3M_1 \ge 17M_1 ,
\end{equation*}
and $c_{12} \ge 1$ then gives $c_{12}LM \ge LM > 17M_1$.

Every floor on $M$ bounds $M$ from below by a number fixed before $M$ in the order of
Remark~\ref{4.5:rem-constants-order}. No floor bounds $M$ from above. Hence if a power of $L$
satisfies all floors on $M$, so does every larger power of $L$.

Since $L \ge 27$, the powers of $L$ are unbounded. So there is a least power of $L$ that satisfies
every floor on $M$; call it $M_\flat$. The powers of $L$ satisfying all floors are then exactly the
powers of $L$ that are $\ge M_\flat$.

\emph{The ceilings.} The ceilings (c7), (c7$'$), (c8), (cf1), (cf4), (cf5), the ceiling
$3C_N^{s}\varepsilon_1' \le \alpha_2$ from clause (AX5), the ceiling
$\mathfrak{K}^2\cdot 4C_fc_9 \le \bar\varepsilon_4$ and the condition $\beta_0 + \varepsilon < 1$
have the form
\begin{equation*}
\varphi \cdot x \le c, \qquad c > 0 ,
\end{equation*}
where $\varphi$ depends only on quantities fixed earlier in the order of choice. The quantity $x$ is
one of $\alpha_0$, $\alpha_1$, a weighted sum of the two, $\varepsilon_1'$, $c_9$ or $\varepsilon$.
This is the form treated by the half-line lemma, which asserts that finitely many conditions of this
form, each imposed after the letters on which its $\varphi$ depends, are met by choosing each letter
in turn as the minimum of finitely many strictly positive numbers. The remaining ceilings, apart
from the two ceilings
of Lemma~\ref{4.5:lem-uniform-increment} that contain $\alpha_2$ in a sum (treated below), bound one
letter from above by numbers fixed before it, and are met below by taking minima.

Four points need a word of explanation.
\begin{itemize}
\item The third member of (c7$'$) and the second member of (cf4) are $b/L \le \alpha_1^{(98)}$ and
$b'/L \le \alpha_1^{(98)}$. Since $b = B_b(1+2\beta)M\alpha_0$ and
$b' = B_b^{f}(1+2\beta)M\alpha_0$, with $B_b$, $B_b^{f}$, $\beta$ and $M$ fixed before $\alpha_0$,
these are ceilings on $\alpha_0$, imposed after $M/L$.
\item The fifth entry of $\bar\varepsilon_1$ is a $(d,L)$-number, so every ceiling with
$\bar\varepsilon_1$ on its right-hand side has $c$ fixed before the amplitudes.
\item A ceiling on a sum, $c\cdot M(\alpha_0 + c'\alpha_1) \le C$, is met by halving. First fix
$\alpha_1$ so that $cc'M\alpha_1 \le C/2$. Then fix $\alpha_0$ so that $cM\alpha_0 \le C/2$.
\item The ceilings $c_Q(\alpha_1+\alpha_2) \le \tfrac12 C_1\bar\varepsilon_1$ and
$c_{12}LMB^{109}\cdot\tfrac12\alpha_0 + \tfrac12\alpha_1 + \alpha_2 \le c_{\sharp}$ contain
$\alpha_2$ in a sum. Since $\alpha_2 \le \alpha_1/8$ is among the ceilings of $\alpha_2$, they hold
once
\begin{equation*}
\tfrac98 c_Q\alpha_1 \le \tfrac12 C_1\bar\varepsilon_1
\quad\text{and}\quad
c_{12}LMB^{109}\cdot\tfrac12\alpha_0 + \tfrac58\alpha_1 \le c_{\sharp};
\end{equation*}
the first is a ceiling on $\alpha_1$, and the second is a ceiling on a sum, met by halving.
\end{itemize}

Now fix any power $M \ge M_\flat$ of $L$ and choose the amplitudes in the following order, each as
the minimum of finitely many strictly positive numbers.
\begin{enumerate}
\item Choose $\alpha_1$. Its ceilings are those bounding $\alpha_1$ alone together with the
$\alpha_1$-halves of every ceiling on a sum, namely (c7), the second member of (c8), (cf1), (cf5)
and the ceiling in $c_\sharp$ in the form just derived.
\item Choose $\alpha_0$. Its ceilings are those bounding $\alpha_0$ alone, the remaining halves of
the sum ceilings, and the ceilings on $\alpha_0$ after $M/L$ from the previous paragraph. We also
require $b \le 1/(8C_{\mathrm{BCH}})$ and $b' \le 1/(8C_{\mathrm{BCH}})$, which are ceilings on
$\alpha_0$ of the same form, so that $\max\{b,b'\} < 1/(4C_{\mathrm{BCH}})$.
Consequently the number $1/(4C_{\mathrm{BCH}}) - \max\{b,b'\}$ is strictly positive.
\item Choose $\alpha_2$. Take $\alpha_2$ at most
\begin{equation*}
\min\bigl\{\alpha_1/8,\ \alpha_0/C_{IV},\ 1/(4C_F'),\
1/(4C_{\mathrm{BCH}}) - \max\{b,b'\}\bigr\};
\end{equation*}
each entry is strictly positive.
\item Choose the fluctuation radius as the minimum of $\tfrac13\alpha_2$ and the bound of clause
(K10) of Theorem~\ref{4.5:thm-chart-existence}, and choose $\varepsilon_1'$ subject to
$3C_N^{s}\varepsilon_1' \le \alpha_2$.
\end{enumerate}
Every ceiling then holds by construction. This construction works for every power $M \ge M_\flat$
of $L$. Therefore the admissible $M$ form the half-line $[M_\flat,\infty)$ among the powers of $L$.

\emph{The order of the amplitudes.} The conditions of this section leave $\alpha_0$ and $\alpha_1$
unordered.
Ba{\l}aban's order $\alpha_0 < \alpha_1/(2B_3^2O(1)M)$ is equivalent to his first displayed
condition. It is compatible with the construction: it is a ceiling on $\alpha_0$ imposed after
$\alpha_1$, and half of the bound $\alpha_1/(2B_3^2O(1)M)$ enters the minimum defining $\alpha_0$,
so that the strict inequality holds. The same holds for his ceilings
$O(1)M\alpha_1 \leqq \beta$ on $\alpha_1$ and $(4B_3^2O(1)M)^2\alpha_0 \leqq \beta$ on $\alpha_0$.

\emph{The cube cap of \cite[Theorem~1]{B-CMP102b}.} That theorem requires its cube side, which we
write $M'$, to satisfy
\begin{equation*}
M' \le R_1\mathfrak{M}a_1^{[15]}/\varepsilon_1^{[15]} .
\end{equation*}
It is read only at $\varepsilon_1^{[15]} = \bar\varepsilon_1$, for the cubes of \cite{B-CMP99b} of
fixed side $12\mathfrak{M}$. There it is a ceiling on $\bar\varepsilon_1$ depending only on
$(d,L)$. It is never read for the cube side $M$ of \cite{B-CMP109}: at the opening step of the proof
of Lemma~\ref{4.5:lem-uniform-increment} the gauge on the averaging operator $Q$ is that of
\cite[(1.12)]{B-CMP109}, so \cite[Theorem~1]{B-CMP102b} is not invoked with cube side $M$. When the
response is taken at base one, in the sense of Lemma~\ref{4.5:lem-literal-frame}, the cap is void.

Since this bound is a decreasing function of $\varepsilon_1^{[15]}$, shrinking $\bar\varepsilon_1$
by its fifth entry only enlarges the bound at $\varepsilon_1^{[15]} =
\bar\varepsilon_1$. Hence no floor arises from this cap.

\emph{The one-power hereditary class.} The clause concerning the one-power hereditary class adds no
floor: its hypothesis that the collars of the tower $\mathbf{\Omega}$ are bounded by $\tfrac12M$ is
met by the maximal tower under the standing floor $2R_1M_1 \le M$.
\end{proof}

\begin{remark}[What the operator bounds at free current do not assert]\label{4.5:rem-not-asserted}
Theorem~\ref{4.5:thm-operator-bounds} asserts none of the following.
\begin{enumerate}
\item Bounds on $|\Delta_U\,\cdot\,|$ or $|D^{*}D\,\cdot\,|$, where $\Delta_U$ is the covariant
lattice Laplacian, $D=D_U$ the covariant lattice derivative and $D^{*}$ its adjoint, applied to
$\mathcal{H}^{J}$, to its variation $\delta\mathcal{H}^{J}$, or to a localised piece of it in the
decoupling parameters $s=\{s(\Delta)\}$ of \cite{B-CMP116}, at pairs $(\mathbf{U},\mathbf{J})$
with free current $\mathbf{J}$ off the diagonal; here the \emph{diagonal} consists of the pairs
$(U_{k+1},J_{k+1})$, in which the current is the one belonging to the configuration. These are
the members of the lists \cite[(3.14), (3.9)]{B-CMP109} as printed that involve these two
operators, and the fourth and fifth members of \cite[(190)]{B-CMP102b} at free current.
\item The bound \cite[(21)]{B-CMP102b} off the diagonal.
\item Any of the three hypotheses that Theorem~\ref{4.5:thm-operator-bounds} assumes beyond
the general conditions of this chapter and the statements it cites; the third of
them is assumed only for the variant of \eqref{4.5:eq-operator-bounds-a} in
the notation of \cite{B-CMP109}.
\item That the individual pieces of the expansion in $s$ coincide with the corresponding
pieces in Ba{\l}aban's construction. What does coincide is their sum at $s\equiv 1$ and
their values on the diagonal.
\item Any of the following as a statement in its own right. First, the statements that
Theorem~\ref{4.5:thm-operator-bounds} cites: the corollary on which it rests
together with its two supporting statements on the analysis of \cite{B-CMP102b}, the
regularity statement that supplies the axial representative on the cube $\square_0$ and its
gauge map, the lemma of which the theorem cites only the part stating that, at level $j$, the
background on $\square_0$ determined by the tower $M^{\cdot}(U_{k+1})$ of block averages is
$U_{k+1}$, and the bound on the Baker--Campbell--Hausdorff remainder; of the last, only the
three lines written out within the frame lemma in Ba{\l}aban's notation are asserted. Second,
the lemma on the cluster expansion of \cite{B-CMP116} in which the far functions are evaluated
on one representative modulo $G$-valued gauge maps, which is not a hypothesis of the theorem.
The identification of the glued background with $U_{k+1}$ at $(U_{k+1},J_{k+1})$ is part of
Theorem~\ref{4.5:thm-operator-bounds} and is asserted.
\item The covariance identity \cite[(181)]{B-CMP102b} for $G^{c}$-valued gauge maps,
beyond what the identity theorem for analytic functions gives at the points where it is
used. Those points lie in the two frame lemmas and in the lemma that carries the
regularity conditions over to the configuration shifted by the chart. In particular, the
proof of Theorem~\ref{4.5:thm-operator-bounds} does not use the passage of \cite{B-CMP102b}
in which the equality is taken as a definition for the remaining $v$.
\item The gauge invariance of the far functional, beyond the class of gauge maps for which
Theorem~\ref{4.5:thm-operator-bounds} states it.
\item The covariance \cite[(3.29)]{B-CMP109}, beyond the scope of the remark accompanying
the fine axial frame lemma.
\item The display \cite[(1.17)]{B-CMP109} with its arguments replaced by a general complex
pair $(\mathbf{U},\mathbf{J})$.
\item Any bound on $\nabla\mathcal{A}_f$, where $\mathcal{A}_f$ is one of the auxiliary
potentials of the fine axial frame lemma; it is distinct from the frame potential
$\mathcal{A}$ of \eqref{4.5:eq-operator-bounds-a}, whose bounds displayed there are asserted.
\item That the axial data $V$ of Ba{\l}aban's data-axial frame, in the notation of
\cite{B-CMP109}, have holonomies whose bound carries the first power of $M$ only, in the
absence of the third hypothesis named above. The frame of the fine axial frame lemma, in which
\eqref{4.5:eq-operator-bounds-a} is obtained without that hypothesis, is field-axial and
not data-axial; see the remark accompanying that lemma.
\item That the tower $\{\square_n\}$ of the cube $\square_0$ in \cite[(2.13)]{B-CMP119}
satisfies the collar bound required in the one-power hereditary class.
\end{enumerate}
\end{remark}

\begin{definition}[The complex domain on which the bounds hold]\label{4.6:not-complex-domain}
Let $k$ be a level and $X$ a region at level $k$. In \cite[p.~390]{B-CMP122b} the terms that the
operation $\mathbf{R}$ produces at step $k$ are stated to be defined, as analytic functions, on the
space $\tilde{U}^{c}_{k}(X,\tilde\alpha_0,\tilde\alpha_1)$. By \cite[p.~191]{B-CMP122a}, the member of
index $n=0$ of the stage family $\{\tilde{U}^{(n)c}_{k}(X,\tilde\alpha_0,\tilde\alpha_1)\}_{n\ge0}$ of
\cite[(1.64)--(1.69)]{B-CMP122a} coincides with $\tilde{U}^{c}_{k}(X,\tilde\alpha_0,\tilde\alpha_1)$. The
latter is the inductive multi-scale space of complex configurations $(\mathbf{U},\mathbf{J})$ on $X$ of
\cite[(2.34)--(2.39)]{B-CMP119}. Write $\Omega_j$, indexed by levels $j$ as in \cite[\S2]{B-CMP119},
for the domains of the inductive construction there. This space is graded over the shells
$(\Omega_j\setminus\Omega_{j+1})\cap X$, with the radii of \cite[(2.28)]{B-CMP119}, which for each
level $j\le k$ are
\begin{equation}\label{4.6:eq-complex-domain-1}
\alpha_{0,j}=g_jC_0\bigl(\log g_j^{-2}\bigr)^{q_0},\qquad
\alpha_{1,j}=g_jC_1\bigl(\log g_j^{-2}\bigr)^{q_1}.
\end{equation}
Here $q_0,q_1$ are the exponents fixed in \cite[(2.28)]{B-CMP119}. We read the tildes as marking the
multi-scale space and its sequences of radii:
\begin{equation*}
\tilde\alpha_0=\{\alpha_{0,j}\}_{j\le k},\qquad \tilde\alpha_1=\{\alpha_{1,j}\}_{j\le k}.
\end{equation*}
The single-scale space, carrying the radii of level $k$ alone, is written
$U^{c}_{k}(X,\alpha_{0,k},\alpha_{1,k})$.

In this chapter the symbol $\tilde{U}^{c}_{k}(X,\tilde\alpha_0,\tilde\alpha_1)$ is used for four
nested domains. The first, $U^{\mathrm{pr}}_k(X)$, is Ba{\l}aban's space as printed. The second is
the one-power hereditary
class $U^{\Sigma^1}_{k}(X;a_0,a_1,\gamma_0)$, with constants $a_0$, $a_1$, $\gamma_0$. This class is
the union of the orbits, under complexified gauge transformations, of the pairs
$(\mathbf{U},\mathbf{J})$ on $X$ that satisfy four conditions: condition (i) of \cite{B-CMP109},
namely \cite[(1.11)]{B-CMP109} at the constant $a_0$ and \cite[(1.12)]{B-CMP109}, imposed in the
reading fixed in Definition~\ref{4.4:def-local-data-classes}, with the constant $B^{109}$ of that
definition; condition (ii) of \cite{B-CMP109} at $a_1$;
condition (iii) of \cite{B-CMP109} at $(a_0,\gamma_0)$; and, for every pair $(i,Z)$ in the tower of
$X$, consisting of a level $i\le k$ and a region $Z$, condition (iv) of \cite{B-CMP109} relative to
$(i,Z)$, taken at the radius $a_0L^{-(k-i)}$. Under the class
convention attached to that class, every complex space $U^{c}_{\cdot}(Y,\dots)$, $Y$ a region, is
taken to be $U^{\Sigma^1}_{\cdot}(Y;\dots)$ with the same constants. The third, $U^{\mathrm{ad}}_k(X)$,
is the domain singled out by the adopted wording of the analyticity condition on the multi-scale
spaces, the wording in which the theorem on the multi-scale spaces states its objects. The fourth,
$U^{\mathrm{an}}_k(X)$, is the complex domain on which clause (i-a) of the statement on the
large-field part of the coupling increment asserts its terms analytic. By the inclusion statement
that accompanies the definition of the one-power hereditary class, these domains form a chain:
\begin{equation*}
U^{\mathrm{pr}}_k(X)\supset U^{\Sigma^1}_{k}(X;a_0,a_1,\gamma_0)\supset U^{\mathrm{ad}}_k(X)
\supseteq U^{\mathrm{an}}_k(X).
\end{equation*}
A statement proved on the one-power hereditary class therefore restricts to $U^{\mathrm{ad}}_k(X)$
and $U^{\mathrm{an}}_k(X)$, and not conversely.

We write
\begin{equation*}
\mathfrak{U}_k(X)=\tilde{U}^{c}_{k}(X,\tilde\alpha_0,\tilde\alpha_1),
\end{equation*}
taken to be the largest member of the chain on which all the inputs of this section hold; since a
sup-bound and uniform convergence on a member pass to every smaller member, as shown below, this is
the smallest of the members carrying the individual inputs. These inputs are three:
the family bound for the island operation $\mathbf{T}'_k$, the bound on the interaction terms over
the admissible interaction domains, and the analyticity statement for the large-field terms. Those
of them that concern the objects of the theorem on the multi-scale spaces hold on
$U^{\mathrm{ad}}_k(X)$ or on $U^{\mathrm{an}}_k(X)$; the others hold on the one-power hereditary
class, under its class convention.

Every statement proved in this section is a bound on the supremum over $\mathfrak{U}_k(X)$, together
with uniform convergence on $\mathfrak{U}_k(X)$. Both properties pass to any smaller domain and through
the restriction maps $\mathfrak{U}_k(X)\to\mathfrak{U}_k(X')$, $X'\subset X$. The arguments of this
section use only the sup-bounds and the uniform convergence of their inputs on $\mathfrak{U}_k(X)$,
never the defining conditions of these spaces or the values of the radii
\eqref{4.6:eq-complex-domain-1}. Consequently the results
hold on Ba{\l}aban's space as printed if the inputs hold there. Otherwise they hold on
$\mathfrak{U}_k(X)$ and, by the chain of inclusions, on every domain of the chain contained in it.
\end{definition}

\begin{proof}
Let $D'\subset D$ be two of the domains above, and let $F$ be a function on $D$ with
$\sup_{D}|F|\le b$. Since the supremum over a subset is at most the supremum over the whole set,
$\sup_{D'}|F|\le\sup_{D}|F|\le b$.

Now let $F=\sum_\nu F_\nu$ converge uniformly on $D$, that is,
$\sup_{D}\bigl|\sum_{\nu>\nu_0}F_\nu\bigr|\to0$ as $\nu_0\to\infty$. By the same inequality,
\begin{equation*}
\sup_{D'}\Bigl|\sum_{\nu>\nu_0}F_\nu\Bigr|\le\sup_{D}\Bigl|\sum_{\nu>\nu_0}F_\nu\Bigr|,
\end{equation*}
so the convergence is uniform on $D'$.

Next let $X'\subset X$, let $r$ denote the restriction map from $\mathfrak{U}_k(X)$ to
$\mathfrak{U}_k(X')$, and let $F$ be a function on $\mathfrak{U}_k(X')$. Since
$r(\mathfrak{U}_k(X))\subset\mathfrak{U}_k(X')$, we have
\begin{equation*}
\sup_{\mathfrak{U}_k(X)}|F\circ r|\le\sup_{\mathfrak{U}_k(X')}|F|.
\end{equation*}
Applying this to the tails of a series, uniform convergence on $\mathfrak{U}_k(X')$ gives uniform
convergence of the composed series on $\mathfrak{U}_k(X)$.

If the inputs hold on Ba{\l}aban's space as printed, the conclusions hold there. If they hold on the
one-power hereditary class, the conclusions hold there and, by the chain of inclusions and the first
two paragraphs, also on $U^{\mathrm{ad}}_k(X)$ and $U^{\mathrm{an}}_k(X)$, which are subsets of that
class. If they hold only on $U^{\mathrm{ad}}_k(X)$, the conclusions hold there and, by the same
argument, on $U^{\mathrm{an}}_k(X)$, which is contained in it. If they hold only on
$U^{\mathrm{an}}_k(X)$, the conclusions hold there.
\end{proof}

\begin{definition}[Cubes, cells, touching and tree graphs at level $k$]\label{4.6:def-cubes-cells}
Throughout, lengths are measured in level-$k$ units, that is, in the units of the unit lattice of
\cite{B-CMP119} after $k$ renormalisation steps. The underlying space is either the level-$k$
torus on which the step is performed or $\mathbb{R}^d$. Every object defined below is local and
does not wrap around the torus. The dimension $d \ge 1$ is kept symbolic; in Theorem~0, $d = 4$.
\begin{enumerate}
\item The \emph{$M$-cubes} are the closed cubes
\begin{equation*}
c_n := Mn + [0,M]^d, \qquad n \in \mathbb{Z}^d .
\end{equation*}
\item Let $\check{R}_k$ be the island scale in cells, fixed in the lemma on the island scale, with
$\check{R}_k \ge L$ and $\check{R}_k \in L\mathbb{N}$. The \emph{cells} are the closed cubes of
side $M\check{R}_k$ whose corners lie in $M\check{R}_k\mathbb{Z}^d$.
\item For a level $i \le k$, still in level-$k$ units, put $u_i := L^{-(k-i)}$. The level-$i$
$M$-cubes are the closed cubes of side $Mu_i$, and the level-$i$ cells are the closed cubes of
side $a_i := M\check{R}_i u_i$. Then $a_i \in Mu_i\mathbb{N}$ and $a_k = M\check{R}_k$.
\item Two closed cubes \emph{touch} if they have a point in common, corners included; this is the
relation ``intersect, or touch'' of \cite[p.~386]{B-CMP122b}.
\item For unions $A$, $B$ of $M$-cubes, $A \sqcap B$ denotes the \emph{cube-wise intersection},
the union of the $M$-cubes common to $A$ and $B$. Thus $A \sqcap B = \emptyset$ means that $A$
and $B$ have no common $M$-cube. Wherever containment and disjointness of such unions are used to
split terms into groups, they are understood cube-wise in the same sense.
\item The \emph{components} of a finite union of closed cubes are its classes under the
equivalence relation generated by touching. A touch-connected finite union of closed cubes is
path-connected.
\item For a union $A$ of $M$-cubes, $N(A)$ denotes the number of $M$-cubes of $A$.
\item A \emph{tree graph} is a connected finite union of closed segments. Its \emph{length} is the
sum of the lengths of its segments. Overlapping segments only increase this length. Every use of
the length in this book is an upper bound on the length of a shortest tree graph with the
required property, and in such a graph no two segments overlap.
\end{enumerate}
\end{definition}

\begin{proof}
For the assertions $a_i \in Mu_i\mathbb{N}$ and $a_k = M\check{R}_k$: $\check{R}_i$ lies in
$L\mathbb{N} \subset \mathbb{N}$, so
$a_i = Mu_i\,\check{R}_i$ lies in $Mu_i\mathbb{N}$. At $i = k$ we have $u_k = L^0 = 1$, and hence
$a_k = M\check{R}_k$.

For the path-connectedness of a touch-connected union: each closed cube is convex and therefore
path-connected. If two closed cubes touch, they share a point $x$. Joining a point of the first
cube to $x$ by a segment inside the first cube, and $x$ to a point of the second cube by a segment
inside the second, gives a path in their union between the two points. Now let
$A = Q_1 \cup \dots \cup Q_n$ be touch-connected and
take $y \in Q_a$, $z \in Q_b$. By the definition of components there is a chain
$Q_a = Q_{j_0}, Q_{j_1}, \dots, Q_{j_r} = Q_b$ in which consecutive cubes touch. Joining the
paths through the successive common points gives a path in $A$ from $y$ to $z$.

For the effect of overlaps on the length: let the segments of a tree graph be
$\sigma_1, \dots, \sigma_m$. By subadditivity of
one-dimensional measure, the sum $\sum_{l} |\sigma_l|$ of their lengths is at least the
one-dimensional measure of $\sigma_1 \cup \dots \cup \sigma_m$, with equality when no two
segments overlap. If two segments of a tree graph overlapped in a sub-segment of positive length,
removing the interior of that sub-segment from one of them would replace that segment by at most
two closed segments, would leave the union unchanged, since the removed points still lie in the
other segment, and hence would leave it connected, while strictly decreasing the length. So a
shortest tree graph has no such overlap.
\end{proof}

\begin{definition}[The step datum and admissible island families]\label{4.6:def-step-datum}
Fix a level $k$ and work inside the large-field operation $\mathbf{R}$ at step $k$, as constructed
in \cite[\S1]{B-CMP122a} and \cite[\S1]{B-CMP122b}, in the level-$k$ units, $M$-cubes and cells of
Definition~\ref{4.6:def-cubes-cells}. Write $\Lambda_k$ for the small-field region at step $k$ and
$\Lambda_k^{c}$ for its complement.

\smallskip\noindent
\emph{The sets $Z$ and $Z_k$.} Let $Z_k$ be the large-field region on which $\mathbf{R}$ acts at
step $k$, as in \cite[\S1]{B-CMP122a}, and let $Z\subset\Lambda_k^{c}$ be the union of those
components of $Z_k$ which satisfy conditions (i) and (ii) of \cite[\S1]{B-CMP122a}.
Each such component is, by the parallelepiped rule of
\cite[(2.3)]{B-CMP119} as used in \cite[\S1]{B-CMP122a}, a box made of cells, each cell of side
$M\check{R}_k$, and it is contained in a cube of side $100M\check{R}_k$. This is the same
parallelepiped rule on
which the lemma on admissibility of a single island in this chapter rests. Following
\cite{B-CMP122a}, the remaining part $Z_k\setminus Z$ is again denoted $Z_k$.

\smallskip\noindent
\emph{The two classes of components.} The components of $Z$ are of two kinds. The
\emph{first-class islands} $X_a$ are called \emph{healed}, which here means that each is summed,
together with its history $\mathfrak{h}_a$, inside the curly bracket $\{\cdots\}$ of the
representation of the operation $\mathbf{R}$ in \cite[\S1]{B-CMP122b}. The
\emph{second-class components} are $Y_1,\dots,Y_m$ for some integer $m\ge0$ (in this definition
$m$ denotes this number, not the exponent of the torus $\mathbb{T}_m$); they are called
\emph{thrown back}, which here means that each is fixed, together with its history
$\mathfrak{h}_i$, and the corresponding island integral operations $\mathbf{T}'_k(Y_i)$ stand
outside the curly bracket.

\smallskip\noindent
\emph{The outer datum.} The \emph{outer datum} of the step is
\begin{equation*}
\mathcal{O}:=\bigl(Z_k;\,(Y_i,\mathfrak{h}_i)_{i\le m}\bigr),
\qquad
|\mathcal{O}|:=Z_k\cup\bigcup_{i=1}^{m}Y_i .
\end{equation*}
By \cite[p.~390]{B-CMP122b}, the set $|\mathcal{O}|$ is the new large-field region produced by
the step.

\smallskip\noindent
\emph{Admissibility.} A finite family $\mathfrak{X}=\{(X_a,\mathfrak{h}_a)\}$ of sets with
histories is \emph{admissible under} $\mathcal{O}$ if and only if
\begin{equation}\label{4.6:eq-step-datum-a}
\begin{aligned}
&\text{(c1)}\quad \text{for every } a,\ X_a \text{ is a first-class island on its own;}\\
&\text{(c2)}\quad X_a\cap X_b=\emptyset \quad\text{for } a\neq b;\\
&\text{(c3)}\quad X_a\cap|\mathcal{O}|=\emptyset \quad\text{for every } a.
\end{aligned}
\end{equation}
In (c1), ``on its own'' means that conditions (i) and (ii) of \cite[\S1]{B-CMP122a}, the maturity
criterion of this chapter, and the run of the construction of this chapter carried out for the
island itself are all read inside $X_a$ alone. Conditions
\eqref{4.6:eq-step-datum-a} coincide with the admissibility convention of the family, read
together with the lemma on admissibility of a single island of this chapter.

\smallskip\noindent
\emph{The conventions of the family.} The statements on admissible families used in the rest of
this chapter are formulated under conventions that include those listed here; none of them is
altered later.
\begin{enumerate}
\item The $M$-cube convention of Definition~\ref{4.6:def-cubes-cells}.
\item The cut-off convention and the attachment convention of the family, as fixed in this
chapter.
\item The coupling profile. Fix once and for all
$\psi_\varphi\in C^{\infty}([0,\infty[;[0,1])$ with $\psi_\varphi\equiv1$ on $[0,\tfrac18]$ and
$\psi_\varphi\equiv0$ on $[\tfrac14,\infty[$. For a level $i\le k$ (the letter $i$ now denoting a
level), let $A$ be the level-$i$
member of the sequence of large-field regions in force, namely the prepared sequence
$\{Z''_i\}$ of \cite[(1.10)--(1.11)]{B-CMP122a}, produced by the operations $\mathbf{T}$ at the
current stage of the preparation in \cite[\S1]{B-CMP122a} and retained. For a cell $D$ of
level $i$, write $D=\prod_{\mu=1}^{d}D_\mu$ with closed intervals $D_\mu$. Then
\begin{equation*}
\varphi_i(x):=\prod_{D}\Bigl(1-\prod_{\mu=1}^{d}
\psi_\varphi\bigl(\operatorname{dist}(x_\mu,D_\mu)/a_i\bigr)\Bigr),
\end{equation*}
the product running over the level-$i$ cells $D$ of the level-$i$ cell hull $[A]_i$ of $A$. At
levels at which the member $A$ is made of $M$-cubes, the cells of $A$ are by definition the
cells of $[A]_i$. The prepared coupling function is $g''_k(\cdot):=g_k^{(N)}(\cdot)$ of
\cite[(1.62)]{B-CMP122a}, where the superscript $(N)$ is the notation of
\cite[(1.62)]{B-CMP122a} for the prepared function and is unrelated to the $N$ of $SU(N)$.
With $\mathbf{1}_S$ the indicator function of a set $S$, the coupling function of
\cite[p.~365]{B-CMP122b} is given by
\begin{equation*}
\frac{1}{g_k(\cdot)^{2}}
=\frac{1}{g''_k(\cdot)^{2}}\,\mathbf{1}_{Z^{c}}+\frac{1}{g_k^{2}}\,\mathbf{1}_{Z}.
\end{equation*}
\item The vacuum convention of the family, as fixed in this chapter.
\item The strip convention, with strip width
$m_{\sharp}:=\bigl\lfloor\tfrac12(\check{R}_k-1)\bigr\rfloor$. One has $m_{\sharp}\ge1$ because
$\check{R}_k\ge L\ge3$ (Definitions~\ref{1.1:def-lattice} and~\ref{4.6:def-cubes-cells}).
\end{enumerate}
An alternative convention for the support $\mathcal{D}$ of the coupling excess on the family is
displayed later in this chapter; it alters none of the conventions above.
\end{definition}

\begin{definition}[Strips, interaction domains and the bracket]\label{4.6:def-strips-bracket}
We work at level $k$, in the units of Definition~\ref{4.6:def-cubes-cells}, with its $M$-cubes
$c_n=Mn+[0,M]^d$, its cells (the closed cubes of side $M\check{R}_k$ cornered in
$M\check{R}_k\mathbb{Z}^d$) and its notion of touching. The sup-distance of two sets is their
distance in the maximum norm; we write it $\operatorname{dist}_\infty$. Throughout, Ba{\l}aban's
enlarged region $Z^{\sim}$ of a region $Z$ is replaced by the strip $Z^{\sharp}$ defined below; we
call this the strip convention. Its number of layers $m_\sharp$ is a fixed non-negative integer,
which satisfies $\check{R}_k-2m_\sharp\ge 1$.

\emph{Strips.} Let $S$ be a union of cells or of $M$-cubes. For $1\le n\le m_\sharp$ the
\emph{$n$-th layer} of $S$ is the set of $M$-cubes at sup-distance $(n-1)M$ from $S$. The
\emph{strip} $S^{\sharp}$ of $S$ is $S$ together with its layers $1,\dots,m_\sharp$, and
$S^{\boxplus}$ is the union of $S$ with the cells touching $S$.
These sets have the properties (i)--(iii) below; the quantitative content of (ii) and (iii) is
the display
\begin{equation}\label{4.6:eq-strips-bracket-b}
\operatorname{dist}_\infty\bigl(S_1^{\sharp},S_2^{\sharp}\bigr)\ge(\check{R}_k-2m_\sharp)M\ge M,
\qquad
S^{\sharp}\subset S^{\boxplus},
\end{equation}
for $S_1,S_2$ as in (ii) and $S$ as in (iii).
\begin{enumerate}
\item[(i)] If $S$ is a box of cells, then $S^{\sharp}$ is a box of $M$-cubes.
\item[(ii)] If $S_1,S_2$ are unions of cells that share no point, then the chain of inequalities
in \eqref{4.6:eq-strips-bracket-b} holds, and $S_1^{\sharp}$, $S_2^{\sharp}$ share no point.
\item[(iii)] If $S$ is a union of cells, then the inclusion in \eqref{4.6:eq-strips-bracket-b}
holds.
\end{enumerate}

\emph{The island union.} Let $\mathfrak{X}$ be a family of islands $X_a$ admissible under the
outer datum $\mathcal{O}$, and let $Y_i$ be the components carried by $\mathcal{O}$
(Definition~\ref{4.6:def-step-datum}). Put
\begin{equation*}
Z(\mathfrak{X}):=\bigcup_a X_a\cup\bigcup_i Y_i,
\qquad
Z(\mathfrak{X})^{\sharp}:=\bigcup_a X_a^{\sharp}\cup\bigcup_i Y_i^{\sharp};
\end{equation*}
the sets $X_a^{\sharp}$, $Y_i^{\sharp}$ are pairwise non-touching boxes of $M$-cubes; we call them
the strips of $Z(\mathfrak{X})^{\sharp}$.

\emph{Interaction domains.} With $\mathbf{D}_k$ the class of localization domains at level $k$
and $\sqcap$ the cube-wise intersection,
\begin{equation}\label{4.6:eq-strips-bracket-a}
\begin{split}
\mathcal{Y}(\mathfrak{X}):=\bigl\{Y\in\mathbf{D}_k:\ {}&Y\sqcap Z(\mathfrak{X})^{\sharp}
\text{ is a nonempty union of whole strips of }Z(\mathfrak{X})^{\sharp},\\
&Y\setminus Z(\mathfrak{X})^{\sharp}\neq\emptyset\bigr\}.
\end{split}
\end{equation}
This is the condition of \cite{B-CMP122b} that the intersection of $Y$ with Ba{\l}aban's enlarged
region $Z^{\sim}$ be a union of components of $Z^{\sim}$ containing at least one component, with
$Z^{\sim}$ replaced by the strip union $Z(\mathfrak{X})^{\sharp}$. A domain contained in a strip is
excluded by the second condition; in Ba{\l}aban's words, \emph{``The sum in the last exponential
does not include terms with localization domains equal to one of the components of $Z^{\sim}$;
these terms are included into the operations $\mathbf{T}'_k$''} \cite[p.~379]{B-CMP122b}.

\emph{The island operations.} For an island $X_a$ of $\mathfrak{X}$, $\mathbf{T}'_k(X_a)$ is the
operation of \cite[(1.71)]{B-CMP122b}, read under the strip convention: the integral operation of
the island, including its sum over the internal admissible histories and the frame average. Its
exponent is, in the notation of \cite[(1.71)]{B-CMP122b},
\begin{equation*}
\begin{split}
&-\tfrac12\bigl\langle DH''B,\zeta DH''B\bigr\rangle
-A\bigl((g'')^{-2}\zeta_1\lceil X_a,U''\bigr)\\
&\qquad+V(X_a^{\sharp},U_k)+E_k(X_a)-E_k(\Lambda\cap X_a),
\end{split}
\end{equation*}
together with the normalisation of \cite[(1.71)]{B-CMP122b}. Here $\zeta$, $\zeta_1$ are
Ba{\l}aban's cut-off functions, $g''$ is the prepared coupling function, $U_k$ is the block field,
and $D$, $H''$, $B$, $U''$, $V$, $E_k$, $\Lambda$ are as in \cite[(1.71)]{B-CMP122b}.
The operation $\mathbf{T}'_k(X_a)$ integrates the variables of $X_a$ and is the identity on
functions of the other variables.

\emph{The bracket.} The curly bracket of \cite[(1.72)]{B-CMP122b}, read under the strip
convention, is
\begin{equation}\label{4.6:eq-strips-bracket-1}
\{\cdots\}:=\sum_{\mathfrak{X}\text{ admissible under }\mathcal{O}}\
\prod_{a\in\mathfrak{X}}\mathbf{T}'_k(X_a)\,
\exp\sum_{Y\in\mathcal{Y}(\mathfrak{X})}V(Y).
\end{equation}
Here $V(Y)$ is the term with localization domain $Y$ in the last exponential of
\cite[(1.72)]{B-CMP122b}.
The index set of the exponent depends on $\mathfrak{X}$ through \eqref{4.6:eq-strips-bracket-a}.
The history sums inside the operations $\mathbf{T}'_k(X_a)$ act on the terms $V(Y)$ with
$Y\supseteq X_a$; in Ba{\l}aban's words, \emph{``the sum in the definition (1.71) acts also on the
last exponential in (1.72)''} \cite[p.~379]{B-CMP122b}.
\end{definition}

\begin{proof}
We prove (i)--(iii) and the assertion on $Z(\mathfrak{X})^{\sharp}$.

(i) Let $S=\prod_{i=1}^d[v_i,v'_i]$ be a box of cells. Since $\check{R}_k$ is an integer,
$v_i,v'_i\in M\check{R}_k\mathbb{Z}\subset M\mathbb{Z}$. The sup-distance of $c_n$ to $S$ is
the maximum over $i$ of the distances of the intervals $[Mn_i,M(n_i+1)]$ and $[v_i,v'_i]$.
Each of these is a multiple of $M$. Hence, for $m_\sharp\ge1$, the $M$-cubes of $S^{\sharp}$ are
exactly the $M$-cubes at sup-distance at most $(m_\sharp-1)M$ from $S$; those of $S$ are at
sup-distance $0$. This condition holds coordinatewise, in the form
\begin{equation*}
v_i-m_\sharp M\le Mn_i\le v'_i+(m_\sharp-1)M,\qquad 1\le i\le d .
\end{equation*}
The union of these cubes is $\prod_i[v_i-m_\sharp M,v'_i+m_\sharp M]$, a box of
$M$-cubes. For $m_\sharp=0$ we have $S^{\sharp}=S$.

(ii) Let $c$ be a cube of the $n$-th layer of $S_j$. Then $c$ is at sup-distance
$(n-1)M\le(m_\sharp-1)M$ from $S_j$, and $c$ has sup-diameter $M$. Hence every point of $c$ is at
sup-distance at most $m_\sharp M$ from $S_j$, and so is every point of $S_j^{\sharp}$.

The sup-distance of two cells is a multiple of $M\check{R}_k$, because their corners lie in
$M\check{R}_k\mathbb{Z}^d$. Now $S_1$ and $S_2$ are unions of closed cells sharing no
point. Their sup-distance is therefore the infimum, over pairs formed by a cell of $S_1$ and a
cell of $S_2$, of the sup-distances of these cells; each of these is positive and a multiple of
$M\check{R}_k$, hence at least $M\check{R}_k$. By the triangle inequality,
\begin{equation*}
\operatorname{dist}_\infty\bigl(S_1^{\sharp},S_2^{\sharp}\bigr)
\ge M\check{R}_k-2m_\sharp M=(\check{R}_k-2m_\sharp)M .
\end{equation*}
This is at least $M$ because $\check{R}_k-2m_\sharp\ge1$. In particular the distance is positive,
so $S_1^{\sharp}$ and $S_2^{\sharp}$ share no point.

(iii) Let $c$ be a cube of $S^{\sharp}$ not contained in $S$. Then
\begin{equation*}
\operatorname{dist}_\infty(c,S)\le(m_\sharp-1)M<M\check{R}_k,
\end{equation*}
since $\check{R}_k\ge 2m_\sharp+1$. Since $\check{R}_k$ is an integer, $c$ lies in a cell
$\mathsf{c}$. We have
$\operatorname{dist}_\infty(\mathsf{c},S)\le\operatorname{dist}_\infty(c,S)<M\check{R}_k$. Since
$S$ is a union of cells, this distance is a multiple of $M\check{R}_k$, hence $0$. Thus
$\mathsf{c}$ touches $S$, and $c\subset\mathsf{c}\subset S^{\boxplus}$.

The islands $X_a$ and the components $Y_i$ are boxes of cells that pairwise share no point, by
Definition~\ref{4.6:def-step-datum}. By (i) their strips are boxes of $M$-cubes. By (ii) any two
of these strips are at sup-distance at least $M$, so they do not touch.
\end{proof}

\begin{definition}[Standing geometric facts on boxes and strips]\label{4.6:def-geometric-facts}
Let $k$ be a level, and let the step datum at level $k$ be as in Definition~\ref{4.6:def-step-datum}. A union of closed $M$-cubes is \emph{touch-connected} if any two of its cubes are joined by a chain of its cubes in which consecutive cubes share at least one point. The following facts hold at every level $k$ and are used below without further mention.
\begin{enumerate}
\item[(G1)] Every first-class and every second-class component $C$ of $\Lambda_k^c$ is a box of cells, and it is contained in a cube $Q_C$ of side $100M\check{R}_k$. Distinct components $C \neq C'$ of $\Lambda_k^c$ share no point. In formulas,
\begin{equation}\label{4.6:eq-geometric-facts-a}
C \subset Q_C, \quad Q_C \text{ a cube of side } 100M\check{R}_k;
\qquad C \cap C' = \emptyset \quad (C \neq C').
\end{equation}
\item[(G2)] The strip facts \eqref{4.6:eq-strips-bracket-b} of Definition~\ref{4.6:def-strips-bracket} hold.
\item[(G3)] We read the class $\mathbf{D}_k$ of localisation domains of \cite{B-CMP109} as consisting of touch-connected unions of closed $M$-cubes: every $Y \in \mathbf{D}_k$ satisfies
\begin{equation}\label{4.6:eq-geometric-facts-b}
Y = \bigcup_{c \in \mathcal{S}_Y} c,
\qquad \mathcal{S}_Y \text{ a touch-connected family of closed } M\text{-cubes}.
\end{equation}
Should $\mathbf{D}_k$ be read as admitting members that are not touch-connected, the family $\mathcal{Y}(\mathfrak{X})$ of admissible interaction domains is restricted to its touch-connected members. These are the domains produced by the construction in \cite[(1.59)--(1.64)]{B-CMP122b}.
\end{enumerate}
\end{definition}

\begin{proof}
(G1) Both assertions about first-class and second-class components are part of Definition~\ref{4.6:def-step-datum}: each such component lies in a cube of side $100M\check{R}_k$ \cite[\S1]{B-CMP122a}, and it is a box of cells by the parallelepiped rule fixed in that definition.

That distinct components of $\Lambda_k^c$ share no point is part of the statement of the relative small-field machine describing the components of the complement $\Lambda_k^c$ of the small-field region, taken together with the separation statement for these components that belongs to the same result.

(G2) This is \eqref{4.6:eq-strips-bracket-b}, established in Definition~\ref{4.6:def-strips-bracket}.

(G3) The cited passages support this reading. In \cite{B-CMP109} the distance function $d_k$ is defined on the class $\mathbf{D}_k$: for $X \in \mathbf{D}_k$ the length $d_k(X)$ is defined through tree graphs contained in $X$ and joining its $M$-cubes. Such a tree graph exists only if $X$ is touch-connected.

In \cite{B-CMP122b} the bound for the interaction terms on a domain $Y$ is derived using that $Y$ is connected, so the domains of $\mathbf{D}_k$ entering there are connected.

For an island $X$ of the admissible family $\mathfrak{X}$ and $Z$ as in Definition~\ref{4.6:def-step-datum}, the interaction domain is, by \cite[p.~373]{B-CMP122b}, the smallest localisation domain containing $X$ and the components of $Z$ that intersect $X$. In particular it belongs to $\mathbf{D}_k$, and the reading supported in the first paragraph of this part of the proof applies to it; this gives \eqref{4.6:eq-geometric-facts-b} for the interaction domains.

Suppose $\mathbf{D}_k$ is read as admitting members that are not touch-connected. Then the restriction stated in (G3) keeps in $\mathcal{Y}(\mathfrak{X})$ exactly the touch-connected domains, namely those produced by \cite[(1.59)--(1.64)]{B-CMP122b}. On these, \eqref{4.6:eq-geometric-facts-b} holds by definition.
\end{proof}

\begin{definition}[Expansion of the bracket into polymers]\label{4.6:def-polymer-expansion}
Let $\mathfrak{X}$ run over the admissible island families of \eqref{4.6:eq-step-datum-a}, with
islands $X_a$, $a\in\mathfrak{X}$. Let $\mathcal{Y}(\mathfrak{X})$ be the family of admissible
interaction domains and $\{\cdots\}$ the bracket of \eqref{4.6:eq-strips-bracket-a}, as fixed in
Definition~\ref{4.6:def-strips-bracket}. Let $V(Y)$ be the interaction term of a domain
$Y\in\mathcal{Y}(\mathfrak{X})$, $\mathbf{T}'_k(X_a)$ the island integral operation of
\cite[(1.71)]{B-CMP122b}, and $Y_i$ the second-class components, with strips $Y_i^{\sharp}$.
Two closed $M$-cubes \emph{touch} if they share a point. The \emph{touch-components} of a union of
closed $M$-cubes are its maximal subfamilies connected under this relation.

\emph{Mayer step.} By the identity
\begin{equation}\label{4.6:eq-polymer-expansion-1}
e^{V}-1=\int_0^1 dt\,V\,e^{tV}
\end{equation}
of \cite[(2.1)]{B-CMP116}, the bracket takes the form
\begin{equation}\label{4.6:eq-polymer-expansion-2}
\{\cdots\}=\sum_{\mathfrak{X}}\ \sum_{\mathbf{D}\subset\mathcal{Y}(\mathfrak{X})}\
\prod_{a\in\mathfrak{X}}\mathbf{T}'_k(X_a)\prod_{Y\in\mathbf{D}}\int_0^1 dt_Y\,V(Y)\,
e^{t_YV(Y)}.
\end{equation}
A pair $(\mathfrak{X},\mathbf{D})$ indexing a term of \eqref{4.6:eq-polymer-expansion-2} is a
\emph{datum} of the expansion.

\emph{Raw domain and core of a datum.} The \emph{raw domain} of $(\mathfrak{X},\mathbf{D})$ is
\begin{equation*}
X'_0:=\bigcup_{Y\in\mathbf{D}}Y\ \cup\ \bigcup_{a\in\mathfrak{X}}X_a^{\sharp}.
\end{equation*}
Its \emph{core} is $X'_0\setminus\bigcup_iY_i^{\sharp}$, the set of $M$-cubes of $X'_0$ that
lie in no second-class strip. This is the set $X'_0\setminus\bigcup_iY_i$ of
\cite{B-CMP122b} with each $Y_i$ replaced by its strip $Y_i^{\sharp}$; the replacement is the
convention under which the tree length over the core, written $d_{\mathrm{out}}$, is taken
relative to $\bigcup_iY_i^{\sharp}$.

\emph{Classes.} The \emph{core cubes} of $Y\in\mathbf{D}$ are the $M$-cubes of $Y$ that lie in
the core. Two touch-components of the core are \emph{glued} if some $Y\in\mathbf{D}$ has
$M$-cubes of the core in both. The \emph{classes} of the datum are the equivalence classes of the
touch-components of the core under the equivalence relation generated by gluing. This is the
connection rule of \cite[p.~388]{B-CMP122b}, printed there as
\begin{quote}
we consider two such components as connected, if they are contained in one localization
domain $Y$ of a term
\end{quote}
and read here on the core. The \emph{core part} of a class $p$ is the union of
its touch-components. The class $p$ carries
\begin{equation*}
\mathfrak{X}_p:=\{a\in\mathfrak{X}:X_a^{\sharp}\subset\text{the core part of }p\},
\end{equation*}
\begin{equation*}
\mathbf{D}_p:=\{Y\in\mathbf{D}:\text{the core cubes of }Y\text{ lie in the core part of }p\},
\end{equation*}
and its \emph{domain}
\begin{equation}\label{4.6:eq-polymer-expansion-3}
X'_p:=\bigcup_{Y\in\mathbf{D}_p}Y\ \cup\ \bigcup_{a\in\mathfrak{X}_p}X_a^{\sharp}.
\end{equation}

\emph{Polymers and activities.} A \emph{polymer} is a set that arises as the domain of a class of
some datum. The set of polymers of the step is denoted $\mathcal{P}$; on a torus it is finite,
and the lattice on $\mathbb{R}^d$ is treated separately below. The \emph{activity} of a
polymer $X'$ is
\begin{equation}\label{4.6:eq-polymer-expansion-a}
F(X'):=\sum_{\substack{(\mathfrak{X},\mathbf{D})\colon\ \text{one class,}\\
\text{domain }X'}}\ \prod_{a\in\mathfrak{X}}\mathbf{T}'_k(X_a)
\prod_{Y\in\mathbf{D}}\int_0^1 dt_Y\,V(Y)\,e^{t_YV(Y)} .
\end{equation}
The sum runs over the admissible families $\mathfrak{X}$ and the sets
$\mathbf{D}\subset\mathcal{Y}(\mathfrak{X})$ whose datum has exactly one class, with domain exactly
$X'$. This is the activity of \cite[(1.91)]{B-CMP122b}, where the summation is over the families of
islands and the sets of localization domains whose connected localization domain equals $X'$, with
second-class components replaced by their strips.

\emph{Core of a polymer; incompatibility.} For a polymer $X'$ set
\begin{equation}\label{4.6:eq-polymer-expansion-b}
\operatorname{core}(X'):=X'\setminus\bigcup_iY_i^{\sharp},\qquad
X'\,\iota\,X''\ :\Longleftrightarrow\ \operatorname{core}(X')\cap\operatorname{core}(X'')
\neq\emptyset .
\end{equation}
The set $\operatorname{core}(X')$ depends on $X'$ alone and not on the datum from which $X'$
arises. Polymers $X',X''$ with $X'\,\iota\,X''$ are \emph{incompatible}. The relation $\iota$ is
reflexive and symmetric.
\end{definition}

\begin{proof}
We show that the classes carry well-defined, disjoint sets $\mathfrak{X}_p$ and $\mathbf{D}_p$
exhausting $\mathfrak{X}$ and $\mathbf{D}$. We then verify the two properties of $\iota$.

Let $a\in\mathfrak{X}$. By definition of the raw domain, $X_a^{\sharp}\subset X'_0$. The strips
of first-class and of second-class components share no point, by (G2), that is by
\eqref{4.6:eq-strips-bracket-b}. Hence $X_a^{\sharp}$ meets no $Y_i^{\sharp}$, and
$X_a^{\sharp}$ is contained in the core. The island $X_a$ is a box of cells, by (G1) of
Definition~\ref{4.6:def-geometric-facts}. The strip of a box of cells is a box of $M$-cubes, by
(G2), so $X_a^{\sharp}$ is a box of $M$-cubes. A box of $M$-cubes is touch-connected; hence
$X_a^{\sharp}$ lies inside one touch-component of the core, and so inside the core part of exactly
one class. Thus $a$ belongs to exactly one $\mathfrak{X}_p$.

Let $Y\in\mathbf{D}$. By (G3) of Definition~\ref{4.6:def-geometric-facts}, $Y$ is a union of
closed $M$-cubes, so its core cubes are defined. Since $Y\in\mathcal{Y}(\mathfrak{X})$,
the set $Y\setminus Z(\mathfrak{X})^{\sharp}$ is non-empty. Its $M$-cubes lie off
$\bigcup_iY_i^{\sharp}$, and they lie in $X'_0$ because $Y\subset X'_0$. Hence they are cubes of
the core, and $Y$ has at least one core cube.

Any two touch-components of the core that contain core cubes of $Y$ are glued. By the definition of
the classes, they lie in the same class. Therefore all core cubes of $Y$ lie in the core part of a
single class $p$, and $Y$ belongs to $\mathbf{D}_p$ for exactly this $p$. Consequently the domains
\eqref{4.6:eq-polymer-expansion-3} are well defined.

Symmetry of $\iota$ is immediate from the symmetry of the intersection in
\eqref{4.6:eq-polymer-expansion-b}. For reflexivity, let $X'=X'_p$ be the domain of a class $p$
of some datum. We show that every $M$-cube $c$ of the core part of $p$ lies in
$\operatorname{core}(X')$.

The cube $c$ lies in $X'_0$ and in no $Y_i^{\sharp}$. Suppose first that $c\subset X_a^{\sharp}$
for some $a\in\mathfrak{X}$. By the first paragraph, $X_a^{\sharp}$ lies in a single
touch-component, which contains $c$ and hence belongs to $p$. Thus $a\in\mathfrak{X}_p$ and
$c\subset X'$. Suppose instead that $c\subset Y$ for some $Y\in\mathbf{D}$. Then $c$ is a core
cube of $Y$, all of whose core cubes lie in the class of $c$, namely $p$. Thus $Y\in\mathbf{D}_p$
and $c\subset X'$.

In either case $c\subset X'\setminus\bigcup_iY_i^{\sharp}=\operatorname{core}(X')$. A class
contains at least one touch-component, so $\operatorname{core}(X')\neq\emptyset$, and
$X'\,\iota\,X'$.
\end{proof}

\begin{remark}[Locality of the interaction terms; printed without proof in {\cite[p.~379]{B-CMP122b}}]\label{4.6:rem-interaction-locality}
Let $Y$ be an interaction domain in the sense of
Definition~\ref{4.6:def-strips-bracket}, and let $V(Y)$ denote the interaction term attached to $Y$.
Let $(\mathbf{U},\mathbf{J})$ be the complex multi-scale configuration, and let $Y_i$ be the
domains of the outer datum $\mathcal{O}$. Let $X_a$ be the islands, and let $\mathfrak{h}_a$ be
the history of the island $X_a$. Then the following holds:
\begin{equation}\label{4.6:eq-interaction-locality-a}
\begin{gathered}
V(Y)\ \text{is a function of}\ (\mathbf{U},\mathbf{J})|_{Y},\
\text{of the (fixed) variables of the}\ Y_i\subset Y,\\
\text{and of the integration variables and the histories}\ \mathfrak{h}_a\
\text{of the islands}\ X_a\subset Y\ \text{only}.
\end{gathered}
\end{equation}
This is stated in \cite[p.~379]{B-CMP122b} without proof, in the following sentence, which is
cited as printed:
\begin{quote}
Notice that the terms $V(Y,U_k)$ depend on the sequence $\{\Omega_j^c,Z_j\}$ restricted to
the components of $Z$ contained in $Y$
\end{quote}
We take this sentence to assert that the dependence is on the listed data alone: $V(Y)$
depends on no data other than those listed in \eqref{4.6:eq-interaction-locality-a}; in
particular, it depends on the configuration only through its restriction
$(\mathbf{U},\mathbf{J})|_{Y}$.
\end{remark}

\begin{lemma}[The polymer-gas representation of the bracket]\label{4.6:lem-polymer-gas}
Strips are taken in the convention of Definition~\ref{4.6:def-strips-bracket}. Data
$(\mathfrak{X},\mathbf{D})$, their cores and classes, the set $\mathcal{P}$ of polymers, the
activities $F(X')$ of \eqref{4.6:eq-polymer-expansion-a} and the incompatibility relation
$\iota$ of \eqref{4.6:eq-polymer-expansion-b} are those of
Definition~\ref{4.6:def-polymer-expansion}. Two polymers are \emph{compatible} when they are
not $\iota$-related. For a datum and a class $p$ of it, $(\mathfrak{X}_p,\mathbf{D}_p)$
denotes the part of the datum in the class $p$, and $X'_p$ its polymer.
\begin{itemize}
\item[(a)] Assume the locality of the interaction terms
(Remark~\ref{4.6:rem-interaction-locality}; printed without proof in \cite[\S3]{B-CMP122b};
cited as printed). Then for every datum, the integrand of its term in the expansion of the
bracket $\{\cdots\}$ factorises over the classes of the datum. The factor of the class $p$
involves, among the integration variables of the islands, only those of the islands in
$\mathfrak{X}_p$.
\item[(b)] The map $(\mathfrak{X},\mathbf{D})\mapsto\{(\mathfrak{X}_p,\mathbf{D}_p)\}_p$ is a
bijection from the set of data onto the set of finite families of single-class data whose
polymers are pairwise compatible. Consequently, under the hypothesis of (a),
\begin{equation}\label{4.6:eq-polymer-gas-a}
\{\cdots\}=\sum_{r\ge 0}\
\sum_{\substack{\{X'_1,\dots,X'_r\}\subset\mathcal{P}\\ \text{pairwise compatible}}}\
\prod_{p=1}^{r}F(X'_p).
\end{equation}
\item[(c)] Every polymer $X'$ is touch-connected, and it is the disjoint union of its core and
of the second-class strips $Y_i^{\sharp}$ contained in it:
\begin{equation}\label{4.6:eq-polymer-gas-b}
X'=\operatorname{core}(X')\sqcup\Bigl(\,\bigsqcup_{i:\,Y_i^{\sharp}\subset X'}Y_i^{\sharp}\Bigr).
\end{equation}
Each second-class strip contained in $X'$ touches $\operatorname{core}(X')$.
Every first-class strip $X_a^{\sharp}$ of a datum lies in the core of the datum. For every
datum $(\mathfrak{X},\mathbf{D})$ and every $Y\in\mathbf{D}$, the set $Y$ is touch-connected
and $Y\setminus Z(\mathfrak{X})^{\sharp}\neq\emptyset$.
\end{itemize}
\end{lemma}

\begin{proof}
\emph{(a).} Fix a datum $(\mathfrak{X},\mathbf{D})$. Its term in the expansion of the bracket
(Definition~\ref{4.6:def-polymer-expansion}) is
\begin{equation*}
\prod_{a\in\mathfrak{X}}\mathbf{T}'_k(X_a)\prod_{Y\in\mathbf{D}}\int_0^1 dt_Y\,
V(Y)\,e^{t_YV(Y)}.
\end{equation*}
Let $a\in\mathfrak{X}$ and $Y\in\mathbf{D}$. We use the locality of the interaction terms
(Remark~\ref{4.6:rem-interaction-locality}; printed without proof in
\cite[\S3]{B-CMP122b}; cited as printed). By \eqref{4.6:eq-interaction-locality-a}, $V(Y)$ is
a function only of:
\begin{itemize}
\item $(\mathbf{U},\mathbf{J})$ restricted to $Y$;
\item the fixed variables of the second-class components $Y_i\subset Y$;
\item the integration variables and histories of the islands $X_a\subset Y$.
\end{itemize}
Hence $V(Y)$ involves the integration variables of $X_a$ only if $X_a\subset Y$.

In that case $X_a^{\sharp}\subset Y$ by condition \eqref{4.6:eq-strips-bracket-a}. The cubes
of $X_a^{\sharp}$ are core cubes, since first-class strips lie in the core
(Definition~\ref{4.6:def-polymer-expansion}; see the end of this proof), and they lie in $Y$.
Thus $Y$ has core cubes in common with $X_a^{\sharp}$. By the definition of the classes, the
island $a$ and the domain $Y$ therefore belong to one class.

The operation $\mathbf{T}'_k(X_a)$ acts on the integration variables of $X_a$ and is the
identity on functions of the other variables. Three facts follow.
\begin{itemize}
\item Each operation $\mathbf{T}'_k(X_a)$, together with every $V(Y)$ on which it does not
act as the identity, lies in one class.
\item Distinct classes share no first-class integration variables.
\item Functions of the fixed variables of the second-class components merely multiply.
\end{itemize}
Consequently the term of the datum equals
\begin{equation}\label{4.6:eq-polymer-gas-1}
\prod_{p}\Bigl[\prod_{a\in\mathfrak{X}_p}\mathbf{T}'_k(X_a)\prod_{Y\in\mathbf{D}_p}
\int_0^1 dt_Y\,V(Y)\,e^{t_YV(Y)}\Bigr],
\end{equation}
and the factor of the class $p$ involves first-class integration variables of
$\mathfrak{X}_p$ only.

\emph{(b).} The classes of one datum have cores sharing no point. Indeed, two core cubes
sharing a point lie in one touch-component of the core, hence in one class. Therefore the
family $\{X'_p\}$ is pairwise compatible, and the map of (b) takes values in the stated set.

Conversely, let $(\mathfrak{X}_p,\mathbf{D}_p)_{p\le r}$ be single-class data with pairwise
compatible polymers $X'_p$. Put
\begin{equation*}
\mathfrak{X}:=\bigcup_p\mathfrak{X}_p,\qquad \mathbf{D}:=\bigcup_p\mathbf{D}_p .
\end{equation*}
We show in three steps that $(\mathfrak{X},\mathbf{D})$ is a datum whose classes are the given
ones.

\emph{(i) $\mathfrak{X}$ is admissible.} The conditions \eqref{4.6:eq-step-datum-a} split
into two kinds. The first and third conditions are conditions on single islands, and they
hold because each $\mathfrak{X}_p$ is admissible. The second condition holds within each
$\mathfrak{X}_p$. It also holds for a pair of islands $a\in\mathfrak{X}_p$,
$b\in\mathfrak{X}_{p'}$ with $p\neq p'$: we have
$X_a^{\sharp}\subset\operatorname{core}(X'_p)$ and
$X_b^{\sharp}\subset\operatorname{core}(X'_{p'})$, and these cores share no point by
compatibility.

\emph{(ii) $\mathbf{D}\subset\mathcal{Y}(\mathfrak{X})$.} Let $Y\in\mathbf{D}_p$, and suppose
a cube of $Y$ lay inside some $X_b^{\sharp}$ with $b\in\mathfrak{X}_{p'}$, $p'\neq p$. That
cube would be a core cube of $X'_{p'}$, since $X_b^{\sharp}\subset\operatorname{core}(X'_{p'})$.
Being a core cube, it lies off $\bigcup_iY_i^{\sharp}$. Being also in $Y\subset X'_p$, it
would be a core cube of $X'_p$ as well. The cores of $X'_p$ and $X'_{p'}$ would then share a
point, which compatibility excludes. Hence
\begin{equation*}
Y\sqcap Z(\mathfrak{X})^{\sharp}=Y\sqcap Z(\mathfrak{X}_p)^{\sharp}.
\end{equation*}
Condition \eqref{4.6:eq-strips-bracket-a} holds for $Y$ relative to $\mathfrak{X}_p$, because
$Y\in\mathcal{Y}(\mathfrak{X}_p)$. By this equality it transfers to $Y$ relative to
$\mathfrak{X}$.

\emph{(iii) The classes of $(\mathfrak{X},\mathbf{D})$ are the given ones.} The core of
$(\mathfrak{X},\mathbf{D})$ is the union of the cores of the $X'_p$. These cores pairwise
share no point, so no touch-component of the union meets two of them. A domain
$Y\in\mathbf{D}_p$ has its core cubes in $\operatorname{core}(X'_p)$, so no gluing joins two
of them either. Within each $p$, the touch-components are glued into one class, because
$(\mathfrak{X}_p,\mathbf{D}_p)$ is a single-class datum.

By (iii), applying the first map to $(\mathfrak{X},\mathbf{D})$ returns the family
$\{(\mathfrak{X}_p,\mathbf{D}_p)\}_p$. Applying the union construction to the parts of a datum
returns the datum. The two maps are therefore mutually inverse.

We now derive \eqref{4.6:eq-polymer-gas-a}. By (a), the term \eqref{4.6:eq-polymer-gas-1} of a
datum is the product of the terms of its single-class parts. By the bijection, the sum over
data of the bracket becomes a sum over finite families of single-class data with pairwise
compatible polymers. We group the single-class data according to their polymers; the
activity \eqref{4.6:eq-polymer-expansion-a} of a polymer collects the single-class data with
that polymer. This gives \eqref{4.6:eq-polymer-gas-a}, the counterpart of
\cite[(1.90)]{B-CMP122b} in the present strip convention. The term $r=0$ is the datum
$(\emptyset,\emptyset)$, whose value is $1$.

The relation $\iota$ is reflexive, and this is consistent with the grouping. Two single-class
data with the same polymer never occur together in a family, since their cores coincide and
so share a point. They are alternatives inside the same activity $F$.

\emph{(c).} Let $X'$ be the polymer of the class $p$ of a datum. Fix a second-class strip
$Y_i^{\sharp}$. Cube-wise,
\begin{equation*}
X'\cap Y_i^{\sharp}=\bigcup_{Y\in\mathbf{D}_p}\bigl(Y\cap Y_i^{\sharp}\bigr).
\end{equation*}
By \eqref{4.6:eq-strips-bracket-a}, each $Y\cap Y_i^{\sharp}$ is either $\emptyset$ or
$Y_i^{\sharp}$, so $X'\cap Y_i^{\sharp}\in\{\emptyset,Y_i^{\sharp}\}$. The core of $X'$ is
$X'$ with the second-class strips removed. Hence $X'$ is the disjoint union of its core and
of whole second-class strips, which is \eqref{4.6:eq-polymer-gas-b}.

\emph{Touch-connectedness of $X'$.} Each $X_a^{\sharp}$ is a box of $M$-cubes by
\eqref{4.6:eq-strips-bracket-b}, hence touch-connected. Each $Y$ is touch-connected by
\eqref{4.6:eq-geometric-facts-b}. The class is chained from touch-components of the core by
gluing along the domains $Y$ that meet the touch-components already chained. Therefore $X'$
is touch-connected.

\emph{Second-class strips touch the core.} Let $Y_i^{\sharp}\subset X'$. By the cube-wise
identity above, $Y_i^{\sharp}\subset Y$ for some $Y\in\mathbf{D}_p$. The set
$Y\setminus Z(\mathfrak{X})^{\sharp}$ is non-empty and $Y$ is touch-connected (see the last
paragraph of this proof). So there is a chain of pairwise touching cubes of $Y$ leading from a
cube of $Y_i^{\sharp}$ to a cube of $Y\setminus Z(\mathfrak{X})^{\sharp}$. The last cube of
the chain lies outside $Y_i^{\sharp}$; let $c$ be the first cube of the chain outside
$Y_i^{\sharp}$.
Then $c$ touches the preceding cube, which lies in $Y_i^{\sharp}$; so $c$ shares a point with
$Y_i^{\sharp}$. By \eqref{4.6:eq-strips-bracket-b}, strips of distinct components share no
point, so $c$ lies in no other second-class strip. Since $c\in Y\subset X'$, it is a core cube
of $X'$ touching $Y_i^{\sharp}$.

\emph{First-class strips lie in the core.} That every $X_a^{\sharp}$ lies in the core is shown
in Definition~\ref{4.6:def-polymer-expansion}, where the core is introduced.

\emph{The domains $Y$.} Let $(\mathfrak{X},\mathbf{D})$ be a datum and $Y\in\mathbf{D}$.
Since $\mathbf{D}\subset\mathcal{Y}(\mathfrak{X})$ and the admissible interaction domains are
localization domains of the class $\mathbf{D}_k$, the set $Y$ is a touch-connected union of
closed $M$-cubes by \eqref{4.6:eq-geometric-facts-b}. The condition
$Y\setminus Z(\mathfrak{X})^{\sharp}\neq\emptyset$ is part of the definition of the admissible
interaction domains $\mathcal{Y}(\mathfrak{X})$ (Definition~\ref{4.6:def-strips-bracket}).
\end{proof}

\begin{definition}[Tree-graph lengths relative to the strips]\label{4.6:def-tree-lengths}
Let $\mathfrak{X}$ be an admissible family of healed islands
(Definition~\ref{4.6:def-step-datum}), let $Y_i$ be the second-class components of the outer datum,
and let $Z(\mathfrak{X})$, $Z(\mathfrak{X})^{\sharp}$ and $\mathcal{Y}(\mathfrak{X})$ be as in
Definition~\ref{4.6:def-strips-bracket}. For unions $Y$ and $Z'$ of $M$-cubes we write
$d_{k,Z'}(Y)$ for
Ba{\l}aban's length \cite[(1.67)]{B-CMP122b}: $M^{-1}$ times the length of a shortest tree graph
contained in $Y$ and intersecting all $M$-cubes of the components of $Y\setminus Z'$. Following
the strip convention of Definition~\ref{4.6:def-strips-bracket}, Ba{\l}aban's region $Z$ is
replaced throughout by its strip
$Z^{\sharp}$. Further, $d_k(X)$ denotes Ba{\l}aban's linear size \cite{B-CMP109}, recalled below.
We set
\begin{equation}\label{4.6:eq-tree-lengths-a}
\begin{aligned}
d_{\sharp}(Y)&:=d_{k,Z(\mathfrak{X})^{\sharp}}(Y),
&&Y\in\mathcal{Y}(\mathfrak{X}),\\
d_{\mathrm{out}}(X)&:=d_{k,\bigcup_iY_i^{\sharp}}(X),\\
d_k(X)&\le d_{\mathrm{out}}(X)
&&\text{if }\operatorname{core}(X)=X,\\
N_a&:=N(X_a^{\sharp})\le(100\check{R}_k+2m_{\sharp})^d\le(101\check{R}_k)^d,\\
d_{k,Z(\mathfrak{X})^{\sharp}}(Y)&=d_{k,Z(\mathfrak{X}_p)^{\sharp}}(Y),
&&Y\in\mathbf{D}_p.
\end{aligned}
\end{equation}
In words:
\begin{itemize}
\item $d_{\sharp}(Y)$ is $M^{-1}$ times the minimal length of a tree graph contained in $Y$ that
meets every $M$-cube of $Y\setminus Z(\mathfrak{X})^{\sharp}$. If no such tree graph exists,
we put $d_{\sharp}(Y)=+\infty$. For such $Y$ we have $V(Y)=0$ by the interaction-term bound
in \eqref{4.6:eq-standing-bounds-b}, assumed below, and such $Y$ are dropped from all sums.
\item For a polymer, or a union $X$ of polymers, $d_{\mathrm{out}}(X)$ is $M^{-1}$ times the
minimal length of a tree graph in $X$ that meets every $M$-cube of $\operatorname{core}(X)$. Here
$\operatorname{core}$ is as in \eqref{4.6:eq-polymer-expansion-b}.
\item $d_k(X)$ is the linear size of \cite{B-CMP109}: $M^{-1}$ times the length of a shortest
tree graph meeting every $M$-cube of $X$. Ba{\l}aban's wording can be read with or without the
requirement that the tree graph be contained in $X$.
\item $N_a$ is the number of $M$-cubes of the strip $X_a^{\sharp}$ of the island $X_a$.
\item The last line of \eqref{4.6:eq-tree-lengths-a} holds for $Y$ in the part $\mathbf{D}_p$ of
$\mathbf{D}$ belonging to class $p$ (Definition~\ref{4.6:def-polymer-expansion}). There
$\mathfrak{X}_p$ is the part of $\mathfrak{X}$ in class $p$, and
$Z(\mathfrak{X}_p):=\bigcup_{a\in\mathfrak{X}_p}X_a\cup\bigcup_iY_i$.
\end{itemize}
\end{definition}

\begin{proof}
We justify the three relations in \eqref{4.6:eq-tree-lengths-a}.

\emph{The inequality $d_k\le d_{\mathrm{out}}$.} Suppose $\operatorname{core}(X)=X$. Then the
tree graphs admitted in $d_{\mathrm{out}}(X)$ are the tree graphs in $X$ meeting every $M$-cube
of $X$. If the tree graphs in $d_k(X)$ are required to be contained in $X$, these are
exactly the tree graphs admitted in $d_k(X)$, and $d_k(X)=d_{\mathrm{out}}(X)$. If that
requirement is not made, the minimum defining $d_k(X)$ is taken over a larger set of tree
graphs, and $d_k(X)\le d_{\mathrm{out}}(X)$. Under either reading the inequality holds. It is
the only relation between the two lengths used in what follows.

\emph{The count $N_a$.} The island $X_a$ is contained in a cube of side $100M\check{R}_k$
(Definition~\ref{4.6:def-step-datum}); by \eqref{4.6:eq-geometric-facts-a} its strip satisfies
$N(X_a^{\sharp})\le(100\check{R}_k+2m_{\sharp})^d$. The inequality
$2m_{\sharp}\le\check{R}_k$, which holds by Definition~\ref{4.6:def-strips-bracket}, then gives
$100\check{R}_k+2m_{\sharp}\le101\check{R}_k$, hence
$N_a\le(101\check{R}_k)^d$.

\emph{The equality on the family.} Let $Y\in\mathbf{D}_p$. The strip of a union of cells is the
union of the strips of its parts, since the layers of $M$-cubes are defined by sup-distance.
Hence $Z(\mathfrak{X})^{\sharp}$ is the union of $Z(\mathfrak{X}_p)^{\sharp}$ with the strips
$X_a^{\sharp}$, $a\notin\mathfrak{X}_p$, of the islands of the other classes. In the proof of
part (b) of Lemma~\ref{4.6:lem-polymer-gas} it is shown that $Y$ meets no strip of another
class. Therefore
\begin{equation*}
Y\setminus Z(\mathfrak{X})^{\sharp}=Y\setminus Z(\mathfrak{X}_p)^{\sharp}.
\end{equation*}
So the $M$-cubes that an admissible tree graph must meet are the same in
$d_{k,Z(\mathfrak{X})^{\sharp}}(Y)$ and in $d_{k,Z(\mathfrak{X}_p)^{\sharp}}(Y)$. The tree graphs
are in both cases contained in $Y$, so the two minima coincide.
\end{proof}

\begin{definition}[The logarithm of the polymer gas and its two groups]
\label{4.6:def-polymer-logarithm}
Fix a level $k$. Let $\mathcal{P}$ be the set of polymers of the step at level $k$ and $F(X')$ the activity of a polymer $X'\in\mathcal{P}$. Let $\iota$ be the incompatibility relation between polymers, under which two polymers are related when their cores share a point. All three are as in Definition~\ref{4.6:def-polymer-expansion}.

\begin{enumerate}
\item For $X',X''\in\mathcal{P}$ put $\zeta(X',X'')=0$ if $X'$ and $X''$ are $\iota$-related, and $\zeta(X',X'')=1$ otherwise.

\item For $n\ge 1$ let $C_n$ be the set of connected graphs on the vertex set $\{1,\dots,n\}$. A graph is identified with its set of edges $\{p,q\}$, $p\neq q$. For $(X'_1,\dots,X'_n)\in\mathcal{P}^n$ the \emph{Ursell coefficient} is
\begin{equation*}
\rho^{T}(X'_1,\dots,X'_n)
:=\sum_{g\in C_n}\ \prod_{\{p,q\}\in g}\bigl(\zeta(X'_p,X'_q)-1\bigr).
\end{equation*}
For $n=1$ the set $C_1$ consists of the graph with one vertex and no edge. The empty product equals $1$, so $\rho^{T}(X'_1)=1$.

\item Let $Z_k$ and the regions $Y_i$ be those of the outer datum $\mathcal{O}$ of the step, as in Definition~\ref{4.6:def-step-datum}, and let $|\mathcal{O}|$ be its support. Let $S\mapsto S^{\boxplus}$ be the enlargement of $S$ by the cells touching it, as in Definition~\ref{4.6:def-strips-bracket}. For a union $X$ of $M$-cubes, $\mathbf{R}'^{(k)}(X)$ is defined by the first relation, and $|\mathcal{O}|^{\boxplus}$ by the third relation, of
\begin{equation}\label{4.6:eq-polymer-logarithm-a}
\begin{aligned}
\mathbf{R}'^{(k)}(X)
&:=\sum_{n\ge 1}\frac{1}{n!}
\sum_{\substack{(X'_1,\dots,X'_n)\in\mathcal{P}^n\\ X'_1\cup\dots\cup X'_n=X}}
\rho^{T}(X'_1,\dots,X'_n)\\
&\qquad\qquad\times\prod_{p=1}^{n}F(X'_p),\\
\{\cdots\}&=\exp\sum_{X}\mathbf{R}'^{(k)}(X),\\
|\mathcal{O}|^{\boxplus}&:=Z_k^{\boxplus}\cup\bigcup_{i}Y_i^{\boxplus}.
\end{aligned}
\end{equation}
The first relation occupies two rows. The second relation is not part of the definition; it is the identity discussed after this definition, the sum running over unions $X$ of $M$-cubes.

\item The term $\mathbf{R}'^{(k)}(X)$ belongs to the \emph{first group} if $X\sqcap|\mathcal{O}|^{\boxplus}\neq\emptyset$, where $\sqcap$ denotes cube-wise intersection. For such $X$ we write $\mathbf{B}'^{(k)}(X):=\mathbf{R}'^{(k)}(X)$; these are the new boundary terms of the step. The term belongs to the \emph{second group} if $X\sqcap|\mathcal{O}|^{\boxplus}=\emptyset$; these are the new $\mathbf{R}$-terms of the step.
\end{enumerate}
\end{definition}

The first relation of \eqref{4.6:eq-polymer-logarithm-a}, which defines $\mathbf{R}'^{(k)}(X)$, is the formula \cite[(2.13)]{B-CMP116}, written for the polymers, activities and incompatibility relation of the present step.

Its role is the following. Assume that the series $\sum_{X}\mathbf{R}'^{(k)}(X)$, with every $\mathbf{R}'^{(k)}(X)$ written out as in the first relation of \eqref{4.6:eq-polymer-logarithm-a}, converges absolutely in the strong sense that it still converges when $\rho^{T}$ and $F$ are replaced by their absolute values. The curly bracket $\{\cdots\}$ of Definition~\ref{4.6:def-strips-bracket} has the polymer-gas representation \eqref{4.6:eq-polymer-gas-a} of Lemma~\ref{4.6:lem-polymer-gas}. The identity of this chapter that expresses the logarithm of a polymer gas as its series of connected terms then gives the second relation of \eqref{4.6:eq-polymer-logarithm-a}, $\{\cdots\}=\exp\sum_{X}\mathbf{R}'^{(k)}(X)$, the sum running over unions $X$ of $M$-cubes. This is the counterpart of \cite[(1.98)]{B-CMP122b}. The absolute convergence assumed here and the logarithm identity are proved later in this chapter; together with Lemma~\ref{4.6:lem-polymer-gas} they give the second relation of \eqref{4.6:eq-polymer-logarithm-a}. Here that relation serves only to explain why the terms defined by the first relation of \eqref{4.6:eq-polymer-logarithm-a} are called the logarithm of the polymer gas.

The division into two groups is that of \cite[p.~390]{B-CMP122b}:
\begin{quote}
To the first group we assign all the terms with the localization domains $X$ intersecting the large field region $Z_k^{\sim}\cup\bigcup_{i=1}^{m}Y_i^{\sim}$, to the second group the terms with the domains disjoint with this region
\end{quote}
In the quoted sentence $m$ is the number of regions $Y_i$, and the tilde denotes the enlargement by one layer of cells of side $M\check{R}_k$, the same operation as in \cite[(2.30), (2.41)]{B-CMP119}. It is therefore the operation $S\mapsto S^{\boxplus}$ of Definition~\ref{4.6:def-strips-bracket}, not the strip $S\mapsto S^{\sharp}$. This is why the group split is made with respect to the set $|\mathcal{O}|^{\boxplus}$ defined by the third relation of \eqref{4.6:eq-polymer-logarithm-a}.

\begin{remark}[Modified coupling profiles and their support]\label{4.6:rem-coupling-support}
We work inside Ba{\l}aban's operation $\mathbf{R}$ at step $k$, with the step datum, the union $Z$
and an admissible family $\mathfrak{X}$ of islands $X_a$ as in
Definition~\ref{4.6:def-step-datum}. The island integral operations $\mathbf{T}'_k(X_a)$ and the
interaction terms $V(Y)$ are those of Definition~\ref{4.6:def-strips-bracket}. Position-dependent
couplings enter these objects at two places.

First, they enter the exponent of $\mathbf{T}'_k(X_a)$ through \cite[(1.1)--(1.48)]{B-CMP122b}.
The pivot of that computation is the identity \cite[(1.37), p.~365]{B-CMP122b}:
\begin{align*}
&-g_k^{-2}A\bigl(\zeta_0,U_{k,Z}(V'_kV_\Lambda)\bigr)-A\bigl((g''_k(\cdot))^{-2},U''_k\bigr)
 +g_k^{-2}A(\zeta_0,U_0)\\
&\quad=-A\bigl((g_k(\cdot))^{-2},U_k\bigr)
 -A\bigl(((g''_k(\cdot))^{-2}-g_k^{-2})\zeta,U''_{k,Z}\bigr)
 -A\bigl((g''_k(\cdot))^{-2}\zeta_1,U''_{k,Z}\bigr)\\
&\qquad-\tfrac12\sum_p\eta^4\zeta(x(p))\,\bigl|(DH''_{1,k,Z}B)(p)\bigr|^2+O(1).
\end{align*}
All letters in this identity are those of \cite{B-CMP122b}; in particular $\zeta,\zeta_0,\zeta_1$
are Ba{\l}aban's cut-off functions there, not the complex source. The term $O(1)$ is, in
Ba{\l}aban's words, ``the sum of all the small terms''.

Second, they enter the stage terms of \cite[(1.62), p.~189]{B-CMP122a}. There
$\mathbb{C}^{(n+1)}_k$ is the sum of the logarithm of the integral in
\cite[(1.60)]{B-CMP122a} and the term
\begin{equation*}
A\Bigl(\frac{1}{(g^{(n+1)}_k(\cdot))^{2}}-\frac{1}{(g^{(n)}_k(\cdot))^{2}},\,
U^{(n+1)}_k\Bigr).
\end{equation*}
The symbols $\mathbb{C}^{(n+1)}_k$, $g^{(n)}_k(\cdot)$ and $U^{(n+1)}_k$ denote Ba{\l}aban's stage
objects of \cite[\S1]{B-CMP122a}.
These stage terms build $\mathbb{B}''_k$, and through it they enter $V(Y)$ and $\mathbf{T}'_k$.

On an admissible family, all of these couplings are products of the coupling profiles
$\varphi_i$ of the sequence of large-field regions in force; this is the product rule for the
couplings on an admissible family. Write $Z''_i$ for the level-$i$
large-field regions of the prepared sequence and of the stage sequences. Write
$\operatorname{dist}_\infty$ for the sup-norm distance.

The prepared and the stage sequences are based on partitions into $M$-cubes at each of their
scales. Ba{\l}aban describes them as ``a new
sequence of the large field regions \dots\ the complements of these sets form an admissible
sequence of domains based on partitions into $M$-cubes in the corresponding scales''
\cite[p.~178]{B-CMP122a}. At every level $i\le k$ of these sequences the profiles $\varphi''_i$
of $\{Z''_i\}$ are built on the level-$i$
cells of the cell hull $[Z''_i]_i$, with ramps of width $\tfrac14 a_i$. Consequently
\begin{equation}\label{4.6:eq-coupling-support-a}
\operatorname{supp}\bigl((g''_k)^{-2}-g_k^{-2}\bigr)\subset\mathcal{D}
:=\bigcup_{i\le k}\bigl\{x:\operatorname{dist}_\infty\bigl(x,[Z''_i]_i\bigr)<\tfrac14 a_i\bigr\}.
\end{equation}
In \cite{B-CMP122b} the profiles $\varphi''$ sit instead on the partitions into the region's own
cubes, and ``The second term on the right-hand side is localized in $\Lambda$ as a function of the
field $U''_k$'' \cite[p.~361]{B-CMP122b}. Here $\Lambda=\Lambda^{(1.73)}$ is $Z''_k$ plus one layer
of $M$-cubes, ``a rectangular parallelepiped contained in a cube of the size $100M$''
\cite[p.~191]{B-CMP122a}. The set $\mathcal{D}$ compares with these regions as follows:
\begin{equation*}
\mathcal{D}\not\subset\Lambda^{(1.73)},\qquad \mathcal{D}\subset X_a .
\end{equation*}
On an admissible family, every occurrence of $\Lambda$ as a support in
\cite[(1.14)--(1.48)]{B-CMP122b} is to be read as $\mathcal{D}$; this is how the representation
and bound of the large-field terms uses the Wilson-action terms with position-dependent couplings
of \cite[(1.14)--(1.48)]{B-CMP122b} at the regions $\{Z''_i\}$.

Likewise, consider the weight difference of \cite[(1.62)]{B-CMP122a} at stage $n+1$, and write
$j+1$ for the level of the stage sequence attached to that stage; it exceeds by $n+1$ the level
attached to stage zero. This weight
difference is supported on the $\tfrac14 a_{j+1}$-collar of $[Z''_{j+1}]_{j+1}$, instead of on
the layer of own cubes used in \cite{B-CMP122a}.

The change enters the bounds of the step only at two places. One is the bound on the island
operations over an admissible family: \cite[(1.73)--(1.75), (1.79)--(1.89)]{B-CMP122b} consume
\cite[(1.1)--(1.48)]{B-CMP122b}. The other is the strip bound on the interaction terms:
\cite[(1.68)]{B-CMP122b} consumes the stage bounds \cite[(1.70)]{B-CMP122b} through
\cite[(1.59)--(1.69)]{B-CMP122b}. No other estimate of the step depends on a profile.

The remaining changes of notation on an admissible family are of the same kind as the one
treated above. For each of them we state what is changed and at which places in the bounds of
the step it is used.

\emph{The completion terms.} After the operation $\mathbf{R}$ at step $k$, the small-field
action over the healed region is completed. In Ba{\l}aban's words: ``Now we add all the missing
terms with localization domains intersecting the domain $Z$, more exactly we add these terms to
complete the effective action \dots\ but these will be completed yet. Denote this action by
$A'_k$'' \cite[pp.~378--379]{B-CMP122b}.

On an admissible family the added terms, namely the $E$-terms and the $R''$-terms in the notation
of \cite{B-CMP122b}, are by definition the members of the family of effective-action terms on the
whole lattice, by the completion items of the definition of that family. Their
negatives are carried by those interaction terms $V(Y)$ that receive the corresponding
brackets. For the subtracted action terms of \cite[pp.~378--379]{B-CMP122b} in particular,
whether such a term is present in $V(Y)$ is decided by the components of $Z$ that $Y$ contains.

The change enters at two places:
\begin{itemize}
\item inside the definition of $V(Y)$, and hence inside the inputs of the strip bound on the
interaction terms, which are \cite[(1.59)--(1.69)]{B-CMP122b};
\item in the completed action $A'_k$ outside the curly bracket $\{\cdots\}$ of
Definition~\ref{4.6:def-strips-bracket}.
\end{itemize}

\emph{The members carrying a profile.} Some members of $V(Y)$, and some members of the exponent
of $\mathbf{T}'_k(X_a)$, carry a coupling profile. These are of two kinds.

First, the couplings of \cite[(1.70), p.~377]{B-CMP122b}. The coupling $1/(g_k(\cdot))^2$ occurs
in the term $A'_k\bigl(1/(g_k(\cdot))^2,U_k\bigr)$ outside the integral. The coupling
$1/(g''_k(\cdot))^2$ occurs in the term
\begin{equation*}
-A\bigl((1/(g''_k(\cdot))^2)\,\zeta_1,\cdot\bigr)
\end{equation*}
inside the integral.

Second, the counterterms $-\beta_iA(h_z\varphi_i,\cdot)$, where $h_z$ is the weight of these
counterterms in Ba{\l}aban's construction. These carry the profiles
$\varphi_i$ of the sequence of large-field regions in force.

Two clauses of the lemma on the coupling profiles locate these members. Write
$\varphi^{(X_a)^c}$ for the profile built from the complement of $X_a$. By clause (f) of that
lemma,
\begin{equation}\label{4.6:eq-coupling-support-2}
\bigl\{x:\varphi^{(X_a)^c}(x)\neq1\bigr\}\subset X_a^{\sharp}.
\end{equation}
By clause (e), the ramp of $\varphi_i$ has radius less than $\tfrac14 a_i$. Hence the members of
$V(Y)$ that carry a profile depend, among the islands, on exactly those $X_a$ with
$X_a^{\sharp}\subset Y$. This agrees with the locality of the interaction terms
(Remark~\ref{4.6:rem-interaction-locality}, display \eqref{4.6:eq-interaction-locality-a}),
which is stated in \cite[p.~379]{B-CMP122b} without proof and is cited as printed.

As weights, the profiles take values in $[0,1]$. The support of $1-\varphi_i$ is placed as
follows:
\begin{itemize}
\item at the levels whose large-field regions are unions of level-$i$ cells, it lies inside the
one-cell layer of \cite[\S1]{B-CMP122b}, by clauses (a) and (b) of the lemma on the coupling
profiles, together with the statement on the one-cell layer among the changes of notation on an
admissible family;
\item at the levels of the prepared and the stage sequences, which are based on partitions into
$M$-cubes as described above, it is the collar of \eqref{4.6:eq-coupling-support-a}.
\end{itemize}

The change enters at two places:
\begin{itemize}
\item inside the inputs of the bound on the island operations over an admissible family and of
the strip bound on the interaction terms, where the profiles act as weights;
\item in the polymer-gas representation of the bracket
(Lemma~\ref{4.6:lem-polymer-gas}(a)), through the locality of the interaction terms
(Remark~\ref{4.6:rem-interaction-locality}, which is stated in \cite[p.~379]{B-CMP122b} without
proof and is cited as printed).
\end{itemize}

\emph{The vacuum constants.} In \cite[(1.70)--(1.72)]{B-CMP122b} a constant independent of the
field is assigned to the same place as the term that produces it. Precisely, take a vacuum
value $E^{(i)}(X',1,z')$ or $R^{(i)}(X',1)$.
\begin{itemize}
\item If $X'\subset X_a^{\sharp}$, the value is assigned to $E_k(X_a)$.
\item Otherwise, it is assigned to the interaction term $V(Y)$ that receives the bracket of
$(X',z')$.
\end{itemize}
At $i=k$ the $R$-value is that of $R''^{(k)}$ only, since these terms ``will be completed
yet'' \cite[pp.~378--379]{B-CMP122b}. The value $\mathbf{R}'^{(k)}(X,1)$ is shifted after
\cite[(1.98)]{B-CMP122b}, term by term. Finally, $E_k(\Lambda\cap X)$ denotes the count term of
\cite[(1.10)]{B-CMP122b} on $\Lambda\cap X$, the letter $\Lambda$ being used here as in
\cite[\S1]{B-CMP122b}. The notation $E^{(i)}$, $R^{(i)}$, $z'$,
$R''^{(k)}$ and $E_k$ is that of \cite{B-CMP122b}.

This change enters at three places.
\begin{itemize}
\item The difference $E_k(X_a)-E_k(\Lambda\cap X_a)$ occurs in the exponent of
$\mathbf{T}'_k(X_a)$. It is the normalisation in the bound on the island operations over an
admissible family.
\item The constants assigned to the interaction terms change the prefactor $\alpha$ of the
strip bound on the interaction terms. The bound \cite[(1.68)]{B-CMP122b} is kept, with $\alpha$
replaced by $Ce^{-\frac12\beta_s\kappa\check{R}_k}$, where $C$ is a constant, $\beta_s$ is the
decay fraction of the lemma that places the vacuum constants and $\kappa$ is the tree-decay rate
of \cite{B-CMP119}. This quantity is at most $\tfrac12\alpha$
under a ceiling on $\gamma$, and, by the lemma that places the vacuum constants,
$\tfrac12\beta_s\kappa\ge\kappa(d)$, where $\kappa(d)$ is a decay rate depending on the
dimension $d$ alone, distinct from the rate $\kappa$.
\item The shift
\begin{equation*}
\mathbf{R}'^{(k)}(X,\cdot)-\mathbf{R}'^{(k)}(X,1)
\end{equation*}
is applied to the logarithm terms $\mathbf{R}'^{(k)}(X)$ of
Definition~\ref{4.6:def-polymer-logarithm}, after the logarithm of the polymer gas has been
formed.
\end{itemize}

\emph{The strip convention.} In \cite[(1.68)--(1.72), (1.90)--(1.100)]{B-CMP122b} the
enlargements marked by a tilde are replaced by strips $S^{\sharp}$, that is, by $S$ plus
$m_{\sharp}$ layers of $M$-cubes. This replacement governs the following objects, as written in
Definitions~\ref{4.6:def-strips-bracket}, \ref{4.6:def-polymer-expansion},
\ref{4.6:def-tree-lengths} and~\ref{4.6:def-polymer-logarithm}:
\begin{itemize}
\item the strips themselves;
\item the interaction domains $\mathcal{Y}(\mathfrak{X})$;
\item the cores $\operatorname{core}(X')$;
\item the lengths $d_{\sharp}$ and $d_{\mathrm{out}}$;
\item the cube counts $N_a$.
\end{itemize}
The replacement is not made at the split into the two groups, where the enlargement
$S^{\boxplus}$ is kept.

The cost of the strip convention lies entirely inside the strip bound on the interaction terms.
It consists of two items:
\begin{itemize}
\item the factor
\begin{equation*}
e^{-\frac12\beta_{\sharp}\kappa(m_{\sharp}-1)}
\end{equation*}
in the prefactor $\alpha^{\sharp}$, where $\beta_{\sharp}$ is the strip decay fraction of that
bound;
\item the floor condition of that bound.
\end{itemize}
The remaining estimates of the step use, of the strips, only the separation property recorded in
Definition~\ref{4.6:def-strips-bracket}. In particular, that property states that the strip of a
box of cells is a box of $M$-cubes, and that strips of components sharing no point share no point.
\end{remark}

\begin{proof}
\emph{The support inclusion.} On $Z$ the telescoping identity \cite[(2.24)]{B-CMP119}, written
for the profiles $\varphi''_i$ of the sequence $\{Z''_i\}$ (the telescoping property of coupling
profiles on an admissible family),
expresses $(g''_k(x))^{-2}$ as $g_k^{-2}$ plus the increments $\beta_i$, $i\le k$, each weighted
by $1-\varphi''_i(x)$. The coupling subtracted in the identity \cite[(1.37)]{B-CMP122b} is this
constant $g_k^{-2}$ \cite[p.~365]{B-CMP122b}. Hence, for $x\in Z$,
\begin{equation}\label{4.6:eq-coupling-support-1}
(g''_k(x))^{-2}-g_k^{-2}=\sum_{i\le k}\beta_i\bigl(1-\varphi''_i(x)\bigr).
\end{equation}
The $i$-th summand on the right of \eqref{4.6:eq-coupling-support-1} vanishes wherever
$\varphi''_i(x)=1$. The profile $\varphi''_i$ is built on the level-$i$
cells of $[Z''_i]_i$ with a ramp of width $\tfrac14 a_i$. It therefore equals $1$ at every $x$
with $\operatorname{dist}_\infty(x,[Z''_i]_i)\ge\tfrac14 a_i$. Hence the support of the $i$-th
summand lies in the $i$-th set of the union in \eqref{4.6:eq-coupling-support-a}. Taking the
union over $i\le k$ gives \eqref{4.6:eq-coupling-support-a}.

\emph{The size of the collars.} At the levels whose cells have side at most $M$, the collar that
$\mathcal{D}$ adds round $Z''_i$ is of order $M$. At the top levels, where $a_i>M$, it is
cell-wide; these levels run up to $a_k=M\check{R}_k$. At $i=k$ the set $\mathcal{D}$ contains
whole level-$k$ cells of $[Z''_k]_k$, each of side $M\check{R}_k$, together with their
quarter-cell collars. Each such cell with its collar is a cube of volume
\begin{equation*}
\bigl(\tfrac32 M\check{R}_k\bigr)^{d}.
\end{equation*}

\emph{Comparison with $\Lambda^{(1.73)}$.} In \cite{B-CMP122b} the profiles are built on the
partitions into the region's own cubes, and, in the sentence quoted above from
\cite[p.~361]{B-CMP122b}, ``the second term on the right-hand side'' is localized in
$\Lambda=\Lambda^{(1.73)}$ as a function of the field $U''_k$. That set is $Z''_k$
plus one layer of $M$-cubes, a rectangular parallelepiped contained in a cube of size $100M$.
By contrast, the level-$k$ part of $\mathcal{D}$ rounds $Z''_k$ up to whole level-$k$ cells, each
of side $M\check{R}_k$, and adds to each a quarter-cell collar, whereas $\Lambda^{(1.73)}$ exceeds
$Z''_k$ by one layer of $M$-cubes only. The cells so added are not contained in
$\Lambda^{(1.73)}$; this is the relation $\mathcal{D}\not\subset\Lambda^{(1.73)}$.

\emph{Containment in the island.} The old large-field region inside the island $X_a$
(Definition~\ref{4.6:def-step-datum}) sits inside $X_a$ by ``at least ten
layers of $M\check{R}_k$-cubes'' \cite[\S1]{B-CMP122a}. That is, it lies at least ten cells inside
$X_a$. The regions $Z''_i$ lie in that old large-field region. Every point of the cell hull
$[Z''_i]_i$ is at sup-distance at most one level-$i$ cell side $a_i$ from $Z''_i$, and the collar
of $\mathcal{D}$ adds less than $\tfrac14 a_i$. Since $a_i\le a_k=M\check{R}_k$, each set of the
union in \eqref{4.6:eq-coupling-support-a} lies within sup-distance less than
$\tfrac54 M\check{R}_k$ of the old large-field region. This is less than the ten cells that
separate that region from the complement of $X_a$; hence $\mathcal{D}\subset X_a$.

\emph{The stage weight difference.} At stage $n+1$ the weight difference of
\cite[(1.62)]{B-CMP122a} is carried by the profile at the level $j+1$ of the stage sequence. The
argument of the first paragraph, applied to that profile alone, places its support in the
$\tfrac14 a_{j+1}$-collar of $[Z''_{j+1}]_{j+1}$, in place of the layer of own cubes.
\end{proof}

\begin{definition}[Standing bounds on island operations and interaction terms]\label{4.6:def-standing-bounds}
Fix the level $k$ of the step. The step datum, the admissible island families $\mathfrak{X}$ and the
histories of their islands are those of Definition~\ref{4.6:def-step-datum}. The strips $S^{\sharp}$,
the admissible interaction domains $\mathcal{Y}(\mathfrak{X})$, the island operations $\mathbf{T}'_k$
and the interaction terms $V(Y)$ are those of Definition~\ref{4.6:def-strips-bracket}. The tree
length $d_{\sharp}(Y)$ relative to the strip union $Z(\mathfrak{X})^{\sharp}$ is that of
Definition~\ref{4.6:def-tree-lengths}, and the complex domains $\mathfrak{U}_k(\cdot)$ of
configurations $(\mathbf{U},\mathbf{J})$ are those of Definition~\ref{4.6:not-complex-domain}. We say
that the \emph{standing bounds hold at level $k$} if the following three statements hold.
\begin{enumerate}
\item[(a)] Let $X$ be a first-class island admissible alone, in the sense of
Definition~\ref{4.6:def-step-datum}, and let $\Phi$ be a bounded function of
the integration variables of $X$ and of other variables. Then, for all
$(\mathbf{U},\mathbf{J})\in\mathfrak{U}_k(X^{\sharp})$, pointwise in the other variables and
uniformly in the history of $X$,
\begin{equation}\label{4.6:eq-standing-bounds-a}
\begin{gathered}
\bigl|\mathbf{T}'_k(X,(\mathbf{U},\mathbf{J}))\Phi\bigr|
\le e^{\mathfrak{s}}\,\bigl(\mathbf{T}'_k(X,(U,0))1\bigr)\,\sup_X|\Phi| ,
\qquad
\mathfrak{s}:=c_\sigma\sup_\omega|\sigma(\omega)|\le 1,\\[2pt]
\mathbf{T}'_k(X,(U,0))1\le \exp\bigl(-2(1+\beta_0)^{-1}p_0(g_k)\bigr),\\[2pt]
p_T:=2(1+\beta_0)^{-1}p_0(g_k)-\mathfrak{s}\ \ge\ 2(1+\beta_0)^{-1}p_0(g_k)-1 .
\end{gathered}
\end{equation}
In the first line, $\sup_X$ is the supremum over the integration variables of $X$ on the supports
of the characteristic functions entering $\mathbf{T}'_k$. The function $\sigma(\omega)$ and the
configuration $(U,0)$ are those of \cite[(1.73)--(1.75)]{B-CMP122b}. The constant
$c_\sigma\in\{1,3\}$ is the constant in the exponent of whichever of the bounds
\cite[(1.73)]{B-CMP122b}, \cite[(1.75)]{B-CMP122b} is used. In the second line, $\beta_0$ is the
constant of \cite{B-CMP122b} entering the exponent of \cite[(1.89)]{B-CMP122b}, and $p_0(\cdot)$
is the degree function.

\item[(b)] Let $\mathfrak{X}$ be an admissible family and let $Y\in\mathcal{Y}(\mathfrak{X})$. Then,
uniformly in the histories, in the first-class integration variables on the supports of the
characteristic functions, and in $(\mathbf{U},\mathbf{J})\in\mathfrak{U}_k(Y)$,
\begin{equation}\label{4.6:eq-standing-bounds-b}
|V(Y)|\le \alpha\,\exp\bigl(-(1+2\beta)\kappa\, d_{\sharp}(Y)\bigr).
\end{equation}
Here $\alpha=\alpha^{\mathrm{fam}}$ can be made arbitrarily small by taking $\gamma$ small
enough. The number $\beta\in\,]0,\tfrac14]$ is the decay fraction of \cite{B-CMP122b}, and
$\kappa$ is the tree-decay rate of \cite[\S2]{B-CMP119}.

\item[(c)] For every island $X_a$ of an admissible family $\mathfrak{X}$, every analytic $\Phi$ and
every $Y\in\mathcal{Y}(\mathfrak{X})$, the maps
\begin{equation}\label{4.6:eq-standing-bounds-c}
\begin{gathered}
(\mathbf{U},\mathbf{J})\longmapsto \mathbf{T}'_k(X_a,(\mathbf{U},\mathbf{J}))\Phi
\quad\text{on } \mathfrak{U}_k(X_a^{\sharp}),\\[2pt]
(\mathbf{U},\mathbf{J})\longmapsto V(Y,(\mathbf{U},\mathbf{J}))
\quad\text{on } \mathfrak{U}_k(Y)
\end{gathered}
\end{equation}
are analytic. The bounds (a) and (b) hold on these domains.
\end{enumerate}
\end{definition}

The statements (a)--(c) are the forms, on an admissible family, of bounds of Ba{\l}aban.

The first line of \eqref{4.6:eq-standing-bounds-a} is the bound of
\cite[(1.73)--(1.75)]{B-CMP122b}. There it reads
\begin{equation*}
|\mathbf{T}'(X,(\mathbf{U},\mathbf{J}))F|
\le\bigl(\mathbf{T}'(X,(U,0))1\bigr)\,e^{3\sup|\sigma|}\,\sup|F| .
\end{equation*}
It holds pointwise in the other variables because the density of the operation at $(U,0)$ is
positive, by the positivity dichotomy of the large-field operations. When Ba{\l}aban applies it
in \cite[p.~388]{B-CMP122b}, he estimates $|\sigma|$ in
\cite[(1.73)]{B-CMP122b} by $1$. The condition $\mathfrak{s}\le 1$ in
\eqref{4.6:eq-standing-bounds-a} is the corresponding requirement on $c_\sigma\sup_\omega|\sigma(\omega)|$,
the supremum being taken over all arguments $\omega$ of $\sigma$. That this requirement is met for
$\gamma$ small enough is the content of the bound on the phase $\sigma$ for small $\gamma$, which
enters here by statement.

The second line of \eqref{4.6:eq-standing-bounds-a} is \cite[(1.89)]{B-CMP122b}, read with the
integer parameter in the assumption of that bound equal to zero. An island admissible alone
satisfies that assumption with this value of the parameter, and the constant
$2(1+\beta_0)^{-1}$ in the exponent is the one printed there. The derivation of
\cite[(1.89)]{B-CMP122b} uses the bound for the interaction term on Ba{\l}aban's enlargement of the
island, which \cite[p.~379]{B-CMP122b} states without proof to be the same bound with a different
constant; it is cited as printed.

The third line of \eqref{4.6:eq-standing-bounds-a} defines $p_T$. Its lower bound is the condition
$\mathfrak{s}\le 1$ of the first line, subtracted from $2(1+\beta_0)^{-1}p_0(g_k)$. Substituting the
second line into the first gives, under the hypotheses of (a),
\begin{equation}\label{4.6:eq-standing-bounds-1}
\bigl|\mathbf{T}'_k(X,(\mathbf{U},\mathbf{J}))\Phi\bigr|
\le e^{\mathfrak{s}}\,e^{-2(1+\beta_0)^{-1}p_0(g_k)}\sup_X|\Phi|
= e^{-p_T}\sup_X|\Phi| .
\end{equation}

The bound \eqref{4.6:eq-standing-bounds-b} is \cite[(1.68)]{B-CMP122b}. There it is stated for
interaction domains not contained in Ba{\l}aban's enlarged large-field region, with
$d_{k,Z}(Y)$ in place of $d_{\sharp}(Y)$, and with $\alpha$ ``arbitrarily small for $\gamma$ small
enough''. On an admissible family the region $Z$ is replaced by the strip union
$Z(\mathfrak{X})^{\sharp}$, and this replacement is what $d_{\sharp}(Y)$ records.

The analyticity (c) is that of \cite[(1.65)]{B-CMP122b} for the operation $\mathbf{T}'_k$ and of
\cite[(1.68)--(1.69)]{B-CMP122b} for the interaction terms.

We call the hypothesis of the joint induction of the step construction at all steps $j\le k$ the
\emph{induction hypothesis below level $k$}; it is a single induction hypothesis. It enters only
inside the statements supplying (a)--(c) and in one further estimate of this chapter, the bound
on the support of the coupling excess on an admissible family, where it is invoked under this
name. None of the inequalities carried by the induction hypothesis below level $k$ is used in the
estimates of this chapter that precede that bound.

\begin{definition}[One-sided conditions on the constants]\label{4.6:def-constant-conditions}
Let $\beta\in\left]0,\tfrac14\right]$ be the decay fraction and $\beta_0$ the exponent parameter of
\cite{B-CMP122b}, and let $\kappa$ be the tree-decay rate of \cite{B-CMP119}. Let $\alpha$ be the
prefactor of the bound on the interaction terms \cite[(1.68)]{B-CMP122b},
let $g_k$ be the coupling at level $k$, let
$p_0(\cdot)$ be the degree function, and let $\check{R}_k$ be the island scale at level $k$, measured
in cells. Let $p_T$ and $\mathfrak{s}$ be as in the standing bounds on island operations and
interaction terms (Definition~\ref{4.6:def-standing-bounds}), so that
\begin{equation*}
p_T = 2(1+\beta_0)^{-1}p_0(g_k)-\mathfrak{s},
\qquad 0\le \mathfrak{s}\le 1 .
\end{equation*}
Put
\begin{align*}
A_d &:= 5^d\bigl[(1+3^d)\log 2+2d\log 3\bigr], &
C_D &:= 4\cdot 3^d e^{2A_d},\\
c_1(d) &:= 3+5\sqrt{d}, &
K_{\mathrm{KP}}(d) &:= 4e^{2(A_d+5^d)},
\end{align*}
and
\begin{equation*}
\bar{c}:=\max\bigl\{\alpha^{1/3},\,e^{-p_0(g_k)}\bigr\}.
\end{equation*}
The constants satisfy the \emph{one-sided conditions on the constants} (at level $k$) if the
following six inequalities hold:
\begin{equation}\label{4.6:eq-constant-conditions-a}
\tfrac12\beta\kappa\ \ge\ A_d+5^d+\log 2 ,
\end{equation}
\begin{equation}\label{4.6:eq-constant-conditions-b}
\alpha^{1/3}\ \le\ \min\Bigl\{\exp\bigl(-(1+\beta)\kappa\cdot 2\sqrt{d}-\tfrac14\beta\kappa\bigr),\
\frac{\beta\kappa}{24\cdot 5^d\,C_D}\Bigr\},
\end{equation}
\begin{equation}\label{4.6:eq-constant-conditions-c}
-p_T+c_1(d)\,\kappa\,(101\check{R}_k)^d\ \le\ -\tfrac32\,p_0(g_k),
\end{equation}
\begin{equation}\label{4.6:eq-constant-conditions-d}
100^d\exp\bigl(-\tfrac12p_0(g_k)\bigr)\ \le\ \frac{\beta\kappa}{16\cdot 5^d},
\end{equation}
\begin{equation}\label{4.6:eq-constant-conditions-e}
3^d\,K_{\mathrm{KP}}(d)\,e^{(1+\frac12\beta)\kappa\cdot 2\sqrt{d}}\;\bar{c}\ \le\ 1,
\end{equation}
\begin{equation}\label{4.6:eq-constant-conditions-f}
\beta\kappa\bigl(5^{-d}\check{R}_k^{\,d}-1\bigr)\ \ge\ 2\log K_{\mathrm{KP}}(d).
\end{equation}
\end{definition}

Each of the six conditions is one-sided. None of them fixes the value of a constant, none of them
is a cap in the sense of Definition~\ref{2.3:def-cap}, and none is a condition on $L$.

The condition \eqref{4.6:eq-constant-conditions-a} is a lower bound on $\kappa$, once $\beta$ is
fixed, by a quantity depending on $d$ only. It plays the role of the requirement in
\cite[p.~389]{B-CMP122b} that $\kappa$ be large enough for $(6\cdot 2^d)^{-1}\beta\kappa\ge 1$. The
condition
\eqref{4.6:eq-constant-conditions-b} is an upper bound on $\alpha$, in the role of the smallness of
$\alpha$ assumed in \cite[p.~389]{B-CMP122b}. It is a condition on $\gamma$: by the bound on the
interaction terms, the prefactor $\alpha$ is arbitrarily small for $\gamma$ small.

The condition \eqref{4.6:eq-constant-conditions-c} is the form taken, for the families of islands
considered in this chapter, by the smallness assumption on $g_k$ displayed in
\cite[p.~389]{B-CMP122b}, namely
\begin{equation*}
-2(1+\beta_0)^{-1}p_0(g_k)+1+2\kappa(100R_k)^d\ \leqq\ -\tfrac32\,p_0(g_k),
\end{equation*}
where $R_k$ is the island scale of that paper. It can hold only if $\beta_0<\tfrac13$. To see this,
insert the
expression for $p_T$ and rearrange \eqref{4.6:eq-constant-conditions-c} into the equivalent form
\begin{equation*}
\Bigl(\frac{2}{1+\beta_0}-\frac32\Bigr)p_0(g_k)\ \ge\ c_1(d)\,\kappa\,(101\check{R}_k)^d+\mathfrak{s}.
\end{equation*}
Since $\kappa>0$, $\check{R}_k>0$ and $\mathfrak{s}\ge 0$, the right-hand side is positive. Since
$p_0(g_k)>0$ in the small-coupling regime, the bracket must therefore be positive, that is,
$2(1+\beta_0)^{-1}>\tfrac32$, which is equivalent to $\beta_0<\tfrac13$. The value
$\beta_0=\tfrac17$ proposed in \cite[p.~389]{B-CMP122b} satisfies this.
Moreover, \eqref{4.6:eq-constant-conditions-c} reduces to the conjunction of two conditions: an
inequality between exponents that does not involve $L$, namely that the exponent $p_0$ of the
degree function $p_0(\cdot)$ exceed $d$ times the exponent governing the growth of the island
scale, which is weaker than the requirement that $p_0$ exceed $d+2$ times that exponent; and a
smallness condition on $\gamma$ that is polynomial in $L$, met by choosing $\gamma$ after $L$.
The condition is carried here in the form \eqref{4.6:eq-constant-conditions-c}, as stated.

The condition \eqref{4.6:eq-constant-conditions-d} is the smallness of $g_k$ used at the same
place in \cite[p.~389]{B-CMP122b}, where it reads
\begin{equation*}
\exp\bigl(-\tfrac12p_0(g_k)\bigr)\,100^dR_k^{-d}\ <\ \exp\bigl(-\tfrac12p_0(g_k)\bigr)
\ \leqq\ \frac{1}{6\cdot 2^d}\,\beta\kappa
\end{equation*}
for $g_k$ small. The condition \eqref{4.6:eq-constant-conditions-e} is an upper bound
on $\bar{c}$, and hence a condition on $\gamma$: its entry $\alpha^{1/3}$ is arbitrarily small for
$\gamma$ small by the bound on the interaction terms, and its entry $e^{-p_0(g_k)}$ is small for
$g_k$ small.

The condition \eqref{4.6:eq-constant-conditions-f} is a lower bound on $\check{R}_k$. Under
\eqref{4.6:eq-constant-conditions-a} it follows from $\check{R}_k^{\,d}\ge 6\cdot 5^d$. Indeed,
\begin{equation*}
2\log K_{\mathrm{KP}}(d)=2\log 4+4(A_d+5^d)=4\bigl(A_d+5^d+\log 2\bigr).
\end{equation*}
By \eqref{4.6:eq-constant-conditions-a} we have $\beta\kappa\ge 2(A_d+5^d+\log 2)$, and therefore
\begin{equation*}
\frac{2\log K_{\mathrm{KP}}(d)}{\beta\kappa}\ \le\ 2\ <\ 5 .
\end{equation*}
If $\check{R}_k^{\,d}\ge 6\cdot 5^d$, then $5^{-d}\check{R}_k^{\,d}-1\ge 5$, and
\eqref{4.6:eq-constant-conditions-f} follows on multiplying by $\beta\kappa$. The bound
$\check{R}_k^{\,d}\ge 6\cdot 5^d$ is met in two independent ways, and the definition does not
prescribe which of them is used:
\begin{itemize}
\item by the lower bound $\check{R}_k\ge L$ on the island scale; hence the bound holds for
every $L$ beyond a threshold depending on $d$ only;
\item for $\gamma$ small, by the lower bound $\check{R}_k\ge(\log\check{x}_k)^{r}$ on the island
scale, where $r$ is the exponent of that bound and $\check{x}_k$ is the inverse-coupling variable
at level $k$; the right-hand side tends to infinity as $\gamma\to0$.
\end{itemize}

The constants are chosen in the following order. First, $d$, $\beta$ and $\beta_0$ are fixed. Next,
$\kappa$ is chosen after $M$, subject to \eqref{4.6:eq-constant-conditions-a}; $\kappa$ grows with
$M$, as in \cite[(0.26)]{B-CMP109}. No condition is placed on $L$. Finally, $\gamma$ is chosen last,
and \eqref{4.6:eq-constant-conditions-b}, \eqref{4.6:eq-constant-conditions-c},
\eqref{4.6:eq-constant-conditions-d} and \eqref{4.6:eq-constant-conditions-e} are conditions to be
met by this choice. The same holds for \eqref{4.6:eq-constant-conditions-f} when it is met in the
second of the two ways above.

The regime condition of \cite[p.~178]{B-CMP122a} does not appear among the six conditions. It
enters through the standing bounds of Definition~\ref{4.6:def-standing-bounds}.

\begin{lemma}[Cubes met by a tree graph]\label{4.6:lem-tree-cubes}
Let $\Gamma\subset\mathbb{R}^d$ be a tree graph (Definition~\ref{4.6:def-cubes-cells}) of length
$\ell M$, where $\ell\ge 0$ is real.
The number of closed $M$-cubes $c_n=Mn+[0,M]^d$ meeting $\Gamma$ is at most $5^d(\ell+1)$:
\begin{equation}\label{4.6:eq-tree-cubes-1}
\#\{\,n\in\mathbb{Z}^d : c_n\cap\Gamma\neq\emptyset\,\}\le 5^d(\ell+1).
\end{equation}
In particular, with the lengths of Definition~\ref{4.6:def-tree-lengths}
(see \eqref{4.6:eq-tree-lengths-a}):
\begin{enumerate}
\item for $Y\in\mathcal{Y}(\mathfrak{X})$ with $d_{\sharp}(Y)<\infty$,
\begin{equation*}
N\bigl(Y\setminus Z(\mathfrak{X})^{\sharp}\bigr)\le 5^d\bigl(d_{\sharp}(Y)+1\bigr);
\end{equation*}
\item for a union $X$ of polymers,
\begin{equation*}
N\bigl(\operatorname{core}(X)\bigr)\le 5^d\bigl(d_{\mathrm{out}}(X)+1\bigr);
\end{equation*}
\item for a union $X$ of $M$-cubes,
\begin{equation*}
N(X)\le 5^d\bigl(d_k(X)+1\bigr).
\end{equation*}
This holds for either reading of the clause defining $d_k(X)$ in \cite{B-CMP109},
since in both the realising tree graph meets every $M$-cube of $X$.
\end{enumerate}
\end{lemma}

\begin{proof}
Put
\begin{equation*}
\mathcal{N}:=\{\,n\in\mathbb{Z}^d : c_n\cap\Gamma\neq\emptyset\,\}.
\end{equation*}
Since $\Gamma$ is a finite union of segments, it is bounded, and $\mathcal{N}$ is finite.
Choose greedily a maximal subset $\mathcal{N}'\subset\mathcal{N}$ such that
$|n-n'|_\infty\ge 3$ for all distinct $n,n'\in\mathcal{N}'$.
Every $n\in\mathcal{N}$ then has some $n'\in\mathcal{N}'$ with $|n-n'|_\infty\le 2$.
If $n\in\mathcal{N}'$, take $n'=n$. If $n\notin\mathcal{N}'$, maximality forbids adjoining $n$ to
$\mathcal{N}'$, so some $n'\in\mathcal{N}'$ satisfies $|n-n'|_\infty\le 2$.
For a fixed $n'$ there are exactly $5^d$ points $n\in\mathbb{Z}^d$ with $|n-n'|_\infty\le 2$. Hence
\begin{equation}\label{4.6:eq-tree-cubes-2}
|\mathcal{N}|\le 5^d\,|\mathcal{N}'|.
\end{equation}

For each $n\in\mathcal{N}'$ pick a point $x_n\in c_n\cap\Gamma$.
Let $n,n'\in\mathcal{N}'$ be distinct. Some coordinate $\mu$ has $|n_\mu-n'_\mu|\ge 3$, and
\begin{equation*}
x_{n,\mu}\in[Mn_\mu,\,M(n_\mu+1)],\qquad x_{n',\mu}\in[Mn'_\mu,\,M(n'_\mu+1)].
\end{equation*}
The two intervals lie at distance at least $3M-M$ from one another. Therefore
\begin{equation*}
|x_{n,\mu}-x_{n',\mu}|\ge 3M-M=2M,
\end{equation*}
and so $|x_n-x_{n'}|\ge 2M$ in the Euclidean norm.

Suppose $|\mathcal{N}'|\ge 2$.
By the previous paragraph, the open balls $B(x_n,M)$, $n\in\mathcal{N}'$, are pairwise disjoint.
Fix $n\in\mathcal{N}'$. The set $\Gamma$ is a connected finite union of segments. It contains $x_n$, and
it contains a point outside the closed ball $\bar B(x_n,M)$, namely any $x_{n'}$ with $n'\neq n$,
whose distance from $x_n$ is at least $2M>M$.

Let $\Gamma_n$ be the connected component of $\Gamma\cap\bar B(x_n,M)$ containing $x_n$.
It is a finite union of segments joining $x_n$ to the sphere $\partial B(x_n,M)$.
Otherwise $\Gamma_n$ would lie in the open ball $B(x_n,M)$, and $\Gamma_n$ and its complement in
$\Gamma$ would separate $\Gamma$, contradicting the connectedness of $\Gamma$.
A connected finite union of segments contains a polygonal curve between any two of its points. So
$\Gamma_n$ contains a rectifiable curve from the centre $x_n$ to a point of the sphere. Any such curve
has length at least $M$, and hence $\Gamma_n$ has length at least $M$.

A segment meets a sphere in at most two points. Hence the length of $\Gamma$ inside
$\bar B(x_n,M)$ equals its length inside the open ball $B(x_n,M)$, and this length is at least $M$.
Summing over the pairwise disjoint open balls $B(x_n,M)$, $n\in\mathcal{N}'$, gives
\begin{equation*}
\ell M\ge |\mathcal{N}'|\,M.
\end{equation*}

Thus $|\mathcal{N}'|\le\ell$ whenever $|\mathcal{N}'|\ge 2$, while otherwise $|\mathcal{N}'|\le 1$.
In all cases
\begin{equation*}
|\mathcal{N}'|\le\max\{1,\ell\}\le\ell+1.
\end{equation*}
Together with \eqref{4.6:eq-tree-cubes-2} this yields $|\mathcal{N}|\le 5^d(\ell+1)$, which is
\eqref{4.6:eq-tree-cubes-1}.

For the particular cases, apply \eqref{4.6:eq-tree-cubes-1} to shortest tree graphs realising the
lengths of Definition~\ref{4.6:def-tree-lengths}.
\begin{itemize}
\item In the first case the tree graph has length $M d_{\sharp}(Y)$; it exists because
$d_{\sharp}(Y)<\infty$.
\item In the second case it has length $M d_{\mathrm{out}}(X)$.
\item In the third case it has length $M d_k(X)$.
\end{itemize}
By definition these tree graphs meet every $M$-cube of $Y\setminus Z(\mathfrak{X})^{\sharp}$, of
$\operatorname{core}(X)$ and of $X$, respectively. Each of
$N(Y\setminus Z(\mathfrak{X})^{\sharp})$, $N(\operatorname{core}(X))$ and $N(X)$ is therefore at most the
number of $M$-cubes met by the corresponding tree graph. The estimate \eqref{4.6:eq-tree-cubes-1},
with $\ell$ equal to $d_{\sharp}(Y)$, $d_{\mathrm{out}}(X)$ and $d_k(X)$ respectively, gives the three
bounds.
\end{proof}

\begin{lemma}[A serpentine path through a box of cubes]\label{4.6:lem-serpentine-path}
Let $d\ge 1$, let $n_1,\dots,n_d\ge 1$ be integers, and let $B$ be a box of $n_1\times\cdots\times n_d$
$M$-cubes in the sense of Definition~\ref{4.6:def-cubes-cells}, that is, the union of the cubes
$c_\nu$ with $\nu\in\bar n+\prod_{\mu=1}^{d}\{0,\dots,n_\mu-1\}$ for some $\bar n\in\mathbb{Z}^d$.
In this lemma only, write $N:=\prod_{\mu=1}^{d}n_\mu$ for the number of cubes of $B$ (not the $N$
of the gauge group $SU(N)$). Then $B$ contains a polygonal path through the centres of all its cubes,
consecutive cubes sharing a face, of length $(N-1)M$, and this path is contained in $B$.
\end{lemma}

\begin{proof}
The centre of $c_\nu$ is $M\nu+\tfrac{M}{2}(1,\dots,1)$. Two cubes $c_\nu$, $c_{\nu'}$ share a face
exactly when $\nu-\nu'$ or $\nu'-\nu$ is a unit coordinate vector, and then their centres are at
distance $M$. We
produce an
enumeration $c^{(1)},\dots,c^{(N)}$ of the cubes of $B$ in which $c^{(i)}$ and $c^{(i+1)}$ share a
face for $1\le i\le N-1$; the polygonal path joining the successive centres then passes through
the centres of all cubes of $B$ and consists of $N-1$ segments of length $M$, so its length is
$(N-1)M$. We argue by induction on $d$.

For $d=1$ the cubes $c_{\bar n},c_{\bar n+1},\dots,c_{\bar n+n_1-1}$ are consecutive intervals of
the line, each sharing an endpoint with the next, and the path is the row of their centres.

Let $d\ge 2$. For $j=0,\dots,n_d-1$ the cubes $c_\nu$ of $B$ with $\nu_d=\bar n_d+j$ form the
layer $j+1$ of $B$, which is $B'\times\bigl(M(\bar n_d+j)+[0,M]\bigr)$, where $B'$ is the
$(d-1)$-dimensional box of $n_1\times\cdots\times n_{d-1}$ $M$-cubes whose cubes are indexed by
$\bar n'+\prod_{\mu=1}^{d-1}\{0,\dots,n_\mu-1\}$, with $\bar n':=(\bar n_1,\dots,\bar n_{d-1})$.
By the induction hypothesis, $B'$ has an enumeration $P$ of its
$N':=n_1\cdots n_{d-1}$ cubes, consecutive cubes sharing a face. Taking the product of each cube of
$P$ with $M(\bar n_d+j)+[0,M]$ gives an enumeration of layer $j+1$ in which consecutive cubes again
share a face (the product of their common $(d-2)$-dimensional face with that interval). For a cube
$c'$ of $B'$, call $c'\times\bigl(M(\bar n_d+j)+[0,M]\bigr)$ its lift to layer~$j+1$. Run
layer~$1$ along $P$; step from the lift to layer~$1$ of the last cube of $P$ to its lift to
layer~$2$, which shares a face with it, their centres being at distance $M$; run
layer~$2$ along $P$ reversed, ending at the lift of the first cube of $P$; step from it to the lift
of that cube to layer~$3$; run layer~$3$ along $P$; and so on, alternating directions up to
layer~$n_d$. Every cube
of $B$ occurs exactly once, consecutive cubes share a face, and the length is
\begin{equation*}
n_d(N'-1)M+(n_d-1)M=(n_dN'-1)M=(N-1)M .
\end{equation*}

Finally, the segment between the centres of two face-adjacent cubes passes through the centre of
their common face, and each of its two halves lies inside one of the two cubes. Hence every
segment of the path, and so the whole path, is contained in $B$.
\end{proof}

\begin{lemma}[Touch-connected families of cubes through a cube]\label{4.6:lem-connected-families}
Let $c$ be an $M$-cube and let $s \ge 1$ be an integer. The number of touch-connected families
$\mathcal{S}$ of $M$-cubes with $c \in \mathcal{S}$ and $|\mathcal{S}| = s$ is at most
$9^{d(s-1)}$.
\end{lemma}

\begin{proof}
We use the $M$-cubes, touching and the touch-graph of
Definition~\ref{4.6:def-cubes-cells}. The vertices of the touch-graph are the $M$-cubes. Two
distinct $M$-cubes are joined by an edge when they touch. A family $\mathcal{S}$ is
touch-connected when the subgraph induced on $\mathcal{S}$ is connected.

The $M$-cube $c_n = Mn + [0,M]^d$ touches $c_{n'}$, $n' \ne n$, exactly when
$\max_{1 \le i \le d} |n'_i - n_i| \le 1$. Hence every $M$-cube has
\begin{equation*}
\Delta := 3^d - 1
\end{equation*}
neighbours in the touch-graph.

Let $\mathcal{S} \ni c$ be touch-connected with $|\mathcal{S}| = s$. Since the induced
subgraph on $\mathcal{S}$ is connected, it has a spanning tree. This tree has $s-1$ edges,
and each of them joins two touching cubes of $\mathcal{S}$. Fix one such tree for each
$\mathcal{S}$.

A depth-first traversal of the tree from $c$ crosses every edge of the tree exactly twice,
once in each direction. It is therefore a closed walk in the touch-graph that starts and
ends at $c$ and has $2(s-1)$ steps. Its vertex set is the vertex set of the tree, namely
$\mathcal{S}$. (For $s = 1$ this is the walk of zero steps at $c$.)

Since $\mathcal{S}$ is the vertex set of its walk, $\mathcal{S}$ is determined by the walk.
The assignment of a walk to each family is therefore injective. Consequently the number of
families in question is at most the number of walks of $2(s-1)$ steps in the touch-graph
that start at $c$. At each step such a walk moves to one of the $\Delta$ neighbours of its
current cube, so there are at most $\Delta^{2(s-1)}$ such walks.

Finally, $\Delta = 3^d - 1 < 3^d$ gives $\Delta^2 < 9^d$. Hence
\begin{equation*}
\Delta^{2(s-1)} \le 9^{d(s-1)},
\end{equation*}
with strict inequality when $s \ge 2$. This proves the bound.
\end{proof}

The counting of interaction domains used in \cite[p.~389]{B-CMP122b}, where it is carried out in the usual way, is made explicit in the following lemma.

\begin{lemma}[Anchored interaction domains and their weighted sum]\label{4.6:lem-anchored-domains}
Fix an admissible family $\mathfrak{X}$. For $Y\in\mathcal{Y}(\mathfrak{X})$, an \emph{anchor} of $Y$
is a pair $(c,S)$ in which $c$ is an $M$-cube of $Y\setminus Z(\mathfrak{X})^{\sharp}$ and
$S\subset Y$ is a strip of $Z(\mathfrak{X})^{\sharp}$ touched by $c$. The cube $c$ is then called an
\emph{anchor cube} of $Y$. Let $A_d$ be the constant of \eqref{4.6:eq-constant-conditions-a}.
\begin{enumerate}
\item[($\alpha$)] Every $Y\in\mathcal{Y}(\mathfrak{X})$ with $d_{\sharp}(Y)<\infty$ has an anchor.
\item[($\beta$)] Let $c$ be an $M$-cube and let $D\in\mathbb{N}=\{0,1,2,\dots\}$. Then
\begin{equation}\label{4.6:eq-anchored-domains-1}
\#\bigl\{Y\in\mathcal{Y}(\mathfrak{X}) : c \text{ is an anchor cube of } Y,\
\lfloor d_{\sharp}(Y)\rfloor=D\bigr\}\le 2e^{A_d(D+2)}.
\end{equation}
\item[($\gamma$)] Assume the condition \eqref{4.6:eq-constant-conditions-a} on $\beta\kappa$. Then
every union $W$ of $M$-cubes satisfies
\begin{equation}\label{4.6:eq-anchored-domains-a}
\sum_{\substack{Y\in\mathcal{Y}(\mathfrak{X}),\ Y\subset W\\ d_{\sharp}(Y)<\infty}}
e^{-\frac12\beta\kappa\, d_{\sharp}(Y)}\le C_D\,N\bigl(W\setminus Z(\mathfrak{X})^{\sharp}\bigr),
\qquad C_D=4\cdot 3^d e^{2A_d}.
\end{equation}
\end{enumerate}
\end{lemma}

\begin{proof}
Throughout, write $Z^{\sharp}=Z(\mathfrak{X})^{\sharp}$. This set is the union of the strips of the
components of $Z(\mathfrak{X})$.

\emph{($\alpha$).} Let $Y\in\mathcal{Y}(\mathfrak{X})$ with $d_{\sharp}(Y)<\infty$. Such a $Y$
contains a whole strip $S$ of $Z^{\sharp}$ and an $M$-cube not contained in $Z^{\sharp}$. By
\eqref{4.6:eq-geometric-facts-b}, $Y$ is a touch-connected union of closed $M$-cubes. Hence there
is a chain $c_{(0)},c_{(1)},\dots,c_{(r)}$ of $M$-cubes of $Y$ in which consecutive cubes touch,
$c_{(0)}\subset S$ and $c_{(r)}\not\subset Z^{\sharp}$. Let $c_{(i)}$ be the first cube of the chain
that is not contained in $S$; then $i\ge1$ and $c_{(i-1)}\subset S$. Suppose $c_{(i)}$ lay in a strip
$S'\neq S$ of $Z^{\sharp}$. Since $c_{(i)}$ touches $c_{(i-1)}$, the strips $S$ and $S'$ would share
a point. This is excluded by \eqref{4.6:eq-strips-bracket-b}, by which distinct strips share no
point. So $c_{(i)}$ lies in no strip, that is, $c_{(i)}$ is a cube of $Y\setminus Z^{\sharp}$. It
touches $S$ through $c_{(i-1)}$, and therefore $(c_{(i)},S)$ is an anchor of $Y$.

\emph{($\beta$).} Fix $c$ and $D$, and let $Y$ belong to the set on the left of
\eqref{4.6:eq-anchored-domains-1}. By Definition~\ref{4.6:def-tree-lengths} and
\eqref{4.6:eq-tree-lengths-a}, there is a shortest tree graph $\Gamma_Y\subset Y$ meeting every
$M$-cube of $Y\setminus Z^{\sharp}$. Its length is
\begin{equation*}
M d_{\sharp}(Y)<(D+1)M,
\end{equation*}
and, since $c$ is a cube of $Y\setminus Z^{\sharp}$, it meets $c$ in particular. Let $\mathcal{S}$
be the set of $M$-cubes meeting $\Gamma_Y$. Then $c\in\mathcal{S}$. By Lemma~\ref{4.6:lem-tree-cubes},
applied with $\ell=d_{\sharp}(Y)<D+1$, we have $|\mathcal{S}|\le 5^d(D+2)$.

We show that $\mathcal{S}$ is touch-connected. The curve $\Gamma_Y$ is connected, and it is covered
by the closed cubes of $\mathcal{S}$. Suppose $\mathcal{S}$ split into two non-empty families
$\mathcal{S}_1,\mathcal{S}_2$ such that no cube of $\mathcal{S}_1$ touches a cube of $\mathcal{S}_2$.
The unions of the two families would then be disjoint closed sets, being finite unions of closed
cubes. They would cover $\Gamma_Y$, and each would meet $\Gamma_Y$. This contradicts the
connectedness of $\Gamma_Y$.

Next, $Y$ is determined by the pair formed by $Y\setminus Z^{\sharp}$ and the set of strips of
$Z^{\sharp}$ contained in $Y$. Indeed, the decomposition \eqref{4.6:eq-polymer-gas-b} of
Lemma~\ref{4.6:lem-polymer-gas} holds for interaction domains by \eqref{4.6:eq-strips-bracket-a}:
$Y$ is the disjoint union of $Y\setminus Z^{\sharp}$ and whole strips. Now consider the two
components of the pair.
\begin{itemize}
\item Every cube of $Y\setminus Z^{\sharp}$ meets $\Gamma_Y$, so $Y\setminus Z^{\sharp}$ is a union of
cubes of $\mathcal{S}$.
\item Every strip contained in $Y$ touches $Y\setminus Z^{\sharp}$. To see this, run the chain
argument of ($\alpha$) starting from that strip. The strip therefore touches a cube of
$\mathcal{S}$.
\item A cube touches or equals at most $3^d$ cubes, and each of these lies in at most one strip,
because strips share no point by \eqref{4.6:eq-strips-bracket-b}. Hence at most $3^d|\mathcal{S}|$
strips are candidates.
\end{itemize}

Put $s=|\mathcal{S}|$. By Lemma~\ref{4.6:lem-connected-families} there are at most $9^{d(s-1)}$
choices of $\mathcal{S}$. Given $\mathcal{S}$, there are at most $2^s$ choices of
$Y\setminus Z^{\sharp}\subset\mathcal{S}$ and at most $2^{3^d s}$ choices of the set of strips. Put
$q=2\cdot 9^d\cdot 2^{3^d}$; note that $q\ge2$ and that $5^d(D+2)$ is an integer. Summing over
$s$ gives
\begin{equation}\label{4.6:eq-anchored-domains-2}
\#\{Y\}\le\sum_{s=1}^{5^d(D+2)}9^{d(s-1)}\,2^s\,2^{3^d s}
\le\sum_{s=1}^{5^d(D+2)}q^s\le 2\,q^{5^d(D+2)}.
\end{equation}
The last step uses that a geometric sum of ratio $q\ge2$ is at most twice its last term:
\begin{equation*}
\sum_{s=1}^n q^s\le q^n\,\frac{q}{q-1}\le 2q^n .
\end{equation*}
Finally,
\begin{equation*}
5^d\log q=5^d\bigl[\log2+2d\log3+3^d\log2\bigr]=5^d\bigl[(1+3^d)\log2+2d\log3\bigr]=A_d ,
\end{equation*}
so the right side of \eqref{4.6:eq-anchored-domains-2} equals $2e^{A_d(D+2)}$. This proves
\eqref{4.6:eq-anchored-domains-1}.

\emph{($\gamma$).} Let $Y$ be a term of the sum in \eqref{4.6:eq-anchored-domains-a}. By ($\alpha$),
$Y$ has an anchor $(c,S)$. Its cube $c$ is a cube of $Y\setminus Z^{\sharp}$, and hence of
$W\setminus Z^{\sharp}$ because $Y\subset W$. Its strip $S$ is one of the at most $3^d$ strips
touching $c$, by the count in ($\beta$). Put $D=\lfloor d_{\sharp}(Y)\rfloor$. Then
$d_{\sharp}(Y)\ge D$, so $e^{-\frac12\beta\kappa d_{\sharp}(Y)}\le e^{-\frac12\beta\kappa D}$. For
given $(c,S,D)$, the number of such $Y$ is bounded by \eqref{4.6:eq-anchored-domains-1}. Therefore
\begin{align*}
\sum_{Y}e^{-\frac12\beta\kappa\, d_{\sharp}(Y)}
&\le\sum_{c}\ \sum_{S}\ \sum_{D\ge0}2e^{A_d(D+2)}e^{-\frac12\beta\kappa D}\\
&\le N(W\setminus Z^{\sharp})\cdot 3^d\cdot 2e^{2A_d}
\sum_{D\ge0}e^{-(\frac12\beta\kappa-A_d)D},
\end{align*}
where $c$ runs over the $M$-cubes of $W\setminus Z^{\sharp}$ and $S$ over the strips touching $c$.
By \eqref{4.6:eq-constant-conditions-a},
\begin{equation*}
\tfrac12\beta\kappa-A_d\ge 5^d+\log2\ge\log2 .
\end{equation*}
The series over $D$ is therefore at most $\sum_{D\ge0}2^{-D}=2$. Hence
\begin{equation}\label{4.6:eq-anchored-domains-3}
\sum_{Y}e^{-\frac12\beta\kappa\, d_{\sharp}(Y)}\le N(W\setminus Z^{\sharp})\cdot3^d\cdot2e^{2A_d}\cdot2
=C_D\,N(W\setminus Z^{\sharp}),
\end{equation}
which is \eqref{4.6:eq-anchored-domains-a}.
\end{proof}

\begin{lemma}[Polymers whose core contains a given cube]\label{4.6:lem-polymer-count}
Let $c$ be an $M$-cube and let $D\in\mathbb{N}$. Then
\begin{equation}\label{4.6:eq-polymer-count-1}
\#\bigl\{X'\in\mathcal{P}\;:\;c\in\operatorname{core}(X'),\
\lfloor d_{\mathrm{out}}(X')\rfloor=D\bigr\}\;\le\;2e^{A_d(D+2)},
\end{equation}
and every polymer $X'$ in the set on the left satisfies
$N(\operatorname{core}X')\le 5^d(D+2)$.
\end{lemma}

\begin{proof}
The argument is the count in the proof of part $(\beta)$ of the lemma on anchored interaction
domains (Lemma~\ref{4.6:lem-anchored-domains}), combined with part $(c)$ of the polymer-gas
representation of the bracket (Lemma~\ref{4.6:lem-polymer-gas}).

Fix a polymer $X'$ in the set on the left of \eqref{4.6:eq-polymer-count-1}. By
Lemma~\ref{4.6:lem-polymer-gas}$(c)$, $X'$ is the disjoint union of its core and of the
second-class strips contained in $X'$, and each of these strips touches the core. In particular
$X'$ is determined by its core together with the set of second-class strips it contains.

Since $\lfloor d_{\mathrm{out}}(X')\rfloor=D$, the number $d_{\mathrm{out}}(X')$ is finite and
$d_{\mathrm{out}}(X')<D+1$. By the definition of $d_{\mathrm{out}}$
(Definition~\ref{4.6:def-tree-lengths}) there is a tree graph $\Gamma$ contained in $X'$,
of length $M\,d_{\mathrm{out}}(X')<(D+1)M$, which meets every $M$-cube of
$\operatorname{core}(X')$; in particular $\Gamma$ meets $c$, since $c\in\operatorname{core}(X')$.
Let $\mathcal{S}$ be the family of closed $M$-cubes met by $\Gamma$. Then:
\begin{itemize}
\item $c\in\mathcal{S}$;
\item $\mathcal{S}$ is touch-connected, because $\Gamma$ is connected and is covered by the
finitely many closed cubes of $\mathcal{S}$;
\item by the lemma on cubes met by a tree graph (Lemma~\ref{4.6:lem-tree-cubes}), applied with
$\ell=d_{\mathrm{out}}(X')<D+1$, one has $|\mathcal{S}|\le 5^d(\ell+1)\le 5^d(D+2)$;
\item every cube of $\operatorname{core}(X')$ is met by $\Gamma$, so
$\operatorname{core}(X')\subset\mathcal{S}$, and therefore
\begin{equation*}
N(\operatorname{core}X')\;\le\;|\mathcal{S}|\;\le\;5^d(D+2),
\end{equation*}
which is the second assertion of the lemma.
\end{itemize}
The second-class strips inside $X'$ touch $\operatorname{core}(X')\subset\mathcal{S}$, so they
range over at most $3^d|\mathcal{S}|$ candidates.

Thus $X'$ is fixed by the choice of a touch-connected family $\mathcal{S}\ni c$ with
$|\mathcal{S}|\le 5^d(D+2)$, whose number for each size is bounded by the lemma on
touch-connected families through a cube (Lemma~\ref{4.6:lem-connected-families}), of the core as
a subfamily of $\mathcal{S}$, and of the strips among at most $3^d|\mathcal{S}|$ candidates. These
are the data, with the same bounds, of the count in the proof of
Lemma~\ref{4.6:lem-anchored-domains}$(\beta)$; that count gives the bound
$2e^{A_d(D+2)}$, which is \eqref{4.6:eq-polymer-count-1}.
\end{proof}

\begin{definition}[Abstract polymer systems and their logarithm]\label{4.6:def-abstract-polymers}
An \emph{abstract polymer system} consists of the following data.
\begin{itemize}
\item A finite set $\mathcal{P}$, whose elements are called polymers.
\item A reflexive and symmetric relation $\iota\subset\mathcal{P}\times\mathcal{P}$. Two polymers
$X',X''$ with $X'\,\iota\,X''$ are called \emph{incompatible}; they are called \emph{compatible} if
they are not incompatible.
\item A function $F\colon\mathcal{P}\to\mathbb{C}$, the activity.
\end{itemize}
The hard-core factor is $\zeta(X',X'')=0$ if $X'\,\iota\,X''$ and $\zeta(X',X'')=1$ otherwise.
For $n\ge 1$ put $[n]:=\{1,\dots,n\}$, and let $C_n$ be the set of edge sets $g$ of graphs on $[n]$
for which $([n],g)$ is connected. For $(X'_1,\dots,X'_n)\in\mathcal{P}^n$ let
$G(X'_1,\dots,X'_n)$ be the graph on $[n]$ in which $\{p,q\}$ is an edge if and only if $p\neq q$
and $X'_p\,\iota\,X'_q$; write $E(G)$ for its edge set. The \emph{Ursell coefficient} is
\begin{equation}\label{4.6:eq-abstract-polymers-1}
\begin{split}
\rho^{T}(X'_1,\dots,X'_n)
&:=\sum_{g\in C_n}\ \prod_{\{p,q\}\in g}\bigl(\zeta(X'_p,X'_q)-1\bigr)\\
&=\sum_{\substack{g\subset E(G)\\ ([n],g)\ \text{connected}}}(-1)^{|g|}.
\end{split}
\end{equation}
In particular $\rho^{T}=1$ when $n=1$, and $\rho^{T}=0$ when $G$ is disconnected.
We write $\tau(G)$ for the number of spanning trees of $G$; thus $\tau(G)=1$ when $n=1$, and
$\tau(G)=0$ when $G$ is disconnected. The \emph{partition function} of the system is
\begin{equation}\label{4.6:eq-abstract-polymers-2}
\Xi:=\sum_{\substack{\mathfrak{F}\subset\mathcal{P}\\ \text{pairwise compatible}}}\
\prod_{X'\in\mathfrak{F}}F(X'),
\end{equation}
the empty family $\mathfrak{F}=\emptyset$ contributing $1$; since $\mathcal{P}$ is finite, this is
a finite sum. Its \emph{logarithm series} is
\begin{equation}\label{4.6:eq-abstract-polymers-3}
\Phi:=\sum_{n\ge 1}\frac{1}{n!}\sum_{(X'_p)\in\mathcal{P}^n}
\rho^{T}(X'_1,\dots,X'_n)\prod_{p=1}^{n}F(X'_p).
\end{equation}
\end{definition}

The expansion of the logarithm of a polymer partition function into the series
\eqref{4.6:eq-abstract-polymers-3} is the formula \cite[(2.11)--(2.13)]{B-CMP116}, which is
stated there as well known, the convergence theorem being referred to the literature cited in
that paper. The convergence of the exponentiated cluster expansion is proved below in this
section, in the abstract setting fixed above.

\begin{proof}
We verify the second equality in \eqref{4.6:eq-abstract-polymers-1} and the values recorded
after it. Let $p\neq q$. If $\{p,q\}\in E(G)$, then $X'_p\,\iota\,X'_q$, so
$\zeta(X'_p,X'_q)-1=-1$; if $\{p,q\}\notin E(G)$, then $X'_p$ and $X'_q$ are compatible, so
$\zeta(X'_p,X'_q)-1=0$. Hence for $g\in C_n$ the product in the first line of
\eqref{4.6:eq-abstract-polymers-1} vanishes unless $g\subset E(G)$, in which case it equals
$(-1)^{|g|}$. This gives the second line.

For $n=1$ the only graph on $[1]$ has empty edge set and is connected, so the sum has the single
term $(-1)^0=1$. If $G$ is disconnected, no subset $g\subset E(G)$ makes $([n],g)$ connected,
since every path in $([n],g)$ is a path in $G$; the sum is empty and $\rho^{T}=0$. In the same
way, for $n=1$ the graph itself is its unique spanning tree, so $\tau(G)=1$, and a disconnected
graph has no spanning tree, so $\tau(G)=0$.
\end{proof}

\begin{lemma}[The exponential formula for compatible families]\label{4.6:lem-exponential-formula}
Let $(\mathcal{P},\iota,F)$ be an abstract polymer system, with partition function $\Xi$ and
logarithm $\Phi$ as in Definition~\ref{4.6:def-abstract-polymers}. If the double series defining
$\Phi$ converges absolutely, that is, if
\begin{equation}\label{4.6:eq-exponential-formula-1}
\sum_{n\ge 1}\frac{1}{n!}\sum_{(X'_1,\dots,X'_n)\in\mathcal{P}^n}
\Bigl|\rho^{T}(X'_1,\dots,X'_n)\prod_{p=1}^{n}F(X'_p)\Bigr|<\infty,
\end{equation}
then $\exp\Phi=\Xi$.
\end{lemma}

\begin{proof}
Throughout, $\zeta(X',X'')$ is the hard-core factor of
Definition~\ref{4.6:def-abstract-polymers}: $\zeta(X',X'')=0$ if $X'\,\iota\,X''$ and
$\zeta(X',X'')=1$ otherwise. A family $\mathfrak{F}\subset\mathcal{P}$ is called compatible if
its elements are pairwise compatible, so that
$\Xi=\sum_{\mathfrak{F}\text{ compatible}}\prod_{X'\in\mathfrak{F}}F(X')$, the empty family
contributing $1$.

For $z\in\mathbb{C}$ let $\Xi(z)$ and $\Phi(z)$ denote the same objects with $F$ replaced by
$zF$. Since $\mathcal{P}$ is finite, $\Xi(z)$ is a polynomial in $z$. Setting
\begin{equation}\label{4.6:eq-exponential-formula-2}
\Psi_n:=\sum_{(X'_1,\dots,X'_n)\in\mathcal{P}^n}\rho^{T}(X'_1,\dots,X'_n)
\prod_{p=1}^{n}F(X'_p)\qquad(n\ge 1),
\end{equation}
we have $\Phi(z)=\sum_{n\ge1}z^n\Psi_n/n!$. By \eqref{4.6:eq-exponential-formula-1},
$\sum_{n\ge1}|\Psi_n|/n!<\infty$, so this series converges absolutely and uniformly on
$|z|\le 1$. Hence $\Phi$ is analytic on $|z|<1$ and continuous on $|z|\le1$, and
$\Phi(1)=\Phi$, $\Xi(1)=\Xi$.

It suffices to prove $\exp\Phi(z)=\Xi(z)$ as an identity of formal power series in $z$.
Indeed, on $|z|<1$ the exponential of the convergent power series $\Phi(z)$ is given by its
formal exponential, so the two sides agree as functions on $|z|<1$. Both sides are continuous on
$|z|\le1$ ($\Xi$ is a polynomial, and $\exp\Phi$ is the composition of continuous maps), so they
also agree at $z=1$, which is the assertion.

\emph{Step (i).} Fix $n\ge0$ and a sequence $(X'_1,\dots,X'_n)\in\mathcal{P}^n$. Let $K_n$ be
the complete graph on $[n]$, with edge set $E(K_n)$, and for $e=\{p,q\}\in E(K_n)$ write
$\zeta_e:=\zeta(X'_p,X'_q)$. Writing each factor $\zeta_e$ as $1+(\zeta_e-1)$ and expanding the
product gives
\begin{equation}\label{4.6:eq-exponential-formula-3}
\prod_{e\in E(K_n)}\zeta_e=\sum_{g\subset E(K_n)}\prod_{e\in g}(\zeta_e-1).
\end{equation}
Let $\operatorname{Part}[n]$ be the set of partitions of $[n]$ into non-empty blocks (for
$n=0$ it consists of the empty partition alone). To each $g\subset E(K_n)$ associate the
partition $\pi\in\operatorname{Part}[n]$ of $[n]$ into the vertex sets of the connected
components of $([n],g)$, together with the restrictions of $g$ to the blocks $B\in\pi$; each
restriction is a connected graph on its block. Conversely, a partition $\pi$ and a connected
graph on each block determine $g$ as the union of these graphs, whose components are exactly
the blocks. This correspondence $g\leftrightarrow(\pi;\text{a connected graph on each block})$
is therefore a bijection, and the product over $e\in g$ splits over the blocks. For a block $B$
we write $\rho^{T}((X'_p)_{p\in B})$ for $\rho^{T}$ evaluated at the subsequence
$(X'_p)_{p\in B}$ listed in increasing order of $p$; the increasing bijection of $B$ onto
$[|B|]$ carries connected graphs on $B$ to connected graphs on $[|B|]$ and preserves the factors
$\zeta_e-1$, so the sum over connected graphs on $B$ of $\prod_{e}(\zeta_e-1)$ equals
$\rho^{T}((X'_p)_{p\in B})$. Grouping the terms of \eqref{4.6:eq-exponential-formula-3}
accordingly yields
\begin{equation}\label{4.6:eq-exponential-formula-4}
\prod_{1\le p<q\le n}\zeta(X'_p,X'_q)
=\sum_{\pi\in\operatorname{Part}[n]}\prod_{B\in\pi}\rho^{T}\bigl((X'_p)_{p\in B}\bigr).
\end{equation}

\emph{Step (ii).} For every $n\ge0$,
\begin{equation}\label{4.6:eq-exponential-formula-5}
\sum_{(X'_p)\in\mathcal{P}^n}\,\prod_{p<q}\zeta(X'_p,X'_q)\prod_{p=1}^{n}F(X'_p)
=n!\sum_{\substack{\mathfrak{F}\text{ compatible}\\|\mathfrak{F}|=n}}\,
\prod_{X'\in\mathfrak{F}}F(X').
\end{equation}
Indeed, if a sequence has a repeated entry $X'_p=X'_q$ with $p<q$, then $X'_p\,\iota\,X'_q$
because $\iota$ is reflexive, so $\zeta(X'_p,X'_q)=0$ and the term vanishes. If the entries are
distinct, the product of the $\zeta$'s is $1$ when they are pairwise compatible and $0$
otherwise. The sequences of $n$ distinct pairwise compatible polymers are exactly the $n!$
orderings of the compatible families of cardinality $n$, and $\prod_pF(X'_p)$ does not depend
on the ordering. Multiplying \eqref{4.6:eq-exponential-formula-5} by $z^n/n!$, summing over
$n\ge0$ and inserting \eqref{4.6:eq-exponential-formula-4}, we obtain
\begin{equation}\label{4.6:eq-exponential-formula-6}
\Xi(z)=\sum_{n\ge0}\frac{z^n}{n!}\sum_{(X'_p)\in\mathcal{P}^n}\,
\sum_{\pi\in\operatorname{Part}[n]}\prod_{B\in\pi}
\Bigl[\rho^{T}\bigl((X'_p)_{p\in B}\bigr)\Bigr]\prod_{p=1}^{n}F(X'_p),
\end{equation}
all sums over $(X'_p)$ and $\pi$ being finite.

\emph{Step (iii).} Fix $n$ and $\pi\in\operatorname{Part}[n]$. The summand in
\eqref{4.6:eq-exponential-formula-6} is
$\prod_{B\in\pi}\bigl[\rho^{T}((X'_p)_{p\in B})\prod_{p\in B}F(X'_p)\bigr]$, and
$\mathcal{P}^n=\prod_{B\in\pi}\mathcal{P}^{B}$, so the sum over $(X'_p)\in\mathcal{P}^n$
factorises over the blocks. Relabelling $B$ increasingly onto $[|B|]$, under which $\rho^{T}$ is
invariant by its definition recalled in Step (i), the factor of the block $B$ is
\begin{equation*}
\sum_{(X'_p)_{p\in B}\in\mathcal{P}^{B}}\rho^{T}\bigl((X'_p)_{p\in B}\bigr)
\prod_{p\in B}F(X'_p)=\Psi_{|B|},
\end{equation*}
which depends on $|B|$ only. Thus the coefficient of $z^n$ in $\Xi(z)$ is
$\frac{1}{n!}\sum_{\pi\in\operatorname{Part}[n]}\prod_{B\in\pi}\Psi_{|B|}$. Each partition of
$[n]$ into $s$ blocks corresponds to exactly $s!$ ordered tuples $(B_1,\dots,B_s)$ of its blocks,
these being non-empty and hence distinct; and the ordered tuples of disjoint blocks covering
$[n]$ with $|B_c|=m_c$ for $c=1,\dots,s$ number $n!/(m_1!\cdots m_s!)$. Therefore the coefficient
of $z^n$ in $\Xi(z)$ equals
\begin{equation}\label{4.6:eq-exponential-formula-7}
\sum_{s\ge0}\frac{1}{s!}\sum_{\substack{(m_1,\dots,m_s):\ m_c\ge1\\ m_1+\dots+m_s=n}}\,
\prod_{c=1}^{s}\frac{\Psi_{m_c}}{m_c!},
\end{equation}
where only $s\le n$ contribute, and for $n=0$ only the empty tuple ($s=0$), with value $1$.
On the other hand, since $m_c\ge1$ forces $s\le n$, the formal series
\begin{equation*}
\sum_{s\ge0}\frac{1}{s!}\Bigl(\sum_{m\ge1}\frac{z^m\Psi_m}{m!}\Bigr)^{s}=\exp\Phi(z)
\end{equation*}
is well defined, and expanding the $s$-th power shows that its coefficient of $z^n$ is exactly
\eqref{4.6:eq-exponential-formula-7}. Hence $\exp\Phi(z)=\Xi(z)$ as formal power series, and the
lemma follows from the reduction above.
\end{proof}

\begin{lemma}[The connected coefficient is bounded by spanning trees]
\label{4.6:lem-spanning-tree-bound}
In the abstract setting of Definition~\ref{4.6:def-abstract-polymers}, for every $n\ge 1$ and
every $(X'_1,\dots,X'_n)\in\mathcal{P}^n$,
\begin{equation*}
\bigl|\rho^{T}(X'_1,\dots,X'_n)\bigr|\le\tau\bigl(G(X'_1,\dots,X'_n)\bigr).
\end{equation*}
\end{lemma}

\begin{proof}
Write $G:=G(X'_1,\dots,X'_n)$ and $E:=E(G)$. By Definition~\ref{4.6:def-abstract-polymers},
\begin{equation}\label{4.6:eq-spanning-tree-bound-1}
\rho^{T}(X'_1,\dots,X'_n)=\sum_{\substack{g\subset E\\ ([n],g)\ \text{connected}}}(-1)^{|g|}.
\end{equation}
If $G$ is disconnected, both sides of the asserted inequality vanish; if $n=1$, both equal $1$.
We may therefore assume that $G$ is connected and $n\ge 2$. We take the vertex $1$ as root. A
subset $g\subset E$ is identified with the graph $([n],g)$; we call it a connected spanning
subgraph of $G$ when this graph is connected. For such $g$ and $u,v\in[n]$ let
$\operatorname{dist}_g(u,v)$ denote the number of edges of a shortest path from $u$ to $v$ in
$([n],g)$.

\emph{The tree attached to a connected subgraph.} For a connected spanning subgraph $g\subset E$
put $\ell_g(v):=\operatorname{dist}_g(1,v)$, and for $v\neq 1$ let $p_g(v)$ be the smallest, in
the order of $[n]$, among the $g$-neighbours $u$ of $v$ with $\ell_g(u)=\ell_g(v)-1$. This set is
non-empty: $\ell_g(v)\ge 1$ is finite because $g$ is connected, and the vertex preceding $v$ on a
shortest $g$-path from $1$ to $v$ belongs to it. Put
\begin{equation*}
T(g):=\bigl\{\{v,p_g(v)\}: v\in[n],\ v\neq 1\bigr\}.
\end{equation*}
Each of these edges lies in $g\subset E$. They are $n-1$ distinct edges: if
$\{v,p_g(v)\}=\{u,p_g(u)\}$ with $u\neq v$, then $u=p_g(v)$ and $v=p_g(u)$, so
$\ell_g(u)=\ell_g(v)-1$ and $\ell_g(v)=\ell_g(u)-1$, which is impossible. The graph $([n],T(g))$
is connected, since the sequence $v,p_g(v),p_g(p_g(v)),\dots$ decreases $\ell_g$ by one at each
step and reaches $1$ after $\ell_g(v)$ steps. A connected graph on $n$ vertices with $n-1$ edges
is a tree, so $T(g)$ is a spanning tree of $G$ contained in $g$. The path just described has
length $\ell_g(v)$, so $\ell_{T(g)}(v)\le\ell_g(v)$; since $T(g)\subset g$, also
$\ell_g\le\ell_{T(g)}$; hence $\ell_{T(g)}=\ell_g$.

For a spanning tree $T$ of $G$ we write $\ell_T$ and $p_T$ for $\ell_{E(T)}$ and $p_{E(T)}$. In a
tree rooted at $1$, a vertex $v\neq 1$ has exactly one tree-neighbour one level up, its parent,
and all its other tree-neighbours lie one level down; thus $p_T(v)$ is the parent of $v$. In
particular, $p_{T(g)}(v)=p_g(v)$, because $\{v,p_g(v)\}\in T(g)$ and
$\ell_{T(g)}(p_g(v))=\ell_g(v)-1$. Define
\begin{equation}\label{4.6:eq-spanning-tree-bound-2}
\begin{aligned}
E^{*}(T):={}&E(T)\cup\bigl\{\{u,v\}\in E:\ \ell_T(u)=\ell_T(v)\bigr\}\\
&\cup\bigl\{\{u,v\}\in E:\ \ell_T(v)=\ell_T(u)+1,\ u>p_T(v)\bigr\}.
\end{aligned}
\end{equation}

\emph{Claim.} For a connected spanning subgraph $g\subset E$ and a spanning tree $T$ of $G$,
\begin{equation*}
T(g)=T\iff E(T)\subset g\subset E^{*}(T).
\end{equation*}

Proof of $\Rightarrow$. Let $T=T(g)$. Then $E(T)=T(g)\subset g$ by construction. Let
$\{u,v\}\in g\setminus E(T)$. By the triangle inequality for $\operatorname{dist}_g$,
$|\ell_g(u)-\ell_g(v)|\le 1$, and $\ell_g=\ell_T$ by the previous paragraph. If
$\ell_T(u)=\ell_T(v)$, the edge lies in $E^{*}(T)$. Otherwise, after exchanging the names of the
endpoints, $\ell_T(v)=\ell_T(u)+1$. Then $u$ is a $g$-neighbour of $v$ one level up, and $u\neq
p_T(v)$, because $\{v,p_T(v)\}\in E(T)$ while $\{u,v\}\notin E(T)$. Since $p_T(v)=p_g(v)$ is the
smallest $g$-neighbour of $v$ one level up, $u>p_T(v)$, and the edge lies in $E^{*}(T)$. Hence
$g\subset E^{*}(T)$.

Proof of $\Leftarrow$. Let $E(T)\subset g\subset E^{*}(T)$. By
\eqref{4.6:eq-spanning-tree-bound-2}, every edge of $g$ joins two vertices whose $T$-depths differ
by at most one. Along a shortest $g$-path $1=v_0,v_1,\dots,v_m=v$, with $m=\ell_g(v)$, we
therefore have $\ell_T(v_{i+1})\le\ell_T(v_i)+1$, so $\ell_T(v)\le\ell_g(v)$. Since
$E(T)\subset g$, also $\ell_g\le\ell_T$, and hence $\ell_g=\ell_T$. Now let $v\neq 1$. The
$g$-neighbours $u$ of $v$ one level up are those with $\{u,v\}\in g$ and
$\ell_T(v)=\ell_T(u)+1$. Such an edge either lies in $E(T)$, in which case $u=p_T(v)$, the only
tree-neighbour of $v$ one level up, or lies in the third set of
\eqref{4.6:eq-spanning-tree-bound-2}, in which case $u>p_T(v)$; and $p_T(v)$ itself is a
$g$-neighbour of $v$, because $E(T)\subset g$. Hence $p_g(v)=p_T(v)$ for every $v\neq 1$, and
\begin{equation*}
T(g)=\bigl\{\{v,p_T(v)\}: v\neq 1\bigr\}=E(T).
\end{equation*}
This proves the claim.

\emph{Conclusion.} By the claim, the connected spanning subgraphs of $G$ are partitioned, according
to the value of $g\mapsto T(g)$, into the Boolean intervals
\begin{equation*}
[E(T),E^{*}(T)]:=\{g:\ E(T)\subset g\subset E^{*}(T)\},
\end{equation*}
where $T$ runs over the image of $g\mapsto T(g)$, a set of spanning trees of $G$. Every $g$ in such
an interval is connected and spanning, because it contains $E(T)$. Writing $g=E(T)\cup S$ with
$S\subset E^{*}(T)\setminus E(T)$, we have $|g|=n-1+|S|$, so
\eqref{4.6:eq-spanning-tree-bound-1} becomes
\begin{equation*}
\rho^{T}(X'_1,\dots,X'_n)=\sum_{T}(-1)^{n-1}\sum_{S\subset E^{*}(T)\setminus E(T)}(-1)^{|S|}.
\end{equation*}
The inner sum equals $1$ if $E^{*}(T)=E(T)$ and $0$ otherwise, since the alternating sum over all
subsets of a non-empty finite set vanishes. Hence
\begin{equation*}
\rho^{T}(X'_1,\dots,X'_n)=(-1)^{n-1}\,\#\bigl\{T:\ E^{*}(T)=E(T)\bigr\},
\end{equation*}
with $T$ again running over the image of $g\mapsto T(g)$. Since this image consists of spanning
trees of $G$, its cardinality is at most $\tau(G)$, and therefore
\begin{equation*}
\bigl|\rho^{T}(X'_1,\dots,X'_n)\bigr|\le\#\{T\}\le\tau(G).
\end{equation*}
\end{proof}

\begin{lemma}[The rooted-tree majorant under a summability condition]
\label{4.6:lem-rooted-tree-majorant}
Let $\mathcal{P}$, the incompatibility relation $\iota$, the graphs $G(X'_1,\dots,X'_n)$ and the
spanning-tree count $\tau(G)$ be as in Definition~\ref{4.6:def-abstract-polymers}. Let
$\mathsf{w}\colon\mathcal{P}\to[0,\infty[$ and $a\colon\mathcal{P}\to[0,\infty[$ satisfy
\begin{equation}\label{4.6:eq-rooted-tree-majorant-a}
\sum_{X'\in\mathcal{P}\,:\,X'\,\iota\,X''}\mathsf{w}(X')\,e^{a(X')}\;\le\;a(X'')
\qquad\text{for every } X''\in\mathcal{P}.
\end{equation}
Then for every $X''\in\mathcal{P}$ the series
\begin{equation}\label{4.6:eq-rooted-tree-majorant-b}
B(X'')\;:=\;\sum_{n\ge 1}\frac{1}{(n-1)!}
\sum_{(X'_2,\dots,X'_n)\in\mathcal{P}^{n-1}}
\tau\bigl(G(X'',X'_2,\dots,X'_n)\bigr)\prod_{q=2}^{n}\mathsf{w}(X'_q)
\end{equation}
converges, and $B(X'')\le e^{a(X'')}$. Here the term $n=1$ of the series is $1$.
\end{lemma}

\begin{proof}
For a statement $P$ we write $\mathbf{1}\{P\}$ for $1$ if $P$ holds and $0$ otherwise.
Since $\mathcal{P}$ is finite, every inner sum in \eqref{4.6:eq-rooted-tree-majorant-b} is finite.
All its terms are nonnegative.

\emph{Step 1: the tree representation.} For $n\ge1$ let $\mathcal{T}_n$ be the set of trees on the
vertex set $[n]$, each regarded as rooted at the vertex $1$. In particular, $\mathcal{T}_1$ consists
of the one-vertex tree. A spanning tree of a graph $G$ on $[n]$ is exactly a tree $T$ on $[n]$ with
$E(T)\subset E(G)$. Hence
\begin{equation*}
\tau(G)=\sum_{T\in\mathcal{T}_n}\mathbf{1}\{E(T)\subset E(G)\}.
\end{equation*}
For $n=1$ both sides equal $1$. If $G$ is disconnected, both sides vanish.

Take $G=G(X'_1,\dots,X'_n)$. By Definition~\ref{4.6:def-abstract-polymers}, a pair $\{p,q\}$ is an
edge of $G$ if and only if $p\ne q$ and $X'_p\,\iota\,X'_q$. Every edge $\{p,q\}$ of a tree has
$p\ne q$, so
\begin{equation*}
\mathbf{1}\{E(T)\subset E(G)\}=\prod_{\{p,q\}\in E(T)}\mathbf{1}\{X'_p\,\iota\,X'_q\}.
\end{equation*}
Put $X'_1:=X''$ and, for $T\in\mathcal{T}_n$,
\begin{equation}\label{4.6:eq-rooted-tree-majorant-1}
W_T(X''):=\sum_{(X'_2,\dots,X'_n)\in\mathcal{P}^{n-1}}
\;\prod_{\{p,q\}\in E(T)}\mathbf{1}\{X'_p\,\iota\,X'_q\}\prod_{q=2}^{n}\mathsf{w}(X'_q)\;\ge 0 .
\end{equation}
For $n=1$ this is $W_T(X'')=1$: the empty tuple contributes one term, and both products are empty.
Interchanging the finite sums over $T$ and over the tuple gives
\begin{equation*}
B(X'')=\sum_{n\ge1}\frac{1}{(n-1)!}\sum_{T\in\mathcal{T}_n}W_T(X'').
\end{equation*}

\emph{Step 2: reduction to a claim.} For $N\ge1$ let $B_N(X'')$ be the partial sum of this series over
$n\le N$. The term $n=1$ equals $\frac{1}{0!}\,W_T(X'')=1$, so $B_1\equiv1$. We shall show: for every
$N\ge2$ and every $X''\in\mathcal{P}$,
\begin{equation}\label{4.6:eq-rooted-tree-majorant-2}
B_N(X'')\;\le\;\exp\Bigl(\sum_{X'\,\iota\,X''}\mathsf{w}(X')\,B_{N-1}(X')\Bigr).
\end{equation}

Granting \eqref{4.6:eq-rooted-tree-majorant-2}, we prove $B_N\le e^{a}$ on $\mathcal{P}$ for all
$N\ge1$, by induction on $N$.
\begin{itemize}
\item For $N=1$: $B_1=1\le e^{a}$, because $a\ge0$.
\item Suppose $B_{N-1}(X')\le e^{a(X')}$ for every $X'\in\mathcal{P}$. Since $\mathsf{w}\ge0$ and the
exponential is increasing, \eqref{4.6:eq-rooted-tree-majorant-2} and then
\eqref{4.6:eq-rooted-tree-majorant-a} give
\begin{equation*}
B_N(X'')\le\exp\Bigl(\sum_{X'\,\iota\,X''}\mathsf{w}(X')\,e^{a(X')}\Bigr)\le e^{a(X'')}.
\end{equation*}
\end{itemize}
The terms of the series are nonnegative, so $N\mapsto B_N(X'')$ is nondecreasing. It is bounded by
$e^{a(X'')}$. Therefore the series converges and
\begin{equation*}
B(X'')=\sup_N B_N(X'')\le e^{a(X'')}.
\end{equation*}
It remains to prove \eqref{4.6:eq-rooted-tree-majorant-2}.

\emph{Step 3: the root-degree decomposition.} Let $n\ge2$ and $T\in\mathcal{T}_n$. The tree $T$ has
edges, so its root $1$ has some degree $j\ge1$. Deleting the root leaves $j$ subtrees. Their vertex
sets $V_1,\dots,V_j$ form an unordered partition of $\{2,\dots,n\}$ into nonempty blocks. Each $V_c$
contains exactly one child $r_c$ of the root, and carries the tree $T_c$ (the restriction of $T$ to
$V_c$), rooted at $r_c$.

Conversely, take an unordered partition $\{V_1,\dots,V_j\}$ of $\{2,\dots,n\}$ into nonempty blocks.
On each $V_c$ take a tree $T_c$ rooted at a chosen vertex $r_c\in V_c$. Joining $1$ to each $r_c$
gives a tree on $[n]$ whose root has degree $j$. These two constructions are inverse to each other.
Hence trees $T\in\mathcal{T}_n$ with root degree $j$ correspond bijectively to such data.

\emph{Step 4: factorisation of the weight.} For a finite set $V$, a tree $T_V$ on $V$ rooted at
$r\in V$, and $X'\in\mathcal{P}$, define $W_{T_V}(X')$ by the formula
\eqref{4.6:eq-rooted-tree-majorant-1}, with three changes:
\begin{itemize}
\item the summation runs over $(X'_q)_{q\in V\setminus\{r\}}$;
\item the edge product runs over $E(T_V)$;
\item the weight product runs over $q\in V\setminus\{r\}$, with $X'_r:=X'$.
\end{itemize}

Now fix $T\in\mathcal{T}_n$ with its decomposition. The edges of $T$ are the root edges $\{1,r_c\}$
and the edges of the trees $T_c$. By symmetry of $\iota$, the root edge $\{1,r_c\}$ contributes
$\mathbf{1}\{X'_{r_c}\,\iota\,X''\}$. The edges of $T_c$ involve only the variables $X'_q$ with
$q\in V_c$. The weight product factorises as
\begin{equation*}
\prod_{q\ge2}\mathsf{w}(X'_q)
=\prod_{c=1}^{j}\Bigl(\mathsf{w}(X'_{r_c})\prod_{q\in V_c\setminus\{r_c\}}\mathsf{w}(X'_q)\Bigr).
\end{equation*}
The blocks are disjoint, so the variables of different blocks are summed independently. By
distributivity,
\begin{equation}\label{4.6:eq-rooted-tree-majorant-3}
W_T(X'')=\prod_{c=1}^{j}\;\sum_{X'_{r_c}\,\iota\,X''}\mathsf{w}(X'_{r_c})\,W_{T_c}(X'_{r_c}).
\end{equation}

\emph{Step 5: relabelling the blocks.} Let $V$ be a block with $|V|=m$, and fix $r\in V$. Choose a
bijection $V\to[m]$ with $r\mapsto1$. It carries trees on $V$ rooted at $r$ bijectively onto
$\mathcal{T}_m$. It leaves the weight unchanged, because the summation variables attached to the
non-root vertices are dummy variables. Hence
\begin{equation*}
\sum_{T_V\text{ on }V\text{ rooted at }r}W_{T_V}(X')
=\sum_{T'\in\mathcal{T}_m}W_{T'}(X')=:\mathcal{W}_m(X').
\end{equation*}
Summing over the $m$ choices of $r\in V$ gives $m\,\mathcal{W}_m(X')$.

Set
\begin{equation*}
b_m:=\sum_{X'\,\iota\,X''}\mathsf{w}(X')\,\mathcal{W}_m(X')\;\ge0 .
\end{equation*}
Sum \eqref{4.6:eq-rooted-tree-majorant-3} over all trees with root degree $j$, using the bijection of
Step~3. The sums over the roots $r_c$ and over the trees $T_c$ factorise block by block, so
\begin{equation*}
\sum_{\substack{T\in\mathcal{T}_n\\ \deg(1)=j}}W_T(X'')
=\sum_{\{V_1,\dots,V_j\}}\;\prod_{c=1}^{j}|V_c|\,b_{|V_c|}.
\end{equation*}
The sum runs over unordered partitions of $\{2,\dots,n\}$ into $j$ nonempty blocks.

\emph{Step 6: counting partitions.} Fix sizes $m_1,\dots,m_j\ge1$ with $\sum_c m_c=n-1$. The number of
ordered $j$-tuples of disjoint nonempty blocks of these sizes exhausting $\{2,\dots,n\}$ is the
multinomial coefficient $(n-1)!/(m_1!\cdots m_j!)$. The blocks of a partition are nonempty and
disjoint, hence distinct, so each unordered partition into $j$ blocks arises from exactly $j!$
ordered tuples. Therefore
\begin{equation*}
\sum_{\substack{T\in\mathcal{T}_n\\ \deg(1)=j}}W_T(X'')
=\frac{1}{j!}\sum_{\substack{(m_1,\dots,m_j):\ m_c\ge1\\ \sum_c m_c=n-1}}
\frac{(n-1)!}{\prod_c m_c!}\prod_{c=1}^{j}m_c\,b_{m_c}.
\end{equation*}
Summing over $j\ge1$ and dividing by $(n-1)!$ gives
\begin{equation}\label{4.6:eq-rooted-tree-majorant-4}
\frac{1}{(n-1)!}\sum_{T\in\mathcal{T}_n}W_T(X'')
=\sum_{j\ge1}\frac{1}{j!}\sum_{\substack{(m_1,\dots,m_j):\ m_c\ge1\\ \sum_c m_c=n-1}}
\;\prod_{c=1}^{j}\frac{b_{m_c}}{(m_c-1)!}.
\end{equation}
Here only $j\le n-1$ contribute.

\emph{Step 7: summing over $n$.} Sum \eqref{4.6:eq-rooted-tree-majorant-4} over $2\le n\le N$. For
$j\ge1$ and $m_c\ge1$, the condition $1\le\sum_c m_c\le N-1$ reduces to $\sum_c m_c\le N-1$. Add the
term $n=1$, which equals $1$. It is the term $j=0$ of the same expression: the empty tuple satisfies
$\sum_c m_c=0\le N-1$, and its product is empty. Thus
\begin{equation*}
B_N(X'')=\sum_{j\ge0}\frac{1}{j!}
\sum_{\substack{(m_1,\dots,m_j):\ m_c\ge1\\ \sum_c m_c\le N-1}}
\;\prod_{c=1}^{j}\frac{b_{m_c}}{(m_c-1)!}.
\end{equation*}

\emph{Step 8: relaxing the constraint.} The condition $\sum_c m_c\le N-1$ implies $m_c\le N-1$ for
each $c$. Replacing the former by the latter only adds nonnegative terms, since every $b_m\ge0$. Hence
\begin{equation*}
B_N(X'')\le\sum_{j\ge0}\frac{1}{j!}
\Bigl(\sum_{m=1}^{N-1}\frac{b_m}{(m-1)!}\Bigr)^{j}.
\end{equation*}
By Step~1 applied at the point $X'$, the $m$-th term of the series defining $B(X')$ is
$\mathcal{W}_m(X')/(m-1)!$. So
\begin{equation*}
\sum_{m=1}^{N-1}\frac{\mathcal{W}_m(X')}{(m-1)!}=B_{N-1}(X'),
\end{equation*}
and therefore
\begin{equation*}
\sum_{m=1}^{N-1}\frac{b_m}{(m-1)!}
=\sum_{X'\,\iota\,X''}\mathsf{w}(X')\,B_{N-1}(X').
\end{equation*}
Summing the exponential series gives \eqref{4.6:eq-rooted-tree-majorant-2}. This proves the claim and
hence the lemma.
\end{proof}

\begin{remark}[Use of the tree bounds; infinite polymer sets]\label{4.6:rem-tree-bounds-use}
Let $\mathcal{P}$ be a finite set of polymers with activities $F(X')$, $X'\in\mathcal{P}$, and let
$w,a:\mathcal{P}\to[0,\infty[$ satisfy the summability condition
\eqref{4.6:eq-rooted-tree-majorant-a} and, termwise, $|F(X')|\le w(X')$ for every $X'\in\mathcal{P}$. Then
\begin{equation}\label{4.6:eq-tree-bounds-use-1}
\sum_{n\ge1}\frac{1}{n!}\sum_{(X'_1,\dots,X'_n)\in\mathcal{P}^n}
\Bigl|\rho^{T}(X'_1,\dots,X'_n)\prod_{p=1}^{n}F(X'_p)\Bigr|
\le\sum_{X'_1\in\mathcal{P}}w(X'_1)\,B(X'_1)<\infty ,
\end{equation}
with $B(X'_1)$ the series \eqref{4.6:eq-rooted-tree-majorant-b}. Hence the double series of $\Phi$
converges absolutely and $\Xi=\exp\Phi$. For the sub-sums pinned at a polymer, which are estimated
later in this chapter, the spanning-tree bound (Lemma~\ref{4.6:lem-spanning-tree-bound}) and the
rooted-tree majorant (Lemma~\ref{4.6:lem-rooted-tree-majorant}) are used directly.

If $\mathcal{P}$ is infinite (polymers in $\mathbb{R}^d$), but the polymers contained in any bounded
region are finite in number, one applies the exponential formula
(Lemma~\ref{4.6:lem-exponential-formula}) and Lemmas~\ref{4.6:lem-spanning-tree-bound}
and~\ref{4.6:lem-rooted-tree-majorant} to the finite systems
$\mathcal{P}_R$, the set of polymers contained in the box of side $R$ (here $R$ denotes only this box
side), and lets $R\to\infty$: every bound derived from them in what follows is uniform in $R$, and the
pinned sums are monotone in $R$.
\end{remark}

\begin{proof}
By the spanning-tree bound (Lemma~\ref{4.6:lem-spanning-tree-bound}),
$|\rho^{T}(X'_1,\dots,X'_n)|\le\tau(G(X'_1,\dots,X'_n))$, and $|F(X'_p)|\le w(X'_p)$ for each $p$.
Hence the $n$-th term on the left of \eqref{4.6:eq-tree-bounds-use-1} is at most
\begin{equation*}
\frac{1}{n!}\sum_{X'_1\in\mathcal{P}}w(X'_1)
\sum_{(X'_2,\dots,X'_n)\in\mathcal{P}^{n-1}}\tau(G(X'_1,X'_2,\dots,X'_n))\prod_{q=2}^{n}w(X'_q),
\end{equation*}
where the sum over $(X'_1,\dots,X'_n)$ has been split into the sum over the first entry and the sum
over the remaining ones. Since $1/n!\le1/(n-1)!$ and all terms are non-negative, summing over
$n\ge1$ and exchanging the order of summation gives the first inequality of
\eqref{4.6:eq-tree-bounds-use-1}, the inner series being exactly
\eqref{4.6:eq-rooted-tree-majorant-b} at $X''=X'_1$. Under \eqref{4.6:eq-rooted-tree-majorant-a},
the rooted-tree majorant (Lemma~\ref{4.6:lem-rooted-tree-majorant}) gives that each $B(X'_1)$
converges and $B(X'_1)\le e^{a(X'_1)}<\infty$; as $\mathcal{P}$ is finite, the right side of
\eqref{4.6:eq-tree-bounds-use-1} is a finite sum of finite terms. Thus the double series of $\Phi$
converges absolutely, which is the hypothesis of the exponential formula
(Lemma~\ref{4.6:lem-exponential-formula}); that lemma gives $\Xi=\exp\Phi$.
\end{proof}

\begin{remark}\label{4.6:prop-activity-bound}
The statement of this proposition is not written out here; parts of its proof are printed below.
\end{remark}

\begin{lemma}[The activity bound: splitting the decay and summation]\label{4.6:prf-activity-summation}
Assume the hypotheses of Proposition~\ref{4.6:prop-activity-bound}:
\begin{itemize}
\item the standing bounds on island operations and interaction terms,
\eqref{4.6:eq-standing-bounds-a} and \eqref{4.6:eq-standing-bounds-b};
\item the localisation hypothesis of Proposition~\ref{4.6:prop-activity-bound};
\item the standing geometric facts on boxes and strips (of which
\eqref{4.6:eq-geometric-facts-a} and \eqref{4.6:eq-strips-bracket-b} are used below);
\item the conditions on the constants of Definition~\ref{4.6:def-constant-conditions}, namely
the floor on $\kappa$ of Definition~\ref{4.6:def-constant-conditions} and the conditions
\eqref{4.6:eq-constant-conditions-b},
\eqref{4.6:eq-constant-conditions-c} and \eqref{4.6:eq-constant-conditions-d}.
\end{itemize}
Let $X'\in\mathcal{P}$ be a polymer. For a family $\mathfrak{X}$ of islands and a set
$\mathbf{D}$ of interaction domains, write $F(X';\mathfrak{X},\mathbf{D})$ for the summand of
$F(X')$ indexed by $(\mathfrak{X},\mathbf{D})$ in the first step of the proof of
Proposition~\ref{4.6:prop-activity-bound}, $q=|\mathfrak{X}|$ for the number of islands and
$|\mathbf{D}|$ for the number of interaction domains. Let the corresponding bound obtained in that first step, and the claims (i) and (ii) obtained in its second step, hold. Then, uniformly in the outer
variables and for $(\mathbf{U},\mathbf{J})\in\mathfrak{U}_k(X')$,
\begin{equation*}
|F(X')|\le c_1(X')\,e^{-(1+\beta)\kappa\, d_{\mathrm{out}}(X')}\,
e^{-\frac14 a_v N(\operatorname{core}(X'))},
\qquad a_v=\frac{\beta\kappa}{4\cdot 5^d},
\end{equation*}
with $c_1(X')=\alpha^{1/3}$ in Case~A and $c_1(X')=e^{-p_0(g_k)}$ in Case~B of
the corresponding display. In particular the corresponding bound of
Proposition~\ref{4.6:prop-activity-bound} holds.
\end{lemma}

\begin{proof}
Throughout, $X'$ is fixed, $a$ runs over the islands of $\mathfrak{X}$, $X_a^{\sharp}$ is the
strip of the island $X_a$, and $N_a$ is the number of $M$-cubes of $X_a^{\sharp}$, as in
\eqref{4.6:eq-tree-lengths-a}; the $Y_i^{\sharp}$ are the strips of the outer datum. Recall
that $0<\beta\le\frac14$. The condition \eqref{4.6:eq-constant-conditions-b} bounds
$\alpha^{1/3}$ by an exponential with non-positive exponent, so $\alpha\le1$. By its definition
$p_0(g')=A_0(\log g'^{-2})^{p_0}$ in the glossary (Definition~\ref{0.2:not-glossary}), the degree
function $p_0(\cdot)$ is non-negative on the couplings $g_k\le1$ of the flow, so $p_0(g_k)\ge0$.

\emph{Splitting the decay} (this is \cite[(1.93)]{B-CMP122b}). Write
$1+2\beta=(1+\beta)+\frac14\beta+\frac34\beta$. By claim (i) of
the corresponding display,
\begin{equation*}
\sum_{Y\in\mathbf{D}}(1+\beta)\kappa\, d_{\sharp}(Y)\ge(1+\beta)\kappa
\Big[d_{\mathrm{out}}(X')-\sum_{a\in\mathfrak{X}}N_a-2\sqrt d\,(|\mathbf{D}|+q)\Big].
\end{equation*}
Since $\frac14\beta\kappa=a_v5^d$, claim (ii) of the corresponding display, which bounds
$5^d\sum_{Y\in\mathbf{D}}d_{\sharp}(Y)$ from below, multiplied by $a_v$ gives
\begin{align*}
\sum_{Y\in\mathbf{D}}\tfrac14\beta\kappa\, d_{\sharp}(Y)
&=a_v5^d\sum_{Y\in\mathbf{D}}d_{\sharp}(Y)
\ge a_v\Big[N(\operatorname{core}(X'))-\sum_aN_a-5^d|\mathbf{D}|\Big]\\
&=a_vN(\operatorname{core}(X'))-a_v\sum_aN_a-\tfrac14\beta\kappa|\mathbf{D}|.
\end{align*}
Every island has $N_a\ge1$, so $q\le\sum_aN_a$. The $q$ junction charges
$(1+\beta)\kappa\,2\sqrt d$ can therefore be placed one per island inside
$(1+\beta)\kappa\,2\sqrt d\,N_a$. Inserting the two lower bounds into
$\alpha^{|\mathbf{D}|}\exp\bigl(-(1+2\beta)\kappa\sum_{Y\in\mathbf{D}}d_{\sharp}(Y)\bigr)$ and
keeping the part $\frac34\beta\kappa\,d_{\sharp}(Y)$ of each exponent untouched, we obtain
\begin{equation}\label{4.6:eq-activity-summation-1}
\begin{aligned}
\prod_{Y\in\mathbf{D}}\alpha e^{-(1+2\beta)\kappa d_{\sharp}(Y)}
&\le e^{-(1+\beta)\kappa d_{\mathrm{out}}(X')}\,e^{-a_vN(\operatorname{core}(X'))}
\prod_{a\in\mathfrak{X}}e^{[(1+\beta)\kappa(1+2\sqrt d)+a_v]N_a}\\
&\quad\times\prod_{Y\in\mathbf{D}}\Bigl[\alpha e^{(1+\beta)\kappa 2\sqrt d+\frac14\beta\kappa}
\,e^{-\frac34\beta\kappa d_{\sharp}(Y)}\Bigr].
\end{aligned}
\end{equation}
We simplify the three kinds of factors on the right.
\begin{itemize}
\item The island exponent. Since $\beta\le\frac14$ we have $1+\beta\le\frac54$, and
$a_v=\beta\kappa/(4\cdot5^d)\le\kappa/16$. The coefficient of $N_a$ is therefore at most
$\kappa\bigl(\frac54+\frac52\sqrt d+\frac1{16}\bigr)\le c_1(d)\kappa$, where
$c_1(d)=3+5\sqrt d$.
\item The domain prefactor. The first member of \eqref{4.6:eq-constant-conditions-b} and its
consequence read
\begin{equation*}
\alpha^{1/3}\le e^{-(1+\beta)\kappa\,2\sqrt d-\frac14\beta\kappa},\qquad
\alpha e^{(1+\beta)\kappa 2\sqrt d+\frac14\beta\kappa}\le\alpha\cdot\alpha^{-1/3}=\alpha^{2/3}.
\end{equation*}
\item The remaining decay. Since $d_{\sharp}(Y)\ge0$, we have
$e^{-\frac34\beta\kappa d_{\sharp}(Y)}\le e^{-\frac12\beta\kappa d_{\sharp}(Y)}$.
\end{itemize}
Put
\begin{equation*}
a_Y=\alpha^{2/3}e^{-\frac12\beta\kappa d_{\sharp}(Y)},\qquad
S(\mathfrak{X},X')=\sum_{\substack{Y\in\mathcal{Y}(\mathfrak{X}):\ Y\subset X',\\
d_{\sharp}(Y)<\infty}}e^{-\frac12\beta\kappa d_{\sharp}(Y)}.
\end{equation*}
The last exponential factor of the corresponding display is at most
$e^{\alpha S(\mathfrak{X},X')}$. This holds because
$\mathbf{D}\subset\{Y\in\mathcal{Y}(\mathfrak{X}):Y\subset X'\}$ and
$(1+2\beta)\kappa\ge\frac12\beta\kappa$. Insert \eqref{4.6:eq-activity-summation-1}, the three
simplifications and this bound on the last exponential factor into
the corresponding display, together with the island factors $e^{-p_T}$ carried there. We
obtain
\begin{equation}\label{4.6:eq-activity-summation-2}
\begin{aligned}
|F(X';\mathfrak{X},\mathbf{D})|
&\le e^{-(1+\beta)\kappa d_{\mathrm{out}}(X')}\,e^{-a_vN(\operatorname{core}(X'))}\\
&\quad\times\prod_{a\in\mathfrak{X}}e^{-p_T+c_1(d)\kappa N_a}
\prod_{Y\in\mathbf{D}}a_Y\cdot e^{\alpha S(\mathfrak{X},X')}.
\end{aligned}
\end{equation}

\emph{Size of the islands} (this is the first factor of \cite[(1.95)]{B-CMP122b}). By
\eqref{4.6:eq-tree-lengths-a}, $N_a\le(101\check{R}_k)^d$. Condition
\eqref{4.6:eq-constant-conditions-c} then gives
\begin{equation*}
e^{-p_T+c_1(d)\kappa N_a}\le e^{-p_T+c_1(d)\kappa(101\check{R}_k)^d}
\le e^{-\frac32p_0(g_k)}.
\end{equation*}

\emph{The sum over the interaction domains at fixed islands} (this is
\cite[(1.94)]{B-CMP122b}). Apply part~$(\gamma)$ of Lemma~\ref{4.6:lem-anchored-domains}, display
\eqref{4.6:eq-anchored-domains-a}, with $W=X'$. This is allowed since the floor on $\kappa$ of
Definition~\ref{4.6:def-constant-conditions} holds. It
gives
\begin{equation*}
S(\mathfrak{X},X')\le C_DN(X'\setminus Z(\mathfrak{X})^{\sharp})
\le C_DN(\operatorname{core}(X')).
\end{equation*}
The second inequality holds because $Z(\mathfrak{X})^{\sharp}\supseteq\bigcup_iY_i^{\sharp}$.
Define
\begin{equation*}
\mathcal{Y}'=\{Y\in\mathcal{Y}(\mathfrak{X}):Y\subset X',\ d_{\sharp}(Y)<\infty\}.
\end{equation*}
The data with islands $\mathfrak{X}$ have $\mathbf{D}\subset\mathcal{Y}'$. Indeed, a domain $Y$
with $d_{\sharp}(Y)=\infty$ has $V(Y)=0$ by \eqref{4.6:eq-standing-bounds-b}, and so contributes
nothing.

A datum of a single class has $\mathbf{D}=\emptyset$ only if $q=1$ and $X'=X_a^{\sharp}$. There
are two reasons. First, the strips of two islands never touch, by
\eqref{4.6:eq-strips-bracket-b}. Second, a pair with $q=0$ and $\mathbf{D}=\emptyset$ does not
form a class of data, so it does not occur.
This is the case in which, in \cite[p.~389]{B-CMP122b}, $X'$ is one of the island domains and
the smallness comes from the bound \cite[(1.89)]{B-CMP122b}.

All terms are non-negative, so we drop every other constraint on $\mathbf{D}$. By the definition
of $a_Y$ and the bound on $S(\mathfrak{X},X')$,
\begin{equation*}
\sum_{Y\in\mathcal{Y}'}a_Y=\alpha^{2/3}S(\mathfrak{X},X')
\le\alpha^{2/3}C_DN(\operatorname{core}(X')).
\end{equation*}
We distinguish two cases.
\begin{itemize}
\item If $\mathfrak{X}=\{a\}$ and $X'=X_a^{\sharp}$, expanding the product gives
\begin{equation*}
\sum_{\mathbf{D}\subset\mathcal{Y}'}\prod_{Y\in\mathbf{D}}a_Y
=\prod_{Y\in\mathcal{Y}'}(1+a_Y)\le\exp\Bigl(\sum_{Y\in\mathcal{Y}'}a_Y\Bigr)
\le\exp\bigl(\alpha^{2/3}C_DN(\operatorname{core}(X'))\bigr).
\end{equation*}
\item Otherwise $\mathbf{D}\ne\emptyset$. The sum over non-empty $\mathbf{D}$ equals
$\prod_{Y\in\mathcal{Y}'}(1+a_Y)-1$, and $e^x-1\le xe^x$ for $x\ge0$, so
\begin{align*}
\sum_{\emptyset\ne\mathbf{D}\subset\mathcal{Y}'}\prod_{Y\in\mathbf{D}}a_Y
&\le\Bigl(\sum_{Y\in\mathcal{Y}'}a_Y\Bigr)\exp\Bigl(\sum_{Y\in\mathcal{Y}'}a_Y\Bigr)\\
&\le\alpha^{1/3}\bigl[\alpha^{1/3}C_DN(\operatorname{core}(X'))\bigr]
e^{\alpha^{2/3}C_DN(\operatorname{core}(X'))}\\
&\le\alpha^{1/3}\exp\bigl(2\alpha^{1/3}C_DN(\operatorname{core}(X'))\bigr).
\end{align*}
In the last line we used $x\le e^x$ and $\alpha^{2/3}\le\alpha^{1/3}$, which holds because
$\alpha\le1$.
\end{itemize}
Set $c_0(\mathfrak{X},X')=1$ if $\mathfrak{X}=\{a\}$ and $X'=X_a^{\sharp}$, and
$c_0(\mathfrak{X},X')=\alpha^{1/3}$ otherwise. Multiply both cases by
$e^{\alpha S(\mathfrak{X},X')}\le e^{\alpha C_DN(\operatorname{core}(X'))}$ and use
$\alpha\le\alpha^{1/3}$ and $\alpha^{2/3}\le\alpha^{1/3}$. This gives
\begin{equation}\label{4.6:eq-activity-summation-3}
\sum_{\mathbf{D}}\prod_{Y\in\mathbf{D}}a_Y\cdot e^{\alpha S(\mathfrak{X},X')}
\le c_0(\mathfrak{X},X')\,e^{3\alpha^{1/3}C_DN(\operatorname{core}(X'))}
\le c_0(\mathfrak{X},X')\,e^{\frac12a_vN(\operatorname{core}(X'))}.
\end{equation}
The last inequality is the second member of \eqref{4.6:eq-constant-conditions-b}:
$\alpha^{1/3}\le\beta\kappa/(24\cdot5^dC_D)$ gives
$3\alpha^{1/3}C_D\le\beta\kappa/(8\cdot5^d)=\frac12a_v$.

\emph{The sum over the islands} (this is \cite[(1.96)]{B-CMP122b}). The islands of a datum
with domain $X'$ are admissible first-class boxes $X_a$ with
$X_a^{\sharp}\subset\operatorname{core}(X')$, by \eqref{4.6:eq-polymer-gas-b}. Such an $X_a$ is
a box of cells with two properties:
\begin{itemize}
\item since $X_a\subset X_a^{\sharp}\subset\operatorname{core}(X')$, every cell of $X_a$, in
particular its lowest cell, the one with the least corner in the lexicographic order, lies in
$\operatorname{core}(X')$;
\item by \eqref{4.6:eq-geometric-facts-a}, its side lengths lie in $\{1,\dots,100\}$ cells.
\end{itemize}
A cell contains $\check{R}_k^d\ge1$ cubes of side $M$. Hence the number of candidates is at most
\begin{equation*}
100^d\cdot\#\{\text{cells}\subset\operatorname{core}(X')\}
\le100^dN(\operatorname{core}(X'))\,\check{R}_k^{-d}\le100^dN(\operatorname{core}(X')).
\end{equation*}
This is the count of \cite[(1.96)]{B-CMP122b}, written there in Ba{\l}aban's notation.

The terms are non-negative, so we drop the pairwise constraints and the single-class
constraint. A sum over $q$-element sets of candidates is at most $1/q!$ times the sum over
$q$-tuples of candidates. Using the island bound from the size step, we get
\begin{equation*}
\sum_{\mathfrak{X}:\,|\mathfrak{X}|=q}\ \prod_{a\in\mathfrak{X}}e^{-\frac32p_0(g_k)}
\le\frac1{q!}\Bigl(100^dN(\operatorname{core}(X'))\,e^{-\frac32p_0(g_k)}\Bigr)^q.
\end{equation*}

\emph{Case A}, where $X'\supseteq Y_i^{\sharp}$ for some $i$. A set containing $Y_i^{\sharp}$ is
not a first-class strip, so $X'\ne X_a^{\sharp}$. Hence $c_0(\mathfrak{X},X')=\alpha^{1/3}$ for
every datum, and $q\ge0$. Summing over $q$ and using $e^{-\frac32p_0(g_k)}\le e^{-\frac12p_0(g_k)}$,
\begin{align*}
\sum_{q\ge0}\frac1{q!}\Bigl(100^dN(\operatorname{core}(X'))e^{-\frac32p_0(g_k)}\Bigr)^q
&=\exp\bigl(100^dN(\operatorname{core}(X'))e^{-\frac32p_0(g_k)}\bigr)\\
&\le\exp\bigl(100^de^{-\frac12p_0(g_k)}N(\operatorname{core}(X'))\bigr)
\le e^{\frac14a_vN(\operatorname{core}(X'))}.
\end{align*}
The last inequality is \eqref{4.6:eq-constant-conditions-d}.

\emph{Case B}, where $X'\sqcap\bigcup_iY_i^{\sharp}=\emptyset$. Here every datum has $q\ge1$
and $c_0(\mathfrak{X},X')\le1$. As in \cite[p.~389]{B-CMP122b}, we extract the factor
$e^{-p_0(g_k)}$ before the sum. Since $q\ge1$ and $p_0(g_k)\ge0$,
\begin{equation*}
\prod_{a\in\mathfrak{X}}e^{-\frac32p_0(g_k)}
\le e^{-p_0(g_k)}\prod_{a\in\mathfrak{X}}e^{-\frac12p_0(g_k)}.
\end{equation*}
Summing over $q$ and using \eqref{4.6:eq-constant-conditions-d},
\begin{equation*}
\sum_{q\ge1}\frac1{q!}\Bigl(100^dN(\operatorname{core}(X'))e^{-\frac12p_0(g_k)}\Bigr)^q
\le\exp\bigl(100^de^{-\frac12p_0(g_k)}N(\operatorname{core}(X'))\bigr)
\le e^{\frac14a_vN(\operatorname{core}(X'))}.
\end{equation*}

\emph{Collecting.} We bound $|F(X')|$ by
$\sum_{\mathfrak{X}}\sum_{\mathbf{D}}|F(X';\mathfrak{X},\mathbf{D})|$ and apply
\eqref{4.6:eq-activity-summation-2}. In it we bound each island factor by the size step; the
island factors do not depend on $\mathbf{D}$. We then carry out the sum over $\mathbf{D}$ by
\eqref{4.6:eq-activity-summation-3}. This gives
\begin{align*}
|F(X')|&\le e^{-(1+\beta)\kappa d_{\mathrm{out}}(X')}\,e^{-a_vN(\operatorname{core}(X'))}\,
e^{\frac12a_vN(\operatorname{core}(X'))}
\Bigl[\sum_{\mathfrak{X}}c_0(\mathfrak{X},X')\prod_{a\in\mathfrak{X}}e^{-\frac32p_0(g_k)}\Bigr]\\
&\le e^{-(1+\beta)\kappa d_{\mathrm{out}}(X')}\,e^{-\frac12a_vN(\operatorname{core}(X'))}\,
c_1(X')\,e^{\frac14a_vN(\operatorname{core}(X'))}.
\end{align*}
In the second line we used the bound on the bracket from the sum over the islands:
\begin{itemize}
\item in Case~A, $c_1(X')=\alpha^{1/3}$, which is the common value of $c_0(\mathfrak{X},X')$;
\item in Case~B, $c_1(X')=e^{-p_0(g_k)}$, which is the extracted factor, since
$c_0(\mathfrak{X},X')\le1$.
\end{itemize}
Since $-\frac12a_v+\frac14a_v=-\frac14a_v$, this is the asserted bound. The two cases exhaust
$\mathcal{P}$, as stated in Proposition~\ref{4.6:prop-activity-bound}.
\end{proof}

\begin{proposition}[Convergence and tree decay of the logarithm]\label{4.6:prop-logarithm-convergence}
Assume the hypotheses of the activity bound for a single polymer,
Proposition~\ref{4.6:prop-activity-bound}, namely the bound on the island integrals in family form,
the strip form of the bound on the interaction terms, the localisation statement and the remaining
hypothesis listed there, together with condition~\eqref{4.6:eq-constant-conditions-a}, the ceiling
on $\alpha^{1/3}$ and the two conditions on the coupling of
Definition~\ref{4.6:def-constant-conditions}; assume the format statement on which the polymer-gas
representation of Lemma~\ref{4.6:lem-polymer-gas} rests; and assume in addition
condition~\eqref{4.6:eq-constant-conditions-e}. Then for every union $X$ of $M$-cubes the
series of Definition~\ref{4.6:def-polymer-logarithm} defining $\mathbf{R}'^{(k)}(X)$ converges
absolutely, uniformly in the outer variables and on $\mathfrak{U}_k(X)$, and
\begin{equation}\label{4.6:eq-logarithm-convergence-a}
\bigl|\mathbf{R}'^{(k)}(X,(\mathbf{U},\mathbf{J}))\bigr|
\le K_{\mathrm{KP}}(d)\, c_1(X)\,
\exp\bigl(-(1+\tfrac12\beta)\kappa\, d_{\mathrm{out}}(X)\bigr),
\end{equation}
where
\begin{equation*}
c_1(X) :=
\begin{cases}
e^{-p_0(g_k)} & \text{if } X\sqcap\bigcup_i Y_i^{\sharp}=\emptyset,\\
\bar c & \text{otherwise,}
\end{cases}
\qquad
K_{\mathrm{KP}}(d) = 4e^{2(A_d+5^d)},
\end{equation*}
with $\bar c$ and $A_d$ as in Definition~\ref{4.6:def-constant-conditions}. Moreover the series
\begin{equation}\label{4.6:eq-logarithm-convergence-1}
\Phi := \sum_{n\ge1}\frac1{n!}\sum_{(X'_1,\dots,X'_n)\in\mathcal{P}^n}
\rho^{T}(X'_1,\dots,X'_n)\prod_{p=1}^n F(X'_p)
\end{equation}
converges absolutely, and consequently the bracket $\{\cdots\}$ of \eqref{4.6:eq-polymer-gas-a}
satisfies $\{\cdots\} = \exp\sum_X \mathbf{R}'^{(k)}(X)$, which is the identity
\eqref{4.6:eq-polymer-logarithm-a}, by the exponential formula for compatible families
(Lemma~\ref{4.6:lem-exponential-formula}).
\end{proposition}

The bound \eqref{4.6:eq-logarithm-convergence-a} is the counterpart, in the present setting, of
\cite[(1.99)]{B-CMP122b}, and the identity \eqref{4.6:eq-polymer-logarithm-a} that of
\cite[(1.98)]{B-CMP122b}.

\begin{proof}
Throughout, $G(X'_1,\dots,X'_n)$ is the incompatibility graph on $\{1,\dots,n\}$, with an edge
$\{p,q\}$ exactly when $X'_p\,\iota\,X'_q$, and $\tau(G)$ is its number of spanning trees.

\emph{Only connected sequences contribute.} By Definition~\ref{4.6:def-polymer-logarithm},
$\rho^{T}(X'_1,\dots,X'_n)$ is a sum over connected graphs $g$ on $\{1,\dots,n\}$ of
$\prod_{\{p,q\}\in g}(\zeta(X'_p,X'_q)-1)$, and $\zeta(X'_p,X'_q)-1$ vanishes unless
$X'_p\,\iota\,X'_q$. Hence only graphs $g$ contained in $G(X'_1,\dots,X'_n)$ contribute, and if
$G(X'_1,\dots,X'_n)$ is not connected no connected graph on $\{1,\dots,n\}$ is contained in it, so
that $\rho^{T}=0$. In the series defining $\mathbf{R}'^{(k)}(X)$ only sequences with
$\bigcup_pX'_p=X$ and $G(X'_1,\dots,X'_n)$ connected therefore contribute. If no such sequence
exists, $\mathbf{R}'^{(k)}(X)=0$ and \eqref{4.6:eq-logarithm-convergence-a} holds trivially; we
assume from now on that one exists.

\emph{Step $(\alpha)$: tree extraction.} Let $(X'_1,\dots,X'_n)$ be a contributing sequence. We show
\begin{equation}\label{4.6:eq-logarithm-convergence-2}
d_{\mathrm{out}}(X)\le\sum_{p=1}^n d_{\mathrm{out}}(X'_p)+2\sqrt d\,(n-1).
\end{equation}
Recall (Definition~\ref{4.6:def-tree-lengths}, display \eqref{4.6:eq-tree-lengths-a}) that
$M\,d_{\mathrm{out}}(X')$ is the length of a shortest tree graph contained in $X'$ and meeting every
$M$-cube of $\operatorname{core}(X')$. For each $p$ choose such a shortest tree graph
$\Gamma_p\subset X'_p$. Choose a spanning tree of $G(X'_1,\dots,X'_n)$; it has $n-1$ edges. For
each of its edges $\{p,q\}$ we have $X'_p\,\iota\,X'_q$, so there are cubes
$c\in\operatorname{core}(X'_p)$ and $c'\in\operatorname{core}(X'_q)$ sharing a point $z$. Since
$\Gamma_p$ meets $c$ and $\Gamma_q$ meets $c'$, a straight segment inside $c$ from a point of
$\Gamma_p\cap c$ to $z$, followed by a straight segment inside $c'$ from $z$ to a point of
$\Gamma_q\cap c'$, joins $\Gamma_p$ to $\Gamma_q$ inside $c\cup c'\subset X$; each segment has
length at most the diameter $\sqrt d\,M$ of an $M$-cube, so the joining path has length at most
$2\sqrt d\,M$. This is the same joining as in the proof of Proposition~\ref{4.6:prop-activity-bound}.
The union of the $\Gamma_p$ and of the $n-1$ joining paths is a connected graph contained in $X$,
of total length at most
\begin{equation*}
M\sum_{p}d_{\mathrm{out}}(X'_p)+2\sqrt d\,M(n-1),
\end{equation*}
and it meets every cube of $\bigcup_p\operatorname{core}(X'_p)=\operatorname{core}(X)$, where
$\operatorname{core}(X)=X\setminus\bigcup_iY_i^{\sharp}$. A spanning tree of this graph is a tree
graph with the same two properties and no greater length. Dividing by $M$ gives
\eqref{4.6:eq-logarithm-convergence-2}. Now put
\begin{equation}\label{4.6:eq-logarithm-convergence-3}
\tilde w(X') := |F(X')|\,e^{(1+\frac12\beta)\kappa\,(d_{\mathrm{out}}(X')+2\sqrt d)},
\qquad X'\in\mathcal{P}.
\end{equation}
Multiplying \eqref{4.6:eq-logarithm-convergence-2} by $(1+\frac12\beta)\kappa$ and exponentiating,
\begin{equation}\label{4.6:eq-logarithm-convergence-4}
e^{(1+\frac12\beta)\kappa\, d_{\mathrm{out}}(X)}\prod_{p=1}^n|F(X'_p)|
\le e^{-(1+\frac12\beta)\kappa\,2\sqrt d}\prod_{p=1}^n\tilde w(X'_p).
\end{equation}
By the activity bound of Proposition~\ref{4.6:prop-activity-bound},
$|F(X')|\le c_1(X')\,e^{-(1+\beta)\kappa d_{\mathrm{out}}(X')}$ once the factor
$\exp(-\frac14a_v N(\operatorname{core}X'))\le1$ is dropped; since
$(1+\beta)-(1+\frac12\beta)=\frac12\beta$, this gives
\begin{equation}\label{4.6:eq-logarithm-convergence-5}
\tilde w(X')\le c_1(X')\,e^{(1+\frac12\beta)\kappa\,2\sqrt d}\,
e^{-\frac12\beta\kappa\, d_{\mathrm{out}}(X')}.
\end{equation}

\emph{Step $(\beta)$: the summability condition for $\tilde w$ with
$a(X'):=N(\operatorname{core}X')$.} Let $X''\in\mathcal{P}$. If $X'\,\iota\,X''$, there are a cube
$c'\in\operatorname{core}(X')$ and a cube $c''\in\operatorname{core}(X'')$ sharing a point; for a
given $c''$ there are at most $3^d$ such cubes $c'$ (equal to $c''$ or touching it). For an $M$-cube
$c'$ and $D\in\mathbb{N}$, the count of polymers through a cube
(Lemma~\ref{4.6:lem-polymer-count}) gives at most $2e^{A_d(D+2)}$ polymers $X'$ with
$c'\in\operatorname{core}(X')$ and $\lfloor d_{\mathrm{out}}(X')\rfloor=D$, each with
$N(\operatorname{core}X')\le5^d(D+2)$; for them $d_{\mathrm{out}}(X')\ge D$. Every polymer has
$c_1(X')\le\bar c$, the two values $\alpha^{1/3}$ and $e^{-p_0(g_k)}$ of
Proposition~\ref{4.6:prop-activity-bound} being at most $\bar c$. With
\eqref{4.6:eq-logarithm-convergence-5} we obtain
\begin{align*}
\sum_{X':\,X'\iota X''}\tilde w(X')\,e^{N(\operatorname{core}X')}
&\le\sum_{c''\in\operatorname{core}X''}\ \sum_{c'}\
\sum_{X':\,\operatorname{core}(X')\ni c'}\tilde w(X')\,e^{N(\operatorname{core}X')}\\
&\le N(\operatorname{core}X'')\,3^d\,\bar c\,e^{(1+\frac12\beta)\kappa\,2\sqrt d}
\sum_{D\ge0}2e^{A_d(D+2)}e^{5^d(D+2)}e^{-\frac12\beta\kappa D}\\
&= N(\operatorname{core}X'')\,3^d\,\bar c\,e^{(1+\frac12\beta)\kappa\,2\sqrt d}\,
2e^{2(A_d+5^d)}\sum_{D\ge0}e^{-(\frac12\beta\kappa-A_d-5^d)D}\\
&\le N(\operatorname{core}X'')\,
\bigl[3^d\,K_{\mathrm{KP}}(d)\,e^{(1+\frac12\beta)\kappa\,2\sqrt d}\,\bar c\,\bigr]
\le N(\operatorname{core}X'').
\end{align*}
In the penultimate step we used condition \eqref{4.6:eq-constant-conditions-a}, which gives
$\frac12\beta\kappa-A_d-5^d\ge\log2$, so that the sum over $D$ is a geometric series of ratio at most
$\frac12$ and is at most $2$, and $2\cdot2e^{2(A_d+5^d)}=K_{\mathrm{KP}}(d)$; in the last step we used
condition \eqref{4.6:eq-constant-conditions-e}, by which the bracket is at most one. Thus
\eqref{4.6:eq-rooted-tree-majorant-a} holds for $(\mathcal{P},\tilde w,a)$ with
$a=N(\operatorname{core}\,\cdot)$.

The same computation, pinned at one cube $c_0$ and restricted to polymers contained in $X$, gives
\begin{equation}\label{4.6:eq-logarithm-convergence-6}
\sum_{X'\subset X:\ \operatorname{core}(X')\ni c_0}\tilde w(X')\,e^{N(\operatorname{core}X')}
\le c_1(X)\,e^{(1+\frac12\beta)\kappa\,2\sqrt d}\,K_{\mathrm{KP}}(d).
\end{equation}
Here there is no sum over $c''$ and no factor $3^d$, the count of
Lemma~\ref{4.6:lem-polymer-count} being applied directly at $c_0$, and $\bar c$ is replaced by
$c_1(X)$ because every $X'\subset X$ has $c_1(X')\le c_1(X)$: if
$X\sqcap\bigcup_iY_i^{\sharp}=\emptyset$, every $X'\subset X$ satisfies
$X'\sqcap\bigcup_iY_i^{\sharp}=\emptyset$, the case of Proposition~\ref{4.6:prop-activity-bound} in
which $c_1(X')=e^{-p_0(g_k)}$, so $c_1(X')=e^{-p_0(g_k)}=c_1(X)$; otherwise
$c_1(X)=\bar c\ge c_1(X')$.

\emph{Step $(\gamma)$: assembly.} Pick $c_0\in\operatorname{core}(X)$; the core is non-empty, since
every polymer has a non-empty core (it contains a first-class strip or the off-cubes of an
interaction domain) and $X$ is the union of the polymers of a contributing sequence. For a contributing
sequence $s=(X'_1,\dots,X'_n)$ let $m(s):=\#\{p:\operatorname{core}(X'_p)\ni c_0\}$; then $m(s)\ge1$,
because $c_0\in\operatorname{core}(X)=\bigcup_p\operatorname{core}(X'_p)$. Both $\tau(G)$ and
$\prod_p\tilde w(X'_p)$ are invariant under permutations of $s$, so that each of the $n$ terms in the
middle sum below equals the term with $p=1$. Writing $\mathbf{1}\{\,\cdot\,\}$ for the indicator,
and writing $\cup=X$ under a sum over $s$ for the constraint $\bigcup_pX'_p=X$, the sum running
also over $n\ge1$, we obtain
\begin{align*}
\sum_{s:\,\cup=X}\frac1{n!}\,\tau(G)\prod_p\tilde w(X'_p)
&\le\sum_{s:\,\cup=X}\frac{m(s)}{n!}\,\tau(G)\prod_p\tilde w(X'_p)\\
&=\sum_{p=1}^n\ \sum_{s:\,\cup=X}\mathbf{1}\{\operatorname{core}(X'_p)\ni c_0\}\,
\frac1{n!}\,\tau(G)\prod_q\tilde w(X'_q)\\
&=\sum_{s:\,\cup=X,\ \operatorname{core}(X'_1)\ni c_0}\frac1{(n-1)!}\,
\tau(G)\prod_p\tilde w(X'_p)\\
&\le\sum_{X'_1\subset X:\ \operatorname{core}(X'_1)\ni c_0}\tilde w(X'_1)
\sum_{n\ge1}\frac1{(n-1)!}\sum_{(X'_2,\dots,X'_n)\in\mathcal{P}_X^{\,n-1}}
\tau(G(X'_1,X'_2,\dots,X'_n))\prod_{q\ge2}\tilde w(X'_q),
\end{align*}
where $\mathcal{P}_X$ is the set of polymers contained in $X$, and in the last step the constraint
$\bigcup_pX'_p=X$ was relaxed to $X'_p\subset X$
for all $p$, which only adds non-negative terms. The summability condition
\eqref{4.6:eq-rooted-tree-majorant-a} holds for $(\mathcal{P}_X,\tilde w,N(\operatorname{core}\,\cdot))$
a fortiori, its left side being a sub-sum of the one estimated in Step $(\beta)$. The rooted-tree
majorant (Lemma~\ref{4.6:lem-rooted-tree-majorant}), applied to this system, bounds the inner series
--- which is the series \eqref{4.6:eq-rooted-tree-majorant-b} for $\mathcal{P}_X$ at $X'_1$ --- by
$e^{N(\operatorname{core}X'_1)}$, and then \eqref{4.6:eq-logarithm-convergence-6} bounds the
remaining sum over $X'_1$ by $c_1(X)e^{(1+\frac12\beta)\kappa\,2\sqrt d}K_{\mathrm{KP}}(d)$.

By the spanning-tree bound (Lemma~\ref{4.6:lem-spanning-tree-bound}),
$|\rho^{T}(X'_1,\dots,X'_n)|\le\tau(G(X'_1,\dots,X'_n))$. Combining it with
\eqref{4.6:eq-logarithm-convergence-4} and with the estimate just proved,
\begin{align*}
\bigl|\mathbf{R}'^{(k)}(X)\bigr|
&\le\sum_{n\ge1}\frac1{n!}\sum_{\cup=X}|\rho^{T}|\prod_p|F(X'_p)|\\
&\le e^{-(1+\frac12\beta)\kappa\, d_{\mathrm{out}}(X)}\,
e^{-(1+\frac12\beta)\kappa\,2\sqrt d}\,c_1(X)\,
e^{(1+\frac12\beta)\kappa\,2\sqrt d}\,K_{\mathrm{KP}}(d)\\
&=K_{\mathrm{KP}}(d)\,c_1(X)\,e^{-(1+\frac12\beta)\kappa\, d_{\mathrm{out}}(X)}.
\end{align*}
The first line bounds the series of absolute values, so the series defining $\mathbf{R}'^{(k)}(X)$
converges absolutely, and \eqref{4.6:eq-logarithm-convergence-a} holds.

\emph{Absolute convergence of $\Phi$ and the exponential identity.} Since
$d_{\mathrm{out}}\ge0$, \eqref{4.6:eq-logarithm-convergence-3} gives $|F|\le\tilde w$. By the
spanning-tree bound, and since $1/n!\le1/(n-1)!$,
\begin{align*}
\sum_{n\ge1}\frac1{n!}\sum_{(X'_p)\in\mathcal{P}^n}|\rho^{T}|\prod_p|F(X'_p)|
&\le\sum_{n\ge1}\frac1{n!}\sum_{(X'_p)\in\mathcal{P}^n}\tau(G)\prod_p\tilde w(X'_p)\\
&\le\sum_{X'_1\in\mathcal{P}}\tilde w(X'_1)\,B(X'_1)
\le\sum_{X'_1\in\mathcal{P}}\tilde w(X'_1)\,e^{N(\operatorname{core}X'_1)}<\infty,
\end{align*}
where $B(X'_1)$ is the series \eqref{4.6:eq-rooted-tree-majorant-b} with weight $\tilde w$, bounded
by the rooted-tree majorant (Lemma~\ref{4.6:lem-rooted-tree-majorant}), which applies by
Step $(\beta)$, and the last sum is finite because $\mathcal{P}$ is finite. Hence the exponential
formula (Lemma~\ref{4.6:lem-exponential-formula}) applies and gives $\Xi=\exp\Phi$, where $\Xi$ is
the sum over pairwise compatible families of products of activities, which equals the bracket
$\{\cdots\}$ by the polymer-gas representation \eqref{4.6:eq-polymer-gas-a} of
Lemma~\ref{4.6:lem-polymer-gas}. Every sequence with $\rho^{T}\ne0$ has a connected graph $G$, and
its union is a union $X$ of polymers; grouping the absolutely convergent series
\eqref{4.6:eq-logarithm-convergence-1} according to $X=\bigcup_pX'_p$ gives
$\Phi=\sum_X\mathbf{R}'^{(k)}(X)$, and therefore $\{\cdots\}=\exp\sum_X\mathbf{R}'^{(k)}(X)$, which is
\eqref{4.6:eq-polymer-logarithm-a}.

\emph{Uniformity.} Every bound used above is uniform on $\mathfrak{U}_k(X)$ and in the outer
variables, because the activity bound is, and so are the hypotheses
of Proposition~\ref{4.6:prop-activity-bound} on the interaction terms and the island integrals from
which it is derived.

\emph{Infinite volume.} When the polymers live in $\mathbb{R}^d$, the exhaustion by the finite
systems $\mathcal{P}_R$ of Remark~\ref{4.6:rem-tree-bounds-use} applies, since every bound above is
independent of the side $R$ of the box defining $\mathcal{P}_R$.
\end{proof}

\begin{proposition}[Representation and bound of the large-field terms]\label{4.6:prop-large-field-bound}
Let $k$ be the level of the step, and let the family of configurations and the spaces
$\mathfrak{U}_k(X)$ be those of Convention~\ref{4.6:not-complex-domain}. Let $Z_k$ be the
large-field region and $Y_i$, $i$ in a finite index set, the regions of the outer datum of the
step, as in Definition~\ref{4.6:def-polymer-expansion}. Assume the
following.
\begin{itemize}
\item The hypotheses of Proposition~\ref{4.6:prop-logarithm-convergence}, that is, those of
Proposition~\ref{4.6:prop-activity-bound} together with the further one-sided condition on
$\bar{c}$ assumed in Proposition~\ref{4.6:prop-logarithm-convergence}.
\item The analyticity statement \eqref{4.6:eq-standing-bounds-c} of
Definition~\ref{4.6:def-standing-bounds}.
\item The following statements, on the family of Convention~\ref{4.6:not-complex-domain}: the
format statement, the statement on the prepared sequences of large-field regions, and the
positivity dichotomy of the large-field operations.
\item The one-sided condition \eqref{4.6:eq-constant-conditions-f}.
\end{itemize}
Let $\mathbf{R}'^{(k)}(X)$ be the logarithm terms of
Definition~\ref{4.6:def-polymer-logarithm}, and let $c_1(X)$ and
$K_{\mathrm{KP}}(d)=4e^{2(A_d+5^d)}$ be as in Proposition~\ref{4.6:prop-logarithm-convergence}.
Then the following hold.
\begin{enumerate}
\item[(a)] For every union $X$ of $M$-cubes and all $(\mathbf{U},\mathbf{J})\in\mathfrak{U}_k(X)$,
\begin{equation}\label{4.6:eq-large-field-bound-1}
\bigl|\mathbf{R}'^{(k)}(X,(\mathbf{U},\mathbf{J}))\bigr|
\le K_{\mathrm{KP}}(d)\,c_1(X)\,
\exp\Bigl(-\bigl(1+\tfrac12\beta\bigr)\kappa\, d_{k,\cup_i Y_i^{\sharp}}(X)\Bigr).
\end{equation}
\item[(b)] For every $X$ in the second group of Definition~\ref{4.6:def-polymer-logarithm},
that is, with
$X\sqcap\bigl(Z_k\cup\bigcup_i Y_i\bigr)^{\boxplus}=\emptyset$, and for all
$(\mathbf{U},\mathbf{J})\in\mathfrak{U}_k(X)$,
\begin{equation}\label{4.6:eq-large-field-bound-a}
\bigl|\mathbf{R}'^{(k)}(X,(\mathbf{U},\mathbf{J}))\bigr|
\le \exp\bigl(-p_0(g_k)\bigr)\exp\bigl(-\kappa\, d_k(X)\bigr).
\end{equation}
\item[(c)] Let $X$ be in the second group. The real unit point $(U,J)=(1,0)$ lies in
$\mathfrak{U}_k(X)$; write $\mathbf{R}'^{(k)}(X,1)$ for the value of
$\mathbf{R}'^{(k)}(X,\cdot)$ there. Then $|\mathbf{R}'^{(k)}(X,1)|$ obeys
\eqref{4.6:eq-large-field-bound-a}, and the shifted member
$\mathbf{R}'^{(k)}(X,\cdot)-\mathbf{R}'^{(k)}(X,1)$ obeys
\eqref{4.6:eq-large-field-bound-a} with an extra factor $2$ on the right-hand side. If
condition \eqref{4.6:eq-constant-conditions-f} holds with $2K_{\mathrm{KP}}(d)$ in place of
$K_{\mathrm{KP}}(d)$, the shifted member obeys \eqref{4.6:eq-large-field-bound-a} itself.
\item[(d)] For every union $X$ of $M$-cubes the function $\mathbf{R}'^{(k)}(X,\cdot)$ depends
only on the variables on $X$, and it is an
analytic function of $(\mathbf{U},\mathbf{J})$ on $\mathfrak{U}_k(X)$.
\end{enumerate}
\end{proposition}

The inductive hypothesis of the construction at the levels not exceeding $k$ enters through
the statements assumed above.
\cite[(1.99)]{B-CMP122b} states the bound \eqref{4.6:eq-large-field-bound-1} with an
unspecified constant $O(1)$ and with the length $d_{k,\cup_i Y_i}(X)$; here the constant is
$K_{\mathrm{KP}}(d)$, which depends on $d$ only, and the length is measured relative to the
strips $\bigcup_i Y_i^{\sharp}$. Estimate \eqref{4.6:eq-large-field-bound-a} is the bound
\cite[(1.100)]{B-CMP122b}, here established on the family and the space of
Convention~\ref{4.6:not-complex-domain}, on which the standing bounds and
\eqref{4.6:eq-standing-bounds-c} are supplied.

\begin{proof}
\emph{Clause (a).} By Proposition~\ref{4.6:prop-logarithm-convergence}, whose hypotheses are
assumed, for every union $X$ of $M$-cubes the series defining $\mathbf{R}'^{(k)}(X)$
converges absolutely, uniformly on $\mathfrak{U}_k(X)$, and
\begin{equation*}
\bigl|\mathbf{R}'^{(k)}(X,(\mathbf{U},\mathbf{J}))\bigr|
\le K_{\mathrm{KP}}(d)\,c_1(X)\exp\Bigl(-\bigl(1+\tfrac12\beta\bigr)\kappa\,
d_{\mathrm{out}}(X)\Bigr).
\end{equation*}
By Definition~\ref{4.6:def-tree-lengths}, $d_{\mathrm{out}}(X)=d_{k,\cup_i Y_i^{\sharp}}(X)$.
This is \eqref{4.6:eq-large-field-bound-1}.

\emph{Clause (b).} Let $X$ be in the second group. By Definition~\ref{4.6:def-strips-bracket},
$Y_i^{\sharp}\subset Y_i^{\boxplus}$ for every $i$. The set $Y_i^{\boxplus}$ consists of $Y_i$
and the cells touching $Y_i$. Each such cell touches $Z_k\cup\bigcup_j Y_j$, so
$Y_i^{\boxplus}\subset(Z_k\cup\bigcup_j Y_j)^{\boxplus}$, and therefore
$X\sqcap Y_i^{\boxplus}=\emptyset$. Hence $X\sqcap\bigcup_i Y_i^{\sharp}=\emptyset$. Three
consequences follow. First, $c_1(X)=e^{-p_0(g_k)}$ by the definition of $c_1(X)$ in
Proposition~\ref{4.6:prop-logarithm-convergence}. Second, $\operatorname{core}(X)=X$, since
no $M$-cube of $X$ lies in a strip $Y_i^{\sharp}$. Third,
$d_{\mathrm{out}}(X)\ge d_k(X)$ by Definition~\ref{4.6:def-tree-lengths}.
This is the remark of \cite[p.~390]{B-CMP122b} that for the terms of the second group the
length in \cite[(1.99)]{B-CMP122b} may be replaced by $d_k(X)$ and $c_1$ is equal to
$\exp(-p_0(g_k))$.

If no tuple of polymers has union $X$, then $\mathbf{R}'^{(k)}(X)=0$ and
\eqref{4.6:eq-large-field-bound-a} holds trivially. Otherwise let $X'\subset X$ be a
polymer. Since $X'\sqcap\bigcup_i Y_i^{\sharp}=\emptyset$, $X'$ is in the second case of
Proposition~\ref{4.6:prop-activity-bound}. It therefore contains the strip $X_a^{\sharp}$ of
an island $X_a$ of its datum, and we obtain
\begin{equation*}
X\supseteq X'\supseteq X_a^{\sharp}\supseteq X_a.
\end{equation*}
The island $X_a$ contains at least one cell of level $k$, which consists of
$\check{R}_k^{d}$ $M$-cubes. A shortest tree graph realising $d_k(X)$ has length
$d_k(X)M$ and meets every $M$-cube of $X$, in particular these $\check{R}_k^{d}$ cubes.
Lemma~\ref{4.6:lem-tree-cubes} therefore gives
$\check{R}_k^{d}\le 5^{d}(d_k(X)+1)$, that is,
\begin{equation}\label{4.6:eq-large-field-bound-2}
d_k(X)\ge 5^{-d}\check{R}_k^{d}-1 .
\end{equation}
By condition \eqref{4.6:eq-constant-conditions-f} and \eqref{4.6:eq-large-field-bound-2},
\begin{equation}\label{4.6:eq-large-field-bound-3}
\log K_{\mathrm{KP}}(d)\le \tfrac12\beta\kappa\bigl(5^{-d}\check{R}_k^{d}-1\bigr)
\le \tfrac12\beta\kappa\, d_k(X).
\end{equation}
Insert $c_1(X)=e^{-p_0(g_k)}$ and $d_{\mathrm{out}}(X)\ge d_k(X)$ into
\eqref{4.6:eq-large-field-bound-1}, and then use \eqref{4.6:eq-large-field-bound-3}:
\begin{align*}
\bigl|\mathbf{R}'^{(k)}(X,(\mathbf{U},\mathbf{J}))\bigr|
&\le K_{\mathrm{KP}}(d)\,e^{-p_0(g_k)}\,e^{-(1+\frac12\beta)\kappa d_k(X)}\\
&= e^{-p_0(g_k)}e^{-\kappa d_k(X)}
\bigl[K_{\mathrm{KP}}(d)\,e^{-\frac12\beta\kappa d_k(X)}\bigr]
\le e^{-p_0(g_k)}e^{-\kappa d_k(X)} .
\end{align*}
This holds for all $(\mathbf{U},\mathbf{J})\in\mathfrak{U}_k(X)$, and is
\eqref{4.6:eq-large-field-bound-a}.

\emph{Clause (c).} Let $X$ be in the second group. The real unit point $(U,J)=(1,0)$ lies in
the real slice of the class of Convention~\ref{4.6:not-complex-domain}, by the real-slice
statement for that class, and hence in $\mathfrak{U}_k(X)$; this is consistent with the fact
that the spaces $\tilde{U}^{c}_{k}(X,\tilde\alpha_0,\tilde\alpha_1)$ of \cite{B-CMP122b}
contain their centre.
Clause (b) applied at this point bounds $|\mathbf{R}'^{(k)}(X,1)|$ by the right-hand side of
\eqref{4.6:eq-large-field-bound-a}.

The shifted member is obtained by subtracting the value
at the unit point term by term after the exponentiation \eqref{4.6:eq-polymer-logarithm-a}.
By the triangle inequality and clause (b), applied at $(\mathbf{U},\mathbf{J})$ and at $1$,
\begin{equation*}
\bigl|\mathbf{R}'^{(k)}(X,(\mathbf{U},\mathbf{J}))-\mathbf{R}'^{(k)}(X,1)\bigr|
\le 2\,e^{-p_0(g_k)}e^{-\kappa d_k(X)} .
\end{equation*}
Suppose instead that condition \eqref{4.6:eq-constant-conditions-f} holds with
$2K_{\mathrm{KP}}(d)$ in place of $K_{\mathrm{KP}}(d)$. The argument of clause (b) then
gives $\log(2K_{\mathrm{KP}}(d))\le\frac12\beta\kappa\,d_k(X)$. The first line of the
display in clause (b), applied at both points and added, gives
\begin{align*}
\bigl|\mathbf{R}'^{(k)}(X,(\mathbf{U},\mathbf{J}))-\mathbf{R}'^{(k)}(X,1)\bigr|
&\le 2K_{\mathrm{KP}}(d)\,e^{-p_0(g_k)}e^{-(1+\frac12\beta)\kappa d_k(X)}\\
&= e^{-p_0(g_k)}e^{-\kappa d_k(X)}
\bigl[2K_{\mathrm{KP}}(d)\,e^{-\frac12\beta\kappa d_k(X)}\bigr]
\le e^{-p_0(g_k)}e^{-\kappa d_k(X)} .
\end{align*}
This is \eqref{4.6:eq-large-field-bound-a} for the shifted member.

\emph{Clause (d).} We first treat a single activity $F(X')$.

\begin{itemize}
\item \emph{Analyticity.} $F(X')$ is analytic on $\mathfrak{U}_k(X')$ by
\eqref{4.6:eq-standing-bounds-c}. Indeed, it is built from two kinds of operations, both of
which preserve analyticity:
\begin{itemize}
\item finite sums over data, the sums over histories inside $\mathbf{T}'_k$ being finite;
\item integrals over $t\in[0,1]$ of integrands that are analytic and continuous in $t$.
\end{itemize}
\item \emph{Locality.} $F(X')$ depends only on the variables on $X'$. Each interaction term
$V(Y)$ entering it has $Y\subset X'$. Each island operation $\mathbf{T}'_k(X_a)$ entering it
has $X_a^{\sharp}\subset X'$. The strip collar of $X_a$, that is, the collar around $X_a$ on
whose variables $\mathbf{T}'_k(X_a)$ may depend, lies inside
$X'$, which contains $X_a^{\sharp}$.
\end{itemize}
The bounds of Propositions~\ref{4.6:prop-activity-bound}
and~\ref{4.6:prop-logarithm-convergence} hold uniformly on $\mathfrak{U}_k(X')$ and in the
outer variables.

We now pass to $\mathbf{R}'^{(k)}(X)$. By Proposition~\ref{4.6:prop-logarithm-convergence},
the series for $\mathbf{R}'^{(k)}(X)$ converges absolutely and uniformly on
$\mathfrak{U}_k(X)$. This set is a domain in the finite-dimensional complex manifold of
configurations on the bonds of $X$. Each term of the series is holomorphic there and depends
only on variables on $X$. A uniform limit of holomorphic functions is holomorphic.
Hence $\mathbf{R}'^{(k)}(X,\cdot)$ depends only on the variables on $X$ and is analytic on
$\mathfrak{U}_k(X)$. In the reading of Convention~\ref{4.6:not-complex-domain} of the space
$\tilde{U}^{c}_{k}(X,\tilde\alpha_0,\tilde\alpha_1)$, this is the assertion of
\cite[p.~390]{B-CMP122b} that $\mathbf{R}'^{(k)}(X)$ depends on the background field
restricted to $X$ and extends analytically in $(\mathbf{U},\mathbf{J})$ to that space.
\end{proof}

\begin{remark}[The bound in the form used by the induction]\label{4.6:rem-inductive-form}
Let $\kappa_0$ be the exponent of \cite[(2.31)]{B-CMP119}, where the large-field terms are bounded in the form
\begin{equation}\label{4.6:eq-inductive-form-1}
|\mathbf{R}^{(k)}(X)| \le g_k^{\kappa_0}\, e^{-\kappa d_k(X)} .
\end{equation}
This form follows from clauses (b)--(c) of Proposition~\ref{4.6:prop-large-field-bound} if and only if
\begin{equation}\label{4.6:eq-inductive-form-2}
2e^{-p_0(g_k)} \le g_k^{\kappa_0},
\end{equation}
that is, if and only if
\begin{equation*}
A_0\bigl(\log g_k^{-2}\bigr)^{p_0} \ge \tfrac12\kappa_0 \log g_k^{-2} + \log 2 .
\end{equation*}
This condition holds when the fixed exponent satisfies $p_0\ge 2$ and $g_k$ lies below a ceiling
depending only on $A_0$ and $\kappa_0$; this ceiling is a condition on $\gamma$. The passage from
\eqref{4.6:eq-large-field-bound-a} to \eqref{4.6:eq-inductive-form-1} is not made in the present
section. The large-field part of the coupling increment uses \eqref{4.6:eq-large-field-bound-a}
itself: the term $\mathfrak{r}'^{(i)}(X)$ of the relative small-field machine is $\mathbf{R}'$
evaluated at the empty admissible label of that machine, at which every $X$ belongs to the second
group of clause~(b) of Proposition~\ref{4.6:prop-large-field-bound}, so
\eqref{4.6:eq-large-field-bound-a} holds for it.
\end{remark}

\begin{proof}
Clauses (b)--(c) of Proposition~\ref{4.6:prop-large-field-bound} together bound the term
$\mathbf{R}^{(k)}(X)$ of \eqref{4.6:eq-inductive-form-1} by the prefactor $2e^{-p_0(g_k)}$ times
the decay factor $e^{-\kappa d_k(X)}$. The decay factor is the same as in
\eqref{4.6:eq-inductive-form-1}. Passing to \eqref{4.6:eq-inductive-form-1} is therefore the
comparison of prefactors \eqref{4.6:eq-inductive-form-2}.

Both sides of \eqref{4.6:eq-inductive-form-2} are positive, so we may take logarithms. Since $\kappa_0\log g_k = -\tfrac12\kappa_0\log g_k^{-2}$, the condition \eqref{4.6:eq-inductive-form-2} reads
\begin{equation*}
\log 2 - p_0(g_k) \le -\tfrac12\kappa_0\log g_k^{-2}.
\end{equation*}
Equivalently, $p_0(g_k)\ge \tfrac12\kappa_0\log g_k^{-2}+\log 2$. The degree function is $p_0(g_k)=A_0(\log g_k^{-2})^{p_0}$, and inserting it gives the displayed inequality of the remark.

For the sufficient condition, write $x_k=g_k^{-2}$. Suppose $p_0\ge 2$ and
\begin{equation*}
\log x_k\ \ge\ \max\{1,\ \kappa_0/A_0,\ (2\log 2/A_0)^{1/2}\},
\end{equation*}
that is, $g_k$ does not exceed the exponential of minus one half of this maximum, a ceiling on the
coupling fixed by $A_0$ and $\kappa_0$. Since $\log x_k\ge 1$, we have
$(\log x_k)^{p_0}\ge(\log x_k)^2$. Since $\log x_k\ge\kappa_0/A_0$, one half of
$A_0(\log x_k)^2$ is at least $\tfrac12\kappa_0\log x_k$; since $(\log x_k)^2\ge 2\log 2/A_0$, the
other half is at least $\log 2$. Hence
\begin{equation*}
A_0(\log x_k)^{p_0}\ \ge\ A_0(\log x_k)^2\ \ge\ \tfrac12\kappa_0\log x_k+\log 2 .
\end{equation*}

The last assertion is the identification made in the remark. At the empty admissible label of the
relative small-field machine every $X$ is in the second group, and on that group
\eqref{4.6:eq-large-field-bound-a} is the bound of Proposition~\ref{4.6:prop-large-field-bound}.
\end{proof}

\begin{remark}[Transfer of the step bounds to the widened support. Proof outstanding]\label{4.6:rem-widened-support}
Throughout, $k$ is the step and $i\le k$ a level. We write $u_i=L^{-(k-i)}$ for the unit at level $i$, $a_i=M\check{R}_iu_i$ for the side of a level-$i$ cell, and $[A]_i$ for the level-$i$ cell hull of a set $A$. The island $X_a$ is an island of the admissible family (Definition~\ref{4.6:def-step-datum}), and $g''_k(\cdot)$ is the prepared coupling function. The set $\mathcal{D}$ is the support of the coupling excess on the family, displayed in \eqref{4.6:eq-coupling-support-a} (Remark~\ref{4.6:rem-coupling-support}). The symbol $\Lambda^{(1.73)}$ denotes Ba{\l}aban's support: the region $Z''_k$ of \cite[\S1]{B-CMP122b} together with one layer of $M$-cubes around it, a set contained in a cube of side $100M$ \cite[\S1]{B-CMP122b}.

The proof of the representation and bound of the large-field terms (Proposition~\ref{4.6:prop-large-field-bound}) reads nothing of the coupling profiles. The support $\mathcal{D}$ enters it only through the two standing bounds of Definition~\ref{4.6:def-standing-bounds}:
\begin{itemize}
\item the bound on the island operations, through the exponent of $\mathbf{T}'_k$: the Wilson-action terms of \cite[(1.37)]{B-CMP122b}, the stage terms $\mathbb{B}''_k$ inside \cite[(1.71)]{B-CMP122b}, and the normalisation;
\item the bound on the interaction terms, through the interaction terms $V(Y)$ of Definition~\ref{4.6:def-standing-bounds} as constructed in \cite[(1.59)--(1.70)]{B-CMP122a}, namely through the nested site sums and the completion terms of that construction.
\end{itemize}
The estimates of \cite{B-CMP122a} and \cite{B-CMP122b} from which these two bounds are drawn are established on Ba{\l}aban's own supports. There the ramps of the coupling profiles are layers of the member's own cubes; here a member is a set $A$ of the prepared or stage sequence, and its own cubes at level $i$ are the level-$i$ $M$-cubes of $A$, each of side $Mu_i$. On these supports the coupling excess lies inside $\Lambda^{(1.73)}$. It is in $\Lambda^{(1.73)}$ that \cite[\S1]{B-CMP122b} localises the coupling excess $(g''_k(\cdot))^{-2}-g_k^{-2}$. In the same part of \cite{B-CMP122b} a standing assumption is made, that an auxiliary integer of that paper, a function of $g_k$, does not exceed a scale of that paper which is also a function of $g_k$; we call it the regime assumption of \cite[\S1]{B-CMP122b}.

The profiles of the admissible family (Remark~\ref{4.6:rem-coupling-support}) take their ramps over level-$i$ cells at the prepared and stage sequences. With them the supports are $\mathcal{D}\supsetneq\Lambda^{(1.73)}$. At the prepared sequence the region $Z''_k$ is rounded up to level-$k$ cells, and these are not contained in $\Lambda^{(1.73)}$. So the localisation in $\Lambda^{(1.73)}$ fails on the family, although the excess stays inside $X_a$. Ba{\l}aban's estimates therefore reach the family only through the following statement, which we call the \emph{widened-support statement}.

\emph{Widened-support statement.} The following estimates hold with every support $\Lambda^{(1.73)}$, and every own-cube layer, replaced by $\mathcal{D}$:
\begin{enumerate}
\item[(i)] the estimates \cite[(1.14)--(1.48)]{B-CMP122b}. In particular this covers the budget of small terms in \cite[(1.37)]{B-CMP122b}, where $O(1)$ is the sum of all the small terms, each small after division by $g_k^2$ \cite[p.~365]{B-CMP122b}. It also covers the regime assumption recalled above, and every part of \cite[(1.38)--(1.48)]{B-CMP122b} that reads the support.
\item[(ii)] the estimates \cite[(1.60)--(1.70)]{B-CMP122a}. In particular this covers the weight-difference members of \cite[(1.62)]{B-CMP122a},
\begin{equation}\label{4.6:eq-widened-support-1}
A\Bigl(\frac{1}{(g_k^{(n+1)}(\cdot))^{2}}-\frac{1}{(g_k^{(n)}(\cdot))^{2}},\;U_k^{(n+1)}\Bigr),
\end{equation}
where $g_k^{(n)}(\cdot)$ and $U_k^{(n)}$ are the coupling function and the configuration of the $n$-th stage in \cite[(1.62)]{B-CMP122a}, and their share of the constant $C_0$ of \cite[(1.70)]{B-CMP122a}.
\end{enumerate}
Here, for each level $i$, the part of $\mathcal{D}$ contributed by the level-$i$ member is the $\tfrac14a_i$-collar of $[Z''_i]_i$; its volume is at most $2^d(3a_i/2)^d$ per level-$i$ cell of $[Z''_i]_i$:
\begin{equation}\label{4.6:eq-widened-support-2}
\bigl|\text{the }\tfrac14a_i\text{-collar of }[Z''_i]_i\bigr|
\le 2^d\Bigl(\frac{3a_i}{2}\Bigr)^{d}\cdot\#\{\text{level-}i\text{ cells of }[Z''_i]_i\}.
\end{equation}
\end{remark}

Proof outstanding for the widened-support statement.

\textit{Route.} There are two ways forward, and either suffices for the use made of the widened-support statement in Proposition~\ref{4.6:prop-large-field-bound}.

The first way is to re-run \cite[(1.14)--(1.48)]{B-CMP122b} and \cite[(1.60)--(1.70)]{B-CMP122a} with the supports replaced by $\mathcal{D}$. The points requiring attention are the budget of small terms in \cite[(1.37)]{B-CMP122b} and the weight differences \eqref{4.6:eq-widened-support-1}. This route would rest on the following inputs.
\begin{itemize}
\item In the notation of \cite{B-CMP122b}, every constant of \cite[(1.14)--(1.37)]{B-CMP122b} is either absolute or an explicit function of $g_k$ alone. The absolute constants are $d,L,M,M_1,A_0,A_1,B_3,B_5,\gamma_0,\delta,\delta_0,\beta_0$. The functions of $g_k$ are $\varepsilon_k,\delta'_k,p_0(g_k),p_1(g_k)$ and the two quantities that enter the regime assumption.
\item \cite[p.~361]{B-CMP122b} notes that the number on the right-hand side of the estimate there, divided by $g_k^2$, is still small if $g_k$ is small enough.
\item Each affected small term carries a positive power of $g_k$, times a power of $\log g_k^{-1}$, times a volume.
\item Per cell of $[Z''_k]_k$ the ratio $|\mathcal{D}|/|\Lambda^{(1.73)}|$ is bounded, up to a constant factor, by $(3\check{R}_k/200)^d$. Since $\check{R}_k$ is at most $L$ times a fixed power of $\log g_k^{-2}$, this ratio is bounded by a power of $\log g_k^{-1}$.
\end{itemize}
What has to be established is that the same kind of ceiling on the constant $\gamma$ which absorbs these terms on $\Lambda^{(1.73)}$ also absorbs them with the volume of $\mathcal{D}$ in place of that of $\Lambda^{(1.73)}$, with the ceiling moved.

The term kept in \cite[(1.37)]{B-CMP122b},
\begin{equation*}
-A\bigl(((g''_k)^{-2}-g_k^{-2})\zeta,\;U''_{k,Z}\bigr),
\end{equation*}
where $\zeta$ is the cut-off function and $U''_{k,Z}$ the configuration of \cite[(1.37)]{B-CMP122b}, needs no such change. It is localised in $\mathcal{D}$, which lies inside $X_a$, so its locality is intact. Its coefficient obeys the graded bound of \cite[p.~364]{B-CMP122b}, namely $O(1)(k-j)$ on $\Omega''_j\setminus\Omega''_{j+1}$, where the $\Omega''_j$ are the regions of \cite[\S1]{B-CMP122b} indexed by the level $j$, for the cell-made profiles and for the own-cube profiles alike.

The second way is to replace the family's profiles, at the levels where this is needed, by \emph{own-cube profiles}. The widened-support statement is then not needed. The verification that follows is complete; these profiles are not adopted in this book, and Remark~\ref{4.6:rem-coupling-support} keeps the cell-made profiles.

Let $i$ be a level at which the member $A$ of the prepared or stage sequence, whichever is in use at this point, is made of $M$-cubes and not of cells \cite[p.~178]{B-CMP122a}. At such a level the product defining the profile is taken over the own cubes of $A$, that is, over the level-$i$ $M$-cubes $D$ of $A$, each of side $Mu_i$. Writing $D_\mu$ for the projection of $D$ on the $\mu$-th coordinate axis and $\psi_\varphi$ for the bump of the profile, the own-cube profile is
\begin{equation}\label{4.6:eq-widened-support-3}
\varphi_i(x)=\prod_{D}\Bigl(1-\prod_{\mu=1}^{d}\psi_\varphi\bigl(\operatorname{dist}(x_\mu,D_\mu)/(Mu_i)\bigr)\Bigr).
\end{equation}

For this profile the set where $\varphi_i$ differs from $1$ is controlled as follows:
\begin{equation}\label{4.6:eq-widened-support-4}
\begin{aligned}
\operatorname{supp}(1-\varphi_i)
&\subset\{x:\operatorname{dist}_\infty(x,A)<\tfrac14Mu_i\}\\
&\subset A\cup(\text{the first layer of own cubes of }A).
\end{aligned}
\end{equation}
This is inside the one-cube layer on which Ba{\l}aban's profile may differ from $1$. That layer comes from the two conditions of \cite[(2.24), p.~259]{B-CMP119}, namely $\varphi_j\in C_0^\infty(\Lambda_j)$ and $\varphi_j=1$ on $\Lambda_j^{\sim-1}$, transported to the member. So the profile equals $0$ on $A$ and $1$ off one layer of the cubes of $A$.

Writing $\varphi''_i$ for the level-$i$ profile of the prepared sequence, built as in \eqref{4.6:eq-widened-support-3} at the levels in question, we obtain, member by member, that $\operatorname{supp}(1-\varphi''_i)$ on the family lies inside the set on which Ba{\l}aban's own $\varphi''_i$ may differ from $1$. Consequently
\begin{equation*}
\operatorname{supp}\bigl((g''_k(\cdot))^{-2}-g_k^{-2}\bigr)
=\bigcup_{i\le k}\operatorname{supp}(1-\varphi''_i)\cap Z
\end{equation*}
lies inside Ba{\l}aban's support of the same excess, which \cite[\S1]{B-CMP122b} localises in $\Lambda^{(1.73)}$. The weights stay in $[0,1]$.

The hypotheses that \cite[(1.14)--(1.48)]{B-CMP122b} and \cite[(1.60)--(1.70)]{B-CMP122a} place on the profiles are the two conditions of \cite[(2.24)]{B-CMP119} and the supports, and no derivative of the profile enters their constants. These estimates therefore apply to the family as stated.

The clauses of the lemma on the coupling profiles hold for \eqref{4.6:eq-widened-support-3} at these levels:
\begin{enumerate}
\item[(a)] $\varphi_i$ is smooth and vanishes within distance $\tfrac18Mu_i$ of $A$.
\item[(b)] $\varphi_i=1$ beyond distance $\tfrac14Mu_i$ of $A$, in particular off the first layer of own cubes of $A$. The remaining inequality of this clause reads $u_i\le\tfrac34Mu_i$, that is, $M\ge4/3$.
\item[(c)] $\varphi_i$ is natural under the motions preserving the level-$i$ cell partition. These motions preserve the partition into own cubes, because $a_i\in Mu_i\mathbb{N}$, so that every cell is a union of own cubes.
\item[(d)] $\varphi_i$ is multiplicative over cube-disjoint families, by the product structure of \eqref{4.6:eq-widened-support-3}.
\item[(e)] The locality clause at radius $\tfrac14a_i$ holds a fortiori. By \eqref{4.6:eq-widened-support-3}, $\varphi_i(x)$ is determined by the own cubes within distance $\tfrac14Mu_i$ of $x$. Since $a_i\in Mu_i\mathbb{N}$ gives $\tfrac14Mu_i\le\tfrac14a_i$, these cubes lie in cells within $\tfrac14a_i$ of $x$.
\item[(f)] For a set $S$, the profile $\varphi^{(S)^{c}}$ of this clause differs from $1$ only on $\{\operatorname{dist}_\infty(\cdot,S)<\tfrac14Mu_i\}$. This set is contained in $S$ together with its first layer of $M$-cubes, hence in $S^{\sharp}$, since $m_\sharp\ge1$.
\end{enumerate}

The lemma on the interaction terms reads the profile only through clauses (e) and (f), through the fact that its total sums do not involve the profile, and through the fact that $\psi_\varphi$ enters none of its constants; each of these holds as stated for the own-cube profile. Its remaining ingredients read only the values of the profile at the cell-made levels retained by the family, which the own-cube modification does not touch.

The own-cube profile changes one clause in the definition of the family's coupling profiles. Proposition~\ref{4.6:prop-large-field-bound} is stated with the cell-made profiles; with them the widened-support statement enters only through the estimates from which its inputs are drawn. With the own-cube profiles it does not enter at all.

\begin{remark}[The remaining modifications leave the bounds unchanged]
\label{4.6:rem-remaining-modifications}
Besides the modified coupling profiles of Remark~\ref{4.6:rem-coupling-support}, the step of
this section differs from Ba{\l}aban's construction in \cite{B-CMP122a,B-CMP122b} by the
following modifications. None of them calls for a new estimate beyond the transfer of the
bounds in Remark~\ref{4.6:rem-widened-support}, whose proof is outstanding.
\begin{enumerate}
\item[(i)] \emph{Completion copies.} The completion copies enter twice. They are subtracted
in the interaction terms $V(Y)$, and they are added in the term $A'_k$ of the step. They
are bounded by the induction hypothesis at the levels $\le k$. Given that hypothesis, they
leave the bounds unchanged.
\item[(ii)] \emph{The profile members of the family} (clause~(f) of the lemma on the
coupling profiles). They enter
clause~(a) of the polymer-gas representation (Lemma~\ref{4.6:lem-polymer-gas}) only
through the locality of the interaction terms (Remark~\ref{4.6:rem-interaction-locality},
printed without proof in \cite{B-CMP122b}; cited as printed). They also enter, as weights,
the estimates from which the standing bounds of Definition~\ref{4.6:def-standing-bounds}
are derived.
At the levels at which Ba{\l}aban's layers are unions of cells, they leave the bounds
unchanged. At the levels at which those layers are unions of $M$-cubes, their effect is
that of the modified coupling profiles of Remark~\ref{4.6:rem-coupling-support}; the
transfer of the bounds in that case is Remark~\ref{4.6:rem-widened-support}, whose proof
is outstanding.
\item[(iii)] \emph{The vacuum normalisation.} It enters exactly three places:
  \begin{itemize}
  \item the normalisation $E_k(X_a)-E_k(\Lambda\cap X_a)$ of
  Definition~\ref{4.6:def-standing-bounds} in the standing bound on the island operations
  $\mathbf{T}'_k(X_a)$;
  \item the prefactor $\alpha$ of the standing bound on the interaction terms;
  \item clause~(c) of Proposition~\ref{4.6:prop-large-field-bound}, through a factor $2$.
  \end{itemize}
\item[(iv)] \emph{The strip convention $\mathfrak{D}^{\sharp}$.} It is carried throughout.
Its price is contained in the standing bound on the interaction terms. The estimates of
the polymer gas use only the properties \eqref{4.6:eq-strips-bracket-b} of the strips.
\item[(v)] \emph{The incompatibility relation.} Two relations between polymers $X'$,
$X''$ can be used. The first is
touch-incompatibility: the cores share a point. The second is wall-incompatibility: the
sets $X'\setminus\bigcup_iY_i$ and $X''\setminus\bigcup_iY_i$ share a cube or a wall of a
cube, as in \cite[(2.11)]{B-CMP116}. Each
gives a valid polymer gas, provided the same relation is used both for the classes and
for the incompatibility $\iota$. The bounds of Proposition~\ref{4.6:prop-large-field-bound}
on the gas cover either choice.
\end{enumerate}
\end{remark}

\begin{proof}
\emph{(i).} By its definition, each completion copy is a member of the whole-lattice
family at its own level $j\le k$. The induction hypothesis at the levels $\le k$
therefore supplies, for each copy, the bounds
\cite[(2.27), (2.31), (2.42)]{B-CMP119} at level $j$.

These are the same inductive bounds on which Ba{\l}aban's estimates of the interaction
terms in \cite[(1.59)--(1.69)]{B-CMP122b} rest. The standing bound on the interaction
terms of
Definition~\ref{4.6:def-standing-bounds} is derived from those estimates, so it rests on
them too.

Hence both the copies subtracted in $V(Y)$ and the copies added in $A'_k$ are estimated
by bounds already in use. No new estimate is needed, and given the induction hypothesis
the bounds are unchanged.

\emph{(ii).} First consider clause~(a) of Lemma~\ref{4.6:lem-polymer-gas}. There the
profiles enter only through the locality of the interaction terms
(Remark~\ref{4.6:rem-interaction-locality}, printed without proof in \cite{B-CMP122b};
cited as printed). Clause~(f) of the lemma on the coupling profiles is what supplies that
locality for the profile members of the whole-lattice family, as these are fixed in the
construction of that family.

Next consider the estimates from which the standing bounds on $\mathbf{T}'_k$ and on
$V(Y)$ are derived. There, clauses~(a) and~(b) of the lemma on the coupling profiles give
two facts, at the levels at which Ba{\l}aban's layers are unions of cells:
\begin{itemize}
\item the profiles act as weights with values in $[0,1]$;
\item their supports lie inside those layers.
\end{itemize}
Since the weights take values in $[0,1]$ and are supported inside the layers on which
those estimates are made, the estimates apply to the weighted terms as they stand. So the
bounds are unchanged at these levels.

At the levels at which the layers are unions of $M$-cubes, the weights are the modified
coupling profiles of Remark~\ref{4.6:rem-coupling-support}. Their effect on the bounds is
the subject of Remark~\ref{4.6:rem-widened-support}, whose proof is outstanding. Outside
that case the profiles leave the bounds unchanged.

\emph{(iii).} The standing bound on the island operations is stated with the normalisation
$E_k(X_a)-E_k(\Lambda\cap X_a)$ of Definition~\ref{4.6:def-standing-bounds} included. The
standing bound on the interaction terms is
stated with its prefactor $\alpha$ as furnished by the lemma on the vacuum normalisation.
Both bounds therefore enter the argument with the normalisation already in them, as
supplied.

The only other place the normalisation enters is clause~(c) of
Proposition~\ref{4.6:prop-large-field-bound}, where it costs a factor $2$. Nothing else
in the argument depends on it.

\emph{(iv).} The price of the strip convention is given by clause~(3) of the strip theorem.
It is the factor $e^{-\frac12\beta_\sharp\kappa(m_\sharp-1)}$ in the prefactor
$\alpha^{\sharp}$, $\beta_\sharp$ being the constant so denoted in the strip theorem,
subject to the side condition attached to it there. This factor is part
of the standing bound on the interaction terms. The estimates of the polymer gas that
follow use only the properties \eqref{4.6:eq-strips-bracket-b} of the strips.

The division of the terms into the first and the second group is made with the hull
$\boxplus$, not with the strip. For every $i$, the set $Y_i^{\sharp}$ is contained in
$Y_i^{\boxplus}$, by \eqref{4.6:eq-strips-bracket-b}. The union of the hulls
$Y_i^{\boxplus}$ is contained in $(Z_k\cup\bigcup_iY_i)^{\boxplus}$. Hence a set $X$ of
the second group, for which
\begin{equation*}
X\sqcap (Z_k\cup\textstyle\bigcup_iY_i)^{\boxplus}=\emptyset,
\end{equation*}
shares no $M$-cube with any strip: $X\sqcap Y_i^{\sharp}=\emptyset$ for every $i$. This is
what gives the second group the prefactor
$e^{-p_0(g_k)}$ in \eqref{4.6:eq-large-field-bound-a}.

Finally, by \eqref{4.6:eq-strips-bracket-b}, the strips are boxes of $M$-cubes, and strips
of components that share no point are at least one $M$-cube apart. This is what makes
the cores well defined and gives $X_a^{\sharp}\subset\operatorname{core}(X')$ in
Lemma~\ref{4.6:lem-polymer-gas}.

\emph{(v).} The modified incompatibility can be worded in two ways. In the first, two
polymers $X'$, $X''$ are incompatible when their cores share a point; this is the relation
$\iota$. In the second, they are incompatible when
\begin{equation*}
(X'\setminus\textstyle\bigcup_iY_i)\cap(X''\setminus\bigcup_iY_i)
\ \text{contains a cube or a wall of a cube;}
\end{equation*}
this is the wall form of \cite[(2.11)]{B-CMP116}, transported to the present setting.

In \cite[p.~386]{B-CMP122b} the components are the classes of the relation
``intersect, or touch''. The identity \cite[(1.90)]{B-CMP122b} requires two things:
\begin{itemize}
\item ``component'' means touch-component;
\item compatibility means that the cores share no point, contact at a corner included.
\end{itemize}

The construction of this section uses touch both for the classes and for $\iota$,
consistently. Clause~(b) of Lemma~\ref{4.6:lem-polymer-gas} is proved for this choice.
The wall relation, used consistently both for the classes and for the incompatibility,
yields an equally valid gas with different activities.

The bounds of Proposition~\ref{4.6:prop-large-field-bound} on the gas cover either choice,
for two reasons:
\begin{itemize}
\item wall-incompatibility implies touch-incompatibility;
\item the polymers are touch-connected in either case.
\end{itemize}
The two relations must only not be mixed.
\end{proof}

\begin{definition}[Block sides, nests in scope and the size of candidate cubes]\label{4.6:not-block-sides}
We use the geometric conventions fixed at the opening of this section: the torus
$\mathbb{T}=\mathbb{R}^d/P\mathbb{Z}^d$ with $d=4$, carrying the $\ell^\infty$ quotient metric; the fixed
step $k$; the scales $\lambda_i=L^{-(k-i)}$; for each $i$ the grid $\pi_i$ of closed cubes of side
$M\lambda_i$ cornered in $M\lambda_i\mathbb{Z}^d$; the large-field cell side
$a=a_k=M\check{R}_k$, with $\check{R}_k$ a positive integer and $P\in a\mathbb{Z}$; and regions of a
given grain, that is, finite families of closed cells of that grain identified with their unions.

\emph{Units.} Since $\lambda_k=1$, the top grid $\pi_k$ consists of cubes of side $M$. For every
integer $g$ we put
\begin{equation}\label{4.6:eq-block-sides-1}
u_g:=M\lambda_g=ML^{-(k-g)} .
\end{equation}
If $\eta$ denotes the spacing of the finest lattice in \cite{B-CMP99b} and lengths are measured in
units in which $L^k\eta=1$, then $u_g$ is the side $ML^g\eta$ of the big blocks of
\cite{B-CMP99b} and of the big blocks of size $ML^g\eta$ of \cite[p.~229]{B-CMP96}; at the unit
scale, $u_k=M$ is the constant $M_0$ of \cite[(1.118)]{B-CMP95}. Of $L$ we use here that it is an
integer with $L\ge2$, the standing condition on the scale factor.

\emph{Nests in scope.} All nest levels are at most $k$. A nest is a sequence
\begin{equation}\label{4.6:eq-block-sides-2}
\mathbf{\Omega}=(\Omega_0\supseteq\Omega_1\supseteq\dots\supseteq\Omega_k),
\qquad \Omega_{k+1}:=\Omega_{k+2}:=\emptyset ,
\end{equation}
and the nests considered are the members of the class $\mathfrak{N}_{\mathrm{rec}}$ defined later in
this section. That each $\Omega_i$ is a region of grain $\pi_i$ and that the sequence is decreasing
is a hypothesis of this section; for the nests that occur it is the conclusion of the geometric
lemma on the nests of $\mathfrak{N}_{\mathrm{rec}}$. For the covariances of the successive
integrations we admit in addition the wavy augmentation of \cite{B-CMP99b}, in which a new sequence
is formed by adding a further set as a new last member; it is taken here at a truncation level $j$,
the added member being $\Omega_{j+1}$, with added level $j+1\le k$. The last preparatory integration
of \cite{B-CMP122a} runs at $j=k-1$; hence its covariance taken componentwise on
$Z_{j+1}\cap\Omega_{j+1}$, in the form \cite[(2.21)]{B-CMP119}, and its covariance on
$\Omega^c_{j+1}\setminus Z''_{j+1}$, where $Z''_{j+1}$ is the region so denoted in
\cite{B-CMP122a}, are taken at levels at most $k$. Accordingly, region data $\mathfrak{P}$ carry a
declared level at most $k-1$.

\emph{Candidate cubes.} The candidate cubes of scale $g$, introduced below, are cubes of side
$2u_g$; every candidate cube occurring under the scope just described therefore has side at most
$2M$.

\emph{Outside the scope.} The top covariance of the operation $\mathbf{T}$ on the nest augmented by
$\Omega_{k+1}$, whose candidate cubes have scale $k+1$ and side $2ML$, is not treated in this
section. That augmented nest is a member of $\mathfrak{N}_{\mathrm{rec}}$ to which the constructions
of this section are not applied; for it, the additional term in the first case of the reach lemma
later in this section would read $(2L-7)^+M$, where $t^+=\max\{t,0\}$, and that case is not
treated.

\emph{The letter $R$.} In the conventions recalled above the cell multiplier $\check{R}_k$ is
abbreviated $R$; the letter $R$ in the remark on seams later in this section denotes a different
quantity.
\end{definition}

The block sides and grids just introduced have the following properties, used throughout the
section: for every integer $g$,
\begin{equation}\label{4.6:eq-block-sides-3}
u_{g+1}=Lu_g ;
\end{equation}
for $g<k$ the grid $\pi_g$ refines $\pi_{g+1}$; $u_g\le M$ for $g\le k$; and every candidate cube in
scope has side at most $2M$.

\begin{proof}
By \eqref{4.6:eq-block-sides-1},
\begin{equation*}
u_{g+1}=ML^{-(k-g-1)}=L\cdot ML^{-(k-g)}=Lu_g ,
\end{equation*}
which is \eqref{4.6:eq-block-sides-3}.

For the nesting of the grids, let $g<k$; a cube of $\pi_{g+1}$ has the form
$u_{g+1}n+[0,u_{g+1}]^d$ with
$n\in\mathbb{Z}^d$. By \eqref{4.6:eq-block-sides-3} and because $L$ is an integer, this cube equals
\begin{equation*}
\bigcup_{e\in\{0,\dots,L-1\}^d}\bigl(u_g(Ln+e)+[0,u_g]^d\bigr),
\end{equation*}
a union of $L^d$ cubes of side $u_g$ cornered in $u_g\mathbb{Z}^d$, that is, of cubes of $\pi_g$.
Hence every cube of $\pi_{g+1}$ is a union of cubes of $\pi_g$, and $\pi_g$ refines $\pi_{g+1}$.
These grids are well defined on $\mathbb{T}$ for $g\le k$: since $P\in a\mathbb{Z}$ and
$a/u_g=\check{R}_kL^{k-g}$ is a positive integer, $P$ is an integer multiple of $u_g$, so
$u_g\mathbb{Z}^d$ is invariant under translation by $P\mathbb{Z}^d$.

For $g\le k$ we have $k-g\ge0$, and since $L\ge2>1$ this gives $L^{-(k-g)}\le1$; by
\eqref{4.6:eq-block-sides-1}, $u_g\le M$, with equality exactly when $g=k$.

Finally, the scales of the candidate cubes in scope are nest levels. By the scope hypothesis these
levels are at most $k$: the members of a nest \eqref{4.6:eq-block-sides-2} carry levels
$0,\dots,k$, and in the wavy augmentation the added level is $j+1\le k$. A candidate cube of scale
$g\le k$ therefore has side $2u_g\le 2M$ by the preceding paragraph. The candidate cubes of scale
$k+1$, of side $2u_{k+1}=2ML$, belong only to the top covariance of $\mathbf{T}$ on the nest
augmented by $\Omega_{k+1}$, which is outside the scope of this section.
\end{proof}

\begin{definition}[Blocks, membership bits, candidate cubes and their enlargements]
\label{4.6:def-candidate-cubes}
Let the step $k$, the torus $\mathbb{T}$ with its $\ell^\infty$ quotient metric, the enlargements
$A^{+s}$, the scales $\lambda_i$ and the grids $\pi_i$ be as in Definition~\ref{4.6:not-block-sides}.
For $0\le g\le k$ write $u_g=M\lambda_g$ for the side of the cells of $\pi_g$.

\emph{Blocks.} A \emph{$\pi_g$-block} is a cell of $\pi_g$, that is, a closed cube of side $u_g$
cornered in $u_g\mathbb{Z}^d$, for $0\le g\le k$. Every $\pi_{g+1}$-block is the union of $L^d$
$\pi_g$-blocks.

\emph{Membership bits.} For a closed union $A$ of blocks and a block $e$ of any scale, the
\emph{membership bit} $[e\subset A]\in\{0,1\}$ equals $1$ if $e\subset A$ and $0$ otherwise. If
$A$ is a union of $\pi_i$-blocks and $e\in\pi_g$, then
\begin{equation}\label{4.6:eq-candidate-cubes-a}
[e\subset A]=
\begin{cases}
[c\subset A]\ \text{ for every } c\in\pi_0 \text{ with } c\subset e, & g\le i,\\[2pt]
\min_{c\in\pi_0,\ c\subset e}\,[c\subset A], & g>i.
\end{cases}
\end{equation}

\emph{Candidate cubes.} For $0\le g\le k$ put
\begin{equation*}
\mathcal{C}_g:=\bigl\{x+[-u_g,u_g]^d:\ x\in u_g\mathbb{Z}^d\bigr\},
\qquad
\mathcal{C}:=\bigsqcup_{g\le k}\mathcal{C}_g .
\end{equation*}
A member of $\mathcal{C}_g$ is a cube of side $2u_g$, the union of the $2^d$ $\pi_g$-blocks having
its centre $x$ as a corner. Since the sides $u_g$, $0\le g\le k$, are distinct, the \emph{scale}
$j(Q)$ of $Q\in\mathcal{C}$ is the index $g$ for which $Q$ has side $2u_g$; it is determined by
the cube itself. Two candidate cubes of the same scale whose centres are neighbours in
$u_g\mathbb{Z}^d$ overlap in half of each. The members of $\mathcal{C}_g$ are cubes of the kind from
which Ba{\l}aban's covers at scale $g$ are drawn: the cubes of size $2M_0$ centred at points of the
block lattice in \cite[p.~36]{B-CMP95}, the cubes of size $2ML^j\eta$, each a sum of $2^d$ big
blocks with a centre $y$, in \cite[p.~229]{B-CMP96}, and the cubes with centres at points of the
block lattice which are unions of big blocks in \cite{B-CMP99b}; here $\eta$ is the lattice spacing
and $j$ the scale of \cite{B-CMP96}, so that the block side $ML^j\eta$ there is the side $u_j$ of
this chapter, and $M_0$ is the block size of \cite{B-CMP95}. The covers themselves use at each scale
only those cubes whose centres lie in the part of the lattice assigned to that scale. The family of
such cubes, denoted $\mathcal{D}$ in \cite{B-CMP96} and \cite{B-CMP99b}, is called in
\cite{B-CMP99b} a partition of the lattice; since cubes of one scale overlap in half, that word is
to be understood as a cover,
$\mathbb{T}=\bigcup_{\square\in\mathcal{D}}\square$, carrying functions $h_\square$ with
$\sum_{\square\in\mathcal{D}}h_\square^2=1$ \cite[(2.36)]{B-CMP96}.

\emph{Enlargements.} Let $Q=x+[-s,s]^d$ be a block-aligned cube of scale $g$, that is,
$x\in u_g\mathbb{Z}^d$ and $s\in u_g\mathbb{N}$. For an integer $n\ge0$,
\begin{equation}\label{4.6:eq-candidate-cubes-1}
Q^{(n)}:=Q^{+nu_g}=x+[-s-nu_g,\ s+nu_g]^d ,
\end{equation}
and for real $t\ge0$ also $Q^{(t)}:=Q^{+tu_g}$. For a candidate cube ($s=u_g$) and real
$t\in[-1,0)$,
\begin{equation*}
Q^{(t)}:=x+[-(1+t)u_g,\ (1+t)u_g]^d ,
\end{equation*}
the concentric cube of side $(2+2t)u_g$; thus
\begin{equation*}
Q^{(-1/3)}=x+\bigl[-\tfrac23u_g,\tfrac23u_g\bigr]^d,
\qquad
Q^{(-2/3)}=x+\bigl[-\tfrac13u_g,\tfrac13u_g\bigr]^d .
\end{equation*}
For a candidate cube $\square$ we write $\square^{n}:=\square^{(n)}$, the cube of side
$(2+2n)u_{j(\square)}$ with the same centre as $\square$; this is the cube $\square^n$ of
\cite{B-CMP99b}, written there, in the same units as in \cite{B-CMP96}, as of size
$(2+2n)ML^j\eta$.

\emph{The double cube.} For a candidate cube $\square$ we put
\begin{equation}\label{4.6:eq-candidate-cubes-b}
\widetilde{\square}:=\square^{1}=\square^{(1)} .
\end{equation}
This is the cube of \cite[p.~235]{B-CMP96} containing $\square$ in the middle and of size $4M$,
where $\square$ has size $2M$ (sizes there being measured in units in which the block side is $M$).
The cube whose characteristic function is written
$\widetilde{\square}$ in \cite[(3.95)]{B-CMP99b}, where the construction is referred to
\cite[(2.82)]{B-CMP96}, is taken to be the same cube; the identification
\eqref{4.6:eq-candidate-cubes-b} is a reading of that notation. In this chapter
$\widetilde{\square}$ enters no statement comparing the labels of two cubes; it occurs as a fixed
cut-off, in the exactness of the level agreement at a cube stated later in this section, in the
telescoping sum
of \cite[(3.95)]{B-CMP99b}, in the notation of one later proposition of this section, and in a
consistency check on the scales of derived cubes, and each of these holds equally with
$\widetilde{\square}=\square^{2}$.

\emph{Derived cubes.} A cube obtained from a candidate cube by enlargement keeps the scale of that
candidate cube, and its further enlargements are measured in the block side of that scale: for
$\square\in\mathcal{C}$ we put $\square^{2}:=\square^{(2)}$ and $\square^{4}:=\square^{(4)}$, of
scale $j(\square)$, and likewise $\square_1^{2}$, $\square_1^{4}$, of scale $j(\square_1)$, for a
second candidate cube $\square_1$. Since $\square^2$ and $\square^4$ are block-aligned cubes of
scale $j(\square)$ with half-sides $3u_{j(\square)}$ and $5u_{j(\square)}$,
\eqref{4.6:eq-candidate-cubes-1} gives
\begin{equation*}
(\square^{2})^{(5)}=\square^{7},\qquad (\square^{4})^{(4)}=\square^{8},\qquad
(\square^{4})^{(5)}=\square^{9}.
\end{equation*}
In \cite{B-CMP99b}, as we read it, the boundary conditions of the operators localised in
$\square^4$ are placed outside $\square^8$; the Dirichlet domain $\Omega_0(\square^4)$ of the local
sequence $\mathbf{\Omega}(\square^4)$ constructed in this section is contained in
$(\square^4)^{(5)}=\square^9$. Only the containment of that domain in the set $F(\mathfrak{a})$
(the $U$-footprint of the corresponding letter $\mathfrak{a}$, defined below) is used, so the
difference between $\square^8$ and $\square^9$ plays no role.

\emph{Domains.} For a finite family $X$ of block-aligned cubes, each of its own scale, and an
integer $n\ge0$,
\begin{equation*}
X^{(n)}:=\bigcup_{Q\in X}Q^{(n)} ,
\end{equation*}
each $Q^{(n)}$ taken in the scale of $Q$. In \cite{B-CMP99b} the enlargement $\square^n$ is defined
for single cubes; unions of enlarged cubes also occur there, the largest localisation domain there,
that of the operators localised in $\square^{4}$ (or $\square_1^{4}$), being a union of $5^d$ cubes
of that family. The notation $X^{(n)}$ is used here for such unions.
\end{definition}

\begin{proof}[Proof of \eqref{4.6:eq-candidate-cubes-a}]
Let $A$ be a union of $\pi_i$-blocks and let $e\in\pi_g$. Iterating the statement that every
$\pi_{g'+1}$-block is the union of $L^d$ $\pi_{g'}$-blocks, $0\le g'<k$, every $\pi_g$-block is the
union of the $L^{dg}$ $\pi_0$-blocks contained in it; in particular there are $\pi_0$-blocks
$c\subset e$.

Case $g\le i$. The same iteration, for $g'=g,\dots,i-1$, shows that every $\pi_i$-block is the union
of $L^{d(i-g)}$ $\pi_g$-blocks, so $A=e_1\cup\dots\cup e_r$ with $e_1,\dots,e_r\in\pi_g$. Let
$c\in\pi_0$ with $c\subset e$. If $e\subset A$, then $c\subset e\subset A$. Conversely, suppose
$c\subset A$. Since $c\subset e$, the interior of $c$ lies in the interior of $e$. The finer grid
$\pi_0$ refines the coarser grid $\pi_g$ (Definition~\ref{4.6:not-block-sides}), and two distinct
cells of $\pi_g$ meet only in their boundaries; hence for $e_j\neq e$ the intersection $e\cap e_j$
lies in the boundary of $e$ and does not meet the interior of $c$. Every point of the interior of
$c$ lies in $A$, hence in some $e_j$, and by what was just said that $e_j$ is $e$. As the interior
of $c$ is non-empty, $e$ is one of $e_1,\dots,e_r$, so $e\subset A$. Thus $e\subset A$ if and only
if $c\subset A$, that is, $[e\subset A]=[c\subset A]$, for every $\pi_0$-block $c\subset e$.

Case $g>i$. Since $e$ is the union of the $\pi_0$-blocks $c\subset e$, we have $e\subset A$ if and
only if $c\subset A$ for every such $c$, that is,
$[e\subset A]=\min_{c\in\pi_0,\ c\subset e}[c\subset A]$.
\end{proof}

\begin{definition}[Nests, region data, traces and boundary separation]\label{4.6:def-nest-traces}
Fix the step $k$. The torus $\mathbb{T}$, its quotient metric $\operatorname{dist}_\infty$ and the
grids $\pi_g$ ($0\le g\le k$) of closed blocks of side $u_g$ are those of
Definition~\ref{4.6:not-block-sides}. Blocks and the membership bits $[e\subset A]\in\{0,1\}$ are those of
Definition~\ref{4.6:def-candidate-cubes}. For a closed set $V\subset\mathbb{T}$ we write
$\pi_0(V):=\{c\in\pi_0 : c\subset V\}$.

\emph{Nests.} A \emph{nest} is a sequence
\begin{equation*}
\mathbf{\Omega}=(\Omega_0\supseteq\Omega_1\supseteq\cdots\supseteq\Omega_k)
\end{equation*}
in which each $\Omega_i$ is a closed union of $\pi_i$-blocks. We put
$\Omega_{k+1}:=\Omega_{k+2}:=\emptyset$. This is the shape of the decreasing sequences of domains in
\cite{B-CMP99b} and in \cite[(2.1), (2.3)]{B-CMP96}. In \cite{B-CMP99b}, where the sequence of
domains is introduced, we read it, in the notation of that paper, as
\begin{equation*}
\Omega_0\supset\Omega_1\supset\cdots\supset\Omega_k,\quad \Omega_{k+1}=\emptyset,\quad
\Lambda_j=\Omega_j^{(j)}\setminus\Omega_{j+1}^{(j)},\quad \Omega_j\setminus\Omega_{j+1}=B_j(\Lambda_j);
\end{equation*}
there $B_j$ is an operation of that paper and not the layer defined below.

The first member $\Omega_0$ is part of the datum; for the nests of the class $\mathfrak{N}$ one has
$\Omega_0=\mathbb{T}$. It is
not determined from $\Omega_1$. In \cite{B-CMP99b} the set $\Omega_0$ is a collar of $\Omega_1$;
nothing below uses this. The separation condition \cite[(2.2)]{B-CMP96} is not part of the definition
of a nest. Condition \eqref{4.6:eq-nest-traces-a} below is not part of the definition either; it is
assumed as an explicit hypothesis wherever it is used.

\emph{Layers.} For $0\le g\le k$ the \emph{layer} of $\mathbf{\Omega}$ at level $g$ is
\begin{equation*}
B_g:=\Omega_g\setminus\operatorname{int}\Omega_{g+1}.
\end{equation*}
It is a closed set, and $B_k=\Omega_k$.

\emph{Boundary separation.} A nest $\mathbf{\Omega}$ satisfies the \emph{boundary separation condition}
if
\begin{equation}\label{4.6:eq-nest-traces-a}
\Omega_{i+1}\cap\operatorname{cl}(\mathbb{T}\setminus\Omega_i)=\emptyset
\qquad\text{for every } 0\le i\le k,
\end{equation}
that is, if $\Omega_{i+1}\subset\operatorname{int}\Omega_i$ for every $0\le i\le k$.

\emph{Region data.} A \emph{region datum} is a finite family $\mathfrak{P}$ of pairwise disjoint closed
unions of blocks. Its members are called \emph{pieces}. Each piece $P$ carries a declared level
$j(P)\in\{0,\dots,k-1\}$. An example is a single domain which is a closed union of blocks, as the
domain of \cite[\S E]{B-CMP99b}; the region data of the class $\mathfrak{R}_K$ are further examples.

\emph{Traces.} Let $V\subset\mathbb{T}$ be closed. The \emph{trace} of $(\mathbf{\Omega};\mathfrak{P})$
on $V$ is the triple
\begin{equation}\label{4.6:eq-nest-traces-1}
\operatorname{tr}_V(\mathbf{\Omega};\mathfrak{P})
:=\bigl(\tau_V,\ \mathcal{P}_V,\ \sim_V\bigr).
\end{equation}
Here $\tau_V$ is the map $(i,c)\mapsto[c\subset\Omega_i]$ on $\{0,\dots,k\}\times\pi_0(V)$. Next,
$\mathcal{P}_V:=\{c\in\pi_0(V): c\subset\bigcup\mathfrak{P}\}$. Finally, $\sim_V$ is the equivalence
relation on $\mathcal{P}_V$ under which $c\sim_V c'$ if and only if $c$ and $c'$ lie in the same piece.

Two pairs $(\mathbf{\Omega};\mathfrak{P})$ and $(\mathbf{\Omega}';\mathfrak{P}')$ at the same step $k$
have \emph{equal traces on $V$} if the two triples \eqref{4.6:eq-nest-traces-1} coincide. This means
equal maps, equal sets, and equal equivalence relations.

\emph{Wavy nests.} Let $P\in\mathfrak{P}$ be a piece of level $j=j(P)$. The \emph{wavy nest} of
$(\mathbf{\Omega};P)$ is the sequence
\begin{equation}\label{4.6:eq-nest-traces-2}
\widetilde{\mathbf{\Omega}}:=\bigl(\Omega_0,\dots,\Omega_j,\ \widetilde{\Omega}_{j+1}:=P,\
\widetilde{\Omega}_{j+2}:=\emptyset\bigr),
\end{equation}
as in \cite{B-CMP99b}. It is a function of $\mathbf{\Omega}$, of $P$ and of the level $j$ only. Its
trace on $V$ is the map $(i,c)\mapsto[c\subset\widetilde{\Omega}_i]$ on
$\{0,\dots,j+2\}\times\pi_0(V)$, where $\widetilde{\Omega}_i:=\Omega_i$ for $i\le j$.

The following assertions hold.
\begin{enumerate}
\item[(a)] The two formulations of \eqref{4.6:eq-nest-traces-a} are equivalent.
\item[(b)] Suppose that for every $0\le i\le k-1$ the set $\Omega_{i+1}$ has positive distance from
$\mathbb{T}\setminus\Omega_i$ in some metric inducing the topology of $\mathbb{T}$, the distance to
the empty set being $+\infty$. Then $\mathbf{\Omega}$ satisfies \eqref{4.6:eq-nest-traces-a}. For
the sequences of \cite[(2.1)]{B-CMP96} the separation
\begin{equation*}
\operatorname{dist}(\mathbb{T}\setminus\Omega_i,\Omega_{i+1})>RML^i\eta>0
\end{equation*}
of \cite[(2.2)]{B-CMP96}, in the distance used there, gives this hypothesis for $1\le i\le k-1$. In
this item only, $R$ is the positive integer, $M$ the size of big blocks and $\eta$ the lattice spacing
of that display. For $i=0$ the hypothesis holds when $\Omega_0=\mathbb{T}$.
\item[(c)] Suppose that for every $1\le i\le k$ every $\pi_i$-block meeting $\Omega_i$ lies in
$\Omega_{i-1}$. Then $\mathbf{\Omega}$ satisfies \eqref{4.6:eq-nest-traces-a}. For nests of the class
$\mathfrak{N}$ the hypothesis of (c) is the separation property stated with that class; for nests of
the class $\mathfrak{N}_{\mathrm{rec}}$ condition \eqref{4.6:eq-nest-traces-a} is part of the
admissibility proposition for that class.
\item[(d)] The trace $\operatorname{tr}_V(\mathbf{\Omega};\mathfrak{P})$ determines, and is determined
by, the following data over all blocks $e\subset V$ of all scales $g\in\{0,\dots,k\}$:
\begin{itemize}
\item the bits $[e\subset\Omega_i]$, $0\le i\le k$;
\item the bit $[e\subset\bigcup\mathfrak{P}]$;
\item for blocks contained in $\bigcup\mathfrak{P}$, the relation ``lie in a common piece''.
\end{itemize}
\item[(e)] Let $V\subset V'$ be closed. If two pairs have equal traces on $V'$, they have equal traces
on $V$.
\item[(f)] Let $(\mathbf{\Omega};\mathfrak{P})$ and $(\mathbf{\Omega}';\mathfrak{P})$ have equal
traces on $V$, and let $P\in\mathfrak{P}$. Then the wavy nests of $(\mathbf{\Omega};P)$ and
$(\mathbf{\Omega}';P)$ have equal traces on $V$. More generally, let $(\mathbf{\Omega};\mathfrak{P})$
and $(\mathbf{\Omega}';\mathfrak{P}')$ have equal traces on $V$, and let $P\in\mathfrak{P}$ and
$P'\in\mathfrak{P}'$ be pieces of the same level $j$ such that
\begin{equation*}
\{c\in\pi_0(V):c\subset P\}=\{c\in\pi_0(V):c\subset P'\}.
\end{equation*}
Then the wavy nests of $(\mathbf{\Omega};P)$ and $(\mathbf{\Omega}';P')$ have equal traces on $V$.
\end{enumerate}
\end{definition}

\begin{proof}
We first note that $B_g$ is closed. Indeed,
$B_g=\Omega_g\cap(\mathbb{T}\setminus\operatorname{int}\Omega_{g+1})$ is the intersection of a closed
set and the complement of an open set. Moreover $B_k=\Omega_k$, because $\Omega_{k+1}=\emptyset$.

(a) We have
$\operatorname{int}\Omega_i=\mathbb{T}\setminus\operatorname{cl}(\mathbb{T}\setminus\Omega_i)$. Hence a
set is disjoint from $\operatorname{cl}(\mathbb{T}\setminus\Omega_i)$ if and only if it is contained in
$\operatorname{int}\Omega_i$.

(b) Let $0\le i\le k-1$. If $\mathbb{T}\setminus\Omega_i=\emptyset$, its closure is empty and
\eqref{4.6:eq-nest-traces-a} holds for this $i$. Otherwise, in a metric inducing the topology of
$\mathbb{T}$, the distance from a point to a set equals its distance to the closure of that set. So
every $x\in\Omega_{i+1}$ has positive distance from $\operatorname{cl}(\mathbb{T}\setminus\Omega_i)$.
Hence $x$ does not lie in that closure, and \eqref{4.6:eq-nest-traces-a} holds for this $i$. For $i=k$
the left side of \eqref{4.6:eq-nest-traces-a} is empty, since $\Omega_{k+1}=\emptyset$. The last two
sentences of (b) are read off from \cite[(2.2)]{B-CMP96} and from
$\mathbb{T}\setminus\mathbb{T}=\emptyset$.

(c) Let $1\le i\le k$ and $x\in\Omega_i$. The grid $\pi_i$ on the compact torus $\mathbb{T}$ is
finite. Hence the union $F'$ of the $\pi_i$-blocks not containing $x$ is closed, and $x\notin F'$. So
some open neighbourhood $\mathcal{U}$ of $x$ misses $F'$. Since the blocks of $\pi_i$ cover
$\mathbb{T}$, $\mathcal{U}$ is contained in the union of the $\pi_i$-blocks
containing $x$. Each of these meets $\Omega_i$ at $x$, so by hypothesis lies in $\Omega_{i-1}$. Thus
$\mathcal{U}\subset\Omega_{i-1}$ and $x\in\operatorname{int}\Omega_{i-1}$. This proves
$\Omega_i\subset\operatorname{int}\Omega_{i-1}$ for $1\le i\le k$, that is,
\eqref{4.6:eq-nest-traces-a} for $0\le i\le k-1$. For $i=k$ it holds because $\Omega_{k+1}=\emptyset$.

(d) We first check that $\sim_V$ is well defined. Let $c\in\mathcal{P}_V$. The sets $c\cap P$,
$P\in\mathfrak{P}$, are finitely many, closed in $c$, pairwise disjoint, and cover $c$. Since $c$ is
connected, exactly one of them is non-empty, and $c$ lies in exactly one piece. The same argument
applies to any block $e\subset\bigcup\mathfrak{P}$.

Every block of scale $g$ is a union of $\pi_0$-blocks. This follows by applying $g$ times the fact that
a $\pi_{g'+1}$-block is the union of $L^d$ blocks of $\pi_{g'}$, $0\le g'<g$. Hence every piece, and
$\bigcup\mathfrak{P}$,
is a union of $\pi_0$-blocks. Moreover each block $e\subset V$ contains at least one $c\in\pi_0(V)$.

The $\pi_0$-blocks $c\subset V$ are blocks of scale $0$. So the listed data contain $\tau_V$,
$\mathcal{P}_V$ and $\sim_V$.

Conversely, let $e\subset V$ be a block of scale $g$. Since $\Omega_i$ is a union of $\pi_i$-blocks,
\eqref{4.6:eq-candidate-cubes-a} gives the following. If $g\le i$, then
$[e\subset\Omega_i]=[c\subset\Omega_i]$ for any $\pi_0$-block $c\subset e$. If $g>i$, then
$[e\subset\Omega_i]=\min_{c\subset e}[c\subset\Omega_i]$. In both cases the right side is read off
from $\tau_V$, because $c\subset e\subset V$.

Applying \eqref{4.6:eq-candidate-cubes-a} with $i=0$ to the union $\bigcup\mathfrak{P}$ of
$\pi_0$-blocks gives $[e\subset\bigcup\mathfrak{P}]=\min_{c\subset e}[c\subset\bigcup\mathfrak{P}]$
for $g>0$. For $g=0$ this is the definition of $\mathcal{P}_V$.

Finally, let $e$ and $e'$ be blocks contained in $\bigcup\mathfrak{P}$. Then $e$ lies in exactly one
piece, and that piece contains every $\pi_0$-block $c\subset e$. Likewise for $e'$. Hence $e$ and $e'$
lie in a common piece if and only if $c\sim_V c'$ for some, equivalently every, pair of $\pi_0$-blocks
$c\subset e$, $c'\subset e'$.

(e) Since $V\subset V'$, we have $\pi_0(V)\subset\pi_0(V')$. The trace on $V$ is obtained from the trace
on $V'$ in three steps: restrict $\tau_{V'}$ to $\{0,\dots,k\}\times\pi_0(V)$; intersect
$\mathcal{P}_{V'}$ with $\pi_0(V)$; and restrict $\sim_{V'}$ to $\mathcal{P}_{V}$. Equal triples on $V'$
therefore give equal triples on $V$.

(f) The first assertion is the case $\mathfrak{P}'=\mathfrak{P}$, $P'=P$ of the second, in which the
hypotheses on the levels and on the sets of $\pi_0$-blocks hold trivially. We prove the second. Write
$(\tau'_V,\mathcal{P}'_V,\sim'_V)$ for the trace of $(\mathbf{\Omega}';\mathfrak{P}')$ on $V$ and
$\widetilde{\mathbf{\Omega}}'=(\widetilde{\Omega}'_i)$ for the wavy nest of $(\mathbf{\Omega}';P')$.
Let $c\in\pi_0(V)$. For $i\le j$ we have
$[c\subset\widetilde{\Omega}_i]=[c\subset\Omega_i]=\tau_V(i,c)$, and
$[c\subset\widetilde{\Omega}'_i]=\tau'_V(i,c)$. These agree because the traces of the two pairs on $V$
are equal. For $i=j+1$, the bit $[c\subset P]$ equals $[c\subset P']$ by the hypothesis on $P$ and
$P'$. For $i=j+2$ both bits vanish. The two wavy nests have the same length $j+3$, since the levels
agree. Hence their traces on $V$ are equal.
\end{proof}

\begin{definition}[The layer covers and the local sequence of a cube]\label{4.6:def-local-sequence}
Let $\mathbf{\Omega}=(\Omega_0\supseteq\dots\supseteq\Omega_k)$ be a nest in the sense of
Definition~\ref{4.6:def-nest-traces}, with $\Omega_{k+1}=\Omega_{k+2}=\emptyset$. Let the
grids $\pi_g$ of blocks of side $u_g$, the membership bits $[e\subset A]$, the candidate cubes
$\mathcal{C}_g$ and the enlargements $Q^{(t)}$ be those of
Definition~\ref{4.6:def-candidate-cubes}. For $0\le g\le k$ we call
$B_g:=\Omega_g\setminus\operatorname{int}\Omega_{g+1}$ the \emph{layer} of scale $g$.
The set $\Omega_g\setminus\Omega_{g+1}$ is written $B_g(\Lambda_g)$ in \cite{B-CMP99b} and
$B^g(\Lambda_g)$ in \cite[(2.4)]{B-CMP96} (see Definition~\ref{4.6:def-nest-traces}), with
$\Lambda_g$
as in the paragraph on levels below. Since $\Omega_{g+1}\subset\Omega_g$ and $\Omega_{g+1}$ is
closed, this set is contained in the layer $B_g$ and differs from it only by points of the
boundary of $\Omega_{g+1}$.

\emph{Membership.} For $0\le g\le k$ put
\begin{equation*}
\mathcal{D}_g(\mathbf{\Omega}):=\bigl\{Q\in\mathcal{C}_g:\ \exists\, e\in\pi_g,\ e\subset Q,\
[e\subset\Omega_g]=1,\ [e\subset\Omega_{g+1}]=0\bigr\},
\end{equation*}
and let $\mathcal{D}(\mathbf{\Omega})$ be the disjoint union of the families
$\mathcal{D}_g(\mathbf{\Omega})$, $0\le g\le k$. A block $e$ as in the display is called a
\emph{witnessing block} of $Q$. We read the corresponding passage of \cite{B-CMP99b} as saying
that the cubes of scale $j$ are those with at least one of their blocks contained in
$B_j(\Lambda_j)$, and that the set $T_1\setminus\Omega_1$, where $T_1$ denotes the set so written
in \cite{B-CMP99b}, is taken there in place of $\Lambda_0=\Omega_0\setminus\Omega_1$. In the
display this containment of a block $e$ is expressed by the membership bits
$[e\subset\Omega_j]=1$, $[e\subset\Omega_{j+1}]=0$. The family
$\mathcal{D}_g(\mathbf{\Omega})$
consists of the same cubes as the cover of $B^g(\Lambda_g)$ in \cite[p.~229]{B-CMP96}, whose cubes
are formed by $2^d$ $\pi_g$-blocks (the big blocks of \cite[p.~229]{B-CMP96}) around a centre in
$\Lambda_g$, centres on the boundary of that
set included: the centre, in the words of \cite[p.~229]{B-CMP96}, ``more exactly it belongs to
the boundary of this set also''. At $g=0$ and for nests with $\Omega_0=\mathbb{T}$ the displayed
formula is the variant of \cite{B-CMP99b} just described; for a general $\Omega_0$ the family
$\mathcal{D}_0(\mathbf{\Omega})$ defined by the display may differ from that variant, and the
difference plays no role in what follows. Every $Q\in\mathcal{D}_g(\mathbf{\Omega})$ has each of its
witnessing
blocks $e$ contained in the layer $B_g$; more precisely $e\subset\Omega_g$ and
$\operatorname{int}e\cap\Omega_{g+1}=\emptyset$.

\emph{Levels.} The \emph{level datum} of a point or bond $z$ is the collection of the
bits $[c\subset\Omega_i]$, $0\le i\le k$, of the (at most $2^d$) $\pi_0$-blocks $c$ containing $z$.
Write $\Omega_n^{(n)}$ for
the subset of the block lattice of level $n$ whose blocks make up $\Omega_n$ (compare
\cite[(2.2)]{B-CMP96}), and
$\Lambda_n:=\Omega_n^{(n)}\setminus\Omega_{n+1}^{(n)}$, as in \cite{B-CMP99b} and
\cite[(2.3)]{B-CMP96}; in the latter $\Omega_0$ does not occur and $\Lambda_0=\Omega_1^c$.
The level $n(z)$ of $z$, and the relation ``$z\in\Lambda_n$'', are given by the
rule of \cite{B-CMP96, B-CMP99b} applied to the level datum of $z$; one and the same rule is used
for all nests.

\emph{The local sequence.} We read the recipe of \cite{B-CMP99b} as assigning to a cube
$\square\in\mathcal{D}_j$ a sequence $\{\Omega_n(\square)\}_{n=0,\dots,j+1}$ as follows. Here
$\square^{n}$ is the notation of \cite{B-CMP99b} for the concentric enlargement of the cube
$\square$ that is written $Q^{(n)}$, for a cube $Q$, in Definition~\ref{4.6:def-candidate-cubes}.
Either
$\square^{3}\subset B_j(\Lambda_j)$, or $\square^{3}$ meets $B_{j+1}(\Lambda_{j+1})$; in the
second case $\Omega_{j+1}(\square)=\square^{3}\cap B_{j+1}(\Lambda_{j+1})$. Further
$\Omega_j(\square)=\square^{4}$. For $0\le n<j$ the sets $\Omega_n(\square)$ are cubes concentric
with $\square$, with
\begin{equation*}
\operatorname{dist}\bigl(\Omega_n(\square)^c,\Omega_{n+1}(\square)\bigr)=2R_0M_0L^n\eta,
\end{equation*}
and $\Omega_0(\square)\subset\square^{5}$. Here $R_0$ and $M_0$ are fixed numbers of
\cite{B-CMP99b} and $\eta$ is the lattice spacing there.

The construction below is applied to the following cubes $Q=x+[-s,s]^d$ of scale $g$ that are
unions of $\pi_g$-blocks: the candidate cubes
in either of the two cases above, the candidate cubes with $Q^{(3)}\not\subset\Omega_g$, and the
derived cubes $\square^{2}$, $\square^{2}_1$, $\square^{4}$, $\square^{4}_1$ of the modified
covers $\mathcal{D}'$, $\mathcal{D}''$ of \cite{B-CMP99b}. In \cite{B-CMP99b} the recipe above is
formulated for a cube $\square\in\mathcal{D}_j$ in the two cases named there; for every cube just
listed, in particular for the candidate cubes with $Q^{(3)}\not\subset\Omega_g$ and for the
derived cubes, we take the display \eqref{4.6:eq-local-sequence-a} as the definition of the local
sequence. For such a cube we put
\begin{equation}\label{4.6:eq-local-sequence-a}
\begin{gathered}
T(Q):=\bigcup\bigl\{e\in\pi_g:\ e\subset Q^{(3)},\ [e\subset\Omega_{g+1}]=1,\
[e\subset\Omega_{g+2}]=0\bigr\},\\[2pt]
\mathbf{\Omega}(Q):=\bigl(\Omega_0(Q)\supset\dots\supset\Omega_{g-1}(Q)\supset\Omega_g(Q)\\
\qquad\supset\Omega_{g+1}(Q)\supset\Omega_{g+2}(Q)\bigr),\\[2pt]
\Omega_g(Q):=Q^{(4)},\qquad \Omega_{g+1}(Q):=T(Q),\qquad \Omega_{g+2}(Q):=\emptyset,\\[2pt]
\Omega_n(Q):=x+[-s_n,s_n]^d,\qquad
s_n:=s+4u_g+\sum_{m=n}^{g-1}2R_0M_0\lambda_m\qquad(0\le n<g).
\end{gathered}
\end{equation}
Here $\lambda_m:=L^{-(k-m)}$ denotes the scale of level $m$, so that $u_m=M\lambda_m$.
The length $L^n\eta$ of \cite{B-CMP99b} corresponds in these units to $\lambda_n$; by the
definition of $s_n$,
\begin{equation*}
s_n-s_{n+1}=2R_0M_0\lambda_n\quad(0\le n<g-1),\qquad
s_{g-1}-(s+4u_g)=2R_0M_0\lambda_{g-1}.
\end{equation*}
The sequence $\mathbf{\Omega}(Q)$ is called the \emph{local sequence} of $Q$. By its definition it
depends only on the cube $Q$, its scale $g$ and the membership bits of blocks in $\mathbf{\Omega}$,
and not on the family of cubes to which $Q$ belongs. With the numbers $R_0$, $M_0$ of
\cite{B-CMP99b} one has $\Omega_0(Q)\subset Q^{(5)}$; this is the counterpart of the inclusion
$\Omega_0(\square)\subset\square^{5}$ of \cite{B-CMP99b}. If $Q^{(3)}\subset B_g$, in particular
in the first of the two cases above, then $T(Q)=\emptyset$.

The \emph{wavy local sequence} $\widetilde{\mathbf{\Omega}}(Q)$ is given by the same formula
\eqref{4.6:eq-local-sequence-a}, with the nest $\mathbf{\Omega}$ replaced by the wavy nest
$\widetilde{\mathbf{\Omega}}$ of Definition~\ref{4.6:def-nest-traces}.
\end{definition}

\begin{proof}[Proof of the two properties stated in the definition]
\emph{Witnessing blocks lie in the layer.} Let $Q\in\mathcal{D}_g(\mathbf{\Omega})$, and let
$e\in\pi_g$ with $e\subset Q$, $[e\subset\Omega_g]=1$ and $[e\subset\Omega_{g+1}]=0$. The first bit
gives $e\subset\Omega_g$.

For $g=k$ the set $\Omega_{g+1}$ is empty, so that $\operatorname{int}e\cap\Omega_{g+1}=\emptyset$;
let now $g<k$.
The set $\Omega_{g+1}$ is a closed union of $\pi_{g+1}$-blocks, and by
Definition~\ref{4.6:def-candidate-cubes} every $\pi_{g+1}$-block is a union of $L^d$
$\pi_g$-blocks. Hence $\Omega_{g+1}$ is a union of $\pi_g$-blocks. Suppose that
$\operatorname{int}e$ met $\Omega_{g+1}$. Then $\operatorname{int}e$ would meet some $\pi_g$-block
$e'\subset\Omega_{g+1}$. A cell of the grid $\pi_g$ that meets the interior of another cell of
$\pi_g$ coincides with it (Definition~\ref{4.6:def-candidate-cubes}), so $e'=e$. This gives
$e\subset\Omega_{g+1}$, contradicting
$[e\subset\Omega_{g+1}]=0$. Therefore $\operatorname{int}e\cap\Omega_{g+1}=\emptyset$.

Now let $p\in e$. If $p$ were in $\operatorname{int}\Omega_{g+1}$, some open neighbourhood of $p$
would be contained in $\Omega_{g+1}$. Since $e$ is the closure of $\operatorname{int}e$, that
neighbourhood meets $\operatorname{int}e$, which is impossible by what has just been shown. Hence
$e\cap\operatorname{int}\Omega_{g+1}=\emptyset$. Together with $e\subset\Omega_g$ this gives
$e\subset\Omega_g\setminus\operatorname{int}\Omega_{g+1}=B_g$.

\emph{The first case.} Assume $Q^{(3)}\subset B_g$, and suppose that some $e\in\pi_g$ satisfies
$e\subset Q^{(3)}$ and $[e\subset\Omega_{g+1}]=1$. Then $e\subset\Omega_{g+1}$, so
$\operatorname{int}e\subset\operatorname{int}\Omega_{g+1}$, and $\operatorname{int}e\neq\emptyset$.
On the other hand $\operatorname{int}e\subset Q^{(3)}\subset B_g$, and $B_g$ is disjoint from
$\operatorname{int}\Omega_{g+1}$. This is a contradiction. Hence no block enters the union defining
$T(Q)$ in \eqref{4.6:eq-local-sequence-a}, and $T(Q)=\emptyset$.
\end{proof}

\begin{definition}[The localized generators of a cube]\label{4.6:def-local-generators}
Let $\mathbf{\Omega}$ be a nest with region datum $\mathfrak{P}$, and let $Q$ be a cube of its cover
$\mathcal{D}(\mathbf{\Omega})$, with local sequence $\mathbf{\Omega}(Q)$ as in
Definition~\ref{4.6:def-local-sequence} and \eqref{4.6:eq-local-sequence-a}. We write the members
of the local sequence as
$\Omega_0(Q)\supseteq\Omega_1(Q)\supseteq\cdots$, and $\widetilde{\mathbf{\Omega}}(Q)$ for its wavy
version, as there. Let $\eta$ be the spacing of the finest lattice of the sequence, so that the
block lattice $\mathbb{T}^{(n)}$ (Definition~\ref{1.5:def-block-average}) has spacing $L^n\eta$. For
a site $y$ of $\mathbb{T}^{(n)}$, $B^n(y)$ is the block of order $n$ with centre $y$. Let $a_n$ be
the constants of \cite[(3.23)--(3.24)]{B-CMP99b}: $a_0=a_1=a$, and $a_{n+1}$ is obtained from the
preceding constants by the recursion given there. Let $U$ be a configuration on the bonds of
$\Omega_0(Q)$.

For a cube $\square$ of his cover, Ba{\l}aban introduces in \cite{B-CMP99b} operators
$G'_\square(U)$, $C_\square(U)$ and $G_\square(U)$, constructed for the sequence of domains
attached to $\square$. We read the display of $C_\square(U)$ there as
\begin{equation*}
  C_\square(U)=\bigl(Q'(U)\,G'_\square(U)^2\,Q'^*(U)\bigr)^{-1}.
\end{equation*}
These operators are given by the definitions made in \cite[Sect.~A]{B-CMP99b} for the sequence
$\{\Omega_j\}_{j\le k}$, with that sequence replaced by the local one. In the present notation the
cube is $Q$ and the local sequence is $\mathbf{\Omega}(Q)$, and the operators are as follows.

\begin{enumerate}
\item[(a)] \emph{Site sets.} The function $\chi_Q$ is the characteristic function of
$\Omega_0(Q)$; as a
multiplication operator it acts on functions on sites and on functions on bonds (through the
bonds contained in $\Omega_0(Q)$). The set $\Lambda_n(Q)$ consists of the points of
$\mathbb{T}^{(n)}$ in
$\Omega_n(Q)\setminus\Omega_{n+1}(Q)$. The site set of the local sequence is the disjoint union
$\mathfrak{B}_Q:=\bigsqcup_n\Lambda_n(Q)$. More generally, for a nest $\mathbf{\Omega}$ with
members $\Omega_n$, $\mathfrak{B}[\mathbf{\Omega}]$ is the disjoint union over $n$ of the sets of
points of $\mathbb{T}^{(n)}$ in $\Omega_n\setminus\Omega_{n+1}$. Thus
$\mathfrak{B}_Q=\mathfrak{B}[\mathbf{\Omega}(Q)]$.

\item[(b)] \emph{Averaging operators.} The operator
$Q'_Q\colon\ell^2(\Omega_0(Q))\to\ell^2(\mathfrak{B}_Q)$ is defined by
\begin{equation}\label{4.6:eq-local-generators-1}
  (Q'_Q\lambda)(y):=(Q'_n(U)\lambda)(y)\qquad\text{for }y\in\Lambda_n(Q),
\end{equation}
where, reading \cite[(3.18), (3.19)]{B-CMP99b},
\begin{equation}\label{4.6:eq-local-generators-2}
  (Q'_n(U)\lambda)(y)=\sum_{x\in B^n(y)}L^{-nd}\,
  R\bigl(U(\Gamma^{(n)}_{y,x})\bigr)\lambda(x).
\end{equation}
Here $R(X)Y=XYX^{-1}$ \cite[(56)]{B-CMP98}, and $U(\Gamma)$ is the product of the bond variables
along the contour $\Gamma$. The contours are those of \cite{B-CMP98}. For $n=1$ they are the axial
contours $\Gamma_{y,x}$, $x\in B(y)$, of \cite[(58)]{B-CMP98}. For $n\ge2$, $\Gamma^{(n)}_{y,x}$ is
the chain through the centres of the intermediate blocks of \cite[(77)]{B-CMP98}. For $n=0$ the
block $B^0(y)$ is $\{y\}$ and the contour is empty, so $Q'_0(U)$ is the identity.

The operator $Q_Q$ is defined in the same way on functions on bonds. Let $A$ be a function on the
bonds of spacing $\eta$ in $\Omega_0(Q)$, and let $c$ be a bond of $\mathbb{T}^{(n)}$ belonging to
$\Lambda_n(Q)$, in the sense in which \cite[(2.20)]{B-CMP96} writes $b\in\Lambda_j$. Then
$(Q_QA)(c):=(Q_n(U)A)(c)$, where $Q_n(U)$ is the averaging operation of order
$n$ on bond functions of \cite[(3.16)]{B-CMP99b} and $Q_0$ is the identity
\cite[(2.20)]{B-CMP96}. The value $(Q_n(U)A)(c)$ is a sum over the bonds
$b\subset B^n(c_-)\cup B^n(c_+)$, where $c_\pm$ are the endpoints of $c$, of terms in the values
$A(b)$ \cite{B-CMP99b}. This
agrees with the locality of the averages of order $n$ in \cite[p.~31]{B-CMP98}.

\item[(c)] \emph{The operators built from $Q'_Q$.} Let $\mathfrak{a}'$ be the diagonal operator on
$\ell^2(\mathfrak{B}_Q)$ that multiplies by $a_n(L^n\eta)^{d-2}$ on $\Lambda_n(Q)$, that is,
$\mathfrak{a}'=\operatorname{diag}\bigl(a_n(L^n\eta)^{d-2}\bigr)$. Let $\Delta_U$ be the covariant
lattice Laplacian of spacing $\eta$ in the background $U$. We put
\begin{align}
  \Delta'_{a,Q}&:=\chi_Q\bigl(\Delta_U+Q'^{*}_Q\,\mathfrak{a}'\,Q'_Q\bigr)\chi_Q,
  \qquad G'_Q:=(\Delta'_{a,Q})^{-1},\label{4.6:eq-local-generators-3}\\
  T'_Q&:=Q'_Q\,G'^{2}_Q\,Q'^{*}_Q,\qquad C_Q:=(T'_Q)^{-1},\label{4.6:eq-local-generators-4}\\
  P_Q&:=G'_Q\,Q'^{*}_Q\,C_Q\,Q'_Q\,G'_Q,\qquad R_Q:=1-P_Q.\label{4.6:eq-local-generators-5}
\end{align}
The cut by $\chi_Q$ on both sides is the Dirichlet restriction to $\Omega_0(Q)$ described after
\cite[(3.24)]{B-CMP99b}, and $G'_Q$ is the inverse on $\ell^2(\Omega_0(Q))$. The
definition of $P_Q$ and $R_Q$ is the local form of \cite[(3.25)]{B-CMP99b}, where $P=I-R$. Put
\begin{equation*}
  \mathcal{G}_Q:=G'^{2}_Q-G'^{2}_QQ'^{*}_Q\,C_Q\,Q'_QG'^{2}_Q.
\end{equation*}
Since $C_QT'_Q=T'_QC_Q=1$, we have
\begin{equation*}
  Q'_Q\mathcal{G}_Q=Q'_QG'^{2}_Q-T'_QC_QQ'_QG'^{2}_Q=0,
\end{equation*}
and likewise $\mathcal{G}_QQ'^{*}_Q=0$. Moreover $\Delta'_{a,Q}G'_Q=G'_Q\Delta'_{a,Q}=1$. Hence
\begin{equation*}
  \Delta'_{a,Q}\,\mathcal{G}_Q\,\Delta'_{a,Q}
  =1-G'_QQ'^{*}_QC_QQ'_QG'_Q=R_Q.
\end{equation*}
Because $\mathcal{G}_Q$ annihilates $Q'^{*}_Q$ on the right and $Q'_Q$ on the left,
$\Delta'_{a,Q}$ may be replaced on both sides by $\chi_Q\Delta_U\chi_Q$. This is the local form
of the identities \cite[(2.17), (2.26)--(2.27)]{B-CMP96}.

\item[(d)] \emph{The operators on bond functions.} Let $\mathfrak{a}$ be the diagonal operator
that multiplies by $a_n(L^n\eta)^{d-2}$ on the values at bonds of $\mathbb{T}^{(n)}$ in
$\Lambda_n(Q)$. Let $\Delta(U)$ be the operator on bond functions of \cite[(3.26)]{B-CMP99b}, and
let $D$ be the covariant derivative, which maps functions on sites to functions on bonds in the
background $U$. We put
\begin{equation}\label{4.6:eq-local-generators-6}
  \Delta_{a,Q}:=\chi_Q\bigl(\Delta(U)+D\,R_Q\,D^*+Q^*_Q\,\mathfrak{a}\,Q_Q\bigr)\chi_Q,
  \qquad G_Q:=(\Delta_{a,Q})^{-1}.
\end{equation}
This is the local form of \cite[(3.26), (3.27)]{B-CMP99b} and \cite[(2.19)]{B-CMP96}. The
projection that enters $\Delta_{a,Q}$ is the local operator $R_Q$ of item~(c). This is the form
in which the localized operator appears in \cite{B-CMP99b}; it is item~(6) of the lemma later in
this section on the localized operators at the cut-off.

\item[(e)] \emph{Further operators and the block matrices.} We put
\begin{equation}\label{4.6:eq-local-generators-7}
  T_Q:=Q_Q\,G_Q\,Q^*_Q,\qquad H_Q:=G_Q\,Q^*_Q\,T_Q^{-1},\qquad S_Q:=G_Q-H_Q\,T_Q\,H^*_Q.
\end{equation}
The operator $H_Q$ is the local form of \cite[(3.126)]{B-CMP99b}. The block matrices are
\begin{equation}\label{4.6:eq-local-generators-8}
  \mathfrak{Q}_Q:=\begin{pmatrix}\Delta_{a,Q}&Q^*_Q\\ Q_Q&0\end{pmatrix},\qquad
  \mathfrak{Q}'_Q:=\begin{pmatrix}\Delta'_{a,Q}&Q'^{*}_Q\\ Q'_Q&0\end{pmatrix}.
\end{equation}

\item[(f)] \emph{General localized generators.} Let $\mathfrak{K}$ be any of the operator
constructions of \cite[Sects.~A, D, E]{B-CMP99b} that assign an operator to a sequence of domains,
a configuration on its largest member, and region data, and that take the sequence only as an
argument. We put
\begin{equation}\label{4.6:eq-local-generators-9}
  K_Q:=\mathfrak{K}\bigl(\mathbf{\Omega}(Q);\,U\restriction\Omega_0(Q);\,
  \operatorname{tr}_{Q^{(5)}}(\mathbf{\Omega};\mathfrak{P});\,x\bigr),
\end{equation}
where $\operatorname{tr}_{Q^{(5)}}(\mathbf{\Omega};\mathfrak{P})$ is the trace on $Q^{(5)}$ of the
nest $\mathbf{\Omega}$ with its region datum $\mathfrak{P}$, and $x\ge0$ is a parameter that does
not depend on the
nest. The parameter is present in those recipes in which a resolvent $(x+\cdot)^{-1}$ replaces an
inverse. The admission of such resolvent recipes is a reading of Ba{\l}aban's construction; it
is part of the tower form of reading convention $(\rho3)$ of
Definition~\ref{4.5:not-reading-conventions}. The wavy
localized generators are
\begin{equation*}
  \widetilde{K}_Q:=\mathfrak{K}\bigl(\widetilde{\mathbf{\Omega}}(Q);\cdot\bigr),
\end{equation*}
with the remaining arguments as in \eqref{4.6:eq-local-generators-9}.

\item[(g)] \emph{Domain.} All the operators of items (b)--(f) are partial functions of
$U\restriction\Omega_0(Q)$, where $\Omega_0(Q)\subset Q^{(5)}$ \cite{B-CMP99b}. Each is defined at
those configurations at which every matrix inverted in its definition is invertible.

\item[(h)] \emph{Transfer convention.} Let $A$ be an operator on $\ell^2(\mathfrak{B}_Q)$, $B$ an
operator on $\ell^2(\mathfrak{B}[\mathbf{\Omega}])$ for a nest $\mathbf{\Omega}$, and $h$ a bump
function. A composition of $A$ and $B$ through $h$ is taken on the common sites of
$\mathfrak{B}_Q$ and $\mathfrak{B}[\mathbf{\Omega}]$: at a site that belongs to one index set but
not to the other, $h$ is set equal to $0$. Suppose the levels of $\mathbf{\Omega}(Q)$ and of
$\mathbf{\Omega}$ agree near the support of $h$, in the sense of the level-agreement condition of
this chapter. Then this composition is the composition of
operators used in \cite{B-CMP99b}.
\end{enumerate}

In the display of $C_\square$ in \cite{B-CMP99b} the averaging operator $Q'$ carries no subscript.
We read it as the operator $Q'_Q$ of item~(b), constructed for the local sequence. This follows
the words ``constructed for this sequence'' that accompany the display, and the same usage for
operators ``constructed for the sequence $\{\Omega'_j\}$'' in \cite[Corollary~3.6]{B-CMP99b}.
Under the
other reading, $Q'$ is built from the level data of the global sequence on $\Omega_0(Q)$. Then
$C_Q$ depends on the membership bits $[e\subset\Omega_g]$ of blocks $e\subset Q^{(5)}$, instead
of on those of blocks $e\subset Q^{(3)}$. In either reading these blocks lie inside the
nest-reading domain fixed later in this section for the localized generators. The only statement
affected by the choice of reading is the
sharpening from $Q^{(5)}$ to $Q^{(3)}$ in the locality statement for $C_Q$ later in this section.
\end{definition}

\begin{definition}[Surgery of the cover at an anchor cube]\label{4.6:def-cover-surgery}
Let $\mathbf{\Omega}$ be a nest, let $\mathcal{D}(\mathbf{\Omega})$ be its mixed-size cover and
$\mathcal{D}_g(\mathbf{\Omega})$ its layer covers (Definition~\ref{4.6:def-local-sequence}). For a
cube $Q$ let $Q^{n}$ denote its concentric enlargements (Definition~\ref{4.6:def-candidate-cubes}),
let $z(Q)$ denote its centre, and let $\operatorname{dist}_\infty$ denote the $\ell^\infty$
distance. An \emph{anchor} is a cube $\square$ of the cover together with its level, the index $j$
with $\square\in\mathcal{D}_j(\mathbf{\Omega})$; in what follows $j$ denotes the level of the anchor
under discussion. The surgery of
the cover at an anchor is the operation of \cite{B-CMP99b}. It is fixed as follows.

\emph{The case index.} Put
\begin{equation*}
\kappa(\square;\mathbf{\Omega}) :=
\begin{cases}
1 & \text{if no member of } \mathcal{D}_{j+1}(\mathbf{\Omega}) \text{ intersects } \square,\\
2 & \text{otherwise.}
\end{cases}
\end{equation*}

\emph{The closest coarser member.} If $\kappa(\square;\mathbf{\Omega})=2$, the set
\begin{equation*}
\mathcal{M}(\square;\mathbf{\Omega}) := \{\, Q\in\mathcal{D}_{j+1}(\mathbf{\Omega}) :
Q\cap\square\neq\emptyset \,\}
\end{equation*}
is non-empty. Let $\square_1(\square;\mathbf{\Omega})$ be the member of
$\mathcal{M}(\square;\mathbf{\Omega})$ whose centre is $\operatorname{dist}_\infty$-closest to the
centre of $\square$, namely
\begin{equation}\label{4.6:eq-cover-surgery-a}
\begin{aligned}
\operatorname{dist}_\infty\bigl(z(\square_1),z(\square)\bigr)
&= \min_{Q\in\mathcal{M}(\square;\mathbf{\Omega})}
\operatorname{dist}_\infty\bigl(z(Q),z(\square)\bigr),\\
z(\square_1)-z(\square)
&= \operatorname{lexmin}\bigl\{\, z(Q)-z(\square) : Q\in\mathcal{M}(\square;\mathbf{\Omega}),\\
&\qquad\quad \operatorname{dist}_\infty\bigl(z(Q),z(\square)\bigr)
= \operatorname{dist}_\infty\bigl(z(\square_1),z(\square)\bigr) \,\bigr\},
\end{aligned}
\end{equation}
where $\operatorname{lexmin}$ of a finite set of vectors is its least element in the
lexicographic order of the coordinates.
In \cite{B-CMP99b} the cube $\square_1$ is the member of $\mathcal{D}_{j+1}$ closest to the centre
of $\square$. We fix the choice among several members at equal distance by the
lexicographic rule in the second line of \eqref{4.6:eq-cover-surgery-a}.

\emph{The inserted cube.} Put
\begin{equation*}
\square_0 :=
\begin{cases}
\square^{2} & \text{if } \kappa(\square;\mathbf{\Omega})=1,\\
\square_1^{2} & \text{if } \kappa(\square;\mathbf{\Omega})=2.
\end{cases}
\end{equation*}
In both cases the cubes that the surgery removes are those that meet $\square_0$. This is the
reading consistent with the description of $\mathcal{D}'$ in \cite{B-CMP99b} as the cover
$\mathcal{D}$ outside $\square_0$. The removed cubes enter what follows only through the test
$C\cap\square_0=\emptyset$.

\emph{The surgered covers.} Following \cite{B-CMP99b}, put
\begin{align*}
\mathcal{D}'(\square;\mathbf{\Omega}) &:= \{\square_0\} \cup
\{\, C\in\mathcal{D}(\mathbf{\Omega}) : C\cap\square_0=\emptyset \,\},\\
\mathcal{D}''(\square;\mathbf{\Omega}) &:= \{\square_0^{2}\} \cup
\{\, C\in\mathcal{D}(\mathbf{\Omega}) : C\cap\square_0^{2}=\emptyset \,\}.
\end{align*}

\emph{Membership tests.} Unwinding the two definitions just given, for every cube $C$ we have
\begin{align*}
C\in\mathcal{D}'(\square;\mathbf{\Omega}) &\iff C=\square_0 \ \text{ or }\
\bigl(C\in\mathcal{D}(\mathbf{\Omega}) \text{ and } C\cap\square_0=\emptyset\bigr),\\
C\in\mathcal{D}''(\square;\mathbf{\Omega}) &\iff C=\square_0^{2} \ \text{ or }\
\bigl(C\in\mathcal{D}(\mathbf{\Omega}) \text{ and } C\cap\square_0^{2}=\emptyset\bigr).
\end{align*}
Indeed, each right-hand side states that $C$ belongs to one of the two sets whose union defines
the corresponding cover.

Once the anchor datum $(\square,\square_0)$ is given, these tests involve the nest
$\mathbf{\Omega}$ only through the membership $C\in\mathcal{D}(\mathbf{\Omega})$. In particular,
they read no set of removed cubes.

\emph{Surgery data and surgered families.} A \emph{surgery datum} $\mathfrak{s}$ is a finite list
of pairs (anchor, inserted cube) produced by iterating the rule of \cite{B-CMP99b} described
above. The \emph{surgered family} of $\mathbf{\Omega}$ along $\mathfrak{s}$ is
\begin{align*}
\mathcal{F}_{\mathfrak{s}}(\mathbf{\Omega}) := {} & \{\text{inserted cubes of } \mathfrak{s}\}\\
& \cup \{\, C\in\mathcal{D}(\mathbf{\Omega}) : C \text{ misses every inserted cube of }
\mathfrak{s} \,\}.
\end{align*}
The inserted cubes are derived cubes in the sense of Definition~\ref{4.6:def-candidate-cubes}, of
side $6u_j$, $6u_{j+1}$, $10u_j$ or $10u_{j+1}$, where $j$ is the level of the corresponding
anchor. They are never candidate cubes.
\end{definition}

\begin{definition}[One-scale bump profiles and the gluing hypothesis; {\cite[(1.118)]{B-CMP95}}, {\cite[p.~229, (2.36)]{B-CMP96}}]\label{4.6:def-bump-profiles}
Throughout, $\mathcal{C}_g$ is the set of candidate cubes of scale $g$ (of side $2u_g$, centred at the points of $u_g\mathbb{Z}^d$), $\mathcal{C}$ is the union of the $\mathcal{C}_g$, and $Q^{(t)}$ is the concentric enlargement of a cube $Q$ of Definition~\ref{4.6:def-candidate-cubes}, which for a candidate $Q=x+[-u_g,u_g]^d$ and a real $t>-1$ reads
\begin{equation*}
Q^{(t)}=x+\bigl[-(1+t)u_g,\,(1+t)u_g\bigr]^d ,
\end{equation*}
so that the double cube $\widetilde{\square}=\square^{1}$ of \eqref{4.6:eq-candidate-cubes-b} is $\square^{(1)}$ in the notation of enlargements. For a cube $Q$ we write $j(Q)$ for its scale. The nest $\mathbf{\Omega}$, the mixed-size family $\mathcal{D}(\mathbf{\Omega})$ and the level datum of a point are those of Definition~\ref{4.6:def-local-sequence}; the surgery datum $\mathfrak{s}$, the surgered family $\mathcal{F}_{\mathfrak{s}}(\mathbf{\Omega})$ and its inserted cube are those of Definition~\ref{4.6:def-cover-surgery}.

\smallskip
\noindent(a) \emph{The one-variable function.} The letter $h$ denotes the function of \cite[(1.118)]{B-CMP95}: $h\in C_0^\infty(\mathopen{]}-2/3,2/3\mathclose{[})$, $h(t)=1$ for $t\in[-1/3,1/3]$, and $h$ is chosen so that
\begin{equation}\label{4.6:eq-bump-profiles-1}
\sum_{n\in\mathbb{Z}}h^2(t-n)=1\qquad\text{for every } t\in\mathbb{R}.
\end{equation}

\smallskip
\noindent(b) \emph{One-scale profiles.} The \emph{basic profile} of a candidate $Q=x+[-u_g,u_g]^d\in\mathcal{C}_g$ is
\begin{equation}\label{4.6:eq-bump-profiles-a}
\eta_Q(y):=\prod_{\mu=1}^{d}h\Bigl(\frac{y_\mu-x_\mu}{u_g}\Bigr).
\end{equation}
The following properties are read off from \eqref{4.6:eq-bump-profiles-a} and from the properties of $h$ in~(a).
\begin{itemize}
\item[(b1)] $\eta_Q(y)$ depends only on $y$, on the centre $x$ and on the scale $g$; in particular it does not depend on the nest $\mathbf{\Omega}$.
\item[(b2)] Since the support of $h$ is a compact subset of $\mathopen{]}-2/3,2/3\mathclose{[}$, the function $\eta_Q$ vanishes outside a compact subset of $x+\mathopen{]}-\tfrac23u_g,\tfrac23u_g\mathclose{[}^{\,d}$, which is the interior of $Q^{(-1/3)}$. Hence
\begin{equation*}
\operatorname{supp}\eta_Q\subset\operatorname{int}Q^{(-1/3)}\subset Q .
\end{equation*}
\item[(b3)] Since $h=1$ on $[-1/3,1/3]$, every factor in \eqref{4.6:eq-bump-profiles-a} equals $1$ when $|y_\mu-x_\mu|\le u_g/3$ for all $\mu$, that is, when $y\in Q^{(-2/3)}$; hence $\eta_Q\equiv1$ on $Q^{(-2/3)}$.
\item[(b4)] For each scale $g$ separately, write the centres as $x=u_g n$ with $n=(n_1,\dots,n_d)\in\mathbb{Z}^d$. The square of \eqref{4.6:eq-bump-profiles-a} is the product over $\mu$ of $h^2(y_\mu/u_g-n_\mu)$, so the sum over $n\in\mathbb{Z}^d$ factorises into the product over $\mu$ of one-variable sums, each of which equals $1$ by \eqref{4.6:eq-bump-profiles-1} applied at $t=y_\mu/u_g$:
\begin{equation}\label{4.6:eq-bump-profiles-2}
\sum_{x\in u_g\mathbb{Z}^d}\eta_{x+[-u_g,u_g]^d}(y)^2
=\prod_{\mu=1}^{d}\sum_{n\in\mathbb{Z}}h^2\Bigl(\frac{y_\mu}{u_g}-n\Bigr)\equiv1
\qquad\text{on }\mathbb{T}.
\end{equation}
\end{itemize}

\smallskip
\noindent(c) \emph{Rescaling.} In \cite[(1.118)]{B-CMP95} the function
\begin{equation*}
h_z(x)=\prod_{\mu=1}^{d}h\Bigl(\frac{x_\mu-z_\mu}{M_0}\Bigr)
\end{equation*}
is attached to the cube $\square_z$ of size $2M_0$ centred at a point $z$ of the lattice of spacing $M_0$, where $M_0$ denotes the block size of \cite[p.~36]{B-CMP95}, and $\sum_z h_z^2=1$ is stated there. In \cite[p.~229]{B-CMP96} the cubes of scale $j$ are ``of the size $2ML^j\eta$, each cube being a sum of $2^d$ big blocks with a center $y$'' ($\eta$ denoting the lattice spacing in \cite{B-CMP96}), and the sentence ``We construct also the corresponding family of functions $h$ described in (1.118), and rescale them to proper scales.'' follows (the reference is to \cite[(1.118)]{B-CMP95}). In the units of this chapter these cubes are candidates of scale $j$, that is, cubes $y+[-u_j,u_j]^d$ with $y\in u_j\mathbb{Z}^d$: the geometry of \cite[(1.118)]{B-CMP95} with $M_0$ replaced by $u_j$. We read ``rescale them to proper scales'' as the dilation $M_0\mapsto u_j$, so that the function corresponding to the cube $\square=y+[-u_j,u_j]^d$ (here, as in \cite{B-CMP96}, $y$ is the centre and $x$ the variable) is
\begin{equation}\label{4.6:eq-bump-profiles-b}
\prod_{\mu=1}^{d}h\Bigl(\frac{x_\mu-y_\mu}{u_j}\Bigr)=\eta_{\square}(x).
\end{equation}
This dilation has two properties. (i) It carries the cubes of \cite[(1.118)]{B-CMP95} to the cubes of scale $j$, that is, $\square_z\mapsto\square$. (ii) The identity \eqref{4.6:eq-bump-profiles-1} becomes a sum over the lattice $u_j\mathbb{Z}^d$ of centres of that scale, which gives the per-scale identity \eqref{4.6:eq-bump-profiles-2}. Any other dilation violates (i) or (ii).

The reading \eqref{4.6:eq-bump-profiles-b} agrees with the following statement of \cite{B-CMP99b}: ``The characteristic function $1-\widetilde{\square}$ at the beginning of the term, and the function $h_\square$ at the end, restrict a kernel of the term to points separated at least by a distance $ML^j\eta$''. With $\widetilde{\square}=\square^{1}=\square^{(1)}$ as in \eqref{4.6:eq-candidate-cubes-b}, the function $1-\widetilde{\square}$ (one minus the characteristic function of $\widetilde{\square}$, in the notation of the quotation) vanishes on $\square^{(1)}$. The distance from $\square$ to the complement of $\square^{(1)}$ is $u_j$, which corresponds to $ML^j\eta$ in the units of \cite{B-CMP99b}. Hence this separation holds if and only if $\operatorname{supp}h_\square\subset\square$. By (b2) this inclusion holds for the profile \eqref{4.6:eq-bump-profiles-b}.

\smallskip
\noindent(d) \emph{The gluing hypothesis.} \cite[p.~229, (2.36)]{B-CMP96} states that the functions $h_\square$ attached to the mixed-size family $\mathcal{D}$ satisfy
\begin{equation*}
\sum_{\square\in\mathcal{D}}h_\square^2=1 .
\end{equation*}
For a surgered family (written $\mathcal{D}'$ in \cite{B-CMP99b}; here $\mathcal{F}_{\mathfrak{s}}(\mathbf{\Omega})$), \cite{B-CMP99b} uses ``the corresponding partition of unity $\{h'_\square\}_{\square\in\mathcal{D}'}$''.

The rule by which the rescaled functions of \cite[(1.118)]{B-CMP95} are completed to the functions $h_\square$ of the mixed-size family, or of a surgered family, is not displayed in \cite[pp.~230--236]{B-CMP96} or in \cite[pp.~37--40]{B-CMP95}; the one modified prescription on those pages, \cite[p.~235, (2.70)]{B-CMP96}, concerns the operator inverted on each cube and not the functions $h_\square$.

The functions attached to the members of these families in this book cannot be taken to be the profiles $\eta_Q$ of the members $Q$ themselves. First, for the family $\mathcal{D}(\mathbf{\Omega})$ the sum of the squares of the profiles of its members is at least $2$ at every seam, for nests satisfying the separation condition on nests. Second, if the functions $h'_\square$ of a surgered family satisfy $\sum_{\square\in\mathcal{D}'}h'^{2}_\square=1$ and its kept members carry functions with the supports given by \cite[(1.118)]{B-CMP95}, then the function of the inserted cube is not of the form \eqref{4.6:eq-bump-profiles-a}.

What this book uses about these functions is the following hypothesis on their type, denoted $(\mathrm{G\text{-}loc})\wedge(\mathrm{G\text{-}supp})$. For every family
\begin{equation*}
\mathcal{F}\in\{\mathcal{D}(\mathbf{\Omega}),\ \mathcal{F}_{\mathfrak{s}}(\mathbf{\Omega})\}
\end{equation*}
and every $Q\in\mathcal{F}$ a function $h_Q[\mathcal{F}]$ is given, and these functions obey the two conditions displayed below.
\begin{equation}\label{4.6:eq-bump-profiles-c}
\begin{aligned}
&(\mathrm{G\text{-}loc})\quad
h_Q[\mathcal{F}](y)=\Phi^h\Bigl(y;\,Q,\,j(Q);\,\text{the inserted-cube datum of }\mathfrak{s};\\
&\qquad\qquad\bigl\{(Q',[Q'\in\mathcal{F}]):Q'\in\mathcal{C},\ y\in Q'\bigr\};\,
\text{the level datum of }y\Bigr),\\
&(\mathrm{G\text{-}supp})\quad
\operatorname{supp}h_Q[\mathcal{F}]\subset Q\quad\text{for every }Q\in\mathcal{F}\cap\mathcal{C},\\
&\qquad\qquad\operatorname{supp}h_I[\mathcal{F}]\subset I^{(5)}\quad\text{for every inserted cube }I.
\end{aligned}
\end{equation}
The arguments of $\Phi^h$ in $(\mathrm{G\text{-}loc})$ are as follows.
\begin{itemize}
\item The inserted-cube datum of $\mathfrak{s}$ is that of the surgery datum $\mathfrak{s}$ of Definition~\ref{4.6:def-cover-surgery}. For $\mathcal{F}=\mathcal{D}(\mathbf{\Omega})$, where no surgery is made, this argument is the empty datum.
\item $[Q'\in\mathcal{F}]\in\{0,1\}$ is the membership bit of the candidate $Q'$ in $\mathcal{F}$.
\item The level datum of $y$ is that of the point $y$ in Definition~\ref{4.6:def-local-sequence}.
\end{itemize}
In $(\mathrm{G\text{-}loc})$, $\Phi^h$ is one function, the same for every nest $\mathbf{\Omega}$, every family $\mathcal{F}$ and every surgery datum $\mathfrak{s}$. It depends on these only through its displayed arguments, and in this sense it is label-free. Thus the value of $h_Q[\mathcal{F}]$ at a point $y$ reads only which of the candidates containing $y$ are members of $\mathcal{F}$, together with the layer of $y$ itself.

\smallskip
\noindent(e) \emph{The sum of squares, a rule of the type, and the windows.} Whether the identity
\begin{equation*}
\sum_{Q\in\mathcal{F}}h_Q[\mathcal{F}]^2=1 ,
\end{equation*}
which is the form of \cite[p.~229, (2.36)]{B-CMP96} for the family $\mathcal{F}$, holds on admissible nests depends on the particular rule $\Phi^h$; it is a property to be verified for a given rule, and it is not a condition in \eqref{4.6:eq-bump-profiles-c}. The exhibited gluing rule $(\mathrm{G}^{\natural})$ is one rule of the type \eqref{4.6:eq-bump-profiles-c} that meets every constraint on the functions $h$ stated in \cite[(1.118)]{B-CMP95}, \cite[p.~229, (2.36)]{B-CMP96} and \cite{B-CMP99b}. For $\mathcal{F}=\mathcal{D}(\mathbf{\Omega})$ and the rule $(\mathrm{G}^{\natural})$, the identity above holds on nests satisfying the separation condition on nests. The arguments of this book use only the type \eqref{4.6:eq-bump-profiles-c}. No identification of $(\mathrm{G}^{\natural})$ with the functions of \cite[p.~229]{B-CMP96} or of \cite{B-CMP99b} is used.

For every rule of the type \eqref{4.6:eq-bump-profiles-c} the reading window of the bumps equals $2$, an absolute constant. Accordingly we put
\begin{equation}\label{4.6:eq-bump-profiles-3}
r_h:=2,\qquad r_*:=\max\{2,r_h\}=2 .
\end{equation}

\smallskip
\noindent(f) \emph{The cut-off $\zeta_\square$.} The function $\zeta_\square$ of \cite[(3.105)]{B-CMP99b} is used here in the following form. \cite{B-CMP99b} states that $\zeta_\square=1$ on a cube containing $\square$ whose boundary is at a distance at least a stated multiple of $M$ from the boundary of $\square$; the numeral is read here as $2/3$ (any value $\ge2/3$ serves below). It further states that this implies that $\operatorname{supp}h_\square$ is separated from $\operatorname{supp}(1-\zeta_\square)$ by a distance $\ge M$, that $\zeta_\square h_\square=h_\square$, and that the ``localizations introduced by $1-\zeta_\square$ and $h_\square$ are separated by a distance $\ge M$''. We read these statements as follows: $\zeta_\square$ is a fixed $C_0^\infty$ cut-off that depends on the pair $(\square,j(\square))$ only, and not on the nest, with
\begin{equation}\label{4.6:eq-bump-profiles-d}
\zeta_\square\equiv1\ \text{on }\square^{(2/3)},\qquad
\operatorname{supp}\zeta_\square\subset\square^{(2)} .
\end{equation}
The properties in \eqref{4.6:eq-bump-profiles-d} have the following consequences.
\begin{itemize}
\item[(f1)] For $\square\in\mathcal{F}\cap\mathcal{C}$, $(\mathrm{G\text{-}supp})$ gives $\operatorname{supp}h_\square[\mathcal{F}]\subset\square\subset\square^{(2/3)}$, where $\zeta_\square=1$; hence $\zeta_\square h_\square[\mathcal{F}]=h_\square[\mathcal{F}]$, the identity stated in \cite{B-CMP99b}.
\item[(f2)] The function $1-\zeta_\square$ vanishes on $\square^{(2/3)}$. The distance from $\square^{(-1/3)}$ to the complement of $\operatorname{int}\square^{(2/3)}$ is
\begin{equation*}
\Bigl(\tfrac23+\tfrac13\Bigr)u_j=u_j ;
\end{equation*}
hence a function supported in $\square^{(-1/3)}$, as the profile $\eta_\square$ is by (b2), is separated from $\operatorname{supp}(1-\zeta_\square)$ by at least $u_j$. This is the distance $M$ of \cite{B-CMP99b} measured in the $L^j\eta$-scaled units of that passage, that is, $u_j$ here. Here $\tfrac13$ is the margin of \cite[(1.118)]{B-CMP95} inside $\square$. With $\operatorname{supp}h_\square\subset\square^{(-1/3)}$ the separation ``$\ge M$'' is therefore exact.
\end{itemize}
The support condition in \eqref{4.6:eq-bump-profiles-d} is our reading of the statements of \cite{B-CMP99b} quoted above; they also admit the reading $\zeta_\square\in C_0^\infty(\widetilde{\square})$, which gives the narrower inclusion $\operatorname{supp}\zeta_\square\subset\widetilde{\square}$. In either case $\zeta_\square$ reads no membership datum, so its reading radius is $0$.

The function $\zeta_\square$ of \eqref{4.6:eq-bump-profiles-d} is distinct from the functions $\zeta_\square$ of the partition of unity subordinate to the cover $\{\square\}$ in \cite[p.~276]{B-CMP119}.
\end{definition}

\begin{lemma}[Block-locality of the averaging operations; {\cite[(58)--(63), (77)--(81), (88)--(92), (124)--(125); p.~26, p.~31]{B-CMP98}}]
\label{4.6:lem-averaging-locality}
Use the setting of \cite{B-CMP98}:

\begin{itemize}
\item $\Omega\supset\Omega^{(1)}\supset\dots\supset\Omega^{(k)}$ are the lattices of the block hierarchy.
\item For a lattice $\Omega'$ and $y\in\Omega'^{(1)}$, $B(y)$ is the block of $\Omega'$ with centre $y$, and $B^{j}(x)$ is the $j$-fold block with centre $x$.
\item $\Gamma_{y,x}$, $x\in B(y)$, are the contours of \cite[(58)]{B-CMP98}.
\item For a bond $c$ (or $b$), $c_-$, $c_+$ (resp.\ $b_-$, $b_+$) are its initial and final points.
\item $R(X)Y=XYX^{-1}$ \cite[(56)]{B-CMP98}.
\item $V_0$ and $U_0$ are the background configurations.
\item $V_1$ and $U_1$ are the configurations on which the operations in (a) and (c) act relative to
the backgrounds $V_0$ and $U_0$; $U_1$ is defined at bonds of the lattice $\Omega$
\cite[p.~31]{B-CMP98}.
\item $\overline{U}^{\,j}$ denotes the average of order $j$ of \cite{B-CMP98}; $\overline{V}$ denotes
the one-step average of the configuration $V$ of \cite[(59)]{B-CMP98}, which is used in the form
defined there and not reproduced; and $\overline{V_1V_0}$ denotes the average of $V_1V_0$.
\end{itemize}

Then the following hold.
\begin{enumerate}
\item[(a)] \emph{Contour operations and the one-step average of bond variables.} For a configuration $V'$, \cite[(58)]{B-CMP98} reads
\begin{equation*}
(R_{0,y}V')(\Gamma_{y,x})=\prod_{b\subset\Gamma_{y,x}}R\bigl(V_0(\Gamma_{y,b_-})\bigr)V'_b=1,
\qquad x\in B(y),\ x\neq y .
\end{equation*}
The one-step average satisfies, with a function $v$ determined below,
\begin{equation*}
\overline{V}_c(\overline{V}_0)_c^{-1}
=v(c_-)\,(\overline{V_1V_0})_c(\overline{V}_0)_c^{-1}\,\overline{R}_{0,c}\,v^{-1}(c_+),
\qquad \overline{R}_{0,c}=R\bigl((\overline{V}_0)_c\bigr).
\end{equation*}
On functions $v$ the operation acts by $(R_{0,y}v)(x)=R(V_0(\Gamma_{y,x}))v(x)$; for the function
$v$ of the preceding display, \cite[(60)]{B-CMP98} reads
\begin{equation*}
(R_{0,y}v)(x)=v(y)\,(R_{0,y}V_1)(\Gamma_{y,x}),\qquad x\in B(y),\ x\neq y,\ y\in\Omega'^{(1)}.
\end{equation*}
The function $v$ is fixed by the condition
\begin{equation*}
(\overline{R_0v})(y)=v(y)\exp\Bigl[\,i\sum_{x\in B(y)}L^{-d}\,\tfrac{1}{i}\log v^{-1}(y)(R_{0,y}v)(x)\Bigr]=1,
\end{equation*}
which gives
\begin{equation*}
v(y)=\exp\Bigl[-i\sum_{x\in B(y)}L^{-d}\,\tfrac{1}{i}\log (R_{0,y}V_1)(\Gamma_{y,x})\Bigr].
\end{equation*}
Consequently
\begin{align*}
\overline{V}_c(\overline{V}_0)_c^{-1}
={}&\exp\Bigl[-i\sum_{x\in B(c_-)}L^{-d}\,\tfrac{1}{i}\log(R_{0,c_-}V_1)(\Gamma_{c_-,x})\Bigr]
(\overline{V_1V_0})_c(\overline{V}_0)_c^{-1}\\
&\cdot\overline{R}_{0,c}\exp\Bigl[\,i\sum_{x'\in B(c_+)}L^{-d}\,\tfrac{1}{i}\log(R_{0,c_+}V_1)(\Gamma_{c_+,x'})\Bigr].
\end{align*}

\item[(b)] \emph{Averages of functions.} Let $v$ be defined on a lattice $\Omega'$ and let $y\in\Omega'^{(1)}$. Then
\begin{align*}
(\overline{R_0v})(y)&=(\overline{R(V_0)v})(y)
=\bigl\{(\overline{R(V_0)v})(x)\bigr\}_{x\in B(y)}\\
&=v(y)\exp\Bigl[\,i\sum_{x\in B(y)}L^{-d}\,\tfrac{1}{i}\log v^{-1}(y)\,R\bigl(V_0(\Gamma_{y,x})\bigr)\,v(x)\Bigr].
\end{align*}
For a function $u$ on $\Omega$ the averages of higher order are defined recursively by
\begin{gather*}
(\overline{R_0u})(x_1)=(\overline{R(U_0)u})(x_1),\qquad x_1\in\Omega^{(1)},\\
(\overline{R_0u}^{\,j+1})(x_{j+1})=\bigl(\overline{R(\overline{U}_0^{\,j})\,\overline{R_0u}^{\,j}}\bigr)(x_{j+1}),
\qquad x_{j+1}\in\Omega^{(j+1)},
\end{gather*}
and \cite[(81)]{B-CMP98} reads, at order $k$,
\begin{equation*}
(\overline{R_0u}^{\,k})(y)=1,\qquad y\in\Omega^{(k)}.
\end{equation*}
The contour operations of order $j+1$ are $(R_{0,x_{j+1}}U_1)(\Gamma^{(j+1)}_{x_{j+1},x})$ for $x\in B^{j+1}(x_{j+1})$. They are built, as in \cite[(77)]{B-CMP98}, along the chain of block centres $x_{j+1},x_j,\dots,x$.

\item[(c)] \emph{The iterated averaging of bond variables.} The one-step operation is
\begin{align*}
(\overline{\overline{R(V_0)V_1}})_c
&=(\overline{R_{0,c_-}V_1})^{-1}(\overline{V_1V_0})_c(\overline{V}_0)_c^{-1}\,\overline{R}_{0,c}\,
\overline{R_{0,c_+}V_1}\\
&=(\overline{R_{0,c_-}V_1})^{-1}\,\widetilde{V}_1\,\overline{R}_{0,c}\,\overline{R_{0,c_+}V_1},
\end{align*}
where the second equality shows that $\widetilde{V}_1$ stands for
$(\overline{V_1V_0})_c(\overline{V}_0)_c^{-1}$.
Its iterates are
\begin{gather*}
\overline{\overline{U}}_1=\overline{\overline{R(U_0)U_1}},\\
\overline{\overline{U}}_1^{\,j+1}=\overline{\overline{R(\overline{U}_0^{\,j})\,\overline{\overline{U}}_1^{\,j}}};
\end{gather*}
that is, the operation of order $j+1$ is the composition of the one-step operation above, taken
with background $V_0=\overline{U}_0^{\,j}$, with the operation $\overline{\overline{U}}_1^{\,j}$ of
order $j$. For every bond $b\subset\Omega^{(k)}$ these satisfy the identity, called the fundamental
equality in \cite[(92)]{B-CMP98},
\begin{equation*}
(\overline{R_{0,b_-}U_1}^{(k)})^{-1}\,\widetilde{U}_1^{\,k}\,\overline{R}^{\,k}_{0,b}\,
\overline{R_{0,b_+}U_1}^{(k)}=(\overline{\overline{U}}_1^{\,k})_b .
\end{equation*}
For every bond $c\subset\Omega^{(k)}$, the variable $(\overline{\overline{U}}_1^{\,k})_c$ (written
$(\overline{\overline{U}}^{\,k})_c$ in the sentence of \cite[p.~31]{B-CMP98} quoted below) depends
only on the variables $U_b$ with
\begin{equation*}
b\subset B^k(c_-)\cup B^k(c_+).
\end{equation*}

\item[(d)] \emph{The linear part.} Here, in (d) and in the remarks on (d) after the lemma, $A$
denotes the field variable of
\cite[(120)--(123)]{B-CMP98}; it is not the action of Definition~\ref{1.1:def-wilson-action}.
The quantity $Q(V_0,A,c)$ is
\begin{equation*}
Q(V_0,A,c)=\tfrac{1}{i}\log(\overline{\overline{V}}_1)_c ,
\end{equation*}
which is \cite[(121)]{B-CMP98}; here $V_1$ is the configuration of \cite[(120)]{B-CMP98}, used in
the form defined there and not reproduced. Then, by \cite[(122)--(123)]{B-CMP98},
\begin{equation*}
Q(V_0,A,c)=L\,(Q(V_0)A)_c+C(V_0,A,c),
\end{equation*}
and the remainder obeys
\begin{equation*}
|C(V_0,A,c)|\le C_1L^2|A|^2<C_1(L\alpha_1)^2 .
\end{equation*}
The linear form is
\begin{align*}
(Q(V_0)A)_c={}&\sum_{x\in B(c_-)}L^{-(d+1)}(R_{0,c_-}A)([x,x'])\\
&+\sum_{x\in B(c_-)}L^{-(d+1)}\bigl[g(-i\operatorname{ad}_Y)g^{-1}(-i\operatorname{ad}_{Y_x})-1\bigr]
(R_{0,c_-}A)(\Gamma_{c_-,x})\\
&+\sum_{x\in B(c_-)}L^{-(d+1)}\bigl[g(-i\operatorname{ad}_Y)g^{-1}(-i\operatorname{ad}_{Y_x})-1\bigr]
(R_{0,c_-}A)([x,x'])\\
&-\sum_{x\in B(c_-)}L^{-(d+1)}\bigl[g(-i\operatorname{ad}_Y)g^{-1}(i\operatorname{ad}_{Y_x})
e^{-i\operatorname{ad}_Y}-1\bigr]\,\overline{R}_{0,c}\,(R_{0,c_+}A)(\Gamma_{c_+,x'})\\
&+\Bigl[e^{i\operatorname{ad}_Y}-g(-i\operatorname{ad}_Y)\sum_{x\in B(c_-)}L^{-d}g^{-1}(i\operatorname{ad}_{Y_x})\Bigr]
L^{-1}(R_{0,c_-}A)(c).
\end{align*}
Its main term is
\begin{equation*}
(Q_0A)_c=(Q_{V_0}A)_c=\sum_{x\in B(c_-)}L^{-(d+1)}(R_{0,c_-}A)([x,x']).
\end{equation*}
\end{enumerate}

In every one of these formulas, each sum runs over the points $x$ (or $x'$) of a single block: $B(y)$, $B(c_-)$ or $B(c_+)$. The operations of order $j+1$ are compositions, through the block centres, of one-step operations with operations of order $j$; the formulas are written on the fixed block hierarchy
$\Omega\supset\Omega^{(1)}\supset\dots\supset\Omega^{(k)}$ and refer to no sequence of domains.
\end{lemma}

This is \cite[(58)--(63), (77)--(81), (88)--(92), (121)--(125)]{B-CMP98}, together with the sentences of \cite[p.~31, p.~26]{B-CMP98} quoted below; the proofs are those of \cite{B-CMP98}.

On the iterated operation, \cite[p.~31]{B-CMP98} states, of the expression on the left-hand side
of the fundamental equality displayed in (c), which is the right-hand side of \cite[(88)]{B-CMP98}:
\begin{quote}
``We may consider this expression as a new averaging operation of $k$th order acting on a configuration $U_1$ defined at bonds of the lattice $\Omega$. A result of the averaging is a configuration defined at bonds of the lattice $\Omega^{(k)}$.''
\end{quote}
The same page states, after \cite[(91)]{B-CMP98}:
\begin{quote}
``Let us remark that these averages have the same locality properties as the averages $\overline{U}^k$, namely $(\overline{\overline{U}}^k)_c$, $c\subset\Omega^{(k)}$, depends on the variables $U_b$ for $b\subset B^k(c_-)\cup B^k(c_+)$.''
\end{quote}

Proposition~2 of \cite{B-CMP98}, whose bound is \cite[(54)]{B-CMP98}, is proved under the smallness assumption
\begin{equation*}
|U(\partial p)-1|<\alpha_0\eta^2,\qquad \eta=L^{-k},
\end{equation*}
which is \cite[(52)]{B-CMP98}; \cite[(53)]{B-CMP98} is the inductive bound of the same page. On
the locality of the result, \cite[p.~26]{B-CMP98} states:
\begin{quote}
``The result is local in the sense that if $p=\langle x,y,z,w\rangle$, then it is enough to assume (52) for $p\subset B^k(x)\cup B^k(y)\cup B^k(z)\cup B^k(w)$.''
\end{quote}

Several quantities in (a), (c) and (d) are used in the form defined in \cite{B-CMP98} and are not reproduced here:
\begin{itemize}
\item the configuration $V$ and its one-step average $\overline{V}$ in (a);
\item the quantities $\widetilde{U}_1^{\,k}$, $\overline{R}^{\,k}_{0,b}$ and
$\overline{R_{0,b_\pm}U_1}^{(k)}$ of order $k$ in \cite[(88)]{B-CMP98};
\item the auxiliary quantities $Y$, $Y_x$ and the function $g$;
\item the configuration $V_1$ of \cite[(120)--(121)]{B-CMP98};
\item the points $x'$ paired with $x\in B(c_-)$ in the bonds $[x,x']$;
\item the constants $C_1$ and $\alpha_1$ of \cite[(123)]{B-CMP98}.
\end{itemize}
The block-locality read from the display in (d) is the locality of its index sets: every sum in the display runs over $x\in B(c_-)$, and the field $A$ appears explicitly only through $(R_{0,c_-}A)$ and $(R_{0,c_+}A)$ evaluated at $[x,x']$, $\Gamma_{c_-,x}$, $\Gamma_{c_+,x'}$ or $c$. The operator brackets are formed from $Y$, $Y_x$ and $g$, which are used in the form defined in \cite{B-CMP98} and not reproduced here, so that the part of this reading that concerns them is taken from \cite{B-CMP98} as stated.

\begin{definition}[Letters, words and footprints of the expansions]\label{4.6:def-letters-words}
Fix the level $k$ and a nest $\mathbf{\Omega}$, with its wavy version
$\widetilde{\mathbf{\Omega}}$, the mixed-size cover $\mathcal{D}(\mathbf{\Omega})$ and its
layers $\mathcal{D}_g$, as in Definition~\ref{4.6:def-local-sequence}. For a cube
$C\in\mathcal{C}_g$ (so in particular for $C\in\mathcal{D}_g$) we write $j(C):=g$ for its
scale. A \emph{candidate} cube is a member of some $\mathcal{C}_g$; the \emph{derived} cubes
$\square^{2}$, $\square^{4}$ and the \emph{inserted} cube $\square_0$ are those produced by
the surgeries of Definition~\ref{4.6:def-cover-surgery}.

\emph{Ingredients.} The \emph{ingredient list} $\mathfrak{J}$ consists of the following
operators and functions.
\begin{itemize}
\item The memberships and level data of Definition~\ref{4.6:def-local-sequence}; the local
sequences of that definition and the localized generators of
Definition~\ref{4.6:def-local-generators}, at candidate and at derived cubes, for
$\mathbf{\Omega}$ and for $\widetilde{\mathbf{\Omega}}$.
\item The bumps and the cut-offs $\zeta_{\square}$ of
Definitions~\ref{4.6:def-cover-surgery} and~\ref{4.6:def-bump-profiles}.
\item The explicit semi-local operators $D$, $D^*$ of \cite[(3.3)--(3.9)]{B-CMP99b},
$\Delta_U$ of \cite[(3.23)]{B-CMP99b}, $\Delta(U)$, $\Delta'(U)$ of
\cite[(3.10)]{B-CMP99b} and $J$ of \cite[(3.11)]{B-CMP99b}.
\item The global averaging operators $Q_n[\mathbf{\Omega}]$, $Q'_n[\mathbf{\Omega}]$ and
their adjoints, with their index sets $\Lambda_n[\mathbf{\Omega}]$, and the weights $a_n$.
\item The commutator $K'(h):=h\Delta'_a-\Delta'_a h$ of \cite[(3.88)]{B-CMP99b}. Up to the
averaging term it is the operator denoted $K(h)$ in \cite[(2.38)]{B-CMP96}, given there by
\begin{equation*}
(K(h)\lambda)(x)=\sum_{b\in\operatorname{st}(x)}(\partial h)(b)(\partial\lambda)(b)
-(\Delta h)(x)\lambda(x),
\end{equation*}
where $\operatorname{st}(x)$ is the set of bonds containing the site $x$, $\Delta$ is the
lattice Laplacian and $\lambda$ is a function on the sites, in the notation of
\cite{B-CMP96}; these two symbols carry this meaning only in the present display.
\item The commutator $K(h):=h\mathfrak{L}-\mathfrak{L}h$ with
$\mathfrak{L}:=\Delta(U)+DD^*+Q^*\mathfrak{a}Q$, as in \cite[(3.100)--(3.104)]{B-CMP99b};
inside $\mathfrak{L}$ the symbol $\mathfrak{a}$ is the diagonal operator built from the
weights $a_n$, and not a letter in the sense defined below.
\item The operator $P_Q(\partial h):=[DP_QD^*,h]$ of \cite[(3.101)]{B-CMP99b}, taken with
the localized operator $P_Q$ of Definition~\ref{4.6:def-local-generators}.
\item The following operators of \cite{B-CMP99b}, which we call the
\emph{tower vertices}: $\Delta'_\pi$ of
\cite[(3.117)--(3.120)]{B-CMP99b}; the vertex $\mathsf{C}^{(2)}$, written $C^{(2)}$ in
\cite{B-CMP99b}; $\Delta^{(2)}$ and
$\widetilde{\Delta}^{(2)}$ of \cite[(3.135)--(3.136)]{B-CMP99b}; the constraint pair
$(Q,Q^*)$ of \cite[(3.110)]{B-CMP99b}; $D^{(2)}$ of \cite{B-CMP99b}; the parametrization
operator of \cite[(3.157)]{B-CMP99b}; $\mu(B)$ of \cite[(3.169)]{B-CMP99b}; $\lambda(A)$ of
\cite[(3.184)]{B-CMP99b}; and $(I+D\mu)Q$ of \cite[(3.185)]{B-CMP99b}. In each of them every
global operator is substituted by its own expansion.
\item The characteristic functions $\widetilde{\square}$, $1-\widetilde{\square}$,
$1-\zeta_{\square}$, $\chi_P$ for $P$ a piece, and the characteristic function $\Delta(y)$.
\end{itemize}

\emph{Spreads.} The \emph{spread} of an operator at a point $y$ is the largest
$\ell^\infty$-distance from $y$ at which its kernel at $y$ reads its argument. A multiplier
has spread $0$. The operators $D$, $D^*$, $\Delta_U$, $\Delta(U)$, $\Delta'(U)$ and $J$
read within one step of the $n$-lattice at a point of level $n$; their spread there is at
most $\lambda_n$, where
\begin{equation*}
\lambda_n\le 1\le M.
\end{equation*}
Fix a scale $g$. At a point of level $n\le g+1$ the spread is at most
\begin{equation*}
\lambda_{g+1}=L\lambda_g\le M\lambda_g=u_g,
\end{equation*}
given $M\ge L$. The operators $Q_n$, $Q'_n$, $Q^*$,
$S(\partial h)$ (the operator of \cite{B-CMP99b} built from the bond differences
$\partial h$ of a bump $h$), the parametrization operator of \cite[(3.157)]{B-CMP99b},
$\mu(\cdot)$ of \cite[(3.169)]{B-CMP99b} and $(I+D\mu)Q$ spread by one averaging block $B^n$
of the level $n$ of the point; the side of this block is $\lambda_n$, which is at most $u_g$
at levels $n\le g+1$, given $M\ge L$ ($M=L^{m'}$, Notation~\ref{0.2:not-glossary}; the
exponent is written $m$ in \cite{B-CMP99b}). Let $\square$ be a cube of scale $g$. Under the
level agreement condition $(\mathrm{LA}_{\square})$ at $\square$, the level $n(y)$ of every
point $y$ of the neighbourhood of $\square$ satisfies
\begin{equation}\label{4.6:eq-letters-words-c}
n(y)\le g+1;
\end{equation}
we call \eqref{4.6:eq-letters-words-c} the \emph{scale condition} at $\square$, and we refer
to it and to $(\mathrm{LA}_{\square})$ together by the single label
\eqref{4.6:eq-letters-words-c}. In particular, at the
points of the neighbourhood of a cube of scale $g$ every explicit operator listed above has
spread at most $u_g$.

\emph{Closure.} The \emph{closure} $\Theta$ of $\mathfrak{J}$ is the smallest class of
operators that contains $\mathfrak{J}$ and is closed under the following six nest-free
operations of \cite{B-CMP99b}:
\begin{itemize}
\item[(o1)] products along chains of cubes \cite{B-CMP99b};
\item[(o2)] insertion of characteristic functions \cite{B-CMP99b};
\item[(o3)] commutation with a bump or with a cut-off $\zeta_{\square}$ \cite{B-CMP99b};
\item[(o4)] differences over nested cubes \cite[(3.97)]{B-CMP99b};
\item[(o5)] substitution of a global operator by its own expansion, or by a word of it: for
$G'$, for $P$ (the global operator whose localized version is $P_Q$), and for the operators
composed of the tower vertices \cite{B-CMP99b};
\item[(o6)] grouping of adjacent factors into one factor \cite{B-CMP99b}.
\end{itemize}

\emph{Letters.} A \emph{letter} $\mathfrak{a}$ consists of four components: its type, its
anchor cubes $E_1,E_2,\dots$, its constituent cubes $C_1,\dots,C_{m_{\mathfrak{a}}}$, and
its $\alpha$-data, a fourth component which together with the other three determines
$\mathfrak{a}$ (the letter $\alpha$ in this name is unrelated to the radii $\alpha$ of
Notation~\ref{0.2:not-glossary}). The type states which of the
operations (o1)--(o6) are applied to which
elements of $\mathfrak{J}$; the anchor cubes are the anchors of all the surgery data of
Definition~\ref{4.6:def-cover-surgery} on which the families and bumps entering
$\mathfrak{a}$ are built, an anchor $E$ being a member of $\mathcal{D}_{j(E)}$; the
constituent cubes are the cubes, candidate, derived or inserted, at which the elements of
$\mathfrak{J}$ entering $\mathfrak{a}$ are localized. The operator of $\Theta$ so described
is the \emph{value} of $\mathfrak{a}$. The \emph{localization domain} (the domain $X$ of
\cite{B-CMP99b}) and the \emph{$U$-footprint} of $\mathfrak{a}$ are
\begin{equation*}
X(\mathfrak{a}):=\bigcup_{i=1}^{m_{\mathfrak{a}}}C_i,\qquad
F(\mathfrak{a}):=\bigcup_{i=1}^{m_{\mathfrak{a}}}C_i^{(5)}.
\end{equation*}
The \emph{index-reading domain} and the \emph{nest-reading domain} of $\mathfrak{a}$ are
\begin{align*}
\rho_\iota(\mathfrak{a})&:=F(\mathfrak{a})\cup
\bigcup_{E\ \text{anchor of}\ \mathfrak{a},\ j(E)<k}E^{+2\min\{u_{j(E)+1},\,M\}},\\
\rho(\mathfrak{a})&:=\rho_\iota(\mathfrak{a})
\cup\bigcup_{C_i\ \text{candidate}}C_i^{+r_hM}
\cup\bigcup_{I\ \text{inserted constituent}}\bigl(I^{(5)}\bigr)^{+r_hM},
\end{align*}
where $r_h$ is the bump window. An anchor $E$ of the top scale $j(E)=k$ contributes no
second member to $\rho_\iota(\mathfrak{a})$: since $\mathcal{D}_{k+1}=\emptyset$, the case
test $\mathrm{CASE}(E;\mathbf{\Omega})$ of Definition~\ref{4.6:def-cover-surgery} returns
case~1 identically and is nest-free there. One has the following inclusion, and the
\emph{reading radius} is the number $r_*$ defined beside it:
\begin{equation}\label{4.6:eq-letters-words-a}
\rho(\mathfrak{a})\subset F(\mathfrak{a})^{+2M},\qquad r_*:=\max\{2,r_h\}=2.
\end{equation}

\emph{Words.} A \emph{word} $\omega$ is either a finite sequence
$(\mathfrak{a}_0,\dots,\mathfrak{a}_\ell)$ of letters with
$X(\mathfrak{a}_{i-1})\cap X(\mathfrak{a}_i)\neq\emptyset$ for $1\le i\le \ell$
\cite{B-CMP99b}, or a finite tree of letters \cite{B-CMP99b}. Its value $K_\omega$ is the
ordered product of the values of its letters if $\omega$ is a sequence, and the contraction
of these values along the tree if $\omega$ is a tree. The sets $X(\omega)$, $F(\omega)$,
$\rho_\iota(\omega)$ and $\rho(\omega)$ are the unions of the corresponding sets of the
letters of $\omega$; they satisfy $\rho(\omega)\subset F(\omega)^{+2M}$. The
\emph{footprint} of $\omega$ in the sense of \cite{B-CMP99b} is $[\omega]:=F(\omega)$.

\emph{Outermost constituents and margin.} $C_{\mathrm{first}}(\omega)$ is the cube of the
leftmost bump of the leftmost letter of $\omega$, and $C_{\mathrm{last}}(\omega)$ is the
cube of the rightmost bump of its rightmost letter; both are members of the cover
$\mathcal{D}(\mathbf{\Omega})$ before surgery, as is shown in the support lemma for words
below. The \emph{margin} of $\omega$ is
\begin{equation}\label{4.6:eq-letters-words-b}
\mu(\omega):=3\cdot\min\bigl\{u_{j(C_{\mathrm{first}}(\omega))},\,
u_{j(C_{\mathrm{last}}(\omega))}\bigr\}.
\end{equation}
The margin $\mu(\omega)$ of a word is not to be confused with the operator $\mu(\cdot)$ of
\cite[(3.169)]{B-CMP99b}.
\end{definition}

\begin{proof}[Proof of the inclusions in Definition~\ref{4.6:def-letters-words}]
Let $\mathfrak{a}$ be a letter with constituent cubes $C_1,\dots,C_{m_{\mathfrak{a}}}$, and
write $F=F(\mathfrak{a})$. The enlargement $A\mapsto A^{+s}$ is monotone: $A\subset B$
implies $A^{+s}\subset B^{+s}$, and $s\le t$ implies $A^{+s}\subset A^{+t}$. Every cube is
contained in its concentric enlargements, so $C_i\subset C_i^{(5)}\subset F$ for every $i$.
The bump window $r_h$ of the bumps of Definitions~\ref{4.6:def-cover-surgery}
and~\ref{4.6:def-bump-profiles} satisfies $r_h\le 2$; hence
$r_*=\max\{2,r_h\}=2$, as in \eqref{4.6:eq-letters-words-a}, and
$r_hM\le 2M$. We bound the members of $\rho(\mathfrak{a})$ one at a time.

First, $F\subset F^{+2M}$.

Second, let $C_i$ be a candidate constituent. By the two monotonicity properties and
$C_i\subset C_i^{(5)}\subset F$,
\begin{equation*}
C_i^{+r_hM}\subset C_i^{+2M}\subset\bigl(C_i^{(5)}\bigr)^{+2M}\subset F^{+2M}.
\end{equation*}

Third, let $I$ be an inserted constituent. Then $I^{(5)}\subset F$ by the definition of
$F(\mathfrak{a})$, and therefore
\begin{equation*}
\bigl(I^{(5)}\bigr)^{+r_hM}\subset\bigl(I^{(5)}\bigr)^{+2M}\subset F^{+2M}.
\end{equation*}

Fourth, let $E$ be an anchor of $\mathfrak{a}$ with $j(E)<k$. Every anchor cube of a letter
is one of its constituent cubes, by part~(a) of the lemma on anchors and constituents
below. Hence $E\subset E^{(5)}\subset F$, and, since
$\min\{u_{j(E)+1},M\}\le M$,
\begin{equation*}
E^{+2\min\{u_{j(E)+1},M\}}\subset E^{+2M}\subset F^{+2M}.
\end{equation*}

The set $\rho_\iota(\mathfrak{a})$ is the union of the members treated in the first and
fourth steps, and $\rho(\mathfrak{a})$ is the union of $\rho_\iota(\mathfrak{a})$ with the
members treated in the second and third steps. Therefore
$\rho_\iota(\mathfrak{a})\subset\rho(\mathfrak{a})\subset F(\mathfrak{a})^{+2M}$, which is
the first half of \eqref{4.6:eq-letters-words-a}.

Let $\omega$ be a word and let $\mathfrak{a}$ run over its letters. The enlargement commutes
with unions, because a point lies within distance $2M$ of a union exactly when it lies
within distance $2M$ of one of its members. Hence
\begin{equation*}
\rho(\omega)=\bigcup_{\mathfrak{a}}\rho(\mathfrak{a})
\subset\bigcup_{\mathfrak{a}}F(\mathfrak{a})^{+2M}
=\Bigl(\bigcup_{\mathfrak{a}}F(\mathfrak{a})\Bigr)^{+2M}=F(\omega)^{+2M}.
\end{equation*}

The factor $3$ in \eqref{4.6:eq-letters-words-b} comes from the following arithmetic. Let
$C$ be $C_{\mathrm{first}}(\omega)$ or $C_{\mathrm{last}}(\omega)$, and let $h_C$ be its
bump. By \eqref{4.6:eq-bump-profiles-c} of Definition~\ref{4.6:def-bump-profiles},
$\operatorname{supp}h_C\subset C$. On each side of this bump stands at most one flanking
explicit operator. When the scale condition \eqref{4.6:eq-letters-words-c} holds at $C$,
the spread of this operator at the points of the neighbourhood of $C$ is at most
$u_{j(C)}$, by the spread bounds above for the one-step and for the one-block operators.
This puts the rows, respectively the columns,
of the outermost factors in
\begin{equation*}
C^{+u_{j(C)}}\subset C^{+2u_{j(C)}}=C^{(2)},
\end{equation*}
the second concentric enlargement of $C$. The distance property of the concentric
enlargements of grid cubes gives, for the cube $C$ of scale $j(C)$,
\begin{equation*}
\operatorname{dist}\bigl(C^{(2)},\mathbb{T}\setminus C^{(5)}\bigr)=3u_{j(C)},
\end{equation*}
and this is the factor $3$ in \eqref{4.6:eq-letters-words-b}. The remaining steps of the
support estimate, in which the margin
\eqref{4.6:eq-letters-words-b} is used, are carried out in the support lemma for words
below.
\end{proof}

\begin{definition}[Locality of atom families under equal nest traces]\label{4.6:def-atom-locality}
Recall that an atom rule, in the sense of the definition of atom rules earlier in this section,
assigns to every operator $K$ of the list of operators of the step enumerated earlier in this
section and to every nest in force $\mathbf{\Omega}\in\mathfrak{N}_{\mathrm{rec}}(\mathcal{E})$,
where $\mathfrak{N}_{\mathrm{rec}}(\mathcal{E})$ is the class of nests in force enumerated earlier
in this section for the data $\mathcal{E}$ of the step:
\begin{itemize}
\item an index set $\mathcal{W}_K(\mathbf{\Omega})$ of words $\omega$;
\item footprints $[\omega]$, which are regular closed sets with connected interior;
\item margins $\mu(\omega)>0$;
\item atoms $K_\omega[\mathbf{\Omega}]$, which are functions of the local datum $(\mathbf{U},\mathbf{J})$
on the footprint $[\omega]$.
\end{itemize}
Nests, region data and traces are those of Definition~\ref{4.6:def-nest-traces}.

In this definition the footprint $[\omega]$ of a word $\omega$ is taken to be the footprint
$F(\omega)$ of Definition~\ref{4.6:def-letters-words}. For a set $W''\subset\mathbb{T}$ we put
\begin{equation*}
K^{W''}[\mathbf{\Omega}]:=\sum_{\substack{\omega\in\mathcal{W}_K(\mathbf{\Omega})\\ F(\omega)\subset W''}}
K_\omega[\mathbf{\Omega}] ,
\end{equation*}
the series being absolutely convergent, since it is a sub-series of the series
$\sum_{\omega\in\mathcal{W}_K(\mathbf{\Omega})}K_\omega[\mathbf{\Omega}]=K[\mathbf{\Omega}]$, which
converges absolutely by the convergence requirement in the definition of an atom rule. For a
set $A$ and $s\ge 0$, $A^{+s}$ denotes its $\ell^\infty$ enlargement by $s$, and $M=L^{m'}$ is the
auxiliary parameter of the glossary (Notation~\ref{0.2:not-glossary}).

Let $r\ge 0$. The atom rule is \emph{local of radius $r$} if the following holds for every
operator $K$ of the list of operators of the step, every region $W$, and every two nests in force
of the same kind, $\mathbf{\Omega}\in\mathfrak{N}_{\mathrm{rec}}(\mathcal{E})$ with region datum
$\mathfrak{P}$ and $\mathbf{\Omega}'\in\mathfrak{N}_{\mathrm{rec}}(\mathcal{E}')$ with region datum
$\mathfrak{P}'$; here $\mathfrak{N}_{\mathrm{rec}}(\mathcal{E})$ denotes the class of nests in force
enumerated earlier in this section for the data $\mathcal{E}$ of the step, and the data
$\mathcal{E},\mathcal{E}'$ may differ. Whenever
\begin{equation*}
\operatorname{tr}_{W^{+rM}}(\mathbf{\Omega};\mathfrak{P})
=\operatorname{tr}_{W^{+rM}}(\mathbf{\Omega}';\mathfrak{P}') ,
\end{equation*}
the following three requirements hold.
\begin{itemize}
\item[(a)] The index sets agree on $W$:
\begin{equation*}
\{\omega\in\mathcal{W}_K(\mathbf{\Omega}) : F(\omega)\subset W\}
=\{\omega\in\mathcal{W}_K(\mathbf{\Omega}') : F(\omega)\subset W\}.
\end{equation*}
\item[(b)] For every $\omega$ in the set of (a), $K_\omega[\mathbf{\Omega}]=K_\omega[\mathbf{\Omega}']$
as functions on the local class of $F(\omega)$ (defined below).
\item[(c)] For every $\omega$ in the set of (a) and every pair $b,b'$ of arguments of the kernel
of $K_\omega[\mathbf{\Omega}]$ with $K_\omega[\mathbf{\Omega}](b,b')\neq 0$, both $b$ and $b'$ lie
in $F(\omega)$ at distance at least $\mu(\omega)$ from $\partial F(\omega)$.
\end{itemize}

For $r=0$, with the footprint taken to be $F(\omega)$, this is the locality property in the
definition of atom rules. Enlarging $r$ weakens the property asked of the rule: for $r'\ge r$,
requirements (a)--(c) are then asked only of pairs of nests whose traces agree on the set
$W^{+r'M}\supseteq W^{+rM}$, and equal traces on $W^{+r'M}$ restrict to equal traces on
$W^{+rM}$ (Definition~\ref{4.6:def-nest-traces}).

An equivalent form of radius zero takes as
footprint of $\omega$ the reading domain $\rho(\omega)$ of Definition~\ref{4.6:def-letters-words}
in place of $F(\omega)$. This form is not adopted here.

The local class is defined as follows. In the local datum $(\mathbf{U},\mathbf{J})$ of an atom,
$\mathbf{U}$ stands for the complexified configuration of \cite[(3.37)]{B-CMP99b}, which has the
form $U'U$, the product of the configurations $U$ and $U'$ of that display, and $\mathbf{J}$
stands for the current $J=D^*\eta^{-2}\operatorname{Im}\partial U$ \cite[(3.11)]{B-CMP99b} of
Definition~\ref{4.6:def-letters-words}, where $\eta$ is the lattice spacing in the letter of
\cite{B-CMP99b} (written $\xi$ in Definition~\ref{4.4:def-regular-backgrounds}) and $\partial U$
denotes the plaquette variables of $U$; the current $J$ is a function of $U$. The
\emph{local class} of $F(\omega)$ is the set of configurations on $F(\omega)$ that obey
\cite[(3.35), (3.37)]{B-CMP99b} on the class cubes contained in $F(\omega)$. Membership of a
configuration in the local class of $F(\omega)$ is decided by the trace of the nest in force on
$F(\omega)$, by clause (c) of the lemma of this section on decidability from traces. Under the
hypothesis of the definition of locality of radius $r$ we have $F(\omega)\subset W\subset W^{+rM}$
for every $\omega$ of the set in (a), so the traces of the two nests on $F(\omega)$ agree and the
local class of $F(\omega)$ is the same for $\mathbf{\Omega}$ and for $\mathbf{\Omega}'$;
requirement (b) thus compares two functions on one domain.

We record the following consequence of locality of radius $r$; its proof is given after the
definition. Suppose the atom rule is local of radius $r$. Let $K$ be an operator of the list of
operators of the step, let $W$ be a region, and let $\mathbf{\Omega},\mathfrak{P}$ and
$\mathbf{\Omega}',\mathfrak{P}'$ be as above. Then
\begin{equation}\label{4.6:eq-atom-locality-a}
\begin{aligned}
&\operatorname{tr}_{W^{+rM}}(\mathbf{\Omega};\mathfrak{P})
=\operatorname{tr}_{W^{+rM}}(\mathbf{\Omega}';\mathfrak{P}')\\
&\qquad\Longrightarrow\quad
K^{W''}[\mathbf{\Omega}]=K^{W''}[\mathbf{\Omega}']
\ \text{term by term, for every } W''\subset W .
\end{aligned}
\end{equation}
\end{definition}

\begin{proof}[Proof of \eqref{4.6:eq-atom-locality-a}]
Assume that the traces of $(\mathbf{\Omega},\mathfrak{P})$ and $(\mathbf{\Omega}',\mathfrak{P}')$ on
$W^{+rM}$ are equal, and let $W''\subset W$. The pair of nests then satisfies the hypothesis of
the definition of locality of radius $r$, with this $W$. Requirements (a) and (b) are therefore
available.

The index set of the series defining $K^{W''}[\mathbf{\Omega}]$ is
\begin{equation*}
\{\omega\in\mathcal{W}_K(\mathbf{\Omega}) : F(\omega)\subset W''\}
=\{\omega\in\mathcal{W}_K(\mathbf{\Omega}) : F(\omega)\subset W\}\cap\{\omega : F(\omega)\subset W''\},
\end{equation*}
because $W''\subset W$. The same holds with $\mathbf{\Omega}'$ in place of $\mathbf{\Omega}$.

The footprint $F(\omega)$ is label-free: it is determined by the word $\omega$ alone, and not by
the nest in force, since the alphabet of the words and the footprint map $F$ of
Definition~\ref{4.6:def-letters-words} do not depend on the nest. The second set on the right is
therefore the same for $\mathbf{\Omega}$ and for $\mathbf{\Omega}'$. By (a), the first set on the
right is also the same for both nests. Hence the series defining $K^{W''}[\mathbf{\Omega}]$ and
$K^{W''}[\mathbf{\Omega}']$ run over one and the same index set. In other words, a word with
$F(\omega)\subset W''\subset W$ is admissible for $\mathbf{\Omega}$ if and only if it is
admissible for $\mathbf{\Omega}'$.

Let $\omega$ be a word of this common index set. Since $F(\omega)\subset W$, requirement (b) applies to it and gives $K_\omega[\mathbf{\Omega}]=K_\omega[\mathbf{\Omega}']$ as functions on the local class of $F(\omega)$. The values of the two terms indexed by $\omega$ therefore agree.

Both series converge absolutely, by the convergence requirement recalled at the beginning of the definition. They have the same index set and equal terms. Hence $K^{W''}[\mathbf{\Omega}]=K^{W''}[\mathbf{\Omega}']$ term by term, which is \eqref{4.6:eq-atom-locality-a}.
\end{proof}

\begin{lemma}[Membership and level tests at radius zero]\label{4.6:lem-membership-tests}
Let $\mathbf{\Omega}=(\Omega_0\supseteq\cdots\supseteq\Omega_k)$ be a $\pi$-grained nest, that is,
a nest in the sense of Definition~\ref{4.6:def-nest-traces}, each $\Omega_i$ a closed union of
$\pi_i$-blocks, carrying a region datum $\mathfrak{P}$, and let $\Lambda_j$ be its layers as in
that definition. We say that a test on a point, a bond or a cube is \emph{decided at radius
zero} if it is determined by the membership bits and the region datum on the blocks containing
that point or bond, or on the blocks of which that cube is the union, and by nothing else.
\begin{itemize}
\item[(a)] For every point or bond $z$, the level datum of $z$
(Definition~\ref{4.6:def-local-sequence}) decides the level $n(z)$ of $z$, that is, the index
with $z\in\Lambda_{n(z)}$, and for every level $n$ decides whether $z\in\Lambda_n$. Hence it
decides the level at which the global operators $Q'^{*}\mathfrak{a}'Q'$ and $Q^{*}\mathfrak{a}Q$
built from the averaging operators $Q'$ and $Q$, the constraint
\cite[(3.110)]{B-CMP99b} and the vertex $C^{(2)}_j$ of \cite[(3.136)]{B-CMP99b} act at $z$. The
memberships $z\in\Lambda'$, when $z$ is a point, and $z\subset\Lambda$, when $z$ is a bond, in
the regions $\Lambda'$ and $\Lambda$ of \cite[\S3]{B-CMP99b} are decided by the region datum on
the block of $z$. All
these tests are decided at radius zero.
\item[(b)] For a candidate cube $Q\in\mathcal{C}_g$, the membership $Q\in\mathcal{D}_g(\mathbf{\Omega})$
is decided by the membership bits of the $2^d$ blocks of $\pi_g$ of which $Q$ is the union.
\item[(c)] Let $\widehat{C}$ be a cube of the class of \cite[(3.35)]{B-CMP99b}, and let $j$,
$0\le j\le k$, be the index attached to it in \cite[p.~396]{B-CMP99b}, which in terms of the
nest is the index with
\begin{equation}\label{4.6:eq-membership-tests-1}
\widehat{C}\subset\Omega_j,\qquad \widehat{C}\cap\Omega_{j+2}=\emptyset,\qquad
\widehat{C}\cap\Omega_j\not\subset\Omega_{j+1}.
\end{equation}
Then $j$ is decided by the trace
$\operatorname{tr}_{\widehat{C}}(\mathbf{\Omega};\mathfrak{P})$. Consequently the local class
of the footprint $[\omega]$ of a word $\omega$, in the sense of the definition of local classes,
is the same set for two nests whose traces on $[\omega]$ coincide.
\end{itemize}
\end{lemma}

\begin{proof}
(a) By Definition~\ref{4.6:def-local-sequence}, the level datum of $z$ consists of membership
bits of the blocks containing $z$ in the sets $\Omega_i$. Since the layer $\Lambda_j$ is formed
from $\Omega_j$ and $\Omega_{j+1}$ alone (Definition~\ref{4.6:def-nest-traces}),
whether $z\in\Lambda_n$ is read off from these bits, and so is the unique $n=n(z)$ with
$z\in\Lambda_n$. \cite[(2.20)]{B-CMP96} states $(QA)(b)=(Q_jA)(b)$ for $b\in\Lambda_j$. Hence
the component of $Q$ that acts at $z$ is $Q_{n(z)}$. The same holds for $Q'$, and the
diagonal weight matrices $\mathfrak{a}'$, $\mathfrak{a}$ carry at $z$ the entry of index $n(z)$;
therefore the level at which $Q'^{*}\mathfrak{a}'Q'$ and $Q^{*}\mathfrak{a}Q$ act at $z$ is
$n(z)$. The constraint \cite[(3.110)]{B-CMP99b}, which we read as
$L^{j}\eta\,Q_j(U)A=B$ on $\Lambda_j$ (symbols as in that display), is imposed on
$\Lambda_j$ level by level, so at $z$ only its component with $j=n(z)$ acts. Likewise, the
vertex $C^{(2)}_j$ of
\cite[(3.136)]{B-CMP99b} acts at $z$ for $j=n(z)$. The memberships $z\in\Lambda'$ (for a point
$z$) and $z\subset\Lambda$ (for a bond $z$) in the regions $\Lambda'$ and $\Lambda$ of
\cite[\S3]{B-CMP99b} are statements about the region on the block of $z$, and
so they are read from the region datum $\mathfrak{P}$ on that block. None of these tests uses
data outside $z$ and its blocks, so each is decided at radius zero.

(b) By Definition~\ref{4.6:def-local-sequence},
\begin{equation*}
\mathcal{D}_g(\mathbf{\Omega})=\bigl\{Q\in\mathcal{C}_g:\ \exists\,e\in\pi_g,\ e\subset Q,\
[e\subset\Omega_g]=1,\ [e\subset\Omega_{g+1}]=0\bigr\}.
\end{equation*}
A cube $Q\in\mathcal{C}_g$ is the union of $2^d$ blocks of $\pi_g$
(Definition~\ref{4.6:def-candidate-cubes}), and these are the only $e\in\pi_g$ with
$e\subset Q$. Since $\Omega_g$ and $\Omega_{g+1}$ are unions of blocks of $\pi_g$ and of
$\pi_{g+1}$ respectively, and $g+1\ge g$, the first clause of \eqref{4.6:eq-candidate-cubes-a},
applied with $i=g$ and with $i=g+1$, gives
\begin{equation*}
[e\subset\Omega_g]=[c\subset\Omega_g],\qquad [e\subset\Omega_{g+1}]=[c\subset\Omega_{g+1}]
\end{equation*}
for every $\pi_0$-block $c\subset e$. Thus the membership $Q\in\mathcal{D}_g(\mathbf{\Omega})$ is
decided by
the bits of the $2^d$ blocks of $Q$ itself.

(c) In the reading of \cite[p.~396]{B-CMP99b} adopted here, for each cube $\square$ of the
class of \cite[(3.35)]{B-CMP99b} there is a unique index $j$, $0\le j\le k$, such that
\begin{itemize}
\item $\square\subset B^j(\Lambda_j)\cup B^{j+1}(\Lambda_{j+1})$;
\item $\square\cap B^j(\Lambda_j)\neq\emptyset$;
\item $\square$ is a union of several big blocks of the lattice $T_{L^{-j}}$ (notation of
\cite[p.~396]{B-CMP99b}).
\end{itemize}
The set $B^j(\Lambda_j)$ of this reading is the set $\Omega_j\setminus\Omega_{j+1}$, in the form
$\Omega_j\setminus\Omega_{j+1}=B_j(\Lambda_j)$ recalled in
Definition~\ref{4.6:def-nest-traces}, with
$\Omega_{k+1}=\Omega_{k+2}=\emptyset$. In terms of the nest, the first two of these conditions
on $\widehat{C}$ read as the three conditions of \eqref{4.6:eq-membership-tests-1}; the third
concerns the shape of $\widehat{C}$ alone.
Each condition in \eqref{4.6:eq-membership-tests-1} is a statement about the sets
$\widehat{C}\cap\Omega_i$, that is, about the bits of the blocks of $\widehat{C}$, and these
are recorded by the trace $\operatorname{tr}_{\widehat{C}}(\mathbf{\Omega};\mathfrak{P})$ of
Definition~\ref{4.6:def-nest-traces}. Since $j$ is unique, it is therefore decided by this
trace. By what was just shown, the index of each class cube is decided by the trace on that
cube; hence the local class, in the sense of the definition of local classes, of the footprint
$[\omega]$ of a word $\omega$ is the same set for
two nests with equal traces on $[\omega]$.
\end{proof}

\begin{lemma}[The local sequence reads only the triple enlargement]\label{4.6:lem-sequence-reach}
Let $g$ be a scale, let $Q\in\mathcal{C}_g$, and let $\mathbf{\Omega}$, $\mathbf{\Omega}'$ be nests. The local
sequences $\mathbf{\Omega}(Q)$, $\mathbf{\Omega}'(Q)$ are those of
Definition~\ref{4.6:def-local-sequence}, given by \eqref{4.6:eq-local-sequence-a}. Call a map
\emph{label-free} if its value depends only on its displayed arguments, and not on the nest from
which these arguments are read.
\begin{itemize}
\item[(a)] There is a label-free map $\Sigma(Q,g;\cdot)$ such that
\begin{equation}\label{4.6:eq-sequence-reach-1}
\mathbf{\Omega}(Q)=\Sigma\Bigl(Q,g;\bigl([e\subset\Omega_{g+1}],\,[e\subset\Omega_{g+2}]\bigr)_{e\in\pi_g,\;
e\subset Q^{(3)}}\Bigr).
\end{equation}
\item[(b)] Suppose that the traces on $Q^{(3)}$ of $\mathbf{\Omega}$ and of $\mathbf{\Omega}'$ at the levels
$g+1$ and $g+2$ coincide; this holds in particular if
$\operatorname{tr}_{Q^{(5)}}(\mathbf{\Omega})=\operatorname{tr}_{Q^{(5)}}(\mathbf{\Omega}')$. Then
$\mathbf{\Omega}(Q)=\mathbf{\Omega}'(Q)$ as tuples of sets, member by member. The derived objects
coincide as well: the layers $\Lambda_n(Q)$, the site set $\mathfrak{B}_Q$, the type map, the blocks,
$\chi_Q$, and the weight matrices $\mathfrak{a}'$, $\mathfrak{a}$ of
Definitions~\ref{4.6:def-local-sequence} and~\ref{4.6:def-local-generators}.
\item[(c)] Every member of $\mathbf{\Omega}(Q)$ is contained in $Q^{(5)}$, as in~\cite{B-CMP99b}.
\item[(d)] Let $\widetilde{\mathbf{\Omega}}$ be the wavy nest formed from $\mathbf{\Omega}$ with a set $P$ at
scale $j$, as in the definition of the wavy nest; here $P\subset\mathbb{T}$ is the set entering that
definition at scale $j$, and the letter does not denote the period of the torus. Write
$\widetilde{T}(Q)$ for the seam datum of the local sequence $\widetilde{\mathbf{\Omega}}(Q)$, that is,
of the local sequence \eqref{4.6:eq-local-sequence-a} read from $\widetilde{\mathbf{\Omega}}$. For a cube
$Q$ of scale $j$,
\begin{equation}\label{4.6:eq-sequence-reach-2}
\widetilde{T}(Q)=\bigcup\bigl\{e : e\in\pi_j,\ e\subset Q^{(3)},\ [e\subset P]=1\bigr\}.
\end{equation}
For a cube $Q$ of scale $j-1$,
\begin{equation*}
\widetilde{T}(Q)=\bigcup\bigl\{e : e\in\pi_{j-1},\ e\subset Q^{(3)},\
[e\subset\Omega_j]=1,\ [e\subset P]=0\bigr\},
\end{equation*}
where $\Omega_j$ is the level-$j$ member of $\mathbf{\Omega}$, not of $\widetilde{\mathbf{\Omega}}$. For a
cube $Q$ of scale at most $j-2$, the datum
$\widetilde{T}(Q)$ coincides with the datum $T(Q)$ of $\mathbf{\Omega}$.
\item[(e)] For $Q\in\mathcal{C}_g$ as above, the outer seam, that is, whether
$Q^{(4)}\not\subset\Omega_g$, is read by none of the objects in \textup{(a)} and \textup{(b)}.
\end{itemize}
\end{lemma}

\begin{proof}
We use the formulas of Definition~\ref{4.6:def-local-sequence}.

\emph{Clause (a).} By \eqref{4.6:eq-local-sequence-a}, the member $\Omega_{g+1}(Q)$ of the local
sequence is the seam datum $T(Q)$. The displayed formula for $T(Q)$ expresses it as a function of
the two bit rows $\bigl([e\subset\Omega_{g+1}]\bigr)$ and $\bigl([e\subset\Omega_{g+2}]\bigr)$, indexed
by the $\pi_g$-blocks $e\subset Q^{(3)}$. Apart from these two bit rows, the formula involves only
$Q$ and $g$; in particular it does not depend on which nest the bits are read from. Every other
member of $\mathbf{\Omega}(Q)$ is a concentric cube $Q^{(t)}$, determined
by $Q$ and $g$ alone, and hence label-free. Let $\Sigma(Q,g;\cdot)$ be the map that assigns to the
two bit rows the tuple whose member at level $g+1$ is the value of $T(Q)$ and whose other members
are these concentric cubes. Then $\Sigma(Q,g;\cdot)$ is label-free, and
\eqref{4.6:eq-sequence-reach-1} holds.

\emph{Clause (b).} By hypothesis, the traces on $Q^{(3)}$ at the levels $g+1$ and $g+2$ coincide.
Hence, for every $\pi_g$-block $e\subset Q^{(3)}$,
\begin{equation*}
[e\subset\Omega_{g+1}]=[e\subset\Omega'_{g+1}],\qquad [e\subset\Omega_{g+2}]=[e\subset\Omega'_{g+2}].
\end{equation*}
Since $Q^{(3)}\subset Q^{(5)}$, the same conclusion holds under the stronger hypothesis
$\operatorname{tr}_{Q^{(5)}}(\mathbf{\Omega})=\operatorname{tr}_{Q^{(5)}}(\mathbf{\Omega}')$. Apply
\eqref{4.6:eq-sequence-reach-1} to $\mathbf{\Omega}$ and to $\mathbf{\Omega}'$. The same label-free
map $\Sigma(Q,g;\cdot)$ is evaluated at the same arguments, so $\mathbf{\Omega}(Q)=\mathbf{\Omega}'(Q)$
member by member, as subsets. By their definitions, the objects $\Lambda_n(Q)$, $\mathfrak{B}_Q$, the
type map, the blocks, $\chi_Q$, $\mathfrak{a}'$ and $\mathfrak{a}$ are label-free functions of the
tuple $\mathbf{\Omega}(Q)$. For $\mathfrak{a}'$ and $\mathfrak{a}$ this uses one further fact: the
coefficients $a_n$ entering them are functions of $n$, of the parameter $a$ of the operators
$\Delta_{a,Q}$, $\Delta'_{a,Q}$, and of $L$ alone~\cite{B-CMP99b}. Equal tuples
therefore give equal derived objects.

\emph{Clause (c).} \cite{B-CMP99b} states this containment for the sequence given by the
same formulas.

\emph{Clause (d).} Apply \eqref{4.6:eq-sequence-reach-1} to the nest $\widetilde{\mathbf{\Omega}}$ in
place of $\mathbf{\Omega}$, with the bits of $\widetilde{\mathbf{\Omega}}$ on the blocks
$e\subset Q^{(3)}$ as listed in the statement of \textup{(d)} for the three ranges of scales.

\emph{Clause (e).} By \eqref{4.6:eq-sequence-reach-1}, the arguments of $\Sigma(Q,g;\cdot)$ are bits at
the levels $g+1$ and $g+2$ only. By \textup{(b)}, every derived object is a function of
$\mathbf{\Omega}(Q)$. Hence no bit at level $g$ is read, and in particular the condition
$Q^{(4)}\not\subset\Omega_g$ is not read.
\end{proof}

\begin{lemma}[Inverses and constrained-inversion blocks of equal matrices]\label{4.6:lem-matrix-algebra}
Let $V$ and $E$ be finite-dimensional inner-product spaces, and write ${}^*$ for the adjoint.
\begin{itemize}
\item[(a)] Let $(\Delta';Q';\mathfrak{a}';\chi)$ be a tuple of linear maps, each acting between two of
the spaces $V$, $E$, and suppose a second tuple on the same spaces is equal to it. Then every
polynomial or rational word in the entries of the tuple takes the same value for both tuples. In
particular this holds for
\begin{gather*}
G':=(\Delta')^{-1},\qquad T':=Q'G'^{\,2}Q'^{*},\qquad C:=T'^{-1},\\
P:=G'Q'^{*}CQ'G',\qquad R:=1-P,
\end{gather*}
and for $(x+\Delta')^{-1}$ at every value of the parameter $x$. The two values have equal domains of
definition, the invertibility loci being the zero-set complements of the same determinants. The
same holds for analytic continuations in a parameter.
\item[(b)] Let $\mathcal{Q}$ be a symmetric ($\mathcal{Q}^*=\mathcal{Q}$) invertible map on $V$ with
$G:=\mathcal{Q}^{-1}$. Let $W\colon V\to E$ be linear with $T:=WGW^{*}$ invertible, and put
$H:=GW^{*}T^{-1}$. Then the block map on $V\oplus E$ in the display below is invertible, with
\begin{equation}\label{4.6:eq-matrix-algebra-1}
\begin{pmatrix}\mathcal{Q} & W^{*}\\ W & 0\end{pmatrix}^{-1}
=\begin{pmatrix}G-HTH^{*} & H\\ H^{*} & -T^{-1}\end{pmatrix}.
\end{equation}
Consequently every block of this constrained inversion, and every Schur complement entering it, is a
rational function of the entries of $\mathcal{Q}$ and $W$ alone.
\end{itemize}
\end{lemma}

\begin{proof}
(a) Each word is built from the entries of the tuple in finitely many steps. Every step is a sum, a
product, an adjoint, the addition of a scalar multiple of the identity, or the inversion of a square
matrix $M$ formed at an earlier step. By Cramer's rule, $M$ is invertible exactly when
$\det M\neq0$, and then $M^{-1}=(\det M)^{-1}\operatorname{adj}M$. Here the entries of the adjugate
$\operatorname{adj}M$ and the number $\det M$ are polynomials in the entries of $M$.

We argue by induction on the number of steps. Suppose the matrices formed at earlier steps coincide
for the two tuples. Then the determinants of the matrices to be inverted coincide, so the two
inversions are defined at the same points, and where they are defined their values coincide.

This applies in turn to $G'$, then to $T'$ and $C=T'^{-1}$, and then to $P$ and $R$. For
$(x+\Delta')^{-1}$, the determinant $\det(x+\Delta')$ is one and the same polynomial in $x$ for both
tuples, so the set of parameters $x$ at which the inverse exists is the same.

Now suppose the two tuples depend analytically on a parameter and coincide as functions of it. Then
any word formed from them is one and the same function on one and the same domain. Hence its
analytic continuations in the parameter coincide.

(b) Since $\mathcal{Q}^*=\mathcal{Q}$, also $G^*=G$, and therefore
\begin{equation*}
T^*=WG^*W^*=T,\qquad H^{*}=(T^{-1})^{*}WG^{*}=T^{-1}WG .
\end{equation*}
We multiply the block map of \eqref{4.6:eq-matrix-algebra-1} from the right by the proposed inverse
and compute the four blocks of the product.

Upper left block. Since $\mathcal{Q}G=I$ and $\mathcal{Q}H=W^{*}T^{-1}$, we have
$\mathcal{Q}HTH^{*}=W^{*}H^{*}=W^{*}T^{-1}WG$. Hence
\begin{equation*}
\mathcal{Q}(G-HTH^{*})+W^{*}H^{*}=I-W^{*}T^{-1}WG+W^{*}T^{-1}WG=I .
\end{equation*}

Upper right block. We have
\begin{equation*}
\mathcal{Q}H+W^{*}(-T^{-1})=W^{*}T^{-1}-W^{*}T^{-1}=0 .
\end{equation*}

Lower left block. We have $WHTH^{*}=(WGW^{*})T^{-1}WG$, so
\begin{equation*}
W(G-HTH^{*})+0\cdot H^{*}=WG-(WGW^{*})T^{-1}WG=WG-WG=0 ,
\end{equation*}
because $WGW^{*}=T$.

Lower right block. We have
\begin{equation*}
WH+0\cdot(-T^{-1})=WGW^{*}T^{-1}=I .
\end{equation*}

The product is therefore the identity of $V\oplus E$. A linear map of a finite-dimensional space that
has a right inverse is invertible, and the right inverse is its inverse. This proves
\eqref{4.6:eq-matrix-algebra-1}.

Each block on the right of \eqref{4.6:eq-matrix-algebra-1} is obtained from $\mathcal{Q}$ and $W$ by
sums, products, adjoints and the inversions $G=\mathcal{Q}^{-1}$ and $T^{-1}$. By the argument of
part (a), each block is therefore a rational function of the entries of $\mathcal{Q}$ and $W$ alone,
and equal pairs $(\mathcal{Q},W)$ give equal blocks. The block inversions in the Lagrange-multiplier
computations of \cite{B-CMP99b} and of \cite[(2.18)--(2.23)]{B-CMP96} are of this form.
\end{proof}

\begin{proposition}[Equal local sequences give equal generators]\label{4.6:prop-generator-equality}
Let $(\mathbf{\Omega},\mathfrak{P})$ and $(\mathbf{\Omega}',\mathfrak{P}')$ be decreasing
$\pi$-grained nests of levels at most $k$ with region data, in the setting of
Definition~\ref{4.6:def-local-sequence}. Augmented nests are allowed; see the last paragraph of
the statement. Let $Q$ be a block-aligned cube of scale $g$ to which the construction of
Definition~\ref{4.6:def-local-sequence} applies. Assume that
\begin{equation}\label{4.6:eq-generator-equality-1}
[e\subset\Omega_i]=[e\subset\Omega'_i]
\end{equation}
for $i\in\{g,g+1,g+2\}$ and every block $e\in\pi_g$ with $e\subset Q^{(3)}$.
For the cubes of an augmented nest, assume in addition that the region bits on these blocks
agree. These hypotheses hold in particular if
\begin{equation*}
\operatorname{tr}_{Q^{(5)}}(\mathbf{\Omega};\mathfrak{P})
=\operatorname{tr}_{Q^{(5)}}(\mathbf{\Omega}';\mathfrak{P}').
\end{equation*}
Write $U|_{\Omega_0(Q)}$ for the restriction of a configuration $U$ to the bonds of $\Omega_0(Q)$.
Then the following hold.
\begin{itemize}
\item[(i)] $\mathbf{\Omega}(Q)=\mathbf{\Omega}'(Q)$; if moreover $Q\in\mathcal{C}_g$, then
$Q\in\mathcal{D}_g(\mathbf{\Omega})$ if and only if $Q\in\mathcal{D}_g(\mathbf{\Omega}')$.
The following objects coincide for the two
nests: the site sets $\mathfrak{B}_Q$ and $\mathfrak{B}_Q^{v}$, the types, the blocks, the
weights, the set $\Omega_0(Q)$ and the spaces. This includes the entries attached to the inner
seams of both scales $g$ and $g+1$.
\item[(ii)] Each of the operators
\begin{equation*}
\begin{gathered}
Q'_Q,\quad Q_Q,\quad \chi_Q\Delta_U\chi_Q,\quad \chi_Q\Delta(U)\chi_Q,\quad \chi_Q D\chi_Q,
\quad \chi_Q D^{*}\chi_Q,\\
\Delta'_{a,Q},\quad T'_Q,\quad \Delta_{a,Q},\quad
\mathfrak{Q}_Q,\quad \mathfrak{Q}'_Q
\end{gathered}
\end{equation*}
is the same matrix-valued function of $U|_{\Omega_0(Q)}$ for $\mathbf{\Omega}$ and for
$\mathbf{\Omega}'$. For $\Delta_{a,Q}$ this includes its term $DR_QD^{*}$.
Consequently the following are the same partial functions of $U|_{\Omega_0(Q)}$, with the same
domain of definition, for the two nests:
\begin{itemize}
\item $G'_Q$, $C_Q$, $P_Q$, $R_Q$, $G_Q$, $T_Q^{-1}$, $H_Q$ and $S_Q$;
\item $\mathfrak{Q}_Q^{-1}$, block by block;
\item the resolvents of all these operators in a parameter $x$;
\item $P_Q(\partial h)$, for every fixed multiplier $h$;
\item every local tower operator $K_Q$ of Definition~\ref{4.6:def-local-generators}.
\end{itemize}
Here $\Omega_0(Q)\subset Q^{(5)}$.
\item[(iii)] Under a gauge transformation $u$ on $\Omega_0(Q)$ the objects of \textup{(ii)}
transform, for both nests, by conjugation with the same $R(u)$, in the sense of
\cite[(3.31)--(3.34)]{B-CMP99b}. The local class is common to the two nests, by
Lemma~\ref{4.6:lem-membership-tests}\textup{(c)}.
\item[(iv)] Parts \textup{(i)}--\textup{(iii)} use no positivity, no admissibility condition and no
bound. The existence of the generators on the regularity class of
\cite[(3.35)--(3.38)]{B-CMP99b} is therefore one and
the same statement for both nests. That statement is \cite[Corollary~3.6]{B-CMP99b}, taken
together with the part of \cite[Theorem~3.11]{B-CMP99b} that concerns these generators.
\end{itemize}
Augmented nests are treated as follows. Let
$\widetilde{\mathbf{\Omega}}=(\Omega_0,\dots,\Omega_j,\widetilde{\Omega}_{j+1})$ be an augmented
nest, whose last member $\widetilde{\Omega}_{j+1}$ is the set adjoined at level $j+1$. Then the
proposition applies to $\widetilde{\mathbf{\Omega}}$
verbatim. For a cube $Q$ of scale $j$, the member of level $j+1$ of its local sequence is
\begin{equation*}
\Omega_{j+1}(Q)=\bigcup\bigl\{e\in\pi_j : e\subset Q^{(3)},\ e\subset\widetilde{\Omega}_{j+1}\bigr\},
\end{equation*}
that is, the intersection of $\widetilde{\Omega}_{j+1}$ with $Q^{(3)}$ taken block by block.
\end{proposition}

\begin{proof}
\emph{Part \textup{(i)}.}
By hypothesis \eqref{4.6:eq-generator-equality-1}, the two nests have the same membership bits
$[e\subset\Omega_g]$ and $[e\subset\Omega_{g+1}]$ on the $\pi_g$-blocks $e\subset Q$. We apply
Lemma~\ref{4.6:lem-membership-tests}\textup{(b)}, which decides membership of $Q$ in
$\mathcal{D}_g$ from these bits. It gives $Q\in\mathcal{D}_g(\mathbf{\Omega})$ if and only if
$Q\in\mathcal{D}_g(\mathbf{\Omega}')$.

Next we use Lemma~\ref{4.6:lem-sequence-reach}. It states that the local sequence
$\mathbf{\Omega}(Q)$ of the display \eqref{4.6:eq-local-sequence-a} is the value of a map
$\Sigma(Q,g;\cdot)$ that does not depend on the nest. The arguments of this map are the bits
$[e\subset\Omega_{g+1}]$ and $[e\subset\Omega_{g+2}]$ for the $\pi_g$-blocks $e\subset Q^{(3)}$.
For an augmented nest, the region bits on these blocks are also arguments. By
\eqref{4.6:eq-generator-equality-1} and the added hypothesis on region bits, the arguments agree
for $\mathbf{\Omega}$ and $\mathbf{\Omega}'$. Hence $\mathbf{\Omega}(Q)=\mathbf{\Omega}'(Q)$ as
tuples of sets.

The same lemma gives the remaining coincidences. The derived objects $\Lambda_n(Q)$,
$\mathfrak{B}_Q$, the type map, the blocks, $\chi_Q$, $\mathfrak{a}'$ and $\mathfrak{a}$ are
functions of $\mathbf{\Omega}(Q)$ alone, so they coincide. All members of $\mathbf{\Omega}(Q)$ lie
in $Q^{(5)}$. The site set $\mathfrak{B}_Q^{v}$, the weights, $\Omega_0(Q)$ and the spaces are
read off from the same tuple, so they coincide as well. This applies to the entries of
$\mathbf{\Omega}(Q)$ attached to the inner seams of scale $g$ and of scale $g+1$ alike, because
both are entries of $\mathbf{\Omega}(Q)$.

The trace condition implies the hypotheses. Indeed, equality of the traces on $Q^{(5)}$ contains
equality of the bits of levels $g$, $g+1$, $g+2$ and of the region bits on all blocks contained in
$Q^{(5)}$. Since $Q^{(3)}\subset Q^{(5)}$, this includes the blocks contained in $Q^{(3)}$.

\emph{Part \textup{(ii)}: the places where the nest enters the definitions.}
By Definition~\ref{4.6:def-local-generators}, the localized generators of $Q$ are the operators of
\cite{B-CMP99b} with the nest $\{\Omega_j\}_{j\le k}$ replaced by the local sequence
$\mathbf{\Omega}(Q)$. We go through the definitions of these operators in \cite{B-CMP99b} and
list every place where the nest enters. There are three kinds of entry.

\smallskip
\noindent($\alpha$) \emph{Index sets and levels of summands.} The nest enters here as follows.
\begin{itemize}
\item The quadratic form of \cite[(3.16)]{B-CMP99b} carries the sums
$\sum_{j=0}^{k}a_j\sum_{b\in\Lambda_j}$.
\item The definition \cite[(3.18)]{B-CMP99b} is made for $y\in\Lambda_j$.
\item The operator $Q'$ is defined on $\mathfrak{B}=\bigcup_j\Lambda_j$.
\item The null space in \cite[(3.21)]{B-CMP99b} is
$N(Q')=\{\lambda : Q'\lambda=0\}$.
\item The nest also enters through the index sets of \cite[(3.24)]{B-CMP99b}.
\end{itemize}
The definition \cite[(2.20)]{B-CMP96} has the same structure: the global averaging acts on
$\Lambda_j$ as the $j$-fold one.

\smallskip
\noindent($\beta$) \emph{The Dirichlet domain.} The nest enters through the domain $\Omega_0$. This
is the domain of the Hilbert space $\ell^2(\Omega_0)$ on which the operators act, with Dirichlet
conditions on its boundary. The choice of $\Omega_0$ made in \cite{B-CMP99b} is an entry of this
kind.

\smallskip
\noindent($\gamma$) \emph{Derived operators.} The nest enters the operators $R$, $P$ and $G'$ of
\cite[(3.25)]{B-CMP99b}, whose localized versions are $R_Q$, $P_Q$ and $G'_Q$, but only through
($\alpha$) and ($\beta$); see also \cite[(2.17), (2.26)--(2.27)]{B-CMP96}.

\smallskip
We now show that the nest enters nowhere else. We take the ingredients in turn.

\emph{The averaging operators.} In \cite{B-CMP99b} the averaging operators are compositions of
$j$ one-step averaging operators. The linearised one-step operator $Q(V)$ is given by the
explicit formula \cite[(124)]{B-CMP98}. Every summand of that formula is indexed by a point
$x\in B(c_-)$ of one block. In the notation of \cite[(124)]{B-CMP98}, $V$ stands for the
background configuration, written $V_0$ in that display, and $(\overline{V}_0)_c$ is its one-step
average on a bond $c$ of the coarser lattice, with endpoints $c_-$ and $c_+$. The letter $A$
denotes the bond field on which the linearised operator $Q(V_0)$ acts, and not the Wilson action
of Definition~\ref{1.1:def-wilson-action}. Finally, $[x,x']$ and $\Gamma_{c_\pm,\cdot}$ are the
bonds and contours of that display. The summands read the following quantities:
\begin{itemize}
\item the values $(R_{0,c_-}A)([x,x'])$, $(R_{0,c_-}A)(\Gamma_{c_-,x})$,
$(R_{0,c_+}A)(\Gamma_{c_+,x'})$ and $(R_{0,c_-}A)(c)$, where the operation $R_{0,\cdot}$ of
\cite[(58)--(60)]{B-CMP98} is taken along contours inside the block;
\item the matrix
\begin{equation*}
\overline{R}_{0,c}=R\bigl((\overline{V}_0)_c\bigr)
\end{equation*}
of \cite[(59)]{B-CMP98};
\item the functions $g(\pm i\operatorname{ad}_Y)$, $g^{-1}(\pm i\operatorname{ad}_{Y_x})$ and
$e^{\pm i\operatorname{ad}_Y}$ of the auxiliary letters $Y$, $Y_x$, where this $g(\cdot)$ is a
fixed function of \cite[(124)]{B-CMP98} unrelated to the scale $g$.
\end{itemize}
The letters $Y$ and $Y_x$ are defined in \cite{B-CMP98} before \cite[(124)]{B-CMP98}. We read
them as objects of the background attached to the block $B(c_-)$. Two features of
\cite[(124)]{B-CMP98} support this reading. First, $Y_x$ carries the index $x\in B(c_-)$.
Second, \cite{B-CMP98} states, of the terms of \cite[(124)]{B-CMP98}, that the operators
occurring in these terms can be estimated
by $O(L^2\alpha_0)$, where $\alpha_0$ is the constant of \cite[(52)]{B-CMP98}.

\emph{The linear parts $Q'_j$.} The operators $Q'_j(U)$ of \cite{B-CMP99b} are the linear parts of
the averaging operations defined by \cite[(78)--(80)]{B-CMP98}, together with
\cite[(81)]{B-CMP98}. Each of these displays sums over $x\in B(y)$ for one block, and composes
through block centres.

The operator $Q'_j$ is displayed in \cite[(3.19)]{B-CMP99b}; its contours are those of
\cite[(58), (60), (77)]{B-CMP98}. They are:
\begin{itemize}
\item the contours $\Gamma_{y,x}$ with $x\in B(y)$;
\item the chains $\Gamma^{(j+1)}_{x_{j+1},x}$ with $x\in B^{j+1}(x_{j+1})$, which run through the
block centres $x_{j+1},x_j,\dots,x_1,x_0=x$.
\end{itemize}
In the chains, each centre $x_i$ lies in the block $B(x_{i+1})$ of the next one. For a chain
descending from a centre $y=x_j$ through $x_{j-1},\dots,x_1$ to $x=x_0$, the operators built
along it read the gauge field variables only on bonds $b\subset B^{j}(y)$.

\emph{The $j$-fold averages.} These are given by \cite[(89)--(92)]{B-CMP98}.
\cite[p.~31]{B-CMP98} states that $(\overline{\overline{U}}{}^k)_c$, for $c\subset\Omega^{(k)}$,
depends on the variables $U_b$ for
\begin{equation*}
b\subset B^k(c_-)\cup B^k(c_+).
\end{equation*}
\cite[p.~26]{B-CMP98} states that, for $p=\langle x,y,z,w\rangle$, it is enough to assume
\cite[(52)]{B-CMP98} for
\begin{equation*}
p\subset B^k(x)\cup B^k(y)\cup B^k(z)\cup B^k(w).
\end{equation*}
Finally, \cite{B-CMP99b} states that the global $Q(U)$ coincides with $Q_j(U)$ on
$B^j(\Lambda_j)$.

\emph{Conclusion for the averaging operators.} All these objects are formulas in $U$ and in the
fixed hierarchy of blocks. In \cite{B-CMP99b} they are written before the sequence of domains is
introduced. In \cite{B-CMP98} they are written on a decreasing chain of lattices
$\Omega,\Omega',\dots,\Omega^{(k)}$ with no sequence of domains at all. By
their form they are block-local and do not depend on the domains. This is the content of
Lemma~\ref{4.6:lem-averaging-locality}, which we use in the form stated there.

\emph{The remaining ingredients.}
\begin{itemize}
\item The weights $a_j$ obey a recursion in $j$ alone \cite{B-CMP99b}.
\item The operator $\Delta_U$ of \cite[(3.23), (3.50)]{B-CMP99b} does not depend on the nest.
\item Neither do $\Delta(U)$ of \cite[(3.10)]{B-CMP99b}, $D$ and $D^{*}$ of
\cite[(3.3)--(3.9)]{B-CMP99b}, or $J$ of \cite[(3.11)]{B-CMP99b}.
\item The regularity class of \cite[(3.35)--(3.38)]{B-CMP99b}, the norms of
\cite[(3.41)]{B-CMP99b}, the distances $d(y,y')$ and the sets written $\Delta(y)$ in
\cite{B-CMP99b} (a notation unrelated to the operator $\Delta(U)$)
enter only hypotheses and bounds. They are local in a cube or at a point, by
Lemma~\ref{4.6:lem-membership-tests}\textup{(c)}, and they enter the value of no operator.
\item The choice of $\Omega_0$ is an entry of kind ($\beta$).
\end{itemize}

\emph{The data tuple.} Consider the tuple consisting of:
\begin{itemize}
\item the set $\Omega_0(Q)$;
\item the site sets $\mathfrak{B}_Q$ and $\mathfrak{B}_Q^{v}$, with their types and blocks;
\item the weights $(a_n)$;
\item for a cube of an augmented nest, the region bits on the $\pi_g$-blocks $e\subset Q^{(3)}$
that are arguments of the map $\Sigma(Q,g;\cdot)$ of Lemma~\ref{4.6:lem-sequence-reach};
\item the restricted matrices of $\Delta_U$, $\Delta(U)$ and $D$;
\item the components of $Q'_n$ and $Q_n$ at the points selected by the index sets.
\end{itemize}
By ($\alpha$)--($\gamma$) and the absence of any other entry, this tuple depends on
$(\mathbf{\Omega},\mathfrak{P})$ only through the pair
\begin{equation*}
\bigl(\mathbf{\Omega}(Q),\ \text{the region bits on the $\pi_g$-blocks } e\subset Q^{(3)}\bigr),
\end{equation*}
where the second entry is present only for a cube of an augmented nest. By part (i),
$\mathbf{\Omega}(Q)=\mathbf{\Omega}'(Q)$; for a cube of an augmented nest the region bits on these
blocks agree by the added hypothesis on region bits. Hence these pairs coincide for
$(\mathbf{\Omega},\mathfrak{P})$ and $(\mathbf{\Omega}',\mathfrak{P}')$. Therefore the data tuples
coincide, and so do the operators of the first list in (ii) that are read directly from them.

\emph{Transport of equality.} All remaining operators are obtained from the data tuple by
algebraic operations. We apply parts (a) and (b) of Lemma~\ref{4.6:lem-matrix-algebra} to the
two equal tuples. These give equality of every inverse, every resolvent, every block of the
constrained-inversion (Schur complement) problems, and every minimiser built from the data. They
also give equality of the domains of definition, because the invertibility loci are the zero sets
of the same determinants. This yields the equality of $G'_Q$, $C_Q$, $P_Q$, $R_Q$, $G_Q$,
$T_Q^{-1}$, $H_Q$, $S_Q$; of the blocks of $\mathfrak{Q}_Q^{-1}$; of their resolvents in $x$; and
of the local tower operators $K_Q$, as partial functions.

\emph{The term $DR_QD^{*}$.} This term is included in the above. By the previous paragraph, $R_Q$
is a word in the data. Then
\begin{equation}\label{4.6:eq-generator-equality-2}
\Delta_{a,Q}=\chi_Q\bigl(\Delta(U)+D(\chi_QR_Q\chi_Q)D^{*}+Q_Q^{*}\mathfrak{a}Q_Q\bigr)\chi_Q
\end{equation}
is again a word in the data. Hence $\Delta_{a,Q}$ is the same function for the two nests, and so
is $G_Q$. That $G_Q$ is built with the localized operator $R_Q$ is part of
Definition~\ref{4.6:def-local-generators}.

\emph{The multiplier term.} For a fixed multiplier $h$,
\begin{equation*}
P_Q(\partial h)=[DP_QD^{*},h]
\end{equation*}
is the commutator of equal matrices with the fixed multiplication operator $h$. Hence it is equal
for the two nests.

\emph{The window in $U$.} It remains to locate the variables of $U$ read by the components of
$Q'_n$, and to see that $\Omega_0(Q)\subset Q^{(5)}$. At a point of $\Lambda_n(Q)$, the components
of $Q'_n$ read variables at most half an averaging block beyond $\Lambda_n(Q)$, as described in
\cite{B-CMP99b}. Under the margins $2R_0M_0$ of \cite{B-CMP99b}, this region lies inside
$Q^{(5)}$. The inclusion $\Omega_0(Q)\subset Q^{(5)}$ follows from part (i), since all members of
$\mathbf{\Omega}(Q)$ lie in $Q^{(5)}$.

\emph{Part \textup{(iii)}.} The gauge transformation laws \cite[(3.31)--(3.32)]{B-CMP99b}, and
\cite[(3.33)--(3.34)]{B-CMP99b}, are identities of the ingredients that do not depend on the nest.
Conjugation by $R(u)$ passes through sums, products, adjoints and inverses. All operators in (ii)
are obtained from the ingredients by these operations together with the data of part (i), which
are equal for the two nests. Hence both families transform by the same $R(u)$ under gauge
transformations on $\Omega_0(Q)$. The local class is common to the two nests by
Lemma~\ref{4.6:lem-membership-tests}\textup{(c)}.

\emph{Part \textup{(iv)}.} The arguments above use only the definitions of the operators and
Lemma~\ref{4.6:lem-matrix-algebra}. No positivity, admissibility condition or bound was used. The
existence of the generators on the regularity class of \cite[(3.35)--(3.38)]{B-CMP99b} is given by
\cite[Corollary~3.6]{B-CMP99b}, together with the part of \cite[Theorem~3.11]{B-CMP99b} that
concerns these generators. By (i) and (ii) its hypotheses and conclusion are the same for the two
nests, so it is one and the same statement for both.

\emph{Augmented nests.} For an augmented nest, the hypotheses and arguments above apply verbatim.
The only difference is that the bits of the last member are the membership bits
$[e\subset\widetilde{\Omega}_{j+1}]$ of the set $\widetilde{\Omega}_{j+1}$, with the region bits
added to the hypothesis. For a cube $Q$ of
scale $j$, Lemma~\ref{4.6:lem-sequence-reach} reads the member of level $j+1$ of the local
sequence block by block on $Q^{(3)}$. Hence
\begin{equation*}
\Omega_{j+1}(Q)=\bigcup\bigl\{e\in\pi_j : e\subset Q^{(3)},\ e\subset\widetilde{\Omega}_{j+1}\bigr\},
\end{equation*}
that is, the intersection of $\widetilde{\Omega}_{j+1}$ with $Q^{(3)}$ taken block by block.
\end{proof}

\begin{remark}[Equal local sequences and the difference of generators]
The proposition is the mechanism of \cite[Theorem~3.14]{B-CMP99b} in its degenerate case. The two
local sequences coincide, so the difference between the corresponding generators is zero. In
particular, the class index of $Q$ (Lemma~\ref{4.6:lem-membership-tests}) and its restricted
operators are determined by the trace of the nest with region datum on $Q^{(5)}$, at every inner
seam.
\end{remark}

In what follows a nest $\mathbf{\Omega}$ is fixed, and the objects attached to it are as follows.

\begin{itemize}
\item $\mathcal{D}(\mathbf{\Omega})$ is its mixed-size cover, and $h_\square$, for $\square\in\mathcal{D}(\mathbf{\Omega})$, are the functions of Definition~\ref{4.6:def-bump-profiles}.
\item For a cube $\square$, $j$ denotes its scale and $\mathbf{\Omega}(\square)$ its local sequence, with members $\Omega_0(\square)\supseteq\dots$.
\item $\chi_\square$ is the characteristic function of $\Omega_0(\square)$.
\item $G'_\square$, $C_\square$, $P_\square$, $R_\square$, $G_\square$ and $\Delta_{a,\square}$ are the localized generators of Definition~\ref{4.6:def-local-generators}.
\item $D$, $D^*$, $\Delta(U)$, $Q$, $\mathfrak{a}$, $\Delta_a$ and the projection $P=I-R$ are the operators of \cite{B-CMP99b} at the nest $\mathbf{\Omega}$. The subscript $\square$ on $Q_\square$, $\mathfrak{L}_\square$ marks the same objects formed for $\mathbf{\Omega}(\square)$.
\end{itemize}

Commutators are $[X,Y]:=XY-YX$.

\medskip
\noindent\emph{Notation for this section.} Put
\begin{equation*}
\mathfrak{L}:=\Delta(U)+DD^*+Q^*\mathfrak{a}Q ,
\end{equation*}
and $\mathfrak{L}_\square:=\mathfrak{L}[\mathbf{\Omega}(\square)]$ for the same expression formed for the local sequence. The operators $D$, $D^*$ and $\Delta(U)$ are formed from the background configuration $U$ alone, so the terms $\Delta(U)$ and $DD^*$ do not depend on the nest. Then
\begin{equation}\label{4.6:eq-commutator-identity-a}
\Delta_a=\mathfrak{L}-DPD^*
\end{equation}
by the definition of $\Delta_a$ and $P=I-R$ in \cite{B-CMP99b}. Compare \cite[(2.19)]{B-CMP96}, which in the notation of that paper reads $\Delta_a=\Delta-\partial P\partial^*+Q^*aQ$, with $\partial,\partial^*,a$ in place of $D,D^*,\mathfrak{a}$.

Put
\begin{equation*}
G_0:=\sum_{\square\in\mathcal{D}(\mathbf{\Omega})}h_\square G_\square h_\square ,
\end{equation*}
as in \cite[(3.87)]{B-CMP99b}; the one-scale forms are \cite[(1.119)]{B-CMP95} and \cite[(2.37)]{B-CMP96}. For a function $h$ put
\begin{equation*}
K(h):=h\mathfrak{L}-\mathfrak{L}h,\qquad P_1(\partial h):=[DPD^*,h],\qquad
P_\square(\partial h):=[DP_\square D^*,h].
\end{equation*}
For each $\square$, $\zeta_\square$ is the cut-off function of \cite[(3.105)]{B-CMP99b}, which satisfies $\zeta_\square h_\square=h_\square\zeta_\square=h_\square$.

The \emph{level function} of a nest assigns to each averaging block the index $n$ of the layer $\Lambda_n$ of that nest which contains it. We say that the pair $(\mathbf{\Omega},\square)$ has \emph{level agreement} if
\begin{equation}\label{4.6:eq-commutator-identity-b}
(\mathrm{LA}_{\square}):\quad
\begin{gathered}
\text{the level function of }\mathbf{\Omega}\text{ and the level function of }
\mathbf{\Omega}(\square)\\
\text{agree block by block on }\square^{(2)}.
\end{gathered}
\end{equation}
Here $\square^{(2)}$ contains $\operatorname{supp}h_\square\subset\square$ together with a margin of two averaging blocks of level at most $j+1$. The condition $(\mathrm{LA}_{\square})$ is a property of the pair $(\mathbf{\Omega},\square)$. It is a hypothesis of part~(3) below and of no locality statement of this section.

\medskip
\noindent\emph{The corresponding passages of \cite{B-CMP99b}.} The commutators of $h$ with $D^*D$, $DD^*$ and $\Delta'$, where $\Delta(U)=D^*D+\Delta'$, are treated around \cite[(3.100)]{B-CMP99b}. There $[D,h]$ is of order zero with coefficients determined by derivatives of $h$, and $\Delta'$ is small and local by \cite[(3.10)]{B-CMP99b}. The commutator $[DRD^*,h]$ is treated through $DRD^*=DD^*-DPD^*$ and \cite[(3.101)]{B-CMP99b}, and the commutator with $Q^*\mathfrak{a}Q$ in \cite[(3.102)--(3.103)]{B-CMP99b}.

The resulting identity \cite[(3.104)]{B-CMP99b}, applied to $\Delta_aG_0$, gives \cite[(3.105)]{B-CMP99b}. After that, $P$ is replaced by its random walk expansion, and only two terms of the second and third sums need to be considered.

The one-scale precursors are \cite[(1.122), (1.123)]{B-CMP95}, the corresponding identity for the one-scale partition of unity of \cite[(1.118), (1.119)]{B-CMP95} and the resulting walk expansion of the propagator of that paper, and \cite[(2.38), (2.50)]{B-CMP96}.

\begin{lemma}[The commutator identity with the local projection]
\label{4.6:lem-commutator-identity}
Fix the nest $\mathbf{\Omega}$ and a cube $\square\in\mathcal{D}(\mathbf{\Omega})$. Write $h:=h_\square$, $G:=G_\square$, $\zeta:=\zeta_\square$, and let $P_\square$ and $\mathfrak{L}_\square$ be as above.
\begin{enumerate}
\item[(1)] The identity
\begin{equation*}
\Delta_ah=h\Delta_a-K(h)-P_1(\partial h)
\end{equation*}
holds with
\begin{equation*}
K(h)=[h,\Delta(U)]+[h,DD^*]+[h,Q^*\mathfrak{a}Q].
\end{equation*}
This is \cite[(3.104)]{B-CMP99b}, the operator $K(h)$ there being the one displayed here.

\item[(2)] Without further hypotheses,
\begin{equation}\label{4.6:eq-commutator-identity-1}
\begin{split}
\Delta_a hGh={}&h(\mathfrak{L}-DP_\square D^*)Gh-K(h)Gh-(1-\zeta)DPD^*hGh\\
&-\zeta(DPD^*-DP_\square D^*)hGh-\zeta P_\square(\partial h)Gh .
\end{split}
\end{equation}

\item[(3)] Assume that $\operatorname{supp}h\subset\square$ and that $(\mathrm{LA}_{\square})$ holds. The support condition is satisfied by the one-scale profiles of Definition~\ref{4.6:def-bump-profiles}, see \eqref{4.6:eq-bump-profiles-a}, and by the glued bumps under the support hypothesis $(\mathrm{G\text{-}supp})$ of that definition. Then
\begin{equation}\label{4.6:eq-commutator-identity-2}
h(\mathfrak{L}-DP_\square D^*)G_\square h=h^2 .
\end{equation}

\item[(4)] Assume that the conclusion of part~(3) holds for every $\square\in\mathcal{D}(\mathbf{\Omega})$ and that $\sum_{\square}h_\square^2=1$, as in \cite[(2.36)]{B-CMP96}. Then
\begin{equation}\label{4.6:eq-commutator-identity-3}
\begin{split}
\Delta_aG_0=I&-\sum_\square K(h_\square)G_\square h_\square
-\sum_\square(1-\zeta_\square)DPD^*h_\square G_\square h_\square\\
&-\sum_\square\zeta_\square(DPD^*-DP_\square D^*)h_\square G_\square h_\square
-\sum_\square\zeta_\square P_\square(\partial h_\square)G_\square h_\square .
\end{split}
\end{equation}
This is \cite[(3.105)]{B-CMP99b} term by term, the projection in the commutator of the fourth sum being read as the local projection $P_\square$; part~(5) shows that this is the only reading under which that display is an identity of the displayed shape.

\item[(5)] Let $\mathbb{F}$ be the free associative $\mathbb{Z}$-algebra on the five letters $A,B,G,h,\zeta$ modulo the two-sided ideal generated by $\zeta h-h$ and $h\zeta-h$; in symbols, $\mathbb{F} := \mathbb{Z}\langle A,B,G,h,\zeta\rangle/(\zeta h-h,\, h\zeta-h)$. The letter $A$ stands for $DPD^*$, the letter $B$ for $DP_\square D^*$, and the letters $G,h,\zeta$ for $G_\square,h_\square,\zeta_\square$; under this assignment the elements of $\mathbb{F}$ become operators. In $\mathbb{F}$ put
\begin{align*}
\mathsf{r}&:=-h(A-B)Gh-[A,h]Gh,\\
\mathsf{r}_{\mathrm{loc}}&:=-(1-\zeta)AhGh-\zeta(A-B)hGh-\zeta[B,h]Gh,\\
\mathsf{r}_{\mathrm{glob}}&:=-(1-\zeta)AhGh-\zeta(A-B)hGh-\zeta[A,h]Gh .
\end{align*}
Then $\mathsf{r}=\mathsf{r}_{\mathrm{loc}}$. Moreover
\begin{equation}\label{4.6:eq-commutator-identity-4}
\mathsf{r}-\mathsf{r}_{\mathrm{glob}}=hBGh+\zeta AhGh-\zeta BhGh-hAGh\neq0
\quad\text{in }\mathbb{F}.
\end{equation}
As operators, summed over the cubes, this difference is
\begin{equation*}
\sum_\square\zeta_\square\,[D(P-P_\square)D^*,h_\square]\,G_\square h_\square .
\end{equation*}
Of this operator only the non-vanishing of the corresponding element of $\mathbb{F}$ is asserted and used. Consequently, of the two readings of the projection in the commutator of the fourth sum of \cite[(3.105)]{B-CMP99b}, the global $P$ and the local $P_\square$, only $P_\square$ makes that display an identity of the displayed shape.

\item[(6)] The second piece of the third sum of \eqref{4.6:eq-commutator-identity-3} is
\begin{equation*}
+\sum_\square\zeta_\square DP_\square D^*h_\square G_\square h_\square ,
\end{equation*}
with the local projection $P_\square$. In the treatment of the sums of \cite[(3.105)]{B-CMP99b}, this term is cancelled by the product
\begin{equation*}
\zeta_\square DG'_\square Q'^*C_\square Q'G'_\square D^*h_\square G_\square h_\square .
\end{equation*}
The operator $Q'$ is written there without subscript, and we read it as $Q'_\square$. With this reading the product equals $\zeta_\square DP_\square D^*h_\square G_\square h_\square$, by the definition
\begin{equation*}
P_\square=G'_\square Q'^*_\square C_\square Q'_\square G'_\square ,
\end{equation*}
which is \cite[(3.25)]{B-CMP99b} for the sequence $\mathbf{\Omega}(\square)$. Hence the operator $G_\square$ there is the one built with the local projection $R_\square$ of Definition~\ref{4.6:def-local-generators}.
\end{enumerate}
\end{lemma}

\begin{proof}
\emph{Part (1).} By \eqref{4.6:eq-commutator-identity-a},
\begin{equation*}
\Delta_ah=\mathfrak{L}h-DPD^*h .
\end{equation*}
By the definition of $K(h)$ we have $\mathfrak{L}h=h\mathfrak{L}-K(h)$. By the definition of $P_1(\partial h)$, which is the commutator of \cite[(3.101)]{B-CMP99b}, we have $DPD^*h=hDPD^*+P_1(\partial h)$. Hence
\begin{equation*}
\Delta_ah=\bigl(h\mathfrak{L}-K(h)\bigr)-\bigl(hDPD^*+P_1(\partial h)\bigr)
=h\Delta_a-K(h)-P_1(\partial h),
\end{equation*}
again by \eqref{4.6:eq-commutator-identity-a}. This is the identity \cite[(3.104)]{B-CMP99b}. The operator $K(h)$ occurring in it is therefore
\begin{equation*}
h\mathfrak{L}-\mathfrak{L}h=[h,\Delta(U)]+[h,DD^*]+[h,Q^*\mathfrak{a}Q],
\end{equation*}
by the definition of $\mathfrak{L}$ and bilinearity of the commutator.

This agrees with the description of $K(h)$ given after \cite[(3.104)]{B-CMP99b}. By \cite[(3.100)]{B-CMP99b} the commutator $[D,h]$ is of order zero with coefficients given by $\partial h$. Hence $[h,D^*D]$ and $[h,DD^*]$ are first order operators with coefficients in $\partial h$. Further, $\Delta(U)=D^*D+\Delta'$ \cite{B-CMP99b}, with $\Delta'$ a local operator of order zero. Finally
\begin{equation*}
[h,Q^*\mathfrak{a}Q]=S^*(\partial h)\mathfrak{a}Q-Q^*\mathfrak{a}S(\partial h)
\end{equation*}
by \cite[(3.103)]{B-CMP99b}, with the operator $S(\partial h)$ defined there. These two terms are the last two terms on the right-hand side of \cite[(3.103)]{B-CMP99b}.

\emph{Part (2).} By \eqref{4.6:eq-commutator-identity-a},
\begin{equation*}
\Delta_ahGh=\mathfrak{L}hGh-DPD^*hGh,
\end{equation*}
and as in part~(1), $\mathfrak{L}hGh=(h\mathfrak{L}-K(h))Gh$. Split
\begin{equation*}
DPD^*h=(1-\zeta)DPD^*h+\zeta(DPD^*-DP_\square D^*)h+\zeta DP_\square D^*h .
\end{equation*}
By the definition of $P_\square(\partial h)$ and $\zeta h=h$,
\begin{equation*}
\zeta DP_\square D^*h=\zeta\bigl(hDP_\square D^*+P_\square(\partial h)\bigr)
=hDP_\square D^*+\zeta P_\square(\partial h).
\end{equation*}
Subtracting the three pieces of $DPD^*hGh$ from $(h\mathfrak{L}-K(h))Gh$ and collecting $h\mathfrak{L}Gh-hDP_\square D^*Gh=h(\mathfrak{L}-DP_\square D^*)Gh$ gives \eqref{4.6:eq-commutator-identity-1}.

\emph{Part (3).} Let $\phi$ be any element of the space on which the operators act, and put $F:=G_\square h\phi$. By the definition of $G_\square$, the function $F$ is supported in $\Omega_0(\square)$, that is $F=\chi_\square F$, and $\Omega_0(\square)\subset\square^{5}$.

(a) We show that $h\mathfrak{L}F=h\mathfrak{L}_\square F$.

The value $(h\mathfrak{L}F)(b)$ is non-zero only for bonds $b\in\operatorname{supp}h\subset\square$. At such $b$, the values $(\Delta(U)F)(b)$ and $(DD^*F)(b)$ read $F$ only at bonds within one lattice step of $b$. Moreover these two operators are the same in $\mathfrak{L}$ and in $\mathfrak{L}_\square$.

The value $(Q^*\mathfrak{a}QF)(b)$ reads two things. First, for each level $n$, whether the bonds $c$ of the $n$-lattice within one $n$-block of $b$ lie in $\Lambda_n$. Second, the values of $F$ within two $n$-blocks of $b$. Under $(\mathrm{LA}_{\square})$ the levels $n$ occurring near $\operatorname{supp}h$ satisfy $n\le j+1$. Since $u_j=ML^j\eta$, with $\eta$ the lattice spacing of \cite{B-CMP99b}, and $M\ge L$, $M$ being a positive integer power of $L$ in \cite{B-CMP99b}, we have $L^{j+1}\eta\le u_j$. Hence everything that $(Q^*\mathfrak{a}QF)(b)$ reads lies inside $\square^{(2)}$.

On $\square^{(2)}$, by $(\mathrm{LA}_{\square})$, the layers $\Lambda_n$ of $\mathbf{\Omega}$ and $\Lambda_n(\square)$ of $\mathbf{\Omega}(\square)$ have the same trace. The averaging operators $Q_n(U)$ and the numbers $a_n$ are given by the same formulas for both. Therefore $Q^*\mathfrak{a}QF=Q_\square^*\mathfrak{a}Q_\square F$ on $\operatorname{supp}h$, and (a) follows.

(b) Since
\begin{equation*}
\operatorname{supp}h\subset\square\subset\square^{4}=\Omega_j(\square)\subset\Omega_0(\square),
\end{equation*}
we have $\chi_\square h=h\chi_\square=h$. Together with $F=\chi_\square F$ this gives
\begin{equation*}
h\mathfrak{L}_\square F=h\chi_\square\mathfrak{L}_\square\chi_\square F,\qquad
hDP_\square D^*F=h\chi_\square DP_\square D^*\chi_\square F .
\end{equation*}

(c) By (a) and (b),
\begin{equation*}
\begin{split}
h(\mathfrak{L}-DP_\square D^*)G_\square h\phi
&=h\chi_\square(\mathfrak{L}_\square-DP_\square D^*)\chi_\square G_\square h\phi\\
&=h\Delta_{a,\square}G_\square h\phi=h\cdot h\phi .
\end{split}
\end{equation*}
In the second equality, $\chi_\square(\mathfrak{L}_\square-DP_\square D^*)\chi_\square$ is the operator $\Delta_{a,\square}$ of Definition~\ref{4.6:def-local-generators}, that is, \eqref{4.6:eq-commutator-identity-a} for $\mathbf{\Omega}(\square)$ on $\Omega_0(\square)$. In the third, $G_\square$ is the inverse of $\Delta_{a,\square}$ there, and $h\phi$ is supported in $\square\subset\Omega_0(\square)$. As $\phi$ was arbitrary, \eqref{4.6:eq-commutator-identity-2} follows.

\emph{Part (4).} By the definition of $G_0$,
\begin{equation*}
\Delta_aG_0=\sum_\square\Delta_ah_\square G_\square h_\square .
\end{equation*}
Apply \eqref{4.6:eq-commutator-identity-1} to each summand. By the hypothesis, \eqref{4.6:eq-commutator-identity-2} turns its first term into $h_\square^2$. Since $\sum_\square h_\square^2=1$, the first terms sum to $I$, and the remaining terms give the four sums of \eqref{4.6:eq-commutator-identity-3}.

\emph{Part (5).} Since $\zeta_\square h_\square=h_\square=h_\square\zeta_\square$, the assignment
\begin{equation*}
A\mapsto DPD^*,\quad B\mapsto DP_\square D^*,\quad G\mapsto G_\square,\quad h\mapsto h_\square,\quad \zeta\mapsto\zeta_\square
\end{equation*}
defines a ring homomorphism from $\mathbb{F}$ to operators. By part~(1), the image of $\mathsf{r}$ is what remains of $\Delta_ahGh$ after $h(\mathfrak{L}-DP_\square D^*)Gh-K(h)Gh$ is subtracted. Indeed
\begin{equation*}
\begin{split}
\Delta_ahGh&=h(\mathfrak{L}-DPD^*)Gh-K(h)Gh-P_1(\partial h)Gh\\
&=h(\mathfrak{L}-DP_\square D^*)Gh-K(h)Gh-h(A-B)Gh-[A,h]Gh
\end{split}
\end{equation*}
under this assignment.

In $\mathbb{F}$, using the relations $\zeta h=h$ and $h\zeta=h$, we compute
\begin{equation*}
\mathsf{r}=-hAGh+hBGh-AhGh+hAGh=hBGh-AhGh
\end{equation*}
and
\begin{equation*}
\begin{split}
\mathsf{r}_{\mathrm{loc}}
&=-AhGh+\zeta AhGh-\zeta AhGh+\zeta BhGh-\zeta BhGh+\zeta hBGh\\
&=-AhGh+hBGh .
\end{split}
\end{equation*}
Hence $\mathsf{r}=\mathsf{r}_{\mathrm{loc}}$. For the global reading,
\begin{equation*}
\begin{split}
\mathsf{r}_{\mathrm{glob}}
&=-AhGh+\zeta AhGh-\zeta AhGh+\zeta BhGh-\zeta AhGh+\zeta hAGh\\
&=-AhGh+\zeta BhGh-\zeta AhGh+hAGh ,
\end{split}
\end{equation*}
so that
\begin{equation*}
\mathsf{r}-\mathsf{r}_{\mathrm{glob}}=hBGh-\zeta BhGh+\zeta AhGh-hAGh ,
\end{equation*}
which is the element in \eqref{4.6:eq-commutator-identity-4}. Using $\zeta h=h$ once more, this element equals
\begin{equation*}
\zeta(A-B)hGh-\zeta h(A-B)Gh=\zeta[A-B,h]Gh .
\end{equation*}
Its image under the homomorphism is $\zeta_\square[D(P-P_\square)D^*,h_\square]G_\square h_\square$. Summed over $\square$, this gives the displayed operator.

It remains to show that the element is non-zero. Consider the rewriting system $\{\zeta h\to h,\ h\zeta\to h\}$ on words in $A,B,G,h,\zeta$. Each rewriting shortens a word by one letter, so the system terminates. There are no inclusion ambiguities. Its overlap ambiguities are $\zeta h\zeta$ and $h\zeta h$, and both resolve:
\begin{itemize}
\item $(\zeta h)\zeta\to h\zeta\to h$ and $\zeta(h\zeta)\to\zeta h\to h$;
\item $(h\zeta)h\to hh$ and $h(\zeta h)\to hh$.
\end{itemize}
A rewriting system that terminates and whose ambiguities all resolve is confluent; hence every element of $\mathbb{Z}\langle A,B,G,h,\zeta\rangle$ has a unique irreducible representative modulo the ideal, and the irreducible words form a $\mathbb{Z}$-basis of $\mathbb{F}$, so normal forms are unique. The words $hBGh$, $\zeta AhGh$, $\zeta BhGh$ and $hAGh$ contain neither $\zeta h$ nor $h\zeta$, so they are irreducible, and they are pairwise distinct. They occur with coefficients $+1,+1,-1,-1$. Hence the element is non-zero in $\mathbb{F}$.

The reading statement follows. Under the local reading, by part~(4) and $\mathsf{r}=\mathsf{r}_{\mathrm{loc}}$, the display \cite[(3.105)]{B-CMP99b} is an identity. The global reading differs from it by the non-zero element \eqref{4.6:eq-commutator-identity-4}.

\emph{Part (6).} By parts~(2) and~(4), the third sum of \eqref{4.6:eq-commutator-identity-3} consists of the terms
\begin{equation*}
-\zeta_\square(DPD^*-DP_\square D^*)h_\square G_\square h_\square
=-\zeta_\square DPD^*h_\square G_\square h_\square
+\zeta_\square DP_\square D^*h_\square G_\square h_\square .
\end{equation*}
Their second piece is $+\zeta_\square DP_\square D^*h_\square G_\square h_\square$, with the local projection $P_\square$. In \cite{B-CMP99b} it is stated that $\zeta_\square h_\square=h_\square$ and that the product displayed in part~(6) of the statement cancels the second term in the third sum. With $Q'=Q'_\square$, the definition of $P_\square$ displayed in part~(6), which is \cite[(3.25)]{B-CMP99b} for $\mathbf{\Omega}(\square)$, shows that this product is $\zeta_\square DP_\square D^*h_\square G_\square h_\square$. This is exactly the second piece above.
\end{proof}

\begin{corollary}[Completeness of the two substitutions in the commutator identity]
\label{4.6:cor-two-substitutions}
Fix a nest $\mathbf{\Omega}$, its cover $\mathcal{D}(\mathbf{\Omega})$, and for
$\square\in\mathcal{D}(\mathbf{\Omega})$ the functions $h_\square$, the local propagators
$G_\square$, the local projections $P_\square$ and the cut-offs $\zeta_\square$ of
Lemma~\ref{4.6:lem-commutator-identity}. Assume the hypotheses of clause~(4) of that lemma, so
that the identity \cite[(3.105)]{B-CMP99b} holds, with $G_0=\sum_\square h_\square G_\square
h_\square$, in the form $\Delta_aG_0=I-R$ and
\begin{equation}\label{4.6:eq-two-substitutions-1}
\begin{split}
R={}&\sum_{\square}K(h_\square)G_\square h_\square
+\sum_{\square}(1-\zeta_\square)DPD^*h_\square G_\square h_\square\\
&+\sum_{\square}\zeta_\square\bigl(DPD^*-DP_\square D^*\bigr)h_\square G_\square h_\square
+\sum_{\square}\zeta_\square P_\square(\partial h_\square)G_\square h_\square .
\end{split}
\end{equation}
Here $Q'$ is the averaging operator of \cite{B-CMP99b}, $\Delta'_a$ the associated operator,
whose one-scale form in \cite[p.~225]{B-CMP96} is $\Delta+Q'^*\mathfrak{a}'Q'$, and
$G'=(\Delta'_a)^{-1}$ its inverse, which enters the expansion \cite[(3.90)]{B-CMP99b}, with local
versions $Q'_\square$ and $G'_\square$ at the cubes $\square$; and $C$ is the operator of
\cite{B-CMP99b} inverting $Q'G'^2Q'^*$, whose expansion in \cite{B-CMP99b} starts from the identity
\cite[(3.95)]{B-CMP99b}, with local inverses $C_\square$. We call an operator local if it is
attached to a single cube $\square$ of the cover (such as $h_\square$, $\zeta_\square$,
multiplication by the indicator function of the cube $\widetilde{\square}$ of
\eqref{4.6:eq-candidate-cubes-b}, $G_\square$, $G'_\square$, $P_\square$,
$P_\square(\partial h_\square)$, $C_\square$, $K(h_\square)$, $K'(h_\square)$) or has bounded
range ($D$, $D^*$, $Q'$, $Q'^*$), and global otherwise; the global operators in question are $P$,
the projection of \eqref{4.6:eq-commutator-identity-a}, $G'$ and $C$.
\begin{enumerate}
\item[(i)] The global projection $P$ occurs in \eqref{4.6:eq-two-substitutions-1} in exactly
two places: in the second sum and in the first piece of the third sum. The second piece of the
third sum and the whole fourth sum contain only the local projections $P_\square$, and the
first sum contains no projection. Both occurrences of $P$ are substituted in the text following
\cite[(3.106)]{B-CMP99b}.
\item[(ii)] The expansion \cite[(3.90)]{B-CMP99b} of $G'$ contains only local operators
$h_\square$, $G'_\square$, $K'(h_\square)$, $\square\in\mathcal{D}(\mathbf{\Omega})$.
\item[(iii)] No global operator occurs in the expansion \cite[(3.98)]{B-CMP99b}.
\item[(iv)] Consequently the two substitutions of (i) are complete: no global $P$, $G'$ or $C$
remains in the expansion of $R$, nor in $G=G_0(I-R)^{-1}$ of \cite[(3.106)]{B-CMP99b}.
\end{enumerate}
\end{corollary}

\begin{proof}
\emph{(i).} The sums in \eqref{4.6:eq-two-substitutions-1} are the sums over
$\square\in\mathcal{D}(\mathbf{\Omega})$ of the terms $K(h)Gh$, $(1-\zeta)DPD^*hGh$,
$\zeta(DPD^*-DP_\square D^*)hGh$ and $\zeta P_\square(\partial h)Gh$ of the per-cube identity in
clause~(2) of Lemma~\ref{4.6:lem-commutator-identity}, written with $h=h_\square$,
$G=G_\square$, $\zeta=\zeta_\square$; the remaining term $h(\mathfrak{L}-DP_\square D^*)G_\square h$
of that identity equals $h^2$ by clause~(3), and these terms add up to $I$ because
$\sum_\square h_\square^2=1$. We inspect the four sums in turn. In the first, clause~(1) of
Lemma~\ref{4.6:lem-commutator-identity} gives
\begin{equation*}
K(h_\square)=[h_\square,\Delta]+[h_\square,DD^*]+[h_\square,Q^*\mathfrak{a}Q],
\end{equation*}
which contains no projection, and the propagator is the local $G_\square$. The second sum
carries the global operator $DPD^*$. The third sum splits into the piece with $DPD^*$, which
carries the global $P$, and the piece with $DP_\square D^*$, which carries only the local
$P_\square$. The fourth sum carries only $P_\square(\partial h_\square)$. Hence the global $P$
occurs exactly in the second sum and in the first piece of the third sum. Both of these
occurrences are substituted in the text following \cite[(3.106)]{B-CMP99b}.

\emph{(ii).} The operator $K'(h)$ of \cite[(3.88)]{B-CMP99b} is the operator $K(h)$ of
\cite[(2.38)]{B-CMP96}. Put, as in \cite[(2.37)]{B-CMP96},
$G'_0=\sum_\square h_\square G'_\square h_\square$. The analogues for $\Delta'_a$ of clauses~(1)
and~(3) of Lemma~\ref{4.6:lem-commutator-identity} are
$\Delta'_ah_\square=h_\square\Delta'_a-K'(h_\square)$ and
$h_\square\Delta'_aG'_\square h_\square=h_\square^2$; together they give
\begin{align*}
\Delta'_ah_\square G'_\square h_\square
&=h_\square\Delta'_aG'_\square h_\square-K'(h_\square)G'_\square h_\square\\
&=h_\square^2-K'(h_\square)G'_\square h_\square ,
\end{align*}
and summing over $\square$ with $\sum_\square h_\square^2=1$ gives
\begin{equation}\label{4.6:eq-two-substitutions-2}
\Delta'_aG'_0=I-\sum_{\square}K'(h_\square)G'_\square h_\square=:I-R'.
\end{equation}
Since $G'$ is the inverse of $\Delta'_a$, \eqref{4.6:eq-two-substitutions-2} gives
$G'=G'_0(I-R')^{-1}=G'_0\sum_{n\ge0}R'^{\,n}$. Inserting the definitions of $G'_0$ and $R'$
into each product $G'_0R'^{\,n}$, we read \cite[(3.90)]{B-CMP99b} as the expansion
\begin{equation}\label{4.6:eq-two-substitutions-3}
G'=\sum_{\omega}h_{\square_0}G'_{\square_0}h_{\square_0}\,
K'(h_{\square_1})G'_{\square_1}h_{\square_1}\cdots
K'(h_{\square_n})G'_{\square_n}h_{\square_n},
\end{equation}
where $\omega=(\square_0,\dots,\square_n)$, $\square_i\in\mathcal{D}(\mathbf{\Omega})$ and
$\square_i\cap\square_{i+1}\neq\emptyset$. This is the form of \cite[(2.50)]{B-CMP96} with
$G'(\square_i)$ replaced by $G'_{\square_i}$ and $K$ by $K'$, and it is also the form of
\cite[(1.123)]{B-CMP95}. Every factor in \eqref{4.6:eq-two-substitutions-3} is one of
$h_{\square_i}$, $G'_{\square_i}$, $K'(h_{\square_i})$, attached to a single cube of
$\mathcal{D}(\mathbf{\Omega})$; no modification of the cover enters. This proves (ii).

\emph{(iii).} Write $\widetilde{\square}$ also for multiplication by the indicator function of
the cube $\widetilde{\square}$ of size $4M$ attached to $\square$ in
\eqref{4.6:eq-candidate-cubes-b}, and let $C_0=\sum_\square h_\square C_\square h_\square$ be the
operator of \cite[(3.95)]{B-CMP99b}, the analogue of the glued inverse of \cite[(2.70)]{B-CMP96}.
We read \cite[(3.95)]{B-CMP99b} as the identity
\begin{equation}\label{4.6:eq-two-substitutions-4}
\begin{split}
Q'G'^2Q'^*C_0={}&I+\sum_{\square}(1-\widetilde{\square})Q'G'^2Q'^*h_\square C_\square h_\square\\
&+\sum_{\square}\widetilde{\square}Q'\bigl(G'^2-G'^2_\square\bigr)Q'^*h_\square C_\square h_\square\\
&+\sum_{\square}\bigl[\widetilde{\square}Q'G'^2_\square Q'^*,h_\square\bigr]C_\square h_\square .
\end{split}
\end{equation}
The operators $C_\square$ here are the multi-scale form of the operators
$C_\square=\bigl((Q'G'(\widetilde{\square})^2Q'^*)\restriction_\square\bigr)^{-1}$ of
\cite[(2.70)]{B-CMP96}, built with the same cube $\widetilde{\square}$ of size $4M$. The
identity \eqref{4.6:eq-two-substitutions-4} is obtained by telescoping: for each $\square$,
\begin{align*}
Q'G'^2Q'^*h_\square C_\square h_\square
={}&(1-\widetilde{\square})Q'G'^2Q'^*h_\square C_\square h_\square\\
&+\widetilde{\square}Q'\bigl(G'^2-G'^2_\square\bigr)Q'^*h_\square C_\square h_\square\\
&+\bigl[\widetilde{\square}Q'G'^2_\square Q'^*,h_\square\bigr]C_\square h_\square
+h_\square\widetilde{\square}Q'G'^2_\square Q'^*C_\square h_\square ,
\end{align*}
and the last term equals $h_\square^2$, because $\widetilde{\square}h_\square=h_\square$ and
$Q'=Q'_\square$ near $\square$ under \eqref{4.6:eq-commutator-identity-b}. Summing over
$\square$ and using $\sum_\square h_\square^2=1$ gives \eqref{4.6:eq-two-substitutions-4}.

In the first and second sums of \eqref{4.6:eq-two-substitutions-4} the operator $G'^2$ is
global. The series \cite[(3.96)]{B-CMP99b} for $C$, obtained from
\eqref{4.6:eq-two-substitutions-4}, is treated in \cite{B-CMP99b} by replacing the
operator $G'$ everywhere in it by an expansion similar to \cite[(3.90)]{B-CMP99b}; thus every
global factor $G'^2$ is replaced by the expansion of the form \eqref{4.6:eq-two-substitutions-3}
over the cover $\mathcal{D}'$, whose factors are again attached to single cubes. The regrouping
carried out in \cite{B-CMP99b} after this replacement produces, besides the terms of
\cite[(3.97)]{B-CMP99b}, a commutator term in which, as in the third sum of
\eqref{4.6:eq-two-substitutions-4}, only operators attached to single cubes occur. By the form
of \eqref{4.6:eq-two-substitutions-3} inserted in the second sum of
\eqref{4.6:eq-two-substitutions-4}, we read the terms of \cite[(3.97)]{B-CMP99b} as
\begin{equation*}
\widetilde{\square}Q'\bigl(G'^2_{\square_0}-G'^2_\square\bigr)Q'^*h_\square C_\square h_\square ,
\end{equation*}
where $\square_0$ is a cube of the expansion over $\mathcal{D}'$ substituted for $G'$. This is
a difference of two operators of local sequences. By the identity $A^2-B^2=(A-B)A+B(A-B)$ with
$A=G'_{\square_0}$, $B=G'_\square$, it factorises as
\begin{equation*}
\begin{split}
&\widetilde{\square}Q'\bigl(G'_{\square_0}-G'_\square\bigr)G'_{\square_0}Q'^*
h_\square C_\square h_\square\\
&\quad+\widetilde{\square}Q'G'_\square\bigl(G'_{\square_0}-G'_\square\bigr)Q'^*
h_\square C_\square h_\square ,
\end{split}
\end{equation*}
in which only the local operators $G'_{\square_0}$, $G'_\square$ occur, together with
$\widetilde{\square}$, $Q'$, $h_\square$ and $C_\square$. The terms of the third sum of
\eqref{4.6:eq-two-substitutions-4} are stated in \cite{B-CMP99b} to be already localized, and
the commutator term produced by the regrouping
contains, as noted above, only operators attached to single cubes. Hence no
global operator occurs among the factors of \cite[(3.98)]{B-CMP99b}. This proves (iii).

\emph{(iv).} By (i), the global $P$ enters \eqref{4.6:eq-two-substitutions-1} only at the two
places substituted in \cite{B-CMP99b}; every other factor of \eqref{4.6:eq-two-substitutions-1}
is one of $h_\square$, $\zeta_\square$, $G_\square$, $P_\square$, $P_\square(\partial h_\square)$,
$K(h_\square)$, $D$, $D^*$. The expansions of $G'$ and of $C$
used in \cite{B-CMP99b} are \cite[(3.90)]{B-CMP99b} and \cite[(3.98)]{B-CMP99b}, which by (ii)
and (iii) contain no global operator. Therefore the two substitutions leave no global $P$,
$G'$ or $C$ in the expansion of $R$. Finally $G=G_0(I-R)^{-1}$ of \cite[(3.106)]{B-CMP99b} is
built from $G_0=\sum_\square h_\square G_\square h_\square$, a sum of local operators, and
from $R$ alone, so no global $P$, $G'$ or $C$ remains in it either.
\end{proof}

\begin{lemma}[Level agreement near an inner seam]\label{4.6:lem-inner-seam}
Let $\mathbf{\Omega}=(\Omega_0\supseteq\Omega_1\supseteq\dots\supseteq\Omega_k)$ be a nest as in
Definition~\ref{4.6:def-nest-traces}, with $\Omega_{k+1}=\Omega_{k+2}=\emptyset$. Let $0\le j\le k$
and let $\square\in\mathcal{D}_j$ be a cube of the cover of Definition~\ref{4.6:def-local-sequence}.
Thus $\square=x+[-u_j,u_j]^d$, and $\square$ contains a $\pi_j$-block $e$ with
$[e\subset\Omega_j]=1$ and $[e\subset\Omega_{j+1}]=0$. Distances are $\ell^\infty$-distances on
$\mathbb{T}$, with $\operatorname{dist}(A,\emptyset)=+\infty$. For $A\subset\mathbb{T}$ and $s\ge0$,
$A^{+s}$ is the set of points at distance at most $s$ from $A$, and $\square^{(2)}=\square^{+2u_j}$.
\begin{enumerate}
\item[(a)] In the local sequence $\mathbf{\Omega}(\square)$, every $\pi_j$-block of $\square^{(2)}$ has
level $j$ or $j+1$, according to its $\Omega_{j+1}$-bit. Consequently the level agreement
$(\mathrm{LA}_{\square})$ of \eqref{4.6:eq-commutator-identity-b} holds if and only if every
$\pi_j$-block $b\subset\square^{(2)}$ satisfies $[b\subset\Omega_j]=1$ and
$[b\subset\Omega_{j+2}]=0$. In particular, $(\mathrm{LA}_{\square})$ holds whenever
$\square^{(2)}\subset\Omega_j\setminus\Omega_{j+2}$.
\item[(b)] \emph{Inner clause.} Let $R$ be a positive integer. Suppose the nest satisfies at level
$j+1$ the separation condition of \cite[(2.2)]{B-CMP96}, which there reads
\begin{equation*}
(L^{j+1}\eta)^{-1}\operatorname{dist}\bigl(\mathbb{T}\setminus\Omega_{j+1},\,\Omega_{j+2}\bigr)>RM,
\end{equation*}
with $\eta$ the lattice spacing of \cite{B-CMP96}. Since a $\pi_{j+1}$-block has side
$u_{j+1}=ML^{j+1}\eta$ in those units, the condition reads
\begin{equation}\label{4.6:eq-inner-seam-1}
\operatorname{dist}\bigl(\mathbb{T}\setminus\Omega_{j+1},\,\Omega_{j+2}\bigr)>R\,u_{j+1},
\end{equation}
and suppose
\begin{equation}\label{4.6:eq-inner-seam-2}
(R+1)L\ge 4 .
\end{equation}
Then $\square^{(2)}\cap\Omega_{j+2}=\emptyset$.
\item[(b$'$)] Let $L\ge2$. If every $\pi_{j+2}$-block that meets $\Omega_{j+2}$ is contained in
$\Omega_{j+1}$, then $\square^{(2)}\cap\Omega_{j+2}=\emptyset$.
\item[(c)] \emph{Outer clause.} Suppose the centre $x$ of $\square$ lies at distance at most $2u_j$
from the boundary $\partial\Omega_j$ of $\Omega_j$ in $\mathbb{T}$. Then
$\square^{(2)}\not\subset\Omega_j$, and $(\mathrm{LA}_{\square})$ fails at $\square$.
\end{enumerate}
\end{lemma}

\begin{proof}
We use one elementary fact repeatedly.

\emph{Fact (F).} Let $g\le g'$, let $b$ be a $\pi_g$-block, and let $A$ be a union of
$\pi_{g'}$-blocks. Then either $b\subset A$ or $\operatorname{int}b\cap A=\emptyset$.

To see this, note that the grids are nested: a set that is $\pi_{g'}$-grained, that is, a union of
$\pi_{g'}$-blocks, is also $\pi_g$-grained. Hence $A$ is a union of $\pi_g$-blocks. Distinct
$\pi_g$-blocks meet $b$ only in
$\partial b$. So if $b$ is not among the $\pi_g$-blocks of $A$, then $A\cap\operatorname{int}b=\emptyset$.
If $b$ is among them, then $b\subset A$.

Fact (F) has a consequence for the closure of a complement. The set
$\operatorname{cl}(\mathbb{T}\setminus A)$ is the union of the $\pi_g$-blocks not contained in $A$.
Indeed, each such block $b$ has $\operatorname{int}b\subset\mathbb{T}\setminus A$, hence
$b=\operatorname{cl}\operatorname{int}b\subset\operatorname{cl}(\mathbb{T}\setminus A)$. Conversely,
$\mathbb{T}\setminus A$ lies in the union of these blocks, and that union is closed.

\emph{(a).} By Definition~\ref{4.6:def-local-sequence}, $\square^{(2)}\subset\square^{(4)}=\Omega_j(\square)$.
On $\square^{(2)}$ the level of a $\pi_j$-block in $\mathbf{\Omega}(\square)$ is $j$ or $j+1$,
according to its $\Omega_{j+1}$-bit, and this bit is its bit in $\mathbf{\Omega}$. By Fact~(F), the
level of a $\pi_j$-block $b\subset\square^{(2)}$ in $\mathbf{\Omega}$ falls into one of four cases.
\begin{itemize}
\item It is $<j$ if $[b\subset\Omega_j]=0$.
\item It is $j$ if $[b\subset\Omega_j]=1$ and $[b\subset\Omega_{j+1}]=0$.
\item It is $j+1$ if $[b\subset\Omega_{j+1}]=1$ and $[b\subset\Omega_{j+2}]=0$.
\item It is $\ge j+2$ if $[b\subset\Omega_{j+2}]=1$.
\end{itemize}
In the middle two cases the local and global levels coincide. In the first case the
$\Omega_{j+1}$-bit is $0$, so the local level is $j$ while the global level is $<j$. In the last
case the $\Omega_{j+1}$-bit is $1$, so the local level is $j+1$ while the global level is $\ge j+2$.
Agreement therefore fails exactly at the blocks of global level $<j$ or $\ge j+2$, which proves the
equivalence.

For the last assertion of (a), suppose $\square^{(2)}\subset\Omega_j\setminus\Omega_{j+2}$. Then
every block $b\subset\square^{(2)}$ lies in $\Omega_j$ and is disjoint from $\Omega_{j+2}$. Hence
$[b\subset\Omega_j]=1$ and $[b\subset\Omega_{j+2}]=0$.

\emph{(b).} If $\Omega_{j+2}=\emptyset$ there is nothing to prove, so assume otherwise.

We first place $e$. The block $e$ is a $\pi_j$-block with $e\not\subset\Omega_{j+1}$, and
$\Omega_{j+1}$ is a union of $\pi_{j+1}$-blocks. By Fact~(F),
$\operatorname{int}e\cap\Omega_{j+1}=\emptyset$. Hence
\begin{equation*}
e=\operatorname{cl}\operatorname{int}e\subset\operatorname{cl}(\mathbb{T}\setminus\Omega_{j+1}).
\end{equation*}

Next we bound the distance between the two grained sets. By Fact~(F), the set
$\operatorname{cl}(\mathbb{T}\setminus\Omega_{j+1})$ is a finite union of $\pi_{j+1}$-blocks. The set
$\Omega_{j+2}$ is $\pi_{j+2}$-grained, hence also $\pi_{j+1}$-grained. The distance between two finite
unions of closed blocks is the minimum of the distances between pairs of blocks. For two closed
$\pi_{j+1}$-blocks, the $\ell^\infty$-distance is the maximum over coordinates of the gap between two
grid intervals of length $u_{j+1}$ on a circle whose length is a multiple of $u_{j+1}$ (the blocks of
$\pi_{j+1}$, of which the members of a nest in Definition~\ref{4.6:def-nest-traces} are unions, tile
$\mathbb{T}$). Each such gap lies in $u_{j+1}\mathbb{Z}_{\ge0}$. Hence
\begin{equation*}
\operatorname{dist}\bigl(\operatorname{cl}(\mathbb{T}\setminus\Omega_{j+1}),\Omega_{j+2}\bigr)
\in u_{j+1}\mathbb{Z}_{\ge0}.
\end{equation*}
The distance from a set equals the distance from its closure. So this number equals the left side of
\eqref{4.6:eq-inner-seam-1}, and it exceeds $Ru_{j+1}$. Since $R$ is an integer, it is therefore at
least $(R+1)u_{j+1}$. Because $e\subset\operatorname{cl}(\mathbb{T}\setminus\Omega_{j+1})$ and
$u_{j+1}=Lu_j$, we obtain
\begin{equation}\label{4.6:eq-inner-seam-3}
\operatorname{dist}(e,\Omega_{j+2})\ge
\operatorname{dist}\bigl(\operatorname{cl}(\mathbb{T}\setminus\Omega_{j+1}),\Omega_{j+2}\bigr)
\ge(R+1)u_{j+1}=(R+1)Lu_j .
\end{equation}

Now we relate $\square^{(2)}$ to $e$. The cube is $\square=x+[-u_j,u_j]^d$, and $e$ is one of its
$2^d$ blocks. So in each coordinate $e$ is $[0,u_j]$ or $[-u_j,0]$ shifted by $x$. Every point of
$[-u_j,u_j]$ lies within $u_j$ of either interval, so $\square\subset e^{+u_j}$. Consequently
\begin{equation*}
\square^{(2)}=\square^{+2u_j}\subset e^{+3u_j}.
\end{equation*}
By the triangle inequality and \eqref{4.6:eq-inner-seam-3},
\begin{equation*}
\operatorname{dist}(\square^{(2)},\Omega_{j+2})
\ge\operatorname{dist}(e,\Omega_{j+2})-3u_j
\ge\bigl((R+1)L-3\bigr)u_j>0 ,
\end{equation*}
where the last inequality is \eqref{4.6:eq-inner-seam-2}. Thus $\square^{(2)}\cap\Omega_{j+2}=\emptyset$.

We add three comments on the hypotheses of (b). First, condition \eqref{4.6:eq-inner-seam-2} holds for
every integer $R\ge1$ as soon as $L\ge2$; in \cite[p.~224]{B-CMP96}, ``$R$ is a big positive integer''.
Second, if the integrality of the distance is not used, \eqref{4.6:eq-inner-seam-1} alone gives
$\operatorname{dist}(e,\Omega_{j+2})>RLu_j$. The same conclusion then follows under $RL>3$, that is,
$RL\ge4$. Third, under the standing choice of \cite{B-CMP99b}, ``$R$, $M$ are sufficiently large. We
take these numbers as powers of $L$'', both forms of the condition hold.

\emph{(b$'$).} Write $u=u_{j+2}$. We first show that $\Omega_{j+2}^{+u}\subset\Omega_{j+1}$. Let $p$
lie at distance at most $u$ from a $\pi_{j+2}$-block $b\subset\Omega_{j+2}$. In each coordinate,
$p_i$ lies within $u$ of the grid interval $b_i$. Hence $p_i$ lies in one of the three grid intervals
of length $u$ that meet $b_i$: the interval $b_i$ itself or one of its two neighbours. Choosing such
an interval in every coordinate gives a $\pi_{j+2}$-block that contains $p$ and meets $b$. By
hypothesis this block lies in $\Omega_{j+1}$, so $p\in\Omega_{j+1}$.

Consequently the open set $\{p:\operatorname{dist}(p,\Omega_{j+2})<u\}$ lies in
$\operatorname{int}\Omega_{j+1}$. It is therefore disjoint from
$\operatorname{cl}(\mathbb{T}\setminus\Omega_{j+1})=\mathbb{T}\setminus\operatorname{int}\Omega_{j+1}$.
This gives
\begin{equation*}
\operatorname{dist}\bigl(\operatorname{cl}(\mathbb{T}\setminus\Omega_{j+1}),\Omega_{j+2}\bigr)
\ge u_{j+2}=L^2u_j\ge4u_j>3u_j .
\end{equation*}
The argument of (b) from the placement of $e$ onward then applies with $(R+1)Lu_j$ replaced by
$L^2u_j$, and yields $\operatorname{dist}(\square^{(2)},\Omega_{j+2})>0$.

For the nests used in the construction, the hypothesis of (b$'$) is the boundary separation
\eqref{4.6:eq-nest-traces-a} at level $j+2$: under that separation, in its form without the cell
multipliers $\check{R}_i$, every $\pi_{j+2}$-block meeting $\Omega_{j+2}$ lies in $\Omega_{j+1}$.

\emph{(c).} Let $z\in\partial\Omega_j$ with $\operatorname{dist}(x,z)\le2u_j$. The open cube
$x+(-3u_j,3u_j)^d$ is a neighbourhood of $z$ contained in $\square^{(2)}$. Because $z$ is a boundary
point of $\Omega_j$, this neighbourhood meets $\mathbb{T}\setminus\Omega_j$. Hence
$\square^{(2)}\not\subset\Omega_j$. Since $x$ is a vertex of the grid $\pi_j$, the set $\square^{(2)}$
is a union of $\pi_j$-blocks. At least one of these blocks $b$ has $[b\subset\Omega_j]=0$, so
$(\mathrm{LA}_{\square})$ fails by (a).

The cover of \cite[p.~229]{B-CMP96} admits the cubes of part~(c). Its cubes are built from big blocks
``with a center $y\in\Lambda_j$ (more exactly it belongs to the boundary of this set also)'', where
$\Lambda_j$ is the set of \cite[(2.3)]{B-CMP96}. At such a cube the hypothesis $(\mathrm{LA}_{\square})$
of the exactness clause of Lemma~\ref{4.6:lem-commutator-identity} is not available. The clause of
that lemma that reproduces \cite[(3.105)]{B-CMP99b} then requires that the functions $h_{\square}$ of
\cite[p.~229]{B-CMP96}, satisfying \cite[(2.36)]{B-CMP96}, vanish where the levels of
$\mathbf{\Omega}$ and $\mathbf{\Omega}(\square)$ differ. Otherwise its identity acquires the additional
term
\begin{equation}\label{4.6:eq-inner-seam-4}
\begin{split}
-\sum_{\square}h_{\square}\bigl(\mathfrak{L}-\mathfrak{L}_{\square}\bigr)G_{\square}h_{\square}
&=-\sum_{\square}h_{\square}\bigl(Q^*\mathfrak{a}Q[\mathbf{\Omega}]
-Q^*\mathfrak{a}Q[\mathbf{\Omega}(\square)]\bigr)G_{\square}h_{\square}.
\end{split}
\end{equation}
The equality in \eqref{4.6:eq-inner-seam-4} holds because, in
$\mathfrak{L}=\Delta(U)+DD^*+Q^*\mathfrak{a}Q$, only the last summand depends on the nest. The
additional term is local: it is a local letter of the alphabet $\mathfrak{A}$ of this chapter, and it
reads the nests only through their traces on $\square^{(2)}$, so the results of this chapter that rest
on that locality are unaffected whether or not the functions $h_{\square}$ vanish there. Neither
part~(c) nor level agreement at the cubes it describes is used in what follows.
\end{proof}

\begin{lemma}[rescaled profiles double-count at every seam]\label{4.6:lem-seam-double-count}
Let $\mathbf{\Omega}=(\Omega_0\supseteq\Omega_1\supseteq\dots\supseteq\Omega_k)$ be a nest in the
sense of Definition~\ref{4.6:def-nest-traces}. Thus each $\Omega_i$ is a closed union of
$\pi_i$-blocks, the $\Omega_i$ decrease, and $\Omega_{k+1}=\Omega_{k+2}=\emptyset$. Let
$\mathcal{C}_g$ be the family of candidate cubes $Q=x+[-u_g,u_g]^d$, $x\in u_g\mathbb{Z}^d$, of
scale $g$, and let $\mathcal{D}_g(\mathbf{\Omega})\subset\mathcal{C}_g$ be the members of scale
$g$. Let $\mathcal{D}(\mathbf{\Omega})$ be the mixed-size cover, the union over $g$ of the
families $\mathcal{D}_g(\mathbf{\Omega})$. All three are as in
Definition~\ref{4.6:def-local-sequence}. Let $B_g=\Omega_g\setminus\operatorname{int}\Omega_{g+1}$
be the layers. For $Q\in\mathcal{C}_g$ let $\eta_Q$ be the profile \eqref{4.6:eq-bump-profiles-a}
of Definition~\ref{4.6:def-bump-profiles}.
\begin{itemize}
\item[(a)] For each $g$, $\sum_{Q\in\mathcal{C}_g}\eta_Q^2\equiv 1$ on $\mathbb{T}$.
\item[(b)] Assume that $\mathbf{\Omega}$ satisfies the boundary separation condition
\eqref{4.6:eq-nest-traces-a} at level $g$. Then the following hold.
\begin{itemize}
\item For every $y\in B_g$, every $Q\in\mathcal{C}_g$ with $\eta_Q(y)\neq0$ belongs to
$\mathcal{D}_g(\mathbf{\Omega})$, and hence $\sum_{Q\in\mathcal{D}_g(\mathbf{\Omega})}\eta_Q^2(y)=1$.
\item The function $\sum_{Q\in\mathcal{D}_g(\mathbf{\Omega})}\eta_Q^2$ vanishes off
$B_g^{+(2/3)u_g}$.
\end{itemize}
Without the separation condition the first conclusion of (b) fails. Take $d=1$ and $u=u_g$, with
$\Omega_g=\Omega_{g+1}=[0,A]$ for some $A\in Lu\mathbb{N}$ with $A\ge Lu$, and
$\Omega_{g+2}=\emptyset$. Then the
point $y=0$ lies in $B_g$, but $\sum_{Q\in\mathcal{D}_g(\mathbf{\Omega})}\eta_Q^2(0)=0$.
\item[(c)] Assume \eqref{4.6:eq-nest-traces-a} at levels $g$ and $g+1$. Let $y\in\partial\Omega_{g+1}$,
a point of the seam between the layers $B_g$ and $B_{g+1}$. Then $y\in B_g\cap B_{g+1}$ and
\begin{equation*}
\sum_{Q\in\mathcal{D}(\mathbf{\Omega})}\eta_Q^2(y)\;\ge\;
\sum_{Q\in\mathcal{D}_g(\mathbf{\Omega})}\eta_Q^2(y)+\sum_{Q\in\mathcal{D}_{g+1}(\mathbf{\Omega})}\eta_Q^2(y)
\;=\;2 .
\end{equation*}
In particular, at no point of any seam does the family
$\{\eta_Q\}_{Q\in\mathcal{D}(\mathbf{\Omega})}$ satisfy the identity
$\sum_{\square\in\mathcal{D}}h_\square^2=1$ of \cite[(2.36)]{B-CMP96}.
\item[(d)] (An explicit two-layer model.) Let $d=1$ and $u:=u_g$, so that $u_{g+1}=Lu$. Let
$\mathbb{T}$ be a circle whose period is a large element of $L^2u\mathbb{N}$, or let
$\mathbb{T}=\mathbb{R}$. Put
\begin{equation*}
\Omega_i:=\mathbb{T}\ (i\le g),\qquad \Omega_{g+1}:=[0,A],\qquad \Omega_{g+2}:=\emptyset,
\end{equation*}
with $A\in Lu\mathbb{N}$ and $A\ge 2Lu$. Then the following hold.
\begin{itemize}
\item $\sum_{Q\in\mathcal{D}(\mathbf{\Omega})}\eta_Q^2(y)=2$ for every $y\in[-Lu/3,\,u/3]$, the
collar of the seam at $0$.
\item $\sum_{Q\in\mathcal{D}(\mathbf{\Omega})}\eta_Q^2(y)=1$ for $y\le-2Lu/3$.
\item $\sum_{Q\in\mathcal{D}(\mathbf{\Omega})}\eta_Q^2(y)=1$ for $(2/3)u\le y\le A-(2/3)u$.
\end{itemize}
\item[(e)] (Sub-selections, in the model of (d).) Let $\mathcal{D}'\subset\mathcal{D}(\mathbf{\Omega})$
be a subfamily. Call $[-u,u]\in\mathcal{D}_g(\mathbf{\Omega})$
and $[-Lu,Lu]\in\mathcal{D}_{g+1}(\mathbf{\Omega})$ the two seam-centred cubes; their centres are
$x=0$ and $x'=0$. Then the following hold.
\begin{itemize}
\item If $\mathcal{D}'$ contains both seam-centred cubes, then
$\sum_{Q\in\mathcal{D}'}\eta_Q^2(0)\ge2$.
\item If $\mathcal{D}'$ does not contain the member $[-Lu,Lu]$ of scale $g+1$, then
$\sum_{Q\in\mathcal{D}'}\eta_Q^2(2u/3)=0$.
\item If $\mathcal{D}'$ contains neither seam-centred cube, then $\sum_{Q\in\mathcal{D}'}\eta_Q^2=0$
on $\left]0,u/3\right[$.
\end{itemize}
Whether some subfamily $\mathcal{D}'$ satisfies $\sum_{Q\in\mathcal{D}'}\eta_Q^2\equiv1$ near the
seam is not decided by (e).
\end{itemize}
Consequently the family of all rescaled profiles $\eta_Q$, $Q\in\mathcal{D}(\mathbf{\Omega})$, does
not satisfy the identity $\sum_{\square\in\mathcal{D}}h_\square^2=1$ of \cite[(2.36)]{B-CMP96}. In
\cite[p.~229]{B-CMP96} the functions $h_\square$ of \cite[(2.36)]{B-CMP96} are introduced by the
words ``construct also the corresponding family of functions $h$ described in (1.118), and rescale
them to proper scales'', the functions $h$ being those of \cite[(1.118)]{B-CMP95}. In this book the
functions $h_\square$ on $\mathcal{D}(\mathbf{\Omega})$ for which \cite[(2.36)]{B-CMP96} is used are
those of the gluing hypothesis of Definition~\ref{4.6:def-bump-profiles}, not the plain profiles
$\eta_Q$.
\end{lemma}

\begin{proof}
Throughout, $h\in C_0^\infty(\left]-2/3,2/3\right[)$ is the one-variable function of
\cite[(1.118)]{B-CMP95} entering \eqref{4.6:eq-bump-profiles-a}. It satisfies $h=1$ on
$[-1/3,1/3]$ and $\sum_{n\in\mathbb{Z}}h^2(t-n)=1$. Since the support of $h$ is a compact subset of
the open interval $\left]-2/3,2/3\right[$, we have $h(t)=0$ whenever $|t|\ge2/3$.

\emph{(a)} Fix $g$ and $y\in\mathbb{T}$. By \eqref{4.6:eq-bump-profiles-a},
$\eta_{x+[-u_g,u_g]^d}(y)=\prod_{\mu=1}^d h((y_\mu-x_\mu)/u_g)$. For fixed $y$ only finitely many
centres give a non-zero term. Summing over $x\in u_g\mathbb{Z}^d$ and writing
$x_\mu=u_gn_\mu$, the sum of products factorises coordinate by coordinate:
\begin{equation*}
\sum_{x\in u_g\mathbb{Z}^d}\prod_{\mu=1}^d h\Bigl(\frac{y_\mu-x_\mu}{u_g}\Bigr)^2
=\prod_{\mu=1}^d\sum_{n\in\mathbb{Z}}h\Bigl(\frac{y_\mu}{u_g}-n\Bigr)^2=1 .
\end{equation*}
The last equality holds because each factor equals $1$ by the identity $\sum_nh^2(t-n)=1$ of
\cite[(1.118)]{B-CMP95}.

\emph{(b)} Write $u:=u_g$. Let $y\in B_g$ and let $Q=x+[-u,u]^d\in\mathcal{C}_g$ satisfy
$\eta_Q(y)\neq0$. Every factor $h((y_\mu-x_\mu)/u)$ is then non-zero, so $|y_\mu-x_\mu|<(2/3)u$
for all $\mu$.

\emph{Every $\pi_g$-block containing $y$ lies in $Q$.} The $\pi_g$-blocks contained in $Q$ are the
products $\prod_\mu I_\mu$ with $I_\mu\in\{[x_\mu-u,x_\mu],[x_\mu,x_\mu+u]\}$. A $\pi_g$-block
containing $y$ is a product of intervals of the grid $u\mathbb{Z}$, each containing $y_\mu$. Since
$x_\mu\in u\mathbb{Z}$ and $|y_\mu-x_\mu|<(2/3)u<u$, the grid intervals containing $y_\mu$ are
among the two intervals $[x_\mu-u,x_\mu]$ and $[x_\mu,x_\mu+u]$. There is one such interval if
$y_\mu\neq x_\mu$, and there are two if $y_\mu=x_\mu$. Hence every $\pi_g$-block containing $y$ is
one of the blocks of $Q$.

\emph{(b1)} Since $y\in\Omega_g$ and $\Omega_g$ is a union of $\pi_g$-blocks, some $\pi_g$-block
$e_0\ni y$ satisfies $e_0\subset\Omega_g$.

\emph{(b2)} Since $y\in B_g$, we have $y\notin\operatorname{int}\Omega_{g+1}$. The union of the
$\pi_g$-blocks containing $y$ is a neighbourhood of $y$. Indeed, the blocks of $\pi_g$ cover
$\mathbb{T}$, only finitely many of them meet a given small ball around $y$, and those among them
not containing $y$ are closed and at positive distance from $y$. Moreover $\Omega_{g+1}$ is a union
of $\pi_{g+1}$-blocks, hence of $\pi_g$-blocks. Were every $\pi_g$-block containing $y$ contained
in $\Omega_{g+1}$, the point $y$ would be interior to $\Omega_{g+1}$. Therefore some $\pi_g$-block
$e_1\ni y$ satisfies $e_1\not\subset\Omega_{g+1}$.

\emph{(b3)} The point $y$ lies in the closed set $\Omega_g$, so either
$y\in\operatorname{int}\Omega_g$ or $y\in\partial\Omega_g$. We treat the two cases in turn.

\emph{First case: $y\in\operatorname{int}\Omega_g$.} We claim that every $\pi_g$-block $e\ni y$ is
contained in $\Omega_g$. The set $\operatorname{int}\Omega_g$ is an open neighbourhood of $y$, and
$y\in e=\operatorname{cl}\operatorname{int}e$. Hence $\operatorname{int}e$ meets $\Omega_g$ at some
point $z$. This $z$ lies in some $\pi_g$-block $e'\subset\Omega_g$. Now
$z\in\operatorname{int}e\cap e'$, and distinct grid cubes meet only in their boundaries, so
$e'=e$. Thus $e\subset\Omega_g$. Take $e:=e_1$ from (b2). Then $e\subset Q$ by the paragraph above,
$e\subset\Omega_g$, and $e\not\subset\Omega_{g+1}$.

\emph{Second case: $y\in\partial\Omega_g$.} Here $y\in\operatorname{cl}(\mathbb{T}\setminus\Omega_g)$.
Since $\Omega_{g+1}$ is closed, the separation condition \eqref{4.6:eq-nest-traces-a} at level $g$
gives $\Omega_{g+1}\cap\operatorname{cl}(\mathbb{T}\setminus\Omega_g)=\emptyset$, hence
$y\notin\Omega_{g+1}$. Every $\pi_g$-block containing $y$ therefore
fails to be contained in $\Omega_{g+1}$. Take $e:=e_0$ from (b1). Then $e\subset Q$,
$e\subset\Omega_g$, and $e\not\subset\Omega_{g+1}$.

In either case $Q$ contains a $\pi_g$-block $e$ with $[e\subset\Omega_g]=1$ and
$[e\subset\Omega_{g+1}]=0$. By the definition of membership in
Definition~\ref{4.6:def-local-sequence}, $Q\in\mathcal{D}_g(\mathbf{\Omega})$. So every term of
the sum in (a) that is non-zero at $y$ belongs to $\mathcal{D}_g(\mathbf{\Omega})$. The remaining
terms vanish at $y$. Hence
\begin{equation*}
\sum_{Q\in\mathcal{D}_g(\mathbf{\Omega})}\eta_Q^2(y)=\sum_{Q\in\mathcal{C}_g}\eta_Q^2(y)=1
\end{equation*}
by (a).

\emph{Support.} Let $Q=x+[-u,u]^d\in\mathcal{D}_g(\mathbf{\Omega})$, with witnessing block $e$:
a $\pi_g$-block $e\subset Q$ with $e\subset\Omega_g$ and $e\not\subset\Omega_{g+1}$. We first show
$e\subset B_g$. Suppose $\operatorname{int}e$ met $\Omega_{g+1}$ at a point $z$. Then $z$ would lie
in a $\pi_g$-block $e'\subset\Omega_{g+1}$, and $z\in\operatorname{int}e\cap e'$ would force
$e'=e\subset\Omega_{g+1}$, a contradiction. So $\operatorname{int}e\cap\Omega_{g+1}=\emptyset$.
The open set $\operatorname{int}\Omega_{g+1}$ is thus disjoint from $\operatorname{int}e$, hence
from $\operatorname{cl}\operatorname{int}e=e$. Together with $e\subset\Omega_g$ this gives
$e\subset B_g$.

As a $\pi_g$-block contained in $Q$, the block $e$ has the centre $x$ as a corner. Since $h$
vanishes outside $\left]-2/3,2/3\right[$,
\begin{equation*}
\operatorname{supp}\eta_Q\subset x+\left]-\tfrac23u,\tfrac23u\right[^d\subset e^{+(2/3)u}
\subset B_g^{+(2/3)u}.
\end{equation*}
Every term of $\sum_{Q\in\mathcal{D}_g(\mathbf{\Omega})}\eta_Q^2$ therefore vanishes off
$B_g^{+(2/3)u_g}$, and so does the sum.

\emph{The example without the separation condition.} Here $d=1$, $u=u_g$,
$\Omega_g=\Omega_{g+1}=[0,A]$ and $\Omega_{g+2}=\emptyset$. Then
$B_g=[0,A]\setminus\left]0,A\right[$ contains $y=0$. A candidate $x+[-u,u]\in\mathcal{C}_g$ has
$\eta(0)=h(-x/u)$, which is non-zero only if $|x|<(2/3)u$. Since $x\in u\mathbb{Z}$, this forces
$x=0$. So the only candidate of scale $g$ whose profile is non-zero at $0$ is $[-u,u]$. Its
$\pi_g$-blocks are $[-u,0]$ and $[0,u]$. The first satisfies $[-u,0]\not\subset\Omega_g$. The
second satisfies $[0,u]\subset[0,A]=\Omega_{g+1}$, because $A\ge Lu\ge u$. Hence neither block is
at the same time contained in $\Omega_g$ and not contained in $\Omega_{g+1}$, so $[-u,u]$ has no
witnessing block and is not a member. Every
member of $\mathcal{D}_g(\mathbf{\Omega})$ therefore has a profile vanishing at $0$, and
$\sum_{Q\in\mathcal{D}_g(\mathbf{\Omega})}\eta_Q^2(0)=0$.

\emph{(c)} Let $y\in\partial\Omega_{g+1}$.

\emph{$y\in B_g$.} The set $\Omega_{g+1}$ is closed, so $y\in\Omega_{g+1}\subset\Omega_g$. Also
$y\notin\operatorname{int}\Omega_{g+1}$. Hence $y\in B_g$.

\emph{$y\in B_{g+1}$.} Since $y\in\operatorname{cl}(\mathbb{T}\setminus\Omega_{g+1})$, the
separation condition \eqref{4.6:eq-nest-traces-a} at level $g+1$ gives
$y\notin\Omega_{g+2}\supseteq\operatorname{int}\Omega_{g+2}$. This is trivially true if
$\Omega_{g+2}=\emptyset$. Together with $y\in\Omega_{g+1}$ this gives $y\in B_{g+1}$.

Part (b) at level $g$ (using \eqref{4.6:eq-nest-traces-a} at level $g$) gives
$\sum_{Q\in\mathcal{D}_g(\mathbf{\Omega})}\eta_Q^2(y)=1$. Part (b) at level $g+1$ (using
\eqref{4.6:eq-nest-traces-a} at level $g+1$) gives
$\sum_{Q\in\mathcal{D}_{g+1}(\mathbf{\Omega})}\eta_Q^2(y)=1$. The families
$\mathcal{D}_g(\mathbf{\Omega})$ and $\mathcal{D}_{g+1}(\mathbf{\Omega})$ consist of cubes of the
different sides $2u_g$ and $2u_{g+1}$, so they are disjoint subfamilies of
$\mathcal{D}(\mathbf{\Omega})$. All summands $\eta_Q^2(y)$ are non-negative. The displayed
inequality follows, and with it the statement about \cite[(2.36)]{B-CMP96}.

\emph{(d)} Write $Q=x+[-u,u]$ with $x\in u\mathbb{Z}$ for the candidates of scale $g$, and
$Q=x'+[-Lu,Lu]$ with $x'\in Lu\mathbb{Z}$ for those of scale $g+1$. Their profiles are
$h((y-x)/u)$ and $h((y-x')/(Lu))$ respectively.

\emph{Members of scale $g$.} Since $\Omega_g=\mathbb{T}$, every block is contained in $\Omega_g$.
A candidate $x+[-u,u]$ is therefore a member exactly when $[x-u,x]\not\subset[0,A]$ or
$[x,x+u]\not\subset[0,A]$. Both blocks lie in $[0,A]$ exactly when $u\le x\le A-u$. So the members
are the centres $x\le0$ and $x\ge A$.

\emph{The scale-$g$ sum on $\left]-\infty,u/3\right]$.} Let $y\le u/3$. A centre $x\ge u$ has
$(y-x)/u\le-2/3$, so its profile $h((y-x)/u)$ vanishes at $y$. Hence every contributor to the full
scale-$g$ sum that is non-zero at $y$ has $x\le0$ and is a member. By (a) with $d=1$,
\begin{equation}\label{4.6:eq-seam-double-count-1}
\sum_{Q\in\mathcal{D}_g(\mathbf{\Omega})}\eta_Q^2(y)=\sum_{x\in u\mathbb{Z}}h^2\Bigl(\frac{y-x}{u}\Bigr)=1
\qquad\text{for } y\le u/3 .
\end{equation}

\emph{Members of scale $g+1$.} Since $\Omega_{g+2}=\emptyset$, no block is contained in
$\Omega_{g+2}$. A candidate $x'+[-Lu,Lu]$ is therefore a member exactly when
$[x'-Lu,x']\subset[0,A]$ or $[x',x'+Lu]\subset[0,A]$. The second condition holds exactly when
$0\le x'\le A-Lu$, and the first exactly when $Lu\le x'\le A$. Since $A\ge Lu$ and $A\in Lu\mathbb{Z}$,
the union of the two conditions over $x'\in Lu\mathbb{Z}$ is $0\le x'\le A$. The non-members are
thus the centres $x'\le-Lu$ and $x'\ge A+Lu$.

\emph{The scale-$(g+1)$ sum on $[-Lu/3,u/3]$.} Let $-Lu/3\le y\le u/3$.
\begin{itemize}
\item A centre $x'\le-Lu$ has $(y-x')/(Lu)\ge-1/3+1=2/3$, so it contributes $0$.
\item Since $A\ge2Lu$, we have $u/3\le(4/3)Lu\le A-(2/3)Lu$, so $y\le A-(2/3)Lu$. A centre
$x'\ge A+Lu$ then has $y-x'\le-(2/3)Lu-Lu$, so it contributes $0$.
\end{itemize}
Every contributor to the full scale-$(g+1)$ sum that is non-zero at $y$ therefore has
$0\le x'\le A$ and is a member. By (a),
\begin{equation}\label{4.6:eq-seam-double-count-2}
\sum_{Q\in\mathcal{D}_{g+1}(\mathbf{\Omega})}\eta_Q^2(y)=1\qquad\text{for } -Lu/3\le y\le u/3 .
\end{equation}

\emph{Other scales.} For a scale $i\le g-1$ we have $\Omega_{i+1}=\mathbb{T}$, so every block is
contained in $\Omega_{i+1}$ and there are no members. For $i\ge g+2$ we have
$\Omega_i\subseteq\Omega_{g+2}=\emptyset$, so no $\pi_i$-block is contained in $\Omega_i$ and there
are no members of scale $i$.
Hence $\mathcal{D}(\mathbf{\Omega})=\mathcal{D}_g(\mathbf{\Omega})\cup\mathcal{D}_{g+1}(\mathbf{\Omega})$.

\emph{The collar.} By \eqref{4.6:eq-seam-double-count-1} and \eqref{4.6:eq-seam-double-count-2},
$\sum_{Q\in\mathcal{D}(\mathbf{\Omega})}\eta_Q^2=1+1=2$ on $[-Lu/3,u/3]$.

\emph{The region $y\le-2Lu/3$.} The members of scale $g+1$ have $x'\ge0$, so
$(y-x')/(Lu)\le-2/3$, and they contribute $0$. The scale-$g$ sum is $1$ by
\eqref{4.6:eq-seam-double-count-1}. Hence the total is $1$.

\emph{The region $(2/3)u\le y\le A-(2/3)u$.} This interval is non-empty since $A\ge2Lu$.
\begin{itemize}
\item A member of scale $g$ with $x\le0$ has $(y-x)/u\ge2/3$.
\item A member of scale $g$ with $x\ge A$ has $(y-x)/u\le-2/3$.
\end{itemize}
So the members of scale $g$ contribute $0$. A contributor $x'$ to the full scale-$(g+1)$ sum that is
non-zero at $y$ has $|y-x'|<(2/3)Lu$.
\begin{itemize}
\item Lower bound:
\begin{equation*}
x'>y-\tfrac23Lu\ \ge\ \tfrac23u-\tfrac23Lu\ >\ -\tfrac23Lu .
\end{equation*}
Since $x'\in Lu\mathbb{Z}$, this gives $x'\ge0$.
\item Upper bound:
\begin{equation*}
x'<y+\tfrac23Lu\ \le\ A-\tfrac23u+\tfrac23Lu\ <\ A+Lu .
\end{equation*}
Since $x',A\in Lu\mathbb{Z}$, this gives $x'\le A$.
\end{itemize}
Such $x'$ are therefore members, and by (a) the scale-$(g+1)$ member sum is $1$. Hence the total is
$1$.

\emph{(e)} By the determination of the members in (d), the cube $[-u,u]$ (centre $x=0\le0$) belongs
to $\mathcal{D}_g(\mathbf{\Omega})$. The cube $[-Lu,Lu]$ has centre $x'=0$, which lies in the
member range $0\le x'\le A$, so it belongs to
$\mathcal{D}_{g+1}(\mathbf{\Omega})$. By (d), every member of $\mathcal{D}(\mathbf{\Omega})$ has
scale $g$ or $g+1$.

\emph{Both seam-centred cubes kept.} At $y=0$ these two cubes alone contribute
$h(0)^2+h(0)^2=2$, because $h=1$ on $[-1/3,1/3]$. All other summands are non-negative, so
$\sum_{Q\in\mathcal{D}'}\eta_Q^2(0)\ge2$.

\emph{Both seam-centred cubes dropped.} Let $y\in\left]0,u/3\right[$. We check every other
candidate of the two scales.
\begin{itemize}
\item Scale $g$, $x=u$: the argument $y/u-1$ lies in $\left]-1,-2/3\right[$, so the profile
vanishes at $y$.
\item Scale $g$, $x=-u$: the argument $y/u+1$ lies in $\left]1,4/3\right[$, so the profile vanishes
at $y$.
\item Scale $g+1$, $x'=Lu$: the argument $(y-Lu)/(Lu)$ lies in
$\left]-1,-1+1/(3L)\right[\subset\left]-1,-2/3\right[$, so the profile vanishes at $y$.
\item Scale $g+1$, $x'=-Lu$: the argument $(y+Lu)/(Lu)$ exceeds $1$, so the profile vanishes at $y$.
\item A candidate of scale $g$ with $|x|\ge2u$ has $|y-x|/u>2-1/3>2/3$, and a candidate of scale
$g+1$ with $|x'|\ge2Lu$ has $|y-x'|/(Lu)>2-1/(3L)>2/3$. Their profiles vanish at $y$.
\end{itemize}
The only candidates of either scale whose profiles can be non-zero on $\left]0,u/3\right[$ are thus
the two seam-centred cubes. If $\mathcal{D}'$ contains neither, then
$\sum_{Q\in\mathcal{D}'}\eta_Q^2=0$ on $\left]0,u/3\right[$.

\emph{The member $[-Lu,Lu]$ dropped.} Let $y=2u/3$. By (d), the members of
$\mathcal{D}(\mathbf{\Omega})$ other than $[-Lu,Lu]$ are the cubes of scale $g$ with centres
$x\le0$ or $x\ge A$, and the cubes of scale $g+1$ with centres $x'\in Lu\mathbb{Z}$,
$Lu\le x'\le A$. We check each of them.
\begin{itemize}
\item Scale $g$, $x\le0$: the argument $(y-x)/u=2/3-x/u$ is at least $2/3$, so the profile
vanishes at $y$.
\item Scale $g$, $x\ge A$: since $A\ge2Lu$, the argument satisfies
\begin{equation*}
\frac{y-x}{u}\ \le\ \frac23-\frac{A}{u}\ \le\ \frac23-2L\ <\ -\frac23 ,
\end{equation*}
so the profile vanishes at $y$.
\item Scale $g+1$, $Lu\le x'\le A$: the argument $(y-x')/(Lu)$ is at most $2/(3L)-1$. The blocking
factor $L$ of Definition~\ref{1.1:def-lattice} is an odd integer greater than $1$, so $L\ge3$ and
$2/(3L)-1<-2/3$; hence the profile vanishes at $y$.
\end{itemize}
Hence, if $\mathcal{D}'\subset\mathcal{D}(\mathbf{\Omega})$ does not contain $[-Lu,Lu]$, every
summand of $\sum_{Q\in\mathcal{D}'}\eta_Q^2(2u/3)$ vanishes, and
$\sum_{Q\in\mathcal{D}'}\eta_Q^2(2u/3)=0$.
\end{proof}

\begin{remark}[The partition of unity on the mixed-size cover]\label{4.6:rem-mixed-cover-unity}
Let $\eta_Q$ be the one-scale profiles of Definition~\ref{4.6:def-bump-profiles}, read with the
scale convention \eqref{4.6:eq-bump-profiles-b}.
\begin{enumerate}
\item[(i)] Let $d=1$, and let the nest $\mathbf{\Omega}$ satisfy the separation condition
of Lemma~\ref{4.6:lem-seam-double-count}, in the two-layer configuration of
clause (d) of that lemma: one seam, placed at $y=0$, between the layers $B_g$, on the side
$y\le0$, and $B_{g+1}$, on the side $y\ge0$. Near the seam, the squared profiles of the members of
$\mathcal{D}_g(\mathbf{\Omega})$ sum to $1$ for $y\le u_g/3$, and those of the members of
$\mathcal{D}_{g+1}(\mathbf{\Omega})$ sum to $1$ for $y\ge-Lu_g/3$. On $[-Lu_g/3,u_g/3]$ the sum of
$\eta_Q^2$ over the mixed-size cover $\mathcal{D}(\mathbf{\Omega})$ equals $2$. This is clause (d)
of Lemma~\ref{4.6:lem-seam-double-count}.
\item[(ii)] Consequently the family $\{h_\square\}_{\square\in\mathcal{D}}$ of \cite[(2.36)]{B-CMP96},
which satisfies $\sum_{\square\in\mathcal{D}}h_\square^2=1$, is not the family of the products
of \cite[(1.118)]{B-CMP95} rescaled scale by scale without change. It includes a normalisation at
the seams, or a convention of equivalent effect. In what follows the gluing hypothesis
$(\mathrm{G\text{-}loc})$, $(\mathrm{G\text{-}supp})$ of \eqref{4.6:eq-bump-profiles-c} is imposed
on that normalisation as a hypothesis; it is consistent with every constraint stated in
\cite{B-CMP96} and \cite{B-CMP95} on the cover and its partition of unity, and it is not identified
with the construction of \cite{B-CMP96}. The identity \cite[(2.36)]{B-CMP96} itself enters only the
convergence statement for a single nest.
\item[(iii)] No statement of this chapter comparing two nests depends on which rule satisfying
the gluing hypothesis $(\mathrm{G\text{-}loc})$, $(\mathrm{G\text{-}supp})$ is used.
\end{enumerate}
\end{remark}

\begin{proof}
\emph{(i).} Here $d=1$. Place the seam at $0$, with $B_g$ on the side $y\le0$ and $B_{g+1}$ on the
side $y\ge0$ near $0$. Put $u=u_g$, so that $u_{g+1}=Lu$.

As in clause (d) of Lemma~\ref{4.6:lem-seam-double-count}, under its separation condition,
the members of $\mathcal{D}_g(\mathbf{\Omega})$ near $0$ are the cubes
$x+[-u,u]$ with $x\in u\mathbb{Z}$, $x\le0$. Those of $\mathcal{D}_{g+1}(\mathbf{\Omega})$ are the cubes
$x+[-Lu,Lu]$ with $x\in Lu\mathbb{Z}$, $x\ge0$.

The cube centred at the seam belongs to $\mathcal{D}_g(\mathbf{\Omega})$, since one of its two
blocks, $[-u,0]$, lies in $B_g$, so that its centre $0$ lies on the boundary of $B_g$. This agrees
with the description of the cover in \cite[p.~229]{B-CMP96}; there $j$ is the scale index,
$\eta$ the lattice spacing, $\Lambda_j$ the union of big blocks of scale $j$ forming the $j$-th layer,
and $B^j(\Lambda_j)$ that layer as a subset of the fine lattice:
\begin{quote}
We cover $B^j(\Lambda_j)$ by a sum of cubes $\square$ of the size $2ML^j\eta$, each cube being a sum
of $2^d$ big blocks with a center $y\in\Lambda_j$ (more exactly it belongs to the boundary of this
set also).
\end{quote}
It also agrees with the sentence of \cite[pp.~229--230]{B-CMP96}
\begin{quote}
Let us assume that $\square$ is a cube connected with a $L^j\eta$-scale (i.e. a cube of the size
$2ML^j\eta$) and intersecting maybe the domain $B^{j+1}(\Lambda_{j+1})$
\end{quote}
according to which a cube of the $j$-th scale may intersect the next layer. For the same
reason the cube $[-Lu,Lu]$ belongs to $\mathcal{D}_{g+1}(\mathbf{\Omega})$.

Let $y\le u/3$. For $x\in u\mathbb{Z}$ with $x\ge u$ we have $(y-x)/u\le-2/3$, so that
$h((y-x)/u)=0$, since the support of $h$ lies in $]-2/3,2/3[$. Hence, by
$\sum_{n}h^2(t-n)=1$,
\begin{equation*}
\sum_{x\in u\mathbb{Z},\,x\le0}h^2\Big(\frac{y-x}{u}\Big)
=\sum_{n\in\mathbb{Z}}h^2\Big(\frac{y}{u}-n\Big)=1 .
\end{equation*}

Let $y\ge-Lu/3$. For $x\in Lu\mathbb{Z}$ with $x\le-Lu$ we have $(y-x)/(Lu)\ge2/3$, so these terms
vanish. In the same way the squared profiles of the scale-$(g+1)$ members sum to $1$.

On $[-Lu/3,u/3]$ both sums equal $1$. In the two-layer configuration of clause (d) the mixed-size
cover $\mathcal{D}(\mathbf{\Omega})$ consists of the members of $\mathcal{D}_g(\mathbf{\Omega})$
and of $\mathcal{D}_{g+1}(\mathbf{\Omega})$, so the sum of $\eta_Q^2$ over
$\mathcal{D}(\mathbf{\Omega})$ equals $2$ there.
This is the configuration and the conclusion of clause (d) of
Lemma~\ref{4.6:lem-seam-double-count}.

\emph{(ii).} By (i), the identity $\sum_{\square\in\mathcal{D}}h_\square^2=1$ of
\cite[(2.36)]{B-CMP96} fails for the unmodified rescaled products of \cite[(1.118)]{B-CMP95}. At
\cite[p.~229]{B-CMP96} the family of \cite[(2.36)]{B-CMP96} is described as the family of functions
of \cite[(1.118)]{B-CMP95}, rescaled to the proper scales. Since the plain rescaling fails
\cite[(2.36)]{B-CMP96}, the family so described
contains a seam normalisation or a convention of equivalent effect.

That rule is not written out on \cite[pp.~230--236]{B-CMP96} or on
\cite[pp.~37--40]{B-CMP95}. The one modification of the prescription on those pages,
\cite[(2.70)]{B-CMP96}, replaces, in the operator inverted on each cube $\square$, the propagator
$G'$ by the propagator $G'(\widetilde{\square})$ of a concentric cube $\widetilde{\square}$ of side
$4M$, the restriction to $\square$ being kept. The gluing there is by the same functions
$h_\square$. The partition of unity $\{h'_\square\}$ of the modified cover in \cite{B-CMP99b} is
introduced by name as the partition corresponding to that cover, without a formula for it. For the
construction of such covers and partitions, \cite{B-CMP119} refers to two earlier papers, one of
which is \cite{B-CMP99b}; no construction from them is used here.

That the gluing hypothesis $(\mathrm{G\text{-}loc})$, $(\mathrm{G\text{-}supp})$ of
\eqref{4.6:eq-bump-profiles-c} is compatible with every constraint stated in \cite{B-CMP95} and
\cite{B-CMP96} on the cover and its partition of unity, in particular with \cite[(2.36)]{B-CMP96},
is shown by the exhibited rule $(\mathrm{G}^{\natural})$ with its covering property
\eqref{4.6:eq-gluing-example-a}.

Two rules satisfying the gluing hypothesis illustrate what such a normalisation may be; neither is
identified with the construction of \cite{B-CMP96}.

The first is the smooth normalisation across seams,
\begin{equation*}
h_\square:=\eta_\square\Big(\sum_{Q\in\mathcal{D}}\eta_Q^2\Big)^{-1/2}.
\end{equation*}
On $\mathcal{D}$ it coincides with $(\mathrm{G}^{\natural})$, and it is smooth. By
\eqref{4.6:eq-gluing-example-a} restricted to $\mathcal{D}$, its squares sum to $1$ for nests
satisfying the separation condition.

The second is the layer restriction $h_\square:=\eta_\square\mathbf{1}_{B_{j(\square)}}$, where
$j(\square)$ is the scale of $\square$. By clauses (a) and (b) of
Lemma~\ref{4.6:lem-seam-double-count}, for nests satisfying the separation condition, its squares sum to
$1$ pointwise on $\bigcup_gB_g$. This rule
reads the level datum of the point, which is why $(\mathrm{G\text{-}loc})$ carries a level-datum
argument. Its differences $\partial h$ are not small on bonds crossing a seam. This does not match
the smallness of $K(h)$ on which \cite[(2.38)--(2.44)]{B-CMP96} relies, and this rule is not
considered further.

\emph{(iii).} Each ingredient of the statements comparing two nests is treated in turn.
\begin{itemize}
\item The bump window is $r_h=2$ by the lemma of this section determining the bump window of the
glued bumps, whose proof uses only $(\mathrm{G\text{-}loc})$ and $(\mathrm{G\text{-}supp})$.
\item The index data do not involve any bump: neither the lemma of this section bounding their
reach nor the statement of this chapter on their locality reads a bump.
\item The value statement for whole atoms $K_\omega$ involves the bumps only through the single
map $\Phi^h$, which does not depend on the nest; by the statement of this section on the single
evaluation map, one and the same map is applied to both nests.
\item The margin in the support property of $K_\omega$ involves only $(\mathrm{G\text{-}supp})$.
\end{itemize}
Hence none of these ingredients depends on the choice of rule satisfying the gluing hypothesis.

A rule not satisfying the gluing hypothesis, with bump window $r_h$, would change only the reading
radius $r_*$, defined in the section of this chapter on the reading radius, to $\max\{2,r_h\}$.
An example is a rule normalising at a point against members that do not contain the point.
\end{proof}

\begin{lemma}[The bump of an inserted cube on its outer shell]\label{4.6:lem-inserted-shell}
Let $\mathbf{\Omega}$ be a nest, let $\mathcal{D}=\mathcal{D}(\mathbf{\Omega})$ be its cover with
the scale classes $\mathcal{D}_g$ of Definition~\ref{4.6:def-local-sequence}, and let
$\square\in\mathcal{D}_j$ be an anchor cube to which the first case of the surgery of
Definition~\ref{4.6:def-cover-surgery} applies. Thus $\square_0:=\square^{2}$ is inserted, every
member of $\mathcal{D}$ that intersects $\square_0$ is removed, and $\mathcal{D}'$ denotes the
surgered cover. Put $u:=u_j$, and assume
\begin{equation}\label{4.6:eq-inserted-shell-1}
\square^{(4)}\subset\operatorname{int}\Omega_j ,
\end{equation}
as holds, for instance, for an anchor lying deep inside the layer $B_j$. Write
$\operatorname{dist}_\infty$ for the $\ell^\infty$ distance on $\mathbb{T}$ and let
\begin{equation*}
S:=\{\,y\in\mathbb{T} : 0<\operatorname{dist}_\infty(y,\square_0)<u\,\}
\end{equation*}
be the outer shell of $\square_0$ of width $u$.
\begin{itemize}
\item[(a)] Suppose that a rule assigns to every kept member $C\in\mathcal{D}'\setminus\{\square_0\}$
a bump $h'_C$ with $\operatorname{supp}h'_C\subset C$ (as the profiles of
\cite[(1.118)]{B-CMP95} satisfy, and as the gluing hypothesis \eqref{4.6:eq-bump-profiles-c}
posits), assigns a function $h'_{\square_0}$ to $\square_0$, and achieves
\begin{equation}\label{4.6:eq-inserted-shell-2}
\sum_{C\in\mathcal{D}'}h'^{\,2}_C=1\quad\text{on }\square_0^{(1)} .
\end{equation}
Then $h'^{\,2}_{\square_0}\equiv1$ on $S$. In particular
$\operatorname{supp}h'_{\square_0}\supsetneq\square_0$, and $h'_{\square_0}$ is not a profile of
the type of \cite[(1.118)]{B-CMP95} for $\square_0$ rescaled to the half-side $3u$ of $\square_0$;
such a profile vanishes outside $\square_0^{+(-(1/3)\cdot3u)}\subset\square_0$.
\item[(b)] Suppose instead that every member $Q\in\mathcal{D}'$ carries a profile $\eta_Q$ of the
type of \cite[(1.118)]{B-CMP95}, with $\eta_{\square_0}$ supported inside $\square_0$, and let
\begin{equation*}
h_Q[\mathcal{D}']:=\eta_Q\Bigl(\sum_{C\in\mathcal{D}'}\eta_C^2\Bigr)^{-1/2},
\end{equation*}
with $h_Q[\mathcal{D}']:=0$ where the sum vanishes. Then
$\sum_{Q\in\mathcal{D}'}h_Q[\mathcal{D}']^2=0$ on $S$. Thus the uniform normalisation with all
profiles of the type of \cite[(1.118)]{B-CMP95} is zero on the shell, and it is not a partition of
unity on $\square_0^{(1)}$.
\end{itemize}
\end{lemma}

\begin{proof}
We first unwind the geometry. The anchor is a candidate cube of scale $j$, so
$\square=x+[-u,u]^d$ with $x\in u\mathbb{Z}^d$
(Definition~\ref{4.6:def-bump-profiles}, display \eqref{4.6:eq-bump-profiles-a}). Hence
$\square_0=\square^{2}=x+[-3u,3u]^d$ is a box whose corners lie in $u\mathbb{Z}^d$. The concentric
enlargements are
\begin{equation*}
\square^{(4)}=\square_0^{(2)}=\square_0^{+2u}=x+[-5u,5u]^d,
\qquad
\square_0^{(1)}=\square_0^{+u}=x+[-4u,4u]^d .
\end{equation*}
In particular $S\subset\square_0^{+u}=\square_0^{(1)}$, so \eqref{4.6:eq-inserted-shell-2} holds
on $S$.

\emph{Kept members of scale $j'\ge j$.} Let $C\in\mathcal{D}'\setminus\{\square_0\}$ be of scale
$j'\ge j$. Then $C=x'+[-u_{j'},u_{j'}]^d$ with $x'\in u_{j'}\mathbb{Z}^d$. Since
$u_{j'}=L^{j'-j}u$ is an integer multiple of $u$, we have $u_{j'}\mathbb{Z}^d\subset u\mathbb{Z}^d$,
so $C$ is also a box with corners in $u\mathbb{Z}^d$. Because $C$ was kept by the surgery,
$C\cap\square_0=\emptyset$.

The $\ell^\infty$ distance between two boxes is the largest, over the coordinates, of the gaps
between their coordinate intervals. For boxes with corners in $u\mathbb{Z}^d$ every such gap lies
in $u\mathbb{Z}_{\ge0}$. Since the two closed boxes are disjoint, some gap is positive. Hence
$\operatorname{dist}_\infty(C,\square_0)\in u\mathbb{Z}_{>0}$, and therefore
$\operatorname{dist}_\infty(C,\square_0)\ge u$. Thus every point of $C$ has distance at least $u$
from $\square_0$, so $C\cap S=\emptyset$. Since $\operatorname{supp}h'_C\subset C$, it follows that
$h'_C=0$ on $S$.

\emph{Kept members of scale $j'<j$.} Let $C'\in\mathcal{D}'\setminus\{\square_0\}$ be a member of
$\mathcal{D}_{j'}$ with $j'<j$. By the membership condition of
Definition~\ref{4.6:def-local-sequence}, there is a block $e'\in\pi_{j'}$ with $e'\subset C'$,
$[e'\subset\Omega_{j'}]=1$ and $[e'\subset\Omega_{j'+1}]=0$. The nest is decreasing and
$j'+1\le j$, so $\Omega_{j'+1}\supseteq\Omega_j$. Since $e'\not\subset\Omega_{j'+1}$, also
$e'\not\subset\Omega_j$. Hence $\operatorname{int}e'\cap\Omega_j=\emptyset$, and so
$e'\subset\operatorname{cl}(\mathbb{T}\setminus\Omega_j)$.

The cube $C'$ has side $2u_{j'}$ and contains the block $e'$ of side $u_{j'}$. In each coordinate,
an interval of length $2u_{j'}$ that contains a subinterval of length $u_{j'}$ lies within distance
$u_{j'}$ of that subinterval, so $C'\subset e'^{\,+u_{j'}}$. Moreover $u_{j'}=L^{j'-j}u\le u/L$.
Hence
\begin{equation*}
C'\subset e'^{\,+u_{j'}}\subset\operatorname{cl}(\mathbb{T}\setminus\Omega_j)^{+u/L}.
\end{equation*}

By hypothesis \eqref{4.6:eq-inserted-shell-1} we have
$\square_0^{(2)}=\square^{(4)}\subset\operatorname{int}\Omega_j$. Since
$\operatorname{cl}(\mathbb{T}\setminus\Omega_j)=\mathbb{T}\setminus\operatorname{int}\Omega_j$,
this gives $\square_0^{(2)}\cap\operatorname{cl}(\mathbb{T}\setminus\Omega_j)=\emptyset$. As
$\square_0^{(2)}=\{y:\operatorname{dist}_\infty(y,\square_0)\le2u\}$, the closed set
$\operatorname{cl}(\mathbb{T}\setminus\Omega_j)$ has no point within distance $2u$ of the compact set
$\square_0$. Since the distance between these two sets is attained, we conclude
$\operatorname{dist}_\infty(\square_0,\operatorname{cl}(\mathbb{T}\setminus\Omega_j))>2u$.

Now let $y\in C'$, and pick $z\in\operatorname{cl}(\mathbb{T}\setminus\Omega_j)$ with
$\operatorname{dist}_\infty(y,z)\le u/L$. By the triangle inequality, and since $L\ge2$,
\begin{equation*}
\operatorname{dist}_\infty(y,\square_0)\ \ge\ \operatorname{dist}_\infty(z,\square_0)-\frac{u}{L}
\ >\ 2u-\frac{u}{L}\ \ge\ \frac{3}{2}\,u\ >\ u .
\end{equation*}
Hence $C'\cap S=\emptyset$. Since $\operatorname{supp}h'_{C'}\subset C'$, it follows that
$h'_{C'}=0$ on $S$.

\emph{Conclusion for (a).} Every scale $j'$ is either at least $j$ or smaller than $j$, so the two
cases cover all kept members. Thus every kept bump vanishes on $S$, and
\eqref{4.6:eq-inserted-shell-2} on $S$ reduces to $h'^{\,2}_{\square_0}=1$ there.

The same two arguments show that every kept member is disjoint from $\square_0$: in the first case
by the surgery, and in the second case because its distance from $\square_0$ exceeds $u$. Hence the
kept bumps also vanish on $\square_0\subset\square_0^{(1)}$, and
\eqref{4.6:eq-inserted-shell-2} gives $h'^{\,2}_{\square_0}=1$ on $\square_0$ as well. Therefore
$\operatorname{supp}h'_{\square_0}$ contains $\square_0$ and the closure of $S$. Since $S$ is
non-empty and disjoint from $\square_0$, we get $S\not\subset\square_0$, and so
$\operatorname{supp}h'_{\square_0}\supsetneq\square_0$.

A profile of the type of \cite[(1.118)]{B-CMP95} for $\square_0$, rescaled to its half-side $3u$,
is built from the one-variable function $h$, which is supported in $]-2/3,2/3[$. Such a profile
therefore vanishes outside
\begin{equation*}
\square_0^{+(-(1/3)\cdot3u)}=x+[-2u,2u]^d\subset\square_0 .
\end{equation*}
In particular it vanishes on $S$, so it cannot equal $h'_{\square_0}$.

\emph{Proof of (b).} Let $y\in S$. Then $y\notin\square_0$, and $\eta_{\square_0}$ is supported
inside $\square_0$, so $\eta_{\square_0}(y)=0$. For a kept member $C$, the profile satisfies
$\operatorname{supp}\eta_C\subset\operatorname{int}C^{(-1/3)}\subset C$
(Definition~\ref{4.6:def-bump-profiles}). We have shown above that $C\cap S=\emptyset$, so
$\eta_C(y)=0$. Hence $\sum_{C\in\mathcal{D}'}\eta_C^2(y)=0$. By the convention fixed in part~(b),
$h_Q[\mathcal{D}'](y)=0$ for every $Q\in\mathcal{D}'$, and therefore
$\sum_{Q\in\mathcal{D}'}h_Q[\mathcal{D}']^2=0$ on $S$. Since $S\subset\square_0^{(1)}$, the family
$\{h_Q[\mathcal{D}']\}$ is not a partition of unity on $\square_0^{(1)}$.
\end{proof}

Part~(a) agrees with Ba{\l}aban's construction. In \cite{B-CMP99b} the Dirichlet box of the derived
cube is $\square_0^{(4)}$ (through the local sequence \eqref{4.6:eq-local-sequence-a}), a box much
larger than $\square_0$. There the bump attached to $\square_0$ is required to equal one only on a
cube $\widetilde{\square}\subset\square_0$. In \cite[p.~38]{B-CMP95} two cubes $\square_{z_1}$,
$\square_{z_2}$ of the cover may be disjoint or overlapping, and the two positions are treated as
separate cases of an estimate: ``We have to consider separately the cases when $\square_{z_1}$,
$\square_{z_2}$ are disjoint, and when they are overlapping''.

\begin{proposition}[An example of an admissible gluing rule]\label{4.6:prop-gluing-example}
For a cube $C$ of scale $g$ and a real number $t$ we write $C^{(t)}$ for the concentric cube
whose half-side is that of $C$ plus $t\,u_g$; thus $C^{(t)}\subset C$ for $t<0$. Fix, once and for
all and independently of the nest, a function $\theta\in C^\infty([0,\infty))$ with
$0\le\theta\le1$, $\theta\equiv1$ on $[0,1/4]$ and $\theta\equiv0$ on $[1/2,\infty)$. For an
inserted cube $I=x_I+[-s,s]^d$ of scale $g$, with $s\in u_g\mathbb{N}$, write $u:=u_g$ and put
\begin{equation}\label{4.6:eq-gluing-example-1}
\psi_I(y):=\prod_{\mu=1}^{d}\theta\Bigl(\frac{\bigl(|y_\mu-x_{I,\mu}|-s-\tfrac43 u\bigr)^{+}}{u}\Bigr).
\end{equation}
Let $\mathcal{F}$ be either the cover $\mathcal{D}(\mathbf{\Omega})$ or a surgered family
$\mathcal{F}_{\mathfrak{s}}(\mathbf{\Omega})$ (Definition~\ref{4.6:def-cover-surgery}); its
inserted cubes are the cubes created by surgery, and $\mathcal{D}(\mathbf{\Omega})$ has none. Let
$\eta_Q$, $Q\in\mathcal{C}$, be the one-scale profiles of Definition~\ref{4.6:def-bump-profiles}.
Define
\begin{equation}\label{4.6:eq-gluing-example-2}
N_{\mathcal{F}}:=\sum_{Q\in\mathcal{F}\cap\mathcal{C}}\eta_Q^2
+\sum_{I\ \text{inserted in}\ \mathcal{F}}\psi_I^2 ,
\end{equation}
and, on the set $\{N_{\mathcal{F}}>0\}$,
\begin{equation}\label{4.6:eq-gluing-example-3}
h^{\natural}_Q[\mathcal{F}]:=\eta_Q\,N_{\mathcal{F}}^{-1/2}\quad(Q\in\mathcal{F}\cap\mathcal{C}),
\qquad
h^{\natural}_I[\mathcal{F}]:=\psi_I\,N_{\mathcal{F}}^{-1/2}\quad(I\ \text{inserted}),
\end{equation}
with $h^{\natural}:=0$ on the complement of $\{N_{\mathcal{F}}>0\}$. Then the following hold.
\begin{enumerate}
\item[(a)] $\psi_I$ is a $C^\infty$ function of $y$ that depends only on $y$, $x_I$, $s$ and
$u_g$, not on the nest; $0\le\psi_I\le1$; $\psi_I\equiv1$ on $I^{(4/3)}$; and
\begin{equation*}
\operatorname{supp}\psi_I\subset I^{(4/3+1/2)}\subset\operatorname{int}I^{(2)}\subset I^{(5)}.
\end{equation*}
\item[(b)] $h^{\natural}$ has the locality property $(\mathrm{G\text{-}loc})$ of
\eqref{4.6:eq-bump-profiles-c}: $N_{\mathcal{F}}(y)$ involves $\eta_{Q'}(y)$ only for candidate
cubes $Q'\ni y$, through their membership bits, and the profiles of the inserted cubes; hence
$N_{\mathcal{F}}(y)$, and with it $h^{\natural}_Q[\mathcal{F}](y)$ and
$h^{\natural}_I[\mathcal{F}](y)$, are label-free functions of the data of
Definition~\ref{4.6:def-cover-surgery}, in which the level datum is not used.
\item[(c)] $h^{\natural}$ has the support property $(\mathrm{G\text{-}supp})$ of
\eqref{4.6:eq-bump-profiles-c}:
\begin{equation*}
\operatorname{supp}h^{\natural}_Q\subset\operatorname{supp}\eta_Q\subset Q^{(-1/3)}\subset Q,
\qquad
\operatorname{supp}h^{\natural}_I\subset I^{(2)}\subset I^{(5)}.
\end{equation*}
\item[(d)] $\sum_{\mathcal{F}}(h^{\natural})^2$ equals $1$ on $\{N_{\mathcal{F}}>0\}$ and vanishes
on its complement.
\item[(e)] Let $\mathbf{\Omega}=(\Omega_0\supseteq\dots\supseteq\Omega_k)$ be $\pi$-grained and
decreasing, as in Lemma~\ref{4.6:lem-seam-double-count}, and satisfy
\eqref{4.6:eq-nest-traces-a} at every level. Then $\mathcal{F}=\mathcal{D}(\mathbf{\Omega})$
satisfies
\begin{equation}\label{4.6:eq-gluing-example-a}
\{N_{\mathcal{F}}>0\}\supseteq\Omega_0 ;
\end{equation}
in particular $N_{\mathcal{F}}>0$ on all of $\mathbb{T}$ when $\Omega_0=\mathbb{T}$.
\item[(f)] Let $\mathcal{F}\in\{\mathcal{D}'(\square;\mathbf{\Omega}),
\mathcal{D}''(\square;\mathbf{\Omega})\}$ be a family obtained by surgery at one anchor cube
$\square$, with the single inserted cube $\square_0$. Then $h^{\natural}_{\square_0}=1$ on
$\square_0$, and hence on every subset of $\square_0$, in particular on the double cube
$\widetilde{\square}\subset\square_0$ on which \cite{B-CMP99b} states $h'_{\square_0}=1$, where
$h'_{\square_0}$ denotes the bump that \cite{B-CMP99b} assigns to the inserted cube.
\item[(g)] In the setting of \textup{(e)}: at a point $y\in B_g$ with
$\operatorname{dist}_\infty(y,B_{g'})>\tfrac23u_{g'}$ for every layer $B_{g'}$ with $g'\neq g$,
one has $N_{\mathcal{D}(\mathbf{\Omega})}(y)=1$ and $h^{\natural}_Q(y)=\eta_Q(y)$ for every
member $Q$ of $\mathcal{D}(\mathbf{\Omega})$, i.e.\ the profile of \cite[(1.118)]{B-CMP95} itself.
At a point
$y\in B_g\cap B_{g'}$ with $g\neq g'$ one has $N_{\mathcal{D}(\mathbf{\Omega})}(y)\ge2$ and
$h^{\natural}_Q(y)=\eta_Q(y)/\sqrt{N_{\mathcal{D}(\mathbf{\Omega})}(y)}$.
\end{enumerate}
The proposition makes no assertion about \eqref{4.6:eq-gluing-example-a} for surgered families,
and it does not identify $h^{\natural}$ with the gluing used in \cite{B-CMP99b} or in any other
of Ba{\l}aban's papers.
\end{proposition}

\begin{proof}
(a) The right-hand side of \eqref{4.6:eq-gluing-example-1} is built from $y$, the centre $x_I$,
the half-side $s$, the scale $u=u_g$ and the fixed function $\theta$ only; the nest does not enter.
Each factor takes values in $[0,1]$ because $\theta$ does, so $0\le\psi_I\le1$.

If $y\in I^{(4/3)}$, then $|y_\mu-x_{I,\mu}|\le s+\tfrac43u$ for every $\mu$, so every argument of
$\theta$ in \eqref{4.6:eq-gluing-example-1} equals $0$, and $\theta(0)=1$ gives $\psi_I(y)=1$.

If $\psi_I(y)\neq0$, every factor is non-zero; since $\theta\equiv0$ on $[1/2,\infty)$, each
argument is $<1/2$, i.e.
\begin{equation*}
|y_\mu-x_{I,\mu}|<s+\bigl(\tfrac43+\tfrac12\bigr)u\qquad(\mu=1,\dots,d).
\end{equation*}
Hence $\{\psi_I\neq0\}\subset\operatorname{int}I^{(4/3+1/2)}$, and its closure lies in
$I^{(4/3+1/2)}=I^{(11/6)}$. Since $11/6<2$, this cube lies in $\operatorname{int}I^{(2)}$, and
$I^{(2)}\subset I^{(5)}$.

For smoothness it suffices to treat one factor, the function
$f(t):=\theta\bigl((|t|-a)^{+}/u\bigr)$ of one real variable, with $a:=s+\tfrac43u>0$; the only
points at which the inner expression fails to be smooth are $t=0$ (the kink of $|\cdot|$) and
$|t|=a$ (the kink of $(\cdot)^{+}$). On the open set $\{|t|<a+u/4\}$ the argument of $\theta$ lies
in $[0,1/4)$, where $\theta\equiv1$, so $f\equiv1$ there; this set contains both $t=0$ and
$|t|=a$. On the open set $\{|t|>a\}$ we have $|t|\neq0$ and $(|t|-a)^{+}=|t|-a$, so $f$ is the
composition of the smooth maps $t\mapsto(|t|-a)/u$ and $\theta$. These two open sets cover
$\mathbb{R}$, so $f$ is $C^\infty$, and so is $\psi_I$, a product of such functions of the
coordinates.

(b) By \eqref{4.6:eq-gluing-example-2},
\begin{equation*}
N_{\mathcal{F}}(y)=\sum_{Q'\in\mathcal{C}}\mathbf{1}[Q'\in\mathcal{F}]\,\eta_{Q'}(y)^2
+\sum_{I\ \text{inserted in}\ \mathcal{F}}\psi_I(y)^2 .
\end{equation*}
By Definition~\ref{4.6:def-bump-profiles}, $\operatorname{supp}\eta_{Q'}\subset Q'$, so
$\eta_{Q'}(y)\neq0$ only for candidate cubes $Q'$ containing $y$; these enter with their membership
bits $\mathbf{1}[Q'\in\mathcal{F}]$. The profiles $\eta_{Q'}$ are nest-free functions of $y$ and of
the centre and scale of $Q'$ (Definition~\ref{4.6:def-bump-profiles}), and the $\psi_I$ are
nest-free by (a). Hence $N_{\mathcal{F}}(y)$ is a label-free function of the membership bits and the
inserted cubes, that is, of the data of Definition~\ref{4.6:def-cover-surgery}, and the level
datum does not occur. The functions \eqref{4.6:eq-gluing-example-3} are obtained from
$N_{\mathcal{F}}(y)$, $\eta_Q(y)$ and $\psi_I(y)$ by the fixed operations of
\eqref{4.6:eq-gluing-example-3}, and so inherit these properties. This is $(\mathrm{G\text{-}loc})$
of \eqref{4.6:eq-bump-profiles-c}.

(c) On $\{N_{\mathcal{F}}>0\}$, $h^{\natural}_Q$ vanishes where $\eta_Q$ does, and off this set
$h^{\natural}_Q=0$; hence $\operatorname{supp}h^{\natural}_Q\subset\operatorname{supp}\eta_Q$, and
Definition~\ref{4.6:def-bump-profiles} gives
\begin{equation*}
\operatorname{supp}\eta_Q\subset\operatorname{int}Q^{(-1/3)}\subset Q^{(-1/3)}\subset Q .
\end{equation*}
In the
same way $\operatorname{supp}h^{\natural}_I\subset\operatorname{supp}\psi_I$, which lies in
$I^{(2)}\subset I^{(5)}$ by (a). This is $(\mathrm{G\text{-}supp})$ of
\eqref{4.6:eq-bump-profiles-c}; for members $Q$ of $\mathcal{D}(\mathbf{\Omega})$ it contains the
support condition $\operatorname{supp}h_{\square}\subset\square^{(-1/3)}$ that \cite{B-CMP99b}
states for the bump $h_{\square}$ it assigns to a member $\square$ of the cover.

(d) At a point $y$ with $N_{\mathcal{F}}(y)>0$, by \eqref{4.6:eq-gluing-example-3} and
\eqref{4.6:eq-gluing-example-2},
\begin{equation*}
\sum_{Q\in\mathcal{F}\cap\mathcal{C}}h^{\natural}_Q(y)^2+\sum_{I}h^{\natural}_I(y)^2
=\frac{1}{N_{\mathcal{F}}(y)}\Bigl(\sum_{Q\in\mathcal{F}\cap\mathcal{C}}\eta_Q(y)^2
+\sum_{I}\psi_I(y)^2\Bigr)=1 ,
\end{equation*}
and at a point with $N_{\mathcal{F}}(y)=0$ every $h^{\natural}$ vanishes by definition.

(e) Here $\mathcal{F}=\mathcal{D}(\mathbf{\Omega})$ has no inserted cubes. Let $y\in\Omega_0$ and
$g:=\max\{i\le k:y\in\Omega_i\}$, where for $g=k$ we read $\Omega_{k+1}=\emptyset$. By maximality
of $g$, $y\notin\Omega_{g+1}$, hence
$y\notin\operatorname{int}\Omega_{g+1}\subseteq\Omega_{g+1}$; since $y\in\Omega_g$, we get
$y\in\Omega_g\setminus\operatorname{int}\Omega_{g+1}=B_g$. The members of
$\mathcal{D}_g(\mathbf{\Omega})$ are members of $\mathcal{D}(\mathbf{\Omega})\cap\mathcal{C}$, and all
terms of \eqref{4.6:eq-gluing-example-2} are non-negative; therefore, by part (b) of
Lemma~\ref{4.6:lem-seam-double-count}, whose hypotheses are those of (e), in particular
\eqref{4.6:eq-nest-traces-a} at level $g$,
\begin{equation*}
N_{\mathcal{D}(\mathbf{\Omega})}(y)\ \ge\ \sum_{Q\in\mathcal{D}_g(\mathbf{\Omega})}\eta_Q(y)^2\ =\ 1 .
\end{equation*}
Thus $\Omega_0\subseteq\{N_{\mathcal{D}(\mathbf{\Omega})}>0\}$, which is
\eqref{4.6:eq-gluing-example-a}; if $\Omega_0=\mathbb{T}$ the set is all of $\mathbb{T}$.

(f) By Definition~\ref{4.6:def-cover-surgery}, the surgery removes every cube intersecting
$\square_0$, so every kept candidate $C\in\mathcal{F}\cap\mathcal{C}$ satisfies
$C\cap\square_0=\emptyset$; as $\operatorname{supp}\eta_C\subset C$, every kept $\eta_C$ vanishes on
$\square_0$. The only inserted cube is $\square_0$, and $\psi_{\square_0}\equiv1$ on
$\square_0^{(4/3)}\supset\square_0$ by (a). Hence $N_{\mathcal{F}}=\psi_{\square_0}^2=1$ on
$\square_0$, and \eqref{4.6:eq-gluing-example-3} gives $h^{\natural}_{\square_0}=1$ there. The
restriction to one inserted cube is used at this point: for a surgery datum $\mathfrak{s}$ with a
second inserted cube $I'$ whose $\psi_{I'}$ is positive at some point of $\square_0$, one has there
$N_{\mathcal{F}}\ge\psi_{\square_0}^2+\psi_{I'}^2>1$, hence $h^{\natural}_{\square_0}<1$, and the
conclusion of (f) fails; such data are not covered by (f).

(g) Group the members of $\mathcal{D}(\mathbf{\Omega})$ by their scale, $\mathcal{D}_{g'}(\mathbf{\Omega})$
being those of scale $g'$. By part (b) of Lemma~\ref{4.6:lem-seam-double-count}, for each $g'$ the sum
$\sum_{Q\in\mathcal{D}_{g'}(\mathbf{\Omega})}\eta_Q^2$ vanishes off $B_{g'}^{+(2/3)u_{g'}}$ and equals $1$
on $B_{g'}$. At a point $y\in B_g$ with $\operatorname{dist}_\infty(y,B_{g'})>\tfrac23u_{g'}$ for all
$g'\neq g$, $y$ lies outside every $B_{g'}^{+(2/3)u_{g'}}$, $g'\neq g$, so only the scale-$g$ terms
survive and
\begin{equation*}
N_{\mathcal{D}(\mathbf{\Omega})}(y)=\sum_{Q\in\mathcal{D}_g(\mathbf{\Omega})}\eta_Q(y)^2=1 ;
\end{equation*}
then \eqref{4.6:eq-gluing-example-3} gives $h^{\natural}_Q(y)=\eta_Q(y)$ for every member $Q$ of
$\mathcal{D}(\mathbf{\Omega})$. At a point
$y\in B_g\cap B_{g'}$ with $g\neq g'$ the same lemma gives
\begin{equation*}
\sum_{Q\in\mathcal{D}_g(\mathbf{\Omega})}\eta_Q(y)^2=\sum_{Q\in\mathcal{D}_{g'}(\mathbf{\Omega})}\eta_Q(y)^2=1 ,
\end{equation*}
so $N_{\mathcal{D}(\mathbf{\Omega})}(y)\ge2$ by non-negativity of the remaining terms, with equality
in the model of part (d) of Lemma~\ref{4.6:lem-seam-double-count}; the formula for $h^{\natural}_Q(y)$
is \eqref{4.6:eq-gluing-example-3}.
\end{proof}

For surgered families the construction gives the following towards
\eqref{4.6:eq-gluing-example-a}. Let $\mathcal{F}$ be a surgered family with inserted cube
$\square_0$. Since $\psi_{\square_0}\equiv1$ on $\square_0^{(4/3)}$, one has $N_{\mathcal{F}}\ge1$
on this cube. In the setting of Lemma~\ref{4.6:lem-inserted-shell}, that is, for a case-1 anchor
$\square\in\mathcal{D}_j$ with $\square^{(4)}\subset\operatorname{int}\Omega_j$,
$\square_0=\square^{2}$ and $u=u_j$, this cube contains the shell
$S=\{y:0<\operatorname{dist}_\infty(y,\square_0)<u_j\}$ of that lemma and the further region up to
$\ell^\infty$-distance $\tfrac43u_j$ from $\square_0$, on which no kept profile is positive: each
kept cube of scale $j$ has its centre at distance at least $5u_j$ from the centre of $\square_0$ in
some coordinate, and, as $\operatorname{supp}\eta_C\subset\operatorname{int}C^{(-1/3)}$, its profile
vanishes at every point whose distance from the centre of $\square_0$ in that coordinate is at most
$5u_j-\tfrac23u_j=\tfrac{13}{3}u_j$. Beyond this region, positivity of
$N_{\mathcal{F}}$ depends on where the one-variable function $h$ of \cite[(1.118)]{B-CMP95} is
non-zero inside $(-2/3,2/3)\setminus[-1/3,1/3]$, which is not specified by the conditions on $h$
listed in Definition~\ref{4.6:def-bump-profiles}. Accordingly $h^{\natural}$ satisfies the gluing
hypothesis \eqref{4.6:eq-bump-profiles-c} (parts (b)--(d) of
Proposition~\ref{4.6:prop-gluing-example}) and the conditions of \cite{B-CMP99b} and
\cite[(1.118)]{B-CMP95} established in parts (c), (f) and (g); only the inclusion
\eqref{4.6:eq-gluing-example-a} for surgered families is left aside.

\begin{lemma}[Gluing rules read the nest within two top blocks]\label{4.6:lem-gluing-window}
Use the torus $\mathbb{T}$, the enlargements $A^{+s}$ and the block sides of
Notation~\ref{4.6:not-block-sides}, the candidate cubes $\mathcal{C}$ and the cover
$\mathcal{D}(\mathbf{\Omega})$ of a nest $\mathbf{\Omega}$, and the surgeries of
Definition~\ref{4.6:def-cover-surgery}. Let the glued bumps obey
$(\mathrm{G\text{-}loc})$ and $(\mathrm{G\text{-}supp})$ of \eqref{4.6:eq-bump-profiles-c},
with scope $u\le M$ in the sense of \eqref{4.6:eq-bump-profiles-c}, so that every candidate
the rule reads has half-side $u'\le M$. Let $\mathcal{F}=\mathcal{F}_{\mathfrak{s}}(\mathbf{\Omega})$
be any family obtained from $\mathcal{D}(\mathbf{\Omega})$ by surgery with surgery datum
$\mathfrak{s}$, the cover $\mathcal{F}=\mathcal{D}(\mathbf{\Omega})$ itself, with no surgery
performed, included; and let $Q$ be a member of $\mathcal{F}$, either a candidate or an inserted
cube. Write $j(Q)$ for its scale index: $j(Q)=g$ when $Q\in\mathcal{C}_g$, and for an inserted
cube the index attached to it by the surgery of Definition~\ref{4.6:def-cover-surgery}. Then the
glued bump $h_Q[\mathcal{F}]$, as a
function on $\mathbb{T}$, together with its lattice differences $\partial h_Q(b)$ at bonds
$b$ and $\Delta h_Q(x)$ at lattice points $x$, is determined by
\begin{enumerate}
\item the pair $(Q,j(Q))$;
\item the surgery datum $\mathfrak{s}$, that is, the anchors and the inserted cubes;
\item the bits of the nest $\mathbf{\Omega}$ on $Q^{+2M}$; when $Q=I$ is an inserted cube,
the bits of $\mathbf{\Omega}$ on $(I^{(5)})^{+2M}$, a set contained in
$F(\mathfrak{a})^{+2M}$, where $F(\mathfrak{a})$ is the footprint, in the sense of
\eqref{4.6:eq-letters-words-a}, of the letter $\mathfrak{a}$ in which the bump of $I$ enters.
\end{enumerate}
More precisely, $h_Q[\mathcal{F}]$ and its lattice differences are determined by
$(Q,j(Q))$, by the membership bits $[Q'\in\mathcal{D}(\mathbf{\Omega})]$ of the candidates
$Q'\subset Q^{+2M}$, each of which is a test at radius $0$ at $Q'$ by
Lemma~\ref{4.6:lem-membership-tests}\,(b), by the label-free incidence of these candidates
with the inserted cubes, and by the level data of the points of $Q$, which are bits of
$\pi_0$-blocks contained in $Q^{+M}$ (with $Q^{+2M}$, $Q^{+M}$ and the points of $Q$
replaced by $(I^{(5)})^{+2M}$, $(I^{(5)})^{+M}$ and the points of $I^{(5)}$ when $Q=I$
is inserted). Thus every gluing rule obeying $(\mathrm{G\text{-}loc})$ and
$(\mathrm{G\text{-}supp})$ has the absolute bump window $r_h=2$, that is, a window $2M$ of two
top blocks, the same for every scale $j(Q)$.
\end{lemma}

\begin{proof}
Write $S:=Q$ if $Q$ is a candidate and $S:=I^{(5)}$ if $Q=I$ is an inserted cube.

\emph{Support.} By $(\mathrm{G\text{-}supp})$ the function $h_Q[\mathcal{F}]$ vanishes off
$S$. This is known from $Q$ and $j(Q)$ alone, without reading any bit of the nest.

\emph{Values at points of $S$.} Let $y\in S$. By $(\mathrm{G\text{-}loc})$ the value
$h_Q[\mathcal{F}](y)$ is determined by the candidates $Q'$ with $y\in Q'$, together with
their bits $[Q'\in\mathcal{F}]$, and by the level datum of $y$. We locate each of these
inputs.

First, by the scope $u\le M$, a candidate $Q'$ read by the rule is a closed cube of side
$2u'$ with $u'\le M$. Every point of a closed cube of side $2u'$ containing $y$ lies at
$\ell^\infty$-distance at most $2u'$ from $y$, so
\begin{equation*}
Q'\subset y^{+2u'}\subset y^{+2M}\subset S^{+2M},
\end{equation*}
the last inclusion because $y\in S$.

Second, the bit $[Q'\in\mathcal{F}]$ of such a candidate is obtained from data already
listed. Inserted cubes are never candidates, and a surgery removes from the cover exactly
the cubes meeting an inserted cube (Definition~\ref{4.6:def-cover-surgery}); hence
\begin{equation*}
[Q'\in\mathcal{F}]
=[Q'\in\mathcal{D}(\mathbf{\Omega})]\wedge[\,Q'\text{ misses every inserted cube}\,].
\end{equation*}
The second factor is the incidence of $Q'$ with the inserted cubes of the surgery datum
$\mathfrak{s}$; it involves no label of the nest. The first factor is, by
Lemma~\ref{4.6:lem-membership-tests}\,(b), a test at radius $0$ at $Q'$, so it is read
from the bits of $\mathbf{\Omega}$ on $Q'\subset S^{+2M}$.

Third, the level datum of $y$ is carried by the $\pi_0$-blocks containing $y$. A
$\pi_0$-block is a closed cube of side $M\lambda_0\le M$ (since $\lambda_0\le1$), so a
$\pi_0$-block containing $y$ lies in $y^{+M}\subset S^{+M}$.

All three inputs are therefore read inside $S^{+2M}$, and the level data inside $S^{+M}$;
this gives the assertion for the values of $h_Q[\mathcal{F}]$.

\emph{Lattice differences.} A difference $\partial h_Q(b)$ at a bond $b$, or
$\Delta h_Q(x)$ at a lattice point $x$, is a fixed linear combination of values of
$h_Q[\mathcal{F}]$ at lattice points within one lattice step of $b$, respectively of $x$.
By the support step, only those of these values at points of $S$ can be non-zero, and the
remaining ones are zero without any reading. Each value at a point of $S$ is determined by
the data listed in the statement, by the preceding step. Hence $\partial h_Q(b)$ and
$\Delta h_Q(x)$ are determined by the same data.

Combining the three steps, $h_Q[\mathcal{F}]$ and its lattice differences are determined by
$(Q,j(Q))$, the surgery datum $\mathfrak{s}$, the membership bits of the candidates in
$S^{+2M}$, their incidence with the inserted cubes, and the level data of the points of
$S$, read from $\pi_0$-blocks inside $S^{+M}$; all of these are bits of the nest on
$S^{+2M}$ or data of the surgery.
\end{proof}

\begin{lemma}[The reading table of the primitives]\label{4.6:lem-reading-table}
Assume the following hypotheses:
\begin{itemize}
\item the profile hypothesis \eqref{4.6:eq-bump-profiles-a} on the one-scale bump profiles, and the
gluing hypothesis of Definition~\ref{4.6:def-bump-profiles};
\item the block-locality of the averaging operations, Lemma~\ref{4.6:lem-averaging-locality};
\item the formula \eqref{4.6:eq-local-sequence-a} for the seam datum $T(Q)$ and the local sequence
$\mathbf{\Omega}(Q)$ of a cube;
\item the display \eqref{4.6:eq-bump-profiles-d} for the cut-off functions $\zeta_{\square}$.
\end{itemize}
Let $\mathbf{\Omega}=(\Omega_0\supseteq\dots\supseteq\Omega_k)$ be a nest with region datum
$\mathfrak{P}$, and with wavy version $\widetilde{\mathbf{\Omega}}$ where this enters, and let $U$ be
the configuration. Cubes of the cover $\mathcal{D}(\mathbf{\Omega})$ are written $C$. For a surgery
anchor $\square\in\mathcal{D}_j$ of scale $j=j(\square)$ we write $\mathcal{D}'(\square)$,
$\mathcal{D}''(\square)$ for the two surgered covers of Definition~\ref{4.6:def-cover-surgery}, with
the cubes $\square_0$, $\square_1$, $\widetilde{\square}_0$ as there; $I$ denotes an inserted cube;
in row~8 and its proof the index $\square$ of $\zeta_{\square}$ and $\widetilde{\square}$ ranges over
all cubes of the cover, as in \eqref{4.6:eq-bump-profiles-d} and \eqref{4.6:eq-candidate-cubes-b};
$\widetilde{\mathcal{D}}_n$ is the set $\mathcal{D}_n$ formed for the wavy nest. For a letter
$\mathfrak{a}$, resp.\ a word $\omega$, of the expansions, $F(\mathfrak{a})$, resp.\ $F(\omega)$,
denotes its $U$-footprint, in the sense of the definition of letters, words and footprints of the
expansions in this section; for an inserted cube $I$ we write $F$ for the $U$-footprint of the letter
or word in which $I$ occurs. Operators and scalars not defined
in this chapter are those of
\cite{B-CMP99b}, in Ba{\l}aban's notation, at the displays named in the proof. In rows 11 to 13 of the
table and in the corresponding parts of the proof the letter $C$ also denotes Ba{\l}aban's operator
$C$ of \cite{B-CMP99b}; there it is always called the operator $C$. Likewise $C^{(2)}$, $C_j^{(2)}$
and $C'^{(k)}(\Lambda)$ in rows 11 and 12 and in their proofs denote operators of \cite{B-CMP99b},
not enlargements of a cube. Likewise $\chi_P$ and $\Delta(y)$ in row~8 and its proof are the objects
of \cite{B-CMP99b} so written, read from the region datum; neither is the operator $P$ of row~12 or
the operator $\Delta$ of row~9.

We say that a primitive attached to a cube $C$ (or to a point or bond $z$) \emph{reads the nest} on
a set $V\subset\mathbb{T}$ if its value is determined by $C$ and its scale (by $z$), by the surgery
datum when a surgered cover is involved, and by $\operatorname{tr}_V(\mathbf{\Omega};\mathfrak{P})$;
it \emph{reads $U$} on $W$ if, the nest being held fixed, it depends on $U$ only through the
restriction of $U$ to the bonds in $W$. The sets in the third column below are absolute: they are
determined by the cube, its scale and $M$, not by the nest. \emph{Radius $0$} means $V=C$ for a
primitive attached to $C$, and $V$ the block containing $z$ for a primitive attached to $z$.

Then the following hold.
\begin{enumerate}
\item[(i)] Each primitive in the table reads the nest at most on the set in the third column and
reads $U$ at most on the set in the fourth column. The operators of row~12 are read on the whole
torus; they do not enter the localized formulas of \cite{B-CMP99b} themselves, each of them being
substituted there by its expansion. The scalars of row~13 cancel and occur in no localized term.
\end{enumerate}

\par\noindent
{\upshape\small
\begin{tabular}{r p{0.28\textwidth} p{0.31\textwidth} p{0.22\textwidth}}
 & primitive & nest read & $U$ read\\ \hline
1 & membership $C\in\mathcal{D}_n$, or $C\in\widetilde{\mathcal{D}}_n$ for the wavy nest
 & $C$ (radius $0$); for $\widetilde{\mathcal{D}}_n$ with the region bits on $C$ & none\\
2 & level datum at a point or bond $z$; the tests $z\in\Lambda$, $y\in\Lambda'$
 & the block of $z$ (radius $0$) & none\\
3 & local sequences $\mathbf{\Omega}(C)$, $\widetilde{\mathbf{\Omega}}(C)$ & $C^{(3)}$ & none\\
4 & generators at $C\in\mathcal{D}$ or at a derived cube $\square^{2}$, $\square_1^{2}$,
 $\square^{4}$, $\square_1^{4}$, including the local tower operators, resolvents and
 the block operators of the associated constrained quadratic problems
 & $C^{(3)}\subset C^{(5)}$; $C^{(5)}$ under the other reading of $Q'$ in
 Definition~\ref{4.6:def-local-generators} & $\Omega_0(C)\subset C^{(5)}$\\
5 & membership of a given cube $C$ in $\mathcal{D}'(\square)$, $\mathcal{D}''(\square)$
 & $C$ (radius $0$) and the anchor data $\square_0$, $\widetilde{\square}_0$; removal sets not read
 & none\\
6 & surgery $\mathcal{D}'(\square)$: case test, $\square_1$, $\square_0$; surgery
 $\mathcal{D}''(\square)$: $\widetilde{\square}_0$
 & memberships of the scale-$(j+1)$ candidates meeting $\square$, all inside
 $\square^{+2u_{j+1}}\subset\square^{+2M}$ if $j<k$, none if $j=k$; the case and the cubes
 $\square_1$, $\square_0$, $\widetilde{\square}_0$ are then fixed by
 \eqref{4.6:eq-cover-surgery-a} & none\\
7 & bumps $h_C[\mathcal{D}]$, $h_C[\mathcal{D}'(\square)]$, $h_C[\mathcal{D}''(\square)]$ and
 their differences $\partial h$, $\Delta h$
 & $C^{+2M}$, resp.\ $(I^{(5)})^{+2M}\subset F^{+2M}$ for an inserted $I$; and the anchor datum
 & none\\
8 & $\zeta_{\square}$, $\widetilde{\square}$, $1-\widetilde{\square}$, $1-\zeta_{\square}$,
 $\chi_P$, $\Delta(y)$ & radius $0$ & none\\
9 & global $Q'_n$, $Q_n$, $a_n$, $S_n(\partial h)$, $D$, $D^*$, $\Delta_U$, $\Delta$, $\Delta'$,
 $J$ at a place $y$ & layer index of $y$ (radius $0$) & the block(s) $B^{n(y)}(\cdot)$ around $y$\\
10 & $K'(h)$, $K(h)$, $P_C(\partial h)$ & as in rows 7 and 9; $C^{(3)}$ for $P_C$
 & one block around $\operatorname{supp}\partial h$; $\Omega_0(C)$ for $P_C$\\
11 & tower vertices $\Delta'_\pi$, $C^{(2)}$, $\Delta^{(2)}$, $\widetilde{\Delta}^{(2)}$, the pair
 $(Q,Q^*)$ of \cite[(3.110)]{B-CMP99b}, $D^{(2)}$, the operator $C$, $\mu(B)$, $\lambda(A)$,
 $(I+D\mu)Q$
 & the tests $b\in\Lambda_j$ and $\Lambda$, $\Lambda'$: radius $0$ & blocks (see the proof)\\
12 & unexpanded globals $G'$, $P$, $R$, the operator $C$, $G$, $G'RD^*$, $H$, $H_1$, $(QGQ^*)^{-1}$,
 $(QG_1Q^*)^{-1}$, the tower operators $\mathfrak{P}$, $\mathfrak{G}$ of
 \cite[(3.147), (3.153)]{B-CMP99b}, $H'$, $C'^{(k)}(\Lambda)$, the tower operator
 $\widetilde{\mathfrak{h}}$, $G_2$, $\widetilde{G}_2$ & $\mathbb{T}$ & $\mathbb{T}$\\
13 & scalars $Z(B)$, $Z$, $Z'$, $Z'(\mu)$, $|\det(\Delta\restriction N(Q'))|$,
 $|\det(\widetilde{\Delta}\restriction N(\widetilde{Q}'))|$, the Jacobians of the
 $\delta$-functions of \cite[(3.166), (3.168)]{B-CMP99b} and of $B=C\overline{B}$, with $C$ the
 operator $C$
 & global & global\\ \hline
\end{tabular}}

\begin{enumerate}
\item[(ii)] Every primitive of rows 1 to 11 reads the nest inside the enlargement $C^{(5)}$ of its own
cube, with two exceptions: the bumps, and through them $K'(h)$, $K(h)$ of row~10, which read within a
window of width $2M$ beyond the support
cube (Lemma~\ref{4.6:lem-gluing-window}); and the tests of the coarser candidates in the surgeries,
which read at most $2u_{j+1}\le 2M$ beyond the anchor, that is, a margin $2L$ in the units of the
anchor's own scale.
\item[(iii)] Under the layer-separation property \eqref{4.6:eq-letters-words-c} (no averaging block
of a level large enough to leave $C^{(5)}$ meets the bump supports of $C$, and, for an inserted cube
$I$, $\operatorname{supp}h_I\subset I^{(3)}$, so that the one-block spreads of the flanking operators
stay in $I^{(5)}$; under the exhibited rule $(\mathrm{G}^{\natural})$ one even has
$\operatorname{supp}h_I\subset I^{(2)}$), the locality in $U$ stated in \cite{B-CMP99b} places the
$U$-reads of a word $\omega$ in $\bigcup_i C_i^{(5)}$, the union over the cubes $C_i$ at which the
letters of $\omega$ are formed. Without
\eqref{4.6:eq-letters-words-c}, the $U$-read set of a word $\omega$ is contained in
$F(\omega)^{+2M}$, which lies inside the window of the reading radius $r_*$, since
$(I^{(5)})^{+2M}\subset F(\omega)^{+2M}$ for every inserted cube $I$ occurring in $\omega$.
\end{enumerate}
\end{lemma}

\begin{proof}
We justify the rows of the table in order, and then parts~(ii) and~(iii).

\emph{Row 1.} Membership of a cube in $\mathcal{D}_n$ is the test of
Definition~\ref{4.6:def-local-sequence}, which takes the cubes of the cover of
\cite[p.~229]{B-CMP96}, the cover used in \cite{B-CMP99b}. The test is a test on $C$ itself, so it
reads the nest with radius $0$; for the wavy nest the test is the same, with the region bits on $C$,
as in \cite{B-CMP99b}. No configuration enters.

\emph{Row 2.} The level datum at a point or bond $z$, and the predicates $z\in\Lambda$,
$y\in\Lambda'$, are those of \cite{B-CMP99b} and of \cite[(2.3), (2.20)]{B-CMP96}; each is decided
by the nest on the block containing $z$, hence has radius $0$, and reads no configuration.

\emph{Row 3.} By the formula \eqref{4.6:eq-local-sequence-a}, which is among the hypotheses,
Lemma~\ref{4.6:lem-sequence-reach} applies: the local sequence $\mathbf{\Omega}(C)$ of the
construction in \cite{B-CMP99b} depends only on $C$, its scale and the membership bits of the blocks
contained in $C^{(3)}$, and not otherwise on the nest, and likewise its wavy version
$\widetilde{\mathbf{\Omega}}(C)$.

\emph{Row 4.} The hypotheses of Proposition~\ref{4.6:prop-generator-equality} are the block-locality
of Lemma~\ref{4.6:lem-averaging-locality} and the formula \eqref{4.6:eq-local-sequence-a}; both are
assumed. By that proposition the generators at $C$, formed as in \cite{B-CMP99b} for the local
sequence $\mathbf{\Omega}(C)$, are determined by the bits of the nest on $C^{(3)}$. Under the other
reading of the operator $Q'$ in Definition~\ref{4.6:def-local-generators} the bits are read on
$C^{(5)}$. The same holds at the derived cubes and for the local tower operators, resolvents and
block operators of the associated constrained quadratic problems formed from these generators in
\cite{B-CMP99b}. The generators are
built from $U$ on $\Omega_0(C)$, and $\Omega_0(C)\subset C^{(5)}$ because all members of the local
sequence lie in $C^{(5)}$ (Lemma~\ref{4.6:lem-sequence-reach}).

\emph{Row 5.} By the construction of \cite{B-CMP99b} recalled in
Definition~\ref{4.6:def-cover-surgery}, whether a given cube $C$ belongs to $\mathcal{D}'(\square)$
or to $\mathcal{D}''(\square)$ is decided by the membership of $C$ itself (radius $0$) and by the
anchor data $\square_0$, $\widetilde{\square}_0$; the sets of removed cubes are not read.

\emph{Row 6.} The case test and the cubes $\square_1$, $\square_0$ of the surgery
$\mathcal{D}'(\square)$, and the cube $\widetilde{\square}_0$ of $\mathcal{D}''(\square)$, are
determined, as in \cite{B-CMP99b} and \eqref{4.6:eq-cover-surgery-a}, by the memberships of the
scale-$(j+1)$ candidates that meet $\square$. If $j<k$, Lemma~\ref{4.6:lem-surgery-surcharge} shows
that every scale-$(j+1)$ candidate meeting $\square$ lies in $\square^{+2u_{j+1}}$. Moreover, since
$u_{j+1}=M\lambda_{j+1}$ and $j+1\le k$,
\begin{equation*}
u_{j+1}=M L^{-(k-j-1)}\le M ,
\end{equation*}
so $\square^{+2u_{j+1}}\subset\square^{+2M}$. If $j=k$, then $\mathcal{D}_{k+1}=\emptyset$, because
the nest has levels at most $k$, and no candidate is tested.

\emph{Row 7.} The profile hypothesis \eqref{4.6:eq-bump-profiles-a}, with the profile of
\cite[(1.118)]{B-CMP95} on the cover of \cite[p.~229]{B-CMP96}, and the gluing hypothesis of
Definition~\ref{4.6:def-bump-profiles} are assumed; under them the bumps of \cite{B-CMP99b} are of
the type treated in Lemma~\ref{4.6:lem-gluing-window}. By that lemma, for each of the families
$\mathcal{F}=\mathcal{D},\mathcal{D}'(\square),\mathcal{D}''(\square)$, the bump $h_C[\mathcal{F}]$
together with $\partial h$ and $\Delta h$ is determined by $C$, its scale, the surgery datum and the
bits of the nest on $C^{+2M}$, and, for an inserted cube $I$, on $(I^{(5)})^{+2M}$, which lies in
$F^{+2M}$ with $F$ the $U$-footprint of the letter or word in which $I$ occurs.

\emph{Row 8.} The cut-off functions $\zeta_{\square}$ and $1-\zeta_{\square}$ have radius $0$ by
\eqref{4.6:eq-bump-profiles-d}; $\widetilde{\square}$ and $1-\widetilde{\square}$ by
\eqref{4.6:eq-candidate-cubes-b}; $\chi_P$ and $\Delta(y)$ are read from the region datum. These
objects are those of \cite{B-CMP99b}.

\emph{Row 9.} The operators $Q'_n$, $Q_n$, $a_n$, $S_n(\partial h)$, $D$, $D^*$, $\Delta_U$,
$\Delta$, $\Delta'$, $J$ are those of \cite[(3.3)--(3.10), (3.15)--(3.19), (3.24),
(3.102)]{B-CMP99b}. At a place $y$ each of them reads the nest only through the layer index $n(y)$,
which has radius $0$ by part~(a) of Lemma~\ref{4.6:lem-membership-tests}. It reads $U$ on the block
or blocks $B^{n(y)}(\cdot)$ around $y$, by \cite[(124), p.~31]{B-CMP98} and \cite{B-CMP99b}, where
these operators are described as semi-local.

\emph{Row 10.} The operators $K'(h)$, $K(h)$ and $P_C(\partial h)$ are those of \cite[(3.88),
(3.100)--(3.104)]{B-CMP99b}; the display \cite[(3.88)]{B-CMP99b} is \cite[(2.38)]{B-CMP96}. They are
composed of the bumps of row~7 and the operators of row~9, so they
read the nest as those rows do; $P_C(\partial h)$ is composed of the generators at $C$ and reads
$C^{(3)}$ by row~4. Their $U$-reads are one block around $\operatorname{supp}\partial h$, and
$\Omega_0(C)$ for $P_C$.

\emph{Row 11.} The tower vertices are those of \cite[(3.117)--(3.120), (3.134)--(3.136),
(3.156)--(3.157), (3.169), (3.184)--(3.185)]{B-CMP99b}, in which the operators $G'$, $R$, $G'RD^*$,
$H$, $H'$, $C'^{(k)}(\Lambda)$, $\widetilde{P}$, $\widetilde{G}'\widetilde{R}$ are substituted by
their expansions. They
read the nest through the tests $b\in\Lambda_j$ and through $\Lambda$, $\Lambda'$, all of radius $0$
by row~2. They read $U$ on blocks, as follows.
\begin{itemize}
\item By \cite{B-CMP99b}, on $\Lambda_j$ the operator $C_j^{(2)}$ coincides with
$C^{(2)}(L_j\eta A)$, a second-order term in the expansion of $Q_j(\eta A)$; this is the term
$C(V_0,A,c)$ of \cite[(122)--(123)]{B-CMP98}, which is indexed by blocks. Its block-locality is that
of Lemma~\ref{4.6:lem-averaging-locality}.
\item By \cite{B-CMP99b}, $D^{(2)}(B)$ is similar to $C^{(2)}(A)$, only restricted to unit blocks.
\item By \cite{B-CMP99b}, the operator $C$ is almost local: a value $(CB)(b)$ depends on $B$
restricted to several blocks surrounding the bond $b$.
\item $\mu(B)$ is given by \cite[(3.169)]{B-CMP99b} on $B(y)$, $y\in\Lambda'$, with operators
$R_y(V)$ of the type of \cite[(58)]{B-CMP98}.
\end{itemize}

\emph{Row 12.} The unexpanded operators $G'$, $P$, $R$, the operator $C$, $G$, $G'RD^*$, $H$, $H_1$,
$(QGQ^*)^{-1}$, $(QG_1Q^*)^{-1}$, the tower operators $\mathfrak{P}$, $\mathfrak{G}$,
$\widetilde{\mathfrak{h}}$, and $H'$, $C'^{(k)}(\Lambda)$, $G_2$, $\widetilde{G}_2$ are read on the
whole torus $\mathbb{T}$, both as to the nest and as to $U$. None of them enters the localized
formulas of \cite{B-CMP99b} itself: each is substituted there by its expansion, and the same
substitution is the one made inside the tower vertices of row~11.

\emph{Row 13.} The scalars $Z(B)$, $Z$, $Z'$, $Z'(\mu)$, $|\det(\Delta\restriction N(Q'))|$, the
wavy determinant $|\det(\widetilde{\Delta}\restriction N(\widetilde{Q}'))|$ of
\cite[(3.160)]{B-CMP99b}, and the Jacobians of the $\delta$-functions of \cite[(3.166),
(3.168)]{B-CMP99b} and of $B=C\overline{B}$, with $C$ the operator $C$ of row~11, are global. They
cancel from \cite[(3.126), (3.129),
(3.130), (3.138), (3.147), (3.153), (3.158), (3.185), (3.186)]{B-CMP99b}, and the final formulas of
\cite{B-CMP99b} are free of them; hence they occur in no localized term. This proves~(i).

\emph{Part (ii).} By the third column, rows 1, 2, 5, 8, 9 and 11 have radius $0$, and rows 3 and 4,
and $P_C$ in row~10, read inside $C^{(3)}\subset C^{(5)}$, or inside $C^{(5)}$ under the other
reading of $Q'$. The remaining entries are the bumps of row~7, which by
Lemma~\ref{4.6:lem-gluing-window} read within $2M$ beyond the support cube and enter $K'(h)$,
$K(h)$ of row~10 through row~7; and the candidate tests of row~6, which enter the anchor data of
row~5 and read at most $2u_{j+1}\le 2M$ beyond the anchor. In the units $u_j=1$, $u_{j+1}=L$ of
Lemma~\ref{4.6:lem-surgery-surcharge}, the margin $2u_{j+1}$ is $2L$.

\emph{Part (iii).} The $U$-reads of the table are the sets $\Omega_0(C)\subset C^{(5)}$ of rows 4
and~10 and the blocks of rows 9, 10 and~11 around points of the relevant supports. The locality in
$U$ stated in \cite{B-CMP99b} places these reads, for a word $\omega$, in $\bigcup_iC_i^{(5)}$, the
union over the cubes $C_i$ at which the letters of $\omega$ are formed; it is used here under the
layer-separation property
\eqref{4.6:eq-letters-words-c}, which keeps the averaging blocks of the levels that would leave
$C^{(5)}$ away from the bump supports of $C$, and keeps the one-block spreads of the operators
flanking $h_I$ inside $I^{(5)}$ through $\operatorname{supp}h_I\subset I^{(3)}$. Without
\eqref{4.6:eq-letters-words-c} the $U$-read set of a word $\omega$ is contained in
$F(\omega)^{+2M}$, and this lies inside the window of the reading radius $r_*$ because
$(I^{(5)})^{+2M}\subset F(\omega)^{+2M}$ for every inserted cube $I$ occurring in $\omega$, by row~7
with $F=F(\omega)$.
\end{proof}

\begin{lemma}[Locality of letters implies locality of words]\label{4.6:lem-word-locality}
Let $\mathfrak{A}$ be a set, and let $\rho$ assign to every $\mathfrak{a}\in\mathfrak{A}$ a closed
subset $\rho(\mathfrak{a})$ of $\mathbb{T}$. Let $\mathfrak{A}^*$ be the free monoid on
$\mathfrak{A}$, with unit the empty word $1$, and extend $\rho$ to $\mathfrak{A}^*$ by
\begin{equation*}
  \rho(uu'):=\rho(u)\cup\rho(u'),\qquad \rho(1):=\emptyset .
\end{equation*}
Let $\mathbb{Z}\langle\langle\mathfrak{A}\rangle\rangle$ be the completion of the free ring
$\mathbb{Z}\langle\mathfrak{A}\rangle$ with respect to word length. Its elements are written as
formal series $\sum_{w\in\mathfrak{A}^*}c_w w$ with $c_w\in\mathbb{Z}$, and the product is induced
by concatenation. For $W\subset\mathbb{T}$ put
\begin{equation*}
  \pi_W\Bigl(\sum_w c_w w\Bigr):=\sum_w c_w\,\mathbf{1}[\rho(w)\subset W]\,w .
\end{equation*}
A \emph{valuation} is a map $v$ from $\mathfrak{A}$ to a unital associative algebra of kernels
that is complete for a norm $\|\cdot\|$, extended multiplicatively to a monoid homomorphism on
$\mathfrak{A}^*$: $v(\mathfrak{a}_1\cdots\mathfrak{a}_\ell)=v(\mathfrak{a}_1)\cdots v(\mathfrak{a}_\ell)$,
and $v(1)$ is the unit of the algebra. Then:
\begin{enumerate}
\item[(i)] $\pi_W$ is a continuous, unital, idempotent algebra endomorphism of
$\mathbb{Z}\langle\langle\mathfrak{A}\rangle\rangle$, and $\pi_{W'}\pi_W=\pi_{W'}$ for
$W'\subset W$.
\item[(ii)] If two valuations $v,v'$ agree on
$\{\mathfrak{a}\in\mathfrak{A}:\rho(\mathfrak{a})\subset W\}$, then $v(w)=v'(w)$ for every word $w$
with $\rho(w)\subset W$. Consequently, for every
$\mathfrak{K}=\sum_w c_w w\in\mathbb{Z}\langle\langle\mathfrak{A}\rangle\rangle$, the word-by-word
evaluations of $\pi_W\mathfrak{K}$ under $v$ and under $v'$ coincide term by term. They also
coincide as operators whenever the sum of $|c_w|\,\|v(w)\|$ over the words $w$ with
$\rho(w)\subset W$ is finite.
\item[(iii)] If $\mathfrak{r}\in\mathbb{Z}\langle\langle\mathfrak{A}\rangle\rangle$ has no constant
term, then $\sum_{n\ge0}\mathfrak{r}^n$ converges in
$\mathbb{Z}\langle\langle\mathfrak{A}\rangle\rangle$. Moreover, for every
$\mathfrak{g}_0\in\mathbb{Z}\langle\langle\mathfrak{A}\rangle\rangle$,
\begin{equation}\label{4.6:eq-word-locality-1}
  \pi_W\Bigl(\mathfrak{g}_0\sum_{n\ge0}\mathfrak{r}^n\Bigr)
  =\pi_W(\mathfrak{g}_0)\sum_{n\ge0}(\pi_W\mathfrak{r})^n .
\end{equation}
\end{enumerate}
\end{lemma}

\begin{proof}
Since $\mathfrak{A}^*$ is free, every word factors uniquely into letters, so the two rules
defining $\rho$ on $\mathfrak{A}^*$ are consistent. Induction on the length $\ell$ then gives
\begin{equation*}
  \rho(\mathfrak{a}_1\cdots\mathfrak{a}_\ell)=\bigcup_{i=1}^{\ell}\rho(\mathfrak{a}_i).
\end{equation*}
Topologically, a sequence in $\mathbb{Z}\langle\langle\mathfrak{A}\rangle\rangle$ converges if
and only if, for every $\ell$, its coefficients on the words of length at most $\ell$ are
eventually constant.

(i) For words $u,u'$ the identity $\rho(uu')=\rho(u)\cup\rho(u')$ gives
\begin{equation}\label{4.6:eq-word-locality-2}
  \mathbf{1}[\rho(uu')\subset W]=\mathbf{1}[\rho(u)\subset W]\,\mathbf{1}[\rho(u')\subset W].
\end{equation}
We apply this word by word. For $a=\sum_u a_u u$ and $b=\sum_{u'} b_{u'} u'$, the coefficient of
$w$ in $ab$ is the finite sum $\sum_{uu'=w}a_ub_{u'}$. By \eqref{4.6:eq-word-locality-2}, the
coefficient of $w$ in $\pi_W(ab)$ is therefore
\begin{equation*}
  \mathbf{1}[\rho(w)\subset W]\sum_{uu'=w}a_ub_{u'}
  =\sum_{uu'=w}\bigl(\mathbf{1}[\rho(u)\subset W]a_u\bigr)
  \bigl(\mathbf{1}[\rho(u')\subset W]b_{u'}\bigr),
\end{equation*}
which is the coefficient of $w$ in $\pi_W(a)\pi_W(b)$. Additivity of $\pi_W$ holds coefficient
by coefficient.

Since $\rho(1)=\emptyset\subset W$, we have $\pi_W(1)=1$.

For $W'\subset W$, the indicators satisfy
$\mathbf{1}[\rho(w)\subset W']\,\mathbf{1}[\rho(w)\subset W]=\mathbf{1}[\rho(w)\subset W']$.
This gives $\pi_{W'}\pi_W=\pi_{W'}$, and the case $W'=W$ is idempotence.

For continuity, note that $\pi_W$ acts on each coefficient separately and does not change the
words. Hence, if the coefficients of a sequence on words of length at most $\ell$ are eventually
constant, so are those of its image.

(ii) By the union formula above, $\rho(w)\subset W$ holds if and only if
$\rho(\mathfrak{a}_i)\subset W$ for every letter $\mathfrak{a}_i$ of $w$. For such a word, each
factor satisfies $v(\mathfrak{a}_i)=v'(\mathfrak{a}_i)$ by hypothesis. A product of equal values
is equal, so
\begin{equation*}
  v(w)=v(\mathfrak{a}_1)\cdots v(\mathfrak{a}_\ell)=v'(\mathfrak{a}_1)\cdots v'(\mathfrak{a}_\ell)=v'(w),
\end{equation*}
and for $w=1$ both sides are the unit of the algebra.

The terms of $\pi_W\mathfrak{K}$ are $c_w w$ with $\rho(w)\subset W$. Their evaluations $c_wv(w)$
and $c_wv'(w)$ therefore agree term by term. If the sum of $|c_w|\,\|v(w)\|$ over these words
is finite, then the sum of $|c_w|\,\|v'(w)\|$ is the same finite number. Both operator series
then converge absolutely in the norm $\|\cdot\|$, hence converge, since the algebra is complete;
they have the same terms, and so they have the same sum.

(iii) Every word occurring in $\mathfrak{r}$ has length at least one, so every word occurring in
$\mathfrak{r}^n$ has length at least $n$. Hence the coefficient of a word of length $\ell$ in
$\mathfrak{r}^n$ vanishes for $n>\ell$. It follows that, for every $\ell$, the coefficients of
the partial sums $\sum_{n\le N}\mathfrak{r}^n$ on words of length at most $\ell$ are constant for
$N\ge\ell$, and so $\sum_{n\ge0}\mathfrak{r}^n$ converges.

Since $\pi_W$ does not change the words, $\pi_W\mathfrak{r}$ also has no constant term, and
$\sum_{n\ge0}(\pi_W\mathfrak{r})^n$ converges by the same argument.

Now apply (i):
\begin{align*}
  \pi_W\Bigl(\mathfrak{g}_0\sum_{n\ge0}\mathfrak{r}^n\Bigr)
  &=\pi_W(\mathfrak{g}_0)\,\pi_W\Bigl(\sum_{n\ge0}\mathfrak{r}^n\Bigr)
  =\pi_W(\mathfrak{g}_0)\sum_{n\ge0}\pi_W(\mathfrak{r}^n)\\
  &=\pi_W(\mathfrak{g}_0)\sum_{n\ge0}(\pi_W\mathfrak{r})^n .
\end{align*}
The first equality is multiplicativity of $\pi_W$, the second is its continuity and additivity,
and the third is multiplicativity again. This is \eqref{4.6:eq-word-locality-1}.

No commutation relation between letters, or between their values, has been used. The kernels
may act on a fibre of colour indices; this fibre enters neither the sets $\rho(w)$ nor the
argument.
\end{proof}

\begin{definition}[Local letters and the supports of their kernels]\label{4.6:def-local-letters}
Let $\mathfrak{a}$ be a candidate letter in the sense of Definition~\ref{4.6:def-letters-words},
with $U$-footprint $F(\mathfrak{a})$, nest-reading domain $\rho(\mathfrak{a})$, valuation
$v_{\mathbf{\Omega}}(\mathfrak{a})$ and admissibility indicator
$\operatorname{adm}_{\mathbf{\Omega}}(\mathfrak{a})\in\{0,1\}$ at a nest $\mathbf{\Omega}$, and let
$\mathfrak{P}$ be the region datum. Write $\mathrm{Rd}(\mathfrak{a})$ for the read set of
$\mathfrak{a}$, the set on which its valuation reads the configuration $U$, and
$U|_{\mathrm{Rd}(\mathfrak{a})}$ for the restriction of $U$ to it. The read set satisfies
$\mathrm{Rd}(\mathfrak{a})\subset F(\mathfrak{a})^{+2M}$ always, and
$\mathrm{Rd}(\mathfrak{a})\subset F(\mathfrak{a})$ under the standing condition for a letter
carrying one label (Definition~\ref{4.6:def-letters-words}).

The letter $\mathfrak{a}$ is \emph{local} if there are a function $\phi_{\mathfrak{a}}$ and one
map $\Phi_{\mathfrak{a}}$, not depending on the label, for which the two conditions
$(\ell\text{-ind})$ and $(\ell\text{-val})$ in the first two lines of
\eqref{4.6:eq-local-letters-a} hold; the last line of the display fixes notation for the kernel
supports described below, the factors $m_L$, $E_L$, $m_R$, $E_R$ being those named there:
\begin{equation}\label{4.6:eq-local-letters-a}
\begin{aligned}
&(\ell\text{-ind})\quad
\operatorname{adm}_{\mathbf{\Omega}}(\mathfrak{a})
=\phi_{\mathfrak{a}}\bigl(\operatorname{tr}_{\rho(\mathfrak{a})}(\mathbf{\Omega};\mathfrak{P})\bigr)
\quad\text{for every nest }\mathbf{\Omega},\\
&(\ell\text{-val})\quad
v_{\mathbf{\Omega}}(\mathfrak{a})
=\Phi_{\mathfrak{a}}\bigl(U|_{\mathrm{Rd}(\mathfrak{a})},\,
\operatorname{tr}_{\rho(\mathfrak{a})}(\mathbf{\Omega};\mathfrak{P})\bigr)
\quad\text{if }\operatorname{adm}_{\mathbf{\Omega}}(\mathfrak{a})=1,\\
&\widehat{X}_L(\mathfrak{a}):=(\operatorname{supp} m_L)^{+\sigma(E_L)},\qquad
\widehat{X}_R(\mathfrak{a}):=(\operatorname{supp} m_R)^{+\sigma(E_R)}.
\end{aligned}
\end{equation}
Condition $(\ell\text{-ind})$ says that the map
$\mathbf{\Omega}\mapsto\operatorname{adm}_{\mathbf{\Omega}}(\mathfrak{a})$ factors through the trace
$\operatorname{tr}_{\rho(\mathfrak{a})}(\mathbf{\Omega};\mathfrak{P})$; condition
$(\ell\text{-val})$ says that the single label-free map $\Phi_{\mathfrak{a}}$ gives the valuation
whenever $\operatorname{adm}_{\mathbf{\Omega}}(\mathfrak{a})=1$, the two sides being read as partial
functions of the configuration on the read set $\mathrm{Rd}(\mathfrak{a})$.

\emph{Kernel supports.} With $\widehat{X}_L(\mathfrak{a})$ and $\widehat{X}_R(\mathfrak{a})$ as in
the last line of \eqref{4.6:eq-local-letters-a}, fix one nest $\mathbf{\Omega}$. Suppose that the
leftmost factor of
$\mathfrak{a}$ is a multiplier $m_L$ (a bump, a cut-off $\zeta_{\square}$, or a characteristic
function $\chi$), preceded by at most one explicit operator $E_L$ of spread $\sigma(E_L)$, and
symmetrically that its rightmost factor is a multiplier $m_R$ followed by at most one explicit
operator $E_R$ of spread $\sigma(E_R)$; when the operator is absent its spread is taken to be $0$,
and $A^{+0}=A$.
The rows of the kernel of $\mathfrak{a}$ lie in $\widehat{X}_L(\mathfrak{a})$ and its columns lie
in $\widehat{X}_R(\mathfrak{a})$. These sets serve the bookkeeping of kernel supports at one nest;
they are not part of the definition of locality.

By \eqref{4.6:eq-bump-profiles-c} of Definition~\ref{4.6:def-bump-profiles}, the support of a
multiplier satisfies $\operatorname{supp} m\subset C$ for a candidate bump at the candidate cube $C$,
$\operatorname{supp} m\subset I^{(5)}$ for a bump inserted at the cube $I$, and
$\operatorname{supp}\zeta_{\square}\subset\square^{(2)}$; and by the spreads fixed in
Definition~\ref{4.6:def-letters-words}, $\sigma(E)$ is at most one lattice step plus one averaging
block of the level at which $E$ acts.
\end{definition}

\begin{proof}[Proof that the rows and columns lie in $\widehat{X}_L(\mathfrak{a})$ and $\widehat{X}_R(\mathfrak{a})$]
Write the operator of $\mathfrak{a}$ at the fixed nest as $E_L\,m_L\,\Pi_R$, where $\Pi_R$ is the
product of the remaining factors and $E_L$ is absent if there is no explicit operator on the left.

A multiplier kills the rows off its support: the kernel of $m_L \Pi_R$ is
$(m_L\Pi_R)(x,y)=m_L(x)\,\Pi_R(x,y)$, which vanishes unless $x\in\operatorname{supp} m_L$. If $E_L$ is
absent this is the claim, since then $\widehat{X}_L(\mathfrak{a})=\operatorname{supp} m_L$.

If $E_L$ is present, an operator of spread $\sigma=\sigma(E_L)$ applied on the outside moves the
rows by at most $\sigma$: its kernel $E_L(x,x')$ vanishes unless $x\in\{x'\}^{+\sigma}$. In
\begin{equation*}
(E_L m_L \Pi_R)(x,y)=\sum_{x'}E_L(x,x')\,m_L(x')\,\Pi_R(x',y)
\end{equation*}
a term is non-zero only if $x'\in\operatorname{supp} m_L$ and $x\in\{x'\}^{+\sigma}$; hence the
left side vanishes unless $x\in(\operatorname{supp} m_L)^{+\sigma(E_L)}=\widehat{X}_L(\mathfrak{a})$.

The columns are treated symmetrically. Write the operator as $\Pi_L\,m_R\,E_R$. Its kernel is
\begin{equation*}
(\Pi_Lm_RE_R)(x,y)=\sum_{y'}\Pi_L(x,y')\,m_R(y')\,E_R(y',y),
\end{equation*}
and a term is non-zero only if $y'\in\operatorname{supp} m_R$ and $y\in\{y'\}^{+\sigma(E_R)}$.
Hence the kernel vanishes unless $y\in\widehat{X}_R(\mathfrak{a})$; without $E_R$ the factor
$m_R(y)$ gives $y\in\operatorname{supp} m_R=\widehat{X}_R(\mathfrak{a})$ directly.
\end{proof}

\begin{lemma}[The closure of the local primitives is local]\label{4.6:lem-closure-local}
Assume the following two hypotheses:
\begin{itemize}
\item the reading of the letters of the alphabet, and its version for the propagator
towers, \eqref{4.6:eq-tower-reading-a};
\item the type hypotheses $(\mathrm{G\text{-}loc})$ and $(\mathrm{G\text{-}supp})$ on the glued
bumps, displayed in \eqref{4.6:eq-bump-profiles-c}.
\end{itemize}
Then the following hold.
\begin{enumerate}
\item[(a)] Every row of the reading table of Lemma~\ref{4.6:lem-reading-table} is a local letter
in the sense of Definition~\ref{4.6:def-local-letters}, that is, it satisfies both conditions of
\eqref{4.6:eq-local-letters-a}.
\item[(b)] The closure $\Theta(\mathfrak{J})$ of the ingredient list $\mathfrak{J}$ of
Definition~\ref{4.6:def-letters-words} consists of local letters. For a letter of
$\Theta(\mathfrak{J})$ built from letters $\mathfrak{a},\mathfrak{a}',\dots$, its reading domain
$\rho$, its footprint $F$ and its read set are the unions of those of
$\mathfrak{a},\mathfrak{a}',\dots$. Its evaluation map $\Phi$ is the corresponding label-free
expression in $\Phi_{\mathfrak{a}},\Phi_{\mathfrak{a}'},\dots$.
\item[(c)] The kernel of a word, whether a product or a tree contraction, has its rows in the row
support $\widehat{X}_L$ of its leftmost letter and its columns in the column support
$\widehat{X}_R$ of its rightmost letter.
\end{enumerate}
\end{lemma}

\begin{proof}
We argue by induction on the construction of $\Theta(\mathfrak{J})$ from $\mathfrak{J}$ by the
operations (o1)--(o6) of Definition~\ref{4.6:def-letters-words}. The base of the induction is
clause~(a). For rows~1 to~11 of the table of Lemma~\ref{4.6:lem-reading-table} it follows from
Lemma~\ref{4.6:lem-membership-tests}, Lemma~\ref{4.6:lem-sequence-reach},
Proposition~\ref{4.6:prop-generator-equality}, Lemma~\ref{4.6:lem-gluing-window}, the reading of
the cut-offs $\zeta_{\square}$ (Definition~\ref{4.6:def-cover-surgery}), the reading of the double
cubes $\widetilde{\square}$, the block-locality of the averaging operations
(Lemma~\ref{4.6:lem-averaging-locality}), and the tie-break rule of
Definition~\ref{4.6:def-cover-surgery}. Row~12 does not occur once the substitution made in that
table is carried out, and row~13 does not occur. The hypotheses $(\mathrm{G\text{-}loc})$ and
$(\mathrm{G\text{-}supp})$ are those under which Lemma~\ref{4.6:lem-gluing-window} applies to the
bump rows; the reading of the letters of the alphabet, together with its version for the propagator
towers \eqref{4.6:eq-tower-reading-a}, is the hypothesis on the letters of the alphabet.
For the inductive step we show that each operation, applied to local
letters, produces a local letter with the reading domain, footprint, read set and evaluation map
described in~(b) and with the kernel supports described in~(c).

\emph{Operation} (o1). Let $\mathfrak{a},\mathfrak{a}'$ be local letters with reading domains
$\rho=\rho(\mathfrak{a})$ and $\rho'=\rho(\mathfrak{a}')$. The trace of a nest with region datum on
$\rho\cup\rho'$ determines, by restriction, its traces
$\operatorname{tr}_{\rho}(\mathbf{\Omega};\mathfrak{P})$ and
$\operatorname{tr}_{\rho'}(\mathbf{\Omega};\mathfrak{P})$. Hence, by the first condition of
\eqref{4.6:eq-local-letters-a} applied to $\mathfrak{a}$ and to $\mathfrak{a}'$, the admissibility
data of both factors are functions of the trace on $\rho\cup\rho'$. By the second condition, applied
to each factor, their values are functions of the trace on $\rho\cup\rho'$ and of $U$ on the union of
the read sets. The product is therefore a local letter with reading domain $\rho\cup\rho'$, footprint
$F(\mathfrak{a})\cup F(\mathfrak{a}')$, read set the union of the read sets, and evaluation map
\begin{equation*}
\Phi:=\Phi_{\mathfrak{a}}\cdot\Phi_{\mathfrak{a}'},
\end{equation*}
which is label-free because both factors are. The kernel supports compose: the rows of
$v(\mathfrak{a})v(\mathfrak{a}')$ lie among the rows of $v(\mathfrak{a})$, and its columns lie among
the columns of $v(\mathfrak{a}')$. Iterating over the letters of a word gives~(c) for products.

\emph{Operation} (o2). A multiplier $m$ is a letter from rows~7 and~8 of the table of
Lemma~\ref{4.6:lem-reading-table}. It is a function of its cube and of the trace of the nest on its
read set. That read set is joined to the reading domain $\rho$. For a bump this is covered by
Definition~\ref{4.6:def-letters-words}: there $\rho$ contains $C^{+2M}$ for a bump attached to a
cube $C$, and $(I^{(5)})^{+2M}$ for a bump attached to an inserted cube $I$. Hence, for a local
letter $\mathfrak{a}$, the letter $m\,\mathfrak{a}$ is local, with reading domain the union of
$\rho(\mathfrak{a})$ and the read set of $m$, and footprint the union of the two footprints. It is
evaluated by the product of the evaluation map of
$m$ with $\Phi_{\mathfrak{a}}$. The kernel $m\,v(\mathfrak{a})$ has its rows in
$\operatorname{supp} m$.

\emph{Operations} (o3) \emph{and} (o4). A difference or a commutator of two local letters is a
function of the data read by the two letters. It is therefore a local letter whose footprint is the
union of the two footprints, with evaluation map the difference, resp.\ the commutator, of the two
evaluation maps. Its rows and its columns lie in the union of the corresponding supports of the two
letters.

\emph{Operation} (o5). Let $E$ be a label-free multilinear expression in local letters, and let $w$
be a word of local letters. Then $E[v(w)]$ is obtained by the operations (o1) and (o2), so the two
preceding cases apply to it. The family of such letters is indexed label-freely by the pair formed
by the symbol of $E$ and the word~$w$.

\emph{Operation} (o6). This operation is the operation (o1), treated above.

Each operation thus preserves the two conditions of \eqref{4.6:eq-local-letters-a}, with reading
domain, footprint and read set the unions and evaluation map the corresponding label-free
expression. Together with the base case this proves~(b). The kernel statements for (o1) and (o2)
give~(c) for products and, since (o6) is an instance of (o1), the same holds for tree contractions.
\end{proof}

\begin{lemma}[The formal expansions contain only local letters]\label{4.6:lem-formal-expansions}
Let $\mathbf{\Omega}$ be a nest, $\mathcal{D}(\mathbf{\Omega})$ its mixed-size cover, and let
$h_\square$, $\square\in\mathcal{D}(\mathbf{\Omega})$, be the bumps, so that
$\sum_{\square}h_\square^2=1$ \cite[(2.36)]{B-CMP96}. Write $\mathfrak{g}'_0$, $\mathfrak{c}_0$ and
$\mathfrak{g}_0$ for the formal block sums over $\square\in\mathcal{D}(\mathbf{\Omega})$ whose summands
are $h_\square G'_\square h_\square$, $h_\square C_\square h_\square$ and $h_\square G_\square h_\square$.
These are the formal counterparts of $G'_0$ of \cite[(2.37)]{B-CMP96}, of $C$ of
\cite[(2.70)]{B-CMP96}, and of $G_0$ of \cite[(3.106)]{B-CMP99b}. Write $\mathfrak{r}'$ for the formal
counterpart of the remainder $\sum_\square K(h_\square)G'_\square h_\square$ of \cite[(2.38)]{B-CMP96},
and $\mathfrak{r}_{95}$ for the formal remainder of \cite[(3.95)--(3.96)]{B-CMP99b}. For a family
$\mathcal{F}$ of cubes, $\mathfrak{K}'[\mathcal{F}]$ denotes the series $\mathfrak{K}'$ below formed on
$\mathcal{F}$ in place of $\mathcal{D}(\mathbf{\Omega})$. For a cube $\square$,
$\mathcal{D}'_\square$ denotes the surgered family $\mathcal{D}'$ attached to $\square$.

Define in $\mathbb{Z}\langle\langle\mathfrak{A}^{\mathrm{cand}}\rangle\rangle$, following
Ba{\l}aban's construction, the following five series.
\begin{enumerate}
\item The expansion of $G'$ \cite[(3.87)--(3.90)]{B-CMP99b}, \cite[(2.37)--(2.38), (2.50)]{B-CMP96}:
\begin{equation*}
\mathfrak{K}':=\sum_{n\ge0}\mathfrak{g}'_0(\mathfrak{r}')^n .
\end{equation*}
\item The expansion of $(Q'G'^2Q'^*)^{-1}$ \cite[(3.95)--(3.96)]{B-CMP99b}:
\begin{equation*}
\mathfrak{K}_C:=\sum_{n\ge0}\mathfrak{c}_0(\mathfrak{r}_{95})^n .
\end{equation*}
It is formed after two operations have been carried out. First, the operation labelled (o5) among
the substitution operations on words of this chapter, namely the substitution of the series
$\mathfrak{K}'[\mathcal{D}'_\square]$ for a factor $G'$, is applied to every factor $G'$ of the first
sum of \cite[(3.95)--(3.96)]{B-CMP99b}, and of the first piece of its second sum. Second, the
regroupings of terms set out in this chapter for these sums are carried out.
\item The expansion of $P$, formed on the families of cubes prescribed in \cite{B-CMP99b}:
\begin{equation*}
\mathfrak{K}_P:=\mathfrak{K}'\,Q'^*\,\mathfrak{K}_C\,Q'\,\mathfrak{K}' .
\end{equation*}
\item The expansion of $R$. Put
$\mathfrak{K}_R:=\sum_{\square}\bigl[\mathfrak{k}_1+\mathfrak{k}_2+\mathfrak{k}_3+\mathfrak{k}_4\bigr](\square)$, where
\begin{align}
\mathfrak{k}_1&:=K(h_\square)G_\square h_\square,\qquad
\mathfrak{k}_2:=(1-\zeta_\square)D\mathfrak{K}_PD^*h_\square G_\square h_\square,
\label{4.6:eq-formal-expansions-1}\\
\mathfrak{k}_3&:=\zeta_\square D\mathfrak{K}_PD^*h_\square G_\square h_\square
-\zeta_\square DP_\square D^*h_\square G_\square h_\square,
\label{4.6:eq-formal-expansions-2}\\
\mathfrak{k}_4&:=\zeta_\square P_\square(\partial h_\square)G_\square h_\square.
\label{4.6:eq-formal-expansions-3}
\end{align}
The element $\mathfrak{k}_3$ contains a word, which we call its \emph{exceptional word}, that is
transformed in \cite{B-CMP99b} before it enters the expansion; in $\mathfrak{k}_3$ this word is taken
in its transformed form. The display of \cite{B-CMP99b} in which this transformation is written is
read here from the sentences that precede and follow it; the steps so read are listed in (c).
\item The expansion of $G$ \cite[(3.106)--(3.107)]{B-CMP99b}:
\begin{equation*}
\mathfrak{K}_G:=\mathfrak{g}_0\sum_{n\ge0}(\mathfrak{K}_R)^n .
\end{equation*}
\end{enumerate}
Then the following hold.
\begin{enumerate}
\item[(a)] Every letter of $\mathfrak{K}'$, $\mathfrak{K}_C$, $\mathfrak{K}_P$, $\mathfrak{K}_R$,
$\mathfrak{K}_G$ lies in $\Theta(\mathfrak{J})$. No letter of any word of the five series is one of
the unexpanded operators $G'$, $C$, $P$, $R$, $G$.
\item[(b)] The global projection $P$ of \eqref{4.6:eq-commutator-identity-a} occurs in
\cite[(3.105)]{B-CMP99b} in exactly two places: in $\mathfrak{k}_2$, and in the first piece of
$\mathfrak{k}_3$.
\item[(c)] In \cite{B-CMP99b} the exceptional word of $\mathfrak{k}_3$ is transformed only by the
following four steps, all inside $\Theta(\mathfrak{J})$:
\begin{itemize}
\item the operation labelled (o2) among the substitution operations on words of this chapter, with
the factor $h''_{\square_0}-\zeta_\square$, where $\square_0$ is the inserted cube and
$h''_{\square_0}$ is its bump in the surgered family $\mathcal{D}''$;
\item the operation labelled (o4);
\item the operation labelled (o3);
\item the local identity of clause (6) of Lemma~\ref{4.6:lem-commutator-identity}.
\end{itemize}
\item[(d)] Every word of $\mathfrak{K}'$, $\mathfrak{K}_C$ and $\mathfrak{K}_G$ begins with a bump
$h_\square$ of a member $\square$ of the unsurgered cover $\mathcal{D}(\mathbf{\Omega})$, and ends with
one. The surgered families $\mathcal{D}'$, $\mathcal{D}''$ occur only inside substituted factors.
\end{enumerate}
\end{lemma}

\begin{proof}
\emph{(b).} By clause (2) of Lemma~\ref{4.6:lem-commutator-identity}, for each cube $\square$, the
operator $\Delta_a h_\square G_\square h_\square$ equals
$h_\square(\mathfrak{L}-DP_\square D^*)G_\square h_\square$
minus the four terms
\begin{equation*}
\begin{gathered}
K(h_\square)G_\square h_\square,\qquad
(1-\zeta_\square)DPD^*h_\square G_\square h_\square,\\
\zeta_\square(DPD^*-DP_\square D^*)h_\square G_\square h_\square,\qquad
\zeta_\square P_\square(\partial h_\square)G_\square h_\square .
\end{gathered}
\end{equation*}
By clauses (4) and (5) of Lemma~\ref{4.6:lem-commutator-identity}, the four sums of
\cite[(3.105)]{B-CMP99b} are, term by term, the sums over $\square\in\mathcal{D}(\mathbf{\Omega})$
of the four terms above. The hypotheses of clause (4) concern only the leading term
$h_\square(\mathfrak{L}-DP_\square D^*)G_\square h_\square$ and play no role in locating the global
$P$. We read off where the
global $P$ stands in these four terms. The first term contains only $K(h_\square)$, $G_\square$ and
$h_\square$. The fourth contains only the local $P_\square$. The third is the difference of a term
with the global $P$ and a term with the local $P_\square$. The second contains the global $P$. Hence
the global $P$ occurs exactly in the second sum and in the first piece of the third sum. In the
definitions \eqref{4.6:eq-formal-expansions-1}--\eqref{4.6:eq-formal-expansions-3} these are exactly
the places of $\mathfrak{k}_2$ and of the first piece of $\mathfrak{k}_3$ where $P$ has been replaced
by $\mathfrak{K}_P$. By Corollary~\ref{4.6:cor-two-substitutions}, both places are substituted in
\cite{B-CMP99b}.

\emph{(c).} We follow the passage of \cite{B-CMP99b} after \cite[(3.105)]{B-CMP99b} in which the
exceptional word of $\mathfrak{k}_3$ is treated, sentence by sentence. Its first step is the
operation (o2), with the factor $h''_{\square_0}-\zeta_\square$. Its next step is the operation (o4).
It then applies the operation (o3). It closes with the identity of clause (6) of
Lemma~\ref{4.6:lem-commutator-identity}, which is an identity between local expressions. No other
step acts on this word.

By Lemma~\ref{4.6:lem-closure-local}, $\Theta(\mathfrak{J})$ consists of local letters. That lemma
is proved as an implication from the reading \eqref{4.6:eq-tower-reading-a} of the propagator towers
and the hypotheses $(\mathrm{G\text{-}loc})$, $(\mathrm{G\text{-}supp})$ on the bumps. The factor
$h''_{\square_0}-\zeta_\square$, and the words produced by (o4), (o3) and clause (6), are formed from
letters of $\Theta(\mathfrak{J})$. Hence each step of the transformation stays inside
$\Theta(\mathfrak{J})$.

\emph{(a).} The letters entering the five definitions are of three kinds:
\begin{itemize}
\item the localized generators at cubes, $G'_\square$, $C_\square$, $G_\square$, $P_\square$,
$K(h_\square)$ and $P_\square(\partial h_\square)$;
\item the bumps and cut-offs $h_\square$, $\zeta_\square$, together with $D$, $D^*$, $Q'$ and $Q'^*$;
\item the letters produced by the operation (o5), by the regroupings of terms set out in this
chapter, and by the transformation treated in (c).
\end{itemize}
By Lemma~\ref{4.6:lem-closure-local}, every ingredient is a local letter and $\Theta(\mathfrak{J})$
consists of local letters. The letters above are ingredients, or are obtained from ingredients by
the substitutions and regroupings of the construction, which are the substitutions treated in
Corollary~\ref{4.6:cor-two-substitutions}, the regroupings named above and the transformation
treated in (c). They therefore lie in $\Theta(\mathfrak{J})$.

It remains to check that no letter of any word is an unexpanded operator. Every factor $G'$ is
replaced by
$\mathfrak{K}'$, or by $\mathfrak{K}'[\mathcal{D}'_\square]$ where the operation (o5) prescribes
this. The factor $(Q'G'^2Q'^*)^{-1}$ is replaced by $\mathfrak{K}_C$. By (b), the global $P$ occurs
only in the two places where it has been replaced by $\mathfrak{K}_P$. The operator $R$ of
\cite[(3.105)]{B-CMP99b} is replaced by $\mathfrak{K}_R$. The operator
$G=G_0(1-R)^{-1}$ of \cite[(3.106)]{B-CMP99b} is replaced by $\mathfrak{K}_G$. By
Corollary~\ref{4.6:cor-two-substitutions}, no global $P$, $G'$ or $C$ survives in the expansion of
$R$, nor in that of $G$. Hence no letter of any word of the five series is one of the unexpanded
operators $G'$, $C$, $P$, $R$, $G$.

\emph{(d).} We inspect the definitions.

Words begin with a bump. A word of $\mathfrak{K}'$ of order $n$ is a summand of $\mathfrak{g}'_0$
followed by $n$ summands of $\mathfrak{r}'$. Its leftmost factor is therefore
$h_\square G'_\square h_\square$ with $\square\in\mathcal{D}(\mathbf{\Omega})$, and it begins with
$h_\square$. In the same way, a word of $\mathfrak{K}_C$ begins with $h_\square C_\square h_\square$,
which is $R'_0(\square)=h_\square C_\square h_\square$ in the notation of \cite{B-CMP99b}. A word of
$\mathfrak{K}_G$ begins with $h_\square G_\square h_\square$. In all three cases
$\square\in\mathcal{D}(\mathbf{\Omega})$, the unsurgered cover over which $\mathfrak{g}'_0$,
$\mathfrak{c}_0$ and $\mathfrak{g}_0$ are summed.

Words end with a bump. For $n=0$ the word is a summand of $\mathfrak{g}'_0$, $\mathfrak{c}_0$ or
$\mathfrak{g}_0$, and it ends with $h_\square$. For $n\ge1$ the word ends with a summand of
$\mathfrak{r}'$, of $\mathfrak{r}_{95}$ or of $\mathfrak{K}_R$ respectively. Each summand of
$\mathfrak{r}'$ ends with $G'_\square h_\square$. Each summand of $\mathfrak{r}_{95}$ ends with
$C_\square h_\square$. Each summand of $\mathfrak{K}_R$ ends with $G_\square h_\square$, by
\eqref{4.6:eq-formal-expansions-1}--\eqref{4.6:eq-formal-expansions-3}. In every case
$\square\in\mathcal{D}(\mathbf{\Omega})$. These are words of the form
\begin{equation*}
h_{\square_0}(\cdots)h_{\square_n},\qquad \square_i\in\mathcal{D}(\mathbf{\Omega}),
\end{equation*}
the form of the summands of \cite[(2.50)]{B-CMP96} and \cite[(1.123)]{B-CMP95}, where the cover is
written $\mathcal{D}$.

Surgered families. The family $\mathcal{D}'_\square$ enters only through the factors
$\mathfrak{K}'[\mathcal{D}'_\square]$ substituted by (o5) inside the sums defining $\mathfrak{K}_C$.
The families prescribed for $\mathfrak{K}_P$ enter only through $\mathfrak{K}_P$, which is
substituted inside $\mathfrak{k}_2$ and $\mathfrak{k}_3$. The family $\mathcal{D}''$ enters only
through the factor $h''_{\square_0}-\zeta_\square$ of the transformed exceptional word of
$\mathfrak{k}_3$, hence inside a summand of $\mathfrak{K}_R$; such a summand stands in a word of
$\mathfrak{K}_G$ of order $n\ge1$ after the leading factor $h_\square G_\square h_\square$, and it
still ends with $G_\square h_\square$, $\square\in\mathcal{D}(\mathbf{\Omega})$. In none of these
cases does a surgered family furnish the first or the last bump of a word of $\mathfrak{K}'$,
$\mathfrak{K}_C$ or $\mathfrak{K}_G$.
\end{proof}

\begin{definition}[The reading of the propagator towers]\label{4.6:def-tower-reading}
In this definition, and wherever it is invoked, the following letters denote the objects of
\cite{B-CMP99b} that bear these names, in Ba{\l}aban's notation there, and not the glossary
symbols of the same shape:
\begin{gather*}
H,\ H_1,\ H',\ G,\ G',\ G_1,\ G_2,\ \widetilde{G}_2,\ \Delta'_\pi,\ \Delta^{(2)},\
\widetilde{\Delta}^{(2)},\ D,\ D^*,\ D_1,\ D^{(2)},\ Q,\ Q^*,\ R,\ I,\\
\Lambda,\ C,\ C^{(2)},\ C'^{(k)}(\Lambda),\ C^{(k)}(\Lambda),\ \widehat{C}^{(k)}(\Lambda),\
\mu,\ \mu(B),\ \lambda(A),\ \mathfrak{P},\ \mathfrak{G},\ \mathfrak{h},\ \widetilde{\mathfrak{h}}.
\end{gather*}
We write $\mathbf{\Omega}$ for the nest and $\widetilde{\mathbf{\Omega}}$ for the wavy nest of
\cite{B-CMP99b}. Letters, words, atoms $K_\omega$ and the closure $\Theta$ of the ingredient list
$\mathfrak{J}$ are those of Definition~\ref{4.6:def-letters-words}; the formal expansions
$\mathfrak{K}'$, $\mathfrak{K}_C$, $\mathfrak{K}_G$ are those of
Lemma~\ref{4.6:lem-formal-expansions}.

The \emph{tower ingredients} are the objects of the first and of the second group below.
\begin{itemize}
\item[] \textit{First group.} The operator $H$ of \cite[(3.109)--(3.112),
(3.121)--(3.126)]{B-CMP99b}; the operator $\Delta'_\pi$ of \cite[(3.117)--(3.120)]{B-CMP99b};
the operator $G$ of \cite[(3.130)]{B-CMP99b}; the quadratic remainder $C^{(2)}$ of
\cite{B-CMP99b}, which is the term written $C(V_0,A,c)$ in \cite[(122)--(123)]{B-CMP98}, in the
notation there, indexed by the blocks $c$; the operator $G_1$ of \cite[(3.128), (3.138)]{B-CMP99b}; the
operators $\Delta^{(2)}$ and
$\widetilde{\Delta}^{(2)}$ of \cite[(3.135)--(3.136)]{B-CMP99b}; the operator $H_1$ of
\cite[(3.129)]{B-CMP99b}; the inverses $(QGQ^*)^{-1}$ and $(QG_1Q^*)^{-1}$ of \cite{B-CMP99b};
the operator $\mathfrak{P}$ of \cite[(3.147)]{B-CMP99b}; the operator $\mathfrak{G}$ of
\cite[(3.153)]{B-CMP99b}.
\item[] \textit{Second group.} The datum $\Lambda$, the operator $D^{(2)}$, the quadratic form of
\cite[(3.156)]{B-CMP99b} on the subspace on which it is considered there, the parametrisation
$C$ and the wavy nest $\widetilde{\mathbf{\Omega}}$, all as in \cite{B-CMP99b}; the operator
$H'$ of \cite[(3.162)--(3.164)]{B-CMP99b}; the operator $\mu(B)$ of \cite[(3.169)]{B-CMP99b};
the operators $\mathfrak{h}$, $\widetilde{\mathfrak{h}}$ and $C'^{(k)}(\Lambda)$ of
\cite[(3.171)--(3.176)]{B-CMP99b}; $\lambda(A)$ of \cite[(3.184)]{B-CMP99b}; $\widetilde{G}_2$
and $G_2$ of \cite[(3.186)]{B-CMP99b}; $C^{(k)}(\Lambda)$ of \cite[(3.185)]{B-CMP99b};
$\widehat{C}^{(k)}(\Lambda)$ of \cite[(3.158)]{B-CMP99b}.
\end{itemize}
A \emph{tower atom} is an atom $K_\omega$ of a tower ingredient after every global operator in it
has been replaced as prescribed in part~(1)(d) below.

The \emph{reading of the propagator towers} is the conjunction of the following three statements.
\begin{itemize}
\item[(1)] \textit{Classification.} Every constituent of the formulas
\cite[(3.109)--(3.186)]{B-CMP99b} for the tower ingredients is of one of the following kinds:
\begin{itemize}
\item[($\emptyset$)] it does not depend on the nest, and it is local in the configuration $U$;
\item[(a)] it reads the nest through tests of levels and of points of the nest;
\item[(b)] it reads the nest through the Dirichlet domain $\Omega_0$;
\item[(c)] it reads the nest through the region $\Lambda$ at radius zero, by way of the
block-local averaging objects of \cite{B-CMP98} and of the companion paper to which
\cite{B-CMP99b} refers for them;
\item[(d)] it is a global operator of $\mathbf{\Omega}$ or of $\widetilde{\mathbf{\Omega}}$; each
such operator is replaced by its formal expansion of Lemma~\ref{4.6:lem-formal-expansions};
\item[(s)] it is a scalar that cancels, namely the thirteenth entry of the reading table of
Lemma~\ref{4.6:lem-reading-table}.
\end{itemize}
The list ($\emptyset$), (a)--(d), (s) is exhaustive over the loci of \cite{B-CMP99b} cited in
the two groups above: at each of them, no constituent depends on the nest in any other way.
\item[(2)] \textit{Identities used before expansion.} The identities that \cite{B-CMP99b} uses
among the operators above, namely
\begin{equation}\label{4.6:eq-tower-reading-1}
RD^*GQ^*=0,\qquad\text{hence}\qquad QGDR=0
\end{equation}
of \cite[(3.124)--(3.125)]{B-CMP99b}, the chain
\begin{equation}\label{4.6:eq-tower-reading-2}
QG_1DR=QDG'R=D_1QG'R=0
\end{equation}
of \cite{B-CMP99b}, the identities \cite[(3.151)--(3.152)]{B-CMP99b}, and the identities
\cite[(3.172)--(3.182)]{B-CMP99b}, together with the Gaussian-integral and Lagrange
computations of \cite{B-CMP99b} that lead to them, are identities among the global operators of
a single nest. They are used only to simplify the global formulas \cite[(3.126), (3.129),
(3.147), (3.153), (3.185)]{B-CMP99b} before any operator is expanded, and they impose no
condition on the localized terms. After the expansion, the only identity imposed on letters is
the local identity \cite[(3.25)]{B-CMP99b}, which is clause (6) of
Lemma~\ref{4.6:lem-commutator-identity}.
\item[(3)] \textit{Residual readings.} The reading includes the following interpretations of
\cite{B-CMP99b}:
\begin{itemize}
\item[--] some of the displays listed above enter only through the type of each operator in
them, that is, through the nest from which the operator is built and the kind in part~(1) to
which it belongs; no other feature of their form is used;
\item[--] the inverses $(QGQ^*)^{-1}$ and $(QG_1Q^*)^{-1}$ are read by analogy with the
preceding operators of \cite{B-CMP99b};
\item[--] for each operator of \cite[(3.160)--(3.186)]{B-CMP99b}, the reading allows its
attribution to the nest $\mathbf{\Omega}$ or to the wavy nest $\widetilde{\mathbf{\Omega}}$ to be
made either way, and which one is made is immaterial below;
\item[--] the quadratic form of \cite[(3.156)]{B-CMP99b} is taken on the subspace on which
\cite{B-CMP99b} considers it; the region $\Lambda$ enters at radius zero through block-local
objects as in kind~(c); and $\mu(B)$ is block-local, of the contour type of the operators
$R_{0,y}$ of \cite[(58), (77)]{B-CMP98} (compare Lemma~\ref{4.6:lem-averaging-locality});
\item[--] the residual convention on the averaging operators under which
Lemma~\ref{4.6:lem-reading-table} is stated is in force.
\end{itemize}
\end{itemize}
\end{definition}

Under the reading of the propagator towers (Definition~\ref{4.6:def-tower-reading}), every tower
atom $K_\omega$ has one of the forms
\begin{equation}\label{4.6:eq-tower-reading-a}
K_\omega\in\bigl\{\,w,\quad (I+D\mu)Q\,w\,Q^*(I+\mu^*D^*),\quad D\,w\,D^*\,\bigr\},
\end{equation}
where $w$ is a word, that is, a finite tree, over letters of the closure $\Theta$ for
$\mathbf{\Omega}$ and for $\widetilde{\mathbf{\Omega}}$ and vertices of radius zero from the
closed list that forms the eleventh entry of the reading table of
Lemma~\ref{4.6:lem-reading-table}, whose outermost
expansion factors are $\mathfrak{K}'$-, $\mathfrak{K}_C$- or $\mathfrak{K}_G$-words; the
second form is the explicit vertex of \cite[(3.185)]{B-CMP99b}. Consequently
Lemma~\ref{4.6:lem-closure-local} applies to every tower atom.

\begin{proof}[Proof of the form \eqref{4.6:eq-tower-reading-a} under the reading]
Let $K_\omega$ be a tower atom, arising from one of the tower ingredients.

First, no global operator survives in $K_\omega$. By part~(1), each constituent of that
ingredient is of one of the kinds ($\emptyset$), (a), (b), (c), (d), (s), and there is no further
way in which it reads the nest. The constituents of kind (d) are global operators of
$\mathbf{\Omega}$ or of
$\widetilde{\mathbf{\Omega}}$, and by the definition of a tower atom each of them has been
replaced by its formal expansion of Lemma~\ref{4.6:lem-formal-expansions}. The constituents of
kind (s) are scalars that cancel, so they do not occur in $K_\omega$. By the third item of
part~(3), for an operator of \cite[(3.160)--(3.186)]{B-CMP99b},
whether it is counted as an operator of
$\mathbf{\Omega}$ or of $\widetilde{\mathbf{\Omega}}$ does not affect this: in either case it is
replaced by the expansion of the nest to which it is assigned.

Second, the expansion imposes no relation among letters other than a local one. By part~(2), the
identities \eqref{4.6:eq-tower-reading-1}, \eqref{4.6:eq-tower-reading-2},
\cite[(3.151)--(3.152)]{B-CMP99b} and \cite[(3.172)--(3.182)]{B-CMP99b} are identities among
the global operators of one nest, applied to the global formulas \cite[(3.126), (3.129),
(3.147), (3.153), (3.185)]{B-CMP99b} before any expansion; they therefore act on that ingredient
before the substitution of part~(1)(d) and impose nothing on the localized terms that the
substitution produces. After the substitution the only identity imposed on letters is the local
identity \cite[(3.25)]{B-CMP99b}, clause (6) of Lemma~\ref{4.6:lem-commutator-identity}. Hence
$K_\omega$ is a product and tree contraction of letters, with no global constraint linking them:
it is a word, a finite tree. By the reading table of Lemma~\ref{4.6:lem-reading-table}, each of
its letters is either a letter of the closure $\Theta$ for $\mathbf{\Omega}$ or for
$\widetilde{\mathbf{\Omega}}$, or a vertex of radius zero from the closed list that forms the
eleventh entry of that table; constituents of kind (s) do not occur.

Third, the outermost factors. The global operators replaced in part~(1)(d) are expanded by the
formal expansions $\mathfrak{K}'$, $\mathfrak{K}_C$, $\mathfrak{K}_G$, so the outermost
expansion factors of the word are $\mathfrak{K}'$-, $\mathfrak{K}_C$- or $\mathfrak{K}_G$-words,
and clause (d) of Lemma~\ref{4.6:lem-formal-expansions} supplies the outermost bumps of such
words. Beyond these factors, in the global formulas \cite[(3.126), (3.129), (3.147), (3.153),
(3.185)]{B-CMP99b}, simplified as in part~(2), each expanded global operator is flanked by at
most one explicit pair of vertices: $(I+D\mu)Q$ on the left with $Q^*(I+\mu^*D^*)$ on the
right, as in \cite[(3.185)]{B-CMP99b}, or $D$ on the left with $D^*$ on the right. This
gives the three forms in \eqref{4.6:eq-tower-reading-a}.

Finally, an atom of the form \eqref{4.6:eq-tower-reading-a} is a word over the alphabet for which
Lemma~\ref{4.6:lem-closure-local} is stated, with outermost bumps supplied by clause (d) of
Lemma~\ref{4.6:lem-formal-expansions}, flanked by at most one explicit vertex; $K_\omega$ is
therefore a word of the class for which Lemma~\ref{4.6:lem-closure-local} is stated, and that
lemma applies to it.
\end{proof}

\begin{definition}[One expansion recipe for all nests]\label{4.6:def-one-recipe}
Let $\mathbf{\Omega}$ be a nest $\Omega_0\supseteq\dots\supseteq\Omega_k$. Let $\mathfrak{K}$ be any one
of the series $\mathfrak{K}'$, $\mathfrak{K}_C$, $\mathfrak{K}_P$, $\mathfrak{K}_R$, $\mathfrak{K}_G$ or a
tower series of Lemma~\ref{4.6:lem-formal-expansions}, an element of
$\mathbb{Z}\langle\langle\mathfrak{A}^{\mathrm{cand}}\rangle\rangle$. Its universal formula in
\cite{B-CMP99b}, recalled with its equation numbers in Lemma~\ref{4.6:lem-formal-expansions}, is a
series $\sum_w c_w\,w$ in the letters of $\mathfrak{A}^{\mathrm{cand}}$; this fixes integers $c_w$,
one for each word $w$, depending on $\mathfrak{K}$ only. To the nest $\mathbf{\Omega}$ we assign
\begin{equation}\label{4.6:eq-one-recipe-a}
\mathfrak{K}[\mathbf{\Omega}]
=\sum_{w}c_w\,\operatorname{adm}_{\mathbf{\Omega}}(w)\,w,
\qquad \operatorname{adm}_{\mathbf{\Omega}}(w)\in\{0,1\},
\end{equation}
the sum running over all words $w$ in the letters of $\mathfrak{A}^{\mathrm{cand}}$, where
$\operatorname{adm}_{\mathbf{\Omega}}$ is the indicator of the set of words admissible for
$\mathbf{\Omega}$. Thus the nest enters $\mathfrak{K}[\mathbf{\Omega}]$ only through the admissibility
indicator $\operatorname{adm}_{\mathbf{\Omega}}$.

Two constructions in \cite{B-CMP99b} are allowed to use an arbitrary partition. The first is the
freedom stated there that the operator \cite[(3.90)]{B-CMP99b} may be constructed for an arbitrary
partition. The second is the statement that the objects constructed there on the cover
$\mathcal{D}$ ``can be constructed for an arbitrary partition instead of $\mathcal{D}$''. Here both
freedoms are exercised by fixed rules:
the local rules producing the covers $\mathcal{D}'$ and $\mathcal{D}''$ obtained by the surgery of
Definition~\ref{4.6:def-cover-surgery}, and the single gluing rule $\Phi^h$
(Definitions~\ref{4.6:def-cover-surgery} and~\ref{4.6:def-bump-profiles}).

Consequently any two nests $\mathbf{\Omega},\mathbf{\Omega}'$ are expanded by the same rule, with
values in $\mathbb{Z}\langle\langle\mathfrak{A}^{\mathrm{cand}}\rangle\rangle$; this is the form in
which the expansions enter Lemma~\ref{4.6:lem-word-locality}.
\end{definition}

\begin{proof}
Let $\mathbf{\Omega}$ and $\mathbf{\Omega}'$ be two nests that are to be compared, and fix one of
the expansions $\mathfrak{K}$.

By the definition \eqref{4.6:eq-one-recipe-a}, $\mathfrak{K}[\mathbf{\Omega}]$ and
$\mathfrak{K}[\mathbf{\Omega}']$ are sums over the same words $w$. The coefficients
$c_w$ are the same in both sums, since they are read off the formula for $\mathfrak{K}$ alone,
before any nest is chosen. The two sums
differ only in the values $\operatorname{adm}_{\mathbf{\Omega}}(w)$ and
$\operatorname{adm}_{\mathbf{\Omega}'}(w)$, each of which lies in $\{0,1\}$.

Hence $\mathfrak{K}[\mathbf{\Omega}]$ and $\mathfrak{K}[\mathbf{\Omega}']$ are produced by one rule
whose only nest-dependent input is the admissibility indicator.
\end{proof}

\begin{lemma}[Anchors and inserted cubes as constituents of words]\label{4.6:lem-anchor-constituents}
Let $\omega$ be a word of the expansions, in the sense of Definition~\ref{4.6:def-letters-words}, with
U-footprint $F(\omega)$. Let $\square\in\mathcal{D}_j$ be the anchor cube of a surgery of the cover
(Definition~\ref{4.6:def-cover-surgery}) used by $\omega$. Let $\square_0=\square_0(\square)$ be the
cube inserted by that surgery, and $\mathcal{D}'(\square)$ the surgered cover. Let $\mathcal{D}''$ be
the second surgered cover, and let $\widetilde{\square}_0$ be the cube inserted in $\mathcal{D}''$ at
$\square_0$; both are fixed together with the surgeries of the cover in
Definition~\ref{4.6:def-cover-surgery}.
\begin{itemize}
\item[(a)] The anchor $\square$ is a constituent of a letter of $\omega$. Hence
$\square^{(5)}\subset F(\omega)$.
\item[(b)] Assume the following:
\begin{itemize}
\item the gluing rule satisfies the kept-member clause of \eqref{4.6:eq-bump-profiles-c}: for every
surgered cover used by $\omega$, in particular for $\mathcal{D}'(\square)$ and $\mathcal{D}''$, and for
every kept member $Q$ of it, $\operatorname{supp} h'_Q\subset Q$, where $h'_Q$ is the bump attached to
$Q$ in the surgered family; here a kept member of a surgered cover is a member of the unsurgered cover
that does not meet the inserted cube, so that a kept member of $\mathcal{D}'(\square)$ is a cube
$Q\in\mathcal{D}$ with $Q\cap\square_0=\emptyset$ (Definition~\ref{4.6:def-cover-surgery});
\item for every anchor of a surgery used by $\omega$, in particular for the anchor $\square$, the
level agreement of \eqref{4.6:eq-commutator-identity-b} at that anchor holds;
\item the display \eqref{4.6:eq-letters-words-c}, among the conditions imposed on letters, words and
footprints in Definition~\ref{4.6:def-letters-words}, holds;
\item $M\ge L$ and $L\ge 2$.
\end{itemize}
The one-scale profiles of Definition~\ref{4.6:def-bump-profiles} satisfy the kept-member clause, so
for them the first hypothesis holds.
Under these hypotheses, $\square_0$ is a constituent of the same word $\omega$, and so is the cube
$\widetilde{\square}_0$ for the cover $\mathcal{D}''$. Hence $\square_0^{(5)}\subset F(\omega)$.
\end{itemize}
For a gluing rule that violates the kept-member clause, part~(b) is not asserted.
\end{lemma}

\begin{proof}
\emph{Part (a).} A word uses a surgery with anchor $\square$ through one of two kinds of summands.

The first kind is a summand of $R$ in \cite[(3.95)]{B-CMP99b}. Each such summand ends with the
factor $h_{\square}C_{\square}h_{\square}$.

The second kind is a summand of the second or third sum in \cite[(3.105)]{B-CMP99b}, which is the
identity of part~(4) of Lemma~\ref{4.6:lem-commutator-identity}. Each such summand ends with the
factor $h_{\square}G_{\square}h_{\square}$.

In both cases $\square$ is the cube to which a factor of a letter of $\omega$ is attached. It is
therefore a constituent of that letter. The inclusion $\square^{(5)}\subset F(\omega)$ follows from
the definition of the U-footprint of a word in Definition~\ref{4.6:def-letters-words}.

\emph{Part (b).} For the one-scale profiles the bump attached to a kept member $Q$ is its profile
$\eta_Q$, and $\operatorname{supp}\eta_Q\subset\operatorname{int}Q^{(-1/3)}\subset Q$
(Definition~\ref{4.6:def-bump-profiles}); so the kept-member clause holds. We treat first the letters
coming from \cite[(3.95)]{B-CMP99b}.

\emph{Step 1: support of the factor to the right of the walk.} In such a letter, a walk in the
surgered cover $\mathcal{D}'(\square)$ stands immediately to the left of the factor
$Q'^{*}h_{\square}$. Two facts bound the support of this factor.
\begin{itemize}
\item The bump $h_{\square}$ of the anchor satisfies $\operatorname{supp} h_{\square}\subset\square$,
by the support clause of \eqref{4.6:eq-bump-profiles-c}.
\item The operator $Q'^{*}$ spreads supports by at most one averaging block of level at most $j+1$;
the bound $j+1$ on the level holds under the level agreement at $\square$ and under
\eqref{4.6:eq-letters-words-c}. An averaging block of level $i$ has side $\lambda_i$, and by the
conventions fixing the scales and block sides, $u_i=M\lambda_i$ and $\lambda_{i+1}=L\lambda_i$ for
every scale $i$. Such a block therefore has side at most $\lambda_{j+1}=L\lambda_j$, while
$u_j=M\lambda_j$; hence, under $M\ge L$, this side is at most $u_j$.
\end{itemize}
Consequently, for every function to which it is applied,
\begin{equation}\label{4.6:eq-anchor-constituents-1}
\operatorname{supp}\bigl(Q'^{*}h_{\square}\,\cdot\,\bigr)\subset\square^{(1)} .
\end{equation}
The walk is multiplied on the right by this factor. It must therefore end at a cube $Q$ of
$\mathcal{D}'(\square)$ whose bump support meets $\operatorname{supp}(Q'^{*}h_{\square}\,\cdot\,)$,
hence meets $\square^{(1)}$ by \eqref{4.6:eq-anchor-constituents-1}.

\emph{Step 2: the inserted cube contains $\square^{(2)}$.} By
Definition~\ref{4.6:def-cover-surgery},
\begin{equation}\label{4.6:eq-anchor-constituents-2}
\mathcal{D}'(\square)=\{\square_0\}\cup\{Q\in\mathcal{D}:\ Q\cap\square_0=\emptyset\} .
\end{equation}
We show $\square_0\supseteq\square^{(2)}$ in each of the two cases of the surgery.

In case~1 we have $\square_0=\square^{2}$, which contains $\square^{(2)}$.

In case~2, $\square_1\in\mathcal{D}_{j+1}$ is the member closest to the centre of $\square$, and
$\square_0=\square_1^{2}$. The anchor $\square$ meets $\square_1$, so
$\square\subset\square_1^{+2u_j}$. Hence $\square^{(2)}\subset\square_1^{+4u_j}$.

Since $u_i=M\lambda_i$ for every scale $i$ and $\lambda_{j+1}=L\lambda_j$, we have $u_{j+1}=Lu_j$;
with $L\ge 2$ this gives $4u_j\le 2u_{j+1}$. Therefore
\begin{equation}\label{4.6:eq-anchor-constituents-3}
\square^{(2)}\subset\square_1^{+4u_j}\subset\square_1^{+2u_{j+1}}=\square_1^{2}=\square_0 .
\end{equation}

\emph{Step 3: only $\square_0$ qualifies as the end cube.} Let $Q$ be a kept member of
$\mathcal{D}'(\square)$, that is, a member other than $\square_0$. By
\eqref{4.6:eq-anchor-constituents-2}, $Q$ is disjoint from $\square_0$. By the kept-member clause,
$\operatorname{supp}h'_Q\subset Q$. Step~2 gives $\square_0\supseteq\square^{(2)}\supseteq\square^{(1)}$,
so $\operatorname{supp}h'_Q$ is disjoint from $\square^{(1)}$. By
\eqref{4.6:eq-anchor-constituents-1} it does not meet the support of $Q'^{*}h_{\square}$.

Hence the only cube at which the walk can end is $Q=\square_0$. This holds wherever the walk starts,
and whether or not commutator factors $K(h)$ or $K'(h)$ occur along it. The instance with no such
factor is the one treated in \cite{B-CMP99b}.

\emph{Step 4: the letters coming from \cite[(3.105)]{B-CMP99b}.} The same analysis applies:
\cite{B-CMP99b} states, for the sums of \cite[(3.105)]{B-CMP99b}, that the analysis is the same as
before. In this case the
$\mathcal{D}'(\square)$-walk stands immediately to the left of
$D^{*}h_{\square}G_{\square}h_{\square}$. Here
$\operatorname{supp}(D^{*}h_{\square}G_{\square}h_{\square}\,\cdot\,)\subset\square^{(1)}$, because
the leftmost factor $h_{\square}$ has support in $\square$ and $D^{*}$ spreads supports by one bond;
so \eqref{4.6:eq-anchor-constituents-1} holds with this factor, and Steps~1--3 are repeated with this
factor in place of $Q'^{*}h_{\square}$.

\emph{Step 5: nested anchors and the cover $\mathcal{D}''$.} The formal expansion
$\mathfrak{K}_C$ of Lemma~\ref{4.6:lem-formal-expansions} is substituted into the letters. Anchors
nested inside it are handled by repeating Steps~1--4 for each of them.

The cube $\widetilde{\square}_0$ is, in the same way, the end cube of the $\mathcal{D}''$-walk at
$\square_0$.

\emph{Conclusion.} By Steps~3--5, $\square_0$ and $\widetilde{\square}_0$ are cubes to which walks
of letters of $\omega$ are attached. They are therefore constituents of $\omega$. The inclusion
$\square_0^{(5)}\subset F(\omega)$ follows from the definition of the U-footprint in
Definition~\ref{4.6:def-letters-words}.

Part~(b) is used in this chapter only for the sharper constant in the bound on the reach of words and
for reading case~2 of the surgery inside the footprint.
\end{proof}

\begin{lemma}[The reading surcharge of the surgeries]\label{4.6:lem-surgery-surcharge}
Let $j<k$, and let $\square=a+[-1,1]^d$ be the anchor cube of a surgery of scale $j$
(Definition~\ref{4.6:def-cover-surgery}). Lengths are measured in units in which $u_j=1$, so that
$u_{j+1}=L$. The grids are placed so that the centres of the candidate cubes of scale $j$ lie
in $\mathbb{Z}^d$ and those of scale $j+1$ lie in $L\mathbb{Z}^d$. Thus $a\in\mathbb{Z}^d$, and a
candidate of scale $j+1$ is $Q'=c+[-L,L]^d$ with $c\in L\mathbb{Z}^d$.

In these units, for a candidate $Q=c+[-u_i,u_i]^d$ of scale $i$ the concentric enlargements are
\begin{equation*}
Q^{n}=c+[-(n+1)u_i,(n+1)u_i]^d ,
\end{equation*}
and for every cube $Q$ carrying the scale $i$, candidates and their derived cubes $Q^{n}$ alike,
\begin{equation*}
Q^{(n)}=Q^{+nu_i}.
\end{equation*}
Here $A^{+s}$ is the closed $\ell^\infty$-enlargement of $A$ by $s$, and $t^+=\max\{t,0\}$.

The case test at $\square$ in Definition~\ref{4.6:def-cover-surgery} asks whether $\square$
meets a cube of $\mathcal{D}_{j+1}$: in case~$1$ it does not, in case~$2$ it does. The test is
decided by the membership in $\mathcal{D}_{j+1}$ of the candidates of scale $j+1$ that meet
$\square$, and the region it reads is the union of these candidates. Let $\omega$ be a word using
the surgery at $\square$, and let $F(\omega)$ be its footprint
(Definition~\ref{4.6:def-letters-words}). The \emph{surcharge} at $\square$ is the least $s\ge0$
such that the region read by the case test lies in $F(\omega)^{+s}$.
\begin{enumerate}
\item[(1)] A candidate $Q'=c+[-L,L]^d$ of scale $j+1$ meets $\square$, as closed sets, if and
only if $|c_\nu-a_\nu|\le L+1$ for $\nu=1,\dots,d$. Hence
\begin{equation}\label{4.6:eq-surgery-surcharge-1}
\bigcup\{Q' : Q'\cap\square\ne\emptyset\}\subset a+[-(2L+1),2L+1]^d=\square^{+2u_{j+1}}.
\end{equation}
In coordinate $\nu$, the \emph{reach} of these candidates beyond $a_\nu$ on the side of
increasing $\nu$-th coordinate is the largest value of $x_\nu-a_\nu$ over the points $x$ of the
candidates of scale $j+1$ that meet $\square$; on the side of decreasing $\nu$-th coordinate it is
the largest value of $a_\nu-x_\nu$ over the same points. The reach equals $2L+1$ on the side of
increasing $\nu$-th coordinate if and only if $a_\nu\equiv-1 \pmod L$. It equals $2L+1$ on the
side of decreasing $\nu$-th coordinate if and only if $a_\nu\equiv 1\pmod L$. For
$a_\nu\equiv 0\pmod L$ the reach is exactly $2L$. Under the open reading, in which two cubes are
said to meet when their interiors intersect, the condition is $|c_\nu-a_\nu|\le L$, and the reach
is bounded by $2L$ in place of $2L+1$.
\item[(2)] In case~$2$, part (b) of Lemma~\ref{4.6:lem-anchor-constituents} implies that the
region read by the case test lies inside $F(\omega)$. Let $\square_1=c_1+[-L,L]^d$ be the cube of
Definition~\ref{4.6:def-cover-surgery}. For every $L\ge 2$, every candidate of scale $j+1$ that
meets $\square$ lies in
\begin{equation*}
\square_1^{3}\subset\square_1^{7}=(\square_1^{2})^{(5)}=\square_0^{(5)}\subset F(\omega).
\end{equation*}
\item[(3)] In case~$1$, part (b) of Lemma~\ref{4.6:lem-anchor-constituents} implies that the
region read by the case test lies inside $(\square^{7})^{+(2L-7)^+u_j}$, where
$\square^{7}\subset F(\omega)$. The surcharge is therefore at most $(2L-7)^+u_j$, and
\begin{equation}\label{4.6:eq-surgery-surcharge-2}
(2L-7)^+u_j=\Bigl(2-\frac{7}{L}\Bigr)^{+}u_{j+1}<2u_{j+1}\le 2M .
\end{equation}
The bound $(2L-7)^+u_j$ is attained under the closed reading at $a_\nu\equiv\pm1\pmod L$, in
words whose only cubes near $\square$ are $\square$ and $\square_0$. Under the open reading the
surcharge is at most $(2L-8)^+u_j$.
\item[(4)] If part (b) of Lemma~\ref{4.6:lem-anchor-constituents} is not used, part (a) alone
gives $\square^{5}\subset F(\omega)$, and this cube reaches $6$ from $a$. Hence at every anchor at
which the case test is made, the surcharge is at most $(2L-5)^+u_j<2u_{j+1}\le 2M$.
\end{enumerate}
In all cases
\begin{equation*}
\rho_\iota(\omega)\subset F(\omega)^{+2M}\qquad\text{and}\qquad
\rho(\omega)\subset F(\omega)^{+r_*M},\quad r_*=2 .
\end{equation*}
\end{lemma}

\begin{proof}
\emph{Part (1).} Two closed cubes meet if and only if their projections meet in every coordinate.
For the $\nu$-th coordinate,
\begin{equation*}
[c_\nu-L,c_\nu+L]\cap[a_\nu-1,a_\nu+1]\ne\emptyset
\iff |c_\nu-a_\nu|\le L+1 .
\end{equation*}
Let $Q'$ meet $\square$ and let $x\in Q'$. Then $|x_\nu-a_\nu|\le|x_\nu-c_\nu|+|c_\nu-a_\nu|$ for
every $\nu$, and the right-hand side is at most $2L+1$. This gives the inclusion in
\eqref{4.6:eq-surgery-surcharge-1}. The equality in \eqref{4.6:eq-surgery-surcharge-1} holds
because $\square^{+2L}=a+[-(1+2L),1+2L]^d$ and $2L=2u_{j+1}$.

For the reach, write $a_\nu=qL+r$ with $q\in\mathbb{Z}$ and $0\le r\le L-1$. On the side of
increasing coordinate we need the largest $c_\nu\in L\mathbb{Z}$ with $c_\nu\le a_\nu+L+1$.
Since $a_\nu+L+1=(q+1)L+r+1$, there are two cases.
\begin{itemize}
\item If $r\le L-2$, this largest value is $(q+1)L$. The reach is
\begin{equation*}
(q+1)L+L-a_\nu=2L-r\le 2L,
\end{equation*}
with equality if and only if $r=0$.
\item If $r=L-1$, the largest value is $(q+2)L$. The reach is
\begin{equation*}
(q+2)L+L-qL-(L-1)=2L+1 .
\end{equation*}
\end{itemize}
In both cases this $c_\nu$ satisfies $c_\nu\ge a_\nu-L-1$. Thus the reach $2L+1$ occurs on this
side exactly when $a_\nu\equiv-1\pmod L$, and the reach is $2L$ when $a_\nu\equiv0\pmod L$.

The side of decreasing coordinate follows by the reflection $x\mapsto-x$. This reflection
preserves $L\mathbb{Z}$ and sends the residue $1$ to $-1$. So the reach $2L+1$ occurs there
exactly when $a_\nu\equiv1\pmod L$, and it is again $2L$ when $a_\nu\equiv0\pmod L$.

A candidate meeting $\square$ with this extremal $c_\nu$ exists. For $\nu'\ne\nu$, let $c_{\nu'}$
be a multiple of $L$ nearest to $a_{\nu'}$. Then $|c_{\nu'}-a_{\nu'}|\le L/2\le L+1$, so the
criterion above holds in every coordinate.

Under the open reading, the interiors $(c_\nu-L,c_\nu+L)$ and $(a_\nu-1,a_\nu+1)$ meet if and only
if $|c_\nu-a_\nu|<L+1$. Since $c_\nu$ and $a_\nu$ are integers, this holds if and only if
$|c_\nu-a_\nu|\le L$. Then $c_\nu+L-a_\nu\le 2L$, and the same bound holds on the other side.

\emph{Part (2).} In case~$2$ some cube of $\mathcal{D}_{j+1}$ meets $\square$; by part (1) its
centre lies at $\ell^\infty$-distance at most $L+1$ from $a$. The cube $\square_1$ is the cube of
$\mathcal{D}_{j+1}$ closest to the centre $a$ of $\square$, closeness being measured by the
$\ell^\infty$ distance between centres; hence $|c_1-a|_\infty\le L+1$.
Let $Q''=c''+[-L,L]^d$ be any candidate of scale $j+1$ meeting $\square$. By part (1),
$|c''-a|_\infty\le L+1$. Hence $|c''-c_1|_\infty\le 2L+2$, and
\begin{equation*}
Q''\subset c_1+[-(3L+2),3L+2]^d\subset c_1+[-4L,4L]^d=\square_1^{3}.
\end{equation*}
The second inclusion holds if and only if $3L+2\le4L$, that is, if and only if $L\ge2$.

Next, $\square_1^{3}\subset\square_1^{7}$. Moreover
\begin{equation*}
\square_1^{7}=c_1+[-8L,8L]^d=(c_1+[-3L,3L]^d)^{+5L}=(\square_1^{2})^{(5)},
\end{equation*}
because $u_{j+1}=L$. In case~$2$ the inserted cube is $\square_0=\square_1^{2}$
(Definition~\ref{4.6:def-cover-surgery}).

By part (b) of Lemma~\ref{4.6:lem-anchor-constituents}, $\square_0$ is a constituent of the same
word $\omega$ as $\square$. Hence $\square_0^{(5)}\subset F(\omega)$, by the same inclusion of the
$(5)$-enlargement of a constituent in the footprint that gives $\square^{(5)}\subset F(\omega)$ in
part (a) of that lemma. This proves part (2).

\emph{Part (3).} In case~$1$ the inserted cube is $\square_0=\square^{2}=a+[-3,3]^d$. Its
enlargement is
\begin{equation*}
\square_0^{(5)}=a+[-8,8]^d=\square^{7}.
\end{equation*}
As in part (2), part (b) of Lemma~\ref{4.6:lem-anchor-constituents} gives
$\square^{7}\subset F(\omega)$.

By \eqref{4.6:eq-surgery-surcharge-1}, the region read lies in $a+[-(2L+1),2L+1]^d$, and
\begin{equation*}
a+[-(2L+1),2L+1]^d\subset\bigl(a+[-8,8]^d\bigr)^{+(2L-7)^+}.
\end{equation*}
Indeed, if $2L+1\le8$ the enlargement is not needed, and otherwise $8+(2L-7)=2L+1$.

This gives the bound on the surcharge in part (3). The equality in
\eqref{4.6:eq-surgery-surcharge-2} holds because $u_j=u_{j+1}/L$. The strict inequality holds
because $(2L-7)^+<2L$. The last inequality compares absolute lengths: $u_{j+1}=M\lambda_{j+1}$
with $\lambda_{j+1}=L^{-(k-j-1)}\le1$, because $j<k$; hence $2u_{j+1}\le2M$.

By part (1), the reach $2L+1$ is attained at $a_\nu\equiv-1\pmod L$ on the increasing side and at
$a_\nu\equiv1\pmod L$ on the decreasing side. This is the attainment asserted in part (3) for words
whose only cubes near $\square$ are $\square$ and $\square_0$.

Under the open reading the region read lies in $a+[-2L,2L]^d$, by part (1). This set is contained
in $(a+[-8,8]^d)^{+(2L-8)^+}$.

\emph{Part (4).} Part (a) of Lemma~\ref{4.6:lem-anchor-constituents} alone gives
\begin{equation*}
\square^{(5)}=a+[-6,6]^d=\square^{5}\subset F(\omega).
\end{equation*}
The argument of part (3) with $6$ in place of $8$ gives
\begin{equation*}
a+[-(2L+1),2L+1]^d\subset\bigl(a+[-6,6]^d\bigr)^{+(2L-5)^+}.
\end{equation*}
Since $(2L-5)^+<2L$, we have $(2L-5)^+u_j<2u_{j+1}$, and $2u_{j+1}\le2M$ as in part (3). So at
every anchor at which the case test is made, in either case, the region read lies in
$F(\omega)^{+2M}$.

\emph{The reading domains.} Let $C$ and $I$ denote the cubes that carry the windows of the bumps
of the letters of $\omega$ in Definition~\ref{4.6:def-letters-words}. For these cubes the
inclusions
\begin{equation*}
C^{+2M}\subset(C^{(5)})^{+2M}
\qquad\text{and}\qquad
(I^{(5)})^{+2M}\subset F(\omega)^{+2M}
\end{equation*}
hold. Together with the
bound at the anchors from part (4), these give $\rho_\iota(\omega)\subset F(\omega)^{+2M}$ and
$\rho(\omega)\subset F(\omega)^{+r_*M}$ with $r_*=2$, the value fixed in
\eqref{4.6:eq-letters-words-a}.
\end{proof}

\begin{theorem}[Locality of the background-field random-walk expansions]
\label{4.6:thm-expansion-locality}
Use the torus $\mathbb{T}$, the block sides $u_g$, the big-block size $M$ and the $\ell^\infty$
enlargements $A^{+s}$ of Notation~\ref{4.6:not-block-sides}, the candidate cubes and their
enlargements $Q^{(t)}$ of Definition~\ref{4.6:def-candidate-cubes}, the nests, region data, traces
and wavy nests of Definition~\ref{4.6:def-nest-traces}, the cover $\mathcal{D}(\mathbf{\Omega})$ and
the local sequences of Definition~\ref{4.6:def-local-sequence}, the surgeries and surgered families
$\mathcal{F}_{\mathfrak{s}}(\mathbf{\Omega})$ of Definition~\ref{4.6:def-cover-surgery}, and the
letters, words, footprints, reading domains and margins of Definition~\ref{4.6:def-letters-words}.
The nests considered are the decreasing nests
$\mathbf{\Omega}=(\Omega_0\supseteq\Omega_1\supseteq\dots\supseteq\Omega_k)$ of levels at most $k$ in
which every $\Omega_i$ is a closed union of $\pi_i$-blocks, each carried together with its region
datum $\mathfrak{P}$. The readings of \cite[Sect.~3]{B-CMP99b} fixed in the definitions referred to
above are in force. Assume the following hypotheses.
\begin{enumerate}
\item[(a)] The one-scale bump profiles \eqref{4.6:eq-bump-profiles-a}, in the form of
Definition~\ref{4.6:def-bump-profiles}.
\item[(b)] The gluing hypothesis \eqref{4.6:eq-bump-profiles-c}, that is, both its locality clause
and its support clause.
\item[(c)] The block-locality of the averaging operations, Lemma~\ref{4.6:lem-averaging-locality},
in full, including the residual reading that goes beyond the block-locality and domain-freeness
quoted there.
\item[(d)] The reading of the operators of \cite[Sect.~3]{B-CMP99b} as words and trees over the
candidate alphabet $\mathfrak{A}^{\mathrm{cand}}$ (Definition~\ref{4.6:def-letters-words}), including
the reading of the propagator towers \eqref{4.6:eq-tower-reading-a}
(Definition~\ref{4.6:def-tower-reading}).
\item[(e)] The reading \eqref{4.6:eq-local-sequence-a} of the local sequence of a cube, and the
reading \eqref{4.6:eq-candidate-cubes-b} of the candidate cubes and their enlargements.
\item[(f)] The reading \eqref{4.6:eq-bump-profiles-d} of the cut-off functions $\zeta_{\square}$.
\item[(g)] The surgery rule \eqref{4.6:eq-cover-surgery-a} at an anchor cube.
\item[(h)] One expansion recipe for all nests, \eqref{4.6:eq-one-recipe-a}
(Definition~\ref{4.6:def-one-recipe}).
\end{enumerate}
The further hypothesis \eqref{4.6:eq-letters-words-c} is assumed only where clauses~(iii) and~(v)
below say so.

Let $K$ be one of the operators $G'$ of \cite[Theorem~3.7]{B-CMP99b},
$(Q'G'^2Q'^*)^{-1}$ of \cite[Theorem~3.9]{B-CMP99b}, and $R$, $P$, $G$ of
\cite[Theorem~3.10]{B-CMP99b}; or, under the reading \eqref{4.6:eq-tower-reading-a}, one of the tree
expansions constructed in \cite[Sect.~3]{B-CMP99b}: those of $G$ (\cite[Sect.~3.D]{B-CMP99b}), of
$G_1$, $H$, $H_1$, of the two tower operators of \cite[(3.147), (3.153)]{B-CMP99b}, of
$(QGQ^*)^{-1}$ and $(QG_1Q^*)^{-1}$, and of the operators $C^{(k\mapsto j)}(\Lambda)$,
$\widehat{C}^{(j)}(\Lambda)$ of \cite[Sect.~3.E]{B-CMP99b}, the last two on a pair consisting of a
nest and a set $\Lambda$, with its wavy nest $\widetilde{\mathbf{\Omega}}$; or an $x$-resolvent of
one of these operators.

Let $\mathfrak{K}_K\in\mathbb{Z}\langle\langle\mathfrak{A}^{\mathrm{cand}}\rangle\rangle$ be the formal
expansion of $K$ over the candidate alphabet $\mathfrak{A}^{\mathrm{cand}}$, which is the same for
every nest (Lemma~\ref{4.6:lem-formal-expansions} and hypothesis~(h)); write $c_\omega$ for the
coefficient of the word $\omega$ in it, $\operatorname{adm}_{\mathbf{\Omega}}(\omega)\in\{0,1\}$ for
the admissibility indicator and $v_{\mathbf{\Omega}}$ for the valuation at the nest
$\mathbf{\Omega}$, extended multiplicatively to words and to tree contractions. Put
\begin{equation}\label{4.6:eq-expansion-locality-1}
\begin{aligned}
\mathcal{W}_K(\mathbf{\Omega})&:=\bigl\{\omega \text{ a word of } \mathfrak{K}_K :
c_\omega\neq0,\ \operatorname{adm}_{\mathbf{\Omega}}(\omega)=1\bigr\},\\
K_\omega[\mathbf{\Omega}]&:=c_\omega\, v_{\mathbf{\Omega}}(\omega),
\end{aligned}
\end{equation}
and, for a word $\omega$ with constituent cubes $C_1,\dots,C_n$, the footprint
\begin{equation}\label{4.6:eq-expansion-locality-2}
[\omega]:=F(\omega):=\bigcup_{i} C_i^{(5)} .
\end{equation}
Write $C_{\mathrm{first}}(\omega)$ and $C_{\mathrm{last}}(\omega)$ for the constituents carried by the
leftmost and the rightmost factor of $\omega$, $j(C)$ for the scale of a cube $C$, $\mu(\omega)$ for
the margin \eqref{4.6:eq-letters-words-b}, and $\rho(\omega)$, $\rho_\iota(\omega)$ for the domains
on which $\omega$ reads the nest and on which it reads its index data. For a closed set
$W'\subset\mathbb{T}$ put
\begin{equation}\label{4.6:eq-expansion-locality-3}
K^{W'}[\mathbf{\Omega}]:=\sum_{\omega\in\mathcal{W}_K(\mathbf{\Omega}),\ F(\omega)\subset W'}
K_\omega[\mathbf{\Omega}] .
\end{equation}
Then for all nests $\mathbf{\Omega}$, $\mathbf{\Omega}'$ as above, with region data $\mathfrak{P}$,
$\mathfrak{P}'$, every word $\omega$ and every closed $W\subset\mathbb{T}$, the following hold.
\begin{enumerate}
\item[(i)] Locality of the index sets. If
$\operatorname{tr}_{\rho(\omega)}(\mathbf{\Omega};\mathfrak{P})
=\operatorname{tr}_{\rho(\omega)}(\mathbf{\Omega}';\mathfrak{P}')$, then
$\omega\in\mathcal{W}_K(\mathbf{\Omega})$ if and only if $\omega\in\mathcal{W}_K(\mathbf{\Omega}')$.
In particular, if the traces are equal on $W^{+2M}$, then
\begin{equation}\label{4.6:eq-expansion-locality-4}
\{\omega\in\mathcal{W}_K(\mathbf{\Omega}) : F(\omega)\subset W\}
=\{\omega\in\mathcal{W}_K(\mathbf{\Omega}') : F(\omega)\subset W\}.
\end{equation}
For the index data alone, which are read on $\rho_\iota(\omega)$, the reach of $\rho_\iota(\omega)$
beyond $F(\omega)$ is less than $2M$ in absolute terms
(Lemma~\ref{4.6:lem-surgery-surcharge}), independently of the bumps, since the index data read no
bump and no removal set; for the words of \cite[(3.90)]{B-CMP99b} this reach is zero.
\item[(ii)] Locality of the atoms. Under the hypothesis of~(i),
$K_\omega[\mathbf{\Omega}]=K_\omega[\mathbf{\Omega}']$ as partial functions of the background
configuration on their common $U$-read set, with the same domain of definition. For the generators
at a cube $Q$ this is Proposition~\ref{4.6:prop-generator-equality}, which requires equality of the
traces on $Q^{(5)}$ only.
\item[(iii)] The $U$-footprint. Under \eqref{4.6:eq-letters-words-c}, the $U$-read set of
$K_\omega$ is contained in $F(\omega)$ \cite[Sect.~3]{B-CMP99b}; without
\eqref{4.6:eq-letters-words-c} it is contained in $F(\omega)^{+2M}$.
\item[(iv)] Shape of the index sets and footprints. $\mathcal{W}_K(\mathbf{\Omega})$ is a set of
finite words or trees over the alphabet $\mathfrak{A}^{\mathrm{cand}}$ \cite[Sect.~3]{B-CMP99b}, in
which the localization domains of consecutive letters intersect; $F(\omega)$ is a finite union of
closed cubes whose interior is connected and dense in $F(\omega)$, so that $F(\omega)$ is regular
closed.
\item[(v)] Support of the kernels. Suppose that $\beta,\beta'$ are arguments of the kernel of
$K_\omega$ with $K_\omega(\beta,\beta')\neq0$. If $\omega$ is unflanked, that
is, its outer factors carry no explicit vertex (this is the case for every word of
\cite[(3.90), (3.98), (3.107)]{B-CMP99b}), then $\beta\in C_{\mathrm{first}}(\omega)$ and
$\beta'\in C_{\mathrm{last}}(\omega)$. If $\omega$ is flanked, and $M\ge L$ holds together with
either \eqref{4.6:eq-letters-words-c} or the level agreement \eqref{4.6:eq-commutator-identity-b} at
the level of each flanking block, then $\beta\in C_{\mathrm{first}}(\omega)^{(2)}$ and
$\beta'\in C_{\mathrm{last}}(\omega)^{(2)}$. The sets $C_{\mathrm{first}}(\omega)^{(2)}$ and
$C_{\mathrm{last}}(\omega)^{(2)}$ lie in $F(\omega)$ at distance at least
\begin{equation}\label{4.6:eq-expansion-locality-5}
\mu(\omega)=3\min\bigl\{u_{j(C_{\mathrm{first}})},\,u_{j(C_{\mathrm{last}})}\bigr\}>0
\end{equation}
from $\mathbb{T}\setminus F(\omega)$.
\item[(vi)] Convergence. Let $\mathbf{\Omega}$ be a nest as above satisfying the standing
hypotheses of \cite{B-CMP99b}: $M$ sufficiently large; the background configuration $U$ satisfies
\cite[(3.35)--(3.36)]{B-CMP99b} or \cite[(3.37)--(3.38)]{B-CMP99b}, according to the case of
\cite[Theorem~3.1]{B-CMP99b}, with $M\alpha_0\le a_0$ in the notation there; the sequence of
domains satisfies
\cite[(2.1)--(2.2)]{B-CMP96} (stated in \cite[p.~224]{B-CMP96}), and for the nests of the class
$\mathfrak{N}_{\mathrm{rec}}$ of Definition~\ref{4.6:def-atom-locality} this condition is supplied
by the admissibility statement for these nests; for the
operators on a pair of a nest and a set $\Lambda$, the condition of \cite[Sect.~3.E]{B-CMP99b} that
``a distance between $A$ and $\Lambda_k^c$ is bigger than $RM$'', in the notation there (there $R$
is Ba{\l}aban's constant, not the operator $R$ above); the level agreement
\eqref{4.6:eq-commutator-identity-b} for every cube of the cover
(Lemma~\ref{4.6:lem-inner-seam}); and the partition property
$\sum_{\square}h_{\square}^2=1$ of the gluing rule (\cite[(2.36)]{B-CMP96}; for the rule of
Proposition~\ref{4.6:prop-gluing-example} this is \eqref{4.6:eq-gluing-example-a}). Then
$\sum_{\omega}K_\omega[\mathbf{\Omega}]$ converges absolutely in the norms of \cite{B-CMP99b}
(\cite[Theorems~3.7, 3.9, 3.10, 3.12, 3.13, 3.15]{B-CMP99b}). Hence every sub-sum
$K^{W'}[\mathbf{\Omega}]$ converges absolutely, and $K^{\mathbb{T}}[\mathbf{\Omega}]=K$.
\end{enumerate}
\end{theorem}

\begin{proof}
The clauses are proved in the order (i)--(vi). Lemma~\ref{4.6:lem-reading-table} and
Lemma~\ref{4.6:lem-closure-local} are proved as implications from hypotheses that are among
(a)--(h); they are therefore available here.

\emph{Clause (i).} Fix a word $\omega$ and suppose that
$\operatorname{tr}_{\rho(\omega)}(\mathbf{\Omega};\mathfrak{P})
=\operatorname{tr}_{\rho(\omega)}(\mathbf{\Omega}';\mathfrak{P}')$. The admissibility of $\omega$
for $\mathbf{\Omega}$ is the conjunction of four kinds of test.
\begin{itemize}
\item The membership tests: the membership of each constituent of $\omega$ in the cover
$\mathcal{D}(\mathbf{\Omega})$ or in a surgered family $\mathcal{F}_{\mathfrak{s}}(\mathbf{\Omega})$;
these are the entries of the reading table of Lemma~\ref{4.6:lem-reading-table} for memberships in
the cover and in the surgered families.
\item The case tests: for each anchor $\square$ of a surgery used by $\omega$, the case
$\mathrm{CASE}(\square;\mathbf{\Omega})$ and the closest coarser member $\square_1$; by the entry of
the same table for the surgery case test, together with the surgery rule
\eqref{4.6:eq-cover-surgery-a}, these are read inside $\square^{+2u_{j(\square)+1}}$.
\item The scale tests: the scale-admissibility of \cite[Sect.~3]{B-CMP99b}, which reads the scales
$j(C_i)$ of the constituents and the level data at the constituents; the latter are given by the
entry of the same table for the level data.
\item The shape tests: the conditions that make $\omega$ a chain or a tree, which do not refer to
the nest.
\end{itemize}
Every membership, case and scale test reads membership bits of the nest and of its region datum
inside $\rho_\iota(\omega)$, and $\rho_\iota(\omega)\subset\rho(\omega)$
(Definition~\ref{4.6:def-letters-words}). Since the traces of $(\mathbf{\Omega};\mathfrak{P})$ and of
$(\mathbf{\Omega}';\mathfrak{P}')$ on $\rho(\omega)$ are equal, these bits are equal for the two
nests, every test returns the same verdict for both, and the shape tests do not involve the nest at
all. Hence
$\operatorname{adm}_{\mathbf{\Omega}}(\omega)=\operatorname{adm}_{\mathbf{\Omega}'}(\omega)$.
The coefficient $c_\omega$ does not depend on the nest, by \eqref{4.6:eq-one-recipe-a}, so the
condition $c_\omega\neq0$ is the same for both nests. By \eqref{4.6:eq-expansion-locality-1},
$\omega\in\mathcal{W}_K(\mathbf{\Omega})$ if and only if $\omega\in\mathcal{W}_K(\mathbf{\Omega}')$.

Now suppose that the traces are equal on $W^{+2M}$ and that $F(\omega)\subset W$. By
\eqref{4.6:eq-letters-words-a}, applied to the letters of $\omega$, and by
Lemma~\ref{4.6:lem-surgery-surcharge}, applied to the surgeries used by $\omega$, one has
$\rho(\omega)\subset F(\omega)^{+2M}$; an anchor of the top scale has no coarser member
$\square_1$ and contributes no further read. Since $F(\omega)\subset W$, every point $y$ satisfies
\begin{equation*}
\operatorname{dist}_\infty(y,W)\le\operatorname{dist}_\infty(y,F(\omega)),
\end{equation*}
so $F(\omega)^{+2M}\subset W^{+2M}$ and therefore $\rho(\omega)\subset W^{+2M}$. Equality of the
traces on $W^{+2M}$ entails their equality on the subset $\rho(\omega)$, and the first part applies:
$\omega\in\mathcal{W}_K(\mathbf{\Omega})$ if and only if $\omega\in\mathcal{W}_K(\mathbf{\Omega}')$.
Applying this to every word with $F(\omega)\subset W$ gives both inclusions in
\eqref{4.6:eq-expansion-locality-4}.

The membership, case, scale and shape tests read no bump and no removal set
(Lemma~\ref{4.6:lem-membership-tests} and Definition~\ref{4.6:def-cover-surgery}). The reach of
$\rho_\iota(\omega)$ beyond $F(\omega)$ is therefore bounded by the surcharge of the surgeries used
by $\omega$, and is independent of the bumps; by Lemma~\ref{4.6:lem-surgery-surcharge} this
surcharge is less than $2M$ in absolute terms, and it is zero for the words of
\cite[(3.90)]{B-CMP99b}.

\emph{Clause (ii).} Assume the hypothesis of~(i). Every letter of $\omega$ belongs to the closure
$\Theta$ of the ingredients of Definition~\ref{4.6:def-letters-words}: this is
Lemma~\ref{4.6:lem-formal-expansions} for the words of the expansions, and the reading
\eqref{4.6:eq-tower-reading-a} for the towers. By Lemma~\ref{4.6:lem-closure-local}, every element of
$\Theta$ is a local letter in the sense of \eqref{4.6:eq-local-letters-a}, with nest-reading domain
$\rho$; thus the valuations $v_{\mathbf{\Omega}}$ and $v_{\mathbf{\Omega}'}$ agree on every letter
$\mathfrak{a}$ for which the traces of the two nests agree on $\rho(\mathfrak{a})$, in particular on
every letter with $\rho(\mathfrak{a})\subset\rho(\omega)$. Since $\rho(\omega)$ is the union of the
domains $\rho(\mathfrak{a})$ over the letters $\mathfrak{a}$ of $\omega$,
Lemma~\ref{4.6:lem-word-locality}(ii), applied with $W:=\rho(\omega)$, gives
$v_{\mathbf{\Omega}}(\omega)=v_{\mathbf{\Omega}'}(\omega)$. Multiplying by the nest-free coefficient
$c_\omega$ gives $K_\omega[\mathbf{\Omega}]=K_\omega[\mathbf{\Omega}']$ by
\eqref{4.6:eq-expansion-locality-1}. For the generators at a cube $Q$ the equality is
Proposition~\ref{4.6:prop-generator-equality}, which reads the traces on $Q^{(5)}$ only. The
$U$-read set itself, and with it the domain of definition of $K_\omega$ as a partial function of the
background configuration, is determined by the level data at the neighbourhoods of the constituents,
which are read at radius zero (the entry of the reading table of
Lemma~\ref{4.6:lem-reading-table} for those neighbourhoods), and by the supports of the bumps (the
entry of the same table for the bumps); both are equal under the two nests, by the equality of the
traces on $\rho(\omega)$. Hence the two partial functions have the same domain and agree on it.

\emph{Clause (iii).} The reading table of Lemma~\ref{4.6:lem-reading-table} places the $U$-read set
of each letter $\mathfrak{a}$ of $\omega$ inside the $U$-footprint $F(\mathfrak{a})$ of
$\mathfrak{a}$ when \eqref{4.6:eq-letters-words-c} holds, and inside $F(\mathfrak{a})^{+2M}$ when it
is not assumed; here $F(\mathfrak{a})\subset F(\omega)$, hence also
$F(\mathfrak{a})^{+2M}\subset F(\omega)^{+2M}$. By Lemma~\ref{4.6:lem-closure-local} the
$U$-read set of a word is the union of those of its letters. This gives both inclusions of~(iii).

\emph{Clause (iv).} That the elements of $\mathcal{W}_K(\mathbf{\Omega})$ are finite words or trees
over $\mathfrak{A}^{\mathrm{cand}}$, with intersecting localization domains for consecutive letters,
is part of Definition~\ref{4.6:def-letters-words}. For the footprint, put $K_i:=C_i^{(5)}$, so that
$F:=F(\omega)=\bigcup_i K_i$ is a finite union of closed cubes and hence closed. Put
$O_\omega:=\bigcup_i\operatorname{int}K_i$; then $O_\omega\subset\operatorname{int}F$. A closed cube
is the closure of its interior, so
\begin{equation*}
\operatorname{cl}O_\omega\supseteq\bigcup_i\operatorname{cl}\operatorname{int}K_i=\bigcup_iK_i=F,
\end{equation*}
and $\operatorname{cl}O_\omega\subset F$ because $F$ is closed; thus $O_\omega$, and a fortiori
$\operatorname{int}F$, is dense in $F$, and $F=\operatorname{cl}\operatorname{int}F$, that is, $F$ is
regular closed. For connectedness, let $C_{i-1}$ and $C_i$ be consecutive constituents (adjacent in
the chain or tree). They share a point $p$. Since $C_{i-1}\subset\operatorname{int}C_{i-1}^{(5)}$ and
$C_i\subset\operatorname{int}C_i^{(5)}$, the point $p$ lies in
$\operatorname{int}K_{i-1}\cap\operatorname{int}K_i$. Each $\operatorname{int}K_i$ is connected, and
a union of connected sets along a chain or a tree in which adjacent members intersect is connected
(by induction along the chain or tree); hence $O_\omega$ is connected. Let
$x\in\operatorname{int}F$, and choose an open $\ell^\infty$-ball $B_x$ with $x\in B_x\subset
\operatorname{int}F$. Since $O_\omega$ is dense in $F\supseteq B_x$ and $B_x$ is open and non-empty,
$B_x\cap O_\omega\neq\emptyset$, so $B_x\cup O_\omega$ is connected. Then
\begin{equation*}
\operatorname{int}F=O_\omega\cup\bigcup_{x\in\operatorname{int}F}B_x
=\bigcup_{x\in\operatorname{int}F}\bigl(B_x\cup O_\omega\bigr)
\end{equation*}
is a union of connected sets all containing the connected set $O_\omega$, and is therefore
connected.

\emph{Clause (v).} By Lemma~\ref{4.6:lem-formal-expansions}(d) and the last clause of the reading of
the towers (Definition~\ref{4.6:def-tower-reading}), the leftmost factor of every word $\omega$ is
$E_L\,h_{C_{\mathrm{first}}}$, where $C_{\mathrm{first}}\in\mathcal{D}(\mathbf{\Omega})$ is a
candidate member and $E_L$ is either absent or a single explicit vertex of spread $\sigma$; the
symmetric statement holds on the right. For a flanked word, under the hypotheses of~(v) for that
case, $\sigma\le u_{j(C_{\mathrm{first}})}$; this bound on the spread is the one of
Definition~\ref{4.6:def-letters-words}: the spread consists of one lattice step and one
averaging block of level at most $j(C_{\mathrm{first}})+1$, each of size at most
$u:=u_{j(C_{\mathrm{first}})}$ under $M\ge L$ together with either \eqref{4.6:eq-letters-words-c}
or \eqref{4.6:eq-commutator-identity-b}. By the support clause of the gluing hypothesis
\eqref{4.6:eq-bump-profiles-c}, $\operatorname{supp}h_{C_{\mathrm{first}}}\subset C_{\mathrm{first}}$.
By the kernel clause of Lemma~\ref{4.6:lem-closure-local}, the rows of $K_\omega$ lie in
\begin{equation*}
\widehat{X}_L=\bigl(\operatorname{supp}h_{C_{\mathrm{first}}}\bigr)^{+\sigma}
\subset C_{\mathrm{first}}^{+u}\subset C_{\mathrm{first}}^{(2)},
\end{equation*}
and $\widehat{X}_L=\operatorname{supp}h_{C_{\mathrm{first}}}\subset C_{\mathrm{first}}$ when $E_L$ is
absent, which is the case for every word of the expansions $\mathfrak{K}'$, $\mathfrak{K}_C$,
$\mathfrak{K}_G$ of Lemma~\ref{4.6:lem-formal-expansions}, that is, of
\cite[(3.90), (3.98), (3.107)]{B-CMP99b}.
In the same way the columns of $K_\omega$ lie in $C_{\mathrm{last}}^{(2)}$, and in $C_{\mathrm{last}}$
for the unflanked words.

It remains to bound the distance from the complement. Let $C$ be $C_{\mathrm{first}}$ or
$C_{\mathrm{last}}$, and $u:=u_{j(C)}$. By \eqref{4.6:eq-expansion-locality-2},
$F(\omega)\supseteq C^{(5)}=C^{+5u}$ (Definition~\ref{4.6:def-candidate-cubes}), so
$\mathbb{T}\setminus F(\omega)\subset
\mathbb{T}\setminus C^{(5)}$. Let $x\in C^{(2)}=C^{+2u}$ and $y\in\mathbb{T}\setminus C^{(5)}$, so
that $\operatorname{dist}_\infty(x,C)\le2u$ and $\operatorname{dist}_\infty(y,C)>5u$. The triangle
inequality for the distance to a set gives
\begin{equation*}
\operatorname{dist}_\infty(x,y)\ge\operatorname{dist}_\infty(y,C)-\operatorname{dist}_\infty(x,C)
>5u-2u=3u .
\end{equation*}
Hence every point of $C^{(2)}$ lies in $F(\omega)$ at distance at least $3u_{j(C)}$ from
$\mathbb{T}\setminus F(\omega)$, and
$3u_{j(C)}\ge3\min\{u_{j(C_{\mathrm{first}})},u_{j(C_{\mathrm{last}})}\}=\mu(\omega)$, which is
positive since the block sides are positive. This proves \eqref{4.6:eq-expansion-locality-5} and
clause~(v).

\emph{Clause (vi).} Under the hypotheses listed in~(vi), which are those standing in
\cite{B-CMP99b}, the convergence of $\sum_\omega K_\omega[\mathbf{\Omega}]$ in the norms of
\cite{B-CMP99b} is the content of \cite[Theorems~3.7, 3.9, 3.10, 3.12, 3.13, 3.15]{B-CMP99b}, in
the proofs of which the random-walk expansions, whose sum is $K$, are summed. The convergence is
absolute, that is, the sum of the norms of the terms is finite, because the bounds of
\cite{B-CMP99b} are obtained from termwise majorants: ``We use part of the exponentials to control
the sum over $y$'s \dots\ the remaining are used to construct an overall exponential factor'', and
``We will use the factor $O(M^{-1/2})$ to control the sum over random walks $\omega$''
\cite[Sect.~3]{B-CMP99b}. The form of the bounds on individual walks in
\cite[(3.94), (3.99), (3.108)]{B-CMP99b} is not used. For a closed
$W'\subset\mathbb{T}$, the series \eqref{4.6:eq-expansion-locality-3} is a sub-series of this
series; the sum of the norms of its terms is bounded by that of the full series, so
$K^{W'}[\mathbf{\Omega}]$ converges absolutely. Every footprint is a subset of $\mathbb{T}$, so
$K^{\mathbb{T}}[\mathbf{\Omega}]$ is the full series, and $K^{\mathbb{T}}[\mathbf{\Omega}]=K$.
\end{proof}

\begin{corollary}[Locality inherited by words in the expanded operators]\label{4.6:cor-operator-words}
Work in the setting and under the hypotheses of Theorem~\ref{4.6:thm-expansion-locality}. Its
nests are $\pi$-grained decreasing nests $\mathbf{\Omega}$ with region data $\mathfrak{P}$. The region
data considered here have level at most $k-1$. Let $K^{\circ}$ be an operator that is
\emph{defined} as a nest-free word in the operators $K$ of that theorem. Being so defined means that
$K^{\circ}$ is obtained from these operators by finitely many of the following operations,
with coefficients that do not depend on $\mathbf{\Omega}$ or on $\mathfrak{P}$:
\begin{enumerate}
\item[(a)] finite sums and products, a product being the grafting of words along a shared
constituent or along an explicit vertex, with footprint the union of the footprints;
\item[(b)] transposition;
\item[(c)] insertion of nest-free multipliers, of multipliers of radius~$0$, and of the
component projections $\chi_P$ read from the region datum $\mathfrak{P}$;
\item[(d)] Taylor coefficients and contour integrals in a nest-free parameter;
\item[(e)] families indexed by a parameter $x$, and integrals over $x$, of the following
two kinds. First, resolvents $(x+K)^{-1}$ expanded as in Theorem~\ref{4.6:thm-expansion-locality},
with the local resolvents $(x+\Delta_{a,\square})^{-1}$, and the other local generators at
$x$, in place of the inverses. Second, the square root
\begin{equation}\label{4.6:eq-operator-words-1}
K^{-1/2}=\frac{1}{\pi}\int_0^\infty x^{-1/2}\,(x+K)^{-1}\,dx ,
\end{equation}
expanded termwise.
\end{enumerate}
Let $(K^{\circ}_\omega)_\omega$ be the atom family of $K^{\circ}$ obtained word by word
from the atom families of its constituents. Then $(K^{\circ}_\omega)_\omega$ satisfies
verbatim the index-locality, valuation-locality and support-locality clauses and the
footprint clause of Theorem~\ref{4.6:thm-expansion-locality}. It satisfies the convergence
clause of that theorem under one further hypothesis: the $x$-integrals and the resolvent
series entering through~(e) converge, in the sense of that convergence clause.
\end{corollary}

\begin{proof}
Write $\pi_W$ for the word-by-word restriction of Lemma~\ref{4.6:lem-word-locality}. For a
closed set $W\subset\mathbb{T}$ it keeps exactly the words $w$ with $\rho(w)\subset W$. The
atom family of $K^{\circ}$ arises from those of the operators $K$ through the operations
(a)--(e). The proof has three parts. We show that each operation commutes with $\pi_W$. We
show that each operation preserves the locality of letters. We then separate off convergence.

\emph{Commutation with the restriction.} By Lemma~\ref{4.6:lem-word-locality}(i), $\pi_W$ is a
continuous unital idempotent algebra endomorphism of
$\mathbb{Z}\langle\langle\mathfrak{A}\rangle\rangle$.

For~(a): $\pi_W$ is linear and multiplicative, so it commutes with finite sums and products.
For a grafted product the footprint is the union of the footprints. This agrees with the
rule $\rho(uv)=\rho(u)\cup\rho(v)$ by which $\pi_W$ reads products.

For~(b): transposition reverses the order of the letters in each word and replaces each
letter by its transpose. By the induction in the proof of Lemma~\ref{4.6:lem-closure-local},
the transpose of a local letter is a local letter with the same domain $\rho$. Hence the
union of the domains of the letters of a word is unchanged, and so is the indicator
$\mathbf{1}[\rho(w)\subset W]$.

For~(c): an inserted multiplier or projection $\chi_P$ enters as one further factor of a
product. The case~(a) applies to it.

For~(d) and~(e): in these operations a parameter, nest-free in~(d) and equal to $x$
in~(e), may enter both the coefficients $c_w$ of an element $\sum c_w w$ and the letters
themselves, as it does in the local resolvents $(x+\Delta_{a,\square})^{-1}$. In the induction
in the proof of Lemma~\ref{4.6:lem-closure-local} the parameter is treated as nest-free, and
the domains $\rho$ of the letters that depend on it, and of the letters obtained from them by
differentiation in the parameter or by forming resolvents, do not depend on the parameter.
Hence for each word $w$ the factor $\mathbf{1}[\rho(w)\subset W]$ by which $\pi_W$ multiplies
its coefficient is the same at every value of the parameter. Taylor coefficients, contour
integrals and the integral over $x$ in~\eqref{4.6:eq-operator-words-1}, taken word by word,
therefore commute with $\pi_W$, the interchange with the infinite sum over words being
justified by the continuity of $\pi_W$.

\emph{Locality of letters.} We use the induction in the proof of
Lemma~\ref{4.6:lem-closure-local}. That induction shows that the operations of~(a)--(d),
applied to local
letters, produce local letters. The operation~(e) introduces letters that depend on $x$. We
treat $x$ as a nest-free parameter. The local resolvents $(x+\Delta_{a,\square})^{-1}$, and the
other local generators at $x$, are then covered by Proposition~\ref{4.6:prop-generator-equality}.
That proposition holds for every value of $x$, so these are local letters for every $x$.
Hence every letter occurring in an atom $K^{\circ}_\omega$ is local.

\emph{Conclusion.} By the two parts above, each atom $K^{\circ}_\omega$ is a word in local
letters, and $\pi_W$ acts on it word by word in the same way as on the atoms of the operators
of Theorem~\ref{4.6:thm-expansion-locality}. The index-locality, valuation-locality,
support-locality and footprint clauses of that theorem are properties of the individual words
of an atom family, read through $\pi_W$ and Lemma~\ref{4.6:lem-word-locality}. They therefore
hold for $(K^{\circ}_\omega)_\omega$ with the same proof.

The convergence clause is different. It requires, in addition, that the integral over $x$
in~\eqref{4.6:eq-operator-words-1} and the resolvent series in~(e) converge. This is not a
consequence of the locality statements used above. It is the one further hypothesis under
which the corollary asserts the convergence clause.
\end{proof}

Operators that carry a region datum of level $k$ are outside the scope of this corollary.
They fall under the scope of Corollary~\ref{4.4:cor-localised-independence}. No statement is made here
about operators of \cite{B-CMP122a} and \cite{B-CMP122b} that cannot be expressed as such words.

\begin{proposition}[A counterexample at reading radius zero: equal traces, different admissible words]
\label{4.6:prop-radius-zero-counterexample}
Let $d\ge 1$ and let $j$ be a level with $j+2\le k$. Measure lengths in units in which $u_j=1$, so
that $u_{j+1}=L$ and $u_{j+2}=L^2$. Let $A,A'\in L^2\mathbb{Z}$ with $A\ge L^2$ and
$A'-A\ge L^2$, and let $\mathbb{T}$ be the torus of period $P\in L^2\mathbb{Z}$ with $P\ge 4A'$; points
of $\mathbb{T}$ are written with coordinates in $[-P/2,P/2)$. Define two sequences of closed sets
$\mathbf{\Omega}=(\Omega_i)_{0\le i\le k}$ and $\mathbf{\Omega}'=(\Omega'_i)_{0\le i\le k}$ by
\begin{equation}\label{4.6:eq-radius-zero-counterexample-1}
\begin{gathered}
\Omega_i=\Omega'_i=\mathbb{T}\quad(i\le j),\qquad
\Omega_{j+2}=\Omega'_{j+2}=\{x:\ A\le x_1\le A'\},\\
\Omega_{j+1}=\{x:\ L\le x_1\le A'\},\qquad
\Omega'_{j+1}=\{x:\ 2L\le x_1\le A'\},\qquad
\Omega_i=\Omega'_i=\emptyset\quad(i>j+2).
\end{gathered}
\end{equation}
Put $\square:=[-1,1]^d\in\mathcal{C}_j$ (centre $0$), $Q'':=[0,2L]^d\in\mathcal{C}_{j+1}$ (centre
$(L,\dots,L)\in L\mathbb{Z}^d$), so that $\square\cap Q''=[0,1]^d$, and put
$\square_0:=\square^{2}=[-3,3]^d$ and $W:=[-8,8]^d$. Then:
\begin{enumerate}
\item[(a)] $\mathbf{\Omega}$ and $\mathbf{\Omega}'$ are decreasing nests, each $\Omega_i$, $\Omega'_i$
a union of $\pi_i$-blocks, and $\square\in\mathcal{D}_j(\mathbf{\Omega})\cap\mathcal{D}_j(\mathbf{\Omega}')$;
\item[(b)] $Q''\in\mathcal{D}_{j+1}(\mathbf{\Omega})$ and $Q''\notin\mathcal{D}_{j+1}(\mathbf{\Omega}')$;
\item[(c)] no member of $\mathcal{D}_{j+1}(\mathbf{\Omega}')$ meets $\square$, whereas $Q''$ meets
$\square$; hence the case index $\kappa(\square;\cdot)\in\{1,2\}$ of the surgery of
Definition~\ref{4.6:def-cover-surgery} (equal to $1$ if $\square$ itself is used, and to $2$ if the
closest coarser member $\square_1$ is used) takes the values $\kappa(\square;\mathbf{\Omega})=2$ and
$\kappa(\square;\mathbf{\Omega}')=1$;
\item[(d)] the case-one grouped symbol
\begin{equation}\label{4.6:eq-radius-zero-counterexample-2}
S:=(1-\widetilde{\square})\,Q'h'_{\square_0}G'_{\square_0}h'_{\square_0}\cdot
h'_{\square_0}G'_{\square_0}h'_{\square_0}\cdot Q'^{*}h_{\square}C_{\square}h_{\square}
\end{equation}
of \cite[(3.95)--(3.98)]{B-CMP99b}, the operators $Q'$, $G'_{\square_0}$, $C_{\square}$ and the
bumps $h'_{\square_0}$, $h_{\square}$ being those of that construction, has footprint
$\square_0^{(5)}\cup\square^{(5)}=W$; and, writing $\mathcal{W}(\mathbf{\Omega})$ for
$\mathcal{W}_K(\mathbf{\Omega})$ with $K$ the operator whose expansion there contains $S$,
$S\in\mathcal{W}(\mathbf{\Omega}')$ but $S\notin\mathcal{W}(\mathbf{\Omega})$;
\item[(e)] if moreover $L\ge 9$, the traces of $\mathbf{\Omega}$ and $\mathbf{\Omega}'$ on $W$ coincide.
\end{enumerate}
Consequently, when $L\ge 9$, the locality property of Definition~\ref{4.6:def-atom-locality}, taken
with the footprint $F(\omega)$ and reading radius $0$, does not hold for the index sets of admissible
words defined through the surgery of the cover (Definition~\ref{4.6:def-cover-surgery}). Consistently with
Theorem~\ref{4.6:thm-expansion-locality}, the traces of $\mathbf{\Omega}$ and $\mathbf{\Omega}'$ differ
on $W^{+2M}$.
\end{proposition}

\begin{proof}
Throughout, candidate cubes of scale $g$ are the cubes $c+[-u_g,u_g]^d$ with
$c\in u_g\mathbb{Z}^d$, membership bits are those of Definition~\ref{4.6:def-candidate-cubes}, and
membership in the layer cover is that of Definition~\ref{4.6:def-local-sequence}: a candidate
$Q\in\mathcal{C}_g$ lies in $\mathcal{D}_g(\mathbf{\Omega})$ if and only if there is a block
$e\in\pi_g$ with $e\subset Q$, $[e\subset\Omega_g]=1$ and $[e\subset\Omega_{g+1}]=0$. All
verifications are made coordinatewise. Since $P\ge 4A'$, the sets in
\eqref{4.6:eq-radius-zero-counterexample-1} and the cubes $\square$, $Q''$, $\square_0$, $W$,
$\square_1^{2}$, $e''$ used below lie in the fundamental domain, so the first coordinate $x_1$ is
unambiguous on them.

(a) The faces $x_1\in\{L,2L\}$ of the slabs $\Omega_{j+1}$, $\Omega'_{j+1}$ lie in $L\mathbb{Z}$, and
the faces $x_1\in\{A,A'\}$ lie in $L^2\mathbb{Z}$; since moreover $P\in L^2\mathbb{Z}$, the sets
$\Omega_{j+1}$ and $\Omega'_{j+1}$ are unions of $\pi_{j+1}$-blocks (side $L$), and $\Omega_{j+2}$ is a
union of $\pi_{j+2}$-blocks (side $L^2$). The levels $i\le j$ equal $\mathbb{T}$ and the levels
$i>j+2$ are empty, so they are unions of blocks of their scales as well. Since $L\ge 2$ we have
$A\ge L^2\ge 2L\ge L$, whence
\begin{equation*}
\Omega_{j+2}\subset\Omega'_{j+1}\subset\Omega_{j+1}\subset\mathbb{T},
\end{equation*}
and both sequences are decreasing. For the membership of $\square$, take the block
$e:=[-1,0]\times[-1,0]^{d-1}\in\pi_j$ (any block of $\pi_j$ inside $\square$ with first factor
$[-1,0]$ serves equally). It lies in $\Omega_j=\Omega'_j=\mathbb{T}$, so its level-$j$ bit is $1$ under
both nests. On $e$ we have $x_1\le 0<L$, so $e\not\subset\{x_1\ge L\}$; since
$\{x_1\ge L\}$ contains both $\Omega_{j+1}$ and $\Omega'_{j+1}$, the level-$(j+1)$ bit of $e$ is $0$
under both nests. Hence $\square\in\mathcal{D}_j(\mathbf{\Omega})\cap\mathcal{D}_j(\mathbf{\Omega}')$.

(b) Consider the $\pi_{j+1}$-block $e'':=[L,2L]\times[0,L]^{d-1}\subset Q''$. On it
$L\le x_1\le 2L$, and $2L\le A'$ (because $A'\ge A+L^2\ge 2L^2\ge 2L$), so
$e''\subset\{L\le x_1\le A'\}=\Omega_{j+1}$. It is not contained in
$\Omega_{j+2}\subset\{x_1\ge A\}$: it contains points with $x_1<2L\le L^2\le A$, for instance points
with $x_1=L<A$. Thus $[e''\subset\Omega_{j+1}]=1$ and $[e''\subset\Omega_{j+2}]=0$, and
$Q''\in\mathcal{D}_{j+1}(\mathbf{\Omega})$. Under $\mathbf{\Omega}'$, the $\pi_{j+1}$-blocks contained in
$Q''=[0,2L]^d$ have $x_1$-range $[0,L]$ or $[L,2L]$, and neither range lies in $\{x_1\ge 2L\}$, which
contains $\Omega'_{j+1}$. So every $\pi_{j+1}$-block in $Q''$ has level-$(j+1)$ bit $0$ under
$\mathbf{\Omega}'$, and $Q''\notin\mathcal{D}_{j+1}(\mathbf{\Omega}')$.

(c) A scale-$(j+1)$ candidate $c+[-L,L]^d$, $c\in L\mathbb{Z}^d$, meets $[-1,1]^d$ if and only if
$|c_\nu|\le L+1$ for every $\nu$. Since $c_\nu\in L\mathbb{Z}$ and $L\ge 2$ (so that $L+1<2L$), this
means $c_\nu\in\{-L,0,L\}$ for every $\nu$. Membership of such a candidate in
$\mathcal{D}_{j+1}(\mathbf{\Omega}')$ would require a $\pi_{j+1}$-block inside it with bit $1$ for
$\Omega'_{j+1}$, hence a block contained in $\{x_1\ge 2L\}$. But the blocks of such a candidate have
$x_1$-range contained in $[c_1-L,c_1+L]\subset[-2L,2L]$, with lower end at most $c_1\le L<2L$; no such
block lies in $\{x_1\ge 2L\}$. Hence no member of $\mathcal{D}_{j+1}(\mathbf{\Omega}')$ meets $\square$.
Under $\mathbf{\Omega}$, the member $Q''$ of $\mathcal{D}_{j+1}(\mathbf{\Omega})$ found in (b) meets
$\square$ in $[0,1]^d$. The surgery of Definition~\ref{4.6:def-cover-surgery} is in case one exactly
when the anchor meets no member of $\mathcal{D}_{j+1}$, and in case two otherwise; therefore
$\kappa(\square;\mathbf{\Omega})=2$ and $\kappa(\square;\mathbf{\Omega}')=1$.

(d) The symbol $S$ of \eqref{4.6:eq-radius-zero-counterexample-2}, with $\widetilde{\square}$ as in
\eqref{4.6:eq-candidate-cubes-b} and $\square_0=\square^{2}=[-3,3]^d$, is the case-one grouped symbol of
\cite[(3.95)--(3.98)]{B-CMP99b}. Its footprint is
\begin{equation*}
\square_0^{(5)}\cup\square^{(5)}=[-8,8]^d\cup[-6,6]^d=[-8,8]^d=W .
\end{equation*}
Under $\mathbf{\Omega}'$ the anchor $\square$ belongs to $\mathcal{D}_j(\mathbf{\Omega}')$ by (a) and is
in case one by (c); case one inserts exactly the cube $\square_0=\square^{2}=[-3,3]^d$. Hence
$S\in\mathcal{W}(\mathbf{\Omega}')$. Under $\mathbf{\Omega}$ the anchor $\square$ is in case two by (c),
and its inserted cube is $\square_1^{2}$, where $\square_1$ is the member of
$\mathcal{D}_{j+1}(\mathbf{\Omega})$ meeting $\square$ closest to $0$. By the computation in (c), a
candidate meeting $\square$ has $c_\nu\in\{-L,0,L\}$; if it belongs to
$\mathcal{D}_{j+1}(\mathbf{\Omega})$, one of its $\pi_{j+1}$-blocks, whose $x_1$-range is
$[c_1-L,c_1]$ or $[c_1,c_1+L]$, lies in $\Omega_{j+1}\subset\{x_1\ge L\}$, which forces $c_1\ge L$ and
hence $c_1=L$. By the tie-break rule \eqref{4.6:eq-cover-surgery-a},
\begin{equation*}
\square_1=(L,0,\dots,0)+[-L,L]^d,\qquad \square_1^{2}=(L,0,\dots,0)+[-3L,3L]^d\neq[-3,3]^d .
\end{equation*}
So under $\mathbf{\Omega}$ no derived cube $[-3,3]^d$ exists, and $S\notin\mathcal{W}(\mathbf{\Omega})$.

(e) Assume $L\ge 9$. By Definition~\ref{4.6:def-nest-traces}, the trace of a nest on $W$ is the bit map
$(i,c)\mapsto[c\subset\Omega_i]$ over the $\pi_0$-blocks $c\subset W$. At levels $i\le j$ we have
$\Omega_i=\Omega'_i=\mathbb{T}$, and all bits equal $1$ under both nests. At level $j+1$, a block
$c\subset W=[-8,8]^d$ has $x_1\le 8<9\le L$ on $c$, so $c\not\subset\{x_1\ge L\}$; since this set
contains $\Omega_{j+1}$ and $\Omega'_{j+1}$, all bits are $0$ under both nests. At level $j+2$,
$x_1\le 8<L^2\le A$ on $c$, so $c\not\subset\Omega_{j+2}=\Omega'_{j+2}$, and all bits are $0$ under both
nests. At levels $i>j+2$ both nests are empty and all bits are $0$. The two traces on $W$ coincide.

Conclusion. Let $L\ge 9$. By (e) the nests $\mathbf{\Omega}$ and $\mathbf{\Omega}'$ have equal traces on
the footprint $W$ of $S$, while by (d) the word $S$ is admissible for $\mathbf{\Omega}'$ and not for
$\mathbf{\Omega}$. This contradicts the locality property of Definition~\ref{4.6:def-atom-locality}
taken with footprint $F(\omega)$ and reading radius $0$, so that property fails for these index sets.
There is no conflict with Theorem~\ref{4.6:thm-expansion-locality}, whose window is $2M$. Indeed,
since $u_g=M\lambda_g$ with $\lambda_k=1$, the side of the top blocks is $u_k=M$, and the block sides
increase with the level, so $u_{j+1}=L\le u_k=M$ in the present units; as $W\supseteq\square$,
\begin{equation*}
W^{+2M}\supseteq\square^{+2u_{j+1}}=[-(2L+1),2L+1]^d\supseteq[L,2L]\times[0,L]^{d-1}=e''.
\end{equation*}
By (b), $[e''\subset\Omega_{j+1}]=1$, while $[e''\subset\Omega'_{j+1}]=0$ since the $x_1$-range
$[L,2L]$ of $e''$ does not lie in $\{x_1\ge 2L\}$. The set $\Omega_{j+1}$ is a union of
$\pi_{j+1}$-blocks and $e''\in\pi_{j+1}$. Hence, by the rule (B0) of
Definition~\ref{4.6:def-candidate-cubes}, every $\pi_0$-block $c\subset e''$ has
$[c\subset\Omega_{j+1}]=1$, and likewise $[c\subset\Omega'_{j+1}]=0$. The level-$(j+1)$ bits of these
blocks differ, so the traces differ on $W^{+2M}$.
\end{proof}

\begin{remark}[On the bump functions]
With the one-scale profiles built from the function $h$ of
\cite[(1.118)]{B-CMP95}, one has
\begin{equation*}
\operatorname{supp}\eta_{Q''}\subset Q''^{(-1/3)}=\,]2L/3,4L/3[^d,
\end{equation*}
which does not meet $\square$. No member of $\mathcal{D}_{j+1}$ of either nest has a profile meeting
$\operatorname{supp}\eta_{\square}\subset\,]-2/3,2/3[^d$, since such members have $c_1\ge L$, resp.
$c_1\ge 2L$. So under the gluing rule $(\mathrm{G}^{\natural})$ of
Proposition~\ref{4.6:prop-gluing-example} one has
$h_{\square}[\mathcal{D}(\mathbf{\Omega})]=h_{\square}[\mathcal{D}(\mathbf{\Omega}')]$ here. Under a
general rule of the type \eqref{4.6:eq-bump-profiles-c}, $h_{\square}$ may differ at points of
$\square\cap Q''$, because it reads the membership bit of $Q''$, which differs between the two nests.
This happens inside the window $2M$ of Theorem~\ref{4.6:thm-expansion-locality} and does not
contradict it; the statement at the level of index sets is unaffected. The numbers $L\ge 9$, $A$,
$A'$ and $P$ are parameters of this example only and condition no other statement.
\end{remark}

\begin{corollary}[Cancellation of words localized in the common region]
\label{4.6:cor-common-region-cancel}
Assume the hypotheses of Theorem~\ref{4.6:thm-expansion-locality}, and let $K$ be one of the
operators covered by that theorem. Let $\mathbf{\Omega}=(\Omega_0\supseteq\dots\supseteq\Omega_k)$
and $\mathbf{\Omega}'=(\Omega'_0\supseteq\dots\supseteq\Omega'_k)$ be two nests in the scope of
that theorem, with their region data, and put
\begin{equation}\label{4.6:eq-common-region-cancel-1}
\Omega:=\Omega_k\cap\Omega'_k .
\end{equation}
Let $\omega$ be a word, in the sense of Definition~\ref{4.6:def-letters-words}, whose
nest-reading domain satisfies $\rho(\omega)\subset\Omega$. Then $\omega$ is admissible for
$\mathbf{\Omega}$
if and only if it is admissible for $\mathbf{\Omega}'$, and, when it is admissible, its values at
the two nests are equal. Consequently, in the difference
$K[\mathbf{\Omega}]-K[\mathbf{\Omega}']$ of the expansions of $K$ at the two nests, the terms
of all words $\omega$ with $\rho(\omega)\subset\Omega$ cancel.
\end{corollary}

\begin{proof}
The set $\Omega$ of \eqref{4.6:eq-common-region-cancel-1} is the common region of
\cite[before (3.154)]{B-CMP99b}. It is contained in $\Omega_k$ and in $\Omega'_k$, the innermost
members of the
two nests, and hence, the nests being decreasing, in every member of each nest. On $\Omega$ the
trace of $\mathbf{\Omega}$ and the trace of $\mathbf{\Omega}'$, each
taken with its region datum, are both trivial; in particular they are equal.

Let $\omega$ be a word with $\rho(\omega)\subset\Omega$. By
Theorem~\ref{4.6:thm-expansion-locality}, admissibility of $\omega$ and its value depend on
the nest only through the trace of the nest, with its region datum, on the nest-reading domain
of $\omega$. The nest-reading domain $\rho(\omega)$ lies in $\Omega$, where the two traces
coincide by the
preceding paragraph. Hence $\omega$ is admissible for both nests or for neither, and when it is
admissible its values at $\mathbf{\Omega}$ and at $\mathbf{\Omega}'$ are equal. Consequently each
such word enters $K[\mathbf{\Omega}]$ and $K[\mathbf{\Omega}']$ with the same term, or enters
neither of them, and in the difference $K[\mathbf{\Omega}]-K[\mathbf{\Omega}']$ these terms
cancel.

The cancellation, in the difference of the expansions at two nests, of the terms with
``localizations contained in $\Omega$'' is stated in
\cite[the sentence following (3.154)]{B-CMP99b}. The corollary gives this
statement with that phrase read as $\rho(\omega)\subset\Omega$. The literal reading of the phrase
is $\bigcup_i X_i\subset\Omega$, the $X_i$ being the localization domains in the terms there;
the two conditions differ by the walks that come within distance $2M+5u$ of
the complement $\Omega^{c}$ (with $u$ the length so denoted there), and these walks are absorbed
there ``after adjusting a definition of $\delta_0$''. The bound \cite[(3.154)]{B-CMP99b} given
in the same place for the remaining terms of the difference, whose treatment there involves
\cite[(3.108)]{B-CMP99b}, is not used in what follows.
\end{proof}

\begin{remark}[Inner and outer seams, surgery near seams, and thin separating rings]
\label{4.6:rem-seams-rings}
Let $\mathbf{\Omega}=(\Omega_0\supseteq\Omega_1\supseteq\dots\supseteq\Omega_k)$ be a nest as in
Definition~\ref{4.6:def-nest-traces}. Let $g$ be a level, and let $Q\in\mathcal{D}_g(\mathbf{\Omega})$
be a cube of scale $g$ of the cover, with local sequence $\mathbf{\Omega}(Q)$
(Definition~\ref{4.6:def-local-sequence}). By an \emph{inner seam} at $Q$ we mean a piece of the
boundary of one of the members $\Omega_{g+1}$, $\Omega_{g+2}$ lying in $Q^{(3)}$. By an
\emph{outer seam} at $Q$ we mean a piece of $\partial\Omega_g$ lying in $Q^{(3)}$.
\begin{itemize}
\item[($\sigma1$)] \emph{Inner seams.} An inner seam at $Q$ enters the construction only through the
seam datum $T(Q)$ of Lemma~\ref{4.6:lem-sequence-reach}: equal data $T(Q)$ give equal local
sequences (Lemma~\ref{4.6:lem-sequence-reach}), and equal local sequences together with equal
level-$g$ bits on the blocks of $Q^{(3)}$ give equal localised generators at $Q$
(Proposition~\ref{4.6:prop-generator-equality}). A cube
$Q'\in\mathcal{D}_{g+1}(\mathbf{\Omega})$ which protrudes across the same seam has, by the definition of
its local sequence, top member $\Omega_{g+1}(Q')=Q'^{(4)}$. For a wavy nest
$\widetilde{\mathbf{\Omega}}$, the boundary $\partial\Lambda$ of the region $\Lambda$ of
\cite{B-CMP99b} that enters its region datum is treated as an inner seam: it enters only through the
wavy datum $\widetilde{T}(Q)$ of Lemma~\ref{4.6:lem-sequence-reach}.
\item[($\sigma2$)] \emph{Outer seams.} An outer seam occurs for cubes of $\mathcal{D}_g(\mathbf{\Omega})$
near $\partial\Omega_g$. Such cubes are admitted by the description of the cover in
\cite[p.~229]{B-CMP96}, where the cubes have a ``center $y\in\Lambda_j$ (more exactly \dots\ the
boundary of this set also)''. In \cite[pp.~229--230, 235]{B-CMP96} two cases are treated for a cube of
scale $j$ of the cover: either it lies in the layer $B^j(\Lambda_j)$ of its own scale, or it also
meets $B^{j+1}(\Lambda_{j+1})$, the next layer inward; a cube of scale $j$ meeting $\Omega_j^c$
across $\partial\Omega_j$ is a third configuration. The outer seam
is not an argument of the map $\Sigma(Q,g;\cdot)$ of
Lemma~\ref{4.6:lem-sequence-reach}.

Three properties of a single nest at such cubes are not settled in this section:
\begin{itemize}
\item the applicability of \cite[Corollary~3.6]{B-CMP99b};
\item the level agreement $(\mathrm{LA}_{\square})$ of \eqref{4.6:eq-commutator-identity-b};
\item the gluing of \cite{B-CMP96}.
\end{itemize}
Suppose that Ba{\l}aban's rule were read as attaching to an outer-seam cube a local sequence
different from that of Definition~\ref{4.6:def-local-sequence}, and suppose, as a hypothesis on that
reading, that this local sequence is determined by membership bits of the $\pi_g$-blocks
$e\subset Q^{(3)}$, now at levels $g-1,\dots,g+2$. Then the bits
$[e\subset\Omega_{g-1}]$, $[e\subset\Omega_g]$ of these blocks would join the datum. The datum would
still consist of bits on $Q^{(3)}$, so that the additional reading radius would be $0$.
\item[($\sigma3$)] \emph{Surgery and bumps near seams.} Their reading surcharge is absolute:
smaller than $2M$ for the surgeries and equal to $2M$ for the glued bumps. The gluing of bumps across
a layer seam, the situation of Lemma~\ref{4.6:lem-seam-double-count}, is the place at which the gluing
hypothesis \eqref{4.6:eq-bump-profiles-c} is used. Under $(\mathrm{G\text{-}loc})$ the glued bump at a
point reads only the candidate cubes containing that point and the level datum at that point. Under
the rule $(\mathrm{G}^{\natural})$ of Proposition~\ref{4.6:prop-gluing-example}, the normalisation
$1/\sqrt{N_{\mathcal{F}}}$ across the seam depends only on data of this kind. The two nests of a
comparison carry the same gluing rule $\Phi^h$.
\item[($\sigma4$)] \emph{Thin rings.} Write $Z''_j$ for the regions produced by the completion in
\cite{B-CMP122a}; that completion separates $Z''_j$ and $Z''_{j-1}$ by a single layer of $M$-cubes. In
this item the letters $\eta$, $M$, $R$ are those of \cite{B-CMP96}: $\eta$ is its lattice spacing,
$M$ is the size of its big blocks, and $R$ is its parameter, ``a big positive integer''; $M$ and $R_k$
denote also the parameters so lettered in \cite{B-CMP122a}. The letters $R$, $R_k$ here denote neither
the ball radius nor the witness constant $R_K$. The condition \cite[(2.2)]{B-CMP96} asks that
\begin{equation*}
(L^j\eta)^{-1}\operatorname{dist}(\Omega_j^c,\Omega_{j+1})>RM .
\end{equation*}
Identify $M$, $R$ of \cite{B-CMP96} with the parameters $M$, $R_k$ of \cite{B-CMP122a}, which are the
parameters of the $M$-cubes and of the $MR_k$-cubes of \cite{B-CMP119}; under this identification the
margin $RM$ of \cite[(2.2)]{B-CMP96} becomes $R_kM$. Whether the nests built
from the regions $Z''_j$
then satisfy \cite[(2.2)]{B-CMP96} is a property of a single nest, and it is not settled here.
That the nests attached to the successive steps of the large-field construction are admissible is
taken here, by statement, from the admissibility statement for those nests, admissibility being
understood in the sense fixed in that statement; this section does not use that statement to decide
the question just stated. For the local
sequences, \cite{B-CMP99b} states that the conditions \cite[(2.1)--(2.2)]{B-CMP96} hold for them;
whether the local sequences of Definition~\ref{4.6:def-local-sequence} meet them is likewise a
property of a single nest and is not settled here. No statement of this section comparing two
nests depends on this question. The boundary separation \eqref{4.6:eq-nest-traces-a} is a
qualitative counterpart of \cite[(2.2)]{B-CMP96}. It is the only separation condition read by
Lemma~\ref{4.6:lem-seam-double-count} and by $(\mathrm{COV}^{\natural})$,
\eqref{4.6:eq-gluing-example-a}, restricted to $\mathcal{D}$, and it holds on every member of the
class $\mathfrak{N}$ of nests with region data, by the boundary-separation statement for that class.
\item[($\sigma5$)] \emph{High-level averaging blocks adjacent to small cubes.} Only the footprint
clause and the flanked case of the support clause of the definition of letters and words take such
blocks into account, and only through the footprint display \eqref{4.6:eq-letters-words-c}.
\item[($\sigma6$)] \emph{The margin $RM$ for region data.} Ba{\l}aban's treatment of region data in
\cite{B-CMP99b} carries a margin $RM$, with $R$, $M$ as in ($\sigma4$). It is a one-sided geometric
lower bound on the region datum.
Whether the region data of this chapter satisfy it is not settled here; it is not imposed on them.
\end{itemize}
\end{remark}

\begin{proof}
We justify in turn the assertions of ($\sigma1$)--($\sigma5$) that are not definitions.

($\sigma1$). By Lemma~\ref{4.6:lem-sequence-reach}, which is stated under the reading convention on
local sequences in force in this chapter, the local sequence is obtained from membership bits by a
label-free map $\Sigma$:
\begin{equation*}
\mathbf{\Omega}(Q)=\Sigma\bigl(Q,g;\,([e\subset\Omega_{g+1}],[e\subset\Omega_{g+2}])_{e}\bigr).
\end{equation*}
Here $e$ runs over the $\pi_g$-blocks $e\subset Q^{(3)}$. These bits are exactly what an inner seam at
$Q$ contributes, and the seam datum $T(Q)$ of that lemma is built from them. Consequently, two nests
whose traces agree on $Q^{(3)}$ at levels $g+1$, $g+2$ have the same local sequence at $Q$, as a tuple
of sets. Proposition~\ref{4.6:prop-generator-equality} then gives equality of the generators at $Q$
whenever the bits agree at levels $g$, $g+1$, $g+2$ on the blocks of $Q^{(3)}$. This is the exactness
asserted. The value $\Omega_{g+1}(Q')=Q'^{(4)}$ for a protruding cube
$Q'\in\mathcal{D}_{g+1}(\mathbf{\Omega})$ holds by Definition~\ref{4.6:def-local-sequence}. For a cube
of a wavy nest, Lemma~\ref{4.6:lem-sequence-reach} gives the datum in its wavy form
$\widetilde{T}(Q)$, built from the region bits of the blocks $e\subset Q^{(3)}$; the seam
$\partial\Lambda$ is encoded there in the same way.

($\sigma2$). The arguments of $\Sigma(Q,g;\cdot)$ displayed above are bits at levels $g+1$ and $g+2$
only. The boundary of $\Omega_g$ is encoded by bits at level $g$, so it is not an argument of
$\Sigma$. Now suppose, as in ($\sigma2$), that a different local sequence were attached to outer-seam
cubes, determined by hypothesis by the bits of the blocks $e\subset Q^{(3)}$ at levels
$g-1,\dots,g+2$. Compared with
the arguments of $\Sigma$, the new arguments are the bits at levels $g-1$ and $g$ of the same blocks.
The datum would therefore remain a family of bits on $Q^{(3)}$,
and the reading region $Q^{(3)}$, and with it every reading radius derived from it, would be
unchanged.

($\sigma3$). The two surcharges are the entries of the reading table,
Lemma~\ref{4.6:lem-reading-table}, for the surgeries and for the glued bumps. For the surgeries see
Lemma~\ref{4.6:lem-surgery-surcharge}. For the bumps see Lemma~\ref{4.6:lem-gluing-window}, by which
the glued bump $h_Q[\mathcal{F}]$ depends on the nest only through its bits on $Q^{+2M}$, the
remaining data being $(Q,j(Q))$ and the surgery datum.
Lemma~\ref{4.6:lem-seam-double-count} shows that across a layer seam the rescaled profiles do not by
themselves form a partition of unity on $\mathcal{D}$. The gluing rule $\Phi^h$ is therefore needed
exactly there, and with it the hypothesis \eqref{4.6:eq-bump-profiles-c}. The description of what
$(\mathrm{G\text{-}loc})$ reads is the content of that hypothesis. The normalisation $1/\sqrt{N_{\mathcal{F}}}$
of $(\mathrm{G}^{\natural})$ is of that type by Proposition~\ref{4.6:prop-gluing-example}. That both
nests carry the same $\Phi^h$ is part of Definition~\ref{4.6:def-one-recipe}; see
\eqref{4.6:eq-one-recipe-a}.

($\sigma4$). Part~(b) of Lemma~\ref{4.6:lem-seam-double-count} is proved under the hypothesis
\eqref{4.6:eq-nest-traces-a} at level $g$ and under no other separation hypothesis. The same
separation condition is the one read by $(\mathrm{COV}^{\natural})$ restricted to $\mathcal{D}$. On
every member of $\mathfrak{N}$ the condition \eqref{4.6:eq-nest-traces-a} holds by the
boundary-separation statement for that class. Theorem~\ref{4.6:thm-expansion-locality}, the statement
of this section comparing two nests, and the statements on which it depends read the two nests only
through their traces and the conditions just named; none of them has a clause of
\cite[(2.2)]{B-CMP96} among its hypotheses. Hence no such statement depends on whether
\cite[(2.2)]{B-CMP96} holds.

($\sigma5$). The letters see averaging blocks only through the footprints of
\eqref{4.6:eq-letters-words-c}. Among the clauses describing letters and words, a high-level
averaging block adjacent to a small cube enters these footprints only in the footprint clause and in
the flanked case of the support clause.
\end{proof}

\begin{lemma}[Placement of the enlarged reading window under the scale floor]\label{4.6:lem-window-placement}
Let $r\ge r_*$, where $r_*=2$ is the reading radius fixed in \eqref{4.6:eq-letters-words-a}. Assume the following.
\begin{itemize}
\item[(a)] The atom rule has the locality property of Definition~\ref{4.6:def-atom-locality} at radius $r$, together with its companion statement for the two nests in force; both are displayed in \eqref{4.6:eq-atom-locality-a}. The rule fixed in this section has this property for every $r\ge r_*$; any atom rule with this property may be used instead.
\item[(b)] The scale floor holds: $\check{R}_s\ge r$ at every step $s\le k$, that is, $rM\le a_s$, with $\gamma$ chosen last.
\end{itemize}
Let $F$ be a term of the expansion at step $k$ attached to the pair $(X,m)$. Let $\mathfrak{L}$ and $\mathfrak{L}'$ be two labels, with nests in force $\mathbf{\Omega}^*[\mathfrak{L}]$ and $\mathbf{\Omega}^*[\mathfrak{L}']$ and region data $\mathfrak{R}$ and $\mathfrak{R}'$, such that
\begin{equation*}
\mathrm{TRACE}^{\flat}_X(\mathfrak{L})=\mathrm{TRACE}^{\flat}_X(\mathfrak{L}').
\end{equation*}
Let $W'\subset X$ be a closed union of cells with $W'=\operatorname{cl}\operatorname{int}W'$. Then the following hold.
\begin{enumerate}
\item[(1)] \emph{Geometry.} Every $\pi$-block $c'\subset W'^{+rM}$ lies in an $\mathfrak{a}_k$-cell $E\in[X]_k^{+1}$. If $m<k$ and statement (T2) below holds, then $E$ is either a cell of a component $C$ of $\Lambda_k^c$ meeting $X$, or a $\Lambda_k$-cell touching such a component $C$.
\end{enumerate}
The following three statements enter as hypotheses: (T2) in the second sentence of part (1), and all three in parts (2) and (3).
\begin{itemize}
\item[(T1)] For every component $C$ of $\Lambda_k^c$ meeting $X$, item (i) of $\mathrm{TRACE}^{\flat}_X$ records in full the datum $\mathfrak{d}(C)$ of the component $C$. This datum determines, on every cell of $C$, every member of every nest in force, and it determines the region data there.
\item[(T2)] For every $\mathfrak{a}_k$-cell $e\subset\Lambda_k$ and every nest $\mathbf{\Omega}^*$ in force, $\Omega^*_i\cap e=e$ for all $i\le k$. This holds independently of the label, and the region data on $e$ are trivial.
\item[(T3)] Suppose $m=k$, and let $e$ be a cell of $[X]_k^{+1}$ which is a cell of a component of $\Lambda_k^c$ not meeting $X$. Then item (ii) of $\mathrm{TRACE}^{\flat}_X$ carries on $e$ the set-traces of the members of the nests in force and of the sets constituting the region data, with their names and marks.
\end{itemize}
Statements (T1) and (T3) describe what the items (i) and (ii) of the trace datum $\mathrm{TRACE}^{\flat}_X$ record, and (T2) is the full-trace property of the cells of $\Lambda_k$ at every nest in force.
Under these additional assumptions:
\begin{enumerate}
\item[(2)] \emph{Traces.} The traces on the enlarged window agree:
\begin{equation}\label{4.6:eq-window-placement-1}
\operatorname{tr}_{W'^{+rM}}\bigl(\mathbf{\Omega}^*[\mathfrak{L}];\mathfrak{R}\bigr)
=\operatorname{tr}_{W'^{+rM}}\bigl(\mathbf{\Omega}^*[\mathfrak{L}'];\mathfrak{R}'\bigr).
\end{equation}
\item[(3)] \emph{Compressions.} For every operator $K$ to which Definition~\ref{4.6:def-atom-locality} applies,
\begin{equation}\label{4.6:eq-window-placement-2}
K^{W'}[\mathfrak{L}]=K^{W'}[\mathfrak{L}'],
\end{equation}
word by word. Every set on which these compressions depend lies in the set
\begin{equation*}
\bigl|[X]_k^{+1}\bigr|\ \cup\ \bigcup\{|C| : C \text{ a component of } \Lambda_k^c \text{ meeting } X\}\ \subset\ \mathfrak{c}(X).
\end{equation*}
In particular, the collar $\mathfrak{c}(X)$ is unchanged.
\end{enumerate}
\end{lemma}

\begin{proof}
\emph{Part} (1). Let $c'\subset W'^{+rM}$ be a $\pi$-block. Let $E$ be the $\mathfrak{a}_k$-cell containing $c'$; it exists by the nesting of the block grid in the grid of $\mathfrak{a}_k$-cells. Choose a point $y\in\operatorname{int}c'\subset\operatorname{int}E$. Since $y\in W'^{+rM}$, there is a point $p\in W'$ with $|y-p|_\infty\le rM$. By the scale floor (b) at $s=k$ we have $rM\le a_k$, hence
\begin{equation*}
|y-p|_\infty\le rM\le a_k .
\end{equation*}

We next show that $p$ lies in a cell of $[W']_k$. A point of $\operatorname{int}W'$ lies in an $\mathfrak{a}_k$-cell whose interior meets $\operatorname{int}W'$. Hence $\operatorname{int}W'\subset|[W']_k|$. The set $|[W']_k|$ is a finite union of closed cells, so it is closed. Therefore
\begin{equation*}
p\in W'=\operatorname{cl}\operatorname{int}W'\subset|[W']_k| .
\end{equation*}
So $p$ lies in some cell $E_p\in[W']_k$. Since $\operatorname{int}W'\subset\operatorname{int}X$, we have $[W']_k\subset[X]_k$, and so $E_p\in[X]_k$.

We now use the arithmetic of the cells. Write $E_p=z+[0,a_k]^d$ with $z\in a_k\mathbb{Z}^d$. Then
\begin{equation*}
E_p^{+a_k}=z+[-a_k,2a_k]^d=\bigcup_{v\in\{-1,0,1\}^d}\bigl(z+a_kv+[0,a_k]^d\bigr).
\end{equation*}
The right-hand side is the union of the $3^d$ cells touching or equal to $E_p$. Since $p\in E_p$ and $|y-p|_\infty\le a_k$, we have $y\in p^{+a_k}\subset E_p^{+a_k}$. Hence $y$ lies in one of these cells; call it $E'$. Now $y\in\operatorname{int}E$, and $E'$ is a closed cell of the same grid containing $y$. A point of the interior of a cell lies in no other cell of the grid, so $E'=E$. Thus $E$ touches or equals $E_p\in[X]_k$, which gives $E\in[X]_k^{+1}$.

Let now $m<k$. By its choice, the cell $E_p$ meets $X$. By the nesting of the $\pi_m$-grid in the grid of $\mathfrak{a}_k$-cells, some $\pi_m$-cube $q$ of $X$ satisfies $q\subset E_p$. Suppose $E_p$ were a $\Lambda_k$-cell. Since $m+1\le k$, statement (T2) would give
\begin{equation*}
q\subset E_p=\Omega^*_{m+1}\cap E_p\subset\Omega^*_{m+1}.
\end{equation*}
This contradicts the cube-wise inclusion $X\subset(\Omega^*_{m+1})^c$ that holds for a region to which a term is attached at level $m$. Hence $E_p$ is a cell of $\Lambda_k^c$. It belongs to a component $C$ of $\Lambda_k^c$, and $C$ meets $X$ because $E_p$ does.

By the first half of this proof, $E$ touches or equals $E_p$. There are two cases.
\begin{itemize}
\item $E$ is a $\Lambda_k$-cell. Then it touches the cell $E_p$ of $C$, that is, it touches $C$.
\item $E$ is a cell of $\Lambda_k^c$, belonging to a component $C'$. Then $C'$ contains a cell, namely $E$, which touches or equals a cell of $C$, namely $E_p$. The components of $\Lambda_k^c$ are the classes generated by the relation ``intersect or touch'' among its cells, so distinct components do not touch. Hence $C'=C$, and $E$ is a cell of the component $C$, which meets $X$.
\end{itemize}
This proves part (1).

\emph{Part} (2). Let $c'\subset W'^{+rM}$ be a $\pi$-block, and let $E$ be the cell given by part (1). We show that the traces of the two labels coincide on $E$. Consider first $m<k$. By part (1), $E$ is then a cell of a component meeting $X$, or a $\Lambda_k$-cell; these are cases (A) and (B) below. Consider next $m=k$. Then $E\in[X]_k^{+1}$, and one of three things holds: $E$ is a cell of a component meeting $X$, $E$ is a $\Lambda_k$-cell, or $E$ is a cell of a component of $\Lambda_k^c$ not meeting $X$. These are cases (A), (B) and (C).
\begin{itemize}
\item[(A)] $E$ is a cell of a component $C$ meeting $X$. By (T1), item (i) of $\mathrm{TRACE}^{\flat}_X$ records in full the datum $\mathfrak{d}(C)$ of $C$. This datum determines on $E$ every member of every nest in force and the region data. The item is the same for $\mathfrak{L}$ and $\mathfrak{L}'$, by the hypothesis on the trace data.
\item[(B)] $E\subset\Lambda_k$. By (T2), every nest in force has the full trace on $E$ for $\mathfrak{L}$ and for $\mathfrak{L}'$ alike, independently of the labels, and the region data on $E$ are trivial.
\item[(C)] $m=k$ and $E$ is a cell of a component not meeting $X$. Since $m=k$, the $\mathfrak{a}_m$-cells are the $\mathfrak{a}_k$-cells, so $E$ is a cell of $[X]_k^{+1}$ not inside a component meeting $X$. By (T3), item (ii) of $\mathrm{TRACE}^{\flat}_X$ carries on $E$ the set-traces of the members of the nests in force and of the sets constituting the region data, with their names and marks. This item is the same for the two labels.
\end{itemize}
In every case the cell-level traces on $E$ of the nests in force and of the region data are the same for $\mathfrak{L}$ and $\mathfrak{L}'$. The $\pi$-bits on $c'\subset E$ are read off these cell-level traces, by the nesting of the $\pi$-grid in the cell grid. Hence the membership bits of $c'$ agree for the two labels. The block $c'\subset W'^{+rM}$ was arbitrary, and this gives \eqref{4.6:eq-window-placement-1}.

When $m<k$, the window $W'^{+rM}$ may contain $\Lambda_k$-cells which are adjacent to components meeting $X$ and lie outside $|[X]_m^{+1}|$. On such cells every nest in force carries the trivial full trace, by (T2). Hence no datum of $\mathrm{TRACE}^{\flat}_X$ beyond items (i) and (ii) is needed on them.

\emph{Part} (3). Apply the locality property (a) with $W:=W'$, together with its companion statement for the two nests in force, \eqref{4.6:eq-atom-locality-a}. The hypothesis of that property is the equality of traces on the enlarged window, and this is \eqref{4.6:eq-window-placement-1}. The property therefore yields \eqref{4.6:eq-window-placement-2} word by word. It coincides, word for word, with the equality of compressions that equal trace data $\mathrm{TRACE}^{\flat}_X$ are used to produce at radius $r_*$; here it is obtained at radius $r$ in place of $r_*$.

The compressions depend only on data on the footprints of \cite[p.~3]{B-CMP116}, which lie in $W'$, enlarged by $rM$, hence only on data inside $W'^{+rM}$. By part (1), every $\pi$-block of $W'^{+rM}$ lies in a cell of $[X]_k^{+1}$, and when $m<k$ this cell is a cell of a component meeting $X$ or a $\Lambda_k$-cell touching such a component. Hence every set on which the compressions depend lies in $|[X]_k^{+1}|$ together with the components meeting $X$, and this set is contained in the collar $\mathfrak{c}(X)$. So the collar $\mathfrak{c}(X)$ is unchanged.
\end{proof}

\begin{remark}[Sharpness of the scale floor in one dimension]\label{4.6:rem-floor-sharp}
The scale floor of Lemma~\ref{4.6:lem-window-placement} is sharp, as the following
one-dimensional example shows. Take $d=1$ and measure lengths in units of $M$, so that the floor reads
$r\le\check{R}$. The cells are the intervals $[a\check{R},(a+1)\check{R}]$, $a\in\mathbb{Z}$, where
$\check{R}$ is an integer. In this example the cell hull $[Y]^{+1}$ of a union $Y$ of cells in
Lemma~\ref{4.6:lem-window-placement} is the union of the cells whose index differs by at most $1$
from the index of a cell of $Y$; for an integer $t\ge1$ we write $[Y]^{+t}$ for the union of the
cells whose index differs by at most $t$ from the index of a cell of $Y$.
Put $X:=[0,\check{R}]$ and $W':=[\check{R}-2,\check{R}]$; the requirement $W'\subset X$ of
Lemma~\ref{4.6:lem-window-placement} forces $\check{R}\ge2$. Then $[X]^{+1}=[-\check{R},2\check{R}]$ and
\begin{equation}\label{4.6:eq-floor-sharp-1}
W'^{+r}=[\check{R}-2-r,\;\check{R}+r]\subset[-\check{R},2\check{R}]
\quad\Longleftrightarrow\quad r\le\check{R}.
\end{equation}
If the floor is dropped, the argument of Lemma~\ref{4.6:lem-window-placement} places every block
$c'\subset W'^{+r}$ only in $[X]^{+\rho}$ with $\rho=\lceil r/\check{R}\rceil$, which for
$r>\check{R}$ is strictly larger than $[X]^{+1}$. This weaker
placement is not used in what follows.
\end{remark}

\begin{proof}
Since $\check{R}\ge2$, the interval $W'$ lies in $X$. It is the union of the unit blocks
$[\check{R}-2,\check{R}-1]$ and $[\check{R}-1,\check{R}]$, and it is the closure of its interior.
The set $X$ is the single cell with index $0$. The cells with index in $\{-1,0,1\}$ therefore have
union $[X]^{+1}=[-\check{R},2\check{R}]$.

In one dimension the $\ell^\infty$ enlargement $W'^{+r}$ is the closed $r$-neighbourhood of $W'$,
which gives the first equality in \eqref{4.6:eq-floor-sharp-1}. The inclusion holds if and only if
both $\check{R}-2-r\ge-\check{R}$ and $\check{R}+r\le2\check{R}$ hold, that is, if and only if
\begin{equation*}
r\le 2\check{R}-2 \quad\text{and}\quad r\le\check{R}.
\end{equation*}
Since $\check{R}\ge2$, we have $\check{R}\le2\check{R}-2$, so the second condition implies the
first. This proves the equivalence in \eqref{4.6:eq-floor-sharp-1}. In particular, if
$r>\check{R}$, then the point $\check{R}+r$ of $W'^{+r}$ lies outside $[X]^{+1}$.

For the placement without the floor, note that in this example
$[X]^{+\rho}=[-\rho\check{R},(1+\rho)\check{R}]$. By the definition of
$\rho=\lceil r/\check{R}\rceil$ we have $\rho\check{R}\ge r$. Hence
$\check{R}+r\le(1+\rho)\check{R}$, and
\begin{equation*}
\check{R}-2-r\ge -r\ge-\rho\check{R}.
\end{equation*}
Therefore $W'^{+r}\subset[X]^{+\rho}$.
\end{proof}

\begin{remark}[Reading radius two needs no new scale floor]\label{4.6:rem-radius-two-floor}
Let $k$ be a level and $X$ a large-field region of level $m\le k$, read at step $k$ as in
Lemma~\ref{4.6:lem-window-placement}. Let $\check{R}_s$ be the cell
multipliers, so that the large-field cells at step $s$ have side $a_s=M\check{R}_s$, and let
the glued bumps obey $(\mathrm{G\text{-}loc})$ and $(\mathrm{G\text{-}supp})$ of
\eqref{4.6:eq-bump-profiles-c} in Definition~\ref{4.6:def-bump-profiles}. Consider the point of
the argument at which, when two nests carry equal trace
data $\mathrm{TRACE}^{\flat}_X$ attached to $X$, the compressions to any closed union of cells
$W'\subset X$ of the operators of the renormalisation step agree term by term. This is the point
at which the
atom hypothesis \eqref{4.6:eq-atom-locality-a} enters. When that agreement is obtained from
Lemma~\ref{4.6:lem-window-placement}, the reading radius is $r_*=2$, and the only condition
the lemma places on the cell multipliers is
\begin{equation}\label{4.6:eq-radius-two-floor-1}
\check{R}_s\ge 2\qquad\text{at every step } s\le k .
\end{equation}
This floor is implied by the floor $\check{R}_s\ge 3$ at every step $s\le k$, which is among
the scale conditions of the theorem on the locality of the compressions in the trace datum
$\mathrm{TRACE}^{\flat}_X$, whose proof uses this agreement. It coincides with the
floor $\check{R}_s\ge r_\zeta$, where
$r_\zeta$ denotes the window of the cut-off functions $\zeta_{\square}$; for the family of
cut-offs used here, $r_\zeta=2$, by the lemma on the window of the cut-offs $\zeta_{\square}$.
No further
floor is needed, and no collar is needed beyond the one on which the trace datum
$\mathrm{TRACE}^{\flat}_X$ is read.
\end{remark}

\begin{proof}
We first fix the radius. By Lemma~\ref{4.6:lem-gluing-window}, every gluing rule of type
$(\mathrm{G\text{-}loc})\wedge(\mathrm{G\text{-}supp})$, that is, every rule satisfying the two
conditions \eqref{4.6:eq-bump-profiles-c} of Definition~\ref{4.6:def-bump-profiles}, has window
$r_h=2$. That is, within
that lemma's scope $u\le M$, each glued
bump $h_Q[\mathcal{F}]$ and its lattice differences are determined by the member $Q$ and its
scale $j(Q)$, by the surgery datum, and by the membership bits of the nest on $Q^{+2M}$ when
$Q$ is a candidate cube, and on the correspondingly enlarged set of
Lemma~\ref{4.6:lem-gluing-window} when $Q$ is an inserted cube. Hence
the reading radius in the atom hypothesis \eqref{4.6:eq-atom-locality-a} may be taken to be
\[
r=r_*,\qquad r_*=\max\{2,r_h\}=2 ;
\]
the hypothesis itself is the statement displayed there and is not derived here.

We next apply Lemma~\ref{4.6:lem-window-placement} with $r=2$. Its hypothesis on the scale is
\eqref{4.6:eq-radius-two-floor-1}, which is the inequality $2M\le a_s$ for $s\le k$.
Since $\check{R}_s\ge 3$ at every step $s\le k$, by the scale conditions recalled in the
statement above, the inequality \eqref{4.6:eq-radius-two-floor-1} holds. Since $r_\zeta=2$ for
the family of cut-offs used here, by the lemma on the window of the cut-offs $\zeta_{\square}$,
it is the
same inequality as $\check{R}_s\ge r_\zeta$. No new floor therefore arises.

We turn to the collar. Let $\omega$ be a word of the compression to $W'$; its footprint
$[\omega]$ lies in $W'$. By clause~(1) of Lemma~\ref{4.6:lem-window-placement} with $r=2$, the
traces on $[\omega]^{+2M}\subset W'^{+2M}$ that the term-by-term agreement requires are part of
the data $\mathrm{TRACE}^{\flat}_X$. No collar is therefore needed beyond the one on which the
trace datum $\mathrm{TRACE}^{\flat}_X$ of Lemma~\ref{4.6:lem-window-placement} is read.

The radius would be $r_*=4$, with the floor $\check{R}_s\ge 4$, if the glued bumps were known
only to satisfy $\operatorname{supp}h_Q\subset Q^{1}$, the double cube of $Q$. That case does
not occur. For
the profiles of Definition~\ref{4.6:def-bump-profiles}, under the reading
\eqref{4.6:eq-bump-profiles-b}, the support of $\eta_Q$ lies
inside $Q$. For the glued bump, $(\mathrm{G\text{-}supp})$ in \eqref{4.6:eq-bump-profiles-c}
states that its support lies inside $Q$.

The support part of the atom hypothesis enters at one further place, the localisation of the
footprints of the atoms. There no reading radius $r$ occurs: the margin flanking the footprint
reads a single label of the data, as the footprint itself does. In the refined
atom families, which are not used in the argument above, the atoms read the names carried by
the trace datum $\mathrm{TRACE}^{\flat}_X$ only within distance $2M\le r_\zeta M$, because
$r_\zeta=2$. They therefore impose
nothing beyond \eqref{4.6:eq-radius-two-floor-1}.

We record, finally, what happens for a gluing rule outside the type
$(\mathrm{G\text{-}loc})\wedge(\mathrm{G\text{-}supp})$, for instance bumps renormalised against
members of the family that do not contain the point. Suppose such a rule has window $r_h$. Then
the same argument gives $r_*=\max\{2,r_h\}$ and the floor $\check{R}_s\ge r_*$ at every step
$s\le k$. This floor is of the same kind as $\check{R}_s\ge r_\zeta$: it bounds $\check{R}_s$
from below only, it bounds the parameter $N_0$ entering the scale conditions from below only,
it is imposed after $L$ with $\gamma$ chosen
last, and it is not a cap in the sense of Definition~\ref{2.3:def-cap}. It is implied by
$\check{R}_s\ge 3$ if and only if $r_h\le 3$. The gluing of
Definition~\ref{4.6:def-bump-profiles} and \eqref{4.6:eq-bump-profiles-c} is of the type
$(\mathrm{G\text{-}loc})\wedge(\mathrm{G\text{-}supp})$, and no rule outside it occurs.
Throughout, the footprint of a word $\omega$ is $[\omega]$ and not its nest-reading domain; the
argument above uses only the former.
\end{proof}

\begin{definition}[The renormalisation step, the level gap and the collars]\label{4.7:not-step-gap}
Let $k$ be the current level and let $i\le k+1$.

\emph{The step.} The step $i-1\to i$ is the step of that name in the multiscale construction of
\cite[\S3]{B-CMP119}. It consists of the
analysis of \cite[\S\S3--5]{B-CMP109} together with that of \cite[\S\S1--2]{B-CMP116}, both
carried out on the multiscale cover of \cite[\S3]{B-CMP119}. In the words of
\cite[\S3]{B-CMP119}, in which the reference [I] is \cite{B-CMP109},
\begin{quote}
we obtain all the formulas and bounds of Sects.~3, 4 [I] with $\eta$ replaced by $L^{-n}$
\end{quote}
(in this sentence the letter $n$ is the layer index, which equals $i-1$ on the layer
$\Lambda_i$), and
\begin{quote}
the terms of the expansion (I.6.41), with localization domains $Y$ contained in
$\Lambda_{k+1}$, satisfy exactly all the conditions of those lemmas.
\end{quote}
\cite[\S3]{B-CMP119} takes over the restrictions that \cite{B-CMP109} imposes on its
constants, and they are in force throughout the step.

\emph{Old levels and the gap.} A level $i'$ with $i'\le i-1$ is called an \emph{old level}.
For an old level $i'$ the \emph{gap} is
\begin{equation}\label{4.7:eq-step-gap-1}
n:=i-1-i' .
\end{equation}
From here to the end of this chapter $n$ denotes the gap \eqref{4.7:eq-step-gap-1}, and not
the layer index
of the quotation above. We put $\eta=L^{-(i-1)}$ and, for every level $m$ (the subscript here
being a level, not the torus exponent of $\mathbb{T}_m$), $\xi_m:=L^{-m}$, the fine spacing
in level-$m$ units. The factors $L^{j}\eta$ and $L^{j-1}\eta$ of \cite[\S3]{B-CMP119}, taken
at $j=i'$, are then expressed through the gap:
\begin{equation}\label{4.7:eq-step-gap-2}
L^{i'}\eta=L^{i'-i+1}=L^{-n},
\qquad
L^{i'-1}\eta=L^{-(n+1)}=\frac{L^{-i}}{L^{-i'}}=\frac{\xi_i}{\xi_{i'}} .
\end{equation}

\emph{Domain classes.} For each level $j$ let $\pi_j$ be the partition into cubes used at
level $j$ in \cite{B-CMP116}, and let $\mathbf{D}_j$ be the class of localisation domains at
level $j$ of \cite{B-CMP116}, which are built from $\pi_j$-cubes. In the step
$i-1\to i$ the new domains are $Y\in\mathbf{D}_{i-1}$, built from $\pi_{i-1}$-cubes, and the
old domains are $X\in\mathbf{D}_{i'}$. For such an $X$ we write $X_0$ for ``the smallest
localization domain from $\mathbf{D}_k$ containing $X$'' \cite[\S\S1--2]{B-CMP116}, where the
letter $k$ of \cite{B-CMP116} stands for the level $i-1$, so that $X_0\in\mathbf{D}_{i-1}$.

\emph{The own-cubes collar.} For a level $p$ and a domain $X\in\mathbf{D}_p$ carrying a space
of level-$p$
configurations, the domain
$\tilde{X}^{-2}$ of \cite[p.~262]{B-CMP109} is $X$ less the two layers of its own
cubes, that is, of $\pi_p$-cubes, nearest to the boundary of $X$; the sentence of
\cite[p.~262]{B-CMP109} reads
\begin{quote}
The domain $\tilde{X}^{-2}$ is obtained from $X$ by taking away two layers of cubes from
$\pi_j$, which are closest to the boundary of $X$.
\end{quote}
(with $j=p$ in the quoted sentence).
For a space of level-$p$ configurations on a domain $Z\in\mathbf{D}_{p-1}$ (the domain $Y$ of
\cite[\S\S1--2]{B-CMP116}), the two layers removed are $\pi_{p-1}$-layers. This is the practice
of the papers: the window $\square_0=\tilde{\square}^{5}$ of \cite[\S3]{B-CMP109},
$\tilde{\square}^{l}$ denoting the $l$-th enlargement of the cube $\square$ in the sense of
\cite[\S3]{B-CMP109}, is cut in \cite[(3.40)]{B-CMP109} to $\tilde{\square}^{3}$, that is, to
$\square_0$ less two layers of cubes of the partition written $\pi_k$ in \cite[\S3]{B-CMP109},
$k$ there being that paper's level; and \cite[\S\S1--2]{B-CMP116} writes
\begin{quote}
we consider this space defined by the conditions I.(i)--(iii) on the domain $Y$.
\end{quote}
The cut operative in this chapter is not $\tilde{X}^{-2}$ but the cut $Z^{\wr p}$ fixed with
the one-power hereditary class: for every pair $(p,Z)$ of a level $p$ and a
domain $Z$ in the finite label set of that class, $Z^{\wr p}$ is the union
of the $\pi_{p-1}$-cubes $\square\subset Z$ with $\tilde{\square}^{2}\subset Z$;
it is cut in $\pi_{p-1}$-layers also when $Z\in\mathbf{D}_p$. For $X\in\mathbf{D}_p$ one has
$\tilde{X}^{-2}\subset X^{\wr p}$, so that the conditions imposed on $X^{\wr p}$ are an
additional demand beyond those imposed on $\tilde{X}^{-2}$. Consequently the one-power
hereditary class is contained in the class defined literally by the collar
$\tilde{X}^{-2}$ of \cite[p.~262]{B-CMP109}.
\end{definition}

\begin{definition}[Running couplings, radii and fluctuation size]\label{4.7:not-running-couplings}
Throughout this section $k$ is the level of the step under consideration, with the step and the
levels as in Definition~\ref{4.7:not-step-gap}. In this section the index $m$ denotes a level, not
the torus exponent.

\emph{Couplings.} For $0\le m\le k$, $g_m$ denotes the running coupling at level $m$ of the
history actually produced by the steps performed, and $x_m:=g_m^{-2}$. In the present induction the
flow constants $\lambda_{\sharp}$ and $B_{\sharp}$ of the glossary (Notation~\ref{0.2:not-glossary})
are
\begin{equation*}
\lambda_{\sharp}=\frac{c_0}{2}-\bar{\varsigma}\ \ge\ \frac{c_0}{4},
\qquad
B_{\sharp}=C_{\beta}+\bar{\varsigma}\ \le\ C_{\beta}+1,
\end{equation*}
where $\bar{\varsigma}$ is the ceiling that the induction on the coupling increment imposes on the
part of the increment descended from large-field regions. The two-sided flow bounds that this
induction takes as its hypothesis at step $k$ read, in the letters just introduced,
\begin{equation}\label{4.7:eq-running-couplings-1}
\lambda_{\sharp}(b-a)-c_5\ \le\ x_a-x_b\ \le\ B_{\sharp}(b-a),
\qquad 0\le a\le b\le k .
\end{equation}

\emph{Radii.} The radii of \cite[(2.28)]{B-CMP119} are
\begin{equation*}
\alpha_{0,j}=g_jC_0\bigl(\log g_j^{-2}\bigr)^{q_0},
\qquad
\alpha_{1,j}=g_jC_1\bigl(\log g_j^{-2}\bigr)^{q_1}.
\end{equation*}
We take $q:=q_0=q_1$, as fixed in the order of choice of the parameters, and set
\begin{equation}\label{4.7:eq-running-couplings-2}
r_m:=\min\{C_0,C_1\}\bigl(\log g_m^{-2}\bigr)^{q}g_m .
\end{equation}
Since $C_0\le C_1$, the minimum in \eqref{4.7:eq-running-couplings-2} is $C_0$, and with
$q_0=q$ the right-hand side is $g_mC_0(\log g_m^{-2})^{q_0}$; thus $r_m=\alpha_{0,m}$.

\emph{Index convention.} A letter carrying a level index $i$ (such as $\alpha_{0,i}$,
$\alpha_{1,i}$, $r_i$, and $\delta_i$ below) denotes the quantity defined in this section computed
from the true coupling $g_i$ whenever $i\le k$. At $i=k+1$ every such letter is computed with the
stand-in
\begin{equation*}
\tilde{g}:=g_k
\end{equation*}
in place of $g_{k+1}$.

\emph{Fluctuation size.} With the degree function $p_0(\cdot)$,
$p_0(g')=A_0(\log g'^{-2})^{p_0}$ \cite{B-CMP119}, put $\varepsilon_m:=g_m\,p_0(g_m)$ and set
\begin{equation}\label{4.7:eq-running-couplings-3}
\delta_m:=\frac{A_1}{A_0}\,\varepsilon_m=A_1\bigl(\log g_m^{-2}\bigr)^{p_0}g_m ,
\end{equation}
the second expression being the first with $\varepsilon_m$ and $p_0(g_m)$ written out. At $m=k$
this is the quantity $\delta_k=g_kA_1A_0^{-1}p_0(g_k)$ of \cite{B-CMP119}, the threshold of
\cite[(3.16)]{B-CMP119}. The complex box of \cite[(3.43)]{B-CMP119} is printed there as
$|\mathbf{A}|<C_1p_1(g_k)$, with $\mathbf{A}$ the fluctuation variable and $p_1(\cdot)$
Ba{\l}aban's second degree function; we write $C_1^{\mathrm{fl}}$ for the constant printed $C_1$
in this box, to keep it apart from the radius constant $C_1$ above. At level $m$ the fluctuation
size on the box is $C_1^{\mathrm{fl}}\delta_m$, where $g_mp_1(g_m)$ is identified with $\delta_m$;
a smaller power of the logarithm in $p_1(\cdot)$ than in $p_0(\cdot)$ only improves the bound
$C_1^{\mathrm{fl}}\delta_m$.
The constant $C_1^{\mathrm{fl}}$ is chosen in the order of choice of the parameters before $C_0$.
The quantity written $\delta'_j$ in \cite{B-CMP122a} belongs to the large-field operation
$\mathbf{R}$ and is not used here.

Finally, the exponents $q$ and $p_0$ are subject to the one-sided condition
\begin{equation}\label{4.7:eq-running-couplings-4}
q-p_0\ \ge\ 3r+2\ >\ 0 ,
\end{equation}
where $r$ is the fixed exponent of Ba{\l}aban's construction that also enters the printed
one-sided condition $p_0\ge 5r$ on the exponent of the degree function $p_0(\cdot)$.
\end{definition}

\begin{definition}[The constants of the step distinguished by letter]\label{4.7:not-step-constants}
Throughout this section the following symbols are fixed. Each denotes exactly one object. Where a
published paper uses the same glyph for a different object, the paper's glyph is kept only inside
quotations of that paper.

\begin{enumerate}
\item \emph{Enlargement and H\"older parameters.}
\begin{itemize}
\item $\beta_s$ denotes the enlargement parameter $\beta$ of \cite[\S3]{B-CMP109} and
\cite[(2.34)--(2.39)]{B-CMP119}; \cite{B-CMP119} observes that one can take $\beta=1/4$.
\item $\beta_0\in\,]0,\tfrac12]$ denotes the parameter of \cite[(2.6)]{B-CMP119}.
\item $\beta_0^H$ denotes the H\"older ceiling written $\beta_0$ in
\cite[(3.31), (4.18)]{B-CMP109}, and $\beta_0^H<1$.
\item $\beta_H\in\,]0,\beta_0^H]$ is the H\"older exponent.
\end{itemize}
The exponents $\beta_0^H$ and $\beta_H$ are fixed in the order of choice before $L$
and before the sup-Lipschitz constant $B_3$ that appears in the third line of
\eqref{4.7:eq-step-constants-a} below. The large-field
part of the coupling increment uses no property of these two exponents other than
$0<\beta_H\le\beta_0^H<1$.

\item \emph{Decay exponents.} We set $\kappa_1:=\kappa_0-5$. Outside quotations, the $\kappa_1$ of
\cite{B-CMP116} is written $\varkappa_1$. The absolute constant $C_1$ of
\cite[(1.20)]{B-CMP116} is written $C_1^{116}$.

\item \emph{The letters $B$ and $A$.} The constant $B$ of the coupling flow
(Definition~\ref{1.2:def-flow-constants}) occurs only through $B_\sharp$. Otherwise $B$, $A$ and
$\mathbf{A}$ are the fields of \cite[(3.30)]{B-CMP109}. The symbol $B^{109}$ is the constant of
Definition~\ref{4.4:def-local-data-classes}, that is, the constant $B$ of
\cite[(1.12)]{B-CMP109}, which that paper requires to be sufficiently large and determines later;
it may be increased. The symbol $A_\square$ denotes the potential on the cube
$\square$ in the same display.

\item \emph{A threshold.} $\nu_{\mathrm{CL}}$ denotes the threshold $\nu$ that enters condition
$(\Gamma3)$ of this chapter. The subscript separates it from every other use of $\nu$.

\item \emph{The frame constant.} Let $c_1$, $C_H$, $C_V$ and $c_Q$ be the constants of the
axial-frame lemma. Set
\begin{equation*}
B_b:=6dLc_1C_HC_V,\qquad B_b^f:=3dLc_1C_H\cdot 448d^2c_Q .
\end{equation*}
The frame constant $\mathfrak{b}$ of the frame statement in the operator bounds at free current is
the constant of the axial-frame lemma read at $\beta=\beta_s$, namely $\mathfrak{b}$ of the first
line of \eqref{4.7:eq-step-constants-a} below. The different product $B_b(1+2\beta_s)$, which
belongs to the lemma on the axial gauge, is not used in this section. Wherever
\cite[(3.27), (3.32)]{B-CMP109} and two further places of that paper carry the
product $B_3\,O(1)$, that product is read as $\mathfrak{b}$; the identification of this product
with $\mathfrak{b}$ is supplied by the axial-frame lemma.

\item \emph{The constant of $Q_j$.} $O_Q$ denotes the constant in the norm bound of
Ba{\l}aban's level-$j$ averaging operator $Q_j$ of \cite{B-CMP109}. This
constant arises because a bond average at level $j$ sums at most $O(d)L^j$ bonds of the finest
lattice.

\item \emph{The constant of \cite[(3.37)]{B-CMP109}.} $O_{(3.37)}$ denotes the constant written
$O(1)$ in \cite[(3.37)]{B-CMP109}. It depends only on $(d,L)$. The quantities $n$, $i$, $i'$ are
those of Definition~\ref{4.7:not-step-gap}. The window cube $\square_0$, the enlarged cubes
$\tilde{\square}^{3}$ and $\tilde{\square}^{4}$, and the level cubes $\square_{i'}$,
$\square_{k+1}$ are written as in \cite[\S3]{B-CMP109}. The radii $\alpha_{0,i}$ are those of
Definition~\ref{4.7:not-running-couplings}, and $\alpha_0'''$ is the triple-primed radius of
Definition~\ref{4.4:def-radius-ceilings}.

Let $\tau\hat{B}$ be the window datum. It is bounded in two regions:
\begin{itemize}
\item on the blocks of $\tilde{\square}^{4}$, by the frame statement in the operator bounds at
free current;
\item off those blocks, by the second clause of the axial-frame lemma, which gives the first
assertion of the second line of \eqref{4.7:eq-step-constants-a} for the window data $B_f$; this
bound is applied layer by layer on $\square_{i'}\setminus\square_{k+1}$ through the layerwise
corollary on the potential map.
\end{itemize}
The data off $\tilde{\square}^{4}$ enter the restriction of $\mathbf{H}_{i'}$ to
$\tilde{\square}^{3}$ through the exponential decay of the potential map used at
\cite[(3.37)]{B-CMP109}.

Feed $\tau\hat{B}$ to the sup-Lipschitz bound of the potential map invoked at
\cite[(3.37)]{B-CMP109}, whose constant is $B_3$.
The third line of \eqref{4.7:eq-step-constants-a} is then the bound that $O_{(3.37)}$ sizes.
The third member of \cite[(3.37)]{B-CMP109}, the distance $|u_{i'}-1|$ of the gauge map from the
identity, is bounded by a constant depending only on $(d,L)$ times
$O_{(3.37)}B_3\mathfrak{b}M\alpha_{0,i}$, with no factor $L^{-n}$. The displays are
\begin{equation}\label{4.7:eq-step-constants-a}
\begin{gathered}
\mathfrak{b}:=B_b^f(1+2\beta_s),\\
\sup_{\mathfrak{B}(\square_0)}|B_f|\le\mathfrak{v}_f,\qquad
\mathfrak{v}_f:=448d^2c_Q\,M\alpha_0''',\\
\sup_{\tilde{\square}^{3}}\bigl\{|\mathbf{H}_{i'}(\square_0,\tau\hat{B})|,\,
|\nabla\mathbf{H}_{i'}(\square_0,\tau\hat{B})|\bigr\}
\le O_{(3.37)}B_3\,\mathfrak{b}\,M\alpha_{0,i}L^{-n}.
\end{gathered}
\end{equation}
Here $\mathfrak{B}(\square_0)$ is the set over which the second clause of the axial-frame lemma
takes the supremum of the window data $B_f$, and $\mathfrak{v}_f$ is $8dLc_Q$ times the size
$b_f$ of those data in that lemma.
The bound $\mathfrak{v}_f$ involves, among the radii, only $\alpha_0'''$, and $M$ only to the
first power. We have
$O_{(3.37)}\ge O_Q$. The constant $O_{(3.37)}$ carries both the factor $O(d)$ of $\|Q_{i'}\|$ and
the share of the data off $\tilde{\square}^{4}$. The factor $M$ stays outside $O_{(3.37)}$,
because both shares are linear in $M\alpha_0$.

\item \emph{A second $\delta_1$.} $\delta_1^{H1}$ denotes the constant $\delta_1$, depending only
on $(d,L)$, in the list of constants $R_1,\delta_0,\delta_1,B_0,\dots$ of the axial-frame lemma.
It is distinct from the quantity $\delta_1=O(M^{-1})$ of \cite[(4.22)]{B-CMP109}. That quantity
is written $\delta_1$ only where this display of Ba{\l}aban is quoted.

\item \emph{The gauge map of \cite[(3.38)]{B-CMP109}.} We set
$g_{\langle 3.38\rangle}:=u_j\bar{u}_j^{-1}$. Here $\bar{u}_j$ is the function of
\cite[(3.38)]{B-CMP109} described there as ``constant on blocks naturally connected with the
function $U_n(X,\cdot)$, and equal to $u_j$ at centers of the blocks''; inside this quotation
$n$ and $X$ are in Ba{\l}aban's notation.

\item \emph{Tower index and gap.} $\mathsf{n}$ denotes the tower index. It is Ba{\l}aban's $n$ in
\cite[(1.15)--(1.16)]{B-CMP109}, where $n=1,\dots,j$. At a label $(p,Z)\in\mathfrak{T}_j(X)$ of
the one-power hereditary class it ranges over $1\le\mathsf{n}\le p$. The plain letter $n$ denotes
the level gap of Definition~\ref{4.7:not-step-gap}; inside quotations Ba{\l}aban's $n$ is kept
as $n$.

\item \emph{Field and domain.} $\mathsf{Z}$ denotes a $\mathfrak{g}^{c}$-valued lattice field,
the argument of the field norms $N_m(\mathsf{Z})$. Where it is specialised, $\mathsf{Z}:=A/\xi_i$.
The letter $Z$ denotes a domain, as in the labels $(p,Z)\in\mathfrak{T}_j(X)$.

\item \emph{Scale, site and radius.} We set
\begin{equation*}
\mathsf{s}_n:=L^{-3(n+1)}\le L^{-3}.
\end{equation*}
This is the scale attached to the current argument $\mathbf{J}$ in the far representation.
We write $\mathcal{S}_i^{\flat}$ for the set of admissible sites at level $i$; a site is denoted
$\mathfrak{s}\in\mathcal{S}_i^{\flat}$, and the letter $\mathfrak{s}$ is used for sites only. The serif letter
\begin{equation*}
s_n:=\tfrac14\,\alpha_{3,i'}^{\sharp}/(b_iL^{-n})
\end{equation*}
is the radius of the circle in the complex variable $\sigma$; here $\alpha_{3,i'}^{\sharp}$ is a
radius at level $i'$ and $b_i$ a constant at level $i$ of the argument in which that circle is
used. It is a different symbol from $\mathsf{s}_n$ and is not that scale.

\item \emph{Abbreviations in one argument.} The symbols $\mathsf{t}$ and $\mathsf{a}$ abbreviate
$\alpha_{0,i'}$ and $\alpha_{0,i}L^{-2(n+1)}$ respectively; they are used in one argument only.
The other related letters are as follows:
\begin{itemize}
\item $t$ and $t_\square$ remain Ba{\l}aban's Duhamel interpolation parameters;
\item $a$ and $b$ are the indices of the lower and upper flow bounds at step $k$;
\item $a_0$ and $a_1$ are the radii in the definition of the one-power hereditary class.
\end{itemize}

\item \emph{The constants written $C_1$.} $C_1^{(172)}$ denotes the constant $C_1$ of
\cite[(172)]{B-CMP102b}, which enters the second ceiling on the radii in the lemma that sets
those ceilings. It is distinct from the
radius constant $C_1$ of \cite{B-CMP119}, from $C_1^{116}$, and from $C_1^{\mathrm{fl}}$.

\item \emph{The potential function and its derivative.} The nonlinear potential function of
\cite{B-CMP109} is written in bold, $\mathbf{H}_j$, throughout. The lightface symbol
$H_j:=D\mathbf{H}_j(0)$ denotes its derivative at $0$.

\item \emph{Floors.} The floors of the step are labelled $(\varphi1)$--$(\varphi6)$ where they are
set; a further floor, $(\varphi\mathrm{B})$, is imposed there as well; it plays no role when the
background configuration is $U=1$.
\end{enumerate}
\end{definition}

\begin{definition}[The shift-field bounds at free current]\label{4.7:not-shift-fields}
In this section the operator bounds at free current are used in the forms listed below, at free
background configuration and current $(\mathbf{U},\mathbf{J})$. Ba{\l}aban's step is taken at
level $i-1$, where $i$ is the level index of Notation~\ref{4.7:not-running-couplings}; indices
that Ba{\l}aban writes in \cite{B-CMP109} and \cite{B-CMP116} with the letter of his step are
written here with $i-1$ in its place. In the exponentials $e^{i\eta(\cdot)}$, $e^{i\xi(\cdot)}$
and in the logarithms below, the factor $i$ is the imaginary unit, not the level index.
Further, $\eta$ is the lattice spacing and $T_{\eta}$ the torus
at that spacing. We write $\mathfrak{U}^{T_{\eta}}$ for the set of
pairs $(\mathbf{U},\mathbf{J})$ that are defined and satisfy conditions (i)--(iii) of
\cite{B-CMP109} on the whole lattice $T_{\eta}$. The constants $\mathfrak{b}$ and
$\mathfrak{v}_f$ are those of \eqref{4.7:eq-step-constants-a}.

\emph{Near part.}
\begin{itemize}
\item[(N-a)] For $(\mathbf{U},\mathbf{J})\in\mathfrak{U}^{T_{\eta}}$, let
$\mathbf{U}^{\mathrm{gl}}$ be the glued background. Then $H_{i-1}(B')$, taken at the background
$(\mathbf{U}^{\mathrm{gl}},\mathbf{J})$,
is analytic in $(\mathbf{U},\mathbf{J},B')$ for $|B'|<c_9/8$, where $c_9$ is the analyticity
constant of the operator bounds at free current. In this domain
$N_{\beta}(H_{i-1}(B'))\le 16C_9^{K}|B'|$, with $N_{\beta}$ the weighted norm of those bounds. The
version at the interpolation parameter $s$ obeys the
same bound multiplied by $e^{c_{\mathbb{K}}\varkappa_1}$; this is \cite[(1.21)]{B-CMP116}. Here and
below $s$ denotes this interpolation parameter of \cite{B-CMP116}, not the depth below the cutoff.
\item[(N-b)] In the frame $\mathfrak{w}$ the background on $\square_0$ is $e^{i\eta\mathcal{A}}$.
Here $\mathcal{A}=L^{-1}\mathcal{H}^{\square_0}(B_f)$ is the potential given by the axial-frame
lemma, with $B_f$ the data of the window. It is the field-axial counterpart of Ba{\l}aban's
$L^{-1}\mathbf{H}_{k+1}(\square_0,(1/i)\log V)$ of \cite[(3.27)]{B-CMP109}, his $k+1$ being the
level written $i$ here. It obeys, on
$\tilde{\square}^{4}$, the first line of the display \eqref{4.7:eq-shift-fields-a} below.
\end{itemize}

Consequently the level-$i$ potential $\mathbf{H}_{i}=L\mathcal{A}$ is bounded by
$L\mathfrak{b}M\alpha_0$ on $\tilde{\square}^{4}$. This is one factor $L$ weaker than the bound
$B_3O(1)M\alpha_0$ displayed in \cite[(3.27)]{B-CMP109}.

Off $\tilde{\square}^{4}$, the operator bounds at free current bound the window's data in the
frame only by $\sup|B|\le\tfrac12 C_1^{(172)}\bar{\varepsilon}_1$ over the base $1$. Here
$C_1^{(172)}$ is the constant of \cite[(172)]{B-CMP102b}, $\bar{\varepsilon}_1$ is the smallness
parameter for the window's data fixed in the operator bounds at free current, and this bound is
not proportional to
$\alpha_0$. We therefore use, in addition, two further bounds of the axial-frame lemma. The first
is
\begin{equation}\label{4.7:eq-shift-fields-1}
\sup_{\mathfrak{B}(\square_0)}|B_f|\;\le\;\mathfrak{v}_f\;=\;8dLc_Qb_f\;=\;448\,d^2c_Q\,M\alpha_0''' .
\end{equation}
Here $\mathfrak{B}(\square_0)$ is the set, and $c_Q$, $b_f$ are the constants, of the axial-frame
lemma in which this bound is stated.
This involves the radius $\alpha_0'''$ only and is linear in $M$. The second is the axial-frame
lemma's bound on the potential $\mathcal{A}$, which holds on all of $\square_{i}$ and not only
on $\tilde{\square}^{4}$.

\emph{Shift field.} Put
\begin{equation*}
\mathbf{A}=\bigl(t\tilde{\zeta}_{\square}+t_{\square}\zeta_{\square}\bigr)\cdot
R(\mathfrak{w})H_{i-1}(B').
\end{equation*}
Here $\zeta_{\square}$, $\tilde{\zeta}_{\square}$ are the cut-off functions of the cube $\square$,
and $R(\mathfrak{w})$ is the passage to the frame $\mathfrak{w}$. Define the field-dependent radius
\begin{equation}\label{4.7:eq-shift-fields-2}
r_{\square}(B')=\frac{\alpha_2}{3C_N^{s}\sup|B'|}.
\end{equation}
Assume $3C_N^{s}\varepsilon_1'\le\alpha_2$, where $\varepsilon_1'$ is a smallness parameter of
the operator bounds at free current, fixed there. Then $\mathbf{A}$ obeys
\cite[(3.31)]{B-CMP109} with plain derivatives, in the form of the second line of
\eqref{4.7:eq-shift-fields-a}. This holds both at $s\equiv 1$ and at $s$.

The linear clause of the axial-frame lemma is the third line of \eqref{4.7:eq-shift-fields-a}. In
it $C_N$ is replaced by $C_N^{s}$ at $s$.

Finally, put
\begin{equation*}
\hat{\mathbf{A}}=\frac{1}{i\eta}\log\bigl(e^{i\eta\mathbf{A}}e^{i\eta\mathcal{A}}\bigr).
\end{equation*}
This is the field $A$ of \cite[(3.30)]{B-CMP109}. It obeys the fourth line of
\eqref{4.7:eq-shift-fields-a}, which replaces \cite[(3.32)]{B-CMP109}. The bounds are
\begin{equation}\label{4.7:eq-shift-fields-a}
\begin{aligned}
&|\mathcal{A}|,\ |\nabla^{\eta}\mathcal{A}|,\ \|\mathcal{A}\|_{1,\beta}\le\mathfrak{b}M\alpha_0
 &&\text{on }\tilde{\square}^{4},\\
&|\mathbf{A}|,\ |\nabla^{\eta}\mathbf{A}|,\ \|\mathbf{A}\|_{1,\beta}<\alpha_2
 &&\text{on }\square_0,\ |t|\le1,\ |t_{\square}|\le r_{\square}(B'),\\
&|\mathbf{A}|,\ |\nabla^{\eta}\mathbf{A}|,\ \|\mathbf{A}\|_{1,\beta}
 \le C_N\bigl(|t|+|t_{\square}|\bigr)\sup|B'|, &&\\
&|\hat{\mathbf{A}}|,\ |\nabla^{\eta}\hat{\mathbf{A}}|,\ \|\hat{\mathbf{A}}\|_{1,\beta}
 <2\bigl(\alpha_2+\mathfrak{b}M\alpha_0\bigr) &&\text{on }\tilde{\square}^{4},\\
&C_s=C_1^{116}\max\{1,C_1^{\mathrm{fl}}\}\,C_{\mathsf{l}}. &&
\end{aligned}
\end{equation}
In the last line $C_1^{\mathrm{fl}}$ is the constant of Notation~\ref{4.7:not-running-couplings},
and $C_{\mathsf{l}}$ is fixed in the paragraph on the constant $C_s$ below.

\emph{Far part.} Let $n$ denote the gap and $i$ the level index, as in
Notation~\ref{4.7:not-running-couplings}. Let $\mathsf{s}_n=L^{-3(n+1)}$, so that
$\mathsf{s}_n\le L^{-3}$, and let
$J=\mathsf{s}_n\mathbf{J}$. The kernel of \cite[(3.13)]{B-CMP109} is replaced by
\begin{equation}\label{4.7:eq-shift-fields-3}
\Phi^{\mathbb{K}}
=E^{(j)}\bigl(X,\;e^{i\xi A(t_{\square})}\mathbf{U},\;J+\mathfrak{J}(B(t),t_{\square})\bigr).
\end{equation}
Here $A$ and $P_1$ denote the field and the operator as they occur in
\cite[(3.13)--(3.14)]{B-CMP109}.
The list of \cite[(3.14)]{B-CMP109} is replaced by
\begin{equation}\label{4.7:eq-shift-fields-4}
|A|,\ |P_1A|,\ |\nabla_{\mathbf{U}}A|,\ |\mathfrak{J}|<\alpha_2\quad\text{on }X .
\end{equation}
The kernel slot $\mathfrak{J}$ takes the place of the member of \cite[(3.14)]{B-CMP109} involving
the Laplacian, whose analogue at free current is not proved and is not used here.

The condition that $(\mathbf{U},\mathbf{J})$ lie in the domain of \cite[(3.16)]{B-CMP109} is read
representative-wise, as in reading convention $(\rho4)$ of
Notation~\ref{4.5:not-reading-conventions}. The far clauses are the following.
\begin{itemize}
\item[(F1)] $\Phi^{\mathbb{K}}$ is analytic on the set of \cite[(3.16)]{B-CMP109} times
$\{|t_{\square}|\le r\}$, where $r$ is the radius of the disc of the parameter $t_{\square}$ in
the far clauses, and $|\Phi^{\mathbb{K}}|\le E_0e^{-\kappa d_j(X)}$, where $E_0$ is the bound
constant of those clauses.
\item[(F2)] The bound \cite[(3.9)]{B-CMP109} holds, and hence \cite[(3.17)]{B-CMP109} holds at
the rate $c_*\delta_0$, where $c_*$ is the rate constant of the far clauses.
\end{itemize}
Clauses (F3)--(F6) are the remaining far clauses of the operator bounds at free current; they are
stated there and are used here in the form stated there. They concern, in order: (F3) the
physical point; (F4) the statement \cite[(1.19)]{B-CMP109}; (F5) the versions at $s$ of the far
clauses; and (F6) the two further far terms.

\emph{Hypotheses.} The hypotheses are those of the operator bounds at free current. These include
its inputs taken from Ba{\l}aban's papers; one further such input is needed only for the variant
of (N-b) in the literal frame. They also include the following conditions.
\begin{itemize}
\item The ceilings on $\alpha_0$ and $\alpha_1$.
\item The ceiling on the ratio:
\begin{equation*}
\alpha_2\le\min\Bigl\{\frac{\alpha_1}{8},\;\frac{\alpha_0}{C_{IV}},\;\frac{1}{4C_F'},\;
\frac{1}{4C_{\mathrm{BCH}}}-\max\{b,b'\}\Bigr\}.
\end{equation*}
Here $b$, $b'$, $C_F'$ and $C_{\mathrm{BCH}}$ are constants of the operator bounds at free
current, fixed there.
\item The smallness of $\varepsilon_1$, a further smallness parameter of the operator bounds at
free current, against $\alpha_2$.
\item The floors, listed next.
\end{itemize}
In addition, the condition $C_{IV}\ge 4(4d+2)$ is imported; here $C_{IV}$ is the constant of the
lemma that places the pairs produced by the far clauses in Ba{\l}aban's class
$U^{c}_j(X,\alpha_0,\alpha_1)$; this lemma is used in the proof of (F1), described below.

In the floors, $c_{12}=\mathsf{o}$ is the number of Definition~\ref{4.4:def-local-data-classes}.
Three further constants occur in them: Ba{\l}aban's block-size constant of \cite{B-CMP99b},
which is chosen before $\kappa$; his constant $M_1$ of \cite{B-CMP119}; and the constant $R_1$ of
the operator bounds at free current. The floors are:
\begin{itemize}
\item $17M_1\le LM$ and $17M_1\le c_{12}LM$; these follow from $M>L^2M_1$, $L\ge17$ and
$c_{12}\ge1$;
\item $c_{12}L\ge2$;
\item the floor on the product of $R_1$ with that block-size constant required by the operator
bounds at free current together with \cite[(2.59)]{B-CMP96}, imposed after $L$;
\item $\delta_0M\ge32$;
\item $2R_1M_1\le M$;
\item $\delta_1^{H1}M\ge\varkappa_1$; this is a floor on $M$ imposed after the constants depending
on $(d,L)$ and after $\varkappa_1$;
\item $L\ge27$ odd, the floor on $M_1$ of the rim-regularity construction, and $M>L^2M_1$;
\item $L^2\ge 12/c_{12}$;
\item $L\ge 4d$, with $R_1\ge L$;
\item $M$ above the absolute floor that makes $|\nabla\zeta_{\square}|\le1$ and
$|\nabla\tilde{\zeta}_{\square}|\le1$.
\end{itemize}
These conditions are hypotheses of the operator bounds at free current. They are assumed here and
are not re-derived. The last floor enters the constant $C_N$ of the linear clause, and hence
$C_s$.

In the proof of (F1), the second member of \cite[(3.16)]{B-CMP109} is read through the ceilings
and the bound on the current slot. The pair produced lies in Ba{\l}aban's class
$U^{c}_j(X,\alpha_0,\alpha_1)$, by the lemma on that class whose constant $C_{IV}$ appears above,
applied to the maximal tower of
$X$; that lemma holds uniformly in $M$. The membership of this pair in the one-power hereditary
class $U^{\Sigma^1}_j$ is the content of the proposition establishing membership in that class.
Expressed in the radii of
\cite[(2.28)]{B-CMP119}, the ceilings become conditions on those radii together with the
definitions of the constants $c_T$ and $c_{\alpha}$.

\emph{The constant $C_s$.} Let $C_{\mathsf{l}}$ be the largest constant by which (N-a), (N-b),
(F1), (F2) and (F5) bound their lists in terms of $|B'|$; here (F5) is the clause on the versions
at $s$ named above. In (F1) this bound is linear. The constant
$C_{\mathsf{l}}$ contains $16C_9^{K}$, $C_N^{s}$, $e^{c_{\mathbb{K}}\varkappa_1}$ and their
products. By \cite[(1.20)]{B-CMP116},
\begin{equation}\label{4.7:eq-shift-fields-5}
|B'|\le C_1^{116}g_{i-1}|B| .
\end{equation}
Let $\delta_{i-1}$ and $C_1^{\mathrm{fl}}$ be as in Notation~\ref{4.7:not-running-couplings}.
The constant $C_s$ is then given by the last line of \eqref{4.7:eq-shift-fields-a}.
With this constant, the near norms are at most $C_s\delta_{i-1}$, and the far lists are at most
$C_s\delta_{i-1}L^{-n}$.
\end{definition}

\begin{proof}
The first, second and fourth lines of \eqref{4.7:eq-shift-fields-a} are clauses of the operator
bounds at free current: (N-b), the bound \cite[(3.31)]{B-CMP109} with plain derivatives, and the
replacement of \cite[(3.32)]{B-CMP109}. The third line is the linear clause of the axial-frame
lemma. Each holds under the hypotheses listed above, with $\mathcal{A}$ the potential of the
axial-frame lemma.

For the consequence drawn from (N-b), note that $\mathbf{H}_{i}=L\mathcal{A}$. Multiplying the
first line of \eqref{4.7:eq-shift-fields-a} by $L$ gives
\begin{equation*}
|\mathbf{H}_{i}|\le L\mathfrak{b}M\alpha_0\quad\text{on }\tilde{\square}^{4}.
\end{equation*}
Off $\tilde{\square}^{4}$ this argument gives nothing proportional to $\alpha_0$. That is why
\eqref{4.7:eq-shift-fields-1} and the bound on $\mathcal{A}$ over all of $\square_{i}$ are used
as well.

For the size lines, first consider the near part. By the definition of $C_{\mathsf{l}}$, which
contains the constants $16C_9^{K}$, $C_N^{s}$ and $e^{c_{\mathbb{K}}\varkappa_1}$, every norm in
(N-a), (N-b) and in the linear clause that is bounded in terms of $|B'|$ is at most
$C_{\mathsf{l}}|B'|$. By
\eqref{4.7:eq-shift-fields-5}, these norms are therefore at most
$C_1^{116}C_{\mathsf{l}}\,g_{i-1}|B|$.

The same holds for the members of the lists of (F1), (F2) and (F5). Their bounds in terms of
$|B'|$ are, by definition, at most $C_{\mathsf{l}}|B'|$.

By the last line of \eqref{4.7:eq-shift-fields-a},
\begin{equation*}
C_s=C_1^{116}C_{\mathsf{l}}\max\{1,C_1^{\mathrm{fl}}\}\;\ge\;C_1^{116}C_{\mathsf{l}} .
\end{equation*}
By Notation~\ref{4.7:not-running-couplings}, the fluctuation field obeys
\begin{align*}
g_{i-1}|B|&\le\max\{1,C_1^{\mathrm{fl}}\}\,\delta_{i-1}&&\text{on the near part},\\
g_{i-1}|B|&\le\max\{1,C_1^{\mathrm{fl}}\}\,\delta_{i-1}L^{-n}&&\text{on the far part}.
\end{align*}
Inserting these two bounds into the estimate $C_1^{116}C_{\mathsf{l}}\,g_{i-1}|B|$ obtained above
and using the last line of \eqref{4.7:eq-shift-fields-a}, we find that the near norms are at most
$C_s\delta_{i-1}$ and the far lists are at most $C_s\delta_{i-1}L^{-n}$.
\end{proof}

\begin{definition}[True pairs lie in the regularity class]\label{4.7:not-true-pairs}
Let $i$ be the level fixing the spacing $\eta=L^{-(i-1)}$, and for every level $m$ put
$\ell_m=L^{m-(i-1)}$. For a $\mathfrak{g}^{c}$-valued field $\mathsf{Z}$ set
\begin{equation}\label{4.7:eq-true-pairs-1}
N_m(\mathsf{Z})=\max\Bigl\{\ell_m\sup|\mathsf{Z}|,\ \ell_m^{2}\sup|\nabla\mathsf{Z}|,\
\ell_m^{3}\sup\bigl(|\Delta\mathsf{Z}|,|\partial^{*}\partial\mathsf{Z}|\bigr)\Bigr\}.
\end{equation}
Let $c_T^{\circ}=c_T^{\circ}(d,L,M)$ be the constant of the lemma asserting that the exponential
of a small $\mathfrak{g}^{c}$-valued field, taken together with its current, lies in Ba{\l}aban's
class $U^{c}_j$. Put
\begin{equation}\label{4.7:eq-true-pairs-a}
c_T=\min\Bigl\{c_T^{\circ},\ \frac18,\ \frac{1}{C_{IV}}\Bigr\},
\end{equation}
where $C_{IV}$ is the constant of that name fixed in Definition~\ref{4.7:not-shift-fields}.
Then for every level $j$, every domain $X$ and every $\mathfrak{g}^{c}$-valued field $\mathsf{Z}$
with $N_j(\mathsf{Z})\le c_T r_j$, where $r_j$ and the radii $\alpha_{0,j},\alpha_{1,j}$ are those
of Notation~\ref{4.7:not-running-couplings}, the true pair
$(e^{\mathrm{i}\eta\mathsf{Z}},J(e^{\mathrm{i}\eta\mathsf{Z}}))$, that is, the configuration
$e^{\mathrm{i}\eta\mathsf{Z}}$ taken together with its own current
$J(e^{\mathrm{i}\eta\mathsf{Z}})$, satisfies:
\begin{enumerate}
\item[(a)] $(e^{\mathrm{i}\eta\mathsf{Z}},J(e^{\mathrm{i}\eta\mathsf{Z}}))
\in U^{c}_j(X,\alpha_{0,j},\alpha_{1,j})$;
\item[(b)] $(e^{\mathrm{i}\eta\mathsf{Z}},J(e^{\mathrm{i}\eta\mathsf{Z}}))\in
U^{\Sigma^1}_j(X,\alpha_{0,j},\alpha_{1,j})$, where $U^{\Sigma^1}_j$ is the one-power
hereditary class.
\end{enumerate}
These memberships are used only at true pairs.
\end{definition}

\begin{proof}
\emph{Clause (a).} By \eqref{4.7:eq-true-pairs-a}, $c_T\le c_T^{\circ}$, so the hypothesis
$N_j(\mathsf{Z})\le c_T r_j$ implies $N_j(\mathsf{Z})\le c_T^{\circ}r_j$. This is the hypothesis of
the lemma on exponentials of small fields named above, and its conclusion is (a).

In that lemma, the curvature condition of the class $U^{c}_j$ is supplied
by \cite[(1.17)]{B-CMP109}. There Ba{\l}aban derives from \cite[Proposition~9]{B-CMP102b}
the bound
\begin{equation}\label{4.7:eq-true-pairs-2}
|\partial U_n(\dot{M}(\mathbf{U}))-1|<B_3\,2\alpha_0'L^{-2n}(L^n\xi)^2=2B_3\alpha_0'\xi^2,
\end{equation}
which we cite as printed from \cite[(1.17)]{B-CMP109}; in it the primed constant $\alpha_0'$ is
left free. The bound uses no tower of
localisation domains, because \cite[Proposition~9]{B-CMP102b} involves none.

\emph{Clause (b).} Clause (b) is the true-pair clause of the proposition placing background pairs
in the one-power hereditary class, which asserts that the conclusion of (a) holds with
$U^{\Sigma^1}_j$ in place of $U^{c}_j$ under the same hypothesis and with the same constant. We
record how the curvature constant of that class is met from \cite[(1.17)]{B-CMP109}.
Let $(p,Z)\in\mathfrak{T}_j(X)$ be a label of the one-power hereditary class,
and put $\mu=L^{-(j-p)}$. In \cite{B-CMP109} one primed constant $\alpha_0'$ is common to
\cite[(1.11)]{B-CMP109}, \cite[(1.12)]{B-CMP109} and \cite[(1.14)]{B-CMP109}, and it is this
constant that appears in \eqref{4.7:eq-true-pairs-2}. Of these, \cite[(1.12)]{B-CMP109} converts
to the label $(p,Z)$, read at cube side $M/L$, with a single factor $\mu$: the primed constant
becomes
\begin{equation*}
\alpha_0''=\mu\alpha_0',
\end{equation*}
and the constant $B^{109}$ becomes $LB^{109}$. The top layer of the tower of the label $(p,Z)$
has cube side $M/L$, by the statement fixing the top layer of each label's tower. Reading
\eqref{4.7:eq-true-pairs-2} at every label $(p,Z)\in\mathfrak{T}_j(X)$ at cube side $M/L$
therefore gives the curvature condition at $(p,Z)$ with constant
$2B_3\alpha_0''=2B_3\mu\alpha_0'$, and, the primed constant $\alpha_0'$ being free,
\begin{equation}\label{4.7:eq-true-pairs-3}
2B_3\mu\alpha_0'\le\mu\alpha_0,
\end{equation}
which is the required constant $\mu\alpha_0$ at that label. The printed bound is thus used at the
labels $(p,Z)$ and at cube side $M/L$, beyond the setting in which it is printed.

Hence, by the proposition placing background pairs in the one-power hereditary class, the
conclusion of (a) holds with $U^{\Sigma^1}_j$ in place of $U^{c}_j$, under the same
hypothesis $N_j(\mathsf{Z})\le c_T r_j$ with the same constant \eqref{4.7:eq-true-pairs-a}. This
is (b).
\end{proof}

\begin{remark}
Let $(p,Z)$ and $\mu$ be as in the proof of clause (b). If the conversion is applied to the
curvature variables alone, with group-valued factor equal to $1$, each of the two powers carries a
factor $\mu$, and one obtains
\begin{equation*}
2B_3\mu^2\alpha_0'\le\mu^2\alpha_0 .
\end{equation*}
Since $0<\mu\le 1$ (as $L>1$ and $p\le j$), we have $\mu^2\alpha_0\le\mu\alpha_0$. Either way the
curvature condition holds at $(p,Z)$ with constant at most $\mu\alpha_0$.
\end{remark}

\begin{definition}[The working radii at a pair of levels]\label{4.7:not-working-radii}
Let $(i,i')$ be a pair of levels, with $i\ge1$. Let $\alpha_{0,i},\alpha_{1,i}$ be the radii at level $i$ of
Lemma~\ref{4.4:lem-twist-stability}, with their constants $C_0\le C_1$. The remaining
constants are the following:
\begin{itemize}
\item $B_3$ is the sup-norm Lipschitz constant among the constants of the step
(Definition~\ref{4.7:not-step-constants}); it enters Ba{\l}aban's condition relating $\alpha_3$ to
$\alpha_0$ in \cite{B-CMP109}.
\item $\beta_s$ is the value taken here by the parameter $\beta$ of that condition, fixed among the
constants of the step (Definition~\ref{4.7:not-step-constants}); it satisfies $\beta_sL^{-2}<1$.
\item $O_{(3.49)}$ and $O_{(4.17)}$ are the constants written $O(1)$ in \cite[(3.49)]{B-CMP109} and in
\cite[(4.17)]{B-CMP109}; the constant $O_{(4.17)}$ includes $O_Q$, the constant in the bound on the norm
of the block-averaging operator $Q_j$ of \cite{B-CMP109}.
\item $C_{IV}$ is the constant of the ratio of radii required in the operator bounds at free current.
\item $\mathfrak{b}$ is the frame constant entering the shift-field bounds
(Definition~\ref{4.7:not-shift-fields}).
\end{itemize}
The \emph{working radii} at the pair $(i,i')$, together with the constant $c_\alpha$ and the quantity
$b_i$ attached to them, are
\begin{equation}\label{4.7:eq-working-radii-a}
\begin{gathered}
\alpha^{\sharp}_{3,i'}:=\frac{\beta_sL^{-2}\alpha_{0,i'}}{2B_3},\qquad
\alpha_{2,i}:=c_\alpha\,\alpha_{0,i-1},\\
c_\alpha:=\min\Bigl\{\frac{\beta_s}{16B_3O_{(3.49)}L^{2}},\ \frac{1}{24},\ \frac{1}{3C_{IV}}\Bigr\},\\
b_i:=2O_{(4.17)}\bigl(\alpha_{2,i}+\mathfrak{b}M\alpha_{0,i}\bigr).
\end{gathered}
\end{equation}
The constants entering $b_i$ are those of the shift-field bounds \eqref{4.7:eq-shift-fields-a} and of
\cite[(4.17)]{B-CMP109}. The quantity $b_i$ is taken to satisfy in addition
$b_i\ge O_{(3.49)}\alpha_{2,i}$; since $\mathfrak{b}M\alpha_{0,i}\ge0$, this holds as soon as
$2O_{(4.17)}\ge O_{(3.49)}$.

In the simple case the radius $\alpha_2$ at
level $i'$ is $c_T r_{i'}$, with $c_T$ as in \eqref{4.7:eq-true-pairs-a} and $r_{i'}$ the value at
level $i'$ of the sequence $r$ fixed with the running couplings.

In \cite{B-CMP109} the radius $\alpha_2$ is only required to be sufficiently small and is otherwise
free; it is fixed here by \eqref{4.7:eq-working-radii-a}. The working radii have the following
properties.
\begin{enumerate}
\item[(a)] The radius $\alpha^{\sharp}_{3,i'}$ is the largest $\alpha_3$ satisfying the assumption
$2B_3\alpha_3\le\beta L^{-2}\alpha_0$ of \cite{B-CMP109} at level $i'$ with $\beta=\beta_s$.
\item[(b)] The condition $B_3\alpha_3<\tfrac12\alpha_1$ of \cite{B-CMP109} holds at level $i'$ for
$\alpha_3=\alpha^{\sharp}_{3,i'}$.
\item[(c)] For every $i\ge1$ one has
\begin{equation*}
\alpha_{2,i}\le\min\Bigl\{\frac{\alpha_{1,i}}{8},\ \frac{\alpha_{0,i}}{C_{IV}}\Bigr\}.
\end{equation*}
\end{enumerate}
\end{definition}

\begin{proof}
(a) Multiplying the definition of $\alpha^{\sharp}_{3,i'}$ by $2B_3$ gives
$2B_3\alpha^{\sharp}_{3,i'}=\beta_sL^{-2}\alpha_{0,i'}$. Thus the assumption holds with equality at
$\alpha_3=\alpha^{\sharp}_{3,i'}$. It fails for every $\alpha_3>\alpha^{\sharp}_{3,i'}$.

(b) By (a), and then by $\beta_sL^{-2}<1$ and $\alpha_{0,i'}\le\alpha_{1,i'}$ (the latter since
$C_0\le C_1$, the two radii being those of Lemma~\ref{4.4:lem-twist-stability}),
\begin{equation*}
B_3\alpha^{\sharp}_{3,i'}=\tfrac12\beta_sL^{-2}\alpha_{0,i'}<\tfrac12\alpha_{0,i'}
\le\tfrac12\alpha_{1,i'}.
\end{equation*}

(c) By $\alpha_{2,i}=c_\alpha\alpha_{0,i-1}$ and the entries $1/24$ and $1/(3C_{IV})$ of the minimum
defining $c_\alpha$,
\begin{equation*}
\alpha_{2,i}\le\min\Bigl\{\frac{\alpha_{0,i-1}}{24},\ \frac{\alpha_{0,i-1}}{3C_{IV}}\Bigr\}.
\end{equation*}
The one-step case of the comparability of radii across a level gap,
\eqref{4.7:eq-radii-comparison-a}, bounds $\alpha_{0,i-1}$ by $3\alpha_{0,i}$; moreover
$\alpha_{0,i}\le\alpha_{1,i}$, since $C_0\le C_1$, as in (b). Hence
\begin{equation*}
\frac{\alpha_{0,i-1}}{24}\le\frac{3\alpha_{0,i}}{24}=\frac{\alpha_{0,i}}{8}\le\frac{\alpha_{1,i}}{8},
\qquad
\frac{\alpha_{0,i-1}}{3C_{IV}}\le\frac{3\alpha_{0,i}}{3C_{IV}}=\frac{\alpha_{0,i}}{C_{IV}},
\end{equation*}
and together with the previous display this gives (c). The bound so obtained is the ratio of radii
required in the operator bounds at free current.
\end{proof}

\begin{lemma}[Comparability of radii across a level gap]\label{4.7:lem-radii-comparison}
Assume condition $(\Gamma4)$ among the smallness conditions
\eqref{4.7:eq-smallness-conditions-a}, and, from the inductive hypothesis at step $k$, the flow
bounds and the bound $g_j\le\gamma$ for $0\le j\le k$. Let
$i'\le i-1\le k$, and put
$n=i-1-i'\ge 0$.

Write $\alpha_{\cdot,j}$ for either of the radii $\alpha_{0,j}$, $\alpha_{1,j}$ of
Definition~\ref{4.7:not-running-couplings}. The same radius letter is used in the
numerator and the denominator of each ratio below. When $i=k+1$, the quantities $g_i$,
$x_i$ and $\alpha_{\cdot,i}$ are read as $g_k$, $x_k$ and $\alpha_{\cdot,k}$. Let $\beta_0$ be
the constant of the smallness conditions \eqref{4.7:eq-smallness-conditions-a}. Then
\begin{align}
\frac{\alpha_{\cdot,i-1}}{\alpha_{\cdot,i'}} &\le 2(1+n)^{1/2},
\label{4.7:eq-radii-comparison-1}\\
\frac{\alpha_{\cdot,i}}{\alpha_{\cdot,i-1}} &\le 2(1+\beta_0),
\label{4.7:eq-radii-comparison-2}\\
\frac{\alpha_{\cdot,i}}{\alpha_{\cdot,i'}} &\le \rho_n,
\qquad \rho_n:=4(1+\beta_0)(1+n)^{1/2}.
\label{4.7:eq-radii-comparison-a}
\end{align}
Moreover,
\begin{equation}\label{4.7:eq-radii-comparison-5}
g_{i'}\le(1+\beta_0)\,g_{i-1},\qquad g_{i-1}\le(1+\beta_0)\,g_i ,
\end{equation}
and, for a single step in the inverse direction,
\begin{equation}\label{4.7:eq-radii-comparison-4}
\frac{\alpha_{\cdot,i-1}}{\alpha_{\cdot,i}}\le 2(1+\beta_0)\le 3 .
\end{equation}
\end{lemma}

\begin{proof}
Recall that $x_j=g_j^{-2}$.

\emph{Comparison of inverse couplings.} Under $(\Gamma4)$, the flow bounds
of the inductive hypothesis at step $k$ give the following
comparison of inverse couplings whenever $i'\le i-1\le k$ and $n=i-1-i'$:
\begin{equation}\label{4.7:eq-radii-comparison-3}
x_{i-1}-c_5\le x_{i'}\le x_{i-1}+B_\sharp n .
\end{equation}
By hypothesis $g_j\le\gamma$ for every $j\le k$, so $x_j\ge\gamma^{-2}$ for $j\le k$; this
lower bound is where the factors $\gamma^2$ and $\log\gamma^{-2}$ below enter. In each chain
below, the last inequality is the smallness supplied by $(\Gamma4)$.

We use two elementary inequalities. The first is
\begin{equation*}
(1-t)^{-1/2}\le 1+t \qquad\text{on } [0,\tfrac12].
\end{equation*}
It holds because
\begin{equation*}
(1+t)^2(1-t)-1=t(1-t-t^2)\ge t\bigl(1-\tfrac12-\tfrac14\bigr)\ge 0 .
\end{equation*}
The second is
\begin{equation*}
\log(1-t)\ge -2t \qquad\text{on } [0,\tfrac12].
\end{equation*}
It holds because $\log(1-t)+2t$ vanishes at $t=0$, and its derivative $2-(1-t)^{-1}$ is
non-negative on $[0,\tfrac12]$. We also use $\log(1+t)\le t$ for $t\ge 0$.

\emph{Across the gap.} Dividing the left inequality of
\eqref{4.7:eq-radii-comparison-3} by $x_{i-1}\ge\gamma^{-2}$ gives
$x_{i'}/x_{i-1}\ge 1-c_5\gamma^2$. Hence, by the first elementary inequality,
\begin{equation*}
g_{i'}\le(1-c_5\gamma^2)^{-1/2}g_{i-1}\le(1+\beta_0)\,g_{i-1}.
\end{equation*}
Dividing the right inequality by $x_{i-1}$ gives
\begin{equation*}
\frac{g_{i-1}}{g_{i'}}=\Bigl(\frac{x_{i'}}{x_{i-1}}\Bigr)^{1/2}
\le(1+B_\sharp n\gamma^2)^{1/2}\le(1+n)^{1/2}.
\end{equation*}
Next,
\begin{equation*}
\log x_{i'}\ge\log x_{i-1}+\log(1-c_5\gamma^2)\ge\log x_{i-1}-2c_5\gamma^2,
\end{equation*}
by the second elementary inequality. Dividing by $\log x_{i-1}\ge\log\gamma^{-2}$ gives
\begin{equation*}
\Bigl(\frac{\log x_{i-1}}{\log x_{i'}}\Bigr)^{q}
\le\Bigl(1-\frac{2c_5\gamma^2}{\log\gamma^{-2}}\Bigr)^{-q}\le 2 .
\end{equation*}

\emph{One step, for $i\le k$.} In this case \eqref{4.7:eq-radii-comparison-3} applies
with $i+1$ in place of $i$ and $i'=i-1$, so that the gap is $1$. It gives
$x_i-c_5\le x_{i-1}\le x_i+B_\sharp$.

From the right-hand bound, divided by $x_i\ge\gamma^{-2}$,
\begin{equation*}
\frac{g_i}{g_{i-1}}\le(1+B_\sharp\gamma^2)^{1/2}\le 1+\beta_0 .
\end{equation*}
The first chain of the previous paragraph, with the same shift of indices, gives
$g_{i-1}\le(1+\beta_0)g_i$.

From $x_i\le x_{i-1}+c_5$ and $\log(1+t)\le t$ we get
$\log x_i\le\log x_{i-1}+c_5\gamma^2$. From $x_{i-1}\le x_i+B_\sharp$ we get
$\log x_{i-1}\le\log x_i+B_\sharp\gamma^2$. Dividing the first by
$\log x_{i-1}\ge\log\gamma^{-2}$ and the second by $\log x_i\ge\log\gamma^{-2}$ gives
\begin{equation*}
\Bigl(\frac{\log x_i}{\log x_{i-1}}\Bigr)^{q}
\le\Bigl(1+\frac{c_5\gamma^2}{\log\gamma^{-2}}\Bigr)^{q}\le 2,
\qquad
\Bigl(\frac{\log x_{i-1}}{\log x_i}\Bigr)^{q}
\le\Bigl(1+\frac{B_\sharp\gamma^2}{\log\gamma^{-2}}\Bigr)^{q}\le 2 .
\end{equation*}
When $i=k+1$, the quantities $g_i$, $x_i$ and $\alpha_{\cdot,i}$ are read as $g_k$, $x_k$
and $\alpha_{\cdot,k}$. All four one-step ratios are then equal to $1$, so the same bounds
hold. Together with the first chain across the gap, the bound
$g_{i-1}\le(1+\beta_0)g_i$ proves \eqref{4.7:eq-radii-comparison-5}.

\emph{The radii.} In a radius $\alpha_{\cdot,j}$ of
Definition~\ref{4.7:not-running-couplings}, the level enters only through the factor
$g_j(\log x_j)^q$. The ratio of one radius letter at two levels is therefore the ratio
of the $g$'s times the $q$-th power of the ratio of the logarithms.

With the bounds across the gap, this gives
\begin{equation*}
\frac{\alpha_{\cdot,i-1}}{\alpha_{\cdot,i'}}\le(1+n)^{1/2}\cdot 2,
\end{equation*}
which is \eqref{4.7:eq-radii-comparison-1}. With the one-step bounds, it gives
\begin{equation*}
\frac{\alpha_{\cdot,i}}{\alpha_{\cdot,i-1}}\le(1+\beta_0)\cdot 2,
\end{equation*}
which is \eqref{4.7:eq-radii-comparison-2}. Multiplying the two gives
\eqref{4.7:eq-radii-comparison-a}:
\begin{equation*}
\frac{\alpha_{\cdot,i}}{\alpha_{\cdot,i'}}
=\frac{\alpha_{\cdot,i}}{\alpha_{\cdot,i-1}}\cdot\frac{\alpha_{\cdot,i-1}}{\alpha_{\cdot,i'}}
\le 4(1+\beta_0)(1+n)^{1/2}=\rho_n .
\end{equation*}

Finally, consider a single step in the inverse direction. The bounds
$g_{i-1}\le(1+\beta_0)g_i$ and $(\log x_{i-1}/\log x_i)^q\le 2$ give
\begin{equation*}
\frac{\alpha_{\cdot,i-1}}{\alpha_{\cdot,i}}\le 2(1+\beta_0)\le 3 ;
\end{equation*}
the last inequality holds because $\beta_0\le\tfrac12$ by the smallness conditions
\eqref{4.7:eq-smallness-conditions-a}. This proves \eqref{4.7:eq-radii-comparison-4}.
This one-step inverse bound is the only inverse
comparison of radii used in this chapter. No bound on the inverse ratio
$\alpha_{\cdot,i'}/\alpha_{\cdot,i-1}$ over a gap $n\ge1$ is asserted.
\end{proof}

\begin{lemma}[Size of the total field across a level gap]\label{4.7:lem-total-field}
Let $(i,i')$ be a pair of levels at gap $n$, as in Lemma~\ref{4.7:lem-radii-comparison}. Let the working
radii $b_i$, $\alpha^{\sharp}_{3,i'}$ and the constant $c_\alpha$ be as in \eqref{4.7:eq-working-radii-a},
let $\mathfrak{b}$ be the frame constant of \eqref{4.7:eq-step-constants-a}, and let $\beta_{\mathrm{s}}$ be
the parameter on which $\mathfrak{b}$ depends through the factor $1+2\beta_{\mathrm{s}}$. Let $r_{i'}$ be the
radius at level $i'$ that enters $\Theta_n$ below through the quantity $c_Tr_{i'}$, and let $c_T>0$ be the
constant multiplying it there. Let $\rho_n$ be as in
\eqref{4.7:eq-radii-comparison-a}, that is, $\rho_n=4(1+\beta_0)(1+n)^{1/2}$, with $\beta_0$ the constant
of that display. Let $B=Q(\eta A)$ be the field of
\cite[(3.30)]{B-CMP109}, $A$ denoting there Ba{\l}aban's composite potential, recalled in the proof; that is,
$B$ is the total field about $\mathbf{U}=1$ in the
fundamental case of \cite[\S3]{B-CMP109}, in which the region there denoted $X$ lies in
$\tilde{\square}^2$. Assume that the fluctuation variable $\mathbf{A}$ of \cite[(3.30)]{B-CMP109} and the
background configuration $\mathbf{U}$, through its potential $\mathcal{A}$ in the axial frame, satisfy the
shift-field bounds at free current \eqref{4.7:eq-shift-fields-a}
(Definition~\ref{4.7:not-shift-fields}). Lengths on $\tilde{\square}^4$ are
measured in units of level $i'$, and $\beta_H$ is the H\"older exponent at which the bounds
\eqref{4.7:eq-shift-fields-a} are read, with
$0\le\beta_H\le\beta^H_0$, where $\beta^H_0<1$ denotes the H\"older ceiling written $\beta_0$ in
\cite[(3.31)]{B-CMP109}. Then
\begin{equation}\label{4.7:eq-total-field-a}
\begin{gathered}
|B|\le b_iL^{-n},\qquad |\partial B|\le b_iL^{-2n},\qquad
|\partial\partial B|\le b_iL^{-(2+\beta_H)n}\quad\text{on }\tilde{\square}^4,\\
\Theta_n:=b_i\cdot\max\Bigl\{\frac{4}{\alpha^{\sharp}_{3,i'}},\,\bigl(c_Tr_{i'}\bigr)^{-1}\Bigr\}
\le\tilde{K}\rho_n .
\end{gathered}
\end{equation}
Here $\partial B$ and $\partial\partial B$ denote the lattice derivatives $\partial_\nu B_\mu$ and
$\partial_\lambda\partial_\nu B_\mu$, bounded pointwise, and
\begin{equation*}
\tilde{K}=\tilde{K}(d,L,M,\beta_{\mathrm{s}},B_3,\mathfrak{b},c_T,c_\alpha)\ge 1 .
\end{equation*}
The constant $\tilde{K}$ is independent of the coupling and of the levels $i$, $i'$ (the gap $n$ entering
only through $\rho_n$), and it contains no supremum
over the data; the radius letters enter the bound on $\Theta_n$ only through ratios.
\end{lemma}

\begin{proof}
We first recall Ba{\l}aban's construction of the field $B$, then transcribe his bounds into the level-indexed
notation of this chapter, and finally bound $\Theta_n$.

\emph{The field $B$ in \cite{B-CMP109}.} We are in the fundamental case recalled in the statement.
After the gauge transformations
\cite[(3.26)--(3.27)]{B-CMP109}, which are explicitly given analytic functions of $\mathbf{U}$ depending only on
the restriction of $\mathbf{U}$ to the cube $\square_0$ \cite[p.~276]{B-CMP109}, Ba{\l}aban considers the
expression obtained from \cite[(3.28)]{B-CMP109} by replacing the term
$(t\tilde{\zeta}_\square+t_\square\zeta_\square)\mathbf{H}_k(B')$, in the notation
there, by a $\mathfrak{g}^{c}$-valued variable $\mathbf{A}$. He then sets
\begin{equation*}
\begin{aligned}
A&=\frac{1}{\mathrm{i}\eta}\log\Bigl(\exp\bigl(\mathrm{i}\eta\mathbf{A}\bigr)\,
\exp\Bigl(\mathrm{i}L^{-1}\eta\,\mathbf{H}_{k+1}\bigl(\square_0,\tfrac{1}{\mathrm{i}}\log V\bigr)\Bigr)\Bigr),\\
B&=Q(\eta A)\quad\text{on }\square_0 ,
\end{aligned}
\end{equation*}
which is \cite[(3.30)]{B-CMP109}, with $V$ the configuration denoted so in \cite[\S3]{B-CMP109}. The function
$B$ is defined on the set of bonds that determine
$U_j(\square_0)$; here $j$ is the level of the block variables $U_j$ of \cite{B-CMP109}, at which the
averaging $Q$ is taken. Thus $\mathbf{A}$ is the fluctuation variable, $A$ is the composite potential, and $B$
is the level-$j$ average $Q(\eta A)$ of $A$. This averaging is the only source of the factor $L^j\eta$ in the
bounds below.

The variable $\mathbf{A}$ is assumed regular and to satisfy
\begin{equation*}
|\mathbf{A}|,\ |\nabla^{\eta}\mathbf{A}|,\ \|\mathbf{A}\|_{1,\beta}<\alpha_2\quad\text{on }\square_0,
\qquad 0\le\beta\le\beta^H_0<1,
\end{equation*}
which is \cite[(3.31)]{B-CMP109}, with the H\"older ceiling written $\beta_0$ there denoted $\beta^H_0$ here.
The function $\mathbf{H}_{k+1}(\square_0,\tfrac{1}{\mathrm{i}}\log V)$ obeys the same
bounds with $\alpha_2$ replaced by $B_3O(1)M\alpha_0$ and with $\square_0$ replaced by $\tilde{\square}^4$; see
\cite[(3.27), p.~275]{B-CMP109} for these bounds and their range.

The assumption \cite[(3.31)]{B-CMP109} implies
\begin{equation*}
|A|,\ |\nabla^{\eta}A|,\ \|A\|_{1,\beta}<2\bigl(\alpha_2+B_3O(1)M\alpha_0\bigr)
\quad\text{on }\tilde{\square}^4 ,
\end{equation*}
which is \cite[(3.32)]{B-CMP109}. The expression is then expanded in $B$ up to the fifth order
\cite[(3.33)--(3.34)]{B-CMP109}. By the definition of $B$ as $Q(\eta A)$ and by \cite[(3.32)]{B-CMP109}, one has
$|B|<O(1)L^j\eta$ \cite[\S3]{B-CMP109}.

For the first derivatives, \cite[(3.32)]{B-CMP109} is combined with \cite[Proposition~5]{B-CMP98}, which bounds
functional derivatives of averaging operations. This gives
\begin{equation*}
|(\partial_\nu B_\mu)(x)|<O(1)\bigl(\alpha_2+B_3O(1)M\alpha_0\bigr)(L^j\eta)^2<\alpha_1(L^j\eta)^2 ,
\end{equation*}
which is \cite[(4.17)]{B-CMP109}. There is no bound on second derivatives of $A$. There are, however, the bounds
on H\"older norms of first derivatives contained in \cite[(3.32)]{B-CMP109}. Arguing as in the proof of
\cite[(4.17)]{B-CMP109}, these give
\begin{equation*}
|(\partial_\lambda\partial_\nu B_\mu)(x)|<\alpha_1(L^j\eta)^{2+\beta},\qquad 0\le\beta\le\beta^H_0<1,
\end{equation*}
which is \cite[(4.18)]{B-CMP109}, with the same renaming of the H\"older ceiling. Both \cite[(4.17)]{B-CMP109}
and \cite[(4.18)]{B-CMP109} hold also for the field $\delta B$ of \cite[p.~285]{B-CMP109}.

In \cite{B-CMP119} the same expansion is carried out with the variable $\mathbf{A}$ of
\cite[(3.30)]{B-CMP109} set equal to zero; there the field is $L^{n'}\eta A$, with $n'$ the step index of
\cite{B-CMP119}, so that $L^{n'}\eta$ takes the place of $L^j\eta$. Hence $B$ is the total field, taken
about $\mathbf{U}=1$.

\emph{The bounds in the notation of the pair of levels $(i,i')$.} We work on $\tilde{\square}^4$ in units of
level $i'$, the pair $(i,i')$ and its gap $n$ being those of Lemma~\ref{4.7:lem-radii-comparison} and
Definition~\ref{4.7:not-working-radii}. Ba{\l}aban's averaging $Q$ at level $j$ is taken over blocks of side
$L^j\eta$ in the lattice of spacing $\eta$; in the units of level $i'$ this block side is $L^{-n}$, that is,
$L^j\eta=L^{-n}$. With this reading we transcribe his bounds below, after fixing their coefficients.

The coefficients of the bounds \cite[(3.31)]{B-CMP109} and \cite[(3.32)]{B-CMP109} are taken in the forms given
in \eqref{4.7:eq-shift-fields-a} (the shift-field bounds at free current,
Definition~\ref{4.7:not-shift-fields}), as assumed in the statement: bound \cite[(3.31)]{B-CMP109} enters in its
linear form and bound \cite[(3.32)]{B-CMP109} in its free-current form, both given there. The bounds
\cite[(4.17)]{B-CMP109} and \cite[(4.18)]{B-CMP109} are derived from \cite[(3.32)]{B-CMP109}; we use them with
\cite[(3.32)]{B-CMP109} replaced by its free-current form, Ba{\l}aban's radii
$\alpha_0$ and $\alpha_2$ carrying the level indices fixed in \eqref{4.7:eq-working-radii-a}. Bound
\cite[(4.18)]{B-CMP109} is applied with $\beta=\beta_H$, which is admissible since $0\le\beta_H\le\beta^H_0$.

Consider first the background's share of $B$. This share is $\hat{B}=Q(\eta\mathcal{A})$, where
$\mathcal{A}$ is the background potential in the axial frame. By the operator bounds at free
current, it satisfies
\begin{equation*}
|Q(\eta\mathcal{A})|\le O_Q L^{-n}\cdot\mathfrak{b}M\alpha_{0,i},
\end{equation*}
with $O_Q$ the constant in the bound on the norm of the averaging operator $Q$ and $\mathfrak{b}$ as in
\eqref{4.7:eq-step-constants-a}
(Definition~\ref{4.7:not-step-constants}). This share is one of the summands that the working radius $b_i$,
defined in \eqref{4.7:eq-working-radii-a} (Definition~\ref{4.7:not-working-radii}), dominates in the bound on
$|B|$ below.

We now transcribe the three bounds on $B$. For the field itself, $B=Q(\eta A)$, and in the units of level $i'$
the averaging carries the factor $L^{-n}$, so that
\begin{equation*}
|B|\le O_QL^{-n}\sup_{\tilde{\square}^4}|A| ,
\end{equation*}
where $\sup_{\tilde{\square}^4}|A|$ is bounded by the free-current form of \cite[(3.32)]{B-CMP109} given in
\eqref{4.7:eq-shift-fields-a}. The definition of $b_i$ in \eqref{4.7:eq-working-radii-a} makes $b_i$ at least
$O_Q$ times the coefficient of that bound, the background's share $O_Q\mathfrak{b}M\alpha_{0,i}$
included; hence $|B|\le b_iL^{-n}$.
For the derivatives, \cite[(4.17)]{B-CMP109} is derived with \cite[(3.32)]{B-CMP109} replaced by its
free-current form in \eqref{4.7:eq-shift-fields-a}. Its middle member then carries, in place of
$\alpha_2+B_3O(1)M\alpha_0$, the coefficient $\alpha_{2,i}+\mathfrak{b}M\alpha_{0,i}$ of that free-current
form, the background's share entering through the bound on $Q(\eta\mathcal{A})$ above; this coefficient is
multiplied by $O_{(4.17)}$, the constant written $O(1)$ in front of it in \cite[(4.17)]{B-CMP109}. By its
definition in \eqref{4.7:eq-working-radii-a}, the working radius $b_i$ is at least
$O_{(4.17)}(\alpha_{2,i}+\mathfrak{b}M\alpha_{0,i})$. The bound \cite[(4.18)]{B-CMP109} is obtained from
\cite[(3.32)]{B-CMP109} by the same considerations as \cite[(4.17)]{B-CMP109}; carried out with the
free-current form, these considerations produce in front of $(L^j\eta)^{2+\beta}$ the coefficient
$\alpha_{2,i}+\mathfrak{b}M\alpha_{0,i}$ multiplied by an absolute constant, and the definition of $b_i$ in
\eqref{4.7:eq-working-radii-a} takes $b_i$ at least this product as well. Since
$L^j\eta=L^{-n}$, we have $(L^j\eta)^2=L^{-2n}$ and
$(L^j\eta)^{2+\beta_H}=L^{-(2+\beta_H)n}$. Hence bound \cite[(4.17)]{B-CMP109} gives
$|\partial B|\le b_iL^{-2n}$, and bound \cite[(4.18)]{B-CMP109} with $\beta=\beta_H$ gives
$|\partial\partial B|\le b_iL^{-(2+\beta_H)n}$. These are the three inequalities in the
first line of \eqref{4.7:eq-total-field-a}.

\emph{The ratio $\Theta_n$.} The working radius $b_i$ belongs to level $i$: by its definition in
\eqref{4.7:eq-working-radii-a} it is the sum $\alpha_{2,i}+\mathfrak{b}M\alpha_{0,i}$ multiplied by a factor
depending only on absolute constants and on the constants listed in $\tilde{K}$. The quantity
$\alpha^{\sharp}_{3,i'}$, defined in \eqref{4.7:eq-working-radii-a}, and the quantity $c_Tr_{i'}$, with
$r_{i'}$ and $c_T$ as in the statement, belong to level $i'$. Hence each of the two ratios
$b_i/\alpha^{\sharp}_{3,i'}$ and $b_i/(c_Tr_{i'})$ that occur in $\Theta_n$ is, summand by summand of $b_i$, a
sum of two terms, each a coefficient depending only on the constants listed in $\tilde{K}$ (for the second
ratio including the factor $c_T^{-1}$) times a quotient $\alpha_{a,i}/\alpha_{b,i'}$, where $a\in\{0,2\}$ and
$\alpha_{b,i'}$ stands for the level-$i'$ radius in the denominator. We write each such quotient as
\begin{equation*}
\frac{\alpha_{a,i}}{\alpha_{b,i'}}=\frac{\alpha_{a,i}}{\alpha_{a,i'}}\cdot\frac{\alpha_{a,i'}}{\alpha_{b,i'}} .
\end{equation*}
The second factor is a ratio of two radii at the single level $i'$; through the definitions in
\eqref{4.7:eq-step-constants-a} and \eqref{4.7:eq-working-radii-a} it depends only on the constants listed in
$\tilde{K}$. This is the sense in which the radius letters enter $\Theta_n$ only through ratios. The first
factor is, in the notation of Lemma~\ref{4.7:lem-radii-comparison}, a ratio
$\alpha_{\cdot,i}/\alpha_{\cdot,i'}$ of radii of one and the same kind at levels $i$ and $i'$.

By the comparability of radii across a level gap (Lemma~\ref{4.7:lem-radii-comparison}), each such first factor
satisfies
\begin{equation*}
\frac{\alpha_{\cdot,i}}{\alpha_{\cdot,i'}}\le\rho_n=4(1+\beta_0)(1+n)^{1/2} ,
\end{equation*}
as stated in \eqref{4.7:eq-radii-comparison-a}. The factors of $\Theta_n$ other than these ratios depend only
on $d$, $L$, $M$, $\beta_{\mathrm{s}}$, $B_3$, $\mathfrak{b}$, $c_T$ and $c_\alpha$, through the definitions in
\eqref{4.7:eq-step-constants-a} and \eqref{4.7:eq-working-radii-a}. We collect these factors, summed over the
finitely many terms above, into a constant
$\tilde{K}=\tilde{K}(d,L,M,\beta_{\mathrm{s}},B_3,\mathfrak{b},c_T,c_\alpha)$, enlarged if necessary so that
$\tilde{K}\ge1$. This gives $\Theta_n\le\tilde{K}\rho_n$, the second line of \eqref{4.7:eq-total-field-a}.

No coupling, no level index and no supremum over the data enters $\tilde{K}$. The radii, which carry the
dependence on the coupling and the level, enter only through the ratios bounded by $\rho_n$.
\end{proof}

\begin{definition}[The small factor from differentiation in the cube parameter]
\label{4.7:not-small-factor}
For a cube $\square$ let $t_{\square}$ denote the complex parameter attached to $\square$ in
the interpolation of the step. Differentiation $\partial/\partial t_{\square}$ at $t_{\square}=0$
replaces exactly one factor by the fluctuation part. The constant of
\cite[(3.54)]{B-CMP109}, in the words that follow that display, ``is much bigger than $E_0$
\dots\ This differentiation yields a factor $O(1)\alpha_2^{-1}\varepsilon_1$'', where $E_0$,
$\alpha_2$ and $\varepsilon_1$ are the constants and radii of that paper. In
\cite[(1.22)]{B-CMP116} the same differentiation is carried out by a Cauchy estimate on a
circle in the variable $t_{\square}$; the fractional factors that accompany it there play no
role on the layer considered here and are not reproduced. We fix the size of this one
factor as follows. For every level $i\ge1$ set
\begin{equation}\label{4.7:eq-small-factor-a}
\vartheta_i := C_{\vartheta}\bigl(\log g_{i-1}^{-2}\bigr)^{-(q-p_0)},
\end{equation}
where
\begin{equation*}
C_{\vartheta} := 2C_sA_1\max\Bigl\{(c_{\alpha}C_0)^{-1},\;
\bigl(c_T\min\{C_0,C_1\}\bigr)^{-1}\Bigr\},
\end{equation*}
with $c_T$ the constant of \eqref{4.7:eq-true-pairs-a}, $c_{\alpha}$ the constant of the
working radii in Notation~\ref{4.7:not-working-radii}, and $C_s$ the constant of the
shift-field bounds at free current (Notation~\ref{4.7:not-shift-fields}). We also set
\begin{equation*}
\vartheta(\gamma) := C_{\vartheta}\bigl(\log \gamma^{-2}\bigr)^{-(q-p_0)}.
\end{equation*}
\end{definition}

\begin{lemma}[Bounds on the small factor]
Let $\vartheta_i$, $C_{\vartheta}$ and $\vartheta(\gamma)$ be as in
Notation~\ref{4.7:not-small-factor}. In the inequalities below, $\delta_{i-1}$ is the
fluctuation size and $\alpha_{2,i}$, $r_{i-1}$ are the radii of
Notations~\ref{4.7:not-running-couplings} and~\ref{4.7:not-working-radii}. For every level
$i\ge1$,
\begin{equation}\label{4.7:eq-small-factor-1}
\frac{2C_s\delta_{i-1}}{\alpha_{2,i}} \le \vartheta_i,
\qquad
\frac{2C_s\delta_{i-1}}{c_T\, r_{i-1}} \le \vartheta_i.
\end{equation}
Assume moreover $q>p_0$. Then, for every level $i$ with $0<g_{i-1}\le\gamma<1$,
\begin{equation}\label{4.7:eq-small-factor-2}
\vartheta_i \le \vartheta(\gamma),
\qquad
\vartheta(\gamma)\to 0 \quad\text{as } \gamma\to 0 .
\end{equation}
The quantities $\vartheta_i$, $C_{\vartheta}$ and $\vartheta(\gamma)$ do not depend on the
gap $n$.
\end{lemma}

\begin{proof}
The constant $C_{\vartheta}$ is $2C_sA_1$ times the larger of the two reciprocals
$(c_{\alpha}C_0)^{-1}$ and $\bigl(c_T\min\{C_0,C_1\}\bigr)^{-1}$. The first entry serves the
denominator $\alpha_{2,i}$ and the second the denominator $c_T r_{i-1}$. Since a maximum
dominates each of its entries,
\begin{equation*}
C_{\vartheta}\ge \frac{2C_sA_1}{c_{\alpha}C_0},
\qquad
C_{\vartheta}\ge \frac{2C_sA_1}{c_T\min\{C_0,C_1\}} .
\end{equation*}
Multiplying the first of these by $(\log g_{i-1}^{-2})^{-(q-p_0)}$, we see that the first
inequality in \eqref{4.7:eq-small-factor-1} follows from the bound
\begin{equation}\label{4.7:eq-small-factor-3}
\frac{\delta_{i-1}}{\alpha_{2,i}}
\le \frac{A_1}{c_{\alpha}C_0}\bigl(\log g_{i-1}^{-2}\bigr)^{-(q-p_0)} ;
\end{equation}
and multiplying the bound
\begin{equation}\label{4.7:eq-small-factor-4}
\frac{\delta_{i-1}}{r_{i-1}}
\le \frac{A_1}{\min\{C_0,C_1\}}\bigl(\log g_{i-1}^{-2}\bigr)^{-(q-p_0)}
\end{equation}
by $2C_s/c_T$ and using the second of these lower bounds on $C_{\vartheta}$, we obtain the
second inequality in \eqref{4.7:eq-small-factor-1}.
The bounds \eqref{4.7:eq-small-factor-3} and \eqref{4.7:eq-small-factor-4} follow from the
sizes of the fluctuation size $\delta_{i-1}$ and of the radii $\alpha_{2,i}$, $r_{i-1}$ as these
are fixed in Notations~\ref{4.7:not-running-couplings} and~\ref{4.7:not-working-radii};
with them the two inequalities \eqref{4.7:eq-small-factor-1} hold.

Both $\vartheta_i$ and $\vartheta(\gamma)$ carry the same constant $C_{\vartheta}$ and the
same exponent $-(q-p_0)$: they are the values of the function
$x\mapsto C_{\vartheta}x^{-(q-p_0)}$ at $x=\log g_{i-1}^{-2}$ and at $x=\log\gamma^{-2}$.
By hypothesis $q>p_0$ and $0<g_{i-1}\le\gamma<1$. Since $q>p_0$, the function above is
decreasing on $(0,\infty)$;
since $0<g_{i-1}\le\gamma<1$, one has
\begin{equation*}
\log g_{i-1}^{-2}\ge\log\gamma^{-2}>0 ,
\end{equation*}
so both arguments lie in $(0,\infty)$. This gives the first relation in
\eqref{4.7:eq-small-factor-2}. As $\gamma\to0$ one has $\log\gamma^{-2}\to\infty$, and since
the exponent $-(q-p_0)$ is negative,
$\vartheta(\gamma)=C_{\vartheta}(\log\gamma^{-2})^{-(q-p_0)}\to0$.

None of $C_s$, $A_1$, $c_{\alpha}$, $c_T$, $C_0$, $C_1$, $q$, $p_0$, $g_{i-1}$ or $\gamma$
involves the gap $n$. Hence $\vartheta_i$, $C_{\vartheta}$ and $\vartheta(\gamma)$ are free
of $n$.
\end{proof}

\begin{definition}[The kernel families and their marking]\label{4.7:not-kernel-families}
Fix a level $i$. The following objects and conventions are in force for the rest of the section.

\emph{The kernel family.} The family of kernels at level $i-1$ is
\begin{equation}\label{4.7:eq-kernel-families-1}
\mathbf{E}_{i-1}=\sum_{j\le i-1}\mathbf{E}^{(j)} .
\end{equation}
Each $\mathbf{E}^{(j)}$ consists of kernels indexed by triples $(j,X,z)$, with $j$ the level,
$X$ a localisation domain and $z\in X$ a base point of the kernel, as in \cite[(2.27)]{B-CMP119}.
The kernels satisfy conditions (i)--(iv) of \cite[(2.27)]{B-CMP119}. For every $j\le i-1$ the
supremum $S(\mathbf{E}^{(j)})$ of the kernels of $\mathbf{E}^{(j)}$, normalised by their decay, is
bounded by one constant $E_0$. Their one-point sums over the whole lattice are covariant, as in
\cite[(2.29)]{B-CMP119}.

\emph{The large-field terms.} The sum
\begin{equation}\label{4.7:eq-kernel-families-2}
\mathbf{R}'_{i-1}=\sum_{j\le k_1}\mathbf{R}^{(j)}
\end{equation}
of Ba{\l}aban's large-field terms $\mathbf{R}^{(j)}$ of the levels $j\le k_1$, for a level $k_1$, is
merged into $\mathbf{E}_{i-1}$; in \cite{B-CMP119}, whose $E_k$ is the family written
$\mathbf{E}_{i-1}$ here, this reads ``we include it into $E_k$''.
The terms $\mathbf{R}^{(j)}$ obey the bound \cite[(2.31), p.~260]{B-CMP119}, with the exponent
$\kappa_0$ of that display:
\begin{equation}\label{4.7:eq-kernel-families-3}
\bigl|\mathbf{R}^{(j)}\bigl(X,(\mathbf{U},\mathbf{J})\bigr)\bigr|
\le g_j^{\kappa_0}\exp\bigl(-\kappa\, d_j(X)\bigr).
\end{equation}
Their summation is controlled as stated in the paragraph of \cite{B-CMP119} following
\cite[(2.31)]{B-CMP119}:
\begin{quote}
After the vacuum energy renormalization we obtain a sum of marginal terms, i.e., terms with bounds
$O(1)(L^j\eta)^4 g_j^{\kappa_0}\exp(-\kappa d_j(X))$. The sum over $X$ is controlled by the exponential
factor, and by the factor $(L^j\eta)^4$. The sum over $j$ is controlled by $g_j^{\kappa_0}$
\end{quote}
Here $\eta=L^{-(i-1)}$ is the spacing fixed in the units of this chapter (not the bare spacing
$L^{-K}$).

\emph{Analyticity domains.} The analyticity domains of condition (ii) of \cite[(2.27)]{B-CMP119} are
read under this chapter's convention on classes of analyticity domains: they are the domains
$U^{\Sigma^1}_j(X;\alpha_{0,j},\alpha_{1,j})$ of the one-power hereditary class, with the radii
$\alpha_{0,j},\alpha_{1,j}$ of Notation~\ref{4.7:not-running-couplings}. Accordingly, at every step
$k$ of the induction of this chapter the induction hypothesis $(\gamma_k)$ is read with analyticity
domains in this class.

\emph{The marking.} The marking is the one fixed in the theorem of this chapter on the marked and
unmarked kernel families. The superscripts $\mathfrak{A}$ and $\mathfrak{R}$ distinguish the regular
and the large-field-descended sub-families of the kernel family \eqref{4.7:eq-kernel-families-1}; for
every level $j$ we put
\begin{equation}\label{4.7:eq-kernel-families-4}
\mathfrak{E}_j:=\mathbf{E}^{(j),\mathfrak{A}},\qquad
\mathfrak{M}_j:=\mathbf{E}^{(j),\mathfrak{R}},\qquad
\mathbf{E}^{(j)}=\mathfrak{E}_j+\mathfrak{M}_j ,
\end{equation}
and call $\mathfrak{E}_j$ the \emph{unmarked} and $\mathfrak{M}_j$ the \emph{marked} kernel family at
level $j$. Kernels and kernel families are written in bold face, as in
\eqref{4.7:eq-kernel-families-1}; the light-face $E_k$ occurs only in the sentence of
\cite{B-CMP119} quoted above.
An activity of the cluster expansion is \emph{marked} if it is born from a term of
$\mathbf{R}'_{i-1}$, or if it descends from a kernel of a marked family $\mathfrak{M}_{j'}$ of an
earlier level $j'$. Let $\mathfrak{t}$ denote the sum over clusters of the cluster expansion performed
in the present renormalisation step, and let $\mathfrak{t}^{\mathfrak{R}}$ be the part of
$\mathfrak{t}$ from clusters with at least one marked activity. We set
\begin{equation}\label{4.7:eq-kernel-families-5}
\mathfrak{t}^{\mathfrak{A}}:=\mathfrak{t}-\mathfrak{t}^{\mathfrak{R}} .
\end{equation}
The coupling increment $\beta_j$ of Definition~\ref{4.2:def-increment} splits correspondingly as
\begin{equation}\label{4.7:eq-kernel-families-6}
\beta_j=\beta_j^{\mathfrak{A}}+\beta_j^{\mathfrak{R}} .
\end{equation}
By the rule of separate renormalisation, each of the two sub-series $\mathfrak{t}^{\mathfrak{A}}$,
$\mathfrak{t}^{\mathfrak{R}}$ is renormalised by its own coefficient of the form
\cite[(5.42)]{B-CMP109}, namely $\mathfrak{t}^{\mathfrak{A}}$ by $\beta_j^{\mathfrak{A}}$ and
$\mathfrak{t}^{\mathfrak{R}}$ by $\beta_j^{\mathfrak{R}}$; compare \cite[\S5]{B-CMP109} and
\cite{B-CMP119}.
\end{definition}

\begin{definition}[The regularity class, its cube gauge condition and restrictions]
\label{4.7:not-class-restrictions}
Throughout the rest of this chapter, for $j\ge1$, $X\in\mathbf{D}_{j-1}$ and constants
$a_0,a_1,\gamma_0>0$, we write $U^{\Sigma^1}_j(X;a_0,a_1,\gamma_0)$ for the one-power
hereditary class. This is the union of the $G^c$-orbits \cite[(1.10)]{B-CMP109} of the
representatives at $(j,X;a_0,a_1,\gamma_0)$, as in the definition of the one-power hereditary class;
$\gamma_0$ is omitted when $\gamma_0=a_0$. The finite label set $\mathfrak{T}_j(X)$, the cut
domains $Z^{\wr p}$ and the tower conditions $(\mathrm{iv})_{p,Z}[c]$ are also as fixed in
that definition.
Here $\mathbf{D}_j$ is the class of localisation domains at level $j$, $\xi_j=L^{-j}$, and
$\mathsf{o}$ and $B^{109}$ are the constants of \cite[(1.12)]{B-CMP109} fixed in
Definition~\ref{4.4:def-local-data-classes}. Since $L$ is odd, $\mathbf{D}_j\subset\mathbf{D}_{j-1}$
as families of sets, so the class at level $j$ is in particular defined for $X\in\mathbf{D}_j$;
this is how it enters $(\mathrm{T}\times2)$ and $(\mathrm{T}\times3)$ below.
The following reading and convention are in force, and the three restriction statements of (c)
hold.

\smallskip
\emph{(a) The cube gauge condition.} In every use of the class at $(j,X;a_0,a_1,\gamma_0)$,
condition \cite[(1.12)]{B-CMP109} is read for every cube $\square\subset X$ of side at most
$\mathsf{o}LM\cdot L^j$ steps of the lattice of spacing $\xi_j$, with one and the same bound for
all such cubes. For each
such cube there is a $G$-valued gauge transformation $u$ on $\square$ such that
$U^u=\exp(i\xi_jA)$ on $\square$, where $U$ is the $G$-valued factor of the representative,
with
\begin{equation}\label{4.7:eq-class-restrictions-1}
|A|,\ |\nabla^{\xi_j}A|<\mathsf{o}LMB^{109}a_0\quad\text{on }\square .
\end{equation}

\smallskip
\emph{(b) The class convention.} Every space $U^c_{\cdot}(Y,\dots)$ of
Ba{\l}aban's, $Y$ a localisation domain, that occurs at the following places in what follows is
read $U^{\Sigma^1}_{\cdot}(Y;\dots)$ with the same constants: in the hypothesis of the lemma in
which the constants $K_{\varepsilon}$ and $K_R$ are fixed, in the two members of the
induction hypothesis $(\gamma_k)$ in which
such a space enters, in the hypotheses $(H_{\cdot})$, in the spaces \cite[(3.16)]{B-CMP109} and
the hypothesis
space $\mathfrak{U}$ of \cite[Lemma~4]{B-CMP109} as they enter the operator bounds at free
current, in the bound for the large-field-descended sub-family, and in the statement of the
large-field part of the coupling increment. The induction hypothesis $(\gamma_k)$ is carried on
$U^{\Sigma^1}$, which is the inductive class of the construction. At level $k+1$ it is
re-established, for the regular sub-family, by the statements of the regular part of the step
that prove clause (i-a) of Theorem~0, and, for the
large-field-descended sub-family, by the bound just named. The operator bounds at free current
and the independence of the localised terms from the small-field region
(Corollary~\ref{4.4:cor-localised-independence}) are read on the same class.

\smallskip
\emph{(c) The restrictions.} Let $\beta$ be the constant in the radii
$(1+\beta)\alpha_0,(1+\beta)\alpha_1$ of \cite{B-CMP116}, let
$\alpha_{0,i},\alpha_{1,i}$ be the radii of \cite[(2.28)]{B-CMP119}, as in
Lemma~\ref{4.4:lem-twist-stability}, let $\beta_s$ be the enlargement constant of these running
radii fixed among the constants of the step in Notation~\ref{4.7:not-step-constants}, and for
integers $n\ge0$ put $\mathsf{s}_n=L^{-3(n+1)}$. For a set $\mathcal{V}$ of pairs
$(\mathbf{U},\mathbf{J})$ on the bonds of a domain $W$ and a domain $V\subset W$, we write
$\mathcal{V}\restriction_V$ for the set of restrictions of its members to the bonds of $V$; the
same sign applied to a single pair denotes its restriction. The
following statements carry the names displayed:
\begin{equation}\label{4.7:eq-class-restrictions-a}
\begin{aligned}
&(\mathrm{T}\times1)\quad U^{\Sigma^1}_i(Y;a_0,a_1,\gamma_0)\restriction_X
\subset U^{\Sigma^1}_i(X;a_0',a_1',\gamma_0'),\\
&(\mathrm{T}\times2)\quad U^{\Sigma^1}_{k+1}(Y;(1+\beta)\alpha_0,(1+\beta)\alpha_1,\alpha_0)
\restriction_X\subset U^{\Sigma^1}_j(X;\tfrac12\alpha_0,\tfrac12\alpha_1,\alpha_0),\\
&\qquad\quad\text{and, in running form,}\\
&\qquad\quad (\mathbf{U},\mathbf{J})\in U^{\Sigma^1}_i(Y;(1+\beta_s)\alpha_{0,i},
(1+\beta_s)\alpha_{1,i},\alpha_{0,i})\\
&\qquad\qquad\Longrightarrow\ (\mathbf{U},\mathsf{s}_n\mathbf{J})\restriction_X\in
U^{\Sigma^1}_{i'}(X;\tfrac12\alpha_{0,i'},\tfrac12\alpha_{1,i'},\alpha_{0,i'}),\\
&(\mathrm{T}\times3)\quad U^{\Sigma^1}_{k+1}(X;\alpha_0,\alpha_1)\restriction_Y
\subset U^{\Sigma^1}_{k+1}(Y;\alpha_0,\alpha_1)\\
&\qquad\qquad\subset U^{\Sigma^1}_{k+1}(Y;(1+\beta)\alpha_0,(1+\beta)\alpha_1,\alpha_0).
\end{aligned}
\end{equation}
The hypotheses of each statement are as follows.
\begin{itemize}
\item[$(\mathrm{T}\times1)$] Here $i\ge1$, $X,Y\in\mathbf{D}_{i-1}$, $X\subset Y$,
$a_0\le a_0'$, $a_1\le a_1'$ and $\gamma_0\le\gamma_0'$. The statement includes its window
member: take $i=k+1$, so that $Y\in\mathbf{D}_k$, and let $\square_0=\tilde{\square}^5$ be
the window cube of \cite[\S3]{B-CMP109}, a domain of $\mathbf{D}_k$ with $\square_0\subset Y$,
on which the hypothesis space $\mathfrak{U}$ is defined. Then every representative at
$(k+1,Y;(1+\beta)\alpha_0,(1+\beta)\alpha_1,\alpha_0)$ satisfies
$(\mathrm{iv})_{k+1,\square_0}[(1+\beta)\alpha_0]$ on
$\square_0^{\wr(k+1)}=\tilde{\square}^3$, for either of the two rules by which the tower of a
domain is constructed in the definition of the class. This is the window condition used in
\cite{B-CMP116}, read on the one-power hereditary class; it is condition
\cite[(3.40)]{B-CMP109} at the radius $(1+\beta)\alpha_0$, which is smaller by the margin
$\beta\alpha_0$ than the radius $(1+2\beta)\alpha_0$ at which the hypothesis space
$\mathfrak{U}$ imposes it.
\item[$(\mathrm{T}\times2)$] In the first line, $L\ge2(1+\beta)$, $j\le k$, $X\in\mathbf{D}_j$,
$Y\in\mathbf{D}_k$ and $X\subset Y$. The restriction of each representative is itself a
representative, so one representative serves both members of \cite[(3.16)]{B-CMP109}. This is
the statement of \cite{B-CMP116}, read on the one-power hereditary class. In the
running form, $i'\le i-1$, $n=i-1-i'$, $X\in\mathbf{D}_{i'}$, $Y\in\mathbf{D}_{i-1}$ and
$X\subset Y$, provided that
\begin{equation}\label{4.7:eq-class-restrictions-2}
(1+\beta_s)\,\frac{\alpha_{\cdot,i}}{\alpha_{\cdot,i'}}\,L^{-(n+1)}\le\frac12,
\qquad \mathsf{s}_n\alpha_{0,i}\le\alpha_{0,i'} ,
\end{equation}
where $\alpha_{\cdot,i}$ stands for either of $\alpha_{0,i},\alpha_{1,i}$, as in
Lemma~\ref{4.4:lem-radii-ratio}. Both conditions are among the standing conditions on the
running radii $\alpha_{0,i},\alpha_{1,i}$ fixed with the constants of the step in
Notation~\ref{4.7:not-step-constants}.
\item[$(\mathrm{T}\times3)$] Here $Y\in\mathbf{D}_k$, $X\in\mathbf{D}_{k+1}$ and $Y\subset X$;
no further condition is imposed. This is the inclusion stated in
\cite[after (2.13)]{B-CMP116}, read on the one-power hereditary class.
\end{itemize}
These are restrictions valid for every such nested pair of domains and every gap $n\ge0$,
wherever the smaller domain lies in the larger one and however thin the larger one is. They
involve no condition on $M$, no split into interior and rim cases, and no separate case at
$n=0$. For a domain $X$ write $\tilde{X}^{-2}$ for the set on which \cite[(1.16)]{B-CMP109}
imposes its bounds in Ba{\l}aban's own definition of the class, namely $X$ less two layers of
the partition cubes of which $X$ is a union (for $X\in\mathbf{D}_j$, its cubes of the partition
$\pi_j$), as opposed to the cut $Z^{\wr p}$, which removes two layers of cubes of the partition
$\pi_{p-1}$. In particular, no
separate treatment is needed when the larger domain $Y$ of a pair $X\subset Y$ has
$\tilde{Y}^{-2}=\emptyset$, nor when, for a pair $Y\subset X$, the set $\tilde{Y}^{-2}$ is not
contained in $\tilde{X}^{-2}$: in each of these configurations the smaller domain of the pair,
paired with the level of the class read on it, is a label in the label set $\mathfrak{T}_{\cdot}$
of the larger domain, and the conditions $(\mathrm{iv})_{p,Z}$ are imposed at that label with
the cut $Z^{\wr p}$.
\end{definition}

\begin{proof}
Each of the three statements is an instance of the restriction theorem for the one-power
hereditary class or of its corollary.
Statement $(\mathrm{T}\times1)$ is the restriction theorem in the case $i'=i$ of equal levels,
where the
scale ratio $\mu=L^{-(i-i')}$ of that theorem equals $1$, so that its conditions
$\mu a_0\le a_0'$, $\mu a_1\le a_1'$, $\gamma_0\le\gamma_0'$ become those listed above; its
window member is the window member of the corollary, which holds because the label
$(k+1,\square_0)$ belongs to $\mathfrak{T}_{k+1}(Y)$. The first line of $(\mathrm{T}\times2)$
is the far-source member of the corollary, and its running form is the running member of the
corollary. The two provisos \eqref{4.7:eq-class-restrictions-2} are the hypotheses of that
running member; as noted after \eqref{4.7:eq-class-restrictions-2}, both are among the standing
conditions on the running radii of Notation~\ref{4.7:not-step-constants}, so the
running form
holds, with no further condition, for every pair of levels $i'\le i-1$ at which the step is run.
Statement $(\mathrm{T}\times3)$ is the domain member of the corollary.
\end{proof}

\begin{lemma}[Linear splitting of the potentials by family]\label{4.7:lem-potential-splitting}
Consider the renormalisation step from level $i-1$ to level $i$
(Definition~\ref{4.7:not-step-gap}). Let its regular input be
\begin{equation}\label{4.7:eq-potential-splitting-1}
-A\Bigl(\frac{1}{g^2_{i-1}(\cdot)},\cdot\Bigr)
+\sum_{j\le i-1}\bigl[\mathfrak{E}_j\oplus\mathfrak{M}_j\bigr]+R'_{i-1},
\end{equation}
together with the further term $R''$ and the boundary terms, as in
\cite[(3.26)--(3.28)]{B-CMP119}. Here $g_{i-1}(\cdot)$ is the position-dependent coupling of
Definition~\ref{1.5:def-domains}, $A(1/g^2_{i-1}(\cdot),\cdot)$ is the Wilson action weighted by it
(Definition~\ref{1.1:def-wilson-action}),
$\mathfrak{E}_j$ and $\mathfrak{M}_j$ are the unmarked and the marked kernel families at level $j$
(Definition~\ref{4.7:not-kernel-families}), and $R'_{i-1}$ is the remainder of the input at
level $i-1$. Assume the following.
\begin{enumerate}
\item[(1)] For each $j\le i-1$ the families $\mathfrak{E}_j$ and $\mathfrak{M}_j$ consist of
localized kernels $(j,X,z)$ with the properties \cite[(2.27)(i)--(iii)]{B-CMP119}, whose
whole-lattice one-point sums are covariant in the sense of \cite[(2.29)]{B-CMP119}.
\item[(2)] For each $j\le i-1$ the two families have decay-normalised sups, written
$S(\mathfrak{E}_j)$ and $S(\mathfrak{M}_j)$, taken on
$U^{\Sigma^1}_j(X;\alpha_{0,j},\alpha_{1,j})$; with $E_0$ the fixed ceiling on the kernel sups
among the constants of the step,
they satisfy $S(\mathfrak{E}_j)\le E_0$ and $S_j:=S(\mathfrak{M}_j)\le E_0$.
\item[(3)] Each of these families is renormalised by its own coefficient of the form
\cite[(5.42)]{B-CMP109}, grouped with it as in \cite[(4.38)]{B-CMP109}.
\item[(4)] For every $j<i$, the members of the torus families indexed by domains $X$ that do not
wrap around the torus are the members of the corresponding family on $\mathbb{Z}^4$; this is the
conclusion of the volume lemma at index $j$; on $\mathbb{Z}^4$ the condition is empty.
\item[(5)] The theorem of this chapter on the operator bounds at free current, whose clauses are
lettered (N) and (F) where it is stated, holds in these two clauses, together with the
axial-frame lemma that accompanies it, in part (a) of its first clause, in its second and third
clauses, in part (ii) of its fourth clause, and in its fifth clause.
\item[(6)] The hereditary corollary holds layer by layer.
\item[(7)] The one-power hereditary class $U^{\Sigma^1}_j$ is the class of
Definition~\ref{4.7:not-class-restrictions}, with the bound defining the class through the constant
$B^{109}$ of Definition~\ref{4.4:def-local-data-classes} read in the form displayed as the first
item of \eqref{4.7:eq-class-restrictions-a}, and with the collars $Z^{\wr p}$, the set $Z$ less two layers
of cubes of the partition $\pi_{p-1}$ (Definition~\ref{4.7:not-step-gap}); the cut to these
collars is made at every pair $(p,Z)$ of the finite label set $\mathfrak{T}_j(X)$ of the class.
\item[(8)] The restriction theorem for the regularity class holds in its one-step clause and its
running clause, with the restrictions of \eqref{4.7:eq-class-restrictions-a}. The accompanying
proposition holds in its three clauses (near backgrounds, far backgrounds, true pairs), the first
taken with Ba{\l}aban's single factorisation, in which $U:=1$.
\item[(9)] The floors (\(\varphi1\))--(\(\varphi6\)) and the smallness conditions
(\(\Gamma4\))--(\(\Gamma6\)) of \eqref{4.7:eq-smallness-conditions-a} hold.
\end{enumerate}
Then the potentials $V'(Y)$, $Y\in\mathbf{D}_{i-1}$, of \cite[Lemma~1]{B-CMP116}, whether
with no factor singled out or with exactly one factor singled out as in
\cite[(3.44)--(3.47)]{B-CMP119}, split linearly into the
contributions of the individual families and of the remainder:
\begin{equation}\label{4.7:eq-potential-splitting-2}
V'(Y)=\sum_{j\le i-1}\bigl[V^{\mathfrak{E}_j}(Y)+V^{\mathfrak{M}_j}(Y)\bigr]
+V^{R'}(Y)+(\text{vertex-free terms}).
\end{equation}
For $j\le i-1$ let $\pi_j(Y)$ be Ba{\l}aban's bound for the full share of a level-$j$ family of
unit sup, namely its kernel terms together with the leftovers of its own counterterm, as displayed
in \eqref{4.7:eq-unit-share-a} (the share bound $\pi_j(Y)$ is a function of the domain $Y$, not
the cube partition $\pi_{p-1}$ of hypothesis (7)); this bound carries the factor
$e^{-(1-2\delta)\kappa d_{i-1}(Y)}$. Then
\begin{equation}\label{4.7:eq-potential-splitting-a}
\begin{gathered}
\bigl|V^{\mathfrak{M}_j}(Y)\bigr|\le S_j\,\pi_j(Y)\qquad(j\le i-1),\\
\sum_{j\le i-1}\pi_j(Y)\le K_{\varepsilon}\,\vartheta_i\,e^{-(1-2\delta)\kappa d_{i-1}(Y)},\\
\bigl|V^{R'}(Y)\bigr|\le K_R\,\vartheta_i\,g_{i-1}^{5}
\sum_{j\le k_1}g_j^{\kappa_0-5}\,e^{-(1-2\delta)\kappa d_{i-1}(Y)}.
\end{gathered}
\end{equation}
Here $\vartheta_i$ is the small factor of Definition~\ref{4.7:not-small-factor}; in the last bound
$k_1$ and $\kappa_0$ are the level and the exponent of the bound on the remainder $R'_{i-1}$
among the constants of the step. The constant
$K_{\varepsilon}$ is independent of the induction step $k$ at which the lemma is applied, of $i$,
of the position, of the site (the torus $T_{\eta}$ or $\mathbb{Z}^4$) and of $\gamma$;
the constants on which it depends are those listed with the bound on $\pi_j$ established below.
The constant $K_R$ is independent of the induction step $k$.
\end{lemma}

The proof of Lemma~\ref{4.7:lem-potential-splitting} is given in the statements that follow it in
this section.

\begin{lemma}[Linearity and degree one in the sup constant]\label{4.7:prf-linearity-degree}
Consider the step $i-1\to i$ in the setting and under the hypotheses of
Lemma~\ref{4.7:lem-potential-splitting} (linear splitting of the potentials by family). For a
localisation domain $Y$ write $V'(Y)$ for the localised potential that the step produces from the
regular input. For $j\le i-1$ write $V^{\mathfrak{M}_j}(Y)$ for the sum of those terms of $V'(Y)$
that are built from the members of $\mathfrak{M}_j$, together with the terms left over, and not
cancelled, from the counterterm of $\mathfrak{M}_j$.
\begin{itemize}
\item[(D1)] Every operation that \cite[\S\S3--5]{B-CMP109} and \cite[\S1]{B-CMP116} perform on the
regular input is linear in the list consisting of the kernel families and their counterterms.
Consequently $V'(Y)$ is the sum over families displayed in the statement of
Lemma~\ref{4.7:lem-potential-splitting}, and its summand for $\mathfrak{M}_j$ is
$V^{\mathfrak{M}_j}(Y)$.
\item[(D2)] Every bound proved in these papers for the terms generated by a family at level $j$ has
the form (sup constant of the family)$\times$(a factor independent of the family). Consequently
$V^{\mathfrak{M}_j}(Y)$ satisfies Ba{\l}aban's level-$j$ bound with the sup constant $E_0$ replaced
by $S_j$; this is the bound for $V^{\mathfrak{M}_j}(Y)$ in \eqref{4.7:eq-potential-splitting-a}.
\end{itemize}
\end{lemma}

\begin{proof}
\emph{Proof of \textup{(D1)}.} We list the operations in the order in which they act and say,
for each, why it is linear in the list of families and counterterms.

The first is the Duhamel formula in the interpolation parameter $t_\square$ attached to each cube
$\square$, applied cube by cube, together with the split into two cases,
\cite[(3.4)--(3.5)]{B-CMP109} and \cite{B-CMP119}; the Duhamel formula writes a function as its
value at one end of the interpolation plus the integral of its derivative, and is therefore linear
in the function to which it is applied, and the split into the two cases is likewise linear in that
function. The
second is the Cauchy representation \cite[(3.15)]{B-CMP109}, a contour integral of the kernel
itself and hence linear in it. The third is the rewriting
\cite[(3.18)--(3.28)]{B-CMP109}, whose gauge maps, in Ba{\l}aban's words, ``are explicitly given
analytic functions of $\mathbf{U}$, depending on $\mathbf{U}$ restricted to the cube $\square_0$''
\cite[p.~276]{B-CMP109}; these maps depend on the configuration only and not on the kernel family
to which they are applied. The fourth is the Taylor expansion \cite[(3.34)]{B-CMP109}, whose terms
are derivatives of the function expanded and hence linear in it. The fifth consists of the Ward
identities \cite[(4.3)--(4.31)]{B-CMP109}, which are identities holding for each gauge-invariant
member separately.

The sixth consists of the extension of the localised functions ``to all $X\in\mathbf{D}_j^0$, where
$\mathbf{D}_j^0$ is the class of localization domains constructed for the lattice
$\xi\mathbb{Z}^4$'' \cite[p.~290]{B-CMP109}, of the sum
\begin{equation*}
\Pi:=\sum_X \mathbf{E}^{(2)}(X)
\end{equation*}
of \cite[(4.37)]{B-CMP109}, and of the coefficient $\beta$ defined in \cite[(5.42)]{B-CMP109}, which
is linear in $\Pi$. For a family living on the torus, this extension together with the cancellation
described next is the content of the volume-cancellation identity at level $j$.

The seventh is the cancellation of the terms of a family against that family's own counterterm; it
is an identity, not an estimate. In Ba{\l}aban's words \cite[\S4]{B-CMP109}: ``the only terms in the
expansion of (4.38), which we do not control yet, more exactly for which the sum over $j$ has not a
uniform bound, are terms (4.42), (4.44) \dots\ We define the $\beta$-function $\beta_j(g_{j-1})$
equal to this coefficient. The corresponding terms from (4.34) are equal to (4.42), (4.44) also,
hence both groups of terms cancel''; the displays cited in this sentence are
\cite[(4.34), (4.38), (4.42), (4.44)]{B-CMP109}. The same cancellation is carried out once more later
in \cite{B-CMP109}. The eighth is the localisation in the decoupling parameters $s$
\cite[(1.9)--(1.10), (1.23)]{B-CMP116}, together with the summation over the terms with the same
localisation domain $Y$ \cite[\S1]{B-CMP116}.

Each of the eight operations is linear in the list of families and their counterterms. Hence
$V'(Y)$ is the sum over families displayed in the statement of
Lemma~\ref{4.7:lem-potential-splitting}, and its summand for
$\mathfrak{M}_j$ is $V^{\mathfrak{M}_j}(Y)$, the terms built from the members of $\mathfrak{M}_j$
together with the non-cancelled leftovers of its counterterm \cite[(4.38)]{B-CMP109}. That the
leftovers of the counterterm are processed by the same operations is stated in \cite[\S4]{B-CMP109}:
``we apply to it the whole procedure of these two sections. As a result we obtain all the well
controlled terms, and the terms in (4.34) with the function $E^{(2)}$ replaced by the corresponding
function calculated for the expression (4.38)''.

Finally, a factor $H'(Z_0)$ attached to a distinguished domain $Z_0$ appears in \cite{B-CMP119},
where Ba{\l}aban writes of it: ``Although different, it satisfies the same bound as the other
factors''. This factor is the derivative of the same linear expression in the direction of the
variable attached to $Z_0$; it is therefore again linear in the list of families and
counterterms. This proves (D1).

\emph{Proof of \textup{(D2)}.} We go through the bounds of \cite{B-CMP109} and \cite{B-CMP116}
that enter the estimate of a family's terms, and read off the power of the family's sup constant
in each.

\emph{The simple case.} The bound \cite[(3.17)]{B-CMP109}, on p.~273, reads
\begin{equation*}
\begin{split}
|(3.15)|&\le \frac{1}{r}E_0\exp\bigl(-\kappa d_j(X)\bigr)\\
&\le E_0\exp\bigl(-\kappa d_j(X)\bigr)\cdot 6\alpha_2^{-1}L^j\eta B_0B_3\\
&\qquad\cdot\exp\Bigl(-\tfrac12\delta_0\operatorname{dist}^{(\xi)}(X,\square)
-\tfrac12\delta_0M(L^j\eta)^{-1}\Bigr)\,
|\zeta_\square\mathbf{H}_k(B')|,
\end{split}
\end{equation*}
in which $E_0$ appears to the first power. In this display every letter other than $E_0$ is
Ba{\l}aban's notation at that place of \cite[\S3]{B-CMP109}: $r$ is the radius of the circle of the
Cauchy representation \cite[(3.15)]{B-CMP109}, $B'$ is the field variable there, and
$\zeta_\square$ is the function attached to the cube $\square$ by his partition of unity
$\{\zeta_\square\}$; it is unrelated to the complex source $\zeta$ listed in the glossary,
Notation~\ref{0.2:not-glossary}. The
corresponding bound here is among the operator bounds at free current, in which the far
representation $\Phi^{\mathbb{K}}$ is bounded by its supremum over a circle. The background pairs
$(\mathbf{U},\mathbf{J})$ parametrised by that circle lie in $U^{\Sigma^1}_j$; this membership is
proved in this section, in part~(b) of the lemma on the
simple share.
The bound is therefore again the family's sup constant times a factor independent of the family.

\emph{The fundamental case.} The bound \cite[(3.54)]{B-CMP109} is a Cauchy estimate over a circle
that lies inside the kernel's own space of analyticity. This is the content of
\cite[Lemma~4]{B-CMP109}: ``The functions in (3.53) are analytic on the above spaces'' (the
functions of \cite[(3.53)]{B-CMP109}). The images
of the circle lie in $U^{\Sigma^1}_j$, uniformly over the ranges of $(\tau,B',\sigma)$ printed in
\cite[\S3]{B-CMP109}; that they lie in $U^{\Sigma^1}_j$ is the content of the lemma of this section
on \cite[Lemma~4]{B-CMP109}, and the uniformity over
these ranges is the content of the lemma of this section on the fifth-order terms.
In \cite[(3.54)]{B-CMP109}, p.~280, the
integrand is evaluated at
$\exp i(\tau B+\sigma B)$, and the right-hand side is
\begin{equation*}
E_0\alpha_3^{-5}|B|^5,
\end{equation*}
again with $E_0$ to the first power; here $B$, $\tau$ and $\sigma$ are Ba{\l}aban's notation at
\cite[(3.53)--(3.54)]{B-CMP109}, $B$ being the field variable and $\tau$, $\sigma$ the parameters
of the Cauchy representation. The bound \cite[(4.5)]{B-CMP109} carries the factor
$E_0\exp(-\kappa d_j(X))$, once more to the first power.

\emph{The bound \cite[(4.22)]{B-CMP109}, on p.~286.} The display immediately before it reads
\begin{equation*}
\begin{split}
&\sum_{x,x_3}\Bigl(8B_3\frac{1}{\alpha_2}\Bigr)^4E_0
\exp\Bigl(-\kappa d_j(X)-\delta_0\operatorname{dist}^{(\xi)}(X,x)
-\delta_0\operatorname{dist}^{(\xi)}(X,x_3)\Bigr)\\
&\qquad\cdot|B|^2|\delta B(x)|\,|(\partial B)(\Gamma_{x,x_3})|\\
&\quad<\sum_{x,x_3}\Bigl(8B_3\frac{\alpha_1}{\alpha_2}\Bigr)^4E_0
\exp\Bigl(-\tfrac13\kappa d_j(X)-\delta_1|x-x_0|-\delta_1|x_3-x|\Bigr)
|x_3-x|\,(L^j\eta)^5 .
\end{split}
\end{equation*}
Here $\delta B(x)$, $\Gamma_{x,x_3}$, $(\partial B)(\Gamma_{x,x_3})$, $x_0$ and $\delta_1$ are
Ba{\l}aban's notation at that place of \cite[\S4]{B-CMP109}.
It carries $E_0$ to the first power on both sides, and it is the display from which the degree in
$E_0$ of the second sum in \cite[(4.21)]{B-CMP109} is read. Ba{\l}aban continues: ``Summing over
$x$, $x_3$ we get finally the following estimate''
\begin{equation*}
|(\text{the second sum in }(4.21))|
\le\Bigl(32B_3\frac{\alpha_1}{\alpha_2}c_0(\delta_1)c_1(\delta_1)\Bigr)^4
\exp\Bigl(-\tfrac13\kappa d_j(X)\Bigr)(L^j\eta)^5 ,
\end{equation*}
where $c_0(\delta_1)$ and $c_1(\delta_1)$ are the constants of \cite[(4.22)]{B-CMP109}.
No factor $E_0$ is printed in \cite[(4.22)]{B-CMP109}; the degree in the sup constant of the
second sum in \cite[(4.21)]{B-CMP109} is read from the display above it, whose summand carries
$E_0$ to the first power, so that this bound too is of degree one in the sup constant. On the same
page Ba{\l}aban states that ``we can
bound the
derivative $\mathbf{E}^{(n)}(X,x_1,\dots,x_n)$ by a constant times the exponential''; this is a
Cauchy estimate from the sup of the member, and is therefore of degree one in that sup.

\emph{The coupling-constant section.} The argument of \cite[\S5]{B-CMP109} uses symmetries and
analyticity in a strip, and the constant of that analyticity is the Cauchy constant of
$\mathbf{E}^{(2)}$, which is proportional to the sup. The terms of the primed sum $\Pi'$ of
\cite[\S5]{B-CMP109} that survive are bounded through \cite[(5.44)]{B-CMP109}, which is proportional
to the sup. The leftovers of the counterterm are proportional to the coefficient
$\beta^{\mathfrak{M}}_j$ of the family $\mathfrak{M}_j$, which, by the bound on the coefficients of
the families in the extraction of the coupling coefficients in this chapter, used here at the levels
$j\le i-1$ under the induction hypothesis $(H_{i-1})$, satisfies
\begin{equation*}
|\beta^{\mathfrak{M}}_j|\le C_{\Pi}S_j(c_Tr_j)^{-2},
\end{equation*}
where $C_{\Pi}$, $c_T$ and $r_j$ are the constants of that bound.

\emph{The cluster expansion of \cite[\S1]{B-CMP116}.} The bound \cite[(1.24)]{B-CMP116}, p.~7,
reads, with Ba{\l}aban's constants $C_1$ and $\kappa_1$ of that paper written $C_1^{116}$ and
$\varkappa_1$,
\begin{equation*}
\begin{split}
|(1.23)|&\le 8B_0C_1^{116}e^{16\varkappa_1}\alpha_2^{-1}g_k|B|\cdot
E_0\Bigl(\frac{\alpha_1}{\alpha_3}\Bigr)^5(L^j\eta)^5\\
&\qquad\cdot\exp\Bigl(-(\varkappa_1-1)M^{-4}\bigl|Y_0\setminus\tilde{\square}^{4}\bigr|\Bigr)
\exp\bigl(-\kappa d_j(X)\bigr),
\end{split}
\end{equation*}
where $Y_0$ is Ba{\l}aban's notation at that place of \cite[\S1]{B-CMP116}.
It contains one fluctuation factor $g_k|B|$ and $E_0$ to the first power; the fifth power
$(\alpha_1/\alpha_3)^5$ is a ratio of a field size to a radius, not a power of the sup constant.
The same holds for \cite[(1.29)]{B-CMP116}. The conclusion of \cite[Lemma~1]{B-CMP116}, the bound
\cite[(1.36)]{B-CMP116}, written with the same convention,
\begin{equation*}
|V'_k(Y,U,J,B)|\le E_0\varepsilon_1C_1^{116}M^q\exp\bigl(C_2\varkappa_1\bigr)
\exp\bigl(-(1-2\delta)\kappa d_k(Y)\bigr),
\end{equation*}
contains $E_0$ to the first power; here $\varepsilon_1$, $C_2$, $q$ and $\delta$ are the constants
of \cite[\S1]{B-CMP116}. The geometric sums \cite[(1.26)--(1.28), (1.31)--(1.32)]{B-CMP116}
carry no weight depending on the family. In these sums the number of domains is counted here by
the domain count of this book,
\begin{equation*}
N(Y)\le 16\bigl(2d_k(Y)+1\bigr),
\end{equation*}
valid for every $d_k(Y)\ge 0$, in which $N(Y)$ is the count that stands in these sums in the place
held by $M^{-4}|Y|$ in \cite[\S1]{B-CMP116}. With this count the sums again carry no weight
depending on the family.

\emph{Thresholds.} The smallness conditions on the radii and on the fluctuation variable,
\cite[\S3, (3.31)]{B-CMP109} and \cite[(1.20)--(1.22)]{B-CMP116}, do not involve the
sup constant; these are conditions $(\Gamma 5)$--$(\Gamma 6)$ of the list of threshold conditions
of this chapter. The conditions on the total potential, in \cite[\S1]{B-CMP116}, among them the one
that Ba{\l}aban states as ``We have used the assumption $\varepsilon_2\le 1$'', where
$\varepsilon_2$ is proportional to $E_0\varepsilon_1$, are met with $E_0$ replaced by $2E_0$,
since $S(\mathfrak{E}_j)\le E_0$ and $S_j\le E_0$ by the hypotheses of
Lemma~\ref{4.7:lem-potential-splitting}; this is condition $(\Gamma 3)$ of the same list.

Every bound entering the estimate of the terms of a level-$j$ family is therefore of degree one
in that family's sup constant, multiplied by a factor that does not depend on the family, and
every threshold involved holds for the family $\mathfrak{M}_j$. Hence the bound of
$V^{\mathfrak{M}_j}(Y)$ is Ba{\l}aban's level-$j$ bound with sup constant $S_j$, which is the bound
for $V^{\mathfrak{M}_j}(Y)$ in \eqref{4.7:eq-potential-splitting-a}. This proves (D2).
\end{proof}

\begin{lemma}[The share of a family of unit sup]\label{4.7:prf-unit-share}
Let $k$ be a level, and write $\eta=L^{-k}$ (this meaning of $\eta$ is local to this lemma and its
proof). Let $j$ be a level with $L^{j}\eta\le 1$. Consider a family of localised kernels of level $j$
whose sup constant equals one, and fix a localisation domain $Y\in\mathbf{D}_k$. The terms of the
family that are assigned (routed) to $Y$ are of two sorts, the \emph{fundamental} terms and the
\emph{simple} terms, described in parts (a) and (b) of the proof below. Let $\pi^{\mathrm{a}}_j(Y)$
be the size of the sum of the fundamental terms routed to $Y$; let $\pi^{\mathrm{b}}_j(Y)$ be the
size of the sum of the simple terms routed to $Y$, together with the two further far terms treated
in part (b); and let $\pi^{\mathrm{lo}}_j(Y)$ be the size of the sum of the leftovers at level $j$ of
the grouping of \cite[(4.38)]{B-CMP109}. Sizes are taken in the sense in which
\cite[(1.29)]{B-CMP116} bounds such sums, and all three are taken at level $j$, before summation
over $j$. Define
\begin{equation}\label{4.7:eq-unit-share-a}
  \pi_j(Y) := \pi^{\mathrm{a}}_j(Y) + \pi^{\mathrm{b}}_j(Y) + \pi^{\mathrm{lo}}_j(Y).
\end{equation}
Then, in the normalisation of \cite{B-CMP109},
\begin{equation}\label{4.7:eq-unit-share-1}
  \pi_j(Y) \le \varepsilon^{I}_1\, c_{\pi}\, (L^{j}\eta)^{\beta_H}\,
  e^{-(1-2\delta)\kappa d_k(Y)} .
\end{equation}
Here $\varepsilon^{I}_1$ is the constant $\varepsilon_1$ of \cite{B-CMP109}. The exponent
$\beta_H>0$ is a fixed positive value of the exponent $\beta$ of \cite[(4.18)]{B-CMP109}, taken
within the range of that display, $0\le\beta\le\beta^{H}_0$ (the upper endpoint being written
$\beta_0$ in \cite[(4.18)]{B-CMP109}). The numbers $\delta$ and $\kappa$ are those
of the decay factors of \cite{B-CMP116}. The constant $c_{\pi}$ is independent of the sup constant
of the family and of $j$.
\end{lemma}

\begin{proof}
Throughout, $L^{j}\eta=L^{j-k}\le 1$.

\emph{(a) The fundamental terms.} These are the terms whose polymer $X$ satisfies
$X\subset\tilde{\square}^{2}$. Such a term carries the data $(\square,\square_0,Y_0,X)$ of
\cite[\S1]{B-CMP116}, consisting of two cubes $\square$ and $\square_0$, a set $Y_0$ and the
polymer $X$, and is routed to $Y=Y_0\cup\square_0$.

The terms of orders zero to four of the expansion \cite[(3.34)]{B-CMP109} are treated by the
analysis of \cite[\S4]{B-CMP109}. The term of order zero is the vacuum energy. The term of order
one vanishes, by \cite[(4.14)]{B-CMP109}. The terms of orders two to four are reduced, by the Ward
identities, to two parts. The first part consists of the marginal monomials
\cite[(4.42), (4.44)]{B-CMP109}, which are cancelled identically. The second part consists of
irrelevant leftovers, of four kinds. For each kind, the power of $L^{j}\eta$ that
\cite{B-CMP109} itself gives is as follows.
\begin{itemize}
\item[($\alpha$)] Terms with one extra first difference carry the power five, by
  \cite[(4.22)]{B-CMP109}.
\item[($\beta$)] The Taylor remainder \cite[(4.30)]{B-CMP109}. Of it,
  \cite[\S4]{B-CMP109} says:
  ``By (4.18) the last term on the right-hand side can be estimated by
  $2|\Gamma_{x,x_3}|^{2}\alpha_1(L^{j}\eta)^{2+\beta}$ with a positive $\beta$. \dots{} it implies
  the bound (4.22) with the power $4+\beta$ instead of $5$ in the last factor'' (the symbols
  $\Gamma_{x,x_3}$ and $\alpha_1$ are those of \cite[\S4]{B-CMP109}). This kind
  therefore carries the power $4+\beta_H$. Here \cite{B-CMP109} itself takes $\beta$ positive,
  inside the range $0\le\beta\le\beta^{H}_0$ of \cite[(4.18)]{B-CMP109}.
\item[($\gamma$)] The third-order terms of \cite[(5.43)]{B-CMP109} containing $\Pi'$.
  Of them, \cite[\S5]{B-CMP109} says: ``we can write them in the form in which
  $\delta B$ is differentiated once, and $B$ twice. The bound (5.44), and the bounds (4.14), (4.15)
  for derivatives of $\delta B$, $B$ imply that all these third order terms in (5.43) are
  irrelevant''. In the second quoted sentence, the bounds for the derivatives of
  $\delta B$ and $B$ (the quantities so denoted in \cite[\S4]{B-CMP109}) used here are
  \cite[(4.17), (4.18)]{B-CMP109}. This kind carries the power
  $2+(2+\beta_H)$.
\item[($\delta$)] The extension differences \cite[(4.5), (4.36)]{B-CMP109}, which
  \cite[\S4]{B-CMP109} calls ``an irrelevant expression again''. This kind carries any power.
\end{itemize}
The class is described alike in \cite{B-CMP119}, as ``terms which can be bounded by
$O((L^{j}L^{-n})^{5-\beta})\exp(-\kappa d_j(X))$, where $\beta$ is a positive number'' (in
\cite{B-CMP119} the current spacing $\eta$ is written $L^{-n}$, $n$ there being the level written
$k$ here).

The remainder of order five is the term \cite[(1.23)]{B-CMP116}. It is bounded by the product of
\cite[(1.24)]{B-CMP116} and \cite[(1.25)]{B-CMP116}, and this bound is summed in three stages.
\begin{itemize}
\item For a cube $\square'\subset\tilde{\square}^{2}$, the sum over polymers $X$ containing
  $\square'$ is controlled by \cite[(1.26)]{B-CMP116}.
\item For the sum over cubes $\square'\subset\tilde{\square}^{2}$, \cite[\S1]{B-CMP116} writes:
  ``we use the factor $(L^{j}\eta)^{5}$ in (1.24). This yields $(6L)^{4}L^{j}\eta$, and the sum
  over $j$ is bounded by $2(6L)^{4}$''.
\item The sum over $Y_0$ and $\square_0$ is controlled by \cite[(1.27)--(1.28)]{B-CMP116}, the
  number of localisation domains being bounded by the lemma of this book on the number of
  localisation domains.
\end{itemize}
The three bounds are gathered in \cite[(1.29)]{B-CMP116}.

For the leftovers of kinds ($\alpha$)--($\delta$), \cite[\S1]{B-CMP116} states the same: ``the
terms in this expansion can be bounded similarly as in (1.24), using the inequalities (I.4.5),
(I.4.22), (I.4.36), (I.5.44) \dots{} The summation over all possible choices of $\square_0$, $Y_0$,
$j$, $X$ yields an expression satisfying (1.29)''.

The count of cubes $\square'\subset\tilde{\square}^{2}$ turns the factor $(L^{j}\eta)^{5}$ into
$(6L)^{4}L^{j}\eta$, as in the quoted computation. After this count, the kinds ($\beta$) and
($\gamma$), whose powers are $4+\beta_H$ and $2+(2+\beta_H)$, carry the factor
$(L^{j}\eta)^{\beta_H}$. Since $L^{j}\eta=L^{-(k-j)}$, writing $n=k-j$ gives their control of the
sum over $j$:
\begin{equation}\label{4.7:eq-unit-share-2}
  \sum_{j\le k}(L^{j}\eta)^{\beta_H}
  \le \sum_{n\ge 0}L^{-\beta_H n}
  = \bigl(1-L^{-\beta_H}\bigr)^{-1} < \infty .
\end{equation}

\emph{(b) The simple terms.} These are the terms whose polymer meets the complement of the cube,
$X\cap(\tilde{\square}^{2})^{c}\neq\emptyset$. As in \cite[\S1]{B-CMP116}, they are routed to
localisation domains $Y$ containing the set $X_0$ that \cite[\S1]{B-CMP116} attaches to the term.
Each such term obeys the bound \cite[(1.30)]{B-CMP116}, which follows from
\cite[(3.17)]{B-CMP109} and \cite[(1.21)]{B-CMP116}. This bound is summed in three stages.
\begin{itemize}
\item For a cube $\square'$, the sum over the polymers $X$ of simple terms containing $\square'$ is
  controlled by \cite[(1.32)]{B-CMP116} and \cite[(1.26)]{B-CMP116}.
\item The sum over such cubes $\square'$ is controlled by the first exponential in
  \cite[(1.30)]{B-CMP116}: ``The first term under the exponential gives also the factor
  $L^{j}\eta$, which controls the sum over $j$''.
\item The sum over cubes $\square$ is controlled as in \cite[\S1]{B-CMP116}, the number of
  localisation domains being bounded by the lemma of this book on the number of localisation
  domains.
\end{itemize}
The result is, in the words of \cite[\S1]{B-CMP116}, ``again a bound of the form (1.29) for the
considered sum, with the last exponential replaced by $\exp(-(1-2\delta)\kappa d_k(Y))$''.

With the simple terms go the two further far terms of \cite{B-CMP109}, which are also treated in
\cite[\S1]{B-CMP116}. They are bounded by the sixth of the estimates on far representations in the
operator bounds at free current.

\emph{Conclusion.} Consider a family of level $j$ with sup constant one. Collect the bounds of
part (a), as gathered in \cite[(1.29)]{B-CMP116}; the bounds of part (b), of the form
\cite[(1.29)]{B-CMP116} with the last exponential replaced by
$\exp(-(1-2\delta)\kappa d_k(Y))$; and the bounds on the leftovers of the grouping
\cite[(4.38)]{B-CMP109}. All are taken at level $j$, before summation over $j$, and expressed in
the normalisation of \cite{B-CMP109}. By \eqref{4.7:eq-unit-share-a}, they bound the three summands
of $\pi_j(Y)$. The factor $(L^{j}\eta)^{\beta_H}$ in \eqref{4.7:eq-unit-share-1} is the one
carried, after the count of cubes, by the kinds ($\beta$) and ($\gamma$), and the exponential in
\eqref{4.7:eq-unit-share-1} is the replaced exponential $e^{-(1-2\delta)\kappa d_k(Y)}$ of part (b);
collecting the remaining constants into $c_{\pi}$, which depends neither on the sup constant nor
on $j$, gives \eqref{4.7:eq-unit-share-1}.
\end{proof}

\begin{lemma}[Ba{\l}aban's Lemma~4 at every level gap]\label{4.7:prf-lemma-four-gaps}
Let $n\ge 0$ and let $(i,i')$ be a pair of levels at gap $n$. We use the following objects:
\begin{itemize}
\item the working radii at a pair of levels (Notation~\ref{4.7:not-working-radii}), namely
$\alpha_{0,i}$, $\alpha_{1,i}$, $\alpha_{0,i'}$, $\alpha_{1,i'}$ and $\alpha^{\sharp}_{3,i'}$;
\item the constants of the step distinguished by letter, $O_{(3.37)}$, $\mathfrak{b}$ and
$\mathfrak{v}_f$, of \eqref{4.7:eq-step-constants-a};
\item the shift-field bounds at free current (Notation~\ref{4.7:not-shift-fields});
\item the regularity class and its restrictions (Notation~\ref{4.7:not-class-restrictions});
\item the ratio $\rho_n=4(1+\beta_0)(1+n)^{1/2}$ of Lemma~\ref{4.7:lem-radii-comparison}.
\end{itemize}
Assume the conditions \eqref{4.7:eq-smallness-conditions-a}.

Let $\square_0$ be the window cube. Let $\mathbf{U}$ lie in the source space of
\cite[Lemma~4]{B-CMP109} written in index-$i$ letters,
\begin{equation*}
\mathfrak{U}_i:=U'^{c}_i\bigl(\square_0,(1+2\beta_s)\alpha_{0,i},(1+2\beta_s)\alpha_{1,i},
\alpha_{0,i}\bigr),
\end{equation*}
read in the one-power hereditary class $U^{\Sigma^1}$. Let $\mathcal{A}$ be the background potential in the
axial frame, and put $\hat{B}:=Q(\eta\mathcal{A})$; this is the field $B$ at vanishing fluctuation
$\mathbf{A}=0$. Let $h_n$, $\theta_{\times}$, $\theta_J$, $\bar{\theta}$ and $\bar{\theta}_J$ be defined as
in the display below. Assume in particular, in the form displayed below, the condition $(\varphi1)$, the
ceiling on $\alpha_{0,i}$ and the exponential inequality; they are among the conditions
\eqref{4.7:eq-smallness-conditions-a} ($(\varphi1)$ itself and two conditions of the kind $(\Gamma6)$).

Then the membership stated in the last line of the display holds, uniformly in $(\tau,B')$, for the
following data: every $\tau\in[0,1]$; every $\mathfrak{g}^c$-valued $B'$ with $|B'|<\alpha^{\sharp}_{3,i'}$
on the bonds of the tower of $\square_0$; and every localization domain $X\subset\tilde{\square}^2$.
\begin{equation}\label{4.7:eq-lemma-four-gaps-a}
\begin{aligned}
&h_n:=O_{(3.37)}B_3\mathfrak{b}M\alpha_{0,i}L^{-n},\\
&C_1>8(1+\beta_0)O_{(3.37)}B_3\mathfrak{b}M\,C_0\quad(\varphi1),\\
&\bigl(4LO_{(3.37)}B_3\mathfrak{b}M\bigr)^2\alpha_{0,i}\le\beta_s,\\
&\exp\Bigl\{28d\tfrac{M}{L}(1+2\beta_s)\bigl(\alpha_{1,i}+\mathsf{o}LMB^{109}\alpha_{0,i}\bigr)
+4dLc_Q\mathfrak{b}M\alpha_{0,i}\\
&\qquad+cO_{(3.37)}B_3\mathfrak{b}M\alpha_{0,i}\Bigr\}
(1+2\beta_s)\le1+3\beta_s,\\
&\theta_{\times}:=\sup_{n\ge0}\bigl[(1+8\beta_s)\rho_nL^{-2(n+1)}+2\beta_sL^{-2}\bigr],\qquad
\bar{\theta}:=\frac{4(1+\beta_0)(1+8\beta_s)+2\beta_s}{4(1+\beta_0)(1+9\beta_s)},\\
&\theta_J:=\theta_{\times}+\beta_sL^{-2},\qquad
\bar{\theta}_J:=\frac{4(1+\beta_0)(1+8\beta_s)+4\beta_s}{4(1+\beta_0)(1+9\beta_s)},\\
&\bigl(U_{i'}(\square_0,\exp i(\tau\hat{B}+B')),\,J_{i'}(\square_0,\exp i(\tau\hat{B}+B'))\bigr)\\
&\qquad\restriction X\in U^{\Sigma^1}_{i'}(X;\alpha_{0,i'},\alpha_{1,i'}).
\end{aligned}
\end{equation}
\end{lemma}

\begin{proof}
\emph{Notation.} Equation numbers of the form (3.xx) and (1.xx) quoted from Ba{\l}aban in this proof are
those of \cite{B-CMP109}. The correspondence with Ba{\l}aban's notation is as follows:
\begin{itemize}
\item his level $k+1$ is our $i$, and his level $j$ is our $i'$;
\item his factor $L^{j-1}\eta$ is our $L^{-(n+1)}$;
\item his level-free radii $\alpha_0$, $\alpha_1$ become level-indexed radii: $\alpha_{0,i}$,
$\alpha_{1,i}$ in the hypothesis space $\mathfrak{U}_i$ and in the size of the background, and
$\alpha_{0,i'}$, $\alpha_{1,i'}$ in the conclusions for the image at level $i'$, as displayed at each step
below; his $\alpha_3$ is $\alpha^{\sharp}_{3,i'}$; his $\beta$ is $\beta_s$.
\end{itemize}
The index $n$ is the gap between the two levels. The tower index of condition~(iv) of the class is
written $\mathsf{n}$, with $1\le\mathsf{n}\le p$ at a label $(p,Z)$ of the one-power hereditary class
(Notation~\ref{4.7:not-class-restrictions}).

\emph{The fluctuation.} Ba{\l}aban's smallness condition $g_k|B|<\varepsilon_1$ on the fluctuation
\cite{B-CMP109} is replaced throughout by the fluctuation size $\delta_{i-1}$. At the point where it
enters the estimate \cite[(3.43)]{B-CMP109}, it is replaced by $C_1^{\mathrm{fl}}\delta_{i-1}$. The shift
norms are those of Notation~\ref{4.7:not-shift-fields}.

\emph{The source space.} \cite[Lemma~4]{B-CMP109} writes the space with the two arguments
$(1+2\beta)\alpha_0$, $(1+2\beta)\alpha_1$ after $\square_0$,
$U'^{c}_{k+1}(\square_0,(1+2\beta)\alpha_0,(1+2\beta)\alpha_1)$. The space $\mathfrak{U}_i$ carries in
addition the third argument $\alpha_{0,i}$, the argument $\gamma_0=\alpha_0$ of the definition of these
spaces in \cite{B-CMP109}. The third argument only makes the space smaller, so
\cite[Lemma~4]{B-CMP109} applies to $\mathfrak{U}_i$ a fortiori.

The potentials have a common domain \cite{B-CMP116},
\begin{equation*}
U^{\Sigma^1}_i\bigl(Y;(1+\beta_s)\alpha_{0,i},(1+\beta_s)\alpha_{1,i},\alpha_{0,i}\bigr).
\end{equation*}
This domain meets the defining conditions of $\mathfrak{U}_i$ on $\square_0\subset Y$, as follows.
\begin{itemize}
\item Conditions (i)--(iii) restrict from $Y$ to $\square_0$, since, in Ba{\l}aban's words, ``the
conditions outside $Y$ are unessential'' \cite{B-CMP116}.
\item The one-power condition $(\mathrm{iv})^{\Sigma^1}$ of $\mathfrak{U}_i$ over the label set
$\mathfrak{T}_i(\square_0)$ follows from the inclusion $\mathfrak{T}_i(\square_0)\subset\mathfrak{T}_i(Y)$,
by restriction $(\mathrm{T}\times1)$ of \eqref{4.7:eq-class-restrictions-a}. This is the heredity
theorem for the one-power class at the value $\mu=1$ of its parameter $\mu$, taken in its window member.
\end{itemize}
The member of this condition at the label $(i,\square_0)$ is the one read through \cite[(3.40)]{B-CMP109}
in the proof of the near part of the proposition on the image of the step in the one-power class, used
below. It is condition~(iv) for the window tower of
$\square_0$ on $\square_0^{\wr}=\tilde{\square}^3$, at the radius
$(1+\beta_s)\alpha_{0,i}<(1+2\beta_s)\alpha_{0,i}$. It holds because $(i,\square_0)\in\mathfrak{T}_i(Y)$,
and it holds for either of the two rules by which the window tower is formed.

\emph{The fluctuation variable.} Let $\mathbf{A}$ be the fluctuation variable furnished by the operator
bounds at free current. We show that it obeys \cite[(3.31)]{B-CMP109} at the radius $\alpha_{2,i}$ (the
radius at which \cite[(3.31)]{B-CMP109} is read in index-$i$ letters), for the parameter $t$ and the
cube parameter $t_\square$ of \cite[(3.31)]{B-CMP109} (Notation~\ref{4.7:not-shift-fields}) in the ranges
\begin{equation*}
|t|\le1,\qquad |t_\square|\le\frac{\alpha_{2,i}}{2C_s\delta_{i-1}}.
\end{equation*}
By the linear clause of the axial-frame lemma in the cube parameter $t_\square$,
\begin{equation*}
|\mathbf{A}|,\ |\nabla^{\eta}\mathbf{A}|,\ \|\mathbf{A}\|_{1,\beta}
\le C_N^{s}\bigl(|t|+|t_\square|\bigr)\sup|B'|,
\qquad\text{with}\qquad C_N^{s}\sup|B'|\le C_s\delta_{i-1}.
\end{equation*}
Hence each of the three quantities is bounded by
\begin{equation*}
(1+|t_\square|)\,C_s\delta_{i-1}\le\tfrac12\vartheta_i\alpha_{2,i}+\tfrac12\alpha_{2,i}<\alpha_{2,i},
\end{equation*}
by condition $(\Gamma5)$ of \eqref{4.7:eq-smallness-conditions-a} together with the range of $t_\square$.

The statement of the operator bounds at free current grants only the field-dependent disc
$|t_\square|\le\alpha_{2,i}/(3C_N^{s}\sup|B'|)$, which may be smaller than the range used here. That
range is covered by the linear clause just quoted, with no constant changed.

\emph{The bound on the nonlinear potential: condition~(ii).} Ba{\l}aban's argument is this. On
$\tilde{\square}^3$ he bounds
\begin{equation*}
|\mathbf{H}_j(\square_0,\tau Q(\dots))|,\ |\nabla^{\xi}\mathbf{H}_j(\square_0,\tau Q(\dots))|
\quad\text{by}\quad B_3^2O(1)M\alpha_0L^{j-1}\eta
\end{equation*}
in \cite[(3.37), p.~277]{B-CMP109}, and then assumes $B_3^2O(1)M\alpha_0<\tfrac12\alpha_1$
\cite[p.~277]{B-CMP109}. The bound rests on the sup-Lipschitz property of
\cite[Proposition~9]{B-CMP102b},
\begin{equation*}
\|\mathbf{H}_j(D)\|_\infty,\ \|\nabla^{\xi}\mathbf{H}_j(D)\|_\infty\le B_3\|D\|_\infty,
\end{equation*}
localised by the decay \cite[(190)]{B-CMP102b}. It is applied at
$D=\tau\hat{B}=\tau Q_j(L^{-1}\eta\mathbf{H}_{k+1})$, where $\|Q_j\|\le O_QL^j$ on fine bonds, by the norm
bound on the averaging operator $Q_j$. Data outside $\tilde{\square}^4$ are cruder. They enter $H_j$
restricted to $\tilde{\square}^3$ only with the factor $e^{-\delta_0M(L^j\eta)^{-1}}$, by
\cite[(190)]{B-CMP102b}.

We run this argument on the statement of the operator bounds at free current, read locally. The field
$\hat{B}=Q(\eta\mathcal{A})$ is taken on the bonds of the tower of $\square_0$, and
$|\mathcal{A}|\le\mathfrak{b}M\alpha_{0,i}$ on $\tilde{\square}^4$; this is a bound on
$\tilde{\square}^4$ only. Off $\tilde{\square}^4$, the window data $B_f$ of the frame are sized by the
window-data bound of the axial-frame lemma,
\begin{equation*}
\sup_{\mathfrak{B}(\square_0)}|B_f|\le\mathfrak{v}_f=448d^2c_Q\,M(1+2\beta_s)\alpha_{0,i}
\end{equation*}
(see \eqref{4.7:eq-step-constants-a}; the radius $\alpha_0'''$ in which that display writes
$\mathfrak{v}_f=448d^2c_Q\,M\alpha_0'''$ is $\alpha_0''':=(1+2\beta_s)\alpha_{0,i}$, the radius
$(1+2\beta)\alpha_0$ read at $\beta=\beta_s$ and at level $i$). Here $\mathfrak{B}(\square_0)$ denotes the
window bond set of the frame. This far share is controlled by two further bounds:
\begin{itemize}
\item the bound of the axial-frame lemma on the potential, which holds on all of the cube $\square_i$
(Ba{\l}aban's $\square_{k+1}$);
\item the layerwise bound on the potential, on $\square_{i'}\setminus\square_i$.
\end{itemize}
It enters $\mathbf{H}_{i'}$ restricted to $\tilde{\square}^3$ through the decay
\cite[(190)]{B-CMP102b}. Hence $|\hat{B}|\le O_Q\mathfrak{b}M\alpha_{0,i}L^{-n}$ on the bonds within
$\tilde{\square}^4$, and
\begin{equation}\label{4.7:eq-lemma-four-gaps-1}
\sup_{\tilde{\square}^3}\bigl\{|\mathbf{H}_{i'}(\square_0,\tau\hat{B})|,\
|\nabla\mathbf{H}_{i'}(\square_0,\tau\hat{B})|\bigr\}\le h_n,
\end{equation}
with $h_n$ as in \eqref{4.7:eq-lemma-four-gaps-a}. The
constant $O_{(3.37)}$ of \eqref{4.7:eq-step-constants-a} is the $O(1)$ of \cite[(3.37)]{B-CMP109}. It
contains the $O(d)$ of $\|Q_{i'}\|$ and the share of the far data, and not only the share from
$\tilde{\square}^4$.

Ba{\l}aban's display carries $L^{j-1}\eta$, that is, $L^{-(n+1)}$. The additional $L^{-1}$ there comes
from his bound $B_3O(1)M\alpha_0$ for $\mathbf{H}_{k+1}=L\mathcal{A}$, which uses a bound on
$|\mathcal{A}|$ by a constant times $L^{-1}$. That bound is not contained in the statement of the operator
bounds at free current. We therefore run on the weaker factor $L^{-n}$.

The condition to be met is $h_n<\tfrac12\alpha_{1,i'}$. By Lemma~\ref{4.7:lem-radii-comparison} and the
working radii of Notation~\ref{4.7:not-working-radii}, $\alpha_{0,i}/\alpha_{1,i'}\le(C_0/C_1)\rho_n$.
Therefore
\begin{equation}\label{4.7:eq-lemma-four-gaps-2}
\begin{aligned}
h_n&\le O_{(3.37)}B_3\mathfrak{b}M\,\frac{C_0}{C_1}\,4(1+\beta_0)(1+n)^{1/2}L^{-n}\,\alpha_{1,i'}\\
&=\Bigl[\frac{8(1+\beta_0)O_{(3.37)}B_3\mathfrak{b}MC_0}{C_1}\Bigr]
\Bigl[\tfrac12(1+n)^{1/2}L^{-n}\Bigr]\alpha_{1,i'}
<\tfrac12\alpha_{1,i'}.
\end{aligned}
\end{equation}
The last inequality holds because the first bracket is less than $1$ by $(\varphi1)$ of
\eqref{4.7:eq-lemma-four-gaps-a}, and because
$\sup_n(1+n)^{1/2}L^{-n}=1$, attained at $n=0$, since $(1+n)^{1/2}\le2^n\le L^n$. Condition $(\varphi1)$
is a lower bound on $C_1$ at fixed $C_0$. The constants $C_0$ and $C_1$ are chosen after $M$, and $C_1$
may be taken as large as needed.

\emph{The ceiling on $\alpha_{0,i}$.} Ba{\l}aban works ``for $\alpha_1$ sufficiently small, e.g.
$O(1)M\alpha_1\leqq\beta$'' \cite[after (3.41)]{B-CMP109}, and states ``We assume that
$(4B_3^2O(1)M)^2\alpha_0\leqq\beta$'' \cite[p.~278]{B-CMP109}. In index-$i$ letters, and with the factor
$L^{-n}$ of \eqref{4.7:eq-lemma-four-gaps-1}, the corresponding ceiling is the inequality
$(4LO_{(3.37)}B_3\mathfrak{b}M)^2\alpha_{0,i}\le\beta_s$ of \eqref{4.7:eq-lemma-four-gaps-a}. It is a
condition of the kind $(\Gamma6)$: a quantity of the form $C(d,L,M)\gamma(\log\gamma^{-2})^q$ is required
to lie below a prescribed constant, and it is met under condition $(\Gamma4\mathrm{f})$ of
\eqref{4.7:eq-smallness-conditions-a}.

\emph{The curvature: first half of condition~(iii), with $B'\ne0$.} Put
\begin{equation*}
\mathsf{a}:=\alpha_{0,i}L^{-2(n+1)},\qquad \mathsf{t}:=\alpha_{0,i'}.
\end{equation*}
At $B'=0$ Ba{\l}aban obtains ``(3.44) with $\tau Q(\dots)$, and the factor $1+6\beta$ on the right-hand
side'' \cite[(3.44)--(3.47)]{B-CMP109}.

With $B'$, \cite[(3.50)]{B-CMP109} defines the correction $\mathbf{A}_2$, with
$|\mathbf{A}_2|,|\nabla^{\xi}\mathbf{A}_2|\leqq B_3|B'|<B_3\alpha_3$. On \cite[p.~280]{B-CMP109} the
left-hand side of the curvature bound \cite[(3.51)]{B-CMP109} is less than
\begin{equation*}
(1+6\beta)\alpha_0(L^{j-1}\eta)^2+2B_3\alpha_3\leqq\dots\leqq(1+7\beta)L^{-2}\alpha_0 .
\end{equation*}
It is followed by \cite[(3.52)]{B-CMP109}, which gives $<(1+8\beta)L^{-2}\alpha_0\xi^2$, ``by the
inequality (3.43) slightly modified for the present situation''.

In the level-indexed notation of this chapter, the linear part is less than
\begin{equation*}
(1+6\beta_s)\mathsf{a}+2B_3\alpha^{\sharp}_{3,i'}=(1+6\beta_s)\mathsf{a}+\beta_sL^{-2}\mathsf{t},
\end{equation*}
by the identity $B_3\alpha^{\sharp}_{3,i'}=\tfrac12\beta_sL^{-2}\alpha_{0,i'}$ of the working radii.
Three remarks on this bound are in order.
\begin{itemize}
\item The factor $(L^{j-1}\eta)^2$ is here the exact change of units
$\xi_i^2=L^{-2(n+1)}\xi_{i'}^2$ applied to the curvature of the window minimiser
$U_{k+1}(\square_0,\dot{M}\mathbf{U})$ of \cite[(3.40)]{B-CMP109}.
\item The passage from $1+2\beta$ to $1+6\beta$, through \cite[(3.43)]{B-CMP109} at $B'=0$ and
\cite[(3.45)]{B-CMP109}, does depend on the size of the background. Each of its steps is closed by the
ceiling on $\alpha_{0,i}$ in \eqref{4.7:eq-lemma-four-gaps-a}, with its factor $L^2$, because $B_3\ge1$.
\item At \cite[(3.41)]{B-CMP109} the factor $1+2\beta_s$ is multiplied by the exponential of the frame
exponent. That exponent is controlled by the frame-exponent estimates of the axial-frame lemma. The
resulting inequality is the exponential line of \eqref{4.7:eq-lemma-four-gaps-a}, which is a member of
$(\Gamma6)$. Ba{\l}aban closes \cite[(3.41)]{B-CMP109} by $O(1)M\alpha_1\leqq\beta$. That is a ceiling on
$\alpha_1$, which no ceiling on $\alpha_0$ provides, so this line is not a consequence of the ceiling on
$\alpha_{0,i}$.
\end{itemize}

The quadratic part is treated by the elementary inequality of \cite[(3.43)]{B-CMP109}, namely
$<8(B_3^2O(1)M\alpha_0L^{j-1}\eta)^2\xi^2\exp(\dots)$. We apply it to $\mathbf{H}_{i'}+\mathbf{A}_2$,
using $|\mathbf{H}_{i'}|\le h_n$ from \eqref{4.7:eq-lemma-four-gaps-1} and
\begin{equation*}
|\mathbf{A}_2|,\ |\nabla^{\xi}\mathbf{A}_2|<h':=B_3\alpha^{\sharp}_{3,i'}=\tfrac12\beta_sL^{-2}\mathsf{t}.
\end{equation*}
The added term is at most
\begin{equation}\label{4.7:eq-lemma-four-gaps-3}
8(h_n+h')^2e^{(\dots)}\le32(h_n^2+h'^2)=32h_n^2+8(\beta_sL^{-2}\mathsf{t})^2
\le2\beta_s\mathsf{a}+\beta_sL^{-2}\mathsf{t}.
\end{equation}
The steps are justified as follows.
\begin{itemize}
\item $e^{(\dots)}\le2$ under $(\Gamma6)$.
\item $(x+y)^2\le2x^2+2y^2$.
\item $32h'^2=8(\beta_sL^{-2}\mathsf{t})^2$ by the definition of $h'$.
\item $8(\beta_sL^{-2}\mathsf{t})^2\le\beta_sL^{-2}\mathsf{t}$, because $8\beta_sL^{-2}\mathsf{t}\le1$
by $(\Gamma6)$.
\item $32h_n^2\le2\beta_s\mathsf{a}$, as the following lines show.
\end{itemize}
For the last item,
\begin{equation*}
32h_n^2=32(O_{(3.37)}B_3\mathfrak{b}M)^2\alpha_{0,i}^2L^{-2n}
\le2\beta_s\alpha_{0,i}L^{-2n}L^{-2}=2\beta_s\mathsf{a},
\end{equation*}
and this is equivalent to $16L^2(O_{(3.37)}B_3\mathfrak{b}M)^2\alpha_{0,i}\le\beta_s$. That is the
ceiling on $\alpha_{0,i}$ in \eqref{4.7:eq-lemma-four-gaps-a}, and it does not depend on $n$.

Adding the linear and the quadratic parts,
\begin{equation}\label{4.7:eq-lemma-four-gaps-4}
|\partial\mathbf{U}-1|<\bigl[(1+8\beta_s)\mathsf{a}+2\beta_sL^{-2}\mathsf{t}\bigr]\xi_{i'}^2
\le\theta_{\times}\,\mathsf{t}\,\xi_{i'}^2 .
\end{equation}
The second inequality uses $\alpha_{0,i}\le\rho_n\alpha_{0,i'}$ (Lemma~\ref{4.7:lem-radii-comparison})
and the definition of $\theta_{\times}$ in \eqref{4.7:eq-lemma-four-gaps-a}. The sequence
$\rho_nL^{-2(n+1)}$ decreases in $n$, so $n=0$ is the worst case and
\begin{equation*}
\theta_{\times}\le\frac{4(1+\beta_0)(1+8\beta_s)+2\beta_s}{L^2}.
\end{equation*}
Now $4(1+\beta_0)(1+9\beta_s)=4(1+\beta_0)(1+8\beta_s)+4(1+\beta_0)\beta_s$ and
$4(1+\beta_0)\beta_s>2\beta_s$. Hence $\theta_{\times}<1$ follows from condition $(\varphi2)$ of
\eqref{4.7:eq-smallness-conditions-a}, namely $L^2\ge4(1+\beta_0)(1+9\beta_s)$.
Under $(\varphi2)$, moreover, $\theta_{\times}\le\bar{\theta}<1$, where $\bar{\theta}$ of
\eqref{4.7:eq-lemma-four-gaps-a} is the value of the right-hand side at equality in $(\varphi2)$. It is a
pure number of $(\beta_0,\beta_s)$, and the margin $1-\bar{\theta}>0$ is used below.

\emph{The current: second half of condition~(iii).} Ba{\l}aban states that this half ``is obtained in the
same way'' \cite{B-CMP109}, with the factor $(L^{j-1}\eta)^3$ of \cite[(3.42)]{B-CMP109}. We use this
sentence as printed for the member containing $\partial^{*}\partial$, which is therefore bounded, as the
corresponding member of \cite[(3.52)]{B-CMP109}, by $\theta_{\times}\mathsf{t}$. The one further term is
the quadratic correction to the current of the composite, estimated as follows.

By the estimate of the current of an exponential configuration, the quadratic correction to the current
of a configuration $e^{i\xi\mathbf{H}}$ is at most $86d$ times the square of the supremum of
$|\mathbf{H}|$. We read this estimate here for the representative $\mathcal{U}$, and obtain
\begin{equation*}
86d\bigl(\sup_{\tilde{\square}^3}|\mathbf{H}_{i'}(\square_0,D)|\bigr)^2\le86d(h_n+h')^2
\le\beta_sL^{-2}\mathsf{t},
\end{equation*}
where $D=\tau\hat{B}+B'$ is the argument of the representative $\mathcal{U}$.
The last inequality follows from member~(vi) of $(\Gamma6)$ in \eqref{4.7:eq-smallness-conditions-a},
which carries the factor $172dL^2$, in the form with the factor $86d$ derived from that member in the
proof that follows the definition of the conditions \eqref{4.7:eq-smallness-conditions-a}.
The further factors
$L^{-(n+1)}$ and
$\mathsf{s}_n=L^{-3(n+1)}$ that appear in this notation are at most $1$ and only improve the bounds. Hence
\begin{equation}\label{4.7:eq-lemma-four-gaps-5}
|\mathcal{J}|<\theta_J\,\mathsf{t},
\end{equation}
with $\theta_J$ as defined in \eqref{4.7:eq-lemma-four-gaps-a}, which by the two preceding estimates
bounds the factor produced by the argument.
Under $(\varphi2)$,
\begin{equation*}
\theta_J\le\theta_{\times}+2\beta_sL^{-2}\le\frac{4(1+\beta_0)(1+8\beta_s)+4\beta_s}{L^2}
\le\bar{\theta}_J<1 .
\end{equation*}
Indeed
\begin{equation*}
1-\bar{\theta}_J=\frac{4\beta_0\beta_s}{4(1+\beta_0)(1+9\beta_s)}>0,
\end{equation*}
as $\beta_0>0$. One term $\beta_sL^{-2}$ in this bound is room to spare.

\emph{Condition~(iv).} Ba{\l}aban writes ``The condition (iv) follows from the corresponding identity
(3.38)'' \cite[p.~280]{B-CMP109}. That is, condition~(iv) ``is a consequence of the condition (iii), with a
bit better constant'' \cite[(3.38)]{B-CMP109}. We carry this out.

The fields entering condition~(iv) for the image are those of condition~(iii), conjugated by
$g_{\langle3.38\rangle}=u_{i'}\bar{u}_{i'}^{-1}$. The third member of \cite[(3.37)]{B-CMP109} bounds
$|u_{i'}-1|$ without the factor $L^{j-1}\eta$. In the level-indexed notation,
\begin{equation*}
|u_{i'}-1|\le c'\bigl(O_{(3.37)}B_3\mathfrak{b}M\alpha_{0,i}+B_3\alpha^{\sharp}_{3,i'}\bigr),
\end{equation*}
where $c'$ is a constant depending on $(d,L)$, of the same kind as the constant $2dLc_Q$ in the frame
estimates of the axial-frame lemma, and is absorbed in $c$ below. The conjugation costs
\begin{equation*}
\|\operatorname{Ad}_{g_{\langle3.38\rangle}}\|,\ \|R(g_{\langle3.38\rangle})\|\le
e^{c(O_{(3.37)}B_3\mathfrak{b}M\alpha_{0,i}+\beta_s\alpha_{0,i'})}.
\end{equation*}
Here $\operatorname{Ad}_{g_{\langle3.38\rangle}}$ and $R(g_{\langle3.38\rangle})$ are the actions of
$g_{\langle3.38\rangle}$ on curvatures and on currents. The letter
$c$ denotes one constant depending on $(d,L)$ that dominates the costs of conjugation by $u_{i'}$ and
$\bar{u}_{i'}$ (from \cite[(3.37)]{B-CMP109}) and by the averaging gauge $v_{i'}$ of
\cite[(97)]{B-CMP98}. It is the same letter as in the three exponential members of $(\Gamma6)$.

The two conditions
\begin{equation*}
\bar{\theta}\,e^{c(O_{(3.37)}B_3\mathfrak{b}M\alpha_{0,i}+\beta_s\alpha_{0,i'})}\le1,\qquad
\bar{\theta}_J\,e^{c(O_{(3.37)}B_3\mathfrak{b}M\alpha_{0,i}+\beta_s\alpha_{0,i'})}\le1
\end{equation*}
are smallness conditions on the radii of the kind collected in $(\Gamma6)$. The two margins
$1-\bar{\theta}$ and $1-\bar{\theta}_J$ pay for the conjugation. The identity used here is part (iii) of
the conjugation lemma for window minimisers.

\emph{The one-power condition $(\mathrm{iv})^{\Sigma^1}$ for the image.} This is the near part of the
proposition on the image of the step in the one-power class. The same identity is used at every label
$(p,Z)\in\mathfrak{T}_{i'}(X)$. The
conjugation lemma for window minimisers, parts (iii) and (iv), is applied at $Z$ and its tower
$(\Omega_q(Z))_{q\le\mathsf{n}}$.
These data are admissible: $Z\subset X\subset\tilde{\square}^2$ lies inside
$\square_{i'}\supset\tilde{\square}^4$ with a margin $2M$ to spare. The lemma gives
$U_{\mathsf{n}}(Z;\dot{M}\mathcal{U})=(\mathcal{U}\restriction Z)^{g_{\langle3.38\rangle}}$. Hence, on all
of $Z$,
\begin{equation*}
\begin{aligned}
|\partial U_{\mathsf{n}}(Z;\dot{M}\mathcal{U})-1|
&\le\|\operatorname{Ad}_{g_{\langle3.38\rangle}}\|\,\theta_{\times}\,\alpha_{0,i'}\xi_{i'}^2
<\alpha_{0,i'}\xi_{i'}^2\\
&=\alpha_{0,i'}L^{-2(i'-p)}\xi_p^2\le\alpha_{0,i'}L^{-(i'-p)}\xi_p^2 .
\end{aligned}
\end{equation*}
The identity
$J_{\mathsf{n}}(Z;\dot{M}\mathcal{U})=(L^{\mathsf{n}}\xi_{i'})^3R(g_{\langle3.38\rangle})\mathcal{J}$
(part (iv) of the same lemma, and \cite[(3.38)]{B-CMP109}) gives
\begin{equation*}
\begin{aligned}
|J_{\mathsf{n}}(Z;\dot{M}\mathcal{U})|
&\le\|R(g_{\langle3.38\rangle})\|\,\theta_J\,\alpha_{0,i'}L^{-3(i'-p)}(L^{\mathsf{n}}\xi_p)^3\\
&\le\alpha_{0,i'}L^{-(i'-p)}(L^{\mathsf{n}}\xi_p)^2 .
\end{aligned}
\end{equation*}
Here $L^{\mathsf{n}}\xi_p\le1$ because $\mathsf{n}\le p$, which is the range of condition~(iv) itself
and imposes no condition on the gap $n$; and $\|R(g_{\langle3.38\rangle})\|\theta_J\le1$ by the preceding
paragraph. The
curvature bound carries two powers of $L^{-(i'-p)}$ where the definition of the one-power hereditary class
asks for one.

For the factor $U$ of Ba{\l}aban's factorisation \cite[p.~277]{B-CMP109} we take $U:=1$. Then
$U_{\mathsf{n}}(Z;\dot{M}1)=1$ and $J_{\mathsf{n}}(Z;\dot{M}1)=0$, so this half of the condition is void.

\emph{The conditions \cite[(1.11), (1.12)]{B-CMP109} for the image.} They are void at $U:=1$. Indeed
$u=1$ and $A_\square=0$ on every cube $\square\subset X$, so the cube gauge condition of
\eqref{4.7:eq-class-restrictions-a} is satisfied vacuously. The floor $(\varphi B)$ of
\eqref{4.7:eq-smallness-conditions-a}, a lower bound on the constant $B^{109}$ of
\cite[(1.12)]{B-CMP109}, which may be taken as large as needed, is therefore not used in this proof.

\emph{Conclusion.} Following Ba{\l}aban, ``We replace $\tau Q'$ by the variable $B'$ with values in
$\mathfrak{g}^c$, and we prove (3.36) for $B'$ satisfying $|B'|<\alpha_3$'' \cite{B-CMP109}; here the
bound is $|B'|<\alpha^{\sharp}_{3,i'}$ on the bonds of the tower of $\square_0$. Take $\tau\in[0,1]$,
$\hat{B}=Q(\eta\mathcal{A})$ and $U:=1$. The conditions defining
$U^{\Sigma^1}_{i'}(X;\alpha_{0,i'},\alpha_{1,i'})$ are met as follows.
\begin{itemize}
\item Condition~(i), condition \cite[(1.12)]{B-CMP109} and the $U$-half of condition~(iv) are void.
\item Condition~(ii) follows from \eqref{4.7:eq-lemma-four-gaps-2} under $(\varphi1)$, together with
$B_3\alpha^{\sharp}_{3,i'}=\tfrac12\beta_sL^{-2}\alpha_{0,i'}<\tfrac12\alpha_{1,i'}$. This last
inequality holds because $C_0\le C_1$ in the working radii and $\beta_sL^{-2}<1$.
\item Condition~(iii) follows from \eqref{4.7:eq-lemma-four-gaps-4} and \eqref{4.7:eq-lemma-four-gaps-5}.
\item The $\mathbf{U}$-half of $(\mathrm{iv})^{\Sigma^1}$ follows from the near part of the proposition on
the image of the step in the one-power class, as shown above.
\end{itemize}
All bounds hold uniformly in $(\tau,B')$ over these ranges. This is the last line of
\eqref{4.7:eq-lemma-four-gaps-a}.
\end{proof}

\begin{lemma}[The fifth-order remainder and the small factor]\label{4.7:prf-fifth-order}
Work at a pair of levels $i'$, $i$ at level gap $n$, in the setting of Ba{\l}aban's Lemma~4 at every
level gap (Lemma~\ref{4.7:prf-lemma-four-gaps}), with the parametrised configuration
$U_{i'}(\square_0,\cdot)$ of \eqref{4.7:eq-lemma-four-gaps-a}, and with the kernel
$E^{(i')}(X,\cdot)$, which is analytic on the one-power hereditary class
$U^{\Sigma^1}_{i'}(X;\alpha_{0,i'},\alpha_{1,i'})$ and bounded there by
$S_{i'}e^{-\kappa d_{i'}(X)}$. Use the working radii
$\alpha^{\sharp}_{3,i'}$, $\alpha_{2,i}$ and $b_i$ of Definition~\ref{4.7:not-working-radii}, and the
total field $B$ and the ratio $\Theta_n$ of Lemma~\ref{4.7:lem-total-field}. Let $\chi_4$ be the
indicator of the bonds in $\tilde{\square}^{4}$. For $\tau\in[0,1]$ let
$D^5_\tau[\chi_4B]^{\otimes 5}$ be the fifth derivative of
$B''\mapsto E^{(i')}(X,U_{i'}(\square_0,e^{\mathrm{i}B''}))$ at $B''=\tau B$, evaluated on five
copies of $\chi_4B$. Let $F$ be any function of the cube parameter $t_\square$ that is analytic on
the closed disc $|t_\square|\le\alpha_{2,i}/(2C_s\delta_{i-1})$; in particular $F$ may be any term
of the expansion \cite[(3.34)]{B-CMP109}, or any combination of its terms of orders two to four
that remains once the marginal part has been cancelled by means of the identity expressing those
terms as a marginal part plus an irrelevant part. Then
\begin{equation}\label{4.7:eq-fifth-order-a}
\begin{gathered}
s_n:=\frac{\tfrac14\,\alpha^{\sharp}_{3,i'}}{b_iL^{-n}},\qquad
\Bigl|\int_0^1 d\tau\,\frac{(1-\tau)^4}{4!}\,D^5_\tau[\chi_4B]^{\otimes 5}\Bigr|
\le S_{i'}\,\Theta_n^5\,L^{-5n}\,e^{-\kappa d_{i'}(X)},\\
\bigl|\partial_{t_\square}F\big|_{t_\square=0}\bigr|
\le \vartheta_i\,\sup_{|t_\square|=\alpha_{2,i}/(2C_s\delta_{i-1})}|F| .
\end{gathered}
\end{equation}
The terms of the fifth-order remainder of \cite[(3.34)]{B-CMP109} containing at least one factor
$(1-\chi_4)B$ are bounded as stated in \cite[\S3]{B-CMP109}; see the fourth step of the proof.
\end{lemma}

\begin{proof}
\emph{First step: the total field and its correction.} By \cite[(3.48)]{B-CMP109} the total field
splits as $B=\hat{B}+Q'$, and by \cite[(3.49)]{B-CMP109} the correction satisfies, on
$\square_0$ and in the notation of that display,
\begin{equation*}
|Q'|\le O_{(3.49)}\,L^{j}\eta\,|\mathbf{A}| < O_{(3.49)}\,\alpha_2\,L^{j}\eta ,
\end{equation*}
with an absolute constant $O_{(3.49)}$. On $\tilde{\square}^{4}$ the total field is bounded by
Lemma~\ref{4.7:lem-total-field}, and the field $\hat{B}$ satisfies the same bound at the present
level gap:
\begin{equation}\label{4.7:eq-fifth-order-1}
\sup_{\tilde{\square}^{4}}|B|\le b_iL^{-n},\qquad \sup_{\tilde{\square}^{4}}|\hat{B}|\le b_iL^{-n}.
\end{equation}
At the present pair of levels $L^{j}\eta$ is $L^{-n}$ and $\alpha_2$ is $\alpha_{2,i}$, so the bound
on $Q'$ reads $\sup|Q'|\le O_{(3.49)}\,\alpha_{2,i}L^{-n}$. By the comparison of working radii with
constant $c_\alpha$ in \eqref{4.7:eq-working-radii-a}, wherever $Q'$ is defined,
\begin{equation}\label{4.7:eq-fifth-order-2}
\sup|Q'|\le O_{(3.49)}\,\alpha_{2,i}L^{-n}
\le \tfrac14\,\alpha^{\sharp}_{3,i'}\cdot\tfrac12
\Bigl[\frac{\alpha_{0,i-1}}{\alpha_{0,i'}}\Bigr]L^{-n}
\le \tfrac14\,\alpha^{\sharp}_{3,i'} .
\end{equation}

\emph{Second step: a circle in the localised direction about each point of the segment.} The
localisation of the inner product in the domains $\tilde{\square}^{4}$ and
$(\tilde{\square}^{4})^{c}$ is the one of \cite[\S3]{B-CMP109}. In \cite[\S4]{B-CMP109} the
corresponding field carries the localising factor $\tilde{\zeta}_\square$ of
\cite[(4.16), p.~285]{B-CMP109}. Here the localisation is effected by $\chi_4$, with $s_n$ as in
\eqref{4.7:eq-fifth-order-a}. Fix $\tau\in[0,1]$ and let $|\sigma|\le s_n$. Then
\begin{equation*}
\tau B+\sigma\chi_4B=\tau\hat{B}+B',\qquad B':=\tau Q'+\sigma\chi_4B .
\end{equation*}
The field $\chi_4B$ is read only on $\tilde{\square}^{4}$, where \eqref{4.7:eq-fifth-order-1}
holds. Hence, by \eqref{4.7:eq-fifth-order-2}, by \eqref{4.7:eq-fifth-order-1} and by the
definition of $s_n$, on every bond of the tower,
\begin{equation*}
|B'|\le \tau\sup|Q'|+|\sigma|\sup_{\tilde{\square}^{4}}|B|
\le \tfrac14\,\alpha^{\sharp}_{3,i'}+s_n\,b_iL^{-n}
=\tfrac12\,\alpha^{\sharp}_{3,i'} .
\end{equation*}
Consider, for each $\tau\in[0,1]$, the function
\begin{equation}\label{4.7:eq-fifth-order-3}
f_\tau(\sigma):=E^{(i')}\bigl(X,U_{i'}(\square_0,e^{\mathrm{i}(\tau B+\sigma\chi_4B)})\bigr),
\qquad |\sigma|\le s_n .
\end{equation}
By the conclusion of Ba{\l}aban's Lemma~4 at every level gap, \eqref{4.7:eq-lemma-four-gaps-a},
a configuration $U_{i'}(\square_0,e^{\mathrm{i}(\tau\hat{B}+B')})$ with the bounds just obtained
lies in the one-power hereditary class $U^{\Sigma^1}_{i'}(X;\alpha_{0,i'},\alpha_{1,i'})$. On that
class the kernel is analytic and bounded by $S_{i'}e^{-\kappa d_{i'}(X)}$. Hence $f_\tau$ is
analytic on $|\sigma|\le s_n$ and satisfies $|f_\tau(\sigma)|\le S_{i'}e^{-\kappa d_{i'}(X)}$
there. This holds for every $\tau\in[0,1]$.

The radius $s_n$ may be smaller than one. No condition is needed on that account, because the
circle is centred at the point $\tau$ of the segment and not at $0$. In particular, no condition
at far levels enters.

\emph{Third step: the Cauchy bound for the localised fifth-order term.} By the definition of
$D^5_\tau$ in the statement, $D^5_\tau[\chi_4B]^{\otimes5}=f_\tau^{(5)}(0)$. Cauchy's inequality on
the circle $|\sigma|=s_n$ gives
\begin{equation*}
\Bigl|\Bigl(\frac{d}{d\sigma}\Bigr)^{5}f_\tau\Big|_{\sigma=0}\Bigr|
\le 5!\,s_n^{-5}\,\sup_{|\sigma|=s_n}|f_\tau|
\le 5!\,s_n^{-5}\,S_{i'}e^{-\kappa d_{i'}(X)} .
\end{equation*}
Since $\int_0^1(1-\tau)^4/4!\,d\tau=1/5!$, integrating this bound against $(1-\tau)^4/4!$ over
$[0,1]$ bounds the fifth-order term of \cite[(3.34)]{B-CMP109} with all five arguments equal to
$\chi_4B$:
\begin{align*}
\Bigl|\int_0^1 d\tau\,\frac{(1-\tau)^4}{4!}\,D^5_\tau[\chi_4B]^{\otimes5}\Bigr|
&\le S_{i'}e^{-\kappa d_{i'}(X)}\,s_n^{-5}
= S_{i'}e^{-\kappa d_{i'}(X)}\Bigl(\frac{4b_iL^{-n}}{\alpha^{\sharp}_{3,i'}}\Bigr)^{5}\\
&\le S_{i'}\,\Theta_n^5\,L^{-5n}\,e^{-\kappa d_{i'}(X)} .
\end{align*}
The last inequality holds because $\Theta_n$ is defined in Lemma~\ref{4.7:lem-total-field} as
$b_i$ times the maximum of $4/\alpha^{\sharp}_{3,i'}$ and a second quantity. Hence
$\Theta_n\ge 4b_i/\alpha^{\sharp}_{3,i'}$. This is the first inequality in
\eqref{4.7:eq-fifth-order-a}.

\emph{Fourth step: the terms carrying a factor outside $\tilde{\square}^{4}$.} Write
$B=\chi_4B+(1-\chi_4)B$ and expand the fivefold multilinear form $D^5_\tau[B]^{\otimes5}$. The
fifth-order term of \cite[(3.34)]{B-CMP109} is then the term bounded in the third step plus terms
with at least one factor $(1-\chi_4)B$. These latter terms are not estimated here. For them we
take Ba{\l}aban's sentence in \cite[\S3]{B-CMP109} as printed: they are bounded by an
``arbitrary power of $L^{j}\eta$ using the exponential decay properties \dots{} we do not repeat
them''.

\emph{Fifth step: the circle in the cube parameter.} Every term of \cite[(3.34)]{B-CMP109} is
analytic in $t_\square$ on the disc of the cube parameter on which
Lemma~\ref{4.7:prf-lemma-four-gaps} is set. So is the identity expressing the terms of orders two
to four as a marginal part plus an irrelevant part. The bounds of the first three steps hold
uniformly on that disc, because the only input on the fluctuation field used in them is
\cite[(3.31)]{B-CMP109}. One may therefore first carry out the cancellation in that identity and
differentiate in $t_\square$ afterwards; the combination that remains after the cancellation is a
finite combination of functions analytic on that disc, hence analytic on it as well. This is the
order in which the argument of \cite{B-CMP119} proceeds.

The derivative $\partial_{t_\square}$ of \cite{B-CMP109} is taken at $t_\square=0$. Let $F$ be as
in the statement. The circle $|t_\square|=\alpha_{2,i}/(2C_s\delta_{i-1})$ lies in the disc of the
cube parameter just described, and $F$ is analytic on the closed disc it bounds. Cauchy's
inequality on this circle gives
\begin{equation*}
\bigl|\partial_{t_\square}F\big|_{t_\square=0}\bigr|
\le \frac{2C_s\delta_{i-1}}{\alpha_{2,i}}\,
\sup_{|t_\square|=\alpha_{2,i}/(2C_s\delta_{i-1})}|F|
\le \vartheta_i\,\sup_{|t_\square|=\alpha_{2,i}/(2C_s\delta_{i-1})}|F| .
\end{equation*}
The second inequality is the bound defining the small factor in
\eqref{4.7:eq-small-factor-a}. This is the second inequality in \eqref{4.7:eq-fifth-order-a}. The
factor $\vartheta_i$ is the one small factor produced by differentiation in the cube parameter.
\end{proof}

\begin{lemma}[Fundamental-case share and counterterm leftovers]\label{4.7:prf-fundamental-share}
Let $i'$ and $i$ be the pair of levels at gap $n$ considered in this section. Let $S_{i'}$ and $S_j$ be
the sup-bounds of the kernel families at levels $i'$ and $j$, where $j$ is the level of the marked
family $\mathfrak{M}_j$ whose own counterterms are treated in (b). Let $b_i$, $\alpha^{\sharp}_{3,i'}$ be
the working radii of \eqref{4.7:eq-working-radii-a}, $c_T$ the constant of
\eqref{4.7:eq-true-pairs-a}, $\Theta_n$, $\tilde{K}$, $\rho_n$ as in
\eqref{4.7:eq-total-field-a}, $r_{i'}$ and $r_j$ the radii at levels $i'$ and $j$ that enter
\eqref{4.7:eq-total-field-a} through the product $c_Tr_{i'}$, and $\vartheta_i$ the small
factor of \eqref{4.7:eq-small-factor-a}. In what follows $C$ denotes a generic constant collecting
counting and decay factors, and $c$ is a constant free of the coupling.
\begin{enumerate}
\item[(a)] The fifth-order part of the level-$i'$ fundamental share is bounded by
\begin{equation}\label{4.7:eq-fundamental-share-1}
S_{i'}\,\vartheta_i\,c\,(6L)^4\,\tilde{K}^5\rho_n^5\,L^{-n}.
\end{equation}
The part of orders two to four is bounded by
\begin{equation}\label{4.7:eq-fundamental-share-2}
C\,S_{i'}\,\vartheta_i\,\tilde{K}^4\rho_n^4\,L^{-\beta_H n}.
\end{equation}
\item[(b)] The leftovers of the own counterterms of the marked family $\mathfrak{M}_j$ are bounded by
\begin{equation}\label{4.7:eq-fundamental-share-3}
C\,S_j\,\vartheta_i\,\tilde{K}^2\rho_n^2\,L^{-\beta_H n}.
\end{equation}
\end{enumerate}
\end{lemma}

\begin{proof}
\emph{(a), fifth order.} Fix a cube $\square$, a window cube $\square_0$, a localisation domain $Y_0$
and a polymer $X$. By the fifth-order remainder bound (Lemma~\ref{4.7:prf-fifth-order},
display~\eqref{4.7:eq-fifth-order-a}), the contribution of this quadruple is at most
\begin{equation*}
C\,S_{i'}\,\vartheta_i\,\Theta_n^5\,L^{-5n}\,e^{-\kappa d_{i'}(X)},
\end{equation*}
multiplied by the decay factor of \cite[(1.25)]{B-CMP116}. The counting of the cubes $\square'$
contributes a factor at most $(6L)^4L^{4n}$. The remaining sums are performed with
\cite[(1.26)--(1.28)]{B-CMP116}, whose constants do not depend on the coupling; they are collected
in $c$. This gives the bound
\begin{equation*}
c\,(6L)^4\,S_{i'}\,\vartheta_i\,\Theta_n^5\,L^{4n}L^{-5n}
=c\,(6L)^4\,S_{i'}\,\vartheta_i\,\Theta_n^5\,L^{-n}.
\end{equation*}
By Lemma~\ref{4.7:lem-total-field}, $\Theta_n\le\tilde{K}\rho_n$, so
$\Theta_n^5\le\tilde{K}^5\rho_n^5$. This is \eqref{4.7:eq-fundamental-share-1}.

\emph{(a), orders two to four.} Let $E^{(m)}$, $2\le m\le4$, denote the kernels of order $m$ at
$\mathbf{U}=1$. They are obtained by Cauchy's formula applied to \cite[(4.4)]{B-CMP109}.

Cauchy's formula may be applied at true pairs for the following reason. By
\eqref{4.7:eq-true-pairs-a} (Notation~\ref{4.7:not-true-pairs}), true pairs lie in the regularity
class, and by the true-pairs clause of the membership proposition for the one-power hereditary
class they lie in the one-power hereditary class
$U^{\Sigma^1}_{i'}$. The expansion is controlled through the bound of the form
$B_3\prod_{\nu}|B_{\nu}|$ of \cite[\S4]{B-CMP109}, and the radius available to Cauchy's formula is
at least $\min\{c_T r_{i'},\tfrac14\alpha^{\sharp}_{3,i'}\}$. Hence
\begin{equation*}
|E^{(m)}|\le C\,S_{i'}\,\bigl(\min\{c_T r_{i'},\tfrac14\alpha^{\sharp}_{3,i'}\}\bigr)^{-m}.
\end{equation*}

A term of order $m$ carries $m\le4$ total fields. Each of them is bounded by $b_i$ times its own
power of $L^{-n}$: the field by $b_iL^{-n}$, its derivative by $b_iL^{-2n}$, and its second
derivatives by $b_iL^{-(2+\beta_H)n}$. This follows from Lemma~\ref{4.7:lem-total-field}, from
\cite[(4.17)--(4.18)]{B-CMP109}, and from the bound of \cite[p.~286]{B-CMP109}, in whose notation
$j$ denotes Ba{\l}aban's level index (not the level of part (b)) and $\alpha_1$ is Ba{\l}aban's
constant,
\begin{equation*}
|(\partial B_{\mu_3})(\Gamma_{x,x_3})|<\alpha_1(L^j\eta)^2|\Gamma_{x,x_3}|
\end{equation*}
along the paths $\Gamma_{x,x_3}$. The path lengths are absorbed as in the proof of
\cite[(4.22)]{B-CMP109}.

Unwinding the definition of $\Theta_n$ in \eqref{4.7:eq-total-field-a} gives
\begin{equation*}
\frac{b_i}{\min\{c_T r_{i'},\tfrac14\alpha^{\sharp}_{3,i'}\}}
=b_i\max\Bigl\{\frac{4}{\alpha^{\sharp}_{3,i'}},\,(c_T r_{i'})^{-1}\Bigr\}=\Theta_n .
\end{equation*}
Therefore each leftover term is at most $C\,S_{i'}\,\Theta_n^m\,L^{-\varpi n}$, where $\varpi$ is the
total power of $L^{-n}$ carried by the fields. The leftover terms are of four kinds
$(\alpha)$--$(\delta)$; $\varpi$ equals $5$ for those of kind $(\alpha)$,
$4+\beta_H$ for those of kinds $(\beta)$ and $(\gamma)$, and is at least $5$ for those of
kind $(\delta)$. The smallest of these powers is $4+\beta_H$, so in every case
$\varpi\ge4+\beta_H$.

One factor $\vartheta_i$ is gained from the Cauchy circle in the cube parameter $t_\square$, as in
\eqref{4.7:eq-small-factor-a}. The number of cubes is at most $L^{4n}$, and
$L^{4n}L^{-\varpi n}\le L^{-\beta_H n}$.

It remains to bound $\Theta_n^m$. Since $\tilde{K}\ge1$, and $\rho_n\ge1$ by its definition in
Lemma~\ref{4.7:lem-total-field}, we
have $\tilde{K}\rho_n\ge1$. Together with Lemma~\ref{4.7:lem-total-field} and $m\le4$ this gives
$\Theta_n^m\le(\tilde{K}\rho_n)^m\le\tilde{K}^4\rho_n^4$. This is
\eqref{4.7:eq-fundamental-share-2}.

\emph{(b).} The coefficient of the own counterterm of the marked family $\mathfrak{M}_j$ satisfies
\begin{equation*}
|\beta^{\mathfrak{M}}_j|\le C_{\Pi}\,S_j\,(c_T r_j)^{-2},
\end{equation*}
by the bound on the marked-family coefficient established in the Cauchy extraction of the beta
coefficient leading to \eqref{4.7:eq-beta-extraction-a}. That bound is proved by induction on the
level; it is used here at the level $j\le i-1$, at which it belongs to the induction hypothesis
$(H_{i-1})$ in force at this stage.

Consider the action $A$ evaluated on the minimising configuration $U_j(\square_0,e^{\mathrm{i}B})$
with data $e^{\mathrm{i}B}$ on the window cube $\square_0$. Its Taylor kernels in $B$ are bounded by a
constant $c_A$ that is independent of $S_j$ and of the working radii. The reason is twofold. First,
$A$ is entire
in the bond variables. Second, the minimiser is analytic in its data on a ball whose size does not
depend on these quantities: by \cite[p.~277]{B-CMP102b}, the constants $a_0$, $a_1$, $B_3$ depend
on $d$ and $L$ only. This is the content of the lemma on the absolute Taylor constant of the action
at the minimiser.

These counterterm contributions are processed by the same procedure as the other terms, with the
same Ward structure. The terms \cite[(4.42)]{B-CMP109} and \cite[(4.44)]{B-CMP109} cancel
identically. The remaining terms are of the kinds $(\alpha)$--$(\delta)$ above and are controlled as
in \cite[\S4]{B-CMP109}: the sum over $j$ of these terms, each carrying a constant coefficient, is
bounded uniformly. By
\cite[(4.39)--(4.40)]{B-CMP109}, the expansion about $\mathbf{U}=1$ starts at order two.

Consequently each uncancelled piece is at most
\begin{equation*}
C_{\Pi}\,S_j\,(c_T r_j)^{-2}\,c_A\,b_i^m\,L^{-\varpi n}\,\vartheta_i,
\qquad m\ge2,\quad \varpi\ge4+\beta_H .
\end{equation*}
We have
\begin{equation*}
(c_T r_j)^{-2}\,b_i^m\le\Theta_n^2\,b_i^{m-2}\le\Theta_n^2 ,
\end{equation*}
where the second inequality holds because $b_i\le1$, which is condition $(\Gamma6)$ of
\eqref{4.7:eq-smallness-conditions-a}.
Lemma~\ref{4.7:lem-total-field} gives $\Theta_n^2\le\tilde{K}^2\rho_n^2$. After the count of at most
$L^{4n}$ cubes, as in (a), we have $L^{4n}L^{-\varpi n}\le L^{-\beta_H n}$. Absorbing $C_{\Pi}$ and
$c_A$ into $C$ gives \eqref{4.7:eq-fundamental-share-3}.
\end{proof}

\begin{lemma}[Simple-case share and restriction to an old level]\label{4.7:prf-simple-share}
Let $i$ be a level, let $n\ge 0$ and put $i'=i-n-1$. Let $X\subset Y$ be domains, and write
$\mathsf{s}_n=L^{-3(n+1)}$. Assume the floor $(\varphi4)$ of
\eqref{4.7:eq-smallness-conditions-a}, that is
\begin{equation}\label{4.7:eq-simple-share-1}
L\ \ge\ 4(1+\beta_0)(1+2\beta_s).
\end{equation}
\begin{enumerate}
\item[(a)] Restriction to the old level $i'$:
\begin{equation}\label{4.7:eq-simple-share-2}
\begin{split}
&U^{\Sigma^1}_{i}\bigl(Y;(1+\beta_s)\alpha_{0,i},(1+\beta_s)\alpha_{1,i},\alpha_{0,i}\bigr)
\upharpoonright X\\
&\qquad\subset\ U^{\Sigma^1}_{i'}\bigl(X;\tfrac12\alpha_{0,i'},\tfrac12\alpha_{1,i'},
\alpha_{0,i'}\bigr),
\end{split}
\end{equation}
where a pair $(\mathbf{U},\mathbf{J})$ of the left-hand space enters the right-hand space with
the current $\mathsf{s}_n\mathbf{J}$ in the current slot.
\item[(b)] Assume in addition the conditions \eqref{4.7:eq-smallness-conditions-a}. In the
simple case, in which $\alpha_2=c_T r_{i'}$, the pairs produced by the operator bounds at free
current from configurations of the left-hand space of \eqref{4.7:eq-simple-share-2},
restricted to $X$, lie in
$U^{\Sigma^1}_{i'}(X;\alpha_{0,i'},\alpha_{1,i'})$, and every term of the bound
\cite[(3.17)]{B-CMP109}, in its running form, is at most
\begin{equation}\label{4.7:eq-simple-share-3}
3S_{i'}\,\vartheta_i\,(1+n)^{1/2}L^{-n}\,e^{-\frac12 c_*\delta_0 M L^{n}}
\end{equation}
times the decay factor of the term. Here $\vartheta_i$ is the small factor of
\eqref{4.7:eq-small-factor-a}, $c_*$ is the rate constant and $\delta_0$ the decay rate of the
operator bounds at free current, and $S_{i'}$ is the sup constant of the marked kernel family
$\mathfrak{M}_{i'}$ at level $i'$.
\end{enumerate}
\end{lemma}

\begin{proof}
\emph{The role of the inclusion.} The bound \cite[(3.17)]{B-CMP109} requires the base
configuration to lie in the space $U^{c}_j(X,\frac12\alpha_0,\frac12\alpha_1,\alpha_0)$ of
\cite[(3.16)]{B-CMP109}; here that space is read as the one-power hereditary class
$U^{\Sigma^1}$. The paper \cite{B-CMP116} states that
\begin{quote}
All these spaces contain the subspace $U^c_{k+1}(Y,(1+\beta)\alpha_0,(1+\beta)\alpha_1,
\alpha_0)$
\end{quote}
with radii that do not depend on the level. For the one-power hereditary class this is the
restriction to an old level stated in
\eqref{4.7:eq-class-restrictions-a} (Notation~\ref{4.7:not-class-restrictions}), and
\eqref{4.7:eq-simple-share-2} is its form with the radii of the two levels $i$ and $i'$. The
second member of the requirement \cite[(3.16)]{B-CMP109} consists of the clauses (i)--(iii)
of the space on the left of \eqref{4.7:eq-simple-share-2} itself, and one and the same
representative configuration serves both members, by the reading convention $(\rho4)$ of
Notation~\ref{4.5:not-reading-conventions}.

\emph{Step 1: clauses (i)--(iii), change of units.} Let $(\mathbf{U},\mathbf{J})$ belong to
the left-hand space of \eqref{4.7:eq-simple-share-2}, and keep the decomposition
$\mathbf{U}=U'U$ which exhibits its membership at level $i$. The spacings satisfy
$\xi_i=L^{-(n+1)}\xi_{i'}$. By \cite[(1.11)]{B-CMP109} and the first bound of
\cite[(1.14)]{B-CMP109}, in the notation of those displays,
\begin{equation}\label{4.7:eq-simple-share-4}
|\partial U-1|,\ |\partial\mathbf{U}-1|\ <\ (1+\beta_s)\alpha_{0,i}\,\xi_i^2
\ =\ (1+\beta_s)\alpha_{0,i}\,L^{-2(n+1)}\,\xi_{i'}^2 .
\end{equation}
By \cite[(1.13)]{B-CMP109} the potential $A'$ of $U'$ satisfies
$e^{i\xi_iA'}=e^{i\xi_{i'}A'_{i'}}$ with $A'_{i'}:=L^{-(n+1)}A'$, and in the units of
level $i'$
\begin{equation}\label{4.7:eq-simple-share-5}
|A'_{i'}|<(1+\beta_s)\alpha_{1,i}L^{-(n+1)},\qquad
|\nabla^{\xi_{i'}}A'_{i'}|<(1+\beta_s)\alpha_{1,i}L^{-2(n+1)} .
\end{equation}
The current slot receives $\mathsf{s}_n\mathbf{J}$ with $\mathsf{s}_n=L^{-3(n+1)}$, by
\cite[(3.12)]{B-CMP109}. The powers of $L^{-(n+1)}$ are those of the change of units: two for
plaquette variables, as in \eqref{4.7:eq-simple-share-4}, one for a potential, as in the
definition of $A'_{i'}$, one more for each difference quotient, as in
\eqref{4.7:eq-simple-share-5}, and three for a current, by \cite[(3.12)]{B-CMP109}.

\emph{Step 2: the gauge condition of \cite[(1.12)]{B-CMP109}.} This condition is read as in
\eqref{4.7:eq-class-restrictions-a}, with the constants $\mathsf{o}$ and $B^{109}$ of
Definition~\ref{4.4:def-local-data-classes}. Let $\square\subset X$ be a cube of side at most
$\mathsf{o}LM$ units of level $i'$. In units of level $i$ its side is at most
$\mathsf{o}LM\cdot L^{-(n+1)}<\mathsf{o}LM$, and $\square\subset X\subset Y$. Hence the gauge
condition for $Y$ at level $i$ supplies a gauge transformation $u$ on $\square$ with
\begin{equation*}
U^{u}=e^{i\xi_iA_\square}=e^{i\xi_{i'}A_{\square,i'}},\qquad
A_{\square,i'}:=L^{-(n+1)}A_\square ,
\end{equation*}
and
\begin{equation}\label{4.7:eq-simple-share-6}
\begin{split}
|A_{\square,i'}|&<\mathsf{o}LMB^{109}(1+\beta_s)\alpha_{0,i}L^{-(n+1)},\\
|\nabla^{\xi_{i'}}A_{\square,i'}|&<\mathsf{o}LMB^{109}(1+\beta_s)\alpha_{0,i}L^{-2(n+1)}.
\end{split}
\end{equation}
A domain $Y$ may be too thin to contain any cube of side $\mathsf{o}L^2M$ in the units of the
current lattice; the cube size
in the gauge condition is therefore taken as in \eqref{4.7:eq-class-restrictions-a}, on the
ground of the statement of \cite{B-CMP116} quoted above.

By \eqref{4.7:eq-simple-share-4}--\eqref{4.7:eq-simple-share-6}, every clause of the first
member of the right-hand space of \eqref{4.7:eq-simple-share-2} other than clause (iv) holds
at the radii $\frac12\alpha_{0,i'}$, $\frac12\alpha_{1,i'}$ as soon as
\begin{equation}\label{4.7:eq-simple-share-7}
(1+\beta_s)\,\frac{\alpha_{\cdot,i}}{\alpha_{\cdot,i'}}\,L^{-(n+1)}\ \le\ \frac12 ,
\end{equation}
for $\alpha_{\cdot,\,\cdot}$ either of $\alpha_{0,\,\cdot}$, $\alpha_{1,\,\cdot}$. Only one
power of $L^{-(n+1)}$ is needed: it is the power in the sup bounds of
\eqref{4.7:eq-simple-share-5} and \eqref{4.7:eq-simple-share-6}. The bounds
\eqref{4.7:eq-simple-share-4} and the gradient bounds carry $L^{-2(n+1)}\le L^{-(n+1)}$.

\emph{Step 3: clause (iv) of the one-power hereditary class.} For a domain $W$ and a level
$j$ let $\mathfrak{T}_j(W)$ be the finite label set of the one-power hereditary class on $W$
at level $j$ (Notation~\ref{4.7:not-class-restrictions}); its labels are pairs $(p,Z)$ with
$p\le j$, $Z\in\mathbf{D}_{p-1}$ and $Z\subset W$. If $(p,Z)\in\mathfrak{T}_{i'}(X)$, then
$p\le i'\le i$, $Z\in\mathbf{D}_{p-1}$ and $Z\subset X\subset Y$. Therefore, by the definition
of the label set,
\begin{equation*}
\mathfrak{T}_{i'}(X)\subset\mathfrak{T}_i(Y).
\end{equation*}
For $(p,Z)\in\mathfrak{T}_{i'}(X)$ the left-hand space of \eqref{4.7:eq-simple-share-2}
asserts clause (iv) at $(p,Z)$ with the constant $(1+\beta_s)\alpha_{0,i}L^{-(i-p)}$. Since
$i-p=(n+1)+(i'-p)$, this constant is
\begin{equation*}
(1+\beta_s)\alpha_{0,i}L^{-(n+1)}\cdot L^{-(i'-p)} .
\end{equation*}
Clause (iv) at $(p,Z)$ is one statement about $\mathbf{U}\upharpoonright Z$ and
$U\upharpoonright Z$, whichever space quotes it (Notation~\ref{4.7:not-class-restrictions}).
If $(1+\beta_s)[\alpha_{0,i}/\alpha_{0,i'}]L^{-(n+1)}\le\frac12$, the constant above is at most
$\frac12\alpha_{0,i'}L^{-(i'-p)}$, the constant of the target clause at $(p,Z)$. This is the
same single power of $L^{-(n+1)}$ as in \eqref{4.7:eq-simple-share-7}.

\emph{Step 4: the floor.} The inclusion \eqref{4.7:eq-simple-share-2} involves the factor
$1+\beta_s$, while the floor \eqref{4.7:eq-simple-share-1} is stated with $1+2\beta_s$. Since
$1+\beta_s\le1+2\beta_s$, it suffices to prove \eqref{4.7:eq-simple-share-7} with
$(1+2\beta_s)$ in place of $(1+\beta_s)$. We use the comparability of radii across a level gap
(Lemma~\ref{4.7:lem-radii-comparison}) and the floor \eqref{4.7:eq-simple-share-1}.

If $n=0$, then $i'=i-1$, and the middle ratio of that lemma gives
$\alpha_{\cdot,i}/\alpha_{\cdot,i-1}\le2(1+\beta_0)$. Hence the left side is at most
\begin{equation*}
\frac{2(1+\beta_0)(1+2\beta_s)}{L}\ \le\ \frac12 .
\end{equation*}

If $n\ge1$, the lemma gives $\alpha_{\cdot,i}/\alpha_{\cdot,i'}\le\rho_n$, with $\rho_n$ as in
\eqref{4.7:eq-radii-comparison-a}. By \eqref{4.7:eq-simple-share-1},
\begin{equation}\label{4.7:eq-simple-share-8}
\rho_n(1+2\beta_s)L^{-(n+1)}
=\frac{4(1+\beta_0)(1+2\beta_s)}{L}\,(1+n)^{1/2}L^{-n}\ \le\ (1+n)^{1/2}L^{-n}.
\end{equation}
The ratio of consecutive terms of $(1+n)^{1/2}L^{-n}$ is $((n+2)/(n+1))^{1/2}L^{-1}<1$.
Hence, for $n\ge1$,
\begin{equation*}
(1+n)^{1/2}L^{-n}\ \le\ \frac{\sqrt2}{L}\ \le\ \frac12 ,
\end{equation*}
the last inequality because $L\ge4$ by \eqref{4.7:eq-simple-share-1}. In both cases
\eqref{4.7:eq-simple-share-7} holds, so Steps 2 and 3 apply.

The current slot requires $\mathsf{s}_n\alpha_{0,i}\le\alpha_{0,i'}$. For $n=0$ the middle
ratio gives $\alpha_{0,i}\le2(1+\beta_0)\alpha_{0,i-1}$, and $2(1+\beta_0)\le\rho_0$ with
$i-1=i'$; for $n\ge1$ the lemma gives $\alpha_{0,i}\le\rho_n\alpha_{0,i'}$. So in both cases
$\alpha_{0,i}\le\rho_n\alpha_{0,i'}$. The inequality of \eqref{4.7:eq-simple-share-8} holds for
every $n\ge0$. Combined with
$(1+n)^{1/2}L^{-n}\le1$, it gives $\rho_nL^{-(n+1)}\le1$. Therefore
\begin{equation}\label{4.7:eq-simple-share-9}
\mathsf{s}_n\alpha_{0,i}\ \le\ \rho_nL^{-3(n+1)}\alpha_{0,i'}\ \le\ L^{-2(n+1)}\alpha_{0,i'}
\ \le\ \alpha_{0,i'} .
\end{equation}
With Steps 1--3 this proves (a).

\emph{Step 5: the produced pairs.} Assume now \eqref{4.7:eq-smallness-conditions-a}. In the
simple case $\alpha_2=c_Tr_{i'}$, with $c_T$ as in \eqref{4.7:eq-true-pairs-a}. Clause (F1)
of the operator bounds at free current produces, from configurations of the left-hand space
of \eqref{4.7:eq-simple-share-2} restricted to $X$, which lie in the right-hand space by (a),
the pairs
\begin{equation*}
\bigl(e^{i\xi_{i'}A(t_\square)}\,\mathbf{U}\upharpoonright X,\ \
\mathsf{s}_n\mathbf{J}+\mathfrak{J}\bigr),
\qquad |A(t_\square)|,\ |\nabla_{\mathbf{U}}A(t_\square)|\le\tfrac23\alpha_2 .
\end{equation*}
Here $A(t_\square)$ is the potential that clause (F1) attaches to the cube $\square$ at the
value $t_\square$ of the cube parameter, and $\nabla_{\mathbf{U}}A(t_\square)$ is its
derivative in the background configuration.
These pairs lie in $U^{\Sigma^1}_{i'}(X;\alpha_{0,i'},\alpha_{1,i'})$ by the far clause of the
proposition that places the produced pairs in the regularity class. That clause applies, by
statement and at every label $(p,Z)$, the admissible-tower lemma, that is, the lemma on towers
of localisation domains which, applied to a tower of $X$, places a pair in the class attached
to that tower; it does so with the parameters $M/L$,
$LB^{109}$ and the third parameter of that lemma as fixed in the far clause, using its clause
for the one-power hereditary class and its uniformity in $M$. Its floors are
\begin{equation*}
\delta_0M\ge32L,\qquad 2LR_1M_1\le M,\qquad \mathsf{o}L\ge2,\qquad L\ge4d,
\end{equation*}
which are among the conditions \eqref{4.7:eq-smallness-conditions-a}. Clause (F1) by itself
applies that lemma to the maximal tower of $X$ only, and so places the pairs in
$U^{c}_{i'}(X,\alpha_{0,i'},\alpha_{1,i'})$; the far clause, applying it at every label
$(p,Z)$, gives membership in the one-power hereditary class.

\emph{Step 6: the bound.} Clause (F1) requires its list of bounds at $\frac13\alpha_2$; this
is the requirement of \cite{B-CMP109} that ``$\varepsilon_1$ is so small that
$\mathbf{H}_j(B(t))$ satisfies (3.14) with $1/3\alpha_2$'', see \cite[(3.14)]{B-CMP109}. With
$C_s$ the constant of \eqref{4.7:eq-shift-fields-a}, $\delta_{i-1}$ the fluctuation size at
level $i-1$ (distinct from the decay rate $\delta_0$ of the operator bounds at free current)
and $\vartheta_i$ the small factor of \eqref{4.7:eq-small-factor-a},
\begin{equation}\label{4.7:eq-simple-share-10}
C_s\delta_{i-1}L^{-n}
\ \le\ \tfrac12\vartheta_i\cdot2(1+n)^{1/2}L^{-n}\,c_Tr_{i'}
\ \le\ \vartheta_i\,c_Tr_{i'}\ \le\ \tfrac13\,c_Tr_{i'} .
\end{equation}
In the first inequality the factor $\frac12\vartheta_i$ is the small factor of
\eqref{4.7:eq-small-factor-a}, and the factor $2(1+n)^{1/2}$ is the first ratio
$\alpha_{\cdot,i-1}/\alpha_{\cdot,i'}\le2(1+n)^{1/2}$ of Lemma~\ref{4.7:lem-radii-comparison},
which passes from level $i-1$ to level $i'$.
The second inequality uses $(1+n)^{1/2}L^{-n}\le1$ from Step 4. The third uses
$\vartheta_i\le\frac13$ from the conditions \eqref{4.7:eq-smallness-conditions-a}. For the
current slot, \eqref{4.7:eq-simple-share-9} and $c_Tr_{i'}\le\frac18\alpha_{0,i'}$, which
follows from the choice of $c_T$ in \eqref{4.7:eq-true-pairs-a} under the conditions
\eqref{4.7:eq-smallness-conditions-a}, give
\begin{equation*}
\mathsf{s}_n\alpha_{0,i}+c_Tr_{i'}
\ \le\ \bigl(\rho_nL^{-3(n+1)}+\tfrac18\bigr)\alpha_{0,i'}\ <\ \alpha_{0,i'} ,
\end{equation*}
since $\rho_nL^{-3(n+1)}\le L^{-2(n+1)}<\frac78$.

Clause (F2) of the operator bounds at free current then gives \cite[(3.17)]{B-CMP109} in its
running form. Each term is at most
\begin{equation*}
S_{i'}\cdot3C_s\delta_{i-1}\,(c_Tr_{i'})^{-1}L^{-n}\,e^{-\frac12c_*\delta_0ML^{n}}
\end{equation*}
times its decay factor. By the first inequality of \eqref{4.7:eq-simple-share-10},
$C_s\delta_{i-1}(c_Tr_{i'})^{-1}\le\vartheta_i(1+n)^{1/2}$, so the term is bounded by
\eqref{4.7:eq-simple-share-3} times its decay factor. For big domains the bound is
\cite[(3.48), p.~280]{B-CMP119}, in the notation of that paper:
\begin{equation*}
\bigl|\mathbf{E}^{(j)}(X,z)\bigr|\ \leqq\
E_0\exp\bigl(-\kappa(L^jL^{-n})^{-1}\bigr)\exp\bigl(-\tfrac12\kappa d_j(X)\bigr).
\end{equation*}
Here $j$, $n$, $\kappa$ and $E_0$ are the letters of \cite{B-CMP119}; in particular this $n$ is
not the level gap $n$ above.
The sums \cite[(1.26)--(1.28)]{B-CMP116} and \cite[(1.31)--(1.32)]{B-CMP116}, through which
these terms are summed, involve no coupling. This proves (b).
\end{proof}

\begin{lemma}[Summation over old levels and the remainder share]\label{4.7:prf-level-summation}
Work at the step $i-1\to i$ in the setting of Lemma~\ref{4.7:lem-potential-splitting}. For an old
level $i'\le i-1$ let $n\ge 0$ be its level gap, as in Lemma~\ref{4.7:lem-radii-comparison}. Let
$\pi_{i'}(Y)$ be the share of the family at level $i'$ defined in \eqref{4.7:eq-unit-share-a}, and
let $\vartheta_i$ be the small factor of \eqref{4.7:eq-small-factor-a}.
\begin{enumerate}
\item[(a)] There is a constant $c_{\pi}''$, having $(6L)^4\tilde{K}^5\bigl(4(1+\beta_0)\bigr)^5$ among
its factors, such that for every old level $i'\le i-1$
\begin{equation}\label{4.7:eq-level-summation-1}
\pi_{i'}(Y)\le c_{\pi}''\,\vartheta_i\,(1+n)^{5/2}L^{-\beta_H n}\,
e^{-(1-2\delta)\kappa d_{i-1}(Y)} .
\end{equation}
\item[(b)] With the constants of \eqref{4.7:eq-level-summation-a} below,
\begin{equation}\label{4.7:eq-level-summation-2}
\sum_{i'\le i-1}\pi_{i'}(Y)\le K_{\varepsilon}\,\vartheta_i\,e^{-(1-2\delta)\kappa d_{i-1}(Y)} .
\end{equation}
\item[(c)] The share of the remainder $R'_{i-1}$ of the regular input obeys the remainder bound of
\eqref{4.7:eq-potential-splitting-a} with the constant $K_R$ of \eqref{4.7:eq-level-summation-a},
and $K_R$ does not depend on $k$.
\end{enumerate}
Here
\begin{equation}\label{4.7:eq-level-summation-a}
\begin{gathered}
K_{\varepsilon}:=c_{\pi}''\,\sigma_{L,\beta_H},\qquad
\sigma_{L,\beta_H}:=\sum_{n\ge 0}(1+n)^{5/2}L^{-\beta_H n},\\
K_R:=\bigl(4(1+\beta_0)^2\tilde{K}\bigr)^{5}\,C_R^{\mathrm{cnt}},
\end{gathered}
\end{equation}
where $C_R^{\mathrm{cnt}}$ denotes the product of the counting and decay constants of Ba{\l}aban's
treatment of the remainder in \cite{B-CMP119}. The constant $K_{\varepsilon}$ is fixed by
$d,L,M,\varkappa_1,\mathfrak{p},\beta_H,\beta_s,\beta_0$, by the constants $B_3$, $O_Q$ and
$O_{(\cdot)}$, which depend on $(d,L)$ only, by the constant $\mathfrak{b}$, which depends on
$(d,L)$ and is fixed after $\beta_s$, and by the constants $c_T$, $c_{\alpha}$, which depend on
$(d,L,M)$ only.
\end{lemma}

\begin{proof}
\emph{(a)} The share $\pi_{i'}(Y)$ is gathered from three contributions: the fundamental-case share
and the counterterm leftovers, bounded in Lemma~\ref{4.7:prf-fundamental-share}, and the
simple-case share, bounded in Lemma~\ref{4.7:prf-simple-share}. Adding these three bounds gives
\eqref{4.7:eq-level-summation-1}, with a constant $c_{\pi}''$ that contains the factor
$(6L)^4\tilde{K}^5\bigl(4(1+\beta_0)\bigr)^5$.

\emph{(b)} Distinct old levels $i'\le i-1$ have distinct gaps $n\ge 0$. Summing
\eqref{4.7:eq-level-summation-1} over $i'\le i-1$ and enlarging the sum to all $n\ge 0$ therefore
gives
\begin{equation*}
\sum_{i'\le i-1}\pi_{i'}(Y)\le c_{\pi}''\,\vartheta_i\,e^{-(1-2\delta)\kappa d_{i-1}(Y)}
\sum_{n\ge 0}(1+n)^{5/2}L^{-\beta_H n}.
\end{equation*}
The series is $\sigma_{L,\beta_H}$, and it is finite for every $L\ge 2$ because $\beta_H>0$: the
ratio of its consecutive terms is $\bigl((2+n)/(1+n)\bigr)^{5/2}L^{-\beta_H}$, which tends to
$L^{-\beta_H}<1$ as $n\to\infty$. With $K_{\varepsilon}=c_{\pi}''\sigma_{L,\beta_H}$ this is
\eqref{4.7:eq-level-summation-2}, which is the bound on the sum of the shares in
\eqref{4.7:eq-potential-splitting-a}. The constants entering $c_{\pi}''$ are those of the three
contributions gathered in (a); they depend only on the quantities listed after
\eqref{4.7:eq-level-summation-a}, and so does $K_{\varepsilon}$.

\emph{(c)} In \cite{B-CMP119} the expression determined by the remainder (written $R_k$ there) is
treated in exactly the same way as the expression determined by the kernels (written $E_k$ there),
except for the second renormalisation, connected with the coupling constant renormalisation
counterterm, which is not needed for the remainder. We follow that treatment. The pieces of the
remainder at level $j\le i-1$ have supremum at most $g_j^{\kappa_0}$ by
\cite[(2.31)]{B-CMP119}, the supremum being taken over the one-power hereditary class
$U^{\Sigma^1}_j$, on which Ba{\l}aban's domains are read throughout this chapter. After the
vacuum-energy renormalisation \cite{B-CMP119} the remainder is marginal: its marginal monomials are
not cancelled, the whole factor $L^{-4n}$ is used to compensate the count of cubes, and the
fluctuation factor is the scale-free small factor $\vartheta_i$.

The constant written $O(1)$ in Ba{\l}aban's bound is a product of $m$ ratios, $m\in\{2,\dots,5\}$,
each of the form
\begin{equation*}
\Theta_{n,j}:=b_i\cdot\max\Bigl\{\frac{4}{\alpha_{3,j}^{\sharp}},\ \frac{1}{c_T r_j}\Bigr\},
\end{equation*}
the analogue at level $j$ of the ratio $\Theta_n$ of \eqref{4.7:eq-total-field-a}. We keep the
ratio of radii exact: for $\alpha_{\cdot,\cdot}$ either of the radii,
\begin{equation*}
\frac{\alpha_{\cdot,i}}{\alpha_{\cdot,j}}
=\Bigl[\frac{\alpha_{\cdot,i}}{\alpha_{\cdot,i-1}}\Bigr]\,
\frac{g_{i-1}}{g_j}\Bigl(\frac{\log x_{i-1}}{\log x_j}\Bigr)^{q}
\le 4(1+\beta_0)\,\frac{g_{i-1}}{g_j},
\end{equation*}
the bracketed factor being at most $2(1+\beta_0)$ by the middle ratio bound of
Lemma~\ref{4.7:lem-radii-comparison}. Using this bound on the ratio of radii in place of the
comparison of radii across the gap that enters the bound on $\Theta_n$ in
Lemma~\ref{4.7:lem-total-field}, we obtain
\begin{equation*}
\Theta_{n,j}\le 4(1+\beta_0)\tilde{K}\,\frac{g_{i-1}}{g_j}.
\end{equation*}
Since $g_{i-1}/g_j\ge(1+\beta_0)^{-1}$ and $m\le 5$, we have
$(g_{i-1}/g_j)^{m}\le(1+\beta_0)^{5-m}(g_{i-1}/g_j)^{5}$, and hence
\begin{equation*}
\Theta_{n,j}^{m}
\le\bigl(4(1+\beta_0)\tilde{K}\bigr)^{m}(1+\beta_0)^{5-m}\Bigl(\frac{g_{i-1}}{g_j}\Bigr)^{5}
\le\bigl(4(1+\beta_0)^2\tilde{K}\bigr)^{5}\Bigl(\frac{g_{i-1}}{g_j}\Bigr)^{5},
\end{equation*}
where in the last step $\tilde{K}\ge 1$ and $\beta_0\ge 0$ are used, so that
$4^m\tilde{K}^m(1+\beta_0)^{5}\le 4^5\tilde{K}^5(1+\beta_0)^{10}$. The irrelevant pieces of the
remainder carry additional powers $L^{-n}\le 1$, which are discarded. Hence the remainder bound of
\eqref{4.7:eq-potential-splitting-a} holds with the constant $K_R$ of
\eqref{4.7:eq-level-summation-a}, the product of $\bigl(4(1+\beta_0)^2\tilde{K}\bigr)^5$ and the
counting and decay constants of \cite{B-CMP119}; none of these factors depends on $k$.

With (b) and (c) the bounds on the sum of the shares and on the remainder share are established,
and the proof of Lemma~\ref{4.7:lem-potential-splitting} is complete.
\end{proof}

\begin{definition}[Smallness of the reference coupling and floors on the constants]
\label{4.7:def-smallness-conditions}
The constants of the step are subject to the conditions displayed below. They are of two kinds.

The \emph{ceilings} $(\Gamma0)$--$(\Gamma6)$ are conditions on the smallness parameter $\gamma$ of the
coupling. Each of them is imposed after $M$ and after every other constant it involves, so that $\gamma$
is the last constant chosen. They are restrictions of the kind that Ba{\l}aban describes as
\emph{stronger restrictions on $g_k$, or $\gamma$}~\cite[p.~355]{B-CMP122b}. The constants of the
coupling flow enter them only through the triple
\begin{equation*}
(\underline{\lambda},\,c_5,\,\bar{B})=(c_0/4,\,c_5,\,C_\beta+1),
\end{equation*}
and none of them involves $\bar{\varsigma}$.

The \emph{floors} $(\varphi1)$--$(\varphi6)$ and $(\varphi\mathrm{B})$ are conditions on $L$, $M$,
$M_1$, $R_1$, $C_0$, $C_1$, $\delta_0$, $\delta_1^{H1}$, $c_{12}$, $\beta_H$ and $B^{109}$. Here $c_{12}$
denotes the constant $\mathsf{o}$ of Definition~\ref{4.4:def-local-data-classes}; both letters are used
below.

In $(\Gamma6)$ and in the condition quoted under $(\varphi1)$, $\alpha_0,\alpha_1,\alpha_2,\alpha_3$
stand for the running radii $\alpha_{0,i},\dots,\alpha_{3,i}$ of
Definitions~\ref{4.7:not-running-couplings} and~\ref{4.7:not-working-radii}, at every level $i$; and $c$
denotes a constant independent of $\gamma$. The exponents $\kappa_0$, $\kappa_1$, the bound $E_0$, the
threshold $\nu_{\mathrm{CL}}$ and the constants $K_{CE}$, $c_Q$, $C_F'$, $\delta_0$, $\delta_1^{H1}$,
$R_1$, $\beta_0$, $\beta_0^H$, $\beta_s$ are constants of the step, fixed in
Definition~\ref{4.7:not-step-constants}, in \eqref{4.7:eq-kernel-power-bound-a} and in the operator
bounds at free current.
\begin{align}
&(\Gamma0)\quad && C_{\mathfrak{R}}\,\gamma^{\kappa_1}\le E_0,
\text{ together with the budget condition described below;}\nonumber\\
&(\Gamma1) && K_{CE}K_{\varepsilon}(1+\beta_0)^{2\kappa_1}\vartheta(\gamma)\le\tfrac12;\nonumber\\
&(\Gamma2) && K_Rc_{\Sigma}'(1+\beta_0)^{\kappa_0-2}\vartheta(\gamma)\gamma^3\le1;\nonumber\\
&(\Gamma3) && \bigl[K_{\pi}+K_{\varepsilon}+K_Rc_{\Sigma}'\gamma^{\kappa_0-2}/E_0\bigr]
E_0\vartheta(\gamma)\le\tfrac12\nu_{\mathrm{CL}};\nonumber\\
&(\Gamma4) && \text{(a) }c_5\gamma^2\le\beta_0,\qquad\text{(b) }(C_\beta+1)\gamma^2\le\beta_0,
\nonumber\\
& && \text{(c) }\Bigl(1-\frac{2c_5\gamma^2}{\log\gamma^{-2}}\Bigr)^{-q}\le2,\qquad
\text{(d) }\Bigl(1+\frac{c_5\gamma^2}{\log\gamma^{-2}}\Bigr)^{q}\le2,\nonumber\\
& && \text{(e) }\gamma\le e^{-q},\qquad
\text{(f) }\Bigl(1+\frac{(C_\beta+1)\gamma^2}{\log\gamma^{-2}}\Bigr)^{q}\le2;\nonumber\\
&(\Gamma5) && 3\vartheta(\gamma)\le1;\nonumber
\displaybreak[0]\\
&(\Gamma6) && \text{(i) the smallness of }\alpha_0,\alpha_1,\alpha_2,\alpha_3
\text{ required in \cite[Lemma~4]{B-CMP109};}\nonumber\\
& && \text{(ii) }\exp\Bigl(28d\,\tfrac{M}{L}(1+2\beta_s)\bigl(\alpha_{1,i}
+c_{12}LMB^{109}\alpha_{0,i}\bigr)\nonumber\\
& && \qquad\quad+4dLc_Q\mathfrak{b}M\alpha_{0,i}
+cO_{(3.37)}B_3\mathfrak{b}M\alpha_{0,i}\Bigr)(1+2\beta_s)\le1+3\beta_s;\nonumber\\
& && \text{(iii) the ceilings of the operator bounds at free current,}\nonumber\\
& && \qquad\quad\text{including } c_T\,r_i\le1/(4C_F');\nonumber\\
& && \text{(iv) }8\beta_sL^{-2}\alpha_0\le1,\qquad b_i\le1;\nonumber\\
& && \text{(v) }(4LO_{(3.37)}B_3\mathfrak{b}M)^2\alpha_0\le\beta_s;\nonumber\\
& && \text{(vi) }172dL^2\bigl[(4(1+\beta_0)O_{(3.37)}B_3\mathfrak{b}M)^2+1\bigr]\alpha_0
\le\beta_s;\nonumber\\
& && \text{(vii) }\bar{\theta}\,e^{c(O_{(3.37)}B_3\mathfrak{b}M\alpha_0+\beta_s\alpha_0)}\le1,
\qquad\bar{\theta}_J\,e^{c(O_{(3.37)}B_3\mathfrak{b}M\alpha_0+\beta_s\alpha_0)}\le1;\nonumber\\
& && \text{(viii) the exponential factor in \cite[(3.43)]{B-CMP109} is at most }2;\nonumber
\displaybreak[0]\\
&(\varphi1) && C_1>8(1+\beta_0)O_{(3.37)}B_3\mathfrak{b}M\cdot C_0,\qquad C_0\le C_1;
\nonumber\\
&(\varphi2) && L^2\ge4(1+\beta_0)(1+9\beta_s);\nonumber\\
&(\varphi3) && \beta_H\in\,]0,\beta_0^H],\qquad\beta_0^H<1;\nonumber\\
&(\varphi4) && L\ge4(1+\beta_0)(1+2\beta_s);\nonumber\\
&(\varphi5) && 17M_1\le LM,\quad17M_1\le c_{12}LM,\quad c_{12}L\ge2,\quad\delta_0M\ge32,
\quad2R_1M_1\le M,\nonumber\\
& && \delta_1^{H1}M\ge\varkappa_1,\quad L\ge27\text{ odd},\quad M>L^2M_1,\quad
L^2\ge12/c_{12},\quad L\ge4d,\quad R_1\ge L,\nonumber\\
& && \text{together with the further floors listed below;}\nonumber\\
&(\varphi6) && \text{(i) }\delta_0M\ge32L,\qquad\text{(ii) }2LR_1M_1\le M;\nonumber\\
&(\varphi\mathrm{B}) && B^{109}\ge\frac{4(1+\beta_0)O_{(3.37)}B_3\mathfrak{b}+1}{\mathsf{o}L}.
\label{4.7:eq-smallness-conditions-a}
\end{align}
The constants and the roles of these conditions are as follows.
\begin{itemize}
\item[$(\Gamma0)$--$(\Gamma3)$] Here $\vartheta(\gamma)$ is the small factor of
\eqref{4.7:eq-small-factor-a}. The constants $K_{\varepsilon}=c_{\pi}''\sigma_{L,\beta_H}$ and $K_R$
are those of \eqref{4.7:eq-level-summation-a}, and $C_{\mathfrak{R}}=2K_{CE}$ and $c_{\Sigma}'$ are
those of \eqref{4.7:eq-kernel-power-bound-a}. Condition $(\Gamma0)$ comprises, besides the displayed
inequality, the budget condition stated in clause (iii) of the lemma that fixes the running form of the
constants of \cite{B-CMP109}. The constant $K_\pi$ in $(\Gamma3)$ is the share
constant of the unmarked kernel families $\mathfrak{E}_j$. Condition $(\Gamma3)$ consists of the
thresholds of the cluster expansion: it ensures that the unmarked total and the marked total are each
at most $\nu_{\mathrm{CL}}/2$. These are Ba{\l}aban's thresholds $\varepsilon_2$ of the cluster
expansion with $E_0$ replaced by $2E_0$.
\item[$(\Gamma4)$] Condition (e) is a ceiling on $\gamma$ that depends on $q$; it is not a bound
on $q$.
\item[$(\Gamma6)$] Member (ii) is the rung of \cite[(3.41)]{B-CMP109} at the running radii. Its
exponent is formed from the constants supplied by the first and the fourth clauses of the axial-frame
lemma of the operator bounds at free current.
In member (iii), $r_i$ is the radius fixed in Definition~\ref{4.7:not-running-couplings}, and $c_T$ is
the constant that multiplies it in the product $c_Tr_{i'}$ of Lemma~\ref{4.7:lem-total-field}.
Member (v) is the condition of \eqref{4.7:eq-lemma-four-gaps-a}. The slacks $\bar\theta$,
$\bar\theta_J$ in (vii) are those of \eqref{4.7:eq-lemma-four-gaps-a}. The constant $b_i$ in (iv) is
that of Definition~\ref{4.7:not-working-radii}. The constants $\mathfrak{b}$ and $O_{(3.37)}$ are those of
\eqref{4.7:eq-step-constants-a}. Members (ii), (v), (vi) and (vii) are the conditions read in the proof
of Ba{\l}aban's Lemma~4 at every level gap (Lemma~\ref{4.7:prf-lemma-four-gaps}); member (iii) is
imposed by the operator bounds at free current. They involve $M/L$, $B^{109}$, $c_{12}$,
$c_Q$, $\mathfrak{b}$, $O_{(3.37)}$, $B_3$, $C_0$ and $C_1$, all of which are chosen before
$\gamma$. At the running radii each member of $(\Gamma6)$ has the form
$C(d,L,M,\dots)\,\gamma(\log\gamma^{-2})^q\le c$.
\item[$(\varphi1)$] This condition is strict. It is Ba{\l}aban's condition
$B_3^2O(1)M\alpha_0<\frac12\alpha_1$ of \cite{B-CMP109}, applied to the frame clause of the
operator bounds at free current in its localised form, with the far data sized by the axial-frame
lemma. The
$O(1)$ is the constant of \cite[(3.37)]{B-CMP109}, that is, $O_{(3.37)}$ of
\eqref{4.7:eq-step-constants-a}, which dominates the constant $O_Q$ of the bound on $\|Q_j\|$.
\item[$(\varphi5)$] These are the floors of the operator bounds at free current, all of which are
adopted here. Besides the members displayed, they comprise:
\begin{itemize}
\item the floor of that theorem that involves $R_1$ and is imposed together with
\cite[(2.59)]{B-CMP96} after $L$;
\item the rim floor on $M_1$;
\item a floor on $M$ making the gradients of Ba{\l}aban's partitions of unity $\{\zeta_\square\}$ and
$\{\tilde\zeta_\square\}$ of \cite{B-CMP119} at most $1$.
\end{itemize}
The last of these floors is used by the constant $C_N$ of the linear clause of the axial-frame lemma,
and hence by $C_s$.
\item[$(\varphi6)$] These are the floors of the one-power hereditary class at the cube side $M/L$:
they are the members $\delta_0M\ge32$ and $2R_1M_1\le M$ of $(\varphi5)$ with $M$ replaced by $M/L$.
They are used
in that class's construction and in its top layer at cube side $M/L$. The
conditions $c_{12}L\ge2$ and $L\ge4d$ needed there are already members of $(\varphi5)$.
\item[$(\varphi\mathrm{B})$] This floor on $B^{109}$ is imposed, but no argument of this chapter uses
it when the group-valued factor of Definition~\ref{4.4:def-group-factor-bounds} equals $1$.
\end{itemize}
All other constants are fixed by the definitions referred to. There is no upper bound on $q$, and the
constant $L_0$ is not evaluated. None of the conditions is a cap, in the sense of
Definition~\ref{2.3:def-cap}, on $K$, on the number of steps, on $L$, on $M$ or on the coupling.
\end{definition}

\begin{proof}
We verify the implications recorded above.

\emph{Step 1: $(\Gamma4)$(e) bounds the radii.}
For $q>0$ the function $x\mapsto(\log x)^qx^{-1/2}$ has, at $x>1$, the derivative
\begin{equation*}
x^{-3/2}(\log x)^{q-1}\bigl(q-\tfrac12\log x\bigr),
\end{equation*}
which is $\le0$ when $\log x\ge2q$. Condition (e) says $\log\gamma^{-2}\ge2q$. Hence
$x\mapsto(\log x)^qx^{-1/2}$ decreases on $x\ge\gamma^{-2}$. Each radius $\alpha_{\cdot,m}$ of
Definition~\ref{4.7:not-running-couplings} is $C_\cdot$ times the value of this function at a point
$x\ge\gamma^{-2}$. Consequently
\begin{equation*}
\alpha_{\cdot,m}\le C_\cdot(\log\gamma^{-2})^q\gamma .
\end{equation*}

\emph{Step 2: $(\Gamma4)$(c), (d), (f) follow from (a), (b), (e).}
Here $\beta_0\le\frac12$, as fixed in the order of choice of the parameters (this is also the value
used in Step~7). The derivation of (c) below uses in addition $q\ge1$; for $q<1$ condition (c) is
imposed as displayed.
By (a) and $\log\gamma^{-2}\ge2q$,
\begin{equation}\label{4.7:eq-smallness-conditions-1}
\frac{c_5\gamma^2}{\log\gamma^{-2}}\le\frac{\beta_0}{2q},
\end{equation}
and by (b) the same bound holds with $c_5$ replaced by $C_\beta+1$.

For (d), the inequality $(1+u)^q\le e^{qu}$ and \eqref{4.7:eq-smallness-conditions-1} give a left
member at most $e^{\beta_0/2}\le e^{1/4}<2$. Condition (f) follows in the same way from (b).

For (c), put $u=2c_5\gamma^2/\log\gamma^{-2}$. By \eqref{4.7:eq-smallness-conditions-1},
$u\le\beta_0/q$, where $\beta_0\le\frac12$. For $0<\beta_0<q$,
\begin{equation*}
-q\log(1-\beta_0/q)=\beta_0\,\psi(\beta_0/q),\qquad
\psi(y)=-\frac{\log(1-y)}{y}=\sum_{j\ge0}\frac{y^j}{j+1}.
\end{equation*}
Since $\psi$ increases in $y$, the quantity $-q\log(1-\beta_0/q)$ decreases in $q$. Therefore
\begin{equation*}
(1-u)^{-q}\le(1-\beta_0/q)^{-q}\le(1-\beta_0)^{-1}\le2 .
\end{equation*}

\emph{Step 3: member (vi) of $(\Gamma6)$ gives the bound on $86d(h_n+h')^2$ used for the running
current factor $\theta_J$ of \eqref{4.7:eq-lemma-four-gaps-a}, where $h_n$ and $h'$ are as there.}
Let $\rho_n=4(1+\beta_0)(1+n)^{1/2}$ be as in Lemma~\ref{4.7:lem-radii-comparison}. For $n\ge0$ and
$L\ge2$ we have $1+n\le2^n\le L^{2n}$, so $\rho_nL^{-n}\le4(1+\beta_0)$. By this, the quantity $h_n$
of \eqref{4.7:eq-lemma-four-gaps-a} satisfies
\begin{equation*}
h_n\le4(1+\beta_0)O_{(3.37)}B_3\mathfrak{b}M\cdot\alpha_{0,i'} ,
\end{equation*}
and there also $h'\le\alpha_{0,i'}$. Using $(x+y)^2\le2x^2+2y^2$ we obtain
\begin{align*}
86d(h_n+h')^2&\le172d\bigl[(4(1+\beta_0)O_{(3.37)}B_3\mathfrak{b}M)^2+1\bigr]\alpha_{0,i'}^2\\
&\le\beta_sL^{-2}\alpha_{0,i'} ,
\end{align*}
where the last inequality is member (vi) of $(\Gamma6)$ at $\alpha_0=\alpha_{0,i'}$, multiplied by
$L^{-2}\alpha_{0,i'}$. The resulting condition does not depend on $n$.

\emph{Step 4: $(\varphi4)$ contains $L\ge2(1+\beta_s)$.}
Since $\beta_0,\beta_s\ge0$, we have $4(1+\beta_0)(1+2\beta_s)\ge4(1+\beta_s)\ge2(1+\beta_s)$.
Hence $(\varphi4)$ implies the condition $L\ge2(1+\beta_s)$ of the one-power hereditary class at the
running radii.

\emph{Step 5: the first members of $(\varphi5)$.}
Suppose $M>L^2M_1$, $L\ge17$ (which follows from $L\ge27$) and $c_{12}\ge1$. Then
\begin{equation*}
LM>L^3M_1\ge17M_1\qquad\text{and}\qquad c_{12}LM\ge LM .
\end{equation*}
So $17M_1\le LM$ and $17M_1\le c_{12}LM$ follow.

\emph{Step 6: the window tower.}
The construction of the window tower implicitly requires $M\ge O(1)R_1M_1L$, so that the
$M$-cubes of $\tilde{\square}^{4}$ belong to the one-power hereditary class. This requirement follows
from $(\varphi6)$(ii): the collar of that class has width at most $R_1M_1$, and $(\varphi6)$(ii) gives
$R_1M_1\le\frac12M/L$. Thus the requirement holds with the explicit constant $O(1)=2$.

\emph{Step 7: compatibility of $(\varphi2)$ and $(\varphi4)$ with $(\varphi5)$.}
Suppose $\beta_0\le\frac12$ and $\beta_s=\frac14$. The right members of $(\varphi2)$ and $(\varphi4)$
satisfy
\begin{equation*}
4(1+\beta_0)(1+9\beta_s)\le4\cdot\tfrac32\cdot\tfrac{13}4=\tfrac{39}2,\qquad
4(1+\beta_0)(1+2\beta_s)\le4\cdot\tfrac32\cdot\tfrac32=9 .
\end{equation*}
Both conditions are therefore implied by the requirement $L>11$ of Ba{\l}aban's construction, as
fixed in the order of choice of the parameters, and by the member $L\ge27$
of $(\varphi5)$. This conclusion concerns $(\varphi2)$ and $(\varphi4)$ only.
\end{proof}

\begin{corollary}[Power bound for large-field-descended kernels at every level]
\label{4.7:cor-kernel-power-bound}
Fix $k$. Assume the following.
\begin{itemize}
\item The standing hypotheses of the linear splitting of the potentials by family
(Lemma~\ref{4.7:lem-potential-splitting}) on the step and its regular input, and the conditions
$(\Gamma0)$--$(\Gamma6)$ of Definition~\ref{4.7:def-smallness-conditions}, displayed in
\eqref{4.7:eq-smallness-conditions-a}. The conditions on the kernel families entering that
input are checked in the proof below.
\item The hypothesis $\mathfrak{i}(k)$ of the induction along the coupling flow at step $k$,
together with $x_0,\dots,x_k\ge\gamma^{-2}$.
\item The exponent $\kappa_0$ satisfies $\kappa_0\ge17$.
\item The cluster-expansion bound for the marked sub-series, with its constant
$K_{CE}=4b^{\mathrm{adm}}/\nu_{CL}$, where $b^{\mathrm{adm}}$ and $\nu_{CL}$ are the two
constants of that bound, in the normalisation by the prefactor of the marked potential.
\item Lemma~\ref{4.7:lem-potential-splitting} itself.
\item Part (a) of the theorem on marked kernel families (whose notation is that of
Notation~\ref{4.7:not-kernel-families}), and the rule attaching to each marked family
its own renormalisation coefficient.
\item Part (b) of the lemma on sums of powers of the couplings along the flow, together with
the lower flow bound of the induction hypothesis at step $k$.
\item The volume lemma, through its volume-cancellation property at every index at most $k$.
\item The independence of the localised terms from the small-field region
(Corollary~\ref{4.4:cor-localised-independence}).
\item The restriction inclusion of \eqref{4.7:eq-class-restrictions-a}
(Notation~\ref{4.7:not-class-restrictions}), and the form of
\cite[Lemma~1, (1.34)]{B-CMP116} valid on the one-power hereditary class $U^{\Sigma^1}$.
\end{itemize}
Set $\kappa_1=\kappa_0-5$, and define the constants $C_{\mathfrak{R}}$ and $c_{\Sigma}'$ by the
first line of \eqref{4.7:eq-kernel-power-bound-a}. Then, for every $i\le k+1$, every
$X\in\mathbf{D}_i$ and every pinning point $z$, the bound in the second line of
\eqref{4.7:eq-kernel-power-bound-a} holds; we refer to it as the bound \textup{(E)}:
\begin{equation}\label{4.7:eq-kernel-power-bound-a}
\begin{gathered}
C_{\mathfrak{R}}:=2K_{CE},\qquad
c_{\Sigma}':=1+\frac{8(1+\beta_0)^{\kappa_0-7}}{c_0(\kappa_0-7)},\\
|\mathbf{E}^{(i),\mathfrak{R}}(X,\cdot)|,\ |\mathbf{E}^{(i),\mathfrak{R}}(X,\cdot,z)|
\le C_{\mathfrak{R}}\,g_i^{\kappa_1}e^{-\kappa d_i(X)}
\quad\text{on } U^{\Sigma^1}_i(X;\alpha_{0,i},\alpha_{1,i}).
\end{gathered}
\end{equation}
The kernels are analytic on $U^{\Sigma^1}_i(X;\alpha_{0,i},\alpha_{1,i})$ and invariant under
$G^{c}$-gauge transformations. For $i\le k$ the couplings and radii in
\eqref{4.7:eq-kernel-power-bound-a} are the true ones. For $i=k+1$ the coupling $g_{k+1}$,
both in the bound and in the radii $\alpha_{0,k+1},\alpha_{1,k+1}$, is the stand-in
$\tilde{g}:=g_k$; that is, the class is taken at the stand-in coupling. The same constants serve
at every site $\mathfrak{s}\in\mathcal{S}_i^{\flat}$.
\end{corollary}

\begin{proof}
We argue by induction on $i$. For $i=1$ there are no large-field-descended kernels:
$\mathbf{E}^{(1),\mathfrak{R}}=0$ by the theorem on marked kernel families, and there is nothing
to prove.

Let $2\le i\le k+1$, and assume the induction hypothesis, which we write $(\mathrm{E})_{<i}$:
the bound \textup{(E)} of \eqref{4.7:eq-kernel-power-bound-a}, with the analyticity and
invariance asserted with it, holds at every $i'\le i-1$ with the constants of the statement.

\emph{The kernels at level $i$.} The kernels $\mathbf{E}^{(i)}(X,\cdot,z)$ with $X\subset\Lambda_i$
are the whole-lattice expressions. In the words of \cite[p.~278]{B-CMP119}, ``the terms
$E_0^{(k+1)}(\Lambda_{k+1},X,z)$ of this representation, for $X\subset\Lambda_{k+1}$, do not
depend on $\Lambda_{k+1}$, and coincide with the corresponding terms arising from the
expressions defined on the whole lattice''. This is the content of
Corollary~\ref{4.4:cor-localised-independence}. The kernels are produced by a single
application of the step operation $\mathbf{T}$ to the regular input of
Lemma~\ref{4.7:lem-potential-splitting}. No vertex processed by boundary terms, by $R''$ or by
the large-field operation $\mathbf{R}$ occurs in them, because those terms are collected in the
boundary terms of that representation, written $B^{(k+1)}$ in \cite{B-CMP119}.

\emph{The input families.} For $i'\le i-1$ we take as marked family
$\mathfrak{M}_{i'}:=\mathbf{E}^{(i'),\mathfrak{R}}$ and as unmarked family
$\mathfrak{E}_{i'}:=\mathbf{E}^{(i'),\mathfrak{A}}$. We now verify the conditions of
Lemma~\ref{4.7:lem-potential-splitting} on these families.
\begin{itemize}
\item By $(\mathrm{E})_{<i}$ and $(\Gamma0)$, the decay-normalised supremum of
$\mathfrak{M}_{i'}$ satisfies $S_{i'}\le C_{\mathfrak{R}}g_{i'}^{\kappa_1}\le E_0$.
\item Each $\mathfrak{M}_{i'}$ carries its own renormalisation coefficient
$\beta^{\mathfrak{R}}_{i'}$, by the rule attaching to each marked family its own coefficient.
\item The volume-cancellation property of $\mathfrak{M}_{i'}$ holds by the volume lemma at
index $i'$.
\item The fresh remainder $R'_{i-1}$ satisfies the bound \cite[(2.31)]{B-CMP119}, which is part
of the induction hypothesis $(\gamma_{i-1})$.
\item The unmarked families satisfy $S(\mathfrak{E}_{i'})\le E_0$ by the bound on the regular
kernel families.
\end{itemize}
By the definition of the marking in the theorem on marked kernel families,
$\mathbf{E}^{(i),\mathfrak{R}}(X)=\mathfrak{t}^{\mathfrak{R}}(X)$. Here
$\mathfrak{t}^{\mathfrak{R}}(X)$, in the notation of Notation~\ref{4.7:not-kernel-families}, is
the sub-series of the step's exponentiated cluster expansion formed by the terms containing at
least one marked activity.

\emph{First layer: the cluster expansion.} Let $v^{\mathrm{m}}$ denote the supremum, over
domains $Y$, of the decay-normalised prefactor of the marked potential. We apply the
cluster-expansion bound for the marked sub-series under $(\Gamma3)$. For the pinned sub-series
this bound carries one additional factor of degree one. It gives
\begin{equation*}
|\mathfrak{t}^{\mathfrak{R}}(X)|\le K_{CE}\,v^{\mathrm{m}}e^{-\kappa_{\mathrm{cl}}d_i(X)},
\end{equation*}
where $\kappa_{\mathrm{cl}}$ is the decay rate of the cluster expansion. By
\cite[(2.41)]{B-CMP116} it satisfies $\kappa_{\mathrm{cl}}\ge\kappa$, so that
$e^{-\kappa_{\mathrm{cl}}d_i(X)}\le e^{-\kappa d_i(X)}$. The constant $K_{CE}$ depends neither
on $k$ nor on the couplings.

\emph{Second layer: the splitting of the potentials.} We apply
Lemma~\ref{4.7:lem-potential-splitting} with the bound
$S_{i'}\le(1+\beta_0)^{\kappa_1}C_{\mathfrak{R}}g_{i-1}^{\kappa_1}$, obtained from
$(\mathrm{E})_{<i}$ and the comparison $g_{i'}\le(1+\beta_0)g_{i-1}$ for $i'\le i-1$ of
Lemma~\ref{4.7:lem-radii-comparison}. With $\vartheta_i$ the small factor of
\eqref{4.7:eq-small-factor-a}, and $K_\varepsilon$, $K_R$ the constants of
\eqref{4.7:eq-level-summation-a}, the lemma gives
\begin{equation*}
v^{\mathrm{m}}\le(1+\beta_0)^{\kappa_1}C_{\mathfrak{R}}g_{i-1}^{\kappa_1}K_\varepsilon\vartheta_i
+K_R\vartheta_i\,g_{i-1}^{5}\sum_{j\le k_1}g_j^{\kappa_0-5}.
\end{equation*}
Here $k_1$ is the index of Lemma~\ref{4.7:lem-potential-splitting}, and $k_1\le i-1$.

\emph{The sum along the flow.} We use part (b) of the lemma on sums of powers of the couplings,
together with the lower flow bound, whose constant shift $c_5$ we carry: for levels $j<l$,
\begin{equation*}
x_j\ge x_l-c_5+\lambda_\sharp(l-j).
\end{equation*}
Let $\varpi>2$. The function $t\mapsto(x_l-c_5+\lambda_\sharp t)^{-\varpi/2}$ is decreasing on
$[0,\infty)$. Hence each summand with $l-j=t\ge1$ is at most the integral of this function over
$[t-1,t]$. Since $g_j^{\varpi}=x_j^{-\varpi/2}$, this gives
\begin{align*}
\sum_{j<l}g_j^{\varpi}
&\le\int_0^\infty(x_l-c_5+\lambda_\sharp t)^{-\varpi/2}\,dt
=\frac{2\,(x_l-c_5)^{-(\varpi-2)/2}}{\lambda_\sharp(\varpi-2)}\\
&\le\frac{2(1+\beta_0)^{\varpi-2}g_l^{\varpi-2}}{\lambda_\sharp(\varpi-2)}.
\end{align*}
We take $\varpi=\kappa_0-5$, which is at least $12$ by the hypothesis $\kappa_0\ge17$, and
$l=i-1\ge k_1$, and use $\lambda_\sharp\ge c_0/4$. The sum
over $j\le k_1$ is at most the sum over $j\le i-1$. The summand $j=i-1$ contributes
$g_{i-1}^{5}g_{i-1}^{\kappa_0-5}=g_{i-1}^{\kappa_0}\le g_{i-1}^{\kappa_0-2}$, since
$g_{i-1}\le\gamma\le1$. The summands
$j<i-1$ contribute at most
\begin{equation*}
g_{i-1}^{5}\cdot\frac{8(1+\beta_0)^{\kappa_0-7}g_{i-1}^{\kappa_0-7}}{c_0(\kappa_0-7)}
=\frac{8(1+\beta_0)^{\kappa_0-7}}{c_0(\kappa_0-7)}\,g_{i-1}^{\kappa_0-2}.
\end{equation*}
By the definition of $c_\Sigma'$ in \eqref{4.7:eq-kernel-power-bound-a}, therefore,
\begin{equation*}
g_{i-1}^{5}\sum_{j\le i-1}g_j^{\kappa_0-5}\le c_\Sigma'\,g_{i-1}^{\kappa_0-2}.
\end{equation*}
We insert this into the second layer and the result into the first layer, and obtain
\begin{equation}\label{4.7:eq-kernel-power-bound-1}
|\mathbf{E}^{(i),\mathfrak{R}}(X)|\le K_{CE}\Bigl[K_\varepsilon(1+\beta_0)^{\kappa_1}\vartheta_i
C_{\mathfrak{R}}g_{i-1}^{\kappa_1}+K_Rc_\Sigma'\vartheta_i\,g_{i-1}^{\kappa_0-2}\Bigr]
e^{-\kappa d_i(X)}.
\end{equation}

\emph{Closing the bound.} For $i\le k$ we use $g_{i-1}\le(1+\beta_0)g_i$. For $i=k+1$ we have
$g_{i-1}=\tilde{g}$, which is the coupling standing at index $i$. We also use
$\kappa_0-2=\kappa_1+3$, and $g_i\le\gamma$; the latter holds because $x_0,\dots,x_k\ge\gamma^{-2}$,
and at $i=k+1$ because $g_i=\tilde{g}=g_k$. Then the right side of
\eqref{4.7:eq-kernel-power-bound-1} is at most
\begin{align*}
&\bigl[K_{CE}K_\varepsilon(1+\beta_0)^{2\kappa_1}\vartheta_i\bigr]C_{\mathfrak{R}}g_i^{\kappa_1}
e^{-\kappa d_i(X)}\\
&\qquad+K_{CE}\bigl[K_Rc_\Sigma'(1+\beta_0)^{\kappa_0-2}\vartheta_i\,g_i^{3}\bigr]g_i^{\kappa_1}
e^{-\kappa d_i(X)}\\
&\qquad\le\bigl(\tfrac12C_{\mathfrak{R}}+K_{CE}\bigr)g_i^{\kappa_1}e^{-\kappa d_i(X)}
=C_{\mathfrak{R}}\,g_i^{\kappa_1}e^{-\kappa d_i(X)}.
\end{align*}
In this chain, by $(\Gamma1)$ and $(\Gamma2)$, with $g_i\le\gamma$, the two brackets are at
most $\tfrac12$ and $1$ respectively. The final equality is the definition
$C_{\mathfrak{R}}=2K_{CE}$.

\emph{Domain.} By the two statements of this chapter on the potentials with the source current
as a slot variable, and by the form of \cite[Lemma~1, (1.34)]{B-CMP116} valid on the one-power
hereditary class, the potentials are analytic on the enlarged class
\begin{equation*}
U^{\Sigma^1}_i(Y;(1+\beta_s)\alpha_{0,i},(1+\beta_s)\alpha_{1,i},\alpha_{0,i}).
\end{equation*}
In the discussion of the activities of
\cite[(2.13)]{B-CMP116} the potentials are shown to be analytic on $U^c_{k+1}(X,\alpha_0,\alpha_1)$.
It follows there that the activities of \cite[(2.13)]{B-CMP116}, and the whole sum
$\mathbf{E}^{(k+1)}(X)$, are analytic in $(\mathbf{U},\mathbf{J})$ on $U^c_{k+1}(X,\alpha_0,\alpha_1)$.
This holds under the hypothesis printed there, ``Of course $Y\subset Z_0$, and
$\tilde{Z}_0\subset Z\subset X$ for the activities in (2.13)'', in which $Z_0$, $\tilde{Z}_0$
and $Z$ are the regions attached in that paper to each activity of \cite[(2.13)]{B-CMP116}.

The passage from $X$ to $Y$ is the inclusion \eqref{4.7:eq-class-restrictions-a}:
for $Y\in\mathbf{D}_{i-1}$ and $Y\subset X\in\mathbf{D}_i$,
\begin{equation*}
U^{\Sigma^1}_i(X;\alpha_{0,i},\alpha_{1,i})\big|_{Y}\subset
U^{\Sigma^1}_i(Y;(1+\beta_s)\alpha_{0,i},(1+\beta_s)\alpha_{1,i},\alpha_{0,i}).
\end{equation*}
It is applied with the parameter $\mu=1$ of \eqref{4.7:eq-class-restrictions-a}.
For this inclusion the label sets satisfy $\mathfrak{T}_i(Y)\subset\mathfrak{T}_i(X)$, and no
further condition on $Y$ or $X$ is needed. The quadratic forms and covariances of the step use
only the whole-lattice conditions (i)--(iii) of the class, as in the discussion of
\cite[(2.13)]{B-CMP116}. Hence every
activity of the step, and therefore $\mathfrak{t}^{\mathfrak{R}}(X)=\mathbf{E}^{(i),\mathfrak{R}}(X)$
together with its pinned version, is analytic on $U^{\Sigma^1}_i(X;\alpha_{0,i},\alpha_{1,i})$. At
$i=k+1$ every quantity carrying the index $i$ is taken at the stand-in coupling $\tilde{g}$. Nothing
is asserted about $U^{\Sigma^1}_{k+1}$ at the true coupling $g_{k+1}$. Once $g_{k+1}$ is
available, the induction along the coupling flow applies the present corollary again at $k+1$,
with the true couplings.

\emph{Invariance.} The kernels are $G^{c}$-gauge invariant by part (a) of the theorem on marked
kernel families. The class $U^{\Sigma^1}$ is a union of $G^{c}$-orbits, so invariance holds on
the whole of $U^{\Sigma^1}_i(X;\alpha_{0,i},\alpha_{1,i})$.

\emph{Extension.} Away from true pairs, the far pieces are extended through the far
representation $\Phi^{\mathbb{K}}$ of the operator bounds at free current. Since
$\Phi^{\mathbb{K}}$ is a representation only, the values at true pairs do not change. Hence
the marked quantity $\Pi^{\mathfrak{R}}$ and the renormalisation coefficients
$\beta^{\mathfrak{R}}$ of the marked families, both evaluated at
$(\mathbf{U},\mathbf{J})=(1,0)$, do not change either.

\emph{Sites and pinning.} Every constant used above depends only on
$(d,L,M,\kappa,\varkappa_1,\mathfrak{p},\beta_H,\beta_s,\beta_0)$. At a torus site the same
argument applies, with the volume-cancellation property at indices below $i$ taken in the form in
which it extends to the torus and cancels there; this is part (c) of the hypothesis $(H_i)$ of
the volume lemma. The pinned kernels
$\mathbf{E}^{(i),\mathfrak{R}}(X,\cdot,z)$ differ from the unpinned ones by one factor of degree
one, in the first layer and in the second. So the bound above holds for them with the same
constants.

\emph{No compounding.} At each step the suprema $S_{i'}$ are taken afresh from
$(\mathrm{E})_{<i}$, with the same constant $C_{\mathfrak{R}}$. The level $k$ enters only
through the couplings $g_j$. Hence the constants do not grow along the induction, and the bound
holds at every level $i\le k+1$.
\end{proof}

\begin{proposition}[The large-field part of the beta coefficient by Cauchy extraction]
\label{4.7:prop-beta-extraction}
Fix $k$ such that the induction hypothesis $\mathfrak{i}(k)$ holds and
$x_0,\dots,x_k\ge\gamma^{-2}$, where $\gamma\le 1$. Assume the hypotheses of
Corollary~\ref{4.7:cor-kernel-power-bound}, the conditions $(\Gamma0)$--$(\Gamma6)$ of
Definition~\ref{4.7:def-smallness-conditions}, and $\kappa_0\ge 17$. Put
$\kappa_1=\kappa_0-5$ and $C_{\mathfrak{R}}:=2K_{CE}$. Assume further the following
statements.
\begin{itemize}
\item[(A1)] The power bound for large-field-descended kernels at every level, that is, the
bound $(E)$ of Corollary~\ref{4.7:cor-kernel-power-bound}, display
\eqref{4.7:eq-kernel-power-bound-a}. At $i=k+1$, $g_{k+1}$ is replaced, in the bound and in
the radii, by the stand-in $\tilde{g}=g_k$.
\item[(A2)] True pairs lie in the regularity class, as stated in
Definition~\ref{4.7:not-true-pairs}.
Here a true pair is a pair consisting of a configuration and its current. The constant
$c_T$ and the radius $r_j$ of the polydisc norm $N_j$ are those of display
\eqref{4.7:eq-true-pairs-a}.
\item[(A3)] The decay bound \cite[Proposition~9, (190)]{B-CMP102b}, taken at the real
background $1$, together with the remark in \cite{B-CMP109} that the same estimate holds
for the derivative $H_j=D\mathbf{H}_j(0)$ of the nonlinear potential function
$\mathbf{H}_j$ of that paper and for the second-order operators applied to
$\mathbf{H}_j$: for every bond $b$ and every $X$,
\begin{equation}\label{4.7:eq-beta-extraction-3}
N_j(H_j\delta_b\lceil X)\le C_{\mathsf{h}}e^{-\delta_0\operatorname{dist}(b,X)},
\end{equation}
where $\delta_b$ is the unit element at the bond $b$ and $\lceil X$ denotes restriction
to $X$.
\item[(A4)] The second-moment bound for whole-lattice sums, which evaluates the coefficient
of \cite[(5.42)]{B-CMP109} (see also \cite[(3.62)]{B-CMP119}). Let $C_A$ be a constant with
\begin{equation*}
\sum_X\bigl(Md_j(X)+c\bigr)^2e^{-\kappa d_j(X)}\le C_A ,
\end{equation*}
where $c$ is the constant of that second-moment formula. There is a constant $C_{\Pi}$,
containing the factor $4C_{\mathsf{h}}^2C_A$, such that
the following holds. If a whole-lattice sum $\Pi(b,b')=\sum_X\Pi_X(b,b')$ at level $j$ has,
for some $S_j>0$ and for all $X$, $b$, $b'$, summands bounded by
\begin{equation*}
|\Pi_X(b,b')|\le\frac{4C_{\mathsf{h}}^2S_j}{(c_Tr_j)^2}\,
e^{-\kappa d_j(X)-\delta_0[\operatorname{dist}(b,X)+\operatorname{dist}(b',X)]},
\end{equation*}
then the coefficient $\beta$ that \cite[(5.42)]{B-CMP109} assigns to $\Pi$ satisfies
$|\beta|\le C_{\Pi}S_j(c_Tr_j)^{-2}$.
\end{itemize}
Let $C'$ be the constant defined in \eqref{4.7:eq-beta-extraction-a} below, and let
$\bar{\varsigma}:=C'\gamma^{\kappa_0-8}$ be the ceiling for $|\beta^{\mathfrak{R}}|$
used in the induction along the coupling flow.

Then for every $j\le k$
\begin{equation}\label{4.7:eq-beta-extraction-1}
|\beta^{\mathfrak{R}}_j|\le C'g_j^{\kappa_1-2}=C'g_j^{\kappa_0-7},
\end{equation}
and at $j=k+1$ the bound \eqref{4.7:eq-beta-extraction-1} holds with $g_{k+1}$ replaced by
the stand-in $\tilde{g}=g_k$.
This bound is free of $k$, $K$, $m$, of position and of the site. At $j=k+1$, with
$C_{\Pi}$ the constant of (A4), which contains the factor $4C_{\mathsf{h}}^2C_A$,
\begin{equation}\label{4.7:eq-beta-extraction-a}
\begin{gathered}
C':=\frac{2C_{\Pi}K_{CE}}{(c_T\min\{C_0,C_1\})^2},\\
|\beta^{\mathfrak{R}}_{k+1}|\le C'g_k^{\kappa_0-7}\le C'\gamma^{\kappa_0-7}
\le C'\gamma^{\kappa_0-8}=\bar{\varsigma}.
\end{gathered}
\end{equation}
\end{proposition}

\begin{proof}
\emph{The coefficient.} Let $E^{(j),\mathfrak{R}}(X,\cdot)$ be the kernels without pin of
the large-field-descended sub-family at level $j$
(Definition~\ref{4.7:not-kernel-families}). Define
\begin{equation}\label{4.7:eq-beta-extraction-2}
\Pi^{(j),\mathfrak{R}}(b,b'):=\sum_X\partial_{B(b)}\partial_{B(b')}
E^{(j),\mathfrak{R}}\bigl(X,U_j(e^{iB})\bigr)\Big|_{B=0}.
\end{equation}
This is the whole-lattice sum of the sub-family. By definition, $\beta^{\mathfrak{R}}_j$ is
the coefficient that \cite[(5.42)]{B-CMP109} assigns to $\Pi^{(j),\mathfrak{R}}$.

\emph{Reduction to a second derivative along true pairs.} By \cite[(4.2)]{B-CMP109},
\begin{equation*}
E\bigl(X,U_j(\square_0,e^{iB})\bigr)=E\bigl(X,\exp i\xi\mathbf{H}_j(\square_0,B)\bigr),
\end{equation*}
with $\mathbf{H}_j$ the nonlinear potential function of \cite{B-CMP109}. For a
$\mathfrak{g}^c$-valued field $\mathsf{Z}$, put
\begin{equation*}
\mathcal{F}_X(\mathsf{Z}):=E^{(j),\mathfrak{R}}\bigl(X;e^{i\xi\mathsf{Z}},
J(e^{i\xi\mathsf{Z}})\bigr),
\end{equation*}
and let $H_j:=D\mathbf{H}_j(0)$. The summand in \eqref{4.7:eq-beta-extraction-2} is the
second derivative at $B=0$ of $B\mapsto\mathcal{F}_X(\mathbf{H}_j(\square_0,B))$. By the
chain rule it equals
\begin{equation*}
D^2\mathcal{F}_X(0)[H_j\delta_b,H_j\delta_{b'}]
+D\mathcal{F}_X(0)\bigl[D^2\mathbf{H}_j(0)[\delta_b,\delta_{b'}]\bigr].
\end{equation*}
By \cite[(4.14)]{B-CMP109} the first variation of the kernel vanishes at the configuration
$1$, that is $D\mathcal{F}_X(0)=0$, so the second term is zero.
The summand is therefore $D^2\mathcal{F}_X(0)[H_j\delta_b,H_j\delta_{b'}]$. The argument of
$\mathcal{F}_X$ is a configuration together with its own current, that is, a true pair.

\emph{Cauchy extraction along $H_j\delta_b$.} By (A3), the estimate
\eqref{4.7:eq-beta-extraction-3} holds for every bond $b$. For complex $s,t$ put
\begin{equation*}
f(s,t):=\mathcal{F}_X(sH_j\delta_b+tH_j\delta_{b'}).
\end{equation*}
On true pairs with $N_j(\cdot)\le c_Tr_j$, the pairs lie in
$U^{\Sigma^1}_j(X;\alpha_{0,j},\alpha_{1,j})$ by (A2). There (A1) holds, so $f$ is analytic
and
\begin{equation*}
|f(s,t)|\le C_{\mathfrak{R}}g_j^{\kappa_1}e^{-\kappa d_j(X)}.
\end{equation*}
Define the radii
\begin{equation*}
\mathsf{r}_b:=\frac{c_Tr_j\,e^{\delta_0\operatorname{dist}(b,X)}}{2C_{\mathsf{h}}},\qquad
\mathsf{r}_{b'}:=\frac{c_Tr_j\,e^{\delta_0\operatorname{dist}(b',X)}}{2C_{\mathsf{h}}}.
\end{equation*}
Let $|s|\le\mathsf{r}_b$ and $|t|\le\mathsf{r}_{b'}$. The triangle inequality for $N_j$ and
\eqref{4.7:eq-beta-extraction-3} give
\begin{equation*}
N_j\bigl((sH_j\delta_b+tH_j\delta_{b'})\lceil X\bigr)
\le\tfrac12c_Tr_j+\tfrac12c_Tr_j=c_Tr_j.
\end{equation*}
So the closed bidisc of radii $(\mathsf{r}_b,\mathsf{r}_{b'})$ lies in the region where the
bound on $f$ holds. The Cauchy estimate for $\partial_s\partial_t f(0,0)$ on this bidisc gives
\begin{equation}\label{4.7:eq-beta-extraction-4}
|\partial_s\partial_tf(0,0)|\le\frac{4C_{\mathsf{h}}^2C_{\mathfrak{R}}g_j^{\kappa_1}}
{(c_Tr_j)^2}\,e^{-\kappa d_j(X)-\delta_0[\operatorname{dist}(b,X)+\operatorname{dist}(b',X)]}.
\end{equation}
Since $\partial_s\partial_tf(0,0)=D^2\mathcal{F}_X(0)[H_j\delta_b,H_j\delta_{b'}]$, this
bounds the summand.

The constants in \eqref{4.7:eq-beta-extraction-4} do not depend on $j$. A ball in the sup
norm in place of the $N_j$-polydisc would cost powers of $\xi_j^{-1}$, with $\xi_j$ the
lattice spacing at level $j$.

At $j=k+1$ the estimate is taken inside the region
$U^{\Sigma^1}_{k+1}(X;\alpha_{0,k+1},\alpha_{1,k+1})$ whose radii are computed with the
stand-in $\tilde{g}=g_k$ in place of $g_{k+1}$.
There (A1) holds with $g_{k+1}$ replaced by $g_k$, in the bound and in the radii.

\emph{The second moment.} By \eqref{4.7:eq-beta-extraction-4}, the summands of
\eqref{4.7:eq-beta-extraction-2} satisfy the hypothesis of (A4) with
$S_j=C_{\mathfrak{R}}g_j^{\kappa_1}$, since the prefactor in
\eqref{4.7:eq-beta-extraction-4} is $4C_{\mathsf{h}}^2S_j(c_Tr_j)^{-2}$. The second-moment
formula \cite[(5.42)]{B-CMP109}, \cite[(3.62)]{B-CMP119}, with the summation bound
$\sum_X(Md_j(X)+c)^2e^{-\kappa d_j(X)}\le C_A$ of (A4), therefore gives
\begin{equation}\label{4.7:eq-beta-extraction-5}
|\beta^{\mathfrak{R}}_j|\le C_{\Pi}C_{\mathfrak{R}}g_j^{\kappa_1}(c_Tr_j)^{-2}.
\end{equation}
In the same way, let $\mathfrak{M}$ be any of the kernel families
of Definition~\ref{4.7:not-kernel-families} whose whole-lattice summands reduce, as above,
to second derivatives along $H_j\delta_b$, $H_j\delta_{b'}$ of kernels that are analytic and
bounded by $S_je^{-\kappa d_j(X)}$ on true pairs with $N_j(\cdot)\le c_Tr_j$. The Cauchy
estimate above, with $C_{\mathfrak{R}}g_j^{\kappa_1}$ replaced by $S_j$, bounds these
summands as in the hypothesis of (A4), and (A4) gives
\begin{equation}\label{4.7:eq-beta-extraction-6}
|\beta^{\mathfrak{M}}_j|\le C_{\Pi}S_j(c_Tr_j)^{-2}.
\end{equation}

\emph{Arithmetic.} By condition $(\Gamma4f)$ of
Definition~\ref{4.7:def-smallness-conditions}, $(\log g_j^{-2})^q\ge1$; with the radius
$r_j$ of \eqref{4.7:eq-true-pairs-a} this gives
\begin{equation*}
(c_Tr_j)^{-2}\le(c_T\min\{C_0,C_1\})^{-2}g_j^{-2}.
\end{equation*}
The factor discarded in this inequality is a negative power of $\log g_j^{-2}$, at most one.
With $C_{\mathfrak{R}}=2K_{CE}$, \eqref{4.7:eq-beta-extraction-5} becomes
\begin{equation*}
|\beta^{\mathfrak{R}}_j|\le\frac{2C_{\Pi}K_{CE}}{(c_T\min\{C_0,C_1\})^2}\,
g_j^{\kappa_1-2}=C'g_j^{\kappa_1-2}.
\end{equation*}
Since $\kappa_1=\kappa_0-5$, we have $g_j^{\kappa_1-2}=g_j^{\kappa_0-7}$. This is
\eqref{4.7:eq-beta-extraction-1}. None of the constants $C_{\Pi}$, $K_{CE}$, $c_T$, $C_0$,
$C_1$ depends on $k$, $K$, $m$, on position or on the site.

\emph{The step $j=k+1$.} Here $g_{k+1}$ is replaced by $g_k$, so the argument above, run
with $g_k$ in place of $g_{k+1}$, gives $|\beta^{\mathfrak{R}}_{k+1}|\le C'g_k^{\kappa_0-7}$.
Since $x_k\ge\gamma^{-2}$, we have $g_k\le\gamma$. Since $\kappa_0\ge17$, the exponent
$\kappa_0-7$ is positive, so $g_k^{\kappa_0-7}\le\gamma^{\kappa_0-7}$. Since $\gamma\le1$,
$\gamma^{\kappa_0-7}\le\gamma^{\kappa_0-8}$. By definition, $C'\gamma^{\kappa_0-8}$ is
$\bar{\varsigma}$. This proves \eqref{4.7:eq-beta-extraction-a}.

The requirement $\bar{\varsigma}\le\min\{c_0/4,1\}$ is the last ceiling that the induction
along the coupling flow places on $\gamma$. The constant $C'$ involves only $K_{CE}$,
$C_{\Pi}$, $c_T$, $C_0$ and $C_1$, and not $\gamma$, so this requirement involves no
circularity.
\end{proof}

\begin{proposition}[The large-field part of the coupling increment]\label{4.7:prop-coupling-increment}
Fix a level $k$ in the induction over levels that constructs the coupling flow (called below the level induction), and assume the following.
\begin{itemize}
\item The induction hypothesis $\mathfrak{i}(k)$ of the level induction holds, together with its hypotheses $(\gamma_k)$, taken with domain the one-power hereditary class $U^{\Sigma^1}$, and its flow hypothesis $(\beta_k)$. Moreover $x_0,\dots,x_k\ge\gamma^{-2}$, and $\gamma$ lies below the ceilings of the level induction.
\item The conditions $(\Gamma0)$--$(\Gamma6)$ and $(\varphi1)$--$(\varphi6)$ of \eqref{4.7:eq-smallness-conditions-a} hold, condition $(\Gamma6)$ being taken with its member (vi), which carries the factor $86d$, as fixed in Definition~\ref{4.7:def-smallness-conditions}. The floor $(\varphi B)$ of that display is adopted, as in Corollary~\ref{4.7:cor-kernel-power-bound}, but is not used below. Moreover $\gamma\le1$ and $\kappa_0\ge17$, and hypothesis (A3) of Proposition~\ref{4.7:prop-beta-extraction} holds.
\item Clauses (N) and (F) of the operator bounds at free current hold. So do the following clauses of the frame lemma, which those two clauses use beyond their own statement: part (a) of its first clause and part (ii) of its fourth clause, which carry the cost of the adjoint action of the frame map; its second clause, the bound on the global supremum of the window data by a multiple of $M\alpha_0$; its third clause, which accompanies that bound; and its fifth clause, the linear clause. So does the corollary supplying the operator bounds at free current, applied layer by layer.
\item The one-power hereditary class $U^{\Sigma^1}$ is the class of the induction. The hypotheses $(\gamma_k)$ and $(H_{\cdot})$, the condition \cite[(3.16)]{B-CMP109}, the hypothesis space $\mathfrak{U}$ of \cite[Lemma~4]{B-CMP109}, and the bound (E) of \eqref{4.7:eq-kernel-power-bound-a} are taken with domain $U^{\Sigma^1}$. Condition \cite[(1.12)]{B-CMP109} is taken in the form fixed in \eqref{4.7:eq-class-restrictions-a}. Collars are the cut $Z^{\wr p}$ of the class at every label, not the collar made of Ba{\l}aban's own cubes.
\item The restrictions \eqref{4.7:eq-class-restrictions-a} of the regularity class with its cube gauge condition (Definition~\ref{4.7:not-class-restrictions}) hold, and so does the proposition on membership in that class with its near, far and true-pair clauses. These rest on three inputs:
\begin{itemize}
\item the lemma on admissible towers accompanying the operator bounds at free current, at scale $M/L$;
\item the nesting lemma for \cite[(3.38)]{B-CMP109}, taken with the single factorisation $U:=1$ used in \cite{B-CMP109};
\item \cite[(1.17)]{B-CMP109}, cited as printed.
\end{itemize}
\item The following further statements are available:
\begin{itemize}
\item the corollary furnishing the constant $K_{CE}$ of \eqref{4.7:eq-kernel-power-bound-a};
\item clauses (a) and (d) of the theorem on the marked and unmarked kernel families, together with the convention extending the kernels to the whole lattice;
\item the statement that true pairs lie in the regularity class (Definition~\ref{4.7:not-true-pairs}), called below the true-pair statement;
\item the independence of the localised terms from the small-field region (Corollary~\ref{4.4:cor-localised-independence});
\item the lemma giving existence of the whole-lattice tensor and site-freeness of the coefficient (the tensor lemma), at all indices $\le k$.
\end{itemize}
\end{itemize}
Then the following hold.
\begin{itemize}
\item[(a)] The coefficient splits as
\begin{equation*}
\beta_{k+1}=\beta^{\mathfrak{A}}_{k+1}+\beta^{\mathfrak{R}}_{k+1}.
\end{equation*}
Each summand is the coefficient \cite[(5.42)]{B-CMP109} of the whole-lattice tensor of its own sub-family: the regular sub-family for $\beta^{\mathfrak{A}}_{k+1}$, the large-field-descended one for $\beta^{\mathfrak{R}}_{k+1}$. Two things are not asserted here: the existence of the $\mathbb{Z}^4$ tensor at every admissible site, and the site-freeness of $\beta_{k+1}$. Both are conclusions of the tensor lemma.
\item[(b)] The large-field part is bounded by
\begin{equation*}
|\beta^{\mathfrak{R}}_{k+1}|\le C'g_k^{\kappa_0-7}\le\bar{\varsigma},
\end{equation*}
with $C'$ and $\bar{\varsigma}$ as in \eqref{4.7:eq-beta-extraction-a}.
\item[(c)] The bound (E) of \eqref{4.7:eq-kernel-power-bound-a}, with $C_{\mathfrak{R}}$ and $\kappa_1=\kappa_0-5$ as there, holds for every $i\le k+1$, every localisation domain $X\in\mathbf{D}_i$ and every pin $z$, for two families of kernels: the unpinned kernels of \cite[(2.13)]{B-CMP116}, and the kernels $E^{(i),\mathfrak{R}}_{\mathfrak{s}}(X;\cdot,z)$ pointed at $z$ of \cite[(3.44)--(3.47)]{B-CMP119}. It holds at every admissible site $\mathfrak{s}\in\mathcal{S}_i^{\flat}$ on $U^{\Sigma^1}_i(X;\alpha_{0,i},\alpha_{1,i})$, where these kernels are analytic and $G^c$-gauge invariant. Hence (E) holds on the polydisc
\begin{equation}\label{4.7:eq-coupling-increment-1}
\bigl\{\mathbf{U}=\exp(\sqrt{-1}\,A)\ :\ \sup_X|A|<a^{\mathfrak{R}}_i,\ \ \sup_X|\mathbf{J}|<a^{\mathfrak{R}}_i\bigr\},
\end{equation}
with
\begin{equation}\label{4.7:eq-coupling-increment-a}
a^{\mathfrak{R}}_i:=\frac{\xi_i^{3}\,c_T\,r_i}{8d}>0,\qquad E^{\mathfrak{R}}_i:=C_{\mathfrak{R}}\,g_i^{\kappa_1},
\end{equation}
where $\xi_i=L^{-i}$ is the spacing at which the lattice differences in the proof below are taken, $c_T$ is the constant of \eqref{4.7:eq-true-pairs-a}, and $r_i$ is the radius with which the true-pair statement is formulated in Definition~\ref{4.7:not-true-pairs}; the quantities $a^{\mathfrak{R}}_i$ and $E^{\mathfrak{R}}_i$ are free of $\mathfrak{s}$, $X$ and $z$.

At $i=k+1$ the coupling $g_{k+1}$ is replaced by the stand-in $\tilde{g}=g_k$. There clause (c), with $\tilde{g}$ in the bound and in the radii, is the hypothesis $(H_{k+1})$(d) of the tensor lemma for the large-field-descended sub-family, namely its bound at the stand-in $\tilde{g}:=g_{j-1}$, taken at $j=k+1$. For $i\le k$ the bound, at the true couplings, re-proves what the hypothesis $(H_i)$(d) received from this proposition at induction index $i-1$.
\end{itemize}
\end{proposition}

\begin{proof}
\emph{Clause (a).} Clause (a) of the theorem on the marked and unmarked kernel families, with the convention extending the kernels to the whole lattice, gives the decomposition $\beta_{k+1}=\beta^{\mathfrak{A}}_{k+1}+\beta^{\mathfrak{R}}_{k+1}$ into the coefficients \cite[(5.42)]{B-CMP109} of the two sub-families. Nothing is used or claimed about the existence of the tensor or the site-freeness of $\beta_{k+1}$.

\emph{Clause (b).} Apply Corollary~\ref{4.7:cor-kernel-power-bound} at $i=k+1$. At this index the coupling $g_{k+1}$, both in the bound and in the radii, is replaced by the stand-in $\tilde{g}=g_k$. The corollary gives (E) of \eqref{4.7:eq-kernel-power-bound-a} for the large-field-descended kernels at level $k+1$, analytic on $U^{\Sigma^1}_{k+1}(X;\alpha_{0,k+1},\alpha_{1,k+1})$ with the radii evaluated at $\tilde{g}$.

Proposition~\ref{4.7:prop-beta-extraction} extracts the large-field part of the coefficient by Cauchy's estimate inside that domain. At $j=k+1$ it gives $|\beta^{\mathfrak{R}}_{k+1}|\le C'g_k^{\kappa_0-7}$, with $C'$ as in \eqref{4.7:eq-beta-extraction-a}. Since $x_k=g_k^{-2}\ge\gamma^{-2}$, we have $g_k\le\gamma$. With $\gamma\le1$ and $\kappa_0\ge17$ as assumed in the statement,
\begin{equation*}
C'g_k^{\kappa_0-7}\le C'\gamma^{\kappa_0-7}\le C'\gamma^{\kappa_0-8}=\bar{\varsigma},
\end{equation*}
where $\bar{\varsigma}$ is the ceiling of \eqref{4.7:eq-beta-extraction-a}.

\emph{Clause (c), the bound on the class.} For every $i\le k+1$, Corollary~\ref{4.7:cor-kernel-power-bound} gives (E) for the unpinned and the pointed large-field-descended kernels. It does so at every site $\mathfrak{s}\in\mathcal{S}_i^{\flat}$, on $U^{\Sigma^1}_i(X;\alpha_{0,i},\alpha_{1,i})$, where the kernels are analytic and $G^c$-gauge invariant. For $i\le k$ the couplings and radii are the true ones $g_i$, $\alpha_{0,i}$, $\alpha_{1,i}$; for $i=k+1$ the coupling is the stand-in $\tilde{g}=g_k$.

\emph{Clause (c), the inclusion of the polydisc.} It remains to show that the polydisc \eqref{4.7:eq-coupling-increment-1} lies in $U^{\Sigma^1}_i(X;\alpha_{0,i},\alpha_{1,i})$. Let $\mathbf{U}=\exp(\sqrt{-1}\,A)$ and $\mathbf{J}$ satisfy
\begin{equation*}
\sup_X|A|<a^{\mathfrak{R}}_i,\qquad \sup_X|\mathbf{J}|<a^{\mathfrak{R}}_i.
\end{equation*}
Extend $A$ by $0$ off $X$. Put $\mathsf{Z}:=A/\xi_i$; this is the $\mathfrak{g}^c$-valued field of the true-pair statement, and is distinct from the large-field domains $Z$ below. The lattice differences are taken at spacing $\xi_i\le1$. Measured in the units of level $i$ (the inequality does not depend on the choice of frame), they give
\begin{equation*}
N_i(\mathsf{Z})\le 8d\,\xi_i^{-3}\sup_X|A|<8d\,\xi_i^{-3}\,a^{\mathfrak{R}}_i=c_T\,r_i .
\end{equation*}
Here $N_i$ is the norm with which the true-pair statement is formulated (Definition~\ref{4.7:not-true-pairs}); it is not the quantity $N_k$ of the glossary (Definition~\ref{0.2:not-glossary}). The last equality is the definition \eqref{4.7:eq-coupling-increment-a} of $a^{\mathfrak{R}}_i$.

The true-pair statement therefore applies. It gives, for the tower of $X$ itself, the conditions (i), (ii) and (iv) on $\mathbf{U}$ of the definition of the class in \cite[(1.10)--(1.16)]{B-CMP109}, in the form of Definition~\ref{4.7:not-class-restrictions}, together with the first condition of \cite[(1.14)]{B-CMP109}. By \cite[(1.15)--(1.16)]{B-CMP109} these are functions of $\mathbf{U}$ alone.

The true-pair clause of the membership proposition gives condition (iv) in the form required by the class $U^{\Sigma^1}$ at every label $(p,Z)$ of the finite label set $\mathfrak{T}_i(X)$ of that class. Here \cite[(1.17)]{B-CMP109} is cited as printed, at the data $(p,Z;M/L)$, under the geometric condition, among the hypotheses of the membership proposition, on which its far and true-pair clauses rest.

Finally,
\begin{equation*}
|\mathbf{J}|<a^{\mathfrak{R}}_i\le\alpha_{0,i},
\end{equation*}
and $\alpha_{0,i}$ is the radius bounding the current in the second condition of \cite[(1.14)]{B-CMP109} at level $i$; so this is that second condition. Thus every point of the polydisc \eqref{4.7:eq-coupling-increment-1} lies in $U^{\Sigma^1}_i(X;\alpha_{0,i},\alpha_{1,i})$, and (E) holds there with the constant $E^{\mathfrak{R}}_i=C_{\mathfrak{R}}g_i^{\kappa_1}$.

Both $a^{\mathfrak{R}}_i$ and $E^{\mathfrak{R}}_i$ are formed from $\xi_i$, $c_T$, $r_i$, $d$, $C_{\mathfrak{R}}$ and $g_i$ only. Hence they are free of $\mathfrak{s}$, $X$ and $z$. Moreover $a^{\mathfrak{R}}_i>0$, since $\xi_i$, $c_T$ and $r_i$ are positive. Together with the identification of the couplings and radii in the first paragraph of the proof of clause (c), this gives the last paragraph of clause (c).

\emph{Why the factor $\xi_i^3$ is needed.} The radius of the polydisc cannot be taken to be $c_Tr_i$ without the factor $\xi_i^{3}/(8d)$. The first condition of \cite[(1.14)]{B-CMP109} reads, in the letters of that display, $|\partial\mathbf{U}-1|<\alpha_0\xi^2$, with $\xi$ the spacing and $\alpha_0$ that display's first radius, so a sup-radius $c_Tr_i$ does not give it; a factor $\xi_i^2$ would suffice. The tensor lemma uses only $a^{\mathfrak{R}}_i>0$.
\end{proof}

\begin{definition}[Ba{\l}aban's regularity spaces and their inclusion sentences
{\cite[(1.10)--(1.17), (3.16), (3.40)]{B-CMP109}, \cite[pp.~7, 9, 15]{B-CMP116}, \cite[p.~375]{B-CMP122b}}]
\label{4.8:not-regularity-spaces}
The spaces defined below are those of \cite[(1.10)--(1.17), pp.~262--263]{B-CMP109}. The
space and the implication collected under (a) and (e) at the end are those of
\cite[(3.16), (3.40)]{B-CMP109}. The inclusion sentences (b), (c), (d) are those of
\cite[p.~9]{B-CMP116}, \cite[p.~7]{B-CMP116} and \cite[p.~15]{B-CMP116} respectively, and
(f) is that of \cite[p.~375]{B-CMP122b}.

\emph{Setting.} Fix a level $j$ and put $\xi=L^{-j}$. Let $G=SU(N)$ be the gauge group.
Let $k$ denote the level of the step of \cite[\S3]{B-CMP109}.

\begin{itemize}
\item $\pi_j$ is Ba{\l}aban's family of closed cubes $\square$ of side $M=L^{m'}$ in the
lattice of spacing $\xi$, with centres as fixed in \cite{B-CMP109}.
\item For a cube $\square$ that is a union of cubes of one of the families $\pi_i$,
$\tilde{\square}^{n}$ denotes $\square$ enlarged by $n$ layers of cubes of $\pi_i$ on every
side; for $\square\in\pi_j$ it is a cube of side $(1+2n)M$.
\item For a cube $\square$ of the cover of \cite[\S3]{B-CMP109}, a union of $2^d$ neighbouring
cubes of $\pi_k$, $\square_0=\tilde{\square}^{5}$ denotes its window, the layers being cubes
of $\pi_k$.
\item $\mathbf{D}_j$ is the family of localization domains, the finite wall-connected unions
of $\pi_j$-cubes.
\item For a bond configuration $V$ we write $\partial V$ for the plaquette function
$p\mapsto V(\partial p)$.
\end{itemize}

The objects of interest are pairs $(\mathbf{U},\mathbf{J})$, where $\mathbf{U}$ is a
complexified bond configuration and $\mathbf{J}$ an external current. The constant $\beta>0$
is Ba{\l}aban's small constant of \cite[\S3]{B-CMP109}. The constant $B^{(109)}$ is the
sufficiently large constant $B$ of \cite[(1.12)]{B-CMP109}, as in
Definition~\ref{4.4:def-local-data-classes}; it is determined later in that paper. The
constant $c_{12}$ is the number written $O(1)$ in the bound of
\cite[(1.12)]{B-CMP109}. Finally, $B_3$ is the constant of
\cite[Proposition~9]{B-CMP102b}.

\emph{The spaces.} Let $X\in\mathbf{D}_j$ and $\alpha_0,\alpha_1,\gamma_0>0$. The space
$U^c_j(X,\alpha_0,\alpha_1,\gamma_0)$ is the union of the gauge orbits
$[(\mathbf{U},\mathbf{J})]$, in the sense of \cite[p.~262]{B-CMP109}, determined by
configurations $\mathbf{U}$, $\mathbf{J}$ that satisfy the following four conditions.
\begin{enumerate}
\item[(i)] $\mathbf{U}=U'U$ with $U$ a $G$-valued configuration such that
\begin{equation*}
|\partial U-1|<\alpha_0\xi^2\quad\text{on }X,
\end{equation*}
which is \cite[(1.11)]{B-CMP109}. Moreover, \cite[(1.12)]{B-CMP109} holds: for each cube
$\square\subset X$ of side $\mathsf{o}LM$, with $\mathsf{o}$ the constant of
Definition~\ref{4.4:def-local-data-classes}, there is a $G$-valued gauge transformation $u$,
defined on $\square$, such that $U^u=\exp i\xi A$ with
\begin{equation*}
|A|,\ |\nabla^\xi A|<c_{12}LMB^{(109)}\alpha_0\quad\text{on }\square .
\end{equation*}
Here $\nabla^\xi$ is the lattice gradient of spacing $\xi$, and $\nabla^\xi_U$, used in (ii),
is its covariant version in the background $U$, as in \cite[(1.12)--(1.13)]{B-CMP109}.
\item[(ii)] The field $A'$ that \cite[(1.10)]{B-CMP109} attaches to the factor $U'$ satisfies
$|A'|,\ |\nabla^\xi_U A'|<\alpha_1$, which is \cite[(1.13)]{B-CMP109}.
\item[(iii)] The bounds
\begin{equation*}
|\partial\mathbf{U}-1|<\alpha_0\xi^2,\qquad |\mathbf{J}|<\gamma_0
\end{equation*}
hold; this is \cite[(1.14)]{B-CMP109}.
\item[(iv)] Regard $X$ as $\Omega_0$. Let $\{\Omega_n\}_{n=1,\dots,j}$ be the sequence of
maximal possible domains satisfying the conditions \cite[(1.3)--(1.6)]{B-CMP99a}, read with
$j,\xi$ in place of $k,\eta$ and with $R=R_1$, where $R_1$ is the constant, a power of $L$,
of \cite[(1.3)--(1.4)]{B-CMP99a}. Let $\Lambda_p$ be the bond sets of \cite[\S1]{B-CMP99a}
carrying the data of this sequence at level $p$, and let
$\bar{\mathbf{U}}^p$ denote Ba{\l}aban's average of $\mathbf{U}$ at level $p$ in
\cite{B-CMP109}.
\begin{itemize}
\item For each subsequence $\{\Omega_0,\Omega_1,\dots,\Omega_n\}$, $n=1,\dots,j$, Ba{\l}aban
constructs in the axial gauges the functions $U_n(V)$: $U_n(V)$ is the minimiser of the
variational problem of depth $n$ on this sequence of domains, with data $V$ on the sets
$\Lambda_p$.
\item The layered average $M^{\boldsymbol{\cdot}}(\mathbf{U})$ is given by
$M^{\boldsymbol{\cdot}}(\mathbf{U},b)=M^p(\mathbf{U},b)=\bar{\mathbf{U}}^p(b)$ for
$b\in\Lambda_p$.
\item $J_n(M^{\boldsymbol{\cdot}}(\mathbf{U}))$ is defined by \cite[(1.8)]{B-CMP109}, with
$U_n(M^{\boldsymbol{\cdot}}(\mathbf{U}))$ in place of $U_j$ and $L^{-n}$ in place of $\xi$.
\end{itemize}
The condition is that, for $n=1,\dots,j$, the pair
\begin{equation*}
\bigl(U_n(M^{\boldsymbol{\cdot}}(\mathbf{U})),\,J_n(M^{\boldsymbol{\cdot}}(\mathbf{U}))\bigr),
\end{equation*}
which is \cite[(1.15)]{B-CMP109}, satisfies \cite[(1.16)]{B-CMP109}:
\begin{equation}\label{4.8:eq-regularity-spaces-1}
|\partial U_n(M^{\boldsymbol{\cdot}}(\mathbf{U}))-1|<\alpha_0\xi^2,\qquad
|J_n(M^{\boldsymbol{\cdot}}(\mathbf{U}))|<\alpha_0(L^n\xi)^2\qquad\text{on }\tilde{X}^{-2}.
\end{equation}
The same bounds are required with the $G$-valued factor $U$ of (i) in place of $\mathbf{U}$.
Here $\tilde{X}^{-2}$ is obtained from $X$ by removing the two layers of cubes of $\pi_j$
that are closest to the boundary of $X$.
\end{enumerate}
In the inclusion sentences below, conditions (i)--(iii) are also read on other domains in
place of $X$, as in \cite[pp.~7, 15]{B-CMP116}.
Condition (iv) is imposed because, in constructions later in \cite{B-CMP109}, the functions
\cite[(1.15)]{B-CMP109} are substituted in place of the variables $\mathbf{U}$,
$\mathbf{J}$, and they must then have the required properties. When $\gamma_0=\alpha_0$ the
last argument is omitted:
\begin{equation*}
U^c_j(X,\alpha_0,\alpha_1)=U^c_j(X,\alpha_0,\alpha_1,\alpha_0).
\end{equation*}

\emph{A sufficient condition.} The following is stated in
\cite[(1.17), p.~263]{B-CMP109}. Let $\alpha_0'$, $\alpha_1'$ be constants smaller than
$\alpha_0$, $\alpha_1$ respectively, and suppose that $(\mathbf{U},\mathbf{J})$ satisfies
(i)--(iii) with $\alpha_0'$, $\alpha_1'$ in place of $\alpha_0$, $\alpha_1$. Recall that
$M^{\boldsymbol{\cdot}}(\mathbf{U})=\bar{\mathbf{U}}^p$ on $\Lambda_p$; for these averages
\cite[p.~263]{B-CMP109} records the bound
$|\partial\bar{\mathbf{U}}^p-1|<2\alpha_0'(L^p\xi)^2$, and \cite[Proposition~9]{B-CMP102b},
applied there to this datum, gives, on $\tilde{X}^{-2}$,
\begin{align*}
|\partial U_n(M^{\boldsymbol{\cdot}}(\mathbf{U}))-1|
&<B_3\,2\alpha_0'L^{-2n}(L^n\xi)^2=2B_3\alpha_0'\xi^2,\\
|J_n(M^{\boldsymbol{\cdot}}(\mathbf{U}))|&<B_3\,2\alpha_0'(L^n\xi)^2 .
\end{align*}
In \cite[p.~263]{B-CMP109} it is concluded that for $\alpha_0'$ sufficiently small these
estimates imply condition (iv), so that $(\mathbf{U},\mathbf{J})\in U^c_j(X,\alpha_0,\alpha_1)$.

As a special case, \cite[p.~263]{B-CMP109} considers configurations that satisfy (i)--(iii)
with $\alpha_0'$, $\alpha_1'$ as above and for which $\mathbf{J}$ is the current
\cite[(1.8)]{B-CMP109} built on $\mathbf{U}$ and satisfies $|\mathbf{J}|<\alpha_0'$, and
requires the same of the configuration constructed in the same way from $U$ in place of
$\mathbf{U}$. It is stated there that $(\mathbf{U},\mathbf{J})$ then belongs to
$U^c_j(X,\alpha_0,\alpha_1)$, and that, in particular, the minimal configurations $U_j$ with
$|\partial U_j-1|<\varepsilon_0\xi^2$, for $\varepsilon_0$ sufficiently small, satisfy these
conditions.

\emph{Inclusion sentences.} Here $k$ is the level of the step, as above.
\begin{enumerate}
\item[(a)] The space of \cite[(3.16)]{B-CMP109}, for $X\in\mathbf{D}_j$, $j\le k$, is
\begin{equation}\label{4.8:eq-regularity-spaces-2}
\begin{split}
&U^c_j\bigl(X,\tfrac12\alpha_0,\tfrac12\alpha_1,\alpha_0\bigr)\cap
\bigl\{(\mathbf{U},\mathbf{J}):\ \mathbf{U},\mathbf{J}\text{ satisfy (i)--(iii) at the
level }k+1,\\
&\qquad\text{that is, with }\xi=L^{-(k+1)},\text{ and with the constants }
(1+\beta)\alpha_0,\ (1+\beta)\alpha_1,\ \alpha_0\bigr\}.
\end{split}
\end{equation}
\item[(b)] Summing the terms with a common localization domain $Y\in\mathbf{D}_k$ takes
place in \cite[p.~9]{B-CMP116}. The terms there carry domains $X\in\mathbf{D}_j$, $j\le k$,
with $X\subset X_0\subset Y$ for the domain $X_0$ of that construction. The following is
stated there: each term is defined and analytic on the corresponding space
\eqref{4.8:eq-regularity-spaces-2} restricted to $Y$; all these spaces contain the subspace
\begin{equation*}
U^c_{k+1}\bigl(Y,(1+\beta)\alpha_0,(1+\beta)\alpha_1,\alpha_0\bigr);
\end{equation*}
and all the terms are defined and analytic on it.
\item[(c)] In \cite[p.~7]{B-CMP116} a term is defined on the space written
\begin{equation*}
U'^c_{k+1}\bigl(\square_0,(1+2\beta)\alpha_0,(1+2\beta)\alpha_1,\alpha_0\bigr)
\end{equation*}
there. It is observed there that, by the construction, the term depends only on
$\mathbf{U}$, $\mathbf{J}$ restricted to $Y_0\cup\square_0=Y$, where $Y_0$ is the domain of
that construction, and that the conditions outside $Y$ are inessential; the space is then
regarded as defined by conditions (i)--(iii) on $Y$, and its subspace
$U^c_{k+1}(Y,(1+\beta)\alpha_0,(1+\beta)\alpha_1,\alpha_0)$ is taken.
\item[(d)] Let $X\in\mathbf{D}_{k+1}$. The following statements are made in
\cite[p.~15]{B-CMP116}.
\begin{itemize}
\item The functions treated on that page, which carry a domain $Y$, extend to functions of
$(\mathbf{U},\mathbf{J})$ that are analytic on
\begin{equation*}
U^c_{k+1}\bigl(Y,(1+\beta)\alpha_0,(1+\beta)\alpha_1,\alpha_0\bigr).
\end{equation*}
\item The domains $Z_0$, $\tilde{Z}_0$ and $Z$ attached to the activities in
\cite[(2.13)]{B-CMP116} satisfy $Y\subset Z_0$ and $\tilde{Z}_0\subset Z\subset X$.
\item The potentials are analytic on the subspace $U^c_{k+1}(X,\alpha_0,\alpha_1)$.
\item The quadratic forms and covariances entering the term $H(Z)$ of \cite[p.~15]{B-CMP116}
are analytic on the space of configurations satisfying (i)--(iii) on $Z$, with constants
much bigger than $\alpha_0$, $\alpha_1$; these constants are written $\alpha_0'$,
$\alpha_1'$ on that page and are unrelated to the constants $\alpha_0'$, $\alpha_1'$ of the
sufficient condition above.
\item Hence the activities in \cite[(2.13)]{B-CMP116}, and the whole sum
$\mathbf{E}^{(k+1)}(X)$, are analytic on $U^c_{k+1}(X,\alpha_0,\alpha_1)$; this is the
analyticity statement of the inductive assumptions.
\item For the pair $(U,0)$ the operators are symmetric and the measure is positive.
\end{itemize}
\item[(e)] By \cite[(3.40)--(3.41)]{B-CMP109}, the assumption that $\mathbf{U}$ belongs to
\begin{equation*}
U^c_{k+1}\bigl(\square_0,(1+2\beta)\alpha_0,(1+\beta)\alpha_1,\alpha_0\bigr)
\end{equation*}
implies in particular the inequalities \cite[(3.41)]{B-CMP109} on $\tilde{\square}^3$. These
are the bounds of condition (iv) for the window $\square_0$, with the constant
$(1+2\beta)\alpha_0$, on $\tilde{\square}^3$, that is, on $\square_0$ with its two outer
layers of cubes of $\pi_k$ removed (layers of $\pi_k$, although the space is indexed by
$k+1$).
\item[(f)] A statement of the same kind, for other spaces, is made in
\cite[p.~375]{B-CMP122b}. It concerns the display \cite[(1.65)]{B-CMP122b}, which for the
composite configuration $U_k^0(\mathbf{U},\mathbf{J})$ of that paper reads
\begin{equation*}
\begin{split}
&(\mathbf{U},\mathbf{J})\in\tilde U^c_k(Y_4,\tilde\alpha_0,\tilde\alpha_1)\\
&\qquad\Longrightarrow\quad
U_k^0(\mathbf{U},\mathbf{J})\in\tilde U''^c_k(Y_1,\tilde\alpha_0,\tilde\alpha_1);
\end{split}
\end{equation*}
here $Y_4$, $Y_1$ are domains, $\tilde U^c_k$, $\tilde U''^c_k$ regularity spaces and
$\tilde\alpha_0$, $\tilde\alpha_1$ constants, all as defined in \cite{B-CMP122b}. The
statement is that this property is simple to see for all the old terms, because they are
connected with the old determining set. For that set the regularity conditions on the
crucial domain $\Omega''_{k_0+1}\setminus\Omega_{k_0+1}$ of that construction, $k_0$ being
the index fixed there, are weaker than those for the new determining set. This concerns the
spaces of \cite{B-CMP122b}, not the space $U^c_j$ defined above.
\end{enumerate}
That \cite[Proposition~9]{B-CMP102b} applies to the complex averages $\bar{\mathbf{U}}^p$,
and with it the conclusion of the sufficient condition above and its special case at complex
pairs $(\mathbf{U},\mathbf{J})$, is stated in \cite[(1.17), p.~263]{B-CMP109} without proof;
cited as printed.
The inclusions in (b) and (d) are stated in \cite[p.~9, p.~15]{B-CMP116} without proof;
cited as printed. The passage to the subspace in (c) is stated in \cite[p.~7]{B-CMP116}
without proof; cited as printed. The sentence reported in (f) is stated in
\cite[p.~375]{B-CMP122b} without proof; cited as printed. No inclusion statement between
the spaces $U^c_j$ occurs in \cite{B-CMP109}. A further regularity class is defined in
\cite[(2.34)--(2.39), p.~261]{B-CMP119}, where it is stated that one can take $\beta=1/4$;
no inclusion statement for that class occurs there.
\end{definition}

\begin{definition}[Cube lattices, localisation domains and two-layer interiors]\label{4.8:not-cube-lattices}
All lattices below are sublattices of one fine lattice, and lengths are measured in its steps.
In this section the letters $m$ and $p$ denote levels, not the torus exponent or a plaquette.
At level $m$ the unit of length is $L^m$ steps, so that one step has length
$\xi_m=L^{-m}$ in the units of level $m$; this is the spacing $\xi$ of the lattice at level $m$
in \cite{B-CMP109}. Here $L$ is the odd blocking factor (in \cite[\S1]{B-CMP109} an odd integer
greater than $11$), and $M=L^{m'}$; below only the oddness of $L$ and $L\ge2$ are used.
\begin{itemize}
\item[(a)] $\pi_m$ is the lattice of closed cubes of size $M$ at level $m$, with centres at points
of the level-$m$ lattice (coordinates in $L^m\mathbb{Z}$, in steps) \cite[\S1]{B-CMP109}; a cube
of $\pi_m$ has side $ML^m$ steps. In the frame of \cite[\S1]{B-CMP109}, in which the cubes have
their corners at lattice points, the walls of
the cubes of $\pi_m$ lie on the hyperplanes on which one coordinate, in steps, belongs to
$ML^m\mathbb{Z}$; in the frame in which cubes are centred at lattice points, on those on which
one coordinate belongs to $ML^m(\mathbb{Z}+\tfrac12)$.
\item[(b)] For a cube $\square$ which is a union of $\pi_m$-cubes and for $n\ge0$,
$\tilde{\square}^{n}$ is the cube with the same centre obtained from $\square$ by adding $n$
layers of $\pi_m$-cubes. For $\square\in\pi_m$ this is the cube of size $(1+2n)M$ with centre at
the centre of $\square$, as in \cite[\S1]{B-CMP109}; for a cover cube of \cite[\S3]{B-CMP109}, a
union of $2^d$ neighbouring cubes of $\pi_m$, of size $2M$ (side $2ML^m$ steps), it has size
$(2+2n)M$. The level of the added
layers is part of the notation and is named at each use.
\item[(c)] A finite family of cubes of $\pi_m$ is \emph{wall-connected}, in the sense of
\cite[\S1]{B-CMP109}, if any two of its cubes
are joined by a chain of cubes of the family in which consecutive cubes share a common face of
dimension $d-1$. $\mathbf{D}_m$ is the family of sets which are unions of finite wall-connected
families of cubes of $\pi_m$ \cite[\S1]{B-CMP109}. A set may belong to the families of several
levels.
\item[(d)] For $p\ge1$ and $Z\in\mathbf{D}_{p-1}$ the \emph{two-layer interior} of $Z$ at $p$ is
\begin{equation}\label{4.8:eq-cube-lattices-b}
Z^{\wr p}:=\bigcup\bigl\{\square\in\pi_{p-1}:\ \square\subset Z,\ \tilde{\square}^{2}\subset Z\bigr\},
\end{equation}
with $\tilde{\square}^{2}$ formed with layers of $\pi_{p-1}$-cubes: two layers of cubes are taken
away, the cut being made in $\pi_{p-1}$, so that $Z^{\wr p}$ depends on $p$.
\item[(e)] For $X\in\mathbf{D}_p$, $\tilde{X}^{-2}$ is the domain of clause (iv) of Ba{\l}aban's
regularity spaces, Definition~\ref{4.8:not-regularity-spaces}: it is obtained from $X$ by taking
away the
two layers of cubes of $\pi_p$ which are closest to the boundary of $X$. We read these layers as
follows: the first layer consists of the cubes of $\pi_p$ in $X$ that meet the boundary of $X$,
the second of the remaining cubes of $\pi_p$ in $X$ that meet a cube of the first layer.
\end{itemize}
Then, for every level $m\ge1$,
\begin{equation}\label{4.8:eq-cube-lattices-a}
\mathbf{D}_m\subset\mathbf{D}_{m-1}\quad\text{as families of sets.}
\end{equation}
Moreover, let $p\ge1$, let $Z,Z'\in\mathbf{D}_{p-1}$, let $\square$ be a cube which is a union
of $\pi_{p-1}$-cubes, let $\tilde{\square}^{n}$ be formed with layers of $\pi_{p-1}$-cubes, put
$\square_0=\tilde{\square}^{5}$ (a \emph{window}), and let $X\in\mathbf{D}_p$. Then
\begin{equation}\label{4.8:eq-cube-lattices-c}
\begin{aligned}
&Z\subset Z'\ \Longrightarrow\ Z^{\wr p}\subset Z'^{\wr p};\\
&\square_0\in\mathbf{D}_{p-1},\qquad \square_0^{\wr p}=\tilde{\square}^{3};\\
&\tilde{X}^{-2}\subset X^{\wr(p+1)}\subset X^{\wr p}.
\end{aligned}
\end{equation}
\end{definition}

\begin{proof}
\emph{The inclusion \eqref{4.8:eq-cube-lattices-a}.} In the corner frame the walls of $\pi_m$ lie
at coordinate values in $ML^m\mathbb{Z}=ML^{m-1}(L\mathbb{Z})\subset ML^{m-1}\mathbb{Z}$. In the
centred frame, for $i\in\mathbb{Z}$,
\begin{equation*}
ML^m\bigl(i+\tfrac12\bigr)=ML^{m-1}\,L\bigl(i+\tfrac12\bigr)
=ML^{m-1}\Bigl(\bigl(Li+\tfrac{L-1}{2}\bigr)+\tfrac12\Bigr),
\end{equation*}
and $Li+\tfrac{L-1}{2}\in\mathbb{Z}$ because $L$ is odd. In either frame every wall of $\pi_m$ is
a wall of $\pi_{m-1}$. Since the side $ML^m$ equals $L$ times the side $ML^{m-1}$, every cube of
$\pi_m$ is the union of $L^d$ cubes of $\pi_{m-1}$, and no cube of $\pi_{m-1}$ crosses a wall of
$\pi_m$. Let $Y\in\mathbf{D}_m$ be the union of a finite wall-connected family $F$ of
$\pi_m$-cubes, and let $F'$ be the family of the $\pi_{m-1}$-cubes contained in cubes of $F$; it
is finite and its union is $Y$. The $L^d$ sub-cubes of one cube of $F$ form an array in which any
two are joined by a chain of sub-cubes with consecutive members sharing a face. If two cubes
$\square_1,\square_2$ of $F$ share a wall $W$, then $W$ is a union of faces of sub-cubes of
$\square_1$; the $\pi_{m-1}$-cube on the other side of such a face lies in $\square_2$, since it
does not cross $W$. Concatenating chains along a chain in $F$ shows that $F'$ is wall-connected,
so $Y\in\mathbf{D}_{m-1}$.

\emph{First line of \eqref{4.8:eq-cube-lattices-c}.} If $\square\in\pi_{p-1}$ satisfies
$\square\subset Z$ and $\tilde{\square}^{2}\subset Z$, then $\square\subset Z'$ and
$\tilde{\square}^{2}\subset Z'$, so each cube in the union \eqref{4.8:eq-cube-lattices-b} for $Z$
occurs in that for $Z'$.

\emph{Second line.} $\square_0=\tilde{\square}^{5}$ is a cube made of $\pi_{p-1}$-cubes, namely
$\square$ and five layers around it. These form a rectangular array, which is finite and
wall-connected, so $\square_0\in\mathbf{D}_{p-1}$. Index the $\pi_{p-1}$-cubes in each coordinate
direction by consecutive integers, and let $\square$ occupy the indices from $a_\mu$ to $b_\mu$ in
direction $\mu$. Then $\square_0$ occupies the indices from $a_\mu-5$ to $b_\mu+5$. A
$\pi_{p-1}$-cube $\Delta\subset\square_0$ with index $i_\mu$ has $\tilde{\Delta}^{2}$ occupying
the indices from $i_\mu-2$ to $i_\mu+2$, hence $\tilde{\Delta}^{2}\subset\square_0$ if and only
if $a_\mu-3\le i_\mu\le b_\mu+3$ for every $\mu$, that is, if and only if
$\Delta\subset\tilde{\square}^{3}$. The union of these cubes is $\tilde{\square}^{3}$, which
proves $\square_0^{\wr p}=\tilde{\square}^{3}$.

\emph{Third line.} By \eqref{4.8:eq-cube-lattices-a}, $X\in\mathbf{D}_p\subset\mathbf{D}_{p-1}$,
so both $X^{\wr(p+1)}$ and $X^{\wr p}$ are defined. Let $\square\in\pi_p$ be a cube of $X$ that
survives in $\tilde{X}^{-2}$, with $\tilde{\square}^{1},\tilde{\square}^{2}$ formed with
$\pi_p$-layers. Since $\square$ is not in the first layer of item (e), it does not meet the
boundary of $X$,
so every $\pi_p$-cube touching $\square$ lies in $X$; since $\square$ is not in the second layer,
it meets no cube of the first layer, so none of these touching cubes meets the boundary of $X$,
and hence each such cube $\blacksquare$ satisfies $\tilde{\blacksquare}^{1}\subset X$. As
$\tilde{\square}^{2}$ is the union of these
$\tilde{\blacksquare}^{1}$, $\tilde{\square}^{2}\subset X$, and $\square\subset X^{\wr(p+1)}$.
This is the first inclusion.

For the second, let $\square\in\pi_p$ with $\square\subset X$ and $\tilde{\square}^{2}\subset X$
($\pi_p$-layers), and let $\square'\subset\square$ be one of its $L^d$ cubes of $\pi_{p-1}$, with
$\tilde{\square}'^{2}$ formed with $\pi_{p-1}$-layers. The cube $\tilde{\square}'^{2}$ consists
of the points within distance $2ML^{p-1}$ of $\square'$ in every coordinate, and
$\tilde{\square}^{1}$ of the points within distance $ML^p$ of $\square$. Since
$\square'\subset\square$ and $2ML^{p-1}\le ML^p$,
\begin{equation*}
\tilde{\square}'^{2}\subset\tilde{\square}^{1}\subset\tilde{\square}^{2}\subset X.
\end{equation*}
Hence $\square'\subset X^{\wr p}$. Since $\square$ is the union of such $\square'$,
$X^{\wr(p+1)}\subset X^{\wr p}$.
\end{proof}

\begin{definition}[The two tower rules and their collars]\label{4.8:def-tower-rules}
Throughout, lengths are counted in steps of the lattice, of spacing $\eta$, on which the sets
below live. $A^{c}$ denotes the complement of a set $A$ in that lattice, and $\operatorname{dist}$
the lattice distance, with $\operatorname{dist}(\cdot,\emptyset)=+\infty$. Cubes are cubes of the
cube lattices of Definition~\ref{4.8:not-cube-lattices}.

\emph{Block letters.} The letter $\mathfrak{M}$ denotes the size of big blocks of the propagator
papers. It depends on $d$ and $L$ only, and is written $M_1$ in \cite{B-CMP99a} and in
\cite[(1)]{B-CMP102b}. The letter $M_1$ is the block letter of \cite{B-CMP119}. It is a power of
$L$, and
\begin{equation*}
M_1=R_\star'\,\mathfrak{M}\ge\mathfrak{M},
\end{equation*}
where $R_\star':=M_1/\mathfrak{M}$ is the ratio of the two block letters. The gap number $R_1$ of
\cite{B-CMP99a} is a sufficiently large positive integer, a power of $L$, such that the propagator
theorems on which \cite{B-CMP99a} relies hold for blocks of side $R\mathfrak{M}$ whenever
$R\ge R_1$. It obeys the one-sided floor $R_1\ge L$, and the ratio of the block letters obeys
$R_\star'\ge R_1$. Since $M_1\ge\mathfrak{M}$, any lower bound stated with $M_1$ implies the same
lower bound with $\mathfrak{M}$ in its place.

\emph{Towers.} Let $Z$ be a set, let $n\ge0$, and let $\mathfrak{b}\in\{\mathfrak{M},M_1\}$. A
\emph{tower} of $Z$ to depth $n$ with block $\mathfrak{b}$ is a sequence
\begin{equation*}
Z=\Omega_0\supset\Omega_1\supset\dots\supset\Omega_n ,
\end{equation*}
in which each $\Omega_q$ is a union of cubes of side $\mathfrak{b}L^{q}$ steps and consecutive
members obey a separation rule. The case in which a member is the whole lattice is admitted. The
\emph{data sets} of the tower are
\begin{equation*}
\Lambda_q:=\Omega_q\setminus\Omega_{q+1}\quad(q<n),\qquad \Lambda_n:=\Omega_n .
\end{equation*}
The data set $\Lambda_0$ carries no averaging, in agreement with the convention
$B^{0}(\Lambda_0)=\Lambda_0$ of \cite{B-CMP99a}, $B^{0}$ denoting the zero-fold block operation.

\emph{The two rules.} In each rule, $\Omega_{q+1}$ (respectively $\Omega_q$) is chosen maximal,
given the preceding member.
\begin{itemize}
\item[$\mathfrak{r}_1$.] This is the rule of condition (iv) in the definition of the spaces
$U^c_j(X,\alpha_0,\alpha_1,\gamma_0)$ of \cite[p.~262]{B-CMP109}, with block $\mathfrak{M}$. The
tower is the maximal sequence obeying \cite[(1.3)--(1.4)]{B-CMP99a} with $R=R_1$; this is the
condition \cite[(1)]{B-CMP102b}, stated there for every $R\ge R_1$. In steps, $\Omega_q$ is a union
of cubes of side $\mathfrak{M}L^{q}$, and
\begin{equation*}
\operatorname{dist}(\Omega_q^{c},\Omega_{q+1})>R_1\mathfrak{M}L^{q}\qquad(0\le q<n).
\end{equation*}
\item[$\mathfrak{r}_2$.] This is the rule of \cite[(2.13), p.~256]{B-CMP119}, with block $M_1$. In
steps, $\Omega_q$ is a union of cubes of side $M_1L^{q}$, and
\begin{equation*}
\operatorname{dist}(\Omega_q,\Omega_{q-1}^{c})\ge M_1L^{q}\qquad(1\le q\le n).
\end{equation*}
A tower built under $\mathfrak{r}_2$ satisfies the conditions of \cite[(1)]{B-CMP102b} with block
$\mathfrak{M}$ and $R:=R_\star'\ge R_1$. Take $Z=\square_0$, the window $\tilde{\square}^{5}$ of a
cover cube (the five-layer enlargement of Definition~\ref{4.8:not-cube-lattices}). Then
$\mathfrak{r}_2$ gives the tower $\{\square_n\}$ on which \cite[(2.13)]{B-CMP119} is read; this is
the dial tower of reading convention $(\rho2)$ of Definition~\ref{4.5:not-reading-conventions},
and the same tower enters clause (iv) in the definition of the class of configurations of the
operator theorem at free current.
\end{itemize}
Put $\mathfrak{R}:=\{\mathfrak{r}_1,\mathfrak{r}_2\}$. For $\mathfrak{r}\in\mathfrak{R}$ write
$\boldsymbol{\Omega}^{\mathfrak{r}}(Z)=(\Omega^{\mathfrak{r}}_0(Z),\dots,\Omega^{\mathfrak{r}}_n(Z))$
for the tower of $Z$ to depth $n$ under $\mathfrak{r}$.

Let $U_n^{\mathfrak{r}}(Z;\cdot)$ be the axial-gauge minimiser of the depth-$n$ problem on
$\boldsymbol{\Omega}^{\mathfrak{r}}(Z)$, as given by \cite[Theorem~1]{B-CMP102b}. Let $\mathbf{U}$
be a complexified bond configuration, and let $M^{\boldsymbol{\cdot}}(\mathbf{U})$ be its layered
datum in the sense of \cite[p.~262]{B-CMP109},
\begin{equation*}
M^{\boldsymbol{\cdot}}(\mathbf{U},b)=M^{p}(\mathbf{U},b)=\bar{\mathbf{U}}^{p}(b)\qquad
(b\in\Lambda_p),
\end{equation*}
where $\bar{\mathbf{U}}^{p}(b)$ is the block average of $\mathbf{U}$ after $p$ averaging steps
(Definition~\ref{1.5:def-block-average}), and $\bar{\mathbf{U}}^{0}=\mathbf{U}$ on $\Lambda_0$.
Let $J_n^{\mathfrak{r}}(Z;M^{\boldsymbol{\cdot}}(\mathbf{U}))$ be defined by \cite[(1.8)]{B-CMP109}.
In that formula, $U_n^{\mathfrak{r}}(Z;M^{\boldsymbol{\cdot}}(\mathbf{U}))$ replaces $U_j$, and
$L^{-n}$ replaces $\xi$. Then, for $\mathfrak{r}\in\mathfrak{R}$, the objects
\begin{equation}\label{4.8:eq-tower-rules-a}
\begin{gathered}
\boldsymbol{\Omega}^{\mathfrak{r}}(Z),\qquad U_n^{\mathfrak{r}}(Z;\cdot),\\
\mathbf{U}\longmapsto\bigl(U_n^{\mathfrak{r}}(Z;M^{\boldsymbol{\cdot}}(\mathbf{U})),\,
J_n^{\mathfrak{r}}(Z;M^{\boldsymbol{\cdot}}(\mathbf{U}))\bigr)
\end{gathered}
\end{equation}
are determined by the set $Z$, the rule $\mathfrak{r}$ and $n$ alone, and read $\mathbf{U}|_{Z}$
only. Neither rule names a level.

\emph{Collars.} For $\mathfrak{r}\in\mathfrak{R}$ define, with $\max\emptyset:=0$,
\begin{equation*}
\operatorname{collar}(\mathfrak{r}):=\sup_{Z,\,n}\;L^{-n}\max\bigl\{\operatorname{dist}(x,Z^{c}):
x\in Z\setminus\Omega^{\mathfrak{r}}_n(Z)\bigr\}.
\end{equation*}
This is the width of the part of $Z$ not covered by the last member, in top-block units $L^{n}$.
If $L\ge4d$ and $R_1\ge4d$, then
\begin{equation}\label{4.8:eq-tower-rules-b}
\operatorname{collar}(\mathfrak{r}_1)\le R_1M_1,\qquad
\operatorname{collar}(\mathfrak{r}_2)\le3dM_1\le R_1M_1 .
\end{equation}
Both collars are therefore at most $\tfrac12M'$ for every number $M'\ge2R_1M_1$.
In \cite[p.~256]{B-CMP119} the sharper bound $2M_1$ is stated for $\mathfrak{r}_2$, in the form
that the two-layer interior of the base lies in the last member; it is not used here.

\emph{Adjoining further rules.} A further rule may be adjoined to $\mathfrak{R}$ provided it names
no level, satisfies the conditions of \cite[(1)]{B-CMP102b}, and has collar at most $R_1M_1$. The
properties of towers that rest on \eqref{4.8:eq-tower-rules-a} and \eqref{4.8:eq-tower-rules-b}
then extend to it. Statements whose source is one of Ba{\l}aban's regularity spaces $\tilde U^c$
of \cite{B-CMP119} do not extend to it. Both rules of
$\mathfrak{R}=\{\mathfrak{r}_1,\mathfrak{r}_2\}$ are used in what follows:
\begin{itemize}
\item $\mathfrak{r}_1$ carries the class $U^{\mathrm{lit}}_j$ given by the conditions of
\cite[p.~262]{B-CMP109} read literally, which contains the one-power hereditary class
$U^{\Sigma^1}_j$, and the window tower of Ba{\l}aban's construction;
\item $\mathfrak{r}_2$ carries the window tower of the operator bounds at free current and clause
(iv) in the form of \cite{B-CMP119}.
\end{itemize}
The set $\{\mathfrak{r}_2\}$ alone would not carry the first of these.
\end{definition}

\begin{proof}
\emph{The towers are well defined.} Let $\mathfrak{r}=\mathfrak{r}_1$, let $\Omega_q$ be given, and
let $\mathcal{C}$ be a family of cubes of side $\mathfrak{M}L^{q+1}$. The distance from a set to a
union is the minimum of the distances to the pieces. Hence the union of $\mathcal{C}$ satisfies
$\operatorname{dist}(\Omega_q^c,\cdot)>R_1\mathfrak{M}L^q$ if and only if each cube of $\mathcal{C}$
does. Any cube at positive distance from $\Omega_q^{c}$ lies in $\Omega_q$.

It follows that the union of all cubes of side $\mathfrak{M}L^{q+1}$ at distance greater than
$R_1\mathfrak{M}L^{q}$ from $\Omega_q^{c}$ is the unique maximal admissible $\Omega_{q+1}$. The
same argument, with the non-strict inequality, applies to $\mathfrak{r}_2$. Each member is thus
fixed by the preceding one, $\mathfrak{b}$, $R_1$ and $q$. No level of the renormalisation group
enters, so $\boldsymbol{\Omega}^{\mathfrak{r}}(Z)$ is determined by $Z$, $\mathfrak{r}$ and $n$.

\emph{The remaining objects of \eqref{4.8:eq-tower-rules-a}.} The minimiser of
\cite[Theorem~1]{B-CMP102b} is a function of the tower and the datum, and
$J_n^{\mathfrak{r}}$ is given by \cite[(1.8)]{B-CMP109} in terms of $U_n^{\mathfrak{r}}$ and $L^{-n}$.
Hence both are determined by $Z$, $\mathfrak{r}$ and $n$, as a function of the datum
$M^{\boldsymbol{\cdot}}(\mathbf{U})$. That datum is formed layer by layer on the data sets
$\Lambda_p$. For $b\in\Lambda_p$ the average $\bar{\mathbf{U}}^{p}(b)$ involves $\mathbf{U}$ only
on bonds of the blocks of side $L^{p}$ containing the endpoints of $b$
(Definition~\ref{1.5:def-block-average}), and $\Lambda_p$ is a union of such blocks
(Definition~\ref{4.8:not-cube-lattices}) contained in $Z$; on $\Lambda_0$ the datum is
$\mathbf{U}$ itself. Hence the datum, and with it both maps, depend on $\mathbf{U}$ only through
$\mathbf{U}|_{Z}$.

\emph{$\mathfrak{r}_2$ satisfies the conditions of \cite[(1)]{B-CMP102b}.} We use block
$\mathfrak{M}$ and $R=R_\star'$, which is at least $R_1$ by the block-letter convention above. For
$0\le q<n$ the gap between consecutive members satisfies
\begin{equation*}
\operatorname{dist}(\Omega_q^{c},\Omega_{q+1})\ge M_1L^{q+1}=L\,R_\star'\mathfrak{M}L^{q}
>R_\star'\mathfrak{M}L^{q}.
\end{equation*}
This is the strict inequality required there.

\emph{Collar of $\mathfrak{r}_1$.} We use one fact about cubes: two points of a cube of side
$\ell$ steps are at distance at most $d\ell$.

Let $0\le q<n$, and let $x\in\Omega_q$ with
$\operatorname{dist}(x,\Omega_q^{c})>(R_1+dL)\mathfrak{M}L^{q}$. Let $\square$ be the cube of side
$\mathfrak{M}L^{q+1}$ containing $x$. Every $y\in\square$ satisfies
$\operatorname{dist}(x,y)\le dL\mathfrak{M}L^{q}$. Hence
\begin{equation*}
\operatorname{dist}(y,\Omega_q^{c})>R_1\mathfrak{M}L^{q}\qquad\text{for all }y\in\square .
\end{equation*}
By maximality $\square\subset\Omega_{q+1}$, so $x\in\Omega_{q+1}$. Equivalently, every
$x\in\Omega_q\setminus\Omega_{q+1}$ has $\operatorname{dist}(x,\Omega_q^{c})\le(R_1+dL)\mathfrak{M}L^q$.
In particular $\Omega_q^{c}\neq\emptyset$ whenever $\Omega_q\setminus\Omega_{q+1}\neq\emptyset$.

Put $D_q:=\sum_{r<q}(R_1+dL)\mathfrak{M}L^{r}$, so that $D_0=0$. We show by induction on $q$ that
every $x\in Z\setminus\Omega_q$ has $\operatorname{dist}(x,Z^c)\le D_q$. For $q=0$ the set
$Z\setminus\Omega_0$ is empty.

Assume the claim for $q$, and let $x\in Z\setminus\Omega_{q+1}$. If $x\notin\Omega_q$, the
hypothesis gives $\operatorname{dist}(x,Z^c)\le D_q\le D_{q+1}$. Otherwise
$x\in\Omega_q\setminus\Omega_{q+1}$. Choose $z\in\Omega_q^{c}$ with
$\operatorname{dist}(x,z)\le(R_1+dL)\mathfrak{M}L^q$. Either $z\in Z^{c}$, or
$z\in Z\setminus\Omega_q$ and then $\operatorname{dist}(z,Z^c)\le D_q$. In both cases the triangle
inequality gives $\operatorname{dist}(x,Z^c)\le D_{q+1}$.

It remains to bound $D_n$. Since $L\ge2$, $L\ge4d$ and $R_1\ge4d$,
\begin{align*}
D_n&=(R_1+dL)\mathfrak{M}\frac{L^{n}-1}{L-1}\le\frac{2(R_1+dL)\mathfrak{M}L^{n}}{L},\\
\frac{2(R_1+dL)}{L}&=\frac{2R_1}{L}+2d\le\frac{R_1}{2d}+\frac{R_1}{2}\le R_1 .
\end{align*}
Hence $D_n\le R_1\mathfrak{M}L^n\le R_1M_1L^n$, which is the first bound of
\eqref{4.8:eq-tower-rules-b}.

\emph{Collar of $\mathfrak{r}_2$.} Let $1\le q\le n$, and let $x\in\Omega_{q-1}$ with
$\operatorname{dist}(x,\Omega_{q-1}^{c})\ge(1+d)M_1L^{q}$. The cube of side $M_1L^q$ containing
$x$ has all its points within $dM_1L^q$ of $x$. All its points are therefore at distance at least
$M_1L^q$ from $\Omega_{q-1}^{c}$. By maximality the cube lies in $\Omega_q$, and so does $x$.

Thus every $x\in\Omega_{q-1}\setminus\Omega_q$ has
$\operatorname{dist}(x,\Omega_{q-1}^{c})<(1+d)M_1L^q$. The induction of the previous paragraph,
with these increments, shows that every point of $Z\setminus\Omega_n$ lies within
$\sum_{q=1}^{n}(1+d)M_1L^{q}$ of $Z^{c}$. Since $L\ge2$ and $d=4$,
\begin{equation*}
\sum_{q=1}^{n}(1+d)M_1L^{q}\le(1+d)M_1\frac{L^{n+1}}{L-1}\le2(1+d)M_1L^{n}\le3dM_1L^{n}.
\end{equation*}
Finally $3d\le4d\le R_1$. This gives the second bound of \eqref{4.8:eq-tower-rules-b}.

The final assertion about $M'$ follows: if $M'\ge2R_1M_1$, then both collars are at most
$R_1M_1\le\tfrac12M'$.
\end{proof}

\begin{definition}[The tower clause and the literal class]\label{4.8:def-tower-clause}
For every level $q\ge0$ we write $\xi_q=L^{-q}$. We use the cube lattices $\pi_q$, the families
$\mathbf{D}_q$ of
localisation domains and the two-layer interiors of
Definition~\ref{4.8:not-cube-lattices}, the set $\mathfrak{R}$ of tower rules and the towers
$\boldsymbol{\Omega}^{\mathfrak{r}}(Z)$ of Definition~\ref{4.8:def-tower-rules}, and
Ba{\l}aban's regularity spaces $U^c_j(X,\alpha_0,\alpha_1,\gamma_0)$ of
Definition~\ref{4.8:not-regularity-spaces}.

\emph{The tower clause.} Let $p\ge1$, $Z\in\mathbf{D}_{p-1}$, $\mathfrak{r}\in\mathfrak{R}$ and
$c>0$. Let $\mathbf{U}=U'U$, with $U$ taking values in the gauge group $G$, be given on a set
of bonds that contains the bonds of $Z$.
For $n=1,\dots,p$ we write:
\begin{itemize}
\item $U_n^{\mathfrak{r}}(Z;\cdot)$ for the minimiser in the axial gauge attached to the first
$n+1$ members $\Omega_0,\dots,\Omega_n$ of the tower $\boldsymbol{\Omega}^{\mathfrak{r}}(Z)$,
as in \cite[(1.15)]{B-CMP109};
\item $J_n^{\mathfrak{r}}(Z;\cdot)$ for its current, given by \cite[(1.8)]{B-CMP109} and
measured in level-$n$ units, as in \cite[(1.15)]{B-CMP109}.
\end{itemize}
We write $M^{\boldsymbol{\cdot}}(\mathbf{U})$ for the layered average of $\mathbf{U}$ defined
in \cite[(1.15)]{B-CMP109}; on each data set $\Lambda_q$ of the tower it equals the block
average $\bar{\mathbf{U}}^q$ of $\mathbf{U}$ at level $q$ (Definition~\ref{1.5:def-block-average}).
The configuration $\mathbf{U}=U'U$ satisfies the \emph{tower clause}
$(\mathrm{iv})^{\mathfrak{r}}_{p,Z}[c]$ if, for $n=1,\dots,p$,
\begin{equation}\label{4.8:eq-tower-clause-a}
\begin{aligned}
\bigl|\partial U_n^{\mathfrak{r}}\bigl(Z;M^{\boldsymbol{\cdot}}(\mathbf{U})\bigr)-1\bigr|
&< c\,\xi_p^2,\\
\bigl|J_n^{\mathfrak{r}}\bigl(Z;M^{\boldsymbol{\cdot}}(\mathbf{U})\bigr)\bigr|
&< c\,(L^n\xi_p)^2
\end{aligned}
\qquad\text{on } Z^{\wr p},
\end{equation}
and if the same bounds hold with $U$ in place of $\mathbf{U}$. Here $Z^{\wr p}$ is the set
of~\eqref{4.8:eq-cube-lattices-b}. The tower clause is condition (iv) of
\cite[(1.15)--(1.16)]{B-CMP109} transported to the label $(p,Z)$: the range $n=1,\dots,j$ of
that condition is taken at level $p$, the tower is the one built by the rule $\mathfrak{r}$,
the bounds are required on $Z^{\wr p}$ in place of the two-layer interior cut in $\pi_j$, and
the bounds for the $G$-valued factor $U$ are retained.

\emph{The literal class.} Let $a_0,a_1,\gamma_0>0$. The class
$U^{\mathrm{lit}}_j(X;a_0,a_1,\gamma_0)$ is Ba{\l}aban's space of \cite[(1.11)--(1.16)]{B-CMP109},
taken word for word with the rule $\mathfrak{r}_1$ and with $a_0,a_1$ in place of
$\alpha_0,\alpha_1$. In both cases below the class is defined by conditions (i)--(iii) of
\cite[(1.11)--(1.14)]{B-CMP109} together with a fourth condition, which depends on the domain.
\begin{itemize}
\item If $X\in\mathbf{D}_j$, the fourth condition consists of the
bounds \cite[(1.16)]{B-CMP109} at the constant $a_0$. These bounds are required on the
two-layer interior $\tilde{X}^{-2}$ of $X$, cut in $\pi_j$. By~\eqref{4.8:eq-cube-lattices-c},
this condition is weaker than $(\mathrm{iv})^{\mathfrak{r}_1}_{j,X}[a_0]$.
\item Suppose $U^{\mathrm{lit}}_j(Y;a_0,a_1,\gamma_0)$ is a level-$j$ space on a domain
$Y\in\mathbf{D}_{j-1}\setminus\mathbf{D}_j$. This is the situation of the domain $Y$ of
\cite{B-CMP116}, and of the window $\square_0$ at $j=k+1$. The two-layer interior is then cut
in $\pi_{j-1}$, as Ba{\l}aban does in this situation. The fourth condition therefore reads
$(\mathrm{iv})^{\mathfrak{r}_1}_{j,Y}[a_0]$: the collar is cut in the domain's own cubes.
\end{itemize}

\emph{The reading of \cite[(1.12)]{B-CMP109}.} Consider a space on a domain $X$ at level $j$
with constant $a_0$. Let $\mathsf{o}$ be the number written $O(1)$ in \cite[(1.12)]{B-CMP109},
as in Definition~\ref{4.4:def-local-data-classes}. That display involves a constant described
there as the
\begin{quote}
sufficiently large constant $B$ (it will be determined later)
\end{quote}
and outside quotations we write $B^{109}$ for it, as in
Definition~\ref{4.4:def-local-data-classes}; this constant may be enlarged as needed.
Condition \cite[(1.12)]{B-CMP109} is read as follows. It is required for every cube
$\square\subset X$ of side at most $\mathsf{o}\,LM\,L^{j}$ lattice steps: there is a gauge
transformation $u$ on $\square$ with values in $G$ such that $U^u=\exp(i\xi_jA)$ on $\square$,
where $A$ here denotes a field with values in the Lie algebra of $G$ (not the Wilson action),
and $A$ and its lattice gradient $\nabla^{\xi_j}A$ at spacing $\xi_j$ obey the single bound
\begin{equation}\label{4.8:eq-tower-clause-b}
|A|,\ |\nabla^{\xi_j}A| < \mathsf{o}\,LM\,B^{109}\,a_0 \qquad\text{on } \square .
\end{equation}
This is the reading convention $(\rho1)$ of Definition~\ref{4.5:not-reading-conventions}.
\end{definition}

\begin{definition}[The one-power hereditary regularity class]\label{4.8:def-one-power-class}
Let $j\ge 1$, and let $X\in\mathbf{D}_{j-1}$ be a localisation domain of level $j-1$ in the
sense of Definition~\ref{4.8:not-cube-lattices}. Let $a_0,a_1,\gamma_0>0$, write
$\xi_j=L^{-j}$, and let $\mathfrak{R}$ be the set of the two tower rules of
Definition~\ref{4.8:def-tower-rules}. The set of labels of $X$ at level $j$ and the one-power
tower condition are
\begin{equation}\label{4.8:eq-one-power-class-a}
\begin{gathered}
\mathfrak{T}_j(X):=\{(p,Z):\ 1\le p\le j,\ Z\in\mathbf{D}_{p-1},\ Z\subset X\},\\
(\mathrm{iv})^{\Sigma^1}:\quad
(\mathrm{iv})^{\mathfrak{r}}_{p,Z}\bigl[a_0L^{-(j-p)}\bigr]\\
\text{for every }(p,Z)\in\mathfrak{T}_j(X)\text{ and every }\mathfrak{r}\in\mathfrak{R},
\end{gathered}
\end{equation}
where $(\mathrm{iv})^{\mathfrak{r}}_{p,Z}[c]$ is the tower clause
\eqref{4.8:eq-tower-clause-a} of Definition~\ref{4.8:def-tower-clause}. The set
$\mathfrak{T}_j(X)$ is finite: $p$ takes $j$ values and, for each of them, only finitely many
unions of $\pi_{p-1}$-cubes lie in $X$.

A \emph{one-power representative} at $(j,X;a_0,a_1,\gamma_0)$ is a pair
$(\mathbf{U},\mathbf{J})$ of configurations on the bonds of $X$, together with one
decomposition $\mathbf{U}=U'U$ with $U$ taking values in the gauge group $G=\mathrm{SU}(N)$,
such that, at level $j$
on $X$ and with this single decomposition throughout, the following hold:
\begin{itemize}
\item[(i)] condition \cite[(1.11)]{B-CMP109} with constant $a_0$, that is
$|\partial U-1|<a_0\xi_j^2$ on $X$, where $\partial U$ denotes the field of plaquette
variables of $U$ (Definition~\ref{1.1:def-wilson-action}); and condition
\cite[(1.12)]{B-CMP109} in the reading
\eqref{4.8:eq-tower-clause-b}, with the bound $c_{12}LMB^{(109)}a_0$;
\item[(ii)] condition \cite[(1.13)]{B-CMP109} with constant $a_1$;
\item[(iii)] condition \cite[(1.14)]{B-CMP109} with constants $(a_0,\gamma_0)$;
\item[(iv)] condition $(\mathrm{iv})^{\Sigma^1}$ of \eqref{4.8:eq-one-power-class-a}.
\end{itemize}
The $G$-valued factor $U$ of this decomposition is the configuration that enters the second
half of each tower clause in (iv), where the bounds are required also with $U$ in place of
$\mathbf{U}$. The \emph{one-power hereditary regularity class}
$U^{\Sigma^1}_j(X;a_0,a_1,\gamma_0)$ is the union of the $G^c$-orbits, in the sense of
\cite[(1.10)]{B-CMP109}, of the one-power representatives at $(j,X;a_0,a_1,\gamma_0)$. When
$\gamma_0=a_0$ the argument $\gamma_0$ is omitted, following Ba{\l}aban's convention stated
after \cite[(1.16)]{B-CMP109}, and one writes $U^{\Sigma^1}_j(X;a_0,a_1)$.
\end{definition}

In words: the bounds of \cite[(1.16)]{B-CMP109} are imposed, under both tower rules, for the
tower of every localisation domain $Z\subset X$ of every level $p-1<j$, on the set $Z^{\wr p}$
obtained from $Z$ by removing the two layers of $\pi_{p-1}$-cubes closest to its boundary;
see \eqref{4.8:eq-cube-lattices-b}. At $p=j$ the constant is Ba{\l}aban's constant $a_0$: the
pair $(j,X)$ itself belongs to $\mathfrak{T}_j(X)$, with the two layers removed in
$\pi_{j-1}$. When $X\in\mathbf{D}_j$, so that Ba{\l}aban's set $\tilde{X}^{-2}$, obtained
from $X$ by removing the two layers of $\pi_j$-cubes closest to its boundary, is defined,
these two layers of $\pi_{j-1}$-cubes lie inside the two layers of $\pi_j$-cubes, so that
$X^{\wr j}\supset\tilde{X}^{-2}$ and at $(j,X)$ the bounds are required on a set containing
$\tilde{X}^{-2}$. At $p<j$ the
constant is $a_0L^{-(j-p)}$, and
the tower clause measures the bounds in level-$p$ units, so that the condition carries one
power of the scale ratio $L^{-(j-p)}$, whereas Ba{\l}aban's bound $a_0\xi_j^2$, rewritten in
level-$p$ units by means of $\xi_j=L^{-(j-p)}\xi_p$, reads
\begin{equation*}
a_0\xi_j^2=a_0L^{-2(j-p)}\xi_p^2
\end{equation*}
and carries two. The single power is exactly what the far-case proposition supplies, in the
bound \eqref{4.8:eq-far-case-b}. The
analogous condition with two powers, $a_0L^{-2(j-p)}$ in place of $a_0L^{-(j-p)}$, is also
hereditary, that is, reproduced by the inductive step, but it is not supplied at the rim
sub-towers.

\begin{remark}[Comparison with the literal and two-power classes]\label{4.8:rem-class-comparison}
Let $j\ge 1$ and $X\in\mathbf{D}_{j-1}$. Let $\mathfrak{T}_j(X)$, the rule set $\mathfrak{R}$, the literal class $U^{\mathrm{lit}}_j$ and the one-power class $U^{\Sigma^1}_j$ be as in Definition~\ref{4.8:def-tower-clause} and Definition~\ref{4.8:def-one-power-class}.

Let $\Sigma^2$, the \emph{two-power twin} of $\Sigma^1$, be defined word for word as $\Sigma^1$ in Definition~\ref{4.8:def-one-power-class}, with one change: the clause \eqref{4.8:eq-one-power-class-a} is replaced by
\begin{equation*}
(\mathrm{iv})^{\mathfrak{r}}_{p,Z}\bigl[a_0L^{-2(j-p)}\bigr]
\quad\text{for every }(p,Z)\in\mathfrak{T}_j(X)\text{ and every }\mathfrak{r}\in\mathfrak{R}.
\end{equation*}
The set of labels $\mathfrak{T}_j(X)$ and the rule set $\mathfrak{R}$ are the same as for $\Sigma^1$. The class $U^{\Sigma^2}_j(X;a_0,a_1,\gamma_0)$ is the union of the $G^c$-orbits of the $\Sigma^2$-representatives, exactly as for $U^{\Sigma^1}_j$.

Two further two-power classes enter; their definitions are not repeated here. The first is the two-power class with Ba{\l}aban's published constants, which we write $U^{[1]}_j(X;\cdot)$. Of it we use only that its tower clause carries the two-power factor $a_0L^{-2(j-p)}$, Ba{\l}aban's constants as published, and the nesting rule $\mathfrak{r}_2$ of \cite[(2.13)]{B-CMP119}.

The second is the two-power class on which the sixth-family lemma is stated, which we write $U^{[2]}_j(X;\cdot)$. Of it we use only the following. The tower clause with the two-power factor $a_0L^{-2(j-p)}$, which is Ba{\l}aban's constant $a_0\xi_j^2$ written in level-$p$ units, is imposed for every triple consisting of a level $p$, a domain $X'$ and a value of the structure label of that class, in the range fixed in the sixth-family lemma. The clause is localised to $X'$ in the manner of Ba{\l}aban's regularity conditions \cite{B-CMP109}, and it is imposed also for the unitary part $U$.

In the last two inclusions of \eqref{4.8:eq-class-comparison-a} below, both two-power classes $U^{[1]}_j$ and $U^{[2]}_j$ are taken under the following reading:
\begin{itemize}
\item[(R)] clause (iv) is read at every label $(p,Z)$ of the tower under both rules $\mathfrak{r}_1$ (clause (iv) of \cite{B-CMP109}) and $\mathfrak{r}_2$ (\cite[(2.13)]{B-CMP119}), each clause at label $(p,Z)$ being imposed on $Z^{\wr p}$, that is, on $Z$ with two layers of $\pi_{p-1}$-cubes removed.
\end{itemize}

Then, the dots standing for the common arguments $(a_0,a_1,\gamma_0)$,
\begin{equation}\label{4.8:eq-class-comparison-a}
\begin{split}
U^{\mathrm{lit}}_j(X;\cdot)&\supset U^{\Sigma^1}_j(X;\cdot)\supset U^{\Sigma^2}_j(X;\cdot)\\
&\supseteq U^{[1]}_j(X;\cdot)\supseteq U^{[2]}_j(X;\cdot).
\end{split}
\end{equation}
In \eqref{4.8:eq-class-comparison-a} the symbols $\supset$ and $\supseteq$ both denote inclusion, not necessarily strict; $\supseteq$ marks the two inclusions that are asserted only under the reading (R). The first two inclusions are inclusions of clause sets and are proved below. The last two hold for the two two-power classes just described, under the reading (R).

Read literally, with the clause imposed at one label and under one rule, the last two inclusions have no ground: neither rule's clause is known to imply the other's, a derivation requiring locality of the tower at complex points, which is not proved.

Consequently every bound proved for all configurations in the larger class $U^{\Sigma^1}_j(X;\cdot)$ holds, by inclusion, for all configurations in the two-power classes read as in (R); no passage in the converse direction is available. The class $U^{\Sigma^1}_j$ is a third class, defined differently from both $U^{[1]}_j$ and $U^{[2]}_j$.
\end{remark}

\begin{proof}
We prove the first two inclusions in \eqref{4.8:eq-class-comparison-a}. In each case we show that every representative of the smaller class, with its decomposition $\mathbf{U}=U'U$, satisfies every clause that defines the larger class.

\emph{First inclusion.} Let $(\mathbf{U},\mathbf{J})$ be a $\Sigma^1$-representative at $(j,X;a_0,a_1,\gamma_0)$. The pair $(p,Z)=(j,X)$ lies in $\mathfrak{T}_j(X)$, since $1\le j\le j$, $X\in\mathbf{D}_{j-1}$ and $X\subset X$. At this label the factor is
\begin{equation*}
a_0L^{-(j-j)}=a_0L^0=a_0 .
\end{equation*}
Hence \eqref{4.8:eq-one-power-class-a} contains the clause $(\mathrm{iv})^{\mathfrak{r}_1}_{j,X}[a_0]$ of \eqref{4.8:eq-tower-clause-a}.

The literal class imposes \cite[(1.16)]{B-CMP109} at $a_0$ on the two-layer interior $\tilde X^{-2}$ cut in $\pi_j$. By the comparison of cube lattices \eqref{4.8:eq-cube-lattices-c}, that region is contained in $X^{\wr j}$, the interior cut in $\pi_{j-1}$ on which $(\mathrm{iv})^{\mathfrak{r}_1}_{j,X}[a_0]$ is imposed. So this condition of the literal class is implied by $(\mathrm{iv})^{\mathfrak{r}_1}_{j,X}[a_0]$.

Clauses (ii) at $a_1$ and (iii) at $(a_0,\gamma_0)$, at level $j$ on $X$, are common to both classes. The same holds for the part \cite[(1.11)]{B-CMP109} at $a_0$ of clause (i).

The remaining part of (i) is \cite[(1.12)]{B-CMP109}. In Definition~\ref{4.8:def-one-power-class} it is imposed, under the reading fixed there, at $c_{12}LMB^{(109)}a_0$. It is imposed on a family of cubes that includes the cubes of \cite[(1.12)]{B-CMP109}, so it implies condition \cite[(1.12)]{B-CMP109} as stated there.

Thus every clause of the literal class holds, and $U^{\mathrm{lit}}_j(X;\cdot)\supset U^{\Sigma^1}_j(X;\cdot)$.

\emph{Second inclusion.} The classes $\Sigma^1$ and $\Sigma^2$ have the same clauses (i)--(iii), the same labels $\mathfrak{T}_j(X)$ and the same rules $\mathfrak{R}$. They differ only in the constant of the tower clause.

For $(p,Z)\in\mathfrak{T}_j(X)$ we have $j-p\ge 0$, and $L\ge 1$. Hence
\begin{equation*}
a_0L^{-2(j-p)}\le a_0L^{-(j-p)} .
\end{equation*}
By \eqref{4.8:eq-tower-clause-a}, the clause $(\mathrm{iv})^{\mathfrak{r}}_{p,Z}[c]$ consists of strict upper bounds whose right-hand sides, $c\,\xi_p^2$ and $c\,(L^n\xi_p)^2$, increase with $c$. So the clause at $c=a_0L^{-2(j-p)}$ implies the clause at $c=a_0L^{-(j-p)}$, for every label and every rule.

Every $\Sigma^2$-representative is therefore a $\Sigma^1$-representative. Both classes are unions of the $G^c$-orbits of their representatives, so $U^{\Sigma^1}_j(X;\cdot)\supset U^{\Sigma^2}_j(X;\cdot)$.
\end{proof}

\begin{lemma}[Elementary properties of the one-power class]\label{4.8:lem-class-properties}
Let $j\ge 1$ and $X\in\mathbf{D}_{j-1}$. Let $U^{\Sigma^1}_j(X;a_0,a_1,\gamma_0)$ be the one-power hereditary regularity class of Definition~\ref{4.8:def-one-power-class}, and let $U^{\mathrm{lit}}_j(X;\cdot)$ be the literal class, the class defined by Ba{\l}aban's conditions (i)--(iv) as printed in \cite{B-CMP109}, which is compared with the one-power class in Remark~\ref{4.8:rem-class-comparison}. A dot in place of the parameters means that the assertion holds at every parameter triple. For a set $S$ of pairs on the bonds of $X$ and for $X'\subset X$, write $S|_{X'}$ for the set of restrictions of the members of $S$ to the bonds of $X'$.
\begin{itemize}
\item[(a)] $U^{\Sigma^1}_j(X;\cdot)\subset U^{\mathrm{lit}}_j(X;\cdot)$.
\item[(b)] If $(U'U,\mathbf{J})$ is a $\Sigma^1$-representative at $(j,X;a_0,a_1,\gamma_0)$, then so is its real point $(U,0)$.
\item[(c)] The class $U^{\Sigma^1}_j(X;\cdot)$ is cut out of the set of pairs satisfying conditions (i)--(iii) of Definition~\ref{4.8:def-one-power-class} by finitely many strict inequalities on functions of the restrictions $\mathbf{U}|_{Z}$, one family for each label $(p,Z)\in\mathfrak{T}_j(X)$ of \eqref{4.8:eq-one-power-class-a}. These functions are continuous, and analytic where they are defined \cite[Proposition~9]{B-CMP102b}; they are the functions used by clause (iv), \cite[(1.16)]{B-CMP109}, of Ba{\l}aban's regularity spaces. Hence $U^{\Sigma^1}_j(X;\cdot)$ is relatively open in that set.
\item[(d$_1$)] $U^{\Sigma^1}_j(X;a_0,a_1,\gamma_0)$ increases, with respect to inclusion, with each of $a_0$, $a_1$, $\gamma_0$ and $B^{109}$, the constant of Definition~\ref{4.4:def-local-data-classes}.
\item[(d$_2$)] If $X'\subset X$ and $X'\in\mathbf{D}_{j-1}$, then
\begin{equation*}
U^{\Sigma^1}_j(X;a_0,a_1,\gamma_0)\big|_{X'}\subset U^{\Sigma^1}_j(X';a_0,a_1,\gamma_0).
\end{equation*}
\end{itemize}
\end{lemma}

\begin{proof}
(a) This is the first inclusion of \eqref{4.8:eq-class-comparison-a}, proved in Remark~\ref{4.8:rem-class-comparison}.

(b) Condition (i) concerns only the real factor $U$, so it holds for $(U,0)$ because it holds for $(U'U,\mathbf{J})$. Let $A'$ denote the algebra-valued field of the complex factor $U'$ that is bounded in condition (ii). For the real point this factor is the identity, so $A'=0$, and condition (ii) at $a_1$ holds. Condition (iii) at $(a_0,\gamma_0)$, evaluated at the pair $(U,0)$, is the bound \cite[(1.11)]{B-CMP109} at $a_0$. That bound is part of condition (i), which has just been verified.

Clause $(\mathrm{iv})^{\Sigma^1}$ for $(U'U,\mathbf{J})$ has a half that concerns the real factor $U$. This half is, word for word, the whole of clause $(\mathrm{iv})^{\Sigma^1}$ for the configuration $\mathbf{U}:=U$. Hence $(U,0)$ satisfies (i)--(iii) and $(\mathrm{iv})^{\Sigma^1}$, and it is a $\Sigma^1$-representative at $(j,X;a_0,a_1,\gamma_0)$.

(c) Once (i)--(iii) hold, the conditions that remain are those of clause $(\mathrm{iv})^{\Sigma^1}$. They are indexed by the triples consisting of:
\begin{itemize}
\item a label $(p,Z)\in\mathfrak{T}_j(X)$;
\item a rule $\mathfrak{r}\in\mathfrak{R}$;
\item a depth $n\le p$ of the tower $\boldsymbol{\Omega}^{\mathfrak{r}}(Z)$.
\end{itemize}
The index set $\mathfrak{T}_j(X)\times\mathfrak{R}\times\{n\le p\}$ is finite. The set $\mathfrak{T}_j(X)$ is finite by Definition~\ref{4.8:def-one-power-class}, $\mathfrak{R}$ consists of the two rules $\mathfrak{r}_1,\mathfrak{r}_2$, and for each label only finitely many depths $n\le p$ occur.

For each index, the corresponding condition is a strict inequality. It requires that the supremum, over finitely many plaquettes or bonds, of the functions of $\mathbf{U}|_{Z}$ used by clause (iv), \cite[(1.16)]{B-CMP109}, of Ba{\l}aban's regularity spaces be less than the given constant. Each of these functions is continuous, and analytic where defined \cite[Proposition~9]{B-CMP102b}. A supremum over finitely many continuous functions is a maximum and is continuous. A strict upper bound on a continuous function defines an open set, and a finite intersection of open sets is open. Hence the class is cut out of the set defined by (i)--(iii) by finitely many strict inequalities, and it is relatively open in that set.

(d$_1$) Every defining inequality is an upper bound by a radius, and each radius increases with $a_0$, $a_1$, $\gamma_0$ and $B^{109}$. This applies to conditions (i)--(iii), including \cite[(1.12)]{B-CMP109} at the radius $c_{12}LMB^{109}a_0$, and to the conditions $(\mathrm{iv})^{\mathfrak{r}}_{p,Z}[a_0L^{-(j-p)}]$ of $(\mathrm{iv})^{\Sigma^1}$. Enlarging any of these four parameters therefore weakens each condition. As a result, every $\Sigma^1$-representative at the smaller parameters is a $\Sigma^1$-representative at the larger ones, and the same inclusion holds for the unions of their $G^c$-orbits.

(d$_2$) This is Theorem~\ref{4.8:thm-restriction-units}, in the case of equal levels, that is, at scale ratio $\mu=L^{0}=1$. The proof of that theorem uses the present lemma only through (d$_1$), so there is no circularity.
\end{proof}

\begin{theorem}[Restriction across levels as a change of units]\label{4.8:thm-restriction-units}
Let $1\le i'\le i$ and put $\mu:=L^{-(i-i')}$. Let $Y\in\mathbf{D}_{i-1}$ and $X\in\mathbf{D}_{i'-1}$ with
$X\subset Y$. Let $(\mathbf{U},\mathbf{J})$ be a $\Sigma^1$-representative at $(i,Y;a_0,a_1,\gamma_0)$
in the sense of Definition~\ref{4.8:def-one-power-class}, with decomposition $\mathbf{U}=U'U$. Then:
\begin{enumerate}
\item[(a)] $(\mathbf{U},\mathbf{J})\restriction X$ is a $\Sigma^1$-representative at
$(i',X;\mu a_0,\mu a_1,\gamma_0)$;
\item[(b)] under the re-normalisation $\mathfrak{s}=\mu^3$ of the current in \cite[(3.12)]{B-CMP109},
$(\mathbf{U},\mu^3\mathbf{J})\restriction X$ is a $\Sigma^1$-representative at
$(i',X;\mu a_0,\mu a_1,\mu^3\gamma_0)$.
\end{enumerate}
In both cases the decomposition is $\mathbf{U}\restriction X=(U'\restriction X)(U\restriction X)$, with the
same $G$-valued factor $U$, and the gauge transformations entering the conditions at
$(i',X)$ are restrictions of those entering them at $(i,Y)$. Consequently
\begin{equation*}
\begin{gathered}
U^{\Sigma^1}_i(Y;a_0,a_1,\gamma_0)\restriction X\subset U^{\Sigma^1}_{i'}(X;a_0',a_1',\gamma_0')\\
\text{whenever } \mu a_0\le a_0',\ \mu a_1\le a_1',\ \gamma_0\le\gamma_0' .
\end{gathered}
\end{equation*}
No condition is imposed on $M$, on the position of $X$ in $Y$, on the shape or thickness of $Y$
or of $X$, or on $i-i'$.
\end{theorem}

\begin{proof}
The bonds of $X$ are bonds of $Y$, and the configuration is not changed; what changes is the
level at which it is read, hence the lattice spacing in which algebra-valued fields and
difference quotients are expressed. We first record this change of units, then verify the
conditions of Definition~\ref{4.8:def-one-power-class} one by one.

\emph{Change of units.} Since $\xi_m=L^{-m}$, we have $\xi_i=L^{-(i-i')}\xi_{i'}=\mu\,\xi_{i'}$.
For an algebra-valued field $A$ on bonds put $A^\flat:=\mu A$; then
$e^{\mathrm{i}\xi_iA}=e^{\mathrm{i}\xi_{i'}A^\flat}$. For a spacing $\xi$ write $\nabla^{\xi}$ for a difference
quotient at spacing $\xi$, plain or covariant. A difference quotient at spacing $\xi_{i'}$
divides the same difference, taken between the same lattice points and, in the covariant
case, with the same parallel transporter, by $\xi_{i'}=\xi_i/\mu$ instead of $\xi_i$; hence
$\nabla^{\xi_{i'}}=\mu\nabla^{\xi_i}$ in both cases. Since $i-i'\ge0$ we have $0<\mu\le1$.
Collecting these facts, for every algebra-valued field $A$ on bonds and every $b\ge0$ we have
\begin{equation}\label{4.8:eq-restriction-units-a}
\begin{gathered}
\xi_i=\mu\,\xi_{i'},\qquad A^\flat:=\mu A,\qquad
e^{\mathrm{i}\xi_iA}=e^{\mathrm{i}\xi_{i'}A^\flat},\qquad
\nabla^{\xi_{i'}}=\mu\,\nabla^{\xi_i},\\
|A^\flat|=\mu|A|,\qquad
|\nabla^{\xi_{i'}}A^\flat|=\mu^2|\nabla^{\xi_i}A|\le\mu|\nabla^{\xi_i}A|,\\
b\,\xi_i^2=\mu^2b\,\xi_{i'}^2\le(\mu b)\,\xi_{i'}^2 .
\end{gathered}
\end{equation}

\emph{Condition \textup{(i)}.} We keep the same factor $U$. The curvature bound at level $i$ on
$Y$ holds on every plaquette of $X$, and by \eqref{4.8:eq-restriction-units-a} with $b=a_0$,
\begin{equation*}
|\partial U-1|<a_0\xi_i^2\le(\mu a_0)\,\xi_{i'}^2\quad\text{on } X,
\end{equation*}
which is the curvature bound at level $i'$ with constant $\mu a_0$. For the gauge bound of
\cite[(1.12)]{B-CMP109}, read under the cube-size convention \eqref{4.8:eq-tower-clause-b}, let
$\square\subset X$ be a cube of side at most $c_{12}LM\cdot L^{i'}$ steps. Since $X\subset Y$ and
$L^{i'}\le L^{i}$, $\square$ is a cube contained in $Y$ of side at most $c_{12}LM\cdot L^{i}$ steps;
it is therefore itself one of the cubes to which the hypothesis at $(i,Y)$ applies, whatever
the thickness of $Y$. That hypothesis provides a gauge transformation $u$ on $\square$ with
\begin{equation*}
U^u=e^{\mathrm{i}\xi_iA}\ \text{on }\square,\qquad |A|,\ |\nabla^{\xi_i}A|<c_{12}LMB^{(109)}a_0 .
\end{equation*}
By \eqref{4.8:eq-restriction-units-a}, $U^u=e^{\mathrm{i}\xi_{i'}A^\flat}$ on $\square$, and
\begin{equation*}
|A^\flat|=\mu|A|<c_{12}LMB^{(109)}(\mu a_0),\qquad
|\nabla^{\xi_{i'}}A^\flat|\le\mu|\nabla^{\xi_i}A|<c_{12}LMB^{(109)}(\mu a_0),
\end{equation*}
which is the gauge bound at level $i'$ on $X$ with constant $\mu a_0$, witnessed by the same
$u$.

\emph{Condition \textup{(ii)}.} Write $U'=e^{\mathrm{i}\xi_iA'}$ as in condition (ii) at level $i$. By
\eqref{4.8:eq-restriction-units-a}, $U'=e^{\mathrm{i}\xi_{i'}\mu A'}$, and the same relations convert the
bounds at $a_1$ on $A'$ into the corresponding bounds at $\mu a_1$ on $\mu A'$; this is
condition (ii) at level $i'$ on $X$ with constant $\mu a_1$.

\emph{Condition \textup{(iii)}.} Its bounds at $a_0$ pass to bounds at $\mu a_0$ exactly as in
condition (i). Its bound $|\mathbf{J}|<\gamma_0$ on the bonds of $Y$ holds on the bonds of $X$,
which gives (a); and $|\mu^3\mathbf{J}|<\mu^3\gamma_0$, which gives (b).

\emph{Condition $(\mathrm{iv})^{\Sigma^1}$.} By the definition of the label sets in
\eqref{4.8:eq-one-power-class-a},
\begin{equation*}
\mathfrak{T}_{i'}(X)\subset\mathfrak{T}_i(Y),
\end{equation*}
since every $(p,Z)\in\mathfrak{T}_{i'}(X)$ has $1\le p\le i'\le i$, $Z\in\mathbf{D}_{p-1}$ and
$Z\subset X\subset Y$. The set $\mathfrak{R}$ of rules is the same at both levels. By
\eqref{4.8:eq-tower-rules-a}, the clause \eqref{4.8:eq-tower-clause-a} at label $(p,Z)$ under a
rule $\mathfrak{r}$ is one and the same statement about $\mathbf{U}\restriction Z$ and
$U\restriction Z$, whichever class it is quoted in; and since $Z\subset X$ these are the
restrictions to $Z$ of $\mathbf{U}\restriction X$ and $U\restriction X$. At $(i,Y)$ the
clause is imposed with constant
\begin{equation*}
c=a_0L^{-(i-p)}=L^{-(i-i')}a_0\,L^{-(i'-p)}=(\mu a_0)\,L^{-(i'-p)},
\end{equation*}
which is the constant that $(\mathrm{iv})^{\Sigma^1}$ requires at level $i'$ on $X$ with first
constant $\mu a_0$, in both (a) and (b). Hence
$(\mathrm{iv})^{\mathfrak{r}}_{p,Z}[(\mu a_0)L^{-(i'-p)}]$ holds for every
$(p,Z)\in\mathfrak{T}_{i'}(X)$ and every $\mathfrak{r}\in\mathfrak{R}$. The clause is a condition on
$\mathbf{U}\restriction Z$ and $U\restriction Z$ and does not involve the current, so it is
unaffected by the re-normalisation $\mathbf{J}\mapsto\mu^3\mathbf{J}$. This proves (a)
and (b).

\emph{Orbits.} A gauge transformation acts on each bond through its values at the sites of
that bond, so restriction to the bonds of $X$ commutes with the action:
\begin{equation*}
\bigl((\mathbf{U},\mathbf{J})^u\bigr)\restriction X=\bigl((\mathbf{U},\mathbf{J})\restriction X\bigr)^{u\restriction X}.
\end{equation*}
Hence the restriction to $X$ of an element of the orbit of a $\Sigma^1$-representative at
$(i,Y;a_0,a_1,\gamma_0)$ lies in the orbit of a $\Sigma^1$-representative at
$(i',X;\mu a_0,\mu a_1,\gamma_0)$ by (a), so
\begin{equation*}
U^{\Sigma^1}_i(Y;a_0,a_1,\gamma_0)\restriction X\subset U^{\Sigma^1}_{i'}(X;\mu a_0,\mu a_1,\gamma_0).
\end{equation*}

\emph{Constants.} By clause (d$_1$) of Lemma~\ref{4.8:lem-class-properties}, the class
$U^{\Sigma^1}_{i'}(X;\cdot)$ increases with each of its three constants; for $\mu a_0\le a_0'$,
$\mu a_1\le a_1'$ and $\gamma_0\le\gamma_0'$ it is therefore contained in
$U^{\Sigma^1}_{i'}(X;a_0',a_1',\gamma_0')$, which gives the stated inclusion. No step used a
condition on $M$, on the position of $X$ in $Y$, on the shape or thickness of $Y$ or of $X$, or
on $i-i'$.
\end{proof}

\begin{corollary}[The inclusion sentences on the one-power class]\label{4.8:cor-inclusion-sentences}
Let $\beta$ and $\alpha_0,\alpha_1$ be the constant and radii of Ba{\l}aban's regularity spaces
(Notation~\ref{4.8:not-regularity-spaces}). The following four statements hold. The first three
are Ba{\l}aban's sentences in \cite[p.~9]{B-CMP116}; in \cite[p.~7]{B-CMP116} together with
\cite[(3.40)]{B-CMP109}; and in \cite[p.~15]{B-CMP116}. They are used here as stated in those
places, with Ba{\l}aban's space $U^c$ replaced throughout by the one-power class $U^{\Sigma^1}$ of
Definition~\ref{4.8:def-one-power-class}. The fourth is the running form of the first at the
radii of \cite[(2.28)]{B-CMP119}. All levels are at least $1$, as in
Theorem~\ref{4.8:thm-restriction-units}.
\begin{enumerate}
\item[(a)] Let $L\ge 2(1+\beta)$, $1\le j\le k$, $X\in\mathbf{D}_j$, $Y\in\mathbf{D}_k$
and $X\subset Y$. Then
\begin{equation}\label{4.8:eq-inclusion-sentences-a}
\begin{gathered}
(\mathbf{U},\mathbf{J})\ \text{a $\Sigma^1$-representative at}\
\bigl(k+1,Y;(1+\beta)\alpha_0,(1+\beta)\alpha_1,\alpha_0\bigr)\\
\Longrightarrow\ (\mathbf{U},\mathbf{J})\restriction X\ \text{a $\Sigma^1$-representative at}\
\bigl(j,X;\tfrac12\alpha_0,\tfrac12\alpha_1,\alpha_0\bigr).
\end{gathered}
\end{equation}
The same pair, with the same decomposition $\mathbf{U}=U'U$, satisfies conditions (i)--(iii)
of Definition~\ref{4.8:def-one-power-class}
at level $k+1$ on $Y$ with the constants $(1+\beta)\alpha_0,(1+\beta)\alpha_1,\alpha_0$. Hence a
single representative serves both members of the intersection defining the space
\cite[(3.16)]{B-CMP109}.
\item[(b)] Let $Y\in\mathbf{D}_k$ and $\square_0\in\mathbf{D}_k$ with
$\square_0\subset Y$. The case in view is the window $\square_0=\tilde{\square}^{5}$ of side
$12M$ of a cube $\square$ of Ba{\l}aban's cover in \cite[\S3]{B-CMP109}, that is, of a union of
$2^d$ neighbouring cubes of $\pi_k$. Every $\Sigma^1$-representative at
$(k+1,Y;(1+\beta)\alpha_0,\cdot,\cdot)$ satisfies
\begin{equation}\label{4.8:eq-inclusion-sentences-b}
(\mathrm{iv})^{\mathfrak{r}}_{k+1,\square_0}\bigl[(1+\beta)\alpha_0\bigr]
\quad\text{on } \square_0^{\wr(k+1)}\ \text{ for both } \mathfrak{r}\in\mathfrak{R}.
\end{equation}
For the window, $\square_0^{\wr(k+1)}=\tilde{\square}^{3}$. Since $1+\beta<1+2\beta$, this
implies, by Lemma~\ref{4.8:lem-class-properties}(d$_1$), the window clause assumed in
\cite[(3.40)]{B-CMP109} at the constant $(1+2\beta)\alpha_0$, with room $\beta$ to spare. It
holds both for Ba{\l}aban's window tower
$\mathfrak{r}_1$ (his tower taken ``as in the condition (iv)'' \cite[\S3]{B-CMP109}) and for
the tower $\mathfrak{r}_2$ (Definition~\ref{4.8:def-tower-rules}).
\item[(c)] Let $Y\in\mathbf{D}_k$, $X\in\mathbf{D}_{k+1}$ and $Y\subset X$. Then
\begin{equation}\label{4.8:eq-inclusion-sentences-c}
\begin{aligned}
U^{\Sigma^1}_{k+1}(X;\alpha_0,\alpha_1)\restriction Y
&\subset U^{\Sigma^1}_{k+1}(Y;\alpha_0,\alpha_1)\\
&\subset U^{\Sigma^1}_{k+1}\bigl(Y;(1+\beta)\alpha_0,(1+\beta)\alpha_1,\alpha_0\bigr).
\end{aligned}
\end{equation}
Here $U^{\Sigma^1}_{k+1}(\cdot\,;\alpha_0,\alpha_1)$ is the two-argument form of the space, as
for Ba{\l}aban's spaces in Notation~\ref{4.8:not-regularity-spaces}: it denotes
$U^{\Sigma^1}_{k+1}(\cdot\,;\alpha_0,\alpha_1,\alpha_0)$.
\item[(d)] Let $\alpha_{0,i},\alpha_{1,i}$ be the radii of \cite[(2.28)]{B-CMP119} at
level $i$, and write $\alpha_{\cdot,i}$ for either of them. Let $\beta_s$ be the constant by
which these radii are enlarged at level $i$ in the large-field part of the coupling increment,
and $\beta_0$ the constant of that statement for which
$\alpha_{\cdot,i}/\alpha_{\cdot,i-1}\le 2(1+\beta_0)$ at consecutive levels. Let $1\le i'\le i-1$
and $\mathfrak{n}:=i-1-i'$. Let $X\in\mathbf{D}_{i'}$,
$Y\in\mathbf{D}_{i-1}$ and $X\subset Y$. Put $\mathfrak{s}:=L^{-3(\mathfrak{n}+1)}$. Assume
the two inequalities
\begin{equation}\label{4.8:eq-inclusion-sentences-1}
(1+\beta_s)\,\frac{\alpha_{\cdot,i}}{\alpha_{\cdot,i'}}\,L^{-(\mathfrak{n}+1)}\le\frac12,
\qquad \mathfrak{s}\,\alpha_{0,i}\le\alpha_{0,i'},
\end{equation}
the first for $\alpha_{\cdot,i}$ standing for either of $\alpha_{0,i}$, $\alpha_{1,i}$
(with the same index in numerator and denominator). Then
\begin{equation}\label{4.8:eq-inclusion-sentences-d}
\begin{aligned}
(\mathbf{U},\mathbf{J})&\in U^{\Sigma^1}_{i}
\bigl(Y;(1+\beta_s)\alpha_{0,i},(1+\beta_s)\alpha_{1,i},\alpha_{0,i}\bigr)\\
&\Longrightarrow\ (\mathbf{U},\mathfrak{s}\mathbf{J})\restriction X\in U^{\Sigma^1}_{i'}
\bigl(X;\tfrac12\alpha_{0,i'},\tfrac12\alpha_{1,i'},\alpha_{0,i'}\bigr).
\end{aligned}
\end{equation}
Under the condition $L\ge 4(1+\beta_0)(1+2\beta_s)$ of the large-field part of the coupling
increment, the inequalities \eqref{4.8:eq-inclusion-sentences-1} are derived in the proof of
that statement: for $\mathfrak{n}=0$ from the ratio bound $2(1+\beta_0)$, and for
$\mathfrak{n}\ge1$ from the ratio bound of Lemma~\ref{4.4:lem-radii-ratio}, which grows like
$(1+\mathfrak{n})^{1/2}$ and is compensated by the factor $L^{-(\mathfrak{n}+1)}$.
\end{enumerate}
\end{corollary}

\begin{proof}
Statements (a), (c) and (d) rest on the restriction theorem,
Theorem~\ref{4.8:thm-restriction-units}; statement (b) is read off the definition of the class.
The theorem is applied at the levels $i\ge i'$ and the ratio $\mu=L^{-(i-i')}$, and it is
combined with
the monotonicity in the constants, Lemma~\ref{4.8:lem-class-properties}(d$_1$). Since every
condition of Definition~\ref{4.8:def-one-power-class} is an upper bound by the corresponding
constant, a $\Sigma^1$-representative at given constants is one at any larger constants; this is
the content of Lemma~\ref{4.8:lem-class-properties}(d$_1$), stated there for the classes.
In each application the lower-level domain has to lie in
$\mathbf{D}_{i'-1}$, as Theorem~\ref{4.8:thm-restriction-units} requires. We use for this that
a domain of $\mathbf{D}_{i'}$ is a domain of $\mathbf{D}_{i'-1}$, by the nesting of the cube
lattices of Notation~\ref{4.8:not-cube-lattices}.

(a). Let $(\mathbf{U},\mathbf{J})$ be a $\Sigma^1$-representative at
$(k+1,Y;(1+\beta)\alpha_0,(1+\beta)\alpha_1,\alpha_0)$. We apply
Theorem~\ref{4.8:thm-restriction-units} with $i=k+1$ and $i'=j$. Its hypotheses hold:
$1\le j\le k+1$, $Y\in\mathbf{D}_k=\mathbf{D}_{i-1}$, $X\in\mathbf{D}_{j-1}$ and $X\subset Y$.
Here $\mu=L^{-(k+1-j)}\le L^{-1}$, because $j\le k$. The theorem gives that
$(\mathbf{U},\mathbf{J})\restriction X$ is a $\Sigma^1$-representative at
\begin{equation*}
\bigl(j,X;\mu(1+\beta)\alpha_0,\mu(1+\beta)\alpha_1,\alpha_0\bigr),
\end{equation*}
with the same group-valued factor $U$. Since $L\ge 2(1+\beta)$,
\begin{equation*}
\mu(1+\beta)\le L^{-1}(1+\beta)\le\tfrac12 .
\end{equation*}
Hence $\mu(1+\beta)\alpha_0\le\frac12\alpha_0$ and $\mu(1+\beta)\alpha_1\le\frac12\alpha_1$,
while the current constant $\alpha_0$ is unchanged. Lemma~\ref{4.8:lem-class-properties}(d$_1$)
then gives \eqref{4.8:eq-inclusion-sentences-a}.

The current constant $\alpha_0$ is also admissible under the other normalisation of the current.
If the current is renormalised by $\mu^3$, the second form of
Theorem~\ref{4.8:thm-restriction-units} gives that $(\mathbf{U},\mu^3\mathbf{J})\restriction X$
is a $\Sigma^1$-representative at
\begin{equation*}
\bigl(j,X;\mu(1+\beta)\alpha_0,\mu(1+\beta)\alpha_1,\mu^3\alpha_0\bigr).
\end{equation*}
Since $\mu^3\alpha_0\le\alpha_0$, Lemma~\ref{4.8:lem-class-properties}(d$_1$) shows that
$(\mathbf{U},\mu^3\mathbf{J})\restriction X$ is again a $\Sigma^1$-representative at
$(j,X;\tfrac12\alpha_0,\tfrac12\alpha_1,\alpha_0)$.

Conditions (i)--(iii) at level $k+1$ on $Y$, with the constants
$(1+\beta)\alpha_0,(1+\beta)\alpha_1,\alpha_0$ and for the given decomposition
$\mathbf{U}=U'U$, are part of the definition of a $\Sigma^1$-representative at
$(k+1,Y;(1+\beta)\alpha_0,(1+\beta)\alpha_1,\alpha_0)$. Definition~\ref{4.8:def-one-power-class}
therefore supplies them directly. These conditions are the second member of the intersection
\cite[(3.16)]{B-CMP109}, and \eqref{4.8:eq-inclusion-sentences-a} supplies the first. Both come
from one and the same representative.

(b). Since $\square_0\in\mathbf{D}_k=\mathbf{D}_{(k+1)-1}$, $\square_0\subset Y$ and
$p=k+1$ lies in the range $1\le p\le k+1$, the label $(k+1,\square_0)$ belongs to the set
$\mathfrak{T}_{k+1}(Y)$ of
\eqref{4.8:eq-one-power-class-a}. Let $(\mathbf{U},\mathbf{J})$ be a $\Sigma^1$-representative
at $(k+1,Y;(1+\beta)\alpha_0,\cdot,\cdot)$. Its clause $(\mathrm{iv})^{\Sigma^1}$ asserts the
tower clause \eqref{4.8:eq-tower-clause-a} at every label of $\mathfrak{T}_{k+1}(Y)$ and for
every $\mathfrak{r}\in\mathfrak{R}$. At the label $(p,Z)=(k+1,\square_0)$ its constant is
\begin{equation*}
(1+\beta)\alpha_0\,L^{-((k+1)-(k+1))}=(1+\beta)\alpha_0 .
\end{equation*}
This is \eqref{4.8:eq-inclusion-sentences-b} for both rules, read on $\square_0^{\wr(k+1)}$ as
in \eqref{4.8:eq-cube-lattices-b}. For the window $\square_0=\tilde{\square}^{5}$, removing two
layers of $\pi_k$-cubes leaves $\tilde{\square}^{3}$, so that
$\square_0^{\wr(k+1)}=\tilde{\square}^{3}$. Since $(1+\beta)\alpha_0\le(1+2\beta)\alpha_0$,
Lemma~\ref{4.8:lem-class-properties}(d$_1$) shows that $(\mathbf{U},\mathbf{J})$ is also a
$\Sigma^1$-representative at $(k+1,Y;(1+2\beta)\alpha_0,\cdot,\cdot)$, the other two constants
unchanged; the same reading at the label $(k+1,\square_0)$ then gives the tower clause at the
constant $(1+2\beta)\alpha_0$ on $\tilde{\square}^{3}$, for both rules, which is the window
clause assumed in \cite[(3.40)]{B-CMP109}.

(c). We apply Theorem~\ref{4.8:thm-restriction-units} with $i=i'=k+1$, so that
$\mu=1$. Its hypotheses hold with the r\^oles of $X$ and $Y$ there taken by $Y$ and $X$ here:
$X\in\mathbf{D}_{i-1}$, since $X\in\mathbf{D}_{k+1}$ and $\mathbf{D}_{k+1}\subset\mathbf{D}_k$;
$Y\in\mathbf{D}_{i'-1}$, since $i'-1=k$;
and $Y\subset X$. The theorem gives the first inclusion
of \eqref{4.8:eq-inclusion-sentences-c}, which is also
Lemma~\ref{4.8:lem-class-properties}(d$_2$). The second inclusion is
Lemma~\ref{4.8:lem-class-properties}(d$_1$). There the first two arguments increase from
$\alpha_0,\alpha_1$ to $(1+\beta)\alpha_0,(1+\beta)\alpha_1$, and the third argument is
$\alpha_0$ on both sides. No condition on $L$ enters.

(d). Let $(\mathbf{U},\mathbf{J})$ belong to the class on the left of
\eqref{4.8:eq-inclusion-sentences-d}. We apply Theorem~\ref{4.8:thm-restriction-units} with
levels $i$ and $i'$. Its hypotheses hold: $1\le i'\le i$, $Y\in\mathbf{D}_{i-1}$,
$X\in\mathbf{D}_{i'-1}$ and $X\subset Y$. The ratio is
\begin{equation*}
\mu=L^{-(i-i')}=L^{-(\mathfrak{n}+1)},\qquad \mu^3=\mathfrak{s}.
\end{equation*}
By Definition~\ref{4.8:def-one-power-class}, the class on the left of
\eqref{4.8:eq-inclusion-sentences-d} is the union of the $G^c$-orbits of its
$\Sigma^1$-representatives; so $(\mathbf{U},\mathbf{J})=(\mathbf{U}_0,\mathbf{J}_0)^u$ for a
$G^c$-valued gauge transformation $u$ on $Y$ and a $\Sigma^1$-representative
$(\mathbf{U}_0,\mathbf{J}_0)$ at
$(i,Y;(1+\beta_s)\alpha_{0,i},(1+\beta_s)\alpha_{1,i},\alpha_{0,i})$. The form of the theorem
with the current renormalised by $\mu^3$ gives that
$(\mathbf{U}_0,\mathfrak{s}\mathbf{J}_0)\restriction X$ is a $\Sigma^1$-representative at
\begin{equation*}
\bigl(i',X;\mu(1+\beta_s)\alpha_{0,i},\mu(1+\beta_s)\alpha_{1,i},
\mathfrak{s}\,\alpha_{0,i}\bigr).
\end{equation*}
By the first inequality of \eqref{4.8:eq-inclusion-sentences-1},
$\mu(1+\beta_s)\alpha_{\cdot,i}\le\frac12\alpha_{\cdot,i'}$ for $\alpha_{\cdot,i}$ standing for
either of $\alpha_{0,i}$, $\alpha_{1,i}$. By the
second, $\mathfrak{s}\,\alpha_{0,i}\le\alpha_{0,i'}$. Lemma~\ref{4.8:lem-class-properties}(d$_1$)
shows that $(\mathbf{U}_0,\mathfrak{s}\mathbf{J}_0)\restriction X$ is a $\Sigma^1$-representative
at $(i',X;\frac12\alpha_{0,i'},\frac12\alpha_{1,i'},\alpha_{0,i'})$. The gauge action on the
current is linear, so $(\mathbf{U},\mathfrak{s}\mathbf{J})=(\mathbf{U}_0,\mathfrak{s}\mathbf{J}_0)^u$;
and restriction to $X$ commutes with the gauge action, by the identity for orbits used in the
proof of Theorem~\ref{4.8:thm-restriction-units}:
\begin{equation*}
(\mathbf{U},\mathfrak{s}\mathbf{J})\restriction X
=\bigl((\mathbf{U}_0,\mathfrak{s}\mathbf{J}_0)\restriction X\bigr)^{u\restriction X}.
\end{equation*}
Hence $(\mathbf{U},\mathfrak{s}\mathbf{J})\restriction X$ lies in the orbit of a
$\Sigma^1$-representative at $(i',X;\frac12\alpha_{0,i'},\frac12\alpha_{1,i'},\alpha_{0,i'})$,
which is \eqref{4.8:eq-inclusion-sentences-d}.
\end{proof}

\begin{remark}\label{4.8:prop-near-case}
The statement of this proposition is not written out here; parts of its proof are printed below.
\end{remark}

\begin{lemma}[Near case: the tower clause at every label]\label{4.8:prf-near-case-labels}
Let $j$, $X$, the datum $D=\tau Q(L^{-1}\eta\mathbf{H}_{k+1})+B'$, Ba{\l}aban's minimiser
$W=U_j(\square_0,e^{iD})$ of \cite[Lemma~4]{B-CMP109} with its window tower $(\square_q)_{q\le j}$
and its gauge $u_j=u_j(D)$, the $u_j$-copy $(\mathcal{U},\mathcal{J})=(W,J_j)^{u_j^{-1}}$ and the
decomposition $U=1$, $U'=\mathcal{U}$ be as in Proposition~\ref{4.8:prop-near-case}. In particular
$X\in\mathbf{D}_j$, $X\subset\tilde{\square}^{2}$ and $\mathcal{U}=\exp i\xi\mathbf{H}_j(\square_0,D)$.
Assume the following.
\begin{enumerate}
\item The sub-tower conjugation lemma for Ba{\l}aban's minimisers, which is also the first antecedent
of Proposition~\ref{4.8:prop-near-case}, in the form applied below. Its
hypotheses are: a domain with a margin inside the top member of an outer tower; an admissible inner
tower; a depth $n$ with $n\le p\le j$; and a regular outer minimiser for which the sub-problems have
unique minimal orbits. Its conclusion is that the depth-$n$ minimiser on the inner tower, with the
layered averages of the $u_j$-copy $\mathcal{U}$ of the outer minimiser as data, is the restriction
of $\mathcal{U}$ to the domain, transformed by a gauge $g=u_j\bar u^{-1}$. This yields identities
for its curvature and its current, stated below. Its application to the window tower of $W$ includes
the check of the margin condition for every $X\in\mathbf{D}_j$ with $X\subset\tilde{\square}^{2}$,
and it is used together with the representation lemma for the gauge $u_j(D)$, built on
\cite[Proposition~9]{B-CMP102b}, which gives $u_j(D)$ through that proposition together with its
Lipschitz property in the supremum norm.
\item The given bound, under its ceiling.
\item The corresponding floor, where $c_{\mathrm{Ad}}$ is the absolute constant of
\eqref{4.8:eq-near-case-labels-a} below; on $\beta\le1$ one may take $c_{\mathrm{Ad}}=\tfrac13$.
\end{enumerate}
Let $B'$ satisfy $|B'|<2\alpha_3$, let $(p,Z)\in\mathfrak{T}_j(X)$ and $\mathfrak{r}\in\mathfrak{R}$,
and let $1\le n\le p$. Then, on all of $Z$,
\begin{equation}\label{4.8:eq-near-case-labels-1}
\begin{aligned}
\bigl|\partial U_n^{\mathfrak{r}}(Z;M^{\boldsymbol{\cdot}}(\mathcal{U}))-1\bigr|
&<(1+c_{\mathrm{Ad}}\beta)(1+8\beta)L^{-2}\alpha_0\,\xi_j^2,\\
\bigl|J_n^{\mathfrak{r}}(Z;M^{\boldsymbol{\cdot}}(\mathcal{U}))\bigr|
&<(1+c_{\mathrm{Ad}}\beta)(1+11\beta)L^{-2}\alpha_0\,(L^n\xi_j)^3,
\end{aligned}
\end{equation}
and
\begin{equation*}
U_n^{\mathfrak{r}}(Z;M^{\boldsymbol{\cdot}}(1))=1,\qquad J_n^{\mathfrak{r}}(Z;M^{\boldsymbol{\cdot}}(1))=0.
\end{equation*}
Consequently, the tower clause
$(\mathrm{iv})^{\mathfrak{r}}_{p,Z}[\tfrac12\alpha_0L^{-(j-p)}]$ of \eqref{4.8:eq-tower-clause-a}
holds for every $(p,Z)\in\mathfrak{T}_j(X)$ and every $\mathfrak{r}\in\mathfrak{R}$. That is, the
clause \eqref{4.8:eq-one-power-class-a} holds with $\tfrac12\alpha_0$ in place of $a_0$, and a
fortiori with $\alpha_0$. Moreover $(\mathcal{U},\mathcal{J})\restriction X$ is a
$\Sigma^1$-representative at $(j,X;\tfrac12\alpha_0,\tfrac12\alpha_1,\alpha_0)$ for every
$|B'|<2\alpha_3$, which gives in particular the half clause.
\end{lemma}

\begin{proof}
Fix $B'$ with $|B'|<2\alpha_3$. Fix $(p,Z)\in\mathfrak{T}_j(X)$, $\mathfrak{r}\in\mathfrak{R}$ and
$n\le p$. We treat the two halves of the tower clause in turn. The half with
$\mathbf{U}=U'U=\mathcal{U}$ comes first; the half with $U=1$ follows.

\emph{The conjugation identity.} We apply the sub-tower conjugation lemma of the first hypothesis
above with the following data:
\begin{itemize}
\item its domain is $Z$;
\item its inner tower is the tower $\boldsymbol{\Omega}^{\mathfrak{r}}(Z)$ truncated at depth $n$;
\item its outer tower is the window tower $(\square_q)_{q\le j}$ of $W$;
\item its top level is $j$.
\end{itemize}
We check its hypotheses in order.

The margin condition requires that the domain lie in the top member $\square_j$ of the outer tower,
with every endpoint of a bond of the domain at distance greater than $2\sqrt{d}\,L^j\delta$ from the
complement of $\square_j$, where $\delta$ is the margin parameter of the sub-tower conjugation lemma.
This condition holds for $X$ by the check included in the first hypothesis above. That check uses
$X\subset\tilde{\square}^{2}$ and
$\square_j\supset\tilde{\square}^{4}$, the latter from the inclusion
$\tilde{\square}^{4}\subset\square_{k+1}$ of \cite[\S3]{B-CMP109}. Since $Z\subset X$, every endpoint
of a bond of $Z$ is an endpoint of a bond of $X$. Hence the margin condition is inherited by $Z$.

The inner tower must be admissible, that is, made of unions of level-$p$ blocks. Under either rule of
Definition~\ref{4.8:def-tower-rules}, the members $\Omega_q$ of $\boldsymbol{\Omega}^{\mathfrak{r}}(Z)$
are built from cubes of side $\mathfrak{M}L^q$, respectively $M_1L^q$. The ordering $n\le p\le j$ holds
by the choice of $n$ and since $(p,Z)\in\mathfrak{T}_j(X)$.

Regularity of $W$, and uniqueness for the sub-problems in the sense of the uniqueness clause of
\cite[Theorem~1]{B-CMP102b}, are part of the restrictions of \cite[Lemma~4]{B-CMP109}. At a complex
datum $D$ these properties are carried over by the identity-theorem sentence of the sub-tower
conjugation lemma. That sentence applies because $M^{\boldsymbol{\cdot}}(\mathcal{U})(D)$ lies in the
polydisc of \cite[Proposition~9]{B-CMP102b} under the same ceilings.

The conclusion is Ba{\l}aban's identity \cite[(3.38)]{B-CMP109}, with $(Z,\mathfrak{r})$ in place of
his domain:
\begin{equation*}
U_n^{\mathfrak{r}}(Z;M^{\boldsymbol{\cdot}}(\mathcal{U}))=(\mathcal{U}\restriction Z)^{g},
\qquad g=u_j\bar u^{-1}.
\end{equation*}
Here $\bar u$ is constant on the blocks naturally associated with the minimiser $U_n(Z,\cdot)$ and
equal to $u_j$ at the centres of these blocks. In the form in which we use it, the identity reads
\begin{equation}\label{4.8:eq-near-case-labels-2}
\partial U_n^{\mathfrak{r}}(Z;M^{\boldsymbol{\cdot}}(\mathcal{U}))=R(g)(\partial\mathcal{U}),\qquad
J_n^{\mathfrak{r}}(Z;M^{\boldsymbol{\cdot}}(\mathcal{U}))=(L^n\xi_j)^3R(g)\mathcal{J}.
\end{equation}
The current on the left is written in level-$n$ units, so that it is $(L^n\xi_j)^3$ times the level-$j$
current. The operator $R(g)$ acts at each point $x$ through $\operatorname{Ad}_{g(x)}$.

\emph{The gauge.} Whatever the blocks are, for every point $x$ there is a point $y$, the centre of the
block containing $x$, such that
\begin{equation*}
g(x)=u_j(x)\,u_j(y)^{-1}.
\end{equation*}
The bound $|u_j-1|<B_3^2O_{(3.37)}M\alpha_0$ of \cite[(3.37)]{B-CMP109} is stated there at
$\mathbf{A}=0$ and $B'=0$. Here $O_{(3.37)}$ denotes the constant written $O(1)$ in
\cite[(3.37)]{B-CMP109}. At the datum $D$ the gauge is $u_j(D)$. By the representation of the gauge
$u_j(D)$ and its Lipschitz property in the supremum norm (the representation lemma for $u_j(D)$
of the first hypothesis above), the passage to $D$ adds at most $B_3|B'|<2B_3\alpha_3$. Put
\begin{equation*}
\delta_u:=B_3^2O_{(3.37)}M\alpha_0+2B_3\alpha_3 .
\end{equation*}
Then every value of $u_j(D)$ lies within $\delta_u$ of $1$.

Next we bound $\delta_u$. Ba{\l}aban's ceiling $(4B_3^2O_{(3.37)}M)^2\alpha_0\le\beta$ of
\cite[\S3]{B-CMP109} gives $B_3^2O_{(3.37)}M\alpha_0\le\beta/16$ and $\alpha_0\le\beta/16$, because
$B_3^2O_{(3.37)}M\ge1$. His ceiling $2B_3\alpha_3\le\beta L^{-2}\alpha_0$ of \cite[\S3]{B-CMP109},
read at the proof radius $2\alpha_3$, is $4B_3\alpha_3\le\beta L^{-2}\alpha_0$. Hence the second
summand of $\delta_u$ satisfies $2B_3\alpha_3\le\tfrac12\beta L^{-2}\alpha_0$. Using
$\alpha_0\le\beta/16$, $\beta\le1$ and $L\ge3$ ($L$ being an odd integer greater than one), which
gives $1+1/(2L^2)\le16/15$, we obtain
\begin{equation*}
\delta_u\le\frac{\beta}{16}+\frac{\beta\alpha_0}{2L^2}
\le\frac{\beta}{16}\Bigl(1+\frac{1}{2L^2}\Bigr)\le\frac{\beta}{15}.
\end{equation*}
In particular $\delta_u\le\tfrac1{15}<1$. For every $x$ we have $|u_j(x)|\le1+\delta_u$. The Neumann
series gives $|u_j(y)^{-1}|\le(1-\delta_u)^{-1}$. Hence $|g(x)^{\pm1}|\le(1+\delta_u)/(1-\delta_u)$ in
the operator norm. Since $\|\operatorname{Ad}_{h}\|\le|h|\,|h^{-1}|$ for every invertible $h$, it
follows that
\begin{equation}\label{4.8:eq-near-case-labels-a}
\bigl\|\operatorname{Ad}_{g}^{\pm1}\bigr\|
\le\Bigl(\frac{1+\delta_u}{1-\delta_u}\Bigr)^2
=1+\frac{4\delta_u}{(1-\delta_u)^2}
\le1+c_{\mathrm{Ad}}\beta ,
\end{equation}
with an absolute constant $c_{\mathrm{Ad}}$. The equality holds because
$(1+\delta_u)^2-(1-\delta_u)^2=4\delta_u$.

To see that one may take $c_{\mathrm{Ad}}=\tfrac13$, note that $\delta_u\le\tfrac1{15}$ gives
$(1-\delta_u)^2\ge(14/15)^2$. Therefore
\begin{equation*}
\frac{4\delta_u}{(1-\delta_u)^2}\le\frac{4\beta}{15}\Bigl(\frac{15}{14}\Bigr)^2
=\frac{4}{15}\cdot\frac{225}{196}\,\beta=\frac{15}{49}\,\beta<\frac{\beta}{3},
\end{equation*}
where the last inequality is $45<49$. With this choice $c_{\mathrm{Ad}}\beta\le\tfrac13<\tfrac12$. The
same bound holds for $\|R(g)^{\pm1}\|$, since $R(g)$ acts pointwise through the adjoint action of the
values of $g$.

\emph{The curvature member.} The corresponding floor and $1+8\beta\le1+11\beta$ give
$(1+c_{\mathrm{Ad}}\beta)(1+8\beta)L^{-2}\le\tfrac12$. In the proof of
Proposition~\ref{4.8:prop-near-case}, the first member of condition (iii) was obtained from
\cite[(3.52)]{B-CMP109}: it states $|\partial\mathcal{U}-1|<(1+8\beta)L^{-2}\alpha_0\xi_j^2$ on
$\tilde{\square}^{3}\supset Z$, at the datum $D$, for every $|B'|<2\alpha_3$. The first identity in
\eqref{4.8:eq-near-case-labels-2} gives $\partial U_n^{\mathfrak{r}}-1=R(g)(\partial\mathcal{U}-1)$.
Together with \eqref{4.8:eq-near-case-labels-a} and $\xi_j=L^{-(j-p)}\xi_p$, this yields, on all
of $Z$,
\begin{equation*}
\begin{aligned}
\bigl|\partial U_n^{\mathfrak{r}}(Z;M^{\boldsymbol{\cdot}}(\mathcal{U}))-1\bigr|
&<(1+c_{\mathrm{Ad}}\beta)(1+8\beta)L^{-2}\alpha_0\,\xi_j^2
\le\tfrac12\alpha_0\,\xi_j^2\\
&=\tfrac12\alpha_0L^{-2(j-p)}\xi_p^2\le\tfrac12\alpha_0L^{-(j-p)}\xi_p^2 .
\end{aligned}
\end{equation*}
The last step uses $j\ge p$. Ba{\l}aban closes his own condition (iv) at $(j,X)$ in the same way: from
the identity \cite[(3.38)]{B-CMP109}, the bound on $u_j$ in \cite[(3.37)]{B-CMP109} and condition
(iii), with a slightly better constant.

\emph{The current member.} Since $Z\subset X\subset\tilde{\square}^{2}\subset\tilde{\square}^{3}$, we
have $\sup_Z|\mathcal{J}|\le\sup_{\tilde{\square}^{3}}|\mathcal{J}|$. The given bound states $\sup_{\tilde{\square}^{3}}|\mathcal{J}|<\theta_J\alpha_0$ with
$\theta_J=(1+11\beta)L^{-2}$. The second identity in \eqref{4.8:eq-near-case-labels-2}, the bound
\eqref{4.8:eq-near-case-labels-a}, the corresponding floor, and $L^n\xi_p\le1$ (valid
since $n\le p$) then give, on all of $Z$,
\begin{equation*}
\begin{aligned}
\bigl|J_n^{\mathfrak{r}}(Z;M^{\boldsymbol{\cdot}}(\mathcal{U}))\bigr|
&\le(1+c_{\mathrm{Ad}}\beta)(L^n\xi_j)^3\sup_Z|\mathcal{J}|
<(1+c_{\mathrm{Ad}}\beta)\,\theta_J\alpha_0\,(L^n\xi_j)^3\\
&=(1+c_{\mathrm{Ad}}\beta)(1+11\beta)L^{-2}\alpha_0\,(L^n\xi_j)^3
\le\tfrac12\alpha_0\,(L^n\xi_j)^3\\
&=\tfrac12\alpha_0L^{-3(j-p)}(L^n\xi_p)^3
\le\tfrac12\alpha_0L^{-(j-p)}(L^n\xi_p)^2 .
\end{aligned}
\end{equation*}
This proves \eqref{4.8:eq-near-case-labels-1}. It also proves the $\mathbf{U}$-half of
$(\mathrm{iv})^{\mathfrak{r}}_{p,Z}[\tfrac12\alpha_0L^{-(j-p)}]$, on all of $Z$ and hence on
$Z^{\wr p}$.

\emph{The half with $U=1$.} Here the data are $M^{\boldsymbol{\cdot}}(1)=1$, since the layered averages
of the identity configuration are the identity. Consider the depth-$n$ problem on
$\boldsymbol{\Omega}^{\mathfrak{r}}(Z)$ truncated at $n$ with data $1$. Its energy vanishes at the
configuration $Y=1$, which is therefore an absolute minimum. By the uniqueness clause of
\cite[Theorem~1]{B-CMP102b}, used as in the uniqueness step of the proof of the sub-tower conjugation
lemma of the first hypothesis above,
the minimal orbit is unique, and its axial representative is $1$. Hence
$U_n^{\mathfrak{r}}(Z;M^{\boldsymbol{\cdot}}(1))=1$ and $|\partial U_n^{\mathfrak{r}}-1|=0$.

The current \cite[(1.8)]{B-CMP109} is built on $\operatorname{im}\partial U_n$, where
$\operatorname{im}U:=\tfrac1{2i}(U-U^{-1})$, and this vanishes at $\partial U_n=1$. Hence
$J_n^{\mathfrak{r}}(Z;M^{\boldsymbol{\cdot}}(1))=0$. Both inequalities of the tower clause therefore hold
trivially for $U$, at every constant.

\emph{The half clause.} We now collect, for every $|B'|<2\alpha_3$, the conditions defining a
$\Sigma^1$-representative at $(j,X;\tfrac12\alpha_0,\tfrac12\alpha_1,\alpha_0)$ for
$(\mathcal{U},\mathcal{J})\restriction X$ with $U=1$, $U'=\mathcal{U}$.
\begin{itemize}
\item Condition (i) is void, because $U=1$.
\item Condition (ii) holds at $\tfrac12\alpha_1$ by the corresponding chain in the proof
of Proposition~\ref{4.8:prop-near-case}.
\item For condition (iii), the first inequality of \cite[(1.14)]{B-CMP109} holds at
$(1+8\beta)L^{-2}\alpha_0\le\tfrac12\alpha_0$, by \cite[(3.52)]{B-CMP109} and the corresponding floor. The second member holds at $\gamma_0=\alpha_0$. Indeed
$|\mathcal{J}|<\theta_J\alpha_0<\alpha_0$, because the floor gives
$L^2\ge2(1+c_{\mathrm{Ad}}\beta)(1+11\beta)>1+11\beta$, so $\theta_J<1$. This is also the bound that
\cite[Lemma~4]{B-CMP109} states.
\item The clause $(\mathrm{iv})^{\Sigma^1}$ holds with constant $\tfrac12\alpha_0$, that is,
$(\mathrm{iv})^{\mathfrak{r}}_{p,Z}[\tfrac12\alpha_0L^{-(j-p)}]$ for every
$(p,Z)\in\mathfrak{T}_j(X)$ and every $\mathfrak{r}\in\mathfrak{R}$. For the $\mathbf{U}$-half this is
the content of the two preceding chains. There the curvature member is reached a fortiori, since
$1+8\beta\le1+11\beta$, and the current member uses $\theta_J$ from the given bound. The $U$-half is void.
\end{itemize}
Hence $(\mathcal{U},\mathcal{J})\restriction X$ is a $\Sigma^1$-representative at
$(j,X;\tfrac12\alpha_0,\tfrac12\alpha_1,\alpha_0)$ for every $|B'|<2\alpha_3$, which gives in
particular the half clause.
\end{proof}

\begin{remark}
The argument above for the tower clause uses nothing about the domain $Z$ beyond $Z\subset X$, and
nothing about the rule $\mathfrak{r}$ beyond the admissibility of the truncated tower. With the
decomposition $U=1$, $U'=\mathcal{U}$ of the near case, it uses no lower bound on the constant
$B^{(109)}$ of \cite[(1.12)]{B-CMP109}. It uses no condition on the ratio $\alpha_1/\alpha_0$ other
than Ba{\l}aban's condition $B_3^2O_{(3.37)}M\alpha_0<\tfrac12\alpha_1$ of \cite[\S3]{B-CMP109}.

At the running radii, the ordering of $\alpha_0$ below $\alpha_1$ expressed by that condition is a
strict lower bound on the constant $C_1$ in terms of $B_3$, $M$ and $C_0$, together with
$C_0\le C_1$. This floor is imposed on $C_1$, which is free upward and chosen after $M$ and $C_0$,
among the conditions of the statement of this chapter on the large-field part of the coupling
increment. It is imposed there and not in the present argument.
\end{remark}

\begin{proposition}[Far case: small perturbations stay in the class]\label{4.8:prop-far-case}
Assume the increment lemma, the statement that bounds the change of the tower quantities
$\partial U_n^{\mathfrak{r}}(Z;\cdot)$ and $J_n^{\mathfrak{r}}(Z;\cdot)$ of
\eqref{4.8:eq-tower-clause-a} when a background in the regularity class is multiplied by
$e^{i\xi_jA}$ with $A$ and its covariant differences small, together with its clause for towers:
the conclusions of the increment lemma hold for every tower that is admissible in the sense of
\cite[(1)]{B-CMP102b} and has collar at most $\frac12 M$, uniformly in $M=L^{m'}$.

Let $j\ge 1$ and $X\in\mathbf{D}_j\subset\mathbf{D}_{j-1}$.
Let $\alpha_0,\alpha_1,\gamma_0>0$, and let
$(\mathbf{U},\mathbf{J})$ be a $\Sigma^1$-representative at
$(j,X;\frac12\alpha_0,\frac12\alpha_1,\gamma_0)$ in the sense of
Definition~\ref{4.8:def-one-power-class}, with its decomposition $\mathbf{U}=U'U$. Let
$\alpha_2>0$ and let $A$ be a $\mathfrak{g}^c$-valued function on the bonds of $X$ with
$|A|<\alpha_2$ and $|\nabla^{\xi_j}_{\mathbf{U}}A|<\alpha_2$, where
$\nabla^{\xi_j}_{\mathbf{U}}$ is the covariant difference at spacing $\xi_j$ in the
background $\mathbf{U}$. Assume the ceilings of the increment lemma,
\begin{align*}
&C_7\alpha_0\le\bar\varepsilon_1,\qquad
c_Q(\alpha_1+\alpha_2)\le\tfrac12 C_1\bar\varepsilon_1,\qquad
c_{12}LMB^{(109)}\cdot\tfrac12\alpha_0+\tfrac12\alpha_1+\alpha_2\le c_\sharp,\\
&\alpha_1\le 1,\qquad
\alpha_2\le\min\{\alpha_1/8,\ \alpha_0/C_{IV}\},
\end{align*}
where $C_7$, $\bar\varepsilon_1$, $c_Q$, $C_1$, $c_\sharp$, $C_{IV}$ are constants of the
increment lemma that depend on $(d,L)$ only, $B^{(109)}$ is the constant $B$ of
\cite[(1.12)]{B-CMP109} and $c_{12}$ the constant
written $O(1)$ there, and assume the floors $c_{12}L\ge 2$, $L\ge 4d$ and
\begin{equation}\label{4.8:eq-far-case-a}
\delta_0 M\ge 32L,\qquad\qquad M\ge 2LR_1M_1,
\end{equation}
where $\delta_0$, $R_1$ and the block size $M_1$ are those of the increment lemma. Then there
is a $\mathfrak{g}^c$-valued function $A''$ on the bonds of $X$ with
$e^{i\xi_jA}\mathbf{U}=e^{i\xi_jA''}U$, and:
\begin{enumerate}
\item[(a)] the configuration $e^{i\xi_jA}\mathbf{U}=e^{i\xi_jA''}U$ satisfies on $X$
condition~(i) of Definition~\ref{4.8:def-one-power-class} at $\alpha_0$ (indeed at
$\frac12\alpha_0$, its $G$-valued factor being the factor $U$ of $\mathbf{U}$),
condition~(ii) at $\alpha_1$, the first inequality of \cite[(1.14)]{B-CMP109} at $\alpha_0$,
and the condition \eqref{4.8:eq-one-power-class-a} at $\alpha_0$;
\item[(b)] hence the pair $(e^{i\xi A(t_\square)}\mathbf{U},\,J+\mathfrak{J})$ produced in
the far case of the operator bounds at free current, with $\xi=\xi_j$, where $A(t_\square)$
denotes the perturbation and $J$, $\mathfrak{J}$ the current and the produced current
increment in the notation of that statement, lies in
$U^{\Sigma^1}_j(X;\alpha_0,\alpha_1)$, its current being bounded by the estimate
$|J+\mathfrak{J}|\le L^{-3}\alpha_0+\alpha_2<\alpha_0$ of that statement; and so does
$(e^{i\xi_jA}\mathbf{U},\mathbf{J}')$ for every current $\mathbf{J}'$ with
$|\mathbf{J}'|<\alpha_0$.
\end{enumerate}
\end{proposition}

\begin{proof}
\emph{Conditions} (i), (ii) \emph{and the first inequality of}
\cite[(1.14)]{B-CMP109}. These three are supplied directly by the corresponding conclusion of
the increment lemma, which is pointwise and does not involve $M$.

\emph{The condition \eqref{4.8:eq-one-power-class-a}.} Fix a label
$(p,Z)\in\mathfrak{T}_j(X)$ and a rule $\mathfrak{r}\in\mathfrak{R}$. Put
\begin{equation*}
\mu:=L^{-(j-p)},\qquad M':=M/L=L^{m'-1};
\end{equation*}
$\mu$ is the scale ratio of Theorem~\ref{4.8:thm-restriction-units} at $(i,i')=(j,p)$, and
$\mu\le 1$ because $p\le j$. By Notation~\ref{4.8:not-cube-lattices}, the cubes of
$\pi_{p-1}$ formed with the cube size $M$ are the cubes of $\pi_p$ formed with the cube
size $M'$: both have sides of $ML^{p-1}=M'L^p$ lattice steps and the same centres. Hence
$Z$, an element of $\mathbf{D}_{p-1}$ formed with $M$, is an element of $\mathbf{D}_p$ formed
with $M'$, and the set $Z^{\wr p}$ of \eqref{4.8:eq-cube-lattices-b} is the two-layer interior
$\tilde{Z}^{-2}$ of $Z$ at level $p$ formed with $M'$. We apply the increment lemma with its
letters replaced according to
\begin{equation}\label{4.8:eq-far-case-1}
\begin{split}
&(j,X,M,B^{(109)};\alpha_0,\alpha_1,\alpha_2;A;\mathfrak{r}_1)\\
&\qquad\longmapsto
(p,Z,M',LB^{(109)};\mu\alpha_0,\mu\alpha_1,\mu\alpha_2;\mu A|_Z;\mathfrak{r}).
\end{split}
\end{equation}

\emph{Licence for the substitution \eqref{4.8:eq-far-case-1}.} It is given by the clause for
towers assumed in the statement, which we read as follows: $M=L^{m'}$ may be any power of $L$
obeying the floors of the increment lemma, with $m'$ the exponent of
Notation~\ref{4.8:not-cube-lattices}, and $B^{(109)}$ may be any value of the constant of
\cite[(1.12)]{B-CMP109}, the constants of the lemma remaining unchanged. This reading rests on
the way the increment lemma depends on these letters. In its hypotheses and its proof, $M$
enters only through the side $2M$ of its recentring cube, with $2M\le c_{12}LM$ because
$c_{12}L\ge 2$; through its parameter $r=M/\lambda$, where $\lambda$ is a further parameter of
the lemma; through its depth comparison, in which
the relevant depth is at least
\begin{equation*}
2M/\lambda-M/(2\lambda)-3\ \ge\ M/\lambda\ \ge\ M\ \ge\ 2R_1M_1\ >\ R_1M_1;
\end{equation*}
through the bound $|\nu|\le dM\alpha_2$ on its parameter $\nu$, which need not be small; and
through the quantity $(1+M)e^{-\delta_0M/32}$, which satisfies
\begin{equation*}
(1+M)e^{-\delta_0M/32}\ \le\ 1+\frac{32}{e\delta_0},
\end{equation*}
since $e^{-\delta_0M/32}\le 1$ and $x e^{-\delta_0x/32}\le 32/(e\delta_0)$ for every
$x\ge 0$; the lemma uses this bound under $\delta_0M\ge 32$. Apart from the floors, in which
$M$ occurs directly, $M$ and $B^{(109)}$ enter the ceilings only through the product
$c_{12}LMB^{(109)}$, in the ceiling on $\alpha_0$. The tower
enters only through the lemma's bound on the collar and through admissibility in the sense of
\cite[(1)]{B-CMP102b}. Here the collar of a tower of $X$ is the depth of $X\setminus\Omega_n$,
where $\Omega_n$ is the member of index $n$ of the tower,
in the notation of Definition~\ref{4.8:def-tower-rules}; it is at most $\frac12 M$ whether it
is measured in the units $L^n$ of the top block, as in \eqref{4.8:eq-tower-rules-b}, or in
level-$j$ units, and the first of these two readings implies the second, because
$L^n\le L^j$. The constants $C_K$, $C_{IV}$, $\bar\varepsilon_1$, $c_Q$, $C_1$, $c_\sharp$ of
the lemma depend on $(d,L)$ only and are free of $M$, $n$, $j$, $k$ and $X$.

Both rules of $\mathfrak{R}$ produce towers admissible in the sense of
\cite[(1)]{B-CMP102b} with collar at most $R_1M_1$, by \eqref{4.8:eq-tower-rules-b}; and
$R_1M_1\le\frac12 M'$ because the second inequality of \eqref{4.8:eq-far-case-a}, divided by
$L$, reads $M'\ge 2R_1M_1$. So the tower of $Z$ under $\mathfrak{r}$ is covered by the clause
at the cube size $M'$.

\emph{Hypotheses at the substituted letters.} By the change of units
\eqref{4.8:eq-restriction-units-a}, with $\xi_j=\mu\xi_p$: writing $U'=e^{i\xi_jA'}$ as in
condition~(ii), the restriction $\mathbf{U}|_Z=e^{i\xi_p\mu A'}U$ satisfies (i)--(iii) at
level $p$ with the radii $(\frac12\mu\alpha_0,\frac12\mu\alpha_1)$. For
\cite[(1.11)]{B-CMP109} and the first inequality of \cite[(1.14)]{B-CMP109} this holds because
a bound $b\xi_j^2$ equals $\mu^2b\xi_p^2$, so that these conditions hold at
$\frac12\mu^2\alpha_0\le\frac12\mu\alpha_0$. For condition~(ii), $|\mu A'|=\mu|A'|$, and a
difference quotient at spacing $\xi_p$ is $\mu$ times the one at spacing $\xi_j$. For
\cite[(1.12)]{B-CMP109} under \eqref{4.8:eq-tower-clause-b}: a cube contained in $Z$ of side
at most
\begin{equation*}
c_{12}LM'\cdot L^p\ \le\ c_{12}LM\cdot L^j
\end{equation*}
lattice steps is a cube to which \cite[(1.12)]{B-CMP109} applies for the representative at
level $j$ on $X$, and after the change of units the bound there becomes
\begin{equation*}
\mu\cdot c_{12}LMB^{(109)}\cdot\tfrac12\alpha_0
\ =\ c_{12}LM'\bigl(LB^{(109)}\bigr)\cdot\tfrac12(\mu\alpha_0),
\end{equation*}
which is the bound demanded at the substituted letters. The hypothesis of the increment lemma
that the condition (iv) holds for the maximal tower of $X$ built with the block size
$\mathfrak{M}$ and the gap number $R_1$,
both for $\mathbf{U}$ and for $U$, becomes at the substituted letters the same condition at
$\frac12\mu\alpha_0$, extended by the clause for towers to the $\mathfrak{r}$-tower of $Z$:
that is,
$(\mathrm{iv})^{\mathfrak{r}}_{p,Z}[\frac12\alpha_0L^{-(j-p)}]$ in the sense of
\eqref{4.8:eq-tower-clause-a}. This is the condition \eqref{4.8:eq-one-power-class-a} at the
label $(p,Z)$ and the rule $\mathfrak{r}$ for the radius $\frac12\alpha_0$, which holds because
$(\mathbf{U},\mathbf{J})$ is a $\Sigma^1$-representative at radius $\frac12\alpha_0$. Finally,
again by \eqref{4.8:eq-restriction-units-a}, $|\mu A|=\mu|A|<\mu\alpha_2$ and
\begin{equation*}
|\nabla^{\xi_p}_{\mathbf{U}}(\mu A)|\ =\ \mu^2|\nabla^{\xi_j}_{\mathbf{U}}A|
\ \le\ \mu|\nabla^{\xi_j}_{\mathbf{U}}A|\ <\ \mu\alpha_2 .
\end{equation*}

\emph{Ceilings at the substituted letters.} They read
\begin{align*}
&C_7\mu\alpha_0\le\bar\varepsilon_1,\qquad
c_Q\mu(\alpha_1+\alpha_2)\le\tfrac12C_1\bar\varepsilon_1,\qquad
\mu\alpha_1\le 1,\\
&c_{12}LM'\bigl(LB^{(109)}\bigr)\cdot\tfrac12\mu\alpha_0+\tfrac12\mu\alpha_1+\mu\alpha_2
=\mu\bigl(c_{12}LMB^{(109)}\cdot\tfrac12\alpha_0+\tfrac12\alpha_1+\alpha_2\bigr)
\le c_\sharp,\\
&\mu\alpha_2\le\min\{\mu\alpha_1/8,\ \mu\alpha_0/C_{IV}\}.
\end{align*}
Each is $\mu\le 1$ times the corresponding assumed ceiling; here $C_{IV}$, which depends on
$(d,L)$ only, is the same number at $M'$ as at $M$.

\emph{Floors at $M'$.} The floor $\delta_0M'\ge 32$ is equivalent to the first inequality of
\eqref{4.8:eq-far-case-a}, and the floor $2R_1M_1\le M'$ to the second. The floors
$c_{12}L\ge 2$ and $L\ge 4d$ do not involve $M$ and are unchanged. The floors of the regularity
class of the increment lemma,
\begin{equation*}
17M_1\le LM',\qquad 17M_1\le c_{12}LM',
\end{equation*}
follow from
\begin{align*}
LM'&=M\ge 2LR_1M_1\ge 2R_1M_1\ge 17M_1,\\
c_{12}L\cdot M'&\ge 2M'\ge 4R_1M_1\ge 17M_1,
\end{align*}
where both lines use $2R_1\ge 2L$, the gap number $R_1$ of the increment lemma being at least
$L$, and $2L\ge 17$, the latter by $L\ge 4d$ with $d=4$, and the
second line uses $c_{12}L\ge 2$ and $M'\ge 2R_1M_1$.

\emph{Conclusion.} All hypotheses of the increment lemma hold at the substituted letters
\eqref{4.8:eq-far-case-1}. Its conclusion that the condition (iv) holds with $\alpha_0$, both
for $e^{i\xi_jA}\mathbf{U}$ and for $U$, reads at these letters
\begin{equation*}
(\mathrm{iv})^{\mathfrak{r}}_{p,Z}[\mu\alpha_0]
=(\mathrm{iv})^{\mathfrak{r}}_{p,Z}[\alpha_0L^{-(j-p)}],
\end{equation*}
which is what \eqref{4.8:eq-one-power-class-a} demands at $(p,Z,\mathfrak{r})$ for the radius
$\alpha_0$. Since $(p,Z)\in\mathfrak{T}_j(X)$ and $\mathfrak{r}\in\mathfrak{R}$ were arbitrary,
\eqref{4.8:eq-one-power-class-a} holds at $\alpha_0$, and (a) is proved.

\emph{Part} (b). By the far case of the operator bounds at free current, the perturbation
$A(t_\square)$ satisfies the hypotheses imposed on $A$ in the present proposition, so part (a)
applies with $A(t_\square)$ in place of $A$. It gives for the configuration
$e^{i\xi A(t_\square)}\mathbf{U}$ of the produced pair a decomposition with the same factor
$U$ and the corresponding combined exponent, and conditions (i), (ii), the first inequality of
\cite[(1.14)]{B-CMP109} and \eqref{4.8:eq-one-power-class-a}, all at $(\alpha_0,\alpha_1)$.
In the class $U^{\Sigma^1}_j(X;\alpha_0,\alpha_1)$ the current radius is $\gamma_0=\alpha_0$,
and the current of the pair obeys $|J+\mathfrak{J}|\le L^{-3}\alpha_0+\alpha_2<\alpha_0$ by the
far case of the operator bounds at free current. Hence the pair is a $\Sigma^1$-representative
at $(j,X;\alpha_0,\alpha_1)$ and lies in the union of orbits $U^{\Sigma^1}_j(X;\alpha_0,\alpha_1)$.
The same argument with (a) as it stands and a current $\mathbf{J}'$ with $|\mathbf{J}'|<\alpha_0$
gives $(e^{i\xi_jA}\mathbf{U},\mathbf{J}')\in U^{\Sigma^1}_j(X;\alpha_0,\alpha_1)$.
\end{proof}

The single power of the scale ratio $L^{-(j-p)}$ in \eqref{4.8:eq-one-power-class-a} is what
the proof of Proposition~\ref{4.8:prop-far-case} supplies. At the label $(p,Z)$ and the tower
member of index $n$, the increment
produced by the increment lemma at the substituted letters is
\begin{equation}\label{4.8:eq-far-case-b}
\tfrac12C_{IV}(\mu\alpha_2)\lambda'^2\ \le\ \tfrac12\alpha_0\mu\lambda'^2,
\qquad \lambda':=L^{n-p},
\end{equation}
the inequality by the ratio $\alpha_2\le\alpha_0/C_{IV}$. This is one factor $\mu$ short of
$\alpha_0\mu^2\lambda'^2$, and is of the size $\alpha_0\mu\lambda'^2$ demanded by
Definition~\ref{4.8:def-one-power-class}. The factor $\lambda'^2$ comes from recentring on the
cube of half-side $M'/\lambda'$; this cube lies in $Z$ because $Z^{\wr p}$ has depth $2M'$ in
$Z$; the role of the second inequality of \eqref{4.8:eq-far-case-a} is to place the near
region in the top layer. A second factor $\mu$ would require a recentring cube of
level-$j$ size inside $Z$, and no such cube exists for the sub-towers at the rim.

\begin{proposition}[Small pairs and true backgrounds lie in the class]\label{4.8:prop-small-pairs}
Let $j\ge 1$ and $X\in\mathbf{D}_{j-1}$. Let $\alpha_0,\alpha_1,\gamma_0$ and $\alpha_0',\alpha_1'$ be positive
numbers with
\begin{equation}\label{4.8:eq-small-pairs-1}
\alpha_1'\le\alpha_1,\qquad 2B_3\alpha_0'\le\alpha_0,
\end{equation}
where $B_3$ is the constant of \cite[Proposition~9]{B-CMP102b}. Let the cube size $M$ satisfy
$M\ge 2LR_1M_1$, the second condition in \eqref{4.8:eq-far-case-a}. Assume the following.

\begin{itemize}
\item For every $(p,Z)\in\mathfrak{T}_j(X)$ (see \eqref{4.8:eq-one-power-class-a}) and every
$\mathfrak{r}\in\mathfrak{R}$, the estimate \cite[(1.17)]{B-CMP109} holds at complex pairs, as
stated on \cite[p.~263]{B-CMP109}. It is applied at level $p$ on $Z$, for the tower of $Z$ built under
the rule $\mathfrak{r}$, with $M/L$ in place of the cube size $M$ and $LB^{(109)}$ in place of
$B^{(109)}$.
\end{itemize}

Let $(\mathbf{U},\mathbf{J})$ be a pair on the bonds of $X$, with $\mathbf{U}=U'U$. Suppose it satisfies
conditions (i), (ii), (iii) of Ba{\l}aban's regularity spaces
(Definition~\ref{4.8:not-regularity-spaces}) at level $j$ on $X$ with the constants
$(\alpha_0',\alpha_1',\gamma_0)$. Within (i), condition \cite[(1.12)]{B-CMP109} is read under
\eqref{4.8:eq-tower-clause-b}. Then
\begin{equation}\label{4.8:eq-small-pairs-2}
(\mathbf{U},\mathbf{J})\in U^{\Sigma^1}_j(X;\alpha_0,\alpha_1,\gamma_0).
\end{equation}
Consequently, clause (a) of the lemma of this chapter whose hypothesis is \cite[(1.17)]{B-CMP109}
and whose statement is formulated with Ba{\l}aban's space $U^c_j$ holds with $U^{\Sigma^1}_j$ in
place of $U^c_j$ and with the same constant, written $c_T$, as in that lemma.

Let $\varepsilon_0$ denote the constant so named in \cite[p.~263]{B-CMP109}. The pairs of
\cite[p.~263]{B-CMP109} built on the minimal configurations $U_j$ with
$|\partial U_j-1|<\varepsilon_0\xi_j^2$, $\varepsilon_0$ sufficiently small, lie in every
$U^{\Sigma^1}_j(X;\alpha_0,\alpha_1)$ (the third constant being omitted when $\gamma_0=\alpha_0$).
\end{proposition}

\begin{proof}
By the definition of the one-power hereditary regularity class
(Definition~\ref{4.8:def-one-power-class}), the pair $(\mathbf{U},\mathbf{J})$, with its
decomposition $\mathbf{U}=U'U$, is a candidate representative. Conditions (i)--(iii) at level $j$
on $X$ are assumed, with the constants $(\alpha_0',\alpha_1',\gamma_0)$ and with
$\alpha_1'\le\alpha_1$. What has to be verified is the clause $(\mathrm{iv})^{\Sigma^1}$ of
\eqref{4.8:eq-one-power-class-a}. That clause requires
$(\mathrm{iv})^{\mathfrak{r}}_{p,Z}[\alpha_0L^{-(j-p)}]$ of \eqref{4.8:eq-tower-clause-a} for every
$(p,Z)\in\mathfrak{T}_j(X)$ and every $\mathfrak{r}\in\mathfrak{R}$.

Fix $(p,Z)\in\mathfrak{T}_j(X)$ and $\mathfrak{r}\in\mathfrak{R}$. Put $\mu:=L^{-(j-p)}$, the scale ratio
of Theorem~\ref{4.8:thm-restriction-units} at the pair of levels $(j,p)$. Then $\xi_j=\mu\xi_p$ and
$\mu\le 1$.

\emph{Step 1: the pair at level $p$ on $Z$.} We apply the change of units
\eqref{4.8:eq-restriction-units-a} at $(j,p)$. It shows that the pair obeys (i)--(iii) at level $p$
on $Z$, with $M/L$ in place of $M$, with $LB^{(109)}$ in place of $B^{(109)}$, and with the constants
\begin{equation*}
(\alpha_0'',\alpha_1''):=(\mu\alpha_0',\mu\alpha_1').
\end{equation*}
In detail, the individual conditions transform as follows.

For \cite[(1.11)]{B-CMP109} and the first inequality of \cite[(1.14)]{B-CMP109}, the bound at level
$j$ reads at level $p$ as
\begin{equation*}
\alpha_0'\xi_j^2=\mu^2\alpha_0'\xi_p^2\le\mu\alpha_0'\xi_p^2 .
\end{equation*}

For \cite[(1.12)]{B-CMP109}, let $\square\subset Z$ be a cube of side at most
$c_{12}L(M/L)L^p=c_{12}ML^p$ steps. Since $p\le j$, we have $c_{12}ML^p\le c_{12}LML^j$, and
$Z\subset X$. Hence $\square$ is a cube contained in $X$ of side at most $c_{12}LML^j$ steps, so
\eqref{4.8:eq-tower-clause-b} at level $j$ applies to it. It supplies a gauge transformation $u$ on
$\square$ with $U^u=e^{i\xi_jA}$ and
\begin{equation*}
|A|,\ |\nabla^{\xi_j}A|<c_{12}LMB^{(109)}\alpha_0'.
\end{equation*}
With $A^\flat:=\mu A$ we have $e^{i\xi_jA}=e^{i\xi_pA^\flat}$. By \eqref{4.8:eq-restriction-units-a},
\begin{equation*}
|A^\flat|=\mu|A|,\qquad |\nabla^{\xi_p}A^\flat|=\mu^2|\nabla^{\xi_j}A|\le\mu|\nabla^{\xi_j}A|.
\end{equation*}
Therefore both $|A^\flat|$ and $|\nabla^{\xi_p}A^\flat|$ are smaller than
\begin{equation*}
c_{12}LMB^{(109)}\cdot\mu\alpha_0'
=c_{12}L\,(M/L)\,(LB^{(109)})\cdot(\mu\alpha_0').
\end{equation*}
Only one factor $\mu$ enters here. For this reason the constant in the place of $\alpha_0$ is
$\mu\alpha_0'$, and not $\mu^2\alpha_0'$.

Condition (ii) holds at $\mu\alpha_1'$, by the same change of units applied to $U'$.

Condition (iii) is treated as (i): its first inequality, that of \cite[(1.14)]{B-CMP109}, was
treated above together with \cite[(1.11)]{B-CMP109}, with the constant $\mu\alpha_0'$; its current
bound $|\mathbf{J}|<\gamma_0$ is unchanged, since no renormalisation of the current is made.

\emph{Step 2: the estimate at level $p$, for $\mathbf{U}$.} At level $p$ on $Z$ we now apply the
estimate \cite[(1.17)]{B-CMP109}, which \cite[p.~263]{B-CMP109} obtains from
\cite[Proposition~9]{B-CMP102b}; at complex pairs it is the hypothesis of the proposition and is
used as stated there. We use it for the tower of $Z$ that enters
\eqref{4.8:eq-tower-clause-a}. Theorem~1 and Proposition~9 of \cite{B-CMP102b} are stated for every
sequence admissible in the sense of \cite[(1)]{B-CMP102b}, and the towers of both rules in
$\mathfrak{R}$ are such sequences; so the estimate is available for either rule. The constant $B_3$
depends on $d$ and $L$ only \cite{B-CMP102b}. It is therefore the same when the cube size is $M/L$.
With the constants of Step~1 the estimate gives, for $n=1,\dots,p$ and on $Z^{\wr p}$
(see \eqref{4.8:eq-cube-lattices-b}),
\begin{equation}\label{4.8:eq-small-pairs-3}
\begin{gathered}
|\partial U_n^{\mathfrak{r}}(Z;M^{\boldsymbol{\cdot}}(\mathbf{U}))-1|<2B_3\mu\alpha_0'\,\xi_p^2,\\
|J_n^{\mathfrak{r}}(Z;M^{\boldsymbol{\cdot}}(\mathbf{U}))|<2B_3\mu\alpha_0'\,(L^n\xi_p)^2 .
\end{gathered}
\end{equation}
By \eqref{4.8:eq-small-pairs-1}, $2B_3\mu\alpha_0'\le\mu\alpha_0=\alpha_0L^{-(j-p)}$. Hence
\eqref{4.8:eq-small-pairs-3} gives the bounds for $\mathbf{U}$ in
$(\mathrm{iv})^{\mathfrak{r}}_{p,Z}[\alpha_0L^{-(j-p)}]$.

\emph{Step 3: the same bounds for $U$.} Clause \eqref{4.8:eq-tower-clause-a} also requires the same
bounds with $U$ in place of $\mathbf{U}$. For these we apply \cite[(1.17)]{B-CMP109}, cited as in
Step~2, at the real point $(U,0)$, for which $U'=1$ and the current vanishes. This pair obeys
(i)--(iii) at level $j$ on $X$ with $\alpha_0'$ in the first slot and $0$ in the second; it is the
passage to the real point made in
Lemma~\ref{4.8:lem-class-properties}(b). The configuration $U$ is $G$-valued, so here the estimate is
used at a real configuration. Steps~1 and~2, repeated for $(U,0)$, give
\eqref{4.8:eq-small-pairs-3} with $U$ in place of $\mathbf{U}$. Hence
$(\mathrm{iv})^{\mathfrak{r}}_{p,Z}[\alpha_0L^{-(j-p)}]$ holds.

\emph{Step 4: the geometry used.} The two bounds in \eqref{4.8:eq-small-pairs-3} express the
membership of $U_n^{\mathfrak{r}}$ in the space of \cite[(1.7)--(1.9)]{B-CMP99a} on its top layer,
the $n$-th member $\Omega^{\mathfrak{r}}_n(Z)$ of the tower $\boldsymbol{\Omega}^{\mathfrak{r}}(Z)$,
on which $U_n^{\mathfrak{r}}$ lives. Membership in that space is the plaquette bound together with
the special second-derivative bound, both imposed pointwise on the bonds of that layer. Consequently
the estimate uses, of the geometry, only the inclusion
\begin{equation*}
Z^{\wr p}\subset\Omega^{\mathfrak{r}}_n(Z).
\end{equation*}
By \eqref{4.8:eq-tower-rules-b} the collar of the tower, between $\Omega^{\mathfrak{r}}_n(Z)$ and
the boundary of $Z$, has width at most $R_1M_1$, while the two layers removed from $Z$ in
$Z^{\wr p}$ have depth $2M/L$. The condition $M\ge 2LR_1M_1$ gives $2M/L\ge 4R_1M_1>R_1M_1$; hence
the inclusion.

Since $(p,Z)\in\mathfrak{T}_j(X)$ and $\mathfrak{r}\in\mathfrak{R}$ were arbitrary,
$(\mathrm{iv})^{\Sigma^1}$ holds with the constant $\alpha_0$. The pair $(\mathbf{U},\mathbf{J})$
therefore lies in $U^{\Sigma^1}_j(X;\alpha_0,\alpha_1,\gamma_0)$, which is
\eqref{4.8:eq-small-pairs-2}.

\emph{Step 5: the consequences.} The condition $2B_3\alpha_0'\le\alpha_0$ in
\eqref{4.8:eq-small-pairs-1} is the requirement of \cite[p.~263]{B-CMP109} that $\alpha_0'$ be
sufficiently small at $(j,X)$. Thus no condition beyond those of \cite{B-CMP109} enters.

Since \eqref{4.8:eq-small-pairs-2} supplies membership in $U^{\Sigma^1}_j$ under the same conditions
on $\alpha_0',\alpha_1'$ under which the lemma named in the statement, whose hypothesis is
\cite[(1.17)]{B-CMP109}, obtains membership in $U^c_j$ from that estimate, its clause (a) holds
with $U^{\Sigma^1}_j$ in place of $U^c_j$ and with the same constant $c_T$.

For the same reason, consider the pairs built on the minimal configurations $U_j$ with
$\varepsilon_0$ sufficiently small. By \cite[p.~263]{B-CMP109} they satisfy the conditions
(i)--(iii) with constants $\alpha_0',\alpha_1'$ as above. By \eqref{4.8:eq-small-pairs-2} they
therefore lie in every $U^{\Sigma^1}_j(X;\alpha_0,\alpha_1)$.
\end{proof}

\begin{remark}
If the right-hand sides of \cite[(1.17)]{B-CMP109} at a complex configuration $\mathbf{U}$ were
$2B_3\alpha_0'+C\alpha_1'$ for some constant $C$, in place of $2B_3\alpha_0'$, then the constants
$\mu\alpha_0'$, $\mu\alpha_1'$ obtained at level $p$ by the change of units in the proof of
Proposition~\ref{4.8:prop-small-pairs} would still give, at level $p$, $\mu$ times the corresponding
bound at $(j,X)$, so the scaling part of the argument is unaffected.
\end{remark}

\begin{definition}[The inductive hypothesis on the one-power class]\label{4.8:def-inductive-hypothesis}
Let $j\ge 1$. Let $\mathbf{D}_j$ be the family of localisation domains at level $j$, that is, the finite
wall-connected unions of cubes of $\pi_j$. Let $E^{(j)}(X,g_{j-1},\cdot)$ be the localised terms of the
effective action at level $j$, in the sense of \cite{B-CMP109}, at the coupling $g_{j-1}$. Let
$\alpha_0,\alpha_1$ be Ba{\l}aban's radii, and let $E_0$ and $\kappa$ be the constants of
\cite[(1.18)]{B-CMP109}. For a domain $X$, $U^{\Sigma^1}_j(X;\alpha_0,\alpha_1)$ denotes the one-power
hereditary regularity class of Definition~\ref{4.8:def-one-power-class}, taken with the parameters
$(\alpha_0,\alpha_1)$ and with its third parameter equal to the first.

The \emph{inductive hypothesis on the one-power class at level $j$} is the following statement:
\begin{equation}\label{4.8:eq-inductive-hypothesis-a}\tag{$H^{\Sigma^1}_j$}
\begin{gathered}
\text{for every } X\in\mathbf{D}_j \text{ the function } E^{(j)}(X,g_{j-1},\cdot)
\text{ is defined and analytic on } U^{\Sigma^1}_j(X;\alpha_0,\alpha_1),\\
\text{it is } G^c\text{-invariant there, and }
|E^{(j)}(X,g_{j-1},\cdot)|\le E_0\,e^{-\kappa d_j(X)} \text{ on } U^{\Sigma^1}_j(X;\alpha_0,\alpha_1).
\end{gathered}
\end{equation}
Here $G^c$-invariance is meant in the sense of \cite[(1.19)]{B-CMP109}, and the bound is the bound
\cite[(1.18)]{B-CMP109}, with $d_j(X)$ as there.
\end{definition}

\begin{proposition}[Reproduction of the one-power class under one step]
\label{4.8:prop-class-reproduction}
Let $k$ be a level. Assume the following.
\begin{itemize}
\item The inductive hypothesis $(H^{\Sigma^1}_j)$ of Definition~\ref{4.8:def-inductive-hypothesis}
holds for every $j\le k$.
\item Condition \eqref{4.8:eq-tower-clause-b} of Definition~\ref{4.8:def-tower-clause} holds.
\item $L\ge 2(1+\beta)$.
\item The floors (f2) and (f3) of \eqref{4.8:eq-far-case-a} hold.
\item The $L^2$-floor holds, together with the bound
$(\mathrm{iii})_2$ of the corresponding display, taken as an input exactly as in
Proposition~\ref{4.8:prop-near-case}.
\item The hypotheses of Corollary~\ref{4.8:cor-inclusion-sentences} and of
Propositions~\ref{4.8:prop-near-case}, \ref{4.8:prop-far-case} and~\ref{4.8:prop-small-pairs}
hold, including the statements from which these propositions are derived.
\end{itemize}
Carry out the arguments of \cite[\S\S3--5]{B-CMP109} together with those of
\cite[\S\S1--2]{B-CMP116}, and the steps of this book that rest on them (the operator bounds at
free current, the large-field part of the coupling increment, and the step that reads the running
sentence \eqref{4.8:eq-inclusion-sentences-d} at the radii of \cite[(2.28)]{B-CMP119}), with the
one-power class
$U^{\Sigma^1}$ in place of Ba{\l}aban's regularity space $U^c$ throughout. For
$Y\in\mathbf{D}_k$ call a pair \emph{entering on $Y$} if it is a $\Sigma^1$-representative at
$(k+1,Y;(1+\beta)\alpha_0,(1+\beta)\alpha_1,\alpha_0)$. Fix $Y\in\mathbf{D}_k$ and a level
$j\le k$; in (b)--(e) the domain $X$ ranges over $\mathbf{D}_j$. Then the memberships in the
regularity classes that these arguments read or prove are the following six, and each of them
holds.
\begin{enumerate}
\item[(a)] Near incoming pair. Let $\square$ be a cover cube of \cite[\S3]{B-CMP109} and
$\square_0=\tilde{\square}^5\subset Y$ the window built on it. For every pair entering on $Y$ the
condition \cite[(3.40)]{B-CMP109} holds on $\tilde{\square}^3$, under each of the two rules
$\mathfrak{r}_1,\mathfrak{r}_2\in\mathfrak{R}$.
\item[(b)] Near product, \cite[(3.24), (3.53)]{B-CMP109}. For $X\in\mathbf{D}_j$ with
$X\subset\tilde{\square}^2$, the pair of \cite[Lemma~4]{B-CMP109}, restricted to $X$, lies in
$U^{\Sigma^1}_j(X;\alpha_0,\alpha_1)$. The membership is that of the whole gauge orbit of the
pair; a representative in the class is
the $u_j$-copy $(\mathcal{U},\mathcal{J})$ of Proposition~\ref{4.8:prop-near-case}, with the
decomposition $U=1$.
\item[(c)] Far incoming pair, \cite[(3.16)]{B-CMP109}. Every pair entering on $Y$ restricts, on
every $X\in\mathbf{D}_j$ with $X\subset Y$, to a $\Sigma^1$-representative at
$(j,X;\tfrac12\alpha_0,\tfrac12\alpha_1,\alpha_0)$. In the running form
\eqref{4.8:eq-inclusion-sentences-d} this holds at every level gap $\mathfrak{n}\ge0$.
\item[(d)] Far product. The products formed at \cite[(3.13)]{B-CMP109}, in the form of the
far-product map $\Phi^{\mathbb{K}}$ and of the produced pairs that the operator bounds at
free current form from it, and the two far-type terms of \cite[(3.21)--(3.22)]{B-CMP109}, lie in
$U^{\Sigma^1}_j(X;\alpha_0,\alpha_1)$.
\item[(e)] Kernels and backgrounds. The trivial pair $(\mathbf{U},\mathbf{J})=(1,0)$, at which
the kernels of \cite[(4.4)]{B-CMP109} are evaluated, and the true and small background pairs
covered by Proposition~\ref{4.8:prop-small-pairs}, lie in $U^{\Sigma^1}_j(X;\alpha_0,\alpha_1)$.
\item[(f)] Domain. For $Y\in\mathbf{D}_k$, $X\in\mathbf{D}_{k+1}$ and $Y\subset X$, the
restriction $U^{\Sigma^1}_{k+1}(X;\alpha_0,\alpha_1)\restriction Y$ is contained in the space of
the potentials,
$U^{\Sigma^1}_{k+1}(Y;(1+\beta)\alpha_0,(1+\beta)\alpha_1,\alpha_0)$.
\end{enumerate}
Consequently \cite[Lemma~1]{B-CMP116} holds with the domain of \cite[(1.34)]{B-CMP116} read as
\begin{equation*}
U^{\Sigma^1}_{k+1}\bigl(Y;(1+\beta)\alpha_0,(1+\beta)\alpha_1,\alpha_0\bigr)\times\{B\},
\end{equation*}
the second factor being that of \cite[(1.34)]{B-CMP116}. Moreover, for every
$X\in\mathbf{D}_{k+1}$, the sum
$\mathbf{E}^{(k+1)}(X)$, taken as a whole, is analytic on
$U^{\Sigma^1}_{k+1}(X;\alpha_0,\alpha_1)$. This is the analyticity clause of
$(H^{\Sigma^1}_{k+1})$. Every estimate of the arguments named above holds unchanged.
\end{proposition}

\begin{proof}
We first name the statement that supplies each of the six memberships. We then show that these
memberships are all that the arguments use. Finally we draw the conclusion.

\emph{Membership (a)} is the inclusion sentence \eqref{4.8:eq-inclusion-sentences-b} of
Corollary~\ref{4.8:cor-inclusion-sentences}. It applies to the window $\square_0$ built on a cover
cube, since $\square_0\in\mathbf{D}_k$ and $\square_0\subset Y$. It gives the tower clause on
$\tilde{\square}^3$ at the constant $(1+\beta)\alpha_0$ under both rules of $\mathfrak{R}$. The
condition \cite[(3.40)]{B-CMP109} there is this window clause at the constant
$(1+2\beta)\alpha_0$, and $(1+\beta)\alpha_0<(1+2\beta)\alpha_0$; so it holds, with room
$\beta\alpha_0$.

\emph{Membership (b)} is Proposition~\ref{4.8:prop-near-case}. Its hypotheses are among ours: the
$L^2$-floor, the bound $(\mathrm{iii})_2$ of
the corresponding display, and the restrictions of \cite[Lemma~4]{B-CMP109}. Under them it states
that the $u_j$-copy $(\mathcal{U},\mathcal{J})$, restricted to $X$, is a $\Sigma^1$-representative
at $(j,X;\alpha_0,\alpha_1)$ with $U=1$. Hence the pair of \cite[Lemma~4]{B-CMP109}, restricted to
$X$, lies in the orbit class $U^{\Sigma^1}_j(X;\alpha_0,\alpha_1)$.

\emph{Membership (c)} is the inclusion sentence \eqref{4.8:eq-inclusion-sentences-a}. Its
hypotheses are $L\ge2(1+\beta)$, $j\le k$, $X\in\mathbf{D}_j$, $Y\in\mathbf{D}_k$ and $X\subset Y$,
and all of them hold. For the running form at every gap $\mathfrak{n}\ge0$ it is the running
sentence \eqref{4.8:eq-inclusion-sentences-d}.

\emph{Membership (d)} is Proposition~\ref{4.8:prop-far-case}. It needs a
$\Sigma^1$-representative at half constants as its base, and the base is obtained in two ways.

For the products at \cite[(3.13)]{B-CMP109}, the base is supplied by (c), that is, by
\eqref{4.8:eq-inclusion-sentences-a}.

For the two far-type terms of \cite[(3.21)--(3.22)]{B-CMP109}, the base is taken at the
representative. It is supplied by the half clause of
Proposition~\ref{4.8:prop-near-case}, which holds for every $|B'|<\alpha_3$. Here $B'$ is the proof
variable of \cite[\S3]{B-CMP109} into which $\tau Q'$ is absorbed. The half clause is needed at
every $|B'|<\alpha_3$, and not only at $B'=0$, because the bases of these two terms carry a
non-zero field $\mathbf{A}$. In detail, let $A$ be the $\mathfrak{g}^c$-valued perturbation in the
exponent of the far-type term, as in Proposition~\ref{4.8:prop-far-case}, let $W$ be the
minimiser of \cite[Lemma~4]{B-CMP109}, let
$(\mathcal{U},\mathcal{J})$ be its $u_j$-copy, and set $A^\sharp:=R(u_j^{-1})A$. Then
\begin{equation*}
e^{i\xi_j A}W=\bigl(e^{i\xi_j A^\sharp}\mathcal{U}\bigr)^{u_j}.
\end{equation*}
The operator bounds at free current, in the bound they place on the perturbation interpolated in
the parameter $t_\square$, give $|A(t_\square)|\le\tfrac23\alpha_2$ and
$|\nabla^{\xi_j}A(t_\square)|\le\tfrac23\alpha_2$, the covariant difference being that of those
bounds, whenever $|t_\square|\le r$. The passage from $A$ to
$A^\sharp=R(u_j^{-1})A$ costs at most the factor $1+c_{\mathrm{Ad}}\beta$, by the adjoint bound
\eqref{4.8:eq-near-case-labels-a}; and $(1+c_{\mathrm{Ad}}\beta)\cdot\tfrac23<1$ since
$c_{\mathrm{Ad}}\beta<\tfrac12$. Hence
\begin{equation}\label{4.8:eq-class-reproduction-1}
\max\bigl\{|A^\sharp|,\;|\nabla^{\xi_j}_{\mathcal{U}}A^\sharp|\bigr\}
\le(1+c_{\mathrm{Ad}}\beta)\cdot\tfrac23\alpha_2<\alpha_2 .
\end{equation}
So $A^\sharp$ satisfies the smallness hypothesis $|A|,|\nabla^{\xi_j}_{\mathbf{U}}A|<\alpha_2$ of
Proposition~\ref{4.8:prop-far-case}. By the half clause, for every
$|B'|<\alpha_3$ the pair $(\mathcal{U},\mathcal{J})\restriction X$ is a $\Sigma^1$-representative at
$(j,X;\tfrac12\alpha_0,\tfrac12\alpha_1,\alpha_0)$. Proposition~\ref{4.8:prop-far-case} therefore
applies to $\mathcal{U}$ with the perturbation $A^\sharp$, and it places
$e^{i\xi_j A^\sharp}\mathcal{U}$ in the class. The orbit of this product is, by the displayed
identity, the orbit of $e^{i\xi_j A}W$, and so the far-type terms lie in
$U^{\Sigma^1}_j(X;\alpha_0,\alpha_1)$.

Within the operator bounds at free current, the far products take their base from
\cite[(3.16)]{B-CMP109}, that is, from \eqref{4.8:eq-inclusion-sentences-a}. There the half clause
is not needed.

\emph{Membership (e)} rests on two statements. For the true and small background pairs it is
Proposition~\ref{4.8:prop-small-pairs}. For the trivial pair $(1,0)$, at which the kernels of
\cite[(4.4)]{B-CMP109} are evaluated, it is either
that proposition, or Proposition~\ref{4.8:prop-far-case} applied at the base $(1,0)$.

\emph{Membership (f)} is the inclusion sentence \eqref{4.8:eq-inclusion-sentences-c}. Its
hypotheses are $Y\in\mathbf{D}_k$, $X\in\mathbf{D}_{k+1}$ and $Y\subset X$.

\emph{The memberships used.} Every bound of \cite[\S\S3--5]{B-CMP109} and \cite[\S\S1--2]{B-CMP116},
and of the steps resting on them, is of one of two kinds.

The first kind is the supremum, over produced pairs, of the bound of $(H^{\Sigma^1}_j)$. The
produced pairs lie in the domain for all values of the interpolation (Cauchy) parameters of
\cite[\S3]{B-CMP109}, of \cite[\S\S1--2]{B-CMP116} and of the operator bounds at free current,
among them $\tau$ and the interpolation parameter $t_\square$ of the operator bounds at free
current. These parameters move $A$ and $B'$
about a fixed incoming point, inside the
region $|A|<\alpha_2$, $|B'|<\alpha_3$. On that region
Propositions~\ref{4.8:prop-near-case} and~\ref{4.8:prop-far-case} hold uniformly. No radius
shrinks, because the ceiling $\alpha_2\le\alpha_0/C_{IV}$ of Proposition~\ref{4.8:prop-far-case}
is the same, with the same $C_{IV}$, for every such point.

The second kind uses only conditions (i)--(iii) of \cite[\S1]{B-CMP109} on the domain, with
constants
$\alpha_0'$, $\alpha_1'$ much larger than $\alpha_0$, $\alpha_1$. This is the domain sentence of
\cite[\S\S1--2]{B-CMP116}, which is membership (f) above.

The real point $(U,0)$ used in the perturbative argument of \cite[\S\S1--2]{B-CMP116} lies in the
class by
Lemma~\ref{4.8:lem-class-properties}(b). No step uses the literal class contrapositively: a
smaller domain of analyticity enters only through memberships, and all of these are among (a)--(f).
Therefore every estimate of these arguments holds unchanged with $U^{\Sigma^1}$ in place of $U^c$.
In particular \cite[Lemma~1]{B-CMP116} holds with the domain of \cite[(1.34)]{B-CMP116} read as
stated.

\emph{Analyticity.} The sum $\mathbf{E}^{(k+1)}(X)$ is obtained by composing the analytic maps of
these arguments. Each composition is taken on a set that is open, by
Lemma~\ref{4.8:lem-class-properties}(c). Hence $\mathbf{E}^{(k+1)}(X)$ is analytic on
$U^{\Sigma^1}_{k+1}(X;\alpha_0,\alpha_1)$.

\emph{Invariance.} The class is a union of $G^c$-orbits, so the $G^c$-invariance required in
$(H^{\Sigma^1}_{k+1})$ holds on it unchanged.

\emph{Base of the induction.} At the base $j=1$ no term $E^{(j')}$ of an earlier level $j'<j$
enters.
\end{proof}

Consider the comparison of the one-power class with the literal and two-power classes
(Remark~\ref{4.8:rem-class-comparison}). Read it with the tower clause imposed at every label
$(p,Z)$ of the tower, under both rules $\mathfrak{r}_1$ and $\mathfrak{r}_2$, and with collars of
two layers of the label's own cubes. Then it places the two-power classes inside the one-power
class. Hence the analyticity and the decay bound \cite[(1.18)]{B-CMP109} asserted in
$(H^{\Sigma^1}_j)$ hold on these smaller spaces as well. The localised boundary terms, which live
on domains of the type of $\tilde U^c$,
are not covered by Proposition~\ref{4.8:prop-class-reproduction}. These terms are treated by
Corollary~\ref{4.4:cor-localised-independence}, not here.

\begin{lemma}[The perturbation bound on the top layer]\label{4.8:lem-top-layer-bound}
Let $0\le j\le k$, and let $\alpha_{0,j},\alpha_{1,j}$ be the radii of \cite[(2.28)]{B-CMP119} at
level $j$, built with the constants written $C_0$ and $C_1$ in that display; we write $C_1'$ for
the second of them, to keep it apart from the constant $C_1$ of Proposition~\ref{4.8:prop-far-case}.
Consider the product
$\mathbf{U}^\flat=e^{i\xi\mathcal{H}_w}\mathbf{U}$, and let $A:=\mathcal{H}_w$ be its twist field on
the top layer $\bar n=k$, written in level-$j$ units; $\nabla_{\mathbf{U}}$ denotes the covariant
lattice derivative in the same units. Let $C_{IV}=C_{IV}(d,L)$ be the constant of the ceiling on the
perturbation size in Proposition~\ref{4.8:prop-far-case} (written there
$\alpha_2\le\min\{\alpha_1/8,\alpha_0/C_{IV}\}$ for a generic size $\alpha_2$). Let $w_k$
be the flow amplitude of the twist field $\mathcal{H}_w$ at level $k$, and $C_1^{\mathrm{fl}}$ the
flow constant entering $w_k$; let $C_H$, $\|C\|$, $A_1$ and $\beta_0$ be the constants entering the
layered amplitude bounds of $\mathcal{H}_w$ and the bound of the ratio of $w_k$ to the radii of
\cite[(2.28)]{B-CMP119}, both recalled in the proof below; and let $c_\flat$, $\theta$ and $\beta$
be the constants of the floor on $C_0$ and $C_1'$ imposed in the same construction as the layered
amplitude bounds of $\mathcal{H}_w$; this floor is recalled in \eqref{4.8:eq-top-layer-bound-1}
below. Then
\begin{equation}\label{4.8:eq-top-layer-bound-a}
\begin{gathered}
|A|,\ |\nabla_{\mathbf{U}}A|\ \le\ L^{-(k-j)}C_Hw_k\ <\
\min\Bigl\{\frac{\alpha_{1,j}}{8},\ \frac{\alpha_{0,j}}{C_{IV}}\Bigr\}\\
\Longleftarrow\quad
C_1'>16\,C_H(1+\beta_0)\|C\|C_1^{\mathrm{fl}}A_1
\quad\text{and}\quad
C_0>2\,C_{IV}C_H(1+\beta_0)\|C\|C_1^{\mathrm{fl}}A_1 .
\end{gathered}
\end{equation}
Neither of the two conditions in \eqref{4.8:eq-top-layer-bound-a} involves $k$ or $j$; they are
lower bounds on $C_1'$ and $C_0$ in terms of $C_{IV}$, $C_H$, $\|C\|$, $C_1^{\mathrm{fl}}$, $A_1$ and
$\beta_0$. Of these, $C_{IV}=C_{IV}(d,L)$ depends only on $d$ and $L$, and $C_1^{\mathrm{fl}}$ is
fixed before $C_0$; and $C_0$, $C_1'$ (the latter written $C_1$ there) are required in
\cite{B-CMP119} only to be sufficiently large positive numbers, so that they may be increased
freely. The two conditions are therefore floors on freely increasable constants, which can be met
uniformly in $k$ and $j$, and not bounds on the depth $k-j$. The first of them is implied by the floor
\begin{equation}\label{4.8:eq-top-layer-bound-1}
C_0,\ C_1'\ \ge\ \frac{8c_\flat(1+\beta_0)\|C\|C_1^{\mathrm{fl}}A_1C_H^2}{(1-\theta)\beta}
\end{equation}
together with the standing bounds $c_\flat\ge4$, $C_H\ge1$ and $(1-\theta)\beta\le\frac14$.
\end{lemma}

\begin{proof}
Write $\alpha_2:=L^{-(k-j)}C_Hw_k$.

\emph{The first inequality.} We apply the first two members of the layered amplitude bounds for
the twist field $\mathcal{H}_w$ on the top layer, that is at $\bar n=k$. There the depth term
vanishes, $D_k=0$, and the length of the layer is $\ell_k=L^{k-j}$ in level-$j$ units. The first
member then gives $|A|\le C_Hw_kL^{-(k-j)}=\alpha_2$. The second member bounds the covariant
derivative $\nabla_{\mathbf{U}}A$ directly and gives
\begin{equation*}
|\nabla_{\mathbf{U}}A|\ \le\ C_Hw_kL^{-2(k-j)}\ \le\ C_Hw_kL^{-(k-j)}\ =\ \alpha_2 ,
\end{equation*}
the middle inequality because $L\ge1$ and $k-j\ge0$.

\emph{The ratio of the amplitude to the radii.} Write $C_{(0)}:=C_0$ and $C_{(1)}:=C_1'$. The
bound of the ratio of the flow amplitude $w_k$ to the radii of \cite[(2.28)]{B-CMP119} states that,
for $\iota\in\{0,1\}$,
\begin{equation}\label{4.8:eq-top-layer-bound-2}
\frac{w_k}{\alpha_{\iota,j}}\ \le\ \frac{2\|C\|C_1^{\mathrm{fl}}A_1}{C_{(\iota)}}\,
(1+\beta_0)\,(1+k-j)^{1/2}.
\end{equation}
The factor $(1+k-j)^{1/2}$ grows with the depth $k-j$, but it is compensated by the factor
$L^{-(k-j)}$ in $\alpha_2$. Indeed, with $n:=k-j\ge0$ we have $1+n\le 2^n$ (Bernoulli's
inequality), hence $(1+n)^{1/2}\le 2^{n/2}\le L^n$ since the blocking factor $L$ of
Definition~\ref{1.1:def-lattice} is an odd integer larger than one, so that $L\ge3\ge2$; thus
\begin{equation}\label{4.8:eq-top-layer-bound-3}
(1+k-j)^{1/2}L^{-(k-j)}\ \le\ 1\qquad\text{for every }k-j\ge0 .
\end{equation}

\emph{The two comparisons.} Multiplying \eqref{4.8:eq-top-layer-bound-2} by $C_HL^{-(k-j)}$ and
using \eqref{4.8:eq-top-layer-bound-3}, we obtain for $\iota\in\{0,1\}$
\begin{equation*}
\frac{\alpha_2}{\alpha_{\iota,j}}
\ =\ \frac{C_HL^{-(k-j)}w_k}{\alpha_{\iota,j}}
\ \le\ \frac{2C_H(1+\beta_0)\|C\|C_1^{\mathrm{fl}}A_1}{C_{(\iota)}} .
\end{equation*}
For $\iota=0$ the right-hand side is strictly smaller than $1/C_{IV}$ exactly when
$C_0>2C_{IV}C_H(1+\beta_0)\|C\|C_1^{\mathrm{fl}}A_1$, the second condition in
\eqref{4.8:eq-top-layer-bound-a}; under it $\alpha_2<\alpha_{0,j}/C_{IV}$. For $\iota=1$ the
right-hand side is strictly smaller than $1/8$ exactly when
$C_1'>16C_H(1+\beta_0)\|C\|C_1^{\mathrm{fl}}A_1$, the first condition in
\eqref{4.8:eq-top-layer-bound-a}; under it $\alpha_2<\alpha_{1,j}/8$. Together with the first
inequality this proves \eqref{4.8:eq-top-layer-bound-a}. The two conditions are uniform in $k$
and $j$ because \eqref{4.8:eq-top-layer-bound-3} holds for every depth; in particular no upper bound
on the depth $k-j$ is imposed.

\emph{The first condition follows from the floor.} By the standing bounds of the statement,
$c_\flat\ge4$, $C_H\ge1$ and $(1-\theta)\beta\le\frac14$.
Hence
\begin{equation*}
\frac{8c_\flat C_H^2}{(1-\theta)\beta}\ \ge\ 8\cdot4\cdot4\,C_H^2\ =\ 128\,C_H^2
\ \ge\ 128\,C_H\ >\ 16\,C_H ,
\end{equation*}
and multiplying by the positive number $(1+\beta_0)\|C\|C_1^{\mathrm{fl}}A_1$ shows that
\eqref{4.8:eq-top-layer-bound-1} for $C_1'$ implies
$C_1'>16C_H(1+\beta_0)\|C\|C_1^{\mathrm{fl}}A_1$. The floor \eqref{4.8:eq-top-layer-bound-1}
does not contain $C_{IV}$, so the second condition in \eqref{4.8:eq-top-layer-bound-a} is an
additional requirement beside it.

\emph{Use of the bound.} Since the second inequality in \eqref{4.8:eq-top-layer-bound-a} is strict,
the open interval
\begin{equation*}
\Bigl(L^{-(k-j)}C_Hw_k,\ \min\bigl\{\alpha_{1,j}/8,\ \alpha_{0,j}/C_{IV}\bigr\}\Bigr)
\end{equation*}
is non-empty, and for every $\alpha_2'$ in it the field $A$ satisfies $|A|,|\nabla_{\mathbf{U}}A|<
\alpha_2'$ and $\alpha_2'$ satisfies the ceiling $\alpha_2'\le\min\{\alpha_{1,j}/8,\alpha_{0,j}/C_{IV}\}$
of Proposition~\ref{4.8:prop-far-case}. This is the form in which the product
$\mathbf{U}^\flat=e^{i\xi\mathcal{H}_w}\mathbf{U}$ is supplied by Proposition~\ref{4.8:prop-far-case}
at $A:=\mathcal{H}_w$ in level-$j$ units. The hypothesis of Proposition~\ref{4.8:prop-far-case}
on the pair, namely that $(\mathbf{U},\mathfrak{s}\mathbf{J})$ restricted to the
domain $X\in\mathbf{D}_j$ is a $\Sigma^1$-representative at level $j$ with the halved radii
$\frac12\alpha_{0,j}$, $\frac12\alpha_{1,j}$, is supplied by the running inclusion
\eqref{4.8:eq-inclusion-sentences-d} at the radii of \cite[(2.28)]{B-CMP119}. Here $\mathfrak{s}$
is the rescaling factor of the current in \eqref{4.8:eq-inclusion-sentences-d}. The other ceilings
of Proposition~\ref{4.8:prop-far-case} and its floors are conditions on the constants and are not
the subject of this lemma.
\end{proof}

The straddling producer of the large-field construction is the configuration treated here. The one-power class $U^{\Sigma^1}_j(X;\cdot)$ asks for clause (iv) under both tower rules $\mathfrak{r}_1,\mathfrak{r}_2$ of Definition~\ref{4.8:def-tower-rules}. Some configurations take their source in one of the regularity spaces of \cite[(2.34)--(2.39)]{B-CMP119} or \cite[(1.64)--(1.69)]{B-CMP122a}. Such a configuration carries the clauses of the second rule only through the definition of that space. Yet it feeds an old term of the expansion of \cite{B-CMP109}, whose domain is now $U^{\Sigma^1}_j(X)$. The straddling producer is one such configuration. The following lemma supplies the perturbation bounds on which the membership argument for this producer's target space is built; that argument applies Proposition~\ref{4.8:prop-far-case} and the uniform-increment lemma on which that proposition rests, and is described after the proof.

\begin{lemma}[the layerwise perturbation bound for straddling sources]
\label{4.8:lem-layerwise-bound}

\emph{Setting.} Let $j\le k$ and let $Y$ be a straddling domain at level $k+1$. By this we mean a domain of that level in the straddling case of the large-field construction; the case $Y\subset\Lambda_{k+1}$ is that of the top layer, treated in Lemma~\ref{4.8:lem-top-layer-bound}. Let $(\mathbf{U},\mathbf{J})$ lie in the post-step space $\mathfrak{U}^{\mathrm{post}}_{k+1}(Y)$ of Definition~\ref{4.4:def-post-step-space}, that is, in the space of \cite[(2.34)--(2.39)]{B-CMP119}. Assume that $(\mathbf{U},\mathbf{J})$ satisfies the clauses of the one-power class $\Sigma^1$ at every sub-domain of an old term of \cite{B-CMP109}, under both rules $\mathfrak{r}_1,\mathfrak{r}_2$.

Recall the nested sets
\begin{equation*}
\Lambda_{k+1}\subset\Omega_{k+1}\subset\Lambda_k\subset\cdots\subset\Lambda_j .
\end{equation*}
Let $X\in\mathbf{D}_j$ be any domain with $X\subset\Lambda_j$. No further restriction is placed on $X$, in accordance with the range of the $X$-sum in \cite[p.~259]{B-CMP119}. The target space is $U^c_j(X,\alpha_{0,j},\alpha_{1,j})$ of \cite[(2.27)(ii)]{B-CMP119}, that is, the space defined by \cite[(1.11)--(1.16)]{B-CMP109} with the radii $\alpha_{0,j},\alpha_{1,j}$.

Write $A:=\mathcal{H}_w|_X$, in level-$j$ units, for the restriction to $X$ of the $\mathfrak{g}^c$-valued field $\mathcal{H}_w$ of the straddling producer.

In the large-field construction the nested sequence above is divided into layers, indexed by $\bar n$ with $j\le\bar n\le k$. The layer of index $\bar n$ is the part of $\Lambda_j$ that the construction assigns to level $\bar n$; in particular $X\cap\Omega_k$ lies in the layer of index $k$. The unit of length on the layer of index $\bar n$ is $\ell_{\bar n}=L^{\bar n-j}$ level-$j$ units. On that layer,
\begin{equation*}
D_{\bar n}:=\delta_H\, \operatorname{dist}\bigl(\text{layer }\bar n,\ \Omega_{k+1}\cup S_{k+1}\bigr).
\end{equation*}
Here $\operatorname{dist}$ is the lattice distance, taken in the units used by the large-field construction, and the set $S_{k+1}$ is that of the construction. The region $\Omega_{k+1}\cup S_{k+1}$ is the one from which the decay of $\mathcal{H}_w$ is measured. The constant $\delta_H$ is the decay rate of the layered decay bound \eqref{4.8:eq-layerwise-bound-4} below, and $\mathsf{w}_\star$ is the constant of the construction for which $D_{\bar n}\ge\delta_H\mathsf{w}_\star(k-\bar n)$ holds. The letter $d$ denotes the dimension throughout.

The source's layer-$j$ factor is $F_j=f_{k+1,j}+\theta\beta 2^{-(k+1-j)}$. Here $\beta$ is the schedule constant of the layered analyticity spaces (Definition~\ref{4.4:def-post-step-space}). The constant $\theta$ and the coefficient $f_{k+1,j}$ are those entering the layer-$j$ factor of those spaces. On the layer of index $j$ the space $\mathfrak{U}^{\mathrm{post}}_{k+1}(Y)$ imposes on the source the clauses of the target space with the radii $F_j\alpha_{0,j}$, $F_j\alpha_{1,j}$ in place of $\alpha_{0,j}$, $\alpha_{1,j}$. The factor satisfies $1-F_j\ge(1-\theta)\beta/2$; the floors below presuppose $(1-\theta)\beta>0$.

Assume the floor of the large-field construction
\begin{equation}\label{4.8:eq-layerwise-bound-1}
2\sqrt2\, e^{-\delta_H\mathsf{w}_\star}\le 1 .
\end{equation}
Assume also the two floors of \eqref{4.8:eq-top-layer-bound-a}:
\begin{equation*}
C_1'>16C_H(1+\beta_0)\|C\|C_1^{\mathrm{fl}}A_1,\qquad
C_0>2C_{IV}C_H(1+\beta_0)\|C\|C_1^{\mathrm{fl}}A_1 .
\end{equation*}
Assume finally that the parameters $q_0,q_1$ of the radii $\alpha_{0,\cdot},\alpha_{1,\cdot}$ of \cite[(2.28)]{B-CMP119} satisfy $q_0,q_1\ge p_0$, the condition of Lemma~\ref{4.4:lem-radii-ratio}.
Then the following hold.
\begin{enumerate}
\item[(a)] On $X$ one has $|A|\le\alpha_2$ and $|\nabla_{\mathbf{U}}A|\le\alpha_2$, where
\begin{equation}\label{4.8:eq-layerwise-bound-a}
\alpha_2:=C_Hw_k\max_{j\le\bar n\le k}e^{-D_{\bar n}}L^{-(\bar n-j)}
<\min\Bigl\{\frac{\alpha_{1,j}}{8},\ \frac{\alpha_{0,j}}{C_{IV}}\Bigr\}.
\end{equation}
More precisely, for $C_\cdot=C_0$ and for $C_\cdot=C_1'$, with $\alpha_{\cdot,j}$ the corresponding radius,
\begin{equation}\label{4.8:eq-layerwise-bound-2}
\frac{\alpha_2}{\alpha_{\cdot,j}}\le\frac{2C_H(1+\beta_0)\|C\|C_1^{\mathrm{fl}}A_1}{C_\cdot}.
\end{equation}
The floors \eqref{4.8:eq-layerwise-bound-1} and those of \eqref{4.8:eq-top-layer-bound-a} used here do not depend on $k$, $j$ or $\bar n$.
\item[(b)] Consider the labels $(p,Z)\in\mathfrak{T}_j(X)$ with $Z\subset X\cap\Omega_{j+1}$. At these labels Proposition~\ref{4.8:prop-far-case} applies with this $A$ and with its perturbation radius equal to
\begin{equation*}
\alpha_2^{\mathrm{far}}:=\min\{\alpha_{1,j}/8,\ \alpha_{0,j}/C_{IV}\},
\end{equation*}
provided its other hypotheses hold. These are the hypothesis that $(\mathbf{U},\mathbf{J})$ is a $\Sigma^1$-representative at $(j,X;\tfrac12\alpha_{0,j},\tfrac12\alpha_{1,j},\cdot)$, its ceilings, as stated there at the radii $\alpha_{0,j},\alpha_{1,j}$, and its floors, namely $c_{12}L\ge2$, $L\ge4d$ and \eqref{4.8:eq-far-case-a}; they are inputs of (b) and are not derived in the proof below.
\item[(c)] Suppose, in addition,
\begin{equation}\label{4.8:eq-layerwise-bound-b}
C_0>\frac{2C_{IV}C_H(1+\beta_0)\|C\|C_1^{\mathrm{fl}}A_1}{(1-\theta)\beta},\qquad
C_1'>\frac{16C_H(1+\beta_0)\|C\|C_1^{\mathrm{fl}}A_1}{(1-\theta)\beta}.
\end{equation}
Then
\begin{equation}\label{4.8:eq-layerwise-bound-3}
F_j\alpha_{0,j}+\tfrac12C_{IV}\alpha_2\le\alpha_{0,j},\qquad
F_j\alpha_{1,j}+4\alpha_2\le\alpha_{1,j}.
\end{equation}
These are the radius inequalities required when the uniform-increment lemma on which Proposition~\ref{4.8:prop-far-case} rests is applied at every label $(p,Z)$ with $p\le j$ and $Z$ meeting $\Omega_j\setminus\Omega_{j+1}$. There that lemma is used in two of its parts, the bound on the increment of the tower quantities uniform over the labels and the pointwise bound for clause (ii), both read additively against $F_j$.
\end{enumerate}
The floors \eqref{4.8:eq-layerwise-bound-b} are independent of $k$. They bound $C_0$ and $C_1'$ from below only, and both constants may be increased freely. They do not interfere with the existence of $L_0$.
\end{lemma}

\begin{proof}
\emph{Proof of (a).} Fix $\bar n$ with $j\le\bar n\le k$. We use, as an input of this lemma, the first two members of the layered decay bound of the large-field construction, stated there for the field $\mathcal{H}$ of the straddling producer with source term $\mathsf{b}$. On a layer of index $n$, with unit $\ell_n$, they read
\begin{equation}\label{4.8:eq-layerwise-bound-4}
\ell_n|\mathcal{H}|(y),\ \ell_n^2|\nabla\mathcal{H}|(y)
\le C_H\sup_{y'}\bigl[|\mathsf{b}(y')|\,e^{-\delta_H \operatorname{dist}(y,y')}\bigr],
\end{equation}
with $C_H=C_H(d,L)\ge1$ and $|\mathsf{b}|\le w_k$. The second member is the covariant derivative $\nabla_{\mathbf{U}}$ itself, so no conversion between derivatives is needed. We apply these bounds with $n=\bar n$.

In the large-field construction the distances $D_{\bar n}$ are measured from the region $\Omega_{k+1}\cup S_{k+1}$ carrying $\mathsf{b}$. Since $|\mathsf{b}|\le w_k$, the supremum in \eqref{4.8:eq-layerwise-bound-4} is therefore at most $w_ke^{-D_{\bar n}}$ for $y$ in the layer of index $\bar n$. On that layer $\ell_{\bar n}=L^{\bar n-j}$ level-$j$ units. Dividing by $\ell_{\bar n}=L^{\bar n-j}$ gives, in level-$j$ units,
\begin{equation*}
|\mathcal{H}_w|\le C_Hw_ke^{-D_{\bar n}}L^{-(\bar n-j)}.
\end{equation*}
Dividing by $\ell_{\bar n}^2=L^{2(\bar n-j)}$ gives
\begin{equation*}
|\nabla_{\mathbf{U}}\mathcal{H}_w|\le C_Hw_ke^{-D_{\bar n}}L^{-2(\bar n-j)}
\le C_Hw_ke^{-D_{\bar n}}L^{-(\bar n-j)} ,
\end{equation*}
where the last step uses $L\ge1$. Since $X\subset\Lambda_j$, the domain $X$ is covered by the layers $\bar n=j,\dots,k$. Taking the maximum over these layers gives $|A|,|\nabla_{\mathbf{U}}A|\le\alpha_2$ on $X$, with $\alpha_2$ as in \eqref{4.8:eq-layerwise-bound-a}.

Two facts about the distances are used; both are part of the layer structure of the large-field construction. First, $D_k=0$. Second, $D_{\bar n}\ge\delta_H\mathsf{w}_\star(k-\bar n)$.

We now invoke Lemma~\ref{4.4:lem-radii-ratio} at $n=j$. It holds for $q_0,q_1\ge p_0$, for $C_\cdot=C_0$ and $C_\cdot=C_1'$, with its constant $C_1$ equal to $C_1^{\mathrm{fl}}$, and gives
\begin{equation}\label{4.8:eq-layerwise-bound-5}
\frac{w_k}{\alpha_{\cdot,j}}\le\frac{2\|C\|C_1^{\mathrm{fl}}A_1}{C_\cdot}\,(1+\beta_0)\,(1+k-j)^{1/2}.
\end{equation}
Put $a:=k-\bar n\ge0$ and $b:=\bar n-j\ge0$; both are integers, and $a+b=k-j$. Since
\begin{equation*}
(1+a)(1+b)=1+a+b+ab\quad\text{and}\quad ab\ge0,
\end{equation*}
we have $1+k-j\le(1+a)(1+b)$.

Next, $1+a\le2^a$, so $(1+a)^{1/2}\le(\sqrt2)^a$. By \eqref{4.8:eq-layerwise-bound-1}, $\sqrt2\,e^{-\delta_H\mathsf{w}_\star}\le1/2$. Combining these facts with $D_{\bar n}\ge\delta_H\mathsf{w}_\star a$, we obtain
\begin{equation*}
(1+a)^{1/2}e^{-D_{\bar n}}
\le(1+a)^{1/2}e^{-\delta_H\mathsf{w}_\star a}
\le\bigl(\sqrt2\,e^{-\delta_H\mathsf{w}_\star}\bigr)^a
\le2^{-a}\le1 .
\end{equation*}
Further, $L\ge2$ gives $1+b\le4^b\le L^{2b}$, and therefore $(1+b)^{1/2}L^{-b}\le1$.

Now insert \eqref{4.8:eq-layerwise-bound-5} into the definition of $\alpha_2$ and split $(1+k-j)^{1/2}$ as above. This yields
\begin{align*}
\frac{\alpha_2}{\alpha_{\cdot,j}}
&\le C_H\,\frac{2\|C\|C_1^{\mathrm{fl}}A_1}{C_\cdot}\,(1+\beta_0)
\max_{j\le\bar n\le k}\Bigl[(1+a)^{1/2}e^{-D_{\bar n}}\,(1+b)^{1/2}L^{-b}\Bigr]\\
&\le C_H\,\frac{2\|C\|C_1^{\mathrm{fl}}A_1}{C_\cdot}\,(1+\beta_0).
\end{align*}
This is \eqref{4.8:eq-layerwise-bound-2}.

For $C_\cdot=C_1'$, the first floor of \eqref{4.8:eq-top-layer-bound-a} states that the right-hand side is strictly less than $1/8$. For $C_\cdot=C_0$, the second floor states that it is strictly less than $1/C_{IV}$. So the strict inequality in \eqref{4.8:eq-layerwise-bound-a} follows from the strict floors.

These are the same floors as in the top-layer bound (Lemma~\ref{4.8:lem-top-layer-bound}). No new floor enters, because \eqref{4.8:eq-layerwise-bound-1} is already a floor of the large-field construction. It is equivalent to $\delta_H\mathsf{w}_\star\ge\log(2\sqrt2)$. It is implied by the lower bound $\delta_HR_HM\ge\log(2\sqrt2)$ on $M$ already imposed in that construction. Here $R_H$ denotes the constant of that construction that enters this lower bound. None of the bounds depends on $k$, $j$ or $\bar n$.

\emph{Proof of (b).} Proposition~\ref{4.8:prop-far-case} is applied at $A=\mathcal{H}_w|_X$. Its perturbation radius is taken to be $\alpha_2^{\mathrm{far}}$, the largest value its ceiling $\alpha_2\le\min\{\alpha_1/8,\alpha_0/C_{IV}\}$ allows at the radii $\alpha_{0,j},\alpha_{1,j}$. The proposition requires the strict bounds
\begin{equation*}
|A|<\alpha_2^{\mathrm{far}},\qquad |\nabla_{\mathbf{U}}A|<\alpha_2^{\mathrm{far}} .
\end{equation*}
Both hold, because by (a)
\begin{equation*}
|A|,\ |\nabla_{\mathbf{U}}A|\le\alpha_2<\alpha_2^{\mathrm{far}} .
\end{equation*}

\emph{Proof of (c).} Let $(p,Z)$ be a label with $p\le j$ and $Z$ meeting $\Omega_j\setminus\Omega_{j+1}$. Write $\mu:=L^{-(j-p)}$. By Theorem~\ref{4.8:thm-restriction-units}, in the form \eqref{4.8:eq-restriction-units-a}, the source's factor $F_j$ survives restriction to $(p,Z)$ in the one-power slots. The restricted source therefore has radii $\mu F_j\alpha_{0,j}$ and $\mu F_j\alpha_{1,j}$ at the label $(p,Z)$. The ceilings of the uniform-increment lemma are read at these source radii.

The uniform-increment lemma on which Proposition~\ref{4.8:prop-far-case} rests, through its bound on the increment of the tower quantities uniform over the labels and its pointwise bound for clause (ii), adds $\tfrac12C_{IV}\alpha_2$ to the $\alpha_0$-radius and $4\alpha_2$ to the $\alpha_1$-radius. At the label $(p,Z)$ the factors $\mu$ and $\lambda'^2$ of the uniform-increment lemma may occur in the inequalities of \eqref{4.8:eq-layerwise-bound-3}. In each inequality in which they occur, they multiply both members and cancel. It therefore suffices to prove \eqref{4.8:eq-layerwise-bound-3} at level $j$.

The pointwise form's increment $4\alpha_2$ is applied at the source radius. Since $1-F_j\ge(1-\theta)\beta/2>0$, we have $F_j\le1$, and hence
\begin{equation*}
F_j\alpha_{1,j}\le\alpha_{1,j}\le\tfrac12 .
\end{equation*}
So that source radius lies within the ceiling of the uniform-increment lemma.

For the first inequality, combine \eqref{4.8:eq-layerwise-bound-2} at $C_\cdot=C_0$ with the first floor of \eqref{4.8:eq-layerwise-bound-b}, and then use $1-F_j\ge(1-\theta)\beta/2$:
\begin{align*}
\tfrac12C_{IV}\alpha_2
&\le\tfrac12C_{IV}\,\alpha_{0,j}\,
\frac{2C_H(1+\beta_0)\|C\|C_1^{\mathrm{fl}}A_1}{C_0}\\
&<\tfrac12\,(1-\theta)\beta\,\alpha_{0,j}
\le(1-F_j)\,\alpha_{0,j}.
\end{align*}
For the second inequality, combine \eqref{4.8:eq-layerwise-bound-2} at $C_\cdot=C_1'$ with the second floor of \eqref{4.8:eq-layerwise-bound-b}:
\begin{align*}
4\alpha_2
&\le4\,\alpha_{1,j}\,
\frac{2C_H(1+\beta_0)\|C\|C_1^{\mathrm{fl}}A_1}{C_1'}\\
&<\tfrac12\,(1-\theta)\beta\,\alpha_{1,j}
\le(1-F_j)\,\alpha_{1,j}.
\end{align*}
Rearranging gives \eqref{4.8:eq-layerwise-bound-3}.

The floors \eqref{4.8:eq-layerwise-bound-b} are of the same form as the lower bound on $C_0$ and $C_1'$ imposed in the large-field construction, namely
\begin{equation*}
C_0,\ C_1'\ge\frac{8c_\flat(1+\beta_0)\|C\|C_1^{\mathrm{fl}}A_1C_H^2}{(1-\theta)\beta},
\end{equation*}
where $c_\flat$ is the numerical constant of that lower bound in the large-field construction, which satisfies $c_\flat\ge4$.
The second of them follows from that bound, since $c_\flat\ge4$ and $C_H\ge1$ give
\begin{equation*}
8c_\flat C_H^2\ge32C_H^2\ge32C_H>16C_H .
\end{equation*}
\end{proof}

The member $\bar n=k$ of \eqref{4.8:eq-layerwise-bound-a} is the bound of Lemma~\ref{4.8:lem-top-layer-bound}. That lemma concerns $Y\subset\Lambda_{k+1}$, that is, the top layer only.

For a general $X\in\mathbf{D}_j$ with $X\subset\Lambda_j$, the maximum over $\bar n$ cannot be replaced by the member $\bar n=j$. Indeed, $e^{-D_j}$ is the smallest of the exponentials. On $X\cap\Omega_k$, however, $D_k=0$, and the layered decay bound gives only $C_Hw_kL^{-(k-j)}$, nothing smaller.

Under the further floor $\delta_HR_HM\ge\log L$ on $M$, one could instead bound $e^{-D_{\bar n}}L^{-(\bar n-j)}$ by $L^{-(k-j)}$ for every $\bar n$. No such floor is imposed above.

The hypothesis that the source satisfies the clauses of the one-power class at every sub-domain is an input of Lemma~\ref{4.8:lem-layerwise-bound}; it is not established in this section. For true backgrounds it is supplied by Proposition~\ref{4.8:prop-small-pairs}, whose argument does not distinguish the two rules. For images of the large-field operation $\mathbf{T}$ it is supplied by the representation and bound of the large-field terms, applied along the segment of data obtained by varying the twist parameter $w$ of the step (Definition~\ref{4.4:def-twisted-argument}). That statement is made for every admissible sequence in the sense of \cite[(1)]{B-CMP102b}.

This section does not establish as a theorem the membership of the straddling producer, nor that of the other producers whose sources lie in the later papers' classes. The adjunction of a tower rule at no cost, stated beside \eqref{4.8:eq-tower-rules-b}, holds only for the producer propositions of this section. These are Propositions~\ref{4.8:prop-far-case} and \ref{4.8:prop-small-pairs}, and the producer proposition for Ba{\l}aban's representative at the minimiser of \cite[(3.37)]{B-CMP109}.

The right-hand two-power cells of the class comparison \eqref{4.8:eq-class-comparison-a} need no producer. Read them in the sense in which clause (iv) holds at every label $(p,Z)$ of the tower under both rules $\mathfrak{r}_1,\mathfrak{r}_2$, with collars of two layers of the label's own cubes. Under that reading both cells are contained in $U^{\Sigma^1}_j(X;\cdot)$ by \eqref{4.8:eq-class-comparison-a} itself (Remark~\ref{4.8:rem-class-comparison}).

A further configuration feeds an old term of the expansion of \cite{B-CMP109} from a class of the later papers. It arises from Ba{\l}aban's construction in \cite[p.~375]{B-CMP122b}, whose notation is used in this paragraph and the next: $Y_4$ and $Y_1$ are the domains, $\tilde\alpha_0,\tilde\alpha_1$ the radii and $U^0_k$ the map of that page. There a source
\begin{equation*}
(\mathbf{U},\mathbf{J})\in\tilde U^c_k(Y_4,\tilde\alpha_0,\tilde\alpha_1)
\end{equation*}
is carried to a configuration $U^0_k(\mathbf{U},\mathbf{J})$. This configuration lies in the space
\begin{equation*}
\tilde U''^c_k(Y_1,\tilde\alpha_0,\tilde\alpha_1)
\end{equation*}
of \cite[(1.64)--(1.69)]{B-CMP122a}, taken at the final value of the stage index. \cite[p.~375]{B-CMP122b} states this membership; it is used here as stated there.

The map $(\mathbf{U},\mathbf{J})\mapsto U^0_k(\mathbf{U},\mathbf{J})$ is a composite of the re-minimisations \cite[(1.61)--(1.63)]{B-CMP122b}. It is not a restriction, so restriction across levels as a change of units (Theorem~\ref{4.8:thm-restriction-units}) does not apply to it.

For an old term $E^{(j)}(X,\cdot)$ of \cite{B-CMP109}, the target of this configuration is $U^{\Sigma^1}_j(X)$. Membership in this target is not established in this section.

It is supplied by the orbit-form membership theorem for the composite of the re-minimisations \cite[(1.61)--(1.63)]{B-CMP122b}. That theorem is stated separately for each old term that receives the configuration, under the right-hand two-power cell of \eqref{4.8:eq-class-comparison-a} read on both rules. This reading on both rules is available because the inputs of that theorem are stated for every admissible set. These inputs are the tower-locality statement for an admissible determining set and \cite[Theorem~1]{B-CMP102b}. At the target, the theorem is combined with the producer propositions of this section:
\begin{itemize}
\item the producer proposition for Ba{\l}aban's representative at the minimiser of \cite[(3.37)]{B-CMP109} and Proposition~\ref{4.8:prop-small-pairs}, neither of which distinguishes the two rules;
\item Proposition~\ref{4.8:prop-far-case}, applied for each rule $\mathfrak{r}\in\mathfrak{R}$.
\end{itemize}
Membership would also follow from locality of clause (iv) in the tower at complex configurations. That locality is not proved: the tower-locality statement for an admissible determining set named above holds at real configurations only.

\begin{remark}[Lower bounds on $L$ and $M$ and amplitude ceilings]\label{4.8:rem-lower-bounds}
We collect the conditions on $L$, $M$ and $B^{(109)}$, and the ceilings on the radii, under
which the statements of this section hold. Every lower bound below is one-sided. The constants
on which it depends are listed after it in brackets, and each of them is fixed before it.
\begin{enumerate}
\item[(1)] $L\ge 2(1+\beta)$ [$\beta$]. Here only one power of $L^{-1}$ is available, and this is
the case that fixes the bound. This is the condition
of the first inclusion sentence of Corollary~\ref{4.8:cor-inclusion-sentences}. It enters, in
running form, the treatment of the large-field part of the coupling increment.
\item[(2)] The corresponding floor of Proposition~\ref{4.8:prop-near-case}
[$\beta$, $c_{\mathrm{Ad}}$], where $c_{\mathrm{Ad}}$ is an absolute constant:
\begin{equation*}
L^2\ge 2(1+c_{\mathrm{Ad}}\beta)(1+11\beta).
\end{equation*}
Here $1+11\beta=L^2\theta_J$, where $\theta_J$ is the current-room ratio of the corresponding hypothesis; at constant radii $\theta_J=\theta_\times+\beta L^{-2}$, where
$\theta_\times$ is the curvature analogue of $\theta_J$ used in the treatment of the large-field
part of the coupling increment. This floor also covers the factor $1+8\beta$ of the curvature
member \cite[(3.52)]{B-CMP109}, for every $|B'|<\alpha_3$.
\item[(3)] $c_{12}L\ge 2$ [$c_{12}$] and $L\ge 4d$ [$d$]. These are the floors of the increment
lemma on which Proposition~\ref{4.8:prop-far-case} rests.
\item[(4)] The following floor [$B_3$, $O_{(3.37)}$, $c_{12}$, $L$], where $O_{(3.37)}$ is the
constant written $O(1)$ in \cite[(3.37)]{B-CMP109}:
\begin{equation*}
B^{(109)}\ge \frac{2\bigl(B_3^2\,O_{(3.37)}+1\bigr)}{c_{12}L}.
\end{equation*}
This condition is not used by any statement of this section: with the factorisation $U:=1$ of
Proposition~\ref{4.8:prop-near-case}, the condition \cite[(1.12)]{B-CMP109} is void. It is
recorded for completeness among the lower bounds on $B^{(109)}$, the constant of
\cite[(1.12)]{B-CMP109}.
\item[(5)] $\delta_0M\ge 32L$ [$\delta_0$, $L$] and $M\ge 2LR_1M_1$ [$L$, $R_1$, $M_1$]. These are
the floors \eqref{4.8:eq-far-case-a}.
\end{enumerate}
The bounds \eqref{4.8:eq-top-layer-bound-a} and \eqref{4.8:eq-layerwise-bound-a} depend only on
floors of the layered large-field bounds. Both use the condition on $C_1'$ and the condition on
$C_0$ in \eqref{4.8:eq-top-layer-bound-a}:
\begin{equation*}
C_1'>16C_H(1+\beta_0)\|C\|C_1^{\mathrm{fl}}A_1,\qquad
C_0>2C_{IV}C_H(1+\beta_0)\|C\|C_1^{\mathrm{fl}}A_1 .
\end{equation*}
The condition on $C_1'$ in \eqref{4.8:eq-top-layer-bound-a} is a consequence of the floor
\begin{equation}\label{4.8:eq-lower-bounds-2}
C_0,\ C_1'\ \ge\ \frac{8c_\flat(1+\beta_0)\|C\|C_1^{\mathrm{fl}}A_1C_H^2}{(1-\theta)\beta},
\end{equation}
as is shown below. The layerwise bound \eqref{4.8:eq-layerwise-bound-a} uses in addition the
floor
\begin{equation}\label{4.8:eq-lower-bounds-1}
2\sqrt{2}\,e^{-\delta_H\mathsf{w}_\star}\le 1.
\end{equation}
The bound used in the straddling case of Lemma~\ref{4.8:lem-layerwise-bound}, where $X$ meets
$\Omega_j\setminus\Omega_{j+1}$, uses in addition the conditions
\eqref{4.8:eq-layerwise-bound-b}:
\begin{equation*}
C_0>\frac{2C_{IV}C_H(1+\beta_0)\|C\|C_1^{\mathrm{fl}}A_1}{(1-\theta)\beta},\qquad
C_1'>\frac{16C_H(1+\beta_0)\|C\|C_1^{\mathrm{fl}}A_1}{(1-\theta)\beta}.
\end{equation*}
The first of these is a new condition beside \eqref{4.8:eq-lower-bounds-2} unless
$C_{IV}\le 4c_\flat C_H$; the second is a consequence of \eqref{4.8:eq-lower-bounds-2}, as is
shown below.
These are conditions of the same kind as \eqref{4.8:eq-lower-bounds-2}. They are one-sided and
bear on the constants $C_0$, $C_1'$, which are taken to be sufficiently large positive numbers in
\cite[\S2]{B-CMP119}, where $C_1'$ is written $C_1$ (it is the constant of the radii
\cite[(2.28)]{B-CMP119}). They depend only on $C_{IV}$, $C_H$, $\beta_0$, $\|C\|$,
$C_1^{\mathrm{fl}}$, $A_1$, $\theta$ and $\beta$, all fixed before $C_0$; here $\theta$ and
$\beta$ are the letters of the layered large-field bounds in the factor $(1-\theta)\beta$.
None of the letters $L$, $M$, $B^{(109)}$, $k$, $j$, $\bar n$ occurs in them.

\emph{Amplitude ceilings.} The ceilings on the radii are the following.
\begin{itemize}
\item The restrictions of \cite[\S3]{B-CMP109} under which \cite[Lemma~4]{B-CMP109} holds, read
at the proof radius $2\alpha_3$, namely
$2B_3\alpha_3<\tfrac12\alpha_1$ and $4B_3\alpha_3\le\beta L^{-2}\alpha_0$. Among ``all the
restrictions'' of that lemma there is also $c\,\alpha_2\le\alpha_3$, where $c$ is the constant
in \cite[(3.49)]{B-CMP109}.
\item The ceilings of the increment lemma, of the form $\alpha\le c(d,L,M,B^{(109)},\dots)$.
\item The ceiling $2B_3\alpha_0'\le\alpha_0$ of Proposition~\ref{4.8:prop-small-pairs}.
\item The ceiling in the corresponding display:
\begin{equation*}
86d\Bigl(\sup_{\tilde{\square}^{3}}\bigl|\mathbf{H}_j\bigl(\square_0,\tau Q(L^{-1}\eta
\mathbf{H}_{k+1})\bigr)\bigr|+\sup_{\tilde{\square}^{3}}|\mathbf{A}_2|\Bigr)^2\le\beta L^{-2}
\alpha_0 .
\end{equation*}
\end{itemize}
All of them are chosen after $(L,B^{(109)},M)$. None of them bounds $L$, $M$ or $B^{(109)}$
from above.

\emph{What the floors do not depend on.} No lower bound on $L$, $M$ or $B^{(109)}$ depends on
any of the following:
\begin{itemize}
\item a level gap: $i-i'$, $j-p$, $k-j$, the gap $\mathfrak{n}=i-1-i'$, or the level gap
$\mathsf{g}$ of a rim sub-tower;
\item the layer index $\bar n$, or $k$, $K$, $N$;
\item the coupling;
\item the ratio $\alpha_1/\alpha_0$.
\end{itemize}
Two running ratios grow like $(1+\text{gap})^{1/2}$: the ratios $\alpha_{\cdot,i}/\alpha_{\cdot,i'}$
of Lemma~\ref{4.4:lem-radii-ratio}, and the ratios $w_k/\alpha_{\cdot,j}$ in
\eqref{4.8:eq-top-layer-bound-a} and \eqref{4.8:eq-layerwise-bound-a}. Each is compensated by an
explicit factor $L^{-\text{gap}}$, and, under \eqref{4.8:eq-lower-bounds-1}, by $e^{-D_{\bar n}}$.
The restriction of \cite[\S3]{B-CMP109} ordering $\alpha_0$ against $\alpha_1$ is used as
printed. At running radii it becomes a lower bound on the constant of the radius schedule in the
treatment of the large-field part of the coupling increment; this constant may be taken as large
as needed, and the bound is imposed after $(M,C_0)$ in that treatment and nowhere else.
Apart from this printed restriction, no condition of this section depends on
$\alpha_1/\alpha_0$. No condition of this section
is a cap in the sense of Definition~\ref{2.3:def-cap}.
\end{remark}

\begin{proof}
Items (1)--(5) are lower bounds of the form ``$\ge$''. The letters each one contains are those
listed after it.

In item (2), $\beta\ge0$ gives $1+8\beta\le1+11\beta$. Hence the corresponding floor implies the same inequality with $1+11\beta$ replaced by $1+8\beta$.

Next we compare the conditions on $C_0$ and $C_1'$. We use three properties of the constants
of the layered large-field bounds: $C_H=C_H(d,L)\ge1$, $c_\flat\ge4$, and
$0<(1-\theta)\beta\le\tfrac14$. From these,
\begin{equation*}
\frac{8c_\flat C_H^2}{(1-\theta)\beta}\ \ge\ \frac{32C_H^2}{(1-\theta)\beta}
\ >\ \frac{16C_H}{(1-\theta)\beta}\ \ge\ 64\,C_H\ \ge\ 16\,C_H .
\end{equation*}
Multiply by the positive number $(1+\beta_0)\|C\|C_1^{\mathrm{fl}}A_1$. Then
\eqref{4.8:eq-lower-bounds-2} for $C_1'$ gives both the condition on $C_1'$ in
\eqref{4.8:eq-top-layer-bound-a} and the second
condition of \eqref{4.8:eq-layerwise-bound-b}, each with strict inequality. Since
$(1-\theta)\beta\le1$, the first condition of \eqref{4.8:eq-layerwise-bound-b} implies the
condition on $C_0$ in \eqref{4.8:eq-top-layer-bound-a}.

We now show that the running ratios force no lower bound to depend on a gap. As shown in the
proof of
Lemma~\ref{4.8:lem-layerwise-bound}, $(1+b)^{1/2}L^{-b}\le1$ for every integer $b\ge0$, and,
under \eqref{4.8:eq-lower-bounds-1}, $(1+a)^{1/2}e^{-D_{\bar n}}\le2^{-a}$ for the integer
$a=k-\bar n\ge0$.
Every factor $(1+\text{gap})^{1/2}$ is thus bounded by $1$ once it is multiplied by the
explicit factor $L^{-\text{gap}}$, or by $e^{-D_{\bar n}}$. Therefore no lower bound needs to
contain a gap, $\bar n$, $k$ or $K$.

Finally, each amplitude ceiling is an upper bound on a radius by a quantity depending on
constants already fixed, among them $L$, $M$ and $B^{(109)}$. Choosing the radii after these
constants therefore imposes no upper bound on $L$, $M$ or $B^{(109)}$.
\end{proof}

\begin{remark}[Remarks on the literal inclusion sentences]\label{4.8:rem-literal-inclusions}
We collect five remarks on the inclusion sentences of \cite{B-CMP116} and \cite{B-CMP122b}
and on the route of \cite{B-CMP109} to them, read on Ba{\l}aban's regularity spaces
$U^c_j(X,\alpha_0,\alpha_1,\gamma_0)$ of Notation~\ref{4.8:not-regularity-spaces} as well
as on the one-power class $U^{\Sigma^1}_j$. None of them is used in the proofs of
Lemma~\ref{4.8:lem-class-properties}, Corollary~\ref{4.8:cor-inclusion-sentences} or
Proposition~\ref{4.8:prop-class-reproduction}. As in condition (iv) of Ba{\l}aban's spaces,
$\tilde X^{-2,j}$ denotes the domain obtained from $X$ by removing the two layers of cubes of
$\pi_j$ closest to the boundary of $X$ (we call these two layers the collar of $X$ in $\pi_j$),
and $\tilde Y^{-2,k}$ is obtained in the same way from
$Y\in\mathbf{D}_k$ with the cubes of $\pi_k$.
\begin{enumerate}
\item[(a)] Let $Y\in\mathbf{D}_k$, $X\in\mathbf{D}_{k+1}$, $Y\subset X$. Read the sentence of
\cite[p.~15]{B-CMP116} as the set inclusion
\begin{equation*}
U^c_{k+1}(X,\alpha_0,\alpha_1)|_{Y}\subset U^c_{k+1}(Y,(1+\beta)\alpha_0,(1+\beta)\alpha_1,\alpha_0),
\end{equation*}
with the collars of condition (iv) cut in the cubes of each domain's own family, which is
the practice of \cite[(3.40)]{B-CMP109}. This set inclusion does not hold for any
$\beta<(L^2-1)/(2L^2+1)$, in particular for any $\beta\le\tfrac14$, the values of $\beta$
considered at \cite[p.~198]{B-CMP122a}; the enlargement factor $1+\beta$ at this threshold
equals $3L^2/(2L^2+1)$ and tends to $\tfrac32$ as $L\to\infty$. Read with both collars cut in
the cubes of $\pi_{k+1}$, the family of the index $k+1$, the sentence is a statement about the
locality of the tower construction of condition (iv); neither a proof nor a counterexample to
it in that reading is given in this chapter. Nor does anything in this chapter contradict the
enlargement $1+2\beta$ of the first radius in \cite[(3.40)]{B-CMP109}, equal to $\tfrac32$ at
$\beta=\tfrac14$, which that display consumes. On the
one-power class the sentence is
the inclusion
\eqref{4.8:eq-inclusion-sentences-c}, which holds without any condition. For the sentence of
\cite[p.~9]{B-CMP116}, read literally on Ba{\l}aban's class for $X\in\mathbf{D}_j$, $j\le k$,
$Y\in\mathbf{D}_k$, $X\subset Y$ (we call the tower of such an $X$, at a level $j\le k$ below
the level $k+1$ of the class on $Y$, a rim sub-tower), on Ba{\l}aban's class it is not proved
here; on the one-power
class it holds, by \eqref{4.8:eq-inclusion-sentences-a}, whenever $L\ge2(1+\beta)$.
\item[(b)] Let $j\le k$ and let $\mathsf{g}=k+1-j\ge1$ be the level gap of a rim sub-tower
(so that $\mathsf{g}=\mathfrak{n}+1$, where $\mathfrak{n}\ge0$ is the gap of the running form
in Corollary~\ref{4.8:cor-inclusion-sentences}).
The only route in \cite{B-CMP109} to the literal sentence of \cite[p.~9]{B-CMP116} at such a
sub-tower is \cite[(1.17)]{B-CMP109}, read through the proof of that display, in which the
curvature of the data carries two powers of the scale ratio. It compares
$2B_3(1+\beta)L^{-2\mathsf{g}}\alpha_0$ with $\tfrac12\alpha_0$, where
$B_3=B_3(d,L)$ is the constant of \cite[Proposition~9]{B-CMP102b}, and it closes exactly in
the case
\begin{equation}\label{4.8:eq-literal-inclusions-1}
4B_3(1+\beta)\le L^{2\mathsf{g}},
\end{equation}
and not otherwise. In a reading of the statement of \cite[(1.17)]{B-CMP109} in which the bound
involving $\alpha_0$ carries one power of the scale ratio, the case would read
$4B_3(1+\beta)\le L^{\mathsf{g}}$. Since
$B_3$ depends on $L$, at fixed $\mathsf{g}$ the case \eqref{4.8:eq-literal-inclusions-1}
relates $L$ to the growth of $B_3(d,L)$ in $L$. It is a case split describing the route of
\cite{B-CMP109}, not a hypothesis, and it is never formed
on the one-power class, where, for $L\ge2(1+\beta)$, \eqref{4.8:eq-inclusion-sentences-a}
holds for every $\mathsf{g}\ge1$.
\item[(c)] Let $X\in\mathbf{D}_j$, $Y\in\mathbf{D}_k$. Suppose
\begin{equation}\label{4.8:eq-literal-inclusions-2}
\delta_H M/L\ge\log\bigl(2C_TB_3(1+\beta)/\beta\bigr),\qquad M\ge 8R_1M_1L .
\end{equation}
Then at $G$-valued representatives and at points $x\in\tilde X^{-2,j}\cap\tilde Y^{-2,k}$,
the quantities entering condition (iv) for the tower built on $X$ as $\Omega_0$ and for the
tower built on $Y$ as $\Omega_0$,
written $(\mathrm{iv})_X$ and $(\mathrm{iv})_Y$, satisfy
\begin{equation}\label{4.8:eq-literal-inclusions-3}
\bigl|(\mathrm{iv})_X-(\mathrm{iv})_Y\bigr|(x)\le C_TB_3\,a\,\xi^2e^{-\delta_H D(x)},
\end{equation}
where $C_T$ is the constant of the lemma of this book that bounds the difference of the
condition-(iv) quantities of two towers at $G$-valued points, $a$ is the radius at which
condition (iv) is read, $\xi$ is the lattice spacing at which that lemma reads condition (iv),
and $D(x)$ is the distance, as a function of $x$, that governs the
exponential decay in that lemma.
\item[(d)] The one-power class is monotone in every radius and under restriction. The
inductive construction of Proposition~\ref{4.8:prop-class-reproduction} uses the class only
through memberships in $U^{\Sigma^1}$ and analyticity on $U^{\Sigma^1}$; it never argues from
non-membership in $U^{\mathrm{lit}}_j$.
\item[(e)] \cite[p.~375]{B-CMP122b} states a third sentence of the same kind: a regularity
space of the stage construction of that paper (written there
$\tilde U''^c_k(Y_1,\tilde\alpha_0,\tilde\alpha_1)$) lies in the domains of the old terms, the
crucial domain, in the words of that page, being a difference of two domains of that
construction (written there $\Omega''_{k_0+1}\setminus\Omega_{k_0+1}$). It is stated in
\cite[p.~375]{B-CMP122b} without proof. It concerns the stage classes of that paper and
belongs to their analysis, not to that of the classes of \cite{B-CMP109}; it is not a
consequence of
Corollary~\ref{4.8:cor-inclusion-sentences}, which concerns the classes of
\cite{B-CMP109}.
\end{enumerate}
\end{remark}

\begin{proof}
(a) The assertions about the own-cube reading of the sentence of \cite[p.~15]{B-CMP116}
are the content of the theorem of this book asserting that the set inclusion displayed in
(a), with the collars cut in each domain's own cubes, fails for every
$\beta<(L^2-1)/(2L^2+1)$; this range contains every
$\beta\le\tfrac14$, since $2L^2\ge18>5$ for every odd $L\ge3$ and
\begin{equation*}
\frac{L^2-1}{2L^2+1}>\frac14\iff 4L^2-4>2L^2+1\iff 2L^2>5 .
\end{equation*}
The threshold $(L^2-1)/(2L^2+1)$ increases to $\tfrac12$ as $L\to\infty$, and the
corresponding enlargement factor is
\begin{equation*}
1+\frac{L^2-1}{2L^2+1}=\frac{3L^2}{2L^2+1}=\frac32-\frac{3}{2(2L^2+1)}\longrightarrow\frac32
\qquad(L\to\infty).
\end{equation*}
The remaining assertions of (a) about that sentence on Ba{\l}aban's class state what that
theorem does not reach: that with both collars cut
in the cubes of $\pi_{k+1}$ the sentence
is a statement about the locality of the tower construction, for which neither a proof nor a
counterexample is given in this chapter, and that nothing in this chapter contradicts the
enlargement $1+2\beta$ consumed by \cite[(3.40)]{B-CMP109}, whose value $\tfrac32$ at
$\beta=\tfrac14$ exceeds $3L^2/(2L^2+1)$ for every $L$, by the last display.
On the one-power class the same sentence, with $U^c$ replaced by $U^{\Sigma^1}$, is the
inclusion \eqref{4.8:eq-inclusion-sentences-c} of Corollary~\ref{4.8:cor-inclusion-sentences}.
It is stated for $Y\in\mathbf{D}_k$, $X\in\mathbf{D}_{k+1}$,
$Y\subset X$, with no further condition. The sentence of \cite[p.~9]{B-CMP116} at the rim
sub-towers concerns $X\in\mathbf{D}_j$, $j\le k$, $Y\in\mathbf{D}_k$, $X\subset Y$. On
Ba{\l}aban's class it is not proved here. On the
one-power class it is the first part of Corollary~\ref{4.8:cor-inclusion-sentences}, namely
\eqref{4.8:eq-inclusion-sentences-a}, whose only condition is $L\ge2(1+\beta)$.

(b) In the proof of \cite[(1.17)]{B-CMP109} the layered averages of data of radius
$\alpha_0'$ satisfy the curvature bound with constant $2\alpha_0'$, and
\cite[Proposition~9]{B-CMP102b} converts it into a bound on the axial-gauge minimisers with a
further factor $B_3$, giving $2B_3\alpha_0'$. The data of the class at $(k+1,Y)$, of radius
$(1+\beta)\alpha_0$, have curvature below $(1+\beta)\alpha_0\xi_{k+1}^2$, and
$\xi_{k+1}^2=L^{-2\mathsf{g}}\xi_j^2$ because $\xi_m=L^{-m}$ and $k+1-j=\mathsf{g}$; thus the
curvature of the data carries two powers of the scale ratio, and in the units $\xi_j^2$ of
level $j$ the display yields at a rim sub-tower of gap $\mathsf{g}$ the
quantity $2B_3(1+\beta)L^{-2\mathsf{g}}\alpha_0$, to be compared with the radius
$\tfrac12\alpha_0$ of the class $U^c_j(X,\tfrac12\alpha_0,\tfrac12\alpha_1,\alpha_0)$ at
$(j,X)$.
Since $\alpha_0>0$, multiplying by $2L^{2\mathsf{g}}/\alpha_0$ gives
\begin{equation*}
2B_3(1+\beta)L^{-2\mathsf{g}}\alpha_0\le\tfrac12\alpha_0
\iff 4B_3(1+\beta)\le L^{2\mathsf{g}},
\end{equation*}
which is \eqref{4.8:eq-literal-inclusions-1}. If the bound involving $\alpha_0$ carries one
power, $L^{-\mathsf{g}}$ replaces $L^{-2\mathsf{g}}$, and the same multiplication, now by
$2L^{\mathsf{g}}/\alpha_0$, gives $4B_3(1+\beta)\le L^{\mathsf{g}}$. The condition
$L\ge2(1+\beta)$ of \eqref{4.8:eq-inclusion-sentences-a} does not involve $\mathsf{g}$, so
that inclusion holds for every $\mathsf{g}\ge1$.

(c) The bound \eqref{4.8:eq-literal-inclusions-3} is part (i) of the lemma bounding the
difference of the condition-(iv) quantities of two towers at $G$-valued points, whose
remaining hypothesis on the base configuration holds at real, that is $G$-valued, bases,
applied at a $G$-valued representative and at a point
$x\in\tilde X^{-2,j}\cap\tilde Y^{-2,k}$ under the two floors
\eqref{4.8:eq-literal-inclusions-2}.

(d) Monotonicity in the radii $\alpha_0,\alpha_1,\gamma_0$ is
Lemma~\ref{4.8:lem-class-properties}$(\mathrm{d}_1)$. Monotonicity under restriction to
$X'\subset X$, $X'\in\mathbf{D}_{j-1}$, is Lemma~\ref{4.8:lem-class-properties}$(\mathrm{d}_2)$.
The conclusions of Proposition~\ref{4.8:prop-class-reproduction} are memberships in the
one-power class, together with analyticity on it. Lemma~\ref{4.8:lem-class-properties}(a)
gives the inclusion $U^{\Sigma^1}_j(X;\cdot)\subset U^{\mathrm{lit}}_j(X;\cdot)$ in one
direction only, and no step proceeds from a configuration's not lying in
$U^{\mathrm{lit}}_j$.

(e) Corollary~\ref{4.8:cor-inclusion-sentences} is an assertion about the one-power class
built on the conditions (i)--(iv) of \cite{B-CMP109}. The regularity space of the stage
construction of \cite{B-CMP122b} is a
different class, and no part of that corollary bears on it.
\end{proof}

\begin{definition}[The lattice, levels and closed cubes]\label{4.9:not-lattice-conventions}
The following conventions are in force throughout this section.
\begin{enumerate}
\item The dimension is $d=4$. The gauge group is $G=SU(N)$ with $N\ge 2$. Its Lie algebra
$\mathfrak{g}$ is realised as the traceless Hermitian $N\times N$ matrices, so that a
$G$-valued variable is written $U=e^{i\xi A}$ with $A$ taking values in $\mathfrak{g}$,
where $\xi$ is the spacing of the fine lattice in the units of the level in force (see the
next item).
The symbol $|\cdot|$ denotes the matrix norm of Ba{\l}aban's papers, the norm in which the
bounds of \cite[(1.12)]{B-CMP109} are stated. It is unitarily
invariant and satisfies $\|\cdot\|_{\mathrm{op}}\le|\cdot|$, where
$\|\cdot\|_{\mathrm{op}}$ is the operator norm.
\item There is one fine lattice, and all lengths are measured in its steps. At level $j$ the
unit of length is $L^j$ steps, and we put $\xi_j:=L^{-j}$ for the spacing of the fine
lattice in level-$j$ units. The partition $\pi_j$ of the lattice into cubes of side $ML^j$
steps and the class $\mathbf{D}_j$ of localisation domains at level $j$ are those of
\cite{B-CMP109}. We write $S:=c_{12}LM$, where $c_{12}$ is the constant written $O(1)$ in
the cube size $O(1)LM$ of \cite[(1.12)]{B-CMP109}, so that at level $j$ the cubes of size
$O(1)LM$ occurring there have side $S$ in level-$j$ units, that is, side $SL^{j}$ steps.
\item A bond is an oriented pair $b=\langle b_-,b_+\rangle=\langle x,x+e_\mu\rangle$, where
$e_\mu$ is the step in direction $\mu$. For a $G$-valued gauge transformation $u$ we set
$U^{u}(b):=u(b_-)U(b)u(b_+)^{-1}$, as in \cite[(1.10)]{B-CMP109}.
\item A \emph{closed cube of side $s$} is a set $\prod_{\mu=1}^{d}[a_\mu,a_\mu+s]$ with
$a_\mu$ real. These are the continuous space cubes of \cite{B-CMP109}.
\item Let $D$ be a closed subset of the continuous space in which the fine lattice and the
closed cubes of the preceding item lie. Then $\operatorname{st}(D)$ is the set of sites lying in $D$.
Next, $\operatorname{bd}(D)$ is the set of bonds with both ends in $\operatorname{st}(D)$.
Further, $\operatorname{pp}(D)$ is the set of pairs $(b,b+e_\nu)$, $\nu=1,\dots,d$, with
$b$ and $b+e_\nu$ both in $\operatorname{bd}(D)$; here
$b+e_\nu:=\langle b_-+e_\nu,\,b_++e_\nu\rangle$ is the bond $b$ translated by one step in
direction $\nu$, so that $\operatorname{pp}(D)$ is the set of parallel adjacent pairs of
bonds of $D$. For $r\ge0$ measured in steps,
\begin{equation*}
D^{+r}:=\{x:\operatorname{dist}_\infty(x,D)\le r\},
\end{equation*}
where $x$ ranges over the same continuous space and $\operatorname{dist}_\infty$ is the
distance in the supremum norm, measured in steps. Since $\operatorname{dist}_\infty(\cdot,D)$
is continuous, $D^{+r}$ is again a closed set.
\item Let $\lambda=(x_0,x_1,\dots,x_\ell)$ be a lattice path, that is, a sequence of sites
in which consecutive sites are nearest neighbours. For a bond field $U$ and nearest
neighbours $x,y$, put $U(x,y):=U(b)$ if $b=\langle x,y\rangle$ is a bond, and
$U(x,y):=U(b)^{-1}$ if $b=\langle y,x\rangle$ is a bond. The holonomy of $U$ along
$\lambda$ is the ordered product
\begin{equation*}
T_\lambda(U):=U(x_0,x_1)\,U(x_1,x_2)\cdots U(x_{\ell-1},x_\ell).
\end{equation*}
It satisfies
\begin{equation}\label{4.9:eq-lattice-conventions-1}
T_\lambda(U^{u})=u(x_0)\,T_\lambda(U)\,u(x_\ell)^{-1}.
\end{equation}
\end{enumerate}
\emph{Attachment of bonds to sets.} In \cite{B-CMP109} a set is assigned the bonds which
intersect it. In \cite{B-CMP102b} it is assigned the bonds with at least one end-point in
it. In this section a set is assigned the bonds with both ends in it, as in the definition
of $\operatorname{bd}(D)$ above. Every proof in this section that uses these sets is a
restriction, an identity, or a citation on a set enlarged by a collar of at least one step.
Each of them holds verbatim under either of the other two conventions, with
$\operatorname{st}(D)$ read as the set of end-points of the bonds of $D$.
\end{definition}

\begin{proof}
We verify \eqref{4.9:eq-lattice-conventions-1}. Let $x,y$ be nearest neighbours. We first
show that
\begin{equation*}
U^{u}(x,y)=u(x)\,U(x,y)\,u(y)^{-1}
\end{equation*}
in both orientations.

If $b=\langle x,y\rangle$ is a bond, then $b_-=x$ and $b_+=y$. The definition of $U^{u}$
then gives $U^{u}(x,y)=U^{u}(b)=u(x)U(b)u(y)^{-1}$, which is the claim.

If $b=\langle y,x\rangle$ is a bond, then $b_-=y$ and $b_+=x$. In this case
\begin{equation*}
\begin{split}
U^{u}(x,y)&=U^{u}(b)^{-1}
=\bigl(u(y)U(b)u(x)^{-1}\bigr)^{-1}\\
&=u(x)\,U(b)^{-1}\,u(y)^{-1}
=u(x)\,U(x,y)\,u(y)^{-1}.
\end{split}
\end{equation*}

We now insert this identity into each factor of $T_\lambda(U^{u})$:
\begin{equation*}
\begin{split}
T_\lambda(U^{u})
&=\prod_{i=1}^{\ell}u(x_{i-1})\,U(x_{i-1},x_i)\,u(x_i)^{-1}\\
&=u(x_0)\,U(x_0,x_1)\,u(x_1)^{-1}\,u(x_1)\,U(x_1,x_2)\,u(x_2)^{-1}\cdots
u(x_{\ell-1})\,U(x_{\ell-1},x_\ell)\,u(x_\ell)^{-1},
\end{split}
\end{equation*}
where the product is ordered from left to right in $i$. For $1\le i\le\ell-1$ the adjacent
factors $u(x_i)^{-1}u(x_i)$ equal the identity matrix. What remains is
$u(x_0)\,T_\lambda(U)\,u(x_\ell)^{-1}$.
\end{proof}

\begin{theorem}[Ba{\l}aban's minimal-orbit theorem with cube gauge {\cite[Theorem~1, (7)--(9), p.~279]{B-CMP102b}}]\label{4.9:thm-minimal-orbit}
Fix the level $k$. The lattice, the levels $j$ with $0\le j\le k$ and the cubes are those
fixed in Definition~\ref{4.9:not-lattice-conventions}. Let $\eta$ be the fine spacing, let
$\Omega_0\supset\Omega_1\supset\dots$ be the nest of small-field regions, with $\Lambda_j$
and $B^{j}(\Lambda_j)$ as in \cite[(2), (6)]{B-CMP102b}, so that
\begin{equation*}
B^{j}(\Lambda_j)\cup B^{j+1}(\Lambda_{j+1})=\Omega_j\setminus\Omega_{j+2},
\end{equation*}
and let $R_1$, $M_1$ be the constants of \cite{B-CMP99a}. Let $M'$ denote a
multiple of $R_1M_1$.

\emph{The class of cubes.} A cube $\square$ belongs to the class if, for some
$j$ with $0\le j\le k$, it is contained in
$B^{j}(\Lambda_j)\cup B^{j+1}(\Lambda_{j+1})$ and is a union of big blocks of
the $L^{-j}$-lattice, in the terminology of \cite{B-CMP102b}. More precisely:
$\square$ has side $2M'L^{j}\eta$, and the
cube $\widetilde{\square}$ with the same centre as $\square$ and side
$(2M'+4R_1M_1)L^{j}\eta$ is contained in
$B^{j}(\Lambda_j)\cup B^{j+1}(\Lambda_{j+1})$ but is not contained in
$B^{j+1}(\Lambda_{j+1})$.

\emph{The theorem.} There exist positive constants $a_0$, $a_1$, $B_3$,
$B_4$ and $M(\varepsilon_1)$, with $B_3a_1\le a_0$, where $B_4$ depends on a further
parameter of \cite[Theorem~1]{B-CMP102b}, such that the
following holds. Let $V$ be a configuration satisfying the regularity condition
\cite[(7)]{B-CMP102b} on its plaquette variables with a parameter
$\varepsilon_1\le a_1$. Then there exists
a minimal orbit in the space
\begin{equation*}
\mathfrak{U}_k(\{\Omega_j\},B_3\varepsilon_1)\cap
\mathfrak{B}_k(\mathfrak{B}_k,V).
\end{equation*}
The minimal configurations $U$ have, among the regularity properties listed in
\cite[(9)]{B-CMP102b}, the following. Let $\square$ be a cube of the class
described above, of side $2M'L^{j}\eta$, with $M'\le M(\varepsilon_1)$. Then
there exists a gauge transformation $u$, defined on a neighbourhood of
$\square$, such that on $\square$ one has $U^{u^{-1}}=e^{i\eta A}$ with a
Lie-algebra-valued field $A$ satisfying
\begin{equation*}
|A|<B_3M'\varepsilon_1\,(L^{j}\eta)^{-1},\qquad
|\nabla^{\eta}A|<B_3M'\varepsilon_1\,(L^{j}\eta)^{-2}.
\end{equation*}
The constants $a_0$, $a_1$, $B_3$ depend on $d$ and $L$ only, and
\begin{equation*}
M(\varepsilon_1)=R_1M_1\,(a_1/\varepsilon_1).
\end{equation*}
\end{theorem}

This is \cite[Theorem~1, (7)--(9), p.~279]{B-CMP102b}. In it, $M'$ is the free
multiple of $R_1M_1$ that \cite{B-CMP102b} writes $M$, and $M_1$ is the block
constant of \cite{B-CMP99a}. The field $A$ is the potential of $U$ on $\square$
in the gauge $u$. The constant $B_3$ is the one through which
\cite[(1.2)]{B-CMP109} passes from the condition on $V$ to the regularity of
the minimal configurations.
Throughout this chapter the letters $a_0$ and $a_1$ denote only the constants
of \cite[Theorem~1]{B-CMP102b} in the theorem above.

\begin{proposition}[Ba{\l}aban's cube gauge for regular fields
{\cite[Proposition~6, (1.135)--(1.136)]{B-CMP99a}}]\label{4.9:prop-cube-gauge-regular}
We use the notation of \cite{B-CMP99a}. Throughout, $d$ is the dimension, $\eta$ is the lattice spacing
and $L^j\eta$ is the length at level $j$. The constants $R_1$, $M_1$, $B_1$ and $c_1$ are those of
\cite[Theorems~2 and~4]{B-CMP99a}. The integer $M^{\ast}$ is a multiple of $R_1M_1$; it is printed $M$
in \cite{B-CMP99a}. The constant $R$ is the constant of that name in \cite{B-CMP99a}, and
$\alpha_0>0$. Cubes are lattice cubes as in the conventions of
Notation~\ref{4.9:not-lattice-conventions}. The setting is the
following.
\begin{enumerate}
\item[(a)] Let $\{\Omega_j\}$ be the sequence of domains of \cite[(1.7)--(1.9)]{B-CMP99a}. Let $U_0$
belong to $\mathfrak{U}_k(\{\Omega_j\},\alpha_0)$, that is, for every $j$,
\begin{equation}\label{4.9:eq-cube-gauge-regular-1}
\begin{aligned}
|U_0(\partial p)-1| &< \alpha_0\eta^2(L^j\eta)^{-2} && \text{for } p\in\Omega_j,\\
|(D^{\eta *}_{U_0}\partial U_0)(b)| &< \alpha_0\eta^2(L^j\eta)^{-3} && \text{for } b\in\Omega_j .
\end{aligned}
\end{equation}
\item[(b)] The cube $\square$ satisfies $\square\subset\Omega_j$ and $\square\not\subset\Omega_{j+1}$.
It is a union of cubes of side $R_1M_1L^j\eta$, and its own side is $M^{\ast}L^j\eta$.
Without loss of generality $j=k$: the domains $\Omega_{j'}$, $j'>j$, may be dropped from the
assumptions and the scale changed. This is assumed from now on.

Let
\begin{equation*}
\square_0\supset\square_1\supset\dots\supset\square_{k-1}\supset\square_k\supset\square
\end{equation*}
be a sequence of cubes. As a notational convenience put $\square_{k+1}:=\square$. For
$0\le j\le k$, the distance between the
boundaries of $\square_j$ and $\square_{j+1}$ equals $R_1M_1L^j\eta$.

Cover $\square_0$ by a smallest family of cubes of side $R_1M_1$. The union of this family is a
cube, denoted $\widetilde{\square}$; the distance of its boundary to $\square$ equals $2R_1M_1$.

It is assumed that $RM^{\ast}$ is bigger than $R_1M_1$, for example that
\begin{equation}\label{4.9:eq-cube-gauge-regular-2}
L^{-1}RM^{\ast}>2dR_1M_1 ;
\end{equation}
under~\eqref{4.9:eq-cube-gauge-regular-2} one has $\widetilde{\square}\subset\Omega_{k-1}$.
\item[(c)] Consider the configuration $U_0$ on the cube $\widetilde{\square}$. The configuration
$U_0'$ on $\widetilde{\square}$ is the one which \cite[Section~F]{B-CMP99a} constructs from it.
The configuration $U_0''$ is defined to equal $U_0'$ on
$\widetilde{\square}$ and $1$ outside $\widetilde{\square}$.
\end{enumerate}
Assume
\begin{equation}\label{4.9:eq-cube-gauge-regular-3}
7dL^2M^{\ast}\alpha_0\le c_1 .
\end{equation}
Then there exists a gauge transformation $u$ defined on $\widetilde{\square}$, and a
$\mathfrak{g}$-valued field $A$ on the bonds of
$\widetilde{\square}$, with the following two properties. First,
\begin{equation}\label{4.9:eq-cube-gauge-regular-4}
U_0'^{\,u^{-1}}=U_1=e^{i\eta A}\qquad\text{on } \widetilde{\square}.
\end{equation}
Second, for every $j\in\{0,\dots,k\}$, each of the four quantities
\begin{equation}\label{4.9:eq-cube-gauge-regular-5}
L^j\eta\,|A|,\qquad (L^j\eta)^2\,|\nabla^{\eta}A|,\qquad
(L^j\eta)^3\,|\partial^{\eta *}\partial^{\eta}A|,\qquad (L^j\eta)^3\,|\Delta^{\eta}A|
\end{equation}
is smaller than $7dL^2B_1M^{\ast}\alpha_0$ on $\square_j$.

Here $\nabla^{\eta}A$ is the ordinary finite-difference gradient \cite[p.~98]{B-CMP99a}. The operators
$\partial^{\eta}$, $\partial^{\eta *}$ and $\Delta^{\eta}$ are the finite-difference operators of
spacing $\eta$ of \cite{B-CMP99a}.
\end{proposition}

This is \cite[Proposition~6, (1.135)--(1.136), pp.~98--99]{B-CMP99a}.

The proof there runs as follows. It shows that $U_0''$ belongs to
$\mathfrak{U}_k(\{\square_j\},L^3\alpha_0)$, intersected with the set of configurations in the axial
gauge of \cite[(1.132)]{B-CMP99a}.
It also shows that the level-$j$ averaged configuration $\bar U_0''^{\,j}$ of \cite{B-CMP99a}
satisfies $|\bar U_0''^{\,j}-1|<6dL^2M^{\ast}\alpha_0$ on $\square_j^{(j)}$, the set of points of the
lattice of level $j$ in $\square_j$ \cite[(1.133)]{B-CMP99a}.
Under~\eqref{4.9:eq-cube-gauge-regular-3}, the assumptions of \cite[Theorem~4]{B-CMP99a} are then
satisfied for the pair of configurations $1$, $U_0''$.

The current bound, the second line of~\eqref{4.9:eq-cube-gauge-regular-1}, is part of the hypothesis
$U_0\in\mathfrak{U}_k(\{\Omega_j\},\alpha_0)$.

\begin{definition}[The cube gauge condition and its readings]\label{4.9:def-cube-gauge-condition}
The lattice, the levels and the closed cubes are those of
Convention~\ref{4.9:not-lattice-conventions}, and sides of cubes are counted in lattice steps.
The gauge group is $G=SU(N)$, and $\mathfrak{g}$ is its Lie algebra, realised as the traceless
Hermitian matrices.
Let $c_{12}$ denote the constant written $O(1)$ in \cite[(1.12)]{B-CMP109}. Let $B_{12}$ denote the
constant $B$ of that display \cite[p.~262]{B-CMP109}. These are the constants written $\mathsf{o}$
and $B^{109}$ in Definitions~\ref{4.4:def-local-data-classes} and~\ref{4.4:def-radius-ceilings}. Put
\begin{equation*}
S:=c_{12}LM,\qquad \xi_j:=L^{-j}.
\end{equation*}
For a set $D$ we use the following notation:
\begin{itemize}
\item $\operatorname{st}(D)$ is the set of sites in $D$;
\item $\operatorname{bd}(D)$ is the set of bonds with both end points in $D$;
\item $\operatorname{pp}(D)$ is the set of pairs $(b,b+e_\nu)$ of parallel adjacent bonds of
$\operatorname{bd}(D)$, where $b+e_\nu$ is the translate of $b$ by one step in direction $\nu$, and every
direction $\nu$ is admitted, the direction of $b$ itself included.
\end{itemize}
For a bond $b$ from $b_-$ to $b_+$ and $u\colon\operatorname{st}(D)\to G$, the gauge transform is
$U^{u}(b)=u(b_-)U(b)u(b_+)^{-1}$.

\emph{The clause.} Let $j$ be a level, $D$ a closed set, $U$ a $G$-valued field on a set of bonds
containing $\operatorname{bd}(D)$, and $\rho>0$. The \emph{cube gauge clause} $\mathsf{g}_j(U;D;\rho)$
is the following condition. There are $u\colon\operatorname{st}(D)\to G$ and
$A\colon\operatorname{bd}(D)\to\mathfrak{g}$ such that
\begin{equation}\label{4.9:eq-cube-gauge-condition-a}
\begin{gathered}
U^{u}=e^{i\xi_jA}\ \text{on }\operatorname{bd}(D),\qquad
|A(b)|<\rho\ \text{for } b\in\operatorname{bd}(D),\\
\xi_j^{-1}\,|A(b+e_\nu)-A(b)|<\rho\ \text{for }(b,b+e_\nu)\in\operatorname{pp}(D).
\end{gathered}
\end{equation}
If $\operatorname{st}(D)=\emptyset$, the clause imposes no condition.

\emph{The readings.} Let $X$ be a closed set, $U$ a $G$-valued field on $\operatorname{bd}(X)$, and
$a>0$. Let $B$ denote the absolute constant of \cite[(2.38)]{B-CMP119}, which is the constant $B$ of
\cite[(1.69)]{B-CMP122a}. The four readings of the clause are the following conditions, in which
$\square$ ranges over closed cubes and $\square\subset X$ is required only where it is written:
\begin{equation}\label{4.9:eq-cube-gauge-condition-b}
\begin{aligned}
&(\ell 0)_j(X;a)\colon
&&\mathsf{g}_j(U;\square;SB_{12}a)\ \text{for every }\square\subset X\text{ of side }SL^j;\\
&(\ell 1)_j(X;a)\colon
&&\mathsf{g}_j(U;\square;SB_{12}a)\ \text{for every }\square\subset X\text{ of side}\le SL^j;\\
&(\ell 2)_j(X;a)\colon
&&\mathsf{g}_j(U;\square\cap X;SB_{12}a)\ \text{for every }\square\text{ of side }SL^j;\\
&(\ell 2)^{\flat}_j(X;a)\colon
&&\mathsf{g}_j(U;\square\cap X;CMBa)\ \text{for every }\square\text{ of side }CML^j,
\ C\in\{1,2\}.
\end{aligned}
\end{equation}

The reading $(\ell 0)$ is the condition \cite[(1.12)]{B-CMP109} at level $j$. The reading $(\ell 1)$ is
the widened reading of \cite[(1.12)]{B-CMP109}, clause (i) of the definition of the one-power
hereditary class. Written out, $(\ell 1)_j(X;a)$ asks for the following for every closed
cube $\square\subset X$ of side at most $SL^j$:
\begin{itemize}
\item there are $u\colon\operatorname{st}(\square)\to G$ and
$A\colon\operatorname{bd}(\square)\to\mathfrak{g}$ with $U^{u}=e^{i\xi_jA}$;
\item $|A(b)|<SB_{12}a$ on $\operatorname{bd}(\square)$;
\item $\xi_j^{-1}|A(b+e_\nu)-A(b)|<SB_{12}a$ for every pair of parallel adjacent bonds of $\square$
and every $\nu$, the direction of the bond included.
\end{itemize}
The constant $SB_{12}a$ is the same for every side at most $SL^j$.

We call a closed cube $\square\subset X$ \emph{thin in $X$ at level $j$} if it lies in no closed cube
contained in $X$ of side exactly $SL^j$. Every cube of a set $X$ with fewer than $\lceil c_{12}L\rceil$
consecutive cells of $\pi_j$ in some direction is thin in $X$ at level $j$.

The reading $(\ell 2)^{\flat}$ is the clause of \cite[(2.38)]{B-CMP119} for the top index, that is, the
clause $m=k$ of \cite[(1.69)]{B-CMP122a}, with the requirement that the cube lie in the small-field
region of that index dropped.

\emph{The nest reading.} Let $\{\Omega_n\}_{n\le k}$ be a nest of small-field regions, let $(a_n)$
be radii, and let $Y$ be a closed set, with $U$ a $G$-valued field on a set of bonds containing
$\operatorname{bd}(Y)$. We say that $U$ satisfies the \emph{nest reading of \cite[(2.38)]{B-CMP119}}
on $Y$ with radii $(a_n)$ if, for every index $n\le k$ of the nest,
\begin{equation}\label{4.9:eq-cube-gauge-condition-c}
\begin{gathered}
\mathsf{g}_n(U;\square\cap Y;BCMa_n)\ \text{for every closed cube }\square\text{ of side }CML^n,
\ C\in\{1,2\},\\
\text{such that }\square\subset\Omega_n\setminus\Omega_{n+2}\text{ and }
\square\cap\Omega^{c}_{n+1}\neq\emptyset\ \text{if }n<k,
\quad \square\subset\Omega_k\ \text{if }n=k.
\end{gathered}
\end{equation}
The condition is required for each such cube and for both values $C\in\{1,2\}$.
\end{definition}

\begin{remark}
On a cube thin in $X$ at level $j$ the reading $(\ell 0)_j(X;a)$ imposes no condition, while
$(\ell 1)_j(X;a)$ does.
The bound on $|A|$ in $(\ell 1)$ follows there from the curvature bound of \cite[(1.11)]{B-CMP109},
by the comb-gauge corollary on the site box of $\square$. On such cubes the bound on the differences
of $A$ is not among the four conditions of \cite[p.~262]{B-CMP109}; it is the subject of the theorem
on the gradient bound on thin cubes.
\end{remark}

\begin{remark}[What one-step reproduction of the condition means]\label{4.9:rem-reproduction-meaning}
Write $E^{(j)}(X,g_{j-1},\cdot)$ for the localised term of level $j$ on the localisation domain
$X$. In
the induction the one-power hereditary class $U^{\Sigma^1}_j(X;\alpha_0,\alpha_1)$ enters only as a
domain, through the hypothesis
\begin{equation}\label{4.9:eq-reproduction-meaning-1}
E^{(j)}(X,g_{j-1},\cdot)\ \text{is defined and analytic on}\ U^{\Sigma^1}_j(X;\alpha_0,\alpha_1).
\end{equation}
Let $(\ell 1)_j(X;a)$ be the reading of the cube gauge condition fixed in
Definition~\ref{4.9:def-cube-gauge-condition} and displayed in
\eqref{4.9:eq-cube-gauge-condition-b}. The assertion that one renormalisation step reproduces
this reading on the one-power hereditary class means exactly the following three statements,
\begin{equation}\label{4.9:eq-reproduction-meaning-a}
\begin{gathered}
\text{(a) gauge clause of the source}\ \Longrightarrow\ (\ell 1)_j(X;a_0)\ \text{for the representative},\\
\text{(b) unit pair, small pairs, real backgrounds}\ \Longrightarrow\ (\ell 1)_j(X;\alpha_0'),\\
\text{(c) one and the same $B_{12}$ in (a) and (b)},
\end{gathered}
\end{equation}
which in full read as follows.
\begin{enumerate}
\item[(a)] \emph{Transport.} At every place in the construction where a membership in
$U^{\Sigma^1}_j(X;a_0,\dots)$ is asserted and is derived from another membership (a
restriction of a source, or an argument substituted into an old term), the representative
$\mathbf{U}=U'U$ exhibited there has a $G$-valued factor $U$ satisfying $(\ell 1)_j(X;a_0)$,
provided the gauge clause of the source's space holds for the source's representative at the
source's level, domain and radius.
\item[(b)] \emph{Seeds.} Every pair whose membership is not derived from another membership,
namely the unit pair, the small pairs and the real backgrounds, satisfies
$(\ell 1)_j(X;\alpha_0')$, where $\alpha_0'$ is the constant of its membership as printed.
\item[(c)] \emph{Constant.} One number $B_{12}$, chosen before $M$ and independent of $k$, $j$,
$X$, the nest of small-field regions, the volume and the coupling, serves both in (a) and
in (b).
\end{enumerate}

Since the class enters only through \eqref{4.9:eq-reproduction-meaning-1}, what is asked of a
reading of its definition is what is asked of any domain of analyticity: a reading of a
definition imposes nothing further only if every argument ever substituted into the functions
defined on that domain satisfies it. This is the requirement Ba{\l}aban formulates when the
functions of his construction are substituted for the variables $\mathbf{U},\mathbf{J}$: one has
to make sure that they have the required properties \cite[discussion following (1.15)]{B-CMP109}.

Every argument substituted into a term whose domain is the one-power hereditary class carries,
at the place where its membership is asserted, either a membership derived from another one or a
membership that is not so derived. In the first case the only condition available at that
place is the condition satisfied by the source's representative, at the source's level, domain
and radius. Reproduction therefore means that this condition is carried to the exhibited
representative, which is statement (a). In the second case there is no source to carry the
condition from. The arguments are then the unit pair, the small pairs and the real backgrounds,
and the condition has to hold for them directly, at the constant $\alpha_0'$ of their printed
membership, which is statement (b).

By Definition~\ref{4.9:def-cube-gauge-condition}, the reading $(\ell 1)_j(X;a)$ requires
$\mathsf{g}_j(U;\square;SB_{12}a)$ on every closed cube $\square\subset X$ of side at most
$SL^j$. The constant $B_{12}$ is thus part of the definition of the domain. Hence (a) and (b)
concern one and the same domain only if they hold with one and the same number $B_{12}$.
Statement (c) asks that this number be fixed before $M$ and be independent of $k$, $j$, $X$,
the nest, the volume and the coupling.

The places covered by (a) are of two kinds. In the first kind the source membership is a
membership in a class of \cite{B-CMP109} or \cite{B-CMP116}; one of them is the product
string in the proof of Corollary~\ref{4.4:cor-localised-independence} (independence of the
localised terms from the small-field region). In the second kind the source is a
space of \cite{B-CMP119} or \cite{B-CMP122a}, and the condition available
for the source's representative is \cite[(2.38)]{B-CMP119}, respectively
\cite[(1.65), (1.67), (1.69)]{B-CMP122a}.
\end{remark}

\begin{lemma}[Restriction and monotonicity of the condition]\label{4.9:lem-restriction-monotone}
Let $j$ be a level, let $D'\subset D$ be closed sets, let $U$ be $G$-valued on a set of bonds
containing $\operatorname{bd}(D)$, and let $0<\rho\le\rho'$. Then
$\mathsf{g}_j(U;D;\rho)$ implies $\mathsf{g}_j(U;D';\rho')$.

In particular, for the readings of Definition~\ref{4.9:def-cube-gauge-condition}
(see~\eqref{4.9:eq-cube-gauge-condition-b}), with $X$ closed, $U$ on $\operatorname{bd}(X)$
and $a>0$:
\begin{enumerate}
\item $(\ell 2)_j(X;a)$ implies $(\ell 1)_j(X;a)$, and $(\ell 1)_j(X;a)$ implies
$(\ell 0)_j(X;a)$;
\item each of $(\ell 1)$, $(\ell 2)$, $(\ell 2)^{\flat}$ that holds on $X$ holds on every
closed $X'\subset X$;
\item each of $(\ell 1)$, $(\ell 2)$, $(\ell 2)^{\flat}$ is monotone in $a$, $B$ and $B_{12}$.
\end{enumerate}
\end{lemma}

\begin{proof}
Since $D'\subset D$, a site of $D'$ is a site of $D$, a bond with both ends in $D'$ has both
ends in $D$, and a parallel adjacent pair of bonds of $D'$ is such a pair of $D$; thus
$\operatorname{st}(D')\subset\operatorname{st}(D)$,
$\operatorname{bd}(D')\subset\operatorname{bd}(D)$ and
$\operatorname{pp}(D')\subset\operatorname{pp}(D)$. In particular $U$ is defined on
$\operatorname{bd}(D')$. If $\operatorname{st}(D')=\emptyset$ the clause
$\mathsf{g}_j(U;D';\rho')$ is void. Otherwise let $u$ and $A$ be as in
$\mathsf{g}_j(U;D;\rho)$ (see~\eqref{4.9:eq-cube-gauge-condition-a}), and take their
restrictions to $\operatorname{st}(D')$ and to $\operatorname{bd}(D')$. For a bond of
$\operatorname{bd}(D')$ both ends lie in $\operatorname{st}(D')$, so $U^{u}=e^{i\xi_jA}$ holds
there with the restricted data. On $\operatorname{bd}(D')$ we have $|A(b)|<\rho\le\rho'$, and on
$\operatorname{pp}(D')$ we have
$\xi_j^{-1}|A(b+e_\nu)-A(b)|<\rho\le\rho'$. This is $\mathsf{g}_j(U;D';\rho')$.

(1) Let $Q\subset X$ be a closed cube of side at most $SL^j$. It lies in a closed cube
$\square$ of space of side $SL^j$, and then $Q\subset\square\cap X$. The clause that
$(\ell 2)_j(X;a)$ gives on $\square\cap X$ passes to $Q$ by the first assertion, applied with
$D=\square\cap X$ and $D'=Q$, the radius given by $(\ell 2)_j(X;a)$ not exceeding the radius
$SB_{12}a$ that $(\ell 1)_j(X;a)$ requires. Hence $(\ell 2)_j(X;a)$ implies
$(\ell 1)_j(X;a)$. The cubes of side equal to $SL^j$ are among those of side at most $SL^j$,
with the same radius $SB_{12}a$; so $(\ell 1)_j(X;a)$ implies $(\ell 0)_j(X;a)$.

(2) Let $X'\subset X$ be closed; then $\operatorname{bd}(X')\subset\operatorname{bd}(X)$ and
$U$ restricts to $\operatorname{bd}(X')$. For $(\ell 1)$: a closed cube contained in $X'$ is a
closed cube contained in $X$, so the required clause holds on it. For $(\ell 2)$ and
$(\ell 2)^{\flat}$: for every cube $\square$ of space, $\square\cap X'\subset\square\cap X$, and
the clause on $\square\cap X$ passes to $\square\cap X'$ by the first assertion with
$\rho=\rho'$.

(3) The constants $a$, $B$, $B_{12}$ enter each reading only through the radius, and this
radius does not decrease when any of them is increased. The first assertion with $D'=D$ and
$\rho\le\rho'$ then shows that each reading persists under an increase of $a$, $B$ or $B_{12}$.
\end{proof}

\begin{lemma}[Change of level lowers the radius]\label{4.9:lem-level-change}
Let $i'\le i$ be levels and put
\begin{equation*}
\mu:=L^{-(i-i')}.
\end{equation*}
Let $D$ be a closed set, let $U$ be a $G$-valued field on a set of bonds containing
$\operatorname{bd}(D)$, and let $\rho>0$. If $\mathsf{g}_i(U;D;\rho)$ holds, then
$\mathsf{g}_{i'}(U;D;\mu\rho)$ holds, with the same gauge transformation
$u\colon\operatorname{st}(D)\to G$.
\end{lemma}

\begin{proof}
Throughout, $\mathrm{i}$ denotes the imaginary unit, and the clauses are those of
Definition~\ref{4.9:def-cube-gauge-condition}, display~\eqref{4.9:eq-cube-gauge-condition-a}.
If $\operatorname{st}(D)=\emptyset$, both clauses are void and there is nothing to prove.
Otherwise, by $\mathsf{g}_i(U;D;\rho)$ there are $u\colon\operatorname{st}(D)\to G$ and
$A\colon\operatorname{bd}(D)\to\mathfrak{g}$ with $U^{u}=e^{\mathrm{i}\xi_iA}$ on
$\operatorname{bd}(D)$, with $|A(b)|<\rho$ for every $b\in\operatorname{bd}(D)$, and with
$\xi_i^{-1}|A(b+e_\nu)-A(b)|<\rho$ for every pair $(b,b+e_\nu)\in\operatorname{pp}(D)$.

Since $\xi_j=L^{-j}$, we have $\xi_i=L^{-(i-i')}L^{-i'}=\mu\,\xi_{i'}$. Define
$A^{\flat}:=\mu A$, again a $\mathfrak{g}$-valued function on $\operatorname{bd}(D)$. Then
\begin{equation*}
e^{\mathrm{i}\xi_iA}=e^{\mathrm{i}\xi_{i'}\mu A}=e^{\mathrm{i}\xi_{i'}A^{\flat}}
\quad\text{on }\operatorname{bd}(D),
\end{equation*}
so $U^{u}=e^{\mathrm{i}\xi_{i'}A^{\flat}}$ on $\operatorname{bd}(D)$ with the same $u$.
Next, $|A^{\flat}(b)|=\mu|A(b)|<\mu\rho$ for every $b\in\operatorname{bd}(D)$. Finally,
$\xi_{i'}^{-1}=\mu\,\xi_i^{-1}$, and since $L>1$ and $i'\le i$ we have $\mu\le1$; hence for
every pair $(b,b+e_\nu)\in\operatorname{pp}(D)$
\begin{equation*}
\begin{split}
\xi_{i'}^{-1}\bigl|A^{\flat}(b+e_\nu)-A^{\flat}(b)\bigr|
&=\mu\,\xi_i^{-1}\cdot\mu\,\bigl|A(b+e_\nu)-A(b)\bigr|\\
&<\mu^{2}\rho\le\mu\rho .
\end{split}
\end{equation*}
The pair $(u,A^{\flat})$ therefore satisfies the three requirements of
$\mathsf{g}_{i'}(U;D;\mu\rho)$.
\end{proof}

\begin{lemma}[Invariance under group-valued gauge maps]\label{4.9:lem-gauge-invariance}
Let $j$, $D$, $U$ and $\rho>0$ be as in Definition~\ref{4.9:def-cube-gauge-condition}, and let
$v\colon\operatorname{st}(D)\to G$. Then
\begin{equation*}
\mathsf{g}_j(U;D;\rho)\quad\Longleftrightarrow\quad\mathsf{g}_j(U^{v};D;\rho).
\end{equation*}
\end{lemma}

\begin{proof}
The clause $\mathsf{g}_j(\,\cdot\,;D;\rho)$ of \eqref{4.9:eq-cube-gauge-condition-a} reads its
configuration only on $\operatorname{bd}(D)$. Both endpoints $b_-,b_+$ of a bond
$b\in\operatorname{bd}(D)$ lie in $\operatorname{st}(D)$, so $U^{v}$ is defined on $\operatorname{bd}(D)$.
If $\operatorname{st}(D)=\emptyset$, both sides are void.

For $u,u'\colon\operatorname{st}(D)\to G$ and $b\in\operatorname{bd}(D)$ we have
\begin{equation*}
(U^{v})^{u'}(b)=u'(b_-)v(b_-)U(b)v(b_+)^{-1}u'(b_+)^{-1}=U^{u'v}(b),
\end{equation*}
where $u'v$ is the pointwise product. Taking $u'=uv^{-1}$ gives
\begin{equation}\label{4.9:eq-gauge-invariance-1}
(U^{v})^{uv^{-1}}=U^{u}\qquad\text{on }\operatorname{bd}(D).
\end{equation}

Suppose $\mathsf{g}_j(U;D;\rho)$ holds, witnessed by $u\colon\operatorname{st}(D)\to G$ and
$A\colon\operatorname{bd}(D)\to\mathfrak{g}$. By \eqref{4.9:eq-gauge-invariance-1}, the map
$uv^{-1}\colon\operatorname{st}(D)\to G$ satisfies $(U^{v})^{uv^{-1}}=e^{i\xi_jA}$ on
$\operatorname{bd}(D)$. The bounds on $A$ in \eqref{4.9:eq-cube-gauge-condition-a} are unchanged,
since $A$ is unchanged. Hence the pair $(uv^{-1},A)$ witnesses $\mathsf{g}_j(U^{v};D;\rho)$.

For the converse, note that $U=(U^{v})^{v^{-1}}$ on $\operatorname{bd}(D)$, by the composition
identity displayed before \eqref{4.9:eq-gauge-invariance-1} with $u'=v^{-1}$, since $v^{-1}v=1$.
Applying
the implication just proved to the configuration $U^{v}$ and the map
$v^{-1}\colon\operatorname{st}(D)\to G$ shows that $\mathsf{g}_j(U^{v};D;\rho)$ implies
$\mathsf{g}_j(U;D;\rho)$.
\end{proof}

\begin{lemma}[The identity field satisfies the condition]\label{4.9:lem-identity-field}
Let $j$ be a level, let $D$ be a closed set, and let $\rho>0$. The configuration $U\equiv 1$,
equal to the identity of $G$ on every bond of a set of bonds containing
$\operatorname{bd}(D)$, satisfies the cube gauge condition $\mathsf{g}_j(1;D;\rho)$ of
\eqref{4.9:eq-cube-gauge-condition-a} (Definition~\ref{4.9:def-cube-gauge-condition}).
\end{lemma}

\begin{proof}
If $\operatorname{st}(D)=\emptyset$, the condition is void. Otherwise define
$u\colon\operatorname{st}(D)\to G$ by $u\equiv 1$ and $A\colon\operatorname{bd}(D)\to\mathfrak{g}$
by $A\equiv 0$; the zero matrix is traceless and Hermitian, so $A$ takes values in
$\mathfrak{g}$. For every bond $b\in\operatorname{bd}(D)$ with initial site $b_-$ and final
site $b_+$,
\begin{equation*}
U^{u}(b)=u(b_-)\,U(b)\,u(b_+)^{-1}=1\cdot 1\cdot 1=1=e^{i\xi_j\cdot 0}=e^{i\xi_jA(b)} .
\end{equation*}
Next, $|A(b)|=0<\rho$ for every $b\in\operatorname{bd}(D)$. Finally, for every pair
$(b,b+e_\nu)\in\operatorname{pp}(D)$ we have $A(b+e_\nu)-A(b)=0$, so that
$\xi_j^{-1}|A(b+e_\nu)-A(b)|=0<\rho$. These are the three requirements of
\eqref{4.9:eq-cube-gauge-condition-a}, and none of them uses $j$, $D$ or $\rho$ beyond
$\rho>0$; hence $\mathsf{g}_j(1;D;\rho)$ holds for every $j$, every closed $D$ and every
$\rho>0$.
\end{proof}

\begin{lemma}[Holonomy bound from the cube gauge]\label{4.9:lem-holonomy-bound}
Let $j$ be a level, $D$ a closed set, $\rho>0$, and let $U$ be a $G$-valued bond field on a set
of bonds containing $\operatorname{bd}(D)$. If $\mathsf{g}_j(U;D;\rho)$ holds in the sense of
\eqref{4.9:eq-cube-gauge-condition-a}, then every closed lattice path $\lambda$ of $\ell$ bonds,
all in $\operatorname{bd}(D)$, satisfies
\begin{equation}\label{4.9:eq-holonomy-bound-1}
\|T_{\lambda}(U)-1\|_{\mathrm{op}}\le \ell\,\xi_j\,\rho .
\end{equation}
\end{lemma}

\begin{proof}
By Definition~\ref{4.9:def-cube-gauge-condition}, there are $u\colon\operatorname{st}(D)\to G$ and
$A\colon\operatorname{bd}(D)\to\mathfrak{g}$ with $U^{u}=e^{i\xi_jA}$ on $\operatorname{bd}(D)$ and
$|A(b)|<\rho$ for $b\in\operatorname{bd}(D)$.

Fix $b\in\operatorname{bd}(D)$. The matrix $A(b)\in\mathfrak{g}$ is Hermitian, so
$A(b)=W\operatorname{diag}(\theta_1,\dots,\theta_N)W^{-1}$ with $W$ unitary and $\theta_n$ real,
and hence $e^{\pm i\xi_jA(b)}-1=W\operatorname{diag}(e^{\pm i\xi_j\theta_n}-1)W^{-1}$. Since
$|e^{i t}-1|=2|\sin(t/2)|\le|t|$ for real $t$, and since the operator norm is dominated by the
norm $|\cdot|$ in which the clause is written,
\begin{equation}\label{4.9:eq-holonomy-bound-2}
\|e^{\pm i\xi_jA(b)}-1\|_{\mathrm{op}}\le \xi_j\max_n|\theta_n|
=\xi_j\|A(b)\|_{\mathrm{op}}<\xi_j\rho .
\end{equation}

Let $\lambda$ pass through the sites $x_0,x_1,\dots,x_\ell=x_0$, the $i$-th step running along a
bond $b_i\in\operatorname{bd}(D)$ from $x_{i-1}$ to $x_i$. Then $T_{\lambda}(U)$ is the ordered
product of the factors $U(b_i)^{\epsilon_i}$, where $\epsilon_i=+1$ if $b_i$ is traversed in its
own direction and $\epsilon_i=-1$ otherwise. Every bond of $\operatorname{bd}(D)$ has both ends in
$D$, so every $x_i$ lies in $\operatorname{st}(D)$, and $u(x_i)$ is defined. From
$U^{u}(b)=u(b_-)U(b)u(b_+)^{-1}$, in either orientation
$U^{u}(b_i)^{\epsilon_i}=u(x_{i-1})U(b_i)^{\epsilon_i}u(x_i)^{-1}$. The product telescopes:
\begin{equation*}
T_{\lambda}(U^{u})=u(x_0)\,T_{\lambda}(U)\,u(x_\ell)^{-1}=u(x_0)\,T_{\lambda}(U)\,u(x_0)^{-1}.
\end{equation*}
On the other hand, $T_{\lambda}(U^{u})=V_1\cdots V_\ell$ with
$V_i=e^{i\epsilon_i\xi_jA(b_i)}$, a product of $\ell$ unitaries of the kind estimated in
\eqref{4.9:eq-holonomy-bound-2}. Since $u(x_0)$ is unitary,
$\|T_{\lambda}(U)-1\|_{\mathrm{op}}=\|V_1\cdots V_\ell-1\|_{\mathrm{op}}$.

Finally, we have the identity
\begin{equation*}
V_1\cdots V_\ell-1=\sum_{i=1}^{\ell}V_1\cdots V_{i-1}(V_i-1).
\end{equation*}
Each $V_1\cdots V_{i-1}$ is unitary, so it has operator norm one. Hence
\begin{equation*}
\|V_1\cdots V_\ell-1\|_{\mathrm{op}}\le\sum_{i=1}^{\ell}\|V_i-1\|_{\mathrm{op}}
\le\ell\,\xi_j\,\rho,
\end{equation*}
by \eqref{4.9:eq-holonomy-bound-2}. This is \eqref{4.9:eq-holonomy-bound-1}.
\end{proof}

\begin{corollary}[Passage to smaller cubes by gluing]\label{4.9:cor-gluing-smaller-cubes}
Let $j$ be a level and $a>0$. Let $S:=c_{12}LM$, and let $B$ and $B_{12}$ be the constants of the
readings of the cube gauge condition in Definition~\ref{4.9:def-cube-gauge-condition}. Let
$C_{\mathrm{gl}}\ge1$ and $r_{\mathrm{gl}}\in(0,1]$ be the constants of Lemma~\ref{4.4:lem-gluing-boxes}.
Assume:
\begin{itemize}
\item the gluing of small gauges on a covered box, Lemma~\ref{4.4:lem-gluing-boxes};
\item the following property of the cover constructed direction by direction in the proof of
Corollary~\ref{4.4:cor-gluing-cubes}: with lengths measured in units of $L^j$ lattice steps, every
lattice run of length $T\ge M$, $T$ a multiple of $\xi_j$, is covered by runs of length $M$ which
satisfy, in that direction, the first four admissibility conditions of
Definition~\ref{4.4:def-admissible-covers}, and whose integer $n$ in the fifth admissibility
condition satisfies $(n+1)M\le 2T$;
\item $M\ge 80$ and $c_{12}L\ge1$;
\item the two conditions
\begin{equation}\label{4.9:eq-gluing-smaller-cubes-a}
(\mathrm{gl})_a:\quad 6S^2Ba\le r_{\mathrm{gl}},\qquad\qquad B_{12}\ge 4C_{\mathrm{gl}}B .
\end{equation}
\end{itemize}
Then for every closed set $X$ and every $G$-valued $U$ on $\operatorname{bd}(X)$, the reading
$(\ell 2)^{\flat}_j(X;a)$ implies the widened reading $(\ell 1)_j(X;a)$ of
\eqref{4.9:eq-cube-gauge-condition-b}.
\end{corollary}

\begin{proof}
We measure lengths in units in which $L^j$ lattice steps have length $1$. The lattice spacing is then
$\xi:=\xi_j=L^{-j}\le1$, a cube of side $SL^j$ lattice steps has side $S$, and the gradient of an
algebra-valued $A$ on pairs of parallel adjacent bonds is $\xi^{-1}(A(b+e_\nu)-A(b))$, as in
\eqref{4.9:eq-cube-gauge-condition-a}. Since $M=L^{m'}$ is an integer, $M$ is a multiple of $\xi$.

Two remarks are used throughout. First, the sets $\operatorname{bd}(D)$ and $\operatorname{pp}(D)$ are
determined by $\operatorname{st}(D)$: they consist of the bonds with both ends in
$\operatorname{st}(D)$ and of the parallel adjacent pairs of such bonds. Hence if
$\operatorname{st}(D')\subset\operatorname{st}(D)$ and $\rho\le\rho'$, restriction of $u$ and $A$,
exactly as in the proof of Lemma~\ref{4.9:lem-restriction-monotone}, gives
$\mathsf{g}_j(U;D;\rho)\Rightarrow\mathsf{g}_j(U;D';\rho')$. Second, since $c_{12}L\ge1$ we have
$M\le S$, and since $C_{\mathrm{gl}}\ge1$ the second condition in
\eqref{4.9:eq-gluing-smaller-cubes-a} gives $B\le B_{12}$; hence $MBa\le SB_{12}a$.

Let $\square^{\ast}\subset X$ be a closed cube of side $S^{\ast}_0\le S$. We must show
$\mathsf{g}_j(U;\square^{\ast};SB_{12}a)$. Its set of sites is a lattice box
$\operatorname{st}(\square^{\ast})=\prod_\mu J_\mu$, whose runs have lengths
$|J_\mu|:=\xi\cdot\#J_\mu\in(S^{\ast}_0-\xi,\,S^{\ast}_0+\xi]$; these are multiples of $\xi$, and
any two of them differ by at most $\xi$.

\emph{Case 1: $|J_\mu|\le M$ for every $\mu$.} A run of at most $M/\xi$ sites spans a closed interval
of length at most $M-\xi$. Hence $\operatorname{st}(\square^{\ast})\subset\operatorname{st}(\square)$
for a closed cube $\square$ of side $M$. As $\square^{\ast}\subset X$, also
$\operatorname{st}(\square^{\ast})\subset\operatorname{st}(\square\cap X)$. The reading
$(\ell 2)^{\flat}_j(X;a)$, with its parameter $C$ of Definition~\ref{4.9:def-cube-gauge-condition}
equal to $1$, taken for this cube $\square$ of side $M$, together with
Lemma~\ref{4.9:lem-restriction-monotone} in the form of the first remark, gives
$\mathsf{g}_j(U;\square^{\ast};MBa)$. Since $MBa\le SB_{12}a$, monotonicity in the radius gives
$\mathsf{g}_j(U;\square^{\ast};SB_{12}a)$.

\emph{Case 2: $|J_\nu|>M$ for some $\nu$.} Both $|J_\nu|$ and $M$ are multiples of $\xi$, so
$|J_\nu|\ge M+\xi$. Any two lengths differ by at most $\xi$, so $|J_\mu|\ge M$ for every $\mu$. In
each direction $\mu$ take the runs of length $M$ given by the second hypothesis for the run
$J_\mu$, of length $T=|J_\mu|$. Together they form an $M$-admissible cover $\{c_m\}$ of
$\prod_\mu J_\mu$ in the sense of Definition~\ref{4.4:def-admissible-covers}. The first four
admissibility conditions bear on one direction at a time and hold because they hold for the runs
of each direction. For the fifth condition put $S^{\ast}:=\max_\mu|J_\mu|$; then
\begin{equation*}
(n_\mu+1)M\le 2|J_\mu|\le 4S^{\ast},\qquad S^{\ast}\le S^{\ast}_0+\xi\le 2S ,
\end{equation*}
where $n_\mu$ is the integer that enters the fifth admissibility condition for the cover in
direction $\mu$, and the first inequality is the property of the second hypothesis applied to
$J_\mu$. Since $|J_\mu|\ge M$ for
all $\mu$, this $S^{\ast}$ is also the quantity $\max\{M,\max_\mu|J_\mu|\}$ with which
Lemma~\ref{4.4:lem-gluing-boxes} is applied to a cover by runs of length $M$. The side $M$ of the
cover satisfies $M\ge80\ge80\xi$ and $M\ge1$.

Each box $c_m$ is a lattice cube of side $M$ with
$c_m\subset\operatorname{st}(\square^{\ast})\subset\operatorname{st}(X)$. As in Case~1, $c_m$ lies in
a closed cube $\square_m$ of side $M$ with $c_m\subset\operatorname{st}(\square_m\cap X)$. The reading
$(\ell 2)^{\flat}_j(X;a)$ with its parameter $C$ equal to $1$ and the first remark then give a
$G$-valued $u_m$ on $c_m$ and a
$\mathfrak{g}$-valued $A_m$ with $U^{u_m}=e^{i\xi A_m}$ and $|A_m|,|\nabla A_m|<\mathsf{a}$ on $c_m$,
where $\mathsf{a}:=MBa$. Using $S^{\ast}\le2S$, $M\le S$ and the first condition in
\eqref{4.9:eq-gluing-smaller-cubes-a},
\begin{equation*}
S^{\ast}\mathsf{a}\le 2S\cdot MBa\le 2S^2Ba\le 6S^2Ba\le r_{\mathrm{gl}} .
\end{equation*}
The hypotheses of Lemma~\ref{4.4:lem-gluing-boxes} hold: $U$ is $G$-valued on the bonds of
$\square^{\ast}$, because these lie in $\operatorname{bd}(X)$. The lemma yields one $G$-valued $u$
on $\operatorname{st}(\square^{\ast})$ and a $\mathfrak{g}$-valued $A$ with $U^u=e^{i\xi A}$ and
\begin{equation*}
|A|,\ |\nabla A|\ \le\ C_{\mathrm{gl}}\frac{S^{\ast}}{M}\,MBa\ \le\ 2C_{\mathrm{gl}}SBa
\ <\ 4C_{\mathrm{gl}}SBa\ \le\ SB_{12}a .
\end{equation*}
Here the strict inequality holds because $C_{\mathrm{gl}},S,B,a>0$, and the last one is the
second condition in \eqref{4.9:eq-gluing-smaller-cubes-a}. These bounds hold on
$\operatorname{bd}(\square^{\ast})$ and $\operatorname{pp}(\square^{\ast})$. Hence
$\mathsf{g}_j(U;\square^{\ast};SB_{12}a)$.

In both cases $\mathsf{g}_j(U;\square^{\ast};SB_{12}a)$ holds for every closed cube
$\square^{\ast}\subset X$ of side at most $S$. This is $(\ell 1)_j(X;a)$.
\end{proof}

\begin{lemma}[Prolonging a cube within partition cubes]\label{4.9:lem-cube-prolongation}
Let $\mathsf{K}>0$ and let $\mathcal{P}$ be a cubic partition of mesh $\mathsf{K}$. Its cubes are
the products $\prod_\mu Q_\mu$ in which, for each coordinate direction $\mu$, the factor
$Q_\mu$ runs over the closed intervals of length $\mathsf{K}$ of a fixed partition of the
$\mu$-th coordinate line. These intervals are the \emph{partition intervals} in direction
$\mu$. Let $\Omega$ be a union of closed cubes of $\mathcal{P}$, and let
$\square=\prod_\mu I_\mu\subset\Omega$ be a closed cube of side $s>0$. For each $\mu$ let
$q_\mu$ be the number of partition intervals in direction $\mu$ whose interior meets
$I_\mu$, and put $q:=\min_\mu q_\mu$. Then $q\ge\max\{1,s/\mathsf{K}\}$. Moreover, for every
$t\in[s,q\mathsf{K}]$ there is a closed cube $\square'$ of side $t$ with
\begin{equation*}
\square\subset\square'\subset\Omega .
\end{equation*}
\end{lemma}

\begin{proof}
Fix a direction $\mu$. Let $V_\mu$ be the union of the $q_\mu$ partition intervals whose
interior meets $I_\mu$.

We first show that $V_\mu\supset I_\mu$. Let $x\in I_\mu$. If $x$ lies in the interior of a
partition interval, that interval's interior meets $I_\mu$. Otherwise $x$ is a common
endpoint of two adjacent partition intervals. Since $I_\mu$ has positive length $s$, it
contains points arbitrarily close to $x$ on at least one side of $x$. Hence the interior of
one of these two intervals meets $I_\mu$, and this interval contains $x$.

Since $I_\mu$ is connected, the intervals making up $V_\mu$ are consecutive. Hence $V_\mu$ is
a closed interval, and $|V_\mu|=q_\mu\mathsf{K}$. From $I_\mu\subset V_\mu$ we get
$q_\mu\mathsf{K}\ge s$, and $q_\mu\ge1$. Taking the minimum over $\mu$ gives
$q\ge\max\{1,s/\mathsf{K}\}$. For $t\in[s,q\mathsf{K}]$ we therefore have
\begin{equation*}
|V_\mu|=q_\mu\mathsf{K}\ \ge\ q\mathsf{K}\ \ge\ t\ \ge\ s .
\end{equation*}

Next we slide an interval of length $t$ inside $V_\mu$ so that it contains $I_\mu$. Write
$V_\mu=[v,v+q_\mu\mathsf{K}]$ and $I_\mu=[a,a+s]$, and set
\begin{equation*}
c:=\min\{a,\ v+q_\mu\mathsf{K}-t\},\qquad I'_\mu:=[c,c+t].
\end{equation*}
We check the required inclusions.
\begin{itemize}
\item $c\ge v$: indeed $a\ge v$ because $I_\mu\subset V_\mu$, and
$v+q_\mu\mathsf{K}-t\ge v$ because $t\le q_\mu\mathsf{K}$.
\item $c+t\le v+q_\mu\mathsf{K}$: this holds by the choice of $c$.
\item $c\le a$: this also holds by the choice of $c$.
\item $c+t\ge a+s$: indeed $a+t\ge a+s$ because $t\ge s$, and
$v+q_\mu\mathsf{K}\ge a+s$ because $I_\mu\subset V_\mu$.
\end{itemize}
Hence $I_\mu\subset I'_\mu\subset V_\mu$, and $I'_\mu$ has length $t$.

Put $\square':=\prod_\mu I'_\mu$. This is a closed cube of side $t$, and
\begin{equation*}
\square\subset\square'\subset\prod_\mu V_\mu .
\end{equation*}
It remains to prove that $\prod_\mu V_\mu\subset\Omega$.

Let $Q=\prod_\mu Q_\mu$ be a cube of $\mathcal{P}$ with $Q\subset\prod_\mu V_\mu$. For each
$\mu$, the partition interval $Q_\mu$ lies in $V_\mu$. Distinct partition intervals meet at
most in an endpoint, so $Q_\mu$ is one of the intervals making up $V_\mu$. Hence
$\operatorname{int}Q_\mu\cap I_\mu\neq\emptyset$ for every $\mu$, and therefore
\begin{equation*}
\operatorname{int}Q\cap\square=\prod_\mu\bigl(\operatorname{int}Q_\mu\cap I_\mu\bigr)\neq\emptyset .
\end{equation*}
Choose a point $y$ in this intersection.

The point $y$ lies in no closed cube of $\mathcal{P}$ other than $Q$. Indeed, such a cube
$\widetilde Q=\prod_\mu\widetilde Q_\mu$ differs from $Q$ in some direction $\nu$. Then the
distinct partition intervals $\widetilde Q_\nu$ and $Q_\nu$ meet at most in an endpoint, so
$\widetilde Q_\nu\cap\operatorname{int}Q_\nu=\emptyset$ and $y\notin\widetilde Q$.

On the other hand, $y\in\square\subset\Omega$, and $\Omega$ is a union of cubes of
$\mathcal{P}$. So $y$ lies in one of the cubes of $\Omega$, and that cube must be $Q$. Thus
$Q$ is one of the cubes of $\Omega$.

Finally, $\prod_\mu V_\mu$ is the union of the cubes $\prod_\mu Q_\mu$ of $\mathcal{P}$
obtained by choosing, for each $\mu$, one of the intervals $Q_\mu$ making up $V_\mu$. Each of
these cubes lies in $\prod_\mu V_\mu$, so by the previous paragraph it is one of the cubes of
$\Omega$. Hence $\prod_\mu V_\mu\subset\Omega$, and so $\square\subset\square'\subset\Omega$.
\end{proof}

The cube gauge condition of Definition~\ref{4.9:def-cube-gauge-condition} is the gauge clause of \cite[(1.12)]{B-CMP109}, \cite[(2.38)]{B-CMP119} and \cite[(1.65), (1.67), (1.69)]{B-CMP122a}, in any of its readings \eqref{4.9:eq-cube-gauge-condition-b}. It is a condition on the $G$-valued factor $U$ alone, and no schedule of radii enters it. The one-power hereditary class $U^{\Sigma^1}_j(X;\alpha_0,\alpha_1)$ is a union of $G^{c}$-orbits. A membership in it is therefore proved by exhibiting one representative with one decomposition $\mathbf{U}=U'U$. The following four moves of the $G$-valued factor $U$ are used in this chapter:
\begin{equation}\label{4.9:eq-unitary-moves-a}
\begin{aligned}
&(\mathrm{T}0)\quad U\mapsto 1;\\
&(\mathrm{T}1)\quad U\mapsto U\quad\text{(same level and domain; the radius not decreased)};\\
&(\mathrm{T}2)\quad U\mapsto U|_{X'},\quad X'\subset X,\quad \text{level } i\mapsto i'\le i;\\
&(\mathrm{T}3)\quad U\mapsto U^{v},\quad v\ \ G\text{-valued}.
\end{aligned}
\end{equation}

\begin{lemma}[Four moves of the unitary factor]\label{4.9:lem-unitary-moves}
Let $\star$ stand for any one of the readings $(\ell 1)$, $(\ell 2)$, $(\ell 2)^{\flat}$ of \eqref{4.9:eq-cube-gauge-condition-b}. In each reading the clause $\mathsf{g}_j(U;D;\rho)$ is understood with its bound $\xi_j^{-1}|A(b+e_\nu)-A(b)|<\rho$ on $\operatorname{pp}(D)$ included, $j$ being the level of the clause.
\begin{itemize}
\item[(T0)] For every level $i$, every closed set $X$ and every $a'>0$, the reading $\star_i(X;a')$ holds for the field $1$.
\end{itemize}
Let now $i$ be a level, $X$ a closed set, $U$ a $G$-valued field on $\operatorname{bd}(X)$ and $a>0$, and suppose that $\star_i(X;a)$ holds for $U$. Then:
\begin{itemize}
\item[(T1)] $\star_i(X;a')$ holds for $U$ at every radius $a'\ge a$, in particular at $a$;
\item[(T2)] for every closed $X'\subset X$ and every level $i'\le i$, the reading $\star_{i'}(X';\mu a)$ holds for $U|_{X'}$, where $\mu:=L^{-(i-i')}$;
\item[(T3)] for every $G$-valued $v$ on $\operatorname{st}(X)$, $\star_i(X;a)$ holds for $U^{v}$.
\end{itemize}
\end{lemma}

\begin{proof}
Each reading $\star_j(X;a)$ is a conjunction of clauses $\mathsf{g}_j(U;D;\rho)$. Each such $D$ is either a closed cube contained in $X$ or the intersection of a closed cube of space with $X$. The radius $\rho$ is the product of $a$ with a constant that does not depend on the level: $SB_{12}$ for $(\ell 1)$ and $(\ell 2)$, and $CMB$ for $(\ell 2)^{\flat}$, where $C$ denotes the cube-size constant of the reading $(\ell 2)^{\flat}$ in \eqref{4.9:eq-cube-gauge-condition-b}. A clause on $D\subset X$ involves $U$ only on $\operatorname{bd}(D)\subset\operatorname{bd}(X)$.

(T0). By Lemma~\ref{4.9:lem-identity-field}, the field $1$ satisfies $\mathsf{g}_j(1;D;\rho)$ for every $j$, $D$ and $\rho>0$. Hence every clause of $\star_i(X;a')$ holds, for every $a'>0$.

(T1). Nothing is changed except the radius. Let $a'\ge a$. Every clause $\mathsf{g}_i(U;D;\rho)$ of $\star_i(X;a)$ yields the clause with the same $D$ and radius $\rho a'/a\ge\rho$, by Lemma~\ref{4.9:lem-restriction-monotone} applied with $D'=D$.

(T2). Let $\square\subset X'$ be a closed cube of side at most $SL^{i'}$, as in $(\ell 1)_{i'}(X';\cdot)$. Since $i'\le i$, we have $SL^{i'}\le SL^{i}$. Hence $\square$ is a closed cube contained in $X$ of side at most $SL^{i}$, and $(\ell 1)_i(X;a)$ gives $\mathsf{g}_i(U;\square;SB_{12}a)$.

For the ambient readings, let $\square$ be a cube of space of side $SL^{i'}$ (of side $CML^{i'}$ for $(\ell 2)^{\flat}$). It lies in a cube of space $\square^{\ast}$ of side $SL^{i}$ (of side $CML^{i}$), and
\begin{equation*}
\square\cap X'\subset\square^{\ast}\cap X .
\end{equation*}
The reading at level $i$ gives the clause on $\square^{\ast}\cap X$ with radius $SB_{12}a$ (with radius $CMBa$). Lemma~\ref{4.9:lem-restriction-monotone} then gives the clause on $\square\cap X'$ with the same radius.

In all three cases we now apply Lemma~\ref{4.9:lem-level-change} with the levels $i'\le i$. It turns $\mathsf{g}_i(U;D;\rho)$ into $\mathsf{g}_{i'}(U;D;\mu\rho)$, with the same $u$. It thus multiplies the bound $SB_{12}a$ (respectively $CMBa$) by $\mu$, which is the bound of the reading at radius $\mu a$. The clauses involve $U$ only on $\operatorname{bd}(D)\subset\operatorname{bd}(X')$, so they hold for $U|_{X'}$. This proves $\star_{i'}(X';\mu a)$.

(T3). For each $D$ occurring in $\star_i(X;a)$ we have $\operatorname{st}(D)\subset\operatorname{st}(X)$. Lemma~\ref{4.9:lem-gauge-invariance}, applied with $v$ restricted to $\operatorname{st}(D)$, gives the equivalence of $\mathsf{g}_i(U;D;\rho)$ and $\mathsf{g}_i(U^{v};D;\rho)$. Hence $\star_i(X;a)$ holds for $U^{v}$.
\end{proof}

Throughout, a pair $(\mathbf{U},\mathbf{J})$ of the one-power hereditary class $U^{\Sigma^1}_j(X;\alpha_0,\alpha_1)$ is exhibited with a decomposition $\mathbf{U}=U'U$. Here $U'$ is a complex factor and $U$ is a $G$-valued factor, the \emph{unitary factor} of the representative. The gauge clause of the class, its clause (i), is the reading $(\ell 1)=(\mathrm{L}\text{-}12)$ of \eqref{4.9:eq-cube-gauge-condition-b}, imposed on $U$. We write $(i,X;a)$ for the level $i$, the closed domain $X$ and the radius $a$ at which a reading of the cube gauge condition of Definition~\ref{4.9:def-cube-gauge-condition} is imposed. For levels $i'\le i$ we put $\mu=L^{-(i-i')}$, as in Lemma~\ref{4.9:lem-level-change}. The four moves $(\mathrm{T}0)$–$(\mathrm{T}3)$ of the unitary factor are those of \eqref{4.9:eq-unitary-moves-a}. The constants $\beta$, $\beta_s$ and $\beta_0$ are the ratio-of-radii constants of the running-shift statement: $\beta$ enters the radii $(1+\beta)\alpha_0$ and $(1+\beta)\alpha_{0,k+1}$ of the one-power hereditary class at level $k+1$ in the items below, $\beta_s$ enters the radius $(1+\beta_s)\alpha_{0,i}$ of the running form of the far-source statement, and $\beta_0$ is the constant of the bound on the ratio of radii in Lemma~\ref{4.4:lem-radii-ratio}.

The loci at which the step exhibits a representative are listed below. Each item gives:
\begin{itemize}
\item the source and the product;
\item the unitary factor exhibited;
\item the move it realises;
\item the one-sided condition under which the radius of the product is reached.
\end{itemize}

By the moves they realise, the loci group as follows:
\begin{equation}\label{4.9:eq-gauge-transport-a}
\begin{aligned}
&(\mathrm{T}0):\ \langle 2\rangle,\ \langle 4'\rangle,\ \langle 5\mathrm{a}\rangle,\
  \langle 5\mathrm{b}\rangle(\mathrm{a});\qquad
 (\mathrm{T}1):\ \langle 4\rangle,\ \langle 5\mathrm{b}\rangle(\mathrm{b}),\ \langle 7\rangle,\
  \langle 10\rangle,\ \langle 11\rangle,\ \langle 13\rangle;\\
&(\mathrm{T}2):\ \langle 3\rangle,\ \langle 4\rangle\ (\text{inner}),\ \langle 6\rangle,\
  \langle 7\rangle,\ \langle 8\rangle,\ \langle 13\rangle;\\
&\text{prolongation}:\ \langle 8\rangle,\ \langle 13\rangle;\qquad
 \text{gluing, after the moves of }\langle 8\rangle:\ \langle 9\rangle;\\
&\text{seeds}:\ \langle 5\mathrm{c}\rangle,\ \langle 13\rangle\ \text{in its real-background clause};\qquad
 \text{no move}:\ \langle 1\rangle,\ \langle 12\rangle.
\end{aligned}
\end{equation}

\begin{enumerate}
\item[$\langle 1\rangle$] \emph{Near source.}
The statement is the near-source statement of the one-power hereditary class.
The source is a representative of the one-power hereditary class at $(k+1,Y;(1+\beta)\alpha_0)$.
The product is the expression \cite[(3.40)]{B-CMP109} on the enlarged cube $\widetilde{\square}^{3}$.
Only clause (iv) of the class is read here.
The gauge clause is not obtained at this locus; it is used at $\langle 10\rangle$.
There is no move and no condition.

\item[$\langle 2\rangle$] \emph{Near product.}
The statement is the near-product lemma of the one-power hereditary class.
Its source is the pair of \cite[Lemma~4]{B-CMP109}; the product is a membership in $U^{\Sigma^1}_j(X;\alpha_0,\alpha_1)$ with $X\subset\widetilde{\square}^{2}$.
The decomposition is $U:=1$, $U':=\mathcal{U}$, where $\mathcal{U}$ denotes the configuration of that pair (a configuration, not a class).
Clause \cite[(1.11)]{B-CMP109} reads $|\partial U-1|=0$.
Clause \cite[(1.12)]{B-CMP109}, in the reading $(\mathrm{L}\text{-}12)$, holds on every cube with $u:=1$, since $U^{u}=1=\exp(i\xi_j\cdot 0)$.
The requirement is therefore vacuous at every radius, and no lower bound on $B_{12}$ is used.
The large-field part of the coupling increment uses the same decomposition, with $u=1$ and vanishing potential on every cube $\subset X$.
The move is $(\mathrm{T}0)$; there is no condition.

\item[$\langle 3\rangle$] \emph{Far source, also in its running form.}
The statement is the far-source statement of the one-power hereditary class.
The source is at $(k+1,Y;(1+\beta)\alpha_0)$.
The product is at $(j,X;\tfrac12\alpha_0)$, for every closed $X\subset Y$ and every $j\le k$.
It is exhibited with the same $G$-valued factor $U$ and the restricted gauges, and the decomposition $\mathbf{U}=U'U$ is the same.
The move is $(\mathrm{T}2)$.
The condition is $\mu(1+\beta)\le\tfrac12$, which follows from $L\ge 2(1+\beta)$.
In the running form, where the source is at a level $i$ and the product at a level $i'\le i$, the condition is
\begin{equation*}
(1+\beta_s)\,\frac{\alpha_{0,i}}{\alpha_{0,i'}}\,L^{-(i-i')}\le\tfrac12 ,
\end{equation*}
where $L^{-(i-i')}$ is the factor $\mu$ of Lemma~\ref{4.9:lem-level-change} for this descent;
the condition follows from the floor of the running form of the far-source statement.

\item[$\langle 4\rangle$] \emph{Far product.}
This is the far-product clause of the operator bounds at free current, obtained through the lemma of
the operator bounds at free current that exhibits the representative displayed below, together with
\cite[(3.13)]{B-CMP109}.
The source is at $(j,X;\tfrac12\alpha_0)$ and the product at $(j,X;\alpha_0)$.
The representative is $e^{i\xi A}\mathbf{U}=e^{i\xi A''}U$ on $X$, where $\xi$ denotes the lattice spacing and $A,A''$ are the algebra-valued fields of that statement.
The $G$-valued factor is the same $U$, and \cite[(1.11)--(1.12)]{B-CMP109} are untouched.

Inside it there is an inner instantiation at a level $p$ and a domain $Z$.
It uses the cube side parameter $M/L$ in place of $M$ and the constant $LB_{12}$ in place of $B_{12}$.
There, \cite[(1.12)]{B-CMP109}, read by $(\rho1)$ of Notation~\ref{4.5:not-reading-conventions}, is required on every cube $\subset Z$ of side at most $c_{12}L(M/L)L^{p}$ steps.
Since $c_{12}L(M/L)L^{p}\le c_{12}LML^{j}$, every such cube is a cube of the source.
The moves are $(\mathrm{T}1)$, and $(\mathrm{T}2)$ for the inner instantiation.
The radius comparison is $\tfrac12\le1$; for the inner instantiation the radius identity is
\begin{equation*}
c_{12}L(M/L)(LB_{12})\cdot\tfrac12\mu\alpha_0=\mu\cdot c_{12}LMB_{12}\cdot\tfrac12\alpha_0 .
\end{equation*}

\item[$\langle 4'\rangle$] \emph{The two far-type terms} \cite[(3.21)--(3.22)]{B-CMP109}.
The source is the near product $\langle 2\rangle$, read at the radius $\tfrac12\alpha_0$ at which
the source of $\langle 4\rangle$ is taken; the product is that of $\langle 4\rangle$.
The representative is $e^{i\xi A}W=(e^{i\xi A^{\sharp}}\mathcal{U})^{u_j}$, where $W$, $A^{\sharp}$ and the gauge $u_j$ are those of the two terms.
This is the representative $e^{i\xi A^{\sharp}}\mathcal{U}$ of the orbit, over $U=1$.
The move is $(\mathrm{T}0)$; there is no condition.

\item[$\langle 5\mathrm{a}\rangle$] \emph{Kernels} of \cite[(4.4)]{B-CMP109} at the pair $(1,0)$.
The product is a membership in $U^{\Sigma^1}_j(X;\cdot)$, and the representative is the unit pair.
The move is $(\mathrm{T}0)$; there is no condition.

\item[$\langle 5\mathrm{b}\rangle$] \emph{Small pairs}, by the small-pairs lemma.
The source condition is $\mathsf{N}_j(\mathsf{Z})\le c_{T}r_j$, where $\mathsf{N}_j$, $c_T$ and $r_j$ are the norm, constant and radius of that lemma.
The product is a membership in $U^{\Sigma^1}_j(X;\alpha_{0,j},\alpha_{1,j})$.
\begin{enumerate}
\item[(a)] The pair is $(e^{i\eta\mathsf{Z}},J(e^{i\eta\mathsf{Z}}))$, with $\mathsf{Z}$ a $\mathfrak{g}^{c}$-valued field, the whole of which is measured by $\mathsf{N}_j$ against $r_j$. The decomposition is $U:=1$, $U':=e^{i\eta\mathsf{Z}}$. For real $\mathsf{Z}$ this decomposition is also admissible, since the class asks for a decomposition, not for a canonical one.
\item[(b)] The pair is $(We^{i\eta\mathsf{Z}},J_W+\mathfrak{j})$. Here $W=U'_WU_W$ is a representative of the class with unitary factor $U_W$ and current $J_W$, and $\mathfrak{j}$ is the current increment of the lemma. Then $We^{i\eta\mathsf{Z}}=(U'_We^{i\eta\operatorname{Ad}_{U_W}\mathsf{Z}})\,U_W$, so the unitary factor is that of $W$.
\end{enumerate}

The moves are $(\mathrm{T}0)$ in (a) and $(\mathrm{T}1)$ in (b); there is no condition.

\item[$\langle 5\mathrm{c}\rangle$] \emph{Real backgrounds.}
The statement is the real-background clause of the one-power hereditary class.
The product is clause (i) at $\alpha_0'$.
Here $\mathbf{U}=U:=U_k|_X$ and $U'=1$.
This is a seed: its membership is not derived from another membership.
The input is the real-background theorem.

\item[$\langle 6\rangle$] \emph{Change of domain.}
The statement is the change-of-domain statement of the one-power hereditary class.
The source is at $(k+1,X;\alpha_0)$ and the product at $(k+1,Y;(1+\beta)\alpha_0)$, with $Y\subset X$.
This is restriction with $\mu=1$.
The move is $(\mathrm{T}2)$; the condition is $1\le1+\beta$.

\item[$\langle 7\rangle$] \emph{The product string} of Corollary~\ref{4.4:cor-localised-independence}.
The source is a representative at $(k+1,Y;(1+\beta)\alpha_{0,k+1})$.
The product is $(\mathbf{U}^{\flat},\mathbf{J}^{\flat})|_X\in U^{\Sigma^1}_j(X;\alpha_{0,j},\alpha_{1,j})$, for $j\le k$.
The representative is $U^{\flat}:=U$, $U'^{\flat}:=e^{i\xi\mathcal{H}}U'=:\exp(i\xi A'^{\flat})$.
The whole twist $\mathcal{H}$, real or complex, is placed in the complex factor, and $U^{\flat}=U$.
Nothing is glued: the cubes required of the product are among those of the source.
The moves are $(\mathrm{T}2)$ and $(\mathrm{T}1)$.
The condition, which follows from Lemma~\ref{4.4:lem-radii-ratio} and \eqref{4.9:eq-gauge-transport-1} below, is
\begin{equation*}
(1+\beta)\,\frac{\alpha_{0,k+1}}{\alpha_{0,j}}\,L^{-(k+1-j)}\le\tfrac12 .
\end{equation*}

\item[$\langle 8\rangle$] \emph{Transport of \eqref{4.9:eq-cube-gauge-condition-c}.}
The statement is the first case of the passage from level $k+1$ to level $j$ in Corollary~\ref{4.4:cor-localised-independence}, the case in which the source condition is \eqref{4.9:eq-cube-gauge-condition-c}.
The source is \eqref{4.9:eq-cube-gauge-condition-c} at level $k+1$ on $Y$; the product is \eqref{4.9:eq-cube-gauge-condition-c} at level $j$ on $X$.
The passage is made for a nest whose partitions satisfy the following condition on their meshes, called below the partition condition: each member $\Omega_{n_0}$ is a union of closed cubes of a cubic partition, compatible with it, of mesh at least $ML^{n_0}\xi$. In the nest, moreover, consecutive members are separated as in \cite[(2.2)]{B-CMP96}:
\begin{equation*}
\operatorname{dist}(\Omega^{c}_{n_0+1},\Omega_{n_0+2})>RML^{n_0+1}\xi .
\end{equation*}
The unitary factor is $U^{\flat}=U$. The comparison of cubes is the claim of the paragraph of Corollary~\ref{4.4:cor-localised-independence} on \eqref{4.9:eq-cube-gauge-condition-c}, used here as stated there; its device is Lemma~\ref{4.9:lem-cube-prolongation}, whose hypothesis that the domain be a union of closed cubes of a cubic partition is supplied by the partition condition, together with the separation displayed above.
The move is Lemma~\ref{4.9:lem-cube-prolongation} followed by $(\mathrm{T}2)$.
For a member $\Omega_{n_0}$ of the nest, the condition is
\begin{equation*}
L^{-(n_0-j)}(1+\beta)\alpha_{0,n_0}\le\alpha_{0,j},
\end{equation*}
which follows from Lemma~\ref{4.4:lem-radii-ratio} and \eqref{4.9:eq-gauge-transport-1}.

\item[$\langle 9\rangle$] \emph{Straddling sources.}
The statement is the second case of that passage in Corollary~\ref{4.4:cor-localised-independence}, the case of straddling sources.
The source is \eqref{4.9:eq-cube-gauge-condition-c} on $Y$.
The product is the class of \cite[(1.11)--(1.14)]{B-CMP109} at $(j,X)$, with $X\subset\Lambda_j$.
The unitary factor is $U^{\flat}=U$.
Clause \cite[(1.12)]{B-CMP109}, in the reading $(\mathrm{L}\text{-}12)$, is obtained through Corollary~\ref{4.4:cor-gluing-cubes} with its side parameter equal to $M$.
The move is $\langle 8\rangle$ followed by Corollary~\ref{4.9:cor-gluing-smaller-cubes}.
The condition is that the hypotheses of Corollary~\ref{4.9:cor-gluing-smaller-cubes} hold,
\eqref{4.9:eq-gluing-smaller-cubes-a} among them.

\item[$\langle 10\rangle$] \emph{The operators' glued background $\mathbf{U}^{\mathrm{gl}}$.}
The statement is the proposition on the operators' glued background.
The source is a pair $(\mathbf{U},\mathbf{J})$ of the space $\mathfrak{U}$ of pairs from which that
proposition starts.
The product is $((\mathbf{U}^{\mathrm{gl}})^{h},\operatorname{Ad}_h\mathbf{J})\in\mathfrak{S}$ of Definition~\ref{4.5:def-admissible-pairs}, with $h$ a $G$-valued gauge transformation.
The unitary factor read by $\mathfrak{S}$ is the source's $U$ (the factor that proposition denotes $U_r$).
In clause (S1) of $\mathfrak{S}$, $U$ is unchanged, and \cite[(1.11)--(1.12)]{B-CMP109} hold
verbatim on cubes straddling the boundaries $\partial\square_0$, $\partial\square_1$ of the two
cubes $\square_0$, $\square_1$ used by that proposition in forming $\mathbf{U}^{\mathrm{gl}}$,
as on any other cube.
The whole displacement of $\mathbf{U}^{\mathrm{gl}}$ from $\mathbf{U}$, real part included, is placed in the twist.
The move is $(\mathrm{T}1)$.
The condition is the ceiling $c_{12}LMB\alpha_0'''\le a_0$ on the radius $\alpha_0'''$ of the
operator class $\mathfrak{S}$ of Definition~\ref{4.5:def-admissible-pairs}.

\item[$\langle 11\rangle$] \emph{The real point}, by the real-point lemma of the one-power hereditary class.
The source is $(U'U,\mathbf{J})$ and the product is $(U,0)$.
Clause (i) of the product is that of $U$.
The move is $(\mathrm{T}1)$; there is no condition.

\item[$\langle 12\rangle$] \emph{Cut-off extensions.}
The statement is that of the cut-off extensions made in Corollary~\ref{4.4:cor-localised-independence}.
The source is \cite[(1.12)]{B-CMP109} localised on a cube $\square$.
The product is
\begin{equation*}
\widehat{U}:=\exp(i\xi\eta_{\square}A)=:\exp(i\xi\widehat{A})\in\mathcal{U}_{\mathbb{R}},
\end{equation*}
with $\eta_{\square}$ the cut-off function at $\square$ (a function on the lattice, not the spacing $\eta$).
This is not a membership in the one-power hereditary class.
The gauge is trivial on every cube, and $|\widehat{A}|<\mathsf{a}_U$ and $|\nabla\widehat{A}|<2\mathsf{a}_U$ everywhere.
Both members of the clause are used.
The class $\mathcal{U}_{\mathbb{R}}$ is taken at $2\mathsf{a}_U$.

\item[$\langle 13\rangle$] \emph{Composites of the large-field operation, direct form.}
The statements are the second step of the theorem on the composites of the large-field operation and the theorem on their direct form, which exhibits the composites together with their cube gauge clause.
The source is the space $\widetilde{U}^{c}_k(Y_4)$ of the new nest.
The product is the data domain $D_E(X')$ of every statement $E$ that uses the composites.
One sets $\bar{U}_i:=\bar{U}$, the source's unitary factor.
The cube clause of $\mathcal{W}_i$ is that of $\bar{U}$, inherited unchanged, so no loss of radius occurs.
In the real-background clause of that theorem the unitary factor is that of its real background
$U_2$, which is real by that theorem.
No re-minimisation, no application of the minimising operator of that theorem, and no
block-constant gauge factor acts on $\bar{U}$.
The move is $(\mathrm{T}1)$ or $(\mathrm{T}2)$, followed by Lemma~\ref{4.9:lem-cube-prolongation};
in the real-background clause the factor is a seed.
The input for that seed is the seed theorem for the composites of the large-field operation.
\end{enumerate}

The floor used at $\langle 7\rangle$ and $\langle 8\rangle$ is
\begin{equation}\label{4.9:eq-gauge-transport-1}
L\ge 2\sqrt{2}\,(1+\beta)(1+\beta_0).
\end{equation}

\begin{proposition}[The step transports the gauge condition]\label{4.9:prop-gauge-transport}
Assume the following.
\begin{itemize}
\item[(D)] At each locus $\langle 2\rangle$–$\langle 13\rangle$ above, the statement cited in that item exhibits the representative of its product with the unitary factor displayed in that item; at $\langle 1\rangle$ only clause (iv) is read.
\end{itemize}
Let $\star$ be any of $(\ell 1)=(\mathrm{L}\text{-}12)$, $(\ell 2)$, $(\ell 2)^{\flat}$ of \eqref{4.9:eq-cube-gauge-condition-b}, read with its $\nabla$-member and with one pair of constants $(B,B_{12})$. Then:
\begin{itemize}
\item[(a)] At $\langle 2\rangle$, $\langle 4'\rangle$, $\langle 5\mathrm{a}\rangle$ and $\langle 5\mathrm{b}\rangle$(a) the exhibited representative satisfies $\star$ at every radius, unconditionally.
\item[(b)] At $\langle 3\rangle$, $\langle 4\rangle$, $\langle 5\mathrm{b}\rangle$(b), $\langle 6\rangle$, $\langle 7\rangle$ and $\langle 11\rangle$ it satisfies $\star$ at the product's level, domain and radius, whenever the source's representative satisfies $\star$ at the source's level, domain and radius. The only conditions are those in the last entry of each item. There is no condition on $M$, on the shape or thickness of $X$ or $Y$, on the position of $X$ in $Y$, on $i-i'$, on $k$, on the nest, on the volume or on the coupling.
\item[(c)] At $\langle 8\rangle$, for a nest as in $\langle 8\rangle$, satisfying the partition
condition and the separation displayed there,
\eqref{4.9:eq-cube-gauge-condition-c} is transported from level $k+1$ on $Y$ to level $j$ on $X$
by the claim on \eqref{4.9:eq-cube-gauge-condition-c} in
Corollary~\ref{4.4:cor-localised-independence}, whose device is
Lemma~\ref{4.9:lem-cube-prolongation}, followed by $(\mathrm{T}2)$, under the last entry of
$\langle 8\rangle$. At $\langle 9\rangle$:
\begin{itemize}
\item $\star=(\ell 1)$ holds by Corollary~\ref{4.9:cor-gluing-smaller-cubes}, under its
hypotheses, \eqref{4.9:eq-gluing-smaller-cubes-a} among them;
\item $\star=(\ell 2)^{\flat}$ holds by part (iv) of the proposition on the ambient readings at
straddling sources, stated later in this section;
\item $\star=(\ell 2)$ fails, by part (iii) of that proposition.
\end{itemize}
\item[(d)] The loci $\langle 10\rangle$ and $\langle 12\rangle$ yield no membership in the one-power hereditary class; they use the clause. At $\langle 10\rangle$ the same $U$ is handed to the operator class $\mathfrak{S}$.
\end{itemize}
Consequently, at the four loci $\langle 2\rangle$, $\langle 4\rangle$, $\langle 5\mathrm{b}\rangle$, $\langle 7\rangle$ the condition is not created, for $(\mathrm{L}\text{-}12)$ and for $(\ell 2)$ alike. At two of them the requirement is vacuous. At the other two the set of cube clauses required of the product coincides with a set of cube clauses already required of the source. Moreover, the seed corollary of the one-power hereditary class holds in clause (i):
\begin{itemize}
\item at five of its six loci, with no input;
\item at the sixth, $\langle 5\mathrm{c}\rangle$, if and only if the seed requirement of Remark~\ref{4.9:rem-reproduction-meaning}, \eqref{4.9:eq-reproduction-meaning-a}, holds there.
\end{itemize}
\end{proposition}

\begin{proof}
The proof applies Lemma~\ref{4.9:lem-unitary-moves} item by item.
In each item, the unitary factor exhibited, given by (D), identifies the move.
The last entry compares the factor $\mu$ of Lemma~\ref{4.9:lem-level-change} with the ratio of the radii.
Lemma~\ref{4.9:lem-unitary-moves} gives the following:
\begin{itemize}
\item after $(\mathrm{T}0)$, each of $(\ell 1)$, $(\ell 2)$, $(\ell 2)^{\flat}$, $\nabla$-member included, holds at every radius;
\item after $(\mathrm{T}1)$ and $(\mathrm{T}3)$, it holds at the same radius;
\item after $(\mathrm{T}2)$ from $(i,X;a)$ to $(i',X')$, it holds at the radius $\mu a$.
\end{itemize}

Lemma~\ref{4.9:lem-unitary-moves} carries no hypothesis on $M$, on the shape or position of the
domains, on $i-i'$, on $k$, on the nest, on the volume or on the coupling.
By Lemma~\ref{4.9:lem-restriction-monotone}, each reading is monotone in the radius.
Hence the product's radius is reached as soon as the radius delivered by the move does not exceed it.

(a) By (D), the unitary factor is $1$ at all four loci:
\begin{itemize}
\item at $\langle 2\rangle$ the decomposition is $U:=1$;
\item at $\langle 4'\rangle$ it is the orbit's representative over $U=1$;
\item at $\langle 5\mathrm{a}\rangle$ it is the unit pair;
\item at $\langle 5\mathrm{b}\rangle$(a) the decomposition is $U:=1$.
\end{itemize}

These are therefore $(\mathrm{T}0)$, and $\star$ holds at every radius.

(b) The six loci are treated in turn.

At $\langle 3\rangle$ the factor is the source's $U$, restricted to $X\subset Y$ and read at level $j\le k$. This is $(\mathrm{T}2)$ with $\mu=L^{-(k+1-j)}\le L^{-1}$, and it delivers the radius $\mu(1+\beta)\alpha_0$. If $L\ge2(1+\beta)$, then
\begin{equation*}
\mu(1+\beta)\alpha_0\le\frac{(1+\beta)\alpha_0}{L}\le\tfrac12\alpha_0 .
\end{equation*}
In the running form, the source radius $(1+\beta_s)\alpha_{0,i}$ becomes $\mu(1+\beta_s)\alpha_{0,i}$ with $\mu=L^{-(i-i')}$. This is at most $\tfrac12\alpha_{0,i'}$ by the displayed condition of $\langle 3\rangle$.

At $\langle 4\rangle$ the factor is the same $U$ at the same $(j,X)$. This is $(\mathrm{T}1)$, since the radius is not decreased: $\tfrac12\alpha_0\le\alpha_0$.
For the inner instantiation at $(p,Z)$, with cube side parameter $M/L$ and constant $LB_{12}$, two facts hold.
First, its cubes have side at most $c_{12}L(M/L)L^{p}\le c_{12}LML^{j}$, so they are cubes of the source.
Second, the cube side parameter $M/L$ and the constant $LB_{12}$ leave the product unchanged:
\begin{equation*}
c_{12}L(M/L)\cdot LB_{12}=c_{12}LMB_{12}.
\end{equation*}
So the radius demanded, $c_{12}L(M/L)(LB_{12})\cdot\tfrac12\mu\alpha_0$, equals $\mu$ times the source's radius $c_{12}LMB_{12}\cdot\tfrac12\alpha_0$, with $\mu=L^{-(j-p)}$. The inner instantiation is therefore exactly $(\mathrm{T}2)$ from level $j$ to level $p$.
For $(\ell 2)^{\flat}$ the same holds: with $C$ the constant in the side of its ambient cubes, $C(M/L)(LB)=CMB$.

At $\langle 5\mathrm{b}\rangle$(b), multiplying the relation $U_We^{i\eta\mathsf{Z}}U_W^{-1}=e^{i\eta\operatorname{Ad}_{U_W}\mathsf{Z}}$ on the left by $U'_W$ and on the right by $U_W$ gives, bondwise,
\begin{equation*}
U'_WU_We^{i\eta\mathsf{Z}}=U'_We^{i\eta\operatorname{Ad}_{U_W}\mathsf{Z}}U_W .
\end{equation*}
So the unitary factor is the source's $U_W$, and the move is $(\mathrm{T}1)$ at the same radius.

At $\langle 6\rangle$ the move is $(\mathrm{T}2)$ with $i=i'$, so $\mu=1$. The delivered radius $\alpha_0$ is at most $(1+\beta)\alpha_0$.

At $\langle 7\rangle$ the factor is $U^{\flat}=U$. Restriction from $(k+1,Y)$ to $(j,X)$ is $(\mathrm{T}2)$, and the change of the complex factor to $e^{i\xi\mathcal{H}}U'$ is $(\mathrm{T}1)$. The delivered radius is $\mu(1+\beta)\alpha_{0,k+1}$ with $\mu=L^{-n}$, where $n:=k+1-j\ge1$.
Lemma~\ref{4.4:lem-radii-ratio} gives $\alpha_{0,k+1}\le(1+\beta_0)(1+n)^{1/2}\alpha_{0,j}$. Together with \eqref{4.9:eq-gauge-transport-1} this yields
\begin{equation*}
(1+\beta)\,\frac{\alpha_{0,k+1}}{\alpha_{0,j}}\,L^{-n}
\le(1+\beta)(1+\beta_0)(1+n)^{1/2}L^{-n}
\le\frac{(1+n)^{1/2}L^{1-n}}{2\sqrt{2}}\le\tfrac12 .
\end{equation*}
The last inequality holds because $(1+n)^{1/2}\le\sqrt{2}\,L^{n-1}$ for $n\ge1$ and $L\ge2$. Hence the delivered radius is at most $\tfrac12\alpha_{0,j}\le\alpha_{0,j}$.

At $\langle 11\rangle$ the factor of $(U,0)$ is the source's $U$, so the move is $(\mathrm{T}1)$.

(c) At $\langle 8\rangle$ the factor is $U^{\flat}=U$ by (D). The cubes required at level $j$ are compared with those of the source by the claim of the paragraph of Corollary~\ref{4.4:cor-localised-independence} on \eqref{4.9:eq-cube-gauge-condition-c}, used here as stated there. Its device is Lemma~\ref{4.9:lem-cube-prolongation}, which applies because, by the partition condition of $\langle 8\rangle$, each member $\Omega_{n_0}$ of the nest is a union of closed cubes of a cubic partition of mesh at least $ML^{n_0}\xi$; the comparison also uses the separation of consecutive members stated in $\langle 8\rangle$. The move $(\mathrm{T}2)$ then delivers the radius $L^{-(n_0-j)}(1+\beta)\alpha_{0,n_0}$. This is at most $\alpha_{0,j}$ by the last entry of $\langle 8\rangle$, which Lemma~\ref{4.4:lem-radii-ratio} and \eqref{4.9:eq-gauge-transport-1} give as at $\langle 7\rangle$.

At $\langle 9\rangle$ the factor is again $U$.
\begin{itemize}
\item $(\ell 1)$: after $\langle 8\rangle$, Corollary~\ref{4.9:cor-gluing-smaller-cubes} applies under its hypotheses, \eqref{4.9:eq-gluing-smaller-cubes-a} among them. Its gluing step is Corollary~\ref{4.4:cor-gluing-cubes} with side parameter $M$.
\item $(\ell 2)^{\flat}$ and $(\ell 2)$: the two assertions are parts (iv) and (iii) of the
proposition on the ambient readings at straddling sources, stated later in this section.
\end{itemize}

(d) At $\langle 10\rangle$ the product is a member of $\mathfrak{S}$, not of the one-power hereditary class. Its unitary factor is the source's $U$ unchanged, which is $(\mathrm{T}1)$, under the ceiling $c_{12}LMB\alpha_0'''\le a_0$.
At $\langle 12\rangle$ the product is a member of $\mathcal{U}_{\mathbb{R}}$. It is obtained from both members of the clause, with the trivial gauge on every cube.

For the consequences, consider first the four loci $\langle 2\rangle$, $\langle 4\rangle$, $\langle 5\mathrm{b}\rangle$, $\langle 7\rangle$.
\begin{itemize}
\item At $\langle 2\rangle$ and at $\langle 5\mathrm{b}\rangle$ in form (a), the factor is $1$ and the requirement is vacuous, by (a). At $\langle 5\mathrm{b}\rangle$ in form (b), the cube clauses of the product are those of $W$, by (b).
\item At $\langle 4\rangle$ and $\langle 7\rangle$, the factor is the source's $U$, and the cube clauses required of the product are cube clauses of the source, by (b). Since (a) and (b) hold for each of the three readings, this holds for $(\mathrm{L}\text{-}12)$ and for $(\ell 2)$ alike.
\end{itemize}

Now consider the seed corollary.
At $\langle 5\mathrm{c}\rangle$ the factor $U_k|_X$ is not derived from another membership. Its clause (i) at $\alpha_0'$ is therefore exactly what the seed requirement of Remark~\ref{4.9:rem-reproduction-meaning} asserts there.
At each of the five other loci of that corollary, clause (i) follows from (a) or (b), with no input.
\end{proof}

\begin{remark}[The twist belongs to the complex factor]\label{4.9:rem-twist-complex-factor}
Write the variables of the step as $\mathbf{U}=U'U$, with $U'$ the complex factor and $U$ the unitary
factor, as in \cite[(1.11)--(1.14)]{B-CMP109}. Throughout this chapter the whole twist of the step,
real or complex, is placed in the complex factor $U'$, and the unitary factor $U$ is not moved. Let
$\star$ denote any of the readings $(\ell 1)$, $(\ell 2)$, $(\ell 2)^{\flat}$ of the gauge clause, as
in Proposition~\ref{4.9:prop-gauge-transport}. The clause $\star$ is asked of one representative of
a gauge orbit. A $G^{c}$-valued map that is not $G$-valued acts on the orbit and never on the
exhibited factor.
\end{remark}

\begin{proof}
The placement of the twist is forced by the form in which the step controls its functions. There are
three points.

\emph{The covariant first difference.} Clause~(ii) of the conditions (i), (ii) imposed on the pair
$(U',U)$ in \cite{B-CMP109} measures
the complex factor $U'$ against the radius $\alpha_1$ through a covariant first difference. Every
response function of the step satisfies a bound of this kind. This holds for $\mathcal{H}$, for
$\mathbf{H}_j$, and for the response function written $A(t_{\square})$ in \cite{B-CMP109}, by
\cite[(3.9), (3.14), (3.37)]{B-CMP109} together with the response-function bound of the step.
So a twist placed in $U'$ is measured by the very quantity the step bounds.

Now suppose instead that the twist moved the unitary factor. Clause~(ii) would then need a bound on a
first difference of a potential of the moved unitary factor. Nothing in the step provides such a
bound. Moreover, no estimate of a first difference of such a potential in terms of first-order
quantities of the moved factor can hold in general; this is a classical non-inequality for difference
operators. The twist therefore cannot be carried by $U$.

Ba{\l}aban proceeds in the same way at the one place where he exhibits the representative, in
\cite{B-CMP109}: there the orbit of the level-$j$ configuration considered at that place is shown to
contain the configuration $\exp i\xi\mathbf{H}_j$, and this configuration satisfies the conditions (i)
and (ii) above. In that configuration the whole of $\mathbf{H}_j$ is measured against $\alpha_1$. That
is, the unitary factor is taken to be $U=1$.

\emph{The re-minimised representative of \cite{B-CMP122b}.} One construction in Ba{\l}aban's papers
does move the unitary factor. It is the re-minimised representative of
\cite[(1.61)--(1.65)]{B-CMP122b}: the configuration in the minimal orbit determined by the
constraint space of \cite[Theorem~1]{B-CMP102b} and the averaged data $\dot{M}(\mathbf{U})$. This
representative is not used in this chapter.

\emph{Orbits.} The clause $\star$ is a condition on one representative, namely the exhibited one.
Let $u$ be a $G^{c}$-valued map that is not $G$-valued. For a bond $b$, the gauge transform is
\begin{equation*}
  \mathbf{U}^{u}(b) = u(b_-)\,\mathbf{U}(b)\,u(b_+)^{-1} .
\end{equation*}
Replacing $\mathbf{U}$ by $\mathbf{U}^{u}$ stays within the orbit of $\mathbf{U}$. Every such map is
read as acting on the orbit of $\mathbf{U}$; the factor $U$ of the exhibited representative, for
which $\star$ is asserted, is left unchanged. Accordingly, whenever a non-unitary gauge map is applied in
the proof of Proposition~\ref{4.9:prop-gauge-transport}, the clause $\star$ is checked on the
exhibited representative, and not on its image under that map.
\end{proof}

The gradient bound of the cube gauge condition (Definition~\ref{4.9:def-cube-gauge-condition})
is asked of pairs of a configuration and a current at real backgrounds. A real background is not a
function on the closed set $X$ of Definition~\ref{4.9:def-cube-gauge-condition} alone.
Ba{\l}aban's minimal configuration $U_k(V)$ lives on the whole
lattice in \cite{B-CMP109} and \cite{B-CMP116}, and on the support of its determining set in
\cite{B-CMP119}, \cite{B-CMP122a} and \cite{B-CMP122b}. Its curvature and its current are both
bounded there \cite[(2), (8)]{B-CMP102b}.

What cannot be done is to derive the gradient bound from data of first order on $X$. For a
field whose curvature and current are both bounded on a collar round a cube, in the sense of
\cite[(1.7)--(1.9)]{B-CMP99a}, the gauge with the gradient bound is a theorem of Ba{\l}aban
\cite[Proposition~6]{B-CMP99a}. Whether a cube $\square\subset X$ leaves room for a collar inside
$X$ is a property relative to $X$ only: the collar is taken in the ambient lattice, not in $X$.
The following condition states the two bounds in the units of a level.

\begin{definition}[Regularity on a collar and change of level]\label{4.9:def-collar-regularity}
Let $R$ be a closed set, let $j$ be a level and $a>0$. Let $U$ be a $G$-valued configuration on
the bonds meeting $R^{+1}$. Lengths are measured in the units of level $j$, so that the lattice
spacing is $\xi_j=L^{-j}$. We say that $U$ satisfies $\mathrm{Reg}_j[R;a]$ if
\begin{equation}\label{4.9:eq-collar-regularity-a}
\begin{aligned}
|\partial U(p)-1| &< a\xi_j^2 &&\text{for the plaquettes } p \text{ meeting } R,\\
|(D^{\xi_j *}_U\partial U)(b)| &< a\xi_j^2 &&\text{for the bonds } b \text{ meeting } R.
\end{aligned}
\end{equation}
This is the pair of conditions \cite[(1.7)--(1.9)]{B-CMP99a}, which coincides with
\cite[(2)]{B-CMP102b}, at index $j$, written in level-$j$ units. The covariant divergence
$D^{\xi *}_U$ carries one factor $\xi^{-1}$. In the units of a level $k\ge j$, put $\eta=L^{-k}$
and $\ell_j=L^j\eta$. Then $\eta/\ell_j=\xi_j$ and $D^{\xi_j*}_U=\ell_j D^{\eta *}_U$, so the
two bounds read $a\eta^2\ell_j^{-2}$ and $a\eta^2\ell_j^{-3}$ respectively.

For matrices put $\operatorname{Im}W=(W-W^*)/(2i)$ and $\operatorname{Re}W=(W+W^*)/2$. Let
$\operatorname{Tr}$ denote the ordinary matrix trace, so $\operatorname{Tr}=N\operatorname{tr}$.
Let $\pi$ denote the projection $W\mapsto W-N^{-1}\operatorname{Tr}(W)\,1$ onto traceless
matrices. There is a constant $c_G$, depending only on $N$, with the following property. For $U$
as above, the first two lines of the display below hold without further hypothesis; its third
line holds whenever the following three conditions are met:
\begin{itemize}
\item $|\partial U(p)-1|<a\xi_j^2$ for the plaquettes $p$ meeting $R$;
\item $|J(b)|<a$ for the bonds $b$ meeting $R$, where
$J=D^{\xi_j*}_U\xi_j^{-2}\pi\operatorname{Im}\partial U$ is the current of
\cite[(1.2), (1.8)]{B-CMP109};
\item $c_G a\xi_j^2\le\frac12$.
\end{itemize}
\begin{equation}\label{4.9:eq-collar-regularity-b}
\begin{gathered}
\mathrm{Reg}_n[R;a]\ \Longrightarrow\ \mathrm{Reg}_j[R;L^{-2(n-j)}a]
\qquad (j\le n),\\
\mathrm{Reg}_{j-1}[R;a]\ \Longrightarrow\ \mathrm{Reg}_j[R;L^{3}a]\qquad (j\ge 1),\\
|(D^{\xi_j *}_U\partial U)(b)|<\bigl(1+2(d-1)a\bigr)a\xi_j^2
\quad\text{for the bonds } b \text{ meeting } R.
\end{gathered}
\end{equation}
The third conversion is not needed for minimal configurations.
\end{definition}

\begin{proof}
A change of level is a change of the unit of length. The configuration $U$, the set $R$, and the
plaquettes and bonds meeting $R$ are unchanged. In level-$i$ units the spacing is $\xi_i=L^{-i}$.
The plaquette variable $\partial U$ does not depend on the unit, and $D^{\xi*}_U$ is $\xi^{-1}$
times an operator independent of the unit. Hence, for any two levels $j,n$,
\begin{equation*}
\xi_n=L^{-(n-j)}\xi_j,\qquad
D^{\xi_j*}_U=\frac{\xi_n}{\xi_j}D^{\xi_n*}_U=L^{-(n-j)}D^{\xi_n*}_U .
\end{equation*}

\emph{First line.} Let $j\le n$ and assume $\mathrm{Reg}_n[R;a]$. For a plaquette $p$ meeting
$R$ we have $|\partial U(p)-1|<a\xi_n^2=L^{-2(n-j)}a\xi_j^2$. For a bond $b$ meeting $R$,
\begin{equation*}
|(D^{\xi_j*}_U\partial U)(b)|=L^{-(n-j)}|(D^{\xi_n*}_U\partial U)(b)|
<L^{-(n-j)}a\xi_n^2=L^{-3(n-j)}a\xi_j^2 .
\end{equation*}
Since $L\ge1$ and $n-j\ge0$, the right-hand side is at most $L^{-2(n-j)}a\xi_j^2$. Thus the current
bound gains $L^{-3(n-j)}$, more than required. This is the remark following
\cite[(1.7)--(1.9)]{B-CMP99a}: if the conditions hold for one index, they hold for all smaller
indices.

\emph{Second line.} Read the same two identities upward, with $n=j-1$: $\xi_{j-1}=L\xi_j$ and
$D^{\xi_j*}_U=LD^{\xi_{j-1}*}_U$. Assume $\mathrm{Reg}_{j-1}[R;a]$. Then, for the plaquettes
$p$ and the bonds $b$ meeting $R$,
\begin{align*}
|\partial U(p)-1|&<a\xi_{j-1}^2=L^2a\xi_j^2\le L^3a\xi_j^2,\\
|(D^{\xi_j*}_U\partial U)(b)|&<L\,a\xi_{j-1}^2=L^3a\xi_j^2 .
\end{align*}
This is the factor $L^3\alpha_0$ of \cite[(1.132)]{B-CMP99a}.

\emph{Third line.} Abbreviate $\xi=\xi_j$. Fix a bond $b$ meeting $R$. Every plaquette $p$
containing $b$ meets $R$, so $W=\partial U(p)$ satisfies $|W-1|<a\xi^2$. Since $W$ is unitary,
$W^{-1}=W^*$, and
\begin{equation*}
W-1=i\operatorname{Im}W+(\operatorname{Re}W-1),\qquad
\operatorname{Re}W-1=-\tfrac12(W-1)(W^{-1}-1).
\end{equation*}
The second identity holds because the right-hand side equals $-\frac12(2-W-W^*)$. Here $|\cdot|$
is the matrix norm in which the bounds \eqref{4.9:eq-collar-regularity-a} are posed; it is
submultiplicative and invariant under taking adjoints, and its value $|1|$ at the identity is a
number depending on $N$ only. Moreover $W^{-1}-1=(W-1)^*$. Hence
\begin{equation*}
|\operatorname{Re}W-1|\le\tfrac12|W-1|^2<\tfrac12(a\xi^2)^2 .
\end{equation*}

Next, $(1-\pi)\operatorname{Im}W=N^{-1}\operatorname{Tr}(\operatorname{Im}W)\,1$. Let
$e^{i\theta_r}$, $r=1,\dots,N$, be the eigenvalues of $W$, each argument $\theta_r$ chosen of
smallest absolute value; those of $\operatorname{Im}W$ are $\sin\theta_r$. Each
$|e^{i\theta_r}-1|$ is at most the operator norm of $W-1$. All norms on $N\times N$ matrices are
equivalent, so each $|\theta_r|$ is bounded by a constant depending on $N$ times $a\xi^2$. The
constant $c_G$ of the statement is fixed, depending on $N$ only, as the larger of two numbers. The
first is so large that $c_Ga\xi^2\le\frac12$ forces each $|\theta_r|$ to be less than $1/N$ of a
full turn. Since $\det W=1$, the sum $\sum_r\theta_r$ is an integer multiple of a full turn; its
absolute value is less than a full turn, and therefore $\sum_r\theta_r=0$. Consequently
\begin{equation*}
|\operatorname{Tr}\operatorname{Im}W|=\Bigl|\sum_r(\sin\theta_r-\theta_r)\Bigr|
\le c_G'(a\xi^2)^3 ,
\end{equation*}
with $c_G'$ depending on $N$ only, by $|\sin\theta-\theta|\le|\theta|^3/6$.

The covariant divergence annihilates the constant function $1$. It also satisfies
$|(D^{\xi*}_UF)(b)|\le2(d-1)\xi^{-1}\sup|F|$, the supremum being over the plaquettes containing
$b$. Since $\xi^2J=D^{\xi*}_U\pi\operatorname{Im}\partial U$,
\begin{equation*}
D^{\xi*}_U\partial U-i\xi^2J
=D^{\xi*}_U\bigl[(\operatorname{Re}\partial U-1)+i(1-\pi)\operatorname{Im}\partial U\bigr].
\end{equation*}
The bracket is bounded by $\frac12(a\xi^2)^2+N^{-1}|1|\,c_G'(a\xi^2)^3$. The second number
defining $c_G$ is so large that $c_Ga\xi^2\le\frac12$ gives $N^{-1}|1|\,c_G'a\xi^2\le\frac12$;
then the bracket is
at most $(a\xi^2)^2$. Since $\xi=L^{-j}\le1$, at $b$ we obtain
\begin{equation*}
|D^{\xi*}_U\partial U-i\xi^2J|\le2(d-1)\xi^{-1}(a\xi^2)^2\le2(d-1)a\cdot a\xi^2 .
\end{equation*}
Finally $|\xi^2J(b)|<a\xi^2$, so
\begin{equation*}
|(D^{\xi*}_U\partial U)(b)|<a\xi^2+2(d-1)a\cdot a\xi^2=\bigl(1+2(d-1)a\bigr)a\xi^2 .
\end{equation*}
\end{proof}

\begin{lemma}[Cube gauge from regularity on a collar]\label{4.9:lem-seed-cube-gauge}
Let $j\ge 0$ and $a>0$. Let $\square^{a}$ be a closed cube which is a union of cubes of the grid
of mesh $R_1M_1L^{j}$ lattice steps, this grid being compatible with the block structure. Let the
side of $\square^{a}$ be $M^{\ast}L^{j}$ lattice steps, where $M^{\ast}$ is a multiple of $R_1M_1$. Put
\begin{equation*}
\mathcal{N}(\square^{a}):=(\square^{a})^{+4R_1M_1L^{j}},
\end{equation*}
and assume that $\mathcal{N}(\square^{a})$ is embedded in the torus. Let $U$ be a $G$-valued field on
the bonds meeting $\mathcal{N}(\square^{a})^{+1}$. If $\mathrm{Reg}_j[\mathcal{N}(\square^{a});a]$
holds (Definition~\ref{4.9:def-collar-regularity}) and $7dL^{2}M^{\ast}a\le c_1$, then
\begin{equation}\label{4.9:eq-seed-cube-gauge-1}
\mathsf{g}_j\bigl(U;\square^{a};7dL^{2}B_1M^{\ast}a\bigr)
\end{equation}
holds, in the sense of Definition~\ref{4.9:def-cube-gauge-condition}. Here $B_1$ and $c_1$ are the
constants of Theorems~2 and~4 of \cite{B-CMP99a}. They depend on $d$ and $L$ only.
\end{lemma}

\begin{proof}
The proof applies the statement of Ba{\l}aban's cube gauge for regular fields
(Proposition~\ref{4.9:prop-cube-gauge-regular}, which is \cite[Proposition~6]{B-CMP99a}) to an
auxiliary nest of domains built around $\square^{a}$.

\emph{Units.} We change the scale so that level $j$ becomes the top level $k$ of \cite{B-CMP99a}.
That is, we put $k:=j$, let $\eta:=\xi_j$ be the fine lattice spacing, and write $\ell_i:=L^{i}\eta$
for $0\le i\le k$. Then $\ell_k=1$. In these units the grid of the statement has mesh
$R_1M_1=R_1M_1\ell_k$, the side of $\square^{a}$ is $M^{\ast}$, and $\mathcal{N}(\square^{a})$ is the
closed neighbourhood of $\square^{a}$ of sup-distance $4R_1M_1$.

\emph{The enlarged cube.} Let $\widetilde{\square}\supset\square^{a}$ be the cube that
\cite[pp.~98--99]{B-CMP99a} constructs from the cube there denoted $\square$, which is our
$\square^{a}$. Its boundary lies at distance $2R_1M_1$ from $\square^{a}$. Hence
$\widetilde{\square}\subset(\square^{a})^{+2R_1M_1}$.

\emph{The auxiliary nest.} Put $\Omega_k:=\widetilde{\square}$. Downward, for $i=k-1,\dots,0$, let
$\Omega_i$ be the smallest union of cubes of side $M_1\ell_i$ of the block grid that contains
$(\Omega_{i+1})^{+(R_1+1)M_1\ell_i}$. Then, for $i<k$, $\Omega_i$ is a union of cubes of side
$M_1\ell_i$. Moreover $\Omega_i\supset(\Omega_{i+1})^{+(R_1+1)M_1\ell_i}$, so that
\begin{equation*}
\ell_i^{-1}\operatorname{dist}(\Omega_i^{c},\Omega_{i+1})\ge(R_1+1)M_1>R_1M_1 .
\end{equation*}
Thus the conditions \cite[(1.3)--(1.4)]{B-CMP99a} on a nest of domains hold, with Ba{\l}aban's
separation parameter $R$ equal to $R_1$.

\emph{The nest stays inside the collar.} Every cube of side $M_1\ell_i$ that meets a set lies in the
$M_1\ell_i$-neighbourhood of that set. The passage from $\Omega_{i+1}$ to $\Omega_i$ therefore
enlarges by at most $(R_1+2)M_1\ell_i$. Consequently $\Omega_0\subset\widetilde{\square}^{+r}$ with
\begin{equation*}
r\le\sum_{i=0}^{k-1}(R_1+2)M_1\ell_i\le\frac{(R_1+2)M_1}{L-1}\le R_1M_1 .
\end{equation*}
The middle inequality holds because $\ell_i=L^{i-k}$, so that
$\sum_{i<k}\ell_i\le\sum_{n\ge1}L^{-n}=(L-1)^{-1}$. The last inequality uses $L\ge4$ and
$R_1\ge1$, which give
\begin{equation*}
R_1+2\le 3R_1\le(L-1)R_1 .
\end{equation*}
Now consider $\Omega_0$ together with the bonds and plaquettes which \cite{B-CMP99a} attaches to it,
and the bonds and plaquettes read by the covariant divergence $D^{\eta *}_{U}\partial U$ at bonds of
$\Omega_0$; the latter reach one further lattice step. All of these lie in
\begin{equation*}
(\square^{a})^{+(2R_1M_1+R_1M_1+2\eta)}\subset\mathcal{N}(\square^{a}).
\end{equation*}

\emph{The configuration.} Let $U_0$ be the configuration on all bonds of the fine lattice on the
torus defined as follows: $U_0:=U$ on the bonds meeting $\mathcal{N}(\square^{a})$, and $U_0:=1$ on
all other bonds.

By the first of the conversions \eqref{4.9:eq-collar-regularity-b} in
Definition~\ref{4.9:def-collar-regularity}, the hypothesis $\mathrm{Reg}_k[\mathcal{N}(\square^{a});a]$
implies $\mathrm{Reg}_i[\mathcal{N}(\square^{a});L^{-2(k-i)}a]$ for every $i\le k$. Since
$L^{-2(k-i)}\le1$ and the inequalities in \eqref{4.9:eq-collar-regularity-a} are strict, this gives
$\mathrm{Reg}_i[\mathcal{N}(\square^{a});a]$ for every $i\le k$.

By Definition~\ref{4.9:def-collar-regularity}, $\mathrm{Reg}_i$ is the pair of conditions
\cite[(1.7), (1.9)]{B-CMP99a} at index $i$, written in level-$i$ units. Moreover
$\Omega_i\subset\Omega_0$, and every bond and plaquette read on $\Omega_0$ lies in
$\mathcal{N}(\square^{a})$, where $U_0=U$. Hence \cite[(1.7), (1.9)]{B-CMP99a} hold for $U_0$ at every
index $i\le k$ on $\Omega_i$, with constant $a$. That is,
\begin{equation*}
U_0\in\mathfrak{U}_k(\{\Omega_i\},a).
\end{equation*}
The space $\mathfrak{U}_k(\{\Omega_i\},a)$ constrains $U_0$ on $\Omega_0$ only, so the values $1$
assigned outside $\mathcal{N}(\square^{a})$ play no role.

\emph{The setting of the cited proposition.} The following hold.
\begin{itemize}
\item $\square^{a}\subset\Omega_k$, and there is no domain $\Omega_{k+1}$.
\item $\square^{a}$ is a union of cubes of size $R_1M_1\ell_k$.
\item $\square^{a}$ has size $M^{\ast}$, a multiple of $R_1M_1$.
\item $\widetilde{\square}=\Omega_k\subset\Omega_{k-1}$.
\end{itemize}
This is the setting in which Proposition~\ref{4.9:prop-cube-gauge-regular} is stated, with
$j=k$, with Ba{\l}aban's $M$ equal to $M^{\ast}$ and with $\alpha_0=a$. Ba{\l}aban's auxiliary
assumption $L^{-1}RM>2dR_1M_1$ serves there only to obtain $\widetilde{\square}\subset\Omega_{k-1}$,
and here that inclusion holds by construction. Finally, the smallness condition
$7dL^{2}M\alpha_0\le c_1$ of Proposition~\ref{4.9:prop-cube-gauge-regular} is the hypothesis
$7dL^{2}M^{\ast}a\le c_1$.

\emph{Application.} Proposition~\ref{4.9:prop-cube-gauge-regular} provides the following. Let
$U_0'$ be the configuration constructed there, which is itself a gauge transform of $U_0$. There is
a gauge transformation $u$ on $\widetilde{\square}$ taking $U_0'$ to $e^{i\eta A}$ on
$\widetilde{\square}$, and the bounds \cite[(1.136)]{B-CMP99a} hold on the cube $\square_k$ of
\cite[pp.~98--99]{B-CMP99a}, which contains $\square^{a}$. At index $k$ we have $L^{k}\eta=\ell_k=1$,
so the first two of these bounds read
\begin{equation}\label{4.9:eq-seed-cube-gauge-2}
|A|<7dL^{2}B_1M^{\ast}a,\qquad|\nabla^{\eta}A|<7dL^{2}B_1M^{\ast}a\qquad\text{on }\square_k .
\end{equation}
Here $\nabla^{\eta}$ is the ordinary finite-difference gradient \cite[p.~98, (1.127)]{B-CMP99a}.

All gauge transformations in \cite{B-CMP99a} are $G$-valued, and $A$ is $\mathfrak{g}$-valued. Let
$u_{\ast}$ be the composite of the gauge transformation taking $U_0$ to $U_0'$ with $u$, restricted
to $\operatorname{st}(\square^{a})$. Then $U_0^{u_{\ast}}=e^{i\xi_jA}$ on $\operatorname{bd}(\square^{a})$,
because $\eta=\xi_j$.

On $\operatorname{bd}(\square^{a})$ we have $U_0=U$, so $U^{u_{\ast}}=e^{i\xi_jA}$ there. By
\eqref{4.9:eq-seed-cube-gauge-2}, restricted to $\operatorname{bd}(\square^{a})$ and to
$\operatorname{pp}(\square^{a})$, the two bounds required in \eqref{4.9:eq-cube-gauge-condition-a}
hold with $\rho=7dL^{2}B_1M^{\ast}a$. This is \eqref{4.9:eq-seed-cube-gauge-1}.
\end{proof}

\begin{proposition}[Real backgrounds satisfy the ambient gauge condition]
\label{4.9:prop-real-backgrounds}
Let $d=4$, and let $T$ be the torus on which the step is performed. Write $S:=c_{12}LM$,
where $c_{12}$ is the constant written $O(1)$ in the cube side $O(1)LM$ of
\cite[(1.12)]{B-CMP109}, the number $\mathsf{o}$ of Definition~\ref{4.4:def-local-data-classes}.
Let $R_1,M_1,a_1,B_3$ and
$M(\varepsilon_1)=R_1M_1a_1/\varepsilon_1$ be the constants of Ba{\l}aban's minimal-orbit
theorem (Theorem~\ref{4.9:thm-minimal-orbit}), and $B_1,c_1$ those of
Lemma~\ref{4.9:lem-seed-cube-gauge}, depending on $d$ and $L$ only.
Let $\mathsf{g}_j(U;D;\rho)$ be the cube gauge condition \eqref{4.9:eq-cube-gauge-condition-a},
and let $(\ell 1)_j$, $(\ell 2)_j$, $(\ell 2)^{\flat}_j$ be its readings
\eqref{4.9:eq-cube-gauge-condition-b}, with their constants $B_{12}$ and $B$ and the positive
integer $C$ of the cube side $CML^j$ (Definition~\ref{4.9:def-cube-gauge-condition}), taken
with $C\le 2$, as \cite[following (1.69)]{B-CMP122a} allows. Let
$\mathrm{Reg}_j[R;a]$ be the regularity hypothesis \eqref{4.9:eq-collar-regularity-a}. Let
$m_0\in\{0,1\}$ be the lowest index of the nest of part~(c) below, and $\kappa_G$ the group
constant of the comb-gauge hypothesis of part~(c)(ii$_0$). Each of the following floors is
used only in the assertion that names it:
\begin{equation}\label{4.9:eq-real-backgrounds-b}
\begin{gathered}
B_{12}\ge 2B_3,\qquad B_{12}\ge 14dL^2B_1,\\
B\ge 2B_3,\qquad B\ge 14dL^2B_1,\qquad B\ge 14dL^5B_1,\qquad
B\ge \pi L^{m_0}(d-1)\kappa_G .
\end{gathered}
\end{equation}

\smallskip\noindent
\textup{(a)} \emph{Minimal configurations.} Let $k\ge1$ and let
$U=U_k(V)\in\mathfrak{U}_k(\varepsilon_0)$ be a configuration in the minimal orbit, with
$\mathfrak{U}_k(\varepsilon_0)$ the space of \cite[(1.2)]{B-CMP109}; this is the situation of
\cite[\S\S1--5]{B-CMP109} and \cite[\S\S1--2]{B-CMP116}, in which the nest is trivial.
Assume the ceilings
$2\varepsilon_0\le a_1$ and $2S\varepsilon_0\le R_1M_1a_1$; the floors $S\ge 4R_1M_1$ and
$M\ge 4R_1M_1$ (both of which hold when $c_{12}L\ge2$ and $M\ge 2LR_1M_1$); and
\begin{equation}\label{4.9:eq-real-backgrounds-a}
\text{the side of } T\text{, measured in level-}k\text{ units, is at least } 3S .
\end{equation}
Then for every closed cube $\square\subset T$ of side $s_\square L^k$ with
$0<s_\square\le S$,
\begin{equation}\label{4.9:eq-real-backgrounds-1}
\mathsf{g}_k\bigl(U;\square;2B_3(\tfrac12 s_\square+2R_1M_1)\varepsilon_0\bigr).
\end{equation}
Hence, for every closed set $X\subset T$, every $1\le j\le k$ and $\mu:=L^{-(k-j)}$, the
restriction $U|_{X}$ obeys $(\ell 2)_j(X;\mu\varepsilon_0)$ and $(\ell 1)_j(X;\mu\varepsilon_0)$
as soon as $B_{12}\ge 2B_3$, and $(\ell 2)^{\flat}_j(X;\mu\varepsilon_0)$ as soon as
$B\ge 2B_3$ and $CM\le S$ (which holds when $c_{12}L\ge 2$, since $C\le 2$).

\smallskip\noindent
\textup{(b)} \emph{Real fields regular on a collar.} Let $U$ be $G$-valued, $j\ge0$, $X$
closed, $s\ge 3R_1M_1$, $a>0$ with $14dL^2sa\le c_1$, all sets below being embedded in $T$.
\begin{itemize}
\item[(in)] If $\mathrm{Reg}_j[X^{+7R_1M_1L^j};a]$, then for every closed cube
$\square\subset X$ of side at most $sL^j$
\begin{equation*}
\mathsf{g}_j(U;\square;14dL^2B_1sa).
\end{equation*}
\item[(amb)] If $\mathrm{Reg}_j[X^{+(s+7R_1M_1)L^j};a]$, then for every closed cube
$\square$ of side at most $sL^j$ in $T$
\begin{equation*}
\mathsf{g}_j(U;\square\cap X;14dL^2B_1sa).
\end{equation*}
\end{itemize}
So, with $s:=S$ respectively $s:=CM$ (the hypotheses of~(b) holding at that value of $s$):
the floor $B_{12}\ge 14dL^2B_1$ gives $(\ell 1)_j(X;a)$ from~(in) and $(\ell 2)_j(X;a)$
from~(amb); the floor $B\ge 14dL^2B_1$ gives $(\ell 2)^{\flat}_j(X;a)$ from~(amb) at
$s=CM$.

\smallskip\noindent
\textup{(c)} \emph{Instances.} Let $\{\Omega_n\}_{n\le k}$ be a nest, and let $U$ be real
with $\mathrm{Reg}_n[\Omega_n\setminus\Omega_{n+1};a_n]$ for $n<k$ and
$\mathrm{Reg}_k[\Omega_k;a_k]$. Put $\bar a_j:=\sup_{j\le n\le k}L^{-2(n-j)}a_n$; then
$\mathrm{Reg}_j[\Omega_j;\bar a_j]$. In~(i) and in the first assertion of~(ii), part~(b) is
applied with the amplitude named there, and its hypotheses are assumed at the values used:
in~(i), $S\ge 3R_1M_1$, $CM\ge 3R_1M_1$, $14dL^2S\bar a_j\le c_1$ and
$14dL^2CM\bar a_j\le c_1$; in~(ii), $CM\ge 3R_1M_1$ and
\begin{equation*}
14dL^2CM\max\{L^3\bar a_{m-1},\bar a_m\}\le c_1 ;
\end{equation*}
in both, all collars are embedded in $T$. In~(i), moreover, the floor
$B_{12}\ge 14dL^2B_1$ is assumed for $(\ell 1)_j$ and $(\ell 2)_j$, and the floor
$B\ge 14dL^2B_1$ for $(\ell 2)^{\flat}_j$.
\begin{itemize}
\item[(i)] Let $j$ be an index at which $\Omega_j$ and $\Lambda_j$ are made by the operation
$\mathbf{T}$, and let $X\subset\Lambda_j$ be closed. Put
$\mathsf{m}_j:=L^{-j}\operatorname{dist}_\infty(\Lambda_j,\Omega_j^{c})$, and let $R_j$
be the number of \cite[(2.5)]{B-CMP119}. Then $\mathsf{m}_j\ge MR_j$: at such an index
\cite{B-CMP119} requires the boundaries of these domains to be at distance at least $2MR_j$,
and for $d=4$ this gives $\mathsf{m}_j\ge MR_j$ whether the distance is Euclidean or taken in
the sup norm. The reading
$(\ell 1)_j(X;\bar a_j)$ holds if $7R_1M_1\le MR_j$, and $(\ell 2)^{\flat}_j(X;\bar a_j)$
holds if $2M+7R_1M_1\le MR_j$; both conditions follow from $M\ge 2LR_1M_1$ together with
$R_j\ge L$ ($L\ge 4$ holding under the condition $L>11$ of \cite{B-CMP109}). Further,
$(\ell 2)_j(X;\bar a_j)$ holds if $S+7R_1M_1\le MR_j$, which, given $M\ge 2LR_1M_1$,
$L\ge 4$ and $c_{12}L\ge 1$, follows from
\begin{equation}\label{4.9:eq-real-backgrounds-c}
R_j\ge 2c_{12}L .
\end{equation}
By \cite[(2.5)]{B-CMP119}, \eqref{4.9:eq-real-backgrounds-c} is a ceiling on $\gamma$,
imposed after $L$; it is needed for $(\ell 2)_j$ only.
\item[(ii)] (The cubes of the class in \eqref{4.9:eq-cube-gauge-condition-c}, the thin shells
included.) Let $m$ and $m-1$ be indices of the nest, and assume
$L^{-m}\operatorname{dist}_\infty(\Omega_{m-1}^{c},\Omega_m)\ge 7R_1M_1$. Where the shell
$\Omega_{m-1}\setminus\Omega_m$ is one layer of $M$-cubes of level $m-1$ thick, as in
\cite{B-CMP122a}, this condition is
\begin{equation}\label{4.9:eq-real-backgrounds-d}
M\ge 7LR_1M_1 .
\end{equation}
Then every closed cube $\square\subset\Omega_m$ of side $CML^m$ obeys
\begin{equation}\label{4.9:eq-real-backgrounds-2}
\mathsf{g}_m\bigl(U;\square;14dL^5B_1CM\max\{\bar a_{m-1},L^{-3}\bar a_m\}\bigr).
\end{equation}
Consequently, given amplitudes $a^{\mathrm{tgt}}_n>0$, if the distance hypothesis above and
the bounds $\bar a_{m-1}\le a^{\mathrm{tgt}}_m$, $\bar a_m\le a^{\mathrm{tgt}}_m$ hold at every
index $m>m_0$ of the nest, a real background obeys
\eqref{4.9:eq-cube-gauge-condition-c} for the nest $\{\Omega_n\}$, every set $Y$ and the
amplitudes $(a^{\mathrm{tgt}}_n)$ at every index $m>m_0$, as soon as $B\ge 14dL^5B_1$,
under the ceiling $28dL^5M\,a^{\mathrm{tgt}}_m\le c_1$.
\item[(ii$_0$)] At the lowest index $m_0$ of the nest (in \cite{B-CMP119}, $m_0=1$) there is
no $\Omega_{m_0-1}$: the large-field region begins there, no collar is available, and none
is needed. Let $a^{\mathrm{curv}}_{m_0}>0$ be such that the curvature member of
$\mathrm{Reg}_{m_0}[\Omega_{m_0};a^{\mathrm{curv}}_{m_0}]$ holds, and assume the
comb-gauge statement in the following form: for every closed cube $\square\subset\Omega_{m_0}$
of side $CML^{m_0}$, under the ceiling
\begin{equation*}
(d-1)\cdot 2M\cdot a^{\mathrm{curv}}_{m_0}\le 2\sin\bigl(\pi/(2N)\bigr),
\end{equation*}
there are $u$ on $\operatorname{st}(\square)$ and a $\mathfrak{g}$-valued $A$ on
$\operatorname{bd}(\square)$ with $U^{u}=e^{i\xi_{m_0}A}$ on $\operatorname{bd}(\square)$ and
\begin{equation*}
|A|<\tfrac{\pi}{2}(d-1)\kappa_G\,CM\,a^{\mathrm{curv}}_{m_0},
\end{equation*}
$\kappa_G$ being a constant of the group. Then every such cube obeys
\begin{equation}\label{4.9:eq-real-backgrounds-3}
\mathsf{g}_{m_0}\bigl(U;\square;\pi L^{m_0}(d-1)\kappa_G\,CM\,a^{\mathrm{curv}}_{m_0}\bigr),
\end{equation}
from the curvature bound alone; the floor $B\ge\pi L^{m_0}(d-1)\kappa_G$ then makes this
radius at most $BCMa^{\mathrm{curv}}_{m_0}$, which, after restriction to $\square\cap Y$ for
an arbitrary set $Y$, gives \eqref{4.9:eq-cube-gauge-condition-c} at index $m_0$ with
$a^{\mathrm{curv}}_{m_0}$ in place of $a_{m_0}$. This floor involves the fixed index
$m_0\in\{0,1\}$, not a depth.
\end{itemize}
\end{proposition}

\begin{proof}
\emph{Part (a).} Put $\varepsilon_1:=2\varepsilon_0$; by the first ceiling
$\varepsilon_1\le a_1$. By \cite[Proposition~2]{B-CMP98}, in the form in which
\cite[(1.2)]{B-CMP109} uses it, the membership $U\in\mathfrak{U}_k(\varepsilon_0)$ implies
$|V(\partial p')-1|<2\varepsilon_0$ for the plaquettes $p'$ of the block lattice; hence $V$
satisfies condition \cite[(7)]{B-CMP102b} of Ba{\l}aban's minimal-orbit theorem
(Theorem~\ref{4.9:thm-minimal-orbit}) with $\varepsilon_1=2\varepsilon_0$.

Consider the trivial nest $\Omega_i=T$ for $i\le k$, which \cite{B-CMP99a} admits. For it,
$\Lambda_k$ is the whole level-$k$ lattice and $B^k(\Lambda_k)=T$, and the class of cubes of
Theorem~\ref{4.9:thm-minimal-orbit} at index $k$ consists of all cubes $\square^a$ which are
unions of big blocks, that is of cubes of the grid of mesh $M_1$ in level-$k$ units (in which
$L^k\eta=1$), of side $2M'$ with $M'$ a multiple of $R_1M_1$, and whose concentric cube
$\widetilde{\square}$ of side $2M'+4R_1M_1$ embeds in $T$.

Write $\square=\prod_\mu I_\mu$ with $|I_\mu|=s_\square$. The smallest union of grid
intervals containing $I_\mu$ has length at most $s_\square+2M_1$. Put
\begin{equation*}
M':=R_1M_1\Bigl\lceil\frac{s_\square+2M_1}{2R_1M_1}\Bigr\rceil .
\end{equation*}
Then $2M'\ge s_\square+2M_1$ and $2M'$ is a multiple of $M_1$, so there is a grid interval
$I^a_\mu\supset I_\mu$ of length $2M'$. Since $M_1\le R_1M_1$, since $s_\square\le S$ and
$4R_1M_1\le S$ by the first floor, and since the second ceiling says
$S\le R_1M_1a_1/(2\varepsilon_0)=M(\varepsilon_1)$,
\begin{equation*}
M'<\tfrac12(s_\square+2M_1)+R_1M_1\le\tfrac12 s_\square+2R_1M_1\le S\le M(\varepsilon_1).
\end{equation*}
The cube $\widetilde{\square}$ concentric with $\square^a:=\prod_\mu I^a_\mu$ has side
\begin{equation*}
2M'+4R_1M_1<s_\square+8R_1M_1\le 3S,
\end{equation*}
using again $s_\square\le S$ and $8R_1M_1\le 2S$; by \eqref{4.9:eq-real-backgrounds-a} it
embeds in $T$. Thus $\square^a$ belongs to the class at index $k$, with $M'\le M(\varepsilon_1)$.

In the notation of \cite{B-CMP109}, $U=U_k(V)$ denotes a configuration in the minimal orbit
of $V$. The regularity property \cite[(9)]{B-CMP102b} of
Theorem~\ref{4.9:thm-minimal-orbit}, applied to $\square^a$, gives a
$G$-valued gauge transformation on a neighbourhood of $\square^a$ taking $U$ to $e^{i\eta A}$
on $\square^a$, with
\begin{equation*}
|A|,\ |\nabla^{\eta}A|<B_3M'\varepsilon_1<2B_3\bigl(\tfrac12 s_\square+2R_1M_1\bigr)
\varepsilon_0 ,
\end{equation*}
where the last inequality is $M'<\frac12 s_\square+2R_1M_1$ and $\varepsilon_1=2\varepsilon_0$.
This is $\mathsf{g}_k(U;\square^a;2B_3(\frac12 s_\square+2R_1M_1)\varepsilon_0)$, and
Lemma~\ref{4.9:lem-restriction-monotone} restricts it to $\square\subset\square^a$, which is
\eqref{4.9:eq-real-backgrounds-1}.

For the consequences: by the floors, $\frac12 S+2R_1M_1\le S$ and
$\frac12 CM+2R_1M_1\le CM$ (as $CM\ge M\ge 4R_1M_1$). A cube in $T$ of side $SL^j$
(respectively $CML^j$) has side at most $SL^k$ (respectively $CML^k$, which is at most
$SL^k$ as $CM\le S$), since $j\le k$; so
\eqref{4.9:eq-real-backgrounds-1} and the monotonicity in the radius of
Lemma~\ref{4.9:lem-restriction-monotone} give $\mathsf{g}_k(U;\square;2B_3S\varepsilon_0)$
(respectively $\mathsf{g}_k(U;\square;2B_3CM\varepsilon_0)$) for every closed cube of side at
most $SL^j$ (respectively $CML^j$). The change of level
(Lemma~\ref{4.9:lem-level-change}) gives the same at level $j$ with radius $\mu$ times the
bound, and Lemma~\ref{4.9:lem-restriction-monotone} restricts to $\square\cap X$ (for the
ambient readings) or to cubes $\square\subset X$ (for $(\ell 1)_j$). Under
$B_{12}\ge 2B_3$ the radius $2B_3S\mu\varepsilon_0$ is at most $SB_{12}\mu\varepsilon_0$, which
gives $(\ell 2)_j(X;\mu\varepsilon_0)$ and $(\ell 1)_j(X;\mu\varepsilon_0)$; under
$B\ge 2B_3$ the radius $2B_3CM\mu\varepsilon_0$ is at most $BCM\mu\varepsilon_0$, which gives
$(\ell 2)^{\flat}_j(X;\mu\varepsilon_0)$.

\emph{Part (b).} Let $\square=\prod_\mu I_\mu$ have side $s_\square L^j$, $s_\square\le s$. Put
\begin{equation*}
M^{\ast}:=R_1M_1\Bigl\lceil\frac{s_\square+2R_1M_1}{R_1M_1}\Bigr\rceil
<s_\square+3R_1M_1\le 2s,
\end{equation*}
the last step by $s_\square\le s$ and $3R_1M_1\le s$. As in part~(a), coordinate by
coordinate: the smallest union of intervals of the grid of mesh $R_1M_1L^j$ containing
$I_\mu$ has length at most $(s_\square+2R_1M_1)L^j\le M^{\ast}L^j$ and lies in
$I_\mu^{+R_1M_1L^j}$; the difference between $M^{\ast}L^j$ and its length is a multiple of
$R_1M_1L^j$ smaller than $3R_1M_1L^j$, hence at most $2R_1M_1L^j$, and adding that many grid
intervals on one side gives a grid interval of length $M^{\ast}L^j$ containing $I_\mu$ and
contained in $I_\mu^{+3R_1M_1L^j}$. The product is a cube $\square^a\supset\square$ of the
$R_1M_1L^j$-grid, of side $M^{\ast}L^j$, with $\square^a\subset\square^{+3R_1M_1L^j}$, and
therefore
\begin{equation*}
\mathcal{N}(\square^a)=(\square^a)^{+4R_1M_1L^j}\subset\square^{+7R_1M_1L^j}.
\end{equation*}
In case (in), $\square\subset X$, so $\square^{+7R_1M_1L^j}\subset X^{+7R_1M_1L^j}$. In
case (amb), either $\operatorname{st}(\square\cap X)=\emptyset$, and the clause is void, or
$\square$ meets $X$; since $\square$ has side at most $sL^j$, then $\square\subset X^{+sL^j}$
and $\square^{+7R_1M_1L^j}\subset X^{+(s+7R_1M_1)L^j}$. In both cases the hypothesis gives
$\mathrm{Reg}_j[\mathcal{N}(\square^a);a]$, the conditions of
\eqref{4.9:eq-collar-regularity-a} on a set being conditions on the plaquettes and bonds
meeting it. Moreover $7dL^2M^{\ast}a\le 14dL^2sa\le c_1$. The seed lemma
(Lemma~\ref{4.9:lem-seed-cube-gauge}) gives
$\mathsf{g}_j(U;\square^a;7dL^2B_1M^{\ast}a)$, and $7dL^2B_1M^{\ast}a\le 14dL^2B_1sa$;
Lemma~\ref{4.9:lem-restriction-monotone} then restricts to $\square$ in case (in) and to
$\square\cap X$ in case (amb).

For the consequences: with $s:=S$, (in) gives
$\mathsf{g}_j(U;\square;14dL^2B_1Sa)$ for every closed cube $\square\subset X$ of side at
most $SL^j$, and (amb) gives $\mathsf{g}_j(U;\square\cap X;14dL^2B_1Sa)$ for every closed
cube of side at most $SL^j$; under $B_{12}\ge 14dL^2B_1$ these radii are at most $SB_{12}a$,
and Lemma~\ref{4.9:lem-restriction-monotone} yields $(\ell 1)_j(X;a)$ and $(\ell 2)_j(X;a)$.
With $s:=CM$, (amb) gives radius $14dL^2B_1CMa\le BCMa$ under $B\ge 14dL^2B_1$, whence
$(\ell 2)^{\flat}_j(X;a)$.

\emph{Part (c).} The set $\Omega_j$ is the union of the layers $\Omega_n\setminus\Omega_{n+1}$,
$j\le n<k$, and of $\Omega_k$. On the layer of index $n$ the hypothesis
$\mathrm{Reg}_n[\,\cdot\,;a_n]$ and the first conversion in
\eqref{4.9:eq-collar-regularity-b} give $\mathrm{Reg}_j[\,\cdot\,;L^{-2(n-j)}a_n]$, and
$L^{-2(n-j)}a_n\le\bar a_j$; the same holds on $\Omega_k$ with $n=k$. Since every plaquette
and every bond meeting $\Omega_j$ meets one of these sets, $\mathrm{Reg}_j[\Omega_j;\bar a_j]$
holds.

(i) Since $X\subset\Lambda_j$, one has $X^{+rL^j}\subset\Omega_j$ for $r\le\mathsf{m}_j$, and
$\mathsf{m}_j\ge MR_j$. Part~(b) is applied with $a:=\bar a_j$: for $(\ell 1)_j$ through (in)
and the floor $B_{12}\ge 14dL^2B_1$, with the collar $X^{+7R_1M_1L^j}$, which lies in
$\Omega_j$ when $7R_1M_1\le MR_j$; for $(\ell 2)^{\flat}_j$ through (amb) at $s:=CM$ and the
floor $B\ge 14dL^2B_1$, with the collar $X^{+(CM+7R_1M_1)L^j}$, which lies in
$\Omega_j$ when $2M+7R_1M_1\le MR_j$, since $CM\le 2M$; for
$(\ell 2)_j$ through (amb) at $s:=S$ and the floor $B_{12}\ge 14dL^2B_1$, with the collar
$X^{+(S+7R_1M_1)L^j}$, which lies in $\Omega_j$
when $S+7R_1M_1\le MR_j$. In each case $\mathrm{Reg}_j[\Omega_j;\bar a_j]$ gives the required
regularity on the collar. The condition $M\ge 2LR_1M_1$ gives $M\ge 8R_1M_1$, as $L\ge4$
(which lies inside the condition $L>11$ of \cite{B-CMP109}); hence
\begin{equation*}
7R_1M_1\le M\le MR_j,\qquad 2M+7R_1M_1\le 3M\le MR_j ,
\end{equation*}
the latter because $R_j\ge L\ge4$. Further $S+7R_1M_1\le 2S$, as $7R_1M_1\le M$ by the above
and $M\le S$ by $c_{12}L\ge1$; and $2S=2c_{12}LM\le MR_j$ under
\eqref{4.9:eq-real-backgrounds-c}. Finally,
$R_j$ is the least power of $L$ with $R_j\ge(\log g_j^{-2})^r$, with the exponent $r$ of
\cite[(2.5)]{B-CMP119}; it exceeds $1$ for $g_j$ small, and hence $R_j\ge L$.

(ii) Let $\square\subset\Omega_m$ have side $CML^m$. By the distance hypothesis,
$\square^{+7R_1M_1L^m}\subset\Omega_{m-1}$. There $\mathrm{Reg}_{m-1}[\,\cdot\,;\bar a_{m-1}]$
holds, hence $\mathrm{Reg}_m[\,\cdot\,;L^3\bar a_{m-1}]$ by the second conversion in
\eqref{4.9:eq-collar-regularity-b}, and on $\Omega_m$ even
$\mathrm{Reg}_m[\,\cdot\,;\bar a_m]$. So
$\mathrm{Reg}_m[\square^{+7R_1M_1L^m};a]$ with $a:=\max\{L^3\bar a_{m-1},\bar a_m\}$.
Part~(b)(in) at level $m$, with $X:=\square$ and $s:=CM\ge 3R_1M_1$, gives
$\mathsf{g}_m(U;\square;14dL^2B_1CMa)$, and
\begin{equation*}
14dL^2B_1CM\max\{L^3\bar a_{m-1},\bar a_m\}
=14dL^5B_1CM\max\{\bar a_{m-1},L^{-3}\bar a_m\},
\end{equation*}
which is \eqref{4.9:eq-real-backgrounds-2}. For an arbitrary $Y$,
Lemma~\ref{4.9:lem-restriction-monotone} restricts to $\square\cap Y$. If
$\bar a_{m-1}\le a^{\mathrm{tgt}}_m$ and $\bar a_m\le a^{\mathrm{tgt}}_m$, the maximum is at
most $a^{\mathrm{tgt}}_m$, and under $B\ge 14dL^5B_1$ the radius is at most
$BCMa^{\mathrm{tgt}}_m$, which gives \eqref{4.9:eq-cube-gauge-condition-c}.

(ii$_0$) The comb-gauge statement on the sites of $\square$ gives a $\mathfrak{g}$-valued $A$
with $|A|<\frac{\pi}{2}(d-1)\kappa_G\,CM\,a^{\mathrm{curv}}_{m_0}$. Since
$\xi_{m_0}=L^{-m_0}$,
\begin{equation*}
\xi_{m_0}^{-1}|A(b+e_\nu)-A(b)|\le 2L^{m_0}\sup|A|
<\pi L^{m_0}(d-1)\kappa_G\,CM\,a^{\mathrm{curv}}_{m_0},
\end{equation*}
and also $|A|$ is below the right-hand side, as $L^{m_0}\ge1$: at a fixed finite level the
gradient member is the trivial bound. This is \eqref{4.9:eq-real-backgrounds-3}, and under
$B\ge\pi L^{m_0}(d-1)\kappa_G$ its radius is at most $BCMa^{\mathrm{curv}}_{m_0}$. For an
arbitrary $Y$, Lemma~\ref{4.9:lem-restriction-monotone} restricts to $\square\cap Y$, which
gives \eqref{4.9:eq-cube-gauge-condition-c} at index $m_0$ with $a^{\mathrm{curv}}_{m_0}$ in
place of $a_{m_0}$.
\end{proof}

\begin{corollary}[Gauge condition for the printed pairs]\label{4.9:cor-printed-pairs}
Assume $B_{12}\ge 2B_3$ and the hypotheses of Proposition~\ref{4.9:prop-real-backgrounds}.
Consider the small pairs, and the pairs at real backgrounds, whose membership in
Ba{\l}aban's class is established in \cite[\S\S3--5]{B-CMP109} and
\cite[\S\S1--2]{B-CMP116}. They satisfy the gauge condition \cite[(1.12)]{B-CMP109}
in the reading $(\ell 1)$ of \eqref{4.9:eq-cube-gauge-condition-b} at their own
curvature constant. Explicitly, writing $S=c_{12}LM$, the condition holds on every cube
whose side is at most $S$ in the units of the level, with the single bound
\begin{equation*}
c_{12}LMB_{12}\,\alpha_0',
\end{equation*}
where $\alpha_0'$ is the constant of their membership. The constant $\alpha_0'$ is as
follows.
\begin{itemize}
\item For the small pairs $(e^{i\eta\mathsf{Z}},J(e^{i\eta\mathsf{Z}}))$, with $\mathsf{Z}$
a $\mathfrak{g}^{c}$-valued field, and for the unit pair, whose unitary factor is $U:=1$,
the condition is void: it holds at every value of $\alpha_0'$.
\item For the real backgrounds $U_k(V)\in\mathfrak{U}_k(\varepsilon_0)$, one has
$\alpha_0'=\varepsilon_0$ at their own level $k$. At a level $j$ with $1\le j\le k$ one
has $\alpha_0'=L^{-(k-j)}\varepsilon_0$, with a single factor $\mu=L^{-(k-j)}$.
\end{itemize}
The step $k\to k+1$ reproduces this statement: Proposition~\ref{4.9:prop-real-backgrounds}
applies at level $k+1$, its hypotheses being read at that level, with the same
constants, and these depend on $d$ and $L$ only.

Consequently two hypotheses are supplied. The first is hypothesis~(f) of the statement
that one renormalisation step closes on the one-power hereditary class. The second is
hypothesis~(i), which is the reading $(\ell 1)$, of the statement that the minimal
configurations $U_j$ with $\varepsilon_0$ sufficiently small lie in every
$U^{\Sigma^1}_j(X;\alpha_0,\alpha_1)$. For the latter the only link to the radius is the
condition $2B_3\varepsilon_0\le\alpha_0$.
\end{corollary}

\begin{proof}
\emph{The small pairs and the unit pair.}
The unit pair is decomposed with unitary factor $U:=1$. A small pair
$(e^{i\eta\mathsf{Z}},J(e^{i\eta\mathsf{Z}}))$, with $\mathsf{Z}$ a
$\mathfrak{g}^{c}$-valued field, is decomposed as
\begin{equation*}
U:=1,\qquad U':=e^{i\eta\mathsf{Z}}.
\end{equation*}
These decompositions are the ones listed for these two pairs in
\eqref{4.9:eq-gauge-transport-a}. The gauge condition bears on the unitary factor
$U$, which here is $1$. By Lemma~\ref{4.9:lem-identity-field}, $U\equiv1$ obeys
$\mathsf{g}_j(1;D;\rho)$ for every level $j$, every set $D$ and every $\rho>0$; it
suffices to take $u:=1$ and $A:=0$. Hence $(\ell 1)$ holds for these pairs at every
radius, and so at every value of $\alpha_0'$.

\emph{Real backgrounds at their own level.}
Let $U=U_k(V)\in\mathfrak{U}_k(\varepsilon_0)$ be a minimal configuration. The
hypotheses of Proposition~\ref{4.9:prop-real-backgrounds} hold by assumption. That
proposition gives, for every closed cube $\square$ of side $s_\square L^k$ with
$0<s_\square\le S$,
\begin{equation*}
\mathsf{g}_k\bigl(U;\square;2B_3(\tfrac12 s_\square+2R_1M_1)\varepsilon_0\bigr).
\end{equation*}
The floor $S\ge4R_1M_1$ is among its hypotheses, so $2R_1M_1\le\tfrac12S$. Together
with $s_\square\le S$ this gives
\begin{equation}\label{4.9:eq-printed-pairs-1}
\begin{aligned}
2B_3\bigl(\tfrac12 s_\square+2R_1M_1\bigr)\varepsilon_0
&\le 2B_3\,S\,\varepsilon_0\\
&= c_{12}LM\cdot 2B_3\,\varepsilon_0\\
&\le c_{12}LMB_{12}\,\varepsilon_0 ,
\end{aligned}
\end{equation}
where the last step uses $B_{12}\ge2B_3$. The cubes of
Proposition~\ref{4.9:prop-real-backgrounds} are exactly the cubes of side at most
$S=c_{12}LM$ in level-$k$ units. Since the cube gauge clause at a radius $\rho$ implies
it at every radius $\rho'\ge\rho$, $U$ satisfies $(\ell 1)$ on all of them with the
single bound $c_{12}LMB_{12}\alpha_0'$, where $\alpha_0'=\varepsilon_0$. This
$\alpha_0'$ is the constant of the membership
$U\in\mathfrak{U}_k(\varepsilon_0)$ of \cite[(1.2)]{B-CMP109}.

\emph{Real backgrounds at a lower level.}
Let $1\le j\le k$ and $\mu=L^{-(k-j)}$. The second conclusion of
Proposition~\ref{4.9:prop-real-backgrounds} applies to every closed set $X$ and every
such $j$. It states that $U|_X$ obeys the reading $(\ell 1)$ of
\eqref{4.9:eq-cube-gauge-condition-b} at level $j$, with the constant $\varepsilon_0$ of
level $k$ carrying exactly one factor $\mu$. Thus, on every cube of side at most $S$ in
level-$j$ units, $U|_X$ satisfies the cube gauge condition with the single bound
\begin{equation*}
c_{12}LMB_{12}\,\mu\varepsilon_0=c_{12}LMB_{12}\,L^{-(k-j)}\varepsilon_0 .
\end{equation*}
The constant of the membership at level $j$ is therefore
$\alpha_0'=\mu\varepsilon_0=L^{-(k-j)}\varepsilon_0$.

\emph{Reproduction under the step.}
Proposition~\ref{4.9:prop-real-backgrounds} is stated for every level $k\ge1$. Its
ceilings and floors involve the constants $R_1$, $M_1$, $a_1$ and $B_3$, which depend
on $d$ and $L$ only and not on the level. It therefore applies at level $k+1$, with the
same constants, to $U_{k+1}(V)\in\mathfrak{U}_{k+1}(\varepsilon_0)$, its hypotheses being
read at level $k+1$; in particular the torus-side condition is then read in
level-$(k+1)$ units. The argument above then applies verbatim to $U_{k+1}(V)$, and to
the unit pair and the small pairs considered above at level $k+1$.

\emph{The supplied hypotheses.}
The preceding parts of the proof establish the condition in the form stated, which is
hypothesis~(f); the parts on real backgrounds establish it for the minimal
configurations, which is hypothesis~(i) of the statement on minimal configurations.
\end{proof}

\begin{remark}[Choice of the constant, locality, small tori]\label{4.9:rem-seed-remarks}
We make three observations on Proposition~\ref{4.9:prop-real-backgrounds}, which states that real
backgrounds satisfy the ambient gauge condition.
\begin{enumerate}
\item \emph{Choice of the constant.} Ba{\l}aban's papers themselves defer the choice of the
constant of the gauge condition to the theorem that supplies the background configuration,
\cite[Theorem~1]{B-CMP102b}, whose constant $B_3$ is the one used in
Proposition~\ref{4.9:prop-real-backgrounds}. In \cite{B-CMP109} the gauge condition is imposed
``with a sufficiently large constant $B$
(it will be determined later)'', and in \cite{B-CMP122a} one reads ``for example we can take the
constant $B = B_3$ from Theorem 1 [16]'' (reference~16 of that paper is \cite{B-CMP102b}). The
factor $2$ in the radius
\begin{equation*}
2B_3\bigl(\tfrac12 s_{\square}+2R_1M_1\bigr)\varepsilon_0
\end{equation*}
of the conclusion $\mathsf{g}_k(U;\square;\cdot)$ of Proposition~\ref{4.9:prop-real-backgrounds}
is not introduced here: it is inherited from the proposition of Ba{\l}aban's that
\cite{B-CMP122a} invokes beside this choice of the constant.
\item \emph{Locality.} In the clause of Proposition~\ref{4.9:prop-real-backgrounds} that carries the
layerwise
constants $a_n$, these are the constants of the other conditions defining the real configuration's
membership in the space of \cite[\S2]{B-CMP119}, namely the first
inequality of \cite[(2.34)]{B-CMP119} and \cite[(2.36)]{B-CMP119}, which enter through the third
condition of \eqref{4.9:eq-collar-regularity-b} and are taken by statement; nothing about them is
asserted here. A single constant $\varepsilon_1$ for all layers, as in \cite[(8)]{B-CMP102b}, would not
suffice when the depth $k-j$, $j$ being the level at which Lemma~\ref{4.9:lem-seed-cube-gauge} is
read, tends to infinity. The cube gauge lemma
(Lemma~\ref{4.9:lem-seed-cube-gauge}), however, is local: it reads the constants on the collar
$\mathcal{N}(\square^{a})$ only, so that no dependence on the depth $k-j$ enters.
\item \emph{Small tori.} The condition $(\mathrm{T}_{\mathrm{low}})^{S}$ of
\eqref{4.9:eq-real-backgrounds-a} is a one-sided lower bound on the side of the torus. On a torus
too small to contain the enlarged cube $\widetilde{\square}$ the ambient gauge construction is not
available, and Proposition~\ref{4.9:prop-real-backgrounds} claims nothing there.
\end{enumerate}
\end{remark}

\begin{lemma}[A flat field winding around a hole]\label{4.9:lem-flat-ring}
Fix a level $k$ and put $M_s:=ML^k$. Let $C_i$, $i\in\mathbb{Z}^4$, be the cells of $\pi_k$,
with the lattice and closed cubes as in the conventions of~\ref{4.9:not-lattice-conventions}.
Every site lies in exactly one cell. For each direction $\mu$, the $\mu$-th coordinates of
the sites of the cells of index $i_\mu$ form a run of $M_s$ consecutive integers, called
\emph{block} $i_\mu$, and consecutive indices give consecutive runs. After a translation,
block $0$ begins at the same coordinate $p$ in directions $1$ and $2$. Thus, for
$\mu\in\{1,2\}$, block $i_\mu$ is $[p+i_\mu M_s,\,p+(i_\mu+1)M_s-1]$. Put
$H:=[p,p+2M_s-1]$, the union of blocks $0$ and $1$, and let $c:=p+M_s-1$ be the last
coordinate of block $0$.

Let
\begin{equation*}
\mathsf{R}:=\bigl\{i\in\mathbb{Z}^4:\ (i_1,i_2)\in\{-1,0,1,2\}^2\setminus\{0,1\}^2,\
i_3=i_4=0\bigr\}.
\end{equation*}
This is a wall-connected set of twelve cells, and
$X_{\mathsf{R}}:=\bigcup_{i\in\mathsf{R}}C_i\in\mathbf{D}_k$. The \emph{hole} consists of
the four cells with $(i_1,i_2)\in\{0,1\}^2$ and $i_3=i_4=0$. Its sites are those with
$x_1,x_2\in H$ and $x_3,x_4$ in block $0$.

Let $Y'$ be any union of cells, all of index $i_1\le-2$, such that
$Y:=X_{\mathsf{R}}\cup Y'$ is wall-connected. An example is a tail of cells leaving the wall
of $C_{(-1,0,0,0)}$ that faces $-e_1$; the choice $Y'=\emptyset$ is allowed. Then
$Y\in\mathbf{D}_k$.

Let $N\ge2$ and $h_0:=\operatorname{diag}(-1,-1,1,\dots,1)\in SU(N)$. A bond
$\langle x,x+e_2\rangle$ is \emph{cut} if $x_2=c$, $x_1$ lies in block $2$, and $x_3,x_4$ lie
in block $0$. Let $U_{\mathsf{R}}$ be the field that equals $h_0$ on the cut bonds and $1$ on
all other bonds of $\operatorname{bd}(Y)$.

Let $\Gamma$ be the square lattice loop, traversed in either orientation, that satisfies the
following. It lies in the plane of directions $1,2$ at a fixed $(x_3,x_4)$ with $x_3,x_4$ in
block $0$, and it runs through the coordinates $p-1$ and $p+2M_s$. Equivalently, its sites
are the points of the boundary of the square $[p-1,p+2M_s]^2$ in the coordinates
$(x_1,x_2)$. Its sides have $2M_s+1$ steps, and it has $n_{\Gamma}=4(2M_s+1)$ bonds. Let
$c_{12}$ be the constant written $O(1)$ in \cite[(1.12)]{B-CMP109}, and put $S:=c_{12}LM$.
Then:
\begin{itemize}
\item[(W1)] $U_{\mathsf{R}}$ is flat on $Y$: $\partial U_{\mathsf{R}}=1$ on every plaquette
whose four corners are sites of $Y$.
\item[(W2)] $\Gamma\subset X_{\mathsf{R}}$ and $T_{\Gamma}(U_{\mathsf{R}})=h_0^{\pm1}$, the
sign depending on the orientation of $\Gamma$.
\item[(W3)] Let $\square$ be a closed cube of side at most $2M_s$. Then $U_{\mathsf{R}}$ is a
pure gauge on $\operatorname{bd}(\square\cap Y)$: there is
$u:\operatorname{st}(\square\cap Y)\to G$ with $U_{\mathsf{R}}^{u}=1$ on
$\operatorname{bd}(\square\cap Y)$. Consequently
$\mathsf{g}_n(U_{\mathsf{R}};\square\cap Y;\rho)$ holds for every level $n$ and every
$\rho>0$, with $A\equiv0$.
\item[(W4)] $\Gamma$ lies in a closed cube of side $c_{12}LM_s$ if and only if
$2M_s+1\le c_{12}LM_s$. This inequality follows from the one-sided floor
\begin{equation}\label{4.9:eq-flat-ring-a}
(c_{12}L-2)M\ge1 .
\end{equation}
\end{itemize}
\end{lemma}

\begin{proof}
Since every site lies in exactly one cell, the sites of $Y$ are exactly the sites of the
cells $C_i$ with $i\in\mathsf{R}$ and of the cells of $Y'$.

\emph{(W1).} Consider first a plaquette whose corners are sites of $Y$ and which contains no
cut bond. All four of its bonds lie in $\operatorname{bd}(Y)$ and carry $1$, so its plaquette
variable is $1$.

Now consider a plaquette whose corners are sites of $Y$ and which contains a cut bond
$b=\langle x,x+e_2\rangle$. It lies in a plane of directions $(2,\nu)$ with
$\nu\in\{1,3,4\}$, and its second bond in direction $2$ is $b+e_\nu$ or $b-e_\nu$. The
endpoints of the shifted bond are $x\pm e_\nu$ and $x+e_2\pm e_\nu$. Their second coordinates
are $c$ (block $0$) and $c+1$ (block $1$).

Suppose first that the shifted coordinate $x_\nu\pm1$ leaves the block of $x_\nu$.
\begin{itemize}
\item For $\nu=1$ it lies in block $1$ or in block $3$. The two endpoints then lie in the
cells $(1,0,0,0)$ and $(1,1,0,0)$, which belong to the hole, or in $(3,0,0,0)$ and
$(3,1,0,0)$.
\item For $\nu=3$ they lie in $(2,0,\pm1,0)$ and $(2,1,\pm1,0)$.
\item For $\nu=4$ they lie in $(2,0,0,\pm1)$ and $(2,1,0,\pm1)$.
\end{itemize}
None of these cells is in $\mathsf{R}$: hole cells are excluded, $i_1=3$ lies outside
$\{-1,\dots,2\}$, and $i_3\ne0$ or $i_4\ne0$ in the remaining cases. None of them has
$i_1\le-2$, so none belongs to $Y'$. Hence these endpoints are not sites of $Y$, and no
plaquette of this kind has all its corners in $Y$.

Otherwise the shifted coordinate stays in its block. The shifted bond then starts at a site
with second coordinate $c$, first coordinate in block $2$, and third and fourth coordinates
in block $0$, so it is cut as well and carries $h_0$. The two remaining bonds have direction
$\nu\ne2$; they are not cut and carry $1$. The two bonds in direction $2$ are traversed in
opposite senses around the plaquette. In either orientation, the ordered product is therefore,
up to cyclic order, $h_0\cdot1\cdot h_0^{-1}\cdot1=1$.

\emph{(W2).} Every site of $\Gamma$ has one of $x_1,x_2$ in $\{p-1,p+2M_s\}$. These
coordinates are the last coordinate of block $-1$ and the first of block $2$. The other of
$x_1,x_2$ lies in $[p-1,p+2M_s]$, that is, in blocks $-1,\dots,2$. Moreover $x_3,x_4$ lie in
block $0$. Hence the cell of each site of $\Gamma$ has $(i_1,i_2)\in\{-1,\dots,2\}^2$ with at
least one entry in $\{-1,2\}$, and $i_3=i_4=0$; it belongs to $\mathsf{R}$. Thus
$\Gamma\subset X_{\mathsf{R}}\subset Y$, and $U_{\mathsf{R}}$ is defined on every bond of
$\Gamma$.

The bonds of $\Gamma$ in direction $2$ lie on its two sides $x_1=p-1$ and $x_1=p+2M_s$. Each
of these sides runs through all values of $x_2$ from $p-1$ to $p+2M_s$. So $\Gamma$ contains
exactly two bonds from $x_2=c$ to $x_2=c+1$:
\begin{itemize}
\item the one at $x_1=p+2M_s$ (block $2$) is cut;
\item the one at $x_1=p-1$ (block $-1$) is not cut.
\end{itemize}
Every other bond of $\Gamma$ either has direction $1$ or does not go from $x_2=c$ to $c+1$, so
it is not cut. Hence $T_{\Gamma}(U_{\mathsf{R}})$ is an ordered product of factors $1$ and a
single factor $h_0^{\pm1}$. The exponent is $+1$ or $-1$ according to the sense in which
$\Gamma$ traverses the cut bond. So $T_{\Gamma}(U_{\mathsf{R}})=h_0^{\pm1}$.

\emph{(W3).} Write $\operatorname{st}(\square)=\prod_\mu J_\mu$, where each $J_\mu$ is an
interval of integers with $\#J_\mu\le2M_s+1$.

If $\operatorname{bd}(\square\cap Y)$ contains no cut bond, put $u:=1$. Every bond of
$\operatorname{bd}(\square\cap Y)$ carries $1$, so $U_{\mathsf{R}}^u=1$ there.

Otherwise $J_1$ meets block $2=[p+2M_s,p+3M_s-1]$, since the endpoints of a cut bond have
first coordinate in block $2$. As $J_1$ consists of at most $2M_s+1$ consecutive integers and
contains an integer $\ge p+2M_s$, we get $\min J_1\ge p$. Every site of $\square\cap Y$ thus
has $x_1\ge p$ and lies in a cell of index $i_1\ge0$. Such a cell is not in $Y'$, so it is a
cell of $\mathsf{R}$ with $i_1\in\{0,1,2\}$. In particular, $x_2$ ranges over blocks
$-1,\dots,2$ on $\operatorname{st}(\square\cap Y)$. Define
\begin{equation*}
u(x):=1\quad(x_2\le c),\qquad u(x):=h_0\quad(x_2\ge c+1),
\qquad x\in\operatorname{st}(\square\cap Y).
\end{equation*}

Let $b$ be a bond of $\operatorname{bd}(\square\cap Y)$ from $x_2=c$ to $x_2=c+1$. Its
endpoints lie in cells $(i_1,0,0,0)$ and $(i_1,1,0,0)$ of $\mathsf{R}$ with
$i_1\in\{0,1,2\}$. Since the cells with $(i_1,i_2)\in\{0,1\}^2$ form the hole, $i_1=2$. So $b$
is cut, and
\begin{equation*}
U_{\mathsf{R}}^u(b)=1\cdot h_0\cdot h_0^{-1}=1 .
\end{equation*}
Every other bond of $\operatorname{bd}(\square\cap Y)$ has equal values of $u$ at its two ends.
Indeed, either it has a direction other than $2$, or it is a bond in direction $2$ whose
endpoints both have $x_2\le c$ or both have $x_2\ge c+1$. Such a bond is not cut, so it
carries $1$, and $U_{\mathsf{R}}^u=u\,1\,u^{-1}=1$ on it. Thus $U_{\mathsf{R}}^u=1$ on all
of $\operatorname{bd}(\square\cap Y)$.

Take this $u$ and $A\equiv0$ in the cube gauge condition
(Definition~\ref{4.9:def-cube-gauge-condition}), with $D=\square\cap Y$. We have
$\operatorname{bd}(\square\cap Y)\subset\operatorname{bd}(Y)$, so $U_{\mathsf{R}}$ is defined
there. The identity $U_{\mathsf{R}}^u=e^{i\xi_nA}$ reads $1=1$. The bounds $|A(b)|<\rho$ and
$\xi_n^{-1}|A(b+e_\nu)-A(b)|<\rho$ read $0<\rho$. Hence
$\mathsf{g}_n(U_{\mathsf{R}};\square\cap Y;\rho)$ of \eqref{4.9:eq-cube-gauge-condition-a}
holds for every $n$ and every $\rho>0$.

\emph{(W4).} In directions $1$ and $2$ the sites of $\Gamma$ fill the coordinate range
$[p-1,p+2M_s]$, an extent of $2M_s+1$ steps. In directions $3$ and $4$ they have a single
coordinate. A closed cube of side $c_{12}LM_s$ has extent $c_{12}LM_s$ in each direction.
Hence it can contain $\Gamma$ if and only if $2M_s+1\le c_{12}LM_s$.

Now assume \eqref{4.9:eq-flat-ring-a}. Multiplying by $L^k\ge1$ gives
\begin{equation*}
c_{12}LM_s-2M_s=(c_{12}L-2)ML^k\ge L^k\ge1,
\end{equation*}
that is, $2M_s+1\le c_{12}LM_s$.
\end{proof}

The floor \eqref{4.9:eq-flat-ring-a} is one branch of a case split; it is not a bound on $L$.
If $c_{12}L\le2$, then $S\le2M$. In that case the reading $(\ell 2)$ of the cube gauge
condition (Definition~\ref{4.9:def-cube-gauge-condition}) is implied, up to the constant, by
the cubes of the reading $(\ell 2)^{\flat}$. Ba{\l}aban's data $c_{12}\ge10$
\cite{B-CMP99a} and $L>11$ \cite{B-CMP109} put us in the branch $c_{12}L>2$ of the floor.

\begin{proposition}[Ambient versus inner cube gauge conditions]\label{4.9:prop-ambient-condition}
We use the notation of Definition~\ref{4.9:def-cube-gauge-condition}. Thus $S=c_{12}LM$ and
$\xi_j=L^{-j}$, and $(\ell 0)$, $(\ell 1)=(\mathrm{L}\text{-}12)$, $(\ell 2)$, $(\ell 2)^{\flat}$ are the
readings \eqref{4.9:eq-cube-gauge-condition-b} of the cube gauge condition. Each reading is taken with
its gradient member and with one pair $(B,B_{12})$. We write $\operatorname{dist}_{\infty}$ for the
sup-distance in lattice steps, and $D^{+r}$ for the closed sup-distance $r$-neighbourhood of $D$.
At a level $k$ write $M_s=ML^k$ and
\begin{equation*}
h_0=\operatorname{diag}(-1,-1,1,\dots,1)\in G .
\end{equation*}
Assume the following four statements.
\begin{enumerate}
\item[(a)] $(\ell 0)$ does not imply $(\ell 1)$.
\item[(b)] At a level $k$ there are a closed set $X$ with a hole of one
cell and a flat $G$-valued field $U$ on $\operatorname{bd}(X)$ with the following properties:
for every $a>0$, $U$ obeys $(\ell 1)_k(X;a)$, with $A\equiv 0$; and $X$ contains a closed lattice
path around the hole
which has $4(M_s+1)$ bonds, lies in a closed cube of side $M_s+1$, and along which the holonomy of $U$
is $h_0^{\pm1}$.
\item[(c)] The half of the independence statement, Corollary~\ref{4.4:cor-localised-independence},
that concerns cubes in the small-field region. In the setting of that corollary, let $j$ be a level and
let $X\subset\Lambda_j$ be closed. Then the $G$-valued factor $U$ of
$(\mathbf{U}^{\flat},\mathbf{J}^{\flat})|_{X}$ obeys
$\mathsf{g}_j(U;\square\cap X;BCM\alpha_{0,j})$ for every closed cube $\square\subset\Omega_j$ of
side $CML^j$ with $C\in\{1,2\}$.
\item[(d)] Identities for the step.
\begin{itemize}
\item The two identities stating that, when the minimal orbit of a localised variational problem is
the orbit of a flat field, the plaquette variables of its minimal configuration are $1$ and its current
is $0$. Precisely: let $Y$ be a domain at level $k+1$ as in Definition~\ref{4.4:def-post-step-space},
let $X''\subset Y$ be an admissible domain, and let $\iota$ be an index of the localised problems on
$X''$ over which the clauses \cite[(2.35)]{B-CMP119} and \cite[(2.37)]{B-CMP119} range. Write
$\partial U_{\iota,X''}$ and $\mathbf{J}_{\iota,X''}$ for the plaquette variables and the current of
the corresponding localised minimal configuration. If the minimal orbit of that problem is the orbit
of a flat field, then $\partial U_{\iota,X''}-1=0$ and $\mathbf{J}_{\iota,X''}=0$.
\item The identity $\mathfrak{J}(0)=0$ for the current of the step at $w=0$.
\end{itemize}
\end{enumerate}
Then the following hold.

\smallskip
(i) \emph{Hierarchy.}
\begin{itemize}
\item $(\ell 2)_j(X;a)\Rightarrow(\ell 1)_j(X;a)\Rightarrow(\ell 0)_j(X;a)$.
\item $(\ell 2)^{\flat}_j(X;a)\Rightarrow(\ell 1)_j(X;a)$ under the conditions of
Corollary~\ref{4.9:cor-gluing-smaller-cubes}.
\item If $c_{12}L\ge 2$, then $(\ell 2)_j(X;a)$ gives $\mathsf{g}_j(U;\square\cap X;SB_{12}a)$ for
every closed cube $\square$ of side $CML^j$, $C\in\{1,2\}$. These are the cubes of $(\ell 2)^{\flat}$,
all at the one radius $SB_{12}a$.
\end{itemize}
None of these implications reverses:
\begin{itemize}
\item $(\ell 0)$ does not imply $(\ell 1)$;
\item under condition $(\mathrm{gl})_a$ of \eqref{4.9:eq-gluing-smaller-cubes-a}, $(\ell 1)$ does not
imply $(\ell 2)^{\flat}$: the field of (b) obeys
$(\ell 1)_k(X;a)$ but not $(\ell 2)^{\flat}_k(X;a)$;
\item if $(c_{12}L-2)M\ge 1$ and $4S\cdot SB_{12}a<2$, then $(\ell 2)^{\flat}$ does not imply
$(\ell 2)$: the flat field $U_{\mathsf{R}}$ of Lemma~\ref{4.9:lem-flat-ring} obeys
$(\ell 2)^{\flat}_k(X_{\mathsf{R}};a)$ but not $(\ell 2)_k(X_{\mathsf{R}};a)$.
\end{itemize}

\smallskip
(ii) \emph{What the ambient cube adds.} In the fixed-size reading $(\ell 0)$, which is
Ba{\l}aban's, replace the requirement ``every closed cube $\square\subset X$ of the single side
$SL^j$, with the gauge $u$ defined on $\square$'' by the requirement ``every closed cube $\square$ of
the ambient lattice, with the gauge $u$ defined on $\square\cap X$''.
This adds the three properties below, and on cubes $\square\subset X$ it adds nothing.
\begin{enumerate}
\item[(R1)] \emph{Heredity independent of shape.} The condition passes to every closed
$X'\subset X$ and to every lower level, whatever the thickness of $X$ and of $X'$. The widened reading
$(\ell 1)$ shares this property. The fixed-size reading $(\ell 0)$ does not, since it is void on sets
$X'$ too thin to contain a cube of side $SL^j$.
\item[(R2)] \emph{Seeds without geometry.} Here the seeds are the real backgrounds of
Proposition~\ref{4.9:prop-real-backgrounds}. Let a $G$-valued field obey
$\mathsf{g}_j(\cdot\,;\square;\rho)$
on the whole ambient cube $\square$. Then its restriction to $X$ obeys
$\mathsf{g}_j(\cdot\,;\square\cap X;\rho)$. Hence real backgrounds obey the ambient readings by
Ba{\l}aban's theorems on cubes, applied with a collar in the ambient lattice
(Proposition~\ref{4.9:prop-real-backgrounds}).
\item[(R3)] \emph{No holonomy inside $\square$.} If $\mathsf{g}_j(U;\square\cap X;\rho)$ holds, then
every closed lattice path $\lambda$ of $\ell$ bonds in $\square\cap X$ satisfies
$\|T_{\lambda}(U)-1\|_{\mathrm{op}}\le \ell\xi_j\rho$. This is a condition on the holes and the
re-entrant parts of $X$ inside $\square$. On these, the conditions \cite[(1.11)]{B-CMP109} and
$(\ell 1)$ together say nothing, as the second non-reversal in (i) shows. It is this property that
separates the ambient readings from $(\ell 1)$.
\end{enumerate}

\smallskip
(iii) \emph{The reading $(\ell 2)$ is not reproduced by the straddling-sources half of the
independence statement.} Let $X_{\mathsf{R}}$, $U_{\mathsf{R}}$, $Y$, the loop $\Gamma$ and its number
of bonds $n_{\Gamma}$ be those of Lemma~\ref{4.9:lem-flat-ring} at level $k$. Assume
$(c_{12}L-2)M\ge 1$ and the amplitude ceiling
\begin{equation}\label{4.9:eq-ambient-condition-a}
(\mathrm{c\text{-}W})\qquad 4S\bigl(a_1+SB_{12}a_0\bigr)\le \tfrac12 .
\end{equation}
Let $\{\Omega_n\}_{n\le k+1}$, $\{\Lambda_n\}$ be a nest with $X_{\mathsf{R}}\subset\Lambda_k$ and
$X_{\mathsf{R}}\cap\Omega_{k+1}=\emptyset$. Nests produced by the large-field operation $\mathbf{T}$
leave room for this: there $\Lambda_k\setminus\Omega_{k+1}$ contains a collar of at least
$R_k\ge L$ cells. Then:
\begin{enumerate}
\item[(iii-1)] The pair $(U_{\mathsf{R}},0)$ satisfies on $Y$ every clause of
\cite[(2.34)--(2.39)]{B-CMP119}. Here every cube of the family \eqref{4.9:eq-cube-gauge-condition-c}
of index $n$ (side $CML^n$, $C\in\{1,2\}$) is required to
satisfy its clause. The pair also satisfies the hereditary conditions $(\Sigma)$ and
$(\Sigma\text{-iv})'$ of the post-step space of Definition~\ref{4.4:def-post-step-space}. All of this
holds at all positive radii, because every left-hand member is $0$. Hence
\begin{equation*}
(U_{\mathsf{R}},0)\in\mathfrak{U}^{\mathrm{post},U}_{k+1}(Y).
\end{equation*}
\item[(iii-2)] At $w=0$ one has $(\mathbf{U}^{\flat},\mathbf{J}^{\flat})=(U_{\mathsf{R}},0)$.
\item[(iii-3)] No point of the $G^{c}$-orbit of $(U_{\mathsf{R}},0)|_{X_{\mathsf{R}}}$ has a
decomposition $e^{i\xi_kA'}\widetilde U$ of the form $\mathbf{U}=U'U$ with $|A'|<a_1$ and with
$\widetilde U$ obeying $(\ell 2)_k(X_{\mathsf{R}};a_0)$.
\end{enumerate}
Consequently, consider the straddling-sources half of Corollary~\ref{4.4:cor-localised-independence}.
That half holds with its target gauge condition read as $(\mathrm{L}\text{-}12)$: this is its entry in
the table \eqref{4.9:eq-gauge-transport-a}, obtained from the half that concerns cubes in the
small-field region together with Corollary~\ref{4.9:cor-gluing-smaller-cubes}, under the conditions of
that corollary. By (iii-1)--(iii-3), it fails with that target read as $(\ell 2)$.

\smallskip
(iv) \emph{The reading $(\ell 2)^{\flat}$ is reproduced everywhere.} We work in the setting of
Corollary~\ref{4.4:cor-localised-independence}. Let $X\subset\Lambda_j$ be closed, with
\begin{equation*}
L^{-j}\operatorname{dist}_{\infty}(X,\Omega_j^{c})\ge 2M .
\end{equation*}
For nests produced by $\mathbf{T}$, this margin follows from the margin conditions on the domains
$\Omega_j$, $\Lambda_j$ in \cite{B-CMP119}. Then the $G$-valued factor $U$ of
$(\mathbf{U}^{\flat},\mathbf{J}^{\flat})|_{X}$ obeys $(\ell 2)^{\flat}_j(X;\alpha_{0,j})$.
Together with parts (a) and (b) of Proposition~\ref{4.9:prop-gauge-transport} and with
Proposition~\ref{4.9:prop-real-backgrounds}, this gives the following. The reading
$(\ell 2)^{\flat}$ holds at every locus of the table \eqref{4.9:eq-gauge-transport-a}, and at the
seeds, namely the real backgrounds of Proposition~\ref{4.9:prop-real-backgrounds}.
\end{proposition}

\begin{proof}
\emph{(i), the implications.} We begin with $(\ell 2)\Rightarrow(\ell 1)\Rightarrow(\ell 0)$, which is
Lemma~\ref{4.9:lem-restriction-monotone}.
\begin{itemize}
\item Let $\square'\subset X$ be a closed cube of side at most $SL^j$. Then $\square'$ lies in a closed
cube $\square$ of side $SL^j$ of the ambient lattice, and $\square'\subset\square\cap X$. So
$\mathsf{g}_j(U;\square\cap X;\rho)$ gives $\mathsf{g}_j(U;\square';\rho)$, by restricting $u$ and $A$.
This proves $(\ell 2)\Rightarrow(\ell 1)$.
\item The cubes of side exactly $SL^j$ are among those of side at most $SL^j$. This proves
$(\ell 1)\Rightarrow(\ell 0)$.
\end{itemize}
The implication $(\ell 2)^{\flat}\Rightarrow(\ell 1)$ is Corollary~\ref{4.9:cor-gluing-smaller-cubes}.
For the third implication, let $c_{12}L\ge 2$. Every closed cube $\square'$ of side $CML^j$,
$C\in\{1,2\}$, satisfies
\begin{equation*}
CML^j\le 2ML^j\le c_{12}LML^j=SL^j .
\end{equation*}
So $\square'$ lies in a closed cube $\square$ of side $SL^j$, and $\square'\cap X\subset\square\cap X$.
By Lemma~\ref{4.9:lem-restriction-monotone}, $(\ell 2)_j(X;a)$ then gives
$\mathsf{g}_j(U;\square'\cap X;SB_{12}a)$.

\emph{(i), the first non-reversal.} That $(\ell 0)$ does not imply $(\ell 1)$ is hypothesis (a).

\emph{(i), the annulus.} Let $X$, $U$ and the closed path be as in hypothesis (b); call the path
$\lambda$. There $(\ell 1)_k(X;a)$ holds with $A\equiv 0$. The path $\lambda$ has $4(M_s+1)$ bonds. It
lies in a closed cube of side $M_s+1\le 2M_s$, hence in a closed cube $\square$ of side
$2M_s=2ML^k$. Its holonomy is $h_0^{\pm1}$.

Suppose $(\ell 2)^{\flat}_k(X;a)$ held. For the cube $\square$, with $C=2$, it would give
$\mathsf{g}_k(U;\square\cap X;2MBa)$. The bonds of $\lambda$ lie in $\operatorname{bd}(\square\cap X)$.
Lemma~\ref{4.9:lem-holonomy-bound} would then give the chain
\begin{align*}
2=\|h_0^{\pm1}-1\|_{\mathrm{op}}
&\le 4(M_s+1)\xi_k\cdot 2MBa\le 16M^2Ba\\
&\le 6S^2Ba\le r_{\mathrm{gl}}\le 1 .
\end{align*}
We justify each step in turn.
\begin{itemize}
\item The equality holds because $h_0^{-1}=h_0$ and
$h_0-1=\operatorname{diag}(-2,-2,0,\dots,0)$.
\item The second inequality uses $M_s+1\le 2M_s$ and $M_s\xi_k=M$.
\item The third uses $S=c_{12}LM\ge 2M$, that is $c_{12}L\ge 2$, as in the third implication of (i).
\item The fourth is $(\mathrm{gl})_a$ of \eqref{4.9:eq-gluing-smaller-cubes-a}.
\item The last holds because $r_{\mathrm{gl}}\in(0,1]$ (Lemma~\ref{4.4:lem-gluing-boxes}).
\end{itemize}
The chain gives $2\le 1$, which is absurd. So $(\ell 2)^{\flat}_k(X;a)$ fails while $(\ell 1)_k(X;a)$
holds.

\emph{(i), the ring.} Let $(c_{12}L-2)M\ge 1$ and $4S\cdot SB_{12}a<2$. Property (W3) of the ring of
Lemma~\ref{4.9:lem-flat-ring} gives $(\ell 2)^{\flat}_k(X_{\mathsf{R}};a)$.

Suppose $(\ell 2)_k(X_{\mathsf{R}};a)$ held. By (W4) there is a closed cube $\square\supset\Gamma$ of
side $SL^k$, and $n_{\Gamma}\xi_k\le 4S$. Then $(\ell 2)$ gives
$\mathsf{g}_k(U_{\mathsf{R}};\square\cap X_{\mathsf{R}};SB_{12}a)$. Lemma~\ref{4.9:lem-holonomy-bound},
applied to the $n_{\Gamma}$ bonds of $\Gamma$, gives
\begin{equation*}
\|T_{\Gamma}(U_{\mathsf{R}})-1\|_{\mathrm{op}}\le n_{\Gamma}\xi_k\,SB_{12}a\le 4S\cdot SB_{12}a<2 .
\end{equation*}
By (W2), however, $T_{\Gamma}(U_{\mathsf{R}})$ has the eigenvalue $-1$. So the left-hand side is at least
$2$, which is a contradiction.

\emph{(ii).} The three properties are given by the lemmas on the cube gauge condition.
\begin{itemize}
\item (R1) follows from two facts. First, Lemma~\ref{4.9:lem-restriction-monotone}: for $X'\subset X$
and every ambient cube $\square$ we have $\square\cap X'\subset\square\cap X$, and for $(\ell 1)$ a
cube contained in $X'$ is a cube contained in $X$. Second, Lemma~\ref{4.9:lem-level-change}, for the
passage to lower levels. The fixed-size reading $(\ell 0)$ constrains only the cubes of side $SL^j$
contained in the set. It is therefore void on a set that contains no such cube.
\item (R2) is Lemma~\ref{4.9:lem-restriction-monotone} with $D=\square$ and $D'=\square\cap X$.
\item (R3) is Lemma~\ref{4.9:lem-holonomy-bound} with $D=\square\cap X$.
\end{itemize}
On a cube $\square\subset X$ one has $\square\cap X=\square$. So the two requirements coincide there.

\emph{(iii), the clauses of \cite[(2.34)--(2.39)]{B-CMP119}.}
\begin{itemize}
\item Clause \cite[(2.34)]{B-CMP119} is property (W1) of the ring of Lemma~\ref{4.9:lem-flat-ring}.
\item For \cite[(2.36)]{B-CMP119} and \cite[(2.39)]{B-CMP119}: the pair has $\mathbf{J}=0$ and
$A'=0$.
\item For \cite[(2.38)]{B-CMP119}, take first a cube of that family of index $n\ge k+1$. It lies in
$\Omega_{k+1}$. Since $X_{\mathsf{R}}\cap\Omega_{k+1}=\emptyset$, no cut bond of the ring lies there,
so $U_{\mathsf{R}}=1$ on its bonds. A cube of that family of index $n\le k$ has side at most
$2ML^n\le 2M_s$, and for it (W3) applies.
\item For \cite[(2.35)]{B-CMP119}, \cite[(2.37)]{B-CMP119}, $(\Sigma)$ and $(\Sigma\text{-iv})'$, fix an
admissible $X''\subset Y$ and an index $\iota$ as in hypothesis (d). The flat field
$U_{\mathsf{R}}|_{X''}$ satisfies the
constraints with data $\dot{M}U_{\mathsf{R}}$. Its action is $0$, which is the absolute minimum. It lies
in the space \cite[(6)]{B-CMP102b}. The data obey \cite[(7)]{B-CMP102b} with $\varepsilon_1$
arbitrarily small, by \cite[Proposition~2]{B-CMP98} at zero curvature. By the uniqueness clause of
\cite[Theorem~1]{B-CMP102b}, the minimal orbit is the orbit of $U_{\mathsf{R}}|_{X''}$. Hypothesis (d)
then gives $\partial U_{\iota,X''}-1=0$ and $\mathbf{J}_{\iota,X''}=0$.
\end{itemize}
Every left-hand member is therefore $0$. This proves (iii-1).

\emph{(iii-2).} At $w=0$ one has $\mathcal{H}_0=0$, and $\mathfrak{J}(0)=0$ by hypothesis (d). Hence
$(\mathbf{U}^{\flat},\mathbf{J}^{\flat})=(U_{\mathsf{R}},0)$.

\emph{(iii-3), the orbit.} Suppose that $U_{\mathsf{R}}^{g}=e^{i\xi_kA'}\widetilde U$ on
$\operatorname{bd}(X_{\mathsf{R}})$. Here $g$ is $G^{c}$-valued, $|A'|<a_1$, and $\widetilde U$ obeys
$(\ell 2)_k(X_{\mathsf{R}};a_0)$.

\emph{Step 1: the gauge on a cube around $\Gamma$.} By (W4) there is a closed cube $\square\supset\Gamma$
of side $SL^k$. Put $\rho_S=SB_{12}a_0$. The reading $(\ell 2)$ gives a $u$ on
$\operatorname{st}(\square\cap X_{\mathsf{R}})$ and an $\widetilde A$ with
$\widetilde U^{u}=e^{i\xi_k\widetilde A}$ and $|\widetilde A|<\rho_S$. For a bond $b$ with initial site
$b_-$, writing out the gauge transform gives
\begin{equation*}
(U_{\mathsf{R}}^{g})^{u}(b)=e^{i\xi_k\operatorname{Ad}_{u(b_-)}A'(b)}\,e^{i\xi_k\widetilde A(b)} .
\end{equation*}

\emph{Step 2: each factor is close to $1$.} The value $u(b_-)$ is unitary, so conjugation by it
preserves $\|\cdot\|_{\mathrm{op}}$. Moreover $\|\cdot\|_{\mathrm{op}}\le|\cdot|$. For arbitrary
$N\times N$ matrices $\Phi_1,\Phi_2$, expanding the exponentials in power series shows that
$\|e^{\Phi_1}e^{\Phi_2}-1\|_{\mathrm{op}}$ is at most
$e^{\|\Phi_1\|_{\mathrm{op}}+\|\Phi_2\|_{\mathrm{op}}}-1$. So this matrix and its inverse are both within
$e^{\xi_k(a_1+\rho_S)}-1$ of $1$ in $\|\cdot\|_{\mathrm{op}}$.

\emph{Step 3: the product around $\Gamma$ is close to $1$.} For $N\times N$ matrices
$\Psi_1,\dots,\Psi_{n_{\Gamma}}$ we have
\begin{equation*}
\|\Psi_1\cdots\Psi_{n_{\Gamma}}-1\|_{\mathrm{op}}\le\prod_i\bigl(1+\|\Psi_i-1\|_{\mathrm{op}}\bigr)-1 .
\end{equation*}
Apply this to the product
around $\Gamma$. Use $n_{\Gamma}\xi_k\le 4S$, then \eqref{4.9:eq-ambient-condition-a}, and finally
$e^{1/2}<2$, which holds because $e<4$. This gives that the product is within
\begin{equation*}
e^{n_{\Gamma}\xi_k(a_1+\rho_S)}-1\le e^{4S(a_1+\rho_S)}-1\le e^{1/2}-1<1
\end{equation*}
of $1$.

\emph{Step 4: the contradiction.} The same product equals
$(ug)(x_0)\,T_{\Gamma}(U_{\mathsf{R}})\,(ug)(x_0)^{-1}$, where $x_0$ is the base point of $\Gamma$.
By (W2) this matrix has the eigenvalue $-1$. The operator norm dominates the spectral radius, so
this matrix lies at operator-norm distance at least $2$ from $1$. This contradicts Step 3.

\emph{(iii), the consequence.} The level is $j=k$ and the domain is $X_{\mathsf{R}}\subset\Lambda_k$.
By (iii-1) and (iii-2), $(U_{\mathsf{R}},0)$ is a source of the straddling-sources half with
$(\mathbf{U}^{\flat},\mathbf{J}^{\flat})=(U_{\mathsf{R}},0)$ at $w=0$. By (iii-3), the target
decomposition does not exist when it is read under $(\ell 2)$.

\emph{(iv).} Hypothesis (c) applies to every closed cube $\square\subset\Omega_j$ of side $CML^j$,
$C\in\{1,2\}$. For each such cube it concludes $\mathsf{g}_j(U;\square\cap X;BCM\alpha_{0,j})$.

Now let $\square$ be a closed cube of the ambient lattice of side $CML^j$.
\begin{itemize}
\item If $\operatorname{st}(\square\cap X)\neq\emptyset$, then
$\square\subset X^{+2ML^j}\subset\Omega_j$ by the margin assumption. So hypothesis (c) applies to
$\square$.
\item If $\operatorname{st}(\square\cap X)=\emptyset$, the condition on $\square$ is void.
\end{itemize}
This proves $(\ell 2)^{\flat}_j(X;\alpha_{0,j})$. The last sentence of (iv) combines this with parts (a)
and (b) of Proposition~\ref{4.9:prop-gauge-transport}, for the loci of the table, and with
Proposition~\ref{4.9:prop-real-backgrounds}, for the seeds.
\end{proof}

\begin{corollary}[A single gauge on small localisation domains]\label{4.9:cor-single-gauge-domain}
Let $\|\cdot\|$ be the norm on algebra-valued fields on bonds, $c_{\mathrm{norm}}>0$ the norm
constant and $\varpi>0$ the smallness ceiling of the gauge theorem on thin cubes. We assume the
first step of that theorem in the following form, called the gauge-fixing step in this corollary.
\begin{itemize}
\item[(GF)] Let $n$ be a level, $X$ a closed set and $\rho,a_1,\gamma_0>0$, and write
$\mathfrak{r}:=c_{\mathrm{norm}}\rho$ and $\mathfrak{a}_1:=c_{\mathrm{norm}}a_1$. Let
$\mathbf{U}=e^{i\xi_nA'}U$ and $\mathbf{J}$ be given on $\operatorname{bd}(X)$, with complex
factor $A'$ of radius $a_1$ and current $\mathbf{J}$ of radius $\gamma_0$, and let the unitary
factor $U$ satisfy the cube gauge condition $\mathsf{g}_n(U;X;\rho)$ of
Definition~\ref{4.9:def-cube-gauge-condition} with a gauge $u$. Assume
$\xi_n(\mathfrak{a}_1+\mathfrak{r})\le\varpi$. Take the construction of the first step of the
gauge theorem on thin cubes with the trivial initial gauge, with second gauge $u$, and with the
one-bond covariant difference of \cite[(3.3), (3.37)]{B-CMP99b}, whose length is one bond and
whose amplitude is $\mathfrak{d}:=\mathfrak{a}_1$. Then the three estimates of that step hold,
the first of them with right member $\xi_n(\mathfrak{a}_1+2\mathfrak{r}\mathfrak{a}_1)$, and
together they give a function $\widetilde{A}$ on $\operatorname{bd}(X)$ with
$\mathbf{U}^{u}=e^{i\xi_n\widetilde{A}}$ on $\operatorname{bd}(X)$ and
\begin{equation}\label{4.9:eq-single-gauge-domain-1}
\begin{aligned}
\|\widetilde{A}\| &\le \tfrac{9}{8}\,(\mathfrak{a}_1+\mathfrak{r}),\\
\xi_n^{-1}\,\|\widetilde{A}(b+e_\nu)-\widetilde{A}(b)\|
&\le \tfrac{7}{4}\,(\mathfrak{a}_1+2\mathfrak{r}\mathfrak{a}_1+\mathfrak{r})
\quad\text{on } \operatorname{pp}(X).
\end{aligned}
\end{equation}
\end{itemize}

Let $j\le n$ be levels and let $X\in\mathbf{D}_j$. Let $\square$ be a closed cube of side at
most $2ML^n$ with
$X\subset\square$, so that $X$ has diameter at most $2M$ in the units of level $n$. Let
$Y\supset X$, and let $(\mathbf{U},\mathbf{J})=(e^{i\xi_nA'}U,\mathbf{J})$ on $Y$ be a
representative, at level $n$ and with radii $(a_0,a_1,\gamma_0)$ (the amplitudes of the unitary
factor $U$, of the complex factor $A'$ and of the current $\mathbf{J}$), of one of the following:
\begin{itemize}
\item[$(\alpha)$] the one-power hereditary class $U^{\Sigma^1}_n(Y;\cdot)$, read with the reading
$(\ell 2)^{\flat}$ of Definition~\ref{4.9:def-cube-gauge-condition};
\item[$(\beta)$] a space of configurations with a nest of small-field regions as in
\cite[(2.38)]{B-CMP119} or \cite[(1.65), (1.67), (1.69)]{B-CMP122a}, whose nest
$\{\Omega_i\}_i$ satisfies $\square\subset\Omega_n$, to which clause (ii) of
Proposition~\ref{4.9:prop-large-field-cubes} applies.
\end{itemize}
Put $P:=(\mathbf{U},\mathbf{J})|_{X}$. Set
\begin{equation*}
\rho:=2MBa_0,\qquad \mathfrak{r}:=c_{\mathrm{norm}}\rho,\qquad \mathfrak{a}_1:=c_{\mathrm{norm}}a_1,
\end{equation*}
and assume $\xi_n(\mathfrak{a}_1+\mathfrak{r})\le\varpi$.

Then there are a single $G$-valued function $u$ on $\operatorname{st}(X)$ and a
function $\widetilde{A}$ on $\operatorname{bd}(X)$ such that
$\mathbf{U}^{u}=e^{i\xi_n\widetilde{A}}$ on $\operatorname{bd}(X)$,
$|\operatorname{Ad}_u\mathbf{J}|<\gamma_0$, and
\begin{align*}
\|\widetilde{A}\| &\le \tfrac{9}{8}\,(\mathfrak{a}_1+\mathfrak{r}),\\
\xi_n^{-1}\,\|\widetilde{A}(b+e_\nu)-\widetilde{A}(b)\|
&\le \tfrac{7}{4}\,(\mathfrak{a}_1+2\mathfrak{r}\mathfrak{a}_1+\mathfrak{r})
\quad\text{on } \operatorname{pp}(X).
\end{align*}
Consequently $P$ satisfies, in the units of level $n$ and with the potential and its differences
measured in the norm $\|\cdot\|$, the gauge hypothesis of the boundary lemma of the large-field
bounds for
domains of diameter at most $2M$ at level $n$. The amplitude of that hypothesis is
\begin{equation*}
\mathfrak{a}:=\max\bigl\{\tfrac{9}{8}(\mathfrak{a}_1+\mathfrak{r}),\;
\tfrac{7}{4}(\mathfrak{a}_1+2\mathfrak{r}\mathfrak{a}_1+\mathfrak{r}),\;\gamma_0\bigr\}.
\end{equation*}
\end{corollary}

\begin{proof}
Let $U$ be the unitary factor of the representative. We first show that in either case $U$
satisfies the cube gauge condition $\mathsf{g}_n(U;\square\cap Y;\rho)$ of
Definition~\ref{4.9:def-cube-gauge-condition}, with $\rho=2MBa_0$. In case $(\alpha)$ this is
the reading $(\ell 2)^{\flat}$ of the class at level $n$, applied to the cube $\square$. In
case $(\beta)$ it is clause (ii) of Proposition~\ref{4.9:prop-large-field-cubes}, which applies
because $\square\subset\Omega_n$.

Next we restrict to $X$. Since $X\subset\square$ and $X\subset Y$, we have
\begin{equation*}
X=\square\cap X\subset\square\cap Y .
\end{equation*}
Lemma~\ref{4.9:lem-restriction-monotone}, applied with $D=\square\cap Y$, $D'=X$ and
$\rho'=\rho$, therefore gives $\mathsf{g}_n(U;X;\rho)$. By
Definition~\ref{4.9:def-cube-gauge-condition} this means the following: there are
$u\colon\operatorname{st}(X)\to G$ and $A\colon\operatorname{bd}(X)\to\mathfrak{g}$ with
$U^{u}=e^{i\xi_nA}$ on $\operatorname{bd}(X)$, with $|A(b)|<\rho$ on $\operatorname{bd}(X)$,
and with $\xi_n^{-1}|A(b+e_\nu)-A(b)|<\rho$ on $\operatorname{pp}(X)$.

We now apply the gauge-fixing step (GF) to the restriction of the representative to $X$, with
$\rho=2MBa_0$, with the radii $a_1$ and $\gamma_0$ of the representative, and with the following
choices:
\begin{itemize}
\item the trivial initial gauge;
\item as its second gauge, the gauge $u$ just obtained from $\mathsf{g}_n(U;X;\rho)$;
\item as difference operator, the one-bond covariant difference of \cite[(3.3), (3.37)]{B-CMP99b},
so that the difference has length one bond and its amplitude is $\mathfrak{d}:=\mathfrak{a}_1$.
\end{itemize}
The amplitudes entering the step are $\mathfrak{r}=c_{\mathrm{norm}}\rho$ for the unitary factor
and $\mathfrak{a}_1=c_{\mathrm{norm}}a_1$ for the complex factor $e^{i\xi_nA'}$. Its smallness
condition is the assumption $\xi_n(\mathfrak{a}_1+\mathfrak{r})\le\varpi$.

With these choices all hypotheses of (GF) hold: $U$ satisfies $\mathsf{g}_n(U;X;\rho)$ with the
gauge $u$, the factors $A'$ and $\mathbf{J}$ have the radii $a_1$ and $\gamma_0$ of the
representative, and the smallness condition is assumed. Hence the first of the three estimates
of that step has right member $\xi_n(\mathfrak{a}_1+2\mathfrak{r}\mathfrak{a}_1)$, and the step
gives, in the gauge $u$, a function $\widetilde{A}$ on $\operatorname{bd}(X)$ with
$\mathbf{U}^{u}=e^{i\xi_n\widetilde{A}}$ and the two bounds
\eqref{4.9:eq-single-gauge-domain-1}, which are the two bounds of the conclusion. Moreover,
$\mathbf{J}$ has radius $\gamma_0$, so $|\mathbf{J}|<\gamma_0$; since $u$ is $G$-valued,
$\operatorname{Ad}_u$ conjugates each value of $\mathbf{J}$ by a unitary matrix, which leaves the
norm $|\cdot|$ unchanged, so that $|\operatorname{Ad}_u\mathbf{J}|=|\mathbf{J}|<\gamma_0$. The
gauge $u$ of the cube gauge condition is thus the single gauge of the statement.

Each of the three quantities $\|\widetilde{A}\|$,
$\xi_n^{-1}\|\widetilde{A}(b+e_\nu)-\widetilde{A}(b)\|$ and $|\operatorname{Ad}_u\mathbf{J}|$ is
therefore at most $\mathfrak{a}$. This is the gauge hypothesis of the boundary lemma of the
large-field bounds, in the units of level $n$, with the potential and its differences measured in
the norm $\|\cdot\|$.
\end{proof}

For domains that do not lie in a cube of side at most $2ML^n$ the corollary is not used, the
trivial decay bound of clause (ii) of the large-field decay lemma serving instead, and the
gauge $u$ is asserted only to exist, no particular choice of it being fixed.

\emph{The families of cubes.} Throughout, $\bar U$ denotes the unitary factor of the configuration on
which the large-field operation $\mathbf{R}$ acts; we call it the input unitary factor. Lengths are
measured in lattice steps, so that a cube of side $ML^m$ is a cube of side $M$ in the units of
level $m$. We write $\operatorname{dist}_\infty$ for the distance in the supremum norm and
$\Omega^c$ for the complement of a set $\Omega$.

In \cite[(1.64)--(1.69)]{B-CMP122a} and \cite[(1.65)]{B-CMP122b}, a gauge of $\bar U$ is
required on cubes of five families:
\begin{itemize}
\item[(c1)] The stage families. Let $\mathfrak{U}_n(X)=\widetilde{U}^{(n)c}_k(X,\tilde\alpha_0,
\tilde\alpha_1)$ be the spaces of complexified configurations of the stages of the large-field
operation, with layer radii $\tilde\alpha_{0,m}$, $\tilde\alpha_{1,m}$. For every cube $\square$ of
side $CML^m$ with $\square\subset\Omega_m$ and, if $m<k$, $\square\cap\Omega^c_{m+1}\ne\emptyset$
(a cube attached to the shell $\Omega_m\setminus\Omega_{m+1}$ of index $m$) there
is a $G$-valued $u$ on $\square\cap X$ with
\begin{equation}\label{4.9:eq-large-field-cubes-b}
\bar U^{u}=e^{i\eta A},\qquad
\ell_m|A|,\ \ell_m^2|\nabla^{\eta}A|<\lambda(m)\,BCM\,\tilde\alpha_{0,m}.
\end{equation}
Here $\eta$ is the spacing of the finest lattice in the units of the top level, and
$\ell_m=L^m\eta$. Thus the side $CML^m$ equals $C\ell_mM$ in the units of the top level.
\item[(c2)] The old operations $\mathbf{B}^{(j)}$: the condition
\eqref{4.9:eq-cube-gauge-condition-c} with respect to the old nest.
\item[(c3)] The old operations $\mathbf{E}^{(j)}$, $\mathbf{R}^{(j)}$: the widened reading
$(\mathrm{L}\text{-}12)_j(X';\cdot)$ of \eqref{4.9:eq-cube-gauge-condition-b}, at level $j$ on
$X'$.
\item[(c4)] The cut-off cube: one cube, attached to some level $n$, of side at most $2ML^n$
(that is, at most $2M$ in the units of level $n$), containing $X$.
\item[(c5)] The hulls and local cubes of the operators, of sides at most $17\mathfrak{M}\ell_m$
and $M/12$; these are sub-cubes of the cubes of the preceding families.
\end{itemize}
In (c1), $C$ is a positive integer, which the text following \cite[(1.69)]{B-CMP122a} allows one
to take at most $2$. Further, $B$ is an absolute constant, for which that text proposes $B_3$; in
Proposition~\ref{4.9:prop-large-field-cubes} below, $B$ is subject to the floors of hypothesis (c).
The gain $\lambda(m)$ is one of $L_\ast^{2\max\{0,m-k_0\}}$, $L_\ast^{2(j-m)}$ and $1$. Here
$L_\ast$ is the constant written $L_0$ in \cite[(1.65), (1.67)]{B-CMP122a}, and $k_0$, $j$ are the
indices so named there. In particular $\lambda(m)\ge1$. For $k_0<m<k$ the shell
$\Omega_m\setminus\Omega_{m+1}$ is one layer of $M$-cubes of its level thick, that is, of
thickness $ML^m$ in lattice steps \cite{B-CMP122a}. The configuration on which
\cite[(1.65)]{B-CMP122b} is imposed lies in the standard space of
\cite[(2.34)--(2.39)]{B-CMP119} with respect to the new nest of small-field regions. That nest
is defined on the domain $Y_4$ of that standard space, and the new level is $k$ on the union of the
large-field region $Z$ and the top region $\Omega_k^{\mathrm{old}}$ of the old nest.

\begin{proposition}[Cube conditions inside the large-field spaces]\label{4.9:prop-large-field-cubes}
Let $\sigma=\{\Omega_n\}_{n\le k}$ be a nest of small-field regions, so that
$\Omega_{n+1}\subset\Omega_n$. Let $(a_n)_{n\le k}$ be positive radii. Assume that every
$\Omega_n$ is a union of closed cubes of a cubic partition $\mathcal{P}_n$ of mesh
$\mathsf{K}_n$; that is, $\mathcal{P}_n$ is compatible with $\Omega_n$. Assume further that
\begin{equation}\label{4.9:eq-large-field-cubes-a}
\begin{gathered}
(\mathrm{N\text{-}par})\quad \mathsf{K}_n\in ML^n\mathbb{N};\\
(\mathrm{adm})^{\flat}\quad\operatorname{dist}_\infty(\Omega^c_{n+1},\Omega_{n+2})>2ML^n;\\
(\mathrm{rat})\quad L^{-(n-m)}a_n\le a_m\quad\text{for }m\le n\le k.
\end{gathered}
\end{equation}
Here $(\mathrm{N\text{-}par})$ stands for the partition property together with the mesh condition.
It is required for every $n\le k$, and $(\mathrm{adm})^{\flat}$ for every $n$ with $n+2\le k$.

Let $Y$ be closed, and let $\bar U$ be $G$-valued on $\operatorname{bd}(Y)$ and satisfy
\eqref{4.9:eq-cube-gauge-condition-c} for the nest $\sigma$, the set $Y$ and the radii $(a_n)$.
Assume moreover:
\begin{enumerate}
\item[(a)] The real backgrounds of the stages, namely $U^{(n)}_k$, $U''_k$, $U_{k,\Lambda}$, and
the real background $U_2$ outside the band $\mathfrak{K}_c$, satisfy their layerwise regularity
conditions $\mathrm{Reg}_n[\cdot\,;\cdot]$ (regularity at scale $n$) of
\eqref{4.9:eq-collar-regularity-a}, with constants at most the radii of the corresponding layers.
\item[(b)] The composites formed by the large-field step keep the input unitary factor, with one
exception. That is, every composite of the step has
unitary factor $\bar U$, both in the direct form of the composites of the step (each composite
written as one map from the input space of the step to the domain on which a cube condition is
read) and along the
nesting $\mathfrak{U}_{n+1}\to\mathfrak{U}_n$ of the stages. The exception is the one composite
started from the real background $U_2$ outside the band $\mathfrak{K}_c$, whose unitary factor
is that of $U_2$. This is what the statement on the composites of the large-field step in direct
form and the statement on the nesting of the stage spaces assert.
\item[(c)] The constant $B$ of \eqref{4.9:eq-cube-gauge-condition-c} satisfies the floors
\eqref{4.9:eq-real-backgrounds-b}, in particular
\begin{equation*}
B\ge\max\{14dL^5B_1,\ \pi L^{m_0}(d-1)\kappa_G\},
\end{equation*}
where $m_0\in\{0,1\}$ is the lowest index of the nest; and $M$ satisfies the floor
\eqref{4.9:eq-real-backgrounds-d}.
\end{enumerate}
We say that a nest $\{\Omega'_n\}_{n\le k}$ is dominated by $\sigma$ if $\Omega'_n\subset\Omega_n$
for every $n\le k$; calling the level of a nest at a point $x$ the largest index $n$ with $x$ in
its $n$-th member, this means that its level nowhere exceeds that of $\sigma$.
Then the following hold.
\begin{enumerate}
\item[(i)] The conditions \cite[(2.38)]{B-CMP119} and \cite[(1.65), (1.67), (1.69)]{B-CMP122a}
are, by their definitions, clauses $\mathsf{g}_m(\bar U;\square\cap X;\cdot)$ of
Definition~\ref{4.9:def-cube-gauge-condition}, read with the gradient member
$(L^m\eta)^2|\nabla^{\eta}A|$. They are imposed on cubes $\square$ that lie in the small-field
region $\Omega_m$ of their index $m$ and are not required to lie in $X$.
\item[(ii)] For every $m\le k$ and every closed cube $\square\subset\Omega_m$ of side $CML^m$
with $C\in\{1,2\}$, the clause $\mathsf{g}_m(\bar U;\square\cap Y;BCMa_m)$ holds. For every
closed cube $\square\subset\Omega_m$ of side at most $2ML^m$, the clause
$\mathsf{g}_m(\bar U;\square\cap Y;2BMa_m)$ holds.
\item[(iii)] Let $\square$ be a closed cube of side $CML^m$, $C\in\{1,2\}$, with
$\square\subset\Omega_m$. Then, for every $X'\subset Y$, every gain $\lambda\ge1$ and every
target radius $a^{\tau}_m\ge a_m$,
\begin{equation*}
\mathsf{g}_m(\bar U;\square\cap X';\lambda BCMa^{\tau}_m)
\end{equation*}
holds. Here a target family is a family of cubes on which a space entered after the step imposes
a cube condition, and $a^{\tau}_m$ are its radii. Thus the cube condition of a target family holds
for the same $\bar U$, with no perturbation of $\bar U$.

The inclusion $\square\subset\Omega_m$ holds for every cube of level $m$ of a family attached to a
nest dominated by $\sigma$. When $\sigma$ is the new nest of the step, the stage nests and the old
nest are dominated by $\sigma$.
\item[(iv)] Let $j\le k$, let $X'\subset Y\cap\Omega_j$ be closed, and assume the hypotheses of
Corollary~\ref{4.9:cor-gluing-smaller-cubes} with $a=a_j$, among them
\eqref{4.9:eq-gluing-smaller-cubes-a} at $a=a_j$. Then $\bar U|_{X'}$ satisfies
$(\ell 1)_j(X';a_j)$, which is $(\mathrm{L}\text{-}12)_j(X';a_j)$. If moreover
$X'^{+2ML^j}\subset\Omega_j$, then it also satisfies $(\ell 2)^{\flat}_j(X';a_j)$.
\item[(v)] The real backgrounds listed in (a) satisfy \eqref{4.9:eq-cube-gauge-condition-c},
respectively the condition of (c1), on their own nests, thin shells included.
\item[(vi)] Under hypotheses (a), (b) and (c), for the regions $X'$ of family (c3) subject to the
hypotheses of (iv), and for the families (c1), (c2), (c4) and (c5) attached to nests dominated by
$\sigma$, with domains contained in $Y$ and with radii at each level $m$ not below $a_m$, the gains
being at least $1$ (so that the radius of the clause on a cube of side $CML^m$ is
$\lambda BCMa^{\tau}_m$ with $\lambda\ge1$ and $a^{\tau}_m\ge a_m$, and that on the cube of (c4),
attached to level $n$, is not below $2BMa_n$), every cube
condition that the large-field step imposes is contained either in the condition
\eqref{4.9:eq-cube-gauge-condition-c} satisfied by $\bar U$ or, for the composite started from
$U_2$, in the condition that $U_2$ satisfies by (v); in either case it has, by (i), the form of
\cite[(1.69)]{B-CMP122a}. Hence the step requires no cube condition that is not already
satisfied. The condition \eqref{4.9:eq-cube-gauge-condition-c} does not contain the ambient
reading $(\ell 2)$ of \eqref{4.9:eq-cube-gauge-condition-b} on cubes of side greater than
$2ML^j$ and at most $c_{12}LML^j$ at level $j$, with $c_{12}$ the constant written $O(1)$ in
\cite[(1.12)]{B-CMP109}. None of the families (c1)--(c5) requires it.
\end{enumerate}
\end{proposition}

\begin{proof}
\emph{Proof of (i).} This is the content of the definitions. The conditions
\cite[(2.38)]{B-CMP119} and \cite[(1.64)--(1.69), pp.~190--191]{B-CMP122a} ask the following
for every cube $\square$ of the family they describe: there is a gauge transformation $u$ defined
on $\square\cap X$, and a $\mathfrak{g}$-valued $A$ with $\bar U^u=\exp(i\eta A)$, such that
$L^m\eta|A|$ and $(L^m\eta)^2|\nabla^\eta A|$ are below the radius attached to the index $m$ of
the cube. The cubes of the family are described through the nest alone, by
$\square\subset\Omega_m$, together with $\square\cap\Omega^c_{m+1}\ne\emptyset$ if $m<k$. They
are not described through $X$.

\emph{Proof of (ii).} Put $\Omega_{k+1}:=\emptyset$. Let $\square\subset\Omega_m$ be a closed
cube. Since the $\Omega_n$ decrease, $\Omega_m$ is the union of the sets
$\Omega_n\setminus\Omega_{n+1}$, $m\le n\le k$. Hence
\begin{equation*}
n_0:=\min\{n\ge m:\ \square\cap(\Omega_n\setminus\Omega_{n+1})\ne\emptyset\}
\end{equation*}
exists and $n_0\le k$. For $m\le n<n_0$, the inclusion $\square\subset\Omega_n$ and
$\square\cap(\Omega_n\setminus\Omega_{n+1})=\emptyset$ give $\square\subset\Omega_{n+1}$.
Therefore $\square\subset\Omega_{n_0}$, and $\square\cap\Omega^c_{n_0+1}\ne\emptyset$.

We call a closed cube $\square$ of side $CML^n$, $C\in\{1,2\}$, a cube of index $n$ if
$\square\subset\Omega_n$ and, in case $n<k$, also $\square\cap\Omega^c_{n+1}\ne\emptyset$ and
$\square\cap\Omega_{n+2}=\emptyset$. Every cube of index $n$ belongs to the family on which
\eqref{4.9:eq-cube-gauge-condition-c} imposes its clause at index $n$, since that family is
described, as recalled in the proof of (i), by $\square\subset\Omega_n$ together with
$\square\cap\Omega^c_{n+1}\ne\emptyset$ if $n<k$. On such cubes the hypothesis
\eqref{4.9:eq-cube-gauge-condition-c} therefore gives $\mathsf{g}_n(\bar U;\square\cap Y;BCMa_n)$.

Case $n_0=m$, side $CML^m$. Then $\square\subset\Omega_m$. If $m=k$, $\square$ is a cube of
index $k$. If $m<k$, then $\square\cap\Omega^c_{m+1}\ne\emptyset$. Every point of $\square$ lies
within sup-distance $CML^m\le2ML^m$ of a point of that intersection. A point of
$\square\cap\Omega_{m+2}$ would therefore contradict $(\mathrm{adm})^{\flat}$ at $n=m$ (if
$m=k-1$, then $\Omega_{m+2}=\Omega_{k+1}=\emptyset$ and there is no such point). So
$\square\cap\Omega_{m+2}=\emptyset$, $\square$ is a cube of index $m$, and
$\mathsf{g}_m(\bar U;\square\cap Y;BCMa_m)$ holds.

Case $n_0>m$, side $s\le2ML^m$. Since $L\ge2$ and $\mathsf{K}_{n_0}\in ML^{n_0}\mathbb{N}$,
\begin{equation*}
s\le 2ML^m\le ML^{n_0}\le\mathsf{K}_{n_0}.
\end{equation*}
Apply Lemma~\ref{4.9:lem-cube-prolongation} in $\Omega_{n_0}$, with the partition
$\mathcal{P}_{n_0}$, mesh $\mathsf{K}_{n_0}$ and $t:=ML^{n_0}$. Let $q$ be the number defined in
that lemma for the cube $\square$ and the partition $\mathcal{P}_{n_0}$. Since $q\ge1$, this value
of $t$ lies in $[s,q\mathsf{K}_{n_0}]$. The lemma gives a closed cube $\square'$ of side
$ML^{n_0}$ with $\square\subset\square'\subset\Omega_{n_0}$. If $n_0<k$, then
\begin{equation*}
\square'\cap\Omega^c_{n_0+1}\supset\square\cap\Omega^c_{n_0+1}\ne\emptyset,
\end{equation*}
and $\square'\cap\Omega_{n_0+2}=\emptyset$ by $(\mathrm{adm})^{\flat}$ at $n=n_0$, as in the
first case (if $n_0=k-1$, because $\Omega_{n_0+2}=\Omega_{k+1}=\emptyset$). If $n_0=k$, the
inclusion $\square'\subset\Omega_k$ suffices. Either way $\square'$
is a cube of index $n_0$ with $C=1$, and $\mathsf{g}_{n_0}(\bar U;\square'\cap Y;BMa_{n_0})$
holds.

Since $\square\cap Y\subset\square'\cap Y$, Lemma~\ref{4.9:lem-restriction-monotone} gives
$\mathsf{g}_{n_0}(\bar U;\square\cap Y;BMa_{n_0})$. Lemma~\ref{4.9:lem-level-change}, with
$i=n_0$ and $i'=m$, then gives
\begin{equation*}
\mathsf{g}_m\bigl(\bar U;\square\cap Y;L^{-(n_0-m)}BMa_{n_0}\bigr).
\end{equation*}
By $(\mathrm{rat})$ the radius is at most $BMa_m$. Since $BMa_m\le BCMa_m$ and $BMa_m\le2BMa_m$,
Lemma~\ref{4.9:lem-restriction-monotone} raises the radius to either value.

This proves the first assertion of (ii), and the second when $n_0>m$. It remains to treat a cube
$\square\subset\Omega_m$ of side $s\le2ML^m$ with $n_0=m$. We prolong it by
Lemma~\ref{4.9:lem-cube-prolongation} in $\Omega_m$, with partition $\mathcal{P}_m$ and mesh
$\mathsf{K}_m$.
\begin{itemize}
\item If $s\le ML^m$, take $t:=ML^m$. Then $t\le\mathsf{K}_m\le q\mathsf{K}_m$.
\item If $s>ML^m$, take $t:=2ML^m\ge s$. Either $\mathsf{K}_m\ge2ML^m$, and then
$t\le q\mathsf{K}_m$. Or $\mathsf{K}_m=ML^m<s$; then $q\ge s/\mathsf{K}_m>1$, so $q\ge2$ and
$q\mathsf{K}_m\ge t$.
\end{itemize}
The resulting cube $\square'\supset\square$ lies in $\Omega_m$ and has side $CML^m$ with
$C\in\{1,2\}$. It meets $\Omega_m\setminus\Omega_{m+1}$, because $\square$ does, so its $n_0$ is
$m$. The first case gives $\mathsf{g}_m(\bar U;\square'\cap Y;BCMa_m)$. Since
$BCMa_m\le2BMa_m$, Lemma~\ref{4.9:lem-restriction-monotone} gives
$\mathsf{g}_m(\bar U;\square\cap Y;2BMa_m)$.

\emph{Proof of (iii).} By (ii), $\mathsf{g}_m(\bar U;\square\cap Y;BCMa_m)$ holds. Since
$X'\subset Y$, we have $\square\cap X'\subset\square\cap Y$. Since $\lambda\ge1$ and
$a^\tau_m\ge a_m$, we have $BCMa_m\le\lambda BCMa^\tau_m$.
Lemma~\ref{4.9:lem-restriction-monotone} now gives the assertion.

For the second paragraph of (iii), let $\{\Omega'_n\}_{n\le k}$ be dominated by $\sigma$. By (i)
a cube of level $m$ of a family attached to it lies in $\Omega'_m$, hence in $\Omega_m$. When
$\sigma$ is the new nest of the step, the levels of the large-field step are ordered, by the
statement on the ordering of the levels of the large-field step: the new
level is nowhere below the stage levels, which are nowhere
below the old level, and the region $Z\cup\Omega_k^{\mathrm{old}}$, on which the new level
is $k$, lies in $\Omega_n$ for every $n\le k$; so the stage nests and the old nest are dominated
by $\sigma$.

\emph{Proof of (iv).} Let $\square^\ast\subset X'$ be a closed cube. Since $X'\subset\Omega_j$,
every closed cube of side $ML^j$ contained in $\square^\ast$ lies in $\Omega_j$. The proof of
Corollary~\ref{4.9:cor-gluing-smaller-cubes} uses the cube condition, in its two steps, on cubes
of side $ML^j$: in the first step on a cube $\square\supset\square^\ast$, and in the second, for
each of the covering cubes $c$ of that proof, on a cube $\square_c\supset c$. The cubes
$\square^\ast$ and $c$ of these two steps are contained in the cubes $\square$ and $\square_c$ of
side $ML^j$, so their side $s$ is at most $ML^j$, which is the condition $s\le t$ of
Lemma~\ref{4.9:lem-cube-prolongation} for $t:=ML^j$; and they lie in $X'\subset\Omega_j$. Since
$q\ge1$, the value $t:=ML^j$ lies in $[s,\mathsf{K}_j]\subset[s,q\mathsf{K}_j]$, and
Lemma~\ref{4.9:lem-cube-prolongation}, applied in $\Omega_j$ with the partition
$\mathcal{P}_j$, places the cubes $\square\supset\square^\ast$ and $\square_c\supset c$ of side
$ML^j$ inside $\Omega_j$. On each of them, (ii) with $m=j$ and $C=1$ gives the clause
with radius $MBa_j$. These are the inputs of the proof of
Corollary~\ref{4.9:cor-gluing-smaller-cubes}, which therefore yields $(\ell 1)_j(X';a_j)$ under
its hypotheses with $a=a_j$. If $X'^{+2ML^j}\subset\Omega_j$, the cubes on which the ambient
reading $(\ell 2)^{\flat}_j(X';a_j)$ imposes the clause lie in $\Omega_j$. On them (ii) gives the
clause, restricted to their intersections with $X'$ by Lemma~\ref{4.9:lem-restriction-monotone}.

\emph{Proof of (v).} This is the clause of Proposition~\ref{4.9:prop-real-backgrounds} for
general nests, together with its variant for thin shells under
\eqref{4.9:eq-real-backgrounds-d}, applied to each of the real backgrounds listed in (a) on its
own nest. The regularity constants of each of these
backgrounds are at most the radii by hypothesis (a). By hypothesis (c), the constant $B$ satisfies
the floors \eqref{4.9:eq-real-backgrounds-b}, and $M$ satisfies
\eqref{4.9:eq-real-backgrounds-d}, under which the thin shells are covered.

\emph{Proof of (vi).} By hypothesis (b), the unitary factor of each composite of the step other
than the one started from $U_2$ is $\bar U$. This is the entry of
\eqref{4.9:eq-gauge-transport-a} in Proposition~\ref{4.9:prop-gauge-transport} concerning the
composites of the large-field step in direct form.
The cube condition of each such composite is
therefore a condition on $\bar U$, and no perturbation of $\bar U$ enters. The one place where
the unitary factor is not $\bar U$ is the real background $U_2$, and its condition holds by (v).
For the families (c1) and (c2), attached to nests dominated by $\sigma$, every cube of level $m$
of side $CML^m$, $C\in\{1,2\}$, lies in $\Omega_m$ by the second paragraph of (iii). The domain
$X'$ of the family lies in $Y$, and its radius at level $m$ is some $a^{\tau}_m\ge a_m$, with a gain
$\lambda\ge1$, by the hypotheses of (vi); for (c1), the gain is $\lambda(m)\ge1$ by the remark
following the list of families, the radius is $\tilde\alpha_{0,m}\ge a_m$ and the domain is
$X\subset Y$, which are the hypotheses just named. The first paragraph of (iii) then gives its
clause on $\square\cap X'$ with the gain and the radius of the family. The cube of (c4), of side
at most $2ML^n$ and attached to level $n$ of a
nest dominated by $\sigma$, lies in $\Omega_n$ in the same way; the second assertion of (ii)
gives the clause on its intersection with $Y$ with radius $2BMa_n$, and
Lemma~\ref{4.9:lem-restriction-monotone} passes it to its intersection with every
$X'\subset Y$ and to its radius, which is not below $2BMa_n$ by the hypotheses of (vi). The cubes
of (c5) are sub-cubes of cubes of
the preceding families, and Lemma~\ref{4.9:lem-restriction-monotone} passes the clause to them.
For (c3), with $X'$ subject to the hypotheses of (iv), the condition is the conclusion of (iv).
Both (iii) and (iv), and the treatment of (c4) and (c5), derive their conclusions from
\eqref{4.9:eq-cube-gauge-condition-c} through (ii).

It remains to read the last two sentences of (vi). As recalled in the proof of (i), the condition
\eqref{4.9:eq-cube-gauge-condition-c} imposes its clause at index $m$ only on cubes of side
$CML^m$ with $C\in\{1,2\}$, that is, of side at most $2ML^m$; it imposes none on cubes of side
greater than $2ML^j$ at level $j$, and in this sense it does not contain $(\ell 2)$ on those cubes.
By the list of families, (c1), (c2) and (c4) consist of cubes of side at most $2ML^m$ at their
level $m$, the cubes of (c5) are sub-cubes of cubes of these families, and (c3) asks the reading
$(\ell 1)$, whose cubes lie in $X'$, and not the ambient reading $(\ell 2)$. Hence none of the
families requires $(\ell 2)$ on cubes of side greater than $2ML^j$. This proves (vi).
\end{proof}

\emph{Sufficient conditions for $(\mathrm{adm})^{\flat}$ and $(\mathrm{rat})$.} Condition
$(\mathrm{adm})^{\flat}$ holds when $L>2$, both for the nests made by the operation $\mathbf{T}$
and at the thin shells of (c1). For the nests made by $\mathbf{T}$ the distance exceeds
$RML^{n+1}$ by \cite[(2.2)]{B-CMP96}, with $R$ the constant appearing there, and
$RML^{n+1}\ge2ML^n$ when $R\ge1$ and $L>2$; at the thin shells
the distance equals $ML^{n+1}$, which exceeds $2ML^n$ when $L>2$. Condition $(\mathrm{rat})$
holds when $L\ge\sqrt{2}(1+\beta_0)$ and the $a_n$ are the radii $\alpha_{0,n}$, or the radii
$\alpha_{1,n}$, along the coupling flow, as bounded by the ratio of radii along the coupling flow
(Lemma~\ref{4.4:lem-radii-ratio}). The condition $L\ge\sqrt{2}(1+\beta_0)$ is contained in the
floor $L\ge2\sqrt{2}(1+\beta)(1+\beta_0)$ on $L$.

We verify the statement on $(\mathrm{rat})$. We use the ratio of radii along the coupling flow
(Lemma~\ref{4.4:lem-radii-ratio}), for either family of radii:
\begin{equation}\label{4.9:eq-large-field-cubes-1}
\frac{\alpha_{\cdot,n'}}{\alpha_{\cdot,n}}\le(1+\beta_0)(1+n'-n)^{1/2}
\qquad\text{for } n\le n'.
\end{equation}
Taking $n'=n$ gives $1\le 1+\beta_0$. Let $m\le n\le k$. If $n=m$, then $(\mathrm{rat})$ is an
equality. If $n-m\ge1$, then \eqref{4.9:eq-large-field-cubes-1} gives
$a_n\le(1+\beta_0)(1+n-m)^{1/2}a_m$. On the other hand, $L\ge\sqrt2(1+\beta_0)$, together with
$1+\beta_0\ge1$ and $2^{n-m}\ge1+n-m$, gives
\begin{equation*}
L^{n-m}\ge 2^{(n-m)/2}(1+\beta_0)^{n-m}\ge(1+\beta_0)(1+n-m)^{1/2}.
\end{equation*}
Hence $L^{-(n-m)}a_n\le a_m$.

\begin{remark}[Arithmetic of the stage nesting]\label{4.9:rem-stage-nesting}
Let $L_{\ast}$ denote the constant of \cite[(1.65), (1.67)]{B-CMP122a}, let $\alpha_{0,j}$ be the
radius at level $j$, and let $U$ be any configuration as in Lemma~\ref{4.9:lem-level-change}.
Consider a cube $\square$ of the family of \cite[(1.67)]{B-CMP122a} at stage $n$.
It lies at level $j$ in $\Omega_m$ and carries the gain $L_{\ast}^{2(j-m)}$, so that its gauge clause is
$\mathsf{g}_j(U;\square;L_{\ast}^{2(j-m)}\alpha_{0,j})$. Suppose that $\square$ is contained in a cube
$\square_{1}$ of the same family at stage $n+1$. The cube $\square_{1}$ lies at level $j+1$ and carries
the gain $L_{\ast}^{2(j+1-m)}$, so that its gauge clause is
$\mathsf{g}_{j+1}(U;\square_{1};L_{\ast}^{2(j+1-m)}\alpha_{0,j+1})$.

Then the clause of $\square_{1}$ implies the clause of $\square$, with the same gauge, provided
\begin{equation}\label{4.9:eq-stage-nesting-1}
L^{-1}L_{\ast}^{2(j+1-m)}\alpha_{0,j+1}\le L_{\ast}^{2(j-m)}\alpha_{0,j}.
\end{equation}
By the bound on the ratio of radii along the coupling flow (Lemma~\ref{4.4:lem-radii-ratio}), the
consecutive radii satisfy $\alpha_{0,j+1}\le\sqrt{2}\,(1+\beta_0)\,\alpha_{0,j}$. Hence
\eqref{4.9:eq-stage-nesting-1} holds as soon as
\begin{equation}\label{4.9:eq-stage-nesting-2}
\sqrt{2}\,(1+\beta_0)\,L_{\ast}^{2}\le L .
\end{equation}
Up to the ratio of radii, this is Ba{\l}aban's condition $L_{\ast}^{2}\le\frac13 L$ of
\cite{B-CMP122a}.

Now take the class of configurations in which the gauge clause is required on each cube separately.
For that class, the membership statement needed is clause~(ii) of
Proposition~\ref{4.9:prop-large-field-cubes}, and that clause is obtained without aligning the cube
with the partition.

The remark concerns the gauge clause only. It does not treat the statements that use the band
$\mathfrak{K}_c$, nor Ba{\l}aban's re-minimised representative and the uses made of it, nor the
condition to which that representative is carried. Once the
unitary factor $\bar U$ of Proposition~\ref{4.9:prop-large-field-cubes} is kept, none of them is a
statement about the gauge clause.
\end{remark}

\begin{proof}
Apply the change of level lemma (Lemma~\ref{4.9:lem-level-change}) to the clause of $\square_{1}$. We
take $i=j+1$ and $i'=j$, so that $\mu=L^{-(i-i')}=L^{-1}$, and
$\rho=L_{\ast}^{2(j+1-m)}\alpha_{0,j+1}$. The lemma gives
\begin{equation*}
\mathsf{g}_{j+1}\bigl(U;\square_{1};L_{\ast}^{2(j+1-m)}\alpha_{0,j+1}\bigr)
\Longrightarrow
\mathsf{g}_{j}\bigl(U;\square_{1};L^{-1}L_{\ast}^{2(j+1-m)}\alpha_{0,j+1}\bigr),
\end{equation*}
and the gauge $u$ is the same on both sides.

The conditions of the clause are strict upper bounds, by the radius, on the gauge-transformed potential
and on its finite differences. Hence they persist when the radius is increased. They also persist when
they are read only on the bonds of the subcube $\square\subset\square_{1}$. Under
\eqref{4.9:eq-stage-nesting-1}, therefore, the right-hand side implies
$\mathsf{g}_j(U;\square;L_{\ast}^{2(j-m)}\alpha_{0,j})$, which is the clause of $\square$.

It remains to deduce \eqref{4.9:eq-stage-nesting-1} from \eqref{4.9:eq-stage-nesting-2}. We divide
\eqref{4.9:eq-stage-nesting-1} by $L_{\ast}^{2(j-m)}>0$ and multiply it by $L$. This shows that
\eqref{4.9:eq-stage-nesting-1} is equivalent to
$L_{\ast}^{2}\alpha_{0,j+1}\le L\,\alpha_{0,j}$. By the ratio-of-radii bound and by
\eqref{4.9:eq-stage-nesting-2},
\begin{equation*}
L_{\ast}^{2}\alpha_{0,j+1}
\le\sqrt{2}\,(1+\beta_0)\,L_{\ast}^{2}\,\alpha_{0,j}
\le L\,\alpha_{0,j}.
\end{equation*}
\end{proof}

\begin{lemma}[One step reproduces the thin-cube gradient bound]\label{4.9:lem-thin-cube-gradient}
Fix, in this order: $d$ and $N$; the blocking factor $L$, subject to the floors on $L$ in the
first line of \eqref{4.9:eq-thin-cube-gradient-2} below and to the floor on $L$ under which
the one-power hereditary class is defined; the constants $R_1$, $M_1$, $B_1$,
$c_1$ of \cite[Theorems~2, 4 and Proposition~6]{B-CMP99a} and $a_1$, $B_3$ of
\cite[Theorem~1]{B-CMP102b},
which depend on $d$ and $L$ only; the constants $C_{\mathrm{gl}}$, $r_{\mathrm{gl}}$ of the
gluing lemma (Lemma~\ref{4.4:lem-gluing-boxes}) and the group constant $\kappa_G$, which
depend on $d$ and $G$ only. Let $B^{\circ}$ denote the absolute constant of
\cite[(2.38)]{B-CMP119} as Ba{\l}aban chooses it, and let $B_{12}^{\Sigma^1}$ denote the
lower bound on $B_{12}$ imposed with the one-power hereditary class
$U^{\Sigma^1}_j(X;\alpha_0,\alpha_1)$. Put
\begin{equation}\label{4.9:eq-thin-cube-gradient-1}
\begin{aligned}
B&:=\max\bigl\{B^{\circ},\ 2B_3,\ 14dL^5B_1,\ \pi L^{m_0}(d-1)\kappa_G\bigr\},\\
B_{12}&:=\max\bigl\{4C_{\mathrm{gl}}B,\ 2B_3,\ 14dL^2B_1,\ B_{12}^{\Sigma^1}\bigr\},
\end{aligned}
\end{equation}
where $m_0\in\{0,1\}$ is the lowest index of the nest; the last entry of the second maximum
is not used in what follows. Put $S:=c_{12}LM$. The floors of the statement are
\begin{equation}\label{4.9:eq-thin-cube-gradient-2}
\begin{gathered}
c_{12}L\ge 2,\qquad L\ge 2(1+\beta),\qquad L\ge 2\sqrt{2}\,(1+\beta)(1+\beta_0),\\
(\mathrm{f}3)\ \ M\ge 2LR_1M_1,\qquad M\ge 80;
\end{gathered}
\end{equation}
the first line consists of floors on $L$, imposed when $L$ is fixed, the second of floors on
$M$. Then choose $M$ subject to the second line of \eqref{4.9:eq-thin-cube-gradient-2}, and
then the volume, the amplitudes and $\gamma$, subject to the one-sided conditions below; here
$\varepsilon_0$ is the radius of the space $\mathfrak{U}_k(\varepsilon_0)$ of regular real
configurations to which the real backgrounds belong, and $a$ is the radius at which the cube
clause is read:
\begin{equation}\label{4.9:eq-thin-cube-gradient-a}
\begin{aligned}
&(\varphi\text{-gl}) && B_{12}\ge 4C_{\mathrm{gl}}B,\\
&(\varphi\text{-seed}) && B_{12}\ge 2B_3\ \text{(trivial nests)},
\qquad B_{12}\ge 14dL^2B_1\ \text{(general nests)},\\
&(\varphi\text{-B}) && B\ge 2B_3,\quad B\ge 14dL^2B_1,\quad B\ge 14dL^5B_1,
\quad B\ge \pi L^{m_0}(d-1)\kappa_G,\\
&(\mathrm{f}3)^{\sharp} && M\ge 7LR_1M_1,\\
&(\mathrm{C\text{-}}1) && 2\varepsilon_0\le a_1,\qquad 2S\varepsilon_0\le R_1M_1a_1,\\
&(\mathrm{C\text{-}}2) && 14dL^2Sa\le c_1\quad(\text{resp. } 28dL^5Ma\le c_1),\\
&(\mathrm{gl}) && 6S^2Ba\le r_{\mathrm{gl}},
\end{aligned}
\end{equation}
the floor $(\mathrm{T}_{\mathrm{low}})^{S}$ of \eqref{4.9:eq-real-backgrounds-a} on the volume,
read in the units of every level at which a cube clause is read, that is, the torus side is at
least $3S$ in those units,
and the conditions $(\mathrm{adm})^{\flat}$, $(\mathrm{N\text{-}par})$ of
\eqref{4.9:eq-large-field-cubes-a} on the nests. In $(\varphi\text{-B})$ the floor
$14dL^2B_1$ serves $(\ell 2)^{\flat}$ on regions made by the operation $\mathbf{T}$ inside
$\Lambda_j$; the floor $14dL^5B_1$ serves the class cubes without margin and the thin shells.
Assume:
\begin{enumerate}
\item[(a)] for trivial nests, the bounds \cite[Theorem~1, (8)--(9)]{B-CMP102b}; for general
nests, those of \cite[Proposition~6]{B-CMP99a};
\item[(b)] \cite[Proposition~2]{B-CMP98}, in the form in which it is quoted in
\cite{B-CMP109};
\item[(c)] the gluing lemma on boxes (Lemma~\ref{4.4:lem-gluing-boxes}) and its form on
cubes (Corollary~\ref{4.4:cor-gluing-cubes});
\item[(d)] the factorisations $\mathbf{U}=U'U$, into a complex factor and a $G$-valued
factor, of the representatives exhibited at the loci of \eqref{4.9:eq-gauge-transport-a},
and, on the side of the large-field operation $\mathbf{R}$, the factorisation of the
representative of the source of $\mathbf{R}$ in direct form, that is, with its $G$-valued
factor $\bar U$ equal to the source's unitary factor;
\item[(e)] the nest of small-field regions $\Lambda_j$ satisfies the margin conditions of
\cite[\S2]{B-CMP119}, with the margins $R_j$ of \cite[(2.5)]{B-CMP119};
\item[(f)] for the real backgrounds at the stages of the large-field operation, the constants
of the regularity bounds $\mathrm{Reg}_j[R;a]$ at scale $j$ in
\eqref{4.9:eq-collar-regularity-a}, layer by layer, are bounded by the radii.
\end{enumerate}
Then the three clauses of \eqref{4.9:eq-reproduction-meaning-a} hold:
\begin{enumerate}
\item[(i)] the transport clause holds at every locus of \eqref{4.9:eq-gauge-transport-a};
\item[(ii)] the clause on initial pairs holds;
\item[(iii)] the clause on the constant holds: $B$ and $B_{12}$ depend on $d$, $L$, $G$ only
and are chosen before $M$; no floor imposed in this section (on $L$, $M$, $B$, $B_{12}$ or
the volume) depends on $k$, $j$, the level difference $i-i'$ of a change of level, the layer
index $\bar n$, $K$, the nest, $X$, the coupling or the ratio $\alpha_1/\alpha_0$; the
ceilings bound the amplitudes,
hence $\gamma$, from above, are chosen after $(L,M,B,B_{12})$, and hold uniformly in $k$
and $j$.
\end{enumerate}
That is, the step from level $k$ to level $k+1$ reproduces the gauge condition
\cite[(1.12)]{B-CMP109} in the widened reading
$(\mathrm{L}\text{-}12)=(\ell 1)$ of \eqref{4.9:eq-cube-gauge-condition-b} on
$U^{\Sigma^1}$, its gradient member included, with the single constant $B_{12}$. The same
holds for $(\ell 2)^{\flat}$. For $(\ell 2)$ it holds at the loci $\langle 1\rangle$--$\langle 7\rangle$,
with all their primed and lettered sub-items, and $\langle 10\rangle$, $\langle 11\rangle$
of \eqref{4.9:eq-gauge-transport-a} and at the initial
pairs, there under $(\mathrm{R\text{-}}c_{12})$ of \eqref{4.9:eq-real-backgrounds-c}, and it
fails at $\langle 9\rangle$. The conditions $(c_{12}L-2)M\ge 1$ of \eqref{4.9:eq-flat-ring-a}
and $(\mathrm{c\text{-}W})$ of
\eqref{4.9:eq-ambient-condition-a} enter only the example of this failure and are not
hypotheses of the lemma.
\end{lemma}

\begin{proof}
By Remark~\ref{4.9:rem-reproduction-meaning}, reproduction of the gauge condition by one
step means exactly the three clauses of \eqref{4.9:eq-reproduction-meaning-a}; we verify
them in turn.

\emph{Clause (i).} Let a locus of \eqref{4.9:eq-gauge-transport-a} be given at which a
membership in $U^{\Sigma^1}_j(X;\alpha_0,\alpha_1)$ is derived from another one, and assume
that the representative of the source satisfies $(\mathrm{L}\text{-}12)$ at the source's level,
domain and radius. At every locus of \eqref{4.9:eq-gauge-transport-a} other than
$\langle 8\rangle$, $\langle 9\rangle$, $\langle 13\rangle$ we apply the transport proposition
(Proposition~\ref{4.9:prop-gauge-transport}) with
$\star=(\ell 1)=(\mathrm{L}\text{-}12)$, read with its gradient member and with the one pair
$(B,B_{12})$ of \eqref{4.9:eq-thin-cube-gradient-1}. Its inputs are the factorisation of the
exhibited representative, which is hypothesis (d), and the floors listed in that proposition,
which are among \eqref{4.9:eq-thin-cube-gradient-2} and \eqref{4.9:eq-thin-cube-gradient-a};
no other condition enters. It gives $(\mathrm{L}\text{-}12)$ for the exhibited representative
at the product's level, domain and radius. At $\langle 8\rangle$ and $\langle 9\rangle$ the source
clause is the clause \eqref{4.9:eq-cube-gauge-condition-c} of \cite[(2.38)]{B-CMP119} in the
source's nest; it is carried to the product under the conditions
\eqref{4.9:eq-large-field-cubes-a} on the nest, and the gradient bound of
$(\mathrm{L}\text{-}12)$ is reassembled from its gradient member by the passage to smaller
cubes by gluing (Corollary~\ref{4.9:cor-gluing-smaller-cubes}). The hypotheses of that
corollary are $M\ge 80$ and $c_{12}L\ge 1$, both implied by
\eqref{4.9:eq-thin-cube-gradient-2}, the ceiling $(\mathrm{gl})$ and the floor
$(\varphi\text{-gl})$ of \eqref{4.9:eq-thin-cube-gradient-a}; its input is hypothesis (c).
At $\langle 13\rangle$, on the side of the large-field operation, clause (i) is the proposition
on cube conditions inside the large-field spaces
(Proposition~\ref{4.9:prop-large-field-cubes}): its hypotheses are the conditions
\eqref{4.9:eq-large-field-cubes-a}, the radii condition $(\mathrm{rat})$ among them following
from the ratio of radii along the coupling flow (Lemma~\ref{4.4:lem-radii-ratio}) under the
floor $L\ge 2\sqrt{2}\,(1+\beta)(1+\beta_0)$ of \eqref{4.9:eq-thin-cube-gradient-2}, inside
which the floor $L\ge\sqrt{2}\,(1+\beta_0)$ used there is read; the floor
$(\varphi\text{-gl})$ used in its fourth clause;
the factorisation on the side of $\mathbf{R}$ of hypothesis (d); the margins of
hypothesis (e); and, for
its fifth clause, the layerwise bounds of hypothesis (f); the real stage backgrounds are
treated there through the bounds of \cite[Proposition~6]{B-CMP99a} in hypothesis (a), thin
shells included under $(\mathrm{f}3)^{\sharp}$.

\emph{Clause (ii).} The pairs whose membership is not derived from another membership are
the unit pair, the small pairs and the real backgrounds. For them the clause is
Proposition~\ref{4.9:prop-real-backgrounds} together with
Corollary~\ref{4.9:cor-printed-pairs}, at the constant $\alpha_0'$ of the printed membership.
The inputs of that proposition are the two parts of hypothesis (a):
\cite[Theorem~1, (8)--(9)]{B-CMP102b} in its first clause (trivial nests) and
\cite[Proposition~6]{B-CMP99a} in its second and third clauses. Its conditions are met as
follows: the
ceilings $(\mathrm{C\text{-}}1)$, used in its first clause (trivial nests), and
$(\mathrm{C\text{-}}2)$, used in its second and third clauses, are in
\eqref{4.9:eq-thin-cube-gradient-a}; the floors $S\ge 4R_1M_1$ and $M\ge 4R_1M_1$ follow from
$c_{12}L\ge 2$, $(\mathrm{f}3)$ and $L\ge 2(1+\beta)$ in \eqref{4.9:eq-thin-cube-gradient-2};
$(\mathrm{f}3)^{\sharp}$ is used in its third clause at shells of one cube of side $M$ only;
$(\mathrm{T}_{\mathrm{low}})^{S}$ guarantees that the ambient cubes are embedded in the torus;
and the floor $(\varphi\text{-seed})$ on $B_{12}$ converts the radius produced there into the
radius $SB_{12}\alpha_0'$ of $(\mathrm{L}\text{-}12)$, for trivial and for general nests.

\emph{Clause (iii).} By \eqref{4.9:eq-thin-cube-gradient-1}, $B$ is a maximum of $B^{\circ}$,
which is absolute, and of expressions in $d$, $L$, $B_1$, $B_3$, $\kappa_G$ and $m_0$;
since $m_0\in\{0,1\}$, the entry $\pi L^{m_0}(d-1)\kappa_G$ is one of the two numbers
$\pi(d-1)\kappa_G$, $\pi L(d-1)\kappa_G$, each of which depends on $d$, $L$, $G$ only;
and $B_{12}$ is a maximum of expressions in $B$, $C_{\mathrm{gl}}$, $d$, $L$, $B_1$, $B_3$
and of $B_{12}^{\Sigma^1}$, the floor imposed with the one-power hereditary class, which
depends on $d$, $L$, $G$ only and is fixed before $M$.
Hence both depend on $d$, $L$, $G$ only and are fixed before $M$. The floors
\eqref{4.9:eq-thin-cube-gradient-2}, the floors $(\varphi\text{-gl})$, $(\varphi\text{-seed})$,
$(\varphi\text{-B})$, $(\mathrm{f}3)^{\sharp}$ of \eqref{4.9:eq-thin-cube-gradient-a} and the
floor $(\mathrm{T}_{\mathrm{low}})^{S}$ on the volume are written in terms of $d$, $L$, $M$,
$R_1$, $M_1$, $\beta$, $\beta_0$, $c_{12}$ and the constants of
\eqref{4.9:eq-thin-cube-gradient-1} only;
none of them contains $k$, $j$, $i-i'$, $\bar n$, $K$, the nest, $X$, the coupling or
$\alpha_1/\alpha_0$, and neither does the floor on $L$ under which the one-power hereditary
class is defined. The
ceilings $(\mathrm{C\text{-}}1)$, $(\mathrm{C\text{-}}2)$ and $(\mathrm{gl})$ bound $\varepsilon_0$
and $a$ from above by quantities in $d$, $L$, $M$, $B$ and the constants of
\cite{B-CMP99a}, \cite{B-CMP102b} and of the gluing lemma, and the condition
$(\mathrm{R\text{-}}c_{12})$ on the margins is met by the choice of $\gamma$ after $L$; all of
them are imposed after $(L,M,B,B_{12})$, and none contains $k$ or $j$, so they hold uniformly
in $k$ and $j$. Finally, $B$ and $B_{12}$ grow with $L$ through $B_1$, $B_3$ and the displayed
powers of $L$; the only conditions that read them are ceilings on amplitudes chosen after
them, among them $(\mathrm{gl})$ and the ceiling $c_{12}LMB_{12}\alpha_0\le a_0$ on the radius
of the one-power hereditary class. Hence no condition on $L$ reads $B$ or $B_{12}$, and the
existential constant $L_0$ is not affected.

\emph{The concluding assertions.} Clauses (i)--(iii) are, by
Remark~\ref{4.9:rem-reproduction-meaning}, the statement that the step from level $k$ to
level $k+1$ reproduces \cite[(1.12)]{B-CMP109} in the reading $(\mathrm{L}\text{-}12)$ on
$U^{\Sigma^1}$, gradient member included, with the single constant $B_{12}$. For
$(\ell 2)^{\flat}$ the same three steps apply: at the loci other than $\langle 8\rangle$,
$\langle 9\rangle$, $\langle 13\rangle$ by
Proposition~\ref{4.9:prop-gauge-transport} with $\star=(\ell 2)^{\flat}$; at $\langle 8\rangle$ and
$\langle 9\rangle$ because $(\ell 2)^{\flat}$ is the source clause
\eqref{4.9:eq-cube-gauge-condition-c} of \cite[(2.38)]{B-CMP119} itself, carried to the
product as in the proof of clause (i), so that no reassembly by gluing is needed; at
$\langle 13\rangle$ by Proposition~\ref{4.9:prop-large-field-cubes};
and at the initial pairs by Proposition~\ref{4.9:prop-real-backgrounds} read for
$(\ell 2)^{\flat}$, under the floor $(\varphi\text{-B})$ on $B$. For $(\ell 2)$,
Proposition~\ref{4.9:prop-gauge-transport} with $\star=(\ell 2)$ gives the clause at
$\langle 1\rangle$--$\langle 7\rangle$, with all their primed and lettered sub-items, and at
$\langle 10\rangle$, $\langle 11\rangle$, and
Proposition~\ref{4.9:prop-real-backgrounds} gives it at the initial pairs under
$(\mathrm{R\text{-}}c_{12})$. At $\langle 9\rangle$ it fails: by clause (i) of the proposition on
ambient versus inner cube gauge conditions (Proposition~\ref{4.9:prop-ambient-condition}),
$(\ell 2)^{\flat}$ does not imply $(\ell 2)$, the flat field winding around a hole
\eqref{4.9:eq-flat-ring-a} satisfying the source clause of $\langle 9\rangle$ while its whole
$G^{c}$-orbit misses the class read by $(\ell 2)$, under $(c_{12}L-2)M\ge 1$ of
\eqref{4.9:eq-flat-ring-a} and
$(\mathrm{c\text{-}W})$ of \eqref{4.9:eq-ambient-condition-a}.
\end{proof}

\begin{remark}[Statements not asserted here]\label{4.9:rem-not-asserted}
The results of this section concern one thing only: the gauge condition in clause~(i) of
the one-power hereditary class $U^{\Sigma^1}_j(X;\alpha_0,\alpha_1)$. No other clause of that
class is touched. In particular, the following are not asserted here.
\begin{itemize}
\item Any clause of the one-power hereditary class other than the gauge condition in
clause~(i). This includes its remaining clauses and the reading of \cite[(1.17)]{B-CMP109}.
It also includes the fourth of the hereditary conditions attached to the post-step space
$\mathfrak{U}^{\mathrm{post},U}_{k+1}(Y)$ of Definition~\ref{4.4:def-post-step-space}, in
its primed form.
\item The layerwise constants in the regularity conditions
\eqref{4.9:eq-collar-regularity-a} satisfied by the true background configurations in the
nests of small-field regions of \cite{B-CMP119}, \cite{B-CMP122a} and \cite{B-CMP122b}.
These constants are quoted from the cited papers and are not re-derived here.
\item The representative that Ba{\l}aban's construction obtains by minimising the
configuration once more.
\item Any use of configurations lying in the band $\mathfrak{K}_c$ around the real background
that accompanies the composite configurations of this chapter.
\item The margins at those indices $j$ produced by the large-field operation $\mathbf{R}$
for which $\Lambda_j=\Omega_j$ \cite{B-CMP119}. At such indices, clause (c)(i) of
Proposition~\ref{4.9:prop-real-backgrounds} is replaced by clause (c)(ii) of that
proposition, combined with the passage to smaller cubes by gluing
(Corollary~\ref{4.9:cor-gluing-smaller-cubes}). This combination is applied on cubes
contained in a closed set $X\subset\Omega_j$. It yields the reading $(\ell 1)_j(X;a)$, that
is, the widened reading $(\mathrm{L}\text{-}12)_j(X;a)$ of
\eqref{4.9:eq-cube-gauge-condition-b}. It does so under the constant condition
$B_{12}\ge 4C_{\mathrm{gl}}B$ of \eqref{4.9:eq-gluing-smaller-cubes-a} and under condition
$(\varphi\text{-}\mathrm{B})$ among \eqref{4.9:eq-thin-cube-gradient-a}. The ambient gauge
condition of Proposition~\ref{4.9:prop-real-backgrounds}, with its constants, is asserted only
for cubes contained in $\Omega_j$.
\item Small tori, namely those whose side, in level-$k$ units, does not satisfy the lower
bound required in the hypotheses of Proposition~\ref{4.9:prop-real-backgrounds}.
\end{itemize}
Finally, two proofs are not reproduced in this section: that of the gluing lemma
(Lemma~\ref{4.4:lem-gluing-boxes}), and that of the statement which establishes the widened
reading $(\mathrm{L}\text{-}12)_j(X;a)$ of \eqref{4.9:eq-cube-gauge-condition-b} as clause~(i)
of the one-power hereditary class $U^{\Sigma^1}_j(X;\alpha_0,\alpha_1)$. Both statements are
used here as stated.
\end{remark}

\begin{definition}[Lattices, units and regions]\label{4.10:not-lattices-regions}
Throughout, $L$ is a fixed odd integer, $m'\ge 1$ is an integer, and
\begin{equation}\label{4.10:eq-lattices-regions-1}
M:=L^{m'}.
\end{equation}
No condition on $L$ other than oddness is used in what follows.

\emph{The fine lattice.} All objects live on one fine periodic lattice $\mathbb{T}$, a
four-dimensional discrete torus, whose spacing is one \emph{fine unit}. The \emph{sites} are
the points of $\mathbb{T}$, and $e_1,\dots,e_4$ denote the unit vectors. A
\emph{positively oriented fine bond}, or simply a \emph{bond}, is an ordered pair
$b=\langle x,x+e_\mu\rangle$ of sites, with $\mu\in\{1,\dots,4\}$; we write
\begin{equation*}
b_-:=x,\qquad b_+:=x+e_\mu,\qquad b^{-1}:=\langle x+e_\mu,x\rangle,
\end{equation*}
and call $b^{-1}$ the \emph{reversed bond}. A \emph{plaquette} $P$ is an elementary unit
square of $\mathbb{T}$ together with a base corner and an orientation. Its boundary is the
closed path of four bonds $b_1,b_2,b_3,b_4$, of which the last two are traversed against
their positive orientation.

\emph{Units.} Lengths in \emph{level-$i$ units} are counted in steps of $L^i$ fine units.
Accordingly the fine spacing, measured in level-$i$ units, is
\begin{equation}\label{4.10:eq-lattices-regions-2}
\xi_i:=L^{-i}.
\end{equation}

\emph{Coarse sites and bonds.} For $p\ge 0$ the \emph{level-$p$ sites} are the centres of
the blocks of side $L^p$ fine units of the block partition used by Ba{\l}aban. For $p=0$
these are all the sites of $\mathbb{T}$; for every $p$ they are sites of $\mathbb{T}$,
because $L$ is odd, so that a block of side $L^p$ fine units has a lattice point at its
centre. A \emph{level-$p$ bond} is a pair $c=\langle y,y+L^p e_\mu\rangle$ of level-$p$
sites, and we write
\begin{equation*}
c_-:=y,\qquad c_+:=y+L^p e_\mu .
\end{equation*}
For $p=0$ the level-$p$ bonds are the fine bonds.

\emph{Cube partitions and domains.} For each level $i$, $\pi_i$ denotes the partition of
$\mathbb{T}$ into cubes of side $M$ level-$i$ units, that is, of side $L^{i+m'}$ fine units
by \eqref{4.10:eq-lattices-regions-1}, centred at the level-$(i+m')$ sites. We write
$\mathcal{D}_i$ for the set of connected finite unions of $\pi_i$-cubes. These classes
satisfy
\begin{equation*}
\mathcal{D}_i\subset\mathcal{D}_{i-1}.
\end{equation*}

\emph{Thickenings and two-layer interiors.} Let $p\ge 1$ and let $\square$ be a
$\pi_{p-1}$-cube. For an integer $r\ge 0$, $\widetilde{\square}^{\,r}$ denotes the union of
the $\pi_{p-1}$-cubes lying within $r$ layers of $\square$. For $Z\in\mathcal{D}_{p-1}$,
the \emph{two-layer interior} $Z^{\wr}$ of $Z$ is the union of those $\pi_{p-1}$-cubes
$\square\subset Z$ for which $\widetilde{\square}^{\,2}\subset Z$.

\emph{Finiteness.} All sets of sites, bonds and plaquettes occurring in what follows are
finite.
\end{definition}

\begin{definition}[Complex configurations and holomorphic maps]\label{4.10:def-holomorphic-maps}
\emph{Matrices and groups.} $M_N(\mathbb{C})$ is the algebra of complex $N\times N$ matrices,
identified with $\mathbb{C}^{N^2}$ through the entries, and $|\cdot|$ is the operator norm on
$M_N(\mathbb{C})$ with respect to the Hermitian norm on $\mathbb{C}^N$. Here and in what follows
$\operatorname{Tr}$ denotes the ordinary matrix trace, $\operatorname{Tr}\mathbf{1}=N$, where
$\mathbf{1}$ is the identity matrix. For $Y\in M_N(\mathbb{C})$, $\mathrm{spec}(Y)$ is the set of
eigenvalues of $Y$, and every $\lambda\in\mathrm{spec}(Y)$ satisfies $|\lambda|\le|Y|$. We write
\begin{equation*}
G:=SU(N),\qquad G^c:=SL(N,\mathbb{C}),\qquad \mathfrak{g}:=\mathfrak{su}(N),\qquad
\mathfrak{g}^c:=\mathfrak{sl}(N,\mathbb{C}),
\end{equation*}
so that $\mathfrak{g}$ consists of the traceless anti-Hermitian matrices and $\mathfrak{g}^c$ of the
traceless matrices. As real vector spaces $\mathfrak{g}^c=\mathfrak{g}\oplus i\mathfrak{g}$, and
$\mathfrak{g}^c$ is a complex vector space of dimension $N^2-1$. The exponential map
$\exp:M_N(\mathbb{C})\to M_N(\mathbb{C})$,
\begin{equation*}
\exp Y=e^{Y}=\sum_{k\ge0}\frac{Y^{k}}{k!},
\end{equation*}
is entire (given by an everywhere convergent power series) and satisfies
$\det e^{Y}=e^{\operatorname{Tr}Y}$; hence $e^{Y}\in G^c$ for $Y\in\mathfrak{g}^c$.

\emph{Configurations.} Let $S$ be a finite set of bonds, fine or coarse, in the conventions of
Definition~\ref{4.10:not-lattices-regions}. A \emph{configuration} on $S$ is an element
$W\in\mathrm{Conf}(S):=(G^c)^{S}$, $b\mapsto W(b)$; it is extended to reversed bonds by
$W(b^{-1}):=W(b)^{-1}$. A configuration is \emph{$G$-valued} if $W\in G^{S}$. A \emph{potential}
on $S$ is an element of the complex vector space $(\mathfrak{g}^c)^{S}$. For a potential $B$ and a
configuration $V_0$ on $S$, $e^{iB}V_0$ denotes the configuration $c\mapsto e^{iB(c)}V_0(c)$. For a
configuration $W$ on $S$ and a plaquette $P$ whose four bonds lie in $S$ up to orientation, the
\emph{holonomy} of $W$ around $P$ is
\begin{equation*}
\partial W(P):=W(b_1)W(b_2)W(b_3)W(b_4),
\end{equation*}
the ordered product around $P$, reversed bonds entering through inverses. The set
$\mathrm{Conf}(S)$ is regarded as a subset of the complex vector space $M_N(\mathbb{C})^{S}$ and
carries the induced topology.

\emph{Holomorphy.} Let $\mathcal{O}$ be an open subset of $\mathrm{Conf}(S)$, or of a complex vector
space, or of a product of such; its ambient complex vector space is $M_N(\mathbb{C})^{S}$, the
vector space itself, or the product of the ambient spaces, respectively. A map $F$ from
$\mathcal{O}$ to a complex vector space $E$ is \emph{holomorphic} if every point $Q$ of
$\mathcal{O}$ has an open neighbourhood $\mathcal{N}$ in the ambient complex vector space and a
holomorphic map $\tilde F:\mathcal{N}\to E$ (given locally by convergent power series, equivalently
complex differentiable) with $\tilde F=F$ on $\mathcal{N}\cap\mathcal{O}$. A map with values in
$\mathrm{Conf}(S')$ is holomorphic if it is holomorphic as a map into $M_N(\mathbb{C})^{S'}$. For
open subsets of complex submanifolds this agrees with holomorphy in local holomorphic coordinates;
that equivalence is not used.

\emph{Facts.} The following five facts are used:
\begin{equation}\label{4.10:eq-holomorphic-maps-a}
\begin{aligned}
&\mathrm{(Hol1)}\quad G^c=\{\det=1\}\ \text{is a closed complex submanifold};\\
&\mathrm{(Hol2)}\quad \text{polynomial maps are holomorphic};\\
&\mathrm{(Hol3)}\quad B\mapsto e^{i\xi B}V_0\ \text{holomorphic};\
 (U,A')\mapsto e^{i\xi A'}U\ \text{continuous};\\
&\mathrm{(Hol4)}\quad F_1\circ F\ \text{holomorphic};\ F\ \text{continuous};\\
&\mathrm{(Hol5)}\quad \text{holomorphy is local}.
\end{aligned}
\end{equation}
The precise statements, to which \eqref{4.10:eq-holomorphic-maps-a} refers, are the following.
\begin{itemize}
\item[] \textup{(Hol1)}\ $G^c$ is a closed complex submanifold of $M_N(\mathbb{C})$ of complex
dimension $N^2-1$, namely the level set $\{\det=1\}$. Hence $\mathrm{Conf}(S)$ is a closed complex
submanifold of $M_N(\mathbb{C})^{S}$, carrying the induced topology, and $G^{S}$ is a compact
subset of it.
\item[] \textup{(Hol2)}\ The restriction to $\mathcal{O}$ of a map that is polynomial in the matrix
entries is holomorphic. In particular $W\mapsto W(b)W(b')$, $W\mapsto W(b)^{-1}$,
$W\mapsto\partial W(P)$ and the restriction maps $W\mapsto W|_{S'}$, $S'\subset S$, are
holomorphic on $\mathrm{Conf}(S)$.
\item[] \textup{(Hol3)}\ For fixed $V_0\in\mathrm{Conf}(S)$ and a fixed real number $\xi>0$, the map
$B\mapsto e^{i\xi B}V_0$ from $(\mathfrak{g}^c)^{S}$ to $\mathrm{Conf}(S)$ is holomorphic, and the
map $(U,A')\mapsto e^{i\xi A'}U$ from $G^{S}\times(\mathfrak{g}^c)^{S}$ to $\mathrm{Conf}(S)$ is
continuous.
\item[] \textup{(Hol4)}\ If $F:\mathcal{O}\to\mathcal{O}'$ and $F_1:\mathcal{O}'\to E$ are
holomorphic, then $F_1\circ F$ is holomorphic. Holomorphic maps are continuous, and preimages of
open sets under them are open.
\item[] \textup{(Hol5)}\ Holomorphy is a local property: a map $F$ on $\mathcal{O}$ is holomorphic
if and only if every point $Q$ of $\mathcal{O}$ has an open neighbourhood
$\mathcal{O}_Q\subset\mathcal{O}$ in $\mathcal{O}$ such that $F|_{\mathcal{O}_Q}$ is holomorphic.
\end{itemize}
\end{definition}

\begin{proof}
\emph{Eigenvalues.} If $Yv=\lambda v$ with $v\neq0$, then $|\lambda|\,|v|=|Yv|\le|Y|\,|v|$, and
$|\lambda|\le|Y|$ follows on dividing by $|v|$.

\emph{The decomposition of $\mathfrak{g}^c$.} Let $Y\in\mathfrak{g}^c$ and put
$Y_1=(Y-Y^*)/2$, $Y_2=(Y+Y^*)/(2i)$, so that $Y=Y_1+iY_2$. Then $Y_1^*=(Y^*-Y)/2=-Y_1$ and
$Y_2^*=(Y^*+Y)/(-2i)=-Y_2$, and both are traceless because
$\operatorname{Tr}Y^*=\overline{\operatorname{Tr}Y}=0$; thus $Y_1,Y_2\in\mathfrak{g}$. Conversely
$\mathfrak{g}+i\mathfrak{g}$ consists of traceless matrices. The sum is direct: if
$Y_0\in\mathfrak{g}$ and $Y_0=iY'$ with $Y'\in\mathfrak{g}$, then $Y_0^*=-iY'^*=iY'=Y_0$, so $Y_0$
is both Hermitian and anti-Hermitian, hence $Y_0=0$. Finally $\mathfrak{g}^c$ is the kernel of the
complex-linear form $\operatorname{Tr}$ on $M_N(\mathbb{C})$, which is non-zero since
$\operatorname{Tr}\mathbf{1}=N$; so $\mathfrak{g}^c$ is a complex subspace of dimension $N^2-1$.

\emph{The exponential.} Since $|Y^k/k!|\le|Y|^k/k!$, the series converges absolutely for every $Y$,
and each entry of $e^{Y}$ is an everywhere convergent power series in the entries of $Y$; thus
$\exp$ is entire. For $Y\in\mathfrak{g}^c$ we have $\operatorname{Tr}Y=0$, so
$\det e^{Y}=e^{0}=1$ and $e^{Y}\in G^c$. In particular, for a potential $B$ and a configuration
$V_0$, each $iB(c)$ lies in the complex space $\mathfrak{g}^c$, so $e^{iB(c)}\in G^c$, and
$e^{iB(c)}V_0(c)$ has determinant $1$; thus $e^{iB}V_0\in\mathrm{Conf}(S)$.

\emph{(Hol1).} The map $\det$ is a polynomial in the matrix entries, hence holomorphic on
$M_N(\mathbb{C})$, and at an invertible $Y$ its differential is
$d\det_Y(Y')=\det(Y)\operatorname{Tr}(Y^{-1}Y')$. For $Y\in G^c$ and $Y'=Y$ this equals
$\det(Y)\operatorname{Tr}\mathbf{1}=N\neq0$, so $d\det_Y$ is a non-zero complex-linear form at
every point of $\{\det=1\}$. By the holomorphic implicit function theorem, $\{\det=1\}=G^c$ is a
complex submanifold of $M_N(\mathbb{C})$ of complex codimension one, that is, of dimension
$N^2-1$; it is closed as the preimage of $\{1\}$ under the continuous map $\det$. Consequently
$\mathrm{Conf}(S)=(G^c)^S$, a finite product of closed complex submanifolds, is a closed complex
submanifold of $M_N(\mathbb{C})^{S}$. For $Y\in G$ one has $Y^*Y=\mathbf{1}$, so $|Y|=1$; and $G$
is the zero set of the continuous maps $Y\mapsto Y^*Y-\mathbf{1}$ and $Y\mapsto\det Y-1$. Hence $G$
is closed and bounded in $M_N(\mathbb{C})\cong\mathbb{C}^{N^2}$, thus compact, and the finite
product $G^{S}$ is a compact subset of $\mathrm{Conf}(S)$.

\emph{(Hol2).} A map $F_0$ from the ambient space to $E$ whose components are polynomials in the
matrix entries is a finite power series, hence holomorphic on the whole ambient space; with
$\mathcal{N}$ the ambient space and $\tilde F=F_0$ the definition is satisfied at every point. The
entries of $W(b)W(b')$ are sums of products of entries of $W(b)$ and $W(b')$. For invertible $Y$,
$Y^{-1}=\operatorname{adj}(Y)/\det Y$, where the entries of the adjugate $\operatorname{adj}(Y)$
are signed minors of order $N-1$, hence polynomials; on $G^c$ this gives
$Y^{-1}=\operatorname{adj}(Y)$, so $W\mapsto W(b)^{-1}$ is the restriction to $\mathrm{Conf}(S)$ of
the polynomial map $W\mapsto\operatorname{adj}(W(b))$. Each factor of $\partial W(P)$ is $W(b_i)$ or
$\operatorname{adj}(W(b_i))$, so $\partial W(P)$ is the restriction of a product of four polynomial
maps, again polynomial. The restriction $W\mapsto W|_{S'}$ is the restriction of the linear
coordinate projection $M_N(\mathbb{C})^{S}\to M_N(\mathbb{C})^{S'}$, and it maps
$\mathrm{Conf}(S)$ into $\mathrm{Conf}(S')$.

\emph{(Hol3).} The domain $(\mathfrak{g}^c)^S$ is a complex vector space. For each $c\in S$, the
component $B\mapsto e^{i\xi B(c)}V_0(c)$ is the composition of the complex-linear map
$B\mapsto i\xi B(c)$ into $M_N(\mathbb{C})$, the entire map $\exp$, and the map $Y\mapsto YV_0(c)$,
which is polynomial in the entries and hence holomorphic by (Hol2). A composition of holomorphic
maps between complex vector spaces is holomorphic, and a map into $M_N(\mathbb{C})^{S}$ is
holomorphic when each of its components is.
The values lie in $\mathrm{Conf}(S)$ by the paragraph on the exponential. Hence
$B\mapsto e^{i\xi B}V_0$ is holomorphic into $\mathrm{Conf}(S)$. The map
$(W,A')\mapsto(e^{i\xi A'(c)}W(c))_{c\in S}$ on $M_N(\mathbb{C})^{S}\times(\mathfrak{g}^c)^{S}$ is
continuous, because $\exp$ and matrix multiplication are continuous. Its restriction to
$G^{S}\times(\mathfrak{g}^c)^{S}$, $(U,A')\mapsto e^{i\xi A'}U$, is therefore continuous, and it
takes values in $\mathrm{Conf}(S)$ because $U(c)\in G\subset G^c$ and $e^{i\xi A'(c)}\in G^c$ by the
paragraph on the exponential.

\emph{(Hol4).} Let $Q\in\mathcal{O}$, let $\tilde F$ on $\mathcal{N}\ni Q$ be a local extension of
$F$, and let $\tilde F_1$ on $\mathcal{N}'\ni F(Q)$ be a local extension of $F_1$. Since $\tilde F$
is continuous and $\tilde F(Q)=F(Q)\in\mathcal{N}'$, the set
$\mathcal{N}_0:=\mathcal{N}\cap\tilde F^{-1}(\mathcal{N}')$ is an open neighbourhood of $Q$ with
$\tilde F(\mathcal{N}_0)\subset\mathcal{N}'$. On $\mathcal{N}_0$ the map
$\tilde F_1\circ\tilde F$ is holomorphic, by the chain rule for complex differentiable maps. For
$Q'\in\mathcal{N}_0\cap\mathcal{O}$ we have $\tilde F(Q')=F(Q')\in\mathcal{N}'\cap\mathcal{O}'$, so
$\tilde F_1(\tilde F(Q'))=F_1(F(Q'))$. Thus $\tilde F_1\circ\tilde F$ is a holomorphic extension of
$F_1\circ F$ near $Q$. Next, $F$ coincides on $\mathcal{N}\cap\mathcal{O}$, an open neighbourhood
of $Q$ in $\mathcal{O}$, with the continuous map $\tilde F$; so $F$ is continuous at $Q$. Since $Q$
is arbitrary, $F$ is continuous, and preimages of open sets under $F$ are open in $\mathcal{O}$.

\emph{(Hol5).} If $F$ is holomorphic, then so is its restriction to any open
$\mathcal{O}_Q\subset\mathcal{O}$, with the same local extensions. Conversely, let $Q\in\mathcal{O}$
and let $F|_{\mathcal{O}_Q}$ be holomorphic. Then there are an open $\mathcal{N}\ni Q$ and a
holomorphic $\tilde F$ on $\mathcal{N}$ with $\tilde F=F$ on $\mathcal{N}\cap\mathcal{O}_Q$. As
$\mathcal{O}_Q$ is open in $\mathcal{O}$, and $\mathcal{O}$ carries the induced topology,
$\mathcal{O}_Q=\mathcal{O}\cap\mathcal{N}_1$ for an open set $\mathcal{N}_1$ of the ambient space.
With $\mathcal{N}_0:=\mathcal{N}\cap\mathcal{N}_1$ we have
$\mathcal{N}_0\cap\mathcal{O}=\mathcal{N}\cap\mathcal{O}_Q$, so $\tilde F|_{\mathcal{N}_0}$ is a
local extension of $F$ at $Q$. Hence $F$ is holomorphic.
\end{proof}

\begin{definition}[Truncated towers, constraint bonds and the data space]\label{4.10:not-truncated-towers}
The lattice $\mathbb{T}$, fine units, level-$q$ bonds and the classes $\mathcal{D}_i$ are those of
Definition~\ref{4.10:not-lattices-regions}. Configurations $\mathrm{Conf}(S)=(G^c)^S$ on a finite
bond set $S$ and their complex structure are those of Definition~\ref{4.10:def-holomorphic-maps}.

\emph{Towers.} Let $Z$ be a region of the fine lattice to which the tower rules of this chapter
apply. In the application made later in this section, $Z\in\mathcal{D}_{p-1}$ for some $p\ge 1$.
\begin{itemize}
\item[] \emph{Rule} $\mathrm{I}$ is the maximal rule of \cite{B-CMP99a} and
\cite[pp.~277--278]{B-CMP102b}. Put $\Omega^{\mathrm{I}}_0(Z):=Z$. For $q\ge 1$, let
$\Omega^{\mathrm{I}}_q(Z)$ be the maximal union of blocks of side $M_1L^q$ fine units whose
distance from the complement of $\Omega^{\mathrm{I}}_{q-1}(Z)$ exceeds $R_1M_1L^{q-1}$ fine units.
Here $R_1$ and $M_1$ are Ba{\l}aban's parameters \cite[pp.~277--278]{B-CMP102b}.
\item[] \emph{Rule} $\mathrm{II}$ is the rule \cite[(2.13)]{B-CMP119}, with its distances measured
in fine units. Like rule $\mathrm{I}$, it produces a decreasing sequence of regions
$\Omega^{\mathrm{II}}_q(Z)$, $q\ge 0$, determined by $Z$.
\end{itemize}
For a rule $\varrho\in\{\mathrm{I},\mathrm{II}\}$ and a depth $n\ge 1$, the \emph{truncated tower}
of depth $n$ is
\begin{equation*}
\mathcal{T}=\mathcal{T}_{\varrho}(Z,n):=(\Omega_0,\Omega_1,\dots,\Omega_n),
\qquad \Omega_q:=\Omega^{\varrho}_q(Z).
\end{equation*}
The tuple $\mathcal{T}$ is determined by $(Z,n,\varrho)$ alone. It is a sequence of domains of the
kind considered in \cite[pp.~277--278]{B-CMP102b}. This is the setting in which the quantities
attached to a tower were defined earlier in this chapter, and those quantities are used here without
change.

\emph{Bonds of a region.} $\mathrm{bd}(Z)$ denotes the finite set of positively oriented fine
bonds of $Z$. The convention is the one under which the tower quantities are fields on the bonds
of a domain, which follows \cite[pp.~277--278]{B-CMP102b}. Let $Z\subset X$. The tower
quantities read a configuration on $X$ only through its values on the bonds of $Z$. These bonds
are bonds of $X$, and restriction from $X$ to $Z$ is understood in this sense.

\emph{Constraint bonds.} The \emph{constraint bond sets} of $\mathcal{T}$ are the following.
\begin{itemize}
\item[] $\Lambda_0(\mathcal{T})$ is the set of fine bonds of the skin $\Omega_0\setminus\Omega_1$
at which the tower construction prescribes the field.
\item[] For $1\le q<n$, $\Lambda_q(\mathcal{T})$ is the set of level-$q$ bonds of
$\Omega_q\setminus\Omega_{q+1}$ at which the tower construction prescribes the $q$-fold block
averages.
\item[] $\Lambda_n(\mathcal{T})$ is the set of level-$n$ bonds of $\Omega_n$.
\end{itemize}
The boundary conventions are those of \cite[pp.~277--278]{B-CMP102b}. Of these sets, only two
facts are used below: they are finite, and they are determined by $\mathcal{T}$. If some
$\Lambda_q(\mathcal{T})$ is empty, the corresponding factor in \eqref{4.10:eq-truncated-towers-1}
is a point.

\emph{The data space.} Put
$\mathfrak{B}(\mathcal{T}):=\bigsqcup_{q=0}^{n}\Lambda_q(\mathcal{T})$. The \emph{data space} of
$\mathcal{T}$ is
\begin{equation}\label{4.10:eq-truncated-towers-1}
\mathbb{D}(\mathcal{T}):=\mathrm{Conf}(\mathfrak{B}(\mathcal{T}))
=\prod_{q=0}^{n}\mathrm{Conf}(\Lambda_q(\mathcal{T})).
\end{equation}
The equality holds because $\mathrm{Conf}(S)=(G^c)^S$ and $\mathfrak{B}(\mathcal{T})$ is the
disjoint union of the sets $\Lambda_q(\mathcal{T})$. Each $\Lambda_q(\mathcal{T})$ is finite, so
$\mathfrak{B}(\mathcal{T})$ is a finite bond set. Hence $\mathbb{D}(\mathcal{T})$ is a complex
manifold in the sense of Definition~\ref{4.10:def-holomorphic-maps}.
\end{definition}

\begin{lemma}[Holomorphy of the averaged data along a tower]\label{4.10:lem-averaged-data}
Let $\mathcal{T}=\mathcal{T}_{\varrho}(Z,n)$ be a truncated tower of depth $n$ of the region $Z$, with constraint bonds
$\Lambda_q(\mathcal{T})$, $0\le q\le n$, and data space $\mathbb{D}(\mathcal{T})$, as in
Definition~\ref{4.10:not-truncated-towers}. Assume the following statement about Ba{\l}aban's block-averaging
operation \cite{B-CMP98}. It is taken in the form in which this operation prescribes the data in the constraints
\cite[(3)]{B-CMP102b}, and in which averages of $G^c$-valued configurations enter \cite[(1.15)]{B-CMP109}.
\begin{itemize}
\item[] For every $p\ge1$ and every level-$p$ bond $c$ there are a finite set $\mathrm{supp}_p(c)$ of fine bonds and
an open subset $\mathcal{A}_c\subset\mathrm{Conf}(\mathrm{supp}_p(c))$ with the following two properties.
First, the $p$-fold block average $(\mathcal{M}^pW)(c)\in G^c$ of a configuration $W$ is defined exactly when
$W|_{\mathrm{supp}_p(c)}\in\mathcal{A}_c$. Second, $W\mapsto(\mathcal{M}^pW)(c)$ is a holomorphic map
$\mathcal{A}_c\to G^c$ in the sense of Definition~\ref{4.10:def-holomorphic-maps}. Moreover,
$\mathrm{supp}_p(c)\subset\mathrm{bd}(Z)$ for every $c\in\Lambda_p(\mathcal{T})$ with $1\le p\le n$.
\end{itemize}
Put
\begin{equation}\label{4.10:eq-averaged-data-1}
\mathcal{A}(\mathcal{T}):=\bigl\{W\in\mathrm{Conf}(\mathrm{bd}(Z)):\ W|_{\mathrm{supp}_p(c)}\in\mathcal{A}_c
\text{ for all } c\in\Lambda_p(\mathcal{T}),\ 1\le p\le n\bigr\}.
\end{equation}
This is the set of fields on $Z$ whose averages along $\mathcal{T}$ are defined. For
$W\in\mathcal{A}(\mathcal{T})$ define the \emph{averaged data}
\begin{equation}\label{4.10:eq-averaged-data-2}
\dot{\mathcal{M}}_{\mathcal{T}}W:=\Bigl(W|_{\Lambda_0(\mathcal{T})},\,
(\mathcal{M}^1W)|_{\Lambda_1(\mathcal{T})},\,\dots,\,(\mathcal{M}^nW)|_{\Lambda_n(\mathcal{T})}\Bigr)
\in\mathbb{D}(\mathcal{T}).
\end{equation}
Then $\mathcal{A}(\mathcal{T})$ is open in $\mathrm{Conf}(\mathrm{bd}(Z))$, and the map
$\dot{\mathcal{M}}_{\mathcal{T}}:\mathcal{A}(\mathcal{T})\to\mathbb{D}(\mathcal{T})$ is holomorphic.
\end{lemma}

\begin{proof}
The index set consists of the pairs $(p,c)$ with $1\le p\le n$ and $c\in\Lambda_p(\mathcal{T})$. It is finite,
since each $\Lambda_p(\mathcal{T})$ is a finite set of bonds.

Fix such a pair $(p,c)$. By the hypothesis, $\mathrm{supp}_p(c)\subset\mathrm{bd}(Z)$. Hence the restriction
\begin{equation*}
r_c:\ \mathrm{Conf}(\mathrm{bd}(Z))\to\mathrm{Conf}(\mathrm{supp}_p(c)),\qquad W\mapsto W|_{\mathrm{supp}_p(c)},
\end{equation*}
is defined. It is holomorphic on $\mathrm{Conf}(\mathrm{bd}(Z))$ by fact (H2) of
\eqref{4.10:eq-holomorphic-maps-a}. The hypothesis also gives that $\mathcal{A}_c$ is open in
$\mathrm{Conf}(\mathrm{supp}_p(c))$. Therefore, by fact (H4) of \eqref{4.10:eq-holomorphic-maps-a}, the preimage
$r_c^{-1}(\mathcal{A}_c)$ is open in $\mathrm{Conf}(\mathrm{bd}(Z))$. By its definition
\eqref{4.10:eq-averaged-data-1},
\begin{equation*}
\mathcal{A}(\mathcal{T})=\bigcap_{p=1}^{n}\ \bigcap_{c\in\Lambda_p(\mathcal{T})} r_c^{-1}(\mathcal{A}_c).
\end{equation*}
This is a finite intersection of open sets, so it is open.

Now consider the components of $\dot{\mathcal{M}}_{\mathcal{T}}$. The data space $\mathbb{D}(\mathcal{T})$ is the
configuration space of the disjoint union of the sets $\Lambda_q(\mathcal{T})$, $0\le q\le n$. A point of it is
therefore the family of its values on the bonds of these sets, and \eqref{4.10:eq-averaged-data-2} lists the
components of $\dot{\mathcal{M}}_{\mathcal{T}}$. They are of two kinds.

The first kind is the skin component $W\mapsto W|_{\Lambda_0(\mathcal{T})}$. This is a restriction map, and it is
holomorphic by fact (H2).

The second kind consists of the components indexed by $c\in\Lambda_p(\mathcal{T})$ with $1\le p\le n$. Each is the
composite $W\mapsto(\mathcal{M}^p\,\cdot\,)(c)\circ r_c(W)$. On $\mathcal{A}(\mathcal{T})$ the map $r_c$ is
holomorphic and takes values in $\mathcal{A}_c$, by the definition \eqref{4.10:eq-averaged-data-1}. The hypothesis
makes the average $V\mapsto(\mathcal{M}^pV)(c)$ holomorphic on $\mathcal{A}_c$. Hence the composite is defined on
$\mathcal{A}(\mathcal{T})$, and it is holomorphic there by facts (H2) and (H4) of
\eqref{4.10:eq-holomorphic-maps-a} together with the holomorphy of the average on $\mathcal{A}_c$ given by the
hypothesis.

Finally, a map into a product is holomorphic when each of its components is. Since every component of
$\dot{\mathcal{M}}_{\mathcal{T}}$ is holomorphic on $\mathcal{A}(\mathcal{T})$, the map
$\dot{\mathcal{M}}_{\mathcal{T}}:\mathcal{A}(\mathcal{T})\to\mathbb{D}(\mathcal{T})$ is holomorphic.
\end{proof}

\begin{definition}[Ba{\l}aban's analytic continuation of the constrained minimiser
{\cite[Theorem~1, \S G, Proposition~9]{B-CMP102b}}]
\label{4.10:def-minimiser-continuation}
Let $Z$, the tower rule $\varrho\in\{\mathrm{I},\mathrm{II}\}$, the depth $n$, the truncated tower
$\mathcal{T}=\mathcal{T}_{\varrho}(Z,n)$, its set of constraint bonds $\mathfrak{B}(\mathcal{T})$, the
data space $\mathbb{D}(\mathcal{T})=\mathrm{Conf}(\mathfrak{B}(\mathcal{T}))$ and the set
$\mathrm{bd}(Z)$ of positively oriented fine bonds of $Z$ be as in
Definition~\ref{4.10:not-truncated-towers}. Let $M_N(\mathbb{C})$, its operator norm $|\cdot|$, the
groups $G=SU(N)$ and $G^c=SL(N,\mathbb{C})$, the Lie algebra $\mathfrak{g}^c$ and holomorphy of maps
between complex configuration spaces be as in Definition~\ref{4.10:def-holomorphic-maps}. For a bond
$c$, coarse or fine, we write $c_-$ and $c_+$ for its initial and its final site.

\emph{The minimiser at $G$-valued data.} By \cite[Theorem~1]{B-CMP102b} there are positive
constants, depending on $d$ and $L$ only, among them a threshold which is called $a_1$ there and
which we write $\bar{\varepsilon}$ (the letter $a_1$ being reserved for a radius of the one-power
class), with the following property. Let $V_0$ be a $G$-valued configuration on the constraint
bonds which satisfies condition \cite[(7)]{B-CMP102b} with a constant $\varepsilon_1\le\bar{\varepsilon}$,
that is,
\begin{equation*}
|(\partial V_0)(p')-1|<\varepsilon_1
\end{equation*}
for every plaquette $p'$ of the constraint lattices, with the convention of \cite{B-CMP102b} at
plaquettes meeting two levels. Then there is a minimal orbit of the Wilson action over the fields on
the bonds of $\Omega_0$ whose averaged data are $V_0$. This orbit is unique as a critical orbit in
the space considered in \cite[Theorem~1]{B-CMP102b}, and its elements have the regularity
\cite[(9)--(10)]{B-CMP102b}. The minimal configuration in the axial gauge is denoted $U_k(V_0)$ in
\cite{B-CMP102b}.

\emph{Admitted data.} We write
$\mathcal{R}(\mathcal{T})\subset G^{\mathfrak{B}(\mathcal{T})}$ for the set of $G$-valued data admitted
in \cite[\S G]{B-CMP102b}, namely those which satisfy \cite[(7)]{B-CMP102b} with $\varepsilon_1$ below
the threshold of that section; this threshold is a constant depending on $d$ and $L$ only.

\begin{enumerate}
\item[(a)] \emph{The continuation at a fixed datum.} Let $V_0\in\mathcal{R}(\mathcal{T})$. We write
$\mathcal{B}(V_0)$ for the set of potentials
$B\in(\mathfrak{g}^c)^{\mathfrak{B}(\mathcal{T})}$ which satisfy \cite[(172)]{B-CMP102b}. It is an
open subset of the complex vector space $(\mathfrak{g}^c)^{\mathfrak{B}(\mathcal{T})}$, it contains
$0$, and $|B(c)|<\pi$ for every $B\in\mathcal{B}(V_0)$ and every $c\in\mathfrak{B}(\mathcal{T})$.
We write
\begin{equation*}
\Phi_{V_0}:\mathcal{B}(V_0)\to\mathrm{Conf}(\mathrm{bd}(Z))
\end{equation*}
for the map which sends $B$ to the axial-gauge configuration that \cite[\S G]{B-CMP102b} constructs
from the datum $V'V_0$, $V'=e^{iB}$; of it Ba{\l}aban writes that
``we obtain a gauge field configuration in the axial gauge, which we denote also by
$U_k(V'V_0)$''. It is stated there to be an analytic function of $B$, so that $\Phi_{V_0}$ is
holomorphic. Two features used in this item are not stated by Ba{\l}aban in these words but are
read from condition \cite[(172)]{B-CMP102b} itself: that this condition consists of strict
inequalities, so that $\mathcal{B}(V_0)$ is open; and that it forces $|B(c)|<\pi$ in the operator
norm. In \cite[\S G]{B-CMP102b} these potentials are described by saying that $V'$ is ``small''
and lies in a ``small neighborhood of the identity''.

\item[(b)] \emph{One function of the datum.} The maps $\Phi_{V_0}$ are restrictions of one function
of the datum. Suppose $V_0,V_0'\in\mathcal{R}(\mathcal{T})$, $B\in\mathcal{B}(V_0)$,
$B'\in\mathcal{B}(V_0')$ and $e^{iB}V_0=e^{iB'}V_0'$ in $\mathbb{D}(\mathcal{T})$. Then
$\Phi_{V_0}(B)=\Phi_{V_0'}(B')$. Ba{\l}aban states that the configuration of (a) is ``an analytic
function of $V$'' \cite[\S G]{B-CMP102b}, and that it ``has an extension to an analytic function of
$G^c$-valued small configurations $V$'' \cite[Proposition~9]{B-CMP102b}. Accordingly we put
\begin{equation}\label{4.10:eq-minimiser-continuation-1}
\begin{gathered}
\mathcal{S}(\mathcal{T}):=\bigl\{e^{iB}V_0:\ V_0\in\mathcal{R}(\mathcal{T}),\
B\in\mathcal{B}(V_0)\bigr\}\subset\mathbb{D}(\mathcal{T}),\\
\mathcal{U}_{\mathcal{T}}(e^{iB}V_0):=\Phi_{V_0}(B).
\end{gathered}
\end{equation}
By the consistency just stated, \eqref{4.10:eq-minimiser-continuation-1} defines a map
$\mathcal{U}_{\mathcal{T}}:\mathcal{S}(\mathcal{T})\to\mathrm{Conf}(\mathrm{bd}(Z))$. The set
$\mathcal{S}(\mathcal{T})$ is the neighbourhood of $G$-valued data of
\cite[Proposition~9]{B-CMP102b}; it is this set that is meant wherever this chapter refers to the
neighbourhood of $G$-valued data of that proposition. Since $0\in\mathcal{B}(V_0)$ and
$e^{i0}V_0=V_0$, we have $\mathcal{R}(\mathcal{T})\subset\mathcal{S}(\mathcal{T})$.

\item[(c)] \emph{Agreement with the minimiser.} For $V_0\in\mathcal{R}(\mathcal{T})$ the
configuration $\mathcal{U}_{\mathcal{T}}(V_0)=\Phi_{V_0}(0)$ is the axial-gauge minimal
configuration $U_k(V_0)$ of \cite[Theorem~1]{B-CMP102b}: in the words of
\cite[\S G]{B-CMP102b}, ``for $V$ with values in $G$ it coincides with the minimal configuration
constructed in the previous sections''. Thus $\mathcal{U}_{\mathcal{T}}$ is the continued minimiser
of \cite[Proposition~9]{B-CMP102b} to which this chapter refers elsewhere.

\item[(d)] \emph{Gauge extension.} Let $\Gamma(\mathcal{T})$ be the group of $G^c$-valued gauge
transformations of the data considered in \cite[(181)]{B-CMP102b}, that is, of $G^c$-valued
functions $v$ on the sites of the constraint lattices. Such a $v$ acts on $\mathbb{D}(\mathcal{T})$
by $V\mapsto V^v$, and
it determines the fine-lattice gauge transformation $\tilde{v}$ on $Z$ of \cite[\S G]{B-CMP102b},
which is constant on blocks, and thereby the map $U\mapsto U^{\tilde{v}}$ of
$\mathrm{Conf}(\mathrm{bd}(Z))$; here
\begin{equation*}
\begin{aligned}
(V^v)(c)&:=v(c_-)\,V(c)\,v(c_+)^{-1}, && c\in\mathfrak{B}(\mathcal{T}),\\
(U^{\tilde{v}})(b)&:=\tilde{v}(b_-)\,U(b)\,\tilde{v}(b_+)^{-1}, && b\in\mathrm{bd}(Z).
\end{aligned}
\end{equation*}
For $v\in\Gamma(\mathcal{T})$ write
$\mathcal{S}(\mathcal{T})^v:=\{V^v:\ V\in\mathcal{S}(\mathcal{T})\}$, and put
\begin{equation*}
\mathcal{O}(\mathcal{T}):=\bigcup_{v\in\Gamma(\mathcal{T})}\mathcal{S}(\mathcal{T})^v .
\end{equation*}
On $\mathcal{O}(\mathcal{T})$ the prescription
\begin{equation*}
\mathcal{U}_{\mathcal{T}}(V^v):=\mathcal{U}_{\mathcal{T}}(V)^{\tilde{v}},
\qquad V\in\mathcal{S}(\mathcal{T}),\ v\in\Gamma(\mathcal{T}),
\end{equation*}
defines one function, which for $v$ the identity is the map of (b). Ba{\l}aban writes: ``We can
extend it to all $G^c$-valued $v$ treating the equality as a definition \dots\ It is an analytic
function on the union of the set of orbits, and it satisfies (181)'' \cite[\S G, (181)]{B-CMP102b};
see also \cite[Proposition~9]{B-CMP102b}.
\end{enumerate}
The content of items (a)--(d) is collected, in this order, in the four clauses of the following
display; its first clause is item (a), its second item (b), its third item (c) and its fourth
item (d). For all $V_0,V_0'\in\mathcal{R}(\mathcal{T})$:
\begin{equation}\label{4.10:eq-minimiser-continuation-a}
\begin{aligned}
&\text{(a)}\quad 0\in\mathcal{B}(V_0),\ \ \mathcal{B}(V_0)\ \text{open in}\
(\mathfrak{g}^c)^{\mathfrak{B}(\mathcal{T})},\\
&\qquad |B(c)|<\pi\ \ \text{for}\ B\in\mathcal{B}(V_0),\ c\in\mathfrak{B}(\mathcal{T}),\\
&\qquad \Phi_{V_0}:\mathcal{B}(V_0)\to\mathrm{Conf}(\mathrm{bd}(Z))\ \text{holomorphic};\\
&\text{(b)}\quad \text{for}\ B\in\mathcal{B}(V_0),\ B'\in\mathcal{B}(V_0'):\quad
e^{iB}V_0=e^{iB'}V_0'\ \Longrightarrow\ \Phi_{V_0}(B)=\Phi_{V_0'}(B'),\\
&\qquad \mathcal{U}_{\mathcal{T}}(e^{iB}V_0)=\Phi_{V_0}(B)\ \ \text{on}\ \mathcal{S}(\mathcal{T});\\
&\text{(c)}\quad \mathcal{U}_{\mathcal{T}}(V_0)=\Phi_{V_0}(0)=U_k(V_0);\\
&\text{(d)}\quad \mathcal{U}_{\mathcal{T}}(V^v)=\mathcal{U}_{\mathcal{T}}(V)^{\tilde{v}}\ \
\text{for}\ V\in\mathcal{S}(\mathcal{T}),\ v\in\Gamma(\mathcal{T}),\\
&\qquad \text{one function on}\ \mathcal{O}(\mathcal{T})\ \text{extending}\
\mathcal{U}_{\mathcal{T}}\ \text{of (b)}.
\end{aligned}
\end{equation}
\end{definition}

The existence and properties of the minimiser at $G$-valued data are
\cite[Theorem~1 with condition~(7)]{B-CMP102b}. The construction of the continuation, its
analyticity, its independence of the decomposition $e^{iB}V_0$ of the datum, its agreement with the
minimiser at $G$-valued data and its gauge extension are \cite[\S G, (172), (181),
Proposition~9]{B-CMP102b}. Apart from the two features read from condition \cite[(172)]{B-CMP102b}
in item (a), they are stated there in the words quoted. The proofs are Ba{\l}aban's and are not
reproduced.

\begin{definition}[The tower quantities at complex data]\label{4.10:def-tower-quantities}
Fix a tower rule $\varrho\in\{\mathrm{I},\mathrm{II}\}$, a domain $Z$ and a depth $n$. Let $\mathcal{T}=\mathcal{T}_{\varrho}(Z,n)$ be the truncated tower of depth $n$, and let $W\in\mathrm{Conf}(\mathrm{bd}(Z))$.

The following objects are taken from two earlier statements.
\begin{itemize}
\item The set $\mathcal{A}(\mathcal{T})$ and the averaged data $\dot{\mathcal{M}}_{\mathcal{T}}W$ are those of Lemma~\ref{4.10:lem-averaged-data}.
\item The neighbourhood $\mathcal{S}(\mathcal{T})$ of $G$-valued data, and the continuation $\mathcal{U}_{\mathcal{T}}$ regarded as one function on it, are those of Definition~\ref{4.10:def-minimiser-continuation}.
\end{itemize}

For complex data the field $U_n(Z;\dot{\mathcal{M}}W)$ for the rule $\varrho$ is the analytic continuation of the constrained minimiser given by \cite[Proposition~9]{B-CMP102b}, wherever the averaged data lie in that proposition's neighbourhood of $G$-valued data.

In this sense we say that $U_n(Z;\dot{\mathcal{M}}W)$ \emph{exists} if and only if
\begin{equation*}
W\in\mathcal{A}(\mathcal{T})
\quad\text{and}\quad
\dot{\mathcal{M}}_{\mathcal{T}}W\in\mathcal{S}(\mathcal{T}).
\end{equation*}
In that case we set
\begin{equation}\label{4.10:eq-tower-quantities-1}
U_n(Z;\dot{\mathcal{M}}W)
:=\mathcal{U}_{\mathcal{T}}\bigl(\dot{\mathcal{M}}_{\mathcal{T}}W\bigr)
\in\mathrm{Conf}(\mathrm{bd}(Z)).
\end{equation}

The \emph{existence set} of the tower $\mathcal{T}$ is
\begin{equation}\label{4.10:eq-tower-quantities-2}
\mathcal{V}(\mathcal{T})
:=\bigl\{W\in\mathcal{A}(\mathcal{T}):\
\dot{\mathcal{M}}_{\mathcal{T}}W\in\mathcal{S}(\mathcal{T})\bigr\}.
\end{equation}

The tower quantities may instead be read with the orbit extension of \cite[(181)]{B-CMP102b}. In that reading $\mathcal{S}(\mathcal{T})$ is replaced by the set $\mathcal{O}(\mathcal{T})$ that this orbit extension attaches to the tower $\mathcal{T}$, in the existence condition and in \eqref{4.10:eq-tower-quantities-2}; here the letter $\mathcal{O}$ does not denote a generic open set. A corollary stated below shows that nothing in what follows changes under this replacement.

The dependence on the rule $\varrho$ is carried by the tower $\mathcal{T}$. It is suppressed in the notation $U_n(Z;\cdot)$. Clause (iv) of this chapter, which is distinct from clause (iv) of Theorem~0, is imposed for both rules, $\varrho=\mathrm{I}$ and $\varrho=\mathrm{II}$.
\end{definition}

\begin{definition}[The current as a holomorphic local function {\cite[(1.8), (1.15)]{B-CMP109}}]
\label{4.10:def-current-holomorphic}
Let $\xi>0$ be a spacing parameter and let $\Omega$ be a region (in this definition only, $\Omega$
denotes such a region). The definition of the current in
\cite[(1.8)]{B-CMP109} assigns to a configuration $U\in\mathrm{Conf}(\mathrm{bd}(\Omega))$ a
$\mathfrak{g}^c$-valued function
\begin{equation*}
\mathcal{J}_{\xi}(U)\colon \mathrm{bd}_J(\Omega)\to\mathfrak{g}^c .
\end{equation*}
Here $\mathrm{bd}_J(\Omega)\subset\mathrm{bd}(\Omega)$ is the set of bonds $b$ for which the stencil
of the defining formula at $b$, that is, the set of bonds whose variables enter the value at $b$,
lies in $\Omega$. It contains the bonds at the points of $\Omega$ lying two layers of cubes inside
$\Omega$, and these are the bonds at which the current is evaluated in what follows.

The value $\mathcal{J}_{\xi}(U)(b)$ at a bond $b\in\mathrm{bd}_J(\Omega)$ has the following structure.
\begin{itemize}
\item It is computed from finitely many bond variables $U(b')$ and $U(b')^{-1}$, with $b'$ near $b$.
\item From these variables a polynomial expression with $\xi$-dependent coefficients is formed,
namely covariant difference quotients applied to the matrices
\begin{equation*}
\operatorname{im}\partial U(P)=\frac{1}{2i}\bigl(\partial U(P)-\partial U(P)^{-1}\bigr),
\end{equation*}
where $\partial U(P)$ is the holonomy of $U$ around the plaquette $P$.
\item The projection $\Pi_{\mathfrak{g}^c}$ of $M_N(\mathbb{C})$ onto $\mathfrak{g}^c$ used in
\cite[(1.8)]{B-CMP109} is then applied to this
expression; $\Pi_{\mathfrak{g}^c}$ is taken to be complex-linear.
\end{itemize}
Since $\mathcal{J}_{\xi}(U)(b)$ is thus a polynomial expression in the bond variables and their
inverses followed by a complex-linear map, facts $\mathrm{(H2)}$ and $\mathrm{(H4)}$ of the
list~\eqref{4.10:eq-holomorphic-maps-a} in Definition~\ref{4.10:def-holomorphic-maps} show that
\begin{equation}\label{4.10:eq-current-holomorphic-1}
\mathrm{Conf}(\mathrm{bd}(\Omega))\ni U\longmapsto \mathcal{J}_{\xi}(U)
\in(\mathfrak{g}^c)^{\mathrm{bd}_J(\Omega)}
\end{equation}
is a holomorphic map.

Now let $Z$, $n$, the rule $\varrho$ and the tower $\mathcal{T}=\mathcal{T}_{\varrho}(Z,n)$ be as in
Definition~\ref{4.10:def-tower-quantities}. For $W\in\mathcal{V}(\mathcal{T})$ the tower field
$U_n(Z;\dot{\mathcal{M}}W)$ of Definition~\ref{4.10:def-tower-quantities} exists, and the
\emph{current} at $W$ is
\begin{equation}\label{4.10:eq-current-holomorphic-2}
J_n(Z;\dot{\mathcal{M}}W):=\mathcal{J}_{L^{-n}}\bigl(U_n(Z;\dot{\mathcal{M}}W)\bigr)
\in(\mathfrak{g}^c)^{\mathrm{bd}_J(Z)} .
\end{equation}
Thus $J_n(Z;\dot{\mathcal{M}}W)$ is obtained from the defining formula of
\cite[(1.8)]{B-CMP109} by inserting $U_n(Z;\dot{\mathcal{M}}W)$ as the configuration and $L^{-n}$
as the spacing parameter $\xi$. This is the current of \cite[(1.8), (1.15)]{B-CMP109}.
\end{definition}

\begin{definition}[Decomposed one-power representatives and their ambient spaces]
\label{4.10:def-decomposed-representatives}
This definition makes explicit the one-power hereditary class of this chapter, together with the
spaces on which it lives. The lattice, the units, the classes
$\mathcal{D}_i$, the two-layer interior $Z^{\wr}$ and the spacing $\xi_i=L^{-i}$ are those of
Notation~\ref{4.10:not-lattices-regions}. The groups $G$, $G^c$, the algebras $\mathfrak{g}$,
$\mathfrak{g}^c$, the configuration spaces $\mathrm{Conf}(\cdot)$, the holonomy $\partial W(P)$ and
the holomorphy facts \textup{(H1)}--\textup{(H5)} of \eqref{4.10:eq-holomorphic-maps-a} are those
of Definition~\ref{4.10:def-holomorphic-maps}.

Let $j\ge 1$, $X\in\mathcal{D}_{j-1}$, and let $a_0,a_1,\gamma_0>0$ be radii; when $\gamma_0$ is
not written, $\gamma_0=a_0$. Consider a triple $(U,A',\mathbf{J})$ consisting of
\begin{itemize}
\item a $G$-valued configuration $U$ on $\mathrm{bd}(X)$,
\item a potential $A'\in(\mathfrak{g}^c)^{\mathrm{bd}(X)}$,
\item a free current $\mathbf{J}\in(\mathfrak{g}^c)^{\mathrm{bd}(X)}$,
\end{itemize}
and put $U':=e^{i\xi_jA'}$ and $\mathbf{U}:=e^{i\xi_jA'}U$ (bondwise product). The clauses at level
$j$ on $X$ for $(U,A',\mathbf{J})$ are the following:
\begin{equation}\label{4.10:eq-decomposed-representatives-b}
\begin{aligned}
\text{(R1)}\quad & |\partial U(P)-1|<a_0\xi_j^2;\\
& \text{for every cube } \square\subset X \text{ of side at most } c_{\square}LM
\text{ there are}\\
& u\colon\mathrm{s}(\square)\to G \text{ and } A\in\mathfrak{g}^{\mathrm{bd}(\square)}
\text{ with } U^u=e^{i\xi_jA}\ \text{on }\mathrm{bd}(\square),\\
& |A|,\ |\nabla^{\xi_j}A|<c_{\square}LMB_0a_0\ \text{on }\square;\\
\text{(R2)}\quad & |A'(b)|<a_1,\qquad |\nabla^{\xi_j}_UA'|<a_1\ \text{on }X;\\
\text{(R3)}\quad & |\partial\mathbf{U}(P)-1|<a_0\xi_j^2,\qquad |\mathbf{J}(b)|<\gamma_0 .
\end{aligned}
\end{equation}
Here $c_{\square}$ and $B_0$ are the constants of \cite[(1.12)]{B-CMP109}, the latter written $B$
there; the subscript is added because in this section the letter $B$ is reserved for potentials.
The quantifiers are as follows. In (R1), the first inequality is required for every plaquette $P$
of $X$ (the curvature member of (R1)); the remainder of (R1) (its gauge member,
\cite[(1.12)]{B-CMP109}) requires that for every cube $\square\subset X$ of the fine lattice of
side at most $c_{\square}LM$ level-$j$ units there be a $G$-valued gauge transformation $u$ on the
set $\mathrm{s}(\square)$ of sites of $\square$ and a potential $A\in\mathfrak{g}^{\mathrm{bd}(\square)}$
such that
\begin{equation*}
U^u(b):=u(b_-)U(b)u(b_+)^{-1}=e^{i\xi_jA(b)}\qquad (b\in\mathrm{bd}(\square)),
\end{equation*}
and such that the displayed bounds on $A$ and $\nabla^{\xi_j}A$ hold on $\square$. In (R2), the
first inequality holds
for every $b\in\mathrm{bd}(X)$ and the second everywhere on $X$. In (R3), the first inequality
holds for every plaquette $P$ of $X$ and the second for every $b\in\mathrm{bd}(X)$.

Here $\nabla^{\xi_j}$ is $\xi_j^{-1}$ times the fine lattice difference, applied to the function
$A$ on parallel neighbouring bonds of $\square$, and $\nabla^{\xi_j}_U$ is its $U$-covariant
version. Of these operators only the following property enters: each component of
$\nabla^{\xi_j}A$ is a continuous function of $A$, and each component of $\nabla^{\xi_j}_UA'$ is a
continuous function of $(U,A')$; indeed they are finite differences, with conjugations by bond
variables of $U$ in the covariant case.

\emph{Sub-towers.} The index set is
\begin{equation*}
\mathfrak{T}_j(X):=\{(p,Z):\ 1\le p\le j,\ Z\in\mathcal{D}_{p-1},\ Z\subset X\},
\end{equation*}
a finite set. For $(p,Z)\in\mathfrak{T}_j(X)$, a constant $c>0$ and a configuration $\mathbf{W}$ on
a set of bonds containing $\mathrm{bd}(Z)$, the clause $(\mathrm{iv})_{p,Z}[c]$ for $\mathbf{W}$
requires the following: for both rules $\varrho\in\{\mathrm{I},\mathrm{II}\}$ and every
$n=1,\dots,p$, the tower field $U_n(Z;\dot{\mathcal{M}}(\mathbf{W}|_Z))$ and the tower current
$J_n(Z;\dot{\mathcal{M}}(\mathbf{W}|_Z))$ exist in the sense of
Definition~\ref{4.10:def-tower-quantities} (the current being that of
Definition~\ref{4.10:def-current-holomorphic}), and
\begin{equation}\label{4.10:eq-decomposed-representatives-a}
\begin{aligned}
\bigl|\partial U_n(Z;\dot{\mathcal{M}}(\mathbf{W}|_Z))(P)-1\bigr|&<c\,\xi_p^2
&&\text{for } P\in\mathrm{Pl}(Z^{\wr}),\\
\bigl|J_n(Z;\dot{\mathcal{M}}(\mathbf{W}|_Z))(b)\bigr|&<c\,(L^n\xi_p)^2
&&\text{for } b\in\mathrm{Bd}(Z^{\wr}),
\end{aligned}
\end{equation}
where $\mathrm{Pl}(Z^{\wr})$ and $\mathrm{Bd}(Z^{\wr})$ are the finite sets of plaquettes and of
bonds at the points of $Z^{\wr}$.

\emph{Representatives.} A \emph{decomposed one-power representative} at
$(j,X;a_0,a_1,\gamma_0)$ is a triple $(U,A',\mathbf{J})$ as above which satisfies (R1), (R2) and
(R3) of \eqref{4.10:eq-decomposed-representatives-b} and, for every $(p,Z)\in\mathfrak{T}_j(X)$,
the clause $(\mathrm{iv})_{p,Z}[a_0L^{-(j-p)}]$ both for $\mathbf{W}=\mathbf{U}=e^{i\xi_jA'}U$
(the \emph{$\mathbf{U}$-half}) and for $\mathbf{W}=U$ (the \emph{$U$-half}). A \emph{one-power
representative} is a pair $(\mathbf{U},\mathbf{J})$ admitting such a decomposition, that is, a pair
of the form $(e^{i\xi_jA'}U,\mathbf{J})$ with $(U,A',\mathbf{J})$ a decomposed one-power
representative. The complex gauge group consists of the maps $g$ from the sites of $X$ to $G^c$;
it acts by $(\mathbf{U},\mathbf{J})\mapsto(\mathbf{U}^g,\mathbf{J}^g)$, where, as in
\cite[(1.10)]{B-CMP109},
\begin{equation*}
\mathbf{U}^g(b)=g(b_-)\mathbf{U}(b)g(b_+)^{-1},\qquad
\mathbf{J}^g(b)=g(b_-)\mathbf{J}(b)g(b_-)^{-1}.
\end{equation*}
The set $U^{\Sigma}_j(X;a_0,a_1,\gamma_0)$ is the union of the orbits of
the one-power representatives under the complex gauge group.

\emph{Ambient spaces.} The space of decomposed pairs is
\begin{equation*}
\mathfrak{P}(X):=G^{\mathrm{bd}(X)}\times(\mathfrak{g}^c)^{\mathrm{bd}(X)}
\times(\mathfrak{g}^c)^{\mathrm{bd}(X)},
\end{equation*}
whose points are the triples $(U,A',\mathbf{J})$, with the product topology of the compact group
$G^{\mathrm{bd}(X)}$ and the two finite-dimensional vector spaces. The space of pairs is the complex
manifold $\mathfrak{Q}(X):=\mathrm{Conf}(\mathrm{bd}(X))\times(\mathfrak{g}^c)^{\mathrm{bd}(X)}$.
Define $\rho:\mathfrak{P}(X)\to\mathfrak{Q}(X)$ by $\rho(U,A',\mathbf{J}):=(e^{i\xi_jA'}U,\mathbf{J})$,
and, for $Z\subset X$, define $\sigma_Z,\sigma'_Z:\mathfrak{P}(X)\to\mathrm{Conf}(\mathrm{bd}(Z))$ by
\begin{equation*}
\sigma_Z(U,A',\mathbf{J}):=(e^{i\xi_jA'}U)|_{\mathrm{bd}(Z)},\qquad
\sigma'_Z(U,A',\mathbf{J}):=U|_{\mathrm{bd}(Z)} .
\end{equation*}
Then $\rho$ is continuous, $\sigma_Z$ and $\sigma'_Z$ are
continuous, and for fixed $U$ the map $A'\mapsto\sigma_Z(U,A',\mathbf{J})$ is holomorphic.
\end{definition}

\begin{proof}
Finiteness of $\mathfrak{T}_j(X)$: the index $p$ ranges over the finitely many integers
$1,\dots,j$, and for each $p$ the second entry $Z$ is a subset of $X$; since $X$ is a subset of
the fine lattice $\mathbb{T}$, which is a finite discrete torus, $X$ has finitely many subsets.
Hence $\mathfrak{T}_j(X)$ is finite.

Continuity of $\rho$: the second component of $\rho$ is the coordinate projection
$(U,A',\mathbf{J})\mapsto\mathbf{J}$, which is continuous for the product topology. The first
component is $(U,A')\mapsto e^{i\xi_jA'}U$, the bondwise product of the exponential of the
potential $A'$ with the configuration $U\in G^{\mathrm{bd}(X)}\subset\mathrm{Conf}(\mathrm{bd}(X))$;
it is continuous by \textup{(H3)} of \eqref{4.10:eq-holomorphic-maps-a}. A map into a product is
continuous when its components are, so $\rho$ is continuous.

Continuity of $\sigma_Z$ and $\sigma'_Z$: the restriction map
$\mathrm{Conf}(\mathrm{bd}(X))\to\mathrm{Conf}(\mathrm{bd}(Z))$, $W\mapsto W|_{\mathrm{bd}(Z)}$, is a
coordinate projection of $(G^c)^{\mathrm{bd}(X)}$ onto $(G^c)^{\mathrm{bd}(Z)}$, hence continuous.
The map $\sigma_Z$ is this restriction composed with the first component of $\rho$, which is
continuous by the preceding paragraph; so $\sigma_Z$ is continuous. The map $\sigma'_Z$ is the
restriction composed with the coordinate projection $(U,A',\mathbf{J})\mapsto U$ and the inclusion
$G^{\mathrm{bd}(X)}\subset\mathrm{Conf}(\mathrm{bd}(X))$; so $\sigma'_Z$ is continuous.

Holomorphy of $\sigma_Z$ in $A'$ for fixed $U$: for fixed $U$, the map $A'\mapsto e^{i\xi_jA'}U$
is holomorphic by \textup{(H2)} and \textup{(H3)} of \eqref{4.10:eq-holomorphic-maps-a}, and the
restriction to $\mathrm{bd}(Z)$ is a coordinate projection, hence holomorphic; the composition
$A'\mapsto\sigma_Z(U,A',\mathbf{J})$ is therefore holomorphic.
\end{proof}

\begin{lemma}[The principal logarithm of a matrix]\label{4.10:lem-principal-logarithm}
We use the algebra $M_N(\mathbb{C})$, its operator norm $|\cdot|$ and the spectrum
$\mathrm{spec}$ as in Definition~\ref{4.10:def-holomorphic-maps}. In this lemma
$\operatorname{Tr}$ denotes the ordinary trace of a matrix, the sum of its diagonal entries.
Let
\begin{equation*}
\mathcal{L}:=\{g\in M_N(\mathbb{C}):\ \mathrm{spec}(g)\cap(-\infty,0]=\emptyset\},
\end{equation*}
and let $\log z$ denote the principal branch of the logarithm on $\mathbb{C}\setminus(-\infty,0]$.
\begin{enumerate}
\item[(i)] $\mathcal{L}$ is open in $M_N(\mathbb{C})$. Let $g\in\mathcal{L}$, and let $\eta$ be a cycle
in $\mathbb{C}\setminus(-\infty,0]$ winding once around each eigenvalue of $g$. The map
$\log:\mathcal{L}\to M_N(\mathbb{C})$ given by
\begin{equation}\label{4.10:eq-principal-logarithm-1}
\log g:=\frac{1}{2\pi i}\oint_{\eta}(\log z)\,(z-g)^{-1}\,dz
\end{equation}
is well defined and holomorphic.
\item[(ii)] $\exp(\log g)=g$ for every $g\in\mathcal{L}$.
\item[(iii)] If $Y\in M_N(\mathbb{C})$ and $\mathrm{spec}(Y)\subset\{z:|\operatorname{Im}z|<\pi\}$,
then $e^{Y}\in\mathcal{L}$ and $\log(e^{Y})=Y$.
\item[(iv)] For $g\in\mathcal{L}$ one has $e^{\operatorname{Tr}\log g}=\det g$; hence
$\operatorname{Tr}\log g\in 2\pi i\mathbb{Z}$ when $\det g=1$.
\item[(v)] If $B\in M_N(\mathbb{C})$ and $|B|<\pi$, then
$\mathrm{spec}(iB)\subset\{z:|\operatorname{Im}z|<\pi\}$.
\item[(vi)] If $g\in\mathcal{L}$ is unitary, then $\log g$ is anti-Hermitian.
\end{enumerate}
\end{lemma}

\begin{proof}
All six assertions are properties of the holomorphic functional calculus. For $g\in M_N(\mathbb{C})$
and $f$ holomorphic on an open set $\mathcal{O}\supset\mathrm{spec}(g)$ we set
\begin{equation}\label{4.10:eq-principal-logarithm-2}
f(g):=\frac{1}{2\pi i}\oint_{\eta}f(z)\,(z-g)^{-1}\,dz,
\end{equation}
where $\eta$ is a cycle in $\mathcal{O}\setminus\mathrm{spec}(g)$ winding once around each eigenvalue
of $g$ and not at all around any point of $\mathbb{C}\setminus\mathcal{O}$. Definition
\eqref{4.10:eq-principal-logarithm-1} is the case $f=\log$,
$\mathcal{O}=\mathbb{C}\setminus(-\infty,0]$. Since $\mathcal{O}$ is then simply connected, every cycle in
it has winding number zero around every point of $(-\infty,0]$, so a cycle as in (i) satisfies the
requirements just listed.

\emph{(i).} The eigenvalues of a matrix are the roots of its characteristic polynomial, whose
coefficients are polynomials in the entries; hence the eigenvalues depend continuously on the
matrix. Let $g\in\mathcal{L}$ and let $\eta$ be as in (i). The image of $\eta$ and the half-line
$(-\infty,0]$ are closed sets disjoint from the finite set $\mathrm{spec}(g)$. Therefore there is
a neighbourhood $\mathcal{N}$ of $g$ such that for $g'\in\mathcal{N}$ every eigenvalue of $g'$ lies
close to an eigenvalue of $g$, off the image of $\eta$ and off $(-\infty,0]$. The winding number
of $\eta$ is locally constant on the complement of its image, so $\eta$ winds once around each
eigenvalue of $g'$: the spectrum of $g'$ stays inside $\eta$ and away from $(-\infty,0]$. Thus
$\mathcal{N}\subset\mathcal{L}$, so $\mathcal{L}$ is open, and the same $\eta$ serves in
\eqref{4.10:eq-principal-logarithm-1} for every $g'\in\mathcal{N}$.

The value of \eqref{4.10:eq-principal-logarithm-1} does not depend on the admissible cycle. If
$\eta$ and $\eta'$ are two admissible cycles for $g$, the function
$z\mapsto(\log z)(z-g)^{-1}$ is holomorphic on $\mathbb{C}\setminus((-\infty,0]\cup\mathrm{spec}(g))$.
The cycle $\eta-\eta'$ lies in this open set and winds zero times around each eigenvalue of $g$ and
around each point of $(-\infty,0]$. By Cauchy's theorem the integrals over $\eta$ and $\eta'$
coincide.

For $g'\in\mathcal{N}$ and $z$ on the image of $\eta$, Cramer's rule shows that each entry of
$(z-g')^{-1}$ is a rational function of the entries of $g'$ whose denominator
$\det(z-g')$ does not vanish. The integrand is therefore holomorphic in $g'$ and jointly continuous
in $(z,g')$ on the compact image of $\eta$ times a compact neighbourhood of $g$ in $\mathcal{N}$.
Differentiation under the integral sign then shows that $g'\mapsto\log g'$ is holomorphic on
$\mathcal{N}$, and hence on $\mathcal{L}$.

\emph{(ii) and (iii).} The calculus \eqref{4.10:eq-principal-logarithm-2} has two further
properties. First, $(h\circ f)(g)=h(f(g))$ whenever $f$ is holomorphic near $\mathrm{spec}(g)$ and
$h$ is holomorphic near $f(\mathrm{spec}(g))=\mathrm{spec}(f(g))$. Second, $p(g)$ has its
elementary meaning when $p$ is given by a power series converging on a neighbourhood of
$\mathrm{spec}(g)$; in particular the identity function gives $g$ and $\exp$ gives the exponential
series $e^{g}$.

For (ii), let $g\in\mathcal{L}$ and take $f=\log$ and $h=\exp$. Since $\exp$ is entire, the
composition rule applies, and since $\exp(\log z)=z$ on $\mathbb{C}\setminus(-\infty,0]$, which is
a neighbourhood of $\mathrm{spec}(g)$, we obtain
\begin{equation*}
\exp(\log g)=(\exp\circ\log)(g)=\mathrm{id}(g)=g.
\end{equation*}

For (iii), let $Y$ be as stated. The exponential maps the strip $\{|\operatorname{Im}z|<\pi\}$ into
$\mathbb{C}\setminus(-\infty,0]$. Indeed, write $z=x+iy$ with $|y|<\pi$. Then $e^{z}=e^{x}e^{iy}$
is non-zero, and it is a negative real number only if $e^{iy}=-1$, that is, only if
$y\in\pi+2\pi\mathbb{Z}$, which is excluded. Hence, by the identity
$\mathrm{spec}(f(g))=f(\mathrm{spec}(g))$ recorded above, applied with $f=\exp$ and $g=Y$
(it is proved under (iv) below),
\begin{equation*}
\mathrm{spec}(e^{Y})=e^{\mathrm{spec}(Y)}\subset\mathbb{C}\setminus(-\infty,0],
\end{equation*}
so $e^{Y}\in\mathcal{L}$. On the strip one has $\log(e^{z})=z$, because
$\log(e^{x}e^{iy})=x+iy$ when $y$ lies in $(-\pi,\pi)$, the range of the principal argument. The
strip is an open neighbourhood of $\mathrm{spec}(Y)$. We apply the composition rule with $f=\exp$,
which is entire, and $h=\log$, which is holomorphic near $\exp(\mathrm{spec}(Y))$. This gives
\begin{equation*}
\log(e^{Y})=(\log\circ\exp)(Y)=\mathrm{id}(Y)=Y.
\end{equation*}

\emph{(iv).} The eigenvalues of $f(g)$, counted with algebraic multiplicity, are the numbers
$f(\lambda)$, with $\lambda$ running through the eigenvalues of $g$ counted in the same way. To see
this, triangularise $g$: write $g=STS^{-1}$ with $S$ invertible and $T$ upper triangular, with
diagonal entries $\lambda_1,\dots,\lambda_N$, the eigenvalues of $g$ with multiplicity. Then
$(z-g)^{-1}=S(z-T)^{-1}S^{-1}$, so $f(g)=Sf(T)S^{-1}$. The matrix $(z-T)^{-1}$ is upper triangular
with diagonal entries $(z-\lambda_a)^{-1}$. Hence $f(T)$ is upper triangular, and by Cauchy's
integral formula its diagonal entries are
\begin{equation*}
\frac{1}{2\pi i}\oint_{\eta}f(z)\,(z-\lambda_a)^{-1}\,dz=f(\lambda_a),
\end{equation*}
because $\eta$ winds once around $\lambda_a$. Taking $f=\log$ gives
$\operatorname{Tr}\log g=\sum_{a=1}^{N}\log\lambda_a$, and therefore
\begin{equation*}
e^{\operatorname{Tr}\log g}=\prod_{a=1}^{N}e^{\log\lambda_a}=\prod_{a=1}^{N}\lambda_a=\det g.
\end{equation*}
If $\det g=1$, then $e^{\operatorname{Tr}\log g}=1$, and so $\operatorname{Tr}\log g\in 2\pi i\mathbb{Z}$.

\emph{(v).} An eigenvalue of $iB$ has the form $i\lambda$ with $\lambda\in\mathrm{spec}(B)$. By
Definition~\ref{4.10:def-holomorphic-maps}, $|\lambda|\le|B|<\pi$. Consequently
$\operatorname{Im}(i\lambda)=\operatorname{Re}\lambda$ satisfies
$|\operatorname{Re}\lambda|\le|\lambda|<\pi$, that is, $\operatorname{Im}(i\lambda)\in(-\pi,\pi)$.

\emph{(vi).} A unitary $g$ is normal, so $g=S\,\mathrm{diag}(\mu_a)\,S^{*}$ with $S$ unitary and
$|\mu_a|=1$; moreover $\mu_a\notin(-\infty,0]$ because $g\in\mathcal{L}$. Since
$(z-g)^{-1}=S\,(z-\mathrm{diag}(\mu_a))^{-1}\,S^{*}$ and $(z-\mathrm{diag}(\mu_a))^{-1}$ is
diagonal with entries $(z-\mu_a)^{-1}$, Cauchy's integral formula in
\eqref{4.10:eq-principal-logarithm-1} gives $\log g=S\,\mathrm{diag}(\log\mu_a)\,S^{*}$.

The principal branch satisfies $\overline{\log z}=\log\bar z$ off the cut, since
$\log z=\log|z|+i\operatorname{Arg}z$ with $\operatorname{Arg}z\in(-\pi,\pi)$ and
$\operatorname{Arg}\bar z=-\operatorname{Arg}z$. It also satisfies $\log(1/z)=-\log z$ off the
cut, since $\operatorname{Arg}(1/z)=-\operatorname{Arg}z\in(-\pi,\pi)$. Using these two facts and
$\bar\mu_a=\mu_a^{-1}$, we obtain
\begin{align*}
(\log g)^{*}&=S\,\mathrm{diag}\bigl(\overline{\log\mu_a}\bigr)\,S^{*}
=S\,\mathrm{diag}(\log\bar\mu_a)\,S^{*}\\
&=S\,\mathrm{diag}(-\log\mu_a)\,S^{*}=-\log g .
\end{align*}
Hence $\log g$ is anti-Hermitian.
\end{proof}

\begin{lemma}[Local holomorphic parametrisation of data by potentials]
\label{4.10:lem-data-parametrisation}
Let $S$ be a finite set of bonds, let $V_0\in\mathrm{Conf}(S)$, and let
$B_*\in(\mathfrak{g}^c)^S$ satisfy $|B_*(c)|<\pi$ for all $c\in S$. Put
$V_*:=e^{iB_*}V_0$. Then there exist a connected open neighbourhood $\mathcal{W}$ of
$V_*$ in $\mathrm{Conf}(S)$ and a holomorphic map
$\beta:\mathcal{W}\to(\mathfrak{g}^c)^S$ such that $\beta(V_*)=B_*$ and
\begin{equation}\label{4.10:eq-data-parametrisation-1}
e^{i\beta(V)}V_0=V\qquad\text{for every }V\in\mathcal{W}.
\end{equation}
\end{lemma}

\begin{proof}
Throughout, products of configurations with exponentials of potentials are taken bond
by bond, as in Definition~\ref{4.10:def-holomorphic-maps}, and $\mathcal{L}$ and $\log$
are the set of matrices without spectrum on $(-\infty,0]$ and the principal logarithm of
Lemma~\ref{4.10:lem-principal-logarithm}.

Define $\kappa:\mathrm{Conf}(S)\to M_N(\mathbb{C})^S$ by
$\kappa(V)(c):=V(c)V_0(c)^{-1}$ for $c\in S$. Each component is the product of the
variable matrix $V(c)$ with the fixed matrix $V_0(c)^{-1}$, so $\kappa$ is holomorphic by
the holomorphy fact (H2) listed in \eqref{4.10:eq-holomorphic-maps-a}. In particular
$\kappa$ is continuous. By Lemma~\ref{4.10:lem-principal-logarithm}(i) the set
$\mathcal{L}$ is open in $M_N(\mathbb{C})$, hence $\mathcal{L}^S$ is open in
$M_N(\mathbb{C})^S$. It follows that the set
$\mathcal{W}_1:=\kappa^{-1}(\mathcal{L}^S)$ is open in $\mathrm{Conf}(S)$.

Unwinding the definitions, for every $c\in S$,
\begin{equation*}
\kappa(V_*)(c)=e^{iB_*(c)}V_0(c)V_0(c)^{-1}=e^{iB_*(c)}.
\end{equation*}
The matrix $iB_*(c)$ has operator norm $|B_*(c)|<\pi$. Hence, by clauses (v) and (iii) of
Lemma~\ref{4.10:lem-principal-logarithm}, $e^{iB_*(c)}\in\mathcal{L}$. Thus
$\kappa(V_*)\in\mathcal{L}^S$, that is, $V_*\in\mathcal{W}_1$.

Let $\mathcal{W}$ be the connected component of $\mathcal{W}_1$ that contains $V_*$. The
space $\mathrm{Conf}(S)=(G^c)^S$ is a manifold and is therefore locally connected. Hence
every connected component of its open subset $\mathcal{W}_1$ is open, and $\mathcal{W}$
is a connected open neighbourhood of $V_*$ in $\mathrm{Conf}(S)$.

For $V\in\mathcal{W}$ and $c\in S$ put
\begin{equation}\label{4.10:eq-data-parametrisation-2}
\beta(V)(c):=\frac1i\,\log\kappa(V)(c).
\end{equation}
This is well defined because $\kappa(V)(c)\in\mathcal{L}$ for $V\in\mathcal{W}$. The
restriction of $\kappa$ to $\mathcal{W}$ is holomorphic, with values in $\mathcal{L}^S$.
By Lemma~\ref{4.10:lem-principal-logarithm}(i), $\log$ is holomorphic on $\mathcal{L}$.
Composing these maps and multiplying by the constant $1/i$, we obtain from the holomorphy
facts (H2) and (H4) listed in \eqref{4.10:eq-holomorphic-maps-a} that $\beta$ is
holomorphic on $\mathcal{W}$ as a map into $M_N(\mathbb{C})^S$.

By Lemma~\ref{4.10:lem-principal-logarithm}(ii), for $V\in\mathcal{W}$ and $c\in S$,
\begin{equation*}
e^{i\beta(V)(c)}V_0(c)=\exp\bigl(\log\kappa(V)(c)\bigr)V_0(c)
=V(c)V_0(c)^{-1}V_0(c)=V(c).
\end{equation*}
This is \eqref{4.10:eq-data-parametrisation-1}. Next, by clauses (iii) and (v) of
Lemma~\ref{4.10:lem-principal-logarithm}, applied to $iB_*(c)$, whose norm is less than
$\pi$,
\begin{equation*}
\beta(V_*)(c)=\frac1i\log e^{iB_*(c)}=\frac1i\,iB_*(c)=B_*(c).
\end{equation*}
Hence $\beta(V_*)=B_*$.

It remains to show that $\beta(V)(c)$ is traceless for every $V\in\mathcal{W}$ and
$c\in S$. Then $\beta$ takes values in $(\mathfrak{g}^c)^S$, and it is holomorphic as a
map into this linear subspace of $M_N(\mathbb{C})^S$. Fix $c\in S$ and define
$\varphi_c:\mathcal{W}\to\mathbb{C}$ by
$\varphi_c(V):=\operatorname{Tr}\log\kappa(V)(c)$. The function $\varphi_c$ is continuous
on $\mathcal{W}$, since it is the composition of the continuous maps $\kappa$, $\log$ and
$\operatorname{Tr}$. Since $V(c)$ and $V_0(c)$ lie in $G^c=SL(N,\mathbb{C})$, we have
\begin{equation*}
\det\kappa(V)(c)=\det V(c)\,\bigl(\det V_0(c)\bigr)^{-1}=1.
\end{equation*}
Therefore Lemma~\ref{4.10:lem-principal-logarithm}(iv) gives
$\varphi_c(V)\in2\pi i\mathbb{Z}$ for every $V\in\mathcal{W}$. On the other hand, by the
computation of $\beta(V_*)$ above,
\begin{equation*}
\varphi_c(V_*)=\operatorname{Tr}\bigl(iB_*(c)\bigr)=0,
\end{equation*}
because $B_*(c)\in\mathfrak{g}^c$ is traceless. A continuous function on a connected set
that takes values in the discrete set $2\pi i\mathbb{Z}$ is constant. Since $\mathcal{W}$
is connected, $\varphi_c\equiv0$ on $\mathcal{W}$. Consequently, for every
$V\in\mathcal{W}$,
\begin{equation*}
\operatorname{Tr}\beta(V)(c)=\frac1i\varphi_c(V)=0,
\end{equation*}
that is, $\beta(V)(c)\in\mathfrak{g}^c$. As $c\in S$ was arbitrary, the proof is complete.
\end{proof}

\begin{proposition}[The continuation is holomorphic on an open domain]
\label{4.10:prop-continuation-open}
Let $\mathcal{T}=\mathcal{T}_{\varrho}(Z,n)$ be a truncated tower of depth $n$ under the rule
$\varrho\in\{\mathrm{I},\mathrm{II}\}$. Let $\mathcal{R}(\mathcal{T})$ be the set of admitted
$G$-valued data and $\mathbb{D}(\mathcal{T})=\mathrm{Conf}(\mathfrak{B}(\mathcal{T}))$ the data space.
Assume the following parts of properties (P9.1) and (P9.2) of Ba{\l}aban's analytic continuation
of the constrained minimiser (Definition~\ref{4.10:def-minimiser-continuation},
display~\eqref{4.10:eq-minimiser-continuation-a}); item~\textup{(a)} is taken from (P9.1) and
item~\textup{(b)} from (P9.2).
\begin{itemize}
\item[\textup{(a)}] For every $V_0\in\mathcal{R}(\mathcal{T})$ the set $\mathcal{B}(V_0)$ of
potentials satisfying \cite[(172)]{B-CMP102b} is an open subset of
$(\mathfrak{g}^c)^{\mathfrak{B}(\mathcal{T})}$ with $0\in\mathcal{B}(V_0)$. Every
$B\in\mathcal{B}(V_0)$ satisfies $|B(c)|<\pi$
for every $c\in\mathfrak{B}(\mathcal{T})$. The map
$\Phi_{V_0}:\mathcal{B}(V_0)\to\mathrm{Conf}(\mathrm{bd}(Z))$ is holomorphic.
\item[\textup{(b)}] The set $\mathcal{S}(\mathcal{T})$ consists of the data $e^{iB}V_0$ with
$V_0\in\mathcal{R}(\mathcal{T})$ and $B\in\mathcal{B}(V_0)$. The map $\mathcal{U}_{\mathcal{T}}$ is
the function on $\mathcal{S}(\mathcal{T})$ whose value at $V=e^{iB}V_0$ is $\Phi_{V_0}(B)$, and this
value may be computed in any such representation of $V$.
\end{itemize}
Then $\mathcal{S}(\mathcal{T})$ is an open subset of $\mathbb{D}(\mathcal{T})$, and
\begin{equation*}
\mathcal{U}_{\mathcal{T}}:\mathcal{S}(\mathcal{T})\to\mathrm{Conf}(\mathrm{bd}(Z))
\end{equation*}
is holomorphic.
\end{proposition}

\begin{proof}
Let $V_*\in\mathcal{S}(\mathcal{T})$. By the description of $\mathcal{S}(\mathcal{T})$ in
hypothesis~\textup{(b)}, there are $V_0\in\mathcal{R}(\mathcal{T})$ and $B_*\in\mathcal{B}(V_0)$ with
$V_*=e^{iB_*}V_0$. By the bound on $\mathcal{B}(V_0)$ in hypothesis~\textup{(a)},
$|B_*(c)|<\pi$ for all $c\in\mathfrak{B}(\mathcal{T})$.

We apply the local holomorphic parametrisation of data by potentials
(Lemma~\ref{4.10:lem-data-parametrisation}) with bond set $S=\mathfrak{B}(\mathcal{T})$, with the
$G$-valued configuration $V_0\in\mathcal{R}(\mathcal{T})\subset\mathrm{Conf}(\mathfrak{B}(\mathcal{T}))$
and with the potential $B_*$. The hypotheses of that lemma hold: $S$ is a finite set of bonds, and
$|B_*(c)|<\pi$ for all $c\in S$. Since $\mathrm{Conf}(S)=\mathbb{D}(\mathcal{T})$, the lemma gives a
connected open neighbourhood $\mathcal{W}$ of $V_*$ in $\mathbb{D}(\mathcal{T})$ and a holomorphic map
$\beta:\mathcal{W}\to(\mathfrak{g}^c)^{\mathfrak{B}(\mathcal{T})}$ such that
\begin{equation*}
\beta(V_*)=B_*,\qquad e^{i\beta(V)}V_0=V\quad\text{for every }V\in\mathcal{W}.
\end{equation*}

By hypothesis~\textup{(a)}, $\mathcal{B}(V_0)$ is open. The map $\beta$ is holomorphic, hence
continuous. Therefore the set
\begin{equation*}
\mathcal{W}':=\mathcal{W}\cap\beta^{-1}\bigl(\mathcal{B}(V_0)\bigr)
\end{equation*}
is open in $\mathbb{D}(\mathcal{T})$. It contains $V_*$, because $\beta(V_*)=B_*\in\mathcal{B}(V_0)$.

Let $V\in\mathcal{W}'$. Then $V=e^{i\beta(V)}V_0$, with $V_0\in\mathcal{R}(\mathcal{T})$ and
$\beta(V)\in\mathcal{B}(V_0)$. By the description of $\mathcal{S}(\mathcal{T})$ in
hypothesis~\textup{(b)}, this gives $V\in\mathcal{S}(\mathcal{T})$. Hence
$\mathcal{W}'\subset\mathcal{S}(\mathcal{T})$. Since $V_*$ was an arbitrary point,
$\mathcal{S}(\mathcal{T})$ is a neighbourhood of each of its points, that is, it is open.

Moreover, by hypothesis~\textup{(b)} the value of $\mathcal{U}_{\mathcal{T}}$ at $V$ may be computed
in the representation $V=e^{i\beta(V)}V_0$. This gives
\begin{equation}\label{4.10:eq-continuation-open-1}
\mathcal{U}_{\mathcal{T}}(V)=\Phi_{V_0}\bigl(\beta(V)\bigr)\qquad\text{for every }V\in\mathcal{W}'.
\end{equation}
By \eqref{4.10:eq-continuation-open-1}, on $\mathcal{W}'$ the map $\mathcal{U}_{\mathcal{T}}$ is the
composition of two holomorphic maps:
\begin{itemize}
\item the map $\beta|_{\mathcal{W}'}$, which takes values in $\mathcal{B}(V_0)$;
\item the map $\Phi_{V_0}$, which is holomorphic on $\mathcal{B}(V_0)$ by hypothesis~\textup{(a)}.
\end{itemize}
A composition of holomorphic maps is holomorphic; this is the composition rule among the facts on
holomorphic maps collected in~\eqref{4.10:eq-holomorphic-maps-a}. Hence $\mathcal{U}_{\mathcal{T}}$
is holomorphic on the open neighbourhood $\mathcal{W}'$ of $V_*$.

Every point of the open set $\mathcal{S}(\mathcal{T})$ therefore has an open neighbourhood on which
$\mathcal{U}_{\mathcal{T}}$ is holomorphic. Holomorphy is a local property; this is a further fact
collected in~\eqref{4.10:eq-holomorphic-maps-a}. Therefore $\mathcal{U}_{\mathcal{T}}$ is holomorphic
on $\mathcal{S}(\mathcal{T})$.
\end{proof}

No analyticity in the $G$-valued factor $V_0$ is needed in this argument. At fixed $V_0$, the
potentials $B$ near $B_*$ already sweep out a full neighbourhood of $e^{iB_*}V_0$ in the data space,
because $B\mapsto e^{iB}V_0$ is locally biholomorphic where $|B|<\pi$.

\begin{theorem}[Openness of the existence set and holomorphy]\label{4.10:thm-existence-set}
Let $Z$ be a region, let $n\ge 1$, let $\varrho\in\{\mathrm{I},\mathrm{II}\}$, and let
$\mathcal{T}=\mathcal{T}_{\varrho}(Z,n)$ be the truncated tower of depth $n$ for the rule
$\varrho$. Assume the following:
\begin{itemize}
\item (P9.1) and (P9.2): the statements so tagged in Ba{\l}aban's analytic continuation of
the constrained minimiser, \eqref{4.10:eq-minimiser-continuation-a}, used here as stated in
\cite[Proposition~9]{B-CMP102b};
\item the averaging hypothesis: the block averages are holomorphic on their open domains of
definition; this is the hypothesis from which Lemma~\ref{4.10:lem-averaged-data} (holomorphy
of the averaged data along a tower) is proved;
\item the current hypothesis: the current of \cite[(1.8)]{B-CMP109} is a holomorphic map of
the configuration, as stated in Definition~\ref{4.10:def-current-holomorphic}.
\end{itemize}
Then the following hold.
\begin{enumerate}
\item[(i)] The existence set $\mathcal{V}(\mathcal{T})$ of
Definition~\ref{4.10:def-tower-quantities} is an open subset of
$\mathrm{Conf}(\mathrm{bd}(Z))$, the complex manifold of all $G^c$-valued configurations on
the bonds of $Z$. No condition is imposed on the elements of $\mathcal{V}(\mathcal{T})$ other
than membership in it.
\item[(ii)] The map $W\mapsto U_n(Z;\dot{\mathcal{M}}W)$ is holomorphic from
$\mathcal{V}(\mathcal{T})$ to $\mathrm{Conf}(\mathrm{bd}(Z))$.
\item[(iii)] The map $W\mapsto J_n(Z;\dot{\mathcal{M}}W)$ is holomorphic from
$\mathcal{V}(\mathcal{T})$ to $(\mathfrak{g}^c)^{\mathrm{bd}_J(Z)}$.
\item[(iv)] For every plaquette $P$ of $Z$ the map
$W\mapsto\partial U_n(Z;\dot{\mathcal{M}}W)(P)\in M_N(\mathbb{C})$ is holomorphic on
$\mathcal{V}(\mathcal{T})$, and for every bond $b\in\mathrm{bd}_J(Z)$ the map
$W\mapsto J_n(Z;\dot{\mathcal{M}}W)(b)\in\mathfrak{g}^c$ is holomorphic on
$\mathcal{V}(\mathcal{T})$. In particular the real-valued functions
\begin{equation*}
W\mapsto\bigl|\partial U_n(Z;\dot{\mathcal{M}}W)(P)-1\bigr|,
\qquad
W\mapsto\bigl|J_n(Z;\dot{\mathcal{M}}W)(b)\bigr|
\end{equation*}
are continuous on $\mathcal{V}(\mathcal{T})$.
\end{enumerate}
\end{theorem}

\begin{proof}
We recall from Definition~\ref{4.10:def-tower-quantities} that the existence set is
\begin{equation}\label{4.10:eq-existence-set-1}
\mathcal{V}(\mathcal{T})
=\bigl\{W\in\mathcal{A}(\mathcal{T}):\ \dot{\mathcal{M}}_{\mathcal{T}}W\in
\mathcal{S}(\mathcal{T})\bigr\}
=\dot{\mathcal{M}}_{\mathcal{T}}^{-1}\bigl(\mathcal{S}(\mathcal{T})\bigr),
\end{equation}
and that for $W\in\mathcal{V}(\mathcal{T})$ the tower field is
\begin{equation}\label{4.10:eq-existence-set-2}
U_n(Z;\dot{\mathcal{M}}W)=\mathcal{U}_{\mathcal{T}}\bigl(\dot{\mathcal{M}}_{\mathcal{T}}W\bigr)
\in\mathrm{Conf}(\mathrm{bd}(Z)).
\end{equation}
We shall use the holomorphy facts \eqref{4.10:eq-holomorphic-maps-a}, in particular those on
polynomial maps and on compositions.

(i) By Lemma~\ref{4.10:lem-averaged-data}, which uses the averaging hypothesis, the set
$\mathcal{A}(\mathcal{T})$ is open in $\mathrm{Conf}(\mathrm{bd}(Z))$ and the averaged data
map $\dot{\mathcal{M}}_{\mathcal{T}}:\mathcal{A}(\mathcal{T})\to\mathbb{D}(\mathcal{T})$ is
holomorphic; in particular it is continuous. By Proposition~\ref{4.10:prop-continuation-open},
whose hypotheses are (P9.1) and (P9.2), the set $\mathcal{S}(\mathcal{T})$ is open in
$\mathbb{D}(\mathcal{T})$. The preimage of an open set under a continuous map is open; hence,
by \eqref{4.10:eq-existence-set-1}, $\mathcal{V}(\mathcal{T})$ is open in
$\mathcal{A}(\mathcal{T})$. Since $\mathcal{A}(\mathcal{T})$ is itself open in
$\mathrm{Conf}(\mathrm{bd}(Z))$, an open subset of $\mathcal{A}(\mathcal{T})$ is open in
$\mathrm{Conf}(\mathrm{bd}(Z))$, and so $\mathcal{V}(\mathcal{T})$ is open in
$\mathrm{Conf}(\mathrm{bd}(Z))$. The set is defined by \eqref{4.10:eq-existence-set-1} alone,
so no condition beyond membership is imposed on its elements.

(ii) The restriction $\dot{\mathcal{M}}_{\mathcal{T}}|_{\mathcal{V}(\mathcal{T})}$ of the
holomorphic map $\dot{\mathcal{M}}_{\mathcal{T}}$ to the open subset
$\mathcal{V}(\mathcal{T})$ is holomorphic, and by \eqref{4.10:eq-existence-set-1} it takes
its values in $\mathcal{S}(\mathcal{T})$. By Proposition~\ref{4.10:prop-continuation-open}
the map $\mathcal{U}_{\mathcal{T}}:\mathcal{S}(\mathcal{T})\to\mathrm{Conf}(\mathrm{bd}(Z))$
is holomorphic on the open set $\mathcal{S}(\mathcal{T})$. By
\eqref{4.10:eq-existence-set-2}, on $\mathcal{V}(\mathcal{T})$ the map
$W\mapsto U_n(Z;\dot{\mathcal{M}}W)$ is the composition
$\mathcal{U}_{\mathcal{T}}\circ\dot{\mathcal{M}}_{\mathcal{T}}|_{\mathcal{V}(\mathcal{T})}$,
and a composition of holomorphic maps is holomorphic
\eqref{4.10:eq-holomorphic-maps-a}.

(iii) By Definition~\ref{4.10:def-current-holomorphic}, for $W\in\mathcal{V}(\mathcal{T})$,
\begin{equation}\label{4.10:eq-existence-set-3}
J_n(Z;\dot{\mathcal{M}}W)=\mathcal{J}_{L^{-n}}\bigl(U_n(Z;\dot{\mathcal{M}}W)\bigr),
\end{equation}
so that $W\mapsto J_n(Z;\dot{\mathcal{M}}W)$ is the composition of the map of (ii) with
$\mathcal{J}_{L^{-n}}$. By the current hypothesis, in the form stated in
Definition~\ref{4.10:def-current-holomorphic} with spacing parameter $\xi=L^{-n}$ and region
$\Omega=Z$, the map
$\mathcal{J}_{L^{-n}}:\mathrm{Conf}(\mathrm{bd}(Z))\to(\mathfrak{g}^c)^{\mathrm{bd}_J(Z)}$
is holomorphic. The map of (ii) is holomorphic by (ii). Their composition is therefore
holomorphic \eqref{4.10:eq-holomorphic-maps-a}.

(iv) Let $P$ be a plaquette of $Z$. The holonomy map $U\mapsto\partial U(P)$ from
$\mathrm{Conf}(\mathrm{bd}(Z))$ to $M_N(\mathbb{C})$ is a product of the bond variables on the
boundary of $P$ and of their inverses; since an element of $G^c$ has determinant one, its
inverse is its adjugate matrix, and so the holonomy map is a polynomial in the matrix entries
of the bond variables. It is therefore holomorphic \eqref{4.10:eq-holomorphic-maps-a}, and its
composition with the holomorphic map of (ii), which is
$W\mapsto\partial U_n(Z;\dot{\mathcal{M}}W)(P)$, is holomorphic on
$\mathcal{V}(\mathcal{T})$. Let $b\in\mathrm{bd}_J(Z)$. The coordinate projection
$(\mathfrak{g}^c)^{\mathrm{bd}_J(Z)}\to\mathfrak{g}^c$ onto the component at $b$ is
complex-linear, hence holomorphic, and its composition with the holomorphic map of (iii),
which is $W\mapsto J_n(Z;\dot{\mathcal{M}}W)(b)$, is holomorphic on
$\mathcal{V}(\mathcal{T})$. Finally, holomorphic maps are continuous, and the functions
$Y\mapsto|Y-1|$ on $M_N(\mathbb{C})$ and $Y\mapsto|Y|$ on $\mathfrak{g}^c$ are continuous,
since a norm is continuous. The two real-valued functions of the statement are compositions
of these continuous functions with the two holomorphic maps just obtained, and are therefore
continuous on $\mathcal{V}(\mathcal{T})$.
\end{proof}

\begin{corollary}[Analyticity at complex one-power representatives]
\label{4.10:cor-representative-analyticity}
Assume the following statements:
\begin{itemize}
\item the first two parts of Ba{\l}aban's analytic continuation of the constrained minimiser,
\eqref{4.10:eq-minimiser-continuation-a}, read from \cite[Proposition~9]{B-CMP102b};
\item the statement on block averages at complex configurations and the statement on the
current at complex configurations, which together with the first item are the inputs of
Theorem~\ref{4.10:thm-existence-set}.
\end{itemize}
Let $j\ge1$, $X\in\mathcal{D}_{j-1}$ and $a_0,a_1,\gamma_0>0$. Let
$(U_*,A'_*,\mathbf{J}_*)\in\mathfrak{P}(X)$ be a decomposed one-power representative at
$(j,X;a_0,a_1,\gamma_0)$ in the sense of Definition~\ref{4.10:def-decomposed-representatives},
and put $\mathbf{U}_*:=e^{i\xi_jA'_*}U_*$. Let $(p,Z)\in\mathfrak{T}_j(X)$, $1\le n\le p$ and
$\varrho\in\{\mathrm{I},\mathrm{II}\}$, and let $\mathcal{T}=\mathcal{T}_{\varrho}(Z,n)$. Then the
following hold.
\begin{enumerate}
\item[(a)] $\mathbf{U}_*|_{\mathrm{bd}(Z)}\in\mathcal{V}(\mathcal{T})$ and
$U_*|_{\mathrm{bd}(Z)}\in\mathcal{V}(\mathcal{T})$. Consequently the open set
$\mathcal{V}:=\mathcal{V}(\mathcal{T})\subset\mathrm{Conf}(\mathrm{bd}(Z))$ is a neighbourhood both
of $\mathbf{U}_*|_{\mathrm{bd}(Z)}$ and of $U_*|_{\mathrm{bd}(Z)}$; no regularity condition is
imposed on its members; the quantities $U_n(Z;\dot{\mathcal{M}}W)$ and
$J_n(Z;\dot{\mathcal{M}}W)$ exist for every $W\in\mathcal{V}$; and the maps
\begin{equation*}
W\mapsto U_n(Z;\dot{\mathcal{M}}W)\in\mathrm{Conf}(\mathrm{bd}(Z)),\qquad
W\mapsto J_n(Z;\dot{\mathcal{M}}W)\in(\mathfrak{g}^c)^{\mathrm{bd}_J(Z)}
\end{equation*}
are holomorphic on $\mathcal{V}$.
\item[(b)] (The $\mathbf{U}$-half as a function of the decomposed pair.) On the open
neighbourhood $\sigma_Z^{-1}(\mathcal{V}(\mathcal{T}))$ of $(U_*,A'_*,\mathbf{J}_*)$ in
$\mathfrak{P}(X)$ the maps
\begin{equation*}
(U,A',\mathbf{J})\mapsto U_n\bigl(Z;\dot{\mathcal{M}}\sigma_Z(U,A',\mathbf{J})\bigr),\qquad
(U,A',\mathbf{J})\mapsto J_n\bigl(Z;\dot{\mathcal{M}}\sigma_Z(U,A',\mathbf{J})\bigr)
\end{equation*}
are defined and continuous, and for fixed $U$ they are holomorphic in $A'$.
\item[(b$'$)] (The $U$-half as a function of the $G$-valued factor.) On the open
neighbourhood $\sigma_Z'^{-1}(\mathcal{V}(\mathcal{T}))$ of $(U_*,A'_*,\mathbf{J}_*)$ in
$\mathfrak{P}(X)$ the maps
\begin{equation*}
(U,A',\mathbf{J})\mapsto U_n\bigl(Z;\dot{\mathcal{M}}(U|_{\mathrm{bd}(Z)})\bigr),\qquad
(U,A',\mathbf{J})\mapsto J_n\bigl(Z;\dot{\mathcal{M}}(U|_{\mathrm{bd}(Z)})\bigr)
\end{equation*}
are defined and continuous; they depend on $U$ alone, and as functions of
$U\in G^{\mathrm{bd}(X)}$ they are restrictions to the subset $G^{\mathrm{bd}(X)}$ of maps
holomorphic on an open subset of $\mathrm{Conf}(\mathrm{bd}(X))$.
\item[(c)] (Pairs.) On the open neighbourhood
\begin{equation*}
\bigl\{(\mathbf{U},\mathbf{J})\in\mathfrak{Q}(X):
\mathbf{U}|_{\mathrm{bd}(Z)}\in\mathcal{V}(\mathcal{T})\bigr\}
\end{equation*}
of $(\mathbf{U}_*,\mathbf{J}_*)$ in $\mathfrak{Q}(X)$ the maps
\begin{equation*}
(\mathbf{U},\mathbf{J})\mapsto U_n\bigl(Z;\dot{\mathcal{M}}(\mathbf{U}|_{\mathrm{bd}(Z)})\bigr),
\qquad
(\mathbf{U},\mathbf{J})\mapsto J_n\bigl(Z;\dot{\mathcal{M}}(\mathbf{U}|_{\mathrm{bd}(Z)})\bigr)
\end{equation*}
are holomorphic; they do not depend on $\mathbf{J}$.
\end{enumerate}
\end{corollary}

\begin{proof}
Under the hypotheses of the corollary, Theorem~\ref{4.10:thm-existence-set} is available for
the region $Z$, the depth $n$ and the rule $\varrho$, that is, for the truncated tower
$\mathcal{T}=\mathcal{T}_{\varrho}(Z,n)$. Of the representative $(U_*,A'_*,\mathbf{J}_*)$ only
the defining clauses of Definition~\ref{4.10:def-decomposed-representatives} are used.

\emph{Part (a).} Since $(U_*,A'_*,\mathbf{J}_*)$ is a decomposed one-power representative at
$(j,X;a_0,a_1,\gamma_0)$ and $(p,Z)\in\mathfrak{T}_j(X)$,
Definition~\ref{4.10:def-decomposed-representatives} gives the clause
$(\mathrm{iv})_{p,Z}[a_0L^{-(j-p)}]$ both for $\mathbf{W}=\mathbf{U}_*$ (the $\mathbf{U}$-half) and
for $\mathbf{W}=U_*$ (the $U$-half). That clause is required for both rules and for every depth
from $1$ to $p$; in particular it holds for the rule $\varrho$ and the depth $n$, since
$1\le n\le p$. Its first requirement is the existence of the tower field
$U_n(Z;\dot{\mathcal{M}}(\mathbf{W}|_{\mathrm{bd}(Z)}))$ for the rule $\varrho$. By
Definition~\ref{4.10:def-tower-quantities}, this field exists exactly when
$\mathbf{W}|_{\mathrm{bd}(Z)}$ belongs to the existence set $\mathcal{V}(\mathcal{T})$. Hence
$\mathbf{U}_*|_{\mathrm{bd}(Z)}\in\mathcal{V}(\mathcal{T})$ and
$U_*|_{\mathrm{bd}(Z)}\in\mathcal{V}(\mathcal{T})$. The value of the constant
$a_0L^{-(j-p)}$ in the clause plays no role here; only the existence requirement is used.

The remaining assertions of (a) are clauses (i), (ii) and (iii) of
Theorem~\ref{4.10:thm-existence-set}, applied with $\mathcal{V}=\mathcal{V}(\mathcal{T})$. By
clause (i), $\mathcal{V}$ is open in $\mathrm{Conf}(\mathrm{bd}(Z))$, and no condition is imposed
on its members other than membership. Since $\mathcal{V}$ is open and contains the two points
just found, it is a neighbourhood of each of them. By clauses (ii) and (iii),
$U_n(Z;\dot{\mathcal{M}}W)$ and $J_n(Z;\dot{\mathcal{M}}W)$ exist for every $W\in\mathcal{V}$, and
the two maps of (a) are holomorphic on $\mathcal{V}$.

\emph{Part (b).} By Definition~\ref{4.10:def-decomposed-representatives}, the map
$\sigma_Z:\mathfrak{P}(X)\to\mathrm{Conf}(\mathrm{bd}(Z))$ is continuous. Since
$\mathcal{V}(\mathcal{T})$ is open by part (a), the preimage
$\sigma_Z^{-1}(\mathcal{V}(\mathcal{T}))$ is open in $\mathfrak{P}(X)$. We have
$\sigma_Z(U_*,A'_*,\mathbf{J}_*)=\mathbf{U}_*|_{\mathrm{bd}(Z)}$, and this lies in
$\mathcal{V}(\mathcal{T})$ by part (a); so the preimage contains $(U_*,A'_*,\mathbf{J}_*)$. On this
preimage the two maps of (b) are the compositions of $\sigma_Z$ with the maps
$W\mapsto U_n(Z;\dot{\mathcal{M}}W)$ and $W\mapsto J_n(Z;\dot{\mathcal{M}}W)$ of
Theorem~\ref{4.10:thm-existence-set}(ii),(iii). These maps are defined on
$\mathcal{V}(\mathcal{T})$ and are holomorphic, hence continuous. The compositions are
therefore defined and continuous.

Fix $U$ and $\mathbf{J}$. The set of $A'$ with $(U,A',\mathbf{J})$ in the preimage is open in
$(\mathfrak{g}^c)^{\mathrm{bd}(X)}$, since it is the preimage of an open set under the continuous
map $A'\mapsto(U,A',\mathbf{J})$. By Definition~\ref{4.10:def-decomposed-representatives}, the
map $A'\mapsto\sigma_Z(U,A',\mathbf{J})$ is holomorphic. A composition of holomorphic maps is
holomorphic (one of the holomorphy facts collected in
\eqref{4.10:eq-holomorphic-maps-a}). Hence, for fixed $U$, both maps of (b) are holomorphic in
$A'$.

\emph{Part (b$'$).} The argument of part (b) applies with $\sigma'_Z$ in place of $\sigma_Z$.
The map $\sigma'_Z$ is continuous by Definition~\ref{4.10:def-decomposed-representatives}, so
$\sigma_Z'^{-1}(\mathcal{V}(\mathcal{T}))$ is open in $\mathfrak{P}(X)$. It contains
$(U_*,A'_*,\mathbf{J}_*)$, because $\sigma'_Z(U_*,A'_*,\mathbf{J}_*)=U_*|_{\mathrm{bd}(Z)}$ lies in
$\mathcal{V}(\mathcal{T})$ by part (a). On it the two maps of (b$'$) are the compositions of
$\sigma'_Z$ with the continuous maps of Theorem~\ref{4.10:thm-existence-set}(ii),(iii). They
are therefore defined and continuous. Since $\sigma'_Z(U,A',\mathbf{J})=U|_{\mathrm{bd}(Z)}$
involves neither $A'$ nor $\mathbf{J}$, they depend on $U$ alone.

As functions of $U$, the maps factor through $U\mapsto U|_{\mathrm{bd}(Z)}$. This map is the
restriction to $G^{\mathrm{bd}(X)}$ of the restriction map
$\mathrm{Conf}(\mathrm{bd}(X))\to\mathrm{Conf}(\mathrm{bd}(Z))$, which is holomorphic by
\eqref{4.10:eq-holomorphic-maps-a}. Let
\begin{equation}\label{4.10:eq-representative-analyticity-1}
\mathcal{O}':=\bigl\{\mathbf{V}\in\mathrm{Conf}(\mathrm{bd}(X)):
\mathbf{V}|_{\mathrm{bd}(Z)}\in\mathcal{V}(\mathcal{T})\bigr\}.
\end{equation}
Then $\mathcal{O}'$ is the preimage of the open set $\mathcal{V}(\mathcal{T})$ under this
holomorphic, hence continuous, restriction map, and so it is open. It contains $U_*$ by
part (a).

On $\mathcal{O}'$ consider the maps
$\mathbf{V}\mapsto U_n(Z;\dot{\mathcal{M}}(\mathbf{V}|_{\mathrm{bd}(Z)}))$ and
$\mathbf{V}\mapsto J_n(Z;\dot{\mathcal{M}}(\mathbf{V}|_{\mathrm{bd}(Z)}))$. They are compositions
of the holomorphic restriction map with the holomorphic maps of
Theorem~\ref{4.10:thm-existence-set}(ii),(iii), and so they are holomorphic on $\mathcal{O}'$.
For $U\in G^{\mathrm{bd}(X)}\cap\mathcal{O}'$ their values are those of the two maps of (b$'$).
These are therefore restrictions to $G^{\mathrm{bd}(X)}$ of maps holomorphic on the open subset
$\mathcal{O}'$ of $\mathrm{Conf}(\mathrm{bd}(X))$.

\emph{Part (c).} The map $(\mathbf{U},\mathbf{J})\mapsto\mathbf{U}|_{\mathrm{bd}(Z)}$ from
$\mathfrak{Q}(X)$ to $\mathrm{Conf}(\mathrm{bd}(Z))$ is holomorphic by
\eqref{4.10:eq-holomorphic-maps-a}. The displayed set in (c) is the preimage of the open set
$\mathcal{V}(\mathcal{T})$ under this map, and so it is open. It contains
$(\mathbf{U}_*,\mathbf{J}_*)$, because $\mathbf{U}_*|_{\mathrm{bd}(Z)}\in\mathcal{V}(\mathcal{T})$
by part (a). On it the two maps of (c) are the compositions of this restriction map with the
holomorphic maps of Theorem~\ref{4.10:thm-existence-set}(ii),(iii). Since a composition of
holomorphic maps is holomorphic (\eqref{4.10:eq-holomorphic-maps-a}), they are holomorphic.
They do not depend on $\mathbf{J}$, because the restriction map does not.
\end{proof}

\begin{remark}[What is assumed about the representative]\label{4.10:rem-representative-assumptions}
Let $j\ge1$, $X\in\mathcal{D}_{j-1}$, $a_0,a_1,\gamma_0>0$, and let $(U_*,A'_*,\mathbf{J}_*)$ be a
decomposed one-power representative at $(j,X;a_0,a_1,\gamma_0)$, with
$\mathbf{U}_*=e^{i\xi_jA'_*}U_*$. For $(p,Z)\in\mathfrak{T}_j(X)$, $1\le n\le p$ and
$\varrho\in\{\mathrm{I},\mathrm{II}\}$ put $\mathcal{T}=\mathcal{T}_{\varrho}(Z,n)$, and let
$\mathbf{W}$ stand for either of the two halves $\mathbf{U}_*$, $U_*$. Then the following hold.
\begin{enumerate}
\item[(i)] Part~(a) of Corollary~\ref{4.10:cor-representative-analyticity} uses only the
presupposition of existence carried by the clause $(\mathrm{iv})_{p,Z}[a_0L^{-(j-p)}]$ of the
definition of a one-power representative, namely that, for each $(p,Z)\in\mathfrak{T}_j(X)$, each
$n$, each rule $\varrho$ and both halves,
\begin{equation}\label{4.10:eq-representative-assumptions-1}
\mathbf{W}|_{\mathrm{bd}(Z)}\in\mathcal{V}(\mathcal{T}).
\end{equation}
It uses nothing else.
\item[(ii)] For the $G$-valued half $U_*$, the inclusion
\begin{equation*}
\dot{\mathcal{M}}_{\mathcal{T}}(U_*|_{\mathrm{bd}(Z)})\in\mathcal{R}(\mathcal{T})\subset\mathcal{S}(\mathcal{T})
\end{equation*}
follows from \cite[Proposition~2]{B-CMP98} and \cite[(7)]{B-CMP102b} when $a_0$ lies below a
threshold depending on $d$ and $L$ only. This fact is not used.
\item[(iii)] For the complex half $\mathbf{U}_*$, no such statement is printed in
Ba{\l}aban's papers, and none is used. There, \eqref{4.10:eq-representative-assumptions-1} is part
of the definition of a one-power representative.
\end{enumerate}
\end{remark}

\begin{proof}
By the definition of a one-power representative, the clause
$(\mathrm{iv})_{p,Z}[a_0L^{-(j-p)}]$ is imposed at the representative for every
$(p,Z)\in\mathfrak{T}_j(X)$. It is imposed for the complex half $\mathbf{U}_*$. The same definition
imposes it with the $G$-valued factor $U_*$ in place of $\mathbf{U}_*$. In both cases it is imposed
for the towers of both rules $\varrho=\mathrm{I},\mathrm{II}$.

That clause is a statement about the quantities $\partial U_n(Z;\dot{\mathcal{M}}\mathbf{W})$ and
$J_n(Z;\dot{\mathcal{M}}\mathbf{W})$. Imposing it therefore presupposes that these quantities exist.
By Definition~\ref{4.10:def-tower-quantities}, existence means two things:
$\mathbf{W}|_{\mathrm{bd}(Z)}$ lies in $\mathcal{A}(\mathcal{T})$, so that the averaged data
$\dot{\mathcal{M}}_{\mathcal{T}}(\mathbf{W}|_{\mathrm{bd}(Z)})$ are defined, and these averaged data
lie in the neighbourhood $\mathcal{S}(\mathcal{T})$ of $G$-valued
data on which Ba{\l}aban's analytic continuation of the constrained minimiser
(Definition~\ref{4.10:def-minimiser-continuation}) is defined. In the notation of
Definition~\ref{4.10:def-tower-quantities}, this is precisely
\eqref{4.10:eq-representative-assumptions-1}. Part~(a) of
Corollary~\ref{4.10:cor-representative-analyticity} asserts
\eqref{4.10:eq-representative-assumptions-1} for $\mathbf{W}=\mathbf{U}_*$ and for
$\mathbf{W}=U_*$. It uses this presupposition of the clause, and nothing else. This proves item~(i).

For item~(ii), let $a_0$ lie below a threshold depending on $d$ and $L$ only. The configuration
$U_*$ is $G$-valued and satisfies $|\partial U_*-1|<a_0\xi_j^2$. By
\cite[Proposition~2]{B-CMP98}, as recalled in \cite{B-CMP102b}, its averages
are $G$-valued and satisfy the regularity condition \cite[(7)]{B-CMP102b}, whose curvature
parameter, there denoted $\varepsilon_1$, can be taken $O(a_0)$. Hence
$\dot{\mathcal{M}}_{\mathcal{T}}(U_*|_{\mathrm{bd}(Z)})\in\mathcal{R}(\mathcal{T})$, and
$\mathcal{R}(\mathcal{T})\subset\mathcal{S}(\mathcal{T})$. This conclusion enters none of the
arguments of this section.

For item~(iii), no corresponding statement for a complex configuration
$\mathbf{U}_*=e^{i\xi_jA'_*}U_*$ is printed in Ba{\l}aban's papers, and none is invoked. At the
complex half, \eqref{4.10:eq-representative-assumptions-1} is one of the defining conditions of a
one-power representative: that definition imposes
$(\mathrm{iv})_{p,Z}[a_0L^{-(j-p)}]$, whose quantities exist only in the sense of
Definition~\ref{4.10:def-tower-quantities}.
\end{proof}

\begin{corollary}[The gauge-orbit extension of the continuation]\label{4.10:cor-orbit-extension}
Let $Z$ be a region, $n\ge 1$, $\varrho\in\{\mathrm{I},\mathrm{II}\}$ and
$\mathcal{T}=\mathcal{T}_{\varrho}(Z,n)$. Assume conditions \textup{(P9.1)}, \textup{(P9.2)} and
\textup{(P9.4)} of Ba{\l}aban's analytic continuation of the constrained minimiser
(Definition~\ref{4.10:def-minimiser-continuation}, display~\eqref{4.10:eq-minimiser-continuation-a};
\cite[(181), Proposition~9]{B-CMP102b}), and assume moreover the hypotheses under which
Theorem~\ref{4.10:thm-existence-set} and Corollary~\ref{4.10:cor-representative-analyticity} are
proved. Let $\Gamma(\mathcal{T})$ and the induced fine transformations $\tilde v$ be as in
\textup{(P9.4)}, and let
\begin{equation*}
\mathcal{O}(\mathcal{T})=\bigcup_{v\in\Gamma(\mathcal{T})}\mathcal{S}(\mathcal{T})^{v}
\end{equation*}
be the orbit union defined there. Then:
\begin{enumerate}
\item $\mathcal{O}(\mathcal{T})$ is open in $\mathbb{D}(\mathcal{T})$;
\item $\mathcal{U}_{\mathcal{T}}$, extended to $\mathcal{O}(\mathcal{T})$ by \textup{(P9.4)}, is
holomorphic on $\mathcal{O}(\mathcal{T})$;
\item Theorem~\ref{4.10:thm-existence-set} and Corollary~\ref{4.10:cor-representative-analyticity}
hold verbatim when $\mathcal{S}(\mathcal{T})$ is replaced by $\mathcal{O}(\mathcal{T})$ in the
definition of the existence set $\mathcal{V}(\mathcal{T})$.
\end{enumerate}
\end{corollary}

\begin{proof}
Fix $v\in\Gamma(\mathcal{T})$. The map $V\mapsto V^{v}$ on $\mathbb{D}(\mathcal{T})$, given bondwise by
$(V^{v})(c)=v(c_-)V(c)v(c_+)^{-1}$, and the map $U\mapsto U^{\tilde v}$ on
$\mathrm{Conf}(\mathrm{bd}(Z))$, given bondwise by
$(U^{\tilde v})(b)=\tilde v(b_-)U(b)\tilde v(b_+)^{-1}$, are polynomial in the entries: for fixed $v$
the matrices $v(c_\pm)$, $\tilde v(b_\pm)$ and their inverses are constant, so each entry of $V^{v}$,
respectively $U^{\tilde v}$, is a polynomial in the entries of $V$, respectively $U$. By the holomorphy
facts~\eqref{4.10:eq-holomorphic-maps-a} (polynomial maps in the entries
are holomorphic) both maps are holomorphic. Their inverses are the maps of the same form built from
$v^{-1}$ and from the pointwise inverse $\tilde v^{-1}$, since
$(V^{v})^{v^{-1}}=V=(V^{v^{-1}})^{v}$ bondwise and likewise on the fine bonds; these inverses are again
polynomial, hence holomorphic. Thus both maps are biholomorphisms.

By Proposition~\ref{4.10:prop-continuation-open}, which uses (P9.1) and (P9.2),
$\mathcal{S}(\mathcal{T})$ is open in $\mathbb{D}(\mathcal{T})$ and
$\mathcal{U}_{\mathcal{T}}\colon\mathcal{S}(\mathcal{T})\to\mathrm{Conf}(\mathrm{bd}(Z))$ is
holomorphic. The set $\mathcal{S}(\mathcal{T})^{v}$ is the image of $\mathcal{S}(\mathcal{T})$ under
the biholomorphism $V\mapsto V^{v}$, hence open; therefore the union $\mathcal{O}(\mathcal{T})$ is open.
This is the first assertion.

Let $V\in\mathcal{S}(\mathcal{T})^{v}$. Then $V^{v^{-1}}\in\mathcal{S}(\mathcal{T})$ and
$V=(V^{v^{-1}})^{v}$, so the prescription of (P9.4) assigns to $V$ the value
$\mathcal{U}_{\mathcal{T}}(V^{v^{-1}})^{\tilde v}$. Hence
on the open set $\mathcal{S}(\mathcal{T})^{v}$ the function of (P9.4) is
\begin{equation*}
V\longmapsto \mathcal{U}_{\mathcal{T}}\bigl(V^{v^{-1}}\bigr)^{\tilde v}.
\end{equation*}
By (P9.4) this prescription defines one function on $\mathcal{O}(\mathcal{T})$: the values obtained
from different $v$ at a point of an overlap
$\mathcal{S}(\mathcal{T})^{v}\cap\mathcal{S}(\mathcal{T})^{v'}$ coincide. On
$\mathcal{S}(\mathcal{T})^{v}$ the function is the composition of $V\mapsto V^{v^{-1}}$ (holomorphic by
the first paragraph), of $\mathcal{U}_{\mathcal{T}}$ on $\mathcal{S}(\mathcal{T})$ (holomorphic by
Proposition~\ref{4.10:prop-continuation-open}) and of $U\mapsto U^{\tilde v}$ (holomorphic by the first
paragraph); by the composition fact among~\eqref{4.10:eq-holomorphic-maps-a} it is holomorphic there.
The sets $\mathcal{S}(\mathcal{T})^{v}$, $v\in\Gamma(\mathcal{T})$, form an open cover of
$\mathcal{O}(\mathcal{T})$, and the extended $\mathcal{U}_{\mathcal{T}}$ is holomorphic on each member;
by the fact among~\eqref{4.10:eq-holomorphic-maps-a} that holomorphy is a local property,
$\mathcal{U}_{\mathcal{T}}$ is holomorphic on $\mathcal{O}(\mathcal{T})$. This is the second assertion.

The proofs of Theorem~\ref{4.10:thm-existence-set} and
Corollary~\ref{4.10:cor-representative-analyticity}, run under the hypotheses assumed above, use, of
the set entering the definition of the
existence set $\mathcal{V}(\mathcal{T})$, only that it is open in $\mathbb{D}(\mathcal{T})$ and that
$\mathcal{U}_{\mathcal{T}}$ is holomorphic on it. Both properties hold for $\mathcal{O}(\mathcal{T})$ by
the two assertions just proved, so both proofs apply verbatim with $\mathcal{O}(\mathcal{T})$ in place
of $\mathcal{S}(\mathcal{T})$. This is the third assertion.
\end{proof}

\begin{remark}[Single-valuedness of the continuation near $G$-valued data]
\label{4.10:rem-single-valuedness}
Nothing in this remark is used in the proofs of this book. It concerns the single-valuedness
clause of Ba{\l}aban's analytic continuation of the constrained minimiser
(Definition~\ref{4.10:def-minimiser-continuation}, the second clause of
\eqref{4.10:eq-minimiser-continuation-a}). That clause is stated in substance in
\cite[Proposition~9]{B-CMP102b}. In Ba{\l}aban's construction the continuation is obtained by
solving the equations \cite[(174)--(175)]{B-CMP102b}, or \cite[(179)--(180)]{B-CMP102b}, for a
function $\mathcal{H}(B)$ of the potential $B$. Ba{\l}aban notes that these equations determine
an analytic function $\mathcal{H}$ \cite{B-CMP102b}, and single-valuedness is the uniqueness of
that solution.

Independently of this, near $G$-valued data the single-valuedness clause follows from
\cite[Theorem~1]{B-CMP102b} and the identity theorem. Put
\begin{equation*}
E:=(\mathfrak{g}^c)^{\mathfrak{B}(\mathcal{T})},\qquad
E_{\mathbb{R}}:=(i\mathfrak{g})^{\mathfrak{B}(\mathcal{T})}\subset E .
\end{equation*}
Let $V_0,V_0'\in\mathcal{R}(\mathcal{T})$. Suppose that $V_0$ satisfies \cite[(7)]{B-CMP102b}
with a threshold $\varepsilon_1<\bar\varepsilon$. Suppose further that $B'_*\in\mathcal{B}(V_0')$
satisfies $e^{iB'_*}V_0'=V_0$. Since $B'_*\in\mathcal{B}(V_0')$, the first clause of
\eqref{4.10:eq-minimiser-continuation-a} gives $|B'_*(c)|<\pi$ for every
$c\in\mathfrak{B}(\mathcal{T})$, so Lemma~\ref{4.10:lem-data-parametrisation} applies with
$S=\mathfrak{B}(\mathcal{T})$, the datum $V_0'$ and the potential $B'_*$. Let $\mathcal{W}$ and
$\beta$ be the neighbourhood of $V_0$ and the holomorphic map it provides. Put
$\beta'(B):=\beta(e^{iB}V_0)$ on the set
\begin{equation}\label{4.10:eq-single-valuedness-1}
D_0:=\bigl\{B\in\mathcal{B}(V_0):\ e^{iB}V_0\in\mathcal{W},\
\beta(e^{iB}V_0)\in\mathcal{B}(V_0')\bigr\},
\end{equation}
which is a set on which both representations are admitted. Let $D$ be the connected component of
$D_0$ containing $0$. Then, for every $B\in D$,
\begin{equation}\label{4.10:eq-single-valuedness-2}
e^{i\beta'(B)}V_0'=e^{iB}V_0,\qquad \Phi_{V_0}(B)=\Phi_{V_0'}\bigl(\beta'(B)\bigr).
\end{equation}
On overlaps not connected to $G$-valued data in this way, the single-valuedness clause is taken
from \cite[Proposition~9]{B-CMP102b} as stated in the second clause of
\eqref{4.10:eq-minimiser-continuation-a}; its proof is not reproduced.
\end{remark}

\begin{proof}
\emph{The real form.} Since $\mathfrak{g}^c=\mathfrak{g}\oplus i\mathfrak{g}$ and
$i(i\mathfrak{g})=\mathfrak{g}$, we have $\mathfrak{g}^c=i\mathfrak{g}\oplus i(i\mathfrak{g})$.
Taking this bond by bond gives $E=E_{\mathbb{R}}\oplus iE_{\mathbb{R}}$. In particular
$E_{\mathbb{R}}$ spans $E$ over $\mathbb{C}$. Every $B\in E_{\mathbb{R}}$ has the form $B=iY$ with
$Y\in\mathfrak{g}^{\mathfrak{B}(\mathcal{T})}$. Then $e^{iB}=e^{-Y}$ bondwise, and
$e^{-Y(c)}\in G$ because $Y(c)\in\mathfrak{g}$. Hence the datum $e^{iB}V_0=e^{-Y}V_0$ is
$G$-valued.

\emph{The set $D_0$.} The maps $B\mapsto e^{iB}V_0$ and $\beta$ are continuous, and
$\mathcal{B}(V_0)$, $\mathcal{W}$, $\mathcal{B}(V_0')$ are open. Hence $D_0$ is open in $E$.
It contains $0$, because $0\in\mathcal{B}(V_0)$, $e^{i0}V_0=V_0\in\mathcal{W}$, and
$\beta(V_0)=B'_*\in\mathcal{B}(V_0')$. On $D_0$ the map $\beta'$ is holomorphic, as the composite
of the holomorphic maps $B\mapsto e^{iB}V_0$ and $\beta$. Moreover
$e^{i\beta'(B)}V_0'=e^{iB}V_0$, because $e^{i\beta(V)}V_0'=V$ for $V\in\mathcal{W}$. This is the
first identity of \eqref{4.10:eq-single-valuedness-2}. Consequently, regarding $\Phi_{V_0}$ and
$\Phi_{V_0'}$ as maps into the vector space $M_N(\mathbb{C})^{\mathrm{bd}(Z)}$, the difference
\begin{equation*}
F:=\Phi_{V_0}-\Phi_{V_0'}\circ\beta'
\end{equation*}
is a holomorphic map on $D_0$, and hence on $D$.

\emph{Agreement at real points.} By hypothesis, $V_0$ satisfies \cite[(7)]{B-CMP102b} with the
threshold $\varepsilon_1<\bar\varepsilon$. Hence, for $Y$ small, the data $e^{-Y}V_0$ satisfy
\cite[(7)]{B-CMP102b} with a threshold not exceeding $\bar\varepsilon$. Choose a neighbourhood of
$0$ in $E_{\mathbb{R}}$ consisting of such $B=iY$, and let $N_{\iota}$ be its intersection with
$D$. Then $N_{\iota}$ is a nonempty, relatively open subset of $D\cap E_{\mathbb{R}}$ containing
$0$. For $B=iY\in N_{\iota}$ we apply the third clause of
\eqref{4.10:eq-minimiser-continuation-a} in each admitted representation: in the one
based at $V_0$, at the potential $iY$, and in the one based at $V_0'$, at the potential
$\beta'(iY)$. Both represent the same $G$-valued datum $e^{-Y}V_0$. Together with the uniqueness
assertion of \cite[Theorem~1]{B-CMP102b} for that datum, this gives
$\Phi_{V_0}(iY)=\Phi_{V_0'}(\beta'(iY))$, that is, $F=0$ on $N_{\iota}$.

\emph{The identity theorem along a real form.} Let $F$ be a holomorphic map on a connected open
subset $D$ of $E$. Suppose $F$ vanishes on a nonempty relatively open subset $N_{\iota}$ of
$D\cap(e_0+E_{\mathbb{R}})$ for some $e_0\in E$. We show that $F$ vanishes on $D$.

Take $e_1\in N_{\iota}$. Expand
\begin{equation*}
F(e_1+h)=\sum_{n\ge0}P_n(h),
\end{equation*}
where each $P_n$ is a homogeneous polynomial of degree $n$ on $E$ and the series converges for
$h$ in a ball about $0$. Since $N_{\iota}$ is relatively open, $e_1+h\in N_{\iota}$ for every
sufficiently small $h\in E_{\mathbb{R}}$. Fix such an $h$. For real $t$ near $0$,
\begin{equation*}
0=F(e_1+th)=\sum_{n\ge0}t^nP_n(h).
\end{equation*}
A convergent power series in $t$ that vanishes for all small real $t$ has all its coefficients
zero, so $P_n(h)=0$ for all $n$. By homogeneity, $P_n$ vanishes on all of $E_{\mathbb{R}}$.

Now choose a basis of $E_{\mathbb{R}}$ over $\mathbb{R}$. Because $E=E_{\mathbb{R}}\oplus
iE_{\mathbb{R}}$, it is also a basis of $E$ over $\mathbb{C}$. In the corresponding coordinates,
$P_n$ is a polynomial in finitely many complex variables that vanishes at every real point. By
induction on the number of variables, all its coefficients are zero, since a polynomial in one
variable with infinitely many zeros is zero. Hence every coefficient of the power series of $F$
at $e_1$ vanishes, and $F=0$ on a neighbourhood of $e_1$. Since $D$ is connected, the identity
theorem for holomorphic maps gives $F=0$ on $D$.

\emph{Conclusion.} Apply this with $e_0=0$ to the map $F$ and the set $N_{\iota}$ constructed
above. It follows that $F=0$ on the connected component $D$ of $D_0$ containing the real points of
$N_{\iota}$. This is the second identity of \eqref{4.10:eq-single-valuedness-2}.
\end{proof}

\begin{definition}[Regularity sets and gauge witnesses on a cube]\label{4.10:def-gauge-witnesses}
Let $j\ge1$, $X\in\mathcal{D}_{j-1}$ and the radii $a_0,a_1,\gamma_0>0$ be as in
Definition~\ref{4.10:def-decomposed-representatives}, and let $\mathfrak{P}(X)$ be the space of
triples introduced there. Let
\begin{equation*}
\mathfrak{R}_{123}\subset\mathfrak{P}(X)
\end{equation*}
be the set of triples satisfying the conditions (R1), (R2) and (R3) of
\eqref{4.10:eq-decomposed-representatives-b}. Let $\mathfrak{R}\subset\mathfrak{R}_{123}$ be the
set of decomposed one-power representatives of
Definition~\ref{4.10:def-decomposed-representatives}.

Let $\square$ be a cube of the kind occurring in (R1), and let $\mathrm{s}(\square)$ denote the
set of its sites. For $c>0$ put
\begin{equation}\label{4.10:eq-gauge-witnesses-1}
\begin{split}
\mathcal{C}_{\square}(c):=\bigl\{A\in\mathfrak{g}^{\mathrm{bd}(\square)}:\ {}&|A(b)|<c
\text{ for all } b\in\mathrm{bd}(\square),\\
&\text{and every component of }\nabla^{\xi_j}A\text{ has norm }<c\bigr\}.
\end{split}
\end{equation}
This is an open subset of the real vector space $\mathfrak{g}^{\mathrm{bd}(\square)}$.

Define the bondwise exponential
\begin{equation*}
\mathcal{E}_{\square}:\mathfrak{g}^{\mathrm{bd}(\square)}\to G^{\mathrm{bd}(\square)},
\qquad \mathcal{E}_{\square}(A)(b):=e^{i\xi_jA(b)} .
\end{equation*}
Then set
\begin{equation}\label{4.10:eq-gauge-witnesses-2}
\begin{split}
\mathcal{G}_{\square}(c):=\bigl\{U\in G^{\mathrm{bd}(\square)}:\ {}&\text{there are }
u\in G^{\mathrm{s}(\square)}\text{ and }A\in\mathcal{C}_{\square}(c)\\
&\text{with }U^{u}=\mathcal{E}_{\square}(A)\bigr\},
\end{split}
\end{equation}
where $U^{u}$ is the gauge transform of $U$ by $u$. A pair $(u,A)$ as in
\eqref{4.10:eq-gauge-witnesses-2} is called a \emph{witness} for $U$.

Let $c_{\square}$ and $B$ be the constants of \cite[(1.12)]{B-CMP109}, and put
\begin{equation*}
c_j:=c_{\square}LMBa_0 .
\end{equation*}
The gauge condition contained in (R1) reads, in these terms, as follows. For a triple whose first
entry is $U$, one has
\begin{equation*}
U|_{\mathrm{bd}(\square)}\in\mathcal{G}_{\square}(c_j)
\end{equation*}
for every cube $\square\subset X$ of side at most $c_{\square}LM$ level-$j$ units.
\end{definition}

\begin{proof}[Proof that $\mathcal{C}_{\square}(c)$ is open]
The bond set $\mathrm{bd}(\square)$ is finite, so $\mathfrak{g}^{\mathrm{bd}(\square)}$ is a
finite-dimensional real vector space. For each $b\in\mathrm{bd}(\square)$ the map
$A\mapsto|A(b)|$ is continuous on this space. The scaled lattice difference $\nabla^{\xi_j}$ is
linear, and $\nabla^{\xi_j}A$ has finitely many components. Hence each map sending $A$ to the norm
of one component of $\nabla^{\xi_j}A$ is continuous as well. Every condition in
\eqref{4.10:eq-gauge-witnesses-1} is a strict inequality $<c$ on one of these finitely many
continuous functions. It therefore defines the preimage of the open interval $(-\infty,c)$, which is
an open set. So $\mathcal{C}_{\square}(c)$ is a finite intersection of open sets, and hence open.
\end{proof}

\begin{lemma}[Openness of the non-gauge regularity clauses]\label{4.10:lem-nongauge-openness}
Let $j\ge 1$, $X\in\mathcal{D}_{j-1}$ and the radii $a_0,a_1,\gamma_0>0$ be as in
Definition~\ref{4.10:def-decomposed-representatives}. The subset of $\mathfrak{P}(X)$
defined by the curvature condition of (R1), by (R2) and by (R3) of
\eqref{4.10:eq-decomposed-representatives-b} is open.
\end{lemma}

\begin{proof}
The conditions in question are the curvature condition of (R1), (R2) and (R3) of
\eqref{4.10:eq-decomposed-representatives-b}. Each of them is a strict inequality on the
value at a triple $(U,A',\mathbf{J})\in\mathfrak{P}(X)$ of one of the following functions:
\begin{itemize}
\item $(U,A',\mathbf{J})\mapsto|\partial U(P)-1|$, for a plaquette $P$;
\item $(U,A',\mathbf{J})\mapsto|A'(b)|$, for a bond $b$;
\item the norm of a component of $\nabla^{\xi_j}_U A'$;
\item $(U,A',\mathbf{J})\mapsto|\partial(e^{i\xi_j A'}U)(P)-1|$, for a plaquette $P$;
\item $(U,A',\mathbf{J})\mapsto|\mathbf{J}(b)|$, for a bond $b$.
\end{itemize}
The region $X$ has finitely many plaquettes and finitely many bonds. Hence the set
described in the lemma is cut out by finitely many such strict inequalities.

Each of the listed functions is continuous on $\mathfrak{P}(X)$. This follows from the
holomorphy facts (H2) and (H3) of Definition~\ref{4.10:def-holomorphic-maps}, displayed in
\eqref{4.10:eq-holomorphic-maps-a}, together with the description of the space
$\mathfrak{P}(X)$ and of the maps on it in
Definition~\ref{4.10:def-decomposed-representatives}. The operator norm $|\cdot|$ on
$M_N(\mathbb{C})$ is continuous, and a composition of continuous maps is continuous.

Let $f$ be a continuous real function on $\mathfrak{P}(X)$. The set of points at which $f$
satisfies a given strict inequality is the preimage under $f$ of an open half-line, so it
is open in $\mathfrak{P}(X)$. The set described in the lemma is the intersection of
finitely many such sets. A finite intersection of open sets is open, and therefore this
set is open.
\end{proof}

\begin{lemma}[Openness of the gauge condition below $\pi$]\label{4.10:lem-gauge-openness}
Let $\square$ be a cube of condition (R1) in \eqref{4.10:eq-decomposed-representatives-b}, and
let $c>0$. Let $\mathcal{G}_{\square}(c)\subset G^{\mathrm{bd}(\square)}$ be the gauge-condition
set of Definition~\ref{4.10:def-gauge-witnesses}.
\begin{enumerate}
\item[(i)] If $\xi_j c<\pi$, then $\mathcal{G}_{\square}(c)$ is open in $G^{\mathrm{bd}(\square)}$.
\item[(ii)] For every $c>0$, the set $\mathcal{G}_{\square}(c)$ is a neighbourhood in
$G^{\mathrm{bd}(\square)}$ of each of its points $U_0$ that admits a witness $(u_0,A_0)$ with
$\xi_j|A_0(b)|<\pi$ for all $b\in\mathrm{bd}(\square)$.
\end{enumerate}
\end{lemma}

\begin{proof}
Recall from Definition~\ref{4.10:def-gauge-witnesses} that a witness of a point
$U\in\mathcal{G}_{\square}(c)$ is a pair $(u,A)$, where $u$ is a $G$-valued gauge transformation
on the sites $\mathrm{s}(\square)$ of the cube and $A\in\mathcal{C}_{\square}(c)$, such that the
gauge transform $U^{u}$ equals the bondwise exponential $\mathcal{E}_{\square}(A)$, that is,
$U^{u}(b)=e^{i\xi_j A(b)}$ for every $b\in\mathrm{bd}(\square)$.

\emph{Part (i) from part (ii).} Suppose $\xi_j c<\pi$, and let $U_0\in\mathcal{G}_{\square}(c)$ with
any witness $(u_0,A_0)$. Since $A_0\in\mathcal{C}_{\square}(c)$, we have $|A_0(b)|<c$ for every
bond $b\in\mathrm{bd}(\square)$, and therefore
\begin{equation*}
\xi_j|A_0(b)|<\xi_j c<\pi \qquad\text{for all } b\in\mathrm{bd}(\square).
\end{equation*}
Hence every point of $\mathcal{G}_{\square}(c)$ admits a witness of the kind required in part (ii),
and part (ii) shows that $\mathcal{G}_{\square}(c)$ is a neighbourhood of each of its points;
that is, $\mathcal{G}_{\square}(c)$ is open.

\emph{Part (ii).} Let $U_0\in\mathcal{G}_{\square}(c)$ have a witness $(u_0,A_0)$ with
$\xi_j|A_0(b)|<\pi$ for all $b\in\mathrm{bd}(\square)$.

The map $U\mapsto U^{u_0}$ is a homeomorphism of $G^{\mathrm{bd}(\square)}$ onto itself: on each
bond it is multiplication on the left and on the right by fixed elements of $G$, hence continuous,
and its inverse is $U\mapsto U^{u_0^{-1}}$, which is continuous for the same reason.

Fix $b\in\mathrm{bd}(\square)$. By the witness property, $U_0^{u_0}(b)=e^{i\xi_j A_0(b)}$, and since
$\xi_j|A_0(b)|<\pi$, parts (v) and (iii) of the principal-logarithm lemma
(Lemma~\ref{4.10:lem-principal-logarithm}) give
\begin{equation*}
U_0^{u_0}(b)=e^{i\xi_j A_0(b)}\in\mathcal{L},
\qquad
\log U_0^{u_0}(b)=i\xi_j A_0(b).
\end{equation*}

Put
\begin{equation*}
N_1:=\bigl\{U\in G^{\mathrm{bd}(\square)}:\ U^{u_0}(b)\in\mathcal{L}
\text{ for all } b\in\mathrm{bd}(\square)\bigr\}.
\end{equation*}
By part (i) of Lemma~\ref{4.10:lem-principal-logarithm}, $\mathcal{L}$ is open in
$M_N(\mathbb{C})$. The map $U\mapsto U^{u_0}(b)$ is continuous for each of the finitely many
bonds $b$, so $N_1$ is open. By the display above, $N_1$ contains $U_0$. Let $N_2$ be the connected
component of $N_1$ that contains $U_0$. Since $G^{\mathrm{bd}(\square)}$ is a manifold, it is
locally connected, so the connected components of its open subsets are open; thus $N_2$ is an open
connected neighbourhood of $U_0$.

For $U\in N_2$ and $b\in\mathrm{bd}(\square)$ define
\begin{equation*}
A_U(b):=\frac{1}{i\xi_j}\,\log U^{u_0}(b).
\end{equation*}
The map $U\mapsto A_U$ is continuous on $N_2$. Indeed, it is the composition of the continuous
map $U\mapsto U^{u_0}(b)$, which takes $N_2$ into $\mathcal{L}$, with the logarithm, which is
holomorphic on $\mathcal{L}$ by part (i) of Lemma~\ref{4.10:lem-principal-logarithm}.

We next show that each $A_U(b)$ is an admitted value of a potential. By part (vi) of
Lemma~\ref{4.10:lem-principal-logarithm}, $\log U^{u_0}(b)$ is anti-Hermitian, since
$U^{u_0}(b)$ is unitary and lies in $\mathcal{L}$. Hence $i\xi_j A_U(b)=\log U^{u_0}(b)$ is
anti-Hermitian; equivalently, $iA_U(b)$ is anti-Hermitian. Potentials are written throughout in
the convention $U=e^{i\xi A}$ of \cite{B-CMP109}, with $A$ taking values in $\mathfrak{g}$. In this
convention, the values admitted for a potential are exactly the traceless matrices $A$ for which
$iA$ is anti-Hermitian. It therefore remains to show that $A_U(b)$ is traceless.

Since $U^{u_0}(b)\in G$, part (iv) of Lemma~\ref{4.10:lem-principal-logarithm} gives
$\operatorname{Tr}\log U^{u_0}(b)\in 2\pi i\mathbb{Z}$. The function
$U\mapsto\operatorname{Tr}\log U^{u_0}(b)$ is continuous on $N_2$, takes values in the discrete set
$2\pi i\mathbb{Z}$, and $N_2$ is connected; so it is constant on $N_2$. At $U_0$ its value is
\begin{equation*}
\operatorname{Tr}\log U_0^{u_0}(b)=\operatorname{Tr}\bigl(i\xi_j A_0(b)\bigr)=0,
\end{equation*}
because $A_0(b)$ is the value of a potential and is therefore traceless. Hence
$\operatorname{Tr}A_U(b)=0$ for all $U\in N_2$ and all $b$. Consequently $A_U\in\mathfrak{g}^{\mathrm{bd}(\square)}$
for every $U\in N_2$, and $U\mapsto A_U$ is a continuous map from $N_2$ into the real vector space
$\mathfrak{g}^{\mathrm{bd}(\square)}$.

Two identities about $A_U$ complete the argument. First, by part (ii) of
Lemma~\ref{4.10:lem-principal-logarithm}, $\exp(\log g)=g$ for $g\in\mathcal{L}$. Applied to
$g=U^{u_0}(b)$ on each bond, this gives
\begin{equation*}
\mathcal{E}_{\square}(A_U)(b)=e^{i\xi_j A_U(b)}=\exp\bigl(\log U^{u_0}(b)\bigr)=U^{u_0}(b),
\end{equation*}
that is, $\mathcal{E}_{\square}(A_U)=U^{u_0}$. Second, by the identity
$\log U_0^{u_0}(b)=i\xi_j A_0(b)$ obtained above, $A_{U_0}=A_0$, and $A_0\in\mathcal{C}_{\square}(c)$.

By Definition~\ref{4.10:def-gauge-witnesses}, $\mathcal{C}_{\square}(c)$ is open in
$\mathfrak{g}^{\mathrm{bd}(\square)}$. Since $U\mapsto A_U$ is continuous, the set
\begin{equation*}
N_3:=\bigl\{U\in N_2:\ A_U\in\mathcal{C}_{\square}(c)\bigr\}
\end{equation*}
is open in $N_2$, hence in $G^{\mathrm{bd}(\square)}$, and it contains $U_0$ because
$A_{U_0}=A_0\in\mathcal{C}_{\square}(c)$. For every $U\in N_3$, the pair $(u_0,A_U)$ is a witness
for $U$: $u_0$ is a $G$-valued gauge transformation on $\mathrm{s}(\square)$,
$A_U\in\mathcal{C}_{\square}(c)$, and $U^{u_0}=\mathcal{E}_{\square}(A_U)$. Therefore
$N_3\subset\mathcal{G}_{\square}(c)$. Thus $\mathcal{G}_{\square}(c)$ contains the open
neighbourhood $N_3$ of $U_0$, and is a neighbourhood of $U_0$ in $G^{\mathrm{bd}(\square)}$.
\end{proof}

\begin{proposition}[Openness of the gauge condition beyond $\pi$. Proof outstanding]\label{4.10:prop-gauge-openness-large}
Let $\square$ be a cube of clause (R1) in \eqref{4.10:eq-decomposed-representatives-b}, at level
$j$, and let $c>0$ satisfy $\xi_jc\ge\pi$. Let $\mathcal{C}_{\square}(c)$, $\mathcal{E}_{\square}$
and $\mathcal{G}_{\square}(c)$ be as in Definition~\ref{4.10:def-gauge-witnesses}; thus
$\mathcal{E}_{\square}$ is the bondwise exponential
$A\mapsto(e^{i\xi_jA(b)})_{b\in\mathrm{bd}(\square)}$, and
\begin{equation}\label{4.10:eq-gauge-openness-large-1}
\mathcal{G}_{\square}(c)=\bigl\{U\in G^{\mathrm{bd}(\square)}:\ \text{there are }
u\in G^{\mathrm{s}(\square)}\text{ and }A\in\mathcal{C}_{\square}(c)
\text{ with }U^{u}=\mathcal{E}_{\square}(A)\bigr\}.
\end{equation}
Then $\mathcal{G}_{\square}(c)$ is open in $G^{\mathrm{bd}(\square)}$.
\end{proposition}

Proof outstanding.

\emph{Route.} By \eqref{4.10:eq-gauge-openness-large-1},
\begin{equation}\label{4.10:eq-gauge-openness-large-2}
\mathcal{G}_{\square}(c)=\bigcup_{u\in G^{\mathrm{s}(\square)}}
\bigl(\mathcal{E}_{\square}(\mathcal{C}_{\square}(c))\bigr)^{u^{-1}} .
\end{equation}
Each map $U\mapsto U^{u^{-1}}$ is a homeomorphism of $G^{\mathrm{bd}(\square)}$. By
\eqref{4.10:eq-gauge-openness-large-2}, it would therefore suffice to show that
$\mathcal{E}_{\square}(\mathcal{C}_{\square}(c))$ is open in $G^{\mathrm{bd}(\square)}$. In the
application the constant is $c=c_{\square}LMBa_0$, and the levels concerned are those with
$\xi_jc_{\square}LMBa_0\ge\pi$.

The openness of $\mathcal{G}_{\square}(c_{\square}LMBa_0)$ is also the first conclusion of the
lemma of this chapter stating that the gauge condition of
Definition~\ref{4.4:def-local-data-classes} defines an open set. Its proof rests on the statement
that, for a fixed gauge transformation, the set of fields it carries
into the condition is open, so that the condition defines a union of open sets. For a given $u$
this is the assertion that $(\mathcal{E}_{\square}(\mathcal{C}_{\square}(c)))^{u^{-1}}$ is open,
which is equivalent to the openness of $\mathcal{E}_{\square}(\mathcal{C}_{\square}(c))$ asked for
above. At the levels with $\xi_jc_{\square}LMBa_0\ge\pi$ that lemma's first conclusion holds at
the points covered by Lemma~\ref{4.10:lem-gauge-openness} and is otherwise the content of the
present proposition with $c=c_{\square}LMBa_0$.

What is established is the following. By Lemma~\ref{4.10:lem-gauge-openness}, clause (ii),
$\mathcal{G}_{\square}(c)$ is a neighbourhood of each of its points that admits a witness $(u,A)$
with $\xi_j|A(b)|<\pi$ for all $b\in\mathrm{bd}(\square)$. At exponents of norm below $\pi$ the
map $\exp:\mathfrak{g}\to G$ is a local diffeomorphism. By clause (i) of the same lemma, for every
$c'>0$ with $\xi_jc'<\pi$ the set $\mathcal{G}_{\square}(c')$ is open; in the application these
are the levels with $\xi_jc_{\square}LMBa_0<\pi$.

What remains is openness at the other points, where every witness has a potential with an
exponent of norm at least $\pi$ on some bond. There $\exp:\mathfrak{g}\to G$ need not be locally
surjective. For $N=2$ it is not an open map at potentials with $|\xi_jA(b)|=2\pi$.

The condition in \cite[(1.12)]{B-CMP109} can also be read differently: the potential is taken to
be the principal logarithm, $A:=(i\xi_j)^{-1}\log U^{u}$. Under that reading the corresponding
set is open for every $c$. This follows from the proof of Lemma~\ref{4.10:lem-gauge-openness}(ii),
without the hypothesis on the witness. The two readings coincide when $\xi_jc<\pi$.

The reading of \cite[(1.12)]{B-CMP109} adopted in Definition~\ref{4.4:def-local-data-classes}
and in Definition~\ref{4.10:def-gauge-witnesses} is the existential one: there are
$u$ and $A$ with $U^{u}=\exp(i\xi_jA)$ and $|A|,|\nabla^{\xi_j}A|<c_{\square}LMBa_0$. The
proposition is stated for that reading.

\begin{proposition}[Openness of the tower clause near representatives]
\label{4.10:prop-tower-clause-open}
Let $j\ge1$, $X\in\mathcal{D}_{j-1}$ and $a_0,a_1,\gamma_0>0$. Let $\mathfrak{P}(X)$, the maps
$\sigma_Z,\sigma'_Z$ and the index set $\mathfrak{T}_j(X)$ be as in
Definition~\ref{4.10:def-decomposed-representatives}, and let
$\mathfrak{R}\subset\mathfrak{R}_{123}\subset\mathfrak{P}(X)$ be as in
Definition~\ref{4.10:def-gauge-witnesses}. Assume the following statements:
\begin{itemize}
\item the first two of the statements collected in \eqref{4.10:eq-minimiser-continuation-a}, on
Ba{\l}aban's analytic continuation of the constrained minimiser;
\item the statement on the block averages and the averaged data map
$\dot{\mathcal{M}}_{\mathcal{T}}$;
\item the statement on the current $\mathcal{J}_{\xi}$ and the tower currents
$J_n(Z;\dot{\mathcal{M}}W)$.
\end{itemize}
Then every $t_*=(U_*,A'_*,\mathbf{J}_*)\in\mathfrak{R}$ has an open neighbourhood $N'$ in
$\mathfrak{P}(X)$ such that $N'\cap\mathfrak{R}_{123}\subset\mathfrak{R}$. In particular
$\mathfrak{R}$ is open in the subspace topology of $\mathfrak{R}_{123}$.
\end{proposition}

In this proposition and its proof the letters $N$, $N'$, $N_\iota$ denote neighbourhoods in
$\mathfrak{P}(X)$; they have nothing to do with the rank of the gauge group.

\begin{proof}
The hypotheses listed in the statement are those from which
Theorem~\ref{4.10:thm-existence-set} (openness of the existence set and holomorphy) and
Corollary~\ref{4.10:cor-representative-analyticity} (analyticity at complex one-power
representatives) are proved as implications. Under them we may therefore use both results.

Let
\begin{equation*}
I:=\bigl\{(p,Z,n,\varrho):\ (p,Z)\in\mathfrak{T}_j(X),\ 1\le n\le p,\
\varrho\in\{\mathrm{I},\mathrm{II}\}\bigr\}.
\end{equation*}
The set $\mathfrak{T}_j(X)$ is finite because $X$ is finite. For each of its members $n$ ranges
over finitely many values, and $\varrho$ over two. Hence $I$ is finite. For
$\iota=(p,Z,n,\varrho)\in I$ put $\mathcal{T}_\iota:=\mathcal{T}_{\varrho}(Z,n)$ and
\begin{equation*}
N_\iota:=\sigma_Z^{-1}\bigl(\mathcal{V}(\mathcal{T}_\iota)\bigr)\cap
\sigma_Z'^{-1}\bigl(\mathcal{V}(\mathcal{T}_\iota)\bigr).
\end{equation*}
By clause (i) of Theorem~\ref{4.10:thm-existence-set}, $\mathcal{V}(\mathcal{T}_\iota)$ is an
open subset of $\mathrm{Conf}(\mathrm{bd}(Z))$. The maps $\sigma_Z$ and $\sigma'_Z$ are continuous.
Hence both preimages are open in $\mathfrak{P}(X)$, and so is their intersection $N_\iota$.

By clause (a) of Corollary~\ref{4.10:cor-representative-analyticity}, applied to the
representative $t_*$ and to $(p,Z)\in\mathfrak{T}_j(X)$, $1\le n\le p$,
$\varrho\in\{\mathrm{I},\mathrm{II}\}$, both halves of $t_*$ restricted to $\mathrm{bd}(Z)$ lie in
$\mathcal{V}(\mathcal{T}_\iota)$. That is, $\sigma_Z(t_*)\in\mathcal{V}(\mathcal{T}_\iota)$ and
$\sigma'_Z(t_*)\in\mathcal{V}(\mathcal{T}_\iota)$, so $t_*\in N_\iota$.

Put $N:=\bigcap_{\iota\in I}N_\iota$. As a finite intersection of open sets containing $t_*$, it
is an open neighbourhood of $t_*$. Let $t\in N$ and $\iota\in I$. Then $\sigma_Z(t)$ and
$\sigma'_Z(t)$ belong to $\mathcal{V}(\mathcal{T}_\iota)$. By the definition of the existence
set, the tower fields $U_n(Z;\dot{\mathcal{M}}\sigma_Z(t))$ and
$U_n(Z;\dot{\mathcal{M}}\sigma'_Z(t))$ therefore exist, together with their currents
$J_n(Z;\dot{\mathcal{M}}\sigma_Z(t))$ and $J_n(Z;\dot{\mathcal{M}}\sigma'_Z(t))$.

On $N$ consider, for $\iota=(p,Z,n,\varrho)\in I$, $P\in\mathrm{Pl}(Z^{\wr})$ and
$b\in\mathrm{Bd}(Z^{\wr})$, the real functions
\begin{align*}
f_{\iota,P}(t)&:=\bigl|\partial U_n\bigl(Z;\dot{\mathcal{M}}\sigma_Z(t)\bigr)(P)-1\bigr|,\\
g_{\iota,b}(t)&:=\bigl|J_n\bigl(Z;\dot{\mathcal{M}}\sigma_Z(t)\bigr)(b)\bigr|,
\end{align*}
and let $f'_{\iota,P}$, $g'_{\iota,b}$ be defined in the same way with $\sigma'_Z$ in place of
$\sigma_Z$. Since $I$ is finite and $Z^{\wr}$ has finitely many plaquettes and bonds, there are
finitely many such functions.

Each of them is continuous on $N$. Indeed, it is the composition of the continuous map $\sigma_Z$
(respectively $\sigma'_Z$), which maps $N$ into $\mathcal{V}(\mathcal{T}_\iota)$, with a function
that is continuous on $\mathcal{V}(\mathcal{T}_\iota)$ by clause (iv) of
Theorem~\ref{4.10:thm-existence-set}.

At $t_*$ each of these functions satisfies the corresponding strict inequality of
\eqref{4.10:eq-decomposed-representatives-a} with $c=a_0L^{-(j-p)}$, because $t_*\in\mathfrak{R}$.
Explicitly, for all $\iota$, $P$ and $b$ as above,
\begin{align*}
f_{\iota,P}(t_*),\ f'_{\iota,P}(t_*)&<a_0L^{-(j-p)}\,\xi_p^2,\\
g_{\iota,b}(t_*),\ g'_{\iota,b}(t_*)&<a_0L^{-(j-p)}\,(L^n\xi_p)^2.
\end{align*}
Each strict inequality between a continuous function and a constant defines an open subset of $N$
containing $t_*$. The intersection of these finitely many open sets is therefore an open
neighbourhood $N'\subset N$ of $t_*$ on which all the inequalities hold.

Let $t\in N'$ and $(p,Z)\in\mathfrak{T}_j(X)$. For both halves of $t$ and both rules $\varrho$,
and for every $1\le n\le p$, two facts hold. First, the tower fields and their currents exist,
because $t\in N$. Second, they satisfy the inequalities of
\eqref{4.10:eq-decomposed-representatives-a} with $c=a_0L^{-(j-p)}$, because $t\in N'$. Hence the
clause $(\mathrm{iv})_{p,Z}[a_0L^{-(j-p)}]$ holds at $t$ for both halves and both rules. If
moreover $t\in\mathfrak{R}_{123}$, then $t\in\mathfrak{R}$. This proves
$N'\cap\mathfrak{R}_{123}\subset\mathfrak{R}$.

For the last assertion write $N'(t_*)$ for the neighbourhood just constructed at
$t_*\in\mathfrak{R}$. We claim that
\begin{equation*}
\mathfrak{R}=\mathfrak{R}_{123}\cap\bigcup_{t_*\in\mathfrak{R}}N'(t_*).
\end{equation*}
The inclusion of the left side in the right side holds because $\mathfrak{R}\subset\mathfrak{R}_{123}$
and $t_*\in N'(t_*)$ for every $t_*\in\mathfrak{R}$. The reverse inclusion is the first part,
$N'(t_*)\cap\mathfrak{R}_{123}\subset\mathfrak{R}$, applied to each $t_*$. The union of the open
sets $N'(t_*)$ is open in $\mathfrak{P}(X)$. Hence $\mathfrak{R}$ is the intersection of
$\mathfrak{R}_{123}$ with an open set, that is, it is relatively open in $\mathfrak{R}_{123}$.
\end{proof}

\begin{lemma}[Openness of the triple-to-pair map]\label{4.10:lem-triple-pair-map}
Let $j\ge 1$ and $X\in\mathcal{D}_{j-1}$, and let $\mathfrak{P}(X)$, $\mathfrak{Q}(X)$ and the map
$\rho$ be as in Definition~\ref{4.10:def-decomposed-representatives}, so that
$\rho(U,A',\mathbf{J})=(e^{i\xi_jA'}U,\mathbf{J})$. Let $(U,A',\mathbf{J})\in\mathfrak{P}(X)$
satisfy $\xi_j|A'(b)|<\pi$ for all $b\in\mathrm{bd}(X)$. Then $\rho$ maps every neighbourhood of
$(U,A',\mathbf{J})$ in $\mathfrak{P}(X)$ onto a neighbourhood of $\rho(U,A',\mathbf{J})$ in
$\mathfrak{Q}(X)$.
\end{lemma}

\begin{proof}
Let $\mathcal{N}$ be a neighbourhood of $(U,A',\mathbf{J})$ in $\mathfrak{P}(X)$. It contains a
product $N_U\times N_{A'}\times N_{\mathbf{J}}$ of open neighbourhoods $N_U$ of $U$, $N_{A'}$ of $A'$
and $N_{\mathbf{J}}$ of $\mathbf{J}$, and since $\rho(\mathcal{N})$ contains
$\rho(N_U\times N_{A'}\times N_{\mathbf{J}})$ it suffices to show that the latter set is a
neighbourhood of $\rho(U,A',\mathbf{J})$.

We apply the local holomorphic parametrisation of data by potentials
(Lemma~\ref{4.10:lem-data-parametrisation}) with $S=\mathrm{bd}(X)$, $V_0=U$ and $B_*=\xi_jA'$.
Its hypotheses hold: $U$ is $G$-valued on $\mathrm{bd}(X)$ and $G\subset G^c$, so
$U\in\mathrm{Conf}(\mathrm{bd}(X))$; $B_*\in(\mathfrak{g}^c)^{\mathrm{bd}(X)}$; and, since
$\xi_j=L^{-j}>0$, $|B_*(b)|=\xi_j|A'(b)|<\pi$ for every $b\in\mathrm{bd}(X)$ by assumption. Put
$\mathbf{U}:=e^{i\xi_jA'}U$, which is the point $V_*=e^{iB_*}V_0$ of that lemma. The lemma gives an
open neighbourhood $\mathcal{W}$ of $\mathbf{U}$ in $\mathrm{Conf}(\mathrm{bd}(X))$ and a holomorphic,
in particular continuous, map $\beta:\mathcal{W}\to(\mathfrak{g}^c)^{\mathrm{bd}(X)}$ with
$\beta(\mathbf{U})=\xi_jA'$ and $e^{i\beta(\mathbf{V})}U=\mathbf{V}$ for every
$\mathbf{V}\in\mathcal{W}$.

Define
\begin{equation*}
\mathcal{W}':=\{\mathbf{V}\in\mathcal{W}:\ \xi_j^{-1}\beta(\mathbf{V})\in N_{A'}\}.
\end{equation*}
This set is the preimage of the open set $N_{A'}$ under the continuous map
$\mathbf{V}\mapsto\xi_j^{-1}\beta(\mathbf{V})$ on the open set $\mathcal{W}$, hence open in
$\mathrm{Conf}(\mathrm{bd}(X))$; it contains $\mathbf{U}$ because
$\xi_j^{-1}\beta(\mathbf{U})=A'\in N_{A'}$. For $\mathbf{V}\in\mathcal{W}'$ put
$A'':=\xi_j^{-1}\beta(\mathbf{V})$. Then $A''\in N_{A'}$ and
\begin{equation*}
e^{i\xi_jA''}U=e^{i\beta(\mathbf{V})}U=\mathbf{V}.
\end{equation*}
Consequently, for every $\mathbf{V}\in\mathcal{W}'$ and every $\mathbf{J}'\in N_{\mathbf{J}}$, the
triple $(U,A'',\mathbf{J}')$ lies in $N_U\times N_{A'}\times N_{\mathbf{J}}$, because $U\in N_U$, and
$\rho(U,A'',\mathbf{J}')=(\mathbf{V},\mathbf{J}')$. Hence
\begin{equation*}
\rho(N_U\times N_{A'}\times N_{\mathbf{J}})\supset\mathcal{W}'\times N_{\mathbf{J}}.
\end{equation*}
The right-hand side is the product of an open neighbourhood of $\mathbf{U}$ in
$\mathrm{Conf}(\mathrm{bd}(X))$ and an open neighbourhood of $\mathbf{J}$, hence a neighbourhood of
$(\mathbf{U},\mathbf{J})=\rho(U,A',\mathbf{J})$ in $\mathfrak{Q}(X)$. Therefore
$\rho(N_U\times N_{A'}\times N_{\mathbf{J}})$, and with it $\rho(\mathcal{N})$, is a neighbourhood of
$\rho(U,A',\mathbf{J})$.
\end{proof}

\begin{corollary}[Openness of the one-power class and its gauge saturation]
\label{4.10:cor-class-openness}
Let $\mathfrak{R}_{123}\subset\mathfrak{P}(X)$ and $\mathfrak{R}\subset\mathfrak{R}_{123}$ be the
sets of Definition~\ref{4.10:def-gauge-witnesses}, formed with the conditions (R1)--(R3) of
\eqref{4.10:eq-decomposed-representatives-b}, and put $c_j:=c_{\square}LMBa_0$.
Consider the following hypothesis:
\begin{itemize}
\item[(G)] for every cube $\square$ occurring in (R1), the set $\mathcal{G}_{\square}(c_j)$ is open
in $G^{\mathrm{bd}(\square)}$.
\end{itemize}
By Lemma~\ref{4.10:lem-gauge-openness}(i), hypothesis (G) holds at every level $j$ with
$\xi_jc_{\square}LMBa_0<\pi$; at the remaining levels it is the content of
Proposition~\ref{4.10:prop-gauge-openness-large}, whose proof is outstanding.
\begin{enumerate}
\item If (G) holds, then $\mathfrak{R}_{123}$ is open in $\mathfrak{P}(X)$. Without any proviso,
$\mathfrak{R}_{123}$ is a neighbourhood in $\mathfrak{P}(X)$ of each of its points
$(U,A',\mathbf{J})$ such that, for every cube $\square$ of (R1), the restriction
$U|_{\mathrm{bd}(\square)}$ admits a witness whose exponents have norm below $\pi$, that is, a
witness $(u,A)$ with $\xi_j|A(b)|<\pi$ for all $b\in\mathrm{bd}(\square)$.
\item Assume the first two statements of Ba{\l}aban's analytic continuation of the constrained
minimiser, \eqref{4.10:eq-minimiser-continuation-a} (\cite{B-CMP102b}), together with the
remaining hypotheses under which Proposition~\ref{4.10:prop-tower-clause-open} is proved. Then
$\mathfrak{R}$ is open in the subspace topology of $\mathfrak{R}_{123}$. Hence, whenever
$\mathfrak{R}_{123}$ is open in $\mathfrak{P}(X)$, in particular under (G), the set
$\mathfrak{R}$ is open in $\mathfrak{P}(X)$.
\item Assume that $\mathfrak{R}$ is open in $\mathfrak{P}(X)$ and that $\xi_ja_1<\pi$; the latter
is implied by $a_1\le\pi$, since $\xi_j\le L^{-1}<1$. Then the set $\rho(\mathfrak{R})$ of
one-power representatives is open in $\mathfrak{Q}(X)$, and
\begin{equation*}
U^{\Sigma}_j(X;a_0,a_1,\gamma_0)=\bigcup_{g}\rho(\mathfrak{R})^{g},
\end{equation*}
the union over all $G^c$-valued gauge transformations $g$, is open in $\mathfrak{Q}(X)$.
\end{enumerate}
\end{corollary}

\begin{proof}
\emph{Part (1).} The membership of a triple $(U,A',\mathbf{J})\in\mathfrak{P}(X)$ in
$\mathfrak{R}_{123}$ consists of two groups of conditions. The first group is formed by the
curvature member of (R1), by (R2) and by (R3); by Lemma~\ref{4.10:lem-nongauge-openness} these
conditions define an open subset $\mathcal{O}$ of $\mathfrak{P}(X)$. The second group is the
gauge member of (R1), which is the condition
\begin{equation*}
U|_{\mathrm{bd}(\square)}\in\mathcal{G}_{\square}(c_j)
\end{equation*}
for each of the finitely many cubes $\square$ of (R1). For each such cube the map
$(U,A',\mathbf{J})\mapsto U|_{\mathrm{bd}(\square)}$ is continuous on $\mathfrak{P}(X)$: it is
the projection $(U,A',\mathbf{J})\mapsto U$ followed by the restriction
$U\mapsto U|_{\mathrm{bd}(\square)}$, which is continuous on $G^{\mathrm{bd}(X)}$. Hence, when
$\mathcal{G}_{\square}(c_j)$ is open, which is hypothesis (G), the preimage
$\mathcal{P}_{\square}$ of $\mathcal{G}_{\square}(c_j)$ under this map is open in
$\mathfrak{P}(X)$. Since
\begin{equation*}
\mathfrak{R}_{123}=\mathcal{O}\cap\bigcap_{\square}\mathcal{P}_{\square},
\end{equation*}
the intersection running over the finitely many cubes of (R1), $\mathfrak{R}_{123}$ is a finite
intersection of open sets and therefore open.

For the statement without proviso, let $t_0=(U_0,A'_0,\mathbf{J}_0)\in\mathfrak{R}_{123}$ be such
that, for every cube $\square$ of (R1), the restriction $U_0|_{\mathrm{bd}(\square)}$ admits a
witness $(u,A)$ with $\xi_j|A(b)|<\pi$ for all $b\in\mathrm{bd}(\square)$. By
Lemma~\ref{4.10:lem-gauge-openness}(ii), applied with $c=c_j$, the set
$\mathcal{G}_{\square}(c_j)$ is a neighbourhood of $U_0|_{\mathrm{bd}(\square)}$ in
$G^{\mathrm{bd}(\square)}$, so it contains an open set $N_{\square}$ containing
$U_0|_{\mathrm{bd}(\square)}$. Arguing as above, the preimage of $N_{\square}$ under the continuous
map $(U,A',\mathbf{J})\mapsto U|_{\mathrm{bd}(\square)}$ is an open subset of $\mathfrak{P}(X)$
containing $t_0$ and contained in $\mathcal{P}_{\square}$. The intersection of these finitely many
preimages with $\mathcal{O}$ is an open set containing $t_0$ and contained in
$\mathfrak{R}_{123}$; thus $\mathfrak{R}_{123}$ is a neighbourhood of $t_0$.

\emph{Part (2).} By Proposition~\ref{4.10:prop-tower-clause-open}, which is proved as an
implication from \eqref{4.10:eq-minimiser-continuation-a} and the further hypotheses assumed
in part (2), $\mathfrak{R}$ is open in the subspace topology of $\mathfrak{R}_{123}$, that is,
$\mathfrak{R}=\mathcal{N}\cap\mathfrak{R}_{123}$ for some open set $\mathcal{N}$ of
$\mathfrak{P}(X)$. If $\mathfrak{R}_{123}$ is open in $\mathfrak{P}(X)$, then $\mathfrak{R}$ is
the intersection of two open sets of $\mathfrak{P}(X)$ and is open: a relatively open subset of an
open set is open. Under (G), $\mathfrak{R}_{123}$ is open by part (1).

\emph{Part (3).} By Lemma~\ref{4.10:lem-triple-pair-map}, $\rho$ maps every neighbourhood of a
point $(U,A',\mathbf{J})\in\mathfrak{P}(X)$ with $\xi_j|A'(b)|<\pi$ for all
$b\in\mathrm{bd}(X)$ onto a neighbourhood of $\rho(U,A',\mathbf{J})$ in $\mathfrak{Q}(X)$. Let
$t=(U,A',\mathbf{J})\in\mathfrak{R}$. Every element of $\mathfrak{R}$ satisfies $|A'(b)|<a_1$ for
all $b\in\mathrm{bd}(X)$; since $\xi_j>0$ and $\xi_ja_1<\pi$ by assumption,
\begin{equation*}
\xi_j|A'(b)|<\xi_ja_1<\pi\qquad\text{for all } b\in\mathrm{bd}(X),
\end{equation*}
so Lemma~\ref{4.10:lem-triple-pair-map} applies at $t$. As $\mathfrak{R}$ is open in
$\mathfrak{P}(X)$ by assumption, it is a neighbourhood of $t$; hence $\rho(\mathfrak{R})$ is a
neighbourhood of $\rho(t)$ in $\mathfrak{Q}(X)$. Since every point of $\rho(\mathfrak{R})$ is of
the form $\rho(t)$ with $t\in\mathfrak{R}$, the set $\rho(\mathfrak{R})$ is a neighbourhood of each
of its points, and is therefore open in $\mathfrak{Q}(X)$.

Each $G^c$-valued gauge transformation $g$ acts on $\mathfrak{Q}(X)$ by the map
$(\mathbf{U},\mathbf{J})\mapsto(\mathbf{U}^{g},\mathbf{J}^{g})$. This map is continuous, and its
inverse is the action of $g^{-1}$, which is continuous as well; so the action of $g$ is a
homeomorphism of $\mathfrak{Q}(X)$. Consequently each image $\rho(\mathfrak{R})^{g}$ of the open
set $\rho(\mathfrak{R})$ is open, and the union
$U^{\Sigma}_j(X;a_0,a_1,\gamma_0)=\bigcup_{g}\rho(\mathfrak{R})^{g}$, as a union of open sets, is
open in $\mathfrak{Q}(X)$.
\end{proof}

\begin{remark}[Scope of the statements and two radius ceilings]\label{4.10:rem-radius-ceilings}
The following are not asserted in this section, and none of them is used or needed in it.
\begin{itemize}
\item The existence of the tower field $U_n(Z;\cdot)$ on the whole set cut out by the
conditions (R1)--(R3).
\item The deduction of the tower clause (iv) from the preceding clauses (i)--(iii) of the
definition of the one-power representatives, taken at smaller constants,
as in \cite[(1.17)]{B-CMP109}. It enters the construction as a separate input and is not
derived here.
\item The comparison of tower fields along two towers.
\item Any block-law or comparison statement.
\item Any inclusion of classes across widths.
\item Any estimate of $U_n$ or of $J_n$.
\item Any estimate of the size of the existence sets $\mathcal{V}(\mathcal{T})$.
\item Any uniformity of the sets $\mathcal{V}(\mathcal{T})$ in $Z$, $n$, $p$, $j$, $M$ or the
rule $\varrho$.
\end{itemize}
Theorem~\ref{4.10:thm-existence-set} and Corollary~\ref{4.10:cor-representative-analyticity}
are qualitative statements of openness and holomorphy. What is used of them is only the
openness of the sets on which the analytic maps are composed, and
analyticity by composition on those sets. The Cauchy estimates of the induction do not take
their radii from these two statements. They take them from the fact that the pairs produced by
the step lie in the domain uniformly in the Cauchy parameters. No numerical value is used
anywhere in this section. The quantities $L$, $M=L^{m'}$, $N$, $R_1$, $M_1$, $j$, $p$, $n$ and
the coupling remain symbolic and unrestricted.

Theorem~\ref{4.10:thm-existence-set}, Corollary~\ref{4.10:cor-representative-analyticity} and
Proposition~\ref{4.10:prop-tower-clause-open} carry no side condition.

The printed input defining the admitted data carries Ba{\l}aban's smallness threshold
$\varepsilon_1\le\bar\varepsilon(d,L)$ of \cite[Theorem~1 and \S G]{B-CMP102b}, where
$\varepsilon_1$ is the smallness parameter of that theorem. This threshold enters only the
definition of the set $\mathcal{R}(\mathcal{T})$ of
admitted $G$-valued data. It is not a condition on the radii $a_0$, $a_1$, $\gamma_0$ of this
section.

Corollary~\ref{4.10:cor-class-openness}, in its third conclusion, and
Lemma~\ref{4.10:lem-triple-pair-map} use condition (E1) below. The first conclusion of
Corollary~\ref{4.10:cor-class-openness} is proved under the sufficient condition (E2):
\begin{equation}\label{4.10:eq-radius-ceilings-a}
\text{(E1)}\quad \xi_j\,a_1<\pi,
\qquad\qquad
\text{(E2)}\quad \xi_j\,c_{\square}LMB\,a_0<\pi .
\end{equation}
Condition (E1) is implied by the one-sided ceiling $a_1\le\pi$ on the radius $a_1$, a ceiling
that depends on nothing.

By the identity \eqref{4.10:eq-radius-ceilings-1} below, condition (E2) reads
$L^{m'+1-j}c_{\square}B\,a_0<\pi$. At the levels $j\ge m'+1$ it is implied by the one-sided
ceiling $c_{\square}B\,a_0<\pi$ on $a_0$, in terms of the constants $c_{\square}$ and $B$ of
\cite[(1.12)]{B-CMP109}.

At the levels $1\le j\le m'$, condition (E2) is not adopted. At fixed $a_0$ it would bound
$L^{m'+1-j}$ from above, which is a cap in the sense of Definition~\ref{2.3:def-cap}. At those
levels the openness of $\mathcal{G}_{\square}(c_j)$ that (E2) would deliver through
Lemma~\ref{4.10:lem-gauge-openness} falls, whenever $\xi_jc_j\ge\pi$, under
Proposition~\ref{4.10:prop-gauge-openness-large}, whose proof is outstanding.
\end{remark}

\begin{proof}
We verify the identity behind the second form of (E2), the two implications from the ceilings,
and the reading of (E2) as a cap at the levels $1\le j\le m'$.

Throughout, $j\ge1$ and the radii $a_0$, $a_1$ are positive, as in
Corollary~\ref{4.10:cor-representative-analyticity}. The constants $c_{\square}$ and $B$ are
positive.

Since $\xi_j=L^{-j}$ and $M=L^{m'}$, we have
\begin{equation}\label{4.10:eq-radius-ceilings-1}
\xi_j\,c_{\square}LMB\,a_0
=L^{-j}\cdot L\cdot L^{m'}\,c_{\square}B\,a_0
=L^{m'+1-j}\,c_{\square}B\,a_0 .
\end{equation}
This is the second form of (E2).

For (E1): since the blocking factor satisfies $L>1$ and since $j\ge1$, we have
$\xi_j=L^{-j}<1$. Since $a_1>0$,
the ceiling $a_1\le\pi$ gives
\begin{equation*}
\xi_j\,a_1<a_1\le\pi .
\end{equation*}

For (E2) at $j\ge m'+1$: the exponent $m'+1-j$ is at most $0$, so $L^{m'+1-j}\le1$. Since
$c_{\square}B\,a_0>0$, the identity \eqref{4.10:eq-radius-ceilings-1} and the ceiling
$c_{\square}B\,a_0<\pi$ give
\begin{equation*}
\xi_j\,c_{\square}LMB\,a_0
=L^{m'+1-j}\,c_{\square}B\,a_0
\le c_{\square}B\,a_0<\pi .
\end{equation*}

For the levels $1\le j\le m'$: the exponent $m'+1-j$ is at least $1$. Hence at fixed $a_0$ the
inequality $L^{m'+1-j}c_{\square}B\,a_0<\pi$ is an upper bound on the positive power
$L^{m'+1-j}$ of $L$. At $j=1$ it reads $Mc_{\square}B\,a_0<\pi$.
\end{proof}

\begin{definition}[Lattice, complex configurations, decompositions and block averages]
\label{4.11:not-lattice-configurations}
The following conventions are in force in this section.

\emph{Lattice and units.} Let $d\ge 2$ ($d=4$ in the application), and let $L\ge 2$ be an odd
integer. All levels live on one finite periodic lattice $\mathbb{L}=(\mathbb{Z}/\ell_{\mathbb{L}}\mathbb{Z})^d$,
with period a positive integer $\ell_{\mathbb{L}}$ and spacing $1$; lengths measured on
$\mathbb{L}$ are said to be in \emph{fine units}. Lengths in
level-$i$ units are fine lengths multiplied by $L^{-i}$, and $\xi_i:=L^{-i}$ is the fine spacing
in level-$i$ units. Bonds are $b=\langle x,x+e_\mu\rangle$ together with their reverses $\bar b$;
for a bond $b$ we write $b_-$ for its initial point. Plaquettes $p$ are the unit squares of
$\mathbb{L}$; for a bond $b\in p$, $\partial_b p$ denotes the boundary contour of $p$ that starts at
$b_-$ and runs through $b$. For $v\ge 0$ the \emph{$v$-blocks} are the cubes of side $L^v$ of the
$v$-th $L$-adic partition of $\mathbb{L}$, so that each $(v+1)$-block is the union of $L^d$
$v$-blocks; the $v$-sites are the $v$-blocks, and the $v$-bonds are the ordered pairs of adjacent
$v$-blocks. The numbers $M_1,R_1\ge 1$ are Ba{\l}aban's tower parameters, and a
\emph{$v$-superblock} is a cube of side $M_1L^v$ of a fixed partition compatible with the
$v$-blocks. The cube parameter is $M=L^{m'}$, and $\pi_i$ is the partition of $\mathbb{L}$ into
cubes of side $ML^i$ (in fine units). For a union $Z$ of $\pi_c$-cubes, $Z^{\wr,c}$ denotes $Z$
less the $\pi_c$-cubes lying within two layers of $\partial Z$.

\emph{Group, configurations, gauge.} Let $G\subset U(N)$ be a compact matrix group with Lie
algebra $\mathfrak{g}$ ($G=SU(N)$ in the application), let $G^{c}$ and $\mathfrak{g}^{c}$ be the
complexifications, and let $|\cdot|$ be the operator norm. A configuration on a set of bonds is a
map $\mathbf{U}\colon b\mapsto\mathbf{U}(b)\in G^{c}$ with $\mathbf{U}(\bar b)=\mathbf{U}(b)^{-1}$.
For a plaquette $p$ and a fixed first bond $b\in p$, $\partial\mathbf{U}(p)$ is the ordered product
of the values of $\mathbf{U}$ around $\partial_b p$. A change of the first bond replaces
$\partial\mathbf{U}(p)$ by a conjugate $W\,\partial\mathbf{U}(p)\,W^{-1}$, where $W$ is a product
of values of $\mathbf{U}$. For unitary $W$ one has $|WXW^{-1}-1|=|X-1|$ exactly, the operator norm
being unitarily invariant; in particular, for $G$-valued $\mathbf{U}$ the quantity
$|\partial\mathbf{U}(p)-1|$ does not depend on the choice of the first bond.

A gauge transformation is a map $g$ from sites to $G^{c}$; it acts by
\begin{equation*}
\mathbf{U}^{g}(\langle x,y\rangle)=g(x)\,\mathbf{U}(\langle x,y\rangle)\,g(y)^{-1}.
\end{equation*}
Let $\mathbf{J}$ be a $\mathfrak{g}^{c}$-valued bond field, transforming by
$\mathbf{J}\mapsto\operatorname{Ad}_g\mathbf{J}$. The $G^{c}$-orbit of the pair
$(\mathbf{U},\mathbf{J})$ is denoted $[(\mathbf{U},\mathbf{J})]$ (cf.~\cite{B-CMP109}).

A \emph{decomposition} of $\mathbf{U}$ is a pair $(U',U)$ with $U$ $G$-valued and
$U'(b)=\exp(i\mathsf{A}'(b))$ for a $\mathfrak{g}^{c}$-valued bond potential $\mathsf{A}'$, such
that $\mathbf{U}=U'U$ bondwise; $U$ is called the \emph{$G$-valued factor of the decomposition}. A
decomposition is not canonical: the membership conditions imposed in this section are asked of a
chosen decomposition, called the \emph{witness}. A \emph{real point} is a pair $(U,\mathbf{J})$
with $U$ $G$-valued, presented with the trivial witness $(1,U)$.

\emph{Block averages.} For $v\ge 1$ and a $v$-bond $c=\langle B,B'\rangle$,
$M^{v}(\mathbf{U})(c)\in G^{c}$ denotes Ba{\l}aban's $v$-fold block average \cite{B-CMP98},
\cite[(1.5)]{B-CMP99a}. The following properties are used, and only these.
\begin{itemize}
\item[(Mloc)] $M^{v}(\mathbf{U})(c)$ depends only on the restriction of $\mathbf{U}$ to the bonds
with both endpoints in $B\cup B'$; it depends analytically on $\mathbf{U}$ near $G$-valued
$\mathbf{U}$, and it is $G$-valued for $G$-valued $\mathbf{U}$.
\item[(Mnest)] $M^{v}(\mathbf{U})(c)$ is a fixed function of the values
$M^{v-1}(\mathbf{U})(c')$ for the $(v-1)$-bonds $c'$ with endpoints in $B\cup B'$ (iterated
averaging, \cite[(1.5)]{B-CMP99a}).
\item[(M2)] (\cite[Proposition~2]{B-CMP98}.) There is a positive threshold depending only on $d$
and $L$ such that for $G$-valued $U$ with $\sup|\partial U-1|$ below this threshold
\begin{equation*}
|\partial(M^{v}U)-1|\le 2L^{2v}\sup|\partial U-1|
\end{equation*}
for every plaquette of the lattice of $v$-blocks, lengths being in fine units.
\item[(M1)] (\cite[Propositions~6--7]{B-CMP98}.) There is a constant $c_Q=c_Q(d,L)$ with the
following property. Let $U$ be $G$-valued and regular in the sense of the cited propositions, and
let $\mathsf{A}$ be a small $\mathfrak{g}^{c}$-valued bond field. Then, blockwise,
\begin{equation*}
M^{v}(e^{i\mathsf{A}}U)=\bigl[e^{i\mathcal{B}}M^{v}(U)\bigr]^{u},
\end{equation*}
where $u$ is a block gauge, that is, a gauge transformation on the $v$-sites, and $\mathcal{B}$ is a
$\mathfrak{g}^{c}$-valued field on $v$-bonds, both analytic in $\mathsf{A}$, and $|\mathcal{B}|$ and
$|\log u|$ are each at most $c_Q L^{v}\sup|\mathsf{A}|$, the supremum being taken over the two
blocks of the $v$-bond.
\end{itemize}
\end{definition}

\begin{definition}[The large-field tower, tower rules and constrained minimisers]
\label{4.11:def-tower-rules}
The lattice $\mathbb{L}$, the spacings $\xi_i=L^{-i}$, the $p$-blocks, $p$-bonds and
$p$-superblocks (cubes of side $M_1L^p$), the tower parameters $M_1,R_1$, the
configurations $\mathbf{V}$ and the block averages $M^{q}(\mathbf{V})(c)$ are those of
Definition~\ref{4.11:not-lattice-configurations}.

\emph{The tower and its layers.} Fixed once and for all are the regions produced by the
first $k+1$ steps of the construction of \cite{B-CMP119},
\begin{equation*}
\mathbb{L}\supset\Omega_1\supset\Omega_2\supset\dots\supset\Omega_{k+1},
\end{equation*}
each $\Omega_n$ a union of $n$-superblocks, together with the small-field regions
$\Lambda_n\subset\Omega_n$ of \cite[(2.1)]{B-CMP119}. The \emph{layers} are
$\mathfrak{L}_n:=\Omega_n\setminus\Omega_{n+1}$ for $1\le n\le k$, together with
$\Omega_{k+1}$. At a reading level $i\le k+1$ the \emph{layer index} of a point
$x\in\Omega_1$ is $n_i(x):=n$ if $x\in\mathfrak{L}_n$ with $n<i$, and $n_i(x):=i$ if
$x\in\Omega_i$. For $j\le i$ and $x\in\Omega_1$,
\begin{equation}\label{4.11:eq-tower-rules-1}
n_j(x)=\min\bigl(n_i(x),j\bigr).
\end{equation}

\emph{Tower rules.} A \emph{tower rule} $\mathcal{T}$ assigns to a region
$Z\subset\mathbb{L}$ and an index $t\ge1$ a decreasing sequence
\begin{equation*}
Z=\Omega_0^{\mathcal{T}}(Z)\supset\Omega_1^{\mathcal{T}}(Z)\supset\dots\supset
\Omega_t^{\mathcal{T}}(Z),
\end{equation*}
each $\Omega_q^{\mathcal{T}}(Z)$ with $q\ge1$ a union of $q$-superblocks, determined by
the set $Z$, by $q$ and by the fixed regions of the tower alone. Two rules occur.
The \emph{maximal rule} $\mathcal{T}=\mathrm{max}$ is that of
\cite[(1.3)--(1.6)]{B-CMP99a} and \cite{B-CMP102b}: $\Omega_q^{\mathrm{max}}(Z)$ is the
maximal union of $q$-superblocks at distance greater than $R_1M_1L^{q-1}$ from the
complement of $\Omega_{q-1}^{\mathrm{max}}(Z)$. The second rule, the \emph{layered rule}
$\mathcal{T}=\mathrm{lay}$, is that of \cite[(2.13)]{B-CMP119}; its determining set is
\begin{equation*}
\Gamma_t(Z)\cup\bigcup_{i<t}\Gamma_i,
\end{equation*}
where, for $i<t$, $\Gamma_i$ is the set of $i$-bonds of the layer $\mathfrak{L}_i$, the
layers being those met by $Z$, and $\Gamma_t(Z)$ is the set of $t$-bonds of
$Z\cap\Omega_t$.

\emph{Data.} Let $\mathcal{T}$ be a rule, $Z$ a region, $t$ an index and $\mathbf{V}$ a
configuration on the bonds of $Z$. The \emph{fixed bonds} are the bonds of $Z$ not having
both endpoints in $\Omega_1^{\mathcal{T}}(Z)$; $\Gamma_q^{\mathcal{T}}(Z)$ is the set of
$q$-bonds with both endpoint blocks in
$\Omega_q^{\mathcal{T}}(Z)\setminus\Omega_{q+1}^{\mathcal{T}}(Z)$; and
$\Gamma_t^{\mathcal{T},\mathrm{top}}(Z)$ is the set of $t$-bonds with both endpoint blocks
in $\Omega_t^{\mathcal{T}}(Z)$. The \emph{data} of $\mathbf{V}$ are
\begin{align*}
D_{\mathcal{T}}^{t,Z}(\mathbf{V}):=\Bigl(\;&\mathbf{V}\text{ on the fixed bonds};\;
M^{q}(\mathbf{V})(c)\text{ for }c\in\Gamma_q^{\mathcal{T}}(Z),\ 1\le q<t;\\
&M^{t}(\mathbf{V})(c)\text{ for }c\in\Gamma_t^{\mathcal{T},\mathrm{top}}(Z)\Bigr).
\end{align*}

\emph{The constrained problem.} For data $D$, the problem $P_{\mathcal{T}}(Z,t;D)$ is: minimise
the Wilson action
\begin{equation*}
S_Z(V)=\sum_{q\subset Z}\bigl(1-\operatorname{Re}\operatorname{tr}\partial V(q)\bigr),
\end{equation*}
the sum running over the plaquettes $q$ contained in $Z$, over $G$-valued $V$ on the bonds
of $Z$ with $D_{\mathcal{T}}^{t,Z}(V)=D$. By \cite[Theorem~1]{B-CMP102b}, for $G$-valued
data obeying the small-field condition \cite[(7)]{B-CMP102b} with parameter
$\varepsilon_1\le a_1(d,L)$, where $a_1(d,L)$ is the constant of that theorem and depends on
$d$ and $L$ only, this problem has a unique minimal orbit;
$U_t^{\mathcal{T}}(Z;D)$ denotes its element in the axial gauge. For $G^{c}$-valued data in
the neighbourhood of $G$-valued data in which the axial-gauge minimiser continues
analytically \cite[Proposition~9]{B-CMP102b}, $U_t^{\mathcal{T}}(Z;D)$ denotes the analytic
continuation there.

\emph{Tower quantities.} We write
\begin{equation*}
U_t^{\mathcal{T}}(Z;\dot{\mathcal{M}}\mathbf{V}):=
U_t^{\mathcal{T}}\bigl(Z;D_{\mathcal{T}}^{t,Z}(\mathbf{V})\bigr),
\end{equation*}
and $J_t^{\mathcal{T}}(Z;\dot{\mathcal{M}}\mathbf{V})$
for the current of \cite[(1.8)]{B-CMP109} built on
$U_t^{\mathcal{T}}(Z;\dot{\mathcal{M}}\mathbf{V})$ with $\xi$ replaced by $\xi_t=L^{-t}$
(the \emph{level-$t$ normalisation}). The pair
$\bigl(U_t^{\mathcal{T}}(Z;\dot{\mathcal{M}}\mathbf{V}),
J_t^{\mathcal{T}}(Z;\dot{\mathcal{M}}\mathbf{V})\bigr)$ is the pair of \emph{tower
quantities} of $\mathbf{V}$ at $(\mathcal{T},t,Z)$; they are \emph{defined at}
$\mathbf{V}$ when $D_{\mathcal{T}}^{t,Z}(\mathbf{V})$ lies in the domain just described
($G$-valued data under the small-field condition, or $G^{c}$-valued data in the
neighbourhood of analytic continuation). Data lying in this domain are called
\emph{admissible data} for $(\mathcal{T},t,Z)$; thus the tower quantities are defined at
$\mathbf{V}$ exactly when $D_{\mathcal{T}}^{t,Z}(\mathbf{V})$ are admissible data for
$(\mathcal{T},t,Z)$. By construction,
\begin{equation}\label{4.11:eq-tower-rules-a}
\parbox{0.86\linewidth}{the tower quantities at $(\mathcal{T},t,Z)$ depend on
$(\mathcal{T},t,Z)$ and on $\mathbf{V}$ restricted to the bonds of $Z$, and on nothing
else; in particular not on the level at which a configuration class (these classes are
defined later in this chapter) is read, nor on any ambient domain containing $Z$.}
\end{equation}
Indeed, the sequence $\Omega_q^{\mathcal{T}}(Z)$ is determined by $Z$, $q$ and the fixed
tower; the data $D_{\mathcal{T}}^{t,Z}(\mathbf{V})$ are computed from $\mathbf{V}$, which
is a configuration on the bonds of $Z$; and $U_t^{\mathcal{T}}$ and $J_t^{\mathcal{T}}$ are
defined from $(\mathcal{T},t,Z)$ and these data alone.

\emph{Position conventions.} For a plaquette $q$ (respectively a bond $b$) we say that $q$
(respectively $b$) is \emph{at} $x$ if $x$ is its lowest corner (respectively its initial
point). A bound on $|\partial V-1|$ \emph{at} $x$ is a bound on $|\partial V(q)-1|$ for all
plaquettes $q$ at $x$.
\end{definition}

\begin{proof}[Proof of \eqref{4.11:eq-tower-rules-1}]
Since $\Omega_1\supset\Omega_2\supset\dots\supset\Omega_{k+1}$, the set $\Omega_1$ is the
disjoint union of $\mathfrak{L}_1,\dots,\mathfrak{L}_{i-1}$ and $\Omega_i$, so $n_i(x)$ is
well defined for $x\in\Omega_1$; moreover $x\in\Omega_{n_i(x)}$ in either case of the
definition, because $\mathfrak{L}_n\subset\Omega_n$. Let $j\le i$.

If $n_i(x)<j$, then $n_i(x)<i$, so $x\in\mathfrak{L}_n$ with $n=n_i(x)<j$; by the
definition of the layer index at level $j$, $n_j(x)=n=\min(n_i(x),j)$.

If $n_i(x)\ge j$, then $x\in\Omega_{n_i(x)}\subset\Omega_j$, so $n_j(x)=j=\min(n_i(x),j)$.
\end{proof}

\begin{definition}[Radii, schedules, index sets and reading sets]\label{4.11:def-radii-schedules}
Let $g_n$, $n\ge 1$, be the running couplings.

\emph{Radii.} For $n\ge 1$ we use the radii of \cite[(2.28)]{B-CMP119}:
\begin{equation}\label{4.11:eq-radii-schedules-1}
\alpha_{0,n}:=g_n C_0\bigl(\log g_n^{-2}\bigr)^{q_0},\qquad
\alpha_{1,n}:=g_n C_1\bigl(\log g_n^{-2}\bigr)^{q_1}.
\end{equation}

\emph{Schedule parameters.} Fix $\beta\in(0,1)$ and $\theta\in(0,1)$. Ba{\l}aban's construction
takes, for example, $\beta=1/4$. These two letters are local to this section. They are not the
coefficients $\beta_{k+1}$ and not the parameter $\vartheta$ of the glossary.

\emph{Schedules.} For $n\le i$ put
\begin{equation}\label{4.11:eq-radii-schedules-2}
f_{i,n}:=1-\beta\bigl(1-2^{-(i-n)}\bigr)=1-\beta+\beta\,2^{-(i-n)} .
\end{equation}
Then $f_{i,i}=1-\beta(1-1)=1$. For fixed $n$, the number $2^{-(i-n)}$ decreases as $i$
increases. Since $\beta>0$, the second expression in \eqref{4.11:eq-radii-schedules-2} shows
that $f_{i,n}$ is non-increasing in $i$.

Let $k$ be the level. Put
\begin{equation}\label{4.11:eq-radii-schedules-3}
F_n:=f_{k+1,n}+\theta\beta\,2^{-(k+1-n)}\quad (n\le k),\qquad
F_{k+1}:=F_{\mathrm{top}}:=1+\beta .
\end{equation}
A \emph{schedule at level $i$} is a map $\varphi\colon\{1,\dots,i\}\to(0,\infty)$.
\begin{itemize}
\item The \emph{standard schedule} at level $i$ is $\varphi^{\mathrm{st}}_i(n):=f_{i,n}$.
\item The \emph{post schedule} at level $k+1$ is $\varphi^{\mathrm{post}}(n):=F_n$ for $n\le k$,
and $\varphi^{\mathrm{post}}(k+1):=1+\beta$.
\end{itemize}

\emph{Admissible regions.} For each $p\ge 1$ a family $\mathcal{A}_p$ of \emph{admissible
regions of level $p$} is fixed. Two families serve:
\begin{itemize}
\item the finite unions of cubes of the partition $\pi_{p-1}$ of the lattice into cubes of
side $ML^{p-1}$;
\item the localisation domains of \cite{B-CMP119}.
\end{itemize}
What is used about $\mathcal{A}_p$ in what follows is only this: whether $Z$ is admissible
depends on the pair $(p,Z)$ alone.

\emph{Index sets.} For a region $X$ and $i\ge 1$ put
\begin{equation}\label{4.11:eq-radii-schedules-4}
\mathfrak{T}_i(X):=\bigl\{(t,p,Z):\ 1\le t\le p\le i,\ Z\in\mathcal{A}_p,\ Z\subset X\bigr\}.
\end{equation}
In a triple $(t,p,Z)$, $t$ is the tower index and $p$ is the region level. These sets are
monotone:
\begin{equation}\label{4.11:eq-radii-schedules-a}
X\subset Y,\quad i'\le i\quad\Longrightarrow\quad
\mathfrak{T}_{i'}(X)\subset\mathfrak{T}_i(Y).
\end{equation}
Indeed, let $(t,p,Z)\in\mathfrak{T}_{i'}(X)$. Then $1\le t\le p\le i'\le i$, $Z\in\mathcal{A}_p$
and $Z\subset X\subset Y$. So the triple satisfies every condition in
\eqref{4.11:eq-radii-schedules-4} with $(i,Y)$ in place of $(i',X)$.

\emph{Reading sets.} Write $Z^{\wr,c}$ for $Z$ with the $\pi_c$-cubes lying within two layers
of $\partial Z$ removed. To each pair $(p,Z)$ is attached the \emph{reading set}
\begin{equation*}
P(p,Z):=Z^{\wr,p-1}\subset Z .
\end{equation*}
Thus $P(p,Z)$ is $Z$ less two layers of $\pi_{p-1}$-cubes. It plays the role of the set on which
bounds are read, written $X^{\sim 2}$ in \cite[(iv)]{B-CMP109}.

To each triple $(t,p,Z)$ is also attached a set $\mathcal{R}(t,p,Z)\subset\{\max,\,\mathrm{lay}\}$.
It lists the tower rules of Definition~\ref{4.11:def-tower-rules} that are read at the triple:
\begin{itemize}
\item both rules are read for the interior class;
\item the layered rule is read at straddling sources.
\end{itemize}
The interior class and the straddling sources are those fixed in the definition of the
configuration classes later in this section.
The sets $P$ and $\mathcal{R}$ depend on $(t,p,Z)$ only.
\end{definition}

\begin{definition}[Bounds on the tower quantities in two normalisations]\label{4.11:def-tower-bounds}
Let $X$ be a region and let $\mathbf{V}$ be a configuration on $X$.

The following objects are those of Definition~\ref{4.11:def-radii-schedules}:
\begin{itemize}
\item the index sets $\mathfrak{T}_i(X)$ of triples $(t,p,Z)$;
\item the reading sets $P(p,Z)$;
\item the sets $\mathcal{R}(t,p,Z)$ of tower rules read at a triple;
\item the schedules and the radii $\alpha_{0,n}$.
\end{itemize}
The following objects are those of Definition~\ref{4.11:def-tower-rules}:
\begin{itemize}
\item the tower regions $\Omega_n$ and the layer index $n_i(x)$;
\item the tower rules $\mathcal{T}$;
\item the tower quantities of $\mathbf{V}$ at $(\mathcal{T},t,Z)$, namely the constrained minimiser $U_t^{\mathcal{T}}(Z;\dot{\mathcal{M}}\mathbf{V})$ and its current $J_t^{\mathcal{T}}(Z;\dot{\mathcal{M}}\mathbf{V})$.
\end{itemize}
Here $\partial$ denotes the plaquette holonomy.

\begin{itemize}
\item[(a)] \emph{Tower normalisation.} Let $i\le k+1$ be a level and let $\varphi$ be a schedule at level $i$. Then $\mathsf{D}_i(X;\varphi)[\mathbf{V}]$, the \emph{tower-type predicate}, is the following statement. For every $(t,p,Z)\in\mathfrak{T}_i(X)$ and every $\mathcal{T}\in\mathcal{R}(t,p,Z)$, the tower quantities of $\mathbf{V}$ at $(\mathcal{T},t,Z)$ are defined. Moreover, for every $x\in P(p,Z)\cap\Omega_1$, with $n:=n_i(x)$, the two bounds \eqref{4.11:eq-tower-bounds-a} and \eqref{4.11:eq-tower-bounds-b} hold:
\begin{equation}\label{4.11:eq-tower-bounds-a}
\bigl|\partial U_t^{\mathcal{T}}(Z;\dot{\mathcal{M}}\mathbf{V})-1\bigr|
<\varphi(n)\,\alpha_{0,n}\,L^{-2n}\quad\text{at }x,
\end{equation}
\begin{equation}\label{4.11:eq-tower-bounds-b}
\bigl|J_t^{\mathcal{T}}(Z;\dot{\mathcal{M}}\mathbf{V})\bigr|
<\varphi(n)\,\alpha_{0,n}\,L^{\,n-\min(t,n)}\,L^{-3(n-t)}\quad\text{at }x.
\end{equation}
\item[(b)] \emph{One-power normalisation.} Let $j\ge1$ and $a_0>0$. Then $\mathsf{D}^{\Sigma}_j(X;a_0)[\mathbf{V}]$, the \emph{one-power predicate}, is the following statement. For every $(t,p,Z)\in\mathfrak{T}_j(X)$ and every $\mathcal{T}\in\mathcal{R}(t,p,Z)$, the tower quantities of $\mathbf{V}$ at $(\mathcal{T},t,Z)$ are defined. Moreover, at every $x\in P(p,Z)$ the bounds \eqref{4.11:eq-tower-bounds-1} hold:
\begin{equation}\label{4.11:eq-tower-bounds-1}
\begin{gathered}
\bigl|\partial U_t^{\mathcal{T}}(Z;\dot{\mathcal{M}}\mathbf{V})-1\bigr|
<a_0L^{-(j-p)}L^{-2p},\\
\bigl|J_t^{\mathcal{T}}(Z;\dot{\mathcal{M}}\mathbf{V})\bigr|
<a_0L^{-(j-p)}L^{2(t-p)}.
\end{gathered}
\end{equation}
\end{itemize}
The tower-type predicate $\mathsf{D}_i(X;\varphi)$ and the one-power predicate $\mathsf{D}^{\Sigma}_j(X;a_0)$ are together called the \emph{derived-slot predicates}.

The bounds \eqref{4.11:eq-tower-bounds-a} and \eqref{4.11:eq-tower-bounds-b} are \cite[(2.35), (2.37)]{B-CMP119}. In item (a) they are required for every admissible sub-region, that is, for every triple of $\mathfrak{T}_i(X)$, and not for $X$ alone.

With the tower index written $t$, the right members of \cite[(2.35), (2.37)]{B-CMP119} take the following form:
\begin{equation*}
\begin{gathered}
\bigl(1-\beta(1-2^{-(j-n)})\bigr)\,\alpha_{0,n}\,L^{-2t}\bigl(L^{n}L^{-t}\bigr)^{-2},\\
\varphi(n)\,\alpha_{0,n}\,L^{\,n-\min(t,n)}\bigl(L^{n}L^{-t}\bigr)^{-3}.
\end{gathered}
\end{equation*}
The factor $1-\beta(1-2^{-(j-n)})$ in the first of these stands where the right member of \eqref{4.11:eq-tower-bounds-a} has the schedule value $\varphi(n)$. The powers of $L$ simplify as follows:
\begin{equation*}
L^{-2t}\bigl(L^{n}L^{-t}\bigr)^{-2}=L^{-2t}L^{-2n+2t}=L^{-2n},
\qquad
\bigl(L^{n}L^{-t}\bigr)^{-3}=L^{-3(n-t)}.
\end{equation*}
Hence these are the right members of \eqref{4.11:eq-tower-bounds-a} and \eqref{4.11:eq-tower-bounds-b}.

The predicate of item (b) is the clause of the one-power hereditary class that bounds the tower quantities at $(p,Z)$, taken with constant $a_0L^{-(j-p)}$, in the half of that clause stated for the full configuration, here $\mathbf{V}$, and not in the half stated for its $G$-valued factor. To see the normalisation, recall that $\xi_p=L^{-p}$, so that
\begin{equation*}
(L^{t}\xi_p)^2=L^{2(t-p)}.
\end{equation*}
The two bounds in \eqref{4.11:eq-tower-bounds-1} therefore read as follows:
\begin{equation*}
\begin{gathered}
\bigl|\partial U_t^{\mathcal{T}}(Z;\dot{\mathcal{M}}\mathbf{V})-1\bigr|
<a_0L^{-(j-p)}\,\xi_p^{2},\\
\bigl|J_t^{\mathcal{T}}(Z;\dot{\mathcal{M}}\mathbf{V})\bigr|
<a_0L^{-(j-p)}\,(L^{t}\xi_p)^{2}.
\end{gathered}
\end{equation*}
\end{definition}

\begin{definition}[Configuration classes and the clause on the $G$-valued factor]
\label{4.11:def-configuration-classes}
Throughout, $G$ is the gauge group, $G^{c}$ its complexification, $\mathfrak{g}^{c}$ the complexified Lie
algebra, and
$\xi_j=L^{-j}$ is the fine spacing in level-$j$ units. For a configuration $V$, $\partial V(q)$ is its
holonomy around the plaquette $q$, and $\nabla_U$ is the covariant derivative in the background $U$.
The radii $\alpha_{0,n},\alpha_{1,n}$ and the schedules $\varphi$ (among them $\varphi^{\mathrm{st}}_i$
and $\varphi^{\mathrm{post}}$) are those of Definition~\ref{4.11:def-radii-schedules}. The tower regions
$\Omega_n$, the layers $\mathfrak{L}_n=\Omega_n\setminus\Omega_{n+1}$ and the layer index $n_i(x)$ of a
point $x$ at reading level $i$ are those fixed with the tower rules in this section. The derived-slot
predicate of tower type $\mathsf{D}_i(X;\varphi)$ is that of Definition~\ref{4.11:def-tower-bounds}.
A configuration class is the union of the orbits of its
representatives under $G^{c}$-valued gauge transformations on $X$. This holds for each of the classes
below. In the exponentials below, $\mathrm{i}$ is the imaginary unit.

\smallskip
\emph{(a) The one-power class, recalled from the definition of the one-power hereditary class, after
\cite{B-CMP109}.}
Let $a_0,a_1,\gamma_0>0$. A \emph{$\Sigma$-representative} at $(j,X;a_0,a_1,\gamma_0)$ is a triple
$(U',U,\mathbf{J})$ on $X$ satisfying the conditions below. In them, $A'$ denotes a $\mathfrak{g}^{c}$-valued
bond potential on $X$, and $\mathsf{D}^{\Sigma}_j(X;a_0)$ is the derived-slot predicate of one-power type,
fixed in the definition of the one-power hereditary class.
\begin{equation}\label{4.11:eq-configuration-classes-a}
\begin{aligned}
&\text{(R1)}\quad U \text{ is } G\text{-valued},\quad |\partial U-1|<a_0\xi_j^2 \text{ on } X,
 \quad U \text{ satisfies the gauge condition \cite[(1.12)]{B-CMP109}};\\
&\text{(R2)}\quad U'=\exp(\mathrm{i}\,\xi_jA'),\qquad |A'|<a_1,\qquad |\nabla_U A'|<a_1;\\
&\text{(R3)}\quad |\partial(U'U)-1|<a_0\xi_j^2,\qquad |\mathbf{J}|<\gamma_0;\\
&\text{(R4)}^{\Sigma}\quad \mathsf{D}^{\Sigma}_j(X;a_0)[U'U]\quad\text{and}\quad
 \mathsf{D}^{\Sigma}_j(X;a_0)[U].
\end{aligned}
\end{equation}
The second half of (R4)$^{\Sigma}$ is the last clause of \cite[(1.16)]{B-CMP109}, ``and the same bounds
hold for $U$ instead of $\mathbf{U}$''. The one-power class $U^{\Sigma}_j(X;a_0,a_1,\gamma_0)$ is the
union of the orbits $[(U'U,\mathbf{J})]$ of the $\Sigma$-representatives.

\smallskip
\emph{(b) The classes of \cite{B-CMP119}.}
Let $\varphi$ be a schedule. A \emph{representative} at $(i,X;\varphi)$ is a triple $(U',U,\mathbf{J})$ on
$X$ with $U'=\exp(\mathrm{i}\,\mathsf{A}')$, where $\mathsf{A}'$ is a $\mathfrak{g}^{c}$-valued bond
potential. It must
satisfy the conditions below at every $x\in X\cap\Omega_1$, with $n=n_i(x)$.
\begin{equation}\label{4.11:eq-configuration-classes-b}
\begin{aligned}
&|\partial U-1|,\ |\partial(U'U)-1|<\varphi(n)\,\alpha_{0,n}\,L^{-2n}
 &&\text{\cite[(2.34)]{B-CMP119}};\\
&|\mathbf{J}|<\varphi(n)\,\alpha_{0,n}\,L^{3(i-n)}
 &&\text{\cite[(2.36)]{B-CMP119}};\\
&U \text{ satisfies the cube gauge condition}
 &&\text{\cite[(2.38)]{B-CMP119}};\\
&L^{n}|\mathsf{A}'|<\varphi(n)\,\alpha_{1,n},\qquad L^{2n}|\nabla_U\mathsf{A}'|<\varphi(n)\,\alpha_{1,n}
 &&\text{\cite[(2.39)]{B-CMP119}}.
\end{aligned}
\end{equation}
The current $\mathbf{J}$ is measured in the level-$i$ normalisation. The potential $\mathsf{A}'$ and its
covariant derivative are measured in fine units. In addition, the conditions of
\cite[(2.35), (2.37)]{B-CMP119} are required, read through the derived-slot predicate of tower type, in
which reading they take the form
\[
\mathsf{D}_i(X;\varphi)[U'U],
\]
where $\mathsf{D}_i(X;\varphi)$ is the derived-slot predicate of tower type. The class
$\tilde{U}_i(X;\varphi)$ is the union of the orbits of the representatives at $(i,X;\varphi)$.

We abbreviate $\tilde{U}^{c}_i(X):=\tilde{U}_i(X;\varphi^{\mathrm{st}}_i)$; this is the space written
$\tilde{U}^{c}_i(X,\tilde\alpha_0,\tilde\alpha_1)$ in \cite[(2.34)--(2.39)]{B-CMP119}, with
$\tilde\alpha_0,\tilde\alpha_1$ Ba{\l}aban's radius data.
The post-step class $\mathfrak{U}^{\mathrm{post}}_{k+1}(Y)$ of
Definition~\ref{4.4:def-post-step-space} is the class
\[
\mathfrak{U}^{\mathrm{post}}_{k+1}(Y)=\tilde{U}_{k+1}(Y;\varphi^{\mathrm{post}}).
\]

\smallskip
\emph{(c) The clause on the $G$-valued factor.}
The \emph{$U$-clause} at $(i,X;\varphi)$ is the condition
\[
\mathsf{D}_i(X;\varphi)[U]
\]
on the $G$-valued factor $U$ of a representative $(U',U,\mathbf{J})$. At $(k+1,Y;\varphi^{\mathrm{post}})$
it reads $\mathsf{D}_{k+1}(Y;\varphi^{\mathrm{post}})[U]$; in this form it is the condition on the
$G$-valued factor at the post-step level. The \emph{$U$-augmented classes} are defined as
follows. First, $\tilde{U}^{U}_i(X;\varphi)$ is the union of the orbits of those representatives at
$(i,X;\varphi)$ that satisfy, in addition, the $U$-clause. Then
\[
\mathfrak{U}^{\mathrm{post},U}_{k+1}(Y):=\tilde{U}^{U}_{k+1}(Y;\varphi^{\mathrm{post}}),
\qquad
\tilde{U}^{c,U}_i(X):=\tilde{U}^{U}_i(X;\varphi^{\mathrm{st}}_i).
\]
The configuration class of \cite[(1.12), (1.16)]{B-CMP109} is used as defined there and is not one of the
classes fixed in this definition; its last clause is the sentence quoted after
\eqref{4.11:eq-configuration-classes-a}. In that class and in the one-power class $U^{\Sigma}_j$, the
corresponding clause on the $G$-valued factor is already part of the definition: in the one-power class
it is $\mathsf{D}^{\Sigma}_j(X;a_0)[U]$, the second half of
(R4)$^{\Sigma}$ in \eqref{4.11:eq-configuration-classes-a}.

\smallskip
\emph{(d) Straddling sources.}
A region $Y$ is a \emph{straddling source} if $Y\cap\Omega_{k+1}\neq\emptyset$ and
$Y\cap\mathfrak{L}_n\neq\emptyset$ for some $n\le k$. At a straddling source $Y$, \cite{B-CMP119} states
that the new boundary terms $\mathbf{B}^{(k+1)}(Y)$ are analytic on the space
\cite[(3.43)]{B-CMP119}. That space is $\mathfrak{U}^{\mathrm{post}}_{k+1}(Y)$; its representatives satisfy
the predicate $\mathsf{D}_{k+1}(Y;\varphi^{\mathrm{post}})$ on the product $U'U$ and are not subject to
the $U$-clause.
\end{definition}

\begin{definition}[The maps of the step: restriction, twist, substitution, composites]
\label{4.11:def-step-maps}
The \emph{maps of the step} are the maps by which a point of one of the source classes of the
step becomes an argument of an older term. Let $X\subset Y$, and let $(\mathbf{U},\mathbf{J})$ be
a complex pair on $Y$: $\mathbf{U}$ is $G^{c}$-valued and $\mathbf{J}$ is
$\mathfrak{g}^{c}$-valued. The maps of the step are the following four.
\begin{enumerate}
\item[(a)] The \emph{restriction} to $X$, which sends $(\mathbf{U},\mathbf{J})$ to
$(\mathbf{U},\mathbf{J})|_X$.
\item[(b)] The \emph{twist}, which is the map
\begin{equation}\label{4.11:eq-step-maps-1}
(\mathbf{U},\mathbf{J})\longmapsto(\mathbf{U}^{\flat},\mathbf{J}^{\flat})|_X,
\qquad
\mathbf{U}^{\flat}:=e^{i\mathsf{H}}\,\mathbf{U},
\qquad
\mathbf{J}^{\flat}:=\mathbf{J}+\mathfrak{J}.
\end{equation}
Here $\varpi$ denotes the fluctuation data of the step, and $\mathsf{H}$ is the response
$\mathsf{H}=\mathsf{H}(\mathbf{U},\mathbf{J};\varpi)$ to the fluctuation data $\varpi$. This
response is the one of \cite[(179)--(180)]{B-CMP102b} and \cite[(3.10)--(3.13)]{B-CMP109}.
In \cite{B-CMP119} it is introduced as in \cite[(3.13)]{B-CMP109}. The term $\mathfrak{J}$ is
the shift of the current that accompanies $\mathsf{H}$. The twisted pair is taken together
with its \emph{witness}
\begin{equation}\label{4.11:eq-step-maps-2}
\bigl(e^{i\mathsf{H}}U',\,U\bigr).
\end{equation}
In the witness, $U$ is the $G$-valued factor of the configuration classes of
Definition~\ref{4.11:def-configuration-classes}. The factor $U'=\exp(i\mathsf{A}')$ is the
twist factor, and $\mathsf{A}'$ is a $\mathfrak{g}^{c}$-valued bond potential.
\item[(c)] The \emph{substitution} of \cite[(1.15)]{B-CMP109}. Fix a tower rule
$\mathcal{T}$, a level $t$ and a region $Z$. Let $\dot{\mathcal{M}}\mathbf{U}$ denote the
data of $\mathbf{U}$, that is, its fixed bonds and its block averages. The substitution is
the map
\begin{equation}\label{4.11:eq-step-maps-3}
\Phi_{\mathrm{sub}}(\mathbf{U},\mathbf{J})
:=\Bigl(U_t^{\mathcal{T}}\bigl(Z;\dot{\mathcal{M}}\mathbf{U}\bigr),\,
J_t^{\mathcal{T}}\bigl(Z;\dot{\mathcal{M}}\mathbf{U}\bigr)\Bigr).
\end{equation}
Here $U_t^{\mathcal{T}}(Z;\dot{\mathcal{M}}\mathbf{U})$ is the constrained minimiser of the
tower of rule $\mathcal{T}$ on $Z$ with data $\dot{\mathcal{M}}\mathbf{U}$, and
$J_t^{\mathcal{T}}(Z;\dot{\mathcal{M}}\mathbf{U})$ is its current, both in the normalisation
of level $t$.
\item[(d)] The \emph{composites} of the chain of operations $\mathbf{R}$ of
\cite{B-CMP122b}.
\end{enumerate}
\end{definition}

\begin{definition}[Floors on $L$ and $M$ and ceilings on the coupling]\label{4.11:not-floors-ceilings}
Let $\beta\in(0,1)$ be the schedule parameter and $\alpha_{0,n},\alpha_{1,n}$ Ba{\l}aban's radii at
level $n$ \cite[(2.28)]{B-CMP119}, as in Definition~\ref{4.11:def-radii-schedules}; by that
definition,
\begin{equation*}
\alpha_{0,n}=C_0\,g_n(\log g_n^{-2})^{q_0},\qquad \alpha_{1,n}=C_1\,g_n(\log g_n^{-2})^{q_1},
\end{equation*}
with constants $C_0,C_1$ and exponents $q_0,q_1$. Let
$\varrho=\exp(q_0c_5\gamma^2)$ be the constant of the lemma of this section on the comparability
of the radii along the coupling flow. In what follows the side conditions below are
in force. Each of them is one-sided. None of them involves the level $k$, the step index $j$, the
number of layers, or a lower bound on the reference coupling $g$.

\emph{Floors.} The blocking factor and the cube side satisfy
\begin{align}
L^2&\ge\sqrt{2}\,(1+\beta)\,\varrho,\label{4.11:eq-floors-ceilings-a}\\
M&\ge 2\bigl(R_1+\sqrt{d}\,L\bigr)M_1.\label{4.11:eq-floors-ceilings-b}
\end{align}
The floor \eqref{4.11:eq-floors-ceilings-a} involves only $\beta$ and $\varrho$. The floor
\eqref{4.11:eq-floors-ceilings-b} involves only $R_1$, $M_1$, $L$ and $d$; here $R_1$ is a
constant of the tower construction, the one place where this floor is used, and it is fixed
before $M$.

\emph{Ceilings.} For every level $n$,
\begin{align}
2(1+\beta)\,\alpha_{0,n}&\le\min\{a_1(d,L),\,c_2(d,L)\},\label{4.11:eq-floors-ceilings-c}\\
c_Q(1+\beta)\,\alpha_{1,n}&\le b_H,\label{4.11:eq-floors-ceilings-d}
\end{align}
where $a_1(d,L)$ and $c_2(d,L)$ are Ba{\l}aban's thresholds, which depend only on $d$ and
$L$. They are always written with their arguments: $c_2(d,L)$ is not the constant $c_2$ of the
glossary (the one with $c_5=2c_2$), and $a_1(d,L)$ is not the parameter $a_1$ of the one-power
class. In \eqref{4.11:eq-floors-ceilings-d}, $b_H$ is the radius of the polydisc on
which the analytic response is controlled, and $c_Q$ is the constant of the perturbative route
to the tower quantities, in which the response on that polydisc is used. For
the lemma on the comparability of the radii we require moreover
\begin{equation*}
B_\sharp\gamma^2\le 1,\qquad \gamma\le e^{-1/2}.
\end{equation*}

The ceiling \eqref{4.11:eq-floors-ceilings-c} holds for all $n$ as soon as
$g_n\le\gamma\le e^{-q_0}$ for all $n$ and
\begin{equation}\label{4.11:eq-floors-ceilings-1}
2(1+\beta)\,C_0\,\gamma\,(\log\gamma^{-2})^{q_0}\le\min\{a_1(d,L),\,c_2(d,L)\}.
\end{equation}
It is therefore a ceiling on $\gamma$, chosen after $C_0$, $q_0$ and $L$, of the kind imposed
in \cite{B-CMP109}.

The ceiling \eqref{4.11:eq-floors-ceilings-d} holds for all $n$ as soon as
$g_n\le\gamma\le e^{-q_1}$ for all $n$ and
\begin{equation}\label{4.11:eq-floors-ceilings-2}
c_Q(1+\beta)\,C_1\,\gamma\,(\log\gamma^{-2})^{q_1}\le b_H.
\end{equation}
It is therefore a ceiling on $\gamma$, chosen after $C_1$, $q_1$, $L$ and $M$, of the same
kind. The parameter $M$ enters because in \cite{B-CMP109} the radius $\alpha_1$ depends on $M$.
\end{definition}

\begin{proof}
We verify that \eqref{4.11:eq-floors-ceilings-1} together with $g_n\le\gamma\le e^{-q_0}$
implies \eqref{4.11:eq-floors-ceilings-c}, and that \eqref{4.11:eq-floors-ceilings-2} together
with $g_n\le\gamma\le e^{-q_1}$ implies \eqref{4.11:eq-floors-ceilings-d}.
Fix an exponent $\mu>0$ and put $h_\mu(t)=t(\log t^{-2})^\mu$ for $t\in(0,e^{-\mu}]$. On this
interval $\log t^{-2}=-2\log t\ge 2\mu>0$. Differentiating,
\begin{align*}
h_\mu'(t)&=(\log t^{-2})^\mu+t\,\mu\,(\log t^{-2})^{\mu-1}\cdot\bigl(-2t^{-1}\bigr)\\
&=(\log t^{-2})^{\mu-1}\bigl(\log t^{-2}-2\mu\bigr)\ \ge\ 0 .
\end{align*}
Hence $h_\mu$ is non-decreasing on $(0,e^{-\mu}]$.

By the formulas for the radii recalled in the definition,
$\alpha_{0,n}=C_0\,h_{q_0}(g_n)$ and $\alpha_{1,n}=C_1\,h_{q_1}(g_n)$.
Suppose $g_n\le\gamma\le e^{-q_0}$. Then $g_n$ and $\gamma$ both lie in $(0,e^{-q_0}]$, so
$h_{q_0}(g_n)\le h_{q_0}(\gamma)$. Multiplying by $2(1+\beta)C_0$ and using
\eqref{4.11:eq-floors-ceilings-1} gives
\begin{equation*}
2(1+\beta)\alpha_{0,n}\le 2(1+\beta)C_0\,h_{q_0}(\gamma)\le\min\{a_1(d,L),\,c_2(d,L)\},
\end{equation*}
which is \eqref{4.11:eq-floors-ceilings-c}.

In the same way, $g_n\le\gamma\le e^{-q_1}$ gives $h_{q_1}(g_n)\le h_{q_1}(\gamma)$.
Multiplying by $c_Q(1+\beta)C_1$ and using \eqref{4.11:eq-floors-ceilings-2} yields
\eqref{4.11:eq-floors-ceilings-d}.
\end{proof}

\begin{proposition}[The constrained-minimiser clause at straddling sources. Stated; proof
incomplete at the step marked $(\ast)$]
\label{4.11:prop-straddling-clause}
Let $k\ge 0$. Call a region $Y$ at level $k+1$ \emph{straddling} if
$Y\cap\Omega_{k+1}\neq\emptyset$ and $Y\cap\mathfrak{L}_n\neq\emptyset$ for some $n\le k$.
For every straddling region $Y$ and every representative $(U',U,\mathbf{J})$ of the
post-step class $\mathfrak{U}^{\mathrm{post}}_{k+1}(Y)$ of
Definition~\ref{4.11:def-configuration-classes}, the $G$-valued factor $U$ satisfies
the clause
\begin{equation*}
\mathsf{D}_{k+1}(Y;\varphi^{\mathrm{post}})[U]
\end{equation*}
of Definition~\ref{4.11:def-configuration-classes}, with the post schedule
$\varphi^{\mathrm{post}}$.
\end{proposition}

\begin{proof}
The clause is the last condition of \cite[(1.16)]{B-CMP109}, taken at level $k+1$ with
the post schedule. It is the condition on the $G$-valued factor that enters the
analyticity of the new boundary terms $\mathbf{B}^{(k+1)}(Y)$ at a straddling region
on the post-step space, which is asserted at \cite[p.~261]{B-CMP119}.

Fix a straddling region $Y$ and a representative $(U',U,\mathbf{J})$ of
$\mathfrak{U}^{\mathrm{post}}_{k+1}(Y)$. What has to be shown is the following step.

$(\ast)$ The $G$-valued factor $U$ of this representative satisfies
$\mathsf{D}_{k+1}(Y;\varphi^{\mathrm{post}})[U]$.

Step $(\ast)$ is the entire assertion of the proposition.
\end{proof}

\begin{lemma}[Comparability of the radii along the coupling flow]\label{4.11:lem-radii-comparable}
Let $k\ge 0$. For $1\le i\le k+1$ write $x_i=g_i^{-2}$, and let $\alpha_{0,i}$ be Ba{\l}aban's radius
at level $i$ (Definition~\ref{4.11:def-radii-schedules}; \cite[(2.28)]{B-CMP119}). Let $q_0$
denote the exponent of $\log x_i$ in that radius. Assume the following hypotheses.
\begin{itemize}
\item[(a)] The two-sided slope of clause~(0) (Theorem~\ref{4.1:thm-clause-zero}) holds at all levels
$\le k+1$: for all $1\le j\le n\le k+1$,
\begin{equation*}
x_j-x_n\le B_\sharp\,(n-j),\qquad x_n\le x_j+c_5 .
\end{equation*}
\item[(b)] $x_i\ge\gamma^{-2}$ for $1\le i\le k+1$.
\item[(c)] $B_\sharp\gamma^2\le 1$ and $\gamma\le e^{-1/2}$.
\end{itemize}
Put $C_{\mathrm{rad}}:=\exp(q_0c_5\gamma^2)$, so that $C_{\mathrm{rad}}\le e^{q_0c_5}$. Then, for all
$1\le j\le n\le k+1$,
\begin{equation}\label{4.11:eq-radii-comparable-a}
\frac{\alpha_{0,n}}{\alpha_{0,j}}\le C_{\mathrm{rad}}\,(1+n-j)^{1/2}.
\end{equation}
The constant $C_{\mathrm{rad}}$ does not depend on $j$, $n$ or $k$.
\end{lemma}

\begin{proof}
Fix $1\le j\le n\le k+1$.

\emph{The coupling factor.} By the first inequality of (a), $x_j\le x_n+B_\sharp(n-j)$. Dividing
by $x_n>0$ and using $g_n^2/g_j^2=x_j/x_n$, we obtain
\begin{equation*}
\frac{g_n^2}{g_j^2}=\frac{x_j}{x_n}\le 1+\frac{B_\sharp(n-j)}{x_n}.
\end{equation*}
By (b), $1/x_n\le\gamma^2$. By (c), $B_\sharp\gamma^2\le1$. Hence
\begin{equation*}
\frac{g_n^2}{g_j^2}\le 1+B_\sharp\gamma^2(n-j)\le 1+(n-j).
\end{equation*}
Taking square roots gives
\begin{equation}\label{4.11:eq-radii-comparable-1}
\frac{g_n}{g_j}\le(1+n-j)^{1/2}.
\end{equation}

\emph{The logarithmic factor.} By (b) and (c), $x_j\ge\gamma^{-2}\ge e$, so
$\log x_j\ge\log\gamma^{-2}\ge1$; likewise $\log x_n\ge1$. The second inequality of (a) reads
$x_n\le x_j(1+c_5/x_j)$. Taking logarithms and using $\log(1+u)\le u$ for $u\ge0$, together with
$1/x_j\le\gamma^2$ from (b), we obtain
\begin{equation*}
\log x_n\le\log x_j+\log\Bigl(1+\frac{c_5}{x_j}\Bigr)\le\log x_j+\frac{c_5}{x_j}
\le\log x_j+c_5\gamma^2 .
\end{equation*}
Divide by $\log x_j>0$ and raise both sides to the power $q_0$. Then use $1+u\le e^u$ and
$\log x_j\ge\log\gamma^{-2}\ge1$. This gives
\begin{equation}\label{4.11:eq-radii-comparable-2}
\Bigl(\frac{\log x_n}{\log x_j}\Bigr)^{q_0}
\le\Bigl(1+\frac{c_5\gamma^2}{\log x_j}\Bigr)^{q_0}
\le\exp\Bigl(\frac{q_0c_5\gamma^2}{\log\gamma^{-2}}\Bigr)
\le\exp(q_0c_5\gamma^2).
\end{equation}
The last inequality uses $\log\gamma^{-2}\ge1$, which holds because $\gamma\le e^{-1/2}$.

\emph{Conclusion.} By the definition of the radius (Definition~\ref{4.11:def-radii-schedules}), the
ratio of the radii factorises as
\begin{equation*}
\frac{\alpha_{0,n}}{\alpha_{0,j}}=\frac{g_n}{g_j}\Bigl(\frac{\log x_n}{\log x_j}\Bigr)^{q_0}.
\end{equation*}
Inserting \eqref{4.11:eq-radii-comparable-1} and \eqref{4.11:eq-radii-comparable-2} yields
\eqref{4.11:eq-radii-comparable-a} with $C_{\mathrm{rad}}=\exp(q_0c_5\gamma^2)$. This constant involves
only $q_0$, $c_5$ and $\gamma$, so it does not depend on $j$, $n$ or $k$. Finally,
$\gamma^2\le e^{-1}\le1$ by (c), so $C_{\mathrm{rad}}\le e^{q_0c_5}$.
\end{proof}

\begin{corollary}[The level-transfer inequality for schedule factors]\label{4.11:cor-level-transfer}
Assume the comparability of the radii along the coupling flow,
\eqref{4.11:eq-radii-comparable-a}, with its constant $\varrho$, and the floor on $L$,
\eqref{4.11:eq-floors-ceilings-a}. Then for every integer $m\ge 1$
\begin{equation}\label{4.11:eq-level-transfer-a}
(1+\beta)\,\varrho\,(1+m)^{1/2}L^{-2m}
\le \frac{1}{L^{2}}\,(1+\beta)\,\varrho\,\sqrt{2}\,
\Bigl(\frac{1+m}{2}\Bigr)^{1/2}L^{-2(m-1)}
\le 1 .
\end{equation}
Moreover, let $\beta,\theta\in(0,1)$ be the schedule parameters, let $\alpha_{0,n}$ be the
radius at level $n$ of \cite[(2.28)]{B-CMP119}, and let $\varphi$ be any of the factors of the
standard and post schedules: $f_{i,n}$, $F_n$ with $n\le k$, or $F_{\mathrm{top}}$. Then
$\varphi\le 1+\beta$, and
\begin{equation}\label{4.11:eq-level-transfer-b}
\varphi\cdot\frac{\alpha_{0,n}}{\alpha_{0,j}}\cdot L^{-2(n-j)}\le 1
\qquad(1\le j<n\le k+1).
\end{equation}
\end{corollary}

\begin{proof}
The first inequality in \eqref{4.11:eq-level-transfer-a} holds with equality, because
$(1+m)^{1/2}=\sqrt{2}\,((1+m)/2)^{1/2}$ and $L^{-2m}=L^{-2}L^{-2(m-1)}$.

For the second inequality we use that the blocking factor satisfies $L\ge 2$. First,
$(1+m)/2\le 2^{m-1}$ for every integer $m\ge1$. For $m=1$ both sides equal $1$. If the
inequality holds for $m$, then
\begin{equation*}
2^{m}=2\cdot 2^{m-1}\ge 1+m\ge \frac{2+m}{2}.
\end{equation*}
Next, $L\ge 2$ and $m-1\ge 0$ give
\begin{equation*}
2^{m-1}\le 2^{4(m-1)}\le L^{4(m-1)} .
\end{equation*}
Hence $(1+m)/2\le L^{4(m-1)}$. Taking square roots of these positive numbers yields
\begin{equation}\label{4.11:eq-level-transfer-1}
\Bigl(\frac{1+m}{2}\Bigr)^{1/2}L^{-2(m-1)}\le 1 .
\end{equation}
The remaining factor satisfies $(1+\beta)\varrho\sqrt{2}\,L^{-2}\le 1$: this is the floor
\eqref{4.11:eq-floors-ceilings-a}, $L^{2}\ge\sqrt{2}\,(1+\beta)\varrho$, divided by $L^{2}$.
The middle member of \eqref{4.11:eq-level-transfer-a} is the product of these two
non-negative quantities, each at most $1$. It is therefore at most $1$.

We now prove the second assertion. Each of the factors $\varphi$ listed there is at
most $1+\beta$. Indeed, $f_{i,n}\le 1$ and $F_{\mathrm{top}}=1+\beta$. For $F_n$, the bound
$2^{-(k+1-n)}\le 1$ gives
\begin{equation*}
F_n=1-\beta+\beta(1+\theta)2^{-(k+1-n)}\le 1+\beta\theta\le 1+\beta .
\end{equation*}
Let $1\le j<n\le k+1$ and put $m=n-j\ge 1$. We apply \eqref{4.11:eq-radii-comparable-a} to
the pair $j<n$, in the form that bounds $\alpha_{0,n}/\alpha_{0,j}$ by
$\varrho\,(1+(n-j))^{1/2}$. Combining this with $\varphi\le 1+\beta$ and with
\eqref{4.11:eq-level-transfer-a} gives
\begin{equation}\label{4.11:eq-level-transfer-2}
\varphi\,\frac{\alpha_{0,n}}{\alpha_{0,j}}\,L^{-2(n-j)}
\le (1+\beta)\,\varrho\,(1+m)^{1/2}L^{-2m}\le 1 ,
\end{equation}
which is \eqref{4.11:eq-level-transfer-b}.
\end{proof}

\begin{lemma}[Real points: augmented and plain classes agree]\label{4.11:lem-real-points}
In the setting of Definition~\ref{4.11:def-configuration-classes}, let $U$ be a $G$-valued
configuration on $X$ and let $\mathbf{J}$ be a real current on $X$.
Consider the real point $(U,\mathbf{J})$ presented with the trivial witness $(1,U)$, that is,
the triple $(1,U,\mathbf{J})$. This triple is a representative of the augmented class
$\tilde{U}^{U}_i(X;\varphi)$ if and only if it is a representative of the class
$\tilde{U}_i(X;\varphi)$. Likewise, in the class of \cite{B-CMP109} fixed in
Definition~\ref{4.11:def-configuration-classes} and in the one-power class
$U^{\Sigma}_j(X;a_0,a_1,\gamma_0)$ defined by \eqref{4.11:eq-configuration-classes-a}, the triple
$(1,U,\mathbf{J})$ satisfies the defining conditions in full if and only if it satisfies them
with the clause on the $G$-valued factor omitted.
\end{lemma}

\begin{proof}
We use the classes and the clause on the $G$-valued factor as fixed in
Definition~\ref{4.11:def-configuration-classes}. For the witness $(1,U)$ the twist factor is
$U'=1$, so that its potential is $\mathsf{A}'=0$, and the configuration presented is
\begin{equation*}
\mathbf{U}:=U'U=U .
\end{equation*}

We list the conditions that membership of $(1,U,\mathbf{J})$ in $\tilde{U}_i(X;\varphi)$ imposes
for this witness; they are the conditions \eqref{4.11:eq-configuration-classes-b}, together with
the derived-slot predicate of tower type $\mathsf{D}_i(X;\varphi)$.
\begin{itemize}
\item The plaquette condition is required for $U$ and for $U'U$. Since $U'U=U$, it is the same
inequality, at every point $x\in X\cap\Omega_1$, written twice.
\item The current bound is required for $\mathbf{J}$.
\item The cube gauge condition is required for $U$.
\item The bound on the twist potential is required for $\mathsf{A}'$. Since $\mathsf{A}'=0$,
both $\mathsf{A}'$ and $\nabla_U\mathsf{A}'$ vanish identically, and the condition is void.
\item The derived-slot predicate is required for the configuration presented, and
$\mathsf{D}_i(X;\varphi)[\mathbf{U}]=\mathsf{D}_i(X;\varphi)[U]$ because $\mathbf{U}=U$.
\end{itemize}

The augmented class $\tilde{U}^{U}_i(X;\varphi)$ imposes, in addition to these conditions, the
clause on the $G$-valued factor. For the witness $(1,U)$ the $G$-valued factor is $U$, so this
clause is $\mathsf{D}_i(X;\varphi)[U]$. This is precisely the last item of the list above.

If $(1,U,\mathbf{J})$ is a representative of $\tilde{U}^{U}_i(X;\varphi)$, it satisfies in
particular all the conditions of $\tilde{U}_i(X;\varphi)$, so it is a representative of
$\tilde{U}_i(X;\varphi)$. Conversely, if it is a representative of $\tilde{U}_i(X;\varphi)$, it
satisfies $\mathsf{D}_i(X;\varphi)[U]$, which is the additional clause. Hence it is a
representative of $\tilde{U}^{U}_i(X;\varphi)$.

In the class of \cite{B-CMP109} and in the one-power class $U^{\Sigma}_j(X;a_0,a_1,\gamma_0)$,
the clause on the $G$-valued factor is part of the definition. In each it is the second half of
the last defining condition, which requires the derived-slot predicate of that class both for
$U'U$ and for the $G$-valued factor $U$. For the one-power class this is the last condition in
\eqref{4.11:eq-configuration-classes-a}, which requires
\begin{equation*}
\mathsf{D}^{\Sigma}_j(X;a_0)[U'U]\quad\text{and}\quad\mathsf{D}^{\Sigma}_j(X;a_0)[U] .
\end{equation*}
At $U'=1$ we have $U'U=U$, so in each class the two halves of this condition coincide. Requiring
both is therefore the same as requiring the first alone, and this proves the equivalence for the
class of \cite{B-CMP109} and for the one-power class.
\end{proof}

\begin{remark}
The lemma concerns the trivial witness only. Suppose instead that a $G$-valued configuration
$\mathbf{U}$ is presented with a non-trivial witness $(U',U_1)$, so that
\begin{equation*}
\mathbf{U}=U'U_1,\qquad U'=e^{i\mathsf{A}'}\neq 1 .
\end{equation*}
In this case membership in $\tilde{U}_i(X;\varphi)$ via $(U',U_1)$ is not shown to give the
clause on the $G$-valued factor for $U_1$. Nor does it, in general, give the cube gauge
condition of \eqref{4.11:eq-configuration-classes-b} for the trivial witness $(1,\mathbf{U})$.

To see this, write $u$ for the gauge transformation of the cube gauge condition for $U_1$, and
$A$ for the potential of $U_1$ in that gauge, so that
\begin{equation*}
\mathbf{U}^u=e^{i\operatorname{Ad}_u\mathsf{A}'}e^{i\xi A},
\end{equation*}
where $\xi$ is the lattice spacing. Write $B$ and $C$ for the constants in the bound of the cube
gauge condition. At a point $x\in X$, with $n=n_i(x)$, the potential of $\mathbf{U}^u$ has size
up to
\begin{equation*}
BCM\alpha_{0,n}L^{-n}+(1+\beta)\alpha_{1,n}L^{-n} .
\end{equation*}
This exceeds the bound $BCM\alpha_{0,n}L^{-n}$ that the cube gauge condition imposes.

The real base points of the inductive construction are the background configurations $U_i$, whose
membership in the classes rests on their own regularity. They are presented with the trivial
witness; in particular, the real point of a representative of the one-power class is $(1,U,0)$.
For these base points Lemma~\ref{4.11:lem-real-points} applies.
\end{remark}

\begin{lemma}[Cauchy families over a fixed $G$-valued factor]\label{4.11:lem-cauchy-families}
Let the classes $\tilde{U}_i(X;\varphi)$ and $\tilde{U}^{U}_i(X;\varphi)$, their representatives with
witnesses, and the clause $\mathsf{D}_i(X;\varphi)[\,\cdot\,]$ on the $G$-valued factor be as in
Definition~\ref{4.11:def-configuration-classes}. Fix a $G$-valued configuration $U$, and let
\begin{equation*}
\mathcal{F}=\{(U'_\zeta U,\ \mathbf{J}_\zeta,\ \mathsf{a}_\zeta)\}_{\zeta}
\end{equation*}
be a non-empty family over $U$, indexed by a parameter $\zeta$; in the Cauchy estimates $\zeta$
ranges over a polydisc. Here $U'_\zeta$ is a twist
factor, $\mathbf{J}_\zeta$ is a current, and the third entry $\mathsf{a}_\zeta$ is an auxiliary
variable of an old term, in the role of the variable written $A$ in \cite[(3.13)]{B-CMP109}.
Each member of $\mathcal{F}$ is taken with witness $(U'_\zeta,U)$. Then the following are
equivalent:
\begin{enumerate}
\item[(a)] every member of $\mathcal{F}$ is a representative of $\tilde{U}^{U}_i(X;\varphi)$;
\item[(b)] every member of $\mathcal{F}$ is a representative of $\tilde{U}_i(X;\varphi)$, and $U$
satisfies the clause $\mathsf{D}_i(X;\varphi)[U]$.
\end{enumerate}
The condition on $U$ in \textup{(b)} is one condition, independent of $\zeta$.
\end{lemma}

\begin{proof}
Every member of $\mathcal{F}$ has a witness of the form $(U'_\zeta,U)$, and the second component
of the witness is the same configuration $U$ for every value of $\zeta$. By
Definition~\ref{4.11:def-configuration-classes}, the conditions defining a representative of
$\tilde{U}^{U}_i(X;\varphi)$ with a given witness are the conditions defining a representative of
$\tilde{U}_i(X;\varphi)$ together with the clause $\mathsf{D}_i(X;\varphi)$ evaluated at the
$G$-valued factor of the witness. For the member indexed by $\zeta$ that clause is
$\mathsf{D}_i(X;\varphi)[U]$, which involves only $U$ and does not contain $\zeta$.

Assume \textup{(a)}. For each $\zeta$ the member indexed by $\zeta$ satisfies all conditions
defining a representative of $\tilde{U}^{U}_i(X;\varphi)$, hence in particular those defining a
representative of $\tilde{U}_i(X;\varphi)$. Since $\mathcal{F}$ is not empty, there is at least one
member, and for it the clause on the $G$-valued factor of its witness reads
$\mathsf{D}_i(X;\varphi)[U]$; so $U$ satisfies this clause, and \textup{(b)} holds.

Assume \textup{(b)}. Fix $\zeta$. The member indexed by $\zeta$ satisfies the conditions defining
a representative of $\tilde{U}_i(X;\varphi)$. The clause on the $G$-valued factor of its witness is
$\mathsf{D}_i(X;\varphi)[U]$, which holds by assumption. By the description of
$\tilde{U}^{U}_i(X;\varphi)$ recalled above, this member is a representative of
$\tilde{U}^{U}_i(X;\varphi)$. Since $\zeta$ was arbitrary, \textup{(a)} holds.
\end{proof}

\begin{remark}
At the following places of \cite{B-CMP109} and \cite{B-CMP119}, the families over which Cauchy
estimates are taken are either of the form of Lemma~\ref{4.11:lem-cauchy-families} or of the form
\begin{equation*}
\{(\mathbf{U},\ \mathbf{J}+s\mathfrak{J},\ \mathsf{a}+t\mathsf{b}')\}_{s,t},
\end{equation*}
where $\mathfrak{J}$ is a fixed current shift and $\mathsf{b}'$ a fixed increment of the auxiliary
variable, so that only the current and the auxiliary variable move, affinely in $s,t$, and the first
entry does not vary with the parameters $s,t$; nothing is asserted about Cauchy families at other
places. Specifically:
\begin{itemize}
\item the function \cite[(3.13)]{B-CMP109} is a function of $(\mathbf{U},\mathbf{J},\mathsf{a})$,
analytic for $(\mathbf{U},\mathbf{J})$ in a class of \cite{B-CMP109} and for auxiliary variables
of modulus less than the radius $\alpha_2$ of \cite[\S3]{B-CMP109};
\item in \cite[\S3]{B-CMP109} a Cauchy integral is taken in parameters $t_\square$ attached to the
cubes $\square$; these parameters move the auxiliary variable;
\item the substitution $\mathbf{U}'=\exp(i\xi\,\mathsf{a})\mathbf{U}$ of \cite[(3.10)]{B-CMP109}
is a twist; it is carried in the twist factor, with the $G$-valued factor fixed;
\item in \cite{B-CMP119} the same device is applied to the new boundary terms at a straddling
region $Y$, with a reference back to \cite[(3.13)]{B-CMP109} (``as in (I.3.13)'' in the wording
of \cite{B-CMP119}, where [I] is \cite{B-CMP109}).
\end{itemize}
In the constructions used in this chapter, all Cauchy polydiscs move the twist potentials
$\mathsf{A}'$, the current $\mathbf{J}$ and the auxiliary variable over a fixed $G$-valued
configuration. The complex data of the composite maps of Definition~\ref{4.11:def-step-maps}
produce twists, and every composite leaves the $G$-valued factor of its argument unchanged; the
decoupling parameters of the step, which are distinct from the parameter $s$ above, act on
operators, by interpolation. None of
the parameters at the places listed deforms the $G$-valued factor.
Lemma~\ref{4.11:lem-cauchy-families} does not apply to a family in which the $G$-valued factor
varies with the parameter.
\end{remark}

\begin{lemma}[The transfer lemma: the $G$-valued factor on the real slice]\label{4.11:lem-transfer}
Let $\mathcal{C}^{U}_{\mathrm{src}}$ be a $U$-augmented class in the sense of
Definition~\ref{4.11:def-configuration-classes}. Let $\mathcal{C}_{\mathrm{tgt}}$ be a class of
either type, tower type or one-power type, and let $\mathsf{D}_{\mathrm{tgt}}[\cdot]$ be its
derived-slot predicate. Let $\mathcal{C}$ be a class of either type, with derived-slot predicate
$\mathsf{D}[\cdot]$. For a witness $(U',U,\mathbf{J})$ of a pair in $\mathcal{C}$:
\begin{itemize}
\item the condition $\mathsf{D}[U'U]$ is called its $\mathbf{U}$-half;
\item the condition $\mathsf{D}[U]$ is called its $U$-half.
\end{itemize}
Let $\Phi$ be a map on representatives of $\mathcal{C}^{U}_{\mathrm{src}}$, and write
$\Phi(U'U,\mathbf{J})$ for its value at the representative $(U',U,\mathbf{J})$. Let $\Psi$ be a
map on $G$-valued configurations. Assume the following two hypotheses:
\begin{equation}\label{4.11:eq-transfer-a}
\begin{aligned}
&\text{(T1)}\quad &&\text{for every representative $(U',U,\mathbf{J})$ of
$\mathcal{C}^{U}_{\mathrm{src}}$, the orbit of $\Phi(U'U,\mathbf{J})$}\\
&&&\text{has a witness $(U'_{\Phi},\Psi(U),\mathbf{J}_{\Phi})$ satisfying the pointwise}\\
&&&\text{conditions of $\mathcal{C}_{\mathrm{tgt}}$ and
$\mathsf{D}_{\mathrm{tgt}}[U'_{\Phi}\Psi(U)]$;}\\
&\text{(T2)}\quad &&\text{for a real representative $(1,U,0)$, that is, one with $U$
$G$-valued,}\\
&&&\text{twist factor $1$ and current $0$ (Lemma~\ref{4.11:lem-real-points}),}\\
&&&\text{the witness in (T1) has $U'_{\Phi}=1$.}
\end{aligned}
\end{equation}
Then, for every representative $(U',U,\mathbf{J})$ of $\mathcal{C}^{U}_{\mathrm{src}}$, the
witness of (T1) satisfies $\mathsf{D}_{\mathrm{tgt}}[\Psi(U)]$, the $U$-half of the target.
\end{lemma}

\begin{proof}
Let $(U',U,\mathbf{J})$ be a representative of $\mathcal{C}^{U}_{\mathrm{src}}$. Write
$\mathsf{D}_{\mathrm{src}}$ for the derived-slot predicate of $\mathcal{C}^{U}_{\mathrm{src}}$. It
is $\mathsf{D}^{\Sigma}_j(X;a_0)$ if the class is of one-power type and $\mathsf{D}_i(X;\varphi)$
if it is of tower type.

\emph{Step one: $(1,U,0)$ is a representative of $\mathcal{C}^{U}_{\mathrm{src}}$.} This follows
from the lemma on real points, Lemma~\ref{4.11:lem-real-points}, and from the argument of its
proof. We carry out that argument here. First we check the pointwise conditions at $(1,U,0)$,
deducing each from the corresponding condition at $(U',U,\mathbf{J})$.
\begin{itemize}
\item Condition (R1) of \eqref{4.11:eq-configuration-classes-a}, respectively the first
inequality of \cite[(2.34)]{B-CMP119} together with \cite[(2.38)]{B-CMP119}, concerns the
$G$-valued factor $U$ alone. It therefore holds at $(1,U,0)$ unchanged.
\item Condition (R2) of \eqref{4.11:eq-configuration-classes-a}, respectively
\cite[(2.39)]{B-CMP119}, is void at $\mathsf{A}'=0$. Indeed $1=\exp(i\cdot 0)$, and the potential
and its covariant derivative $\nabla_U\mathsf{A}'$ both vanish, so the required bounds hold.
\item At $(1,U,0)$ the full configuration is $1\cdot U=U$. Hence the first inequality of (R3) of
\eqref{4.11:eq-configuration-classes-a}, respectively the second inequality of
\cite[(2.34)]{B-CMP119}, evaluated at $\mathbf{U}=U$, is the first inequality of (R1) of
\eqref{4.11:eq-configuration-classes-a}, respectively the first inequality of
\cite[(2.34)]{B-CMP119}. That inequality holds at $(U',U,\mathbf{J})$ by hypothesis and involves
$U$ alone, so it holds at $(1,U,0)$.
\item The current vanishes, $|\mathbf{J}|=0$. This is below $\gamma_0$ in the one-power case, and
below the radius on the right-hand side of \cite[(2.36)]{B-CMP119} in the tower case.
\end{itemize}
Next we check the derived-slot conditions at $(1,U,0)$.
\begin{itemize}
\item The $\mathbf{U}$-half at $(1,U,0)$ is $\mathsf{D}_{\mathrm{src}}[1\cdot U]
=\mathsf{D}_{\mathrm{src}}[U]$. This is exactly the $U$-clause of $(U',U,\mathbf{J})$. It holds
because $\mathcal{C}^{U}_{\mathrm{src}}$ is $U$-augmented, so every representative satisfies the
$U$-clause.
\item The $U$-half at $(1,U,0)$ is again $\mathsf{D}_{\mathrm{src}}[U]$, so it holds for the same
reason.
\end{itemize}
Hence $(1,U,0)$ is a representative of $\mathcal{C}^{U}_{\mathrm{src}}$. It is real, since $U$ is
$G$-valued, the twist factor is $1$ and the current is $0$.

\emph{Step two: the $U$-half holds at the real point.} Apply (T1) to the representative
$(1,U,0)$. The value of $\Phi$ there is $\Phi(1\cdot U,0)=\Phi(U,0)$. Hypothesis (T1) of
\eqref{4.11:eq-transfer-a} gives a witness $(U'_{\Phi},\Psi(U),\mathbf{J}_{\Phi})$ for the orbit
of $\Phi(U,0)$ that satisfies $\mathsf{D}_{\mathrm{tgt}}[U'_{\Phi}\Psi(U)]$. Since $(1,U,0)$ is a
real representative, hypothesis (T2) of \eqref{4.11:eq-transfer-a} gives $U'_{\Phi}=1$ for this
witness. Substituting $U'_{\Phi}=1$ into $\mathsf{D}_{\mathrm{tgt}}[U'_{\Phi}\Psi(U)]$ yields
$\mathsf{D}_{\mathrm{tgt}}[\Psi(U)]$.

\emph{Step three: return to the original representative.} By (T1), the witness for
$\Phi(U'U,\mathbf{J})$ at the original representative $(U',U,\mathbf{J})$ has $G$-valued factor
$\Psi(U)$. This is the same configuration as the one obtained in step two. The $U$-half of the
target for this witness is, by definition, exactly $\mathsf{D}_{\mathrm{tgt}}[\Psi(U)]$. That
condition holds by step two.
\end{proof}

\begin{lemma}[Heredity of the clause under the twist]\label{4.11:lem-twist-heredity}
Let $Y$ be a region, let $1\le j\le k$, let $X\subset Y$, and let $(U',U,\mathbf{J})$ be a
representative of $\mathfrak{U}^{\mathrm{post},U}_{k+1}(Y)$. Let
$(\mathbf{U}^{\flat},\mathbf{J}^{\flat})$ be the twisted pair with witness
$(e^{i\mathsf{H}}U',U)$ (Definition~\ref{4.11:def-step-maps}).
Assume the following three statements.
\begin{itemize}
\item[(a)] The comparability of the radii along the coupling flow, that is, the inequalities
\eqref{4.11:eq-radii-comparable-a} of Lemma~\ref{4.11:lem-radii-comparable}, with the constant
$\varrho$ given there.
\item[(b)] Where the two sets of rules read at a triple differ: every rule read at a triple of
$\mathfrak{T}_j(X)$ in the $\mathsf{D}$-predicates at level $j$ on $X$ is also read at that
triple in the $\mathsf{D}$-predicate at level $k+1$ on $Y$.
\item[(c)] The floors on $L$ and $M$ and the ceilings on the coupling,
\eqref{4.11:eq-floors-ceilings-a}.
\end{itemize}
Then:
\begin{enumerate}
\item[(i)] $\mathsf{D}_j(X;\varphi^{\mathrm{st}}_j)[U|_X]$ holds; that is, the clause on the
$G$-valued factor of $\tilde{U}^{c,U}_j(X)$ (Definition~\ref{4.11:def-configuration-classes})
holds for the twisted pair restricted to $X$;
\item[(ii)] if moreover $X\subset\Omega_j$, then also
$\mathsf{D}^{\Sigma}_j(X;\alpha_{0,j})[U|_X]$ holds; that is, the clause on the $G$-valued
factor of the one-power class $U^{\Sigma}_j(X;\alpha_{0,j},\cdot)$ holds for the twisted pair
restricted to $X$.
\end{enumerate}
\end{lemma}

\begin{proof}
Recall from Definition~\ref{4.11:def-radii-schedules} that, for $n\le l$,
$f_{l,n}=1-\beta+\beta 2^{-(l-n)}$, so that $f_{l,l}=1$ and $f_{l,n}$ is non-increasing in $l$
for fixed $n$; that $\varphi^{\mathrm{st}}_l(n)=f_{l,n}$; and that $\varphi^{\mathrm{post}}(n)=F_n$
for $n\le k$, $\varphi^{\mathrm{post}}(k+1)=1+\beta$, where for $n\le k$
\begin{equation*}
F_n=f_{k+1,n}+\theta\beta 2^{-(k+1-n)}=1-\beta+\beta(1+\theta)2^{-(k+1-n)}.
\end{equation*}
By the level-transfer inequality \eqref{4.11:eq-level-transfer-b} of
Corollary~\ref{4.11:cor-level-transfer} (whose hypotheses, the comparability
\eqref{4.11:eq-radii-comparable-a} and the floors and ceilings
\eqref{4.11:eq-floors-ceilings-a}, are hypotheses (a) and (c)) we have, for $j<n\le k+1$,
\begin{equation}\label{4.11:eq-twist-heredity-1}
\varphi^{\mathrm{post}}(n)\,\frac{\alpha_{0,n}}{\alpha_{0,j}}\,L^{-2(n-j)}\le 1 .
\end{equation}

\emph{Part (i).} Since $(U',U,\mathbf{J})$ represents $\mathfrak{U}^{\mathrm{post},U}_{k+1}(Y)$,
its $G$-valued factor satisfies $\mathsf{D}_{k+1}(Y;\varphi^{\mathrm{post}})[U]$; we show
$\mathsf{D}_j(X;\varphi^{\mathrm{st}}_j)[U|_X]$. Let $(t,p,Z)\in\mathfrak{T}_j(X)$ and
$\mathcal{T}\in\mathcal{R}(t,p,Z)$. By \eqref{4.11:eq-radii-schedules-a},
$(t,p,Z)\in\mathfrak{T}_{k+1}(Y)$; by hypothesis~(b) the rule $\mathcal{T}$ is read at this
triple in $\mathsf{D}_{k+1}(Y;\varphi^{\mathrm{post}})[U]$. By \eqref{4.11:eq-tower-rules-a} the
tower quantities of $U$ at $(\mathcal{T},t,Z)$ depend on $U$ only through its restriction to
$Z$; they are therefore the same objects whether read from $(k+1,Y)$ or from $(j,X)$, and in
particular they are defined. Let $x\in P(p,Z)\cap\Omega_1$ and $n:=n_{k+1}(x)$; by the
definition of the layer index, $n_j(x)=\min(n,j)$.

\emph{Curvature member.} By \eqref{4.11:eq-tower-bounds-a} at level $k+1$ we are given, at $x$,
\begin{equation*}
\bigl|\partial U_t^{\mathcal{T}}(Z;\dot{\mathcal{M}}U)-1\bigr|
<\varphi^{\mathrm{post}}(n)\,\alpha_{0,n}L^{-2n},
\end{equation*}
and the required bound is
\begin{equation*}
\bigl|\partial U_t^{\mathcal{T}}(Z;\dot{\mathcal{M}}U)-1\bigr|
<f_{j,\min(n,j)}\,\alpha_{0,\min(n,j)}\,L^{-2\min(n,j)}.
\end{equation*}
There are three cases.

Case $n<j$. The required right member is $f_{j,n}\alpha_{0,n}L^{-2n}$. Since $1+\theta\le 2$,
\begin{equation*}
\varphi^{\mathrm{post}}(n)=F_n\le 1-\beta+\beta 2^{-(k-n)}=f_{k,n}\le f_{j,n},
\end{equation*}
the last step because $j\le k$ and $f_{\cdot,n}$ is non-increasing.

Case $n=j$ (so $n\le k$). The required right member is $f_{j,j}\alpha_{0,j}L^{-2j}
=\alpha_{0,j}L^{-2j}$. The computation of the previous case, which used only $n\le k$, gives
$F_j\le f_{k,j}$, and $f_{k,j}\le 1$ by the definition of $f_{k,j}$.

Case $j<n\le k+1$. The required right member is $\alpha_{0,j}L^{-2j}$, and
$\varphi^{\mathrm{post}}(n)\alpha_{0,n}L^{-2n}\le\alpha_{0,j}L^{-2j}$ by
\eqref{4.11:eq-twist-heredity-1}.

\emph{Current member.} By \eqref{4.11:eq-tower-bounds-b} at level $k+1$ we are given, at $x$,
\begin{equation*}
\bigl|J_t^{\mathcal{T}}(Z;\dot{\mathcal{M}}U)\bigr|
<\varphi^{\mathrm{post}}(n)\,\alpha_{0,n}L^{n-\min(t,n)}L^{-3(n-t)},
\end{equation*}
and the required bound is the same with $n$ replaced by $n_j=\min(n,j)$ and
$\varphi^{\mathrm{post}}$ by $\varphi^{\mathrm{st}}_j$. For $n\le j$ the two right members are the
same function of $n$ apart from the schedule factor, and the comparison of the factors is the one
made in the first two cases above. For $n>j$ we have $t\le p\le j<n$, so $\min(t,n)=t$ and the
given right member is $\varphi^{\mathrm{post}}(n)\alpha_{0,n}L^{-2(n-t)}$; the required one, with
$n_j=j\ge t$ and $f_{j,j}=1$, is
\begin{equation*}
\alpha_{0,j}L^{j-t}L^{-3(j-t)}=\alpha_{0,j}L^{-2(j-t)} .
\end{equation*}
The ratio of the given right member to the required one is
$\varphi^{\mathrm{post}}(n)(\alpha_{0,n}/\alpha_{0,j})L^{-2(n-j)}$, which is at most $1$ by
\eqref{4.11:eq-twist-heredity-1}.

In every case the given right member is at most the required right member, so the strict given
inequality implies the strict required inequality. This proves
$\mathsf{D}_j(X;\varphi^{\mathrm{st}}_j)[U|_X]$. Since the twisted pair restricted to $X$ has
witness $G$-valued factor $U|_X$ (Definition~\ref{4.11:def-step-maps}), this is its clause on the
$G$-valued factor in $\tilde{U}^{c,U}_j(X)$.

\emph{Part (ii).} Let $X\subset\Omega_j$. Let $(t,p,Z)\in\mathfrak{T}_j(X)$ and
$\mathcal{T}\in\mathcal{R}(t,p,Z)$; as in part (i), $\mathcal{T}$ is read at this triple at level
$k+1$ on $Y$ by hypothesis~(b), and the tower quantities are the same objects and defined. Let
$x\in P(p,Z)\subset X\subset\Omega_j$; since the tower regions are nested, $x\in\Omega_1$, and
$n:=n_{k+1}(x)\ge j$. Recall $p\le j\le k$, so that $F_j\le 1$ by part (i).
By the definition of the predicate $\mathsf{D}^{\Sigma}_j$
(Definition~\ref{4.11:def-tower-bounds}), $\mathsf{D}^{\Sigma}_j(X;\alpha_{0,j})[U|_X]$
asks at $x$ for the curvature bound and the current bound with the constants displayed below;
the given bounds are those of \eqref{4.11:eq-tower-bounds-a} and
\eqref{4.11:eq-tower-bounds-b} at level $k+1$, as in part (i).

\emph{Curvature member.} The required constant is
\begin{equation*}
\alpha_{0,j}L^{-(j-p)}L^{-2p}=\alpha_{0,j}L^{-(j+p)},
\end{equation*}
and the given one is $\varphi^{\mathrm{post}}(n)\alpha_{0,n}L^{-2n}$. For $n=j$,
$F_j\alpha_{0,j}L^{-2j}\le\alpha_{0,j}L^{-(j+p)}$, since $F_j\le 1$ and $p\le j$. For $n>j$,
\begin{equation*}
\varphi^{\mathrm{post}}(n)\alpha_{0,n}L^{-2n}\le\alpha_{0,j}L^{-2j}\le\alpha_{0,j}L^{-(j+p)},
\end{equation*}
by \eqref{4.11:eq-twist-heredity-1} and $p\le j$.

\emph{Current member.} The required constant is
\begin{equation*}
\alpha_{0,j}L^{-(j-p)}L^{2(t-p)}=\alpha_{0,j}L^{2t-j-p},
\end{equation*}
and, since $t\le j\le n$, the given one is $\varphi^{\mathrm{post}}(n)\alpha_{0,n}L^{-2(n-t)}$.
The ratio of the given constant to the required one satisfies
\begin{equation*}
\varphi^{\mathrm{post}}(n)\,\frac{\alpha_{0,n}}{\alpha_{0,j}}\,L^{-2n+j+p}
\le\varphi^{\mathrm{post}}(n)\,\frac{\alpha_{0,n}}{\alpha_{0,j}}\,L^{-2(n-j)}\le 1,
\end{equation*}
where the first inequality uses $p\le j$, and the second holds for $n>j$ by
\eqref{4.11:eq-twist-heredity-1} and for $n=j$ because the middle member is then $F_j\le 1$.

As in part (i), the strict given inequalities imply the strict required ones, which proves
$\mathsf{D}^{\Sigma}_j(X;\alpha_{0,j})[U|_X]$, the clause on the $G$-valued factor of
$U^{\Sigma}_j(X;\alpha_{0,j},\cdot)$ for the twisted pair restricted to $X$.
\end{proof}

\begin{lemma}[Restriction of the clause across levels and widths]\label{4.11:lem-restriction-heredity}
Let $k\ge 1$, let $Y$ be a region, let $X\subset Y$, and let $U$ be the $G$-valued factor of a
configuration on $Y$, with $U|_X$ its restriction to $X$.
\begin{enumerate}
\item[(i)] Assume that the radii satisfy \eqref{4.11:eq-radii-comparable-a}, the conclusion of
Lemma~\ref{4.11:lem-radii-comparable}, with the constant $\varrho$ given there, that the floors
and ceilings \eqref{4.11:eq-floors-ceilings-a} are in force, and that, where the sets of tower
rules read by the hypothesis and by the conclusion differ, the compatibility condition on rule
sets in the hypotheses of Lemma~\ref{4.11:lem-twist-heredity} holds. Then for all levels $i,i'$
with $1\le i'\le i\le k+1$
\begin{equation*}
\mathsf{D}_i(Y;\varphi^{\mathrm{st}}_i)[U]\ \Longrightarrow\
\mathsf{D}_{i'}(X;\varphi^{\mathrm{st}}_{i'})[U|_X],
\end{equation*}
and for all levels $i'$ with $1\le i'\le k$
\begin{equation*}
\mathsf{D}_{k+1}(Y;\varphi^{\mathrm{post}})[U]\ \Longrightarrow\
\mathsf{D}_{i'}(X;\varphi^{\mathrm{st}}_{i'})[U|_X].
\end{equation*}
The second implication is asserted only for $i'\le k$.
\item[(ii)] Let $1\le i'\le i$, and put $\mu:=L^{-(i-i')}$. Then for every value $a_0$ of the
radius argument
\begin{equation}\label{4.11:eq-restriction-heredity-1}
\mathsf{D}^{\Sigma}_i(Y;a_0)[U]\ \Longrightarrow\ \mathsf{D}^{\Sigma}_{i'}(X;\mu a_0)[U|_X],
\end{equation}
with the same $G$-valued factor on both sides; no two minimisers are compared.
\end{enumerate}
\end{lemma}

\begin{proof}
\emph{Part} (i), \emph{first implication.} The argument is that of the proof of the heredity of
the clause under the twist (Lemma~\ref{4.11:lem-twist-heredity}, which is proved as an
implication from \eqref{4.11:eq-radii-comparable-a} and, where the rule sets differ, from the
compatibility condition on rule sets), with the pair
$(k+1,\varphi^{\mathrm{post}})$ used there replaced by the pair $(i,\varphi^{\mathrm{st}}_i)$.
That argument reduces the implication to a comparison of the schedule factor carried by the
hypothesis with the one required by the conclusion, at a point of layer index $n$. With the
replacement just made there are three cases.
\begin{itemize}
\item[$n<i'$.] The factor of the hypothesis is $f_{i,n}$ and that of the conclusion is
$f_{i',n}$. The standard factors are non-increasing in the level, and $i'\le i$, so
$f_{i,n}\le f_{i',n}$.
\item[$n=i'<i$.] The factor of the hypothesis is $f_{i,i'}$; the factor of the conclusion is
$f_{i',i'}=1$; and $f_{i,i'}\le 1$.
\item[$i'<n\le i$.] Here the comparison is the level-transfer inequality
\eqref{4.11:eq-level-transfer-b} of Corollary~\ref{4.11:cor-level-transfer}, applied with the
factor $\varphi$ there equal to $f_{i,n}$ or to $1$. That corollary holds under
\eqref{4.11:eq-radii-comparable-a} and \eqref{4.11:eq-floors-ceilings-a}, which are
hypotheses of part~(i).
\end{itemize}
In each case the factor of the hypothesis is at most that of the conclusion, which proves the
first implication.

\emph{Part} (i), \emph{second implication.} This is Lemma~\ref{4.11:lem-twist-heredity} with
$j:=i'$, which is admissible because $1\le i'\le k$. The hypotheses of that lemma hold:
\eqref{4.11:eq-radii-comparable-a} and, where the rule sets differ, the compatibility condition
on rule sets are among the hypotheses of part~(i). The clause concluded in that lemma concerns
only the $G$-valued factor $U$, on which no twist acts. Its conclusion at $j=i'$ is therefore
the statement $\mathsf{D}_{i'}(X;\varphi^{\mathrm{st}}_{i'})[U|_X]$ derived from
$\mathsf{D}_{k+1}(Y;\varphi^{\mathrm{post}})[U]$.

\emph{Part} (ii). Let $(t,p,Z)\in\mathfrak{T}_{i'}(X)$. Then $(t,p,Z)\in\mathfrak{T}_i(Y)$, since
$\mathfrak{T}_{i'}(X)\subset\mathfrak{T}_i(Y)$, so the hypothesis
$\mathsf{D}^{\Sigma}_i(Y;a_0)[U]$ supplies bounds at this triple. Three things have to be
compared at the triple.
\begin{itemize}
\item The rules read at the triple are the same in the hypothesis and in the conclusion, because
the one-power class reads both tower rules at every triple.
\item The objects bounded are the same: they are the tower quantities defined by
\eqref{4.11:eq-tower-rules-a} at $(t,p,Z)$, built from the same $G$-valued factor.
\item The constants agree at every $x\in P(p,Z)$, by the two identities below.
\end{itemize}
For the constants, the curvature constant of the hypothesis is $a_0L^{-(i-p)}L^{-2p}$. Since
$\mu L^{-(i'-p)}=L^{-(i-i')}L^{-(i'-p)}=L^{-(i-p)}$,
\begin{equation*}
a_0L^{-(i-p)}L^{-2p}=(\mu a_0)\,L^{-(i'-p)}L^{-2p},
\end{equation*}
and the right-hand side is the curvature constant of the conclusion, whose radius is $\mu a_0$.
Likewise the current constant of the hypothesis is
\begin{equation*}
a_0L^{-(i-p)}L^{2(t-p)}=(\mu a_0)\,L^{-(i'-p)}L^{2(t-p)},
\end{equation*}
and the right-hand side is the current constant of the conclusion.

Thus at every triple of $\mathfrak{T}_{i'}(X)$ the conclusion asks for the same bounds, on the
same objects, under the same rules, as the hypothesis already gives. Nothing else enters the
predicate $\mathsf{D}^{\Sigma}_{i'}(X;\mu a_0)[U|_X]$, and \eqref{4.11:eq-restriction-heredity-1}
follows.
\end{proof}

\begin{proposition}[The natural witness of the substituted pair]\label{4.11:prop-natural-witness}
Fix a tower rule $\mathcal{T}$, levels $t$ and $p$ and a region $Z$, as in
Definition~\ref{4.11:def-tower-rules}. Let $\mathcal{C}^{U}_{\mathrm{src}}$ be a $U$-augmented
source class (Definition~\ref{4.11:def-configuration-classes}), namely either
$U^{\Sigma}_i(Y;a_0,a_1,\gamma_0)$ or $\tilde{U}^{U}_i(Y;\varphi)$, with
$(t,p,Z)\in\mathfrak{T}_i(Y)$ and $\mathcal{T}\in\mathcal{R}(t,p,Z)$. Let
\begin{equation*}
\mathcal{C}_{\mathrm{tgt}}:=U^{\Sigma}_t(X;a_0',a_1',\gamma_0')
\end{equation*}
be a target class on a region $X\subset P(p,Z)$. In the case treated in
\cite[\S3]{B-CMP109}, with $j$ the step of that section and $\square_0$, $\square_2$,
$\square_0^{\wr}$ its cubes, one has $(t,Z)=(j,\square_0)$ and
$X\subset\square_2\subset\square_0^{\wr}$. Assume the floors and ceilings
\eqref{4.11:eq-floors-ceilings-c} and \eqref{4.11:eq-floors-ceilings-d} of
Definition~\ref{4.11:not-floors-ceilings}, and assume the following four statements of
Ba{\l}aban.
\begin{itemize}
\item[(a)] \emph{The constrained minimiser} \cite[Theorem~1]{B-CMP102b}. Consider $G$-valued
data on a determining set obeying \cite[(7)]{B-CMP102b} with its constant
$\varepsilon_1\le a_1(d,L)$. For such data the constrained Wilson action has a minimal orbit.
When $B_3\varepsilon_1\le\varepsilon_0\le a_0(d,L)$, where $\varepsilon_0$ is the parameter of
the small-field space \cite[(6)]{B-CMP102b}, this orbit is the unique critical orbit in that
space. Its axial-gauge element is the minimal configuration $U(V)$ at the data $V$; in the tower
notation of Definition~\ref{4.11:def-tower-rules}, $U(V)=U_t^{\mathcal{T}}(Z;V)$. The constants
$a_0(d,L)$, $a_1(d,L)$ and $B_3$ depend on $d$ and $L$ only; $a_0(d,L)$ and $a_1(d,L)$ are not
the parameters $a_0,a_1$ of the class $U^{\Sigma}_i(Y;a_0,a_1,\gamma_0)$.
\item[(b)] \emph{Analytic extension of the minimiser} \cite[Proposition~9]{B-CMP102b}. The
axial-gauge minimal configuration $U(V)$ of (a) extends analytically to $G^{c}$-valued data $V$
near $G$-valued data; we write $b_H$ for the radius of this neighbourhood (the constant of
\eqref{4.11:eq-floors-ceilings-d}). For $V=V'V_0$ with $V_0$
$G$-valued and admissible and $B=(1/i)\log V'$, the configuration $U(V'V_0)U(V_0)^{-1}$, brought
to the Landau gauge, equals $\exp(i\eta\mathcal{H}(B))$. Here $\eta$ is the parameter of
\cite[Proposition~9]{B-CMP102b} multiplying the response; in the level-$t$ normalisation of the
tower quantities it takes the value $L^{-t}$. The map $\mathcal{H}$ is analytic in $B$ with
$\mathcal{H}(0)=0$. By
\cite[(181)]{B-CMP102b} this extends to $G^{c}$-orbits of data: $U(V^{v})=U(V)^{\tilde v}$,
with $\tilde v$ block-constant.
\item[(c)] \emph{Block averages of a twisted configuration} \cite[Propositions~6--7]{B-CMP98}.
Let $U$ be $G$-valued and regular, and let $\mathsf{A}$ be a small $\mathfrak{g}^{c}$-valued
bond field. Then, blockwise, $M^{q}(e^{i\mathsf{A}}U)=[e^{iB}M^{q}(U)]^{u}$, where $u$ is a
block gauge and
\begin{equation*}
|B|,\ |\log u|\le c_QL^{q}\sup|\mathsf{A}|,
\end{equation*}
the supremum being taken over the two blocks. Both $B$ and $u$ are analytic in $\mathsf{A}$, and
$c_Q=c_Q(d,L)$.
\item[(d)] \emph{Curvature of block averages} \cite[Proposition~2]{B-CMP98}. Let $U$ be
$G$-valued with $\sup|\partial U-1|$ below the threshold of \cite[Proposition~2]{B-CMP98}, which
we denote $c_2(d,L)$ (this $c_2(d,L)$ is not the constant $c_2$ of the coupling flow). Then, in
fine units, on $q$-block plaquettes
\begin{equation*}
|\partial(M^{q}U)-1|\le 2L^{2q}\sup|\partial U-1|.
\end{equation*}
\end{itemize}
Then the following holds for every representative $(U',U,\mathbf{J})$ of
$\mathcal{C}^{U}_{\mathrm{src}}$, where $U'=e^{i\mathsf{A}'}$.
\begin{enumerate}
\item The tower quantities of $U$ at $(\mathcal{T},t,Z)$ are defined and $G$-valued.
\item The $G^{c}$-orbit of $U_t^{\mathcal{T}}(Z;\dot{\mathcal{M}}(U'U))$ contains
$e^{i\mathsf{H}_{\mathrm{sub}}}\,U_t^{\mathcal{T}}(Z;\dot{\mathcal{M}}U)$. Here
$\mathsf{H}_{\mathrm{sub}}=L^{-t}\mathcal{H}(B)$, with $B$ the data twist given by (c) and
$\mathcal{H}$ the analytic response of (b). Moreover $\mathsf{H}_{\mathrm{sub}}=0$ when $U'=1$.
\end{enumerate}
Consequently, the substituted pair $\Phi_{\mathrm{sub}}(U'U,\mathbf{J})$ of
\cite[(1.15)]{B-CMP109} (Definition~\ref{4.11:def-step-maps}) admits the \emph{natural witness}
\begin{equation*}
\bigl(e^{i\mathsf{H}_{\mathrm{sub}}},\ U_t^{\mathcal{T}}(Z;\dot{\mathcal{M}}U)\bigr).
\end{equation*}
That is, its $G$-valued factor is the constrained minimiser determined by the averages of the
$G$-valued factor $U$ itself.
\end{proposition}

\begin{proof}
Let $(U',U,\mathbf{J})$ be a representative of $\mathcal{C}^{U}_{\mathrm{src}}$. Write
$\eta=L^{-t}$ for the parameter $\eta$ of (b) in the level-$t$ reading.

\emph{The tower quantities of $U$.} The clause on the $G$-valued factor in the definition of
$\mathcal{C}^{U}_{\mathrm{src}}$ at $(t,p,Z)$ (Definition~\ref{4.11:def-configuration-classes})
includes the statement that the tower quantities of $U$ at $(\mathcal{T},t,Z)$ are defined.

Independently, the first bound of \eqref{4.11:eq-configuration-classes-a} (for the class
$U^{\Sigma}_i$), resp.\ the first of \eqref{4.11:eq-configuration-classes-b} (for the class
$\tilde{U}^{U}_i$), gives at every $x$
\begin{equation*}
|\partial U-1|<(1+\beta)\sup_n\alpha_{0,n}L^{-2n(x)},
\end{equation*}
where $n(x)=n_i(x)$ is the layer index of $x$ (Definition~\ref{4.11:def-tower-rules}).
The part of \eqref{4.11:eq-floors-ceilings-c} involving the threshold $c_2(d,L)$ places $U$ under
hypothesis (d). Let $M^{q}U$ be a $q$-block datum on the determining set, lying in the layer
$\mathfrak{L}_n$. By (d),
\begin{equation*}
\begin{split}
|\partial(M^{q}U)-1|&\le 2(1+\beta)\alpha_{0,n}L^{2(q-n)}\\
&\le 2(1+\beta)\sup_n\alpha_{0,n}\le a_1(d,L).
\end{split}
\end{equation*}
The second inequality holds because $q\le n$ at every data block of either rule inside its layer.
For the rule of \cite{B-CMP119}, the data on $\mathfrak{L}_n$ are $n$-averages by construction
(Definition~\ref{4.11:def-tower-rules}). The last inequality is \eqref{4.11:eq-floors-ceilings-c}.

Hence $V_0:=D_{\mathcal{T}}^{t,Z}(U)$ obeys \cite[(7)]{B-CMP102b} with
$\varepsilon_1\le a_1(d,L)$. Its data are $G$-valued, $U$ being the $G$-valued factor. By (a),
\begin{equation*}
W_U:=U_t^{\mathcal{T}}(Z;V_0)
\end{equation*}
is the $G$-valued axial-gauge minimiser. This proves the first assertion.

\emph{The data of the twisted configuration.} Put $V:=D_{\mathcal{T}}^{t,Z}(U'U)$. On the fixed
bonds, $V=e^{i\mathsf{A}'}U$ bondwise; there put $B:=\mathsf{A}'$ and $u:=1$. On a $q$-bond
datum in layer $n$, hypothesis (c) gives $M^{q}(e^{i\mathsf{A}'}U)=[e^{iB}M^{q}(U)]^{u}$, with
\begin{equation*}
\begin{split}
|B|,\ |\log u|&\le c_QL^{q}\sup|\mathsf{A}'|\le c_Q(1+\beta)\alpha_{1,n}L^{q-n}\\
&\le c_Q(1+\beta)\alpha_{1,n}\le b_H.
\end{split}
\end{equation*}
The second inequality uses the bound $|\mathsf{A}'|<(1+\beta)\alpha_{1,n}L^{-n}$ on layer $n$,
from \eqref{4.11:eq-configuration-classes-b}. The third uses $q\le n$, and the last is
\eqref{4.11:eq-floors-ceilings-d}.

Thus $V=(e^{iB}V_0)^{u}$, with $u$ a block gauge on the determining set that is $G^{c}$-valued
and near $1$. Moreover $e^{iB}V_0$ lies in the neighbourhood of hypothesis (b).

\emph{The minimiser at the twisted data.} By \cite[(181)]{B-CMP102b}, part of hypothesis (b),
\begin{equation*}
U_t^{\mathcal{T}}(Z;V)=U_t^{\mathcal{T}}\bigl(Z;(e^{iB}V_0)^{u}\bigr)
=\bigl[U_t^{\mathcal{T}}(Z;e^{iB}V_0)\bigr]^{\tilde u}.
\end{equation*}
The main part of (b), applied with $V'=e^{iB}$, gives
\begin{equation*}
U_t^{\mathcal{T}}(Z;e^{iB}V_0)=\bigl[e^{i\eta\mathcal{H}(B)}W_U\bigr]^{v}.
\end{equation*}
Here $v$ is the change from the Landau gauge to the axial gauge; it is $G^{c}$-valued and depends
on the configuration.

Hence $U_t^{\mathcal{T}}(Z;\dot{\mathcal{M}}(U'U))$ lies in the $G^{c}$-orbit of
$e^{i\mathsf{H}_{\mathrm{sub}}}W_U$, where $\mathsf{H}_{\mathrm{sub}}:=\eta\mathcal{H}(B)$.

If $U'=1$, then $\mathsf{A}'=0$. In that case $B=0$, by the bound on $|B|$ above and by the
definition on fixed bonds, and $\mathcal{H}(0)=0$ gives $\mathsf{H}_{\mathrm{sub}}=0$. This
proves the second assertion.

\emph{The witness.} By the orbit clause of Definition~\ref{4.11:def-configuration-classes},
applied to the target class $\mathcal{C}_{\mathrm{tgt}}$ (the orbits of
\cite[(1.16)]{B-CMP109}), membership in $\mathcal{C}_{\mathrm{tgt}}$ is a condition on
$G^{c}$-orbits. Therefore
$(e^{i\mathsf{H}_{\mathrm{sub}}},W_U)$ is an admissible witness of the substituted field.
\end{proof}

\begin{lemma}[The target's curvature bound read from the source clause]\label{4.11:lem-curvature-member}
Assume the setting and the hypotheses of Proposition~\ref{4.11:prop-natural-witness}.

\emph{Setting.} A triple $(\mathcal{T},t,Z)$ is fixed, together with a $U$-augmented source class $\mathcal{C}^{U}_{\mathrm{src}}$. This class is, for some level $i$ and region $Y$, either $U^{\Sigma}_i(Y;a_0,a_1,\gamma_0)$ or $\tilde{U}^{U}_i(Y;\varphi)$, and in both cases there is a level $p$ with $(t,p,Z)\in\mathfrak{T}_i(Y)$; this $p$ is fixed. The tower rule satisfies $\mathcal{T}\in\mathcal{R}(t,p,Z)$. The target class is
\begin{equation*}
\mathcal{C}_{\mathrm{tgt}}=U^{\Sigma}_t(X;a_0',a_1',\gamma_0'),
\qquad X\subset P(p,Z).
\end{equation*}

\emph{Notation.} Let $(U',U,\mathbf{J})$ be a representative of $\mathcal{C}^{U}_{\mathrm{src}}$ and write $\mathbf{U}=U'U$. The predicate of $\mathcal{C}^{U}_{\mathrm{src}}$ imposes the same bounds on $\mathbf{U}=U'U$ and on the $G$-valued factor $U$ (Definition~\ref{4.11:def-configuration-classes}); we call these its $\mathbf{U}$-half and its $U$-half. For $x\in P(p,Z)$ let $c_{\mathrm{src}}(x)$ denote the constant of the curvature bound in the predicate of $\mathcal{C}^{U}_{\mathrm{src}}$ (of one-power type in the first case, of tower type in the second) at the index $(t,p,Z)$:
\begin{equation}\label{4.11:eq-curvature-member-1}
c_{\mathrm{src}}(x)=
\begin{cases}
a_0L^{-(i-p)}L^{-2p}
  & \text{if } \mathcal{C}^{U}_{\mathrm{src}}=U^{\Sigma}_i(Y;a_0,a_1,\gamma_0),\\[2pt]
\varphi(n_i(x))\,\alpha_{0,n_i(x)}\,L^{-2n_i(x)}
  & \text{if } \mathcal{C}^{U}_{\mathrm{src}}=\tilde{U}^{U}_i(Y;\varphi).
\end{cases}
\end{equation}
Put
\begin{equation*}
W_U:=U_t^{\mathcal{T}}(Z;\dot{\mathcal{M}}U),
\qquad
\mathbf{W}:=U_t^{\mathcal{T}}(Z;\dot{\mathcal{M}}(U'U)).
\end{equation*}
Then the following hold.
\begin{enumerate}
\item[(a)] Consider the curvature bound of the first condition in \eqref{4.11:eq-configuration-classes-a}, that is \cite[(1.11)]{B-CMP109}, for the $G$-valued factor $W_U$ of the natural witness, in the form
\begin{equation}\label{4.11:eq-curvature-member-2}
|\partial W_U-1|<a_0'\xi_t^2 \quad\text{at every } x\in X.
\end{equation}
This bound is the curvature bound of the $U$-half of the predicate of $\mathcal{C}^{U}_{\mathrm{src}}$ at the index $(t,p,Z)$, read on $X\subset P(p,Z)$, up to the comparison
\begin{equation}\label{4.11:eq-curvature-member-3}
c_{\mathrm{src}}(x)\le a_0'L^{-2t}\qquad (x\in X).
\end{equation}
\item[(b)] The comparison \eqref{4.11:eq-curvature-member-3} is exactly the comparison needed to obtain, from the $\mathbf{U}$-half of the same predicate at the same index, the curvature bound of the third condition in \eqref{4.11:eq-configuration-classes-a}. That bound is the first inequality of \cite[(1.14)]{B-CMP109}, here for the substituted field $\mathbf{W}$. The reason is twofold: both bounds of the target class carry the constant $a_0'$, and both halves of the predicate of $\mathcal{C}^{U}_{\mathrm{src}}$ carry the same constant $c_{\mathrm{src}}$.
\end{enumerate}
\end{lemma}

\begin{proof}
We take $W_U$ first.

\emph{What the target class requires of $W_U$.} The predicate of $\mathcal{C}^{U}_{\mathrm{src}}$ at $(t,p,Z)$ includes that the tower quantities of $U$ at $(\mathcal{T},t,Z)$ are defined (Definition~\ref{4.11:def-tower-bounds}). The $G$-valued factor of the natural witness is $U_t^{\mathcal{T}}(Z;\dot{\mathcal{M}}U)=W_U$. The first condition in \eqref{4.11:eq-configuration-classes-a}, written for $\mathcal{C}_{\mathrm{tgt}}=U^{\Sigma}_t(X;a_0',a_1',\gamma_0')$, asks that at every $x\in X$
\begin{equation*}
|\partial W_U-1|<a_0'\xi_t^2=a_0'L^{-2t},
\end{equation*}
since $\xi_t=L^{-t}$.

\emph{What the source class gives for $W_U$.} The $U$-half of the predicate of $\mathcal{C}^{U}_{\mathrm{src}}$ is taken at the triple $(t,p,Z)$ and the tower rule $\mathcal{T}$. By Definition~\ref{4.11:def-tower-bounds}, it gives at every $x\in P(p,Z)$
\begin{equation*}
|\partial U_t^{\mathcal{T}}(Z;\dot{\mathcal{M}}U)-1|<c_{\mathrm{src}}(x),
\end{equation*}
with $c_{\mathrm{src}}(x)$ as in \eqref{4.11:eq-curvature-member-1}. In the first case the constant comes from the one-power predicate $\mathsf{D}^{\Sigma}_i(Y;a_0)[U]$ at the triple $(t,p,Z)$. In the second case it comes from the tower-type predicate $\mathsf{D}_i(Y;\varphi)[U]$.

\emph{Conclusion of (a).} We have $P(p,Z)\supset X$ and $U_t^{\mathcal{T}}(Z;\dot{\mathcal{M}}U)=W_U$. Hence \eqref{4.11:eq-curvature-member-2} holds on $X$ as soon as \eqref{4.11:eq-curvature-member-3} holds, because then
\begin{equation*}
|\partial W_U-1|<c_{\mathrm{src}}(x)\le a_0'L^{-2t}\qquad (x\in X).
\end{equation*}
This proves (a).

\emph{The substituted field.} The third condition in \eqref{4.11:eq-configuration-classes-a}, written for $\mathcal{C}_{\mathrm{tgt}}$ and the substituted field $\mathbf{W}$, asks that at every $x\in X$
\begin{equation*}
|\partial\mathbf{W}-1|<a_0'L^{-2t}.
\end{equation*}
The $\mathbf{U}$-half of the predicate of $\mathcal{C}^{U}_{\mathrm{src}}$ at the same triple gives, at every $x\in P(p,Z)$,
\begin{equation*}
|\partial\mathbf{W}-1|<c_{\mathrm{src}}(x),
\end{equation*}
with the same constant $c_{\mathrm{src}}(x)$. Indeed, by Definition~\ref{4.11:def-configuration-classes} the two halves carry one constant. In the one-power case the same predicate $\mathsf{D}^{\Sigma}_i(Y;a_0)$ is imposed on $U'U$ and on $U$; this is the last condition in \eqref{4.11:eq-configuration-classes-a}. In the tower-type case the same $\varphi$ enters both halves of $\tilde{U}^{U}_i(Y;\varphi)$.

\emph{Conclusion of (b).} The inequality between constants that delivers the curvature bound of the third condition for $\mathbf{W}$ from the $\mathbf{U}$-half is therefore \eqref{4.11:eq-curvature-member-3}. By the argument for (a), the same inequality delivers the curvature bound of the first condition for $W_U$ from the $U$-half. Conversely, the inequality that delivers the bound for $W_U$ delivers the bound for $\mathbf{W}$. This proves (b).

\emph{Where the comparison is made in \cite{B-CMP109}.} In \cite[\S3]{B-CMP109} the source class is taken at halved radii at equal level. Across levels, the one power of the last condition in \eqref{4.11:eq-configuration-classes-a} gives
\begin{equation*}
c_{\mathrm{src}}\le a_0L^{-(i-t)}L^{-2t}.
\end{equation*}
In either situation the comparison is the one needed for the $\mathbf{U}$-half, whatever form it takes, and no further condition arises for the $G$-valued factor $U$.
\end{proof}

\begin{proposition}[Heredity of the clause under substitution]\label{4.11:prop-substitution-heredity}
Work in the setting and under the assumptions of
Proposition~\ref{4.11:prop-natural-witness}. Thus $(\mathcal{T},t,Z)$ is fixed,
$\mathcal{C}^{U}_{\mathrm{src}}$ is the $U$-augmented source class,
$\Phi_{\mathrm{sub}}$ is the substitution map, and the target class is
$\mathcal{C}_{\mathrm{tgt}}=U^{\Sigma}_t(X;a_0',a_1',\gamma_0')$ on the region $X$.
For a representative $(V',V,\mathbf{K})$ of $\mathcal{C}^{U}_{\mathrm{src}}$ write
$\mathbf{J}_{\mathrm{sub}}$ for the current component of the natural witness of
$\Phi_{\mathrm{sub}}(V'V,\mathbf{K})$, so that this witness is
\begin{equation*}
\bigl(e^{i\mathsf{H}_{\mathrm{sub}}(V'V)},\;
U_t^{\mathcal{T}}(Z;\dot{\mathcal{M}}V),\;\mathbf{J}_{\mathrm{sub}}\bigr).
\end{equation*}
Assume the following substitution hypothesis on the full complex factor. For every
representative $(V',V,\mathbf{K})$ of
$\mathcal{C}^{U}_{\mathrm{src}}$, the real representatives $(1,V,0)$ included, the natural
witness of $\Phi_{\mathrm{sub}}(V'V,\mathbf{K})$ satisfies two sets of conditions.
It satisfies the pointwise members of the target class, that is, the first three members
listed in~\eqref{4.11:eq-configuration-classes-a}. It also satisfies the half of the
derived-slot member of the target evaluated at the full complex factor,
\begin{equation*}
\mathsf{D}^{\Sigma}_t(X;a_0')\bigl[e^{i\mathsf{H}_{\mathrm{sub}}}\,
U_t^{\mathcal{T}}(Z;\dot{\mathcal{M}}V)\bigr].
\end{equation*}
Then for every representative $(U',U,\mathbf{J})$ of $\mathcal{C}^{U}_{\mathrm{src}}$
the natural witness of $\Phi_{\mathrm{sub}}(U'U,\mathbf{J})$ also satisfies the half
of the derived-slot member evaluated at its $G$-valued factor,
\begin{equation*}
\mathsf{D}^{\Sigma}_t(X;a_0')\bigl[U_t^{\mathcal{T}}(Z;\dot{\mathcal{M}}U)\bigr].
\end{equation*}
Hence $\Phi_{\mathrm{sub}}$ maps $\mathcal{C}^{U}_{\mathrm{src}}$ into the $U$-augmented
target class.
\end{proposition}

\begin{proof}
We apply the transfer lemma (Lemma~\ref{4.11:lem-transfer}) with the following choices:
\begin{itemize}
\item the source class is $\mathcal{C}^{U}_{\mathrm{src}}$;
\item the target class is $\mathcal{C}_{\mathrm{tgt}}$, with derived-slot predicate
$\mathsf{D}_{\mathrm{tgt}}=\mathsf{D}^{\Sigma}_t(X;a_0')$;
\item $\Phi:=\Phi_{\mathrm{sub}}$;
\item $\Psi(U):=U_t^{\mathcal{T}}(Z;\dot{\mathcal{M}}U)$ on $G$-valued configurations $U$;
\item for each representative, the witness required by the lemma is the natural witness.
\end{itemize}
The natural witness has twist factor
$U'_{\Phi}=e^{i\mathsf{H}_{\mathrm{sub}}(U'U)}$ and current
$\mathbf{J}_{\Phi}=\mathbf{J}_{\mathrm{sub}}$. The two hypotheses of the lemma are displayed
in~\eqref{4.11:eq-transfer-a}, and we check them in turn.

The first hypothesis asks the following of every representative $(U',U,\mathbf{J})$
of $\mathcal{C}^{U}_{\mathrm{src}}$: the orbit of $\Phi(U'U,\mathbf{J})$ must contain
a witness $(U'_{\Phi},\Psi(U),\mathbf{J}_{\Phi})$ that satisfies the target's
pointwise members and $\mathsf{D}_{\mathrm{tgt}}[U'_{\Phi}\Psi(U)]$. With the choices
above, this is word for word the substitution hypothesis on the full complex factor.

The second hypothesis asks that at a real representative $(1,U,0)$ this witness have trivial
twist factor. At $U'=1$, Proposition~\ref{4.11:prop-natural-witness} gives $B=0$ for the
argument $B$ of the analytic response $\mathcal{H}$ entering $\mathsf{H}_{\mathrm{sub}}$.
By the property of the analytic response fixed where $\mathcal{H}$ is introduced, namely
that it vanishes at zero argument, $\mathcal{H}(0)=0$. Hence $\mathsf{H}_{\mathrm{sub}}=0$ and
$U'_{\Phi}=e^{0}=1$. The natural witness of $\Phi_{\mathrm{sub}}(U,0)$ is therefore
\begin{equation*}
\bigl(1,\;U_t^{\mathcal{T}}(Z;\dot{\mathcal{M}}U),\;
J_t^{\mathcal{T}}(Z;\dot{\mathcal{M}}U)\bigr),
\end{equation*}
which is a real point. This is the second hypothesis.

The transfer lemma now yields, for every representative $(U',U,\mathbf{J})$ of
$\mathcal{C}^{U}_{\mathrm{src}}$,
\begin{equation*}
\mathsf{D}^{\Sigma}_t(X;a_0')\bigl[U_t^{\mathcal{T}}(Z;\dot{\mathcal{M}}U)\bigr].
\end{equation*}
This is the half of the derived-slot member of~\eqref{4.11:eq-configuration-classes-a}
evaluated at the $G$-valued factor of the natural witness. The substitution hypothesis on the
full complex factor supplies the pointwise members and the half evaluated at the full complex
factor. With all these members the natural witness is a representative of the $U$-augmented
version of $\mathcal{C}_{\mathrm{tgt}}$.

Of the clause of the source class on the $G$-valued factor, only two things were used. The
first is that it makes $(1,U,0)$ a representative of $\mathcal{C}^{U}_{\mathrm{src}}$. This is
the retention of real points by $U$-augmented classes, established in the first step of the
proof of the transfer lemma (Lemma~\ref{4.11:lem-transfer}), on which that lemma rests. The
second is its single member at the index triple $(t,p_0,Z)$. This is the member from which
Lemma~\ref{4.11:lem-curvature-member} reads the target's curvature bound for the $G$-valued
factor, one of the pointwise members that the substitution hypothesis supplies.
\end{proof}

\begin{lemma}[Idempotence of the constrained minimiser at equal index]\label{4.11:lem-idempotence}
Let $\mathcal{T}$ be a tower rule, $Z\subset\mathbb{L}$ a region and $t\ge 1$ an index, and let $U$
be a $G$-valued configuration on the bonds of $Z$ whose data $D_{\mathcal{T}}^{t,Z}(U)$ are
admissible, so that the tower quantities of Definition~\ref{4.11:def-tower-rules} are defined
at $U$. Then
\begin{equation}\label{4.11:eq-idempotence-1}
U_t^{\mathcal{T}}\bigl(Z;\dot{\mathcal{M}}\bigl(U_t^{\mathcal{T}}(Z;\dot{\mathcal{M}}U)\bigr)\bigr)
=U_t^{\mathcal{T}}(Z;\dot{\mathcal{M}}U).
\end{equation}
\end{lemma}

\begin{proof}
Put $D:=D_{\mathcal{T}}^{t,Z}(U)$ and $W:=U_t^{\mathcal{T}}(Z;\dot{\mathcal{M}}U)$, so that
$W=U_t^{\mathcal{T}}(Z;D)$ by Definition~\ref{4.11:def-tower-rules}. Since $U$ is $G$-valued and
$D$ is admissible, the problem $P_{\mathcal{T}}(Z,t;D)$ has a unique minimal orbit
(\cite[Theorem~1]{B-CMP102b}, as recorded in Definition~\ref{4.11:def-tower-rules}), and $W$ is
the axial-gauge element of that orbit. In particular $W$ is a $G$-valued configuration on the
bonds of $Z$ which is admissible for its own problem, that is, $W$ satisfies the three
constraints defining $P_{\mathcal{T}}(Z,t;D)$:
\begin{enumerate}
\item $W=U$ on the fixed bonds, the bonds of $Z$ not having both endpoints in
$\Omega_1^{\mathcal{T}}(Z)$;
\item $M^{q}(W)(c)=M^{q}(U)(c)$ for every $c\in\Gamma_q^{\mathcal{T}}(Z)$ and every $q$ with
$1\le q<t$;
\item $M^{t}(W)(c)=M^{t}(U)(c)$ for every $c\in\Gamma_t^{\mathcal{T},\mathrm{top}}(Z)$.
\end{enumerate}

Now consider the problem $P_{\mathcal{T}}(Z,t;D_{\mathcal{T}}^{t,Z}(W))$. By the definition of
the data of a configuration, $D_{\mathcal{T}}^{t,Z}(W)$ is the tuple consisting of $W$ on the
fixed bonds, of $M^{q}(W)$ on $\Gamma_q^{\mathcal{T}}(Z)$ for $1\le q<t$, and of $M^{t}(W)$ on
$\Gamma_t^{\mathcal{T},\mathrm{top}}(Z)$. The sets on which these entries are read, namely the
fixed bonds, the sets $\Gamma_q^{\mathcal{T}}(Z)$ and the set
$\Gamma_t^{\mathcal{T},\mathrm{top}}(Z)$, are built from the regions
$\Omega_q^{\mathcal{T}}(Z)$, $0\le q\le t$, which the tower rule determines from $Z$, $q$ and
the fixed regions alone; hence these sets depend on $(\mathcal{T},t,Z)$ only, and are the same
for $W$ as for $U$. Entry by entry, the three constraints above therefore say that the data of
$W$ coincide with those of $U$:
\begin{equation*}
D_{\mathcal{T}}^{t,Z}(W)=D=D_{\mathcal{T}}^{t,Z}(U).
\end{equation*}
In particular the data of $W$ are admissible, so that the left-hand side of
\eqref{4.11:eq-idempotence-1} is defined. Since $U_t^{\mathcal{T}}(Z;\cdot\,)$ is a function of
the data alone, it takes the same value at $D_{\mathcal{T}}^{t,Z}(W)$ as at
$D_{\mathcal{T}}^{t,Z}(U)$, that is,
\begin{equation*}
U_t^{\mathcal{T}}(Z;\dot{\mathcal{M}}W)=U_t^{\mathcal{T}}(Z;D)=W,
\end{equation*}
which is \eqref{4.11:eq-idempotence-1}.
\end{proof}

The identity extends to complex configurations $\mathbf{U}$ whose data lie in the neighbourhood
of the real data on which $U_t^{\mathcal{T}}(Z;\cdot\,)$ is defined by analytic continuation
(Definition~\ref{4.11:def-tower-rules}): there the three constraints are analytic identities
which hold on the real slice, hence throughout the connected domain. This extension is not used
in what follows.

\begin{remark}[No composition identity across tower indices]\label{4.11:rem-no-composition}
Tower rules, the data $\dot{\mathcal{M}}V$ of a configuration, the constrained problems
$P_{\mathcal{T}}(Z,t;D)$ and the constrained minimisers $U_t^{\mathcal{T}}(Z;D)$ are those of
Definition~\ref{4.11:def-tower-rules}.
\begin{enumerate}
\item[$(\alpha)$] For every tower rule $\mathcal{T}$, every region $Z$, every index $t$ and every
$G$-valued configuration $U$ on the bonds of $Z$ with admissible data, the idempotence identity of
Lemma~\ref{4.11:lem-idempotence} holds: the constrained minimiser under $\mathcal{T}$ with index
$t$ on $Z$ whose data are the data of $U_t^{\mathcal{T}}(Z;\dot{\mathcal{M}}U)$ is
$U_t^{\mathcal{T}}(Z;\dot{\mathcal{M}}U)$ itself.
\item[$(\beta)$] Let $\mathcal{T}$, $Z$, $t$ and $U$ be as in $(\alpha)$, let $\mathcal{T}'$ be a
tower rule, let $Z'\subsetneq Z$ and $t'\le t$, and assume that there are indices $1\le r<t'$ and
$r<q<t$ and a $q$-bond $c\in\Gamma_q^{\mathcal{T}}(Z)$ whose endpoint blocks meet the interface
$\partial\Omega_r^{\mathcal{T}'}(Z')$ of the tower of $Z'$ under $\mathcal{T}'$ (such $r$, $q$, $c$
exist wherever, near $\partial Z'$ inside $Z$, the tower of $Z$ under $\mathcal{T}$ is deeper than
the tower of $Z'$ under $\mathcal{T}'$). Put
$W:=U_t^{\mathcal{T}}(Z;\dot{\mathcal{M}}U)$. In general the
definitions supply no identity between $U_{t'}^{\mathcal{T}'}(Z';\dot{\mathcal{M}}W)$ and the
restriction $W|_{Z'}$. The two problems
\begin{equation*}
P_{\mathcal{T}'}(Z',t';\dot{\mathcal{M}}W)
\qquad\text{and}\qquad
P_{\mathcal{T}}(Z,t;\dot{\mathcal{M}}U)
\end{equation*}
impose different conditions on $Z'$: the data $\dot{\mathcal{M}}W$ and $\dot{\mathcal{M}}U$ of the
problem on $Z'$ differ in general on the fixed bonds of $Z'$ that are not fixed bonds of $Z$, and on
those bonds of $\Gamma_q^{\mathcal{T}'}(Z')$, $q<t'$, at which the problem on $Z$ constrains only
coarser averages or none. Further, the admissible
variations of the first problem need not be contained in those of the second. A statement relating
$U_{t'}^{\mathcal{T}'}(Z';\dot{\mathcal{M}}W)$ to $W|_{Z'}$ is a comparison of minimisers along
two towers; no such comparison is used in this section.
\end{enumerate}
\end{remark}

\begin{proof}
Clause $(\alpha)$ is the idempotence of the constrained minimiser at equal index,
Lemma~\ref{4.11:lem-idempotence}; it follows from the definitions.

For clause $(\beta)$ we use two properties of the block averages $M^q$ of
Definition~\ref{1.5:def-block-average}, which follow from their definition by iterated averaging in
\cite[(1.5)]{B-CMP99a}. Let $c=\langle B,B'\rangle$ be a $q$-bond joining the
$q$-blocks $B$ and $B'$. First,
$M^q(\mathbf{U})(c)$ depends only on the restriction of $\mathbf{U}$ to the bonds with both
endpoints in $B\cup B'$. Second, $M^q(\mathbf{U})(c)$ is a fixed function of the
values $M^{q-1}(\mathbf{U})(c')$ on the $(q-1)$-bonds $c'$ with endpoints in $B\cup B'$
(iterated averaging).

\emph{(a) The data differ.} The data of $P_{\mathcal{T}'}(Z',t';\dot{\mathcal{M}}W)$ consist in
part of the values of $W$ itself on the fixed bonds of $Z'$, that is, on the bonds of $Z'$ not
having both endpoints in $\Omega_1^{\mathcal{T}'}(Z')$. Those of these bonds that are not fixed
bonds of $Z$ are varied in $P_{\mathcal{T}}(Z,t;\dot{\mathcal{M}}U)$, and on them the minimiser
$W$ differs from $U$ in general. Further, for $q<t'$ the constraints of
$P_{\mathcal{T}'}(Z',t';\dot{\mathcal{M}}W)$ on the constraint bond set
$\Gamma_q^{\mathcal{T}'}(Z')$ of the shell
$\Omega_q^{\mathcal{T}'}(Z')\setminus\Omega_{q+1}^{\mathcal{T}'}(Z')$ read the
averages $M^q(W)(c)$, $c\in\Gamma_q^{\mathcal{T}'}(Z')$; at those of these $q$-bonds at which the
problem on $Z$ constrains only coarser averages, or none, the averages $M^q(W)(c)$ and
$M^q(U)(c)$ differ in general. Hence, in general, $P_{\mathcal{T}'}(Z',t';\dot{\mathcal{M}}W)$ and
$P_{\mathcal{T}'}(Z',t';\dot{\mathcal{M}}U)$ have different data. In particular the identity
\begin{equation*}
\dot{\mathcal{M}}\bigl(U_t^{\mathcal{T}}(Z;\dot{\mathcal{M}}U)\bigr)=\dot{\mathcal{M}}U
\quad\text{on the constraint sets of } Z'
\end{equation*}
is not available.

\emph{(b) The admissible variations.} The restriction $W|_{Z'}$ satisfies the constraints of
$P_{\mathcal{T}'}(Z',t';\dot{\mathcal{M}}W)$ tautologically. Suppose every admissible variation
of $P_{\mathcal{T}'}(Z',t';\dot{\mathcal{M}}W)$, extended to the bonds of $Z$ by $W$, were an
admissible variation of $P_{\mathcal{T}}(Z,t;\dot{\mathcal{M}}U)$. Then the difference of the
Wilson actions would localise in $Z'$, and the minimality of $W$ would transfer to $W|_{Z'}$.
By the uniqueness of the minimal orbit, \cite[Theorem~1]{B-CMP102b},
$U_{t'}^{\mathcal{T}'}(Z';\dot{\mathcal{M}}W)$ would then be the axial-gauge representative of the
gauge orbit of $W|_{Z'}$.

At equal index this inclusion of admissible sets holds; this is the content of
Lemma~\ref{4.11:lem-idempotence}. Under the hypothesis of $(\beta)$, with $Z'\subsetneq Z$, it is
not supplied by the two properties
of the block averages and the nesting of the regions, for the following reason. Let $r$, $q$ and
the $q$-bond $c\in\Gamma_q^{\mathcal{T}}(Z)$ be as in the hypothesis of $(\beta)$: $c$ is a coarse
$q$-bond constrained in the
problem on $Z$ whose endpoint blocks meet the interface $\partial\Omega_r^{\mathcal{T}'}(Z')$ of
the tower of $Z'$. By the nesting property, for a configuration $V$ on the bonds of $Z$
the value $M^q(V)(c)$ is computed from $r$-averages of $r$-bonds
crossing that interface and from finer averages of blocks inside the shell
\begin{equation*}
\Omega_r^{\mathcal{T}'}(Z')\setminus\Omega_{r+1}^{\mathcal{T}'}(Z').
\end{equation*}
None of these is a constraint of the problem on $Z'$: $Z'$ constrains the $r$-averages of the
$r$-bonds having both endpoints in that shell, and only those. A configuration $V$
obtained from $W$ by a variation admissible for the problem on $Z'$ therefore need not satisfy
$M^q(V)(c)=M^q(W)(c)$, and when it does not, it is not admissible for the problem on $Z$, since
$c\in\Gamma_q^{\mathcal{T}}(Z)$; the inclusion of admissible sets is thus not a consequence of the
definitions. Bonds $c$ of
this kind exist as soon as $L^q$ exceeds the width of the shells of $Z'$ near
$c$; the hypothesis of $(\beta)$ provides one.

Consequently the definitions show $W|_{Z'}$ to be critical for
$P_{\mathcal{T}'}(Z',t';\dot{\mathcal{M}}W)$ only along
the directions this problem shares with $P_{\mathcal{T}}(Z,t;\dot{\mathcal{M}}U)$. Whether
$U_{t'}^{\mathcal{T}'}(Z';\dot{\mathcal{M}}W)$ coincides with the axial-gauge representative of the
gauge orbit of $W|_{Z'}$ is thus a genuine question of comparing
minimisers along two towers, as asserted in $(\beta)$.
\end{proof}

\begin{lemma}[Heredity of the clause under large-field composites]\label{4.11:lem-composite-heredity}
Let $\mathcal{W}$ be a composite of the $\mathbf{R}$-chain of \cite[(1.50)--(1.65)]{B-CMP122b},
one of the maps of the step of Definition~\ref{4.11:def-step-maps}. The class from which its
argument is taken is called the \emph{source class}. The argument is a point of the source class
augmented by the clause on the $G$-valued factor of
Definition~\ref{4.11:def-configuration-classes}, and we write $\bar{U}$ for the $G$-valued factor
of this point; in particular $\bar{U}$ satisfies the source's clause. A \emph{consumer domain} of
the chain is the domain of an older term that receives a value of $\mathcal{W}$ as its argument.
In \cite{B-CMP122b} the chain comes with a distinguished set of domains, called the
\emph{cut band}; only consumer domains not belonging to it are considered in hypothesis~(a)
below and in the conclusion. The
class of the older term, augmented by the clause on the $G$-valued factor of
Definition~\ref{4.11:def-configuration-classes}, is called the \emph{target class}. Assume the
following two statements.
\begin{enumerate}
\item[(a)] \emph{Composite representation.} Let $X'$ be any consumer domain of the chain outside
the cut band. Then, up to a $G^{c}$-valued gauge transformation depending holomorphically on the
argument in the source class,
\begin{equation*}
\mathcal{W}|_{X'}=e^{i\mathsf{A}''}U^{\natural}
\end{equation*}
for some $\mathfrak{g}^{c}$-valued bond potential $\mathsf{A}''$ on $X'$,
where either $U^{\natural}=\bar{U}$, or $U^{\natural}=U_2$ with $U_2$ a fixed $G$-valued
configuration independent of the argument in the source class. Moreover, the derived-slot
predicates of source and target (the predicates $\mathsf{D}_i(X;\varphi)[\,\cdot\,]$ entering
the clause on the $G$-valued factor) are read at the same triples and with the same rules
for sub-domains of $X'$.
This reading on sub-domains is the one used for the one-power class. At every such
triple the target constants dominate the source constants.
\item[(b)] \emph{Slot property of $U_2$.} The configuration $U_2$ satisfies the derived-slot
predicates of every target class it meets.
\end{enumerate}
Then, for every consumer domain $X'$ outside the cut band, $\mathcal{W}|_{X'}$ satisfies the
clause on the $G$-valued factor of the target class attached to $X'$, with witness
$G$-valued factor $U^{\natural}$.
\end{lemma}

\begin{proof}
Let $X'$ be a consumer domain of the chain outside the cut band. By hypothesis~(a), after a
$G^{c}$-valued gauge transformation depending holomorphically on the argument in the source class,
$\mathcal{W}|_{X'}$ equals
$e^{i\mathsf{A}''}U^{\natural}$, with $U^{\natural}\in\{\bar{U},U_2\}$.

The target class is a union of orbits under $G^{c}$-valued gauge transformations: this is the
orbit property of the configuration classes of Definition~\ref{4.11:def-configuration-classes},
from which the target class is built. Membership
in the target
class is therefore unchanged when the gauge transformation of hypothesis~(a) is applied. Hence,
after this gauge transformation, the pair $(e^{i\mathsf{A}''},U^{\natural})$ is an admissible
witness: a twist factor together with a $G$-valued factor whose product is the restricted
composite. It remains to show that the $G$-valued factor $U^{\natural}$ satisfies the target's half
of the clause, that is, the target predicate $\mathsf{D}_{\mathrm{tgt}}[U^{\natural}|_{X'}]$. We
distinguish the two cases allowed by hypothesis~(a).

\emph{Case $U^{\natural}=\bar{U}$.} The target's half of the clause is
$\mathsf{D}_{\mathrm{tgt}}[\bar{U}|_{X'}]$. This is a family of inequalities, one for each triple
$(t,p,Z)$ in the index set $\mathfrak{T}_i(X')$ of triples of the target class, $i$ denoting the
level at which the target class is read, and each rule
$\mathcal{T}$ that the target reads at that triple.

By hypothesis~(a), every such triple and rule is also read by the source class, on the same
sub-domain of $X'$. This holds because the reading of derived slots on sub-domains is the same for
source and target, namely the reading used for the one-power class. Also by hypothesis~(a), at each
such triple the target constant is at least the source constant.

By the tower rules \eqref{4.11:eq-tower-rules-a}, the tower quantities
$U_t^{\mathcal{T}}(Z;\dot{\mathcal{M}}\bar{U})$ and $J_t^{\mathcal{T}}(Z;\dot{\mathcal{M}}\bar{U})$
of $\bar{U}$ entering the target inequality are the same objects as those entering the source
inequality at the same triple and rule.

Since the argument lies in the augmented source class, its $G$-valued factor $\bar{U}$ satisfies
the source's clause on the $G$-valued factor. In particular it satisfies the source inequality at
$(t,p,Z)$ and
$\mathcal{T}$. The two inequalities bound the same tower quantities, and the target inequality
carries a constant at least as large. Hence the target inequality at that triple follows from the
source inequality at the same triple.

This holds for every triple of $\mathfrak{T}_i(X')$ and every rule read there.
Therefore $\mathsf{D}_{\mathrm{tgt}}[\bar{U}|_{X'}]$ holds. The comparison of constants used here
is exactly the one asserted in hypothesis~(a).

\emph{Case $U^{\natural}=U_2$.} The witness $G$-valued factor is the fixed configuration $U_2$.
The target's half of the clause is $\mathsf{D}_{\mathrm{tgt}}[U_2|_{X'}]$. The target class is
met by $U_2$, so this is precisely what hypothesis~(b) asserts.

In both cases $(e^{i\mathsf{A}''},U^{\natural})$ is, after the gauge transformation, an admissible
witness whose $G$-valued factor satisfies the target's derived-slot predicate on $X'$. Therefore
the clause
on the $G$-valued factor of the target class attached to $X'$ holds for $\mathcal{W}|_{X'}$, with
witness $G$-valued factor $U^{\natural}$.
\end{proof}

\begin{theorem}[Reproduction of the $G$-valued-factor clause by the step]\label{4.11:thm-clause-reproduction}
Declare the $U$-augmented post-step class $\mathfrak{U}^{\mathrm{post},U}_{k+1}(Y)$ to be the
analyticity domain of the new boundary terms $\mathbf{B}^{(k+1)}(Y)$ of \cite{B-CMP119} at a
straddling region $Y$. Assume the following.
\begin{enumerate}
\item[(a)] The comparability of the radii along the coupling flow \eqref{4.11:eq-radii-comparable-a} holds with constant $\varrho$, and the floor \eqref{4.11:eq-floors-ceilings-a} holds, namely
\begin{equation*}
L^2\ \ge\ \sqrt{2}\,(1+\beta)\,\varrho .
\end{equation*}
The running couplings satisfy the ceilings on the coupling fixed in
Notation~\ref{4.11:not-floors-ceilings}.
\item[(b)] The substitution hypothesis of Proposition~\ref{4.11:prop-substitution-heredity} holds. It says the following, in the setting of that proposition. Let $(V',V,\mathbf{K})$ be any representative of the $U$-augmented source class $\mathcal{C}^{U}_{\mathrm{src}}$. Then the natural witness of the substituted pair $\Phi_{\mathrm{sub}}(V'V,\mathbf{K})$ satisfies the first three defining members of the target class $\mathcal{C}_{\mathrm{tgt}}$ and the $\mathbf{U}$-half of its derived-slot member, that is, that predicate evaluated at the complex configuration itself:
\begin{equation*}
\mathsf{D}^{\Sigma}_t(X;a_0')\bigl[e^{i\mathsf{H}_{\mathrm{sub}}}U_t^{\mathcal{T}}(Z;\dot{\mathcal{M}}V)\bigr].
\end{equation*}
\item[(c)] The two hypotheses of Lemma~\ref{4.11:lem-composite-heredity} hold.
\begin{itemize}
\item The first concerns every consumer domain $X'$ of Ba{\l}aban's large-field chain, that is,
every region on which a later term of the chain takes the composite as its argument, lying outside
the cut band of the chain. On each such $X'$ the composite $\mathcal{W}$ of the chain is a twist $e^{i\mathsf{A}''}U^{\natural}$, up to a $G^{c}$-valued gauge transformation holomorphic on the source. Here $U^{\natural}$ is either the source's $G$-valued factor $\bar{U}$, or a fixed $G$-valued configuration $U_2$ independent of the source. The derived slots of source and target are read at the same triples and rules for sub-domains of $X'$, and at every such triple the target constants dominate the source constants.
\item The second says that $U_2$ satisfies the derived slots of every target it meets.
\end{itemize}
\item[(d)] The hypothesis of Lemma~\ref{4.11:lem-twist-heredity} on the rules read at a triple,
which is needed only where the rule sets of source and target differ, holds.
\item[(e)] The four published statements on which the existence of the natural witness of a
substituted pair rests hold in the form in which that existence statement uses them. They are
taken from \cite[Theorem~1, Proposition~9]{B-CMP102b} and \cite[Propositions~2, 6--7]{B-CMP98};
among them are the statement supplying the data twist $B$ and the statement supplying the
analytic response $\mathcal{H}(B)$ of the constrained minimiser to that twist.
\end{enumerate}
Then the following hold.
\begin{enumerate}
\item[(i)] At real points the clause on the $G$-valued factor imposes no condition. A real point
carrying the trivial witness lies in a $U$-augmented class exactly when it lies in the
corresponding plain class.
\item[(ii)] On Cauchy families over a fixed $G$-valued factor $U$ the clause imposes exactly one
further condition. For every level $i$, region $X$ and schedule $\varphi$, membership of every
member, taken with the witness $(U'_\zeta,U)$ formed from its twist factor $U'_\zeta$ and the fixed
factor $U$, in the $U$-augmented class
$\tilde{U}^{U}_i(X;\varphi)$ is equivalent to two things together: membership of every member in
the plain class $\tilde{U}_i(X;\varphi)$, and the single condition $\mathsf{D}_i(X;\varphi)[U]$,
which does not depend on the parameter $\zeta$ of the family.
\item[(iii)] The clause is hereditary. Let $\Phi$ be a production map of the step: a restriction, a twist, a substitution, or a composite of the large-field chain. If $\Phi$ maps the $U$-augmented source class into the target class without the clause on the $G$-valued factor, then $\Phi$ maps it into the $U$-augmented target class.
\end{enumerate}
\end{theorem}

\begin{proof}
Part (i) is Lemma~\ref{4.11:lem-real-points}. That lemma covers each of the classes
$\tilde{U}_i(X;\varphi)$, the analyticity class of \cite{B-CMP109}, and $U^{\Sigma}_j$, and it uses
no hypothesis.

Part (ii) is Lemma~\ref{4.11:lem-cauchy-families}, which likewise uses no hypothesis.

For part (iii) we fix a production map $\Phi$ such that $\Phi$ maps the $U$-augmented source class into the target class without the clause. For each kind of production map this antecedent is the
content of the reproduction statement of the step for that kind of map without the clause on the
$G$-valued factor; in what follows it enters only as an antecedent. We treat the kinds of $\Phi$ in turn.

\emph{Restriction.} Suppose $\Phi$ is a restriction across levels and widths. The $G$-valued factor of the image is the restriction of the source's $G$-valued factor. We must pass the derived-slot predicate from the source region and level to the target region and level.
\begin{itemize}
\item For the tower-type predicate $\mathsf{D}_i$ this is Lemma~\ref{4.11:lem-restriction-heredity}(i),
whose hypothesis is the comparability \eqref{4.11:eq-radii-comparable-a} assumed in~(a).
\item For the one-power predicate $\mathsf{D}^{\Sigma}_i$ it is Lemma~\ref{4.11:lem-restriction-heredity}(ii). That part needs no hypothesis.
\end{itemize}
Together with the antecedent, this places the restricted pair in the $U$-augmented target class,
with the restricted $G$-valued factor.

\emph{Twist.} Suppose $\Phi$ is the twist of a representative $(U',U,\mathbf{J})$ of
$\mathfrak{U}^{\mathrm{post},U}_{k+1}(Y)$, followed by restriction to $X\subset Y$, read at a level
$1\le j\le k$. The twisted pair carries the witness $(e^{i\mathsf{H}}U',U)$, so its $G$-valued factor is again $U$. Lemma~\ref{4.11:lem-twist-heredity} then gives two conclusions.
\begin{itemize}
\item The predicate $\mathsf{D}_j(X;\varphi^{\mathrm{st}}_j)[U|_X]$ holds.
\item When $X\subset\Omega_j$, the predicate $\mathsf{D}^{\Sigma}_j(X;\alpha_{0,j})[U|_X]$ holds as well.
\end{itemize}
The hypotheses of that lemma are the comparability \eqref{4.11:eq-radii-comparable-a} assumed
in~(a) and, where the rule sets differ, hypothesis~(d). The first conclusion is the clause on the
$G$-valued factor of the target class $\tilde{U}^{c,U}_j(X)$; the second is the $U$-half of the
derived-slot member of the one-power class $U^{\Sigma}_j(X;\alpha_{0,j},\cdot)$. Hence the twisted,
restricted pair lies in the $U$-augmented target class.

\emph{Substitution.} Suppose $\Phi=\Phi_{\mathrm{sub}}$, in the setting of
Proposition~\ref{4.11:prop-substitution-heredity} (its fixed triple $(\mathcal{T},t,Z)$, its source
and target classes and its standing conditions): $X$ is
the target region and $a_0'$ the target constant. The image of a representative $(U',U,\mathbf{J})$ of $\mathcal{C}^{U}_{\mathrm{src}}$ carries its natural witness
\begin{equation*}
\bigl(e^{i\mathsf{H}_{\mathrm{sub}}},\,U_t^{\mathcal{T}}(Z;\dot{\mathcal{M}}U)\bigr).
\end{equation*}
The existence of this witness rests on the four hypotheses listed in~(e) and on the ceilings on
the coupling assumed in~(a). The hypothesis of
Proposition~\ref{4.11:prop-substitution-heredity} is hypothesis~(b); under~(b) it shows that this
natural witness also satisfies the $U$-half, that is, the derived-slot predicate evaluated at the
$G$-valued factor of the witness:
\begin{equation*}
\mathsf{D}^{\Sigma}_t(X;a_0')\bigl[U_t^{\mathcal{T}}(Z;\dot{\mathcal{M}}U)\bigr].
\end{equation*}
Hence $\Phi_{\mathrm{sub}}$ maps $\mathcal{C}^{U}_{\mathrm{src}}$ into the $U$-augmented target class.

\emph{Composites of the large-field chain.} Suppose $\Phi$ is such a composite $\mathcal{W}$. The
hypotheses of Lemma~\ref{4.11:lem-composite-heredity} are the two statements in~(c); under them,
for every consumer domain $X'$ outside the cut band, the clause on the $G$-valued factor of the
corresponding $U$-augmented class holds for $\mathcal{W}|_{X'}$. The witness's $G$-valued factor is $U^{\natural}$.

These four cases exhaust the production maps. This proves~(iii).

Consequently the clause on the $G$-valued factor is a declared clause of the analyticity domain $\mathfrak{U}^{\mathrm{post},U}_{k+1}(Y)$, and the step reproduces it under hypotheses (a)--(e).
\end{proof}

\begin{remark}
The following ingredients use neither the comparability of the radii nor the floor assumed
in~(a), nor hypotheses (b)--(d); they use at most the published statements listed in~(e) and, for
the existence of the natural witness of a substituted pair, the ceilings on the coupling assumed
in~(a): parts (i) and (ii); the restriction of the
one-power predicate;
the existence of the natural witness of a substituted pair; the identification of the member of the
target class that, for the natural witness's $G$-valued factor, bounds
\begin{equation*}
\bigl|\partial U_t^{\mathcal{T}}(Z;\dot{\mathcal{M}}U)-1\bigr| \quad\text{on } X
\end{equation*}
with the corresponding member of the source's clause at the same index; the transfer lemma on
which the proof of Proposition~\ref{4.11:prop-substitution-heredity} rests; and the composition
identity
\begin{equation*}
U_t^{\mathcal{T}}\bigl(Z;\dot{\mathcal{M}}(U_t^{\mathcal{T}}(Z;\dot{\mathcal{M}}U))\bigr)
=U_t^{\mathcal{T}}(Z;\dot{\mathcal{M}}U),
\end{equation*}
valid for every rule $\mathcal{T}$, region $Z$, index $t$ and $G$-valued $U$ with admissible data,
which follows from the definition of the constrained minimiser.

The argument nowhere compares the constrained minimisers of two different towers.
\end{remark}

\begin{proposition}[The clause up to a constant factor inside one layer]
\label{4.11:prop-one-layer-clause}
Assume the following two statements.
\begin{enumerate}
\item[(a)] \textup{(}\cite[Proposition~2]{B-CMP98}\textup{)} There is a threshold
$c_2=c_2(d,L)>0$ such that every $G$-valued configuration $V$ with
$\sup\lvert\partial V-1\rvert<c_2$ satisfies, on every plaquette of $q$-blocks and in fine
units,
\begin{equation*}
\lvert\partial(M^{q}V)-1\rvert\le 2L^{2q}\sup\lvert\partial V-1\rvert .
\end{equation*}
\item[(b)] \textup{(}\cite[Theorem~1]{B-CMP102b}, in the form in which it is read in
\cite{B-CMP109}\textup{)} There are constants $a_1$ and $B_3$, depending only on $d$ and
$L$, with the following property. Let $t$ be the top level of a constrained minimisation
problem $P_{\max}(Z,t;D)$, and let $\xi>0$, $\alpha'>0$. Suppose that the data
$D=\dot{\mathcal{M}}V$ are the fixed bonds and block averages of a $G$-valued configuration
$V$ on its determining set, and that they obey the regularity condition
\cite[(7)]{B-CMP102b} in the form
that the $q$-block data satisfy $\lvert\partial(M^{q}V)-1\rvert\le 2\alpha'(L^{q}\xi)^2$ at
every level $q\le t$, and that $\varepsilon_1:=2\alpha'(L^{t}\xi)^2\le a_1$. Then the
constrained Wilson action has a minimal orbit, whose element in the axial gauge is the
minimal configuration $U_t^{\max}(Z;D)$, with current $J_t^{\max}(Z;D)$ in the level-$t$
normalisation; and at every point at the top block scale $L^{t}$ of the tower that is covered
by one of the cubes on which the regularity estimates \cite[(9)--(10)]{B-CMP102b} hold
(a \emph{regularity cube}) one has
\begin{equation*}
\lvert\partial U_t^{\max}(Z;D)-1\rvert<2B_3\alpha'\xi^2,\qquad
\lvert J_t^{\max}(Z;D)\rvert<2B_3\alpha'(L^{t}\xi)^2 .
\end{equation*}
\end{enumerate}
Assume the floor \eqref{4.11:eq-floors-ceilings-b} of
Definition~\ref{4.11:not-floors-ceilings},
\begin{equation*}
M\ge 2(R_1+\sqrt{d}\,L)M_1,
\end{equation*}
where $R_1$ is the constant of the maximal tower rule of
Definition~\ref{4.11:def-tower-rules}, and the ceiling \eqref{4.11:eq-floors-ceilings-c},
\begin{equation*}
2(1+\beta)\alpha_{0,n}\le\min(a_1,c_2)\qquad\text{for all } n,
\end{equation*}
with $a_1$ from \textup{(b)} and $c_2$ from \textup{(a)}. Let $U$ be $G$-valued on $Y$ and
satisfy the curvature bound $\lvert\partial U-1\rvert<\varphi^{\mathrm{post}}(n)\alpha_{0,n}L^{-2n}$
at points of layer index $n$, the first inequality of the class condition
\eqref{4.11:eq-configuration-classes-b}
of Definition~\ref{4.11:def-configuration-classes} at $(k+1,Y;\varphi^{\mathrm{post}})$. Let
$(t,p,Z)\in\mathfrak{T}_{k+1}(Y)$ with $Z$ contained in one layer, that is, either
$Z\subset\mathfrak{L}_n$ with $t\le n\le k$, or $Z\subset\Omega_{k+1}$, in which case we put
$n:=k+1$; and let $\mathcal{T}=\max$. Then the tower quantities of $U$ at $(\max,t,Z)$ are
defined, and at every $x\in P(p,Z)$ that is covered by a regularity cube of \textup{(b)} at the
top level $t$,
\begin{equation}\label{4.11:eq-one-layer-clause-1}
\begin{aligned}
\lvert\partial U_t^{\max}(Z;\dot{\mathcal{M}}U)-1\rvert
&<2B_3\,\varphi^{\mathrm{post}}(n)\,\alpha_{0,n}L^{-2n},\\
\lvert J_t^{\max}(Z;\dot{\mathcal{M}}U)\rvert
&<2B_3\,\varphi^{\mathrm{post}}(n)\,\alpha_{0,n}L^{-2(n-t)}.
\end{aligned}
\end{equation}
That is, \eqref{4.11:eq-tower-bounds-a} and \eqref{4.11:eq-tower-bounds-b} hold for $U$ with
the factor $2B_3\varphi^{\mathrm{post}}(n)$ in place of $\varphi^{\mathrm{post}}(n)$.
\end{proposition}

\begin{proof}
Distances are measured in fine units and with the sup-distance. We write $\Omega_q^{\max}(Z)$
for the tower region of level $q$ of $Z$ under the maximal rule.

\emph{Step 1: the reading set lies at the top of the tower.} For $q\ge1$ put
\begin{equation*}
\delta_q:=\sup\{\operatorname{dist}(x,Z^{c}) : x\in Z\setminus\Omega_q^{\max}(Z)\},
\end{equation*}
and $\delta_0=0$. Let $q\ge1$ and $x\in Z\setminus\Omega_q^{\max}(Z)$, and let $S$ be the
$q$-superblock containing $x$. By the maximality property of the maximal rule
(Definition~\ref{4.11:def-tower-rules}; \cite[(1.3)--(1.6)]{B-CMP99a}),
$\operatorname{dist}(S,\Omega_{q-1}^{\max}(Z)^{c})\le R_1M_1L^{q-1}$. Hence there is a point
$y\notin\Omega_{q-1}^{\max}(Z)$ with
\begin{equation*}
\lvert x-y\rvert\le R_1M_1L^{q-1}+\sqrt{d}\,M_1L^{q},
\end{equation*}
the second term bounding the diameter of $S$. Since
$\operatorname{dist}(y,Z^{c})\le\delta_{q-1}$, the triangle inequality gives
\begin{equation*}
\operatorname{dist}(x,Z^{c})\le\delta_{q-1}+R_1M_1L^{q-1}+\sqrt{d}\,M_1L^{q}
=\delta_{q-1}+(R_1+\sqrt{d}\,L)M_1L^{q-1},
\end{equation*}
and taking the supremum over $x$ bounds $\delta_q$ by the right-hand side. Summing these
inequalities over $q=1,\dots,t$ and using $\delta_0=0$,
\begin{equation*}
\delta_t\le(R_1+\sqrt{d}\,L)M_1\sum_{q=1}^{t}L^{q-1}
\le\frac{(R_1+\sqrt{d}\,L)M_1L^{t}}{L-1}
\le 2(R_1+\sqrt{d}\,L)M_1L^{t-1},
\end{equation*}
the middle step because $\sum_{q=1}^{t}L^{q-1}=(L^{t}-1)/(L-1)$, the last because
$L/(L-1)\le2$ for $L\ge2$. The reading set $P(p,Z)=Z^{\wr,p-1}$ consists of points at
distance at least $2ML^{p-1}$ from $Z^{c}$. By the floor \eqref{4.11:eq-floors-ceilings-b}
and $t\le p$,
\begin{equation*}
2ML^{p-1}\ge 4(R_1+\sqrt{d}\,L)M_1L^{p-1}>2(R_1+\sqrt{d}\,L)M_1L^{t-1}\ge\delta_t .
\end{equation*}
Hence no point of $P(p,Z)$ lies in $Z\setminus\Omega_t^{\max}(Z)$, so
$P(p,Z)\subset\Omega_t^{\max}(Z)$: every point of the reading set is at the top of the tower
of $Z$, at block scale $L^{t}$.

\emph{Step 2: the data.} Since $Z$ lies in one layer, with index $n$, the first inequality of
the class condition \eqref{4.11:eq-configuration-classes-b} says that on $Z$
\begin{equation*}
\lvert\partial U-1\rvert<\varphi^{\mathrm{post}}(n)\,\alpha_{0,n}L^{-2n}.
\end{equation*}
The threshold condition of \textup{(a)} is supplied by the ceiling
\eqref{4.11:eq-floors-ceilings-c}. Applying \textup{(a)}, a $q$-block datum with
$q\le t\le n$ satisfies
\begin{equation*}
\begin{aligned}
\lvert\partial(M^{q}U)-1\rvert
&\le 2\varphi^{\mathrm{post}}(n)\,\alpha_{0,n}L^{2(q-n)}
\le\varepsilon_1:=2\varphi^{\mathrm{post}}(n)\,\alpha_{0,n}L^{2(t-n)}\\
&\le 2(1+\beta)\alpha_{0,n}\le a_1,
\end{aligned}
\end{equation*}
the second inequality because $q\le t$, the last by the ceiling
\eqref{4.11:eq-floors-ceilings-c}.

\emph{Step 3: the minimiser.} We apply \textup{(b)}, in the form in which \cite{B-CMP109}
reads it, with the following identification of its quantities with ours: its fine spacing
$\xi$ is $L^{-n}$, its radius $\alpha'$ is $\varphi^{\mathrm{post}}(n)\alpha_{0,n}$, its class
level is $n$ and its tower index is $t$. Then
$2\alpha'(L^{q}\xi)^2=2\varphi^{\mathrm{post}}(n)\alpha_{0,n}L^{2(q-n)}$ is the bound on the
$q$-block data obtained in Step~2, and $2\alpha'(L^{t}\xi)^2=\varepsilon_1\le a_1$. Hence the
constrained problem has a minimal orbit and the tower quantities of $U$ at $(\max,t,Z)$ are
defined. Let $x\in P(p,Z)$ be covered by a regularity cube of \textup{(b)} (this is the
restriction written ``on $X^{\sim 2}$'' in \cite{B-CMP109}); by Step~1, $x$ is at the top block
scale $L^{t}$. The conclusion of \textup{(b)} gives at $x$
\begin{equation*}
\lvert\partial U_t^{\max}(Z;\dot{\mathcal{M}}U)-1\rvert<2B_3\alpha'\xi^2
=B_3\varepsilon_1L^{-2t}=2B_3\,\varphi^{\mathrm{post}}(n)\,\alpha_{0,n}L^{-2n},
\end{equation*}
and
\begin{equation*}
\begin{aligned}
\lvert J_t^{\max}(Z;\dot{\mathcal{M}}U)\rvert
&<2B_3\alpha'(L^{t}\xi)^2=2B_3\,\varphi^{\mathrm{post}}(n)\,\alpha_{0,n}L^{2(t-n)}\\
&=2B_3\,\varphi^{\mathrm{post}}(n)\,\alpha_{0,n}L^{\,n-\min(t,n)}L^{-3(n-t)},
\end{aligned}
\end{equation*}
where the last equality holds because $t\le n$, so that $\min(t,n)=t$ and
$L^{n-t}L^{-3(n-t)}=L^{-2(n-t)}=L^{2(t-n)}$. These are the two bounds
\eqref{4.11:eq-one-layer-clause-1}.
\end{proof}

For $Z$ meeting several layers the same argument gives no uniform bound. The regularity
condition \cite[(7)]{B-CMP102b} in \textup{(b)} is a supremum over the whole determining set,
and a datum on a shallower layer $n'<n$ contributes
$2\varphi^{\mathrm{post}}(n')\alpha_{0,n'}L^{2(\min(t,n')-n')}$, which, compared with the
target $\varphi^{\mathrm{post}}(n)\alpha_{0,n}L^{-2n}$ of the deep layer at block scale
$L^{\min(t,n)}$, loses the factor $L^{2(n-n')}$ when $t\ge n>n'$. Only a local form of the
regularity estimate, the input of the proposition that follows, follows the layers.

\begin{proposition}[The clause with a perturbative correction. Stated; proof
incomplete at the step marked $(\ast)$]\label{4.11:prop-perturbative-clause}
Assume \eqref{4.11:eq-floors-ceilings-c} and the ceiling \eqref{4.11:eq-floors-ceilings-d}, namely
\begin{equation*}
c_Q(1+\beta)\alpha_{1,n}\le b_H\qquad\text{for all } n,
\end{equation*}
where $b_H$ is the radius of the polydisc on which the axial-gauge minimal configuration extends
analytically to $G^{c}$-valued data \cite[Proposition~9]{B-CMP102b}. Assume also the response
bounds at real minimiser bases, in the layered form whose version without the layer structure is
\cite[(190)]{B-CMP102b}. Let $Y$ be a domain at level
$k+1$, let $(U',U,\mathbf{J})$ be a
representative of
$\mathfrak{U}^{\mathrm{post}}_{k+1}(Y)$, let $(t,p,Z)\in\mathfrak{T}_{k+1}(Y)$, let $\mathcal{T}$ be
the tower rule of \cite{B-CMP119}, let $x\in P(p,Z)\cap\Omega_1$, and put $n:=n_{k+1}(x)$. Then the
tower quantity $U_t^{\mathcal{T}}(Z;\dot{\mathcal{M}}U)$ of $U$ at $(\mathcal{T},t,Z)$ is defined,
and there is a number $\kappa_{\mathrm{pc}}$, depending only on $C_H$, $C_u$, $c_Q$, $d$ and $N$,
such that for every plaquette $q$ at $x$
\begin{equation}\label{4.11:eq-perturbative-clause-1}
\bigl|\partial U_t^{\mathcal{T}}(Z;\dot{\mathcal{M}}U)(q)-1\bigr|
<\Bigl[\varphi^{\mathrm{post}}(n)\bigl(1+\kappa_{\mathrm{pc}}\,\alpha_1^{\ast}(x)\bigr)
+\frac{\kappa_{\mathrm{pc}}\,\alpha_1^{\ast}(x)}{\alpha_{0,n}}\Bigr]\alpha_{0,n}L^{-2n},
\end{equation}
where
\begin{equation}\label{4.11:eq-perturbative-clause-2}
\alpha_1^{\ast}(x):=\sup_{y'}\varphi^{\mathrm{post}}\bigl(n(y')\bigr)\,
\alpha_{1,n(y')}\,e^{-\delta_H \operatorname{dist}_H(x,y')},
\end{equation}
the supremum running over the data blocks $y'$ of $(\mathcal{T},t,Z)$, $n(y')$ being the layer
index of the data block $y'$ and $\operatorname{dist}_H$ the scaled block distance in which the
response bounds at real minimiser bases are stated; here $C_H$, $C_u$, $\delta_H$ are the constants
of the response bounds at real minimiser bases, and $c_Q=c_Q(d,L)$ is the constant of
\cite[Propositions~6--7]{B-CMP98}. The analogous statement holds for the current
$J_t^{\mathcal{T}}(Z;\dot{\mathcal{M}}U)$, with $L^{-2n}$ replaced by the right member of
\eqref{4.11:eq-tower-bounds-b}.
\end{proposition}

\begin{proof}
\emph{The objects.} We use the notation of the proof of Proposition~\ref{4.11:prop-natural-witness}.
The construction carried out there uses only the following inputs: the first bound on the class of
$U$ and the bound on the twist potential $\mathsf{A}'$ among the bounds on the configuration classes
\eqref{4.11:eq-configuration-classes-b}; the factorisation of block averages of twisted
configurations, \cite[Propositions~6--7]{B-CMP98}; the plaquette bound for block averages near the
identity; the analytic extension of the axial-gauge minimal configuration
\cite[Proposition~9]{B-CMP102b}; and
\eqref{4.11:eq-floors-ceilings-c}, \eqref{4.11:eq-floors-ceilings-d}. It does not use the clause of
the source class asserting that the tower quantities of $U$ are defined. It therefore applies to the
present representative, and it yields the following facts.

First, $W_U:=U_t^{\mathcal{T}}(Z;\dot{\mathcal{M}}U)$ exists, and it is the $G$-valued axial-gauge
minimiser with data $D_{\mathcal{T}}^{t,Z}(U)$. This is the first assertion of the proposition.

Second, write $\mathbf{U}:=U'U$. Then
\begin{equation*}
U_t^{\mathcal{T}}(Z;\dot{\mathcal{M}}\mathbf{U})=\bigl[e^{i\mathsf{H}_{\mathrm{sub}}}W_U\bigr]^{\Xi},
\qquad \mathsf{H}_{\mathrm{sub}}=\eta\,\mathcal{H}(B),\qquad \Xi:=v\,\tilde{u}.
\end{equation*}
Here $B$ is the $\mathfrak{g}^{c}$-valued twist of the data produced by
the factorisation of block averages, with $B:=\mathsf{A}'$ on the fixed bonds; $\tilde{u}$ is the
block-constant gauge induced by the block gauge of that factorisation through
\cite[(181)]{B-CMP102b}; $v$ is the change from the Landau gauge to the axial gauge; and
$\eta=L^{-t}$.

The factor $\Xi$ is $G^{c}$-valued. By the gauge-factor clause of the response bounds at real
minimiser bases, assumed in the statement,
$|\log\Xi|\le C_u\sup|B|$ near the point $x$. At a data block $y'$ the twist
obeys
\begin{equation*}
|B(y')|\le b(y'):=c_Q\,\varphi^{\mathrm{post}}\bigl(n(y')\bigr)\,\alpha_{1,n(y')}.
\end{equation*}

\emph{Size of the response.} We apply the response bounds at real minimiser bases, assumed in the
statement, at the real base
$W_U$. In units of layer $n$ at $x$ they give
\begin{equation*}
L^{n}|\mathsf{H}_{\mathrm{sub}}|\le\varepsilon(x),\qquad
L^{2n}|\nabla_{W_U}\mathsf{H}_{\mathrm{sub}}|\le\varepsilon(x),
\end{equation*}
where
\begin{equation*}
\varepsilon(x):=C_H\sup_{y'}b(y')\,e^{-\delta_H \operatorname{dist}_H(x,y')}=C_Hc_Q\,\alpha_1^{\ast}(x).
\end{equation*}
The last equality is the definition of $b(y')$ inserted into \eqref{4.11:eq-perturbative-clause-2}.

\emph{The twist of a plaquette.} Let $q$ be a plaquette at $x$. We transport the four
exponentials of $e^{i\mathsf{H}_{\mathrm{sub}}}$ along $\partial q$ to the base corner of $q$ and
combine them by the Baker--Campbell--Hausdorff formula. This gives
\begin{equation*}
\partial\bigl(e^{i\mathsf{H}_{\mathrm{sub}}}W_U\bigr)(q)
=\exp\bigl(i(\operatorname{curl}_{W_U}\mathsf{H}_{\mathrm{sub}})(q)+\varsigma(q)\bigr)\,
\partial W_U(q).
\end{equation*}
Here $\operatorname{curl}_{W_U}\mathsf{H}_{\mathrm{sub}}$ is the lattice curl of
$\mathsf{H}_{\mathrm{sub}}$ around $q$ taken covariantly with respect to $W_U$. It satisfies
\begin{equation*}
|(\operatorname{curl}_{W_U}\mathsf{H}_{\mathrm{sub}})(q)|
\le 2a\,|\nabla_{W_U}\mathsf{H}_{\mathrm{sub}}|,
\end{equation*}
where $a$ is the spacing of the lattice on which the plaquette $q$ lies, in the units in force.
The remainder $\varsigma(q)$ is bounded at this point only in the incomplete form
\begin{equation*}
|\varsigma(q)|\le c_{\mathrm{BCH}}\bigl(L^{-n}\cdot L^{n}|\mathsf{H}_{\mathrm{sub}}|\bigr)^{2}
\cdot(\cdots),
\qquad(\ast)
\end{equation*}
with $c_{\mathrm{BCH}}=c_{\mathrm{BCH}}(d,N)$ and with the last factor not determined. This is the
step marked $(\ast)$; it comprises the bound on $\varsigma(q)$ and the determination of the
constant $\kappa_{\mathrm{pc}}$ below. In fine units
this gives
\begin{equation*}
\bigl|\partial\bigl(e^{i\mathsf{H}_{\mathrm{sub}}}W_U\bigr)(q)-\partial W_U(q)\bigr|
\le\bigl(2\varepsilon+c_{\mathrm{BCH}}\varepsilon^{2}\bigr)L^{-2n}
\bigl(1+|\partial W_U(q)-1|\bigr).
\end{equation*}
This is the plaquette bound for a configuration multiplied by an exponential
$e^{i\mathsf{H}_{\mathrm{sub}}}$, with its size parameter equal to
$\varepsilon=\varepsilon(x)$; being derived from the preceding expansion, it rests on the step
$(\ast)$.

Next, for a configuration $V$ on which $\Xi$ acts, the gauge transformation by $\Xi$ acts on
the holonomy of $q$ by conjugation with the value
of $\Xi$ at the base corner $z$ of $q$. Thus $\partial([V]^{\Xi})(q)-1$ equals
$\Xi(z)\bigl(\partial V(q)-1\bigr)\Xi(z)^{-1}$. Writing conversely
\begin{equation*}
\partial V(q)-1=\Xi(z)^{-1}\bigl(\partial([V]^{\Xi})(q)-1\bigr)\Xi(z)
\end{equation*}
and estimating the two conjugating factors by $e^{|\log\Xi(z)|}$ each, we obtain
\begin{equation*}
\bigl|\partial([\,\cdot\,]^{\Xi})-1\bigr|\ge e^{-2|\log\Xi|}\,\bigl|\partial(\,\cdot\,)-1\bigr|.
\end{equation*}

We now write
\begin{equation*}
\partial W_U(q)-1=\bigl(\partial(e^{i\mathsf{H}_{\mathrm{sub}}}W_U)(q)-1\bigr)
+\bigl(\partial W_U(q)-\partial(e^{i\mathsf{H}_{\mathrm{sub}}}W_U)(q)\bigr).
\end{equation*}
The first term is bounded by the conjugation inequality with $V=e^{i\mathsf{H}_{\mathrm{sub}}}W_U$,
using $[e^{i\mathsf{H}_{\mathrm{sub}}}W_U]^{\Xi}=U_t^{\mathcal{T}}(Z;\dot{\mathcal{M}}\mathbf{U})$,
together with the bound on $|\log\Xi|$ above. The second term is bounded by the bound on
$|\partial(e^{i\mathsf{H}_{\mathrm{sub}}}W_U)(q)-\partial W_U(q)|$ obtained above. Together they
give
\begin{equation}\label{4.11:eq-perturbative-clause-3}
\begin{split}
|\partial W_U(q)-1|
&\le e^{2C_uc_Q(1+\beta)\alpha_1^{\ast}}\,
\bigl|\partial U_t^{\mathcal{T}}(Z;\dot{\mathcal{M}}\mathbf{U})(q)-1\bigr|\\
&\quad+\bigl(2\varepsilon+c_{\mathrm{BCH}}\varepsilon^{2}\bigr)L^{-2n}
\bigl(1+|\partial W_U(q)-1|\bigr).
\end{split}
\end{equation}

The first term on the right of \eqref{4.11:eq-perturbative-clause-3} is estimated by the bound
\eqref{4.11:eq-tower-bounds-a} for the tower quantity of $\mathbf{U}$ at $x$. Its exponential factor
is estimated by the inequality $e^{y}\le1+y\,e^{y}$ for $y\ge0$, applied with
$y=2C_uc_Q(1+\beta)\alpha_1^{\ast}(x)$.

For the second term we use $\varepsilon\le1$, and $|\partial W_U(q)-1|\le1$, which follows from
\eqref{4.11:eq-floors-ceilings-c}. They give
\begin{equation*}
\bigl(2\varepsilon+c_{\mathrm{BCH}}\varepsilon^{2}\bigr)\bigl(1+|\partial W_U(q)-1|\bigr)
\le2\bigl(2+c_{\mathrm{BCH}}\bigr)C_Hc_Q\,\alpha_1^{\ast}(x)
\le4C_Hc_Q\bigl(1+c_{\mathrm{BCH}}\bigr)\alpha_1^{\ast}(x).
\end{equation*}
Moreover
$\kappa_{\mathrm{pc}}\,\alpha_1^{\ast}L^{-2n}
=(\kappa_{\mathrm{pc}}\,\alpha_1^{\ast}/\alpha_{0,n})\,\alpha_{0,n}L^{-2n}$.

Combining these estimates gives
\begin{equation*}
|\partial W_U(q)-1|
<\Bigl[\varphi^{\mathrm{post}}(n)\bigl(1+\kappa_{\mathrm{pc}}\,\alpha_1^{\ast}\bigr)
+\frac{\kappa_{\mathrm{pc}}\,\alpha_1^{\ast}}{\alpha_{0,n}}\Bigr]\alpha_{0,n}L^{-2n},
\end{equation*}
where $\kappa_{\mathrm{pc}}$ is the maximum of two numbers. The second number is
$4C_Hc_Q(1+c_{\mathrm{BCH}})$.
The first is $e^{2C_uc_Q(1+\beta)}\cdot2C_uc_Q(1+\beta)$ multiplied by a supremum of the radii
$\alpha_{1,\cdot}$ whose normalisation the argument does not fix. Thus, as far as the argument
carries it, $\kappa_{\mathrm{pc}}$ depends on $C_H$, $C_u$, $c_Q$, $c_{\mathrm{BCH}}(d,N)$ and
$\beta$ only; fixing that normalisation, and removing the dependence on $\beta$, which the statement
asserts to be absent, belong to the step $(\ast)$. Since
$W_U=U_t^{\mathcal{T}}(Z;\dot{\mathcal{M}}U)$, this is \eqref{4.11:eq-perturbative-clause-1} with
this $\kappa_{\mathrm{pc}}$.

\emph{The current.} The member of the proposition concerning the current
$J_t^{\mathcal{T}}(Z;\dot{\mathcal{M}}U)$ is obtained by the same computation, with
\eqref{4.11:eq-tower-bounds-b} in place of \eqref{4.11:eq-tower-bounds-a}. For the explicit
remainder of the current the computation uses Lemma~\ref{4.4:lem-current-remainder}, the quadratic
remainder in the current and its Lipschitz bound;
this use belongs to the step marked $(\ast)$. This member
rests on the step $(\ast)$
of the plaquette computation as well, and is obtained only as far as that step carries it.
\end{proof}

\begin{remark}[Neither weakening yields the clause]\label{4.11:rem-weakenings-insufficient}
Neither of the following yields the constrained-minimiser clause at straddling sources,
Proposition~\ref{4.11:prop-straddling-clause} (whose proof is incomplete at one step), in the
form in which that clause is used:
\begin{itemize}
\item the clause up to a constant factor inside one layer,
Proposition~\ref{4.11:prop-one-layer-clause};
\item the clause with a perturbative correction, Proposition~\ref{4.11:prop-perturbative-clause}
(whose proof is incomplete at one step).
\end{itemize}
More precisely:
\begin{enumerate}
\item[(a)] The use of the clause reads the curvature bound for the $G$-valued factor of the
substituted pair at the target class's own radius. A factor $2B_3$, as in
Proposition~\ref{4.11:prop-one-layer-clause}, is a growth of constants across widths, and that
growth is precisely what the one-power hereditary class is built to avoid.
\item[(b)] At every layer $n\le k$, the layer $n=k$ included, the correction in
Proposition~\ref{4.11:prop-perturbative-clause} exceeds the whole room that the post schedule
$\varphi^{\mathrm{post}}$ leaves above the standard factor. The bound of
Proposition~\ref{4.11:prop-perturbative-clause} implies the clause of
Proposition~\ref{4.11:prop-straddling-clause} at a point $x$ only if
$\kappa_{pc}\alpha_1^{\ast}(x)\le 0$.
\end{enumerate}
\end{remark}

\begin{proof}
\emph{Part (a).} The target's curvature bound is read from the clause of
Proposition~\ref{4.11:prop-straddling-clause}, as explained in
Lemma~\ref{4.11:lem-curvature-member}. That lemma identifies the member of the target class
which bounds $|\partial U_t^{\mathcal{T}}(Z;\dot{\mathcal{M}}U)-1|$ for the $G$-valued factor
with the curvature member of the source's clause. The two members are compared at the target
class's own radius, and no constant factor is allowed in the comparison.

The bound of Proposition~\ref{4.11:prop-one-layer-clause} carries the factor $2B_3$ in front of
$\varphi^{\mathrm{post}}$. This factor cannot be absorbed at the target's own radius. Absorbing it
would require letting the constants grow from one width to the next, and the one-power
hereditary class was introduced to exclude exactly that growth. Hence
Proposition~\ref{4.11:prop-one-layer-clause} does not supply the member that
Lemma~\ref{4.11:lem-curvature-member} reads.

\emph{The room.} Recall the schedules. For $n\le i$,
\begin{equation*}
f_{i,n}=1-\beta\bigl(1-2^{-(i-n)}\bigr).
\end{equation*}
For $n\le k$,
\begin{equation*}
\varphi^{\mathrm{post}}(n)=F_n=f_{k+1,n}+\theta\beta 2^{-(k+1-n)},
\end{equation*}
and $\varphi^{\mathrm{post}}(k+1)=1+\beta$, where $\beta,\theta\in(0,1)$. On layer $n\le k$ the
room above the standard factor $f_{k+1,n}$ is therefore
\begin{equation}\label{4.11:eq-weakenings-insufficient-1}
F_n-f_{k+1,n}=\theta\beta\,2^{-(k+1-n)}\le\theta\beta .
\end{equation}
The whole factor is at most $1+\beta$. Indeed $f_{k+1,n}\le 1$ and $\theta\beta<\beta$ for
$n\le k$, while $F_{k+1}=1+\beta$.

\emph{Part (b): the correction.} Fix $x$ with $n=n_{k+1}(x)$. The bound of
Proposition~\ref{4.11:prop-perturbative-clause} has the form of the clause of
Proposition~\ref{4.11:prop-straddling-clause} at $x$, with the factor $\varphi^{\mathrm{post}}(n)$
in front of $\alpha_{0,n}L^{-2n}$ replaced by the bracket
$\varphi^{\mathrm{post}}(n)(1+\kappa_{pc}\alpha_1^{\ast}(x))+\kappa_{pc}\alpha_1^{\ast}(x)/\alpha_{0,n}$.
Subtracting the factor from the bracket gives
\begin{equation}\label{4.11:eq-weakenings-insufficient-2}
\Bigl[\varphi^{\mathrm{post}}(n)\bigl(1+\kappa_{pc}\alpha_1^{\ast}(x)\bigr)
+\frac{\kappa_{pc}\alpha_1^{\ast}(x)}{\alpha_{0,n}}\Bigr]-\varphi^{\mathrm{post}}(n)
=\kappa_{pc}\alpha_1^{\ast}(x)\Bigl(\varphi^{\mathrm{post}}(n)+\frac{1}{\alpha_{0,n}}\Bigr).
\end{equation}
The second factor on the right is positive. Hence the bracket is at most
$\varphi^{\mathrm{post}}(n)$ if and only if $\kappa_{pc}\alpha_1^{\ast}(x)\le 0$. This proves the
last assertion of (b).

Now call the correction the term $\kappa_{pc}\alpha_1^{\ast}(x)/\alpha_{0,n}$ of the bracket.
Recall that $\alpha_1^{\ast}(x)$ is the supremum over data blocks $y'$ of
$\varphi^{\mathrm{post}}(n(y'))\,\alpha_{1,n(y')}\,e^{-\delta_H\operatorname{dist}_H(x,y')}$.
The correction is therefore at least
\begin{equation*}
\kappa_{pc}\,\inf_{y'}\alpha_{1,n(y')}e^{-\delta_H\operatorname{dist}_H(x,y')}\big/\alpha_{0,n}.
\end{equation*}
When the data twist sits on layer $n$ near $x$, the correction is of order
$\kappa_{pc}\alpha_{1,n}/\alpha_{0,n}$. This configuration is admissible, because
\cite[(2.39)]{B-CMP119} lets the twist potential $\mathsf{A}'$ saturate its radius anywhere.

\emph{Part (b): the order of the radii.} The radii $\alpha_{0,n},\alpha_{1,n}$ are those of
\cite[(2.28)]{B-CMP119}, with $C_1$ large against $C_0$. Moreover \cite{B-CMP109} orders them so
that $\alpha_1$ dominates $\alpha_0$: ``Assume now that
$B_3O(1)M\alpha_0<\tfrac12\alpha_1$''. Let $c$ denote the number written $O(1)$ in this
assumption. Then
\begin{equation*}
\alpha_{1,n}\ge 2B_3cM\,\alpha_{0,n}.
\end{equation*}
Thus $\alpha_{1,n}\gg\alpha_{0,n}2^{-(k+1-n)}$ is not only admissible but forced, and
\begin{equation}\label{4.11:eq-weakenings-insufficient-3}
\frac{\kappa_{pc}\alpha_{1,n}}{\alpha_{0,n}}\ge 2\kappa_{pc}B_3cM\gg\theta\beta\,2^{-(k+1-n)} .
\end{equation}
By \eqref{4.11:eq-weakenings-insufficient-1}, the right-hand side is the room on layer $n$.
Hence the correction exceeds the room at every layer $n\le k$. Since the room itself is at most
$\theta\beta$, this holds at $n=k$ as well. The failure of this route is thus not a loss of
uniformity in $k$: it occurs at a single layer already. This proves the first assertion of (b).

\emph{Why produced pairs are unaffected.} The route through
Proposition~\ref{4.11:prop-perturbative-clause} measures the data twist by the class radius
$\alpha_{1,n}$. For pairs produced by the step, the twist size on layer $n$
(Definition~\ref{4.4:def-post-step-space}) satisfies
\begin{equation*}
h_n\ll\theta\beta\,2^{-(k+1-n)}\alpha_{0,n}.
\end{equation*}
This is why the corresponding clause for pairs produced by the step costs nothing. The clause of
Proposition~\ref{4.11:prop-straddling-clause}, by contrast, is out of reach of this route,
because it quantifies over all members of $\mathfrak{U}^{\mathrm{post}}_{k+1}(Y)$ and not only
over the produced ones.
\end{proof}

\begin{remark}[What is not asserted about the $G$-valued factor]\label{4.11:rem-not-asserted}
The following are not asserted in this section.
\begin{enumerate}
\item The constrained-minimiser clause at straddling sources
(Proposition~\ref{4.11:prop-straddling-clause}, whose proof is incomplete at one step), and
its negation.
\item The clause on the $G$-valued factor in any membership in a target class, that is, the
clause on the $G$-valued factor of the definition of the configuration classes, whose pointwise
conditions are listed in \eqref{4.11:eq-configuration-classes-b}, other than as the conclusion
of the heredity clause of Theorem~\ref{4.11:thm-clause-reproduction}. The four memberships of
the production maps $\Phi$ in target classes without that clause enter the heredity clause of
Theorem~\ref{4.11:thm-clause-reproduction} as hypotheses. They are not asserted in this section,
and the clause on the $G$-valued factor is obtained there only from them.
\item Any comparison of the constrained minimisers
$U_t^{\mathcal{T}}(Z;\dot{\mathcal{M}}\mathbf{V})$ belonging to two different towers, whether
these differ in the tower rule or in the tower regions. Such a
comparison is not needed in the proof of the reproduction of the clause on the $G$-valued
factor (Theorem~\ref{4.11:thm-clause-reproduction}).
\item Any statement concerning the displays \cite[(1.12)]{B-CMP109}, \cite[(2.38)]{B-CMP119}
and \cite[(1.17)]{B-CMP109}, or concerning the inclusions across widths in \cite{B-CMP116};
this section contains no such statement.
\item The existence of the tower quantities
$U_t^{\mathcal{T}}(Z;\dot{\mathcal{M}}\mathbf{V})$ and
$J_t^{\mathcal{T}}(Z;\dot{\mathcal{M}}\mathbf{V})$ at complex data beyond the data at which
membership in the classes already presupposes it.
\item Completeness of the enumeration of Cauchy families in the proof of clause~(ii) of
Theorem~\ref{4.11:thm-clause-reproduction}, beyond the cases listed there.
\end{enumerate}
What is asserted about the blocking factor has the following existence form. There is a
threshold $L_{\mathrm{fl}}=L_{\mathrm{fl}}(\beta,\varrho)$, depending only on $\beta$ and
$\varrho$, such that
Theorem~\ref{4.11:thm-clause-reproduction} holds for every odd $L\ge L_{\mathrm{fl}}$,
under the ceilings \eqref{4.11:eq-floors-ceilings-c} and \eqref{4.11:eq-floors-ceilings-d} on
the coupling. Any positive real number whose square is at least $\sqrt{2}(1+\beta)\varrho$ is
such a threshold.
\end{remark}

\begin{proof}
Items 1--6 list statements that are not made, so they require no argument. We prove the
existence form.

Let $L_{\mathrm{fl}}>0$ satisfy $L_{\mathrm{fl}}^2\ge\sqrt{2}(1+\beta)\varrho$. For instance,
$L_{\mathrm{fl}}=(\sqrt{2}(1+\beta)\varrho)^{1/2}$ is such a number, and it depends only on
$\beta$ and $\varrho$. Let $L$ be odd with $L\ge L_{\mathrm{fl}}$. Since $t\mapsto t^2$ is
increasing on $[0,\infty)$, we have
\begin{equation*}
L^2\ \ge\ L_{\mathrm{fl}}^2\ \ge\ \sqrt{2}(1+\beta)\varrho ,
\end{equation*}
and this is the floor \eqref{4.11:eq-floors-ceilings-a}.

We next check which of the side conditions of Definition~\ref{4.11:not-floors-ceilings} (the
floors on $L$ and $M$ and the ceilings on the coupling) bound $L$ from below.
\begin{itemize}
\item The floor \eqref{4.11:eq-floors-ceilings-a} is the only such condition. It reads only
$\beta$ and $\varrho$, and it enters the proof of
Theorem~\ref{4.11:thm-clause-reproduction} in the restriction and twist cases of its heredity
clause.
\item The ceilings \eqref{4.11:eq-floors-ceilings-c} and \eqref{4.11:eq-floors-ceilings-d} are
upper bounds on the coupling. They are imposed after $L$ has been fixed, and they impose no
lower bound on $L$.
\end{itemize}
Hence, for every odd $L\ge L_{\mathrm{fl}}$ and under \eqref{4.11:eq-floors-ceilings-c} and
\eqref{4.11:eq-floors-ceilings-d}, the floor on $L$ that
Theorem~\ref{4.11:thm-clause-reproduction} uses is satisfied, and that theorem applies.
\end{proof}

\begin{definition}[Lattice, gauge group, configurations and the current]
\label{4.12:not-lattice-configurations}
The following conventions are in force throughout this section.

\emph{(i) Lattice.} The dimension is any $d\ge 2$ in this section (in the rest of the book $d=4$).
The step is $k\ge 1$. Put $\eta:=L^{-k}$, and for $0\le m\le k$
put $\ell_m:=L^m\eta=L^{m-k}$, so that $\ell_k=1$ and $\ell_0=\eta$. Sites are the points of
$\eta\mathbb{Z}^d$. Bonds are ordered pairs $\langle x,y\rangle$ of nearest-neighbour sites.
The unit vectors are $e_1,\dots,e_d$, and $x+\mu$ abbreviates $x+\eta e_\mu$.

\emph{(ii) Group.} The group $G$ is a compact Lie group realised as a closed subgroup of the
unitary group $\mathrm{U}(\mathfrak{n})$. Put
\begin{equation*}
\mathfrak{g}:=\{A:\ A=A^{*},\ e^{itA}\in G\ \text{for all real }t\}.
\end{equation*}
We write $|\cdot|$ for the operator norm on $\mathfrak{n}\times\mathfrak{n}$ matrices. For $g\in G$
and any matrix $A$ we have $|gA|=|Ag|=|A|$ and $|\operatorname{Ad}_gA|=|A|$, where
$\operatorname{Ad}_gA:=gAg^{-1}$. For $u\in G$ with $|u-1|\le 1$ the principal logarithm $\log u$
is defined, $\frac1i\log u$ is hermitian, and
\begin{equation}\label{4.12:eq-lattice-configurations-a}
\Bigl|\tfrac1i\log u\Bigr|\le\frac{\pi}{3}\,|u-1| .
\end{equation}
There is a number $\epsilon_G\in(0,1]$, depending only on $G$, such that
$\frac1i\log u\in\mathfrak{g}$ for every $u\in G$ with $|u-1|\le\epsilon_G$.

\emph{(iii) Configurations and gauge transformations.} A configuration $U$ on a set of bonds
assigns to each bond $\langle x,y\rangle$ an element $U(\langle x,y\rangle)\in G$, with
$U(\langle y,x\rangle)=U(\langle x,y\rangle)^{-1}$. For a path
$\gamma=\langle x_0,x_1,\dots,x_r\rangle$ put
\begin{equation*}
U(\gamma):=U(\langle x_0,x_1\rangle)\,U(\langle x_1,x_2\rangle)\cdots U(\langle x_{r-1},x_r\rangle).
\end{equation*}
A gauge transformation $\sigma$ is a $G$-valued function on sites, and
$U^{\sigma}(\langle x,y\rangle):=\sigma(x)U(\langle x,y\rangle)\sigma(y)^{-1}$. Then
$U^{\sigma}(\gamma)=\sigma(x_0)U(\gamma)\sigma(x_r)^{-1}$ and $(U^{\sigma})^{\tau}=U^{\tau\sigma}$,
with $\tau\sigma$ the pointwise product.

\emph{(iv) Plaquettes.} The plaquette $\mathfrak{p}=\mathfrak{p}(x;\mu,\nu)$, $\mu<\nu$, has corners
$x$, $x+\mu$, $x+\mu+\nu$, $x+\nu$, and
\begin{equation*}
\partial U(\mathfrak{p}):=U(\langle x,x+\mu\rangle)\,U(\langle x+\mu,x+\mu+\nu\rangle)\,
U(\langle x+\mu+\nu,x+\nu\rangle)\,U(\langle x+\nu,x\rangle).
\end{equation*}
One has $\partial U^{\sigma}(\mathfrak{p})=\sigma(x)\,\partial U(\mathfrak{p})\,\sigma(x)^{-1}$, hence
$|\partial U^{\sigma}(\mathfrak{p})-1|=|\partial U(\mathfrak{p})-1|$.

\emph{(v) One-forms and covariant derivatives.} A $\mathfrak{g}$-valued $1$-form $\mathcal{V}$ assigns
$\mathcal{V}_\mu(x)\in\mathfrak{g}$ to each positively oriented bond $\langle x,x+\mu\rangle$. Given a
configuration $Y$, $e^{i\eta\mathcal{V}}Y$ denotes the configuration with
\begin{equation*}
(e^{i\eta\mathcal{V}}Y)(\langle x,x+\mu\rangle):=e^{i\eta\mathcal{V}_\mu(x)}\,Y(\langle x,x+\mu\rangle),
\end{equation*}
and with the inverse of this element on the reversed bond $\langle x+\mu,x\rangle$. The covariant
forward derivative is
\begin{equation*}
(\nabla^Y_\mu\mathcal{V}_\nu)(x):=\eta^{-1}\bigl[\operatorname{Ad}_{Y(\langle x,x+\mu\rangle)}
\mathcal{V}_\nu(x+\mu)-\mathcal{V}_\nu(x)\bigr].
\end{equation*}
For $Y\equiv 1$ it is the ordinary forward difference quotient, written $\nabla$. We put
$|\nabla^Y\mathcal{V}(x)|:=\max_{\mu,\nu}|(\nabla^Y_\mu\mathcal{V}_\nu)(x)|$. The letter $1$ also
denotes the configuration identically equal to the identity.

\emph{(vi) The current.} The current $J(U)$ is the $\mathfrak{g}$-valued quantity attached by
\cite[(1.8)]{B-CMP109} to a configuration $U$ at each bond $\mathfrak{b}$ all of whose neighbouring
plaquettes carry $U$. The quantities $E_3$ and $E_4$ introduced later in this section are both
sizes of this $J$, in the normalisation of the lattice of spacing $\eta$ fixed there. Only two
properties of $J$ are used. Both are assumed in what follows as hypotheses and are not proved in
this section:
\begin{itemize}
\item[(J1)] the $\operatorname{Ad}$-covariance of the current: for every $G$-valued $\sigma$,
\begin{equation}\label{4.12:eq-lattice-configurations-b}
|J(U^{\sigma})(\mathfrak{b})|=|J(U)(\mathfrak{b})| ;
\end{equation}
\item[(J2)] the current lemma at the trivial background, used here by statement (it rests on
\cite[(3.11)--(3.13)]{B-CMP109}): there is a constant $C_J=C_J(d)$ with the following property.
Let $\ell\in[\eta,1]$ and $\varpi\le 1/64$, and let $\mathcal{V}$ be a $\mathfrak{g}$-valued
$1$-form satisfying
\begin{equation*}
\ell\,|\mathcal{V}|\le\varpi,\qquad \ell^2\,|\nabla\mathcal{V}|\le\varpi,\qquad
\ell^3\,|D^{*}D\,\mathcal{V}|\le\varpi
\end{equation*}
at all bonds within two lattice steps of $\mathfrak{b}$. Then
\begin{equation}\label{4.12:eq-lattice-configurations-c}
|J(e^{i\eta\mathcal{V}}1)(\mathfrak{b})|\le(2+2C_J\varpi)\,\varpi\,\ell^{-3}.
\end{equation}
\end{itemize}
Here $D^{*}D$ is the second-order operator of \cite[(3.12)]{B-CMP109} at the trivial background.
It enters only through (J2) and through the hypothesis on the response $1$-form stated below in
this section.
\end{definition}

\begin{proof}
We first prove \eqref{4.12:eq-lattice-configurations-a}. The matrix $u\in G\subset
\mathrm{U}(\mathfrak{n})$
is unitary. Hence it is unitarily diagonalisable, with eigenvalues $e^{i\theta_j}$,
$\theta_j\in(-\pi,\pi]$. Its operator norms are computed on the eigenvalues:
\begin{equation*}
|u-1|=\max_j|e^{i\theta_j}-1|=\max_j 2\sin(|\theta_j|/2).
\end{equation*}
Since $|u-1|\le 1$, every $j$ satisfies $\sin(|\theta_j|/2)\le 1/2$ with
$|\theta_j|/2\in[0,\pi/2]$. Therefore $|\theta_j|/2\le\pi/6$, that is, $|\theta_j|\le\pi/3$. In
particular no eigenvalue equals $-1$, and the principal logarithm is defined. It acts on the same
eigenvectors with eigenvalues $i\theta_j$. Thus $\frac1i\log u$ is hermitian with eigenvalues
$\theta_j$, and $|\frac1i\log u|=\max_j|\theta_j|$.

The sine is concave on $[0,\pi/2]$ and vanishes at $0$. Hence on $[0,\pi/6]$ it lies above the
chord from $0$ to $\pi/6$, that is,
\begin{equation*}
\sin t\ge\frac{\sin(\pi/6)}{\pi/6}\,t=\frac{3}{\pi}\,t \qquad (0\le t\le\pi/6).
\end{equation*}
Taking $t=|\theta_j|/2$ gives $2\sin(|\theta_j|/2)\ge(3/\pi)|\theta_j|$ for every $j$. Consequently
\begin{equation*}
\Bigl|\tfrac1i\log u\Bigr|=\max_j|\theta_j|\le\frac{\pi}{3}\max_j 2\sin(|\theta_j|/2)
=\frac{\pi}{3}\,|u-1|,
\end{equation*}
which is \eqref{4.12:eq-lattice-configurations-a}.

We next prove the membership statement of (ii). By the definition of $\mathfrak{g}$, the set
$i\mathfrak{g}=\{iA:\ A\in\mathfrak{g}\}$ consists of the skew-hermitian matrices $X$ with
$e^{tX}\in G$ for all real $t$. Since $G$ is a closed subgroup of $\mathrm{U}(\mathfrak{n})$, the
closed-subgroup theorem shows that $i\mathfrak{g}$ is the Lie algebra of $G$ and that the
exponential map is a homeomorphism of an open neighbourhood $V_0$ of $0$ in $i\mathfrak{g}$ onto an
open neighbourhood of $1$ in $G$. Put $V:=\{X\in V_0:\ |X|<\pi\}$. It is open in $V_0$ and
contains $0$, so $\exp(V)$ is an open neighbourhood of $1$ in $G$, and there is
$\epsilon_G\in(0,1]$ with $\{u\in G:\ |u-1|\le\epsilon_G\}\subset\exp(V)$. Let $u\in G$ with
$|u-1|\le\epsilon_G$, and write $u=e^{X}$ with $X\in V$. The matrix $X$ is skew-hermitian, hence
unitarily diagonalisable with eigenvalues $i\varphi_j$, and $|X|<\pi$ gives
$\varphi_j\in(-\pi,\pi)$. On the eigenspace of $X$ for $i\varphi_j$ the matrix $u$ acts as
$e^{i\varphi_j}$. Since $\varphi\mapsto e^{i\varphi}$ is injective on $(-\pi,\pi)$, each
eigenspace of $u$ is the eigenspace of $X$ for the unique $\varphi_j$ giving that eigenvalue, and
the principal logarithm acts on it as $i\varphi_j$. Hence $\log u=X$, and
$\frac1i\log u=X/i\in\mathfrak{g}$.

The norm identities in (ii) hold because left or right multiplication by a unitary matrix preserves
the operator norm, and $\operatorname{Ad}_gA=gAg^{-1}$ is such a product.

The identities in (iii) are definitional unwindings. In the product defining $U^{\sigma}(\gamma)$,
each inner factor $\sigma(x_i)^{-1}\sigma(x_i)$, for $0<i<r$, cancels. What remains is
$\sigma(x_0)U(\gamma)\sigma(x_r)^{-1}$. For a bond $\langle x,y\rangle$ we have
\begin{equation*}
(U^{\sigma})^{\tau}(\langle x,y\rangle)=\tau(x)\sigma(x)\,U(\langle x,y\rangle)\,
\sigma(y)^{-1}\tau(y)^{-1}=U^{\tau\sigma}(\langle x,y\rangle).
\end{equation*}

For (iv), $\partial U(\mathfrak{p})=U(\gamma)$ for the closed path
$\gamma=\langle x,x+\mu,x+\mu+\nu,x+\nu,x\rangle$. The path formula of (iii) with $x_0=x_r=x$
therefore gives $\partial U^{\sigma}(\mathfrak{p})=\sigma(x)\,\partial U(\mathfrak{p})\,\sigma(x)^{-1}$.
Writing
\begin{equation*}
\partial U^{\sigma}(\mathfrak{p})-1=\operatorname{Ad}_{\sigma(x)}\bigl(\partial U(\mathfrak{p})-1\bigr),
\end{equation*}
the norm identity of (ii) yields $|\partial U^{\sigma}(\mathfrak{p})-1|=|\partial U(\mathfrak{p})-1|$.

Finally, in (v) the element $e^{i\eta\mathcal{V}_\mu(x)}$ lies in $G$ by the definition of
$\mathfrak{g}$, applied with $t=\eta$. Hence $e^{i\eta\mathcal{V}}Y$ is again a configuration.
\end{proof}

\begin{definition}[Running radii and fixed exponents]\label{4.12:not-running-radii}
Let $k$ be the level of the step under consideration, with lattice spacing $\eta=L^{-k}$ and
running couplings $g_0,\dots,g_k$, and with the remaining notation as in
Notation~\ref{4.12:not-lattice-configurations}.

\emph{Running radii.} For $0\le m\le k$, $\tilde{\alpha}_{0,m}$ denotes the radius of
Ba{\l}aban's construction in \cite{B-CMP119} at level $m$:
\begin{equation}\label{4.12:eq-running-radii-1}
\tilde{\alpha}_{0,m}=g_m\,C_0\,\bigl(\log g_m^{-2}\bigr)^{q_0}.
\end{equation}
Here $C_0$ is the auxiliary parameter of the same name, and $q_0$ is the fixed exponent of
the radii in \cite{B-CMP119}. For $g_m$ sufficiently small this radius is an increasing
function of $g_m$; see the proof. The ratios of these radii obey the bound
\cite[(2.6)--(2.7)]{B-CMP119}: for all levels $m\le m'\le k$,
\begin{equation}\label{4.12:eq-running-radii-a}
\frac{\tilde{\alpha}_{0,m'}}{\tilde{\alpha}_{0,m}}\le(1+\beta_0)\,(1+m'-m)^{1/2}.
\end{equation}

\emph{Fixed exponents.} A number $\beta\in(0,\tfrac14)$ is fixed once and for all in the
standing conditions of this chapter on the auxiliary parameters, and
$\theta_*:=1-4\beta>0$. The constant $\beta_0$ of \eqref{4.12:eq-running-radii-a} and an
auxiliary integer $L_{\mathrm{a}}$ satisfy, as part of those standing conditions,
\begin{equation*}
\beta_0\le\tfrac12,\qquad L_{\mathrm{a}}\ge2,\qquad L_{\mathrm{a}}^2\le L/3,\qquad L\ge15.
\end{equation*}

\emph{The level of the large-field operation.} The integer $N_0=N_0(k)$ is the depth
parameter, and $k_0:=k-N_0$ is the level of \cite{B-CMP122a}, $N_0$ steps below $k$. Let
$R_m$ denote Ba{\l}aban's scale at level $m$ in \cite[(1.70)\,ff.]{B-CMP122a}. Under the
standing conditions of this chapter on the auxiliary parameters there are constants
$c_R>0$ and $r_R>0$ such that the depth $N_0$ satisfies
\begin{equation}\label{4.12:eq-running-radii-b}
L^{N_0-1}=R_{k_0+1},\qquad R_m\ge c_R\bigl(\log g_m^{-2}\bigr)^{r_R}.
\end{equation}
These two relations are used here as stated in \cite[(1.70)\,ff.]{B-CMP122a}.
Consequently $\inf_k N_0(k)\to\infty$ as $\sup_m g_m\to0$.
\end{definition}

\begin{proof}
Display \eqref{4.12:eq-running-radii-a} is the bound of \cite[(2.6)--(2.7)]{B-CMP119}, and
display \eqref{4.12:eq-running-radii-b} is used as stated above.
Only the monotonicity of the radius and the growth of $N_0$ need an argument.

\emph{Monotonicity.} Regard $g_m$ as a real variable in $(0,1)$ and write
$u=\log g_m^{-2}$, so that $u>0$. Then
\begin{equation*}
\frac{d}{dg_m}\Bigl[g_m\,C_0\,u^{q_0}\Bigr]
=C_0\Bigl(u^{q_0}+g_m\,q_0\,u^{q_0-1}\cdot\bigl(-\tfrac{2}{g_m}\bigr)\Bigr)
=C_0\,u^{q_0-1}\,(u-2q_0).
\end{equation*}
Since $C_0>0$, this derivative is positive as soon as $u>\max\{0,2q_0\}$, that is, for all
$g_m$ below a threshold depending only on $q_0$.

\emph{Growth of $N_0$.} Put $\bar g=\sup_m g_m$, and assume $\bar g<1$; this holds once
$\bar g$ is small. For every $k$ we have $g_{k_0+1}\le\bar g$. Hence
$\log g_{k_0+1}^{-2}\ge\log\bar g^{-2}>0$. Since $r_R>0$ and $c_R>0$, the two relations
\eqref{4.12:eq-running-radii-b} at $m=k_0+1$ give
\begin{equation*}
L^{N_0(k)-1}=R_{k_0+1}\ge c_R\bigl(\log g_{k_0+1}^{-2}\bigr)^{r_R}
\ge c_R\bigl(\log\bar g^{-2}\bigr)^{r_R}.
\end{equation*}
Taking $\log_L$ of both ends, we obtain for every $k$
\begin{equation}\label{4.12:eq-running-radii-2}
N_0(k)\ge1+\log_L c_R+r_R\,\log_L\log\bar g^{-2}.
\end{equation}
The right side of \eqref{4.12:eq-running-radii-2} does not depend on $k$. It tends to
$+\infty$ as $\bar g\to0$. Hence $\inf_k N_0(k)\to\infty$ as $\sup_m g_m\to0$.
\end{proof}

\begin{definition}[Blocks, determining sets and the averaging map]\label{4.12:def-determining-sets}
Fix the level $k$. We use the lattice, the gauge group $G$, its Lie algebra $\mathfrak{g}$ (in the
hermitian convention, so that $e^{ib}\in G$ for $b\in\mathfrak{g}$) and the configurations of
Notation~\ref{4.12:not-lattice-configurations}. The fine lattice has spacing $\eta=L^{-k}$, and
we put $\ell_m=L^{m-k}$ for $0\le m\le k$; in this definition the letter $m$ denotes such a block
level and not the exponent of the torus $\mathbb{T}_m$, and the letter $p$ denotes a level,
fine-lattice plaquettes being written $\mathfrak{p}$ and block-lattice plaquettes $P$. For a
configuration $U$ and a gauge
transformation $\sigma$, $U^{\sigma}$
is the gauge transform of $U$, and $\partial U(\mathfrak{p})$ is the plaquette variable of a
plaquette $\mathfrak{p}$ of the fine lattice.
\begin{enumerate}
\item[(i)] \emph{Blocks.} For $0\le m\le k$ the \emph{level-$m$ blocks} are the cubes of side
$\ell_m$ of the $L$-adic family anchored at the origin. Every level-$m$ block is the union of $L^d$
level-$(m-1)$ blocks.
\item[(ii)] \emph{Determining sets and block lattices.} A \emph{domain} is a union $X$ of blocks. A
\emph{determining set} $\mathfrak{B}$ on $X$ is a partition of $X$ into blocks, and $\ell(B)$
denotes the side of $B\in\mathfrak{B}$. Each block $B\in\mathfrak{B}$ carries a \emph{base point}
$z_B\in B$, and $B(z)$ denotes the block whose base point is $z$. The \emph{block lattice}
$\mathbb{L}(\mathfrak{B})$ has the base points as vertices and one bond for each pair of adjacent
blocks, that is, blocks sharing a $(d-1)$-dimensional piece of face.

\emph{Data} on $\mathbb{L}(\mathfrak{B})$ are $G$-valued functions $D$ on its oriented bonds such
that $D(\langle z',z\rangle)=D(\langle z,z'\rangle)^{-1}$. A \emph{gauge transformation} of
$\mathbb{L}(\mathfrak{B})$ is a $G$-valued function $v$ on the base points. It acts by
\begin{equation*}
D^{v}(\langle z,z'\rangle):=v(z)\,D(\langle z,z'\rangle)\,v(z')^{-1},
\end{equation*}
so that $(D^{v})^{v'}=D^{v'v}$.

For a bond $\mathbf{b}=\langle z,z'\rangle$ of $\mathbb{L}(\mathfrak{B})$ we put
$\ell(\mathbf{b}):=\max\bigl(\ell(B(z)),\ell(B(z'))\bigr)$. A plaquette $P$ of
$\mathbb{L}(\mathfrak{B})$ is an elementary closed circuit of $\mathbb{L}(\mathfrak{B})$, and
$\partial D(P)$ is the ordered product of $D$ along it. We write $\ell(P)$ for the largest side of
a block through which $P$ passes.
\item[(iii)] \emph{The determining set at a prescribed level.} Let $X$ be a domain, let $c$ be the
\emph{level function} on $X$ of one of the spaces of this chapter, that is, the function assigning
to every site $x$ of $X$ the level $c(x)$ that the space carries at $x$, and let $p$
be a level. The determining set $\mathfrak{B}(p,X)$, fixed by the determining-set hypothesis
stated after this list, is a partition of $X$ into blocks of the nested $L$-adic families of
item~(i). The
block containing a site $x$ has level
\begin{equation}\label{4.12:eq-determining-sets-a}
b_{p,X}(x)=\min\bigl(p,\,c(x),\,t_X(x)\bigr),
\end{equation}
where $c$ is suppressed from the notation $b_{p,X}$, as from $\mathfrak{B}(p,X)$, and where the
\emph{boundary tower} $t_X$, which assigns a level $t_X(x)$ to every site $x$ of $X$,
depends only on $X$. Two blocks meeting at a site have levels
differing by at most one. Each region on which the level is constant is a union of blocks of that
level. Every block $B$ has a base point $z_B\in B$, and $\mathbb{L}(\mathfrak{B}(p,X))$ is the
block lattice of item~(ii). The boundary tower satisfies, for every $m$ with $0\le m\le k$ and
every site $x$ of $X$,
\begin{equation}\label{4.12:eq-determining-sets-b}
t_X(x)\ge m\qquad\text{whenever}\qquad \operatorname{dist}(x,\partial X)\ge 2M_1\ell_m ,
\end{equation}
where $\operatorname{dist}(x,\partial X)$ denotes the distance of $x$ from the boundary $\partial X$
of $X$, measured in the same units as the side $\ell_m$.
\item[(iv)] \emph{The smallest nearby block side.} For a bond or plaquette $\xi$ of $X$ let
$\xi^{+}$ denote the set of sites that are sites of $\xi$ or nearest neighbours of such sites, and
put
\begin{equation}\label{4.12:eq-determining-sets-1}
\ell_{\mathfrak{B}}(\xi):=\min\bigl\{\ell(B):\ B\in\mathfrak{B},\ B\cap\xi^{+}\neq\emptyset\bigr\}.
\end{equation}
\item[(v)] \emph{The averaging map.} The map $M_{\mathfrak{B}}$ assigns to every configuration $U$
on the bonds of $X$ data $M_{\mathfrak{B}}(U)$ on $\mathbb{L}(\mathfrak{B})$. These data are the
averages of \cite[(1.5)]{B-CMP98}, composed at interfaces between blocks as in
\cite[(2.14)]{B-CMP119}.

First, $M_{\mathfrak{B}}$ is normalised, gauge covariant and local:
\begin{equation}\label{4.12:eq-determining-sets-c}
\begin{gathered}
M_{\mathfrak{B}}(1)=1,\qquad M_{\mathfrak{B}}(U^{\sigma})=M_{\mathfrak{B}}(U)^{\bar\sigma},\\
\text{$M_{\mathfrak{B}}(U)(\langle z,z'\rangle)$ depends only on the restriction of $U$}\\
\text{to the bonds with both ends in $B(z)\cup B(z')$},
\end{gathered}
\end{equation}
where $\bar\sigma$ is the restriction of $\sigma$ to the base points, $U$ is a configuration on
the bonds of $X$, and $\langle z,z'\rangle$ is a bond of $\mathbb{L}(\mathfrak{B})$.

Second, there are constants $c_Q\ge1$ and $c_{\mathrm{av}}>0$, depending only on $d$ and $L$, with
the following property. Let $\mathfrak{a}$ be a $\mathfrak{g}$-valued $1$-form on the bonds of $X$
with $\sup|\mathfrak{a}|\le c_{\mathrm{av}}$, and let $e^{i\eta\mathfrak{a}}\cdot 1$ be the
configuration whose
variable on a bond $\mathfrak{b}$ is $e^{i\eta\mathfrak{a}(\mathfrak{b})}$. Then there exist a
$G$-valued function $u$ on the base points and a $\mathfrak{g}$-valued function $b$ on the bonds of
$\mathbb{L}(\mathfrak{B})$ such that
\begin{equation}\label{4.12:eq-determining-sets-d}
\begin{gathered}
M_{\mathfrak{B}}(e^{i\eta\mathfrak{a}}\cdot 1)=(e^{ib})^{u},\\
|b(\langle z,z'\rangle)|\le c_Q\,\ell(\langle z,z'\rangle)\,
\sup\bigl\{|\mathfrak{a}(\mathfrak{b})|:\ \mathfrak{b}\subset B(z)\cup B(z')\bigr\}.
\end{gathered}
\end{equation}

Third, suppose $\sup|\partial U-1|\,\eta^{-2}\le c_{\mathrm{av}}$. For a plaquette $P$ of
$\mathbb{L}(\mathfrak{B})$ let $\mathcal{P}_P$ denote the set of plaquettes $\mathfrak{p}$ of the
fine lattice contained in the blocks through which $P$ passes. Then every plaquette $P$ of
$\mathbb{L}(\mathfrak{B})$ satisfies
\begin{equation}\label{4.12:eq-determining-sets-e}
|\partial M_{\mathfrak{B}}(U)(P)-1|
\le 2\,\Bigl(\frac{\ell(P)}{\eta}\Bigr)^{2}\sup\bigl\{|\partial U(\mathfrak{p})-1|:\
\mathfrak{p}\in\mathcal{P}_P\bigr\}.
\end{equation}
\end{enumerate}

It is a hypothesis of this chapter that, for every domain $X$, a boundary tower $t_X$ and, for every
level function $c$ on $X$ of a space of this chapter and every level $p$, a determining set
$\mathfrak{B}(p,X)$ with base points are fixed once and for all with all the properties listed in
item~(iii), including
\eqref{4.12:eq-determining-sets-a} and \eqref{4.12:eq-determining-sets-b}. We call it the
\emph{determining-set hypothesis} and invoke it under that name. Item~(iii) abstracts, as
hypotheses, the properties of three constructions of Ba{\l}aban: the determining sets
$\mathbb{B}_p(X)\cup\{\Gamma_i\}_{i<p}$ of \cite{B-CMP122a}, which enter the minimisers of the
large-field operation; the minimal determining sets and towers of \cite{B-CMP119}; and the maximal
possible domains of \cite{B-CMP109}.

The properties \eqref{4.12:eq-determining-sets-c}, \eqref{4.12:eq-determining-sets-d} and
\eqref{4.12:eq-determining-sets-e} of item~(v) enter this chapter as inputs. The normalisation and
the locality in \eqref{4.12:eq-determining-sets-c}, the estimate \eqref{4.12:eq-determining-sets-d}
for bonds of $\mathbb{L}(\mathfrak{B})$ joining blocks of equal level, and
\eqref{4.12:eq-determining-sets-e} are Ba{\l}aban's \cite[Propositions~2, 6 and~7]{B-CMP98}, taken
here by statement, in the form of the averaging bounds of this chapter. In
particular:
\begin{itemize}
\item \eqref{4.12:eq-determining-sets-e} is taken under \cite[hypothesis (52)]{B-CMP98};
\item \eqref{4.12:eq-determining-sets-d} is the averaging estimate at the identity base
configuration and at $\mathfrak{g}$-valued argument.
\end{itemize}
The covariance in \eqref{4.12:eq-determining-sets-c} is taken by statement from \cite{B-CMP109},
in the form of the covariance lemma for the averaging map.
For bonds of $\mathbb{L}(\mathfrak{B})$ joining blocks of different levels,
\eqref{4.12:eq-determining-sets-d} is the extension of its equal-level case to interfaces composed
as in \cite[(2.14)]{B-CMP119}, stated in the form of the averaging bounds. In that case it is not
taken by statement from \cite{B-CMP98}: it is a hypothesis of this chapter, which we call the
\emph{interface averaging hypothesis} and invoke under that name.
\end{definition}

\begin{definition}[Constrained minimisers and their response to small data]
\label{4.12:def-constrained-minimisers}
Fix the level $k$ and write $\eta=L^{-k}$ for the lattice spacing. Let $\mathfrak{B}$ be a
determining set on a domain $X$, with block lattice $\mathbb{L}(\mathfrak{B})$ and averaging
map $M_{\mathfrak{B}}$, as in Definition~\ref{4.12:def-determining-sets}. For a plaquette or
bond $\xi$ of $X$, $\ell_{\mathfrak{B}}(\xi)$ denotes the local block side of $\mathfrak{B}$ at
$\xi$ fixed in that definition. Configurations, the plaquette variables $\partial U(\mathfrak{p})$,
the Lie algebra $\mathfrak{g}$ and the current $J(U)$ are those fixed in the
conventions~\ref{4.12:not-lattice-configurations}.

\begin{enumerate}
\item[(a)] For $G$-valued data $D$ on $\mathbb{L}(\mathfrak{B})$, the \emph{constraint manifold}
is
\begin{equation*}
C(\mathfrak{B},D):=\{\,U\ \text{$G$-valued on the bonds of } X : M_{\mathfrak{B}}(U)=D\,\}.
\end{equation*}
\item[(b)] The \emph{Wilson action on $X$} is
\begin{equation*}
\mathcal{A}(U):=\sum_{\mathfrak{p}\subset X}\operatorname{Re}\operatorname{tr}
\bigl(1-\partial U(\mathfrak{p})\bigr).
\end{equation*}
\item[(c)] A configuration $U\in C(\mathfrak{B},D)$ is \emph{critical} if
$\frac{d}{ds}\mathcal{A}(U_s)\big|_{s=0}=0$ for every differentiable curve $s\mapsto U_s$ in
$C(\mathfrak{B},D)$ with $U_0=U$.
\end{enumerate}

\emph{The condition on the data.} Hypothesis \cite[(7)]{B-CMP102b}, with parameter
$\varepsilon_1$, is used in the following form (compare \cite[(172)]{B-CMP102b}). It holds for
$G$-valued data $D$ on $\mathbb{L}(\mathfrak{B})$ as soon as
\begin{equation}\label{4.12:eq-constrained-minimisers-a}
|\partial D(P)-1|\le\varepsilon_1\qquad\text{for every plaquette $P$ of }
\mathbb{L}(\mathfrak{B}).
\end{equation}
Moreover, it is preserved under $G$-valued gauge transformations of $\mathbb{L}(\mathfrak{B})$.

\emph{Existence and uniqueness.} The following is \cite[Theorem~1]{B-CMP102b}, with the spaces
\cite[(2), (6), (8)]{B-CMP102b} taken at the local block scale $\ell_{\mathfrak{B}}(\cdot)$ of
$\mathfrak{B}$; this block-scale form of the theorem is used here by statement. It is applied
with the domain of that theorem taken to be the localisation domain $X$, as in
\cite{B-CMP109}, and with hypothesis \cite[(7)]{B-CMP102b} taken in the form
\eqref{4.12:eq-constrained-minimisers-a}.
There are constants $a_0$, $a_1$, $B_3\ge1$ and $\varepsilon_0\in(0,a_0]$, depending on $d$
and $L$ only, with the following property. Let $\mathfrak{B}$ be any determining set, of the kind
fixed in Definition~\ref{4.12:def-determining-sets}, on a domain $X$. Let $D$ be $G$-valued
data on $\mathbb{L}(\mathfrak{B})$ satisfying
\eqref{4.12:eq-constrained-minimisers-a} with
\begin{equation*}
\varepsilon_1\le a_1 \qquad\text{and}\qquad B_3\varepsilon_1\le\varepsilon_0 .
\end{equation*}
Then the following hold.
\begin{enumerate}
\item[(i)] There is a critical $G$-valued configuration $U(\mathfrak{B},D)\in C(\mathfrak{B},D)$
such that, for all plaquettes $\mathfrak{p}\subset X$ and all bonds $\mathfrak{b}\subset X$, the
two bounds in the first row of the following display hold; the last row of the display is the
threshold used in (ii):
\begin{equation}\label{4.12:eq-constrained-minimisers-b}
\begin{gathered}
|\partial U(\mathfrak{B},D)(\mathfrak{p})-1|<B_3\varepsilon_1\eta^2
\ell_{\mathfrak{B}}(\mathfrak{p})^{-2},\qquad
|J(U(\mathfrak{B},D))(\mathfrak{b})|<B_3\varepsilon_1\ell_{\mathfrak{B}}(\mathfrak{b})^{-3},\\
|\partial U(\mathfrak{p})-1|<\varepsilon_0\eta^2\ell_{\mathfrak{B}}(\mathfrak{p})^{-2}.
\end{gathered}
\end{equation}
\item[(ii)] Let $U\in C(\mathfrak{B},D)$ be a critical $G$-valued configuration that satisfies,
for all plaquettes $\mathfrak{p}\subset X$, the bound in the last row of
\eqref{4.12:eq-constrained-minimisers-b}. Then $U=U(\mathfrak{B},D)^{\sigma}$ for some
$G$-valued gauge transformation $\sigma$.
\end{enumerate}
The configuration $U(\mathfrak{B},D)$ is the \emph{constrained minimiser} with data $D$.

\emph{Response to small data.} The response of $U(\mathfrak{B},D)$ to data near the identity is
\cite[Proposition~9 and (190)]{B-CMP102b}, taken at the trivial background $V_0=1$ and for
$G$-valued data. It is used here by statement in the following form, in which these bounds are
expressed through a block-sum norm $\|b\|(x)$ whose decay is measured in the scaled block
distance of \cite[(1.46)--(1.47)]{B-CMP122a}. Several objects enter it.
\begin{itemize}
\item $\delta>0$ is the decay constant of \cite[(1.47)]{B-CMP122a}.
\item For a site $x\in X$ and a bond $\mathbf{b}$ of $\mathbb{L}(\mathfrak{B})$, $d(x,\mathbf{b})$
is the scaled block distance of \cite[(1.46)--(1.47)]{B-CMP122a}. It is measured in
$M_1$-cubes on the corresponding scales, and is taken from $x$ to the nearer of the two blocks
of $\mathbf{b}$.
\item For a $\mathfrak{g}$-valued $b$ on the bonds of $\mathbb{L}(\mathfrak{B})$, the block-sum
norm is
\begin{equation*}
\|b\|(x):=\sum_{\mathbf{b}\in\mathbb{L}(\mathfrak{B})}e^{-\delta d(x,\mathbf{b})}\,|b(\mathbf{b})|.
\end{equation*}
\item $\mathbf{1}$ is the identity configuration. For a $\mathfrak{g}$-valued $1$-form
$\mathcal{H}$ on $X$, $e^{i\eta\mathcal{H}}\mathbf{1}$ is the configuration whose value on a
bond $\mathfrak{b}$ is $e^{i\eta\mathcal{H}(\mathfrak{b})}$.
\end{itemize}
There are constants $b_H>0$, $B_H\ge1$ and $c_d<\infty$ with the following property. They
depend on $d$, on $L$ and on the constants fixed in \cite{B-CMP99b, B-CMP102b}. Let
$\mathfrak{B}$ be as above, and let $b$ be a $\mathfrak{g}$-valued function on the bonds of
$\mathbb{L}(\mathfrak{B})$ such that $\sup|b|<b_H$ and the data $e^{ib}$ satisfy the
hypotheses of the existence and uniqueness statement above. Then
\begin{equation*}
U(\mathfrak{B},e^{ib})=\bigl(e^{i\eta\mathcal{H}}\mathbf{1}\bigr)^{u'}
\end{equation*}
for some $G$-valued $u'$ and a $\mathfrak{g}$-valued $1$-form
$\mathcal{H}=\mathcal{H}_{\mathfrak{B}}(b)$ on $X$. This $1$-form satisfies a bound at every
site $x\in X$, for every plaquette or bond $\xi$ at $x$ (that is, with lower corner, resp.\
initial point, $x$), and for every bond $\mathfrak{b}'$ of the fine lattice that is either a
bond of $\xi$ or a nearest neighbour of a bond of $\xi$:
\begin{equation}\label{4.12:eq-constrained-minimisers-c}
\begin{gathered}
\ell_{\mathfrak{B}}(\xi)\,|\mathcal{H}(\mathfrak{b}')|\le B_H\|b\|(x),\qquad
\ell_{\mathfrak{B}}(\xi)^2\,|\nabla\mathcal{H}(x)|\le B_H\|b\|(x),\\
\ell_{\mathfrak{B}}(\xi)^3\,|D^{*}D\,\mathcal{H}(\mathfrak{b}')|\le B_H\|b\|(x).
\end{gathered}
\end{equation}
Here $\nabla\mathcal{H}(x)$ denotes the forward lattice derivative of $\mathcal{H}$ at $x$,
that is, the covariant forward derivative $\nabla^{Y}$ taken at the identity configuration
$Y=\mathbf{1}$, and $D^{*}D$ is the second-order operator of \cite[(3.12)]{B-CMP109}. The
constant $c_d$ satisfies
\begin{equation}\label{4.12:eq-constrained-minimisers-d}
\sum_{\mathbf{b}\in\mathbb{L}(\mathfrak{B})}e^{-\delta d(x,\mathbf{b})/2}\le c_d
\qquad\text{for every } x\in X,
\end{equation}
uniformly in $\mathfrak{B}$, $X$ and $k$.
\end{definition}

\begin{definition}[Large-field geometry, admissible boxes and the source datum]
\label{4.12:not-admissible-boxes}
Fix the step $k$, and write $\eta=L^{-k}$ for the lattice spacing at step $k$. We set out only
those features of the geometry of Ba{\l}aban's large-field operation $\mathbf{R}$ that are used in what
follows. Throughout, $Y_4$ is the domain of the source space $\tilde{U}^c_k(Y_4,\tilde{\alpha}_0,
\tilde{\alpha}_1)$ of item~(c), and $\tilde{\alpha}_{0,k}$ is the running radius of
Notation~\ref{4.12:not-running-radii}. The source space is the one constructed in the source-space
theorem of this chapter; the one clause of that source space used below is
\eqref{4.12:eq-admissible-boxes-a} of item~(c). The geometric input \eqref{4.12:eq-admissible-boxes-b}
of item~(b) is taken from \cite[(1.73)]{B-CMP122a} as stated in item~(b).

\smallskip
\noindent(a) \emph{Large-field geometry.} Let $Z$ be the large-field component treated by $\mathbf{R}$ at
step $k$, and let $\Omega_m^{\mathrm{old}}$ be the old domains, indexed by the levels $h\le m\le k$.
The \emph{old level}
$o(x)$ equals $k$ on $\Omega_k^{\mathrm{old}}$ and equals $m$ on the collar
$Z\cap(\Omega_m^{\mathrm{old}}\setminus\Omega_{m+1}^{\mathrm{old}})$. The \emph{healed level}
$\mathsf{n}(x)$ equals $k$ on $Z\cup\Omega_k^{\mathrm{old}}$. Let $Z''_k\subset Z$ be the sub-region of
$Z$ singled out in \cite[(1.73)]{B-CMP122a}. Let $\mathfrak{B}''_k$ be the determining set of the step,
in the sense of
Definition~\ref{4.12:def-determining-sets}, attached to the large-field geometry just described;
$\mathsf{b}(x)$ denotes the level of its block containing $x$.
With $k_0=k-N_0$, $N_0$ being the depth parameter of the large-field construction, with $\Gamma''_m$
the thin shells, and with $Z''_{k_0+1}$ the corresponding region at level $k_0+1$,
\begin{equation*}
\mathsf{b}(x)=
\begin{cases}
k & \text{on } \Omega_k^{\mathrm{old}}\cup(Z\setminus Z''_k),\\
m & \text{on } \Gamma''_m,\quad k_0<m<k,\\
j & \text{on } \Gamma''_j\cap Z''_{k_0+1},\quad j\le k_0 .
\end{cases}
\end{equation*}
Let $\Omega_{k_0+1}$ be the domain of level $k_0+1$ of the large-field construction, and
$\Omega^{\sim 5}_{k_0+1}$ its enlargement. Put
\begin{equation*}
B_5:=Z\cap\bigl(\Omega^{\sim 5}_{k_0+1}\setminus\Omega_{k_0+1}\bigr),
\end{equation*}
on which $\mathsf{b}=k$ and $o=k_0$. Let $\Omega''^{\sim}_{h+1}$ be the enlarged domain of level $h+1$ of
the large-field construction, whose boundary $\partial\Omega''^{\sim}_{h+1}$ is the cut surface. The
\emph{core} is
\begin{equation*}
\mathcal{C}:=Z\cap\bigl(\Omega''^{\sim}_{h+1}\bigr)^{c}.
\end{equation*}
On $\mathcal{C}$ the $G$-valued factor entering the composite configurations of the step is not the
factor $\bar{U}$ of item~(c) below but a second $G$-valued factor, denoted $U_2$, carried by those
composite configurations on the core. Finally, $\mathfrak{K}_c$ is
the band formed by two layers,
each of width $MR_{h+1}$, about the cut surface $\partial\Omega''^{\sim}_{h+1}$, where $R_{h+1}$ is the
level-$(h+1)$ radius of the large-field construction. Two properties of this geometry are used below;
the first concerns the old levels and the levels of the blocks of $\mathfrak{B}''_k$, the second the
region $Z''_k$; both are taken by statement:
\begin{equation*}
o(x)\le\mathsf{b}(x)\le k-1\quad\text{for } x\in Z''_k\setminus\mathcal{C},
\qquad
\mathfrak{K}_c\setminus\mathcal{C}\subset Z''_k .
\end{equation*}

\smallskip
\noindent(b) \emph{Admissible boxes.} An \emph{admissible box} is a box
\begin{equation*}
\mathfrak{Q}=\prod_{\mu}\,[a_\mu,\,a_\mu+D_\mu]\ \cap\ \eta\mathbb{Z}^d
\end{equation*}
such that $\mathfrak{Q}\subset Y_4$, $\mathsf{n}\equiv k$ on $\mathfrak{Q}$, and
$\max_\mu D_\mu\le D_{\mathfrak{Q}}:=204M$. Since $\mathfrak{Q}\subset Y_4$ and $\mathsf{n}\equiv k$ on
$\mathfrak{Q}$, the bound \eqref{4.12:eq-admissible-boxes-a} of item~(c) below holds at every plaquette
of $\mathfrak{Q}$,
with radius $\tilde{\alpha}_{0,k}\eta^2$.

Let $\Lambda$ be the set of \cite[(1.73)]{B-CMP122a} that contains $Z''_k$, and let $\Lambda^{\sim}$ be
$\Lambda$ with one
layer of unit $M$-cubes added, so that $\Lambda\subset\Lambda^{\sim}$. With $\Lambda$, $\Lambda^{\sim}$
and $B_5$ as above, the following properties are used here as stated in \cite[(1.73)]{B-CMP122a}:
\begin{equation}\label{4.12:eq-admissible-boxes-b}
\begin{gathered}
Z''_k\subset\Lambda,\qquad \Lambda^{\sim}\subset Z''_k\cup B_5,\qquad \Lambda^{\sim}\subset Y_4,\\
\Lambda^{\sim}\ \text{is a box with edges}\ \le 102M .
\end{gathered}
\end{equation}

\smallskip
\noindent(c) \emph{The source datum.} The source datum is a point
\begin{equation*}
(\mathbb{U},\mathbb{J})\in\tilde{U}^c_k(Y_4,\tilde{\alpha}_0,\tilde{\alpha}_1)
\end{equation*}
of the source space constructed in the source-space theorem of this chapter. By clause~(i) of
\cite[(2.34)]{B-CMP119}, $\mathbb{U}=U'\bar{U}$ with $\bar{U}$ $G$-valued; $\bar{U}$ is the $G$-valued
factor of $\mathbb{U}$. The following is a clause of the definition of the source space, and therefore a
hypothesis on the datum (cf.\ \cite[(2.34)]{B-CMP119}): on the part of $Y_4$ where $\mathsf{n}=k$, the
factor $\bar{U}$ is $G$-valued and
\begin{equation}\label{4.12:eq-admissible-boxes-a}
\bigl|\partial\bar{U}(\mathfrak{p})-1\bigr|<\tilde{\alpha}_{0,k}\,\eta^2
\quad\text{for every plaquette } \mathfrak{p} \text{ there.}
\end{equation}
The only property of $\bar{U}$ used in what follows is \eqref{4.12:eq-admissible-boxes-a} on an admissible
box.

\smallskip
\noindent\emph{Claim.} $\Lambda^{\sim}$ is an admissible box, and $Z''_k\subset\Lambda^{\sim}$.
\end{definition}

\begin{proof}[Proof of the claim]
By \eqref{4.12:eq-admissible-boxes-b}, $\Lambda^{\sim}$ is a box whose edges are at most $102M$. Since
$102M\le 204M=D_{\mathfrak{Q}}$, the side condition $\max_\mu D_\mu\le D_{\mathfrak{Q}}$ holds. Again by
\eqref{4.12:eq-admissible-boxes-b}, $\Lambda^{\sim}\subset Y_4$.

It remains to show that $\mathsf{n}\equiv k$ on $\Lambda^{\sim}$. By \eqref{4.12:eq-admissible-boxes-b},
$\Lambda^{\sim}\subset Z''_k\cup B_5$. By item~(a), $Z''_k\subset Z$, and $B_5\subset Z$ by its
definition as an intersection with $Z$. Hence $\Lambda^{\sim}\subset Z$. Since $\mathsf{n}=k$ on
$Z\cup\Omega_k^{\mathrm{old}}$, we get $\mathsf{n}\equiv k$ on $\Lambda^{\sim}$. Thus $\Lambda^{\sim}$
satisfies every condition of item~(b) and is an admissible box.

Finally, $Z''_k\subset\Lambda$ by \eqref{4.12:eq-admissible-boxes-b}. Since $\Lambda^{\sim}$ is
$\Lambda$ with a layer of cubes added, $\Lambda\subset\Lambda^{\sim}$. Hence $Z''_k\subset\Lambda^{\sim}$.

In particular, by the remark following the definition of admissible boxes,
\eqref{4.12:eq-admissible-boxes-a} holds at every plaquette of $\Lambda^{\sim}$ with radius
$\tilde{\alpha}_{0,k}\eta^2$.
\end{proof}

\begin{definition}[Consumers, their determining sets and radii]\label{4.12:def-consumer-radii}
The following notation is used throughout this definition:
\begin{itemize}
\item the spacing $\eta=L^{-k}$, the block sides $\ell_m=L^{m-k}$, the running radii
$\tilde{\alpha}_{0,m}$, and the letters $\theta_*=1-4\beta$, $\beta$, $L_0$, $k_0$, $N_0$ of
Definition~\ref{4.12:not-running-radii};
\item the determining sets $\mathfrak{B}$, their block-averaging maps $M_{\mathfrak{B}}$ and the
boundary tower $t_X$ of Definition~\ref{4.12:def-determining-sets};
\item the constrained minimisers $U(\mathfrak{B},D)$ of
Definition~\ref{4.12:def-constrained-minimisers};
\item the admissible boxes $\mathfrak{Q}$, the regions $Z''_k$ and $B_5$, the core $\mathcal{C}$,
the old and block level functions $o$ and $\mathsf{b}$, and the $G$-valued factor $\bar{U}$ of the source
datum, all from Definition~\ref{4.12:not-admissible-boxes}; here and below $G$ denotes the gauge
group.
\end{itemize}
The consumers' spaces are those of \cite[(1.64)--(1.69)]{B-CMP122a},
\cite[(2.34)--(2.39)]{B-CMP119} and \cite[(1.11)--(1.16)]{B-CMP109}, grouped in the five classes
(K1)--(K5) of the consumer classification.

\begin{enumerate}
\item[(a)] \emph{Consumers.} A \emph{consumer} is a term $E$ of one of the classes (K1)--(K5), with
domain $X'$. By the convention on derived slots, $E$ reads derived slots for every sub-domain
$X''\subset X'$ of its class of domains, and for every slot level $p$ of its space. Here
$p\le j$ for the spaces $\tilde{U}_j$ and $\mathfrak{U}$ of level $j$, and $p\le k$ for the
classes (K3) and (K4). The five classes are listed in (b). The space of $E$ on $X'$, as given
for its class in the consumer classification, is denoted $D_E(X')$. In the consumer
classification the radius vector of $D_E(X')$ carries, at each point, the pointwise entries
$E_1$, $E_3$ and the derived slots $E_2$, $E_4$; for the $G$-valued factor the derived slots
$E_2$ and $E_4$ are the two conditions of (g), on the plaquette variable and on the current
respectively. The subscripts distinguish these entries from the consumer $E$ itself.

\item[(b)] \emph{Level function.} The \emph{level function} $c_E$ of $E$ on $X'$ is defined as
follows.
\begin{itemize}
\item For class (K1) (the old terms $E^{(j)}$, $R^{(j)}$ of \cite[(1.11)--(1.16)]{B-CMP109}, in
the one-scale space there; the superscript $(j)$ distinguishes them from the consumer $E$),
$c_E\equiv j$.
\item For class (K2) (the old terms $B^{(j)}$), $c_E=\min(o,j)$.
\item For class (K3) (the last-step terms $B^{(k)}$), $c_E=o$.
\item For class (K4) (the terms $C^{(\cdot)}$ with target space $\mathfrak{U}_N$), $c_E=\mathsf{b}$.
\item Terms of class (K5) are read through the classes (K1)--(K4).
\end{itemize}

\item[(c)] \emph{Determining set and minimiser.} By
\eqref{4.12:eq-determining-sets-a}, the consumer's determining set for the pair $(p,X'')$ is the
set $\mathfrak{B}_E(p,X'')$ whose block level at $x$ is
\begin{equation}\label{4.12:eq-consumer-radii-1}
b_E(x)=\min\bigl(p,\,c_E(x),\,t_{X''}(x)\bigr).
\end{equation}
We put
\begin{equation}\label{4.12:eq-consumer-radii-2}
U_E(p,X''):=U\bigl(\mathfrak{B}_E(p,X''),\,M_{\mathfrak{B}_E(p,X'')}(\bar{U})\bigr)
\end{equation}
whenever the right-hand side exists by \eqref{4.12:eq-constrained-minimisers-b}.

\item[(d)] \emph{Consumer level.} For a plaquette or bond $\xi\subset X''$ put
\[
n_E(\xi):=\max\{c_E(y): y \text{ a site of } \xi\}.
\]

\item[(e)] \emph{Two-layer interior.} The derived slots of $E$ on a domain or sub-domain $X''$
are imposed at the plaquettes and bonds of
\begin{equation}\label{4.12:eq-consumer-radii-b}
X''^{\sim 2}:=\bigl\{x\in X'':\ \operatorname{dist}(x,\partial X'')\ge 2M\ell_{c_E(x)}\bigr\}.
\end{equation}
Here $\partial X''$ is the boundary of $X''$ and $\operatorname{dist}$ is the distance of
Definition~\ref{4.4:not-units-distances}.
This is the consumers' form of the exclusion, in \cite[\S1]{B-CMP109} and in
\cite[(2.35)]{B-CMP119}, of the two layers of cubes closest to the boundary of the domain.

\item[(f)] \emph{Radii.} Let $p$ be a slot level and $X''$ a sub-domain as in (a), let
$\xi\subset X''^{\sim 2}$ be a plaquette or a bond, and put
$n:=n_E(\xi)$. Let $\vartheta_E(\xi)>0$ be the coefficient-times-gain of the space of $E$ at the
region of level $n$. The consumer's radii at $\xi$ for the slot $p$ are
\begin{align}
\operatorname{tgt}_2^E(\xi)&:=\vartheta_E(\xi)\,\tilde{\alpha}_{0,n}\,\eta^{2}\,\ell_n^{-2},
\label{4.12:eq-consumer-radii-3}\\
\operatorname{tgt}_4^E(\xi)&:=\vartheta_E(\xi)\,\tilde{\alpha}_{0,n}\,\ell_n^{-3}\,L^{(n-p)_+}.
\label{4.12:eq-consumer-radii-c}
\end{align}
The radii decrease with $n$. Hence the choice of the maximum in (d) means that, at an interface
between regions of different level, the more demanding of the two assignments is taken.

\item[(g)] \emph{The $G$-valued half of the fourth clause.} In the space $D_E(X')$ of the
consumer, the part of the fourth clause that concerns the $G$-valued factor consists of two
conditions. These are imposed for every slot level $p$ and every sub-domain $X''$ as in (a), for
every plaquette $\xi\subset X''^{\sim 2}$ and every bond $\xi\subset X''^{\sim 2}$ respectively:
\begin{equation}\label{4.12:eq-consumer-radii-4}
\begin{aligned}
\bigl|\partial U_E(p,X'')(\xi)-1\bigr|&\le \operatorname{tgt}_2^E(\xi),\\
\bigl|J\bigl(U_E(p,X'')\bigr)(\xi)\bigr|&\le \operatorname{tgt}_4^E(\xi).
\end{aligned}
\end{equation}
In these conditions $U_E(p,X'')$ is regarded as a configuration on the lattice of spacing $\eta$;
plaquette variables are unchanged by this rescaling.

\item[(h)] \emph{Lower bounds on the coefficient-times-gain.} With $n=n_E(\xi)$ we have
\begin{equation}\label{4.12:eq-consumer-radii-d}
\vartheta_E(\xi)\ge\theta_*\,\hat{\vartheta}(n),\qquad
\hat{\vartheta}(n):=
\begin{cases}
1, & n\le k_0,\\
L_0^{\,2(n-k_0)}, & k_0<n<k,
\end{cases}
\end{equation}
and, if $E$ is of class (K4) and $\xi\subset B_5$,
\begin{equation}\label{4.12:eq-consumer-radii-e}
\vartheta_E(\xi)\ge(1-2\beta)\,L_0^{\,2N_0}\qquad\text{for } n=k.
\end{equation}

\item[(i)] \emph{Points of level $k$.} Let $\mathfrak{Q}\subset Z''_k\cup B_5$ be an admissible
box; this is where the sub-domains $X''$ lie. Then the only consumers having points of level
$c_E=k$ in $\mathfrak{Q}$ are those of class (K4), and their points of level $k$ lie in
$B_5\cap\mathfrak{Q}$.

\item[(j)] \emph{Scaled distance and separation of levels.} $d(\cdot,\cdot)$ denotes the scaled
block distance of \cite[(1.46)]{B-CMP122a}. It is the length of a shortest lattice path, measured
along the path in units of $M_1$-cubes of the local block scale. For the consumer geometries
inside an admissible box, the following holds. If $c_E(x)=n$ and $c_E(y)\ge m$ with $m\ge n+2$,
then
\begin{equation}\label{4.12:eq-consumer-radii-a}
d(x,y)\ \ge\ (m-n-1)\,\frac{M}{M_1}.
\end{equation}
In what follows, $d$ enters only through \eqref{4.12:eq-consumer-radii-a} and through its
appearance in the hypothesis on the response one-form $\mathcal{H}_{\mathfrak{B}}(b)$.
\end{enumerate}

The consumers' spaces enter what follows only through three properties.
\begin{itemize}
\item The level function $c_E$ and the region of each consumer class are as in (b).
\item The derived slots $E_2$, $E_4$ at a point carry the same coefficient-times-gain
$\vartheta_E$ as the pointwise entries $E_1$, $E_3$ there. Moreover $\vartheta_E$ obeys the lower
bound $\theta_*\hat{\vartheta}(c_E)$, as in \eqref{4.12:eq-consumer-radii-d}, with the level-$k$
bound \eqref{4.12:eq-consumer-radii-e}.
\item Off the core $\mathcal{C}$, the $G$-valued factor of every composite $\mathcal{W}_i$ is
$\bar{U}$; this property is taken from the theorem on the $G$-valued factor of the composites.
\end{itemize}
\end{definition}

\begin{proof}
We justify, in turn, the radii $\operatorname{tgt}_2^E$ and $\operatorname{tgt}_4^E$ in (f), the
bounds in (h), the assertion (i), and the separation estimate in (j).

\emph{The radius $\operatorname{tgt}_2^E$.} This is the radius of the derived slot $E_2$ in the
$G$-valued half of the fourth clause, read on the lattice of spacing $\eta$ as in (g). In the
radius vector of the consumer classification, the first entry at level $n$, read on the lattice
of spacing $\eta$, is $\vartheta_E\,\tilde{\alpha}_{0,n}\,\eta^{2}\,\ell_n^{-2}$; this is the
radius of the pointwise entry $E_1$. By the
second of the three properties listed at the end of the definition, the slot $E_2$ carries the
same coefficient-times-gain $\vartheta_E$ as the pointwise entry $E_1$ at the same point; this
gives \eqref{4.12:eq-consumer-radii-3}.

\emph{The radius $\operatorname{tgt}_4^E$.} Consider the normalisation on the lattice of spacing
$\eta$ in which the radius of the entry $E_3$ at level $n$ is
$\vartheta_E\,\tilde{\alpha}_{0,n}\,\ell_n^{-3}$. In this normalisation the radius of the entry
$E_4$ of slot $p$, at a bond of consumer level $n$, is
$\vartheta_E\,\tilde{\alpha}_{0,n}\,\ell_n^{-3}L^{(n-p)_+}$; this is
\eqref{4.12:eq-consumer-radii-c}. This radius is the one of \cite[(2.37)]{B-CMP119}, namely
\[
\alpha_{0,n}\,L^{\,n-\min(p,n)}\,\bigl(L^{n}L^{-p}\bigr)^{-3},
\]
transported from the lattice of spacing $L^{-p}$ to the lattice of spacing $\eta$. The exponent
$n-\min(p,n)$ equals $(n-p)_+$.

\emph{The bounds \eqref{4.12:eq-consumer-radii-d} and \eqref{4.12:eq-consumer-radii-e}.} They are
read off from the coefficient-times-gain of the consumers' spaces in
\cite[(1.64)--(1.69)]{B-CMP122a}, \cite[(2.34)--(2.39)]{B-CMP119} and
\cite[(1.11)--(1.16)]{B-CMP109}. We write $\theta_N(m)$ and $\kappa_j$ for the coefficients that
enter this coefficient-times-gain at the levels $m$ and $j$ respectively, and $h$ for the level
of the consumers' spaces that bounds from below, strictly, the range $h<m<k$ of the levels $m$ at
which the coefficients $\theta_N(m)$ enter below level $k$. These coefficients satisfy
\[
\theta_N(m)\ge\theta_* \quad (h<m<k),\qquad
\theta_N(k)\ge 1-2\beta,\qquad
\kappa_j\ge 1-\beta\ge\theta_*.
\]
The coefficient-times-gain is a coefficient multiplied by a gain. For $n<k$ the gain is bounded
below by $\hat{\vartheta}(n)$, that is, by $1$ for $n\le k_0$ and by $L_0^{2(n-k_0)}$ for
$k_0<n<k$; together with the lower bound $\theta_*$ on the coefficients this gives
\eqref{4.12:eq-consumer-radii-d}. At level $k$, on $B_5$ and for the class (K4), the coefficient
bound $1-2\beta$ multiplied by the gain $L_0^{2N_0}$ of the space gives
\eqref{4.12:eq-consumer-radii-e}.

\emph{Assertion (i).} Let $E$ be of class (K1), (K2) or (K3). On $Z''_k\cup B_5$ the old level
satisfies $o\le k_0$, a property of the large-field geometry of
Definition~\ref{4.12:not-admissible-boxes}. Hence $c_E=\min(o,j)\le o\le k_0$ there for the
class (K2), and $c_E=o\le k_0$ there for the class (K3). Since $k_0<k$, such consumers have no
point of level $k$ in $\mathfrak{Q}$. For the class (K1) the level function is the constant level
$j<k$ of an old term, so such a consumer has no point of level $k$ either. Terms of class (K5)
are read through (K1)--(K4). Hence only
consumers of class (K4) remain. For these $c_E=\mathsf{b}$, and their points of level $k$ lie in
$B_5\cap\mathfrak{Q}$.

\emph{The separation \eqref{4.12:eq-consumer-radii-a}.} This is the separation estimate
\cite[(1.47)]{B-CMP122a}, applied to the consumer geometries inside an admissible box. Each
region of constant intermediate level is at least one layer of cubes of side $M\ell$ wide, with
$\ell$ the block side of its own scale: for the thin shells, and for the shells $\Gamma''_j$ and
the old collars, this width is part of their construction in \cite[\S1]{B-CMP122a}.
\end{proof}

\begin{remark}[Insufficiency of the unchanged real factor]\label{4.12:rem-unchanged-factor}
Let $E$ be a consumer in the sense of Definition~\ref{4.12:def-consumer-radii}, with domain
$X'$ and sub-domain $X''$, fix a level $p$, and let $(\mathbb{U},\mathbb{J})$ be a point of the
source space of Definition~\ref{4.12:not-admissible-boxes}, with $G$-valued factor $\bar{U}$.
Write $\mathfrak{B}_{\nu}(p,X'')$ for the determining set that the source space attaches to
$X''$ at level $p$, built with the healed level $\nu$, and $\mathfrak{B}_E(p,X'')$ for the
determining set of the consumer (Definition~\ref{4.12:def-determining-sets}). Inside $Z''_k$
these two determining sets differ. Consequently, the fact that $\bar{U}$ is left unchanged, as
the clause of the source space on its $G$-valued factor provides, does not by itself deliver the
consumer's slots $E_2$ and $E_4$ for $\bar{U}$. That these slots nevertheless hold for
$\bar{U}$ is the content of part~(c) of the proposition on the consumer's derived slots for the
unchanged factor and of its corollary, both stated below this remark, for the consumers within
the scope stated in them.
\end{remark}

\begin{proof}
Inside $Z''_k$ the consumer's determining set $\mathfrak{B}_E(p,X'')$ consists of the shells
$\Gamma''_i$, or of old collars, of levels $i<p$.

For the set $\mathfrak{B}_{\nu}(p,X'')$ of the source space, the healed level satisfies $\nu=k$
on $Z''_k\setminus\mathcal{C}$, and $k\ge p$. Hence on $Z''_k\setminus\mathcal{C}$ the set
$\mathfrak{B}_{\nu}(p,X'')$ is made of blocks of level $\min(p,t_{X''})$, where $t_{X''}$ is
the boundary tower of $X''$.

The two sets are therefore built from different block structures inside $Z''_k$. The bounds
which the source space carries for $\bar{U}$ are bounds along $\mathfrak{B}_{\nu}(p,X'')$. The
derived slots $E_2$ and $E_4$ of the consumer $E$, by contrast, are
conditions along $\mathfrak{B}_E(p,X'')$, with the consumer's radii $\operatorname{tgt}_2^E$
and $\operatorname{tgt}_4^E$. Leaving $\bar{U}$ unchanged transfers the former and does not
by itself give the latter.
\end{proof}

\begin{definition}[Cube-size floors and coupling thresholds]\label{4.12:def-coupling-thresholds}
The letters $\theta_*$, $\beta$, $\beta_0$, $L_0$ and $N_0$ are those of
Definition~\ref{4.12:not-running-radii}. The constants $c_2$ and $c_Q$ are those of the averaging
map of Definition~\ref{4.12:def-determining-sets}. The constants $a_1$, $B_3$ and
$\varepsilon_0$ are those of Definition~\ref{4.12:def-constrained-minimisers}. The constants
$b_H$, $B_H$, $c_d$ and $\delta$ are those of \eqref{4.12:eq-constrained-minimisers-c}. The
constant $C_J$ is that of \eqref{4.12:eq-lattice-configurations-c}. The number $204M$ is the bound
on the side of an admissible box in Definition~\ref{4.12:not-admissible-boxes}.
The constants $C_0$ and $q_0$ are those entering the running radii of
Definition~\ref{4.12:not-running-radii}.

The conditions below are all one-sided. None of them is an upper bound on $L$, $M$, $N$, $K$, $k$
or $N_0$, and none of them is a lower bound on the couplings.

\emph{Order of choice.} First, $L_0\ge 2$ and
\begin{equation*}
L\ge L_*(\beta):=\max\{15,\ 6/\theta_*,\ 2e^2,\ 3L_0^2\}.
\end{equation*}
Then $M_1$ and $\delta$ are fixed as in \cite{B-CMP99b,B-CMP102b}. Then $M\ge M_*(L)$, with
$M_*(L)$ as below. Finally, the couplings are taken small.

\emph{Floors on $M$.} These are the conditions
\begin{gather}
\delta M/M_1\ge 2, \tag{$\Phi 0$}\label{4.12:eq-coupling-thresholds-a}\\
M\ge M_1, \tag{$\Phi 1$}\label{4.12:eq-coupling-thresholds-b}\\
2L\,e^{-\delta M/(2M_1)}\le 1. \tag{$\Phi 2$}\label{4.12:eq-coupling-thresholds-c}
\end{gather}
Condition \eqref{4.12:eq-coupling-thresholds-a} is the choice made in \cite{B-CMP122a}; it is
needed for a separation property of the large-field geometry used later in this chapter. We write
$M_*(L)$ for a floor such that \eqref{4.12:eq-coupling-thresholds-a}--\eqref{4.12:eq-coupling-thresholds-c}
hold for every $M\ge M_*(L)$.

\emph{Thresholds on the couplings.} Set
\begin{equation*}
\mathfrak{a}_{\max}:=\tfrac{\pi}{3}(d-1)\cdot 204M\cdot\tilde{\alpha}_{0,k},
\qquad
K_b:=\tfrac{\pi}{3}(d-1)\,c_Q\,c_d\,(1+3L^2)\cdot 204M .
\end{equation*}
The following conditions are required to hold at the step $k$:
\begin{equation}\tag{$\Gamma_0$}\label{4.12:eq-coupling-thresholds-d}
\begin{aligned}
&\text{(i)} && \tilde{\alpha}_{0,k}\le c_2;\\
&\text{(ii)} && 2\tilde{\alpha}_{0,k}\le a_1,\qquad 2B_3\tilde{\alpha}_{0,k}\le\varepsilon_0;\\
&\text{(iii)} && \mathfrak{a}_{\max}\le\min(c_2,1);\\
&\text{(iv)} && c_Q\mathfrak{a}_{\max}\le\min(b_H/2,\,1/10),\quad 5c_Q\mathfrak{a}_{\max}\le a_1,
\quad 5B_3c_Q\mathfrak{a}_{\max}\le\varepsilon_0;\\
&\text{(v)} && 64\,L\,B_HK_b\,\tilde{\alpha}_{0,k}\le 1,\qquad 2C_JB_HK_b\,\tilde{\alpha}_{0,k}\le 1;
\end{aligned}
\end{equation}
\begin{gather}
4L^3B_3(1+\beta_0)\,W(N_0)\le\theta_*, \tag{$\Gamma_1$}\label{4.12:eq-coupling-thresholds-e}\\
6L^3B_HK_b(1+\beta_0)\,W(N_0)\le\theta_*, \tag{$\Gamma_2$}\label{4.12:eq-coupling-thresholds-f}\\
4L^3B_3\le(1-2\beta)L_0^{2N_0}, \tag{$\Gamma_3$}\label{4.12:eq-coupling-thresholds-g}
\end{gather}
where
\begin{equation}\label{4.12:eq-coupling-thresholds-h}
W(N_0):=\max\bigl\{L_0^{-2N_0},\ (1+N_0)^{1/2}L^{-N_0}\bigr\}.
\end{equation}

The assertion of this definition has two parts. First, there is a floor
$N_*=N_*(d,L,L_0,M,\beta,\beta_0)$ such that
\eqref{4.12:eq-coupling-thresholds-e}--\eqref{4.12:eq-coupling-thresholds-g} hold for every
$N_0\ge N_*$. Second, there is a threshold
\begin{equation*}
\gamma_*=\gamma_*(d,L,L_0,M,\beta,\beta_0,C_0,q_0)>0
\end{equation*}
such that, for every $k$, all of
\eqref{4.12:eq-coupling-thresholds-d}--\eqref{4.12:eq-coupling-thresholds-g} hold at the step $k$
whenever $\sup_m g_m\le\gamma_*$.
\end{definition}

\begin{proof}
\emph{The floors on $M$.} The constants $M_1>0$ and $\delta>0$ are fixed once $L$ is fixed, so
each condition can be solved for $M$.
\begin{itemize}
\item Condition \eqref{4.12:eq-coupling-thresholds-b} is the inequality $M\ge M_1$.
\item Condition \eqref{4.12:eq-coupling-thresholds-a} is equivalent to $M\ge 2M_1/\delta$.
\item Taking logarithms, condition \eqref{4.12:eq-coupling-thresholds-c} is equivalent to
$\delta M/(2M_1)\ge\log(2L)$, that is, to $M\ge 2M_1\delta^{-1}\log(2L)$.
\end{itemize}
Each of these is a lower bound on $M$. Hence all three conditions hold for every $M$ at least
\begin{equation*}
M_*(L):=\max\bigl\{M_1,\ 2M_1/\delta,\ 2M_1\delta^{-1}\log(2L)\bigr\}.
\end{equation*}

\emph{The conditions on $N_0$.} Since $L_0\ge 2$, we have $L_0^{-2N_0}\to 0$ and
$L_0^{2N_0}\to\infty$ as $N_0\to\infty$.

Since $L\ge 15$, the term $(1+N_0)^{1/2}L^{-N_0}$ tends to $0$. It is also decreasing in $N_0$: the
ratio of consecutive terms is $\bigl((2+N_0)/(1+N_0)\bigr)^{1/2}L^{-1}$, which is less than $1$.
The term $L_0^{-2N_0}$ is decreasing as well. Hence $W(N_0)$ is decreasing in $N_0$ and
\begin{equation*}
W(N_0)\to 0\qquad(N_0\to\infty).
\end{equation*}

The left-hand sides of \eqref{4.12:eq-coupling-thresholds-e} and
\eqref{4.12:eq-coupling-thresholds-f} are $W(N_0)$ times factors that do not depend on $N_0$. These
factors are built from $L$, $B_3$, $B_H$, $\beta_0$ and $K_b$, and $K_b$ depends on $d$, $L$ and
$M$. The right-hand side $\theta_*=1-4\beta$ is positive because $\beta<\tfrac14$
(Definition~\ref{4.12:not-running-radii}). So each of these two inequalities holds for all $N_0$
beyond a finite value.

In \eqref{4.12:eq-coupling-thresholds-g} the left-hand side does not depend on $N_0$. The factor
$1-2\beta$ is positive because $\beta<\tfrac14$ (Definition~\ref{4.12:not-running-radii}), and
$L_0^{2N_0}$ increases to infinity. So this inequality, too, holds for all $N_0$ beyond a finite
value.

Let $N_*(d,L,L_0,M,\beta,\beta_0)$ be the largest of the three values just found. By the
monotonicity in $N_0$ established above,
\eqref{4.12:eq-coupling-thresholds-e}--\eqref{4.12:eq-coupling-thresholds-g} hold for every
$N_0\ge N_*$.

\emph{The threshold on the couplings.} Each item of \eqref{4.12:eq-coupling-thresholds-d} is an
upper bound on $\tilde{\alpha}_{0,k}$: through the definition of $\mathfrak{a}_{\max}$, items
(iii) and (iv) are upper bounds on $\tilde{\alpha}_{0,k}$ as well. Each bound is by a positive
quantity depending on $d$, $L$, $M$ and the fixed constants listed at the start of the
definition. The dependence on $M$, which enters through
$\mathfrak{a}_{\max}$ and $K_b$ via the box-side bound $204M$, is therefore part of the smallness
required of $\tilde{\alpha}_{0,k}$.

By \eqref{4.12:eq-running-radii-a}, these bounds on the radius follow from $\sup_m g_m\le\gamma_*$
for a suitable $\gamma_*>0$. By \eqref{4.12:eq-running-radii-b}, the floors $N_0\ge N_*$ of the
previous paragraph, and with them
\eqref{4.12:eq-coupling-thresholds-e}--\eqref{4.12:eq-coupling-thresholds-g}, hold for all $k$ once
$\sup_m g_m$ is below a threshold. Taking $\gamma_*$ to be the smaller of the two thresholds gives
a $\gamma_*(d,L,L_0,M,\beta,\beta_0,C_0,q_0)>0$ with the stated property.

No new condition on $L$ or $k$ is introduced. The dependence of $K_b$ on $M$, through
the bound $204M$ on the side of an admissible box, is absorbed into the threshold $\gamma_*$.
\end{proof}

\begin{lemma}[Holonomies along two monotone paths]\label{4.12:lem-monotone-paths}
Let $\varphi\ge 0$, let $\mathfrak{Q}$ be a box, and let $U$ be a $G$-valued configuration on
the bonds of $\mathfrak{Q}$ such that
\begin{equation*}
|\partial U(\mathfrak{p})-1|\le\varphi\qquad\text{for every plaquette }\mathfrak{p}\subset\mathfrak{Q}.
\end{equation*}
Let $\Pi,\Pi'$ be two monotone lattice paths in $\mathfrak{Q}$ from a site $u$ to a site $v$,
that is, paths every step of which is $+\eta e_\mu$ for some $\mu$. Let $n_{\mathrm{tr}}$ be the
least number of adjacent transpositions that transform the step word of $\Pi'$ into that
of $\Pi$. Then
\begin{equation}\label{4.12:eq-monotone-paths-1}
|U(\Pi)U(\Pi')^{-1}-1|\le n_{\mathrm{tr}}\,\varphi,
\qquad
n_{\mathrm{tr}}\le\sum_{\mu<\nu}n_\mu n_\nu,
\end{equation}
where $n_\mu$ is the number of $e_\mu$-steps, which is the same for $\Pi$ and $\Pi'$.
\end{lemma}

Here $U(\Pi)$ is the ordered product of the bond variables along $\Pi$, the variable of the
first step standing on the left, and $U(\mathsf{a})$ denotes the same product for the
monotone path starting at $u$ with step word $\mathsf{a}$. The conventions are those of
Definition~\ref{4.12:not-lattice-configurations}: the group $G$ is contained in the group of
unitary $\mathfrak{n}\times\mathfrak{n}$ matrices, and $|\cdot|$ is the operator norm.

\begin{proof}
A monotone path from $u$ to $v$ makes $(v-u)_\mu/\eta$ steps in direction $e_\mu$, so the step
words of $\Pi$ and $\Pi'$ are arrangements of the same multiset of letters, $n_\mu$ copies of
$e_\mu$ for each $\mu$. Hence one is obtained from the other by adjacent transpositions, and
their least number is at most the number of inversions, which is at most
$\sum_{\mu<\nu}n_\mu n_\nu$. Indeed, number the copies of each letter in order of occurrence
in both words and call an inversion a pair of occurrences of distinct letters whose relative
order in the current word is opposite to that in the target word. As long as the two words
differ, some adjacent pair of the current word is an inversion; exchanging it lowers the
number of inversions by exactly one. The number of inversions never exceeds the number of pairs
of occurrences of distinct letters, $\sum_{\mu<\nu}n_\mu n_\nu$. This proves the second
inequality in \eqref{4.12:eq-monotone-paths-1}. Every intermediate word is again the step word
of a monotone path from $u$ to $v$; such a path lies in the box spanned by $u$ and $v$, hence
in $\mathfrak{Q}$.

Consider one transposition, replacing consecutive steps $e_\mu e_\nu$ taken at a site $y$ by
$e_\nu e_\mu$. If $\mu<\nu$, put $\mathfrak{p}=\mathfrak{p}(y;\mu,\nu)$, the plaquette with
corner $y$ spanned by $e_\mu$ and $e_\nu$ in this order. From the definition of
the plaquette variable (Definition~\ref{4.12:not-lattice-configurations}),
\begin{equation}\label{4.12:eq-monotone-paths-2}
\begin{split}
&U(\langle y,y+\eta e_\mu\rangle)\,U(\langle y+\eta e_\mu,y+\eta e_\mu+\eta e_\nu\rangle)\\
&\qquad=\partial U(\mathfrak{p})\,U(\langle y,y+\eta e_\nu\rangle)\,
U(\langle y+\eta e_\nu,y+\eta e_\nu+\eta e_\mu\rangle).
\end{split}
\end{equation}
If $\nu<\mu$, the same identity holds with $\partial U(\mathfrak{p})^{-1}$ in place of
$\partial U(\mathfrak{p})$ and $\mathfrak{p}=\mathfrak{p}(y;\nu,\mu)$. In both cases the four
corners of $\mathfrak{p}$ lie on the two intermediate paths, so $\mathfrak{p}\subset\mathfrak{Q}$.

Let $\mathsf{a}$ and $\mathsf{a}'$ be the words before and after the transposition, and let
$\mathsf{a}_0$ be their common prefix, the path with word $\mathsf{a}_0$ ending at $y$. The two
paths coincide after $y+\eta e_\mu+\eta e_\nu$. Writing $U(\mathsf{a})$ as $U(\mathsf{a}_0)$,
times the left side of \eqref{4.12:eq-monotone-paths-2}, times the common remainder, and inserting
$U(\mathsf{a}_0)^{-1}U(\mathsf{a}_0)$ to the left of the right side of
\eqref{4.12:eq-monotone-paths-2}, we obtain
\begin{equation}\label{4.12:eq-monotone-paths-3}
U(\mathsf{a})=\bigl[U(\mathsf{a}_0)\,\partial U(\mathfrak{p})^{\pm1}\,U(\mathsf{a}_0)^{-1}\bigr]\,
U(\mathsf{a}').
\end{equation}
Since \eqref{4.12:eq-monotone-paths-2} holds in the form stated for $\mu<\nu$ and with the
inverse plaquette variable for $\nu<\mu$, \eqref{4.12:eq-monotone-paths-3} holds in both cases
with the corresponding sign.

Now let $\mathsf{a}^{(0)},\dots,\mathsf{a}^{(n_{\mathrm{tr}})}$ be the words of a chain of
$n_{\mathrm{tr}}$ adjacent transpositions from the word of $\Pi'$ to that of $\Pi$. Applying
\eqref{4.12:eq-monotone-paths-3} with $\mathsf{a}=\mathsf{a}^{(i)}$ and
$\mathsf{a}'=\mathsf{a}^{(i-1)}$ gives $U(\mathsf{a}^{(i)})=\Xi_i\,U(\mathsf{a}^{(i-1)})$, where
$\Xi_i$ denotes the bracketed factor of \eqref{4.12:eq-monotone-paths-3} for the $i$-th
transposition, and iterating,
\begin{equation*}
U(\Pi)\,U(\Pi')^{-1}=\Xi_{n_{\mathrm{tr}}}\cdots\Xi_1 .
\end{equation*}
Each $\Xi_i$ is the conjugate, by the element $U(\mathsf{a}_0)$ of $G$, of a plaquette variable
$\partial U(\mathfrak{p})$ with $\mathfrak{p}\subset\mathfrak{Q}$, or of its inverse. Since
conjugation by an element of $G$ is an isometry of the operator norm and fixes the identity, it
does not change $|\,\cdot\,-1|$. Moreover
\begin{equation*}
|\partial U(\mathfrak{p})^{-1}-1|
=|\partial U(\mathfrak{p})^{-1}(1-\partial U(\mathfrak{p}))|
=|\partial U(\mathfrak{p})-1|.
\end{equation*}
Hence $|\Xi_i-1|\le\varphi$ for every $i$. For unitary matrices $\Xi,\Xi'$, writing
$\Xi\Xi'-1=(\Xi-1)\Xi'+(\Xi'-1)$,
\begin{equation*}
|\Xi\Xi'-1|\le|\Xi-1|\,|\Xi'|+|\Xi'-1|=|\Xi-1|+|\Xi'-1|.
\end{equation*}
By induction on the number of factors, applied with $\Xi=\Xi_j$ and
$\Xi'=\Xi_{j-1}\cdots\Xi_1$ for $j=2,\dots,n_{\mathrm{tr}}$, we obtain
$|\Xi_{n_{\mathrm{tr}}}\cdots\Xi_1-1|\le n_{\mathrm{tr}}\varphi$. This is the first
inequality in \eqref{4.12:eq-monotone-paths-1}.
\end{proof}

\begin{lemma}[The axial gauge of a box]\label{4.12:lem-axial-gauge}
We use the gauge group $G$, its Lie algebra $\mathfrak{g}$ (hermitian convention), the holonomy
$U(\gamma)$ of a path $\gamma$ and the gauge transform $U^{\sigma}$ of
Definition~\ref{4.12:not-lattice-configurations}. Let $\mathfrak{Q}$ be a box of the lattice of
spacing $\eta$ whose edges have length at most $\Delta$, and let $\hat{z}$ be its corner with minimal
coordinates. Let $\varphi\ge 0$, and let $U$ be $G$-valued on the bonds of $\mathfrak{Q}$ with
$|\partial U(\mathfrak{p})-1|\le\varphi$ for every plaquette $\mathfrak{p}\subset\mathfrak{Q}$.
Assume
\begin{equation*}
(d-1)\,\frac{\Delta}{\eta}\,\varphi\le 1 .
\end{equation*}
For $x\in\mathfrak{Q}$ let $\gamma_x$ be the lexicographic path from $\hat{z}$ to $x$, which takes
first all its $e_1$-steps, then all its $e_2$-steps, and so on up to its $e_d$-steps, and put
$a(x):=U(\gamma_x)$. Then $a$ is $G$-valued, $U^{a}=1$ on every bond of every path $\gamma_x$, and
for every bond $\mathfrak{b}\subset\mathfrak{Q}$
\begin{equation}\label{4.12:eq-axial-gauge-1}
|U^{a}(\mathfrak{b})-1|\le (d-1)\,\frac{\Delta}{\eta}\,\varphi ,
\end{equation}
and $\mathfrak{a}(\mathfrak{b}):=\frac{1}{i\eta}\log U^{a}(\mathfrak{b})$ is $\mathfrak{g}$-valued with
\begin{equation}\label{4.12:eq-axial-gauge-2}
|\mathfrak{a}(\mathfrak{b})|\le \frac{\pi}{3}\,(d-1)\,\Delta\,\varphi\,\eta^{-2}.
\end{equation}
\end{lemma}

\begin{proof}
For every $x\in\mathfrak{Q}$, $a(x)=U(\gamma_x)$ is a product of the $G$-valued variables $U$ on the
bonds of $\gamma_x$ (the empty product $1$ when $x=\hat{z}$), so $a$ is $G$-valued.

Let $\mathfrak{b}\subset\mathfrak{Q}$ be a bond, written with its positive orientation as
$\mathfrak{b}=\langle x,x+\eta e_\mu\rangle$ with $x,x+\eta e_\mu\in\mathfrak{Q}$. By the
definition of the gauge transform and of $a$,
\begin{equation*}
U^{a}(\mathfrak{b})=U(\gamma_x)\,U(\mathfrak{b})\,U(\gamma_{x+\eta e_\mu})^{-1}
=U(\gamma_x\cdot\mathfrak{b})\,U(\gamma_{x+\eta e_\mu})^{-1},
\end{equation*}
where $\gamma_x\cdot\mathfrak{b}$ is the path $\gamma_x$ followed by the step $\mathfrak{b}$. Both
$\gamma_x\cdot\mathfrak{b}$ and $\gamma_{x+\eta e_\mu}$ are monotone lattice paths from $\hat{z}$ to
$x+\eta e_\mu$. They lie in $\mathfrak{Q}$, because along a monotone path from $\hat{z}$ every
coordinate lies between the corresponding coordinates of $\hat{z}$ and of the endpoint. For
$\nu=1,\dots,d$ let $n_\nu$ be the number of $e_\nu$-steps of $\gamma_x$, that is, the $\nu$-th
coordinate of $x-\hat{z}$ divided by $\eta$. Since the edges of $\mathfrak{Q}$ have length at most
$\Delta$, we have $0\le n_\nu\le \Delta/\eta$. The step word of $\gamma_x\cdot\mathfrak{b}$ is
$e_1^{n_1}\cdots e_d^{n_d}\,e_\mu$, and that of $\gamma_{x+\eta e_\mu}$ is
$e_1^{n_1}\cdots e_\mu^{n_\mu+1}\cdots e_d^{n_d}$. Moving the final letter $e_\mu$ of the first word
to the left past the block $e_{\mu+1}^{n_{\mu+1}}\cdots e_d^{n_d}$ turns it into the second word,
using
\begin{equation*}
\sum_{\nu>\mu}n_\nu\le (d-1)\,\frac{\Delta}{\eta}
\end{equation*}
adjacent transpositions. Hence the number $N$ of adjacent transpositions needed is at most
$(d-1)\Delta/\eta$. The lemma on holonomies along two monotone paths
(Lemma~\ref{4.12:lem-monotone-paths}), applied to $\gamma=\gamma_x\cdot\mathfrak{b}$ and
$\gamma'=\gamma_{x+\eta e_\mu}$, gives
\begin{equation*}
|U^{a}(\mathfrak{b})-1|\le N\varphi\le (d-1)\,\frac{\Delta}{\eta}\,\varphi ,
\end{equation*}
which is \eqref{4.12:eq-axial-gauge-1}.

By \eqref{4.12:eq-axial-gauge-1} and the hypothesis, $|U^{a}(\mathfrak{b})-1|\le 1$. Applied to the
element $U^{a}(\mathfrak{b})$ of $G$, \eqref{4.12:eq-lattice-configurations-a} yields two facts.
First, $\log U^{a}(\mathfrak{b})\in i\mathfrak{g}$, so that $\mathfrak{a}(\mathfrak{b})$ is
$\mathfrak{g}$-valued. Second,
\begin{equation*}
|\log U^{a}(\mathfrak{b})|\le\frac{\pi}{3}\,(d-1)\,\frac{\Delta}{\eta}\,\varphi .
\end{equation*}
Dividing by $\eta$ gives \eqref{4.12:eq-axial-gauge-2}.

Finally, let $\mathfrak{b}$ be a bond of some path $\gamma_y$. Since $\gamma_y$ is monotone, it
traverses $\mathfrak{b}$ in the positive direction, so we may write
$\mathfrak{b}=\langle x,x+\eta e_\mu\rangle$ as above. The part of $\gamma_y$ preceding
$\mathfrak{b}$ ends at $x$. It consists of all the $e_1$-steps, then all the $e_2$-steps, and so on,
up to some of the $e_\mu$-steps, so it is the lexicographic path $\gamma_x$. Moreover $n_\nu=0$ for
$\nu>\mu$. The count above then gives no transposition, so
$\gamma_x\cdot\mathfrak{b}=\gamma_{x+\eta e_\mu}$, and the identity displayed at the beginning of the
proof gives
\begin{equation*}
U^{a}(\mathfrak{b})=U(\gamma_{x+\eta e_\mu})\,U(\gamma_{x+\eta e_\mu})^{-1}=1 .
\end{equation*}
\end{proof}

\begin{lemma}[Plaquettes of a left-perturbed configuration]\label{4.12:lem-perturbed-plaquettes}
We use the lattice, the gauge group $G$ and the configurations of
Definition~\ref{4.12:not-lattice-configurations}. Let $Y$ be a $G$-valued configuration and
$\mathcal{W}$ a $\mathfrak{g}$-valued $1$-form. Let $\mathfrak{p}=\mathfrak{p}(x;\mu,\nu)$ be a
plaquette, and let $w_0$ be the maximum of $|\mathcal{W}|$ over the four bonds of $\mathfrak{p}$.
Assume $4\eta w_0\le 1$. Then
\begin{equation}\label{4.12:eq-perturbed-plaquettes-1}
\begin{split}
\bigl|\partial(e^{i\eta\mathcal{W}}Y)(\mathfrak{p})-1\bigr|
&\le \bigl|\partial Y(\mathfrak{p})-1\bigr|
+\eta^{2}\Bigl(\bigl|(\nabla^{Y}_{\mu}\mathcal{W}_{\nu})(x)\bigr|
+\bigl|(\nabla^{Y}_{\nu}\mathcal{W}_{\mu})(x)\bigr|\Bigr)\\
&\quad+4\eta\,\bigl|\partial Y(\mathfrak{p})-1\bigr|\,w_0+12\,\eta^{2}w_0^{2}.
\end{split}
\end{equation}
\end{lemma}

\begin{proof}
Throughout the proof we use the following properties of the matrix norm $|\cdot|$:
it is submultiplicative, and it is unchanged by left or right multiplication by an
element of $G$, these elements being unitary. In particular $|\operatorname{Ad}_g A|=|A|$
for $g\in G$ and every matrix $A$.

Write
\begin{align*}
y_1&:=Y(\langle x,x+\mu\rangle), & y_2&:=Y(\langle x+\mu,x+\mu+\nu\rangle),\\
y_3&:=Y(\langle x+\nu,x+\nu+\mu\rangle), & y_4&:=Y(\langle x,x+\nu\rangle),
\end{align*}
and
\begin{align*}
\mathcal{W}_1&:=\mathcal{W}_{\mu}(x), & \mathcal{W}_2&:=\mathcal{W}_{\nu}(x+\mu),\\
\mathcal{W}_3&:=\mathcal{W}_{\mu}(x+\nu), & \mathcal{W}_4&:=\mathcal{W}_{\nu}(x),
\end{align*}
so that $\partial Y(\mathfrak{p})=y_1y_2y_3^{-1}y_4^{-1}$ and $|\mathcal{W}_i|\le w_0$ for
$i=1,\dots,4$. By the definition of the left-perturbed configuration $e^{i\eta\mathcal{W}}Y$
and of the plaquette variable,
\begin{equation}\label{4.12:eq-perturbed-plaquettes-2}
\partial(e^{i\eta\mathcal{W}}Y)(\mathfrak{p})
=e^{i\eta\mathcal{W}_1}y_1\,e^{i\eta\mathcal{W}_2}y_2\,y_3^{-1}e^{-i\eta\mathcal{W}_3}\,
y_4^{-1}e^{-i\eta\mathcal{W}_4}.
\end{equation}

\emph{Moving the group elements to the right.} For $g\in G$ and any matrix $A$ we have
$ge^{A}g^{-1}=e^{\operatorname{Ad}_gA}$, that is, $ge^{A}=e^{\operatorname{Ad}_gA}g$. Applied
with $g=y_1$ to the factor $y_1e^{i\eta\mathcal{W}_2}$, then with $g=y_1y_2y_3^{-1}$ to the
factor $y_1y_2y_3^{-1}e^{-i\eta\mathcal{W}_3}$, and finally with
$g=y_1y_2y_3^{-1}y_4^{-1}=\partial Y(\mathfrak{p})$ to the factor
$\partial Y(\mathfrak{p})e^{-i\eta\mathcal{W}_4}$, the identity
\eqref{4.12:eq-perturbed-plaquettes-2} becomes
\begin{equation}\label{4.12:eq-perturbed-plaquettes-3}
\partial(e^{i\eta\mathcal{W}}Y)(\mathfrak{p})
=e^{i\eta X_1}e^{i\eta X_2}e^{i\eta X_3}e^{i\eta X_4}\,\partial Y(\mathfrak{p}),
\end{equation}
where
\begin{align*}
X_1&=\mathcal{W}_1, & X_2&=\operatorname{Ad}_{y_1}\mathcal{W}_2,\\
X_3&=-\operatorname{Ad}_{y_1y_2y_3^{-1}}\mathcal{W}_3, &
X_4&=-\operatorname{Ad}_{\partial Y(\mathfrak{p})}\mathcal{W}_4.
\end{align*}
Since $\operatorname{Ad}$ preserves the norm, $|X_i|=|\mathcal{W}_i|\le w_0$ for each $i$.

\emph{The sum of the exponents.} Since $y_1y_2y_3^{-1}=\partial Y(\mathfrak{p})y_4$, we have
$X_3=-\operatorname{Ad}_{\partial Y(\mathfrak{p})}\operatorname{Ad}_{y_4}\mathcal{W}_3$, hence
\begin{equation*}
X_3+X_4=-\bigl(\operatorname{Ad}_{y_4}\mathcal{W}_3+\mathcal{W}_4\bigr)
+\bigl(1-\operatorname{Ad}_{\partial Y(\mathfrak{p})}\bigr)
\bigl(\operatorname{Ad}_{y_4}\mathcal{W}_3+\mathcal{W}_4\bigr).
\end{equation*}
By the definition of the covariant forward derivative,
\begin{equation*}
(\nabla^{Y}_{\nu}\mathcal{W}_{\mu})(x)=\eta^{-1}\bigl(\operatorname{Ad}_{y_4}\mathcal{W}_3
-\mathcal{W}_1\bigr),\qquad
(\nabla^{Y}_{\mu}\mathcal{W}_{\nu})(x)=\eta^{-1}\bigl(\operatorname{Ad}_{y_1}\mathcal{W}_2
-\mathcal{W}_4\bigr).
\end{equation*}
Adding $X_1+X_2=\mathcal{W}_1+\operatorname{Ad}_{y_1}\mathcal{W}_2$ to the expression for
$X_3+X_4$ therefore gives
\begin{equation}\label{4.12:eq-perturbed-plaquettes-4}
\begin{split}
\Sigma:=X_1+X_2+X_3+X_4
&=-\eta(\nabla^{Y}_{\nu}\mathcal{W}_{\mu})(x)+\eta(\nabla^{Y}_{\mu}\mathcal{W}_{\nu})(x)\\
&\quad+\bigl(1-\operatorname{Ad}_{\partial Y(\mathfrak{p})}\bigr)
\bigl(\operatorname{Ad}_{y_4}\mathcal{W}_3+\mathcal{W}_4\bigr).
\end{split}
\end{equation}
For $g\in G$ and any matrix $A$,
\begin{equation*}
|(1-\operatorname{Ad}_g)A|\le|A-gA|+|gA-gAg^{-1}|\le 2|g-1|\,|A|,
\end{equation*}
because $|A-gA|=|(1-g)A|\le|g-1|\,|A|$ and
$|gA-gAg^{-1}|=|gA(g-1)g^{-1}|\le|A|\,|g-1|$. Applied with $g=\partial Y(\mathfrak{p})$
and $A=\operatorname{Ad}_{y_4}\mathcal{W}_3+\mathcal{W}_4$, for which
$|A|\le|\mathcal{W}_3|+|\mathcal{W}_4|\le 2w_0$, this and
\eqref{4.12:eq-perturbed-plaquettes-4} give
\begin{equation}\label{4.12:eq-perturbed-plaquettes-5}
|\Sigma|\le\eta\Bigl(\bigl|(\nabla^{Y}_{\mu}\mathcal{W}_{\nu})(x)\bigr|
+\bigl|(\nabla^{Y}_{\nu}\mathcal{W}_{\mu})(x)\bigr|\Bigr)
+4\,\bigl|\partial Y(\mathfrak{p})-1\bigr|\,w_0.
\end{equation}

\emph{Expansion of the four exponentials.} Put $s:=\eta\sum_{i=1}^{4}|X_i|$; then
$s\le 4\eta w_0\le 1$ by hypothesis. Multiplying out the four exponential series, the terms of
total degree $0$ and $1$ give $1+i\eta\Sigma$, so that
\begin{equation*}
e^{i\eta X_1}e^{i\eta X_2}e^{i\eta X_3}e^{i\eta X_4}=1+i\eta\Sigma+R,
\end{equation*}
where $R$ is the sum of the terms of total degree at least $2$. By submultiplicativity, each such
term is dominated in norm by the corresponding term of the product of the scalar series of
$e^{\eta|X_1|}\cdots e^{\eta|X_4|}=e^{s}$, in which the terms of total degree $0$ and $1$ sum to
$1+s$. Hence $|R|\le e^{s}-1-s$. Since $(e^{s}-1-s)/s^{2}=\sum_{n\ge2}s^{n-2}/n!$ is
non-decreasing in $s\ge0$ and equals $e-2$ at $s=1$, we have $e^{s}-1-s\le(e-2)s^{2}\le s^{2}$ for
$0\le s\le1$. Consequently, using $s\le 4\eta w_0$ and $16(e-2)\le 12$,
\begin{equation}\label{4.12:eq-perturbed-plaquettes-6}
\bigl|e^{i\eta X_1}e^{i\eta X_2}e^{i\eta X_3}e^{i\eta X_4}-1\bigr|
\le\eta|\Sigma|+16\,\eta^{2}w_0^{2}\,(e-2)\le\eta|\Sigma|+12\,\eta^{2}w_0^{2}.
\end{equation}

\emph{Conclusion.} Write $\Pi:=e^{i\eta X_1}e^{i\eta X_2}e^{i\eta X_3}e^{i\eta X_4}$. Since
$\partial Y(\mathfrak{p})$ is unitary, $|\partial Y(\mathfrak{p})|=1$, and
\begin{equation*}
|\Pi\,\partial Y(\mathfrak{p})-1|
\le|\Pi-1|\,|\partial Y(\mathfrak{p})|+|\partial Y(\mathfrak{p})-1|
=|\Pi-1|+|\partial Y(\mathfrak{p})-1|.
\end{equation*}
By \eqref{4.12:eq-perturbed-plaquettes-3} the left side is
$|\partial(e^{i\eta\mathcal{W}}Y)(\mathfrak{p})-1|$. Inserting
\eqref{4.12:eq-perturbed-plaquettes-6}, and then \eqref{4.12:eq-perturbed-plaquettes-5}
multiplied by $\eta$, yields \eqref{4.12:eq-perturbed-plaquettes-1}.
\end{proof}

\begin{lemma}[Gauge covariance of the constrained minimiser]\label{4.12:lem-minimiser-covariance}
Let $\mathfrak{B}$ and $X$ be as in Definition~\ref{4.12:def-determining-sets}, see
\eqref{4.12:eq-determining-sets-a}. Let $D$ be data satisfying the hypotheses of the statement on
constrained minimisers in Definition~\ref{4.12:def-constrained-minimisers}, see
\eqref{4.12:eq-constrained-minimisers-b}. In this lemma only, let $g$ be a $G$-valued function on
the base points $z_B$ of the blocks $B$ of $\mathfrak{B}$, and let $D^{g}$ denote the gauge
transform of $D$ by $g$. Then $D^{g}$ satisfies the same hypotheses, and there is a $G$-valued
function $\sigma$ on the sites of $X$ such that
\begin{equation*}
U(\mathfrak{B},D^{g})=U(\mathfrak{B},D)^{\sigma}.
\end{equation*}
\end{lemma}

\begin{proof}
The plaquette variables of $D^{g}$ are conjugates, by elements of $G$, of the plaquette variables
of $D$. The condition \eqref{4.12:eq-constrained-minimisers-a} is unchanged under such
conjugation. Hence it holds for $D^{g}$ with the same constant $\varepsilon_1$, and $D^{g}$
satisfies the hypotheses of the statement on constrained minimisers.

Let $\tilde g$ be the $G$-valued function on the sites of $X$ that equals $g(z_B)$ at every site of
the block $B$, for every block $B$ of $\mathfrak{B}$. By \eqref{4.12:eq-determining-sets-c},
\begin{equation*}
M_{\mathfrak{B}}(U^{\tilde g})=M_{\mathfrak{B}}(U)^{g}\qquad\text{for every configuration } U .
\end{equation*}
Thus $M_{\mathfrak{B}}(U)=D$ implies $M_{\mathfrak{B}}(U^{\tilde g})=D^{g}$. The function
$\tilde g^{-1}$ is obtained from $g^{-1}$ by the same construction. Applying
\eqref{4.12:eq-determining-sets-c} to it shows that $M_{\mathfrak{B}}(U)=D^{g}$ implies
$M_{\mathfrak{B}}(U^{\tilde g^{-1}})=D$.

Therefore $U\mapsto U^{\tilde g}$ maps $C(\mathfrak{B},D)$ bijectively onto
$C(\mathfrak{B},D^{g})$, with inverse $U\mapsto U^{\tilde g^{-1}}$. Each bond variable is
multiplied on both sides by fixed elements of $G$. Hence both maps send differentiable curves to
differentiable curves.

The action $\mathcal{A}$ is gauge invariant, so $\mathcal{A}(U^{\tilde g})=\mathcal{A}(U)$.

Put $U_0:=U(\mathfrak{B},D)^{\tilde g}$. Then $U_0\in C(\mathfrak{B},D^{g})$. Moreover $U_0$ is a
critical point of $\mathcal{A}$ on $C(\mathfrak{B},D^{g})$, for the following reason. Let $t\mapsto
U(t)$ be a differentiable curve in $C(\mathfrak{B},D^{g})$ with $U(0)=U_0$. Then $t\mapsto
U(t)^{\tilde g^{-1}}$ is a differentiable curve in $C(\mathfrak{B},D)$ through $U(\mathfrak{B},D)$.
Along it $\mathcal{A}$ takes the same values as along $t\mapsto U(t)$. Its derivative at $t=0$
vanishes because $U(\mathfrak{B},D)$ is critical on $C(\mathfrak{B},D)$.

The plaquette variables of $U_0$ are conjugates of those of $U(\mathfrak{B},D)$. Hence, for every
plaquette $\mathfrak{p}$ of $X$,
\begin{equation}\label{4.12:eq-minimiser-covariance-1}
\begin{split}
|\partial U_0(\mathfrak{p})-1|
&=|\partial U(\mathfrak{B},D)(\mathfrak{p})-1|\\
&<B_3\varepsilon_1\eta^{2}\ell_{\mathfrak{B}}(\mathfrak{p})^{-2}
\le\varepsilon_0\eta^{2}\ell_{\mathfrak{B}}(\mathfrak{p})^{-2}.
\end{split}
\end{equation}
Here the strict inequality is the bound of clause~(i) of \eqref{4.12:eq-constrained-minimisers-b}
for the minimiser with data $D$. The last inequality is $B_3\varepsilon_1\le\varepsilon_0$.

Since $D^{g}$ satisfies the hypotheses, clause~(ii) of \eqref{4.12:eq-constrained-minimisers-b}
applies to $D^{g}$ and to the critical point $U_0$. By \eqref{4.12:eq-minimiser-covariance-1} it
gives a $G$-valued function $\sigma'$ with
\begin{equation*}
U_0=U(\mathfrak{B},D^{g})^{\sigma'}.
\end{equation*}
Applying the gauge transformation $\sigma'^{-1}$ to both sides, and using
$(U^{\tilde g})^{\sigma'^{-1}}=U^{\sigma'^{-1}\tilde g}$ with the pointwise product, we obtain
\begin{equation*}
U(\mathfrak{B},D^{g})=U(\mathfrak{B},D)^{\sigma'^{-1}\tilde g}.
\end{equation*}
This is the assertion with $\sigma=\sigma'^{-1}\tilde g$.
\end{proof}

\begin{lemma}[An arithmetic bound on the gain factors]\label{4.12:lem-gain-arithmetic}
Let $L$ be the blocking factor, let $L_0\ge 2$ be a number with $L\ge 3L_0^2$, let $k$ be the
level, let $N_0\ge 1$ be an integer, and put $k_0=k-N_0$. For integers $m\ge 1$ put
\begin{equation*}
f_1(m)=(1+m)^{1/2}L^{-m},\qquad f_2(m)=(1+m)^{1/2}L^{-2m}.
\end{equation*}
Then the following hold for $i=1,2$.
\begin{enumerate}
\item For every integer $n$ with $k_0<n\le k-1$,
\begin{equation}\label{4.12:eq-gain-arithmetic-1}
\frac{f_i(k-n)}{L_0^{2(n-k_0)}}\le L_0^{-2N_0}.
\end{equation}
\item For every integer $n\le k_0$, so that $k-n\ge N_0$,
\begin{equation}\label{4.12:eq-gain-arithmetic-2}
f_i(k-n)\le (1+N_0)^{1/2}L^{-N_0}.
\end{equation}
\end{enumerate}
\end{lemma}

\begin{proof}
Since $L\ge 1$, we have $L^{-2m}\le L^{-m}$ for every $m\ge 1$, hence $f_2(m)\le f_1(m)$.
It therefore suffices in both clauses to prove the bound for $f_1$.

Clause (1). Let $k_0<n\le k-1$ and put $m=k-n$, so that $m\ge 1$. From $L\ge 3L_0^2$ we get
$L^{-m}\le 3^{-m}L_0^{-2m}$. Moreover
\begin{equation*}
m+(n-k_0)=(k-n)+(n-k_0)=k-k_0=N_0 .
\end{equation*}
Consequently
\begin{equation*}
\begin{aligned}
f_1(m)\,L_0^{-2(n-k_0)}
&\le (1+m)^{1/2}3^{-m}\,L_0^{-2m-2(n-k_0)}\\
&=(1+m)^{1/2}3^{-m}\,L_0^{-2N_0}.
\end{aligned}
\end{equation*}
It remains to see that $(1+m)^{1/2}3^{-m}\le \sqrt{2}/3$ for $m\ge 1$; since $\sqrt2/3<1$, this gives
\eqref{4.12:eq-gain-arithmetic-1} for $i=1$. For $m\ge1$ the ratio of consecutive values is
\begin{equation*}
\frac{(2+m)^{1/2}3^{-m-1}}{(1+m)^{1/2}3^{-m}}
=\frac13\Bigl(\frac{2+m}{1+m}\Bigr)^{1/2}\le\frac13\Bigl(\frac32\Bigr)^{1/2}<1,
\end{equation*}
because $(2+m)/(1+m)$ decreases in $m$ and equals $3/2$ at $m=1$. Hence the sequence
$(1+m)^{1/2}3^{-m}$ is decreasing on $m\ge1$, and its largest value is the value $\sqrt2/3$ at
$m=1$. The bound for $f_2$ follows from $f_2\le f_1$.

Clause (2). The function $f_1$ is decreasing on the integers $m\ge 1$: for such $m$,
\begin{equation*}
\frac{f_1(m+1)}{f_1(m)}=\frac1L\Bigl(\frac{2+m}{1+m}\Bigr)^{1/2}\le\frac{(3/2)^{1/2}}{L}<1,
\end{equation*}
where the first inequality uses again that $(2+m)/(1+m)\le 3/2$ for $m\ge1$, and the last uses
$L\ge 3L_0^2\ge 12$, a consequence of the hypotheses $L\ge 3L_0^2$ and $L_0\ge2$. Now let
$n\le k_0$. Then $k-n\ge k-k_0=N_0\ge 1$, and monotonicity gives
\begin{equation*}
f_1(k-n)\le f_1(N_0)=(1+N_0)^{1/2}L^{-N_0},
\end{equation*}
which is \eqref{4.12:eq-gain-arithmetic-2} for $i=1$. The bound for $f_2$ follows from
$f_2\le f_1$.
\end{proof}

\begin{proposition}[The re-blocking step inside admissible boxes.
Stated; proof incomplete at the step marked $(\ast)$]
\label{4.12:prop-reblocking-inside}
Let $d\ge 2$, and let $G$ be a compact gauge group. Let the parameters
$L, L_0, \beta, \beta_0, M, M_1, C_0, q_0$ satisfy two sets of conditions:
\begin{itemize}
\item the two standing conditions under which the source space
$\tilde{U}^c_k(Y_4,\tilde{\alpha}_0,\tilde{\alpha}_1)$ of Definition~\ref{4.12:not-admissible-boxes}
is constructed, as they are fixed with the theorem that constructs it;
\item the one-sided floors $(\Phi 0)$--$(\Phi 2)$ of Definition~\ref{4.12:def-coupling-thresholds}.
\end{itemize}
Then there is a threshold
\begin{equation*}
N_* = N_*(d, L, L_0, M, \beta, \beta_0) < \infty
\end{equation*}
(equivalently, by the relation between $N_0$ and the running couplings \cite{B-CMP122a}, a threshold
$\gamma_*>0$ on the running couplings) with the following property. Suppose
the parameter $N_0$ of the construction satisfies $N_0\ge N_*$ at the step $k$ under consideration
(equivalently, $\sup_m g_m\le\gamma_*$). Then the conditions $(\Gamma_0)$--$(\Gamma_3)$ of
Definition~\ref{4.12:def-coupling-thresholds} are in force, as shown there, and under them the
statements below are true.

Let $(\mathbb{U},\mathbb{J})$ be any point of the source space
$\tilde{U}^c_k(Y_4,\tilde{\alpha}_0,\tilde{\alpha}_1)$, and let $\bar{U}$ be its $G$-valued factor. Let
$\mathfrak{Q}\subset Y_4$ be an admissible box, and let $\mathcal{C}$ be the core. All of these are as in
Definition~\ref{4.12:not-admissible-boxes}. Let $E$ be a consumer of one of the classes
$(K1)$--$(K5)$ of Definition~\ref{4.12:def-consumer-radii}, and suppose its domain $X'$ satisfies
\begin{equation*}
X'\subset\mathfrak{Q}\setminus\mathcal{C}.
\end{equation*}
Fix a slot level $p$ of the space of $E$ and a sub-domain $X''\subset X'$ read by $E$ in the sense of
Definition~\ref{4.12:def-consumer-radii}. Write
$\mathfrak{B}_E(p,X'')$ for the determining set of $E$ (Definition~\ref{4.12:def-determining-sets}).
Write $U_E=U_E(p,X'')$ for the constrained minimiser along it
(Definition~\ref{4.12:def-constrained-minimisers}):
\begin{equation*}
U_E(p,X'') := U\bigl(\mathfrak{B}_E(p,X''),\, M_{\mathfrak{B}_E(p,X'')}(\bar{U})\bigr).
\end{equation*}
Then the following hold for every such $p$ and $X''$.
\begin{enumerate}
\item[(a)] $U_E(p,X'')$ exists. It is obtained from Ba{\l}aban's theorem on constrained minimisers
\cite[Theorem~1]{B-CMP102b}, in the form \eqref{4.12:eq-constrained-minimisers-b}, with
$\varepsilon_1 = 2\tilde{\alpha}_{0,k}\ell_p^2$.
\item[(b)] Let $\mathfrak{p}\subset X''^{\sim 2}$ be a plaquette and $\mathfrak{b}\subset X''^{\sim 2}$ a bond.
Put $n:=n_E(\mathfrak{p})$ or $n:=n_E(\mathfrak{b})$ respectively. If $n\ge p$, then
\begin{equation*}
|\partial U_E(\mathfrak{p})-1|\le 2L^2B_3\tilde{\alpha}_{0,k}\eta^2,
\qquad
|J(U_E)(\mathfrak{b})|\le 2L^3B_3\tilde{\alpha}_{0,k}\ell_p^{-1}.
\end{equation*}
If $n<p$, then
\begin{equation*}
|\partial U_E(\mathfrak{p})-1|\le 3L^2B_HK_b\tilde{\alpha}_{0,k}\eta^2\ell_n^{-1},
\qquad
|J(U_E)(\mathfrak{b})|\le 3L^3B_HK_b\tilde{\alpha}_{0,k}\ell_n^{-2}.
\end{equation*}
Here the constant $K_b$ is
\begin{equation*}
K_b:=\frac{\pi}{3}(d-1)\,c_Q\,c_d\,(1+3L^2)\,D_{\mathfrak{Q}}.
\end{equation*}
\item[(c)] Consequently, for every plaquette $\mathfrak{p}$ and every bond $\mathfrak{b}$ contained in
$X''^{\sim 2}$,
\begin{equation*}
|\partial U_E(p,X'')(\mathfrak{p})-1|<\operatorname{tgt}_2^E(\mathfrak{p}),
\qquad
|J(U_E(p,X''))(\mathfrak{b})|<\operatorname{tgt}_4^E(\mathfrak{b}).
\end{equation*}
Here $\operatorname{tgt}_2^E$ and $\operatorname{tgt}_4^E$ are the radii of $E$ for the two derived slots of the
$G$-valued factor (Definition~\ref{4.12:def-consumer-radii}).
\end{enumerate}
\end{proposition}

The proof establishes two further bounds. Put $\mathfrak{B}:=\mathfrak{B}_E(p,X'')$, and let $b$ be the datum
on the bonds of $\mathbb{L}(\mathfrak{B})$ constructed in a later step of the proof of
Proposition~\ref{4.12:prop-reblocking-inside}.

The first bound holds for every bond $\mathbf{b}$ of $\mathbb{L}(\mathfrak{B})$ with
$\ell(\mathbf{b})=\ell_i$:
\begin{equation}\label{4.12:eq-reblocking-inside-c}
|b(\mathbf{b})|\le c_Q\,\mathfrak{a}_{\max}\,\ell(\mathbf{b})
=\frac{\pi}{3}(d-1)\,c_Q\,D_{\mathfrak{Q}}\cdot\tilde{\alpha}_{0,k}\cdot\ell_i .
\end{equation}
The second holds for every plaquette or bond $\xi\subset X''^{\sim 2}$ with $n_E(\xi)=i<p$. The first line
refers to a plaquette $\xi$, the second to a bond $\xi$:
\begin{equation}\label{4.12:eq-reblocking-inside-d}
\begin{aligned}
\frac{|\partial U_E(\xi)-1|}{\operatorname{tgt}_2^E(\xi)}
&\le \frac{3L^2B_HK_b(1+\beta_0)(1+k-i)^{1/2}L^{-(k-i)}}{\vartheta_E(\xi)},\\
\frac{|J(U_E)(\xi)|}{\operatorname{tgt}_4^E(\xi)}
&\le \frac{3L^3B_HK_b(1+\beta_0)(1+k-i)^{1/2}L^{-(k-i)}}{\vartheta_E(\xi)}.
\end{aligned}
\end{equation}

\begin{proof}[Proof of clause \textup{(a)} and of the block-scale bound]
Fix $E$, $p$ and $X''$ with $X''\subset X'\subset\mathfrak{Q}\setminus\mathcal{C}$, as in the statement. Write
\begin{equation*}
\mathfrak{B}:=\mathfrak{B}_E(p,X''),\qquad V:=M_{\mathfrak{B}}(\bar{U}).
\end{equation*}

\emph{Existence of the minimiser.} All blocks of $\mathfrak{B}$ lie in $X''$, and $X''\subset\mathfrak{Q}$.
The box $\mathfrak{Q}$ is admissible, so the level $\nu$ is identically equal to $k$ on $\mathfrak{Q}$.
At this level, the bound \eqref{4.12:eq-admissible-boxes-a} on the $G$-valued factor of the source datum
gives
\begin{equation*}
|\partial\bar{U}(\mathfrak{p})-1|<\tilde{\alpha}_{0,k}\eta^2
\end{equation*}
for every plaquette $\mathfrak{p}$ there.

We apply the averaging bound \eqref{4.12:eq-determining-sets-e} to $\bar{U}$. This is permitted by
condition~(i) of \eqref{4.12:eq-coupling-thresholds-d}. It gives, for every plaquette $P$ of
$\mathbb{L}(\mathfrak{B})$,
\begin{equation}\label{4.12:eq-reblocking-inside-1}
|\partial V(P)-1|
\le 2\tilde{\alpha}_{0,k}\eta^2\Bigl(\frac{\ell(P)}{\eta}\Bigr)^2
= 2\tilde{\alpha}_{0,k}\ell(P)^2
\le 2\tilde{\alpha}_{0,k}\ell_p^2 =:\varepsilon_1 .
\end{equation}
Every block of $\mathfrak{B}$ has level at most $p$, so $\ell(P)\le\ell_p\le 1$, which gives the last
inequality.

By \eqref{4.12:eq-constrained-minimisers-a}, which is the form in which hypothesis
\cite[(7)]{B-CMP102b} is read here, \eqref{4.12:eq-reblocking-inside-1} says that $V$ satisfies this
hypothesis with the value of $\varepsilon_1$ just defined. By condition~(ii) of
\eqref{4.12:eq-coupling-thresholds-d},
\begin{equation*}
\varepsilon_1\le a_1,\qquad B_3\varepsilon_1\le\varepsilon_0 .
\end{equation*}
So the hypotheses of clause~(i) of \eqref{4.12:eq-constrained-minimisers-b} are met. That clause is
Ba{\l}aban's theorem \cite[Theorem~1]{B-CMP102b}. It yields the constrained minimiser
\begin{equation*}
U_E:=U_E(p,X'')=U(\mathfrak{B},V).
\end{equation*}
This minimiser is $G$-valued, and it satisfies, for all plaquettes $\mathfrak{p}$ and all bonds
$\mathfrak{b}$ in $X''$,
\begin{equation}\label{4.12:eq-reblocking-inside-a}
\begin{aligned}
|\partial U_E(\mathfrak{p})-1|
&<2B_3\tilde{\alpha}_{0,k}\ell_p^2\cdot\eta^2\ell_{\mathfrak{B}}(\mathfrak{p})^{-2},\\
|J(U_E)(\mathfrak{b})|
&<2B_3\tilde{\alpha}_{0,k}\ell_p^2\cdot\ell_{\mathfrak{B}}(\mathfrak{b})^{-3}.
\end{aligned}
\end{equation}
This proves clause~(a).

\emph{The block scale on $X''^{\sim 2}$.} Let $\xi$ be a plaquette or a bond contained in $X''^{\sim 2}$,
and let $y$ be a site of $\xi$. By \eqref{4.12:eq-consumer-radii-b},
\begin{equation*}
\operatorname{dist}(y,\partial X'')\ge 2M\ell_{c_E(y)} .
\end{equation*}
By $(\Phi 1)$, that is by \eqref{4.12:eq-coupling-thresholds-b}, we have
$2M\ell_{c_E(y)}\ge 2M_1\ell_{c_E(y)}$. Therefore \eqref{4.12:eq-determining-sets-b} gives
\begin{equation*}
t_{X''}(y)\ge c_E(y).
\end{equation*}
Hence, by \eqref{4.12:eq-determining-sets-a}, the block of $\mathfrak{B}$ at $y$ has level
\begin{equation*}
\min\bigl(p,\,c_E(y),\,t_{X''}(y)\bigr)=\min\bigl(p,\,c_E(y)\bigr).
\end{equation*}

Now consider the sites of $\xi$ together with their nearest neighbours. By
\eqref{4.12:eq-determining-sets-a}, they lie in blocks whose levels differ by at most one from the levels
of the blocks at the sites of $\xi$. For a neighbour $y'$ of a site of $\xi$, the inequality
$c_E(y')\ge n_E(\xi)-1$ cannot be asserted. The following two facts suffice instead:
\begin{itemize}
\item the largest level among the blocks at the sites of $\xi$ is $\min(p,n_E(\xi))$;
\item every block meeting a site of $\xi$, or a neighbour of one, has level at least
$\min(p,n_E(\xi))-1$.
\end{itemize}
The smallest block side near $\xi$ is therefore at least
\begin{equation*}
\ell_{\min(p,n_E(\xi))-1}=\ell_{\min(p,n_E(\xi))}/L ,
\end{equation*}
since $\ell_m=L^{m-k}$. Thus
\begin{equation}\label{4.12:eq-reblocking-inside-b}
\ell_{\mathfrak{B}}(\xi)\ge \ell_{\min(p,\,n_E(\xi))}/L
\qquad\text{for every plaquette or bond } \xi\subset X''^{\sim 2}.
\end{equation}

The bounds \eqref{4.12:eq-reblocking-inside-a} and \eqref{4.12:eq-reblocking-inside-b} are the two
inputs from which clauses~(b) and~(c), and the bounds \eqref{4.12:eq-reblocking-inside-c} and
\eqref{4.12:eq-reblocking-inside-d}, are derived in what follows.
\end{proof}

\begin{remark}
In particular, clause~(c) of Proposition~\ref{4.12:prop-reblocking-inside} yields one half of clause~(iv) of the space clause $D_E(X')$ of the consumer $E$, namely the half that concerns the $G$-valued factor. It does so for the composites $\mathcal{W}_1,\dots,\mathcal{W}_4$ of the construction of the source space.

The reason is that the $G$-valued factor of each $\mathcal{W}_i$ off the core $\mathcal{C}$ is $\bar{U}$. For every consumer $E$ of the classes $(K1)$--$(K5)$ whose domain satisfies
\begin{equation*}
X'\subset\mathfrak{Q}\setminus\mathcal{C},
\end{equation*}
the bounds of clause~(c) are therefore exactly the bounds that this half of clause~(iv) of $D_E(X')$ requires.
\end{remark}

\begin{lemma}[Points at or above the slot level]\label{4.12:prf-above-slot-level}
Assume the hypotheses of Proposition~\ref{4.12:prop-reblocking-inside}, whose proof is incomplete
at the step marked $(\ast)$ there. Let $E$ be a consumer as there, let $p$ be a slot level of $E$ and let
$X''\subset X'$ be a sub-domain read by $E$. Let $U_E=U_E(p,X'')$ be the minimiser provided by
clause~(a) of that proposition. Let $\mathfrak{p}\subset X''^{\sim 2}$ be a plaquette and
$\mathfrak{b}\subset X''^{\sim 2}$ a bond. Put $n:=n_E(\mathfrak{p})$, respectively
$n:=n_E(\mathfrak{b})$, and suppose that $n\ge p$. Then
\begin{equation}\label{4.12:eq-above-slot-level-1}
|\partial U_E(\mathfrak{p})-1|<2L^2B_3\tilde{\alpha}_{0,k}\eta^2,
\qquad
|J(U_E)(\mathfrak{b})|<2L^3B_3\tilde{\alpha}_{0,k}\ell_p^{-1},
\end{equation}
and
\begin{equation}\label{4.12:eq-above-slot-level-2}
\frac{|\partial U_E(\mathfrak{p})-1|}{\operatorname{tgt}_2^E(\mathfrak{p})}\le\frac12,
\qquad
\frac{|J(U_E)(\mathfrak{b})|}{\operatorname{tgt}_4^E(\mathfrak{b})}\le\frac12 .
\end{equation}
\end{lemma}

\begin{proof}
Since $n\ge p$, the minimum of $p$ and $n$ equals $p$. The bounds
\eqref{4.12:eq-reblocking-inside-a} and \eqref{4.12:eq-reblocking-inside-b} of
Proposition~\ref{4.12:prop-reblocking-inside} then give
\begin{align*}
|\partial U_E(\mathfrak{p})-1|
&<2B_3\tilde{\alpha}_{0,k}\ell_p^2\cdot\eta^2L^2\ell_p^{-2}=2L^2B_3\tilde{\alpha}_{0,k}\eta^2,\\
|J(U_E)(\mathfrak{b})|
&<2B_3\tilde{\alpha}_{0,k}\ell_p^2\cdot L^3\ell_p^{-3}=2L^3B_3\tilde{\alpha}_{0,k}\ell_p^{-1}.
\end{align*}
This is \eqref{4.12:eq-above-slot-level-1}, that is, clause~(b) of
Proposition~\ref{4.12:prop-reblocking-inside} for $n\ge p$.

We compare these bounds with the radii of $E$. By Definition~\ref{4.12:def-consumer-radii}, for
$\xi=\mathfrak{p}$ or $\xi=\mathfrak{b}$ and $n=n_E(\xi)$,
\begin{equation*}
\operatorname{tgt}_2^E(\xi)=\vartheta_E(\xi)\,\tilde{\alpha}_{0,n}\,\eta^2\ell_n^{-2},
\qquad
\operatorname{tgt}_4^E(\xi)=\vartheta_E(\xi)\,\tilde{\alpha}_{0,n}\,\ell_n^{-3}L^{(n-p)_+}.
\end{equation*}
Since $\ell_m=L^{m-k}$ and $n\ge p$, we have $L^{(n-p)_+}=L^{n-p}=\ell_n/\ell_p$, so that
$\operatorname{tgt}_4^E(\mathfrak{b})$ equals
$\vartheta_E(\mathfrak{b})\tilde{\alpha}_{0,n}\ell_n^{-2}\ell_p^{-1}$. Dividing the bounds
\eqref{4.12:eq-above-slot-level-1} by the radii therefore gives
\begin{equation}\label{4.12:eq-above-slot-level-3}
\frac{|\partial U_E(\mathfrak{p})-1|}{\operatorname{tgt}_2^E(\mathfrak{p})}
<2L^2B_3\,\frac{\tilde{\alpha}_{0,k}}{\tilde{\alpha}_{0,n}}\,
\frac{\ell_n^2}{\vartheta_E(\mathfrak{p})},
\qquad
\frac{|J(U_E)(\mathfrak{b})|}{\operatorname{tgt}_4^E(\mathfrak{b})}
<2L^3B_3\,\frac{\tilde{\alpha}_{0,k}}{\tilde{\alpha}_{0,n}}\,
\frac{\ell_n^2}{\vartheta_E(\mathfrak{b})}.
\end{equation}
By \eqref{4.12:eq-running-radii-a} of Notation~\ref{4.12:not-running-radii},
\begin{equation*}
\frac{\tilde{\alpha}_{0,k}}{\tilde{\alpha}_{0,n}}\le(1+\beta_0)(1+k-n)^{1/2}.
\end{equation*}
Since $L\ge 1$, the factor $2L^2B_3$ in \eqref{4.12:eq-above-slot-level-3} may be replaced by
$2L^3B_3$. We distinguish two cases.

\emph{Case $n\le k-1$.} Here $\ell_n^2=L^{-2(k-n)}$. The lower bound
\eqref{4.12:eq-consumer-radii-d} of Definition~\ref{4.12:def-consumer-radii} gives
$\vartheta_E(\xi)\ge\theta_*\hat{\vartheta}(n)$. Put
\begin{equation*}
\Theta:=\sup_{n\le k-1}\frac{(1+k-n)^{1/2}L^{-2(k-n)}}{\hat{\vartheta}(n)}.
\end{equation*}
Then both right-hand sides of \eqref{4.12:eq-above-slot-level-3} are at most
$2L^3B_3(1+\beta_0)\Theta/\theta_*$.

In the notation of Lemma~\ref{4.12:lem-gain-arithmetic}, the numerator in $\Theta$ is $f_2(k-n)$,
with $k-n\ge1$. That lemma holds under its hypotheses $L\ge 3L_0^2$ and $L_0\ge 2$, and it gives
two bounds:
\begin{itemize}
\item for $k_0<n\le k-1$, where $\hat{\vartheta}(n)=L_0^{2(n-k_0)}$, the quotient
$f_2(k-n)/\hat{\vartheta}(n)$ is at most $L_0^{-2N_0}$;
\item for $n\le k_0$, where $\hat{\vartheta}(n)=1$ and $k-n\ge N_0$, the quotient is at most
$(1+N_0)^{1/2}L^{-N_0}$.
\end{itemize}
Hence $\Theta\le W(N_0)$, with $W(N_0)$ as in \eqref{4.12:eq-coupling-thresholds-h}.
Condition \eqref{4.12:eq-coupling-thresholds-e} of
Definition~\ref{4.12:def-coupling-thresholds} states
$4L^3B_3(1+\beta_0)W(N_0)\le\theta_*$. Therefore
\begin{equation*}
\frac{2L^3B_3(1+\beta_0)\Theta}{\theta_*}
\le\frac{2L^3B_3(1+\beta_0)W(N_0)}{\theta_*}\le\frac12,
\end{equation*}
and both ratios in \eqref{4.12:eq-above-slot-level-2} are at most $\tfrac12$.

\emph{Case $n=k$.} By the observation recorded with Definition~\ref{4.12:def-consumer-radii},
inside an admissible box $\mathfrak{Q}$, which contains $X''$, the only consumers having points of
level $k$ are those of class (K4), and they have them only at points of $B_5\cap\mathfrak{Q}$.
So this case occurs only when $E$ is of class (K4) and $\xi\subset B_5\cap\mathfrak{Q}$.
There \eqref{4.12:eq-consumer-radii-e} of Definition~\ref{4.12:def-consumer-radii} applies and
gives $\vartheta_E(\xi)\ge(1-2\beta)L_0^{2N_0}$. Moreover
$\tilde{\alpha}_{0,k}/\tilde{\alpha}_{0,n}=1$ and $\ell_n=\ell_k=1$. By
\eqref{4.12:eq-above-slot-level-3}, both ratios are therefore less than
$2L^3B_3/\bigl((1-2\beta)L_0^{2N_0}\bigr)$. Condition \eqref{4.12:eq-coupling-thresholds-g} of
Definition~\ref{4.12:def-coupling-thresholds} states $4L^3B_3\le(1-2\beta)L_0^{2N_0}$, so this
quantity is at most $\tfrac12$.

In both cases \eqref{4.12:eq-above-slot-level-2} holds, which completes the proof.
\end{proof}

\begin{lemma}[Points below the slot level: direct expansion]\label{4.12:prf-below-slot-level}
Assume the hypotheses of Proposition~\ref{4.12:prop-reblocking-inside}.
Let $(\mathbb{U},\mathbb{J})$,
its $G$-valued factor $\bar{U}$, the admissible box $\mathfrak{Q}$ and the consumer $E$ with
$X'\subset\mathfrak{Q}\setminus\mathcal{C}$ be as there. Fix a slot level $p$ of $E$ and a
sub-domain $X''\subset X'$ read by $E$. Write
\begin{equation*}
\mathfrak{B}:=\mathfrak{B}_E(p,X''),\qquad V:=M_{\mathfrak{B}}(\bar{U}),\qquad
U_E:=U_E(p,X''),
\end{equation*}
and put
\begin{equation*}
\mathfrak{a}_{\max}:=\tfrac{\pi}{3}(d-1)D_{\mathfrak{Q}}\tilde{\alpha}_{0,k},\qquad
K_b:=\tfrac{\pi}{3}(d-1)c_Qc_d(1+3L^2)D_{\mathfrak{Q}}.
\end{equation*}
Then the following hold.
\begin{enumerate}
\item[(a)] There are a $\mathfrak{g}$-valued $b$ on the bonds of $\mathbb{L}(\mathfrak{B})$ and a
$G$-valued $\varsigma$ on the base points with $V=(e^{ib})^{\varsigma}$. This $b$ obeys
\eqref{4.12:eq-reblocking-inside-c}, and $\|b\|(x)\le K_b\tilde{\alpha}_{0,k}\ell_{c_E(x)}$ for
every site $x$ of $X''$; this is \eqref{4.12:eq-below-slot-level-a} below.
\item[(b)] Let $\mathfrak{p}\subset X''^{\sim 2}$ be a plaquette with $n:=n_E(\mathfrak{p})<p$.
Then
\begin{equation*}
|\partial U_E(\mathfrak{p})-1|\le 3L^2B_HK_b\tilde{\alpha}_{0,k}\eta^2/\ell_n .
\end{equation*}
Let $\mathfrak{b}\subset X''^{\sim 2}$ be a bond with $n:=n_E(\mathfrak{b})<p$. Then
\begin{equation*}
|J(U_E)(\mathfrak{b})|\le 3L^3B_HK_b\tilde{\alpha}_{0,k}\ell_n^{-2}.
\end{equation*}
\item[(c)] For such plaquettes and bonds the bounds \eqref{4.12:eq-reblocking-inside-d} hold,
and both quotients estimated there are at most $\tfrac12$.
\end{enumerate}
\end{lemma}

\begin{proof}
Recall that $\eta=L^{-k}$ and $\ell_m=L^{m-k}$. Hence $\ell_{m+1}=L\ell_m$ and $\ell_0=\eta$.
The map $m\mapsto\ell_m$ is increasing, and $\ell_m\le 1$ for $m\le k$; in particular
$\eta\le 1$.

Let $\xi\subset X''^{\sim 2}$ be a plaquette or a bond with $n:=n_E(\xi)<p$. Let $x_0$ be its
lower corner if $\xi$ is a plaquette, and its initial point if $\xi$ is a bond. Recall that
$c_E(x_0)\le n$. The proof is a direct expansion at the trivial background and proceeds in
eight steps.

(i) The box gauge. We apply the axial gauge of a box (Lemma~\ref{4.12:lem-axial-gauge}) to
$\bar{U}$ on the admissible box $\mathfrak{Q}$, with $D:=D_{\mathfrak{Q}}$ and
$\varphi:=\tilde{\alpha}_{0,k}\eta^2$.

The hypotheses of that lemma hold. The plaquette bound \eqref{4.12:eq-admissible-boxes-a}
gives $|\partial\bar{U}(\mathfrak{p})-1|\le\tilde{\alpha}_{0,k}\eta^2=\varphi$ for all
$\mathfrak{p}\subset\mathfrak{Q}$. The edges of $\mathfrak{Q}$ are at most $D_{\mathfrak{Q}}$.
For the smallness condition, note that \eqref{4.12:eq-coupling-thresholds-d} contains
$\tfrac{\pi}{3}(d-1)D_{\mathfrak{Q}}\tilde{\alpha}_{0,k}\le 1$, hence
$(d-1)D_{\mathfrak{Q}}\tilde{\alpha}_{0,k}\le 3/\pi<1$. Together with $\eta\le1$ this
gives
\begin{equation*}
(d-1)(D_{\mathfrak{Q}}/\eta)\,\tilde{\alpha}_{0,k}\eta^2
=(d-1)D_{\mathfrak{Q}}\tilde{\alpha}_{0,k}\,\eta\le 1 .
\end{equation*}

The lemma provides a $G$-valued $a$ on the sites of $\mathfrak{Q}$ such that
$\bar{U}^{a}=e^{i\eta\mathfrak{a}}\cdot 1$ on $\mathfrak{Q}$, with $\mathfrak{a}$ a
$\mathfrak{g}$-valued $1$-form. With the present $\varphi$ and $D$, its bound on $\mathfrak{a}$
reads
\begin{equation}\label{4.12:eq-below-slot-level-1}
\sup_{\mathfrak{Q}}|\mathfrak{a}|\le\mathfrak{a}_{\max}
=\tfrac{\pi}{3}(d-1)D_{\mathfrak{Q}}\tilde{\alpha}_{0,k}.
\end{equation}

(ii) The data in the box gauge. Since $\bar{U}^{a}=e^{i\eta\mathfrak{a}}\cdot1$ on
$\mathfrak{Q}$, we have $\bar{U}=(e^{i\eta\mathfrak{a}})^{a^{-1}}$ there. Let $\bar{a}$ denote
the restriction of $a$ to the base points. The gauge covariance
\eqref{4.12:eq-determining-sets-c} of the averaging map, applied with $\sigma=a^{-1}$, gives
\begin{equation*}
V=M_{\mathfrak{B}}(\bar{U})
=M_{\mathfrak{B}}\bigl((e^{i\eta\mathfrak{a}})^{a^{-1}}\bigr)
=M_{\mathfrak{B}}(e^{i\eta\mathfrak{a}})^{\bar{a}^{-1}} .
\end{equation*}

Next we use the estimate \eqref{4.12:eq-determining-sets-d}; for interface bonds it is used in
the reading for interface bonds given in Definition~\ref{4.12:def-determining-sets}. It is
applicable because
\eqref{4.12:eq-coupling-thresholds-d} contains $\mathfrak{a}_{\max}\le c_2$. It yields a
$G$-valued $u$ on the base points and a $\mathfrak{g}$-valued $b$ on the bonds of
$\mathbb{L}(\mathfrak{B})$ such that $M_{\mathfrak{B}}(e^{i\eta\mathfrak{a}})=(e^{ib})^{u}$ and
\begin{equation}\label{4.12:eq-below-slot-level-2}
|b(\mathbf{b})|\le c_Q\,\mathfrak{a}_{\max}\,\ell(\mathbf{b})
\qquad\text{for every bond $\mathbf{b}$ of $\mathbb{L}(\mathfrak{B})$.}
\end{equation}
This is \eqref{4.12:eq-reblocking-inside-c}. Hence $V=(e^{ib})^{\varsigma}$, where
$\varsigma:=\bar{a}^{-1}u$ is $G$-valued on the base points.

Every block of $\mathbb{L}(\mathfrak{B})$ has side at most $\ell_p$, and $p\le k$, as every
slot level is at most $k$ (Definition~\ref{4.12:def-consumer-radii}). So
$\ell(\mathbf{b})\le\ell_p\le 1$, and \eqref{4.12:eq-below-slot-level-2} gives
$\sup|b|\le c_Q\mathfrak{a}_{\max}$. By \eqref{4.12:eq-coupling-thresholds-d},
$c_Q\mathfrak{a}_{\max}\le b_H/2<b_H$.

Now fix a plaquette $P$ of $\mathbb{L}(\mathfrak{B})$, and let $v_1,\dots,v_4$ be the factors
$e^{\pm ib(\mathbf{b})}$, $\mathbf{b}\in P$, of $\partial(e^{ib})(P)$, in the order in which
they enter. Then
\begin{equation*}
v_1v_2v_3v_4-1=(v_1-1)v_2v_3v_4+(v_2-1)v_3v_4+(v_3-1)v_4+(v_4-1).
\end{equation*}
Each $v_i$ is unitary, so the products of unitaries on the right have norm one; moreover
$|v^{-1}-1|=|v-1|$ for unitary $v$, since $v^{-1}-1=-v^{-1}(v-1)$. From the power series,
$|e^{ib(\mathbf{b})}-1|\le|b(\mathbf{b})|e^{|b(\mathbf{b})|}$. Hence
\begin{equation*}
|\partial(e^{ib})(P)-1|
\le\sum_{\mathbf{b}\in P}|e^{ib(\mathbf{b})}-1|
\le 4c_Q\mathfrak{a}_{\max}e^{c_Q\mathfrak{a}_{\max}}
\le 5c_Q\mathfrak{a}_{\max}.
\end{equation*}
The last step uses $c_Q\mathfrak{a}_{\max}\le 1/10$, which is part of
\eqref{4.12:eq-coupling-thresholds-d}, and $4e^{1/10}\le 5$.

We apply \eqref{4.12:eq-constrained-minimisers-a}, using the conditions
$5c_Q\mathfrak{a}_{\max}\le a_1$ and $5B_3c_Q\mathfrak{a}_{\max}\le\varepsilon_0$ of
\eqref{4.12:eq-coupling-thresholds-d}. It follows that the data $e^{ib}$ satisfy the hypotheses
of \eqref{4.12:eq-constrained-minimisers-b} with
$\varepsilon_1':=5c_Q\mathfrak{a}_{\max}$.

(iii) The minimiser in the box gauge. We apply the gauge covariance of the constrained minimiser
(Lemma~\ref{4.12:lem-minimiser-covariance}) with data $e^{ib}$ and gauge function $\varsigma$.
It shows that $V=(e^{ib})^{\varsigma}$ satisfies the same hypotheses, and that
\begin{equation*}
U_E=U(\mathfrak{B},V)=U(\mathfrak{B},e^{ib})^{\sigma}
\qquad\text{for a $G$-valued $\sigma$.}
\end{equation*}

The hypotheses of \eqref{4.12:eq-constrained-minimisers-c} on $b$ are $\sup|b|<b_H$ and that
$e^{ib}$ satisfies the hypotheses of \eqref{4.12:eq-constrained-minimisers-b}. Both were
established in step (ii). Hence
\begin{equation*}
U(\mathfrak{B},e^{ib})=(e^{i\eta\mathcal{H}}\cdot1)^{u'},
\end{equation*}
with $u'$ $G$-valued, and $\mathcal{H}:=\mathcal{H}_{\mathfrak{B}}(b)$ obeys the bounds of
\eqref{4.12:eq-constrained-minimisers-c}. Therefore
$U_E=(e^{i\eta\mathcal{H}}\cdot1)^{\sigma u'}$.

Two facts convert this into bounds on $U_E$: the invariance of $|\partial U(\mathfrak{p})-1|$
under gauge transformations, recorded in Notation~\ref{4.12:not-lattice-configurations}, and the
covariance \eqref{4.12:eq-lattice-configurations-b} of the current. Together they give
\begin{equation}\label{4.12:eq-below-slot-level-3}
|\partial U_E(\mathfrak{p})-1|=|\partial(e^{i\eta\mathcal{H}}\cdot1)(\mathfrak{p})-1|,
\qquad
|J(U_E)(\mathfrak{b})|=|J(e^{i\eta\mathcal{H}}\cdot1)(\mathfrak{b})| .
\end{equation}

(iv) The size of $\|b\|$. Let $x$ be a site of $X''$ and put $n':=c_E(x)$.

Suppose first that $n'<p$. Split
$\|b\|(x)=\sum_{\mathbf{b}}e^{-\delta d(x,\mathbf{b})}|b(\mathbf{b})|$ according to the value of
$\ell(\mathbf{b})$.

Bonds with $\ell(\mathbf{b})\le\ell_{n'+1}$. By \eqref{4.12:eq-below-slot-level-2}, each has
$|b(\mathbf{b})|\le c_Q\mathfrak{a}_{\max}\ell_{n'+1}$. Use
$e^{-\delta d}\le e^{-\delta d/2}$ and the summability \eqref{4.12:eq-constrained-minimisers-d}.
The contribution of these bonds is then at most
\begin{equation*}
c_Q\mathfrak{a}_{\max}\ell_{n'+1}\sum_{\mathbf{b}}e^{-\delta d(x,\mathbf{b})/2}
\le c_Q\mathfrak{a}_{\max}c_d\,L\ell_{n'} .
\end{equation*}

Bonds with $\ell(\mathbf{b})=\ell_m$ for some $m\ge n'+2$. Necessarily $m\le p$. By the
construction of $\mathfrak{B}_E(p,X'')$ in Definition~\ref{4.12:def-consumer-radii}, the larger
block of $\mathbf{b}$ has level $m$ at its sites, this level being the minimum of $p$, of $c_E$
and of the level given there by the boundary tower. Hence those sites have $c_E\ge m$. The
smaller block of $\mathbf{b}$ has level at
least $m-1$, so its sites have $c_E\ge m-1$.

For $m\ge n'+3$, the separation \eqref{4.12:eq-consumer-radii-a} therefore gives
$d(x,\mathbf{b})\ge(m-n'-2)M/M_1$. For $m=n'+2$ we use $d(x,\mathbf{b})\ge 0$. In both cases
$d(x,\mathbf{b})\ge(m-n'-2)_+M/M_1$, and so
\begin{equation*}
e^{-\delta d(x,\mathbf{b})}
\le e^{-\delta d(x,\mathbf{b})/2}\,e^{-\delta(m-n'-2)_+M/(2M_1)} .
\end{equation*}
We use \eqref{4.12:eq-below-slot-level-2} and apply
\eqref{4.12:eq-constrained-minimisers-d} to the bonds of each size $\ell_m$ separately. With
$\ell_m=L^{m-n'}\ell_{n'}$, the substitution $j:=m-n'$, and then $i:=j-2$, the contribution of
these bonds is at most
\begin{align*}
c_Q\mathfrak{a}_{\max}c_d\sum_{m\ge n'+2}\ell_m\,e^{-\delta(m-n'-2)_+M/(2M_1)}
&=c_Q\mathfrak{a}_{\max}c_d\,\ell_{n'}
\Bigl[L^2+\sum_{j\ge3}L^{j}e^{-\delta(j-2)M/(2M_1)}\Bigr]\\
&=c_Q\mathfrak{a}_{\max}c_d\,\ell_{n'}
\Bigl[L^2+L^2\sum_{i\ge1}\bigl(Le^{-\delta M/(2M_1)}\bigr)^{i}\Bigr]\\
&\le 2L^2c_Q\mathfrak{a}_{\max}c_d\,\ell_{n'} .
\end{align*}
In the last step, \eqref{4.12:eq-coupling-thresholds-c} makes the geometric series at most one.

Adding the two contributions, and using $L+2L^2\le 1+3L^2$ (that is, $L^2-L+1>0$) together with
the definition of $\mathfrak{a}_{\max}$, we obtain
\begin{equation}\label{4.12:eq-below-slot-level-a}
\|b\|(x)\le c_Q\mathfrak{a}_{\max}c_d(L+2L^2)\,\ell_{c_E(x)}
\le K_b\,\tilde{\alpha}_{0,k}\,\ell_{c_E(x)},
\qquad
K_b=\tfrac{\pi}{3}(d-1)c_Qc_d(1+3L^2)D_{\mathfrak{Q}} .
\end{equation}

Suppose now that $c_E(x)\ge p$. There are no blocks coarser than $\ell_p$, so every bond has
$\ell(\mathbf{b})\le\ell_p$. The same computation applies, with the first class of bonds taken
at $\ell(\mathbf{b})\le\ell_p$ and the second class empty. It gives
\begin{equation*}
\|b\|(x)\le c_Q\mathfrak{a}_{\max}c_d\,\ell_p
\le K_b\tilde{\alpha}_{0,k}\,\ell_{\min(p,c_E(x))}.
\end{equation*}
Since $\ell_{\min(p,c_E(x))}\le\ell_{c_E(x)}$, the bound \eqref{4.12:eq-below-slot-level-a}
holds at every site $x$ of $X''$. This completes the proof of (a).

(v) The size of $\mathcal{H}$ near $\xi$. We apply the bounds
\eqref{4.12:eq-constrained-minimisers-c} at the site $x_0$, at which $\xi$ sits. Three facts
enter:
\begin{itemize}
\item By \eqref{4.12:eq-reblocking-inside-b}, $\ell_{\mathfrak{B}}(\xi)\ge\ell_n/L$, since
$\min(p,n)=n$.
\item By \eqref{4.12:eq-below-slot-level-a} at $x_0$, $\|b\|(x_0)\le K_b\tilde{\alpha}_{0,k}\ell_n$.
Here we use $c_E(x_0)\le n<p$ and the monotonicity of $m\mapsto\ell_m$.
\item The bounds \eqref{4.12:eq-constrained-minimisers-c} themselves, in the form written out
below.
\end{itemize}
These give, for the bonds $\mathfrak{b}'$ of $\xi$ and their nearest neighbours in the first and
third lines, and at the site $x_0$ in the second line,
\begin{align*}
|\mathcal{H}(\mathfrak{b}')|&\le B_H\|b\|(x_0)/\ell_{\mathfrak{B}}(\xi)
\le LB_HK_b\tilde{\alpha}_{0,k},\\
|\nabla\mathcal{H}(x_0)|&\le B_H\|b\|(x_0)/\ell_{\mathfrak{B}}(\xi)^2
\le L^2B_HK_b\tilde{\alpha}_{0,k}/\ell_n,\\
|D^{*}D\mathcal{H}(\mathfrak{b}')|&\le B_H\|b\|(x_0)/\ell_{\mathfrak{B}}(\xi)^3
\le L^3B_HK_b\tilde{\alpha}_{0,k}/\ell_n^2 .
\end{align*}

(vi) The plaquette bound. Let $\xi=\mathfrak{p}$ be a plaquette. We apply the plaquettes of a
left-perturbed configuration (Lemma~\ref{4.12:lem-perturbed-plaquettes}) with $Y\equiv1$ and
$\mathcal{W}:=\mathcal{H}$, at $\mathfrak{p}=\mathfrak{p}(x_0;\mu,\nu)$.

With $Y\equiv 1$ we have $|\partial Y(\mathfrak{p})-1|=0$ and $\nabla^{Y}=\nabla$. The four bonds
of $\mathfrak{p}$ are bonds of $\xi$, so step (v) gives
$w_0\le LB_HK_b\tilde{\alpha}_{0,k}$. The condition $64LB_HK_b\tilde{\alpha}_{0,k}\le1$ of
\eqref{4.12:eq-coupling-thresholds-d}, together with $\eta\le1$, gives
$4\eta w_0\le 4LB_HK_b\tilde{\alpha}_{0,k}\le 1$. So the lemma applies.

Its first and third terms vanish. Each of the two covariant-derivative terms is bounded by
$|\nabla\mathcal{H}(x_0)|$, which step (v) controls. Hence
\begin{align*}
|\partial(e^{i\eta\mathcal{H}}\cdot1)(\mathfrak{p})-1|
&\le 2\eta^2\,L^2B_HK_b\tilde{\alpha}_{0,k}/\ell_n
+12\eta^2L^2B_H^2K_b^2\tilde{\alpha}_{0,k}^2\\
&\le 3L^2B_HK_b\tilde{\alpha}_{0,k}\,\eta^2/\ell_n .
\end{align*}
The last step holds because
$12L^2B_H^2K_b^2\tilde{\alpha}_{0,k}^2\le L^2B_HK_b\tilde{\alpha}_{0,k}/\ell_n$. This in turn
follows from
\begin{equation*}
12B_HK_b\tilde{\alpha}_{0,k}\ell_n\le 12B_HK_b\tilde{\alpha}_{0,k}\le 1,
\end{equation*}
where we used $\ell_n\le1$ and the condition $64LB_HK_b\tilde{\alpha}_{0,k}\le1$ of
\eqref{4.12:eq-coupling-thresholds-d}.

By \eqref{4.12:eq-below-slot-level-3}, this is the plaquette bound in (b). It is the plaquette
part of clause (b) of Proposition~\ref{4.12:prop-reblocking-inside} in the case $n<p$.

(vii) The current bound. Let $\xi=\mathfrak{b}$ be a bond. We apply the current bound
\eqref{4.12:eq-lattice-configurations-c} at $\mathfrak{b}$ with
\begin{equation*}
\ell:=\ell_n/L,\qquad w:=B_H\|b\|(x_0).
\end{equation*}

First, $\ell\ge\eta$. Indeed $n\ge1$, as every consumer level is at least one, so
\begin{equation*}
\ell=\ell_n/L=\ell_{n-1}\ge\ell_0=\eta .
\end{equation*}

Second, $w\le 1/64$. By \eqref{4.12:eq-below-slot-level-a} at $x_0$,
$w\le B_HK_b\tilde{\alpha}_{0,k}\ell_n\le B_HK_b\tilde{\alpha}_{0,k}$, and this is at most $1/64$
by the condition $64LB_HK_b\tilde{\alpha}_{0,k}\le1$ of \eqref{4.12:eq-coupling-thresholds-d}.

Third, the hypotheses
$\ell|\mathcal{H}|\le w$, $\ell^2|\nabla\mathcal{H}|\le w$ and $\ell^3|D^{*}D\mathcal{H}|\le w$
near $\mathfrak{b}$ are exactly the bounds \eqref{4.12:eq-constrained-minimisers-c} at $x_0$,
because $\ell=\ell_n/L\le\ell_{\mathfrak{B}}(\xi)$.

Hence
\begin{align*}
|J(e^{i\eta\mathcal{H}}\cdot1)(\mathfrak{b})|
&\le(2+2C_Jw)\,w\,(L/\ell_n)^3\\
&\le 3B_HK_b\tilde{\alpha}_{0,k}\ell_n\cdot L^3\ell_n^{-3}
=3L^3B_HK_b\tilde{\alpha}_{0,k}\,\ell_n^{-2}.
\end{align*}
Here we used
\begin{equation*}
2C_Jw\le 2C_JB_HK_b\tilde{\alpha}_{0,k}\ell_n\le 2C_JB_HK_b\tilde{\alpha}_{0,k}\le1,
\end{equation*}
the last inequality being the condition $2C_JB_HK_b\tilde{\alpha}_{0,k}\le1$ of
\eqref{4.12:eq-coupling-thresholds-d}.

By \eqref{4.12:eq-below-slot-level-3}, this is the current bound in (b). It is the current part
of clause (b) of Proposition~\ref{4.12:prop-reblocking-inside} in the case $n<p$. This
completes the proof of (b).

(viii) Comparison with the radii. For $n<p$ the factor $L^{(n-p)_+}$ in the consumer's radii
$\operatorname{tgt}_2^E$ and $\operatorname{tgt}_4^E$ of Definition~\ref{4.12:def-consumer-radii}
equals one. Dividing the bounds of steps (vi) and (vii) by these radii therefore gives
\begin{align*}
\frac{|\partial U_E(\mathfrak{p})-1|}{\operatorname{tgt}_2^E(\mathfrak{p})}
&\le 3L^2B_HK_b\,\frac{\tilde{\alpha}_{0,k}}{\tilde{\alpha}_{0,n}}\,
\frac{\ell_n}{\vartheta_E(\mathfrak{p})},\\
\frac{|J(U_E)(\mathfrak{b})|}{\operatorname{tgt}_4^E(\mathfrak{b})}
&\le 3L^3B_HK_b\,\frac{\tilde{\alpha}_{0,k}}{\tilde{\alpha}_{0,n}}\,
\frac{\ell_n}{\vartheta_E(\mathfrak{b})}.
\end{align*}
Insert \eqref{4.12:eq-running-radii-a} for the ratio
$\tilde{\alpha}_{0,k}/\tilde{\alpha}_{0,n}$, and write $\ell_n=L^{-(k-n)}$. The result is
\eqref{4.12:eq-reblocking-inside-d}.

Since $n<p\le k$, we have $n\le k-1$, that is, $k-n\ge1$. We now combine three inputs:
\begin{itemize}
\item the lower bound \eqref{4.12:eq-consumer-radii-d} on $\vartheta_E$;
\item the arithmetic bound on the gain factors (Lemma~\ref{4.12:lem-gain-arithmetic}), applied
with $m=k-n\ge1$;
\item the inequality $L^2\le L^3$ for the plaquette quotient.
\end{itemize}
With $W(N_0)$ as in \eqref{4.12:eq-coupling-thresholds-h}, they show that both quotients are at
most
\begin{equation*}
3L^3B_HK_b(1+\beta_0)\,W(N_0)/\theta_*\le\tfrac12 .
\end{equation*}
The final inequality is \eqref{4.12:eq-coupling-thresholds-f}, namely
$6L^3B_HK_b(1+\beta_0)W(N_0)\le\theta_*$. This proves (c).
\end{proof}

\begin{lemma}[Conclusion and uniformity of the constants]\label{4.12:prf-uniform-constants}
Work in the setting of Proposition~\ref{4.12:prop-reblocking-inside}, stated above with proof
incomplete at the step marked $(\ast)$; every reference below to its clauses is subject to this
proviso. Let $E$ be a consumer with $X'\subset\mathfrak{Q}\setminus\mathcal{C}$, let $p$ be a slot
level of $E$, and let $X''\subset X'$ be a sub-domain read by $E$. Then:
\begin{itemize}
\item[(1)] the bounds of clause (b) of Proposition~\ref{4.12:prop-reblocking-inside} hold for every
plaquette $\mathfrak{p}\subset X''^{\sim 2}$ and every bond $\mathfrak{b}\subset X''^{\sim 2}$;
\item[(2)] in every case the ratio of the bounded quantity to the corresponding radius of the
consumer $E$ is at most $\tfrac12<1$, which is clause (c) of that proposition;
\item[(3)] the constants entering these bounds, namely $B_3$, $B_H$, $c_Q$, $c_d$, $C_J$ and
$\delta$, which are fixed by the inputs of Proposition~\ref{4.12:prop-reblocking-inside}, the bound
$D_{\mathfrak{Q}}\le 204M$ on the side of the admissible box, and $L$, do not depend on $k$, on the
consumer $E$, on $p$ or on $X''$.
\end{itemize}
The dependence on $k$ is confined to $N_0=N_0(k)$ and $\tilde{\alpha}_{0,k}$. These are controlled
uniformly by the coupling thresholds $(\Gamma_0)$--$(\Gamma_3)$ of
Definition~\ref{4.12:def-coupling-thresholds}, through \eqref{4.12:eq-running-radii-b} and the
bound on the running radii $\tilde{\alpha}_{0,k}$.
\end{lemma}

\begin{proof}
Let $\mathfrak{p}\subset X''^{\sim 2}$ be a plaquette or $\mathfrak{b}\subset X''^{\sim 2}$ a bond.
Put $n:=n_E(\mathfrak{p})$, respectively $n:=n_E(\mathfrak{b})$. Either $n\ge p$ or $n<p$, so these
two cases exhaust the plaquettes and bonds of $X''^{\sim 2}$.

Suppose first that $n\ge p$. Lemma~\ref{4.12:prf-above-slot-level} then gives the bounds of
clause (b) at $\mathfrak{p}$, respectively at $\mathfrak{b}$. It also shows that the ratio of each
bounded quantity to the corresponding radius of $E$ is at most $\tfrac12$.

Suppose next that $n<p$. The direct expansion of Lemma~\ref{4.12:prf-below-slot-level} gives the
bounds of clause (b) at $\mathfrak{p}$, respectively at $\mathfrak{b}$. It also shows that the same
ratio is at most $\tfrac12$.

Together the two cases give (1). They also give (2), since $\tfrac12<1$ in each case.

For (3), the constants entering the bounds of the two lemmas are $B_3$, $B_H$, $c_Q$, $c_d$, $C_J$
and $\delta$, taken from the inputs, together with the bound $D_{\mathfrak{Q}}\le 204M$ and the
blocking factor $L$. None of them depends on $k$, on the consumer $E$, on the slot level $p$ or on
the sub-domain $X''$. The level $k$ enters the bounds only through $N_0=N_0(k)$ and
$\tilde{\alpha}_{0,k}$. Both quantities are controlled uniformly by the conditions
$(\Gamma_0)$--$(\Gamma_3)$ of Definition~\ref{4.12:def-coupling-thresholds}, through
\eqref{4.12:eq-running-radii-b} and the bound on the running radii. This concludes the argument for
clauses (b) and (c). With these clauses established, the proof of
Proposition~\ref{4.12:prop-reblocking-inside} is complete, except at the step marked $(\ast)$.
\end{proof}

\begin{corollary}[Real-factor slot bounds inside admissible boxes]\label{4.12:cor-real-factor-slots}
Assume the following two statements:
\begin{itemize}
\item[(a)] the conclusion of Proposition~\ref{4.12:prop-reblocking-inside} (the re-blocking step
inside admissible boxes; stated; proof incomplete at the step marked $(\ast)$), together with the
conditions under which it is stated;
\item[(b)] the statement that, off the core $\mathcal{C}$, the $G$-valued factor of every composite
$\mathcal{W}_i$, $i=1,\dots,4$, of the theorem that fixes these four composites and their consumers
is the $G$-valued factor $\bar{U}$ of the source datum $(\mathbb{U},\mathbb{J})$ of
Definition~\ref{4.12:not-admissible-boxes}.
\end{itemize}
Let $E$ be a consumer of that theorem (Definition~\ref{4.12:def-consumer-radii}) whose domain satisfies
$X'\subset\mathfrak{Q}\setminus\mathcal{C}$ for some admissible box $\mathfrak{Q}$
(Definition~\ref{4.12:not-admissible-boxes}). By \eqref{4.12:eq-admissible-boxes-b} this includes every
consumer with $X'\subset Z''_k\setminus\mathcal{C}$, whether or not $X'$ meets the cut band
$\mathfrak{K}_c$.
Then, for each $i=1,\dots,4$, the half of clause (iv) of the consumer's space $D_E(X')$ that concerns
the $G$-valued factor is satisfied by the $G$-valued factor
\begin{equation*}
\bar{U}_i:=\bar{U}
\end{equation*}
of the composite $\mathcal{W}_i$. Consumers whose domain meets $\mathcal{C}$ are not covered; there the
$G$-valued factor of the composites is $U_2$.
\end{corollary}

\begin{proof}
Fix $i\in\{1,\dots,4\}$. Since $X'\subset\mathfrak{Q}\setminus\mathcal{C}$, the domain $X'$ lies
off the core $\mathcal{C}$. Hypothesis (b) therefore applies on $X'$, and the $G$-valued factor of
$\mathcal{W}_i$ on $X'$ is $\bar{U}$.

Since $X'\subset\mathfrak{Q}\setminus\mathcal{C}$, the consumer $E$ and the box $\mathfrak{Q}$ satisfy
the hypotheses of Proposition~\ref{4.12:prop-reblocking-inside} (stated; proof incomplete at the step
marked $(\ast)$).
By hypothesis (a) we may apply that proposition to $E$, to the box $\mathfrak{Q}$ and to the
$G$-valued factor $\bar{U}$ of the source datum $(\mathbb{U},\mathbb{J})$. The conclusion of that
proposition gives, for $\bar{U}$, the bounds required of the $G$-valued factor in clause (iv) of
$D_E(X')$, with the consumer's own radii $\operatorname{tgt}_2^E$, $\operatorname{tgt}_4^E$ for the two
derived slots; that is, $\bar{U}$ satisfies the half of clause (iv) of $D_E(X')$ that concerns the
$G$-valued factor. By the
first paragraph, $\bar{U}$ is the $G$-valued factor $\bar{U}_i$ of $\mathcal{W}_i$ on $X'$. This proves
the assertion for $\mathcal{W}_i$, and hence for each $i=1,\dots,4$.
\end{proof}

\begin{remark}[The minimiser on a refined determining set]\label{4.12:rem-refined-minimiser}
Let $(\mathbb{U},\mathbb{J})$, $\bar{U}$ and $\mathfrak{Q}$ be as in Proposition~\ref{4.12:prop-reblocking-inside},
whose proof is incomplete at one step. Let $E$ be a consumer with
$X'\subset\mathfrak{Q}\setminus\mathcal{C}$, let $p$ be a slot level of $E$, and let $X''\subset X'$ be a
sub-domain read by $E$ (Definition~\ref{4.12:def-consumer-radii}). Since $\nu=k\ge p$ on $X''$, the
determining set attached to $(\mathbb{U},\mathbb{J})$ on $X''$,
\begin{equation*}
\mathfrak{B}_\nu:=\mathfrak{B}_\nu(p,X''),
\end{equation*}
has block level $\min(p,t_{X''})$ (Definition~\ref{4.12:def-determining-sets}). Let
\begin{equation*}
U_\nu:=U\bigl(\mathfrak{B}_\nu,M_{\mathfrak{B}_\nu}(\bar{U})\bigr)
\end{equation*}
be the corresponding constrained minimiser (Definition~\ref{4.12:def-constrained-minimisers}).
Write $\mathfrak{B}:=\mathfrak{B}_E(p,X'')$ for the consumer's determining set and
$U_E:=U_E(p,X'')$ for the consumer's minimiser.

On $X''$ the set $\mathfrak{B}$ refines $\mathfrak{B}_\nu$. Where $c_E<\min(p,t_{X''})$, a block $Q$
of $\mathfrak{B}_\nu$ of level $\min(p,t_{X''}(Q))$ is subdivided into blocks of $\mathfrak{B}$; we say
that $Q$ is \emph{re-blocked}. Elsewhere the two sets coincide.

Put $D_2:=M_{\mathfrak{B}}(U_\nu)$. Write
\begin{equation*}
D_1:=M_{\mathfrak{B}}(\bar{U})=\bigl(e^{iB_{\mathrm{rb}}}D_2\bigr)^{v}
\end{equation*}
with a $G$-valued gauge transformation $v$. The $\mathfrak{g}$-valued datum $B_{\mathrm{rb}}$ on the
bonds of $\mathbb{L}(\mathfrak{B})$ is called the \emph{re-block datum}.

Grant the composition property of the averages under refinement. Then there are gauge
transformations $\tau$ and $\lambda$ such that
\begin{equation}\label{4.12:eq-refined-minimiser-1}
U_E=\Bigl[e^{i\eta\mathcal{H}_{\mathfrak{B}}(B_{\mathrm{rb}};U_\nu^{\tau})}\,U_\nu^{\tau}\Bigr]^{\lambda}.
\end{equation}
Here $\mathcal{H}_{\mathfrak{B}}(b;U_0)$ denotes the response $1$-form of the constrained minimiser
along $\mathfrak{B}$ to a datum $b$ around a base configuration $U_0$.
In words: the minimiser along the consumer's set is the minimiser along the set attached to
$(\mathbb{U},\mathbb{J})$, gauge transformed, and twisted by the response to the re-block datum.

The identity \eqref{4.12:eq-refined-minimiser-1} is not used in the proof of
Proposition~\ref{4.12:prop-reblocking-inside}, for two reasons.

\emph{First, it is unnecessary inside an admissible box.}
\begin{itemize}
\item At points of consumer level $n\le k-1$, the consumer's two radii
(Definition~\ref{4.12:def-consumer-radii}) exceed what the unit-scale smoothness of $\bar{U}$
delivers, by factors of order $L^{k-n}$ and $L^{2(k-n)}$ respectively. Under
$(\Gamma_1)$--$(\Gamma_2)$ these factors absorb the constants of
\eqref{4.12:eq-constrained-minimisers-b} and \eqref{4.12:eq-constrained-minimisers-c}.
\item The only points of consumer level $k$ inside $\mathfrak{Q}$ are those of consumers of class (K4)
on the band $B_5$. There the gain $L_0^{2N_0}$ absorbs the same constants under $(\Gamma_3)$.
\end{itemize}
Thus nowhere inside $\mathfrak{Q}$ is the sharp radius attached to $(\mathbb{U},\mathbb{J})$ needed.

\emph{Second, with the published tools, the route through \eqref{4.12:eq-refined-minimiser-1} does
not close.} It would require, for a level $i<p$ of the blocks of $\mathfrak{B}$ that subdivide a
re-blocked level-$p$ block, the datum bound
\begin{equation*}
|B_{\mathrm{rb}}|\le\ell_i\bigl(\ell_p\tilde{\alpha}_{0,k}+\dotsb\bigr)
\end{equation*}
on every level-$i$ bond of $\mathbb{L}(\mathfrak{B})$ inside that block. Here the dots denote further
terms inside the bracket; what matters below is only that the bound is of first order in $\ell_i$.
The datum bound in
turn requires a gauge $v$ in which the level-$i$ averages of $\bar{U}$ and of $U_\nu$ differ at first
order in $\ell_i$.

Suppose one uses, blockwise, the sup-norm form of the estimates of
\cite[Propositions~6, 7]{B-CMP98}, the estimate \eqref{4.12:eq-determining-sets-e} and the axial gauge
of a box (Lemma~\ref{4.12:lem-axial-gauge}), with one axial gauge per re-blocked block or per pair of
adjacent re-blocked blocks. This gives the following.
\begin{itemize}
\item On bonds interior to a re-blocked block, the discrepancy is of first order in $\ell_i$.
\item On bonds crossing between adjacent re-blocked blocks, the discrepancy is only of order
$\ell_p^2$ times the curvature.
\end{itemize}
The reason for the second item is that the discrepancy on a crossing bond is a flux of the difference
of the two curvatures through a surface of area of order $\ell_p^2$. The equality of the level-$p$ data
constrains this flux only on average.

A datum of order $\tilde{\alpha}_{0,k}\ell_p^2$ at a shell point of level $i$ gives a ratio of
response to target of order
\begin{equation*}
B_Hc_d(1+\beta_0)(1+k-i)^{1/2}\,\ell_p^2/\vartheta_E .
\end{equation*}
For $p$ near $k$ and $i$ low, this ratio is not small under any admissible condition. The factor
$(1+k-i)^{1/2}$ is unbounded along the flow, and no gain or decay accompanies it.

In the global axial gauge of the box $\mathfrak{Q}$, the first-order bound in $\ell_i$ does hold: it
is \eqref{4.12:eq-reblocking-inside-c} of Proposition~\ref{4.12:prop-reblocking-inside}, whose proof
is incomplete at one step. There, however, the comparison with $U_\nu$ is superfluous.
The same gauge makes the data themselves small, once they are read in that gauge, and
\eqref{4.12:eq-constrained-minimisers-c} applies at the trivial background. This is the route taken in
the direct expansion at points below the slot level, \ref{4.12:prf-below-slot-level}.
\end{remark}

\begin{proof}[Proof of the identity \eqref{4.12:eq-refined-minimiser-1}]
The argument has three steps.

\emph{Step one: $U_\nu$ is admissible for the refined constraints.} By the definition of $D_2$,
$M_{\mathfrak{B}}(U_\nu)=D_2$. Hence $U_\nu$ lies on the constraint manifold
$C(\mathfrak{B},D_2)$.

\emph{Step two: $U_\nu$ is critical for them.} The configuration $U_\nu$ is critical for
$\mathcal{A}$ on $C(\mathfrak{B}_\nu,M_{\mathfrak{B}_\nu}(\bar{U}))$. The set $\mathfrak{B}$ refines
$\mathfrak{B}_\nu$ on $X''$ in the way just described. By the composition property of the averages
under refinement, which we grant, the averages of a configuration over the blocks of
$\mathfrak{B}_\nu$ are determined by its averages over the blocks of $\mathfrak{B}$: on a re-blocked
block through the averages over the blocks of $\mathfrak{B}$ that subdivide it, and trivially where
the two sets coincide. Hence every $U\in C(\mathfrak{B},D_2)$, which has the same averages along
$\mathfrak{B}$ as $U_\nu$, satisfies
\begin{equation*}
M_{\mathfrak{B}_\nu}(U)=M_{\mathfrak{B}_\nu}(U_\nu)=M_{\mathfrak{B}_\nu}(\bar{U}),
\end{equation*}
the second equality because $U_\nu$ lies on $C(\mathfrak{B}_\nu,M_{\mathfrak{B}_\nu}(\bar{U}))$.
Thus $C(\mathfrak{B},D_2)\subset C(\mathfrak{B}_\nu,M_{\mathfrak{B}_\nu}(\bar{U}))$, and by Step one
$U_\nu$ lies in the smaller set. Every tangent direction to $C(\mathfrak{B},D_2)$ at $U_\nu$ is
tangent to the larger manifold, along which the first variation of $\mathcal{A}$ at $U_\nu$ vanishes;
so $U_\nu$ is critical for $\mathcal{A}$ on $C(\mathfrak{B},D_2)$. Clause~(ii) of
\eqref{4.12:eq-constrained-minimisers-b}, applied to this critical point, gives a gauge
transformation $\tau$ with
\begin{equation*}
U(\mathfrak{B},D_2)=U_\nu^{\tau}.
\end{equation*}

\emph{Step three: the response around the base.} The consumer's minimiser is the constrained minimiser
along $\mathfrak{B}$ with data $D_1=M_{\mathfrak{B}}(\bar{U})$. By the choice of $v$,
\begin{equation*}
D_1=\bigl(e^{iB_{\mathrm{rb}}}D_2\bigr)^{v}.
\end{equation*}
So $D_1$ is, up to the gauge transformation $v$, the datum $D_2$ twisted by $B_{\mathrm{rb}}$. Apply
the response statement for the constrained minimiser around a general base, with base
$U(\mathfrak{B},D_2)=U_\nu^{\tau}$ and datum $B_{\mathrm{rb}}$. It gives a gauge transformation
$\lambda$ such that
\begin{equation*}
U_E=\Bigl[e^{i\eta\mathcal{H}_{\mathfrak{B}}(B_{\mathrm{rb}};U_\nu^{\tau})}\,U_\nu^{\tau}\Bigr]^{\lambda},
\end{equation*}
which is \eqref{4.12:eq-refined-minimiser-1}.
\end{proof}

\begin{remark}[What the re-blocking step does not assert]\label{4.12:rem-reblocking-not-asserted}
Let $E$ be a consumer with domain $X'$, in the sense of Definition~\ref{4.12:def-consumer-radii}.
Let $p$ be a slot level of $E$, and let $X''\subset X'$ be a sub-domain read by $E$.
The \emph{re-block region} of the pair $(p,X'')$ is the set of points $x\in X''$ with $c_E(x)<p$.
The scope of Proposition~\ref{4.12:prop-reblocking-inside} (stated; proof incomplete at the step
marked $(\ast)$) and of this section is delimited as follows.
\begin{enumerate}
\item[(1)] Nothing is asserted for consumers whose domain $X'$ meets the core $\mathcal{C}$. On the
core the group-valued factor of the composites $\mathcal{W}_i$ is $U_2$, by the construction of
the composites.

Nothing is asserted either for consumers whose domain is not contained in any admissible box
in the sense of Notation~\ref{4.12:not-admissible-boxes}. Two instances are:
\begin{itemize}
\item the consumers of class (K4) whose domain meets both $Z''_k$ and $\Omega^{\mathrm{old}}_k$;
\item the old terms $B^{(j)}$, $R^{(j)}$, $E^{(j)}$ with $k_0<j<k$ that are localised on the
cleaned collars or on the band $B_5$ and read slots $p>o(x)$ there.
\end{itemize}
For such consumers there are pairs $(p,X'')$ whose re-block region is non-empty, and for these
no box gauge in the sense of Lemma~\ref{4.12:lem-axial-gauge} is available.
\item[(2)] Clause~(iv) of the definition of the analyticity spaces of this section, among them the
source space $\tilde{U}^c_k(Y_4,\tilde{\alpha}_0,\tilde{\alpha}_1)$, is the clause that controls
the minimisers along determining sets; it is not clause~(iv) of Theorem~0. Nothing is asserted
about its complex half at the base of the induction over levels by which that clause is
established for the spaces of the consumers. That half consists of the complex
minimisers $U(\mathfrak{B}_E(p,X''),M_{\mathfrak{B}_E(p,X'')}(\mathcal{W}_i))$
along the determining sets of the
consumer. The term of order $\tilde{\alpha}_{1,k}$ that arises in the comparison with the
source's minimiser, that is, the constrained minimiser with datum $\bar{U}$ along the
determining sets of the source, belongs to that half. For the group-valued half, on the
other hand, the bounds of Proposition~\ref{4.12:prop-reblocking-inside} (stated; proof incomplete
at the step marked $(\ast)$) are consistent with the radii of the consumer in
Definition~\ref{4.12:def-consumer-radii}, in particular with \eqref{4.12:eq-consumer-radii-c},
at every slot level $p$.
\item[(3)] No holomorphy statement is made, and none is needed.
\item[(4)] Clause~(iv) of the source space, for $\bar{U}$ along the determining sets of the source
datum, is not used. Of the conditions on the source datum, only
\eqref{4.12:eq-admissible-boxes-a} is used.
\item[(5)] No numerical value is asserted for any threshold. The thresholds $N_*$, $M_*$, $\gamma_*$
are existential.
\end{enumerate}
\end{remark}

\begin{proof}
We justify items~(3) and~(4) and the consistency assertion of item~(2), and then describe where
the argument stops for the consumers excluded in item~(1).

The group-valued half of clause~(iv) concerns the $G$-valued factor at each point
$(\mathbb{U},\mathbb{J})$ of the source space $\tilde{U}^c_k(Y_4,\tilde{\alpha}_0,\tilde{\alpha}_1)$.
It is therefore a statement about real constrained minimisers, namely the minimisers
$U_E(p,X'')$ of the Wilson action on the constraint manifold. This is item~(3).

Of the clauses of the source space bearing on $\bar{U}$, only the bound
\eqref{4.12:eq-admissible-boxes-a} on admissible boxes enters
Proposition~\ref{4.12:prop-reblocking-inside} (stated; proof incomplete at the step marked
$(\ast)$). Hence, inside an admissible box, that proposition does not depend on reproducing the
group-valued half of clause~(iv) for the source space. This is item~(4).

For the group-valued half, the bounds of Proposition~\ref{4.12:prop-reblocking-inside} (stated;
proof incomplete at the step marked $(\ast)$) are consistent with the radii of the consumer. Fix a
plaquette $\mathfrak{p}$ or a bond $\mathfrak{b}$ in $X''^{\sim 2}$, and put $n=n_E(\mathfrak{p})$,
resp.\ $n=n_E(\mathfrak{b})$. We treat the slot levels $p\le n$ and $p>n$ separately.
\begin{itemize}
\item Let $p\le n$. The bound \eqref{4.12:eq-reblocking-inside-a} is taken at the coarsest
scale. The radius of the slot $E_4$ in \eqref{4.12:eq-consumer-radii-c} carries the weakening
factor $L^{(n-p)_+}$. The bound taken at the coarsest scale, together with this weakening factor,
gives the consistency for the slot $E_4$.
\item Let $p>n$. Then the block scale is the scale of the consumer.
\end{itemize}

We now turn to item~(1). Let $E$ be a consumer with $X'\cap\mathcal{C}=\emptyset$ and $X'$
not contained in any admissible box. Let $(p,X'')$ be a pair whose re-block region
$\{x\in X'' : c_E(x)<p\}$ is non-empty. For such a pair the argument of
Proposition~\ref{4.12:prop-reblocking-inside} (stated; proof incomplete at the step marked
$(\ast)$) stops at the step at which Lemma~\ref{4.12:lem-axial-gauge} is invoked to supply a box
gauge.

The lemma needs a box that contains $X''$
and on which \eqref{4.12:eq-admissible-boxes-a} holds with radius $\tilde{\alpha}_{0,k}\eta^2$.
The constant $K_b$ of the datum grows linearly with the side of the box. Since no bound on the
size of $X''$ is available, a box large enough to contain it need not exist inside the healed
region. A bound on the side
of the box would be a bound on $\operatorname{diam}X'$, and such a bound is not admitted.

The alternative route compares $U_E(p,X'')$ with the source's minimiser, that is, the
constrained minimiser with datum $\bar{U}$ along the determining sets of the source. This route
stops at the datum. It would require a bound on the re-blocking term
of that comparison, on a block of a level-$i$ shell, by $\ell_i$ times the sum of
$\ell_p\tilde{\alpha}_{0,k}$ and further terms.

Consider the bonds of $\mathbb{L}(\mathfrak{B}_E(p,X''))$ that cross between adjacent subdivided
blocks of the determining sets of the source datum. On these bonds the required bound does not
follow from two inputs taken together:
\cite[Propositions~6, 7]{B-CMP98}, in their sup-norm form, and the axial gauge of
Lemma~\ref{4.12:lem-axial-gauge}. What these two inputs give on such bonds is a bound of order
$\ell_p^2\tilde{\alpha}_{0,k}$. For $p$ near $k$ a bound of this order is weaker than the
required one on the shells of low level, so the comparison cannot be completed.

Either of the two lemmas stated below would remove this obstruction.
\end{proof}

\begin{proposition}[The re-blocking step outside admissible boxes. Proof outstanding]\label{4.12:prop-reblocking-outside}
Let the standing hypotheses and the data $(\mathbb{U},\mathbb{J})$ and $\bar{U}$ be those of
Proposition~\ref{4.12:prop-reblocking-inside}, whose proof is incomplete at one step; the hypothesis
there that the domain of the consumer lie in an admissible box is not assumed. Let $E$ be a
consumer in the sense of Definition~\ref{4.12:def-consumer-radii}, with
domain $X'$ and level function $c_E$. Assume the following:
\begin{itemize}
\item $X'\cap\mathcal{C}=\emptyset$;
\item $X'$ is not contained in any box of side at most $D_{\mathfrak{Q}}=204M$ in the healed region.
\end{itemize}
Let $p$ be a slot level of $E$, and let $X''\subset X'$ be a sub-domain read by $E$ such that
\begin{equation*}
\{x\in X'' : c_E(x)<p\}\neq\emptyset .
\end{equation*}
Then the conclusions (a) and (b) of Proposition~\ref{4.12:prop-reblocking-inside} hold for $E$ at
the pair $(p,X'')$.
\end{proposition}

Examples of such consumers and pairs are given in
Remark~\ref{4.12:rem-reblocking-not-asserted}.

Proof outstanding.

\emph{Route.} A proof would carry the argument for Proposition~\ref{4.12:prop-reblocking-inside},
whose proof is itself incomplete at one step, to the consumers and pairs of the statement, with the
remaining bookkeeping of that argument unchanged. Either of the following two statements would
suffice.

\begin{enumerate}
\item[(i)] \emph{A re-block datum lemma.} For a level $m$, write $\mathcal{M}_m$ for
Ba{\l}aban's block-averaging map onto level $m$ (Definition~\ref{1.5:def-block-average}). Let
$\square$ be a level-$p$ block. Let $U_1,U_2$ be $G$-valued
configurations on $\square$ and its neighbouring blocks, and let $F_1,F_2\ge0$. Suppose that
$|\partial U_r(\mathfrak{p})-1|\le F_r\eta^2$ for $r=1,2$ and every plaquette $\mathfrak{p}$ there.
Suppose further that $\mathcal{M}_p(U_1)=\mathcal{M}_p(U_2)$ on the bonds of $\square$. The
statement required is the following. There is a constant $C_{\mathrm{rb}}$, depending on $d$ only,
such that for every level $j$ there exist a $G$-valued gauge transformation $\sigma$ on the
level-$j$ base points in $\square$ and its neighbouring blocks and a $\mathfrak{g}$-valued function
$B$ on the level-$j$ bonds in $\square$ and its neighbouring blocks such that
\begin{equation}\label{4.12:eq-reblocking-outside-1}
\mathcal{M}_j(U_1)=\bigl(e^{iB}\mathcal{M}_j(U_2)\bigr)^{\sigma},
\end{equation}
and such that, for every level-$j$ bond $\mathbf{b}$ in $\square$, the bonds crossing
$\partial\square$ included,
\begin{equation}\label{4.12:eq-reblocking-outside-2}
|B(\mathbf{b})|\le C_{\mathrm{rb}}\,\ell_j\,\ell_p\,(F_1+F_2).
\end{equation}
It is a discrete Coulomb-gauge estimate under the averaged constraint, with inputs
\cite[Propositions~6--7]{B-CMP98} and \cite[(1.7)--(1.19)]{B-CMP99a}.

\item[(ii)] \emph{Tower localisation at real data.} Let $\mathfrak{Q}$ be an admissible box, and
put $X''':=X''\cap\mathfrak{Q}$. Let $\Delta\ge0$, and let $x\in X''$ satisfy
$\operatorname{dist}(x,\partial\mathfrak{Q})\ge\Delta$. Write $U(\mathfrak{B},D)$ for the
constrained minimiser with determining set $\mathfrak{B}$ and datum $D$, and take as datum
$\mathcal{M}\cdot\bar{U}$, the averages of $\bar{U}$
(Proposition~\ref{4.12:prop-reblocking-inside}) over the blocks of the determining set concerned.
The statement required is that,
for every plaquette $\mathfrak{p}$ at $x$,
\begin{equation}\label{4.12:eq-reblocking-outside-3}
\begin{split}
\bigl|\partial U(\mathfrak{B}_E(p,X''),\mathcal{M}\cdot\bar{U})(\mathfrak{p})-1\bigr|
&\le\bigl|\partial U(\mathfrak{B}_E(p,X'''),\mathcal{M}\cdot\bar{U})(\mathfrak{p})-1\bigr|\\
&\quad+C_{\mathrm{loc}}\,B_3\,\tilde{\alpha}_{0,k}\,e^{-\delta\Delta/M_1}\,
\eta^2\,\ell_{\mathfrak{B}}(\mathfrak{p})^{-2},
\end{split}
\end{equation}
with some constant $C_{\mathrm{loc}}$. For every bond $\mathfrak{b}$ at $x$ the analogous
inequality is required, with $|J(U(\cdot))(\mathfrak{b})|$ in place of
$|\partial U(\cdot)(\mathfrak{p})-1|$ on both sides and $\ell_{\mathfrak{B}}(\mathfrak{b})^{-3}$ in
place of $\ell_{\mathfrak{B}}(\mathfrak{p})^{-2}$. This route requires in particular that the coarse
constraints across $\partial X'''\setminus\partial X''$ be compatible with the finer constraints of
the determining set $\mathfrak{B}_E(p,X''')$; this compatibility is part of what would have to be
proved and is not proved here. The inputs are Ba{\l}aban's tower analysis in
\cite{B-CMP119} and \cite[(1.8)--(1.10)]{B-CMP116}, a localisation lemma for the cluster
expansions of \cite{B-CMP116}, and a lemma on the boundary towers of the determining sets.
\end{enumerate}

\begin{lemma}[{Uniqueness and locality of the constrained minimizer, completing
Ba{\l}aban's Lemma~4 of \cite{B-CMP109}. Proof outstanding}]\label{4.13:lem-minimizer-unique}
Let the constants satisfy the restrictions of \cite[Lemma~4]{B-CMP109}, and let the
fields be those admitted there.
\begin{enumerate}
\item[(a)] \emph{Uniqueness.} At every field admitted by \cite[Lemma~4]{B-CMP109},
the constrained variational problem of that lemma has exactly one minimizer, so that
the minimizer map of that lemma is well defined.
\item[(b)] \emph{Locality.} The minimizer is local, in the sense furnished by the
averaging operations of \cite{B-CMP98}. The property required of it is that the
nesting identity \cite[(3.38)]{B-CMP109} hold when it is read for an arbitrary
sub-domain of the domain of \cite[Lemma~4]{B-CMP109} in place of that domain.
\end{enumerate}
\end{lemma}

Proof outstanding.

\emph{Route.} Statement (a) is to be derived from \cite[Theorem~1]{B-CMP102b}.
For configurations satisfying the curvature bound of its hypothesis, that theorem
gives a minimal orbit in a first space of configurations. Under its smallness
conditions on the constants, it also gives that this orbit is the unique critical
orbit in a second, explicitly described space of configurations. A proof of (a)
would have to establish three things:
that every field admitted by \cite[Lemma~4]{B-CMP109} satisfies the hypothesis of
\cite[Theorem~1]{B-CMP102b}, with constants meeting that theorem's smallness
conditions; that every minimizer of the constrained variational problem of
\cite[Lemma~4]{B-CMP109} is a critical point lying in the second space, the one in
which \cite[Theorem~1]{B-CMP102b} asserts uniqueness of the critical orbit; and that
uniqueness of the orbit yields uniqueness of the minimizer in the form in which
\cite[Lemma~4]{B-CMP109} uses it.
Statement (b) is to be derived from the averaging operations of \cite{B-CMP98}.
A proof would have to extract from \cite{B-CMP98} a locality property of the
minimizer and show that this property, together with (a), allows the nesting identity
\cite[(3.38)]{B-CMP109} to be read with an arbitrary sub-domain
$Z\in\mathcal{D}_j$ (Definition~\ref{1.5:def-domains}) of the domain of
\cite[Lemma~4]{B-CMP109} in place of that domain.

The lemma is an input of the statement on the one-power hereditary class. It
enters that statement through the nesting identity \cite[(3.38)]{B-CMP109}, read
for an arbitrary sub-domain $Z$; that reading uses both parts of the lemma.

\begin{lemma}[The gauge condition at the small pairs' and real-background pairs' own curvature
constant, reproduced by one step. Proof outstanding]\label{4.13:lem-own-constant}
Let $B$ be the constant of the gauge condition \cite[(1.12)]{B-CMP109}. Let $c_\square$ be the
cube-size multiple fixed in this chapter's definition of the widened reading of that condition,
the reading of \cite[(1.12)]{B-CMP109} on every cube of side at most a fixed multiple of $LM$.

Consider two kinds of pairs of Ba{\l}aban's one-step analysis in \cite[\S\S3--5]{B-CMP109}
and \cite[\S\S1--2]{B-CMP116}:
\begin{itemize}
\item the small pairs, and
\item the pairs at real backgrounds,
\end{itemize}
in each case those whose membership in Ba{\l}aban's class, the class of pairs on which that
analysis is carried out, is printed in \cite[\S\S3--5]{B-CMP109} and \cite[\S\S1--2]{B-CMP116}.
For each such pair, let $\alpha_0'$ denote the curvature constant of its printed membership.

Then every such pair satisfies the gauge condition \cite[(1.12)]{B-CMP109} in the widened reading at its own curvature constant. That is, on every cube of side at most $c_\square LM$, condition \cite[(1.12)]{B-CMP109} holds with the single bound
\begin{equation*}
c_\square\, L\, M\, B\, \alpha_0' .
\end{equation*}
Moreover, this property is reproduced by one step of the renormalization of \cite[\S\S3--5]{B-CMP109} and \cite[\S\S1--2]{B-CMP116}.
\end{lemma}

Proof outstanding.

\paragraph{Route.}
The statement is to be derived together with the thin-cube gradient bound of this chapter.
A proof would have to establish two things: for the small pairs and the pairs at real
backgrounds whose membership in Ba{\l}aban's class is printed in the cited sections, the
gauge condition \cite[(1.12)]{B-CMP109} in the widened reading at their own printed curvature
constant $\alpha_0'$; and the reproduction of this property under the one-step argument of
\cite[\S\S3--5]{B-CMP109} and \cite[\S\S1--2]{B-CMP116}.

\begin{definition}[Levels, territories, covers and the form at free current]
\label{4.14:not-free-current-notation}
The following notation is fixed for Ba{\l}aban's step $k$. Objects lettered here and used in
other sections carry the meaning given here.

\emph{Levels and units.} We write $t:=k+1$. Lengths are measured in level-$k$ units, in which
the lattice spacing of level $j$ is $\ell_j:=L^{j-k}$; thus $\ell_k=1$ and $\ell_t=L$. There are
three block sizes $M_1$, $M_1^{\sharp}$ and $M$, each a power of $L$, with
$M_1\mid M_1^{\sharp}\mid M$. We also write
$\ell:=M/(12d)$, the real side of the unit boxes; here, and everywhere below, the letter $d$
standing alone denotes the dimension, $d=4$.

\emph{Nest, layers and territories.} The regions of the step form the nest
$\{\Omega_j\}_{j\le k}$, together with the new region $\Omega_t$. The layers are
\begin{equation*}
\Lambda_j:=\Omega_j^{(j)}\setminus\Omega_{j+1}^{(j)}\quad(j<k),\qquad
\Lambda_k:=\Omega_k^{(k)},
\end{equation*}
as in \cite{B-CMP96}, where $\Omega_j^{(j)}$ and $\Omega_{j+1}^{(j)}$ are the sets so written.
The small-field region of \cite{B-CMP119} is written
$\Lambda_j^{119}$, and its complement, the large-field region, is
$Z_j:=(\Lambda_j^{119})^{c}$, as in \cite[(2.3)]{B-CMP119}; the letters $\Lambda_j$ and
$\Lambda_j^{119}$ denote different sets. With $T$ the lattice on which the step is taken, the
territories are
\begin{equation*}
\begin{gathered}
X_n:=\operatorname{closure}(\Omega_n\setminus\Omega_{n+1})\quad(n<k),\qquad X_k:=\Omega_k,\\
X_0:=\operatorname{closure}(T\setminus\Omega_1).
\end{gathered}
\end{equation*}
In the bookkeeping of the second cover below and of the hull construction, and there only, the
index $k+1$ is read with $\Omega_{k+1}=\emptyset$, as in \cite{B-CMP122a}; this reading does not
concern the new region $\Omega_t$ of the bond sets below.

\emph{Distances.} We use the block distance $d(\cdot,\cdot)$ of \cite{B-CMP96}, always written
with its arguments, the weighted distance $\tau$ of \cite{B-CMP96} through listed sets, and its
variant $\tau^{pr}$ used in proofs, with $\tau\le\tau^{pr}$.

\emph{Frames.} We use the unit frame $\mathfrak{F}^{u}$ and the many-scale frame
$\mathfrak{F}^{ms}$, with their block sets $\mathfrak{B}_s$ and $\mathfrak{B}$. In them each
coarse bond $c$ carries an anchor $a(c)$; the anchor distance of two coarse bonds $c,c'$ is
$d(a(c),a(c'))$, and the frame constant is $c_\tau:=2L$.

\emph{Domains.} All domains are regular closed unions of closed cubes of the stated
partitions. The regularisation is applied without further mention; the distances and
inclusions stated below are insensitive to it.

\emph{Bonds.} $\mathbb{B}^{+}$ is the set of free unit bonds of $\Omega_t$ from which the
pivot bonds are removed (the pivot bond $b_0(c)$ of a coarse bond $c$ is its middle bond);
$\mathbb{B}^{+}$ is the set written $\Omega_t^{(k)\ast}$ in
\cite{B-CMP119}. The set $\mathbb{B}\subset\mathbb{B}^{+}$ is arbitrary, and
$\mathbb{B}':=\mathbb{B}^{+}\setminus\mathbb{B}$. Further, $\Lambda_t$ is the region at level
$t$ fixed in the statements that use the following bond sets. The set $\mathbb{B}^{K1}$ is the
set of
bonds with both end-points in $\Lambda_t$; $\mathbb{B}^{\ast}$ is the set of
bonds with at least
one end-point in $\Lambda_t$, and $\mathfrak{x}:=\mathbb{B}^{\ast}\setminus\mathbb{B}^{K1}$.

\emph{The form at free current.} The form
\begin{equation*}
\mathsf{Q}:=C_{\flat}^{\top}\,\Delta^{(k)}\,C_{\flat}\qquad\text{on }\mathfrak{g}^{\mathbb{B}^{+}}
\end{equation*}
is taken at free current $J$, with $\Delta^{(k)}$ the operator of the step so written in
\cite{B-CMP109}, and is given by four defining clauses. These are obtained by the
extension of \cite{B-CMP109}, in which the function $J_{k+1}$ is replaced in all propagators
and operators by the variable $J$, and the configuration $U_{k+1}$ in all other places by the
variable $U$. Its
blocks are
\begin{equation*}
\mathsf{A}:=\mathsf{Q}_{\mathbb{B}\mathbb{B}},\qquad \mathsf{C}:=\mathsf{A}^{-1},\qquad
\mathsf{J}:=\mathsf{Q}_{\mathbb{B}\mathbb{B}'},
\end{equation*}
and $\mathsf{G}:=\|\mathfrak{g}(\cdot)\|^{2}$ is one of the terms of the form; here
$\mathfrak{g}(\cdot)$ is the unit-lattice quantity of \cite{B-CMP109} so written, distinct from
the Lie algebra $\mathfrak{g}$ in $\mathfrak{g}^{\mathbb{B}^{+}}$. The graph of
the current is $\mathfrak{G}:=\{J=J(\mathbf{U})\}$. On the complex fibre of pairs $(A',J)$ we
use the scale-weighted norm $\|(A',J)\|_{\star}$, and $\mathcal{U}(r)$ denotes the complex ball
of radius $r$ in this fibre. The cores are the sets $\mathsf{Z}_0\subset\mathfrak{c}^{\flat}$;
a core is a set of level-$k$ $M$-cubes.

\emph{Constants.} Let $\gamma_0>0$ be the lower bound of the base form, and let $\mu_1$ be a
constant depending on $d$ and $L$ only. Let
$B_{\mathsf{Q}}$ and $\delta_{\mathsf{Q}}$ be the constant and the decay rate of the walk
expansion of the form $\mathsf{Q}$, as in Definition~\ref{4.4:def-form-walk-hypothesis}; let
$e$ be the base of the natural logarithm, and
let $S_u(\cdot)$ be the function of the decay rate attached to the unit frame
$\mathfrak{F}^{u}$; it is not the sheet format $\mathcal{S}_{u}$ of that frame. Set
\begin{gather*}
\gamma_1:=e\,B_{\mathsf{Q}}\,S_u(\delta_{\mathsf{Q}}),\qquad \gamma_{+}:=2/\gamma_0,\qquad
\gamma_{-}:=1/\gamma_1,\qquad \varepsilon_1:=\gamma_0/(8\gamma_1),\\
m_1:=\gamma_0/4,\qquad \delta_0:=\mu_1/4,\qquad B_0:=12\cdot 2^{d}/m_1,\qquad
B_{\mathsf{J}}:=B_{\mathsf{Q}}.
\end{gather*}

\emph{The second cover.} The second cover is
$\mathcal{D}^{\sharp}=\bigcup_j\mathcal{D}^{\sharp}_j$. A box
$\square\in\mathcal{D}^{\sharp}_j$ has side $2s_j$, where
\begin{equation*}
s_j:=M_1^{\sharp}\,\ell_j,
\end{equation*}
and the boxes of $\mathcal{D}^{\sharp}_j$ are centred at the points of the $s_j$-lattice in
$X_j$. For $n\ge 0$, $\square^{n}$ is the box concentric with $\square$ of half-side
$(1+n)s_j$; in particular $\square^{4}$ has side $10s_j$. The hull of $\square$ is
$\widehat{\square}:=\square^{3}$.

The cover carries profiles $g_\square$. Set $h_\square:=g_\square\big(\sum g^{2}\big)^{-1/2}$,
where $\sum g^{2}$ is the sum of the squared profiles, and $h_\square(c):=h_\square(x_c)$,
where $x_c$ is the point attached to the coarse bond $c$. The
index set
of a box is $\mathfrak{X}(\square''):=\{c:\ x_c\in\square^{1}\}$ (the argument is written
$\square''$ as in the statements that use these index sets; the set depends on $\square$ alone).

\emph{Modified objects.} The statements named in $(\alpha)$--$(\delta)$ below are used with the
following four objects.

\smallskip
\noindent(\emph{$\alpha$}) \emph{Reading hulls.} The following sets and radii are fixed:
\begin{equation}\label{4.14:eq-free-current-notation-a}
\begin{gathered}
\square^{\mathsf{r}}:=\square^{2L+1},\qquad
\widehat{X}:=\bigcup_{\square\subset X}\square^{\max\{5,2L\}},\\
\rho_0:=c_X(d)\,L^{2}M_1,\qquad \rho_T=\rho_{\mathsf{Q}}:=8d\,L^{2}M_1^{\sharp}.
\end{gathered}
\end{equation}
Here $c_X(d)$ is a constant depending on $d$ only. In $\square^{\mathsf{r}}$, $\square$ is a box
of the second cover;
in $\widehat{X}$, $X$ is a domain at level $j$ and the union runs over the boxes $\square$ of
the first cover $\mathcal{D}_j$, described below, contained in $X$.
In every domain list of the terms of the inversion lemma attached to a single box of the second
cover, and of the steps of that lemma, the hull $\widehat{\square}=\square^{3}$ is replaced, as
a listed set only, by
$\square^{\mathsf{r}}$, which has half-side $(2L+2)s_j$. The per-box localisation
$T^{\square}$, the projection $1''$, the profiles $h_\square$, the index sets
$\mathfrak{X}(\square'')$ and the four box conditions attached to the boxes of the second cover
are unchanged.

In the first cover, that is the cover $\mathcal{D}_j$ of \cite{B-CMP99b} at block size
$M_1$, the domains of the elementary operators of the walk expansion are the enlarged sets
$\widehat{X}$ of
\eqref{4.14:eq-free-current-notation-a}. For a box $\square$ of $\mathcal{D}_j$, $\square^{n}$
denotes the concentric enlargement of $\square$ used for the first cover; each box
$\square^{\max\{5,2L\}}$ entering
$\widehat{X}$ has half-side $(2L+1)M_1\ell_j$.
By \cite[(1.118)]{B-CMP95}, the functions $h_\square$ attached to the boxes
$\square\in\mathcal{D}_j$ of the first cover satisfy
$\operatorname{supp}h_\square\subset\square$; these are not the profiles of the second cover
defined above. With this inclusion, $\widehat{X}$ is a reading
set, in the sense of the reading property of the operator and sheet formats, for the elementary
operators of the walk expansion. Under the weaker inclusion
$\operatorname{supp}h_\square\subset\square^{1}$ the same holds for the set defined with
$\square^{\max\{5,2L+1\}}$ in place of $\square^{\max\{5,2L\}}$, with only the constant
$c_X(d)$ changed.

\smallskip
\noindent(\emph{$\beta$}) \emph{Symmetric cell weights.} In the localisation of the local terms
and in the branching lemma for the local terms, the terms are $N_z:=w_z\,N$. Here
$N$ is the local term of those two statements (in this item the letter $N$ has this meaning
only, and not that of the rank in $SU(N)$ or of $N_k$), $\Delta(z)$ is the block indexed by
$z$, and for a bond $\iota$
\begin{equation}\label{4.14:eq-free-current-notation-b}
w_z(\iota):=\frac{\#\{\text{end-points of }\iota\text{ lying in }\Delta(z)\}}
{\#\{\text{end-points of }\iota\text{ lying in }\Omega_0\}}.
\end{equation}
Thus $w_z(\iota)=\tfrac12$ on each of the two end-point blocks of an interior bond, that is, of
a bond with both end-points in $\Omega_0$. For a bond
with one end-point outside $\Omega_0$, as in \cite{B-CMP96}, $w_z(\iota)=1$ on the one block
containing its other end-point. Hence, for every bond $\iota$,
\begin{equation*}
\sum_z w_z(\iota)=1,\qquad 0\le w_z\le 1.
\end{equation*}
Moreover $w_{rz}(r\iota)=w_z(\iota)$ for every $r$ in the group $\mathfrak{r}$ of lattice
symmetries.

The domain list of $z$ consists of the single set $X_N^{+}(z)$. Here $X_N(z)$ is the set
attached to $z$ in the two statements just named, $c_N$ is the constant of those statements so
written, and $X_N^{+}(z)$ is the set $X_N(z)$ taken at radius $c_N+dL$. The junction
indicators $\chi_y$
of the product lemma are unchanged; they are indicators, unrelated to the cut-off $\chi$
of $(\delta)$.

\smallskip
\noindent(\emph{$\gamma$}) \emph{The half of the symmetrising operator.} In the first defining
clause of the form with free current, $W_\pi[J]$ is the symmetric operator whose quadratic form
is
\begin{equation}\label{4.14:eq-free-current-notation-c}
A\longmapsto\langle A,\,M_J\,\mathfrak{S}\,\mathsf{w}A\rangle
+\tfrac12\,\langle\mathfrak{S}\,\mathsf{w}A,\,M_J\,D\,\mathsf{w}A\rangle .
\end{equation}
Here $\mathfrak{S}$ is the symmetrising operator, $D$ the covariant derivative and $M_J$ the
multiplication vertex at the current $J$, and $\mathsf{w}$ is the operator of that first
defining clause so written, acting on $A$ (the letter $\mathsf{w}$ is distinct from the complex
source $w$ and from the cell weights $w_z$ of $(\beta)$).

\smallskip
\noindent(\emph{$\delta$}) \emph{The cut-off at scale $M/(2L)$.} In the cut-off lemma,
the cut-off $\chi$ attached to an index $\nu$ with support $[\nu]$ satisfies
\begin{equation}\label{4.14:eq-free-current-notation-d}
\chi=1\ \text{on }[\nu],\qquad
\chi=0\ \text{at weighted distance greater than }\frac{M}{2L}\ \text{from }[\nu].
\end{equation}
It is Lipschitz with constant $2L/M$ with respect to the weighted length.
\end{definition}

\begin{definition}[The class of real backgrounds]\label{4.14:def-real-backgrounds}
Use the levels, territories and covers of Notation~\ref{4.14:not-free-current-notation}. There
$\Omega_0\supset\dots\supset\Omega_k$ is the nest, $X_n$ are its territories, $\mathbb{B}^{+}$ is
the set of free unit bonds of the step with the pivots removed, and $\mathcal{D}^{\sharp}$ is the
second cover, at block size $M_1^{\sharp}$. For $\square\in\mathcal{D}^{\sharp}$, $\square^{4}$ is
its fourth concentric hull, and $\ell_n$ is the lattice spacing of level $n$ in level-$k$ units.
Moreover $G$ is the gauge group, and $\eta$ is the lattice spacing in the representation
$U^{u}=e^{i\eta A}$, as in the sentence preceding clause (i) on \cite[p.~190]{B-CMP122a}.
Further, $\nabla$ is the lattice gradient, and $\alpha_0$ is the small-field radius of the step,
one of the radii among the auxiliary parameters $\mathfrak{p}$. Gauge transformations act on
configurations by $U\mapsto U^{u}$.

The \emph{class of real backgrounds} $\mathcal{U}_{\mathbb{R}}$ is the set of $G$-valued
configurations $U$ on $\Omega_0$ that satisfy the following three conditions.
\begin{enumerate}
\item[(u1)] Let $\mathcal{C}_1$ be the seam-enlarged class of cubes at block size $M_1$, fixed with
the geometry of the step. For every cube of $\mathcal{C}_1$ there is a $G$-valued gauge
transformation $u$ on the cube for which the first line of \eqref{4.14:eq-real-backgrounds-a}
holds on it. This is condition \cite[(3.35)]{B-CMP99b}, read at block size $M_1$, on the class
$\mathcal{C}_1$, and with constant $O(1)M_1\alpha_0$; we write $\mathsf{b}$ for the absolute
constant written $O(1)$ there, fixed with the step, so that this constant is $\mathsf{b}M_1\alpha_0$.
\item[(reg)] Let $Q\subset\Omega_k$ be a cube of side $M$ that contains a bond of $\mathbb{B}^{+}$.
For every such $Q$ there is a $G$-valued gauge transformation $u$ on $Q$ for which the second line
of \eqref{4.14:eq-real-backgrounds-a} holds on $Q$, where $a_Q$ is the gauge constant attached to
the cube $Q$. Compare the condition
\cite[(1.12)]{B-CMP109} and the clause of \cite[(2.38)]{B-CMP119} for the cubes contained in the
innermost region of the nest there, written there as of the size $CM$ with $C$ a positive integer
constant; (reg) is the form of these conditions used here.
\item[(reg)$^{\sharp}$] For every $\square\in\mathcal{D}^{\sharp}$ there is a $G$-valued gauge
transformation $u$ on $\square^{4}\cap\Omega_0$ with the following property. Let $X_n$ be any
territory that $\square^{4}$ meets; there are at most two such territories, and they are adjacent.
Then the third line of \eqref{4.14:eq-real-backgrounds-a} holds on $\square^{4}\cap X_n$.
\end{enumerate}
The three lines referred to are
\begin{equation}\label{4.14:eq-real-backgrounds-a}
\begin{aligned}
&U^{u}=e^{i\eta A},\qquad |A|,\ |\nabla A|<\mathsf{b}M_1\alpha_0 \qquad\text{on the cube},\\
&U^{u}=e^{i\eta A},\qquad |A|,\ |\nabla A|<a_Q \qquad\text{on } Q,\\
&U^{u}=e^{i\eta A},\qquad \ell_n|A|,\ \ell_n^{2}|\nabla A|<a^{\sharp}
\qquad\text{on } \square^{4}\cap X_n .
\end{aligned}
\end{equation}
In all three lines $A$ takes values in the Lie algebra of $G$.

The constant $a^{\sharp}$ in (reg)$^{\sharp}$ is one number, and it enters (reg)$^{\sharp}$ only
as an upper bound. A gauge transformation whose constants on $\square^{4}$ are smaller than
$a^{\sharp}$ satisfies (reg)$^{\sharp}$. Where $a^{\sharp}$ is used in what follows, it enters only
through the ceiling $a^{\sharp}\le a^{\sharp}_{\mathrm{ceil}}$ imposed in the lemma of this chapter
in which (reg)$^{\sharp}$ is used.

Condition (reg)$^{\sharp}$ is supplied in either of two forms:
\begin{itemize}
\item[$(\alpha')$] as \cite[(3.35)]{B-CMP99b} at the second block size $M_1^{\sharp}$;
\item[$(\beta')$] as Ba{\l}aban's nest-indexed clauses \cite[(2.38)]{B-CMP119},
\cite[(1.65), (1.67), (1.69)]{B-CMP122a}.
\end{itemize}
For the form $(\beta')$, the geometry, the class of cubes and the value of $a^{\sharp}$ at every
nest are fixed in the statement of this chapter on the nest-indexed form of (reg)$^{\sharp}$. The
condition \cite[(3.36)]{B-CMP99b} is not among the
defining conditions of $\mathcal{U}_{\mathbb{R}}$.

The data space of the step has the pair of radii $(\alpha_1,\alpha_0^{J})$, fixed with that space.
Let $a_1^{\oplus}$ and $a_0^{\oplus}$ denote the pair of reference constants, fixed with the data
space, against which these radii are compared. The requirement placed on the data space, on the
side of the statements that use it, is, with the inequality between pairs read componentwise,
\begin{equation}\label{4.14:eq-real-backgrounds-b}
(\alpha_1,\alpha_0^{J})\le \tfrac14\,\varepsilon_1\,(a_1^{\oplus},a_0^{\oplus}),
\qquad \text{and the real parts of its points lie in } \mathcal{U}_{\mathbb{R}} .
\end{equation}
At the levels weighted by an $\mathbf{R}$-operation, the constant $\alpha_0$ of (u1) and the radii
in \eqref{4.14:eq-real-backgrounds-b} carry Ba{\l}aban's factor
$\mathsf{L}_0^{2e_n}$,
where $\mathsf{L}_0$ is the auxiliary integer of \cite[(1.24)]{B-CMP122a},
$e_n$ is the $\mathsf{L}_0$-exponent attached to the level $n$ at which the condition is
read and $R_n$ is the block count of \cite[(2.5)]{B-CMP119} at that level; this factor satisfies
$\mathsf{L}_0^{2e_n}\le\mathsf{c}_1R_n$. The exponent $e_n$, the constant
$\mathsf{c}_1$ and this bound are fixed in the statement of this chapter on the powers of
$\mathsf{L}_0$ at the levels weighted by an $\mathbf{R}$-operation. The conditions are otherwise
unchanged.

It is not asserted that the data space of the step is contained in $\mathcal{U}_{\mathbb{R}}$;
condition \eqref{4.14:eq-real-backgrounds-b} requires only that the real parts of its points lie in
$\mathcal{U}_{\mathbb{R}}$.

Finally, $\mathcal{U}_{\mathbb{R}}$ is invariant under $G$-valued gauge transformations. Compare the
corresponding sentence following \cite[(2.38)]{B-CMP119}.
\end{definition}

\begin{proof}[Proof of the invariance]
Let $U\in\mathcal{U}_{\mathbb{R}}$, and let $v$ be a $G$-valued gauge transformation on
$\Omega_0$. Then $U^{v}$ is $G$-valued.

Each of (u1), (reg) and (reg)$^{\sharp}$ is an existence-of-gauge clause. It asserts that, for
every set $S$ of a family fixed by the geometry of the step, there is a $G$-valued gauge
transformation $u$ on $S$ such that $U^{u}=e^{i\eta A}$, with a bound on $A$ and $\nabla A$. For
(u1), (reg) and (reg)$^{\sharp}$ these families are
\begin{itemize}
\item the cubes of the seam-enlarged class at block size $M_1$, on which (u1) is stated,
\item the cubes $Q\subset\Omega_k$ of side $M$ that contain a bond of $\mathbb{B}^{+}$, and
\item the sets $\square^{4}\cap\Omega_0$ with $\square\in\mathcal{D}^{\sharp}$, together with
their intersections with the territories $X_n$.
\end{itemize}
These families are fixed with the nest, the bond set $\mathbb{B}^{+}$ and the cover
$\mathcal{D}^{\sharp}$ of Notation~\ref{4.14:not-free-current-notation}, which are the same for
$U$ and for $U^{v}$. Hence the same families serve for $U$ and for $U^{v}$.

Fix such a set $S$ and the gauge transformation $u$ given for $U$. The action of a gauge
transformation $\tilde u$ on a bond $b$ from $x$ to $y$ is
$U^{\tilde u}(b)=\tilde u(x)\,U(b)\,\tilde u(y)^{-1}$; hence $(U^{v})^{\tilde u}=U^{\tilde u v}$,
where $\tilde u v$ is the pointwise product. Put $u'=u\,v^{-1}$, restricted to $S$;
it is $G$-valued, and $(U^{v})^{u'}=U^{u v^{-1}v}=U^{u}$ on $S$. Therefore
\begin{equation*}
(U^{v})^{u'}=U^{u}=e^{i\eta A}\qquad\text{on } S,
\end{equation*}
with the same $A$. The bound required of $A$ and $\nabla A$ on $S$ is therefore unchanged. This
holds for each of (u1), (reg) and (reg)$^{\sharp}$, and so $U^{v}\in\mathcal{U}_{\mathbb{R}}$.
\end{proof}

\begin{proposition}[Admissibility of the prepared large-field data; printed without proof in
{\cite[pp.~178--179]{B-CMP122a}}]
\label{4.14:prop-prepared-data-admissible}
Let $k\ge 1$. Levels, lengths and territories are as in
Definition~\ref{4.14:not-free-current-notation}. In particular,
$\ell_j=L^{j-k}$ is the lattice spacing of level $j$ measured in units of level $k$, and
$\operatorname{dist}_\infty$ is the distance in the maximum norm in the same units. Let
$\Omega_0,\Omega_1,\dots,\Omega_k$
be domains, and let $\Omega_t$ be a further domain (in the application, the next domain
$\Omega_{k+1}$ of the nest, carved on the lattice of level $k$ by the renormalization
transformation from level $k$ to level $k+1$). Conditions (A1)
and (A2) on these data involve the relations
\begin{equation}\label{4.14:eq-prepared-data-admissible-a}
\begin{aligned}
&\Omega_0\supset\Omega_1\supset\dots\supset\Omega_k,
\qquad
\ell_j^{-1}\operatorname{dist}_\infty\bigl(\Omega_j^{c},\Omega_{j+1}\bigr)\ge M
\quad(1\le j\le k-1),\\
&\Omega_t\subset\Omega_k,
\qquad
\operatorname{dist}_\infty\bigl(\Omega_t,\Omega_k^{c}\bigr)\ge 6M ,
\end{aligned}
\end{equation}
and read as follows.
\begin{itemize}
\item[(A1)] The nesting and the inequality in the first line of
\eqref{4.14:eq-prepared-data-admissible-a} hold. For every $j\ge 1$ the domain
$\Omega_j$ is a union of closed $M$-cubes of level $j$, that is, of closed cubes of side
$M\ell_j$ ($M$ lattice spacings of level $j$), all taken from
partitions compatible with one another under $L$-adic refinement, as in \cite[\S3]{B-CMP119}.
Finally, $\Omega_0$ is
the Dirichlet hull of $\Omega_1$ in the sense of \cite[\S E]{B-CMP99b}.
\item[(A2)] The inclusion and the inequality in the second line of
\eqref{4.14:eq-prepared-data-admissible-a} hold, and $\Omega_t$ is a union of closed
$LM$-cubes of level $k$.
\end{itemize}
For the nest of the prepared large-field data, $\Omega_j=\Omega''_j$ of
\cite[(1.12)]{B-CMP122a} for $1\le j\le k$ and $\Omega_0$ the Dirichlet hull of $\Omega''_1$,
condition (A1) holds.
\end{proposition}

The last sentence of Proposition~\ref{4.14:prop-prepared-data-admissible} is stated in
\cite[pp.~178--179]{B-CMP122a} without proof; cited as printed. There the sets
$Z''_k$ and $Z''_{k-N_0+1}$ are defined in \cite[(1.10)]{B-CMP122a}, where $N_0$ is the
integer introduced immediately before \cite[(1.10)]{B-CMP122a}: the smallest positive integer
with $L^{-N_0+1}MR_{k-N_0+1}=M$, the numbers $R_j$ being the block-size factors of the
$MR_j$-cubes of \cite[p.~256]{B-CMP119}. The new domains $\Omega''_j$ are defined in
\cite[(1.12)]{B-CMP122a}. These are the domains $\Omega''_j$ of the prepared nest named in the
last sentence of Proposition~\ref{4.14:prop-prepared-data-admissible}. The sentences stated
without proof are the following.
\begin{quote}
\dots{} and complete these two sets to a sequence
$Z''_k, Z''_{k-1},\dots,Z''_{k-N_0+2}, Z''_{k-N_0+1}$ in such a way that the complements of
these sets form an admissible sequence of domains based on partitions into $M$-cubes in the
corresponding scales.
\end{quote}
\begin{quote}
The sequence $\{\Omega''_j\}$ is an admissible sequence [A], and the corresponding generating
set is denoted by $\mathbb{B}''_k$.
\end{quote}
The reference [A] in the second quoted sentence is the one \cite{B-CMP122a} gives for
admissible sequences of domains. Nested sequences of domains of this kind, each $\Omega_j$
read on the lattice of level $j$, together with their layers $\Omega_j\setminus\Omega_{j+1}$,
are introduced in \cite[(2.1)--(2.2)]{B-CMP96}.

Condition (A2) is not among the sentences quoted above. Conditions (A1) and (A2) together
enter the operator theorem at free current as a hypothesis. In that
theorem they are assumed, not derived.

The same conditions are supplied, for every nest of domains produced by the succession of
renormalization transformations \cite[\S3]{B-CMP119} and large-field operations
\cite{B-CMP122a}, by the admissibility lemma for the produced nests. That lemma gives them
with the stronger margin
\begin{equation}\label{4.14:eq-prepared-data-admissible-b}
\ell_j^{-1}\operatorname{dist}_\infty\bigl(\Omega_j^{c},\Omega_{j+1}\bigr)\ge LM
\qquad(1\le j\le k-1)
\end{equation}
in place of the inequality in the first line of \eqref{4.14:eq-prepared-data-admissible-a}.

\begin{definition}[One-sided conditions on block sizes and couplings]\label{4.14:def-block-size-conditions}
The conditions below are floors on the block sizes $L$, $M_1$, $M_1^{\sharp}$, $M$ and ceilings on
the couplings of the individual levels. Every condition is one-sided. The numbers depending only
on $(d,L)$ are fixed before $M$. No constant is evaluated.

\emph{Notation.} The form $\mathsf{Q}$, the levels, the territories and the second cover
$\mathcal{D}^{\sharp}_j$ are those of Notation~\ref{4.14:not-free-current-notation}, and the
regularity conditions on $M$-cubes and on the boxes of the second cover are those of the class of
real backgrounds (Definition~\ref{4.14:def-real-backgrounds}).

We write $\mathsf{F}_{u}$, $\mathsf{F}^{P}_{u}$ and $\mathsf{F}_{ms}$, $\mathsf{F}^{P}_{ms}$ for
the floor conditions $\mathsf{F}_{\mathfrak{F}}$, $\mathsf{F}^{P}_{\mathfrak{F}}$ of the inversion
lemma and of the positivity lemma in the unit frame $\mathfrak{F}^{u}$ and in the many-scale frame
$\mathfrak{F}^{ms}$. In each of them the first three arguments are a constant, a decay rate and a
third constant of the operator estimates to which the lemma is applied, and the last three are
numbers attached to the frame; the arguments are written in the order those lemmas use.

The quantities $B_T$, $\delta_T$, $m^{\mathrm{flat}}$, $\varepsilon_T$, $a_1^{\mathrm{el}}$ and
$S_{ms}(\cdot)$ are the $(d,L)$-constants of the many-scale frame. The rate-dependent constant
$c_1(\cdot)$ is that of \cite{B-CMP96}; it is distinct from the number $\mathsf{c}_1$ defined
below. The numbers $B_{\mathsf{Q}}$ and
$\delta_{\mathsf{Q}}$ are the constant and the decay rate of the walk expansion of $\mathsf{Q}$.
The constants $\gamma_0$, $\gamma_1$ and $m_1$ are the $(d,L)$-constants of the operators at free
current. The
radius $r_{\flat}$ and the constant $C'_1$ are the $(d,L)$-constants of the unit-lattice averaging
bounds. The numbers $c_X=c_X(d)$ and $C_d$ depend on $d$ only. We set
\begin{equation*}
\rho_0:=c_XL^2M_1,\qquad \rho_{\mathsf{Q}}:=8dL^2M_1^{\sharp},\qquad
\lambda^{\sharp}:=\frac{C_dL^2}{M_1^{\sharp}},\qquad w^{\sharp}:=\frac{M_1^{\sharp}}{2L},
\qquad N_d:=2\cdot 5^d .
\end{equation*}
We write $M=L^{m'}$; this is an odd integer.

\emph{Floors.} The floors are
\begin{equation}\label{4.14:eq-block-size-conditions-a}
\begin{gathered}
M_1\ge M_F(d,L);\\
M_1^{\sharp}\ge L^2,\qquad M_1^{\sharp}>2L\rho_0,\\
\mathsf{F}_{ms}\bigl(B_T,\delta_T,\tfrac14 m^{\mathrm{flat}};N_d,\lambda^{\sharp},w^{\sharp}\bigr),\\
\mathsf{F}^{P}_{ms}\bigl(B_T,\delta_T,m^{\mathrm{flat}};N_d,\lambda^{\sharp},w^{\sharp}\bigr);\\
M\ge 13L^2M_1^{\sharp},\qquad \frac{M}{L}\ge R_E,\qquad M\ge 144\,d\,\rho_{\mathsf{Q}},\\
\ell=\frac{M}{12d}\ge 6,\qquad M>2L,\\
\mathsf{F}_{u}\Bigl(B_{\mathsf{Q}},\delta_{\mathsf{Q}},m_1;2^d,\frac{36d^2}{M},\frac{M}{48d}\Bigr),\\
\mathsf{F}^{P}_{u}\Bigl(B_{\mathsf{Q}},\delta_{\mathsf{Q}},\gamma_0;2^d,\frac{36d^2}{M},
\frac{M}{48d}\Bigr);\\
L\ \text{odd},\qquad L\ge 3,\qquad L\ge 8,
\end{gathered}
\end{equation}
together with the two floor conditions on $M$ at rate $\underline{a}$ described in item (c) below.
The four groups of lines of~\eqref{4.14:eq-block-size-conditions-a}, separated by semicolons, are
the floors of items (a), (b), (c) and (d) below, in this order, and each group is referred to by
the letter of its item.
In detail:
\begin{enumerate}
\item[(a)] \emph{The floor on $M_1$.} Here $M_F(d,L)$ is a number depending only on $(d,L)$. It
dominates the following thresholds:
\begin{itemize}
\item the block-size thresholds of
\cite[Theorems~3.1--3.4, 3.7, 3.9--3.11]{B-CMP99b};
\item the block-size thresholds of \cite[Propositions~2.6--2.7]{B-CMP96};
\item a further $(d,L)$-threshold $M_E$;
\item the threshold attached to the collar of the enlarged domain built around a region on which
the construction of the background is carried out;
\item the threshold below which the inequality $b(d,L)C_{\mathrm{el}}M_1^{-1/2}\le\frac12$ fails;
here $b(d,L)$ and $C_{\mathrm{el}}$ are numbers not depending on $M_1$, the former depending only
on $(d,L)$.
\end{itemize}
\item[(b)] \emph{The floor on $M_1^{\sharp}$.} The conditions
of~\eqref{4.14:eq-block-size-conditions-a} on $M_1^{\sharp}$, namely the two inequalities
$M_1^{\sharp}\ge L^2$, $M_1^{\sharp}>2L\rho_0$ and the two conditions $\mathsf{F}_{ms}$,
$\mathsf{F}^{P}_{ms}$, form a floor on $M_1^{\sharp}$. It is read after $M_1$.
\item[(c)] \emph{The floors on $M$.} These are the conditions
of~\eqref{4.14:eq-block-size-conditions-a} involving $M$, together with two conditions at a rate
$\underline{a}$. Let $\underline{a}$ be one
quarter of the smallest decay rate used in this chapter; these rates include
$\delta_4^{\ast}/8$, where $\delta_4^{\ast}$ is one of the decay rates of the estimates of this
chapter. The floors on $M$ are the following five conditions.
\begin{enumerate}
\item[(i)] $M\ge 13L^2M_1^{\sharp}$.
\item[(ii)] $M$ satisfies, at rate $\underline{a}$, the floor which, once it holds, makes the
explicit expressions $S_{ms}(\underline{a})$ and $c_1(\underline{a})$ valid bounds, and a second
floor on $M$ at the same rate; moreover $M/L\ge R_E$. Here $R_E$ is a number fixed before $M$, so
that $M/L\ge R_E$ is a floor on $M$.
\item[(iii)] $M\ge 144\,d\,\rho_{\mathsf{Q}}$, $\ell=M/12d\ge 6$ and $M>2L$.
\item[(iv)] The condition $\mathsf{F}_{u}$ of~\eqref{4.14:eq-block-size-conditions-a}.
\item[(v)] The condition $\mathsf{F}^{P}_{u}$ of~\eqref{4.14:eq-block-size-conditions-a}.
\end{enumerate}

The number $\ell=M/12d$ is a real side length. By definition $M/\ell=12d\in\mathbb{N}$. Since $M$
is odd and $12d$ is even, $\ell$ is not an integer, and no condition below requires it to be one.
The boxes of side $\ell$ used below are centred at the points of $\ell\mathbb{Z}^d$, cut-off
functions at scale $\ell$ take real arguments, and $M\in\ell\mathbb{Z}$.

All these are floors on $M$ alone. They are read after $d$, $L$, $M_1$, $M_1^{\sharp}$.
\item[(d)] \emph{The floor on $L$.} The conditions of~\eqref{4.14:eq-block-size-conditions-a} on
$L$ include the condition $L\ge 8$ required by the uniform localisation. They form a floor on $L$
only, of the same kind as the conditions through which the constant $L_0=L_0(N)$ of Theorem~0 is
chosen.
\end{enumerate}

\emph{Ceilings.} Let $R_n$ be the block count of \cite[(2.5)]{B-CMP119} and $\alpha_{0,n}$ the
radius of \cite[(2.28)]{B-CMP119} at level $n$. Let $\mathsf{L}_0$ be Ba{\l}aban's auxiliary
integer of \cite[(1.24)]{B-CMP122a}. Let $B$ be Ba{\l}aban's absolute constant of
\cite[(2.38)]{B-CMP119} and \cite[(1.65)]{B-CMP122a}, and let $\beta_0$ be Ba{\l}aban's exponent
in \cite[(2.6)--(2.9)]{B-CMP119}. Set
\begin{equation}\label{4.14:eq-block-size-conditions-1}
c(\beta_0):=\sup_{t\ge1}3^{-t}t^{\beta_0},\qquad \mathsf{c}_1:=1+c(\beta_0)L(L+1).
\end{equation}
The number $c(\beta_0)$ is finite, and it is fixed by $\beta_0$. Let $a_{\mathrm{pl}}$ be the
constant bounding the plaquette variables that is implied by the regularity condition on
$M$-cubes, and let $\bar a_0\le r_{\flat}$ be the bound on the unit plaquettes of the $k$-fold
block average of $U$. The numbers $\mathsf{o}_1$ and $\mathsf{o}_2$ (the latter used
in~\eqref{4.14:eq-block-size-conditions-d}) are two absolute constants of the unit-lattice
averaging bounds. The ceilings are
\begin{equation}\label{4.14:eq-block-size-conditions-b}
\begin{gathered}
C(d)L^dM_1\,\mathsf{c}_1R_n\alpha_{0,n}\le a_F(d,L)\quad\text{for every layer }n;\\
a_Q\le\frac{a_1^{\oplus}\gamma_0}{24\gamma_1},\qquad 2\,\mathsf{o}_1a_{\mathrm{pl}}\le r_{\flat},
\qquad C'_1L^2\bar a_0\le\frac12;\\
a^{\sharp}\le a^{\sharp}_{\mathrm{ceil}}
:=\frac{a_1^{\mathrm{el}}\,m^{\mathrm{flat}}}{24\,e\,L^2B_TS_{ms}(\delta_T)} ,
\end{gathered}
\end{equation}
where $C(d)$ depends on $d$ only and $a_F(d,L)$ on $(d,L)$ only; the constant $C(d)$ need not
coincide with the constant $C_d$ above, and $e$ is the base of the natural logarithm. The three
lines of~\eqref{4.14:eq-block-size-conditions-b} are the ceilings of items (e), (f) and (g) below,
in this order, and each is referred to by the letter of its item. In detail:
\begin{enumerate}
\item[(e)] \emph{The first line.} Its basic form is $C(d)L^dM_1\alpha_0\le a_F(d,L)$, where
$\alpha_0$ is the constant with which the condition \cite[(3.35)]{B-CMP99b} holds at block size
$M_1$. The condition is read layer by layer, with $\alpha_0$ replaced by
$\mathsf{c}_1R_n\alpha_{0,n}$. The reason is as follows.
\begin{itemize}
\item At the levels weighted by an $\mathbf{R}$-operation, the bounds (i), (ii) defining
Ba{\l}aban's analyticity space in \cite{B-CMP122a} carry the factor $\mathsf{L}_0^{2e(Q)}$ on all
seven of their members, including the bound on $|U(\partial p)-1|$.
\item Hence the constant with which that space supplies \cite[(3.35)]{B-CMP99b} there is
$\mathsf{L}_0^{2e(Q)}\alpha_{0,n}$.
\item This constant satisfies $\mathsf{L}_0^{2e(Q)}\alpha_{0,n}\le\mathsf{c}_1R_n\alpha_{0,n}$, by
the value lemma of this chapter and its per-operation form.
\end{itemize}
The condition keeps its form and its place in the order of choice, and it is not a cap in the
sense of Definition~\ref{2.3:def-cap}. At such levels neither the regularity condition on
$M$-cubes nor the second line of~\eqref{4.14:eq-block-size-conditions-b} is changed. The reason
is that the cubes on which they are read lie at distance at least $8LMR_t-2M$ from the post-step
large-field region $Z_k^{\mathrm{post}}$.
\item[(f)] \emph{The second line.} These are ceilings on the regularity constant $a_Q$ on
$M$-cubes, that is, on a quantity of the form $O(1)LMB\alpha_0$. They are read after $M$.
\item[(g)] \emph{The third line.} This is the ceiling read by the estimates in the many-scale
frame. The factor $L^2$ in its denominator leaves room for the weights attached to the
territories.
\end{enumerate}

\emph{The layerwise ceiling.} For every level $n\in\{1,\dots,k\}$,
\begin{equation}\label{4.14:eq-block-size-conditions-c}
L^2\,\mathsf{c}_1\,B\,M\,R_n\alpha_{0,n}\le a^{\sharp}_{\mathrm{ceil}} .
\end{equation}
This is a ceiling on each layer's own quantity $MR_n\alpha_{0,n}$, which is $M$ times a
polylogarithmic factor in $g_n$ times $g_n$. It is read after $d$, $L$, $M_1$, $M_1^{\sharp}$,
$M$. It has the form of Ba{\l}aban's smallness conditions in \cite[p.~192]{B-CMP122a} and of the
smallness of $M\alpha_0$ required in \cite{B-CMP99b}, with one more logarithm. It
implies the third line of~\eqref{4.14:eq-block-size-conditions-b} for the value
$a^{\sharp}:=\max_{\square}\mathfrak{a}(\square)$ at which the regularity condition on the boxes
$\square$ of the second cover is supplied, where $\mathfrak{a}(\square)$ is the boxwise value of
its constant.

\emph{Radii and order of choice.} The radius $a_1^{\oplus}$ and the order in which
the constants are chosen are
\begin{equation}\label{4.14:eq-block-size-conditions-d}
\begin{gathered}
a_1^{\oplus}:=\min\Bigl\{\varepsilon_Ta_1^{\mathrm{el}},\ \frac{r_{\flat}}{2\mathsf{o}_2},\
\frac{r_{\flat}^2}{16L^{d+1}\mathsf{o}_2},\ \frac{\log 2}{4dL\,\mathsf{o}_2}\Bigr\},\\
d,L\ \to\ M_1\ \to\ M_1^{\sharp}\ \to\ (d,L)\text{-constants}\ \to\ \kappa_1\ \to\ M\\
\to\ \text{layer couplings}\ \to\ \text{radii}\ \to\ g .
\end{gathered}
\end{equation}
The second radius $a_0^{\oplus}$ is the companion radius fixed together with $a_1^{\oplus}$; it is
the only place at which the constant $c_{\mathfrak{g}}$ enters.
The $(d,L)$-constants of this chain include
\begin{equation*}
B_T,\ \delta_T,\ m^{\mathrm{flat}},\ S_{ms}(\cdot),\ c_1(\cdot),\ r_{\flat},\
a_1^{\mathrm{el}},\ c(\beta_0)L(L+1).
\end{equation*}
The constant $\kappa_1$ of the chain is fixed after these $(d,L)$-constants and before $M$.
The layer couplings are constrained through $a_Q$ and $a^{\sharp}$ by the first two lines
of~\eqref{4.14:eq-block-size-conditions-b} and by~\eqref{4.14:eq-block-size-conditions-c}.
Each of these is a condition $g_n\le g^{\flat}(d,L,M)$ on the layer's own coupling $g_n$, with
$g^{\flat}(d,L,M)>0$ depending only on $d$, $L$, $M$. No quantity in the chain is chosen in
terms of a quantity standing to its right.
\end{definition}

\begin{definition}[Ba{\l}aban's small-coupling regime for large-field operations]\label{4.14:def-small-coupling-regime}
Let $g_j$ be the running coupling at level $j$, and let $r>0$ and $\beta_0>0$ be the fixed
exponents of \cite[(2.5), (2.6)]{B-CMP119}. Let $R_j$ be the number of \cite[(2.5)]{B-CMP119}:
the smallest number of the form $L^{a}$, with $a$ a non-negative integer, such that
$R_j\ge(\log g_j^{-2})^{r}$. Let $\mathsf{L}_0$ be Ba{\l}aban's auxiliary integer of
\cite[(1.24)]{B-CMP122a}; it is not the constant $L_0$ of Theorem~0. For an operation
$\mathbf{R}$ performed at level $k$, let $N_k$ be its number of preliminary integrations.
Let $N_0(k)$ be the smallest positive integer $n$ with $L^{-n+1}R_{k-n+1}=1$, whenever such an
integer exists \cite[p.~179]{B-CMP122a}.

\emph{Ba{\l}aban's small-coupling regime} consists of the following four requirements: the size
condition on $R$, the size condition on $N_k$, the size condition on $\mathsf{L}_0$ and the size
condition on $\beta_0$, displayed in this order. The two inequalities relating $R_i$ and $R_j$
hold for every pair of levels $j<i$ of
one sequence of renormalisation steps; the conditions on $R_k$, $N_0(k)$ and $N_k$ are imposed at
every level $k$ at which an operation $\mathbf{R}$ acts.
\begin{equation}\label{4.14:eq-small-coupling-regime-a}
\begin{aligned}
&\text{on $R$:}\quad R_j\in\{1,L,L^2,\dots\},\qquad 1\le R_j<L\max\{1,(\log g_j^{-2})^{r}\};\\
&\qquad\qquad R_i\le LR_j,\qquad R_j\le(L+1)(i-j)^{\beta_0}R_i;\\
&\qquad\qquad R_k>1;\\
&\text{on $N_k$:}\quad N_0(k)\ \text{exists, and}\ k>N_k>N_0(k);\\
&\text{on $\mathsf{L}_0$:}\quad \mathsf{L}_0\ge1,\qquad \mathsf{L}_0^2\le\tfrac13L;\\
&\text{on $\beta_0$:}\quad \beta_0>0\ \text{is fixed}.
\end{aligned}
\end{equation}
No operation $\mathbf{R}$ is performed at a level at which the size condition on $N_k$ fails.

The two inequalities relating $R_i$ and $R_j$ are \cite[(2.9)]{B-CMP119}. The condition
$N_k>N_0(k)$ is
assumed in \cite[p.~179]{B-CMP122a}. In \cite{B-CMP122a} the number $N_k$ is taken to be a
positive power of $\log g_k^{-2}$. The bound on $\mathsf{L}_0^2$ is that of
\cite[pp.~191--192]{B-CMP122a}. Ba{\l}aban also requires
$2\le\mathsf{L}_0<\tfrac12L$ \cite[(1.24)]{B-CMP122a}.

The number $N_0$ is introduced through this equality in \cite[p.~179]{B-CMP122a}; in the regime,
the existence of $N_0(k)$ at a level $k$ where an $\mathbf{R}$ acts is one of the requirements.
Of $\beta_0$, Ba{\l}aban notes that it ``can be chosen arbitrarily small,
if $g$ is sufficiently small'' \cite[p.~255]{B-CMP119}. No use of the regime in what follows
requires $\beta_0$ to be small.

Put
\begin{equation*}
c(\beta_0):=\sup_{t\ge1}3^{-t}t^{\beta_0}.
\end{equation*}
In the regime \eqref{4.14:eq-small-coupling-regime-a} the following hold.
\begin{enumerate}
\item[(a)] $R_{j+1}\ge R_j/L$ for every level $j$.
\item[(b)] At a level $k$ where an $\mathbf{R}$ acts, $R_k>1$ holds if and only if
$g_k^2<e^{-1}$. Moreover $R_k\ge L$, and $R_k\ge3$ under the floor on $L$ in
Definition~\ref{4.14:def-block-size-conditions}.
\item[(c)] For $1\le n<k$, the ratio of the $(n+1)$-st to the $n$-th term of the sequence
$n\mapsto L^{-n+1}R_{k-n+1}$ is a power of $L$ lying in $[L^{-2},1]$.
\item[(d)] At a level $k$ where an $\mathbf{R}$ acts, $N_0(k)\ge2$ and
\begin{equation}\label{4.14:eq-small-coupling-regime-1}
N_0(k)\le1+\log_L\bigl[(L+1)(N_0(k)+1)^{\beta_0}R_k\bigr].
\end{equation}
Consequently $N_0(k)$ is at most of order $\log\log g_k^{-2}$, and, with $N_k$ a positive power
of $\log g_k^{-2}$ as in \cite{B-CMP122a}, the requirement
$N_k>N_0(k)$ is a one-sided smallness condition on $g_k$ at fixed $k$, which a larger $L$ only
relaxes.
\item[(e)] The requirements $\mathsf{L}_0\ge2$ and $\mathsf{L}_0^2\le\frac13L$ of
\cite{B-CMP122a} together force $L\ge12$.
\item[(f)] $c(\beta_0)<\infty$ for every $\beta_0>0$, and $c(\beta_0)\le1$ when $\beta_0\le1$.
\end{enumerate}
\end{definition}

\begin{proof}
(a) Apply the inequality $R_j\le(L+1)(i-j)^{\beta_0}R_i$ of the size condition on $R$ to the
pair of levels $j<j+1$.
Since
$(i-j)^{\beta_0}=1$ for $i=j+1$, this gives $R_j\le(L+1)R_{j+1}$. Both numbers are powers of
$L$, so $R_j/R_{j+1}=L^{a}$ for some integer $a$. Then $L^{a}\le L+1<L^2$, hence $a\le1$, that is
$R_{j+1}\ge R_j/L$.

(b) Suppose first that $(\log g_k^{-2})^r>1$. Then the defining inequality of $R_k$ gives
$R_k\ge(\log g_k^{-2})^r>1$. Conversely, suppose $R_k>1$. Then $R_k=L^{a}$ with $a\ge1$. The
number $L^{a-1}$ is of the admitted form and is smaller than $R_k$. By the minimality of $R_k$ it
fails the defining inequality, so $(\log g_k^{-2})^r>L^{a-1}\ge1$. The definition of $R_k$
through \cite[(2.5)]{B-CMP119} requires $\log g_k^{-2}>0$, that is $g_k<1$, for the power
$(\log g_k^{-2})^{r}$ to be defined.
Since $r>0$, we have
$(\log g_k^{-2})^r>1$ if and only if $\log g_k^{-2}>1$, that is, if and only if $g_k^2<e^{-1}$.
This condition does not involve $L$ or any other parameter, and it is one-sided. Finally,
$R_k$ is a power of $L$ exceeding $1$, so $R_k\ge L$. The floor on $L$ in
Definition~\ref{4.14:def-block-size-conditions} then gives $R_k\ge3$.

(c) Let $1\le n<k$. The ratio of the $(n+1)$-st to the $n$-th term is
$R_{k-n}/(LR_{k-n+1})$, which is a power of $L$ because the $R_j$ are. By (a) with $j=k-n$ we
have $R_{k-n}\le LR_{k-n+1}$, so the ratio is at most $1$. By the inequality $R_i\le LR_j$ of
the size condition on $R$, applied to the pair $k-n<k-n+1$, we have $R_{k-n+1}\le LR_{k-n}$,
so the ratio is at least
$L^{-2}$. Under the size condition on $R$ alone, these ratio bounds therefore allow a step from
the value $L$ to the value $L^{-1}$, which passes over the value $1$; for this reason the
existence of $N_0(k)$ is listed as a separate requirement in the size condition on $N_k$.

(d) For $n=1$ the defining equality of $N_0(k)$ reads $R_k=1$, which contradicts $R_k>1$.
Hence $N_0(k)\ge2$, and the defining equality gives $L^{N_0(k)-1}=R_{k-N_0(k)+1}$. Apply the
inequality $R_j\le(L+1)(i-j)^{\beta_0}R_i$ of the size condition on $R$ to the levels
$j=k-N_0(k)+1<i=k$, whose difference
is $N_0(k)-1$.
Since $N_0(k)-1<N_0(k)+1$ and $\beta_0>0$, this gives
\begin{equation*}
L^{N_0(k)-1}\le(L+1)(N_0(k)-1)^{\beta_0}R_k\le(L+1)(N_0(k)+1)^{\beta_0}R_k.
\end{equation*}
Taking $\log_L$ gives \eqref{4.14:eq-small-coupling-regime-1}.

By (b) we have $(\log g_k^{-2})^r>1$. The bound $R_k<L\max\{1,(\log g_k^{-2})^r\}$ therefore
reads $R_k<L(\log g_k^{-2})^r$, so that $\log_LR_k<1+r\log_L\log g_k^{-2}$. Hence
\begin{equation*}
N_0(k)<2+\log_L(L+1)+\beta_0\log_L(N_0(k)+1)+r\log_L\log g_k^{-2}.
\end{equation*}
The term $\beta_0\log_L(N_0(k)+1)$ grows more slowly than $N_0(k)$. So $N_0(k)$ is at most of
order $\log\log g_k^{-2}$. As $L$ grows, $\log_L(L+1)=\log(L+1)/\log L$ decreases, and so do
$\beta_0\log_L(N_0(k)+1)$ and $r\log_L\log g_k^{-2}$, whose arguments exceed $1$ (the second by
(b)); so a larger $L$ only improves the bound. In the
choice of \cite{B-CMP122a}, $N_k$ is a positive power of $\log g_k^{-2}$. Therefore the
requirement $N_k>N_0(k)$ of the size condition on $N_k$ is met once $g_k$ is small enough. It
is thus a one-sided
smallness condition on $g_k$ at the level $k$, which a larger $L$ only relaxes.

(e) From $\mathsf{L}_0\ge2$ we get $4\le\mathsf{L}_0^2\le\frac13L$, hence $L\ge12$. This is a
one-sided lower bound on $L$ implicit in the choices of \cite{B-CMP122a}. Such a lower bound is
absorbed in the requirement $L\ge L_0$ of Theorem~0; the number $12$ is not used as a threshold
anywhere below.

(f) The function $t\mapsto3^{-t}t^{\beta_0}$ is continuous on $[1,\infty)$ and tends to $0$ as
$t\to\infty$. Hence its supremum over $[1,\infty)$ is finite. Now let $\beta_0\le1$ and
$t\ge1$. Then $t^{\beta_0}\le t$, and
\begin{equation*}
3^{t}=e^{t\log3}\ge1+t\log3>t.
\end{equation*}
So $3^{-t}t^{\beta_0}<1$ for every $t\ge1$, and therefore $c(\beta_0)\le1$.
\end{proof}

\begin{definition}[The components a large-field operation acts on]\label{4.14:def-operation-components}
Let an operation $\mathbf{R}$ act at level $k$, let $N=N_k$ be the number of its preliminary
integrations, fixed in Definition~\ref{4.14:def-small-coupling-regime}, and put $h:=k-N_k$.
As in Ba{\l}aban's construction, $Z_k$ denotes the large-field
region at level $k$ (the complement of the small-field region), $M$ and $R_k$ are the parameters
of Definition~\ref{4.14:def-small-coupling-regime}, and $\mathbf{T}_{j\to j+1}$ denotes the one-step
renormalisation from level $j$ to level $j+1$.

The class is defined by induction on the level: for an operation $\mathbf{R}$ at an earlier level
$k'$, its class is the set defined below at the level $k'$. During the operation at level $k'$
this class is redefined at each of its $N_{k'}$ preliminary integrations: a component in which a
new large-field characteristic function is introduced is removed from the class; such a component
is called a \emph{drop-out} of the operation at level $k'$. The class is redefined without it, and
Ba{\l}aban treats the drop-out like the remaining large-field
regions \cite[\S1]{B-CMP122a}; if this happens when the integrations in it are finished at some
$n$, the terms $\mathbb{B}_{k'}^{(n)}$ are left as a part of all boundary terms
\cite[p.~191]{B-CMP122a}. Each drop-out $X'$ carries a state $s(X')$, the integer with
$0\le s(X')\le N_{k'}$ assigned to it in the lemma on the components excluded during an operation.
The \emph{second-class components} of the operation at level $k'$ are the components $Y$ of the
class that remains after all $N_{k'}$ preliminary integrations which carry a large-field
characteristic function of the final step of that operation; on them the new operation has the
form \cite[(1.102)]{B-CMP122b}.

The class $\mathsf{Z}\subset Z_k$ on which the operation $\mathbf{R}$ at level $k$ acts is
\begin{equation}\label{4.14:eq-operation-components-a}
\mathsf{Z}=\bigcup\bigl\{X:\ X \text{ a connected component of } Z_k
  \text{ satisfying (i), (ii-a), (ii-b)}\bigr\},
\end{equation}
where:
\begin{enumerate}
\item[(i)] $X$ is contained in a cube of size $100MR_k$;
\item[(ii-a)] $X$ contains no new large-field region created by one of the operations
$\mathbf{T}_{j\to j+1}$, $h\le j\le k-1$, and every previous region contained in $X$ satisfies (i)
on its own scale;
\item[(ii-b)] $X$ contains no component $X'$ of the class of an earlier operation $\mathbf{R}$ at a
level $k'\in(k-N_k,k)$ in which a large-field function was introduced during that operation.
\end{enumerate}
Condition (ii-b) excludes both kinds of such components: every drop-out $X'$ of
state $0\le s(X')\le N_{k'}$ of such an operation, and every second-class component $Y$
of such an operation.

Condition (i) is Ba{\l}aban's condition (i), and conditions (ii-a), (ii-b) render his condition
(ii) on the components of the large-field region \cite[p.~177]{B-CMP122a}. Condition (ii-a) is his
own description of condition (ii): in the last $N$ one-step operations no large-field
characteristic functions are included \cite[\S1]{B-CMP122a}; condition (ii-b) is a reading of
condition (ii), explained below. In the analysis of the
operation $\mathbf{R}$ on its class later in this chapter only condition (ii-b) is used. It is
used, moreover, only for components whose nest of domains was modified by the earlier operation
and not filled, a component being \emph{filled} when, as in \cite[(1.33)]{B-CMP122b}, the new
determining set on the whole component consists of the component itself at level $k'$. These are,
first, the drop-outs $X'$ with $1\le s(X')\le N_{k'}$. Second, where the bookkeeping of the
large-field terms adopted in this chapter keeps a second-class component $Y$ on its stage nest
$\mathcal{N}^{(N)}$, the nest of domains after $N$ preliminary integrations, instead of filling it,
they are those components $Y$. The second of these is Ba{\l}aban's own instance
of condition (ii): ``The components with the second, large field function are included again into
the large field domain'' \cite[\S1]{B-CMP122b}. Condition (i), condition (ii-a), the exclusion of
the components of state $0$ and the condition on $N_k$ of
Definition~\ref{4.14:def-small-coupling-regime} are retained because they are part of Ba{\l}aban's
definition.

The window $(k-N_k,k)$ for the level $k'$ is the one fixed here. Ba{\l}aban's phrase ``in the
preceding $N$ renormalization steps'' may also be read as covering the operation
$\mathbf{R}^{(h)}$ at level $h$. Under that wider reading Ba{\l}aban's condition excludes more
components than (ii-b) does, so his class is contained in $\mathsf{Z}$. The geometric arguments
based on (ii-b) need only $k'\ge h+1$.

Condition (ii-b) requires that a component into which an operation $\mathbf{R}$ introduced a large
field remains excluded for $N$ steps. This requirement is a reading of Ba{\l}aban's condition (ii),
not a statement printed by him. Ba{\l}aban defines the class by condition (ii), and he applies that
condition to his own components carrying a new large-field function at the moment they arise, by
excluding them from the class. This reading is the rule adopted in this book.
\end{definition}

\begin{theorem}[The operator theorem at free current]\label{4.14:thm-operator-theorem}
Use the levels, territories, covers, bond sets, the cores $\mathsf{Z}_0$, the form $\mathsf{Q}$ at
free current, the operators $\mathsf{A}$, $\mathsf{C}$, $\mathsf{J}$, the balls $\mathcal{U}(1)$ and
$\mathcal{U}(\varepsilon_1)$ and
the constants $\gamma_0$, $\gamma_1$, $\gamma_{+}$, $\gamma_{-}$, $\varepsilon_1$, $m_1$, $\delta_0$,
$B_0$, $B_{\mathsf{J}}$, $B_{\mathsf{Q}}$, $\delta_{\mathsf{Q}}$, $\rho_{\mathsf{Q}}$
and $\ell=M/12d$ of Definition~\ref{4.14:not-free-current-notation}, and the class
$\mathcal{U}_{\mathbb{R}}$ of real backgrounds of Definition~\ref{4.14:def-real-backgrounds}.
Assume the following.
\begin{itemize}
\item The nest $\{\Omega_j\}_{j\le k}$ and $\Omega_t$ satisfy conditions (A1) and (A2) of
\eqref{4.14:eq-prepared-data-admissible-a}.
\item The floors on $M_1$, on $M_1^{\sharp}$, on $M$ and on $L$ of
\eqref{4.14:eq-block-size-conditions-a} hold
(Definition~\ref{4.14:def-block-size-conditions}).
\item The ceilings of \eqref{4.14:eq-block-size-conditions-b} hold
(Definition~\ref{4.14:def-block-size-conditions}).
\end{itemize}
Let $\mathbb{B}\subset\mathbb{B}^{+}$ be arbitrary, let $\mathbb{B}'=\mathbb{B}^{+}\setminus\mathbb{B}$,
$\mathsf{A}=\mathsf{Q}_{\mathbb{B}\mathbb{B}}$, $\mathsf{J}=\mathsf{Q}_{\mathbb{B}\mathbb{B}'}$ and,
where it exists, $\mathsf{C}=\mathsf{A}^{-1}$. Then the following hold; no clause is lettered (J),
so that clause letters are not confused with the operator $\mathsf{J}$.
\begin{enumerate}
\item[(A)] \emph{Objects.} The series defining $\mathsf{Q}$ in
Definition~\ref{4.14:not-free-current-notation} converge, and $\mathsf{Q}$ is analytic on
$\mathcal{U}(1)$; on $\mathcal{U}(\varepsilon_1)$ the inverse $\mathsf{C}=\mathsf{A}^{-1}$ exists;
the operators $\mathsf{Q}$, $\mathsf{A}$, $\mathsf{C}$, $\mathsf{J}$ are analytic, and $\mathsf{Q}$,
$\mathsf{A}$, $\mathsf{C}$ are complex symmetric; $\mathsf{A}\mathsf{C}=1$ on all of
$\mathcal{U}(\varepsilon_1)$. On $\mathfrak{G}\cap\mathcal{U}(\varepsilon_1)$, where
$\mathfrak{G}=\{J=J(\mathbf{U})\}$, let $A_k\in\mathfrak{g}^{\mathbb{B}^{+}}$ be the field of the
fluctuation integral of Ba{\l}aban's construction, written $A_k=\varphi\oplus A'$ with $\varphi$
supported on $\mathbb{B}$ and $A'$ on $\mathbb{B}'$; then the quadratic form in $A_k$ of that
integral is
\begin{equation}\label{4.14:eq-operator-theorem-a}
\tfrac12\langle \varphi,\mathsf{A}\varphi\rangle+\langle \varphi,\mathsf{J}A'\rangle
+\tfrac12\langle A',\mathsf{Q}_{\mathbb{B}'\mathbb{B}'}A'\rangle ,
\end{equation}
so that $\mathsf{C}$ and $\mathsf{J}$ are the covariance and the source operator of the
conditional integral over $A_k|_{\mathbb{B}}$. At $\mathbb{B}=\mathbb{B}^{\ast}$ these are the
operators of \cite[(3.23), (3.25)]{B-CMP119}; at $\mathbb{B}=\mathbb{B}^{K1}$ they are the
covariance and source operator of the conditional integral over $A_k|_{\mathbb{B}^{K1}}$.
\item[(B)] \emph{Bounds at every base point.} At every point $(U,0)$ with
$U\in\mathcal{U}_{\mathbb{R}}$, $\mathsf{Q}$ is real symmetric with
$\gamma_0/2\le\mathsf{Q}\le\gamma_1$, and $\mathsf{C}$ is real symmetric with
$\gamma_{-}\le\mathsf{C}\le\gamma_{+}$. At every real point of $\mathcal{U}(\varepsilon_1)$,
$\gamma_0/4\le\mathsf{Q}\le\gamma_1$ and $\gamma_{-}\le\mathsf{C}\le2\gamma_{+}$. On
$\mathcal{U}(\varepsilon_1)$, $\operatorname{Re}\mathsf{Q}\ge\gamma_0/4$.
\item[(C)] \emph{The source operator.} $\mathsf{J}$ is real at real points,
$\mathsf{J}^{T}=\mathsf{Q}_{\mathbb{B}'\mathbb{B}}$, and $\|\mathsf{J}\|\le\gamma_1$ on
$\mathcal{U}(1)$.
\item[(D)] \emph{Sheets.} In the sheet format of the unit frame,
$\mathsf{Q}\in\mathcal{S}_{u}[\rho_{\mathsf{Q}};B_{\mathsf{Q}},\delta_{\mathsf{Q}};1]$, and for every
$x\ge0$
\begin{equation*}
(x+\mathsf{A})^{-1}\in\mathcal{S}_{u}\Bigl[\frac{M}{4};\,\frac{B_0m_1}{m_1+2x},\,\delta_0;\,
\varepsilon_1\Bigr],
\end{equation*}
with walks and domains independent of $x$.
\item[(E)] \emph{Expansion properties of $\mathsf{C}$.} $\mathsf{C}$ obeys the walk-expansion properties
$(\mathrm{W}1)$, $(\mathrm{W}2)$, $(\mathrm{W}3)'$, $(\mathrm{W}5)^{\mathrm{cov}}$,
$(\mathrm{W}6)^{\mathrm{geo}}$ with the constants $(B_0,\delta_0,\varepsilon_1)$. Here
$(\mathrm{W}1)$--$(\mathrm{W}6)$ are the properties of the walk-expansion format fixed with the
operator theorem at free current with the hulls $\widehat{\square}$, the statement whose domain
lists of the one-box terms carry the hulls $\widehat{\square}$ where the domain lists of the
present theorem carry the reading hulls $\square^{\mathsf{r}}$
(clause (B) above is the property $(\mathrm{W}4)$ for
$\mathsf{Q}$ and $\mathsf{C}$), and $(\mathrm{W}3)'$, $(\mathrm{W}5)^{\mathrm{cov}}$,
$(\mathrm{W}6)^{\mathrm{geo}}$ are the variants of $(\mathrm{W}3)$, $(\mathrm{W}5)$,
$(\mathrm{W}6)$ fixed with that format. Further,
$(\mathrm{W}6)^{\mathrm{geo}}$ is asserted for labels coinciding on an open neighbourhood of
$[\nu]$, the list $[\nu]$ consisting of the reading hulls $\square^{\mathsf{r}}$ of its boxes.
\item[(F)] \emph{The source operator's expansion.} $\mathsf{J}$ satisfies clause (F) of the
operator theorem at free current with the hulls $\widehat{\square}$, read with the enlarged domain
lists described after the present theorem.
\item[(G)] \emph{Resolvents.} Clause (G) of the operator theorem at free current with the hulls
$\widehat{\square}$ holds, read with the same enlarged domain lists.
\item[(H)] \emph{Local balls.} Let $U\in\mathcal{U}_{\mathbb{R}}$ be given globally, and let
$(A',J)$ be given only on
\begin{equation*}
[\nu]^{\approx}=\{x:\ \tau(\{x\},[\nu])<M/2L\},
\end{equation*}
a subset of the one-cube collar $[\nu]^{\sim}$ by the localisation lemma with cut-off distance
$M/2L$ (so that data given on $[\nu]^{\sim}$ suffice a
fortiori), with each of the quantities entering the norm $\|(A',J)\|_{\star}$ less than
$\varepsilon_1/2$ there. Then, for each operator $T$ of (E)--(G), the term $T_\nu$ is defined and
analytic, and the properties $(\mathrm{W}3)'$, $(\mathrm{W}5)^{\mathrm{cov}}$,
$(\mathrm{W}6)^{\mathrm{geo}}$ of (E)--(G) hold for it at the radius $\tfrac12\varepsilon_1$.
\item[(I)] \emph{Uniformity.} The constants $B_0$, $B_{\mathsf{J}}$, $\delta_0$, $\delta_{\mathsf{Q}}$,
$\gamma_0$, $\gamma_1$, $\varepsilon_1$ and the radii $a_1^{\oplus}$, $a_0^{\oplus}$ of
Definition~\ref{4.14:def-block-size-conditions} depend on $d$ and $L$ only;
they are independent of $M$, $k$, $K$, the nest, $\mathbb{B}$, $\Lambda_t$, $\mathsf{Z}_0$, the
volume and $g$.
\item[(K)] \emph{The instance $\mathbb{B}=\mathbb{B}^{K1}$.} The three statements of clause (K)
of the operator theorem at free current with the hulls $\widehat{\square}$ hold, read with the
enlarged domain lists described after the present theorem.
\end{enumerate}
Moreover:

\emph{The covariance statement.}\ For every lattice symmetry $r$ of the group $\mathfrak{r}$ (the $2^d$ reflections
and the $d!$ permutations preserving the partitions) the index sets transport by
$\nu\mapsto r\nu$, and
\begin{equation*}
T_{r\nu}[r\cdot\mathrm{data}](r\cdot\mathsf{u})
=\mathcal{R}_r\,T_\nu[\mathrm{data}](\mathsf{u})\,\mathcal{R}_r^{-1}
\end{equation*}
for $T\in\{\mathsf{C},\mathsf{J},(x+\mathsf{A})^{-1},\mathcal{K}\}$, where
$\mathsf{u}=(\mathbf{U},J)$, $\mathcal{K}$ is the Schur complement defined below, $\mathrm{data}$
stands for the geometric data on which the operators depend, and $r\cdot$ denotes transport by
$r$.

\emph{Dependence on $\Lambda_t$.}\ Write $\mathsf{C}^{K1}$ for $\mathsf{C}$ at $\mathbb{B}=\mathbb{B}^{K1}$. Every term of
the walk expansion of $\mathsf{C}^{K1}$ that depends on $\Lambda_t$ has a domain containing a
point of $\Lambda_t^{c}$ together with its $\ell/3$-collar.

\emph{The Schur complement.}\ Let $B_{\mathcal{K}}$ be the constant of the corollary on the Schur complement.
The Schur complement
$\mathcal{K}=\mathsf{Q}_{\mathbb{B}'\mathbb{B}'}-\mathsf{J}^{T}\mathsf{C}\mathsf{J}$ satisfies
\begin{equation*}
\mathcal{K}\in\mathcal{S}_{u}[M/4;\,B_{\mathcal{K}},\,\delta_0/2;\,\varepsilon_1],
\end{equation*}
and $\gamma_0/2\le\mathcal{K}\le\gamma_1$ at every $(U,0)$ with $U\in\mathcal{U}_{\mathbb{R}}$.
\end{theorem}

Wherever the walk-expansion properties $(\mathrm{W}1)$--$(\mathrm{W}6)$ are read, in particular in
clauses (E)--(H) and (K), they are read with
the distance $\mathsf{d}(\omega,b,b')=\tau(\operatorname{dom}\nu,b,b')$ computed on the enlarged
domain lists, that is, with the reading hulls $\square^{\mathsf{r}}$ in place of the hulls
$\widehat{\square}$. Enlarging a list lowers $\tau$; hence every upper bound of the form
$e^{-\delta\mathsf{d}}$ obtained with the smaller lists persists with the enlarged ones. The only
lower bounds on $\tau$ that are used are the lower bound on $\tau$ that holds for every list, and
the lower bound that the inverted sheet $(x+\mathsf{A})^{-1}$ of clause (D) carries at its own
radius $M/4$. The packing estimate uses in addition only the
inequality $\rho_{\mathsf{Q}}\le M/4$, which follows from the floor $M\ge144d\rho_{\mathsf{Q}}$
of \eqref{4.14:eq-block-size-conditions-a}.

The proof of Theorem~\ref{4.14:thm-operator-theorem} occupies the remainder of this section; it
rests on four lemmas, which are proved first.

\begin{lemma}[Reading hulls of second-cover boxes]\label{4.14:lem-reading-hulls}
Let $\mathcal{D}^{\sharp}=\bigcup_j\mathcal{D}^{\sharp}_j$ be the second cover of
Definition~\ref{4.14:not-free-current-notation}. A box $\square\in\mathcal{D}^{\sharp}_j$ has
side $2s_j$, where $s_j=M_1^{\sharp}\ell_j$, and is centred at a point of the $s_j$-lattice
lying in the territory $X_j$. For $n\ge 1$, $\square^{n}$ is the concentric cube of half-side
$(1+n)s_j$, and $\widehat{\square}=\square^{3}$. The profiles are $g_{\square}$ and
\[
h_{\square}=g_{\square}\Bigl(\sum_{\square'\in\mathcal{D}^{\sharp}}g_{\square'}^{2}\Bigr)^{-1/2},
\]
with $h_{\square}(c)=h_{\square}(x_c)$ and $\mathfrak{X}(\square'')=\{c:\ x_c\in\square^{1}\}$.
Put
\begin{equation}\label{4.14:eq-reading-hulls-1}
\square^{\mathsf{r}}:=\square^{2L+1}\quad(\text{half-side }(2L+2)s_j),
\qquad
\rho_T:=8dL^{2}M_1^{\sharp}.
\end{equation}
We say that an object is \emph{determined by the nest near} a set $S$ if it is determined by
the nest on an arbitrarily small neighbourhood of $S$. In other words, two nests that coincide
on some open neighbourhood of $S$ give the same object.
\begin{enumerate}
\item[(i)] The restriction $h_{\square}|_{\square^{1}}$ is determined by the nest near
$\square^{\mathsf{r}}$. Hence so are the values $h_{\square}(c)$ for
$c\in\mathfrak{X}(\square'')$, the set $\mathfrak{X}(\square'')$ and the projection $1''$.
\item[(ii)] $\square^{\mathsf{r}}$ meets no territory other than $X_{j-1},X_j,X_{j+1}$, and at
most two of these, two only if they are adjacent. The $\tau$-diameter of
$\square^{\mathsf{r}}$, that is, the largest weighted distance $\tau$ between two of its
points, is at most $\rho_T$.
\item[(iii)] In clauses (iv)--(v) of the inversion lemma for one-box words, replace
$\widehat{\square}$ by $\square^{\mathsf{r}}$ in the domain lists $\operatorname{dom}\sigma$
of the words and $\operatorname{dom}\varpi$ of the steps, and nowhere else. Then:
\begin{itemize}
\item every line of that lemma expressing $(\mathrm{Sh}1)$--$(\mathrm{Sh}4)$ holds with the
same constants;
\item the admissibility conditions of the words are unchanged, since they read
$\widehat{\square}$;
\item $(\mathrm{Sh}5)$ holds factor by factor.
\end{itemize}
Hence $(\mathrm{Sh}5)$ holds for $\widetilde{T}^{-1}$ and for $T^{-1}\in\mathfrak{L}$, where $T$
denotes a sheet, $\widetilde{T}^{-1}$ the inverted sheet and $\mathfrak{L}$ the class of
operators to which $T^{-1}$ is shown to belong. It also
holds for the operators of the lemma on the form $\mathsf{Q}$, one operator at a time, and for
the core lemma, term by term. Consequently the geometric locality statement contained in
clauses (E)--(G) of the operator theorem at free current
(Theorem~\ref{4.14:thm-operator-theorem}) holds with the support $[\nu]$ listing the hulls
$\square^{\mathsf{r}}$.
\item[(iv)] Consider the first cover, that is, Ba{\l}aban's cover $\mathcal{D}_j$ of
\cite{B-CMP99b} at block size $M_1$, with its normalised profiles $\tilde h_{\square}$. For a
cube $\square$ of $\mathcal{D}_j$, of half-side $M_1\ell_j$, let $\square^{n}$ denote the
concentric cube of half-side $(1+n)M_1\ell_j$. Then:
\begin{itemize}
\item $\tilde h_{\square}$ is determined by the nest near $\square^{2L}$;
\item taking
\[
\widehat{X}:=\bigcup_{\square\subset X}\square^{\max\{5,2L\}}
\]
as the listed domain of the elementary walk term attached to the elementary domain $X$ restores
$(\mathrm{Sh}5)$ for the
operators of the elementary sheets below the top layer (the top layer is the subject of the
top-layer proposition for the first cover);
\item the $\tau$-diameters of these enlarged domains are at most $\rho_0:=c_X(d)L^{2}M_1$,
with the constant $c_X(d)$ of \eqref{4.14:eq-free-current-notation-a}.
\end{itemize}
\item[(v)] Set $\rho_{\mathsf{Q}}:=\rho_T$. The conditions (fl-M)$'$(iii) and
(fl-$\sharp$)$'$ of \eqref{4.14:eq-block-size-conditions-a} are read with these values of
$\rho_{\mathsf{Q}}$ and $\rho_0$. They remain floors, that is, one-sided lower bounds, only.
\end{enumerate}
\end{lemma}

\begin{proof}
The proof uses two groups of conditions.
\begin{itemize}
\item The admissibility conditions (A1) and (A2) of
Proposition~\ref{4.14:prop-prepared-data-admissible}, displayed in
\eqref{4.14:eq-prepared-data-admissible-a}. That proposition is printed without proof in
\cite[pp.~178--179]{B-CMP122a} and cited as printed.
\item The floors (fl-M)$'$(i), (fl-M)$'$(iii) and (fl-$\sharp$)$'$ of
Definition~\ref{4.14:def-block-size-conditions}, displayed in
\eqref{4.14:eq-block-size-conditions-a}.
\end{itemize}
By (A1), for $1\le i\le k-1$,
$\operatorname{dist}_\infty(\Omega_i^{c},\Omega_{i+1})\ge M\ell_i$; all indices $i$ used below
lie in this range. By (fl-M)$'$(i),
$M\ge 13L^{2}M_1^{\sharp}$. Combining the two, and using $s_{i+1}=M_1^{\sharp}\ell_{i+1}$ and
$\ell_{i+1}=L\ell_i$,
\begin{equation}\label{4.14:eq-reading-hulls-2}
\operatorname{dist}_\infty(\Omega_i^{c},\Omega_{i+1})\ \ge\ M\ell_i\ \ge\
13L^{2}M_1^{\sharp}\ell_i\ =\ 13L\,s_{i+1}.
\end{equation}

\emph{(i)} Fix $\square\in\mathcal{D}^{\sharp}_j$ and let $\square'\in\mathcal{D}^{\sharp}_{j'}$.
The profile $g_{\square'}$ is not identically zero on $\square^{1}$ (the cube of half-side
$2s_j$) exactly when $\square'^{1}\cap\square^{1}\neq\emptyset$. In that case
\[
|\operatorname{centre}(\square')-\operatorname{centre}(\square)|_\infty\le 2s_j+2s_{j'},
\]
with $\operatorname{centre}(\square')\in X_{j'}$ and $\operatorname{centre}(\square)\in X_j$.

We first exclude the levels $j'\notin\{j-1,j,j+1\}$.

Let $j'\ge j+2$. Since the sets $\Omega_i$ decrease, their complements increase, so
\[
X_j\subset\operatorname{closure}(\Omega^{c}_{j+1})\subset
\operatorname{closure}(\Omega^{c}_{j'-1}),
\qquad X_{j'}\subset\Omega_{j'}.
\]
A segment from $X_j$ to $X_{j'}$ therefore crosses the pair $(j'-1,j')$: one end-point lies in
$\operatorname{closure}(\Omega^{c}_{j'-1})$ and the other in $\Omega_{j'}$. By
\eqref{4.14:eq-reading-hulls-2} with $i=j'-1$, and since $s_j\le s_{j'}$,
\begin{align*}
|\operatorname{centre}(\square')-\operatorname{centre}(\square)|_\infty
&\ge\operatorname{dist}_\infty(\Omega^{c}_{j'-1},\Omega_{j'})
\ge M\ell_{j'-1}\ge 13L^{2}M_1^{\sharp}\ell_{j'-1}\\
&=13L\,s_{j'}>2s_{j'}+2s_j .
\end{align*}

Let $j'\le j-2$. Symmetrically,
\[
X_{j'}\subset\operatorname{closure}(\Omega^{c}_{j'+1})\subset
\operatorname{closure}(\Omega^{c}_{j-1}),
\qquad X_j\subset\Omega_j .
\]
Since $s_{j'}\le s_j$,
\[
|\operatorname{centre}(\square')-\operatorname{centre}(\square)|_\infty
\ge M\ell_{j-1}\ge 13L\,s_j>2s_j+2s_{j'} .
\]

In both cases this contradicts the bound above. So $j'\in\{j-1,j,j+1\}$, and then
$s_{j'}\le Ls_j$. Consequently
\[
|\operatorname{centre}(\square')-\operatorname{centre}(\square)|_\infty\le(2+2L)s_j,
\]
that is, $\operatorname{centre}(\square')\in\square^{2L+1}=\square^{\mathsf{r}}$.

Whether a point $y$ of the $s_{j'}$-lattice is a centre of a box of $\mathcal{D}^{\sharp}_{j'}$
is the statement $y\in X_{j'}$. Since each territory is defined from the sets of the nest alone,
by closure and set difference (Definition~\ref{4.14:not-free-current-notation}), this is decided
by the nest
on any neighbourhood of $y$. The geometry-locality property $(\mathrm{Sh}5)$ is formulated with
an open neighbourhood of the support $[\omega]$. An open neighbourhood of $\square^{\mathsf{r}}$
is therefore also a neighbourhood of every $y\in\partial\square^{\mathsf{r}}$, which settles the
centres on the boundary of $\square^{\mathsf{r}}$. The profiles $g$ themselves are fixed
functions.

It follows that the restriction to $\square^{1}$ of $\sum_{\square'}g_{\square'}^{2}$, and hence
$h_{\square}|_{\square^{1}}$, is a function of the nest near $\square^{\mathsf{r}}$. The
remaining objects follow:
\begin{itemize}
\item The condition $x_c\in\square^{1}$ decides whether $c\in\mathfrak{X}(\square'')$.
\item $1''$ projects on $\mathfrak{X}_1\cap\mathfrak{X}(\square'')$, where $\mathfrak{X}_1$ is
the index set entering the definition of $1''$.
\item Off $\square^{1}$, $h_{\square}$ is $0$ whatever the normalisation.
\item The per-box localisation
\[
T^{\square}=1''\sum\{T_\omega:\ [\omega]\subset\widehat{\square}\}\,1''
\]
and its index set $\{\omega:\ [\omega]\subset\widehat{\square}\}$ are determined by the nest
near $\widehat{\square}$, by the property $(\mathrm{Sh}5)$ of the sheet $T$ itself; and
$\widehat{\square}\subset\square^{\mathsf{r}}$.
\end{itemize}
This proves (i).

\emph{(ii)} We have $\operatorname{centre}(\square)\in X_j\subset
\Omega_j\cap\operatorname{closure}(\Omega^{c}_{j+1})$, and
$\operatorname{diam}_\infty(\square^{\mathsf{r}})=(4L+4)s_j$. We show that $\square^{\mathsf{r}}$
meets no territory of level outside $\{j-1,j,j+1\}$ and cannot meet both $X_{j-1}$ and
$X_{j+1}$.
\begin{itemize}
\item Reaching $X_{j+2}\subset\Omega_{j+2}$ costs at least
\[
\operatorname{dist}_\infty(\Omega^{c}_{j+1},\Omega_{j+2})\ge M\ell_{j+1}\ge 13L^{3}s_j ,
\]
by \eqref{4.14:eq-reading-hulls-2}. The same bound applies to every $X_n$ with $n\ge j+2$, since
$X_n\subset\Omega_n\subset\Omega_{j+2}$.
\item Reaching $X_{j-2}\subset\operatorname{closure}(\Omega^{c}_{j-1})$ costs at least
$M\ell_{j-1}\ge 13Ls_j$. The same bound applies to every $X_n$ with $n\le j-2$, since
\[
X_n\subset\operatorname{closure}(\Omega^{c}_{n+1})\subset\operatorname{closure}(\Omega^{c}_{j-1}).
\]
\item Meeting both $X_{j-1}\subset\operatorname{closure}(\Omega^{c}_j)$ and
$X_{j+1}\subset\Omega_{j+1}$ requires spanning the layer $j$, whose width is
\[
\operatorname{dist}_\infty(\Omega^{c}_j,\Omega_{j+1})\ge M\ell_j\ge 13L^{2}s_j .
\]
\end{itemize}
Each of these three quantities exceeds $(4L+4)s_j$, because $13L>4L+4$. Hence
$\square^{\mathsf{r}}$ meets no territory of level outside $\{j-1,j,j+1\}$, and at most two
adjacent ones among $X_{j-1},X_j,X_{j+1}$.

For the diameter, join two points of $\square^{\mathsf{r}}$ by a segment. Its Euclidean length is
at most $\sqrt{d}\,(4L+4)s_j$. The segment stays in $\square^{\mathsf{r}}$, hence in the
territories of levels $j-1,j,j+1$, where lengths are measured in units at least
$\ell_{j-1}=\ell_j/L$. Since $s_j/\ell_{j-1}=LM_1^{\sharp}$, its weighted length, which bounds
$\tau$ between its end-points, satisfies
\[
\sqrt{d}\,(4L+4)s_j\,\ell_{j-1}^{-1}=\sqrt{d}\,(4L+4)LM_1^{\sharp}\le 8dL^{2}M_1^{\sharp}=\rho_T .
\]
This proves (ii).

\emph{(iii)} We have $\square^{\mathsf{r}}\supset\widehat{\square}$. The weighted distances $\tau$
and $\tau^{pr}$ do not increase when a listed set is enlarged.

Every line of clauses (iv)--(v) of the inversion lemma for one-box words that expresses the
decay property $(\mathrm{Sh}3)$ was proved with
\[
\mathsf{d}^{\mathrm{old}}=\tau(\text{old list},\cdot,\cdot)\ \ge\ \mathsf{d}^{\mathrm{new}},
\]
where $\mathsf{d}^{\mathrm{new}}$ is the same distance for the enlarged lists. Hence
$e^{-\delta\mathsf{d}^{\mathrm{old}}}\le e^{-\delta\mathsf{d}^{\mathrm{new}}}$, and these lines
hold verbatim.

The supports $[\sigma]$ and $[\varpi]$ stay connected, because
$\square^{\mathsf{r}}\supset\widehat{\square}$ and $\widehat{\square}$ contains the sets that made
them connected. They still contain the indices. The reading property $(\mathrm{Sh}2)$ and the
covariance property $(\mathrm{Sh}4)$ hold factor by factor as before. The hull estimate of the
weighted distance is stated for an arbitrary hull $K$, so it applies with the hulls
$\square^{\mathsf{r}}$. The admissibility conditions of the words read $\widehat{\square}$ and are
not touched by the replacement.

We turn to $(\mathrm{Sh}5)$. A factor of a word reads the point $\mathsf{u}$ on
$\widehat{\square}$, by $(\mathrm{Sh}2)$. It reads the geometry through $h_{\square}$,
$\mathfrak{X}(\square'')$, $1''$ and the admissibility set of $T^{\square}$. By (i), all of these
are determined by the nest near $\square^{\mathsf{r}}$. Therefore two nests that coincide on an
open neighbourhood of $[\varpi]\supseteq\bigcup_i\square_i^{\mathsf{r}}$ give the same factor.
This is $(\mathrm{Sh}5)$ factor by factor, and hence $(\mathrm{Sh}5)$ for $\widetilde{T}^{-1}$, for
$T^{-1}\in\mathfrak{L}$, for the operators of the lemma on the form $\mathsf{Q}$ one at a time,
and for the terms of the core lemma. This yields the geometric locality statement in clauses
(E)--(G) of Theorem~\ref{4.14:thm-operator-theorem} with $[\nu]$ listing the hulls
$\square^{\mathsf{r}}$.

The hull $\widehat{\square}$ alone would not suffice. Consider the nests $\{x_1\ge 0\}$ and
$\{x_1\ge LM\ell_j\}$, and a box $\square$ centred at $x_1=-5s_j$. The hull
$\widehat{\square}=\square^{3}$ reaches only $x_1=-s_j<0$. The two nests nevertheless differ inside
$\square^{2L+1}$, which reaches $x_1=(2L-3)s_j>0$. So the list $\square^{\mathsf{r}}$ separates
them, while the list $\widehat{\square}$ does not.

It remains to control the lower bounds on the distance that these lines use. They are two: the
lower bound on the weighted distance valid for any list, and the lower bound (b3) of the inverted
sheet at its own radius, one of the box conditions (b1)--(b4) of the one-box words, which the
conventions on reading hulls \eqref{4.14:eq-free-current-notation-a} leave unchanged.
\begin{itemize}
\item In the many-scale frame, the inversion lemma reads first-cover sets at radius $\rho_0$
against (b3). This requires $M_1^{\sharp}>2L\rho_0$, which is part of (fl-$\sharp$)$'$, now
taken at the new value of $\rho_0$. The constants $B_T$, $\delta_T$ of the inverted sheet
involve $\rho_0$ and not $\rho_T$, so the condition is not circular.
\item In the unit frame, the inversion lemma reads the lists of $\mathsf{Q}$ at radius
$\rho_{\mathsf{Q}}=8dL^{2}M_1^{\sharp}$. With $\ell=M/12d$, condition (fl-M)$'$(iii), namely
$M\ge 144d\rho_{\mathsf{Q}}$, gives $\rho_{\mathsf{Q}}\le\ell/12$. Since
\[
\rho_{\mathsf{Q}}\le\frac{M}{144d}\le\frac{M}{4},
\]
the condition $\rho\le M/4$ of the core lemma holds.
\end{itemize}
This proves (iii).

\emph{(iv)} In the cover of \cite{B-CMP99b}, let $\square'\in\mathcal{D}_{j'}$ be a cube with
$h_{\square'}\not\equiv 0$ on $\square$. By the same exclusion as in (i),
$j'\in\{j-1,j,j+1\}$. By \cite[(1.118)]{B-CMP95}, $\operatorname{supp}h_{\square'}\subset\square'$,
and $\square'$ has half-side at most $LM_1\ell_j$. So $\square'$ lies within $(2L+1)M_1\ell_j$ of
the centre of $\square$.

Membership $\square'\in\mathcal{D}_{j'}$ is read on $\square'$ itself. Indeed, in
\cite{B-CMP99b} the cubes of the cover at level $j$ are unions of $2^{d}$ big blocks, at least one
of which is contained in the blocked level-$j$ region of \cite{B-CMP99b}. Hence $\tilde h_{\square}$
is determined by the nest
near $\square^{2L}$. The tower of domains attached to $\square$ and the commutator of
$\tilde h_{\square}$ with the block-averaging term of the form read $\square^{5}$, by the reading
statements for the towers
and for the normalised profiles of the first cover.

Listing $\square^{\max\{5,2L\}}$ keeps $(\mathrm{Sh}1)$--$(\mathrm{Sh}4)$, by the monotonicity of
$\tau$ and $\tau^{pr}$ under enlargement of a listed set, used as in (iii). It also gives
$(\mathrm{Sh}5)$, since every object the term reads is determined by the nest near this set.

The $\tau$-diameter of $\widehat{X}$ is at most $c_X(d)L^{2}M_1$. Indeed, the
$\ell_\infty$-diameter of $\widehat{X}$ is at most the bound of \cite{B-CMP99b} on the number of
cubes of an elementary domain times the side $2(2L+1)M_1\ell_j$ of the cubes
$\square^{\max\{5,2L\}}$; passing to Euclidean length costs the factor $\sqrt{d}$, and measuring
in units at least $\ell_{j-1}=\ell_j/L$, as in (ii), costs the factor $L$. Since
$2(2L+1)\le 6L$, the result is at most $c_X(d)L^{2}M_1$, where $c_X(d)$ is $6\sqrt{d}$ times
that bound on the number of cubes; in particular the $\ell_\infty$-diameter of $\widehat{X}$ is
at most $c_X(d)LM_1\ell_j$.

The enlarged $\widehat{X}$ meets at most two adjacent layers. The argument is that of (ii): the
layer widths are the same, at least $13LM_1^{\sharp}\ell_j$ by \eqref{4.14:eq-reading-hulls-2},
and they exceed the $\ell_\infty$-diameter $c_X(d)LM_1\ell_j$ of $\widehat{X}$, since
(fl-$\sharp$)$'$ gives $M_1^{\sharp}>2L\rho_0=2c_X(d)L^{3}M_1$. This proves (iv).

\emph{(v)} The radii $\rho_T$, $\rho_{\mathsf{Q}}$ and $\rho_0$ enter only the following places:
\begin{itemize}
\item condition (fl-$\sharp$)$'$, through $M_1^{\sharp}>2L\rho_0$ and through the floor conditions
$\mathsf{F}_{ms}$, $\mathsf{F}^{P}_{ms}$ of the many-scale format $\mathcal{S}_{ms}$;
\item the box conditions (b3)--(b4) of the inverted sheet, through the $\tau$-diameter of
$\widehat{\square}$, which by the computation of (ii) applied to the side $8s_j$ is at most
$8\sqrt{d}\,LM_1^{\sharp}\le\rho_T$;
\item condition (fl-M)$'$(iii), namely $M\ge 144d\rho_{\mathsf{Q}}$;
\item the condition $\rho\le M/4$ of the core lemma, that is,
$8dL^{2}M_1^{\sharp}\le M/4$, which follows from (fl-M)$'$(iii).
\end{itemize}
These are floors on $M_1^{\sharp}$, imposed after $M_1$, and floors on $M$, imposed after
$M_1^{\sharp}$. This proves (v).
\end{proof}

\begin{remark}[No reading hull in the unit frame]
In the unit frame no hull is needed. The grid and the partition of unity $\{\theta\}$ of that
frame do not depend on the nest, and $\sum\theta^{2}\equiv 1$ requires no normalisation. Moreover
$\widehat{\square}\subset\Omega_k$, by (A2) of
Proposition~\ref{4.14:prop-prepared-data-admissible}, printed without proof in
\cite[pp.~178--179]{B-CMP122a} and cited as printed.
\end{remark}

We use the notation of Definition~\ref{4.14:not-free-current-notation}, together with the following conventions.

\emph{Bonds, parities and pairings.} Bonds are oriented, $b=(b_-,b_+)$. The reversed bond is
$\bar b:=(b_+,b_-)$, and $\mathbf{U}(\bar b)=\mathbf{U}(b)^{-1}$. For an unoriented bond $\iota$ we write
$\partial\iota$ for the set of its two end-points, and for a block $z$ we write $\Delta(z)$ for its set
of sites; this is not to be confused with the Laplacians $\Delta_U$, $\Delta^{(2)}$, $\Delta'_a$ below.
$R$ is the representation through which $\mathbf{U}(b)$ acts on the values of site fields; the letter
$R(\cdot)$ here denotes this representation and not the ball radius $R$ of the observables.
$\eta$ is the lattice spacing of the bond sums. A bond field $F$ is called \emph{even}
(\emph{odd}) if $F(\bar b)=+R(\mathbf{U}(b))^{-1}F(b)$ (respectively $F(\bar b)=-R(\mathbf{U}(b))^{-1}F(b)$)
for every $b$. On the Lie-algebra-valued arguments of the trace pairings below, the same
transformation acts by conjugation, $X\mapsto R(\mathbf{U}(b))^{-1}XR(\mathbf{U}(b))$, which we write
$\operatorname{Ad}R(\mathbf{U}(b))^{-1}X$. When the site field is Lie-algebra-valued, as
$\mathsf{w}A$ is, the factor $R(\mathbf{U}(b))^{-1}$ in the parity convention and in
\eqref{4.14:eq-symmetry-covariance-4} acts as this conjugation. The pairing with the multiplication
vertex is
\begin{equation*}
\langle X,M_{\mathsf{f}}Y\rangle=\sum_b\eta^d\operatorname{tr}\bigl(\mathsf{f}(b)\cdot i[X(b),Y(b)]\bigr).
\end{equation*}
We write $\mathsf{w}$ for the map of \eqref{4.14:eq-free-current-notation-c} from bond fields to site
fields that enters the operator $W_\pi$.

\emph{Symmetries.} $\mathfrak{r}$ is the group of lattice symmetries composed of the $2^d$ reflections and
the $d!$ permutations of the coordinate axes. It acts on bond spaces by $\mathcal{R}_r$, $r\in\mathfrak{r}$.
The letter $r$ inside the brackets of the format $\mathcal{O}^{\alpha}[\,\cdot\,;\,\cdot\,,\cdot\,;r]$
denotes the analyticity radius and not a symmetry.

\begin{lemma}[Symmetric cell weights and lattice-symmetry covariance]
\label{4.14:lem-symmetry-covariance}
We use the notation of Definition~\ref{4.14:not-free-current-notation} and the conventions above.
Then the following hold.
\begin{enumerate}
\item[(a)] Let $N$ be an operator as in the local-term lemma of the operator construction: local
with range $c_N$, with
$\|\chi_yN\|\le\nu\ell_y^{\alpha}$ for every block $y$ and with analyticity radius $r$; let
$\delta>0$. Here $\chi_y$ are the junction indicators
of the product construction, and $\nu$ is a constant; in clause~(d) the letter $\nu$ denotes an
index of a term, not this constant. Put $N_z:=w_zN$, with $w_z$ as in
\eqref{4.14:eq-free-current-notation-b}, and put $\operatorname{dom}z:=\{X_N^{+}(z)\}$, where
$X_N^{+}(z)$ is $X_N(z)$ at radius $c_N+dL$. Then the local-term lemma of the operator construction
holds in the form
\begin{equation}\label{4.14:eq-symmetry-covariance-1}
N\in\mathcal{O}^{\alpha}\bigl[\rho_N^{+};\,n_w\nu e^{\delta(c_N+dL)},\,\delta;\,r\bigr],
\qquad \rho_N^{+}:=2(c_N+dL)+2d,
\end{equation}
where
\begin{equation}\label{4.14:eq-symmetry-covariance-2}
n_w(d,L):=\sup_y\#\{z\in\mathfrak{B}_s:\ d(y,z)\le dL\}\le 3(2dL+1)^d .
\end{equation}
The branching lemma holds with $n_N$ replaced by $n_Nn_w$ and $c_N$ replaced by $c_N+dL$. The
family $(N_z)$ has the sheet properties $(\mathrm{Sh}1)$, $(\mathrm{Sh}2)$, $(\mathrm{Sh}4)$ and
$(\mathrm{Sh}5)$; here $(\mathrm{Sh}1)$--$(\mathrm{Sh}5)$ denote the five properties of the sheet
format of the operator construction: in this order, the conditions on domains, on reading, on
decay, on covariance and on geometry-locality. Finally, $w_{rz}(r\iota)=w_z(\iota)$ for every
symmetry $r\in\mathfrak{r}$, every block $z$ and every bond $\iota$.
\item[(b)] For every site field $\lambda$ and every $\mathsf{f}$,
\begin{equation}\label{4.14:eq-symmetry-covariance-3}
\langle\mathfrak{S}_{-}\lambda,M_{\mathsf{f}}D\lambda\rangle
=\tfrac12\langle\mathfrak{S}\lambda,M_{\mathsf{f}}D\lambda\rangle .
\end{equation}
More is true. Consider the two bilinear forms
\begin{equation*}
(A,A')\mapsto\langle\mathfrak{S}_{-}\mathsf{w}A,M_JD\mathsf{w}A'\rangle
\quad\text{and}\quad
(A,A')\mapsto\tfrac12\langle\mathfrak{S}\mathsf{w}A,M_JD\mathsf{w}A'\rangle .
\end{equation*}
They differ by an antisymmetric form, so their symmetrisations coincide. Consequently the operator
$W_\pi$ of \eqref{4.14:eq-free-current-notation-c} coincides with the symmetric operator of
\begin{equation*}
A\mapsto\langle A,M_J\mathfrak{S}\mathsf{w}A\rangle
+\langle\mathfrak{S}_{-}\mathsf{w}A,M_JD\mathsf{w}A\rangle ,
\end{equation*}
that is, with the operator $W_\pi$ built with the letter $\mathfrak{S}_-$. The two have the same
matrix, and their words after polarisation are the same up to the replacement of the letter
$\mathfrak{S}_-$ by $\tfrac12\mathfrak{S}$; orders and weights are unmoved.
\item[(c)] Under bond reversal the letters have the parities
\begin{equation}\label{4.14:eq-symmetry-covariance-4}
(D\lambda)(\bar b)=-R(\mathbf{U}(b))^{-1}(D\lambda)(b),\qquad
(\mathfrak{S}\lambda)(\bar b)=+R(\mathbf{U}(b))^{-1}(\mathfrak{S}\lambda)(b).
\end{equation}
Thus $D\lambda$ is odd and $\mathfrak{S}\lambda$ is even. Take $A$ and $J$ odd, which is the
convention of \cite{B-CMP119} for the components on reflected bonds. Then both trilinear summands
of $W_\pi$,
\begin{equation}\label{4.14:eq-symmetry-covariance-5}
\operatorname{tr}\bigl(J(b)\cdot i[A(b),(\mathfrak{S}\mathsf{w}A)(b)]\bigr),\qquad
\operatorname{tr}\bigl(J(b)\cdot i[(\mathfrak{S}\mathsf{w}A)(b),(D\mathsf{w}A)(b)]\bigr),
\end{equation}
are independent of the orientation of $b$. Hence the quadratic form that defines $W_\pi$ is natural
under $\mathfrak{r}$.
\item[(d)] Let $r\in\mathfrak{r}$. Then every index set of the construction admits a bijection
$\nu\mapsto r\nu$. The index sets meant are: the tuples of the product construction; the indices
$z$ of the local-term lemma and of the branching lemma; the words $\varpi$ of the inversion
construction,
with the hulls of Lemma~\ref{4.14:lem-reading-hulls}; the core classes; and the words of the
word lemma for the form. Write $T_\nu[\mathfrak{d}](\mathsf{u})$ for the term with index $\nu$,
built from the data $\mathfrak{d}$ and evaluated at $\mathsf{u}=(\mathbf{U},J)$. These bijections
satisfy $[r\nu]=r[\nu]$, $\operatorname{dom}(r\nu)=r(\operatorname{dom}\nu)$, and
\begin{equation}\label{4.14:eq-symmetry-covariance-6}
T_{r\nu}[r\mathfrak{d}](r\mathsf{u})
=\mathcal{R}_r\,T_\nu[\mathfrak{d}](\mathsf{u})\,\mathcal{R}_r^{-1},
\qquad T\in\{\mathsf{C},\ \mathsf{J},\ (x+\mathsf{A})^{-1},\ \mathcal{K}\}.
\end{equation}
Here $x$ is a scalar shift parameter, not to be confused with the kernel arguments below, and
$(x+\mathsf{A})^{-1}$ is the inverse of the shifted operator $x+\mathsf{A}$. This clause rests on
the covariance of the elementary letters under $\mathfrak{r}$, in the following two senses.
\begin{itemize}
\item The averaging scheme used in this chapter is covariant,
$(r\cdot\mathcal{W})\langle ry,rx\rangle=\mathcal{W}\langle y,x\rangle$. Here $\mathcal{W}$ stands for
any of the averaging letters $M(\cdot)$, $\mathfrak{g}$, $Q^{\flat}$, $\mathfrak{a}$, $C^{(2)}$, $h$,
$C_{\flat}$, regarded as kernels.
\item Cubes, cut-offs, Laplacians, $D$, $\mathfrak{S}$, $M_{\mathsf{f}}$,
and the walks of \cite[\S C]{B-CMP99b}, which are built from $\Delta'_a$, $Q'$ and the operator
denoted $R$ there, are manifestly covariant.
\end{itemize}
\end{enumerate}
Clause \textup{(d)} is the covariance clause of the operator theorem at free current,
Theorem~\ref{4.14:thm-operator-theorem}.
\end{lemma}

\begin{proof}
\emph{Clause \textup{(a)}: the distance from the bookkeeping block.} Let $\iota$ be a bond with
$w_z(\iota)\neq0$. The numerator of $w_z(\iota)$ in \eqref{4.14:eq-free-current-notation-b} counts
the end-points of $\iota$ in $\Delta(z)$, so an end-point of $\iota$ lies in $\Delta(z)$. The
bookkeeping block $y(\iota)$ of $\iota$ (the block of its anchor, as in the connector lemma) also
holds an end-point of $\iota$. The two end-point blocks of a bond are joined by the bond $\iota$
itself, an admissible connector of weighted length at most $L$ to the anchor by clause~(i) of the
connector lemma, followed by at most $(d-1)(L-1)/2+1$ grid steps inside an interface face to the
block point.
In this count one uses that a bond of $\mathfrak{B}$ has no end-point in the set $\Omega_{n+1}$ that
follows the layer $n$ in which it lies \cite[(2.3)]{B-CMP96}. The total weighted length of this
path is at most $dL$, and therefore $d(y(\iota),z)\le dL$.

\emph{Clause \textup{(a)}: the decay property.} By the preceding paragraph,
$\chi_yN_z\chi_{y'}\neq0$ forces $d(y,z)\le dL$. Since $N$ is local with range $c_N$, it also
forces $d(y,y')\le c_N$. Then $d(y',z)\le c_N+dL$, so both $y$ and $y'$ lie in $X_N^{+}(z)$.

Put $\mathsf{d}(z;y,y'):=\tau^{pr}(X_N^{+}(z),\{y\},\{y'\})$. This weighted distance through the
listed set $X_N^{+}(z)$ equals $d(y,y')$, because the geodesic from $y$ starts inside the listed
set; hence $\mathsf{d}(z;y,y')\le c_N$. By \eqref{4.14:eq-free-current-notation-b} we have
$0\le w_z\le1$. This gives the first inequality below; the second is the bound on $N$; the third
uses $\mathsf{d}(z;y,y')\le c_N$:
\begin{equation*}
\|\chi_yN_z\chi_{y'}\|\le\|\chi_yN\|\le\nu\ell_y^{\alpha}
\le\nu e^{\delta c_N}\ell_y^{\alpha}e^{-\delta\mathsf{d}(z;y,y')} .
\end{equation*}
This is $(\mathrm{Sh}3)$ with the coefficients
$a_z(y):=\nu e^{\delta c_N}\mathbf{1}\{d(y,z)\le dL\}$. Their row sums satisfy
\begin{equation}\label{4.14:eq-symmetry-covariance-7}
\sum_za_z(y)\le n_w\nu e^{\delta c_N}\le n_w\nu e^{\delta(c_N+dL)} .
\end{equation}

\emph{Clause \textup{(a)}: the count $n_w$.} The number $n_w$ depends on $(d,L)$ only. A block $z$
with $d(y,z)\le dL$ lies in one of the layers $n(y)-1$, $n(y)$, $n(y)+1$, where $n(y)$ is the index
of the layer containing $y$. Indeed, a path that traverses an entire layer has weighted length at
least $M/L$, and $M/L>dL$ under the condition on $M$ among the one-sided conditions
\eqref{4.14:eq-block-size-conditions-a}; so a path of weighted length at most $dL$ from $y$ cannot
reach beyond the two layers adjacent to the layer $n(y)$. Within each of these layers, $z$ lies
within Euclidean
distance $dL\ell$ of $y$, measured in that layer's own unit, where $\ell$ is the side of the unit
boxes. This leaves at most
$(2dL+1)^d$ blocks per layer, hence at most $3(2dL+1)^d$ in all, which is
\eqref{4.14:eq-symmetry-covariance-2}.

\emph{Clause \textup{(a)}: the remaining sheet properties.}
\begin{itemize}
\item $(\mathrm{Sh}1)$: $\operatorname{dom}z$ consists of the one set $X_N^{+}(z)$, a connected union
of blocks holding $y$ and $y'$.
\item $(\mathrm{Sh}2)$: $N_z$ reads $\mathsf{u}$ on $X_N(y)$, and $X_N(y)\subset X_N^{+}(z)$.
\item $(\mathrm{Sh}4)$: $w_z$ is a scalar.
\item $(\mathrm{Sh}5)$: $X_N(y)$ and $w_z$ are determined by the geometry of the nest.
\end{itemize}
For the branching lemma, its proof goes through with the same two changes: the constant $n_N$
becomes $n_Nn_w$, by \eqref{4.14:eq-symmetry-covariance-7}, and the range $c_N$ becomes
$c_N+dL$, by the radius of $X_N^{+}(z)$.

\emph{Clause \textup{(a)}: naturality of the weights.} A symmetry $r\in\mathfrak{r}$ maps blocks to
blocks, because the blocks are centred and $L$ is odd. It also maps end-points of bonds to
end-points. Therefore
\begin{equation*}
\#\bigl(\partial(r\iota)\cap\Delta(rz)\bigr)=\#\bigl(\partial\iota\cap\Delta(z)\bigr),\qquad
\#\bigl(\partial(r\iota)\cap\Omega_0\bigr)=\#\bigl(\partial\iota\cap\Omega_0\bigr).
\end{equation*}
These are the numerator and the denominator of \eqref{4.14:eq-free-current-notation-b}, so
$w_{rz}(r\iota)=w_z(\iota)$. By contrast, the indicator $\mathbf{1}\{y(r\iota)=rz\}$ of the
bookkeeping block would not have this property for a reflection that reverses $\iota$.

\emph{Clause \textup{(b)}: the scalar identity.} By the definitions of the letters,
\begin{equation*}
(\mathfrak{S}-2\mathfrak{S}_{-})\lambda(b)=R(\mathbf{U}(b))\lambda(b_+)-\lambda(b_-)
=\eta(D\lambda)(b).
\end{equation*}
Inserting this into the pairing gives
\begin{equation}\label{4.14:eq-symmetry-covariance-8}
\langle(\mathfrak{S}-2\mathfrak{S}_{-})\lambda,M_{\mathsf{f}}D\lambda\rangle
=\eta\sum_b\eta^d\operatorname{tr}\bigl(\mathsf{f}(b)\cdot i[(D\lambda)(b),(D\lambda)(b)]\bigr)=0,
\end{equation}
since the commutator is antisymmetric and so vanishes on equal arguments. Dividing
\eqref{4.14:eq-symmetry-covariance-8} by $2$ gives \eqref{4.14:eq-symmetry-covariance-3}.

\emph{Clause \textup{(b)}: the bilinear forms.} By the same identity with $\lambda=\mathsf{w}A$,
\begin{equation*}
\tfrac12\langle\mathfrak{S}\mathsf{w}A,M_JD\mathsf{w}A'\rangle
-\langle\mathfrak{S}_{-}\mathsf{w}A,M_JD\mathsf{w}A'\rangle
=\tfrac12\eta\,\langle D\mathsf{w}A,M_JD\mathsf{w}A'\rangle .
\end{equation*}
The right-hand side is antisymmetric in $(A,A')$, by the antisymmetry of the commutator in the
pairing. It is therefore annihilated by the symmetrisation that defines $W_\pi$, and the two
symmetric operators coincide.

\emph{Clause \textup{(b)}: the words.} In the word lemma for the form the words change only by the
scalar $\tfrac12$ and the name of the letter. The admissible letter triples (left letter, vertex,
right letter) of the word lemma for the form
\begin{equation*}
(I,M_J,\mathfrak{S}\mathsf{w}),\quad(\mathfrak{S}_{-}\mathsf{w},M_J,D\mathsf{w})
\end{equation*}
become
\begin{equation*}
(I,M_J,\mathfrak{S}\mathsf{w}),\quad(\tfrac12\mathfrak{S}\mathsf{w},M_J,D\mathsf{w}).
\end{equation*}
Their orders are $0$, and their weights are at most the old ones. The letter $\mathfrak{S}_-$
occurs nowhere else in the first two of the defining series of $\mathsf{Q}$ in
Theorem~\ref{4.14:thm-operator-theorem}. In particular it does not occur in
\begin{equation*}
W_1=W_\pi+(I-D\mathsf{w})^{T}\Delta^{(2)}(I-D\mathsf{w}).
\end{equation*}

\emph{Clause \textup{(c)}: the parities.} We use $\bar b_+=b_-$, $\bar b_-=b_+$ and
$R(\mathbf{U}(\bar b))=R(\mathbf{U}(b))^{-1}$. For $D$,
\begin{align*}
\eta(D\lambda)(\bar b)&=R(\mathbf{U}(\bar b))\lambda(\bar b_+)-\lambda(\bar b_-)
=R(\mathbf{U}(b))^{-1}\lambda(b_-)-\lambda(b_+)\\
&=-R(\mathbf{U}(b))^{-1}\bigl[R(\mathbf{U}(b))\lambda(b_+)-\lambda(b_-)\bigr]
=-\eta R(\mathbf{U}(b))^{-1}(D\lambda)(b).
\end{align*}
For $\mathfrak{S}$,
\begin{align*}
(\mathfrak{S}\lambda)(\bar b)&=\lambda(\bar b_-)+R(\mathbf{U}(\bar b))\lambda(\bar b_+)
=\lambda(b_+)+R(\mathbf{U}(b))^{-1}\lambda(b_-)\\
&=+R(\mathbf{U}(b))^{-1}\bigl[R(\mathbf{U}(b))\lambda(b_+)+\lambda(b_-)\bigr]
=+R(\mathbf{U}(b))^{-1}(\mathfrak{S}\lambda)(b).
\end{align*}
This is \eqref{4.14:eq-symmetry-covariance-4}.

\emph{Clause \textup{(c)}: orientation independence.} On $\bar b$ each field transforms as
$F\mapsto\varepsilon_F\operatorname{Ad}R(\mathbf{U}(b))^{-1}F$, with
\begin{equation*}
\varepsilon_A=\varepsilon_J=\varepsilon_{D\lambda}=-1,\qquad \varepsilon_{\mathfrak{S}\lambda}=+1 .
\end{equation*}
The common transformation $\operatorname{Ad}R(\mathbf{U}(b))^{-1}$ is a conjugation: it preserves
the commutator and the trace, so it cancels in $\operatorname{tr}\circ[\cdot,\cdot]$. Each summand
$\operatorname{tr}(J\cdot i[X,Y])$ of \eqref{4.14:eq-symmetry-covariance-5} is therefore multiplied
by $\varepsilon_J\varepsilon_X\varepsilon_Y$. This product is $(-)(-)(+)=+$ for
$(X,Y)=(A,\mathfrak{S}\mathsf{w}A)$, and $(-)(+)(-)=+$ for
$(X,Y)=(\mathfrak{S}\mathsf{w}A,D\mathsf{w}A)$. Hence the sum over bonds may be taken over
unoriented bonds; it is well defined and natural under $\mathfrak{r}$.

\emph{Clause \textup{(c)}: the remaining letters.} The vertex $M_{\mathsf{f}}$ with $\mathsf{f}$ odd
exchanges parities. The action $\mathcal{R}_r$ acts on each bond space with the parity of the
corresponding letter, and reads $\mathbf{U}(b)$. For a coarse bond $c$ we write $-c$ for the
reversed coarse bond. Likewise
$\Phi^{\flat}(-c)=-R(M(V)(c))^{-1}\Phi^{\flat}(c)$, where $M(V)$ denotes the averaged
configuration, written with its argument to distinguish it from the block size $M$. So
$Q^{\flat}$ and $C^{(2)}$ are odd, and their pairing with the odd current $\mathfrak{j}$ is
independent of orientation.

\emph{Clause \textup{(d)}: the objects in the definition of a term.} Clause~(d) follows from (a),
(b), (c) and the covariance of the elementary letters stated there. The proof is transport of
structure: the bijection $\nu\mapsto r\nu$ applies $r$ to the geometric objects that label $\nu$. The
argument is valid because every object entering the \emph{definition} of a term is natural under
$\mathfrak{r}$. These objects are the following.
\begin{itemize}
\item The covers and their product cut-offs (even profiles, lattice-centred cubes).
\item The midpoint rules $h_{\square}(x_c)$, the index sets $\mathfrak{X}(\square'')$, and the
concentric hulls $\widehat{\square}$ and $\square^{\mathsf{r}}$ of
Lemma~\ref{4.14:lem-reading-hulls}.
\item The weights $w_z$, by clause~(a).
\item The letters $\mathfrak{S}$, $D$, $M_{\mathsf{f}}$ with their parities, by clause~(c).
\item The bond $b_0(c)$, which is the middle bond of $c$ ($L$ odd; \cite{B-CMP109}).
\item The averaging letters $M(\cdot)$, $Q_j$, $Q^{\flat}$, $\mathfrak{g}$, $Q'$. These are covariant
by the transformation properties of the averaging scheme, including the reversal rule
$\bar U(-c)=\bar U(c)^{-1}$ satisfied by the averaged bond variables $\bar U(c)$ of that scheme.
\item The towers and the unit grid ($M\in\ell\mathbb{Z}$).
\item The lengths $\ell_c$ as they enter the flat construction, the geometric admissibility
conditions, the core classes, and the Neumann words.
\end{itemize}
The following enter the \emph{bounds} only, not the definition of any term: the anchors $a(c)$ and
the parameter $\varrho_{\mathrm{con}}$ of the connector lemma, the bookkeeping block $y(\cdot)$,
$c_\omega(b):=a_\omega(b_-)$, the auxiliary functions of the branching lemma indexed by blocks $y$,
and the row constants. For $y(\cdot)$ this holds because, by clause~(a), the terms $N_z$ are defined
through the weights $w_z$.

\emph{Clause \textup{(d)}: the elementary letters and the induction.} The elementary averaging letters
are covariant by the covariance statement of the averaging scheme used in this chapter,
$(r\cdot\mathcal{W})\langle ry,rx\rangle=\mathcal{W}\langle y,x\rangle$; cubes, cut-offs, Laplacians,
$D$, $\mathfrak{S}$, $M_{\mathsf{f}}$ and the walks of \cite[\S C]{B-CMP99b} are manifestly covariant.
This covariance is not derived
from the contours $\Gamma_{y,x}$ of \cite{B-CMP98}, which are built with one fixed ordering; a
definition of this type is not symmetric with respect to lattice Euclidean transformations, as
noted in \cite{B-CMP109}. For $\Delta_U$, $D$ and the Dirichlet inverses of the towers, covariance
is transport of structure. We now proceed by induction along the construction:
\begin{enumerate}
\item[(1)] the elementary letters;
\item[(2)] the words of the product construction, of the local-term lemma and of the branching
lemma;
\item[(3)] the words of the inversion construction;
\item[(4)] the core classes;
\item[(5)] the Schur complement $\mathcal{K}$ of the Schur-complement corollary.
\end{enumerate}
At each stage the terms indexed by $r\nu$ are built by the same operations from the images under
$r$ of the objects indexing $\nu$. Hence $[r\nu]=r[\nu]$ and
$\operatorname{dom}(r\nu)=r(\operatorname{dom}\nu)$, and \eqref{4.14:eq-symmetry-covariance-6}
holds for $\mathsf{C}$, $\mathsf{J}$, $(x+\mathsf{A})^{-1}$ and $\mathcal{K}$. This is the covariance
clause of Theorem~\ref{4.14:thm-operator-theorem}. The only published statement on which this
covariance rests is the convention of \cite{B-CMP119} for the components on reflected bonds.
\end{proof}

\begin{lemma}[Local balls at cut-off radius $M/2L$]\label{4.14:lem-local-balls}
Let $T\in\mathcal{S}_{u}[\rho;B,\delta;r]$, or a core form of such a $T$ (its form relative to a
core $\mathsf{Z}_0\subset\mathfrak{c}^{\flat}$, a set of level-$k$ $M$-cubes), and let $\nu$ be an
index, with support $[\nu]$. Let $U\in\mathcal{U}_{\mathbb{R}}$ be given on the whole lattice, and let
$(A',J)$ be given on
\begin{equation*}
[\nu]^{\approx}:=\bigl\{x:\ \tau(\{x\},[\nu])<M/2L\bigr\},
\end{equation*}
with every member of $\|(A',J)\|_{\star}$ smaller than $r/2$ there. The set of such data $(A',J)$
is called the \emph{local ball} (of radius $r/2$ about $[\nu]$). Let $[\nu]^{\sim}$ be the
one-cube collar of $[\nu]$, that is, $[\nu]$ together with the cubes of $\mathfrak{c}^{\flat}$
that touch a cube of $[\nu]$. Put
\begin{equation}\label{4.14:eq-local-balls-1}
\chi(x):=\max\Bigl\{0,\ 1-\frac{2L}{M}\,\tau(\{x\},[\nu])\Bigr\}.
\end{equation}
Then the following hold.
\begin{itemize}
\item[$(\alpha)$] $\{\chi>0\}\subset[\nu]^{\approx}\subset[\nu]^{\sim}$.
\item[$(\beta)$] $(e^{i\eta\chi A'}U,\ \chi J)\in\mathcal{U}(r)$.
\item[$(\gamma)$] The value of $T_\nu$ at the point of $(\beta)$ is independent of the extension of
the data, and it defines $T_\nu$ on the local ball; so defined, $T_\nu$ is analytic and has the
properties $(\mathrm{W}3)'$, $(\mathrm{W}5)^{\mathrm{cov}}$, $(\mathrm{W}6)^{\mathrm{geo}}$ listed in
clause~(E) of the operator theorem at free current (Theorem~\ref{4.14:thm-operator-theorem},
\eqref{4.14:eq-operator-theorem-a}).
\end{itemize}
If the cut-off is placed at weighted distance $M/2$ instead of $M/2L$, the argument for $(\alpha)$
no longer applies at a seam with a coarser side, and the inclusion $(\alpha)$ is no longer
guaranteed.
\end{lemma}

\begin{proof}
\emph{$(\alpha)$.} The first inclusion holds by \eqref{4.14:eq-local-balls-1}: $\chi(x)>0$ if and
only if $\tau(\{x\},[\nu])<M/2L$.

For the second inclusion, let $x\in[\nu]^{\approx}$. Then there are a cube $Q\in\mathfrak{c}^{\flat}$
of $[\nu]$ and a path $\gamma$ from $Q$ to $x$ of weighted length $<M/2L$. Let $n$ be the level of
$Q$: $Q$ is a level-$n$ $M$-cube in layer $n$ (the territory $X_n$), of side $M\ell_n$, with its
faces on the level-$n$ grid. In layer $m$ a lattice step has Euclidean length $\ell_m$ and weight
one, so the weighted length of a piece of path in layer $m$ is its Euclidean length divided by
$\ell_m$.

The path $\gamma$ cannot traverse a whole layer. Passing from $\Omega_{m+1}$ to $\Omega_m^c$ costs
weighted length at least
\begin{equation*}
\ell_m^{-1}\operatorname{dist}_\infty(\Omega_m^c,\Omega_{m+1})\ \ge\ M .
\end{equation*}
This lower bound is condition (A1) of \eqref{4.14:eq-prepared-data-admissible-a}, part of the
admissibility of the prepared large-field data (Proposition~\ref{4.14:prop-prepared-data-admissible},
which is stated in \cite[pp.~178--179]{B-CMP122a} without proof). Since $M>M/2L$, such a crossing
exceeds the length of $\gamma$. Passing from the inner to the outer side of layer $n$ through $Q$
itself also costs weighted length at least $M$, because $Q$ is $M\ell_n$ wide. Hence $\gamma$ lies
in the layers $\{n,n+1\}$ or in the layers $\{n-1,n\}$.

The Euclidean length of a path is at most its weighted length times the largest spacing of the
layers it meets. In the first case the Euclidean length of $\gamma$ is therefore smaller than
\begin{equation*}
\frac{M}{2L}\,\ell_{n+1}=\frac{M}{2}\,\ell_n .
\end{equation*}
In the second case it is smaller than
\begin{equation*}
\frac{M}{2L}\,\ell_n=\frac{M}{2}\,\ell_{n-1}.
\end{equation*}
Here $\ell_{j\pm1}=L^{\pm1}\ell_j$ was used.

Let $C\in\mathfrak{c}^{\flat}$ be a cube containing $x$. Its level $m$ is that of the layer
containing $x$, so $m\in\{n-1,n,n+1\}$, and $m=n+1$ occurs only in the first case. The faces of $C$
lie on the level-$\min\{m,n\}$ grid, of mesh $\mu:=M\ell_{\min\{m,n\}}$, and so do the faces of $Q$,
the partitions being compatible. Two closed grid-aligned boxes on a common grid of mesh $\mu$ are at
$\ell^\infty$-distance either $0$ or at least $\mu$. We have
\begin{equation*}
\operatorname{dist}_\infty(C,Q)\le|x-Q|,
\end{equation*}
and $|x-Q|$ is at most the Euclidean length of $\gamma$.

If $m\ge n$, then $\mu=M\ell_n$, and, whichever of the two layer alternatives $\gamma$ falls under,
its Euclidean length is smaller than
\begin{equation*}
\frac{M}{2}\,\ell_n<\mu ,
\end{equation*}
since $\frac{M}{2}\ell_{n-1}<\frac{M}{2}\ell_n$.

If $m=n-1$, then $\gamma$ is in the second case and $\mu=M\ell_{n-1}$. The Euclidean length of
$\gamma$ is $<\frac{M}{2}\ell_{n-1}<\mu$.

In both situations $\operatorname{dist}_\infty(C,Q)<\mu$, hence $\operatorname{dist}_\infty(C,Q)=0$.
So $C$ touches the cube $Q$ of $[\nu]$, and $x\in C\subset[\nu]^{\sim}$.

\emph{$(\beta)$.} At a point of $X_m$ the members of $\|\cdot\|_{\star}$ are $\ell_m|A'|$,
$\ell_m^2|\nabla A'|$ and $\ell_m^3|J|$. Since $0\le\chi\le1$,
\begin{equation*}
\ell_m|\chi A'|\le\ell_m|A'|,\qquad \ell_m^3|\chi J|\le\ell_m^3|J| .
\end{equation*}

Now let $y$ be joined to $x$ by one lattice step in $X_m$, of length $\ell_m$ and weight one. Since
$\tau(\cdot,[\nu])$ is a path distance,
\begin{equation*}
|\tau(\{y\},[\nu])-\tau(\{x\},[\nu])|\le1,
\end{equation*}
and $t\mapsto\max\{0,t\}$ is $1$-Lipschitz. Hence one lattice step changes $\chi$ by at most $2L/M$.
The discrete Leibniz rule gives
\begin{equation*}
\chi(y)A'(y)-\chi(x)A'(x)=\chi(y)\bigl(A'(y)-A'(x)\bigr)+\bigl(\chi(y)-\chi(x)\bigr)A'(x).
\end{equation*}
It follows that
\begin{equation*}
\ell_m^2|\nabla(\chi A')|\le\ell_m^2|\nabla A'|+\frac{2L}{M}\,\ell_m|A'| .
\end{equation*}

Each member of $\|(\chi A',\chi J)\|_{\star}$ is thus bounded by the corresponding member of
$\|(A',J)\|_{\star}$, except the gradient member, which is bounded by
$\ell_m^2|\nabla A'|+\frac{2L}{M}\ell_m|A'|$. Therefore
\begin{equation*}
\|(\chi A',\chi J)\|_{\star}\le\Bigl(1+\frac{2L}{M}\Bigr)\frac{r}{2}<r .
\end{equation*}
The last inequality holds as soon as $M>2L$, which is guaranteed by item (iii) of the conditions on
the block size $M$ in \eqref{4.14:eq-block-size-conditions-a}.

By $(\alpha)$, $\chi A'$ and $\chi J$ vanish off $[\nu]^{\approx}$. The point
$(e^{i\eta\chi A'}U,\chi J)$ is therefore defined from the given data alone, $U$ being given on the
whole lattice. Hence it lies in $\mathcal{U}(r)$.

\emph{$(\gamma)$.} By the reading property of the sheet format $\mathcal{S}_{u}$, $T_\nu$ reads
$\mathsf{u}=(\mathbf{U},J)$ only on $[\nu]$. On $[\nu]$ we have $\tau(\{x\},[\nu])=0$, hence
$\chi=1$, so the point of $(\beta)$ coincides there with the given data $(e^{i\eta A'}U,J)$.
Consequently $T_\nu$ takes the same value at every point of $\mathcal{U}(r)$ whose data agree with
the given ones on $[\nu]$. In particular its value at the point of $(\beta)$ does not depend on the
extension, and this value defines $T_\nu$ on the local ball.

The map $(A',J)\mapsto(\chi A',\chi J)$ is linear. By $(\beta)$ it takes the local ball into the
domain $\mathcal{U}(r)$ on which $T_\nu$ is analytic. So $T_\nu$ is analytic on the local ball.
The properties $(\mathrm{W}3)'$, $(\mathrm{W}5)^{\mathrm{cov}}$ and $(\mathrm{W}6)^{\mathrm{geo}}$,
listed in clause~(E) of Theorem~\ref{4.14:thm-operator-theorem}, are those of $T$.

\emph{The cut-off at $M/2$.} Let $Q$ be a cube of $[\nu]$ of level $n$ at a seam with the coarser
layer $n+1$, and let $C$ be a cube of level $n+1$, so that $\mu=M\ell_n$ in the notation of the
proof of $(\alpha)$. A path of weighted length smaller than $M/2$ in layer $n+1$ has Euclidean
length smaller than
\begin{equation*}
\frac{M}{2}\,\ell_{n+1}=\frac{LM}{2}\,\ell_n ,
\end{equation*}
and $\frac{LM}{2}\ell_n\ge M\ell_n=\mu$ because $L\ge2$. The resulting bound on
$\operatorname{dist}_\infty(C,Q)$ is therefore not smaller than $\mu$, and the alternative between
distance $0$ and distance at least $\mu$ no longer forces $C$ to touch $Q$. The argument for
$(\alpha)$ thus does not apply to the weighted ball of radius $M/2$ about $[\nu]$.
\end{proof}

Clause~(H) of the operator theorem at free current
(Theorem~\ref{4.14:thm-operator-theorem}, \eqref{4.14:eq-operator-theorem-a}) is
Lemma~\ref{4.14:lem-local-balls} at $r:=\varepsilon_1$, and at $r:=1$ for $\mathsf{J}$. Data given
on $[\nu]^{\approx}$ suffice; since $[\nu]^{\approx}\subset[\nu]^{\sim}$ by $(\alpha)$, data given on
the one-cube collar $[\nu]^{\sim}$ suffice a fortiori.

\begin{lemma}[Unit-lattice operators at complex configurations]\label{4.14:lem-unit-lattice-operators}
Fix the step $k$, and use Notation~\ref{4.14:not-free-current-notation} for the form
$\mathsf{Q}$ at free current. Write $\mathsf{M}$ for Ba{\l}aban's group average, which assigns to
a configuration $V$ on the unit lattice and a bond $c$ of the lattice of spacing $L$ an element
$\mathsf{M}(V)(c)$, and $\mathsf{M}^{k}$ for the $k$-fold iterate of the block-averaging map of
Definition~\ref{1.5:def-block-average}. For a point $y$ of
the lattice of spacing $L$ let $\mathcal{B}(y)$ be the unit-lattice block centred at $y$, and for
$x\in\mathcal{B}(y)$ let $V\langle y,x\rangle$ be the parallel transport of $V$ along Ba{\l}aban's
contour from $y$ to $x$. For a bond $c$ of the lattice of spacing $L$ let $c_-$, $c_+$ be its
initial and final points and $e_\mu$ the unit vector of its direction, and let $b_0(c)$ be the
unit bond contained in $c$ in its middle ($L$ is odd); we call $b_0(c)$ the pivot of $c$. For a
unit bond $b$ let $b_-$ be its initial point. For a field $B$ on unit bonds with values in the
complexified Lie algebra $\mathfrak{g}^{c}$ put
\begin{equation}\label{4.14:eq-unit-lattice-operators-1}
\begin{aligned}
\Phi^{\mathfrak{g}}(V,B;y,x)&:=\tfrac{1}{i}\log\bigl[(e^{iB}V)\langle y,x\rangle\,
V\langle y,x\rangle^{-1}\bigr],\\
\Phi^{\flat}(V,B;c)&:=\tfrac{1}{i}\log\bigl[\mathsf{M}(e^{iB}V)(c)\,\mathsf{M}(V)(c)^{-1}\bigr],
\end{aligned}
\end{equation}
and for $\Phi\in\{\Phi^{\mathfrak{g}},\Phi^{\flat}\}$ let $\Phi_n(V;B)$ be the part of $\Phi(V,B)$
homogeneous of degree $n$ in $B$. The local unit-lattice operators of $\mathsf{Q}$ at step $k$
are functions of $V:=\mathsf{M}^{k}(\mathbf{U})$ through \eqref{4.14:eq-unit-lattice-operators-1},
by the definition of the unit-lattice operators and \cite{B-CMP109}:
\begin{equation}\label{4.14:eq-unit-lattice-operators-2}
LQ^{\flat}B=\Phi^{\flat}_1(V;B),\qquad \mathfrak{g}(B)=\Phi^{\mathfrak{g}}_1(V;B),\qquad
C^{(2)}(V;B)=\Phi^{\flat}_2(V;B),
\end{equation}
with $C^{(2)}(V;B_1,B_2)$ the symmetric bilinear form whose diagonal is $C^{(2)}(V;B)$. The
remaining operators are the operator $\mathfrak{a}$ of the definition of the unit-lattice
operators; the kernel $Q^{\flat}(V;c,b)$ of $Q^{\flat}$, defined by
$(Q^{\flat}B)(c)=\sum_b L^{-d}Q^{\flat}(V;c,b)B(b)$; the pivot coefficient
$\pi_c:=Q^{\flat}(V;c,b_0(c))$; the operators
\begin{equation}\label{4.14:eq-unit-lattice-operators-3}
C_{\flat}(b_0(c),b)=-\pi_c^{-1}Q^{\flat}(V;c,b),\qquad \mathsf{h}(c)=L^{d-1}\pi_c^{-1},\qquad
\mathsf{G}:=\|\mathfrak{g}(\cdot)\|^2,
\end{equation}
where $\mathsf{G}$ is the term of \cite[(2.5)]{B-CMP109}; and the source vertex
$\mathcal{J}[\mathsf{f}]\colon B\mapsto -2\langle \mathsf{h}C^{(2)}(B),\mathsf{f}\rangle$.

There is a number $r_{\flat}>0$ depending only on $d$ and $L$ with the following property. Let
$W$ be $G$-valued with $|W(\partial p)-1|\le r_{\flat}$ for the unit plaquettes $p$ of the (at
most two) blocks which a given operator reads, let $\bar a_0:=\sup_p|W(\partial p)-1|$ over these
plaquettes, so that $\bar a_0\le r_{\flat}$, and put
\begin{equation*}
\mathsf{P}:=\Bigl\{(A'',B):\ \sup_b|A''_b|<r_{\flat},\ \sup_b|B_b|<r_{\flat}\Bigr\}.
\end{equation*}
\begin{enumerate}
\item For $\Phi\in\{\Phi^{\mathfrak{g}},\Phi^{\flat}\}$ the map
$(A'',B)\mapsto\Phi(e^{iA''}W,B)$ is jointly analytic on $\mathsf{P}$, and $|\Phi|\le1$ there.
\item On $\mathsf{P}$, for every $n$,
\begin{equation}\label{4.14:eq-unit-lattice-operators-4}
\begin{aligned}
|\Phi_n(e^{iA''}W;B)|&\le(\sup|B|/r_{\flat})^n,\\
|\Phi_n(e^{iA''}W;B)-\Phi_n(W;B)|&\le 2(\sup|A''|/r_{\flat})(\sup|B|/r_{\flat})^n.
\end{aligned}
\end{equation}
\item Assume $C'_1L^2\bar a_0\le\tfrac12$. Then $\pi_c(W)=\mathsf{R}(1+E)$, where
$\mathsf{R}:=\operatorname{Ad}W([c_-,b_0(c)_-])$ is an isometry of $\mathfrak{g}^{c}$ and
\begin{equation*}
\|E\|\le C'_1L^2\bar a_0\le\tfrac12;
\end{equation*}
here $C'_1=C\,L^{d-1}$ is the constant of the second member of \cite[(139)]{B-CMP98} in the form
this display takes in the analyticity theorem for the averaging operations, $C$
depending on $d$ and on the average only, so that $C'_1$ depends only on $d$ and $L$, and
$C'_1L^2\bar a_0\le C'_1L^2r_{\flat}$ since $\bar a_0\le r_{\flat}$.
Consequently $\|\pi_c(W)^{-1}\|\le2$, and if moreover
$8L^{d-1}\sup|A''|\le r_{\flat}^2$, then $\pi_c(e^{iA''}W)$ is invertible with
$\|\pi_c(e^{iA''}W)^{-1}\|\le4$.
\item In particular, let $\mathsf{u}=(\mathbf{U},J)\in\mathcal{U}(1)$ with
$\mathbf{U}=e^{i\eta A'}U$, where $\eta$
is as in Definition~\ref{4.14:def-real-backgrounds}, and
$U\in\mathcal{U}_{\mathbb{R}}$, the class of that definition. Assume the ceiling on the regularity
constants on cubes of side $M$ in~\eqref{4.14:eq-block-size-conditions-b}, and let the radius of
$\mathcal{U}(1)$ in the $A'$-direction be $a_1^{\oplus}$
of~\eqref{4.14:eq-block-size-conditions-d}, so that $\alpha_1:=\sup_n\ell_n|A'|<a_1^{\oplus}$. Then
$W:=\mathsf{M}^{k}(U)$ satisfies the hypotheses above: it is
$G$-valued with unit plaquettes bounded by $r_{\flat}$, and $V=e^{iA''}W$ with
$\sup|A''|<r_{\flat}$. The
operators $LQ^{\flat}$, the kernel of $Q^{\flat}$, $\mathfrak{g}$, $\mathfrak{a}$, $C^{(2)}$,
$\pi_c^{-1}$, $C_{\flat}$, $\mathsf{G}$ and $\mathcal{J}[\mathsf{f}]$ are analytic in $A'$ ($J$
does not enter), keep their supports, the gauge covariance of \cite[(0.6)]{B-CMP109} and their
locality with respect to
the geometry, and obey the bounds
\begin{equation}\label{4.14:eq-unit-lattice-operators-5}
\begin{gathered}
|Q^{\flat}(V;c,b)|\le L^{d-1}/r_{\flat},\qquad |\mathfrak{g}(B)(y,x)|\le\sup|B|/r_{\flat},\\
|C^{(2)}(V;B_1,B_2)(c)|\le 4\sup|B_1|\,\sup|B_2|/r_{\flat}^2,\qquad \|\pi_c^{-1}\|\le4,\\
|C_{\flat}(b_0(c),b)|\le 4L^{d-1}/r_{\flat},\qquad
\|\chi_y\mathcal{J}[\mathsf{f}]\|\le C_{\mathcal{J}}(d,L)\sup|\mathsf{f}|,
\end{gathered}
\end{equation}
with $\chi_y$ the Lipschitz cut-off at $y$ and $C_{\mathcal{J}}(d,L)$ depending only on $d$ and
$L$.
\item In the setting of item (4), the multi-level averaging operators $Q_j(\mathbf{U})$ of all
layers $j$, which are compositions of the one-step operators of \cite[(124)]{B-CMP98} (see
\cite[(3.15)]{B-CMP99b}), satisfy
\begin{equation}\label{4.14:eq-unit-lattice-operators-6}
|Q_j(\mathbf{U};c,b)|\le 2(1+\theta_k)\le4,\qquad \theta_k:=6C'_1\alpha_0+\kappa\alpha_1,\qquad
\kappa:=4\bigl(3c'''+2c^{\natural}/L\bigr),
\end{equation}
where $\alpha_1$ is as in item (4), $\alpha_0$ is the radius entering condition (u1)
in~\eqref{4.14:eq-real-backgrounds-a}, and $c'''$, $c^{\natural}$ are the constants of the
complex averaging bounds of the proposition on the unit-lattice averaging operators at complex
configurations.
Here $\kappa$ is not the auxiliary parameter $\kappa$ of $\mathfrak{p}$. The one-step operator of
\cite[(124)]{B-CMP98} is analytic, its analyticity being obtained, as in the analyticity theorem
for the averaging operations, from $L^{-1}$ times the derivative at the value zero of the
parameter of the generating function of that theorem; the bounds
\eqref{4.14:eq-unit-lattice-operators-5} enter that argument as values.
\end{enumerate}
\end{lemma}

\begin{proof}
\emph{Item (1) for $\Phi^{\mathfrak{g}}$.} Joint analyticity and the bound are clause (c) of the
analyticity theorem for the averaging operations at complex backgrounds: $\Phi^{\mathfrak{g}}$ is
the logarithm of a composition of products of group elements near the identity, the group average
$\mathsf{M}$ is analytic on its disc, and the logarithm is analytic near $1$. The final assertion
of that clause concerns backgrounds within $e^{\tau}$ of unitary ones, among them the background
$V'V_0$ of
\cite[Proposition~7]{B-CMP98} with $\tau=O(L\alpha_1)$: the same statements hold with the
parameter $\rho$ of that clause replaced by $\rho+C\tau$ and with the isometry factors replaced by
$e^{2\tau}$.

\emph{Item (1) for $\Phi^{\flat}$.} By \cite[(89), (62)]{B-CMP98}, in the form these displays
take in the analyticity theorem for the averaging operations (there they incorporate also
\cite[(120)]{B-CMP98}), in the notation of those displays, in which $R(u)$ denotes the adjoint
action of $u$ on $\mathfrak{g}^{c}$ (not the ball radius $R$) and $e^{iQ}$ is the middle factor of
\cite[(89)]{B-CMP98} (not the operator $Q^{\flat}$),
\begin{equation*}
e^{i\Phi^{\flat}(V,B;c)}=(\bar R_{0,c_-}V_1)\cdot e^{iQ}\cdot
\bigl[R(\mathsf{M}(V)(c))(\bar R_{0,c_+}V_1)\bigr]^{-1},
\end{equation*}
\begin{equation*}
\bar R_{0,y}V_1=\exp\Bigl[i\sum_{x\in\mathcal{B}(y)}L^{-d}\Phi^{\mathfrak{g}}(V,B;y,x)\Bigr].
\end{equation*}
This is a product of three factors, each jointly analytic and near the identity. The two outer
factors are such by the case of $\Phi^{\mathfrak{g}}$ just treated, including that final
assertion. The
middle factor is such by \cite[Proposition~3]{B-CMP98} at the background $V'V_0$; Ba{\l}aban
states in \cite{B-CMP98} that this proposition holds uniformly for $V'V_0$ in place of $V_0$, with
a change only in \cite[(126)]{B-CMP98}, and in the present setting this is supplied by the
uniformity lemma for
Ba{\l}aban's averaging propositions and by the clause of the analyticity theorem for the averaging
operations that places those propositions here. The adjoint factor obeys
$\|\operatorname{Ad}\mathsf{M}(V)(c)\|\le e^{O(1)r_{\flat}}$ by \cite[(164)]{B-CMP98}. The
logarithm of a product within $\tfrac12$ of $1$ is analytic; we decrease $r_{\flat}$ until
$|\Phi|\le1$ on $\mathsf{P}$. Every quantity used depends only on $d$ and $L$, so $r_{\flat}$ does.

\emph{Item (2).} Since $e^{i0}V=V$, we have $\Phi(\cdot;0)=0$. Fix $(A'',B)\in\mathsf{P}$ with
$B\ne0$ (for $B=0$ both sides of \eqref{4.14:eq-unit-lattice-operators-4} vanish when $n\ge1$).
By homogeneity $\Phi_n(V;sB)=s^n\Phi_n(V;B)$, so $s\mapsto\Phi(e^{iA''}W,sB)$ has the power
series $\sum_n s^n\Phi_n(e^{iA''}W;B)$. It is analytic for $|s|<r_{\flat}/\sup|B|$, because then
$(A'',sB)\in\mathsf{P}$, and it is bounded by $1$ there by item (1). The Cauchy estimate on
circles of radius below $r_{\flat}/\sup|B|$ gives the first inequality. The homogeneous part
$\Phi_n(e^{iA''}W;B)$ is a Cauchy integral of the analytic function of item (1), hence analytic
in $A''$. For fixed $B$ and $A''\ne0$ put
\begin{equation*}
g(z):=\Phi_n\bigl(e^{iz(r_{\flat}/\sup|A''|)A''}W;B\bigr),\qquad |z|<1 .
\end{equation*}
For $|z|<1$ the argument $z(r_{\flat}/\sup|A''|)A''$ has supremum below $r_{\flat}$, so $g$ is
analytic with $|g|\le K:=(\sup|B|/r_{\flat})^n$ by the first inequality. The function
$(g-g(0))/2K$ is analytic on the unit disc, bounded by $1$ and vanishes at $0$; the Schwarz lemma
gives $|g(z)-g(0)|\le2K|z|$. At $z_0:=\sup|A''|/r_{\flat}<1$ one has
$g(z_0)=\Phi_n(e^{iA''}W;B)$ and $g(0)=\Phi_n(W;B)$, which is the second inequality.

\emph{Item (3).} By the support property in the definition of the unit-lattice operators,
$\mathfrak{g}$ and $\mathfrak{a}$ read only bonds inside the blocks $\mathcal{B}(y)$, and none of
these bonds is a pivot. Hence the coefficient of $b_0(c)$ in $Q^{\flat}$ equals the coefficient
of $b_0(c)$ in $Q_1:=Q^{\flat}+\bar D_1\mathfrak{a}$, where $\bar D_1$ is the difference operator
of the definition of the unit-lattice operators; $Q_1$ is the operator of
\cite[(124)]{B-CMP98} at the background $W$, in the form this display takes in the analyticity
theorem for the averaging operations. In that form the main term, applied to a field $B$ on unit
bonds, is
\begin{equation*}
Q'_0B=\sum_{x\in\mathcal{B}(c_-)}L^{-(d+1)}\sum_{b\in[x,x']}
R\bigl(W\langle c_-,x\rangle\,W([x,b_-])\bigr)B(b),\qquad x':=x+Le_\mu,
\end{equation*}
with $R$ again the adjoint action,
where $[x,x']$ is the contour obtained by parallel transport of $c$ to the point $x$
\cite{B-CMP109}. The pivot $b_0(c)\subset c$ is the bond of the unit lattice contained in $c$ and
belonging to the corridor \cite{B-CMP109}; as $L$ is odd it is the middle bond of $c$. Write
$x=c_-+te_\mu+x^{\perp}$ with $x^{\perp}\perp e_\mu$ and $|t|\le(L-1)/2$ (the blocks are centred).
Then $b_0(c)\in[x,x+Le_\mu]$ if and only if $x^{\perp}=0$: for $x^{\perp}\ne0$ the contour does
not meet the line through $c$, while for $x^{\perp}=0$ the contour runs from $c_-+te_\mu$ to
$c_-+(t+L)e_\mu$ and contains
\begin{equation*}
b_0(c)=\bigl[c_-+\tfrac{L-1}{2}e_\mu,\ c_-+\tfrac{L+1}{2}e_\mu\bigr]
\end{equation*}
because $t\le(L-1)/2$ and $t+L\ge(L+1)/2$; this holds for all $L$ values of $t$. For each of
these $L$ points $x$ the transporter from $c_-$ to $x$ and then to $b_0(c)_-$ runs back and forth
on the one line through $c$: by the contour convention of the analyticity theorem for the
averaging operations the contour on that line is the single straight run, and the group average of
a one-element family is that element. Since $W$ is $G$-valued, reversing a bond inverts its
variable, so the retraced segments cancel and the transporter is $W([c_-,b_0(c)_-])$; its adjoint
action is $\mathsf{R}:=\operatorname{Ad}W([c_-,b_0(c)_-])$, an isometry of $\mathfrak{g}^{c}$.
Thus the coefficient of $Y:=B(b_0(c))$ in the main term is $L\cdot L^{-(d+1)}\mathsf{R}Y
=L^{-d}\mathsf{R}Y$.

The remainder $Q''(W)$ of the same display acts on $X:=Y\delta_{b_0(c)}$, and by the second
member of \cite[(139)]{B-CMP98} and by \cite[(140)]{B-CMP98}, in the form used in the analyticity
theorem for the averaging operations,
\begin{equation*}
|Q''(W)X|\le C'_1L^2\bar a_0\,(Q''|X|)_c=C'_1L^2\bar a_0\,L^{-d}|Y|.
\end{equation*}
Removing the factor $L^{-d}$ of the kernel convention, $\pi_c(W)=\mathsf{R}(1+E)$, where
$E:=\mathsf{R}^{-1}L^{d}Q''(W)(\cdot\,\delta_{b_0(c)})_c$ satisfies $\|E\|\le C'_1L^2\bar a_0$.
Since $\bar a_0\le r_{\flat}$ this is at most $C'_1L^2r_{\flat}$, and it is at most $\tfrac12$ by
the hypothesis $C'_1L^2\bar a_0\le\tfrac12$ of item (3). As $\mathsf{R}$ is an isometry, the
Neumann series gives
\begin{equation*}
\|\pi_c(W)^{-1}\|=\|(1+E)^{-1}\mathsf{R}^{-1}\|\le\|(1+E)^{-1}\|\le\frac{1}{1-\|E\|}\le2.
\end{equation*}
At $W=1$ the result agrees with \cite[(0.12)]{B-CMP109}, where the coefficient of the pivot in
$LQ^{\flat}$ is $L^{1-d}$. For $d=1$, $L=3$: $c=[0,3]$, $b_0(c)=[1,2]$, and $b_0(c)$ lies on each
of the three contours $[-1,2]$, $[0,3]$, $[1,4]$.

\emph{The increment.} Let $V=e^{iA''}W$ with $(A'',0)\in\mathsf{P}$. For $B:=Y\delta_{b_0(c)}$ the
convention for the kernel and \eqref{4.14:eq-unit-lattice-operators-2} give
$\Phi^{\flat}_1(V;B)(c)=L\cdot L^{-d}\pi_c(V)Y$. Item (2) at $n=1$ (both sides are linear in $B$,
so it extends from $\mathsf{P}$ to all $Y$) reads
\begin{equation*}
\bigl|L^{1-d}(\pi_c(V)-\pi_c(W))Y\bigr|\le2\,\frac{\sup|A''|}{r_{\flat}}\,\frac{|Y|}{r_{\flat}},
\quad\text{that is,}\quad
\|\pi_c(V)-\pi_c(W)\|\le\frac{2L^{d-1}\sup|A''|}{r_{\flat}^2}.
\end{equation*}
With $\|\pi_c(W)^{-1}\|\le2$,
\begin{equation*}
\bigl\|\pi_c(W)^{-1}(\pi_c(V)-\pi_c(W))\bigr\|\le\frac{4L^{d-1}\sup|A''|}{r_{\flat}^2}\le\frac12
\end{equation*}
under $8L^{d-1}\sup|A''|\le r_{\flat}^2$. Writing
$\pi_c(V)=\pi_c(W)\bigl(1+\pi_c(W)^{-1}(\pi_c(V)-\pi_c(W))\bigr)$, the Neumann series shows that
$\pi_c(V)$ is invertible with $\|\pi_c(V)^{-1}\|\le2\cdot2=4$, and $\pi_c(V)^{-1}$ is analytic in
$A''$ because $\pi_c(V)$ is (item (2)) and the Neumann series converges uniformly.

\emph{Item (4).} The configuration $W:=\mathsf{M}^{k}(U)$ is $G$-valued. The sets read by the
operators lie within distance $L+1<M$ of a bond of $\mathbb{B}^{+}$, so each of them lies,
together with that bond, inside a cube of side $M$ holding a bond of $\mathbb{B}^{+}$, on which the
regularity condition on cubes of side $M$ of~\eqref{4.14:eq-real-backgrounds-a} holds. Let
$\mathsf{o}_1$ and $\mathsf{o}_2$ be the constants appearing in~\eqref{4.14:eq-block-size-conditions-b}
and~\eqref{4.14:eq-block-size-conditions-d}; $\mathsf{o}_1$ is the constant of
\cite[(53)]{B-CMP98} and $\mathsf{o}_2$ the constant of \cite[(164)]{B-CMP98}. By
\cite[Proposition~2, (53)]{B-CMP98} the unit-lattice
plaquettes of $W$ there are bounded by $2\mathsf{o}_1$ times the plaquette constant of that
regularity condition, and this is at most $r_{\flat}$
by the member of the ceiling on the regularity constants
in~\eqref{4.14:eq-block-size-conditions-b} which bounds $2\mathsf{o}_1$ times that plaquette
constant by $r_{\flat}$. The member $C'_1L^2\bar a_0\le\tfrac12$ of that ceiling, for the
unit-plaquette bound $\bar a_0$ of $W$, is the
hypothesis of item (3). Next, $V:=\mathsf{M}^{k}(e^{i\eta A'}U)=e^{iA''}W$, where
\begin{equation*}
|A''_b|\le 2\mathsf{o}_2\sup_n\ell_n|A'|\ \text{ for every bond } b,\qquad
\sup_n\ell_n|A'|<a_1^{\oplus}.
\end{equation*}
This follows from
\cite[Proposition~6, (164)]{B-CMP98}, which states that the averaged complex configuration
multiplied by the inverse of the averaged real one is close to $1$, at a distance proportional to
$\sup_n\ell_n|A'|$, together with \cite[(27)]{B-CMP98}, $|\log u|\le2|u-1|$ for $u$ near $1$; by
\cite[Proposition~6]{B-CMP98} $A''$ is analytic in $A'$. The second member of the minimum defining
$a_1^{\oplus}$ in~\eqref{4.14:eq-block-size-conditions-d} is $r_{\flat}/(2\mathsf{o}_2)$, whence
$\sup|A''|<2\mathsf{o}_2a_1^{\oplus}\le r_{\flat}$; the third is
$r_{\flat}^2/(16L^{d+1}\mathsf{o}_2)$, whence
\begin{equation*}
8L^{d-1}\sup|A''|\le16L^{d-1}\mathsf{o}_2\,a_1^{\oplus}\le r_{\flat}^2/L^2\le r_{\flat}^2 .
\end{equation*}
Items (1)--(3) and the increment therefore apply. The operators are
$\Phi^{\flat}_1$, $\Phi^{\mathfrak{g}}_1$, $\Phi^{\flat}_2$, $\pi_c^{-1}$ and the formulas
\eqref{4.14:eq-unit-lattice-operators-3}. Item (2) with $n=1$ and $B=Y\delta_b$ gives
$|L^{1-d}Q^{\flat}(V;c,b)Y|\le|Y|/r_{\flat}$, which is the bound on the kernel of $Q^{\flat}$; with
$n=1$ for $\Phi^{\mathfrak{g}}$ it gives the bound on $\mathfrak{g}$; with $n=2$ for
$\Phi^{\flat}$ it gives $|C^{(2)}(V;B)(c)|\le(\sup|B|/r_{\flat})^2$, and polarisation, which costs
at most the factor $4$, gives the bound on $C^{(2)}(V;B_1,B_2)$. The increment gives
$\|\pi_c^{-1}\|\le4$, and then $|C_{\flat}(b_0(c),b)|\le4\cdot L^{d-1}/r_{\flat}$ by
\eqref{4.14:eq-unit-lattice-operators-3}. Supports, gauge covariance \cite[(0.6)]{B-CMP109} and
locality
are properties of $\Phi^{\mathfrak{g}}$ and $\Phi^{\flat}$, hence of their homogeneous parts. The
constant $C_{\mathcal{J}}(d,L)$ collects the bound $4$ on $\pi_c^{-1}$, the factor $4/r_{\flat}^2$ of
$C^{(2)}$ and the number of pivots per block. Analyticity in $A'$ follows from the joint
analyticity of item (1) composed with the analytic map $A'\mapsto A''$; $J$ enters none of these
formulas.

\emph{Item (5).} This is the first clause of the proposition on the unit-lattice averaging
operators at complex configurations, whose first item is the complex form of
\cite[(139)]{B-CMP98}; the clause of the analyticity theorem for the averaging operations that
places Ba{\l}aban's propositions in the present setting supplies its inputs. The multi-level
operators $Q_j$ are compositions of one-step operators of the type \cite[(124)]{B-CMP98}
(\cite[(3.15)]{B-CMP99b}, and the composition statement in the definition of the unit-lattice
operators), so that proposition applies to them. As in the analyticity theorem for the averaging
operations, the analyticity of the one-step operator of \cite[(124)]{B-CMP98} is read from $L^{-1}$
times the derivative at the value zero of the parameter of the generating function of that
theorem, and the bounds \eqref{4.14:eq-unit-lattice-operators-5} enter that reading as values.

\emph{Two remarks on dependence.} First, the complex form of \cite[(139)]{B-CMP98} used in item
(5) does not depend on the analyticity
theorem for the averaging operations; its inputs enter that theorem through the clause placing
Ba{\l}aban's propositions in the present setting. Second, the bounds of items (3)--(4) depend only
on $d$
and $L$, not on $k$, and the lemma bounding the form $\mathsf{Q}$ uses them only through the
constant of its bound and the radii $a_1^{\oplus}$, $a_0^{\oplus}$.
\end{proof}

\begin{lemma}[Walk expansion of the form and base-point bound]\label{4.14:prf-form-walk-expansion}
Assume the hypotheses of the operator theorem at free current,
Theorem~\ref{4.14:thm-operator-theorem}. These are: conditions (A1), (A2) of
\eqref{4.14:eq-prepared-data-admissible-a}, the conditions of
Proposition~\ref{4.14:prop-prepared-data-admissible} (printed without proof in
\cite[pp.~178--179]{B-CMP122a}; cited as printed); the floors (fl-99), (fl-$\sharp$)$'$,
(fl-$M$)$'$, (fl-$L$) of \eqref{4.14:eq-block-size-conditions-a}; and the ceilings
(cl-99), (cl-reg)$^{+}$, (cl-reg)$^{\sharp}$ of \eqref{4.14:eq-block-size-conditions-b}.
Then the following hold.
\begin{enumerate}
\item[(a)] On $\mathcal{U}(1)$ the series (q1) and (q2) of clause (A) of
Theorem~\ref{4.14:thm-operator-theorem} converge absolutely, and
\begin{equation*}
G_{\pi}=(G_0^{-1}+W_{\pi})^{-1},\qquad G_1=(G_0^{-1}+W_1)^{-1};
\end{equation*}
the operators $T$, $T_{\pi}$, $T_1$ are invertible.
\item[(b)] The form $\mathsf{Q}$ is a sheet of the unit frame in the class
\begin{equation}\label{4.14:eq-form-walk-expansion-1}
\mathsf{Q}\in\mathcal{S}_{u}[\rho_{\mathsf{Q}};B_{\mathsf{Q}},\delta_{\mathsf{Q}};1],
\qquad \rho_{\mathsf{Q}}=\rho_T=8dL^2M_1^{\sharp}.
\end{equation}
\item[(c)] $\mathsf{Q}$ is complex symmetric, and real symmetric at real points; it has
the covariance property $(\mathrm{Sh}4)$, and the geometry-locality property
$(\mathrm{Sh}5)$ with respect to the lists enlarged by the reading hulls.
\item[(d)] With $S_{u}(\cdot)$ the Schur sum of the unit frame,
\begin{equation*}
\|\mathsf{Q}\|\le e\,B_{\mathsf{Q}}\,S_{u}(\delta_{\mathsf{Q}})=\gamma_1 .
\end{equation*}
\item[(e)] At the base point, $\mathsf{Q}(\mathbf{1},0)\ge\gamma_0$.
\end{enumerate}
\end{lemma}

This is the first step of the proof of Theorem~\ref{4.14:thm-operator-theorem}.

\begin{proof}
We supply the inputs (i)--(vii) in the order in which the assembly of step (vi)
uses them. After the inputs we list the hypotheses of
Theorem~\ref{4.14:thm-operator-theorem} that are used. We then explain why the proofs
of the cited lemmas apply unchanged under the conventions
\eqref{4.14:eq-free-current-notation-a}--\eqref{4.14:eq-free-current-notation-c} of
Notation~\ref{4.14:not-free-current-notation}.

(i) Elementary sheets. By the elementary-sheet lemma, the operators $G'$, $G_0$, $T'$
and the products
\begin{equation*}
(DG'),\qquad (G'D^{\ast}),\qquad (D^{\ast}G_0),\qquad (G_0D)
\end{equation*}
have walk expansions at complex $\mathbf{U}$ that satisfy the norm clauses (n1)--(n3)
of that lemma.

That lemma rests on \cite[Theorems~3.4, 3.7, 3.9, 3.10]{B-CMP99b}, which require of the
background only \cite[(3.35)]{B-CMP99b}, that is, condition (u1) of
\eqref{4.14:eq-real-backgrounds-a}. Walk by walk, clause (n2) is the second inequality of
\cite[(3.43)]{B-CMP99b} with the parameter there called $\beta$ equal to $\tfrac12$.
Clause (n3) is \cite[(3.44)]{B-CMP99b} with the parameter there called $\varepsilon$
equal to $\tfrac12$; both are on \cite[p.~398]{B-CMP99b}. The bounds
\cite[(3.94), (3.108)]{B-CMP99b} carry an additional factor $M_1^{-|\omega|/2}\le1$,
the block size of \cite{B-CMP99b} being the present $M_1$, which only strengthens them.

The domains are taken as in convention \eqref{4.14:eq-free-current-notation-a}:
\begin{equation*}
\widehat{X}:=\bigcup_{\square\subset X}\square^{\max\{5,2L\}},\qquad
\rho_0:=c_X(d)\,L^2M_1 .
\end{equation*}
With these domains three things hold. First, part (iv) of the lemma on reading hulls
of second-cover boxes (Lemma~\ref{4.14:lem-reading-hulls}) gives $(\mathrm{Sh}5)$.
Second, the diameter in $(\mathrm{Sh}1)$ becomes $\rho_0$. Third, the constants
$B_{\mathrm{el}}$, $\delta_{\mathrm{el}}$ of the elementary-sheet lemma are unchanged.
The reason is that the walk weight $a_{\omega}(y)$ is read on $X_0$ and not on
$\widehat{X}$, and condition $(\tau1)$ of the many-scale frame, one of the conditions
$(\tau0)$--$(\tau3)$ recalled in step (t1) below, keeps the decay property
$(\mathrm{Sh}3)$.

(ii) Local operators at complex $\mathbf{U}$. The locality lemma is applied with the
symmetric cell weights $w_z$ of convention \eqref{4.14:eq-free-current-notation-b}.
Part (a) of Lemma~\ref{4.14:lem-symmetry-covariance} (symmetric cell weights and
lattice-symmetry covariance) shows that this changes only the values of the constants
$\nu$, $c_N$ of the locality lemma, whose decay rate is $\delta$, and of the constant
$n_N$ of the branching lemma, with $n_w\le3(2dL+1)^d$ the number defined in that part:
\begin{equation*}
\nu e^{\delta c_N}\;\longmapsto\; n_w\,\nu\,e^{\delta(c_N+dL)},\qquad
c_N\;\longmapsto\; c_N+dL,
\end{equation*}
and, in the branching lemma, $n_N\mapsto n_Nn_w$. The properties
$(\mathrm{Sh}1)$--$(\mathrm{Sh}5)$ are kept.

The local operators so expanded are the following.
\begin{itemize}
\item The operators $Q'$, $Q'^{T}$, $\mathfrak{S}$, the weight $a$ and $M_{\mathsf{f}}$.
These come from the elementary-sheet lemma and from Ba{\l}aban's construction.
\item The operators $Q$ and $Q^{T}$, in all layers. These come from the first clause of
the complex localisation proposition, which is the complex form of
\cite[(147)]{B-CMP98}.
\item The operator $\Delta^{(2)}[\mathsf{f}]$ at complex $\mathbf{U}$, obtained by the
chain below.
\end{itemize}

The chain for $\Delta^{(2)}[\mathsf{f}]$ has five links, each feeding the next.
\begin{enumerate}
\item[(1)] The first lemma of the complex localisation analysis, the complex form of the
bound \cite[(139)]{B-CMP98}. It uses no expansion scheme.
\item[(2)] The complex localisation proposition. The radius $\alpha_0$ is placed as in
the complex averaging bounds, and the floor $L\ge 8$ contained in (fl-$L$) of
\eqref{4.14:eq-block-size-conditions-a} is imposed.
\item[(3)] Its first corollary, the complex form of \cite[(148)]{B-CMP98}, obtained from
the proposition by the Cauchy estimate.
\item[(4)] Part (iii$'$) of that corollary, combined with the lemma on the complex
averaging operation.
\item[(5)] Part (iv) of the complex projection statement, the complex form of the
corresponding real-point bound of Ba{\l}aban's construction, which gives, at complex
$\mathbf{U}$,
\begin{equation*}
\|\chi_y\,\Delta^{(2)}[\mathsf{f}]\|\le C_2\,\ell_y^{-2}\,\sup_{b}\,\ell_b^{3}|\mathsf{f}(b)| .
\end{equation*}
Here $\chi_y$, $\ell_y$ and $C_2$ are, as fixed in the complex projection statement,
the Lipschitz cut-off localising at the point $y$, the lattice spacing at $y$ in
level-$k$ units, and the constant of the bound; the supremum is taken over the
arguments $b$ of $\mathsf{f}$, with $\ell_b$ the lattice spacing at $b$ in level-$k$
units.
\end{enumerate}
The last bound rests on \cite[(142), (159)--(163)]{B-CMP98}. The constants
\begin{equation*}
\theta_k:=6C_1'\alpha_0+\kappa\alpha_1,\qquad \kappa:=4\bigl(3c'''+2c^{\natural}/L\bigr),
\end{equation*}
where $\kappa$ is the constant so written in the complex averaging bounds and not the
auxiliary parameter $\kappa$ of the glossary, and the differentiation identity relating
two kernels of the complex averaging bounds, enter only through their values.

(iii) Unit-lattice operators at complex $\mathbf{U}$. The operators $P_{\Lambda}$,
$C_{\flat}$, $C_{\flat}^{T}$, $\mathsf{G}$ and $\mathcal{J}[\mathsf{f}]$ at complex
$\mathbf{U}$ are supplied by Lemma~\ref{4.14:lem-unit-lattice-operators} (unit-lattice
operators at complex configurations). That lemma uses the $r_{\flat}$-members of the
ceiling (cl-reg)$^{+}$ in \eqref{4.14:eq-block-size-conditions-b}. It also uses the
second, third and fourth members of the minimum defining the radius
$a_1^{\oplus}$ in \eqref{4.14:eq-block-size-conditions-d}.

(iv) The singular word. Write
\begin{equation*}
\mathfrak{b}:=(DG')(D^{\ast}G_0D)(G'D^{\ast})
\end{equation*}
for the one singular word. The singular-word lemma expands it with orders $(1,0,1)$
in its three factors and with smooth junctions between them.

(v) The walk sheet of $T^{-1}=(QG_0Q^{T})^{-1}$. This is the walk-sheet lemma for
$T^{-1}$. The steps of its proof that are used here are (t1), (t2), (t3), (t5), (t6),
carried out as follows.

(t1) The product lemma and part (F1) of the flattening lemma are applied to
$QG_0Q^{T}$ in the many-scale frame $\mathfrak{F}^{ms}$. The frame has the following
ingredients:
\begin{itemize}
\item the anchor points of coarse bonds $c$, the distance between the anchors of two
coarse bonds $c$ and $c'$, and the connection lemma of the many-scale frame;
\item the conditions $(\tau0)$--$(\tau3)$ with $c_{\tau}=2L$;
\item the sum $S_{ms}(\cdot)$, bounded by \cite[Lemma~2.1]{B-CMP96}, which applies
under the condition of (fl-$M$)$'$(ii) of \eqref{4.14:eq-block-size-conditions-a}.
\end{itemize}
The frame is used together with the box system of the second cover: its properties
(b1)--(b4) and (Sc), and the constants
\begin{equation*}
N_d=2\cdot5^{d},\qquad \lambda^{\sharp}=C_dL^2/M_1^{\sharp},\qquad
w^{\sharp}=M_1^{\sharp}/2L .
\end{equation*}
Property (b3) is used at the value $\rho_0$ of step (i). It follows from
$M_1^{\sharp}>2L\rho_0$, which is part of the floor (fl-$\sharp$)$'$ of
\eqref{4.14:eq-block-size-conditions-a}, read with $\rho_0$ as in
\eqref{4.14:eq-free-current-notation-a}. In (b4) the radius is
$\rho_{\square}=\rho_T$.

The hulls $\square^{\mathsf{r}}$ are listed in $\operatorname{dom}\sigma$ and
$\operatorname{dom}\varpi$, as in convention \eqref{4.14:eq-free-current-notation-a}.
By parts (i)--(iii) of Lemma~\ref{4.14:lem-reading-hulls}, properties
$(\mathrm{Sh}1)$--$(\mathrm{Sh}4)$ of the inversion lemma then hold exactly as stated
there, and $(\mathrm{Sh}5)$ holds with respect to the enlarged lists.

(t2) Here $\widetilde{T}$ denotes the operator of the many-scale frame attached to
$T=QG_0Q^{T}$ in the walk-sheet lemma; its inverse is produced by the inversion lemma in
step (t5), and part (F2) of the flattening lemma recovers $T^{-1}$ from it in step (t6).
We show that
\begin{equation*}
\widetilde{T}(U)\ge m_0:=\tfrac12 m^{\mathrm{flat}}\qquad\text{for every }
U\in\mathcal{U}_{\mathbb{R}} .
\end{equation*}
This follows from the positivity lemma at the identity combined with the positivity
lemma on real backgrounds, as follows.

The positivity lemma at the identity applies \cite[Proposition~2.7]{B-CMP96} at
$U=\mathbf{1}$, with the region $\Omega$ there enlarged by a collar. Combined with the
Schur bound, this gives $\widetilde{T}(\mathbf{1})\ge m^{\mathrm{flat}}$.

The positivity lemma on real backgrounds proceeds in seven steps.
\begin{itemize}
\item[($\alpha$)] An IMS localisation over the boxes of the second cover, using the box
properties (b1)--(b2), gives the bound
$\|\mathsf{D}\|\le m^{\mathrm{flat}}/6$ on the IMS error $\mathsf{D}$. This bound
follows from the floor $\mathsf{F}^{P}_{ms}$, which is part of (fl-$\sharp$)$'$.
\item[($\beta$)] For each box $\square$ let $\widetilde{T}^{\square}$ be the partial walk
sum of $\widetilde{T}$. Its tail $\mathsf{E}_{\square}$ satisfies
$\|\mathsf{E}_{\square}\|\le m^{\mathrm{flat}}/12$.
\item[($\gamma$)] On $\square^{4}\cap\Omega_0$ take the gauge $\tilde u$ supplied by
(reg)$^{\sharp}$ of \eqref{4.14:eq-real-backgrounds-a}, with $U^{\tilde u}=e^{i\eta A}$.
Let $\varphi_{\square}$ be a cut-off function with values in $[0,1]$, interpolating
between $\square^{3}$ and $\square^{4}$ and equal to $1$ on $\square^{3}$. The path
\begin{equation*}
\mathsf{u}_z=\bigl(e^{iz\eta\varphi_{\square}A}\mathbf{1},\,0\bigr),\qquad
|z|<r^{\sharp}:=\frac{a_1^{\mathrm{el}}}{2a^{\sharp}},
\end{equation*}
lies in the fibre of the sheet; here $a_1^{\mathrm{el}}$ is the analyticity radius of
the elementary sheets. Indeed, $a^{\sharp}$ is the upper bound that (reg)$^{\sharp}$
places on $A$ on $\square^{4}\cap\Omega_0$; since $0\le\varphi_{\square}\le1$, the
field $z\varphi_{\square}A$ is bounded by $|z|\,a^{\sharp}<a_1^{\mathrm{el}}/2$ for
$|z|<r^{\sharp}$.
\item[($\delta$)] Write $\Gamma_{\tilde u}$ for the operator implementing the gauge
transformation $\tilde u$ on bond functions. By $(\mathrm{Sh}4)$,
\begin{equation*}
\widetilde{T}^{\square}(U,0)=\Gamma_{\tilde u}^{-1}\,
\widetilde{T}^{\square}(U^{\tilde u},0)\,\Gamma_{\tilde u} .
\end{equation*}
By $(\mathrm{Sh}2)$, since the support $[\omega]$ of each walk counted in
$\widetilde{T}^{\square}$ lies in $\square^{3}$, where $\varphi_{\square}=1$,
\begin{equation*}
\widetilde{T}^{\square}(U^{\tilde u},0)=\widetilde{T}^{\square}(\mathsf{u}_1) .
\end{equation*}
Together,
$\widetilde{T}^{\square}(U,0)=\Gamma_{\tilde u}^{-1}\,\widetilde{T}^{\square}
(\mathsf{u}_1)\,\Gamma_{\tilde u}$.
\item[($\varepsilon$)] The Schwarz lemma, applied to the function
$z\mapsto\widetilde{T}^{\square}(\mathsf{u}_z)$ on the disc $|z|<r^{\sharp}$, bounds
the difference between $\widetilde{T}^{\square}(\mathsf{u}_1)$ and its value at the
base point $\mathsf{u}_0$ by the share of $m^{\mathrm{flat}}$ that the assembly in
($\eta$) allots to one box, as soon as
\begin{equation*}
r^{\sharp}\ge\frac{12\,e\,B_T\,S_{ms}(\delta_T)}{m^{\mathrm{flat}}} ,
\end{equation*}
that is, by the definition of $r^{\sharp}$, as soon as
\begin{equation*}
a^{\sharp}\le\frac{a_1^{\mathrm{el}}\,m^{\mathrm{flat}}}
{24\,e\,B_T\,S_{ms}(\delta_T)} .
\end{equation*}
This holds, with a factor $L^2$ to spare, by the ceiling (cl-reg)$^{\sharp}$ of
\eqref{4.14:eq-block-size-conditions-b}. Steps ($\gamma$) and ($\varepsilon$) are the
only place in the proof of Theorem~\ref{4.14:thm-operator-theorem} where
(reg)$^{\sharp}$ and the constant $a^{\sharp}$ enter, and the ceiling
(cl-reg)$^{\sharp}$ is used only here.
\item[($\zeta$)] At $z=0$ the path is at the base point $\mathsf{u}_0=(\mathbf{1},0)$.
This is the base at which the bound $\widetilde{T}(\mathbf{1})\ge m^{\mathrm{flat}}$ of
the positivity lemma at the identity is taken.
\item[($\eta$)] Assembling the box estimates with the IMS decomposition of ($\alpha$)
gives $\widetilde{T}(U)\ge\tfrac12m^{\mathrm{flat}}$.
\end{itemize}

(t3) The sectoriality lemma is applied in the variable $A'$. It gives
\begin{equation*}
\operatorname{Re}\widetilde{T}\ge\tfrac14m^{\mathrm{flat}}\quad\text{for }
\|A'\|<\varepsilon_Ta_1^{\mathrm{el}},\qquad
\varepsilon_T:=\frac{m_0}{4\,e\,B_T\,S_{ms}(\delta_T)} .
\end{equation*}
This holds on $\mathcal{U}(1)$, since $a_1^{\oplus}\le\varepsilon_Ta_1^{\mathrm{el}}$
by the definition of $a_1^{\oplus}$ recalled in step (vi).

(t5) The inversion lemma is applied in $\mathfrak{F}^{ms}$, with index set
$\mathfrak{X}_1:=\mathfrak{B}$ and coercivity constant $\tfrac14m^{\mathrm{flat}}$. It
requires the floor
\begin{equation*}
\mathsf{F}_{ms}\bigl(B_T,\delta_T,\tfrac14m^{\mathrm{flat}};N_d,\lambda^{\sharp},
w^{\sharp}\bigr),
\end{equation*}
which is part of (fl-$\sharp$)$'$. It gives
\begin{equation*}
\widetilde{T}^{-1}\in\mathcal{S}_{ms}\bigl[\rho_T;\,48N_d/m^{\mathrm{flat}},\,
\mu_T/4;\,1\bigr];
\end{equation*}
here $\mu_T$ is the rate entering the output of the inversion lemma, as fixed in the
walk-sheet lemma.

(t6) Part (F2) of the flattening lemma gives
\begin{equation*}
T^{-1}\in\mathcal{O}^{-2}\bigl[\rho_T;B_{T^{-1}},\delta_{T^{-1}};1\bigr].
\end{equation*}

(vi) Assembly. The assembly of the form proceeds in five stages. The weights and tails
of the assembly are chosen so that all tails lie in $[-5,5]$, which is contained in the
admissible tail range $[-\mathsf{e},\mathsf{e}]$ of the assembly, where
$\mathsf{e}=8+d$.
The alphabet $\mathfrak{L}$ collects the operators of (i)--(v), with one radius
$\rho_T$, one constant $B_{\mathfrak{L}}$ and one rate $\delta_{\mathfrak{L}}$. Its
vertex operators are $M_J$, with vertex constant $\nu_J=2c_{\mathfrak{g}}a_0^{\oplus}$
for a constant $c_{\mathfrak{g}}$ of the assembly, together with
$\Delta^{(2)}$ and $\mathcal{J}$.

The words of $W_{\pi}$ carry $\tfrac12\mathfrak{S}$, as in convention
\eqref{4.14:eq-free-current-notation-c}. By part (b) of
Lemma~\ref{4.14:lem-symmetry-covariance}, this is the same symmetric operator, and the
words are the same up to the scalar $\tfrac12$. Indeed, for a generic vector $v$ the
two factorised words of $W_{\pi}$, written as triples (left factor, vertex, right
factor), are
\begin{equation*}
(\mathbf{1},\,M_J,\,\mathfrak{S}v),\qquad (\tfrac12\mathfrak{S}v,\,M_J,\,Dv),
\end{equation*}
of orders $0$. Their weights are not larger than those of the same words carrying
$\mathfrak{S}$.

The five stages, numbered $1$ to $5$, are the following.
\begin{enumerate}
\item[1.] The operators $G_{\pi}$, $T_{\pi}^{-1}$ and the remaining operator of this
stage.
\item[2.] The branching lemma, applied to $\Delta^{(2)}[\mathfrak{j}_0]$, the operator
of (ii) with the argument $\mathfrak{j}_0$ inserted at this stage in the place of
$\mathsf{f}$.
\item[3.] The operators $G_1$, $T_1^{-1}$ and the remaining operator of this stage.
\item[4.] The branching lemma, applied to $\mathcal{J}[\mathfrak{j}]$, the operator of
(iii) with the argument $\mathfrak{j}$ inserted at this stage in the place of
$\mathsf{f}$.
\item[5.] The form itself,
\begin{equation*}
\mathsf{Q}=C_{\flat}^{T}\bigl(P_{\Lambda}T_1^{-1}P_{\Lambda}-a+\mathcal{J}[\mathfrak{j}]
+\mathsf{G}\bigr)C_{\flat},
\end{equation*}
treated by the summation lemma, by the product lemma with the unit-lattice operators of
(iii), and by part (F3) of the flattening lemma.
\end{enumerate}
Stage 5 yields \eqref{4.14:eq-form-walk-expansion-1} with
$\delta_{\mathsf{Q}}:=\delta_{\mathfrak{L}}/32$. In particular, on $\mathcal{U}(1)$ the
series (q1), (q2) converge absolutely; Stages 1 and 3 give
$G_{\pi}=(G_0^{-1}+W_{\pi})^{-1}$ and $G_1=(G_0^{-1}+W_1)^{-1}$ together with
$T_{\pi}^{-1}$ and $T_1^{-1}$; and (v) gives $T^{-1}$. This is (a).

The radii are those of \eqref{4.14:eq-block-size-conditions-d}. The radius
$a_1^{\oplus}$ is the minimum of $\varepsilon_Ta_1^{\mathrm{el}}$ and of the three radii
of Lemma~\ref{4.14:lem-unit-lattice-operators}. The radius $a_0^{\oplus}$ is the largest
number satisfying the four smallness conditions displayed there; these conditions are
the only place where the size of $c_{\mathfrak{g}}$ is constrained.

The sheet properties entering (b), and the properties in (c) and (d), are obtained as
follows.
\begin{itemize}
\item Symmetry and reality hold operator by operator, hence for $\mathsf{Q}$.
\item $(\mathrm{Sh}2)$ follows from the Weierstrass theorem on uniform limits of
analytic functions.
\item $(\mathrm{Sh}4)$ holds operator by operator.
\item $(\mathrm{Sh}5)$ holds operator by operator, with respect to the enlarged lists:
for the elementary walks by part (iv) of Lemma~\ref{4.14:lem-reading-hulls}; for the
words of $T^{-1}$ produced by the inversion lemma by part (iii) of
Lemma~\ref{4.14:lem-reading-hulls}; and for the local and unit-lattice operators
because their supports are determined by the geometry of the nest.
\item The bound (d) is the Schur bound for a sheet in the class
\eqref{4.14:eq-form-walk-expansion-1}.
\end{itemize}

(vii) The base point. First, $(\mathbf{1},0)\in\mathfrak{G}$. By the identification of
the form on the graph, $\mathsf{Q}(\mathbf{1},0)$ is the form $\mathcal{B}$ of the
unit-lattice averaging analysis evaluated at the data
\begin{equation*}
\bigl(\Omega_t^{(k)},\,(\Omega_j)_{j\le k},\,1,\,M_1,\,h\,C^{(2)}\bigr),
\end{equation*}
in which $C^{(2)}$ is the operator so written in
Lemma~\ref{4.14:lem-unit-lattice-operators}, and the fifth entry is the one formed from
$C^{(2)}$ and written $h\,C^{(2)}$ in the unit-lattice averaging analysis.
Part (b) of the coercivity proposition for this form gives
\begin{equation*}
\mathcal{B}\ge\gamma_0^{X}/c_{\Phi}^{2}=:\gamma_0 ,
\end{equation*}
where $\gamma_0^{X}$ and $c_{\Phi}$ are the constants of that proposition.
The proposition gives this bound under four conditions. The first is the pair of
conditions (A1), (A2) of \eqref{4.14:eq-prepared-data-admissible-a}, the conditions of
Proposition~\ref{4.14:prop-prepared-data-admissible} (printed without proof in
\cite[pp.~178--179]{B-CMP122a}; cited as printed). The second holds trivially at
$U=\mathbf{1}$. The third is the proposition's remaining condition on the data
displayed above; it is met through the floors (fl-99) and (fl-$M$)$'$(i)--(ii) of
\eqref{4.14:eq-block-size-conditions-a}, which, together with (A1) and (A2), are the
hypotheses of Theorem~\ref{4.14:thm-operator-theorem} that this step uses. The fourth,
the criticality condition, holds trivially at $U=\mathbf{1}$.

What the proposition uses at $(\mathbf{1},0)$ is the following. It uses the expansion of
the elementary operators at $(\mathbf{1},0)$, where every $J$-vertex vanishes. It also
uses the positivity input at $U=\mathbf{1}$, which is \cite[Proposition~2.7]{B-CMP96} on
the extended sequence of domains.

A second proof of the same lower bound, at level $j=k$ with the region $\Omega_t$, uses
Jensen's inequality and a two-block Poincar\'e inequality and is independent of the
analysis of \cite{B-CMP99b} on which the first proof rests; it is not used here. The
first argument gives (e).

Hypotheses used. The hypotheses of Theorem~\ref{4.14:thm-operator-theorem} enter as
follows.
\begin{itemize}
\item Condition (A1) of \eqref{4.14:eq-prepared-data-admissible-a} is used in three
ways. It is used as \cite[(2.1)--(2.2)]{B-CMP96}, at the pair of block sizes
$(M/LM_1,M_1)$ and at the pair $(M/LM_1^{\sharp},M_1^{\sharp})$, one pair for each of
the two covers. It is also used in \cite[Lemma~2.1]{B-CMP96}, applied with $M/L$ in
place of the product written $RM$ there. Finally, together with (fl-$M$)$'$(i), it
gives property (Sc): a box $\square^{5}$ meets at most two, adjacent, layers.
\item Conditions (A1), (A2) are used in step (vii), as the hypotheses of the analysis of
\cite{B-CMP99b} that enters the first proof there and as the first hypothesis of the
coercivity proposition, at $\Lambda:=\Omega_t^{(k)}$.
\item The floors (fl-99), (fl-$\sharp$)$'$ and (fl-$M$)$'$(i)--(ii) are used.
\item The floor (fl-$L$) is used: $L\ge8$ in the complex localisation proposition, and
$L$ odd, so that the blocks are centred and the pivot $b_0(c)$ is the middle bond of a
coarse bond $c$.
\item The ceiling (cl-99) is used.
\item The members of the ceiling (cl-reg)$^{+}$ read by
Lemma~\ref{4.14:lem-unit-lattice-operators} are used.
\item Condition (reg)$^{\sharp}$ of \eqref{4.14:eq-real-backgrounds-a} and the ceiling
(cl-reg)$^{\sharp}$ are used only in steps ($\gamma$) and ($\varepsilon$) of (t2).
\item Condition (u1) of \eqref{4.14:eq-real-backgrounds-a} is used: the elementary sheets
and the tower bounds of \cite{B-CMP99b} are taken on the class of backgrounds fixed by
(u1) in Definition~\ref{4.14:def-real-backgrounds}.
\end{itemize}

Applicability of the cited proofs. The proofs of the lemmas cited in (i)--(vi) apply
under the conventions
\eqref{4.14:eq-free-current-notation-a}--\eqref{4.14:eq-free-current-notation-c}. These
conventions enter the proofs only in four ways:
\begin{itemize}
\item through the diameters in $(\mathrm{Sh}1)$, namely $\rho_0$ and $\rho_T$, and the
floors required at these values are among the hypotheses;
\item through the constant $B_{\mathfrak{L}}$ of the alphabet and the constant entering
the smallness conditions on $a_0^{\oplus}$ in \eqref{4.14:eq-block-size-conditions-d},
and only through their values;
\item through the scalar $\tfrac12$;
\item through the $(d,L)$-numbers of Lemma~\ref{4.14:lem-unit-lattice-operators}.
\end{itemize}
No inequality changes direction; in particular condition $(\tau1)$ of step (t1) holds at
the new values exactly as before.

The identity $\widetilde{T}^{\square}(U^{\tilde u},0)=\widetilde{T}^{\square}
(\mathsf{u}_1)$, used in step ($\delta$) of (t2), needs only that the support $[\omega]$
of each walk $\omega$ counted in $\widetilde{T}^{\square}$ lies in
$\widehat{\square}=\square^{3}$, where $\varphi_{\square}=1$. This requirement is
unaffected by the conventions. When $[\omega]$ grows, fewer walks qualify, and the tail
is the one controlled by (b3) at the value $\rho_0$.
\end{proof}

\begin{lemma}[Positivity, sector bound and inversion of the block]\label{4.14:prf-positivity-inversion}
Assume the hypotheses of the operator theorem at free current, Theorem~\ref{4.14:thm-operator-theorem}.
Let $\mathbb{B}\subset\mathbb{B}^{+}$ be arbitrary, and let $\mathsf{Q}$ and its $\mathbb{B}\mathbb{B}$-block
$\mathsf{A}=\mathsf{Q}_{\mathbb{B}\mathbb{B}}$ be as there. Let $\gamma_0$, $\gamma_1$, $B_{\mathsf{Q}}$,
$\delta_{\mathsf{Q}}$, $\rho_{\mathsf{Q}}$ be the constants of the walk expansion of the form and base-point bound
(Lemma~\ref{4.14:prf-form-walk-expansion}), let $S_u(\cdot)$ be the summation constant of the unit frame
$\mathfrak{F}^{u}$, and put
\begin{equation*}
m_1:=\frac{\gamma_0}{4},\qquad \varepsilon_1:=\frac{\gamma_0}{8\gamma_1},\qquad
B_0:=\frac{12\cdot 2^{d}}{m_1},
\end{equation*}
\begin{equation*}
\mu_1:=\min\Bigl\{\frac{\delta_{\mathsf{Q}}}{4},\ \frac{m_1}{8e\,B_{\mathsf{Q}}\,S_u(\delta_{\mathsf{Q}}/2)}\Bigr\},
\qquad \delta_0:=\frac{\mu_1}{4}.
\end{equation*}
Then the following hold.
\begin{enumerate}
\item[(a)] $\mathsf{Q}(U,0)\ge \gamma_0/2$ for every $U\in\mathcal{U}_{\mathbb{R}}$.
\item[(b)] $\operatorname{Re}\mathsf{Q}\ge m_1$ on $\mathcal{U}(\varepsilon_1)$.
\item[(c)] For every $x\ge 0$ the operator $x+\mathsf{A}$ is invertible on $\mathcal{U}(\varepsilon_1)$, and
\begin{equation*}
(x+\mathsf{A})^{-1}\in
\mathcal{S}_{u}\Bigl[\frac{M}{4};\ \frac{B_0\,m_1}{m_1+2x},\ \delta_0;\ \varepsilon_1\Bigr],
\end{equation*}
with walks and domains independent of $x$. In particular $\mathsf{C}=\mathsf{A}^{-1}$ exists on
$\mathcal{U}(\varepsilon_1)$.
\end{enumerate}
Assertion (a) is the first clause of clause (B), and (c) is the second half of clause (D), of
Theorem~\ref{4.14:thm-operator-theorem} (see \eqref{4.14:eq-operator-theorem-a}).
\end{lemma}

\begin{proof}
Throughout, $S_u(\cdot)$ is the summation constant of the unit frame $\mathfrak{F}^{u}$, and
$\mathcal{S}_{u}[\rho;B,\delta;r]$ is the sheet format of that frame (domain diameter $\rho$, constant $B$,
decay rate $\delta$, analyticity radius $r$). By Lemma~\ref{4.14:prf-form-walk-expansion} we have
$\mathsf{Q}\in\mathcal{S}_{u}[\rho_{\mathsf{Q}};B_{\mathsf{Q}},\delta_{\mathsf{Q}};1]$ with
$\rho_{\mathsf{Q}}=\rho_T=8dL^2M_1^{\sharp}$; $\mathsf{Q}$ is complex symmetric, and real symmetric at real
points; it has the covariance property $(\mathrm{Sh}4)$ and the reading property $(\mathrm{Sh}2)$; by the Schur
bound, $\|\mathsf{Q}\|\le eB_{\mathsf{Q}}S_u(\delta_{\mathsf{Q}})=\gamma_1$ on $\mathcal{U}(1)$; and
$\mathsf{Q}(1,0)\ge\gamma_0$. Evaluating both bounds on a unit vector of the space on which $\mathsf{Q}$ acts
gives $\gamma_0\le\gamma_1$; this is used below to see that $r_{\mathrm{reg}}>1$.

\emph{Step 1: positivity at real backgrounds.} Put $\ell:=M/12d$. We use the unit box system: the cubes
$\mathsf{b}$ of side $\ell$ centred at the points of $\ell\mathbb{Z}^d$ for which $\mathsf{b}''$, the
concentric enlargement of $\mathsf{b}$ that selects the boxes of the unit box system, contains a bond
of $\mathbb{B}^{+}$, with cut-offs $h_{\mathsf{b}}$ that are products of one-dimensional cut-offs and satisfy
$\sum_{\mathsf{b}} h_{\mathsf{b}}^2=1$, and $1''$ the projection onto the index set $\mathfrak{X}(\mathsf{b}'')$,
which carries the support of $h_{\mathsf{b}}$. Its four properties, which we label (b1)--(b4), hold here
with overlap number $2^d$ in (b1),
Lipschitz constant $3d/\ell=36d^2/M$ in (b2), width $\ell/4=M/48d$ in (b3), and in (b4) the convex hull
$\mathsf{b}^{+\ell}$, the concentric cube of side $3\ell$, with $\rho_{\mathsf{b}}=3d\ell=M/4$.

These require the following. Property (b3) at width $\ell/4$ requires $\rho_{\mathsf{Q}}\le\ell/12$, that is,
\begin{equation*}
M\ \ge\ 144\,d\,\rho_{\mathsf{Q}}\ =\ 144\cdot 8\,d^2L^2M_1^{\sharp}.
\end{equation*}
This is the first member of item (iii) of the floor on $M$ in \eqref{4.14:eq-block-size-conditions-a}. The
same item gives $\ell\ge 6$ and $M/\ell\in\mathbb{N}$, which the construction of the unit box system requires.

Property (b4) requires $\mathsf{b}^{+\ell}\subset\Omega_k$. This follows from condition (A2) of
\eqref{4.14:eq-prepared-data-admissible-a}, a hypothesis of Theorem~\ref{4.14:thm-operator-theorem}. That
condition is part of the statement of Proposition~\ref{4.14:prop-prepared-data-admissible}, which is stated in
\cite[pp.~178--179]{B-CMP122a} without proof. Condition (A2) states that
$\operatorname{dist}_\infty(\Omega_t,\Omega_k^c)\ge 6M$. Since $\mathsf{b}$ holds a bond of $\mathbb{B}^{+}$, the
cube $(\mathsf{b}^{+\ell})^{+\ell}$, whose side is $5\ell=5M/12d\le M$, lies within distance $5\ell+1<6M$ of
$\Omega_t$. Hence it lies inside $\Omega_k$, and so does $\mathsf{b}^{+\ell}$. For the same reason there is a
cube $Q\supset(\mathsf{b}^{+\ell})^{+\ell}$ of side $M$, contained in $\Omega_k$ and holding a bond of
$\mathbb{B}^{+}$.

Fix $U\in\mathcal{U}_{\mathbb{R}}$ (Definition~\ref{4.14:def-real-backgrounds}). We apply the positivity lemma
for sheets in the unit frame with the operator $\mathsf{Q}$, the lower bound $\gamma_0$ and the unit box system
above. Its floor condition is
\begin{equation*}
\mathsf{F}^{P}_{u}\bigl(B_{\mathsf{Q}},\delta_{\mathsf{Q}},\gamma_0;\,2^d,\,36d^2/M,\,M/48d\bigr),
\end{equation*}
which is item (v) of the floor on $M$ in \eqref{4.14:eq-block-size-conditions-a}. The argument of that lemma, in
the present instance, runs as follows. It proves, in the present setting, the positivity indicated in
\cite{B-CMP99b} by the sentence ``Localizing the operators in $\Delta_k$ and using the
methods of Sect.~B''.

Put $\mathsf{D}:=\tfrac12\sum_{\mathsf{b}}[h_{\mathsf{b}},[h_{\mathsf{b}},\mathsf{Q}]]$. Expanding the double
commutator and using $\sum_{\mathsf{b}} h_{\mathsf{b}}^2=1$ gives the identity
\begin{equation}\label{4.14:eq-positivity-inversion-1}
\mathsf{Q}=\sum_{\mathsf{b}} h_{\mathsf{b}}\,\mathsf{Q}\,h_{\mathsf{b}}+\mathsf{D}.
\end{equation}
By (b1), (b2) and the sheet bound on $\mathsf{Q}$,
\begin{equation*}
\|\mathsf{D}\|\ \le\ 2^d\Bigl(\frac{3d}{\ell}\Bigr)^2 eB_{\mathsf{Q}}S_u(\delta_{\mathsf{Q}})\ \le\ \frac{\gamma_0}{6},
\end{equation*}
the last inequality by the floor condition $\mathsf{F}^{P}_{u}$.

For each box let
$\mathsf{Q}^{\mathsf{b}}:=1''\sum\{\mathsf{Q}_\omega:[\omega]\subset\mathsf{b}^{+\ell}\}\,1''$ be the
partial walk sum over walks whose domain lies in $\mathsf{b}^{+\ell}$. Let
$\mathsf{E}_{\mathsf{b}}:=1''\mathsf{Q}1''-\mathsf{Q}^{\mathsf{b}}$ be the tail. Then
$h_{\mathsf{b}}\mathsf{Q}h_{\mathsf{b}}=h_{\mathsf{b}}(\mathsf{Q}^{\mathsf{b}}+\mathsf{E}_{\mathsf{b}})h_{\mathsf{b}}$.
By (b3),
\begin{equation*}
\|\mathsf{E}_{\mathsf{b}}\|\ \le\ eB_{\mathsf{Q}}\,e^{-\delta_{\mathsf{Q}}\ell/8}\,S_u(\delta_{\mathsf{Q}}/2)
\ \le\ \frac{\gamma_0}{12},
\end{equation*}
the last inequality by the first line of $\mathsf{F}^{P}_{u}$.

Next we use the regularity condition of Definition~\ref{4.14:def-real-backgrounds} on the cube $Q$. This is
exactly a cube of side $M$ inside $\Omega_k$ holding a bond of $\mathbb{B}^{+}$. It gives a $G$-valued $u$ on $Q$
with $U^u=e^{i\eta A}$ and $|A|,|\nabla A|<a_Q$. Let $\tilde u:=u$ on $Q$, and let $\tilde u$ equal the identity
of $G$ off $Q$.
Let $\chi$ equal $1$ on $\mathsf{b}^{+\ell}$ and $0$ off $(\mathsf{b}^{+\ell})^{+\ell}$, with Lipschitz constant
$1/\ell$. Since $(\mathsf{b}^{+\ell})^{+\ell}\subset Q$, $\chi A$ is defined on the whole lattice. Let
$a_1^{\oplus}$ be the radius fixed in Lemma~\ref{4.14:prf-form-walk-expansion} that bounds, in the direction
of the background potential, the points of the ball $\mathcal{U}(1)$. Put
$\mathsf{u}_z:=(e^{iz\eta\chi A},0)$ and $r_{\mathrm{reg}}:=a_1^{\oplus}/(2a_Q)$. By the bounds on $A$ and
$\nabla A$ and the Lipschitz bound on $\chi$, we have $\mathsf{u}_z\in\mathcal{U}(1)$ for $|z|<r_{\mathrm{reg}}$.

The class $\mathcal{U}_{\mathbb{R}}$ is gauge invariant, so $U^{\tilde u}\in\mathcal{U}_{\mathbb{R}}$, and the
covariance $(\mathrm{Sh}4)$ gives
\begin{equation*}
\mathsf{Q}^{\mathsf{b}}(U,0)=\Gamma(\tilde u)^{-1}\,\mathsf{Q}^{\mathsf{b}}(U^{\tilde u},0)\,\Gamma(\tilde u).
\end{equation*}
Here
$\Gamma(\tilde u)$ is the action of $\tilde u$ on the bond space, which is unitary. Every walk in
$\mathsf{Q}^{\mathsf{b}}$ has $[\omega]\subset\mathsf{b}^{+\ell}$. On $\mathsf{b}^{+\ell}$ we have $\chi=1$, so
$U^{\tilde u}$ and the configuration of $\mathsf{u}_1$ agree there. By the reading property
$(\mathrm{Sh}2)$, therefore,
\begin{equation}\label{4.14:eq-positivity-inversion-2}
\mathsf{Q}^{\mathsf{b}}(U,0)=\Gamma(\tilde u)^{-1}\,\mathsf{Q}^{\mathsf{b}}(\mathsf{u}_1)\,\Gamma(\tilde u).
\end{equation}

The walk terms of $\mathsf{Q}$ are analytic on $\mathcal{U}(1)$ and their series converges uniformly there, so
$z\mapsto\mathsf{Q}^{\mathsf{b}}(\mathsf{u}_z)$ is analytic on $|z|<r_{\mathrm{reg}}$. By the Schur bound it is
bounded there by $eB_{\mathsf{Q}}S_u(\delta_{\mathsf{Q}})=\gamma_1$. The first member of the ceiling on the
regularity constants $a_Q$ in \eqref{4.14:eq-block-size-conditions-b} reads
\begin{equation*}
a_Q\ \le\ \frac{a_1^{\oplus}\gamma_0}{24eB_{\mathsf{Q}}S_u(\delta_{\mathsf{Q}})}=\frac{a_1^{\oplus}\gamma_0}{24\gamma_1}.
\end{equation*}
It gives $r_{\mathrm{reg}}\ge 12\gamma_1/\gamma_0\ge 12>1$.

For unit vectors $v,v'$, the function
\begin{equation*}
g(z):=\bigl\langle v,\bigl(\mathsf{Q}^{\mathsf{b}}(\mathsf{u}_z)-\mathsf{Q}^{\mathsf{b}}(\mathsf{u}_0)\bigr)v'\bigr\rangle
\end{equation*}
is analytic on
$|z|<r_{\mathrm{reg}}$, vanishes at $0$ and is bounded by $2\gamma_1$. The Schwarz lemma therefore gives
$|g(1)|\le 2\gamma_1/r_{\mathrm{reg}}$. Hence
\begin{equation*}
\|\mathsf{Q}^{\mathsf{b}}(\mathsf{u}_1)-\mathsf{Q}^{\mathsf{b}}(\mathsf{u}_0)\|
\ \le\ \frac{2}{r_{\mathrm{reg}}}\,eB_{\mathsf{Q}}S_u(\delta_{\mathsf{Q}})
=\frac{4a_Q\gamma_1}{a_1^{\oplus}}\ \le\ \frac{\gamma_0}{6}.
\end{equation*}

At the base point $\mathsf{u}_0=(1,0)$ we have
$\mathsf{Q}^{\mathsf{b}}(\mathsf{u}_0)=1''\mathsf{Q}(1,0)1''-\mathsf{E}_{\mathsf{b}}(1,0)$, which is
$\ge\gamma_0-\gamma_0/12$ on the range of $1''$. The point $\mathsf{u}_1$ is real, so
$\mathsf{Q}^{\mathsf{b}}(\mathsf{u}_1)$ is real symmetric and
$\mathsf{Q}^{\mathsf{b}}(\mathsf{u}_1)\ge\gamma_0-\gamma_0/12-\gamma_0/6$ there. By
\eqref{4.14:eq-positivity-inversion-2} the same bound holds for $\mathsf{Q}^{\mathsf{b}}(U,0)$. Adding the tail,
\begin{equation*}
h_{\mathsf{b}}\mathsf{Q}(U,0)h_{\mathsf{b}}\ \ge\
\Bigl(\gamma_0-\frac{\gamma_0}{12}-\frac{\gamma_0}{6}-\frac{\gamma_0}{12}\Bigr)h_{\mathsf{b}}^2
=\frac{2\gamma_0}{3}\,h_{\mathsf{b}}^2 .
\end{equation*}
Summing over the boxes with $\sum_{\mathsf{b}} h_{\mathsf{b}}^2=1$ and subtracting $\|\mathsf{D}\|\le\gamma_0/6$ in
\eqref{4.14:eq-positivity-inversion-1} gives $\mathsf{Q}(U,0)\ge\gamma_0/2$. This is (a).

\emph{Step 2: the sector bound.} We apply the sector lemma with the operator $\mathsf{Q}$ and lower bound
$\gamma_0/2$ at real points (Step 1). Its inputs are the analyticity of $\mathsf{Q}$ on the complex fibre and
the matrix-element bound
$|\mathsf{Q}(\iota,\iota')|\le eB_{\mathsf{Q}}e^{-\delta_{\mathsf{Q}}\,d(\iota,\iota')}$, with $d$ the block
distance of the unit frame between the bonds $\iota,\iota'$, which
follows from the sheet bound and the distance property of the unit frame. The lemma applies the
Schwarz lemma to matrix elements and the Schur bound, and uses nothing specific to the frame. Its radius is
\begin{equation*}
\varepsilon=\frac{\gamma_0/2}{4eB_{\mathsf{Q}}S_u(\delta_{\mathsf{Q}})}=\frac{\gamma_0}{8\gamma_1}=\varepsilon_1,
\end{equation*}
and its conclusion is $\operatorname{Re}\mathsf{Q}\ge\gamma_0/4=m_1$ on $\mathcal{U}(\varepsilon_1)$. This is (b).
No further hypothesis is used.

\emph{Step 3: inversion of the block.} A block of a sheet is a sheet with the same constants, so
$\mathsf{A}=\mathsf{Q}_{\mathbb{B}\mathbb{B}}\in\mathcal{S}_{u}[\rho_{\mathsf{Q}};B_{\mathsf{Q}},\delta_{\mathsf{Q}};1]$.
For $v$ supported on $\mathbb{B}$,
$\operatorname{Re}\langle v,\mathsf{A}v\rangle=\operatorname{Re}\langle v,\mathsf{Q}v\rangle\ge m_1\|v\|^2$ by
(b). Hence $\operatorname{Re}\mathsf{A}\ge m_1$ on $\mathcal{U}(\varepsilon_1)$.

We apply the inversion lemma for sheets in the unit frame to $\mathsf{A}$, with index set $\mathbb{B}$ and the
unit box system of Step 1. Properties (b1)--(b4) hold at $\rho_{\mathsf{Q}}$ by item (iii) of the floor on $M$
and by condition (A2), exactly as in Step 1. The floor condition of the inversion lemma is
\begin{equation*}
\mathsf{F}_{u}\bigl(B_{\mathsf{Q}},\delta_{\mathsf{Q}},m_1;\,2^d,\,36d^2/M,\,M/48d\bigr),
\end{equation*}
which is item (iv) of the floor on $M$ in \eqref{4.14:eq-block-size-conditions-a}. The lemma gives, for every
$x\ge0$,
\begin{equation*}
(x+\mathsf{A})^{-1}\in\mathcal{S}_{u}\Bigl[\max\Bigl\{\rho_{\mathsf{Q}},\frac M4\Bigr\};\,
\frac{12\cdot 2^d}{m_1+2x},\,\frac{\mu_1}{4};\,\varepsilon_1\Bigr],
\end{equation*}
with walks and domains independent of $x$. Since $\rho_{\mathsf{Q}}\le M/144d\le M/4$, the maximum equals $M/4$.
Moreover $12\cdot2^d/(m_1+2x)=B_0m_1/(m_1+2x)$ and $\mu_1/4=\delta_0$. This is (c); at $x=0$ it gives the
existence of $\mathsf{C}=\mathsf{A}^{-1}$.

The set $\mathbb{B}$ is arbitrary. Nothing in the inversion lemma depends on the shape of its index set, and
irregular subsets of $\mathbb{B}^{+}$ are included.

In none of the three steps do the enlarged domain lists of Theorem~\ref{4.14:thm-operator-theorem} enter. The
unit frame has a single level, its grid and cut-offs carry no labels, and $\sum_{\mathsf{b}} h_{\mathsf{b}}^2=1$
holds without normalisation. Since $\mathsf{b}^{+\ell}\subset\Omega_k$ by (A2), no reading hull
(Lemma~\ref{4.14:lem-reading-hulls}) is needed. The gauge used is the one given by the regularity condition on
cubes of side $M$, not the boxwise regularity condition on the second cover.
\end{proof}

\begin{lemma}[Block identities, resummed expansions and uniformity]\label{4.14:prf-block-identities}
Assume the hypotheses of the operator theorem at free current,
Theorem~\ref{4.14:thm-operator-theorem}. Let $\mathbb{B}\subset\mathbb{B}^{+}$ be arbitrary and
$\mathbb{B}'=\mathbb{B}^{+}\setminus\mathbb{B}$, and write $\mathsf{A}=\mathsf{Q}_{\mathbb{B}\mathbb{B}}$,
$\mathsf{J}=\mathsf{Q}_{\mathbb{B}\mathbb{B}'}$. Suppose that the following hold; they are
established in the preceding parts of the proof of the theorem,
Lemmas~\ref{4.14:prf-form-walk-expansion} and~\ref{4.14:prf-positivity-inversion}.
\begin{enumerate}
\item[(a)] $\mathsf{Q}\in\mathcal{S}_{u}[\rho_{\mathsf{Q}};B_{\mathsf{Q}},\delta_{\mathsf{Q}};1]$, and
$\|\mathsf{Q}\|\le\gamma_1$ on $\mathcal{U}(1)$. The block $\mathsf{J}=\mathsf{Q}_{\mathbb{B}\mathbb{B}'}$
of this sheet, with $b\in\mathbb{B}$ and $b'\in\mathbb{B}'$, is a sheet with constants
$(B_{\mathsf{J}},\delta_{\mathsf{Q}},1)$. Property $(\mathrm{Sh}5)$ of the terms of $\mathsf{Q}$
holds with respect to lists in which the many-scale factors inside the words of $\mathsf{Q}$ carry
the reading hulls $\square^{\mathsf{r}}$ and the elementary walks carry the enlarged domains
$\widehat{X}$ with enlargement $\max\{5,2L\}$.
\item[(b)] At every base point $(U,0)$ with $U\in\mathcal{U}_{\mathbb{R}}$ the matrix $\mathsf{Q}$ is
real symmetric with spectrum in $[\gamma_0/2,\gamma_1]$.
\item[(c)] On $\mathcal{U}(\varepsilon_1)$ one has $\operatorname{Re}\mathsf{Q}\ge\gamma_0/4$.
\item[(d)] Clause (D) of Theorem~\ref{4.14:thm-operator-theorem} holds.
\end{enumerate}
Then clauses (A), (B), (C), (E), (F), (G), (H), (I) and (K) of
Theorem~\ref{4.14:thm-operator-theorem} hold, together with the covariance statement, the
dependence statement and the Schur-complement statement listed there. Consequently
Theorem~\ref{4.14:thm-operator-theorem} holds, with the covariance, dependence and
Schur-complement statements.
\end{lemma}

\begin{proof}
Clause (A). By (a), $\mathsf{Q}$ is analytic on $\mathcal{U}(1)$; it is complex symmetric, and real
symmetric at real points, by the symmetry lemma for the form $\mathsf{Q}$. The diagonal block
$\mathsf{A}$ is therefore analytic and complex symmetric, and the off-diagonal block $\mathsf{J}$ is
analytic. By (d) at $x=0$ the inverse $\mathsf{C}=\mathsf{A}^{-1}$ exists on
$\mathcal{U}(\varepsilon_1)$, so that $\mathsf{A}\mathsf{C}=1$ on all of $\mathcal{U}(\varepsilon_1)$
by definition. The inverse of an analytic, invertible, complex symmetric
finite matrix is analytic and complex symmetric, so $\mathsf{C}$ is analytic and complex symmetric.

On the graph $\mathfrak{G}=\{J=J(\mathbf{U})\}$ the graph identity for the form identifies
$\mathsf{Q}$ with Ba{\l}aban's operator $C^{\ast}\Delta^{(k)}C$ on $\mathbb{B}^{+}$; here the italic
$C$ and $\Delta^{(k)}$ are in Ba{\l}aban's notation, distinct from the block inverse $\mathsf{C}$. That
identity uses \cite[(3.117$'$), (3.126), (3.129)]{B-CMP99b} on $\mathfrak{G}$ only, takes the object
with the current substituted as defined in \cite{B-CMP109}, and reads it in the frame of
\cite{B-CMP116}.
Let $\langle\cdot,\cdot\rangle$ denote the bilinear pairing, and write $A_k=\varphi\oplus A'$ with
$\varphi$ on $\mathbb{B}$ and $A'$ on $\mathbb{B}'$. The cross terms of
$\tfrac12\langle A_k,\mathsf{Q}A_k\rangle$ are $\tfrac12\langle\varphi,\mathsf{J}A'\rangle$ and
$\tfrac12\langle A',\mathsf{Q}_{\mathbb{B}'\mathbb{B}}\varphi\rangle$. Since $\mathsf{Q}$ is symmetric,
$\mathsf{Q}_{\mathbb{B}'\mathbb{B}}=\mathsf{J}^{T}$, and the two cross terms are equal. Hence on
$\mathcal{U}(\varepsilon_1)$
\begin{equation}\label{4.14:eq-block-identities-1}
\tfrac12\langle A_k,\mathsf{Q}A_k\rangle
=\tfrac12\langle\varphi,\mathsf{A}\varphi\rangle+\langle\varphi,\mathsf{J}A'\rangle
+\tfrac12\langle A',\mathsf{Q}_{\mathbb{B}'\mathbb{B}'}A'\rangle .
\end{equation}
Expand the square below and use $\mathsf{A}\mathsf{C}=1$ together with the symmetry of $\mathsf{A}$
and $\mathsf{C}$. This gives
\begin{equation}\label{4.14:eq-block-identities-2}
\begin{split}
\tfrac12\langle\varphi,\mathsf{A}\varphi\rangle&+\langle\varphi,\mathsf{J}A'\rangle
+\tfrac12\langle A',\mathsf{Q}_{\mathbb{B}'\mathbb{B}'}A'\rangle\\
&=\tfrac12\bigl\langle\varphi+\mathsf{C}\mathsf{J}A',\,\mathsf{A}(\varphi+\mathsf{C}\mathsf{J}A')\bigr\rangle
+\tfrac12\langle A',\mathcal{K}A'\rangle,\\
&\qquad\text{with}\quad
\mathcal{K}:=\mathsf{Q}_{\mathbb{B}'\mathbb{B}'}-\mathsf{J}^{T}\mathsf{C}\mathsf{J}.
\end{split}
\end{equation}
The identities \eqref{4.14:eq-block-identities-1} and \eqref{4.14:eq-block-identities-2} hold on
all of $\mathcal{U}(\varepsilon_1)$. On $\mathfrak{G}\cap\mathcal{U}(\varepsilon_1)$ the left side of
\eqref{4.14:eq-block-identities-1} is Ba{\l}aban's form, by the identification above, so that
$\mathsf{C}$ and $\mathsf{J}$ are the covariance and the source operator of the conditional
integral over the restriction of $A_k$ to $\mathbb{B}$. For $\mathbb{B}=\mathbb{B}^{\ast}$ these are
the operators of \cite[(3.23), (3.25)]{B-CMP119}. For $\mathbb{B}=\mathbb{B}^{K1}$ they are the
operators of the same conditional integral taken over $\mathbb{B}^{K1}$. For arbitrary $\mathbb{B}$
the statement is the linear algebra of \eqref{4.14:eq-block-identities-2}.

Clause (B). At $(U,0)$, $U\in\mathcal{U}_{\mathbb{R}}$, the matrix $\mathsf{Q}$ is real symmetric with
spectrum in $[\gamma_0/2,\gamma_1]$ by (b). For a real vector $v$ supported on $\mathbb{B}$,
$\langle v,\mathsf{A}v\rangle=\langle v,\mathsf{Q}v\rangle$. By the min--max characterisation, the
spectrum of the compression $\mathsf{A}$ therefore also lies in $[\gamma_0/2,\gamma_1]$. Consequently
$\mathsf{C}=\mathsf{A}^{-1}$ is real symmetric with spectrum in
\begin{equation*}
[\gamma_1^{-1},\,2/\gamma_0]=[\gamma_-,\gamma_+].
\end{equation*}
At a real point of $\mathcal{U}(\varepsilon_1)$ the matrix $\mathsf{Q}$ is real symmetric, by the
symmetry lemma for the form $\mathsf{Q}$. Hence $\operatorname{Re}\mathsf{Q}=\mathsf{Q}$, and by (c)
and (a) we get $\gamma_0/4\le\mathsf{Q}\le\gamma_1$. The same min--max argument gives
$\gamma_0/4\le\mathsf{A}\le\gamma_1$, hence
\begin{equation*}
\gamma_-\le\mathsf{C}\le 4/\gamma_0=2\gamma_+ .
\end{equation*}
On all of $\mathcal{U}(\varepsilon_1)$, $\operatorname{Re}\mathsf{Q}\ge\gamma_0/4$ is (c).

Clause (C). At real points $\mathsf{J}$ is a block of a real matrix, hence real. By the symmetry of
$\mathsf{Q}$, $\mathsf{J}^{T}=\mathsf{Q}_{\mathbb{B}'\mathbb{B}}$. Since $\mathsf{J}$ is $\mathsf{Q}$
composed on both sides with coordinate projections of norm at most one,
$\|\mathsf{J}\|\le\|\mathsf{Q}\|\le\gamma_1$ on $\mathcal{U}(1)$, by the norm bound in (a). None of
the notation of \eqref{4.14:eq-free-current-notation-a}--\eqref{4.14:eq-free-current-notation-d}
enters clauses (A)--(C).

Clauses (E), (F), (G). The core lemma states the following. A sheet
$T\in\mathcal{S}_{u}[\rho;B,\delta;r]$ with $\rho\le M/4$ obeys the walk-expansion properties (W1),
(W2), (W3)$'$, (W5)$^{\mathrm{cov}}$ and (W6)$^{\mathrm{geo}}$ listed in clause (E) of
Theorem~\ref{4.14:thm-operator-theorem}, with the constants $(B,\delta,r)$ at every core.
Property (W2) is a theorem there: the expansions at two cores regroup the same series. That lemma
requires that those cubes of $\mathsf{Z}_0^{\sim}$, the collar of the set $\mathsf{Z}_0$ of cores,
which lie within distance $3M$ of a core be level-$k$ $M$-cubes of one aligned partition. This is
supplied by condition (A2) of
\eqref{4.14:eq-prepared-data-admissible-a}.

Clause (E) follows from (D) at $x=0$. There $\mathsf{C}=(0+\mathsf{A})^{-1}$ lies in
$\mathcal{S}_u[M/4;B_0,\delta_0;\varepsilon_1]$, and $\rho=M/4$ is admissible in the core lemma.

Clause (F) follows from the block $\mathsf{J}=\mathsf{Q}_{\mathbb{B}\mathbb{B}'}$ of the sheet of
$\mathsf{Q}$ in (a), with $b\in\mathbb{B}$ and $b'\in\mathbb{B}'$. By (a) this block is a sheet with
constants $(B_{\mathsf{J}},\delta_{\mathsf{Q}},1)$, hence with $(B_{\mathsf{J}},\delta_0,\varepsilon_1)$,
at the same cores. The core lemma applies because $\rho_{\mathsf{Q}}\le M/4$, which follows from the
floor $M\ge 144d\,\rho_{\mathsf{Q}}$ in \eqref{4.14:eq-block-size-conditions-a}.

Clause (G) follows from (D) for $x\ge0$, together with the remark accompanying the inversion lemma for
sheets that complex $x$ is handled through $\operatorname{Re}x$. Each walk of $n$ steps in the
expansion of $(x+\mathsf{A})^{-1}$ carries the factor
\begin{equation}\label{4.14:eq-block-identities-3}
\Bigl(1+\frac{2x}{m_1}\Bigr)^{-(n+1)} .
\end{equation}

The reading hulls of \eqref{4.14:eq-free-current-notation-a} enter the proof at this point, and only
through (W6)$^{\mathrm{geo}}$. The core lemma passes $(\mathrm{Sh}5)$ to its classes termwise.
$(\mathrm{Sh}5)$ for the unit words of the inversion lemma follows factorwise from $(\mathrm{Sh}5)$ for
the terms of $\mathsf{Q}$. By (a), the latter holds with respect to lists in which the many-scale
factors inside the words of $\mathsf{Q}$ (the words of the inversion lemma for $T^{-1}$) carry the
reading hulls $\square^{\mathsf{r}}$, and the elementary walks carry $\widehat{X}$ with enlargement
$\max\{5,2L\}$. This is parts (iii)--(iv) of the lemma on reading hulls,
Lemma~\ref{4.14:lem-reading-hulls}. Accordingly, in (W6)$^{\mathrm{geo}}$ the phrase ``labels
coinciding on an open neighbourhood of $[\nu]$'' refers to a set $[\nu]$ containing those hulls.
This is the reading of clause (E) in which $[\nu]$ lists the reading hulls $\square^{\mathsf{r}}$ of
its boxes.

By parts (ii) and (v) of Lemma~\ref{4.14:lem-reading-hulls}, the $\tau$-diameters of all domains stay
at most $\rho_{\mathsf{Q}}\le M/4$. Hence the smallness of domains in (W1) and the input of the packing
estimate for sheets are unaffected. The distances $d(\omega,b,b')=\tau(\operatorname{dom}\nu,b,b')$ are
computed on the enlarged lists, and enlarging a list lowers $\tau$. Write $d(\nu,\cdot)$ for the
distance $\tau$ computed through the list of $\nu$, which lists the reading hulls of its boxes; this
list contains the list through which $d(\omega,\cdot)$ is computed. Therefore
$d(\nu,\cdot)\le d(\omega,\cdot)$, as the core lemma requires. The notations
\eqref{4.14:eq-free-current-notation-b} and \eqref{4.14:eq-free-current-notation-c} do not enter
this step. Property (W5)$^{\mathrm{cov}}$ follows from $(\mathrm{Sh}4)$,
applied termwise.

Clause (H). Apply the lemma on local balls, Lemma~\ref{4.14:lem-local-balls}, to the core forms of
$\mathsf{C}$ and of $(x+\mathsf{A})^{-1}$ with $r:=\varepsilon_1$, and to that of $\mathsf{J}$ with
$r:=1$. Two of its inputs are that crossing a layer costs at least $M$, and that the cubes of
$\mathfrak{c}^{\flat}$ are grid-aligned at each level; both are condition (A1) of
\eqref{4.14:eq-prepared-data-admissible-a}. It also uses $M>2L$, one of the floors on $M$ in
\eqref{4.14:eq-block-size-conditions-a}.

The cut-off of \eqref{4.14:eq-free-current-notation-d}, at radius $M/2L$, is exactly the cut-off of
that lemma, and this is the only place where it enters. By part $(\alpha)$ of
Lemma~\ref{4.14:lem-local-balls}, $[\nu]^{\approx}\subset[\nu]^{\sim}$. Hence the statement of (H)
with data given on the larger set $[\nu]^{\sim}$ holds a fortiori.

Clause (I). Each constant introduced in Lemmas~\ref{4.14:prf-form-walk-expansion}
and~\ref{4.14:prf-positivity-inversion} and above is of one of two kinds. Either it is a number
depending on $d$ and $L$ alone, or it depends on $M_1$ or on $M_1^{\sharp}$ as displayed in
\eqref{4.14:eq-free-current-notation-a} and in Definition~\ref{4.14:def-block-size-conditions};
for instance $\rho_0$ depends on $M_1$ and
$\rho_T=\rho_{\mathsf{Q}}$ on $M_1^{\sharp}$. No constant depends on $M$, except through
$\ell=M/12d$ inside the floors on $M$. None depends on $k$, $K$, the nest, $\mathbb{B}$,
$\Lambda_t$, $\mathsf{Z}_0$, the volume or $g$. The constant $c_{\mathfrak{g}}$, which appears in the
radius $a_0^{\oplus}$ of \eqref{4.14:eq-block-size-conditions-d}, enters no other constant.

The constants are chosen in the order \eqref{4.14:eq-block-size-conditions-d}. The constants named
in it are those introduced in Lemmas~\ref{4.14:prf-form-walk-expansion}
and~\ref{4.14:prf-positivity-inversion} and in Definition~\ref{4.14:def-block-size-conditions}.
Along the steps of the proof this order reads
\begin{equation}\label{4.14:eq-block-identities-4}
\begin{gathered}
d,\,L\;\to\;M_1\;\to\;B_{\mathrm{el}},\,\delta_{\mathrm{el}},\,a_1^{\mathrm{el}},\,\rho_0,\,
B_T,\,\delta_T,\,m^{\mathrm{flat}},\,r_{\flat}\;\to\;M_1^{\sharp}\\
\to\;\varepsilon_T,\,\mu_T,\,B_{T^{-1}},\,\delta_{T^{-1}},\,\rho_T,\,B_{\mathfrak{L}},\,
\delta_{\mathfrak{L}},\,x,\,a_1^{\oplus},\,a_0^{\oplus},\,B_{\mathsf{Q}},\,\delta_{\mathsf{Q}},\\
\gamma_0,\,\gamma_1,\,\varepsilon_1,\,m_1,\,\mu_1,\,\delta_0,\,B_0,\,\mathsf{c}_1
\;\to\;\kappa_1\;\to\;M\;\to\;\text{layer couplings}\;\to\;\text{radii}\;\to\;g.
\end{gathered}
\end{equation}
In this order, $M_1$ is fixed by its floor in \eqref{4.14:eq-block-size-conditions-a}. The block
size $M_1^{\sharp}$ is fixed by its floor there, which is a floor on $M_1^{\sharp}$ after $M_1$. The
cube side $M$ is fixed by the floors on $M$ there. The layer couplings are fixed by the ceilings of
Definition~\ref{4.14:def-block-size-conditions}. The radii are fixed by the condition on the radii
imposed where the theorem is used. No arrow in \eqref{4.14:eq-block-identities-4} returns: no constant
is used before it has been fixed.

In this order $M_1$ and $M_1^{\sharp}$ are chosen after $d$ and $L$ alone, and before $M$, the layer
couplings, the radii and $g$; they are fixed once and for all as functions of $d$ and $L$ satisfying
their floors in \eqref{4.14:eq-block-size-conditions-a}. Hence every constant that depends on $M_1$
or $M_1^{\sharp}$ is itself a number depending on $d$ and $L$ only. So the
constants $B_0$, $B_{\mathsf{J}}$, $\delta_0$, $\delta_{\mathsf{Q}}$, $\gamma_0$, $\gamma_1$,
$\varepsilon_1$, $a_1^{\oplus}$, $a_0^{\oplus}$ depend on $d$ and $L$ only, and are free of the
quantities listed in clause (I).

Clause (K). This clause concerns the choice $\mathbb{B}:=\mathbb{B}^{K1}$, for which the following
hold.
First, the pivot rows of $C_{\flat}$ have entries in both $\mathbb{B}$ and $\mathbb{B}'$. The
$\mathbb{B}\mathbb{B}'$-block of $C_{\flat}^{T}(-a)C_{\flat}$ is non-zero, where $-a$ denotes one of
the unit-lattice terms of the form. The
$\mathbb{B}\mathbb{B}'$-block of $C_{\flat}^{T}\mathsf{G}C_{\flat}$ vanishes, by the unit-lattice
identity for the term $\mathsf{G}$. Second, the constants are free of $\Lambda_t$, by clause (I).
Hence every constant used later that is built from them is fixed before any choice depending on
$\Lambda_t$ is made, and no circularity arises.

The covariance statement. This is part (d) of the lemma on symmetric cell weights and lattice-symmetry
covariance, Lemma~\ref{4.14:lem-symmetry-covariance}. Parts (a)--(c) of that lemma are what make every
object occurring in the definition of a term natural under each $r\in\mathfrak{r}$. These are the
symmetric cell weights $w_z$ of \eqref{4.14:eq-free-current-notation-b}, the symmetrising letter
$\tfrac12\mathfrak{S}$ of \eqref{4.14:eq-free-current-notation-c}, and the parities of part (c).

The dependence statement. This is the statement that a word of $\mathsf{C}^{K1}$, the operator
$\mathsf{C}$ at $\mathbb{B}=\mathbb{B}^{K1}$, that depends on $\Lambda_t$ has a domain containing a
point of $\Lambda_t^{c}$ together with its $\ell/3$-collar. It follows from the dependence lemma for
the words of $\mathsf{C}^{K1}$; in applying that lemma, note that the cell weights $w_z$ re-weight the
local terms and do not act on the index $1''$ of the inversion lemma.

The Schur-complement statement. This is the Schur-complement corollary. The operator $\mathcal{K}$ of
\eqref{4.14:eq-block-identities-2} is the Schur complement of $\mathsf{A}$ in $\mathsf{Q}$. It is a
sheet by the product and sum rules for sheets, applied with $\rho\le M/4$ to the sheet of
$\mathsf{Q}$ in (a), to the sheet of (D) and to the expansions of clauses (E) and (F).

For the bounds, take a base point $(U,0)$. There all matrices are real symmetric and
$\mathsf{A}>0$ by (B). Hence, for real $A'$, the right side of \eqref{4.14:eq-block-identities-2} is
minimised over real $\varphi$ at $\varphi=-\mathsf{C}\mathsf{J}A'$, with minimum
$\tfrac12\langle A',\mathcal{K}A'\rangle$.
By (B), $\langle A_k,\mathsf{Q}A_k\rangle\ge(\gamma_0/2)(|\varphi|^2+|A'|^2)$. Taking the minimiser for
$\varphi$ gives $\langle A',\mathcal{K}A'\rangle\ge(\gamma_0/2)|A'|^2$. Taking $\varphi=0$ instead gives
\begin{equation*}
\langle A',\mathcal{K}A'\rangle\le\langle A',\mathsf{Q}_{\mathbb{B}'\mathbb{B}'}A'\rangle
\le\gamma_1|A'|^2 .
\end{equation*}
Thus $\gamma_0/2\le\mathcal{K}\le\gamma_1$ at every $(U,0)$.

Conclusion. Assume the admissibility of the prepared large-field data,
\eqref{4.14:eq-prepared-data-admissible-a}, at every nest of Ba{\l}aban's history, and the one-sided
conditions of Definition~\ref{4.14:def-block-size-conditions}. Then clauses (A)--(K) of
Theorem~\ref{4.14:thm-operator-theorem} hold, together with the covariance, dependence and
Schur-complement statements: clause (D) and statements (a)--(c) are established in
Lemmas~\ref{4.14:prf-form-walk-expansion} and~\ref{4.14:prf-positivity-inversion}, and the remaining
clauses and statements were proved above.
\end{proof}

\begin{definition}[Conventions at a large-field operation]\label{4.14:not-operation-conventions}
Fix a level $k$ at which Ba{\l}aban's large-field operation $\mathbf{R}$ acts, within Ba{\l}aban's
small-coupling regime for large-field operations
(Definition~\ref{4.14:def-small-coupling-regime}, display
\eqref{4.14:eq-small-coupling-regime-a}). Let $\mathsf{Z}$ be the class of components on which the
operation acts (Definition~\ref{4.14:def-operation-components}, display
\eqref{4.14:eq-operation-components-a}). The following conventions are in force.
\begin{enumerate}
\item \emph{Units.} Lengths are measured in the units of the level at which the analysis takes
place. At an operation $\mathbf{R}$ at level $k$ these are the level-$k$ units, and the lattice
spacing of level $j$ is $\ell_j = L^{j-k}$.
\item \emph{Partitions.} For every level $j$ the partition into level-$j$ $M$-cubes, of side
$M\ell_j$, and the partition into level-$j$ $MR_j$-cubes, of side $MR_j\ell_j$, where $R_j$ is a
power of $L$, together form one compatible $L$-adic family of partitions
\cite[p.~256]{B-CMP119}.
\item \emph{Enlargements.} Let $\mathcal{P}$ be one of these partitions, with cubes of side $s$,
and let $X$ be a union of closed cubes of $\mathcal{P}$. For $r\ge 0$ put
\begin{equation*}
X^{(+r)} := \{x : \operatorname{dist}_\infty(x,X)\le r\},
\qquad X^{\sim n} := X^{(+ns)}\quad (n\ge 0\ \text{an integer}).
\end{equation*}
If $r$ is a non-negative integer multiple of $s$, then $X^{(+r)}$ is again a union of closed cubes
of $\mathcal{P}$.
\item \emph{Metric enlargement of the domains.} In \cite[(1.10)--(1.11)]{B-CMP122a} the enlarged
domains are read as $\Omega_j^{\sim n} := \Omega_j^{(+nMR_j\ell_j)}$. On every part of $\Omega_j$
whose trace on $\mathsf{Z}$ enters the analysis, this coincides with the enlargement by $n$ layers
of $MR_j\ell_j$-cubes of \cite{B-CMP122a}; this coincidence is the content of the alignment
statement for the two enlargements.
\item \emph{Components.} Three elementary facts about unions of closed $\mathcal{P}$-cubes of side
$s$ are used in the sequel: the separation of distinct components by at least $s$, the exit from
the first component of a segment joining it to another component, and the location of the exit
point of such a segment.
\item \emph{Scales at the operation.} For every level $i$ of the history let $\sigma_i$ be the side
$MR_i\ell_i$ of the level-$i$ $MR_i$-cubes; its definition is recorded in
\eqref{4.14:eq-operation-conventions-a} below. For every integer $n\ge1$ for which $k-n+1$ is a
level of the history put
\begin{equation*}
f(n) := R_{k-n+1}\ell_{k-n+1} = \sigma_{k-n+1}/M .
\end{equation*}
By \eqref{4.14:eq-small-coupling-regime-a} the set $\{n : f(n) = 1\}$ is non-empty (item (b)
below); let $N_0$ be its least element, and let $k_0 := k-N_0$. The integer $k_0$ is Ba{\l}aban's
index in \cite[(1.64)--(1.67)]{B-CMP122a} and \cite[p.~192]{B-CMP122a}; $N_0$ is the integer of
\cite[p.~179]{B-CMP122a}. Put ${\ast} := k_0+1$. This definition, the definition of $\sigma_i$,
and the two relations proved in (b) and (d) below are collected in one display:
\begin{equation}\label{4.14:eq-operation-conventions-a}
{\ast} := k_0+1 = k-N_0+1,\qquad R_\ast = L^{N_0-1},\qquad
\sigma_i := MR_i\ell_i,\qquad N_0\ge 2 .
\end{equation}
Then the following hold.
\begin{itemize}
\item[(a)] $f$ is non-increasing and takes values in integer powers of $L$, and $f(1) = R_k\ge L$.
\item[(b)] $N_0$ exists, and $\sigma_\ast = MR_\ast\ell_\ast = M$. Equivalently $R_\ast = L^{N_0-1}$,
which is the identity $L^{N_0-1} = R_{k-N_0+1}$ of \cite[p.~192]{B-CMP122a}.
\item[(c)] $\sigma_j\le M$ for $j\le{\ast}$, and $\sigma_j\ge LM$ for ${\ast}<j\le k$.
\item[(d)] $N_0\ge2$; hence $\ell_\ast\le L^{-1}$ and ${\ast}\le k-1$.
\item[(e)] $\sigma_j$ is non-decreasing in $j$. Hence, for levels $j\le k'$ of the history, with
$k'$ as in the item on symbols below, the side $\sigma_{k'}$ is a power-of-$L$ multiple of
$\sigma_j$, so that the partition into cubes of side $\sigma_{k'}$ is compatible with that into
cubes of side $\sigma_j$ (item on partitions).
\item[(f)] With $N_k$ the number of preliminary integrations of the operation at level $k$
(Definition~\ref{4.14:def-small-coupling-regime}) and $h := k-N_k$, we have $h<{\ast}-1$.
\end{itemize}
\item \emph{Symbols.} The symbol $k'$ denotes the level of an earlier, or of any other, operation
$\mathbf{R}$. Ba{\l}aban's $k_0 = k-N_0$ keeps its name.
\end{enumerate}
\end{definition}

\begin{proof}
\emph{Enlargements.} Since $\operatorname{dist}_\infty(x,X)$ is the minimum of
$\operatorname{dist}_\infty(x,Q)$ over the closed cubes $Q\subset X$ of $\mathcal{P}$, we have
$X^{(+r)} = \bigcup_{Q} Q^{(+r)}$. Write $Q = \prod_{i=1}^{4}[a_i,a_i+s]$ with the $a_i$ in the grid
of $\mathcal{P}$. Then $Q^{(+r)} = \prod_{i=1}^{4}[a_i-r,\,a_i+s+r]$. When $r\in s\mathbb{N}$ the
endpoints $a_i-r$ and $a_i+s+r$ again lie in the grid, so $Q^{(+r)}$ is a union of closed
$\mathcal{P}$-cubes, and so is $X^{(+r)}$.

\emph{The gap-one bound.} Apply the second inequality of \cite[(2.9)]{B-CMP119}, as displayed in
\eqref{4.14:eq-small-coupling-regime-a}, at the pair of levels $(m+1,m)$. It gives
$R_m\le (L+1)R_{m+1}$. Both $R_m$ and $R_{m+1}$ are powers of $L$, so $R_m/R_{m+1} = L^{e}$ for an
integer $e$, and $L^e\le L+1<L^2$ forces $e\le1$. Hence $R_{m+1}\ge R_m/L$.

\emph{Item (a).} Since $\ell_{k-n} = L^{-1}\ell_{k-n+1}$,
\begin{equation*}
\frac{f(n+1)}{f(n)} = \frac{R_{k-n}\ell_{k-n}}{R_{k-n+1}\ell_{k-n+1}}
= \frac{R_{k-n}}{L\,R_{k-n+1}} \le 1,
\end{equation*}
by the gap-one bound at $m = k-n$. Each $R_{k-n+1}$ is a power of $L$ and
$\ell_{k-n+1} = L^{-n+1}$, so $f(n)$ is an integer power of $L$. Finally $f(1) = R_k\ell_k = R_k$.
At a level where an operation $\mathbf{R}$ acts, $R_k>1$ by
\eqref{4.14:eq-small-coupling-regime-a}; since $R_k$ is a power of $L$, this gives $R_k\ge L$.

\emph{Item (b).} By \eqref{4.14:eq-small-coupling-regime-a}, at an operation $\mathbf{R}$ at level
$k$ there is a positive integer $n$ with $L^{-n+1}R_{k-n+1} = 1$, that is, with $f(n)=1$. Hence the
minimum $N_0$ exists. With ${\ast} = k-N_0+1$ we get $\sigma_\ast = Mf(N_0) = M$, that is,
$R_\ast\ell_\ast = 1$. Since $\ell_\ast = L^{{\ast}-k}$, this gives
$R_\ast = L^{k-{\ast}} = L^{N_0-1}$.

\emph{Item (c).} Write $j = k-n+1$, so that $\sigma_j = Mf(n)$.

If $j\le{\ast}$, then $n\ge N_0$, and by (a) $f(n)\le f(N_0) = 1$; thus $\sigma_j\le M$.

If ${\ast}<j\le k$, then $1\le n<N_0$, and by (a) $f(n)\ge f(N_0) = 1$. Moreover $f(n)\neq1$ by the
minimality of $N_0$. Since $f(n)$ is an integer power of $L$, it follows that $f(n)\ge L$, and
$\sigma_j\ge LM$.

\emph{Item (d).} By (a), $f(1)\ge L>1$, so $N_0\neq1$, that is, $N_0\ge2$. Therefore
\begin{equation*}
\ell_\ast = L^{-(N_0-1)}\le L^{-1}, \qquad {\ast} = k-N_0+1\le k-1 .
\end{equation*}

\emph{Item (e).} The ratio computed in (a) equals $\sigma_{j-1}/\sigma_j$ with $j = k-n+1$. This
ratio is at most $1$ by the gap-one bound, so $\sigma$ is non-decreasing in $j$. Each ratio
$\sigma_{j+1}/\sigma_j = LR_{j+1}/R_j$ is an integer power of $L$, and by monotonicity it is at
least $1$. Hence, for $j\le k'$, the quotient $\sigma_{k'}/\sigma_j$ is a product of non-negative
integer powers of $L$, which is a power of $L$.

\emph{Item (f).} By \eqref{4.14:eq-small-coupling-regime-a}, $N_k>N_0$. Therefore
\begin{equation*}
h+1 = k-N_k+1 < k-N_0+1 = {\ast},
\end{equation*}
that is, $h<{\ast}-1$.
\end{proof}

\begin{lemma}[Small boxes in unions of grid cubes]\label{4.14:lem-small-boxes}
Let $S>0$ and $o=(o_1,\dots,o_d)\in\mathbb{R}^d$, and call a \emph{grid cube} any closed cube
$\prod_{i=1}^d[\,o_i+n_iS,\;o_i+(n_i+1)S\,]$ with $n_1,\dots,n_d\in\mathbb{Z}$; these are the
cubes of the aligned partition of side $S$ whose faces lie on the hyperplanes
$x_i\in o_i+S\mathbb{Z}$. Let $\Omega\subset\mathbb{R}^d$ be a union of grid cubes, and let
$H=\prod_{i=1}^d I_i\subset\Omega$ be a closed box with every side satisfying $|I_i|\le S$.
Then there is a closed cube $Q$ of side $S$ with $H\subset Q\subset\Omega$.
\end{lemma}

\begin{proof}
Each $I_i$ is a closed interval of positive length, $H$ being a closed box; write
$\operatorname{int} I_i$ for its interior. A \emph{grid point} in the $i$-th coordinate is a
point of $o_i+S\mathbb{Z}$, and a \emph{grid interval} is a closed interval
$[\,o_i+nS,\;o_i+(n+1)S\,]$ with $n\in\mathbb{Z}$; it has length $S$.

\emph{Construction of $Q$.} Fix $i$. We say that $I_i$ \emph{straddles} a grid point $v_i$ if
$v_i\in\operatorname{int}I_i$. There is at most one such point: two distinct points of the open
interval $\operatorname{int}I_i$ are at distance less than $|I_i|\le S$, whereas distinct grid
points are at distance at least $S$. If $I_i=[a_i,b_i]$ straddles $v_i$, put
$J_i:=[a_i,a_i+S]$. This is a closed interval of length $S$ with
$I_i\subset J_i\subset[v_i-S,v_i+S]$: indeed $b_i\le a_i+|I_i|\le a_i+S$, so $I_i\subset J_i$;
moreover $a_i\ge v_i-|I_i|\ge v_i-S$ because $v_i\in I_i$, and $a_i+S<v_i+S$ because
$a_i<v_i$. If $I_i$ straddles no grid point, then $I_i$ lies in one grid interval $K_i$, namely
the one containing its interior, and we put $J_i:=K_i$. Set
\begin{equation*}
Q:=\prod_{i=1}^d J_i .
\end{equation*}
Each $J_i$ is a closed interval of length $S$ containing $I_i$, so $Q$ is a closed cube of
side $S$ with $H\subset Q$.

\emph{A cover of $Q$ by grid cubes.} Let $\sigma$ be the set of indices $i$ for which $I_i$
straddles a grid point. For $i\in\sigma$ put
\begin{equation*}
K_i^{+}:=[v_i,\,v_i+S],\qquad K_i^{-}:=[v_i-S,\,v_i],
\end{equation*}
which are grid intervals since $v_i$ is a grid point. For $\varepsilon\in\{+,-\}^{\sigma}$ put
$K_i^{\varepsilon}:=K_i^{\varepsilon_i}$ for $i\in\sigma$, $K_i^{\varepsilon}:=K_i$ for
$i\notin\sigma$, and $C_\varepsilon:=\prod_{i=1}^d K_i^{\varepsilon}$. Each $C_\varepsilon$ is a
grid cube. Moreover $Q\subset\bigcup_\varepsilon C_\varepsilon$: if $x\in Q$, then for
$i\in\sigma$ the coordinate $x_i$ lies in $J_i\subset[v_i-S,v_i+S]=K_i^-\cup K_i^+$, and we
choose $\varepsilon_i$ with $x_i\in K_i^{\varepsilon_i}$; for $i\notin\sigma$,
$x_i\in J_i=K_i$. With this $\varepsilon$ we have $x\in C_\varepsilon$.

\emph{Each $C_\varepsilon$ lies in $\Omega$.} Fix $\varepsilon$. We first find a point
$q\in H\cap\operatorname{int}C_\varepsilon$. For $i\in\sigma$ with $\varepsilon_i=+$ choose
$q_i\in(v_i,b_i)$; this interval is non-empty, lies in $\operatorname{int}I_i$, and lies in
$(v_i,v_i+S)=\operatorname{int}K_i^{+}$ since $b_i\le a_i+S<v_i+S$. For $i\in\sigma$ with
$\varepsilon_i=-$ choose $q_i\in(a_i,v_i)$, which lies in $\operatorname{int}I_i$ and in
$(v_i-S,v_i)=\operatorname{int}K_i^{-}$ since $a_i>v_i-S$ by $b_i-a_i\le S$ and $b_i>v_i$.
For $i\notin\sigma$ choose $q_i\in I_i\cap\operatorname{int}K_i$, which is non-empty because
$I_i\subset K_i$ has positive length and $K_i\setminus\operatorname{int}K_i$ consists of two
points. Then $q=(q_1,\dots,q_d)\in H\cap\operatorname{int}C_\varepsilon$.

Since $H\subset\Omega$, the point $q$ lies in some grid cube $C$ among those whose union is
$\Omega$. A grid cube $C\ne C_\varepsilon$ does not meet $\operatorname{int}C_\varepsilon$: in
some coordinate $i$ the grid intervals of $C$ and $C_\varepsilon$ are distinct, two distinct
grid intervals meet at most in a common endpoint, and no endpoint of the grid interval of
$C_\varepsilon$ lies in its interior. Hence $\Omega\cap\operatorname{int}C_\varepsilon$ equals
$\operatorname{int}C_\varepsilon$ if $C_\varepsilon$ is one of the grid cubes composing
$\Omega$, and is empty otherwise. As $q\in\Omega\cap\operatorname{int}C_\varepsilon$, the cube
$C_\varepsilon$ is one of the grid cubes composing $\Omega$, and so $C_\varepsilon\subset\Omega$.

Consequently $Q\subset\bigcup_\varepsilon C_\varepsilon\subset\Omega$, which together with
$H\subset Q$ proves the lemma.
\end{proof}

The cube $Q$ is in general not a cube of the partition: it has been slid off the grid in
every coordinate in which $H$ straddles a grid point. This is admissible for the use made of
the lemma, since the families of cubes in \cite[(2.38)]{B-CMP119},
\cite[(1.65), (1.67), (1.69)]{B-CMP122a} and \cite[(1.12)]{B-CMP109} do not refer to any
partition. If $Q$ were required to be a cube of the partition, a single cube would not suffice
whenever some side $I_i$ of $H$ straddles a grid point, since no grid interval contains such an
$I_i$.

\begin{lemma}[The hull of a second-cover box in one cube]\label{4.14:lem-hull-cube}
Let $\{\widehat{\Omega}_i\}_{0\le i\le \hat{k}}$ be a nest satisfying condition (A1) of
\eqref{4.14:eq-prepared-data-admissible-a}, and let
$M\ge 13L^{2}M_1^{\sharp}$, which is the first of the conditions on $M$ in
\eqref{4.14:eq-block-size-conditions-a}. Put
$\widehat{\Omega}_{\hat{k}+1}:=\widehat{\Omega}_{\hat{k}+2}:=\emptyset$, as in \cite{B-CMP122a}.
Let the territories $\widehat{X}_n$ of the nest be as in
Convention~\ref{4.14:not-free-current-notation}: for $n\ge 1$, $\widehat{X}_n$ is the closure of
$\widehat{\Omega}_n\setminus\widehat{\Omega}_{n+1}$, and $\widehat{X}_0$ is the closure of the
complement of $\widehat{\Omega}_1$. Let $j\ge 2$ and let $\square\in\mathcal{D}^{\sharp}_j$ be a
box of the second cover of this nest; it has side $2s_j$, where $s_j=M_1^{\sharp}\ell_j$, and its
centre $c$ lies in $\widehat{X}_j$. Let $\square^{4}$ be the concentric closed cube of half-side
$5s_j$, hence of side $10s_j$, and define
\begin{equation*}
m(\square):=\max\bigl\{\,i\le \hat{k}\;:\;\square^{4}\subset\widehat{\Omega}_i\,\bigr\},
\qquad m:=m(\square).
\end{equation*}
Then the following hold.
\begin{enumerate}
\item[(1)] $m\in\{j-1,\,j\}$.
\item[(2)] $\square^{4}\subset\widehat{\Omega}_m\setminus\widehat{\Omega}_{m+2}$ and
$\square^{4}\not\subset\widehat{\Omega}_{m+1}$. Consequently, among the territories of the nest,
$\square^{4}$ meets only $\widehat{X}_m$ and $\widehat{X}_{m+1}$, and in addition
$\widehat{X}_{m-1}$ at most along $\partial\widehat{\Omega}_m$.
\item[(3)] The side of $\square^{4}$ satisfies $10s_j\le M\ell_m$.
\item[(4)] There is a closed cube $Q=Q(\square)$ of side exactly $M\ell_m$, where $m=m(\square)$,
such that
\begin{equation}\label{4.14:eq-hull-cube-a}
\square^{4}\subset Q\subset \widehat{\Omega}_m\setminus\widehat{\Omega}_{m+2},
\qquad Q\cap\widehat{\Omega}^{c}_{m+1}\neq\emptyset ,
\end{equation}
that is, $Q$ contains a point outside the closed set $\widehat{\Omega}_{m+1}$.
\end{enumerate}
\end{lemma}

\begin{proof}
We use condition (A1) in the following form: (a) the sets are nested,
$\widehat{\Omega}_0\supset\widehat{\Omega}_1\supset\dots\supset\widehat{\Omega}_{\hat{k}}$, and this
persists after appending $\widehat{\Omega}_{\hat{k}+1}=\widehat{\Omega}_{\hat{k}+2}=\emptyset$;
(b) for $1\le i\le \hat{k}$, $\widehat{\Omega}_i$ is a union of closed level-$i$ $M$-cubes, that
is, of cubes of side $M\ell_i$ of one compatible family of aligned partitions; in particular each
$\widehat{\Omega}_i$ is closed; (c) for $1\le i\le \hat{k}-1$,
$\operatorname{dist}_\infty(\widehat{\Omega}^{c}_i,\widehat{\Omega}_{i+1})\ge M\ell_i$. The
inequality in (c) holds trivially, the distance to the empty set being $+\infty$, whenever
$\widehat{\Omega}_{i+1}=\emptyset$. Recall that $\ell_i=L^{i-k}$ is increasing in $i$ and that
$M\ell_{i-1}=(M/L)\ell_i$. The territories are indexed by $j\le\hat{k}$, since
$\widehat{X}_{\hat{k}+1}=\emptyset$; thus $2\le j\le \hat{k}$.

(1) Since $\widehat{\Omega}_j$ is closed, $\widehat{X}_j\subset\widehat{\Omega}_j$, so
$c\in\widehat{\Omega}_j$. As $1\le j-1\le \hat{k}-1$, (c) applies with $i=j-1$, and with
$M\ge 13L^2M_1^{\sharp}$ it gives
\begin{equation*}
\begin{split}
\operatorname{dist}_\infty(c,\widehat{\Omega}^{c}_{j-1})
&\ge\operatorname{dist}_\infty(\widehat{\Omega}_j,\widehat{\Omega}^{c}_{j-1})
\ge M\ell_{j-1}=\frac{M}{L}\,\ell_j\\
&\ge 13LM_1^{\sharp}\ell_j=13Ls_j>5s_j .
\end{split}
\end{equation*}
Every $q\in\square^{4}$ satisfies $|q-c|_\infty\le 5s_j$, hence
$q\notin\widehat{\Omega}^{c}_{j-1}$. Thus $\square^{4}\subset\widehat{\Omega}_{j-1}$. The set over
which $m(\square)$ is the maximum is therefore non-empty, and $m\ge j-1$.

Next, $\widehat{\Omega}_j\setminus\widehat{\Omega}_{j+1}\subset\widehat{\Omega}^{c}_{j+1}$, so $c$
lies in the closure of $\widehat{\Omega}^{c}_{j+1}$, and every neighbourhood of $c$ meets
$\widehat{\Omega}^{c}_{j+1}$. The cube $\square^{4}$ contains the open $|\cdot|_\infty$-ball of
radius $5s_j>0$ about $c$, which is such a neighbourhood. Hence
$\square^{4}\not\subset\widehat{\Omega}_{j+1}$. By (a), $\widehat{\Omega}_i\subset\widehat{\Omega}_{j+1}$
for every $i\ge j+1$, so $\square^{4}\not\subset\widehat{\Omega}_i$ for all such $i$, and $m\le j$.

(2) By definition of $m$, $\square^{4}\subset\widehat{\Omega}_m$. If $m<\hat{k}$, maximality of $m$
gives $\square^{4}\not\subset\widehat{\Omega}_{m+1}$. If $m=\hat{k}$, then
$\widehat{\Omega}_{m+1}=\emptyset$ and the same holds. Choose
$p\in\square^{4}\setminus\widehat{\Omega}_{m+1}$.

Let $q\in\square^{4}$. By (1), $m+1\ge j$, so $\ell_{m+1}\ge\ell_j$. If $m+2\le\hat{k}$, (c) applies
with $i=m+1\ge 1$; otherwise $\widehat{\Omega}_{m+2}=\emptyset$ and the last inequality below is
trivial. We obtain
\begin{equation*}
\begin{split}
|p-q|_\infty&\le 10s_j<13L^2s_j=13L^2M_1^{\sharp}\ell_j\le M\ell_j\\
&\le M\ell_{m+1}\le\operatorname{dist}_\infty(\widehat{\Omega}^{c}_{m+1},\widehat{\Omega}_{m+2}).
\end{split}
\end{equation*}
Since $p\in\widehat{\Omega}^{c}_{m+1}$, this gives $q\notin\widehat{\Omega}_{m+2}$. Hence
$\square^{4}\subset\widehat{\Omega}_m\setminus\widehat{\Omega}_{m+2}$.

For the territories, note first that
$\widehat{X}_n\subset\overline{\widehat{\Omega}^{c}_{n+1}}$ for every $n$. If $n\ge m+2$, then
$\widehat{X}_n\subset\widehat{\Omega}_n\subset\widehat{\Omega}_{m+2}$, which is disjoint from
$\square^{4}$. For $n=m-1$, which is at least $0$ because $m\ge j-1\ge 1$, we have
\begin{equation*}
\widehat{X}_{m-1}\cap\square^{4}\subset\overline{\widehat{\Omega}^{c}_{m}}\cap\widehat{\Omega}_m
=\partial\widehat{\Omega}_m ,
\end{equation*}
because $\widehat{\Omega}_m$ is closed. Now let $0\le n\le m-2$. Then $1\le n+1\le m-1\le\hat{k}-1$,
and (c) with $i=n+1$ gives $\operatorname{dist}_\infty(\widehat{\Omega}^{c}_{n+1},\widehat{\Omega}_{n+2})>0$.
Hence $\overline{\widehat{\Omega}^{c}_{n+1}}$ is disjoint from $\widehat{\Omega}_{n+2}$. Since
$n+2\le m$, (a) gives $\widehat{\Omega}_{n+2}\supset\widehat{\Omega}_m\supset\square^{4}$, so
$\widehat{X}_n\cap\square^{4}=\emptyset$. The only territories that remain are $\widehat{X}_m$ and
$\widehat{X}_{m+1}$.

(3) By $m\ge j-1$ and the chain displayed in the proof of (1),
\begin{equation*}
M\ell_m\ge M\ell_{j-1}\ge 13Ls_j\ge 10s_j .
\end{equation*}

(4) Since $1\le j-1\le m\le\hat{k}$, (b) applies with $i=m$. Thus $\widehat{\Omega}_m$ is a union of
closed cubes of an aligned partition of side $S=M\ell_m$. The closed cube $\square^{4}$ lies in
$\widehat{\Omega}_m$ by (2), and each of its sides has length $10s_j\le M\ell_m$ by (3). Lemma~\ref{4.14:lem-small-boxes}
(small boxes in unions of grid cubes) therefore yields a closed cube $Q$ of side $M\ell_m$ with
$\square^{4}\subset Q\subset\widehat{\Omega}_m$.

The point $p$ chosen in (2) lies in $\square^{4}\subset Q$ and not in $\widehat{\Omega}_{m+1}$, so
$Q\cap\widehat{\Omega}^{c}_{m+1}\neq\emptyset$. For $q\in Q$ we have $|p-q|_\infty\le M\ell_m$,
because $Q$ has side $M\ell_m$. Since $L>1$, and with the last inequality justified exactly as in (2),
\begin{equation*}
|p-q|_\infty\le M\ell_m<LM\ell_m=M\ell_{m+1}
\le\operatorname{dist}_\infty(\widehat{\Omega}^{c}_{m+1},\widehat{\Omega}_{m+2}).
\end{equation*}
As $p\in\widehat{\Omega}^{c}_{m+1}$, it follows that $q\notin\widehat{\Omega}_{m+2}$. Hence
$Q\cap\widehat{\Omega}_{m+2}=\emptyset$ and $Q\subset\widehat{\Omega}_m\setminus\widehat{\Omega}_{m+2}$,
which is \eqref{4.14:eq-hull-cube-a}.

For $j\in\{0,1\}$ the hulls may meet
$\widehat{X}_0$, and the backgrounds on them fall under the alternative $(\alpha')$ of
\eqref{4.14:eq-real-backgrounds-a}, that is, the condition \cite[(3.35)]{B-CMP99b} taken at block
size $M_1^{\sharp}$ at the indices $0,1,2$. That condition does not depend on the nest, so no question
about the large-field operation $\mathbf{R}$ arises there. Indeed, the pairs produced by
$\mathbf{R}$ sit at levels strictly above $h$, the lowest level touched by that operation, and
$h\ge 1$.
\end{proof}

In what follows lengths are measured in the units of the level at which the nest under consideration is read, as in Convention~\ref{4.14:not-operation-conventions}. In these units the spacing of level $i$ is $\ell_i$. The lengths written $L^{i}\xi$ in \cite{B-CMP119} and $L^{i}\eta$ in \cite{B-CMP122a}, where $\xi$ and $\eta$ are Ba{\l}aban's lattice spacings in the two papers, are both equal to $\ell_i$.

The families of cubes of \cite[(2.38)]{B-CMP119} and of \cite[(1.65), (1.67), (1.69)]{B-CMP122a} all have the following form. A cube $\square$ of such a family has a prescribed size, $CM\ell_i$ or $CM$; in \cite{B-CMP122a} the constant $C$ is a positive integer \cite[p.~191]{B-CMP122a}, and below only $C=1$ occurs. For such a cube there must exist a gauge transformation $u$, defined on the intersection of $\square$ with the domain of the configuration $U$, such that $U^{u}=\exp(i\eta A)$, with $\xi$ in place of $\eta$ in \cite{B-CMP119}, $A$ takes values in the Lie algebra $\mathfrak{g}$, and on that intersection
\begin{equation*}
\ell_i\,|A|<\mathsf{L}_0^{2e}\,BCM\alpha_{0,i},\qquad
\ell_i^{2}\,|\nabla^{\eta}A|<\mathsf{L}_0^{2e}\,BCM\alpha_{0,i} .
\end{equation*}
Here $\nabla^{\eta}$, with $\nabla^{\xi}$ in \cite{B-CMP119}, is Ba{\l}aban's lattice derivative at that spacing, $B$ is Ba{\l}aban's absolute constant, $\alpha_{0,i}$ is the radius of index $i$ of \cite[(2.28)]{B-CMP119}, $\mathsf{L}_0$ is Ba{\l}aban's auxiliary integer of \cite[(1.24)]{B-CMP122a}, and $e\ge 0$ is an integer; $e=0$ in \cite[(2.38)]{B-CMP119} and in \cite[(1.69)]{B-CMP122a}. We call $\ell_i$ the \emph{scale} of the family and the right-hand side $\mathsf{L}_0^{2e}BCM\alpha_{0,i}$ its \emph{printed bound}. We call $i$ the \emph{$\alpha$-index} and $e$ the \emph{$\mathsf{L}_0$-exponent} of the printed bound. When the printed bound is the one attached to a cube $Q$, we write $e=e(Q)$.

\begin{lemma}[The printed gauge family of the hull's cube]\label{4.14:lem-gauge-family}
\begin{enumerate}
\item[(a)] Let the nest $\{\widehat{\Omega}_i\}_{i\le\hat{k}}$, the cube $\square$, $H=\square^{4}$,
$m=m(\square)$ and $Q=Q(\square)$ be as in Lemma~\ref{4.14:lem-hull-cube}. Then $Q$ is a cube of
the family of \cite[(2.38)]{B-CMP119}, read on the nest $\{\widehat{\Omega}_i\}$ with top index
$\hat{k}$, as follows.
\begin{itemize}
\item Suppose $m\le\hat{k}-1$. Then $Q$ is a member with Ba{\l}aban's index $n$ of that family equal
to $m$, and $C:=1$:
\begin{equation*}
Q\subset\widehat{\Omega}_m\setminus\widehat{\Omega}_{m+2},\qquad
Q\cap\widehat{\Omega}_{m+1}^{c}\neq\emptyset,\qquad \text{side } M\ell_m=CM\ell_m .
\end{equation*}
At $m=\hat{k}-1$ the set $\widehat{\Omega}_{\hat{k}+1}$ is empty.
\item Suppose $m=\hat{k}$. Then $Q$ is a member of the top alternative: $Q\subset\widehat{\Omega}_{\hat{k}}$, of size $CM$ with $C=1$.
\end{itemize}
A fortiori, $Q$ belongs to the one-sided families of \cite[(1.65), (1.69)]{B-CMP122a} read on
$\{\widehat{\Omega}_i\}$, namely $Q\subset\widehat{\Omega}_m$ and $Q\cap\widehat{\Omega}_{m+1}^{c}\neq\emptyset$.
\item[(b)] Let $\mathbf{R}^{(k)}$ be a large-field operation at level $k$ acting on the class
$\mathsf{Z}$ of Definition~\ref{4.14:def-operation-components}. Attached to it are:
\begin{itemize}
\item its nest $\{\Omega_i\}_{i\le k}$ before the operation (the \emph{old nest});
\item the family $Z''$ and the new domains $\Omega''_i$ of \cite[(1.10)--(1.12)]{B-CMP122a}, in the
reading of Convention~\ref{4.14:not-operation-conventions};
\item the integer $k_0=k-N_0$ of \eqref{4.14:eq-operation-conventions-a} and the lowest level
touched, $h=k-N_k$, of Definition~\ref{4.14:def-operation-components}.
\end{itemize}
Let $X$ be a component that leaves $\mathsf{Z}$ at stage $n\ge1$, that is, a component in which
the preliminary integrations of the operation are finished after $n$ of them \cite[p.~191]{B-CMP122a},
by any of the exits described in the lemma on exits
and stage nests. Put $p:=h+n\le k$. Let $\{\widehat{\Omega}_i\}_{i\le k}$ be the \emph{stage nest}
$\mathcal{N}^{(n)}$:
\begin{equation*}
\widehat{\Omega}_i:=\Omega''_i\quad(i\le p),\qquad \widehat{\Omega}_i:=\Omega_i\quad(p<i\le k).
\end{equation*}
For $i\le h$ one has $\Omega''_i=\Omega_i$ by \cite[(1.12)]{B-CMP122a}. This is the nest on which
\cite[(1.64)--(1.69)]{B-CMP122a} define the space of stage $n$. The level of that stage is
written $j$ there, so that $j=p$ throughout (b). The index of the second cover to which $\square$
belongs in Lemma~\ref{4.14:lem-hull-cube} is written $j_\square$ in this lemma and its proof; it
enters below only through $m\in\{j_\square-1,j_\square\}$. By the lemma on exits
and stage nests, applied with the rule of Definition~\ref{4.14:def-operation-components} for the
components on which an operation acts, the stage nest is the nest of $X$ at every later moment, until a component
containing $X$ is filled, in the sense of that lemma. By the lemma asserting that Ba{\l}aban's
stage nests are admissible, $\{\widehat{\Omega}_i\}$ satisfies condition (A1) of
\eqref{4.14:eq-prepared-data-admissible-a}, for every assignment of states.
Let $\square$, $H$, $m:=m(\square)$ and $Q:=Q(\square)$ be as in Lemma~\ref{4.14:lem-hull-cube} on
this nest, with $\hat{k}=k$. Then exactly one of the following holds.
\begin{itemize}
\item[(G1)] Case $m<p$. Then
\begin{equation*}
Q\subset\Omega''_m,\qquad Q\cap(\Omega''_{m+1})^{c}\neq\emptyset,\qquad
\text{side } M\ell_m=CM\ell_m\ (C=1),
\end{equation*}
so $Q$ is in the family of point (i) of \cite[(1.65)]{B-CMP122a} at the index $m$. Its printed bound
is $\mathsf{L}_0^{2\max\{0,m-k_0\}}BCM\alpha_{0,m}$ at scale $\ell_m$.
\item[(G2)] Case $m=p$ and $Q\cap\Omega_{k_0+1}^{c}\neq\emptyset$, with $\Omega_{k_0+1}$ from the old
nest. Then $Q$ is in the family of point (i) of \cite[(1.65)]{B-CMP122a} at the index $m=j=p$. This is
the alternative
\begin{equation*}
Q\subset\Omega''_p,\qquad Q\cap\Omega_{k_0+1}^{c}\neq\emptyset,\qquad
\text{side } M\ell_p=CM\ell_p\ (C=1).
\end{equation*}
Its printed bound is $\mathsf{L}_0^{2\max\{0,p-k_0\}}BCM\alpha_{0,p}$ at scale $\ell_p$. This case
always occurs when $p\le k_0$, for then
\begin{equation*}
\Omega_{k_0+1}\subset\Omega_{p+1}=\widehat{\Omega}_{p+1}\quad\text{and}\quad
Q\not\subset\widehat{\Omega}_{p+1}.
\end{equation*}
\item[(G3)] Case $m=p$ and $Q\subset\Omega_{k_0+1}$, with $\Omega_{k_0+1}$ from the old nest. Then
$p\ge k_0+1$. Define, with respect to the old nest,
\begin{equation*}
\bar{m}:=\max\{\,i\le p:\ Q\subset\Omega_i\,\}.
\end{equation*}
Then $k_0+1\le \bar{m}\le p$, and
\begin{equation*}
Q\subset\Omega_{\bar{m}},\qquad Q\cap\Omega_{\bar{m}+1}^{c}\neq\emptyset .
\end{equation*}
Hence $Q$ is in the family of point (ii) of \cite[(1.67)]{B-CMP122a} at the index $\bar{m}$, with
Ba{\l}aban's $m$ equal to $\bar{m}$, his $j$ equal to $p$, and $C=1$. Its printed bound is
$\mathsf{L}_0^{2(p-\bar{m})}BCM\alpha_{0,p}$ at scale $\ell_p=\ell_m$, and
\begin{equation*}
p-\bar{m}\le p-k_0-1<m-k_0 .
\end{equation*}
\item[(G4)] Case $m>p$. Then $Q\subset\Omega_m=\widehat{\Omega}_m$, with side $M\ell_m$. Moreover
$Q\cap\Omega_{m+1}^{c}\neq\emptyset$ if $m<k$, while $Q\subset\Omega_k$ if $m=k$. Hence $Q$ is in the family
of point (iii) of \cite[(1.69)]{B-CMP122a}, with $C=1$. Its printed bound is $BCM\alpha_{0,m}$ at
scale $\ell_m$, with no power of $\mathsf{L}_0$.
\end{itemize}
\end{enumerate}
In every case the $\alpha$-index of the printed bound is $m=m(\square)$ and its scale is $\ell_m$.
Its $\mathsf{L}_0$-exponent $e(Q)$ satisfies
\begin{equation}\label{4.14:eq-gauge-family-a}
0\le e(Q)\le\max\{0,\,m-k_0\}.
\end{equation}
More precisely, $e(Q)=\max\{0,m-k_0\}$ in (G1) and (G2), $e(Q)=p-\bar{m}$ in (G3), and $e(Q)=0$ in (a)
and (G4). Finally, $e(Q)>0$ forces $k_0+1\le m\le k$.
\end{lemma}

\begin{proof}
\emph{Part (a).} Assertion \eqref{4.14:eq-hull-cube-a} of Lemma~\ref{4.14:lem-hull-cube} gives
\begin{equation*}
Q\subset\widehat{\Omega}_m\setminus\widehat{\Omega}_{m+2},\qquad
Q\cap\widehat{\Omega}_{m+1}^{c}\neq\emptyset,
\end{equation*}
and says that $Q$ has side $M\ell_m=1\cdot M\ell_m$. If $m=\hat{k}$, it says that
$Q\subset\widehat{\Omega}_{\hat{k}}$ and that $Q$ has side $M$. We read these facts as memberships.
\begin{itemize}
\item For $1\le m\le\hat{k}-1$ they are the defining conditions of the family of
\cite[(2.38)]{B-CMP119} at Ba{\l}aban's index $n=m$ with $C=1$. At $n=\hat{k}-1$ the set $\widehat{\Omega}_{n+2}=\widehat{\Omega}_{\hat{k}+1}$
is empty, by the convention of Lemma~\ref{4.14:lem-hull-cube}.
\item For $m=\hat{k}$ they are the top alternative of that family, a cube contained in $\widehat{\Omega}_{\hat{k}}$
of size $CM$ with $C=1$.
\end{itemize}
Here $m\ge1$, because $m\in\{j_\square-1,j_\square\}$ with $j_\square\ge2$ by
Lemma~\ref{4.14:lem-hull-cube}, $j_\square$ being the index of the second cover to which $\square$ belongs.

The one-sided families of \cite[(1.65), (1.69)]{B-CMP122a} impose only the conditions
$Q\subset\widehat{\Omega}_m$ and $Q\cap\widehat{\Omega}_{m+1}^{c}\neq\emptyset$. They are obtained by dropping the
condition of avoiding $\widehat{\Omega}_{m+2}$, and dropping a condition only enlarges a family. So
$Q$ belongs to them a fortiori.

\emph{Part (b): exhaustiveness.} By Lemma~\ref{4.14:lem-hull-cube}, $m\in\{j_\square-1,j_\square\}$,
where $j_\square\ge2$ is the index of the second cover; and $m\le\hat{k}=k$ by the definition of
$m(\square)$. Hence $1\le m\le k$.
The cases $m<p$, $m=p$ and $m>p$ exhaust the possibilities. In the case $m=p$, either
$Q\cap\Omega_{k_0+1}^{c}\neq\emptyset$ or $Q\subset\Omega_{k_0+1}$. So the four cases (G1)--(G4) are exhaustive
and mutually exclusive.

\emph{Case (G1), $m<p$.} By the definition of the stage nest, $\widehat{\Omega}_m=\Omega''_m$. For $h<m<p$
this is the definition itself. For $m\le h$ one moreover has $\Omega''_m=\Omega_m=\widehat{\Omega}_m$ by
\cite[(1.12)]{B-CMP122a}. Since $m+1\le p$, likewise $\widehat{\Omega}_{m+1}=\Omega''_{m+1}$, with the same
remark if $m+1\le h$. Assertion \eqref{4.14:eq-hull-cube-a} of Lemma~\ref{4.14:lem-hull-cube} therefore gives
\begin{equation*}
Q\subset\widehat{\Omega}_m=\Omega''_m,\qquad
Q\cap\widehat{\Omega}_{m+1}^{c}=Q\cap(\Omega''_{m+1})^{c}\neq\emptyset .
\end{equation*}
Moreover $1\le m\le p-1=j-1$. These are the conditions of point (i) of \cite[(1.65)]{B-CMP122a} at the
index $m$. The size $CM\ell_m$ required there holds with $C=1$, since the side of $Q$ is $M\ell_m$. The
bound attached to the index $m$ in \cite[(1.65)]{B-CMP122a} is the one displayed in (G1), at scale
$\ell_m$.

\emph{Case (G2), $m=p$ and $Q\cap\Omega_{k_0+1}^{c}\neq\emptyset$.} By assertion \eqref{4.14:eq-hull-cube-a}
of Lemma~\ref{4.14:lem-hull-cube},
\begin{equation*}
Q\subset\widehat{\Omega}_p=\Omega''_p=\Omega''_j .
\end{equation*}
The hypothesis of the case is the second condition of the alternative $m=j$ of point (i) of
\cite[(1.65)]{B-CMP122a}. The size is $M\ell_p=CM\ell_j$ with $C=1$. So $Q$ is in that family, and the
attached bound is $\mathsf{L}_0^{2\max\{0,p-k_0\}}BCM\alpha_{0,p}$ at scale $\ell_p$.

Now suppose $p\le k_0$. Then $k_0+1\ge p+1$. Since the old nest is decreasing, this gives
$\Omega_{k_0+1}\subset\Omega_{p+1}$. Since $p+1>p$, the definition of the stage nest gives
$\Omega_{p+1}=\widehat{\Omega}_{p+1}$. Assertion \eqref{4.14:eq-hull-cube-a} of Lemma~\ref{4.14:lem-hull-cube} gives
$Q\not\subset\widehat{\Omega}_{m+1}=\widehat{\Omega}_{p+1}$. Hence $Q\not\subset\Omega_{k_0+1}$, that is,
$Q\cap\Omega_{k_0+1}^{c}\neq\emptyset$. Thus (G2) applies and (G3) cannot occur. This agrees with the
statement in \cite[p.~191]{B-CMP122a} that for $j\le k_0$ there are no powers of $\mathsf{L}_0$ in point
(i) and point (ii) is empty: indeed $\max\{0,p-k_0\}=0$ in that case.

\emph{Case (G3), $m=p$ and $Q\subset\Omega_{k_0+1}$.} By the last paragraph, $p\le k_0$ would force
$Q\not\subset\Omega_{k_0+1}$; hence $p\ge k_0+1$. The set $\{i\le p:\ Q\subset\Omega_i\}$ therefore contains
$k_0+1$. So $\bar{m}$ is defined, and $k_0+1\le \bar{m}\le p$. We show $Q\not\subset\Omega_{\bar{m}+1}$ in two cases.
\begin{itemize}
\item If $\bar{m}<p$, the maximality of $\bar{m}$ gives $Q\not\subset\Omega_{\bar{m}+1}$.
\item If $\bar{m}=p$, then $\Omega_{p+1}=\widehat{\Omega}_{p+1}$ by the definition of the stage nest. When $p=k$,
both sets are empty, by the convention of Lemma~\ref{4.14:lem-hull-cube} and the convention
$\Omega_{k+1}=\emptyset$ of \cite{B-CMP122a}. Assertion \eqref{4.14:eq-hull-cube-a} of
Lemma~\ref{4.14:lem-hull-cube} gives $Q\not\subset\widehat{\Omega}_{p+1}$, hence $Q\not\subset\Omega_{p+1}$.
\end{itemize}
Either way $Q\cap\Omega_{\bar{m}+1}^{c}\neq\emptyset$. Further, $Q\subset\Omega_{\bar{m}}$ by the definition of $\bar{m}$. The
side of $Q$ is $M\ell_p=CM\ell_j$ with $C=1$, and $\bar{m}$ lies in the range $k_0+1,\dots,j$. These are the
conditions of point (ii) of \cite[(1.67)]{B-CMP122a}, with Ba{\l}aban's $m$ equal to $\bar{m}$.

The bound attached there, with $j=p=m$, reads
\begin{equation*}
\ell_m|A|,\ \ell_m^{2}|\nabla^{\eta}A|<\mathsf{L}_0^{2(p-\bar{m})}BCM\alpha_{0,m}.
\end{equation*}
Its scale is $\ell_m$ and its $\alpha$-index is $m$. Its $\mathsf{L}_0$-exponent satisfies
\begin{equation*}
p-\bar{m}\le p-(k_0+1)=m-k_0-1 .
\end{equation*}

\emph{Case (G4), $m>p$.} Both indices $m$ and $m+1$ exceed $p$. By the definition of the stage nest,
$\widehat{\Omega}_m=\Omega_m$, and $\widehat{\Omega}_{m+1}=\Omega_{m+1}$ when $m<k$. Assertion
\eqref{4.14:eq-hull-cube-a} of Lemma~\ref{4.14:lem-hull-cube} therefore reads as follows:
\begin{itemize}
\item for $m<k$, $Q\subset\Omega_m$ and $Q\cap\Omega_{m+1}^{c}\neq\emptyset$;
\item for $m=k$, $Q\subset\Omega_k$.
\end{itemize}
In both cases the side is $M\ell_m$. Since $m\ge p+1=j+1$, these are the conditions of point (iii) of
\cite[(1.69)]{B-CMP122a}, with size $CM\ell_m$ and $C=1$. The attached bound is $BCM\alpha_{0,m}$ at scale
$\ell_m$, without any power of $\mathsf{L}_0$.

\emph{The exponent statement.} The $\alpha$-indices, scales and $\mathsf{L}_0$-exponents are read off the
bounds identified in (a) and (G1)--(G4).
\begin{itemize}
\item In (G1) and (G2), $e(Q)=\max\{0,m-k_0\}$, with $m=p$ in (G2), so \eqref{4.14:eq-gauge-family-a} holds
with equality on the right.
\item In (G3), $e(Q)=p-\bar{m}\ge0$ because $\bar{m}\le p$. Also $e(Q)\le m-k_0-1<m-k_0$, and $m-k_0\ge1$ there.
\item In (a) and (G4), $e(Q)=0$.
\end{itemize}
Now suppose $e(Q)>0$. In (G1) and (G2) this requires $m>k_0$. In (G3) it requires $p\ge \bar{m}+1\ge k_0+2$,
that is, $m=p>k_0$. In all cases $m\le k$, since $m$ is one of the levels $i\le k$ of the operation.
Hence $e(Q)>0$ forces $k_0+1\le m\le k$.
\end{proof}

\begin{remark}[Tiling, later steps and earlier drop-outs]\label{4.14:rem-tiling-later-steps}
We work in the setting of part (b) of the lemma on the printed gauge family of the hull's cube
(Lemma~\ref{4.14:lem-gauge-family}). An operation $\mathbf{R}^{(k)}$ acts on the class
$\mathsf{Z}$ of \eqref{4.14:eq-operation-components-a}, with its old nest
$\{\Omega_i\}_{i\le k}$, its family $Z''$ and its new domains $\Omega''_i$. The numbers
$k_0=k-N_0$ and $\ast=k_0+1$ are those of \eqref{4.14:eq-operation-conventions-a}, and
$h=k-N_k$. A component $X$ of $\mathsf{Z}$ leaves $\mathsf{Z}$ at stage $n\ge1$, and $p=h+n\le k$.
The nest $\{\widehat{\Omega}_i\}$ is the stage nest $\mathcal{N}^{(n)}$, that is,
$\widehat{\Omega}_i=\Omega''_i$ for $i\le p$ and $\widehat{\Omega}_i=\Omega_i$ for $i>p$.
We write $\widehat{X}_j$ for the closure of $\widehat{\Omega}_j\setminus\widehat{\Omega}_{j+1}$
and call it the territory of index $j$.
For a set $S$ and $r>0$, $S^{(+r)}$ denotes the set of points within distance $r$
of $S$, lengths being measured in level-$k$ units; the only property of it used below is
$S\subset S^{(+r)}$. For a
cube $\square$ of the second cover, $Q=Q(\square)$ and $m=m(\square)$ are as in
\eqref{4.14:eq-hull-cube-a}.
\begin{enumerate}
\item[(1)] \emph{Tiling on $\mathsf{Z}$.} Clauses (i), (ii), (iii) of
\cite[(1.64), (1.66), (1.68)]{B-CMP122a} impose, besides the conditions on the cubes of the
families displayed in \cite[(1.65), (1.67), (1.69)]{B-CMP122a}, further conditions, each on a
region of its own. On $\mathsf{Z}$ these regions tile $\Omega''_1$
consistently with the case list of Lemma~\ref{4.14:lem-gauge-family}. To be precise, let
$p>k_0$. Then the old layers $\Omega_{i}\setminus\Omega_{i+1}$,
$k_0+1\le i\le p$, meet
$\mathsf{Z}$ inside
\begin{equation}\label{4.14:eq-tiling-later-steps-1}
\begin{aligned}
(\Omega_{k_0+1}\setminus\Omega_{p+1})\cap\mathsf{Z}
&\subset(\Omega''_p\setminus\widehat{\Omega}_{p+1})\cap\mathsf{Z}\\
&\subset\overline{\widehat{\Omega}_p\setminus\widehat{\Omega}_{p+1}}=\widehat{X}_p ,
\end{aligned}
\end{equation}
that is, inside the territory of index $p$ of $\mathcal{N}^{(n)}$, at the scale $\ell_p$.
Cases (G2) and (G3) of Lemma~\ref{4.14:lem-gauge-family} put $Q$ in exactly this territory.
Moreover, $\ell_p$ is exactly the scale at which the regularity condition bounded by
$a^{\sharp}$ weighs this territory. Off $\mathsf{Z}$ one has
$\Omega''_p=\Omega_p\subset\Omega_{k_0+1}$. This is the reverse inclusion, and it does not
give the tiling. Off $\mathsf{Z}$ the old layers with $i<p$ are the old territories at
their own scales, and nothing is asserted about them here.
\item[(2)] \emph{Later steps.} Consider a later renormalisation step $\mathbf{T}$, and write
$\{\Omega^{\mathrm{now}}_i\}$ for the nest carried at that moment by the term of the expansion
produced when $X$ left $\mathsf{Z}$. This nest
agrees with $\mathcal{N}^{(n)}$ on $X$ at all levels $\le k$ until a fill of a component
containing $X$ (the operation that puts the whole component into the domains at all levels
$\le k$) erases the inner structure of the nest on $X$. After such a fill, part (a)
holds at the filled levels.

Off $X$, but within $M\ell_m$ of a hull, the nest may differ from $\mathcal{N}^{(n)}$ by
later fills and by the removal of new large-field regions from the domains at later steps
$\mathbf{T}$. Ba{\l}aban's
generalised space, which extends the spaces defined in \cite[(1.64)--(1.69)]{B-CMP122a} to
later steps, imposes a condition on this region.
That space is introduced by the sentence of \cite{B-CMP122a} beginning
``it is easy to generalize'', printed without proof there; cited as printed.
The condition it imposes on this region admits two natural readings:
\begin{itemize}
\item the families of the type of \cite[(2.38)]{B-CMP119} read on the current nest;
\item the families (i), (ii), (iii) of $\mathcal{N}^{(n)}$ kept on the levels $\le k$ of $X$,
on which the nest is that of $\mathcal{N}^{(n)}$.
\end{itemize}
Under either reading $Q(\square)$ is a member of the corresponding family. Under both
readings the $\mathsf{L}_0$-exponent,
in the sense of \eqref{4.14:eq-gauge-family-a}, that \cite{B-CMP122a} can attach at index $m$
is at most $\max\{0,m-k_0\}$.
Which of these families the space of the later step uses at such a cube is settled by the
lemma on the families of that space; the bound on the value of the printed
member holds under all of them.
\item[(3)] \emph{Earlier drop-outs.} A component that dropped out at an earlier stage $n'<n$ of
the same operation carries $\mathcal{N}^{(n')}$. For it, part (b) holds with $p':=h+n'$ in
place of $p$ and with the same $k_0$. Now let $Y$ be a component of the second class, that
is, one that remains large-field after the operation, and which carries the stage nest
$\mathcal{N}^{(N_k)}$ of the last stage of the operation. For
it, part (b) holds with $p=k$. Then case (G4) is void, and cases (G2) and (G3) are read at
$m=k$ with the convention $\Omega_{k+1}=\emptyset$ of Lemma~\ref{4.14:lem-hull-cube}.
\item[(4)] \emph{The componentwise nest.} At the moment of stage $n$ of the operation the
components of $\mathsf{Z}$ are survivors at stage $n$ (components still in $\mathsf{Z}$) or
earlier drop-outs (components that left $\mathsf{Z}$ at stages $n'<n$), and off
$\mathsf{Z}$ the nest is the old one.
In this item $X'$ ranges over all components of $\mathsf{Z}$.
For $X'$ let $n(X')$ be $n$ if $X'$ is a survivor and $n'$ if $X'$
dropped out at stage $n'<n$, and put $p(X'):=h+n(X')$. The global nest is
defined componentwise: it is $\Omega''_i$ on $X'$ for $h<i\le p(X')$, and $\Omega_i$ otherwise.
It satisfies (A1) of \eqref{4.14:eq-prepared-data-admissible-a}. The case analysis at a hull
reads the state of at most one component. Hence the stage nest $\mathcal{N}^{(n)}$ in part
(b) may be read as this global componentwise nest, with $p:=p(X')$ for the component $X'$
that $Q$ meets; a cube $Q$ meeting no component of $\mathsf{Z}$ falls under part (a), read
on the old structure.
\end{enumerate}
\end{remark}

\begin{proof}
\emph{Part (1).} Let $p>k_0$. Since the old nest is decreasing, the union of the layers
$\Omega_{i}\setminus\Omega_{i+1}$ over $k_0+1\le i\le p$ is
$\Omega_{k_0+1}\setminus\Omega_{p+1}$. It therefore suffices to prove
\eqref{4.14:eq-tiling-later-steps-1}.

On $\mathsf{Z}$ one has $\Omega''_p\cap\mathsf{Z}=\mathsf{Z}\setminus Z''_p$. By the first
clause of the statement on the new families $Z''_j$,
\begin{equation*}
Z''_p\subset Z''_k=\mathsf{Z}\setminus\Omega_\ast^{(+5M)} .
\end{equation*}
Taking complements in $\mathsf{Z}$, and using $\ast=k_0+1$ together with
$\Omega_\ast\subset\Omega_\ast^{(+5M)}$, we obtain
\begin{equation*}
\Omega''_p\cap\mathsf{Z}=\mathsf{Z}\setminus Z''_p\supseteq\mathsf{Z}\cap\Omega_\ast^{(+5M)}
\supseteq\mathsf{Z}\cap\Omega_{k_0+1}.
\end{equation*}
Since $p+1>p$, the definition of the stage nest gives $\widehat{\Omega}_{p+1}=\Omega_{p+1}$.
Removing this set from both sides yields the first inclusion of
\eqref{4.14:eq-tiling-later-steps-1}. Taking $i=p$ in the definition of the stage nest gives
$\widehat{\Omega}_p=\Omega''_p$.
Hence the middle set lies in $\widehat{\Omega}_p\setminus\widehat{\Omega}_{p+1}$, and so in
its closure $\widehat{X}_p$. This is the second inclusion.

Off $\mathsf{Z}$ the operation does not change the domains, so $\Omega''_p=\Omega_p$. Since
$p\ge k_0+1$ and the old nest is decreasing, $\Omega_p\subset\Omega_{k_0+1}$. This inclusion
runs in the direction opposite to the first inclusion of
\eqref{4.14:eq-tiling-later-steps-1}, which is why nothing is asserted off $\mathsf{Z}$.

\emph{Part (2).} Two facts give the agreement of $\{\Omega^{\mathrm{now}}_i\}$ with
$\mathcal{N}^{(n)}$ on $X$ at all levels $\le k$. The first is Ba{\l}aban's statement in
\cite{B-CMP122a} about the domains $\Omega''_i$, ``They are not changed by the subsequent
operations.'' The second is
the persistence statement for the nests of components that have left $\mathsf{Z}$, in its
clauses (a)
and (b). The agreement lasts until a fill of a component containing $X$. After such a fill,
part (a) of Lemma~\ref{4.14:lem-gauge-family} applies at the filled levels, because part (a)
holds for every nest satisfying (A1).

Under the first reading, membership of $Q(\square)$ follows directly from
\eqref{4.14:eq-hull-cube-a}, as in part (a), and the member of \cite[(2.38)]{B-CMP119}
carries no power of $\mathsf{L}_0$. Under the second reading, part (b) applies verbatim,
because on $X$ the nest at the levels $\le k$ is that of $\mathcal{N}^{(n)}$. The exponents
of the members are then as follows:
\begin{itemize}
\item the member of \cite[(1.65)]{B-CMP122a} has exponent $\max\{0,m-k_0\}$;
\item the member of \cite[(1.67)]{B-CMP122a}, read as in case (G3) of
Lemma~\ref{4.14:lem-gauge-family} with its old-nest index $m'\ge k_0+1$ and top index
$j=p=m$, has
exponent $j-m'\le j-k_0-1$;
\item the members of \cite[(1.69)]{B-CMP122a} and \cite[(2.38)]{B-CMP119} carry no power of
$\mathsf{L}_0$.
\end{itemize}
Each of these $\mathsf{L}_0$-exponents is at most $\max\{0,m-k_0\}$. This agrees with
\cite[p.~192]{B-CMP122a}, where the highest exponent appearing is $k-k_0=N_0$, since
$m\le k$.

\emph{Part (3).} A component that dropped out at stage $n'<n$ keeps $\mathcal{N}^{(n')}$, by
the same persistence as in part (2). Part (b) of Lemma~\ref{4.14:lem-gauge-family} applies
to it with $n'$ in place of $n$, hence with $p'=h+n'$. The number $k_0=k-N_0$ depends only on
the operation, so it is unchanged.

For $Y$ the number of stages is $N_k$, so $p=h+N_k=k$. Case (G4) requires $m>p=k$, which
is impossible because $m\le k$ at the levels of the operation. When $m=k$, the index $p+1=k+1$
enters, and cases (G2) and (G3) are read at $m=k$ with the convention $\Omega_{k+1}=\emptyset$
of Lemma~\ref{4.14:lem-hull-cube}.

\emph{Part (4).} The componentwise nest is obtained by assigning to each component of
$\mathsf{Z}$ its state, so it is one of the assignments of states to the components of
$\mathsf{Z}$. By the proposition establishing conditions (A1), (A2) of
\eqref{4.14:eq-prepared-data-admissible-a} for the nests of the operation, the nest of every
assignment of states to the components of $\mathsf{Z}$ satisfies condition (A1) of that
display. The componentwise
nest is such an assignment; hence it satisfies (A1).

Next we show that the case analysis reads at most one state. Since $m\le k$, we have
$\ell_m=L^{m-k}\le1$, so the $\ell^\infty$-diameter of $Q(\square)$ satisfies
$M\ell_m\le M$. Distinct components of $Z_k$ are at $\ell^\infty$-distance at least $MR_k$,
the complement $\Omega_k^c$ is contained in $Z_k$ at level $k$, and $R_k\ge3$; hence distinct
components of $Z_k$ are at $\ell^\infty$-distance at least $MR_k\ge3M$. Suppose $Q$ meets $X'$
at a point $x$. Writing $|\cdot|$ for the
$\ell^\infty$-distance, for $y\in Q$ and $z$ in another component,
\begin{equation*}
|y-z|\ge|x-z|-|x-y|\ge MR_k-M\ell_m\ge3M-M=2M .
\end{equation*}
So $Q$ stays at distance at least $2M$ from every other component, and only the state of $X'$
is read. If $Q$ meets no component of $\mathsf{Z}$, the nest at $Q$ is the old structure,
made by $\mathbf{T}$ or by earlier fills. Part (a) then applies, since it holds on any nest
satisfying (A1). Hence, when $Q$ meets a component $X'$, the nest read at $Q$ is the global
componentwise nest with $p=p(X')$; when it meets none, part (a) is read on the old structure.
This is the final assertion.
\end{proof}

\begin{lemma}[The auxiliary integer's power against the block count]\label{4.14:lem-integer-power-bound}
Work in Ba{\l}aban's small-coupling regime \eqref{4.14:eq-small-coupling-regime-a} at a large-field
operation $\mathbf{R}^{(k)}$. Thus every block count $R_i$ is a power of $L$; for every pair of
levels $n>m$ of the history the two inequalities
\begin{equation*}
R_n\le L R_m,\qquad R_m\le (L+1)(n-m)^{\beta_0}R_n
\end{equation*}
of \cite[(2.9)]{B-CMP119} hold, $\beta_0>0$ being Ba{\l}aban's exponent; the integer $N_0$ exists;
and Ba{\l}aban's auxiliary integer $\mathsf{L}_0$ satisfies $\mathsf{L}_0\ge 1$ and
$\mathsf{L}_0^2\le \tfrac13 L$. Put $k_0:=k-N_0$ and $\ast:=k_0+1$. Then
$R_{\ast}=L^{N_0-1}$, as in the conventions \eqref{4.14:eq-operation-conventions-a}: this is the
equality $L^{-N_0+1}R_{k-N_0+1}=1$ by which $N_0$ is defined in \cite[p.~192]{B-CMP122a}. For every
level $m$ with $k_0+1\le m\le k$ and every integer $e$ with $0\le e\le m-k_0$,
\begin{equation}\label{4.14:eq-integer-power-bound-a}
\begin{gathered}
\mathsf{L}_0^{2e}\le \mathsf{L}_0^{2(m-k_0)}\le c(\beta_0)\,L(L+1)\,R_m ,\\
c(\beta_0):=\sup_{v\ge 1}3^{-v}v^{\beta_0}<\infty .
\end{gathered}
\end{equation}
In particular, taking $m=k$ and $e=N_0$,
\begin{equation}\label{4.14:eq-integer-power-bound-1}
\mathsf{L}_0^{2N_0}\le 3^{-N_0}N_0^{\beta_0}L(L+1)R_k\le c(\beta_0)\,L(L+1)\,R_k .
\end{equation}
This is a form, free of $N_0$ on the right, of the chain
\begin{equation*}
\mathsf{L}_0^{2N_0}\le \Bigl(\frac13\Bigr)^{N_0}L(L+1)(N_0+1)^{\beta_0}R_k
< L(L+1)N_0^{-1}R_k
\end{equation*}
displayed in \cite[p.~192]{B-CMP122a}.
\end{lemma}

\begin{proof}
Since $\mathsf{L}_0\ge 1$ and $e\le m-k_0$, we have $\mathsf{L}_0^{2e}\le\mathsf{L}_0^{2(m-k_0)}$.
This is the first inequality of \eqref{4.14:eq-integer-power-bound-a}.

For the second inequality put $a:=m-k_0=m-\ast+1$. This is an integer. Since
$k_0+1\le m\le k$, it satisfies $1\le a\le k-k_0=N_0$. By the condition
$\mathsf{L}_0^2\le\tfrac13 L$ and since $a-1=m-\ast$,
\begin{equation*}
\mathsf{L}_0^{2a}=(\mathsf{L}_0^2)^a\le\Bigl(\frac L3\Bigr)^a
=3^{-a}\,L\,L^{a-1}=3^{-a}\,L\,L^{m-\ast}.
\end{equation*}
Because $m\le k$ and $L\ge 1$, we have $L^{m-k}\le 1$. Moreover $k-\ast=N_0-1$, and
$L^{N_0-1}=R_{\ast}$ by the defining equality of $N_0$. Hence
\begin{equation*}
L^{m-\ast}=L^{k-\ast}\,L^{m-k}\le L^{k-\ast}=L^{N_0-1}=R_{\ast}.
\end{equation*}
We now bound $R_{\ast}$ by $R_m$, distinguishing two cases.

If $m=\ast$, then $a=1$, so $a^{\beta_0}=1$. Since $L+1\ge 1$,
\begin{equation*}
R_{\ast}=R_m\le (L+1)\,a^{\beta_0}R_m .
\end{equation*}

If $m>\ast$, apply the second inequality of \cite[(2.9)]{B-CMP119} to the pair of levels
$(n,m)$ of that display taken to be $(m,\ast)$; it applies because $m>\ast$. It reads
\begin{equation*}
R_{\ast}\le (L+1)(m-\ast)^{\beta_0}R_m=(L+1)(a-1)^{\beta_0}R_m\le (L+1)\,a^{\beta_0}R_m .
\end{equation*}
The last inequality holds because $0<a-1<a$ and $\beta_0>0$.

In both cases, therefore,
\begin{equation}\label{4.14:eq-integer-power-bound-2}
L^{m-\ast}\le (L+1)\,a^{\beta_0}R_m .
\end{equation}
Inserting \eqref{4.14:eq-integer-power-bound-2} into the bound for $\mathsf{L}_0^{2a}$, and using
the definition of $c(\beta_0)$ at the point $v=a\ge 1$, gives
\begin{equation*}
\mathsf{L}_0^{2a}\le 3^{-a}a^{\beta_0}\,L(L+1)\,R_m\le c(\beta_0)\,L(L+1)\,R_m .
\end{equation*}
Since $a=m-k_0$, this is the second inequality of \eqref{4.14:eq-integer-power-bound-a}.

For \eqref{4.14:eq-integer-power-bound-1}, take $m=k$. Then $a=k-k_0=N_0$, and the integer
$e=N_0$ is admissible. The middle member of the chain bounding $\mathsf{L}_0^{2a}$ above becomes
$3^{-N_0}N_0^{\beta_0}L(L+1)R_k$, which gives both inequalities of
\eqref{4.14:eq-integer-power-bound-1}.

It remains to show that $c(\beta_0)$ is finite. The function $v\mapsto 3^{-v}v^{\beta_0}$ is
continuous on $[1,\infty)$ and tends to $0$ as $v\to\infty$, so its supremum is finite. For
$\beta_0\le 1$ one has $3^{-v}v^{\beta_0}\le 3^{-v}v$ for $v\ge 1$. Hence in that case
$c(\beta_0)\le\sup_{v\ge 1}3^{-v}v=\tfrac13$, the supremum being attained at $v=1$. Ba{\l}aban's
choice of $\beta_0$ permits this case, since by the remark after \cite[(2.6)]{B-CMP119}
$\beta_0>0$ can be chosen arbitrarily small if $g$ is sufficiently small; the lemma does not
require $\beta_0\le 1$.
\end{proof}

\smallskip\noindent\emph{Remarks.}
\begin{enumerate}
\item Nothing beyond Ba{\l}aban's printed relations enters the proof, and only three of them. The
first is the second inequality of \cite[(2.9)]{B-CMP119}, applied once, to the pair of levels
$(m,\ast)$, which differ by $m-\ast$, in the form with the factor $(n-m)^{\beta_0}$, exactly as
Ba{\l}aban applies it in \cite[p.~192]{B-CMP122a} to a pair of levels differing by $N_0-1$.
The second is the defining equality of $N_0$. The third is $\mathsf{L}_0^2\le L/3$.

The consequence $R_{i+1}\ge R_i/L$ of \cite[(2.9)]{B-CMP119} for consecutive levels would not
suffice. Iterating it from $\ast$ to $k$ gives only
\begin{equation*}
R_k\ge R_{\ast}/L^{N_0-1}=1 .
\end{equation*}
No direction of the coupling flow is used. Neither the number $K$ of renormalisation steps, nor the
value of the level $k$ itself, nor the number $N_k$ of preliminary integrations of the operation,
nor the number of steps elapsed since any large-field region was created enters. The constants
$\beta_0$ and $\mathsf{L}_0$ are Ba{\l}aban's.

\item Suppose the upper bound $\mathsf{L}_0^2\le L/3$ were replaced by Ba{\l}aban's weaker upper
bound $\mathsf{L}_0<\tfrac12L$ of \cite{B-CMP122a}. Then the argument above, with
\eqref{4.14:eq-integer-power-bound-2}, would give
\begin{equation*}
\mathsf{L}_0^{2a}<\Bigl(\frac{L^2}{4}\Bigr)^a=4^{-a}\bigl(L\,L^{m-\ast}\bigr)^2
\le 4^{-a}\bigl[L(L+1)\,a^{\beta_0}R_m\bigr]^2 .
\end{equation*}
This is a bound of the same kind, but squared. By the upper bound on the block counts in
\eqref{4.14:eq-small-coupling-regime-a}, $R_m$ is less than $L$ times the larger of $1$ and a fixed
power of $\log g_m^{-2}$; since $a\le N_0$, both right-hand sides are therefore bounded by a factor
depending on $L$, $N_0$ and $\beta_0$ times a power of $\max\{1,\log g_m^{-2}\}$, the second being
the square of such a bound. The role of the lemma is only to exhibit a bound of this kind, and
Ba{\l}aban's condition $\mathsf{L}_0^2\le L/3$ is kept.

\item There is an alternative route. It takes Ba{\l}aban's bound at $m=k$ from
\cite[p.~192]{B-CMP122a} and then passes to a level $n<k$ by the first inequality
$R_k\le LR_n$ of \cite[(2.9)]{B-CMP119}. This is the case $m=k$ of the lemma followed by one more
factor $L$. The lemma instead applies the second inequality of \cite[(2.9)]{B-CMP119} at the level
$m$ at which the bound is needed. For this reason the bounds that use
\eqref{4.14:eq-integer-power-bound-a} hold level by level, without a maximum over levels.
\end{enumerate}

\begin{lemma}[The top power against every lower level]\label{4.14:lem-top-power-bound}
Let the hypotheses of Lemma~\ref{4.14:lem-integer-power-bound} hold. Thus a large-field
operation $\mathbf{R}^{(k)}$ acts at level $k$ in Ba{\l}aban's small-coupling regime
\eqref{4.14:eq-small-coupling-regime-a} of Definition~\ref{4.14:def-small-coupling-regime}.
Let $N_0$, $k_0=k-N_0$ and $\ast=k_0+1$ be as in Convention~\ref{4.14:not-operation-conventions},
so that $R_\ast=L^{N_0-1}$ by \eqref{4.14:eq-operation-conventions-a}. Put
\begin{equation*}
c(\beta_0):=\sup_{t\ge1}3^{-t}t^{\beta_0},
\qquad
\mathsf{c}_1:=1+c(\beta_0)L(L+1),
\end{equation*}
where $c(\beta_0)$ is the constant of \eqref{4.14:eq-integer-power-bound-a}, and $\mathsf{c}_1$ is the
constant of Definition~\ref{4.14:def-block-size-conditions}. Then for every level $n\le k$
\begin{equation}\label{4.14:eq-top-power-bound-a}
\mathsf{L}_0^{2N_0}\le \mathsf{c}_1R_n .
\end{equation}
Consequently, for every level $n\le k$,
\begin{equation}\label{4.14:eq-top-power-bound-1}
\mathsf{L}_0^{2N_0}\,L^2BM\alpha_{0,n}\le L^2\mathsf{c}_1BM\,R_n\alpha_{0,n}.
\end{equation}
The left-hand side of \eqref{4.14:eq-top-power-bound-1} is the member at the layer $n$ of the
per-operation constant $a^{\sharp}_F$ of the operation $\mathbf{R}^{(k)}$, where $B$ is the constant of
the operator format in which that constant is written; the right-hand side is the boxwise value
$L^2\mathsf{c}_1BMR_n\alpha_{0,n}$ at level $n$. In words: the member of $a^{\sharp}_F$ at the layer $n$
is at most the boxwise value at level $n$. This holds for the fixed operation $\mathbf{R}^{(k)}$ and
for the layers $n\le k$ of that operation.
\end{lemma}

\begin{proof}
By \eqref{4.14:eq-small-coupling-regime-a} we have $\mathsf{L}_0\ge1$ and
$\mathsf{L}_0^2\le L/3$. Raising the second inequality, between non-negative numbers, to the
power $N_0$ and using $R_\ast=L^{N_0-1}$ from \eqref{4.14:eq-operation-conventions-a} gives
\begin{equation}\label{4.14:eq-top-power-bound-2}
\mathsf{L}_0^{2N_0}=(\mathsf{L}_0^2)^{N_0}\le\Big(\frac L3\Big)^{N_0}
=3^{-N_0}\,L\,L^{N_0-1}=3^{-N_0}\,L\,R_\ast .
\end{equation}
It remains to compare $R_\ast$ with $R_n$. The tool is \cite[(2.9), p.~256]{B-CMP119}, which is part
of \eqref{4.14:eq-small-coupling-regime-a}. For every pair of levels $j>i$ of the history it states
\begin{equation}\label{4.14:eq-top-power-bound-3}
R_{j}\le L\,R_{i}\qquad\text{and}\qquad R_{i}\le (L+1)(j-i)^{\beta_0}R_{j} .
\end{equation}
Here $R_{j}\ge1$ for every level $j$, by \eqref{4.14:eq-small-coupling-regime-a}, and $\beta_0>0$ is
the exponent of \cite[(2.9)]{B-CMP119}.
Two elementary facts about $c(\beta_0)$ are used below. First, for every integer $t\ge1$ we have
$3^{-t}t^{\beta_0}\le c(\beta_0)$, by the definition of the supremum. Second, the value at $t=1$
gives $c(\beta_0)\ge3^{-1}\cdot1^{\beta_0}=\tfrac13$. We distinguish three ranges of $n$.

\emph{Range $\ast<n\le k$.} Apply the second inequality of \eqref{4.14:eq-top-power-bound-3} to the
pair $(j,i)=(n,\ast)$. This is admissible, since $n>\ast$. It gives
\begin{equation*}
R_\ast\le (L+1)(n-\ast)^{\beta_0}R_n ,
\end{equation*}
where
\begin{equation*}
1\le n-\ast\le k-\ast=N_0-1<N_0 .
\end{equation*}
Since $\beta_0>0$, we get $(n-\ast)^{\beta_0}\le N_0^{\beta_0}$. Inserting this into
\eqref{4.14:eq-top-power-bound-2} gives
\begin{align*}
\mathsf{L}_0^{2N_0}
&\le 3^{-N_0}N_0^{\beta_0}\,L(L+1)\,R_n\\
&\le c(\beta_0)L(L+1)R_n\\
&\le \mathsf{c}_1R_n .
\end{align*}
The second inequality holds because $N_0$ is an integer with $N_0\ge1$, so that
$3^{-N_0}N_0^{\beta_0}\le c(\beta_0)$. The third holds because
$\mathsf{c}_1=1+c(\beta_0)L(L+1)$ and $R_n\ge0$.

\emph{Range $n=\ast$.} Here $R_\ast=R_n$. Since $N_0\ge1$ and $c(\beta_0)\ge\frac13$, we have
\begin{equation*}
3^{-N_0}L\le\frac L3\le c(\beta_0)L(L+1)<\mathsf{c}_1 .
\end{equation*}
By \eqref{4.14:eq-top-power-bound-2} and $R_n\ge1$ this gives
$\mathsf{L}_0^{2N_0}\le 3^{-N_0}LR_n<\mathsf{c}_1R_n$.

\emph{Range $n<\ast$.} Apply the first inequality of \eqref{4.14:eq-top-power-bound-3} to the pair
$(j,i)=(\ast,n)$. This is admissible, since $\ast>n$. It gives $R_\ast\le LR_n$. Inserting this
into \eqref{4.14:eq-top-power-bound-2} gives
\begin{align*}
\mathsf{L}_0^{2N_0}
&\le 3^{-N_0}L^2R_n\le\frac{L^2}{3}R_n\\
&\le c(\beta_0)L(L+1)R_n<\mathsf{c}_1R_n .
\end{align*}
The third inequality uses $c(\beta_0)\ge\frac13$ and $L(L+1)\ge L^2$. The last one uses
$R_n\ge1$.

Each of the three arguments holds for every integer $N_0\ge1$. Here, in fact, $N_0\ge2$ by
Convention~\ref{4.14:not-operation-conventions}. The three ranges exhaust all levels $n\le k$,
which proves \eqref{4.14:eq-top-power-bound-a}.

Finally, $L^2BM\alpha_{0,n}\ge0$. Multiplying \eqref{4.14:eq-top-power-bound-a} by this number gives
\eqref{4.14:eq-top-power-bound-1}.
\end{proof}

We add two observations, on the scope of Lemma~\ref{4.14:lem-top-power-bound} and on the inputs
of its proof.

The pairing of the power $\mathsf{L}_0^{2N_0}$, which belongs to the operation $\mathbf{R}^{(k)}$,
with the radius $\alpha_{0,n}$ is asserted only for $n\le k$, that is, for the operation's own
levels and the levels below them. This is exactly the range in which
\cite[(1.64)--(1.67)]{B-CMP122a} pairs the powers $\mathsf{L}_0^{2e}$, $0\le e\le m-k_0$, of
Lemma~\ref{4.14:lem-integer-power-bound} with the radii $\alpha_{0,m},\alpha_{1,m}$ for
$m\in(k_0,k]$, together with the levels below.

For a later layer $n>k$ of a later nest, the second inequality of
\eqref{4.14:eq-top-power-bound-3} gives only
\begin{equation*}
R_\ast\le(L+1)(n-\ast)^{\beta_0}R_n ,
\end{equation*}
whose factor $(n-\ast)^{\beta_0}$ grows with the gap $n-\ast$. No bound of the form
\eqref{4.14:eq-top-power-bound-a} is claimed for such $n$.

The value adopted for the gauge constant is the boxwise value $\mathfrak{a}(\square)$. It reads the
block count $R_{m(\square)}$ of each hull's own level, and so it requires no pairing across
operations or across layers.

The proof above uses only the following inputs:
\begin{itemize}
\item the second inequality of \cite[(2.9)]{B-CMP119}, at gaps smaller than $N_0$;
\item the first inequality of \cite[(2.9)]{B-CMP119}, once;
\item the bound $\mathsf{L}_0^2\le L/3$;
\item the defining equality $R_\ast=L^{N_0-1}$ of $N_0$.
\end{itemize}
It uses no monotonicity of the coupling flow. It reads neither $k$, nor $K$, nor the number $N_k$
of preliminary integrations, nor the age of any large-field region, nor the number of large-field
operations.

\begin{corollary}[The gauge constants on a hull]\label{4.14:cor-hull-gauge-value}
Work in the setting of case~(a) or case~(b) of Lemma~\ref{4.14:lem-gauge-family}, with
$\square$, $H=\square^{4}$, $m=m(\square)$, $Q=Q(\square)$ and the territories
$\widehat{X}_n$ of the nest $\{\widehat{\Omega}_i\}$ as there. Let $u$ be the $G$-valued gauge
on $Q$ that the clause of \cite[(2.38)]{B-CMP119} (case~(a)), resp.\ of
\cite[(1.65), (1.67), (1.69)]{B-CMP122a} (case~(b)), provides for the printed member of the
family to which $Q$ belongs. When the analyticity space in which that member is read carries a
localisation domain, which we denote $Y_{\mathrm{loc}}$, $u$ is given on
$Q\cap Y_{\mathrm{loc}}$ only. Write $U^{u}=e^{i\eta A}$ as in
Definition~\ref{4.14:def-real-backgrounds}. Let $B$ be the constant of the printed members, and
let $C$ be the size constant of the printed families, taken equal to $1$. Put
\begin{equation}\label{4.14:eq-hull-gauge-value-a}
\mathfrak{a}(\square):=L^{2}\mathsf{c}_1 B M R_{m(\square)}\alpha_{0,m(\square)},
\qquad \mathsf{c}_1:=1+c(\beta_0)L(L+1),
\end{equation}
where $c(\beta_0):=\sup_{t\ge1}3^{-t}t^{\beta_0}$. Then, for each $n\in\{m-1,m,m+1\}$ such that
$H$ meets $\widehat{X}_n$,
\begin{equation}\label{4.14:eq-hull-gauge-value-1}
\ell_n|A|<\mathfrak{a}(\square),\qquad \ell_n^{2}|\nabla A|<\mathfrak{a}(\square)
\qquad\text{on } H\cap\widehat{X}_n,
\end{equation}
and on $H\cap Y_{\mathrm{loc}}\cap\widehat{X}_n$ in the localised case. Consequently
$(\mathrm{reg})^{\sharp}$
of Definition~\ref{4.14:def-real-backgrounds} holds at $\square$ for every
$a^{\sharp}\ge\mathfrak{a}(\square)$.

Moreover, in case~(b) (in case~(a) one has $e(Q)=0$) let $\mathbf{R}^{(k)}$ be the operation,
with its $N_0$ and $k_0=k-N_0$. Consider the per-operation constant $a^{\sharp}_F$, whose value
at the level $n$ is $\mathsf{L}_0^{2N_0}L^{2}BM\alpha_{0,n}$. Then:
\begin{itemize}
\item[(i)] it dominates the printed member of every cube $Q=Q(\square)$ of the operation, since
$e(Q)\le N_0$;
\item[(ii)] for every level $n\le k$ of the operation its value at the level $n$ is at most the
boxwise value $L^{2}\mathsf{c}_1BMR_n\alpha_{0,n}$ at the level $n$;
\item[(iii)] at the hull, with $m=m(\square)\le k$,
\begin{equation}\label{4.14:eq-hull-gauge-value-2}
\mathsf{L}_0^{2e(Q)}L^{2}BM\alpha_{0,m}\le \mathsf{L}_0^{2N_0}L^{2}BM\alpha_{0,m}
\le L^{2}\mathsf{c}_1BMR_m\alpha_{0,m}=\mathfrak{a}(\square).
\end{equation}
\end{itemize}
\end{corollary}

\begin{proof}
Write $e=e(Q)$. By Lemma~\ref{4.14:lem-gauge-family}, in each of its cases the printed member
at $Q$ has $\alpha$-index $m=m(\square)$, scale $\ell_m$ and $\mathsf{L}_0$-exponent $e$. Thus
the gauge $u$ satisfies
\begin{equation}\label{4.14:eq-hull-gauge-value-3}
\ell_m|A|<\mathsf{L}_0^{2e}BCM\alpha_{0,m},\qquad
\ell_m^{2}|\nabla A|<\mathsf{L}_0^{2e}BCM\alpha_{0,m}
\end{equation}
on $Q$, resp.\ on $Q\cap Y_{\mathrm{loc}}$. Since $H\subset Q$ by
Lemma~\ref{4.14:lem-hull-cube}\,(4), the bounds~\eqref{4.14:eq-hull-gauge-value-3} hold on $H$,
resp.\ on $H\cap Y_{\mathrm{loc}}$.

\emph{The weights.} By Lemma~\ref{4.14:lem-hull-cube}\,(2), $H$ meets only the territories
$\widehat{X}_m$ and $\widehat{X}_{m+1}$, together with $\widehat{X}_{m-1}$ at most along
$\partial\widehat{\Omega}_m$. Hence the levels $n$ to be treated lie in $\{m-1,m,m+1\}$.

For $n=m+1$ we have $\ell_{m+1}=L\ell_m$, because $\ell_j=L^{j-k}$. Therefore
\begin{equation*}
\ell_{m+1}|A|=L\,\ell_m|A|\le L^{2}\,\ell_m|A|,\qquad
\ell_{m+1}^{2}|\nabla A|=L^{2}\,\ell_m^{2}|\nabla A|.
\end{equation*}
For $n\in\{m-1,m\}$ we have $\ell_n\le\ell_m$, so no factor is lost. In all three cases
\eqref{4.14:eq-hull-gauge-value-3} gives
\begin{equation*}
\ell_n|A|<L^{2}\mathsf{L}_0^{2e}BCM\alpha_{0,m},\qquad
\ell_n^{2}|\nabla A|<L^{2}\mathsf{L}_0^{2e}BCM\alpha_{0,m}
\end{equation*}
on $H\cap\widehat{X}_n$, resp.\ on $H\cap Y_{\mathrm{loc}}\cap\widehat{X}_n$.

\emph{The power.} Suppose first that $e=0$. Then $\mathsf{L}_0^{2e}=1\le\mathsf{c}_1$, since
$c(\beta_0)\ge0$. Also $\mathsf{c}_1\le\mathsf{c}_1R_m$, since $R_m\ge1$.

Suppose next that $e\ge1$. In case~(a) the exponent is $e=0$, so we are in case~(b). By the
exponent bound~\eqref{4.14:eq-gauge-family-a} of Lemma~\ref{4.14:lem-gauge-family}, $e>0$
forces $k_0+1\le m\le k$, and moreover $e\le\max\{0,m-k_0\}=m-k_0$. These are the conditions on
$m$ and $e$ in Lemma~\ref{4.14:lem-integer-power-bound}, whose regime is that of the operation
$\mathbf{R}^{(k)}$. Its bound \eqref{4.14:eq-integer-power-bound-a} gives
\begin{equation*}
\mathsf{L}_0^{2e}\le c(\beta_0)L(L+1)R_m\le\mathsf{c}_1R_m .
\end{equation*}

In both cases, therefore, the two quantities in~\eqref{4.14:eq-hull-gauge-value-1} are strictly
smaller than
\begin{equation*}
L^{2}\mathsf{c}_1BCMR_m\alpha_{0,m}=\mathfrak{a}(\square)\qquad(C=1).
\end{equation*}
This proves~\eqref{4.14:eq-hull-gauge-value-1}. Clause $(\mathrm{reg})^{\sharp}$ of
Definition~\ref{4.14:def-real-backgrounds} asks for such a gauge with these two quantities
bounded by the single upper bound $a^{\sharp}$ on each territory met by $\square^{4}$. The strict
bounds~\eqref{4.14:eq-hull-gauge-value-1} and $\mathfrak{a}(\square)\le a^{\sharp}$ give that
clause at $\square$.

\emph{The relation with $a^{\sharp}_F$.} In case~(b) we argue as follows.

For (i): the exponent bound~\eqref{4.14:eq-gauge-family-a} gives $e\le\max\{0,m-k_0\}$. Every
level of the operation satisfies $m\le k$, so $m-k_0\le k-k_0=N_0$ and hence $e\le N_0$. Since
$\mathsf{L}_0\ge1$, this gives $\mathsf{L}_0^{2e}\le\mathsf{L}_0^{2N_0}$. Multiplying this
inequality by $L^{2}BM\alpha_{0,m}$ gives the first inequality
in~\eqref{4.14:eq-hull-gauge-value-2}. Multiplying it by $BM\alpha_{0,m}$ alone and using
$L^{2}\ge1$ gives
\begin{equation*}
\mathsf{L}_0^{2e}BM\alpha_{0,m}\le\mathsf{L}_0^{2N_0}L^{2}BM\alpha_{0,m},
\end{equation*}
which is (i).

For (ii): Lemma~\ref{4.14:lem-top-power-bound} gives $\mathsf{L}_0^{2N_0}\le\mathsf{c}_1R_n$
for every $n\le k$, by \eqref{4.14:eq-top-power-bound-a}. Multiplying by $L^{2}BM\alpha_{0,n}$
gives (ii).

For (iii): taking $n:=m\le k$ in (ii) gives the second inequality
in~\eqref{4.14:eq-hull-gauge-value-2}. The final equality is the
definition~\eqref{4.14:eq-hull-gauge-value-a}.
\end{proof}

\begin{proposition}[The second gauge clause on admissible nests]\label{4.14:prop-second-gauge-clause}
Assume the floor $M\ge 13L^2M_1^{\sharp}$ of Definition~\ref{4.14:def-block-size-conditions}.
Let $\{\widehat{\Omega}_i\}_{i\le \hat k}$ be a nest of one of the kinds covered by the admissibility
lemma for Ba{\l}aban's nests, described in the terms of that lemma (by that lemma, each such nest
satisfies the nesting condition of Lemma~\ref{4.14:lem-hull-cube}):
\begin{itemize}
\item a nest made by the operations $\mathbf{T}$;
\item a hybrid stage nest $\mathcal{N}^{(n)}$ of a drop-out at the stage, or at any later step
$\mathbf{T}$ before a fill;
\item a nest with filled islands.
\end{itemize}
Let $j\ge 2$ and
$\square\in\mathcal{D}^{\sharp}_j$, put $H:=\square^{4}$,
and let $m(\square)$ and $Q(\square)$ be as in \eqref{4.14:eq-hull-cube-a}. Then the following hold.
\begin{itemize}
\item[(G)] The cube $Q(\square)$ of side $M\ell_{m(\square)}$ of Lemma~\ref{4.14:lem-hull-cube} exists.
\item[(C)] $Q(\square)$ is a cube, with size constant $C=1$, of each family that Ba{\l}aban displays
at the index $m(\square)$ for this nest: the family of \cite[(2.38)]{B-CMP119} read on
$\{\widehat{\Omega}_i\}$; and, writing $m=m(\square)$, the families of
\cite[(1.65)(i)]{B-CMP122a} in the case $m<j$ and in
the case $m=j$, of \cite[(1.67)(ii)]{B-CMP122a} at the old-nest index $m'$, and of
\cite[(1.69)(iii)]{B-CMP122a} at the stage nest. For the operation that made the level, the
$\mathsf{L}_0$-exponent obeys
\begin{equation*}
e(Q)\le \max\{0,\,m(\square)-k_0\}.
\end{equation*}
\item[(V)] Let $u$ be the gauge, with $U^{u}=e^{i\eta A}$ as in clause $(\mathrm{reg})^{\sharp}$ of
\eqref{4.14:eq-real-backgrounds-a}, that \cite[(2.38)]{B-CMP119}, resp.\
\cite[(1.65), (1.67), (1.69)]{B-CMP122a}, provides on $Q(\square)$ as a cube of its family.
Weighted to the territories met by $H$, its potential $A$ obeys, on $H\cap\widehat{X}_n$ for each
territory $\widehat{X}_n$ that $H$ meets,
\begin{equation}\label{4.14:eq-second-gauge-clause-1}
\ell_n|A|<\mathfrak{a}(\square),\qquad \ell_n^{2}|\nabla A|<\mathfrak{a}(\square),\qquad
\mathfrak{a}(\square)=L^{2}\mathsf{c}_1BM\,R_{m(\square)}\,\alpha_{0,m(\square)} ,
\end{equation}
with $\mathsf{c}_1$ and $B$ as in \eqref{4.14:eq-hull-gauge-value-a}. Hence
$(\mathrm{reg})^{\sharp}$ holds at $\square$ with
\begin{equation}\label{4.14:eq-second-gauge-clause-2}
\begin{split}
a^{\sharp}&:=\max\{\mathfrak{a}(\square'):\ \square'\in\mathcal{D}^{\sharp}_{j'},\ j'\ge 2\}\\
&\le L^{2}\mathsf{c}_1BM\,\max_{m\le\hat k}R_m\alpha_{0,m}.
\end{split}
\end{equation}
The ceiling $a^{\sharp}\le a^{\sharp}_{\mathrm{ceil}}$ of \eqref{4.14:eq-block-size-conditions-b},
assumed in the operator theorem at free current (Theorem~\ref{4.14:thm-operator-theorem}), is met
whenever
\begin{equation}\label{4.14:eq-second-gauge-clause-3}
L^{2}\mathsf{c}_1B\cdot M\cdot R_m\alpha_{0,m}\le a^{\sharp}_{\mathrm{ceil}}
\qquad\text{for every level } m\in\{1,\dots,\hat k\} ,
\end{equation}
which is the ceiling \eqref{4.14:eq-block-size-conditions-c}, read with the top index $\hat k$ of
the nest; conversely, if the ceiling is met,
the inequality in \eqref{4.14:eq-second-gauge-clause-3} holds at every level $m$ that occurs as
$m(\square')$ for some $\square'$ in the range of the maximum in
\eqref{4.14:eq-second-gauge-clause-2}.
\end{itemize}
For $\square\in\mathcal{D}^{\sharp}_j$ with $j\in\{0,1\}$, the clause $(\mathrm{reg})^{\sharp}$ is
supplied in the form $(\alpha')$ of \eqref{4.14:eq-real-backgrounds-a}.
\end{proposition}

\begin{proof}
\emph{(G).} The nest satisfies the hypotheses of the hull lemma (Lemma~\ref{4.14:lem-hull-cube}):
its nesting condition holds by the admissibility lemma for Ba{\l}aban's nests, the floor
$M\ge 13L^2M_1^{\sharp}$ is
assumed, and $\square\in\mathcal{D}^{\sharp}_j$ with $j\ge 2$. The hull lemma therefore
provides the level $m(\square)$ and the cube $Q(\square)$ of side $M\ell_{m(\square)}$ of
\eqref{4.14:eq-hull-cube-a}.

\emph{(C).} Under the same hypotheses, Lemma~\ref{4.14:lem-gauge-family} applies to $Q(\square)$.
Its case~(a) places $Q(\square)$, at $C=1$, in the family of \cite[(2.38)]{B-CMP119} read on
$\{\widehat{\Omega}_i\}$. Its cases (b1)--(b4) place $Q(\square)$ in the families of
\cite[(1.65)(i), (1.67)(ii), (1.69)(iii)]{B-CMP122a} at the indices listed in (C). The exponent bound
$e(Q)\le\max\{0,m(\square)-k_0\}$ is \eqref{4.14:eq-gauge-family-a}.

\emph{(V).} Corollary~\ref{4.14:cor-hull-gauge-value}, applied in the setting of case~(a) or~(b) of
Lemma~\ref{4.14:lem-gauge-family} at $C=1$, gives \eqref{4.14:eq-second-gauge-clause-1} with the
value \eqref{4.14:eq-hull-gauge-value-a}. The quantity $\mathfrak{a}(\square)$ depends on $\square$
only through $m(\square)\in\{1,\dots,\hat k\}$, so the maximum in \eqref{4.14:eq-second-gauge-clause-2}
is taken over finitely many values, each of which, by the formula for $\mathfrak{a}(\square)$ in
\eqref{4.14:eq-second-gauge-clause-1}, is at most $L^{2}\mathsf{c}_1BM\max_{m\le\hat k}R_m\alpha_{0,m}$;
this is the inequality in \eqref{4.14:eq-second-gauge-clause-2}. Thus $a^{\sharp}$ is a number with
$a^{\sharp}\ge\mathfrak{a}(\square)$ for every $\square$, and by \eqref{4.14:eq-second-gauge-clause-1}
the gauge on each hull satisfies
\begin{equation*}
\ell_n|A|<a^{\sharp},\qquad \ell_n^{2}|\nabla A|<a^{\sharp}.
\end{equation*}
By Definition~\ref{4.14:def-real-backgrounds}, $a^{\sharp}$ in $(\mathrm{reg})^{\sharp}$ is one
number and an upper bound: a gauge whose constants on $H$ are smaller than $a^{\sharp}$ satisfies the
clause. Hence $(\mathrm{reg})^{\sharp}$ holds at $\square$ with the value
\eqref{4.14:eq-second-gauge-clause-2}. Finally, $a^{\sharp}=\max_{\square}\mathfrak{a}(\square)$ by
definition, and
\begin{align*}
\max_{\square}\mathfrak{a}(\square)\le a^{\sharp}_{\mathrm{ceil}}
&\iff \mathfrak{a}(\square)\le a^{\sharp}_{\mathrm{ceil}}\ \text{for every }\square\\
&\iff L^{2}\mathsf{c}_1B\cdot M\cdot R_m\alpha_{0,m}\le a^{\sharp}_{\mathrm{ceil}}
\ \text{for every } m \text{ occurring as some } m(\square).
\end{align*}
The first equivalence holds because a maximum of finitely many numbers is at most
$a^{\sharp}_{\mathrm{ceil}}$ exactly when each of them is. The second is the formula for
$\mathfrak{a}(\square)$ in \eqref{4.14:eq-second-gauge-clause-1}. The levels that occur as
$m(\square)$ form a subset of $\{1,\dots,\hat k\}$. Hence \eqref{4.14:eq-second-gauge-clause-3},
which asks the inequality at every level $m\in\{1,\dots,\hat k\}$, implies the condition on the last
line and
so the ceiling $a^{\sharp}\le a^{\sharp}_{\mathrm{ceil}}$; and the ceiling implies the inequality of
\eqref{4.14:eq-second-gauge-clause-3} at every level that occurs as some $m(\square)$. This is the
relation between the ceiling $a^{\sharp}\le a^{\sharp}_{\mathrm{ceil}}$ of
\eqref{4.14:eq-block-size-conditions-b} and \eqref{4.14:eq-block-size-conditions-c} asserted in (V).

For $j\in\{0,1\}$ there is nothing to prove: at such boxes $(\mathrm{reg})^{\sharp}$ is supplied in
the form $(\alpha')$ of \eqref{4.14:eq-real-backgrounds-a}.
\end{proof}

\begin{remark}[The layerwise ceiling as a small-coupling condition]\label{4.14:rem-layerwise-ceiling}
Let $\beta_0>0$ be the exponent of \cite[(2.6)--(2.9)]{B-CMP119}. Let $r$ be the exponent of
\cite[(2.5)]{B-CMP119}, let $q_0>1$ be the integer exponent of \cite[(2.28)]{B-CMP119}, and let
$C_0$ be the constant of \cite[(2.28)]{B-CMP119}. Let $B$ be Ba{\l}aban's absolute constant of
\cite[(2.38)]{B-CMP119} and \cite[(1.65)]{B-CMP122a}. Let $\mathsf{c}_1=1+c(\beta_0)L(L+1)$ with
$c(\beta_0)=\sup_{t\ge1}3^{-t}t^{\beta_0}$, and let $a^{\sharp}_{\mathrm{ceil}}$ be the ceiling of
Definition~\ref{4.14:def-block-size-conditions}, a positive number depending on
$d,L,M_1,M_1^{\sharp}$ only. For a real variable $\eta$ with $0<\eta<1$ put
\begin{equation}\label{4.14:eq-layerwise-ceiling-1}
\Phi(\eta):=L^{3}\mathsf{c}_1BC_0\,\eta\,(\log \eta^{-2})^{q_0}\max\{1,(\log \eta^{-2})^{r}\}.
\end{equation}
\begin{enumerate}
\item[(a)] For every level $m$ with $g_m<1$,
\begin{equation}\label{4.14:eq-layerwise-ceiling-2}
L^2\mathsf{c}_1BM\,R_m\alpha_{0,m}<M\,\Phi(g_m).
\end{equation}
Since $\Phi(\eta)\to0$ as $\eta\to0$ (the factor $\eta$ against powers of $\log\eta^{-2}$), there
is a number $g^{\sharp}(d,L,M)\in(0,1)$ such that $M\Phi(\eta)\le a^{\sharp}_{\mathrm{ceil}}$ for
all $0<\eta\le g^{\sharp}(d,L,M)$. For any such number, the layerwise ceiling
\eqref{4.14:eq-block-size-conditions-c} holds whenever $g_m\le g^{\sharp}(d,L,M)$ for every
level $m$.
\item[(b)] The layerwise ceiling is one-sided: it is an upper bound on each level's own coupling
$g_m$. It enters the order of constants \eqref{4.14:eq-block-size-conditions-d} after $d$, $L$,
$M_1$, $M_1^{\sharp}$, all $(d,L)$-constants and $M$, among the other conditions on the couplings
of the levels. It involves none of $k$, $K$, $N_k$, $N_0$, the number of steps since a large-field
region was created, a depth, the number $|\mathsf{Z}_0|$ of cubes in a core $\mathsf{Z}_0$, or the
number of operations or of drop-outs. For every
fixed $L$ and $M$ it can be met by taking the couplings small.
\item[(c)] The condition is uniform over histories. Assume Ba{\l}aban's small-coupling regime of
Definition~\ref{4.14:def-small-coupling-regime} at each operation $\mathbf{R}^{(k')}$ considered.
A nest of domains at a late renormalisation
step may contain large-field regions left by several operations $\mathbf{R}^{(k')}$, each with its
own integer $N_0(k')$. The gauge constant $\mathfrak{a}(\square)$ of
\eqref{4.14:eq-hull-gauge-value-a} depends only on the level $m(\square)$ of the hull, through
$R_{m(\square)}\alpha_{0,m(\square)}$.
\item[(d)] Assume the regime of Definition~\ref{4.14:def-small-coupling-regime} at an operation
$\mathbf{R}^{(k)}$, with its integer $N_0$ and with $k_0=k-N_0$ and $\ast=k_0+1$ as there, and put
\begin{equation}\label{4.14:eq-layerwise-ceiling-3}
a^{\sharp}_F:=\mathsf{L}_0^{2N_0}L^2BM\max_{n\le k}\alpha_{0,n}.
\end{equation}
Then $a^{\sharp}_F\ge\mathsf{L}_0^{2e}L^2BM\alpha_{0,n}$ for all $0\le e\le N_0$ and all $n\le k$.
Under the layerwise ceiling, moreover, $a^{\sharp}_F\le a^{\sharp}_{\mathrm{ceil}}$. Nothing is
asserted about products of $\mathsf{L}_0^{2N_0}$ with $\alpha_{0,n}$ for levels $n>k$.
\end{enumerate}
\end{remark}

\begin{proof}
(a) By \cite[(2.5)]{B-CMP119}, $R_m$ is the smallest power of $L$ with exponent in $\mathbb{N}$
such that $R_m\ge(\log g_m^{-2})^{r}$. We first show that $R_m<L\max\{1,(\log g_m^{-2})^r\}$.
\begin{itemize}
\item If $R_m=1$, then $R_m<L$.
\item If $R_m\ge L$, then $R_m/L$ is again a power of $L$ with exponent in $\mathbb{N}$. By
minimality, $R_m/L<(\log g_m^{-2})^r$, so $R_m<L(\log g_m^{-2})^{r}$.
\end{itemize}
In both cases $R_m<L\max\{1,(\log g_m^{-2})^{r}\}$.
By \cite[(2.28)]{B-CMP119}, $\alpha_{0,m}=C_0g_m(\log g_m^{-2})^{q_0}$, which is positive for
$0<g_m<1$. Multiplying the bound on $R_m$ by $L^2\mathsf{c}_1BM\alpha_{0,m}>0$ gives
\begin{align*}
L^2\mathsf{c}_1BM\,R_m\alpha_{0,m}
&<L^3\mathsf{c}_1BC_0M\,g_m(\log g_m^{-2})^{q_0}\max\{1,(\log g_m^{-2})^{r}\}\\
&=M\,\Phi(g_m),
\end{align*}
which is \eqref{4.14:eq-layerwise-ceiling-2}.

As $\eta\to0$, the factor $\eta$ in \eqref{4.14:eq-layerwise-ceiling-1} beats every power of
$\log \eta^{-2}$, so $M\Phi(\eta)\to0$. Since $a^{\sharp}_{\mathrm{ceil}}>0$, there is
$\delta\in(0,1)$ with $M\Phi(\eta)\le a^{\sharp}_{\mathrm{ceil}}$ for $0<\eta<\delta$. Take
$g^{\sharp}(d,L,M):=\delta/2$.
If now $g_m\le g^{\sharp}(d,L,M)$ for every level $m$, then \eqref{4.14:eq-layerwise-ceiling-2}
gives $L^2\mathsf{c}_1BMR_m\alpha_{0,m}<a^{\sharp}_{\mathrm{ceil}}$ at every level. This is
\eqref{4.14:eq-block-size-conditions-c}.

This is the condition of \cite{B-CMP99b} that $M\alpha_0$ be sufficiently small, with the
additional logarithmic factor $R_m$. Ba{\l}aban writes in \cite[p.~192]{B-CMP122a}:
\begin{quote}
Thus the powers of $\mathsf{L}_0$ make only unessential logarithmic contributions to the constants
$\alpha_{0,j}$, $\alpha_{1,j}$, $\varepsilon_j$ in the definitions.
\end{quote}
This sentence is printed without proof in \cite[p.~192]{B-CMP122a}. Its layerwise form
\eqref{4.14:eq-block-size-conditions-c} is proved here from \cite[(2.9)]{B-CMP119}: those powers
are bounded level by level by the logarithmic factor $\mathsf{c}_1R_m$
(Lemmas~\ref{4.14:lem-integer-power-bound} and~\ref{4.14:lem-top-power-bound}).

(b) The condition $g_m\le g^{\sharp}(d,L,M)$ bounds each level's coupling from above only. The
numbers $\beta_0,r,q_0,C_0,B,\mathsf{L}_0$ are Ba{\l}aban's fixed numbers. The factor $L$ enters
only through
$\mathsf{c}_1$, the weight $L^2$ and the factor $L$ in the bound on $R_m$, all fixed before the
couplings. Likewise $a^{\sharp}_{\mathrm{ceil}}$ is fixed before $M$. Hence $g^{\sharp}(d,L,M)$
depends on $d,L,M$ and the constants fixed before them, and on nothing read from a history. It
bounds no block size and no number of steps, whether of the whole flow or since the creation of a
large-field region.

The only inequalities between levels used here and in the lemmas cited are those of
\cite[(2.9)]{B-CMP119}, displayed in \eqref{4.14:eq-small-coupling-regime-a}:
\begin{equation*}
R_n\le L\,R_m,\qquad R_m\le(L+1)(n-m)^{\beta_0}R_n\qquad(n>m),
\end{equation*}
for any two levels of one history. They serve as comparisons between the block counts $R$. No
rate of decrease of the couplings along the flow is used.

(c) Consider a hull $\square$ whose level $m(\square)$ lies among the levels of an operation
$\mathbf{R}^{(k')}$. Lemma~\ref{4.14:lem-integer-power-bound} is applied with that operation's
own indices: put
$k_0':=k'-N_0(k')$ and $\ast':=k_0'+1$. Its conclusion \eqref{4.14:eq-integer-power-bound-a}
involves only $R_{m(\square)}$, and hence the bound of Corollary~\ref{4.14:cor-hull-gauge-value}
at $\square$ involves only $R_{m(\square)}$. Hence, whenever $g_{m(\square)}\le g^{\sharp}(d,L,M)$,
part (a) gives
\begin{equation*}
\mathfrak{a}(\square)=L^2\mathsf{c}_1BM\,R_{m(\square)}\alpha_{0,m(\square)}
< a^{\sharp}_{\mathrm{ceil}}
\end{equation*}
for every hull, whatever operation its level belongs to. No supremum over operations is taken.
No maximum is taken over levels carrying
different couplings, and no direction of the flow enters.

(d) By the regime, $\mathsf{L}_0\ge1$; since $e\le N_0$, we have
$\mathsf{L}_0^{2e}\le\mathsf{L}_0^{2N_0}$. The
first claim follows by multiplying by $L^2BM\alpha_{0,n}\le L^2BM\max_{n'\le k}\alpha_{0,n'}$.

One route to a ceiling for $a^{\sharp}_F$ goes as follows. Lemma~\ref{4.14:lem-integer-power-bound}
at $m=k$, where $k-k_0=N_0$, gives $\mathsf{L}_0^{2N_0}\le c(\beta_0)L(L+1)R_k$. The first
inequality of \cite[(2.9)]{B-CMP119} with the pair of levels $k>n$ gives
$R_k\le LR_n$. Together these yield a ceiling of the same small-coupling kind, with the maximum
over the operation's levels.

Lemma~\ref{4.14:lem-top-power-bound} gives a sharper route. It
gives \eqref{4.14:eq-top-power-bound-a}, $\mathsf{L}_0^{2N_0}\le\mathsf{c}_1R_n$ for every
$n\le k$, with no further factor $L$. Multiply by $L^2BM\alpha_{0,n}\ge0$, take the maximum over
$n\le k$, and use \eqref{4.14:eq-block-size-conditions-c} at each $n\le k$. This gives
\begin{equation*}
a^{\sharp}_F
\le L^2\mathsf{c}_1BM\max_{n\le k}R_n\alpha_{0,n}
\le a^{\sharp}_{\mathrm{ceil}}.
\end{equation*}
Thus, for each operation and for its levels $n\le k$, the per-operation constant
$a^{\sharp}_F$ leads to the same conclusion as $\mathfrak{a}(\square)$. In either form, the value
of the gauge constant required in the second gauge clause
(Proposition~\ref{4.14:prop-second-gauge-clause}) is supplied at hulls made by a large-field
operation, under a ceiling on the coupling of Ba{\l}aban's kind.

For a later level $n>k$, the second inequality of \cite[(2.9)]{B-CMP119} gives only
$R_{\ast}\le(L+1)(n-\ast)^{\beta_0}R_n$, where $\ast=k_0+1$. The factor $(n-\ast)^{\beta_0}$
grows with the number $n-\ast$ of levels between them, so no bound uniform in the history is
asserted there under either form. The layerwise constant
$\mathfrak{a}(\square)$ is the form in which the ceiling is uniform in the history.
\end{proof}

\begin{remark}[Scope of the second gauge clause's supply]\label{4.14:rem-supply-scope}
The five points below concern the application of the operator theorem at free current
(Theorem~\ref{4.14:thm-operator-theorem}) to the inductive large-field expansion. We call
this expansion simply \emph{the expansion}. The points are conditions that the application
has to meet. They concern condition~\eqref{4.14:eq-real-backgrounds-b} and its supply, and lie
outside the implication asserted by Theorem~\ref{4.14:thm-operator-theorem}.
\begin{enumerate}
\item[(i)] \emph{The inclusion.} The real background of the expansion at the nest under
consideration must lie in a space that carries, at every index, one of the gauge clauses
displayed by Ba{\l}aban. These clauses are \cite[(1.12)]{B-CMP109}, \cite[(2.38)]{B-CMP119}
and \cite[(1.65), (1.67), (1.69)]{B-CMP122a}, the families set out in
Definition~\ref{4.14:def-real-backgrounds} and Lemma~\ref{4.14:lem-gauge-family}. The gauges
are taken on the sets
$Q\cap\Omega_0$, for the cubes $Q$ of the clause. The radii of the space must be at most
$\tfrac14\varepsilon_1(a_1^{\oplus},a_0^{\oplus})$, where $a_1^{\oplus},a_0^{\oplus}$ are the
radii fixed in Definition~\ref{4.14:def-block-size-conditions}. This is
condition~\eqref{4.14:eq-real-backgrounds-b}.

Ba{\l}aban's clauses provide a gauge transformation $u$ on $Q\cap X$ for each domain
$X$ of the representation \cite[(1.63)]{B-CMP122a}. Reading them globally, on $Q\cap\Omega_0$,
is a convention of the expansion.

\item[(ii)] \emph{The constant $C$.} The expansion must state its gauge clause at $C=1$, or
with $C\in\{1,2\}$ and the value $1$ allowed. In \cite[p.~191]{B-CMP122a} the constant $C$
in the size $CML^m\eta$ of the cubes is left open:
\begin{quote}
The constant $C$ above is a positive integer. It is usually small, because of geometric
constraints on the cubes, for example we can assume that $C\le 2$.
\end{quote}

With $C=2$ alone, a cube of side $2M\ell_m$ need not fit beside a shell created by the
operation $\mathbf{R}$ whose width is one cube of side $M\ell_m$. The reason is that near
$\partial Z''_m$ the set $\Omega''_m$ is a union of cubes of side $M\ell_m$, but need not be a
union of cubes of side $2M\ell_m$. The choice $C=1$ lies within the latitude left by the
quoted sentence.

\item[(iii)] \emph{The space at later steps on drop-out nests.} A \emph{drop-out} is a
member of $\mathsf{Z}$ in which the preliminary integrations stop at some stage because a new
large field has been introduced. For the
later steps on such a component, Ba{\l}aban's construction rests on the following sentence,
printed without proof in \cite[p.~191]{B-CMP122a}; cited as printed:
\begin{quote}
Finally let us remark that if in a component of $Z$ we finish the integrations for some
$n$, because some new large field has been introduced, then we leave the new terms
$\mathbb{B}_k^{(n)}$ as a part of all boundary terms, and it is easy to generalize the
definition of the spaces in such a case for the next steps.
\end{quote}
(Here $n$ is the number of preliminary integrations.)
Lemma~\ref{4.14:lem-gauge-family} and Proposition~\ref{4.14:prop-second-gauge-clause}
supply two things under every family displayed by Ba{\l}aban: the membership of the hull's
cube $Q(\square)$ in the family, and the value~\eqref{4.14:eq-hull-gauge-value-a} of its
gauge constant. On drop-out nests the expansion therefore has to fix only one thing: which
of these families its space carries.

\item[(iv)] \emph{Hulls at the two lowest levels.} Let $\square$ be a box of the second cover
at level $0$ or $1$, that is, $\square\in\mathcal{D}^{\sharp}_0\cup\mathcal{D}^{\sharp}_1$.
For such $\square$ the second gauge clause in~\eqref{4.14:eq-real-backgrounds-a}
is supplied in the first of its two admitted forms. This form is
\cite[(3.35)]{B-CMP99b} at the second block size $M_1^{\sharp}$.

\item[(v)] \emph{The powers of $\mathsf{L}_0$.} Let $k_0:=k-N_0$. Write $\alpha_{\cdot,n}$ for
either of the radii $\alpha_{0,n},\alpha_{1,n}$ at a level $n$ of a stage nest. The two
conditions below are then read as follows.
\begin{itemize}
\item The ceiling in~\eqref{4.14:eq-block-size-conditions-b} on the constant of
\cite[(3.35)]{B-CMP99b} at block size $M_1$ is to be read, in every layer $n$, at
$M_1\mathsf{c}_1R_n\alpha_{0,n}$. The power of $\mathsf{L}_0$ appears only in the layers with
$k_0<n\le k$; in the layers $n\le k_0$ its exponent is zero, and the condition so read
implies the one at $M_1\alpha_{0,n}$, because $\mathsf{c}_1R_n\ge1$.
\item The radii in~\eqref{4.14:eq-real-backgrounds-b} are to be read, in every layer $n$, at
$\mathsf{c}_1R_n\alpha_{\cdot,n}$. Again the power of $\mathsf{L}_0$ appears only in the
layers with $k_0<n\le k$, and in the layers $n\le k_0$ the condition so read implies the one
at $\alpha_{\cdot,n}$.
\end{itemize}
These are one-sided conditions of the same kind as the ceiling
\eqref{4.14:eq-block-size-conditions-c}. Each bounds a layer's own quantity
$R_n\alpha_{\cdot,n}$, which is $g_n$ times a factor polylogarithmic in $g_n$. They are
imposed after $L$ and $M$, at the same position of the order of choice of
Definition~\ref{4.14:def-block-size-conditions}. None of them is a cap in the sense of
Definition~\ref{2.3:def-cap}.

Two conditions are not affected: the regularity clause of~\eqref{4.14:eq-real-backgrounds-a}
on cubes of side $M$ holding a bond of $\mathbb{B}^{+}$, and the ceiling
in~\eqref{4.14:eq-block-size-conditions-b} on its constants $a_Q$. Point~(v) moves no constant
of Theorem~\ref{4.14:thm-operator-theorem}; in particular, clause~(I) is unchanged.
\end{enumerate}
None of the points (i)--(v) belongs to the implication of
Theorem~\ref{4.14:thm-operator-theorem}. None of them bounds $L$, $K$, $k$ or $N$ from below.
\end{remark}

\begin{proof}
Points (i)--(iv) state which conditions the application has to verify, together with the
reasons given in them. We prove the assertions of~(v). Let $j$ denote the level of the stage
of the operation $\mathbf{R}$ under consideration; it satisfies $j\le k$. If $j\le k_0$, the
conditions \cite[(1.64), (1.65)]{B-CMP122a} carry no power of $\mathsf{L}_0$ and the
conditions \cite[(1.66), (1.67)]{B-CMP122a} are void \cite[p.~191]{B-CMP122a}, so there is
nothing to prove; we assume $j>k_0$.

\emph{The position of the factor.} Consider the definition of the stage space in
\cite[p.~190]{B-CMP122a}. The condition \cite[(1.64)]{B-CMP122a} applies on
$(\Omega''_m\setminus\Omega''_{m+1})\cap X$ for $1\le m\le j-1$, and on
$(\Omega''_j\setminus\Omega_{k_0+1})\cap X$ for $m=j$. There it bounds the seven quantities
\begin{equation*}
|\partial\mathbb{U}-1|,\quad |\partial U_{p,X}-1|,\quad |\mathbb{J}|,\quad |J_{p,X}|,\quad
|\partial U-1|,\quad |A'|,\quad |\nabla^{\eta}_U A'| .
\end{equation*}
Here $U_{p,X}$ and $J_{p,X}$, $p=1,\dots,k$, are the auxiliary configurations of
\cite[p.~190]{B-CMP122a} entering these bounds; the letter $p$ is used in this sense only in
the present proof.
The bound has three factors:
\begin{itemize}
\item a numerical factor independent of $\mathsf{L}_0$,
\begin{equation*}
1-\beta\sum_{i=h+1}^{j}2^{-|m-i|}-\beta\sum_{q=m+1}^{k}2^{-(q-m)} ,
\end{equation*}
where $\beta$ is Ba{\l}aban's constant in the definition of these spaces and $h$ is the lowest
level touched by the operation;
\item the power $\mathsf{L}_0^{2\max\{0,m-k_0\}}$;
\item a bracket of seven entries, one for each quantity. The first five entries carry
$\alpha_{0,m}$ and the last two carry $\alpha_{1,m}$.
\end{itemize}
The condition \cite[(1.66)]{B-CMP122a} applies on $(\Omega_m\setminus\Omega_{m+1})\cap X$ for
$k_0+1\le m\le j$. It bounds the same seven quantities by the analogous product, with the
power $\mathsf{L}_0^{2(j-m)}$ and entries carrying $\alpha_{0,j}$ and $\alpha_{1,j}$.

In both displays the power of $\mathsf{L}_0$ stands outside the bracket. It therefore
multiplies every one of the seven entries, not only the gauge clauses
\cite[(1.65), (1.67)]{B-CMP122a}. In particular it multiplies the entries of
$|\partial U-1|$, of $|\mathbb{J}|$ and $|J_{p,X}|$, and of $|A'|$ and
$|\nabla^{\eta}_UA'|$.

\emph{The exponent.} Let $n$ be the index of the radii in an entry, and let $\mathsf{L}_0^{2e}$
be the power multiplying it.
\begin{itemize}
\item In \cite[(1.64)]{B-CMP122a} we have $n=m$ and $e=\max\{0,m-k_0\}$. Hence $e=0$ unless
$k_0<n$, and in that case $e=n-k_0$.
\item In \cite[(1.66)]{B-CMP122a} we have $n=j$ and $e=j-m$, with $k_0+1\le m\le j$. Hence
\begin{equation*}
0\le e\le j-k_0-1<n-k_0 .
\end{equation*}
\item For levels $n$ with $j<n\le k$ the condition that applies is
\cite[(1.68)]{B-CMP122a}, on $(\Omega_n\setminus\Omega_{n+1})\cap X$ for $j<n<k$ and on
$\Omega_k\cap X$ for $n=k$. It carries no power of $\mathsf{L}_0$, so $e=0$.
\end{itemize}
Thus whenever $k_0<n\le k$,
\begin{equation}\label{4.14:eq-supply-scope-1}
0\le e\le n-k_0\le k-k_0=N_0 .
\end{equation}

\emph{The bound on the power.} Let $k_0<n\le k$ and let $e$ satisfy
\eqref{4.14:eq-supply-scope-1}. The auxiliary integer's power against the block count
(Lemma~\ref{4.14:lem-integer-power-bound}, display~\eqref{4.14:eq-integer-power-bound-a})
gives
\begin{equation*}
\mathsf{L}_0^{2e}\le c(\beta_0)L(L+1)R_n .
\end{equation*}
Since $\mathsf{c}_1=1+c(\beta_0)L(L+1)\ge c(\beta_0)L(L+1)$ and $R_n>0$, it follows that
\begin{equation}\label{4.14:eq-supply-scope-2}
\mathsf{L}_0^{2e}\le \mathsf{c}_1R_n .
\end{equation}
The same bound also follows from the top power against every lower level
(Lemma~\ref{4.14:lem-top-power-bound}). Indeed $\mathsf{L}_0\ge1$ and $e\le N_0$ by
\eqref{4.14:eq-supply-scope-1}, so
\begin{equation*}
\mathsf{L}_0^{2e}\le\mathsf{L}_0^{2N_0}\le\mathsf{c}_1R_n
\end{equation*}
by \eqref{4.14:eq-top-power-bound-a}, which holds for every $n\le k$.

\emph{Consequences for the first clause and the radii.} Fix a level $n$ with $k_0<n\le k$ of
the stage nest. Consider a cube of the class of \cite[(3.35)]{B-CMP99b}. The space supplies
the first clause of~\eqref{4.14:eq-real-backgrounds-a} there through its bound on
$|\partial U-1|$, and that entry carries the power $\mathsf{L}_0^{2e}$. So the only gauge that
Ba{\l}aban's space provides on such a cube has constant $\mathsf{L}_0^{2e}\alpha_{0,n}$, not
$\alpha_{0,n}$, and this is the constant with which the space supplies the first clause
of~\eqref{4.14:eq-real-backgrounds-a}. By
\eqref{4.14:eq-supply-scope-2} this constant is at most $\mathsf{c}_1R_n\alpha_{0,n}$.

The ceiling in \eqref{4.14:eq-block-size-conditions-b} on that clause's constant is
formulated with $\alpha_0$. It is therefore read, in each layer $n$ with $k_0<n\le k$, with
$\alpha_0$ replaced by $\mathsf{c}_1R_n\alpha_{0,n}$, that is, at
$M_1\mathsf{c}_1R_n\alpha_{0,n}$. At a layer $n\le k_0$ the exponent $e$ is zero, by the
cases above, so the constant supplied is $\alpha_{0,n}$; since $\mathsf{c}_1\ge1$ and
$R_n\ge1$ ($R_n\in L^{\mathbb{N}}$ in the regime of Lemma~\ref{4.14:lem-integer-power-bound}),
the ceiling read at $M_1\mathsf{c}_1R_n\alpha_{0,n}$ implies the ceiling at $M_1\alpha_{0,n}$.
The reading at $M_1\mathsf{c}_1R_n\alpha_{0,n}$ therefore serves every layer $n$.

The same argument applies to the radii. The entries bounding $|A'|$ and $|\nabla^{\eta}_UA'|$
(radius $\alpha_{1,n}$), and those bounding $|\mathbb{J}|$ and $|J_{p,X}|$, also carry
$\mathsf{L}_0^{2e}$. Hence the radii in~\eqref{4.14:eq-real-backgrounds-b} are read at
$\mathsf{c}_1R_n\alpha_{\cdot,n}$ in each layer $n$ with $k_0<n\le k$. At a layer $n\le k_0$
these entries carry no power of $\mathsf{L}_0$, and, as $\mathsf{c}_1R_n\ge1$, the condition
read at $\mathsf{c}_1R_n\alpha_{\cdot,n}$ implies the one at $\alpha_{\cdot,n}$; this reading
therefore serves every layer $n$.

It remains to compare these readings with \eqref{4.14:eq-block-size-conditions-c}. That
condition is a ceiling on each layer's $L^2\mathsf{c}_1BMR_m\alpha_{0,m}$, imposed after $L$
and $M$. The readings above are ceilings on the same layer quantities
$\mathsf{c}_1R_n\alpha_{\cdot,n}$, times constants fixed before them. So they are one-sided
conditions of the same kind, at the same position of the order of choice. They bound none of
the quantities listed in Definition~\ref{2.3:def-cap} from above; they are conditions on each
layer's own coupling, and are therefore not caps.

\emph{The regularity clause on cubes of side $M$.} The cubes $Q$ of this clause have side $M$,
hold a bond of $\mathbb{B}^{+}$ and lie within $M$ of $\Omega_t$, $t$ being the index in the
definition of $\mathbb{B}^{+}$. By the separation statement for these cubes, they lie at
distance at least $8LMR_t-2M$ from the large-field region $Z^{\mathrm{post}}_k$ at level $k$
after the operation. The set
$Z^{\mathrm{post}}_k$ contains every drop-out and every region kept as large-field. Hence the
only clause of Ba{\l}aban's spaces that applies on these cubes is the top clause, whose
exponent of $\mathsf{L}_0$ is zero. The clause and its ceiling
in~\eqref{4.14:eq-block-size-conditions-b} on the constants $a_Q$ are therefore unchanged.

On filled islands, in the sense of Proposition~\ref{4.14:prop-second-gauge-clause}, the
required clause is the sentence that the configuration there ``satisfies the induction
hypothesis'', printed without proof in \cite{B-CMP122b}; cited as printed.

\emph{The constants of the operator theorem.} The modification concerns $\alpha_0$, a
constant of the first clause of~\eqref{4.14:eq-real-backgrounds-a}. The supply of that
clause is part of condition~\eqref{4.14:eq-real-backgrounds-b}, by
Definition~\ref{4.14:def-real-backgrounds}. No constant of
Theorem~\ref{4.14:thm-operator-theorem} enters, in particular none of its clause~(I), so none
is changed.
\end{proof}

\begin{definition}[The cube sizes and the waiting depth]\label{4.15:not-cube-sizes}
The following letters and standing hypotheses are in force in this section.

\emph{Letters.} $N$ is the rank of the gauge group. $L$ is an odd integer with $L\ge L_0$,
where $L_0=L_0(N)$ is the constant of hypothesis (H2) of
Definition~\ref{2.2:not-hypotheses}; that hypothesis contains the condition $L\ge 5$.
$M$ is the side of the partition cubes. $M_1$ and $M_2:=L^2M_1<M$ are the smaller cube
parameters; $M$, $M_1$ and $M_2$ are all powers of $L$. $K\ge 0$ is the last level, levels
are integers $j,k,s,t\in\{0,\dots,K\}$, and $d=4$ is the dimension.

\emph{Couplings.} The coupling $g\in\,]0,\gamma_H]$ and the running couplings
$g_0,\dots,g_K$ are those of Theorem~0. The thresholds $\gamma$ and then $\gamma_H$ are the
last ones fixed in the order of choice, in hypotheses (H3) and (H4). We write
$x_j:=g_j^{-2}$.

\emph{Standing hypotheses on the parameters.} In this section, no upper bound is imposed on
$L$, $M$, $K$ or $k$, on a number of
steps, of iterations or of components, or on an age. No lower bound is imposed on any
coupling. The only hypothesis on the couplings that is used is the following part of
clause~(0) of Theorem~0 (Theorem~\ref{4.1:thm-clause-zero}):
\begin{equation}\label{4.15:eq-cube-sizes-a}
x_j>g^{-2}\ge\gamma_H^{-2}>1\qquad(0\le j\le K).
\end{equation}

\emph{Cube sizes.} For $0\le j\le K$ put
\begin{equation*}
\check{x}_j:=\min_{i\le j}x_i .
\end{equation*}
Let $\check{R}_j$ be the least power $L^b$, with $b\in\mathbb{Z}_{\ge 0}$, such that
\begin{equation}\label{4.15:eq-cube-sizes-1}
L^b\ge\bigl(\log\check{x}_j\bigr)^{r},
\end{equation}
where $r>0$ is the exponent fixed in hypothesis (H5) of
Definition~\ref{2.2:not-hypotheses}.

The lemma on the cube sizes shows that \eqref{4.15:eq-cube-sizes-a}, together
with $L\ge 2$ and $r>0$, has three consequences: $\check{R}_j$ is a well-defined power of
$L$; it is at least $1$; and the map $j\mapsto\check{R}_j$ is non-increasing.

Wherever Ba{\l}aban's size parameter $R_j$ occurs, we read $R_j:=\check{R}_j$. This
convention concerns that size parameter only; it does not apply to sets that Ba{\l}aban
denotes by the letter $R$.

\emph{Waiting depth.} The \emph{waiting depth} at level $k$ is
\begin{equation*}
\mathsf{N}=\mathsf{N}_k:=\check{R}_k .
\end{equation*}
\end{definition}

\begin{definition}[Two formulations of the no-creation condition]\label{4.15:def-no-creation}
Let $1\le k\le K$ and let $\sigma'$ be a term of $\mathbf{T}\rho_{k-1}$. Let $\Lambda_k$ be the
small-field region of scale $k$ of $\sigma'$ and
\begin{equation*}
Z_k=\Lambda_k^{c}=\mathbb{T}\setminus\Lambda_k
\end{equation*}
its large-field region of scale $k$. Here $\rho_{k-1}$ is the effective density of
Definition~\ref{1.5:def-densities}; its terms, their small-field and large-field regions and
their histories are made precise later in this chapter. Let $C$ be a component of $Z_k$, in the
sense of components
fixed later in this chapter. For $1\le t\le k$ let $\mathcal{W}^{\wedge}_t$ be the creation set
of the step $t-1\to t$ of the history of $\sigma'$; this set is also defined later in this chapter.

A \emph{creation at level $t$ lies inside $C$} if $C$ meets $\mathcal{W}^{\wedge}_t$. This is
meant in the widest sense: the point sets of $C$ and of $\mathcal{W}^{\wedge}_t$ have a common
point.

Let $\mathsf{N}_k=\check{R}_k$ be the waiting depth at level $k$
(Definition~\ref{4.15:not-cube-sizes}). Condition~(ii) in the definition of the level at which
$\mathbf{R}$ acts on a component of a large-field region is Ba{\l}aban's condition~(ii) of
\cite[p.~176]{B-CMP122a}. Its no-creation clause is formulated for $C$ at
level $k$ in the following two ways.
\begin{itemize}
\item[(ii$'$)] The inequality $k\ge\mathsf{N}_k$ holds, and no creation at any level $t$ with
$k-\mathsf{N}_k<t\le k$ lies inside $C$; that is,
\begin{equation}\label{4.15:eq-no-creation-a}
k\ge\mathsf{N}_k
\quad\text{and}\quad
C\cap\mathcal{W}^{\wedge}_t=\emptyset
\ \text{ for every integer } t \text{ with } k-\mathsf{N}_k<t\le k .
\end{equation}
\item[(ii$''$)] No creation at any level $t$ with $\max\{1,k-\mathsf{N}_k+1\}\le t\le k$ lies
inside $C$; that is,
\begin{equation}\label{4.15:eq-no-creation-b}
C\cap\mathcal{W}^{\wedge}_t=\emptyset
\ \text{ for every integer } t \text{ with } \max\{1,k-\mathsf{N}_k+1\}\le t\le k .
\end{equation}
Under (ii$''$) it is not required that $\mathsf{N}_k$ steps precede level $k$.
\end{itemize}
Reading (ii$'$) is the principal reading and (ii$''$) is the alternative one; the arguments of
this chapter presuppose only the no-creation clause, in whichever of the two readings is adopted.
\end{definition}

\begin{proof}[Justification of the window and of the two readings]
Ba{\l}aban's condition~(ii) reads ``in the preceding $N$ renormalization steps no new large
field regions were created inside this component'' \cite[p.~176]{B-CMP122a}. We take
$N=\mathsf{N}_k$.

The $\mathsf{N}_k$ steps preceding level $k$ are the steps $t-1\to t$ with
$k-\mathsf{N}_k<t\le k$. The step $k-1\to k$ is included among them.

The clause therefore says the following: if the last creation inside $C$ occurred at a level
$t\ge1$, then $k-t\ge\mathsf{N}_k$. Indeed, $k-t\ge\mathsf{N}_k$ fails exactly when
$k-\mathsf{N}_k<t$.

In particular, a creation at level $t=k$ inside $C$ is barred. To see this, note that
\begin{equation*}
k-k=0<1\le\mathsf{N}_k ,
\end{equation*}
where $\mathsf{N}_k=\check{R}_k\ge1$ by Definition~\ref{4.15:not-cube-sizes}.

This window agrees with \cite{B-CMP122a}. There, with $h=k-\mathsf{N}_k$, the product of
the last $\mathsf{N}_k$ one-step operations consists of the operations $\mathbf{T}_j$ indexed by
$j=h,\dots,k-1$. These operations produce the levels $h+1,\dots,k$. These are the levels whose
large-field regions the operation $\mathbf{R}$ of \cite{B-CMP122a} renews: by
\cite[(1.11)]{B-CMP122a}, new regions $Z''_j$ appear for $j\ge h+1$, while
\begin{equation*}
Z''_j=Z_j \quad\text{for } j=h,h-1,\dots,1 .
\end{equation*}
The levels $h+1,\dots,k$ are exactly the integers $t$ with
$k-\mathsf{N}_k<t\le k$.

For integers, $k-\mathsf{N}_k<t$ is the same as $t\ge k-\mathsf{N}_k+1$. When
$k\ge\mathsf{N}_k$ one has $k-\mathsf{N}_k+1\ge1$. Hence, whenever $k\ge\mathsf{N}_k$, the
windows in \eqref{4.15:eq-no-creation-a} and \eqref{4.15:eq-no-creation-b} coincide. The two
readings differ only at levels $k<\mathsf{N}_k$. At such levels (ii$'$) fails, while (ii$''$)
requires $C\cap\mathcal{W}^{\wedge}_t=\emptyset$ for every $t=1,\dots,k$.

We now explain why (ii$''$) is the alternative reading. Besides (ii$'$), the definition of the
level at which $\mathbf{R}$ acts admits a second reading of condition~(ii): one under which the
inequality $k\ge\mathsf{N}_k$ is not part of condition~(ii), and under which that inequality,
at a level at which $\mathbf{R}$ acts, is obtained instead from a structural statement about
creations. Under (ii$''$) the inequality $k\ge\mathsf{N}_k$ is indeed not part of the condition.

A statement about creations can yield $k\ge\mathsf{N}_k$ only if the clause on creations keeps
bearing on those of the preceding steps that exist. At a level $k<\mathsf{N}_k$ it must
therefore bear on all the steps $1,\dots,k$. This is what \eqref{4.15:eq-no-creation-b} does,
since there $\max\{1,k-\mathsf{N}_k+1\}=1$.

A third formalization would have condition~(ii) hold outright whenever fewer than
$\mathsf{N}_k$ steps precede level $k$. It is not admitted as a reading of condition~(ii):
under it, no statement about creations bounds the level of an action from below, and
nothing below is claimed under it.

Whenever an argument of this chapter needs a creation inside $C$, it produces a cube of the
coarse partition of the step $t-1\to t$ (fixed later in this chapter) that is contained both
in $\mathcal{W}^{\wedge}_t$ and in $C$. These arguments therefore
hold verbatim for every narrower sense of having a common point. Examples are a common cube of
the finest partition used, or a common interior point.

Ba{\l}aban's condition~(ii) continues ``and the previous regions contained in it satisfy the
condition (i) on the corresponding scales'' \cite[p.~176]{B-CMP122a}. The definition of the level
at which $\mathbf{R}$ acts adds a further condition~(iii). These clauses, like condition~(i),
only restrict further the levels at which $\mathbf{R}$ acts.

The arguments that follow use only one property of an action: that it presupposes the
no-creation clause of condition~(ii), in the reading adopted. They hold whether or not the
further clauses are imposed.
\end{proof}

\begin{definition}[Compatible partitions of the torus into cubes]\label{4.15:def-torus-partitions}
All regions of all levels are subsets of one torus $\mathbb{T}$. This includes, for every level $j$, the large-field region $Z_j=\mathbb{T}\setminus\Lambda_j$ of scale $j$. It also includes the creation set $\mathcal{W}^{\wedge}_t$ of the step $t-1\to t$, for every level $t\le k$, with which a component of $\Lambda_k^c$ is compared. We write
\begin{equation*}
\mathbb{T}=\mathbb{R}^d/p\,\mathbb{Z}^d ,
\end{equation*}
where $p>0$ is the period in a fixed unit of length. Here and in the rest of this section the letter $p$ denotes this period and not a plaquette. We write $\pi\colon\mathbb{R}^d\to\mathbb{T}$ for the covering map.

\emph{Partitions.} Let $a>0$ be such that $p/a$ is a positive integer. The \emph{partition of side $a$} is the set $\mathcal{P}$ of the closed cubes
\begin{equation}\label{4.15:eq-torus-partitions-1}
\pi\Bigl(\prod_{\nu=1}^{d}\bigl[a n_\nu,\,a(n_\nu+1)\bigr]\Bigr),\qquad n=(n_1,\dots,n_d)\in\mathbb{Z}^d .
\end{equation}
Thus, when $a$ is a power of $L$, the partition into cubes of side $a$ is the one whose cubes have their corners in $\pi(a\mathbb{Z}^d)$. The partition $\mathcal{P}$ has $(p/a)^d$ distinct members, called the \emph{$\mathcal{P}$-cubes}, and their union is $\mathbb{T}$.

\emph{Refinement.} A partition $\mathcal{P}'$ \emph{refines} a partition $\mathcal{P}$ if two conditions hold:
\begin{itemize}
\item every $\mathcal{P}'$-cube is contained, as a point set, in exactly one $\mathcal{P}$-cube;
\item every $\mathcal{P}$-cube is the union of the $\mathcal{P}'$-cubes it contains.
\end{itemize}
Let $\mathcal{P}'$ and $\mathcal{P}$ be partitions of sides $a'$ and $a$ of the form \eqref{4.15:eq-torus-partitions-1}, that is, with common origin. Then $\mathcal{P}'$ refines $\mathcal{P}$ as soon as $a/a'$ is a positive integer.

\emph{Compatibility.} The construction uses finitely many partitions at each of the finitely many levels $\le K$. These are the partitions into coarse cubes, into test cubes, into $M$-cubes, into $M\check{R}_j$-cubes and into unit cubes. They are \emph{pairwise compatible}: of any two of them, one refines the other. The reason is that their sides are, in the common unit, powers of $L$ times one length, and their corners lie on lattices through a common origin. Compare the compatibility of each new partition with all the previous partitions in \cite[p.~265]{B-CMP119}.

Refinement is a partial order, and these finitely many partitions are pairwise comparable in it. We write $\mathcal{P}_{\bullet}$ for the finest of them. The number of $\mathcal{P}_{\bullet}$-cubes of $\mathbb{T}$ is finite.

In what follows no relation between $p$ and any cube side is assumed beyond these divisibilities. In particular, no statement that cubes of a given side fit in the torus is used.
\end{definition}

\begin{definition}[Regions, touching, chains and components]\label{4.15:def-regions-components}
Let $\mathbb{T}$ be the torus and $\mathcal{P}_{\bullet}$ the finest of the pairwise compatible
partitions of Definition~\ref{4.15:def-torus-partitions}.

\emph{Regions.} A \emph{region} is a set $X$ of $\mathcal{P}_{\bullet}$-cubes. The same letter
$X$ also denotes the union of its members as a point set of $\mathbb{T}$, and we speak of the
points of $X$. Inclusion and union of regions are inclusion and union of sets of
$\mathcal{P}_{\bullet}$-cubes. They imply the same relations between the point sets: the point
set of a union of regions is the union of their point sets, and if every member of $X$ is a
member of $X'$, then the union of the members of $X$ is contained in the union of the members
of $X'$.

\emph{Cubes of other partitions.} Let $\mathcal{P}$ be any of the partitions of
Definition~\ref{4.15:def-torus-partitions}. A $\mathcal{P}$-cube $c$ is regarded as the region
consisting of the $\mathcal{P}_{\bullet}$-cubes contained in $c$. Since
$\mathcal{P}_{\bullet}$ refines $\mathcal{P}$, every $\mathcal{P}$-cube is the union of the
$\mathcal{P}_{\bullet}$-cubes it contains, so the point set of this region is $c$ itself. A
\emph{$\mathcal{P}$-region} is a union of $\mathcal{P}$-cubes so regarded; its
\emph{$\mathcal{P}$-cubes} are the $\mathcal{P}$-cubes all of whose $\mathcal{P}_{\bullet}$-cubes
belong to it. All regions of the construction are regions in this sense, being unions of cubes
of partitions used in it. This applies to the large-field region $Z_j$ of scale $j$, to the
regions $P'$, $Q'$ and $R'$ introduced below, to the complement $\Omega_t^{c}$ of the domain
$\Omega_t$ at level $t$, to the creation set $\mathcal{W}^{\wedge}_t$ of the step from $t-1$
to $t$, and to the stage regions of the cube-test recursion introduced below.

\emph{Touching.} Two cubes, of the same or of different partitions, \emph{touch} if their point
sets have a common point. This includes meeting at a corner only, and it includes containment of
one cube in the other. A cube \emph{meets} a region if it touches a member of the region. Two
regions \emph{touch} if some member of the one touches some member of the other. Regions that do
not touch have disjoint point sets. Indeed, a common point of $X$ and $X'$ lies in a member of
$X$ and in a member of $X'$, and these two members then touch.

\emph{Chains and connectedness.} A \emph{chain in $X$} is a finite sequence of members of $X$ in
which consecutive members touch. The region $X$ is \emph{connected} if it is non-empty and any
two of its members are the ends of a chain in $X$.

\emph{Components.} On the members of a region $X$, the relation ``joined by a chain in $X$'' is
an equivalence relation. It is reflexive, since a single member is a chain. It is symmetric, since
a chain read backwards is a chain. It is transitive, since two chains, the second beginning
where the first ends, concatenate to a chain. The \emph{components} of $X$ are the classes of
this relation, each regarded as a region. They are thus the classes generated by the relation
``intersect or touch'', corners included. The empty region has no component.
\end{definition}

\begin{lemma}[Components do not depend on the subdivision]\label{4.15:lem-subdivision-independence}
Let $\mathcal{P}$ be one of the partitions of Definition~\ref{4.15:def-torus-partitions}, let $X$ be
a $\mathcal{P}$-region in the sense of Definition~\ref{4.15:def-regions-components}, and let $S$ be
its set of $\mathcal{P}$-cubes. Let $a\subset c$ and $a'\subset c'$ be members of $X$, where
$c,c'\in S$ are the $\mathcal{P}$-cubes containing them. Then $a$ and $a'$ are joined by a chain
in $X$ if and only if $c$ and $c'$ are joined by a chain in $S$, that is, by a finite sequence of
members of $S$ with consecutive members touching. Hence:
\begin{itemize}
\item the components of $X$ are the unions of the classes of the relation ``joined by a chain
in $S$'' on $S$;
\item every component of a $\mathcal{P}$-region is a $\mathcal{P}$-region;
\item a single cube of any partition, regarded as a region, is connected.
\end{itemize}
\end{lemma}

\begin{proof}
Throughout, $\mathcal{P}_{\bullet}$ is the finest partition, which refines $\mathcal{P}$. Write
$\ell_{\mathcal{P}}$ and $\ell_{\bullet}$ for the sides of $\mathcal{P}$ and of
$\mathcal{P}_{\bullet}$, and $q_{\mathcal{P}}:=\ell_{\mathcal{P}}/\ell_{\bullet}$, a positive integer.
We first record two facts.

(\(\alpha\)) The $\mathcal{P}_{\bullet}$-cubes contained in one $\mathcal{P}$-cube $c$ form a
connected region. Indeed, write
\begin{equation*}
c=\pi\Bigl(\prod_{\nu=1}^{d}\bigl[\ell_{\mathcal{P}}n_\nu,\ \ell_{\mathcal{P}}(n_\nu+1)\bigr]\Bigr),
\qquad \mathbf{n}=(n_1,\dots,n_d)\in\mathbb{Z}^d .
\end{equation*}
The $\mathcal{P}_{\bullet}$-cubes contained in $c$ are the cubes
\begin{equation*}
\pi\Bigl(\prod_{\nu=1}^{d}\bigl[\ell_{\bullet}z_\nu,\ \ell_{\bullet}(z_\nu+1)\bigr]\Bigr)
\quad\text{with}\quad
q_{\mathcal{P}}n_\nu\le z_\nu\le q_{\mathcal{P}}n_\nu+q_{\mathcal{P}}-1
\quad\text{for every }\nu,
\end{equation*}
where $\mathbf{z}=(z_1,\dots,z_d)\in\mathbb{Z}^d$. The set of such index vectors $\mathbf{z}$ is a
box in $\mathbb{Z}^d$, and it is non-empty. Two of these cubes whose index vectors differ by a unit
coordinate vector share a face, hence touch. Any two index vectors in the box are joined by a
sequence of such unit moves that stays inside the box: change the first coordinate one unit at a
time until it agrees, then the second, and so on; every intermediate vector has each coordinate
between the corresponding coordinates of the two ends, hence lies in the box. The corresponding
sequence of cubes is a chain, which proves (\(\alpha\)).

(\(\beta\)) Two $\mathcal{P}$-cubes touch if and only if some $\mathcal{P}_{\bullet}$-cube of the
one touches some $\mathcal{P}_{\bullet}$-cube of the other. This holds because each
$\mathcal{P}$-cube is the union, as a point set, of its $\mathcal{P}_{\bullet}$-cubes: a common
point of the two $\mathcal{P}$-cubes lies in a $\mathcal{P}_{\bullet}$-cube of each, and these two
then touch; conversely, a common point of a $\mathcal{P}_{\bullet}$-cube of the one and a
$\mathcal{P}_{\bullet}$-cube of the other is a common point of the two $\mathcal{P}$-cubes.

Now let $c=c_0,c_1,\dots,c_l=c'$ be a chain in $S$. For $1\le i\le l$ the cubes $c_{i-1}$ and
$c_i$ touch, so by (\(\beta\)) we may choose touching members $b_i\subset c_{i-1}$ and
$b'_i\subset c_i$. By (\(\alpha\)), $a$ is joined to $b_1$ by a chain inside $c_0$; $b'_i$ is
joined to $b_{i+1}$ by a chain inside $c_i$ for $1\le i<l$; and $b'_l$ is joined to $a'$ by a
chain inside $c_l$. Each $c_i$ belongs to $S$, so all its $\mathcal{P}_{\bullet}$-cubes belong to
$X$, and each of these chains is a chain in $X$. Since $b_i$ and $b'_i$ touch, the concatenation
of these chains, in the order listed, is a chain in $X$ from $a$ to $a'$. (If $l=0$, then
$c=c'$ and (\(\alpha\)) alone joins $a$ to $a'$ inside $c$.)

Conversely, let $a=a_0,a_1,\dots,a_l=a'$ be a chain in $X$, and consider the sequence of the
$\mathcal{P}$-cubes containing the $a_i$. These $\mathcal{P}$-cubes belong to $S$. Indeed, a member
of the $\mathcal{P}$-region $X$ lies in exactly one $\mathcal{P}$-cube: since $X$ is a union of
$\mathcal{P}$-cubes, the member lies in some $\mathcal{P}$-cube all of whose
$\mathcal{P}_{\bullet}$-cubes belong to $X$, and since $\mathcal{P}_{\bullet}$ refines
$\mathcal{P}$ that $\mathcal{P}$-cube is the only one containing it. The $\mathcal{P}$-cubes of
$X$ are those all of whose $\mathcal{P}_{\bullet}$-cubes belong to $X$, so the $\mathcal{P}$-cube
containing each $a_i$ is a member of $S$. Consecutive members $a_{i-1},a_i$ touch, so by
(\(\beta\)) the $\mathcal{P}$-cubes containing them are equal or touch. Deleting repetitions gives a
chain in $S$ from $c$ to $c'$.

The remaining assertions follow. By the equivalence just proved, a member $a'$ of $X$ lies in the
class of $a$ in $X$ if and only if the $\mathcal{P}$-cube $c'$ containing it lies in the class of
$c$ in $S$; hence the class of $a$ in $X$ is the union of the $\mathcal{P}$-cubes in the class of
$c$ in $S$. Thus the components of $X$ are the unions of the classes of ``joined by a chain in
$S$'', and each of them, being a union of $\mathcal{P}$-cubes, is a $\mathcal{P}$-region. Finally,
a single cube $c$ of a partition $\mathcal{P}$, regarded as a region, is a $\mathcal{P}$-region
whose set of $\mathcal{P}$-cubes is the singleton $\{c\}$; this set consists of one class, so the
region is non-empty and has exactly one component, that is, it is connected.
\end{proof}

\begin{definition}[Enlargement, hull and the stabilised boxing map]\label{4.15:def-enlargement-hull}
We use the torus $\mathbb{T}=\mathbb{R}^d/p\mathbb{Z}^d$, its covering map
$\pi:\mathbb{R}^d\to\mathbb{T}$, the compatible partitions of
Definition~\ref{4.15:def-torus-partitions}, the finest of them $\mathcal{P}_{\bullet}$, and the
regions, touching, chains and components of Definition~\ref{4.15:def-regions-components}.
Fix a renormalisation step from level $t-1$ to level $t$, and let $\mathcal{P}_t$ be its coarse
partition, the partition of side $W_t=LM\check{R}_t$; its cubes are the coarse cubes of the
step. For a region $X$ we write $X^c$ for its complement, the region of the
$\mathcal{P}_{\bullet}$-cubes not in $X$.

\emph{Enlargement and removal of layers.} For a $\mathcal{P}_t$-region $X$, let $X^{\sim}$ be the
$\mathcal{P}_t$-region whose $\mathcal{P}_t$-cubes are the $\mathcal{P}_t$-cubes touching some
$\mathcal{P}_t$-cube of $X$; thus $\sim$ adds one layer of touching coarse cubes, as in the
layers of \cite[pp.~248, 265]{B-CMP119}. Put $X^{\sim 0}:=X$ and
$X^{\sim(n+1)}:=(X^{\sim n})^{\sim}$ for $n\ge 0$. Let $X^{\sim -1}$ be the
$\mathcal{P}_t$-region whose $\mathcal{P}_t$-cubes are those $\mathcal{P}_t$-cubes of $X$ that do not
touch $X^c$; it is $X$ with one layer of its cubes removed, as in \cite{B-CMP119}.

\emph{Hull.} For an arbitrary region $Y$, the \emph{$\mathcal{P}_t$-hull} $\widehat{Y}$ is the
$\mathcal{P}_t$-region whose $\mathcal{P}_t$-cubes are the $\mathcal{P}_t$-cubes meeting $Y$; when $Y$
is the large-field region $Z_k$, this is the union of all $LMR_{k+1}$-cubes intersecting $Z_k$
of \cite{B-CMP119}, written $\bar{Z}_k$ there. Ba{\l}aban's word \emph{intersecting} may be read
as sharing a point or as sharing a $\mathcal{P}_{\bullet}$-cube; this is immaterial below, since
under either reading $\widehat{Y}\supseteq Y$ and every $\mathcal{P}_t$-cube of $\widehat{Y}$
touches a member of $Y$, and nothing else about $\widehat{Y}$ is used.

\emph{Boxing.} Put $N_t:=p/W_t$, a positive integer, and for $n=(n_1,\dots,n_d)\in\mathbb{Z}^d$
write
\begin{equation*}
c_n:=\pi\Bigl(\prod_{\nu=1}^{d}\,[\,W_t n_\nu,\;W_t(n_\nu+1)\,]\Bigr).
\end{equation*}
Then $c_n=c_{n'}$ if and only if $n\equiv n'\pmod{N_t}$ coordinatewise, so the
$\mathcal{P}_t$-cubes are indexed by $(\mathbb{Z}/N_t\mathbb{Z})^d$. For a non-empty connected
$\mathcal{P}_t$-region $C$ and each $\nu\in\{1,\dots,d\}$, let $S_\nu(C)\subset\mathbb{Z}/N_t\mathbb{Z}$,
the \emph{$\nu$-th coordinate range} of $C$, be the set of the residues $n_\nu \bmod N_t$ of the
indices $n$ for which $c_n$ is a $\mathcal{P}_t$-cube of $C$. Let $\operatorname{box}(C)$ be the
$\mathcal{P}_t$-region consisting of the cubes $c_n$ with $n_\nu \bmod N_t\in S_\nu(C)$ for every
$\nu$, the product of the coordinate ranges. Then $C\subset\operatorname{box}(C)$, since every
$\mathcal{P}_t$-cube $c_n$ of $C$ has $n_\nu \bmod N_t\in S_\nu(C)$ for every $\nu$. It is shown
in the lemma on boxing later in this section that each $S_\nu(C)$ is a cyclic interval and that
$\operatorname{box}(C)$ is connected. When $C$ is contained in the image under $\pi$ of a cube of
$\mathbb{R}^d$ on which $\pi$ is injective, $\operatorname{box}(C)$ is the image of the smallest
axis-parallel parallelepiped containing the lift of $C$; this formalises the smallest rectangular
parallelepiped containing a component in \cite[(3.20)]{B-CMP119}.

\emph{Small components and one application of the boxing map.} A component $C$ of a
$\mathcal{P}_t$-region is
\emph{small} if its point set is contained in $\pi\bigl(y+[0,100W_t]^d\bigr)$ for some
$y\in\mathbb{R}^d$; here $100W_t=100LM\check{R}_t$, and this is Ba{\l}aban's condition of being
contained in a cube of side $100LMR_{k+1}$, that is, one hundred coarse-cube sides,
\cite[(3.20)]{B-CMP119}. By
Lemma~\ref{4.15:lem-subdivision-independence}, every component of a $\mathcal{P}_t$-region is a
connected $\mathcal{P}_t$-region, so $\operatorname{box}(C)$ is defined for it. For a
$\mathcal{P}_t$-region $X$ put
\begin{equation}\label{4.15:eq-enlargement-hull-1}
\mathsf{M}_1X:=X\cup\bigcup\bigl\{\operatorname{box}(C):\ C\ \text{a small component of}\ X\bigr\},
\end{equation}
a $\mathcal{P}_t$-region, being a union of $\mathcal{P}_t$-regions.

\emph{The stabilised boxing map.} For a $\mathcal{P}_t$-region $X$ let $X_q:=\mathsf{M}_1^{\,q}X$
($q\ge 0$ applications of $\mathsf{M}_1$). There is a least integer $q_0=q_0(X)\ge 0$ with
$X_{q_0+1}=X_{q_0}$, and we put
\begin{equation}\label{4.15:eq-enlargement-hull-2}
\mathsf{M}X:=X_{q_0(X)}.
\end{equation}
Then $X_q=\mathsf{M}X$ for all $q\ge q_0$, $\mathsf{M}X\supseteq X$, $\mathsf{M}_1(\mathsf{M}X)=\mathsf{M}X$
and $\mathsf{M}(\mathsf{M}X)=\mathsf{M}X$. No bound on $q_0$ is used.

\emph{Replacement maps.} A map $\Phi$ from $\mathcal{P}_t$-regions to $\mathcal{P}_t$-regions is a
\emph{replacement map} if for every $\mathcal{P}_t$-region $X$ there are components
$C_1,\dots,C_m$ of $X$ ($m\ge 0$) and connected $\mathcal{P}_t$-regions
$B_1\supseteq C_1,\dots,B_m\supseteq C_m$ with
\begin{equation*}
\Phi X=X\cup B_1\cup\dots\cup B_m .
\end{equation*}
The name is explained by the identity $(X\setminus C_i)\cup B_i=X\cup B_i$, valid because
$C_i\subseteq B_i$: replacing $C_i$ by $B_i\supseteq C_i$ in $X$ gives $X\cup B_i$. Of
$\mathsf{M}$, only its structure as a composite of replacement maps is used (in clauses (d) and
(e) of the lemma on boxing later in this section); hence every formalisation of Ba{\l}aban's
modification in which a component is replaced by a connected union of coarse cubes containing it
yields the same conclusions.
\end{definition}

\begin{proof}
We prove the assertions made about $q_0$ and $\mathsf{M}$.

First, $\mathsf{M}_1Y\supseteq Y$ for every $\mathcal{P}_t$-region $Y$, since $Y$ is one of the sets
whose union is taken in \eqref{4.15:eq-enlargement-hull-1}; and $\mathsf{M}_1Y$ is again a
$\mathcal{P}_t$-region, so $\mathsf{M}_1$ can be iterated. Applying this to $Y=X_q$ gives
$X_q\subseteq X_{q+1}$ for every $q\ge 0$, so $(X_q)_{q\ge 0}$ is a non-decreasing sequence of
$\mathcal{P}_t$-regions.

Second, the $\mathcal{P}_t$-cubes of $\mathbb{T}$ are finitely many, namely $N_t^d$. A
$\mathcal{P}_t$-region is determined by its set of $\mathcal{P}_t$-cubes, and along a
non-decreasing sequence these sets are non-decreasing subsets of a set with $N_t^d$ elements, so
the inclusion $X_q\subseteq X_{q+1}$ can be strict for at most $N_t^d$ values of $q$. Hence the
set of $q\ge 0$ with $X_{q+1}=X_q$ is non-empty, and its least element $q_0=q_0(X)$ exists.

Third, $X_q=X_{q_0}$ for all $q\ge q_0$, by induction on $q$. For $q=q_0$ there is nothing to
prove. If $X_q=X_{q_0}$ for some $q\ge q_0$, then
\begin{equation*}
X_{q+1}=\mathsf{M}_1X_q=\mathsf{M}_1X_{q_0}=X_{q_0+1}=X_{q_0},
\end{equation*}
the last equality being the defining property of $q_0$. Thus $X_q=\mathsf{M}X$ for all $q\ge q_0$,
by \eqref{4.15:eq-enlargement-hull-2}.

Fourth, $\mathsf{M}X=X_{q_0}\supseteq X_0=X$, since the sequence is non-decreasing. Moreover
\begin{equation*}
\mathsf{M}_1(\mathsf{M}X)=\mathsf{M}_1X_{q_0}=X_{q_0+1}=X_{q_0}=\mathsf{M}X .
\end{equation*}

Finally, put $Y:=\mathsf{M}X$. The sequence $Y_q=\mathsf{M}_1^{\,q}Y$ has $Y_0=Y$ and, by the
previous display, $Y_1=\mathsf{M}_1Y=Y=Y_0$. Hence $q_0(Y)=0$, and
$\mathsf{M}Y=Y_{q_0(Y)}=Y_0=Y$ by \eqref{4.15:eq-enlargement-hull-2}; that is,
$\mathsf{M}(\mathsf{M}X)=\mathsf{M}X$.
\end{proof}

\begin{definition}[The small-field complement as a union of enlargements
{\cite[(3.2)--(3.5), p.~265; (1.2), (1.4), (1.8), (1.10), pp.~246--248]{B-CMP119}}]
\label{4.15:def-small-field-complement}
Let $1\le t\le K$ and put $k:=t-1$. The step $k\to t$ of the operation $\mathbf{T}$ acts
on a term $\sigma$ of $\rho_k$ inside $\mathbf{T}\rho_k$. At the first step the same construction
applies, with the first-step decompositions of unity of \cite[\S1]{B-CMP119} in place of those of
\cite[\S3]{B-CMP119}.

\emph{Coarse cubes and test cubes.} In the units of the bond lattice of the fluctuation field of
step $k$, put $W_t:=LM\check{R}_t$. The \emph{coarse cubes} of the step are the cubes of side $W_t$
of the partition $\mathcal{P}_t$ of the whole lattice (Definition~\ref{4.15:def-torus-partitions}).
They are Ba{\l}aban's $LMR_{k+1}$-cubes of \cite[\S3]{B-CMP119}, with $R_{k+1}$ read as
$\check{R}_t$. At $t=1$ they are his $LMR_1$-cubes of \cite[\S1]{B-CMP119}. The \emph{test cubes}
are the cubes of side $L^2M_2\check{R}_t$. Regions, touching and components are meant with respect
to $\mathcal{P}_t$, as in Definition~\ref{4.15:def-regions-components}. The enlargement $X^{\sim a}$
by $a$ layers of coarse cubes and the $\mathcal{P}_t$-hull $\widehat{X}$ are as in
Definition~\ref{4.15:def-enlargement-hull}.

\emph{The input region of the step.} The step takes as input a region $Y\subset\mathbb{T}$, namely the
large-field region $Z_k=\Lambda_k^c$ of scale $k$ of the term $\sigma$ entering it. Which region
$Y$ is, for the terms of a history, is fixed later in this chapter. At $t=1$ there are no
large-field regions, and $Y=\emptyset$.

\emph{Primed sets.} For each of the sets $P_t$, $Q_t$, $P_0$, $P_1$, $Q_1$ below, the primed set
($P'_t$, $Q'_t$, $P'_0$, $P'_1$, $Q'_1$) denotes the union of all coarse cubes meeting it, that is,
its $\mathcal{P}_t$-hull. Likewise Ba{\l}aban's $\bar Z_k$, the union of all $LMR_{k+1}$-cubes
intersecting $Z_k$, is the hull $\widehat{Y}$.

\emph{Two properties of enlargements.} Two properties of enlargements are used below. Both follow
from Definition~\ref{4.15:def-enlargement-hull}, where $X^{\sim}$ is $X$ together with every coarse
cube touching $X$, and $X^{\sim a}$ is this operation applied $a$ times.

First, a coarse cube touches $X_1\cup X_2$ if and only if it touches $X_1$ or touches $X_2$. Hence
$(X_1\cup X_2)^{\sim}=X_1^{\sim}\cup X_2^{\sim}$. Applying this $a$ times, and writing a
$\mathcal{P}_t$-region $X$ as the union of its coarse cubes $c$, gives
\begin{equation}\label{4.15:eq-small-field-complement-1}
(X_1\cup X_2)^{\sim a}=X_1^{\sim a}\cup X_2^{\sim a},\qquad
X^{\sim a}=\bigcup_{c\subset X}c^{\sim a}\qquad(a\ge 0).
\end{equation}
Second, $(X^{\sim a})^{\sim b}=X^{\sim(a+b)}$ for $a,b\ge 0$.

\emph{The general step, $2\le t\le K$} \cite[(3.2)--(3.5), p.~265]{B-CMP119}.
\begin{itemize}
\item[(a)] \emph{The set $P_t$.} Ba{\l}aban surrounds $\widehat{Y}$ by four layers of coarse cubes.
In the complement $(\widehat{Y}^{\sim 4})^c$ he introduces the decomposition of unity
\cite[(3.2)]{B-CMP119}. Its summation runs over domains $P_t\subset(\widehat{Y}^{\sim4})^c$ that are
unions of test cubes.
\item[(b)] \emph{The axial gauge fixing and the set $B_t(P''_t)$.} With $P'_t$ the union of the
coarse cubes meeting $P_t$, he introduces the axial gauge fixing in the complement of
$P_t'^{\sim1}\cup\widehat{Y}^{\sim5}$. Here and at the first step, for a subset $P''$ of the unit
lattice of level $t$ (Definition~\ref{1.5:def-block-average}), $B_t(P'')$ denotes the union of the
blocks of level $t$ over the points of $P''$, in the notation of \cite[\S3]{B-CMP119}. He takes the
subset $P''_t$ of the unit lattice of level $t$ determined by
\begin{equation*}
B_t(P''_t)=(P_t'^{\sim1})^c\cap(\widehat{Y}^{\sim5})^c .
\end{equation*}
Only the complement of this region is used below, namely
\begin{equation}\label{4.15:eq-small-field-complement-2}
B_t(P''_t)^c=P_t'^{\sim1}\cup\widehat{Y}^{\sim5}.
\end{equation}
\item[(c)] \emph{The set $Q_t$.} In the domain
\begin{equation*}
(P_t'^{\sim2}\cup\widehat{Y}^{\sim6})^c=\bigl(B_t(P''_t)\bigr)^{\sim-1}
\end{equation*}
he introduces the decomposition of unity \cite[(3.3)]{B-CMP119}, which restricts the approximate
fluctuation fields. Its summation runs over subdomains $Q_t\subset(B_t(P''_t))^{\sim-1}$ that are
unions of test cubes. Let $Q'_t$ be the union of the coarse cubes meeting $Q_t$.
\item[(d)] \emph{The domain $\Omega_t$.} Surrounding $Q'_t\cup((B_t(P''_t))^c)^{\sim}$ by two
layers of coarse cubes, he sets
\begin{equation*}
\Omega_t=\bigl(Q_t'^{\sim2}\cup(B_t(P''_t)^c)^{\sim3}\bigr)^c
\qquad\text{\cite[(3.5)]{B-CMP119}.}
\end{equation*}
On $\Omega_t$ only the small-field characteristic functions of these decompositions of unity
occur.
\end{itemize}
Insert \eqref{4.15:eq-small-field-complement-2} into this formula. Then apply the first identity of
\eqref{4.15:eq-small-field-complement-1} and $(X^{\sim a})^{\sim b}=X^{\sim(a+b)}$. This gives,
with $P':=P'_t$ and $Q':=Q'_t$,
\begin{equation}\label{4.15:eq-small-field-complement-a}
\Omega_t^c=Q'^{\sim2}\cup P'^{\sim4}\cup\widehat{Y}^{\sim8}\qquad(2\le t\le K).
\end{equation}
Thus \eqref{4.15:eq-small-field-complement-a} is Ba{\l}aban's definition \cite[(3.5)]{B-CMP119},
rewritten by means of his definition of $P''_t$ in the form
\eqref{4.15:eq-small-field-complement-2}.

\emph{The first step, $t=1$} \cite[(1.2), (1.4), (1.8), (1.10), pp.~246--248]{B-CMP119}. Here
$Y=\emptyset$. The decompositions of unity preceding $\Omega_1$ have three large-field sets.
\begin{itemize}
\item[(a)] \emph{The sets $P_0$ and $P_1$.} Both restrict plaquette variables. The set $P_0$ is the
large-field set of \cite[(1.2)]{B-CMP119}. The set $P_1$ is that of \cite[(1.4)]{B-CMP119}: the sum
there runs over sets $P_1\subset(P_0'^{\sim})^c$ that are unions of test cubes. Here $P'_0$ and
$P'_1$ are the hulls of $P_0$ and $P_1$, as above.
\item[(b)] \emph{The region $B_1(P''_1)$.} Once $P_0$ and $P_1$ are fixed, Ba{\l}aban takes a
subset $P''_1$ of the unit lattice of level $1$ and introduces the axial gauge fixing in the blocks
of $B_1(P''_1)$. This is the first-step instance of the general step, where the gauge fixing is
introduced in the complement of enlargements of the hulls of the large-field sets already fixed.
Here those sets are $P_0$ and $P_1$, and there is no incoming region. Accordingly, we read
Ba{\l}aban's definition of $P''_1$, given after \cite[(1.4)]{B-CMP119}, as
\begin{equation}\label{4.15:eq-small-field-complement-3}
B_1(P''_1)^c=P_0'^{\sim a_0}\cup P_1'^{\sim a_1}
\end{equation}
with the integers $a_0,a_1\ge 0$ that Ba{\l}aban's definition after \cite[(1.4)]{B-CMP119}
prescribes; their values play no role below. Display
\eqref{4.15:eq-small-field-complement-3} is not a formula of \cite{B-CMP119}; it is the reading of
that definition adopted here.
\item[(c)] \emph{The set $Q_1$.} The set $Q_1$ restricts the approximate fluctuation field, as at
the general step. It is the large-field set of \cite[(1.8)]{B-CMP119}. The sum there runs over sets
\begin{equation*}
Q_1\subset B_1(P''_1)^{\sim-1}=\bigl((B_1(P''_1)^c)^{\sim}\bigr)^c
\end{equation*}
that are unions of test cubes. Here $Q'_1$ is the hull of $Q_1$, as above.
\item[(d)] \emph{The domain $\Omega_1$.} Surrounding $Q'_1$ and $(B_1(P''_1)^c)^{\sim}$ by two
layers of coarse cubes, Ba{\l}aban sets
\begin{equation*}
\Omega_1=\bigl(Q_1'^{\sim2}\cup(B_1(P''_1)^c)^{\sim3}\bigr)^c
\qquad\text{\cite[(1.10)]{B-CMP119}.}
\end{equation*}
Thus $\Omega_1$ is a union of coarse cubes. Its distance from the union of the large-field regions
is at least $2LM\check{R}_1$, and on $\Omega_1$ only the small-field characteristic functions occur.
The resummation is then done with respect to $P_0$, $P_1$ and $Q_1$, with $\Omega_1$ fixed.
\end{itemize}
Insert \eqref{4.15:eq-small-field-complement-3}, and use
\eqref{4.15:eq-small-field-complement-1} and $(X^{\sim a})^{\sim 3}=X^{\sim(a+3)}$. This gives
\begin{equation}\label{4.15:eq-small-field-complement-b}
\Omega_1^c=Q_1'^{\sim2}\cup P_0'^{\sim(a_0+3)}\cup P_1'^{\sim(a_1+3)},\qquad Y=\emptyset .
\end{equation}
Thus \eqref{4.15:eq-small-field-complement-b} is Ba{\l}aban's definition \cite[(1.10)]{B-CMP119},
rewritten by means of the reading \eqref{4.15:eq-small-field-complement-3}.
At the first step the primed sets of the step are, by definition,
\begin{equation*}
P':=P'_0\cup P'_1,\qquad Q':=Q'_1 .
\end{equation*}
Thus $P'$ is the union of the coarse cubes meeting $P_0\cup P_1$. At the general step,
$P':=P'_t$ and $Q':=Q'_t$.

\emph{The form of the small-field complement.} For every step $t-1\to t$ with $1\le t\le K$, and
every term the step produces, $\Omega_t^c$ is a $\mathcal{P}_t$-region of the form
\begin{equation}\label{4.15:eq-small-field-complement-c}
\Omega_t^c=\bigcup_{i\in\mathsf{I}}G_i^{\sim a_i}\;\cup\;\widehat{Y}^{\sim a_Y}.
\end{equation}
Here:
\begin{itemize}
\item $\mathsf{I}$ is finite;
\item each $G_i$ is a $\mathcal{P}_t$-region contained in $P'\cup Q'$;
\item $a_i\ge0$ and $a_Y\ge0$ are integers;
\item $\widehat{Y}$ is the $\mathcal{P}_t$-hull of the input region $Y$ of the step;
\item $Y=\emptyset$ at $t=1$, so that the last member is then empty.
\end{itemize}
Concretely, \eqref{4.15:eq-small-field-complement-c} is \eqref{4.15:eq-small-field-complement-a}
for $2\le t\le K$, and \eqref{4.15:eq-small-field-complement-b} at $t=1$. The numbers of layers
occurring there play no role in what follows. It is in the form
\eqref{4.15:eq-small-field-complement-c}, with the conditions just listed, that the small-field
complement of a step enters as a hypothesis the statements of this section that follow; at $t=1$
this form rests on the reading \eqref{4.15:eq-small-field-complement-3}.

By the second identity of \eqref{4.15:eq-small-field-complement-1},
\eqref{4.15:eq-small-field-complement-c} says that $\Omega_t^c$ is the union of
$\widehat{Y}^{\sim a_Y}$ and of sets $c^{\sim a(c)}$, where $c$ ranges over a set of coarse cubes of
$P'\cup Q'$ and $a(c)\ge0$ is an integer (the largest $a_i$ with $c\subset G_i$).
In particular, $\Omega_t^c$ is the union of enlargements of $Q'$, of $P'$, and of the large-field
region of the previous step; at $t=1$ there is no such region. At $t=1$ this holds by
\eqref{4.15:eq-small-field-complement-b} with $P'=P'_0\cup P'_1$. Under the narrower choice
$P':=P'_1$ it can fail when $P_0\ne\emptyset$, since the member $P_0'^{\sim(a_0+3)}$ of
\eqref{4.15:eq-small-field-complement-b} need not be covered by enlargements of $Q'_1\cup P'_1$.

The sets $P_t$, $Q_t$, $P''_t$ and $\Omega_t$ are those of \cite[(3.2)--(3.5), p.~265]{B-CMP119}.
At the first step $P_0$, $P_1$, $Q_1$, $P''_1$ and $\Omega_1$ are those of
\cite[(1.2), (1.4), (1.8), (1.10), pp.~246--248]{B-CMP119}.
\end{definition}

\begin{definition}[The large-field region of a step and its creation set]\label{4.15:def-creation-set}
Let $1\le t\le K$, and consider the step $t-1\to t$ of the renormalisation transformation
$\mathbf{T}$, acting on a term of $\rho_{t-1}$. Let $\mathcal{P}_t$ be the coarse partition of
the step, $\Omega_t$ Ba{\l}aban's small-field domain of the step, and $P'$, $Q'$ the
$\mathcal{P}_t$-hulls of the large-field sets of Ba{\l}aban's decompositions of unity, as in
Definition~\ref{4.15:def-small-field-complement}; at $t=1$ these are $P'=P_0'\cup P_1'$ and
$Q'=Q_1'$. Enlargements $X^{\sim n}$, the removal of one layer $X^{\sim -1}$, $\mathcal{P}_t$-hulls
and the stabilised boxing map $\mathsf{M}$ are those of
Definition~\ref{4.15:def-enlargement-hull}. A set of $\mathcal{P}_t$-cubes is identified with
the union of its members whenever it enters a union or an enlargement of regions.

\emph{The decomposition.} Ba{\l}aban's decomposition \cite[(3.16)]{B-CMP119} (at the first step
\cite[(1.18)]{B-CMP119}) is replaced by the decomposition fixed in the definition of the
large-field recursion of the step; its outcomes are labels
$\ell=(R,I,\mathsf{R}^{\mathrm{lo}})$, and the bond test, the cube test and the strict growth of
the cube sets $\mathsf{R}_n$ used below are those of that definition. Here $R$ is the large-field
set of the outcome, a union of test cubes (the cubes of side $L^2M_2\check{R}_t$, in the units in
which the $\mathcal{P}_t$-cubes have side $LM\check{R}_t$) contained in
$\Omega_t^{\sim -1}$, and $R'$ is the $\mathcal{P}_t$-hull of $R$.

\emph{Cubes attached to bonds.} Let $\mathsf{K}$ be the set of the coarse cubes of
$\Omega_t^{\sim -1}$; of $\mathsf{K}$ only the fact that its members are $\mathcal{P}_t$-cubes is
used. For a set $J$ of bonds of the step, $\mathsf{K}(J)$ is the set of the cubes of
$\mathsf{K}$ holding an end-point of a bond of $J$. Thus
$\mathsf{K}(J\cup J')=\mathsf{K}(J)\cup\mathsf{K}(J')$ for any two sets $J$, $J'$ of bonds.

\emph{Stage sets.} For a set $\mathsf{R}$ of $\mathcal{P}_t$-cubes put
\begin{equation}\label{4.15:eq-creation-set-1}
Z[\mathsf{R}]:=(\Omega_t^c)^{\sim 2}\cup\mathsf{R}^{\sim},
\end{equation}
which for $\mathsf{R}=R'$ is the complement of the region written, following
\cite[(3.20)]{B-CMP119},
\begin{equation*}
\Lambda_{k+1}=\bigl((\Omega^c_{k+1})^{\sim 2}\cup R'^{\sim}_{k+1}\bigr)^c,
\end{equation*}
with $k+1=t$ and $R'_{k+1}=R'$; the region denoted $\Lambda_t$ in
\eqref{4.15:eq-creation-set-3} below is contained in it.

\emph{The stages.} Start with $\mathsf{R}_0:=R'$ and $Z^{[0]}:=\mathsf{M}Z[R']$. At stage
$n\ge0$, given $(\mathsf{R}_n,Z^{[n]})$, put $\Lambda^n:=(Z^{[n]})^c$. A bond test on a set
$\mathfrak{B}_n$ of bonds near $\Lambda^n$ produces a subset $I_n\subset\mathfrak{B}_n$. Let
$\mathfrak{C}_n$ be the set of the cubes of $\mathsf{K}$ in the rim of $Z^{[n]}$, that is, the
cubes of $\mathsf{K}$ contained in $Z^{[n]}$ and touching its complement, that have not been
tested before; a cube test on $\mathfrak{C}_n$ produces a subset
$\mathsf{R}^{\mathrm{lo}}_n\subset\mathfrak{C}_n$. The two tests are those of the definition of
the large-field recursion; their nature is immaterial here, and only the facts that
$I_n$ is a set of bonds
and $\mathsf{R}^{\mathrm{lo}}_n$ a set of cubes of $\mathsf{K}$ are used. Then
\begin{equation}\label{4.15:eq-creation-set-2}
\mathsf{R}_{n+1}:=\mathsf{R}_n\cup\mathsf{K}(I_n)\cup\mathsf{R}^{\mathrm{lo}}_n,
\qquad
Z^{[n+1]}:=\mathsf{M}\bigl(Z^{[n]}\cup Z[\mathsf{R}_{n+1}]\bigr).
\end{equation}
The run stops at the first $n_*$ with
$I_{n_*}=\mathsf{R}^{\mathrm{lo}}_{n_*}=\emptyset$. For every outcome that labels a term of the
step such an $n_*$ exists: before the stop the sets $\mathsf{R}_n$, all contained in the finite
set of $\mathcal{P}_t$-cubes of $\mathbb{T}$, grow strictly, as shown where the large-field
recursion is defined. No bound on $n_*$ is used.

\emph{Output.} The outcome has
$I:=\bigcup_{n<n_*}I_n$ and $\mathsf{R}^{\mathrm{lo}}:=\bigcup_{n<n_*}\mathsf{R}^{\mathrm{lo}}_n$,
and
\begin{equation}\label{4.15:eq-creation-set-3}
\Lambda_t:=\bigl(Z^{[n_*]}\bigr)^c,
\qquad\text{so that}\qquad
Z_t:=\mathbb{T}\setminus\Lambda_t=Z^{[n_*]}
\end{equation}
is the \emph{large-field region of scale $t$} of the resulting term.

\emph{The creation set.} The \emph{creation set} of the step $t-1\to t$ is
\begin{equation}\label{4.15:eq-creation-set-a}
\mathcal{W}^{\wedge}_t:=P'\cup Q'\cup R'\cup\mathsf{K}(I)\cup\mathsf{R}^{\mathrm{lo}}.
\end{equation}
A component of $\Lambda_t^c$ that meets
$\mathcal{W}^{\wedge}_t$ is said to have a \emph{new large-field region created in it at step
$t$}.

With these definitions, every stage set $Z^{[n]}$ is a $\mathcal{P}_t$-region ($\Omega_t^c$ being
one by \eqref{4.15:eq-small-field-complement-c}),
\begin{equation}\label{4.15:eq-creation-set-b}
\begin{split}
\mathsf{R}_0\subset\mathsf{R}_1\subset\dots\subset\mathsf{R}_{n_*}
&=R'\cup\bigcup_{n<n_*}\bigl(\mathsf{K}(I_n)\cup\mathsf{R}^{\mathrm{lo}}_n\bigr)\\
&=R'\cup\mathsf{K}(I)\cup\mathsf{R}^{\mathrm{lo}},\\
P'\cup Q'\cup\mathsf{R}_{n_*}&=\mathcal{W}^{\wedge}_t,
\end{split}
\end{equation}
the stage sets increase, $Z^{[n+1]}\supseteq Z^{[n]}$ for $n<n_*$, and $\mathcal{W}^{\wedge}_t$
is a $\mathcal{P}_t$-region.
\end{definition}

\begin{proof}[Proof of \eqref{4.15:eq-creation-set-b} and of the closing assertions of
Definition~\ref{4.15:def-creation-set}]
\emph{Monotonicity of the stage sets.} By Definition~\ref{4.15:def-enlargement-hull},
$\mathsf{M}_1X\supseteq X$ for every region $X$, and hence
also $\mathsf{M}X\supseteq X$ for the stabilised iterate. Applied to
$X=Z^{[n]}\cup Z[\mathsf{R}_{n+1}]$,
the second relation of \eqref{4.15:eq-creation-set-2} gives
\begin{equation*}
Z^{[n+1]}\supseteq Z^{[n]}\cup Z[\mathsf{R}_{n+1}]\supseteq Z^{[n]}.
\end{equation*}

\emph{The stage sets are $\mathcal{P}_t$-regions.} For every set $\mathsf{R}$ of
$\mathcal{P}_t$-cubes, $Z[\mathsf{R}]$ is a $\mathcal{P}_t$-region, being by
\eqref{4.15:eq-creation-set-1} a union of enlargements of $\Omega_t^c$, which is a
$\mathcal{P}_t$-region by \eqref{4.15:eq-small-field-complement-c}, and of $\mathsf{R}$, and an
enlargement of a $\mathcal{P}_t$-region adds layers of touching $\mathcal{P}_t$-cubes. The map
$\mathsf{M}_1$ adds to a region the sets $\operatorname{box}(C)$, each a union of
$\mathcal{P}_t$-cubes, so $\mathsf{M}_1$ and its stabilised iterate $\mathsf{M}$ map
$\mathcal{P}_t$-regions to $\mathcal{P}_t$-regions (Definition~\ref{4.15:def-enlargement-hull}).
Hence $Z^{[0]}=\mathsf{M}Z[R']$ is a $\mathcal{P}_t$-region, and if $Z^{[n]}$ is one, so is
$Z^{[n+1]}$ by the second relation of \eqref{4.15:eq-creation-set-2}; by induction on $n$ every
stage set is a $\mathcal{P}_t$-region.

\emph{The chain of cube sets.} For $n<n_*$ the first relation of
\eqref{4.15:eq-creation-set-2} reads
$\mathsf{R}_{n+1}=\mathsf{R}_n\cup\mathsf{K}(I_n)\cup\mathsf{R}^{\mathrm{lo}}_n$, so
$\mathsf{R}_n\subset\mathsf{R}_{n+1}$; this is the chain of inclusions in
\eqref{4.15:eq-creation-set-b}. Unwinding the same relation from $n=0$, where
$\mathsf{R}_0=R'$, up to $n=n_*-1$ gives by induction on $n$
\begin{equation*}
\mathsf{R}_{n}=R'\cup\bigcup_{n'<n}\bigl(\mathsf{K}(I_{n'})\cup\mathsf{R}^{\mathrm{lo}}_{n'}\bigr),
\qquad 0\le n\le n_*,
\end{equation*}
and at $n=n_*$ this is the first equality of \eqref{4.15:eq-creation-set-b}. Since $\mathsf{K}$
is union-preserving, $\mathsf{K}(J\cup J')=\mathsf{K}(J)\cup\mathsf{K}(J')$, and the union
defining $I$ is finite,
\begin{equation*}
\bigcup_{n<n_*}\mathsf{K}(I_n)=\mathsf{K}\Bigl(\bigcup_{n<n_*}I_n\Bigr)=\mathsf{K}(I),
\end{equation*}
while $\bigcup_{n<n_*}\mathsf{R}^{\mathrm{lo}}_n=\mathsf{R}^{\mathrm{lo}}$ by definition. This is
the second equality of \eqref{4.15:eq-creation-set-b}. Taking the union with
$P'\cup Q'$ and comparing with \eqref{4.15:eq-creation-set-a} gives
$P'\cup Q'\cup\mathsf{R}_{n_*}=\mathcal{W}^{\wedge}_t$.

\emph{The creation set is a $\mathcal{P}_t$-region.} By the last line of
\eqref{4.15:eq-creation-set-b}, $\mathcal{W}^{\wedge}_t$ is the union of $P'$, $Q'$ and the
members of $\mathsf{R}_{n_*}=R'\cup\mathsf{K}(I)\cup\mathsf{R}^{\mathrm{lo}}$. The sets $P'$, $Q'$
and $R'$ are $\mathcal{P}_t$-hulls, hence unions of $\mathcal{P}_t$-cubes; $\mathsf{K}(I)$ and
$\mathsf{R}^{\mathrm{lo}}$ are sets of cubes of $\mathsf{K}$, whose members are
$\mathcal{P}_t$-cubes. A union of $\mathcal{P}_t$-cubes is a $\mathcal{P}_t$-region, and so
$\mathcal{W}^{\wedge}_t$ is one.
\end{proof}

\begin{definition}[Histories of a term and the components acted upon]\label{4.15:def-histories}
\emph{Terms and labels.} By Definition~\ref{1.5:def-densities}, for $0\le k\le K$ the measure
$\rho_k\,dV_k$ is a finite sum of terms.
Each term carries a label $\mathcal{S}=(\Omega_j,\Lambda_j)_{1\le j\le k}$, the sequence of
small-field regions created at the steps up to $k$. The set $Z_j:=\mathbb{T}\setminus\Lambda_j$
is the large-field region of scale $j$ of the term. For a term $\sigma$ we write $\Lambda_j(\sigma)$
and $Z_j(\sigma):=\Lambda_j(\sigma)^c$.

The density $\rho_0$ has one term, without characteristic functions, and there is no large-field
region at level $0$. For $1\le j\le K$ one has $\rho_j=\mathbf{R}\mathbf{T}\rho_{j-1}$
(Definition~\ref{1.5:def-densities}). The step
$j-1\to j$ of $\mathbf{T}$ replaces each term $\sigma$ of $\rho_{j-1}$ by a finite sum of terms
$\sigma'$, one for each outcome of the step's decompositions of unity. Each $\sigma'$ comes with its
region $Z_j(\sigma')=\Lambda_j(\sigma')^c$, produced by the recursion of
Definition~\ref{4.15:def-creation-set}, and with its creation set $\mathcal{W}^{\wedge}_j(\sigma')$ of
\eqref{4.15:eq-creation-set-a}. Then $\mathbf{R}$ acts on every component of the large-field region
that has become eligible for it, at the level fixed by the rule for the level of action, which we
recall next.

\emph{The level at which $\mathbf{R}$ acts.} Let $j\ge1$, let $\sigma'$ be a term of
$\mathbf{T}\rho_{j-1}$, and let $C$ be a component of $Z_j(\sigma')$. The operation $\mathbf{R}$ acts
on $C$ at level $j$ if the following three conditions hold for $C$ at level $j$ and $j$ is the first
level at which they hold:
\begin{itemize}
\item[(i)] $C$ is contained in a cube of side $100M\check{R}_j$;
\item[(ii)] condition (ii) of Definition~\ref{4.15:def-no-creation} holds for $C$ at level $j$, in
the reading adopted there. In both formulations recorded there it contains the no-creation clause:
no large-field region created at a level $t\in\{1,\dots,j\}$ with $t>j-\mathsf{N}_j$ lies inside
$C$, where $\mathsf{N}_j$ is the waiting depth at level $j$. The formulation which also requires
that $\mathsf{N}_j$ steps precede level $j$ adds the condition $j\ge\mathsf{N}_j$;
\item[(iii)] no component nested strictly inside $C$ is still waiting, that is, not yet acted upon
by $\mathbf{R}$.
\end{itemize}
The level of an action is the scale $j$ of the region $\Lambda_j^c$ acted upon. If instead the action
on a component of the large-field region of a term of $\mathbf{T}\rho_k$ is attached to the index $k$
of that density, then, since a component of $Z_j(\sigma')$ with $\sigma'$ a term of
$\mathbf{T}\rho_{j-1}$ carries the index $k=j-1$, the index so attached differs by one from the level
$j$.

Let $\mathcal{A}_j(\sigma')$ be the set of the components of $Z_j(\sigma')$ on which $\mathbf{R}$ acts
at level $j$, that is, of those components $C$ of $Z_j(\sigma')$ for which $j$ is the first level at
which conditions (i), (ii), (iii) hold; it may be empty. Membership in $\mathcal{A}_j(\sigma')$ is
decided by these conditions alone, and the phrase ``$\mathbf{R}$ acts on $C$ at level $j$'' means
$C\in\mathcal{A}_j(\sigma')$. This is Ba{\l}aban's class, whose union he denotes by $Z$
\cite[pp.~176--177]{B-CMP122a}; we denote that union by $\mathcal{Z}$. The actions of level $j$
replace $\sigma'$ by a finite sum of terms $\tau$, the outcomes of the decompositions of unity inside
$\mathbf{R}$ \cite{B-CMP122a}, \cite{B-CMP122b}. Each $\tau$ is a term of $\rho_j$ with a label
$(\Omega_i(\tau),\Lambda_i(\tau))_{i\le j}$.

By the definition of the level of action the following two statements hold; they are proved below.
First, the conditions of the level of action hold at that level:
\begin{equation}\label{4.15:eq-histories-a}
C\in\mathcal{A}_j(\sigma')\ \Longrightarrow\ \text{conditions (i), (ii), (iii) hold for $C$ at
level $j$.}
\end{equation}
In particular, the no-creation clause of condition (ii) holds for $C$ at level $j$, in the reading
adopted; the clause is the same in both formulations recorded in
Definition~\ref{4.15:def-no-creation}. Second, if $\mathcal{A}_j(\sigma')=\emptyset$, no operation is
applied to $\sigma'$ at level $j$. Hence $\sigma'$ is itself a term $\tau$ of $\rho_j$, with its
label unchanged, and
\begin{equation}\label{4.15:eq-histories-b}
\begin{aligned}
\mathcal{A}_j(\sigma')=\emptyset\ \Longrightarrow\ {}&\tau=\sigma'\quad\text{and}\\
&\Lambda_j(\tau)^c=\Lambda_j(\sigma')^c=Z_j(\sigma')=Z^{[n_*]},
\end{aligned}
\end{equation}
where $Z^{[n_*]}$ is the final region of the recursion of Definition~\ref{4.15:def-creation-set}.

\emph{Histories.} A \emph{history} of a term $\tau$ of $\rho_k$ is a sequence
$\tau_0,\tau_1,\dots,\tau_k=\tau$ with the following properties: $\tau_0$ is the term of $\rho_0$,
and, for $1\le j\le k$, $\tau_j$ is one of the terms produced from $\tau_{j-1}$ by the step
$j-1\to j$ followed by the actions of level $j$. A history thus fixes the outcomes of the
decompositions of unity at the steps up to $k$.

A history comes with the following data, for $1\le j\le k$:
\begin{itemize}
\item the term $\sigma'_j$ of $\mathbf{T}\rho_{j-1}$ from which $\tau_j$ results;
\item the region $Z_j:=Z_j(\sigma'_j)$;
\item the creation set $\mathcal{W}^{\wedge}_j:=\mathcal{W}^{\wedge}_j(\sigma'_j)$;
\item the class $\mathcal{A}_j:=\mathcal{A}_j(\sigma'_j)$;
\item the region $Z_j^{\mathrm{post}}:=\Lambda_j(\tau_j)^c$.
\end{itemize}
Along a history the symbols $Z_j$, $\mathcal{W}^{\wedge}_j$ and $\mathcal{A}_j$ without argument
refer to $\sigma'_j$; the large-field region of scale $j$ of $\tau_j$ itself is
$Z_j^{\mathrm{post}}$.
The region $Z_j^{\mathrm{post}}$ is the large-field region of scale $j$ of $\tau_j$ after the actions
of level $j$. For $j<K$ it is the region $Y$ read by the step $j\to j+1$
\cite[p.~265]{B-CMP119}. We set $Z_0^{\mathrm{post}}:=\emptyset$. A history of a term $\sigma'$ of
$\mathbf{T}\rho_{k-1}$ is defined in the same way, ending with the step $k-1\to k$.

Conditions (i)--(iii) at level $j$ for a component of $Z_j$ refer to the creation sets
$\mathcal{W}^{\wedge}_t$, $t\le j$, of the history. In particular, the pair
$(I,\mathsf{R}^{\mathrm{lo}})$ of the label of an outcome of a step enters these conditions through
the creation set, and hence through the waiting rule, condition (ii). By \eqref{4.15:eq-histories-b},
\begin{equation}\label{4.15:eq-histories-c}
\mathcal{A}_j=\emptyset\ \Longrightarrow\ \tau_j=\sigma'_j\quad\text{and}\quad
Z_j^{\mathrm{post}}=Z_j .
\end{equation}

The history is \emph{pure-$\mathbf{T}$ below level $k$} if
\begin{equation}\label{4.15:eq-histories-1}
\mathcal{A}_j=\emptyset\qquad\text{for every $j$ with $1\le j<k$.}
\end{equation}
Every term has at least one history; this too is proved below.
\end{definition}

\begin{proof}
We prove \eqref{4.15:eq-histories-a}, \eqref{4.15:eq-histories-b}, the existence of histories, and
\eqref{4.15:eq-histories-c}, in this order.

\emph{Proof of \eqref{4.15:eq-histories-a}.} Let $C\in\mathcal{A}_j(\sigma')$, so that
$\mathbf{R}$ acts on $C$ at level $j$. By the definition of the level at which $\mathbf{R}$ acts,
$j$ is the first level at which conditions (i), (ii), (iii) hold for $C$; in particular they hold at
level $j$. Condition (ii) is taken in the reading adopted, and in both formulations of
Definition~\ref{4.15:def-no-creation} it contains the no-creation clause. Hence, in particular, the
no-creation clause holds for $C$ at level $j$.

\emph{Proof of \eqref{4.15:eq-histories-b}.} By Definition~\ref{1.5:def-densities} and
\cite[pp.~176--177]{B-CMP122a}, the operation $\mathbf{R}$ consists of actions on
components of the large-field region. On a component, and at the level fixed by the definition
above, $\mathbf{R}$ performs its action. While a component waits, no operation that changes the
density acts on it. In Ba{\l}aban's construction the operation of a level is set up on the union of
the class of components acted upon, the set we call $\mathcal{Z}$. Outside that union nothing is
changed: in the notation of \cite[pp.~176--177]{B-CMP122a}, where the union is written $Z$ and the
large-field region of scale $k$ is written $Z_k$,
``we do not make any changes in the operation $T_k(Z_k\cap Z^c)$''.

Now let $\mathcal{A}_j(\sigma')=\emptyset$. Then $\mathcal{Z}=\emptyset$ and $\mathbf{R}$ acts at
level $j$ on no component of $Z_j(\sigma')$. Since $\mathbf{R}$ consists of actions on components,
no operation is applied to $\sigma'$ at level $j$. The actions of level $j$ therefore produce from
$\sigma'$ the single term $\sigma'$ itself. It is a term of $\rho_j$ with its label unchanged, so
$\tau=\sigma'$, and consequently
\begin{equation*}
\Lambda_j(\tau)^c=\Lambda_j(\sigma')^c=Z_j(\sigma').
\end{equation*}
Finally, by Definition~\ref{4.15:def-creation-set} the region $Z_j(\sigma')$ is the region produced
by the recursion of the step $j-1\to j$. This is the region $Z^{[n_*]}$ at the stage $n_*$ at which
the recursion stops.

\emph{Existence of histories.} Let $\tau$ be a term of $\rho_k$; we descend level by level. If
$k=0$, the one-member sequence $\tau_0=\tau$ is a history, since $\rho_0$ has exactly one term. If
$k\ge1$, then $\rho_k=\mathbf{R}\mathbf{T}\rho_{k-1}$. Every term of $\rho_k$ is therefore produced
by the step $k-1\to k$, followed by the actions of level $k$, from some term of $\rho_{k-1}$; choose
such a term and call it $\tau_{k-1}$. Applying the same argument to $\tau_{k-1}$, and so on, yields
terms $\tau_{k-1},\tau_{k-2},\dots,\tau_0$. Here $\tau_0$ is a term of $\rho_0$, hence the term of
$\rho_0$, and each $\tau_j$ is produced from $\tau_{j-1}$ as the definition requires. Thus
$\tau_0,\dots,\tau_k=\tau$ is a history of $\tau$. The same argument, with the last step taken in
$\mathbf{T}$ only, gives a history of every term of $\mathbf{T}\rho_{k-1}$.

\emph{Proof of \eqref{4.15:eq-histories-c}.} Let $\mathcal{A}_j=\mathcal{A}_j(\sigma'_j)=\emptyset$.
The term $\tau_j$ is one of the terms produced from $\sigma'_j$ by the actions of level $j$. Apply
\eqref{4.15:eq-histories-b} with $\sigma'=\sigma'_j$ and $\tau=\tau_j$. It gives $\tau_j=\sigma'_j$
and
\begin{equation*}
Z_j^{\mathrm{post}}=\Lambda_j(\tau_j)^c=Z_j(\sigma'_j)=Z_j. \qedhere
\end{equation*}
\end{proof}

\begin{lemma}[The cube sizes do not increase]\label{4.15:lem-cube-sizes-monotone}
Let $L\ge 2$ be an integer, let $r>0$, let $0\le k\le K$, and let $x_0,\dots,x_k$ be real
numbers with $x_j>1$ for every $j\le k$. Then for every $j\le k$ the number $\check{R}_j$ of
Definition~\ref{4.15:not-cube-sizes} is well defined and is a power $L^{b_j}$ with
$b_j\in\mathbb{Z}_{\ge 0}$, so that $\check{R}_j\ge 1$. Moreover $\check{R}_{j+1}\le\check{R}_j$
for every $j<k$; hence $\check{R}_k\le\check{R}_j$ for every $j\le k$.
\end{lemma}

\begin{proof}
Fix $j\le k$. The number $\check{x}_j=\min_{i\le j}x_i$ is the minimum of finitely many real
numbers, each of which exceeds $1$ by hypothesis; hence $\check{x}_j>1$. Therefore
$\log\check{x}_j>0$, and since $r>0$ the power $(\log\check{x}_j)^r$ is a well-defined
positive real number.

Since $L\ge 2$, Bernoulli's inequality gives $L^b\ge 2^b\ge 1+b$ for every
$b\in\mathbb{Z}_{\ge 0}$. For $j\le k$ let
\begin{equation*}
\mathcal{B}_j:=\bigl\{\,b\in\mathbb{Z}_{\ge 0}\;:\;L^b\ge(\log\check{x}_j)^r\,\bigr\}.
\end{equation*}
Any integer $b\ge(\log\check{x}_j)^r$ satisfies $L^b\ge 1+b>(\log\check{x}_j)^r$, so
$\mathcal{B}_j$ is non-empty. Being a non-empty set of non-negative integers, it has a least
element, which we call $b_j$. By Definition~\ref{4.15:not-cube-sizes}, $\check{R}_j$ is the
least power $L^b$ with $b\in\mathbb{Z}_{\ge 0}$ and $L^b\ge(\log\check{x}_j)^r$. Since
$b\mapsto L^b$ is strictly increasing, this least power is $L^{b_j}$. Thus
$\check{R}_j=L^{b_j}$ is well defined, and $\check{R}_j\ge L^0=1$.

Now let $j<k$. By the definition of the minimum,
\begin{equation*}
\check{x}_{j+1}=\min_{i\le j+1}x_i=\min\{\check{x}_j,\,x_{j+1}\}\le\check{x}_j .
\end{equation*}
By the first paragraph applied at $j+1$ and at $j$, both numbers exceed $1$. Hence
$0<\log\check{x}_{j+1}\le\log\check{x}_j$, since the logarithm is increasing. The map
$u\mapsto u^r$ is increasing on $]0,\infty[$ because $r>0$. Together with $b_j\in\mathcal{B}_j$
this gives
\begin{equation*}
(\log\check{x}_{j+1})^r\le(\log\check{x}_j)^r\le L^{b_j}.
\end{equation*}
Thus $b_j\in\mathcal{B}_{j+1}$. Since $b_{j+1}$ is the least element of $\mathcal{B}_{j+1}$,
we get $b_{j+1}\le b_j$. As $b\mapsto L^b$ is increasing, it follows that
$\check{R}_{j+1}=L^{b_{j+1}}\le L^{b_j}=\check{R}_j$.

Finally, let $j\le k$. If $j=k$ there is nothing to prove. If $j<k$, we apply the inequality
just proved at $j,j+1,\dots,k-1$ in turn and obtain
\begin{equation*}
\check{R}_k\le\check{R}_{k-1}\le\dots\le\check{R}_j .
\end{equation*}
\end{proof}

\begin{remark}
Within Theorem~0 the hypotheses of Lemma~\ref{4.15:lem-cube-sizes-monotone} are met as
follows. The inequality $x_j>1$ is \eqref{4.15:eq-cube-sizes-a}, the part of clause~(0) of
Theorem~0 (Theorem~\ref{4.1:thm-clause-zero}) used in this section. The condition $L\ge 2$ is
contained in the lower bound that hypothesis (H2) imposes on the odd integer $L$, which
contains $L\ge 5$. The condition $r>0$ is part of hypothesis (H5), which fixes the exponent
$r$. None of the bounds on the coupling flow enters the proof.

This is the monotonicity of the cube sizes that is used in what follows; the bound
$\check{R}_j\le L^{k-j}\check{R}_k$ rests on part~(a) of the lemma on cube sizes, which
concerns the same numbers $\check{R}_j$. Here the monotonicity has been derived
from Definition~\ref{4.15:not-cube-sizes} and the inequality $\check{x}_j>1$ alone.
\end{remark}

\begin{lemma}[Elementary properties of components]\label{4.15:lem-component-properties}
Let $X$ be a region in the sense of Definition~\ref{4.15:def-regions-components}.
\begin{enumerate}
\item[(a)] The components of $X$ are finitely many, non-empty, connected regions. They are
pairwise non-touching, and hence have pairwise disjoint point sets. Their union is $X$.
\item[(b)] A non-empty connected region $A\subset X$ is contained in exactly one component of $X$.
\item[(c)] If $X\subset X'$, then every component of $X$ is contained in exactly one component
of $X'$.
\item[(d)] Let $A,A_1,\dots,A_m$ be connected regions such that each $A_i$ touches $A$. Then
$A\cup A_1\cup\dots\cup A_m$ is connected.
\item[(e)] Let $X$ be the union of a finite family $(A_i)_{i\in\mathsf{I}}$ of non-empty
connected regions. Then every component of $X$ is the union of those $A_i$ that it contains,
and it contains at least one of them.
\item[(f)] Let $X'$ be a region and $\mathcal{C}'$ a non-empty set of components of $X'$, and
suppose $X=\bigcup\mathcal{C}'$. Then the components of $X$ are exactly the members of
$\mathcal{C}'$.
\end{enumerate}
\end{lemma}

\begin{proof}
Throughout we use the notions of region, touching, chain, connectedness and component of
Definition~\ref{4.15:def-regions-components}.

(a) By definition the components of $X$ are the classes of the equivalence relation ``joined
by a chain in $X$'' on the set of members of $X$. These classes partition the finite set
$X$ of $\mathcal{P}_{\bullet}$-cubes. There are therefore finitely many of them, each is
non-empty, they are pairwise disjoint as sets of cubes, and their union is $X$.

Suppose a member $b$ of one class touched a member $b'$ of another class. Then the
two-term sequence $(b,b')$ would be a chain in $X$, so $b$ and $b'$ would be joined, and
they would lie in the same class, which is a contradiction. Hence distinct components do
not touch.

Next let $E$ be a class and let $b,b'\in E$. By definition there is a chain
$(b=c_0,c_1,\dots,c_r=b')$ in $X$. Every $c_l$ is joined to $b$ by the initial segment
$(c_0,\dots,c_l)$, which is a chain in $X$, so every $c_l$ belongs to $E$. The chain
therefore lies in $E$, and it is a chain in $E$. Since $E$ is non-empty, $E$ is connected.

Finally, regions that do not touch have disjoint point sets, by
Definition~\ref{4.15:def-regions-components}. Applied to distinct components, this gives the
last assertion of (a).

(b) Let $A\subset X$ be non-empty and connected, and let $a,a'\in A$. There is a chain in
$A$ with ends $a$ and $a'$. Since $A\subset X$, this is also a chain in $X$, so $a$ and $a'$
are joined by a chain in $X$. Consequently all members of $A$ lie in one class, that is, $A$
is contained in one component of $X$. This component is unique: $A\neq\emptyset$, and by (a)
the components are pairwise disjoint, so no two of them can both contain a member of $A$.

(c) Let $E$ be a component of $X$. By (a), $E$ is a non-empty connected region, and
$E\subset X\subset X'$. Part (b), applied to $E$ with $X'$ in place of $X$, shows that $E$
is contained in exactly one component of $X'$.

(d) Put $B=A\cup A_1\cup\dots\cup A_m$. The region $A$ is connected, hence non-empty, so $B$
is non-empty.

For each $i$, since $A_i$ touches $A$, there are a member $b_i$ of $A_i$ and a member $a_i$
of $A$ that touch. Let $b\in A_i$. Because $A_i$ is connected, $b$ is joined to $b_i$ by a
chain in $A_i$. Appending $a_i$, which touches $b_i$, gives a chain in $B$ from $b$ to
$a_i\in A$. Any two members of $A$ are joined by a chain in $A$, and hence by a chain in $B$.

Now let $b,b'\in B$. Each of them is either a member of $A$ or is joined by a chain in $B$
to a member of $A$. Concatenating such a chain from $b$ to $A$, a chain in $A$, and the
reverse of such a chain from $b'$ to $A$ gives a chain in $B$ with ends $b$ and $b'$. Hence
$B$ is connected.

(e) Let $E$ be a component of $X$ and let $a\in E$. Since $X=\bigcup_{i\in\mathsf{I}}A_i$,
we have $a\in A_i$ for some $i$. The region $A_i$ is connected, non-empty and contained in
$X$. By (b) it lies in exactly one component of $X$. This component contains $a$, and by (a)
the component of $X$ containing $a$ is $E$. Thus $A_i\subset E$.

Applying this to every member $a$ of $E$ shows that $E$ is contained in the union of those
$A_i$ with $A_i\subset E$. The reverse inclusion holds because each such $A_i$ is contained
in $E$. Hence $E$ is the union of the $A_i$ it contains. Since $E$ is non-empty by (a), it
contains at least one of them.

(f) By (a) applied to $X'$, the members of $\mathcal{C}'$ are finitely many, non-empty,
connected, pairwise disjoint as sets of cubes, and pairwise non-touching. Hence
$X=\bigcup\mathcal{C}'$ is the union of a finite family of non-empty connected regions, and
(e) applies. Each component $E$ of $X$ is the union of the members of $\mathcal{C}'$ it
contains, and it contains at least one of them, say $D$.

Suppose $E$ contained a second member $D'\neq D$ of $\mathcal{C}'$. Choose $b\in D$ and
$b'\in D'$. Since $b$ and $b'$ lie in the same component of $X$, they are joined by a chain
$(b=c_0,c_1,\dots,c_r=b')$ in $X$. Each $c_l$ is a member of $X$, and so lies in exactly one
member of $\mathcal{C}'$, because these members are pairwise disjoint and their union is $X$.
The chain starts in $D$ and ends in $D'\neq D$. Hence there is an index $l$ such that $c_l$
and $c_{l+1}$ lie in distinct members of $\mathcal{C}'$. Since $c_l$ and $c_{l+1}$ touch,
these two members of $\mathcal{C}'$ touch, which contradicts (a) for $X'$. Therefore $E$
contains only $D$, and $E=D$. In particular every component of $X$ is a member of
$\mathcal{C}'$.

Conversely, let $D\in\mathcal{C}'$. It is a non-empty connected region contained in $X$, so
by (b) it is contained in a component $E$ of $X$. By the preceding paragraph $E$ is itself a
member of $\mathcal{C}'$. So $D$ and $E$ are both members of $\mathcal{C}'$ with $D\subset E$;
since distinct members of $\mathcal{C}'$ are disjoint as sets of cubes and $D$ is non-empty,
$D=E$. Thus every member of $\mathcal{C}'$ is a component of $X$.
\end{proof}

\begin{lemma}[Enlargement and hull preserve connectedness]\label{4.15:lem-enlargement-connected}
Fix a step $t-1\to t$ and let $\mathcal{P}_t$ be its coarse partition. Enlargement
$X\mapsto X^{\sim}$, its iterates $X^{\sim n}$ and the hull $Y\mapsto\widehat{Y}$ are those of
Definition~\ref{4.15:def-enlargement-hull}.
\begin{enumerate}
\item[(a)] For $\mathcal{P}_t$-regions $X$, $X'$ and integers $m,n\ge 0$:
\begin{itemize}
\item $X\subset X^{\sim}$;
\item $(X\cup X')^{\sim}=X^{\sim}\cup X'^{\sim}$;
\item $(X^{\sim m})^{\sim n}=X^{\sim(m+n)}$;
\item $X^{\sim n}=\bigcup_c c^{\sim n}$, the union over the $\mathcal{P}_t$-cubes $c$ of $X$.
\end{itemize}
\item[(b)] If $X$ is a connected $\mathcal{P}_t$-region, then so is $X^{\sim n}$ for every
$n\ge 0$. In particular, $c^{\sim n}$ is connected for every $\mathcal{P}_t$-cube $c$.
\item[(c)] Let $Y$ be a region. Then $Y\subset\widehat{Y}$, and $\widehat{Y}=\bigcup_D\widehat{D}$,
the union over the components $D$ of $Y$. For a non-empty connected region $D$, the sets
$\widehat{D}$ and $\widehat{D}^{\sim n}$ $(n\ge 0)$ are connected $\mathcal{P}_t$-regions
containing $D$.
\end{enumerate}
\end{lemma}

\begin{proof}
(a) A cube touches itself. Hence every $\mathcal{P}_t$-cube of $X$ touches a $\mathcal{P}_t$-cube
of $X$, and is therefore a $\mathcal{P}_t$-cube of $X^{\sim}$; thus $X\subset X^{\sim}$.

The $\mathcal{P}_t$-cubes of $X\cup X'$ are those of $X$ together with those of $X'$. A
$\mathcal{P}_t$-cube therefore touches a cube of $X\cup X'$ if and only if it touches a cube of
$X$ or a cube of $X'$. This is the asserted distributivity
$(X\cup X')^{\sim}=X^{\sim}\cup X'^{\sim}$. By induction on the number of regions, the same
identity holds for every finite union of $\mathcal{P}_t$-regions.

The third relation unwinds the recursive definition $X^{\sim 0}=X$,
$X^{\sim(n+1)}=(X^{\sim n})^{\sim}$. It holds for $n=0$. If it holds for $n$, then
\begin{equation*}
(X^{\sim m})^{\sim(n+1)}=\bigl((X^{\sim m})^{\sim n}\bigr)^{\sim}
=\bigl(X^{\sim(m+n)}\bigr)^{\sim}=X^{\sim(m+n+1)}.
\end{equation*}

The fourth relation follows from distributivity by induction on $n$. For $n=0$ it says that
$X$ is the union of its $\mathcal{P}_t$-cubes $c$. If it holds for $n$, then distributivity over
the finite family $\{c^{\sim n}\}$ gives
\begin{equation*}
X^{\sim(n+1)}=\Bigl(\bigcup_c c^{\sim n}\Bigr)^{\sim}=\bigcup_c\bigl(c^{\sim n}\bigr)^{\sim}
=\bigcup_c c^{\sim(n+1)}.
\end{equation*}

(b) By definition $X^{\sim}=X\cup\bigcup c'$, the union over the $\mathcal{P}_t$-cubes $c'$ that
touch a $\mathcal{P}_t$-cube of $X$. This family is finite, since there are finitely many
$\mathcal{P}_t$-cubes. Each $c'$ is connected, because a single cube of a partition is connected
(Lemma~\ref{4.15:lem-subdivision-independence}), and each $c'$ touches $X$. Since $X$ is
connected, Lemma~\ref{4.15:lem-component-properties}(d), applied with $A=X$ and the $A_i$ the cubes
$c'$, shows that $X^{\sim}$ is connected; it is a $\mathcal{P}_t$-region by definition.

We now induct on $n$. The case $n=0$ is the hypothesis. If $X^{\sim n}$ is a connected
$\mathcal{P}_t$-region, the case just proved, applied to $X^{\sim n}$, shows that
$X^{\sim(n+1)}=(X^{\sim n})^{\sim}$ is one as well. For the last assertion, a
$\mathcal{P}_t$-cube $c$ is a connected $\mathcal{P}_t$-region by
Lemma~\ref{4.15:lem-subdivision-independence}, so the statement applies to $X=c$.

(c) Let $a$ be a member of $Y$. Since the finest partition $\mathcal{P}_{\bullet}$ refines
$\mathcal{P}_t$, $a$ lies in exactly one $\mathcal{P}_t$-cube $c$. This $c$ meets $Y$, because it
contains $a$. Hence $c$ is a $\mathcal{P}_t$-cube of $\widehat{Y}$, and $a\in\widehat{Y}$. Thus
$Y\subset\widehat{Y}$.

By Lemma~\ref{4.15:lem-component-properties}(a), $Y$ is the union of its components.
Consequently a $\mathcal{P}_t$-cube meets $Y$ if and only if it meets some component $D$ of $Y$.
Since $\widehat{Y}$ and each $\widehat{D}$ are the unions of the $\mathcal{P}_t$-cubes meeting
$Y$, respectively $D$, this gives $\widehat{Y}=\bigcup_D\widehat{D}$.

Let now $D$ be a non-empty connected region. By the first assertion of (c), $D\subset\widehat{D}$,
so $\widehat{D}=D\cup\bigcup c$, the union over the $\mathcal{P}_t$-cubes $c$ that meet $D$. This
is a finite family. Each $c$ is connected by Lemma~\ref{4.15:lem-subdivision-independence}, and
each touches $D$, since it meets $D$. Lemma~\ref{4.15:lem-component-properties}(d), applied with
$A=D$, therefore shows that $\widehat{D}$ is connected; it is a $\mathcal{P}_t$-region by definition.
Part (b), applied to $X=\widehat{D}$, shows that $\widehat{D}^{\sim n}$ is a connected
$\mathcal{P}_t$-region for every $n\ge 0$.

Finally, $D\subset\widehat{D}$ was shown above. Applying the first relation of (a) $n$ times gives
$\widehat{D}\subset\widehat{D}^{\sim n}$. Hence $D\subset\widehat{D}\subset\widehat{D}^{\sim n}$.
\end{proof}

\begin{lemma}[Boxing and components under replacement maps]\label{4.15:lem-replacement-maps}
Fix a step $t$ and its coarse partition $\mathcal{P}_t$ of the torus
$\mathbb{T}=\mathbb{R}^d/p\mathbb{Z}^d$, whose cubes have side $W:=W_t$. Let $\pi$ be the
covering map $\mathbb{R}^d\to\mathbb{T}$ and put $N_t:=p/W$, a positive integer. For
$n\in\mathbb{Z}^d$ put
\begin{equation*}
c_n:=\pi\Bigl(\prod_{\nu}\,[\,Wn_\nu,\;W(n_\nu+1)\,]\Bigr),
\end{equation*}
so that $c_n=c_{n'}$ if and only if $n\equiv n'\pmod{N_t}$ coordinatewise, and the
$\mathcal{P}_t$-cubes are indexed by $(\mathbb{Z}/N_t\mathbb{Z})^d$. The coordinate ranges
$S_\nu(C)$, the region $\operatorname{box}(C)$, the maps $\mathsf{M}_1$ and $\mathsf{M}$, the
number $q_0(X)$ and the notion of a replacement map are those of
Definition~\ref{4.15:def-enlargement-hull}.
\begin{itemize}
\item[(a)] The cubes $c_n$ and $c_{n'}$ touch if and only if for every $\nu$ there is
$\varepsilon_\nu\in\{-1,0,1\}$ with $n'_\nu\equiv n_\nu+\varepsilon_\nu\pmod{N_t}$.
\item[(b)] Call a non-empty subset $A$ of $\mathbb{Z}/N_t\mathbb{Z}$ \emph{linked} if any two of
its elements are the ends of a finite sequence in $A$ whose consecutive terms differ by $0$,
$1$ or $-1$ modulo $N_t$. A linked set is a cyclic interval: there are $a\in\mathbb{Z}$ and
$\ell\in\{0,\dots,N_t-1\}$ with
\begin{equation}\label{4.15:eq-replacement-maps-1}
A=\{\,a+i \ (\mathrm{mod}\ N_t) \;:\; 0\le i\le \ell\,\}.
\end{equation}
\item[(c)] For a non-empty connected $\mathcal{P}_t$-region $C$ each coordinate range
$S_\nu(C)$ is linked, hence a cyclic interval, and $\operatorname{box}(C)$ is a connected
$\mathcal{P}_t$-region containing $C$.
\item[(d)] If $\Phi$ is a replacement map, then for every $\mathcal{P}_t$-region $X$ the set
$\Phi X\supseteq X$ is a $\mathcal{P}_t$-region and every component of $\Phi X$ contains a
component of $X$. The same holds for every composite of finitely many replacement maps.
\item[(e)] $\mathsf{M}_1$ is a replacement map and $\mathsf{M}X=\mathsf{M}_1^{q_0(X)}X$. Hence
$\mathsf{M}X\supseteq X$ is a $\mathcal{P}_t$-region, $\mathsf{M}(\mathsf{M}X)=\mathsf{M}X$, and
every component of $\mathsf{M}X$ contains a component of $X$.
\end{itemize}
\end{lemma}

\begin{proof}
\emph{(a)} The point sets of $c_n$ and $c_{n'}$ share a point if and only if there are
$y\in\prod_\nu[Wn_\nu,W(n_\nu+1)]$ and $y'\in\prod_\nu[Wn'_\nu,W(n'_\nu+1)]$ with
$y-y'\in p\mathbb{Z}^d$, that is, if and only if for every $\nu$ there are
$y_\nu\in[Wn_\nu,W(n_\nu+1)]$, $y'_\nu\in[Wn'_\nu,W(n'_\nu+1)]$ and $m\in\mathbb{Z}$ with
\begin{equation*}
y_\nu-y'_\nu=pm=N_tWm .
\end{equation*}
For fixed $\nu$ and $m$ this asks that the intervals $[Wn_\nu,W(n_\nu+1)]$ and
$[W(n'_\nu+N_tm),\,W(n'_\nu+N_tm+1)]$ share a point. Two closed intervals $[Wa,W(a+1)]$ and
$[Wb,W(b+1)]$ with integers $a,b$ share a point if and only if $|a-b|\le1$. Hence the
condition is that $n_\nu-n'_\nu-N_tm\in\{-1,0,1\}$ for some $m\in\mathbb{Z}$, that is,
$n'_\nu\equiv n_\nu+\varepsilon_\nu\pmod{N_t}$ with $\varepsilon_\nu\in\{-1,0,1\}$, for every
$\nu$.

\emph{(b)} If $A=\mathbb{Z}/N_t\mathbb{Z}$, take $a=0$ and $\ell=N_t-1$. Otherwise pick a residue
$u\notin A$ and identify $(\mathbb{Z}/N_t\mathbb{Z})\setminus\{u\}$ with the set of integers
$\{u+1,\dots,u+N_t-1\}$ by sending the residue of $u+i$, $1\le i\le N_t-1$, to the integer
$u+i$. Under this identification two elements of $A$ that differ by $\pm1$ modulo $N_t$ differ
by $\pm1$ as integers: a step between $u+N_t-1$ and $u+N_t\equiv u$, or between $u+1$ and $u$,
would involve the residue $u$, which is not in $A$. (Equivalently: two integers of
$\{u+1,\dots,u+N_t-1\}$ differ by at most $N_t-2$ in absolute value, and the only integers of
absolute value at most $N_t-2$ congruent to $\pm1$ modulo $N_t$ are $\pm1$.) So the image $A'$
of $A$ is a set of integers in which any two elements are joined by a sequence in $A'$ with
steps of size at most $1$. Put $a:=\min A'$ and $b:=\max A'$. A sequence in $A'$ from $a$ to
$b$ with steps of size at most $1$ passes through every integer between $a$ and $b$, so
$A'=\{a,a+1,\dots,b\}$. Returning to residues, $A$ has the form
\eqref{4.15:eq-replacement-maps-1} with $\ell=b-a\le N_t-2$.

\emph{(c)} Since $C$ is connected, any two $\mathcal{P}_t$-cubes of $C$ are joined by a chain
of $\mathcal{P}_t$-cubes of $C$, that is, a finite sequence of such cubes with consecutive
members touching (Lemma~\ref{4.15:lem-subdivision-independence}). By (a), the $\nu$-th
residues of consecutive cubes of such a chain differ by $0$ or $\pm1$ modulo $N_t$. Hence
$S_\nu(C)$ is linked, and by (b) there are $a_\nu\in\mathbb{Z}$ and
$\ell_\nu\in\{0,\dots,N_t-1\}$ with
\begin{equation*}
S_\nu(C)=\{\,a_\nu+i \ (\mathrm{mod}\ N_t) \;:\; 0\le i\le\ell_\nu\,\}.
\end{equation*}
By definition, $\operatorname{box}(C)$ is the $\mathcal{P}_t$-region formed by the cubes $c_n$
such that, for every $\nu$, $n_\nu\equiv a_\nu+i_\nu\pmod{N_t}$ for some
$i_\nu\in\{0,\dots,\ell_\nu\}$. It contains $C$, since the residues of a cube of $C$ lie in
the coordinate ranges. It is connected: two index vectors $i,i'\in\prod_\nu\{0,\dots,\ell_\nu\}$
are joined by a path that changes one coordinate at a time by $\pm1$ within its range, and by
(a), applied with $\varepsilon_\nu=\pm1$ in the changed coordinate and $\varepsilon_\nu=0$ in
the others, consecutive cubes along such a path touch. Thus any two cubes of
$\operatorname{box}(C)$ are joined by a chain in $\operatorname{box}(C)$, and
$\operatorname{box}(C)$ is connected by Lemma~\ref{4.15:lem-subdivision-independence}.

\emph{(d)} By the definition of a replacement map there are components $C_1,\dots,C_m$ of $X$
($m\ge0$) and connected $\mathcal{P}_t$-regions $B_1\supseteq C_1,\dots,B_m\supseteq C_m$ with
\begin{equation*}
\Phi X=X\cup B_1\cup\dots\cup B_m .
\end{equation*}
This is a $\mathcal{P}_t$-region containing $X$. It is the union of the components of $X$,
which are non-empty and connected by
Lemma~\ref{4.15:lem-component-properties}\,(a), and of the sets $B_1,\dots,B_m$, which are
non-empty and connected (each contains the non-empty set $C_i$). By
Lemma~\ref{4.15:lem-component-properties}\,(e), applied to $\Phi X$ and this finite family,
every component of $\Phi X$ contains one of these sets; each of them contains a component of
$X$, a component of $X$ trivially and $B_i$ because $B_i\supseteq C_i$. This proves the claim
for a single replacement map.

For a composite, induct on the number of factors, using
Lemma~\ref{4.15:lem-component-properties}\,(c) to pass a component of $X$ through the
intermediate regions. Let $\Phi'$ be a composite of fewer factors for which the claim holds,
and $\Phi$ a replacement map. Then $\Phi'X\supseteq X$ is a $\mathcal{P}_t$-region, so by the
case of one map $\Phi(\Phi'X)\supseteq\Phi'X\supseteq X$ is a $\mathcal{P}_t$-region. If every
component of $\Phi'X$ contains a component of $X$, and every component of $\Phi(\Phi'X)$
contains a component of $\Phi'X$, then every component of $\Phi(\Phi'X)$ contains a component
of $X$.

\emph{(e)} For a $\mathcal{P}_t$-region $X$,
\begin{equation*}
\mathsf{M}_1X=X\cup\bigcup\{\operatorname{box}(C) : C \text{ a small component of } X\},
\end{equation*}
and each such $C$ is a non-empty connected $\mathcal{P}_t$-region
(Lemma~\ref{4.15:lem-component-properties}\,(a)), so that by (c) each $\operatorname{box}(C)$
is a connected $\mathcal{P}_t$-region containing $C$. Taking for $C_1,\dots,C_m$ the small
components of $X$ and $B_i=\operatorname{box}(C_i)$, this shows that $\mathsf{M}_1$ is a
replacement map. By Definition~\ref{4.15:def-enlargement-hull},
$\mathsf{M}X=\mathsf{M}_1^{q_0(X)}X$, a composite of $q_0(X)$ replacement maps (for
$q_0(X)=0$ it is the identity, which is the replacement map with $m=0$); (d) gives
$\mathsf{M}X\supseteq X$, that $\mathsf{M}X$ is a $\mathcal{P}_t$-region, and that every
component of $\mathsf{M}X$ contains a component of $X$. Finally, by the definition of $q_0(X)$
as the least $q$ with $\mathsf{M}_1^{q+1}X=\mathsf{M}_1^qX$, one has
$\mathsf{M}_1(\mathsf{M}X)=\mathsf{M}X$; hence $q_0(\mathsf{M}X)=0$ and
$\mathsf{M}(\mathsf{M}X)=\mathsf{M}X$.
\end{proof}

No bound on $q_0$, on the number of components or on the size of any region is used in this
proof, and no relation between $W_t$ and the period $p$; in particular the fact that cubes of
side $100M\check{R}_t$ fit in the torus, which rests on the full bounds on the coupling flow
and on hypothesis (H5), is not invoked.

\begin{proposition}[Components of the large-field region after one step]\label{4.15:prop-one-step-components}
Let $1\le t\le K$, and consider the step $t-1\to t$ of the operation $\mathbf{T}$ applied to a
term whose large-field region of scale $t-1$, as read by the step, is the region $Y$; here
$Y=\emptyset$ if $t=1$. Let $\sigma'$ be one of the terms produced, with
$P'$, $Q'$, $R'$, $\mathsf{K}$, $I_n$, $\mathsf{R}^{\mathrm{lo}}_n$, $\mathsf{R}_n$, $Z^{[n]}$,
$n_*$, $\Lambda_t$ and $\mathcal{W}^{\wedge}_t$ as in Definition~\ref{4.15:def-creation-set}.
Assume, as hypothesis, the description of the small-field complement in
Definition~\ref{4.15:def-small-field-complement}: for this step and this term, $\Omega_t^c$ is
a $\mathcal{P}_t$-region of the form \eqref{4.15:eq-small-field-complement-c}, concretely of the
form \eqref{4.15:eq-small-field-complement-a}, or at $t=1$ of the form
\eqref{4.15:eq-small-field-complement-b}. Then:
\begin{enumerate}
\item[(a)] $Y\subset\Lambda_t^c$ and $\mathcal{W}^{\wedge}_t\subset\Lambda_t^c$;
\item[(b)] every component of $\Lambda_t^c$ contains a component of $Y$ or a
$\mathcal{P}_t$-cube contained in $\mathcal{W}^{\wedge}_t$;
\item[(c)] every component of $Y$ is contained in exactly one component of $\Lambda_t^c$.
\end{enumerate}
\end{proposition}

\begin{proof}
Recall from Definition~\ref{4.15:def-creation-set} that for a set $\mathsf{R}$ of
$\mathcal{P}_t$-cubes we write $Z[\mathsf{R}]=(\Omega_t^c)^{\sim 2}\cup\mathsf{R}^{\sim}$, where
$\mathsf{R}^{\sim}$ is the enlargement of the union of the cubes of $\mathsf{R}$. By the
hypothesis \eqref{4.15:eq-small-field-complement-c} and by
Lemma~\ref{4.15:lem-enlargement-connected}(a), for every set $\mathsf{R}$ of
$\mathcal{P}_t$-cubes
\begin{equation}\label{4.15:eq-one-step-components-1}
\begin{split}
Z[\mathsf{R}]
&=(\Omega_t^c)^{\sim 2}\cup\mathsf{R}^{\sim}
=\bigcup_{i\in\mathsf{I}}G_i^{\sim(a_i+2)}\cup\widehat{Y}^{\sim(a_Y+2)}
\cup\mathsf{R}^{\sim 1}\\
&=\bigcup_{(i,c)}c^{\sim(a_i+2)}\cup\bigcup_{D}\widehat{D}^{\sim(a_Y+2)}
\cup\bigcup_{c'\in\mathsf{R}}c'^{\sim 1}.
\end{split}
\end{equation}
The second equality uses the distributivity
$(X\cup X')^{\sim}=X^{\sim}\cup X'^{\sim}$ over the finite union
\eqref{4.15:eq-small-field-complement-c}, applied twice, together with
$(X^{\sim m})^{\sim n}=X^{\sim(m+n)}$, both from
Lemma~\ref{4.15:lem-enlargement-connected}(a). In the third member, $(i,c)$ runs over the pairs
with $i\in\mathsf{I}$ and $c$ a $\mathcal{P}_t$-cube of $G_i$, so that $c$ is a
$\mathcal{P}_t$-cube of $P'\cup Q'$ because $G_i\subset P'\cup Q'$; this uses
$X^{\sim n}=\bigcup_c c^{\sim n}$ over the $\mathcal{P}_t$-cubes $c$ of $X$
(Lemma~\ref{4.15:lem-enlargement-connected}(a)). Next, $D$ runs over the components of $Y$: by
Lemma~\ref{4.15:lem-enlargement-connected}(c), $\widehat{Y}=\bigcup_D\widehat{D}$, and then
$\widehat{Y}^{\sim n}=\bigcup_D\widehat{D}^{\sim n}$ by distributivity. Finally, $c'$ runs over
the members of $\mathsf{R}$, again by $X^{\sim n}=\bigcup_c c^{\sim n}$ with $n=1$.

We call the sets on the right of \eqref{4.15:eq-one-step-components-1} the \emph{pieces} of
$Z[\mathsf{R}]$. They form a finite family: $\mathsf{I}$ is finite, there are finitely many
$\mathcal{P}_t$-cubes, and $Y$ has finitely many components by
Lemma~\ref{4.15:lem-component-properties}(a). Each piece is a non-empty connected
$\mathcal{P}_t$-region: for $c^{\sim(a_i+2)}$ and $c'^{\sim 1}$ this is
Lemma~\ref{4.15:lem-enlargement-connected}(b), and for $\widehat{D}^{\sim(a_Y+2)}$ it is
Lemma~\ref{4.15:lem-enlargement-connected}(c), since $D$ is a non-empty connected region by
Lemma~\ref{4.15:lem-component-properties}(a). Moreover the pieces contain, respectively, the cube
$c$ of $P'\cup Q'$ (because $X\subset X^{\sim}$, iterated), the component $D$ of $Y$ (because
$D\subset\widehat{D}\subset\widehat{D}^{\sim(a_Y+2)}$, by
Lemma~\ref{4.15:lem-enlargement-connected}(c) and (a)), and the cube $c'$ of $\mathsf{R}$. At
$t=1$ there are no pieces of the second kind, since $Y=\emptyset$.

For $0\le n\le n_*$ let $(\mathrm{C}_n)$ be the assertion: every component of $Z^{[n]}$
contains a component of $Y$ or a $\mathcal{P}_t$-cube contained in
$P'\cup Q'\cup\mathsf{R}_n$. We prove $(\mathrm{C}_n)$ by induction on $n$.

For $n=0$: $Z[\mathsf{R}_0]$ is the union of the finite family of its pieces, which are
non-empty connected regions, so by Lemma~\ref{4.15:lem-component-properties}(e) every component
of $Z[\mathsf{R}_0]$ contains a piece, hence, by what was just shown, a component of $Y$ or a
cube of $P'\cup Q'\cup\mathsf{R}_0$. By Lemma~\ref{4.15:lem-replacement-maps}(e), every
component of $Z^{[0]}=\mathsf{M}Z[\mathsf{R}_0]$ contains a component of $Z[\mathsf{R}_0]$, and
therefore contains a component of $Y$ or a cube of $P'\cup Q'\cup\mathsf{R}_0$. This is
$(\mathrm{C}_0)$.

From $n$ to $n+1\le n_*$: the region $Z^{[n]}\cup Z[\mathsf{R}_{n+1}]$ is the union of the
finite family consisting of the components of $Z^{[n]}$ and of the pieces of
$Z[\mathsf{R}_{n+1}]$. Each component of $Z^{[n]}$ is a non-empty connected region
(Lemma~\ref{4.15:lem-component-properties}(a)) and contains, by $(\mathrm{C}_n)$, a component of
$Y$ or a cube of $P'\cup Q'\cup\mathsf{R}_n$, which is a cube of $P'\cup Q'\cup\mathsf{R}_{n+1}$
because $\mathsf{R}_n\subset\mathsf{R}_{n+1}$ by \eqref{4.15:eq-creation-set-b}. Each piece of
$Z[\mathsf{R}_{n+1}]$ is a non-empty connected region containing a component of $Y$ or a cube of
$P'\cup Q'\cup\mathsf{R}_{n+1}$. By Lemma~\ref{4.15:lem-component-properties}(e), every
component of $Z^{[n]}\cup Z[\mathsf{R}_{n+1}]$ contains a member of this family, hence a
component of $Y$ or a cube of $P'\cup Q'\cup\mathsf{R}_{n+1}$. By
Lemma~\ref{4.15:lem-replacement-maps}(e), every component of
$Z^{[n+1]}=\mathsf{M}(Z^{[n]}\cup Z[\mathsf{R}_{n+1}])$ contains a component of
$Z^{[n]}\cup Z[\mathsf{R}_{n+1}]$, and so has the same property. This is $(\mathrm{C}_{n+1})$.

Thus $(\mathrm{C}_{n_*})$ holds. Since $\Lambda_t^c=Z^{[n_*]}$ (Definition~\ref{4.15:def-creation-set})
and $P'\cup Q'\cup\mathsf{R}_{n_*}=\mathcal{W}^{\wedge}_t$ by \eqref{4.15:eq-creation-set-b},
$(\mathrm{C}_{n_*})$ is assertion (b).

(a) By Lemma~\ref{4.15:lem-enlargement-connected}(c) and (a), by
\eqref{4.15:eq-one-step-components-1} with $\mathsf{R}=\mathsf{R}_0$, and by
Lemma~\ref{4.15:lem-replacement-maps}(e), which gives $\mathsf{M}X\supseteq X$,
\begin{equation*}
Y\subset\widehat{Y}\subset\widehat{Y}^{\sim(a_Y+2)}\subset Z[\mathsf{R}_0]
\subset\mathsf{M}Z[\mathsf{R}_0]=Z^{[0]}.
\end{equation*}
For $n<n_*$, applying $\mathsf{M}X\supseteq X$ to $X=Z^{[n]}\cup Z[\mathsf{R}_{n+1}]$ gives
$Z^{[n]}\subset Z^{[n+1]}$. Hence $Y\subset Z^{[n_*]}=\Lambda_t^c$.

For the second inclusion we first show
\begin{equation*}
P'\cup Q'\subset\Omega_t^c\subset(\Omega_t^c)^{\sim 2}\subset Z[\mathsf{R}_0]
\subset Z^{[0]}\subset\Lambda_t^c.
\end{equation*}
For the first inclusion: by \eqref{4.15:eq-small-field-complement-c}, each $G_i$ satisfies
$G_i\subset G_i^{\sim a_i}\subset\Omega_t^c$, and every cube of $P'\cup Q'$ lies in some $G_i$;
indeed, in \eqref{4.15:eq-small-field-complement-a} the regions $G_i$ are $Q'$ and $P'$, and in
\eqref{4.15:eq-small-field-complement-b}, where $P'=P_0'\cup P_1'$ and $Q'=Q_1'$, they are
$Q_1'$, $P_0'$ and $P_1'$. The second
inclusion is $X\subset X^{\sim}$, iterated; the third is the definition of $Z[\mathsf{R}_0]$; the
fourth and fifth were shown in the preceding paragraph. Next,
\begin{equation*}
\mathsf{R}_{n_*}\subset\mathsf{R}_{n_*}^{\sim}\subset Z[\mathsf{R}_{n_*}]\subset Z^{[n_*]}.
\end{equation*}
The first inclusion is $X\subset X^{\sim}$ and the second the definition of
$Z[\mathsf{R}_{n_*}]$. For the third: if $n_*\ge 1$, then
\begin{equation*}
Z^{[n_*]}=\mathsf{M}(Z^{[n_*-1]}\cup Z[\mathsf{R}_{n_*}]),
\end{equation*}
which contains $Z[\mathsf{R}_{n_*}]$ by $\mathsf{M}X\supseteq X$; if $n_*=0$, then
$\mathsf{R}_0=R'$ and
$Z^{[0]}=\mathsf{M}Z[R']\supseteq Z[R']$. Combining the two chains with
\eqref{4.15:eq-creation-set-b}, we obtain
$\mathcal{W}^{\wedge}_t=P'\cup Q'\cup\mathsf{R}_{n_*}\subset\Lambda_t^c$.

(c) Apply Lemma~\ref{4.15:lem-component-properties}(c) with $X=Y$ and $X'=\Lambda_t^c$; its
hypothesis $Y\subset\Lambda_t^c$ is the first inclusion of (a).
\end{proof}

\begin{remark}
Part~(b) of Proposition~\ref{4.15:prop-one-step-components} gives more than the statement that
every component of $\Lambda_t^c$ either contains a component of the region $Y$ read by the step or
meets $P'\cup Q'\cup\mathsf{R}_{n_*}\subset\mathcal{W}^{\wedge}_t$: such a component contains a
whole $\mathcal{P}_t$-cube of $\mathcal{W}^{\wedge}_t$, and this includes the case $t=1$. The
proof uses nothing about the tests, the thresholds or the fields: only the shape of
the recursion and the form \eqref{4.15:eq-small-field-complement-c} of the small-field complement.
\end{remark}

\begin{proposition}[Creation cubes along histories without large-field actions]
\label{4.15:prop-pure-histories}
Assume the reading of the small-field complement as a union of enlargements,
\eqref{4.15:eq-small-field-complement-c}. Let $1\le k\le K$, and let $\tau$ be a term of $\rho_k$,
or a term $\sigma'_k$ of $\mathbf{T}\rho_{k-1}$, whose history in the sense of
Definition~\ref{4.15:def-histories} satisfies
\begin{equation*}
\mathcal{A}_j=\emptyset\qquad\text{for every } j \text{ with } 1\le j<k .
\end{equation*}
Let $Z_j$, $1\le j\le k$, be the large-field regions of the history, as in
Definition~\ref{4.15:def-histories}. Then for every $j$ with $1\le j\le k$, every component of $Z_j$
contains a $\mathcal{P}_s$-cube contained in $\mathcal{W}^{\wedge}_s$ for some $s$ with $1\le s\le j$.
\end{proposition}

\begin{proof}
We argue by induction on $j$. The statement to be proved for a given $j$ is: every component
of $Z_j$ contains a $\mathcal{P}_s$-cube contained in $\mathcal{W}^{\wedge}_s$ for some $s$
with $1\le s\le j$.

Let $j=1$. The step $0\to1$ of $\mathbf{T}$ reads the region $Y=Z_0^{\mathrm{post}}=\emptyset$,
which is the case $t=1$ of Proposition~\ref{4.15:prop-one-step-components}. That proposition is
proved above as an implication from the reading \eqref{4.15:eq-small-field-complement-c}, which
is assumed here. By its clause~(b), every component of $Z_1$ contains either a component of $Y$ or
a $\mathcal{P}_1$-cube contained in $\mathcal{W}^{\wedge}_1$. Since $Y=\emptyset$ has no
component, every component of $Z_1$ contains a $\mathcal{P}_1$-cube contained in
$\mathcal{W}^{\wedge}_1$. This is the statement for $j=1$, with $s=1$.

Now let $2\le j\le k$, and assume the statement for $j-1$. Since $j-1<k$, the hypothesis gives
$\mathcal{A}_{j-1}=\emptyset$. By \eqref{4.15:eq-histories-c} this implies
$\tau_{j-1}=\sigma'_{j-1}$ and $Z_{j-1}^{\mathrm{post}}=Z_{j-1}$. Hence the step $j-1\to j$
reads the region
\begin{equation*}
Y=Z_{j-1}^{\mathrm{post}}=Z_{j-1}.
\end{equation*}
Let $C$ be a component of $Z_j$. We apply clause~(b) of
Proposition~\ref{4.15:prop-one-step-components} with $t=j$ and $Y=Z_{j-1}$. It shows
that one of the following two cases occurs.

In the first case $C$ contains a component $D$ of $Z_{j-1}$. By the induction hypothesis, $D$
contains a $\mathcal{P}_s$-cube $\square\subset\mathcal{W}^{\wedge}_s$ for some $s$ with
$1\le s\le j-1$. Then $\square\subset D\subset C$, so $C$ contains a $\mathcal{P}_s$-cube
contained in $\mathcal{W}^{\wedge}_s$ with $1\le s\le j-1\le j$.

In the second case $C$ contains a $\mathcal{P}_j$-cube contained in $\mathcal{W}^{\wedge}_j$,
which is the assertion with $s=j$.

In both cases the statement holds for $j$, and this completes the induction.
\end{proof}

\begin{lemma}[The waiting rule]\label{4.15:lem-waiting-rule}
Let $1\le k\le K$. Let $\mathbf{R}$ act at level $k$ on the component $C$ of the region
$Z_k(\sigma')$ of a term $\sigma'$ of $\mathbf{T}\rho_{k-1}$, and fix a history of $\sigma'$.
Suppose that, for some $s$ with $1\le s\le k$, the component $C$ contains a whole cube $w$
of $\mathcal{W}^{\wedge}_s$ (in particular, $C$ meets $\mathcal{W}^{\wedge}_s$). Here
$\mathcal{W}^{\wedge}_s$ is the creation set of the step $s-1\to s$ of the history. Then,
under the no-creation condition in the formulation \eqref{4.15:eq-no-creation-a} as well as
in the formulation \eqref{4.15:eq-no-creation-b} of Definition~\ref{4.15:def-no-creation},
\begin{equation}\label{4.15:eq-waiting-rule-a}
k-s\ \ge\ \mathsf{N}_k,
\qquad\text{hence}\qquad
k\ \ge\ s+\mathsf{N}_k\ \ge\ 1+\mathsf{N}_k\ >\ \mathsf{N}_k .
\end{equation}
Here $\mathsf{N}_k$ is the waiting depth at level $k$.
\end{lemma}

\begin{proof}
Let $\mathcal{T}_C$ be the set of those $t\in\{1,\dots,k\}$ for which a creation at level $t$
lies inside $C$. The phrase ``lies inside'' may be given any sense ranging from ``meets'' to
``contains a whole coarse cube of''.
Containing a whole coarse cube is the strongest of these senses and meeting is the weakest.
Since $C$ contains the coarse cube $w$ of $\mathcal{W}^{\wedge}_s$, it meets
$\mathcal{W}^{\wedge}_s$, and so $s\in\mathcal{T}_C$ under every one of these senses.
Hence $\mathcal{T}_C$ is not empty. Its largest element $t_*:=\max\mathcal{T}_C$
satisfies $s\le t_*\le k$.

Since $\mathbf{R}$ acts at level $k$ on $C$, statement \eqref{4.15:eq-histories-a} of
Definition~\ref{4.15:def-histories} shows that Ba{\l}aban's conditions (i)--(iii) on the
action of $\mathbf{R}$ hold for $C$ at level $k$. In particular, the no-creation condition (ii)
holds in the formulation adopted. We treat the two
formulations in turn.

Under \eqref{4.15:eq-no-creation-a}, no $t$ with $k-\mathsf{N}_k<t\le k$ belongs to
$\mathcal{T}_C$. Now $t_*$ belongs to $\mathcal{T}_C$ and satisfies $t_*\le k$. Therefore the
inequality $k-\mathsf{N}_k<t_*$ fails, that is, $t_*\le k-\mathsf{N}_k$.

Under \eqref{4.15:eq-no-creation-b}, no $t$ with
\begin{equation*}
\max\{1,\,k-\mathsf{N}_k+1\}\ \le\ t\ \le\ k
\end{equation*}
belongs to $\mathcal{T}_C$. Now $t_*$ belongs to $\mathcal{T}_C$ and satisfies
$1\le t_*\le k$. Therefore the lower inequality fails for $t=t_*$:
\begin{equation*}
t_*\ <\ \max\{1,\,k-\mathsf{N}_k+1\}.
\end{equation*}
Since $t_*\ge1$, the maximum exceeds $1$. Hence it equals $k-\mathsf{N}_k+1$, and
$k-\mathsf{N}_k+1>1$. The displayed strict inequality between integers then reads
$t_*<k-\mathsf{N}_k+1$, that is, $t_*\le k-\mathsf{N}_k$.

In both cases, therefore, $t_*\le k-\mathsf{N}_k$. Together with $s\le t_*$ this gives
\begin{equation*}
k-s\ \ge\ k-t_*\ \ge\ \mathsf{N}_k ,
\end{equation*}
which is the first assertion of \eqref{4.15:eq-waiting-rule-a}. Rearranging gives
$k\ge s+\mathsf{N}_k$. The hypothesis $s\ge1$ then yields $s+\mathsf{N}_k\ge1+\mathsf{N}_k$,
and $1+\mathsf{N}_k>\mathsf{N}_k$ holds trivially. This proves the remaining inequalities
of \eqref{4.15:eq-waiting-rule-a}.
\end{proof}

The two formulations \eqref{4.15:eq-no-creation-a} and \eqref{4.15:eq-no-creation-b} differ
only on components inside which no creation of any level at most $k$ lies. Along a history
that consists of the operation $\mathbf{T}$ alone below the level considered, there is no
such component, by Proposition~\ref{4.15:prop-pure-histories}.

\begin{lemma}[Large-field actions occur only above the cube size]\label{4.15:lem-action-level}
Assume the following two statements.
\begin{itemize}
\item[(a)] The part of clause~(0) of Theorem~0 (Theorem~\ref{4.1:thm-clause-zero}) displayed in
\eqref{4.15:eq-cube-sizes-a}. It gives
\[
x_j=g_j^{-2}>g^{-2}\ge\gamma_H^{-2}>1\qquad(0\le j\le K),
\]
and at a level $k$ it is used only for $j\le k$.
\item[(b)] The description \eqref{4.15:eq-small-field-complement-c} of the complement $\Omega_t^c$
of the domain $\Omega_t$ of every step $t-1\to t$, $1\le t\le K$
(Definition~\ref{4.15:def-small-field-complement}).
\end{itemize}
Let $\mathbf{R}$ act on a component of a large-field region only at the level prescribed for it by
the definition of the acting level of $\mathbf{R}$, conditions~(i)--(iii) of
Definition~\ref{4.15:def-histories}.
Condition~(ii) of that definition is taken in either of the two
formulations of the no-creation condition, \eqref{4.15:eq-no-creation-a} or
\eqref{4.15:eq-no-creation-b} (Definition~\ref{4.15:def-no-creation}).

Let $1\le k\le K$, and let $\sigma'$ be a term of $\mathbf{T}\rho_{k-1}$. Write $\Lambda_k(\sigma')$
for its small-field region of scale $k$ and $Z_k(\sigma')=\Lambda_k(\sigma')^c$ for its large-field
region. Then for every component $C$ of $Z_k(\sigma')$ on which $\mathbf{R}$ acts at level $k$,
\begin{equation}\label{4.15:eq-action-level-a}
k\ge 1+\mathsf{N}_k,\qquad\text{that is,}\qquad k>\mathsf{N}_k=\check{R}_k .
\end{equation}
Equivalently, at no level $k$ with $k\le\check{R}_k$ does $\mathbf{R}$ act on any component of the
large-field region of any term.

The quantifiers are those of Theorem~0. There exists $L_0=L_0(N)$, the constant of Theorem~0, such
that the following holds for every odd $L\ge L_0$, with the auxiliary constants fixed after $L$ in
their order and the thresholds $\gamma$ and then $\gamma_H$ fixed last:
\eqref{4.15:eq-action-level-a} holds for every $g\in\,]0,\gamma_H]$, every admissible torus and
initial coupling, every $K$, every $k\in\{1,\dots,K\}$, every term of $\mathbf{T}\rho_{k-1}$ and every
component of its large-field region acted upon at level $k$. No condition on $L$, $M$, $K$, $k$ or on
a coupling beyond those of Theorem~0 is introduced.
\end{lemma}

\begin{proof}
Let $1\le k\le K$, let $\sigma'$ be a term of $\mathbf{T}\rho_{k-1}$, and let $C$ be a component of
$Z_k(\sigma')$ on which $\mathbf{R}$ acts at level $k$. Fix a history
$\tau_0,\dots,\tau_{k-1},\sigma'_k=\sigma'$ of $\sigma'$ in the sense of
Definition~\ref{4.15:def-histories}, together with its classes $\mathcal{A}_j$, $1\le j\le k$, of
components acted upon at level $j$. Since $\mathbf{R}$ acts on $C$ at level $k$, we have
$C\in\mathcal{A}_k$, and in particular $\mathcal{A}_k\neq\emptyset$. Put
\begin{equation}\label{4.15:eq-action-level-1}
k_1:=\min\bigl\{\,j\in\{1,\dots,k\}:\ \mathcal{A}_j\neq\emptyset\,\bigr\}.
\end{equation}
The set on the right contains $k$, so it is non-empty and $1\le k_1\le k$.

By the choice of $k_1$ we have $\mathcal{A}_j=\emptyset$ for $1\le j<k_1$. Thus the history contains
no action of $\mathbf{R}$ below level $k_1$. By \eqref{4.15:eq-histories-b}, if
$\mathcal{A}_j=\emptyset$ then no operation is applied to the member of the history at level $j$,
since $\mathbf{R}$ consists of actions on components. Hence below level $k_1$ only the steps
$\mathbf{T}$ have acted on the ancestors of $\sigma'$. By \eqref{4.15:eq-histories-c}, each of these
steps takes as its incoming large-field region the region $\Lambda_j^c$ produced by the preceding one.
The first step takes the empty region, in accordance with the clause $Y=\emptyset$ at $t=1$ of
\eqref{4.15:eq-small-field-complement-c}.

Choose $C_1\in\mathcal{A}_{k_1}$, and take $C_1:=C$ if $k_1=k$. Then $C_1$ is a component of
$Z_{k_1}=Z_{k_1}(\sigma'_{k_1})$ on which $\mathbf{R}$ acts at level $k_1$, where $\sigma'_{k_1}$ is
the term of $\mathbf{T}\rho_{k_1-1}$ attached to level $k_1$ of the history
(Definition~\ref{4.15:def-histories}), with $\sigma'_{k_1}=\sigma'$ if $k_1=k$. Its history is
$\tau_0,\dots,\tau_{k_1-1},\sigma'_{k_1}$, the initial segment of the fixed one, and this history
satisfies $\mathcal{A}_j=\emptyset$ for $1\le j<k_1$.

We apply Proposition~\ref{4.15:prop-pure-histories}, the statement on creation cubes along histories
without large-field actions, with $k_1$ in place of $k$ and $\sigma'_{k_1}$ in place of the term. It
is an implication from the description \eqref{4.15:eq-small-field-complement-c}, which is
hypothesis~(b). Its remaining hypotheses are $1\le k_1\le K$ and the absence of actions of
$\mathbf{R}$ below level $k_1$, and both have just been established. Taking $j=k_1$ in its
conclusion, every component of $Z_{k_1}$, in particular $C_1$, contains a $\mathcal{P}_s$-cube
contained in $\mathcal{W}^{\wedge}_s$ for some $s$ with $1\le s\le k_1$.

Next we apply the waiting rule, Lemma~\ref{4.15:lem-waiting-rule}, at level $k_1$. Its hypotheses are
the following, and each holds here.
\begin{itemize}
\item $\mathbf{R}$ acts at level $k_1$ on the component $C_1$ of $Z_{k_1}(\sigma'_{k_1})$, where
$\sigma'_{k_1}$ is a term of $\mathbf{T}\rho_{k_1-1}$ and a history of it is fixed.
\item $C_1$ contains a whole cube of $\mathcal{W}^{\wedge}_s$ for some $1\le s\le k_1$, by the
preceding paragraph.
\end{itemize}
The lemma holds under either formulation \eqref{4.15:eq-no-creation-a}, \eqref{4.15:eq-no-creation-b}
of condition~(ii). By \eqref{4.15:eq-waiting-rule-a} it gives $k_1\ge s+\mathsf{N}_{k_1}$. Since
$s\ge1$ and $\mathsf{N}_{k_1}=\check{R}_{k_1}$, this yields
\begin{equation}\label{4.15:eq-action-level-2}
k_1\ge 1+\mathsf{N}_{k_1}=1+\check{R}_{k_1}.
\end{equation}

Suppose first that $k=k_1$. Then \eqref{4.15:eq-action-level-2} is
\eqref{4.15:eq-action-level-a}.

Suppose now that $k>k_1$. We apply Lemma~\ref{4.15:lem-cube-sizes-monotone}, the statement that the
cube sizes do not increase. Its hypotheses are the following.
\begin{itemize}
\item $L\ge2$ is an integer. This is contained in hypothesis~(H2) of
Definition~\ref{2.2:not-hypotheses}.
\item $r>0$. This is part of hypothesis~(H5), in which the exponent $r$ is fixed and positive.
\item $0\le k\le K$, and $x_0,\dots,x_k$ are real numbers. This holds since $1\le k\le K$ and
$x_j=g_j^{-2}$ with $g_j>0$.
\item $x_j>1$ for every $j\le k$. This is hypothesis~(a), used only for $j\le k$.
\end{itemize}
The lemma gives $\check{R}_k\le\check{R}_j$ for every $j\le k$, in particular
$\check{R}_k\le\check{R}_{k_1}$. Combining this with \eqref{4.15:eq-action-level-2}, we obtain
\begin{equation*}
k>k_1\ge 1+\check{R}_{k_1}\ge 1+\check{R}_k .
\end{equation*}
By the same lemma $\check{R}_k$ is a power of $L$, hence an integer, and $k$, $k_1$ are integers. The
strict inequality $k>k_1$ therefore gives $k\ge k_1+1$, and so
\begin{equation*}
k\ge 2+\check{R}_k>1+\mathsf{N}_k .
\end{equation*}

In both cases $k\ge1+\mathsf{N}_k$. Since $\mathsf{N}_k=\check{R}_k$ by the definition of the waiting
depth (Definition~\ref{4.15:not-cube-sizes}), this is \eqref{4.15:eq-action-level-a}.

For the equivalent form, suppose $k\le\check{R}_k$. If some component of the large-field region of
some term of $\mathbf{T}\rho_{k-1}$ were acted upon by $\mathbf{R}$ at level $k$, then
\eqref{4.15:eq-action-level-a} would give $k>\check{R}_k$, a contradiction. Hence no component of any
term is acted upon at such a level.
\end{proof}

\begin{corollary}[No action at or below the cube size; the first action]\label{4.15:cor-first-action}
Assume the following.
\begin{itemize}
\item The part of clause~(0) of Theorem~0 (Theorem~\ref{4.1:thm-clause-zero}) that gives
the cube sizes and the waiting depth \eqref{4.15:eq-cube-sizes-a}.
\item The description of the complement of the small-field region adopted in this section,
which is the hypothesis of Proposition~\ref{4.15:prop-pure-histories}.
\item The operation $\mathbf{R}$ acts on components of large-field regions only at the levels
described in Definition~\ref{4.15:def-histories}, its no-creation condition being read either
as \eqref{4.15:eq-no-creation-a} or as \eqref{4.15:eq-no-creation-b}.
\end{itemize}
Call a history of a term \emph{pure-$\mathbf{T}$ below level $k$} if $\mathcal{A}_j=\emptyset$
for every level $j$ with $1\le j<k$. Then the following hold.
\begin{enumerate}
\item[(a)] For every level $k\ge 1$ with $k\le\check{R}_k$ and every term $\sigma'$ of
$\mathbf{T}\rho_{k-1}$, one has $\mathcal{A}_k(\sigma')=\emptyset$: the operation $\mathbf{R}$
acts on no component at that level.
\item[(b)] The history of every term is pure-$\mathbf{T}$ below the first level at which
$\mathbf{R}$ acts in it. At that level the region acted upon has been produced from the empty
region by $\mathbf{T}$-steps alone, and every component of it contains a whole coarse cube of
some creation set, that is, a $\mathcal{P}_s$-cube contained in $\mathcal{W}^{\wedge}_s$ for
some $s$ with $1\le s$ and $s$ not exceeding that level
(Proposition~\ref{4.15:prop-pure-histories}).
\item[(c)] At every level $k$ at which $\mathbf{R}$ acts, the index
$h:=k-\mathsf{N}_k$ of the preceding operation $\mathbf{T}_h$ (in the notation of
\cite{B-CMP122a}, $h=k-N$) satisfies $h\ge 1$.
\end{enumerate}
\end{corollary}

\begin{proof}
(a) This is the equivalent form of the statement that large-field actions occur only above the
cube size (Lemma~\ref{4.15:lem-action-level}, proved as an implication from hypotheses that
are among those assumed here). By \eqref{4.15:eq-action-level-a}, if $\mathbf{R}$ acts at
level $k$ on a component of $Z_k(\sigma')$, then $k>\mathsf{N}_k=\check{R}_k$. Hence, if
$k\le\check{R}_k$, no component of the large-field region of any term of
$\mathbf{T}\rho_{k-1}$ is acted upon at level $k$, which is $\mathcal{A}_k(\sigma')=\emptyset$.

(b) Let $k_1$ be the first level at which $\mathbf{R}$ acts in the history of the term, that
is, the least $j\ge 1$ with $\mathcal{A}_j\neq\emptyset$. By the minimality of $k_1$,
$\mathcal{A}_j=\emptyset$ for $1\le j<k_1$, so the history is pure-$\mathbf{T}$ below level
$k_1$. By Definition~\ref{4.15:def-histories}, at a level $j$ with $\mathcal{A}_j=\emptyset$ no
operation is applied to the term, so only $\mathbf{T}$-steps have acted on the ancestors. By
\eqref{4.15:eq-histories-c} each of these steps reads the large-field region produced by the
previous one, and the first reads $Z_0^{\mathrm{post}}=\emptyset$; thus the region acted upon at
level $k_1$ has been produced from the empty region by $\mathbf{T}$-steps alone. Finally,
Proposition~\ref{4.15:prop-pure-histories}, which is proved as an implication from the
description of the complement of the small-field region assumed here, applies to the member
of the history at level $k_1$ with its
pure-$\mathbf{T}$ history, and gives that every component of $Z_{k_1}$ contains a
$\mathcal{P}_s$-cube contained in $\mathcal{W}^{\wedge}_s$ for some $s$ with $1\le s\le k_1$.

(c) If $\mathbf{R}$ acts at level $k$, then \eqref{4.15:eq-action-level-a} gives
$k\ge 1+\mathsf{N}_k$, that is, $h=k-\mathsf{N}_k\ge 1$.
\end{proof}

\begin{remark}[What the lower bound on acting levels uses; no cycle in the induction]
\label{4.15:rem-non-circularity}
Notation is that of Definition~\ref{4.15:def-histories} and Lemma~\ref{4.15:lem-action-level};
in particular $\mathcal{A}_j$ is the class of components of $Z_j(\sigma')=\Lambda_j(\sigma')^c$
on which $\mathbf{R}$ acts at level $j$.

\smallskip
\noindent(1) \emph{What the proof uses.} The proof of Lemma~\ref{4.15:lem-action-level}
(large-field actions occur only above the cube size) uses exactly two properties of
$\mathbf{R}$.
\begin{itemize}
\item The property \eqref{4.15:eq-histories-a}: an action at level $k$ on a component
presupposes, for that component, the no-creation clause of condition (ii) of the
definition of the levels of action at level $k$.
\item The property \eqref{4.15:eq-histories-b}: at a level at which $\mathbf{R}$ acts on
no component of a term, nothing is done to the term. Hence the next $\mathbf{T}$-step
reads the region $\Lambda_j^c$ produced by the recursion.
\end{itemize}
The proof does not use any of the following:
\begin{itemize}
\item what $\mathbf{R}$ does to a component on which it acts;
\item Ba{\l}aban's preparatory decompositions of unity in \cite{B-CMP122a}, and the
components they exclude, which concern only components of a non-empty class
$\mathcal{A}_j$;
\item the new sequence of regions $Z''_j$ of \cite{B-CMP122a};
\item the division of components into two classes in \cite{B-CMP122b};
\item what the large-field region becomes after an action;
\item any heredity of creations through steps of $\mathbf{R}$;
\item conditions (i) and (iii) of the definition of the levels of action, or the second
clause of condition (ii) in the form stated in \cite{B-CMP122a}; these restrict the
actions further and play no role;
\item Ba{\l}aban's integers $N_0(k)$;
\item any bound on the coupling flow beyond \eqref{4.15:eq-cube-sizes-a};
\item the lemma on the levels of action that accompanies their definition;
\item any analytic estimate;
\item the block-law comparison lemma;
\item the operator theorem at free current;
\item any conclusion of clause (i) of Theorem~0.
\end{itemize}

\smallskip
\noindent(2) \emph{Non-circularity.} Clauses (0) and (i) of Theorem~0 are proved by one
induction on the level; clause (0) is Theorem~\ref{4.1:thm-clause-zero}. At stage $k$ of
that induction, the couplings $g_0,\dots,g_k$ are fixed first. With them the numbers
$\check{x}_0,\dots,\check{x}_k$ and $\check{R}_0,\dots,\check{R}_k$ are fixed as well.
All of these are fixed before the action of step $k$ is rewritten with the profile, and
before the $\mathbf{R}$-half of step $k$ is performed. In the same way, at the next stage
$g_{k+1}$, and then $\check{R}_{k+1}$, are fixed before the $\mathbf{R}$-half of step
$k+1$. The induction is therefore well founded.

The regime clause at level $k$ enters the induction at stage $k$. There it is used by
the operator theorem at free current, after $g_0,\dots,g_k$ have been fixed.

Now consider the proof of the inequality $k\ge 1+\mathsf{N}_k$ of
Lemma~\ref{4.15:lem-action-level}, for an action at level $k$. It invokes only the
following:
\begin{itemize}
\item the numbers $\check{R}_{k_1}$ and $\check{R}_k$, where $k_1\le k$ is the least level
of the fixed history at which $\mathbf{R}$ acts on some component, as in the proof of
Lemma~\ref{4.15:lem-action-level};
\item the inequality \eqref{4.15:eq-cube-sizes-a} for $j\le k$, which is available at
stage $k$;
\item the definitions of the $\mathbf{T}$-steps from level $0$ to level $1$, \dots, from
level $k-1$ to level $k$, all of which have already been performed at stage $k$;
\item the definition of the levels of action at levels $\le k$.
\end{itemize}
None of the following enters that proof: any statement of clause (i) of Theorem~0; the
operator theorem at free current; the representation and bound of the large-field
terms; the independence of the localised terms from the small-field region
(Corollary~\ref{4.4:cor-localised-independence}); the large-field part of the coupling
increment; the regime clause itself.

Consequently, inserting Lemma~\ref{4.15:lem-action-level} at stage $k$ of the induction
creates no cycle.
\end{remark}

\begin{remark}[The chain of inequalities $k>\mathsf{N}_k>N_0(k)$]\label{4.15:rem-level-inequalities}
Throughout, $\mathbf{R}$ acts on components of large-field regions at the levels fixed by the
definition of the level of action, whose condition~(ii) is read in either of its two forms, and
$\mathsf{N}_k=\check{R}_k$.
\begin{itemize}
\item[(a)] \emph{The regime clause in full.} Assume the hypotheses of
Lemma~\ref{4.15:lem-action-level}, and the hypotheses of part~(c) of the lemma on the cube sizes
$\check{R}_j$ and Ba{\l}aban's integer $N_0$: the flow bounds of clause~(0)
(Theorem~\ref{4.1:thm-clause-zero}), the condition on $L$ of that lemma ($L\ge 5$, $L$ odd), and
its one-sided lower bound on $g^{-2}$ (equivalently, a ceiling on $\gamma$ fixed before
$\gamma_H$). Under the three conditions just listed, part~(c) of that lemma states: for
every level $k$ with $k\ge\check{R}_k$,
the integer $N_0(k)$ of the regime condition of \cite{B-CMP122a} exists and
\begin{equation}\label{4.15:eq-level-inequalities-2}
N_0(k)\le 2+\log_L\check{R}_k<\check{R}_k .
\end{equation}
Let $k$ be a level at which $\mathbf{R}$ acts. Then $N_0(k)$ exists, and
\begin{equation}\label{4.15:eq-level-inequalities-1}
k>\mathsf{N}_k>N_0(k).
\end{equation}
Thus Ba{\l}aban's regime clause holds at every acting level. It holds as an implication from
three statements: clause~(0), the description of the complements $\Omega_t^c$ of the
small-field regions assumed in Lemma~\ref{4.15:lem-action-level}, and part~(c) of the lemma on
the cube sizes $\check{R}_j$ and Ba{\l}aban's integer $N_0$.
\item[(b)] \emph{The attachment of the level.} Suppose the level is attached differently: an
action on a component of $\Lambda_{k'+1}^c$, that is, an action at the first step $k'$ on a
component of the large-field region of a term of $\mathbf{T}\rho_{k'}$, receives the index $k'$.
Take $\mathsf{N}=\check{R}_{k'}$, and take the waiting rule in the form $k'-t\ge\check{R}_{k'}$
for creations at levels $t\le k'+1$. Then every acting index $k'$ still satisfies
$k'>\check{R}_{k'}$.
\item[(c)] \emph{Levels $k\le\check{R}_k$.} Part~(a) of Corollary~\ref{4.15:cor-first-action}
is a derived lower bound on the level of an action in terms of that level's own cube size. It
bounds nothing from above, and it imposes nothing on $K$, on a number of steps or on a coupling.
\item[(d)] \emph{Ba{\l}aban's own preparatory large fields.} The validity of
Lemma~\ref{4.15:lem-action-level} does not depend on whether the large-field functions introduced
inside a component during Ba{\l}aban's preparation are counted as creations.
\end{itemize}
\end{remark}

\begin{proof}
(a) Let $k$ be a level at which $\mathbf{R}$ acts on a component $C$ of the large-field region
$Z_k(\sigma')$ of a term $\sigma'$ of $\mathbf{T}\rho_{k-1}$. By
Lemma~\ref{4.15:lem-action-level}, inequality~\eqref{4.15:eq-action-level-a} holds under either
form of condition~(ii), that is,
\begin{equation*}
k\ge 1+\mathsf{N}_k=1+\check{R}_k .
\end{equation*}
In particular $k\ge\check{R}_k$, which is the hypothesis on the level in part~(c) of the lemma on
the cube sizes $\check{R}_j$ and Ba{\l}aban's integer $N_0$.

The remaining hypotheses of that part are the three conditions listed in~(a):
\begin{itemize}
\item the flow bounds of clause~(0), which Theorem~\ref{4.1:thm-clause-zero} supplies;
\item the condition on $L$ of that lemma, $L\ge 5$, $L$ odd;
\item its one-sided lower bound on $g^{-2}$, equivalently a ceiling on $\gamma$ fixed before
$\gamma_H$.
\end{itemize}
These hold by assumption. Part~(c) of that lemma therefore gives the
existence of $N_0(k)$, together with \eqref{4.15:eq-level-inequalities-2}.
Combining \eqref{4.15:eq-level-inequalities-2} and the preceding display with
$\mathsf{N}_k=\check{R}_k$ gives
\begin{equation*}
k>\check{R}_k=\mathsf{N}_k>N_0(k),
\end{equation*}
which is \eqref{4.15:eq-level-inequalities-1}.

The only inputs used are Lemma~\ref{4.15:lem-action-level}, which rests on clause~(0) and on the
description of the complements $\Omega_t^c$ of the small-field regions assumed in
Lemma~\ref{4.15:lem-action-level}, and part~(c) of the lemma on the cube sizes $\check{R}_j$ and
Ba{\l}aban's integer $N_0$. No description of the creation cubes along histories that contain
$\mathbf{R}$-steps is used. Part~(c) of that lemma is not re-proved here, and it is not
an input of the proof of Lemma~\ref{4.15:lem-action-level}.

(b) Under the shifted attachment, let $k'_1$ be the first acting index. The argument of
Lemma~\ref{4.15:lem-action-level} applies in the same way. Since no action precedes the index
$k'_1$, the history is pure-$\mathbf{T}$ below the level $k'_1+1$ of the region acted upon, so
Proposition~\ref{4.15:prop-pure-histories} applies: the component of $\Lambda_{k'_1+1}^c$ acted
upon at $k'_1$ contains a whole coarse cube of $\mathcal{W}^{\wedge}_t$ for some $t$ with
$1\le t\le k'_1+1$, that is, a creation at the level $t$. The waiting rule in the form
$k'-t\ge\check{R}_{k'}$, applied at $k'=k'_1$, then gives
\begin{equation*}
k'_1\ge t+\check{R}_{k'_1}\ge 1+\check{R}_{k'_1}.
\end{equation*}
For a later acting index $k'>k'_1$ we have $\check{R}_{k'}\le\check{R}_{k'_1}$, since
$\check{x}_j=\min_{i\le j}x_i$ is non-increasing in $j$ and $\check{R}_j$ is the least power of
$L$ not smaller than $(\log\check{x}_j)^{r}$. Together with
the preceding display this yields
\begin{equation*}
\check{R}_{k'}\le\check{R}_{k'_1}<k'_1<k'.
\end{equation*}
Hence every acting index $k'$ satisfies $k'>\check{R}_{k'}$ under this attachment as well.

The attachment used in this section is the one of \cite{B-CMP122a}: the level of an action is the
scale $j$ of the region $\Lambda_j^c$ acted upon. It is also the attachment in which the regime
clause is applied, $N_0(k)$ entering its applications through
$\check{R}_{k-N_0(k)+1}$.

(c) By part~(a) of Corollary~\ref{4.15:cor-first-action}, $\mathcal{A}_k=\emptyset$ whenever
$1\le k\le\check{R}_k$; this constrains only the acting level in terms of its own cube size, and
no other quantity occurs in it.

(d) Ba{\l}aban regards the large-field functions introduced inside a component during the
preparation as new large fields for the purpose of condition~(ii), so that they do not satisfy
condition~(ii) \cite{B-CMP122a}. The creation sets $\mathcal{W}^{\wedge}_t$ of
\eqref{4.15:eq-creation-set-a}, by contrast, are attached to $\mathbf{T}$-steps.

Lemma~\ref{4.15:lem-action-level} rests on the waiting rule, Lemma~\ref{4.15:lem-waiting-rule}.
Of the creation sets, the proof of the waiting rule uses only that an action presupposes
condition~(ii) and that condition~(ii) bars creations of the sets $\mathcal{W}^{\wedge}_t$ in its
window; both remain true if further sets are counted as creations. Hence counting the preparatory
large fields as creations, or not counting them, leaves Lemma~\ref{4.15:lem-action-level}
unaffected.
\end{proof}

\begin{definition}[The top large-field region after the actions of a level {\cite[p.~390]{B-CMP122b}}]
\label{4.15:def-post-action-region}
Let $1\le k\le K$, let $\sigma'$ be a term of $\mathbf{T}\rho_{k-1}$, and write
$Z_k(\sigma')=\Lambda_k(\sigma')^c$ for its large-field region of scale $k$. Let
$\mathcal{A}_k(\sigma')$ be the class of components of $Z_k(\sigma')$ on which $\mathbf{R}$ acts
at level $k$, as in Definition~\ref{4.15:def-histories}, and let
\begin{equation*}
\mathcal{Z}=\bigcup_{X\in\mathcal{A}_k(\sigma')}X
\end{equation*}
be the union of this class; it is the set written $Z$ in \cite[pp.~176--177]{B-CMP122a} before any
component is removed from it during the preparation described below, and we count the components
so removed as members of $\mathcal{A}_k(\sigma')$. Let $\tau$
be one of the terms produced from $\sigma'$ by the actions of level $k$, with $\tau=\sigma'$ if
$\mathcal{A}_k(\sigma')=\emptyset$. The \emph{top large-field region after the actions of level $k$}
is the set
\begin{equation*}
Z_k^{\mathrm{post}}(\tau)=\Lambda_k(\tau)^c .
\end{equation*}
It is the large-field region of scale $k$ that the step $k\to k+1$ applied to $\tau$ reads
\cite[p.~265]{B-CMP119}. For $X\in\mathcal{A}_k(\sigma')$ the set $S(X)$ defined in the following
display is the \emph{residual} of $X$, and
\begin{equation}\label{4.15:eq-post-action-region-a}
\Lambda_k(\tau)^c=\bigl(Z_k(\sigma')\setminus\mathcal{Z}\bigr)\cup
\bigcup_{X\in\mathcal{A}_k(\sigma')}S(X),
\qquad S(X):=\Lambda_k(\tau)^c\cap X\subset X .
\end{equation}
To see this, recall that by \cite[(1.1)]{B-CMP122a} the integral operation $T_k(Z_k(\sigma'))$
over the large-field region, $T_k$ denoting Ba{\l}aban's integral operation of level $k$,
factorises into two parts:
\begin{itemize}
\item the operation on $Z_k(\sigma')\cap\mathcal{Z}^c$;
\item the product of the operations on the disjoint components of $\mathcal{Z}$.
\end{itemize}
Ba{\l}aban states there that the operation on $Z_k(\sigma')\cap\mathcal{Z}^c$ is not changed by the
constructions of that section. Hence outside $\mathcal{Z}$ the region is unchanged:
\begin{equation*}
\Lambda_k(\tau)^c\cap\mathcal{Z}^c=Z_k(\sigma')\setminus\mathcal{Z}.
\end{equation*}
Inside $\mathcal{Z}$ the region is, by definition, the union of the residuals, since
$\Lambda_k(\tau)^c\cap\mathcal{Z}$ is the union over $X\in\mathcal{A}_k(\sigma')$ of
$\Lambda_k(\tau)^c\cap X$. Together these give \eqref{4.15:eq-post-action-region-a}.

\smallskip
\noindent The following condition is a hypothesis on this region; it is adopted where it is used
and is not derived here. Its negation is recorded after it.
\begin{itemize}
\item[] \emph{The whole-component condition.} For every $1\le k\le K$, every term $\sigma'$ of
$\mathbf{T}\rho_{k-1}$ and every term $\tau$ produced from $\sigma'$ by the actions of level $k$,
each residual $S(X)$, $X\in\mathcal{A}_k(\sigma')$, in \eqref{4.15:eq-post-action-region-a} is
either $\emptyset$ or $X$. Equivalently, $\Lambda_k(\tau)^c$ is the union of a set of components of
$Z_k(\sigma')$, and this set contains every component not in $\mathcal{A}_k(\sigma')$.
\item[] \emph{The residual alternative.} For some $k$, $\sigma'$, $\tau$ as above and some
$X\in\mathcal{A}_k(\sigma')$, the residual $S(X)$ is a non-empty proper sub-region of $X$.
\end{itemize}

The two formulations of the whole-component condition are equivalent. By
Lemma~\ref{4.15:lem-component-properties}(a), the components of $Z_k(\sigma')$ are pairwise
disjoint and have union $Z_k(\sigma')$. Hence $Z_k(\sigma')\setminus\mathcal{Z}$ is the union of the
components of $Z_k(\sigma')$ that are not in $\mathcal{A}_k(\sigma')$.

Suppose first that each $S(X)$ is $\emptyset$ or $X$. Then the right-hand side of
\eqref{4.15:eq-post-action-region-a} is a union of two kinds of components: all components not in
$\mathcal{A}_k(\sigma')$, and those $X\in\mathcal{A}_k(\sigma')$ with $S(X)=X$.

Conversely, suppose $\Lambda_k(\tau)^c$ is a union of components of $Z_k(\sigma')$. Let
$X\in\mathcal{A}_k(\sigma')$. Since $X$ is disjoint from every other component, the set
$S(X)=\Lambda_k(\tau)^c\cap X$ is $X$ if $X$ is one of these components and $\emptyset$ otherwise.

When $\mathcal{A}_k(\sigma')=\emptyset$ we have $\mathcal{Z}=\emptyset$ and $\tau=\sigma'$. Then
\eqref{4.15:eq-post-action-region-a} reduces to \eqref{4.15:eq-histories-c}. That identity is a
consequence of the definitions, since by Definition~\ref{4.15:def-histories} no operation is
applied to $\sigma'$ at level $k$. In this case the instance of the whole-component condition at
$(k,\sigma',\tau)$ holds vacuously, and $\sigma'$ furnishes no instance of the residual
alternative.

\smallskip
\noindent\emph{The class and the sets it is compared with.} By \cite[pp.~176--177]{B-CMP122a},
each component in the class on which $\mathbf{R}$ acts at level $k$ has two properties:
\begin{itemize}
\item it is contained in a cube of the size $100MR_k$;
\item in the preceding $\mathsf{N}_k$ renormalisation steps no new large-field regions were created
inside it, and the earlier large-field regions contained in it satisfy the cube condition on their
own scales.
\end{itemize}
Here $R_k$ is Ba{\l}aban's large-field size parameter of level $k$, as it enters the cube condition
of \cite[pp.~176--177]{B-CMP122a}.

During the preparation of \cite[\S1]{B-CMP122a}, decompositions of unity are introduced on each
component of the current set written $Z$. A component on which a large-field characteristic
function is introduced no longer satisfies the second property, and it is removed from that set; it
remains a member of $\mathcal{A}_k(\sigma')$. For the components removed at the first
decomposition, Ba{\l}aban states that they are not changed by the subsequent operations and
are treated in the same way as the remaining large-field regions. For the later removals he states
that the components are removed. We write $\mathcal{Z}^{\mathrm{fin}}$ for the union of the
components of $\mathcal{Z}$ that remain after all removals; it is the set written $Z$ in the
constructions that follow, and $\mathcal{Z}\setminus\mathcal{Z}^{\mathrm{fin}}$ is the union of the
removed components.

For the construction of the new action, Ba{\l}aban defines a new sequence of large-field regions
inside $\mathcal{Z}$ \cite[(1.10)--(1.12)]{B-CMP122a}. Let $N_0$ be the smallest positive integer
such that $L^{-N_0+1}MR_{k-N_0+1}=M$; it is assumed that $\mathsf{N}_k>N_0$. Put
$h=k-\mathsf{N}_k$.
\begin{itemize}
\item The two top members $Z''_k$ and $Z''_{k-1}$ are defined in \cite[(1.10)]{B-CMP122a}.
\item These are completed to a sequence $Z''_k,Z''_{k-1},\dots,Z''_{k-N_0+1}$ whose complements form
an admissible sequence of domains based on partitions into $M$-cubes on the corresponding scales.
Thus $Z''_k$ and $Z''_k\setminus Z''_{k-1}$ are unions of $M$-cubes, $Z''_k$ and $Z''_{k-1}$ are
separated by one layer of $M$-cubes, and so on down the sequence.
\item The remaining members are
\begin{equation}\label{4.15:eq-post-action-region-b}
\begin{aligned}
Z''_j&=\bigl(\Omega_j^{\sim 5}\bigr)^c\cap\mathcal{Z},
&\qquad j&=k-N_0,\,k-N_0-1,\,\dots,\,k-\mathsf{N}_k+1=h+1,\\
Z''_j&=Z_j,&\qquad j&=h,\,h-1,\,\dots,\,1 .
\end{aligned}
\end{equation}
\end{itemize}
New domains $\Omega''_j$, $j=k,\dots,1$, are defined in \cite[(1.12)]{B-CMP122a}, with
$\Omega''_j=\Omega_j$ for $j=h,\dots,1$. The integrations of that construction are with respect to
the variables localised in $Z_{j+1}\setminus Z''_{j+1}$, $j=h,\dots,k-1$.

In \cite[(1.33)--(1.35)]{B-CMP122b} a new determining set $B_k$ is introduced. It is equal to
Ba{\l}aban's determining set $B''_k$ on $(\mathcal{Z}^{\mathrm{fin}})^c$ and to $\{Z^{(k)}\}$ on
$\mathcal{Z}^{\mathrm{fin}}$; more precisely $B_k=\{\Gamma_j\}$, with
\begin{equation*}
\Gamma_j=\Gamma''_j\cap(\mathcal{Z}^{\mathrm{fin}})^c\quad(j<k),
\qquad \Gamma_k=\Gamma''_k\cup Z^{(k)} ,
\end{equation*}
together with new gauge field variables $V$ on it and the new background field
$U_k=U_{B_k}(V)$. Decompositions of unity are then introduced on the components of
$\mathcal{Z}^{\mathrm{fin}}$, and these components are divided into two classes
\cite[(1.71)--(1.72)]{B-CMP122b}.
\begin{itemize}
\item First-class components $X_1,\dots,X_n$ are those with the small-field function of the
decomposition; on them the characteristic functions $\chi_{k,\Lambda}$, $\chi$ are removed, because
they are equal to $1$. They carry only the characteristic functions $\chi_k(X_i^{\sim-6})$, which
are combined with the function $\chi_k(\Omega_k^{\sim-4})$, $\Omega_k$ being the old $k$th domain;
these functions are gauge invariant. On these components the density is integrated over the gauge
field variables $V_k$, and the gauge-fixing $\delta$-functions $\delta_{G\cap X_i}(V'_k)$ are
removed by a reversed Faddeev--Popov procedure; this is possible only because the density is
integrated over $V_k$ at least on the components $X_i$. The density without the $\delta$-functions
connected with these components is equivalent, but not equal, to the density with them.
\item Second-class components $Y_1,\dots,Y_m$ carry the large-field characteristic functions
$\chi^c_k(Y_i^{\sim-6})$ together with the earlier functions $\chi_{k,\Lambda_i}\chi_i$ and the
$\delta$-functions $\delta_{G_i}(V'_k)$.
Ba{\l}aban states that they are included again into the large-field domain of the new action.
\end{itemize}
The letter $G$ in the subscripts of these $\delta$-functions is Ba{\l}aban's and does not denote the
gauge group. The top-integrated operation $\mathbf{T}'_k(X)$ of \cite[(1.71)]{B-CMP122b} is
defined for a component $X$ of either class, and all the expressions in it are localised in $X$.
In the definition \cite[(1.72)]{B-CMP122b} of the $\mathbf{R}$-operation, $\Omega_k$ is the new
$k$th domain defined by the determining set of $U_k$. By \cite[(1.33)]{B-CMP122b}, a first-class
component $X_i$ passes into the top member $\Gamma_k=\Gamma''_k\cup Z^{(k)}$ of that set, and so
into the new $k$th domain $\Omega_k$. At the next
step, a component of $Z_{k+1}$ that is not small is obtained by joining whole components of the
region of level $k$, as it stands after level $k$, together with new large-field regions
\cite[\S1]{B-CMP122b}. This concerns that region whatever it is; it does not by itself say what
becomes of a first-class component.

The sentence following \cite[(1.101), p.~390]{B-CMP122b} states that the new large field region is
equal to $Z_k\cup\bigcup_{i=1}^{m}Y_i$. The sentences of \cite{B-CMP122a} and \cite{B-CMP122b}
cited above are compatible with two readings of the top-level set $Z_k$ in it.
\begin{itemize}
\item[] \emph{First reading.} The top-level set is the unchanged part
$Z_k(\sigma')\cap(\mathcal{Z}^{\mathrm{fin}})^c$ of the old region. Under the convention of
\cite[pp.~176--177]{B-CMP122a}, intersections with the set written $Z$ are written explicitly only
where omitting them would lead to a misunderstanding. If the set were read as
$Z_k\cap\mathcal{Z}^{\mathrm{fin}}=\mathcal{Z}^{\mathrm{fin}}\supseteq\bigcup_iX_i$, the sentence
would return the first-class components to the large-field region, although their block fields
have been integrated over and moved into $\Gamma_k$. It would also make the term $\bigcup_iY_i$
redundant.
\item[] \emph{Second reading.} The top-level set is the top member $Z''_k$ of the new sequence, with
the unchanged part outside $\mathcal{Z}^{\mathrm{fin}}$ understood.
\end{itemize}
The set $\mathcal{Z}^{\mathrm{fin}}$ is the disjoint union of the first-class and the second-class
components. Under the two readings the new region is, respectively,
\begin{equation}\label{4.15:eq-post-action-region-c}
\begin{gathered}
\Lambda_k(\tau)^c=\bigl(Z_k(\sigma')\setminus\mathcal{Z}^{\mathrm{fin}}\bigr)
\cup\bigcup_{i=1}^{m}Y_i
=Z_k(\sigma')\setminus\bigcup_{i=1}^{n}X_i,\\
\Lambda_k(\tau)^c=\bigl(Z_k(\sigma')\setminus\mathcal{Z}^{\mathrm{fin}}\bigr)\cup Z''_k
\cup\bigcup_{i=1}^{m}Y_i .
\end{gathered}
\end{equation}
Under the first reading, the components removed during the preparation lie outside
$\mathcal{Z}^{\mathrm{fin}}$, hence inside $Z_k(\sigma')\setminus\mathcal{Z}^{\mathrm{fin}}$. So the
new region is $Z_k(\sigma')$
minus the first-class components, and $S(X)=\emptyset$ for first-class $X$ while $S(X)=X$ for every
other member of the class. The whole-component condition therefore holds under the first reading.

Under the second reading the outcome depends on $Z''_k$.
\begin{itemize}
\item If $Z''_k=\mathcal{Z}$, every component of the class stays whole, $S(X)=X$ for all
$X\in\mathcal{A}_k(\sigma')$, and the whole-component condition again holds.
\item If $Z''_k$ is a proper part of $\mathcal{Z}$, a first-class component $X_i$ leaves in the
top region exactly the residual $S(X_i)=Z''_k\cap X_i$. When $Z''_k\cap X_i$ is non-empty and
different from $X_i$, this is an instance of the residual alternative.
\end{itemize}
A residual of another shape, a shell, would arise if the characteristic functions
$\chi_k(X_i^{\sim-6})$ and $\chi^c_k(Y_i^{\sim-6})$ delimited the new regions. The sentences cited
here do not say that they do. Whether $Z''_k=\mathcal{Z}$ is not settled by these sentences.

The description of $\mathbf{R}$ on a treated component, in the outline of
Definition~\ref{1.5:def-densities}, does not fix
the set $\Lambda_k(\tau)^c$. That description consists of:
\begin{itemize}
\item the integration over the block field inside the component, the extension of the result as a
function independent of $V_k$ there, and the multiplication by the normalised reference density;
\item the top-integrated operation $\mathbf{T}'_k(X)$ of \cite[(1.71)]{B-CMP122b}, which integrates
the variables on $X$ and contains the sum over the sequences $\{\Omega_j^c\cap X,\,Z_j\cap X\}$ of
regions inside $X$;
\item the factors that keep a treated component live, that is, still carrying its large-field
factor: the gauge-fixing factor on a second-class component and the reference law on a first-class
one;
\item the division into the two classes.
\end{itemize}
All of these are consistent with each of the following outcomes:
\begin{itemize}
\item a treated component remains whole in $\Lambda_k(\tau)^c$, carrying $\mathbf{T}'_k(X)$;
\item a first-class component leaves $\Lambda_k(\tau)^c$ whole, its block field having passed under
the reference law into the top-level domain;
\item a treated component is replaced by a proper sub-region, either a core delimited by the new
sequence $Z''_j$ or a shell.
\end{itemize}

The region $\Lambda_k(\tau)^c$ of this definition is the new large-field region of the sentence
following \cite[(1.101), p.~390]{B-CMP122b}. It is built on the data of
\cite[\S1, (1.1), (1.10)--(1.12), pp.~176--179]{B-CMP122a} and of
\cite[(1.33)--(1.35), (1.71)--(1.72)]{B-CMP122b}. The whole-component condition is that sentence
under the first reading, or under the second reading with $Z''_k=\mathcal{Z}$. It enters this book
as a hypothesis. Every statement of this book that uses it lists it explicitly among its
hypotheses. The residual
alternative is the case excluded by it.
\end{definition}

\begin{proposition}[Every top large-field component contains a creation cube]
\label{4.15:prop-creation-cube}
Assume the following two hypotheses.
\begin{enumerate}
\item[(a)] The description of the small-field complement in
Definition~\ref{4.15:def-small-field-complement} holds: for every step $t-1\to t$,
$1\le t\le K$, and every term the step produces, the complement $\Omega_t^c$ of the
domain of the step is the union of enlargements displayed in
\eqref{4.15:eq-small-field-complement-c}.
\item[(b)] The description of the top large-field region after the actions of a level
in Definition~\ref{4.15:def-post-action-region} holds: for every $t$ with
$1\le t\le K$, every term $\sigma'$ of $\mathbf{T}\rho_{t-1}$ and every term $\tau'$
produced from $\sigma'$ by the actions of level $t$, every residual $S(X)$,
$X\in\mathcal{A}_t(\sigma')$, in \eqref{4.15:eq-post-action-region-a} is either empty
or equal to $X$; equivalently, $\Lambda_t(\tau')^c$ is a union of components of
$Z_t(\sigma')$ that contains every component on which $\mathbf{R}$ does not act.
\end{enumerate}
Let $k$ be a level with $1\le k\le K$, let $\tau$ be a term of $\rho_k$, and fix a
history of $\tau$ in the sense of Definition~\ref{4.15:def-histories}, that is, any
succession of steps of $\mathbf{T}$ and of actions of Ba{\l}aban's large-field
operation $\mathbf{R}$ on components, at the levels fixed in
Definition~\ref{4.15:def-histories}, that produces $\tau$. Let $\sigma'_k$ be the term
of $\mathbf{T}\rho_{k-1}$ from which $\tau$ is produced by the actions of level $k$.
Consider
\begin{itemize}
\item the region $Z_k=\Lambda_k(\sigma'_k)^c$ produced by the step $k-1\to k$, which is
the large-field region of scale $k$ before the actions of $\mathbf{R}$ at level $k$;
\item the region $Z_k^{\mathrm{post}}=\Lambda_k(\tau)^c$ after the actions of
$\mathbf{R}$ at level $k$; for $k<K$ this is the set which the step $k\to k+1$ reads as
the large-field region of scale $k$ \cite[p.~265]{B-CMP119}.
\end{itemize}
Then every component of $Z_k$ and every component of $Z_k^{\mathrm{post}}$ contains a
whole $\mathcal{P}_s$-cube contained in the creation set $\mathcal{W}^{\wedge}_s$ for
some $s$ with $1\le s\le k$, where $\mathcal{P}_s$ is the coarse partition of the step
$s-1\to s$. In particular every such component contains a point of
\begin{equation*}
\bigcup_{1\le s\le k}\mathcal{W}^{\wedge}_s .
\end{equation*}
\end{proposition}

The part of Proposition~\ref{4.15:prop-creation-cube} that concerns the levels up to
and including the first level at which $\mathbf{R}$ acts in the history, and in
particular the whole of it along histories without actions of $\mathbf{R}$, is
Proposition~\ref{4.15:prop-pure-histories}, which assumes hypothesis~(a) only.

\begin{proof}
Fix $k$ with $1\le k\le K$, a term $\tau$ of $\rho_k$ and a history
$\tau_0,\dots,\tau_k=\tau$ of $\tau$, together with its data as in
Definition~\ref{4.15:def-histories}: for $1\le j\le k$, the term $\sigma'_j$ of
$\mathbf{T}\rho_{j-1}$ produced from $\tau_{j-1}$ by the step $j-1\to j$, the region
$Z_j=\Lambda_j(\sigma'_j)^c$, the creation set $\mathcal{W}^{\wedge}_j$ of the step,
the class $\mathcal{A}_j$ of components of $Z_j$ on which $\mathbf{R}$ acts at
level $j$, and the region
\begin{equation*}
Z_j^{\mathrm{post}}=\Lambda_j(\tau_j)^c .
\end{equation*}
We prove by induction on $j$, $1\le j\le k$, the following assertion, which we call the
assertion at level $j$: every component of $Z_j$, and every component of
$Z_j^{\mathrm{post}}$, contains a $\mathcal{P}_s$-cube contained in
$\mathcal{W}^{\wedge}_s$ for some $s$ with $1\le s\le j$.

Let $1\le j\le k$, and, if $j\ge2$, assume the assertion at level $j-1$. The step
$j-1\to j$, applied to $\tau_{j-1}$, reads as its incoming large-field region the set
$Y=Z_{j-1}^{\mathrm{post}}$ (Definition~\ref{4.15:def-histories} and
\cite[p.~265]{B-CMP119}), with the convention $Z_0^{\mathrm{post}}=\emptyset$.

We first treat the components of $Z_j$. Hypothesis~(a) is the hypothesis of
Proposition~\ref{4.15:prop-one-step-components}; we apply its clause~(b) to the step
$j-1\to j$, to the term $\tau_{j-1}$ with incoming region $Y=Z_{j-1}^{\mathrm{post}}$,
and to the produced term $\sigma'_j$, whose large-field region of scale $j$ is
$Z_j=\Lambda_j(\sigma'_j)^c$. It gives: a component $C$ of $Z_j$ contains either a
component $D$ of $Z_{j-1}^{\mathrm{post}}$, or a $\mathcal{P}_j$-cube contained in
$\mathcal{W}^{\wedge}_j$.
\begin{itemize}
\item In the first case $j\ge2$, since $Z_0^{\mathrm{post}}=\emptyset$ has no
component. The induction hypothesis, the assertion at level $j-1$, applies to the
component $D$ of $Z_{j-1}^{\mathrm{post}}$ and gives a $\mathcal{P}_s$-cube $w$ with
$w\subset\mathcal{W}^{\wedge}_s$ for some $s\le j-1$. Then $w\subset D\subset C$, so
$C$ contains a $\mathcal{P}_s$-cube contained in $\mathcal{W}^{\wedge}_s$ with
$1\le s\le j-1$.
\item In the second case $C$ contains a $\mathcal{P}_s$-cube contained in
$\mathcal{W}^{\wedge}_s$ with $s=j$.
\end{itemize}
So every component of $Z_j$ contains a $\mathcal{P}_s$-cube contained in
$\mathcal{W}^{\wedge}_s$ for some $s$ with $1\le s\le j$.

We now treat the components of $Z_j^{\mathrm{post}}$. By hypothesis~(b), applied at
level $j$ to the term $\sigma'_j$ of $\mathbf{T}\rho_{j-1}$ and to the term $\tau_j$
produced from it by the actions of level $j$, the region
$Z_j^{\mathrm{post}}=\Lambda_j(\tau_j)^c$ is the union of a set $\mathcal{C}'$ of
components of $Z_j=Z_j(\sigma'_j)$. If $\mathcal{C}'=\emptyset$, then
$Z_j^{\mathrm{post}}=\emptyset$ has no component, and there is nothing to prove.
Otherwise clause~(f) of Lemma~\ref{4.15:lem-component-properties} shows that the
components of $Z_j^{\mathrm{post}}$ are exactly the members of $\mathcal{C}'$. Each of
them is a component of $Z_j$, and so, by the preceding paragraph, contains a
$\mathcal{P}_s$-cube contained in $\mathcal{W}^{\wedge}_s$ for some $s$ with
$1\le s\le j$.

This proves the assertion at level $j$, and the induction is complete. The assertion
at level $k$ is the conclusion of the proposition for $\tau$ and the fixed history:
the components before the actions of $\mathbf{R}$ at level $k$ are those of $Z_k$, and
the components after these actions are those of $Z_k^{\mathrm{post}}$. The final
sentence of the proposition follows, since a $\mathcal{P}_s$-cube contained in
$\mathcal{W}^{\wedge}_s$ is non-empty and its points lie in
$\bigcup_{1\le s\le k}\mathcal{W}^{\wedge}_s$.
\end{proof}

Thus, under hypothesis~(b), every component of the modified large-field region after
an action of $\mathbf{R}$ still contains a whole coarse cube of some creation set.
The assertion about $Z_k^{\mathrm{post}}$ is obtained only as an implication from
hypotheses~(a) and~(b): under the alternative description of the region after the
actions of a level recorded in Definition~\ref{4.15:def-post-action-region}, in which
a treated component may leave in that region a non-empty proper residual, it is not
proved.

\begin{remark}[The case of a proper residual after a large-field action]
\label{4.15:rem-proper-residual}
Let $1\le k\le K$, let $\sigma'$ be a term of $\mathbf{T}\rho_{k-1}$, let $\tau$ be one of the terms
produced from $\sigma'$ by the actions of $\mathbf{R}$ at level $k$, and write
$Z_k(\sigma')=\Lambda_k(\sigma')^c$. For a component $X\in\mathcal{A}_k(\sigma')$ of $Z_k(\sigma')$
acted upon at level $k$, let
\begin{equation*}
S(X)=\Lambda_k(\tau)^c\cap X
\end{equation*}
be its residual. Definition~\ref{4.15:def-post-action-region} describes $\Lambda_k(\tau)^c$ in the case
in which every $S(X)$ is either $\emptyset$ or $X$. Suppose instead that there are
$\sigma'$, $\tau$ and $X\in\mathcal{A}_k(\sigma')$ such that $S(X)$ is a non-empty proper sub-region
of $X$. For such $\sigma'$, $\tau$, $X$ the following hold.
\begin{enumerate}
\item[(a)] The region $\Lambda_k(\tau)^c$ is not a union of whole components of $Z_k(\sigma')$, and
clause (f) of Lemma~\ref{4.15:lem-component-properties} does not apply to it.
\item[(b)] The residual $S(X)$ need not be connected, although $X$ is connected.
\end{enumerate}
Accordingly, for the conclusion of Proposition~\ref{4.15:prop-creation-cube} to persist through the
action of $\mathbf{R}$ on $X$, two things would have to be shown: that $S(X)$ contains one of the
cubes of the partition $\mathcal{P}_s$ contained in the creation set $\mathcal{W}^{\wedge}_s$ of the
step $s-1\to s$, for some $s\le k$, and contained in $X$; and that every component of $S(X)$ contains
such a cube. In that case neither is proved, and the passage of the conclusion of
Proposition~\ref{4.15:prop-creation-cube} through the large-field actions is not proved.

The set $\Lambda_k(\tau)^c$ is determined by Ba{\l}aban's definition of the region $Z''_k$ in
\cite[(1.10), p.~179]{B-CMP122a} and by the description of the new action in
\cite[(1.101) and the sentence following it, p.~390]{B-CMP122b}. The display
\cite[(1.11), p.~179]{B-CMP122a} contains the relation
\begin{equation*}
Z''_j=(\Omega_j^{\sim 5})^c\cap\mathcal{Z}, \qquad h+1\le j\le k-N_0,
\end{equation*}
where $\mathcal{Z}$ is the union of the class of components acted upon, $h=k-\mathsf{N}_k$ and $N_0$
is Ba{\l}aban's integer of that paper; this relation concerns only the lower members $j\le k-N_0$ and
does not decide the case $j=k$. If $\Lambda_k(\tau)^c$ is a union of whole
components of $\Lambda_k(\sigma')^c$, then the description of
Definition~\ref{4.15:def-post-action-region} holds and the proof of
Proposition~\ref{4.15:prop-creation-cube} applies verbatim. If a proper residual $S(X)$ occurs, the
statement to be proved is: every component of $S(X)$ contains a $\mathcal{P}_s$-cube of
$\mathcal{W}^{\wedge}_s$ for some $s\le k$.

The statement on the creation of large-field regions by the renormalisation steps alone is unaffected
by this alternative, since no $\mathbf{R}$-step enters it.
\end{remark}

\begin{proof}
(a) Suppose, to the contrary, that $\Lambda_k(\tau)^c$ is the union of a set $\mathcal{C}'$ of
components of $Z_k(\sigma')$. The region $X$ is itself a component of $Z_k(\sigma')$. By
Lemma~\ref{4.15:lem-component-properties}(a), distinct components of $Z_k(\sigma')$ have pairwise
disjoint point sets. If $X\in\mathcal{C}'$, then $X\subset\Lambda_k(\tau)^c$, and so $S(X)=X$. If
$X\notin\mathcal{C}'$, then $X$ is disjoint from every member of $\mathcal{C}'$, hence from their
union $\Lambda_k(\tau)^c$, and so $S(X)=\emptyset$. Both alternatives contradict the assumption that
$S(X)$ is a non-empty proper sub-region of $X$. Hence $\Lambda_k(\tau)^c$ is not a union of whole
components of $Z_k(\sigma')$, and clause (f) of Lemma~\ref{4.15:lem-component-properties} does not
apply to it.

(b) The region $X$ is connected, being a component of $Z_k(\sigma')$
(Lemma~\ref{4.15:lem-component-properties}(a)). The residual $S(X)$ is a proper sub-region of $X$,
and a proper sub-region of a connected region --- for instance one obtained by removing layers of
cubes --- need not be connected. Hence no connectivity of $S(X)$ follows from that of $X$.
\end{proof}

\begin{remark}[The creation statement concerns top-scale regions only]\label{4.15:rem-top-scale-scope}
Proposition~\ref{4.15:prop-creation-cube} concerns the top large-field regions of a level only:
the region $Z_k$ before the actions of $\mathbf{R}$ at level $k$, and the region
$Z_k^{\mathrm{post}}$ after them. It claims nothing about the lower regions of a treated
component. More precisely, the following hold.
\begin{itemize}
\item[(a)] Let $X$ be a treated component, and consider its new lower regions
$Z''_j\cap X$ of \eqref{4.15:eq-post-action-region-b}. They are in general proper subregions
of the old ones. Their definition does not require them to contain the creation cubes of $X$.
Proposition~\ref{4.15:prop-creation-cube} says nothing about them.
\item[(b)] Proposition~\ref{4.15:prop-creation-cube} does not involve condition~(ii) of
\cite{B-CMP122a}, in either of the readings of it under which
Corollary~\ref{4.15:cor-first-action} is stated.
\item[(c)] Fix a history. At every level $j$ up to and including the first level at which
$\mathbf{R}$ acts, the conclusion of Proposition~\ref{4.15:prop-creation-cube} for $Z_j$
follows from the reading assumed in
Proposition~\ref{4.15:prop-pure-histories} alone. The reading of the post-action region fixed in
Definition~\ref{4.15:def-post-action-region} is used only at the first action of
$\mathbf{R}$ and after it.
\end{itemize}
\end{remark}

\begin{proof}
(a) Let $X$ be a treated component at level $k$. Let $\mathcal{Z}$ be the union of the
components acted upon, and let $h=k-\mathsf{N}_k$. For $h+1\le j\le k-N_0$ the new lower
regions are defined in \cite[(1.11)]{B-CMP122a} by
$Z''_j=(\Omega_j^{\sim 5})^c\cap\mathcal{Z}$. Since $\Omega_j\subset\Omega_j^{\sim 5}$,
passing to complements reverses the inclusion. Moreover, in the admissible sequence of
domains the region $\Omega_j^c$ is contained in the large-field region $Z_j=\Lambda_j^c$.
Together these give
\begin{equation*}
Z''_j=(\Omega_j^{\sim 5})^c\cap\mathcal{Z}
\subset\Omega_j^c\cap\mathcal{Z}
\subset Z_j\cap\mathcal{Z},
\qquad h+1\le j\le k-N_0 .
\end{equation*}
The stated purpose of the construction, in \cite[\S1, before (1.10)]{B-CMP122a}, is to
obtain ``a new, much smaller large field region''; the inclusions may therefore
be proper. The definition does not require $Z''_j\cap X$ to contain the creation cubes of $X$,
that is, the $\mathcal{P}_s$-cubes contained in $\mathcal{W}^{\wedge}_s$ that
Proposition~\ref{4.15:prop-creation-cube} places in $X$.

No statement used together with Proposition~\ref{4.15:prop-creation-cube} is quantified
over these lower regions:
\begin{itemize}
\item The structural statement on the components of the top large-field region is a
statement about every component
of the top region $\Lambda_k^c$, which is $Z_k$ or $Z_k^{\mathrm{post}}$ according as it is
read on $\sigma'_k$ or on $\tau$.
\item The waiting rule reads components of $\Lambda_k^c$.
\item The regime clause that uses the waiting rule reads only the component acted upon at
its own level.
\item The second clause of condition~(ii) of \cite{B-CMP122a} concerns the previous regions
contained in a component. Of these it asks only the size condition~(i) there.
\end{itemize}
Hence Proposition~\ref{4.15:prop-creation-cube} is, as asserted, a statement about $Z_k$ and
$Z_k^{\mathrm{post}}$ only.

(b) Inspection of its statement shows that Proposition~\ref{4.15:prop-creation-cube} makes
no reference to condition~(ii). Consequently the remark of this chapter on Ba{\l}aban's
preparatory large fields, which concerns that condition, has no bearing on it.

(c) Fix a term and one of its histories, and let $k_1$ be the first level at which
$\mathbf{R}$ acts in it. By the definition of $k_1$, $\mathcal{A}_j=\emptyset$ for
$1\le j<k_1$. Thus the history is pure-$\mathbf{T}$ below level $k_1$; this is the first
assertion of Corollary~\ref{4.15:cor-first-action}(b).

Now apply Proposition~\ref{4.15:prop-pure-histories}, with $k=k_1$, to the term
$\sigma'_{k_1}$ of $\mathbf{T}\rho_{k_1-1}$ in this history. Its only hypotheses are the
reading assumed there and the purity just established. It gives
the following: for every $j$
with $1\le j\le k_1$, every component of $Z_j$ contains a $\mathcal{P}_s$-cube contained in
$\mathcal{W}^{\wedge}_s$, for some $s$ with $1\le s\le j$. This is the conclusion of
Proposition~\ref{4.15:prop-creation-cube} for $Z_j$. At level $k_1$ this yields the last
assertion of Corollary~\ref{4.15:cor-first-action}(b), since the components acted upon are
components of $Z_{k_1}$.

Now consider a level $j<k_1$, where $\mathcal{A}_j(\sigma'_j)=\emptyset$.
Definition~\ref{4.15:def-post-action-region} then takes $\tau_j=\sigma'_j$. Hence
$Z_j^{\mathrm{post}}=\Lambda_j(\tau_j)^c=Z_j$, and that definition imposes nothing. It
first carries content at level $k_1$, where $\mathcal{A}_{k_1}\ne\emptyset$ and
$Z_{k_1}^{\mathrm{post}}$ is fixed by that definition.
The regions of later levels are produced from such post-action regions. Hence the reading of
Definition~\ref{4.15:def-post-action-region} enters at the first action of $\mathbf{R}$ and
after it, and not before.
\end{proof}

\begin{remark}[What is not asserted about creation of large-field regions]
\label{4.15:rem-not-asserted}
We record here what the results of this section do not assert.
\begin{enumerate}
\item[(1)] Nothing is asserted under a formulation of the no-creation condition in which that
condition holds outright whenever fewer than $\mathsf{N}_k$ steps precede level $k$, that
is, whenever $k<\mathsf{N}_k$. Such a formulation is neither of the two formulations of
Definition~\ref{4.15:def-no-creation}. Under it nothing bars an action of $\mathbf{R}$ at a
level $k<\mathsf{N}_k$ on a component that satisfies the two conditions, other than the
no-creation condition, under which $\mathbf{R}$ acts at a level (those assumed in
Lemma~\ref{4.15:lem-action-level} besides the no-creation condition), and
\eqref{4.15:eq-action-level-a} is not derivable.
\item[(2)] It is not asserted that $\mathbf{R}$ ever acts. Nothing is asserted about the values
of the stopping stage $n_*$ of the recursion that builds the large-field region of a step
and its creation set \eqref{4.15:eq-creation-set-a}, of the
stabilisation index $q_0$ of the modification map, of Ba{\l}aban's integer $N_0(k)$, of
$\mathsf{N}_k$, or of the first level at which $\mathbf{R}$ acts in a history; none of
these numbers is bounded.
\item[(3)] Proposition~\ref{4.15:prop-creation-cube} is not asserted under the competing
description of the top large-field region that Definition~\ref{4.15:def-post-action-region}
states beside the one assumed in that proposition. Nor is it asserted for the lower
entries $Z''_j\cap X$ of a component $X$ on which $\mathbf{R}$ acts, where $Z''_j$ are the
members of Ba{\l}aban's new sequence of regions; these lower entries are considered
separately in this section.
\item[(4)] It is not asserted that the description of $\mathbf{R}$ in
Definition~\ref{1.5:def-densities} determines the region read by the next
$\mathbf{T}$-step. That description concerns the operator and the density on a component on
which $\mathbf{R}$ acts; the first of the two descriptions of the top region in
Definition~\ref{4.15:def-post-action-region}, the one assumed in
Proposition~\ref{4.15:prop-creation-cube}, is an assumption there and is not derived.
\item[(5)] Let $C$ be a component, at level $k$, of either of the two regions of
Proposition~\ref{4.15:prop-creation-cube} (the large-field region produced by the step
$k-1\to k$, or the large-field region after the actions of $\mathbf{R}$ at level $k$), and
let $\square$ be a cube of a creation set $\mathcal{W}^{\wedge}_t$, $t\le k$, found in $C$
by the proof of that proposition. It is not asserted that $\square$ lies in the top
large-field region at every level between $t$ and $k$.
\item[(6)] Nothing is asserted about the analytic content of $\mathbf{R}$. This covers the
positivity of its denominators, the bounds \cite[(1.79)--(1.89)]{B-CMP122b} and the
reference density. Nothing is asserted about the operator theorem at free current, about
clause~(i) of Theorem~0, or about $N_0(k)$ beyond the remark on Ba{\l}aban's integer
$N_0(k)$ among the results of this section and part~(c) of the lemma on which that remark
rests.
\item[(7)] Nothing is asserted about the particular numbers of layers of enlargement that appear
in \eqref{4.15:eq-small-field-complement-a} and \eqref{4.15:eq-small-field-complement-b}.
Only the form of Definition~\ref{4.15:def-small-field-complement} is used, in which these
numbers are arbitrary non-negative integers. No value of the constant $L_0$ is used.
\item[(8)] Consider, at the first step, the narrower choice $P':=P_1'$ of the primed set of
Definition~\ref{4.15:def-small-field-complement} from which the creation set
\eqref{4.15:eq-creation-set-a} is built, $P_0$ not being counted. Nothing is asserted under
this choice: under it neither Proposition~\ref{4.15:prop-creation-cube} nor
\eqref{4.15:eq-action-level-a} is derivable by the argument of this section, as the proof
below shows. Here $P_1'$ is the union of the coarse
cubes meeting the large-field set $P_1$ of \cite[(1.4)]{B-CMP119}, and $P_0$ is the
large-field set of \cite[(1.2)]{B-CMP119}.
\end{enumerate}
\end{remark}

\begin{proof}
We give the arguments behind items (1), (5) and (8). The remaining items state the scope of
Lemma~\ref{4.15:lem-action-level} and Proposition~\ref{4.15:prop-creation-cube} and need no
argument.

\emph{Item (1).} Let $k<\mathsf{N}_k$. The first formulation of
Definition~\ref{4.15:def-no-creation} requires $k\ge\mathsf{N}_k$, so it fails at such $k$.
In the second formulation we have $k-\mathsf{N}_k+1\le 0$, so the range of levels $t$ is
$1\le t\le k$. That formulation therefore bears on every step $1,\dots,k$ and does not hold
outright. Hence a formulation that holds outright at $k<\mathsf{N}_k$ is neither of the two.

Under that formulation, the no-creation condition at such $k$ holds whatever creations lie
in the component. An action at level $k$ is then barred only by the two other conditions.
Suppose a component at a level $k<\mathsf{N}_k$ satisfies these two conditions. Then
an action on it is permitted, and at that level $k\ge 1+\mathsf{N}_k$ fails. So
\eqref{4.15:eq-action-level-a} cannot be derived from the conditions under that
formulation.

\emph{Item (5).} The proofs of Lemma~\ref{4.15:lem-action-level} and
Proposition~\ref{4.15:prop-creation-cube} pass from $C$ back to the step $t$ through a chain
of components, namely the components that the argument uses. For each member $D$ of this
chain they give
\begin{equation*}
\square\subset D\subset C .
\end{equation*}
The chain is not asserted to contain, at every level $j$ with $t\le j\le k$, a component of
the top large-field region at level $j$; these inclusions therefore do not give the
persistence of $\square$ that item (5) leaves unasserted. They give what the no-creation
condition needs. In both formulations of
Definition~\ref{4.15:def-no-creation}, that condition asks only whether a creation at a
level in the prescribed range lies inside $C$.

\emph{Item (8).} At $t=1$, Definition~\ref{4.15:def-small-field-complement} takes
$P'=P_0'\cup P_1'$, $Q'=Q_1'$ and $Y=\emptyset$, where $P_0'$, $P_1'$, $Q_1'$ are the unions
of the $\mathcal{P}_1$-cubes meeting the plaquette large-field sets $P_0$, $P_1$ of
\cite[(1.2), (1.4)]{B-CMP119} and the large-field set $Q_1$ of the approximate fluctuation
field; with this choice \eqref{4.15:eq-small-field-complement-b} exhibits $\Omega_1^c$ as
the union of an enlargement of $Q_1'$, the enlargement $P_0'^{\sim(a_0+3)}$ of $P_0'$, and
an enlargement of $P_1'$.

The creation set $\mathcal{W}^{\wedge}_1$ is built, by \eqref{4.15:eq-creation-set-a}, from
the primed sets of the step, among them $P'$.

Now take the narrower choice $P':=P_1'$ and assume $P_0\neq\emptyset$. Then
$P_0'\neq\emptyset$. The member $P_0'^{\sim(a_0+3)}$ enters
\eqref{4.15:eq-small-field-complement-b} as the enlargement of $P_0'$, which under this
choice is not one of the sets from which $\mathcal{W}^{\wedge}_1$ is built. Consider the form
of Definition~\ref{4.15:def-small-field-complement}. At $t=1$ there are no large-field
regions at level $0$ and $Y=\emptyset$, so the form presents $\Omega_1^c$ as a union of
enlargements of regions contained in $P'\cup Q'$, that is, under this choice, in
$P_1'\cup Q_1'$. The display \eqref{4.15:eq-small-field-complement-b} does not exhibit
$\Omega_1^c$ in this form, since it gives its member $P_0'^{\sim(a_0+3)}$ as an enlargement
of $P_0'$, a region which in general is not contained in $P_1'\cup Q_1'$; under this choice
the form of Definition~\ref{4.15:def-small-field-complement} at $t=1$ is therefore not
supplied by \eqref{4.15:eq-small-field-complement-b}.

There are no large-field regions at level $0$. So a component of $\Lambda_1^c$ consisting
of enlargements of $P_0'$ alone need contain no point of any creation set. The proof of
Proposition~\ref{4.15:prop-creation-cube} finds a cube of a creation set in a component by
following the chain of components back to the step whose creation set contains that cube.
On the descendants of
such a component there may be no such cube to find. Hence neither
Proposition~\ref{4.15:prop-creation-cube} nor \eqref{4.15:eq-action-level-a} is derivable
by the present argument.

Indeed, suppose $k$ is a level with $k=\mathsf{N}_k=\check{R}_k$, and $C$ is such a
descendant at level $k$ inside which no creation at any level $t$ with $1\le t\le k$ lies.
In the first formulation we then have
$k\ge\mathsf{N}_k$, and the prescribed range $k-\mathsf{N}_k<t\le k$ is $1\le t\le k$. In
the second formulation the range
\begin{equation*}
\max\{1,k-\mathsf{N}_k+1\}\le t\le k
\end{equation*}
is the same. So the no-creation condition holds for $C$ in both formulations, and an action
at level $k=\check{R}_k$ is not excluded by it.

The choice $P'=P_0'\cup P_1'$ of Definition~\ref{4.15:def-small-field-complement} is the one
under which the description of $\Omega_1^c$ above holds, by
\eqref{4.15:eq-small-field-complement-b}. That choice is part of the statement assumed in
Lemma~\ref{4.15:lem-action-level} and Proposition~\ref{4.15:prop-creation-cube}.
\end{proof}

\begin{remark}[The two lowest tower levels: Ba{\l}aban's regularity conditions, carried as a
two-member clause on the datum at levels $0$ and $1$; printed without proof in
{\cite[p.~396, (3.35), (3.37)]{B-CMP99b}}; cited as printed]\label{4.16:rem-lowest-levels}
The following two sentences are printed without proof in
\cite[p.~396, (3.35), (3.37)]{B-CMP99b}; they are cited as printed. The first is
\cite[(3.35)]{B-CMP99b}:
\begin{quotation}
for an arbitrary cube $\square$ of the described above class, and for a configuration $U$
there exists a gauge transformation $u$ on $\square$ such that $U^u = e^{i\eta A}$.
[sic] and if the index of $\square$ is $j$, then
$|A| < O(1)M\alpha_0(L^j\eta)^{-1}$, $|\nabla^\eta A| < O(1)M\alpha_0(L^j\eta)^{-2}$
on $\square$, where $O(1)M$ is a size of $\square$ in $T_{L^{-j}}$
\end{quotation}
The second is \cite[(3.37)]{B-CMP99b}:
\begin{quotation}
$U$ satisfies the condition (3.35), and $|A'| < \alpha_1(L^j\eta)^{-1}$,
$|\nabla^\eta_U A'| < \alpha_1(L^j\eta)^{-2}$ on $\Omega_j$, $j = 0,\dots,k$
\end{quotation}
Here, in both sentences, the class of cubes $\square$, the domains $\Omega_j$, the lattice
spacing $\eta$, the potentials $A$, $A'$ and the radii $\alpha_0$, $\alpha_1$ are those of
\cite[p.~396]{B-CMP99b}. All other letters inside the two quotations are also
Ba{\l}aban's on that page: $M$ is his big-block parameter, the cubes of the class being
unions of big blocks of the lattice written $T_{L^{-j}}$ there; $k$ is the top index of his
sequence of domains $\{\Omega_j\}$, $0\le j\le k$; $\nabla^\eta$ and $\nabla^\eta_U$ are
his lattice derivative and his lattice derivative covariant with respect to $U$; $U^u$ is
the gauge transform of $U$ by $u$; and $O(1)$ denotes a number $\ge 10$. The letters
$M$ and $k$ inside the quotations are read as on \cite[p.~396]{B-CMP99b}; the block size
and the level at which the clause below is applied are fixed where it is used.

These two sentences are definitions of regularity classes in \cite{B-CMP99b}; writing the
configuration of the datum as $U'U$, with $U$ taking values in the gauge group and
$U' = e^{i\eta A'}$ as on \cite[p.~396]{B-CMP99b}, what this book carries as printed is
the following two-member clause on the datum at the two lowest levels:
\begin{itemize}
\item[(i)] the real part $U$ of the configuration satisfies the gauge condition of
\cite[(3.35)]{B-CMP99b} on the boxes of the two lowest levels $j\in\{0,1\}$ that meet
the large-field region $Z$;
\item[(ii)] the potential $A'$ obeys the level-$0$ radii of \cite[(3.37)]{B-CMP99b}, and,
in the lettering of the background pair (Definition~\ref{4.5:def-admissible-pairs}), in
which the current $J$ is carried beside $A'$, the current $J$ obeys the same level-$0$ radii;
\end{itemize}
this clause is used by the operator theorem at free current, and so by clause~(0) of
Theorem~0, proved as Theorem~\ref{4.1:thm-clause-zero}, and by clause~(i) of Theorem~0,
which is proved in one induction with clause~(0); it is also an input of the statement that
the data of the correlation theorem lie in the domain of the operator theorem at free
current, which clause~(iii) of Theorem~0 uses for its local random-walk bounds.

These two conditions together form one of the four items of Ba{\l}aban's text that the
book counts as printed without proof; all four are listed in the index of statements
printed without proof that closes Part~III.
\end{remark}

\begin{definition}[Positions of fine indices and weighted bounds]\label{4.16:def-fine-positions}
Lengths are measured in level-$k$ units, in which $\ell_n=L^{n-k}$ is the lattice spacing of level $n$
and $\eta=\ell_0=L^{-k}$; the finest standard partition $\pi^{(0)}$ is the standard partition at
mesh $\eta$, and the shells $X_n$ are the territories of Definition~\ref{4.14:not-free-current-notation}.

\emph{Fine sites, bonds and plaquettes.} The \emph{fine sites} are the centres of the cells of
$\pi^{(0)}$. The standard partitions are nested, so every cube of every one of them is a union of
cells of $\pi^{(0)}$. A centre lies in the interior of its own cell and in no other cell.
Hence no fine site lies on the boundary of a cube of any standard partition.
A \emph{fine bond} $b=\langle x,x+\eta e_\mu\rangle$ is the closed segment joining two
neighbouring fine sites, where $e_\mu$ is the $\mu$-th unit vector, as in \cite{B-CMP109}.
We write $b_-:=x$ and $b_+:=x+\eta e_\mu$ for its end-points. The \emph{position} of $b$ is its
midpoint $x_b$. A fine plaquette is positioned at its centre.

\emph{Pairs and difference quotients.} A \emph{pair} $(b,b')$ is a pair of parallel
nearest-neighbour fine bonds, that is, $b'=b+\eta e_{\mu'}$ for some direction $\mu'$. The
collinear pair, $\mu'=\mu$, is admitted. For a field $A$ on the fine bonds, a pair carries the
plain difference quotient
\begin{equation}\label{4.16:eq-fine-positions-1}
\nabla A(b,b') := \eta^{-1}\bigl(A(b')-A(b)\bigr).
\end{equation}
These are the difference quotients between parallel neighbouring bonds that enter the
regularity condition on the data. The \emph{position} of the pair is the midpoint of the two
midpoints,
\begin{equation}\label{4.16:eq-fine-positions-2}
x_{(b,b')} := \tfrac12\,\bigl(x_b+x_{b'}\bigr).
\end{equation}
Neither the admission of the collinear pair nor this choice of position affects the estimates
made with pairs later in this chapter; the reason is given where those estimates are made.

\emph{Carried indices and gauge transformations on a set.} A closed set $S$ \emph{carries} a
fine bond $b$ if $b\subset S$, the bond being the closed segment. It carries a pair $(b,b')$ if
it carries both $b$ and $b'$. A map $u$ \emph{on $S$} is a $G$-valued map on the end-points of
the bonds that $S$ carries. For such $u$ and every bond $b$ carried by $S$ we set
\begin{equation}\label{4.16:eq-fine-positions-3}
U^{u}(b) := u(b_-)\,U(b)\,u(b_+)^{-1}.
\end{equation}

\emph{The reading $(\mathrm{W})$ of a weighted bound.} Let $S$ be a closed set and $a>0$.
Consider a clause of the form
\begin{equation*}
\text{``}\ \ell_n|A|,\ \ell_n^2|\nabla A| < a \ \text{ on } S\cap X_n
\ \text{ for each territory } X_n \text{ that } S \text{ meets''},
\end{equation*}
the territories being the shells $X_n$.
Such a clause is read index by index. For each $n$ with $S\cap X_n\neq\emptyset$ it means the
following two families of bounds:
\begin{equation}\label{4.16:eq-fine-positions-a}
\begin{aligned}
\ell_n\,|A(b)| &< a
  &&\text{for every bond $b$ carried by $S$ with $x_b\in X_n$},\\
\ell_n^2\,|\nabla A(b,b')| &< a
  &&\text{for every pair $(b,b')$ carried by $S$ with $x_{(b,b')}\in X_n$}.
\end{aligned}
\tag{W}
\end{equation}
In particular, consider a bond or a pair carried by $S$ whose position lies in
$X_n\cap X_{n+1}$. It is subject to both the bound of index $n$ and the bound of index $n+1$ in
\eqref{4.16:eq-fine-positions-a}; for a bond these carry the weights $\ell_n$ and $\ell_{n+1}$,
for a pair the weights $\ell_n^2$ and $\ell_{n+1}^2$.
\end{definition}

\begin{definition}[The data region, box regularity and the hull condition]\label{4.16:def-data-region}
We use the scales $\eta=L^{-k}$ and $\ell_n=L^{n-k}$, the block sizes $M_1\mid M_1^{\sharp}\mid M$, the
admissible nest $(\Omega_n)_{0\le n\le k}$ together with $\Omega_t\subset\Omega_k$, and the shells
$X_n$ of Definition~\ref{4.14:not-free-current-notation}. We also use the carrying of bonds and pairs by a closed set,
site maps $u$ on a closed set, the gauge-transformed variable $U^{u}$, the plain difference quotient
$\nabla A(b,b')$ and the index-by-index reading of a weighted bound of
Definition~\ref{4.16:def-fine-positions}.

\begin{enumerate}
\item \emph{The data region.} A closed set $\Omega_0^{\mathrm{dat}}\subset\mathbb{T}$ is fixed, the
\emph{data region}. The configuration $U$ is $G$-valued on the fine bonds carried by
$\Omega_0^{\mathrm{dat}}$. The domain on which the operator theorem at free current takes its data is
this region. That theorem's real configuration space $\mathcal{U}_{\mathbb{R}}^{\mathrm{OP}}$ consists of the
$G$-valued $U$ on the fine bonds carried by $\Omega_0^{\mathrm{dat}}$ that satisfy that theorem's hypotheses
$(\mathrm{u}1)$ and $(\mathrm{reg})$ together with condition $(\mathrm{reg})^{\sharp}$ of the third item below.

\item \emph{The hull condition.} The perturbation fields entering the operators of the step live on bond sets
$\mathsf{B}\subset\hat{\mathsf{B}}$ built on a base set $\mathsf{S}\subseteq\Omega_0$, where $\Omega_0$ is the
largest domain of the admissible nest; we call $\Omega_0$ their \emph{hull}. This hull satisfies
\begin{equation}\label{4.16:eq-data-region-b}
\Omega_0^{+2\eta}\subset\Omega_0^{\mathrm{dat}},
\end{equation}
where $\Omega_0^{+2\eta}$ denotes the set of points at sup-norm distance at most $2\eta$ from $\Omega_0$.
We call \eqref{4.16:eq-data-region-b} the \emph{hull condition} $(\mathrm{col})$.

\item \emph{Box regularity.} For $0\le j\le k$ put
\begin{equation}\label{4.16:eq-data-region-1}
s_j:=M_1^{\sharp}\ell_j .
\end{equation}
Let $\mathcal{D}^{\sharp}_j$ be the family of boxes of side $2s_j$ centred at the points of the lattice
$s_j\mathbb{Z}^{d}$ that lie in $X_j$, and put $\mathcal{D}^{\sharp}:=\bigcup_j\mathcal{D}^{\sharp}_j$. For a box
$\square\in\mathcal{D}^{\sharp}_j$ with centre $y$, let $\square^{4}$ be the closed cube of half-side
$5s_j$ about $y$:
\begin{equation}\label{4.16:eq-data-region-2}
\square^{4}:=\{x:\ |x-y|_\infty\le 5s_j\}.
\end{equation}
The regularity clause of the operator theorem at free current is read in the following form, called
condition $(\mathrm{reg})^{\sharp}$. For every $\square\in\mathcal{D}^{\sharp}$ there are a $G$-valued $u$ on
$\square^{4}\cap\Omega_0^{\mathrm{dat}}$ and a function $A$ on the bonds carried by
$\square^{4}\cap\Omega_0^{\mathrm{dat}}$ such that
\begin{equation}\label{4.16:eq-data-region-a}
\begin{gathered}
U^{u}=e^{i\eta A},\\
\ell_n|A|<a^{\sharp}\quad\text{and}\quad \ell_n^{2}|\nabla A|<a^{\sharp}
\quad\text{on } \square^{4}\cap X_n,
\end{gathered}
\end{equation}
for each shell $X_n$ that $\square^{4}$ meets (in the clause as the operator theorem at free current
states it, there are at most two such shells, and they are adjacent).
The bounds in \eqref{4.16:eq-data-region-a} are read index by index, in the sense of
Definition~\ref{4.16:def-fine-positions}. The territories of the theorem are read as the shells $X_n$, as in
Definition~\ref{4.14:not-free-current-notation}; in particular the level-$0$ territory is all of $X_0$.
\end{enumerate}

The constant $a^{\sharp}$ is the regularity constant of the operator theorem at free current. It is subject
to the ceilings imposed on it there, among them $a^{\sharp}\le a_{\mathrm{reg}}^{\oplus}$. In
\eqref{4.16:eq-data-region-a}, $\nabla$ is the plain difference quotient
$\nabla A(b,b')=\eta^{-1}(A(b')-A(b))$ between parallel neighbouring bonds. Its modulus $|\nabla A|$ is the
quantity written $|\nabla^{\eta}A|$ in \cite[(3.35)]{B-CMP99b}.
\end{definition}

We record two consequences of the definition and one remark on the size of $\square^{4}$.

First, the nest is decreasing, so $\Omega_1\subseteq\Omega_0$. Moreover
$\Omega_0\subset\Omega_0^{+2\eta}$. The hull condition \eqref{4.16:eq-data-region-b} therefore gives
\begin{equation}\label{4.16:eq-data-region-3}
\Omega_1\subseteq\Omega_0\subset\Omega_0^{+2\eta}\subset\Omega_0^{\mathrm{dat}} .
\end{equation}

Second, the bond sets $\mathsf{B}\subset\hat{\mathsf{B}}$ on which the perturbation fields of the step live
contain only bonds lying within sup-norm distance $2\eta$ of their base set $\mathsf{S}$. Since
$\mathsf{S}\subseteq\Omega_0$, every such bond lies in $\Omega_0^{+2\eta}$, hence in $\Omega_0^{\mathrm{dat}}$ by
\eqref{4.16:eq-data-region-b}; so $\mathsf{B}$ and $\hat{\mathsf{B}}$ are carried by $\Omega_0^{\mathrm{dat}}$.
Hence $U$ is given on every bond of $\mathsf{B}$ and of $\hat{\mathsf{B}}$.

Finally, in the clause as the operator theorem at free current states it, the superscript of $\square^{4}$
counts enlargements by $s_j$: $\square^{4}$ is the cube concentric with $\square$ of half-side
$(1+4)s_j=5s_j$, hence of side $10s_j$, and it is the cube of \eqref{4.16:eq-data-region-2}. If the clause
were read with a larger cube instead, for instance an enlarged cube attached to $\square$ and the integer $4$
in the manner of \cite{B-CMP109} (``of the size \dots\ and with a center at the center of'' the given cube)
with side $1+2\cdot 4$ times that of the box, hence of half-side $9s_j$, then every inclusion of a set into
$\square^{4}$ used in what follows would hold a fortiori, and the fitting of cubes into boxes used in what
follows would still hold under the same one-sided floor on the block sizes.

\begin{definition}[Ba{\l}aban's cube class and its bounded sub-class]\label{4.16:def-cube-class}
Let the dimension $d$, the lengths $\ell_n=L^{n-k}$ with $\eta=\ell_0=L^{-k}$, the block sizes
$M_1\mid M_1^{\sharp}\mid M$ (powers of $L$), the nest
$\Omega_0\supseteq\Omega_1\supseteq\dots\supseteq\Omega_k$ and the shells $X_0,\dots,X_k$, with
$X_n=\emptyset$ for $n>k$, be as in Definition~\ref{4.16:def-data-region}, and let
$0\le j\le k$.
\begin{itemize}
\item[(i)] A \emph{big block of level $j$} is a level-$j$ $M_1$-cube, that is, a cube of the
standard partition at mesh $M_1\ell_j$, a union of $M_1^{d}$ cells of side $\ell_j$. These are the
``big blocks of the lattice
$T_{L^{-j}}$'' of \cite[\S3]{B-CMP99b}, whose block size $M$ is our $M_1$.
\item[(ii)] For an integer $\mathsf{N}\ge 1$, an \emph{aligned $\mathsf{N}$-block cube of level $j$}
is a closed cube that is the union of $\mathsf{N}^d$ big blocks of level $j$; its side is
$\mathsf{N}M_1\ell_j$. The lower bound $\mathsf{N}\ge 10$ placed on the class cubes below (in
$\mathcal{C}_j$ and in Ba{\l}aban's class) is Ba{\l}aban's convention ``$O(1)$ will
mean a number $\geqq 10$'' \cite[\S3]{B-CMP99b}; the same paper adds that the number $O(1)$ in the
condition \cite[(3.35)]{B-CMP99b} ``can be taken as equal to $12$''.
\item[(iii)] Put $\mathsf{N}^{\sharp}:=4M_1^{\sharp}/M_1$. The \emph{bounded sub-class at level $j$} is
\begin{equation}\label{4.16:eq-cube-class-a}
\begin{split}
\mathcal{C}_j:=\bigl\{\mathfrak{E}:\ &\mathfrak{E}\text{ an aligned }\mathsf{N}\text{-block cube of
level }j\text{ with }10\le\mathsf{N}\le\mathsf{N}^{\sharp},\\
&\mathfrak{E}\subset X_j\cup X_{j+1},\ \mathfrak{E}\cap X_j\neq\emptyset\bigr\}.
\end{split}
\end{equation}
\item[(iv)] \emph{Ba{\l}aban's class.} Ba{\l}aban introduces a class of cubes of which he says:
``For each cube $\square$ of this class there exists a unique index $j$,
$0\leqq j\leqq k$'' \cite[\S3]{B-CMP99b} (in his sentences quoted here $\square$ is a class cube,
our $\mathfrak{E}$; outside quotations, $\square$ and $\square^{4}$ denote the boxes of
Definition~\ref{4.16:def-data-region}). At the index $0$ the region is
``$\Lambda_0=\Omega_0\setminus\Omega_1$, where $\Omega_0=\{$a union of big blocks of the lattice
$T_1$, such that their distances to $\Omega_1$ are $\leqq RM\}$'' \cite[\S3]{B-CMP99b}; here $R$ is
a constant of \cite[\S3]{B-CMP99b}, unrelated to the ball radius $R$ of the glossary
(Definition~\ref{0.2:not-glossary}), and his $M$ is our $M_1$, so that the distance bound reads
$\leqq RM_1$. This set is to
be distinguished from the hull $\Omega_0$ of the nest; we write $\Omega_0^{\mathrm{cl}}$ for it and
call it the \emph{class hull}. It is the union of the level-$0$ big blocks satisfying the distance
bound in the quotation, and it is closed. For $i\ge 0$ write $B^{i}(\Lambda_i)$ for the closed region
of index $i$ of \cite[\S3]{B-CMP99b} (his $\Lambda_i$ is a region of that paper, unrelated to the
regions written $\Lambda$ elsewhere in this chapter). In the notation of this book these regions are
\begin{equation}\label{4.16:eq-cube-class-1}
B^{0}(\Lambda_0)=\operatorname{cl}\bigl(\Omega_0^{\mathrm{cl}}\setminus\Omega_1\bigr),\qquad
B^{1}(\Lambda_1)=X_1,\qquad B^{2}(\Lambda_2)=X_2 ,
\end{equation}
and, in the translation used in this book, $B^{i}(\Lambda_i)=X_i$ for $1\le i\le k$ and
$B^{i}(\Lambda_i)=\emptyset$ for $i>k$, in accordance with $X_i=\emptyset$ for $i>k$.
A cube of Ba{\l}aban's class \emph{at index $j$} is an aligned cube of level-$j$ big blocks with at
least ten blocks a side, contained in $B^{j}(\Lambda_j)\cup B^{j+1}(\Lambda_{j+1})$ and meeting
$B^{j}(\Lambda_j)$. Thus at an index $j\ge1$ a cube of Ba{\l}aban's class is an aligned cube of
level-$j$ big blocks with at least ten blocks a side, contained in $X_j\cup X_{j+1}$ and meeting
$X_j$ \cite[\S3]{B-CMP99b}.
\end{itemize}
These objects have the following properties.
\begin{itemize}
\item[(a)] $\mathsf{N}^{\sharp}$ is an integer, four times a power of $L$. The range
$10\le\mathsf{N}\le\mathsf{N}^{\sharp}$ in \eqref{4.16:eq-cube-class-a} is non-empty as soon as
$\mathsf{N}^{\sharp}\ge 10$. It contains $\mathsf{N}=12$ as soon as $M_1^{\sharp}$ is at least three
times $M_1$, which holds under the floor \eqref{4.16:eq-lowest-level-regularity-b}.
\item[(b)] Every cube of Ba{\l}aban's class at index $1$ with at most $\mathsf{N}^{\sharp}$ blocks a
side belongs to $\mathcal{C}_1$.
\item[(c)] Every cube $\mathfrak{E}$ of Ba{\l}aban's class at index $0$ with at most
$\mathsf{N}^{\sharp}$ blocks a side belongs to $\mathcal{C}_0$ and lies in $\Omega_0^{\mathrm{cl}}$.
Membership in $\mathcal{C}_0$ itself imposes no condition relative to a hull.
\item[(d)] Suppose that Ba{\l}aban's $O(1)$ is read as the single number $12$, so that the cubes of
his class are exactly the aligned $12$-block cubes meeting the conditions of (iv), and that the
floor \eqref{4.16:eq-lowest-level-regularity-b} holds. Then $\mathcal{C}_1$ contains the whole of
Ba{\l}aban's class at index $1$, and $\mathcal{C}_0$ contains the whole of his class at index $0$.
\item[(e)] Ba{\l}aban's regularity bound \cite[(3.35)]{B-CMP99b} reads ``$|A|<O(1)M\alpha_0(L^j\eta)^{-1}$,
$|\nabla^{\eta}A|<O(1)M\alpha_0(L^j\eta)^{-2}$ on $\square$, where $O(1)M$ is a size of $\square$''.
Here $A$ is the field defined, as in the regularity condition $(\mathrm{reg})^{\sharp}$ of
Definition~\ref{4.16:def-data-region}, by $U^{u}=e^{i\eta A}$ for a $G$-valued gauge $u$ on the
cube, and $\nabla^{\eta}$ is the plain difference quotient $\nabla$ between parallel neighbouring fine
bonds. On a class cube $\mathfrak{E}$ with $\mathsf{N}$ blocks a side this bound is the pair of
bounds
\begin{equation}\label{4.16:eq-cube-class-2}
\ell_j|A|<\mathsf{N}M_1\alpha_0,\qquad \ell_j^{2}|\nabla A|<\mathsf{N}M_1\alpha_0
\quad\text{on }\mathfrak{E}.
\end{equation}
These bounds use the single weight $\ell_j$ throughout $\mathfrak{E}$, including at points of
$X_{j+1}$. Ba{\l}aban notes that ``these regularity conditions do not impose any constraints on $U$
outside the domain $\Omega_0$'' \cite[\S3]{B-CMP99b}; his $\Omega_0$ there is the set of the
quotation in (iv), our $\Omega_0^{\mathrm{cl}}$.
\end{itemize}
\end{definition}

\begin{proof}
(a) Both $M_1$ and $M_1^{\sharp}$ are powers of $L$, and $M_1\mid M_1^{\sharp}$. Hence
$M_1^{\sharp}/M_1$ is a power of $L$, and $\mathsf{N}^{\sharp}=4M_1^{\sharp}/M_1$ is four times a power
of $L$, in particular an integer. The integers $\mathsf{N}$ with $10\le\mathsf{N}\le\mathsf{N}^{\sharp}$
form a non-empty set exactly when $\mathsf{N}^{\sharp}\ge 10$. If $M_1^{\sharp}\ge 3M_1$, then
$\mathsf{N}^{\sharp}=4M_1^{\sharp}/M_1\ge 12$, so that $\mathsf{N}=12$ is in this set.

(b) Let $\mathfrak{E}$ be a cube of Ba{\l}aban's class at index $1$ with $\mathsf{N}$ blocks a side,
where $\mathsf{N}\le\mathsf{N}^{\sharp}$. It is an aligned cube of level-$1$ big blocks, and by the
class condition $\mathsf{N}\ge 10$. Thus it is an aligned $\mathsf{N}$-block cube of level $1$ with
$10\le\mathsf{N}\le\mathsf{N}^{\sharp}$. By \eqref{4.16:eq-cube-class-1},
\begin{equation*}
\mathfrak{E}\subset B^{1}(\Lambda_1)\cup B^{2}(\Lambda_2)=X_1\cup X_2,
\qquad
\mathfrak{E}\cap X_1=\mathfrak{E}\cap B^{1}(\Lambda_1)\neq\emptyset .
\end{equation*}
These are the conditions of \eqref{4.16:eq-cube-class-a} at $j=1$, so $\mathfrak{E}\in\mathcal{C}_1$.

(c) We first record three inclusions.

First, the set $\Omega_0^{\mathrm{cl}}\setminus\Omega_1$ is contained in $\mathbb{T}\setminus\Omega_1$.
Taking closures gives
$\operatorname{cl}(\Omega_0^{\mathrm{cl}}\setminus\Omega_1)\subset\operatorname{cl}(\mathbb{T}\setminus\Omega_1)=X_0$.
Since $\Omega_0^{\mathrm{cl}}$ is a finite union of closed cubes, it is closed, and the same closure
lies in $\Omega_0^{\mathrm{cl}}$. Hence
\begin{equation*}
B^{0}(\Lambda_0)\subset X_0\cap\Omega_0^{\mathrm{cl}} .
\end{equation*}

Second, by condition (A1) of Proposition~\ref{4.14:prop-prepared-data-admissible},
$\Omega_1$ is a union of level-$1$ $M$-cubes. These are closed and, the torus being compact, finite
in number, so $\Omega_1$ is closed. If $k\ge 2$, then $X_1=\operatorname{cl}(\Omega_1\setminus\Omega_2)$
is therefore contained in $\Omega_1$. If $k=1$, then $X_1=\Omega_1$. In either case
$X_1\subset\Omega_1$.

Third, the partitions are nested, $\ell_1=L\ell_0$ and $M_1\mid M$. Hence every level-$1$ $M$-cube is
a union of level-$0$ $M_1$-cubes, so $\Omega_1$ is a union of level-$0$ big blocks. Each of these
blocks lies at distance $0$ from $\Omega_1$, so each belongs to the class hull, and
$\Omega_1\subset\Omega_0^{\mathrm{cl}}$.

Now let $\mathfrak{E}$ be a cube of Ba{\l}aban's class at index $0$ with $\mathsf{N}$ blocks a side,
where $\mathsf{N}\le\mathsf{N}^{\sharp}$. As in (b), $\mathfrak{E}$ is an aligned $\mathsf{N}$-block
cube of level $0$ with $10\le\mathsf{N}\le\mathsf{N}^{\sharp}$. By \eqref{4.16:eq-cube-class-1} and the
three inclusions above,
\begin{equation*}
\mathfrak{E}\subset B^{0}(\Lambda_0)\cup B^{1}(\Lambda_1)\subset X_0\cup X_1 .
\end{equation*}
Moreover, $\mathfrak{E}\cap X_0$ contains $\mathfrak{E}\cap B^{0}(\Lambda_0)$, which is non-empty.
Hence $\mathfrak{E}\in\mathcal{C}_0$. Finally, $B^{0}(\Lambda_0)\subset\Omega_0^{\mathrm{cl}}$ and,
by \eqref{4.16:eq-cube-class-1} and the second and third inclusions,
\begin{equation*}
B^{1}(\Lambda_1)=X_1\subset\Omega_1\subset\Omega_0^{\mathrm{cl}},
\end{equation*}
so $\mathfrak{E}\subset\Omega_0^{\mathrm{cl}}$. The conditions of \eqref{4.16:eq-cube-class-a} at $j=0$
involve only $X_0$ and $X_1$, so membership in $\mathcal{C}_0$ does not require inclusion in
$\Omega_0^{\mathrm{cl}}$.

(d) Under this reading every class cube has exactly $12$ blocks a side. By (a) and the floor,
$12\le\mathsf{N}^{\sharp}$. Hence every class cube at index $1$, respectively at index $0$, has at
most $\mathsf{N}^{\sharp}$ blocks a side, and (b), respectively (c), places it in $\mathcal{C}_1$,
respectively $\mathcal{C}_0$.

(e) We have $L^j\eta=L^{j}L^{-k}=\ell_j$. Let $\mathfrak{E}$ have $\mathsf{N}$ blocks a side. Its side
is $\mathsf{N}M_1\ell_j$, which is $\mathsf{N}M_1$ in units of $L^j\eta=\ell_j$. Ba{\l}aban's $M$ is
our $M_1$, so the size written $O(1)M$ is $\mathsf{N}M_1$. The quoted bounds therefore read
$|A|<\mathsf{N}M_1\alpha_0\ell_j^{-1}$ and $|\nabla A|<\mathsf{N}M_1\alpha_0\ell_j^{-2}$ on
$\mathfrak{E}$. Multiplying the first by $\ell_j$ and the second by $\ell_j^{2}$ gives
\eqref{4.16:eq-cube-class-2}. The weight in these bounds is the fixed number $\ell_j$, whichever of
$X_j$, $X_{j+1}$ the position lies in.
\end{proof}

The constraint $\mathsf{N}\le\mathsf{N}^{\sharp}$ is what makes every $\mathfrak{E}\in\mathcal{C}_j$
fit into a single enlarged box $\square^{4}$ of half-side $5s_j$ of the regularity condition
$(\mathrm{reg})^{\sharp}$ (Definition~\ref{4.16:def-data-region}); this is the inclusion
\eqref{4.16:eq-cube-fitting-a}. Under the reading ``at least ten blocks a side'' of Ba{\l}aban's
$O(1)$, class cubes with more than $\mathsf{N}^{\sharp}$ blocks a side, if there are any, lie
outside $\mathcal{C}_j$, and they are not placed in a single box $\square^{4}$ by the fitting
\eqref{4.16:eq-cube-fitting-a}. They are not treated in this section, and no gluing of the gauges of
several boxes is attempted for them.

The two uses of the class named below are at $\mathsf{N}=10$. Under the floor
\eqref{4.16:eq-lowest-level-regularity-b}, (a) gives $\mathsf{N}^{\sharp}\ge 12$, so
$\mathcal{C}_j$ admits $\mathsf{N}=10$. In the
proof of part~(b) of the supremum proposition, hats (the cubes of side
$3\mathsf{m}_4\ell_n$ of that proof) are fitted into aligned ten-block cubes. The clause of the
definition of the class of hats that concerns the part of a hat lying in a closed union $Z$ of
level-$k$ $LM$-cubes involves no class cube itself; what it requires would be obtained in the same
way as in the proof of part~(b) of the supremum proposition, again with $\mathsf{N}=10$.

\begin{definition}[The weighted norm of perturbation and current]\label{4.16:def-weighted-norm}
Let $\mathfrak{g}^{\mathbb{C}}$ denote the complexified Lie algebra of $G$.
Let $A'$ be a $\mathfrak{g}^{\mathbb{C}}$-valued function on the bond set $\hat{\mathsf{B}}$ and $J$ a
$\mathfrak{g}^{\mathbb{C}}$-valued function on the bond set $\mathsf{B}\subset\hat{\mathsf{B}}$ of
Definition~\ref{4.16:def-data-region}; alternatively, $J$ may be given on the fine bonds carried by the
data region $\Omega_0^{\mathrm{dat}}$, as in the formulation of the operator theorem at free current.
Let $\nabla_U$ be the covariant difference quotient $\nabla^{\eta}_U$ of \cite[(3.37)]{B-CMP99b}, taken
in the background $U$ of Definition~\ref{4.16:def-data-region}, and let $a_1^{\oplus}$, $a_0^{\oplus}$ be
the radii of the space of the operator theorem at free current. The \emph{weighted norm} of the pair
$(A',J)$ is
\begin{equation}\label{4.16:eq-weighted-norm-a}
\|(A',J)\|_{\star}
:=\max\Bigl\{\,\sup\frac{\ell\,|A'|}{a_1^{\oplus}},\
\sup\frac{\ell^{2}\,|\nabla_U A'|}{a_1^{\oplus}},\
\sup\frac{\ell^{3}\,|J|}{a_0^{\oplus}}\,\Bigr\}.
\end{equation}
Each of the three members is read index by index, at the position of its index: the index is a fine bond
$b$, positioned at its midpoint $x_b$, or a pair $(b,b')$ of parallel neighbouring fine bonds, positioned at
$x_{(b,b')}$, as in Definition~\ref{4.16:def-fine-positions}. The weight attached to a fine index $\iota$
with position $x_\iota$ is the scale
\begin{equation*}
\ell_\iota:=\ell(x_\iota),
\end{equation*}
where $\ell(\cdot)$ is the pointwise weight, $\ell(x):=\min\{\ell_n:\ x\in X_n\}$, the least
$\ell_n$ over the shells $X_n$ that contain $x$. In \eqref{4.16:eq-weighted-norm-a} the first supremum
runs over the bonds of $\hat{\mathsf{B}}$, the second over the pairs of parallel neighbouring fine bonds
both of which belong to $\hat{\mathsf{B}}$, and the third over the bonds on which $J$ is given. The radius
parameter $r$ used together with this norm ranges over $(0,1]$.

A variant takes in each of the three members the supremum over the layers $n$, with the weight $\ell_n$
on all of $\Omega_n$, as in \cite[(3.37)]{B-CMP99b}, where the bounds are imposed on $\Omega_j$. This is
not a rewriting of \eqref{4.16:eq-weighted-norm-a}: on the boundary $\partial\Omega_n$, which lies in the
shell $X_{n-1}$, the pointwise weight is at most $\ell_{n-1}<\ell_n$, so that the layered bound on each
member is a stronger hypothesis on $\partial\Omega_n$ than the bound on the corresponding member of
\eqref{4.16:eq-weighted-norm-a}; clause~(b) of the statement below that derives the regularity bounds of
\cite[(3.37)]{B-CMP99b} from this norm holds under it a fortiori, without the powers of $L$.
\end{definition}

\begin{definition}[Tiles and unions of tiles]\label{4.16:def-tiles}
Let the level $k$, the lengths $\ell_n$ and the cube size $M$ of the nest be as in
Definition~\ref{4.16:def-data-region}, with lengths measured in level-$k$ units,
so that $\ell_k=1$. A \emph{tile} is a level-$k$ $LM$-cube, that is, a closed cube of side
$LM$ belonging to the standard partition of $\mathbb{T}$ into cubes of that side. The tiles
are exactly the cells of the partition on which the expansion of the correlation theorem
at step $k+1$ is built; these cells have side $LM$ in level-$k$ units.

In this section $Z$ denotes an arbitrary closed union of tiles. In what follows no property
of $Z$ is used other than the statement that a given object is carried by $Z$.
\end{definition}

\begin{proposition}[Lowest-level regularity of the initial data]\label{4.16:prop-lowest-level-regularity}
Let $\mathcal{C}_j$ be the bounded sub-class \eqref{4.16:eq-cube-class-a} of Ba{\l}aban's class of cubes
of \cite{B-CMP99b}, in the form fixed in Definition~\ref{4.16:def-cube-class}: the class cubes
of $\mathsf{N}$ big blocks a side with $\mathsf{N}\ge 10$ and $\mathsf{N}\le\mathsf{N}^{\sharp}$, where
$\mathsf{N}^{\sharp}=4M_1^{\sharp}/M_1$. Assume the following.
\begin{itemize}
\item The decreasing nest of domains $\Omega_0,\Omega_1,\dots,\Omega_k$ satisfies the admissibility
condition \textup{(A1)} of Proposition~\ref{4.14:prop-prepared-data-admissible}, with the shells $X_n$
and the pointwise
weight $\ell(\cdot)$ of Definition~\ref{4.16:def-weighted-norm}.
\item The block sizes, powers of $L$, satisfy $M_1\mid M_1^{\sharp}\mid M$, and the side of
$\mathbb{T}$ is a multiple of $LM$.
\item The hull condition $(\mathrm{col})$ of \eqref{4.16:eq-data-region-b} holds. Of it and of
the nesting of the domains only
$\Omega_1\subseteq\Omega_0\subset\Omega_0^{\mathrm{dat}}$ is used.
\item The floors
\begin{equation}\label{4.16:eq-lowest-level-regularity-a}
M\ge 13L^2M_1^{\sharp}
\end{equation}
and
\begin{equation}\label{4.16:eq-lowest-level-regularity-b}
M_1^{\sharp}\ge 3M_1
\end{equation}
hold; we refer to them as \textup{(F1)} and \textup{(F2)}, respectively. The first floor
\eqref{4.16:eq-lowest-level-regularity-a} is a floor of the operator theorem at
free current. It is used through $10M_1^{\sharp}<M$ in the proof below, and through
$5LM_1^{\sharp}<M$ and $20M_1<M$ in the statements that follow this proposition; each of these
three inequalities follows from \eqref{4.16:eq-lowest-level-regularity-a}, since
$M_1\le M_1^{\sharp}$ and $L\ge 3$. The second floor
\eqref{4.16:eq-lowest-level-regularity-b} is a floor on $M_1^{\sharp}$ imposed after $M_1$. Since
$M_1$ and $M_1^{\sharp}$ are powers of the odd integer $L>1$ with $M_1\mid M_1^{\sharp}$,
\eqref{4.16:eq-lowest-level-regularity-b} is equivalent to $M_1^{\sharp}>M_1$. It gives
$\mathsf{N}^{\sharp}=4M_1^{\sharp}/M_1\ge 12\ge 10$, so that $\mathcal{C}_j$ contains the cubes of ten
and of twelve big blocks a side used in \cite{B-CMP99b}.
\item In addition, the floor $M_1^{\sharp}\ge 5L\,M_1$ holds; we refer to it as
\textup{(F2$^{\sharp}$)}. It is a floor on $M_1^{\sharp}$ imposed after $M_1$ and $L$, of the same
kind as \eqref{4.16:eq-lowest-level-regularity-b} and as the floors on $M_1^{\sharp}$ of the operator
theorem at free current; it places no upper bound on any parameter and no condition on a coupling.
Since $M_1$ and $M_1^{\sharp}$ are powers of $L$, it asks that $M_1^{\sharp}/M_1$ be at least the
least power of $L$ that is $\ge 5L$. It gives $\mathsf{N}^{\sharp}\ge 20L$, the block count of the
seam-enlarged class of cubes treated in the corollary that follows this proposition; for the cubes
of that class at the lower seams, $M_1^{\sharp}\ge 3L\,M_1$, that is $\mathsf{N}^{\sharp}\ge 12L$,
already suffices. Clauses \textup{(a)} and \textup{(b)} below and the two facts established in the
proof do not use \textup{(F2$^{\sharp}$)}; they rest on the floors
\eqref{4.16:eq-lowest-level-regularity-a} and \eqref{4.16:eq-lowest-level-regularity-b} alone.
Whether the floors of the operator theorem at free current imply \textup{(F2$^{\sharp}$)} depends
on a constant of that theorem that is not evaluated, and is not asserted here.
\item $U$ is $G$-valued on the fine bonds carried by $\Omega_0^{\mathrm{dat}}$ and satisfies
$(\mathrm{reg})^{\sharp}$ in the form \eqref{4.16:eq-data-region-a}. The weighted bounds there are read
index by index as in \eqref{4.16:eq-fine-positions-a}, and the territories there are the shells $X_n$.
In particular the boxes of $\mathcal{D}^{\sharp}_0$ are centred at the $s_0$-lattice points of the whole
shell $X_0$, as part~(a) of the supremum proposition reads them.
\item $(A',J)$ satisfies $\|(A',J)\|_{\star}<r$, with the norm of \eqref{4.16:eq-weighted-norm-a}.
\item $Z$ is a union of tiles (Definition~\ref{4.16:def-tiles}), and $j\in\{0,1\}$, $j\le k$.
\end{itemize}
Then the following hold.

\smallskip
\textup{(a)} \emph{The bound \cite[(3.35)]{B-CMP99b} at index $j$, on $Z$, for class cubes of at most
$\mathsf{N}^{\sharp}$ big blocks a side.} For every $\mathfrak{E}\in\mathcal{C}_j$ there is a box
$\square\in\mathcal{D}^{\sharp}_j$ with $\mathfrak{E}\subset\square^{4}$. Let $u$ be the $G$-valued map
that $(\mathrm{reg})^{\sharp}$ gives on $\square^{4}\cap\Omega_0^{\mathrm{dat}}$, restricted to the
end-points of the fine bonds carried by $\mathfrak{E}\cap Z\cap\Omega_0^{\mathrm{dat}}$. Then
$U^{u}=e^{i\eta A}$ on those bonds, with
\begin{equation}\label{4.16:eq-lowest-level-regularity-c}
\ell_j|A(b)|<a^{\sharp},\qquad \ell_j^2|\nabla A(b,b')|<a^{\sharp}
\end{equation}
for every such bond $b$ and every parallel nearest-neighbour pair $(b,b')$ of such bonds.

This is \cite[(3.35)]{B-CMP99b} at index $j$, on the cubes of Ba{\l}aban's class with at most
$\mathsf{N}^{\sharp}$ big blocks a side. These are all the cubes of that class when the number written
$O(1)$ there is fixed equal to $12$. The constant written $O(1)M\alpha_0$ there, that is
$\mathsf{N}M_1\alpha_0$, is replaced by $a^{\sharp}$.

The bound can be sharpened. Under the index-by-index convention \eqref{4.16:eq-fine-positions-a}, a bond
or pair
positioned in $X_{j+1}$ obeys the same bounds with $\ell_{j+1}=L\ell_j$ in place of $\ell_j$.
Throughout $\mathfrak{E}\cap Z\cap\Omega_0^{\mathrm{dat}}$ one has, at the pointwise weight,
\begin{equation}\label{4.16:eq-lowest-level-regularity-1}
\ell(x_b)|A(b)|<a^{\sharp},\qquad \ell(x_{(b,b')})^2|\nabla A(b,b')|<a^{\sharp}.
\end{equation}

Every $\mathfrak{E}\in\mathcal{C}_1$ lies in $\Omega_0^{\mathrm{dat}}$, because
$\mathfrak{E}\subset X_1\cup X_2\subset\Omega_1$ and
$\Omega_1\subseteq\Omega_0\subset\Omega_0^{\mathrm{dat}}$. Hence at $j=1$ the intersection with
$\Omega_0^{\mathrm{dat}}$ removes nothing. At $j=0$ it removes nothing from any cube of Ba{\l}aban's class
as soon as $\Omega_0^{\mathrm{cl}}\subset\Omega_0^{\mathrm{dat}}$; this is the case when the data region
is taken, as in the operator theorem at free current, to be the hull of $\Omega_1$ of
\cite{B-CMP99b} on whose cube class \cite[(3.35)]{B-CMP99b} is imposed.

\smallskip
\textup{(b)} \emph{The bound \cite[(3.37)]{B-CMP99b} at index $j$, on $Z$.} Put $\mathsf{q}_j(x):=1$
if $j=1$ and
$x\in X_0$, and $\mathsf{q}_j(x):=0$ otherwise. On $\Omega_1$ one has $\mathsf{q}_1=1$ exactly on
$\Omega_1\cap X_0=\partial\Omega_1$. For every fine bond $b$ carried by $Z$ with $x_b\in\Omega_j$ (at
$j=0$: instead, every fine bond carried by $Z$ on which $A'$, respectively $J$, is given) and every pair
$(b,b')$ carried by $Z$ and positioned in $\Omega_j$,
\begin{equation}\label{4.16:eq-lowest-level-regularity-d}
\begin{aligned}
\ell_j|A'(b)|&<L^{\mathsf{q}_j(x_b)}\,r\,a_1^{\oplus}\le L\,r\,a_1^{\oplus},\\
\ell_j^2|(\nabla_UA')(b,b')|&<L^{2\mathsf{q}_j(x_{(b,b')})}\,r\,a_1^{\oplus}\le L^2\,r\,a_1^{\oplus},\\
\ell_j^3|J(b)|&<L^{3\mathsf{q}_j(x_b)}\,r\,a_0^{\oplus}\le L^3\,r\,a_0^{\oplus}.
\end{aligned}
\end{equation}
The powers of $L$ are incurred only at indices positioned on $\partial\Omega_1$, and only at $j=1$:
\begin{itemize}
\item for bonds, exactly the bonds crossing $\partial\Omega_1$ (one end-point in $\Omega_1$, one
outside), positioned at the crossing point;
\item for pairs, those positioned at the centre of a plaquette crossing $\partial\Omega_1$.
\end{itemize}
No factor is paid by a bond, pair or plaquette carried by $\Omega_j$, nor by any index at $j=0$.

In the present notation the condition \cite[(3.37)]{B-CMP99b} is
\begin{equation*}
|A'|<\alpha_1(L^j\eta)^{-1},\qquad |\nabla_UA'|<\alpha_1(L^j\eta)^{-2}\qquad\text{on }\Omega_j,
\quad j=0,\dots,k,
\end{equation*}
where $L^j\eta=\ell_j$. Clause~(b) is this bound at $j\in\{0,1\}$, with its one constant $\alpha_1$
given the value $L^2ra_1^{\oplus}$ for both members (using $L\le L^2$). On bonds and pairs carried
by $\Omega_j$, and at $j=0$, the replacement is $ra_1^{\oplus}$. The member for $J$ is a bound of the
kind of \cite[(1.14)]{B-CMP109}, namely $|J|<\gamma_0$ on a region; \cite{B-CMP99b} has no current as a
variable.

Nothing is asserted here for class cubes of more than $\mathsf{N}^{\sharp}$ big blocks a side.
\end{proposition}

\begin{proof}[Proof of Proposition~\ref{4.16:prop-lowest-level-regularity}: two facts]
We first establish two facts about the shells and the boxes of $(\mathrm{reg})^{\sharp}$.

\smallskip
\emph{Fact (i).} For every $j\in\{0,\dots,k\}$, $X_j$ is a union of cubes of side $s_j$ of the standard
partition, that is,
\begin{equation}\label{4.16:eq-lowest-level-regularity-e}
X_j=\bigcup\{\Delta:\ \Delta\ \text{a cube of side } s_j \text{ of the standard partition},\ 
\Delta\subset X_j\}.
\end{equation}
Suppose first $j\ge 1$. By the definition of the shells and condition (A1) of
Proposition~\ref{4.14:prop-prepared-data-admissible}, $X_j$ is a union of level-$j$
$M$-cubes. The standard partitions are nested, and $M_1^{\sharp}$, $M$ are powers of $L$ with
$M_1^{\sharp}\mid M$. So each
level-$j$ $M$-cube is a union of level-$j$ $M_1^{\sharp}$-cubes, whose side is
$M_1^{\sharp}\ell_j=s_j$. Now take $j=0$. By the same definition, $X_0$ is a union of level-$0$ $LM$-cubes.
Since $M_1^{\sharp}\mid LM$, each of these is a union of level-$0$ $M_1^{\sharp}$-cubes, that is, of
$s_0$-cubes. This proves \eqref{4.16:eq-lowest-level-regularity-e}.

\smallskip
\emph{Fact (ii).} Let $j\in\{0,1\}$, let $\square\in\mathcal{D}^{\sharp}_j$ have centre $y\in X_j$, and
write $\square^{4}=\{x:\ |x-y|_\infty\le 5s_j\}$. Then $\square^{4}$ meets $X_j$ and at most one further
shell, and that shell is adjacent to $X_j$; in particular
\begin{equation}\label{4.16:eq-lowest-level-regularity-f}
\#\{n:\ \square^{4}\cap X_n\neq\emptyset\}\le 2 .
\end{equation}

\emph{The size of $\square^4$.} The sup-diameter of $\square^{4}$ is $10s_j=10M_1^{\sharp}\ell_j$.
By \eqref{4.16:eq-lowest-level-regularity-a},
\begin{equation*}
M\ge 13L^2M_1^{\sharp}>10M_1^{\sharp},\qquad\text{so}\qquad 10s_j=10M_1^{\sharp}\ell_j<M\ell_j\le M.
\end{equation*}
The last inequality holds because $\ell_j=L^{j-k}\le 1$ for $j\le k$. The side of $\mathbb{T}$ is a
positive multiple of $LM$, hence at least $M$. So $\square^{4}$ does not wrap around the torus and is a
genuine cube of $\mathbb{T}$.

\emph{The case $j=0$.} Here $y\in X_0=\operatorname{cl}(\mathbb{T}\setminus\Omega_1)$. Suppose
$k\ge 2$. Condition (A1) at $n=1$ gives
$\operatorname{dist}_\infty(\Omega_1^c,\Omega_2)\ge M\ell_1$. By
\eqref{4.16:eq-lowest-level-regularity-a},
\begin{equation*}
M\ell_1=LM\eta>5M_1^{\sharp}\eta=5s_0.
\end{equation*}
Passing to the closure of $\Omega_1^c$ does not decrease this distance. So every $x\in\square^{4}$ has
$|x-y|_\infty\le 5s_0<\operatorname{dist}_\infty(y,\Omega_2)$, that is, $\square^{4}\cap\Omega_2=\emptyset$.
Every shell $X_n$ with $n\ge 2$ satisfies
\begin{equation*}
X_n=\operatorname{cl}(\Omega_n\setminus\Omega_{n+1})\subset\Omega_n\subset\Omega_2,
\end{equation*}
since $\Omega_n$ is a closed union of cubes (with $X_k=\Omega_k\subset\Omega_2$). Hence $\square^{4}$
meets $X_0$ and $X_1$ only. If $k=1$, there are no shells other than $X_0$ and $X_1$.

\emph{The case $j=1$.} If $k=1$, the only shells are $X_0$ and $X_1$ and there is nothing to prove. Let
$k\ge 2$. Then
\begin{equation*}
y\in X_1=\operatorname{cl}(\Omega_1\setminus\Omega_2)\subset\operatorname{cl}(\Omega_2^c).
\end{equation*}
First, suppose $k\ge 3$. Condition (A1) at $n=2$ gives
$\operatorname{dist}_\infty(y,\Omega_3)\ge M\ell_2$. By \eqref{4.16:eq-lowest-level-regularity-a},
\begin{equation*}
M\ell_2=LM\ell_1>5M_1^{\sharp}\ell_1=5s_1.
\end{equation*}
So $\square^{4}$ misses $\Omega_3$, and therefore misses every $X_n\subset\Omega_3$ with $n\ge 3$.

Second, $X_0\subset\operatorname{cl}(\Omega_1^c)$ and $X_2\subset\Omega_2$. Condition (A1) at $n=1$ (here
$k\ge 2$) shows that these two shells are at sup-distance at least $M\ell_1$, and by
\eqref{4.16:eq-lowest-level-regularity-a}
\begin{equation*}
M\ell_1>10M_1^{\sharp}\ell_1=10s_1,
\end{equation*}
the sup-diameter of $\square^{4}$. So $\square^{4}$ cannot meet both $X_0$ and $X_2$.

Since $y\in X_1$, it follows that $\square^{4}$ meets $X_1$ and at most one of $X_0$, $X_2$. This proves
\eqref{4.16:eq-lowest-level-regularity-f} together with the assertions accompanying it.

The form \eqref{4.16:eq-data-region-a} of $(\mathrm{reg})^{\sharp}$ states this property in the
words ``at most two, adjacent''. We have derived it here from (A1) and
\eqref{4.16:eq-lowest-level-regularity-a}, so that the index-by-index convention
\eqref{4.16:eq-fine-positions-a} of the weighted bounds applies to $\square^{4}$ without ambiguity.

Clauses \textup{(a)} and \textup{(b)} are derived from Facts~(i) and~(ii) in the lemma below on fitting
a class cube into one box, and in the two proofs that follow that lemma.
\end{proof}

\begin{lemma}[Fitting a class cube into one box]\label{4.16:prf-cube-fitting}
Assume the hypotheses of Proposition~\ref{4.16:prop-lowest-level-regularity}, and let
$j\in\{0,1\}$ with $j\le k$. Let $\mathfrak{E}\in\mathcal{C}_j$, the class of
\eqref{4.16:eq-cube-class-a}. Then there is a box $\square\in\mathcal{D}^{\sharp}_j$, with
centre $y\in X_j$, such that
\begin{equation}\label{4.16:eq-cube-fitting-a}
\mathfrak{E}\subset\square^{4},
\end{equation}
where $\square^{4}=\{x:\ |x-y|_\infty\le 5s_j\}$ is the closed enlargement of $\square$ of
half-side $5s_j$. At $j=0$ the boxes of $\mathcal{D}^{\sharp}_0$ are those centred at the
$s_0$-lattice points of the whole shell $X_0$, the level-$0$ territory of the regularity
hypothesis $(\mathrm{reg})^{\sharp}$ of \eqref{4.16:eq-data-region-a} being read as all of
$X_0$, as in part~(a) of the supremum proposition. For any other placement of the centres of
the boxes of $\mathcal{D}^{\sharp}_0$ (for instance at the $s_0$-lattice points of
$X_0\cap\Omega_0^{\mathrm{dat}}$ only) nothing is asserted at $j=0$.
\end{lemma}

\begin{proof}
Since $\mathfrak{E}\in\mathcal{C}_j$, the definition \eqref{4.16:eq-cube-class-a} gives
$\mathfrak{E}\cap X_j\neq\emptyset$; pick a point $z\in\mathfrak{E}\cap X_j$.

By the first of the two facts established with Proposition~\ref{4.16:prop-lowest-level-regularity},
\eqref{4.16:eq-lowest-level-regularity-e}, the shell $X_j$ is a union of $s_j$-cubes of the
standard partition. Hence some $s_j$-cube of the standard partition that is contained in $X_j$
contains $z$. (The point $z$ may lie on $\partial X_j$, where several $s_j$-cubes of the
partition contain it; since $z\in X_j$ and $X_j$ is a union of such cubes, one of them lies in
$X_j$.) Fix such a cube. The proof of part~(a) of the supremum proposition shows that the
vertices and the centre of this cube lie in $X_j$, and that there is a centre $y$ of a box
$\square\in\mathcal{D}^{\sharp}_j$
with $|y-z|_\infty\le s_j$. Since the boxes of $\mathcal{D}^{\sharp}_j$ are centred at points of
$X_j$, we have $y\in X_j$.

Let $x\in\mathfrak{E}$. Since $z\in\mathfrak{E}$ and $\mathfrak{E}$ is a cube of side
$\mathsf{N}M_1\ell_j$, we have $|x-z|_\infty\le\mathsf{N}M_1\ell_j$. The definition
\eqref{4.16:eq-cube-class-a} of $\mathcal{C}_j$ requires
$\mathsf{N}\le\mathsf{N}^{\sharp}=4M_1^{\sharp}/M_1$, so that
$\mathsf{N}M_1\ell_j\le 4M_1^{\sharp}\ell_j=4s_j$. By the triangle inequality,
\begin{equation}\label{4.16:eq-cube-fitting-1}
|x-y|_\infty\le|y-z|_\infty+|x-z|_\infty\le s_j+\mathsf{N}M_1\ell_j
\le s_j+4s_j=5s_j .
\end{equation}
Thus every point of $\mathfrak{E}$ lies in $\square^{4}$, which is
\eqref{4.16:eq-cube-fitting-a}.

The factor $4$ in the definition of $\mathsf{N}^{\sharp}$ is accounted for as follows.
Part~(a) of the supremum proposition fitted, about a point of a cell of level $n$, a cube of
side $3\mathsf{m}_4\ell_n$, where $\mathsf{m}_4$ is the power of $L$ fixed there as the
multiplier of its cubes, and for this needed
$M_1^{\sharp}\ge\mathsf{m}_4$. Here the point $z$ may sit at a corner of $\mathfrak{E}$; this is
why the full side $\mathsf{N}M_1\ell_j$ of $\mathfrak{E}$ enters
\eqref{4.16:eq-cube-fitting-1}, and with it the factor $4$. A cube of side larger than $4s_j$
fits in no $\square^{4}$ by this construction, which is why the class $\mathcal{C}_j$ stops at
$\mathsf{N}^{\sharp}$ blocks a side.

At $j=0$ the argument uses the placement of the centres of $\mathcal{D}^{\sharp}_0$ fixed in the
statement, at the $s_0$-lattice points of the whole shell $X_0$: this is the placement under
which part~(a) of the supremum proposition supplies the centre $y$. It agrees with the
definitions under which $(\mathrm{reg})^{\sharp}$ of \eqref{4.16:eq-data-region-a} is imposed,
namely $X_0$ the closure of $\mathbb{T}\setminus\Omega_1$ and the boxes of
$\mathcal{D}^{\sharp}_j$ centred at the $s_j$-lattice points of $X_j$. Under any other placement of
these centres the argument above does not supply the fitting at $j=0$.
\end{proof}

\begin{proof}[Proof of clause~(a) of Proposition~\ref{4.16:prop-lowest-level-regularity}]
Let $\mathfrak{E}\in\mathcal{C}_j$. The first assertion of clause~(a), the existence of a box
$\square\in\mathcal{D}^{\sharp}_j$ with $\mathfrak{E}\subset\square^{4}$, is
Lemma~\ref{4.16:prf-cube-fitting}; the centre $y$ of this box lies in $X_j$.

The regularity hypothesis $(\mathrm{reg})^{\sharp}$ of \eqref{4.16:eq-data-region-a}, applied
at the box $\square$, gives a $G$-valued $u$ on $\square^{4}\cap\Omega_0^{\mathrm{dat}}$ with
$U^{u}=e^{i\eta A}$ on the fine bonds carried by $\square^{4}\cap\Omega_0^{\mathrm{dat}}$. By the
second fact established with Proposition~\ref{4.16:prop-lowest-level-regularity},
\eqref{4.16:eq-lowest-level-regularity-f}, which applies because $j\in\{0,1\}$ and $y\in X_j$,
the cube $\square^{4}$ meets at most two shells, adjacent, one of them $X_j$. Reading the
weighted bounds of $(\mathrm{reg})^{\sharp}$ by positions as in
\eqref{4.16:eq-fine-positions-a}, we obtain, for each shell $X_n$ that $\square^{4}$ meets,
\begin{equation}\label{4.16:eq-cube-fitting-2}
\ell_n|A(b)|<a^{\sharp},\qquad \ell_n^{2}|\nabla A(b,b')|<a^{\sharp},
\end{equation}
the first for every bond $b$ carried by $\square^{4}\cap\Omega_0^{\mathrm{dat}}$ and positioned
in $X_n$, the second for every pair $(b,b')$ carried by $\square^{4}\cap\Omega_0^{\mathrm{dat}}$
and positioned in $X_n$.

We place the bonds and pairs carried by $\mathfrak{E}\cap\Omega_0^{\mathrm{dat}}$. Since
$\mathfrak{E}\cap\Omega_0^{\mathrm{dat}}\subset\square^{4}\cap\Omega_0^{\mathrm{dat}}$ by
\eqref{4.16:eq-cube-fitting-a}, a bond carried by $\mathfrak{E}\cap\Omega_0^{\mathrm{dat}}$ is
carried by $\square^{4}\cap\Omega_0^{\mathrm{dat}}$; its position, its midpoint, lies in the
convex cube $\mathfrak{E}$, hence in $X_j\cup X_{j+1}$ because $\mathfrak{E}\subset X_j\cup
X_{j+1}$ by \eqref{4.16:eq-cube-class-a}. A pair carried by
$\mathfrak{E}\cap\Omega_0^{\mathrm{dat}}$ is positioned at the midpoint of the two midpoints of
its bonds, that is, at the midpoint of two points of the convex cube $\mathfrak{E}$; this
position again lies in $\mathfrak{E}\subset X_j\cup X_{j+1}$. Finally, a shell that contains a
point of $\square^{4}$ is a shell that $\square^{4}$ meets. Hence each such bond or pair is
positioned in a shell $X_n$ with $n\in\{j,j+1\}$ that $\square^{4}$ meets, and it obeys the
$\ell_n$-weighted bound \eqref{4.16:eq-cube-fitting-2}.

Since $\ell_j\le\ell_n$ for $n\in\{j,j+1\}$, the $\ell_j$-weighted bounds of
\eqref{4.16:eq-lowest-level-regularity-c} follow. When $n=j+1$ they hold with a factor to spare:
as $\ell_{j+1}=L\ell_j$,
\begin{equation}\label{4.16:eq-cube-fitting-3}
\ell_j|A(b)|=L^{-1}\ell_{j+1}|A(b)|<L^{-1}a^{\sharp},\qquad
\ell_j^{2}|\nabla A(b,b')|<L^{-2}a^{\sharp};
\end{equation}
in particular a bond or pair positioned in $X_{j+1}$ obeys the same bounds with $\ell_{j+1}$ in
place of $\ell_j$, which is the sharper assertion of clause~(a). Since the pointwise weight
satisfies $\ell(x)\le\ell_n$ for $x\in X_n$, the pointwise bounds
\begin{equation}\label{4.16:eq-cube-fitting-4}
\ell(x_b)|A(b)|<a^{\sharp},\qquad \ell(x_{(b,b')})^{2}|\nabla A(b,b')|<a^{\sharp}
\end{equation}
follow from \eqref{4.16:eq-cube-fitting-2} as well. A position on the lower seam
$\mathfrak{E}\cap X_{j-1}\cap X_j$ lies in $X_j$ and costs nothing: both the $\ell_j$-weighted
and the pointwise bound hold there under the reading \eqref{4.16:eq-fine-positions-a}. On
$\mathfrak{E}\cap X_j$ the bound is the bound of \cite[(3.35)]{B-CMP99b} at its weight
$\ell_j$, with Ba{\l}aban's constant $\mathsf{N}M_1\alpha_0$ replaced by $a^{\sharp}$; on
$\mathfrak{E}\cap X_{j+1}$ it is stronger than that bound. No loss occurs anywhere.

Now restrict $u$ to the end-points of the fine bonds carried by
$\mathfrak{E}\cap Z\cap\Omega_0^{\mathrm{dat}}$. These bonds, and the parallel nearest-neighbour
pairs formed from them, are among the bonds and pairs carried by
$\mathfrak{E}\cap\Omega_0^{\mathrm{dat}}$ just bounded; so the restricted $u$ satisfies
$U^{u}=e^{i\eta A}$ on them, together with \eqref{4.16:eq-lowest-level-regularity-c}, its
sharper form at $\ell_{j+1}$ and the pointwise bounds \eqref{4.16:eq-cube-fitting-4}.

If the hypothesis $(\mathrm{reg})^{\sharp}$ is read as binding only those boxes whose enlargement
$\square^{4}$ carries a bond of the data region (a restriction of the same kind as the one under
which the regularity of the fine data is required only on those members of the supremum
proposition's cube family whose enlargements hold a bond of $\hat{\mathsf{B}}$), nothing
changes: when $\mathfrak{E}$ carries a bond of $\Omega_0^{\mathrm{dat}}$, so does
$\square^{4}\supset\mathfrak{E}$, and otherwise clause~(a) is vacuous.

Finally, the cubes of $\mathcal{C}_1$ lie in
\[
X_1\cup X_2\subset\Omega_1\subseteq\Omega_0\subset\Omega_0^{\mathrm{dat}},
\]
by the nesting of the domains and the hull condition \eqref{4.16:eq-data-region-b};
so at $j=1$ the intersection with $\Omega_0^{\mathrm{dat}}$ is idle. The cubes of Ba{\l}aban's
class at index $0$ lie in the class hull $\Omega_0^{\mathrm{cl}}$
(Definition~\ref{4.16:def-cube-class}), so at $j=0$ the intersection is idle for every such cube
as soon as $\Omega_0^{\mathrm{cl}}\subset\Omega_0^{\mathrm{dat}}$.
\end{proof}

\begin{lemma}[Boundary factors in the weighted bounds]\label{4.16:prf-boundary-factors}
This is clause (b) of Proposition~\ref{4.16:prop-lowest-level-regularity}. Let the data be as
there, with $\|(A',J)\|_{\star}<r$ for some $r\in(0,1]$, and let $Z$ be as there. For
$j\in\{0,1\}$ and $x\in\mathbb{T}$ put $\mathsf{q}_j(x):=1$ if $j=1$ and $x\in X_0$, and
$\mathsf{q}_j(x):=0$ otherwise. Then the three bounds \eqref{4.16:eq-lowest-level-regularity-d}
of that clause hold at every index carried by $Z$ and positioned in $\Omega_j$; for $j=0$ the
bounds on bonds hold at every fine bond carried by $Z$ on which $A'$, resp.\ $J$, is given.
The powers of $L$
in them are incurred only at indices positioned on $\partial\Omega_1$, and only for $j=1$. For
bonds these are exactly the bonds with one end-point in $\Omega_1$ and one outside. For pairs
these are exactly the pairs positioned at the centre of a plaquette that meets $\Omega_1$ without
being carried by it.
\end{lemma}

\begin{proof}
Indices are positioned as in Definition~\ref{4.16:def-fine-positions}. A fine bond $b$ sits at its
midpoint $x_b$. A pair $(b,b')$ sits at $x_{(b,b')}=\tfrac12(x_b+x_{b'})$. A plaquette sits at its
centre. The weight $\ell(\cdot)$ is the pointwise weight of
Definition~\ref{4.16:def-weighted-norm}, with $\ell_n=L^{n-k}$.

\emph{The weight on $\Omega_j$.} Let $j=0$. For every $n\ge0$ we have
$\ell_n=L^{n-k}\ge L^{-k}=\ell_0$. Since $\ell(x)$ is the least $\ell_n$ over the shells $X_n$
containing $x$, we get $\ell(x)\ge\ell_0=\eta$ at every point of $\mathbb{T}$.

Let $j=1$. The set $\Omega_1$ is closed, being a finite union of closed cubes, and
$X_0=\operatorname{cl}(\mathbb{T}\setminus\Omega_1)$. Hence
$\Omega_1\cap X_0=\Omega_1\cap\operatorname{cl}(\mathbb{T}\setminus\Omega_1)=\partial\Omega_1$.

Every point of $\Omega_1$ lies in some shell $X_n$ with $n\ge1$. Indeed, $\Omega_1$ is the union of
the sets $\Omega_n\setminus\Omega_{n+1}$ for $1\le n<k$ and of $\Omega_k$, and each of these is
contained in the corresponding shell. A point $x\in\Omega_1\setminus\partial\Omega_1$ does not lie
in $X_0$, because $X_0\setminus\partial\Omega_1\subset\mathbb{T}\setminus\Omega_1$. Since the shells
cover $\mathbb{T}$, such an $x$ lies only in shells $X_n$ with $n\ge1$.

Because $\ell_n=\ell_0$ only for $n=0$, it follows that on $\Omega_1$ we have
$\ell(x)\ge\ell_0=\ell_1/L$. Moreover $\ell(x)=\ell_1/L$ if and only if $x\in X_0$, that is, if and
only if $x\in\partial\Omega_1$; otherwise $\ell(x)\ge\ell_1$.

The same holds for a general $j\ge2$, with $\partial\Omega_j$ in place of $\partial\Omega_1$ and
$\ell_{j-1}=\ell_j/L$ in place of $\ell_0$. The reason is that $X_n\cap\Omega_j=\emptyset$ for
$n\le j-2$, by the positive distances in condition (A1) of \eqref{4.14:eq-prepared-data-admissible-a}.

For $j=0$ and every $x\in\mathbb{T}$, for $j=1$ and every $x\in\Omega_1$, and for
$p\in\{1,2,3\}$, these facts give
\begin{equation}\label{4.16:eq-boundary-factors-1}
\Bigl(\frac{\ell_j}{\ell(x)}\Bigr)^{p}\le L^{p\,\mathsf{q}_j(x)} .
\end{equation}
To check this, consider the cases in turn.
\begin{itemize}
\item For $j=0$ the left side is at most $1$ and $\mathsf{q}_0=0$.
\item For $j=1$ and $x\in\partial\Omega_1$ the left side equals $L^{p}$, and $\mathsf{q}_1(x)=1$
because $x\in X_0$.
\item For $j=1$ and $x\in\Omega_1\setminus\partial\Omega_1$ the left side is at most $1$, and
$\mathsf{q}_1(x)=0$ because $x\notin X_0$.
\end{itemize}

\emph{The bounds.} For an index $\iota$, write $F_\iota$ for the member attached to it, $p$ for its
weight power and $a_\iota$ for its radius:
\begin{itemize}
\item $F_\iota:=A'(b)$ for $\iota=b$, with $p=1$ and $a_\iota:=a_1^{\oplus}$;
\item $F_\iota:=(\nabla_UA')(b,b')$ for $\iota=(b,b')$, with $p=2$ and $a_\iota:=a_1^{\oplus}$;
\item $F_\iota:=J(b)$ for $\iota=b$, with $p=3$ and $a_\iota:=a_0^{\oplus}$.
\end{itemize}
The norm $\|(A',J)\|_{\star}$ of \eqref{4.16:eq-weighted-norm-a}
(Definition~\ref{4.16:def-weighted-norm}) is the largest of three suprema. Each member is read at
the position $x_\iota$ of its index with the weight $\ell(x_\iota)$. Hence
$\|(A',J)\|_{\star}<r$ gives $\ell(x_\iota)^{p}|F_\iota|<r\,a_\iota$ for every index $\iota$ of
each of the three kinds.

For $\iota$ positioned in $\Omega_j$, and for $j=0$ also at every bond index at which the member
is given, wherever it is positioned, multiply this by the positive number
$(\ell_j/\ell(x_\iota))^{p}$ and use \eqref{4.16:eq-boundary-factors-1}. This gives
\begin{equation}\label{4.16:eq-boundary-factors-2}
\begin{aligned}
\ell_j^{p}|F_\iota|
&=\Bigl(\frac{\ell_j}{\ell(x_\iota)}\Bigr)^{p}\,\ell(x_\iota)^{p}|F_\iota|\\
&<L^{p\,\mathsf{q}_j(x_\iota)}\,r\,a_\iota
\le L^{p}\,r\,a_\iota ,
\end{aligned}
\end{equation}
where the last step uses $\mathsf{q}_j\le1$ and $L>1$. Restricting to indices carried by $Z$
changes nothing. Read for the three kinds of index in turn, \eqref{4.16:eq-boundary-factors-2}
gives the three bounds \eqref{4.16:eq-lowest-level-regularity-d}.

\emph{Where the factor is paid ($j=1$).} By \eqref{4.16:eq-boundary-factors-2}, a factor is paid at
$\iota$ exactly when $x_\iota\in\Omega_1\cap X_0=\partial\Omega_1$.

We first note the coordinates involved.
\begin{itemize}
\item Sites have all coordinates in $\eta\mathbb{Z}+\tfrac12\eta$, being centres of the cells of
$\pi^{(0)}$, whose corners lie on $\eta\mathbb{Z}^d$.
\item By condition (A1), $\Omega_1$ is a union of level-$1$ $M$-cubes, that is, of cubes of side
$M\ell_1=LM\eta$ of the standard partition at that mesh, whose corners lie on $LM\eta\mathbb{Z}^d$.
Its faces therefore lie on hyperplanes $\{x_\mu=c\}$ with $c\in LM\eta\mathbb{Z}\subset\eta\mathbb{Z}$.
\end{itemize}

\emph{Bonds.} A bond $\langle x,x+\eta e_\mu\rangle$ has its midpoint with $\mu$-th coordinate in
$\eta\mathbb{Z}$ and all other coordinates in $\eta\mathbb{Z}+\tfrac12\eta$. The midpoint lies on
$\partial\Omega_1$ if and only if the bond crosses a face of $\Omega_1$ transversally at its
midpoint, that is, if and only if exactly one end-point lies in $\Omega_1$.

Indeed, the midpoint $x_b$ is the centre of the common face of the two cells centred at the
end-points, and it lies in exactly those two cells. Since $\Omega_1$ is a union of cells by (A1),
$x_b\in\Omega_1$ if and only if at least one of the two cells is a cell of $\Omega_1$. A cell is a
cell of $\Omega_1$ if and only if its centre lies in $\Omega_1$, so this happens if and only if at
least one end-point lies in $\Omega_1$.

Next, a small neighbourhood of $x_b$ lies in the union of the two cells. A cell that is not a cell
of $\Omega_1$ has interior disjoint from $\Omega_1$. Hence
$x_b\in\operatorname{cl}(\mathbb{T}\setminus\Omega_1)$ if and only if at least one of the two cells
is not a cell of $\Omega_1$. Combining the two statements, $x_b\in\partial\Omega_1$ if and only if
exactly one end-point lies in $\Omega_1$.

\emph{Pairs of non-collinear bonds.} Let $(b,b')$ be a pair with $b=\langle x,x+\eta e_\mu\rangle$
and $b'=b+\eta e_{\mu'}$, where $\mu'\ne\mu$. Its position is
\begin{equation*}
x_{(b,b')}=x+\tfrac12\eta e_\mu+\tfrac12\eta e_{\mu'} ,
\end{equation*}
which is the centre of the plaquette bounded by $b$ and $b'$. A plaquette centre has two
coordinates in $\eta\mathbb{Z}$, namely those of its two directions. It lies on $\partial\Omega_1$
only if the plaquette crosses a face of $\Omega_1$ normal to one of its directions.

The converse also holds. The centre lies in exactly the four cells centred at the plaquette's
corners, and the plaquette lies in their union. Hence the centre lies in
$\Omega_1\cap\operatorname{cl}(\mathbb{T}\setminus\Omega_1)=\partial\Omega_1$ if and only if
between one and three of the four cells are cells of $\Omega_1$. This happens if and only if the
plaquette meets $\Omega_1$ without being carried by it.

\emph{Collinear pairs.} For a collinear pair ($\mu'=\mu$) the position $x_{(b,b')}$ is the common
end-point of $b$ and $b'$, a site. A site has all coordinates in $\eta\mathbb{Z}+\tfrac12\eta$.
It therefore lies on none of the hyperplanes carrying the faces of $\Omega_1$, and so not on
$\partial\Omega_1$. A collinear pair pays nothing.

\emph{Conclusion.} A bond or a plaquette carried by $\Omega_1$ is positioned off
$\partial\Omega_1$ and pays no factor. At $j=0$ no factor is paid, since $\mathsf{q}_0=0$. The
factors $L$, $L^2$, $L^3$ are paid exactly at the indices described in the statement.

Suppose a pair is positioned at one of its two bonds rather than at the midpoint of their
midpoints. Then $\mathsf{q}_j$ is read at that bond, and the factor $L^2$ is paid at the pairs
containing a bond that crosses $\partial\Omega_1$. The form of the bounds is unchanged, since the
same position is read in the norm and in the bound.
\end{proof}

\begin{remark}[Placing a ceiling on the regularity constant]\label{4.16:rem-regularity-ceiling}
The constant in clause~(a) of Proposition~\ref{4.16:prop-lowest-level-regularity}, that is, the
constant bounding the two quantities in \eqref{4.16:eq-lowest-level-regularity-c}, is $a^{\sharp}$.
This is the regularity constant of the condition $(\mathrm{reg})^{\sharp}$ of
Definition~\ref{4.16:def-data-region}, chosen after $M$
under the ceilings on $a^{\sharp}$ recorded there.
\begin{enumerate}
\item[(1)] Suppose a statement uses the regularity of
Proposition~\ref{4.16:prop-lowest-level-regularity}
with a constant of its own, and requires that constant to lie below a radius. Through
Proposition~\ref{4.16:prop-lowest-level-regularity}, this is a ceiling on $a^{\sharp}$, of the same
kind as the ceilings under which $a^{\sharp}$ is placed, one-sided, and imposed after $M$.
\item[(2)] Matching the constant of \cite[(3.35)]{B-CMP99b} would require the condition
\begin{equation}\label{4.16:eq-regularity-ceiling-1}
a^{\sharp}\le \mathsf{N}M_1\alpha_0 .
\end{equation}
No statement of this chapter reads this condition, and
Proposition~\ref{4.16:prop-lowest-level-regularity}, stated with $a^{\sharp}$, does not impose it.
Condition \eqref{4.16:eq-regularity-ceiling-1} is one-sided only if $a^{\sharp}$ is chosen after
$\alpha_0$. If it is ever used, $a^{\sharp}$ is to be placed after $\alpha_0$ in the order of choice.
\end{enumerate}
\end{remark}

\begin{proof}
For (1): the bounds \eqref{4.16:eq-lowest-level-regularity-c} hold with right-hand side
$a^{\sharp}$, and so does every other bound concluded in clause~(a). A statement
that uses clause~(a) with a constant of its own can therefore take that constant to be
$a^{\sharp}$. The
requirement that its constant lie below a given radius is then the requirement that $a^{\sharp}$ lie
below that radius. This is an upper bound on $a^{\sharp}$ alone, imposed at the place of
$a^{\sharp}$ in the order of choice, after $M$; it is therefore a one-sided condition, of
the same kind as the ceilings on $a^{\sharp}$ recorded in
Definition~\ref{4.16:def-data-region}.

For (2): in \cite[(3.35)]{B-CMP99b} the regularity bounds are stated with the constant written
there $O(1)M\alpha_0$, in which $M$ is the block size of \cite{B-CMP99b}, that is, the present
$M_1$, and not the cube size $M$ of this chapter. On a class cube of $\mathsf{N}$ blocks a side the
factor $O(1)$ is $\mathsf{N}$, so this constant reads
$\mathsf{N}M_1\alpha_0$. In clause~(a) of Proposition~\ref{4.16:prop-lowest-level-regularity} it is
replaced by $a^{\sharp}$. Recovering the constant as it appears in \cite[(3.35)]{B-CMP99b} from
\eqref{4.16:eq-lowest-level-regularity-c} therefore requires \eqref{4.16:eq-regularity-ceiling-1}.

Suppose first that $a^{\sharp}$ is chosen after $\alpha_0$. Then $M_1$ and $\alpha_0$
are fixed when $a^{\sharp}$ is chosen, and for each class cube, with its block count $\mathsf{N}$,
\eqref{4.16:eq-regularity-ceiling-1} is an upper bound on $a^{\sharp}$, hence one-sided.

Suppose instead that $a^{\sharp}$ is chosen before $\alpha_0$. Then \eqref{4.16:eq-regularity-ceiling-1}
is equivalent to
\begin{equation*}
\alpha_0\ge \frac{a^{\sharp}}{\mathsf{N}M_1},
\end{equation*}
a lower bound on $\alpha_0$, and such a bound is not admitted; hence, if
\eqref{4.16:eq-regularity-ceiling-1} is ever used, $a^{\sharp}$ is placed after $\alpha_0$.
\end{proof}

\begin{remark}[All indices, zero steps and the direct route]\label{4.16:rem-direct-route}
We add three complements to the lowest-level regularity of the initial data
(Proposition~\ref{4.16:prop-lowest-level-regularity}).

\begin{enumerate}
\item[(1)] All indices. The proof of clause~(a) of
Proposition~\ref{4.16:prop-lowest-level-regularity}, that is, of the bounds
\eqref{4.16:eq-lowest-level-regularity-c}, serves every index $j$ with $0\le j\le k$, on the
class $\mathcal{C}_j$ of \eqref{4.16:eq-cube-class-a}. Indeed, the fitting
step in the proof of
Proposition~\ref{4.16:prop-lowest-level-regularity}, which gives $\mathfrak{E}\subset\square^{4}$
as in \eqref{4.16:eq-cube-fitting-a}, makes no use of the value of the index. The sentence of
part~(a) of the supremum proposition that the fitting step invokes is stated at every level. The
fact \eqref{4.16:eq-lowest-level-regularity-f}, that $\square^{4}$ meets at
most two shells, which are adjacent and one of which is $X_j$, holds at index $j$ as soon as
\begin{equation}\label{4.16:eq-direct-route-1}
5s_j<M\ell_{j-1}=\frac{M\ell_j}{L}
\qquad\text{and}\qquad
10s_j<M\ell_j .
\end{equation}
Both inequalities follow from the floor \eqref{4.16:eq-lowest-level-regularity-a}. Since
$s_j=M_1^{\sharp}\ell_j$, they read $5LM_1^{\sharp}<M$ and $10M_1^{\sharp}<M$, and both are
implied by $M\ge 13L^2M_1^{\sharp}$, because $13L^2\ge 13>10$ and $13L^2>5L$ for every blocking
factor $L$.

\item[(2)] Zero steps. Let $k=0$, so that, with the convention $\Omega_{k+1}=\emptyset$
of \cite{B-CMP99b}, $\Omega_1=\emptyset$. Then there is one shell,
$X_0=\mathbb{T}$, the weight is constant, $\ell\equiv\eta=1$, and only the index $j=0$
occurs. The two conditions in \eqref{4.16:eq-cube-class-a} hold for every cube, so
$\mathcal{C}_0$ is the set of aligned $\mathsf{N}$-block cubes with
$10\le\mathsf{N}\le\mathsf{N}^{\sharp}$. Ba{\l}aban's class at index~$0$
(Definition~\ref{4.16:def-cube-class}) is the set of aligned
cubes contained in
$\Omega_0^{\mathrm{cl}}$. Suppose the condition $(\mathrm{reg})^{\sharp}$ of
\eqref{4.16:eq-data-region-a} is imposed at $k=0$, with boxes $\mathcal{D}^{\sharp}_0$ on
$\mathbb{T}$. Then the fitting step in the proof of
Proposition~\ref{4.16:prop-lowest-level-regularity} applies word for word. The fact
\eqref{4.16:eq-lowest-level-regularity-f} is void, since there is only one shell. The bounds
\eqref{4.16:eq-lowest-level-regularity-d} of clause~(b) carry no power of $L$, because
$\mathsf{q}_0\equiv 0$.

\item[(3)] The direct route through hypothesis $(\mathrm{u}1)$. Hypothesis
$(\mathrm{u}1)$ of the operator theorem at free current is the bound
\cite[(3.35)]{B-CMP99b} itself. It is imposed at a constant $\alpha_0$ on the cube class of the
nest $\{\Omega_j\}_{j\le k}$ with block size $M_1$; the operator theorem at free current states it
on an enlargement of that cube class across the seams between adjacent shells, and it is read
there.
That theorem's real configuration space
$\mathcal{U}_{\mathbb{R}}$ is cut out by the three conditions $(\mathrm{u}1)$,
$(\mathrm{reg})$ and $(\mathrm{reg})^{\sharp}$.

On $\mathcal{U}_{\mathbb{R}}$, clause~(a) holds with Ba{\l}aban's own constant.
That is, for every index~$j$, every admitted $\mathsf{N}$ and every cube $\mathfrak{E}$ of
Ba{\l}aban's class at index~$j$ with $\mathsf{N}$ blocks a side, the gauge $u$ that hypothesis
$(\mathrm{u}1)$ furnishes on $\mathfrak{E}$ gives $U^{u}=e^{i\eta A}$ on the bonds carried by
$\mathfrak{E}$, with
\begin{equation}\label{4.16:eq-direct-route-2}
\ell_j|A(b)|<\mathsf{N}M_1\alpha_0,
\qquad
\ell_j^2|\nabla A(b,b')|<\mathsf{N}M_1\alpha_0 ,
\end{equation}
at the weight $\ell_j$, for the bonds $b$ and the parallel nearest-neighbour pairs $(b,b')$
carried by $\mathfrak{E}\cap Z$, with $Z$ the union of tiles of
Proposition~\ref{4.16:prop-lowest-level-regularity}. This is hypothesis $(\mathrm{u}1)$
restricted to the bonds
carried by $\mathfrak{E}\cap Z$. Being a restriction of a hypothesis, it needs no floor and no
bound on $\mathsf{N}$.

At the pointwise weight there is a cost. Every cube $\mathfrak{E}$ of Ba{\l}aban's class at
index~$j$ lies in $X_j\cup X_{j+1}$ (Definition~\ref{4.16:def-cube-class}); hence
$\ell(x)\le\ell_{j+1}=L\ell_j$ on $\mathfrak{E}$.
Hence \eqref{4.16:eq-direct-route-2} gives
\begin{equation}\label{4.16:eq-direct-route-3}
\ell(x_b)|A(b)|<L\,\mathsf{N}M_1\alpha_0,
\qquad
\ell(x_{(b,b')})^2|\nabla A(b,b')|<L^2\,\mathsf{N}M_1\alpha_0 .
\end{equation}
The factors $L$ and $L^2$ are paid only on $\mathfrak{E}\cap(X_{j+1}\setminus X_j)$, since
$\ell\le\ell_j$ on $X_j$. At $\mathsf{N}=10$ both right-hand sides in
\eqref{4.16:eq-direct-route-3} are at most $10L^2M_1\alpha_0$. This is the form of the bound in
part~(b) of the supremum proposition:
\begin{equation*}
\ell|A|<10L^2M_1\alpha_0,\qquad \ell^2|\nabla A|<10L^2M_1\alpha_0 .
\end{equation*}
At index $0$, the bound \eqref{4.16:eq-direct-route-2} lives on the class hull
$\Omega_0^{\mathrm{cl}}$. The relation of this hull to the data region
$\Omega_0^{\mathrm{dat}}$ is fixed by the hypotheses of the operator theorem at free current and
is not re-examined here.

Consequently, at $j=1$ the bound \cite[(3.35)]{B-CMP99b} is obtained from either of two inputs:
from $(\mathrm{reg})^{\sharp}$, on the cubes of $\mathcal{C}_1$ (the route of
Proposition~\ref{4.16:prop-lowest-level-regularity}), or from hypothesis $(\mathrm{u}1)$, on
Ba{\l}aban's whole class (the direct route). Neither route asks any
regularity of $U$ beyond its own input.

Where it applies, the route of Proposition~\ref{4.16:prop-lowest-level-regularity} is the
sharper of the two, for two reasons. First, it gives pointwise weights with no factor of~$L$.
Second, the gauge it uses is postulated on the trace $\square^{4}\cap\Omega_0^{\mathrm{dat}}$
itself, re-entrant corners of the data region included. In particular, the cubes of side
$3\mathsf{m}_4\ell_n$ at a level $n$, the hats of the supremum proposition, that meet $X_0$ near
the boundary of the data region, re-entrant corners
included, are served in the supremum proposition by its part~(a), through the boxes of
$(\mathrm{reg})^{\sharp}$ as in Proposition~\ref{4.16:prop-lowest-level-regularity}, and not by
the bound of its part~(b). For these reasons it is the route stated in
Proposition~\ref{4.16:prop-lowest-level-regularity}.

For class cubes of more than $\mathsf{N}^{\sharp}$ blocks a side, should such cubes be admitted,
only the direct route is given here: for them the bounds \eqref{4.16:eq-direct-route-2} are
hypothesis $(\mathrm{u}1)$ restricted to the bonds carried by $\mathfrak{E}\cap Z$, and nothing is
derived from $(\mathrm{reg})^{\sharp}$, since no such cube fits in a single box $\square^{4}$ by
the construction of the fitting step in the proof of
Proposition~\ref{4.16:prop-lowest-level-regularity}. Hypothesis $(\mathrm{u}1)$ is
an assumption of the operator theorem at free current.
\end{enumerate}
\end{remark}

\begin{remark}[Class cubes too large for one box. Proof outstanding]\label{4.16:rem-large-class-cubes}
Let $0\le j\le k$, let $Z$ be a closed union of tiles, and let $\mathfrak{E}$ be a cube of index $j$
in Ba{\l}aban's cube class (Definition~\ref{4.16:def-cube-class}) with $\mathsf{N}$ big blocks a side, where
\begin{equation*}
\mathsf{N}>\mathsf{N}^{\sharp}=4M_1^{\sharp}/M_1 .
\end{equation*}
Such a cube lies outside the sub-class $\mathcal{C}_j$ of \eqref{4.16:eq-cube-class-a}.
The fitting step \eqref{4.16:eq-cube-fitting-a} in the proof of
Proposition~\ref{4.16:prop-lowest-level-regularity} does not place such a cube in a single box
$\square^{4}$ of the box-regularity hypothesis $(\mathrm{reg})^{\sharp}$ of
\eqref{4.16:eq-data-region-a}, and no gauges of different boxes are glued.
The claim is that, under the hypotheses of
Proposition~\ref{4.16:prop-lowest-level-regularity}, clause (a) of that proposition holds on
$\mathfrak{E}$ with Ba{\l}aban's constant, that is,
\begin{equation*}
\ell_j|A|<\mathsf{N}M_1\alpha_0,\qquad \ell_j^2|\nabla A|<\mathsf{N}M_1\alpha_0
\end{equation*}
on the bonds carried by $\mathfrak{E}\cap Z$, at the weight $\ell_j$.

Proof outstanding.

\emph{Route.}
One way is to read the bounds as a restriction of a hypothesis.
Besides $(\mathrm{reg})$ and $(\mathrm{reg})^{\sharp}$, the operator theorem at free current has a
hypothesis $(\mathrm{u}1)$: the bound \cite[(3.35)]{B-CMP99b} at the constant $\alpha_0$, at block
size $M_1$, on Ba{\l}aban's cube class of the nest $\{\Omega_j\}_{j\le k}$, taken in the enlarged form
that theorem prescribes near the boundaries $\partial\Omega_j$ (called below the seam-enlarged
class of $(\mathrm{u}1)$); the cubes of that class are unions of big blocks and have at most $20L$
big blocks a side of the finest level they meet.
On the real configuration space $\mathcal{U}_{\mathbb{R}}^{\mathrm{OP}}$ of that theorem, whose
constraints are $(\mathrm{u}1)\wedge(\mathrm{reg})\wedge(\mathrm{reg})^{\sharp}$, for every cube
$\mathfrak{E}$ that the seam-enlarged class of $(\mathrm{u}1)$ admits, the displayed
bounds are exactly $(\mathrm{u}1)$ restricted to the bonds carried by $\mathfrak{E}\cap Z$, and
nothing is to be proved; this reading needs no floor and no bound on $\mathsf{N}$.
At the pointwise weight it costs the factors $L$ and $L^2$ on
$\mathfrak{E}\cap(X_{j+1}\setminus X_j)$, since the pointwise weight satisfies
$\ell(x)\le\ell_{j+1}=L\ell_j$ for $x\in\mathfrak{E}\subset X_j\cup X_{j+1}$.
At index $0$ the bounds live on the class hull $\Omega_0^{\mathrm{cl}}$; they concern the part of
$\mathfrak{E}$ inside $\Omega_0^{\mathrm{dat}}$, all of $\mathfrak{E}$ when
$\Omega_0^{\mathrm{cl}}\subset\Omega_0^{\mathrm{dat}}$ as in the hypotheses of the operator theorem
at free current, and nothing is asserted on $\mathfrak{E}\setminus\Omega_0^{\mathrm{dat}}$, where $U$
is not given; this is consistent with \cite[\S3]{B-CMP99b}, whose regularity conditions impose no
constraint on $U$ outside the domain written $\Omega_0$ there (here $\Omega_0^{\mathrm{cl}}$).
What remains on this route is whether the seam-enlarged class of $(\mathrm{u}1)$, whose cubes have
at most $20L$ big blocks a side, admits a cube with $\mathsf{N}>\mathsf{N}^{\sharp}$, and the supply
of hypothesis $(\mathrm{u}1)$ itself on such a cube.
The other way is to show that every cube on which clause (a) is applied lies in $\mathcal{C}_j$.
Under the floor $(\mathrm{F}2)$, $M_1^{\sharp}\ge 3M_1$, of
Proposition~\ref{4.16:prop-lowest-level-regularity} one has $\mathsf{N}^{\sharp}\ge12$.
So, by \eqref{4.16:eq-cube-class-a}, the ten-block cubes of Ba{\l}aban's class of index $j$ lie in
$\mathcal{C}_j$.
Every use of clause (a) other than through hypothesis $(\mathrm{u}1)$ of the operator
theorem at free current is on ten-block class cubes, hence inside $\mathcal{C}_j$.
For the cubes of the seam-enlarged class of $(\mathrm{u}1)$, impose the further one-sided floor
$M_1^{\sharp}\ge 5L\cdot M_1$, on $M_1^{\sharp}$ after $L$ and $M_1$; then
$\mathsf{N}^{\sharp}\ge 20L$.
Hence every cube of that class with at least ten big blocks a side lies in $\mathcal{C}_i$, where $i$
is the finest level it meets, every cube of that class, whatever its block count, is placed in one
box $\square^{4}$ by the fitting step \eqref{4.16:eq-cube-fitting-a} verbatim, and
Proposition~\ref{4.16:prop-lowest-level-regularity} and
Remark~\ref{4.16:rem-direct-route} apply to them unchanged.
Under that floor, therefore, no cube of the seam-enlarged class of $(\mathrm{u}1)$ is a cube of the
present remark; without it, from $(\mathrm{F}1)$ and $(\mathrm{F}2)$ alone, nothing is obtained here
for the cubes of that class with $12L$ to $20L$ big blocks a side.
\end{remark}

\begin{remark}[Transfer of the one-step analysis of \cite{B-CMP109} (its Lemmas~5,~6) to the cover of
{\cite[\S3]{B-CMP119}}; printed without proof in {\cite[p.~276]{B-CMP119}}, cited as printed]
\label{4.17:rem-transfer}
The following sentence is stated in \cite[p.~276]{B-CMP119} without proof; cited as printed. In it,
$\{\square\}$ is the cover of \cite[\S3]{B-CMP119}, $\{\zeta_\square\}$ is the partition of
unity associated with it in \cite[\S3]{B-CMP119}, and Ba{\l}aban's reference [I] is
\cite{B-CMP109}. His equation
numbers (I.6.34) and (I.6.41) therefore denote \cite[(6.34)]{B-CMP109} and
\cite[(6.41)]{B-CMP109}, and his ``Lemmas 5, 6 [I]'' are \cite[Lemmas~5 and~6]{B-CMP109}.
The ellipsis marks text of the same page omitted here. The expression
$\mathbf{P}^{(k)}$ and the letters $Y$, $\Lambda_{k+1}$, $\Lambda^{(k)*}_{k+1}$,
$\tilde U^c_{k+1}(Y,\cdot,\cdot)$, $\beta$, $\tilde\alpha_0$, $\tilde\alpha_1$, $C_1$, $p_1$
and $A$ are those of \cite[\S3]{B-CMP119}.
\begin{quotation}
We repeat the analysis of Sect.~3 [I] using the above cover $\{\square\}$, and the partition
of unity $\{\zeta_\square\}$. \dots{} The same remarks apply to the expression
$\mathbf{P}^{(k)}$, and we obtain Lemmas 5, 6 [I] with the corresponding changes. More
precisely, the terms of the expansion (I.6.41), with localization domains $Y$ contained in
$\Lambda_{k+1}$, satisfy exactly all the conditions of those lemmas. The remaining terms
satisfy such conditions also, but with analyticity space (I.6.34) replaced by the space
\begin{multline*}
\tilde U^c_{k+1}\bigl(Y,(1+\beta)\tilde\alpha_0,(1+\beta)\tilde\alpha_1\bigr)\\
\times\bigl\{A : \operatorname{supp} A\subset Y\cap\Lambda^{(k)*}_{k+1},\
|A|<C_1p_1(g_k)\bigr\}. \tag*{(3.43)}
\end{multline*}
\end{quotation}

The passage declares that the one-step analysis of \cite{B-CMP109}, which begins with
\cite[\S3]{B-CMP109} and ends in \cite[Lemmas~5 and~6]{B-CMP109}, carries over to the cover
$\{\square\}$ of \cite[\S3]{B-CMP119}: the terms of the expansion \cite[(6.41)]{B-CMP109}
whose localization domains lie in $\Lambda_{k+1}$ satisfy the conditions of those lemmas
exactly, and the remaining terms satisfy them with the analyticity space
\cite[(6.34)]{B-CMP109} replaced by the space \cite[(3.43)]{B-CMP119} displayed above.

In this chapter the passage is used as an input of the induction that proves the
inductive representation of the effective densities (Theorem~\ref{4.20:thm-representation}),
at the seam between consecutive layers of the cover. It is also an input of clause~(0)
(Theorem~\ref{4.1:thm-clause-zero}), through the bound on the large-field part of the
coupling increment.

It is one of the sentences of Ba{\l}aban's papers that this book uses as printed without
proof; all of them are listed together in the index that closes Part~III.
\end{remark}

\begin{theorem}[The transfer at the seam: Ba{\l}aban's one-step lemmas re-obtained at step $k+1$
on the layered cover of \cite{B-CMP119}]\label{4.17:thm-transfer-seam}
\emph{Setting.} Fix the level $k$. Measure lengths so that the lattice of the $k$-th
fluctuation field has spacing one. For $0\le n\le k+1$ put $\ell_n=L^{n-k}$, and let $\pi_n$ be
the partition of the torus $\mathbb{T}_m$ into closed aligned cubes of side $M\ell_n$; these
partitions are nested. Let
\begin{equation*}
\mathbb{T}_m=\Omega_0\supset\Omega_1\supset\dots\supset\Omega_{k+1}
\end{equation*}
and $\{\Lambda_j\}_{j\le k+1}$ be the sequences of domains of \cite{B-CMP119} at step $k+1$;
they are closed unions of lattice cubes. Let $U_{k+1}$ be the background configuration. For
$1\le n\le k$ write $\theta_n$ for the distance in the maximum norm from $\Omega_{n+1}$ to
$\mathbb{T}_m\setminus\Omega_n$.

Consider one term of the expansion \cite[(3.1)]{B-CMP119} with these data, and in it the last
logarithm of \cite[(3.28)]{B-CMP119}. Formally, this logarithm is given by the formula of the
term $E^{(k+1)}$ of \cite{B-CMP109}, except that the integral is restricted to $\Lambda_{k+1}$
and the operators are determined by the sequence $\{\Omega_j\}$. The part of its exponent that
depends on the fluctuation field $A$, with $\operatorname{supp}A\subset\Lambda^{(k)*}_{k+1}$, is
\begin{equation}\label{4.17:eq-transfer-seam-1}
\begin{split}
\mathcal{E}(A)={}&\mathbf{P}^{(k)}(g_k,U_{k+1},A)\\
&+\bigl\{E_k\bigl(U_k(e^{iB'}V^{(k)}_*)\bigr)-E_k(U_{k+1})\bigr\},
\end{split}
\end{equation}
where
\begin{equation}\label{4.17:eq-transfer-seam-2}
B'=g_kCA-hD(g_kCA).
\end{equation}
Here $C$, $V^{(k)}_*$ and the expression $hD(\cdot)$ are as in \cite[(3.2)]{B-CMP109}, and
$U_k(\cdot)$ denotes, as in \cite[\S3]{B-CMP109}, the level-$k$ configuration as a function of
its argument, here evaluated at $e^{iB'}V^{(k)}_*$. The expression $\mathbf{P}^{(k)}$ is the one
occurring in \cite[(3.28)]{B-CMP119}. Moreover $E_k=\sum_{j\le k}E^{(j)}$, where $E^{(j)}$ is the
expression produced at the $j$-th step, whose terms satisfy \cite[(2.26)--(2.27)]{B-CMP119}.
Finally, $\Lambda^{(k)*}_{k+1}$ is the set appearing in \cite[(3.43)]{B-CMP119} that carries the
support of the fluctuation field.

The operators $H_0$, $H$, $G$, $H_1$, $H_k$ are those of \cite{B-CMP102b}; in this theorem the
letters $H$ and $G$ denote these operators only. They are constructed, at the background
$U_{k+1}$, for the sequence of domains $\{\Omega_j\}_{j\le k}$, whose innermost domain is
$\Omega_k$, with the random walk expansions of \cite[Theorems~3.7, 3.9--3.11 and
Corollary~3.8]{B-CMP99b}. The
first expansion is \cite[(3.4)]{B-CMP109}. It is taken over the layered cover $\{\square\}$ and
the partition of unity $\{\zeta_\square\}$ of \cite[pp.~275--276]{B-CMP119}, with, at each cube,
the per-cube inputs listed with the proof. Each cube
$\square$ of this cover is a union of $2^d$ cubes of $\pi_s$, where $s=s(\square)$ is its
scale.

\emph{Hypotheses.}
\begin{enumerate}
\item Nesting: $\Omega_{j+1}\subset\Lambda_j\subset\Omega_j$ for each $j\le k$, as in
\cite[(2.1)]{B-CMP119}.
\item Grid: each $\Omega_n$ is a union of cubes of $\pi_n$.
\item Thickness: $\theta_n\ge 2M\ell_n$ for $1\le n\le k$, and $\theta_n\ge LM\ell_n$ for
$1\le n\le k-1$. In addition, $\theta_n>(7+7/L)M\ell_n$ for $1\le n\le k$; this bound is used
only in the construction of the layered cover.
\item The set $\Lambda_{k+1}$ is a union of cubes of $\pi_k$ and is contained in
$\Omega_{k+1}$. Its distance in the maximum norm from the complement of $\Omega_{k+1}$ satisfies
\begin{equation*}
\operatorname{dist}_\infty\bigl(\Lambda_{k+1},\mathbb{T}_m\setminus\Omega_{k+1}\bigr)
\ge LM\mathsf{R}_{k+1},
\end{equation*}
as in \cite[(3.20)]{B-CMP119}; here $\mathsf{R}_{k+1}$, a power of $L$, is the scale factor of
\cite{B-CMP119} attached to the $(k+1)$-th step.
\item The layered cover obeys the cover rule, fixed with the proof, that assigns to each cube
$\square$ of the cover its scale $s(\square)$.
\item Growth of the radii: there are fixed constants $\bar A$ and $\bar p$, independent of
$k$, $j$, $v$ and of the couplings, with
\begin{equation*}
\alpha_{\cdot,j+v}\le \bar A(1+v)^{\bar p}\,\alpha_{\cdot,j}
\qquad\text{for all } j \text{ and all } v\ge0.
\end{equation*}
Here $\alpha_{\cdot,j}$ stands for each of the radii $\alpha_{0,j}$, $\alpha_{1,j}$ and of the
constants $\varepsilon_j$ of \cite[\S2]{B-CMP119}. The constants in the conclusion may depend on
$\bar A$ and $\bar p$.
\item Consider an index $n$ at which $\Omega_n$ is determined by the large-field operation
$\mathbf{R}$ and not by the $n$-th renormalization transformation. At every such index,
$\Omega_n$ is a union of cubes of $\pi_n$, and $\theta_n\ge LM\ell_n$.
\item The residual printed inputs listed with the proof.
\end{enumerate}

\emph{Conclusion.} The conclusions of \cite[Lemmas~1 and~2]{B-CMP116} hold for this term with
the following changes. (These are Lemmas~5 and~6 in the continuous numbering of \cite{B-CMP109}
and \cite{B-CMP116} used in \cite{B-CMP119}.) Consider the terms of the expansion
\cite[(1.41)]{B-CMP116}. Those whose localization domains $Y$ are contained in $\Lambda_{k+1}$
satisfy exactly all the conditions of those lemmas. The remaining terms satisfy them as well,
but with the analyticity space \cite[(1.34)]{B-CMP116} replaced by the space
\cite[(3.43)]{B-CMP119}:
\begin{equation}\label{4.17:eq-transfer-seam-3}
\begin{split}
&\tilde{\mathcal{U}}^{c}_{k+1}\bigl(Y,(1+\beta)\tilde\alpha_0,(1+\beta)\tilde\alpha_1\bigr)\\
&\qquad\times\bigl\{A:\ \operatorname{supp}A\subset Y\cap\Lambda^{(k)*}_{k+1},\
|A|<C_1p_1(g_k)\bigr\}.
\end{split}
\end{equation}
Here $\tilde{\mathcal{U}}^{c}_{k+1}$ is the layered complex analyticity space of
\cite[(2.34)--(2.39)]{B-CMP119}, read at the top index $k+1$, with radii
$\tilde\alpha_0,\tilde\alpha_1$; $\beta$ is the schedule constant of
Definition~\ref{4.4:def-post-step-space}; and $C_1$ and $p_1(\cdot)$ are the constant and the
function of the coupling of \cite[(3.43)]{B-CMP119}.
\end{theorem}

The proof is given later in this section. It derives the conclusion as an implication
from the hypotheses above and from published statements cited by statement.

\begin{lemma}[The seam lemma for the one family of cells at all traces]\label{4.17:lem-seam}
Work in the setting of the one-family operator theorem. Its hypotheses are in force, and its
number set is chosen in its order of choice at the seam constants, that is, at one choice of the
values of its auxiliary numbers, subject to the following two readings.
\begin{enumerate}
\item[(a)] The cube-size exponents $\mathsf{m}_2,\mathsf{m}_4$ are chosen under the standing
floor $\mathsf{F}$ on the cube-size exponents, under which all four are chosen in the order of
choice. The exponents
$\mathsf{m}_1,\mathsf{m}_3$ are chosen under $\mathsf{F}$ and, in addition, under one further
floor $\mathsf{F}^\sharp_{\mathrm{sm}}$, a one-sided lower bound on $\mathsf{m}_1$ and
$\mathsf{m}_3$.
\item[(b)] The radius ceilings of the one-family operator theorem are re-read at the seam
constants. In particular $a_1^\sharp$, $a_0^\sharp$ and $a^\oplus$ denote the re-read values,
each an upper bound on a field or source radius fixed before $M$ and before the coupling.
The common ceiling $B^\sharp_{\mathrm{sum}}$ of the sum-clause constants of the operators of the
list is read at the exponents $\mathsf{m}_i$ of~(a).
\end{enumerate}
Let $\mathsf{K}$ be an operator of the list of the one-family operator theorem. For a walk
$\omega$, of any trace, of the walk system of $\mathsf{K}$ (Definition~\ref{4.4:def-walk-systems}),
write
$X_-(\omega)$ and $X_+(\omega)$ for its two footprints and $\#\omega$ for its number of steps.
Let $(\kappa_\omega)_\omega$ denote the weights of the norm clause of the one-family operator
theorem for $\mathsf{K}$, read at the same exponents $\mathsf{m}_i$ and at the single choice of
weighting made when the seam constants are fixed. Then there is a family of weights
$(\kappa^{\mathrm{sm}}_\omega)_\omega$, the seam-typed weights, indexed by the walks of
$\mathsf{K}$, fixed together with the number set at the seam constants and independent of the
datum, with $\kappa^{\mathrm{sm}}_\omega=\kappa_\omega$ for every coarse-valued operator
$\mathsf{K}$, such that the following hold. In~(i) and~(ii), $\omega$ is any walk, of any trace,
of the walk system of $\mathsf{K}$; $\mathsf{u}$ is any datum in
$\mathcal{U}^{\flat\nabla}_{[\omega]}(a_1^\sharp,a_0^\sharp)$, the data class of $\omega$ at the
re-read radii; and $\mathsf{K}_\omega=\mathsf{K}_\omega(\mathsf{u})$ is the term of $\omega$
evaluated at $\mathsf{u}$.
\begin{enumerate}
\item[(i)] $\kappa^{\mathrm{sm}}_\omega\ge\kappa_\omega$.
\item[(ii)] The weights $\kappa^{\mathrm{sm}}_\omega$ are valid $\ell^2$-weights for the norm
clause, that is,
\begin{equation}\label{4.17:eq-seam-1}
\|\mathsf{K}_\omega\|\le\kappa^{\mathrm{sm}}_\omega\,\mathsf{l}(X_-(\omega))^{\eta_{\mathsf{K}}}.
\end{equation}
Here $\|\cdot\|$ is the $\ell^2$ operator norm, and $\mathsf{l}$ and $\eta_{\mathsf{K}}$ are the
size function and the exponent of that clause.
\item[(iii)] Let $\Xi$ be the weight, per step of a walk, at which the statement of the
one-family operator theorem for $\mathsf{K}$ is formed; $\Xi$ is at least the rung
$\Xi_{\mathrm{r}}$ fixed for $\mathsf{K}$ in that theorem.
Then for every block $\Delta$ of the family $\mathfrak{Y}$ of blocks of the one-family operator
theorem, and for each of the two footprints separately,
\begin{equation}\label{4.17:eq-seam-2}
\sum\bigl\{\kappa^{\mathrm{sm}}_\omega\,\Xi^{\#\omega}:\ X_{\mp}(\omega)\cap\Delta\neq\emptyset\bigr\}
\le B^{\mathrm{sm}}_{\mathsf{K}}\le B^\sharp_{\mathrm{sum}} .
\end{equation}
Here $B^{\mathrm{sm}}_{\mathsf{K}}$ is the constant of the sum clause for $\mathsf{K}$ at the
seam-typed weights.
\end{enumerate}
Thus one and the same family of weights carries both the norm clause and the sum clauses of the
one-family operator theorem.
\end{lemma}

Two inputs stand beside Lemma~\ref{4.17:lem-seam}.

The first input is the inclusion of data classes
\begin{equation}\label{4.17:eq-seam-3}
\mathcal{U}(1)\subset\mathcal{U}^{\flat\nabla}(a_1^\sharp,a_0^\sharp).
\end{equation}
Here $\mathcal{U}(1)$ is the data class on which the statements that apply the one-family
operator theorem are formulated, and $\mathcal{U}^{\flat\nabla}(a_1^\sharp,a_0^\sharp)$ is the data
class of the same kind as in Lemma~\ref{4.17:lem-seam}, at the same radii, without the index of
a walk; both classes are defined through the radius $a^\oplus$. The inclusion is taken
by statement from the inclusion lemma for data classes, read at the re-read value of
$a^\oplus$. It is the first link of the chain of inclusions through which those statements are
applied, and it bears on that chain, not on
Lemma~\ref{4.17:lem-seam}.

The second input is the floor $\mathsf{F}^\sharp_{\mathrm{sm}}$ of~(a).

A data class becomes smaller when its radii decrease. The domain
$\mathcal{U}^{\flat\nabla}_{[\omega]}(a_1^\sharp,a_0^\sharp)$ is therefore no larger than at the
radius ceilings before they are re-read, and Lemma~\ref{4.17:lem-seam} claims no more there.
None of the re-readings in~(a) and~(b) is a cap (Definition~\ref{2.3:def-cap}) on $L$, $M$, the
number of renormalisation steps, the level $k$ or the volume, a lower bound on $a_1^\sharp$, a
floor on the coupling, a window, or a two-sided condition.

The proof of Lemma~\ref{4.17:lem-seam} is given in this section.

\begin{definition}[Scales, partitions and cells of one large-field operation]\label{4.18:not-scales-partitions}
Throughout, $d=4$ and the gauge group is $SU(\mathfrak{n})$ with $\mathfrak{n}\ge 2$. The rank is
written $\mathfrak{n}$, and not $N$ as elsewhere in the book, because in this section the letter
$N$ denotes the number of preparatory steps of the large-field operation. One application of
Ba{\l}aban's large-field operation $\mathbf{R}$ at a step $k\ge 1$ is fixed.

\emph{Units.} Lengths are measured in the units of the step-$k$ block lattice. In these units
the fine lattice is $T_{\eta}$, the lattice of spacing $\eta=L^{-k}$. For $i\le k$ the level-$i$
lattice has spacing $u_i=L^{-(k-i)}$. We write $M=L^{m'}$ with $m'\ge 1$.

\emph{Partitions.} For $a>0$, $\mathcal{Q}_a$ is the partition of $\mathbb{R}^d$ into the closed
cubes
\begin{equation*}
\prod_{\mu=1}^{d}\,[a n_\mu,\,a(n_\mu+1)],\qquad n\in\mathbb{Z}^d .
\end{equation*}
For $a,a'\in L^{\mathbb{Z}}$ we write $a'\mid a$ if $a/a'$ is a positive integer. For such $a$
and $a'$ the partitions $\mathcal{Q}_{a'}$ and $\mathcal{Q}_a$ are \emph{compatible}: every cube
of $\mathcal{Q}_a$ is a union of cubes of $\mathcal{Q}_{a'}$, and every cube of
$\mathcal{Q}_{a'}$ lies in a cube of $\mathcal{Q}_a$.

\emph{Distances.} The sup-norm distance of $x,y\in\mathbb{R}^d$ is
$\operatorname{dist}_\infty(x,y)=\max_\mu|x_\mu-y_\mu|$. Distances from a point to a set, and
between two sets, are the corresponding infima. For $S\subset\mathbb{R}^d$ and $r>0$ we put
\begin{equation*}
S^{+r}=\{x\in\mathbb{R}^d:\ \operatorname{dist}_\infty(x,S)<r\}.
\end{equation*}

\emph{Cells.} Let $R_i$ be the cell parameter of level $i$, and put $a_i=MR_iu_i$. The
\emph{level-$i$ cells} are the cubes of $\mathcal{Q}_{a_i}$. The \emph{level-$i$ $M$-cubes} are
the cubes of $\mathcal{Q}_{Mu_i}$. Given the coupling flow of Theorem~\ref{4.1:thm-clause-zero} at
levels $\le k$, and under the ceiling $g_k\le e^{-1}$, the following hold:
\begin{itemize}
\item[(R1)] $R_i$ is a power of $L$, and $R_i\ge L$;
\item[(R2)] $R_k<L(\log g_k^{-2})^{r}$, where $r$ is the fixed exponent entering the cell
parameter;
\item[(R3)] $R_{i-1}\in\{R_i,\,LR_i\}$.
\end{itemize}
Consequently $Mu_i\mid a_i$, $a_i\mid a_{i'}$ for $i\le i'$, and $i\mapsto a_i$ is
non-decreasing.

\emph{Cell hulls.} Let $A$ be a union of cubes of $\mathcal{Q}_b$ with $b\mid a_i$. Then $[A]_i$
is the union of those level-$i$ cells that contain a cube of $A$. One has $A\subset[A]_i$, and
$[A]_i=A$ when $A$ is a union of level-$i$ cells. Moreover every point of $[A]_i$ lies within
sup-distance $a_i$ of $A$:
\begin{equation}\label{4.18:eq-scales-partitions-a}
\operatorname{dist}_\infty(x,A)\le a_i\qquad\text{for every }x\in[A]_i .
\end{equation}

\emph{Touching.} Two closed cubes of one partition \emph{touch} if they share a point. Two cubes
of $\mathcal{Q}_a$ that do not touch are at sup-distance at least $a$.
\end{definition}

\begin{proof}
\emph{Compatibility.} Let $q=a/a'\in\mathbb{Z}_{>0}$. For each $\mu$ the interval
$[an_\mu,a(n_\mu+1)]$ is the union of the intervals $[a'(qn_\mu+j),\,a'(qn_\mu+j+1)]$ with
$0\le j\le q-1$. Taking products over $\mu$, every cube of $\mathcal{Q}_a$ is the union of $q^d$
cubes of $\mathcal{Q}_{a'}$. Conversely, the cube of $\mathcal{Q}_{a'}$ with index $n'$ lies in
the cube of $\mathcal{Q}_a$ whose index has components $\lfloor n'_\mu/q\rfloor$.

\emph{The cell parameters.} Properties (R1)--(R3) are not proved here. They are the content of
the lemma on the cell parameter, which derives them from the coupling flow
(Theorem~\ref{4.1:thm-clause-zero}) under the ceiling $g_k\le e^{-1}$.

We derive the consequences. First, $a_i/(Mu_i)=R_i$ is a positive power of $L$ by (R1), so
$Mu_i\mid a_i$. Second, $u_{i-1}=u_i/L$, so by (R3)
\begin{equation*}
a_{i-1}=MR_{i-1}u_{i-1}=\frac{MR_{i-1}u_i}{L}\in\Bigl\{\frac{a_i}{L},\ a_i\Bigr\}.
\end{equation*}
Hence $a_{i-1}\mid a_i$ and $a_{i-1}\le a_i$. Induction on $i'-i$ gives $a_i\mid a_{i'}$ and
$a_i\le a_{i'}$ for all $i\le i'$.

\emph{Cell hulls.} Since $b\mid a_i$, compatibility applies to $\mathcal{Q}_b$ and
$\mathcal{Q}_{a_i}$.

Let $x\in A$. Then $x$ lies in a cube $Q$ of $\mathcal{Q}_b$ with $Q\subset A$, and $Q$ lies in a
level-$i$ cell $C$. This cell contains the cube $Q$ of $A$, so $C\subset[A]_i$. Therefore
$x\in[A]_i$, and $A\subset[A]_i$.

Now let $A$ be a union of level-$i$ cells, and let $C$ be a level-$i$ cell containing a cube $Q$
of $A$. Choose a point $y$ in the interior of $Q$; it is also an interior point of $C$. Distinct
cubes of $\mathcal{Q}_{a_i}$ meet only along their boundaries, so $C$ is the only level-$i$ cell
containing $y$. Since $y\in A$ and $A$ is a union of level-$i$ cells, the cell of $A$ that
contains $y$ must be $C$. Hence $C\subset A$, so $[A]_i\subset A$, and with the previous
inclusion $[A]_i=A$.

For \eqref{4.18:eq-scales-partitions-a}, let $x\in[A]_i$. Then $x$ lies in a level-$i$ cell $C$
that contains a cube $Q$ of $A$; pick any $y\in Q$. Both $x$ and $y$ lie in $C$, a cube of side
$a_i$, and two points of one cell are at sup-distance at most $a_i$. Hence
\begin{equation*}
\operatorname{dist}_\infty(x,A)\le\operatorname{dist}_\infty(x,y)\le a_i .
\end{equation*}

\emph{Touching.} Let the cubes of $\mathcal{Q}_a$ with indices $n$ and $n'$ not touch. Suppose
$|n_\mu-n'_\mu|\le 1$ for every $\mu$. Then for each $\mu$ the intervals
$[an_\mu,a(n_\mu+1)]$ and $[an'_\mu,a(n'_\mu+1)]$ have a common value $t_\mu$, and the point
$(t_1,\dots,t_d)$ lies in both cubes, which is excluded. So there is a direction $\mu$ with
$|n_\mu-n'_\mu|\ge 2$, say $n'_\mu\ge n_\mu+2$. Then for $x$ in the first cube and $y$ in the
second,
\begin{equation*}
\operatorname{dist}_\infty(x,y)\ge y_\mu-x_\mu\ge an'_\mu-a(n_\mu+1)\ge a .
\end{equation*}
Taking the infimum over $x$ and $y$ gives a sup-distance of at least $a$.
\end{proof}

\begin{definition}[The component, the prepared sequence and the collars]\label{4.18:def-prepared-sequence}
We use the units, partitions and cells of Notation~\ref{4.18:not-scales-partitions}. In
particular lengths are measured in the units of the step-$k$ block lattice, $u_i=L^{-(k-i)}$ is
the spacing of the level-$i$ lattice, $a_i=MR_iu_i$ is the side of a level-$i$ cell, and a
level-$i$ $M$-cube is a cube of the partition $\mathcal{Q}_{Mu_i}$. For a set $S$ and $r>0$,
$S^{+r}$ is the open $r$-neighbourhood of $S$ in the metric $\operatorname{dist}_\infty$, and
$S^{\sim \ell}$ is Ba{\l}aban's notation of \cite{B-CMP122a} for the enlargement of $S$ by $\ell$
layers.

\begin{enumerate}
\item \emph{The component.} $Z$ is one component of the class of Ba{\l}aban's large-field
operation $\mathbf{R}$ at step $k$ \cite{B-CMP122a}, that is, a component with the two properties
``(i) it is contained in a cube of the size $100MR_k$, (ii) in the preceding $N$ renormalization
steps no new large field regions were created inside this component, and the previous regions
contained in it satisfy the condition (i) on the corresponding scales''. Ba{\l}aban works ``in
each component of $Z$ separately'' \cite{B-CMP122a}; accordingly, inside the component, every
set $Z_j$ below is identified with $Z_j\cap Z$.

\item \emph{The label in force.} The \emph{label in force} at the start of the operation consists
of the nested sequences
\begin{equation*}
\Omega_1\supset\Lambda_1\supset\Omega_2\supset\dots\supset\Omega_k\supset\Lambda_k
\end{equation*}
of \cite[(2.1)]{B-CMP119}, and $Z_i:=\Lambda_i^c$. The integer $N=k-h$ is the number of
preparatory integrations, which are those producing the levels $h+1,\dots,k$; the members of the
levels $i\le h$ are carried through the operation unchanged. $N$ is subject to Ba{\l}aban's
conditions ``We assume that $N>N_0$'' \cite{B-CMP122a} and ``$N\leqq R_k$'' \cite{B-CMP122b}.

\item \emph{The integer $N_0$.} $N_0$ is ``the smallest positive integer $N_0$ such, that
$L^{-N_0+1}MR_{k-N_0+1}=M$'' \cite{B-CMP122a}; that is, with $k_0:=k-N_0$,
\begin{equation}\label{4.18:eq-prepared-sequence-a}
\begin{gathered}
L^{N_0-1}=R_{k_0+1},\qquad a_{k_0+1}=M,\\
a_i\le M\quad(i\le k_0+1),\qquad a_i\ge M\quad(i\ge k_0+1).
\end{gathered}
\end{equation}
By the lemma on the range of the integer $N_0$ one has $N_0\ge 2$ and $N_0\le R_k-1$. We write
$I_M:=\{k_0+1,\dots,k\}$.

\item \emph{The prepared sequence.} The prepared sequence $\{Z''_i\}_{i\le k}$ of
\cite{B-CMP122a} is as follows.
\begin{itemize}
\item For $i\in I_M$ the member $Z''_i$ is a union of level-$i$ $M$-cubes (the construction is
``based on partitions into $M$-cubes in the corresponding scales''); these members are nested,
and $Z''_i$ and $Z''_{i-1}$ are ``separated by one layer'' of level-$i$ $M$-cubes. The last member
is $Z''_k=(\Omega^{\sim5}_{k_0+1})^c\cap Z$.
\item For $h<i\le k_0$, $Z''_i=(\Omega_i^{\sim5})^c\cap Z$.
\item For $i\le h$, $Z''_i=Z_i$, as in \cite[(1.11)]{B-CMP122a}.
\end{itemize}
The set
\begin{equation*}
\Lambda:=(\Omega^{\sim4}_{k_0+1})^c\cap Z
\end{equation*}
is, in Ba{\l}aban's words \cite[(1.73)]{B-CMP122a}, ``obtained by adding one layer of $M$-cubes
to $Z''_k$ \dots\ a rectangular parallelepiped contained in a cube of the size $100M$''. The
construction of \cite[\S1]{B-CMP122a} makes the prepared sequence nested,
\begin{equation*}
Z''_1\subset\dots\subset Z''_{h+1}\subset\dots\subset Z''_k\subset\Lambda\subset Z,
\qquad Z''_{h+1}\supset Z_h,
\end{equation*}
and $\Lambda$ satisfies
\begin{equation}\label{4.18:eq-prepared-sequence-b}
\begin{gathered}
\{x\in Z:\operatorname{dist}_\infty(x,Z''_k)<M\}\subset\Lambda
\subset\overline{(Z''_k)^{+M}},\\
\Lambda\subset\Omega^c_{k_0+1}\cap Z\subset Z_{k_0+1}\cap Z .
\end{gathered}
\end{equation}

\item \emph{The collars.} For $0\le n\le N-1$ put $j:=h+n$. The integration of stage $n+1$ is over
the variables localised in the \emph{collar}
\begin{equation}\label{4.18:eq-prepared-sequence-c}
R_n:=Z_{j+1}\setminus Z''_{j+1}\qquad(\text{inside }Z)
\end{equation}
\cite{B-CMP122a}. (The collar carries the stage index $n$; the scale factors $R_i$ entering the
cell sides $a_i$ carry level indices.) After the integration of stage $n+1$, the member of level
$j+1$ in force inside $Z$ is $Z''_{j+1}$, the members of the levels $i>j+1$ are still $Z_i$, and
those of the levels $i\le j$ are $Z''_i$.

\item \emph{The levels above $h$.} At every level $i>h$ the member $Z_i$ is produced by
Ba{\l}aban's one-step operations $\mathbf{T}$ of \cite{B-CMP119}, no new large-field region being
created inside $Z$ at these levels by condition (ii); hence, by the cell structure of the regions
produced by the operations $\mathbf{T}$, $Z_i$ is a union of level-$i$ cells.
For the levels $h<t\le k$ Ba{\l}aban's construction of the operations $\mathbf{T}$ in
\cite{B-CMP119} gives the margins
\begin{equation*}
Z_t\supset(Z'_{t-1})^{\sim10}\supset(Z_{t-1})^{+10a_t},\qquad
\Omega_t\subset\bigl((Z'_{t-1})^{\sim5}\bigr)^c,
\end{equation*}
and therefore
\begin{equation*}
\Omega_t^c\supset(Z_{t-1})^{+5La_t},
\end{equation*}
where $Z'_{t-1}$ is the auxiliary region of that name in Ba{\l}aban's construction of the
operation $\mathbf{T}$ in \cite{B-CMP119}.

\item \emph{The inner core.} $Z^{\sim-8}$ denotes $Z$ less eight successive layers of level-$k$
cells.

\item \emph{Hypothesis \textup{(G)}.} Every member $Z_i\cap Z$, $i\le h$, of the label in force
is a union of level-$i$ cells. Hypothesis~(G) is not assumed in what follows unless this is said.
\end{enumerate}
\end{definition}

\begin{proof}
We derive the equalities of \eqref{4.18:eq-prepared-sequence-a} from the definition of $N_0$,
and the inclusions \eqref{4.18:eq-prepared-sequence-b} from the description of $\Lambda$.

\emph{The equalities in \eqref{4.18:eq-prepared-sequence-a}.} By definition $N_0$ satisfies
$L^{-N_0+1}MR_{k-N_0+1}=M$. Dividing by $M$ and using $k-N_0+1=k_0+1$ gives
$R_{k_0+1}=L^{N_0-1}$. Since $k-(k_0+1)=N_0-1$, the level-$(k_0+1)$ spacing is
$u_{k_0+1}=L^{-(N_0-1)}$, and therefore
\begin{equation*}
a_{k_0+1}=MR_{k_0+1}u_{k_0+1}=M\,L^{N_0-1}L^{-(N_0-1)}=M .
\end{equation*}
The two one-sided inequalities of \eqref{4.18:eq-prepared-sequence-a}, $a_i\le M$ for
$i\le k_0+1$ and $a_i\ge M$ for $i\ge k_0+1$, are stated for this integer $N_0$ in
\cite{B-CMP122a}, at the place where Ba{\l}aban uses it; we take them from there, together with
the range $N_0\ge2$, $N_0\le R_k-1$, which is the content of the lemma on the range of the
integer $N_0$.

\emph{The first line of \eqref{4.18:eq-prepared-sequence-b}.} At level $k$ one has $u_k=1$, so a
level-$k$ $M$-cube is a cube of $\mathcal{Q}_M$, of side $M$. We read Ba{\l}aban's description
\cite[(1.73)]{B-CMP122a} of $\Lambda$ as ``obtained by adding one layer of $M$-cubes to $Z''_k$''
in the $\operatorname{dist}_\infty$ geometry of Notation~\ref{4.18:not-scales-partitions}: the
layer adjoined to the union $Z''_k$ of such cubes consists of the cubes of $\mathcal{Q}_M$ that
meet $Z''_k$, and since $\Lambda\subset Z$, the set $\Lambda$ consists of the points of $Z$ lying
in a cube of $\mathcal{Q}_M$ that meets $Z''_k$.

Let $x\in Z$ and $\operatorname{dist}_\infty(x,Z''_k)<M$. Choose $y\in Z''_k$ with
$|x-y|_\infty<M$. Since $Z''_k$ is a union of cubes of $\mathcal{Q}_M$, there is a cube
$Q'=\prod_\mu[Mn'_\mu,M(n'_\mu+1)]\subset Z''_k$ containing $y$; let
$Q=\prod_\mu[Mn_\mu,M(n_\mu+1)]$ be a cube of $\mathcal{Q}_M$ containing $x$. If
$n_\mu\ge n'_\mu+2$ for some $\mu$, then
\begin{equation*}
x_\mu-y_\mu\ge Mn_\mu-M(n'_\mu+1)\ge M,
\end{equation*}
contradicting $|x-y|_\infty<M$; symmetrically $n'_\mu\ge n_\mu+2$ is impossible. Hence
$|n_\mu-n'_\mu|\le1$ for every $\mu$, and the point with coordinates
$M\max\{n_\mu,n'_\mu\}$ lies in $Q\cap Q'$. Thus $Q$ meets $Z''_k$, and $x\in\Lambda$.

Conversely, let $x\in\Lambda$. Then $x$ lies in a cube $Q$ of $\mathcal{Q}_M$ containing a point
$z\in Z''_k$, and $|x-z|_\infty\le M$ because $Q$ has side $M$; so
$\operatorname{dist}_\infty(x,Z''_k)\le M$. The set of points at $\operatorname{dist}_\infty$ at
most $M$ from $Z''_k$ is the closure of $(Z''_k)^{+M}$: it is closed and contains $(Z''_k)^{+M}$,
and a point $x$ with $\operatorname{dist}_\infty(x,Z''_k)=M$, attained at $z\in Z''_k$, is the
limit of the points $z+t(x-z)$, $t\uparrow1$, whose distance to $Z''_k$ is at most $tM<M$. Hence
$\Lambda\subset\overline{(Z''_k)^{+M}}$.

\emph{The second line of \eqref{4.18:eq-prepared-sequence-b}.} The enlargement
$\Omega^{\sim4}_{k_0+1}$ contains $\Omega_{k_0+1}$, so its complement is contained in
$\Omega^c_{k_0+1}$, and $\Lambda=(\Omega^{\sim4}_{k_0+1})^c\cap Z\subset\Omega^c_{k_0+1}\cap Z$.
Finally $\Omega_{k_0+1}\supset\Lambda_{k_0+1}$ in the label in force, so
$\Omega^c_{k_0+1}\subset\Lambda^c_{k_0+1}=Z_{k_0+1}$, which gives
$\Omega^c_{k_0+1}\cap Z\subset Z_{k_0+1}\cap Z$.
\end{proof}

\begin{definition}[The two profiles of a member and the admissible class]
\label{4.18:def-profiles-admissible}
We work with the units, partitions and cells of
Convention~\ref{4.18:not-scales-partitions} and with the prepared sequence of
Definition~\ref{4.18:def-prepared-sequence}. For $x\in\mathbb{R}^d$ and
$S\subset\mathbb{R}^d$ we write $\operatorname{dist}_\infty(x,S)$ for the distance
from $x$ to $S$ in the sup norm, and $S^{+r}=\{x:\operatorname{dist}_\infty(x,S)<r\}$
for the open $r$-neighbourhood of $S$.

\emph{The smoothed indicator.} Fix $\psi_{\varphi}\in C^\infty([0,\infty);[0,1])$
with $\psi_{\varphi}=1$ on $[0,\tfrac18]$ and $\psi_{\varphi}=0$ on $[\tfrac14,\infty)$.
For $\sigma>0$ and a locally finite family $\mathcal{P}$ of closed cubes
$D=\prod_{\mu}D_\mu$ of side $\sigma$ put
\begin{equation}\label{4.18:eq-profiles-admissible-a}
\Phi_{\mathcal{P},\sigma}(x):=\prod_{D\in\mathcal{P}}
\Bigl(1-\prod_{\mu=1}^{d}\psi_{\varphi}\bigl(\operatorname{dist}(x_\mu,D_\mu)/\sigma\bigr)\Bigr).
\end{equation}

\emph{The two profiles.} Let $A$ be a member of the sequence in force at level $i$.
It is a union of cubes of its defining partition $\mathcal{Q}_b$, where $b=a_i$ at a
cell-made level and $b=Mu_i$ at an $M$-cube-made level
(Definition~\ref{4.18:def-prepared-sequence}).
The \emph{family's profile} of $A$ is
$\varphi_A:=\Phi_{\mathcal{P},a_i}$ with $\mathcal{P}$ the family of level-$i$ cells
of which $[A]_i$ is the union; the \emph{companion profile} of $A$ is
$\varphi^{\circ}_A:=\Phi_{\mathcal{P},b}$ with $\mathcal{P}$ the family of cubes of
$\mathcal{Q}_b$ of which $A$ is the union.

\emph{The admissible class.} A cube of $\mathcal{Q}_b$ \emph{touches} $A$ if it meets
$A$. The class $\mathcal{P}(A)$ consists of the $C^\infty$ functions $\varphi$ on
$\mathbb{R}^d$ with $0\le\varphi\le1$, $\varphi=0$ on a neighbourhood of $A$, and
$\varphi=1$ off the union of $A$ and the cubes of $\mathcal{Q}_b$ touching $A$. These
are the two conditions
$\varphi_j\in C_0^\infty(\Lambda_j)$ and $\varphi_j=1$ on $\Lambda_j^{\sim -1}$ of
\cite[p.~259]{B-CMP119}, with Ba{\l}aban's $\Lambda_j$ in the role of the complement of
$A$ and $\Lambda_j^{\sim-1}$ in the role of that complement less one layer of cubes of
$\mathcal{Q}_b$; the layer is read, as in \cite{B-CMP119}, as a layer ``of cubes taken
from this family of cubes, in terms of which the sets are defined''. The letter
$\mathcal{P}$ for a family of cubes in
\eqref{4.18:eq-profiles-admissible-a} and the class $\mathcal{P}(A)$ are distinguished
by the argument $(A)$.

Then the following hold.
\begin{enumerate}
\item[(a)] $\Phi_{\mathcal{P},\sigma}$ is $C^\infty$ with values in $[0,1]$;
$\Phi_{\mathcal{P},\sigma}(x)=0$ whenever $\operatorname{dist}_\infty(x,D)\le\sigma/8$
for some $D\in\mathcal{P}$, in particular on $\bigcup\mathcal{P}$; and
$\Phi_{\mathcal{P},\sigma}(x)\ne1$ only if $\operatorname{dist}_\infty(x,D)<\sigma/4$
for some $D\in\mathcal{P}$.
\item[(b)] $\varphi^{\circ}_A\in\mathcal{P}(A)$ at every level, and
$\varphi^{\circ}_A=\varphi_A$ at a cell-made level.
\item[(c)] At an $M$-cube-made level, $\varphi_A=0$ on all of $[A]_i$, and indeed at
every point $x$ with $\operatorname{dist}_\infty(x,[A]_i)\le a_i/8$. Hence $\varphi_A$
fails the second condition defining $\mathcal{P}(A)$ at every such point lying neither
in $A$ nor in a cube of $\mathcal{Q}_{Mu_i}$ touching $A$, and so as soon as such a
point exists.
\item[(d)] Since $a_i=MR_iu_i$ with $R_i\ge2$
(Convention~\ref{4.18:not-scales-partitions}), so that
$\tfrac14Mu_i\le\tfrac18a_i$, at every level
\begin{equation}\label{4.18:eq-profiles-admissible-b}
0\le\varphi_A\le\varphi^{\circ}_A\le1,\qquad
\varphi^{\circ}_A-\varphi_A\ne0\ \text{only on}\ \bigl([A]_i\bigr)^{+\frac14a_i}\setminus A.
\end{equation}
\end{enumerate}

\emph{Restriction to a neighbourhood of $Z$.} The members of the other components of
the step-$k$ large-field region, and their factors, lie at sup-distance at least $a_k$
from $Z$, and their profiles are identically $1$ on $Z^{+\frac34a_k}$.
Accordingly, all functions in what
follows are restricted to $Z^{+\frac34a_k}$. There the only factors different from $1$
are those of the members inside $Z$ and, at level $k$ and at the last stage only, the
common factor of the cells of $Z$ itself, which enters with the stages of the operation
introduced below.
\end{definition}

\begin{proof}
\emph{(a).} Let $D=\prod_\mu D_\mu\in\mathcal{P}$ with $D_\mu=[s_\mu,s_\mu+\sigma]$.
On $\mathbb{R}$ we have
$\operatorname{dist}(t,D_\mu)=\max\{s_\mu-t,\,0,\,t-s_\mu-\sigma\}$. This function
is affine on each of the open sets $\{t<s_\mu\}$ and $\{t>s_\mu+\sigma\}$. On the open
set $\{\operatorname{dist}(t,D_\mu)<\sigma/8\}$, which contains $D_\mu$, the composite
$\psi_{\varphi}(\operatorname{dist}(t,D_\mu)/\sigma)$ is identically $1$. Hence
$t\mapsto\psi_{\varphi}(\operatorname{dist}(t,D_\mu)/\sigma)$ is $C^\infty$ with values
in $[0,1]$. The product
$f_D(x)=\prod_{\mu}\psi_{\varphi}(\operatorname{dist}(x_\mu,D_\mu)/\sigma)$ is therefore
$C^\infty$ with values in $[0,1]$. For a product of intervals,
$\operatorname{dist}_\infty(x,D)=\max_\mu\operatorname{dist}(x_\mu,D_\mu)$. So
$f_D(x)=1$ if $\operatorname{dist}_\infty(x,D)\le\sigma/8$, and $f_D(x)\ne0$ only if
$\operatorname{dist}_\infty(x,D)<\sigma/4$.

Let a bounded open set be given. Its $\sigma/4$-neighbourhood is bounded, so by local
finiteness it meets only finitely many $D\in\mathcal{P}$. Consequently, on this set all
factors $1-f_D$ in \eqref{4.18:eq-profiles-admissible-a} except finitely many are
identically $1$. Thus on it the product is a finite product of $C^\infty$ functions
with values in $[0,1]$, and so is $\Phi_{\mathcal{P},\sigma}$.

If $\operatorname{dist}_\infty(x,D)\le\sigma/8$ for some $D\in\mathcal{P}$, then the
factor $1-f_D(x)$ vanishes; this holds on $\bigcup\mathcal{P}$, where the distance is
$0$. If $\Phi_{\mathcal{P},\sigma}(x)\ne1$, some factor differs from $1$, that is,
$f_D(x)\ne0$ for some $D$. Hence $\operatorname{dist}_\infty(x,D)<\sigma/4$.

\emph{(b).} Let $\mathcal{P}$ be the family of cubes of $\mathcal{Q}_b$ whose union is
$A$. It is locally finite because $\mathcal{Q}_b$ is a partition into cubes of side $b$.
By (a), $\varphi^{\circ}_A$ is $C^\infty$ with $0\le\varphi^{\circ}_A\le1$. Again by
(a), it vanishes on the union over $D\in\mathcal{P}$ of the open sets
$\{\operatorname{dist}_\infty(\cdot,D)<b/8\}$, which is a neighbourhood of $A$.

Now let $\varphi^{\circ}_A(x)\ne1$. By (a), $\operatorname{dist}_\infty(x,D)<b/4$ for
some $D\in\mathcal{P}$. Let $E\in\mathcal{Q}_b$ contain $x$, and write
$E=\prod_\mu[bn_\mu,b(n_\mu+1)]$ and $D=\prod_\mu[bn'_\mu,b(n'_\mu+1)]$. If
$|n_\mu-n'_\mu|\ge2$ for some $\mu$, then
$\operatorname{dist}(x_\mu,D_\mu)\ge b>b/4$, which is impossible. Hence
$|n_\mu-n'_\mu|\le1$ for every $\mu$, and $E\cap D\ne\emptyset$. So either $E\subset A$
or $E$ is a cube of $\mathcal{Q}_b$ touching $A$, and in both cases $x$ lies in the
union of $A$ and the cubes touching $A$. This is the second condition, so
$\varphi^{\circ}_A\in\mathcal{P}(A)$.

At a cell-made level, $b=a_i$ and $A$ is a union of level-$i$ cells. The level-$i$
cells holding a cube of $A$ are then exactly the cubes of $A$, so $[A]_i=A$. The two
families in the definitions of $\varphi_A$ and $\varphi^{\circ}_A$ therefore coincide,
and so do the side lengths. Hence $\varphi^{\circ}_A=\varphi_A$.

\emph{(c).} Let $\mathcal{P}$ be the family of level-$i$ cells composing $[A]_i$. By (a),
$\varphi_A=\Phi_{\mathcal{P},a_i}$ vanishes at every $x$ with
$\operatorname{dist}_\infty(x,D')\le a_i/8$ for some $D'\in\mathcal{P}$, in particular on
$\bigcup\mathcal{P}=[A]_i$. If $\operatorname{dist}_\infty(x,[A]_i)\le a_i/8$, then,
since $[A]_i$ is the union of the locally finite family $\mathcal{P}$ of closed cells,
only finitely many of which meet a bounded neighbourhood of $x$, the distance from $x$
to $[A]_i$ equals $\operatorname{dist}_\infty(x,D')$ for one $D'\in\mathcal{P}$; so
$\varphi_A(x)=0$. A point of this kind outside $A$ and outside every cube of
$\mathcal{Q}_{Mu_i}$ touching $A$ is a point off that union at which
$\varphi_A=0\ne1$. At such a point the second condition fails.

\emph{(d).} At a cell-made level, (b) gives $\varphi_A=\varphi^{\circ}_A$, and (a) gives
values in $[0,1]$, so \eqref{4.18:eq-profiles-admissible-b} holds. Now let the level be
$M$-cube-made, so that $b=Mu_i$. Since $[A]_i$ is the union of the level-$i$ cells
holding the cubes of $A$, we have $A\subset[A]_i$, and each cube $D$ of $A$ lies in a
cell of $[A]_i$.

Suppose $\varphi^{\circ}_A(x)\ne1$. By (a), $\operatorname{dist}_\infty(x,D)<Mu_i/4$
for some cube $D$ of $A$. Since $a_i=MR_iu_i$ and $R_i\ge2$,
\begin{equation}\label{4.18:eq-profiles-admissible-1}
\tfrac14Mu_i\le\tfrac18a_i .
\end{equation}
So $\operatorname{dist}_\infty(x,D')\le\operatorname{dist}_\infty(x,D)<a_i/8$ for the
cell $D'$ of $[A]_i$ containing $D$. By (a), $\varphi_A(x)=0\le\varphi^{\circ}_A(x)$.
Where $\varphi^{\circ}_A(x)=1$, (a) gives $\varphi_A(x)\le1=\varphi^{\circ}_A(x)$.
Together with (a) for both functions, this proves the chain of inequalities in
\eqref{4.18:eq-profiles-admissible-b}.

Let now $\varphi^{\circ}_A(x)\ne\varphi_A(x)$. Since $\varphi^{\circ}_A=\varphi_A=0$
on $A\subset[A]_i$ by (a), we have $x\notin A$. There are two cases.
\begin{itemize}
\item If $\varphi^{\circ}_A(x)\ne1$, then, as above, $x$ is within sup-distance
$Mu_i/4\le a_i/8$ of $A\subset[A]_i$. Hence $x\in([A]_i)^{+\frac14a_i}$.
\item If $\varphi^{\circ}_A(x)=1$, then $\varphi_A(x)\ne1$. By (a),
$\operatorname{dist}_\infty(x,D')<a_i/4$ for some cell $D'$ of $[A]_i$, so again
$x\in([A]_i)^{+\frac14a_i}$.
\end{itemize}
This proves the support statement in \eqref{4.18:eq-profiles-admissible-b}.
\end{proof}

\begin{definition}[Coupling functions of the two systems and the profile excess]
\label{4.18:def-coupling-excess}
Work at step $k$ in the setting of Definition~\ref{4.18:def-prepared-sequence}, with the profiles of
Definition~\ref{4.18:def-profiles-admissible}. The increments $\beta_i$ of
Definition~\ref{1.2:def-couplings} satisfy, with $x_0:=g_0^{-2}$, the recursion
\begin{equation*}
g_{i-1}^{-2}=g_i^{-2}+\beta_i\qquad (i\le k),
\end{equation*}
and we use the bound $0\le \beta_i\le b_0$ $(i\le k)$ of the coupling-flow statement for the levels
$i\le k$ that is in force at step $k$. Therefore, by telescoping,
\begin{equation*}
g_i^{-2}=x_0-\sum_{\ell\le i}\beta_\ell .
\end{equation*}

\emph{Coupling function of a system.} A \emph{system} $\Pi=(A_i,\varphi_i)_{i\le k}$ consists of a member
$A_i$ and a profile $\varphi_i$ for every level $i\le k$. Its \emph{coupling function} is
\begin{equation*}
x_{\Pi}:=x_0-\sum_{i\le k}\beta_i\varphi_i ,
\end{equation*}
which is \cite[(2.24)]{B-CMP119} in telescoped form.

\emph{The two stage systems.} Fix a stage $n$ and put $j=h+n$. For $i\le k$ let $A^{(n)}_i$ be the member of
level $i$ of the sequence in force at stage $n$. The \emph{family's system} $\Pi^{(n)}$ consists of the
members $A^{(n)}_i$ with the family's profiles $\varphi_{A^{(n)}_i}$; the \emph{companion system}
$\Pi^{(n),\circ}$ consists of the same members with the companion profiles $\varphi^{\circ}_{A^{(n)}_i}$. We
write
\begin{equation}\label{4.18:eq-coupling-excess-a}
\begin{gathered}
x^{(n)}:=x_{\Pi^{(n)}},\qquad x^{(n),\circ}:=x_{\Pi^{(n),\circ}},\\
e_n:=x^{(n)}-x^{(n),\circ}
=\sum_{i\le k}\beta_i\bigl(\varphi^{\circ}_{A^{(n)}_i}-\varphi_{A^{(n)}_i}\bigr)\ \ge\ 0 .
\end{gathered}
\end{equation}
The function $e_n$ is the \emph{profile excess} at stage $n$. For $i\le h$ the member $A^{(n)}_i$ does not
depend on $n$, and we write $A_i:=A^{(n)}_i$. The current and old parts of $e_n$ are
\begin{equation}\label{4.18:eq-coupling-excess-b}
\begin{aligned}
e_n&=e_n^{\mathrm{cur}}+\mathfrak{d}_{\mathrm{old}},\qquad
e_n^{\mathrm{cur}}:=\sum_{i\in I_M,\ i\le j}\beta_i\bigl(\varphi^{\circ}_{Z''_i}-\varphi_{Z''_i}\bigr),\\
\mathfrak{d}_{\mathrm{old}}&:=\sum_{\substack{i\le h\\ A_i\ M\text{-cube-made}}}
\beta_i\bigl(\varphi^{\circ}_{A_i}-\varphi_{A_i}\bigr).
\end{aligned}
\end{equation}
We also put
\begin{equation*}
\mathcal{E}':=\bigcup_{i\in I_M}\operatorname{supp}\bigl(\varphi^{\circ}_{Z''_i}-\varphi_{Z''_i}\bigr),
\qquad
\mathcal{D}:=\bigcup_{i\le k}\bigl([Z''_i]_i\bigr)^{+\frac14 a_i}.
\end{equation*}

With these definitions the following hold.
\begin{enumerate}
\item The identity and the inequality in \eqref{4.18:eq-coupling-excess-a}, and the decomposition
\eqref{4.18:eq-coupling-excess-b}.
\item The bounds
\begin{equation*}
0\le e_n^{\mathrm{cur}}\le b_0\cdot\#\bigl(I_M\cap[1,j]\bigr)\le b_0N_0,
\qquad
0\le \mathfrak{d}_{\mathrm{old}};
\end{equation*}
the pointwise upper bound on $\mathfrak{d}_{\mathrm{old}}(x)$ in terms of the number $n_{\mathrm{old}}(x)$
of contributing older members at $x$ is stated together with the definition of $n_{\mathrm{old}}$, in the
lemma of this chapter on the count of older members; and, under hypothesis~(G) of
Definition~\ref{4.18:def-prepared-sequence}, $\mathfrak{d}_{\mathrm{old}}=0$.
\item The support inclusions
\begin{equation*}
\operatorname{supp}e_n^{\mathrm{cur}}\subset\mathcal{E}'
\subset(\Lambda\cup\mathcal{D})\setminus Z''_{k_0+1}\subset Z^{\sim -8}.
\end{equation*}
\item The stage increments $w_{n+1}$ and $w^{\circ}_{n+1}$ of the two coupling functions are defined by the
first equalities in
\begin{equation}\label{4.18:eq-coupling-excess-c}
\begin{aligned}
w_{n+1}&:=x^{(n+1)}-x^{(n)}=-\beta_{j+1}\bigl(\varphi_{Z''_{j+1}}-\varphi_{Z_{j+1}}\bigr),\\
w^{\circ}_{n+1}&:=x^{(n+1),\circ}-x^{(n),\circ}
=-\beta_{j+1}\bigl(\varphi^{\circ}_{Z''_{j+1}}-\varphi_{Z_{j+1}}\bigr),
\end{aligned}
\end{equation}
the second equalities holding on $Z^{+\frac34 a_k}$; and on $Z^{+\frac34 a_k}$
\begin{equation}\label{4.18:eq-coupling-excess-d}
e_{n+1}-e_n=w_{n+1}-w^{\circ}_{n+1}=
\begin{cases}
\beta_{j+1}\bigl(\varphi^{\circ}_{Z''_{j+1}}-\varphi_{Z''_{j+1}}\bigr), & j+1\in I_M,\\[2pt]
0, & j+1\le k_0 .
\end{cases}
\end{equation}
Here $\varphi_{Z_{j+1}}$ denotes the profile of the whole member of level $j+1$ in force at stage $n$; its
factors other than the one belonging to $Z_{j+1}\cap Z$ equal $1$ on $Z^{+\frac34 a_k}$.
\end{enumerate}
\end{definition}

\begin{proof}
\emph{The identity in \eqref{4.18:eq-coupling-excess-a}.} By the definition of the coupling function of a
system, applied to $\Pi^{(n)}$ and to $\Pi^{(n),\circ}$,
\begin{equation*}
x^{(n)}-x^{(n),\circ}
=\Bigl(x_0-\sum_{i\le k}\beta_i\varphi_{A^{(n)}_i}\Bigr)
-\Bigl(x_0-\sum_{i\le k}\beta_i\varphi^{\circ}_{A^{(n)}_i}\Bigr)
=\sum_{i\le k}\beta_i\bigl(\varphi^{\circ}_{A^{(n)}_i}-\varphi_{A^{(n)}_i}\bigr).
\end{equation*}

\emph{The decomposition \eqref{4.18:eq-coupling-excess-b}.} By the definition of the two profiles of a member
(Definition~\ref{4.18:def-profiles-admissible}), only $M$-cube-made members contribute to the sum just
obtained. These are of two kinds. The current ones are the prepared members $Z''_i$ with $i\in I_M$ and
$i\le j$; each of them is present from the stage at which level $i$ is prepared. The older ones, if there are
any, are members $A_i$ at levels $i\le h$; they do not change with $n$. Collecting the summands of the first
kind gives $e_n^{\mathrm{cur}}$, and those of the second kind give $\mathfrak{d}_{\mathrm{old}}$, which is
\eqref{4.18:eq-coupling-excess-b}. Under hypothesis~(G) of Definition~\ref{4.18:def-prepared-sequence},
$\mathfrak{d}_{\mathrm{old}}=0$.

\emph{The bounds, and the inequality in \eqref{4.18:eq-coupling-excess-a}.} By the comparison of the two
profiles in \eqref{4.18:eq-profiles-admissible-b}, together with $0\le\beta_i\le b_0$, the sum
$e_n^{\mathrm{cur}}$, which has one summand for each $i\in I_M\cap[1,j]$, satisfies
\begin{equation*}
0\le e_n^{\mathrm{cur}}\le b_0\cdot\#\bigl(I_M\cap[1,j]\bigr)\le b_0N_0 ,
\end{equation*}
and likewise $0\le\mathfrak{d}_{\mathrm{old}}$, each summand of $\mathfrak{d}_{\mathrm{old}}$ being
non-negative. Adding the two lower bounds and
using \eqref{4.18:eq-coupling-excess-b} gives $e_n\ge0$.

\emph{The support inclusions.} The support of the finite sum $e_n^{\mathrm{cur}}$ lies in the union of the
supports of its summands, and the support of $\beta_i(\varphi^{\circ}_{Z''_i}-\varphi_{Z''_i})$ lies in
$\operatorname{supp}(\varphi^{\circ}_{Z''_i}-\varphi_{Z''_i})$; hence
$\operatorname{supp}e_n^{\mathrm{cur}}\subset\mathcal{E}'$. The inclusions
$\mathcal{E}'\subset(\Lambda\cup\mathcal{D})\setminus Z''_{k_0+1}\subset Z^{\sim -8}$ are taken from two
statements of this chapter that localise the supports of the profile differences
$\varphi^{\circ}_{Z''_i}-\varphi_{Z''_i}$, $i\in I_M$, of the prepared members: the support proposition for
the prepared members and the support lemma for the profile differences. Together these two statements give
the chain of inclusions displayed in item~3.

\emph{The increments \eqref{4.18:eq-coupling-excess-c}.} In the transition from stage $n$ to stage $n+1$ the
only member that changes is the one of level $j+1$: it is $Z_{j+1}$, which is cell-made, at stage $n$, and
$Z''_{j+1}$ at stage $n+1$; the latter is cell-made if $j+1\le k_0$ and $M$-cube-made if $j+1\in I_M$. All
other summands of the coupling functions therefore cancel in the differences, and
\begin{equation*}
w_{n+1}=-\beta_{j+1}\bigl(\varphi_{A^{(n+1)}_{j+1}}-\varphi_{A^{(n)}_{j+1}}\bigr),\qquad
w^{\circ}_{n+1}=-\beta_{j+1}\bigl(\varphi^{\circ}_{A^{(n+1)}_{j+1}}-\varphi^{\circ}_{A^{(n)}_{j+1}}\bigr).
\end{equation*}
It remains to identify these profiles on $Z^{+\frac34 a_k}$. At stage $n+1$ the member of level $j+1$ is
$Z''_{j+1}$, so that $\varphi_{A^{(n+1)}_{j+1}}=\varphi_{Z''_{j+1}}$ and
$\varphi^{\circ}_{A^{(n+1)}_{j+1}}=\varphi^{\circ}_{Z''_{j+1}}$. At stage $n$ it is $Z_{j+1}$, and
$\varphi_{Z_{j+1}}$ is the profile of the whole member $A^{(n)}_{j+1}$ of level $j+1$ in force at stage $n$,
whose factors other than the one belonging to $Z_{j+1}\cap Z$ equal $1$ on $Z^{+\frac34 a_k}$. Since
$Z_{j+1}$ is cell-made, its companion profile coincides with its family profile
(Definition~\ref{4.18:def-profiles-admissible}), so that on $Z^{+\frac34 a_k}$
\begin{equation*}
\varphi^{\circ}_{A^{(n)}_{j+1}}=\varphi_{A^{(n)}_{j+1}}=\varphi_{Z_{j+1}} .
\end{equation*}
Inserting these identifications into the two differences gives \eqref{4.18:eq-coupling-excess-c} on
$Z^{+\frac34 a_k}$. There, by the statement on the factors, the increments involve only the part of the
member of level $j+1$ inside the component $Z$.

\emph{The excess increment \eqref{4.18:eq-coupling-excess-d}.} By the definition of $e_n$ and of the
increments,
\begin{equation*}
e_{n+1}-e_n=\bigl(x^{(n+1)}-x^{(n)}\bigr)-\bigl(x^{(n+1),\circ}-x^{(n),\circ}\bigr)
=w_{n+1}-w^{\circ}_{n+1}.
\end{equation*}
Inserting \eqref{4.18:eq-coupling-excess-c}, the terms in $\varphi_{Z_{j+1}}$ cancel on $Z^{+\frac34 a_k}$:
\begin{equation*}
w_{n+1}-w^{\circ}_{n+1}
=-\beta_{j+1}\bigl(\varphi_{Z''_{j+1}}-\varphi_{Z_{j+1}}\bigr)
+\beta_{j+1}\bigl(\varphi^{\circ}_{Z''_{j+1}}-\varphi_{Z_{j+1}}\bigr)
=\beta_{j+1}\bigl(\varphi^{\circ}_{Z''_{j+1}}-\varphi_{Z''_{j+1}}\bigr).
\end{equation*}
If $j+1\in I_M$, this is the first case of \eqref{4.18:eq-coupling-excess-d}. If $j+1\le k_0$, then $Z''_{j+1}$
is cell-made, so that $\varphi^{\circ}_{Z''_{j+1}}=\varphi_{Z''_{j+1}}$
(Definition~\ref{4.18:def-profiles-admissible}) and the right side vanishes; this is the second case.
\end{proof}

\begin{definition}[The plaquette functional and the recursion of the preparatory integrations]
\label{4.18:def-stage-recursion}
\emph{The weighted Wilson functional.} Let $T_{\eta}$ be the fine lattice, of spacing $\eta$, let $w$ be a weight
on $T_{\eta}$, and let $\mathbb{U}$ be a bond field on $T_{\eta}$, possibly complexified. We set
\begin{equation}\label{4.18:eq-stage-recursion-1}
A(w,\mathbb{U}) := \sum_{p} w(x_-(p))\,\mathcal{A}_p(\mathbb{U}),
\end{equation}
where the sum runs over the fine plaquettes $p$, $x_-(p)\in T_{\eta}$ is the base point of $p$, and $\mathcal{A}_p$ is
Ba{\l}aban's non-abelian plaquette function \cite[(0.1)]{B-CMP109}. We write $A(\mathbb{U}) := A(1,\mathbb{U})$
for the unweighted functional. The function $\mathcal{A}_p$ satisfies $\mathcal{A}_p\ge 0$ when $\mathbb{U}$ takes
values in the gauge group $G=SU(\mathfrak{n})$, and it is entire in the plaquette variable; by
\eqref{4.18:eq-stage-recursion-1} the map
$w\mapsto A(w,\mathbb{U})$ is linear. The functional is the non-abelian one throughout. The only comparisons made
in what follows are comparisons of scalar weights.

\emph{Where the profiles enter.} The integration structure of Ba{\l}aban's operation $\mathbf{R}$ does not
depend on the profiles of Definition~\ref{4.18:def-coupling-excess}. This is the case for the determining sets
$\mathfrak{B}^{(n)}_k$,
$\mathfrak{B}^{(n)}_k(Z)$ and the backgrounds $U^{(n)}_k$ of \cite[(1.14)--(1.21)]{B-CMP122a}; for the collars
$R_n$; for the characteristic functions $\chi^{(n+1)}_k$, $\chi^{(j)}$, $\chi'^{(j)}$ of
\cite[(1.22)--(1.27), (1.55), (1.59)]{B-CMP122a}, with their thresholds $\delta_j$, $\delta'_j$; for the change of
variables and the response map
$b\mapsto\bigl(e^{i\eta\mathbb{H}^{(n)}_k(b)}\mathbb{U},\ \mathbb{J}+\mathcal{J}(b)\bigr)$; for the Gaussian forms
$\tfrac12\langle A_j, C^*\Delta^{(j)}_{\Lambda_{j+1}}CA_j\rangle$ and
$\tfrac12\langle B_j, C^*\Delta^{(j)}CB_j\rangle$ in the fluctuation fields $A_j$, $B_j$ of the collar $R_n$;
for the remainder $\mathbb{P}^{(j)}(g_j,A_j,B_j)$ of the
constant-coupling Wilson term; for the normalisation constants $z^{(j)}$ and $E^{(j)}_0$ of
\cite[(1.60)]{B-CMP122a}; for the spaces
$\tilde{U}^{(n)c}_k(X,\tilde\alpha_0,\tilde\alpha_1)$ of \cite[(1.64)--(1.69)]{B-CMP122a}; for the domains
$\mathbb{D}_m$; for the tree size $d_m$; and for the filing rule of \cite[(1.63)]{B-CMP122a}. Each of these objects
is defined from the regions $Z_i$, $Z''_i$, $\Omega_i$, from the constants and from the fields, and none of them
involves a profile $\varphi$. This is the profile-freeness of the integration structure; its proof is an
inspection of Ba{\l}aban's definitions of the objects
just listed. The profiles enter only as weights of Wilson terms, and only in
three places:
\begin{itemize}
\item the coupling function in $-A(x^{(n)},\cdot)$;
\item the counterterms $-\beta_\ell A(\mathsf{t}_z\varphi_\ell,\cdot)$ inside $\mathbb{E}_k$, where $\varphi_\ell$ is
the profile of the member with index $\ell$ (Definition~\ref{4.18:def-coupling-excess}) and $\beta_\ell$,
$\mathsf{t}_z$ are as in \cite[(2.25), (3.65)]{B-CMP119};
\item the weight-difference members of \cite[(1.62)]{B-CMP122a}, that is, the terms $\Delta W_{n+1}$ below.
\end{itemize}

\emph{Stage $0$.} The action in force at stage $0$, denoted $A^{(0)}$, is the effective action at level $k$ of
\cite[(2.23)]{B-CMP119}. Following that display, we write
\begin{equation}\label{4.18:eq-stage-recursion-2}
A^{(0)} = -A(x^{(0)},\cdot) + \mathfrak{E}^{(0)},
\qquad
\mathfrak{E}^{(0)} := \mathbb{E}_k + \mathbb{R}_k + \mathbb{B}_k - E .
\end{equation}
Here $x^{(0)}$ is the coupling function at stage $0$ of Definition~\ref{4.18:def-coupling-excess}, and $E$ is the
constant of that display. The counterterms $-\beta_\ell A(\mathsf{t}_z\varphi_\ell,\cdot)$ of $\mathbb{E}_k$
carry the stage-$0$ profiles.

\emph{Stage $n+1$.} Let $j=h+n$. The density at stage $n$ is $\exp A^{(n)}(U^{(n)}_k)$, where
$A^{(n)} = -A(x^{(n)},\cdot)+\mathfrak{E}^{(n)}$. Following \cite[\S1]{B-CMP122a}, write
\begin{equation}\label{4.18:eq-stage-recursion-3}
-A(x^{(n)},\cdot) = -g_j^{-2}A(\cdot) - W^{(n)}(\cdot),
\qquad
W^{(n)} := A(x^{(n)}-g_j^{-2},\cdot).
\end{equation}
Express $U^{(n)}_k$ through the new background $U^{(n+1)}_k$ and the fluctuation $(A_j,B_j)$ on the collar $R_n$,
and call the resulting configuration the \emph{shifted background} $\mathbb{U}^{\mathrm{sh}}$. Expanding
$-g_j^{-2}A$ to second order produces the Gaussian forms, the remainder $\mathbb{P}^{(j)}$ and the normalisations.
By \cite[(1.60)]{B-CMP122a}, the resulting integral takes the form
\begin{equation}\label{4.18:eq-stage-recursion-a}
\begin{split}
&\chi^{(n+1)}_k\exp\bigl[A^{(n)}(U^{(n+1)}_k, A_j=0)+E^{(j)}_0\bigr]
\cdot\int dA_j\,\chi^{(j)}\,
e^{-\frac12\langle A_j,\,C^*\Delta^{(j)}_{\Lambda_{j+1}}CA_j\rangle}\\
&\qquad\cdot z^{(j)}\int dB_j\,\chi'^{(j)}
\exp\Bigl[-\tfrac12\langle B_j,C^*\Delta^{(j)}CB_j\rangle
+\mathbb{P}^{(j)}+\mathfrak{K}^{(n)}\Bigr].
\end{split}
\end{equation}
The interaction bracket $\mathfrak{K}^{(n)}$ is the increment between the two backgrounds of the part of the stage-$n$
action that is not the constant-coupling Wilson term. Precisely, it is the value of
$\mathfrak{E}^{(n)}-W^{(n)}$ at $\mathbb{U}^{\mathrm{sh}}$ minus its value at $(U^{(n+1)}_k, A_j=0)$. One has,
with the functional $\boldsymbol{\mathcal{R}}^{(n)}$ defined in the second line,
\begin{equation}\label{4.18:eq-stage-recursion-b}
\begin{split}
\mathfrak{K}^{(n)}
&=\boldsymbol{\mathcal{R}}^{(n)}(\mathbb{U}^{\mathrm{sh}})
-\boldsymbol{\mathcal{R}}^{(n)}(U^{(n+1)}_k, A_j=0),\\
\boldsymbol{\mathcal{R}}^{(n)} &:= A^{(n)} + g_j^{-2}A(\cdot).
\end{split}
\end{equation}

\emph{Stage terms and the action at stage $n+1$.} In \eqref{4.18:eq-stage-recursion-a}, we call the
\emph{fluctuation integral} the iterated integral over $A_j$ and $B_j$, including the factor $z^{(j)}$ and the
characteristic functions $\chi^{(j)}$, $\chi'^{(j)}$. By
\cite[(1.61)]{B-CMP122a},
\begin{equation*}
\mathbb{B}^{(n)}_k=\sum_{i=1}^{n}\mathbb{C}^{(i)}_k .
\end{equation*}
By \cite[(1.62)]{B-CMP122a}, the stage term is
\begin{equation}\label{4.18:eq-stage-recursion-5}
\begin{split}
\mathbb{C}^{(n+1)}_k &:= \log\bigl(\text{the fluctuation integral}\bigr) + \Delta W_{n+1},\\
\Delta W_{n+1} &:= A(x^{(n+1)}-x^{(n)},U^{(n+1)}_k) = A(w_{n+1},U^{(n+1)}_k),
\end{split}
\end{equation}
with $w_{n+1}=x^{(n+1)}-x^{(n)}$ the stage increment of Definition~\ref{4.18:def-coupling-excess}, and
$\mathbb{B}^{(n+1)}_k=\mathbb{B}^{(n)}_k+\mathbb{C}^{(n+1)}_k$. The action in force at stage $n+1$ is
\begin{equation}\label{4.18:eq-stage-recursion-c}
\begin{split}
A^{(n+1)} &:= A^{(n)}(U^{(n+1)}_k, A_j=0)+E^{(j)}_0+\log\bigl(\text{the fluctuation integral}\bigr)\\
&\phantom{:}= -A(x^{(n+1)},\cdot)+\mathfrak{E}^{(n+1)},\\
\mathfrak{E}^{(n+1)} &:= \mathfrak{E}^{(n)}+\mathbb{C}^{(n+1)}_k+E^{(j)}_0 .
\end{split}
\end{equation}
All terms of $\mathfrak{E}^{(n)}$ are functions on the descending spaces
$\tilde{U}^{(n)c}_k(X,\tilde\alpha_0,\tilde\alpha_1)$ of \cite[(1.64)--(1.69)]{B-CMP122a}, and they are evaluated
at the current background \cite[\S1]{B-CMP122a}.

\emph{Representation of the stage terms.} By \cite[(1.63)]{B-CMP122a},
$\mathbb{C}^{(n+1)}_k=\sum_X\mathbb{C}^{(n+1)}_k(X)$, where the sum runs over the sets $X$ such that, with
$m=m(X)$ the least $m\ge j+1$ for which $X\subset\Omega^c_{m+1}$, one has $X\in\mathbb{D}_m$,
$X\cap\Omega_m\ne\emptyset$ and $X\cap R_n\ne\emptyset$.
Here $\mathbb{D}_m$ is the set of face-connected unions of cubes of the partition
$\pi_m$ into cubes of side $Mu_m$. The stage terms $\mathbb{C}^{(n+1)}_k$ are bounded uniformly in $n$ and $k$
by a theorem that we call the uniform bound on the stage terms; its weighted norm is, for a family
$C=(C(X))_X$,
\begin{equation}\label{4.18:eq-stage-recursion-6}
\|C\| := \sup_m\ \sup_{\square\in\pi_m}\ \sum_{X\ni\square,\ m(X)=m}\|C(X)\|_\infty\, e^{a_m d_m(X)},
\end{equation}
where $a_m$ are the decay rates of the uniform bound on the stage terms, and $\|C(X)\|_\infty$ is the supremum of
$|C(X)|$ over the space $\tilde{U}^{(\cdot)c,\Sigma}_k(X)$ on which its bound is stated. Then, for every $X$,
with $m=m(X)$ the filing level of $X$, one has pointwise
\begin{equation*}
|C(X)|\le\|C\|\,e^{-a_m d_m(X)} .
\end{equation*}
The constant $c_G(d,L,M)+1$ and the constants $\vartheta$, $\mathfrak{s}$, $\mathfrak{p}_j$,
$\mathfrak{l}_j$,
$K_S$, $K_L$ that accompany this norm, together with the conditions $(K)$, $(Q^+)$ and $\gamma\le\gamma_0$, are those
of the uniform bound on the stage terms. They are used only through the inequalities displayed in the uniform
bound on the stage terms.
\end{definition}

\begin{proof}
Two identities have to be checked: \eqref{4.18:eq-stage-recursion-b}, and the second equality in
\eqref{4.18:eq-stage-recursion-c}. Together with \eqref{4.18:eq-stage-recursion-2}, the second of these also shows,
by induction on $n$, that $A^{(n)}=-A(x^{(n)},\cdot)+\mathfrak{E}^{(n)}$ at every stage. This is the form of the
stage-$n$ action used in the description of stage $n+1$.

\emph{Identity \eqref{4.18:eq-stage-recursion-b}.} In the construction of \cite[(1.60)]{B-CMP122a}, only the
constant-coupling Wilson term $-g_j^{-2}A$ is expanded to second order in the fluctuation: in
\cite[\S1]{B-CMP122a} the coupling function $x^{(n)}$ is replaced by $g_j^{-2}$ in the Wilson term, and the term
with the difference, $-W^{(n)}$, is put into the interaction. The Gaussian forms and the
remainder $\mathbb{P}^{(j)}$ come from this term alone. The rest of $A^{(n)}$ is not expanded; this rest is
\begin{equation*}
\boldsymbol{\mathcal{R}}^{(n)}=A^{(n)}+g_j^{-2}A(\cdot).
\end{equation*}
Its value at $(U^{(n+1)}_k,A_j=0)$ joins the value of $-g_j^{-2}A$ there in the prefactor
$\exp[A^{(n)}(U^{(n+1)}_k,A_j=0)+E^{(j)}_0]$, and its increment between the two backgrounds is the interaction
bracket $\mathfrak{K}^{(n)}$: the bracket was described above as the increment of $\mathfrak{E}^{(n)}-W^{(n)}$,
and this function coincides with $\boldsymbol{\mathcal{R}}^{(n)}$, as we now compute. By the stage-$n$ form
$A^{(n)}=-A(x^{(n)},\cdot)+\mathfrak{E}^{(n)}$ and the definition of $W^{(n)}$ in
\eqref{4.18:eq-stage-recursion-3},
\begin{equation*}
\mathfrak{E}^{(n)}-W^{(n)} = A^{(n)}+A(x^{(n)},\cdot)-A(x^{(n)}-g_j^{-2},\cdot).
\end{equation*}
By linearity of $w\mapsto A(w,\mathbb{U})$,
\begin{equation*}
A(x^{(n)},\cdot)-A(x^{(n)}-g_j^{-2},\cdot)=A(g_j^{-2},\cdot)=g_j^{-2}A(\cdot).
\end{equation*}
Hence $\mathfrak{E}^{(n)}-W^{(n)}=A^{(n)}+g_j^{-2}A(\cdot)=\boldsymbol{\mathcal{R}}^{(n)}$ by the definition in
the second line of \eqref{4.18:eq-stage-recursion-b}. Taking the difference of the values at
$\mathbb{U}^{\mathrm{sh}}$ and at
$(U^{(n+1)}_k,A_j=0)$ gives \eqref{4.18:eq-stage-recursion-b}.

The constant $E$ enters $\mathfrak{E}^{(n)}$ through $\mathfrak{E}^{(0)}$ in \eqref{4.18:eq-stage-recursion-2},
and the normalisation constants $E^{(j')}_0$ with $j'<j$ enter it through \eqref{4.18:eq-stage-recursion-c}; these
constants cancel in this difference. Consequently the bracket is unchanged if $\mathfrak{E}^{(n)}$ is replaced by
$\mathbb{E}_k+\mathbb{R}_k+\mathbb{B}_k+\mathbb{B}^{(n)}_k$.

\emph{Identity \eqref{4.18:eq-stage-recursion-c}.} All functions below are evaluated at the new background
$U^{(n+1)}_k$ with $A_j=0$. By the stage-$n$ form of the action,
\begin{equation*}
A^{(n)}(U^{(n+1)}_k,A_j=0)=-A(x^{(n)},U^{(n+1)}_k)+\mathfrak{E}^{(n)} .
\end{equation*}
By linearity in the weight and by the definition of $\Delta W_{n+1}$ in \eqref{4.18:eq-stage-recursion-5},
\begin{equation*}
-A(x^{(n)},\cdot)
=-A(x^{(n+1)},\cdot)+A(x^{(n+1)}-x^{(n)},\cdot)
=-A(x^{(n+1)},\cdot)+\Delta W_{n+1}.
\end{equation*}
This identity expresses the contribution to $\mathbb{C}^{(n+1)}_k$ that, in \cite[\S1]{B-CMP122a}, comes from the
renormalisation of the Wilson action. By \eqref{4.18:eq-stage-recursion-5}, the logarithm of the fluctuation
integral equals $\mathbb{C}^{(n+1)}_k-\Delta W_{n+1}$.
Adding the three contributions in the definition of $A^{(n+1)}$ gives
\begin{align*}
A^{(n+1)}
&=-A(x^{(n+1)},\cdot)+\Delta W_{n+1}+\mathfrak{E}^{(n)}+E^{(j)}_0
+\mathbb{C}^{(n+1)}_k-\Delta W_{n+1}\\
&=-A(x^{(n+1)},\cdot)+\mathfrak{E}^{(n)}+\mathbb{C}^{(n+1)}_k+E^{(j)}_0 ,
\end{align*}
which is $-A(x^{(n+1)},\cdot)+\mathfrak{E}^{(n+1)}$ by the definition of $\mathfrak{E}^{(n+1)}$ in
\eqref{4.18:eq-stage-recursion-c}. This proves the second equality in \eqref{4.18:eq-stage-recursion-c}.

The induction is now complete. At $n=0$ the form $A^{(0)}=-A(x^{(0)},\cdot)+\mathfrak{E}^{(0)}$ is
\eqref{4.18:eq-stage-recursion-2}, and the step from $n$ to $n+1$ is the computation just made.
\end{proof}

\begin{remark}[The bound required of the localised boundary terms]\label{4.18:rem-required-bound}
Fix a stage $n+1$ with $n+1\in(N-N_0,N]$; equivalently, $j+1=h+n+1\in I_M$. Let
$e^{\mathrm{cur}}_n$ be the current part of the profile excess of
Definition~\ref{4.18:def-coupling-excess}. It is a $C^\infty$ weight with
$0\le e^{\mathrm{cur}}_n\le b_0N_0$, where $b_0$ is the ceiling on the coupling increments
$\beta_i$ of that definition, and its support lies in the set
\begin{equation}\label{4.18:eq-required-bound-1}
\mathcal{E}'_{(j)}
:=\bigcup_{i\in I_M,\ i\le j}
\operatorname{supp}\bigl(\varphi^{\circ}_{Z''_i}-\varphi_{Z''_i}\bigr)
\subset\mathcal{E}' .
\end{equation}
Its stage increment is
\begin{equation}\label{4.18:eq-required-bound-2}
e^{\mathrm{cur}}_{n+1}-e^{\mathrm{cur}}_n
=\beta_{j+1}\bigl(\varphi^{\circ}_{Z''_{j+1}}-\varphi_{Z''_{j+1}}\bigr),
\qquad
0\le e^{\mathrm{cur}}_{n+1}-e^{\mathrm{cur}}_n\le\beta_{j+1}.
\end{equation}
Write $U^{(n)}_k(\cdot)$ for the background field of stage $n$, as a function of the variables
integrated in \cite[(1.60)]{B-CMP122a}, and $U^{(n+1)}_k$ for the corresponding field of
stage $n+1$.

The bound required of the localised boundary terms at stage $n+1$ is the following
assertion. Suppose the stage-$(n+1)$ integral \cite[(1.60)]{B-CMP122a} is formed with the
data of the family's system rather than with those of the companion system
(Definition~\ref{4.18:def-coupling-excess}). Relative to the companion's data, this means
two things.
\begin{itemize}
\item[(i)] The action carries the additional Wilson term $-A(e^{\mathrm{cur}}_n,U^{(n)}_k(\cdot))$,
in the notation of Definition~\ref{4.18:def-stage-recursion}.
\item[(ii)] The member of \cite[(1.62)]{B-CMP122a} gains the term
$A(e^{\mathrm{cur}}_{n+1}-e^{\mathrm{cur}}_n,U^{(n+1)}_k)$.
\end{itemize}
Then the localised pieces $\mathbb{C}^{(n+1)}_k(X)$ of \cite[(1.63)]{B-CMP122a} extend
analytically to the spaces \cite[(1.64)--(1.69)]{B-CMP122a}, and there is an absolute
constant $c_\Omega$ such that for every $X\in\mathbb{D}_m$
\begin{equation}\label{4.18:eq-required-bound-3}
\bigl|\mathbb{C}^{(n+1)}_k\bigl(X,(\mathbb{U},\mathbb{J})\bigr)\bigr|
\le c_\Omega\,C_0\exp\bigl(-(1+3\beta)\kappa\,d_m(X)\bigr).
\end{equation}
Here $\beta$ is the schedule constant of the layered analyticity spaces
(Definition~\ref{4.4:def-post-step-space}), $\kappa$ is the decay parameter among the
auxiliary parameters $\mathfrak{p}$, and $C_0$ is the constant of Ba{\l}aban's bound
\cite[(1.70)]{B-CMP122a}.
The bound obtained in this chapter has $c_\Omega=1$, so that $C_0$ is unchanged, and carries
the decay rates described in (b) below.

Two observations on this formulation are used in what follows.
\begin{itemize}
\item[(a)] In the family's own construction, hypothesis (i) is not an independent datum.
The family's boundary term $\mathbb{B}^{(n)}_k$ also differs from the companion's, by exactly
$+A(e^{\mathrm{cur}}_n,\cdot)$ (the proposition comparing the boundary terms of the two
systems). Hence the additional Wilson term of (i) never occurs as an uncompensated extra
term. Accordingly the assertion above, in (i) and (ii), records only the departure in the
Wilson term; the matching departure of $\mathbb{B}^{(n)}_k$ is understood, and (i) is always
to be read together with it.
\item[(b)] At filing levels $m<k$, the uniform rate $(1+3\beta)\kappa$ in
\eqref{4.18:eq-required-bound-3} is obtained from the theorem supplying the weighted bounds
of the localised terms only under that theorem's rate condition $(F_m)$ at filing level $m$.
The bound established in this chapter carries, at those levels, that theorem's decay rates
$a_m$ in place of $(1+3\beta)\kappa$.
\end{itemize}
\end{remark}

\begin{proposition}[Independence of the stage integrals from the choice of profiles]
\label{4.18:prop-profile-independence}
Assume the following two statements:
\begin{itemize}
\item[($\alpha$)] the profile-independence of the integration structure of the preparatory integrations,
as stated in Definition~\ref{4.18:def-stage-recursion};
\item[($\beta$)] Ba{\l}aban's stage integral and its boundary terms \cite[(1.60)--(1.62)]{B-CMP122a}
take the form of the recursion of Definition~\ref{4.18:def-stage-recursion}, that is, of the displays
\eqref{4.18:eq-stage-recursion-a}, \eqref{4.18:eq-stage-recursion-b} and
\eqref{4.18:eq-stage-recursion-c}, the integral of \cite[(1.60)]{B-CMP122a} being
\eqref{4.18:eq-stage-recursion-a}.
\end{itemize}
Let $\Pi=(\Pi^{(n)})_{0\le n\le N}$ and $\Pi'=(\Pi'^{(n)})_{0\le n\le N}$ be two stage-indexed
systems in the sense of Definition~\ref{4.18:def-coupling-excess}. Suppose that they have the same
members at every stage, namely the members of the sequence in force at that stage, and possibly
different profiles. Write $x_{\Pi^{(n)}}$ and $x_{\Pi'^{(n)}}$ for their coupling functions.
Suppose that both systems are used to write one and the same stage-$0$ action:
\begin{equation*}
A^{(0)} = -A(x_{\Pi^{(0)}},\cdot)+\mathfrak{E}^{(0)}_{\Pi}
= -A(x_{\Pi'^{(0)}},\cdot)+\mathfrak{E}^{(0)}_{\Pi'} ,
\end{equation*}
that is, $\mathfrak{E}^{(0)}_{\Pi'}$ is defined by
\begin{equation}\label{4.18:eq-profile-independence-1}
\mathfrak{E}^{(0)}_{\Pi'} := \mathfrak{E}^{(0)}_{\Pi} + A(x_{\Pi'^{(0)}}-x_{\Pi^{(0)}},\cdot).
\end{equation}
Run the recursion of Definition~\ref{4.18:def-stage-recursion} once with each system. Both runs use
the same couplings $g_j$, collars, fields, characteristic functions and Gaussian measures. Then the
following hold for every $n$:
\begin{itemize}
\item[(a)] $A^{(n)}_{\Pi}=A^{(n)}_{\Pi'}$ as functions of the background and of the exterior
variables;
\item[(b)] the brackets $\mathfrak{K}^{(n)}$ of \eqref{4.18:eq-stage-recursion-b} coincide for
the two systems, the integrands of the stage integral \eqref{4.18:eq-stage-recursion-a} coincide,
and the logarithm $\log\int^{(n+1)}$ of the integral that leads from stage $n$ to stage $n+1$ is
one and the same function for both systems;
\item[(c)] the stage terms of the two systems differ by a sum of single-plaquette terms:
\begin{equation*}
\mathbb{C}^{(n+1)}_{\Pi'}-\mathbb{C}^{(n+1)}_{\Pi}
= A\Bigl(\bigl(x_{\Pi'^{(n+1)}}-x_{\Pi'^{(n)}}\bigr)-\bigl(x_{\Pi^{(n+1)}}-x_{\Pi^{(n)}}\bigr),
\,U^{(n+1)}_k\Bigr);
\end{equation*}
\item[(d)] the non-Wilson parts differ by the Wilson functional of the difference of the coupling
functions:
\begin{equation*}
\mathfrak{E}^{(n)}_{\Pi'}-\mathfrak{E}^{(n)}_{\Pi}=A\bigl(x_{\Pi'^{(n)}}-x_{\Pi^{(n)}},\cdot\bigr).
\end{equation*}
\end{itemize}
Here $0\le n\le N$, and in (b) and (c) $n<N$.
\end{proposition}

\begin{proof}
We argue by induction on $n$. At $n=0$, statement (a) holds by hypothesis, because both systems
write the same action $A^{(0)}$. Statement (d) holds at $n=0$ by the definition
\eqref{4.18:eq-profile-independence-1}.

Assume (a) and (d) at stage $n$. Stage $n$ is the preparatory integration of step $j=h+n$.
Throughout, the displays \eqref{4.18:eq-stage-recursion-a}--\eqref{4.18:eq-stage-recursion-c} are
available by hypothesis ($\beta$).

We first prove (b) at stage $n$. By hypothesis ($\alpha$), every ingredient of the stage integral
\eqref{4.18:eq-stage-recursion-a} is common to $\Pi$ and $\Pi'$, with two possible exceptions: the
stage action $A^{(n)}$, and the split of the stage action into the Wilson part $-g_j^{-2}A$ and the
remainder. The coefficient $g_j^{-2}$ of that split is the same for both systems. Indeed, the
couplings $g_j$ are common to the two runs by hypothesis, and they do not depend on the profiles.
By \eqref{4.18:eq-stage-recursion-b}, the bracket $\mathfrak{K}^{(n)}$ is the difference of
$\boldsymbol{\mathcal{R}}^{(n)}$ at the shifted argument and at the new argument. Here
\begin{equation*}
\boldsymbol{\mathcal{R}}^{(n)}=A^{(n)}+g_j^{-2}A .
\end{equation*}
This function is common to the two systems: $A^{(n)}$ is common by (a) at stage $n$, and $g_j^{-2}$
is common as just observed. Hence the two brackets coincide. The prefactor of the integral is
\begin{equation*}
\exp\bigl[A^{(n)}(U^{(n+1)}_k,A_j=0)+E^{(j)}_0\bigr],
\end{equation*}
where $A_j=0$ means that the integration variable $A_j$ of step $j$ in
\eqref{4.18:eq-stage-recursion-a} is set to zero; this variable is not the member $A_j$ of a system
in the sense of Definition~\ref{4.18:def-coupling-excess}.
It is common to the two systems for two reasons: its first term is common by (a) at stage $n$, and
the normalisation $E^{(j)}_0$ is among the common ingredients by hypothesis ($\alpha$). So every factor of
the integrand of \eqref{4.18:eq-stage-recursion-a} is the same for $\Pi$ and $\Pi'$. The integrands
therefore coincide, and so does the logarithm $\log\int^{(n+1)}$ of the integral. This is (b) at
stage $n$.

Next we prove (c) at stage $n$. By the second step of the recursion of
Definition~\ref{4.18:def-stage-recursion}, the stage terms of the two systems satisfy
\begin{equation*}
\mathbb{C}^{(n+1)}_{\Pi'}-\mathbb{C}^{(n+1)}_{\Pi}
= A\bigl(x_{\Pi'^{(n+1)}}-x_{\Pi'^{(n)}},\cdot\bigr)-A\bigl(x_{\Pi^{(n+1)}}-x_{\Pi^{(n)}},\cdot\bigr),
\end{equation*}
evaluated at $U^{(n+1)}_k$. The functional $A(\cdot,\mathbb{U})$ is linear in the weight
(Definition~\ref{4.18:def-stage-recursion}). The right-hand side is therefore $A$ evaluated at the
difference of the two weights, which is the identity in (c). The value of $A$ at a weight is the
sum over fine plaquettes $p$ of the weight at $x(p)$ times $\mathcal{A}_p$. So the difference is a
sum of single-plaquette terms. This is (c) at stage $n$.

Next we prove (d) at stage $n+1$. By \eqref{4.18:eq-stage-recursion-c}, the difference of the
non-Wilson parts at stage $n+1$ is the difference at stage $n$ plus the difference of the stage
terms. We insert (d) at stage $n$ and (c) at stage $n$, and then use linearity in the weight once
more:
\begin{align*}
\mathfrak{E}^{(n+1)}_{\Pi'}-\mathfrak{E}^{(n+1)}_{\Pi}
&=\bigl[\mathfrak{E}^{(n)}_{\Pi'}-\mathfrak{E}^{(n)}_{\Pi}\bigr]
+\bigl[\mathbb{C}^{(n+1)}_{\Pi'}-\mathbb{C}^{(n+1)}_{\Pi}\bigr]\\
&=A\bigl(x_{\Pi'^{(n)}}-x_{\Pi^{(n)}},\cdot\bigr)
+A\Bigl(\bigl(x_{\Pi'^{(n+1)}}-x_{\Pi'^{(n)}}\bigr)
-\bigl(x_{\Pi^{(n+1)}}-x_{\Pi^{(n)}}\bigr),\cdot\Bigr)\\
&=A\bigl(x_{\Pi'^{(n+1)}}-x_{\Pi^{(n+1)}},\cdot\bigr).
\end{align*}
This is (d) at stage $n+1$.

Finally we prove (a) at stage $n+1$. As at stage $0$, the stage action of each system is its
Wilson part with the system's coupling function as weight, plus its non-Wilson part:
\begin{equation*}
A^{(n+1)}_{\Pi}=-A(x_{\Pi^{(n+1)}},\cdot)+\mathfrak{E}^{(n+1)}_{\Pi},
\end{equation*}
and likewise for $\Pi'$. By linearity in the weight and by (d) at stage $n+1$,
\begin{equation*}
A^{(n+1)}_{\Pi'}-A^{(n+1)}_{\Pi}
=-A\bigl(x_{\Pi'^{(n+1)}}-x_{\Pi^{(n+1)}},\cdot\bigr)
+\bigl[\mathfrak{E}^{(n+1)}_{\Pi'}-\mathfrak{E}^{(n+1)}_{\Pi}\bigr]=0 .
\end{equation*}
This is (a) at stage $n+1$, and the induction is complete.
\end{proof}

\begin{corollary}[The stage integrals of the family are those of its companion]
Take for $\Pi$ the family's systems and for $\Pi'$ the companion systems $\Pi^{\circ}$ of
Definition~\ref{4.18:def-coupling-excess}. For the stage-$0$ action take $A^{(0)}$, the effective
action at level $k$ of the family's construction, written in both ways, the companion's non-Wilson
part at stage $0$ being
\begin{equation*}
\mathfrak{E}^{(0),\circ}:=\mathfrak{E}^{(0)}-A(e_0,\cdot)
=\mathfrak{E}^{(0)}-A(\mathfrak{d}_{\mathrm{old}},\cdot).
\end{equation*}
Let $\int^{(n+1)}$ denote the integral of \eqref{4.18:eq-stage-recursion-a} (the integral of
\cite[(1.60)]{B-CMP122a}), and let the superscript $\circ$ mark the companion's quantities. Write
$\mathbb{B}^{(n),\mathrm{old}}_k$ for the share, in the family's boundary term $\mathbb{B}^{(n)}_k$,
of the older counterterms rewritten with the companion's profiles. Then
at every stage the family's integral is the companion's, and the boundary and stage terms are
related as follows:
\begin{equation}\label{4.18:eq-profile-independence-a}
\begin{aligned}
\int^{(n+1)}(\text{family})
&=\int^{(n+1),\circ},\\
\mathbb{B}^{(n)}_k(\text{family})
&=\mathbb{B}^{(n),\circ}_k+A(e_n-e_0,\cdot)+\mathbb{B}^{(n),\mathrm{old}}_k,\\
\mathbb{C}^{(n+1)}_k(\text{family})
&=\mathbb{C}^{(n+1),\circ}_k+A\bigl(e_{n+1}-e_n,U^{(n+1)}_k\bigr),
\end{aligned}
\end{equation}
where the share $\mathbb{B}^{(n),\mathrm{old}}_k$ is zero under the condition of the chapter's setting under
which the old excess weight vanishes (referred to below as the vanishing condition for the old
excess). In particular, the additional term $-A(e^{\mathrm{cur}}_n,\cdot)$ in the exponent of
\eqref{4.18:eq-stage-recursion-a}, beyond the companion's terms, that hypothesis~(i) of the lemma
on the stage integrals with the extra weight introduces, is compensated identically, before
any expansion, by the additional term $+A(e_n-e_0,\cdot)=+A(e^{\mathrm{cur}}_n,\cdot)$ inside the
family's $\mathbb{B}^{(n)}_k$.
\end{corollary}

\begin{proof}
The second equality in the definition of $\mathfrak{E}^{(0),\circ}$ holds because the current part
$e^{\mathrm{cur}}_0$ of $e_0$ vanishes, as noted at the end of this proof. Under the vanishing
condition for the old excess, $e_0=0$, and then $\mathfrak{E}^{(0),\circ}=\mathfrak{E}^{(0)}$.

Let $(\mathsf{t}_z)_z$ be the partition of unity by tent functions, so that $\sum_z\mathsf{t}_z=1$.
For an older member $\ell$ made of $M$-cubes, linearity in the weight gives, tent by tent,
\begin{equation*}
-\beta_\ell A(\mathsf{t}_z\varphi_\ell,\cdot)
-\beta_\ell A\bigl(\mathsf{t}_z(\varphi^{\circ}_\ell-\varphi_\ell),\cdot\bigr)
=-\beta_\ell A(\mathsf{t}_z\varphi^{\circ}_\ell,\cdot).
\end{equation*}
Summing over $z$ and using $\sum_z\mathsf{t}_z=1$, we see that $\mathfrak{E}^{(0),\circ}$ is
$\mathfrak{E}^{(0)}$ with one change: the counterterm profiles of the older members made of
$M$-cubes are replaced by those of the companion. In other words, $\mathfrak{E}^{(0),\circ}$ is the
non-Wilson part written in the same standard form as $\mathfrak{E}^{(0)}$, with the companion's
profiles in these counterterms; the rewriting rests on the telescoped expression
\eqref{4.18:eq-coupling-excess-a} of the coupling function $x_{\Pi}$ of
Definition~\ref{4.18:def-coupling-excess}.

Proposition~\ref{4.18:prop-profile-independence} now applies to the family's systems and the
companion systems, with the stage-$0$ action written in both ways as above, and yields at every
stage the relations \eqref{4.18:eq-profile-independence-a}.

For the last assertion, the additional term inside the family's $\mathbb{B}^{(n)}_k$ consists of
the weight-difference members $A(e_i-e_{i-1},\cdot)$ of the earlier stages $i\le n$, summed, by
\eqref{4.18:eq-coupling-excess-d}; by linearity in the weight this sum is $A(e_n-e_0,\cdot)$.
Since $e_i=e^{\mathrm{cur}}_i+\mathfrak{d}_{\mathrm{old}}$ for every $i$
(Definition~\ref{4.18:def-coupling-excess}) and $e^{\mathrm{cur}}_0=0$, we have
$e_n-e_0=e^{\mathrm{cur}}_n$, so this term is $+A(e^{\mathrm{cur}}_n,\cdot)$ and compensates the
additional term $-A(e^{\mathrm{cur}}_n,\cdot)$ identically. Here $e^{\mathrm{cur}}_0=0$, because at
stage $0$ no current level is prepared: $h<k_0+1$, which follows from $N>N_0$.
\end{proof}

\begin{lemma}[Support and size of the weight-difference term]\label{4.18:lem-weight-difference}
Work in the setting of Definitions~\ref{4.18:def-prepared-sequence},
\ref{4.18:def-profiles-admissible} and~\ref{4.18:def-coupling-excess}. Fix a stage $n+1$ of the
preparatory integrations, write $j=h+n$, and suppose $j+1\in I_M$. Assume the following.
\begin{enumerate}
\item[(i)] The flow relations of the couplings in Definition~\ref{4.18:def-coupling-excess}; in
particular $0\le\beta_i\le b_0$ for $i\le k$.
\item[(ii)] The properties of the smoothed indicators $\Phi_{\mathcal{P},\sigma}$ stated in
Definition~\ref{4.18:def-profiles-admissible}: $\Phi_{\mathcal{P},\sigma}$ is smooth with values in
$[0,1]$, vanishes at every $x$ with $\operatorname{dist}_\infty(x,D)\le\sigma/8$ for some
$D\in\mathcal{P}$ (in particular on $\bigcup\mathcal{P}$), and differs from $1$ only at points $x$
with $\operatorname{dist}_\infty(x,D)<\sigma/4$ for some $D\in\mathcal{P}$.
\item[(iii)] The relations \eqref{4.18:eq-prepared-sequence-a} and
\eqref{4.18:eq-prepared-sequence-b} of Definition~\ref{4.18:def-prepared-sequence}.
\item[(iv)] The ten-layer growth of the large-field regions at the levels made by the operation
$\mathbf{T}$: at every such level $i>h$ one has $Z_i\supset(Z_{i-1})^{+10a_i}$.
\item[(v)] The nesting of the prepared sequence: $Z''_{j+1}\subset Z''_k$.
\item[(vi)] The position of $Z''_k$ inside $Z$: one has
$\operatorname{dist}_\infty(Z''_k,Z^c)\ge10a_k$; and, writing
$\mathcal{N}:=(Z''_k)^{+\frac54a_k}$, the intersection of $\mathcal{N}$ with the union of the
level-$k$ cells of $Z$ lies in $Z^{\sim-8}$.
\item[(vii)] The chain of inclusions at the levels made by $\mathbf{T}$: if a level $i+1\le k$ is
made by $\mathbf{T}$ inside $Z$, then $\Omega_{i+1}\subset(Z'^{\sim5}_{i})^c$, where $Z'_{i}$
denotes the union of the cubes of side $La_{i+1}$ meeting $Z_{i}$.
\item[(viii)] The filing of terms by level, modelled on \cite[(1.63)]{B-CMP122a}: at every stage
$n+1$, with $j=h+n$, a cube $C\in\pi_{j+1}$ that meets $R_n$ and is contained in
$\Omega^c_{j+2}$ is filed at level $m(C)=j+1$.
\end{enumerate}
Put
\begin{equation}\label{4.18:eq-weight-difference-1}
D_{j+1}:=\bigl\{x:\ \bigl(\varphi^{\circ}_{Z''_{j+1}}-\varphi_{Z''_{j+1}}\bigr)(x)\neq0\bigr\}.
\end{equation}
Then the following hold.
\begin{enumerate}
\item[(a)] One has
\begin{equation}\label{4.18:eq-weight-difference-2}
w_{n+1}=w^{\circ}_{n+1}+(e_{n+1}-e_n),\qquad
0\le e_{n+1}-e_n\le\beta_{j+1}\mathbf{1}_{D_{j+1}},
\end{equation}
and, pointwise, $|w_{n+1}|\le\beta_{j+1}$ and $|w^{\circ}_{n+1}|\le\beta_{j+1}$.
\item[(b)] One has
\begin{equation}\label{4.18:eq-weight-difference-3}
D_{j+1}\subset\bigl([Z''_{j+1}]_{j+1}\bigr)^{+\frac14a_{j+1}}\setminus Z''_{j+1}
\subset(Z_{j+1}\cap Z)\setminus Z''_{j+1}=R_n\cap Z,
\end{equation}
and $D_{j+1}\subset Z^{\sim-8}$.
\item[(c)] Consequently
$\operatorname{supp}w_{n+1}\subset\operatorname{supp}w^{\circ}_{n+1}\cup D_{j+1}$. Inside $Z$ both
supports lie in the union of the closure of $R_n$ with the set where
$\varphi_{Z_{j+1}}\neq1$. Outside $Z$ one has $w_{n+1}=w^{\circ}_{n+1}$; this function is
non-zero there only at the last stage, and then only on the ramp, common to the two systems, of
the factor attached to the cells of $Z$ itself.
\item[(d)] Every cube $C\in\pi_{j+1}$ with $C\cap D_{j+1}\neq\emptyset$ satisfies
$C\cap R_n\neq\emptyset$ and $C\subset\Omega^{c}_{j+2}$. Hence, in the filing of
hypothesis~(viii), $m(C)=j+1$ and $d_{j+1}(C)=0$.
\end{enumerate}
\end{lemma}

\begin{proof}
\emph{(a)} The identity in \eqref{4.18:eq-weight-difference-2} is
\eqref{4.18:eq-coupling-excess-d} at stage $n+1$. The sign of $e_{n+1}-e_n$ and the bound
$e_{n+1}-e_n\le\beta_{j+1}\mathbf{1}_{D_{j+1}}$ follow from \eqref{4.18:eq-profiles-admissible-b},
applied to the member $Z''_{j+1}$, together with $0\le\beta_{j+1}\le b_0$, which is part of
hypothesis~(i). Further,
\begin{equation*}
|w_{n+1}|=\beta_{j+1}\bigl|\varphi_{Z''_{j+1}}-\varphi_{Z_{j+1}}\bigr|\le\beta_{j+1},
\end{equation*}
because both profiles take values in $[0,1]$, so that their difference has absolute value at
most $1$, and $\beta_{j+1}\ge0$. The same argument, with the companion profiles in place of the
family's profiles, gives $|w^{\circ}_{n+1}|\le\beta_{j+1}$.

\emph{(b), the first inclusion.} By hypothesis~(ii), the family's profile
$\varphi_{Z''_{j+1}}$ vanishes on $[Z''_{j+1}]_{j+1}$, a set containing $Z''_{j+1}$, and equals
$1$ off $([Z''_{j+1}]_{j+1})^{+\frac14a_{j+1}}$. Again by hypothesis~(ii), the companion profile
$\varphi^{\circ}_{Z''_{j+1}}$ vanishes on $Z''_{j+1}$ and equals $1$ off
$(Z''_{j+1})^{+\frac14Mu_{j+1}}$. Since $Z''_{j+1}\subset[Z''_{j+1}]_{j+1}$ and
$Mu_{j+1}\le a_{j+1}$, one has
\begin{equation*}
(Z''_{j+1})^{+\frac14Mu_{j+1}}\subset\bigl([Z''_{j+1}]_{j+1}\bigr)^{+\frac14a_{j+1}}.
\end{equation*}
Hence both profiles vanish on $Z''_{j+1}$ and both equal $1$ off
$([Z''_{j+1}]_{j+1})^{+\frac14a_{j+1}}$. Their difference therefore vanishes on both sets, and
this gives the first inclusion in \eqref{4.18:eq-weight-difference-3}.

\emph{(b), the second inclusion.} A cell of $[Z''_{j+1}]_{j+1}$ holds a cube of $Z''_{j+1}$, and
$Z''_{j+1}\subset Z''_k$ by hypothesis~(v). By \eqref{4.18:eq-scales-partitions-a} such a cell
has side $a_{j+1}$. It contains a point of $Z''_k$, so each of its points lies at
$\operatorname{dist}_\infty$-distance at most $a_{j+1}$ from $Z''_k$. Therefore
\begin{equation}\label{4.18:eq-weight-difference-4}
\bigl([Z''_{j+1}]_{j+1}\bigr)^{+\frac14a_{j+1}}\subset(Z''_k)^{+\frac54a_{j+1}}.
\end{equation}
We now distinguish two cases.

\emph{Case $j+1\ge k_0+2$.} Here
$Z''_k=(\Omega^{\sim5}_{k_0+1})^c\cap Z$ (Definition~\ref{4.18:def-prepared-sequence}). Since
$\Omega_{k_0+1}\supset\Lambda_{k_0+1}$,
\begin{equation*}
Z''_k=(\Omega^{\sim5}_{k_0+1})^c\cap Z\subset\Omega^{c}_{k_0+1}\subset Z_{k_0+1}.
\end{equation*}
The level $j+1$ exceeds $h$ and is made by $\mathbf{T}$, so hypothesis~(iv) applies at level
$j+1$. Moreover $Z_j\supset Z_{k_0+1}$ because $j\ge k_0+1$. Together these give
\begin{align*}
Z_{j+1}&\supset(Z_j)^{+10a_{j+1}}\supset(Z_{k_0+1})^{+10a_{j+1}}\\
&\supset(Z''_k)^{+10a_{j+1}}\supset(Z''_k)^{+\frac54a_{j+1}}.
\end{align*}
With \eqref{4.18:eq-weight-difference-4}, the set $([Z''_{j+1}]_{j+1})^{+\frac14a_{j+1}}$ lies in
$Z_{j+1}$.

\emph{Case $j+1=k_0+1$.} By \eqref{4.18:eq-prepared-sequence-a}, $a_{k_0+1}=M$. The
level-$(k_0+1)$ cells are therefore the cubes of $\mathcal{Q}_M$, and $Z''_k$ is a union of such
cubes \cite{B-CMP122a}. A cell holding a cube of $Z''_{k_0+1}\subset Z''_k$ meets the interior of
$Z''_k$. Being a cube of $\mathcal{Q}_M$, it is then one of the cubes of which $Z''_k$ is the
union. Hence $[Z''_{k_0+1}]_{k_0+1}\subset Z''_k$, and by \eqref{4.18:eq-prepared-sequence-b}
\begin{align*}
\bigl([Z''_{k_0+1}]_{k_0+1}\bigr)^{+\frac14M}\cap Z
&\subset\bigl\{x:\ \operatorname{dist}_\infty(x,Z''_k)<\tfrac14M\bigr\}\cap Z\\
&\subset\Lambda\subset Z_{k_0+1}\cap Z.
\end{align*}

\emph{Both cases.} The set $([Z''_{j+1}]_{j+1})^{+\frac14a_{j+1}}$ lies in $Z$, and indeed in
$Z^{\sim-8}$. Since $a_{j+1}\le a_k$, one has
\begin{equation*}
(Z''_k)^{+\frac54a_{j+1}}\subset(Z''_k)^{+\frac54a_k}=\mathcal{N}.
\end{equation*}
By hypothesis~(vi), the intersection of $\mathcal{N}$ with the cells of $Z$ lies in
$Z^{\sim-8}$. Every point of $\mathcal{N}$ lies at
$\operatorname{dist}_\infty$-distance less than $\frac54a_k$ from $Z''_k$, and $\frac54a_k<10a_k$;
hence the bound $\operatorname{dist}_\infty(Z''_k,Z^c)\ge10a_k$ of hypothesis~(vi) gives
$\mathcal{N}\subset Z$.

By \eqref{4.18:eq-weight-difference-4} the set $([Z''_{j+1}]_{j+1})^{+\frac14a_{j+1}}$ lies in
$\mathcal{N}$, hence in $Z$. Since $Z$ is the union of its level-$k$ cells, the set lies in the
intersection of $\mathcal{N}$ with the cells of $Z$, and therefore in $Z^{\sim-8}$. In the first
case it also lies in $Z_{j+1}$. In the second case, since the set lies in $Z$, it coincides with
its intersection with $Z$, which lies in $Z_{k_0+1}\cap Z=Z_{j+1}\cap Z$. In both
cases
\begin{equation*}
\bigl([Z''_{j+1}]_{j+1}\bigr)^{+\frac14a_{j+1}}\setminus Z''_{j+1}
\subset(Z_{j+1}\cap Z)\setminus Z''_{j+1}.
\end{equation*}
Since both profiles vanish on $Z''_{j+1}$, $D_{j+1}\subset(Z_{j+1}\cap Z)\setminus Z''_{j+1}$.
The last set equals $R_n\cap Z$ because $R_n=Z_{j+1}\setminus Z''_{j+1}$. This proves
\eqref{4.18:eq-weight-difference-3}. Finally, $D_{j+1}$ lies in
$([Z''_{j+1}]_{j+1})^{+\frac14a_{j+1}}$, which lies in $Z^{\sim-8}$.

\emph{(c)} The first statement follows from (a). Where $w_{n+1}(x)\neq0$, either
$w^{\circ}_{n+1}(x)\neq0$, or $(e_{n+1}-e_n)(x)\neq0$, and then $x\in D_{j+1}$ by
\eqref{4.18:eq-weight-difference-2}.

\emph{Inside $Z$.} Here $w_{n+1}=-\beta_{j+1}(\varphi_{Z''_{j+1}}-\varphi_{Z_{j+1}})$. On
$Z_{j+1}\cap Z$ the profile $\varphi_{Z_{j+1}}$ vanishes, because $Z_{j+1}\cap Z$ is a union of
its defining cells. There, therefore, $w_{n+1}=-\beta_{j+1}\varphi_{Z''_{j+1}}$. This function
vanishes on $[Z''_{j+1}]_{j+1}$ and is supported in the closure of
$(Z_{j+1}\cap Z)\setminus Z''_{j+1}=R_n\cap Z$, hence in the closure of $R_n$. Off $Z_{j+1}$ and
within $Z$ one has
$\varphi_{Z''_{j+1}}=1$, since by (b) its transition region lies inside $Z_{j+1}$. There
\begin{equation*}
w_{n+1}=-\beta_{j+1}\bigl(1-\varphi_{Z_{j+1}}\bigr),
\end{equation*}
which is supported where $\varphi_{Z_{j+1}}\neq1$, a set common to the family and the companion.
The same description, with $Z''_{j+1}$ in place of $[Z''_{j+1}]_{j+1}$, holds for
$w^{\circ}_{n+1}$.

\emph{Outside $Z$, when $j+1<k$.} The factor attached to $Z_{j+1}\cap Z$ and both factors
attached to $Z''_{j+1}$ equal $1$ outside $Z$. Their transition regions lie within
$a_{j+1}\le a_k$ of sets lying at least eight cell layers inside $Z$. Hence
$w_{n+1}=w^{\circ}_{n+1}=0$ outside $Z$.

\emph{Outside $Z$, when $j+1=k$, the last stage.} Both factors attached to $Z''_k$ equal $1$
outside $Z$. Let $\mathcal{P}$ be the family of level-$k$ cells of $Z$. The factor
$\Phi_{\mathcal{P},a_k}$ of the member of stage $N-1$ is common to the two systems, and it ramps
on $Z^{+\frac14a_k}\setminus Z$. Write $\Psi$ for the product of the remaining factors in the
expression of $w_N$ there; these factors are likewise common to the two systems. There,
therefore,
\begin{equation*}
w_N=w^{\circ}_N=-\beta_k\bigl(1-\Phi_{\mathcal{P},a_k}\bigr)\Psi.
\end{equation*}

\emph{(d)} Let $C\in\pi_{j+1}$ with $C\cap D_{j+1}\neq\emptyset$. Since $D_{j+1}\subset R_n$ by
(b), $C\cap R_n\neq\emptyset$ (the cubes of $\pi_{j+1}$ being closed). The cube $C$ has side
$Mu_{j+1}\le a_{j+1}$ and meets $Z_{j+1}$, so $C$ lies in the closure of
$(Z_{j+1})^{+a_{j+1}}$.

If $j+1=k$, then $\Omega_{k+1}=\emptyset$, and $C\subset\Omega^c_{k+1}$ holds trivially.

If $j+1<k$, the level $j+2\le k$ is made by $\mathbf{T}$ inside $Z$. With $Z'_{j+1}$ as in
hypothesis~(vii), applied with $i=j+1$, that hypothesis gives
$\Omega_{j+2}\subset(Z'^{\sim5}_{j+1})^c$, hence
$\Omega^c_{j+2}\supset Z'^{\sim5}_{j+1}$. The set $Z'^{\sim5}_{j+1}$ contains the closure of
$(Z_{j+1})^{+5La_{j+2}}$. Since $a_{j+1}\le a_{j+2}<5La_{j+2}$,
\begin{equation*}
\Omega^c_{j+2}\supset\overline{(Z_{j+1})^{+5La_{j+2}}}\supset\overline{(Z_{j+1})^{+a_{j+1}}}
\supset C.
\end{equation*}

In both cases $C$ meets $R_n$ and lies in $\Omega^c_{j+2}$. Hypothesis~(viii) then assigns
$m(C)=j+1$. A single cube has
$d_{j+1}(C)=0$, since its spanning tree is a point.
\end{proof}

\begin{remark}[The size of one cube's share of the weight-difference term]\label{4.18:rem-per-cube-size}
Fix the stage $n+1$ and the level $j+1$ as in Lemma~\ref{4.18:lem-weight-difference}. Let $C\in\pi_{j+1}$, and let $(\mathbb{U},\mathbb{J})$ lie in Ba{\l}aban's analyticity space $\tilde{U}^{(n+1)c}_k(C,\tilde\alpha_0,\tilde\alpha_1)$ of \cite{B-CMP122a}. Let $\alpha_{0,j+1},\alpha_{1,j+1}$ be the radii at level $j+1$ (Lemma~\ref{4.4:lem-radii-ratio}). Let $C_{\mathrm{reg}}$ and $\hat{L}_0$ be the constants of the regularity bound on the stage analyticity space $\tilde{U}^{(n+1)c}_k(C,\tilde\alpha_0,\tilde\alpha_1)$, and let $\beta_0$ be the exponent of $N_0$ there (it is unrelated to the couplings $\beta_{j+1}$). Put
\begin{equation}\label{4.18:eq-per-cube-size-1}
\varrho:=C_{\mathrm{reg}}\,\hat{L}_0^{2N_0}\,(\alpha_{0,j+1}+\alpha_{1,j+1}),
\end{equation}
a letter local to this remark, distinct from the $\varrho$ of Notation~\ref{0.2:not-glossary}.
By \cite{B-CMP122a},
\begin{equation*}
\varrho\le C_{\mathrm{reg}}\,L(L+1)\,N_0^{\beta_0}\,R_k\,(\alpha_{0,j+1}+\alpha_{1,j+1}).
\end{equation*}
Assume the following two hypotheses.
\begin{itemize}
\item[(i)] The regularity bound on the stage analyticity space: for every plaquette $p\subset C$, $|\mathbb{U}(\partial p)-1|\le\varrho$.
\item[(ii)] The quadratic bound on Ba{\l}aban's plaquette function $\mathcal{A}_p$ of \cite[(0.1)]{B-CMP109}, with constant $c_{\mathfrak{a}}$: if $|\mathbb{U}(\partial p)-1|\le\frac12$, then $|\mathcal{A}_p(\mathbb{U})|\le c_{\mathfrak{a}}\,|\mathbb{U}(\partial p)-1|^2$.
\end{itemize}
Then, provided $\varrho\le\frac12$,
\begin{equation}\label{4.18:eq-per-cube-size-2}
\bigl|A(w_{n+1}\mathbf{1}_C,\mathbb{U})\bigr|
\le\beta_{j+1}\sum_{p\subset C}|\mathcal{A}_p(\mathbb{U})|
\le\beta_{j+1}\cdot 6c_{\mathfrak{a}}(LM)^4\varrho^2 .
\end{equation}
The right-hand side of \eqref{4.18:eq-per-cube-size-2} is $g_{j+1}^2$ times powers of $\log g^{-2}$, $L$ and $M$. It is therefore smaller than $\vartheta$ --- here the constant of the bound on the stage term $\mathbb{C}^{(n+1)}_k$, which is proportional to $(\log g_k^{-2})^{p_1-q_*}$, and not the separation-window parameter of the glossary --- as soon as $g_k$ lies below a ceiling, which it does under the condition $\gamma\le\gamma_0$. This corroborates, with different constants, the bound by $\vartheta$ on each single-cube term of the weight-difference term made in bounding the stage term $\mathbb{C}^{(n+1)}_k$; both are upper bounds for the same quantity.
\end{remark}

\begin{proof}
The weighted Wilson functional $A(w,\mathbb{U})$ is the sum over plaquettes $p$ of the weight at $p$ times $\mathcal{A}_p(\mathbb{U})$. The weight $w_{n+1}\mathbf{1}_C$ vanishes on plaquettes not contained in $C$. On plaquettes contained in $C$ its modulus is at most $\beta_{j+1}$, by the pointwise bound $|w_{n+1}|\le\beta_{j+1}$ of part (a) of Lemma~\ref{4.18:lem-weight-difference}. The triangle inequality now gives the first inequality in \eqref{4.18:eq-per-cube-size-2}.

For the second inequality, let $p\subset C$. By hypothesis (i) and the proviso, $|\mathbb{U}(\partial p)-1|\le\varrho\le\frac12$, so hypothesis (ii) applies to $p$, and
\begin{equation*}
|\mathcal{A}_p(\mathbb{U})|\le c_{\mathfrak{a}}\,|\mathbb{U}(\partial p)-1|^2\le c_{\mathfrak{a}}\varrho^2 .
\end{equation*}

It remains to count the plaquettes in $C$, measuring $|C|$ in units of the scale $u_j$. The cube $C\in\pi_{j+1}$ has side $Mu_{j+1}=LM\,u_j$. Hence it consists of $(LM)^4$ unit cells of the lattice of spacing $u_j$. In dimension four each unit cell carries six plaquettes at its base point, one for each of the $\binom{4}{2}=6$ coordinate planes. So, counting each plaquette at its base point, $C$ contains $6(LM)^4$ plaquettes.

Summing the bound $c_{\mathfrak{a}}\varrho^2$ over these plaquettes and multiplying by $\beta_{j+1}$ gives the second inequality in \eqref{4.18:eq-per-cube-size-2}.
\end{proof}

\begin{remark}[An uncompensated extra plaquette term would also be absorbed]\label{4.18:rem-extra-term}
Let $e$ be a weight with $0\le e\le b_0N_0$ and $\operatorname{supp}e\subset\mathcal{E}'$, where
$\mathcal{E}'\subset Z^{\sim -8}\cap Z_{j+1}$. In what follows, \emph{the stage argument} is the
argument from which the stage estimate of this chapter is obtained, with its constants $K_S$,
$C_2$, $\kappa_1$, $\vartheta$, its ceiling $\gamma\le\gamma_0$, and the constants $C_0$, $E_0$,
$B_0$, which together with $2$ make up the factor $C_0+E_0+B_0+2$ in the condition fixing that
ceiling. Suppose the
exponent contained, in addition, the uncompensated term $-A(e,\cdot)$. Then the stage argument
absorbs this term as an additional family of single-cube constituents
\begin{equation}\label{4.18:eq-extra-term-1}
  T_{\square}:=-A(e\mathbf{1}_{\square},\cdot),
\end{equation}
where $\square$ runs over the cubes, each in its proper scale, that meet $\mathcal{E}'$. Put
\begin{equation}\label{4.18:eq-extra-term-2}
  \tau:=b_0N_0\cdot 6c_{\mathfrak{a}}(LM)^4\varrho^2 ,
\end{equation}
where $\varrho$ is the local regularity parameter of Remark~\ref{4.18:rem-per-cube-size}
evaluated with the enlarged radii of the spaces on which the stage argument works; the smallness
condition on $\varrho$ stated in that remark is assumed.
The conclusion of the stage argument then persists with two changes. First, the constant bounding
its local sums over the blocks of the cover acquires the additional summand
$K_Se^{C_2\kappa_1}\cdot 6^d\tau\cdot\vartheta$. Second, the factor $C_0+E_0+B_0+2$ in the
condition fixing the ceiling on $\gamma$ is replaced by $C_0+E_0+B_0+2+6^d\tau$. Under $\tau\le 1$
this is a strengthened ceiling on $\gamma$, and no other input of the argument is changed.

No term of this kind is present in the exponent treated in this chapter, and the remark is not
used anywhere.
\end{remark}

\begin{proof}
There are three points to check.

\emph{Size of one constituent.} Each $T_{\square}$ is an entire function of the complexified bond
field, because $A(e\mathbf{1}_{\square},\cdot)$ is a finite weighted sum of the plaquette functions
$\mathcal{A}_p$. To bound it, repeat the computation of Remark~\ref{4.18:rem-per-cube-size} with
the weight $w_{n+1}\mathbf{1}_C$ replaced by $e\mathbf{1}_{\square}$. In that computation
$\varrho$ is the local regularity parameter evaluated with the enlarged radii of the spaces on which
the stage argument works, and the smallness condition on $\varrho$ stated there is assumed. Since
$0\le e\le b_0N_0$, the factor $\beta_{j+1}$ of that computation is replaced by $b_0N_0$. This
gives, for every $\mathbb{U}$ in those spaces,
\begin{equation}\label{4.18:eq-extra-term-3}
  |T_{\square}(\mathbb{U})|\le b_0N_0\sum_{p\subset\square}|\mathcal{A}_p(\mathbb{U})|
  \le b_0N_0\cdot 6c_{\mathfrak{a}}(LM)^4\varrho^2=\tau ,
\end{equation}
whence $\|T_{\square}\|_\infty\le\tau$ on those spaces.

\emph{Dependence on the fluctuation field.} The term $T_{\square}$ depends on the fluctuation field
only through the response map. This dependence is therefore controlled by the Cauchy estimate of
the stage argument, in exactly the way the dependence of a constituent of $\mathbb{B}^{(n)}_k$ or
of $\mathbb{C}^{(n)}_k$ is controlled.

\emph{Local sums.} In the local sums over the blocks of the cover, at most $6^d$ of the cubes
$\square$ meet any one block. By \eqref{4.18:eq-extra-term-3}, the new constituents therefore
contribute at most $K_Se^{C_2\kappa_1}\cdot 6^d\tau\cdot\vartheta$ to the constant bounding these
sums. Correspondingly, the factor in the condition on $\gamma$ becomes
$C_0+E_0+B_0+2+6^d\tau$ in place of $C_0+E_0+B_0+2$.

Under $\tau\le 1$ the new factor is bounded by $C_0+E_0+B_0+2+6^d$, so the condition is a
strengthened ceiling on $\gamma$ and nothing else in the stage argument is affected.
\end{proof}

\begin{theorem}[Bounds on the boundary terms of the preparatory integrations]
\label{4.18:thm-preparatory-boundary}
Fix the step $k$. Assume the following.
\begin{itemize}
\item The inductive theorem on one stage of the preparatory integrations of
\cite[(1.60)--(1.70)]{B-CMP122a}, as stated for the profiles of those displays, holds under its
standing hypotheses. In what follows this theorem is
called \emph{the stage theorem}. It supplies the weighted norm $\|\cdot\|$, the analyticity
spaces, the constants $c_G$, $K_L$, $\mathfrak{s}$, $\mathfrak{p}_j$, $\mathfrak{l}_j$, $q$,
$C_2$, $\kappa_1$, $\vartheta$, $K_S'$, the decay rates $a_m$, and the conditions $(K)$,
$(Q^+)$ and $\gamma\le\gamma_0$.
\item The \emph{profile-class hypothesis}: Ba{\l}aban's statements about
\cite[(1.60)--(1.70)]{B-CMP122a}, and the stage theorem, hold for every system whose
profiles are admissible at every stage.
\item The hypotheses, other than the integration structure of the next item, from which
Proposition~\ref{4.18:prop-profile-independence} is derived.
\item The integration structure fixed in Definition~\ref{4.18:def-stage-recursion}.
\item Hypothesis (ii) of Lemma~\ref{4.18:lem-weight-difference}, one of the inputs from which
that lemma is derived.
\end{itemize}
Assume further that the conditions $(K)$, $(Q^+)$ and $\gamma\le\gamma_0$ of the stage
theorem hold at step $k$. Let $Z$ be any component of the class for which the prepared
sequence and the collars \eqref{4.18:eq-prepared-sequence-c} are defined, with its prepared
sequence and its collars as there. Run the
construction with the
profiles $\varphi_A$ at every stage; this is the \emph{prepared system} $\Pi$. Then, for
every stage $n+1$ with $0\le n\le N-1$, and with $j=h+n$, the following hold.
\begin{itemize}
\item[(a)] The boundary term of the prepared system decomposes as
$\mathbb{C}^{(n+1)}_k=\sum_X\mathbb{C}^{(n+1)}_k(X)$, and the sets $X$ satisfy the conditions
of \cite[(1.63)]{B-CMP122a}.
\item[(b)] Each $\mathbb{C}^{(n+1)}_k(X,\cdot)$ is the restriction of a function holomorphic
on $\tilde{U}^{(n+1)c,\Sigma}_k(X)$. It depends on $(\mathbb{U},\mathbb{J})$, and on the
exterior fluctuation fields, only inside $X$.
\item[(c)] The norm satisfies $\|\mathbb{C}^{(n+1)}_k\|\le C_0$, where
$C_0=c_G(d,L,M)+1$ is the constant of the stage theorem. The rates are those of the stage
theorem, where $\beta$ is the schedule constant of the layered analyticity spaces
(Definition~\ref{4.4:def-post-step-space}) and $\kappa$ the decay parameter of the glossary:
$a_k=(1+3\beta)\kappa$, and, for $j+1\le m\le k-1$, $a_m=(1+3\beta)\kappa$ under
$(F_m)$ and $a_m=(1-\beta)\kappa$ otherwise.
\item[(d)] Consequently the bound of \cite[(1.70)]{B-CMP122a} holds in the form
\begin{equation}\label{4.18:eq-preparatory-boundary-1}
\bigl|\mathbb{C}^{(n+1)}_k(X,(\mathbb{U},\mathbb{J}))\bigr|
\le C_0\exp\bigl(-a_m d_m(X)\bigr),\qquad X\in\mathbb{D}_m .
\end{equation}
At $m=k$, and at every $m$ with $j+1\le m\le k-1$ for which $(F_m)$ holds, the rate is
$(1+3\beta)\kappa$ and this is Ba{\l}aban's bound
$C_0e^{-(1+3\beta)\kappa d_m(X)}$. It is the bound required in
Remark~\ref{4.18:rem-required-bound}, with constant equal to $1$.
\item[(e)] The boundary terms of all stages up to $n+1$ satisfy
\begin{equation}\label{4.18:eq-preparatory-boundary-2}
\sup_x\sum_{i\le n+1}\ \sum_{X\ni x}\bigl\|\mathbb{C}^{(i)}_k(X)\bigr\|_\infty
\le 3(N_0+3)C_0 .
\end{equation}
\item[(f)] There is a decomposition $\mathbb{C}^{(n+1)}_k=\mathbb{C}_G+\mathbb{C}_I$. It obeys
the bounds of the stage theorem:
\begin{equation}\label{4.18:eq-preparatory-boundary-3}
\|\mathbb{C}_G\|\le c_G,\qquad
\|\mathbb{C}_G-\mathbb{C}_G\restriction_{(1,0)}\|\le\vartheta,\qquad
\|\mathbb{C}_I\|\le K_S'(N_0+8)C_0\vartheta\le 1 .
\end{equation}
\end{itemize}
For $n+1\le N-N_0$ the boundary term $\mathbb{C}^{(n+1)}_k$ of the prepared system coincides
term by term with that of the companion system. For $n+1>N-N_0$ it differs from the
companion's term by the single-cube terms
$A\bigl((e_{n+1}-e_n)\mathbf{1}_{\square},U^{(n+1)}_k\bigr)$ of
Lemma~\ref{4.18:lem-weight-difference}.
\end{theorem}

\begin{proof}
\emph{The companion system.} Let $\Pi^{\circ}$ be the companion system of
\eqref{4.18:eq-coupling-excess-a}. It is built as follows.
\begin{itemize}
\item At every stage, every member carries a profile $\varphi^{\circ}_A$ of Ba{\l}aban's
class $\mathcal{P}(A)$.
\item The counterterms of the companion's stage-$0$ split carry the companion's profiles, as
in Proposition~\ref{4.18:prop-profile-independence}.
\item Every other ingredient is the construction's own, by the integration structure of
Definition~\ref{4.18:def-stage-recursion}.
\end{itemize}

\emph{The stage theorem applies to the companion at every stage.} We invoke the
profile-class hypothesis. It says that Ba{\l}aban's statements about
\cite[(1.60)--(1.70)]{B-CMP122a}, and the stage theorem, hold for every admissible profile
system. The proof of the stage theorem reads the profiles at two places only. The first is
the counterterms and the support of the weight of the functional $W$. The second is the
supremum and the support of the weight of the increment $\Delta W$. At neither place is the
shape of a profile seen beyond Ba{\l}aban's two conditions on the class $\mathcal{P}$.

Hence the stage theorem applies to the construction of $\Pi^{\circ}$ at every stage. For
$\Pi^{\circ}$, its inductive hypothesis at stage $n$ implies its inductive hypothesis at
stage $n+1$, with the constant $C_0$. Starting from stage $0$, we obtain it inductively at
every stage. This step uses only the stage theorem under its standing hypotheses and the
profile-class hypothesis.

\emph{The stage integral of the prepared system.} Apply
Proposition~\ref{4.18:prop-profile-independence} to the pair $(\Pi,\Pi^{\circ})$. It shows
that the integral of \cite[(1.60)]{B-CMP122a} at stage $n+1$ is, for the prepared system,
the companion's integral. Hence its logarithm has the companion's representation
$\sum_X[\mathbb{C}_G(X)+\mathbb{C}_L(X)]$. This representation is produced by the Gaussian
normalisation lemma and the cluster-expansion proposition of the stage theorem, and it comes
with their bounds and their holomorphy.

\emph{The weight term of the prepared system.} The prepared system's term of
\cite[(1.62)]{B-CMP122a} is $A(w_{n+1},U^{(n+1)}_k)$. Assign each fine plaquette to the cube
of $\pi_{j+1}$ that contains its base point. Since $A(\cdot,\mathbb{U})$ is linear in the
weight (Definition~\ref{4.18:def-stage-recursion}), this gives
\begin{equation}\label{4.18:eq-preparatory-boundary-4}
A\bigl(w_{n+1},U^{(n+1)}_k\bigr)=\sum_{C\in\pi_{j+1}}A\bigl(w_{n+1}\mathbf{1}_C,U^{(n+1)}_k\bigr).
\end{equation}
At the last stage the weight $w_{n+1}$ has a part common to the prepared and the companion
system, its outer-ramp part. It is filed as the stage theorem files it.

\emph{Definition of the terms.} For each $X$ set
\begin{equation}\label{4.18:eq-preparatory-boundary-5}
\mathbb{C}^{(n+1)}_k(X):=\mathbb{C}_G(X)+\mathbb{C}_L(X)
+\begin{cases}A(w_{n+1}\mathbf{1}_C,\cdot)&\text{if }X=C\in\pi_{j+1},\\
0&\text{otherwise.}\end{cases}
\end{equation}
This is exactly the definition of the stage theorem, with the prepared system's weight
$w_{n+1}$ in place of Ba{\l}aban's.

\emph{Clause} (a). The filing of $\mathbb{C}_G$ and $\mathbb{C}_L$ is that of the stage
theorem. The cube terms that carry the part specific to the prepared system are filed in
$\mathbb{D}_{j+1}$ as single cubes $C\in\pi_{j+1}$, that is with $d_{j+1}(C)=0$. By
Lemma~\ref{4.18:lem-weight-difference}(d) they meet the collar $R_n$.
The part common to both systems is filed as in the stage theorem. Hence every $X$ in the
decomposition satisfies the conditions of \cite[(1.63)]{B-CMP122a}.

\emph{Clause} (b). For $\mathbb{C}_G$ and $\mathbb{C}_L$, holomorphy and locality are the
conclusions of the Gaussian normalisation lemma and the cluster-expansion proposition. The
cube term $A(w_{n+1}\mathbf{1}_C,\cdot)$ is a sum of the plaquette functions
$\mathcal{A}_p$ over plaquettes of $C$. These are entire in the plaquette variables
(Definition~\ref{4.18:def-stage-recursion}). So the cube term is entire in the plaquette
variables of $C$, and it depends on $\mathbb{U}$ only inside $C$.

\emph{Clause} (c), \emph{the per-cube bound.} The per-cube bound of the stage theorem uses
only a bound for the supremum of the weight, multiplied by a factor independent of the
weight.
\begin{itemize}
\item Ba{\l}aban's weight $-\beta_{j+1}(\varphi_{Z''_{j+1}}-\varphi_{Z_{j+1}})$, with
admissible profiles, has supremum at most $\beta_{j+1}$.
\item The prepared system's weight $w_{n+1}$ also has supremum at most $\beta_{j+1}$, by
Lemma~\ref{4.18:lem-weight-difference}(a).
\end{itemize}
The factor independent of the weight is $(LM)^4\cdot 4(2\tilde\alpha_0+\tilde\alpha_1)^2$.
Therefore each cube term of the prepared system obeys the same estimate as in the stage
theorem, the supremum being taken over
$(\mathbb{U},\mathbb{J})\in\tilde{U}^{(n+1)c}_k(C,\tilde\alpha_0,\tilde\alpha_1)$:
\begin{equation}\label{4.18:eq-preparatory-boundary-6}
\bigl\|A(w_{n+1}\mathbf{1}_C,\cdot)\bigr\|_\infty
\le\beta_{j+1}(LM)^4\cdot 4(2\tilde\alpha_0+\tilde\alpha_1)^2\le\vartheta .
\end{equation}
The last inequality is the stage theorem's.

\emph{Clause} (c), \emph{the norm bound.} The norm $\|\cdot\|$ is subadditive over families
summed $X$ by $X$. Applied to \eqref{4.18:eq-preparatory-boundary-5}, and using
\eqref{4.18:eq-preparatory-boundary-6} for the third summand, it gives
\begin{align*}
\|\mathbb{C}^{(n+1)}_k\|
&\le\|\mathbb{C}_G\|+\|\mathbb{C}_L\|
+\sup_{C\in\pi_{j+1}}\bigl\|A(w_{n+1}\mathbf{1}_C,\cdot)\bigr\|_\infty\\
&\le c_G+K_L(\mathfrak{s}+\mathfrak{p}_j+\mathfrak{l}_j)M^q
e^{C_2\kappa_1+4\sqrt{d}(2+4\beta)\kappa}+\vartheta\\
&\le c_G+1=C_0 .
\end{align*}
This is the chain of inequalities of the stage theorem. Its second line uses that theorem's
bounds for $\|\mathbb{C}_G\|$ and $\|\mathbb{C}_L\|$. Its last inequality is the one the
stage theorem establishes under $(K)$, $(Q^+)$ and $\gamma\le\gamma_0$. The rates are those
fixed in the step of the cluster-expansion proposition that determines the decay rates. The
cube terms are single cubes, $d_{j+1}(C)=0$, so the decay factor
$e^{-a_{j+1}d_{j+1}(C)}$ equals $1$ for them whatever the rate.

\emph{Clause} (d). By the definition of the weighted norm of the stage theorem, for every
$X\in\mathbb{D}_m$
\begin{equation*}
\bigl|\mathbb{C}^{(n+1)}_k(X,(\mathbb{U},\mathbb{J}))\bigr|
\le\|\mathbb{C}^{(n+1)}_k\|\,e^{-a_m d_m(X)} .
\end{equation*}
Combined with the bound $\|\mathbb{C}^{(n+1)}_k\|\le C_0$ of (c), this gives
\eqref{4.18:eq-preparatory-boundary-1}.

\emph{Clause} (e). The clause of the cumulation lemma of the stage theorem that bounds the
pointwise sum of sup norms uses only two inputs:
\begin{itemize}
\item the bounds $\le C_0$ on the norms of the boundary terms at all stages;
\item the filing condition $X\cap R_{i-1}\neq\emptyset$.
\end{itemize}
Both are available for the prepared system. For every $i\le n+1$, the norm bound holds by
(c) at stage $i$, and the filing condition holds by (a) at stage $i$. Hence that clause
gives \eqref{4.18:eq-preparatory-boundary-2}.

\emph{Clause} (f). The term $\mathbb{C}_G$ is the companion's. Its two bounds in
\eqref{4.18:eq-preparatory-boundary-3} are those of the stage theorem applied to
$\Pi^{\circ}$. Set $\mathbb{C}_I:=\mathbb{C}_L$ plus the cube terms of
\eqref{4.18:eq-preparatory-boundary-5}. The cube terms are bounded by
\eqref{4.18:eq-preparatory-boundary-6}, the same bound they have in the stage theorem, and
they enter in the same place in its estimate. Therefore $\mathbb{C}_I$ obeys the stage
theorem's bound $\|\mathbb{C}_I\|\le K_S'(N_0+8)C_0\vartheta\le 1$.

\emph{The last assertion.} For $j+1\le k_0$ the increment vanishes, $e_{n+1}-e_n=0$, by
\eqref{4.18:eq-coupling-excess-d}. Then \eqref{4.18:eq-preparatory-boundary-5} coincides
term by term with the companion's definition. For $j+1\in I_M$, which are the stages with
$n+1>N-N_0$ (Remark~\ref{4.18:rem-required-bound}), the difference from the companion's
terms is given by \eqref{4.18:eq-profile-independence-a}. It consists of the single-cube
terms $A\bigl((e_{n+1}-e_n)\mathbf{1}_{\square},U^{(n+1)}_k\bigr)$ of
Lemma~\ref{4.18:lem-weight-difference}.
\end{proof}

\begin{remark}[The decay rate at levels below the top]\label{4.18:rem-decay-rates}
The boundary-term lemma on which Theorem~\ref{4.18:thm-preparatory-boundary} rests is stated with
Ba{\l}aban's uniform decay rate $(1+3\beta)\kappa$. Here $\beta$ is the schedule constant of the layered
analyticity spaces (Definition~\ref{4.4:def-post-step-space}) and $\kappa$ is Ba{\l}aban's auxiliary
parameter. The theorem on the stage terms of the preparatory integrations may deliver a weaker rate
at filing levels below the top. At the top level $k$ its rate is $(1+3\beta)\kappa$. At a filing
level $m<k$ its rate is $(1+3\beta)\kappa$ under the condition of that theorem denoted $(F_m)$, and
$(1-\beta)\kappa$ otherwise.
Theorem~\ref{4.18:thm-preparatory-boundary} holds with these rates. It takes this qualification over
from the theorem on the stage terms and adds none of its own.
\end{remark}

\begin{proof}
Let $m<k$ and suppose that condition $(F_m)$ of the theorem on the stage terms fails, so that the
delivered rate is $(1-\beta)\kappa$; we check that this rate suffices for every later use of a term
of level $m$.

The estimate in \cite{B-CMP122b} through which terms of level below $k$ enter the bounds for the
operation $\mathbf{R}$, and every estimate after it, use a term of level $m<k$ only through the
right-hand member of a chain of inequalities that Ba{\l}aban proves at that place. They never use the
rate $(1+3\beta)\kappa$ in any other way. This is
established in the statements accompanying the theorem on the stage terms (two propositions and a
lemma) that identify what the later estimates take from a term of level below $k$.

The same statements show that the rate $(1-\beta)\kappa$ also supplies this right-hand member as soon
as $L\ge 4$. This condition is implied by the requirement $L\ge 8$ on the blocking factor under which
those statements are made. Hence a term of level $m<k$ carrying the
rate $(1-\beta)\kappa$ enters every later estimate exactly as one carrying $(1+3\beta)\kappa$ would.

Since Theorem~\ref{4.18:thm-preparatory-boundary} is proved from the theorem on the stage terms, its
stage terms carry these rates, and the dependence on $(F_m)$ is inherited, not introduced there.
\end{proof}

\begin{remark}[Scope with respect to older members made of $M$-cubes]
\label{4.18:rem-scope-older-members}
Call the cell hypothesis the condition that every member $Z_i\cap Z$, $i\le h$, of the label in
force (Definition~\ref{4.18:def-prepared-sequence}) is a union of level-$i$ cells.
Theorem~\ref{4.18:thm-preparatory-boundary} does not assume the cell hypothesis.
\begin{enumerate}
\item Under the cell hypothesis, the companion system of profiles $\varphi^{\circ}_A$ coincides
with the family of profiles $\varphi_A$ (the profiles with which the construction of
Theorem~\ref{4.18:thm-preparatory-boundary} is run) at stage $0$ and at the first
$N-N_0$ stages. Consequently, the theorem on the stage terms, the first of the statements that
Theorem~\ref{4.18:thm-preparatory-boundary} takes as hypotheses, is applied to a construction
that differs from the
actual one only in the bookkeeping of the last $N_0$ stages. Moreover, the admissibility
condition on the profiles is used only at the current levels whose members are unions of
$M$-cubes.
\item If the label in force has older members inside $Z$ that are unions of
$M$-cubes (the case treated in the discussion of the older $\mathbf{R}$-operations later in this
chapter), then the companion's stage-$0$ split, that is, the decomposition
$A^{(0)}=-A(x_{\Pi^{(0)}},\cdot)+\mathfrak{E}^{(0)}_{\Pi}$ of the stage-$0$ action of
Proposition~\ref{4.18:prop-profile-independence} for the companion system $\Pi$, rewrites the
counterterm pieces of these members
with admissible profiles, and this rewriting falls within the admissibility condition on the
profiles. No stage increment $w_{n+1}$ of the coupling functions involves those members.
\end{enumerate}
\end{remark}

\begin{proof}
We take the two items in turn.

For the first item: under the cell hypothesis the companion and the family coincide at stage $0$
and at the first $N-N_0$ stages, so the construction to which the theorem on the stage terms, the
first of the statements assumed in Theorem~\ref{4.18:thm-preparatory-boundary}, is applied differs
from the actual construction only in the bookkeeping of the last $N_0$ stages, and the
admissibility condition on the profiles enters only through the current levels whose members are
unions of $M$-cubes.

The second item is about the admissibility condition. That condition requires admissibility
relative to
the prepared members, and a member $Z_i\cap Z$ lying at a top level $i$ of an older operation is
a union of $M$-cubes of scale $i$. The companion and the family have the same members at every
stage and write the same stage-$0$ action, so
Proposition~\ref{4.18:prop-profile-independence} gives equality of their stage integrals. On the
counterterm pieces of such members the companion's profiles are admissible relative to the
prepared members, which is what the admissibility condition asks.

It remains to see why no stage increment $w_{n+1}$ involves the older members. For a stage $n+1$
with
$0\le n\le N-1$ put $j=h+n$. By \eqref{4.18:eq-coupling-excess-c}, the stage increment $w_{n+1}$
of the coupling functions involves only level $j+1$, and
\begin{equation*}
j+1=h+n+1>h.
\end{equation*}
The older members carry levels $i\le h$ and do not change during the preparation. Hence no
increment $w_{n+1}$ involves them.
\end{proof}

\begin{remark}[The final-stage exponent up to older members made of $M$-cubes]
\label{4.18:rem-final-exponent}
Use the notation of Proposition~\ref{4.18:prop-profile-independence}, at the final stage
$n=N$, with the coupling function $x''$, the background $U''_k$ and the boundary terms
$\mathbb{B}''_k$ after the preparation of \cite{B-CMP122b}. Write $\mathfrak{E}^{(N)}$ and
$\mathbb{B}''_k$ for the family's stage term and boundary terms, and $\mathfrak{E}^{(N),\circ}$
and $\mathbb{B}''^{\circ}_k$ for the stage term and boundary terms of the companion
construction. Let $\mathbb{O}''_k$ denote the share of the older members' counterterms. Then
\begin{equation}\label{4.18:eq-final-exponent-2}
-A(x'',U''_k)+\mathbb{B}''_k
=-A\bigl(x''-e^{\mathrm{cur}}_N,\,U''_k\bigr)+\mathbb{B}''^{\circ}_k+\mathbb{O}''_k .
\end{equation}
The family's numerator exponent in \cite[(1.1)]{B-CMP122b} equals the exponent of the
construction on the right-hand side of \eqref{4.18:eq-final-exponent-2}, whose Wilson weight
departs from a weight in Ba{\l}aban's admissible class only at the older members made of
$M$-cubes, and there by $\mathfrak{d}_{\mathrm{old}}$; under the chapter's condition on those
older members there is no departure at all. The equality holds term by term once the
single-plaquette pieces
\begin{equation*}
\pm A\bigl((e_i-e_{i-1})\mathbf{1}_{\square},\,U''_k\bigr)
\end{equation*}
have been moved, all of them evaluated at the final background $U''_k$. Consequently no
estimate of the Wilson term fails for the family's density regarded as a function of the
fields. The two sides of \eqref{4.18:eq-final-exponent-2}, the family's construction and the
companion construction, differ only in which localised term carries the single-plaquette
pieces, and each piece has the size computed in Remark~\ref{4.18:rem-per-cube-size}.
\end{remark}

\begin{proof}
Apply clause~(d) of Proposition~\ref{4.18:prop-profile-independence} at $n=N$. It gives
\begin{equation}\label{4.18:eq-final-exponent-1}
\mathfrak{E}^{(N)}-\mathfrak{E}^{(N),\circ}=A(e_N,\cdot\,).
\end{equation}
With the notation of the remark, identity \eqref{4.18:eq-final-exponent-1}, written out at the
final background, reads \eqref{4.18:eq-final-exponent-2}.

The left-hand side of \eqref{4.18:eq-final-exponent-2} is the family's numerator exponent in
\cite[(1.1)]{B-CMP122b}. On the right-hand side the Wilson weight is $x''-e^{\mathrm{cur}}_N$.
This weight departs from a weight of Ba{\l}aban's admissible class only at the older members
made of $M$-cubes, and the departure there is $\mathfrak{d}_{\mathrm{old}}$. Under the
chapter's condition on those older members there is no departure at all, and the weight is then
itself of Ba{\l}aban's admissible class.

Identity \eqref{4.18:eq-final-exponent-2} compares the two exponents term by term once the
single-plaquette pieces $\pm A((e_i-e_{i-1})\mathbf{1}_{\square},U''_k)$ are moved to the
localised term that holds them. Every term in \eqref{4.18:eq-final-exponent-2},
including these pieces, is evaluated at the same background $U''_k$. This proves the identity
asserted in the remark.

For the consequence, the two sides of \eqref{4.18:eq-final-exponent-2} are the same function of
the fields. Hence an estimate of the Wilson term that holds for the construction on the
right-hand side holds for the family's density as a function. The two sides of
\eqref{4.18:eq-final-exponent-2}, the family's construction and the companion construction,
differ only in which localised term carries the single-plaquette pieces, and the size of each
such piece is the one computed in Remark~\ref{4.18:rem-per-cube-size}.
\end{proof}

\begin{remark}[Older large-field operations inside the component]\label{4.18:rem-older-operations}
Let $Z$ be a component at step $k$ of the class of Ba{\l}aban's construction (a component of this
class, at any step, is called here a \emph{first-class component}), with its label in
force and its prepared sequence, as in Definition~\ref{4.18:def-prepared-sequence}; write $N$ for
the number of preceding steps in the class conditions, put $h=k-N$, and identify $Z_j$ with
$Z_j\cap Z$. Call a \emph{retained region} of an $\mathbf{R}$-operation performed at step $s$ any
component $Y$ of the class of \cite[(1.72)]{B-CMP122a} on which a large-field factor remains
after that operation.
\begin{enumerate}
\item[(a)] Condition (ii) of the class excludes new large-field regions created inside $Z$ at
the levels $h+1,\dots,k$, and requires the previous regions $Z_j\cap Z$ to satisfy condition (i)
on their scales. It imposes nothing on the way in which the large-field regions of levels $\le h$
arose. In particular an $\mathbf{R}$-operation performed at a step $s\le h$ on a first-class
component of that step lying inside what is at step $k$ the component $Z$ is admitted by
conditions (i) and (ii); this includes the case $s=h$.
\item[(b)] The retained regions of such an operation persist as large-field regions of level $s$.
A retained region of a step $s\le h$ lying inside $Z$ is an admissible part of $Z_h\cap Z$.
\item[(c)] Let $\ell$ denote the label in force, let $Y$ be a retained region of an
$\mathbf{R}$-operation at step $s\le h$ lying inside $Z$, let $Z''^{(s)}_i$ be the prepared
sequence and $N_0(s)$ the integer $N_0$ of that operation, and let $Z^{\mathbf{T}}_i$ denote the
members of $\ell$ produced by the operations $\mathbf{T}$. For the levels $i<s$ three
descriptions of the members of $\ell$ inside $Y$ are in question:
\begin{enumerate}
\item[(c1)] no member of $\ell$ lies inside $Y$, that is, $\Lambda_i(\ell)\supseteq Y$ for $i<s$;
\item[(c2)] the members of $\ell$ inside $Y$ are the sets $Z^{\mathbf{T}}_i\cap Y$;
\item[(c3)] the members of $\ell$ inside $Y$ at the levels $i<s$ are the sets
$Z''^{(s)}_i\cap Y$ of the prepared sequence of that operation; for $s-N_0(s)<i<s$ these are
unions of level-$i$ $M$-cubes and not unions of level-$i$ cells.
\end{enumerate}
Under (c1) and under (c2) the cell-geometry condition holds at every step: every member of the
label in force which lies inside $Z$ at a level $\le h$ is a union of cells of its level. In
particular the case of a retained region of an older operation carrying members that are not
unions of cells is vacuous. Under (c3) this case is not vacuous, and the cell-geometry condition
fails in a configuration admitted by conditions (i) and (ii), namely a retained region of an
$\mathbf{R}$-operation at step $s=h$ lying inside $Z$; the passages of \cite{B-CMP122a} and
\cite{B-CMP122b} quoted in the proof below do not exclude it.
\end{enumerate}
\end{remark}

\begin{proof}
(a) Ba{\l}aban introduces the class of components in \cite{B-CMP122a} in the following words:
\begin{quote}
we consider the class of components such that each satisfies the following two properties:
(i) it is contained in a cube of the size $100MR_k$, (ii) in the preceding $N$ renormalization
steps no new large field regions were created inside this component, and the previous regions
contained in it satisfy the condition (i) on the corresponding scales. According to our rule of
construction of the large field regions, for such a component all the regions connected with the
last $N$ steps are rectangular parallelepipeds.
\end{quote}
The preceding $N$ renormalization steps are the one-step operations $\mathbf{T}^{(j)}$,
$j=h,\dots,k-1$, written in \cite[(1.2)]{B-CMP122a}, of which Ba{\l}aban says that they are
``the last $N$ one-step operations \dots\ by the condition (ii) no large field characteristic
functions are included in them''; they are the steps producing the levels $h+1,\dots,k$. The
older operation $\mathbf{T}_h(Z_h)$ of \cite[(1.2)]{B-CMP122a}, acting on the region of level $h$,
is carried unchanged, and \cite[(1.11)]{B-CMP122a} takes
``$Z''_j = Z_j$ for $j = h, h-1,\dots,1$''. Ba{\l}aban's construction moreover contemplates
components with mixed histories: ``if in a component of $Z$ we finish the integrations for some
$n$, because some new large field has been introduced, then we leave the new terms
$\mathbb{B}^{(n)}_k$ as a part of all boundary terms'' \cite{B-CMP122a}. Condition (ii) therefore
speaks of the levels $h+1,\dots,k$ only: it excludes new large-field regions at these levels
inside $Z$, and it requires the previous regions $Z_j\cap Z$ to satisfy (i) on their scales. It
says nothing about how the regions of levels $\le h$ arose. An $\mathbf{R}$-operation at a step
$s\le h$, whose large-field regions are of level $s$, on a first-class component of that step
lying inside $Z$, therefore violates neither (i) nor (ii). For $s=h$ this remains so, because the
step at which the operation is performed is not among the $N=k-h$ preceding steps, which produce
the levels $h+1,\dots,k$.

(b) After an $\mathbf{R}$-operation performed at a step that Ba{\l}aban writes $k$, he states in
\cite[p.~390]{B-CMP122b}:
\begin{quote}
The new large field region is equal to $Z_k\cup\bigcup_{i=1}^{m}Y_i$.
\end{quote}
Here the $Y_i$ form the class of \cite[(1.72)]{B-CMP122a}, and $m$ and $i$ are Ba{\l}aban's
enumeration of its components: each $Y_i$ is a whole first-class component of that step, and it
keeps its large-field characteristic functions and the further factors attached to it in
\cite[(1.72)]{B-CMP122a}, together with the operation $\mathbf{T}'_k(Y_i)$. Hence the retained
regions of
the operation are large-field regions of the level at which it is performed. Applied at a step
$s\le h$, a retained region lying inside $Z$ is a previous region of level $s\le h$ contained in
$Z$, and by (a) its presence is compatible with conditions (i) and (ii); it is thus an admissible
part of $Z_h\cap Z$.

(c) The members of $\ell$ at level $i$ that carry a profile are the complements
$\Lambda_i(\ell)^c$, the profile being formed from the cubes $D$ over the level-$i$ cells of
$[\Lambda_i(\ell)^c]_i$ (Definition~\ref{4.18:def-profiles-admissible}). Which of the three
descriptions holds inside $Y$ for $i<s$ is fixed by the description, in \cite{B-CMP122b}, of the
label after an $\mathbf{R}$-operation in the form \cite[(2.18)]{B-CMP119}; that description is
not derived in this section, and the three descriptions are treated as alternatives.

Description (c1) agrees with Ba{\l}aban's formula for the coupling of the action after the
operation,
\begin{quote}
the new function
\begin{equation*}
1/(g_k(\cdot))^2 = 1/(g''_k(\cdot))^2 Z^c + 1/(g_k^2) Z
\end{equation*}
\end{quote}
\cite{B-CMP122b}; here Ba{\l}aban's $k$ is the step $s$ of the operation, and $Z$ and $Z^c$
are the region of the operation and its complement, read as indicator functions; this function
is constant on
$Y$. Description (c1) also agrees with the rule by which the profiles of the label cease to act
at a retained region: the profile $\varphi^{\ell}_s$ of the member of $\ell$ at level $s$
(Definition~\ref{4.18:def-profiles-admissible}) vanishes on and near $Y$. Description (c2) counts
only the members produced by the operations $\mathbf{T}$; under it the profile
$\varphi^{\ell}_k$ of the member of $\ell$ at level $k$ vanishes on a
retained region $Y\subset\Lambda_k(\ell)^c$, and the coupling function of the label is there the
one in force before the operation. Description (c3) counts among the members of the label the
sets made by the older
$\mathbf{R}$-operation, read as sets, that is, the members of its own prepared sequence; for
$s-N_0(s)<i<s$ these sets are unions of level-$i$ $M$-cubes and not unions of level-$i$ cells
\cite{B-CMP122a}, and the sequences made by the $\mathbf{R}$-operation lie inside $Z$.

For a large-field region of an older $\mathbf{R}$-operation that has been healed, that is, whose
large-field factor has been absorbed back into the small-field format by a later step, the three
descriptions agree: no member of $\ell$ remains inside it, the retained profile is free of that
region, and the coupling function equals the new constant $g_k^{-2}$ on and around that region.

Now assume (c1) or (c2). The members of $\ell$ lying inside $Z$ at a level $\le h$ are then of
two kinds only: sets produced by the operations $\mathbf{T}$, which are unions of cells of their
level by the cell geometry of the sets made by the operations $\mathbf{T}$, and
retained regions of older operations, which are closed unions of cells of their level by the cell
geometry of the retained regions of an $\mathbf{R}$-operation.
Under (c1) no further member lies inside a retained region $Y$ below its level $s$; under (c2)
the members inside $Y$ are the sets $Z^{\mathbf{T}}_i\cap Y$, again unions of cells. Hence every
member of the label in force inside $Z$ at a level $\le h$ is a union of cells of its level. This
holds at every step $k$, so the cell-geometry condition holds at every step, and no member inside
$Z$ is a union of $M$-cubes without being a union of cells: the case of a retained region of an
older operation carrying such members does not occur.

Assume (c3). Then inside a retained region $Y$ of a step $s\le h$ the members of $\ell$ at the
levels $s-N_0(s)<i<s$ are the sets $Z''^{(s)}_i\cap Y$, which are unions of level-$i$ $M$-cubes
and not unions of level-$i$ cells. By (a) and (b), a retained region of an $\mathbf{R}$-operation
at step $s=h$ lying inside $Z$ is admitted by conditions (i) and (ii). Hence the case is not
vacuous, the cell-geometry condition fails in this admitted configuration, and none of the quoted
passages excludes it.
\end{proof}

\begin{lemma}[The window length of an older step against the current cell parameter]\label{4.18:lem-older-window}
Let $r>0$. Assume the coupling-flow statement at all levels $\le k$. Of it only the following
consequences are used, in which $\beta_i$ are the coupling increments of
Definition~\ref{1.2:def-couplings} and $\bar\beta\ge0$ is the constant bounding them in the
coupling-flow statement:
\begin{equation}\label{4.18:eq-older-window-1}
\begin{gathered}
g_m^{-2}\ge g_k^{-2}\quad (m\le k),\\
g_m^{-2}=g_n^{-2}+\sum_{m<i\le n}\beta_i,\qquad \sum_{m<i\le n}\beta_i\le \bar\beta(n-m)\quad (m\le n\le k).
\end{gathered}
\end{equation}
Assume the ceiling $g_k\le e^{-1}$. For every level $m\le k$ let $R_m$ be the least power of $L$ with
$R_m\ge(\log g_m^{-2})^{r}$, where $\log$ is the natural logarithm; this is the cell parameter $R_m$
of Notation~\ref{4.18:not-scales-partitions}. Let $s\le h$, where $h<k$ is
the integer of the large-field operation at step $k$ fixed in
Definition~\ref{4.18:def-prepared-sequence}, and let $N_0(s)\ge 1$ be an integer with
$L^{N_0(s)-1}=R_{s-N_0(s)+1}$; this is the relation satisfied by the integer $N_0$ of the
large-field operation at step $s$ in \cite{B-CMP122a}, whose construction is uniform in the
step index. Then
\begin{equation}\label{4.18:eq-older-window-a}
N_0(s)<c_1(k-s)+2R_k,\qquad c_1:=4+2r(1+\bar\beta)+3r^2\bar\beta.
\end{equation}
The same conclusion holds if $R_m$ is replaced, in the hypothesis on $N_0(s)$ and in
\eqref{4.18:eq-older-window-a}, by numbers $\check{R}_m$ $(m\le k)$ satisfying
\begin{equation}\label{4.18:eq-older-window-2}
(\log g_m^{-2})^{r}\le\check{R}_m<L(\log g_m^{-2})^{r}.
\end{equation}
\end{lemma}

\begin{proof}
Write $y_m:=\log g_m^{-2}$, $\nu:=N_0(s)$, $u:=k-s$ and $s_0:=s-\nu$; then $u\ge 1$, since
$s\le h<k$. By the first line of \eqref{4.18:eq-older-window-1} and the ceiling $g_k\le e^{-1}$,
$g_m^{-2}\ge g_k^{-2}\ge e^{2}$ for $m\le k$, so that
\begin{equation*}
y_m\ge y_k\ge 2,\qquad g_m^2\le g_k^2\le e^{-2}<1 .
\end{equation*}

\emph{Step 1 (the cell parameter).} Since $r>0$, we have $y_m^{r}\ge 2^{r}>1$. Write
$R_m=L^{\sigma}$. Since $R_m\ge y_m^{r}>1=L^{0}$, we have $\sigma\ge1$; then $L^{\sigma-1}$
is a power of $L$ smaller than $R_m$, so by minimality $L^{\sigma-1}<y_m^{r}$, that is,
$R_m<Ly_m^{r}$. For the numbers $\check{R}_m$ the same inequality is the right-hand half of
\eqref{4.18:eq-older-window-2}. Applied at the level $m=s_0+1\le s\le k$ together with the
hypothesis on $\nu$, this gives
\begin{equation*}
L^{\nu-1}=R_{s_0+1}<Ly_{s_0+1}^{r},\qquad\text{that is,}\qquad \nu-2<r\log_L y_{s_0+1}.
\end{equation*}

\emph{Step 2 (the flow).} Let $m\le n\le k$. By the second line of \eqref{4.18:eq-older-window-1}
and $\bar\beta\ge0$, and since $g_n^{-2}\ge1$,
\begin{equation*}
g_m^{-2}=g_n^{-2}+\sum_{m<i\le n}\beta_i\le g_n^{-2}+\bar\beta(n-m)\le g_n^{-2}\bigl(1+\bar\beta(n-m)\bigr).
\end{equation*}
Taking logarithms, $y_m\le y_n+\log(1+\bar\beta(n-m))$. We apply this first with $m=s_0+1$, $n=s$,
so that $n-m=\nu-1$, and then with $m=s$, $n=k$, so that $n-m=u$. Since $\nu\ge1$ and the
logarithm is increasing, we obtain
\begin{equation*}
\begin{aligned}
y_{s_0+1}&\le y_s+\log\bigl(1+\bar\beta(\nu-1)\bigr)\\
&\le y_k+\log(1+\bar\beta u)+\log(1+\bar\beta\nu).
\end{aligned}
\end{equation*}

\emph{Step 3 (splitting the logarithm).} For $Y\ge1$ and $a,b\ge0$,
\begin{equation*}
Y(1+a)(1+b)=Y+Ya+Yb+Yab\ge Y+a+b,
\end{equation*}
hence $\log_L(Y+a+b)\le\log_LY+\log_L(1+a)+\log_L(1+b)$. We use this with $Y=y_k\ge2$,
$a=\log(1+\bar\beta u)\ge0$ and $b=\log(1+\bar\beta\nu)\ge0$.

\emph{Step 4 (elementary bounds).} For $t\ge1$ we have
\begin{equation*}
\log_Lt=\frac{\log t}{\log L}\le\frac{\log t}{\log2}\le\tfrac32\log t,
\end{equation*}
because the blocking factor $L$ is an odd integer greater than $1$
(Definition~\ref{1.1:def-lattice}), so that $L\ge3$, $\log L>\log2$, and
$1/\log2<\frac32$.
For $c\ge0$, the inequality $\log(1+x)\le x$ at $x=\log(1+c)\ge0$ gives
$\log(1+\log(1+c))\le\log(1+c)$. Moreover $\log(1+c)\le\sqrt c$: the function
$f(c)=\sqrt c-\log(1+c)$ vanishes at $c=0$, and for $c>0$
\begin{equation*}
f'(c)=\frac{1}{2\sqrt c}-\frac{1}{1+c}=\frac{(1-\sqrt c)^2}{2\sqrt c\,(1+c)}\ge0.
\end{equation*}
Finally $\sqrt c\le(1+c)/2$, by the inequality of arithmetic and geometric means. With $t=1+\log(1+\bar\beta u)$
and $c=\bar\beta u$ these give
\begin{equation*}
\log_L\bigl(1+\log(1+\bar\beta u)\bigr)\le\tfrac32\sqrt{\bar\beta u}\le\tfrac34(1+\bar\beta u).
\end{equation*}
With $t=1+\log(1+\bar\beta\nu)$ and $c=\bar\beta\nu$ they give
$\log_L(1+\log(1+\bar\beta\nu))\le\frac32\sqrt{\bar\beta\nu}$. The inequality of arithmetic and
geometric means, applied to the numbers $\frac94r^2\bar\beta$ and $\nu$,
reads $\frac32r\sqrt{\bar\beta\nu}\le\frac\nu2+\frac98r^2\bar\beta$; dividing by $r>0$,
\begin{equation*}
\log_L\bigl(1+\log(1+\bar\beta\nu)\bigr)\le\tfrac32\sqrt{\bar\beta\nu}\le\frac{\nu}{2r}+\tfrac98r\bar\beta.
\end{equation*}

\emph{Collecting.} By Step 1, the monotonicity of $\log_L$ and the bound of Step 2, then Step 3,
and then the two bounds of Step 4, each multiplied by $r$,
\begin{equation*}
\begin{aligned}
\nu-2&<r\log_L\bigl(y_k+\log(1+\bar\beta u)+\log(1+\bar\beta\nu)\bigr)\\
&\le r\log_Ly_k+r\log_L\bigl(1+\log(1+\bar\beta u)\bigr)+r\log_L\bigl(1+\log(1+\bar\beta\nu)\bigr)\\
&\le r\log_Ly_k+\tfrac34r(1+\bar\beta u)+\frac\nu2+\tfrac98r^2\bar\beta.
\end{aligned}
\end{equation*}
Subtracting $\nu/2$ and multiplying by $2$,
\begin{equation*}
\nu<4+\tfrac32r+\tfrac94r^2\bar\beta+2r\log_Ly_k+\tfrac32r\bar\beta u.
\end{equation*}

\emph{Step 5 (the current cell parameter).} By definition $R_k\ge y_k^{r}>1$ (for
$\check{R}_k$ this is the left-hand half of \eqref{4.18:eq-older-window-2}). Hence
$r\log_Ly_k=\log_Ly_k^{r}\le\log_LR_k$. Moreover $\log_LR_k\le\log_2R_k\le R_k$, since $L\ge3$
and $t\le 2^{t}$ for $t\ge1$ (Bernoulli's inequality $(1+1)^t\ge1+t$), applied with $t=R_k$.
Next, since $u\ge1$,
\begin{equation*}
4+\tfrac32r+\tfrac94r^2\bar\beta\le\bigl(4+2r+3r^2\bar\beta\bigr)u,
\end{equation*}
and $\frac32r\bar\beta u\le2r\bar\beta u$. Inserting these bounds into the last display of the collecting
step gives
\begin{equation*}
\nu<\bigl(4+2r+3r^2\bar\beta+2r\bar\beta\bigr)u+2R_k=c_1(k-s)+2R_k,
\end{equation*}
because $4+2r+3r^2\bar\beta+2r\bar\beta=4+2r(1+\bar\beta)+3r^2\bar\beta=c_1$. This is
\eqref{4.18:eq-older-window-a}, both for $R_m$ and for $\check{R}_m$.

The constant $c_1$ of \eqref{4.18:eq-older-window-a} depends on $r$ and $\bar\beta$ only.
Of the blocking factor only $L\ge3$ enters, and of the couplings only the ceiling $g_k\le e^{-1}$
and the flow in the form \eqref{4.18:eq-older-window-1}; no value of the couplings $g_m$ beyond
that ceiling is used.
\end{proof}

\begin{lemma}[The graded count of older members made of $M$-cubes]\label{4.18:lem-graded-count}
Let $Z$ be the component and $h$ the integer fixed in Definition~\ref{4.18:def-prepared-sequence},
and let $Z''_i$ be the members of the prepared sequence and $\Omega_i,\Lambda_i,Z_i,\Lambda$ the sets
of the label in force at step $k$, as in that definition. Let $u_i=L^{-(k-i)}$ be the level-$i$
spacing and $a_i=MR_iu_i$ the level-$i$ cell side, as in \eqref{4.18:eq-scales-partitions-a}.

Let $\mathcal{S}$ be the set of steps $s\le h$ at which an operation $\mathbf{R}$ was performed on
a component $Y^{(s)}\subset Z$ of the class considered in \cite{B-CMP122a} at step $s$ (the class
of Definition~\ref{4.18:def-prepared-sequence}, read at step $s$). For $s\in\mathcal{S}$ let
$N(s)>N_0(s)$ be the two integers of that operation \cite{B-CMP122a}, and put
\begin{equation*}
h(s):=s-N(s),\qquad s_0(s):=s-N_0(s).
\end{equation*}
For $i\le s$ let $Z''^{(s)}_i$ be the members of the prepared sequence of step $s$ inside
$Y^{(s)}$. Assume the following.
\begin{enumerate}
\item[(a)] The flow relations of Definition~\ref{4.18:def-coupling-excess} hold at all levels
$\le k$, and $g_k\le e^{-1}$.
\item[(b)] The cell parameters $R_m$, $m\le k$, satisfy the hypothesis of
Lemma~\ref{4.18:lem-older-window} on them, in either of the two forms stated there.
\item[(c)] For every $s\in\mathcal{S}$, $N_0(s)\ge 1$ is an integer with
$L^{N_0(s)-1}=R_{s-N_0(s)+1}$.
\item[(d)] The following clauses hold for all $s,s'\in\mathcal{S}$.
\begin{enumerate}
\item[(P1)] The members of the label in force at step $k$ which lie inside $Z$ at a level
$i\le h$ and are unions of level-$i$ $M$-cubes but not of level-$i$ cells are exactly the sets
\begin{equation*}
A^{(s)}_i:=Z''^{(s)}_i\cap Y^{(s)},\qquad s\in\mathcal{S},\quad s_0(s)<i<s .
\end{equation*}
They are nested,
\begin{equation*}
A^{(s)}_{s_0(s)+1}\subset\cdots\subset A^{(s)}_{s-1}\subset Z''^{(s)}_s\subset Y^{(s)},
\end{equation*}
and $Z''^{(s)}_{h(s)+1}\subset A^{(s)}_i$ for every such $i$, with
$Z''^{(s)}_{h(s)+1}\supset Z_{h(s)}\cap Y^{(s)}$.
\item[(P2)] $\operatorname{dist}_\infty\bigl(Z''^{(s)}_s,(Y^{(s)})^c\bigr)\ge 10a_s$.
\item[(P3)] If $s'<s$ both lie in $\mathcal{S}$ and $Y^{(s')}\cap Y^{(s)}\neq\emptyset$, then
$s'\le h(s)$ and $Y^{(s')}\subset Z_{h(s)}\cap Y^{(s)}$.
\item[(P4)] $Y^{(s)}\subset\Lambda_s^c$ for the label in force at step $k$, and
$\Omega_{s+1}\subset\Lambda_s$; moreover $Y^{(s)}\subset Z_s\cap Z\subset Z_h\cap Z$.
\item[(P5)] The level-$s$ member containing $Y^{(s)}$ is made of level-$s$ cells: it is
$Y^{(s)}$ itself, a closed union of level-$s$ cells.
\end{enumerate}
\item[(e)] $Z''_{h+1}\subset\Lambda\subset Z^{\sim -8}$, where $Z^{\sim -8}$ is $Z$ less eight
layers of cells.
\end{enumerate}

For $x\in Z$ the \emph{shell index} $j(x)\in\{0,\dots,k\}$ is the index with
$x\in\Omega_{j(x)}\setminus\Omega_{j(x)+1}$ for the label in force, where $\Omega_0$ is the whole
space and $\Omega_{k+1}:=\emptyset$. For a set $A$ made of level-$i$ $M$-cubes, $[A]_i$ is the
union of the level-$i$ cells holding a cube of $A$, and $S^{+r}$ is the open
$r$-neighbourhood of a set $S$. Put
\begin{equation*}
n_{\mathrm{old}}(x):=\#\Bigl\{(s,i):\ s\in\mathcal{S},\ s_0(s)<i<s,\
x\in\bigl([A^{(s)}_i]_i\bigr)^{+\frac14 a_i}\setminus A^{(s)}_i\Bigr\}.
\end{equation*}
By \eqref{4.18:eq-profiles-admissible-b} and \eqref{4.18:eq-coupling-excess-b}, the old part
$\mathfrak{d}_{\mathrm{old}}$ of the profile excess of
Definition~\ref{4.18:def-coupling-excess} satisfies
\begin{equation}\label{4.18:eq-graded-count-a}
0\le\mathfrak{d}_{\mathrm{old}}(x)\le b_0\, n_{\mathrm{old}}(x).
\end{equation}
Then for every $x\in Z$, with the constants $c_1,c_2$ of \eqref{4.18:eq-older-window-a},
\begin{equation}\label{4.18:eq-graded-count-b}
\begin{gathered}
n_{\mathrm{old}}(x)\le c_1\bigl(k-j(x)\bigr)+c_2R_k,\\
\{n_{\mathrm{old}}>0\}\subset\bigcup_{s\in\mathcal{S}}Y^{(s)}\subset Z_h\cap Z\subset Z''_{h+1},\\
\operatorname{dist}_\infty\bigl(\{n_{\mathrm{old}}>0\},\,Z\setminus Z''_{h+1}\bigr)
\ge\operatorname{dist}_\infty\bigl(Z_h\cap Z,\,Z\setminus Z''_{h+1}\bigr)\ge 5a_{h+1},\\
\{n_{\mathrm{old}}>0\}\subset\Lambda\subset Z^{\sim -8}.
\end{gathered}
\end{equation}
\end{lemma}

\begin{proof}
Clauses (P1)--(P5) of hypothesis (d) and hypothesis (e) are hypotheses of the lemma, and the
argument below uses them only as stated. For orientation we record the statements of
Ba{\l}aban's construction to which they correspond. Clause (P1) corresponds to the description
of the prepared sequence in \cite{B-CMP122a}, applied at step $s$. Its last inclusion
corresponds to \cite[(1.11)]{B-CMP122a} and \cite[(2.1)]{B-CMP119} applied at step $s$, that is,
with $Z$, $h$ and the label in force replaced by the component $Y^{(s)}$, the integer $h(s)$ and
the label of step $s$, whose sets we write $\Omega^{(s)}_i$, $\Lambda^{(s)}_i$; in this form they
read
\begin{equation*}
Z''^{(s)}_{h(s)+1}=\bigl((\Omega^{(s)}_{h(s)+1})^{\sim 5}\bigr)^c\cap Y^{(s)}
\supset(\Omega^{(s)}_{h(s)+1})^c\cap Y^{(s)}\supset(\Lambda^{(s)}_{h(s)})^c\cap Y^{(s)},
\end{equation*}
and $s_0(s)>h(s)$ is the inequality $N(s)>N_0(s)$. Clause (P2) corresponds to
$Z''^{(s)}_s\subset Z_{s-1}\cap Y^{(s)}$ together with the ten layers of cells of the ten-layer
geometry of the large-field regions, applied at step $s$ inside the component $Y^{(s)}$.
Clause (P3) corresponds to the bound on the steps of older operations and to clause (ii) of the
closure property, both applied at step $s$: $Y^{(s')}$ is connected and lies in the large-field
region of step $s$, whose components are the classes of the relation ``intersect, or touch'' of
\cite{B-CMP122b}, so that, meeting the component $Y^{(s)}$, it lies in it. Clause (P4)
corresponds to \cite[p.~390]{B-CMP122b} applied at step $s$: $Y^{(s)}$ is a level-$s$
large-field region, and the later operations $\mathbf{T}$ append levels and keep
\cite[(2.1)]{B-CMP119}. Hypothesis (e) corresponds to the nesting of the prepared sequence in
the small-field region, together with the eight-layer margin of $\Lambda$ in $Z$.

Say that a pair $(s,i)$ with $s\in\mathcal{S}$ and $s_0(s)<i<s$ \emph{contributes} at $x$ if
$x\in([A^{(s)}_i]_i)^{+\frac14 a_i}\setminus A^{(s)}_i$; thus $n_{\mathrm{old}}(x)$ is the number
of pairs contributing at $x$.

\emph{Localisation of one contribution.} Let $(s,i)$ contribute at $x$. Then there is a cell
$D$ of $[A^{(s)}_i]_i$ with $\operatorname{dist}_\infty(x,D)<\frac14 a_i$. By definition of
$[A^{(s)}_i]_i$, the cell $D$ holds a cube of $A^{(s)}_i$, and $A^{(s)}_i\subset Z''^{(s)}_s$ by
the nesting in (P1). By \eqref{4.18:eq-scales-partitions-a} the cell $D$ has side $a_i$, so every
point of $D$ lies at sup-distance at most $a_i$ from that cube; with $a_i\le a_s$ this gives
\begin{equation*}
\operatorname{dist}_\infty\bigl(x,Z''^{(s)}_s\bigr)<a_i+\tfrac14 a_i\le\tfrac54 a_s<10a_s .
\end{equation*}
By (P2), every point at sup-distance less than $10a_s$ from $Z''^{(s)}_s$ lies in $Y^{(s)}$;
hence $x\in Y^{(s)}$. Moreover $x\notin A^{(s)}_i$ by the definition of a contribution, and
$A^{(s)}_i\supset Z''^{(s)}_{h(s)+1}$ by (P1), since $s_0(s)<i<s$; hence
$x\notin Z''^{(s)}_{h(s)+1}$.

\emph{At most one step contributes.} Let $(s,i)$ and $(s',i')$ contribute at $x$, with $s'<s$.
The localisation step applied to $(s',i')$ and to $(s,i)$ gives
$x\in Y^{(s')}\cap Y^{(s)}$, so this intersection is not empty. By (P3),
$Y^{(s')}\subset Z_{h(s)}\cap Y^{(s)}$. By (P1), $Z_{h(s)}\cap Y^{(s)}\subset Z''^{(s)}_{h(s)+1}$
and $Z''^{(s)}_{h(s)+1}\subset A^{(s)}_i$. Therefore
\begin{equation*}
x\in Y^{(s')}\subset Z_{h(s)}\cap Y^{(s)}\subset Z''^{(s)}_{h(s)+1}\subset A^{(s)}_i ,
\end{equation*}
which contradicts $x\notin A^{(s)}_i$, obtained in the localisation step for $(s,i)$. Hence all
pairs contributing at $x$ share one step $s=s(x)$, and their second entries $i$ range over the
integers with $s-N_0(s)<i<s$, of which there are $N_0(s)-1$. Thus, if some pair contributes at
$x$,
\begin{equation*}
n_{\mathrm{old}}(x)\le N_0(s)-1 .
\end{equation*}
If no pair contributes at $x$, then $n_{\mathrm{old}}(x)=0$ and the first line of
\eqref{4.18:eq-graded-count-b} holds, its right side being nonnegative because $j(x)\le k$,
$c_1,c_2>0$ and $R_k>0$. In what follows $x$ is a point at which some pair contributes, and
$s=s(x)$.

\emph{The shell index.} By the localisation step, $x\in Y^{(s)}$, and $Y^{(s)}\subset\Lambda_s^c$
by (P4); so $x\notin\Lambda_s$. Since $\Omega_{s+1}\subset\Lambda_s$, again by (P4),
$x\notin\Omega_{s+1}$. By the definition of the shell index, $x\in\Omega_{j(x)}$; since the sets
$\Omega_j$ of the label decrease in $j$, $j(x)\ge s+1$ would give $x\in\Omega_{s+1}$. Hence
$j(x)\le s$, that is, $k-s\le k-j(x)$.

\emph{The count.} The step $s$ lies in $\mathcal{S}$, so $s\le h$. Hypothesis (a) gives the flow
relations at all levels $\le k$ and the ceiling $g_k\le e^{-1}$, hypothesis (b) gives the
hypothesis of Lemma~\ref{4.18:lem-older-window} on the cell parameters, and hypothesis (c) gives
the integer $N_0(s)\ge1$ with $L^{N_0(s)-1}=R_{s-N_0(s)+1}$. Lemma~\ref{4.18:lem-older-window}
therefore applies at this $s$, and \eqref{4.18:eq-older-window-a} gives
$N_0(s)<c_1(k-s)+c_2R_k$. Since $c_1>0$ and $k-s\le k-j(x)$,
\begin{equation*}
n_{\mathrm{old}}(x)\le N_0(s)-1<c_1(k-s)+c_2R_k\le c_1\bigl(k-j(x)\bigr)+c_2R_k ,
\end{equation*}
which is the first line of \eqref{4.18:eq-graded-count-b}.

\emph{The support.} By the localisation step, a point $x$ at which some pair $(s,i)$ contributes
lies in $Y^{(s)}$; so $\{n_{\mathrm{old}}>0\}\subset\bigcup_{s\in\mathcal{S}}Y^{(s)}$. By (P4),
$Y^{(s)}\subset Z_h\cap Z$ for every $s\in\mathcal{S}$. By \cite[(1.11)]{B-CMP122a},
$Z''_{h+1}=(\Omega^{\sim 5}_{h+1})^c\cap Z$; by \cite[(1.11)]{B-CMP122a} and
\cite[(2.1)]{B-CMP119}, this set contains $\Omega^c_{h+1}\cap Z$ together with the points of $Z$
in the five layers of level-$(h+1)$ cells surrounding $\Omega^c_{h+1}\cap Z$, and
$\Omega^c_{h+1}\cap Z\supset Z_h\cap Z$. Each such layer has width $a_{h+1}$, so every point of
$Z$ at sup-distance less than $5a_{h+1}$ from $\Omega^c_{h+1}\cap Z$ lies in $\Omega^c_{h+1}\cap Z$
or in one of these five layers; that is,
\begin{equation*}
Z''_{h+1}\supset\bigl\{y\in Z:\ \operatorname{dist}_\infty\bigl(y,\Omega^c_{h+1}\cap Z\bigr)
<5a_{h+1}\bigr\},\qquad \Omega^c_{h+1}\cap Z\supset Z_h\cap Z .
\end{equation*}
Hence
\begin{equation*}
Z_h\cap Z\subset Z''_{h+1},\qquad
\operatorname{dist}_\infty\bigl(Z_h\cap Z,\ Z\setminus Z''_{h+1}\bigr)\ge 5a_{h+1}.
\end{equation*}
This proves the second line of \eqref{4.18:eq-graded-count-b}, and the third line follows
because $\{n_{\mathrm{old}}>0\}$ is a subset of $Z_h\cap Z$. Finally, by hypothesis (e),
$Z''_{h+1}\subset\Lambda\subset Z^{\sim -8}$, and hence
$\{n_{\mathrm{old}}>0\}\subset\Lambda\subset Z^{\sim -8}$, which is the fourth line of
\eqref{4.18:eq-graded-count-b}.
\end{proof}

\begin{remark}[Why the count must be graded by level]\label{4.18:rem-graded-form}
Let $n_{\mathrm{old}}(x)$ denote the count of contributing older members at $x$ and $j(x)$ the shell index of $x$. The graded count of older members made of $M$-cubes (Lemma~\ref{4.18:lem-graded-count}) bounds $n_{\mathrm{old}}(x)$ in a form graded by the difference $k-j(x)$. A bound depending on the running coupling at level $k$ alone is not available.

To see this, consider the supremum over $Z$ of $n_{\mathrm{old}}$. It is of the order of $\log_L R_s$, where $s$ is the oldest step that contributes and $R_s$ is the cell parameter at level $s$. The parameter $R_s$ is controlled by the coupling $g_s$ of level $s$, not by $g_k$. For fixed $g_k$ the difference $k-s$ is unbounded, so no bound of $n_{\mathrm{old}}$ by a quantity depending on $g_k$ alone is available.

The graded form is of the same kind as Ba{\l}aban's own graded bound in \cite{B-CMP122b}. That bound is a bound by $O(1)(k-j)$ on the $j$-th shell $\Omega''_j\setminus\Omega''_{j+1}$ of the sequence of domains $\Omega''_j$ of the $\mathbf{R}$-operation treated there. Ba{\l}aban's construction suppresses that growth through the scaled decay of the terms $H_{k,Z}$ that he attaches to the level $k$ and the large-field region $Z$.

Accordingly, every subsequent use in this chapter of the bound on $n_{\mathrm{old}}$ is through the graded form, and never through a supremum over $Z$.
\end{remark}

\begin{proposition}[The graded excess bound in the large-field estimates]
\label{4.18:prop-graded-excess}
Fix the level $k$ and work in the setting of Definition~\ref{4.18:def-coupling-excess}, of
Lemma~\ref{4.18:lem-older-window} and of Lemma~\ref{4.18:lem-graded-count}. Thus:
\begin{itemize}
\item $Z$ is the large-field region in which the current $\mathbf{R}$-operation is performed.
\item $N_0$, $N$ and $h$ are Ba{\l}aban's integers of that operation, with $N>N_0$ and $h<k$
\cite[\S1]{B-CMP122a}, and with $N\le R_k$ \cite[\S1]{B-CMP122b}.
\item $Z''_i$ are the prepared members, and $\Omega_i$, $\Omega''_i$ are the domains of
\cite[\S1]{B-CMP122a}.
\item $R_i$ and $\check{R}_i$ are the cell parameters of Lemma~\ref{4.18:lem-older-window},
$c_1,c_2$ are the constants of \eqref{4.18:eq-older-window-a}, and $r_0>0$ is the exponent in
the cell parameter there.
\item The ceiling $g_k\le e^{-1}$ of Lemma~\ref{4.18:lem-older-window} is in force.
\item $u_i=L^{-(k-i)}$ and $a_i=MR_iu_i$.
\item $b_0$ is the constant in the bound $0\le\beta_i\le b_0$ ($i\le k$) of
Definition~\ref{4.18:def-coupling-excess}.
\item $\mathfrak{d}_{\mathrm{old}}$ is the old part of the excess.
\end{itemize}
For $j\le h$ put
\begin{equation}\label{4.18:eq-graded-excess-1}
\mathrm{zone}_j:=(\Omega_j\setminus\Omega_{j+1})\cap Z ,
\end{equation}
which equals $(\Omega''_j\setminus\Omega''_{j+1})\cap Z$ \cite[\S1]{B-CMP122a}. For
$x\in\mathrm{zone}_j$ write $j(x):=j$.

The following objects are taken from \cite{B-CMP122b}:
\begin{itemize}
\item the small functions $\mathfrak{m},\mathfrak{m}'\ge0$;
\item the cutoff functions $\zeta_0,\zeta_1$, and $\zeta=\zeta_0(1-\zeta_1)$;
\item the coefficient $G$ of Ba{\l}aban's class of coupling functions on $Z$, and the sets
$\Gamma''_j\subset\mathrm{zone}_j$ ($j\le h$) of \cite[(1.47)]{B-CMP122b};
$c_{\Gamma}$ denotes the constant implicit in the term $O(k-j)$ of the factor
$(g_k^{-2}+O(k-j))$ in \cite[(1.47)]{B-CMP122b}, so that this factor is written
$(g_k^{-2}+c_{\Gamma}(k-j))$ below, whereas $O(1)$ below denotes a factor that does not depend on
$k$, $j$ and $N$;
\item the coupling function $x''=(g''_k)^{-2}$ of the family's system after the preparation, and
the decay constant $\delta$;
\item the free integer $\nu$ of the regime $L^{-N}\le(\log g_k^{-2})^{-\nu}$.
\end{itemize}

The remaining objects are as follows.
\begin{itemize}
\item $\mathfrak{d}=\mathfrak{d}^{(n)}:=e_n$ is the profile excess at the current stage $n$ of
Definition~\ref{4.18:def-coupling-excess}.
\item $T_{\mathfrak{d}}$ is the remainder of the $Z$-part of the large-field estimates
attached to $\mathfrak{d}$, and $T_{\mathfrak{d}_{\mathrm{old}}}$ is the remainder attached to
$\mathfrak{d}_{\mathrm{old}}$.
\item (C1) denotes the ceiling on $g_k$ under which the remainder $T_{\mathfrak{d}}$ is absorbed
in the large-field estimates; it strengthens ceilings on $g_k$ stated in
\cite[\S1]{B-CMP122b}.
\item $x^{\circ}$ is the coupling function of the companion system
(Definition~\ref{4.18:def-coupling-excess}).
\item $\mathcal{E}'$ is the support set of the current part of the excess.
\item $w_{n+1}$ and $w^{\circ}_{n+1}$ are the stage increments of the coupling functions of the
family's system and of the companion system; the difference $w_{n+1}-w^{\circ}_{n+1}$ is the
quantity bounded in the companion estimate, and hypothesis (V) below bounds it by a part
supported in $\mathcal{E}'$ and the old part $\mathfrak{d}_{\mathrm{old}}$.
\item $\mathfrak{t}_p\ge0$ are the per-plaquette coefficients of hypothesis (III) below.
\item $\mathfrak{M}_j$ ($j\le h$), $\mathfrak{M}^{\mathrm{gr}}$ and $\mathfrak{M}^{\zeta_0}$ are
the zone-wise, the graded and the $\zeta_0$-restricted majorants of the coefficients
$\mathfrak{t}_p$ in the large-field estimates; the properties of them that are used are stated
in hypotheses (III) and (IV) below.
\item $\mathrm{Rem}_{\mathrm{abs}}$ denotes the absorbed remainder of the companion estimate.
\end{itemize}

Assume the following.
\begin{enumerate}
\item[(I)] The hypotheses of Lemma~\ref{4.18:lem-graded-count} hold, so that its conclusions
\eqref{4.18:eq-graded-count-a} and \eqref{4.18:eq-graded-count-b} hold on $Z$.
\item[(II)] $Z''_k\subset\{\zeta_0=1\}$, and $\mathrm{zone}_j\subset Z''_k$ for $j\le h$.
Moreover, $\mathfrak{d}_{\mathrm{old}}$ vanishes outside $\bigcup_{j\le h}\mathrm{zone}_j$, and its
support lies in $Z''_{h+1}$.
\item[(III)] The remainder of the old part obeys the termwise bound
$|T_{\mathfrak{d}_{\mathrm{old}}}|\le\sum_p|\mathfrak{d}_{\mathrm{old}}(x(p))|\,\mathfrak{t}_p$,
and, for every $j\le h$, $\sum_{p:\,x(p)\in\mathrm{zone}_j}\mathfrak{t}_p\le\mathfrak{M}_j$.
\item[(IV)] The majorant bounds for the small functions $\mathfrak{m},\mathfrak{m}'$ of
\cite[\S1]{B-CMP122b} hold, namely
$\mathfrak{M}^{\zeta_0}\le g_k^2\mathfrak{m}$ and $\mathfrak{M}^{\mathrm{gr}}\le\mathfrak{m}'$.
The majorants satisfy
\begin{equation*}
\sum_{j\le h}(k-j+1)\mathfrak{M}_j\le\mathfrak{M}^{\mathrm{gr}},
\end{equation*}
and $\sum_{j\le h}\mathfrak{M}_j\le\mathfrak{M}^{\zeta_0}$ whenever
$\mathrm{zone}_j\subset\{\zeta_0=1\}$ for all $j\le h$.
Under the graded majorant bound, moreover, $\mathrm{Rem}_{\mathrm{abs}}$ of a weight bounded by
$b_0c_1(k-j)$ on $\mathrm{zone}_j$ ($j\le h$) is at most $b_0c_1\mathfrak{m}'$.
\item[(V)] The remainder bound of the large-field estimates has the form
$|T_{\mathfrak{d}}|\le b_0N_0g_k^2\mathfrak{m}+|T_{\mathfrak{d}_{\mathrm{old}}}|$. Pointwise,
\begin{equation*}
|\mathfrak{d}^{(n)}|\le b_0N_0+\mathfrak{d}_{\mathrm{old}},\qquad
|w_{n+1}-w^{\circ}_{n+1}|\le b_0N_0\mathbf{1}_{\mathcal{E}'}+\mathfrak{d}_{\mathrm{old}},
\end{equation*}
and $\mathcal{E}'\subset\{\zeta_0=1\}$.
\item[(VI)] The companion estimate absorbs a constant weight part $b_0N_0$ on
$\{\zeta_0=1\}$ under a ceiling on $g_k$, denoted (C4), which is defined through a quantity
$t_4$ carrying the factor $b_0L$. It equally absorbs the constant part $b_0(1+c_2)R_k$ on
$\{\zeta_0=1\}$ under the ceiling (C4$'$), which is (C4) with $b_0L$ replaced by $b_0(1+c_2)L$
in the definition of $t_4$.
\item[(VII)] The function $\zeta$ vanishes on $Z''_{h+1}$ except on ``a layer of the width $2$ at
the boundary, in $L^{-h}$-scale'' \cite[\S1]{B-CMP122b}; that is, except at points of $Z''_{h+1}$
at sup-distance less than $2u_h$ from $Z\setminus Z''_{h+1}$.
\item[(VIII)] On $\Gamma''_j$ ($j\le h$) one has $(g''_k)^{-2}=G+\mathfrak{d}_{\mathrm{old}}$, with
$G\le g_k^{-2}+c_{\Gamma}(k-j)$, in the setting of \cite[(1.46)--(1.48)]{B-CMP122b}.
\item[(IX)] Items (IX.1)--(IX.3) give Ba{\l}aban's absorption of the factor
$(g_k^{-2}+c_{\Gamma}(k-j))$ of \cite[(1.47)]{B-CMP122b}, in the form in which it is used here;
items (IX.4)--(IX.6) concern the regime $L^{-N}\le(\log g_k^{-2})^{-\nu}$, which
\cite[\S1]{B-CMP122b} states for sufficiently large $\nu$. All six items are assumed in the
form stated.
\begin{enumerate}
\item[(IX.1)] For $j\le h$,
\begin{equation}\label{4.18:eq-graded-excess-3}
g_k^{-2}+c_{\Gamma}(k-j)\le g_k^{-2}O(1)(1+k-j)\le g_k^{-2}O(1)(N+2)(1+h-j),
\end{equation}
with factors $O(1)$ that do not depend on $k$, $j$ and $N$.
\item[(IX.2)] The polynomial in $N$ in \eqref{4.18:eq-graded-excess-3} goes into $R_k^2$, using
$N\le R_k$.
\item[(IX.3)] The polynomial in $h-j$ in \eqref{4.18:eq-graded-excess-3} goes into
$e^{-\frac12\delta M(h-j)}$.
\item[(IX.4)] The right member of \cite[(1.47)]{B-CMP122b} so obtained carries $R_k^7$;
\cite[\S1]{B-CMP122b} describes it as ``small under the usual conditions on $N$''. It is small
in the regime above as soon as $\nu\ge p_0+\max\{p_0,q\}+7r_0+1$.
\item[(IX.5)] A right member carrying one further factor $R_k$, that is
$R_k^8$, is small in the same regime as soon as $\nu\ge p_0+\max\{p_0,q\}+8r_0+1$.
\item[(IX.6)] These floors on $\nu$ do not depend on $L$.
\end{enumerate}
\item[(X)] $(g''_k)^{-2}\zeta_1\ge g_k^{-2}\zeta_1\ge0$, and
$x''\zeta_1=(x^{\circ}+\mathfrak{d}_{\mathrm{old}})\zeta_1$; the latter is the identity
$x''\zeta_1=x^{\circ}\zeta_1$ of the companion estimate in the form it takes when
$\mathfrak{d}_{\mathrm{old}}$ does not vanish.
\end{enumerate}

Then the following hold.
\begin{enumerate}
\item[(1)] On $\mathrm{zone}_j$, for $j\le h$,
\begin{equation}\label{4.18:eq-graded-excess-a}
0\le\mathfrak{d}_{\mathrm{old}}\le b_0c_1(k-j)+b_0c_2R_k\le b_0c_1(k-j+1)+b_0c_2R_k .
\end{equation}
\item[(2)] The old part of the remainder obeys
$|T_{\mathfrak{d}_{\mathrm{old}}}|\le b_0(c_1\mathfrak{m}'+c_2R_kg_k^2\mathfrak{m})$, and the full
remainder obeys
\begin{equation}\label{4.18:eq-graded-excess-b}
|T_{\mathfrak{d}}|\le b_0\bigl(N_0g_k^2\mathfrak{m}+c_1\mathfrak{m}'+c_2R_kg_k^2\mathfrak{m}\bigr)
\le b_0\bigl((1+c_2)R_kg_k^2\mathfrak{m}+(1+c_1)\mathfrak{m}'\bigr).
\end{equation}
The ceiling (C1) on $g_k$ under which the remainder $T_{\mathfrak{d}}$ is absorbed in the
large-field estimates is replaced by the ceiling \eqref{4.18:eq-graded-excess-4}, referred to as
(C1$'$):
\begin{equation}\label{4.18:eq-graded-excess-4}
b_0\bigl((1+c_2)\check{R}_kg_k^2\mathfrak{m}+(1+c_1)\mathfrak{m}'\bigr)\le1 ;
\end{equation}
here $\check{R}_k$ is the cell parameter of the alternative in
Lemma~\ref{4.18:lem-older-window}; when Ba{\l}aban's cell parameter is used, (C1$'$) is read
with $R_k$ in place of $\check{R}_k$, and under (C1$'$) so read the right member of
\eqref{4.18:eq-graded-excess-b} is at most $1$.
\item[(3)] For $x\in Z$,
\begin{equation}\label{4.18:eq-graded-excess-5}
\begin{split}
|w_{n+1}(x)-w^{\circ}_{n+1}(x)|
&\le b_0c_1\,(k-j(x))\,\mathbf{1}_{\bigcup_{j\le h}\mathrm{zone}_j}(x)\\
&\quad+b_0(N_0+c_2R_k)\,\mathbf{1}_{\{\zeta_0=1\}}(x),
\end{split}
\end{equation}
the first term being read as zero outside $\bigcup_{j\le h}\mathrm{zone}_j$.
The constant part is absorbed under (C4$'$), and $\mathrm{Rem}_{\mathrm{abs}}$ of the graded part
is at most $b_0c_1\mathfrak{m}'$.
\item[(4)] $\mathfrak{d}_{\mathrm{old}}\,\zeta\equiv0$; hence $\mathfrak{d}_{\mathrm{old}}$
contributes nothing to the term of the large-field estimates that carries the cutoff $\zeta$,
and the ceilings on $g_k$ under which that term is estimated are unchanged.
\item[(5)] On $\Gamma''_j$ ($j\le h$),
\begin{equation*}
(g''_k)^{-2}\le g_k^{-2}+c'_{\Gamma}(k-j)+b_0c_2R_k,\qquad c'_{\Gamma}:=c_{\Gamma}+2b_0c_1 .
\end{equation*}
The bound that takes the place of \cite[(1.47)]{B-CMP122b} in the present setting carries
$R_k^8$ where \cite[(1.47)]{B-CMP122b} has $R_k^7$, and it is small when
$\nu\ge p_0+\max\{p_0,q\}+8r_0+1$. The positivity used in
\cite[(1.48)]{B-CMP122b} is unchanged. On $\Gamma''_j$ the same coefficient bound and the same
absorption hold for $x''\zeta_1=(x^{\circ}+\mathfrak{d}_{\mathrm{old}})\zeta_1$.
\item[(6)] On $\mathrm{zone}_{j'}$ ($j'\le h$),
$|\mathfrak{d}^{(n)}|\le b_0\bigl(N_0+c_1(k-j'+1)+c_2R_k\bigr)$.
\end{enumerate}
\end{proposition}

\begin{proof}
\emph{The pointwise bound \eqref{4.18:eq-graded-excess-a}.} Let $j\le h$ and
$x\in\mathrm{zone}_j$. By hypothesis (I), the bounds \eqref{4.18:eq-graded-count-a} and
\eqref{4.18:eq-graded-count-b} of Lemma~\ref{4.18:lem-graded-count} hold on $\mathrm{zone}_j$.
They give $0\le\mathfrak{d}_{\mathrm{old}}(x)\le b_0c_1(k-j)+b_0c_2R_k$. The second inequality
of \eqref{4.18:eq-graded-excess-a} holds because $b_0c_1\ge0$: here $b_0\ge0$, and $c_1$ is the
sum of non-negative terms displayed in Lemma~\ref{4.18:lem-older-window}. This proves (1).

\emph{The old part of the remainder.} By the termwise bound (III),
$|T_{\mathfrak{d}_{\mathrm{old}}}|\le\sum_p|\mathfrak{d}_{\mathrm{old}}(x(p))|\,\mathfrak{t}_p$.
By (II), $\mathfrak{d}_{\mathrm{old}}$ vanishes outside
$\bigcup_{j\le h}\mathrm{zone}_j$. Hence only plaquettes with $x(p)$ in some $\mathrm{zone}_j$
with $j\le h$ contribute; since all terms are non-negative, their sum is at most the sum over
$j\le h$ of the sums over the plaquettes with $x(p)\in\mathrm{zone}_j$, and on each such zone
(1) applies. Split the bound \eqref{4.18:eq-graded-excess-a} into
its graded part $b_0c_1(k-j+1)$ and its constant part $b_0c_2R_k$, and insert it zone by zone.
Using $\sum_{p:\,x(p)\in\mathrm{zone}_j}\mathfrak{t}_p\le\mathfrak{M}_j$ from (III), this gives
\begin{align*}
|T_{\mathfrak{d}_{\mathrm{old}}}|
&\le b_0c_1\sum_{j\le h}(k-j+1)\mathfrak{M}_j+b_0c_2R_k\sum_{j\le h}\mathfrak{M}_j\\
&\le b_0c_1\mathfrak{M}^{\mathrm{gr}}+b_0c_2R_k\mathfrak{M}^{\zeta_0}.
\end{align*}
The first sum is at most $\mathfrak{M}^{\mathrm{gr}}$ by (IV). For the second sum, the zones
with $j\le h$ lie, by (II), in $Z''_k\subset\{\zeta_0=1\}$, so that (IV) gives
$\sum_{j\le h}\mathfrak{M}_j\le\mathfrak{M}^{\zeta_0}$.

The majorant bounds (IV) give $\mathfrak{M}^{\mathrm{gr}}\le\mathfrak{m}'$ and
$\mathfrak{M}^{\zeta_0}\le g_k^2\mathfrak{m}$. Hence
\begin{equation*}
|T_{\mathfrak{d}_{\mathrm{old}}}|\le b_0\bigl(c_1\mathfrak{m}'+c_2R_kg_k^2\mathfrak{m}\bigr).
\end{equation*}

\emph{The full remainder.} Inserting the last bound into the remainder bound (V),
$|T_{\mathfrak{d}}|\le b_0N_0g_k^2\mathfrak{m}+|T_{\mathfrak{d}_{\mathrm{old}}}|$, gives the first
inequality of \eqref{4.18:eq-graded-excess-b}. For the second inequality, $N_0\le R_k$ gives
\begin{equation*}
b_0N_0g_k^2\mathfrak{m}+b_0c_2R_kg_k^2\mathfrak{m}\le b_0(1+c_2)R_kg_k^2\mathfrak{m},
\end{equation*}
and $b_0c_1\mathfrak{m}'\le b_0(1+c_1)\mathfrak{m}'$ because $b_0\mathfrak{m}'\ge0$. This proves
(2).

The ceiling \eqref{4.18:eq-graded-excess-4} replaces (C1); it is a ceiling on $g_k$ of the same
kind as the ceilings of \cite[\S1]{B-CMP122b} that it strengthens. No power of $R_k$ multiplies
$\mathfrak{m}'$ in it. When it is read with $R_k$ in place of $\check{R}_k$, its left member is
the right member of \eqref{4.18:eq-graded-excess-b}, which is therefore at most $1$.

\emph{The companion system.} By (V),
$|w_{n+1}-w^{\circ}_{n+1}|\le b_0N_0\mathbf{1}_{\mathcal{E}'}+\mathfrak{d}_{\mathrm{old}}$. By (II),
$\mathfrak{d}_{\mathrm{old}}$ vanishes outside $\bigcup_{j\le h}\mathrm{zone}_j$, and on
$\mathrm{zone}_j$ ($j\le h$) the first inequality of \eqref{4.18:eq-graded-excess-a} bounds
$\mathfrak{d}_{\mathrm{old}}$ by $b_0c_1(k-j(x))+b_0c_2R_k$. These zones lie in
$\{\zeta_0=1\}$ by (II), and so does $\mathcal{E}'$ by (V). Collecting the constant terms on
$\{\zeta_0=1\}$ gives \eqref{4.18:eq-graded-excess-5}.

Since $N_0\le R_k$, the constant part satisfies
\begin{equation*}
b_0(N_0+c_2R_k)\le b_0(1+c_2)R_k .
\end{equation*}
By hypothesis (VI), a constant part of this size on $\{\zeta_0=1\}$ is absorbed under
(C4$'$); this is the absorption of the companion estimate with $b_0N_0$ replaced by
$b_0(N_0+c_2R_k)$. The graded part is a weight bounded by $b_0c_1(k-j)$ on $\mathrm{zone}_j$, so
by the graded majorant bound in (IV) its absorbed remainder satisfies
$\mathrm{Rem}_{\mathrm{abs}}\le b_0c_1\mathfrak{m}'$. Thus the conclusion of the companion
estimate, that the absorbed remainder is small, keeps its form; the graded majorant bound is
added to its inputs. This proves (3).

\emph{The term carrying $\zeta$.} By \eqref{4.18:eq-graded-count-b}, the support of
$\mathfrak{d}_{\mathrm{old}}$ keeps sup-distance at least $5a_{h+1}$ from $Z\setminus Z''_{h+1}$.
Since $u_{h+1}=Lu_h$ and $M,R_{h+1}\ge1$, we have
\begin{equation*}
5a_{h+1}=5MR_{h+1}Lu_h>2u_h .
\end{equation*}
By (VII), $\zeta=\zeta_0(1-\zeta_1)$ vanishes at every point of $Z''_{h+1}$ whose sup-distance from
$Z\setminus Z''_{h+1}$ is at least $2u_h$. The support of $\mathfrak{d}_{\mathrm{old}}$ lies in
$Z''_{h+1}$ by (II), so $\zeta$ vanishes in particular on the support of
$\mathfrak{d}_{\mathrm{old}}$. Outside that support $\mathfrak{d}_{\mathrm{old}}=0$. Hence
$\mathfrak{d}_{\mathrm{old}}\zeta\equiv0$. The share of $\mathfrak{d}_{\mathrm{old}}$ in the term
of the large-field estimates that carries $\zeta$ enters through the product
$\mathfrak{d}_{\mathrm{old}}\zeta$ and therefore vanishes, so the ceilings under which that term is
estimated are unchanged. This proves (4).

\emph{The coefficient of \cite[(1.47)]{B-CMP122b}.} Let $j\le h$, so that $k-j\ge1$ as
$j\le h<k$. On $\Gamma''_j\subset\mathrm{zone}_j$, hypothesis (VIII) and
\eqref{4.18:eq-graded-excess-a} give
\begin{align*}
(g''_k)^{-2}=G+\mathfrak{d}_{\mathrm{old}}
&\le\bigl[g_k^{-2}+c_{\Gamma}(k-j)\bigr]+b_0c_1(k-j+1)+b_0c_2R_k\\
&\le g_k^{-2}+c'_{\Gamma}(k-j)+b_0c_2R_k .
\end{align*}
The second line uses $k-j+1\le2(k-j)$. Thus the factor $(g_k^{-2}+c_{\Gamma}(k-j))$ of
\cite[(1.47)]{B-CMP122b} becomes $(g_k^{-2}+c'_{\Gamma}(k-j)+b_0c_2R_k)$, with
$c'_{\Gamma}=c_{\Gamma}+2b_0c_1$.

The absorption \eqref{4.18:eq-graded-excess-3} of hypothesis (IX.1) bounds the old factor by
$g_k^{-2}O(1)(N+2)(1+h-j)$. Under this absorption the polynomial in $N$ goes into $R_k^2$ and the
polynomial in $h-j$ goes into $e^{-\frac12\delta M(h-j)}$. The same line holds with
$c'_{\Gamma}$ in place of $c_{\Gamma}$, the factor $O(1)$ now depending on $c'_{\Gamma}$: since
$g_k\le e^{-1}\le1$ by the ceiling of Lemma~\ref{4.18:lem-older-window}, in whose setting we work,
\begin{equation*}
g_k^{-2}+c'_{\Gamma}(k-j)\le g_k^{-2}\bigl(1+c'_{\Gamma}(k-j)\bigr)
\le g_k^{-2}\max\{1,c'_{\Gamma}\}(1+k-j),
\end{equation*}
and the second inequality of \eqref{4.18:eq-graded-excess-3}, which passes from the factor
$1+k-j$ to $(N+2)(1+h-j)$, does not involve $c_{\Gamma}$. The factor $O(1)$ may be taken at
least $1$, since enlarging it preserves the inequality; as $g_k\le e^{-1}$ and $1+h-j\ge1$, the
summand $b_0c_2R_k$ is then at most $g_k^{-2}O(1)\,b_0c_2R_k(1+h-j)$. With the new summand this
gives
\begin{align*}
g_k^{-2}+c'_{\Gamma}(k-j)+b_0c_2R_k
&\le g_k^{-2}O(1)\bigl(N+2+b_0c_2R_k\bigr)(1+h-j)\\
&\le g_k^{-2}O(1)(2+b_0c_2)R_k(1+h-j).
\end{align*}
The last step uses $N\le R_k$ and $2\le R_k$. Here $R_k\ge2$: by
Lemma~\ref{4.18:lem-older-window}, $R_k$ is the least power of $L$ with
$R_k\ge(\log g_k^{-2})^{r_0}$, and $(\log g_k^{-2})^{r_0}\ge2^{r_0}>1$ because $g_k\le e^{-1}$
and $r_0>0$; hence $R_k\ge L\ge3$, $L$ being odd and greater than $1$. Thus one more factor $R_k$
enters the bound that
takes the place of \cite[(1.47)]{B-CMP122b} in the present setting: where the right member of
\cite[(1.47)]{B-CMP122b} carries $R_k^7$, this bound carries $R_k^8$. By (IX.5), this bound is
small in the regime $L^{-N}\le(\log g_k^{-2})^{-\nu}$ as soon as
$\nu\ge p_0+\max\{p_0,q\}+8r_0+1$, and by (IX.6) this floor on $\nu$ does not depend on $L$.

The positivity used in \cite[(1.48)]{B-CMP122b} is
$(g''_k)^{-2}\zeta_1\ge g_k^{-2}\zeta_1\ge0$, which is hypothesis (X). It does not involve the
bound on $\mathfrak{d}_{\mathrm{old}}$ and is unchanged. For the companion system, (X) gives
$x''\zeta_1=(x^{\circ}+\mathfrak{d}_{\mathrm{old}})\zeta_1$. Where $\mathfrak{d}_{\mathrm{old}}$
vanishes this is $x''\zeta_1=x^{\circ}\zeta_1$. On $\Gamma''_j$ one has $x''=(g''_k)^{-2}$, and
$\zeta_1\ge0$ by (X), since $g_k^{-2}\zeta_1\ge0$ and $g_k^{-2}>0$; hence the first display of (5)
gives
\begin{equation*}
x''\zeta_1\le\bigl(g_k^{-2}+c'_{\Gamma}(k-j)+b_0c_2R_k\bigr)\zeta_1 ,
\end{equation*}
so the coefficient bound and the absorption with $R_k^8$ just proved apply to $x''\zeta_1$ as they
apply to $(g''_k)^{-2}$. This proves (5).

\emph{The pointwise bound on the excess.} Let $j'\le h$ and $x\in\mathrm{zone}_{j'}$. By (V) and
by the second inequality of \eqref{4.18:eq-graded-excess-a} with $j=j'$,
\begin{equation*}
|\mathfrak{d}^{(n)}(x)|\le b_0N_0+\mathfrak{d}_{\mathrm{old}}(x)
\le b_0\bigl(N_0+c_1(k-j'+1)+c_2R_k\bigr),
\end{equation*}
which is (6).
\end{proof}

\begin{theorem}[The domain of the plaquette-action excess in the large-field bounds]
\label{4.18:thm-action-domain}
Let $d=4$, and let the parameters satisfy
\begin{gather*}
L\ge 8,\qquad M=L^{m'}>\max\{M_1,M_0(\kappa,L,d)\},\\
\kappa\ge\kappa_0(d),\qquad \beta\le\tfrac14,\qquad \ell_0^2\le L/3,
\end{gather*}
where $\kappa_0(d)$ and $M_0(\kappa,L,d)$ are the floors in the hypotheses of the theorem on the
stage terms assumed in Theorem~\ref{4.18:thm-preparatory-boundary}, $\beta$ is the schedule
constant of the layered
analyticity spaces (Definition~\ref{4.4:def-post-step-space}), and $\ell_0$ is one of the
fixed auxiliary parameters of the construction, namely the one constrained by the inequality
$\ell_0^2\le L/3$ among the conditions of the theorem on the stage terms assumed in
Theorem~\ref{4.18:thm-preparatory-boundary}. Assume further the conditions on the exponents in
Ba{\l}aban's windows,
\begin{equation*}
q_*\ge p_1+1,\qquad p_0\ge 5r,
\end{equation*}
and the lower bound on the exponent $\nu$ of Ba{\l}aban's regime in the form fixed in
Proposition~\ref{4.18:prop-graded-excess}, which makes explicit the condition that $\nu$ be
sufficiently large under which the bounds \cite[(1.46)--(1.48)]{B-CMP122b} are stated, with the
summand $8r_0$ in the place of the summand $7r_0$ of the unshifted bound; here $r_0$ is the
constant so denoted in that proposition. Let Ba{\l}aban's absolute constants and the
constant $b_0$ be fixed. Let $c_1,c_2$ be the constants of \eqref{4.18:eq-older-window-a}.

Assume the following statements.
\begin{itemize}
\item The first Wilson-term theorem: the large-field estimate of the Wilson term of
\cite{B-CMP122b} for components satisfying its geometric hypothesis, with its support
statement for the excess and the cell count of $[Z''_k]_k$, the volume bound on the domain
$\mathcal{D}$ of \eqref{4.18:eq-action-domain-2} and the non-containment in $\Lambda$ that
accompany it, its locality statement and its bounds on the excess.
\item The second Wilson-term theorem: the same estimates for all components of Ba{\l}aban's
class, with its support statement, the cell count of $[Z''_k]_k$, the volume bound on
$\mathcal{D}$ and the non-containment in $\Lambda$ that accompany it, its locality statement
and its bounds on the excess,
excepting its bound on the old part $\mathfrak{d}_{\mathrm{old}}$ of the excess on the zones
$\mathrm{zone}_j=(\Omega''_j\setminus\Omega''_{j+1})\cap Z$ with $j\le h$.
\item Theorem~\ref{4.18:thm-preparatory-boundary} (bounds on the boundary terms of the
preparatory integrations), together with its standing hypotheses.
\item Lemma~\ref{4.18:lem-graded-count} (the graded count of older members made of
$M$-cubes), together with Proposition~\ref{4.18:prop-graded-excess} (the graded excess bound
in the large-field estimates).
\item The statements of Ba{\l}aban's papers on which the two Wilson-term theorems rest, in the
forms stated with those two theorems, namely:
\begin{itemize}
\item the statement fixing the small function $\mathfrak{m}$;
\item its graded form, fixing the small function $\mathfrak{m}'$;
\item the statement fixing the small function $\Xi$;
\item the bound of \cite{B-CMP122b} for the second remainder of the excess term of
\cite[(1.32)]{B-CMP122b}, whose right member is the quantity $\bar\sigma$ of clause~(ii);
\item the statement of \cite{B-CMP122b} on its cutoff functions;
\item three further statements of Ba{\l}aban's papers entering the proofs of the two
Wilson-term theorems;
\item the regularity statement for the stage terms.
\end{itemize}
\end{itemize}
Then there is $\gamma_W>0$, depending only on the data above and on the small functions
$\mathfrak{m}$, $\mathfrak{m}'$, $\Xi$, with the following property.

Let $k$ be a step at which the flow hypothesis on the running couplings (the hypothesis so
named in Theorem~\ref{4.18:thm-preparatory-boundary}) holds at all levels $\le k$, at which
$0<g_k\le\gamma_W$, and at which Ba{\l}aban's regime
\begin{equation}\label{4.18:eq-action-domain-1}
N_0<N\le R_k,\qquad L^{-N}\le\bigl(\log g_k^{-2}\bigr)^{-\nu}
\end{equation}
is in force. Let $Z$ be any component of the class of components treated by Ba{\l}aban's
operation $\mathbf{R}$ \cite{B-CMP122a}, the construction being run with the profiles
$\varphi_A$ of the family of prepared members at every preparatory stage. Put
\begin{equation}\label{4.18:eq-action-domain-2}
\mathcal{D}:=\bigcup_{i\le k}\bigl([Z''_i]_i\bigr)^{+\frac14 a_i},\qquad
\mathcal{N}:=(Z''_k)^{+\frac54 a_k}.
\end{equation}
Then the following hold.
\begin{enumerate}
\item[(i)] The support of the excess satisfies
\begin{equation*}
\operatorname{supp}\bigl((g''_k)^{-2}-g_k^{-2}\bigr)\cap Z\subset\mathcal{D}
\subset\mathcal{N}\subset Z^{\sim-8}.
\end{equation*}
The number of level-$k$ cells in $[Z''_k]_k$ is at most
$(\lceil 100/R_k\rceil+1)^d$, and
\begin{equation*}
|\mathcal{D}|\le\bigl(100M+\tfrac52 MR_k\bigr)^d .
\end{equation*}
In general $\mathcal{D}$ is not contained in the set $\Lambda$ of \cite[(1.73)]{B-CMP122a}.
\item[(ii)] Every estimate among \cite[(1.14)--(1.48)]{B-CMP122b} that reads the support or
the values of the excess holds with $\Lambda$ replaced by $\mathcal{D}$, in Ba{\l}aban's
form, as follows.
\begin{itemize}
\item The statement \cite[(1.14)]{B-CMP122b} reads: localized in
$\Lambda\cup\mathcal{D}\subset Z^{\sim-8}$.
\item The $Z$-part of the excess term of \cite[(1.32)]{B-CMP122b} has two remainders. The
first is bounded by Ba{\l}aban's bound for the companion, the construction run with the
companion profiles $\varphi^{\circ}_A$ of his admissible class, plus
\begin{equation*}
b_0\bigl((1+c_1)\mathfrak{m}'+(1+c_2)R_kg_k^2\mathfrak{m}\bigr).
\end{equation*}
The second is bounded by
\begin{equation*}
\bar\sigma+b_0c_1\mathfrak{m}'+\Xi\, e^{-\frac14\delta M_1R_k},
\end{equation*}
where $\bar\sigma$ denotes the right member of the bound for this remainder stated in
\cite{B-CMP122b}, in the form in which that statement is assumed above, and
$\delta>0$ is the decay constant of \cite{B-CMP122b}.
\item For every $j$, on $(\Omega''_j\setminus\Omega''_{j+1})\cap Z$ the graded bound reads
\begin{equation}\label{4.18:eq-action-domain-3}
0\le(g''_k)^{-2}-g_k^{-2}\le\bigl(O(1)+b_0c_1\bigr)(k-j)+b_0(1+c_2)R_k .
\end{equation}
\item The bound \cite[(1.37)]{B-CMP122b} holds in the form stated with the two Wilson-term
theorems; its second term is
Ba{\l}aban's term for the companion plus an increment at most $1$, and this increment is
included in the constant $O(1)$ of that bound.
\item The bounds \cite[(1.46)--(1.48)]{B-CMP122b} hold with $g_k^{-2}+O(k-j)$ replaced by
\begin{equation*}
g_k^{-2}+O'(k-j)+b_0c_2R_k ,
\end{equation*}
where $O'$ is the coefficient fixed in Proposition~\ref{4.18:prop-graded-excess}, and with
the factor $R_k^7$ of \cite[(1.47)]{B-CMP122b} replaced by $R_k^8$; the bound
\cite[(1.48)]{B-CMP122b} holds verbatim.
\item The quantity $C_k$ of \cite[(1.14)--(1.48)]{B-CMP122b}, into which the small terms of
these estimates are collected, remains a finite sum of small terms localized in $Z$.
\item Every estimate of \cite[(1.14)--(1.48)]{B-CMP122b} that reads one of the sets $\Lambda$
(the set of clause~(i)), $\Lambda_0$, $B_0$, $\Lambda^{(k)}$ as an
integration domain, and the regime $N\le R_k$, are unchanged.
\end{itemize}
For components satisfying the geometric hypothesis of the first Wilson-term theorem the
constants $c_1,c_2$ may be dropped, and clause (ii) is then the conclusion of the first
Wilson-term theorem on the excess, equivalently that of the second with
$\mathfrak{d}_{\mathrm{old}}\equiv0$. Every component satisfies that geometric hypothesis
under the first two of the three readings of the description of the older
$\mathbf{R}$-operations that is the hypothesis of
Lemma~\ref{4.18:lem-graded-count}. Under the
third description, clause (ii) is the second Wilson-term theorem with its bound on
$\mathfrak{d}_{\mathrm{old}}$ on the zones of index $j\le h$ replaced by
\eqref{4.18:eq-graded-excess-a} and with the three substitutions of
Proposition~\ref{4.18:prop-graded-excess}.
\item[(iii)] At every preparatory stage $n+1$, $0\le n\le N-1$, with $j=h+n$, the family's
term $\mathbb{C}^{(n+1)}_k$ has the following properties:
\begin{itemize}
\item it has the representation \cite[(1.63)]{B-CMP122a};
\item it is holomorphic on the spaces \cite[(1.64)--(1.69)]{B-CMP122a}, under the reading of
these spaces adopted in Theorem~\ref{4.18:thm-preparatory-boundary};
\item it obeys \cite[(1.70)]{B-CMP122a} with Ba{\l}aban's constant $C_0$, which the theorem
on the stage terms assumed in Theorem~\ref{4.18:thm-preparatory-boundary} gives as
$C_0=c_G+1$, and with that theorem's decay rates (Remark~\ref{4.18:rem-decay-rates}).
\end{itemize}
The weight-difference member of the family in \cite[(1.62)]{B-CMP122a} is supported, beyond
the support of the companion's member, on the $\frac14 a_{j+1}$-collar of
$[Z''_{j+1}]_{j+1}$ inside $R_n$. Its share of $C_0$ equals the companion's share, so the
constant $C_0$ is unchanged.
\item[(iv)] The set $\mathcal{D}$, the support $\mathcal{E}'$ of the current part of the
excess weight, and the supports of all modified terms
lie in $Z^{\sim-8}$, and $Z^{\sim-8}\subset X_a$, where $X_a$ is the locality set of the
locality statement of the two Wilson-term theorems.
\end{enumerate}
\end{theorem}

\begin{proof}
\emph{Clauses} (i) \emph{and} (iv). Clause (i) consists of the support statements of the two
Wilson-term theorems, which are hypotheses of the theorem, together with the cell count of
$[Z''_k]_k$, the volume bound on $\mathcal{D}$ and the non-containment in $\Lambda$ that
accompany them among those hypotheses. Clause (iv) consists of the locality
statements of the two Wilson-term theorems, also hypotheses, together with two further inputs.
The first is part~(b) of Lemma~\ref{4.18:lem-weight-difference}, which places the set where
the two profiles of $Z''_{j+1}$ differ inside $Z^{\sim-8}$. The second is the bound
\eqref{4.18:eq-graded-count-b} of Lemma~\ref{4.18:lem-graded-count}, which controls the older
members made of $M$-cubes.

\emph{Clause} (ii). Consider first a component satisfying the geometric hypothesis of the
first Wilson-term theorem. For it, clause (ii) with $c_1,c_2$ dropped is the conclusion of
that theorem on the excess. Equivalently, it is the conclusion of the second Wilson-term
theorem with $\mathfrak{d}_{\mathrm{old}}\equiv0$. In that case the bound of the second
theorem on $\mathfrak{d}_{\mathrm{old}}$ on the zones of index $j\le h$, excepted from the
hypothesis, is not used, since $\mathfrak{d}_{\mathrm{old}}$ vanishes there.

Consider next an arbitrary component under the third description of the older
$\mathbf{R}$-operations. We apply the second Wilson-term theorem, which is a
hypothesis for everything except its bound on $\mathfrak{d}_{\mathrm{old}}$ on the zones
$\mathrm{zone}_j$ with $j\le h$. We replace that bound by the graded bound
\eqref{4.18:eq-graded-excess-a} of Proposition~\ref{4.18:prop-graded-excess}, which holds on
$\mathrm{zone}_j$ as a consequence of Lemma~\ref{4.18:lem-graded-count}. We then carry out the
three substitutions of
Proposition~\ref{4.18:prop-graded-excess} in the estimates that read this bound. These
substitutions are made under the two strengthened ceilings on $g_k$ fixed in that
proposition, and under the shifted lower bound on $\nu$ of that proposition. Both are
among the hypotheses of the theorem: the ceilings through $g_k\le\gamma_W$, as arranged
below, and the bound on $\nu$ directly. The result is the displayed list of clause (ii),
including the graded bound \eqref{4.18:eq-action-domain-3}, the two increments displayed in
clause (ii) for the remainders of \cite[(1.32)]{B-CMP122b}, and the replacements in
\cite[(1.46)--(1.48)]{B-CMP122b}.

\emph{Clause} (iii). This is Theorem~\ref{4.18:thm-preparatory-boundary}, a hypothesis of the
theorem together with its standing hypotheses. The statement on the support of the
weight-difference member of
\cite[(1.62)]{B-CMP122a} is Lemma~\ref{4.18:lem-weight-difference}, on which that theorem
rests.

\emph{The ceiling} $\gamma_W$. Define $\gamma_W$ as the least of the following numbers:
\begin{itemize}
\item the ceilings on $g_k$ under which the second Wilson-term theorem is stated (three of
them), one of them taken in the strengthened form of
Proposition~\ref{4.18:prop-graded-excess};
\item the ceilings on $g_k$ under which the first Wilson-term theorem is stated (four of
them), one of them taken in the strengthened form of that proposition;
\item the number $e^{-1}$;
\item the ceiling $\gamma_0$ in the hypotheses of the theorem on the stage terms assumed in
Theorem~\ref{4.18:thm-preparatory-boundary}.
\end{itemize}
Each of these ceilings is the threshold of a one-sided condition, displayed with the statement
to which it belongs, whose left member is a continuous function of $g$ that decreases to
$0$ with $g$. This left member is either a power of $g$ times a power of $\log g^{-2}$, or
Ba{\l}aban's $\mathfrak{m}'$, or $\Xi\,e^{-cM_1R(g)}$ with a constant $c>0$, where $R(g)$ is the
cell parameter as a function of the coupling. Hence each condition holds for
all sufficiently small $g>0$, so the least of these finitely many
thresholds is positive and $\gamma_W>0$ exists. It depends only on the data listed in the
statement and on $\mathfrak{m}$, $\mathfrak{m}'$, $\Xi$. No value of $\gamma_W$ is asserted.
\end{proof}

\begin{lemma}[Two-sided flow bounds from per-step bounds]\label{4.19:lem-flow-bounds}
Let $n\ge 1$. Let $x_0$ and $\beta_l=\beta^{\mathfrak{A}}_l+\beta^{\mathfrak{R}}_l$, $1\le l\le n$, be real
numbers, and let $\beta^0_l$, $1\le l\le n$, be real numbers. Put $x_l=x_{l-1}-\beta_l$ for
$1\le l\le n$. Let $c_0>0$, $c_2\ge 0$, $C_\beta\ge 0$ and $\bar\varsigma\ge 0$ be constants,
and assume that for all $1\le l\le n$ and all $0\le i\le j\le n$
\begin{gather}
\sum_{l=i+1}^{j}\beta^0_l\ \ge\ (j-i)\,c_0-2c_2,\label{4.19:eq-flow-bounds-1}\\
\begin{split}
&|\beta^{\mathfrak{A}}_l-\beta^0_l|\le \tfrac12 c_0,\qquad |\beta^{\mathfrak{A}}_l|\le C_\beta,\\
&|\beta^{\mathfrak{R}}_l|\le\bar\varsigma\le\min\{\tfrac14 c_0,1\}.
\end{split}\label{4.19:eq-flow-bounds-2}
\end{gather}
Put $\lambda_\sharp:=\tfrac12 c_0-\bar\varsigma$, $B_\sharp:=C_\beta+\bar\varsigma$ and
$c_5:=2c_2$. Then $\lambda_\sharp\ge c_0/4$, and for all $0\le i\le j\le n$
\begin{equation}\label{4.19:eq-flow-bounds-3}
\lambda_\sharp\,(j-i)-c_5\ \le\ x_i-x_j\ \le\ B_\sharp\,(j-i).
\end{equation}
\end{lemma}

When the lemma is applied at level $k+1$ of the induction over the levels, that is with
$n=k+1$, the inequality \eqref{4.19:eq-flow-bounds-1} for all $0\le i\le j\le k+1$ is supplied
by the finite part of the one-loop law (Lemma~\ref{4.2:lem-one-loop}, whose proof is incomplete
at one step), since by that lemma
\begin{equation*}
\sum_{l=i+1}^{j}\beta^0_l=\sum_{l\le j}\beta^0_l-\sum_{l\le i}\beta^0_l
\ge(jc_0-c_2)-(ic_0+c_2)=(j-i)\,c_0-2c_2 ;
\end{equation*}
the per-step bounds \eqref{4.19:eq-flow-bounds-2} at $l=k+1$, in which
$\beta^{\mathfrak{A}}_{k+1}$ is the unmarked part and $\beta^{\mathfrak{R}}_{k+1}$ the
large-field part of the coupling increment, are supplied by the bounds
on the unmarked part of the increment (Lemma~\ref{4.2:lem-unmarked}, whose proof is incomplete
at one step) together with the ceiling $\varepsilon_1^{\mathrm{I}}\le c_0/(2C_{\mathrm{rem}})$,
where $C_{\mathrm{rem}}$ is the constant of that lemma in its bound
$|\beta^{\mathfrak{A}}_{k+1}-\beta^0_{k+1}|\le C_{\mathrm{rem}}\varepsilon_1^{\mathrm{I}}$,
and by the bound in the statement on the large-field part of the coupling increment,
together with the ceiling $C'\gamma^{\kappa_0-8}\le\min\{c_0/4,1\}$ on the constant
$C'$ and the exponent $\kappa_0$ of that bound, fixed in the order of choice of the constants.

\begin{proof}
By the definition $x_l=x_{l-1}-\beta_l$, telescoping gives, for $0\le i\le j\le n$,
\begin{equation*}
x_i-x_j=\sum_{l=i+1}^{j}(x_{l-1}-x_l)=\sum_{l=i+1}^{j}\beta_l .
\end{equation*}

For the upper bound, each step satisfies, by \eqref{4.19:eq-flow-bounds-2},
\begin{equation*}
\beta_l\le|\beta^{\mathfrak{A}}_l|+|\beta^{\mathfrak{R}}_l|\le C_\beta+\bar\varsigma=B_\sharp .
\end{equation*}
Summing this over the $j-i$ indices $l=i+1,\dots,j$ gives $x_i-x_j\le B_\sharp\,(j-i)$, which
is the right inequality of \eqref{4.19:eq-flow-bounds-3}.

For the lower bound, write each step as
\begin{equation*}
\beta_l=\beta^0_l+(\beta^{\mathfrak{A}}_l-\beta^0_l)+\beta^{\mathfrak{R}}_l
\ge\beta^0_l-\tfrac12 c_0-\bar\varsigma,
\end{equation*}
where the inequality uses $|\beta^{\mathfrak{A}}_l-\beta^0_l|\le\tfrac12 c_0$ and
$|\beta^{\mathfrak{R}}_l|\le\bar\varsigma$ from \eqref{4.19:eq-flow-bounds-2}. Summing over
$l=i+1,\dots,j$ and using \eqref{4.19:eq-flow-bounds-1},
\begin{equation*}
\begin{split}
x_i-x_j&\ge\sum_{l=i+1}^{j}\beta^0_l-(j-i)\bigl(\tfrac12 c_0+\bar\varsigma\bigr)\\
&\ge(j-i)c_0-2c_2-(j-i)\bigl(\tfrac12 c_0+\bar\varsigma\bigr)
=\lambda_\sharp\,(j-i)-c_5 ,
\end{split}
\end{equation*}
which is the left inequality of \eqref{4.19:eq-flow-bounds-3}.

Finally, $\bar\varsigma\le c_0/4$ from \eqref{4.19:eq-flow-bounds-2} gives
$\lambda_\sharp=\tfrac12 c_0-\bar\varsigma\ge c_0/4$.
\end{proof}

\begin{remark}
No per-step sign of $x_{l-1}-x_l$ is asserted in Lemma~\ref{4.19:lem-flow-bounds}, and none is
used in what follows.
\end{remark}

\begin{proposition}[Clause (0) concluded]\label{4.19:prop-clause-zero}
Assume (H1)--(H6) of Notation~\ref{2.2:not-hypotheses}. Let $(x_k)$, $x_k=g_k^{-2}$, be the running
inverse couplings of Definition~\ref{1.2:def-couplings}, started at $g_0\in(0,g_-(g)]$. Assume
further the following statements:
\begin{enumerate}
\item[(a)] Ba{\l}aban's one-step renormalization theorem, \cite[Theorem~3]{B-CMP109};
\item[(a$'$)] the widened form of condition \cite[(1.12)]{B-CMP109}, on cubes smaller than those
of that condition, and the statement that the step of the construction reproduces this widened
form;
\item[(b)] the finite part of the one-loop law, Lemma~\ref{4.2:lem-one-loop}, and the bounds on
the unmarked part of the increment, Lemma~\ref{4.2:lem-unmarked}, each stated with its proof
incomplete at the step marked $(\ast)$ there;
\item[(c)] the bound on the large-field part of the coupling increment;
\item[(d)] the independence of the localised terms from the small-field region,
Corollary~\ref{4.4:cor-localised-independence};
\item[(e)] the operator bounds at free current;
\item[(f)] the representation and bound of the large-field terms;
\item[(g)] the one-power hereditary class;
\item[(h)] the torus-independence of the terms localised in non-wrapping domains.
\end{enumerate}
Assume these jointly with the conclusion of Theorem~\ref{4.20:thm-representation}, the inductive
representation of the effective densities, at the scales already reached: for every $n\ge1$ with
$x_i\ge\gamma^{-2}$ for all $i<n$, that conclusion holds at every level $k<n$. Then
$K=K(g_0,g)$ of Definition~\ref{1.2:def-flow-constants} is a well-defined non-negative integer
satisfying
\begin{equation}\label{4.19:eq-clause-zero-1}
\frac{g_0^{-2}-g^{-2}}{B_\sharp}-1\;\le\;K\;\le\;
1+\frac{g_0^{-2}-g^{-2}-B_\sharp+c_5}{\lambda_\sharp}.
\end{equation}
Consequently $K\to\infty$ as $g_0\downarrow0$, and the set of realised cutoffs is unbounded.
Moreover:
\begin{itemize}
\item $g_k\in(0,g)\subset(0,\gamma_H]$ for every $k\le K$;
\item $g_K\in[g_-(g),g)$;
\item for $0\le i\le j\le K$,
\begin{equation}\label{4.19:eq-clause-zero-2}
\lambda_\sharp(j-i)-c_5\;\le\;x_i-x_j\;\le\;B_\sharp(j-i);
\end{equation}
\item the $K$ renormalization steps on $\mathbb{T}_m$ at lattice spacing $L^{-K}$ produce exactly
the couplings $g_0,\dots,g_K$.
\end{itemize}
\end{proposition}

\begin{proof}
\emph{Structure of the argument.} Clause (0) and Theorem~\ref{4.20:thm-representation} are proved
by one induction on the level, which is carried out in the proof of
Theorem~\ref{4.20:thm-representation}; the present proof concerns the couplings. Its inputs at each
level are the statements (a), (a$'$) and (b)--(h), among them Lemmas~\ref{4.2:lem-one-loop}
and~\ref{4.2:lem-unmarked}, each stated with its proof incomplete at the step marked $(\ast)$ there.
Together with the bounds of clause (i), that is, with
Theorem~\ref{4.20:thm-representation} at all earlier levels, these statements give, for every
$n\ge1$ such that $x_i\ge\gamma^{-2}$ for all $i<n$, the two facts below; for the second they supply
the bounds on the parts of the increments at the levels up to $n$ that the lemma on the two-sided
flow bounds, Lemma~\ref{4.19:lem-flow-bounds}, the main step of the induction, converts into the
flow bounds:
\begin{itemize}
\item the increment $\beta_n$ of Definition~\ref{4.2:def-increment} is defined and is a function
of $g_0,\dots,g_{n-1}$ only;
\item by Lemma~\ref{4.19:lem-flow-bounds}, the two-sided bounds \eqref{4.19:eq-clause-zero-2} hold
for $0\le i\le j\le n$, with the
constants $\lambda_\sharp$, $B_\sharp$, $c_5$ of Definition~\ref{1.2:def-flow-constants}.
\end{itemize}
Here Lemma~\ref{4.19:lem-flow-bounds} is applied with $c_0$, $c_2$ and $C_\beta$ the constants of
Definition~\ref{1.2:def-flow-constants}, and with $\bar\varsigma$ the bound on the large-field part
of the increment supplied by (c), so that the constants $\lambda_\sharp$, $B_\sharp$, $c_5$ which
that lemma defines from them are those fixed in Definition~\ref{1.2:def-flow-constants}.

\emph{All members but the last.} By (H4), $g^{-2}\ge\gamma_H^{-2}$. By the ceiling on $\gamma_H$
in Remark~\ref{2.3:rem-thresholds}, $\gamma_H^{-2}\ge\gamma^{-2}+c_5+1>\gamma^{-2}$. Hence
\begin{equation*}
g^{-2}\;\ge\;\gamma_H^{-2}\;>\;\gamma^{-2}.
\end{equation*}
Since $g_0\le g_-(g)=(g^{-2}+B_\sharp)^{-1/2}$,
we have $t:=g_0^{-2}\ge g^{-2}+B_\sharp$. Now let $n\ge1$ be such that
$x_i>g^{-2}+B_\sharp$ for all $i<n$. Then $x_i>\gamma^{-2}$ for all $i<n$, and the second fact
of the preceding paragraph gives \eqref{4.19:eq-clause-zero-2} for $0\le i\le j\le n$. Thus the
sequence $(x_k)$, started at $x_0=t\ge g^{-2}+B_\sharp$, obeys the two-sided flow bounds on every
initial segment $[0,n]$ with $x_i>g^{-2}+B_\sharp$ for all $i<n$. This is
the hypothesis of the first-passage arithmetic,
Lemma~\ref{4.2:lem-first-passage}, with this $t$. That lemma yields the following:
\begin{itemize}
\item $K$ is a well-defined non-negative integer obeying \eqref{4.19:eq-clause-zero-1};
\item $g_k\in(0,g)$ for $k\le K$;
\item $g_K\in[g_-(g),g)$.
\end{itemize}
If $K\ge1$, minimality of $K$ gives $x_i>g^{-2}+B_\sharp$ for $i<K$, so the preceding applies with
$n=K$ and gives \eqref{4.19:eq-clause-zero-2} for $0\le i\le j\le K$. If $K=0$, the only pair is
$i=j=0$, for which
\eqref{4.19:eq-clause-zero-2} reads $-c_5\le0\le0$; this holds since $c_5=2c_2\ge0$.
The inclusion $(0,g)\subset(0,\gamma_H]$ is $g\le\gamma_H$, which is (H4). As $g_0\downarrow0$,
the quantity $g_0^{-2}$ tends to infinity. Hence so does the lower bound in
\eqref{4.19:eq-clause-zero-1}. Thus $K\to\infty$, and the set of realised cutoffs is unbounded.
This gives every member of the proposition except the last.

\emph{The last member.} It rests on two facts. The first is that the increments depend neither on
$K$ nor on the torus: not on $K$, because by the first fact of the structure paragraph $\beta_n$
is a function of $g_0,\dots,g_{n-1}$ only; not on the torus, by (h). The second concerns the side
lengths of the cubes used by the step at level $j$, for $j\le K$. Through the bounds
\eqref{4.19:eq-clause-zero-2} on the generated sequence, these side lengths are bounded in terms of
$R(g_-(g))$ and $L^{K-j}$. Through these bounds, (H5) places the declared pair $(m,K)$ in the set of
pairs admissible for the step at level $j$, for every $j\le K$. Together the two facts show that
the $K$ steps performed on $\mathbb{T}_m$ at spacing $L^{-K}$ subtract exactly the increments of the
sequence $(x_k)$. That sequence was generated before any lattice was declared. So defining $K$
afterwards, as the first passage of that sequence, is consistent, and the $K$ steps produce
exactly $g_0,\dots,g_K$. No upper bound on $K$, $m$ or $L$ is imposed.
\end{proof}

These are the assertions of clause (0) of Theorem~\ref{3.1:thm-main} as restated in
Theorem~\ref{4.1:thm-clause-zero}; the paragraph on fixed-depth windows stated there takes them as
its hypothesis.

\begin{proposition}[Containment: Ba{\l}aban's one coupling hypothesis holds along the flow]
\label{4.19:prop-containment}
Assume hypothesis (H4) of Notation~\ref{2.2:not-hypotheses}, clause (0) of Theorem~\ref{3.1:thm-main}
in the form of Proposition~\ref{4.19:prop-clause-zero}, and the terminal ceiling
$\gamma_H\le(\gamma^{-2}+c_5+1)^{-1/2}$ of hypothesis (H3) (see
Remark~\ref{2.3:rem-thresholds}). Then for all $k\le K$
\begin{equation}\label{4.19:eq-containment-1}
x_k>g^{-2}\ge\gamma_H^{-2}\ge\gamma^{-2}+c_5+1,
\end{equation}
so $0<g_k<\gamma$. In particular:
\begin{itemize}
\item every lower bound $x_k\ge x_*$ of the construction with $x_*\le\gamma^{-2}+c_5+1$ holds
at every level $k\le K$;
\item the one hypothesis of Ba{\l}aban's ultraviolet stability theorem
(\cite[Theorem~1]{B-CMP122b}), namely that all
effective couplings lie in $(0,\gamma]$, is satisfied.
\end{itemize}
\end{proposition}

\begin{proof}
Clause (0) (Proposition~\ref{4.19:prop-clause-zero}) gives $0<g_k<g$ for every $k\le K$.
Since $x_k=g_k^{-2}$ and $t\mapsto t^{-2}$ is strictly decreasing on $(0,\infty)$, this is
$x_k>g^{-2}$, the first inequality of \eqref{4.19:eq-containment-1}. By (H4),
$g\le\gamma_H$, that is, $g^{-2}\ge\gamma_H^{-2}$, the second inequality. The terminal ceiling
$\gamma_H\le(\gamma^{-2}+c_5+1)^{-1/2}$ of (H3), both sides being positive, is equivalent to
$\gamma_H^{-2}\ge\gamma^{-2}+c_5+1$, the third inequality; and
$\gamma^{-2}+c_5+1>\gamma^{-2}$. Hence $x_k>\gamma^{-2}$, and since $x_k>0$ we obtain
\begin{equation*}
g_k=x_k^{-1/2}<\gamma .
\end{equation*}
Together with $g_k>0$ this is $0<g_k<\gamma$.

For the first consequence, let $x_*\le\gamma^{-2}+c_5+1$. Then \eqref{4.19:eq-containment-1}
gives $x_k>\gamma^{-2}+c_5+1\ge x_*$ for every $k\le K$. For the second, $0<g_k<\gamma$ for
every $k\le K$ places every effective coupling in $(0,\gamma]$, which is the hypothesis of
Ba{\l}aban's ultraviolet stability theorem, \cite[Theorem~1]{B-CMP122b}.
\end{proof}

\begin{corollary}[Coupling windows at fixed depth and subsequential limit couplings]
\label{4.19:cor-windows}
Assume clause~(0) of Theorem~\ref{3.1:thm-main} (Proposition~\ref{4.19:prop-clause-zero}). Let
$g_0^{(i)}\downarrow 0$ be a sequence of bare couplings, and put $K_i:=K(g_0^{(i)},g)$. Then, for
every $s\ge 0$ and every $i$ with $K_i\ge s$,
\begin{equation}\label{4.19:eq-windows-1}
x_{K_i-s}\in\bigl(g^{-2}+\max\{0,\lambda_\sharp s-c_5\},\ g^{-2}+B_\sharp(s+1)\bigr].
\end{equation}
Along a subsequence, $g_{K_i-s}\to g_{\infty,s}$ for every $s\ge 0$, where the limits
$g_{\infty,s}$ satisfy
\begin{equation}\label{4.19:eq-windows-2}
g_{\infty,s}^{-2}\in\bigl[g^{-2}+\max\{0,\lambda_\sharp s-c_5\},\ g^{-2}+B_\sharp(s+1)\bigr];
\end{equation}
in particular $g_{\infty,0}\in[g_-(g),g]$.
\end{corollary}

\begin{proof}
Recall $x_k=g_k^{-2}$. Clause~(0) contains the two-sided flow bound
\begin{equation*}
\lambda_\sharp(l-j)-c_5\le x_j-x_l\le B_\sharp(l-j),\qquad 0\le j\le l\le K.
\end{equation*}
Fix $s\ge0$ and $i$ with $K_i\ge s$, and apply this bound at the bare coupling $g_0^{(i)}$
with $j=K_i-s$ and $l=K_i$. It gives
\begin{equation*}
\lambda_\sharp s-c_5\le x_{K_i-s}-x_{K_i}\le B_\sharp s .
\end{equation*}
Clause~(0) also gives $g_{K_i}\in[g_-(g),g)$. Since $g_{K_i}<g$, we have $x_{K_i}>g^{-2}$.
Since $g_{K_i}\ge g_-(g)=(g^{-2}+B_\sharp)^{-1/2}$, we have $x_{K_i}\le g^{-2}+B_\sharp$.
Adding these bounds on $x_{K_i}$ to the two preceding inequalities gives
\begin{equation}\label{4.19:eq-windows-3}
g^{-2}+\lambda_\sharp s-c_5<x_{K_i-s}\le g^{-2}+B_\sharp(s+1).
\end{equation}
Clause~(0) further gives $g_k\in(0,g)$ for every $k\le K$. Taking $k=K_i-s$ yields
$x_{K_i-s}>g^{-2}$. Together with \eqref{4.19:eq-windows-3} this is \eqref{4.19:eq-windows-1}.

By the lower bound on $K$ in clause~(0), $K_i\to\infty$ as $g_0^{(i)}\downarrow0$. Hence, for
each $s$, the condition $K_i\ge s$ holds for all but finitely many $i$. By
\eqref{4.19:eq-windows-1}, the numbers $x_{K_i-s}$ with $K_i\ge s$ all lie in the compact
interval
\begin{equation*}
I_s:=\bigl[g^{-2}+\max\{0,\lambda_\sharp s-c_5\},\ g^{-2}+B_\sharp(s+1)\bigr]\subset(0,\infty).
\end{equation*}

We now make a diagonal extraction. By compactness of $I_0$, choose a subsequence along which
$x_{K_i}$ converges. Inductively, given the subsequence for $s-1$, discard its finitely many
terms with $K_i<s$. By compactness of $I_s$, choose a further subsequence along which
$x_{K_i-s}$ converges. Let the diagonal sequence have as its $n$-th term the $n$-th term of the
$n$-th subsequence. For each $s$, its terms from the $s$-th onward form a subsequence of the
$s$-th subsequence. Hence $x_{K_i-s}$ converges along the diagonal sequence for every
$s\ge0$, and the limit lies in the closed interval $I_s$.

The map $x\mapsto x^{-1/2}$ is continuous on $(0,\infty)$, and $I_s\subset(0,\infty)$.
Therefore every $g_{K_i-s}=x_{K_i-s}^{-1/2}$ converges along the diagonal sequence. Call its
limit $g_{\infty,s}$. Then $g_{\infty,s}^{-2}$ is the limit of $x_{K_i-s}$, so it lies in
$I_s$; this is \eqref{4.19:eq-windows-2}.

For $s=0$, since $c_5\ge0$ (Definition~\ref{1.2:def-flow-constants}), the interval is
\begin{equation*}
I_0=[g^{-2},\ g^{-2}+B_\sharp].
\end{equation*}
Thus $g_{\infty,0}^{-2}\in I_0$. Applying $x\mapsto x^{-1/2}$, which is
decreasing, gives $g_{\infty,0}\in[(g^{-2}+B_\sharp)^{-1/2},g]=[g_-(g),g]$.
\end{proof}

\begin{remark}[What clause (0) does not assert]\label{4.19:rem-not-asserted}
Clause (0) of Theorem~\ref{3.1:thm-main}, in the form of Theorem~\ref{4.1:thm-clause-zero}
and as concluded in Proposition~\ref{4.19:prop-clause-zero}, asserts none of the following.
\begin{itemize}
\item That for a prescribed $K$ some bare coupling $g_0$ gives $g_K=g$.
\item That every integer $K$ is realised as $K(g_0,g)$ for some $g_0$, or that $K(\cdot,g)$ is
monotone; nor the monotonicity of a tuned bare coupling $g_0$.
\item Any per-step sign, either of the differences $x_k-x_{k+1}$ or of the increments
$\beta_k$.
\item Convergence of $g_{K-s}$, for fixed $s$, as $g_0\to0$ along first passage
$K=K(g_0,g)$, without passing to subsequences; only convergence of $g_{K_i-s}$ along a
subsequence, with limit in the window given by the corollary on fixed-depth windows of
clause (0) in this section, is asserted.
\item Any regularity of the increments $\beta_j$ in the couplings.
\end{itemize}
\end{remark}

\begin{proof}
We explain why each item lies outside clause (0).

The existence and uniqueness of a bare coupling giving $g_K=g$ for a prescribed $K$ is one
of the conclusions of clause (v) of Theorem~\ref{3.1:thm-main} (Chapter~\ref{ch:9}), which
is stated for $N=2$ under hypothesis (H7). The
two-sided flow bounds of Theorem~\ref{4.1:thm-clause-zero} alone do not yield
it. Indeed, taking $i=0$ and $j=K$ in them gives
\begin{equation*}
\lambda_\sharp K-c_5\le x_0-x_K\le B_\sharp K ,
\end{equation*}
so that for a prescribed $K$ they place $x_K$ only in the interval
\begin{equation*}
\bigl[x_0-B_\sharp K,\ x_0-\lambda_\sharp K+c_5\bigr],
\end{equation*}
whose length $(B_\sharp-\lambda_\sharp)K+c_5$ grows with $K$. Without continuity of the
map $x_0\mapsto x_K$, no choice of $x_0$ is guaranteed to put $x_K$ in a prescribed window
about $g^{-2}$.

Continuity of the increments in the whole history of couplings would supply that
continuity. Regularity in the last coupling is stated in
\cite[after (1.18), after (1.22)]{B-CMP109}; no proof is given there. Namely, it is stated
for the terms denoted $E^{(j)}$ as functions of $g_{j-1}$ (the inductive assumption
following \cite[(1.18)]{B-CMP109}) and for $\beta_{j+1}$ as a function of $g_j$ (the
sentence following \cite[(1.22)]{B-CMP109}). Analyticity in all effective couplings is
stated there only for the alternative small-field threshold introduced after
\cite[(2.9)]{B-CMP109}, which depends on $K-k$ and is not used here. For the construction
used in this book, continuity of $\beta_{k+1}$ in $(g_0,\dots,g_k)$ is not stated in
\cite{B-CMP109} for this threshold and is not used here. This accounts for the first and the
last item.

No per-step sign follows from these flow bounds: the lower member
$\lambda_\sharp(j-i)-c_5\le x_i-x_j$ holds only over stretches of steps and carries the
defect $c_5$.

Finally, the corollary on fixed-depth windows of clause (0) in this section confines, along a
sequence $g_0^{(i)}\downarrow0$ with $K_i=K(g_0^{(i)},g)$ and for every $s$ with $K_i\ge s$,
the value $x_{K_i-s}$ to a bounded interval; this yields convergence of $g_{K_i-s}$ along a
subsequence only, and nothing in clause (0) identifies the limits of different
subsequences.
\end{proof}

\begin{definition}[Standing data of the construction: numbers, partitions and cut-offs]
\label{4.38:def-standing-data}
Throughout, $j$ denotes the level of the renormalisation step at hand, and lengths are measured in
level-$j$ units. For $i\le j$ the \emph{level-$i$ cells} are the cubes of the standard partition at
level $i$ whose side is $M\check{R}_i$ in level-$i$ units. In level-$j$ units their side is
\begin{equation}\label{4.38:eq-standing-data-1}
  a_i := M\check{R}_i\,L^{-(j-i)},\qquad\text{so that } a_j = M\check{R}_j .
\end{equation}
For a set $W$ we write $[W]_i$ for the smallest union of level-$i$ cells containing $W$. For a
label $\ell$ we write $\Omega_i(\ell)\supset\Lambda_i(\ell)$ for its sequences of localisation
domains of \cite[(2.1)]{B-CMP119}, and $S_i\subset\Omega_i(\ell)\cap\Lambda_i(\ell)^c$, $i\le j$,
for its third sequence of \cite[p.~254]{B-CMP119}, each $S_i$ empty or a union of
$LM_2R_i$-cubes, with $R_i$ as in Lemma~\ref{4.4:lem-collar-avoidance}.
At an $\mathbf{R}$-step, Ba{\l}aban's index $k$ in \cite{B-CMP122a,B-CMP122b} is the level of that
step, so that $k=j$. The prepared objects $Z''_k$, $g''_k$, $U''_k$ below carry this index $k$,
and so does the integer $m_\sharp$ of \eqref{4.38:eq-standing-data-b}. We write $\mathsf{N}$ for
the number of preparatory integrations of an $\mathbf{R}$-step, which is written $N$ in
\cite{B-CMP122a}.

The setting of the construction, its family of small-field and large-field
functions, the walk-expansion rule of Definition~\ref{4.4:def-expansion-rule}, and the split
$\mathbb{B}$ of the data of each step of the family enter the torus-independence lemma of
clause~(0) only through the following data and properties \textup{(L0)}--\textup{(L5)}.

\smallskip
\noindent\textup{(L0)} \emph{The scheme.} The scheme $\mathcal{M}$, that is, the items
\textup{(E-a)}--\textup{(E-f)} of its definition in the setting of the construction.

\smallskip
\noindent\textup{(L1)} \emph{Numbers.} Every threshold, width, cube size and counter (such as
$k_1$ and $\mathsf{N}$ below) of the construction is a function of the auxiliary parameters
$\mathfrak{p}$ and of the couplings already generated. In particular the following hold.
\begin{itemize}
\item The threshold $\varepsilon_k=g_k\,p_0(g_k)$ is given by the degree function $p_0(\cdot)$,
  and $\varepsilon_0$ is also given by this formula, as assumed in \cite{B-CMP119}.
\item The threshold $\delta_k=(A_1/A_0)\,\varepsilon_k$ is that of \cite{B-CMP119}.
\item The widths $\delta_k^{\flat}$ and $\theta_c\delta_k$, with the factor $\theta_c$, are those
  fixed by the family's rules at the split $\mathbb{B}$.
\item The cell-size multiplier $\check{R}_k$ is a function of its argument $\check{x}_k$, fixed in
  the definition of the cell-size multiplier; by the lemma on the cell-size multiplier,
  $\check{R}_k\ge L$.
\item The index $k_1$ is the largest integer satisfying the inequality of \cite{B-CMP119},
  which is read with the index $k_1$ on $R$. Here $R_i$ is as in
  Lemma~\ref{4.4:lem-collar-avoidance}; it is not the ball radius $R$. The inequality reads
  \begin{equation*}
    L^{-(k-k_1)}\,M R_{k_1}\le L .
  \end{equation*}
\item The number of preparatory integrations is $\mathsf{N}:=\check{R}_k$.
\end{itemize}
None of these numbers depends on the number of renormalisation steps. The choices of
\cite{B-CMP109} and \cite{B-CMP119} that make a threshold depend on that number are not used.

\smallskip
\noindent\textup{(L2)} \emph{Geometry.}
\begin{itemize}
\item \emph{Partitions.} Every partition of the construction is either the standard partition
  into cubes with corners in $a\mathbb{Z}^4$, for some $a\in L^{\mathbb{Z}}$
  \cite{B-CMP109}, or a fixed subdivision of it.
\item \emph{Cut-offs are natural.} Every auxiliary cut-off and partition function is
  \emph{natural}, that is, invariant under the simultaneous action on its argument and on its
  data of the lattice motions that preserve the partitions it reads. Except for the coupling
  profile below, each smooth cut-off is a fixed function of the sup-distances to the faces of
  its domain.
\item \emph{Cut-offs attached to standard cubes or sites.} This applies to the functions of
  the partitions of unity attached to standard cubes, to the tent functions attached to
  sites, and to the interpolation profiles. Such a cut-off is a fixed profile of $x$ minus
  the centre of its cube or site.
\item \emph{Cut-offs attached to a label.} Such a cut-off is a fixed natural function of the
  position of $x$ relative to the label's cells. Its value at $x$ is determined by the cells
  at sup-distance less than one cell side from $x$.
\item \emph{Cut-offs attached to an island.} Such a cut-off is, in addition, read inside its
  island.
\end{itemize}
The coupling profile of \cite[(2.24)--(2.25)]{B-CMP119} is taken in the following product form.
Fix $\psi_{\varphi}\in C^\infty([0,\infty[;[0,1])$ with $\psi_{\varphi}\equiv1$ on
$[0,\tfrac18]$ and $\psi_{\varphi}\equiv0$ on $[\tfrac14,\infty[$. For $i\le j$ put
\begin{equation}\label{4.38:eq-standing-data-a}
  \varphi_i(x):=\prod_{D}\Bigl(1-\prod_{\mu=1}^{4}
  \psi_{\varphi}\bigl(\operatorname{dist}(x_\mu,D_\mu)/a_i\bigr)\Bigr).
\end{equation}
Here $D$ runs over the level-$i$ cells of $[\Lambda_i(\ell)^c]_i$, and $D_\mu$ is the $\mu$-th
coordinate interval of $D$. The profile \eqref{4.38:eq-standing-data-a} is natural, by part (c)
of the coupling-profile lemma. It is not,
in general, a function of the sup-distances to the faces of $\Lambda_i(\ell)$: along a flat
wall it varies with the period of the cells. It is part of the setting of the construction. No
number of \textup{(L1)}, that is, no threshold, width, cube size or counter, depends on the choice
of $\psi_{\varphi}$; the localised terms built with $\varphi_i$ do depend on it.

\emph{The profile at an $\mathbf{R}$-step.} At an $\mathbf{R}$-step, $\varphi_i$ ($i\le j$) is
the product profile \eqref{4.38:eq-standing-data-a} of the sequence of large-field regions in
force, with $[\Lambda_i(\ell)^c]_i$ replaced by the round-up $[\,\cdot\,]_i$ of that sequence's
level-$i$ large-field set. The sequences in force are of four kinds.
\begin{enumerate}
\item The sequence made by $\mathbf{T}$, before the $\mathbf{R}$-step.
\item The intermediate sequences of the preparatory integrations, at the stages $n$ with
  $0<n<\mathsf{N}$. Their coupling functions are the $g_k^{(n)}(\cdot)$ of
  \cite[(1.60)--(1.62)]{B-CMP122a}.
\item The prepared sequence $\{Z''_i\}$ of \cite[(1.10)--(1.11)]{B-CMP122a}. Its complements
  form an admissible sequence of domains based on partitions into $M$-cubes in the
  corresponding scales.
\item The sequence of the retained label $\ell$.
\end{enumerate}

\emph{How Ba{\l}aban's coupling function changes at an $\mathbf{R}$-step.} In
\cite[(1.60)--(1.62)]{B-CMP122a} the coupling function is changed at each preparatory
integration. In the notation of \cite[(1.60)]{B-CMP122a}, whose index $j$ is Ba{\l}aban's
and is not the level $j$ of this chapter, the term
$A\bigl(1/(g_k^{(n)}(\cdot))^2-1/g_j^2,\,\cdot\,\bigr)$ appears twice in
\cite[(1.60)]{B-CMP122a}; here $A(x(\cdot),U)$ denotes Ba{\l}aban's Wilson action of $U$ with a
position-dependent weight $x(\cdot)$ in place of the constant inverse coupling
(Definition~\ref{1.5:def-domains}). By \cite[(1.62)]{B-CMP122a},
\begin{equation}\label{4.38:eq-standing-data-2}
\begin{split}
  \mathbf{C}_k^{(n+1)} &= \log\bigl(\text{the integral in \cite[(1.60)]{B-CMP122a}}\bigr)
  + A\Bigl(\frac{1}{(g_k^{(n+1)}(\cdot))^2}-\frac{1}{(g_k^{(n)}(\cdot))^2},\,
  U_k^{(n+1)}\Bigr),\\
  \mathbf{B}_k^{(n+1)} &= \mathbf{B}_k^{(n)}+\mathbf{C}_k^{(n+1)} .
\end{split}
\end{equation}
Here the first term of $\mathbf{C}_k^{(n+1)}$ is the main contribution, determined by the
integral in \cite[(1.60)]{B-CMP122a}, and the second comes from the renormalisation of the
Wilson action \cite[p.~189]{B-CMP122a}.
The superscript $(\mathsf{N})$ is then replaced by a double prime \cite{B-CMP122a}, and
$g''_k(\cdot):=g_k^{(\mathsf{N})}(\cdot)$ is the \emph{prepared} coupling function. Its
excess over the constant is the second term of \cite[(1.14)]{B-CMP122b}:
\begin{equation*}
  A\bigl(1/(g''_k(\cdot))^2-1/g_k^2,\,U''_k\bigr).
\end{equation*}
This term is localised, as a function of $U''_k$, in the box $\Lambda^{(1.73)}$
\cite[p.~359]{B-CMP122b}. The box $\Lambda^{(1.73)}$ of \cite[(1.73)]{B-CMP122a} is obtained by
adding one layer of $M$-cubes to $Z''_k$; it is a union of $M$-cubes and, by condition~(i)
there, a rectangular parallelepiped contained in a cube of size $100M$
\cite[p.~192]{B-CMP122a}. It lies inside $Z$: in \cite[p.~177]{B-CMP122a} every region is
intersected with $Z$, and $Z_k$ is identified with $Z$. The new
function of \cite[p.~365]{B-CMP122b} is
\begin{equation}\label{4.38:eq-standing-data-3}
  \frac{1}{(g_k(\cdot))^2}=\frac{1}{(g''_k(\cdot))^2}\,Z^c+\frac{1}{g_k^2}\,Z ,
\end{equation}
in which a set and its characteristic function are denoted alike.

\emph{Consequences at islands and at the regions $Y_b$ of \cite[p.~390]{B-CMP122b}.}
\begin{itemize}
\item By \eqref{4.38:eq-standing-data-3}, the new function equals $1/g_k^2$ on $Z$. Since the
  excess of $1/(g''_k(\cdot))^2$ over $1/g_k^2$ is localised in $\Lambda^{(1.73)}\subset Z$,
  the new function equals $1/g_k^2$ also around $Z$.
\item At a healed island, that is, an island whose large-field factor has been absorbed back
  into the small-field format by a later step, the island being treated from then on in the
  regular expansion, the new function is the position-dependent coupling function of
  \cite[(2.24)--(2.25)]{B-CMP119} built with the profile \eqref{4.38:eq-standing-data-a} of
  the retained label, which contains no island factor. This is the form in which Ba{\l}aban's
  functional $A'_k\bigl(1/(g_k(\cdot))^2,U_k\bigr)$ of \cite[(1.72)]{B-CMP122b} stands outside
  the sum over islands in \cite[p.~379]{B-CMP122b}.
\item At the regions $Y_b\subset Z$ adjoined to $Z_k$ to form the new large-field region
  $Z_k\cup\bigcup_{b}Y_b$ of \cite[p.~390]{B-CMP122b}, the coupling function used here, built
  with $\varphi_j$ of the label $\ell$, differs from \eqref{4.38:eq-standing-data-3}. Since
  $Y_b$ is part of the new large-field region, $Y_b\subset\Lambda_j(\ell)^c$, and
  $\varphi_j$ of the label $\ell$ vanishes on $Y_b$ and near it. The function
  \eqref{4.38:eq-standing-data-3}, by contrast, has the value $1/g_k^2$ there.
\item The coupling function and the counterterms are written with one and the same profile.
  Hence, by the family's representation proposition for the coupling function and the
  counterterms, $A'_k$ is the same function under both choices, and the difference is one of
  representation only.
\end{itemize}

\smallskip
\noindent\textup{(L3)} \emph{Rules.} This concerns every rule that makes or changes a label:
the rules \cite[(3.2)--(3.5), (3.20)]{B-CMP119}, the family's form of
\cite[(3.16)]{B-CMP119}, the one-cell enlargement $\sim$ (one layer of cells; see
\textup{(L4)} below), the parallelepiped rule, and the two rules by which the family forms
labels. Such a rule reads
only two things. First, it reads the cubes, bonds and fields of the component $W$ it acts on,
together with its round-up $[W]_j$ and the level-$j$ cells touching $[W]_j$. Second, it reads
numbers of \textup{(L1)}.

\smallskip
\noindent\textup{(L4)} \emph{Vacuum convention, components, strips and the rule.}

\emph{Vacuum convention.} The family's convention on the vacuum terms of the step, fixed with
its rules at the split $\mathbb{B}$, is in force.

\emph{Components.} The components of $\Lambda_k^c$ are the classes of the equivalence relation
generated by the relation \emph{intersect, or touch} \cite{B-CMP122b}, contact at corners
included.

\emph{The strip.} In \cite[(1.68)--(1.72), (1.90)--(1.100)]{B-CMP122b} the one-cell
enlargements $Z^{\sim}$, $X^{\sim}$ (one layer of cells) are written. At those places we use
instead the strip
\begin{equation}\label{4.38:eq-standing-data-b}
  Z^{\sharp}:=Z\ \text{plus } m_\sharp \text{ layers of } M\text{-cubes},\qquad
  m_\sharp:=\bigl\lfloor \tfrac12(\check{R}_k-1)\bigr\rfloor .
\end{equation}
The first layer consists of the $M$-cubes touching $Z$. For $1\le n\le m_\sharp$, the $n$-th
layer consists of the $M$-cubes at sup-distance $(n-1)M$ from $Z$. The strips $X^{\sharp}$,
$X_a^{\sharp}$, $Y_b^{\sharp}$, and in general $W^{\sharp}$ for a union $W$ of $M$-cubes, are
formed in the same way. At \cite[(2.30), (2.41)]{B-CMP119}
and at one further passage of \cite{B-CMP122b}, the operation $\sim$ keeps its meaning,
namely enlargement by one layer of the cells in terms of which the enlarged set is defined
($M\check{R}$-cubes of its level).

The strip convention presupposes $m_\sharp\ge1$. The following properties hold.
\begin{itemize}
\item[(s1)] $m_\sharp\ge1$ if and only if $\check{R}_k\ge3$. This floor is implied by
  $\check{R}_k\ge L$, which holds by the lemma on the cell-size multiplier, $L$ being an odd
  integer greater than one.
\item[(s2)] $m_\sharp M<\tfrac12 M\check{R}_k$, and $m_\sharp<\check{R}_k$.
\item[(s3)] Let $Z$ be a union of $M$-cubes of the standard partition. Then every point of
  $Z^{\sharp}$ lies at sup-distance at most $m_\sharp M$ from $Z$.
\item[(s4)] Let $V,W$ be unions of $M$-cubes at sup-distance at least $\check{R}_kM$. Then
  $V^{\sharp}$ and $W^{\sharp}$ are at sup-distance at least
  \begin{equation*}
    (\check{R}_k-2m_\sharp)M\ \ge\ M .
  \end{equation*}
\item[(s5)] If $Z$ is a union of level-$k$ cells, then every point of $Z^{\sharp}$ lies in a
  level-$k$ cell touching $Z$, that is, in the one-cell enlargement $Z^{\sim}$.
\item[(s6)] Let $Z$ be a union of components of $\Lambda_k^c$ at the $\mathbf{R}$-step. Then
  the strips $C^{\sharp}$ of the components $C$ of $Z$ are boxes, they
  are pairwise at
  sup-distance at least $(\check{R}_k-2m_\sharp)M$, they are the components of $Z^{\sharp}$, and
  $C^{\sharp}\subset C^{\sim}$. This is the strip theorem together with its corollary.
\end{itemize}

\emph{The rule.} Let $\mathbb{L}$ be the list of operators of
Theorem~\ref{4.4:thm-uniform-agreement}. For an operator $K\in\mathbb{L}$
(Definition~\ref{4.4:def-walk-systems}) the rule
$\mathcal{R}$ writes $K=\sum_\omega K_\omega$, and has three properties.
\begin{itemize}
\item It is exact at interpolation parameters identically equal to one
  (Proposition~\ref{4.4:prop-rule-properties}).
\item It is local in the sense of Definition~\ref{4.4:def-local-rules}. Here the footprint of a
  walk $\omega$ (Definition~\ref{4.4:def-walk-systems}) is the set read by its term. For a
  localised term with enlarged cover sets $\tilde X_i^{5}$ it is their union
  $\bigcup_i\tilde X_i^{5}$, and for a term attached to a union $W$ of $M$-cubes it is the
  strip $W^{\sharp}$ of \eqref{4.38:eq-standing-data-b}.
  For the terms of $\mathsf{A}$ this is the support $[\omega]$ itself, with reading radius $0$.
  For
  the walk terms of the operators of the class admitted in
  Definition~\ref{4.4:def-expansion-rule} other than $\mathsf{A}$ it is the
  $M/4$-neighbourhood of $[\omega]$. The term
  indexed by $X$ collects the walks whose footprints have hull $X$. A rule that takes instead
  the hull of $[\omega]$ reads $X$ with a collar of width $M/4$.
\item It is covariant under the translations that preserve the partitions it reads.
\end{itemize}

\smallskip
\noindent\textup{(L5)} \emph{Interior.} Let $\Lambda_t$ denote the family's small-field region
at the step. Consider the free bonds of the step, the bond sets of the step being those of
Notation~\ref{4.4:not-units-distances}, with both end-points in
$\Lambda_t$. On these bonds the integrand of the step carries Ba{\l}aban's characteristic
function $\chi_b$ at the threshold $\delta_k$ \cite{B-CMP119}, and nothing else.

\smallskip
\emph{Scope.} Covariance under reflections and permutations of the coordinate axes is not
among \textup{(L0)}--\textup{(L5)}. It is read only at a single step of the proof of the
torus-independence lemma of clause~(0), where it is required; there it is supplied by the
reflection covariance of
the walk expansion of $\mathsf{A}$ and, for the operators of $\mathbb{L}_1$ other than
$\mathsf{A}$, by item (f3) of Definition~\ref{4.4:def-one-family-hypothesis}.

The proof of the torus-independence lemma of clause~(0) uses the family, the rule and the split
only through \textup{(L0)}--\textup{(L5)}; any family, rule and split satisfying
\textup{(L0)}--\textup{(L5)} serves, where \textup{(L0)}--\textup{(L5)} include the strip
convention \eqref{4.38:eq-standing-data-b}, the profile \eqref{4.38:eq-standing-data-a} with its
reading at an $\mathbf{R}$-step, and the clause of \textup{(L2)} on cut-offs attached to a label.

For Ba{\l}aban's construction with the family, rule and split fixed in this chapter:
\begin{itemize}
\item \textup{(L0)}--\textup{(L2)} are part of its setting. This includes the profile
  \eqref{4.38:eq-standing-data-a} and its reading at an $\mathbf{R}$-step.
\item The locality statements of \textup{(L4)} are supplied by parts (a) and (b) of
  Corollary~\ref{4.4:cor-localised-independence} on each torus. For two tori of different
  size, the localised linear objects agree on their common domain by Ba{\l}aban's definitions
  \cite[p.~394 and Corollary~3.8, p.~410]{B-CMP99b}, the convergence of the expansions on the
  joined localisation domains being cited as printed from \cite[(1.70)]{B-CMP122a} and
  \cite[(1.68)--(1.69)]{B-CMP122b}; for the marked terms the agreement follows from the
  construction of the terminal step. For the terms of
  $\mathsf{A}$ they are supplied by Lemma~\ref{4.4:lem-walk-expansion-locality}, on local data
  given on the walk's support, and by its two corollaries. The strip properties are supplied
  by the strip theorem and its corollary.
\item Covariance under translations holds by transport of structure.
\item \textup{(L3)} and \textup{(L5)} hold by the lemma on the family's small-field and
  large-field functions, \textup{(L5)} by its part~(b).
\end{itemize}
\end{definition}

\begin{proof}[Proof of \textup{(s1)}--\textup{(s5)}]
\emph{(s1).} For a real number $y$, $\lfloor y\rfloor\ge1$ holds if and only if $y\ge1$. With
$y=\tfrac12(\check{R}_k-1)$ this reads $\check{R}_k\ge3$. Since $L$ is an odd integer greater
than one, $L\ge3$. Hence $\check{R}_k\ge L$, which holds by the lemma on the cell-size
multiplier, gives $\check{R}_k\ge3$.

\emph{(s2).} By definition of the integer part,
\begin{equation*}
  m_\sharp\le\tfrac12(\check{R}_k-1)<\tfrac12\check{R}_k .
\end{equation*}
Multiplying by $M$ gives $m_\sharp M<\tfrac12M\check{R}_k$. Since $\check{R}_k>0$, also
$m_\sharp<\tfrac12\check{R}_k<\check{R}_k$.

\emph{(s3).} Let $x\in Z^{\sharp}\setminus Z$. Then $x$ lies in an $M$-cube $c$ of the $n$-th
layer for some $1\le n\le m_\sharp$, so $\operatorname{dist}_\infty(c,Z)\le(n-1)M$. Choose
$y\in c$ with
\begin{equation*}
  \operatorname{dist}_\infty(y,Z)=\operatorname{dist}_\infty(c,Z).
\end{equation*}
The cube $c$ has sup-diameter $M$, so $\operatorname{dist}_\infty(x,y)\le M$. Hence
\begin{equation*}
  \operatorname{dist}_\infty(x,Z)\le(n-1)M+M=nM\le m_\sharp M .
\end{equation*}

\emph{(s4).} Let $x\in V^{\sharp}$ and $y\in W^{\sharp}$. By (s3) there are $x'\in V$ and
$y'\in W$ with $\operatorname{dist}_\infty(x,x')\le m_\sharp M$ and
$\operatorname{dist}_\infty(y,y')\le m_\sharp M$. By the triangle inequality,
\begin{equation*}
  \operatorname{dist}_\infty(x,y)\ \ge\ \operatorname{dist}_\infty(x',y')-2m_\sharp M
  \ \ge\ (\check{R}_k-2m_\sharp)M .
\end{equation*}
From $m_\sharp\le\tfrac12(\check{R}_k-1)$ we get $2m_\sharp\le\check{R}_k-1$, that is,
$\check{R}_k-2m_\sharp\ge1$.

\emph{(s5).} Write $a_k=M\check{R}_k$ for the side of a level-$k$ cell
(\eqref{4.38:eq-standing-data-1} with $k=j$). Let
$x\in Z^{\sharp}$. By (s3) and (s2), $\operatorname{dist}_\infty(x,Z)\le m_\sharp M<a_k$. Since
$Z$ is a finite union of closed cells, there is a level-$k$ cell $D\subset Z$ with
$\operatorname{dist}_\infty(x,D)<a_k$. Let $D'$ be a level-$k$ cell containing $x$. Fix a
coordinate $\mu$.
\begin{itemize}
\item $D_\mu$ and $D'_\mu$ are closed intervals of length $a_k$ of the grid $a_k\mathbb{Z}$.
\item Since the sup-distance is the maximum of the coordinate distances,
  $\operatorname{dist}(x_\mu,D_\mu)\le\operatorname{dist}_\infty(x,D)<a_k$.
\item Suppose $D_\mu\cap D'_\mu=\emptyset$. Then a whole grid interval of length $a_k$ lies
  between them, and $x_\mu\in D'_\mu$ would give $\operatorname{dist}(x_\mu,D_\mu)\ge a_k$,
  which is impossible.
\end{itemize}
Hence $D_\mu\cap D'_\mu\ne\emptyset$ for every $\mu$. Therefore $D\cap D'\ne\emptyset$, that is,
$D'$ touches $Z$.
\end{proof}

\begin{definition}[Sites, tori in level units, and admissible sites]\label{4.38:def-admissible-sites}
The numbers, partitions and cut-offs of the construction are those of
Definition~\ref{4.38:def-standing-data}.

\emph{Sites.} A \emph{site} is a pair $\mathfrak{s}=(m,K')$. It designates Wilson's action
\cite[(0.1)]{B-CMP109} on the torus of side $2L^{m}$ with lattice spacing $\varepsilon=L^{-K'}$
(Ba{\l}aban's $\varepsilon$, not to be confused with the thresholds $\varepsilon_k$ of
Definition~\ref{4.38:def-standing-data}).
For $j\le K'$ we measure lengths in \emph{level-$j$ units}, whose unit is the physical length
$L^{j-K'}$. In these units the torus of $\mathfrak{s}$ is
$\mathbb{T}_j(\mathfrak{s})=\mathbb{R}^4/s_j(\mathfrak{s})\mathbb{Z}^4$, where
\begin{equation}\label{4.38:eq-admissible-sites-a}
  s_j(\mathfrak{s}) = 2L^{m+K'-j},
\end{equation}
and $p_{\mathfrak{s}}\colon\mathbb{R}^4\to\mathbb{T}_j(\mathfrak{s})$ denotes the covering map.
By the geometry clause of Definition~\ref{4.38:def-standing-data}, every lattice and every
partition used on $\mathbb{T}_j(\mathfrak{s})$ with side $a$ dividing
$\tfrac12 s_j(\mathfrak{s})$ is the image under $p_{\mathfrak{s}}$ of the standard one on
$\mathbb{R}^4$, whose corners lie in $a\mathbb{Z}^4$.

Let $E$ denote the constant which, as described below, is fixed only at the end of the
construction. During the induction $E$ is set equal to $0$ at every site. This is consistent
because, by part~(iii) of the lemma of this chapter on the large-field operations, $\mathbf{T}$ is
linear, $\mathbf{R}$ is homogeneous of degree one, and no term of the
construction contains $E$. The constant $E$ is fixed afterwards, and only at the sites $(m,K)$ at
which Theorem~0 is stated. The condition that fixes $E$ is not one of the inductive hypotheses at
any other site.

\emph{Admissible sites.} Let $j\ge 0$, and suppose that cell-size multipliers
$\check{R}_0,\dots,\check{R}_j$ of levels $0,\dots,j$ are given which do not depend on the site.
The level-$i$ cells have side $M\check{R}_i$ in level-$i$ units. Put
\begin{equation}\label{4.38:eq-admissible-sites-b}
\begin{gathered}
  \mathcal{S}_j = \bigl\{\mathfrak{s}=(m,K') :\ K'\ge j,\ 
     L^{m+K'-i}\ge L^{11}M\check{R}_i \text{ for } 0\le i\le j\bigr\},\\
  \mathcal{S}_j^{\flat} = \bigl\{\mathfrak{s}=(m,K')\in\mathcal{S}_{j-1} :\ K'\ge j\bigr\}
  \qquad (j\ge 1),
\end{gathered}
\end{equation}
and call the elements of $\mathcal{S}_j$ \emph{admissible sites at level $j$}. Put also
$\mathcal{A}_j=\{s_j(\mathfrak{s}):\mathfrak{s}\in\mathcal{S}_j^{\flat}\}$.

These objects have the following properties.
\begin{enumerate}
\item[(i)] For $j\ge1$ we have $\mathcal{S}_j\subset\mathcal{S}_j^{\flat}$. The sets
  $\mathcal{S}_j$ and $\mathcal{S}_j^{\flat}$ are site-free: their defining conditions involve
  only $j$, $L$, $M$ and $\check{R}_0,\dots,\check{R}_j$, not the site. They are also cofinal,
  with $m$ free upward: if
  $(m,K')$ belongs to one of them, so does $(m'',K')$ for every $m''\ge m$, and for every
  $K'\ge j$ some $(m,K')$ belongs to it. In particular $\mathcal{A}_j$ is unbounded. Any cofinal
  family of torus sides, that is, one unbounded above, serves in
  Theorem~\ref{4.1:thm-clause-zero}, and $\mathcal{A}_j$ is such a family.
\item[(ii)] Let $K=K(g_0,g)$ be the first-passage level of
  Definition~\ref{1.2:def-flow-constants}, under the hypotheses \textup{(H1)}--\textup{(H7)} of
  Notation~\ref{2.2:not-hypotheses}, and let
  $m\ge m_{\mathbb{T}}(g_-(g))$. The sites $(m,K)$ of Theorem~0 lie in
  $\bigcap_{j\le K}\mathcal{S}_j$.
\item[(iii)] For $\mathfrak{s}\in\mathcal{S}_j^{\flat}$,
  \begin{equation}\label{4.38:eq-admissible-sites-1}
    \tfrac12 s_j(\mathfrak{s})\ \ge\ L^{10}M\check{R}_{j-1}\ \ge\ L^{10}M .
  \end{equation}
  Hence the level-$j$ $M$-cubes fit in $\mathbb{T}_j(\mathfrak{s})$ although $\check{R}_j$ is
  not yet known. This is why the coupling increment $\beta_j$ is read off on
  $\mathcal{S}_j^{\flat}$.
\item[(iv)] The pair $(m,K')$ enters the steps of levels $\le j$ only through
  $s_j(\mathfrak{s})$.
\end{enumerate}
\end{definition}

\begin{proof}
(i) Let $\mathfrak{s}=(m,K')\in\mathcal{S}_j$. Then $K'\ge j\ge j-1$. The inequalities
$L^{m+K'-i}\ge L^{11}M\check{R}_i$ for $0\le i\le j-1$ are among those that define
$\mathcal{S}_j$. Hence $\mathfrak{s}\in\mathcal{S}_{j-1}$, and together with $K'\ge j$ this gives
$\mathfrak{s}\in\mathcal{S}_j^{\flat}$.

The conditions in \eqref{4.38:eq-admissible-sites-b} involve only $m$, $K'$, $L$, $M$ and the
site-free numbers $\check{R}_i$, so both sets are site-free.

Each defining inequality is preserved when $m$ is increased, since $L>1$. For fixed $K'\ge j$,
the $i$-th inequality holds as soon as
\[
  m \ge 11 + m' + \log_L\check{R}_i - K' + i ,
\]
because $M=L^{m'}$. There are finitely many $i$, so every sufficiently large $m$ satisfies all of
them. The same argument applies to $\mathcal{S}_j^{\flat}$, whose conditions are those of
$\mathcal{S}_{j-1}$ together with $K'\ge j$.

Finally, $s_j(m,K')=2L^{m+K'-j}$ is unbounded as $m\to\infty$ with $K'$ fixed. Hence
$\mathcal{A}_j$ is unbounded.

(ii) By definition $m_{\mathbb{T}}(g')=11+m'+\log_L R(g')$. The hypothesis
$m\ge m_{\mathbb{T}}(g_-(g))$ therefore reads $L^{m}\ge L^{11}M\,R(g_-(g))$. Thus the defining
inequality of $\mathcal{S}_j$ at each level $i$ is the inequality that defines the torus exponent
$m_{\mathbb{T}}$, read at level $i$.

Two bounds on the cell-size multipliers are used, both taken by statement:
\[
  \check{R}_K\le R(g_-(g)), \qquad \check{R}_i\le L^{K-i}\check{R}_K \quad (0\le i\le K).
\]
The first holds for $K=K(g_0,g)$ under the hypotheses \textup{(H1)}--\textup{(H7)} of
Notation~\ref{2.2:not-hypotheses}, by the comparison of the last cell-size multiplier with the
radius function $R$; the second is part~(b) of the lemma on the cell-size multipliers.
Let $j\le K$ and $0\le i\le j$. Then
\begin{align*}
  L^{m+K-i} = L^{K-i}L^{m}
  &\ge L^{11}M\,L^{K-i}R(g_-(g))\\
  &\ge L^{11}M\,L^{K-i}\check{R}_K \ \ge\ L^{11}M\check{R}_i .
\end{align*}
Since also $K\ge j$, this shows $(m,K)\in\mathcal{S}_j$ for every $j\le K$.

(iii) Let $\mathfrak{s}=(m,K')\in\mathcal{S}_j^{\flat}$. Then $\mathfrak{s}\in\mathcal{S}_{j-1}$.
The defining inequality of $\mathcal{S}_{j-1}$ at $i=j-1$ gives
$L^{m+K'-j+1}\ge L^{11}M\check{R}_{j-1}$. Dividing by $L$ and using
\eqref{4.38:eq-admissible-sites-a}, we get
\[
  \tfrac12 s_j(\mathfrak{s}) = L^{m+K'-j} \ge L^{10}M\check{R}_{j-1}.
\]
The second inequality in \eqref{4.38:eq-admissible-sites-1} follows from $\check{R}_{j-1}\ge 1$,
which holds because $\check{R}_{j-1}\ge L$ by the lemma on the cell-size multipliers.

The side of the torus is therefore at least $2L^{10}M$ in level-$j$ units. This bound uses
$\check{R}_{j-1}$ only, not $\check{R}_j$.

(iv) In dimension four the bare action carries no dependence on the spacing $\varepsilon$:
Ba{\l}aban writes ``For $d = 4$ we drop the superscript $\varepsilon$'' \cite{B-CMP109}. By the
numbers clause of Definition~\ref{4.38:def-standing-data}, all parameters of the steps are
functions of $\mathfrak{p}$ and of the couplings already generated.

For $i\le j$ the level-$i$ torus of $\mathfrak{s}$ has side
$s_i(\mathfrak{s})=L^{j-i}s_j(\mathfrak{s})$, by \eqref{4.38:eq-admissible-sites-a}. So the tori
on which the steps of levels $\le j$ act, together with their lattices and partitions (images
under $p_{\mathfrak{s}}$ of the standard ones), are determined by $s_j(\mathfrak{s})$. Hence
$(m,K')$ enters those steps only through $s_j(\mathfrak{s})$.
\end{proof}

\begin{lemma}[Rounding up to coarser cubes, then touching]\label{4.38:lem-round-up-touch}
Let $c_{\mathsf{K}}>0$, and let $\mathsf{K}$ be the partition of $\mathbb{R}^4$ into the closed cubes
$c_{\mathsf{K}}(n+[0,1]^4)$, $n\in\mathbb{Z}^4$. Let $c_{\mathsf{K}'},c_{\mathsf{k}}>0$, and let
$\mathsf{K}'$ and $\mathsf{k}$ be the partitions of $\mathbb{R}^4$ into the closed cubes
$c_{\mathsf{K}'}(n+[0,1]^4)$ and $c_{\mathsf{k}}(n+[0,1]^4)$,
$n\in\mathbb{Z}^4$. Assume that they are compatible with $\mathsf{K}$ and with each other, with
$\mathsf{K}'$ equal to or finer than $\mathsf{K}$ and $\mathsf{k}$ finer than $\mathsf{K}'$; that is,
$c_{\mathsf{K}}/c_{\mathsf{K}'}\in\{1,2,\dots\}$ and $c_{\mathsf{K}'}/c_{\mathsf{k}}\in\{2,3,\dots\}$.

For $\mathsf{P}\in\{\mathsf{K},\mathsf{K}',\mathsf{k}\}$ the cubes of $\mathsf{P}$ are called
$\mathsf{P}$-cubes. Two cubes \emph{touch} if they share a point. For a union $X$ of
$\mathsf{P}$-cubes, the $\mathsf{P}$-cubes of $X$ are the $\mathsf{P}$-cubes contained in $X$; two
such unions \emph{share no $\mathsf{P}$-cube} if no $\mathsf{P}$-cube is a $\mathsf{P}$-cube of both.
For $\mathsf{P}\in\{\mathsf{K},\mathsf{K}'\}$ and a union $W$ of $\mathsf{k}$-cubes put
\begin{itemize}
\item[] $[W]_{\mathsf{P}}$ is the union of the $\mathsf{P}$-cubes containing a $\mathsf{k}$-cube of $W$;
\item[] $[W]^{+1}_{\mathsf{P}}$ is the union of $[W]_{\mathsf{P}}$ and the $\mathsf{P}$-cubes touching
$[W]_{\mathsf{P}}$.
\end{itemize}
We write $[W]=[W]_{\mathsf{K}}$ and $W^{+1}=[W]^{+1}_{\mathsf{K}}$.
Let $Z$ be a union of $\mathsf{K}$-cubes, and let $Z^{\sim}$ be the union of $Z$ and the
$\mathsf{K}$-cubes touching $Z$. Let $W'\subset W$ be a further union of $\mathsf{k}$-cubes. Write
$|y|_\infty=\max_\mu|y_\mu|$, let $\operatorname{dist}_{\infty}$ be the sup-distance between subsets of
$\mathbb{R}^4$, and let $\tau_t A=A+t$ for $t\in\mathbb{R}^4$. Then:
\begin{enumerate}
\item[(a)] The following three statements are equivalent: $W$ and $Z^{\sim}$ share no
$\mathsf{k}$-cube; $[W]$ and $Z^{\sim}$ share no $\mathsf{K}$-cube; $W^{+1}$ and $Z$ share no
$\mathsf{K}$-cube.
\item[(b)] A $\mathsf{K}$-cube at sup-distance $<c_{\mathsf{K}}$ from a $\mathsf{k}$-cube of $W$ lies in
$W^{+1}$.
\item[(c)] $[W']^{+1}_{\mathsf{K}'}\subset[W]^{+1}_{\mathsf{K}}$.
\item[(d)] Every point of $W^{+1}$ lies within sup-distance $2c_{\mathsf{K}}$ of $W$; and
$\tau_t[W]=[\tau_t W]$ for $t\in c_{\mathsf{K}}\mathbb{Z}^4$.
\item[(e)] If $W^{+1}$ and $Z$ share no $\mathsf{K}$-cube, then
$\operatorname{dist}_{\infty}(x,Z)\ge c_{\mathsf{K}}$ for every $x\in W$.
\end{enumerate}
\end{lemma}

\begin{proof}
We first record three elementary facts, for $\mathsf{P}\in\{\mathsf{K},\mathsf{K}',\mathsf{k}\}$, whose
cubes have side $c_{\mathsf{P}}$.

First, the sup-distance between two products of compact intervals is the maximum over $\mu$ of the
distances between the $\mu$-th factors, and the intervals $[n_\mu,n_\mu+1]$, $[n'_\mu,n'_\mu+1]$ are at
distance $\max(|n_\mu-n'_\mu|-1,0)$. Hence
\begin{equation}\label{4.38:eq-round-up-touch-1}
\operatorname{dist}_{\infty}\bigl(c_{\mathsf{P}}(n+[0,1]^4),\,c_{\mathsf{P}}(n'+[0,1]^4)\bigr)
= c_{\mathsf{P}}\,\max\bigl(|n-n'|_\infty-1,\,0\bigr),
\end{equation}
a multiple of $c_{\mathsf{P}}$; in particular two $\mathsf{P}$-cubes are equal or touch if and only if
$|n-n'|_\infty\le1$.

Second, distinct $\mathsf{P}$-cubes meet only in boundary points, so the interior of a
$\mathsf{P}$-cube meets no other $\mathsf{P}$-cube. Consequently, if a cube $D$ lies in a
$\mathsf{P}$-cube $C$ and in a union $X$ of $\mathsf{P}$-cubes, then the centre of $D$, an interior point
of $C$, lies in a $\mathsf{P}$-cube of $X$, which must be $C$; so $C$ is a $\mathsf{P}$-cube of $X$. Taking
$D=C$: a $\mathsf{P}$-cube contained in a union of $\mathsf{P}$-cubes is one of them.

Third, let $q=c_{\mathsf{K}}/c_{\mathsf{k}}\in\mathbb{Z}$. Writing the index of a $\mathsf{k}$-cube as
$qn+\nu$ with $n\in\mathbb{Z}^4$ and $\nu\in\{0,\dots,q-1\}^4$, we get
$[qn_\mu+\nu_\mu,qn_\mu+\nu_\mu+1]\subset[qn_\mu,qn_\mu+q]$, so
$c_{\mathsf{k}}(qn+\nu+[0,1]^4)\subset c_{\mathsf{K}}(n+[0,1]^4)$. By the second fact this $\mathsf{K}$-cube is
unique. Thus a $\mathsf{k}$-cube lies in exactly one $\mathsf{K}$-cube, its $\mathsf{K}$-cube; likewise
for $\mathsf{k}$ in $\mathsf{K}'$ and $\mathsf{K}'$ in $\mathsf{K}$. Accordingly $[W]$ is the union of
the $\mathsf{K}$-cubes of the $\mathsf{k}$-cubes of $W$. By the second fact, the $\mathsf{K}$-cubes of
$W^{+1}$ are exactly the $\mathsf{K}$-cubes that equal or touch a $\mathsf{K}$-cube of $[W]$, and the
$\mathsf{K}$-cubes of $Z^{\sim}$ are exactly those that equal or touch a $\mathsf{K}$-cube of $Z$.

(a) Suppose a $\mathsf{k}$-cube $c$ of $W$ is contained in $Z^{\sim}$, and let $C$ be its
$\mathsf{K}$-cube. By the second fact, $C$ is a $\mathsf{K}$-cube of $Z^{\sim}$. It is also a
$\mathsf{K}$-cube of $[W]$. Conversely, a $\mathsf{K}$-cube of $[W]$ contained in $Z^{\sim}$ contains a
$\mathsf{k}$-cube of $W$, and that $\mathsf{k}$-cube is then contained in $Z^{\sim}$. This proves the
first equivalence.

For the second, we use that touching is symmetric. Some $\mathsf{K}$-cube $C$ of $[W]$ is a
$\mathsf{K}$-cube of $Z^{\sim}$ if and only if some $\mathsf{K}$-cube $C'$ of $Z$ equals or touches
some $\mathsf{K}$-cube $C$ of $[W]$. This holds if and only if some $\mathsf{K}$-cube $C'$ of $Z$ is a
$\mathsf{K}$-cube of $W^{+1}$.

(b) Let $c$ be a $\mathsf{k}$-cube of $W$, and let $C'$ be a $\mathsf{K}$-cube with
$\operatorname{dist}_{\infty}(c,C')<c_{\mathsf{K}}$. Let $C$ be the $\mathsf{K}$-cube of $c$; it is a
$\mathsf{K}$-cube of $[W]$. Since $c\subset C$,
\begin{equation*}
\operatorname{dist}_{\infty}(C,C')\le\operatorname{dist}_{\infty}(c,C')<c_{\mathsf{K}}.
\end{equation*}
By \eqref{4.38:eq-round-up-touch-1} this distance is a multiple of $c_{\mathsf{K}}$, hence zero, and
$|n-n'|_\infty\le1$ for the indices of $C$ and $C'$. So $C'$ equals or touches $C\subset[W]$, and
therefore lies in $W^{+1}$.

(c) A $\mathsf{K}'$-cube containing a $\mathsf{k}$-cube $c$ of $W'$ lies in a $\mathsf{K}$-cube,
which then contains $c$; hence $[W']_{\mathsf{K}'}\subset[W']_{\mathsf{K}}$. Next, every
$\mathsf{k}$-cube of $W'$ is contained in $W$, so by the second fact it is a $\mathsf{k}$-cube of $W$;
hence $[W']_{\mathsf{K}}\subset[W]_{\mathsf{K}}$. Thus
\begin{equation*}
[W']_{\mathsf{K}'}\subset[W']_{\mathsf{K}}\subset[W]_{\mathsf{K}}\subset[W]^{+1}_{\mathsf{K}}.
\end{equation*}
Now let $D$ be a $\mathsf{K}'$-cube touching $[W']_{\mathsf{K}'}$ at a point $p$. Then
$p\in[W]_{\mathsf{K}}$. The $\mathsf{K}$-cube $C$ containing $D$ contains $p$, so $C$ shares a point with
$[W]_{\mathsf{K}}$. Hence $C$ equals or touches a $\mathsf{K}$-cube of $[W]_{\mathsf{K}}$, and
$D\subset C\subset[W]^{+1}_{\mathsf{K}}$.

(d) Let $y\in W^{+1}$. Then $y$ lies in a $\mathsf{K}$-cube $C'$ that equals or touches a
$\mathsf{K}$-cube $C$ of $[W]$; let $p\in C\cap C'$. The cube $C$ contains a $\mathsf{k}$-cube $c$ of
$W$. Since a $\mathsf{K}$-cube has sup-diameter $c_{\mathsf{K}}$,
\begin{equation*}
|y-p|_\infty\le c_{\mathsf{K}}
\quad\text{and}\quad
\operatorname{dist}_{\infty}(p,c)\le c_{\mathsf{K}},
\end{equation*}
so $\operatorname{dist}_{\infty}(y,W)\le2c_{\mathsf{K}}$. In words: every cube of $[W]$ is within
sup-distance $c_{\mathsf{K}}$ of $W$, and $W^{+1}$ adds one further layer of $\mathsf{K}$-cubes.

For $t\in c_{\mathsf{K}}\mathbb{Z}^4$, the translation $\tau_t$ permutes the $\mathsf{K}$-cubes. Since
$c_{\mathsf{K}}\in c_{\mathsf{k}}\mathbb{Z}$, it also permutes the $\mathsf{k}$-cubes, and it preserves
inclusion. Hence $\tau_tW$ is a union of $\mathsf{k}$-cubes. Moreover, a $\mathsf{K}$-cube $C$ contains a
$\mathsf{k}$-cube of $W$ if and only if $\tau_tC$ contains a $\mathsf{k}$-cube of $\tau_tW$. Thus
$\tau_t[W]=[\tau_tW]$.

(e) Let $x\in W$. Then $x$ lies in a $\mathsf{k}$-cube of $W$, which is contained in its
$\mathsf{K}$-cube $C=c_{\mathsf{K}}(n+[0,1]^4)$, a $\mathsf{K}$-cube of $[W]$. Let $z\in Z$. Then $z$ lies
in a $\mathsf{K}$-cube $C'=c_{\mathsf{K}}(n'+[0,1]^4)$ of $Z$. By hypothesis $C'$ is not a
$\mathsf{K}$-cube of $W^{+1}$, so $C'$ neither equals nor touches $C$. By the first fact,
$|n-n'|_\infty\ge2$, so $|n_\mu-n'_\mu|\ge2$ for some $\mu$. Then the $\mu$-intervals
$c_{\mathsf{K}}[n_\mu,n_\mu+1]$ and $c_{\mathsf{K}}[n'_\mu,n'_\mu+1]$ of $C$ and $C'$ are at least one
full interval apart. Since $x_\mu$ and $z_\mu$ lie in them, $|x_\mu-z_\mu|\ge c_{\mathsf{K}}$, and so
$|x-z|_\infty\ge c_{\mathsf{K}}$. As $z\in Z$ was arbitrary,
$\operatorname{dist}_{\infty}(x,Z)\ge c_{\mathsf{K}}$.
\end{proof}

\begin{definition}[Non-wrapping domains, lifts and the wrapping bound]\label{4.38:def-non-wrapping}
Let $\mathfrak{s}=(m,K')$ be a site with $K'\ge j$. Let $\mathbb{T}_j(\mathfrak{s})=\mathbb{R}^4/s_j(\mathfrak{s})\mathbb{Z}^4$ be the torus of Definition~\ref{4.38:def-admissible-sites}, with side $s_j(\mathfrak{s})$ given by \eqref{4.38:eq-admissible-sites-a} and covering map $p_{\mathfrak{s}}$. We write $s_j$ for $s_j(\mathfrak{s})$. For $t\in\mathbb{R}^4$, $\tau_t$ denotes the translation $x\mapsto x+t$ of $\mathbb{R}^4$ and the map it induces on $\mathbb{T}_j(\mathfrak{s})$. All cubes in $\mathbb{R}^4$ are closed with edges parallel to the axes. A \emph{closed cube of side $a$} in $\mathbb{T}_j(\mathfrak{s})$ is the $p_{\mathfrak{s}}$-image of a closed cube of side $a$ in $\mathbb{R}^4$. Distances are sup-distances $\operatorname{dist}_{\infty}$ unless stated otherwise. Statements about $M$-cubes are made for $\mathfrak{s}\in\mathcal{S}_j^{\flat}$. We write $D_j(\mathfrak{s})$ for the class of localisation domains of \cite{B-CMP109} in $\mathbb{T}_j(\mathfrak{s})$, built from the $M$-cubes of $\mathbb{T}_j(\mathfrak{s})$; at a fixed site this is the class written $\mathbf{D}_j$.
\begin{enumerate}
\item[(a)] A set $X\subset\mathbb{T}_j(\mathfrak{s})$ is \emph{non-wrapping} if it lies in a closed cube of side at most $\tfrac14 s_j$; otherwise it is \emph{wrapping}.
\item[(b)] Let $X$ be non-wrapping and $0\le r\le s_j/16$. Then the set of points at distance at most $r$ from $X$ lies in a closed cube of side less than $\tfrac12 s_j$. The admissible widths include $r=M/4$ on $\mathcal{S}_j^{\flat}$ and $r=L^{11}M\check{R}_j/8$ on $\mathcal{S}_j$. On $\mathcal{S}_j$ they therefore include a width covering the following two sets:
\begin{itemize}
\item $X^{+1}$, which lies within $2M\check{R}_j$ of $X$;
\item $(X^{\sharp})^{+1}$, which lies within $2M\check{R}_j+m_{\sharp}M<3M\check{R}_j$ of $X$.
\end{itemize}
Here $X^{\sharp}$ denotes the strip of $m_{\sharp}$ layers of $M$-cubes around $X$ fixed in \eqref{4.38:eq-standing-data-b}.
\item[(c)] Let $\hat Q\subset\mathbb{R}^4$ be a closed cube of side less than $\tfrac12 s_j$. Then $p_{\mathfrak{s}}$ restricted to $\hat Q$ is injective, and its inverse $p_{\mathfrak{s}}^{-1}$ on $p_{\mathfrak{s}}(\hat Q)$ is called a \emph{lift}. The image $\hat X$ of a set $X\subset p_{\mathfrak{s}}(\hat Q)$ under a lift is also called a lift of $X$, and a point $\hat z$ with $p_{\mathfrak{s}}(\hat z)=z$ is called a lift of $z$. A lift is an isometric homeomorphism. It preserves lattices, partitions, incidences of cubes (touching included), the coordinate axes and tree lengths. Two lifts of the same non-wrapping set differ by a translation $\tau_t$ with $t\in s_j\mathbb{Z}^4$, and $s_j\mathbb{Z}^4\subset M\mathbb{Z}^4$ for $\mathfrak{s}\in\mathcal{S}_j^{\flat}$. For a point $z$ and a point $\hat z$ with $p_{\mathfrak{s}}(\hat z)=z$, the \emph{lift of $X\ni z$ through $\hat z$} is the lift of $X$ that maps $z$ to $\hat z$.
\item[(d)] The quantity $d_j(X)$ is defined in \cite{B-CMP109} through tree graphs contained in $X$. It is intrinsic to $X$: if $\hat X$ is a lift of a non-wrapping $X$, then $d_j(X)=d_j(\hat X)$.
\item[(e)] A sub-domain of a non-wrapping domain is non-wrapping at its own level $i$, where $s_i=L^{j-i}s_j$.
\item[(f)] Every wrapping $X\in D_j(\mathfrak{s})$ satisfies
\begin{equation}\label{4.38:eq-non-wrapping-a}
d_j(X)>\frac{s_j}{4M}-2 .
\end{equation}
\item[(g)] Let $D_j^{\infty}$ denote the localisation domains of \cite{B-CMP109} in $\mathbb{R}^4$, built from the standard $M$-cubes. Let $z$ be a point of the unit lattice and $\hat z$ a lift of $z$. Then the map sending $X$ to the lift of $X$ through $\hat z$ is a bijection of
\[
\{X\in D_j(\mathfrak{s}) :\ X \text{ non-wrapping},\ z\in X\}
\]
onto
\[
\{\hat X\in D_j^{\infty} :\ \hat z\in\hat X,\ \hat X \text{ lies in a closed cube of side}\le\tfrac14 s_j\}.
\]
\item[(h)] On $\mathcal{S}_j$ the following sets are non-wrapping:
\begin{itemize}
\item every first-class component $Z$ in the sense of \cite[p.~175]{B-CMP122a}, which by the same page lies in a cube of side $100M\check{R}_j$;
\item the touching layer of such a component (the cells of side $M\check{R}_j$ touching $Z$) and its strip $Z^{\sharp}$;
\item the cubes $\square^{\sim2}$ of \cite[(3.5)]{B-CMP109}.
\end{itemize}
\end{enumerate}
\end{definition}

\begin{proof}
(b) Let $X$ lie in $p_{\mathfrak{s}}(\hat Q_0)$, where $\hat Q_0$ has side at most $\tfrac14 s_j$. Let $\hat Q$ be the cube obtained by enlarging $\hat Q_0$ by $r$ on every side. Then $p_{\mathfrak{s}}(\hat Q)$ contains every point at distance at most $r$ from $X$. Its side satisfies
\[
\tfrac14 s_j+2r\le\tfrac14 s_j+\tfrac18 s_j=\tfrac38 s_j<\tfrac12 s_j .
\]
We now check the stated widths. On $\mathcal{S}_j^{\flat}$, Definition~\ref{4.38:def-admissible-sites} gives $\tfrac12 s_j\ge L^{10}M$. Hence $s_j/16\ge L^{10}M/8\ge M/4$. On $\mathcal{S}_j$, the defining inequality of Definition~\ref{4.38:def-admissible-sites} with $i=j$ gives
\[
s_j=2L^{m+K'-j}\ge 2L^{11}M\check{R}_j ,
\]
so that $s_j/16\ge L^{11}M\check{R}_j/8$. By Lemma~\ref{4.38:lem-round-up-touch}(d), $X^{+1}$ lies within $2M\check{R}_j$ of $X$. The strip $X^{\sharp}$ lies within $m_{\sharp}M$ of $X$, and, again by Lemma~\ref{4.38:lem-round-up-touch}(d), $(X^{\sharp})^{+1}$ lies within $2M\check{R}_j$ of $X^{\sharp}$. Hence $(X^{\sharp})^{+1}$ lies within $2M\check{R}_j+m_{\sharp}M$ of $X$. The number $m_{\sharp}$ of strip layers fixed in \eqref{4.38:eq-standing-data-b} satisfies $m_{\sharp}<\check{R}_j$, so $2M\check{R}_j+m_{\sharp}M<3M\check{R}_j$. Since $L$ is an odd integer larger than one, $3M\check{R}_j\le L^{11}M\check{R}_j/8$, so both sets are covered.

(c) \emph{Injectivity.} Let $a,b\in\hat Q$ with $p_{\mathfrak{s}}(a)=p_{\mathfrak{s}}(b)$. Then $a-b\in s_j\mathbb{Z}^4$ and $|a_\mu-b_\mu|<\tfrac12 s_j$ for every $\mu$. Hence $a=b$.

\emph{Isometry.} For $a,b\in\hat Q$ and every coordinate $\mu$, the minimum over $n_\mu\in\mathbb{Z}$ of $|a_\mu-b_\mu+s_jn_\mu|$ is attained at $n_\mu=0$, because $|a_\mu-b_\mu|<\tfrac12 s_j$. Therefore the torus distance between $p_{\mathfrak{s}}(a)$ and $p_{\mathfrak{s}}(b)$ equals $|a-b|$, for the sup-distance, the $\ell^1$-distance and the Euclidean distance alike. Thus the lift is an isometry, hence a homeomorphism onto its image. Since $p_{\mathfrak{s}}$ is locally a translation, the lift maps coordinate directions to coordinate directions.

\emph{A uniqueness fact.} Let $A,B\subset\mathbb{R}^4$ each lie in a cube of side less than $\tfrac12 s_j$, share a point, and satisfy $p_{\mathfrak{s}}(A)=p_{\mathfrak{s}}(B)$. Each coordinate range of $A\cup B$ is contained in the union of two intervals of length less than $\tfrac12 s_j$ with a common point, hence in an interval of length less than $s_j$. The argument used for injectivity shows that $p_{\mathfrak{s}}$ is injective on $A\cup B$, so $A=B$.

\emph{Two lifts of the same set.} Let $\hat X,\hat X'$ be two lifts of a non-wrapping $X$, and let $x_0\in X$. Their preimages of $x_0$ differ by some $t\in s_j\mathbb{Z}^4$. The sets $\tau_t\hat X$ and $\hat X'$ then satisfy the hypotheses of the uniqueness fact, so $\tau_t\hat X=\hat X'$. On $\mathcal{S}_j^{\flat}$ we have $\tfrac12 s_j=L^{m+K'-j}\ge L^{10}M$. Since $M=L^{m'}$, it follows that $M$ divides $s_j$, and so $s_j\mathbb{Z}^4\subset M\mathbb{Z}^4$.

\emph{Partitions and lattices.} By Definition~\ref{4.38:def-admissible-sites}, the lattices and partitions of $\mathbb{T}_j(\mathfrak{s})$ of side $a$ dividing $\tfrac12 s_j$ are the $p_{\mathfrak{s}}$-images of the standard ones. Let $C\subset p_{\mathfrak{s}}(\hat Q)$ be a cube of such a partition, so that $C=p_{\mathfrak{s}}(\hat C)$ with $\hat C$ standard. Translate $\hat C$ by an element of $s_j\mathbb{Z}^4$ so that it meets the lift of $C$. The translate is still standard, because $a$ divides $s_j$. By the uniqueness fact it equals the lift of $C$. Lattice points are treated in the same way.

\emph{Incidences.} Let $A,B\subset\hat Q$, let $a\in A$, $b\in B$ with $p_{\mathfrak{s}}(a)=p_{\mathfrak{s}}(b)$, and write $b=a+s_jn$ with $n\in\mathbb{Z}^4$. If $n\ne0$, then $a$ and $a+s_jn$ both lie in $\hat Q$ and differ in the coordinate $\mu$ with $n_\mu\ne0$ by at least $s_j>\tfrac12 s_j$, which is impossible. Hence $n=0$, and
\[
p_{\mathfrak{s}}(A)\cap p_{\mathfrak{s}}(B)=p_{\mathfrak{s}}(A\cap B),
\]
so intersections, and in particular touching, are preserved.

\emph{Tree lengths.} These are preserved by the lemma on lifting tree lengths in a torus, stated in this section for cubes of side at most half the torus side; it applies here because $\hat Q$ has side less than $\tfrac12 s_j$.

(d) A tree contained in $X\subset p_{\mathfrak{s}}(\hat Q)$ is mapped by the lift to a tree contained in $\hat X$, and the projection $p_{\mathfrak{s}}$ maps trees contained in $\hat X$ to trees contained in $X$. Both maps are isometric, hence preserve lengths of curves. The two infima defining $d_j(X)$ and $d_j(\hat X)$ therefore coincide.

(e) Passing from level-$j$ units to level-$i$ units multiplies all lengths by $L^{j-i}$. A cube of side at most $\tfrac14 s_j$ thus becomes a cube of side at most $\tfrac14 L^{j-i}s_j=\tfrac14 s_i$.

(f) Let $\Gamma$ be a shortest tree of $X$. It is contained in $X$, meets every $M$-cube of $X$, and has length $Md_j(X)$. A tree is simply connected, so $\Gamma$ lifts through the covering $p_{\mathfrak{s}}$ to a tree $\hat\Gamma\subset\mathbb{R}^4$ of the same length. Each coordinate projection of the connected set $\hat\Gamma$ is an interval of length at most the length of $\hat\Gamma$. Hence $\hat\Gamma$ lies in a closed cube $\hat Q$ of side $Md_j(X)$.

Let $C$ be an $M$-cube of $X$. It meets $\Gamma$ at a point $p_{\mathfrak{s}}(y)$ with $y\in\hat\Gamma$. Write $C=p_{\mathfrak{s}}(\hat C)$ with $\hat C$ a standard $M$-cube, translated by an element of $s_j\mathbb{Z}^4$ so that $y\in\hat C$. Then $\hat C$ lies within distance $M$ of $y$, hence in the cube $\hat Q'$ obtained by enlarging $\hat Q$ by $M$ on each side. Consequently $X\subset p_{\mathfrak{s}}(\hat Q')$, which is a closed cube of side $Md_j(X)+2M$.

If $Md_j(X)+2M\le\tfrac14 s_j$, then $X$ is non-wrapping. Hence for wrapping $X$ we have $Md_j(X)+2M>\tfrac14 s_j$, which is \eqref{4.38:eq-non-wrapping-a}.

(g) \emph{The map is well defined.} Let $X$ be non-wrapping with $X\subset p_{\mathfrak{s}}(\hat Q_0)$, where $\hat Q_0$ has side at most $\tfrac14 s_j$. Take the lift on $\hat Q_0$, translated by the element of $s_j\mathbb{Z}^4$ that sends the preimage of $z$ to $\hat z$. By (c), the resulting set $\hat X$ does not depend on the choice of $\hat Q_0$. It contains $\hat z$ and lies in a closed cube of side at most $\tfrac14 s_j$. By (c), the lift carries the $M$-cubes of $X$ to standard $M$-cubes and preserves their incidences and tree lengths. Hence $\hat X\in D_j^{\infty}$.

\emph{Injectivity.} If two sets have the same lift $\hat X$, both equal $p_{\mathfrak{s}}(\hat X)$.

\emph{Surjectivity.} Let $\hat X$ belong to the target set, and let $\hat Q$ be a closed cube of side at most $\tfrac14 s_j$ containing it. Then $X:=p_{\mathfrak{s}}(\hat X)$ contains $z$ and is non-wrapping. By (c) applied to $p_{\mathfrak{s}}$ on $\hat Q$, the set $X$ lies in $D_j(\mathfrak{s})$. Since $\hat z\in\hat Q$, the inverse of $p_{\mathfrak{s}}$ on $\hat Q$ maps $z$ to $\hat z$. Its image of $X$ is $\hat X$, so $\hat X$ is the lift of $X$ through $\hat z$.

(h) Let $Z$ be a first-class component, contained in a cube of side $100M\check{R}_j$. Its touching layer consists of the cells of side $M\check{R}_j$ touching $Z$, each of which lies within $M\check{R}_j$ of $Z$, and its strip satisfies $Z^{\sharp}\subset(Z^{\sharp})^{+1}$. As in the proof of (b), Lemma~\ref{4.38:lem-round-up-touch}(d) and the width $m_{\sharp}M$ of the strip show that $Z^{+1}$ and $(Z^{\sharp})^{+1}$ lie within $3M\check{R}_j$ of $Z$; no non-wrapping hypothesis enters this. All these sets therefore lie in a closed cube of side $106M\check{R}_j$. On $\mathcal{S}_j$ we have $\tfrac14 s_j\ge\tfrac12 L^{11}M\check{R}_j$, and $\tfrac12 L^{11}>106$ because $L$ is an odd integer larger than one. Hence all these sets are non-wrapping. A cube $\square^{\sim2}$ of \cite[(3.5)]{B-CMP109} has side at most $\tfrac14 s_j$ on $\mathcal{S}_j$, and by (a) a cube of side at most $\tfrac14 s_j$ is non-wrapping.
\end{proof}

\begin{definition}[Intrinsic objects and the trace of a label]\label{4.38:def-intrinsic-trace}
Fix a level $j$ and a domain $W$. We use the standard lattices, partitions and axes, and the numbers fixed
in Definition~\ref{4.38:def-standing-data}.

\emph{Label cubes, rounding up, touching.} The \emph{label cubes} at the level concerned are the cubes of
the label's partition at that level. At level $j$ these are the $M\check{R}_j$-cubes, called \emph{cells}.
The cells of side $\mathsf{m}_4$ that occur elsewhere in this chapter are a different object and are
always called $\mathsf{m}_4$-cells. Two sets \emph{touch} if they share a point. For a label $\ell$, whose
members are defined below, we put
\begin{equation}\label{4.38:eq-intrinsic-trace-a}
\begin{aligned}
{[W]}&:=\text{the smallest union of label cubes containing } W,\\
W^{+1}&:=[W]\cup\bigl(\text{the label cubes touching } [W]\bigr),\\
\text{the trace of } \ell \text{ on } W^{+1}&:=\text{the family, over the members } A \text{ of } \ell,\\
&\qquad\text{of the cubes of the partition of } A \text{ lying in } W^{+1},\ \text{marks kept}.
\end{aligned}
\end{equation}
Equivalently, $[W]$ is the union of the label cubes that hold a cube of $W$; thus $[W]=W$ when $W$ is a
union of label cubes. At a level $i\le j$ we write $[\,\cdot\,]_i$ and $(\cdot)_i^{+1}$ for the same
operations performed with the label cubes of level $i$, the level-$i$ cells, of side
$a_i=M\check{R}_iL^{-(j-i)}$ in level-$j$ units. The members of a label $\ell$, which enter the
last two lines of \eqref{4.38:eq-intrinsic-trace-a}, are defined next.

\emph{The members of a label.} A \emph{label} $\ell$ at level $j$ is given by the following data, called
its \emph{members}.
\begin{enumerate}
\item The complements $\{\Omega_i(\ell)^c\}$ and $\{\Lambda_i(\ell)^c\}$ of the first two sequences, and
the third sequence $\{S_i\}$, for $i\le j$; in this item we write $\Omega_i=\Omega_i(\ell)$ and
$\Lambda_i=\Lambda_i(\ell)$. These are the sequences of \cite[p.~254]{B-CMP119}. There
$\Omega_i,\Lambda_i\in\mathbf{D}_i$, and the two sequences are nested as in \cite[(2.1)]{B-CMP119}:
\begin{equation*}
\Omega_1\supset\Lambda_1\supset\Omega_2\supset\Lambda_2\supset\dots\supset\Omega_j\supset\Lambda_j .
\end{equation*}
Moreover $S_i\subset\Omega_i\cap\Lambda_i^c$, and each $S_i$ is either empty or a union of
$LM_2R_i$-cubes in the lattice of spacing $L^{-i}$, with $M_2$ and $R_i$ as in
\cite[p.~254]{B-CMP119}.
\item For the family of the construction fixed in Definition~\ref{4.38:def-standing-data}: at each level
$i\le j$ the label carries the datum $(R,I,\mathsf{R}^{\mathrm{lo}})$ of a step of the family at level
$i$, and the final large-field set of that step is $\Lambda_i(\ell)^c$ (for a label retained after the
operation $\mathbf{R}$, the step is that of the label's retained sequence with the healed islands
removed). Here $R$
and $\mathsf{R}^{\mathrm{lo}}$ are sets, $\mathsf{R}^{\mathrm{lo}}$ being a
rejected set of the recursion inside that step, and $I$ is a set of bonds.
\item The sets made by the operation $\mathbf{R}$. These are the histories of the regions $Y_b$ of $\ell$
that $\mathbf{R}$ adds to the large-field region \cite[p.~390]{B-CMP122b}, and the histories
$\mathfrak{h}_a$ of those islands $X_a$ (the components of the large-field region on which the
operation $\mathbf{R}$ acts) that are contained in such a region $Y_b$. Each history is read as a
collection of sets, namely the members of its own older sequences, and these sets lie in the island or
region to which the history belongs, the sets made by $\mathbf{R}$ being formed, in the first line of
\cite[(1.11)]{B-CMP122a}, as intersections with the region on which $\mathbf{R}$ acts.
\item The marks that distinguish, within the large-field region, the cubes of the large-field sets $Z_i$
(the complements of the domains of the second sequence, before the regions added by $\mathbf{R}$ are
joined to them) from those of the regions $Y_b$ added by $\mathbf{R}$.
\end{enumerate}
The sequences $\{S_i\}$ are counted among the members because in \cite{B-CMP119} the summation over
these sequences is part of the operation $\mathbf{T}_k$ at each level $k$ (at the first level, see
\cite[(1.29), p.~254]{B-CMP119}), and the expression $\mathbf{B}_k$ depends also on the sequence
$\{S_i\}$.

\emph{Sets made of $M$-cubes.} At a level $i$ a member $A$ may be a union of $M$-cubes without being a
union of cells. This is the case for the older sets made by $\mathbf{R}$, which in \cite{B-CMP122a} are
based on partitions into $M$-cubes in the corresponding scales. At such a level, ``the cells of $A$''
means the cells of $[A]_i$.

\emph{The trace.} The \emph{trace of $\ell$ on $W^{+1}$} is defined in
\eqref{4.38:eq-intrinsic-trace-a}.
\begin{itemize}
\item For each member $A$ it records the cubes of $A$'s own partition that lie in $W^{+1}$. For the set
$I$ these cubes are the bonds. The marks are kept.
\item The trace is read cube-wise: a cube of a member whose only contact with $W^{+1}$ is along a wall
or at a corner contributes nothing.
\end{itemize}
A trace that records only which cells of $W^{+1}$ lie in the large-field sets does not distinguish two
labels whose data differ only in the bond set $I$, for instance $I=\{b_1\}$ and $I=\{b_2\}$ with
$b_1\neq b_2$. The trace defined here records each member's own cubes, and the bonds for $I$.

\emph{Intrinsic objects.} An object attached to $W$ is \emph{intrinsic} if it is determined by the
following data:
\begin{itemize}
\item the standard lattices, partitions and axes restricted to $W$;
\item the fields on $W$;
\item the gauge group $G$ and the parameters $\mathfrak{p}$;
\item the numbers of Definition~\ref{4.38:def-standing-data} available at the stage concerned;
\item where a label $\ell$ is present, its trace on $W^{+1}$.
\end{itemize}

With these definitions the following hold.
\begin{itemize}
\item[(a)] The trace of $\ell$ on $W^{+1}$ is empty if and only if no cell of $W^{+1}$ is a
cell of $[\Lambda_j(\ell)^c]$.
\item[(b)] An intrinsic object attached to $W$ is carried by every translation $\tau_t$ that preserves
the partitions on which the object depends. If $W$ is non-wrapping in the sense of
Definition~\ref{4.38:def-non-wrapping}, it is also carried by every lift. This is transport of structure.
\item[(c)] Prescriptions that depend on coordinates, such as axial trees, the coordinate-dependent
choice made for each cube $c$, and orderings of the axes, are intrinsic in the above sense.
\end{itemize}
\end{definition}

\begin{proof}
(a) Every member of $\ell$ lies in $\Lambda_j(\ell)^c$. For each kind of member this is seen as
follows; throughout, $i\le j$.
\begin{itemize}
\item For the third sequence, \cite[p.~254]{B-CMP119} gives $S_i\subset\Omega_i(\ell)\cap\Lambda_i(\ell)^c$,
so $S_i\subset\Lambda_i(\ell)^c$.
\item The nesting \cite[(2.1)]{B-CMP119} gives $\Omega_i(\ell)\supset\Lambda_i(\ell)$, hence
$\Omega_i(\ell)^c\subset\Lambda_i(\ell)^c$.
\item The same nesting gives $\Lambda_i(\ell)\supset\Lambda_j(\ell)$ for $i\le j$, hence
$\Lambda_i(\ell)^c\subset\Lambda_j(\ell)^c$.
\end{itemize}
Combining these inclusions,
\begin{equation*}
S_i\subset\Lambda_i(\ell)^c\subset\Lambda_j(\ell)^c,\qquad
\Omega_i(\ell)^c\subset\Lambda_i(\ell)^c\subset\Lambda_j(\ell)^c .
\end{equation*}
For the family of the construction, let $i'\le j$ and consider the datum carried by $\ell$ at level $i'$.
Parts (b) and (d) of the lemma on the family's small-field function $\chi^{\wedge}$ give that $R$,
$\mathsf{R}^{\mathrm{lo}}$ and the set $\mathsf{K}(I)$ attached to $I$ lie in the final large-field set
produced by the recursion inside that step, and that this set is $\Lambda_{i'}^c$, the complement of the
domain of the second sequence at level $i'$ produced by that step. By the description of the datum
carried by $\ell$ among the members, this final large-field set is $\Lambda_{i'}(\ell)^c$; as $i'\le j$,
the nesting gives $\Lambda_{i'}(\ell)^c\subset\Lambda_j(\ell)^c$. So the data carried by $\ell$ at the
levels $i'\le j$ also lie in $\Lambda_j(\ell)^c$, the member $I$ entering this inclusion through the set
$\mathsf{K}(I)$ attached to it. For the sets made
by $\mathbf{R}$, let $Y_b$ be a region of $\ell$ added by $\mathbf{R}$ at a step taken at level
$i'\le j$. By \cite[p.~390]{B-CMP122b} the new large-field region produced by that step contains $Y_b$;
this region is the complement $\Lambda_{i'}(\ell)^c$ of the label's domain at level $i'$, so
$Y_b\subset\Lambda_{i'}(\ell)^c\subset\Lambda_j(\ell)^c$ by the nesting. Every set of the history of
$Y_b$, and of the history of an island $X_a\subset Y_b$, lies in $Y_b$ by the description, among the
members, of the sets made by $\mathbf{R}$ (which rests on \cite[(1.11)]{B-CMP122a}), hence in
$\Lambda_j(\ell)^c$. The marks are not sets: they only label cubes of
members as belonging to a set $Z_i$ or to a region $Y_b$, and contribute no cubes of their own.

Suppose first that some cell $C$ of $W^{+1}$ is a cell of $[\Lambda_j(\ell)^c]$. The set
$\Lambda_j(\ell)^c$ is itself a member of $\ell$ (take $i=j$). Since $C$ is a cell of
$[\Lambda_j(\ell)^c]$, it holds a cube $q$ of the partition belonging to $\Lambda_j(\ell)^c$ itself (a
cell, or an $M$-cube when $\Lambda_j(\ell)^c$ is made of $M$-cubes). Then $q\subset C\subset W^{+1}$, so
$q$ lies in $W^{+1}$ cube-wise, and the trace is non-empty.

Conversely, suppose the trace is non-empty. Then some member has a cube $q$ of its own partition (a
bond, for $I$) lying in $W^{+1}$. A cube contained in the boundary of $W^{+1}$ meets $W^{+1}$ only along
walls or at corners, and by the cube-wise reading it contributes nothing; hence $q$ is not contained in
the boundary of $W^{+1}$, and $q$ contains a point $y$ of the interior of $W^{+1}$. Since every member
lies in $\Lambda_j(\ell)^c$, we have $y\in\Lambda_j(\ell)^c\subset[\Lambda_j(\ell)^c]$, so $y$ lies in a
cell $C$ of $[\Lambda_j(\ell)^c]$. Every neighbourhood of $y$ meets the interior of $C$; choosing a
neighbourhood of $y$ contained in $W^{+1}$, we obtain a non-empty open set $O$ with
$O\subset\operatorname{int}C$ and $O\subset W^{+1}$. The set $W^{+1}$ is a locally finite union of
cells, and the union of their boundaries is a closed set with empty interior; hence $O$ meets the
interior of some cell $C'\subset W^{+1}$. Two distinct cells of the partition have disjoint interiors,
so $C'=C$. Thus $C$ is a cell of $W^{+1}$ and a cell of $[\Lambda_j(\ell)^c]$.

(b) Let $W$ be non-wrapping. By Definition~\ref{4.38:def-non-wrapping}, the set $W^{+1}$ lies inside the
collar of $W$ on which a lift is defined. By the same definition, a lift is an isometric homeomorphism on
that collar which preserves the lattices, the partitions, the incidences (touching included) and the
axes.
\begin{itemize}
\item The lift therefore carries the lattices, partitions and axes restricted to $W$ to those restricted
to the image of $W$ under the lift. It carries the fields on $W$ to the fields there.
\item Since it preserves label cubes and touching, it carries $[W]$ and $W^{+1}$ onto the sets
$[\,\cdot\,]$ and $(\cdot)^{+1}$ of the image of $W$, which are their images under the lift.
\item Since it also preserves the partitions of the members and the marks, it carries the trace of
$\ell$ on $W^{+1}$ onto its image: for each member, the images under the lift of its cubes lying in
$W^{+1}$, with their marks.
\item The gauge group $G$, the parameters $\mathfrak{p}$ and the numbers of
Definition~\ref{4.38:def-standing-data} are not affected.
\end{itemize}
An intrinsic object is determined by these data, so the object attached to the image of $W$ under the
lift is the image of the object attached to $W$.

The same argument, with no condition on $W$, applies to a translation $\tau_t$ that preserves the
partitions on which the object depends. Such a translation preserves lattices and touching, and it does
not move the axes. It therefore maps the restricted lattices, partitions and axes, the fields, the sets
$[W]$ and $W^{+1}$, and the trace on $W^{+1}$ to the corresponding data at $\tau_t W$.

(c) A prescription such as an axial tree, the coordinate-dependent choice made for each cube $c$, or an
ordering of the axes is determined by the axes and the lattice restricted to the domain. Both are among
the data listed in the definition, so such a prescription is intrinsic. By (b) it is carried by every
lift and by every translation preserving the partitions on which it depends.
\end{proof}

\begin{lemma}[Tree lengths on the torus equal lifted tree lengths]\label{4.38:lem-tree-lengths}
Let $s>0$, let $\mathbb{T}=\mathbb{R}^4/s\mathbb{Z}^4$, and let $p\colon\mathbb{R}^4\to\mathbb{T}$ be the
covering map. Let $c=(c_\mu)_{\mu=1}^{4}\in\mathbb{R}^4$, let $a\le s/2$, and put
\begin{equation*}
\hat Q=\prod_{\mu=1}^{4}[c_\mu,c_\mu+a].
\end{equation*}
Let $\hat D_1,\dots,\hat D_n\subset\hat Q$ be closed, and put $D_i:=p(\hat D_i)$. Let $\ell$ denote
length measured in the $\ell^1$-distance of \cite[p.~223]{B-CMP96} on $\mathbb{R}^4$, and on
$\mathbb{T}$ in the quotient distance
\begin{equation*}
\operatorname{dist}(p(x),p(y))=\min_{z\in\mathbb{Z}^4}|x-y+sz|_1,
\end{equation*}
where $|\cdot|_1$ is the $\ell^1$-norm; for this distance $p$ is $1$-Lipschitz and a local isometry.
A tree is a finite connected graph without cycles whose edges are polygonal arcs, and $\ell(\tau)$ is
the sum of the lengths of the edges of the tree $\tau$. Then
\begin{equation}\label{4.38:eq-tree-lengths-1}
\begin{split}
&\inf\{\ell(\tau):\ \tau\subset\mathbb{T}\text{ a tree meeting every }D_i\}\\
&\qquad=\inf\{\ell(\hat\tau):\ \hat\tau\subset\mathbb{R}^4\text{ a tree meeting every }\hat D_i\},
\end{split}
\end{equation}
and the infimum on the right may be taken over trees $\hat\tau\subset\hat Q$.
\end{lemma}

\begin{proof}
Write $I_{\mathbb{T}}$ for the left side of \eqref{4.38:eq-tree-lengths-1}, $I$ for its right side,
and $I_{\hat Q}$ for the infimum of $\ell(\hat\tau)$ over trees $\hat\tau\subset\hat Q$ meeting every
$\hat D_i$. Trivially $I\le I_{\hat Q}$. We prove $I_{\mathbb{T}}\le I$ and $I_{\hat Q}\le I_{\mathbb{T}}$;
together these give $I_{\mathbb{T}}=I=I_{\hat Q}$, which is the assertion.

\emph{The inequality $I_{\mathbb{T}}\le I$.} Let $\hat\tau\subset\mathbb{R}^4$ be a tree meeting every
$\hat D_i$. The map $p$ increases no length, since the distance of $\mathbb{T}$ is the quotient of the
$\ell^1$-distance of $\mathbb{R}^4$; hence $\ell(p(\hat\tau))\le\ell(\hat\tau)$. The set $p(\hat\tau)$
is connected, as the continuous image of a connected set, and it meets every $D_i$: if
$x\in\hat\tau\cap\hat D_i$, then $p(x)\in p(\hat\tau)\cap D_i$. Choose one such meeting point for each
$i$. Subdivide the arcs $p(e)$, $e$ an edge of $\hat\tau$, at the chosen meeting points and at the
endpoints of their mutual intersections and self-intersections (two polygonal arcs meet in finitely
many points and segments), so that $p(\hat\tau)$ becomes a finite connected graph whose edges are
polygonal arcs and whose vertices include the chosen meeting points. As long as this graph is not a
tree, delete an edge lying on a cycle.
This keeps the set connected, keeps every chosen meeting point (each is a vertex), and does not
increase the length. After finitely many deletions one obtains a tree $\tau\subset\mathbb{T}$ meeting
every $D_i$ with $\ell(\tau)\le\ell(\hat\tau)$. Taking the infimum over $\hat\tau$ gives
$I_{\mathbb{T}}\le I$.

\emph{The inequality $I_{\hat Q}\le I_{\mathbb{T}}$.} Let $\tau\subset\mathbb{T}$ be a tree meeting every
$D_i$. Since $\tau$ is simply connected, it lifts through the covering $p$ to a tree
$\hat\tau\subset\mathbb{R}^4$ with $p(\hat\tau)=\tau$. Since $p$ is a local isometry,
$\ell(\hat\tau)=\ell(\tau)$. Since $\tau$ meets $D_i=p(\hat D_i)$, there is $z_i\in\mathbb{Z}^4$ such
that $\hat\tau$ meets $\hat D_i+sz_i$.

Let $\mathsf{f}\colon\mathbb{R}\to[0,a]$ be the $s$-periodic function given on $[0,s]$ by
\begin{equation*}
\mathsf{f}(u)=
\begin{cases}
u, & u\in[0,a],\\
a, & u\in[a,s-a],\\
s-u, & u\in[s-a,s].
\end{cases}
\end{equation*}
The middle interval is non-empty because $a\le s/2$ gives $a\le s-a$. The three formulas agree at the
junctions: both adjacent formulas give $a$ at $u=a$ and at $u=s-a$. Moreover $\mathsf{f}(s)=0=\mathsf{f}(0)$,
so the periodic extension is continuous. On each piece the slope is $0$ or $\pm1$, so $\mathsf{f}$ is
$1$-Lipschitz. Define $\mathsf{F}\colon\mathbb{R}^4\to\mathbb{R}^4$ by
\begin{equation*}
\mathsf{F}(x)_\mu:=c_\mu+\mathsf{f}(x_\mu-c_\mu),\qquad \mu=1,\dots,4.
\end{equation*}
The map $\mathsf{F}$ takes values in $\hat Q$. It maps $\mathbb{R}^4$ onto $\hat Q$: for $x\in\hat Q$
one has $x_\mu-c_\mu\in[0,a]$, where $\mathsf{f}(u)=u$, so $\mathsf{F}(x)=x$.

The $\mu$-th component of $\mathsf{F}$ depends only on $x_\mu$, and it is $1$-Lipschitz in that
variable. Hence $|\mathsf{F}(x)_\mu-\mathsf{F}(y)_\mu|\le|x_\mu-y_\mu|$ for every $\mu$. It follows that
$\mathsf{F}$ increases no $\ell^1$-length; it increases no $\ell^2$- or $\ell^\infty$-length either.

By the $s$-periodicity of $\mathsf{f}$, $\mathsf{F}(x+sz)=\mathsf{F}(x)=x$ for $x\in\hat Q$ and
$z\in\mathbb{Z}^4$. Hence $\mathsf{F}$ maps $\hat D_i+sz_i$ onto $\hat D_i$.

Consequently $\mathsf{F}(\hat\tau)\subset\hat Q$ has the following properties:
\begin{itemize}
\item it is connected, as a continuous image of the connected set $\hat\tau$;
\item its length is at most $\ell(\hat\tau)=\ell(\tau)$;
\item it meets every $\hat D_i$.
\end{itemize}
Applying to $\mathsf{F}(\hat\tau)$ the subdivision and edge deletion used in the first inequality (the
arcs $\mathsf{F}(e)$ are polygonal, since $\mathsf{F}$ is piecewise affine), with the meeting
points kept as vertices, yields a tree in $\hat Q$ meeting every $\hat D_i$ whose length is at most
$\ell(\tau)$. Taking the infimum over $\tau$ gives $I_{\hat Q}\le I_{\mathbb{T}}$.
\end{proof}

\begin{remark}[Walk distances read on the torus and on the lift]
Let $\omega$ be a walk of a walk system in the sense of Definition~\ref{4.4:def-walk-systems}, and let
$|\omega|$ be its footprint. Let $d(\omega)$ be the distance of $\omega$ entering the walk
expansions, namely the infimum of the lengths $\ell(\tau)$ of trees $\tau$ meeting every one of the
finitely many closed cells that make up $|\omega|$, and, for two bonds $b,b'$, let $d(\omega;b,b')$ be
the same infimum with the closed cells of $b$ and $b'$ adjoined to those of $|\omega|$. These
distances are read in the metric of $\mathbb{T}$, and so is the
predicate $d(\omega)\le r'$, with $r'$ the near--far threshold of
Definition~\ref{4.4:def-auxiliary-parameters}, through which the localised averaged Laplacian, the
operators of the cell-resolvent construction and the cell resolvents of
Lemma~\ref{4.4:lem-cell-resolvent} are defined. Where the metric is to be stressed we write
$d_{\mathbb{T}}(\omega)$. A lift of $\omega$ is a walk $\hat\omega$ in $\mathbb{R}^4$ whose image under
$p$ is $\omega$ term by term, and we write $d_{\mathbb{R}^4}(\hat\omega)$ for the same distance
computed in $\mathbb{R}^4$ for a lift $\hat\omega$. For walks whose footprints
lie in the image under $p$ of a cube of side at most $\tfrac12 s$, these distances are, by
Lemma~\ref{4.38:lem-tree-lengths}, those of the lift $\hat\omega$ inside that cube.

The predicate itself is decided by the lift. Put $\rho_w:=2\rho_0$, that is, $\sqrt{4}\,\rho_0$ with
$4$ the dimension, where $\rho_0$ is as in Definition~\ref{4.4:def-form-walk-hypothesis}, and suppose
$r'+2\rho_w<\tfrac12 s$. If $d_{\mathbb{T}}(\omega)\le r'$, then the footprint $|\omega|$ lies in the
image of one cube of side $r'+2\rho_w<\tfrac12 s$, and, by the last sentence of the preceding
paragraph, $d_{\mathbb{T}}(\omega)=d_{\mathbb{R}^4}(\hat\omega)$ for the lift $\hat\omega$ inside that
cube. Otherwise, since $p$ is $1$-Lipschitz, every lift $\hat\omega$ of $\omega$ satisfies
\begin{equation*}
d_{\mathbb{R}^4}(\hat\omega)\ge d_{\mathbb{T}}(\omega)>r'.
\end{equation*}
\end{remark}

\begin{lemma}[Growth recursions factorise over non-touching components]
\label{4.38:lem-growth-factorise}
Let $\mathsf{K}$ be a partition into cubes of side $c_{\mathsf{K}}>1$; call its cubes
$\mathsf{K}$-cubes, and say that two sets \emph{touch} if, regarded as closed subsets of
$\mathbb{R}^4$, they share a point (corners included). Consider a
recursion on unions of $\mathsf{K}$-cubes,
\begin{equation}\label{4.38:eq-growth-factorise-1}
Z_{n+1}=\mathsf{M}\bigl(Z_n\cup(\mathsf{R}_{n+1})^{\sim}\bigr),\qquad
\mathsf{R}_{n+1}=\mathsf{R}_n\cup\mathsf{T}_n ,
\end{equation}
with the following properties.
\begin{itemize}
\item[(r1)] For a union $W$ of $\mathsf{K}$-cubes, $W^{\sim}$ is $W$ with one layer of
$\mathsf{K}$-cubes added.
\item[(r2)] $\mathsf{M}$ replaces the connected components of its argument by boxes
containing them, and this is iterated until the result is stable.
\item[(r3)] The test set of stage $n$ consists of the bonds with an end-point in $Z_n$
within distance $\varpi$ of $Z_n^c$, where $\varpi<c_{\mathsf{K}}$, and of the
$\mathsf{K}$-cubes of $Z_n$ touching $Z_n^c$, minus what was tested at earlier stages.
The test set and the set $\mathsf{T}_n$ of cubes at which tests fail are functions of
$(\mathsf{R}_n,Z_n)$, read within distance $\varpi$, and of the answers there.
\end{itemize}
A \emph{datum} $\ell$ is an answer pattern on the test sets of its own run, its base part
included; the \emph{run} of $\ell$ is the sequence $(Z_n,\mathsf{R}_n)_{n\ge0}$ produced by
\eqref{4.38:eq-growth-factorise-1} from the base part, each test being answered as $\ell$
prescribes, and the run \emph{ends at} $Z$ if $Z_n=Z$ for all large $n$. Then the
following hold.
\begin{itemize}
\item[(a)] (Gluing.) If the runs of two data $\ell^1,\ell^2$ end at regions $Z^1,Z^2$
that do not touch, then the run of $\ell^1\cup\ell^2$ has $Z_n=Z_n^1\cup Z_n^2$ for every
$n$, its test sets are the disjoint unions of those of the two runs, and admissibility,
the terms $E^{T}$ and $E^{H}$ that the ruled family of this chapter attaches to a run,
and the tested factors of the small-field function $\chi^{\wedge}$ of that family
factorise.
\item[(b)] (Restriction.) If a run ends at $Z=\bigsqcup_i C_i$, the $C_i$ being the
connected components of $Z$, then the data $\ell\cap C_i$, with base part equal to the
base part intersected with $C_i$, have the runs $Z_n\cap C_i$.
\end{itemize}
\end{lemma}

\begin{proof}
We prove (a) by induction on $n$, starting from the base parts; the base part of
$\ell^1\cup\ell^2$ is the union of those of $\ell^1$ and $\ell^2$. Assume that at stage
$n$ the region, the set $\mathsf{R}_n$ and the test sets of all earlier stages of the run
of $\ell^1\cup\ell^2$ are the unions of those of the runs of $\ell^1$ and $\ell^2$.

Since $\mathsf{M}$ only enlarges its argument, \eqref{4.38:eq-growth-factorise-1} gives
$Z_n^i\subset Z_{n+1}^i$, hence $Z_n^i\subset Z^i$ for $i=1,2$. As $Z^1$ and $Z^2$ do not
touch, neither do $Z_n^1$ and $Z_n^2$; being unions of cubes of the one partition
$\mathsf{K}$, they are at sup-distance at least $c_{\mathsf{K}}$. Since
$\varpi<c_{\mathsf{K}}$, a point of $(Z_n^1)^c$ within distance $\varpi$ of $Z_n^1$ is not
in $Z_n^2$; and since $c_{\mathsf{K}}>1$ exceeds the length of a bond, a bond with an
end-point in $Z_n^1$ has no end-point in $Z_n^2$. Consequently, for a bond with an
end-point in $Z_n^1$, the distance of that end-point to $(Z_n^1)^c$ and its distance to
$(Z_n^1\cup Z_n^2)^c$ agree whenever one of them is at most $\varpi$; if the other
end-point lies outside $Z_n^1$, it lies in $(Z_n^1)^c$ and not in $Z_n^2$, so its
distances to $(Z_n^1)^c$ and to $(Z_n^1\cup Z_n^2)^c$ are both zero. Thus the bonds tested
at $Z_n^1$ are the same whether $Z_n^1$ is taken alone or as part of the union.

A cube of $Z_n^1$ touches $(Z_n^1)^c$ through a $\mathsf{K}$-cube $c'$ with
$c'\not\subset Z_n^1$. If $c'\subset Z_n^2$, the regions $Z_n^1$ and $Z_n^2$ would touch,
which is excluded; so $c'$ lies outside $Z_n^1\cup Z_n^2$, and the cubes of $Z_n^1$ touching
the complement are the same for $Z_n^1$ and for the union. The same holds with the roles
of $1$ and $2$ exchanged. Together with the induction hypothesis on what was tested
before, and since by (r3) the test sets and $\mathsf{T}_n$ are read within distance
$\varpi$, the test sets, the answers on them and the sets $\mathsf{T}_n$ of the run of
$\ell^1\cup\ell^2$ are the unions of those of the two runs; the test sets are disjoint by
the preceding paragraphs. Hence
$\mathsf{R}_{n+1}=\mathsf{R}_{n+1}^1\cup\mathsf{R}_{n+1}^2$.

By (r1), $\sim$ commutes with unions. Put
\begin{equation*}
W^i:=Z_n^i\cup(\mathsf{R}_{n+1}^i)^{\sim},\qquad i=1,2 .
\end{equation*}
By the induction hypothesis, the identity
$\mathsf{R}_{n+1}=\mathsf{R}^1_{n+1}\cup\mathsf{R}^2_{n+1}$ and (r1), the argument of
$\mathsf{M}$ in \eqref{4.38:eq-growth-factorise-1} for the run of $\ell^1\cup\ell^2$ is
$W^1\cup W^2$.
Then $W^i\subset Z_{n+1}^i\subset Z^i$, so $W^1$ and $W^2$ do not touch, and the connected
components of $W^1\cup W^2$ are those of $W^1$ together with those of $W^2$. By (r2),
every stage of the iteration defining $\mathsf{M}$ applied to $W^i$ stays inside
$\mathsf{M}W^i\subset Z^i$; the boxes grown from $W^1$ and from $W^2$ therefore never
touch, and
\begin{equation*}
\mathsf{M}(W^1\cup W^2)=\mathsf{M}W^1\cup\mathsf{M}W^2 ,
\end{equation*}
that is, $Z_{n+1}=Z_{n+1}^1\cup Z_{n+1}^2$. A run that has stopped has empty test sets at
all later stages, so the argument continues when one of the two runs stops before the
other. This completes the induction; admissibility, $E^{T}$, $E^{H}$ and the tested
factors of $\chi^{\wedge}$ factorise because the test sets are the disjoint unions and
the answers on them are those of $\ell^1$ and of $\ell^2$.

For (b), run the same induction with $Z_n\cap C_i$ in place of $Z_n^i$. Every box
produced by $\mathsf{M}$ is connected and contained in $Z$, hence contained in one
component $C_i$; so the components $C_i$ play the role of the non-touching regions $Z^i$
above, and the run of $\ell\cap C_i$ is $Z_n\cap C_i$.
\end{proof}

\begin{remark}
The recursion by which the regions of the small-field function $\chi^{\wedge}$ of the
ruled family of this chapter at the split $\mathbb{B}$ are grown is of the form
\eqref{4.38:eq-growth-factorise-1}, with $c_{\mathsf{K}}=LM\check{R}_t$, where $t$ is the
level of the step at hand; this side exceeds $\varpi_T+1$, where $\varpi_T$ is the reach
of the tests of that recursion, by the reach condition of the torus-independence lemma of
clause~(0). Any recursion obtained from it by a change of notation that keeps
(r1)--(r3) is again of this form. Deciding, as in (r3), which $\mathsf{K}$-cubes of $Z_n$
touch $Z_n^c$ already reads the touching layer of $Z_n$, that is, the $\mathsf{K}$-cubes
outside $Z_n$ that touch $Z_n$.
\end{remark}

\begin{lemma}[Properties of the coupling profile]\label{4.38:lem-coupling-profile}
Lengths are measured in level-$j$ units. Let $i\le j$ be levels. Call \emph{cells} the level-$i$
label cubes; they have side $a_i=M\check{R}_iL^{-(j-i)}$, and the level-$j$ cells have side
$a_j=M\check{R}_j$. Let $\ell$ be a label, taken before or after the large-field operation
$\mathbf{R}$, with or without its islands adjoined, and put $\Lambda:=\Lambda_i(\ell)$. At a step of
$\mathbf{R}$, $\ell$ denotes the retained label, and $\Lambda$ denotes the level-$i$ domain of the
sequence in force at that step, as fixed in the standing data (Definition~\ref{4.38:def-standing-data}),
which need not be $\Lambda_i(\ell)$. Let $\mathcal{D}$ be the set of cells of
$[\Lambda^c]_i$ (Definition~\ref{4.38:def-intrinsic-trace}), and let $\varphi:=\varphi_i$ be the
coupling profile \eqref{4.38:eq-standing-data-a}:
\begin{equation*}
\varphi(x)=\prod_{D\in\mathcal{D}}\Bigl(1-\prod_{\mu=1}^{4}
\psi_{\varphi}\bigl(\operatorname{dist}(x_\mu,D_\mu)/a_i\bigr)\Bigr),
\end{equation*}
where $D_\mu$ is the $\mu$-th coordinate interval of $D$ and $\psi_{\varphi}\in
C^\infty([0,\infty);[0,1])$ is the fixed function with $\psi_{\varphi}\equiv1$ on $[0,\tfrac18]$ and
$\psi_{\varphi}\equiv0$ on $[\tfrac14,\infty)$. We write $\varphi(\cdot\,;\ell)$ when the label is
to be shown, and, for a set $\Lambda'$, $\varphi^{\Lambda'}$ for the same product taken over the cells
of $[\Lambda'^c]_i$; thus $\varphi=\varphi^{\Lambda}$, and $\varphi^{\ell}:=\varphi^{\Lambda_i(\ell)}$
is the profile of the label $\ell$ alone.
$\operatorname{dist}_{\infty}$ is the sup-distance; two distinct cubes of one partition \emph{touch}
if they intersect. Then the following hold.
\begin{enumerate}
\item[(a)] $\varphi\in C^\infty(\mathbb{R}^4;[0,1])$, and
\begin{equation*}
\varphi=0\ \text{ on }\ \bigl\{x:\ \operatorname{dist}_{\infty}(x,D)\le\tfrac18a_i
\text{ for some }D\in\mathcal{D}\bigr\}\supset\Lambda^c;
\end{equation*}
hence, on every torus, $\varphi\in C_0^\infty(\Lambda)$, which is the condition
``$\varphi_j\in C_0^\infty(\Lambda_j)$'' of \cite[p.~259]{B-CMP119}.
\item[(b)] $\varphi(x)=1$ whenever $\operatorname{dist}_{\infty}(x,D)\ge\tfrac14a_i$ for every
$D\in\mathcal{D}$, that is, whenever $\operatorname{dist}_{\infty}(x,[\Lambda^c]_i)\ge\tfrac14a_i$;
and $\varphi(x)\neq1$ implies $\operatorname{dist}_{\infty}(x,D)<\tfrac14a_i$ for some
$D\in\mathcal{D}$. In particular $\varphi=1$ on $\Lambda^{\sim-1}$ and on the set of points at
sup-distance at most $\tfrac34a_i$ from it; and $h_z\varphi=h_z$ for every level-$i$ site $z$ of
$\Lambda^0$, the set $\Lambda$ less one layer of cells \cite[p.~259]{B-CMP119} (cube-wise
$\Lambda^{\sim-1}$), where $h_z$ is the tent function of \cite[(3.40), (3.65)]{B-CMP119}. At a level
where $\Lambda^c$ is a union of cells these are the condition ``$\varphi_j=1$ on
$\Lambda_j^{\sim-1}$'' of \cite[p.~259]{B-CMP119} and the identity ``$h_z\varphi_j=h_z$ for
$z\in\Lambda_j^0$'' of \cite[p.~283]{B-CMP119}. At the levels $i$ with $k-N_0<i\le k$ ($k$ the
level of the $\mathbf{R}$-step, $N_0$ the integer of \cite[(1.10)--(1.11)]{B-CMP122a} bounding this
range of levels) at which the label's set is the prepared set $Z''_i$ of
\cite[(1.10)--(1.11)]{B-CMP122a}, a union of $M$-cubes which is not a union of cells, the assertions
above hold with $\Lambda^{\sim-1}$ formed from the cells, $[\Lambda^c]_i$ taking the place of
$\Lambda^c$, and $\varphi$ vanishes on all of $[Z''_i]_i$.
\item[(c)] (Naturality.) Let $\mathfrak{G}_i$ be the group of Euclidean motions of the lattice which
preserve the partition into cells. For $r\in\mathfrak{G}_i$,
$\varphi(rx;r\ell)=\varphi(x;\ell)$.
\item[(d)] (Multiplicativity.) If $[\Lambda^c]_i=A\sqcup B$ with $A$, $B$ unions of cells sharing
no cell, then $\varphi^{\Lambda}=\varphi^{A^c}\cdot\varphi^{B^c}$. In particular, at an
$\mathbf{R}$-step,
\begin{equation}\label{4.38:eq-coupling-profile-1}
\varphi^{\mathrm{in\ force}}=\bigl(\text{the non-island factor}\bigr)\cdot
\prod_a\varphi^{(S_a)^c},
\end{equation}
where $S_a$ is the level-$i$ large-field set of the island $X_a$ under the sequence in force (for the
sequence made by the operations $\mathbf{T}$, in force before $\mathbf{R}$: $S_a\subset X_a$, with
$S_a=X_a$ at $i=j$; for the prepared sequence: $S_a=Z''_i\cap X_a$), and the non-island factor is
$\varphi^{A^c}$ for $A$ the set of cells of $[\Lambda^c]_i$ lying in no island $X_a$; under the
sequence made by the operations $\mathbf{T}$ the non-island factor is $\varphi^{\ell}$. Here the
islands' level-$i$ sets are cell-disjoint from the retained label's and from each other's. The sets
$S_a$ are unrelated to the third sequence $\{S_i\}$ of localization domains of
\cite[p.~254]{B-CMP119}.
\item[(e)] (Radius.) $\varphi(x)$ is determined by the cells $D\in\mathcal{D}$ with
$\operatorname{dist}_{\infty}(x,D)<\tfrac14a_i$. If $W$ is a domain and $x$ lies in an $M$-cube of
$W$ (or is a point of $W$), these cells lie in $[W]_i^{+1}\subset W^{+1}$. If $W^{+1}$ contains no
cell of $[\Lambda^c]_i$, then $\varphi\equiv1$ on $W$.
\item[(f)] (The ramp lies in the strip.) At level $j$, let $S$ be a union of level-$j$ cells (an
island $X_a$, or any such $S$), and let $x$ satisfy $\operatorname{dist}_{\infty}(x,S)<\tfrac14a_j$.
Then every closed $M$-cube $c$ containing $x$ lies in $S^{\sharp}$. Hence, with $i=j$,
\begin{equation*}
\bigl\{x:\ \varphi^{(X_a)^c}(x)\neq1\bigr\}\subset X_a^{\sharp};
\end{equation*}
and for $i\le j$ and a set $S'\subset X_a$ which is a union of level-$i$ cells,
\begin{equation*}
\bigl\{x:\ \operatorname{dist}_{\infty}(x,S')<\tfrac14a_i\bigr\}\subset X_a^{\sharp}
\end{equation*}
likewise, because $a_i\le a_j$. The sets $S_a$ and $[Z''_i\cap X_a]_i$ are such sets $S'$, and so
are the round-ups $[\,\cdot\,]_i$ of the large-field sets of the intermediate stages of the
preparatory integrations of \cite[(1.60)--(1.62)]{B-CMP122a} lying inside $X_a$.
\end{enumerate}
\end{lemma}

\begin{proof}
Two elementary facts on a partition of $\mathbb{R}^4$ into closed cubes of side $c$ with corners in
$c\mathbb{Z}^4$ are used throughout; both appear in the proof of Lemma~\ref{4.38:lem-round-up-touch}.
Index the cubes by $n\in\mathbb{Z}^4$. Two cubes are equal or touch if and only if their index
vectors satisfy $|n-n'|_\infty\le1$; if they neither are equal nor touch, then
$|n_\mu-n'_\mu|\ge2$ for some $\mu$, the $\mu$-th coordinate intervals are a full interval $c$ apart,
and every point of the one cube is at sup-distance at least $c$ from the other. Moreover the
sup-distance between two unions of cubes of the partition is an integer multiple of $c$. For a cube
$D=\prod_\mu D_\mu$ and a point $x$ we have
$\operatorname{dist}_{\infty}(x,D)=\max_\mu\operatorname{dist}(x_\mu,D_\mu)$. For $D\in\mathcal{D}$ put
\begin{equation*}
\eta_D(x):=1-\prod_{\mu=1}^{4}\psi_{\varphi}\bigl(\operatorname{dist}(x_\mu,D_\mu)/a_i\bigr),
\qquad\text{so that}\qquad \varphi=\prod_{D\in\mathcal{D}}\eta_D .
\end{equation*}

(a) Let $I=[\alpha,\beta]$ be an interval. The function $t\mapsto\operatorname{dist}(t,I)$ vanishes
on $I$ and is smooth off $\{\alpha,\beta\}$, being $\alpha-t$ for $t<\alpha$ and $t-\beta$ for
$t>\beta$. Since $\psi_{\varphi}(\cdot/a_i)\equiv1$ on $[0,a_i/8]$, the function
$t\mapsto\psi_{\varphi}(\operatorname{dist}(t,I)/a_i)$ is identically $1$ on
$[\alpha-a_i/8,\beta+a_i/8]$, whose interior contains $\alpha$ and $\beta$, and it is smooth
elsewhere; although $t\mapsto\operatorname{dist}(t,I)$ is not smooth at $\alpha$ and $\beta$, the
composite is, because $\psi_{\varphi}$ is constant near $0$. Hence it is $C^\infty$ on $\mathbb{R}$.
Consequently each $\eta_D$ is
$C^\infty$, with values in $[0,1]$ because $\psi_{\varphi}$ takes values in $[0,1]$.

The factor $\eta_D$ equals $1$ at $x$ unless $\operatorname{dist}(x_\mu,D_\mu)<\tfrac14a_i$ for all
$\mu$, that is, unless $\operatorname{dist}_{\infty}(x,D)<\tfrac14a_i$. Let $C$ be a cell containing
$x$. If $|y-x|_\infty<\tfrac34a_i$ and $\operatorname{dist}_{\infty}(y,D)<\tfrac14a_i$, then
$\operatorname{dist}_{\infty}(x,D)<a_i$, so $D$ is at sup-distance less than $a_i$ from $C$ and
therefore equals $C$ or is one of the $3^4-1$ cells touching $C$. Thus on a neighbourhood of each
point all but at most $3^4$ factors are identically $1$: the product is locally finite, so
$\varphi\in C^\infty$ and $0\le\varphi\le1$.

If $\operatorname{dist}_{\infty}(x,D)\le\tfrac18a_i$ for some $D\in\mathcal{D}$, then
$\operatorname{dist}(x_\mu,D_\mu)\le\tfrac18a_i$ for every $\mu$, so every factor
$\psi_{\varphi}(\operatorname{dist}(x_\mu,D_\mu)/a_i)$ equals $1$, $\eta_D(x)=0$ and $\varphi(x)=0$.
Every point of $\Lambda^c$ lies in a cell of $[\Lambda^c]_i$, so $\Lambda^c\subset\bigcup\mathcal{D}$
and the set on which $\varphi$ was just shown to vanish contains $\Lambda^c$. Hence
\begin{equation*}
\operatorname{supp}\varphi\subset\bigl\{x:\ \operatorname{dist}_{\infty}(x,\Lambda^c)
\ge\tfrac18a_i\bigr\},
\end{equation*}
a closed set contained in the interior of $\Lambda$, and compact on a torus. This is
$\varphi\in C_0^\infty(\Lambda)$.

(b) If $\operatorname{dist}_{\infty}(x,D)\ge\tfrac14a_i$, then some $\mu$ has
$\operatorname{dist}(x_\mu,D_\mu)\ge\tfrac14a_i$; $\psi_{\varphi}$ vanishes there, so $\eta_D(x)=1$.
Hence $\operatorname{dist}_{\infty}(x,D)\ge\tfrac14a_i$ for all $D\in\mathcal{D}$ implies
$\varphi(x)=1$; since $\operatorname{dist}_{\infty}(x,[\Lambda^c]_i)$ is the infimum of
$\operatorname{dist}_{\infty}(x,D)$ over $D\in\mathcal{D}$, this is the second formulation.
Contrapositively, $\varphi(x)\ne1$ implies $\eta_D(x)\ne1$ for some $D$, hence
$\operatorname{dist}_{\infty}(x,D)<\tfrac14a_i$ for that $D$. Nothing about the values of
$\psi_{\varphi}$ on $(\tfrac18,\tfrac14)$ is used.

A point of $\Lambda^{\sim-1}$ lies in a cell $C\subset\Lambda$ all of whose touching cells lie in
$\Lambda$. A cell $D\in\mathcal{D}$ meets $\Lambda^c$, so it neither equals nor touches $C$; by the
facts recalled at the start, their index vectors differ by at least $2$ in some coordinate, and
$\operatorname{dist}_{\infty}(x,D)\ge a_i$ for every $x\in C$ (this is the step of the proof of
Lemma~\ref{4.38:lem-round-up-touch}(e)). A point at sup-distance at most $\tfrac34a_i$ from $C$
therefore has $\operatorname{dist}_{\infty}(\cdot,D)\ge a_i-\tfrac34a_i=\tfrac14a_i$ for every
$D\in\mathcal{D}$, and $\varphi=1$ there by the first assertion.

The tent. Let $z$ be a level-$i$ site of $\Lambda^0$. The set $\Lambda^0$ is $\Lambda$ less one
layer of cells \cite[p.~259]{B-CMP119}, that is, cube-wise $\Lambda^{\sim-1}$; so $z$ lies in a
cell $C\subset\Lambda^{\sim-1}$, and, as just shown, $\operatorname{dist}_{\infty}(z,[\Lambda^c]_i)\ge
a_i$. The tent $h_z$ of \cite[(3.65)]{B-CMP119} is the product \cite[(3.40)]{B-CMP119} of the
functions $h(t)=\max\{1-|t|,0\}$ on the level-$i$ unit lattice (the lattice of
\cite[p.~283]{B-CMP119}, on which $\sum_zh_z=1$). Hence
\begin{equation*}
\operatorname{supp}h_z\subset\bigl\{x:\ |x-z|_\infty\le L^{-(j-i)}\bigr\},
\end{equation*}
$L^{-(j-i)}$ being one level-$i$ unit in level-$j$ units. Now $L^{-(j-i)}\le\tfrac34a_i$ is
equivalent to $M\check{R}_i\ge\tfrac43$, which is true since $M\ge2$. Therefore
$\operatorname{dist}_{\infty}(x,[\Lambda^c]_i)\ge a_i-\tfrac34a_i=\tfrac14a_i$ on
$\operatorname{supp}h_z$, the first assertion gives $\varphi=1$ there, and $h_z\varphi=h_z$. In
\cite[p.~283]{B-CMP119} this identity is stated for Ba{\l}aban's $\varphi_j$; here it is a property of
the profile \eqref{4.38:eq-standing-data-a}.

(c) Every $r\in\mathfrak{G}_i$ is the composition of a translation by a vector $t\in a_i\mathbb{Z}^4$
with a signed permutation $s$ of the coordinates about the centre of a cell. Such an $r$ permutes the
cells, and $r\mathcal{D}$ is the set of cells of $[r\Lambda^c]_i$, the set $\mathcal{D}$ of the label
$r\ell$. If $s$ sends the coordinate $\mu$ to the coordinate $\sigma(\mu)$ with sign $\epsilon_\mu$,
then $(rx)_{\sigma(\mu)}$ and $(rD)_{\sigma(\mu)}$ are obtained from $x_\mu$ and $D_\mu$ by one and the
same map $u\mapsto\epsilon_\mu u+b_\mu$, $b_\mu$ a constant; since
$\operatorname{dist}(-u,-I)=\operatorname{dist}(u,I)$ and the distance from a point to an interval
is translation invariant, the multiset
$\{\operatorname{dist}((rx)_\nu,(rD)_\nu)\}_\nu$ equals
$\{\operatorname{dist}(x_\mu,D_\mu)\}_\mu$. The product over $\mu$ is symmetric in its factors, so
$\eta_{rD}(rx)=\eta_D(x)$, and
\begin{equation*}
\varphi(rx;r\ell)=\prod_{D\in\mathcal{D}}\eta_{rD}(rx)=\prod_{D\in\mathcal{D}}\eta_D(x)
=\varphi(x;\ell),
\end{equation*}
the first equality being the re-indexing of the product over the cells of $[r\Lambda^c]_i$ by
$D\mapsto rD$.

(d) The set $\mathcal{D}$ is the disjoint union of the cells of $A$ and the cells of $B$. Since $A$
and $B$ are unions of cells, $[(A^c)^c]_i=A$ and $[(B^c)^c]_i=B$, so the product over $\mathcal{D}$
splits as $\varphi^{\Lambda}=\varphi^{A^c}\cdot\varphi^{B^c}$. For every sequence in force
$S_a\subset X_a$, and $X_a$ is a union of level-$j$ cells, hence of level-$i$ cells (the partitions
are compatible), so $[S_a]_i\subset X_a$. The cells of $A$ lie in no island and the islands pairwise
share no cell, so $A$ and the sets $[S_a]_i$ are pairwise cell-disjoint. Applying the first assertion
repeatedly to the decomposition of $[\Lambda^c]_i$ into $A$ and the sets $[S_a]_i$ gives
\eqref{4.38:eq-coupling-profile-1}.

For the sequence made by the operations $\mathbf{T}$, the complement of its level-$i$ domain
$\Lambda$ is, at an $\mathbf{R}$-step,
\begin{equation*}
\Lambda^{c}=\Lambda_i(\ell)^c\cup\bigcup_aS_a ,
\end{equation*}
$S_a$ the level-$i$ large-field set of the island $X_a$. By the nesting $\Lambda_j\subset\Lambda_i$ of
the retained sequence \cite[(2.1)]{B-CMP119}, and of the sequences glued at the islands, by the
gluing statement for the islands' sequences at an $\mathbf{R}$-step in this chapter,
$\Lambda_i(\ell)^c\subset\Lambda_j(\ell)^c$ and $S_a\subset X_a$. The partitions into level-$i$ and
into level-$j$ cells are compatible and $a_i$ divides $a_j$ (by the lemma on the cell-size
multipliers $\check{R}_i$), so every level-$i$ cell lies inside a level-$j$ cell; since
$\Lambda_j(\ell)^c$ and $X_a$ are unions of level-$j$ cells, $[\Lambda_i(\ell)^c]_i\subset
\Lambda_j(\ell)^c$ and $[S_a]_i\subset X_a$. As $\Lambda_j(\ell)^c$ and the $X_a$ share no cell, these
round-ups are cell-disjoint; the cells of $[\Lambda^c]_i$ lying in no island are exactly those of
$[\Lambda_i(\ell)^c]_i$, and the non-island factor is $\varphi^{\ell}$.

(e) By the proof of (a), $\eta_D(x)=1$ unless $\operatorname{dist}_{\infty}(x,D)<\tfrac14a_i$, so only
such $D$ enter $\varphi(x)$. Apply Lemma~\ref{4.38:lem-round-up-touch}(b) with $\mathsf{K}=\mathsf{K}'$
the partition into level-$i$ cells ($c_{\mathsf{K}}=a_i$) and $\mathsf{k}$ the partition into
level-$i$ $M$-cubes, of side $ML^{-(j-i)}$, of which $W$ is a union and which is finer than the
cells, as that lemma requires. A cell $D$ with $\operatorname{dist}_{\infty}(x,D)<\tfrac14a_i$ is at
sup-distance less than $a_i$ from the level-$i$ $M$-cube of $W$ containing $x$, hence lies in
$[W]_i^{+1}$. Alternatively, from the point $x$: such a $D$ equals or touches the cell of $x$, so
$D\subset[x]_i^{+1}\subset[W]_i^{+1}$. Finally $[W]_i^{+1}\subset[W]_j^{+1}=W^{+1}$ by
Lemma~\ref{4.38:lem-round-up-touch}(c), applied with $\mathsf{K}$ the level-$j$ cells, $\mathsf{K}'$
the level-$i$ cells and $W'=W$.

For the last assertion, suppose $W^{+1}$ contains no cell of $[\Lambda^c]_i$. Then neither does
$[W]_i^{+1}\subset W^{+1}$, and Lemma~\ref{4.38:lem-round-up-touch}(e), applied with $\mathsf{K}$ the
level-$i$ cells and $Z=[\Lambda^c]_i$, gives $\operatorname{dist}_{\infty}(x,[\Lambda^c]_i)\ge
a_i>\tfrac14a_i$ for every $x\in W$. By (b), $\varphi\equiv1$ on $W$.

(f) Recall from \eqref{4.38:eq-standing-data-b} that $S^{\sharp}$ is $S$ together with $m_\sharp$
layers of $M$-cubes, the $n$-th layer consisting of the $M$-cubes at sup-distance $(n-1)M$ from $S$
($1\le n\le m_\sharp$), with $m_\sharp=\lfloor\tfrac12(\check{R}_k-1)\rfloor$ and $k$ the level of
the $\mathbf{R}$-step, here $k=j$. Let $c$ be a closed $M$-cube containing $x$. Then
\begin{equation*}
\operatorname{dist}_{\infty}(c,S)\le\operatorname{dist}_{\infty}(x,S)<\tfrac14a_j
=\tfrac14M\check{R}_j .
\end{equation*}
Both $c$ and $S$ are unions of cubes of the $M$-partition, so
$\operatorname{dist}_{\infty}(c,S)\in M\mathbb{Z}_{\ge0}$, and
$n_0:=\operatorname{dist}_{\infty}(c,S)/M$ is an integer smaller than $\check{R}_j/4$. The
multiplier $\check{R}_j$ is odd, being a power of the odd integer $L$ (by the lemma on the
cell-size multipliers), so $\check{R}_j/4$ is not an integer and $n_0\le(\check{R}_j-1)/4$; also
$m_\sharp-1=(\check{R}_j-3)/2$. If $\check{R}_j\ge5$, then
$(\check{R}_j-1)/4\le(\check{R}_j-3)/2$, so $n_0\le m_\sharp-1$. If $\check{R}_j=3$, then
$n_0<\tfrac34$ forces $n_0=0=m_\sharp-1$. In all cases $c$ lies in $S$ or in one of its first
$m_\sharp$ layers, that is, $c\subset S^{\sharp}$. The case distinction covers all cases because
$\check{R}_j\ge3$, which follows from the floor $\check{R}_j\ge L$ of the standing data
(Definition~\ref{4.38:def-standing-data}).

With (b), applied at level $j$ with $\Lambda=(X_a)^c$: if $\varphi^{(X_a)^c}(x)\ne1$, then
$\operatorname{dist}_{\infty}(x,D)<\tfrac14a_j$ for a cell $D\subset X_a$, hence
$\operatorname{dist}_{\infty}(x,X_a)<\tfrac14a_j$, and $c\subset X_a^{\sharp}$ by what was just
proved. For $S'\subset X_a$ a union of level-$i$ cells: the lemma on the cell-size multipliers
gives $\check{R}_i\le L^{j-i}\check{R}_j$, so
\begin{equation*}
a_i=M\check{R}_iL^{-(j-i)}\le M\check{R}_j=a_j ,
\end{equation*}
and $\operatorname{dist}_{\infty}(x,S')<\tfrac14a_i$ implies
$\operatorname{dist}_{\infty}(x,X_a)\le\operatorname{dist}_{\infty}(x,S')<\tfrac14a_j$, hence
$c\subset X_a^{\sharp}$.
\end{proof}

\begin{remark}
Assertions (a) and (b) show that, at levels where $\Lambda^c$ is a union of cells, the profile
\eqref{4.38:eq-standing-data-a} has the two properties required of $\varphi_j$ in
\cite[p.~259]{B-CMP119} and the property stated in \cite[p.~283]{B-CMP119}; the arguments of
\cite[Sects.~2--3]{B-CMP119} which use $\varphi_j$ through these properties apply to it unchanged.
At the levels of the prepared sets $Z''_i$ the operation $\sim$ of \cite{B-CMP119} adds one layer of
the $M$-cubes in terms of which the sets are defined; that $\varphi=1$ one such layer away from
$Z''_i$ is not part of (b), and at those levels the cells read by the coupling of the Wilson term are
those fixed in the representation and bound of the large-field terms. Assertion (c) is the naturality
of the profile read by the covariance statement of the torus-independence lemma of clause~(0);
(d) and (e) are used in part (iii)(b) of that lemma and in the factorisation of the profile over the
islands at an $\mathbf{R}$-step; and (f) gives the consistency of that factorisation with
\cite{B-CMP122b}.
\end{remark}

\begin{lemma}[The choice of coupling profile moves no total]\label{4.38:lem-profile-telescopes}
For $i\ge 0$ write $x_i(x)=1/g_i^2(x)$ for the position-dependent inverse couplings of
\cite[(2.24)]{B-CMP119}, and $\beta_i=\beta_i(g_{i-1})$ for $i\ge1$; let $\varphi_i$, $i\ge1$, be any
coupling profiles entering \cite[(2.24)]{B-CMP119} that satisfy the two conditions imposed on them in
\cite{B-CMP119} (see Lemma~\ref{4.38:lem-coupling-profile}).
\begin{itemize}
\item[(a)] For every level $j$ and every point $x$,
\begin{equation}\label{4.38:eq-profile-telescopes-a}
x_j(x)+\sum_{i=1}^{j}\beta_i\,\varphi_i(x)=x_0 ,
\end{equation}
a constant. Consequently the two members of the exponent of \cite[(2.23)]{B-CMP119} that carry
the profiles, namely the action $-A(x_j(\cdot),U)$ and the counterterms
$-\sum_{i=1}^{j}\beta_i A(\varphi_i,U)$ of \cite[(2.25)]{B-CMP119}, have the sum $-x_0A(U)$.
\item[(b)] The effective action at level $j$ (the exponent of $\rho_j$), which we denote $A_j$, as a
function of the configuration, the effective density $\rho_j$, its integral $\int\rho_j$, every total
(every quantity formed from the density $\rho_j$ as a whole rather than from its individual localised
terms), the polarisation tensors $\Pi^{(j)}_{\mathfrak{s}}$ (second variations of a total), the
increments $\beta_j$ and the couplings $g_j$ are the same for every such choice of profiles; in
particular they are not altered by taking for $\varphi_i$ the product profile of
\eqref{4.38:eq-standing-data-a}.
\item[(c)] At a step at which the operation $\mathbf{R}$ acts, the same cancellation holds plaquette by
plaquette, on the collar of each island, for the difference of the counterterm members before and
after the step:
\begin{equation}\label{4.38:eq-profile-telescopes-1}
-\beta_jA\bigl(h_z(1-\varphi_j^{\mathrm{pre}}),\cdot\bigr)-\beta_jA\bigl(h_z\varphi_j^{\mathrm{pre}},\cdot\bigr)
+\beta_jA(h_z,\cdot)=0 ,
\end{equation}
where $\varphi_j^{\mathrm{pre}}$ denotes the level-$j$ profile of the sequence of localisation domains
in force at that step (Lemma~\ref{4.38:lem-coupling-profile}) and the $h_z$, $z$ ranging over points,
are the tent functions with which the counterterm is localised.
\item[(d)] What the choice of profile does bear on is the value of individual terms that read the label,
namely $V(Y)$, $R'(X)$ and $\mathfrak{r}'(X)$, also for $X\in\mathcal{X}_j(\ell)$; it decides no total
and no increment $\beta_j$, since $\beta_j=\beta_j^{\mathfrak{A}}+\beta_j^{\Re}$, where
$\beta_j^{\mathfrak{A}}$ and $\beta_j^{\Re}$ are the two parts of the coupling increment,
$\beta_j^{\mathfrak{A}}$ involves no island and $\beta_j^{\Re}$ is a difference of totals. It also
fixes the radius within which a value $\varphi_i(x)$ reads
the label: for the product profile of \eqref{4.38:eq-standing-data-a} this radius is less than a
quarter of a level-$i$ cell, so that for $x$ in a region $W$ all cells read lie in $W^{+1}$.
\end{itemize}
\end{lemma}

Here $\mathcal{X}_j(\ell)$ is the set of domains $X\in\mathbf{D}_j$ on whose one-cell enlargement
$X^{+1}$ the label $\ell$ has empty trace, and $W^{+1}$ is the one-cell enlargement of $W$.

\begin{proof}
(a) The profile $\varphi_j$ enters through \cite[(2.24), p.~259]{B-CMP119}, which reads
\begin{equation*}
\frac{1}{g_{j-1}^2(x)}=\frac{1}{g_j^2(x)}+\beta_j(g_{j-1})\,\varphi_j(x);
\end{equation*}
in the present notation, $x_{i-1}(x)=x_i(x)+\beta_i\varphi_i(x)$ for every $i=1,\dots,j$. Summing these
identities over $i=1,\dots,j$, the intermediate terms
$x_1(x),\dots,x_{j-1}(x)$ cancel, leaving
\begin{equation*}
x_0(x)=x_j(x)+\sum_{i=1}^{j}\beta_i\varphi_i(x).
\end{equation*}
Here $x_0(x)=x_0$ is the bare inverse coupling, a constant; this is
\eqref{4.38:eq-profile-telescopes-a}.

The weighted action $A(c,U)$ of Definition~\ref{1.1:def-wilson-action} is linear in the weight $c$,
and $A(1,U)=A(U)$. The action member of the exponent of \cite[(2.23)]{B-CMP119} is
$-A(x_j(\cdot),U)$. The counterterms are those of \cite[(2.25)]{B-CMP119}, which, written at level $j$
with summation index $i$ and with $U$ for the level-$j$ field, is
\begin{equation*}
\mathbf{E}_j(U)=\sum_{i=1}^{j}\Bigl[\mathbf{E}^{(i)}(\Lambda_i,g_{i-1},U)
-\mathbf{E}^{(i)}(\Lambda_i,g_{i-1},1)-\beta_i(g_{i-1})A(\varphi_i,U)\Bigr].
\end{equation*}
By linearity in the weight and by \eqref{4.38:eq-profile-telescopes-a},
\begin{align*}
-A(x_j(\cdot),U)-\sum_{i=1}^{j}\beta_iA(\varphi_i,U)
&=-A\Bigl(x_j(\cdot)+\sum_{i=1}^{j}\beta_i\varphi_i(\cdot),\,U\Bigr)\\
&=-A(x_0,U)=-x_0A(U).
\end{align*}
This is the add-and-subtract by which the coupling constant renormalization is performed in
\cite{B-CMP119}: ``The coupling constant renormalization is performed by subtracting
$\beta_{k+1}(g_k)A(\varphi_{k+1},U_{k+1})$ from $E^{(k+1)}-E^{(k+1)}(1)$, and by replacing the
function $g_k^{-2}(\cdot)$ in the action \dots\ by $g_{k+1}^{-2}(\cdot)$ defined by the equation
(2.24)''. The subtracted counterterm and the change of the function in the action therefore
compensate each other exactly.

(b) The profiles appear in the representation of $\rho_j$ by \cite[(2.23)--(2.25)]{B-CMP119} because
of the coupling constant renormalisation quoted in the proof of (a). At step $i$, $1\le i\le j$, the
counterterm $\beta_iA(\varphi_i,U)$ is subtracted from
$\mathbf{E}^{(i)}(\Lambda_i,g_{i-1},U)-\mathbf{E}^{(i)}(\Lambda_i,g_{i-1},1)$, and the function
$x_{i-1}(\cdot)$ in the action is replaced by $x_i(\cdot)$. By \cite[(2.24)]{B-CMP119}, in the form
$x_{i-1}(x)=x_i(x)+\beta_i\varphi_i(x)$ used in (a), and by linearity of $A$ in the weight,
\begin{equation*}
-A(x_{i-1}(\cdot),U)=-A(x_i(\cdot),U)-\beta_iA(\varphi_i,U).
\end{equation*}
The subtraction therefore lowers the exponent by $\beta_iA(\varphi_i,U)$ and the replacement raises it
by the same amount, so that the exponent, and with it the integrand, is unchanged: the coupling
constant renormalisation changes the way the density is written and not the density. Summed over
$i=1,\dots,j$, these compensations are the identity of (a), by which the two members carrying the
profiles together give back $-x_0A(U)$. At a step at which $\mathbf{R}$ acts, the counterterm members
on the collar of an island before and after the step differ, and by (c), whose proof below does not use
(b), their difference vanishes plaquette by plaquette, so that this change of writing does not change
the density either. Hence the density $\rho_j$ is the same function of the configuration for every
choice of profiles satisfying the two conditions recalled in Lemma~\ref{4.38:lem-coupling-profile};
so are the effective action $A_j$ as a function of the configuration and the integral $\int\rho_j$, and
so is the value of every total, since a total is formed from the density as a whole. The polarisation tensors
$\Pi^{(j)}_{\mathfrak{s}}$ are second variations of a total and therefore agree. The increment
$\beta_j$ is the whole-lattice second-variation coefficient of Definition~\ref{4.2:def-increment} and
so agrees, and with it the coupling $g_j$ of Definition~\ref{1.2:def-couplings}.

(c) At a step at which $\mathbf{R}$ acts, the profiles entering the counterterms on the collar of an
island, $\varphi_j^{\mathrm{pre}}$ among them, are those of the sequence of localisation domains
$\{\Lambda_i(\ell)\}$ in force at that step:
the prepared sequence during the operation, the retained one after it. In the notation of
\cite{B-CMP122b}, where the level at which $\mathbf{R}$ acts is written $k$ (here $k=j$), the
position-dependent coupling
$1/(g_k(\cdot))^2$ entering the completed action $A'_k$ at that step is the function redefined on
\cite[p.~365]{B-CMP122b}, formed with the profiles of the sequence in force at that step; the action
$A'_k$ of the construction here coincides with it.
By linearity of $A$ in the weight, the left side of \eqref{4.38:eq-profile-telescopes-1} equals
$\beta_jA(w,\cdot)$ with
\begin{equation*}
w=-h_z(1-\varphi_j^{\mathrm{pre}})-h_z\varphi_j^{\mathrm{pre}}+h_z=0 ,
\end{equation*}
pointwise, hence on each plaquette separately. This proves \eqref{4.38:eq-profile-telescopes-1}.

(d) The terms $V(Y)$, $R'(X)$, $\mathfrak{r}'(X)$ contain localised pieces of the two profile-bearing
members of (a): restrictions $-A(x_j(\cdot),U)\restriction W$ of the action to regions $W$, and
localised counterterms $-\beta_iA(h_z\varphi_i,U)$ with $\operatorname{supp}h_z\subset W$. The
identity \eqref{4.38:eq-profile-telescopes-a} cancels the profiles only in the sum of these members over
the whole lattice, not piece by piece; this holds on $\mathcal{X}_j(\ell)$ as well. The choice
of profile decides no total, by (b). It decides no increment $\beta_j$ either: in the decomposition
$\beta_j=\beta_j^{\mathfrak{A}}+\beta_j^{\Re}$ of the increment into its two parts, the part
$\beta_j^{\mathfrak{A}}$ involves no island, and $\beta_j^{\Re}$ is a difference of totals, which by
(b) does not depend on the choice of profile.
Finally, by Lemma~\ref{4.38:lem-coupling-profile} (properties of the coupling profile), the value
$\varphi_i(x)$ of the product profile of \eqref{4.38:eq-standing-data-a} is determined by the level-$i$
cells of $[\Lambda_i(\ell)^c]_i$ within distance $a_i/4$ of $x$, where $a_i$ is the side of a level-$i$
cell and $[\,\cdot\,]_i$ denotes the round-up to level-$i$ cells; all of these cells lie in
$[x]_i^{+1}$, the level-$i$ cell of $x$ enlarged by one cell, hence in $W^{+1}$ when $x\in W$. The two
conditions recalled in Lemma~\ref{4.38:lem-coupling-profile} alone admit other profiles, such as a
ramp, anchored on the domain obtained from a localisation domain by removing one layer of cubes, whose
value at $x$ reads cells at two cells' distance from $x$; the product profile does not.
\end{proof}

The torus-independence theorem below is used inside the induction on the level by which clause (0) of Theorem~0 is proved (Theorem~\ref{4.1:thm-clause-zero}). It supplies the site-freeness of the coupling increments and does not assume it.

\begin{theorem}[Torus-independence of the localized terms]\label{4.38:thm-torus-independence}
Assume the following statements of this book:
\begin{itemize}
\item[(1)] independence of the localized terms from the small-field region
(Corollary~\ref{4.4:cor-localised-independence}), together with the interior sufficiency of the
whole-lattice setting (Lemma~\ref{4.4:lem-interior-sufficiency});
\item[(2)] the lemma on the walk expansion of the averaged Laplacian $\mathsf{A}$ on the base
class of data, with its statement that this expansion is determined by local data given on
the footprint of each walk, and with the accompanying lemma on single cells;
\item[(3)] the reflection and permutation covariance of the walk expansion of
$\mathsf{A}$, and, for the operators of $\mathbb{L}_1$ other than $\mathsf{A}$, the covariance
clause of the hypothesis of Definition~\ref{4.4:def-one-family-hypothesis};
\item[(4)] the operator bounds at free current;
\item[(5)] the large-field part of the coupling increment;
\item[(6)] the representation and bound of the large-field terms, together with the statements
used with it: the six-clause format of that representation; the description of the new
large-field region
in \cite[p.~390]{B-CMP122b}; the clause on points of the small-field region near an island
that lies outside the domain considered, an island being a first-class component of the
large-field region (Definition~\ref{4.38:def-non-wrapping}); the statement on the rounding of the
prepared large-field sequence up to cells; a lemma on the format of that representation, used
together with it; and those clauses of a further statement on that representation which concern
its format, the remaining clauses of that statement not being used. We assume also the
convergence and analytic extension of the expansions on the joined localization domains, used
here as stated in \cite[p.~192, (1.70)]{B-CMP122a} and
\cite[p.~377, (1.68)--(1.69)]{B-CMP122b}; the locality of the pieces of the renormalisation
operation is derived in the proof from the definitions of \cite[\S\S A, C]{B-CMP99b} and the
statements of this list;
\item[(7)] the locality of the rules that make or change a label and the clause on the
integrand at free bonds in the interior, both of Definition~\ref{4.38:def-standing-data}.
\end{itemize}

Let $j\ge1$, and assume the following hypothesis $(\mathrm{H}_j)$, which is the form in which the
induction on the level enters (it is not to be confused with the hypotheses (H1)--(H7) of
Theorem~0, Notation~\ref{2.2:not-hypotheses}).
\begin{itemize}
\item[(a)] The numbers $(g_i,\check R_i)_{0\le i<j}$ and $(\beta_i)_{1\le i<j}$ are defined
and site-free, that is, the same at every site $\mathfrak{s}=(m,K')$, and $x_i\ge\gamma^{-2}$
for these $i$.
\item[(b)] The renormalisation steps $1,\dots,j-1$ have been performed with these numbers at
every site $\mathfrak{s}\in\mathcal{S}_{j-1}$
(Definition~\ref{4.38:def-admissible-sites}), with two clauses of the clause set
$\mathfrak{i}(j-1)$ of the induction at level $j-1$: the flow clause $x_{j-1}>0$, and the clause
that the
step is taken at every site of $\mathcal{S}_{j-1}$ with the constant $E$ of
Definition~\ref{4.38:def-admissible-sites} set to $0$.
\item[(c$^{+}$)] Clauses (i), (ii)(a) and (iii) below hold at every level $i<j$ in place of $j$.
Moreover $f^{\mathfrak{A}}:=(f_i^{\mathfrak{A}},\beta_i^{\mathfrak{A}})_{i<j}$, where
$f_i^{\mathfrak{A}}$ is the family $\{E_{\mathfrak{s}}^{(i),\mathfrak{A}}(X;\cdot)\}$ of
$\mathfrak{A}$-marked localized terms at level $i$, belongs to the
class $\mathfrak{F}_{j-1}$ of volume-consistent families; in particular the bound
\cite[(2.29)]{B-CMP119} holds for the $\mathfrak{A}$-terms at every site, this property being
hereditary along the steps. For $\natural\in\{\cdot,\Re\}$ the clause asserted at the sites is
the last sentence of clause (ii)(a) at level $i$, the Euclidean clause.
\item[(d)] At every $\mathfrak{s}\in\mathcal{S}_j^{\flat}$ the whole-lattice step $j$
converges on the domain of \cite[(2.27)(ii)]{B-CMP119}. This step is the last logarithm in
\cite[(3.28)]{B-CMP119}, taken in the case in which the localization domains $\Omega_j$ and
$\Lambda_j$ are both the whole lattice torus, pointed as in
\cite[(3.37)--(3.42)]{B-CMP119} and expanded as in \cite[(3.44)--(3.47)]{B-CMP119} under the
walk-resummation rule $\mathcal{R}$. The radii are absolute, that is, not evaluated at any
coupling, for the $\mathfrak{A}$-terms; for the $\Re$-terms that domain is read as the tube at
the stand-in coupling $\tilde g:=g_{j-1}$ fixed in the induction on the level. Its terms
$E_{\mathfrak{s}}^{(j),\natural}(X;\cdot,z)$,
$X\in D_j(\mathfrak{s})$, $z$ a unit-lattice point of $X$, split, by the marking of the terms
that defines the two parts $\beta_j^{\mathfrak{A}}$ and $\beta_j^{\Re}$ of the coupling increment
(Definition~\ref{4.2:def-increment}),
into the two sub-families $\natural\in\{\mathfrak{A},\Re\}$, are analytic there. They are real
at real $(\mathbf{U},\mathbf{J})$. They are invariant under the
complexified gauge group $G^{c}$, in the sense of \cite[(1.19)]{B-CMP119}, which is
\cite[(2.27)(iii)]{B-CMP119}. Finally, on the polydisc
\begin{equation*}
\bigl\{\mathbf{U}=e^{iA}:\ \sup_X|A|<a^{\natural},\ \sup_X|\mathbf{J}|<a^{\natural}\bigr\}
\end{equation*}
about $(\mathbf{U},\mathbf{J})=(1,0)$, where $A$ denotes a $\mathfrak{g}^{\mathbb{C}}$-valued
lattice field on $X$ and is not the action functional, they
obey
\begin{equation*}
\bigl|E_{\mathfrak{s}}^{(j),\natural}(X;\cdot,z)\bigr|\le E_0^{\natural}e^{-\kappa d_j(X)},
\end{equation*}
with $E_0^{\natural}>0$, $a^{\natural}>0$ and $\kappa\ge2\kappa(d)$ independent of
$\mathfrak{s}$, $X$ and $z$.
\end{itemize}
For $j=1$, hypothesis $(\mathrm{H}_1)$ consists of: $g_0\in(0,\gamma]$; the cell-size
multiplier $\check R_0$ is given and site-free; and (d), with \cite[\S3]{B-CMP119} read at
$k=0$, where, by the lemma of this book
identifying the zeroth step of that analysis, it coincides with \cite[\S1]{B-CMP119}.

Then the following hold.
\begin{itemize}
\item[(i)] Let $\mathfrak{s}\in\mathcal{S}_j^{\flat}$, and let $X\in D_j(\mathfrak{s})$ be
non-wrapping (Definition~\ref{4.38:def-non-wrapping}) with lift $\hat X$. Then the function
$E_{\mathfrak{s}}^{(j)}(X;\cdot,z)$, read on $\hat X$, is intrinsic
(Definition~\ref{4.38:def-intrinsic-trace}) as an analytic function on the domain of (d): it
depends neither on $\mathfrak{s}$ nor on the lift. No operator, minimiser, partition function
or scalar of $\mathbb{T}_j(\mathfrak{s})$ occurs in it; the periodised objects occur only in
its argument. The functions $E^{(j)}(\hat X;\Phi,\hat z)$, $\hat X\in D_j^{\infty}$, so
defined satisfy
\begin{equation*}
E^{(j)}(\hat X+t;\tau_t\Phi,\hat z+t)=E^{(j)}(\hat X;\Phi,\hat z),\qquad t\in M\mathbb{Z}^4.
\end{equation*}
The same holds for $E^{(j),\mathfrak{A}}$ and for $E^{(j),\Re}$.
\item[(ii)] (a) Let $\mathsf{h}_{\mathfrak{s}}$ be the linearised minimiser-and-current map at
the site $\mathfrak{s}$, and let $\mathsf{h}_\infty$ be its walk-expansion limit. For points
$x,y,z$ of the unit lattice of $\mathbb{R}^4$ (read on $\mathbb{T}_j(\mathfrak{s})$ through
$p_{\mathfrak{s}}$ where the torus tensors are concerned), the series
\begin{equation}\label{4.38:eq-torus-independence-a}
\begin{split}
&\Pi_\infty^{(j)}(x,y,z)\\
&\quad:=\sum_{\hat X\in D_j^{\infty},\ \hat X\ni z}
\bigl\langle\partial^2E^{(j)}(\hat X;(1,0),z);\ \mathsf{h}_\infty\delta_x,\
\mathsf{h}_\infty\delta_y\bigr\rangle
\end{split}
\end{equation}
converges absolutely and decays exponentially in $|x-z|+|y-z|$.
Here $\delta_x$ denotes the function on the unit lattice equal to $1$ at $x$ and to $0$ elsewhere.

It is the pointwise limit, as $s_j(\mathfrak{s})\to\infty$ in $\mathcal{S}_j^{\flat}$, of the
torus tensors $\Pi_{\mathfrak{s}}^{(j)}(x,y,z)$. These are the tensors of
\cite[(3.50)]{B-CMP119}, with $E^{(j)}(X,U_j,z)$ replaced by the total
$E_{\mathfrak{s}}^{(j)}(U_j,z)$.

The two-point tensors $\sum_z\Pi_{\mathfrak{s}}^{(j)}(x,y,z)$ of \cite[(1.20)]{B-CMP109}
converge to
\begin{equation*}
\sum_z\Pi_\infty^{(j)}(x,y,z)=:\Pi_\infty^{(j)}(x-y),
\end{equation*}
which satisfies \cite[(5.10)]{B-CMP109}; it is the limit of the
expression \cite[(5.1)]{B-CMP109}.

For every Euclidean transformation $r$ of the unit lattice of $\mathbb{R}^4$,
\begin{equation*}
\Pi_\infty^{(j)}(rx,ry,rz)=(r\otimes r)\,\Pi_\infty^{(j)}(x,y,z);
\end{equation*}
in particular $\Pi_\infty^{(j)}$ is translation invariant.

All of this holds for $\mathfrak{A}$ and for $\Re$ separately. No rate of convergence is part
of the statement.

At each site $\mathfrak{s}$, let $\mathcal{E}^{\natural}_{\mathfrak{s}}(V,z)$ denote the total
of the terms of the sub-family $\natural$ at $\mathfrak{s}$ pointed at $z$ (all terms for
$\natural=\cdot$), evaluated at the minimiser determined by a regular configuration $V$ (a real
regular background, Definition~\ref{4.4:def-regular-backgrounds}). Then for every Euclidean
transformation $r$ of the unit lattice of $\mathbb{T}_j(\mathfrak{s})$,
\begin{equation*}
\mathcal{E}^{\natural}_{\mathfrak{s}}(rV,rz)=\mathcal{E}^{\natural}_{\mathfrak{s}}(V,z),\qquad
\natural\in\{\cdot,\mathfrak{A},\Re\}.
\end{equation*}
This is the Euclidean clause of a volume-consistent family at level $j$.

(b) The number $\beta_j$, defined as the coefficient \cite[(3.61)]{B-CMP119}, that is
\cite[(1.22)]{B-CMP109}, of $\Pi_\infty^{(j)}$, is well defined and real. It equals
$\beta_j^{\mathfrak{A}}+\beta_j^{\Re}$, and it is a function of $(g_0,\dots,g_{j-1})$ and of
the fixed parameters of the construction alone. Hence $x_j$ is site-free. Given $x_j>0$, which
is the flow
clause of the induction at level $j$, the numbers $g_j$, $\check R_j$, $\varepsilon_j$,
$\delta_j$, $k_1(j)$ of Definition~\ref{4.38:def-standing-data} and the class $\mathcal{S}_j$
are also site-free: they are the same at every site, and free of $m$ and $K'$.
\item[(iii)] Assume moreover that $x_j\ge\gamma^{-2}$ and that step $j$ has been completed at
every $\mathfrak{s}\in\mathcal{S}_j$. Then every term or operation of the representation
\cite[(2.18)]{B-CMP119} of $\rho_j$ that is localized in a non-wrapping $W$ is intrinsic
(Definition~\ref{4.38:def-intrinsic-trace}).

This includes the members carrying the coupling profile, namely
\begin{equation*}
-A(x_j(\cdot),\mathbf{U})\restriction W
\quad\text{and}\quad
-\beta_iA(h_z\varphi_i,\mathbf{U})\ \ (z\ \text{with}\ \operatorname{supp}h_z\subset W),
\end{equation*}
of \cite[(2.23)--(2.25)]{B-CMP119} and \cite[(3.65)]{B-CMP119}. Here
$x_j(\cdot):=1/g_j^2(\cdot)$ is the position-dependent inverse coupling of
\cite[(2.24)]{B-CMP119}, $h_z$ is the cut-off function attached to $z$ in these members,
and $\varphi_i$ is the coupling profile of the
label $\ell$
carried by the term. By
Lemma~\ref{4.38:lem-coupling-profile}, for $x$ the barycentre of a plaquette of $W$,
$\varphi_i(x)$ reads only the cells of
$[\Lambda_i(\ell)^c]_i$ within $\tfrac14a_i$ of $x$, and these lie in
$[x]_i^{+1}\subset W^{+1}$. In particular:
\begin{itemize}
\item[(a)] the terms $E^{(j)}(X;\cdot,z)$, $X\subset\Lambda_j$, $z\in\Lambda_j^0$, of every
label at $\mathfrak{s}$ are the functions of clause (i); here $\Lambda_j^0$ is the range of
$z$ in \cite[(2.26)]{B-CMP119};
\item[(b)] let $\ell$ be any label at $\mathfrak{s}$, and let $X\in\mathbf{D}_j$ be any
non-wrapping union of $M$-cubes, unions of cells included, with lift $\hat X$ and with
$X\subset\Lambda_j(\ell)^{\sim-1}$ (the union of the cells of $\Lambda_j(\ell)$ that touch no
cell of $\Lambda_j(\ell)^c$), the range of localization domains of
\cite[(2.30), p.~260]{B-CMP119}. Let $R^{(j)}(X;\cdot)$ be the term of
\cite[(2.30)]{B-CMP119} with localization domain $X$; it is the sum of
$R'^{(j),\ell}(X;\cdot)$, the term with localization domain $X$ in the logarithm of the
polymer expansion \cite[(1.90)]{B-CMP122b} at the label $\ell$, written as the series
\cite[(2.13)]{B-CMP116}, and of the remaining part
$R''^{(j)}(X;\cdot)$. Then
\begin{equation}\label{4.38:eq-torus-independence-b}
\begin{gathered}
R^{(j)}(X;\cdot)=R''^{(j)}(X;\cdot)+R'^{(j),\ell}(X;\cdot)=\mathfrak{r}^{(j)}(\hat X;\cdot),\\
\mathfrak{r}^{(j)}(\hat X;\cdot):=R''^{(j)}(\hat X;\cdot)+\mathfrak{r}'^{(j)}(\hat X;\cdot),\\
B'_{\mathrm{ann}}(X):=R'^{(j),\ell}(X)-\mathfrak{r}'^{(j)}(\hat X),
\end{gathered}
\end{equation}
where $\mathfrak{r}'^{(j)}$ is the series of $R'^{(j)}$ taken over the catalogue of the empty
trace, as in \eqref{4.38:eq-label-remainder-b}, read on the lift in the last line. The functions
$\mathfrak{r}^{(j)}(\hat X;\cdot)$ form one family, independent of $\mathfrak{s}$, of the lift
and of the label, and covariant under $\tau_t$ for
$t\in M\check R_j\mathbb{Z}^4$. The difference $B'_{\mathrm{ann}}(X)$ vanishes whenever the
trace of $\ell$ on $X^{+1}=[X]\cup(\text{the cells touching }[X])$ is empty, that is,
whenever $X$ belongs to the class $\mathcal{X}_j(\ell)$ of domains with empty trace of
\eqref{4.38:eq-label-remainder-b}. Since $\mathcal{X}_j(\ell)$ is exactly the range of
\cite[(2.30)]{B-CMP119}, as is shown later in this section, $B'_{\mathrm{ann}}(X)$ vanishes on the
whole range;
\item[(c)] the objects
\begin{equation*}
B^{(j)}(X;\cdot),\quad \mathbf{B}'^{(j)}(X;\cdot),\quad \mathbf{T}^{(j-1)}(\cdot\cap X)
\quad\text{with}\ W:=X,
\end{equation*}
and $\mathbf{T}'_j(X)$ with $W:=X^{\sharp}$ depend on $(\mathfrak{s},\text{label})$ only
through the trace of the label on $W^{+1}$; here $B^{(j)}(X;\cdot)$ is the term of
\cite[(2.42)]{B-CMP119}, and $X^{\sharp}$ is the strip of $X$
(Definition~\ref{4.38:def-standing-data}) on which the potential term of
\cite[(1.91)]{B-CMP122b} with domain $X^{\sharp}$, contained in $\mathbf{T}'_j(X)$, lives.
\end{itemize}
Consequently, under the additional hypotheses of this clause, and with the flow clause at
level $j$ and the step statement for volume-consistent
families at level $j-1$, parts (a)--(c$^{+}$) of hypothesis $(\mathrm{H}_{j+1})$ hold.
\end{itemize}
\end{theorem}

The proof of Theorem~\ref{4.38:thm-torus-independence}, from the statements (1)--(7) and
hypothesis $(\mathrm{H}_j)$, is given in the remainder of this section.

\begin{lemma}[Torus-independence of the non-wrapping localized terms]\label{4.38:prf-proof-non-wrapping}
Let $j\ge 1$ and put $k:=j-1$. Assume the following.
\begin{enumerate}
\item Hypothesis $(\mathrm{H}_j)$ of Theorem~\ref{4.38:thm-torus-independence}.
\item The walk expansion of the averaged Laplacian $\mathsf{A}=C^{*}\Delta^{(k)}C$ satisfies the
hypothesis on the quadratic form's walk expansion, Definition~\ref{4.4:def-form-walk-hypothesis}.
\item The operators of $\mathbb{L}_1$ other than $\mathsf{A}$ form one walk family in the sense of
Definition~\ref{4.4:def-one-family-hypothesis}.
\item The locality of the rule and agreement on uniform regions, Theorem~\ref{4.4:thm-uniform-agreement},
holds for the whole-lattice settings of any two sites $\mathfrak{s},\mathfrak{s}'$.
\end{enumerate}
Then part~(i) of Theorem~\ref{4.38:thm-torus-independence} holds at level $j$. That is, let
$\mathfrak{s}$ be a site and let $X\subset\mathbb{T}_j(\mathfrak{s})$ be non-wrapping in the sense of
Definition~\ref{4.38:def-non-wrapping}. Consider the whole-lattice setting of $\mathfrak{s}$, in which
$\Lambda_i$ is the whole torus for every $i\le k$. Consider the term of index $X$ of
\cite[(3.47)]{B-CMP119} in this setting, each of the three parts of \cite[(3.37)]{B-CMP119} being pointed
as in \cite[(3.41)]{B-CMP119}. This term is one and the same series in intrinsic ingredients at every
site. Consequently the terms of index $X$ at two sites coincide, under a common lift of $X$, on the whole
domain of clause~(d) of $(\mathrm{H}_j)$. The same holds separately for each of the two marked
sub-families of terms.
\end{lemma}

\begin{proof}
Throughout, $k=j-1$, and $\hat X$ denotes a lift of $X$ in the sense of
Definition~\ref{4.38:def-non-wrapping}.

\emph{The series $\Psi_X$.}
We apply the first step of the localisation lemma for the terms of \cite[(3.47)]{B-CMP119} and the
proposition on the cancellation of the normalising partition function. The $s$-interpolation integrals
are turned into vertex values by the vertex formula for interpolated functionals,
Lemma~\ref{4.4:lem-vertex-formula}. In the whole-lattice setting of $\mathfrak{s}$, the term of index $X$
of \cite[(3.47)]{B-CMP119} is then a series $\Psi_X$; this holds for each of the three parts of
\cite[(3.37)]{B-CMP119}, pointed as in \cite[(3.41)]{B-CMP119}. The series $\Psi_X$ converges by
clause~(d) of $(\mathrm{H}_j)$.

Its index set and its coefficients are functions of the incidences among the standard sub-cubes of $X$.
These are the Ursell and Mayer coefficients, the rules assigning a walk to the hull of its footprint, and
the tent functions $h_z$. The series $\Psi_X$ is a series in the following ingredients.
\begin{itemize}
\item[$(\alpha)$] The operators $K^{W'}$ of \eqref{4.38:eq-proof-non-wrapping-1} below, with
$W'\subset X$. Here $K$ runs over $\mathbb{L}$: the propagators, the minimiser operators,
$\mathsf{A}=C^{*}\Delta^{(k)}C$, its compressions, their resolvents, and the square root.
\item[$(\beta)$] Local densities on $X$.
  \begin{itemize}
  \item The plaquette terms with weight $x_k$.
  \item The monomials in Ba{\l}aban's polynomials $\mathbf{P}^{(k)}$ of \cite{B-CMP119}.
  \item For $X'\subset X$ and $i\le k$, the inherited differences
  $E^{(i)}(X';\cdot,z')-E^{(i)}(X';1,z')-\beta_iA(h_{z'},\cdot)$.
  \item For $X'\subset X$ and $i\le k_1(k)$, the differences $R^{(i)}(X';\cdot)-R^{(i)}(X';1)$, which
  \cite{B-CMP119} includes in $E_k$.
  \end{itemize}
\item[$(\gamma)$] The factors $\chi(|A(b)|<\delta_k/g_k)$ of \cite{B-CMP119}, in which $A(b)$ denotes
Ba{\l}aban's field variable on the bond $b$ (not the action functional $A$), and unit Gaussian
measures on the bonds $b$ of $X$.
\item[$(\delta)$] The $t(Y)$- and $\tau$-parameters attached to the sets $Y\subset X$. Also the one-bond
surface insertion $\partial_t\chi_t$ of the second part of \cite[(3.37)]{B-CMP119}; it arises in the
first step of the localisation lemma named above. This item is distinct from the condition on the free
interior bonds in Definition~\ref{4.38:def-standing-data}.
\end{itemize}
The operators of $(\alpha)$ are
\begin{equation}\label{4.38:eq-proof-non-wrapping-1}
K^{W'}:=\sum_{|\omega|\subset W'}K_{\omega},\qquad K\in\mathbb{L},
\end{equation}
where $K=\sum_\omega K_\omega$ is the walk expansion of $K$ and $|\omega|$ the footprint of the walk
$\omega$ (Definition~\ref{4.4:def-walk-systems}).

In the whole-lattice setting we have, for every $i\le k$,
\begin{equation}\label{4.38:eq-proof-non-wrapping-2}
\Lambda_i=\mathbb{T}_i(\mathfrak{s}),\qquad [\Lambda_i^{c}]_i=\emptyset .
\end{equation}
Hence $\varphi_i\equiv1$: in \eqref{4.38:eq-standing-data-a} the product over the level-$i$ cells of
$[\Lambda_i^{c}]_i$ is empty (Lemma~\ref{4.38:lem-coupling-profile}). So $x_k(\cdot)\equiv x_k$, and the
plaquette terms of $(\beta)$ carry the constant weight $x_k$.

\emph{Each ingredient is intrinsic} in the sense of Definition~\ref{4.38:def-intrinsic-trace}. We take
the four classes in turn.

$(\alpha)$. Let $\mathfrak{s},\mathfrak{s}'$ be two sites. Let $W$ be a common lift of $X$ together with
its collar; it exists by Definition~\ref{4.38:def-non-wrapping}, because $X$ is non-wrapping. We apply
the locality of the walk-expansion rule of Definition~\ref{4.38:def-standing-data} to the whole-lattice
settings of $\mathfrak{s}$ and $\mathfrak{s}'$. We use it in the second case of the locality definition,
Definition~\ref{4.4:def-local-rules}, together with the corollary of that locality for two settings, and
with the agreement hypothesis above.

By this locality, a walk term $K_\omega$ whose footprint lies in $W$ is a word in three kinds of factor:
\begin{itemize}
\item cut-offs of fixed profile;
\item block-local maps;
\item inverses with Dirichlet conditions on cubes concentric with a cover cube, as in
\cite{B-CMP99b}.
\end{itemize}
The operators concerned are those of the one walk family of the one-family hypothesis above. This is the
form in which these operators are written in \cite{B-CMP99b}.

For $\mathsf{A}$ and its compressions, the terms $\mathsf{A}_\omega$ are those of the walk-expansion
hypothesis above, local on their footprints. For their resolvents and the square root, $K_\omega$ is
instead a word in cell resolvents. In the notation of part~(iv) of
Lemma~\ref{4.4:lem-cell-resolvent}, these are cell resolvents on a cell $q$, built from the terms
$\mathsf{A}_\omega$ with $d(\omega)\le r'$ whose footprints satisfy $|\omega|\subset q^{\sharp}$; this is
part~(iv) of that lemma together with Lemma~\ref{4.38:lem-tree-lengths}.

In either case $K_\omega$ is determined by the lattices, partitions and fields on $W$ and by the fixed
data. Since $W$ is a lift common to both sites, $K_\omega$ is the same at both sites. The global operator
$K$ is the sum over all walks, the wrapping ones included, and it is not an ingredient of $\Psi_X$.

$(\beta)$. We use clause~(a) of $(\mathrm{H}_j)$ and the conclusions of
Theorem~\ref{4.38:thm-torus-independence} at the levels $i<j$, which are part of clause~(c$^{+}$) of
$(\mathrm{H}_j)$. They apply because every $X'\subset X$ is non-wrapping at its own level $i$, with
$\hat X'\subset\hat X$ (Definition~\ref{4.38:def-non-wrapping}).

The inherited terms $R^{(i)}(X';\cdot)$ are read at the empty label of the whole-lattice setting. There
$X'\in\mathcal{X}_i(\emptyset)$ holds trivially, and the terms are $\mathfrak{r}^{(i)}(\hat X')$. This
follows from part~(iii)(b) of Theorem~\ref{4.38:thm-torus-independence} at level $i$ and from the
identification of the remainder terms at the empty label.

$(\gamma)$, $(\delta)$ and the coefficients. These are intrinsic by three facts:
\begin{itemize}
\item the numbers of Definition~\ref{4.38:def-standing-data} (every threshold, width, cube size and clock
is a function of $\mathfrak{p}$ and of the couplings already generated);
\item the geometry of that definition (standard partitions, and cut-offs borne by standard cubes that are
fixed profiles of the position relative to the centre);
\item Definition~\ref{4.38:def-non-wrapping}, by which lifts preserve lattices, partitions, incidences,
axes and tree lengths.
\end{itemize}

So the ingredients of $\Psi_X$ at two sites coincide wherever they exist. Hence the identity of the
lemma holds on the whole domain of clause~(d) of $(\mathrm{H}_j)$.

\emph{Translations.} Two lifts of $X$ differ by a translation in $s_j(\mathfrak{s})\mathbb{Z}^4\subset
M\mathbb{Z}^4$ (Definition~\ref{4.38:def-non-wrapping}); we show that every ingredient of $\Psi_X$ is
carried by the translations $\tau_t$, $t\in M\mathbb{Z}^4$. Let $t\in M\mathbb{Z}^4$ and let $\tau_t$ be
translation by $t$ (in this paragraph $t$ denotes a translation vector, not the interpolation
parameters $t$, $\tau$ of $(\delta)$, nor the parameter $t$ of $\langle\cdot\rangle_t$).

At level $k$, $\tau_t$ acts by $LM\mathbb{Z}^4$. It therefore preserves the following partitions: the
$M_1$-cubes; the $\mathsf{m}_4$-cells of the expansion of $\mathsf{A}$, whose side $\mathsf{m}_4$ divides $M$
by the choice of these cells in that expansion; the $L$-blocks; and the $M$- and $LM$-cubes.

At a level $i\le k_1(k)$, $\tau_t$ acts by $ML^{j-i}\mathbb{Z}^4$, and we claim
\begin{equation}\label{4.38:eq-proof-non-wrapping-3}
ML^{j-i}\mathbb{Z}^4\subset M\check{R}_i\mathbb{Z}^4 .
\end{equation}
This rests on the inequality $M\check{R}_i\le L^{j-i}$, equivalently $ML^{i}\check{R}_i\le L^{j}$.

At $i=k_1$ the inequality is the defining property of $k_1$ among the numbers of
Definition~\ref{4.38:def-standing-data}; see \cite{B-CMP119}. The map $i\mapsto L^{i}\check{R}_i$ is
non-decreasing, by the monotonicity property of the cell-size multipliers stated in the lemma on the
multipliers $\check{R}_i$. Hence for $i\le k_1$
\begin{equation*}
ML^{i}\check{R}_i\le ML^{k_1}\check{R}_{k_1}\le L^{j}.
\end{equation*}
Since the $\check{R}_i$ are powers of $L$, by the same lemma, the inequality $\check{R}_i\le L^{j-i}$
means that $\check{R}_i$ divides $L^{j-i}$. This gives \eqref{4.38:eq-proof-non-wrapping-3}.

Thus $\tau_t$ preserves, at level $k$ and at every level $i\le k_1(k)$, the partitions that the
ingredients read. By Definition~\ref{4.38:def-intrinsic-trace}, an intrinsic object is carried by every
translation preserving the partitions it reads, by transport of structure; hence every ingredient of
$\Psi_X$ is carried by $\tau_t$.

\emph{Sub-families.} The marking of the two sub-families $\natural\in\{\mathfrak{A},\Re\}$ depends on
the levels only, by the definition of the marking. The identity therefore holds cluster by cluster, and
in particular for each marked sub-family separately.

\emph{Objects that do not enter $\Psi_X$.} Three kinds of object are not among the ingredients.

(1) The partition function $\Xi(\mathbb{T})$ normalises the expectation
$\langle\cdot\rangle_t$. It is the integral inside the logarithm on the right-hand side of
\cite[(3.28)]{B-CMP119}, and it is extensive in $|\mathbb{T}|$. It cancels cluster by cluster. Let $b$ be
the bond at which the term is pointed, and let $N/D$ and $H'(Z_0)$ be the pointed quotient and the
cluster function of the proposition on the cancellation of the normalising partition function;
$Z_0^{\hat c}$ denotes the set of that name there. Then
\begin{equation}\label{4.38:eq-proof-non-wrapping-4}
\frac{N}{D}=\sum_{Z_0\ni b}H'(Z_0)\,
\exp\bigl[\log\Xi(Z_0^{\hat c})-\log\Xi(\mathbb{T})\bigr].
\end{equation}
This sum is reduced by \cite[(3.45)--(3.46)]{B-CMP119}, through the proposition on the cancellation of
the normalising partition function.

(2) The global minimiser and the map $\mathsf{h}_{\mathfrak{s}}$ are the argument at which $\Psi_X$ is
evaluated. They are not ingredients.

(3) The objects of the infinite lattice $\mathbb{Z}^4\subset\mathbb{R}^4$ enter only by addition and
subtraction.
Fix $i\le k$ and a cover cube $\square$. Let $\square_0=\tilde\square^{5}$ be the cube in which the
determining set of the local minimiser map $H_i(\square_0)$ is built, let $\square^{\sim2}$ be the cube
of \cite[(3.5)]{B-CMP109}, and let $\square^{3}$ be the cube so denoted in the rewriting below. Put
\begin{equation}\label{4.38:eq-proof-non-wrapping-5}
\mathcal{T}_{\square}:=\sum_{X''\subset\square^{\sim2}}E^{(2)}(X'')\bigl[H_i(\square_0)\,\cdot,\,
H_i(\square_0)\,\cdot\,\bigr]
-\beta_i\bigl\langle\partial H_i(\square_0)\,\cdot,\,\zeta_{\square}\,\partial H_i(\square_0)\,
\cdot\,\bigr\rangle .
\end{equation}
The rewriting carried out in \cite{B-CMP109}, and identically in \cite{B-CMP119}, expresses
\eqref{4.38:eq-proof-non-wrapping-5} as
\begin{equation}\label{4.38:eq-proof-non-wrapping-6}
\begin{split}
\mathcal{T}_{\square}
&=(\Pi_{\infty}-\beta_i\Delta_i)[\cdot,\cdot]
-\sum_{\hat X\not\subset\square^{\sim2}}E^{(2)}(\hat X)\bigl[H_{\infty}\,\cdot,\,H_{\infty}\,\cdot\,\bigr]\\
&\quad+\bigl[\text{terms in }H_i(\square_0)-H_{\infty}\bigr]
+\bigl[\text{tails in }y\text{ off }\square^{3}\bigr].
\end{split}
\end{equation}
Here $\Pi_\infty$ is the infinite-volume polarisation form, and $\Delta_i$ and $H_\infty$ are the
quadratic form paired with $\beta_i$ and the infinite-lattice minimiser map in the rewriting of
\cite{B-CMP109}; the last bracket collects the tails of the sums over points $y$ outside $\square^{3}$.
This is an identity between polynomials in the local field on $\square_0$. Every difference is kept on
the right-hand side; as \cite{B-CMP109} states, this change increases the sum of the irrelevant terms.
Nothing is dropped.

The $s$-parameters enter only through the argument $\mathbf{H}_k(s)$, the small solution of
\cite[(1.3)--(1.4)]{B-CMP116}. Moreover
\cite[(1.9)--(1.10)]{B-CMP116} is linear in the function substituted, by
Lemma~\ref{4.4:lem-vertex-formula}. Hence the terms with index $Y_0$ of \cite[(1.9)--(1.10)]{B-CMP116}
on the right-hand side of
\eqref{4.38:eq-proof-non-wrapping-6} add up to the $Y_0$-term of $\mathcal{T}_{\square}$. By
\eqref{4.38:eq-proof-non-wrapping-5}, this term contains only $E^{(i)}(X'')$ with
$X''\subset\square^{\sim2}$, the map $H_i(\square_0)$, and the number $\beta_i$. The objects
$\Pi_\infty$, $\Delta_i$ and $H_\infty$ therefore do not enter $\Psi_X$.

The individual operators of the ingredients $(\alpha)$--$(\delta)$ are listed one by one in the table of
the ingredients of $\Psi_X$ that accompanies this lemma.
\end{proof}

\begin{proof}[Proof of Lemma~\ref{4.38:prf-proof-non-wrapping}, continued (the ingredients of the localized series)]
\label{4.38:prf-series-ingredients}
We go through the objects that occur in the localized series $\Psi_X$. For each object we
say why its occurrence in $\Psi_X$ is intrinsic in the sense of the preceding part of the
proof, and we name the statement or the passage of Ba{\l}aban's papers on which this rests.
Throughout, $K^{W'}=\sum_{|\omega|\subset W'}K_{\omega}$ denotes the sum of the walk terms
of an operator $K$ whose footprint lies in $W'$.
\begin{enumerate}
\item The operators $G'$, $(Q'G'^{2}Q'^{*})^{-1}$, $P$, $R$, $G_0$, $G$. The family used is
the one walk family of Definition~\ref{4.4:def-one-family-hypothesis}, taken on the cells of
the hypothesis of Definition~\ref{4.4:def-form-walk-hypothesis} and together with the
clauses (F1)--(F3) that accompany that family. Ba{\l}aban's walk expansions
\cite[(3.90), (3.98), (3.107)]{B-CMP99b} are expansions in words such as
$h_{\square}G'_{\square}h_{\square}$, $K(h_{\square})G'_{\square}h_{\square}$,
$h_{\square}C_{\square}h_{\square}$, \dots, that is, in operators attached to the sequence
$\{\Omega_n(\square)\}$ of cubes concentric with $\square$, with Dirichlet conditions off
$\Omega_0(\square)\subset\tilde\square^{5}$
(Lemma~\ref{4.4:lem-walk-expansion-locality}). These expansions are used only in the
treatment of the map $H_{\mathfrak{s}}$ at the trivial background $U=1$.
\item The operators $(QGQ^{*})^{-1}$, $H$, $H_0$, $\tilde G$, $G_1$, $H_1$,
$\Delta^{(2)}$, $\Delta^{(k)}$ and $\mathsf{A}$. They belong to the same family
(Definition~\ref{4.4:def-one-family-hypothesis}): they are obtained from the operators of
the first item by products, inverses, sums and branchings, written in the words of the
first item, and by clause (f3) of that family, in the notation fixed there, the decay
constant of each walk term so obtained equals that of the term it comes from,
$\kappa_{r\omega}=\kappa_{\omega}$. The averaged Laplacian
$\mathsf{A}=C^{*}\Delta^{(k)}C$ is the form $\mathsf{Q}$ of the hypothesis of
Definition~\ref{4.4:def-form-walk-hypothesis}. The random walk expansions of these
operators in Ba{\l}aban's construction are again used only in the treatment of
$H_{\mathfrak{s}}$ at $U=1$.
\item The operators $C^{(k)}(Z_0)$, $C^{(k)}(\Lambda)$, $(C^{(k)})^{1/2}$,
$(x+\mathsf{A})^{-1}$ and $(\mathsf{A}+\lambda)^{-1}$. They enter through their
cell-resolvent walks: the local piece $(x+T_q)^{-1}$ on $q^{*}$ is built from the walk
terms $\mathsf{A}_{\omega}$ with $d(\omega)\le r'$ and $|\omega|\subset q^{\sharp}$, in the
notation of Lemma~\ref{4.4:lem-cell-resolvent}; see the construction of these local pieces
in \cite{B-CMP119}, clause (iv) of Lemma~\ref{4.4:lem-cell-resolvent}, and the corollary on
the local pieces of the cell resolvents that accompanies it.

The operators of this item must be taken through the rule of
Definition~\ref{4.4:def-expansion-rule}, and not as
they stand, for the following reason. Un-expanded, they are periodised: with
$\mathsf{A}_{\mathbb{T}}$ the operator on the torus $\mathbb{T}_j(\mathfrak{s})$ of side
$s_j(\mathfrak{s})$ and $\mathsf{A}_{\infty}$ the operator on $\mathbb{Z}^4$,
\begin{equation}\label{4.38:eq-series-ingredients-1}
\chi_{Z_0}\mathsf{A}_{\mathbb{T}}\chi_{Z_0}
=\chi_{Z_0}\Bigl[\mathsf{A}_{\infty}
+\sum_{n\neq0}\mathsf{A}_{\infty}(\cdot,\cdot+s_j(\mathfrak{s})n)\Bigr]\chi_{Z_0},
\end{equation}
and $\mathsf{A}_{\mathbb{T}}$ at the cut background
$(\mathbf{U},\mathbf{J})\restriction Z\oplus(1,0)$ is a function of
$(\mathbf{U},\mathbf{J})\restriction Z$ which still depends on $s_j(\mathfrak{s})$. The
background locality stated in \cite{B-CMP116}, ``$H(Z)$ is localized in the interior of $Z$
with respect to the external gauge fields'', concerns the dependence on the external gauge
fields and does not by itself exclude this dependence on $s_j(\mathfrak{s})$. The operators
on two tori would agree exactly if the image part in \eqref{4.38:eq-series-ingredients-1},
the sum over $n\neq0$, vanished on $Z_0$; no such vanishing is used here.
\item The decoupling parameters $s(\Delta)$. They enter through vertex
values $K\mapsto K^{W'}$ with $W'\subset X$, by the vertex formula
(Lemma~\ref{4.4:lem-vertex-formula}) applied to the interpolation of
Definition~\ref{4.4:def-interpolation-rule}; a walk carries the parameter $s(\Delta)$ for
``all cubes from $\sigma_0$ which intersect this localization domain'' \cite{B-CMP116}.
\item The objects $\mathbf{H}_k(B')$, $D$, $A_0$. They are the small solutions of the
equations \cite[(1.3)--(1.4)]{B-CMP116}, which, in the words of \cite{B-CMP116},
``separate equations in components''.
\item The minimisers $U_j(\square_0,\cdot)$ and $U^{c}_j(X)$. They are the minimisers
for the determining set built in $\square_0=\tilde\square^{5}$, respectively in $X$
(``We consider $X$ as $\Omega_0$''), as in \cite{B-CMP109}.
\item The objects written $Q$, $C$, $h$, $\tilde D$, $b_0(c)$, $z^{(k)}(c)$ in
\cite{B-CMP109}. They are block-local algebra; ``the operator $C$ is almost local''
\cite{B-CMP109}; see also \cite{B-CMP99b}.
\item The functions $\zeta_{\square}$, $h_{\square}$, $h_z$. They are fixed profiles of
$x$ minus the centre. The coupling profile satisfies $\varphi_i\equiv1$ in $\Psi_X$, since
on the whole lattice it is the empty product; at labels it is the product profile
\eqref{4.38:eq-standing-data-a}, whose properties are those of
Lemma~\ref{4.38:lem-coupling-profile}. See \cite{B-CMP109} and
\cite[(3.40); (2.24), p.~259]{B-CMP119}.
\item The factors $\chi^{(k)}$, the signs of \cite[(2.3)]{B-CMP116}, and $d\mu_0$. They
are products over bonds; off $Z$ the Gaussian integrates to $1$.
\item The Mayer, polymer and cluster sums. Their index sets are families of sub-domains of
$X$, and $\zeta(Z,Z')$ is an incidence \cite{B-CMP116}. The partition function
$\Xi(\mathbb{T})$ cancels cluster by cluster, as in the preceding part of the proof.
\item The inherited terms, the numbers $\beta_i$ and the cell-size multipliers
$\check R_i$. They are intrinsic by the statement being proved, at the levels $i<j$
(induction on the level), and by clause (a) of the hypothesis $(H_j)$.
\item The objects on $\mathbb{Z}^4$. They enter by addition and subtraction, as in the
preceding part of the proof.
\end{enumerate}
The following do not occur in $\Psi_X$: the global propagators, which are sums over all
walks; the global minimiser and $\mathsf{h}_{\mathfrak{s}}$, which are the argument;
the partition function $\Xi(\mathbb{T})$; and the constants $E$, $E_k$.
\end{proof}

\begin{lemma}[Infinite-volume limit of the polarisation tensor]
\label{4.38:prf-polarisation-limit}
Let $j\ge 1$. For a site $\mathfrak{s}$, let $H_{\mathfrak{s}}$ be the operator $H$ of
\cite[(3.126)]{B-CMP99b} at $U=1$, with all the domains $\Omega_i$ equal to
$\mathbb{T}_j(\mathfrak{s})$, and let $\mathsf{D}$ be the finite-range, translation-invariant
linearisation of the current map of \cite[(1.8)]{B-CMP109}. Assume the following.
\begin{itemize}
\item Hypothesis $(\mathrm{H}_j)$ of Theorem~\ref{4.38:thm-torus-independence}.
\item Equality of walk terms on two tori. For two sites
$\mathfrak{s},\mathfrak{s}'\in\mathcal{S}_j^{\flat}$ and every cube $\hat\square\subset\mathbb{R}^4$
of side less than $\frac12\min\{s_j(\mathfrak{s}),s_j(\mathfrak{s}')\}$, the terms of the walk
expansions of $H_{\mathfrak{s}}$ and $H_{\mathfrak{s}'}$ given by
\cite[Theorems~3.10 and~3.12]{B-CMP99b} that are indexed by walks with footprint in
$p_{\mathfrak{s}}\hat\square$, respectively in $p_{\mathfrak{s}'}\hat\square$, and that correspond
to each other under the lifts of Definition~\ref{4.38:def-non-wrapping}, are equal.
\item Part~(b) of the covariance lemma for the closed expression
\cite[(3.37), (3.41)]{B-CMP119} of the total localized term at a background configuration:
under its three hypotheses (the first is the fifth of the conditions defining the scheme
$\mathcal{M}$ of the step, together with the attachment condition of the setting; the second
consists of the bounds \cite[(2.29), (2.32)]{B-CMP119}; the third is stated with that lemma),
that total is covariant under the Euclidean transformations of
the unit lattice.
\item For reflections: the walk expansions of the operators of $\mathbb{L}_1$ other than
$\mathsf{A}$, in the family of Definition~\ref{4.4:def-one-family-hypothesis}, are covariant
under the reflections of the unit lattice, and clause (f3) of that definition holds.
\end{itemize}
Let $\natural\in\{\cdot,\mathfrak{A},\Re\}$; write $E^{(j),\natural}_{\mathfrak{s}}(X;\cdot,z)$
for the corresponding pointed localized terms at the site $\mathfrak{s}$ (with
$E^{(j),\cdot}_{\mathfrak{s}}=E^{(j)}_{\mathfrak{s}}$), and
$\Pi^{(j),\natural}_{\mathfrak{s}}$, $\Pi^{(j),\natural}_{\infty}$ for the tensors built from
them as in Theorem~\ref{4.38:thm-torus-independence}. Let $E^{\natural}$ and $a^{\natural}$ be
the bound and the radius of analyticity that clause~(d) of $(\mathrm{H}_j)$ provides for
these terms. For a site $\mathfrak{s}$, let $\mathsf{h}_{\mathfrak{s}}$ be the linearisation at
$B=0$ of the map $B\mapsto(U_j,J_j)(e^{iB})$ on $\mathbb{T}_j(\mathfrak{s})$, where $U_j(V)$ is
the minimal configuration over $V$ and $J_j(V)$ its current. Then for every
$\mathfrak{s}\in\mathcal{S}_j^{\flat}$ (see \eqref{4.38:eq-admissible-sites-b}) and all points
$x,y,z$ of the unit lattice of $\mathbb{T}_j(\mathfrak{s})$,
\begin{equation}\label{4.38:eq-polarisation-limit-a}
\Pi^{(j)}_{\mathfrak{s}}(x,y,z)
=\sum_{X\in D_j(\mathfrak{s}),\,X\ni z}
\bigl\langle \partial^2 E^{(j)}_{\mathfrak{s}}(X;(1,0),z);\,
\mathsf{h}_{\mathfrak{s}}\delta_x,\,\mathsf{h}_{\mathfrak{s}}\delta_y\bigr\rangle ,
\end{equation}
whatever representative of the minimal orbit is used to define $\mathsf{h}_{\mathfrak{s}}$, and
likewise with $E^{(j),\natural}_{\mathfrak{s}}$ and $\Pi^{(j),\natural}_{\mathfrak{s}}$. Moreover,
for each $\natural$:
\begin{itemize}
\item[(a)] $\mathsf{h}_{\mathfrak{s}}=(H_{\mathfrak{s}},\mathsf{D}H_{\mathfrak{s}})$, and there
are $C_{\mathsf{h}}=C_{\mathsf{h}}(j)$ and $\delta_0>0$ with
$|\mathsf{h}_{\mathfrak{s}}\delta_x(b)|\le C_{\mathsf{h}}e^{-\delta_0\operatorname{dist}(x,b)}$
for every site $\mathfrak{s}\in\mathcal{S}_j^{\flat}$ and every bond $b$; for
$\hat x\in\mathbb{Z}^4$ and $x=p_{\mathfrak{s}}\hat x$, as
$s_j(\mathfrak{s})\to\infty$ in $\mathcal{S}_j^{\flat}$, the lifts of
$\mathsf{h}_{\mathfrak{s}}\delta_x$ converge on every finite set to
$\mathsf{h}_{\infty}\delta_{\hat x}$, the sum over the walks of $\mathbb{R}^4$, which obeys the
same bound.
\item[(b)] With $\delta':=\tfrac14\min\{\delta_0,\kappa/M\}$ and a constant $C$ free of
$\mathfrak{s}$, in torus distances,
\begin{equation*}
|\Pi^{(j),\natural}_{\mathfrak{s}}(x,y,z)|\le C e^{-\delta'(|x-z|+|y-z|)} .
\end{equation*}
\item[(c)] For all $\hat x,\hat y,\hat z\in\mathbb{Z}^4$, as $s_j(\mathfrak{s})\to\infty$ in
$\mathcal{S}_j^{\flat}$,
\begin{equation*}
\Pi^{(j),\natural}_{\mathfrak{s}}(p_{\mathfrak{s}}\hat x,p_{\mathfrak{s}}\hat y,p_{\mathfrak{s}}\hat z)
\longrightarrow \Pi^{(j),\natural}_{\infty}(\hat x,\hat y,\hat z),
\end{equation*}
the series \eqref{4.38:eq-torus-independence-a} converging absolutely and obeying the bound
of~(b) with distances in $\mathbb{R}^4$.
\item[(d)] On minimal configurations only: for every Euclidean transformation $r$ of the unit
lattice of $\mathbb{R}^4$ and every regular $V$, the function
\begin{equation*}
\Phi^{\natural}_{\mathfrak{s}}(V,z):=\sum_{X\ni z}
E^{(j),\natural}_{\mathfrak{s}}\bigl(X;(U_j,J_j)(V),z\bigr)
\end{equation*}
satisfies $\Phi^{\natural}_{\mathfrak{s}}(rV,rz)=\Phi^{\natural}_{\mathfrak{s}}(V,z)$; hence
\begin{equation*}
\Pi^{(j),\natural}_{\mathfrak{s}}(rx,ry,rz)=(r\otimes r)\Pi^{(j),\natural}_{\mathfrak{s}}(x,y,z),
\end{equation*}
and the same holds for $\Pi^{(j),\natural}_{\infty}$.
\item[(e)] For $\hat x,\hat y\in\mathbb{Z}^4$ and $x=p_{\mathfrak{s}}\hat x$,
$y=p_{\mathfrak{s}}\hat y$, as $s_j(\mathfrak{s})\to\infty$ in $\mathcal{S}_j^{\flat}$,
\begin{equation*}
\sum_z\Pi^{(j),\natural}_{\mathfrak{s}}(x,y,z)\longrightarrow
\sum_{\hat z\in\mathbb{Z}^4}\Pi^{(j),\natural}_{\infty}(\hat x,\hat y,\hat z),
\end{equation*}
and the limit satisfies the bound \cite[(5.10)]{B-CMP109}, with the decay constant
$\delta_1$ appearing there.
\end{itemize}
\end{lemma}

\begin{proof}
\emph{The representation.} The first derivative at $(1,0)$ of a $G^{c}$-invariant analytic
function of $(\mathbf{U},\mathbf{J})$ vanishes in both slots; hence its second derivative
annihilates the linearised gauge orbit $(d\lambda,0)$ \cite[(4.13)--(4.15)]{B-CMP109}; of the
relations \cite[(4.15)]{B-CMP109} only the first is used. The terms
$E^{(j),\natural}_{\mathfrak{s}}(X;\cdot,z)$ are such
functions, so \cite[(4.35)]{B-CMP109} gives \eqref{4.38:eq-polarisation-limit-a} at every
$\mathfrak{s}\in\mathcal{S}_j^{\flat}$. The sum over $X\ni z$ of these terms, evaluated at
$(U_j,J_j)(e^{iB})$, does not depend on the representative of the minimal orbit, the terms
being gauge invariant. In its second derivative at $B=0$, computed by the chain rule along an
analytic representative, the contribution of the second derivative of the representative is
multiplied by the first derivative at $(1,0)$, which vanishes; hence
\eqref{4.38:eq-polarisation-limit-a} holds with the linearisation $\mathsf{h}_{\mathfrak{s}}$
of any analytic representative of the minimal orbit.

\emph{(a) The argument converges.} With $H_{\mathfrak{s}}$ and $\mathsf{D}$ as in the statement
(so that $H_{\mathfrak{s}}\in\mathbb{L}_1$), we have
$\mathsf{h}_{\mathfrak{s}}=(H_{\mathfrak{s}},\mathsf{D}H_{\mathfrak{s}})$. At $U=1$
the regularity hypotheses of \cite[Theorems~3.10 and~3.12]{B-CMP99b} hold trivially; these
theorems give $H_{\mathfrak{s}}=\sum_\omega H_\omega$ with the termwise bound
\cite[(3.108)]{B-CMP99b}, whose constant depends, as stated there, on $d$ and $L$ only.
Therefore the walk series converges uniformly in $\mathfrak{s}$, and
$|\mathsf{h}_{\mathfrak{s}}\delta_x(b)|\le C_{\mathsf{h}}e^{-\delta_0\operatorname{dist}(x,b)}$ with
$C_{\mathsf{h}}=C_{\mathsf{h}}(j)$. By the second hypothesis, two sites share, with equal terms,
all walks whose footprints lie in a common lifted cube; this is the cancellation step of the
proof of \cite[Theorem~3.14]{B-CMP99b}, applied to two tori. The remaining walks between two
fixed points form a tail of the uniformly convergent series. Hence, for fixed $\hat x$ and a
fixed finite set, the lifts of $\mathsf{h}_{\mathfrak{s}}\delta_x$ on that set form a Cauchy
family as $s_j(\mathfrak{s})\to\infty$, and they converge to $\mathsf{h}_{\infty}\delta_{\hat x}$,
the sum over the walks of $\mathbb{R}^4$; this is the operator that \cite{B-CMP109} calls $H_j$
with free boundary conditions, and it obeys the same bound. Since $\mathsf{h}_{\mathfrak{s}}$ is
defined without reference to any expansion, so is its limit. The periodicity of the torus
enters the argument only here and through the wrapping domains in~(c). The walk expansions of
\cite[(3.90), (3.98), (3.107)]{B-CMP99b} are used in part~(a) only. Outside part~(a), wherever
termwise identities or reflections of walk terms of the operators of $\mathbb{L}_1$ other than
$\mathsf{A}$ are used, as through the fourth hypothesis in part~(d), the walk family is the one
of Definition~\ref{4.4:def-one-family-hypothesis}.

\emph{(b) Domination.} Clause~(d) of $(\mathrm{H}_j)$ gives analyticity of
$E^{(j),\natural}_{\mathfrak{s}}(X;\cdot,z)$ on the fields whose deviation from $(1,0)$ on $X$
is below $a^{\natural}$, with bound $E^{\natural}e^{-\kappa d_j(X)}$. Put
$d:=d_j(X)$, $d_x:=\operatorname{dist}(x,X)$, $d_y:=\operatorname{dist}(y,X)$. By~(a) the sup norms
on $X$ of $\mathsf{h}_{\mathfrak{s}}\delta_x$ and $\mathsf{h}_{\mathfrak{s}}\delta_y$ are at most
$C_{\mathsf{h}}e^{-\delta_0d_x}$ and $C_{\mathsf{h}}e^{-\delta_0d_y}$. Cauchy's estimate for the
mixed second derivative at $s=t=0$ of the function
\begin{equation*}
(s,t)\mapsto E^{(j),\natural}_{\mathfrak{s}}\bigl(X;(1,0)+s\,\mathsf{h}_{\mathfrak{s}}\delta_x
+t\,\mathsf{h}_{\mathfrak{s}}\delta_y,\,z\bigr),
\end{equation*}
on the polydisc where $s$ and $t$ times these norms are
below $a^{\natural}/2$, bounds each summand of \eqref{4.38:eq-polarisation-limit-a} by
\begin{equation}\label{4.38:eq-polarisation-limit-1}
4E^{\natural}(a^{\natural})^{-2}C_{\mathsf{h}}^2\,e^{-\kappa d}\,e^{-\delta_0(d_x+d_y)}
=:C_{\mathrm{sum}}\,e^{-\kappa d-\delta_0(d_x+d_y)} .
\end{equation}
By \cite[(0.26)]{B-CMP109} and the counting lemma for localization domains,
$\sum_{X\ni z}e^{-\frac12\kappa d_j(X)}\le K_0$ on every torus and on $\mathbb{R}^4$. For
$z\in X$,
\begin{equation*}
|x-z|\le \operatorname{dist}(x,X)+Md_j(X)+2c_dM ,
\end{equation*}
where $c_d$ is the diameter of a unit cube in the norm used; the factor $2$ is needed, since
the sixteen $M$-cubes around a vertex form a domain with $d_j=0$. Since
$\delta'\le\delta_0/4$ and $M\delta'\le\kappa/4$,
\begin{align*}
\delta'(|x-z|+|y-z|)
&\le \delta'(d_x+d_y)+2M\delta'd+4c_dM\delta'\\
&\le \delta_0(d_x+d_y)+\tfrac12\kappa d+4c_dM\delta' ,
\end{align*}
so that
\begin{equation*}
e^{-\kappa d-\delta_0(d_x+d_y)}\le e^{4c_dM\delta'}\,e^{-\delta'(|x-z|+|y-z|)}\,
e^{-\frac12\kappa d} .
\end{equation*}
Summing over $X\ni z$ gives~(b), in torus distances, with
$C:=C_{\mathrm{sum}}K_0e^{4c_dM\delta'}$, which is free of $\mathfrak{s}$.

\emph{(c) Limit.} Fix $\hat x,\hat y,\hat z\in\mathbb{Z}^4$ and put
$x=p_{\mathfrak{s}}\hat x$, $y=p_{\mathfrak{s}}\hat y$, $z=p_{\mathfrak{s}}\hat z$. By
Definition~\ref{4.38:def-non-wrapping}, lifting through $\hat z$ is a bijection from the
non-wrapping $X\in D_j(\mathfrak{s})$ with $X\ni z$ onto the $\hat X\in D_j^{\infty}$ with
$\hat X\ni\hat z$ lying in a closed cube of side at most $\frac14 s_j(\mathfrak{s})$, and
$d_j(X)=d_j(\hat X)$. For these $X$, by
Lemma~\ref{4.38:prf-proof-non-wrapping}, the form
$\partial^2E^{(j),\natural}_{\mathfrak{s}}(X;(1,0),z)$, transported by the lift, is
$\partial^2E^{(j),\natural}(\hat X;(1,0),\hat z)$, a fixed form on a finite-dimensional space,
and its arguments, the lifts of $\mathsf{h}_{\mathfrak{s}}\delta_x$ and
$\mathsf{h}_{\mathfrak{s}}\delta_y$, converge to $\mathsf{h}_{\infty}\delta_{\hat x}$ and
$\mathsf{h}_{\infty}\delta_{\hat y}$ by~(a). By \eqref{4.38:eq-polarisation-limit-1} each such
term is bounded by $C_{\mathrm{sum}}e^{-\kappa d_j(\hat X)}$, which is summable over
$\hat X\ni\hat z$ in $D_j^{\infty}$ with sum at most $C_{\mathrm{sum}}K_0$. A wrapping $X$ has
$d_j(X)>s_j/4M-2$ by \eqref{4.38:eq-non-wrapping-a}, so by
\eqref{4.38:eq-polarisation-limit-1} its term is at most
$C_{\mathrm{sum}}e^{-\frac12\kappa(s_j/4M-2)}e^{-\frac12\kappa d_j(X)}$, and all wrapping $X\ni z$
together contribute at most $C_{\mathrm{sum}}K_0e^{-\frac12\kappa(s_j/4M-2)}$, which tends to
zero. By dominated convergence, $\Pi^{(j),\natural}_{\mathfrak{s}}(x,y,z)$ tends to the series
\eqref{4.38:eq-torus-independence-a}; this series converges absolutely, and the argument of~(b)
applies to it verbatim on $\mathbb{R}^4$, with $\mathsf{h}_{\infty}$ in place of
$\mathsf{h}_{\mathfrak{s}}$ (same bound by~(a)) and $K_0$ on $\mathbb{R}^4$, so it keeps the
decay of~(b).

\emph{(d) Symmetry, on minimal configurations only.} Every Euclidean transformation $r$ of the
unit lattice of $\mathbb{R}^4$ normalises the deck group $s_j\mathbb{Z}^4$, because its linear
part is a signed permutation; so $r$ descends to every $\mathbb{T}_j(\mathfrak{s})$. There, for
regular $V$, $\Phi_{\mathfrak{s}}(rV,rz)=\Phi_{\mathfrak{s}}(V,z)$ \cite[(3.58)]{B-CMP119}, for
the following reasons. At the background $(U_j,J_j)(V)$ the total is the closed expression
\cite[(3.37), (3.41)]{B-CMP119}, by \cite[(1.9)]{B-CMP119}, by clause~(c1) of
Proposition~\ref{4.4:prop-rule-properties}, and by clause~(d) of $(\mathrm{H}_j)$. This
expression is covariant by the third hypothesis. The first hypothesis of the covariance lemma
is the fifth of the conditions defining the scheme $\mathcal{M}$ of the step, together with
the attachment condition, both part of the setting in which the step is taken; its third
hypothesis holds in this setting with no further input; its second
hypothesis is the pair of bounds \cite[(2.29), (2.32)]{B-CMP119}, which hold at $\mathfrak{s}$
by clause~(b) of $(\mathrm{H}_j)$, together with the merging criterion of \cite{B-CMP119}.
When $r$ is a reflection, the bound \cite[(2.32)]{B-CMP119} at reflections requires in
addition the fourth hypothesis: it follows from the covariance of the walk expansions under
reflections together with clause (f3) of Definition~\ref{4.4:def-one-family-hypothesis}.
Further, $r$ maps the minimal orbit over $V$ onto the
minimal orbit over $rV$: the Wilson action satisfies $A(rU)=A(U)$, the composite
block-averaging map to level $j$ commutes with $r$, and the minimal orbit over a given regular
configuration is unique among regular configurations \cite[Theorem~1]{B-CMP102b}.
Differentiating twice at $V=1$ gives
$\Pi^{(j)}_{\mathfrak{s}}(rx,ry,rz)=(r\otimes r)\Pi^{(j)}_{\mathfrak{s}}(x,y,z)$
\cite[(3.59)]{B-CMP119}. Letting $s_j\to\infty$ and using~(c) gives the same identity for
$\Pi^{(j)}_{\infty}$.

No symmetry is claimed off the minimal configurations. There the extension of
\cite{B-CMP109} is built on the cover by cubes of side $2M$: on each cube $\square_0$ of this
cover the configuration is replaced by one defined on $\square_0$ and extended as equal to $U$
outside $\square_0$, and unit translations do not preserve this cover. The torus-independence
of non-wrapping terms alone
gives invariance only under translations by $t\in M\mathbb{Z}^4$; in \cite{B-CMP119} the
invariance is stated for functions with the unrestricted summation.

For $\natural=\mathfrak{A}$: by the proposition expressing the total of the terms
$E^{(j),\mathfrak{A}}_{\mathfrak{s}}$ as the same closed expression on the unmarked family,
which holds under the hypothesis stated with it, this total is covariant by
the same lemma. Its second hypothesis for this
family, the bound \cite[(2.29)]{B-CMP119} for $E^{(i),\mathfrak{A}}$, $i<j$, at regular
minimisers at the site, is the Euclidean condition contained in the membership
$f^{\mathfrak{A}}\in\mathfrak{F}_{j-1}$ of clause~$(\mathrm{c}^+)$ of $(\mathrm{H}_j)$. For
$\natural=\Re$, the statement follows as the difference of the two preceding ones.

\emph{(e) Two-point tensor.} Apply \cite[(2.26)]{B-CMP119} with $\Lambda_j=\mathbb{T}_j(\mathfrak{s})$,
and use the translation covariance in~(d) at the site with the translation by $-z$, writing
$w=y-z$:
\begin{equation}\label{4.38:eq-polarisation-limit-2}
\sum_z\Pi^{(j),\natural}_{\mathfrak{s}}(x,y,z)
=\sum_{w}\Pi^{(j),\natural}_{\mathfrak{s}}(x-y+w,w,0),
\end{equation}
with $w$ ranging over the box $-\frac12 s_j<w_\mu\le\frac12 s_j$, where the torus distance from
$w$ to $0$ is $|w|$. By~(b) the summand on the right of \eqref{4.38:eq-polarisation-limit-2} is
dominated by $Ce^{-\delta'|w|}$, which is summable over $\mathbb{Z}^4$ uniformly in
$\mathfrak{s}$; by~(c) it converges for each fixed $w$. By dominated convergence the right-hand
side of \eqref{4.38:eq-polarisation-limit-2}, for $x=p_{\mathfrak{s}}\hat x$ and
$y=p_{\mathfrak{s}}\hat y$, converges to
\begin{equation*}
\sum_{w\in\mathbb{Z}^4}\Pi^{(j),\natural}_{\infty}(\hat x-\hat y+w,w,0),
\end{equation*}
which equals
$\sum_{\hat z\in\mathbb{Z}^4}\Pi^{(j),\natural}_{\infty}(\hat x,\hat y,\hat z)$ by the
translation invariance of
$\Pi^{(j),\natural}_{\infty}$ from~(d) \cite[(3.63)]{B-CMP119}. The bound
\cite[(5.10)]{B-CMP109}, with Ba{\l}aban's $\delta_1$, for this limit is the content of the
lemma on pointed polarisation sums, applied to the pointed terms over $X\ni z$.
\end{proof}

\begin{lemma}[The coupling increment does not depend on the site]\label{4.38:prf-increment-site-free}
Let $j\ge 1$. The following is one step, at level $j$, of the proof of
Theorem~\ref{4.38:thm-torus-independence}. Assume the following.
\begin{itemize}
\item[(1)] Hypothesis $(\mathrm{H}_j)$ of Theorem~\ref{4.38:thm-torus-independence}, with its
clauses (a), (b), (c$^{+}$) and (d).
\item[(2)] The conclusion at level $j$ of Lemma~\ref{4.38:prf-polarisation-limit} (infinite-volume
limit of the polarisation tensor). In it, the pointed polarisation tensor
$\Pi^{(j)}_{\infty}$ of \eqref{4.38:eq-torus-independence-a} is the infinite-volume limit of the
tensors $\Pi^{(j)}_{\mathfrak{s}}$ of \eqref{4.38:eq-polarisation-limit-a}.
\item[(3)] The tensor \cite[(3.57)]{B-CMP119} reduces, under the reflections and permutations of
the coordinate axes, to the single coefficient \cite[(3.61)]{B-CMP119}, and this coefficient
coincides with \cite[(1.22)]{B-CMP109} through \cite[(3.62)--(3.64)]{B-CMP119}. These steps use
only translation invariance, the first identity of \cite[(4.15)]{B-CMP109},
\cite[(3.63)]{B-CMP119} and the decay of the tensor.
\end{itemize}
Let $\beta_j$ be the coefficient \cite[(3.61)]{B-CMP119}, equivalently \cite[(1.22)]{B-CMP109},
of $\Pi^{(j)}_{\infty}$. Then $\beta_j$ is well defined and real. It satisfies
$\beta_j=\beta_j^{\mathfrak{A}}+\beta_j^{\Re}$. It is a function of $(g_0,\dots,g_{j-1})$ and of
$\mathfrak{p}$ alone. Hence $x_j:=x_{j-1}-\beta_j$ is site-free. If moreover $x_j>0$, which is
clause $(\beta_j)$ of the induction's clause set $\mathfrak{i}(j)$, then
the numbers $g_j$, $\check{R}_j$, $\varepsilon_j$, $\delta_j$, $k_1(j)$ of clause (L1) of
Definition~\ref{4.38:def-standing-data} and the site class $\mathcal{S}_j$ are site-free as well.
That is, they are the same at every site $\mathfrak{s}=(m,K')$, and they do not depend on $m$ or
$K'$.
\end{lemma}

\begin{proof}
\emph{Independence of the pointing.} Clause (d) of $(\mathrm{H}_j)$ supplies the translation
invariance on which the corresponding step of \cite{B-CMP119} rests. There the step reads:
``This function is translation invariant, hence $\Pi^{(j)}_{\mu\nu,\kappa\lambda}$ is
independent of $z$.'' The same conclusion therefore holds for $\Pi^{(j)}_{\infty}$: its
components $\Pi^{(j)}_{\infty,\mu\nu,\kappa\lambda}$ do not depend on the point $z$ at which the
tensor is pointed.

\emph{Reduction to one coefficient.} The reflections and permutations of the coordinate axes used
in \cite{B-CMP119} reduce the tensor \cite[(3.57)]{B-CMP119} to the single coefficient
\cite[(3.61)]{B-CMP119}. By \cite[(3.62)--(3.64)]{B-CMP119}, this coefficient coincides with
\cite[(1.22)]{B-CMP109}. These steps use only translation invariance, the first identity of
\cite[(4.15)]{B-CMP109}, \cite[(3.63)]{B-CMP119} and decay, as stated in hypothesis (3).
Translation invariance is available for $\Pi^{(j)}_{\infty}$ by the previous paragraph. Decay
enters only through the second moments of the translation-invariant kernel of
$\Pi^{(j)}_{\infty}$ on which \cite[(3.62)--(3.64)]{B-CMP119} operate, and these second moments
converge by parts (b) and (c) of Lemma~\ref{4.38:prf-polarisation-limit}, which is
hypothesis (2). Hence $\beta_j$ is well defined.

\emph{Dependence on the couplings.} Three objects enter $\Pi^{(j)}_{\infty}$:
\begin{itemize}
\item[(i)] the domains $D^{\infty}_j$;
\item[(ii)] the site-free family $E^{(j)}(\hat X;\,\cdot\,)$;
\item[(iii)] the map $\mathsf{h}_{\infty}$.
\end{itemize}
They contain the numbers $(g_i,\beta_i,\check{R}_i)_{i<j}$ and the parameters $\mathfrak{p}$, and
nothing else. Here $\beta_i=x_{i-1}-x_i$ with $x_i=g_i^{-2}$. Moreover $\check{R}_i$ is the
cell-size multiplier of level $i$, supplied by the lemma that defines it; like every cube size of
the construction it is, by clause (L1) of Definition~\ref{4.38:def-standing-data}, a function of
$\mathfrak{p}$ and of the couplings $g_0,\dots,g_i$ generated up to level $i$. Consequently
$\beta_j$ is a function of $(g_0,\dots,g_{j-1})$ and $\mathfrak{p}$ alone.

\emph{Reality.} Clause (d) of $(\mathrm{H}_j)$ gives reality at real $(\mathbf{U},\mathbf{J})$, so
that the function of which $\Pi^{(j)}_{\mathfrak{s}}$ is the Hessian is real at real fields. Each
$\Pi^{(j)}_{\mathfrak{s}}$ is therefore real. This matches
the statement of \cite{B-CMP109} that $\Pi_{j+1,\mu\nu}(g_j,x)$ is a real valued function. By
hypothesis (2), $\Pi^{(j)}_{\infty}$ is the limit of the real tensors $\Pi^{(j)}_{\mathfrak{s}}$.
It is therefore real, and so is $\beta_j$.

\emph{Splitting.} The tensor $\Pi^{(j)}_{\infty}$ is linear in the family $E$ of localised terms.
That family is the sum of its two marked sub-families $\natural\in\{\mathfrak{A},\Re\}$. Taking the
coefficient \cite[(3.61)]{B-CMP119} of each part gives
\begin{equation}\label{4.38:eq-increment-site-free-1}
\beta_j=\beta_j^{\mathfrak{A}}+\beta_j^{\Re}.
\end{equation}

\emph{The next level.} We put $x_j:=x_{j-1}-\beta_j$. The quantity $x_{j-1}$ is site-free by
clause (a) of $(\mathrm{H}_j)$, and $\beta_j$ is site-free by the above; hence $x_j$ is
site-free. Suppose moreover that $x_j>0$, which is clause $(\beta_j)$ of $\mathfrak{i}(j)$. Then
$g_j=x_j^{-1/2}$ is defined, and $\check{R}_j$, $\varepsilon_j$, $\delta_j$ and $k_1(j)$ are fixed
from $g_j$ and the couplings already generated by clause (L1) of
Definition~\ref{4.38:def-standing-data}; for instance $\varepsilon_j=g_j\,p_0(g_j)$ with
$p_0(\cdot)$ the degree function. Hence all of these quantities are
site-free. The site class $\mathcal{S}_j$ is fixed by these numbers, and is therefore site-free
as well.

\emph{Independence of the bookkeeping choices.} The coefficient $\beta_j$ does not depend on the
bookkeeping choices collected in clause (L4) of Definition~\ref{4.38:def-standing-data}. The reason
is that $\Pi^{(j)}_{\mathfrak{s}}$ is the second variation of a total. No total is moved by the
walk-resummation rule $\mathcal{R}$, by the vacuum and strip normalisations of that clause, or by
the remaining choice collected there. For the rule $\mathcal{R}$ this is
Proposition~\ref{4.4:prop-rule-properties}; for the vacuum and strip normalisations it is the
theorem stating that these two normalisations change no total; for the remaining choice it is the
statement, proved alongside that theorem, that this choice changes no total.

\emph{Independence of the profile.} Nor does $\beta_j$ depend on the choice of the coupling profile
in clause (L2) of Definition~\ref{4.38:def-standing-data}. By
Lemma~\ref{4.38:lem-profile-telescopes}, through the telescoping identity
\eqref{4.38:eq-profile-telescopes-a}, every total is independent of that choice.
\end{proof}

\begin{lemma}[Label-independence: inherited terms and the nesting chain]\label{4.38:prf-label-inherited}
This lemma is a step of the proof of Theorem~\ref{4.38:thm-torus-independence}. It treats the
members produced by the operation $\mathbf{T}$; the nesting chain is established both at labels
produced by $\mathbf{T}$ and at labels produced by an $\mathbf{R}$-operation. Let $j\ge1$. Assume
hypothesis $(H_j)$ of
Theorem~\ref{4.38:thm-torus-independence}, that $x_j\ge\gamma^{-2}$, and that step $j$ has been
completed at every $\mathfrak{s}\in\mathcal{S}_j$. Assume moreover the following four statements.
\begin{itemize}
\item[(I)] Corollary~\ref{4.4:cor-localised-independence} at every site, together with the
statements by which the two interior series compared there exist on local data:
Lemma~\ref{4.4:lem-interior-sufficiency}, the local half of clause (b2) of the quadratic-form
theorem on the sites of $\mathcal{S}_0$, Lemma~\ref{4.4:lem-cell-coercivity}, and the positivity
statement at level zero.
\item[(II)] The representation and bound of the large-field terms, namely: its clause on
completion terms
(at a label produced by an $\mathbf{R}$-operation, the terms of the new action whose localization
domains meet a healed island are, together with the convention on vacuum terms, members of the
whole-lattice family and free of the label); and its account of the domains of the large-field
step (the domains $Y_b$ introduced by the $\mathbf{R}$-operation of step $j-1$ lie in
$Z_{j-1}=\Lambda_{j-1}^c$; the result of the
$\mathbf{R}$-operation is again of the form \cite[(2.18)]{B-CMP119}; the sequences of a label
produced by the $\mathbf{R}$-operation obey the nesting \cite[(2.1)]{B-CMP119}).
\item[(III)] Clause (d) of the lemma on the small-field function $\chi^{\wedge}$ of the family of
labels constructed in this book: for every label of that family and every level $t$, the set
$\Lambda_t^c$ contains $(\Omega_t^c)^{\sim2}$, the set $\Omega_t^c$ enlarged by two layers of
cells.
\item[(IV)] The strip separation statement: for a label produced by $\mathbf{T}$ before the
$\mathbf{R}$-operation, distinct connected components of the complement of its small-field set
$\Lambda_j$ share no point.
\end{itemize}
Fix $\mathfrak{s}\in\mathcal{S}_j$ and a label $\ell$ at $\mathfrak{s}$, with sequences
$\Lambda_i(\ell)$, $\Omega_i(\ell)$. For $i\le j$ the level-$i$ label cubes (cells) have side
$a_i=M\check{R}_iL^{-(j-i)}$ in level-$j$ units; $[\,\cdot\,]_i$ and $(\cdot)_i^{+1}$ are the
round-up and the one-cell enlargement of Definition~\ref{4.38:def-intrinsic-trace} formed with
these cubes, and $[\,\cdot\,]=[\,\cdot\,]_j$, $(\cdot)^{+1}=(\cdot)_j^{+1}$. Then the following
hold.
\begin{itemize}
\item[(a)] For every non-wrapping domain $W$ and every $i\le j$, the member
$-A(x_j(\cdot),\mathbf{U})\restriction W$ and the members $-\beta_iA(h_z\varphi_i,\mathbf{U})$
with $\operatorname{supp}h_z\subset W$, where $h_z$ are the tent functions of the counterterms
in \cite[(2.23)--(2.25)]{B-CMP119}, are functions of the fields on $W$ and of the trace of $\ell$
on $W^{+1}$.
\item[(b)] For every non-wrapping $X\in\mathbf{D}_j$ with $X\subset\Lambda_j(\ell)$, the localized
term attached to $X$ in the actual setting at $\mathfrak{s}$ is the one attached to $X$ in the
whole-lattice setting of $\mathfrak{s}$; if $\ell$ is produced by an $\mathbf{R}$-operation and
$X$ meets a healed island, this term is a completion term, a member of the whole-lattice family.
\item[(c)] For $X$ as in \textup{(b)}, every $i<j$ and every domain $X'\subset X$ at level $i$,
\begin{equation}\label{4.38:eq-label-inherited-a}
[X']_i^{+1}\subset[X]_j^{+1}\subset\Lambda_i(\ell).
\end{equation}
Consequently $X'\in\mathcal{X}_i(\ell)$, and the inherited terms $R^{(i)}(X';\cdot)$ with
$i\le k_1(k)<j$, where $k_1(k)$ is Ba{\l}aban's lower index for the inherited terms in
\cite{B-CMP119}, $k$ denoting there the index of the step of that representation, are present in
the whole-lattice setting and equal $\mathfrak{r}^{(i)}(\hat X';\cdot)$ there.
\item[(d)] The terms $E_{\mathfrak{s}}^{(j)}(X;\cdot,z)$ of the label $\ell$, with
$X\subset\Lambda_j(\ell)$ and
$z\in\Lambda_j^0$, the set over which the base point $z$ of the pointed terms attached to
$\Lambda_j=\Lambda_j(\ell)$ ranges in Ba{\l}aban's representation, are the functions described by
Theorem~\ref{4.38:thm-torus-independence} for the
pointed localized terms; since $\ell$ is arbitrary, this holds for every label at $\mathfrak{s}$.
\end{itemize}
\end{lemma}

\begin{proof}
\emph{Cells at different levels.} Since $M\check{R}_j/a_i=L^{j-i}\check{R}_j/\check{R}_i$ and this
ratio is a non-negative integer power of $L$ by the choice of the cell-size multipliers, $a_i$
divides $M\check{R}_j$; the partitions into level-$j$ cells, into level-$i$ cells and into
level-$i$ $M$-cubes (of side $ML^{-(j-i)}$ in level-$j$ units) are therefore compatible and
successively finer, each level-$i$ cell lying in one level-$j$ cell and every $M$-cube of the
level-$j$ lattice being a union of level-$i$ $M$-cubes. This is
what allows Lemma~\ref{4.38:lem-round-up-touch} to be applied below with $\mathsf{K}$ the level-$j$
cells, $\mathsf{K}'$ the level-$i$ cells and $\mathsf{k}$ the level-$i$ $M$-cubes. Every further
member of the representation is shown, in the course of the proof of
Theorem~\ref{4.38:thm-torus-independence}, to be intrinsic in the sense of
Definition~\ref{4.38:def-intrinsic-trace}, the label entering through its trace on $W^{+1}$; here
we treat the members carrying the coupling
profile, then the terms of the whole-lattice comparison, then the chain.

\emph{Step 1: the members carrying the profile (part (a)).} By clause (e) of
Lemma~\ref{4.38:lem-coupling-profile}, on a plaquette $p\subset W$ the value $\varphi_i(x(p))$ of
the profile \eqref{4.38:eq-standing-data-a} is determined by the cells of $[\Lambda_i(\ell)^c]_i$
within sup-distance $\tfrac14a_i$ of $x(p)$, all of which lie in $[p]_i^{+1}\subset W^{+1}$. By
Lemma~\ref{4.38:lem-profile-telescopes}, the running inverse coupling satisfies the telescoping
identity \eqref{4.38:eq-profile-telescopes-a}, and $A(c,\mathbf{U})$ is linear in the weight $c$;
hence
\begin{equation*}
-A(x_j(\cdot),\mathbf{U})\restriction W
=-x_0A(\mathbf{U})\restriction W+\sum_{i\le j}\beta_iA(\varphi_i,\mathbf{U})\restriction W ,
\end{equation*}
and the right-hand side is a function of the fields on $W$ and of the values of the $\varphi_i$
on the plaquettes of $W$, hence, by the first sentence of this step, of the trace of $\ell$ on
$W^{+1}$. For a member
$-\beta_iA(h_z\varphi_i,\mathbf{U})$ the tent $h_z$ overhangs $z$ by one level-$i$ lattice unit
$L^{-(j-i)}$, so that the condition $\operatorname{supp}h_z\subset W$, which is the hypothesis
made, is needed for the member to involve only plaquettes of $W$; under it the member is a
function of the fields on $W$. The cells read by $\varphi_i$ at
points of $\operatorname{supp}h_z$ lie at sup-distance less than $\tfrac14a_i+L^{-(j-i)}<a_i$
from $z$, hence inside $[W]_i^{+1}$ by Lemma~\ref{4.38:lem-round-up-touch}(b); and
$[W]_i^{+1}\subset[W]_j^{+1}=W^{+1}$ by
Lemma~\ref{4.38:lem-round-up-touch}(c) with $W'=W$. This proves (a).

\emph{Step 2: the whole-lattice comparison (part (b)).} By Corollary~\ref{4.4:cor-localised-independence}
at $\mathfrak{s}$, for $X\subset\Lambda_j$ the term of the actual setting is that of the
whole-lattice setting of $\mathfrak{s}$. The two interior series compared there exist on local
data: this is Lemma~\ref{4.4:lem-interior-sufficiency}, whose argument is a local datum, read
together with the local half of clause (b2) of the quadratic-form theorem and
Lemma~\ref{4.4:lem-cell-coercivity}, as listed in (I). This matters because the comparison
covers in particular domains $X\subset\Lambda_j$ that encircle a hole, for which there need be no
site of $\mathcal{S}_0$ carrying a label with the same trace on $X^{+1}$. At a label
produced by an $\mathbf{R}$-operation, for $X\subset\Lambda_j(\ell)$ meeting a healed island, the
term is a completion term: \cite[p.~377]{B-CMP122b} states ``Now we add all the missing terms with
localization domains intersecting the domain $Z$, more exactly we add these terms to complete the
effective action''. By the clause on completion terms of (II), with the convention on vacuum
terms, such a
term is by definition the member of the whole-lattice family and is free of the label. This proves
(b).

\emph{Step 3: the inherited terms, given the chain (part (c)).} The first inclusion of
\eqref{4.38:eq-label-inherited-a} is Lemma~\ref{4.38:lem-round-up-touch}(c) with $W'=X'$, $W=X$
and the partitions named above. Assume the second inclusion, proved in Steps~4 and~5. Every
level-$i$ cell of $[X']_i^{+1}$ then lies in $\Lambda_i(\ell)$; a level-$i$ cell inside
$\Lambda_i(\ell)$ holds no $M$-cube of $\Lambda_i(\ell)^c$, hence is not a cell of
$[\Lambda_i(\ell)^c]_i$. So the trace of $\ell$ on $[X']_i^{+1}$ is empty, that is,
$X'\in\mathcal{X}_i(\ell)$, and $\mathcal{X}_i(\ell)$ is the range at level $i$ of the
identification of the remainder terms with the label-free family. By hypothesis $(H_j)$, clause
$(c^{+})$, that identification holds at level $i<j$; hence the inherited terms
$R^{(i)}(X';\cdot)$, $i\le k_1(k)<j$, are present in the whole-lattice setting and equal
$\mathfrak{r}^{(i)}(\hat X';\cdot)$ there.

\emph{Step 4: the chain at a label produced by $\mathbf{T}$.} Let $\ell$ be produced by the
renormalization transformation, before any $\mathbf{R}$-operation. Since $X\subset\Lambda_j$ and
$\Lambda_j$ is a union of cells, $[X]\subset\Lambda_j$, and the cells touching $[X]$ lie in
$\Lambda_j^{\sim}$, the union of $\Lambda_j$ with the cells touching it; so
\begin{equation*}
[X]_j^{+1}\subset\Lambda_j^{\sim}\subset\Omega_j\subset\Lambda_{j-1}\subset\Lambda_i ,
\end{equation*}
where the last three inclusions are justified as follows.

$\Lambda_j^{\sim}\subset\Omega_j$. On the family of labels constructed in this book this is read
from (III): $\Lambda_j^c\supset(\Omega_j^c)^{\sim2}$, that is, $\Lambda_j\subset\Omega_j^{\sim-2}$,
of which one layer is used. Indeed $\Lambda_j\subset\Omega_j^{\sim-1}$ says that no cell of
$\Lambda_j$ equals or touches a cell of $\Omega_j^c$, and $\Lambda_j^{\sim}\subset\Omega_j$ says
that no cell of $\Omega_j^c$ equals or touches a cell of $\Lambda_j$; touching being symmetric,
the two are equivalent. For the labels of Ba{\l}aban's construction the same follows from the
text accompanying \cite[(3.20)]{B-CMP119}, with Ba{\l}aban's $k+1$ equal to $j$: ``We surround the
domain $(\Omega^c_{k+1})^{\sim}\cup R'_{k+1}$ by a layer of $LMR_{k+1}$-cubes, and we denote the
complement of the so-obtained domain by $\Lambda_{k+1}$''. For the domains produced by the
renormalization transformation \cite{B-CMP119} states: ``The domains $\Omega_j$, $\Lambda_j$,
which are determined by the $j$-th renormalization transformation, but not by the
$\mathbf{R}$-operation, are unions of $MR_j$-cubes \dots{} the distance between their boundaries
is at least equal to $2MR_j$''; the quoted sentence concerns only domains not produced by an
$\mathbf{R}$-operation, which is the case of this step; labels produced by an
$\mathbf{R}$-operation are treated in Step~5.

$\Omega_j\subset\Lambda_{j-1}$. By \cite{B-CMP119},
$\Omega_j\subset\bigl((Z'_{j-1})^{\sim5}\bigr)^c$, where $Z'_{j-1}$ is ``the union of all
$LMR_{k+1}$-cubes intersecting the region $Z_k$'' with $k=j-1$, and $Z_{j-1}=\Lambda_{j-1}^c$,
which by (II) holds the sets $Y_b$ of step $j-1$. Since
$Z_{j-1}\subset Z'_{j-1}\subset(Z'_{j-1})^{\sim5}$, taking complements gives
$\bigl((Z'_{j-1})^{\sim5}\bigr)^c\subset\Lambda_{j-1}$.

$\Lambda_{j-1}\subset\Lambda_i$ for $i\le j-1$. This is the nesting \cite[p.~254, (2.1)]{B-CMP119}
of the admissible sequences of $\rho_{j-1}$, with Ba{\l}aban's $k$ equal to $j-1$:
\begin{equation*}
\Omega_1\supset\Lambda_1\supset\Omega_2\supset\dots\supset\Omega_k\supset\Lambda_k .
\end{equation*}

\emph{Step 5: the chain at a label produced by an $\mathbf{R}$-operation.} Here \cite{B-CMP119}
states: ``we admit the possibility that $\Lambda_j=\Omega_j$ for some indices $j$. Such a situation
may arise as a result of an $\mathbf{R}$-operation''; so the inclusion
$\Lambda_j^{\sim}\subset\Omega_j$ is not used at such a label. What is used is only
\begin{equation}\label{4.38:eq-label-inherited-1}
[X]_j^{+1}\subset\Lambda_i(\ell)\qquad(i<j,\ X\subset\Lambda_j(\ell)).
\end{equation}
Write $\Lambda_i^{\mathbf{T}}$, $\Omega_i^{\mathbf{T}}$ for the sequences of the label at
$\mathfrak{s}$ produced by $\mathbf{T}$ before the $\mathbf{R}$-operation, and
$Z_j^{\mathbf{T}}=(\Lambda_j^{\mathbf{T}})^c$. Three facts are used.
\begin{itemize}
\item[(A)] Distinct components of $Z_j^{\mathbf{T}}$ share no point; this is (IV).
\item[(B)] Healing never enlarges a large-field set: $\Lambda_i(\ell)\supset\Lambda_i^{\mathbf{T}}$
for every $i\le j$. At level $j$, with Ba{\l}aban's $k$ equal to $j$, \cite[p.~390]{B-CMP122b}
states ``The new large field region is equal to $Z_k\cup\bigcup_{i=1}^{m}Y_i$'' (the $Y_i$ there
are the sets $Y_b$, $m$ their number), where by \cite{B-CMP122a}, ``We denote $Z_k\setminus Z$ by
$Z_k$'', that is,
$Z_k$ there is the old region less the union $Z$ of the islands. At levels $i<j$ the sequences
produced by the $\mathbf{R}$-operation live inside $Z$, by the first line of
\cite[p.~179, (1.11)]{B-CMP122a}, ``$Z''_j=(\Omega_j^{\sim5})^c\cap Z$'', and by
\cite{B-CMP122b} the result is in the form \cite[(2.18)]{B-CMP119}; these are the statements of
(II).
\item[(C)] The sequences of $\ell$ obey the nesting \cite[(2.1)]{B-CMP119}; this is (II).
\end{itemize}
Now $\Lambda_j(\ell)=\Lambda_j^{\mathbf{T}}\cup\bigcup_aX_a$, the $X_a$ being the healed islands.
Let $C'$ be a cell of $[X]^{+1}$; it equals or touches a cell $C$ of $[X]\subset\Lambda_j(\ell)$.
\begin{itemize}
\item If $C'$ lies in a healed $X_a$, then
\begin{equation*}
C'\subset X_a\subset\Lambda_j(\ell)\subset\Lambda_i(\ell)
\end{equation*}
by (C).
\item If $C'$ equals or touches a cell of $\Lambda_j^{\mathbf{T}}$, then, by the chain of Step~4
applied to the label produced by $\mathbf{T}$ and by (B),
\begin{equation*}
C'\subset\Lambda_j^{\mathbf{T},\sim}\subset\Omega_j^{\mathbf{T}}\subset\Lambda_i^{\mathbf{T}}
\subset\Lambda_i(\ell).
\end{equation*}
\item Otherwise $C'$ neither equals nor touches a cell of $\Lambda_j^{\mathbf{T}}$, so $C'$ is a
cell of $Z_j^{\mathbf{T}}$, and $C'$ lies in no healed island. The cell $C$ is not in
$\Lambda_j^{\mathbf{T}}$, since $C'$ equals or touches it; as $C\subset\Lambda_j(\ell)$, it lies in
some healed $X_a$, and $C'\ne C$, so $C'$ touches $C$. Then $C'$ lies in a component of
$Z_j^{\mathbf{T}}$ sharing a point with $X_a$, that is, in $X_a$ by (A); this is excluded.
\end{itemize}
Hence every cell of $[X]^{+1}$ lies in $\Lambda_i(\ell)$, which is \eqref{4.38:eq-label-inherited-1}.
Together with Step~4 this proves the second inclusion of \eqref{4.38:eq-label-inherited-a} for
every label, and Step~3 then gives (c).

\emph{Step 6: conclusion (part (d)).} By \eqref{4.38:eq-label-inherited-a}, every $X'\subset X$
has $[X']_i^{+1}\subset\Lambda_i(\ell)$, that is, $X'$ lies in $\Lambda_i(\ell)^{\sim-1}$ formed
with level-$i$ cells, and the identification of Step~3 applies at level $i$. At a level $i$ at
which $\Lambda_i(\ell)^c$ is a union of $M$-cubes and not of cells, the enlargement is formed with
$[\Lambda_i(\ell)^c]_i$, as in Definition~\ref{4.38:def-intrinsic-trace} and in the hypotheses of
Lemma~\ref{4.38:lem-round-up-touch}, and the conclusion holds a fortiori, since a cell inside
$\Lambda_i(\ell)$ is not a cell of $[\Lambda_i(\ell)^c]_i$. By (b) the term attached to $X$ is
the one of the whole-lattice setting, and by (c) its inherited terms are the label-free terms
$\mathfrak{r}^{(i)}(\hat X';\cdot)$; the description of the pointed localized terms of the
whole-lattice setting at level $j$, given in Theorem~\ref{4.38:thm-torus-independence}, then
applies to it, together with (a). Hence the terms
$E_{\mathfrak{s}}^{(j)}(X;\cdot,z)$ of the label $\ell$, $X\subset\Lambda_j(\ell)$,
$z\in\Lambda_j^0$, are the functions described by
Theorem~\ref{4.38:thm-torus-independence} for the pointed localized terms; since $\ell$ is
arbitrary, this holds for every label at $\mathfrak{s}$, which is (d).
\end{proof}

\begin{lemma}[Label-independence: the remainder series and its catalogues]
\label{4.38:prf-label-remainder}
Fix a level $j$, a site $\mathfrak{s}=(m,K')\in\mathcal{S}_j$ and a label $\ell$ at $\mathfrak{s}$,
the label retained after the operation $\mathbf{R}$. We work in the setting of
Lemma~\ref{4.38:prf-label-inherited}; in particular $x_j\ge\gamma^{-2}$, and step $j$ has been
completed at every site of $\mathcal{S}_j$. The label has the sequences $\Lambda_i(\ell)$,
$\Omega_i(\ell)$, $i\le j$, the islands healed by $\mathbf{R}$, with their histories, and its sets
$Y_b(\ell)$. Write $\Lambda_j^{\mathbf{T}}$ for the small-field region at level $j$ of the label made
by $\mathbf{T}$ at $\mathfrak{s}$ before the operation $\mathbf{R}$; then $\Lambda_j(\ell)$ is the union
of $\Lambda_j^{\mathbf{T}}$ and the healed islands (Lemma~\ref{4.38:prf-label-inherited}).
\emph{Cells} are the level-$j$ label cubes, of side $M\check{R}_j$. At level $i\le j$
the label cubes have side $a_i=M\check{R}_iL^{-(j-i)}$ in level-$j$ units; $a_i$ divides
$M\check{R}_j$, the quotient $L^{j-i}\check{R}_j/\check{R}_i$ being a power of $L$ by the choice of
the cell-size multipliers. The operations $[\,\cdot\,]$,
$(\cdot)^{+1}$, $[\,\cdot\,]_i$, $(\cdot)_i^{+1}$ and the trace of a label are those of
Definition~\ref{4.38:def-intrinsic-trace} and of \eqref{4.38:eq-intrinsic-trace-a}. We write
$X^{\sharp}$ for the strip of $m_\sharp$ layers of $M$-cubes attached to a union $X$ of $M$-cubes
by \eqref{4.38:eq-standing-data-b}.

Assume the following.
\begin{itemize}
\item[($\alpha$)] The representation and bound of the large-field terms. This includes its lemma
on the series $R'$ and the format conditions on polymers stated there, among them the strip
condition $X_a^{\sharp}\subset X'$ for an island $X_a$ of a polymer $X'$.
\item[($\beta$)] The gluing and restriction statements for admissible labels, in the form of
Lemma~\ref{4.38:lem-growth-factorise} applied to the growth recursion by which the large-field
regions of the family are produced.
\item[($\gamma$)] The localisation corollary for the part $R''^{(j)}$ of the remainder, by which
that part, on domains contained in the small-field region made by $\mathbf{T}$, is the term of the
whole-lattice setting.
\item[($\delta$)] The strip normalisation \eqref{4.38:eq-standing-data-b}.
\item[($\varepsilon$)] The vacuum normalisation convention, which attaches vacuum values to strips.
\end{itemize}
Then the following hold for every $X\in\mathbf{D}_j$.
\begin{enumerate}
\item[(1)] If $X\subset\Lambda_j(\ell)$ and $X$ is non-wrapping, then $R''^{(j)}(X;\cdot)$ is
intrinsic in the sense of Definition~\ref{4.38:def-intrinsic-trace}, the label being read through
its trace on $X^{+1}$, and it is the member $R''^{(j)}(\hat X;\cdot)$ of the label- and
site-free family of \eqref{4.38:eq-torus-independence-b}: by ($\gamma$) when
$X\subset\Lambda_j^{\mathbf{T}}$, and, when $X$ meets a healed island, as the term added by
$\mathbf{R}$ to complete the effective action \cite[p.~377]{B-CMP122b}. The translations $\tau_t$,
$t\in M\check{R}_j\mathbb{Z}^4$, preserve every partition it reads.
\item[(2)] The term $R'^{(j),\ell}(X;\cdot)$ is the series
\begin{equation}\label{4.38:eq-label-remainder-a}
R'^{(j),\ell}(X;\cdot)=\sum_{n\ge1}\frac{1}{n!}
\sum_{\substack{X'_1,\dots,X'_n\in\mathcal{C}_j(X;\ell)\\ X'_1\cup\dots\cup X'_n=X}}
\rho^{T}(X'_1,\dots,X'_n)\prod_{q=1}^{n}F(X'_q;\cdot).
\end{equation}
Here $F(X'_q;\cdot)$ is the activity of the polymer $X'_q$ with its contents, and $\rho^{T}$, in the
notation of \cite[(2.13)]{B-CMP116}, is the
truncated coefficient of the polymer gas, a function of the incidences of the polymers only.
\item[(3)] The catalogue $\mathcal{C}_j(X;\ell)$ consists of the polymers $X'\subset X$ together with
their contents: islands $\{(X_a,\mathfrak{h}_a)\}$ with $X_a\subset X'$, and a finite set
$\mathcal{Y}(X')$ of domains. For labels of the family produced by the growth recursion of
($\beta$), membership is the conjunction of
\begin{equation}\label{4.38:eq-label-remainder-c}
\begin{aligned}
&\text{(c1)}\quad \text{each }(X_a,\mathfrak{h}_a)\text{ is, by itself, a first-class island;}\\
&\text{(c2)}\quad X_a\cap X_{a'}=\emptyset\ \ (a\ne a'),\qquad X_a^{\sharp}\subset X';\\
&\text{(c3)}\quad X_a\cap C=\emptyset\ \text{ for every cell } C \text{ of }\Lambda_j(\ell)^c;\\
&\text{(c4)}\quad \text{every } Y\in\mathcal{Y}(X') \text{ has } Y\in\mathbf{D}_j,\ Y\subset X',\\
&\qquad\quad\ \, Y\cap Z^{\sharp}\ \text{a union of components of } Z^{\sharp},
 \text{ at least one};\\
&\text{(c5)}\quad \text{vacuum values are attached at the strips } X_a^{\sharp},\ Y_b(\ell)^{\sharp}
 \text{ as in }(\varepsilon),
\end{aligned}
\end{equation}
where, in (c4), $Z:=\bigcup_aX_a\cup\bigcup_bY_b(\ell)$ is formed with the islands $X_a$ of the
contents of $X'$ and the sets $Y_b(\ell)$ of the label.
Every polymer satisfying \eqref{4.38:eq-label-remainder-c} occurs in the series.
\item[(4)] Let
\begin{equation}\label{4.38:eq-label-remainder-b}
\begin{aligned}
&\mathcal{C}_j(X;\ell)=\mathcal{C}_j(X;\ell')\quad
\text{whenever }\ell,\ell'\text{ have the same trace on }X^{+1},\\
&\mathcal{X}_j(\ell):=\{X''\in\mathbf{D}_j:\ \text{the trace of }\ell\text{ on }X''^{+1}
\text{ is empty}\},\\
&\mathfrak{r}'^{(j)}(X;\cdot):=\sum_{n\ge1}\frac{1}{n!}
\sum_{\substack{X'_1,\dots,X'_n\in\mathcal{C}_j(X;\emptyset)\\ X'_1\cup\dots\cup X'_n=X}}
\rho^{T}(X'_1,\dots,X'_n)\prod_{q=1}^{n}F(X'_q;\cdot).
\end{aligned}
\end{equation}
The identity in the first line holds for all labels $\ell,\ell'$ at $\mathfrak{s}$ of the family
produced by the growth recursion of ($\beta$), the empty label included. If $\ell$ belongs to that
family, then for
$X\in\mathcal{X}_j(\ell)$ we have
$\mathcal{C}_j(X;\ell)=\mathcal{C}_j(X;\emptyset)$. Moreover, for such $X$ the set $X^{+1}$ contains
no cell of $[\Lambda_i(\ell)^c]_i$ or of $[\Omega_i(\ell)^c]_i$, for any $i\le j$.
\end{enumerate}
\end{lemma}

\begin{proof}
\emph{The part $R''$.}
Let $X\subset\Lambda_j(\ell)$ be non-wrapping; then every $X'\subset X$ is non-wrapping. The term
$R''^{(j)}(X;\cdot)$ is treated by hypothesis ($\gamma$).
We use it together with the construction of $R''^{(j)}$ in \cite{B-CMP119}. In that construction
the terms $R^{(i)}(X';\cdot)$ with $X'\subset X$ enter as additional vertices, for the levels $i$ in
the range written $k_1(k)<i\le k$ in Ba{\l}aban's indexing of levels (his current level $k$ and
the lower end $k_1(k)$ fixed there). The step there passes from level $k$ to level $k+1$, and the
level $k+1$ produced by it is level $j$ here; so these are levels $i\le j-1<j$, and for them the
nesting chain
\eqref{4.38:eq-label-inherited-a} of Lemma~\ref{4.38:prf-label-inherited} shows that these vertices
are present. Each of them is the term $\mathfrak{r}^{(i)}(\hat X';\cdot)$ of
\eqref{4.38:eq-torus-independence-b}, because, for these $i<j$,
\begin{equation*}
[X']_i^{+1}\subset[X]_j^{+1}\subset\Lambda_i(\ell).
\end{equation*}
Level-$i$ label cubes contained in $\Lambda_i(\ell)$ contain no $M$-cube of $\Lambda_i(\ell)^c$.
Hence $[X']_i^{+1}$ contains no level-$i$ label cube of $[\Lambda_i(\ell)^c]_i$, so that $X'$ belongs
to the set $\mathcal{X}_i(\ell)$, defined at level $i$ as in \eqref{4.38:eq-label-remainder-b}.
Moreover, since $[X']_i^{+1}\subset\Lambda_i(\ell)$, the domain $X'$ lies in $\Lambda_i(\ell)$ with
one layer of level-$i$ label cubes removed, which is the condition of the level-$i$ case of the
identification \eqref{4.38:eq-torus-independence-b}; that case therefore gives
$R^{(i)}(X';\cdot)=\mathfrak{r}^{(i)}(\hat X';\cdot)$.
The argument is therefore an induction on the level. Through these additional vertices it uses
hypothesis ($\alpha$), on which the terms $\mathfrak{r}^{(i)}$ depend.

We distinguish the position of $X$ relative to the healed islands. Recall that $\Lambda_j(\ell)$ is the
union of $\Lambda_j^{\mathbf{T}}$ and the healed islands. By the Restriction part of
Lemma~\ref{4.38:lem-growth-factorise}, the retained sequence with the data of its healed islands
split off is again a run of the recursion, and $\Lambda_j(\ell)$ is its small-field region at level
$j$ made by $\mathbf{T}$.
Hence hypothesis ($\gamma$) applies to every $X\subset\Lambda_j^{\mathbf{T}}$, and on such $X$ it
identifies $R''^{(j)}(X;\cdot)$, through the additional vertices $R^{(i)}=\mathfrak{r}^{(i)}$ above,
with the member $R''^{(j)}(\hat X;\cdot)$ of the family of \eqref{4.38:eq-torus-independence-b}.

Suppose instead that $X$ meets a healed island (Lemma~\ref{4.38:prf-label-inherited}). Then
$R''^{(j)}(X;\cdot)$ is the completion term.
In \cite[p.~377]{B-CMP122b} the terms whose localization domains meet the large-field region are
added to complete the effective action, and the action so completed is denoted there by $A'_k$,
the level being written $k$ in that paper. Such a term is, by definition, the member
$R''^{(j)}(\hat X;\cdot)$ of the family of the whole-lattice setting. This follows from the
completion terms of the representation and bound of the large-field terms, hypothesis ($\alpha$),
and from the vacuum normalisation convention ($\varepsilon$).

Finally, let $t\in M\check{R}_j\mathbb{Z}^4$. We have
\begin{equation*}
M\check{R}_j=a_i\cdot\frac{L^{j-i}\check{R}_j}{\check{R}_i},
\end{equation*}
and $L^{j-i}\check{R}_j/\check{R}_i$ is a power of $L$, as recalled in the statement.
Hence $t$ is an integer multiple of $a_i$ in each coordinate. So $\tau_t$ preserves the partition into
level-$i$ label cubes for every $i\le j$, and with it every partition that these terms read. This
proves~(1).

\emph{The part $R'$.}
We apply the lemma on the series $R'$ contained in hypothesis ($\alpha$), with two substitutions.
First, the datum called the label there is taken to be the pair formed by the label and the site.
Second, the trace of the label on $X$ in the supplementary clause of that lemma is replaced by the
trace on $X^{+1}$.
Read in this way, the lemma represents $R'^{(j),\ell}(X;\cdot)$ as a series whose terms and index set
depend on the label and the site only through the trace on $X^{+1}$; the series is identified next,
and its index set in the paragraphs that follow.

In \cite[p.~390]{B-CMP122b} the term $R'^{(j)}(X;\cdot)$ is described as given, with notational
changes, by the convergent series \cite[(7.13)]{B-CMP122a}. Under the format conditions of ($\alpha$)
this series is the series \cite[(2.13)]{B-CMP116}, that is, \eqref{4.38:eq-label-remainder-a}. It is
the logarithm of the polymer gas \cite[(1.90)]{B-CMP122b}, and the terms $R'^{(j)}(X;\cdot)$ are
indexed by domains $X\in\mathbf{D}_j$, as on \cite[p.~390]{B-CMP122b}. The coefficient $\rho^{T}$
depends only on which
of the polymers $X'_1,\dots,X'_n$ are incident with which. This proves~(2), once the index set
$\mathcal{C}_j(X;\ell)$ is identified. We do this now.

\emph{The catalogue.}
A polymer $X'\subset X$ with contents $\{(X_a,\mathfrak{h}_a)\}$ and $\mathcal{Y}(X')$ belongs to
$\mathcal{C}_j(X;\ell)$ when it is admissible under the retained label $\ell$. For the family
considered here we now check that admissibility is the conjunction of the clauses
\eqref{4.38:eq-label-remainder-c}.

\emph{Clause (c1).} Each $X_a$ is a closed union of cells. It satisfies the two conditions of the
definition of a first-class domain in \cite{B-CMP122a}, the first of which requires it to be
contained in a cube of side $100M\check{R}_j$. It also satisfies the further domain condition imposed on
islands in this construction, and it carries its own run of the growth recursion. All these
conditions are evaluated on data inside $X_a$.

The construction in \cite{B-CMP119} also restricts, at each step $t$, the regions in which two sets
of cubes selected at that step may lie; we call them the seed sets of step $t$. The first seed set is
kept outside the enlargement by four layers of cubes of
the region $Z'_{t-1}$, and the second is kept inside a region that the construction determines from
the first, less its outer layer of cubes. Here $Z'_{t-1}$ is the union of the cubes of side
$LM\check{R}_t$, in the units of level $t$, intersecting the
large-field region $Z_{t-1}=\Lambda_{t-1}^c$ of step $t-1$, as in the proof of
Lemma~\ref{4.38:prf-label-inherited}, and $A^{\sim n}$ denotes a set $A$ enlarged by $n$ layers of
such cubes. These are conditions between different components, of four to six layers. They are
implied by (c3).
Indeed, suppose a seed of $X_a$ lies inside $(Z'_{t-1}(\ell))^{\sim6}$. The construction of
\cite{B-CMP119} adds in succession $4,1,1,2,1$ and $1$ layers, ten in all, so that
\begin{equation*}
Z_t(\ell)\supset\bigl(Z'_{t-1}(\ell)\bigr)^{\sim10}.
\end{equation*}
Hence the seed lies in $Z_t(\ell)=\Lambda_t(\ell)^c$. For $t\le j$ the nesting of the sequence
gives $\Lambda_t(\ell)^c\subset\Lambda_j(\ell)^c$. The seed would then be a point of $X_a$ in
$\Lambda_j(\ell)^c$, which contradicts (c3).

\emph{Clause (c2).} The islands pairwise share no point, and $X_a^{\sharp}\subset X'$ by the strip
condition of ($\alpha$). This clause is read inside $X$: no $X_a$ reaches the boundary of $X$.

\emph{Clause (c3).} No island shares a point with a cell of $\Lambda_j(\ell)^c$. This reflects three
facts from Ba{\l}aban's papers. A first-class domain is a component of the large-field region of the
configuration before the operation $\mathbf{R}$, written $\Lambda_k^c$ in \cite{B-CMP122a}, the level
written $k$ there being $j$ here; of that region the paper states ``It is a union of connected
components''. Components are the classes of the relation of intersecting or touching
\cite{B-CMP122b}. The smallest union of cubes of the cell size containing the large-field region is a
union of components, in the words of \cite{B-CMP122b}: ``This domain is a union of components''. The
passage to cells is made through the strip normalisation ($\delta$).

\emph{Clause (c4).} Each $Y\in\mathcal{Y}(X')$ is a domain carrying the term written $V(Y,U_k)$ in
\cite[p.~379]{B-CMP122b}. In
\cite[p.~377]{B-CMP122b} the intersection of $Y$ with the enlargement of $Z$ used there, written
$Y\cap Z^{\sim}$ in that paper, is required to be a union of components of that enlargement,
containing at least one component. The replacement of that enlargement by the strip $Z^{\sharp}$
is made by the strip normalisation ($\delta$).

\emph{Clause (c5).} This is the attachment of vacuum values at the strips $X_a^{\sharp}$ and
$Y_b(\ell)^{\sharp}$ prescribed by ($\varepsilon$).

\emph{The clauses are exactly the admissibility conditions.}
That (c1), (c2) and (c3) together are equivalent to the statement that the label $\ell$ glued with
the $(X_a,\mathfrak{h}_a)$ is an admissible label before the operation $\mathbf{R}$, with the $X_a$
first-class, is hypothesis ($\beta$): the Gluing and Restriction parts of
Lemma~\ref{4.38:lem-growth-factorise}.

Next, by the format conditions of ($\alpha$), read under ($\delta$) and ($\varepsilon$), the clauses
(c4) and (c5) are the only further conditions of the format.

Finally, every polymer satisfying \eqref{4.38:eq-label-remainder-c} occurs. The reason is that the
sums over labels in \cite[(2.18)]{B-CMP119} run over arbitrary unions of cubes. This proves~(3).

\emph{The reading set.}
We show that each clause of \eqref{4.38:eq-label-remainder-c} is a function of the trace of $\ell$
on $X^{+1}$.

Clause (c3) involves only cells sharing a point with some $X_a$. Since $X_a\subset[X]$, such a cell
lies in $[X]$ or touches $[X]$, so it lies in $X^{+1}$.

For (c4) and (c5) we use that the strips are thin. By \eqref{4.38:eq-standing-data-b},
\begin{equation*}
m_\sharp M<\tfrac12 M\check{R}_j .
\end{equation*}
Let $C'$ be a cell of $\bigcup_bY_b(\ell)$ whose strip $C'^{\sharp}$ contains an $M$-cube $c\subset Y$.
If the strips of \cite[p.~377]{B-CMP122b} are read as closed sets, the same applies when the strip
merely touches such a cube. In either case $C'$ is at sup-distance at most $m_\sharp M<M\check{R}_j$
from $c$. Since $c\subset Y\subset X$, Lemma~\ref{4.38:lem-round-up-touch}(b) shows that $C'$ lies in
$X^{+1}$. The distance here is measured from the cube $c$; the argument does not say that $C'$ touches
$X$.

The same argument applies to a whole set. If $Y_b(\ell)^{\sharp}\subset Y$, then $Y_b(\ell)$, together
with the cells touching it, lies in $X^{+1}$. The half of (c4) concerning the sets $Y_b(\ell)$ says
that every cell of $\bigcup_bY_b(\ell)$ whose strip meets $Y$ belongs to such a set $Y_b(\ell)$.

Both sets entering here are determined by the trace. Let $K_0$ be a class, under the equivalence
relation generated by touching, of cells of $\bigcup_bY_b(\ell)$ lying in $X^{+1}$, and suppose
$K_0^{\sharp}\subset Y$. Then $K_0\subset Y\subset X$, so $K_0\subset[X]$. Hence every cell touching
$K_0$ lies in $X^{+1}$ and is seen by the trace, whether or not it is marked as a cell of
$\bigcup_bY_b(\ell)$. Suppose a touch-connected set $Y_b(\ell)$ strictly contained $K_0$. It would
then contain a cell touching $K_0$, lying outside $K_0$ and marked as a cell of
$\bigcup_bY_b(\ell)$, which is impossible because $K_0$ is a whole class. So $K_0$ is a whole set
$Y_b(\ell)$. Conversely, $Y_b(\ell)^{\sharp}\subset Y\subset X$ gives $Y_b(\ell)\subset[X]$. Then
$Y_b(\ell)$ and the cells touching it lie in $X^{+1}$.

Hence every clause of \eqref{4.38:eq-label-remainder-c} is a function of the trace of $\ell$ on
$X^{+1}$, and two labels with the same trace on $X^{+1}$ have the same catalogue. This gives the
first identity of \eqref{4.38:eq-label-remainder-b}.

The strip clause does involve cells outside $[X]$: a set $Y_b(\ell)$ with
$Y_b(\ell)^{\sharp}\subset Y$ may have cells touching it that lie outside $[X]$ but inside $X^{+1}$,
and (c4) reads them. So the catalogue depends on the trace on $X^{+1}$, and the clause of the
statement on torus-independence of the localized terms, which contains
\eqref{4.38:eq-torus-independence-b}, asserting that $B^{(j)}(X;\cdot)$,
$\mathbf{B}'^{(j)}(X;\cdot)$ and $\mathbf{T}^{(j-1)}(\cdot\cap X)$ read the label only through
its trace on $X^{+1}$ (and $\mathbf{T}'_j(X)$ only through its trace on $(X^{\sharp})^{+1}$),
cannot be stated with $[X]$ in place of $X^{+1}$.

\emph{The set $\mathcal{X}_j(\ell)$ and the series $\mathfrak{r}'^{(j)}$.}
By Definition~\ref{4.38:def-intrinsic-trace}, the trace of $\ell$ on $X^{+1}$ is empty exactly when
no cell of $X^{+1}$ is a cell of $[\Lambda_j(\ell)^c]$. Let $X\in\mathcal{X}_j(\ell)$. Then every cell
of $X^{+1}$ lies in $\Lambda_j(\ell)$. The nesting \cite[(2.1)]{B-CMP119} of the retained sequence
alone, as used in the proof of Lemma~\ref{4.38:prf-label-inherited}, then gives
\begin{align*}
X^{+1}&\subset\Lambda_j(\ell)\subset\Omega_j(\ell)\subset\Lambda_i(\ell)\subset\Omega_i(\ell)
&&\text{for } i<j,\\
X^{+1}&\subset\Lambda_j(\ell)\subset\Omega_j(\ell)&&\text{for } i=j.
\end{align*}
Hence $X^{+1}$ contains no cell of $[\Lambda_i(\ell)^c]_i$ or of $[\Omega_i(\ell)^c]_i$ for any
$i\le j$.

Let now $\ell$ belong to the family produced by the growth recursion of ($\beta$).
The trace of the empty label on $X^{+1}$ is empty, so for $X\in\mathcal{X}_j(\ell)$ the labels $\ell$
and $\emptyset$ have the same trace on $X^{+1}$. Therefore the first identity of
\eqref{4.38:eq-label-remainder-b} gives, for $X\in\mathcal{X}_j(\ell)$,
\begin{equation*}
\mathcal{C}_j(X;\ell)=\mathcal{C}_j(X;\emptyset).
\end{equation*}
The series $\mathfrak{r}'^{(j)}(X;\cdot)$ of \eqref{4.38:eq-label-remainder-b} is the series
\eqref{4.38:eq-label-remainder-a} taken over the catalogue $\mathcal{C}_j(X;\emptyset)$. This
proves~(4).
\end{proof}

\begin{lemma}[The second-group range is the set of empty trace]\label{4.38:lem-range-empty-trace}
Let $j$ be a level and let $\ell$ be a label at level $j$, with level-$j$ localization domain
$\Lambda_j(\ell)$, whose complement is a union of $M$-cubes. Cells are the label cubes at level
$j$, of side $M\check{R}_j$, and $[\,\cdot\,]$, $(\cdot)^{+1}$ are the round-up to cells and the
one-cell enlargement of Definition~\ref{4.38:def-intrinsic-trace}, see
\eqref{4.38:eq-intrinsic-trace-a}.

Two unions of cubes of one partition \emph{share} a cube if some cube of that partition lies in
both; a union of $M$-cubes is \emph{contained} in a union of cells if each of its $M$-cubes lies in
one of those cells. A cell \emph{touches} a union $Z$ of cells if it is not a cell of $Z$ and meets
$Z$, along a face of any dimension down to a single corner.

For a union $Z$ of cells put $Z^{\sim}:=Z\cup\{\text{cells touching }Z\}$. For a union $A$ of
$M$-cubes write $A^{\sim}$ for $[A]^{\sim}$. For a domain $A$ whose complement is a union of
$M$-cubes let
\begin{equation}\label{4.38:eq-range-empty-trace-1}
A^{\sim-1}:=\text{the union of the cells that are not cells of }\bigl([A^{c}]\bigr)^{\sim}.
\end{equation}
Write the large-field region at level $j$ as in \cite[p.~390]{B-CMP122b}:
\begin{equation}\label{4.38:eq-range-empty-trace-2}
\Lambda_j(\ell)^{c}=Z_j\cup\bigcup_{b}Y_b ,
\end{equation}
where $Z_j$ is the large-field region inherited from the earlier steps with the domain on which
Ba{\l}aban's operation $\mathbf{R}$ acts removed, the set written $Z_k$ in \cite{B-CMP122a} after
``We denote $Z_k\setminus Z$ by $Z_k$'', and the $Y_b$, finitely many, are the large-field regions
created at step $j$ by the operation $\mathbf{R}$; all are unions of $M$-cubes.

Then for every non-wrapping $X\in\mathbf{D}_j$ which is a union of $M$-cubes the following
are equivalent:
\begin{itemize}
\item[(a)] $X\subset\Lambda_j(\ell)^{\sim-1}$, that is, $X$ lies in the range of the sum
\cite[(2.30)]{B-CMP119};
\item[(b)] $X$ shares no $M$-cube with $Z_j^{\sim}\cup\bigcup_{b}Y_b^{\sim}$, that is, in the
cube-wise reading of disjointness, a term with localization domain $X$ belongs to the second
group of \cite[p.~390]{B-CMP122b};
\item[(c)] $X^{+1}$ shares no cell with $[\Lambda_j(\ell)^{c}]$;
\item[(d)] $X\in\mathcal{X}_j(\ell)$.
\end{itemize}
If they hold, then $\operatorname{dist}_{\infty}(x,\Lambda_j(\ell)^{c})\ge M\check{R}_j$ for every
$x\in X$.
\end{lemma}

\begin{proof}
We first unwind the definition \eqref{4.38:eq-range-empty-trace-1}.

A cell that is not a cell of $[A^c]$ contains no $M$-cube of $A^c$. The cells are unions of
$M$-cubes (Definition~\ref{4.38:def-standing-data}). Every $M$-cube of such a cell therefore lies
in $A$. Hence $A^{\sim-1}$ is the union of the cells of $A$ that touch no cell of $[A^c]$. When
$A^c$ is itself a union of cells, $[A^c]=A^c$ and $A^{\sim-1}$ is the union of the cells of $A$
touching no cell of $A^c$.

This is the operation of \cite{B-CMP119}, carried out with the cells as cubes. There ${}^{\sim}$
adds one layer of cubes of the partition in force touching a domain (one layer of
$M_2R_0$-cubes), and $\Omega_{k+1}^{\sim-1}$ denotes the domain obtained by taking off one layer of
$LMR_{k+1}$-cubes from $\Omega_{k+1}$ \cite[p.~268]{B-CMP119}; here both operations are taken with
the cells of side $M\check{R}_j$. Touching includes meeting at a corner only. So the layer added by
$\sim$ consists of all cells that are not cells of $Z$ but touch one of them, in the same way as,
for an $M$-cube $\square$, the cube $\square^{n}$ of \cite{B-CMP109}, of side $(1+2n)M$ and
centred at the centre of $\square$, contains the $M$-cubes that meet $\square$ at a corner only.

On \cite[p.~259]{B-CMP119} the set $\Lambda_j^{0}$, obtained by removing one layer of
$MR_j$-cubes from the domain $\Lambda_j$ (written $\Lambda_j^{(j)}$ there, as a subset of the
level-$j$ lattice), is the range of $z$ in \cite[(2.26)]{B-CMP119}; the same page
has $\varphi_j=1$ on $\Lambda_j^{\sim-1}$, which agrees with the reading of ${}^{\sim-1}$ as the
removal of one layer.

\emph{(a) is a statement about $M$-cubes.} Put $W_0:=([\Lambda_j(\ell)^c])^{\sim}$; this is a union
of cells. Each $M$-cube of $X$ lies in exactly one cell. That cell is either a cell of $W_0$ or,
by \eqref{4.38:eq-range-empty-trace-1}, a cell of $\Lambda_j(\ell)^{\sim-1}$. Hence (a) holds if
and only if
\begin{equation}\label{4.38:eq-range-empty-trace-3}
X \text{ and } \bigl([\Lambda_j(\ell)^{c}]\bigr)^{\sim} \text{ share no } M\text{-cube}.
\end{equation}

\emph{(a) $\Leftrightarrow$ (b).} The region \eqref{4.38:eq-range-empty-trace-2} is the new
large-field region of \cite[p.~390]{B-CMP122b}, written there as the union of the old region and
the new ones, with no layer added. Here $Z_j$ is the region that \cite{B-CMP122a} writes $Z_k$
after removing from the old large-field region the domain on which the operation $\mathbf{R}$ acts
(``We denote $Z_k\setminus Z$ by $Z_k$''), with $k$ there equal to $j$ here.

Both operations distribute over unions. The round-up of a union of $M$-cubes is the union of the
cells containing one of its $M$-cubes; hence $[A\cup B]=[A]\cup[B]$ for unions $A,B$ of $M$-cubes.
A cell equals or touches a cell of $A\cup B$ exactly when it equals or touches a cell of $A$ or a
cell of $B$; hence $(A\cup B)^{\sim}=A^{\sim}\cup B^{\sim}$ for unions $A,B$ of cells. Applying
both to \eqref{4.38:eq-range-empty-trace-2} gives
\begin{equation}\label{4.38:eq-range-empty-trace-4}
\bigl([\Lambda_j(\ell)^{c}]\bigr)^{\sim}
=[Z_j]^{\sim}\cup\bigcup_{b}[Y_b]^{\sim}
=Z_j^{\sim}\cup\bigcup_{b}Y_b^{\sim}.
\end{equation}

In \cite[p.~390]{B-CMP122b} the second group consists of the terms whose domains are disjoint
from the region written there $Z_k^{\sim}\cup\bigcup_{i=1}^{m}Y_i^{\sim}$; in the present letters
this is the right-hand side of
\eqref{4.38:eq-range-empty-trace-4}. Under the strip convention of
\eqref{4.38:eq-standing-data-b}, $\sim$ in that sentence remains the addition of one layer of
cells. We read ``disjoint'' there cube-wise, as ``shares no $M$-cube''. With this reading,
\eqref{4.38:eq-range-empty-trace-4} shows that (b) is \eqref{4.38:eq-range-empty-trace-3}, hence
equivalent to (a). A point-wise reading of ``disjoint'' would only make the second group's range
smaller than the range of \cite[(2.30)]{B-CMP119}.

\emph{(a) $\Leftrightarrow$ (c).} A non-wrapping domain is read on one of its lifts to
$\mathbb{R}^4$, on which the cells and the $M$-cubes of Definition~\ref{4.38:def-standing-data}
form partitions of $\mathbb{R}^4$.
We apply Lemma~\ref{4.38:lem-round-up-touch} (rounding up to
coarser cubes, then touching) with the following data:
\begin{itemize}
\item $\mathsf{K}$ is the partition of $\mathbb{R}^4$ into closed cells, with
$c_{\mathsf{K}}=M\check{R}_j$;
\item $\mathsf{K}':=\mathsf{K}$;
\item $\mathsf{k}$ is the partition into $M$-cubes, which is compatible with $\mathsf{K}$ and
finer than it, since the cells are unions of $M$-cubes;
\item $Z:=[\Lambda_j(\ell)^{c}]$, a union of $\mathsf{K}$-cubes;
\item $W:=X$, a union of $\mathsf{k}$-cubes.
\end{itemize}
With these choices the round-up and the enlargement of that lemma are those of
Definition~\ref{4.38:def-intrinsic-trace}, and its $Z^{\sim}$ is $([\Lambda_j(\ell)^c])^{\sim}$.
Part~(a) of that lemma states that $X$ and $Z^{\sim}$ share no $M$-cube if and only if $[X]$ and
$Z^{\sim}$ share no cell, if and only if $X^{+1}$ and $Z$ share no cell. The first condition is
\eqref{4.38:eq-range-empty-trace-3} and the last is (c). Hence (a) $\Leftrightarrow$ (c).

Sharing is cube-wise here as well. A cell of $X^{+1}$ that meets $[\Lambda_j(\ell)^c]$ only along a
face, an edge or a corner is not shared with it. To see what this means, index the cells by
$n\in\mathbb{Z}^4$, and for unions $A,A'$ of cells let $\delta(A,A')$ be the least value of
$|n-n'|_{\infty}$ over the indices $n$ of cells of $A$ and $n'$ of cells of $A'$. Two cells are
equal or touch exactly when
their indices differ by at most one in $|\cdot|_{\infty}$. Consequently (c) says
$\delta([X],[\Lambda_j(\ell)^c])\ge2$. A point-wise reading of (c) would instead say
$\delta([X],[\Lambda_j(\ell)^c])\ge3$. That condition defines a set different from the range of
\cite[(2.30)]{B-CMP119}, and under it the range would not be the set of empty trace.

\emph{(c) $\Leftrightarrow$ (d).} By the definition of $\mathcal{X}_j(\ell)$ in
\eqref{4.38:eq-label-remainder-b}, $X\in\mathcal{X}_j(\ell)$ means that the trace of $\ell$ on
$X^{+1}$ is empty. At level $j$ this says that no cell of $[\Lambda_j(\ell)^c]$ lies in $X^{+1}$,
which is (c). The conditions at the levels $i\le j$, namely that $X^{+1}$ holds no cell of
$[\Lambda_i(\ell)^c]_i$ or of $[\Omega_i(\ell)^c]_i$, follow from (c) alone. Indeed, a cell that
holds an $M$-cube of $\Lambda_j(\ell)^c$ is a cell of $[\Lambda_j(\ell)^c]$, and by (c) no cell of
$X^{+1}$ is such a cell. Since the cells are unions of $M$-cubes and $\Lambda_j(\ell)^c$ is a union
of $M$-cubes, every $M$-cube of $X^{+1}$ lies in $\Lambda_j(\ell)$, that is,
$X^{+1}\subset\Lambda_j(\ell)$. The nesting \cite[(2.1)]{B-CMP119} of the label's sequences of
localization domains, as recorded with the definition of $\mathcal{X}_j(\ell)$, then gives
\begin{equation*}
X^{+1}\subset\Lambda_j(\ell)\subset\Omega_j(\ell)\subset\Lambda_i(\ell)\subset\Omega_i(\ell),
\qquad i<j,
\end{equation*}
and for $i=j$ only the first two inclusions are used. A level-$i$ cell lying in $X^{+1}$ therefore
lies in $\Lambda_i(\ell)$ and in $\Omega_i(\ell)$; such a cell holds no $M$-cube of
$\Lambda_i(\ell)^c$ or of $\Omega_i(\ell)^c$, so it is no cell of $[\Lambda_i(\ell)^c]_i$ or of
$[\Omega_i(\ell)^c]_i$. Hence (c) $\Leftrightarrow$ (d).

\emph{The distance bound.} Suppose (c) holds, so that $X^{+1}$ and $Z=[\Lambda_j(\ell)^c]$ share
no cell. Part~(e) of Lemma~\ref{4.38:lem-round-up-touch}, with the same data, gives
$\operatorname{dist}_{\infty}(x,[\Lambda_j(\ell)^c])\ge M\check{R}_j$ for every $x\in X$. Since
$\Lambda_j(\ell)^c\subset[\Lambda_j(\ell)^c]$, it follows that
$\operatorname{dist}_{\infty}(x,\Lambda_j(\ell)^c)\ge M\check{R}_j$ for every $x\in X$.
\end{proof}

\begin{remark}
By the lemma, any statement proved for every $X\in\mathcal{X}_j(\ell)$ holds on the whole range of
\cite[(2.30)]{B-CMP119}. In particular this applies to the vanishing of the annulled difference
$B'_{\mathrm{ann}}(X)$ between the remainder term of label $\ell$ and the series
$\mathfrak{r}'^{(j)}(X)$ over the catalogue of the empty trace in
\eqref{4.38:eq-label-remainder-b}. That difference is considered only on this range, and the set
$\Lambda_j(\ell)^{\sim-2}$, obtained by removing two layers of cells, is not needed.

Neither the identification of the range nor the vanishing of $B'_{\mathrm{ann}}$ on it uses a hull
property of $\Lambda_j(\ell)$; what they use are three facts:
\begin{itemize}
\item cells are adjacent when they touch, corners included, and sharing, containment and
disjointness of domains are read cube-wise;
\item the strips of \eqref{4.38:eq-standing-data-b} are thinner than half a cell,
$m_{\sharp}M<\tfrac12 M\check{R}_j$;
\item the ramps of the coupling profiles $\varphi_i$ are thinner than a quarter cell.
\end{itemize}
\end{remark}

\begin{lemma}[Label-independence: activities read on the one-cell enlargement]
\label{4.38:prf-label-activities}
Fix a level $j$, a site $\mathfrak{s}$, and the label $\ell$ retained after the
$\mathbf{R}$-operation at level $j$. Its data are:
\begin{itemize}
\item the sequences $\Lambda_i(\ell)$, $\Omega_i(\ell)$, $i\le j$;
\item its healed islands, with their histories: the first-class islands whose large-field
factor is absorbed by the $\mathbf{R}$-operation into the small-field format;
\item the sets $Y_b$ adjoined by the $\mathbf{R}$-operation to the large-field region.
\end{itemize}
The index $k$ of \cite{B-CMP122a} and \cite{B-CMP122b} denotes the present level $j$
throughout.

For $X\in\mathbf{D}_j$ let $R'^{(j),\ell}(X)$ be the series
\eqref{4.38:eq-label-remainder-a}. It is taken over the catalogue $\mathcal{C}_j(X;\ell)$
of \eqref{4.38:eq-label-remainder-b}, whose members are described by
\eqref{4.38:eq-label-remainder-c}. Write $F(X')$ for its activities. Let
$\mathfrak{r}'^{(j)}(X)$ be the same series over $\mathcal{C}_j(X;\emptyset)$. Let
$\mathcal{X}_j(\ell)$ be the set of $X\in\mathbf{D}_j$ at which the trace of $\ell$ on
$X^{+1}$ is empty, as in \eqref{4.38:eq-label-remainder-b}.

Assume the following.
\begin{itemize}
\item[(a)] The representation and bound of the large-field terms, in the form used in this
chapter. This includes three things. First, the activity $F(X')$ of a polymer $X'$ is
given by \cite[(1.91)]{B-CMP122b}. Second, the new large-field region equals
$Z_k\cup\bigcup_bY_b$, as stated in \cite[p.~390]{B-CMP122b}, where $Z_k$ (with $k=j$) is
the large-field region made by $\mathbf{T}$ at level $j$ with the healed islands of $\ell$
removed. Third, it includes the two
sentences of \cite[p.~379]{B-CMP122b} (in which Ba{\l}aban's domain $X$ is the polymer
$X'$ here, and the subscript $j$ runs over the levels):
\begin{quotation}
All the expressions in the above integral operation are localized in the domain $X$
\end{quotation}
\begin{quotation}
Notice that the terms $V(Y,U_k)$ depend on the sequence $\{\Omega_j^c,Z_j\}$ restricted
to the components of $Z$ contained in $Y$
\end{quotation}
Of the second sentence only the words ``restricted to the components of $Z$ contained in
$Y$'' are used.
\item[(b)] The strip convention \eqref{4.38:eq-standing-data-b}, with $Z$ as in clause (c4)
of \eqref{4.38:eq-label-remainder-c} for a polymer with contents
$\{(X_a,\mathfrak{h}_a)\}$. The components of
$Z^{\sharp}$ are the strips $X_a^{\sharp}$ and $Y_b^{\sharp}$, and distinct strips share
no point.
\item[(c)] The vacuum convention of Definition~\ref{4.38:def-standing-data}. The vacuum
values entering a term $V(Y)$ are those of
terms whose domains are contained in $Y$. They are attached at the strips $X_a^{\sharp}$
and $Y_b^{\sharp}$ as in clause (c5) of \eqref{4.38:eq-label-remainder-c}.
\item[(d)] The walk family of the operators of the step. The operators entering $V(Y)$ are
built through the walk-resummation rule $\mathcal{R}$ of
Definition~\ref{4.38:def-standing-data}, with footprints in $Y$. The
Dirichlet piece at a cover cube $\square$ is that of the sequence $\{\Omega_n(\square)\}$,
which is fixed by the trace of $\{\Omega_i\}$ on $\tilde\square^{5}$, as in
\cite{B-CMP99b}.
\item[(e)] Label-independence at lower levels. At every level $i<j$ and for every label,
the following terms are intrinsic in the sense of
Definition~\ref{4.38:def-intrinsic-trace}:
\begin{itemize}
\item the localised terms born of $\mathbf{T}$;
\item the remainder terms $R''^{(i)}$ and $R'^{(i)}$;
\item the terms carrying the label.
\end{itemize}
The reading set of such a term with domain $X''$ is $[X'']_i^{+1}$.
\item[(f)] Label-independence of the terms born of $\mathbf{T}$ at level $j$. At level $j$
and for every label, the localised terms born of $\mathbf{T}$, the remainder terms
$R''^{(j)}$ and the label-carrying terms born of $\mathbf{T}$ are intrinsic in the sense of
Definition~\ref{4.38:def-intrinsic-trace}, a term with domain $X''$ having reading set
$[X'']_j^{+1}$; these are assertions, at level $j$, of part (iii)(b) of the statement on
torus-independence of the localized terms, which carries the display
\eqref{4.38:eq-torus-independence-b}.
\end{itemize}
Here $R^{(j)}$ and $R''^{(j)}$ are the remainder terms at level $j$ of part (iii)(b) of that
statement, so that $R^{(j)}=R''^{(j)}+R'^{(j),\ell}$.
Then the following hold.
\begin{enumerate}
\item $R'^{(j),\ell}(X)$ is intrinsic, with reading set $X^{+1}$.
\item For $X\in\mathcal{X}_j(\ell)$ one has $R'^{(j),\ell}(X)=\mathfrak{r}'^{(j)}(X)$, term by
term.
\item On $\mathcal{X}_j(\ell)$ the annulled difference $B'_{\mathrm{ann}}$ of
\eqref{4.38:eq-torus-independence-b} vanishes, and
\begin{equation}\label{4.38:eq-label-activities-1}
R^{(j)}=R''^{(j)}+\mathfrak{r}'^{(j)}=\mathfrak{r}^{(j)} ,
\end{equation}
where $\mathfrak{r}^{(j)}$ is the function at the empty label.
\item For every $X\in\mathbf{D}_j$ that is non-wrapping in the sense of
Definition~\ref{4.38:def-non-wrapping}, $\mathfrak{r}^{(j)}(X)$ depends neither on the site
nor on the lift. It is covariant under the translations $\tau_t$, $t\in M\check{R}_j\mathbb{Z}^4$.
\end{enumerate}
\end{lemma}

\begin{proof}
\emph{The ingredients of an activity.}
Let $X'\subset X$ be a polymer of $\mathcal{C}_j(X;\ell)$ with contents
$\{(X_a,\mathfrak{h}_a)\}$ and its family $\mathbf{D}$ of $V$-domains $Y$, as in
\eqref{4.38:eq-label-remainder-c}. By \cite[(1.91)]{B-CMP122b}, which enters through
hypothesis~(a), $F(X')$ is assembled from the operations $\mathbf{T}'_j(X_a)$ and the terms
$V(Y)$, with $X_a,Y\subset X'$.

The first quoted sentence of hypothesis~(a) says that all expressions of the integral
operation defining $F(X')$ are localised in $X'$. The second says that $V(Y,U_k)$ depends on
the sequence of large-field data only through the components of $Z$ contained in $Y$.
Accordingly $F(X')$ is built from the seven ingredients below. For each we determine what it
reads of the label.

\smallskip
\emph{(1) Integrations.} These are integrations over the fields in the islands
$X_a\subset X'$.

\smallskip
\emph{(2) Operators.} These are built through the rule $\mathcal{R}$ with footprints in $Y$,
by hypothesis~(d). The Dirichlet piece at a cover cube $\square$ is that of
$\{\Omega_n(\square)\}$, and it is fixed by the trace of $\{\Omega_i\}$ on
$\tilde\square^{5}$.

\smallskip
\emph{(3) The members carrying the coupling profile.} The island integrals and the
re-expansions carry the position-dependent couplings $1/(g''_k(\cdot))^2$ and
$1/(g_k(\cdot))^2$. In \cite[(1.70)]{B-CMP122b}, the completed-action term
$A'_k(1/(g_k(\cdot))^2,U_k)$ of \cite[p.~377]{B-CMP122b} stands
outside the integral, and the term $-A((1/(g''_k(\cdot))^2)\zeta_1,\cdot)$ stands inside it,
$\zeta_1$ being the cut-off function of \cite[(1.70)]{B-CMP122b}.
The counterterms $-\beta_iA(h_z\varphi_i,\cdot)$ are also carried, $h_z$ being the partition
of unity of \cite[(3.65)]{B-CMP119}.

The two couplings are identified as follows.
\begin{itemize}
\item The function $1/(g''_k(\cdot))^2$ is the prepared function of \cite[(1.62)]{B-CMP122a}.
Writing $n_{\mathrm{pr}}$ for the number of preparatory integrations there, one has
$g''_k=g_k^{(n_{\mathrm{pr}})}$.
\item The function $1/(g_k(\cdot))^2$ is the new function defined in \cite[p.~365]{B-CMP122b}.
\end{itemize}
For $0<n<n_{\mathrm{pr}}$ we call the sequence of large-field sets in force after the $n$-th
preparatory integration, with coupling function $g_k^{(n)}(\cdot)$ of
\cite[(1.60)--(1.62)]{B-CMP122a}, the stage-$n$ sequence.
In these members each profile is the product profile \eqref{4.38:eq-standing-data-a},
formed with the prepared sequence for $1/(g''_k(\cdot))^2$ and with the retained sequence for
$1/(g_k(\cdot))^2$.

The profile made by $\mathbf{T}$ is
\begin{equation}\label{4.38:eq-label-activities-2}
\varphi_i^{\mathrm{pre}}=\varphi_i^{\ell}\cdot\prod_a\varphi_i^{(S_a)^c}
\end{equation}
by Lemma~\ref{4.38:lem-coupling-profile}(d). Here $S_a$ is the level-$i$ large-field set of
the island $X_a$. For a union $A$ of level-$i$ cells, $\varphi_i^{A^c}$ denotes the profile
\eqref{4.38:eq-standing-data-a} with the cells of $A$ in place of those of
$[\Lambda_i(\ell)^c]_i$, and $\varphi_i^{\ell}$ denotes the profile
\eqref{4.38:eq-standing-data-a} formed with the cells of $[\Lambda_i(\ell)^c]_i$ themselves.
This profile enters through the weight differences
\begin{equation}\label{4.38:eq-label-activities-3}
A\Bigl(\frac{1}{(g_k^{(n+1)}(\cdot))^2}-\frac{1}{(g_k^{(n)}(\cdot))^2},\,\cdot\,\Bigr)
\end{equation}
of \cite[(1.62)]{B-CMP122a}. These are contained in $\mathbf{B}''_k$, the sum of the new
boundary terms of \cite[p.~192]{B-CMP122a}.

Consider a plaquette of $Y$ and split each profile into its non-island factor and its island
factors.

\emph{The non-island factor.} For the sequence made by $\mathbf{T}$ and for the retained
sequence, this factor is $\varphi_i^{\ell}$. For a stage-$n$ sequence and for the prepared
sequence it also involves the cells of $[Z''_i\cap Y_b]_i\subset Y_b$. The corresponding
factors differ from $1$ only on $Y_b^{\sharp}$, by Lemma~\ref{4.38:lem-coupling-profile}(f).
By clause (c4) of \eqref{4.38:eq-label-remainder-c} and hypothesis~(b), $Y_b^{\sharp}$ lies
in $Y$ when it meets $Y$.

By Lemma~\ref{4.38:lem-coupling-profile}(e), applied with $W=Y$, the non-island factor at a
point of $Y$ is determined by the cells of $[\Lambda_i(\ell)^c]_i$ within sup-distance
$\tfrac14a_i$ of that point. These cells lie in $[Y]_i^{+1}\subset[Y]^{+1}$.

\emph{The island factors.} An island factor is one of the following:
\begin{itemize}
\item the factor $\varphi_i^{(S_a)^c}$ of the sequence made by $\mathbf{T}$;
\item the island factor of the prepared sequence;
\item the island factor of a stage-$n$ sequence.
\end{itemize}
The set it is formed from is contained in $X_a$; for the prepared one it is
$[Z''_i\cap X_a]_i\subset X_a$. By Lemma~\ref{4.38:lem-coupling-profile}(f), the ramp lies
in the strip. So every island factor differs from $1$ only on $M$-cubes of $X_a^{\sharp}$.
It therefore affects a plaquette of $Y$ only if $Y$ meets $X_a^{\sharp}$.

By clause (c4) of \eqref{4.38:eq-label-remainder-c}, $Y\cap Z^{\sharp}$ is a union of
components of $Z^{\sharp}$. By hypothesis~(b), $X_a^{\sharp}$ is one such component and
shares no point with the others. Hence $Y$ meets $X_a^{\sharp}$ only if
$X_a^{\sharp}\subset Y$.

\emph{Conclusion for the profile members.} The members carrying the profile on $Y$ read $\ell$
on $[Y]^{+1}$. Of the islands, they read exactly those contained in $Y$. This agrees with
the second sentence quoted in hypothesis~(a).

The inclusion of the ramp in the strip is what this argument uses. A profile whose ramp
extended over a full cell could read an island not contained in $Y$ at cell-index distance
$3$ from $[Y]$, the cell-index distance of two unions of cells being the least sup-norm
distance between the index vectors of their cells (two cells strictly between).

\smallskip
\emph{(4) Older terms.} These are the older terms of each kind listed in hypotheses~(e)
and~(f), with domains $X''\subset Y$ and labels $\mathfrak{h}_a$. Those of levels $i<j$ are
intrinsic by hypothesis~(e); those born of $\mathbf{T}$ at level $j$ are intrinsic by
hypothesis~(f).

Such a term reads $[X'']_i^{+1}$, where $i\le j$ is its level. By
Lemma~\ref{4.38:lem-round-up-touch}(c) this set is contained in $[Y]_j^{+1}$. So the reaches
do not compound: all these reading sets are contained in $[Y]_j^{+1}$, the one-cell
enlargement being taken once, at level $j$, and the enlargements at the several levels do not
add up.

\smallskip
\emph{(5) Counting terms.} These are the counts $\mathsf{N}^{(i)}$. Their coefficients
are numbers of Definition~\ref{4.38:def-standing-data}.

\smallskip
\emph{(6) Averaged trees.} These are the trees, which depend on the chosen coordinate system
\cite{B-CMP122a}, averaged over the hyperoctahedral group $W_4$. The averaged trees are
intrinsic by Definition~\ref{4.38:def-intrinsic-trace}.

\smallskip
\emph{(7) Vacuum values.} By hypothesis~(c), these are the vacuum values of terms whose
domains are contained in $Y$, and of those only.

\smallskip
\emph{Assertion (1).} Each ingredient (1)--(7) of $F(X')$ is determined by:
\begin{itemize}
\item the fields, lattices, partitions and axes on $X'$;
\item numbers of Definition~\ref{4.38:def-standing-data};
\item the trace of $\ell$ on $[Y]^{+1}$ for domains $Y\subset X'\subset X$.
\end{itemize}
Since $Y\subset X$, we have $[Y]\subset[X]$, and every cell touching $[Y]$ lies in $[X]$ or
touches it. Hence $[Y]^{+1}\subset X^{+1}$.

The index set of the series \eqref{4.38:eq-label-remainder-a} is
$\mathcal{C}_j(X;\ell)$. By \eqref{4.38:eq-label-remainder-b} it equals the catalogue formed
with the trace of $\ell$ on $X^{+1}$ in place of $\ell$, a function of that trace.
The weights of the series are functions of incidences. By
Definition~\ref{4.38:def-intrinsic-trace}, $R'^{(j),\ell}(X)$ is therefore intrinsic, with
reading set $X^{+1}$. This proves assertion~(1).

\smallskip
\emph{Assertion (2): equality on the set of empty trace.}
Let $X\in\mathcal{X}_j(\ell)$. By \eqref{4.38:eq-label-remainder-b},
$\mathcal{C}_j(X;\ell)=\mathcal{C}_j(X;\emptyset)$, so the two series have the same index
set. We show that the activities of every polymer also agree.

\emph{Non-island factors.} The cells of $[\Lambda_i(\ell)^c]_i$ lie in $\Lambda_j(\ell)^c$,
by the nesting $\Lambda_j(\ell)\subset\Lambda_i(\ell)$ of the retained sequence and since
level-$i$ cells lie inside level-$j$ cells.
By the definition of $\mathcal{X}_j(\ell)$ in \eqref{4.38:eq-label-remainder-b}, $X^{+1}$
holds no cell of $[\Lambda_j(\ell)^c]$. Hence, for every $i\le j$, $X^{+1}$ holds no cell of
$[\Lambda_i(\ell)^c]_i$, and neither does $[X]_i^{+1}$, which is contained in $X^{+1}$ by
Lemma~\ref{4.38:lem-round-up-touch}(c). The cells of $[Z''_i\cap Y_b]_i$ lie in
$Y_b\subset\Lambda_j(\ell)^c$, and so they also lie outside $X^{+1}$.

Lemma~\ref{4.38:lem-coupling-profile}(e), together with
Lemma~\ref{4.38:lem-round-up-touch}(e), then gives that every non-island factor, of whichever
sequence, is identically $1$ on $X$ under $\ell$; under the empty label there are no such
cells at all, and the non-island factors are identically $1$ as well.

\emph{Island factors.} These are the same functions under $\ell$ and under $\emptyset$. They
are formed from the islands of the polymer, which are the same contents in both catalogues.

\emph{The cut-off.} Write $\eta_\ell$ for the small-field cut-off function attached to the
label $\ell$ in the series \eqref{4.38:eq-label-remainder-a}, and $\eta_\emptyset$ for the one
attached to the empty label. They coincide on $X$, because
$X\cap\bigcup_bY_b=\emptyset$: indeed $Y_b\subset\Lambda_j(\ell)^c$, while $X^{+1}\supset X$
holds no cell of $[\Lambda_j(\ell)^c]$, as shown in the preceding paragraphs.

Hence $R'^{(j),\ell}(X)=\mathfrak{r}'^{(j)}(X)$ term by term: the index sets are the same, and
the weights are the same, the members carrying the profile included. This proves
assertion~(2).

\smallskip
\emph{Assertion (3): the range.}
By Lemma~\ref{4.38:lem-range-empty-trace}, the second-group range is $\mathcal{X}_j(\ell)$.
On this set the annulled difference $B'_{\mathrm{ann}}$ vanishes, by
\eqref{4.38:eq-torus-independence-b}. Assertion~(2) gives
$R'^{(j),\ell}=\mathfrak{r}'^{(j)}$ there. Therefore
\begin{equation*}
R^{(j)}=R''^{(j)}+\mathfrak{r}'^{(j)}=\mathfrak{r}^{(j)}
\end{equation*}
on $\mathcal{X}_j(\ell)$: the first equality is the decomposition
$R^{(j)}=R''^{(j)}+R'^{(j),\ell}$ of part (iii)(b) of the statement on torus-independence of
the localized terms together with assertion~(2), and the second is
\eqref{4.38:eq-torus-independence-b} with $B'_{\mathrm{ann}}=0$. This is
\eqref{4.38:eq-label-activities-1}: the function at the
empty label, one family for all labels. This proves assertion~(3).

\smallskip
\emph{Assertion (4): site, lift, translations.}
Let $X\in\mathbf{D}_j$ be non-wrapping in the sense of Definition~\ref{4.38:def-non-wrapping}.
The term $\mathfrak{r}^{(j)}(X)$ is intrinsic with the empty label, by assertion~(1) applied
with the empty label, hypothesis~(f) for $R''^{(j)}$, and \eqref{4.38:eq-label-activities-1}.
By Lemma~\ref{4.38:lem-round-up-touch}(d), $X^{+1}$ lies within $2M\check{R}_j$ of $X$. Since
the strips consist of $m_{\sharp}$ layers of $M$-cubes and $m_{\sharp}M<\tfrac12M\check{R}_j$
by \eqref{4.38:eq-standing-data-b}, $(X^{\sharp})^{+1}$ lies within
$2M\check{R}_j+m_{\sharp}M<3M\check{R}_j$ of $X$.
So $\mathfrak{r}^{(j)}(X)$ is determined by the lattices, partitions, axes and fields
on $X$ with the collar $2M\check{R}_j+m_{\sharp}M$.

By Definition~\ref{4.38:def-non-wrapping}, a lift is an isometric homeomorphism on a cube
containing $X$ with this collar. It preserves lattices, partitions, incidences (touching
included), axes and tree lengths. By Definition~\ref{4.38:def-intrinsic-trace}, every
ingredient is transported by the lift. By Lemma~\ref{4.38:lem-coupling-profile}(c), the
profile is transported as well. Transport of structure therefore shows that
$\mathfrak{r}^{(j)}(X)$ depends neither on the site nor on the lift. The same transport shows
that it is covariant under the translations $\tau_t$, $t\in M\check{R}_j\mathbb{Z}^4$; these
preserve the level-$i$ cell partitions for all $i\le j$, since, by the choice of the
cell-size multipliers $\check{R}_i$ among the numbers of Definition~\ref{4.38:def-standing-data},
$L^{j-i}\check{R}_j/\check{R}_i$ is a power of $L$, so that $a_i$ divides $M\check{R}_j$.
This proves assertion~(4).
\end{proof}

\begin{lemma}[Label-independence: boundary terms and large-field operations]
\label{4.38:prf-label-boundary}
Fix a level $j$ with $x_j\ge\gamma^{-2}$ such that step $j$ has been completed at every site
of $\mathcal{S}_j$; fix a site $\mathfrak{s}\in\mathcal{S}_j$ and a label $\ell$ at $\mathfrak{s}$. Let
$[\,\cdot\,]_i$ and $(\cdot)_i^{+1}$ be formed with the level-$i$ label cubes, of side
$a_i=M\check{R}_iL^{-(j-i)}$, as in
Lemma~\ref{4.38:lem-round-up-touch}, and let $W^{+1}$ and the trace of a label on $W^{+1}$ be
as in \eqref{4.38:eq-intrinsic-trace-a}. Assume the following.
\begin{enumerate}
\item[(1)] The representation and bound of the large-field terms holds at level $j$;
among its constituents are the following:
\begin{itemize}
\item all expressions of the integral operation of the large-field step are localised in
the domain of that operation \cite{B-CMP122b};
\item the terms $V(Y,U_j)$ (the terms $V(Y,U_k)$ of
\cite[p.~379]{B-CMP122b}, whose $k$ is the present $j$) depend on the large-field sequence
only through its restriction to the components of $Z$ contained in $Y$;
\item the vacuum values entering a term with domain $Y$ are those of terms with domains
contained in $Y$ only;
\item the result of the large-field step has again the form of the representation
\cite[(2.18)]{B-CMP119}, as stated in \cite{B-CMP122b};
\item the items indexed by points of the small-field region of the pre-step label near an
island foreign to the label are assigned to the boundary terms or to the pointed localised
terms $E^{(j)}$ as in \cite{B-CMP119}, read with \cite[(1.59), (1.60)]{B-CMP122a};
\item at the prepared sequence $\{Z''_i\}$ the members built with the coupling profile involve
the level-$j$ cells covering $Z''_j$ (in the notation of \cite{B-CMP122a}, whose $k$ is the
present $j$), which need not lie in the box $\Lambda^{(1.73)}$ of
\cite[(1.73)]{B-CMP122a}; these members are localised inside the island $X_a$ of the
large-field step on which they are built.
\end{itemize}
\item[(2)] The gluing and restriction statement for labels: the clauses (c1), (c2), (c3) of
\eqref{4.38:eq-label-remainder-c} hold if and only if $\ell$ glued with the pairs
$(X_a,\mathfrak{h}_a)$ is an admissible pre-$\mathbf{R}$ label with the $X_a$ first-class.
\item[(3)] The bound on the large-field part of the coupling increment $\beta_j$ (the increment
being that of Definition~\ref{4.2:def-increment}) holds at level $j$, and the conclusion of the
renormalization step holds at level $j-1$.
\end{enumerate}
Let $X\in\mathbf{D}_j$ be a non-wrapping localisation domain at $\mathfrak{s}$. Then
$B^{(j)}(X;\cdot)$, $\mathbf{B}'^{(j)}(X;\cdot)$ and $\mathbf{T}^{(j-1)}(\cdot\cap X)$, with
$W:=X$, and $\mathbf{T}'_j(X)$, with $W:=X^{\sharp}$, depend on $(\mathfrak{s},\ell)$ only
through the trace of $\ell$ on $W^{+1}$. Together with Lemma~\ref{4.38:prf-label-activities}
and Lemma~\ref{4.38:prf-label-remainder}, the induction hypothesis at level $j+1$ holds in
its clauses (a) to (c$^{+}$).
\end{lemma}

\begin{proof}
The four objects are built from the same lists of ingredients as the remainder series
$R'^{(j),\ell}(X)$ of \eqref{4.38:eq-label-remainder-a}. These lists were set out in the proof
of Lemma~\ref{4.38:prf-label-activities}:
\begin{itemize}
\item integrations over fields in islands;
\item operators with footprints in the domains $Y$;
\item the members bearing the coupling profile;
\item older terms;
\item the counts and the averaged trees;
\item vacuum values.
\end{itemize}
It was shown there that each of them reads the label only on the one-cell enlargement of the
domain in which it is localised. Here the label is carried along, and we determine, object by
object, the set $W$ on which it is read.

\emph{The boundary term $B^{(j)}$.} By \cite[(2.42)]{B-CMP119}, the term $B^{(j)}(X;\cdot)$ is
written as
\begin{equation*}
B^{(j)}\bigl(X,(\mathbf{U},\mathbf{J}),A,\{S_i\cap X\}\bigr).
\end{equation*}
It is assembled from the lists above with the label carried. The sets $S_i\cap X$ are read on
$X$, and hence on $X^{+1}$; the members bearing the coupling profile read the label on
$X^{+1}$ by part (e) of Lemma~\ref{4.38:lem-coupling-profile}, as recalled below.

\emph{The operation $\mathbf{T}^{(j-1)}$.} The displays \cite[(2.19)--(2.21)]{B-CMP119}
factorise the operation $\mathbf{T}$ over components. Hence $\mathbf{T}^{(j-1)}(\cdot\cap X)$
is built from the ingredients of the lists above, localised in $W=X$.

\emph{The term $\mathbf{B}'^{(j)}$.} The term $\mathbf{B}'^{(j)}(X)$ is the first group of terms
displayed in \cite{B-CMP119}. For $X\notin\mathcal{X}_j(\ell)$ it coincides
with the remainder series $R'^{(j),\ell}(X)$. By the catalogue identity
\eqref{4.38:eq-label-remainder-b},
\begin{equation*}
\mathcal{C}_j(X;\ell)=\mathcal{C}_j\bigl(X;\operatorname{tr}_{X^{+1}}\ell\bigr).
\end{equation*}
That identity was obtained from hypothesis (2) together with
Lemma~\ref{4.38:lem-round-up-touch}. By the proof of Lemma~\ref{4.38:prf-label-activities}, every
term of the series over this catalogue reads the label on $X^{+1}$. Hence
$\mathbf{B}'^{(j)}(X)$ depends on $(\mathfrak{s},\ell)$ only through the trace of $\ell$ on
$X^{+1}$. For $X\in\mathcal{X}_j(\ell)$ no term is assigned to $\mathbf{B}'^{(j)}(X)$: the
difference $R'^{(j),\ell}(X)-\mathfrak{r}'^{(j)}(X)$ vanishes by
Lemma~\ref{4.38:prf-label-activities}.

\emph{The operation $\mathbf{T}'_j$.} The operation $\mathbf{T}'_j(X)$ contains $V(X^{\sharp})$,
so here $W:=X^{\sharp}$. The older terms of $V(X^{\sharp})$ have domains $X''\subset
X^{\sharp}$. By part (c) of Lemma~\ref{4.38:lem-round-up-touch}, applied to
$X''\subset X^{\sharp}$, each of them reads the label on
\begin{equation*}
[X'']_i^{+1}\subset (X^{\sharp})^{+1}.
\end{equation*}
When such an operation occurs inside $R'^{(j),\ell}(X)$ with an island $X_a$, clause (c2) of
\eqref{4.38:eq-label-remainder-c} gives $X_a^{\sharp}\subset X'\subset X$. The same part of
Lemma~\ref{4.38:lem-round-up-touch} then gives
\begin{equation*}
(X_a^{\sharp})^{+1}\subset X^{+1}.
\end{equation*}
So in that position it reads the label on $X^{+1}$.

\emph{The members bearing the coupling profile.} By part (e) of
Lemma~\ref{4.38:lem-coupling-profile}, on a plaquette $p\subset W$ the value $\varphi_i(x(p))$ is
determined by the cells of $[\Lambda_i(\ell)^c]_i$ within $\tfrac14 a_i$ of $x(p)$. All of these
cells lie in
\begin{equation*}
[p]_i^{+1}\subset W^{+1}.
\end{equation*}

In each of the four cases, then, every ingredient reads the label on $W^{+1}$. This proves the
first assertion.

The proof of this lemma uses the following inputs:
\begin{itemize}
\item the series representation \eqref{4.38:eq-label-remainder-a} of the remainder, established in
Lemma~\ref{4.38:prf-label-remainder};
\item the catalogue identity \eqref{4.38:eq-label-remainder-b}, which rests on hypothesis (2)
and on Lemma~\ref{4.38:lem-round-up-touch};
\item for the lists of ingredients, hypothesis (1) with all of its constituents, among them
the part of the statement on the one-power hereditary class that fixes the form of the terms,
as these appear in \cite[(1.60)--(1.100)]{B-CMP122a} and \cite[(1.14)--(1.48)]{B-CMP122b}.
\end{itemize}
The coupling profile of Lemma~\ref{4.38:lem-coupling-profile}, and its reading at the
large-field step, belong to the setting and add no hypothesis.

Parts (i) and (ii) of the label-independence assertion at level $j$ are
Lemma~\ref{4.38:prf-label-activities} and Lemma~\ref{4.38:prf-label-remainder}. Part (iii) is
the first assertion proved above. Together with hypothesis (3), these give clauses (a) to
(c$^{+}$) of the induction hypothesis at level $j+1$.
\end{proof}

\begin{remark}[The trace needs the one-cell enlargement: a counterexample]
\label{4.38:rem-enlargement-counterexample}
In Definition~\ref{4.38:def-intrinsic-trace} a label is read through its trace on
$W^{+1}=[W]\cup(\text{the label cubes touching }[W])$; see \eqref{4.38:eq-intrinsic-trace-a}.
For a domain $W$ that is a union of $M$-cubes, the smaller set
$W\cup(\text{the cells touching }W)$ does not suffice. The following configuration at
level $j$ shows this.

We index the cells of level $j$ (the cubes of side $M\check{R}_j$) by $n\in\mathbb{Z}^4$,
writing $Q(n)$ for the cell $M\check{R}_j n+[0,M\check{R}_j]^4$, and we put $e:=(1,0,0,0)$.
For a union $A$ of cells, $A^{\sim}$ denotes $A$ together with the cells touching $A$.
As a choice made for the example, we take $\check{R}_j\ge 5$ odd and $m_{\sharp}\ge 2$.
\begin{itemize}
\item[(a)] $\mathsf{C}$ is a box which is a first-class island, with history $\mathfrak{h}$,
contained in $\{n_1\le -2\}$ and containing $Q(-2e)$.
\item[(b)] $\ell$ is a label whose only large-field set at level $j$ is a single set
$Y_1=Y_1(\ell)$ among the sets $Y_i(\ell)$ of clause~(c4) of
\eqref{4.38:eq-label-remainder-c}. Here $Y_1$ is a union of cells with
$Q(e)\subset Y_1\subset\{n_1\ge 1\}$.
\item[(c)] $Y$ is the union of $\mathsf{C}^{\sharp}$ and of the row of $M$-cubes of
$Q(-e)\cup Q(0)$ on the central $x_1$-axis of these two cells, up to and including the
$M$-cube $c\subset Q(0)$ at distance $M$ from $Q(e)$. Then $Y\in\mathbf{D}_j$, and we put
$X:=Y$.
\end{itemize}
Then the following hold.
\begin{itemize}
\item[(1)] No cell of $\ell$ shares a point with $X$. Hence the trace of $\ell$ on
$X\cup(\text{the cells touching }X)$ is empty.
\item[(2)] The following membership holds:
\begin{equation*}
(\{(\mathsf{C},\mathfrak{h})\},\{Y\})\in\mathcal{C}_j(X;\emptyset)\setminus
\mathcal{C}_j(X;\ell).
\end{equation*}
In particular $\mathcal{C}_j(X;\ell)\neq\mathcal{C}_j(X;\emptyset)$.
This is an inequality between catalogues. Nothing is asserted about the values of the
corresponding series.
\item[(3)] $Q(0)\subset[X]$ and $Q(e)\subset X^{+1}$, so the trace of $\ell$ on $X^{+1}$ sees
$Y_1$. Moreover $Q(0)\subset(\Lambda_j(\ell)^{c})^{\sim}$, and $X\notin\mathcal{X}_j(\ell)$.
\item[(4)] (Margin of zero cells.) Keep $\mathsf{C}$ and $Y_1$ three cells apart, as above.
Let $X_3$ be the union of $\mathsf{C}^{\sharp}$ and of the central row up to the wall
$\{x_1=0\}$. Then $X_3\in\mathcal{X}_j(\ell)$, and the last $M$-cube of the row shares a wall
with $Q(0)\subset Y_1^{\sim}$. Further,
\begin{equation*}
\min_{x\in X_3}\operatorname{dist}_{\infty}(x,Y_1)=M\check{R}_j .
\end{equation*}
\end{itemize}
\end{remark}

\begin{proof}
A first-class island is never a single cell: \cite{B-CMP122a} requires of a first-class
domain that ``the size must be greater than $20MR_k$''.
Suppose instead that the class of first-class domains were empty. Then the remainder series
would vanish identically, and there would be nothing to compare. This is why (a) takes
$\mathsf{C}$ to be a box larger than one cell, placed by
$Q(-2e)\subset\mathsf{C}\subset\{n_1\le-2\}$.

The label $\ell$ is admissible beside $\mathsf{C}$, because $\mathsf{C}\subset\{n_1\le-2\}$ and
$Y_1\subset\{n_1\ge 1\}$ share no point. This is the gluing statement, part~(a) of
Lemma~\ref{4.38:lem-growth-factorise}.

Since $\check{R}_j$ is odd, the $M$-cubes of $Q(-e)\cup Q(0)$ meeting the central $x_1$-axis
in their interior form a single row. In the coordinates above, $c$ is the $M$-cube of this row
with $x_1$-range $[M\check{R}_j-2M,\,M\check{R}_j-M]$.

(1) The cells of $\ell$ lie in $Y_1\subset\{x_1\ge M\check{R}_j\}$. The set $X=Y$ lies in
$\{x_1\le M\check{R}_j-M\}$. So no cell of $\ell$ shares a point with $X$, and the trace of
$\ell$ on $X$ together with the cells touching $X$ is empty.

(2) The $M$-cube $c$ lies at distance $M$ from $Q(e)\subset Y_1$. It therefore belongs to the
second layer of $M$-cubes next to $Y_1$, and since $m_{\sharp}\ge2$ we have
$c\subset Y_1^{\sharp}$. On the other hand $Y_1^{\sharp}\not\subset Y$.

Consider clause~(c4) of \eqref{4.38:eq-label-remainder-c} for the domain $Y$, with
$Z:=\bigcup X_a\cup\bigcup Y_i(\ell)$.
\begin{itemize}
\item Under the empty label, $Z=\mathsf{C}$. Then $Y\cap Z^{\sharp}=\mathsf{C}^{\sharp}$, which
is a union of components of $Z^{\sharp}$, at least one. So $Y$ satisfies~(c4).
\item Under $\ell$, $Z=\mathsf{C}\cup Y_1$. Since $m_{\sharp}M<\tfrac12 M\check{R}_j$ and
$\mathsf{C}$, $Y_1$ are three cells apart, $Z^{\sharp}$ has the two components
$\mathsf{C}^{\sharp}$ and $Y_1^{\sharp}$. The set $Y$ meets $Y_1^{\sharp}$, at $c$, but does
not contain it. So $Y\cap Z^{\sharp}$ is not a union of components of $Z^{\sharp}$, and $Y$
violates~(c4).
\end{itemize}
Clauses (c1)--(c3) of \eqref{4.38:eq-label-remainder-c} hold for this element under both
labels: $\mathsf{C}$ is a first-class island by~(a); $\mathsf{C}^{\sharp}\subset Y=X$; and
$\mathsf{C}$ shares no point with a cell of $Y_1$, being three cells away from it. The two
labels are separated by clause~(c4). Hence the
element lies in $\mathcal{C}_j(X;\emptyset)$ and not in $\mathcal{C}_j(X;\ell)$. By~(1), the
trace on $X\cup(\text{the cells touching }X)$ does not distinguish $\ell$ from the empty label.
So that set cannot replace $X^{+1}$ in the catalogue identity of
\eqref{4.38:eq-label-remainder-b}: clause~(c4), through the strips $Y_i^{\sharp}$, reads the
label beyond $X\cup(\text{the cells touching }X)$.

(3) The set $X$ contains $M$-cubes of $Q(0)$, for instance $c$. Hence $Q(0)\subset[X]$. The
cell $Q(e)$ touches $Q(0)$, so $Q(e)\subset X^{+1}$. Since $Q(e)\subset Y_1$ is a cell of
$[\Lambda_j(\ell)^{c}]$, the trace of $\ell$ on $X^{+1}$ is not empty. This agrees with the
identity of \eqref{4.38:eq-label-remainder-b}, by which $\mathcal{C}_j(X;\ell)$ depends on
$\ell$ only through the trace of $\ell$ on $X^{+1}$. It also gives $X\notin\mathcal{X}_j(\ell)$,
and equally $Q(0)\subset(\Lambda_j(\ell)^{c})^{\sim}$. So $X$ lies outside the range of the
second group of terms, which by Lemma~\ref{4.38:lem-range-empty-trace} is
$\mathcal{X}_j(\ell)$. The example therefore does
not contradict the equality of catalogues on $\mathcal{X}_j(\ell)$ stated in
\eqref{4.38:eq-label-remainder-b}.

(4) The strips are thinner than half a cell, $m_{\sharp}M<\tfrac12 M\check{R}_j$. Hence
$\mathsf{C}^{\sharp}\subset\{n_1\le-1\}$, and the row defining $X_3$ lies in $Q(-e)$. Thus
$[X_3]$ consists of cells with $n_1\le -1$, and the cells touching them have $n_1\le 0$. So
$X_3^{+1}$ holds no cell of $Y_1\subset\{n_1\ge1\}$, the trace of $\ell$ on $X_3^{+1}$ is
empty, and $X_3\in\mathcal{X}_j(\ell)$.

The last $M$-cube of the row has its face on $\{x_1=0\}$, and so shares a wall with $Q(0)$.
The cell $Q(0)$ touches $Q(e)\subset Y_1$, so $Q(0)\subset Y_1^{\sim}$. Every point of $X_3$
has $x_1\le 0$, while $Y_1\subset\{x_1\ge M\check{R}_j\}$. Hence
$\operatorname{dist}_{\infty}(x,Y_1)\ge M\check{R}_j$ for all $x\in X_3$. Equality holds at
the points of that face which lie at sup-distance $M\check{R}_j$ from $Q(e)$.

Thus the distance from $X_3$ to $Y_1$ is exactly one cell side, with nothing to spare: every
dependence on the label reaches strictly less than one cell from $X_3$, and so just fails to
reach $Y_1$. Every such dependence of the terms on $X_3$
is read from a point or an $M$-cube of $X_3$ and reaches less than one cell from there:
\begin{itemize}
\item the clauses (c3)--(c5), which by the discussion following
\eqref{4.38:eq-label-remainder-c} read the label only on $X_3^{+1}$;
\item the clauses read at a level $i$, by part~(e) of Lemma~\ref{4.38:lem-round-up-touch}
applied at that level;
\item the coupling profile, by part~(e) of Lemma~\ref{4.38:lem-coupling-profile}, which reads
less than a quarter of a cell.
\end{itemize}
The distance $M\check{R}_j$ to $Y_1$ is not less than one cell, so none of these reaches it.

Both conventions of Definition~\ref{4.38:def-intrinsic-trace} enter at this configuration,
and both are part of that definition:
\begin{itemize}
\item to touch means to share a point, so that contact at a corner counts;
\item the trace is read cube-wise, so that contact along a wall or at a corner contributes no
member.
\end{itemize}
\end{proof}

\begin{remark}[Unchanged remainder split and unchanged coupling increments]
\label{4.38:rem-split-unchanged}
Fix a level $j$, a site $\mathfrak{s}\in\mathcal{S}_j$ and a label $\ell$ at $\mathfrak{s}$. By the
\emph{range} we mean the set of non-wrapping localization domains $X\in\mathbf{D}_j$, unions of
$M$-cubes (cell-unions included), with $X\subset\Lambda_j(\ell)^{\sim-1}$, as in
\eqref{4.38:eq-torus-independence-b}. Write
\[
R^{(j)}(X;\cdot)=R''^{(j)}(X;\cdot)+R'^{(j)}(X;\cdot)
\]
for the remainder terms of the representation \cite[(2.18)]{B-CMP119}, and $R'^{(j),\ell}(X)$ for
the second family read at the label $\ell$.
\begin{itemize}
\item[(a)] For every label $\ell$ the identity
\begin{equation}\label{4.38:eq-split-unchanged-a}
\begin{split}
\sum_{X}R'^{(j),\ell}(X)
={}&\sum_{X\in\mathrm{range}}\mathfrak{r}'^{(j)}(X)
+\sum_{X\in\mathrm{range}}B'_{\mathrm{ann}}(X)\\
&+\sum_{X\notin\mathrm{range}}\mathbf{B}'^{(j)}(X)
\end{split}
\end{equation}
holds, and the middle sum vanishes identically. Consequently the action $A_j$, the
density $\rho_j$, its integral $\int\rho_j$, the polarisation tensors
$\Pi_{\mathfrak{s}}^{(j)}$, every total and every bound are the same for every label $\ell$; in
particular they coincide with their values computed without labels.
The bounds in question are those of \cite[(1.97)--(1.100)]{B-CMP122b} and those expressed
through $d_j$ and $\kappa$.
\item[(b)] The numbers $\Pi_{\infty}^{(j)}$, $\beta_j^{\mathfrak{A}}$, $\beta_j^{\Re}$, $x_j$,
$g_j$, $\check{R}_j$ and the quantity $k_1(j)$ fixed together with $g_j$ and $\check{R}_j$ from
$x_j$ are the same numbers for every label, and the site class $\mathcal{S}_j$ and the integer $K$
are likewise unchanged: they do not
depend on the label $\ell$, on the round-up $[W]$, on the one-cell enlargement $W^{+1}$, on the
class $\mathcal{X}_j(\ell)$, on the split of the range terms into $\mathfrak{r}'^{(j)}$ and
$B'_{\mathrm{ann}}$, or on the cutoff $\psi_{\varphi}$ of the coupling profile.
\end{itemize}
\end{remark}

\begin{proof}
(a) The partition of the terms $R'^{(j),\ell}(X)$ into those with $X$ in the range and those with
$X$ off the range is the split of \cite[(1.98)--(1.101), p.~390]{B-CMP122b}. In that split the
off-range terms are the terms $\mathbf{B}'^{(j)}(X)$, which gives the last sum in
\eqref{4.38:eq-split-unchanged-a}. For $X$ in the range, the difference $B'_{\mathrm{ann}}$ is
defined in \eqref{4.38:eq-torus-independence-b} by
\[
B'_{\mathrm{ann}}(X)=R'^{(j),\ell}(X)-\mathfrak{r}'^{(j)}(X).
\]
Hence $R'^{(j),\ell}(X)=\mathfrak{r}'^{(j)}(X)+B'_{\mathrm{ann}}(X)$ term by term on the range. This
gives the first two sums in \eqref{4.38:eq-split-unchanged-a}, and the identity follows.

By \eqref{4.38:eq-torus-independence-b}, $B'_{\mathrm{ann}}(X)$ vanishes whenever the trace of
$\ell$ on $X^{+1}$ is empty. Here $X^{+1}$ is $[X]$ together with the cells touching $[X]$, and
the condition is that $X\in\mathcal{X}_j(\ell)$. By the lemma stating that the range consists
exactly of the domains $X$ on whose enlargement $X^{+1}$ the trace of $\ell$ is empty, the class
$\mathcal{X}_j(\ell)$ is exactly the
range. Therefore $B'_{\mathrm{ann}}(X)=0$ for every $X$ in the range, and the middle sum of
\eqref{4.38:eq-split-unchanged-a} is identically zero.

Thus the right-hand side of \eqref{4.38:eq-split-unchanged-a} is Ba{\l}aban's own split. No term
is created, none is dropped, and none is moved from one family to another. In particular
\cite[(2.41)]{B-CMP119} is used as stated there. The quantities listed in (a) are defined from
this representation and from the totals built on it, and the bounds listed there are bounds on
the terms of this representation. Hence these quantities and bounds do not depend on the label
$\ell$ and coincide with their values computed without labels.

(b) By Lemma~\ref{4.38:prf-increment-site-free}, the increment $\beta_j$ is read off the
site-free family $\{E^{(j)}(\hat X;\cdot)\}$ together with the walk limit $\mathsf{h}_{\infty}$.
In the localized series $\Psi_X$ a label enters only through the remainder terms $R^{(i)}(X')$,
$X'$ a localization domain at some level $i$, occurring in the activities of $\Psi_X$, and these
are read at the empty label. These terms are
\[
R^{(i)}(X')=R''^{(i)}(X')+\mathfrak{r}'^{(i)}(X'),
\]
a series over the catalogue $\mathcal{C}_i(X';\emptyset)$.

At $\ell=\emptyset$ the conditions defining the catalogue $\mathcal{C}_i(X';\ell)$ that read the
label are void, on whatever set they are read. These are the third of the conditions defining
the catalogue in Lemma~\ref{4.38:prf-label-remainder} and the parts of the fourth and fifth of
those conditions that concern the sets $Y_b$.
Moreover, every profile $\varphi_i^{\emptyset}$ read at the empty label is identically $1$,
because the empty label prescribes
no cells for the profiles to read. The profiles of the virtual islands that occur inside the
activities of
$\mathfrak{r}'^{(i)}$ are given by the same choice of coupling profile under $\ell$ and under
$\emptyset$ alike. Hence neither the index set nor the weights of this series involve any of the
following: the round-up $[W]$, the one-cell enlargement $W^{+1}$, the classes
$\mathcal{X}_i(\cdot)$, or the split into $\mathfrak{r}'^{(i)}$ and $B'_{\mathrm{ann}}$ at any
level.

The cutoff $\psi_{\varphi}$ of the profile enters no number. The individual weights of
$\mathfrak{r}'^{(i)}$ carry island factors that are not identically $1$, and termwise
independence of $\psi_{\varphi}$ is not claimed; only the independence of totals is used, as
follows. By Lemma~\ref{4.38:prf-increment-site-free},
$\Pi_{\mathfrak{s}}^{(j)}$ is the second variation of a total. By the lemma that the choice of
coupling profile moves no total (Lemma~\ref{4.38:lem-profile-telescopes}), that total is free of
$\varphi$. The identity used there is \eqref{4.38:eq-profile-telescopes-a}: the two
$\varphi$-bearing members of the exponent sum to $-x_0A(U)$ plus a $\varphi$-free term.

It follows that $\Pi_{\infty}^{(j)}$ and the two parts $\beta_j^{\mathfrak{A}}$, $\beta_j^{\Re}$
of the increment are the same numbers, whatever the label and whatever the cutoff
$\psi_{\varphi}$. Then $x_j=x_{j-1}-\beta_j$ is the same number, and so are $g_j$,
$\check{R}_j$ and $k_1(j)$, which are fixed from $x_j$; the site class $\mathcal{S}_j$ and the
integer $K$ are likewise unchanged.

Independently of the analysis of index sets and weights above, the independence of the label
also follows from Lemma~\ref{4.38:prf-increment-site-free} and part (a):
$\Pi_{\mathfrak{s}}^{(j)}$ is the second variation of a total, and by (a) no total depends on the
label. On this route, too, the independence of $\psi_{\varphi}$ rests on
\eqref{4.38:eq-profile-telescopes-a}.
\end{proof}

\begin{remark}[The locality induction at one-cell reach]\label{4.38:rem-locality-induction}
Fix $j\ge 1$. The induction of this section proceeds in three stages.
\begin{itemize}
\item From the induction hypothesis at level $j$ it obtains three statements at level $j$: the
intrinsic character of the whole-lattice terms $E^{(j)}_{\mathfrak{s}}(X;\cdot,z)$, the convergence of
the tensors $\Pi^{(j)}_{\mathfrak{s}}$ to $\Pi^{(j)}_{\infty}$, and the site-freeness of $\beta_j$.
\item It then performs the steps of the induction that advance the flow and fix $\check{R}_j$ and
$\mathcal{S}_j$. If moreover $x_j\ge\gamma^{-2}$, the operations $\mathbf{T}$ and $\mathbf{R}$ are
carried out at every admissible site, and from this it obtains the locality statement at level $j$:
every term or operation of the representation of $\rho_j$ that is localized in a non-wrapping $W$ is
intrinsic and depends on the label $\ell$ only through its trace on the one-cell enlargement $W^{+1}$.
\item From there it obtains the induction hypothesis at level $j+1$.
\end{itemize}
The locality statements at the levels $i<j$ are used in this induction at exactly six places:
\begin{enumerate}
\item[(a)] in the torus-independence of the non-wrapping localized terms, at the empty label, where
the older remainder terms $\mathfrak{r}^{(i)}$ must be carried by lifts and by the translations
$\tau_t$;
\item[(b)] in the whole-lattice terms, for the inherited terms $R^{(i)}(X';\cdot)$ at the vertices
$X'\subset X\subset\Lambda_j(\ell)$;
\item[(c)] in the part $R''$ of the remainder, for the same inherited terms;
\item[(d)] for the older terms carried by the islands' histories $\mathfrak{h}_a$;
\item[(e)] for the older terms inside $\mathbf{B}'^{(j)}(X)$;
\item[(f)] for the older terms inside the functional $V(X^{\sharp})$ contained in $\mathbf{T}'_j(X)$
(Lemma~\ref{4.38:prf-label-boundary}).
\end{enumerate}
At each of these places the term depends on the label only through its trace on a set of cells, which
we call its \emph{reading set}; the \emph{reach} of the term is the width, in cells, of the layer by
which its reading set extends beyond the rounded-up domain. The reach is the same at every level: it
is one cell, and reaches taken at different levels combine by a maximum and not by a sum. The coupling
profile $\varphi_i$ has reach a quarter of a cell at its own level, and this lies inside the same
$W^{+1}$. Hence the clause of the induction hypothesis at level $j+1$ that asserts the locality
statements at all levels $i\le j$ is reproduced. The only condition used is $m_{\sharp}\ge 1$, i.e.\
$\check{R}_j\ge 3$; it is implied by the lower bound $\check{R}_j\ge L$ on the cell-size multiplier,
a floor which the construction already has.

Finally, the remainder $\mathfrak{r}'^{(i)}(X)$ is Ba{\l}aban's term $\mathbf{R}'^{(i)}(X)$ at the
empty admissible label. Every localization domain there belongs to the second group of terms of the
$\mathbf{R}$-operation in \cite[p.~390]{B-CMP122b}, so the bound \cite[p.~390, (1.100)]{B-CMP122b}
holds for it.
\end{remark}

\begin{proof}
Throughout, $[\,\cdot\,]_i$ and $(\cdot)^{+1}_i$ denote rounding up to the level-$i$ label cubes and
the enlargement by the cubes touching the result. The level-$i$ label cubes have side $a_i$, and
$W^{+1}=[W]^{+1}_j$.

\emph{The six places.}

(a) In the proof of Lemma~\ref{4.38:prf-proof-non-wrapping}, the locality statement at levels $i<j$ is
applied at the empty label. There it is needed in the form that each $\mathfrak{r}^{(i)}$ of
\eqref{4.38:eq-torus-independence-b} is one family, carried by lifts and covariant under the
translations $\tau_t$. This is part of the locality statement at level $i$.

(b), (c) In Lemma~\ref{4.38:prf-label-inherited}, take a vertex $X'$ with
$X'\subset X\subset\Lambda_j(\ell)$ and $i<j$. Then
\begin{equation*}
[X']^{+1}_i\subset[X]^{+1}_j\subset\Lambda_i(\ell).
\end{equation*}
The first inclusion is clause (c) of Lemma~\ref{4.38:lem-round-up-touch}. The second is the chain
\eqref{4.38:eq-label-inherited-a} at a label made by $\mathbf{T}$, which rests on the nesting of the
domains $\Omega_1,\Lambda_1,\Omega_2,\dots$ determined by the renormalization transformations, each
containing the next; at a label retained after $\mathbf{R}$ it follows from the three
facts established in Lemma~\ref{4.38:prf-label-inherited}. Hence $X'$ lies in the range
$\mathcal{X}_i(\ell)$. On that range
the locality statement at level $i$ identifies the inherited term $R^{(i)}(X';\cdot)$ with
$\mathfrak{r}^{(i)}(\hat X';\cdot)$. This is used once for the whole-lattice terms and once for the
part $R''$.

(d) In Lemma~\ref{4.38:prf-label-activities}, the activity $F(X')$ of the series
\eqref{4.38:eq-label-remainder-a} contains older terms with domains
$X''\subset Y\subset X'$ and labels $\mathfrak{h}_a$. The locality statement at level $i$ for boundary
terms and operations applies to them with reading set $[X'']^{+1}_i$, and
\begin{equation*}
[X'']^{+1}_i\subset[Y]^{+1}_j.
\end{equation*}

(e), (f) In Lemma~\ref{4.38:prf-label-boundary}, the older terms inside $\mathbf{B}'^{(j)}(X)$ (where
$W=X$) and inside $V(X^{\sharp})$ (where $W=X^{\sharp}$) are treated as follows. The locality
statements at level $i$ for boundary terms and for the remainder term apply to them with reading set
$[X'']^{+1}_i$, and
\begin{equation*}
[X'']^{+1}_i\subset W^{+1}.
\end{equation*}

\emph{The reach.} At each place the inclusion of the level-$i$ reading set in the level-$j$
enlargement is an instance of Lemma~\ref{4.38:lem-round-up-touch}. For $X''\subset W$, rounding up to
the coarser level-$j$ cubes and then adding the touching cubes contains rounding up at level $i$ and
then touching. So each level contributes a reach of one cell of its own, and the combined reach is one
cell of the coarsest level read: reaches at successive levels are not added. In (d) we have
$Y\subset X'\subset X$, so
$[Y]^{+1}_j\subset[X]^{+1}_j$ by monotonicity of $[\,\cdot\,]_j$ and of the enlargement.

The members carrying $\varphi_i$ are treated by Lemma~\ref{4.38:lem-coupling-profile}. On a plaquette
$p\subset W$, the value $\varphi_i(x(p))$ is determined by the cells of $[\Lambda_i(\ell)^c]_i$ within
$\tfrac14 a_i$ of $x(p)$. These lie in $[p]^{+1}_i$, which is contained in $W^{+1}$.

Thus every term at level $j$ depends on the label only through its trace on $W^{+1}$, which is the
reading set asserted at every level $i<j$. The locality statement at level $j$ therefore has the same
form as those below it. Together with them it is the locality clause of the induction hypothesis at
level $j+1$.

This applies in particular to the locality of boundary terms and operations at levels $i<j$, which is
part of the induction hypothesis at level $j$. Since the argument of
Lemma~\ref{4.38:prf-label-boundary} reads the label on $W^{+1}$ at level $j$ exactly as the hypothesis
asserts at the levels $i<j$, the induction closes for the boundary terms and operations as well. The
enlargement $W^{+1}$, rather than $[W]$, is needed because the clause on strips of the
catalogue $\mathcal{C}_j(X;\ell)$ (Lemma~\ref{4.38:prf-label-remainder}) reads cells of $X^{+1}$
outside $[X]$.

\emph{Compatibility with the admissible-site condition.} Rounding up and enlarging each add at most
one cell of side
$M\check{R}_j$. Hence $W^{+1}$ lies within distance $2M\check{R}_j$ of $W$, and by the definition of
$\mathcal{S}_j$ in \eqref{4.38:eq-admissible-sites-b} this is at most $s_j(\mathfrak{s})/16$ at every
admissible site. Likewise $(X^{\sharp})^{+1}$ lies within $3M\check{R}_j$ of $X$. The factor $L^{11}$
in the definition of admissible sites therefore suffices for reading the label on $W^{+1}$ and on
$(X^{\sharp})^{+1}$.

The cells of side $M\check{R}_j$ are fixed at the step of the induction at which $\check{R}_j$ is
defined. This precedes the operations $\mathbf{T}$ and $\mathbf{R}$ at level $j$, after which the
locality statements reading $W^{+1}$ are established.

For $t\in M\check{R}_j\mathbb{Z}^4$, the translation $\tau_t$ maps cells to cells. Hence
$\tau_t[W]=[\tau_tW]$.

The number $m_{\sharp}$ of layers of $M$-cubes in the strip, fixed in
\eqref{4.38:eq-standing-data-b}, satisfies $m_{\sharp}\ge 1$ exactly when $\check{R}_j\ge 3$; at
$\check{R}_j=2$ one has $m_{\sharp}=\lfloor\tfrac12\rfloor=0$. The condition $\check{R}_j\ge 3$ is
implied by the floor $\check{R}_j\ge L$ of the cell-size multiplier, which the construction already
has, since $L$ is an odd blocking factor larger than one.

\emph{The remainder at the empty label.} At the empty label the trace on $X^{+1}$ is empty for every
$X\in\mathbf{D}_i$, so the series
\eqref{4.38:eq-label-remainder-a} runs over the catalogue $\mathcal{C}_i(X;\emptyset)$ and is
Ba{\l}aban's $\mathbf{R}'^{(i)}(X)$ at the empty label, which is admissible and at which every
localization domain in $\mathbf{D}_i$
belongs to the second group of terms of the $\mathbf{R}$-operation in \cite[p.~390]{B-CMP122b}. For
such terms
\cite[p.~390, (1.100)]{B-CMP122b} states
\begin{equation}\label{4.38:eq-locality-induction-1}
\bigl|\mathbf{R}'^{(k)}(X,(\mathbf{U},\mathbf{J}))\bigr|
\le \exp\bigl(-p_0(g_k)\bigr)\exp\bigl(-\kappa\,d_k(X)\bigr),
\end{equation}
where $p_0(\cdot)$ is the degree function. Applied with $k=i$, \eqref{4.38:eq-locality-induction-1}
bounds $\mathfrak{r}'^{(i)}(X)$.

In the large-field part of the coupling increment the family
$\mathfrak{r}^{(i)}=R''^{(i)}+\mathfrak{r}'^{(i)}$ enters as follows: the part $R''^{(i)}$ enters
unchanged, and the difference $B'_{\mathrm{ann}}$ of \eqref{4.38:eq-torus-independence-b} vanishes
identically on the range $\mathcal{X}_i(\emptyset)=\mathbf{D}_i$ and contributes nothing.
\end{proof}

\begin{definition}[Families at level $j$, the site-free family and the class $\mathfrak{F}_k$]
In what follows $j\ge 0$ is a level, sites $\mathfrak{s}=(m,K')$, the tori $\mathbb{T}_j(\mathfrak{s})$ of side $s_j(\mathfrak{s})$ and their covering maps $p_{\mathfrak{s}}$ are those of Definition~\ref{4.38:def-admissible-sites}, and non-wrapping domains and their lifts are those of Definition~\ref{4.38:def-non-wrapping}. We write $D_j(\mathfrak{s})$ for the class of localization domains on $\mathbb{T}_j(\mathfrak{s})$, $D_j^{\infty}$ for the class of localization domains on $\mathbb{R}^4$, and $\tau_t$ for translation by $t$.

A \emph{family at level $j$} is a collection
\begin{equation*}
f_j=\bigl\{E^{(j)}_{\mathfrak{s}}(X,\cdot)\bigr\},
\end{equation*}
in which $\mathfrak{s}$ runs over the sites of any fixed cofinal exhaustion of tori with $\tfrac12 s_j(\mathfrak{s})\ge L^{10}M$, and $X\in D_j(\mathfrak{s})$. On the torus of each such $\mathfrak{s}$ the collection is required to have the following properties:
\begin{itemize}
\item the property \cite[(1.7)]{B-CMP109};
\item analyticity of $E^{(j)}_{\mathfrak{s}}(X,\cdot)$ on the complex neighbourhood $U^{c}_j(X,\alpha_0,\alpha_1)$ of \cite{B-CMP109};
\item the bound \cite[(1.18)]{B-CMP109} with the constants $(E_0,\kappa)$;
\item the property \cite[(1.19)]{B-CMP109};
\item reality at real points;
\item the Euclidean clause: the total $\sum_X E^{(j)}_{\mathfrak{s}}(X,\cdot)$, evaluated at the minimal configuration and current determined by a regular block configuration, is invariant under the unit translations and signed permutations of the axes of the torus;
\item the \emph{volume-consistency clause at level $j$}: for non-wrapping $X$, the function $E^{(j)}_{\mathfrak{s}}(X,\cdot)$, read on a lift, is independent of $\mathfrak{s}$ and of the lift, and is covariant under the translations $\tau_t$, $t\in M\mathbb{Z}^4$.
\end{itemize}
By the volume-consistency clause, for each $\hat X\in D_j^{\infty}$ the function $E^{(j)}_{\mathfrak{s}}(p_{\mathfrak{s}}\hat X,\cdot)$, read on the lift $\hat X$, is the same for every site $\mathfrak{s}$ of the exhaustion on whose torus $p_{\mathfrak{s}}\hat X$ is non-wrapping; we denote it by $E^{(j)}(\hat X;\cdot)$, and call the collection $\{E^{(j)}(\hat X;\cdot):\hat X\in D_j^{\infty}\}$ the \emph{site-free family} of $f_j$.

These clauses define the coefficient $\beta[f_j]$.

\begin{itemize}
\item For the pointed terms $E^{(j)}_{\mathfrak{s}}(X;\cdot,z)$, $z\in X$, whose sum over $z\in X$ is $E^{(j)}_{\mathfrak{s}}(X,\cdot)$, the coefficient $\beta[f_j]$ is defined through Lemma~\ref{4.38:prf-polarisation-limit} and Lemma~\ref{4.38:prf-increment-site-free}.
\item For the unpointed terms $E^{(j)}_{\mathfrak{s}}(X,\cdot)$ of \cite{B-CMP109}, the same clauses define it by the construction of \cite[(5.1)--(5.10)]{B-CMP109}. The construction is not applied verbatim: a domain $X$ is anchored at its cube nearest the point $x$ at which the variation is taken, and it is lifted through a fixed lift of that point. Under these clauses, together with the existence, the uniqueness among regular configurations and the analyticity of the minimal configuration near the identity block configuration \cite[Theorem~1]{B-CMP102b} and the bounds on its linearisation used in Lemma~\ref{4.38:prf-polarisation-limit}, the second derivatives, at the identity block configuration, of the map that sends a block configuration to the torus total evaluated at its minimal configuration and current converge along the exhaustion to an infinite-volume polarisation tensor with the symmetry properties \cite[(5.4), (5.6)--(5.8)]{B-CMP109}, the transversality \cite[(5.9)]{B-CMP109}, the colour-diagonal form \cite[(1.21)]{B-CMP109}, the Lie algebra of $SU(N)$ being simple, and the decay \cite[(5.10)]{B-CMP109}; and $\beta[f_j]$, defined from this tensor by \cite[(1.22)]{B-CMP109}, converges absolutely, is real, and does not depend on the pair of directions $\mu\ne\nu$ used in \cite[(1.22)]{B-CMP109}.
\end{itemize}

We put
\begin{equation*}
\mathfrak{F}_k:=\bigl\{(f_j,\beta_j)_{j\le k}:\ \beta_j=\beta[f_j]\ \text{for } j\le k\bigr\}.
\end{equation*}

For $f\in\mathfrak{F}_k$ we write $\mathbf{E}_f$ for the collection of its terms. For $g'\in(0,\gamma]$ we write $\mathsf{S}_{k+1}(g'^{-2},\mathbf{E}_f)$ for the sum of the localized terms that the $(k+1)$-st renormalization step produces from the terms $\mathbf{E}_f$ at inverse coupling $g'^{-2}$. We write $Z^{(k)}$ for the normalising integral of that step, a function of the block configuration, and $Z^{(k)}(1)$ for its value at the identity configuration.
\end{definition}

\begin{corollary}[Volume-consistent families are propagated by the step]
\label{4.38:cor-volume-families}
Assume:
\begin{itemize}
\item[(a)] the operator bounds at free current;
\item[(b)] the first part of the corollary of the layer bookkeeping on the term $\log Z^{(k)}-\log Z^{(k)}(1)$, which states that, for non-wrapping $X$, the localized term of index $X$ of $\log Z^{(k)}-\log Z^{(k)}(1)$, read on a lift, is independent of the site and of the lift, and is covariant under $\tau_t$, $t\in M\mathbb{Z}^4$;
\item[(c)] the lemma on propagation of the family clauses through the step and the second lemma of the step on the family clauses, which, under hypothesis (a), carry on the torus of each site every clause of a family other than volume-consistency from the levels $j\le k$ to level $k+1$, with the same constants; the clause of the propagation lemma that extends the sum over $X$ from the torus to an infinite lattice reads the volume-consistency clauses at the levels $j\le k$, and the propagation lemma carries the Euclidean clause from level to level by an induction;
\item[(d)] part (b) of the covariance lemma for the step, on which that induction rests.
\end{itemize}
Let $f\in\mathfrak{F}_k$ and $g'\in(0,\gamma]$. Form, on the torus of every site of the exhaustion,
\begin{equation}\label{4.38:eq-volume-families-1}
E^{(k+1)}:=\bigl[\log Z^{(k)}-\log Z^{(k)}(1)\bigr]+\mathsf{S}_{k+1}(g'^{-2},\mathbf{E}_f).
\end{equation}
On the torus of each site $\mathfrak{s}$ of the exhaustion the right side of \eqref{4.38:eq-volume-families-1} is a sum
$\sum_{X\in D_{k+1}(\mathfrak{s})}E^{(k+1)}_{\mathfrak{s}}(X,\cdot)$ of localized terms; put
$f_{k+1}:=\{E^{(k+1)}_{\mathfrak{s}}(X,\cdot)\}$, with $\mathfrak{s}$ running over the sites of the exhaustion
with $\tfrac12 s_{k+1}(\mathfrak{s})\ge L^{10}M$, which again form a cofinal exhaustion.
Then $f_{k+1}$ is a family at level $k+1$ with the same constants. Moreover
\begin{equation*}
\bigl(f,(f_{k+1},\beta[f_{k+1}])\bigr)\in\mathfrak{F}_{k+1}.
\end{equation*}
\end{corollary}

\begin{proof}
Let $f\in\mathfrak{F}_k$ and $g'\in(0,\gamma]$, and form \eqref{4.38:eq-volume-families-1} on the torus of every site of the exhaustion. The clauses other than volume-consistency are treated first; volume-consistency at level $k+1$ is treated second.

\emph{All clauses except volume-consistency at level $k+1$.} These clauses are supplied on each torus by the two lemmas of hypothesis (c): under hypothesis (a), the lemma on propagation of the family clauses through the step and the second lemma of the step on the family clauses together carry every clause of a family other than volume-consistency through one step.
From them we obtain, with the same constants, the property \cite[(1.7)]{B-CMP109}, analyticity on $U^{c}_{k+1}(X,\alpha_0,\alpha_1)$, the bound \cite[(1.18)]{B-CMP109} with $(E_0,\kappa)$, the property \cite[(1.19)]{B-CMP109}, and reality at real points.

The clause of the propagation lemma that extends the sum over $X$ from the torus to an infinite lattice uses the volume-consistency clauses at the levels $j\le k$ once, as follows. In \cite[p.~290]{B-CMP109} the sum over $X$ is extended ``to all $X\in\mathbf{D}^0_j$, where $\mathbf{D}^0_j$ is the class of localization domains constructed for the lattice $\xi\mathbb{Z}^4$''. Here $\xi=L^{-j}$ is the spacing of the fine lattice in level-$j$ units, and the class $\mathbf{D}^0_j$ is the class $D_j^{\infty}$ of localization domains on $\mathbb{R}^4$. The same extension is made in \cite{B-CMP119}.

Since $f\in\mathfrak{F}_k$, each $f_j$ with $j\le k$ satisfies the volume-consistency clause at level $j$. This clause identifies the torus terms indexed by $X\subset\square^{\sim 2}$, where $\square$ is a cube of the cover $\{\square\}$ of \cite[p.~276]{B-CMP119} and $\square^{\sim 2}$ is $\square$ enlarged by two layers, with the site-free terms $E^{(j)}(\hat X;\cdot)$, $\hat X$ a lift of $X$. The coefficient of the site-free family is $\beta[f_j]$, which is $\beta_j$ since $f\in\mathfrak{F}_k$. The expression \cite[(4.36)]{B-CMP109} and the cancellation that follows it in \cite{B-CMP109} are written for the terms indexed by $\mathbf{D}^0_j$, that is, for the site-free family; this identification is what allows them to be read for the torus terms.

The Euclidean clause at level $k+1$ follows by the induction of the propagation lemma that carries the Euclidean clause from level to level, hypothesis (c). That induction uses part (b) of the covariance lemma for the step, hypothesis (d).

\emph{Volume-consistency at level $k+1$.} The argument of Lemma~\ref{4.38:prf-proof-non-wrapping} applies without change at level $k+1$; three points need comment.
\begin{itemize}
\item For the unpointed terms, the localized series $\Psi_X$ of that lemma is the series \cite[(2.13)]{B-CMP116}.
\item That lemma requires its ingredient consisting of the local densities on $X$ to be intrinsic. This ingredient includes the inherited terms $E^{(i)}$, $i\le k$. It is intrinsic by the volume-consistency clauses at the levels $\le k$, which hold because $f\in\mathfrak{F}_k$.
\item The part $\log Z^{(k)}-\log Z^{(k)}(1)$ of \eqref{4.38:eq-volume-families-1} is supplied by hypothesis (b).
\end{itemize}
Hence, for non-wrapping $X$, $E^{(k+1)}_{\mathfrak{s}}(X,\cdot)$ read on a lift is independent of $\mathfrak{s}$ and of the lift, and is covariant under $\tau_t$, $t\in M\mathbb{Z}^4$.

\emph{Conclusion.} Thus $f_{k+1}$ has every clause of a family at level $k+1$, with the same constants. These clauses define $\beta[f_{k+1}]$. Since $f\in\mathfrak{F}_k$, the definition of $\mathfrak{F}_{k+1}$ gives $\bigl(f,(f_{k+1},\beta[f_{k+1}])\bigr)\in\mathfrak{F}_{k+1}$. The argument uses the coupling $g'$ only through the condition $g'\in(0,\gamma]$.
\end{proof}

\begin{remark}[Order of the inductive step with torus-independence]\label{4.38:rem-step-order}
In this remark we use the following names for the conclusions of
Theorem~\ref{4.38:thm-torus-independence} at a level $i$. The \emph{localized-term conclusion} at
level $i$ is its part~(i). The \emph{site conclusion} at level $i$ is the conclusion whose last
sentence is the site clause for $\natural\in\{\cdot,\Re\}$. The \emph{accompanying conclusion} at
level $i$ is the remaining one of the three conclusions which
Theorem~\ref{4.38:thm-torus-independence} derives jointly from hypothesis $(\mathrm{H}_i)$, the
other two being the localized-term conclusion and the site conclusion. The \emph{remainder
conclusion} at level $i$ is part~(iii)(b), with its items (i), (ii), (ii)$''$ and (iii).

The clause set $\mathfrak{i}(k)$ of the flow induction of Theorem~\ref{4.1:thm-clause-zero} is
read as follows, and a further clause
$(\vartheta_k)$, not part of $\mathfrak{i}(k)$, is introduced:
\begin{itemize}
\item[] clause $(\gamma_k)$ of $\mathfrak{i}(k)$ is read at every site $\mathfrak{s}\in\mathcal{S}_k$,
with the quantity denoted $E$ in that clause taken to be $0$;
\item[] clause $(\alpha_k)$ of $\mathfrak{i}(k)$ is taken to be the accompanying conclusion at all
levels $i\le k$;
\item[] the further clause $(\vartheta_k)$ is the conjunction of the localized-term conclusion, the
site conclusion and the remainder conclusion at all levels $i\le k$, and of
$f^{\mathfrak{A}}\in\mathfrak{F}_k$.
\end{itemize}

The step from level $k$ to level $k+1$ is carried out in the following order.
\begin{itemize}
\item[(1)] Clauses (a), (b) and (c$^{+}$) of hypothesis $(\mathrm{H}_{k+1})$ of
Theorem~\ref{4.38:thm-torus-independence} are the conjunction of $\mathfrak{i}(k)$, of
$(\vartheta_k)$, and of the bounds $x_i\ge\gamma^{-2}$ for $0\le i\le k$, under the condition
$\tfrac12 s_i(\mathfrak{s})\ge L^{10}M$ for the sites $\mathfrak{s}$ at every level $i\le k$, as in
Corollary~\ref{4.38:cor-volume-families}.
\item[(2)] The second and third steps of the flow induction are performed at every site
$\mathfrak{s}\in\mathcal{S}^{\flat}_{k+1}$. Clause $(\vartheta_k)$ enters them through
Corollary~\ref{4.38:cor-volume-families}. This is clause~(d) of $(\mathrm{H}_{k+1})$.
\item[(3)] The localized-term, site and accompanying conclusions are established at level $k+1$.
They replace the assertion of the first step of the flow induction that neither the number of
steps $K$ nor the torus exponent $m$ occurs.
\item[(4)] The fourth and fifth steps of the flow induction are performed: the coupling flow
(clause $(\beta_{k+1})$), the multiplier $\check{R}_{k+1}$ and the set of sites $\mathcal{S}_{k+1}$.
\item[(5a)] If $x_{k+1}\ge\gamma^{-2}$, the operation $\mathbf{T}$ is applied at every site
$\mathfrak{s}\in\mathcal{S}_{k+1}$.
\item[(6a)] For the terms produced by $\mathbf{T}$, the remainder conclusion is established in
its items (i), (ii) in the form for the remainder terms $R''$, and (iii). Here and below, $R''$ and
$R'$ are the two families of remainder terms of the large-field step which the remainder
conclusion treats in its item (ii) and in its item (ii)$''$, respectively.
\item[(5b)] The operation $\mathbf{R}$ is applied. It is applied under the vacuum condition and
the strip condition, under the hypothesis on the coupling profile as it applies during the
operation $\mathbf{R}$, and under Ba{\l}aban's bounds for old terms beside a new region, printed
without proof in \cite[p.~375]{B-CMP122b}; cited as printed. It also uses the corollary
through which the bounds \cite[(1.68), (1.79), (1.89)]{B-CMP122b} re-enter.
\item[(6b)] The remainder conclusion is established for $R'$ (in the form of item (ii)$''$), for
$\mathbf{B}'^{(k+1)}$ and $\mathbf{T}'_{k+1}$, and for the completion terms. By definition, the
completion terms are the families of the localized-term conclusion and of item (ii)$''$, as in the
completion part of the representation and bound of the large-field terms.
\end{itemize}
The order contains no circle: no item uses an object that is fixed only at a later item.
\end{remark}

\begin{proof}
We go through the objects that are fixed in the course of the step.

The coupling increment $\beta_{k+1}$ is determined in item (4), through clause $(\beta_{k+1})$. It
does not enter the localized terms $E^{(k+1)}_{\mathfrak{s}}$. Hence items (2) and (3), which
concern these terms at the sites of $\mathcal{S}^{\flat}_{k+1}$, do not use it.

The multiplier $\check{R}_{k+1}$ and Ba{\l}aban's $\varepsilon_{k+1}$ of level $k+1$ are not used
before item (4); this follows the order of Ba{\l}aban's construction in \cite{B-CMP119}, where the
region of level $k+1$ formed from the large-field region $Z_k$ is
\begin{quote}
the union of all $LMR_{k+1}$-cubes intersecting the region $Z_k$.
\end{quote}
Here $R_{k+1}$ in the quotation is Ba{\l}aban's multiplier of level $k+1$, written
$\check{R}_{k+1}$ above.
This region is formed only after the cell size of level $k+1$ is known. Items (5a)--(6b) come
after item (4). They may therefore use $\check{R}_{k+1}$ and $\varepsilon_{k+1}$, and the set of
sites $\mathcal{S}_{k+1}$ fixed at item (4).

In Ba{\l}aban's construction \cite{B-CMP119}, the coupling profile $\varphi_{k+1}$ is set after the
operation $\mathbf{T}$ and before the operation $\mathbf{R}$. It depends only on
$[\Lambda_{k+1}^{c}]$, which is known at that point, after item (5a) and before item (5b). Through
the operation $\mathbf{R}$, $\varphi_{k+1}$ is the profile of the sequence in force, in the sense of
the hypothesis on the coupling profile: first the prepared sequence of the $\mathbf{R}$-step and
then the retained one, by that hypothesis as it applies during the operation $\mathbf{R}$. Thus
item (5b), and item (6b) after it, read $\varphi_{k+1}$ only
after it has been fixed.

Finally, item (1) uses only $\mathfrak{i}(k)$, $(\vartheta_k)$ and the bounds
$x_i\ge\gamma^{-2}$ for $0\le i\le k$,
all available at level $k$. Item (2) uses $(\vartheta_k)$ only through
Corollary~\ref{4.38:cor-volume-families}. Each item therefore uses only objects fixed at level $k$
or at an earlier item of the same step, and the order contains no circle.
\end{proof}

\begin{corollary}[The number of steps does not depend on the torus]\label{4.38:cor-step-count}
Let $g_0>0$ and $g>0$ be as in Theorem~0, and write $x_k=g_k^{-2}$.
Assume the statements listed as hypotheses \textup{(1)}--\textup{(7)} of
Theorem~\ref{4.38:thm-torus-independence}, from which that theorem is proved.
\begin{enumerate}
\item[(a)] The running couplings $g_1,g_2,\dots$ are generated from $g_0$ by the increments
$\beta_1,\beta_2,\dots$ (Definition~\ref{4.2:def-increment}) without reference to any site
$\mathfrak{s}=(m,K')$ (Definition~\ref{4.38:def-admissible-sites}), for as long as
$x_k\ge\gamma^{-2}$.
\item[(b)] Consequently the first-passage number $K=K(g_0,g)$ is defined without reference to any
site, and it is the same for every torus exponent $m$.
\item[(c)] Every site $\mathfrak{s}=(m,K')$ with $K'=K$ and $m\ge m_{\mathbb{T}}(g_-(g))$ (these are
the sites of Theorem~0) satisfies, with $\mathcal{S}_j$ the set of admissible sites
of \eqref{4.38:eq-admissible-sites-b},
\begin{equation}\label{4.38:eq-step-count-1}
\mathfrak{s}\in\bigcap_{j\le K}\mathcal{S}_j .
\end{equation}
Hence the $K$ renormalisation steps performed at $\mathfrak{s}$ subtract, at step $j$, exactly
the increment $\beta_j$ of part~(a).
\item[(d)] After the $K$ steps of~(c), the quantity $E$ at $\mathfrak{s}$, which is set equal to
$0$ at every site during the induction over the steps
(Definition~\ref{4.38:def-admissible-sites}), is fixed by the condition of Theorem~0 at the
sites $(m,K)$, as stated in Definition~\ref{4.38:def-admissible-sites}.
\item[(e)] $K$ is a well-defined non-negative integer, which is the last sentence of the statement
on the coupling increments; and the last member of clause~(0)
(Theorem~\ref{4.1:thm-clause-zero}) holds.
\item[(f)] The auxiliary tori of the statement on the cell-size multipliers $\check{R}_j$,
item~(A7), fall under the volume-independence statement for sites.
\end{enumerate}
\end{corollary}

\begin{proof}
Part~(a) is the input of the joint induction of Theorem~\ref{4.38:thm-torus-independence} that
concerns the couplings; it is discharged by Lemma~\ref{4.38:prf-increment-site-free}, as follows.

The bare coupling $g_0$ is given, and it involves no site. Let $j\ge1$, and suppose that
$g_0,\dots,g_{j-1}$ have been defined without reference to any site and that $x_i\ge\gamma^{-2}$
for $i<j$. By Lemma~\ref{4.38:prf-increment-site-free}, the coefficient $\beta_j$ is then well
defined, real, and a function of $(g_0,\dots,g_{j-1})$ and of the parameters alone; hence so is
$x_j$. Moreover $x_j>0$ by clause $(\beta_j)$ of the clause set $\mathfrak{i}(j)$ of the joint
induction of Theorem~\ref{4.38:thm-torus-independence}, so that $g_j=x_j^{-1/2}$, the cell-size
multiplier $\check{R}_j$ and the set $\mathcal{S}_j$ are likewise site-free: they are the same at
every site and free of $m$ and $K'$. The recursion therefore continues, with no site declared,
for as long as $x_k\ge\gamma^{-2}$, which proves~(a).

For~(b), the first-passage number $K(g_0,g)$ of Definition~\ref{1.2:def-flow-constants} is a
function of $g$ and of the sequence $(g_k)$ of~(a) alone.
It is therefore defined without reference to any site, and in particular it does not depend on $m$.

For~(c), by the inclusion stated in Definition~\ref{4.38:def-admissible-sites} after
\eqref{4.38:eq-admissible-sites-b}, the sites $(m,K)$ of
Theorem~0, those with $m\ge m_{\mathbb{T}}(g_-(g))$, lie in $\bigcap_{j\le K}\mathcal{S}_j$. This
is~\eqref{4.38:eq-step-count-1}; so the $K$ steps can be performed at the site, and at step $j$
they subtract the increment $\beta_j$, which by~(a) is the same at every site. This proves~(c).

For~(d): by~(c) the $K$ steps are performed at every site $\mathfrak{s}=(m,K)$ of Theorem~0;
once they have been performed, $E$ at $\mathfrak{s}$ is fixed by the condition of Theorem~0 at
that site, as Definition~\ref{4.38:def-admissible-sites} states.

For~(e): the last sentence of the statement on the coupling increments, that $K$ is a
well-defined non-negative integer, and the last member of clause~(0) are the conjunction of the
site-freeness of the increments $\beta_j$, given
by Lemma~\ref{4.38:prf-increment-site-free}, with the admissibility of the sites,
\eqref{4.38:eq-step-count-1}, which comes from Definition~\ref{4.38:def-admissible-sites}; in
particular, by~(b), $K=K(g_0,g)$ is determined by $g$ and by the site-free sequence $(g_k)$. This
proves~(e).
Part~(f) is the volume-independence statement for sites, read at the auxiliary tori of item~(A7)
of the statement on the cell-size multipliers $\check{R}_j$.
\end{proof}

\begin{remark}[Termwise hyperoctahedral covariance of the localized terms]
\label{4.38:rem-hyperoctahedral-covariance}
Let $j\ge1$. Assume the following:
\begin{itemize}
\item[(a)] the covariance of the walk expansions under reflections;
\item[(b)] clause (f3) of the hypothesis of Definition~\ref{4.4:def-one-family-hypothesis};
\item[(c)] for every site $\mathfrak{s}\in\mathcal{S}_j$, the covariance \cite[(2.32)]{B-CMP119} of
the localized terms at $\mathfrak{s}$ under every global lattice symmetry of the torus
$\mathbb{T}_j(\mathfrak{s})$ that preserves the cell partitions of every level, the strips and the
coupling profile, the case of reflections being a consequence of (a) and (b);
\item[(d)] the torus-independence of the localized terms,
Theorem~\ref{4.38:thm-torus-independence}.
\end{itemize}
Let $r$ be the lattice motion $x\mapsto\sigma x+t$, where $\sigma\in W_4$ acts about the centre of
a cell of side $M\check{R}_j$ and $t\in M\check{R}_j\mathbb{Z}^4$; we write $r=(\sigma,t)$. Then $r$
descends to every
$\mathbb{T}_j(\mathfrak{s})$, $\mathfrak{s}\in\mathcal{S}_j$, as a global symmetry preserving the
cells of every level, the strips and the coupling profile. Let $E^{(j)}$ and $\mathfrak{r}^{(j)}$ be the two
site-free families of localized terms of Theorem~\ref{4.38:thm-torus-independence}. Let $\hat X$
be a domain and $\Phi$ a field argument of these families such that, at some site
$\mathfrak{s}\in\mathcal{S}_j$, both $\hat X$ and $r\hat X$ are non-wrapping at $\mathfrak{s}$ in the
sense of Theorem~\ref{4.38:thm-torus-independence}. Then
\begin{equation}\label{4.38:eq-hyperoctahedral-covariance-1}
\mathfrak{r}^{(j)}(r\hat X;r\Phi)=\mathfrak{r}^{(j)}(\hat X;\Phi),
\end{equation}
and likewise $E^{(j)}(r\hat X;r\Phi)=E^{(j)}(\hat X;\Phi)$.
\end{remark}

\begin{proof}
First, $r$ descends to every site. The linear part $\sigma$ is a signed permutation of the
coordinate axes. It therefore maps $s_j(\mathfrak{s})\mathbb{Z}^4$ onto itself for every
$\mathfrak{s}\in\mathcal{S}_j$, and the translation part $t$ commutes with this lattice of
translations. Hence $r$ normalises the deck group of $p_{\mathfrak{s}}$. It consequently descends
to $\mathbb{T}_j(\mathfrak{s})$ for every $\mathfrak{s}\in\mathcal{S}_j$, as a global symmetry of
the whole torus, and, since $\sigma$ fixes the centre of a cell of side $M\check{R}_j$ and $t$ is a
translation by a vector of $M\check{R}_j\mathbb{Z}^4$, this symmetry maps each cube of side
$M\check{R}_j$ of the partition onto another, so it preserves the partition into such cubes.
Hence it preserves the cells of every level and the strips. For each level $i$ it is therefore a
motion in $\mathfrak{G}_i$, the group of lattice motions preserving the level-$i$ cell partition of
Lemma~\ref{4.38:lem-coupling-profile}\,(c). By Lemma~\ref{4.38:lem-coupling-profile}\,(c) it then
preserves the coupling profile: for every point $x$ and every label $\ell$,
\begin{equation*}
\varphi_i(rx;r\ell)=\varphi_i(x;\ell).
\end{equation*}

Now fix a site $\mathfrak{s}\in\mathcal{S}_j$ at which $\hat X$ and $r\hat X$ are both
non-wrapping. Since $r$ is a global symmetry of $\mathbb{T}_j(\mathfrak{s})$ preserving the cell
partitions of every level, the strips and the coupling profile, hypothesis (c)
applies at $\mathfrak{s}$. It equates the localized term at $\mathfrak{s}$ on $r\hat X$ with
argument $r\Phi$ to the localized term at $\mathfrak{s}$ on $\hat X$ with argument $\Phi$.
Both domains are non-wrapping at $\mathfrak{s}$. Hypothesis (d),
Theorem~\ref{4.38:thm-torus-independence}, therefore identifies each of these two terms with the
corresponding site-free term. This gives
\eqref{4.38:eq-hyperoctahedral-covariance-1}. The same argument, applied to the family
$E^{(j)}$, gives $E^{(j)}(r\hat X;r\Phi)=E^{(j)}(\hat X;\Phi)$.

Of hypothesis (c) only its statement for a symmetry of the whole torus has been used; no
localized form of its proof enters.
\end{proof}

\begin{definition}[The effective action with its position-dependent coupling, live large-field regions, and the representation of the densities as a sum over labels]\label{4.20:not-representation}
Throughout, the hypotheses (H1)--(H6) of Theorem~0 (Definition~\ref{2.2:not-hypotheses}) are in force.
We use the following objects, fixed in Chapter~1 and in Theorem~\ref{4.1:thm-clause-zero}:
\begin{itemize}
\item the bare density $\rho_0$, fixed in Chapter~1, and Ba{\l}aban's large-field operation
$\mathbf{R}$ (Definition~\ref{1.5:def-densities});
\item the exact push-forward $\mathbf{T}$ under the block-averaging map
(Definition~\ref{1.5:def-block-average});
\item the running couplings $g_k$, $x_k=g_k^{-2}$, $0\le k\le K=K(g_0,g)$
(Definition~\ref{1.2:def-couplings});
\item the induction of clause~(0) of Theorem~0, which bounds these couplings
(Theorem~\ref{4.1:thm-clause-zero}).
\end{itemize}
The effective densities are
\begin{equation}\label{4.20:eq-representation-1}
\rho_k=(\mathbf{R}\mathbf{T})^k\rho_0,\qquad 0\le k\le K,
\end{equation}
where the recursion and the choices that Ba{\l}aban leaves open are those of
Definition~\ref{1.5:def-densities}.
At the step that creates a large-field region, the construction adds to Ba{\l}aban's procedure the
split of the creating step, defined in the section of this chapter devoted to that split; its added
indicators carry two further $g_k$-dependent thresholds, named in Chapter~1. Following Ba{\l}aban
we write $\rho_k$ as a density in the block field $V_k$. That $\rho_k\,dV_k$ is a non-negative
measure of finite total mass is one of the hypotheses of the induction on the effective densities,
listed in a proposition of this chapter, and that
$\rho_k$ need not be a function of $V_k$ is the content of the proposition on pointwise values of the
densities. The effective action \cite[(2.23)]{B-CMP119} is evaluated at
Ba{\l}aban's background configuration $U_k=U_k(V_k)$, the minimal configuration determined by the
block data $V_k$ (Definition~\ref{1.5:def-domains}).

For a weight function $c$ on the torus, $A(c,U)$ is the weighted Wilson action of
Definition~\ref{1.5:def-domains},
\begin{equation*}
A(c,U)=\sum_p c(x(p))\,[1-\operatorname{Re}\operatorname{tr}U(\partial p)],
\end{equation*}
where $x(p)$ is the point attached to the plaquette $p$. Near large fields the coupling depends on
position. Let $\Lambda_j$ be the domain of step $j$ on which only small-field characteristic
functions occur; it is a union of cubes of side $M\check R_j$ in the lattice units of step $j$, where
$\check R_j$ is the cube-size factor fixed in the section of this chapter on the split at a creating
step. Let $\Lambda_j^{\sim-1}$ denote $\Lambda_j$ with one layer of these cubes
removed. Choose $\varphi_j\in C_0^\infty(\Lambda_j)$ with $\varphi_j=1$ on $\Lambda_j^{\sim-1}$.
The coupling profiles are those of Definition~\ref{1.5:def-domains}, as in
\cite[(2.24)]{B-CMP119}: $x_0(\cdot)\equiv g_0^{-2}$ and
\begin{equation}\label{4.20:eq-representation-2}
x_j(\cdot)=x_{j-1}(\cdot)-\beta_j(g_{j-1})\,\varphi_j(\cdot),\qquad 1\le j\le k,
\end{equation}
where $\beta_j=\beta_j(g_{j-1})$ is the coupling increment of step $j$
(Definition~\ref{4.2:def-increment}), so that $x_j=x_{j-1}-\beta_j$ as in
Definition~\ref{1.2:def-couplings}. The profiles satisfy
\begin{equation}\label{4.20:eq-representation-3}
x_k(\cdot)-x_k=\sum_{1\le j\le k}\beta_j(g_{j-1})\bigl(1-\varphi_j(\cdot)\bigr),
\end{equation}
so that $x_k(\cdot)=x_k$ on $\bigcap_{1\le j\le k}\Lambda_j^{\sim-1}$; $g_k^{-2}(\cdot)=x_k(\cdot)$ is
the position-dependent inverse coupling of Definition~\ref{1.5:def-domains}. No nesting of the
$\Lambda_j$ is assumed. When
$\Lambda_k\subset\Lambda_j$ for all $j\le k$, the intersection is $\Lambda_k^{\sim-1}$.

Following \cite[(2.23)]{B-CMP119}, the effective action, as a function of the background
configuration, is
\begin{equation}\label{4.20:eq-representation-4}
\begin{split}
A_k\bigl(g_k^{-2},U_k\bigr)={}&-A\bigl(g_k^{-2}(\cdot),U_k\bigr)+\mathbf{E}_k(U_k)
+\mathbf{R}_k(U_k)\\
&+\mathbf{B}_k(U_k,\underline{\eta})-E_k ;
\end{split}
\end{equation}
the plaquette action carries the minus sign. The boundary part $\mathbf{B}_k(U_k,\underline{\eta})$
depends also on the collection $\underline{\eta}=(\eta_j)_{j<k}$ of fluctuation fields of earlier
steps that are stored on the large-field region, $\eta_j$ denoting the fluctuation field of step $j$;
they are integrated by the operations composing
$\mathbf{T}'_k(\mathcal{S})$ in \eqref{4.20:eq-representation-6}, as described later in this
chapter. Ba{\l}aban writes $A$ for this collection; here it is written $\underline{\eta}$, and its
members $\eta_j$, so that the letter $A$ denotes only actions. Where this argument plays no
role it is suppressed, as Ba{\l}aban does. Here
\begin{equation}\label{4.20:eq-representation-5}
\begin{gathered}
\mathbf{E}_k(U)=\sum_{j\le k}\sum_X\bigl[\mathbf{E}^{(j)}(X,U)-\mathbf{E}^{(j)}(X,1)\bigr],\\
\mathbf{R}_k(U)=\sum_{j\le k}\sum_X\bigl[\mathbf{R}^{(j)}(X,U)-\mathbf{R}^{(j)}(X,1)\bigr],\qquad
\mathbf{B}_k=\sum_{j\le k}\sum_X \mathbf{B}^{(j)}(X),
\end{gathered}
\end{equation}
where $\mathbf{E}^{(j)}(X,U)$, $\mathbf{R}^{(j)}(X,U)$ and $\mathbf{B}^{(j)}(X)$ denote the regular,
$\mathbf{R}$-made and boundary terms produced at step $j$ and localized in the domain $X$, and the
argument $1$ is the unit configuration; $E_k$ is a number, into which the subtracted
unit-configuration values enter. The sums defining
$\mathbf{E}_k$ and $\mathbf{R}_k$ run over $X\in\mathcal{D}_j$ ($\mathcal{D}_j$ as in
Definition~\ref{1.5:def-domains}) with $X\subset\Lambda_j$, the same
domain over which the unit-configuration values enter the vacuum energy defined later in this
chapter; the remaining terms, $\mathbf{R}$-terms included, are not differenced, are collected in
$\mathbf{B}_k$ with the boundary terms, and contribute no unit-configuration value to $E_k$.

For $X\in\mathcal{D}_j$, $d_j(X)$ denotes Ba{\l}aban's linear size \cite{B-CMP109}
(Definition~\ref{1.5:def-domains}): the length of a shortest graph of his class for $X$, measured in
units of the scale-$j$ lattice spacing and divided by $M$. The tree length $d^\theta(X)$ is the one
of Definition~\ref{1.5:def-densities}.

A large-field component of a term is called \emph{live at level $k$} if $\mathbf{R}$ has acted on
it and the term still carries the level-$k$ block field on it as a variable: a second-class
region of \cite{B-CMP122b} carries it through its gauge-fixing factor, written
$\delta_{G_i}(V'_k)$ there; a first-class region carries it through the reference law of
\cite[(1.100)--(1.102)]{B-CMP122a}. The consequence of a live
gauge-fixing factor is derived in the proposition on pointwise values of the densities. For
$X\in\mathcal{D}_j$ and a union $Z$ of such regions, $d_{j,Z}(X)$ is the relative linear size of
\cite[(1.67)]{B-CMP122b} (Definition~\ref{1.5:def-domains}).

Following \cite[(2.18)--(2.22)]{B-CMP119}, $\rho_k\,dV_k$ is a sum over the admissible sequences of
domains $\mathcal{S}$. Each such sequence consists of the large-field and small-field domains created
at steps $\le k$; we call it a \emph{label}. Schematically, as an identity of measures,
\begin{equation}\label{4.20:eq-representation-6}
\rho_k\,dV_k=\sum_{\mathcal{S}}\mathbf{T}'_k(\mathcal{S})
\Bigl[\chi_{k,\mathcal{S}}\,\exp A_{k,\mathcal{S}}\bigl(g_k^{-2},U_k(V_k)\bigr)\Bigr]\,dV_k,
\end{equation}
the $\delta$-factors of live second-class components being included in
$\mathbf{T}'_k(\mathcal{S})$. Here $\chi_{k,\mathcal{S}}$ is a product of characteristic functions
and $A_{k,\mathcal{S}}$ has the form \eqref{4.20:eq-representation-4}. The operator
$\mathbf{T}'_k(\mathcal{S})$ is the composition, over the components of the large-field region of
the label, of their large-field integral operations: on a component on which $\mathbf{R}$ has not
acted, the operation $\mathbf{T}_k(Z_k)$ of \cite[(2.18)]{B-CMP119} with the preparation factors,
written $\mathbf{T}''_k(Z_k)$ in \cite[(1.72)]{B-CMP122b}; on a component $X$ on which $\mathbf{R}$
has acted, the top-integrated operation $\mathbf{T}'_k(X)$ of \cite[(1.71)]{B-CMP122b}, which
contains the preparation tests of \cite[(1.82), (1.88)]{B-CMP122a} and whose bounds, the majorants
\cite[(1.73), (1.79)]{B-CMP122b} and the bound \cite[(1.89)]{B-CMP122b}, are part of what
Theorem~\ref{4.20:thm-representation}\,(a) asserts.
\end{definition}

\begin{proof}
We prove \eqref{4.20:eq-representation-3} and the two assertions that follow it. By definition
$x_0(\cdot)\equiv g_0^{-2}=x_0$. For $1\le j\le k$, subtracting the relation
$x_j=x_{j-1}-\beta_j(g_{j-1})$ of Definition~\ref{1.2:def-couplings} from
\eqref{4.20:eq-representation-2} gives
\begin{equation*}
x_j(\cdot)-x_j=x_{j-1}(\cdot)-x_{j-1}+\beta_j(g_{j-1})\bigl(1-\varphi_j(\cdot)\bigr).
\end{equation*}
Summing over $j=1,\dots,k$, the left-hand terms telescope, and since $x_0(\cdot)-x_0=0$ we obtain
\eqref{4.20:eq-representation-3}. On $\bigcap_{1\le j\le k}\Lambda_j^{\sim-1}$ every $\varphi_j$
equals $1$, because $\varphi_j=1$ on $\Lambda_j^{\sim-1}$; hence every summand on the right of
\eqref{4.20:eq-representation-3} vanishes there, and $x_k(\cdot)=x_k$.

Now suppose $\Lambda_k\subset\Lambda_j$ for all $j\le k$. By the definition of the cube-size
factors in the section of this chapter on the split at a creating step,
$\check R_j\le L^{k-j}\check R_k$. Measured in the lattice units of step $k$,
the cubes of $\Lambda_j$, $j<k$, have side $L^{-(k-j)}M\check R_j$, and
$L^{-(k-j)}M\check R_j\le M\check R_k$ because $\check R_j\le L^{k-j}\check R_k$; thus the cubes of
$\Lambda_j$ are no larger than the cubes of side $M\check R_k$ of $\Lambda_k$, and therefore
$\Lambda_k\subset\Lambda_j$ implies $\Lambda_k^{\sim-1}\subset\Lambda_j^{\sim-1}$. Consequently
$\Lambda_k^{\sim-1}$ is contained in every set of the intersection, while the intersection is
contained in its member $\Lambda_k^{\sim-1}$ (the index $j=k$). The intersection is therefore
$\Lambda_k^{\sim-1}$.
\end{proof}

Ba{\l}aban's ultraviolet stability theorem \cite[Theorem~1]{B-CMP122b}
has one hypothesis: the effective couplings
$g_0,\dots,g_K$ are contained in $(0,\gamma]$. In the following theorem this hypothesis is
supplied by clause~(0) of Theorem~0 (Theorem~\ref{4.1:thm-clause-zero}).

\begin{theorem}[The inductive representation of the effective densities (clause (i-a)), with
the bounds for boundary terms beside live regions and the constants clause]
\label{4.20:thm-representation}
For every $0\le k\le K$ the following hold.
\begin{itemize}
\item[(a)] The density $\rho_k$ has the representation of
Definition~\ref{4.20:not-representation}, which is the representation
\cite[(2.18)]{B-CMP119}. It is a sum over the labels $\mathcal{S}$ of
Definition~\ref{4.20:not-representation}, which index the admissible histories of
large-field operations. Each term is a large-field operation applied to the product of
characteristic functions with $\exp A_k(g_k^{-2},U_k)$. Here the effective action
$A_k(g_k^{-2},U_k)$ of Definition~\ref{4.20:not-representation} is
$-A(g_k^{-2}(\cdot),U_k)$ plus the localised terms of every level $j\le k$.

The density $\rho_k$ satisfies the inductive hypotheses and bounds of \cite[\S2]{B-CMP119}
with the exception of \cite[(2.47)]{B-CMP119}, which is not part of this theorem; in place of
the pointwise upper member of \cite[(2.50)]{B-CMP119} it satisfies the integrated bound
$\int\rho_k\,dV_k\le e^{E_+N_k}$ of Theorem~\ref{4.33:thm-stability}.
The display
\cite[(2.48)]{B-CMP119} and the bounds of \cite[Theorem~2]{B-CMP119} are carried in the
forms stated later in this chapter. Where another display of \cite[\S2]{B-CMP119} is used in
a form differing from the printed one, the form used is stated at the point of use.

In particular, every localised term created at step $j\le k$ has the following properties:
\begin{itemize}
\item it is localised in some $X\in\mathcal{D}_j$ of Definition~\ref{1.5:def-domains}, it is
gauge invariant, and it depends on
the background configuration $U_k$ only through the restriction of $U_k$ to $X$;
\item it is analytic on $U^c_j(X,\alpha_0,\alpha_1)$, in the reading of this class fixed at
the end of part (a), and, except for the boundary terms described in (b), it is bounded there
by $\mathrm{const}\cdot e^{-\kappa d_j(X)}$.
\end{itemize}
New terms obey these bounds on the enlarged analyticity spaces, with rate
$(1+4\beta)\kappa$, where $\beta$, without index and distinct from the coupling increments
$\beta_{k+1}$, is the schedule constant of the layered analyticity spaces
(Definition~\ref{4.4:def-post-step-space}).

The class $U^c_j(X,\alpha_0,\alpha_1)$ is understood as the one-power hereditary class. For
this class the passage to sub-regions built from finer cubes is an identity.
It contains Ba{\l}aban's domains of \cite[(1.11)--(1.16)]{B-CMP109} read hereditarily, that
is, with the regularity condition of \cite[(1.11)--(1.16)]{B-CMP109} on constrained
minimisers imposed, with collars of two layers of each member's own cubes, at every member
of the tower of subdomains of $X$ (the nested subdomains from whose block averages these
minimisers are built \cite{B-CMP99b}), under each of Ba{\l}aban's two rules: the maximal
tower of \cite{B-CMP109} and the sequence of domains of \cite[(2.13)]{B-CMP119}.
Analyticity on $U^c_j(X,\alpha_0,\alpha_1)$ therefore yields analyticity on these domains
read hereditarily. That this class
is reproduced by the step is the statement of the section on the one-power hereditary class.

\item[(b)] The exception to the bound in (a) consists of the boundary terms whose localisation
domain contains a live second-class region, in the sense of \cite{B-CMP122b}, live being
understood as in Definition~\ref{4.20:not-representation}. Let $Z=\bigcup_iY_i$ be the union
of these regions.

Consider first a new boundary term created by $\mathbf{R}$ at step $j$. Suppose that its
localisation domain contains no second-class region that was created at an earlier step and
is still live. For such a term \cite[(1.99)]{B-CMP122b} states the bound below. In the
notation of that display, $\mathbf{B}'^{(j)}(X,\cdot)$ denotes the new boundary term
localised in $X$ (it is not the large-field operation $\mathbf{B}$), $(\mathbf{U},\mathbf{J})$
is the argument of the boundary term as written in that display, and $d_{j,\cup_iY_i}(X)$ is
the relative length $d_{j,Z}(X)$ of Definition~\ref{1.5:def-domains}. The bound reads
\begin{equation}\label{4.20:eq-representation-1-bis}
\bigl|\mathbf{B}'^{(j)}(X,(\mathbf{U},\mathbf{J}))\bigr|
\le O(1)\,c_1\,e^{-(1+\frac12\beta)\kappa\, d_{j,\cup_iY_i}(X)} .
\end{equation}
Here $c_1\in\{e^{-p_0(g_j)},\alpha^{1/3}\}$, where $p_0(\cdot)$ is the degree function. The
$Y_i$ are the second-class regions, $\alpha$ is the parameter so written in
\cite{B-CMP122b} (not the radii $\alpha_0,\alpha_1$), and $\beta$ in this display is the
parameter so written in \cite{B-CMP122b}.

\cite{B-CMP122b} states the bound \eqref{4.20:eq-representation-1-bis} for the terms just
described. For the remaining boundary terms of this kind, which are of two sorts,
\begin{itemize}
\item the terms built on the former by later $\mathbf{T}$- and $\mathbf{R}$-steps while the
region is live;
\item the new terms whose domain contains an older live region,
\end{itemize}
the following bound holds: each of them is bounded by a constant with the dependence stated
in (c), times $e^{-c\,d^\theta(X)}$ with $c>0$.
Here $d^\theta(X)$ is the tree length of Definition~\ref{1.5:def-densities}, and
$d^\theta(X)\le d_{j,Z}(X)$. Part (b) asserts decay in $d_{j,Z}(X)$ itself only for the
instance \eqref{4.20:eq-representation-1-bis}, and it asserts no lower bound of $d^\theta(X)$ in
terms of $d_{j,Z}(X)$.

\item[(c)] Every constant in (a) and (b) depends only on $L$, Ba{\l}aban's auxiliary
parameters $\mathfrak{p}$, $\gamma$ and $N$, and it depends on the scale only through $g_k$.
There is one exception: the ratio
\begin{equation*}
\varpi_{jn}:=\max\{1,(x_j/x_n)^2\}=\max\{1,g_n^4/g_j^4\},\qquad j\le n ,
\end{equation*}
which depends on two scales, through the running couplings $g_j$ and $g_n$. It is the
exception named in the constants clause of Theorem~0. It appears in the bound corresponding
to \cite[(2.44)]{B-CMP119}, among the bounds of \cite[Theorem~2]{B-CMP119} in the form
referred to in (a), and in the bounds derived from it.

The ratio is controlled as follows. By the flow bounds of clause~(0)
(Theorem~\ref{4.1:thm-clause-zero}), $x_n\le x_j+c_5$ for $j\le n$. The threshold
$\gamma_H$ of hypothesis (H3) (Definition~\ref{2.2:not-hypotheses}) is chosen with
$\gamma_H\le(\gamma^{-2}+c_5+1)^{-1/2}$, and $g_j<g$ by clause~(0), so that
\begin{equation*}
x_j=g_j^{-2}>g^{-2}\ge\gamma_H^{-2}\ge\gamma^{-2}+c_5+1>c_5 .
\end{equation*}
This gives $x_n\le 2x_j$, whence
\begin{equation*}
g_j^{\kappa_0}\varpi_{jn}\le 4g_j^{\kappa_0-4}g_n^4 ,
\end{equation*}
where $\kappa_0$ is one of Ba{\l}aban's exponents.

No constant, $\varpi_{jn}$ included, depends on $K$, on $m$ or on position, or on $g_0$ other
than through the running couplings.
\end{itemize}
\end{theorem}

The proof is the induction on the level $k$ carried out in the sections that follow, which is
one induction with that of clause~(0) (Theorem~\ref{4.1:thm-clause-zero}); it is concluded in
the section closing the induction, immediately before Theorem~\ref{4.33:thm-stability}.

\begin{proposition}[The relative small-field machine: after each step every stored boundary-type term carries a label domain and a finite set of live pockets, with the format conditions on support, bound, scaffold-reading and anchoring]\label{4.21:prop-relative-machine}
Work in the setting of the representation over labels of
Definition~\ref{4.20:not-representation}, with the labels $\mathcal{S}$ of
\cite[(2.18)]{B-CMP119}.

A \emph{pocket} is a pair $(Y,m)$ of a large-field region $Y$ and a level $m$ attached to $Y$
by the construction. The
\emph{enlarged pocket} $Y^\sharp$ is the region $Y$ with layers of $M$-cubes added, the number
of layers being fixed by the level at which the large-field operation acted on $Y$. Let
$\Omega_j$ be the small-field region of level $j$, and let $Z_j^{\boxplus}$ be the enlargement
of the large-field region of level $j$ that enters Ba{\l}aban's conditions
\cite[(2.41)]{B-CMP119}. Let $(Z_j^{\boxplus})^{(1)}$ be the first enlargement of
$Z_j^{\boxplus}$. Let $\varepsilon_k$ be the current radius of the analyticity domains at
level $k$.

\emph{Terms concerned.} A \emph{boundary-type term} of level $j$ is one of the following:
\begin{itemize}
\item[($\alpha$)] a term of the boundary part $\mathbf{B}_k$ of the effective action at complex
source $w$: one born in the operation $\mathbf{T}$, written
$\mathbf{B}^{(j)}(\cdot,\mathcal{S})$, or one born in the large-field operation $\mathbf{R}$,
written $\mathbf{B}'^{(j)}$;
\item[($\beta$)] a frozen chain term $\mathbf{C}^{(i)}_{k_1}$ of level $i$, so that $j=i$, of a
piece that is still live, frozen at the earlier level $k_1$ at which that piece was arrested;
\item[($\gamma$)] inside a large-field step $\mathbf{R}_k$, a chain term
$\mathbf{C}^{(n)}_k(\cdot,\mathcal{S})$ of level $n\le k$, so that $j=n$;
\item[($\delta$)] a change-of-variables term of an interrupted preparation.
\end{itemize}
The terms kept under the
operation $\mathbf{T}'$ whose domain is an enlarged pocket $Y^\sharp$ itself are not
boundary-type terms.

\emph{The format at level $k$.} The relative format at level $k$, which we write
$(\mathrm{RF})_k$, asserts the following. In every
label, after $\mathbf{R}_k$, every boundary-type term $F$ of level $j\le k$ carries scaffold
data $(A_F,\mathcal{P}_F)$ with these properties.
\begin{itemize}
\item[(f0)] The domain $A_F\in\mathbf{D}_j$, the \emph{label domain} of $F$, obeys the side
conditions that Ba{\l}aban imposes on the domain $X$, namely those of
\cite[(2.41)]{B-CMP119}: $A_F\cap\Omega_j\neq\emptyset$ and
$A_F\cap Z_j^{\boxplus}\neq\emptyset$, where the second condition is waived when
$\mathcal{P}_F\neq\emptyset$. For chain terms they are the conditions of
\cite[(1.63)]{B-CMP122b}. The set $\mathcal{P}_F$ is a finite set of pockets $(Y,m)$ with
$m\le j$. Each of them is live at level $k$, and each has $Y^\sharp$ touching $A_F$. For
each kind, each domain $A$ and each label $\mathcal{S}$ there is at most one such term.
\item[(f1)] Set
\begin{equation*}
X_F:=A_F\cup\bigcup_{(Y,m)\in\mathcal{P}_F}Y^\sharp .
\end{equation*}
The term $F$ depends on the configuration and current $(\mathbf{U},\mathbf{J})$ and on the
further source variable $\mathbf{A}$ only through their restrictions to $X_F$, and on the
complex source $w$
only through its restriction to a neighbourhood $X_F^+$ of $X_F$. It has an orbit datum of
radius $\varepsilon_k$ on the analyticity domain stored for $F$. That domain lies inside the
analyticity space over $X_F$ of the terms of its kind.
\item[(f2)] The term $F$ obeys the bound that Ba{\l}aban's construction gives for terms of its
kind, with $d_j(X)$ replaced by $d_j(A_F)$. The norm is the one in which that bound is
stated, taken at radius $\varepsilon_k$; we write it $\mathfrak{N}_{\varepsilon_k}$. For
boundary terms,
\begin{equation*}
\mathfrak{N}_{\varepsilon_k}\bigl(\mathbf{B}^{(j)}(A,\cdot)\bigr)\le
B_0e^{-\kappa_{\mathrm{out}}d_j(A)}\quad(\kappa_{\mathrm{out}}\ge 2\kappa)
\end{equation*}
holds for fresh terms, that is, for terms created at the step just performed, with the
constant $B_0$ of the bounds \cite[(2.3), (2.42)]{B-CMP119}.
The same bound with $\kappa$ in place of $\kappa_{\mathrm{out}}$ \cite[(2.3), (2.42)]{B-CMP119}
holds for stored terms, that is, for terms created at an earlier step and carried since
then. Terms born in $\mathbf{R}$ obey
\begin{equation*}
\mathfrak{N}_{\varepsilon_k}\bigl(\mathbf{B}'^{(j)}(A,\cdot)\bigr)
\le C\,c_1(g_j)\,
e^{-(1+\frac14\beta_{\mathrm{s}})\kappa d_j(A)}
\le B_0e^{-\kappa d_j(A)} ,
\end{equation*}
where $C$ is a constant and $c_1(g_j)\in\{e^{-p_0(g_j)},\alpha^{1/3}\}$ is Ba{\l}aban's
factor. The chain terms $\mathbf{C}^{(n)}_k$ have, in the norm of chain terms that enters
\cite[(1.70)]{B-CMP122b}, norm at most a constant $C_0^\sharp$. In that norm
the exponential weight is $a_k=\frac85\kappa$ at the top level and $a_i=\frac45\kappa$ at
levels $i<k$, and the sum and the exponential weight of that norm are taken over the label
domains $A$ in place of the domains $X$. Terms of kind
($\delta$) obey the bound stated for them at their construction.
\item[(f3)] The pair $(A_F,\mathcal{P}_F)$ depends on the scaffold only and not on the
complex source. In particular it is the same for two runs of the construction on one
scaffold and for their differences.
\item[(f4)] The domain $A_F$ has a cube in $(Z_j^{\boxplus})^{(1)}$.
\end{itemize}
For $\mathcal{P}_F=\emptyset$ these conditions are, for boundary terms,
\cite[(2.41)--(2.42)]{B-CMP119}, and, for chain terms, \cite[(1.63), (1.70)]{B-CMP122b},
each with an orbit datum added. Nothing is asserted
about $d_j(X_F)$.

\emph{Assertion.} The statement $(\mathrm{RF})_0$ is empty. Let the parameters obey one-sided
floors on $\kappa$, $\kappa_1$ and $M$ of the same kind as Ba{\l}aban's. Let
$\gamma$ obey one-sided ceilings, that is, upper bounds on $\gamma$. Assume the decoupling
statements for the sources, the
statement on tubes about the live regions, the statement giving the change-of-variables terms
of an interrupted preparation their label domains and bounds, the persistence of the weights
of the near groups of the step, both for one run and for the differences of two runs on one
scaffold, and the statement supplying the orbit data of (f1) on the stored analyticity
domains (the last two are not proved; see the route below). Then, within the single induction
of clauses (0) and (i) of
Theorem~0, the inductive statements at level $k$ of the induction proving
Theorem~\ref{4.20:thm-representation}, together with $(\mathrm{RF})_k$, imply
$(\mathrm{RF})_{k+1}$, for every $k<K$.

Moreover, every bound of the induction that involves $d_j(X)$, or another geometric quantity
of the domain $X$ of a boundary-type term, holds with $X$ replaced by $A_F$ in the norm of
(f2), taken at the output radius of each operation: in the operation $\mathbf{T}$ and in the
large-field operation at level $k+1$ the radius is reduced by a margin, and in the preparation
it is the radius of the preparation's own bound. It holds with the same labels, index
sets, constants, floors
and ceilings. The second group of \cite[(1.99)]{B-CMP122b} is unaffected.
\end{proposition}

Proof outstanding.

\emph{Route.} The proof is an induction on the level $k$. It is carried inside the single
induction of clauses (0) and (i) of Theorem~0 by the inductive representation of
the effective densities (Theorem~\ref{4.20:thm-representation}). At level zero the statement
is empty.

The step from $k$ to $k+1$ consists of four operations, taken in this order:
\begin{enumerate}
\item the operation $\mathbf{T}$ (Definition~\ref{1.5:def-densities});
\item at level $k+1$, the preparation;
\item the large-field operation;
\item Ba{\l}aban's localisation, exponentiation and bounds \cite[(1.90)--(1.99)]{B-CMP122b}.
\end{enumerate}
One must show that each of these operations preserves (f0)--(f4), and that in every bound the
domain $X$ of a boundary-type term may be replaced by its label domain $A_F$.

Three ingredients are to be used. The first is geometric: the relative length $d_{j,Z}(X)$ of
\cite[(1.67)]{B-CMP122b} is the plain length $d_j$ of a skeleton of the domain. The second is
a counting of the terms, each in
its own cubes, in which no field variable enters. The third is a decoupling
of the sources.

The proof depends on two inputs that are themselves outstanding. The first is the persistence
of the weights of the near groups of the step, both for one run and for the differences of two
runs on one scaffold. The second is the orbit datum of (f1) on the stored analyticity domains.

\begin{proposition}[Volume sums on labels and the life of wrapped terms]\label{4.21:prop-label-volume-sums}
Let the boundary-type terms be stored in the relative format of Definition~\ref{4.21:def-relative-format}, with the conditions \eqref{4.21:eq-relative-format-a}. Assume the following counting hypothesis. The number of boundary-type terms charged at one cube, read on labels, is bounded by the count that the correlation theorem uses:
\begin{itemize}
\item for the terms $\mathbf{B}$, $\mathbf{B}'$, the count \cite[(1.26)]{B-CMP116};
\item for the chain terms $\mathbb{C}$, the count the correlation theorem makes through the norm of chain terms and its two counting statements for chain terms.
\end{itemize}
Each of these counts reads a term through its domain in $\mathbf{D}_j$, its bound (condition (f2) of \eqref{4.21:eq-relative-format-a}) and the side conditions (condition (f0) of \eqref{4.21:eq-relative-format-a}). Then the following hold.
\begin{enumerate}
\item[(a)] \emph{Sup-norm sums.} Consider the following sums, each a sum of sup-norms of boundary-type terms charged at a cell of their domains:
\begin{itemize}
\item the sums of \cite[(2.47)--(2.48)]{B-CMP119} as they enter the correlation theorem;
\item the volume bound of the correlation theorem;
\item the per-cell sum $\mathsf{K}_{\mathrm{cell}}$ of the correlation theorem;
\item the $\Gamma$-term of the torus bound.
\end{itemize}
Every such sum holds with the charge taken on the label. By \eqref{4.21:eq-relative-format-a}, the label $A_F$ of a term $F$ of level $j$ contains a cube of level $j$ lying in $(Z_j^{\boxplus})^{(1)}$, and $F$ is charged once, at that cube. The one geometric change is
\begin{equation*}
N_j\bigl((Z_j^{\boxplus})^{(1)}\bigr)\le 3^d N_j(Z_j^{\boxplus}).
\end{equation*}
This costs a factor $3^d$ on the term of $\mathsf{K}_{\mathrm{cell}}$ coming from the summation of boundary terms at finer zones.

At the large-field site, a wrapped term $v$ contributes the following, in place of the total that the correlation theorem uses for one term:
\begin{equation}\label{4.21:eq-label-volume-sums-1}
\sum_{\mathfrak{p}\text{ of }v}|\mathfrak{p}|+|\mathrm{val}_v|\le 4\,\mathfrak{N}^A(v)\,e^{4(d_k(K_v)+1)} .
\end{equation}
Here the factor $e^{4d_k(K_v)}\le e^{4d_j(A_v)}$ is absorbed by four units of the decay rate of $v$. These four units are available in every case by the per-cube bound on stored wrapped terms (Lemma~\ref{4.21:lem-wrapped-cube-bound}). The factor $4e^4$ multiplies the contribution of the summation of boundary terms at finer zones to $\mathsf{K}_{\mathrm{cell}}$. Consequently the volume bound of the correlation theorem holds for $\mathfrak{v}+\mathfrak{v}^{\mathrm{rel}}$.

Interior sums need no depth factor, and the lemma supplying the third, depth-dependent factor is not needed for them. In particular, for every pocket $(Y,m)$ with strip $Y^{\sharp}$,
\begin{equation}\label{4.21:eq-label-volume-sums-2}
\sum_{A:\ Y^{\sharp}\text{ touches }A}\mathfrak{N}^A\bigl(\mathbf{B}'^{(k)}(A)\bigr)
\le e^{\lambda_\theta}C_{99}\,c_1^{\flat}\,K(d)\cdot N_k\bigl((Y^{\sharp})^{(1)}\bigr).
\end{equation}
The surface of a pocket enters only the sums over all terms around one fixed pocket, as in \cite{B-CMP122b}. There it lies inside the volume budget $O(1)\log g_j^{-2}|Z_j|$ of the correlation theorem.

\item[(b)] \emph{A term never outlives its pockets.} Let $(Y,m)\in\mathcal{P}_F$ cease to be live at the operation $\mathbf{R}_k$. Then $F$ is a term of the first part there, in the sense of part (a) of Proposition~\ref{4.21:prop-large-field-site}, and it is resummed into $V^{\sharp}$. It therefore ceases to exist as a stored term. Hence, in the relative format after $\mathbf{R}_k$, all pockets of all stored terms are live.

\item[(c)] \emph{Copies.} The terms that are added and subtracted in \cite{B-CMP122b} are terms of ``the effective action determined by the assumption that $Z_k$ is the only large field region'' \cite{B-CMP122a}. They are therefore terms of another label. The relative format holds for every label. The pockets of these terms lie in the rest of the large-field region, not yet treated at this stage of $\mathbf{R}_k$, so these terms are terms of Proposition~\ref{4.21:prop-large-field-site} like the others.

\item[(d)] \emph{Interruption.} Suppose the preparation of a component is interrupted or excluded, as in \cite{B-CMP122a}. Then the following hold.
\begin{itemize}
\item Its pockets stay live.
\item Its chain terms stay stored, on labels, in the relative format of Definition~\ref{4.21:def-relative-format}.
\item The geometric lemma of the scaffold, in its clauses (c)--(e), (h) and (i), holds for these terms at later sites.
\item Its windows are met at later preparation sites by the first two classes of the grain partition, the cubes (${\flat}$a) and (${\flat}$b) of Definition~\ref{4.21:def-grain-partition}.
\end{itemize}
\end{enumerate}
\end{proposition}

\begin{proof}
(a) Each of the sums listed in (a) is a sum, over domains through a charge point, of a pointwise bound or a norm bound. By \eqref{4.21:eq-relative-format-a}, every label $A_F$ of level $j$ contains a cube of level $j$ in $(Z_j^{\boxplus})^{(1)}$, and $F$ is charged there. The number of terms charged at one such cube is bounded by the counting hypothesis. The number of charge cubes is $N_j((Z_j^{\boxplus})^{(1)})\le 3^dN_j(Z_j^{\boxplus})$.

Next consider a wrapped term $v$ at the large-field site. Its pieces $\mathfrak{p}$ are indexed by four families $\mathsf{y}_1,\dots,\mathsf{y}_4$ of cut cells, one for each link. The value $\mathrm{val}_v$ is the piece in which all four families are empty. The bound used for a piece is the first bound of part (c) of Proposition~\ref{4.21:prop-large-field-site}, namely
\begin{equation*}
|\mathfrak{p}|\le\mathfrak{N}^A(v)\,e^{-(\kappa_1-1)\sum_{t}|\mathsf{y}_t|}.
\end{equation*}
Put $\lambda=\kappa_1-1$. Let $\operatorname{lab}_t$ denote the label after the $t$-th link, with $\operatorname{lab}_0=K_v=[A_v]_k$, the label read by the first link (part (b) of Proposition~\ref{4.21:prop-large-field-site}). The four free sums are bounded by parts (b) and (a) of the summation lemma, Lemma~\ref{4.21:lem-source-anchored-sums}, applied in turn.

We begin with $\mathsf{y}_4$. Part (b) of Lemma~\ref{4.21:lem-source-anchored-sums}, applied to the reader of the fourth link, whose label is $\operatorname{lab}_3$, gives the first inequality below. Counting the empty family as the term $1$, part (a) of the same lemma, which gives $d_k(\operatorname{lab}_3)\le d_k(\operatorname{lab}_2)+c_l|\mathsf{y}_3|$, gives the second:
\begin{equation}\label{4.21:eq-label-volume-sums-3}
1+\sum_{\mathsf{y}_4}e^{-\lambda|\mathsf{y}_4|}\le e^{d_k(\operatorname{lab}_3)+1}
\le e^{d_k(\operatorname{lab}_2)+1}\,e^{c_l|\mathsf{y}_3|}.
\end{equation}

The factor $e^{c_l|\mathsf{y}_3|}$ is then combined with $e^{-\lambda|\mathsf{y}_3|}$. The sum over $\mathsf{y}_3$ is therefore taken at the rate $\lambda-c_l$, and part (b) bounds it, together with the empty family, by $e^{d_k(\operatorname{lab}_2)+1}$. Applying part (a) to $\operatorname{lab}_2$ then leaves
\begin{equation*}
e^{2(d_k(\operatorname{lab}_2)+1)}\le e^{2(d_k(\operatorname{lab}_1)+1)}e^{2c_l|\mathsf{y}_2|}.
\end{equation*}
The sum over $\mathsf{y}_2$ is taken at the rate $\lambda-2c_l$. It gives $e^{3(d_k(\operatorname{lab}_1)+1)}\le e^{3(d_k(K_v)+1)}e^{3c_l|\mathsf{y}_1|}$. The sum over $\mathsf{y}_1$ is taken at the rate $\lambda-3c_l$ and gives the total
\begin{equation*}
e^{4(d_k(K_v)+1)}.
\end{equation*}

Part (b) of Lemma~\ref{4.21:lem-source-anchored-sums} applies at each of the rates used. The smallest of them is
\begin{equation*}
\lambda-3c_l=\kappa_1-1-3c_l\ge\lambda_*,
\end{equation*}
and this inequality holds by \eqref{4.21:eq-format-propagation-a}. For a term of the second part, the pieces come from four brackets. Each bracket is bounded in the same way, which gives the factor $4$. This proves \eqref{4.21:eq-label-volume-sums-1}.

The factor $e^{4d_k(K_v)}\le e^{4d_j(A_v)}$ is absorbed by four units of the decay rate of $v$, which Lemma~\ref{4.21:lem-wrapped-cube-bound} leaves available. The remaining factor $4e^4$ multiplies the contribution of the summation of boundary terms at finer zones to $\mathsf{K}_{\mathrm{cell}}$. Thus the volume bound of the correlation theorem holds for $\mathfrak{v}+\mathfrak{v}^{\mathrm{rel}}$.

For the displayed sum \eqref{4.21:eq-label-volume-sums-2}, part (a) of the skeleton lemma shows that every $A$ touched by $Y^{\sharp}$ contains a cube of $(Y^{\sharp})^{(1)}$. We charge the term $\mathbf{B}'^{(k)}(A)$ at such a cube. The sum over the domains $A$ through one cube is bounded by \cite[(1.26)]{B-CMP116}. The number of charge cubes is $N_k((Y^{\sharp})^{(1)})$.

Charging on labels costs no factor depending on age. A cube of level $j$ in the zone of level $n\le j$ contains $L^{d(j-n)}$ cells of that zone. Hence charging the terms of level $j$ at cubes of level $j$ gives each cell of zone $n$ the weight $L^{-d(j-n)}$. These weights are summable over $j$ with no factor depending on age. This is the summation of \cite[(2.47)]{B-CMP119}.

Parts (b)--(d) hold for the reasons given in their statements. For (b): the term is of the first part in the sense of Proposition~\ref{4.21:prop-large-field-site} and is resummed into $V^{\sharp}$. For (c): the copies are terms of another label, and the relative format holds for every label. For (d): the terms of an interrupted or excluded component stay stored on labels in the relative format. The clauses of the geometric lemma, and the classes (${\flat}$a) and (${\flat}$b) of Definition~\ref{4.21:def-grain-partition}, apply to them at later sites.
\end{proof}

\begin{remark}[Order of the induction over steps and sites]\label{4.21:rem-induction-order}
Under the inputs of the theorem on the relative small-field machine, the statements of the
relative small-field machine are established for all steps $k$ in one
induction. Within step $k$ they are established at eight sites, which we denote
$o_1,\dots,o_8$ and list in the following order:
\begin{equation}\label{4.21:eq-induction-order-a}
\begin{array}{cl}
o_1 & \text{closure of step } k-1,\\
o_2 & \text{the } \mathbf{T}\text{-operation from level } k-1 \text{ to level } k,\\
o_3 & \text{entry into the operation } \mathbf{R}_k,\\
o_4 & \text{preparation, stages } n=0,\dots,\check{R}_k-1,\\
o_5 & \text{orbit data of the class chain},\\
o_6 & \text{the third-factor lemma at step } k,\\
o_7 & \text{localisation in } \mathbf{R}_k,\\
o_8 & \text{the boundary terms of \cite[(1.90)--(1.99)]{B-CMP122a}; closure of step } k.
\end{array}
\end{equation}
The induction runs over the pairs $(k,o_i)$ in lexicographic order: $(k,o_i)$ precedes
$(k',o_{i'})$ if $k<k'$, or if $k=k'$ and $i<i'$. Here the number of stages of the
preparation at step $k$ is the region-size parameter $\check{R}_k$ of step $k$. Three
hypotheses of the theorem on the relative small-field machine are read at individual sites,
and we name them by the sites that read them. The \emph{site hypothesis} is stated
separately at the $\mathbf{T}$-site, at the preparation site and at the localisation site,
and at the localisation site it has two clauses; the \emph{preparation hypothesis} is read
only at the preparation site; the \emph{localisation hypothesis of step $k$} is read only at
the localisation site of step $k$. For each site we state what is established there and by
which statements, and then what the site uses; it uses nothing else.
\begin{itemize}
\item[$o_1$.] At $o_1$ the closure of step $k-1$ is established, that is, the conclusion of
the relative small-field machine at level $k-1$; orbit data of radius $\varepsilon_k$ on the
surviving law sites, and the orbit-data letter $\mathsf{e}$ on the frames created at step
$k-1$; the source-tube bound at the creating step $k-1$, in the form in which it is stated
at $o_1$ and at $o_8$, together with the two-point bound of the proposition on the
relative potential; and the bounds of the kept-member part of the orbit-datum statement at
the later sites. These are established by the proposition on the boundary terms and the
proposition on the relative potential at the closure site, by the closure clause of the
orbit-datum statement, and by the source-tube bound. This site uses only the sites of step
$k-1$.
\item[$o_2$.] At $o_2$ the relative Mayer potential $V^{\mathrm{rel}}$ is established, and the
terms of level $k$ created by $\mathbf{T}$ are shown to be in format and bounded in the norm
$\mathfrak{N}^{A,\mathbb{C}}$. This is done by the proposition on the $\mathbf{T}$-site, with
its tabulated bounds for single terms and for families of terms at that site, and by the
transfer statement for terms of the second class. The site uses $o_1$; the site hypothesis
at the $\mathbf{T}$-site at interpolation time $\mathsf{t}=0$; clause (i) of the further input
of the relative small-field machine that is used at this site; what is established at step
$k$ up to the margin $\ell 1$; Stage B of the hypothesis drawn from the correlation
theorem; and the source-tube bound only through the share of the orbit-datum statement for
kept members.
\item[$o_3$.] At $o_3$ it is established, by the orbit-datum statement, that every stored
term entering the $\mathbf{R}$-site has an orbit datum of radius $\varepsilon_k$ and satisfies
the bound of its class in the variable $\mathsf{d}:=d(A_F)$, where $A_F$ is the label of the
term. This site uses only $o_1$ and $o_2$.
\item[$o_4$.] At $o_4$ the chain terms $\mathbb{C}^{(n+1)}_k$ on labels are established for
the stages $n=0,\dots,\check{R}_k-1$, bounded by $C_0^{\sharp}$ in the norm
$\mathfrak{N}^{A,\mathbb{C}}$. This is done by the proposition on the preparation site with
its resummation statement; by the counting lemmas of the preparation site (the overlap
lemma, the local star lemma, the per-cube lemma, the window lemma, the anchoring lemma for
all old anchors, and the summation lemma for non-resonant groups); and by the corollary of
the preparation proposition for the class chain. The site uses $o_3$; the site hypothesis
at the preparation site; the preparation hypothesis; the inherited-datum
hypothesis, for the third and fourth members of the list of summable members at the
preparation site; the difference bound for the kept members of live regions (regions still
carrying their large-field factor at level $k$), which is used only in the
derivation of the coupling-comparison inequality at this site, through the orbit-datum
statement; the margins $\ell 3$, $\ell 4$; and the totals of the kept-member part of the
orbit-datum statement at the preparation site, for groups that read a kept member.
\item[$o_5$.] At $o_5$ the orbit data of the class chain, which are clauses (b) and (c) of the
orbit-datum statement, are established by that statement. This site uses only $o_4$.
\item[$o_6$.] At $o_6$ the two bounds of the third-factor lemma at step $k$, the oscillation
bound and the bound that accompanies it, are established by the third-factor lemma. The site
uses $o_1$--$o_5$ of step $k$, which together form the induction hypothesis of step $k$; the
reference statement that the third-factor lemma takes as input, for all pairs of levels not
exceeding $k$; the two clauses of the site hypothesis at the localisation site of step $k$;
the chart read by the third-factor lemma, supplied by part (a) of a proposition which we
call the chart proposition, and whose part (b) is used at $o_8$; the
margins $\ell 6$, $\ell 7$; and, of the statements of the relative small-field machine
itself, only the geometric statement at the closure site, the lemma on the decoupling family
$\sigma_0$, clause (ii) of the source-anchoring statement, the localisation clause of the
source-counting statement, and clauses (D1) and (N) of the statement on deep units. No
length, label, sum or output of $o_7$ is used.
\item[$o_7$.] At $o_7$ the localisation of $\mathbf{R}_k$ is carried out: the values, the
pieces and the relative boundary family $\mathfrak{v}^{\mathrm{rel}}$ are established,
together with the statement of the localisation proposition for the combined family
$\mathfrak{v}+\mathfrak{v}^{\mathrm{rel}}$. This is done by the localisation proposition, with
its bound in the starred norm and its statement on deep units, and by its tabulated bounds
for single terms and for families at the localisation site. The site uses $o_1$--$o_6$ and
the localisation hypothesis of step $k$.
\item[$o_8$.] At $o_8$ the boundary terms of \cite[(1.90)--(1.99)]{B-CMP122a} are treated:
the term $\mathbf{B}'^{(k)}$ is shown to be in format, with its weight; the terms of regions
that are no longer live cease to be carried; and the conclusion of the relative small-field
machine at level $k$ is established. This is done by the proposition on the boundary terms
and the proposition on the relative potential, by part (b) of the chart proposition named
at $o_6$, and by the lemma on the frames of the closure. This site
uses only $o_7$.
\end{itemize}
\end{remark}

\begin{proof}
The pairs $(k,o_i)$, with $1\le k\le K$ and $1\le i\le 8$, are totally ordered
by the lexicographic order of the remark, and each pair $(k,o_i)$ has only finitely many
predecessors: the eight sites of each step $k'<k$ and the sites $o_1,\dots,o_{i-1}$ of step
$k$. Induction along this order is therefore legitimate, provided that what is established
at $(k,o_i)$ uses only what is established at strictly preceding pairs, together with
inputs of the relative small-field machine that are not produced by the induction.

We go through the sites in the list of the remark. At $o_1$ only sites of step $k-1$ are
read. At $o_2$ the site $o_1$ is read; the remaining items read at $o_2$ are inputs, or
statements of step $k$ preceding the margin $\ell 1$, or the source-tube bound
through the share of the orbit-datum statement for kept members, which belongs to
$o_1$. At $o_3$ only $o_1$ and $o_2$ are read. At $o_5$ only $o_4$ is read. At $o_6$ the sites
$o_1$--$o_5$ of step $k$ are read; of the statements of the relative small-field machine
itself only the geometric statement at the closure site, the lemma on $\sigma_0$, the
source-anchoring and source-counting clauses and the two clauses of the statement on deep
units listed in the remark are used, and none of these involves a length, label, sum or
output of $o_7$, so that $o_6$ reads nothing established at $o_7$ or later. At $o_7$ the sites
$o_1$--$o_6$ and the localisation hypothesis of step $k$, an input, are read. At $o_8$ only
$o_7$ is read.

It remains to treat $o_4$, where the counting lemmas of the preparation site are used at
every stage. These lemmas, namely the overlap lemma, the local star lemma, the per-cube
lemma, the window lemma, the anchoring lemma for all old anchors and the summation lemma
for non-resonant groups, are field-free: they read the label, the indices $(k,h,j)$ (step,
level and top zone), the zone geometry of the step, which is a hypothesis of the theorem on
the relative small-field machine, the definition of arrested attempts
with its accompanying lemma, and the per-cube letters of the per-cube lemma at stage
$n+1$. The sources of these per-cube letters are of two kinds. For frozen terms they are
the bounds of clause (f2) of the format, in the conclusion of the relative small-field
machine at the earlier steps. Otherwise they are supplied by the induction of $o_4$ itself
over the stages $i\le n$ of the class chain; for a sibling component $C'$, whose number of
stages satisfies $n(C')\le n$, the stages $i\le n(C')$ are among the stages $i\le n$ of this
induction. Hence stage $n+1$ of $o_4$ uses only the stages $i\le n$ at step $k$, the sites of
earlier steps, $o_3$, the inputs listed in the remark and the zone geometry, and no output
of $o_5$--$o_8$.
Within $o_4$ this is an induction over the finite set of stages, and it is contained in the
position of $o_4$ in the order.

Thus every site reads only pairs that strictly precede it, together with inputs, and the
induction over pairs $(k,o_i)$ in lexicographic order is well-founded.
\end{proof}

\begin{corollary}[Stored boundary terms in Ba{\l}aban's relative length]\label{4.21:cor-stored-term-bound}
Let $k\le K'$, where $K'$ denotes the last step of the range of steps for which the relative format
of boundary-type terms (Definition~\ref{4.21:def-relative-format}) is established, and consider any
label of the expansion at level $k$ and any stored boundary term
$F$ of that label, of level $j\le k$. Let $(A_F,\mathcal{P}_F)$ be its data in the relative format
of boundary-type terms (Definition~\ref{4.21:def-relative-format}). Put
\begin{equation*}
W_F:=\bigcup_{(Y,m)\in\mathcal{P}_F}Y^{\sharp},\qquad X:=A_F\cup W_F,
\end{equation*}
and write the term, in Ba{\l}aban's notation, as
$\mathbf{B}^{(j)}(X,(\mathbf{U},\mathbf{J}),\mathbf{A};w)$. Then
\begin{equation}\label{4.21:eq-stored-term-bound-1}
\bigl|\mathbf{B}^{(j)}(X,(\mathbf{U},\mathbf{J}),\mathbf{A};w)\bigr|
\le B_0\exp\bigl(-\kappa\, d_{j,W_F}(X)\bigr),
\end{equation}
where, for a set $W$, $d_{j,W}(X)$ is $M^{-1}$ times the length of a shortest tree graph
contained in $X$ and intersecting all $M$-cubes of $X\setminus W$. The bound and the relative
length $d_{j,W}$ are in the form of \cite[(1.67)]{B-CMP122b}.
\end{corollary}

\begin{proof}
Let $\mathcal{T}$ be an optimal tree graph for $d_j(A_F)$, the quantity fixed in the glossary
(Notation~\ref{0.2:not-glossary}): its length is $M\,d_j(A_F)$ and it intersects all $M$-cubes
of $A_F$.

First, $\mathcal{T}$ lies in $X$: as in the relative length above in the case $W=\emptyset$, the
tree graphs over which $d_j(A_F)$ is minimised are contained in $A_F$, and
$A_F\subset A_F\cup W_F=X$.
Second, $\mathcal{T}$ intersects all $M$-cubes of $X\setminus W_F$. Indeed,
\begin{equation*}
X\setminus W_F=(A_F\cup W_F)\setminus W_F\subset A_F,
\end{equation*}
so every $M$-cube of $X\setminus W_F$ is an $M$-cube of $A_F$, and $\mathcal{T}$ intersects all
$M$-cubes of $A_F$. Hence $\mathcal{T}$ is one of the tree graphs over which the minimum
defining $d_{j,W_F}(X)$ is taken, and
\begin{equation}\label{4.21:eq-stored-term-bound-2}
d_{j,W_F}(X)\le M^{-1}\operatorname{len}(\mathcal{T})=d_j(A_F),
\end{equation}
where $\operatorname{len}(\mathcal{T})$ denotes the length of $\mathcal{T}$.

Clause (f2) of the relative format, in \eqref{4.21:eq-relative-format-a}, implies the bound
$B_0\exp(-\kappa\, d_j(A_F))$ for the stored term. Since $\kappa>0$, the function
$t\mapsto\exp(-\kappa t)$ is decreasing, so \eqref{4.21:eq-stored-term-bound-2} gives
\begin{equation*}
B_0\exp\bigl(-\kappa\, d_j(A_F)\bigr)\le B_0\exp\bigl(-\kappa\, d_{j,W_F}(X)\bigr),
\end{equation*}
and \eqref{4.21:eq-stored-term-bound-1} follows.
\end{proof}

The relative format carries more information than \eqref{4.21:eq-stored-term-bound-1}. In it,
every strip $Y^{\sharp}$, $(Y,m)\in\mathcal{P}_F$, is touched by the optimal tree graph of $A_F$
used in the proof above. Moreover the label $A_F$
is a datum of the term $F$, not a function of its domain $X$. At the step at which the term is
created one has $d_j(A_F)\le d_{j,W_F}(X)$ in addition to \eqref{4.21:eq-stored-term-bound-2},
so that the two lengths are equal there; for the terms derived from $F$ at later steps only the
inequality \eqref{4.21:eq-stored-term-bound-2} remains, and it may be strict.
The bound \eqref{4.21:eq-stored-term-bound-1} is the one displayed in the correlation theorem;
the two-point witness theorem shows that the statements which use the stored terms need the
additional information carried by the relative format.

\begin{definition}[Running decay rate and step-wise constants]\label{4.21:not-running-rates}
Fix the step $k$ of the renormalisation group. Let $g_0,\dots,g_k$ be the running couplings of
Definition~\ref{1.2:def-couplings} at the levels up to $k$. The following constants are used at
the preparation site of the operation $\mathbf{R}_k$.
\begin{enumerate}
\item[(a)] \emph{Running decay rate.} Let $\kappa_A^{\mathrm{pr}}(g')$ be the decay rate that the
setting of the preparation site assigns to a step taken at coupling $g'$. At step $k$ the rate used
is the running minimum $\kappa_A(k)$ of this rate over the couplings already reached. The smallness
constant $\varepsilon_S$ is defined from $\kappa_A(k)$ and from the constant $c_{\flat}$ of the grain
lemma of the preparation site:
\begin{equation}\label{4.21:eq-running-rates-a}
\kappa_A(k):=\min_{i\le k}\kappa_A^{\mathrm{pr}}(g_i),
\qquad
\varepsilon_S:=2e^{-(\kappa_A(k)-c_{\flat})}.
\end{equation}
\item[(b)] \emph{Box containment.} Let $\bar c$ be a component and $R(\bar c)$ its region, where
$j$ is the zone index attached to $R(\bar c)$ in the setting of the preparation site. Then
\begin{equation}\label{4.21:eq-running-rates-b}
R(\bar c)\ \text{lies in a box of side at most}\ 100\,\check{R}_{j+1}.
\end{equation}
The boxes in clauses (f)($\alpha$) and (g) of the geometry lemma of the preparation site are
obtained in the same way as \eqref{4.21:eq-running-rates-b}.
The nesting property of the zones is accordingly an input of this step.
\item[(c)] \emph{The step-wise constant $\gamma_0$.} The constant $\gamma_0$ of the site bound at
step $k$ is
\begin{equation}\label{4.21:eq-running-rates-c}
\gamma_0^{(k)}:=5\max_{j\le k}\bigl(w_*(j)\,\bar\alpha_j\bigr)\le\bar\gamma_0(\mathsf{t}_k).
\end{equation}
Here $w_*(j)$ is the weight attached to step $j$ and $\bar\alpha_j$ the bound on the radii at
level $j$, both as they enter the site bound; $\bar\gamma_0$ is the function of the single variable
$\mathsf{t}$ that the site bound supplies, and $\mathsf{t}_k$ is the value of $\mathsf{t}$ at
step $k$. The right-hand side is thus a bound in the single variable $\mathsf{t}$, evaluated at
$\mathsf{t}_k$.
\item[(d)] \emph{The exponent at the preparation site.} At the preparation site, the input bound
governing the reception of stored terms is used in its displayed form at the exponent $r'$. Here
$r'$ is defined from the decay rate $r_{\mathrm{hull}}$ of the $\pi_k$-hull, which is at most
$\tfrac85\kappa$:
\begin{equation}\label{4.21:eq-running-rates-d}
r':=r_{\mathrm{hull}}+2+1\le\tfrac85\kappa+3\le2\kappa+3 .
\end{equation}
Terms of the kind $(\delta)(\beta)$, formed from the groups $(\beta)$ and $(\delta)$ of the expansion
at the preparation site, are bounded by the fourth clause of the counting lemma of the preparation
site, and not by $A_{\flat}b_{\delta}$. Terms of group $(\delta)$ are stored in the form given by
clauses (a) and (f) of the base-frame statement, with
$\mathfrak{N}^{A}:=\mathfrak{N}^{G}$. No further level condition is imposed on group $(\delta)$.
\item[(e)] \emph{The class $Z$.} In clause (a) of the grain lemma of the preparation site, $Z$ is
the class at the start of the present operation $\mathbf{R}_k$. With this choice a sibling excluded
at an earlier stage of the same operation $\mathbf{R}_k$ lies in $Z$, its grain being given by
clauses (b) and (c) of that lemma at the running index $j$. A sibling component $C'$ keeps its
zones of index at most $j(C')$ inside $Z$, by the statement on sibling components of the
preparation site:
\begin{equation}\label{4.21:eq-running-rates-e}
\text{every zone of } C' \text{ of index } \le j(C') \text{ is contained in } Z .
\end{equation}
\item[(f)] \emph{Per-cube floor.} The per-cube bound of the preparation site remains the floor
required by the correlation theorem for the cubes decoupled at Stage~B of the decoupling expansion
of the preparation site. These are the cubes $\Delta$ that belong to the family $\mathsf{Q}$ of
fluctuation cubes of the per-cube bound, or that are proper-scale cubes in the sense of that bound:
\begin{equation}\label{4.21:eq-running-rates-f}
\Delta\in\mathsf{Q}\quad\text{or}\quad \Delta\ \text{is a proper-scale cube}.
\end{equation}
\item[(g)] \emph{Coarser levels at old blocks.} A term of plain type placed at a block of an old
zone meets a number of coarser levels equal to
\begin{equation}\label{4.21:eq-running-rates-g}
k-n',
\end{equation}
where $n'$ is the level attached to the old zone of that block.
\item[(h)] \emph{Two bounds of the summary.} In the summary of the bounds at the preparation
site, two displayed bounds are inputs of this step: the tube bound $\Sigma^{\mathfrak{q}}$ on the
source functional $\mathfrak{q}_C(\zeta)$ is the input stating the tube bound for
$\mathfrak{q}_C(\zeta)$, and the second displayed bound is the input stating the same bound.
\item[(i)] \emph{The relative comparison for chain terms.} The comparison statement used at this
step is the relative comparison theorem for the chain class $\mathbb{C}$.
\item[(j)] \emph{Grains.} The preliminary step of the resummation at the preparation site declares
no objects; in particular, $\pi^{\flat}$ is not among the objects declared in the opening
paragraph of the summary.
\end{enumerate}
\end{definition}

\begin{proof}
(a) The index set $\{i:i\le k\}$ is finite and non-empty. Hence the minimum in
\eqref{4.21:eq-running-rates-a} exists, and $\kappa_A(k)$ is well defined. With $\kappa_A(k)$
fixed, $\varepsilon_S$ is given by the explicit exponential in \eqref{4.21:eq-running-rates-a}.

(b) The containment \eqref{4.21:eq-running-rates-b} follows from the second part of clause~(ii),
written $(\mathrm{ii})_2$, taken together with
the nesting property of the zones. The boxes in clauses (f)($\alpha$) and (g) of the geometry lemma
come from the same two facts.

(d) By definition, $r'$ is $r_{\mathrm{hull}}+3$. The decay rate of the $\pi_k$-hull satisfies
$r_{\mathrm{hull}}\le\tfrac85\kappa$. Adding $3$ to both sides gives the first inequality in
\eqref{4.21:eq-running-rates-d}. The parameter $\kappa$ is positive, so
$\tfrac85\kappa\le2\kappa$. Adding $3$ to both sides gives the second inequality.

(g) The $k-n'$ coarser levels in \eqref{4.21:eq-running-rates-g} are accounted for in two places.
For terms of the format class $\mathfrak{F}$, they are covered by the bound on anchored totals over
the enlarged set $\mathfrak{O}^{+}$ of old blocks. For the terms entering the correlation theorem,
they are covered by the inherited-datum hypothesis.
\end{proof}

\begin{definition}[Orbit data, complex size functionals and the stored domain]
\label{4.21:not-orbit-data}
In this definition $G^{\mathbb{C}}$ denotes the complexified gauge group and $\|g\|$ the operator
norm of $g\in G^{\mathbb{C}}$ as an $N\times N$ matrix. The labels $A_F$, regions $X_F$, pocket lists
$\mathcal{P}_F$ and strips $Y^{\sharp}$ of boundary-type terms are those of the relative format of
boundary-type terms (Definition~\ref{4.21:def-relative-format}).

\emph{Gauge size and charts.} For $g\in G^{\mathbb{C}}$ we put
\begin{equation*}
\operatorname{pol}(g)=\log\max\{\|g\|,\|g^{-1}\|\},\qquad
\mathfrak{G}_{\varepsilon}=\{g\in G^{\mathbb{C}}:\ \operatorname{pol}(g)<\varepsilon\}.
\end{equation*}
For a region $X$ we write $\mathfrak{G}_{\varepsilon}(X)$ for the chart of radius $\varepsilon$ over
$X$, that is, the set of $G^{\mathbb{C}}$-valued gauge transformations $g$ on $X$ with
$\operatorname{pol}(g(x))<\varepsilon$ for every $x\in X$. The lemma on the size of complex gauge
elements supplies the subadditivity and logarithm bound
$(\mathsf{g}1)$ and the identity principle $(\mathsf{g}3)$ on products
$\prod_s\mathfrak{G}_{\varepsilon_s}$, displayed in \eqref{4.21:eq-orbit-data-a}. Items
$(\mathrm{O}1)$, $(\mathrm{O}2)$ and $(\mathrm{D})$ of that display are explained in the paragraphs
on orbit data and on the stored domain below.

\begin{equation}\label{4.21:eq-orbit-data-a}
\begin{aligned}
&(\mathsf{g}1)\quad \operatorname{pol}(g_1g_2)\le\operatorname{pol}(g_1)+\operatorname{pol}(g_2),
\qquad \operatorname{pol}(e^{Z})\le|Z|;\\
&(\mathsf{g}3)\quad \text{two holomorphic functions on a product }
\textstyle\prod_s\mathfrak{G}_{\varepsilon_s}\ \text{of charts (finite index set)}\\
&\qquad\qquad\text{that agree at all $G$-valued points agree identically};\\
&(\mathrm{O}1)\quad F^{\sharp}\ \text{is holomorphic in }(P,\ \text{weak coordinates},\ g);\\
&(\mathrm{O}2)\quad F^{\sharp}(P;\mathbf{A},\boldsymbol{\Phi},g)
=F(P;\mathbf{A}^{[g]},\boldsymbol{\Phi}^{[g]})\quad\text{at $G$-valued }g;\\
&(\mathrm{D})\quad\mathfrak{U}_F=\bigl\{P\in\mathfrak{U}^{\mathrm{tr}}_F:\ Q_a(y)<\theta_a(y)
-\varpi_{\mathrm{fr}}\,r_a(y)\ \text{for every }a\notin\{5,8\}\\
&\qquad\qquad\text{and every clause location $y$ whose stencil meets some }
\mathsf{B}_W\bigr\}.
\end{aligned}
\end{equation}

\emph{Law sites.} The law sites of a region $S$ form the disjoint union
\begin{equation*}
\mathfrak{s}(S)=\mathfrak{s}^{\mathrm{pt}}(S)\sqcup\bigsqcup_{W}\Sigma_W
\end{equation*}
of the pointwise law sites $\mathfrak{s}^{\mathrm{pt}}(S)$ and the law sites $\Sigma_W$ of the
frames $W$ described below.

\emph{Orbit data.} Let $(F,X_F,\mathcal{S})$ be a stored term with strong set $\mathcal{S}$. An
\emph{orbit datum of radius $\varepsilon$} of $(F,X_F,\mathcal{S})$ is a function $F^{\sharp}$ on
\begin{equation*}
\mathfrak{U}_F(X_F)\times\mathfrak{D}_{\mathcal{S}}(X_F)\times\operatorname{supp}\boldsymbol{\Phi}
\times\mathfrak{G}_{\varepsilon}(X_F)
\end{equation*}
that obeys $(\mathrm{O}1)$ and $(\mathrm{O}2)$ of \eqref{4.21:eq-orbit-data-a}; here
$\mathfrak{D}_{\mathcal{S}}(X_F)$ is the data domain over $X_F$ attached to the strong set
$\mathcal{S}$, in which $\mathbf{A}$ ranges. Further,
$\mathbf{A}^{[g]}$, $\boldsymbol{\Phi}^{[g]}$ are the gauge transforms of $\mathbf{A}$,
$\boldsymbol{\Phi}$ by $g$, and the weak coordinates in $(\mathrm{O}1)$ are those of the data
$\mathbf{A}$ on $\mathfrak{D}_{\mathcal{S}}(X_F)$. The complex size functionals are
\begin{equation*}
\mathfrak{N}^{\mathbb{C}}_{\varepsilon}(F)=\sup\,|F^{\sharp}|,\qquad
\mathfrak{N}^{A,\mathbb{C}}(F)=\sum_{\mathcal{S}}e^{\kappa_A|\mathcal{S}|}\,
\mathfrak{N}^{\mathbb{C}}_{\varepsilon}\bigl(F(\cdot,\mathcal{S})\bigr),
\end{equation*}
the supremum being taken over the domain of $F^{\sharp}$. One has
$\mathfrak{N}^{\mathbb{C}}_{\varepsilon}(F)\ge\mathfrak{N}(F)$.

\emph{What a wrapped term carries.} Let $F$ be a wrapped term, written in the representation
\cite[(1.102)]{B-CMP122b}. Every pocket $(Y,m)\in\mathcal{P}_F$ is a \emph{frame} $W=Y$ of level
$k_W=m$, with real data $\Phi_W$ and law sites
\begin{equation*}
\Sigma_W\subset\partial^{+}\Lambda_W\subset Y,
\end{equation*}
where $\Lambda_W$ is the region of the frame $W$ and $\partial^{+}\Lambda_W$ its outer boundary
layer. A strip holds no pointwise law site: its gauge data are read through $\mathbf{U}$ and
belong to the fifth class of the classification of the variables read by a term. The
pointwise law sites of $F$ lie in $A_F$. Every wrapped term has a law site.

\emph{The stored domain.} Put $\varpi_{\mathrm{fr}}=\tfrac18(1-\theta_S)\beta$, where $\beta$ is
the schedule constant of Definition~\ref{4.4:def-post-step-space} and $\theta_S$ is the constant of
the frame-reserve storage condition. For a frame $W$, let
$\mathsf{B}_W$ be the reserve set of $W$ in the frame-reserve storage condition; it lies within one
unit of level $k_W$ of $\Sigma_W$. The term $F$ is
\emph{stored} on the set $\mathfrak{U}_F\subset\mathfrak{U}^{\mathrm{tr}}_F$ defined in item
$(\mathrm{D})$ of \eqref{4.21:eq-orbit-data-a}. The set
$\mathfrak{U}^{\mathrm{tr}}_F$, the quantities $Q_a(y)$, $\theta_a(y)$, $r_a(y)$, the clause
locations and their stencils are those of the frame-reserve storage condition: $\mathfrak{U}^{\mathrm{tr}}_F$
is the parameter domain on which that condition is imposed, a clause location $y$ is a point at
which a clause of that condition is imposed and its stencil is the set of lattice points that
clause reads, and at a clause location $y$ the
quantity $Q_a(y)$ is the one controlled by the clause of index $a$, with threshold $\theta_a(y)$ and
reserve weight $r_a(y)$.

\emph{The carried gauge map.} Let $\mathsf{u}(\mathsf{s})=\mathfrak{u}[\cdot;\cdot]$ be the gauge map
carried by a link of Definition~\ref{4.21:def-links-readers}. The lemma on the carried gauge map
states the following properties.
\begin{enumerate}
\item $\mathfrak{u}$ is holomorphic in the site's parameters (the decoupling parameters
$\mathsf{s}$ of the link).
\item $\mathfrak{u}\equiv1$ on top blocks (the blocks of the top level of the site) where the
layer vanishes.
\item $\mathfrak{u}(z)$ reads the layer on $B^{\mathrm{top}}(z)$ only.
\item The size of $\mathfrak{u}(z)$ is controlled as follows. Here $\mathcal{K}$ denotes the datum
of the link from which the layer is built, $N_y(\mathcal{K})$ the pointwise size of $\mathcal{K}$ at
the point $y$, and $\Theta_0^{\mathrm{loc}}(\mathcal{K})(z)$ the local size functional entering the
bound on $\mathfrak{u}(z)$, which the second inequality controls by the pointwise sizes over
$B^{\mathrm{top}}(z)$:
\begin{equation}\label{4.21:eq-orbit-data-1}
|\log\mathfrak{u}(z)|\le K_u\,\Theta_0^{\mathrm{loc}}(\mathcal{K})(z),\qquad
\Theta_0^{\mathrm{loc}}(\mathcal{K})(z)\le\sup_{B^{\mathrm{top}}(z)}N_y(\mathcal{K}).
\end{equation}
\item $B^{\mathrm{top}}(z)$ lies in the $M$-cube of the local scale of $z$, inside the term's
domain.
\end{enumerate}
\end{definition}

\begin{proof}[Proof that $\mathfrak{N}^{\mathbb{C}}_{\varepsilon}(F)\ge\mathfrak{N}(F)$]
Recall that $\mathfrak{N}(F)$ is at most the supremum of $|F|$ over
$\mathfrak{U}_F(X_F)\times\mathfrak{D}_{\mathcal{S}}(X_F)\times\operatorname{supp}\boldsymbol{\Phi}$
(Definition~\ref{4.21:def-relative-format}). The identity
$\mathbf{1}$ has norm one, and so does its inverse. Hence $\operatorname{pol}(\mathbf{1})=0<\varepsilon$.
It follows that the constant gauge transformation with value $\mathbf{1}$ lies in
$\mathfrak{G}_{\varepsilon}(X_F)$.

This transformation is $G$-valued and acts trivially, so
$\mathbf{A}^{[\mathbf{1}]}=\mathbf{A}$ and $\boldsymbol{\Phi}^{[\mathbf{1}]}=\boldsymbol{\Phi}$.
Condition $(\mathrm{O}2)$ of \eqref{4.21:eq-orbit-data-a} then gives
\begin{equation*}
F^{\sharp}(P;\mathbf{A},\boldsymbol{\Phi},\mathbf{1})=F(P;\mathbf{A},\boldsymbol{\Phi})
\end{equation*}
for every $(P,\mathbf{A},\boldsymbol{\Phi})$ in
$\mathfrak{U}_F(X_F)\times\mathfrak{D}_{\mathcal{S}}(X_F)\times\operatorname{supp}\boldsymbol{\Phi}$.

The supremum of $|F^{\sharp}|$ is at least its supremum over the slice $g=\mathbf{1}$. By the last
display, that slice supremum is the supremum of $|F|$ over
$\mathfrak{U}_F(X_F)\times\mathfrak{D}_{\mathcal{S}}(X_F)\times\operatorname{supp}\boldsymbol{\Phi}$,
which is at least $\mathfrak{N}(F)$.
\end{proof}

On the circle $|\mathsf{s}(\Delta)|=e^{\kappa_1}$, for a cell $\Delta$ of the link's decoupling
family (Definition~\ref{4.21:def-links-readers}), the carried map $\mathsf{u}(\mathsf{s})$ is not
$G$-valued,
so that $(\mathrm{O}2)$ does not apply there. Hence, for a term with a law site, the holomorphy
statement \eqref{4.21:eq-links-readers-b} holds only with sizes measured by
$\mathfrak{N}^{\mathbb{C}}_{\varepsilon}$. Since every wrapped term has a law site (see above), the
relative format of Definition~\ref{4.21:def-relative-format} is stated with
$\mathfrak{N}^{A,\mathbb{C}}$.

\begin{remark}[The carried statements on the enlarged set of old blocks]
\label{4.21:rem-old-blocks-enlarged}
Let $\mathfrak{O}$ be the set of old blocks with which the preparation site is first set up, and
let $\mathfrak{O}^{+}$ be the enlarged set of old blocks, with its window, as fixed in the definition
of the enlarged set of old blocks and of the window points at the preparation site. For a
block $\square$ of the cover we write $\tilde{\square}$ for its probe. We write $d^{\wedge}(\square,R_n)$
for the scaled distance from $\square$ to the collar region $R_n$ of stage $n$, and $w=M/M_1$. We
write $\check{R}_k$ for the region size parameter at step $k$. For an arrested attempt
$\mathfrak{a}$ of an earlier step and its window we write:
\begin{itemize}
\item $k_a$ for the step of $\mathfrak{a}$;
\item $\mathopen{]}h_a,k_a]$ for the range of levels re-labelled by $\mathfrak{a}$;
\item $\Gamma''_i$ for the strata of the window;
\item $T_{\mathfrak{a}}$ for the top raised zone of the window, and $j_a$ for its level;
\item $j^{+}$ for the level which the definition of an arrested attempt attaches to $\mathfrak{a}$,
with $j_a\le j^{+}$;
\item $X_{\mathfrak{a}}$ for the region of $\mathfrak{a}$.
\end{itemize}
Let $C$ be the component and $Z_h$ the zone of level $h$ of the geometry lemma. Of the definition
of $\mathfrak{O}^{+}$ we use only the following consequence of its second clause: a block whose
probe meets $Z_h\cap C$ belongs to $\mathfrak{O}^{+}$.

At the following places the set of old blocks is read as $\mathfrak{O}^{+}$ in place of
$\mathfrak{O}$, and the window as the one attached to $\mathfrak{O}^{+}$:
\begin{itemize}
\item in the statement and in the proof of the proposition bounding the volume sums over anchor
blocks at the preparation site;
\item in the first item of the statement summarising that site's bounds.
\end{itemize}
These are all the places at which the set of old blocks occurs in the statements and proofs carried
over at the preparation site; no other carried statement reads $\mathfrak{O}$. With this reading,
the readings (ii) and (ii$'$) below are adopted, and the statements (i), (iii), (iv) and (v) hold.
\begin{enumerate}
\item[(i)] The old-block estimate used in the proof of that proposition holds for
every block $\square_0\in\mathfrak{O}^{+}$.
\item[(ii)] The last clause of the arrest property of the domain arrangement of the step reads as
follows.
\begin{itemize}
\item The ranges $\mathopen{]}h_a,k_a]$ of re-labelled levels of distinct arrested attempts are
disjoint.
\item The strata $\Gamma''_i$, $i\le j_a-1$, of an older window lie in zones of level at most
$k_a-1\le h-1$.
\item The top raised zone $T_{\mathfrak{a}}$ of that window, of level $j_a\le k_a\le h$, lies in
$X_{\mathfrak{a}}$, and
\begin{equation}\label{4.21:eq-old-blocks-enlarged-1}
X_{\mathfrak{a}}\subset Z_{k_a}\subset Z_h .
\end{equation}
\item Every block whose probe meets either the strata or the top raised zone lies in
$\mathfrak{O}^{+}$:
\begin{equation}\label{4.21:eq-old-blocks-enlarged-a}
\tilde{\square}\cap\Bigl(\bigcup_{i\le j_a-1}\Gamma''_i\ \cup\ T_{\mathfrak{a}}\Bigr)\neq\emptyset
\ \Longrightarrow\ \square\in\mathfrak{O}^{+}.
\end{equation}
\end{itemize}
\item[(ii$'$)] In the proof of that proposition, the justification given after the old-block
estimate for the membership of blocks in the set of old blocks reads: the strata and the top raised
zone of an older window lie inside
$Z_h\cap C$, so that a block whose probe meets one of them belongs to $\mathfrak{O}^{+}$ by the
second clause of its definition.
\item[(iii)] A probe that meets $Z_h\cap C$ has level at most $h$.
\item[(iv)] The second per-block bound and the depth bound of the preparation site hold on
$\mathfrak{O}^{+}$. In particular, with $n$ the stage at which the depth bound is read,
\begin{equation}\label{4.21:eq-old-blocks-enlarged-2}
d^{\wedge}(\square_0,R_n)\ \ge\ w\,(\check{R}_k-5)\qquad\text{for every }\square_0\in\mathfrak{O}^{+}.
\end{equation}
\item[(v)] A lemma of the preparation site uses, in its second clause, the arrest property of
the domain arrangement only through the statement that no level
above $h$ is re-labelled by an earlier attempt; the reading (ii) leaves that statement, and hence
that clause, unchanged.
\end{enumerate}
\end{remark}

\begin{proof}
We first prove (iv). The second per-block bound and the depth bound of the preparation site are
derived afresh, on the enlarged set $\mathfrak{O}^{+}$, in the all-anchor bound at the preparation
site. The depth bound on $\mathfrak{O}^{+}$ gives in particular the inequality
\eqref{4.21:eq-old-blocks-enlarged-2}. This proves (iv).

For (i), the proof of the old-block estimate uses the fact that a block $\square_0$ is old at one
point only, namely through the lower bound $d^{\wedge}(\square_0,R_n)\ge w(\check{R}_k-5)$. By (iv)
this lower bound holds for every $\square_0\in\mathfrak{O}^{+}$, in the form
\eqref{4.21:eq-old-blocks-enlarged-2}. The proof of the estimate therefore applies verbatim to every
block of $\mathfrak{O}^{+}$, which is (i).

For (ii), consider an arrested attempt $\mathfrak{a}$ of an earlier step $k_a$. Item (ii) is the
reading adopted for the last clause of the arrest property of the domain arrangement of the step.
The disjointness of
the ranges $\mathopen{]}h_a,k_a]$ of re-labelled levels of distinct arrested attempts is the first
part of that clause and is not affected by the reading (ii).
The one point at which the reading must be checked against the definition of an arrested attempt is
the inequality $j_a\le k_a$ for the level of the top raised zone.

This inequality holds, since $j_a\le j^{+}\le k_a$ by the definition of an arrested attempt. On the
fifth to seventh of the branches of that definition, and on its fourth branch at the stage
$n=\check{R}_{k_a}-1$, that definition gives moreover the equality $j_a=k_a$. Together with
$k_a\le h$, which is part of the reading (ii), this gives $j_a\le k_a\le h$.

It remains to account for the last part of (ii). The strata lie in zones of level at most $h-1$,
and the top raised zone, which lies in $X_{\mathfrak{a}}$, lies in $Z_h$ by
\eqref{4.21:eq-old-blocks-enlarged-1}. The reading (ii$'$), adopted for the proof of the
proposition, places both the strata and the top raised zone of the older window inside
$Z_h\cap C$.

A block whose probe meets one of them therefore meets $Z_h\cap C$ through its probe. Such a block
belongs to $\mathfrak{O}^{+}$ by the second clause of its definition. This is
\eqref{4.21:eq-old-blocks-enlarged-a}, and (ii) follows.

For (iii), the proof of the fourth part of the last clause of the geometry lemma is carried out with
the set $Z_h\cap C$. Read with this set, it gives that a probe meeting $Z_h\cap C$ has level at
most $h$.

For (v), the reading (ii) of the last clause of the arrest property of the domain arrangement
retains the disjointness of
the ranges $\mathopen{]}h_a,k_a]$ unchanged, and otherwise concerns only the location of the strata
and of the
top raised zone of an older window and the membership in $\mathfrak{O}^{+}$ of the blocks whose
probes meet them. It does not change which levels are re-labelled by earlier attempts. Hence the
statement that no level above $h$ is re-labelled by an earlier attempt, through which alone the
clause in question uses the arrest property, is the same with and without the reading (ii), and that
clause is unchanged.
\end{proof}

\begin{proposition}[Every step reads a stored term through one supremum]
\label{4.21:prop-majorant-structure}
Assume the following four statements on the second stages of the sites of the relative machine.
\begin{enumerate}
\item[(a)] In the second stage of the correlation theorem, the potential $V$ enters the
expansion only as Mayer factors, and after the passage to the polymer representation no $V$
occurs.
\item[(b)] The second stage at the preparation site reads its elements only through the per-cube
quantities $(\mathsf{b},\mathsf{d},\nu)$.
\item[(c)] The strip theorem of the birth step, in the form carried by
Proposition~\ref{4.21:prop-birth-site}\,(a), holds for any family of numbers (or of functions,
in the supremum norm).
\item[(d)] The second-stage steps at the preparation site on $\hat{\mathfrak{E}}$ are majorant
series in the weights $\tilde{\nu}(\hat{e})$.
\end{enumerate}
Consider the four sites of the relative machine: the fluctuation integral $(\mathrm{T})$ of
Proposition~\ref{4.21:prop-fluctuation-site}, the preparation site $(\mathrm{P})$ of
Proposition~\ref{4.21:prop-preparation-machine}, the large-field site $(\mathrm{L})$ of
Proposition~\ref{4.21:prop-large-field-site}, and the birth of new boundary terms of
Proposition~\ref{4.21:prop-birth-site}. Then each of them has the following properties.
\begin{itemize}
\item The site reads a stored term $F$ only through the supremum of one function, namely
$\tau\mapsto F^{\sharp}\circ\mathsf{R}(\tau)$ over the site's datum, whatever the parameter
manifold. Here $F^{\sharp}$ is the orbit datum of $F$ and $\mathsf{R}$ is the map of the site at
hand (at $(\mathrm{T})$ and $(\mathrm{L})$, the link map $\mathsf{R}(\mathsf{s})$).
\item Its first-stage pieces are linear in $F$, and they pass the parameters on.
\item Its second stage is a majorant series in these suprema.
\item Its index sets, labels, cores, corridors, grains and counts depend on the scaffold and on
the set $\mathcal{R}_F$ on which $F$ reads the response, and on nothing else.
\end{itemize}
Precisely, write $\mathcal{E}(F)$ for the value, or the difference of values, that the site forms
from a stored term $F$ (at $(\mathrm{T})$ this is $F(\mathsf{R}(\mathbf{1}))-F(\mathsf{R}(0))$),
and $\mathfrak{N}(F)$ for the supremum of $|F^{\sharp}\circ\mathsf{R}|$ over the site's datum; at
$(\mathrm{T})$ this is the size functional in the bound of
Lemma~\ref{4.21:lem-first-stage-expansion}\,(ii). Then for every stored term $F$ there are an
index set $\Pi(F)$, pieces $F_\pi$ and numbers $c_\pi\ge0$ ($\pi\in\Pi(F)$) such that
\begin{equation}\label{4.21:eq-majorant-structure-a}
\begin{aligned}
&\text{(M1)}\quad \mathcal{E}(F)=\sum_{\pi\in\Pi(F)}F_\pi;\\
&\text{(M2)}\quad \Pi(F)\ \text{is read from the scaffold and from } \mathcal{R}_F \text{ only};\\
&\text{(M3)}\quad |F_\pi|\le c_\pi\,\mathfrak{N}(F);\\
&\text{(M4)}\quad \text{the counts of the second stage use numbers only}.
\end{aligned}
\end{equation}
Here each $F_\pi$ arises from $\tau\mapsto F^{\sharp}\circ\mathsf{R}(\tau)$ by the operations
listed below under (M1); the data from which $\Pi(F)$ is read, the numbers $c_\pi$ and the
counting statements are listed below under (M2), (M3) and (M4).
The entries are as follows.

\emph{Site $(\mathrm{T})$.}
\begin{itemize}
\item (M1): the fundamental theorem of calculus in $\mathsf{s}$ on the finite set $\pi^{T}$,
\cite[(1.9)]{B-CMP116}. This consists of evaluations at $0$ and integrals
$\int_0^1 d\mathsf{s}\,\partial_{\mathsf{s}}$, as in
Lemma~\ref{4.21:lem-first-stage-expansion}\,(i).
\item (M2): $\Pi(F)=\mathfrak{Y}(F)$, determined by $(\pi^{T},\mathfrak{F})$ and by
$\mathcal{R}_F=X_F$.
\item (M3): $c_{\mathsf{y}}=e^{-(\kappa_1-1)|\mathsf{y}|}$. This follows from Cauchy's formula, as in
Lemma~\ref{4.21:lem-first-stage-expansion}\,(ii).
\item (M4): clauses (c) and (d) of Lemma~\ref{4.21:lem-source-anchored-sums}, clause (iv) of
Lemma~\ref{4.21:lem-probe-mass} and Lemma~\ref{4.21:lem-wrapped-cube-bound}. After these come the
Mayer and polymer majorant steps of the correlation theorem's second stage, in the activities
$\nu(Y)$.
\end{itemize}

\emph{Site $(\mathrm{P})$.}
\begin{itemize}
\item (M1): M\"obius differences in $(t,u)$, the two parameters of the first-stage identity of the
preparation site, followed by the fundamental theorem of calculus in $\mathsf{s}$; this is the
first-stage identity of the preparation site, unchanged. The result is
then summed by fibres, using the fundamental theorem of calculus in $s$ on
$\hat{\sigma}_0(A)$ (Lemma~\ref{4.21:lem-grain-resummation}\,(i)).
\item (M2): $\Pi(F)$ consists of the data $\mathcal{A}_c$, $\mathcal{Q}^{w}_c$ and
$\hat{\mathsf{y}}$. These are determined by the cover, by $\pi^{\flat}$, and by $X_c$ and $S_c$.
\item (M3): the constant of Lemma~\ref{4.21:lem-grain-resummation}\,(ii), namely
\begin{equation*}
c_\pi=\Bigl(\prod K_A\,\vartheta(\square)\Bigr)\Bigl(\prod \frac{2}{r_A}\Bigr)
e^{-(\kappa_1-1)|\hat{\mathsf{y}}|},
\end{equation*}
with the two products taken over the same ranges as in that lemma.
\item (M4): clauses (c)--(e) of Proposition~\ref{4.21:prop-preparation-machine} and clause (iv) of
Lemma~\ref{4.21:lem-grain-resummation}. After these come the second-stage majorant steps of the
preparation site on $\hat{\mathfrak{E}}$, in the weights $\tilde{\nu}(\hat{e})$ (hypothesis (d)).
\end{itemize}

\emph{Site $(\mathrm{L})$.}
\begin{itemize}
\item (M1): the fundamental theorem of calculus in
$(\mathsf{s}^{(1)},\dots,\mathsf{s}^{(4)})$ on the product of four finite sets.
\item (M2): $\Pi(F)$ is the family $\mathfrak{Y}$ at each link, determined by
$(\sigma_0^{(t)},\text{hubs})$ and by $\mathcal{R}_{t-1}$.
\item (M3): $c_\pi=e^{-(\kappa_1-1)\sum_t|\mathsf{y}_t|}$. For the exterior pieces of deep units
this is multiplied by $C_{3F}e^{-\varsigma_3 d^{\wedge}}$, by the minimal form of the third-factor
lemma, \eqref{4.21:eq-third-factor-form-a}.
\item (M4): clauses (c) and (e) of Lemma~\ref{4.21:lem-source-anchored-sums}, applied at each of
the four links, together with Lemma~\ref{4.21:lem-probe-mass},
Lemma~\ref{4.21:lem-wrapped-mass} and Lemma~\ref{4.21:lem-deep-unit-sum}.
\end{itemize}

\emph{Birth of new boundary terms.}
\begin{itemize}
\item (M1): sums, products and the operations $\mathbf{T}'_k(X_j)$ over interior variables.
\item (M2): $\Pi(F)$ consists of the classes $\mathbf{D}_k^{Z}$ and $\mathbf{D}_k^{\sharp}$ and the
skeleton map $\operatorname{sk}(\cdot)$. These are determined by $(\pi_k,\mathsf{I},X)$.
\item (M3): $c_\pi=1$, or $c_\pi=\mathbf{T}'1$ where the operation $\mathbf{T}'_k(X_j)$ enters.
\item (M4): the strip theorem and clauses (u3)--(u7) of the combination lemma, in the form carried
by Proposition~\ref{4.21:prop-birth-site}\,(a), and clause (c) of the skeleton lemma, through which
the skeleton label $\operatorname{sk}$ enters Proposition~\ref{4.21:prop-birth-site}\,(c).
\end{itemize}
\end{proposition}

\begin{proof}
We treat the four components of \eqref{4.21:eq-majorant-structure-a} in turn, for all four sites
at once.

\emph{The identities} (M1). At the four sites these are, respectively:
\begin{itemize}
\item Lemma~\ref{4.21:lem-first-stage-expansion}\,(i);
\item the first-stage identity of the preparation site (M\"obius differences in $(t,u)$, then the
fundamental theorem of calculus in $\mathsf{s}$), together with
Lemma~\ref{4.21:lem-grain-resummation}\,(i);
\item clause (c) of Proposition~\ref{4.21:prop-large-field-site};
\item clause (a) of Proposition~\ref{4.21:prop-birth-site}.
\end{itemize}
Each of these identities holds for every function of $\tau$ that has the locality of
$F\circ\mathsf{R}$. That locality is the one assumed in
Lemma~\ref{4.21:lem-first-stage-expansion} and, at the preparation site, the one defining the
class of Lemma~\ref{4.21:lem-grain-resummation}. The identities therefore use $F$ only through the
function $\tau\mapsto F^{\sharp}\circ\mathsf{R}(\tau)$.

The same observation gives linearity. Write $c_{A,Q}$ for the pieces of an element $c$ at the
preparation site, and $\mathfrak{p}$ for a piece of the decoupling expansion. The maps
\begin{itemize}
\item $F\mapsto F_{\mathsf{y}}$ at $(\mathrm{T})$,
\item $c\mapsto c_{A,Q}\mapsto v_{\hat{e}}$ at $(\mathrm{P})$, and
\item $v\mapsto\mathfrak{p}$ at $(\mathrm{L})$
\end{itemize}
are compositions of evaluations, integrals and derivatives in $\tau$. Hence they are linear in the
stored term.

\emph{The index sets} (M2). At the four sites these are given, respectively, by:
\begin{itemize}
\item the definition of $\mathfrak{Y}(F)$ recalled in
Lemma~\ref{4.21:lem-first-stage-expansion}\,(i);
\item Proposition~\ref{4.21:prop-preparation-machine};
\item clause (c) of Proposition~\ref{4.21:prop-large-field-site};
\item clause (c) of the skeleton lemma (the skeleton label $\operatorname{sk}$ of
Proposition~\ref{4.21:prop-birth-site}\,(c)).
\end{itemize}
Each of these definitions involves only sets of the scaffold and the reading set $\mathcal{R}$.
The reading set of a stored term is fixed by the clause of the format of stored terms that
defines it.

\emph{The constants} (M3). Each $c_\pi$ comes from Cauchy's formula applied to a function bounded
by its supremum. Here the function is $\tau\mapsto F^{\sharp}\circ\mathsf{R}(\tau)$, or a piece
obtained from it by the operations in (M1).

\emph{The counts} (M4). The counting statements listed there involve no field variable:
\begin{itemize}
\item Lemma~\ref{4.21:lem-source-anchored-sums};
\item Lemma~\ref{4.21:lem-probe-mass};
\item Lemma~\ref{4.21:lem-wrapped-cube-bound};
\item Lemma~\ref{4.21:lem-wrapped-mass};
\item Lemma~\ref{4.21:lem-deep-unit-sum};
\item clause (iv) of Lemma~\ref{4.21:lem-grain-resummation}.
\end{itemize}
The remaining second-stage steps are covered by the hypotheses.
\begin{itemize}
\item At $(\mathrm{T})$, hypothesis (a) states that $V$ enters only as Mayer factors and that no
$V$ occurs after the passage to the polymer representation. The majorants are therefore series in
the activities $\nu(Y)$ alone.
\item At $(\mathrm{P})$, hypothesis (b) states that the second stage reads elements only through
the per-cube quantities $(\mathsf{b},\mathsf{d},\nu)$, and hypothesis (d) states that its steps on
$\hat{\mathfrak{E}}$ are majorant series in the weights $\tilde{\nu}(\hat{e})$.
\item At the birth site, hypothesis (c) states that the strip theorem holds for any family of
numbers, or of functions in the supremum norm.
\end{itemize}
In each case the second stage is therefore a majorant series in the suprema supplied by (M3).
\end{proof}

\begin{remark}[Grain of an excluded component at the running level]\label{4.21:rem-excluded-grain}
Fix a stage of the $\mathbf{R}$-operation, let $j$ be its running level, and let $C'$ be an excluded
component of this stage (a sibling, in the sense of the definition of siblings of a component), with
level $j(C')$. Let $Z$ be the large-field region and $\Omega_{j+1}$ the small-field region at level
$j+1$, write $\Omega^c_{j+1}$ for the complement of $\Omega_{j+1}$, and let $\pi_{j+1}$ be the
partition into cubes of level $j+1$. Refer to the lemma on the grain partition of a component, with
its parts (1) and (2) and its alternatives (b) and (c) for the grain; a \emph{member} of a grain is
one of its elements. Then the following hold.
\begin{enumerate}
\item The grain of $C'$ is of the form (b) or (c) of that lemma, taken at the running level $j$ of the
stage. The strata and the top zone of $C'$ are then members, or unions of members, of this grain.
\item In the proof of part (1) of that lemma, a cell $\Delta\subset\Omega^c_{j+1}$ satisfies
\begin{equation}\label{4.21:eq-excluded-grain-a}
\begin{gathered}
\operatorname{lev}(\Delta)\le j,\quad\text{or}\\
\operatorname{lev}(\Delta)=j+1
\ \text{and}\ \Delta\ \text{is a}\ \pi_{j+1}\text{-cube of}\ Z\cap\Omega^c_{j+1},
\end{gathered}
\end{equation}
the second case occurring only on the band of an exit of the current stage with $j(C')=j+1$, where
$\Delta$ is a member of a grain of the form (b). Each member of the grain is a union of cells.
\item In the proof of part (2) of that lemma, the preparation re-labels in $C'$ only levels at most
$j(C')\le j+1$. Every member $G$ of a grain of the form (b) satisfies
\begin{equation}\label{4.21:eq-excluded-grain-1}
\operatorname{lev}(G)=j+1 .
\end{equation}
\end{enumerate}
No other step of the grain lemma or of the resummation lemma at the preparation site uses the level
of a cell inside a member of a grain of the form (b).
\end{remark}

\begin{proof}
For the first assertion: part (2) of the grain lemma concludes $\operatorname{lev}(G)\ge j+1$ for
the members $G$ of the grain. If the grain of $C'$ were formed at a level other than the running
level $j$ of the stage, this conclusion would fail. The grain is therefore taken of the form (b) or
(c) at the running $j$, and with this choice the strata and the top zone of $C'$ are members or
unions of members.

For the second assertion: the proof of part (1) uses that the cells in $\Omega^c_{j+1}$ have level at
most $j$. This holds except on the band of an exit of the current stage. There the sibling relation
gives $j(C')=j+1$, and a cell of level $j+1$ on the band is a $\pi_{j+1}$-cube of
$Z\cap\Omega^c_{j+1}$, that is, a member of a grain of the form (b). This is the alternative recorded
in \eqref{4.21:eq-excluded-grain-a}. In the first case the argument of part (1) is unchanged. In the
second case the member is itself a cell, hence a union of cells. In both cases each member is a union
of cells, which is the conclusion part (1) requires.

For the third assertion: the proof of part (2) uses that the preparation has re-labelled only the
levels at most $j$. Inside $C'$ it has re-labelled only the levels at most $j(C')$, and
$j(C')\le j+1$. Let $G$ be a member of a grain of the form (b). If the cells of $G$ fall under the
first alternative of \eqref{4.21:eq-excluded-grain-a}, the argument of part (2) gives
\eqref{4.21:eq-excluded-grain-1} as before. If $G$ is a band cell under the second alternative, then
$\operatorname{lev}(G)=j+1$ directly. So \eqref{4.21:eq-excluded-grain-1} holds in either case.

Finally, the grain lemma and the resummation lemma at the preparation site use the level of a cell
inside a member of a grain of the form (b) only in the two places just treated. Their remaining steps
are therefore unaffected.
\end{proof}

\begin{proposition}[The link map stays in the stored domain]\label{4.21:prop-link-containment}
Let $F$ be a stored boundary-type term with $\mathcal{P}_F\neq\emptyset$, put
$X_F=A_F\cup\bigcup_{\mathcal{P}_F}Y^{\sharp}$, and let the links, the link maps
$\mathsf{R}(\mathsf{s})$, the decoupling families $\sigma_0$ and the polydiscs, source polydiscs and
output tubes of the links be those of the decoupling clause
(Definition~\ref{4.21:def-decoupling-clause}, display~\eqref{4.21:eq-decoupling-clause-a}).
Here $\mathsf{R}(\mathsf{s})$ denotes the link map on the lower end
$(\mathbf{K},\mathbf{J}_{\mathbf{K}})$ of the link, $\mathsf{R}(\mathsf{s})^{h}$ the link map of the
link, and $\restriction X_F$ the restriction to $X_F$.
Let $\mathfrak{U}_F(X_F)\subset\mathfrak{U}^{\mathrm{tr}}_F(X_F)$ be the stored domain of $F$
(Notation~\ref{4.21:not-orbit-data}, display~\eqref{4.21:eq-orbit-data-a}). It is cut out of
$\mathfrak{U}^{\mathrm{tr}}_F(X_F)$ by the inequalities
$Q_a(y)<\theta_a(y)-\varpi_{\mathrm{fr}}r_a(y)$ between the entries $Q_a$, the thresholds
$\theta_a$ and the margins $r_a$ of \eqref{4.21:eq-orbit-data-a}, imposed at the clause locations
$y$ whose stencil meets some $\mathsf{B}_W$, for $a\notin\{5,8\}$, where
$\varpi_{\mathrm{fr}}=\tfrac18(1-\theta_S)\beta$, with $\beta$ the schedule constant of the layered
analyticity spaces (Definition~\ref{4.4:def-post-step-space}) and $\theta_S$ the constant entering
the reserve of the stored domain (Notation~\ref{4.21:not-orbit-data}). For a member
$(Y,m)$ of $\mathcal{P}_F$ we write $W=Y$
for the frame it defines, so that $k_W=m$. Assume the following.
\begin{enumerate}
\item[(a)] \emph{Containment at Ba{\l}aban's decoupling family.} At the sites (T), (P) and (L) the
containment statements assumed in Lemma~\ref{4.21:lem-wrapped-holomorphy} hold, taken at
Ba{\l}aban's decoupling family (the family in place of which
Definition~\ref{4.21:def-decoupling-clause} puts its families $\sigma_0$). In particular
\eqref{4.21:eq-link-containment-1} below holds at that family.
\item[(b)] \emph{Arithmetic of the stored-domain thresholds.} The following arithmetic, which is the
threshold arithmetic of the stored-domain construction, holds. Let $y$ be a clause location whose
stencil meets $\mathsf{B}_W$, and let $a\notin\{5,8\}$.
\begin{itemize}
\item[(b1)] At the first site (T) after the creation of $W$ (layer $n=m$), the factor $1$ of the
threshold exceeds the level $F_n=f_{k+1,k}+\theta_S\beta/2$ (a number; the subscripted letter is
unrelated to the term $F$) by $\tfrac12(1-\theta_S)\beta$, and the increment
$\pi_n$ at layer $n$ satisfies
\begin{equation}\label{4.21:eq-link-containment-2}
\begin{split}
\frac{\pi_n}{\alpha_{0,n}}
&\le \frac{C_H(1+\beta_0)^2\|C\|\,(1+\eta_p|\mathsf{t}|)\,g\sup|\mathsf{B}|}{\alpha_{0,k}}\\
&\le \frac{C_H(1+\beta_0)^2\|C\|}{K_R}
\le \tfrac14(1-\theta_S)\beta ;
\end{split}
\end{equation}
this is the no-attenuation case $n=k$ of clause (ii) of the layer-increment bound, together with
the homogeneity lemma at the site (T). Moreover, on the polydisc of the link the hypothesis
$K_Rg\sup|\mathsf{B}|\le\alpha_{*}$ of the layer-increment bound holds, that is, the condition
$\theta^{\mathbb{C}}\le\tfrac12$ on the quantity $\theta^{\mathbb{C}}$ of that bound.
\item[(b2)] At every later site (T) the excess is $\beta[1-(1+\theta_S)2^{-\ell'}]$ for an integer
$\ell'\ge2$, the exponent of the level over which the excess is measured at that site (at the first
site (T) this exponent is $1$, which gives the excess $\tfrac12(1-\theta_S)\beta$ of (b1)); the
excess is at least $\beta/2$, and the increment there is attenuated.
\item[(b3)] At the site (P) the margin $r_T$ satisfies
\begin{equation}\label{4.21:eq-link-containment-3}
r_T\ \ge\ \frac{\theta_S\beta}{2}
\ \ge\ \beta\,\theta_S\,2^{-(k-k_W)},
\end{equation}
and the lower bound $r_{\min}(y)\ge\beta e_{*}(m)$ of the margin bound at the site (P) holds. At
the site (L), in an untreated component (a component not yet treated at the site (L), in the sense
of the margin proposition at the site (L)), the lower bound $\min(\theta_S,\tfrac12)\beta$ on the
margins given by the margin proposition at the site (L) holds.
\item[(b4)] On the layer $m$, the threshold factor of the domain created at a site (T) of level
$i+1>m$ is $f_{i+1,m}+\theta_S\beta2^{-(i+1-m)}$. Moreover
\begin{equation}\label{4.21:eq-link-containment-4}
f_{i+1,m}+\theta_S\beta2^{-(i+1-m)}=1-\beta+\beta(1+\theta_S)2^{-(i+1-m)} .
\end{equation}
\end{itemize}
\item[(c)] \emph{The dilation lemma accompanying the third-factor lemma, with its remark on the
stored domain.} On the auxiliary disc in the parameter $\lambda$ of the Schwarz-lemma step of the
third-factor lemma (whose minimal form is Definition~\ref{4.21:def-third-factor-form}), each entry
$Q_a(y)$, evaluated at the configuration--current pair dilated in $\lambda$ on that disc (the
dilated pair) and restricted to $X_F$, is at most $\tfrac3{32}+\tfrac12$ times its threshold.
Moreover $1-\tfrac1{80}\le1-\varpi_{\mathrm{fr}}$.
\item[(d)] \emph{Tube containment for descendants.} The following is the descendant tube
containment of the stored-domain construction. Let a term be built from a boundary term
$\mathbf{B}'$ whose domain contains $Y$. Then its tube over the strip lies in
$\mathfrak{U}^{\mathrm{tr}}_{\mathbf{B}'}$, and the threshold factors of the domains created at the
sites (P) and (L) are at most those of the domains created at the sites (T).
\end{enumerate}
Here $\alpha_{0,n}$ is the radius of Lemma~\ref{4.4:lem-twist-stability}. In (b1), $C_H$ and
$\beta_0$ are constants of the layer-increment bound and of the homogeneity lemma at the site (T);
$K_R$ is the denominator of that bound, subject to its lower bound (d2); $\|C\|$ and
$g\sup|\mathsf{B}|$, with $\sup|\mathsf{B}|$ the supremum of the field entering that bound (not the
set $\mathsf{B}_W$), are the two factors in whose product that bound is linear, $g$ being the scalar
factor of that bound (a number, not a group element as in Notation~\ref{4.21:not-orbit-data});
$\mathsf{t}$ is the
interpolation time and $\eta_p$ the constant multiplying $|\mathsf{t}|$; $\alpha_{*}$ is the
threshold in the hypothesis $K_Rg\sup|\mathsf{B}|\le\alpha_{*}$ of the layer-increment bound. In
(b3), $r_T$ is the margin left below $1-\varpi_{\mathrm{fr}}$ at the site (P), $r_{\min}(y)$ is the
least margin at $y$ and $e_{*}(m)$ the factor in its
lower bound. In (b4), $f_{i+1,m}$ is the part of the threshold factor on the layer $m$ of a domain
created at level $i+1$ that does not carry the factor $\theta_S$.

Then at every link of the decoupling clause, on the whole of its polydisc, its source polydisc and
its output tube,
\begin{equation}\label{4.21:eq-link-containment-1}
\bigl(\mathsf{R}(\mathsf{s})^{h}\bigr)\restriction X_F\in\mathfrak{U}_F(X_F).
\end{equation}
\end{proposition}

\begin{proof}
The membership \eqref{4.21:eq-link-containment-1} at Ba{\l}aban's decoupling family is hypothesis
(a). It is used as it stands and is not proved again. We show that it persists for the families
$\sigma_0$ of the decoupling clause.

\emph{Clause locations away from the sets $\mathsf{B}_W$.} By the definition of the stored domain in
\eqref{4.21:eq-orbit-data-a}, $\mathfrak{U}_F$ is obtained from $\mathfrak{U}^{\mathrm{tr}}_F$ by
adding inequalities only at the clause locations whose stencil meets some $\mathsf{B}_W$. At every
other clause location the two domains therefore impose the same condition. There the membership
is that of Lemma~\ref{4.21:lem-wrapped-holomorphy} (holomorphy of wrapped terms at every link).

From now on let $y$ be a clause location whose stencil meets $\mathsf{B}_W$. Here $W=Y$ for a member
$(Y,m)$ of $\mathcal{P}_F$, and every such member is a frame with $k_W=m$
(Notation~\ref{4.21:not-orbit-data}). Let $a\notin\{5,8\}$. We argue site by site.

\emph{The first site (T) after the creation of $W$ (layer $n=m$).}
By (b1) the threshold factor $1$ exceeds $F_n$ by $\tfrac12(1-\theta_S)\beta$. Under the clause (T)
of the decoupling clause, in \eqref{4.21:eq-decoupling-clause-a}, the increment is that of the full
layer $\mathcal{Y}(\mathsf{s};b)$ of Lemma~\ref{4.21:lem-exempted-set}, taken at amplitude one. We
check that \eqref{4.21:eq-link-containment-2} remains valid for it.

The arithmetic of clause (ii) of the layer-increment bound is linear in $\|C\|\,g\sup|\mathsf{B}|$;
this is entry (d2) of the homogeneity lemma at the site (T). Its constant $C_H$, read on the polydisc
$|\sigma|\le e^{\kappa_1}$, becomes $C_He^{16\kappa_1}$. By
Lemma~\ref{4.21:lem-exempted-set} (decoupled operators are blind to the exempted set), this holds
for every $\sigma_0$, with the same constants. The number $C_He^{16\kappa_1}$ is the one for which
the lower bound (d2) on $K_R$ holds. The factor $(1+\eta_p|\mathsf{t}|)$ is replaced by $1$. The
hypothesis $K_Rg\sup|\mathsf{B}|\le\alpha_{*}$ of the layer-increment bound, which is the condition
$\theta^{\mathbb{C}}\le\tfrac12$ on the quantity $\theta^{\mathbb{C}}$ of that bound, holds by
(b1).

The chain \eqref{4.21:eq-link-containment-2} therefore gives
$\pi_n/\alpha_{0,n}\le\tfrac14(1-\theta_S)\beta$ for the increment of the full layer as well. The
increment and the reserve together use at most
\begin{equation*}
\tfrac14(1-\theta_S)\beta+\varpi_{\mathrm{fr}}
=\bigl(\tfrac14+\tfrac18\bigr)(1-\theta_S)\beta\le\tfrac12(1-\theta_S)\beta .
\end{equation*}
The right-hand side is the excess of the threshold factor $1$ over $F_n$ given by (b1). Hence the
reserve
$\varpi_{\mathrm{fr}}=\tfrac18(1-\theta_S)\beta$ is not touched, and the inequality defining
$\mathfrak{U}_F$ holds at $y$.

\emph{Later sites (T).} By (b2) the excess is $\beta[1-(1+\theta_S)2^{-\ell'}]$ with $\ell'\ge2$, and
it is at least $\beta/2$. The increment is attenuated, hence smaller
than at the first site. The same count as above applies.

\emph{The site (P).} The clause (P) of the decoupling clause, in
\eqref{4.21:eq-decoupling-clause-a}, changes no family of cells. The restriction step at the site
(P) restricts one family. The computation of (b3) therefore applies verbatim.
Inequality \eqref{4.21:eq-link-containment-3} holds, and the lower bound
$r_{\min}(y)\ge\beta e_{*}(m)$ of the margin bound at the site (P) stands.

\emph{The site (L), $W$ live, in an untreated component.} The same computation of (b3) applies, and
the lower bound $\min(\theta_S,\tfrac12)\beta$ on the margins, given by the margin proposition at
the site (L), stands. The increments at this site are the three increments of part ($\beta$) of the
increment statement at (L), the first of which carries the factor $2^{-(k-1-m)}$.
On the families of the decoupling clause they are bounded by clause (c3) of
Lemma~\ref{4.21:lem-wrapped-holomorphy}.

\emph{The site (L), $W$ dying.} Here $W$ dies at the site (L), in the sense of the entrywise bound
for dying frames at (L). With $Z_h$ denoting the region of that bound which contains $W$,
\begin{equation*}
\mathsf{B}_W\subset W\subset Z_h\subset\Lambda^{\natural}.
\end{equation*}
By the entrywise bound for dying frames
at (L), each entry is at most $\tfrac1{16}$ of its threshold on
$X\cap(\Lambda^{\natural}\cup\mathfrak{N}^{\mathrm{rec}})$, with $X$ the domain on which that bound
is stated. This holds for every $\sigma_0$, by
clause (c3) of Lemma~\ref{4.21:lem-wrapped-holomorphy}. On $W$ the threshold and the margin
coincide, $\theta_a=r_a$, by the frame-threshold identity of the stored-domain construction, and
$\tfrac1{16}<1-\varpi_{\mathrm{fr}}$. Hence, at a clause location $y$ in $W$,
\begin{equation*}
Q_a(y)\le\tfrac1{16}\theta_a(y)<(1-\varpi_{\mathrm{fr}})\theta_a(y)
=\theta_a(y)-\varpi_{\mathrm{fr}}r_a(y).
\end{equation*}

\emph{The auxiliary $\lambda$-disc.} Consider the auxiliary disc in the parameter $\lambda$ used in
the Schwarz-lemma step of the third-factor lemma. By hypothesis (c), the dilation lemma accompanying
the third-factor lemma with its remark on the stored domain, each entry is at most
$\tfrac3{32}+\tfrac12$ of its threshold, and
\begin{equation*}
\tfrac3{32}+\tfrac12=\tfrac{19}{32}<1-\tfrac1{80}\le1-\varpi_{\mathrm{fr}} .
\end{equation*}
The dilated pair, restricted to $X_F$, therefore lies in $\mathfrak{U}_F(X_F)$.

\emph{Descendants.} A term built from a boundary term $\mathbf{B}'$ whose domain contains $Y$
carries the real data $\Phi_W$ of the frame through the chart variable of the orbit datum
$\mathbf{B}'^{\sharp}$ of $\mathbf{B}'$, by the second item of the remark on descendants of boundary
terms. The analyticity requirement (f1) of the format of stored
terms asks that the tube of this term over the strip lie inside the stored domain of $\mathbf{B}'$,
not only inside $\mathfrak{U}^{\mathrm{tr}}_{\mathbf{B}'}$. The latter containment is
hypothesis (d).

For the former, consider the layer $m$. By (b4), the domain created at a site (T) of level $i+1>m$
has threshold factor $f_{i+1,m}+\theta_S\beta2^{-(i+1-m)}$, which equals the right-hand side of
\eqref{4.21:eq-link-containment-4}. Since $i+1-m\ge1$, we have $2^{-(i+1-m)}\le\tfrac12$, and
therefore
\begin{equation*}
\begin{split}
1-\beta+\beta(1+\theta_S)2^{-(i+1-m)}
&\le1-\beta+\tfrac12\beta(1+\theta_S)\\
&=1-\tfrac12(1-\theta_S)\beta\\
&\le1-\varpi_{\mathrm{fr}} ,
\end{split}
\end{equation*}
the last step because $\tfrac18\le\tfrac12$ and $(1-\theta_S)\beta\ge0$, the latter because by (b1)
$\tfrac12(1-\theta_S)\beta$ is the amount by which $1$ exceeds $F_n$. By the last part of
hypothesis (d), the threshold
factors of the domains created at the sites (P) and (L) are at most the factors of the domains
created at the sites (T), hence also at most $1-\varpi_{\mathrm{fr}}$. This proves
\eqref{4.21:eq-link-containment-1} at every link.
\end{proof}

\begin{remark}[One class chain per step in the per-stage bound]\label{4.21:rem-one-chain-per-step}
Let $C_0^{\sharp}$ be the constant of the per-stage bound on the format of stored terms, and let
$j\le i$ be the endpoints of the interval of steps at which the chains of that bound were run. Fix
a step $t$ with $j\le t\le i$ and a stage $n$ of the $\mathbf{R}$-operation, and let $\mathbb{C}$
denote the class chain of step $t$ at stage $n$. Its members are the chain terms of all components
of the step's class, of all filing levels. In particular, the terms of every attempt of that step
and stage, arrested or not, are members of $\mathbb{C}$. For a localisation domain $X$ let
$\mathbb{C}_X$ be the part of $\mathbb{C}$ localised in $X$. We recall the norm of the class chain
at stage $n$:
\begin{equation}\label{4.21:eq-one-chain-per-step-a}
  \|\mathbb{C}\|_{(n)} = \sup_x \sum_{X\ni x}
  \bigl(\text{weighted size of } \mathbb{C}_X \text{ at stage } n\bigr),
\end{equation}
one sum over all filing levels and over all components of the step's class. For a component $D$
of the class write $\mathbb{C}\restriction D$ for the part of $\mathbb{C}$ formed by the terms of
$D$.

In the per-stage bound one chain is counted per step of the interval from $j$ to $i$. For every
component $D$ of the class one has
\begin{equation}\label{4.21:eq-one-chain-per-step-1}
  \|\mathbb{C}\restriction D\|_{(n)} \le \|\mathbb{C}\|_{(n)} ,
\end{equation}
and the per-stage bound holds as stated. The constants $c_{\mathrm{att}}$ and $c_{\mathrm{sib}}$
of the counting lemma for attempts and siblings are not needed for it. If the per-stage bound is
applied component by component, then the per-cube constant $\mathsf{b}_i$ of the per-stage bound,
whose index is the endpoint $i$ above, is to be replaced by the per-cube constant $\mathsf{b}^{P}$
of the format class. This replacement differs from $\mathsf{b}_i$ by a number at the top levels
and not in degree, and the constant $C_\Sigma$ of the base comparison changes at most by that
number.
\end{remark}

\begin{proof}
The norm \eqref{4.21:eq-one-chain-per-step-a} is the supremum over $x$ of a single sum. This sum
runs over all domains $X\ni x$ on which a term of $\mathbb{C}$ is localised, whatever the filing
level of that term and whatever the component of the step's class to which it belongs.

Fix a component $D$. The terms of $\mathbb{C}\restriction D$ are among the terms of $\mathbb{C}$.
This holds also when $D$ comes from an arrested attempt, since the terms of every attempt of the
given step and stage, arrested or not, are members of the one class chain $\mathbb{C}$. Hence, for
each $x$, the sum defining $\|\mathbb{C}\restriction D\|_{(n)}$ at $x$ is a sub-sum of the sum
defining $\|\mathbb{C}\|_{(n)}$ at $x$. All summands are non-negative weighted sizes, so the
sub-sum is at most the full sum. Taking the supremum over $x$ gives
\eqref{4.21:eq-one-chain-per-step-1}.

The per-stage bound estimates the chain terms by $C_0^{\sharp}$ per stage and per chain, the chains
having been run at the steps of the interval from $j$ to $i$. The chain that it counts at a step of
this interval is the class chain of that step. One chain is therefore counted per step of the
interval. The bound by $C_0^{\sharp}$ per stage and per
chain applies to it with the norm \eqref{4.21:eq-one-chain-per-step-a}. By
\eqref{4.21:eq-one-chain-per-step-1}, the same bound holds for each component separately. No count
of attempts or of siblings enters. The per-stage bound thus holds as stated, and the constants
$c_{\mathrm{att}}$, $c_{\mathrm{sib}}$ are not used in it. What they permit is to estimate the norm
component by component.

Suppose instead that the per-stage bound is applied component by component. Then the per-cube
constant $\mathsf{b}_i$ is to be replaced by $\mathsf{b}^{P}$. The latter exceeds $\mathsf{b}_i$ by a
number at the top levels and has the same degree. The base comparison uses the same per-cube
constants, namely $\mathsf{b}_k$ at level $k$, and its constant $C_\Sigma$ is a maximum involving
them. Retaking that maximum with $\mathsf{b}^{P}$ in place of $\mathsf{b}_k$ changes $C_\Sigma$ at
most by that number.
\end{proof}

\begin{proposition}[Base-point profile of the carried gauge map]\label{4.21:prop-base-point-profile}
For $u\in G^{\mathbb{C}}$, regarded as an invertible complex $N\times N$ matrix with the operator
norm $\|\cdot\|$, we write
$\operatorname{pol}(u):=\log\max\{\|u\|,\|u^{-1}\|\}$ for its size as a gauge transformation.
The sites (T), (P), (L), the decoupling families $\sigma_0$, the decoupling parameters
$\mathsf{s}$ and the exempted set $\mathsf{K}$ are those of the decoupling clause
(Definition~\ref{4.21:def-decoupling-clause}). The layers $\mathcal{Y}(\mathsf{s};b)$ at a
source datum $b$ are those of the links (Definition~\ref{4.21:def-links-readers}). Assume the
following statements.
\begin{enumerate}
\item[(a)] Parts (i) and (iv) of the lemma on the block-local gauge map $\mathfrak{u}$ of the
relative small-field machine. Part (i) states that, at every point $z$ of the unit's domain,
\begin{equation*}
|\log\mathfrak{u}(z)|\le K_u\,\Theta_0^{\mathrm{loc}}(\mathcal{K})(z),
\qquad
\Theta_0^{\mathrm{loc}}(\mathcal{K})(z)\le\sup_{B^{\mathrm{top}}(z)}N_y(\mathcal{K}),
\end{equation*}
where $\mathcal{K}$ is the layer read by $\mathfrak{u}$, $N_y(\mathcal{K})$ is the local size
of the layer $\mathcal{K}$ at the point $y$, as in that lemma, and $B^{\mathrm{top}}(z)$ is the
top block of $z$, which lies in the $M$-cube of the local scale of $z$, inside the unit's domain;
it states moreover that $\mathfrak{u}(z)$ reads the layer on $B^{\mathrm{top}}(z)$ only, and that
$\mathfrak{u}\equiv1$ on top blocks on which the layer vanishes.
Part (iv) states the base-point profile at the sites (T), (P), (L) for
Ba{\l}aban's choice of decoupling family, with the constant $B_H$ and the rate $\delta_1$.
\item[(b)] Part (iii) of the lemma on the carried gauge map $h\,\mathfrak{Q}^{\mathfrak{b}}$ at
site (L): the factor $h$ equals $1$ at law sites off $Z$.
\item[(c)] The localisation bound of the relative small-field machine at site (L), for
Ba{\l}aban's decoupling family, together with its proof, in which the decoupling family enters at
exactly three places: the polydisc condition of that bound; the decay rate of the kernel of the
operator $\mathbf{G}''(\sigma)$ of that bound, which is given by part (a) of the weighted
localisation statement; and one step, which uses $\sigma_0\cap\Lambda=\emptyset$.
\item[(d)] At site (T) and interpolation time $\mathsf{t}=0$: the source datum $b$ vanishes off
the fluctuation domain $\Omega$ of the step, by the containment statement at site (T) of the
correlation theorem; and
\begin{equation}\label{4.21:eq-base-point-profile-1}
\sup|b|\le C_1\,g\,\sup|\mathsf{B}|\le\frac{C_1}{K_R}\,\alpha_*,
\end{equation}
where the first inequality is the bound \cite[(1.19)]{B-CMP116}, with $g$ the coupling constant
of the step at which that bound is applied (the letter of \cite{B-CMP116}, not the reference
coupling), $C_1$ the constant of that bound and $\mathsf{B}$ the variable whose supremum controls
the source datum $b$ in it; the second inequality is the condition
$g\sup|\mathsf{B}|\le\alpha_*/K_R$ that the correlation theorem imposes at site (T), the radius
$\alpha_*$ and the constant $K_R$ being those fixed in that condition.
\end{enumerate}
Put
\begin{equation}\label{4.21:eq-base-point-profile-a}
B_H^S:=\max\Bigl\{B_H,\ \frac{B_\sharp C_1}{K_R},\ B_\sharp(1+K_\sharp)\Bigr\},
\end{equation}
a number depending only on $(d,L,M,\kappa_1)$ and fixed before $\gamma$. Then, on the families on
which the decoupling clause runs, at every law site $s$ of the chart of a wrapped unit and for
every $\mathsf{s}$ with $|\mathsf{s}|\le e^{\kappa_1}$, the following hold.
\begin{enumerate}
\item[(T)] With $\Omega$ the fluctuation domain of the step,
\begin{equation}\label{4.21:eq-base-point-profile-2}
\begin{aligned}
\operatorname{pol}\bigl(\mathsf{u}(\mathsf{s})(s)\bigr)
&\le K_uB_\sharp\sup|b|\;e^{-\delta_P d^{\wedge}(B^{\mathrm{top}}(s),\Omega)}\\
&\le K_uB_H^S\,\alpha_*\;e^{-\delta_P d^{\wedge}(B^{\mathrm{top}}(s),\Omega)}.
\end{aligned}
\end{equation}
\item[(P)] The bound of part (iv) of the lemma in (a) at site (P) holds unchanged.
\item[(L)] ($\alpha$) If $s\notin Z$, then
\begin{equation}\label{4.21:eq-base-point-profile-3}
\begin{aligned}
\operatorname{pol}\bigl(h\,\mathfrak{Q}^{\mathfrak{b}}(\vec{\mathsf{s}})(s)\bigr)
\le K_uB_\sharp\Bigl[&\sum_{s'}\sup|\mathrm{arg}_{s'}|
+K_\sharp\sup\ell^3|\mathsf{J}_3|\Bigr]\\
&\times e^{-\delta_P d^{\wedge}(B^{\mathrm{top}}(s),\Lambda^{\sim})},
\end{aligned}
\end{equation}
where the sum runs over the layers $\mathbb{H}_{s'}$ of the carried map, $\mathrm{arg}_{s'}$
denotes the arguments of $\mathbb{H}_{s'}$ supported in $\Lambda^{\sim}$, $\mathsf{J}_3$ is
the supported source of the field datum $\mathfrak{j}_3$, and $\ell$ is the length scale in the
bound of part (iii)$'$ of Lemma~\ref{4.21:lem-exempted-set}.
($\beta$) If $s\in Z_h$, then the display of the localisation bound in (c) holds unchanged:
for $s'=3,4$,
\begin{equation}\label{4.21:eq-base-point-profile-4}
K_u\,\Theta_0^{\mathrm{loc}}(\mathbb{H}_{s'})(\tilde{s})
\le C_L\,\varepsilon_{\mathfrak{Q}}(k)\,e^{-\delta_L(d_{\mathrm{sc}}(s,\mathsf{S})-1)},
\end{equation}
with $Z$, $Z_h$, $\tilde{s}$ and $\mathsf{S}$ as in the lemma in (b) and in the localisation
bound in (c).
\end{enumerate}
Moreover, the bounds (T) and (L)($\alpha$) have the form of part (iv) of the lemma in (a) with
the pair $(B_H,\delta_1)$ replaced by
$(B_H^S,\min\{\delta_1,\delta_Pw\})$. Part (iii) of the lemma in (b), and the site-(L) theorem
resting on it, depend on this pair only through
\begin{equation}\label{4.21:eq-base-point-profile-5}
C_\natural=2e^2\max\{C_LK_{\mathfrak{Q}},\,K_uB_H,\,K_u\},
\qquad
\delta_\natural=\min\{\delta_1,\delta_L,1\}\le 1 .
\end{equation}
With these replacements $\delta_\natural$ is unchanged, while $C_\natural$ and $\varepsilon_{\mathfrak{Q}}$
are recomputed once with $B_H$ replaced by $B_H^S$, within the upper bounds, taken in supremum
form, under which $C_\natural$ and $\varepsilon_{\mathfrak{Q}}$ are fixed.
\end{proposition}

\begin{proof}
We use two elementary properties of $\operatorname{pol}$. The first is
$\operatorname{pol}(u_1u_2)\le\operatorname{pol}(u_1)+\operatorname{pol}(u_2)$. It holds because
$\|u_1u_2\|\le\|u_1\|\,\|u_2\|$ and $\|(u_1u_2)^{-1}\|\le\|u_2^{-1}\|\,\|u_1^{-1}\|$, so the
maximum for $u_1u_2$ is at most the product of the maxima. The second is
$\operatorname{pol}(e^{\Xi})\le\|\Xi\|$ for every complex $N\times N$ matrix $\Xi$. It holds
because $\|e^{\pm\Xi}\|\le e^{\|\Xi\|}$.

\emph{Site} (T). The map $\mathsf{u}(\mathsf{s})=\mathfrak{u}[\cdot;\cdot]$ is the block-local
gauge map of the lemma in (a), read at the layer $\mathcal{Y}(\mathsf{s};b)$, so that
$\mathsf{u}(\mathsf{s})(s)=e^{\log\mathfrak{u}(s)}$. The second property of $\operatorname{pol}$
and part (i) of the lemma in (a) give, at the law site $s$,
\begin{equation*}
\operatorname{pol}\bigl(\mathsf{u}(\mathsf{s})(s)\bigr)
\le K_u\Theta_0^{\mathrm{loc}}(s)
\le K_u\sup_{B^{\mathrm{top}}(s)}N_y .
\end{equation*}
Here the layer is $\mathcal{Y}(\mathsf{s};b)$, read on $B^{\mathrm{top}}(s)$ only. By (d), its
source datum $b$ vanishes off $\Omega$. Part (iii) of Lemma~\ref{4.21:lem-exempted-set} holds with
the same constants whichever cells are exempted. Applied to this layer, it bounds
$\sup_{B^{\mathrm{top}}(s)}N_y$ by
$B_\sharp\sup|b|\,e^{-\delta_P d^{\wedge}(B^{\mathrm{top}}(s),\Omega)}$. This gives the first
inequality of \eqref{4.21:eq-base-point-profile-2}. By \eqref{4.21:eq-base-point-profile-1},
$\sup|b|\le(C_1/K_R)\alpha_*$. By \eqref{4.21:eq-base-point-profile-a},
$B_\sharp C_1/K_R\le B_H^S$, which gives the second inequality. At site (T) the exempted set is
empty, $\mathsf{K}=\emptyset$, so at $\mathsf{s}=0$ no cell carries the layer, and
$\mathfrak{u}\equiv1$ by part (i) of the lemma in (a); thus the base point of the profile is
$\mathsf{u}(0)=1$.

\emph{Site} (P). At site (P) the decoupling clause does not change the decoupling family.
Part (iv) of the lemma in (a) at site (P) therefore applies as it stands.

\emph{Site} (L), \emph{case} ($\alpha$). Let $s\notin Z$. By (b), $h=1$ there, so
$\operatorname{pol}(h\,\mathfrak{Q}^{\mathfrak{b}}(\vec{\mathsf{s}})(s))
=\operatorname{pol}(\mathfrak{Q}^{\mathfrak{b}}(\vec{\mathsf{s}})(s))$. The carried map at
site (L) is built layer by layer: with $\mathfrak{b}_{s'}$ the base point after the layer
$\mathbb{H}_{s'}$ and $\mathfrak{w}_{s'}$ the gauge factor contributed by that layer,
\begin{equation*}
\mathfrak{Q}^{\mathfrak{b}_{s'}}=\mathfrak{Q}^{\mathfrak{b}_{s'-1}}\,\mathfrak{w}_{s'},
\qquad
|\log\mathfrak{w}_{s'}|\le K_u\,\Theta_0^{\mathrm{loc}}(\mathbb{H}_{s'}),
\end{equation*}
the bound on $|\log\mathfrak{w}_{s'}|$ being part (i) of the lemma in (a) applied to the layer
$\mathbb{H}_{s'}$.
We apply the product inequality for $\operatorname{pol}$ to each factorisation in turn, and the
exponential inequality to each factor $\mathfrak{w}_{s'}$. This bounds
$\operatorname{pol}(\mathfrak{Q}^{\mathfrak{b}}(\vec{\mathsf{s}})(s))$ by the sum over $s'$ of
$K_u\Theta_0^{\mathrm{loc}}(\mathbb{H}_{s'})(s)$.

Each $\Theta_0^{\mathrm{loc}}(\mathbb{H}_{s'})(s)$ is then bounded through its arguments. The
arguments supported in $\Lambda^{\sim}$ are bounded by part (iii) of
Lemma~\ref{4.21:lem-exempted-set}. This yields the term
\begin{equation*}
B_\sharp\sup|\mathrm{arg}_{s'}|\,e^{-\delta_P d^{\wedge}(B^{\mathrm{top}}(s),\Lambda^{\sim})}.
\end{equation*}
The field datum $\mathfrak{j}_3$, whose supported source is $\mathsf{J}_3$, is bounded by
part (iii)$'$ of Lemma~\ref{4.21:lem-exempted-set}. This yields the term
\begin{equation*}
B_\sharp K_\sharp\sup\ell^3|\mathsf{J}_3|\,
e^{-\delta_P d^{\wedge}(B^{\mathrm{top}}(s),\Lambda^{\sim})}.
\end{equation*}
Adding the two contributions gives \eqref{4.21:eq-base-point-profile-3}.

\emph{Site} (L), \emph{case} ($\beta$). Let $s\in Z_h$. By (c), the proof of the localisation
bound uses the decoupling family only at the three places listed there, so it suffices to check
that each of them holds on the site-(L) families of the decoupling clause.
\begin{itemize}
\item Its polydisc condition is supplied by part (iii) of Lemma~\ref{4.21:lem-exempted-set}.
\item The decay rate of the kernel of $\mathbf{G}''(\sigma)$ is given by part (a) of the
weighted localisation statement. By Lemma~\ref{4.21:lem-exempted-set}, that statement holds with
the same constants whichever cells are exempted.
\item The condition $\sigma_0\cap\Lambda=\emptyset$ holds on the site-(L) families of the
decoupling clause, whose exempted set contains $\Lambda$
(Definition~\ref{4.21:def-decoupling-clause}, \eqref{4.21:eq-decoupling-clause-a}, and
Definition~\ref{4.21:def-links-readers}).
\end{itemize}
All three hold on the families of the decoupling clause. These differ from Ba{\l}aban's families
only by cells of $X$ lying off $\Lambda^{\sim}$, and only at the first operation. Hence the proof
of the localisation bound runs unchanged on these families, and the
display \eqref{4.21:eq-base-point-profile-4} holds unchanged for $s'=3,4$.

\emph{The constants after replacement.} The scaled distances of the relative small-field machine
satisfy $d^{\wedge}\ge w\,d_{\mathrm{sc}}$. Since $B^{\mathrm{top}}(s)$ lies in the $M$-cube of
the local scale of $s$ (part (i) of the lemma in (a)), the scaled distance from it to a set is at
least the scaled distance from $s$ minus one unit. Consequently,
\begin{equation*}
d^{\wedge}(B^{\mathrm{top}}(s),\cdot)\ge w\bigl(d_{\mathrm{sc}}(s,\cdot)-1\bigr).
\end{equation*}
Multiplying by $\delta_P$, the exponential factors
$e^{-\delta_P d^{\wedge}(B^{\mathrm{top}}(s),\cdot)}$ in
\eqref{4.21:eq-base-point-profile-2}--\eqref{4.21:eq-base-point-profile-3} are bounded by
\begin{equation*}
e^{-\min\{\delta_1,\delta_Pw\}(d_{\mathrm{sc}}(s,\cdot)-1)}
\end{equation*}
(when $d_{\mathrm{sc}}(s,\cdot)<1$ the left side is at most $1$ and the right side at least $1$).

The prefactors are handled as follows. At (T), the prefactor is at most $K_uB_H^S\alpha_*$ by
\eqref{4.21:eq-base-point-profile-2}. At (L) ($\alpha$), it carries $B_\sharp$ and
$B_\sharp K_\sharp$, and both are at most $B_\sharp(1+K_\sharp)\le B_H^S$ by
\eqref{4.21:eq-base-point-profile-a}. This is the replacement $(B_H,\delta_1)\mapsto
(B_H^S,\min\{\delta_1,\delta_Pw\})$.

By the level condition on $w=M/M_1$, $\delta_Pw\ge4$. Hence
\begin{equation*}
\min\bigl\{\min\{\delta_1,\delta_Pw\},\delta_L,1\bigr\}=\min\{\delta_1,\delta_L,1\}=\delta_\natural ,
\end{equation*}
so $\delta_\natural$ does not change. The constants of the lemma in (b) and of the site-(L)
theorem involve $B_H$ only through $C_\natural$ of \eqref{4.21:eq-base-point-profile-5} and
through $\varepsilon_{\mathfrak{Q}}$. Both are therefore recomputed once
with $B_H^S$ in place of $B_H$, within the upper bounds, taken in supremum form, under which
$C_\natural$ and $\varepsilon_{\mathfrak{Q}}$ are fixed.

One special case is worth recording. If $\delta_1\le\delta_Pw$ and
$B_H\ge B_\sharp\max\{C_1/K_R,1+K_\sharp\}$, then $B_H^S=B_H$ and
$\min\{\delta_1,\delta_Pw\}=\delta_1$, and no constant moves at all. We do not use these two
inequalities, and work throughout with the constant $B_H^S$ of
\eqref{4.21:eq-base-point-profile-a}.
\end{proof}

\begin{definition}[Cubes, tree lengths, cells, blocks and probes]\label{4.21:not-cubes-cells-probes}
The following conventions are in force in what follows. The space dimension $d$ is kept as a
letter; $L$ is the odd blocking factor of Definition~\ref{1.1:def-lattice} (so $L\ge3$) and $M=L^{m'}$
the cube-size parameter among Ba{\l}aban's auxiliary parameters $\mathfrak{p}$; in the quotations from
\cite{B-CMP119} below, $\eta$ is that paper's lattice spacing.
\begin{enumerate}
\item \emph{Cubes and distances.} For each level $\ell$, $\pi_\ell$ denotes the partition into the
closed $M$-cubes of level $\ell$, called \emph{cubes} (of level $\ell$). The partitions are nested
$L$-adically: every cube of $\pi_j$ with $j\ge \ell$ is a union of cubes of $\pi_\ell$, and every
cube of $\pi_\ell$ lies in exactly one cube of $\pi_j$. At level $\ell$ all lengths are measured in
units of the side of a cube of $\pi_\ell$. For closed sets $A,B$ we write $\operatorname{dist}_\ell(A,B)$
for their distance in the maximum norm, measured in this unit. Two closed sets \emph{touch} if
their distance is $0$. Two families of cubes of one and the same level \emph{meet} if they have a
cube in common; sets of different levels \emph{meet} if they have a common interior point.
\item \emph{Counting and enlargements.} $N_\ell(\cdot)$ is the number of cubes of $\pi_\ell$ in a
set. For a union $A$ of cubes of one level, $A^{(n)}$ is $A$ enlarged by $n$ layers of cubes of
that level. For a set $S$, $[S]_i$ is the union of the cubes of $\pi_i$ that meet $S$.
\item \emph{Domains and linear sizes.} $\mathbf{D}_\ell$ is the class of domains of
\cite{B-CMP109}: a domain is ``a union of a connected, finite family of cubes'' of $\pi_\ell$,
connectedness being understood through chains of cubes in which ``two consecutive cubes have a
common wall''. For $X\in\mathbf{D}_\ell$ one considers, following \cite{B-CMP109}, the ``tree graphs
contained in $X$ and intersecting all the cubes in $X$. A length of a shortest graph in this
class, divided by $M$, is the linear size of $X$'', the graphs being taken ``in the continuous
space''; this linear size is $d_\ell(X)$. Since a cube of $\pi_\ell$ has side $M$ in the units of
\cite{B-CMP109} at level $\ell$, $d_\ell(X)$ is the length of a shortest such tree graph in the
unit of level $\ell$ fixed above.
\item \emph{Relative tree length.} For a closed set $W$ and a union $X$ of cubes of $\pi_\ell$,
$d_{\ell,W}(X)$ is the least length of a tree graph contained in $X$ and intersecting every cube
of $X$ not contained in $W$; $d_{\ell,W}(X):=0$ if there is no such tree graph. For $\ell=k$,
$X\in\mathbf{D}_k$ and $W=Z$ this is the relative linear size $d_{k,Z}(X)$ of
\cite[(1.67), p.~376]{B-CMP122b}.
\item \emph{Cells and zones.} At a given stage of the construction the determining set of the
current sequence of domains $\Omega_j$ cuts space into closed cubes in proper scales; in the words
of \cite{B-CMP119}, ``if $\Delta\subset\Omega_j\setminus\Omega_{j+1}$, then the size of $\Delta$ is
equal to $ML^j\eta$''. These cubes are the \emph{cells}; for such a cell $\Delta$ we put
$\operatorname{lev}(\Delta):=j$, so that a cell of level $j$ is a cube of $\pi_j$. The \emph{zone}
$j$ is the union of the cells of level $j$.
\item \emph{Blocks.} For each layer $\Omega_j\setminus\Omega_{j+1}$, \cite{B-CMP119} constructs ``its
cover by cubes $\square$, which are unions of $2^d$ neighboring cubes from $\pi_j$, i.e.\ cubes of
the size $2ML^j\eta$'' (see also \cite{B-CMP109}). Accordingly, a \emph{block} $\square$ of zone $n$
is a cube which is the union of $2^d$ cubes of $\pi_n$, at least one of which is a cell of zone
$n$; the covers used are such that every cube of $\pi_n$ lies in at most $2^d$ blocks of zone $n$.
We write $\operatorname{zone}(\square)=n$, and $\tilde{\square}$, $\tilde{\square}^{4}$ for $\square$
enlarged by one, respectively four, layers of cubes of $\pi_n$.
\item \emph{Probes.} A \emph{probe of level $n$} is a closed cube $\mathsf{Q}$ which is a union of
cubes of $\pi_n$ and satisfies
\begin{equation}\label{4.21:eq-cubes-cells-probes-1}
\Delta_n\subset\mathsf{Q}\subset\{x:\ \operatorname{dist}_n(x,\Delta_n)\le 2\}
\quad\text{for some cell $\Delta_n$ of level $n$;}
\end{equation}
we put $\operatorname{lev}(\mathsf{Q}):=n$. A probe $\mathsf{Q}$ of level $n$ has side at most $5$,
is the union of at most $5^d$ cubes of $\pi_n$, and meets at most $5^d$ cubes of $\pi_j$ for every
$j\ge n$. A cell is a probe of its level; a block $\square$ and its enlargement $\tilde{\square}$
are probes of level $\operatorname{zone}(\square)$, whereas $\tilde{\square}^{4}$ is not a probe of
that level. The letter $\mathsf{Q}$ is also used, later in the book, for a family of fluctuation
cubes; the two meanings never occur in one sentence.
\item \emph{Constants.} We put $\mathsf{c}_d:=2\cdot 3^d$. The letters $c_d$ and $K(d)$ denote the
constants of \cite[(1.26)]{B-CMP116}, in the form
\begin{equation}\label{4.21:eq-cubes-cells-probes-2}
\sum_{A\in\mathbf{D}_\ell,\ A\ni\square'} e^{-(c_d+1)\,d_\ell(A)}\le K(d)
\qquad\text{for every cube $\square'$ of $\pi_\ell$.}
\end{equation}
Finally,
\begin{equation}\label{4.21:eq-cubes-cells-probes-3}
\lambda_\theta:=(1+3^d)\,\mathsf{c}_d\log 2 .
\end{equation}
\end{enumerate}
\end{definition}

\begin{proof}
We prove the assertions about probes made above, in the item headed \emph{Probes}.
Throughout, lengths at level $n$ are measured in the side of a cube of $\pi_n$, so that the cubes
of $\pi_n$ are the unit cubes of a lattice of unit cubes.

\emph{Side and number of cubes.} Let $\mathsf{Q}$ be a probe of level $n$ and $\Delta_n$ a cell as
in \eqref{4.21:eq-cubes-cells-probes-1}. Since $\Delta_n$ is a cube of $\pi_n$, it is a closed unit
cube, and the set of points at maximum-norm distance at most $2$ from it is the closed cube
$\Delta_n^{(2)}$ of side $1+2\cdot 2=5$ with the same centre. This cube is the union of the $5^d$
cubes of $\pi_n$ lying in it, and these are the only cubes of $\pi_n$ contained in it. By
\eqref{4.21:eq-cubes-cells-probes-1}, $\mathsf{Q}\subset\Delta_n^{(2)}$, so the side of the cube
$\mathsf{Q}$ is at most $5$. Moreover each cube of $\pi_n$ of which $\mathsf{Q}$ is the union is
contained in $\Delta_n^{(2)}$, hence is one of those $5^d$ cubes; thus $\mathsf{Q}$ is the union of
at most $5^d$ cubes of $\pi_n$.

\emph{Cubes of coarser levels met by a probe.} Let $j\ge n$. For $j=n$, a cube of $\pi_n$ meets the
family of cubes of $\mathsf{Q}$ if and only if it is one of them, and there are at most $5^d$ of
these. Let $j>n$, and let $Q'$ be a cube of $\pi_j$ which meets $\mathsf{Q}$, that is, which has a
common interior point $x$ with $\mathsf{Q}$. The set $\operatorname{int}Q'\cap\operatorname{int}\mathsf{Q}$
is a non-empty open set, hence has positive volume, and $\mathsf{Q}$ is a finite union of cubes of
$\pi_n$ whose boundaries have volume zero; therefore
$\operatorname{int}Q'\cap\operatorname{int}q\neq\emptyset$ for some cube $q$ of $\pi_n$ contained in
$\mathsf{Q}$. By the $L$-adic nesting, $q$ lies in exactly one cube $Q''$ of $\pi_j$, and then
$\operatorname{int}q\subset\operatorname{int}Q''$, so $\operatorname{int}Q'\cap\operatorname{int}Q''\ne\emptyset$.
Two distinct cubes of the partition $\pi_j$ have disjoint interiors, hence $Q''=Q'$, that is,
$q\subset Q'$. Choosing for each such $Q'$ one cube $q\subset\mathsf{Q}$ of $\pi_n$ with
$q\subset Q'$, we obtain a map from the cubes of $\pi_j$ meeting $\mathsf{Q}$ into the cubes of
$\pi_n$ of $\mathsf{Q}$; it is injective because every cube of $\pi_n$ lies in only one cube of
$\pi_j$. Hence $\mathsf{Q}$ meets at most $5^d$ cubes of $\pi_j$.

\emph{Cells.} If $\Delta$ is a cell of level $n$, then $\Delta$ is a closed cube, a union of cubes
of $\pi_n$ (namely of itself), and $\Delta\subset\Delta\subset\{x:\operatorname{dist}_n(x,\Delta)\le 2\}$;
so \eqref{4.21:eq-cubes-cells-probes-1} holds with $\Delta_n=\Delta$, and $\Delta$ is a probe of
level $n$.

\emph{Blocks and their first enlargement.} Let $\square$ be a block of zone $n$, and let
$\Delta\subset\square$ be a cell of zone $n$, hence of level $n$, among the $2^d$ cubes of $\pi_n$ of
which $\square$ is the union. Then $\square$ is a closed cube of side $2$ containing the unit cube
$\Delta$. In each coordinate direction the projection of $\square$ is an interval of length $2$
containing the unit-length projection of $\Delta$, so every point of that interval lies within
distance $1$ of the projection of $\Delta$; hence every point of $\square$ is within
maximum-norm distance $1$ of $\Delta$. The cube $\tilde{\square}$ is a closed cube and a union of
cubes of $\pi_n$, and every point of $\tilde{\square}$ lies within distance $1$ of $\square$,
hence within distance $1+1=2$ of $\Delta$. Since $\Delta\subset\square\subset\tilde{\square}$, both
$\square$ and $\tilde{\square}$ satisfy \eqref{4.21:eq-cubes-cells-probes-1} with $\Delta_n=\Delta$,
and they are probes of level $n=\operatorname{zone}(\square)$.

\emph{The fourth enlargement.} The cube $\tilde{\square}^{4}$ is a union of cubes of $\pi_n$ of side
$2+2\cdot 4=10>5$ in the unit of level $n$; by the first paragraph it is therefore not a probe of
level $n=\operatorname{zone}(\square)$.
\end{proof}

\begin{corollary}[Radius loss at one reading of the response]\label{4.21:cor-radius-loss}
Let $\eta^{T}_k(s)$ and $\eta^{L}_k(s)$ be the costs of one reading at the law site $s$, at site
(T) and at site (L) respectively, of the step from level $k-1$ to level $k$. Both are taken as
suprema over the families of decoupling parameters run by the decoupling clause
(Definition~\ref{4.21:def-decoupling-clause}). Put
\begin{equation*}
m(s):=\min\{m:\ s\in Z_m\},
\end{equation*}
with $m(s):=k_W$ at a frame site, $k_W$ denoting the level assigned to the frame, and put
$\mathsf{a}_k(s):=(k-N_k-m(s))_+$.
\begin{enumerate}
\item[(i)] At site (T), and at site (L) for $s\notin Z$ (the case (L)$(\alpha)$), the cost of
one reading, $\eta^{T}_k(s)$ at (T) and $\eta^{L}_k(s)$ at (L)$(\alpha)$, satisfies
\begin{equation}\label{4.21:eq-radius-loss-1}
\begin{aligned}
\eta^{T}_k(s),\ \eta^{L}_k(s)
&\le \mathsf{c}_k\,e^{-\delta_P w[(\check{R}_k-2)+\mathsf{a}_k(s)]}\\
&\le \tfrac12\, C_{\natural}\,\alpha^{+}_k\,e^{-\delta_{\natural}(\check{R}_k-2)}
\,e^{-2\mathsf{a}_k(s)} .
\end{aligned}
\end{equation}
Here $\mathsf{c}_k$ is the prefactor of the bound on the reading costs, with the constant
$B_H^{S}$ of \eqref{4.21:eq-base-point-profile-a} in place of $B_H$. At (T) it is
$\mathsf{c}_k=K_uB_H^{S}\alpha_{*,k}$. At (L)$(\alpha)$ it is the prefactor of that bound
read at $B_H^{S}$, and it is controlled through
$\varepsilon_{\mathfrak{Q}}(k)\le K_{\mathfrak{Q}}\alpha^{+}_k$.
The right-hand side of \eqref{4.21:eq-radius-loss-1} is the majorant of clause (iii) of the
lemma that introduces the reading costs, with $C_{\natural}$ taken at $B_H^{S}$. The bound uses
no lower bound on the width of any shell.
\item[(ii)] At site (L) for $s\in Z_h$ (the case (L)$(\beta)$), the decay rate is
$\delta_L$ per cube of side $M$; it is the rate of the case (L)$(\beta)$ of
Proposition~\ref{4.21:prop-base-point-profile}. The shell factor there is bounded by the shell
bound at frame sites, where $m(s):=k_W$; that bound is a bound at frame sites, not a pointwise
one, and this corollary uses it by statement.
\item[(iii)] Write $\ell 1,\ell 2,\dots$ for the radius allowances of the chart-radius
budget of the step. A wrapped term uses the following allowances:
\begin{itemize}
\item $\ell 1$ at (T), where no auxiliary disc in the interpolation time $\mathsf{t}$
occurs, so that $\ell 2$ is not used;
\item $\ell 3$ and $\ell 4$ at (P);
\item $\ell 5$, $\ell 6$ and $\ell 7$ at (L).
\end{itemize}
For a deep wrapped term, the auxiliary disc in $\lambda$ of the minimal form of the
third-factor lemma uses at most one half of the same allowance $\ell 7$, by
\eqref{4.21:eq-third-factor-form-a}. That disc serves only in the proof of the bound. The
stored piece lives on the chart $\mathfrak{G}_{\varepsilon'}$ of clause (a) of the proposition
fixing the chart of the stored pieces, with $\varepsilon'=\varepsilon-\eta_0-\mu$; here
$\varepsilon$ is the chart radius and $\eta_0$, $\mu$ are the two radius losses fixed in that
proposition. The allowance $\ell 7$ is thus shared. After the decoupling clause each piece
carries at most one auxiliary disc, by \eqref{4.21:eq-decoupling-clause-a}. The readings of
this corollary add no allowance to the chart-radius budget and no amount to any allowance in it.
If the chart-radius budget keeps a further allowance $\ell 8$ in reserve, its total radius
loss, which we write $\Lambda_{\mathrm{rad}}$, still satisfies $\Lambda_{\mathrm{rad}}\le 5\mu$,
the comparison $7/32<\tfrac14$ of that budget is unaffected, and nothing above changes.
\end{enumerate}
\end{corollary}

\begin{proof}
Items (ii) and (iii) rest on the statements named in them, in particular on
\eqref{4.21:eq-third-factor-form-a} and \eqref{4.21:eq-decoupling-clause-a}; it remains to
prove (i), which we do in five steps.

\emph{Location of $s$.} By the nesting of the zones, $Z_{m(s)}\subset\Omega^{c}_{m(s)+1}$.
Since $s\in Z_{m(s)}$, the block $B^{\mathrm{top}}(s)$ lies in cells of level at most
$m(s)$.

\emph{Site} (T). Consider the step from level $k-1$ to level $k$, with fluctuation domain
$\Omega=\Omega_k$. By clause (i) of the lemma that introduces the reading costs,
$s\in Z_{k-1}$.

Take any contour from $B^{\mathrm{top}}(s)$ to $\Omega_k$. It crosses each of the whole zones
$m(s)+1,\dots,k-2$, and there are $(k-2-m(s))_+$ of them. By clause $(g^{\wedge})$ of the
lemma on the geometry of zones, each such crossing contributes at least $w$ to the scaled
distance $d^{\wedge}$, whatever the thickness of the zone.

The contour then traverses the stretch from $Z_{k-1}$ to $\Omega_k$. By the proof of clause
(d) of the same lemma, this stretch consists of at least $8\check{R}_k$ units of level $k$,
through cells of level at most $k-1$.

These arcs of the contour are disjoint, so their contributions add:
\begin{equation}\label{4.21:eq-radius-loss-2}
d^{\wedge}(B^{\mathrm{top}}(s),\Omega)
\ge w\bigl[8\check{R}_k+(k-2-m(s))_+\bigr]
\ge w\bigl[(\check{R}_k-2)+\mathsf{a}_k(s)\bigr].
\end{equation}
The second inequality holds term by term. The separation radius $\check{R}_k$ of the zones
at level $k$ is a positive power of the blocking factor $L$, so $\check{R}_k\ge L\ge 2$; in
particular $8\check{R}_k\ge\check{R}_k-2$. Since
$N_k\ge 2$, we have $(k-2-m(s))_+\ge(k-N_k-m(s))_+=\mathsf{a}_k(s)$.

\emph{Site} (L), \emph{case} $(\alpha)$. Here $s\notin Z$, so $s$ lies either in another
component of $\Lambda_k^{c}$ or in $\Lambda_k$. Take any contour from $B^{\mathrm{top}}(s)$
to $\Lambda^{\sim}\subset Z$. It crosses the zones $m(s)+1,\dots,k-1$ of the component
containing $s$, each contributing at least $w$ by clause $(g^{\wedge})$. Inside $Z$ it
crosses at least $9\check{R}_k$ units of level $k$, by clause (f) of the lemma on the
geometry of zones. Hence
\begin{equation}\label{4.21:eq-radius-loss-3}
d^{\wedge}(B^{\mathrm{top}}(s),\Lambda^{\sim})
\ge w\bigl[9\check{R}_k+(k-1-m(s))_+\bigr]
\ge w\bigl[(\check{R}_k-2)+\mathsf{a}_k(s)\bigr].
\end{equation}
The second inequality follows as in \eqref{4.21:eq-radius-loss-2}, now from
$9\check{R}_k\ge\check{R}_k-2$ and $(k-1-m(s))_+\ge\mathsf{a}_k(s)$.

\emph{The first inequality.} Proposition~\ref{4.21:prop-base-point-profile} holds at every
law site of a wrapped term's chart, for every family run by the decoupling clause. At (T),
with $B_H^{S}$ as in \eqref{4.21:eq-base-point-profile-a}, it bounds the reading by
\begin{equation*}
K_uB_H^{S}\alpha_{*,k}\,e^{-\delta_P d^{\wedge}(B^{\mathrm{top}}(s),\Omega)} .
\end{equation*}
At (L)$(\alpha)$ it bounds the reading by
\begin{equation*}
\mathsf{c}_k\,e^{-\delta_P d^{\wedge}(B^{\mathrm{top}}(s),\Lambda^{\sim})},
\end{equation*}
where $\mathsf{c}_k$ is, as in (i), the prefactor of the bound on the reading costs at
(L)$(\alpha)$ read at $B_H^{S}$; the prefactor of the display of
Proposition~\ref{4.21:prop-base-point-profile} at (L)$(\alpha)$ is converted into it by the
conversion line \eqref{4.21:eq-base-point-profile-a}. In both cases no
width of a shell enters. The reading costs are suprema over these families. Since
$\delta_P>0$, the lower bounds \eqref{4.21:eq-radius-loss-2} and
\eqref{4.21:eq-radius-loss-3} give the first inequality in \eqref{4.21:eq-radius-loss-1}.

\emph{The second inequality.} By the level inequality recorded beside the conversion line
\eqref{4.21:eq-base-point-profile-a}, $\delta_P w\ge 4$; by the definition of $\delta_{\natural}$
in that conversion line, $\delta_{\natural}\le 1$. Hence
\begin{equation*}
\delta_P w\ge 4\ge\max\{\delta_{\natural},2\}.
\end{equation*}
Since $\check{R}_k-2\ge 0$, comparing exponents gives
$e^{-\delta_P w(\check{R}_k-2)}\le e^{-\delta_{\natural}(\check{R}_k-2)}$.
Since $\mathsf{a}_k(s)\ge 0$, also $e^{-\delta_P w\,\mathsf{a}_k(s)}\le e^{-2\mathsf{a}_k(s)}$.

It remains to bound the prefactor. The inequality
$\mathsf{c}_k\le\tfrac12 C_{\natural}\alpha^{+}_k$ is the step of the bound on the reading
costs that defines $C_{\natural}$, namely
\begin{equation*}
C_{\natural}=2e^{2}\max\{C_LK_{\mathfrak{Q}},K_uB_H^{S},K_u\},
\end{equation*}
the constant of the conversion line \eqref{4.21:eq-base-point-profile-a} with $B_H$ replaced
by $B_H^{S}$. This proves the second inequality in \eqref{4.21:eq-radius-loss-1},
and with it (i).
\end{proof}

\begin{lemma}[Elementary facts on tree lengths and packings]\label{4.21:lem-tree-length-facts}
Fix a level $\ell$. Cubes, tree graphs, $\mathbf{D}_\ell$, $d_\ell$, $N_\ell(\cdot)$ and $[S]_i$ are
as in Notation~\ref{4.21:not-cubes-cells-probes}. At level $\ell$ lengths are measured in units of
the side of a $\pi_\ell$-cube, and a cube is a cube of $\pi_\ell$ unless another level is named. We
write $\operatorname{dist}$ for the $\infty$-distance $\operatorname{dist}_\ell$. The length of a tree
graph is the sum of the Euclidean lengths of its segments.

Two cubes of one level are \emph{wall-neighbours} if they share a wall, that is, a whole
$(d-1)$-dimensional face. A \emph{wall-chain} is a finite sequence of cubes of one level in which
consecutive members are wall-neighbours. A \emph{box} is a product of $d$ closed intervals, and its
sides are the lengths of these intervals.

There is a constant $\mathsf{c}_d$, depending only on $d$, such that the following hold:
\begin{equation}\label{4.21:eq-tree-length-facts-a}
\begin{aligned}
&\text{(g1)}\quad &&\text{a tree graph of length $t$ intersects at most $\mathsf{c}_d(t+1)$ cubes,
hence}\\
& &&N_\ell(X)\le \mathsf{c}_d\bigl(d_\ell(X)+1\bigr)\quad\text{for } X\in\mathbf{D}_\ell;\\
&\text{(g2)} &&X\subset Y,\ X,Y\in\mathbf{D}_\ell\ \Longrightarrow\
d_\ell(Y)\le d_\ell(X)+N_\ell(Y\setminus X);\\
&\text{(g3)} &&X\in\mathbf{D}_j,\ i\ge j\ \Longrightarrow\ [X]_i\in\mathbf{D}_i,\\
& &&\qquad d_i([X]_i)\le L^{-(i-j)}d_j(X),\quad N_i([X]_i)\le N_j(X);\\
&\text{(g4)} &&\text{a tree graph intersecting two cubes $Q,Q'$ has length at least
$\operatorname{dist}(Q,Q')$;}\\
&\text{(g5)} &&\text{two cubes at distance $x$ are joined by a wall-chain with fewer than}\\
& &&\text{$d(x+1)$ intermediate cubes, all inside the smallest box containing both;}\\
&\text{(g6)} &&\text{sets pairwise at distance at least $R>0$, each intersecting one box}\\
& &&\text{whose sides are at most $S$, number at most $(S/R+1)^d$;}\\
&\text{(g7)} &&\text{a box all of whose sides are at least the side of a cube $\mathsf{Q}$,}\\
& &&\text{and which intersects $\mathsf{Q}$, contains a vertex of $\mathsf{Q}$.}
\end{aligned}
\end{equation}
\end{lemma}

\begin{proof}
\emph{(g1).} Let $\mathcal{T}$ be a tree graph of length $t$. The closed walk that traverses every
edge of $\mathcal{T}$ twice has length $2t$ and passes through every point of $\mathcal{T}$. We
parametrise it by arclength and cut it into at most $2t+1$ arcs of length at most $1$.

In level-$\ell$ units the cubes are products of grid intervals $[n,n+1]$, $n\in\mathbb{Z}$. The
projection of an arc of length at most $1$ onto a coordinate axis is an interval of length at most
$1$, and such an interval meets at most $3$ grid intervals. A cube is intersected by the arc only if
each of the $d$ projections of the arc meets the corresponding factor of the cube. Hence one arc
intersects at most $3^d$ cubes, and $\mathcal{T}$ intersects at most
\begin{equation*}
3^d(2t+1)\le 2\cdot 3^d\,(t+1)
\end{equation*}
cubes. The first assertion therefore holds with $\mathsf{c}_d=2\cdot 3^d$.

Applied to a tree graph of length $d_\ell(X)$ intersecting every cube of $X$, it gives
$N_\ell(X)\le\mathsf{c}_d(d_\ell(X)+1)$.

\emph{(g2).} Let $n=N_\ell(Y\setminus X)$. Using the connectedness of $Y\in\mathbf{D}_\ell$, we
enumerate the cubes of $Y$ not in $X$ as $\Delta_1,\dots,\Delta_n$ in such a way that, for each $t$,
the cube $\Delta_t$ shares a wall $\mathsf{w}$ with a cube $Q$ of
\begin{equation*}
X_{t-1}:=X\cup\Delta_1\cup\dots\cup\Delta_{t-1}.
\end{equation*}

Let $\mathcal{T}\subset X_{t-1}$ be a tree graph of length $d_\ell(X_{t-1})$ intersecting all cubes
of $X_{t-1}$. Take $p\in\mathcal{T}\cap Q$, and let $p'$ be its orthogonal projection on
$\mathsf{w}$. Since $\mathsf{w}$ is a whole face of $Q$ and $Q$ is convex, $[p,p']\subset Q$. The
two points differ in one coordinate only, by at most the side of $Q$, so $|p-p'|\le 1$. Moreover
$p'\in\mathsf{w}\subset\Delta_t$.

Hence $\mathcal{T}\cup[p,p']$ serves for $X_t$: it intersects every cube of $X_t$, and its length is
at most that of $\mathcal{T}$ plus $1$. Therefore $d_\ell(X_t)\le d_\ell(X_{t-1})+1$. Iterating from
$t=1$ to $t=n$ gives
\begin{equation*}
d_\ell(Y)=d_\ell(X_n)\le d_\ell(X)+n=d_\ell(X)+N_\ell(Y\setminus X).
\end{equation*}

\emph{(g3).} Let $\mathcal{T}$ be an optimal tree graph of $X$, of length $d_j(X)$ in level-$j$
units. It lies in $X\subset[X]_i$.

Every $\pi_i$-cube of $[X]_i$ has an interior point in common with $X$. Hence it has such a point in
common with some $\pi_j$-cube of $X$. Since the partitions are nested $L$-adically, it then
contains that $\pi_j$-cube. Consequently $\mathcal{T}$ intersects every $\pi_i$-cube of $[X]_i$.

The side of a $\pi_i$-cube is $L^{i-j}$ times the side of a $\pi_j$-cube, so the length of
$\mathcal{T}$ in level-$i$ units is $L^{-(i-j)}d_j(X)$. This gives
$d_i([X]_i)\le L^{-(i-j)}d_j(X)$.

Distinct $\pi_i$-cubes have disjoint interiors, so the $\pi_j$-cubes of $X$ they contain are
distinct. This gives $N_i([X]_i)\le N_j(X)$.

Finally, take two $\pi_i$-cubes of $[X]_i$ and $\pi_j$-cubes of $X$ inside them. Join the latter by
a wall-chain in $X$, as the connectedness of $X\in\mathbf{D}_j$ allows. The parents (the
$\pi_i$-cubes containing them) of two wall-neighbours either coincide or are wall-neighbours, since
the common wall lies in both parents. The parents along the chain therefore connect the two given
$\pi_i$-cubes inside $[X]_i$, and $[X]_i\in\mathbf{D}_i$.

\emph{(g4).} Take points $p\in\mathcal{T}\cap Q$ and $p'\in\mathcal{T}\cap Q'$. The path in
$\mathcal{T}$ between them contains both points, so its $\infty$-diameter is at least
\begin{equation*}
|p-p'|_\infty\ge\operatorname{dist}(Q,Q').
\end{equation*}
A connected union of segments has length at least its $\infty$-diameter, and $\mathcal{T}$ contains
this path.

\emph{(g5).} Label the two cubes by their lowest corners $q,q'\in\mathbb{Z}^d$. Their distance is
$x=\max_\mu\bigl(|q_\mu-q'_\mu|-1\bigr)_+$, so $|q_\mu-q'_\mu|\le x+1$ for every $\mu$.

Passing from $q$ to $q'$ by changing one coordinate at a time by $\pm1$ produces a wall-chain. It has
\begin{equation*}
\sum_\mu|q_\mu-q'_\mu|\le d(x+1)
\end{equation*}
steps, hence fewer than $d(x+1)$ intermediate cubes. Each coordinate of every corner along the way
lies between $q_\mu$ and $q'_\mu$. Hence all cubes of the chain lie in the smallest box containing
both given cubes.

\emph{(g6).} Pick in each set a point lying in the box. A side of length $s\le S$ is covered by
$\lfloor s/R\rfloor+1\le S/R+1$ half-open intervals of length $R$. Taking products, the box is
covered by at most $(S/R+1)^d$ half-open boxes of side $R$.

Two points of one such box are at $\infty$-distance less than $R$, whereas points of distinct sets
are at distance at least $R$. So each half-open box holds at most one of the chosen points, and the
number of sets is at most $(S/R+1)^d$.

\emph{(g7).} Write $\mathsf{Q}=\prod_\mu[a_\mu,b_\mu]$ and let the box be $\prod_\mu[c_\mu,e_\mu]$.
On each axis, $[c_\mu,e_\mu]$ intersects $[a_\mu,b_\mu]$ and has length at least $b_\mu-a_\mu$.

Suppose it contained neither $a_\mu$ nor $b_\mu$. Being an interval that meets $[a_\mu,b_\mu]$, it
would then lie in $(a_\mu,b_\mu)$, and its length would be less than $b_\mu-a_\mu$, which is
impossible. So it contains some $v_\mu\in\{a_\mu,b_\mu\}$.

Taking products, $v=(v_\mu)_\mu$ is a vertex of $\mathsf{Q}$ contained in the box.
\end{proof}

\begin{lemma}[Transfer of orbit data through the decoupling expansion]\label{4.21:lem-orbit-transfer}
Let $F$ be a stored term in format, with an orbit datum $F^{\sharp}$ of radius $\varepsilon$ in the
sense of Notation~\ref{4.21:not-orbit-data}. Fix a link of the decoupling clause
\eqref{4.21:eq-decoupling-clause-a}, with link map $\mathsf{R}(\mathsf{s})$ and carried gauge map
$\mathsf{u}(\mathsf{s})$ (Definition~\ref{4.21:def-links-readers}), and let $h$ and
$\mathsf{R}(\mathsf{s})^{h}$ be as in Proposition~\ref{4.21:prop-link-containment}. Put
\begin{equation}\label{4.21:eq-orbit-transfer-1}
f^{\sharp}(\mathsf{s};g') := F^{\sharp}\bigl(\mathsf{R}(\mathsf{s})^{h};\mathbf{A},\boldsymbol{\Phi},
h\,\mathsf{u}(\mathsf{s})\,g'\bigr),
\qquad
\eta(s) := \sup_{\mathsf{s}}\operatorname{pol}\bigl((h\,\mathsf{u}(\mathsf{s}))(s)\bigr),
\end{equation}
where $s$ runs over the law sites of $X_F$ and the supremum is over the link's polydisc
$|\mathsf{s}(\Delta)|\le e^{\kappa_1}$; $\eta$ is bounded by
Proposition~\ref{4.21:prop-base-point-profile}. Assume $\eta(s)<\varepsilon(s)$ at every law site $s$.
Then the following hold.
\begin{enumerate}
\item[(i)] $f^{\sharp}$ is holomorphic on the product of the link's polydisc, the source polydisc,
the output tube, the weak coordinates (the weak slots of the link) together with the parameter $w$
of the weak slots, and $\mathfrak{G}_{\varepsilon-\eta}(X_F)$. On this product $|f^{\sharp}|\le
\mathfrak{N}^{\mathbb{C}}_{\varepsilon}(F)$. This is the holomorphy condition
\eqref{4.21:eq-links-readers-b} read in the complex size functional $\mathfrak{N}^{\mathbb{C}}$.
\item[(ii)] Clauses (i)--(iv) of Lemma~\ref{4.21:lem-first-stage-expansion} hold for $f^{\sharp}$,
with the same index set $\mathfrak{Y}(F)$. For $\mathsf{y}\in\mathfrak{Y}(F)$ put
\begin{equation}\label{4.21:eq-orbit-transfer-2}
F^{\sharp}_{\mathsf{y}} := \int_{[0,1]^{\mathsf{y}}}\prod_{\Delta\in\mathsf{y}}
d\mathsf{s}(\Delta)\,\partial_{\mathsf{s}(\Delta)}f^{\sharp}(\mathsf{s};g')
\Big|_{\mathsf{s}=0\text{ off }\mathsf{y}} ,
\end{equation}
regarded as a function of $g'$ and of the variables of (i) other than $\mathsf{s}$, which are carried
along.
Then, for every $g'$ in the whole chart $\mathfrak{G}_{\varepsilon-\eta}(X_F)$ and at every point of
the factors of the domain of (i) other than the link's polydisc and the chart,
\begin{equation}\label{4.21:eq-orbit-transfer-3}
|F^{\sharp}_{\mathsf{y}}| \le \mathfrak{N}^{\mathbb{C}}_{\varepsilon}(F)\,
e^{-(\kappa_1-1)|\mathsf{y}|}.
\end{equation}
\item[(iii)] $F^{\sharp}_{\mathsf{y}}$ is an orbit datum of radius $\varepsilon-\eta$ of the piece
$F_{\mathsf{y}}$ of Lemma~\ref{4.21:lem-first-stage-expansion}. At real data, $\mathsf{s}=\mathbf{1}$
and $g'=1$, the left member of clause (i) of Lemma~\ref{4.21:lem-first-stage-expansion}, written for
$f^{\sharp}$, is the bracket of Ba{\l}aban's construction.
\end{enumerate}
\end{lemma}

\begin{proof}
(i) By Proposition~\ref{4.21:prop-link-containment}, on the whole of the link's polydisc, source
polydisc and output tube, the restriction of $\mathsf{R}(\mathsf{s})^{h}$ to $X_F$ lies in the
stored domain $\mathfrak{U}_F(X_F)$. This is the domain in the first argument of $F^{\sharp}$.

Next let $g'\in\mathfrak{G}_{\varepsilon-\eta}(X_F)$ and let $s$ be a law site. The subadditivity
of $\operatorname{pol}$ in \eqref{4.21:eq-orbit-data-a} and the definition of $\eta$ in
\eqref{4.21:eq-orbit-transfer-1} give
\begin{equation*}
\operatorname{pol}\bigl((h\,\mathsf{u}(\mathsf{s})\,g')(s)\bigr)
\le \operatorname{pol}\bigl((h\,\mathsf{u}(\mathsf{s}))(s)\bigr)+\operatorname{pol}(g'(s))
\le \eta(s)+\operatorname{pol}(g'(s)) .
\end{equation*}
Since $g'\in\mathfrak{G}_{\varepsilon-\eta}(X_F)$, $\operatorname{pol}(g'(s))<\varepsilon(s)-\eta(s)$,
so the right-hand side is $<\varepsilon(s)$, and the gauge argument of $F^{\sharp}$ therefore lies in
$\mathfrak{G}_{\varepsilon}(X_F)$.

By the holomorphy property of an orbit datum \eqref{4.21:eq-orbit-data-a}, $F^{\sharp}$ is
holomorphic in its first argument $P\in\mathfrak{U}_F(X_F)$, in the weak coordinates and in the
gauge argument. The carried map
$\mathsf{u}(\mathsf{s})$ is holomorphic in the site's parameters, by clause (i) of the lemma on the
carried gauge map $\mathfrak{u}$. The link map $\mathsf{R}(\mathsf{s})$ is holomorphic in the link's
polydisc, the source polydisc and the output tube by Definition~\ref{4.21:def-links-readers}. Hence
$f^{\sharp}$ is a composition of holomorphic maps, and so is holomorphic on the stated
product.

Finally, $\mathfrak{N}^{\mathbb{C}}_{\varepsilon}(F)=\sup|F^{\sharp}|$ by
Notation~\ref{4.21:not-orbit-data}, and $f^{\sharp}$ takes values of $F^{\sharp}$ only. Therefore
$|f^{\sharp}|\le\mathfrak{N}^{\mathbb{C}}_{\varepsilon}(F)$.

(ii) By clause (i) of the lemma on the carried gauge map $\mathfrak{u}$, $\mathsf{u}(\mathsf{s})(s)$
reads the link's layer $\mathcal{Y}(\mathsf{s})$ only on $B^{\mathrm{top}}(s)$, and it equals $1$
where $\mathcal{Y}$ vanishes. Let $\mathcal{R}$ be the reading set of $F$ at the link, as in
Lemma~\ref{4.21:lem-first-stage-expansion}. Then
\begin{equation*}
B^{\mathrm{top}}(s)\subset X_F\subset\mathcal{R}.
\end{equation*}
The first inclusion holds because, by the same lemma, $B^{\mathrm{top}}(s)$ lies in the $M$-cube of
the local scale of $s$, inside the domain $X_F$. The second inclusion is the inclusion
$\mathcal{R}\supseteq X_F$, which holds at every link by the inclusion established for the reading
sets in the construction of the links. Hence $f^{\sharp}$ satisfies the class
condition \eqref{4.21:eq-links-readers-a}, with the same reading set $\mathcal{R}$ as $F$.

By (i), $f^{\sharp}$ satisfies the holomorphy condition \eqref{4.21:eq-links-readers-b}, with
$\mathfrak{N}^{\mathbb{C}}_{\varepsilon}(F)$ as its bound. The hypotheses of
Lemma~\ref{4.21:lem-first-stage-expansion} are therefore met by $f^{\sharp}$, at each fixed value of
$g'$ and of the remaining variables, with the same $\mathcal{R}$. Its proof applies to $f^{\sharp}$
as it stands, so its clauses (i)--(iv) hold for $f^{\sharp}$.

The index set of that lemma consists of the families $\mathsf{y}$ each of whose components holds a
source cell and a member seeing $\mathcal{R}$. It depends on $F$ only through $\mathcal{R}$, and so
is the same set $\mathfrak{Y}(F)$. The bound of clause (ii) of that lemma, with the sup bound
$\mathfrak{N}^{\mathbb{C}}_{\varepsilon}(F)$ of (i), is \eqref{4.21:eq-orbit-transfer-3}. It holds
for every $g'$ in the whole chart $\mathfrak{G}_{\varepsilon-\eta}(X_F)$ and at every point of the
factors of the domain of (i) other than the link's polydisc and the chart.

(iii) Holomorphy. Each derivative $\partial_{\mathsf{s}(\Delta)}f^{\sharp}$ in
\eqref{4.21:eq-orbit-transfer-2} is a Cauchy integral of the holomorphic function $f^{\sharp}$, which
is bounded by the single constant $\mathfrak{N}^{\mathbb{C}}_{\varepsilon}(F)$ of (i). Cauchy
integrals and integrals over $[0,1]^{\mathsf{y}}$ of holomorphic functions under one bound are
holomorphic. Hence $F^{\sharp}_{\mathsf{y}}$ is holomorphic in its first argument, the weak
coordinates and $g'\in\mathfrak{G}_{\varepsilon-\eta}(X_F)$. This is the holomorphy property of an
orbit datum of radius $\varepsilon-\eta$.

Transformation law. Let $v$ and $g$ take values in the gauge group $G$, and let
$P\in\mathfrak{U}_F(X_F)$. Apply the transformation law
of \eqref{4.21:eq-orbit-data-a} once with gauge argument $gv$, and once with $g$ at the data
$(\mathbf{A}^{[v]},\boldsymbol{\Phi}^{[v]})$. Use the composition rule
$(\mathbf{A}^{[v]})^{[g]}=\mathbf{A}^{[gv]}$, which is the rule $(\mathbf{A}^{[a]})^{[b]}=
\mathbf{A}^{[ba]}$ with $a=v$ and $b=g$, and the same rule for $\boldsymbol{\Phi}$. This gives
\begin{equation}\label{4.21:eq-orbit-transfer-4}
\begin{split}
F^{\sharp}(P;\mathbf{A},\boldsymbol{\Phi},gv)
&= F\bigl(P;\mathbf{A}^{[gv]},\boldsymbol{\Phi}^{[gv]}\bigr)\\
&= F\bigl(P;(\mathbf{A}^{[v]})^{[g]},(\boldsymbol{\Phi}^{[v]})^{[g]}\bigr)
= F^{\sharp}\bigl(P;\mathbf{A}^{[v]},\boldsymbol{\Phi}^{[v]},g\bigr).
\end{split}
\end{equation}
Since $v$ is $G$-valued, $\operatorname{pol}(v(s))=0$. By subadditivity, $gv$ lies in the chart
$\mathfrak{G}_{\varepsilon}(X_F)$ whenever $g$ does. The first and last members of
\eqref{4.21:eq-orbit-transfer-4} are therefore holomorphic in $g$ on
$\mathfrak{G}_{\varepsilon}(X_F)$, and they agree at $G$-valued $g$. By the identity principle of
\eqref{4.21:eq-orbit-data-a}, the equality of the first and last members of
\eqref{4.21:eq-orbit-transfer-4} holds for all $g\in\mathfrak{G}_{\varepsilon}(X_F)$.

We use the equality of the first and last members of \eqref{4.21:eq-orbit-transfer-4} with
$P=\mathsf{R}(\mathsf{s})^{h}$ and
$g=h\,\mathsf{u}(\mathsf{s})\,g'$, which lies in $\mathfrak{G}_{\varepsilon}(X_F)$ by (i); since
$\operatorname{pol}(v(s))=0$, subadditivity also gives $g'v\in\mathfrak{G}_{\varepsilon-\eta}(X_F)$.
By \eqref{4.21:eq-orbit-transfer-1}, the identity then states that $f^{\sharp}(\mathsf{s};g'v)$ at
the data $(\mathbf{A},\boldsymbol{\Phi})$ equals $f^{\sharp}(\mathsf{s};g')$ at the data
$(\mathbf{A}^{[v]},\boldsymbol{\Phi}^{[v]})$, for every $\mathsf{s}$ in the link's polydisc.
Differentiating in $\mathsf{s}(\Delta)$, $\Delta\in\mathsf{y}$, and integrating over
$[0,1]^{\mathsf{y}}$ as in \eqref{4.21:eq-orbit-transfer-2}, under the integral, this identity gives
the transformation law for $F^{\sharp}_{\mathsf{y}}$. Hence $F^{\sharp}_{\mathsf{y}}$ is an orbit
datum of radius $\varepsilon-\eta$ of $F_{\mathsf{y}}$.

Real point. At real data, $\mathsf{s}=\mathbf{1}$ and $g'=1$, the left member of clause (i) of
Lemma~\ref{4.21:lem-first-stage-expansion}, written for $f^{\sharp}$, is the bracket of Ba{\l}aban's
construction described in (iii). This follows from clause (ii) of the lemma on the carried gauge map
$\mathfrak{u}$ and from the real-point clause (iv) of the statement that fixes the link's data.
\end{proof}

\begin{lemma}[Relative tree length as the length of a skeleton]\label{4.21:lem-skeleton}
Work at a level $\ell$, with the cubes $\pi_\ell$, the domains $\mathbf{D}_\ell$, the length
$d_\ell$, the counting function $N_\ell(\cdot)$ and the enlargement $A^{(1)}$ of
Notation~\ref{4.21:not-cubes-cells-probes}. Let $\mathsf{I}$ be a finite family of boxes, each
a union of cubes of $\pi_\ell$, no two of which touch, and put $W:=\bigcup\mathsf{I}$. For
$X\in\mathbf{D}_\ell$ let $d_{\ell,W}(X)$ be the relative linear size of
\cite[(1.67)]{B-CMP122b}, in units of the side of a cube of $\pi_\ell$: the least length of a
tree graph contained in $X$ and intersecting every cube of $X\setminus W$. Put
\begin{equation*}
\mathbf{D}_\ell^{\mathsf{I}}:=\bigl\{X\in\mathbf{D}_\ell:\ X\cap W \text{ is a union of
members of } \mathsf{I},\ X\not\subset W\bigr\}.
\end{equation*}
For $X\in\mathbf{D}_\ell^{\mathsf{I}}$ fix a tree graph $\mathcal{T}_X$ attaining
$d_{\ell,W}(X)$, chosen by a rule that reads $(\pi_\ell,\mathsf{I},X)$ only, and put
\begin{equation*}
\operatorname{sk}(X):=\{\Delta\in\pi_\ell:\ \Delta\cap\mathcal{T}_X\neq\emptyset\}.
\end{equation*}
Then the following hold.
\begin{enumerate}
\item[(a)] $\operatorname{sk}(X)\in\mathbf{D}_\ell$ and
$X\setminus W\subset\operatorname{sk}(X)\subset X^{(1)}$; every $I\in\mathsf{I}$ with
$I\subset X$ touches a cube of $X\setminus W$.
\item[(b)] One has
\begin{equation*}
d_\ell(\operatorname{sk}(X))\le d_{\ell,W}(X)\le d_\ell(X),
\end{equation*}
and $d_{\ell,W}(X)\le d_\ell(A)$ for every $A\in\mathbf{D}_\ell$ with
$X\setminus W\subset A\subset X$.
\item[(c)] For every $A\in\mathbf{D}_\ell$,
\begin{equation}\label{4.21:eq-skeleton-1}
\#\bigl\{X\in\mathbf{D}_\ell^{\mathsf{I}}:\ \operatorname{sk}(X)=A\bigr\}
\le 2^{(1+3^d)N_\ell(A)}\le e^{\lambda_\theta(d_\ell(A)+1)},
\end{equation}
where $\lambda_\theta$ denotes a constant chosen subject to the lower bound
$\lambda_\theta\ge(1+3^d)\mathsf{c}_d\log 2$, which is the only property of $\lambda_\theta$
used here, and
$\mathsf{c}_d$ is the constant of part (g1) of \eqref{4.21:eq-tree-length-facts-a}.
\end{enumerate}
\end{lemma}

\begin{proof}
Throughout, cubes are the closed cubes of $\pi_\ell$, and by
Notation~\ref{4.21:not-cubes-cells-probes} a set is in $\mathbf{D}_\ell$ exactly when it is a
non-empty union of cubes any two of which are joined by a chain of its cubes in which
consecutive cubes share a wall (a wall-chain); a family of cubes with this property is called
wall-connected.

(a) For a point $x$ put $\mathcal{E}(x):=\{\Delta\in\pi_\ell:\ x\in\Delta\}$, the cubes
containing $x$. This family is
wall-connected: at each coordinate in which $x$ lies on a grid hyperplane of $\pi_\ell$ a cube
containing $x$ lies on one of the two sides, the cubes of $\mathcal{E}(x)$ correspond to the
choices of side, one per such coordinate, and changing a single choice passes to a cube sharing
a wall with the previous one. Moreover $\mathcal{E}(x')\subset\mathcal{E}(x)$ for every $x'$
close enough to $x$, since the finitely many cubes near $x$ that do not contain $x$ are closed
and at positive distance from it. By definition
$\operatorname{sk}(X)=\bigcup_{x\in\mathcal{T}_X}\mathcal{E}(x)$.

Suppose that $\operatorname{sk}(X)$ were the union of two disjoint non-empty subfamilies $F_1$,
$F_2$ with no cube of $F_1$ sharing a wall with a cube of $F_2$. Since each $\mathcal{E}(x)$,
$x\in\mathcal{T}_X$, is non-empty and wall-connected, it is contained in exactly one of $F_1$,
$F_2$. The sets $O_i:=\{x\in\mathcal{T}_X:\ \mathcal{E}(x)\subset F_i\}$, $i=1,2$, are therefore
disjoint and cover $\mathcal{T}_X$; each is relatively open in $\mathcal{T}_X$ by the inclusion
$\mathcal{E}(x')\subset\mathcal{E}(x)$ for $x'$ near $x$; and each is non-empty, since a cube
$\Delta\in F_i$ intersects $\mathcal{T}_X$ at some point $x$, and then $\Delta\in\mathcal{E}(x)$
forces
$\mathcal{E}(x)\subset F_i$. This contradicts the connectedness of the tree graph
$\mathcal{T}_X$. Hence $\operatorname{sk}(X)$ is wall-connected, and it is non-empty because
$X\not\subset W$ and $\mathcal{T}_X$ intersects the cubes of $X\setminus W$; so
$\operatorname{sk}(X)\in\mathbf{D}_\ell$.

The tree $\mathcal{T}_X$ is contained in $X$ and intersects every cube of $X\setminus W$;
the second property gives $X\setminus W\subset\operatorname{sk}(X)$, and the first shows that
every cube intersecting $\mathcal{T}_X$ intersects $X$, hence touches $X$ and lies in $X^{(1)}$,
which gives $\operatorname{sk}(X)\subset X^{(1)}$.

Let $I\in\mathsf{I}$ with $I\subset X$. Since $X\not\subset W$, $X$ has a cube outside $W$;
take a wall-chain of cubes of $X$ from a cube of $I$ to it, and let $\Delta$ be the first cube
of the chain not contained in $I$. Then $\Delta$ shares a wall with a cube of $I$, so it
touches $I$. It lies in no other member $I'$ of $\mathsf{I}$, for then $I$ and $I'$ would
touch. As $X\cap W$ is a union of members of $\mathsf{I}$, a cube of $X$ contained in $W$ lies in
a member of $\mathsf{I}$; hence $\Delta$ is a cube of $X\setminus W$, and $I$ touches it.

(b) Every point of $\mathcal{T}_X$ lies in a cube, which then intersects $\mathcal{T}_X$; so
$\mathcal{T}_X$ lies in $\bigcup\operatorname{sk}(X)$, and by definition it intersects each
cube of $\operatorname{sk}(X)$. It is thus admissible for $d_\ell(\operatorname{sk}(X))$, and
$d_\ell(\operatorname{sk}(X))\le d_{\ell,W}(X)$. A tree graph admissible for $d_\ell(X)$ lies in
$X$ and intersects every cube of $X$, in particular every cube of $X\setminus W$, so it is
admissible for $d_{\ell,W}(X)$; hence $d_{\ell,W}(X)\le d_\ell(X)$. Finally, if
$A\in\mathbf{D}_\ell$ and $X\setminus W\subset A\subset X$, a tree graph attaining $d_\ell(A)$
lies in $A\subset X$ and intersects all cubes of $A\supset X\setminus W$; it is admissible for
$d_{\ell,W}(X)$, so $d_{\ell,W}(X)\le d_\ell(A)$.

(c) Fix $A\in\mathbf{D}_\ell$ and let $X\in\mathbf{D}_\ell^{\mathsf{I}}$ with
$\operatorname{sk}(X)=A$. By (a), $X\setminus W$ is a subfamily of $A$, which leaves at most
$2^{N_\ell(A)}$ possibilities. Since $X\cap W$ is a union of members of $\mathsf{I}$,
\begin{equation*}
X=(X\setminus W)\cup\bigcup\{I\in\mathsf{I}:\ I\subset X\},
\end{equation*}
and by (a) each $I\in\mathsf{I}$ with $I\subset X$ touches a cube of $X\setminus W$, that is a
cube $\Delta$ of $A$ not contained in $W$. Such an $I$ does not contain $\Delta$ and touches
it, so it contains a cube touching $\Delta$ and different from $\Delta$; distinct members of
$\mathsf{I}$ do not touch, hence are disjoint, and so contain distinct such cubes. There are
fewer than $3^d$ cubes touching a given $\Delta$ other than $\Delta$ itself. Thus the members of
$\mathsf{I}$ contained in $X$ form a subset of a family that depends on $A$ only and has fewer
than $3^dN_\ell(A)$ members, which leaves at most $2^{3^dN_\ell(A)}$ possibilities. Together,
\begin{equation*}
\#\bigl\{X\in\mathbf{D}_\ell^{\mathsf{I}}:\ \operatorname{sk}(X)=A\bigr\}
\le 2^{N_\ell(A)}\,2^{3^dN_\ell(A)}=2^{(1+3^d)N_\ell(A)}.
\end{equation*}
By part (g1) of the facts \eqref{4.21:eq-tree-length-facts-a} in
Lemma~\ref{4.21:lem-tree-length-facts},
$N_\ell(A)\le\mathsf{c}_d(d_\ell(A)+1)$, whence
\begin{equation*}
2^{(1+3^d)N_\ell(A)}\le\exp\bigl((1+3^d)\mathsf{c}_d\log 2\,(d_\ell(A)+1)\bigr),
\end{equation*}
which is the second inequality of \eqref{4.21:eq-skeleton-1}, since
$\lambda_\theta\ge(1+3^d)\mathsf{c}_d\log 2$.
\end{proof}

Thus, by (a) and (b), the relative length $d_{\ell,W}(X)$ lies between the length
$d_\ell(\operatorname{sk}(X))$ of a domain that contains $X\setminus W$ and lies in $X^{(1)}$,
and the least length $d_\ell(A)$ of a domain $A\in\mathbf{D}_\ell$ with
$X\setminus W\subset A\subset X$: up to cubes adjacent to $X$, it is the least length of a
domain of $\mathbf{D}_\ell$ containing every cube of $X$ outside $W$; the cubes a tree crosses
between the components of $X\setminus W$ are counted, the cubes of the members of
$\mathsf{I}$ are not.

\begin{proof}[Proof of Theorem~\ref{4.21:thm-format-propagation}]
\label{4.21:prf-propagation-proof}
The statements of this section are used in the order in which they are established. At each point below only statements that are already available at that point are invoked. The proof has four parts: holomorphy, the numerical bounds, the orbit data, and the radii.

\emph{Holomorphy.}
We prove first the holomorphy statement \eqref{4.21:eq-links-readers-b} in the complex size
functional $\mathfrak{N}^{\mathbb{C}}_{\varepsilon}$. At every link of the decoupling expansion it
is clause~(i) of the transfer of orbit data through the decoupling expansion
(Lemma~\ref{4.21:lem-orbit-transfer}). That lemma needs two inputs:
\begin{itemize}
\item the restriction to $X_F$ of $\mathsf{R}(\mathsf{s})^{h}$ lies in the stored domain
$\mathfrak{U}_F(X_F)$ of \eqref{4.21:eq-orbit-data-a}, at every link and on the whole of its
polydisc, source polydisc and output tube;
\item the quantity $\eta(s)$ is bounded.
\end{itemize}
The first input is Proposition~\ref{4.21:prop-link-containment}. The second is Proposition~\ref{4.21:prop-base-point-profile}, which bounds $\operatorname{pol}(h\mathsf{u}(\mathsf{s})(s))$ at every law site, uniformly over the families on which the expansion runs.
Both propositions precede the holomorphy statement at the links in that order.

At the preparation site $(P)$, holomorphy is supplied instead by the complex form of the
preparation corollary, which precedes the preparation site in that order. This form is obtained
from the real one by replacing the variable $\lambda$ of the real statement by the triple
$(\lambda,\boldsymbol{\Phi},g)$ of variables of the complex one, and the size functional
$\mathfrak{N}$ by $\mathfrak{N}^{\mathbb{C}}$, with the same index sets $\mathcal{A}_c$,
$\mathcal{Q}^{w}_c$ and the same supports $\mathsf{b}(e)$, $\mathsf{d}(e)$; in this triple the
letter $g$ denotes the gauge variable of the complex corollary, not the reference coupling.
In the present setting the index set is $\mathcal{A}_c$ of Proposition~\ref{4.21:prop-preparation-machine}, and the supports are $\mathsf{b}(\hat e)$ and $\mathsf{d}^{\mathrm{rel}}(\hat e)$. The resummation by grains enters through the bound of Lemma~\ref{4.21:lem-grain-resummation}(ii). That bound holds with the parameters carried along, and therefore applies at the complexified parameters.

\emph{Numerical bounds.}
By Proposition~\ref{4.21:prop-majorant-structure}, every estimate of the induction step is a numerical consequence of the suprema of the stored terms it reads. Each such estimate reads a stored term through the supremum of one function over its datum. Fix $k<K$. Assume the relative format \eqref{4.21:eq-relative-format-a} at step $k$, together with the statement at step $k$ of the one induction in which the correlation theorem is proved. The relative format at step $k+1$ is reached through the following chain of implications.
\begin{enumerate}
\item By Proposition~\ref{4.21:prop-fluctuation-site}, the terms of level $k+1$ born from the
operation $\mathbf{T}$ in the passage from level $k$ to level $k+1$ are in the relative format
\eqref{4.21:eq-relative-format-a}, and the two subsequent steps of the machine at that site run
on them.
\item Apply Proposition~\ref{4.21:prop-preparation-machine}, together with
Lemma~\ref{4.21:lem-grain-resummation} and the hypothesis of
Theorem~\ref{4.21:thm-format-propagation} on the filing of stored terms. This gives clause~(a) of the
proposition on the
operation $\mathbf{R}$, read on labels, that is with each domain $X$ of a boundary-type term
replaced by its label $A_F$.
\item At this point of the ordering, the third-factor lemma at step $k+1$ is available.
\item Apply Proposition~\ref{4.21:prop-large-field-site}. The third-factor lemma, assumed in
Theorem~\ref{4.21:thm-format-propagation}, enters through Lemma~\ref{4.21:lem-deep-unit-sum}.
The hypothesis of Theorem~\ref{4.21:thm-format-propagation} for the plain family, stated there
together with its two accompanying conditions, enters as well.
The conclusion is clause~(e) of Proposition~\ref{4.21:prop-large-field-site}; this is the
hypothesis under which Proposition~\ref{4.21:prop-birth-site} is stated.
\item By Proposition~\ref{4.21:prop-birth-site}, clauses~(c)--(e) of the proposition on the
operation $\mathbf{R}$ hold. The boundary terms $\mathbf{B}'^{(k+1)}$ are in the relative format,
in the weighted norms $\mathfrak{N}^{A}$ and $\mathfrak{N}^{A,\mathbb{C}}$.
\item By clause~(b) of Proposition~\ref{4.21:prop-label-volume-sums}, the terms carried over pockets that are no longer live have ceased.
\end{enumerate}
These six steps together give the relative format \eqref{4.21:eq-relative-format-a} at step
$k+1$. The remaining estimates of the step read only the sizes of stored terms, not their
domains. These are clauses~(b), (d) and~(f) of the proposition on the operation $\mathbf{R}$,
and two further statements of the induction of the correlation theorem that bound sizes of
stored terms.
Their inputs are therefore supplied by the bounds just obtained, through
Proposition~\ref{4.21:prop-majorant-structure}.

\emph{Orbit data.}
The pieces of the decoupling expansion carry orbit data by
Lemma~\ref{4.21:lem-orbit-transfer}(iii). In the majorant stage described in
Proposition~\ref{4.21:prop-majorant-structure}, new functions are formed as limits of series. A
uniform limit of orbit data, under a majorant independent of the chart, is again an orbit datum.
The reason is Weierstrass's theorem: a locally uniform limit of holomorphic functions is
holomorphic. This theorem is used here in the way it is used in the complex preparation lemma
and in the preparation lemma of the preparation site. For terms created at the step, the orbit
data on the stored domain of \eqref{4.21:eq-orbit-data-a} are supplied by two statements:
clause~(b) of the proposition that equips terms with orbit data at their creation, and
clause~(ii) of the frame lemma.

\emph{Radii.}
The loss of radius at each reading of the response is bounded by
Corollary~\ref{4.21:cor-radius-loss}. The hypothesis $\eta<\varepsilon$ of
Lemma~\ref{4.21:lem-orbit-transfer} is obtained from this bound together with the lower bound
$\varepsilon_k\ge\tfrac34\mathsf{e}$ on the radii $\varepsilon_k$ of the charts of live terms.
Here $\mathsf{e}$ is the radius of the orbit data on the frames. This lower bound is clause~(j)
of the chart hypothesis of Theorem~\ref{4.21:thm-format-propagation}, the hypothesis from which
that theorem also derives the source-tube hypothesis for the class $\mathfrak{S}$.

For the charts of the units of the wrapped-descendant class $\mathfrak{S}$, clause~(j) alone
does not close the inequality $\eta<\varepsilon$; it is obtained from clause~(j) together with
three inputs:
\begin{itemize}
\item Proposition~\ref{4.21:prop-base-point-profile};
\item the second of the two conditions that make up the third-factor hypothesis of
Theorem~\ref{4.21:thm-format-propagation}, the first being its oscillation condition at plain
and wrapped units;
\item where a kept member of a live pocket is read, the source-tube hypothesis that
Theorem~\ref{4.21:thm-format-propagation} derives for the class $\mathfrak{S}$, in a
strengthened form.
\end{itemize}
With these inputs, $\eta<\varepsilon$ holds for the charts of the units of $\mathfrak{S}$ as well.
This completes the induction step and the proof.
\end{proof}

\begin{definition}[Second-class regions, their strips and the admissible geometry]
\label{4.21:def-strips-geometry}
The partitions $\pi_\ell$ into cubes, the distances $\operatorname{dist}_\ell$, the layer notation
$A^{(n)}$ and cells are those of Notation~\ref{4.21:not-cubes-cells-probes}. For every step $k$,
$\Lambda_k$ is
the small-field region of the step. The sets $\Omega_j$ and $Z_j=\Lambda_j^c$ are the domains of
the label, that is, of the admissible sequence of domains on which a term is written. The
\emph{zone of level $j$} is $\Omega_j\setminus\Omega_{j+1}$. We write $x_i=g_i^{-2}$; throughout,
$x_i\ge1$ for $i\le K$, as the powers $\mathsf{T}_k^{r_{\mathrm{pr}}}$ and
$(\log x)^{r_{\mathrm{pr}}}$ below presuppose. The
following objects are taken from the definition of the region size parameter and of the bounds
on the logarithmic scale variables:
\begin{itemize}
\item Ba{\l}aban's radius function $R(\cdot)$, written here as a function of the inverse squared
coupling $x=g'^{-2}$;
\item the number $N(g)$ of levels re-labelled by an operation $\mathbf{R}$ taken at coupling $g$;
\item the constants $C_N$ and $r_{\mathrm{pr}}>0$;
\item the constants $\lambda$ and $c_2$ of the coupling flow;
\item the logarithmic scale variables $\mathsf{T}_k=\log x_k$, $\check{\mathsf{T}}_k$ and
$\mathsf{T}_\gamma$; the variable $\check{\mathsf{T}}_k$ enters only through the bounds of
item~(7), which tie it to $\mathsf{T}_k$, $\check{R}_k$ and $\mathsf{T}_\gamma$.
\end{itemize}
Here $\mathsf{T}_\gamma\ge0$ is the number determined by the parameter $\gamma$ that bounds every
$\check{\mathsf{T}}_k$ from below (item~(7)).

\emph{Second-class regions and their strips.} A \emph{second-class region} is a pair $(Y,m)$ in
which $Y$ is a component of the second class created at the step $m$ by the operation
$\mathbf{R}_m$ in the label. In other words, $Y$ is a component of the set $\bigcup_i Y_i$ which
is adjoined to the large-field region at that step and on which the operation $\mathbf{T}_m$ has
the form \cite[(1.102)]{B-CMP122b}. The \emph{strip} of $(Y,m)$ is the closed box
\begin{equation}\label{4.21:eq-strips-geometry-1}
  Y^{\sharp}:=Y^{(m_{\sharp}(m))},\qquad
  m_{\sharp}(m):=\Bigl\lfloor \tfrac12\bigl(\check{R}_m-1\bigr)\Bigr\rfloor ,
\end{equation}
obtained by adding to $Y$ the stated number of layers of $\pi_m$-cubes. It is one set, and the
same set at every later level. We write $s_i(Y)$ for the largest side of $Y^{\sharp}$,
measured in the unit of level $i$.

Attempts of an operation $\mathbf{R}$, the interrupted and the excluded ones among them, their
localisation stage and the thin zones they leave are those of the lemma on arrested attempts.
For an attempt at step $k_a$ we write $h_a:=k_a-N(g_{k_a})$, similarly $h_b$, and
$h:=k-N(g_k)$ for the operation $\mathbf{R}_k$. The region $(Y,m)$ is \emph{live} until a later
step passes the
component of $\Lambda^c$ that contains $Y$ through the localisation stage of its operation
$\mathbf{R}$. This holds whatever class that component then takes. The region \emph{dies} at that
step. An attempt that is interrupted or excluded before its localisation stage does not end the
life of any second-class region.

\emph{The admissible geometry.} The following facts are in force at every step $k\le K$ and in
every label.
\begin{enumerate}
\item[(1)] The large-field region $Z$ treated by the operation $\mathbf{R}_k$ is a union of
components $C$ of
$\Lambda_k^c$. Distinct components of $\Lambda_k^c$ are at distance at least $\check{R}_k$, and
each component is a box with sides
\begin{equation}\label{4.21:eq-strips-geometry-a}
  a_\mu\,\check{R}_k,\qquad a_\mu\in\{1,\dots,100\}\quad(\mu=1,\dots,4).
\end{equation}
\item[(2)] The strips $C^{\sharp}:=C^{(m_{\sharp}(k))}$ of these components have the following
properties. They are boxes
with sides at most $101\check{R}_k-1$. They are pairwise at distance at least $1$. Each is at
distance at least $\check{R}_k-m_{\sharp}(k)$ from every other component of $\Lambda_k^c$.
Finally, $C^{\sharp}\setminus C\subset\Lambda_k$ and $C^{\sharp}\subset C^{\boxplus}$, where
$C^{\boxplus}$ is the one-layer enlargement of $C$ fixed in the statement on these strips.
\item[(3)] The region size parameter is $\check{R}_j=R\bigl(\min_{i\le j}x_i\bigr)$, where
Ba{\l}aban's radius function satisfies $R(x)<L\max\{1,(\log x)^{r_{\mathrm{pr}}}\}$. It is a
power of $L$, non-increasing in $j$, and
$\check{R}_j/\check{R}_{j+1}\in\{1,L\}$. Consequently, for all $m\le i$,
\begin{equation}\label{4.21:eq-strips-geometry-2}
  \check{R}_m\,L^{-(i-m)}\le \check{R}_i\le \check{R}_m ;
\end{equation}
and, by item (9), $\check{R}_j\ge L$ for $j\le K$.
\item[(4)] Four statements hold in every label \cite[(3.5), (3.20), p.~269]{B-CMP119}. Of the
domains of an admissible sequence, \cite{B-CMP119} states that ``the distance between their
boundaries is at least equal to $2MR_j$''; here $R_j$ is Ba{\l}aban's notation for the parameter
$\check{R}_j$ of item (3).
 \begin{itemize}
 \item[(a)] $\Omega_{j+1}\subset\Lambda_j\subset\Omega_j$. Hence
 $\Omega_j^c\subset Z_j=\Lambda_j^c\subset\Omega_{j+1}^c$.
 \item[(b)] The $\Omega_j$ are nested.
 \item[(c)] Cells that are neighbours across a wall or at a corner differ in level by at most $1$.
 \item[(d)] Every zone of level $j$ is at least one $\pi_{j+1}$-cube thick:
 \begin{equation}\label{4.21:eq-strips-geometry-3}
   \operatorname{dist}_{j+1}\bigl(\Omega_j^c,\Omega_{j+1}\bigr)\ge 1 .
 \end{equation}
 Consequently, a zone of level $j$ is at least $L$ cells of its own level thick, and at least
 $w$ units of length $M_1$ of its own level thick, where in this chapter $w:=M/M_1$.
 \end{itemize}
\item[(5)] Let $N_0$ be as in \cite[p.~179]{B-CMP122a}, put $k_0:=k-N_0$, and let
$h=k-N(g_k)$. Then
$L^{N_0-1}=\check{R}_{k_0+1}$. Every cube of the partition of which the enlarged domain
$\Omega_{k_0+1}^{\sim n}$ of \cite[p.~179]{B-CMP122a} is a
union is a $\pi_k$-cube. The operation $\mathbf{R}_k$ is taken with $N(g_k)>N_0$
\cite[p.~179]{B-CMP122a}; hence $h+1\le k_0$.
\item[(6)] Let two attempts, at steps $k_a<k_b$, act on components containing a common point,
and suppose the first
is interrupted or excluded. Then $k_a\le h_b=k_b-N(g_{k_b})$. Hence the ranges
$\left]h_a,k_a\right]$ and $\left]h_b,k_b\right]$ of re-labelled levels are disjoint. In a component
treated by $\mathbf{R}_k$, no earlier attempt re-labelled any of the levels $h+1,\dots,k$. The thin
zones left by an earlier interrupted or excluded attempt at step $k_a$ lie at levels
$\le k_a-1\le h-1$.
\item[(7)] For $k\le K$ and $0\le \mathsf{b}\le i$, with $s=i-\mathsf{b}$, one has
\begin{equation}\label{4.21:eq-strips-geometry-4}
  \begin{gathered}
  N(g_k)\le C_N\,\mathsf{T}_k^{r_{\mathrm{pr}}},\qquad
  \check{R}_k\ge \check{\mathsf{T}}_k^{\,r_{\mathrm{pr}}},\qquad
  \mathsf{T}_\gamma\le\check{\mathsf{T}}_k,\qquad
  \mathsf{T}_k\le\check{\mathsf{T}}_k+1,\\
  \mathsf{T}_{\mathsf{b}}\le\mathsf{T}_i+c_s,\qquad c_s:=\log\bigl(1+3\lambda s+2c_2\bigr).
  \end{gathered}
\end{equation}
\item[(8)] Define
\begin{equation}\label{4.21:eq-strips-geometry-5}
  \mathsf{R}^{\uparrow}(\mathsf{T}):=
  L\Bigl(\mathsf{T}+1+\log\bigl(1+3\lambda C_N(\mathsf{T}+1)^{r_{\mathrm{pr}}}+2c_2\bigr)
  \Bigr)^{r_{\mathrm{pr}}} .
\end{equation}
Then
\begin{equation}\label{4.21:eq-strips-geometry-6}
  N(g_k)\le C_N2^{r_{\mathrm{pr}}}\check{R}_k,\qquad
  \mathsf{T}_k^{r_{\mathrm{pr}}}\le 2^{r_{\mathrm{pr}}}\check{R}_k,\qquad
  \check{R}_{h+1}\le\mathsf{R}^{\uparrow}(\check{\mathsf{T}}_k).
\end{equation}
\item[(9)] The parameter $\gamma$ satisfies $\mathsf{T}_\gamma^{r_{\mathrm{pr}}}\ge 38$. Hence
$\check{R}_j\ge 38$ for $j\le K$.
\item[(10)] Put
\begin{equation}\label{4.21:eq-strips-geometry-7}
  \tau_d(t):=\sum_{l\ge1}(2l+215)^d\,e^{-(lt-5)} .
\end{equation}
The parameter $\gamma$ satisfies the condition, depending on $d$ only, that
$4\tau_d(t)\le1$ for all $t\ge\mathsf{T}_\gamma^{r_{\mathrm{pr}}}$. This holds if and only if
$4\tau_d\bigl(\mathsf{T}_\gamma^{r_{\mathrm{pr}}}\bigr)\le1$.
\end{enumerate}
\end{definition}

\emph{Which items are proved here.} Items (1) and (2) are the statements on the geometry of the
large-field region of $\mathbf{R}_k$ and on the strips of its components. Item (3), apart from
its consequences in \eqref{4.21:eq-strips-geometry-2} and the bound $\check{R}_j\ge L$, is the
definition of the region size
parameter, and item (7) is the set of bounds on the logarithmic scale variables. Item (6), apart
from its last two sentences, is the lemma on arrested attempts; compare \cite{B-CMP122a}. The
inequality $N(g_k)>N_0$ of item (5) is the assumption made in \cite[p.~179]{B-CMP122a}. These
are used here as stated there. Item (4)(c) and the first inclusions of (4)(a) are
taken from \cite[(3.5), (3.20)]{B-CMP119}. The remaining assertions are proved now.

\begin{proof}
\emph{The consequences in item (3).} Let $m\le i$. Write $\check{R}_m/\check{R}_i$ as the
telescoping product
$\prod_{j=m}^{i-1}\check{R}_j/\check{R}_{j+1}$. It has $i-m$ factors, each equal to $1$ or $L$.
Its value therefore lies between $1$ and $L^{i-m}$. This is
\eqref{4.21:eq-strips-geometry-2}.

By item (9), proved below, $\check{R}_j\ge38>1$ for $j\le K$. Since $\check{R}_j$ is a power of
$L$ exceeding $1$, its exponent is at least $1$, so $\check{R}_j\ge L$.

\emph{Item (9).} Let $j\le K$. By item (7), $\check{R}_j\ge\check{\mathsf{T}}_j^{\,r_{\mathrm{pr}}}$
and $\check{\mathsf{T}}_j\ge\mathsf{T}_\gamma$. Since $\mathsf{T}_\gamma\ge0$, as stated at the
beginning of this definition,
and $r_{\mathrm{pr}}>0$, the map $t\mapsto t^{r_{\mathrm{pr}}}$ is increasing on $t\ge0$. Hence
\begin{equation*}
  \check{R}_j\ge\check{\mathsf{T}}_j^{\,r_{\mathrm{pr}}}
  \ge\mathsf{T}_\gamma^{r_{\mathrm{pr}}}\ge 38 .
\end{equation*}

\emph{Item (4)(a), second chain.} Take complements in $\Omega_{j+1}\subset\Lambda_j\subset\Omega_j$.
This gives $\Omega_j^c\subset\Lambda_j^c\subset\Omega_{j+1}^c$, and $\Lambda_j^c=Z_j$ by
definition.

\emph{Item (4)(b).} This follows from (4)(a), since $\Omega_{j+1}\subset\Lambda_j\subset\Omega_j$.

\emph{Item (4)(d).} We treat the zones according to the operation that made them.

A zone made by the $\mathbf{T}$-step has thickness $8\check{R}_{j+1}$ in the unit of
$\pi_{j+1}$, by the statement on the geometry of the $\mathbf{T}$-step. This is at least $1$.

Next consider a zone of the prepared label of $\mathbf{R}_k$, as defined in
\cite[p.~179, (1.10)--(1.12)]{B-CMP122a}. Let $h\le j<k_0$. There one has, with $Z'_j$ and the
enlargements $(\cdot)^{\sim}$, $(\cdot)^{\sim3}$ as in \cite[(1.10)]{B-CMP122a},
\begin{equation*}
  (\Omega''_{j+1})^c\cap Z=Z''_{j+1}=(Z'_j)^{\sim3}.
\end{equation*}
By \cite[(1.11), (1.12)]{B-CMP122a}, $(\Omega''_j)^c\subset\Omega_j^c$, and by (4)(a),
$\Omega_j^c\subset Z_j$. The two sets are therefore separated by
the three layers of $LM\check{R}_{j+1}$-cubes added in forming $(Z'_j)^{\sim3}$. Measured in the
unit of $\pi_{j+1}$, these layers have thickness $3\check{R}_{j+1}\ge1$.

For $j=k_0$ one has $Z''_{k_0+1}=(Z'_{k_0})^{\sim}$ and $Z''_{k_0}\subset Z_{k_0}$. One such layer
separates them, of thickness $\check{R}_{k_0+1}\ge1$.

For $k_0<j<k$, the complements in \cite[p.~179]{B-CMP122a} are completed to an admissible
sequence based on $M$-cubes at the corresponding scales. Thus $(Z''_{j+1})^c$ and $Z''_j$ are
separated by one layer of $L^{-(k-j-1)}M$-cubes, which are the $\pi_{j+1}$-cubes. The thickness
is therefore $\ge1$.

Finally, at a running or arrested stage of the integrations the zone of level $j$ is
$\Omega^c_{j+1}\setminus Z''_j$. By the description of the stage sets in
\cite[p.~180]{B-CMP122a}, it is thicker still.

This proves \eqref{4.21:eq-strips-geometry-3} in every case. Since the cubes are nested
$L$-adically, one $\pi_{j+1}$-cube has side $L$ in the unit of level $j$. A thickness
$\ge1$ in the unit of $\pi_{j+1}$ is thus a thickness $\ge L$ cells of level $j$. Each of these
cells has side $M=wM_1$ in lattice spacings of level $j$, so the thickness is at least
$Lw\ge w$ units of length $M_1$ of level $j$. Only the lower bound $1$ of
\eqref{4.21:eq-strips-geometry-3} is used below.

\emph{Item (5).} By \cite[p.~179]{B-CMP122a}, $N_0$ is the smallest positive integer with
$L^{-N_0+1}MR_{k-N_0+1}=M$, where Ba{\l}aban's $R_j$ is the parameter $\check{R}_j$ of
item (3). Hence $\check{R}_{k_0+1}=L^{N_0-1}$, since $k_0+1=k-N_0+1$.

By the same passage, the domains $\Omega_j^{\sim n}$ are unions of $L^{-(k-j)}M\check{R}_j$-cubes
of the finest lattice. For $j=k_0+1$ the side is
\begin{equation*}
  L^{-(N_0-1)}M\check{R}_{k_0+1}=M ,
\end{equation*}
so the cubes are the $M$-cubes of level $k$, that is, $\pi_k$-cubes.

The inequality $N(g_k)>N_0$ is assumed in \cite[p.~179]{B-CMP122a}, where $N(g_k)$ is written
$N$. Since both are integers, $N(g_k)\ge N_0+1$. Therefore
$h+1=k-N(g_k)+1\le k-N_0=k_0$.

\emph{The last two sentences of item (6).} Let the component be treated by $\mathbf{R}_k$, with
$h=k-N(g_k)$. Take an earlier interrupted or excluded attempt at step $k_a<k$ on a component
containing a point of it. The first part of item (6), with $k_b=k$, gives $k_a\le h$. Its range
$\left]h_a,k_a\right]$ therefore lies in $\left]-\infty,h\right]$, which is disjoint from
$\left]h,k\right]$. So none of the levels $h+1,\dots,k$ was re-labelled by it. By the lemma on
arrested attempts, its thin zones lie at levels $\le k_a-1$, and $k_a-1\le h-1$.

\emph{Item (8).} From item (7) and item (9) we have
$\check{\mathsf{T}}_k\ge\mathsf{T}_\gamma>1$, because $\mathsf{T}_\gamma\ge0$ and
$\mathsf{T}_\gamma^{r_{\mathrm{pr}}}\ge38>1$.
Hence $\mathsf{T}_k\le\check{\mathsf{T}}_k+1\le2\check{\mathsf{T}}_k$. Since $x_k\ge1$, we have
$\mathsf{T}_k=\log x_k\ge0$. Raising to the power
$r_{\mathrm{pr}}$ and using $\check{\mathsf{T}}_k^{\,r_{\mathrm{pr}}}\le\check{R}_k$ gives
\begin{equation*}
  \mathsf{T}_k^{r_{\mathrm{pr}}}\le2^{r_{\mathrm{pr}}}\check{\mathsf{T}}_k^{\,r_{\mathrm{pr}}}
  \le2^{r_{\mathrm{pr}}}\check{R}_k .
\end{equation*}
Then $N(g_k)\le C_N\mathsf{T}_k^{r_{\mathrm{pr}}}\le C_N2^{r_{\mathrm{pr}}}\check{R}_k$. These are
the first two bounds of \eqref{4.21:eq-strips-geometry-6}.

For the third bound, put $y=\min_{i\le h+1}x_i$; then $1\le y\le x_{h+1}$, since every
$x_i\ge1$. By item (3), Ba{\l}aban's radius
function satisfies $R(x)<L\max\{1,(\log x)^{r_{\mathrm{pr}}}\}$. The right side is non-decreasing
in $x\ge1$. Hence
\begin{equation*}
  \check{R}_{h+1}=R(y)<L\max\bigl\{1,(\log y)^{r_{\mathrm{pr}}}\bigr\}
  \le L\max\bigl\{1,(\log x_{h+1})^{r_{\mathrm{pr}}}\bigr\}.
\end{equation*}
Apply the last inequality of \eqref{4.21:eq-strips-geometry-4} with $\mathsf{b}=h+1$, $i=k$ and
$s=k-h-1=N(g_k)-1$, which is non-negative by item (5). Since $\mathsf{T}_{h+1}=\log x_{h+1}$,
$s<N(g_k)$ and the logarithm is increasing, this gives
\begin{equation*}
  \log x_{h+1}\le\mathsf{T}_k+\log\bigl(1+3\lambda N(g_k)+2c_2\bigr).
\end{equation*}
Apply item (7) twice: $\mathsf{T}_k\le\check{\mathsf{T}}_k+1$, and
$N(g_k)\le C_N\mathsf{T}_k^{r_{\mathrm{pr}}}\le C_N(\check{\mathsf{T}}_k+1)^{r_{\mathrm{pr}}}$.
Since the logarithm is increasing,
\begin{equation*}
  \log x_{h+1}\le\check{\mathsf{T}}_k+1
  +\log\bigl(1+3\lambda C_N(\check{\mathsf{T}}_k+1)^{r_{\mathrm{pr}}}+2c_2\bigr).
\end{equation*}
The right side is at least $\check{\mathsf{T}}_k+1\ge1$, so the maximum with $1$ is attained by
its $r_{\mathrm{pr}}$-th power. This power is $L^{-1}\mathsf{R}^{\uparrow}(\check{\mathsf{T}}_k)$
by \eqref{4.21:eq-strips-geometry-5}. Hence
$\check{R}_{h+1}\le\mathsf{R}^{\uparrow}(\check{\mathsf{T}}_k)$.

\emph{Item (10).} For each $l\ge1$ the summand $(2l+215)^de^{-(lt-5)}$ of
\eqref{4.21:eq-strips-geometry-7} is strictly decreasing in $t$. Hence $\tau_d$ is
non-increasing. The inequality $4\tau_d(t)\le1$ for all $t\ge\mathsf{T}_\gamma^{r_{\mathrm{pr}}}$
is therefore equivalent to the single inequality at $t=\mathsf{T}_\gamma^{r_{\mathrm{pr}}}$. Only
the consequence $\tau_d(t)\le\frac14$ is used.
\end{proof}

\begin{remark}\label{4.21:lem-live-region-geometry}
The statement of this lemma is not written out here; parts of its proof are printed below.
\end{remark}

\begin{proof}[Proof of Lemma~\ref{4.21:lem-live-region-geometry}, concluded]
\label{4.21:prf-live-region-proof}
It remains to prove clauses (e)--(i). As in \cite{B-CMP122a}, $Z$ denotes the large-field
region at step $k$, $C$ a component of it, and $k_0$ the index entering the definition
\cite[(1.10)]{B-CMP122a} of the new large-field region, which is related to $k$ and $N_0$ by
(N$_0$) in~\eqref{4.21:eq-strips-geometry-a}.

\smallskip
\noindent Clause (e). Three properties of the strip are asserted: it is connected, it is
contained in $Z_k$, and it contains $Y$. All three are part of condition (G0)
in~\eqref{4.21:eq-strips-geometry-a}.

\smallskip
\noindent Clause (f). Following \cite[p.~192]{B-CMP122a}, the new large-field region is
\begin{equation*}
Z''_k=\bigl(\Omega^{\sim 5}_{k_0+1}\bigr)^c\cap Z .
\end{equation*}
By the definition of the index $k_0$ it is therefore a union of $M$-cubes, and it is a
rectangular parallelepiped. The domain
\begin{equation*}
\Lambda=\bigl(\Omega^{\sim 4}_{k_0+1}\bigr)^c\cap Z
\end{equation*}
of \cite[(1.73)]{B-CMP122a} is obtained by adding one layer of $M$-cubes to $Z''_k$, so it is
again a union of $M$-cubes. We write $\Lambda_C=\Lambda\cap C$ for its part in the component
$C$, and $\Lambda_C^{\sim}$ for the one-layer enlargement of $\Lambda_C$. A component of the
large-field region satisfies the two conditions stated in \cite{B-CMP122a}:
\begin{quote}
(i) it is contained in a cube of the size $100MR_k$, (ii) in the preceding $N$ renormalization
steps no new large field regions were created inside this component, and the previous regions
contained in it satisfy the condition (i) on the corresponding scales
\end{quote}
In the notation of this chapter a cube of the size $100MR_k$ is a cube of side
$100\check{R}_k$ units of level $k$, and $N$ is the number of steps so denoted in
\cite{B-CMP122a}, not the rank of the gauge group. We refer to these two conditions as
Ba{\l}aban's conditions (i) and (ii), to distinguish them from the clauses of the lemma.
By Ba{\l}aban's condition (i), $\Lambda$ is a rectangular parallelepiped contained in a cube of
size $100M$.

We bound the size of $\Lambda$ in two ways.

$(\alpha)$ The first way is the one of \cite{B-CMP122a}. We have
\begin{equation*}
\Lambda\subset\Omega_{k_0+1}^c\cap C\subset\Lambda_{k_0+1}^c\cap C=Z_{k_0+1}\cap C .
\end{equation*}
By Ba{\l}aban's condition (ii), this last set satisfies condition (i) on the
corresponding scale. Hence it lies in a cube of side $100\check{R}_{k_0+1}$ units of level
$k_0+1$. By the relation (N$_0$) in~\eqref{4.21:eq-strips-geometry-a}, this side equals $100$
units of level $k$.

$(\beta)$ The second way uses condition (i) at the scale $k_0$ only. Since
$\check{R}_{k_0}\le L\check{R}_{k_0+1}=L^{N_0}$, the set $Z_{k_0}\cap C$ lies in a cube of side
\begin{equation*}
100\,\check{R}_{k_0}L^{-N_0}\le 100
\end{equation*}
units of level $k$. The set $Z'_{k_0}$ adds less than $1$ on each face. Next,
$\Lambda=(Z'_{k_0})^{\sim 4}$ adds four more on each face. Hence every side of $\Lambda$ is at
most $100+2+8=110$.

Passing from $\Lambda_C$ to $\Lambda_C^{\sim}$ adds one unit on each face, hence $2$ to every
side length. Under either bound, therefore, every side of $\Lambda^{\sim}_C$ is at most $112$.

For the distances, note first that $N_0\ge 2$, because $\check{R}_{k_0+1}=L^{N_0-1}>1$. Hence
\begin{equation*}
\Lambda\subset\Omega_{k_0+1}^c\subset Z_{k_0+1}\subset Z_{k-1}.
\end{equation*}
The set $\Omega_k^c$ contains $(Z'_{k-1})^{\sim 8}$, and this accounts for $8\check{R}_k$ cubes
of $\pi_k$. A further $2\check{R}_k$ comes from \cite{B-CMP122a}, which states:
\begin{quote}
Each renormalization step adds at least ten layers of $MR_k$-cubes
\end{quote}
Therefore
\begin{equation*}
\operatorname{dist}_k\bigl(\Lambda,Z^c\bigr)\ge 10\check{R}_k,
\qquad
\operatorname{dist}_k\bigl(\Lambda^{\sim},Z^c\bigr)\ge 10\check{R}_k-1 .
\end{equation*}
Combined with (G0) in~\eqref{4.21:eq-strips-geometry-a}, this shows that the sets
$\Lambda^{\sim}_C$ and $\Lambda^{\sim}_{C'}$ of two distinct components $C\ne C'$ are at distance
at least $21\check{R}_k-2$. Since $\check{R}_k\ge1$, these bounds are at least the bounds
$9\check{R}_k$ and $19\check{R}_k$ asserted in clause (f).

Finally, off $Z''_k$ the set $\mathbb{B}''_k$ of \cite[(1.12)--(1.13)]{B-CMP122a} is of level
$k$ in $Z$.

\smallskip
\noindent Clause (g). A component of the large-field region satisfies Ba{\l}aban's conditions
(i) and (ii), quoted in the proof of clause (f).
Moreover, by \cite{B-CMP122a}, the one-step operations of the preceding steps counted in
Ba{\l}aban's condition (ii) carry no large-field characteristic functions in such a component.
A component thrown back at level $m$ carries the factor $\chi_m^c(Y^{\sim -6})$ there, in
Ba{\l}aban's construction. Since those operations carry no such factor, $m\le h$.

By the definition of the core, taken from the passage preceding \cite[(1.81)]{B-CMP122a} on
p.~195, we have
$\mathsf{C}_C=(Z'_h)^{\sim}\cap C$. This is a rectangular parallelepiped contained in
$Z''_{h+1}=(Z'_h)^{\sim 3}$. We bound its side. By Ba{\l}aban's condition (ii), $Z_h\cap C$ lies
in a cube of side $100\check{R}_h$. The set $Z'_h$ adds less than $L\check{R}_{h+1}$ on each
side, and the operation ${\sim}$ adds one layer of width $L\check{R}_{h+1}$. Since
$\check{R}_h\le L\check{R}_{h+1}$, every side of $\mathsf{C}_C$ is therefore at most
\begin{equation*}
100\check{R}_h+4L\check{R}_{h+1}\le 104\,L\check{R}_{h+1}.
\end{equation*}

Since $\mathsf{C}_C\subset(\Omega''_{h+1})^c$, the cells of $\mathsf{C}_C$ have level at most
$h$. Consider the moat $(Z'_h)^{\sim}\setminus Z'_h$. It lies outside $Z'_h\supset Z_h$, hence
in $Z_h^c$, and inside $(Z'_h)^{\sim 3}=Z''_{h+1}$. By the definition of the domains of the
labels in~\eqref{4.21:eq-strips-geometry-a}, $Z_h^c=\Lambda_h$; and
$Z''_{h+1}=(\Omega''_{h+1})^c$ by \cite[p.~195]{B-CMP122a}. By (A0)
in~\eqref{4.21:eq-strips-geometry-a}, which gives $\Lambda_h\subset\Omega_h$, we obtain
\begin{equation*}
(Z'_h)^{\sim}\setminus Z'_h
\subset Z_h^c\cap Z''_{h+1}
\subset \Lambda_h\cap(\Omega''_{h+1})^c
\subset \Omega_h\setminus\Omega''_{h+1}.
\end{equation*}
There the level is $h$. Consequently the boundary $\partial\mathsf{C}_C$ runs along faces of
$LM\check{R}_{h+1}$-cubes inside zone $h$.

The assertion of clause (g) on bonds is the convention on bonds stated after
\cite[(1.81)]{B-CMP122a}.

Now let $(Y,m)$ be a pocket in $C$.
\begin{itemize}
\item If $m<h$, then $Y^{\sharp}\subset Z_{m+1}\subset Z_h$. Hence $Y^{\sharp}$ lies at distance
at least $L\check{R}_{h+1}\ge\check{R}_h$ from $\mathsf{C}_C^c$.
\item If $m=h$, then $Y\subset Z_h$, and $Y^{\sharp}$ reaches at most $\tfrac12(\check{R}_h-1)$
beyond $Z_h$. The distance from $\mathsf{C}_C^c$ is therefore at least
$\check{R}_h-\tfrac12(\check{R}_h-1)$.
\end{itemize}
Finally, let a cube of side at most $1$ in level-$k$ units touch $Y^{\sharp}\subset Z''_k$. Such
a cube lies in $(Z''_k)^{(1)}=\Lambda_C$.

\smallskip
\noindent Clause (h). We have $Z''_{h+1}=(Z'_h)^{\sim 3}$, and the sets $Z''_j$ increase with
$j$. Hence every collar lies at distance at least $3L\check{R}_{h+1}\ge 3\check{R}_h$ level-$h$
units from $Z_h\supset Y$. In level-$(h+1)$ units this distance is at least $3\check{R}_{h+1}$.
On the other hand, $Y^{\sharp}$ reaches beyond $Z_h$ by at most
$\tfrac12\check{R}_h/L\le\tfrac12\check{R}_{h+1}$ units of level $h+1$. In the second case of
clause (h), the assertion follows from clause (e) together with the inclusion $R_\nu\subset Z$.

\smallskip
\noindent Clause (i). We prove its four parts in turn.

(i1) Let $n:=\operatorname{lev}(\mathsf{Q})>m$. The cell $\Delta_n$ of the probe lies in zone
$n$, which is contained in $\Omega_n$. By clause (c$'$),
\begin{equation*}
\operatorname{dist}_n(\Delta_n,Y^{\sharp})\ge\tfrac52>2,
\end{equation*}
and hence $\mathsf{Q}\cap Y^{\sharp}=\emptyset$.

(i2) By (i1), a strip met by $\mathsf{Q}$ belongs to a pocket of level $m\ge n$. The side of
$\mathsf{Q}$ is at most $5$ level-$n$ units; measured in that strip's own units it is therefore
at most $5$, and at most $5/L\le 1$ if $m>n$. By clause (a), the strip has all sides at least
$38$, which exceeds $5$. By (g7)
in~\eqref{4.21:eq-tree-length-facts-a}, such a strip contains a vertex of $\mathsf{Q}$. By
clause (b), distinct strips contain distinct vertices.

(i3) A point of a block of zone $n$ lies within $1$ level-$n$ unit of the block's own cell
$\Delta_n\subset\Omega_n\setminus\Omega_{n+1}$. Apply (A3) of~\eqref{4.21:eq-strips-geometry-a}
first at $j=n-1$: the set $\Omega_{n-1}^c$ is at distance at least $1$ level-$n$ unit from
$\Omega_n$. Hence no interior point of such a block lies in a zone of index at most $n-2$. Apply
(A3) next at $j=n+1$: the set $\Omega_{n+2}$ is at distance at least $1$ level-$(n+2)$ unit,
that is $L^2$ level-$n$ units, from $\Omega_{n+1}^c$, which contains $\Delta_n$.

Consequently a cell $\Delta$ of level $n'$ meets blocks of the zones $n'-1$, $n'$, $n'+1$ only.
Here ``meet'' means having a common interior point.
\begin{itemize}
\item The blocks of zone $n'$ contain $\Delta$; there are $2^d$ of them.
\item The blocks of zone $n'+1$ contain the parent of $\Delta$; again there are $2^d$ of them.
\item The blocks of zone $n'-1$ are cubes of side $2$ on the $\pi_{n'-1}$-grid that contain a
sub-cube of $\Delta$. Such a block has $L+1$ positions along each axis, so there are at most
$(L+1)^d$ of them.
\end{itemize}

(i4) By (A0), $\Omega_h^c\subset\Lambda_h^c=Z_h$. Hence a cell of level less than $h$ lies in
$\Omega_h^c\subset Z_h$. Now let $n:=\operatorname{lev}(\mathsf{Q})>h$. Then, for the current
label, $\Delta_n\subset\Omega_n$. We claim
\begin{equation*}
\operatorname{dist}_n\bigl(Z_h\cap C,\Omega_n\bigr)\ge\tfrac52 .
\end{equation*}
This follows from the computation in the proof of clause (c$'$), with $Z_h$ in place of
$Y^{\sharp}$. If $n\le k_0$, then $Z_h\subset Z_{n-1}$, and the three layers give distance
$3\check{R}_n$. If $n>k_0$, then $Z_h\subset Z_{k_0}$, and the distance is at least $3$.
\end{proof}

\noindent\emph{Remark on the proof.}
In every comparison made in the proof above, both members scale alike. Neither the age $i-m$,
nor the number of preceding steps counted in Ba{\l}aban's condition (ii), nor $N_0$, nor any
number of steps or of pockets occurs. Clause (c) asserts a collar of
width $\tfrac52\check{R}_m$ only. The collar made by the operation $\mathbf{T}$ is wider, but
after a preparation interrupted at level $m$ only $\tfrac52\check{R}_m$ of it is left, and the
lemma claims no more. A late interruption leaves at most $N_0$ thin strata. These are handled by
(A3) of~\eqref{4.21:eq-strips-geometry-a} and by clause (c$'$), never by a claim on their
thickness.

\begin{lemma}[Sums over cells and blocks against the scaled distance]\label{4.21:lem-cell-block-sums}
Let $\mathfrak{B}$ be a determining set, with its cells, zones and blocks as in
Notation~\ref{4.21:not-cubes-cells-probes}, and put $w:=M/M_1$. For each level $i$ let $\ell_i$
denote the $M_1$-unit of level $i$, so that an $M$-cube of the own scale of level $i$ has side $w$
in that unit. For a contour $\Gamma$ put
\begin{equation*}
\operatorname{len}_{\mathfrak{B}}(\Gamma):=\sum_i \ell_i^{-1}\,\bigl|\Gamma\cap(\text{zone } i)\bigr|,
\end{equation*}
where $|\cdot|$ is Euclidean length, and for sets $A,B$ let the \emph{scaled distance}
$d^{\wedge}(A,B)=d^{\wedge}_{\mathfrak{B}}(A,B)$ be the infimum of $\operatorname{len}_{\mathfrak{B}}$
over all contours joining $A$ to $B$, whatever zones they cross. Put
\begin{equation*}
c_{\boxplus}:=3^d(L+5),\qquad \hat{b}_t:=(L+2)^d+3^d,\qquad
C_{\mathrm{cell}}:=4\,\hat{b}_t^{\,2c_{\boxplus}} .
\end{equation*}
For a closed union $\Sigma$ of cells let $\partial^+\Sigma$ be the set of cells not in $\Sigma$ that
touch $\Sigma$, and $\#\Sigma$ the number of cells of $\Sigma$. Let $C_{\mathrm{blk}}$ be the constant
of the corresponding display. Then the following hold.
\begin{enumerate}
\item[(a)] (A short contour meets few cells.) Let $x\ge 0$. A contour $\Gamma$ with
$\operatorname{len}_{\mathfrak{B}}(\Gamma)\le xw$ meets at most $c_{\boxplus}(x+1)$ cells.
\item[(b)] If $\beta w\ge c_{\boxplus}\log\hat{b}_t+1$, then for every closed union $\Sigma$ of cells
\begin{equation}\label{4.21:eq-cell-block-sums-1}
\sum_{\Delta\not\subset\Sigma} e^{-\beta d^{\wedge}(\Delta,\Sigma)}\le C_{\mathrm{cell}}\,\#\partial^+\Sigma ,
\end{equation}
the sum running over cells $\Delta$.
\item[(c)] For the same $\beta$ and $\Sigma$, the sum over all blocks $\square$ of the cover satisfies
\begin{equation}\label{4.21:eq-cell-block-sums-2}
\sum_{\square} e^{-\beta d^{\wedge}(\square,\Sigma)}\le C_{\mathrm{blk}}\bigl(1+\hat{b}_t C_{\mathrm{cell}}\bigr)\,\#\Sigma .
\end{equation}
\item[(d)] Suppose that $w$ satisfies the lower bound (imposed on $w=M/M_1$ after $L$ is fixed)
\begin{equation}\label{4.21:eq-cell-block-sums-a}
\tfrac{1}{16}\,\delta_P\, w\ \ge\ c_{\boxplus}\log\hat{b}_t+1 ,
\end{equation}
where $\delta_P=\tfrac18\delta_0$ is the decay rate attached to the scaled distance. Then
\eqref{4.21:eq-cell-block-sums-1} and \eqref{4.21:eq-cell-block-sums-2} hold at every
$\beta\ge\delta_P/16$, and $\tfrac1{16}\delta_P w\ge 1$.
\end{enumerate}
No level, age, number or thickness of zones enters these bounds.
\end{lemma}

\begin{proof}
We use the following four properties of the scaled distance $d^{\wedge}_{\mathfrak{B}}$, which are the
content of the lemma of this chapter on the geometry of the scaled distance of a determining set (the
scaled distance is measured in $M_1$-cubes on the corresponding scales, as in
\cite[p.~186]{B-CMP122a}):
\begin{enumerate}
\item[(i)] $d^{\wedge}$ obeys the triangle inequality;
\item[(ii)] a contour pays at least $w$ for every zone it crosses;
\item[(iii)] on a set all of whose cells have level at most $n$, $\operatorname{len}_{\mathfrak{B}}$ is at
least $w$ times the Euclidean length measured in level-$n$ units;
\item[(iv)] the payments of $\operatorname{len}_{\mathfrak{B}}$ on disjoint arcs add.
\end{enumerate}

(a) The scaled length of the initial arc of $\Gamma$ up to a point is a nondecreasing, continuous
function of that point, and by (iv) the scaled lengths of consecutive arcs add. Hence $\Gamma$ can be
cut into at most $\max\{1,\lceil 2x\rceil\}\le 2x+1$ consecutive arcs, each of scaled length at most
$\tfrac12 w$. Fix one such arc $\gamma$ and let $n$ be the highest level met by $\gamma$. By (ii),
crossing a zone costs at least $w$, while $\operatorname{len}_{\mathfrak{B}}(\gamma)\le\tfrac12 w$;
hence $\gamma$ crosses no zone and meets the zones $n$ and $n-1$ only. All units met by $\gamma$ are
therefore at most $\ell_n$, all cells met have level at most $n$, and (iii) gives that the Euclidean
length of $\gamma$, in level-$n$ units, is at most
$\operatorname{len}_{\mathfrak{B}}(\gamma)/w\le\tfrac12$. Thus $|\gamma|$ is at most half the side of a
level-$n$ cube, which, the cubes being nested $L$-adically, equals $\tfrac12 L$ sides of a level-$(n-1)$
cube.

A curve of Euclidean length $t$, measured in units of the side of the cubes of a cubic grid, meets at
most $3^d(t+1)$ cubes of that grid: it is the union of at most $\lfloor t\rfloor+1\le t+1$ consecutive
pieces of length less than $1$, and a piece of length less than $1$ starting in a cube lies in the union
of the $3^d$ cubes touching that cube (the cube itself included) and meets no other cube. Applied with
$t=\tfrac12$ on the grid of level $n$ and with
$t=\tfrac12 L$ on the grid of level $n-1$, this shows that $\gamma$ meets at most $3^d\cdot\tfrac32$
cubes of level $n$ and at most $3^d(\tfrac12 L+1)$ cubes of level $n-1$; these include all cells met
by $\gamma$. Summing over the arcs,
\begin{equation*}
\#\{\text{cells met by }\Gamma\}\le(2x+1)\,3^d\bigl(\tfrac12L+\tfrac52\bigr)
=3^d(L+5)\bigl(x+\tfrac12\bigr)\le c_{\boxplus}(x+1).
\end{equation*}

(b) The cells met by a contour form a family connected by touching: otherwise the family splits into
two nonempty subfamilies, no cell of one touching a cell of the other, and their unions would be two
disjoint closed sets covering the contour and each meeting it, which is impossible since a connected
set is not covered by two disjoint closed sets each meeting it.

Let $x\ge0$ and let $\Delta\not\subset\Sigma$ be a cell with $d^{\wedge}(\Delta,\Sigma)<xw$. Choose a
contour $\Gamma$ from $\Sigma$ to $\Delta$ with $\operatorname{len}_{\mathfrak{B}}(\Gamma)<xw$, and let
$\Gamma'$ be its part after its last point in $\Sigma$ (the last point exists because $\Sigma$ is
closed). By (iv), $\operatorname{len}_{\mathfrak{B}}(\Gamma')<xw$, so by (a) $\Gamma'$ meets at most
$c_{\boxplus}(x+1)$ cells; these form a family connected by touching, which contains $\Delta$ and, since
$\Gamma'$ starts at a point of $\Sigma$ and then leaves $\Sigma$, a cell of $\partial^+\Sigma$. Hence
$\Delta$ is joined to a cell of $\partial^+\Sigma$ by a chain of at most $c_{\boxplus}(x+1)$ cells,
consecutive ones touching. By the admissibility condition of the cell geometry in
Definition~\ref{4.21:def-strips-geometry}, a cell touches at most $3^d$ cells
of its own or of the coarser level and at most $(L+2)^d$ of the finer level, hence at most $\hat{b}_t$
others; so at most $\hat{b}_t^{\,n}$ cells are reached from a given cell by chains of $n$ steps. Since
$\hat{b}_t\ge2$,
\begin{equation*}
\#\{\Delta:\ d^{\wedge}(\Delta,\Sigma)<xw\}
\le\#\partial^+\Sigma\sum_{0\le n\le c_{\boxplus}(x+1)}\hat{b}_t^{\,n}
\le 2\,\#\partial^+\Sigma\;\hat{b}_t^{\,c_{\boxplus}(x+1)} .
\end{equation*}
For $\Delta\not\subset\Sigma$ let $x=\lfloor d^{\wedge}(\Delta,\Sigma)/w\rfloor$, an integer $\ge0$; then
$d^{\wedge}(\Delta,\Sigma)<(x+1)w$ and $e^{-\beta d^{\wedge}(\Delta,\Sigma)}\le e^{-\beta wx}$. Therefore,
applying the last display with $x+1$ in place of $x$,
\begin{align*}
\sum_{\Delta\not\subset\Sigma}e^{-\beta d^{\wedge}(\Delta,\Sigma)}
&\le\sum_{x\ge0}\#\{\Delta:\ d^{\wedge}(\Delta,\Sigma)<(x+1)w\}\,e^{-\beta wx}\\
&\le 2\,\#\partial^+\Sigma\;\hat{b}_t^{\,2c_{\boxplus}}
\sum_{x\ge0}\bigl(\hat{b}_t^{\,c_{\boxplus}}e^{-\beta w}\bigr)^x .
\end{align*}
The hypothesis $\beta w\ge c_{\boxplus}\log\hat{b}_t+1$ gives
$\hat{b}_t^{\,c_{\boxplus}}e^{-\beta w}\le e^{-1}$, so the geometric sum is at most $(1-e^{-1})^{-1}$,
and the left side is at most $2(1-e^{-1})^{-1}\hat{b}_t^{\,2c_{\boxplus}}\#\partial^+\Sigma$. Since
$2(1-e^{-1})^{-1}<4$, this is at most $C_{\mathrm{cell}}\,\#\partial^+\Sigma$, which is
\eqref{4.21:eq-cell-block-sums-1}.

(c) Every contour from $\square$ to $\Sigma$ starts at a point of $\square$, which lies in one of the
finitely many cells $\Delta$ meeting $\square$, and is then a contour from $\Delta$ to $\Sigma$; hence
$d^{\wedge}(\square,\Sigma)\ge\min_{\Delta\text{ meets }\square}d^{\wedge}(\Delta,\Sigma)$, and therefore
\begin{equation*}
e^{-\beta d^{\wedge}(\square,\Sigma)}\le\sum_{\Delta\text{ meets }\square}e^{-\beta d^{\wedge}(\Delta,\Sigma)} .
\end{equation*}
We sum over $\square$ and interchange the sums. By Lemma~\ref{4.21:lem-live-region-geometry} (see
the corresponding display), a cell meets at most $C_{\mathrm{blk}}$ blocks of the cover, so
\begin{equation*}
\sum_{\square}e^{-\beta d^{\wedge}(\square,\Sigma)}
\le C_{\mathrm{blk}}\sum_{\Delta}e^{-\beta d^{\wedge}(\Delta,\Sigma)}
\le C_{\mathrm{blk}}\bigl[\#\Sigma+C_{\mathrm{cell}}\,\#\partial^+\Sigma\bigr],
\end{equation*}
where the cells $\Delta\subset\Sigma$ contribute at most $1$ each, hence at most $\#\Sigma$, and the cells
$\Delta\not\subset\Sigma$ are summed by \eqref{4.21:eq-cell-block-sums-1}. Every cell of
$\partial^+\Sigma$ touches a cell of $\Sigma$, and a cell of $\Sigma$ touches at most $\hat{b}_t$ others,
as shown in (b),
so $\#\partial^+\Sigma\le\hat{b}_t\,\#\Sigma$. Inserting this gives \eqref{4.21:eq-cell-block-sums-2}.

(d) Assume \eqref{4.21:eq-cell-block-sums-a} and let $\beta\ge\delta_P/16$. Then
\begin{equation*}
\beta w\ \ge\ \tfrac1{16}\delta_P w\ \ge\ c_{\boxplus}\log\hat{b}_t+1 ,
\end{equation*}
which is the hypothesis of (b) and (c). Since $\hat{b}_t>1$, we have $c_{\boxplus}\log\hat{b}_t>0$, and
\eqref{4.21:eq-cell-block-sums-a} gives $\tfrac1{16}\delta_P w\ge1$.
\end{proof}

\begin{remark}[Where the relative length arises, and source-anchored index sets]
\label{4.21:rem-relative-length-origin}
Throughout this remark $k$ is the step of the $\mathbf{R}$-operation $\mathbf{R}_k$ of \cite{B-CMP122b},
and $Z$ is the large-field region denoted so on \cite[pp.~373--376]{B-CMP122b} and written $Z_k$ on
\cite[p.~390]{B-CMP122b} (the new large-field region being $Z_k\cup\bigcup_i Y_i$; see the sentence
preceding \cite[(1.102)]{B-CMP122b}). $\mathbf{D}_k$ is the class
of localisation domains of level $k$, and $M$-cubes are the cubes of $\pi_k$. For $Y\in\mathbf{D}_k$ the
relative linear size of \cite[(1.67), p.~376]{B-CMP122b} is
\begin{equation}\label{4.21:eq-relative-length-origin-1}
\begin{split}
d_{k,Z}(Y)=M^{-1}\bigl(&\text{the length of a shortest tree graph contained in } Y\\
&\text{and intersecting all $M$-cubes of the components of } Y\setminus Z\bigr).
\end{split}
\end{equation}
For any union $W$ of $M$-cubes we write $d_{k,W}(Y)$ for the right-hand side of
\eqref{4.21:eq-relative-length-origin-1} with $W$ in place of $Z$.
The following four points concern this quantity.
\begin{enumerate}
\item[(a)] The relative length enters at the step from \cite[(1.66)]{B-CMP122b} to
\cite[(1.67)]{B-CMP122b}. The reason is that the domain $Y_1$ of \cite[p.~373]{B-CMP122b} adds whole
components of $Z$ to the domain of a term, and these components bring no decay with them.
\item[(b)] The quantity \eqref{4.21:eq-relative-length-origin-1} is a tree length. A crossing of the tree
through a component of $Z$ is counted at its full length. Only the obligation to visit the $M$-cubes of
that component is lifted.
\item[(c)] At \cite[(1.92)--(1.97)]{B-CMP122b} the first-class components $X_j$ of the operation
$\mathbf{R}_k$ (in the classification of Definition~\ref{4.21:def-strips-geometry} and
\cite[(1.102)]{B-CMP122b}) are absorbed by the factors $\mathbf{T}'_k(X_j)1$. For the second-class
components $Y_i$
(Definition~\ref{4.21:def-strips-geometry}, \cite[(1.102)]{B-CMP122b}), the factor $\mathbf{T}'_k(Y_i)$
stands outside the bracket; see \cite[(1.72), (1.102)]{B-CMP122b}. Consequently
\cite[(1.97), (1.99)]{B-CMP122b} remain relative to $\bigcup_i Y_i$.
\item[(d)] In the setting of the correlation theorem, the reading convention of this chapter for the
enlarged domains of \cite{B-CMP122b} replaces the enlargement $Y_i^{\sim}$ of
\cite[p.~390]{B-CMP122b} by the strip $Y_i^{\sharp}$ of Definition~\ref{4.21:def-strips-geometry}.
This strip is $Y_i$ together with
$m_{\sharp}=\lfloor\frac12(\check{R}_k-1)\rfloor$ layers of $M$-cubes. Correspondingly, the relative size
$d_{k,\bigcup_i Y_i}$ is read as
\begin{equation}\label{4.21:eq-relative-length-origin-2}
d_{\mathrm{out}}:=d_{k,\bigcup_i Y_i^{\sharp}},
\end{equation}
and the parameter $\alpha$ of \cite[(1.99)]{B-CMP122b} is read as the radius $\alpha^{\sharp}$. We write
$C_{99}$ for the constant written $O(1)$ in \cite[(1.99)]{B-CMP122b}; it is kept in every bound. We
also put
\begin{equation}\label{4.21:eq-relative-length-origin-3}
c_1^{\flat}(g):=\max\bigl\{e^{-p_0(g)},\,(\alpha^{\sharp})^{1/3}\bigr\},
\end{equation}
where $p_0(\cdot)$ is the degree function. In this reading the terms $\mathbf{R}'^{(k)}(X)$ of the
exponentiation \cite[(1.98)]{B-CMP122b} satisfy the following bound on the product of the tube and the
set $\mathcal{W}_k(X)$, with the tube and $\mathcal{W}_k(X)$ as in the combination bound for the
exponentiated terms in its strip form:
\begin{equation}\label{4.21:eq-relative-length-origin-4}
\bigl|\mathbf{R}'^{(k)}(X;w)\bigr|\le C_{99}\,c_1\,
e^{-(1+\frac12\beta_s)\kappa\, d_{\mathrm{out}}(X)} .
\end{equation}
Here $\beta_s$ is the rate constant that takes, in this reading, the place of the constant $\beta$ of
\cite[(1.99)]{B-CMP122b}, and
$c_1$ is the factor of \cite[(1.99)]{B-CMP122b}, equal there to $e^{-p_0(g_k)}$ if the domain
$X'$ (as there) does not intersect $\bigcup_i Y_i$, and to $\alpha^{1/3}$ in the remaining cases; with
$\alpha$ read as $\alpha^{\sharp}$, this gives $c_1\le c_1^{\flat}(g_k)$.
This bound is the conclusion of the combination bound for the exponentiated terms in its strip form,
which holds under its input hypothesis on the exponentiation; it is used here under that hypothesis.
As in \cite[p.~390]{B-CMP122b}, the terms whose domains $X$ meet
$Z_k^{\sim}\cup\bigcup_i Y_i^{\sim}$ (under the strip convention,
$Z_k^{\sim}\cup\bigcup_i Y_i^{\sharp}$), the first group of \cite[p.~390]{B-CMP122b}, become the new
boundary terms $\mathbf{B}'^{(k)}(X)$. In every reading their size is relative, in the sense of
\eqref{4.21:eq-relative-length-origin-1}, to the collar of the step.
\end{enumerate}

The mechanism that produces the relative length is on \cite[pp.~373--374]{B-CMP122b}. The decay it
yields is used on \cite[p.~376]{B-CMP122b}. The mechanism rests on two features.
\begin{itemize}
\item The index sets $Y'_i$ ($i=2,3,4$) of the localisation expansions there are
\emph{source-anchored}: every component of $Y'_i$ contains a component of the region $\Lambda$ of
\cite[p.~373]{B-CMP122b}, inside which the changes of the configuration at this step are made.
\item The term's own domain is not exempt from the cut: the family $\sigma_0$ consists of all cubes
disjoint with $\Lambda$.
\end{itemize}
The argument of \cite[p.~376]{B-CMP122b} spends a second quantity, the third term of the exponential
\cite[(1.66)]{B-CMP122b}, on the sum over the localisation domains $X$. Wherever the term's own cubes
are cut, this third term is replaced by the vanishing of those terms of the localisation expansions
whose index sets have a component containing no component of $\Lambda$ \cite[p.~373]{B-CMP122b}. Inside
$\Lambda$, where the cubes are not cut, this replacement is not available, and the factor for the sum
over $X$ is supplied by the third-factor lemma, through the third part of the lemma on deep units.
\end{remark}

\begin{proof}
We take the points in order.

(a) By \cite[p.~373]{B-CMP122b}, $Y_1$ is the smallest domain in $\mathbf{D}_k$ that contains $X$ and
contains the components of $Z$ meeting $X$. The components of $Z$ added in this way carry no decay
factor. Hence the decay available after the resummations can be extracted only in a form that does not
require the tree to visit the cubes of these components. This is the form
\eqref{4.21:eq-relative-length-origin-1}, and it is extracted, as the factor
$\exp(-(1+3\beta)\kappa\, d_{k,Z}(Y_4))$ of \cite[p.~376]{B-CMP122b}, with $Y_4$ and $\beta$ as there,
from the exponential \cite[(1.66)]{B-CMP122b}.

(b) In \eqref{4.21:eq-relative-length-origin-1} the tree graph is required to be contained in $Y$. The
requirement that it intersect every $M$-cube is imposed only on the components of $Y\setminus Z$. A
segment of the tree lying in a component of $Z$ therefore contributes its full length to the minimum.
What is dropped is only the condition that the cubes of that component be met.

(c) In \cite[(1.92)--(1.97)]{B-CMP122b} the first-class components are absorbed by the factors
$\mathbf{T}'_k(X_j)1$. The operation $\mathbf{T}'_k(Y_i)$ of
a second-class component is applied outside the bracket; see \cite[(1.72), (1.102)]{B-CMP122b}. The
sets $Y_i$ are therefore not absorbed at this step. The bounds \cite[(1.97), (1.99)]{B-CMP122b} keep
their length relative to $\bigcup_i Y_i$.

(d) The displays \eqref{4.21:eq-relative-length-origin-2} and \eqref{4.21:eq-relative-length-origin-3}
are definitions. The bound \eqref{4.21:eq-relative-length-origin-4} is the conclusion of the combination
bound for the exponentiated terms in its strip form, written with
$d_{\mathrm{out}}$ in place of $d_{k,\bigcup_i Y_i}$ and with $C_{99}$ kept.

It remains to see how source-anchoring produces the relative length. Let $C$ be a component of $Z$ that
the tree crosses without visiting its cubes. Consider a component of $Y'_i$ that reaches the far side of
$C$ as seen from $\Lambda$.
\begin{itemize}
\item This component is connected and contains a component of $\Lambda$. Hence it contains at least as
many $M$-cubes as the distance, in units of $M$, from $\Lambda$ to $C$.
\item That distance is a fixed fraction of the size of $C$.
\item Each localisation operation yields, after the resummations with $Y'_i$ fixed, the factor
$\exp(-(\kappa_1-1)M^{-d}|Y'_i|)$ (see the derivation of \cite[(1.66)]{B-CMP122b}), and
$M^{-d}|Y'_i|$ counts the $M$-cubes of $Y'_i$.
\end{itemize}
Half of the exponent $(\kappa_1-1)M^{-d}|Y'_i|$ therefore pays for the crossing of $C$, and nothing
further is required at the face of $C$. This gives the relative length
\eqref{4.21:eq-relative-length-origin-1}.

It remains to justify the last paragraph of the remark. Where the term's own cubes are cut, a cube far
from $\Lambda$ is reached only by a long connected family of cut cubes, and the terms of the localisation
expansions whose index sets have a component containing no component of $\Lambda$ vanish; this is what
replaces the third term of \cite[(1.66)]{B-CMP122b}. The cubes inside $\Lambda$ are cut neither in the
expansions of \cite[pp.~373--374]{B-CMP122b} nor in the proof of the statement on the
wrapped-descendant class of stored terms: in the first operation of \cite[p.~373]{B-CMP122b} the family
$\sigma_0$ consists of the cubes disjoint with $\square_0$, or $X$, and with $\Lambda^{\sim}$ (with
$\square_0$ as there), and in the expansions indexed by $Y'_2,Y'_3,Y'_4$ it consists of the cubes
disjoint with $\Lambda$. Hence the replacement is not available inside $\Lambda$, and there the factor
for the sum over $X$ is supplied by the third-factor lemma, through the third part of the lemma on deep
units.
\end{proof}

\begin{remark}[Configurations the estimates must survive]\label{4.21:rem-test-configurations}
We use the pockets, strips and live pockets of Definition~\ref{4.21:def-strips-geometry} and
the geometry of a live pocket described in Lemma~\ref{4.21:lem-live-region-geometry}. The
integer $N_0$ and the sets $Z''_j$ are those fixed in \cite[(1.10), p.~179]{B-CMP122a}. The
\emph{window} is the family of one-layer strata $Z''_j\setminus Z''_{j-1}$ built there, and its
zones are the zones of the stage labels of the successive integrations that follow
\cite[(1.10)]{B-CMP122a}. The estimates of this chapter must take account of each of the
following ten configurations and failures.
\begin{enumerate}
\item[\textup{(1)}] Let $Y^{\sharp}$ be the strip of a live pocket, let $\Delta_1$ be a cube touching
$Y^{\sharp}$, and put $X:=Y^{\sharp}\cup\Delta_1$. Then, by the bound for a live strip with one
adjoined cube, $d_{\mathrm{out}}(X)=0$ and
\begin{equation*}
d_k(X)\ \ge\ (20\check{R}_k)^d/\mathsf{c}_d-1 ,
\end{equation*}
and the estimate that this bound delivers for $X$ is $O(1)(\alpha^{\sharp})^{1/3}$.
\item[\textup{(2)}] Take $j=k-1$ and a connection of shortest length from the strip
$Y^{\sharp}$ of a live pocket to $Z$.
\item[\textup{(3)}] The estimates may not use, for regions $X$ and $Y$, the inequality
\begin{equation*}
d_{\ell,W}(Y)\ \le\ d_{\ell,W}(X)+\sqrt{d}\,N_{\ell}(Y\setminus X),
\end{equation*}
which fails in general; here $d_{\ell,W}$ is the relative length of
Remark~\ref{4.21:rem-relative-length-origin} and $N_{\ell}(\cdot)$ is the number of level-$\ell$
cubes of a set.
\item[\textup{(4)}] If the strip is never decoupled, the families hung on one face of it sum to at least
\begin{equation*}
\exp\Bigl[\tfrac12\,e^{-(\kappa_1-1)}(20\check{R})^{d-1}\Bigr].
\end{equation*}
\item[\textup{(5)}] Consider a row of $n$ pockets merged into one island.
\item[\textup{(6)}] Consider one block $\square_0$ inside a strip, at the preparation site of the
$\mathbf{R}$-operation: all the terms around the pocket, of all levels $m,\dots,k$, read
$\square_0$.
\item[\textup{(7)}] Consider the pocket lying inside $\Lambda$ at the site \textup{(L)} of the
$\mathbf{R}$-step that treats its component (Definition~\ref{4.21:def-strips-geometry}): a count
of $L^{d(k-\ell)}$ per cell of level $\ell$ would leave a factor $L^{dN_k}$.
\item[\textup{(8)}] Consider the units lying wholly inside the core $Z_{h'}^{\sim}$. There can be as many of
them as the past volume of the core, a number bounded by no power of $\log g_k^{-2}$. The
region $X$ of configuration~(1) is itself one of them, at the step at which its pocket is
healed.
\item[\textup{(9)}] Consider a block $\square$ of zone $i$ whose own cell touches a window of
one-layer strata, in
which zone $i-1$ is one level-$i$ cube thick, zone $i-2$ is one level-$(i-1)$ cube thick, and
so on down to the last, thicker zone of the window. Then the enlargement $\tilde{\square}$
holds cells of all the levels of the window, $(2L^{N_0})^d$ of them. A fattened block is
therefore not a bounded union of cells.
\item[\textup{(10)}] Consider stage $j=k-1$ (or any later moment at which the window survives an
interruption), and a block $\square$ of zone $k-1$ whose own cell touches $Z''_{k-1}$. Fix a
term $c$, a label $A$ one of whose blocks is $\square$, and a probe $\mathsf{Q}$. Let $\sigma_0$
be the family of cubes lying off the union of the enlargements $\tilde{\square}^{4}$ over the
blocks of $A$, and let $n_{\mathrm{f}}$ be the number of cells of $\sigma_0$ that share a wall
with the far face of $\tilde{\square}^{4}$. Put $k_0:=k-N_0$. Then
$n_{\mathrm{f}}\ge(10\check{R}_{k_0+1}/L)^{d-1}$, and the majorant sum for $c$, $A$ and
$\mathsf{Q}$ is at least $(1+e^{-(\kappa_1-1)})^{n_{\mathrm{f}}}$. The majorant inequality of
the preparation site, read as an inequality between majorants, then holds only under the
condition
\begin{equation}\label{4.21:eq-test-configurations-1}
e^{-(\kappa_1-1)}\bigl(10\check{R}_{k_0+1}/L\bigr)^{d-1}\ \le\ C,
\end{equation}
with a constant $C$ not depending on $\check{R}$.
Condition \eqref{4.21:eq-test-configurations-1} is an upper bound on $\check{R}$, that is, a
lower bound on $g_k$: a cap, in the manner of Definition~\ref{2.3:def-cap}, but on the
coupling. This
condition is not imposed; it would bear on plain units (units carrying no pocket) as well. At
an old window the count
reads an old value of $\check{R}$ against a compensating
factor linear in $\check{R}_h$. Likewise, a cell under a window lies in $\tilde{\square}_0^{4}$
for blocks $\square_0$ of every zone of the window, so there is no bounded number of blocks
per cube.
\end{enumerate}
\end{remark}

\begin{proof}
\emph{Configuration} (9). By \cite[(1.10)]{B-CMP122a} the whole window fills two level-$k$
layers. The thin strata extend less than $5/4$ of a level-$(k-1)$ unit beyond the first layer.
The zones of the window are arranged as stated in configuration~(9): zone $i-1$ is one
level-$i$ cube thick, zone $i-2$ one level-$(i-1)$ cube thick, and so on down to the last,
thicker zone. The enlargement $\tilde{\square}$ reaches two level-$i$ units beyond the cell of
$\square$, so it holds cells of all the levels of the window. By the cell count for a fattened
block under a window, their number is $(2L^{N_0})^d$. Hence a fattened
block is not a bounded union of cells.

\emph{Configuration} (10). The enlargement $\tilde{\square}^{4}$ reaches four units beyond
$\square$. Starting from zone $k-1$, it crosses the zones $k-2,\dots,k_0+2$ and ends inside zone
$k_0+1$. The cells of zone $k_0+1$ have side $L^{-(N_0-2)}$ level-$(k-1)$ units; since the
cells of zone $j$ of the window have side $L^{-(k-1-j)}$ level-$(k-1)$ units (the scales of the
successive strata in \cite[p.~179]{B-CMP122a}), this is the zone $j=k-N_0+1=k_0+1$.

Let $\sigma_0$ and $n_{\mathrm{f}}$ be as in configuration~(10). The far face of
$\tilde{\square}^{4}$ has side $10$. Since the cells there have side $L^{-(N_0-2)}$, we have
\begin{equation}\label{4.21:eq-test-configurations-2}
n_{\mathrm{f}}\ \ge\ \bigl(10L^{N_0-2}\bigr)^{d-1}
\ =\ \bigl(10\check{R}_{k_0+1}/L\bigr)^{d-1}.
\end{equation}
The equality in \eqref{4.21:eq-test-configurations-2} holds because
$\check{R}_{k_0+1}=L^{N_0-1}$: this is the relation
$L^{-N_0+1}M\check{R}_{k-N_0+1}=M$ by which $N_0$ is chosen in \cite[p.~179]{B-CMP122a},
with $k-N_0+1=k_0+1$.

By the admissibility statement for families of cut cells, every subset $\mathsf{y}$ of these
$n_{\mathrm{f}}$ cells is an admissible family of cut cells. The activity $\nu$ of the element
of the expansion determined by $\mathsf{y}$
satisfies $\nu\propto\exp[-(\kappa_1-1)|\mathsf{y}|]$. For the fixed term $c$, label $A$ and
probe $\mathsf{Q}$,
the majorant sum therefore contains the sum of $e^{-(\kappa_1-1)|\mathsf{y}|}$ over all these
subsets. By the binomial theorem,
\begin{equation*}
\sum_{\mathsf{y}}e^{-(\kappa_1-1)|\mathsf{y}|}
\ =\ \sum_{r=0}^{n_{\mathrm{f}}}\binom{n_{\mathrm{f}}}{r}e^{-(\kappa_1-1)r}
\ =\ \bigl(1+e^{-(\kappa_1-1)}\bigr)^{n_{\mathrm{f}}} .
\end{equation*}
Hence the majorant sum is at least $(1+e^{-(\kappa_1-1)})^{n_{\mathrm{f}}}$.

This quantity stands against the per-block number in the constant $K_2$ of the
preparation-site bound; we write $n_{\mathrm{b}}$ for that per-block number. The majorant
inequality of the preparation site, read as an inequality between majorants, therefore
requires $(1+e^{-(\kappa_1-1)})^{n_{\mathrm{f}}}\le n_{\mathrm{b}}$, that is,
\begin{equation*}
n_{\mathrm{f}}\,\log\bigl(1+e^{-(\kappa_1-1)}\bigr)\ \le\ \log n_{\mathrm{b}} .
\end{equation*}
Since $\log(1+u)\ge u/(1+u)$ for every $u\ge 0$, this gives, with
\eqref{4.21:eq-test-configurations-2},
\begin{equation*}
\begin{split}
e^{-(\kappa_1-1)}\bigl(10\check{R}_{k_0+1}/L\bigr)^{d-1}
&\ \le\ e^{-(\kappa_1-1)}\,n_{\mathrm{f}}\\
&\ \le\ \bigl(1+e^{-(\kappa_1-1)}\bigr)\log n_{\mathrm{b}} ,
\end{split}
\end{equation*}
which is \eqref{4.21:eq-test-configurations-1} with
$C=(1+e^{-(\kappa_1-1)})\log n_{\mathrm{b}}$. Hence the majorant inequality of the preparation
site holds, as an inequality between majorants, only under
\eqref{4.21:eq-test-configurations-1}.
\end{proof}

\begin{definition}[Stages of the preparation and excluded components]
\label{4.21:def-stages-components}
Fix the step $k$ and consider the preparatory integrations of the operation $\mathbf{R}_k$,
whose stage of level $j$ integrates the variables localized in $Z_{j+1}\setminus Z''_{j+1}$, the
sets $Z''_i$ being those of \cite[(1.10)--(1.11)]{B-CMP122a}. Put
\begin{equation}\label{4.21:eq-stages-components-1}
h:=k-N,\qquad k_0:=k-N_0,
\end{equation}
where $N_0$ is the integer fixed in Definition~\ref{4.21:def-strips-geometry} by the relation
$L^{N_0-1}=\check{R}_{k_0+1}\ge L$ (see \eqref{4.21:eq-strips-geometry-a}), and where $N>N_0\ge 2$.
In particular $N\ge 3$. The number $N$ of stages is the value $N=\check{R}_k$ declared with the
admissible geometry of Definition~\ref{4.21:def-strips-geometry}; of this declaration only the
relation between $N$ and the radii $\check{R}$ in \eqref{4.21:eq-strips-geometry-a} is used in
the proofs of this section. That the declared value also satisfies the lower bound on $N$ imposed
on the number of stages elsewhere in the construction is a separate statement and is not used
here.

\emph{Stages.} For $0\le n\le N-1$, stage $n+1$ of the preparation integrates the level
$j:=h+n$; thus $h\le j\le k-1$. On the running components (defined below) the current label
is $\mathbb{B}^{(n)}_k$. We write $Z$ for the large-field region on whose components
$\mathbf{R}_k$ acts, taken as it stands
at the beginning of this operation $\mathbf{R}_k$; its connected components are the
\emph{class components}. In particular a component excluded (a sibling, in the sense defined
below) at an earlier stage of this
$\mathbf{R}_k$ still lies in $Z$; by the convention on the grains of excluded components, its
grain partition $\pi^{\flat}$ is read at the running level $j$.

\emph{Running components and siblings.} A class component $C'$ is either \emph{running} at
stage $n+1$, which means that its preparation has performed the stages $1,\dots,n$ and is at
level $j$, or a \emph{sibling}, which means that $C'$ has been excluded from the running
large-field region (the union of the class components not yet excluded, which
\cite{B-CMP122a} continues to denote by $Z$; we also call it the running region) in the course
of this $\mathbf{R}_k$ (it remains a class component of $Z$ as
fixed above), at a stage $\le n+1$, after $n(C')\le n$ performed stages, by one of the
exclusions of \cite[pp.~181--183]{B-CMP122a} or by the exclusion at
\cite[(1.55)]{B-CMP122a}. A sibling is an arrested attempt $\mathfrak{a}$ with $k_a=k$, in the
sense of the definition of arrested attempts, namely the tuple
\begin{equation*}
\mathfrak{a}=\bigl(C',\,k,\,n(C'),\,n^+(C'),\,\text{branch}\bigr),
\end{equation*}
whose last entry specifies the exclusion at which the exit takes place. On every branch except the
exclusion at \cite[(1.55)]{B-CMP122a} we have $n^+(C')=n(C')$; on that branch
$n^+(C')=n(C')+1$, since there the change of variables of the next stage is made (its band is
at level $h+n(C')+1$ and its datum is $V_{h+n(C')+1}$) while its fluctuation integral is not.
The preparation of a sibling is frozen at its exit, with top raised level
\begin{equation}\label{4.21:eq-stages-components-2}
j(C'):=h+n^+(C')\le j+1 ,
\end{equation}
and its raised cells, its bands and its chain terms $\mathbb{C}^{(i)}_k[C']$,
$1\le i\le n(C')$, are
kept. The equality $j(C')=j+1$ occurs only for an exit at \cite[(1.55)]{B-CMP122a} during the
current stage, that is with $n(C')=n$; such a sibling enters stage $n+1$ through its geometry
alone. Every sibling $\mathfrak{a}$ whose terms $\mathfrak{Q}(\mathfrak{a})$ are read at stage
$n+1$ satisfies $j(C')\le j$. For a running component $C'$ we put $j(C'):=j$ and $n(C'):=n$.

\emph{Chains.} The chains of this $\mathbf{R}_k$ are the terms
\begin{equation*}
\mathbb{C}^{(i)}_k[C'],\qquad C'\ \text{a class component},\quad 1\le i\le n(C');
\end{equation*}
the \emph{local chain} at a point of $C'$ is the family of the chains of $C'$. The collar $R_n$
of stage $n+1$ lies in the running components only. In what follows, the fact that a component
is running enters only through the decay of the block factor $\vartheta(\square)$, and that
factor is only ever bounded above by $1$.

\emph{Zones.} The sets $Z''_i$, $h+1\le i\le k$, are those of \cite[(1.10)--(1.11)]{B-CMP122a};
they depend on the geometry of $\mathbb{B}_k$ only, and they are defined on every class
component, running or sibling. In a class component $C'$ the current zones are:
\begin{itemize}
\item for $h+1\le i<j(C')$: the stratum $\Gamma''_i=Z''_{i+1}\setminus Z''_i$;
\item for $i=j(C')$: the set
$\bigl(\Omega^c_{j(C')+1}\setminus Z''_{j(C')}\bigr)\cap C'$;
\item for $i>j(C')$: the zones made by the $\mathbf{T}$-steps.
\end{itemize}

\emph{Domains.} All domains $\Omega$, $Z$, $\Lambda$ are those of the admissible geometry of
Definition~\ref{4.21:def-strips-geometry}: $\Omega^{(t)}_m$ denotes the domains after the
complete step $t$; $\Lambda_t$ is the small-field region and $Z_t:=\Lambda_t^c$;
$\Omega^{\mathbf{T}}_t\supset\Lambda^{\mathbf{T}}_t$ are the domain and the small-field region
made by $\mathbf{T}_t$; at $\mathbf{R}_k$, the sets $\Lambda_k$, $Z_k$ are the ones made by
$\mathbf{T}_k$. The conditions \eqref{4.21:eq-strips-geometry-a} of
Definition~\ref{4.21:def-strips-geometry} hold, Lemma~\ref{4.21:lem-live-region-geometry}
applies to every second-class region (its conclusions are displayed in
the corresponding display),
and in addition
\begin{equation}\label{4.21:eq-stages-components-a}
Z_{t-1}\subset Z'^{\sim 8}_{t-1}\subset\bigl(\Omega^{\mathbf{T}}_t\bigr)^c ,
\end{equation}
which follows from \cite[(3.5)]{B-CMP119} and is established in the proof of part (d) of
Lemma~\ref{4.21:lem-live-region-geometry}.

\emph{Completed class decisions.} A component carried through the completed class decision of
$\mathbf{R}_t$, that is, through the division of the components into two classes of
\cite[p.~378]{B-CMP122b} (the step by which, in Definition~\ref{4.21:def-strips-geometry}, a
second-class region is created and by which a live one dies), takes its class as a whole. If it
is of the first class it is
\emph{healed}: all of it joins $\Lambda_t$, at level $t$. If it is of the second class it is a
second-class region $(Y,t)$ in the sense of Definition~\ref{4.21:def-strips-geometry}: all of it
carries, in
$\mathbb{B}_{t'}$ for $t\le t'\le k$ and for as long as this second-class region is live, the one
domain
$Y^{(t)}$ of level $t$. This is \cite[(1.33)]{B-CMP122b}: in the notation of that paper, where
$k$ denotes the step and $Z$ its large-field region, the label $\mathbb{B}_k$ is taken equal to
$\mathbb{B}''_k$ on $Z^c$ and to $\{Z^{(k)}\}$ on $Z$; read at step $t$, this gives the flat
label $\{Y^{(t)}\}$ on $Y$. It is the sequence of domains that the
operation assigns to the class before its division into the two classes. The functions
$\chi^c_t(Y^{\sim-6})$ and $\chi_{t,\Lambda\cap Y}$ of \cite[p.~378]{B-CMP122b} persist as
factors of the density on $Y$, not as levels.

\emph{Levels and base.} For a point $x$ in the interior of a current cell:
\begin{itemize}
\item $\operatorname{lev}(x)$ is the current level of its cell;
\item $m_0(x)$, the \emph{original level}, is the level of its cell in $\mathbb{B}_k$, that is,
after $\mathbf{T}_k$ and before the preparations of this $\mathbf{R}_k$;
\item $b(x)$, the \emph{base}, is the largest $t\le k$ such that either
$x\in\Omega^{\mathbf{T}}_t$, or $x$ lies in a component carried through the completed class
decision of $\mathbf{R}_t$, of either class (healed at level $t$, or thrown back as a
second-class region
$(Y,t)$ with the label $Y^{(t)}$); if there is no such $t$, then $b(x):=0$, the level of the
starting label.
\end{itemize}
\end{definition}

\begin{proof}[Proof of the assertions in Definition~\ref{4.21:def-stages-components}]
We prove, in order, the inequality $N\ge 3$, the range of $j$, the bound
\eqref{4.21:eq-stages-components-2}, the description of its equality case, and the bound
$j(C')\le j$ for the siblings whose terms are read at stage $n+1$.

First, the blocking factor satisfies $L>1$, so the inequality $L^{N_0-1}\ge L$ gives
$N_0-1\ge 1$, that is $N_0\ge 2$. Since $N$ and $N_0$ are integers with $N>N_0$, we get
$N\ge N_0+1\ge 3$.

Second, for $0\le n\le N-1$ we have $j=h+n\ge h$ and, by \eqref{4.21:eq-stages-components-1},
\begin{equation*}
j=h+n\le h+N-1=k-1 .
\end{equation*}

Third, let $C'$ be a sibling at stage $n+1$. It was excluded at a stage $\le n+1$, after
$n(C')\le n$ performed stages. The exit happens either at one of the exclusions of
\cite[pp.~181--183]{B-CMP122a}, where $n^+(C')=n(C')$, or at the exclusion of
\cite[(1.55)]{B-CMP122a}, where $n^+(C')=n(C')+1$. The second value arises because on that
branch the change of variables of the stage following the $n(C')$ performed stages has been made
before the exit, while the fluctuation integral of that stage has not. Moreover, by the order of
the operations on \cite[p.~188]{B-CMP122a}: in the integral obtained after the change of
variables of a stage, the decomposition of unity $1=\chi'^{(j)}+(1-\chi'^{(j)})$ is introduced in
each component of the running region at \cite[(1.55)]{B-CMP122a}, the components carrying the
large-field function $1-\chi'^{(j)}$ are excluded from the running region, and only afterwards,
for the components carrying
the small-field function $\chi'^{(j)}$, is the fluctuation integral of the stage written, at
\cite[(1.60)]{B-CMP122a}. The exit on this branch therefore precedes the fluctuation integral of
the stage in which it occurs; a component excluded in this way during the current stage is seen
at stage $n+1$, so the case $j(C')=j+1$ treated below can occur. In every case
$n^+(C')\le n(C')+1$, hence
\begin{equation*}
j(C')=h+n^+(C')\le h+n(C')+1\le h+n+1=j+1 ,
\end{equation*}
which is \eqref{4.21:eq-stages-components-2}. For a running component, $j(C')=j\le j+1$ by
definition.

Fourth, suppose $j(C')=j+1$ for a sibling $C'$. Then $n^+(C')=n+1$. Since
$n^+(C')\le n(C')+1\le n+1$, both inequalities are equalities: $n(C')=n$ and
$n^+(C')=n(C')+1$. The second equality holds only on the branch of the exclusion at
\cite[(1.55)]{B-CMP122a}, and the first says that the exit took place during the current stage
$n+1$. This proves that $j(C')=j+1$ occurs only for an exit at \cite[(1.55)]{B-CMP122a} during
the current stage.

Fifth, let $\mathfrak{a}=(C',k,n(C'),n^+(C'),\text{branch})$ be a sibling. By the definition of
the frozen terms of an arrested attempt, in its clause on the stages at which the terms of a
sibling are read, the terms $\mathfrak{Q}(\mathfrak{a})$ are read only at the stages
$l'>n^+(C')$ of this $\mathbf{R}_k$; by the statement on the reading of these terms after the
exit, they are read only at the stages $l'>n(C')+1$, and at the stage $n(C')+1$ itself, which on
the branch of \cite[(1.55)]{B-CMP122a} is the stage $n^+(C')$, none of them enters. If these
terms are read at stage $n+1$, then $n+1>n^+(C')$,
that is $n^+(C')\le n$, and therefore
\begin{equation*}
j(C')=h+n^+(C')\le h+n=j .
\end{equation*}
In particular a sibling with $j(C')=j+1$ has $n^+(C')=n+1$, so none of its terms
$\mathfrak{Q}(\mathfrak{a})$ is read at stage $n+1$: it enters that stage through its geometry
alone. The sharper bound $j(C')\le j$ is thus available for the siblings whose terms are read at
stage $n+1$; all statements of this section are proved under the bound
\eqref{4.21:eq-stages-components-2}, which holds for every sibling.
\end{proof}

\begin{lemma}[Base levels never exceed current levels]\label{4.21:lem-base-below-level}
Fix the level $k$ and the large-field operation $\mathbf{R}_k$, with the notation of
Definition~\ref{4.21:def-stages-components}; second-class regions $(Y,m)$, their strips
$Y^{\sharp}$ and the sense in which such a region is live or dies are those of
Definition~\ref{4.21:def-strips-geometry}.
For $t\le k$ let $\Omega^{\mathbf{T}}_t\supset\Lambda^{\mathbf{T}}_t$ be the domain and the
small-field region made by $\mathbf{T}_t$, and $Z^{\mathbf{T}}_t$ the complement of
$\Lambda^{\mathbf{T}}_t$; let $\Lambda_t$ be the small-field region after the complete step $t$
and $Z_t:=\Lambda_t^c$, where at $\mathbf{R}_k$ the sets $\Lambda_k$, $Z_k$ are those made by
$\mathbf{T}_k$. For $t\le k$ write $h(t)$ for the level $h$ of
Definition~\ref{4.21:def-stages-components} at step $t$, so that $h(k)=h$.

At step $t$ the operation $\mathbf{R}_t$ acts on a class component $X$ first by its
preparation, which raises cells of $X$, and then by its class decision on $X$. If the class
decision completes, $X$ takes its class whole: if $X$ is of the first class it is
\emph{healed}, and all of $X$ joins $\Lambda_t$ at level $t$; if $X$ is of the second class it
is \emph{thrown back} as a second-class region $(X,t)$, and all of $X$ carries in
$\mathbb{B}_{t'}$, for $t\le t'\le k$ while $(X,t)$ lives, the one domain $X^{(t)}$ of level $t$
\cite[(1.33), p.~365]{B-CMP122b}, \cite[p.~378]{B-CMP122b}. Arrested attempts and the time
during which they are live are those of the definition of arrested attempts: an attempt
$\mathfrak{a}$ of $\mathbf{R}_t$ on $X$ is arrested when it is interrupted or excluded before its
class decision on $X$ completes, and we then write $X_{\mathfrak{a}}:=X$ and $k_a:=t$; in
particular $\mathfrak{a}$ is live until the first later step whose $\mathbf{R}$-operation
completes, of either class, on the component holding $X_{\mathfrak{a}}$.

For a point $x$ in the interior of a current cell, $\operatorname{lev}(x)$ and $m_0(x)$ are as
in Definition~\ref{4.21:def-stages-components}: the current level of its cell, and the level of
its cell in $\mathbb{B}_k$, that is, after $\mathbf{T}_k$ and before the preparations of this
$\mathbf{R}_k$. The \emph{base}
$b(x)$ is the largest $t\le k$ such that $x\in\Omega^{\mathbf{T}}_t$, or $x$ lies in a
component on which the class decision of $\mathbf{R}_t$ completed, of either class (healed at
level $t$, or thrown back as a second-class region $(Y,t)$ with the domain $Y^{(t)}$); if there
is no such
$t$, then $b(x):=0$, the level of the starting label. We write $\mathcal{W}$ for the set of
cells raised by the preparation of a live arrested attempt $\mathfrak{a}$ of a step $k_a<k$, and
$t(x):=k_a$ for $x$ in a cell of $\mathcal{W}$. Finally, as in
Definition~\ref{4.21:def-stages-components},
$\mathsf{r}(x)=\operatorname{lev}(x)-b(x)$ and $\mathsf{r}_0(x)=\operatorname{lev}(x)-m_0(x)$.

Then the following hold. For every arrested attempt $\mathfrak{a}$ live at $\mathbf{R}_k$,
every second-class region $(Y,m)$ live at $\mathbf{R}_k$, and every point $x$ in the interior of
a current cell,
\begin{equation}\label{4.21:eq-base-below-level-a}
\begin{gathered}
X_{\mathfrak{a}}\subset Z_{t'}\quad(k_a\le t'\le k-1),\qquad
Y\subset Z_{t'}\quad(m\le t'\le k-1),\\
\text{(a)}\quad b(x)\le m_0(x)\le \operatorname{lev}(x);
\end{gathered}
\end{equation}
in particular no point of $X_{\mathfrak{a}}$ lies in $\Omega^{\mathbf{T}}_{t'}$ for
$k_a<t'\le k$, and no point of $Y$ lies in $\Omega^{\mathbf{T}}_{t'}$ for $m<t'\le k$.
Moreover:
\begin{enumerate}
\item[(b)] $m_0(x)>b(x)$ if and only if the cell of $x$ was raised by the preparation of an
attempt $\mathfrak{a}'$ of a step $k_{a'}<k$ that was arrested and is still live;
\item[(c)] on the cells of a live second-class region $(Y,m)$ one has
$\mathsf{r}=\mathsf{r}_0=0$, $Y\cap\mathcal{W}=\emptyset$, and
\begin{equation*}
\operatorname{lev}=m_0=b=m;
\end{equation*}
\item[(d)] $b$ is non-decreasing in time, and on a cell of $\mathcal{W}$ it does not change
while the attempt that raised it is live;
\item[(e)] at every point at most one of $\mathsf{r}_0$ and $m_0-b$ is non-zero; $t(x)$ is
well defined on $\mathcal{W}$, and $\operatorname{lev}=m_0\le t(x)$ there.
\end{enumerate}
\end{lemma}

\begin{proof}
We first record three facts on the way levels are assigned, then prove the inclusions
of \eqref{4.21:eq-base-below-level-a}, and then the remaining assertions in order.

\emph{Step 1: the level-assigning operations.} The cell of $x$ receives a level only by one of
the following four operations.
\begin{enumerate}
\item[(i)] The step $\mathbf{T}_t$, when $x\in\Omega^{\mathbf{T}}_t$: the cell receives level
$t$.
\item[(ii)] A completed first-class decision of $\mathbf{R}_t$ on the class component containing
$x$: level $t$ on all of that component.
\item[(iii)] A completed second-class decision of $\mathbf{R}_t$ on the class component $Y$
containing $x$: level $t$ on all of $Y$, carried while $(Y,t)$ lives
\cite[(1.33), p.~365]{B-CMP122b}.
\item[(iv)] A preparation of $\mathbf{R}_t$ on the class component containing $x$, raising cells
of $Z\setminus Z''_i$ (the sets of step $t$) to levels $i$ with $h(t)<i\le t$, whatever the
later fate of the attempt.
\end{enumerate}
Nothing else assigns a level: excluded components are left unchanged \cite{B-CMP122a}, and the
raises made by the preparation of an arrested attempt persist, not changed by the subsequent
operations, by the statement on the raises of arrested attempts that accompanies the definition
of arrested attempts. That none of the
operations (i)--(iv) re-levels a live second-class region is not part of this list; it is proved
below, in part~(c), from the inclusions of \eqref{4.21:eq-base-below-level-a}. Within one step
$t$ the order is: first $\mathbf{T}_t$, then the preparations of $\mathbf{R}_t$, then the class
decisions of $\mathbf{R}_t$.

\emph{Step 2: monotonicity.} Each of the operations (i)--(iv) assigns a level at least equal to
the old one. For (i), $\Omega^{\mathbf{T}}_t\subset\Lambda_{t-1}$, where levels were at most
$t-1$, by \eqref{4.21:eq-strips-geometry-a}. For (ii) and
(iii), the levels of the component were at most $t$ before the decision. An operation of kind
(iv) raises. Hence $\operatorname{lev}$ is non-decreasing in time: nothing lowers a level.

\emph{Step 3: base events.} We call \emph{base events} of $x$ the operations of kinds (i),
(ii), (iii) acting on the cell of $x$; these are exactly the operations that enter the
definition of $b(x)$, and we count the starting label as a base event at step $0$. By Step~1,
each base event at step $t$ leaves $x$ at level exactly $t$.

\emph{Step 4: the inclusions of \eqref{4.21:eq-base-below-level-a}.} Let $S$ be either
$X_{\mathfrak{a}}$, with $t_0:=k_a$, for an arrested attempt $\mathfrak{a}$ live at
$\mathbf{R}_k$, or $Y$, with $t_0:=m$, for a second-class region $(Y,m)$ live at
$\mathbf{R}_k$. First, $S\subset Z_{t_0}$: the component $X_{\mathfrak{a}}$ of an arrested
attempt, interrupted or excluded, is carried through no completed class decision at step $k_a$;
it is therefore not healed there and stays large-field, so $X_{\mathfrak{a}}\subset Z_{k_a}$;
likewise $Y\subset Z_m$, $Y$ being thrown back at step $m$. Now let $t_0<t'\le k-1$ and
suppose $S\subset Z_{t'-1}$. By the inclusion \eqref{4.21:eq-stages-components-a},
$Z_{t'-1}\subset(\Omega^{\mathbf{T}}_{t'})^c$, and since
$\Lambda^{\mathbf{T}}_{t'}\subset\Omega^{\mathbf{T}}_{t'}$ we have
$(\Omega^{\mathbf{T}}_{t'})^c\subset Z^{\mathbf{T}}_{t'}$; so $S\subset Z^{\mathbf{T}}_{t'}$.
Being connected, $S$ lies in one component $K$ of $Z^{\mathbf{T}}_{t'}$. The component $K$ is
carried through no completed class decision of $\mathbf{R}_{t'}$: such a decision would end the
life of $\mathfrak{a}$, respectively $(Y,m)$ would die there, whereas both are live at
$\mathbf{R}_k$ and $t'<k$. Hence $K\subset Z_{t'}$, and so $S\subset Z_{t'}$. By induction on
$t'$ this proves $S\subset Z_{t'}$ for $t_0\le t'\le k-1$, which is the first line of
\eqref{4.21:eq-base-below-level-a}. Applying \eqref{4.21:eq-stages-components-a} once more,
$S\subset Z_{t'-1}\subset(\Omega^{\mathbf{T}}_{t'})^c$ for $t_0<t'\le k$: no point of
$X_{\mathfrak{a}}$ lies in $\Omega^{\mathbf{T}}_{t'}$ for $k_a<t'\le k$, and no point of $Y$
lies in $\Omega^{\mathbf{T}}_{t'}$ for $m<t'\le k$.

\emph{Step 5: the inequalities $b\le m_0\le\operatorname{lev}$.} Let $E$ be the last base event
of $x$ (the starting label if there is no other); it takes place at step $b=b(x)$. Right after
$E$ the level of the cell of $x$ is $b$, by Step~3. Every later operation of Step~1 on this
cell, up to and including $\mathbf{T}_k$, is not a base event, by the choice of $E$ as the last
one and the order of the operations within a step; hence it is of kind (iv), which raises, by
Step~2. Therefore $m_0(x)\ge b(x)$. During $\mathbf{R}_k$, before its class decisions, only
preparations act, and these raise; hence $\operatorname{lev}(x)\ge m_0(x)$. This proves the
second line of \eqref{4.21:eq-base-below-level-a}.

\emph{Step 6: part (b), the direct implication.} Let $m_0(x)>b(x)$, and let $O$ be the last
operation of Step~1 on the cell of $x$ up to $\mathbf{T}_k$; it exists, since the starting
label is such an operation. After $O$ the level is $m_0(x)$. If $O$ were the starting label,
then $m_0(x)=0\le b(x)$, a contradiction. If $O$ were of kind (i), (ii) or (iii) at a step $t$,
then $m_0(x)=t$ by Step~3, and $O$ would be a base event, so that $b(x)\ge t=m_0(x)$, again a
contradiction. Hence $O$ is of kind (iv): the preparation of $\mathbf{R}_t$ on the class
component $X$ containing $x$, raising the cell of $x$ to the level $m_0(x)\le t$. Here
$t\le k-1$, because the preparations of $\mathbf{R}_k$ act after $\mathbf{T}_k$. The class
decision of the same $\mathbf{R}_t$ on $X$ did not complete: a completed decision, of either
class, would be an operation of kind (ii) or (iii) on the cell of $x$ later than $O$. Hence the
attempt $\mathfrak{a}$ of $\mathbf{R}_t$ on $X$ was arrested, with $k_a=t<k$, and the cell of $x$
is one of its raised cells. No $\mathbf{R}_{t'}$ with $t<t'<k$ completed on the component
holding $X_{\mathfrak{a}}$, since it would be an operation of kind (ii) or (iii) on the cell of
$x$ later than $O$; and the class decisions of $\mathbf{R}_k$ have not yet taken place. Hence
$\mathfrak{a}$ is live.

\emph{Step 7: part (b), the converse implication.} Let the cell of $x$ be raised, to the level
$m'$, by the preparation of a live arrested attempt $\mathfrak{a}$ with $k_a<k$. Since the
preparation acts on the class component $X_{\mathfrak{a}}$, the point $x$ lies in
$X_{\mathfrak{a}}$. Just before the raise the level was at most $m'-1$, by Step~2 and because
the raise is strict. Every base event before the raise, at a step $t$, left the level $t$
(Step~3), and the level did not decrease afterwards (Step~2); hence $t\le m'-1$ for all these
events. After the raise there is no base event:
\begin{itemize}
\item an operation of kind (i) at a step $t'$ with $k_a<t'\le k$ is excluded by Step~4, since
no point of $X_{\mathfrak{a}}$ lies in $\Omega^{\mathbf{T}}_{t'}$, and $\mathbf{T}_{k_a}$
precedes $\mathbf{R}_{k_a}$;
\item an operation of kind (ii) or (iii) at step $k_a$ is excluded because $\mathfrak{a}$ was
arrested; at a step $t'$ with $k_a<t'<k$ it is excluded because $\mathfrak{a}$ is live; at step
$k$ it has not yet taken place.
\end{itemize}
Hence $b(x)\le m'-1$. Moreover $m_0(x)=m'$, because after the raise no operation of Step~1
touched the cell of $x$ up to $\mathbf{T}_k$: none of kinds (i), (ii), (iii), as just shown,
and none of kind (iv), because no cell is raised twice. Indeed, by
Definition~\ref{4.21:def-stages-components} a preparation of step $t$
raises only cells of $\Lambda_{h(t)}$, at least three layers away from every older
$X_{\mathfrak{a}'}$, by the statement on the raises of arrested attempts used in Step~1, and
from every older second-class region, by part~(h) of
Lemma~\ref{4.21:lem-live-region-geometry}; since the cell of $x$ lies in
$X_{\mathfrak{a}}$, no later preparation raises it. Therefore $m_0(x)=m'>b(x)$.

\emph{Step 8: part (c).} Let $(Y,m)$ be a live second-class region, so that $m\le k-1$, and let
$x\in Y$.
The completed second-class decision of $\mathbf{R}_m$ on $Y$ is a base event of $x$, so
$b(x)\ge m$. No base event occurs after it: an operation of kind (i) at a step $t'$ with
$m<t'\le k$ is excluded by Step~4; one of kind (ii) or (iii) at a step $t'$ with $m<t'<k$ would
end the life of $(Y,m)$; at step $k$ none has yet taken place. Hence $b(x)=m$, and right after
the
decision of $\mathbf{R}_m$ the level is $m$, by Step~3. Among the later operations of Step~1
up to $\mathbf{T}_k$ there is none of kinds (i), (ii), (iii), as just shown. Consider one of
kind (iv), the preparation of a step $t'$ with $m<t'<k$. If it belongs to an attempt
$\mathfrak{b}$ whose component holds $Y$, then $m\le h_b$, by part~(g) of
Lemma~\ref{4.21:lem-live-region-geometry}, where $h_b$ is the level $h(t')$ of $\mathfrak{b}$;
and $\mathfrak{b}$ raises only outside the set $Z''^{(\mathfrak{b})}_{h_b+1}$ of its own step,
for which
\begin{equation*}
Z''^{(\mathfrak{b})}_{h_b+1}\supset Z_{h_b}\supset Y,
\end{equation*}
the second inclusion by \eqref{4.21:eq-base-below-level-a}, as $m\le h_b\le k-1$. An attempt on
another component does not touch $Y$. The preparation of step $m$ acted before the decision of
$\mathbf{R}_m$. Hence $m_0(x)=m$. During this $\mathbf{R}_k$: if $Y$ lies in a class component
$C'$, then $m\le h$ by part~(g) of Lemma~\ref{4.21:lem-live-region-geometry}, so
$Y\subset Z_h$ by \eqref{4.21:eq-base-below-level-a}, and $Y\subset Z_h\cap C'\subset Z''_{h+1}$;
as the preparation of this $\mathbf{R}_k$ raises inside $Z\setminus Z''_{h+1}$ only
(Definition~\ref{4.21:def-stages-components}), $Y$ is
never raised. If $Y$ lies in a component not treated by this $\mathbf{R}_k$, nothing is raised
there. Therefore $\operatorname{lev}=m_0=b=m$ on $Y$, so $\mathsf{r}=\mathsf{r}_0=0$; and since
$m_0=b$ on $Y$, part~(b) gives $Y\cap\mathcal{W}=\emptyset$. Distinct live second-class regions
are disjoint,
by part~(b) of Lemma~\ref{4.21:lem-live-region-geometry}, so the statement $b=m$ on $Y$ is
unambiguous.

\emph{Step 9: part (d).} The base $b(x)$ is the maximum of the steps of a set of events that
only grows in time, so it is non-decreasing. On a cell of $\mathcal{W}$, Step~7 showed that no
base event occurs after the raise while the raising attempt $\mathfrak{a}$ is live; hence $b$
does not change there.

\emph{Step 10: part (e).} If $m_0(x)>b(x)$, then the cell of $x$ belongs to $\mathcal{W}$, by
part~(b); it is therefore not raised during this $\mathbf{R}_k$, because no cell is raised twice
(Step~7); hence $\mathsf{r}_0(x)=0$. If $\mathsf{r}_0(x)>0$, then the cell was raised during this
$\mathbf{R}_k$, since during $\mathbf{R}_k$ only preparations act (Step~5); hence it does not
belong to $\mathcal{W}$, again because no cell is raised twice, and then $m_0(x)=b(x)$ by
part~(b) and \eqref{4.21:eq-base-below-level-a}. Thus at most one of $\mathsf{r}_0$ and $m_0-b$
is non-zero. A cell of $\mathcal{W}$ is raised by exactly one preparation, since no cell is
raised twice; so the raising attempt, and with it $t(x)=k_a$, is well defined. On $\mathcal{W}$
we have $\mathsf{r}_0=0$, that is, $\operatorname{lev}=m_0$, and by Step~6 the raise of
$\mathfrak{a}$ gave the level $m_0(x)\le k_a$. This completes the proof.
\end{proof}

\begin{definition}[Raised cells, arrested attempts and their frozen data]
\label{4.21:def-arrested-attempts}
Fix the step $k$ and the operation $\mathbf{R}_k$. The levels $h$, $k_0$, the numbers $N>N_0$,
the class $Z$ at the start of $\mathbf{R}_k$, its class components $C'$ (running or sibling) and the
indices $j(C')$, $n(C')$ are as in Definition~\ref{4.21:def-stages-components} and
Lemma~\ref{4.21:lem-base-below-level}. For a point $x$ in the interior of a current cell,
$\operatorname{lev}(x)$ is the current level of its cell, $m_0(x)$ its level in $\mathbb{B}_k$
(after $\mathbf{T}_k$, before the preparations of $\mathbf{R}_k$) and $b(x)$ its base, as there.
\begin{enumerate}
\item A cell is \emph{raised during $\mathbf{R}_k$} if the preparation of a running component or
of a sibling raises it in the course of $\mathbf{R}_k$.
\item An \emph{arrested attempt} is written
$\mathfrak{a}=(X_{\mathfrak{a}},k_a,n,n^+,\text{branch})$. It is an attempt of step $k_a$ that was
interrupted after $n<N(g_{k_a})$ stages \cite[p.~191]{B-CMP122a}, or excluded before any stage or
after all of them \cite[p.~193]{B-CMP122a}; the branch is its case among (B0)--(B7) in the
classification of interrupted and excluded attempts. Its data are the following.
\begin{itemize}
\item The level $h_a$, which is the level $h$ of Definition~\ref{4.21:def-stages-components}
taken at the step $k_a$, and the level $k_0(k_a):=k_a-N_0(k_a)$.
\item The performed levels $h_a,\dots,h_a+n-1$.
\item The raised levels $i$ with $h_a<i\le j_a$, where
\begin{equation}\label{4.21:eq-arrested-attempts-1}
j_a:=h_a+n^+\le k_a .
\end{equation}
The equality $j_a=k_a$ occurs: on the branches (B5)--(B7), and on the branch (B4) with
$n=N(g_{k_a})-1$.
\item The set $X_{\mathfrak{a}}\subset Z_{k_a}$, which is a component of the class at
$\mathbf{R}_{k_a}$.
\item Its \emph{collar} $\mathfrak{c}_{\mathfrak{a}}:=X_{\mathfrak{a}}\cap\Omega_{k_a}$. This is two
layers of $M\check{R}_{k_a}$-cubes, of level $k_a$, and unraised.
\item Its \emph{piece} $W_{\mathfrak{a}}:=X_{\mathfrak{a}}\cap\Omega^c_{k_a}$.
\item Its \emph{bands}, defined in the display below.
\end{itemize}
The bands are
\begin{equation*}
R^{(i)}_{\mathfrak{a}}:=\bigl(Z_{\ell_i+1}\setminus Z''_{\ell_i+1}\bigr)\cap X_{\mathfrak{a}},
\qquad \ell_i:=h_a+i-1 .
\end{equation*}
Its \emph{frozen terms} are
$\mathfrak{F}(\mathfrak{a})=\mathfrak{F}^{\mathbb{C}}(\mathfrak{a})\cup\mathfrak{Q}(\mathfrak{a})$.
These consist of two kinds of term. The first kind is each chain term
$\mathbb{C}^{(i)}_{k_a}(S,\mathcal{S})$ of stage $i\le n$ of $\mathbf{R}_{k_a}$ (as in
Definition~\ref{4.21:def-stages-components} and Lemma~\ref{4.21:lem-base-below-level}, with $k_a$
in place of $k$) with $S\cap R^{(i)}_{\mathfrak{a}}\ne\emptyset$.
The second kind, on the branch (B4) only, is each stored term $v_e\in\mathfrak{Q}(\mathfrak{a})$ (the
collar terms of the attempt, with their enlarged supports $S^+(e)$, as fixed in the statement on the
stored collar terms of an interrupted attempt) whose $S^+(e)$
contains a collar cube of the band of $W_{\mathfrak{a}}$. The attempt $\mathfrak{a}$ is \emph{live}
until the first later completed $\mathbf{R}$-operation, of either class, on the component holding
$X_{\mathfrak{a}}$.
\item A cell is \emph{inherited raised} if it was raised by the preparation of an attempt of an
earlier step $t<k$ that was arrested and is still live. Here still live means that no completed
$\mathbf{R}$-operation, of either class, has acted on the component of the cell since. We put
$\mathcal{W}$ for the set of inherited raised cells. For $x$ in a cell of $\mathcal{W}$ raised by
the arrested attempt $\mathfrak{a}$ we put $t(x):=k_a$.
\item The \emph{raise} at $x$ and its part made during $\mathbf{R}_k$ are
\begin{equation}\label{4.21:eq-arrested-attempts-2}
\mathsf{r}(x):=\operatorname{lev}(x)-b(x),\qquad \mathsf{r}_0(x):=\operatorname{lev}(x)-m_0(x).
\end{equation}
For a cell $\Delta$, $\mathsf{r}(\Delta)$ and $\mathsf{r}_0(\Delta)$ are the maxima of these
functions over $\Delta$.
\end{enumerate}
With these definitions the following hold.
\begin{enumerate}
\item[(i)] $b\le m_0\le\operatorname{lev}$; in particular $\mathsf{r}\ge0$.
\item[(ii)] $\operatorname{lev}(x)>m_0(x)$ if and only if the cell of $x$ is raised during
$\mathbf{R}_k$. Such raises happen inside $Z\setminus Z''_i$ only, with $i\le j(C')$.
\item[(iii)] $m_0(x)>b(x)$ if and only if the cell of $x$ lies in $\mathcal{W}$. The raises of an
arrested attempt persist.
\item[(iv)] No cell is raised twice. Hence $t(x)$ is well defined on $\mathcal{W}$, and
$\operatorname{lev}=m_0\le t(x)$ there. At most one of $\mathsf{r}_0$, $m_0-b$ is non-zero.
\item[(v)] Let $(Y,m)$ be a live second-class region, as in
Lemma~\ref{4.21:lem-live-region-geometry}. On every cell of $Y$, whether in a
class component or in an untreated one, we have $\operatorname{lev}=m_0=b=m$ and
$\mathsf{r}=\mathsf{r}_0=0$; hence $Y\cap\mathcal{W}=\emptyset$. What is used of $Y$ in what follows
is the set $Y$, its strip $Y^{\sharp}\supseteq Y$ (Lemma~\ref{4.21:lem-live-region-geometry}) and
its level $m$.
\end{enumerate}
\end{definition}

\begin{proof}
\emph{Item} (i). By Lemma~\ref{4.21:lem-base-below-level}, in the form
\eqref{4.21:eq-base-below-level-a}, we have $b\le m_0\le\operatorname{lev}$. The reason is this.
The $\mathbf{T}$-steps and the completed $\mathbf{R}$-operations, of either class, set the level of
a cell and its base together. A preparation, whether of $\mathbf{R}_k$ or of an earlier arrested
attempt, raises a level above its base. Nothing lowers a level. Since $b\le\operatorname{lev}$, the
display \eqref{4.21:eq-arrested-attempts-2} gives $\mathsf{r}\ge0$.

\emph{Item} (ii). By definition, $m_0(x)$ is the level of the cell of $x$ before the preparations
of $\mathbf{R}_k$. By the proof of Lemma~\ref{4.21:lem-base-below-level}, during $\mathbf{R}_k$
before its class decision only preparations act on levels, and each of them strictly raises the
levels it changes. Hence $\operatorname{lev}(x)\ne m_0(x)$ exactly when the cell of $x$ has been
raised during $\mathbf{R}_k$, by a running preparation or by a sibling's; in that case
$\operatorname{lev}(x)>m_0(x)$. By \cite[(1.12)]{B-CMP122a}, these raises happen inside
$Z\setminus Z''_i$ only, with $i\le j(C')$.

\emph{Item} (iii). First we show that $m_0>b$ forces membership in $\mathcal{W}$. By
Lemma~\ref{4.21:lem-base-below-level}, levels are set only by three kinds of operation: the
$\mathbf{T}$-steps, the completed $\mathbf{R}$-operations of either class, and the preparations.
The first two kinds set $b$ together with the level.

Suppose $m_0(x)>b(x)$. Then the level of the cell of $x$ in $\mathbb{B}_k$ was set by the
preparation of an attempt of a step $t<k$, and no completed $\mathbf{R}$-operation has acted on the
component of $x$ since. Suppose that attempt had been completed. Its preparation is then followed,
on its own component, by the class decision of the same $\mathbf{R}_t$. That component is healed
whole or thrown back whole \cite[(1.33)]{B-CMP122b}. This decision sets $b=t$,
and $t$ is at least every level raised there. That would contradict $m_0>b$. Hence the attempt is
arrested, and it is live, so the cell lies in $\mathcal{W}$.

The converse implication, that a cell of $\mathcal{W}$ has $m_0>b$, is part of
Lemma~\ref{4.21:lem-base-below-level}. So this one kind of raiser exhausts $\mathcal{W}$.

The raiser is a live arrested attempt $\mathfrak{a}$ with $t(x)=k_a$. It was interrupted after
$n<N(g_{k_a})$ stages \cite[p.~191]{B-CMP122a}, or excluded before any stage or after all of them
\cite[p.~193]{B-CMP122a}. Its raises persist, in the sense that they are not
changed by the subsequent operations. This is the clause on persistence of raises in the
classification of interrupted and excluded attempts.

\emph{Item} (iv). A preparation of step $t$ raises only cells of two kinds:
\begin{equation*}
Z\cap\Omega^{\sim5}_i\subset\Lambda_{i-1}\quad(i\le k_0(t)),
\qquad\text{or}\qquad
\Omega^{\sim7}_{k_0(t)+1}\subset\Lambda_{k_0(t)} .
\end{equation*}
This is shown in the proof of part~(a) of the lemma on the levels of old regions. These cells lie at
least three layers off every older $X_{\mathfrak{a}}$, by the persistence statement just cited. They
also lie at least three layers off every older second-class region, by part~(h) of
Lemma~\ref{4.21:lem-live-region-geometry}. A cell raised by an earlier arrested attempt
$\mathfrak{a}$ lies in $X_{\mathfrak{a}}$. Therefore no later preparation raises it again, and no
cell is raised twice.

It follows that the raiser of a cell of $\mathcal{W}$ is unique, so $t(x)$ is well defined. Such a
cell is not raised during $\mathbf{R}_k$, so $\operatorname{lev}=m_0$ there by item~(ii). Its level
in $\mathbb{B}_k$ is the level assigned by the raise of $\mathfrak{a}$, which persists. That level
$i$ satisfies $h_a<i\le j_a$, and $j_a\le k_a$ by \eqref{4.21:eq-arrested-attempts-1}. Hence
$m_0\le k_a=t(x)$.

Finally, suppose $m_0>b$. Then the cell lies in $\mathcal{W}$ by item~(iii), so it is not raised
during $\mathbf{R}_k$, and $\mathsf{r}_0=0$ by item~(ii). Conversely, suppose $\mathsf{r}_0>0$. Then
the cell is raised during $\mathbf{R}_k$, hence not in $\mathcal{W}$, hence $m_0=b$ by item~(iii).
So at most one of $\mathsf{r}_0$, $m_0-b$ is non-zero. This is the corollary to
Lemma~\ref{4.21:lem-base-below-level}.

\emph{Item} (v). Let $(Y,m)$ be a live second-class region and $x\in Y$. By
\cite[(1.33)]{B-CMP122b}, a component thrown back at $\mathbf{R}_m$ carries the one domain
$Y^{(m)}$ of level $m$ for as long as it lives. No stratum of its completed preparation, and no
structure inside it, keeps a level below $m$ in $\mathbb{B}_k$. The completed class decision of
$\mathbf{R}_m$ on this component enters the definition of the base, so $b\ge m$ on $Y$.

No later $\mathbf{T}_t$ integrates a point of $Y$, for the following reason. By liveness,
$Y\subset Z_{t'}$ for $m\le t'\le k-1$; this is part of \eqref{4.21:eq-base-below-level-a}. By
\eqref{4.21:eq-stages-components-a}, $Z_{t'-1}\subset(\Omega^{\mathbf{T}}_{t'})^c$.
Hence $Y\cap\Omega^{\mathbf{T}}_{t'}=\emptyset$ for $m<t'\le k$, and no $\mathbf{T}$-step after
step $m$ assigns a level to a point of $Y$ or enters $b$ there.

No later arrested attempt's preparation raises $Y$, for the following reason. For an attempt
$\mathfrak{b}$ of a later step whose component holds $Y$, the level $h_{\mathfrak{b}}$ of
$\mathfrak{b}$ satisfies $m\le h_{\mathfrak{b}}$ by part~(g) of
Lemma~\ref{4.21:lem-live-region-geometry}. Its raises keep off older live second-class regions by
part~(h) of that lemma.

$Y$ is not raised during $\mathbf{R}_k$ either. There are two cases. If $Y$ lies in a class
component $C'$, then $m\le h$ by part~(g) of Lemma~\ref{4.21:lem-live-region-geometry}, and
$Y\subset Z_h$ by liveness. Hence
\begin{equation*}
Y\subset Z_h\cap C'\subset Z''_{h+1},
\end{equation*}
whereas $\mathbf{R}_k$ raises inside $Z\setminus Z''_{h+1}$ only. If instead $Y$ lies in an
untreated component, nothing there is raised.

Since a live second-class region is not re-levelled and no completed $\mathbf{R}$-operation has
acted on its component since $\mathbf{R}_m$, the base stays at $m$. Together these give
$\operatorname{lev}=m_0=b=m$ on every cell of $Y$. Therefore $\mathsf{r}=\mathsf{r}_0=0$ there, and
$Y\cap\mathcal{W}=\emptyset$ by item~(iii), because $m_0=b$ on $Y$.
\end{proof}

\begin{definition}[Old blocks, window blocks and window points]\label{4.21:def-old-window-blocks}
We work at the operation $\mathbf{R}_k$ with the stages, the notation and the running, sibling
and excluded components of Definition~\ref{4.21:def-stages-components}; in particular $h=k-N$.
Pockets $(Y,m)$, their strips $Y^{\sharp}$ and the notion of a live pocket are those of
Definition~\ref{4.21:def-strips-geometry}. A \emph{class component} is a running or a sibling
component, and in this section a treated component means a class component. For a block
$\square$ of the cover, $\tilde{\square}$ denotes its probe.
\begin{enumerate}
\item[(a)] The set of \emph{old blocks} is the set $\mathfrak{O}^{+}$ of those blocks $\square$
whose probe $\tilde{\square}$ meets at least one of the following:
\begin{itemize}
\item the strip $Y^{\sharp}$ of a live pocket;
\item a cell contained in $Z_h\cap C'$, for some class component $C'$;
\item an untreated component of $\Lambda_k^c$.
\end{itemize}
\item[(b)] The \emph{window blocks} are the blocks that do not lie in $\mathfrak{O}^{+}$.
\item[(c)] The \emph{window points} are the points of
\begin{equation*}
\Lambda_k\;\cup\bigcup_{C'\ \text{class component}}\bigl(C'\setminus Z_h\bigr)
\end{equation*}
that lie in no strip of a live pocket.
\end{enumerate}
For comparison let $\mathfrak{O}$ be the smaller set of blocks defined as in (a) with the second
item replaced by: a cell of level $<h$ of a class component. The following two facts hold.
\begin{enumerate}
\item[(i)] $\mathfrak{O}\subset\mathfrak{O}^{+}$, and the inclusion is strict; the witness
exhibited in the proof involves no arrested attempt.
\item[(ii)] For every class component $C'$: $Z_h$ is a union of $\pi_h$-cubes, the cells of
$Z_h\cap C'$ have level at most $h$, and every current cell either lies in $Z_h\cap C'$ or is
interior-disjoint from it. Hence the second item of (a) is a condition on the current cells
met by $\tilde{\square}$.
\end{enumerate}
All statements of this section are made and proved for $\mathfrak{O}^{+}$; none of them uses
$\mathfrak{O}$, which serves for comparison only.
\end{definition}

\begin{proof}
(i) The first and third items in the definitions of $\mathfrak{O}$ and $\mathfrak{O}^{+}$
coincide, so it suffices to compare the second items. Let $C'$ be a class component and let
$\Delta$ be a cell of level $<h$ of $C'$. By the admissible geometry of
Definition~\ref{4.21:def-strips-geometry}, display~\eqref{4.21:eq-strips-geometry-a}, such a
cell lies in $(\Omega^{(k)}_h)^c\cap C'$, and the same display gives
\begin{equation*}
(\Omega^{(k)}_h)^c\cap C'\subset Z_h\cap C'.
\end{equation*}
Thus $\Delta$ is a cell contained in $Z_h\cap C'$, and a probe $\tilde{\square}$ that meets
$\Delta$ satisfies the second item of (a). Hence every block of $\mathfrak{O}$ lies in
$\mathfrak{O}^{+}$.

For strictness, again by display~\eqref{4.21:eq-strips-geometry-a},
\begin{equation*}
Z_h\cap C'\supset(\Omega_h\setminus\Lambda_h)\cap C',
\end{equation*}
and $(\Omega_h\setminus\Lambda_h)\cap C'$ consists of two layers of cubes of side $M\check{R}_h$
made of unraised cells of level $h$. Fix a block $\square$ of the cover over these two layers,
that is, a block whose probe $\tilde{\square}$ is contained in their interior. Every cell met by
$\tilde{\square}$ is then an unraised cell of $C'$ of level exactly $h$, not $<h$; hence
$\tilde{\square}$ meets a cell contained in $Z_h\cap C'$, so that $\square\in\mathfrak{O}^{+}$,
and it meets no cell of level $<h$ of a class component. Since $\tilde{\square}\subset C'$, it
meets no untreated component of $\Lambda_k^c$, these being disjoint from the class component
$C'$. Finally, a probe contained in the two unraised layers of a class component meets no strip
of a live pocket. Therefore $\tilde{\square}$ satisfies none of the three conditions defining
$\mathfrak{O}$, and $\square\in\mathfrak{O}^{+}\setminus\mathfrak{O}$.

Blocks over the raised top zone of an arrested attempt lying inside $C'$, whose step $k_a$ and
top-zone level $j_a$ satisfy $j_a=k_a=h$, would also belong to $\mathfrak{O}^{+}\setminus
\mathfrak{O}$. Whether such attempts occur depends on the rule by which an arrested component is
treated again: under the bound on the step of an arrest used in the statement on the structure
of arrested attempts they occur, whereas in the construction of \cite[pp.~385--386]{B-CMP122b},
after an arrest at step $k_a$ the next treatment of the same component takes place after the
number of steps written $R_{k_a+1}$ there, so that $k_a\le h-1$ and no such attempt exists. The
strictness in (i) rests on the two layers above, and nothing in this section uses these further
blocks.

(ii) By display~\eqref{4.21:eq-strips-geometry-a}, read at step $h$, the set $Z_h$ is a union
of $\pi_h$-cubes. By the same display, $Z_h\subset\Omega_{h+1}^c$. Moreover
$Z_h\cap C'\subset Z''_{h+1}$, and this set is never raised during the present
$\mathbf{R}_k$: the zones raised during $\mathbf{R}_k$ lie in sets $Z\setminus Z''_i$, where $i$
runs over the levels at which zones are raised during $\mathbf{R}_k$, and for each such $i$
\begin{equation*}
Z\setminus Z''_i\subset Z\setminus Z''_{h+1},
\end{equation*}
that is, $Z\cap Z''_{h+1}\subset Z''_i$.
Consequently the cells of $Z_h\cap C'$ keep the levels they have in $\Omega_{h+1}^c$, namely
levels at most $h$. The partition into current cells and the partition $\pi_h$ are nested
partitions. Since $Z_h$ is a union of $\pi_h$-cubes and the cells of $Z_h\cap C'$ have level
at most $h$, a current cell either lies in $Z_h\cap C'$ or is interior-disjoint from it.
Therefore the condition that $\tilde{\square}$ meets a cell contained in $Z_h\cap C'$ is a
condition on the current cells met by $\tilde{\square}$.
\end{proof}

\begin{lemma}[What a probe meeting a live second-class region meets]\label{4.21:lem-probe-live-region}
Work at the step $\mathbf{R}_k$, with the current level $\operatorname{lev}$, the level $m_0$ in
$\mathbb{B}_k$, the base $b$ and the set $\mathcal{W}$ of inherited raised cells as in
Lemma~\ref{4.21:lem-base-below-level}. Let $(Y,m)$ be a second-class region live at
$\mathbf{R}_k$, with strip $Y^{\sharp}$ as in Definition~\ref{4.21:def-strips-geometry}. Distances
are measured in level-$m$ units. For a probe $\mathsf{Q}$ we write $n_0$ for the current level of
the probe's own cell. Then the following hold.
\begin{enumerate}
\item[(i)] A cell of $Y$ is never a window point. It is never raised by this $\mathbf{R}_k$ or by
a live arrested attempt. It is met by the probe of no window block
(Definition~\ref{4.21:def-old-window-blocks}).
\item[(ii)] Every cell within distance $5$ of a point of $Y$ is either a cell of $Y$ or a cell of
$\Lambda^{\mathbf{T}}_m$ beside $Y$, and on all these cells
\begin{equation*}
\operatorname{lev}=m_0=b=m .
\end{equation*}
\item[(iii)] A probe $\mathsf{Q}$ meeting a cell of $Y$ has $n_0=m$ and meets no cell of
$\mathcal{W}$.
\end{enumerate}
\end{lemma}

\begin{proof}
Since $(Y,m)$ is live at $\mathbf{R}_k$, we have $m\le k-1$.

\emph{Proof of \textup{(i)}.} By Definition~\ref{4.21:def-strips-geometry} the strip
$Y^{\sharp}$ contains $Y$. By Definition~\ref{4.21:def-old-window-blocks}, window points lie in no
live strip. Hence no point of $Y$ is a window point.

By part~(c) of Lemma~\ref{4.21:lem-base-below-level}, $\mathsf{r}=\mathsf{r}_0=0$ at every point
of $Y$, and $Y\cap\mathcal{W}=\emptyset$. Since $\mathsf{r}_0=0$, no cell of
$Y$ is raised during this $\mathbf{R}_k$. Since $Y\cap\mathcal{W}=\emptyset$, no cell of $Y$ is
raised by a live arrested attempt.

A window block is, by Definition~\ref{4.21:def-old-window-blocks}, a block not in
$\mathfrak{O}^{+}$. Its probe therefore meets the strip of no live second-class region, so in
particular it meets no cell of $Y\subset Y^{\sharp}$.

\emph{Proof of \textup{(ii)}.} The set $Y$ is a whole component of $Z^{\mathbf{T}}_m$
(Definition~\ref{4.21:def-strips-geometry}). The other
components of $Z^{\mathbf{T}}_m$ lie at distance at least $\check{R}_m$ from $Y$, by the separation
condition of the admissible geometry (Definition~\ref{4.21:def-strips-geometry}). By part~(a) of
Lemma~\ref{4.21:lem-live-region-geometry},
\begin{equation*}
\check{R}_m\ge 38>5 .
\end{equation*}

Let $x$ be a point within distance $5$ of $Y$ with $x\notin Y$. Then $x$ lies in no other
component of $Z^{\mathbf{T}}_m$: such a component would otherwise come within distance $5$ of $Y$.
Hence $x\notin Z^{\mathbf{T}}_m$, that is,
\begin{equation*}
x\in\Lambda^{\mathbf{T}}_m\subset\Omega^{\mathbf{T}}_m .
\end{equation*}
The base $b(x)$ is the largest $t\le k$ with $x\in\Omega^{\mathbf{T}}_t$ or with $x$ in a
component carried through a completed class decision at $\mathbf{R}_t$. Since $m\le k$, this gives
$b(x)\ge m$.

By part~(c) of Lemma~\ref{4.21:lem-live-region-geometry}, every cell within distance
$(5/2)\check{R}_m$ of $Y^{\sharp}$ has current level at most $m$. This applies to the cell
containing $x$, because $5<(5/2)\check{R}_m$ and $Y\subset Y^{\sharp}$. Hence
$\operatorname{lev}(x)\le m$. Part~(a) of Lemma~\ref{4.21:lem-base-below-level} states
$b\le m_0\le\operatorname{lev}$, so
\begin{equation*}
m\le b(x)\le m_0(x)\le\operatorname{lev}(x)\le m .
\end{equation*}
All members are therefore equal to $m$.

On the points of $Y$ itself, $\operatorname{lev}=m_0=b=m$ by part~(c) of
Lemma~\ref{4.21:lem-base-below-level}. This proves~(ii).

\emph{Proof of \textup{(iii)}.} Let the probe $\mathsf{Q}$ meet a cell of $Y$. By part~(i1) of
Lemma~\ref{4.21:lem-live-region-geometry}, the current level of its own cell is at most $m$. A
probe has side at most $5$ units of its own level, and a unit of level at most $m$ is at most one
level-$m$ unit; hence the side of $\mathsf{Q}$ is at most $5$ level-$m$ units. Since $\mathsf{Q}$
meets $Y$, it lies in the neighbourhood of $Y$ described in~(ii).

The probe's own cell lies in that neighbourhood, so by~(ii) it has level $m$; that is, $n_0=m$.

By~(ii), no cell of that neighbourhood has $m_0>b$. By part~(b) of
Lemma~\ref{4.21:lem-base-below-level}, a cell has $m_0>b$ if and only if it was raised by the
preparation of an arrested attempt of an earlier step that is still live, that is, if and only if
it belongs to $\mathcal{W}$. Hence $\mathsf{Q}$ meets no cell of $\mathcal{W}$.
\end{proof}

\begin{lemma}[What a window probe meets]\label{4.21:lem-window-probe}
Let the step $k$, the level $h$, the class components of $\Lambda_k^c$, the live pockets and their
strips, the window blocks and the window points be as in
Definition~\ref{4.21:def-old-window-blocks}, and let the current level $\operatorname{lev}$, the
original level $m_0$ and the base $b$ of a cell, and the raised cells, be as in
Definition~\ref{4.21:def-arrested-attempts}. A \emph{treated} component is a class component,
running or sibling. Let $\square_0$ be a window block and $\tilde{\square}_0$ its probe, and let
$\Delta$ denote a cell met by $\tilde{\square}_0$. Then $\Delta$ consists of window points, and
\begin{equation}\label{4.21:eq-window-probe-a}
\begin{aligned}
&\text{either}\quad \Delta\subset\Lambda_k, \quad\text{and then}\quad
  \operatorname{lev}=b=k \ \text{on } \Delta,\\
&\text{or}\quad \Delta\subset C'\setminus Z_h\subset\Lambda_h\cap C'
  \ \text{for a treated } C', \quad\text{and then}\quad b\ge h .
\end{aligned}
\end{equation}
In the second case $\operatorname{lev}=m_0=b$ on $\Delta$, unless $\Delta$ was raised during the
operation $\mathbf{R}_k$ of the present step. In particular $\tilde{\square}_0$ meets no cell of
$\mathcal{W}$, the set of raised cells inherited from the preparation of a live arrested attempt of
an earlier step (Definition~\ref{4.21:def-arrested-attempts}).
\end{lemma}

\begin{proof}
Since $\square_0$ is a window block, Definition~\ref{4.21:def-old-window-blocks} says that its
probe $\tilde{\square}_0$ meets no strip of a live pocket, no cell contained in $Z_h\cap C'$ for a
treated component $C'$, and no untreated component of $\Lambda_k^c$. Hence a cell $\Delta$ met by
$\tilde{\square}_0$ either lies in $\Lambda_k$, or lies in a treated component $C'$ without lying in
$Z_h\cap C'$. In the latter case $\Delta$ lies in $C'\cap\Lambda_h$, since $Z_h=\Lambda_h^c$. Thus
$\Delta$ lies in $\Lambda_k\cup\bigcup_{C'}(C'\setminus Z_h)$, the union over treated components.
Moreover, a cell meeting the strip of a live pocket lies in that strip, by clause~(b) of
Lemma~\ref{4.21:lem-live-region-geometry}; since the probe meets no live strip, $\Delta$ meets
none, and so $\Delta$ consists of window points.

We show $b\ge h$ in the second case. The set $C'\cap\Lambda_h$ received level at least $h$ at
step $h$, either through the operation $\mathbf{T}_h$ or through a healing at $\mathbf{R}_h$.
Hence $b\ge h$ on every cell of $C'\setminus Z_h$.

We show that no cell of $C'\setminus Z_h$ was raised at an earlier step. First, inside $C'$ no
attempt of a step $t$ with $h<t<k$ took place. An arrest or an ejection at such a step would be new
large field inside $C'$, and this is excluded by condition~(ii) of
\cite[p.~177]{B-CMP122a}. For the same reason Definition~\ref{4.21:def-arrested-attempts} admits
no arrested attempt of such a step inside $C'$. In the same direction, clause~(g) of
Lemma~\ref{4.21:lem-live-region-geometry} gives $m\le h$ for the level $m$ of a live pocket
$(Y,m)$ inside $C'$, so that no live pocket inside $C'$ originates at a step $t$ with $h<t<k$.

Secondly, a cell carrying a raise inherited from an earlier step was raised by a live arrested
attempt of that step (Definition~\ref{4.21:def-arrested-attempts}); if the cell lies in $C'$, that
step is at most $h$ by the first point. So let $\mathfrak{a}$ be a live arrested attempt of a step
$t=k_a\le h$ that raised cells,
and let $\Omega^{\mathbf{T}}_t$ be the domain produced by the operation $\mathbf{T}_t$. By
Definition~\ref{4.21:def-arrested-attempts}, the raised cells of $\mathfrak{a}$ lie in its piece
\begin{equation}\label{4.21:eq-window-probe-1}
W_{\mathfrak{a}}=X_{\mathfrak{a}}\cap(\Omega^{\mathbf{T}}_t)^c\subset X_{\mathfrak{a}}
\subset Z_t\subset Z_h .
\end{equation}
Here the collar $\mathfrak{c}_{\mathfrak{a}}=X_{\mathfrak{a}}\cap\Omega^{\mathbf{T}}_t$ consists of
two layers of unraised level-$t$ cubes of side $M\check{R}_t$, by the construction of the collar
in the preparation of an arrested attempt. That $\mathfrak{a}$ is live means that the
component of $X_{\mathfrak{a}}$ was not completed at $\mathbf{R}_t$.

By \eqref{4.21:eq-window-probe-1}, no cell of $C'\setminus Z_h$ is among the raised cells of such
an $\mathfrak{a}$. Together with the absence of attempts of steps $t$ with $h<t<k$ inside $C'$, this
gives $m_0=b$ on $C'\setminus Z_h$. Consequently, on a cell of $C'\setminus Z_h$ that was not raised
during $\mathbf{R}_k$, we have $\operatorname{lev}=m_0=b$.

In the first case of \eqref{4.21:eq-window-probe-a}, the cells of $\Lambda_k$ lie outside
$\Lambda_k^c=Z_k$, which contains the regions $Z_t$, $t\le k$, in which raised cells lie, and are
therefore never raised; this gives $\operatorname{lev}=b=k$ on them.

Finally, every cell of $\mathcal{W}$ was raised at an earlier step, by a live arrested attempt. By
the preceding discussion of the two cases of \eqref{4.21:eq-window-probe-a}, no cell met by
$\tilde{\square}_0$ was so raised. Hence $\tilde{\square}_0$ meets no cell of $\mathcal{W}$.
\end{proof}

\begin{lemma}[Stored terms cease at a completed large-field operation;
{\cite[p.~369, p.~377]{B-CMP122b}}]\label{4.21:lem-terms-cease}
Let $t$ be a step and let $C$ be a component that is carried through a completed site
(L) of the operation $\mathbf{R}_t$. The component $C$ may belong to either class. The
sites of $\mathbf{R}_t$, the stages, the components and their classes are those of
Definition~\ref{4.21:def-stages-components}. Let
$F$ be a stored term of the format class $\mathfrak{F}$, with reading set $X_F$. Then
\begin{equation}\label{4.21:eq-terms-cease-a}
X_F\cap C\neq\emptyset\quad\Longrightarrow\quad F\ \text{ceases at } t .
\end{equation}
Here a stored term \emph{ceases} at $t$ in the sense of the format of stored terms.
\end{lemma}

For every stored term of $\mathfrak{F}$ other than the two kinds named in the next
sentence, the implication \eqref{4.21:eq-terms-cease-a} is stated in
\cite[p.~369]{B-CMP122b} and \cite[p.~377]{B-CMP122b}, and it is used here as stated
there. The two kinds excepted are the frozen terms of an arrested attempt $\mathfrak{a}$
and the terms of $\mathfrak{Q}(\mathfrak{a})$.
For them \eqref{4.21:eq-terms-cease-a} holds by clauses (O2) and (O5) of the definition
of the frozen structure of an arrested attempt. By the first of these clauses a frozen
term is absorbed at the completed site and ceases there; by the second the terms of
$\mathfrak{Q}(\mathfrak{a})$ are dissolved there.

\emph{Reading the condition that $A_F$ meets a set.} Where the reading set $X_F$ of a
stored term $F$ must be shown to meet a set, this is done through the position of its
label $A_F$, and the condition is read as follows. Two cases arise.

First, let $F$ be a chain term, or a stored term of the other kind that clause (f1) of
the format of stored terms treats together with the chain terms. For these clause (f1)
gives
\begin{equation}\label{4.21:eq-terms-cease-1}
X_F\supseteq A_F .
\end{equation}
Hence if $A_F$ meets a set $E$, so does $X_F$.

Second, let $F$ be a term attached to an element $e$. By the statement fixing the format
of the terms attached to an element, and by the definition of $A_e$ as a hull, its
reading set and its label are
\begin{equation}\label{4.21:eq-terms-cease-2}
X_F:=S^+(e)\subset A_F=A_e=[S^+(e)]_{m_0(e)} ,
\end{equation}
where $[\,\cdot\,]_{m_0(e)}$ is the $\pi_{m_0(e)}$-hull. Inclusion
\eqref{4.21:eq-terms-cease-1} may then fail, and the following is used in its place.
By the definition of the hull, $A_e$ is the union of those cubes of $\pi_{m_0(e)}$ that
meet $S^+(e)$. Hence every cube of $\pi_{m_0(e)}$ contained in $A_e$ holds a point of
$S^+(e)=X_F$.

The second case arises in three places of the proofs of this chapter. In each, a set $E$
must be shown to meet $X_F$, and $E$ is one of the following:
\begin{itemize}
\item a component that contains a point $y$ of $A_F$;
\item a region met by a cube $Q_F$ of $\pi_{m_0(e)}$ contained in $A_F$;
\item a region $Y$ met by a probe cube $Q$ that is a cube of $\pi_{m_0(e)}$ contained
in $A_e$.
\end{itemize}
For a term attached to $e$, one exhibits at the use site a cube $Q$ of $\pi_{m_0(e)}$
with $Q\subset A_e$ that meets $E$. Here $Q$ is, respectively, the cube of
$\pi_{m_0(e)}$ in $A_e$ containing $y$, which meets $E$ because $y\in Q\cap E$; the cube
$Q_F$, which meets $E$ by hypothesis; or the probe cube, which meets $E$ by hypothesis.
By the geometric condition fixing the levels of these sets, each of the three sets $E$
is a union of cubes of $\pi_{t'}$ for a level $t'\ge m_0(e)$. By the nesting of the
partitions, every cube of $\pi_{t'}$ with $t'\ge m_0(e)$ is a union of cubes of
$\pi_{m_0(e)}$. Hence $E$ is a union of cubes of $\pi_{m_0(e)}$. A cube of
$\pi_{m_0(e)}$ that meets $E$ therefore lies inside $E$, so that $Q\subset E$ in each of
the three cases. Since $Q\subset A_e$, the cube
$Q$ holds a point of $S^+(e)$, and this point lies in $E$. Thus $X_F\cap E\neq\emptyset$.
When $E$ is a component carried through a completed site (L) of $\mathbf{R}_t$, this is
the hypothesis of Lemma~\ref{4.21:lem-terms-cease}, which then applies to $F$.

\begin{lemma}[The format class and its confined and top old terms]\label{4.21:lem-format-class}
Work at the step $k$, at the large-field operation $\mathbf{R}_k$, in the notation of
Definition~\ref{4.21:def-stages-components}, Definition~\ref{4.21:def-arrested-attempts} and
Definition~\ref{4.21:def-strips-geometry}; in particular $h=k-N$, and for an arrested attempt
$\mathfrak{a}$ of step $k_a$ we use its region $X_{\mathfrak{a}}$, its collar
$\mathfrak{c}_{\mathfrak{a}}$, its piece $W_{\mathfrak{a}}=X_{\mathfrak{a}}\setminus\mathfrak{c}_{\mathfrak{a}}$,
its bands $R^{(i)}_{\mathfrak{a}}$, its frozen label $\mathfrak{B}_{\mathfrak{a}}$ and its level $h_a$.
\begin{enumerate}
\item[(a)] The \emph{format class} $\mathfrak{F}$ consists of the following stored terms.
\begin{itemize}
\item[$(\alpha)$] The boundary terms $\mathbf{B}^{(\ell)}$, $\mathbf{B}'^{(\ell)}$, plain or wrapped, with label
$A_F\in\mathbf{D}_\ell$ of one level $\ell$.
\item[$(\beta)$] The frozen chain terms $c=\mathbb{C}^{(i)}_{k_a}(S,\mathcal{S})$ of live arrested attempts
$\mathfrak{a}$ of earlier steps ($k_a<k$; these terms are called \emph{old}), the members of part $(\alpha)$
of the statement introducing the terms $\mathfrak{Q}(\mathfrak{a})$ of an arrested attempt included. Their
label $S=S_c\in\mathbb{D}_{m(S)}$ has one
level $m=m(S)$: it is the smallest domain of $\mathbf{D}_m$ containing the region of the term, all regions
determining the same domain being resummed, as in \cite[p.~189]{B-CMP122a}, and by the single-level filing
convention for the labels of chain terms. Moreover
$h_a+1\le m\le k_a$ and $S\cap R^{(i)}_{\mathfrak{a}}\neq\emptyset$.
\item[$(\gamma)$] The terms $c=\mathbb{C}^{(i)}_k[C'](S,\mathcal{S})$ of the chains of $\mathbf{R}_k$
itself, where $C'$ is a component of the class (a running component or a sibling), $i\le n(C')$,
$S\in\mathbb{D}_{m(S)}$ and $h+1\le m(S)\le k$; they are filed, by \cite[(1.63)]{B-CMP122a}, in the
label of stage $i$.
\item[$(\delta)$] The terms $v_e\in\mathfrak{Q}(\mathfrak{a})$ of an attempt $\mathfrak{a}$ arrested by a
sibling exit: either an old such attempt ($k_a<k$), or a sibling of the present step with $n^{+}\le n$
($k_a=k$). The label of $v_e$ is $A_e=[S^{+}(e)]_{m_0(e)}\in\mathbb{D}_{m_0(e)}$, of one level
$m_0(e)\le k_a$, and $v_e$ is born at step $k_a$.
\end{itemize}
The \emph{own cubes} of a term are the cubes of its label at the label's one level: the $\pi_\ell$-cubes of
$A_F$, the $\pi_m$-cubes of $S$, the $\pi_{m_0(e)}$-cubes of $A_e$.
\item[(b)] For $\ell\le k$ put
\begin{equation*}
\Omega^{(\ell)}_\ell:=\Omega^{\mathbf{T}}_\ell\cup\bigl(\text{the components healed at }\mathbf{R}_\ell\bigr),
\end{equation*}
a component being \emph{healed} at $\mathbf{R}_\ell$ when its large-field factor is absorbed there back into
the small-field format and the region is re-filed into the regular expansion;
let $Z_\ell$ be the large-field region after $\mathbf{R}_\ell$, and let $Z_\ell^{\boxplus}$ be $Z_\ell$
enlarged by one layer of cubes of side $M\check{R}_\ell$. For a term $F$ of level $\ell$, born at step
$\ell$, with label $A_F$, and for a term $c$ of kind $(\beta)$ or $(\delta)$ of the attempt $\mathfrak{a}$,
with filing level $m:=m(S)$, resp.\ $m:=m_0(e)$, so that $m\le k_a$, and label $S$ ($S=S_c$, resp.\
$S=A_e$), we consider the conditions, which come from the format statement for stored terms together with
\cite[(2.41)]{B-CMP119} and \cite[(1.63)]{B-CMP122a}, and the
splitting
\begin{equation}\label{4.21:eq-format-class-a}
\begin{aligned}
&(\mathrm{f}0)_\ell:\ \ A_F\cap\Omega^{(\ell)}_\ell\neq\emptyset,\ \text{or }A_F\text{ meets two
components of }Z^{\mathbf{T}}_\ell\\
&\qquad\qquad\text{and crosses }\Lambda^{\mathbf{T}}_\ell\text{ between them};\\
&(\mathrm{f}4)_\ell:\ \ A_F\text{ has a cube }Q_F\text{ (its anchor) in }(Z_\ell^{\boxplus})^{(1)};\\
&c\text{ is confined}\iff m<k_a,\ \text{or }m=k_a\text{ and }S\subset\bigl(\Omega^{(k_a)}_{k_a}\bigr)^{c};\\
&c\text{ is a top term}\iff m=k_a\text{ and }S\cap\Omega^{(k_a)}_{k_a}\neq\emptyset .
\end{aligned}
\end{equation}
The condition $(\mathrm{f}0)_\ell$ is used in the set form displayed; it is not equivalent to the condition
that $A_F$ contain a point of level at least $\ell$ after step $\ell$.
\end{enumerate}
Then the following hold.
\begin{enumerate}
\item[(i)] If $c$ is confined, then $S\subset X_{\mathfrak{a}}\subset Z_{k_a}$, and the base satisfies
$b(x)\le m$ at every point $x\in S$, at the birth of $c$ and at every later site at which $c$ is live.
\item[(ii)] If $c$ is a top term, then $c$ is of level $\ell:=k_a$ and is born at step $\ell$; it satisfies
$(\mathrm{f}0)_\ell$ by definition, and $(\mathrm{f}4)_\ell$ with anchor $Q_c$ a cube of $S$ in the band
$R^{(i)}_{\mathfrak{a}}\subset X_{\mathfrak{a}}\subset Z_{k_a}$ for a chain term, resp.\ a cube of $A_e$
over the collar cube $q\in\mathsf{b}(e)$ of the band of $W_{\mathfrak{a}}$ for a term of kind $(\delta)$.
\end{enumerate}
\begin{enumerate}
\item[(c)] From here on the terms of $\mathfrak{F}$ are grouped as follows.
\begin{itemize}
\item[(c1)] The \emph{terms of boundary type} are the boundary terms $\mathbf{B}^{(\ell)}$,
$\mathbf{B}'^{(\ell)}$ with $\ell\le k$, the top frozen chain terms ($k_a<k$) and the top terms of kind
$(\delta)$ ($k_a\le k$). Each is a term $F$ of a level $\ell$, born at step $\ell$, satisfying
$(\mathrm{f}0)_\ell$ and $(\mathrm{f}4)_\ell$ of \eqref{4.21:eq-format-class-a} and the cessation property
of Lemma~\ref{4.21:lem-terms-cease}.
\item[(c2)] The \emph{confined old terms} are the confined frozen chain terms and the confined terms of kind
$(\delta)$ with $k_a<k$. Their labels lie in
\begin{equation*}
X_{\mathfrak{a}}\subset Z_{k_a}\subset Z_h\cap\bigl(\text{the component of the class holding }
\mathfrak{a}\bigr),
\end{equation*}
or in a component not yet treated.
\item[(c3)] The \emph{local sibling terms} are the confined terms of kind $(\delta)$ of the sibling $C''$ of
the present step that holds the point under consideration ($k_a=k$, $n^{+}\le n$); their labels lie in
$C''$, possibly in $C''\setminus Z_h$.
\end{itemize}
The chain terms of a sibling, whatever their filing level, are of kind $(\gamma)$: they form the local chain
of $C''$, and those filed at level $k$ may reach other components. They are estimated with the chain terms of
the present step and not with the old terms, and they are not top frozen terms.
\end{enumerate}
\end{lemma}

\begin{proof}
\emph{Confined terms with $m<k_a$.} We first show
\begin{equation}\label{4.21:eq-format-class-1}
S\subset\Omega_{m+1}^{c},
\end{equation}
with the unprimed domain $\Omega_{m+1}$ (not $\Omega''_{m+1}$). For a chain term, \eqref{4.21:eq-format-class-1}
is the clause of \cite[(1.63)]{B-CMP122a} itself, written in the label of the stage $i$ at which the term is
born; there the index $m+1$ satisfies $m+1\ge\ell_i+2$, where $\ell_i$ is the level integrated at stage $i$
(the level $h_a+i-1$, in the indexing of Definition~\ref{4.21:def-stages-components}, where stage $n+1$
integrates level $h+n$), and has not yet been relabelled. For a term $v_e$ of kind $(\delta)$, the level
assigned by the frozen label
$\mathfrak{B}_{\mathfrak{a}}$ is at most $m_0(e)=m$ on $S^{+}(e)$, by the definition of $m_0(e)$. Write
$\Omega^{\mathfrak{B}_{\mathfrak{a}}}_{m+1}$ for the set of points of $\mathfrak{B}_{\mathfrak{a}}$-level
at least $m+1$, so that the set of points of $\mathfrak{B}_{\mathfrak{a}}$-level at most $m$ is
$(\Omega^{\mathfrak{B}_{\mathfrak{a}}}_{m+1})^{c}$. On $X_{\mathfrak{a}}$ the set
$\Omega^{\mathfrak{B}_{\mathfrak{a}}}_{m+1}$ is either $\Omega''_{m+1}$ or $\Omega_{m+1}$; off the class it
is $\Omega_{m+1}$. In the first case, by \cite[(1.11)--(1.12)]{B-CMP122a},
\begin{equation*}
\Omega''_{m+1}\cap Z=Z\setminus Z''_{m+1}\supseteq Z\cap\Omega_{m+1},
\end{equation*}
$Z$ the class. Hence in every case $\Omega^{\mathfrak{B}_{\mathfrak{a}}}_{m+1}\supseteq\Omega_{m+1}$, that is
\begin{equation*}
\bigl(\Omega^{\mathfrak{B}_{\mathfrak{a}}}_{m+1}\bigr)^{c}\subset\Omega_{m+1}^{c}.
\end{equation*}
Both complements are unions of $\pi_m$-cubes. Since $S^{+}(e)$ lies in
$(\Omega^{\mathfrak{B}_{\mathfrak{a}}}_{m+1})^{c}$, a union of $\pi_m$-cubes, so does its $\pi_m$-hull
$A_e=[S^{+}(e)]_m$; this gives \eqref{4.21:eq-format-class-1} for $S=A_e$.

Next, $\Omega_{m+1}^{c}\subset Z_{m+1}$ by the first condition of the admissible geometry,
\eqref{4.21:eq-strips-geometry-a}, and $Z_{m+1}\subset Z_{k_a}$ since $m+1\le k_a$. So $S\subset Z_{k_a}$.
The set $S$ is connected: for a chain term it is a domain of $\mathbb{D}_m$, and for a term of kind
$(\delta)$ the set $S^{+}(e)$ is connected by the statement introducing the terms
$\mathfrak{Q}(\mathfrak{a})$, and $A_e$ is a domain of $\mathbb{D}_m$. The set $S$ meets $X_{\mathfrak{a}}$
(for a chain term through $S\cap R^{(i)}_{\mathfrak{a}}\neq\emptyset$ and
$R^{(i)}_{\mathfrak{a}}\subset X_{\mathfrak{a}}$; for a term of kind $(\delta)$ through the cube of $A_e$
over the collar cube $q\in\mathsf{b}(e)$ of the band of $W_{\mathfrak{a}}$, which the statement introducing
the terms $\mathfrak{Q}(\mathfrak{a})$ provides, and which lies in $X_{\mathfrak{a}}$),
and $X_{\mathfrak{a}}$ is a component of $Z_{k_a}$. A connected subset of $Z_{k_a}$ meeting the component
$X_{\mathfrak{a}}$ of $Z_{k_a}$ is contained in it; hence $S\subset X_{\mathfrak{a}}$.

The level of the label is at most $m$ on $\Omega_{m+1}^{c}$: its raised cells have levels at most
$\ell_i+1\le m$ (since $m+1\ge\ell_i+2$) for a chain term, and at most $m$ by definition for a term of kind
$(\delta)$. By \eqref{4.21:eq-format-class-1}, $b\le m$ on $S$ at the birth of $c$. Afterwards, for every
later step $t$ one has $X_{\mathfrak{a}}\subset Z_{t-1}$, and by \eqref{4.21:eq-stages-components-a},
$\Omega^{\mathbf{T}}_t\cap Z_{t-1}=\emptyset$, so no operation $\mathbf{T}_t$ integrates a point of
$X_{\mathfrak{a}}$; and a completed operation $\mathbf{R}$ on the component holding $X_{\mathfrak{a}}$ ends
$c$, by Lemma~\ref{4.21:lem-terms-cease}. Therefore $b$ stays at most $m$ on $S$ while $c$ is live. This
proves~(i) when $m<k_a$.

\emph{Confined terms with $m=k_a$.} The small-field region $\Lambda_{k_a}$ after $\mathbf{R}_{k_a}$, whose
complement is $Z_{k_a}$, lies in $\Omega^{(k_a)}_{k_a}$; hence
$(\Omega^{(k_a)}_{k_a})^{c}\subset Z_{k_a}$, and $S\subset Z_{k_a}$ by the definition of a confined term. As
$S$ is connected and meets $X_{\mathfrak{a}}$, the argument of the preceding paragraph gives
$S\subset X_{\mathfrak{a}}$, and moreover $S\subset X_{\mathfrak{a}}\setminus\mathfrak{c}_{\mathfrak{a}}
=W_{\mathfrak{a}}$, where $b\le k_a-1<m$. The argument for later sites in the preceding paragraph used only
$S\subset X_{\mathfrak{a}}$, and applies unchanged. This proves~(i).

\emph{Top terms.} For a chain term, the clause $X\cap(Z_{j'+1}\setminus Z''_{j'+1})\neq\emptyset$ of
\cite[(1.63)]{B-CMP122a} gives the cube of $S$ in the band $R^{(i)}_{\mathfrak{a}}$. For a term of kind
$(\delta)$, the statement introducing the terms $\mathfrak{Q}(\mathfrak{a})$ gives the collar cube
$q\in\mathsf{b}(e)$. Both are born at $\mathbf{R}_{k_a}$, and $(\mathrm{f}0)_{k_a}$ is the definition of a
top term. This proves~(ii).

\emph{Proof of (c).} We verify $(\mathrm{f}0)_\ell$ for the boundary terms of part~(c1), and first the
remark in part~(b): the set condition is not equivalent to the
condition that $A_F$ contain a point of level at least $\ell$ after step $\ell$, since for an arrested attempt
whose top zone has been raised, that zone has level $k_a$ inside $Z_{k_a}$ and lies outside
$\Omega^{(k_a)}_{k_a}$. For a boundary term $\mathbf{B}^{(\ell)}$ born at the operation $\mathbf{T}$, the
condition is the clause $X\cap\Omega_\ell\neq\emptyset$ of \cite[(2.41)]{B-CMP119}, $X$ the region of the
term. For a
boundary term $\mathbf{B}'^{(\ell)}$ of the class born at the operation $\mathbf{R}_\ell$, with region $X$
hanging on the pockets $Y_i$, one has $X\not\subset\bigcup_i Y_i^{\sharp}$, and a cube of $X$ off the strips
$Y_i^{\sharp}$ lies in a healed component, or in $\Lambda^{\mathbf{T}}_\ell\subset\Omega^{\mathbf{T}}_\ell$,
or in a component of $Z^{\mathbf{T}}_\ell$ other than the one containing the $Y_i$ it hangs on; in the last
case $A_F$, being connected, meets both components and crosses $\Lambda^{\mathbf{T}}_\ell$ between them. The
boundary terms satisfy $(\mathrm{f}4)_\ell$ by the format statement for stored terms, together with
\cite[(2.41)]{B-CMP119} and \cite[(1.63)]{B-CMP122a}. For
the top terms the conditions $(\mathrm{f}0)_\ell$ and $(\mathrm{f}4)_\ell$, with $\ell=k_a$ and birth at
step $\ell$, are part~(ii). Every term of boundary type belongs to $\mathfrak{F}$, so the cessation property
is Lemma~\ref{4.21:lem-terms-cease}. This proves~(c1).

For a confined old term, part~(i) gives $S\subset X_{\mathfrak{a}}\subset Z_{k_a}$, and the inclusion of
$Z_{k_a}$ in $Z_h$ intersected with the component of the class holding $\mathfrak{a}$, or else the location
of $X_{\mathfrak{a}}$ in a component not yet treated, is the statement locating live arrested attempts in
the class, together with its clause~(ii). This proves~(c2). For a local sibling term, part~(i) gives
$S\subset X_{\mathfrak{a}}$, the region of the arrested attempt, which for a sibling of the present step is
the sibling $C''$ (Definition~\ref{4.21:def-arrested-attempts}); this proves~(c3).

Finally, the chains of a sibling $C''$ of the present step are chains of $\mathbf{R}_k$ on a component of the
class, so their terms are of kind $(\gamma)$ by part~(a), whatever their filing level; the top frozen chain
terms are of kind $(\beta)$, with $k_a<k$, so a sibling's chain terms are not among them.
\end{proof}

\begin{remark}\label{4.21:lem-under-block}
The statement of this lemma is not written out here; parts of its proof are printed below.
\end{remark}

\begin{proof}[Proof of Lemma~\ref{4.21:lem-under-block}, concluded]
\label{4.21:prf-under-block-proof}
We prove clauses (d) and (e). Throughout, $\operatorname{dist}_t$ denotes distance measured in
units of level $t$; one unit of level $t+1$ is $L$ units of level $t$.

\emph{Clause} (d). Consider first a chain term $c$ of the present $\mathbf{R}_k$; its label is a
set $S\in\mathbb{D}_m$ of one level $m$, and we treat the case $m<k$. This label satisfies
\begin{equation*}
S\subset\Omega_{m+1}^{c},
\end{equation*}
where the domain is the unprimed one. Here $m$ is the minimal index of \cite[(1.63)]{B-CMP122a},
read in the label of the stage $i\le n(C')$ at which $c$ was born. At that stage the domain
$\Omega_{m+1}$, whose index satisfies $m+1\ge h+i+1$, has not yet been re-labelled. Hence, at
that stage, the level of the label on $S$ is at most $m$, and therefore so is the base.

Between that stage and the present one, no $\mathbf{T}$-step and no completed $\mathbf{R}$-operation
has occurred: neither a healing nor a throw-back. Indeed, we are inside $\mathbf{R}_k$, and the
label of a sibling is frozen. Therefore
\begin{equation*}
b\le m \quad\text{on } S
\end{equation*}
at the present stage.

For confined old terms and for local sibling terms, the assertion of (d) is the clause of
Lemma~\ref{4.21:lem-format-class} on confined old terms and local sibling terms. For a sibling it
is applied with $k_a=k$ and $X_{\mathfrak{a}}=C''$.

\emph{Clause} (e). Put $b:=b(y)$. Let $\ell$ be the level of $F$, and let $Q_F$ be its anchor,
given by condition (f4)$_\ell$ of the format in Lemma~\ref{4.21:lem-format-class}; for top
terms, $Q_F$ is the band cube, which lies inside $Z_\ell$. For $t\ge\ell+1$, and as long as $F$
lives, we show that the component $K_t$ in the first line below is defined, and that the
inequality in the second line holds:
\begin{equation}\label{4.21:eq-under-block-proof-a}
\begin{gathered}
K_t:=\text{the component of } Z^{\mathbf{T}}_t=(\Lambda^{\mathbf{T}}_t)^c
\text{ containing } Q_F,\quad\text{and}\\
\operatorname{dist}_t(Q_F,Z_t^{c})\ge\check{R}_t .
\end{gathered}
\end{equation}

\emph{The components $K_t$, by induction on $t$.} By (f4)$_\ell$ the anchor $Q_F$ lies within
$\check{R}_\ell+1$ level-$\ell$ units of $Z_\ell$. Since $\check{R}_\ell+2<8L\check{R}_{\ell+1}$,
the inclusions \eqref{4.21:eq-stages-components-a} give
\begin{equation*}
Q_F\subset Z'^{\sim 8}_\ell\subset Z^{\mathbf{T}}_{\ell+1}.
\end{equation*}
Hence $K_{\ell+1}$ is defined and $K_{\ell+1}\supset Q_F$. Moreover $X_F\supseteq A_F\supset Q_F$,
so $X_F$ meets $K_{\ell+1}$.

Suppose now that $K_t$ is defined. If $K_t$ is carried through a completed stage (L) of
$\mathbf{R}_t$, then $F$ ceases at $t$ by Lemma~\ref{4.21:lem-terms-cease}. That lemma applies,
because $F$ exists at step $t$: its birth step is $\ell<t$. Otherwise, after $\mathbf{R}_t$ we
have $K_t\subset Z_t$. Then, again by \eqref{4.21:eq-stages-components-a},
\begin{equation*}
Z^{\mathbf{T}}_{t+1}\supseteq Z'^{\sim 8}_t\supseteq K_t ,
\end{equation*}
and therefore $K_{t+1}\supseteq K_t$.

\emph{Proof of the inequality in \eqref{4.21:eq-under-block-proof-a}.} Let $t\ge\ell+1$ be such
that $F$ lives at $t$. The complement $Z_t^c$ is the union of $\Lambda^{\mathbf{T}}_t$ and of the
components healed at $t$. We bound the distance from $Q_F$ to each of these two parts.

Consider first $\Lambda^{\mathbf{T}}_t$. Since $\Lambda^{\mathbf{T}}_t\subset\Omega^{\mathbf{T}}_t$,
it suffices to bound the distance from $Q_F$ to $\Omega^{\mathbf{T}}_t$.
\begin{itemize}
\item Let $t=\ell+1$. By clause (d) of Lemma~\ref{4.21:lem-live-region-geometry},
\begin{equation*}
\operatorname{dist}_\ell(Z_\ell,\Omega^{\mathbf{T}}_{\ell+1})\ge 8L\check{R}_{\ell+1}.
\end{equation*}
The anchor $Q_F$ lies within $\check{R}_\ell+2\le L\check{R}_{\ell+1}+2$ level-$\ell$ units of
$Z_\ell$. Hence
$Q_F$ is at distance at least $7L\check{R}_{\ell+1}-2$ level-$\ell$ units from
$\Omega^{\mathbf{T}}_{\ell+1}$. This is at least $\check{R}_{\ell+1}$ level-$(\ell+1)$ units.
\item Let $t\ge\ell+2$. By the induction above, $Q_F\subset K_{t-1}\subset Z_{t-1}$. By clause
(d) of Lemma~\ref{4.21:lem-live-region-geometry},
\begin{equation*}
\operatorname{dist}_{t-1}(Z_{t-1},\Omega^{\mathbf{T}}_t)\ge 8L\check{R}_t .
\end{equation*}
This is at least $8\check{R}_t$ level-$t$ units.
\end{itemize}

Consider next a component healed at $t$. It is a component of $Z^{\mathbf{T}}_t$ other than
$K_t$: otherwise $F$ would have ceased at $t$, by the alternative just treated in the induction.
Its distance from $K_t\supset Q_F$ is therefore at least $\check{R}_t$, by (G0) in
\eqref{4.21:eq-strips-geometry-a}. This proves \eqref{4.21:eq-under-block-proof-a}.

\emph{Proof of} (L$\downarrow$) \emph{in} the corresponding display. Under the hypothesis of
(L$\downarrow$) we have $\ell\le b-1$.
Thus the cell of $y$ received its base at a step $b\ge\ell+1$, and $F$ was alive at that step.

Suppose first that the base was assigned by a completed $\mathbf{R}_b$ on a component
$C'\ni y$ of either class. This means that either $C'$ is healed, or $C'$ is thrown back as a
second-class region (Definition~\ref{4.21:def-strips-geometry}); the latter is the case in which
the base of a thrown-back component is defined. Then
$X_F\supseteq A_F\ni y$, so $X_F$ meets $C'$. For terms of type $(\delta)$ of the format class
(Lemma~\ref{4.21:lem-format-class}), for which $X_F=S^+(e)\subset A_F$, the fact that $X_F$ meets
$C'$ rests instead on the property, recorded with Lemma~\ref{4.21:lem-terms-cease}, that every own
cube of $A_F$ meets $X_F$. Hence $F$ ceased at $b$, by
Lemma~\ref{4.21:lem-terms-cease}, which covers both classes. This contradicts the fact that $F$
is alive, so this case does not occur.

The base was therefore assigned by $\mathbf{T}_b$. By (A0) in \eqref{4.21:eq-strips-geometry-a},
\begin{equation*}
y\in\Omega^{\mathbf{T}}_b\subset\Lambda_{b-1}=Z_{b-1}^{c}.
\end{equation*}
We distinguish two cases.
\begin{itemize}
\item Let $b-1\ge\ell+1$. Apply \eqref{4.21:eq-under-block-proof-a} at $t=b-1$. It gives
$\operatorname{dist}_{b-1}(Q_F,y)\ge\check{R}_{b-1}$, that is, $\check{R}_{b-1}L^{b-1-\ell}$
level-$\ell$ units. By [D-$\check{R}$] in \eqref{4.21:eq-strips-geometry-a} we have
$\check{R}_\ell\le L^{b-1-\ell}\check{R}_{b-1}$. By (c-$\check{R}$) there we have
$\check{R}_{b-1}\ge 38$. Hence
\begin{equation*}
\check{R}_{b-1}L^{b-1-\ell}\ge\max\{\check{R}_\ell,\;38L^{b-1-\ell}\}
\end{equation*}
level-$\ell$ units.
\item Let $b-1=\ell$. Then $y\in\Omega^{\mathbf{T}}_{\ell+1}$, and $Q_F$ lies within
$\check{R}_\ell+1$ level-$\ell$ units of $Z_\ell$. Clause (d) of
Lemma~\ref{4.21:lem-live-region-geometry} gives
\begin{equation*}
\operatorname{dist}_\ell(Q_F,y)\ge 8L\check{R}_{\ell+1}-\check{R}_\ell-1\ge 7\check{R}_\ell-1 .
\end{equation*}
\end{itemize}

By (g4) of Lemma~\ref{4.21:lem-tree-length-facts}, $d_\ell(A_F)$ is at least the distance from
$Q_F$ to the $\pi_\ell$-cube containing $y$. That distance is at least
$\operatorname{dist}_\ell(Q_F,y)-2$. In both cases, since $\check{R}_\ell\ge\check{R}_k$, this
gives
\begin{equation*}
d_\ell(A_F)\ge\check{R}_k-3 .
\end{equation*}
When $\ell\le b-2$, the first case gives in addition
\begin{equation*}
d_\ell(A_F)\ge 38L^{b-1-\ell}-3 .
\end{equation*}

\emph{Proof of} (L$\uparrow$) \emph{in} the corresponding display. Under the hypothesis of
(L$\uparrow$) we have $\ell\ge b+2$.

At the end of step $\ell-1\ge b+1$ the point $y$ lay in $\Lambda_{\ell-1}^{c}=Z_{\ell-1}$. Indeed,
$\Lambda_{\ell-1}$ is the union of $\Lambda^{\mathbf{T}}_{\ell-1}$ and of the components healed
at $\ell-1$. If $y$ had been in $\Lambda_{\ell-1}$, its base would be at least $\ell-1>b$. Since
$Z^{\mathbf{T}}_\ell\supseteq Z_{\ell-1}$ by \eqref{4.21:eq-stages-components-a}, the point $y$
lies in some component $K$ of $Z^{\mathbf{T}}_\ell$.

By condition (f0)$_\ell$ of the format in Lemma~\ref{4.21:lem-format-class}, one of the following
three alternatives holds.
\begin{itemize}
\item The set $A_F$ has a point $z\in\Omega^{\mathbf{T}}_\ell$. Then, by clause (d) of
Lemma~\ref{4.21:lem-live-region-geometry},
\begin{equation*}
\operatorname{dist}_\ell(y,z)\ge\operatorname{dist}_\ell(Z_{\ell-1},\Omega^{\mathbf{T}}_\ell)
\ge 8\check{R}_\ell .
\end{equation*}
\item The set $A_F$ has a point $z$ in a component $C'$ of $Z^{\mathbf{T}}_\ell$ healed at
$\mathbf{R}_\ell$. Then $y\notin C'$. A first-class component is healed as a whole
\cite{B-CMP122b}, so $y\in C'$ would give $b(y)\ge\ell$. Hence $K\neq C'$, and (G0) in
\eqref{4.21:eq-strips-geometry-a} gives $\operatorname{dist}_\ell(y,z)\ge\check{R}_\ell$.
\item The set $A_F$ crosses $\Lambda^{\mathbf{T}}_\ell$ between two components. Then
$d_\ell(A_F)\ge\check{R}_\ell-2$, by (G0) in \eqref{4.21:eq-strips-geometry-a} and (g4) of
Lemma~\ref{4.21:lem-tree-length-facts}.
\end{itemize}

In each alternative, (g4) of Lemma~\ref{4.21:lem-tree-length-facts} gives
\begin{equation*}
d_\ell(A_F)\ge\check{R}_\ell-3\ge\check{R}_k-3 .
\end{equation*}

The proof of (L$\uparrow$) does not use that $y$ is a window point. This hypothesis is kept in
the statement because (L$\uparrow$) is applied at window points only.
\end{proof}

\begin{remark}[Stage counts, bases and excluded components under a window]
\label{4.21:rem-under-block-remarks}
The following four readings accompany Lemma~\ref{4.21:lem-under-block} (what lies under a block)
and Lemma~\ref{4.21:lem-window-probe} (what a window probe meets).
\begin{enumerate}
\item[(1)] Clauses (a) and (b) of Lemma~\ref{4.21:lem-under-block} restate the geometric clauses
of the grain-partition lemma, which governs the grain partition $\pi^{\flat}$, applied from the
level of the block down to the finer cells beneath it. The common domain of
\cite[p.~191]{B-CMP122a} is the case $i\ge k_0+1$ of clause (a). Two counts give the same
interval. Ba{\l}aban's construction counts the stages whose collars overlap, of which there are
at most $N_0$; here $N_0$ is the bound on the number of stages of $\mathbf{R}_k$ whose collars
overlap. The overlap-counting lemma instead enumerates the $(2L^{N_0})^d$
finer cells lying under one block. Both counts describe the interval of clause (b).
\item[(2)] Clause (e) of Lemma~\ref{4.21:lem-under-block} compares the level $\ell$ with the base of
the point. Clause (b) bounds the difference between zone and base under a window block. The
factor supplied by the overlap-counting lemma is exactly what separates this reading from the
reading in terms of the zone used in the preparation-site bound.
\item[(3)] The following holds:
\begin{equation}\label{4.21:eq-under-block-remarks-a}
\begin{gathered}
\text{when the set of old blocks is the enlarged set } \mathfrak{O}^{+}:\\
\text{for every window block } \square_0 \text{, the probe } \tilde{\square}_0\\
\text{meets no confined old term } c\ (k_a<k),\\
\text{whether } c \text{ is frozen or of class } (\delta).
\end{gathered}
\end{equation}
The assertion concerns confined old terms only. Top old terms of class $(\delta)$ do reach
window points. They have the type of the terms $\mathbf{B}^{(\ell)}$,
$\mathbf{B}'^{(\ell)}$ of class $(\alpha)$ in
Lemma~\ref{4.21:lem-format-class} (the format class and its confined and top old terms).
For this reason the fine-term bound for a window block and clause (a) of the window-counting
lemma below contain no contribution from frozen fine terms.
\item[(4)] The case of a sibling $C''$, that is, a component with $k_a=k$, differs from that
of~\eqref{4.21:eq-under-block-remarks-a}. Its window region $C''\setminus Z_h$ consists of
window points. There the window probes do meet two kinds of terms. The first is its own
chain, the local chain of class $(\gamma)$ in Lemma~\ref{4.21:lem-format-class}. The second
arises if $C''$ is a sibling produced in the fourth branch of the classification of arrested
attempts, with arrested attempt $\mathfrak{a}$, and if moreover the stage datum $n^{+}$ of
$\mathfrak{a}$ satisfies $n^{+}\le n$, the present stage of $\mathbf{R}_k$ being the stage
$n+1$: then its own confined terms of $\mathfrak{Q}(\mathfrak{a})$, the local sibling terms,
are met as well. These terms are counted by the second and the fourth clause of the
probe-counting lemma below, in the same way
as the terms of this chain and as the terms $\mathbf{B}^{(\ell)},\mathbf{B}'^{(\ell)}$ of each
level $\ell$. Moreover, the terms filed at $k$ of every chain of this $\mathbf{R}_k$ may reach
any window point. This is the origin of the constant $c_{\mathrm{sib}}$ in the second clause of
that lemma.
\end{enumerate}
\end{remark}

\begin{proof}
We prove the identification in item (1) and the assertion
\eqref{4.21:eq-under-block-remarks-a} of item (3). Items (2) and (4) compare clauses already
established; the counting referred to in item (4) is carried out in the probe-counting lemma
below.

For item (1), let $\Delta\subset Z\cap C'$ be a cell of current level $i$ with
$h+1\le i\le j(C')$ and $i\ge k_0+1$. Clause (a) of Lemma~\ref{4.21:lem-under-block} then gives
$b\ge\min\{i-1,k_0\}=k_0$ on $\Delta$. This is the case in which the stages involved share
the common domain of \cite[p.~191]{B-CMP122a}.

For item (3), let $c$ be a confined old term, let $\mathfrak{a}$ be the arrested attempt to
which it belongs, with step $k_a<k$ and region $X_{\mathfrak{a}}$, and let $c$ be frozen or of
class $(\delta)$. The label of $c$ lies in one of two places.
\begin{itemize}
\item The first possibility is the region $X_{\mathfrak{a}}$ of $\mathfrak{a}$. There
\begin{equation*}
X_{\mathfrak{a}}\subset Z_{k_a}\subset Z_h\cap C',
\end{equation*}
where $C'$ is a component of a class.
\item The second possibility is an untreated component.
\end{itemize}
The inclusions in the first case hold by the clause of the classification of arrested attempts
that locates the region of an arrested attempt. Now let
$\square_0$ be a window block. By Lemma~\ref{4.21:lem-window-probe}
(display~\eqref{4.21:eq-window-probe-a}), the probe $\tilde{\square}_0$ meets no cell of
$Z_h\cap C'$ for any class component $C'$ and no cell of an untreated component. In either of
the two cases above, therefore, the label of $c$ lies in
cells that $\tilde{\square}_0$ does not meet. This proves
\eqref{4.21:eq-under-block-remarks-a}.
\end{proof}

\begin{definition}[The grain of a stage and its constants]\label{4.21:not-stage-grain}
Fix a stage $n+1$ of the construction at step $k$. Here $Z$, $\Omega_{j+1}$ and the level $j$
are the domains and the level of the current sequence of domains of stage $n+1$, as fixed in
the construction of the stages. For a level $\ell$, $\pi_\ell$ denotes the partition
into cubes of level $\ell$.
\begin{enumerate}
\item The \emph{grain} of stage $n+1$ is the partition $\pi^{\flat}$ whose members are of
three kinds:
\begin{itemize}
\item the cubes of $\pi_k$ off $Z$;
\item the cubes of $\pi_{j+1}$ in $Z\cap\Omega^{c}_{j+1}$;
\item the cells of the current cell partition lying in $Z\cap\Omega_{j+1}$.
\end{itemize}
This partition is fixed in the construction of the stages.
\item For a block $\square$ of the cover (the cover over whose blocks the sum in part~(c) of
Lemma~\ref{4.21:lem-cell-block-sums} runs), with $\tilde{\square}^{4}$ the enlargement of
$\square$ fixed with that cover, the
\emph{grain neighbourhood} of $\square$ is
\begin{equation*}
\hat{\mathsf{K}}_{\square}
:=\bigcup\bigl\{\,G\in\pi^{\flat}:\ G\cap\tilde{\square}^{4}\neq\emptyset\,\bigr\}.
\end{equation*}
\item We put
\begin{equation*}
N^{\flat}:=(10L+2)^{d}.
\end{equation*}
\item The counting constants used with the grain are the constants appearing in
Lemma~\ref{4.21:lem-cell-block-sums} (sums over cells and blocks against the scaled
distance). Their values are
\begin{equation*}
\hat{b}_t=(L+2)^{d}+3^{d},\qquad
c_{\boxplus}=3^{d}(L+5),\qquad
C_{\mathrm{cell}}=4\,\hat{b}_t^{\,2c_{\boxplus}},
\end{equation*}
and
\begin{equation*}
C_{\mathrm{blk}}=2^{d+1}+(L+1)^{d}.
\end{equation*}
\item The unit of length along strata is $w=M/M_1$, with $M$ and $M_1$ among
Ba{\l}aban's auxiliary parameters.
\item The symbol $d^{\wedge}$ denotes the scaled distance of
Lemma~\ref{4.21:lem-cell-block-sums}, computed for the current label, that is, for the label
in force at the stage considered. The collar $R_n$ of stage $n+1$, as fixed in the
construction of the stages, is made of cells of levels $j$ and $j+1$.
\item The constant $\delta_P$ is the decay rate whose sixteenth part, $\beta=\delta_P/16$, is
the rate at which Lemma~\ref{4.21:lem-cell-block-sums} is applied. Throughout, the parameters
are assumed to satisfy
the condition \eqref{4.21:eq-cell-block-sums-a}, namely
\begin{equation*}
\tfrac{1}{16}\,\delta_P\,w\ \ge\ c_{\boxplus}\log\hat{b}_t+1 .
\end{equation*}
This is the hypothesis of part~(b) of Lemma~\ref{4.21:lem-cell-block-sums},
$\beta w\ge c_{\boxplus}\log\hat{b}_t+1$, with $\beta=\delta_P/16$. It is therefore also
the hypothesis of part~(c) of that lemma.
\end{enumerate}
\end{definition}

\begin{lemma}[The blocks around one grain member sum to a constant]\label{4.21:lem-grain-block-sum}
Let $\pi^{\flat}$, $\hat{\mathsf{K}}_{\square}$, $N^{\flat}$, $\hat{b}_t$, $c_{\boxplus}$, $C_{\mathrm{blk}}$, $w=M/M_1$,
$d^{\wedge}$ and $R_n$ be as in Notation~\ref{4.21:not-stage-grain}. For a block $\square_0$ of the cover we call
the cubes $G\in\pi^{\flat}$ that meet $\tilde{\square}_0^{4}$ the members of $\hat{\mathsf{K}}_{\square_0}$.
Assume the condition on $w$ fixed in Notation~\ref{4.21:not-stage-grain},
\begin{equation}\label{4.21:eq-grain-block-sum-1}
\tfrac{1}{16}\,\delta_P\, w\ \ge\ c_{\boxplus}\log\hat{b}_t+1 .
\end{equation}
Then for every $\beta\ge\delta_P/16$ and every $G\in\pi^{\flat}$,
\begin{equation}\label{4.21:eq-grain-block-sum-2}
\sum_{\square_0} e^{-\beta d^{\wedge}(\square_0,R_n)}\ \le\ c_{\mathrm{loc}},
\end{equation}
where the sum runs over the blocks $\square_0$ of the cover, in any zone, in $Z$ or off it, such that
$G\in\hat{\mathsf{K}}_{\square_0}$ or some member of $\hat{\mathsf{K}}_{\square_0}$ touches $G$, and
\begin{equation*}
c_{\mathrm{loc}}:=\hat{b}_t^{\,N^{\flat}+3}\,C_{\mathrm{blk}}
\Bigl[L^d+2\hat{b}_t^{\,2c_{\boxplus}}\sum_{x\ge0}(L^2+2x+6)^d e^{-x}\Bigr].
\end{equation*}
The number $c_{\mathrm{loc}}$ depends on $d$ and $L$ only and on nothing else; in particular it depends on no
region size parameter $\check{R}$, on no age, and not on the number of windows, of components or of attempts.
\end{lemma}

\begin{proof}
We use the following properties of the grain of stage $n+1$, which are clauses of the lemma on the grain of a
stage (the lemma establishing the basic properties of $\pi^{\flat}$ and of $\hat{\mathsf{K}}_{\square}$): for every
block $\square$ of the cover, every cell meeting $\square$ lies in
exactly one member of $\hat{\mathsf{K}}_{\square}$; a member of $\pi^{\flat}$ touches at most $\hat{b}_t$ members of
$\pi^{\flat}$; and $\hat{\mathsf{K}}_{\square}$ is wall-connected, has at most $N^{\flat}$ members and contains
$\tilde{\square}^{4}\supset\square$.

(1) Let $\square_0$ be a block in the sum \eqref{4.21:eq-grain-block-sum-2}. By the last of the properties just
recalled, $\hat{\mathsf{K}}_{\square_0}$ is wall-connected, has at most $N^{\flat}$ members, and contains
$\tilde{\square}_0^{4}\supset\square_0$; by the choice of $\square_0$ it contains $G$ or a member touching $G$.
Every cell $\Delta$ meeting $\square_0$ lies in one member $G''\in\hat{\mathsf{K}}_{\square_0}$. Since
$\hat{\mathsf{K}}_{\square_0}$ is wall-connected with at most $N^{\flat}$ members, $G''$ is joined to the member
that equals or touches $G$ by a chain of at most $N^{\flat}$ touching members, hence to $G$ by a chain of at most
$N^{\flat}+1$ touching members. A member touches at most $\hat{b}_t$ members, so the number of chains of length $r$
starting at $G$ is at most $\hat{b}_t^{\,r}$. Since $\hat{b}_t=(L+2)^d+3^d\ge2$, the members $G''$ that can occur in
this way, over all blocks $\square_0$ of the sum, number at most
\begin{equation*}
\sum_{r=0}^{N^{\flat}+1}\hat{b}_t^{\,r}\ \le\ \hat{b}_t^{\,N^{\flat}+2}\ \le\ \hat{b}_t^{\,N^{\flat}+3}.
\end{equation*}

(2) The scaled distance $d^{\wedge}(\square_0,R_n)$ is attained at a cell meeting $\square_0$, as at the beginning of
the proof of part (c) of Lemma~\ref{4.21:lem-cell-block-sums}. Hence
\begin{equation*}
e^{-\beta d^{\wedge}(\square_0,R_n)}\ \le\ \sum_{\Delta\ \text{meets}\ \square_0} e^{-\beta d^{\wedge}(\Delta,R_n)} .
\end{equation*}
We sum this over the blocks $\square_0$ of \eqref{4.21:eq-grain-block-sum-2} and exchange the order of summation.
Each cell $\Delta$ that occurs lies in one of the members $G''$ of step (1), and it occurs for at most
$C_{\mathrm{blk}}$ blocks, because a cell meets at most $C_{\mathrm{blk}}$ blocks of the cover
(Lemma~\ref{4.21:lem-live-region-geometry}, the third assertion of clause (i)). Therefore the left side of
\eqref{4.21:eq-grain-block-sum-2} is at most
\begin{equation}\label{4.21:eq-grain-block-sum-3}
C_{\mathrm{blk}}\sum_{G''}\ \sum_{\Delta\subset G''} e^{-\beta d^{\wedge}(\Delta,R_n)},
\end{equation}
with $G''$ running over the at most $\hat{b}_t^{\,N^{\flat}+3}$ members of step (1). It remains to show that the inner
sum in \eqref{4.21:eq-grain-block-sum-3} is at most the bracket in the definition of $c_{\mathrm{loc}}$.

(3) Fix one member $G''$ and put $\ell:=\operatorname{lev}(G'')$. If $G''$ is a cell (a member of the grain that is a
cell in $Z\cap\Omega_{j+1}$, or a $\pi_{j+1}$-cube in $Z\cap\Omega^c_{j+1}$ that is a band cell in the sense of
Notation~\ref{4.21:not-stage-grain}, and hence itself a cell), the inner sum has one summand, and it is at most $1$
since $d^{\wedge}\ge0$.

Otherwise $G''$ is a $\pi_k$-cube off $Z$ or a $\pi_{j+1}$-cube in $Z\cap\Omega^c_{j+1}$. Its cells of level
$\ell-1$ number at most $L^d$, because a cell of level $\ell$ contained in $G''$ would be $G''$ itself; each of them
contributes at most $1$. Let $\mathsf{W}$ be the union of all cells of level $\le\ell-2$ on the side of $\partial Z$
on which $G''$ lies: in $Z$ when $G''$ is a $\pi_{j+1}$-cube in $Z\cap\Omega^c_{j+1}$, where $\ell-2=j-1$; off $Z$
when $G''$ is a $\pi_k$-cube off $Z$, where $\ell-2=k-2$. Let $\Sigma$ be the closure of the complement of
$\mathsf{W}$. Then $\Sigma$ is a union of cells and contains $R_n$: when $G''$ lies in $Z$, because the cells of
$R_n$ have levels $j$ and $j+1$, both $\ge\ell-1$; when $G''$ lies off $Z$, because $R_n\subset Z$ and $\mathsf{W}$
lies off $Z$. Moreover no cell of $\mathsf{W}$
touches the other side of $\partial Z$, since the cells on both sides of $\partial Z$ adjacent to it have level $k$,
as is shown in the proof of the bound on touching members in the lemma on the grain of a stage. Every contour from
a cell $\Delta\subset\mathsf{W}$ to $R_n$ meets $\Sigma$, so
\begin{equation*}
d^{\wedge}(\Delta,R_n)\ \ge\ d^{\wedge}(\Delta,\Sigma)\qquad(\Delta\subset\mathsf{W}).
\end{equation*}

Let $x\ge0$ and consider a contour $\Gamma$ from $\Sigma$ to a cell $\Delta\subset\mathsf{W}\cap G''$ with
$\operatorname{len}_{\mathfrak{B}}(\Gamma)<xw$. After its last point in $\Sigma$ it runs through cells of level
$\le\ell-2$,
so by the comparison of scaled and Euclidean lengths in \eqref{4.21:eq-cell-block-sums-a} its Euclidean length is
less than $x$ units of level $\ell-2$. It starts in a cell of $\partial^+\Sigma$; this cell touches a cell of level
$\ge\ell-1$, and therefore has level $\ell-2$ by the admissibility condition on the levels of touching cells. It lies
within $x+1$ units of level $\ell-2$ of $G''$, whose side is $L^2$ such units; hence there are at most
$(L^2+2x+4)^d$ possible start cells. By part (a) of Lemma~\ref{4.21:lem-cell-block-sums}, the contour meets at most
$c_{\boxplus}(x+1)$ cells, chained by touching; there are at most $\hat{b}_t^{\,c_{\boxplus}(x+1)}$ such chains from
a given start cell. Consequently
\begin{equation}\label{4.21:eq-grain-block-sum-4}
\#\{\Delta\subset\mathsf{W}\cap G'':\ d^{\wedge}(\Delta,\Sigma)<xw\}
\ \le\ (L^2+2x+4)^d\,\hat{b}_t^{\,c_{\boxplus}(x+1)} .
\end{equation}
We group the cells $\Delta\subset\mathsf{W}\cap G''$ according to the integer $x\ge0$ with
$xw\le d^{\wedge}(\Delta,\Sigma)<(x+1)w$. By \eqref{4.21:eq-grain-block-sum-4} with $x+1$ in place of $x$, the group
of $x$ has at most $(L^2+2x+6)^d\hat{b}_t^{\,c_{\boxplus}(x+2)}$ members, each contributing at most $e^{-\beta wx}$.
Hence
\begin{align*}
\sum_{\Delta\subset\mathsf{W}\cap G''} e^{-\beta d^{\wedge}(\Delta,R_n)}
&\le \sum_{x\ge0}(L^2+2x+6)^d\,\hat{b}_t^{\,c_{\boxplus}(x+2)}e^{-\beta wx}\\
&= \hat{b}_t^{\,2c_{\boxplus}}\sum_{x\ge0}(L^2+2x+6)^d\bigl(\hat{b}_t^{\,c_{\boxplus}}e^{-\beta w}\bigr)^{x}\\
&\le 2\hat{b}_t^{\,2c_{\boxplus}}\sum_{x\ge0}(L^2+2x+6)^d e^{-x}.
\end{align*}
In the last step we used $\hat{b}_t^{\,c_{\boxplus}}e^{-\beta w}\le e^{-1}$, which is the inequality
$\beta w\ge c_{\boxplus}\log\hat{b}_t+1$; for $\beta\ge\delta_P/16$ it follows from \eqref{4.21:eq-grain-block-sum-1},
since then $\beta w\ge\frac{1}{16}\delta_P w$. (The factor $2$ is not needed for this inequality; it is kept so that
the bound matches the definition of $c_{\mathrm{loc}}$.)

Adding the at most $L^d$ cells of level $\ell-1$, the inner sum in \eqref{4.21:eq-grain-block-sum-3} is, for each
member $G''$, at most
\begin{equation*}
L^d+2\hat{b}_t^{\,2c_{\boxplus}}\sum_{x\ge0}(L^2+2x+6)^d e^{-x};
\end{equation*}
the term $L^d$ also covers the case in which $G''$ is a cell. Inserting this into \eqref{4.21:eq-grain-block-sum-3},
with at most $\hat{b}_t^{\,N^{\flat}+3}$ members $G''$, gives \eqref{4.21:eq-grain-block-sum-2}. All quantities
entering $c_{\mathrm{loc}}$, namely $N^{\flat}$, $\hat{b}_t$, $c_{\boxplus}$ and $C_{\mathrm{blk}}$, are functions of $d$
and $L$ only.
\end{proof}

\begin{remark}
(1) Thin layers of cells of intermediate levels cost nothing extra: a cell many levels below $G''$ is reached only by
crossing the layers of those levels, each of which contributes $w$ to the scaled distance.

(2) The lemma is used in the following form. Let $Q$ be a cube contained in one member $G$ of $\pi^{\flat}$, and let
$\square_0$ be a block whose probe $\tilde{\square}_0$ meets $Q$. A point of $\tilde{\square}_0$ is interior to
$\tilde{\square}_0^{4}$, so $\tilde{\square}_0^{4}$ meets $G$ and $G\in\hat{\mathsf{K}}_{\square_0}$. The blocks in
question are thus among those of \eqref{4.21:eq-grain-block-sum-2}, and
\begin{equation}\label{4.21:eq-grain-block-sum-a}
\sum_{\square_0:\ \tilde{\square}_0\ \text{meets}\ Q} e^{-\beta d^{\wedge}(\square_0,R_n)}\ \le\ c_{\mathrm{loc}} .
\end{equation}
\end{remark}

\begin{definition}[Weights, budgets and incidences]\label{4.21:not-weights-incidences}
Throughout, $k$ is the level of the present step $\mathbf{R}_k$; $N$ and the separation radius
$\check{R}_k$ of that step are as declared in
\eqref{4.21:eq-strips-geometry-a}; $\pi^{\flat}$ is the grain of Notation~\ref{4.21:not-stage-grain};
$\mathfrak{F}$ is the format class of Lemma~\ref{4.21:lem-format-class}; pockets $(Y,m)$ and their
strips $Y^{\sharp}$ are those of Definition~\ref{4.21:def-strips-geometry}. In
\eqref{4.21:eq-strips-geometry-a} the letter $N$ is bounded by $N\le C_N\,2^{r_{\mathrm{pr}}}\check{R}_k$
and is also declared equal to $\check{R}_k$; item (a) below covers both readings.

\emph{Data of a term.} A term $c\in\mathfrak{F}$ has its label $A_c$. For chain terms and for the
terms of clause $(\delta)$ of Lemma~\ref{4.21:lem-format-class} we put $S_c:=A_c$; for a leaf we put
$S_c:=[A_c]_{m_c}$, where $m_c$ is the filing level of the leaf in the sense of
Lemma~\ref{4.21:lem-format-class}. We write $\mathcal{P}_c$ for the list of pockets of $c$
($\mathcal{P}_c=\emptyset$ for plain terms), and define the \emph{region} $X_c$ of $c$ and the set of
\emph{incident blocks} by
\begin{equation}\label{4.21:eq-weights-incidences-1}
X_c:=A_c\cup\bigcup_{(Y,m)\in\mathcal{P}_c}Y^{\sharp},\qquad
\mathcal{A}_c:=\{\square:\ \tilde{\square}\cap X_c\neq\emptyset\},
\end{equation}
except that for a term $c$ of clause $(\delta)$ we put $X_c:=S^{+}(c)\subset A_c$. The following are
not members of $\mathfrak{F}$: the small-field leaves (of every level); the near groups and the pieces
of the block and rim totals of the present step; and the kept members of live pockets.

\emph{Confinement.} A term $c\in\mathfrak{F}$ is \emph{confined} if $X_c$ lies inside a single class
component of the present step, and \emph{not confined} otherwise; in
\eqref{4.21:eq-weights-incidences-3} the word confined always has this meaning. This notion is
distinct from the
notion of a confined old term in Lemma~\ref{4.21:lem-format-class} (a term of an arrest at step
$k_a$ stored at a level below $k_a$, or at level $k_a$ outside $\Omega^{(k_a)}_{k_a}$); in particular a
confined old term of an arrest lying in an untreated component need not be confined in the present
sense, and then $\mathsf{e}(c)=d_k([S_c]_k)/N$, which is non-zero in general.

\emph{Weights and budgets.} With $\mathfrak{N}^{A}$ the size functional of stored terms (read as
$\mathfrak{N}^{G}$ for the terms of clause $(\delta)$), $\kappa_A$ the weight of the strong set
$\mathcal{S}(c)$, and $w(c)$ the hull weight of $c$ at its filing level, in the sense of the weighted
bounds of stored terms, we put
\begin{equation}\label{4.21:eq-weights-incidences-2}
\mathfrak{w}(c):=\mathfrak{N}^{A}(c)\,e^{\kappa_A|\mathcal{S}(c)|}\,e^{w(c)},\qquad
W(c):=\mathfrak{w}(c)\,e^{\frac23\mathsf{e}(c)+\mathsf{f}(c)},
\end{equation}
where
\begin{equation}\label{4.21:eq-weights-incidences-3}
\begin{aligned}
\mathsf{e}(c)&:=
\begin{cases}
d_k([S_c]_k)/N, & c \text{ not confined},\\
0, & c \text{ confined},
\end{cases}\\
\mathsf{f}(c)&:=
\begin{cases}
d_k([S_c]_k), & c \text{ of leaf type, or old and confined, or a local sibling},\\
0, & c \text{ a chain term of } \mathbf{R}_k .
\end{cases}
\end{aligned}
\end{equation}
The factor $e^{\mathsf{e}(c)}$ is the budget that a non-confined chain term draws from its running
rate: if $n+1$ is the present stage and $a^{(i)}_k$ denotes the running rate of the chain terms of
stage $i$, then by the rate schedule of the chain terms, for $i\le n$,
$a^{(i)}_k-a^{(n+1)}_k=(n+1-i)/N\ge 1/N$ per unit of $d_k$. Two thirds of this budget are placed
in $W(c)$; the remaining third is reserved for the summation over the chains of $\mathbf{R}_k$, where
the constant $c_{\mathrm{sib}}$ below enters.

\emph{Grain and incidences.} An own cube $Q$ of $c$ (its cubes are of one level, by
Lemma~\ref{4.21:lem-format-class}) is \emph{grain-fine} if $Q$ lies inside one member $G(Q)$ of
$\pi^{\flat}$, and \emph{grain-coarse} otherwise. An \emph{incidence} is a pair $(\square_0,c)$ with
$\square_0\in\mathcal{A}_c$. With $\operatorname{zone}(\square_0)$ the zone of the block in the sense
of Definition~\ref{4.21:def-strips-geometry}, the incidence is
\begin{itemize}
\item a \emph{strip incidence} if $\tilde{\square}_0$ meets $X_c$ only through strips; such a block
$\square_0$ belongs to $\mathfrak{O}^{+}$;
\item a \emph{coarse incidence} if $\tilde{\square}_0$ meets an own cube $Q$ of $c$ with
$\operatorname{lev}(Q)\ge\operatorname{zone}(\square_0)-1$;
\item a \emph{fine incidence} otherwise, that is, when $\tilde{\square}_0$ meets $A_c$ only through own
cubes of level at most $\operatorname{zone}(\square_0)-2$.
\end{itemize}

\emph{Scaffold functions and constants.} Let $\lambda$ and $c_2$ be the constants of the coupling
comparison in \eqref{4.21:eq-strips-geometry-a}, for which $\lambda\ge0$ and $c_2\ge0$, and recall the
stage-count constant $C_N=2^{r_{\mathrm{pr}}}L$ of \eqref{4.21:eq-strips-geometry-a}. For $s\ge0$ put
$c_s:=\log(1+3\lambda s+2c_2)$, and for $\mathsf{T}\ge0$
\begin{equation}\label{4.21:eq-weights-incidences-4}
\mathsf{R}^{\uparrow}(\mathsf{T}):=L\Bigl(\mathsf{T}+1+\log\bigl(1+3\lambda C_N(\mathsf{T}+1)^{r_{\mathrm{pr}}}
+2c_2\bigr)\Bigr)^{r_{\mathrm{pr}}} .
\end{equation}
The constants of the counts that follow, and the two level ceilings imposed on $\gamma$, are
\begin{equation}\label{4.21:eq-weights-incidences-a}
\begin{aligned}
&\theta_0:=(3\,C_N\,2^{r_{\mathrm{pr}}})^{-1},\qquad
c_{\mathrm{sib}}:=e\sum_{u\ge0}(2u+4)^d e^{-u\theta_0},\qquad c_{\mathrm{att}}:=3\cdot4^d,\\
&N^{\mathrm{fr}}:=C_N(\check{\mathsf{T}}_k+2+c_{N+1})^{r_{\mathrm{pr}}},\qquad
\mathfrak{g}(s):=\bigl(1+\log(1+3\lambda(s+1))\bigr)^{r_{\mathrm{pr}}}\quad(s\ge0),\\
&C_{\mathrm{fr}}:=\sum_{s\ge1}L^{d(s+1)}\mathfrak{g}(s)\,e^{3-38L^{s}},\qquad
c_{\mathrm{fr}}:=c_{\mathrm{att}}(3+C_{\mathrm{fr}}),\\
&\mathfrak{c}:=2K(d)A_{\flat}B_0+c_{\mathrm{att}}A_{\flat}C_0^{\sharp}N^{\mathrm{fr}}+A_{\flat}b_{\delta},
\qquad
\mathfrak{c}^{\circ}:=2K(d)A_{\flat}B_0+A_{\flat}b_{\delta},\\
&\mathfrak{t}^{\mathrm{win}}:=c_{\mathrm{sib}}NC_0^{\sharp}+(N+1)A_{\flat}b_{\delta}
+(3+C_{\mathrm{fr}})\mathfrak{c}\\
&\qquad\;=c_{\mathrm{sib}}NC_0^{\sharp}+2(3+C_{\mathrm{fr}})K(d)A_{\flat}B_0
+c_{\mathrm{fr}}A_{\flat}C_0^{\sharp}N^{\mathrm{fr}}+(N+4+C_{\mathrm{fr}})A_{\flat}b_{\delta},\\
&\mathfrak{l}_j:=3(5L^{4})^d\ \ (j\le k_0+1),\qquad
\mathfrak{l}_j:=3(5L^{N_0+2})^d\ \ (j\ge k_0+2),\qquad
\check{R}^{\uparrow}:=\check{R}_{h+1},\\
&\sup_{\mathsf{T}\ge\mathsf{T}_{\gamma}}\bigl(C_N(\mathsf{T}+1)^{r_{\mathrm{pr}}}+4\bigr)
\mathfrak{g}(0)\,e^{4-\mathsf{T}^{r_{\mathrm{pr}}}}\le1,\\
&\sup_{\mathsf{T}\ge\mathsf{T}_{\gamma}}L^{d}\mathfrak{g}(0)\,
e^{3/2-\mathsf{T}^{r_{\mathrm{pr}}}/2}\le\tfrac14 .
\end{aligned}
\end{equation}
Here $h$, $k_0$ and $N_0$ are as in \eqref{4.21:eq-strips-geometry-a}.
The constant $\mathfrak{c}$ is the per-cube constant of the leaf-type class; $\mathfrak{c}^{\circ}$ is
the same constant without the frozen chains, and does not depend on $\check{\mathsf{T}}_k$. The last
two lines of \eqref{4.21:eq-weights-incidences-a} are the \emph{level ceilings}. The following hold.
\begin{enumerate}
\item[(a)] $\theta_0\le\check{R}_k/(3N)$ and $\theta_0\le\frac13$; $c_{\mathrm{sib}}$ is finite and
depends on $(d,L,r_{\mathrm{pr}})$ only.
\item[(b)] For $s\ge0$,
\begin{equation}\label{4.21:eq-weights-incidences-5}
c_{N+1+s}\le c_N+\log\bigl(1+3\lambda(s+1)\bigr).
\end{equation}
\item[(c)] $\mathfrak{g}$ is non-decreasing, and
\begin{equation}\label{4.21:eq-weights-incidences-8}
\mathfrak{g}(a+s+1)\le\mathfrak{g}(a)\,\mathfrak{g}(s)\qquad(a,s\ge0).
\end{equation}
\item[(d)] The stage count $N(g_t)$ of any attempt at a step $t$ with $k-t\le N+1$ satisfies
$N(g_t)\le N^{\mathrm{fr}}$. More generally, for
$s\ge0$ and every step $t$ with $k-t\le N+1+s$,
\begin{equation}\label{4.21:eq-weights-incidences-6}
N(g_t)\le C_N(\check{\mathsf{T}}_k+1+c_{N+1+s})^{r_{\mathrm{pr}}}\le\mathfrak{g}(s)N^{\mathrm{fr}},
\qquad
\check{R}_t\le R(x_t)\le\mathfrak{g}(s)\,\mathsf{R}^{\uparrow}(\check{\mathsf{T}}_k).
\end{equation}
\item[(e)] $C_{\mathrm{fr}}$ is finite for every $(d,L,\lambda,r_{\mathrm{pr}})$; it is a number fixed
after $(d,L,\lambda,r_{\mathrm{pr}})$ and before $\gamma$.
\item[(f)] The two expressions for $\mathfrak{t}^{\mathrm{win}}$ in
\eqref{4.21:eq-weights-incidences-a} agree.
\item[(g)] $\check{R}^{\uparrow}\le\mathsf{R}^{\uparrow}(\check{\mathsf{T}}_k)$.
\item[(h)] Under the level ceilings,
\begin{equation}\label{4.21:eq-weights-incidences-7}
\begin{aligned}
&(N+4)\mathfrak{g}(0)e^{4-\check{R}_k}\le1,\qquad N\mathfrak{g}(0)e^{3-\check{R}_k}\le e^{-1},\\
&\mathfrak{g}(0)e^{3-\check{R}_k}\le\tfrac14e^{-1},\qquad L^{d}\mathfrak{g}(0)e^{3-\check{R}_k}\le\tfrac14 .
\end{aligned}
\end{equation}
\item[(i)] For every $c\in\mathfrak{F}$,
\begin{equation*}
e^{\mathsf{e}(c)+\mathsf{f}(c)+d_k([S_c]_k)}\le e^{3d_k([S_c]_k)};
\end{equation*}
and for a label $A$ at level $\ell$,
the hull inequality $d_k([A]_k)\le d_{\ell}(A)$ gives $e^{\mathsf{e}+\mathsf{f}}\le e^{2d_{\ell}(A)}$.
\end{enumerate}

\emph{On the constants.} $C_{\mathrm{fr}}$ is not evaluated: no inequality in this section bounds
$C_{\mathrm{fr}}$, or any $\mathsf{T}$-free quantity that grows with $L$, $\lambda$ or $r_{\mathrm{pr}}$,
from above by a numerical constant, except floors on parameters chosen later (the floor on
$\kappa_1$, chosen after $L$; the floor on $M$; the floor on the stratum unit $M/M_1$ in
Notation~\ref{4.21:not-stage-grain}) and ceilings on $\gamma$. The only places where $L^d$ or
$\mathfrak{g}(0)$ stand against a numerical constant are ceilings on $\gamma$, namely the level
ceilings and one further ceiling of the same form imposed later in this section; there they multiply a
function of $\mathsf{T}$ that tends to zero as $\mathsf{T}\to\infty$, so that they are conditions on
$\gamma$ imposed after $L$. The constant $38$ in $C_{\mathrm{fr}}$ is the one in
\eqref{4.21:eq-strips-geometry-a}. The factor $L^{d(s+1)}$ is the number of cubes of $\pi_{\ell-s}$
contained in a cube of $\pi_{\ell+1}$, one level above the level $\ell$ of the cutoff. It serves
three counts: the window positions below the floor $\mathfrak{f}-1$ (the windows and their floor
$\mathfrak{f}$ are introduced later in this section; here $L^{ds}\le L^{d(s+1)}$
suffices), one further count of positions later in this section, and the rim (where $L^{d(s+1)}$ is
exact, a rim cube of an element of level $j$ being a $\pi_{j+1}$-cube of base $j+1$).

In $\mathfrak{t}^{\mathrm{win}}$ the coefficient of $A_{\flat}b_{\delta}$ is accounted for as follows:
old confined terms of an arrest never reach a window; the top ones are leaves and are counted in
$(3+C_{\mathrm{fr}})\mathfrak{c}$; the own confined terms of a local sibling pass its window points on
at most $N+1$ levels, which gives $(N+1)A_{\flat}b_{\delta}$, of the same degree in
$\check{\mathsf{T}}_k$ as $NC_0^{\sharp}$. The
term $(N+1)A_{\flat}b_{\delta}$ is absent outside the window regions of the local siblings.

The factor $L^{d}$ in $\mathfrak{l}_j$ is included once for all, so that $\mathfrak{l}_j$ applies to
the blocks of a local sibling $C'$ whose level $j(C')$ equals $j+1$. Every factor
$(L\check{R}^{\uparrow})^d$ in the counts that follow is taken with the value $\check{R}^{\uparrow}$ of
\eqref{4.21:eq-weights-incidences-a}. The level ceilings are ceilings on $\gamma$ of the kind that
makes a stage count times an exponentially small factor at most one, written in supremum form; they
are imposed after $L$, hold for each $L$ when $\gamma$ is small enough, and are not caps in the sense
of Definition~\ref{2.3:def-cap}.

\emph{Rates.} Under the floor $\sigma_{*}\ge c_d+4$ imposed on the rate $\sigma_{*}$, where $c_d+1$ is
the rate of \cite[(1.26)]{B-CMP116}, the excess of $\sigma_{*}$ over $c_d+1$ is used as follows:
two units pay for $e^{\mathsf{e}+\mathsf{f}}\le e^{2d_{\ell}(A)}$
(item (i)), and one unit is displayed as a factor $e^{d_{\ell}(A)}$ in the per-cube bound of the
format class and spent there on cutoffs, own-cube counts and position drift.
\end{definition}

\begin{proof}[Proof of the assertions of Notation~\ref{4.21:not-weights-incidences}]
(a) By \eqref{4.21:eq-strips-geometry-a}, $N\le C_N\,2^{r_{\mathrm{pr}}}\check{R}_k$, that is
$(3\,C_N\,2^{r_{\mathrm{pr}}})^{-1}\le\check{R}_k/(3N)$, which is $\theta_0\le\check{R}_k/(3N)$. Since
$C_N\,2^{r_{\mathrm{pr}}}=4^{r_{\mathrm{pr}}}L\ge1$, also $\theta_0\le\frac13$; hence, if instead $N$ is
read as $\check{R}_k$, as in the declaration of \eqref{4.21:eq-strips-geometry-a}, then
$\check{R}_k/(3N)=\frac13\ge\theta_0$ as well. As
$\theta_0>0$ and the ratio of consecutive terms of the series defining $c_{\mathrm{sib}}$ is
$\bigl((2u+6)/(2u+4)\bigr)^de^{-\theta_0}\to e^{-\theta_0}<1$, the series converges. Since $\theta_0$
depends only on $C_N\,2^{r_{\mathrm{pr}}}=4^{r_{\mathrm{pr}}}L$, the number $c_{\mathrm{sib}}$ depends
on $(d,L,r_{\mathrm{pr}})$ only.

(b) Since $\lambda$ and $c_2$ are non-negative,
\begin{equation*}
(1+3\lambda N+2c_2)\bigl(1+3\lambda(s+1)\bigr)
=1+3\lambda N+2c_2+3\lambda(s+1)(1+3\lambda N+2c_2)
\ge1+3\lambda(N+1+s)+2c_2 .
\end{equation*}
Taking logarithms gives \eqref{4.21:eq-weights-incidences-5}.

(c) The map $s\mapsto\log(1+3\lambda(s+1))$ is non-decreasing and non-negative, and
$r_{\mathrm{pr}}>0$, so $\mathfrak{g}$ is non-decreasing. For $a,s\ge0$, expanding the product gives
$1+3\lambda(a+s+2)\le(1+3\lambda(a+1))(1+3\lambda(s+1))$. With $x:=\log(1+3\lambda(a+1))\ge0$ and
$y:=\log(1+3\lambda(s+1))\ge0$ this reads $\log(1+3\lambda(a+s+2))\le x+y$, and
$1+x+y\le(1+x)(1+y)$. Raising to the power $r_{\mathrm{pr}}$ gives
$\mathfrak{g}(a+s+1)\le\mathfrak{g}(a)\mathfrak{g}(s)$.

(d) By the stage-count bound and the comparison of logarithmic scales in
\eqref{4.21:eq-strips-geometry-a}, $N(g_t)\le C_N\mathsf{T}_t^{r_{\mathrm{pr}}}$ and
$\mathsf{T}_t\le\mathsf{T}_k+c_{k-t}\le\check{\mathsf{T}}_k+1+c_{k-t}$. Since $c_s$ is non-decreasing in
$s$, for $k-t\le N+1$ this gives $\mathsf{T}_t\le\check{\mathsf{T}}_k+1+c_{N+1}$, hence
$N(g_t)\le C_N(\check{\mathsf{T}}_k+2+c_{N+1})^{r_{\mathrm{pr}}}=N^{\mathrm{fr}}$.

Now let $s\ge0$ and $k-t\le N+1+s$. The same two inequalities give $\mathsf{T}_t\le
\check{\mathsf{T}}_k+1+c_{N+1+s}$, and with $N(g_t)\le C_N\mathsf{T}_t^{r_{\mathrm{pr}}}$ this gives the
first inequality of
\eqref{4.21:eq-weights-incidences-6}. Put $\tau:=\check{\mathsf{T}}_k+1+c_N$, which is at least $1$ since
$\check{\mathsf{T}}_k\ge0$ and $c_N\ge0$, and $y:=\log(1+3\lambda(s+1))\ge0$. By (b),
$\check{\mathsf{T}}_k+1+c_{N+1+s}\le \tau+y\le \tau(1+y)$, so
\begin{equation*}
C_N(\check{\mathsf{T}}_k+1+c_{N+1+s})^{r_{\mathrm{pr}}}\le C_N\tau^{r_{\mathrm{pr}}}(1+y)^{r_{\mathrm{pr}}}
=\mathfrak{g}(s)\,C_N\tau^{r_{\mathrm{pr}}}\le\mathfrak{g}(s)N^{\mathrm{fr}},
\end{equation*}
the last step because $c_N\le c_{N+1}$. For the radius, the definition of the separation radius in
\eqref{4.21:eq-strips-geometry-a} together with the same comparison of scales and the computation
just made gives
\begin{equation*}
\check{R}_t\le R(x_t)\le L(\check{\mathsf{T}}_k+1+c_{N+1+s})^{r_{\mathrm{pr}}}
\le\mathfrak{g}(s)\,L\tau^{r_{\mathrm{pr}}} .
\end{equation*}
Finally, by \eqref{4.21:eq-strips-geometry-a},
$N\le C_N(\check{\mathsf{T}}_k+1)^{r_{\mathrm{pr}}}$, so
$c_N\le\log(1+3\lambda C_N(\check{\mathsf{T}}_k+1)^{r_{\mathrm{pr}}}+2c_2)$ and, by
\eqref{4.21:eq-weights-incidences-4}, $L\tau^{r_{\mathrm{pr}}}\le\mathsf{R}^{\uparrow}(\check{\mathsf{T}}_k)$.
This proves \eqref{4.21:eq-weights-incidences-6}.

(e) For $s\ge1$ the $s$-th term of $C_{\mathrm{fr}}$ equals
\begin{equation*}
\exp\Bigl\{d(s+1)\log L+r_{\mathrm{pr}}\log\bigl(1+\log(1+3\lambda(s+1))\bigr)+3-38L^{s}\Bigr\}.
\end{equation*}
Since $\log(1+\log(1+u))\le u$ for $u\ge0$, the exponent is at most
$(d\log L+3\lambda r_{\mathrm{pr}})(s+1)+3-38L^{s}$. As $L>1$, $L^{s}$ outgrows every multiple of
$s$, so the exponent is at most $-s$ for all large $s$, and the series converges. Its value depends
on $(d,L,\lambda,r_{\mathrm{pr}})$ only, so it is fixed after these and before $\gamma$.

(f) Expanding $(3+C_{\mathrm{fr}})\mathfrak{c}$ with the definition of $\mathfrak{c}$ and using
$c_{\mathrm{att}}(3+C_{\mathrm{fr}})=c_{\mathrm{fr}}$,
\begin{equation*}
(3+C_{\mathrm{fr}})\mathfrak{c}=2(3+C_{\mathrm{fr}})K(d)A_{\flat}B_0
+c_{\mathrm{fr}}A_{\flat}C_0^{\sharp}N^{\mathrm{fr}}+(3+C_{\mathrm{fr}})A_{\flat}b_{\delta},
\end{equation*}
and $(N+1)A_{\flat}b_{\delta}+(3+C_{\mathrm{fr}})A_{\flat}b_{\delta}=(N+4+C_{\mathrm{fr}})A_{\flat}b_{\delta}$.

(g) By the comparison of the inverse couplings at the levels of a label with the inverse coupling at
the present level, $x_{h+1}\le x_k(1+3\lambda N+2c_2)$. Hence the
comparison of logarithmic scales used in (d) holds at the level $h+1$ with $c_N$ in place of
$c_{k-t}$, and, with $\tau=\check{\mathsf{T}}_k+1+c_N$ as in (d),
\begin{equation*}
\mathsf{T}_{h+1}\le\mathsf{T}_k+c_N\le\check{\mathsf{T}}_k+1+c_N=\tau .
\end{equation*}
By the definition of the separation radius in \eqref{4.21:eq-strips-geometry-a}, read as in (d),
$\check{R}_{h+1}\le R(x_{h+1})\le L\tau^{r_{\mathrm{pr}}}$, and
$L\tau^{r_{\mathrm{pr}}}\le\mathsf{R}^{\uparrow}(\check{\mathsf{T}}_k)$ was shown at the end of (d).
Since $\check{R}^{\uparrow}=\check{R}_{h+1}$, this proves (g).

(h) By \eqref{4.21:eq-strips-geometry-a}, $\check{R}_k\ge\check{\mathsf{T}}_k^{r_{\mathrm{pr}}}$, and
$N\le C_N(\check{\mathsf{T}}_k+1)^{r_{\mathrm{pr}}}$. Since $\check{\mathsf{T}}_k\ge\mathsf{T}_{\gamma}$ by
the choice of the threshold $\mathsf{T}_{\gamma}$,
the value $\mathsf{T}=\check{\mathsf{T}}_k$ lies in
the range $\mathsf{T}\ge\mathsf{T}_{\gamma}$ of the suprema. Hence the first level ceiling gives
\begin{equation*}
(N+4)\mathfrak{g}(0)e^{4-\check{R}_k}
\le\bigl(C_N(\check{\mathsf{T}}_k+1)^{r_{\mathrm{pr}}}+4\bigr)\mathfrak{g}(0)
e^{4-\check{\mathsf{T}}_k^{r_{\mathrm{pr}}}}\le1 .
\end{equation*}
Multiplying by $e^{-1}$ and dropping the summand $4$ gives $N\mathfrak{g}(0)e^{3-\check{R}_k}\le e^{-1}$;
keeping only the summand $4$ gives $4\mathfrak{g}(0)e^{4-\check{R}_k}\le1$, that is
$\mathfrak{g}(0)e^{3-\check{R}_k}\le\frac14e^{-1}$. The second level ceiling gives
$L^d\mathfrak{g}(0)e^{3/2-\check{R}_k/2}\le\frac14$; since $L^d\mathfrak{g}(0)\ge1$, also
$e^{3/2-\check{R}_k/2}\le\frac14$, and therefore
\begin{equation*}
L^d\mathfrak{g}(0)e^{3-\check{R}_k}
=L^d\mathfrak{g}(0)e^{3/2-\check{R}_k/2}\cdot e^{3/2-\check{R}_k/2}\le\tfrac1{16}\le\tfrac14 .
\end{equation*}
This proves \eqref{4.21:eq-weights-incidences-7}. For fixed $L$ the functions of $\mathsf{T}$ under the
suprema tend to zero as $\mathsf{T}\to\infty$, which is why the level ceilings hold for $\gamma$ small
enough and bound no parameter other than $\gamma$.

(i) By \eqref{4.21:eq-weights-incidences-3} and $N\ge1$ (\eqref{4.21:eq-strips-geometry-a}),
$\mathsf{e}(c)\le d_k([S_c]_k)/N\le d_k([S_c]_k)$ and $\mathsf{f}(c)\le d_k([S_c]_k)$. Hence
$\mathsf{e}(c)+\mathsf{f}(c)\le2d_k([S_c]_k)$, and adding $d_k([S_c]_k)$ to both sides gives the first
bound of item (i). With $A=S_c$ at level $\ell$ and the hull inequality
$d_k([A]_k)\le d_{\ell}(A)$ this gives $e^{\mathsf{e}+\mathsf{f}}\le e^{2d_{\ell}(A)}$. In particular a
confined old term of an arrest in an untreated component, for which $\mathsf{e}(c)\neq0$, is covered
by the first bound of item (i), which is the form used in the per-cube bound of the
format class.
\end{proof}

\begin{lemma}[Per-cube bounds for the format class]\label{4.21:lem-per-cube-bounds}
Consider stage $n+1$ of the operation $\mathbf{R}_k$. Assume the format hypothesis at rate $\kappa$ at
the preparation site, under which in particular both kinds of boundary terms are stored at rate at
least $\kappa$ in the sense of the format bound (f2) of Lemma~\ref{4.21:lem-format-class}, and assume
the lower bound on the radii $\check{R}_\ell$ among the conditions
\eqref{4.21:eq-strips-geometry-a}. The format class $\mathfrak{F}$ and its terms are those of
Lemma~\ref{4.21:lem-format-class}. The weights $W$, $\mathfrak{w}$ and the budgets
$\mathsf{e},\mathsf{f}$ are those of Notation~\ref{4.21:not-weights-incidences}. Here $N$ denotes the
stage bound of $\mathbf{R}_k$: every class component $C'$ has $n(C')\le N$ stages, and $N$ enters the
running spare $(n+1-i)/N$ of the weights; it is not the rank of the gauge group. Then for every level
$\ell$ and every cube $\square'\in\pi_\ell$ the following hold.
\begin{subequations}\label{4.21:eq-per-cube-bounds-a}
\begin{enumerate}
\item[(pc-1)] (Boundary terms.)
\begin{equation}
\sum\bigl\{W(F)\,e^{d_\ell(A_F)} :\ F\in\mathbf{B}^{(\ell)}\cup\mathbf{B}'^{(\ell)},\ A_F\ni\square'\bigr\}
\le 2K(d)\,A_\flat B_0 .
\end{equation}
\item[(pc-2)] (The chains of this $\mathbf{R}_k$.) Let $\ell<k$, so that the terms concerned are confined
terms, and let $\square'$ lie inside the class component $C'$. Then
\begin{equation}
\begin{split}
&\sum\bigl\{W(c):\ c=\mathbb{C}^{(i)}_k[C'](S),\ i\le n(C'),\ m(S)=\ell,\ S\ni\square'\bigr\}\\
&\qquad\le \sum_{i\le n(C')}\bigl\|\mathbb{C}^{(i)}_k[C']\bigr\|^{A,\mathbb{C}}_{(i)}\le N C_0^{\sharp},
\end{split}
\end{equation}
and only the chain of $C'$ has confined terms through $\square'$. For $\ell=k$, the terms of level $k$ of
all the chains of this $\mathbf{R}_k$ satisfy
\begin{equation}
\sum\bigl\{W(c):\ m(S)=k,\ S\ni\square'\bigr\}\le c_{\mathrm{sib}}\,N C_0^{\sharp}.
\end{equation}
For every point $y$,
\begin{equation}
\begin{split}
&\sum\bigl\{W(c):\ c \text{ a chain term of this } \mathbf{R}_k,\ \text{any level},\ S\ni y\bigr\}\\
&\qquad\le c_{\mathrm{sib}}\,N C_0^{\sharp}.
\end{split}
\end{equation}
\item[(pc-3)] (Frozen chains, $k_a<k$.) For every attempt $\mathfrak{a}$ of an earlier step $k_a<k$
whose frozen chain terms belong to $\mathfrak{F}$, and every stage $i$,
\begin{equation}
\sum\bigl\{W(c)\,e^{d_\ell(S)}:\ c=\mathbb{C}^{(i)}_{k_a}(S),\ m(S)=\ell,\ S\ni\square'\bigr\}
\le A_\flat C_0^{\sharp}\,e^{-(\sigma_*-3)D},
\end{equation}
where $D$ is the least value of $d_\ell(S)$ among the terms of the sum. Hence:
\begin{itemize}
\item[($\alpha$)] for $\ell<k$, the top terms of level $\ell$ with $S\ni\square'$, of all attempts at step
$\ell$ and all their stages, satisfy
\begin{equation}
\sum W(c)\,e^{d_\ell(S)}\le c_{\mathrm{att}}\,N(g_\ell)\,A_\flat C_0^{\sharp};
\end{equation}
\item[($\beta$)] the confined terms of level $\ell$ through $\square'$ satisfy
\begin{equation}
\begin{split}
\sum W(c)\,e^{d_\ell(S)}
&\le (k-\ell+1)\max_{t\in[\ell,k]}N(g_t)\,A_\flat C_0^{\sharp}\\
&\le (k-\ell+1)\,C_N\bigl(\mathsf{T}_k+c_{k-\ell}\bigr)^{r_{\mathrm{pr}}}A_\flat C_0^{\sharp}.
\end{split}
\end{equation}
\end{itemize}
\item[(pc-4)] (The terms $\mathfrak{Q}(\mathfrak{a}')$ of live attempts.) For the live attempts
$\mathfrak{a}'$, that is, the old ones
and the siblings with $n^{+}\le n$,
\begin{equation}
\sum\bigl\{W(F)\,e^{d_\ell(A_F)}:\ F\in\mathfrak{Q}(\mathfrak{a}'),\ \mathfrak{a}' \text{ live},\
m_0(F)=\ell,\ A_F\ni\square'\bigr\}\le A_\flat\, b_\delta ,
\end{equation}
top and confined terms together, one level $\ell$ at a time. At a point of a sibling component $C''$ the
local sibling terms occupy at most $k-h+1=N+1$ levels $\ell\in[h,k]$, where $h$ is the level with
$k-h=N$, and hence weigh at most $(N+1)A_\flat b_\delta$ there.
\end{enumerate}
Consequently the sums in \textup{(pc-1)}, \textup{(pc-3)} and \textup{(pc-4)} admit a cutoff;
\textup{(pc-2)} carries no factor $e^{d_\ell}$. Let $C_\bullet$ denote the right-hand side of the bound
concerned (for \textup{(pc-3)}, any one of its three bounds), the sums below running over the same
terms, and let $d_\ell$ stand for $d_\ell(A_F)$, respectively $d_\ell(S)$. Then for every
$D\ge0$
\begin{equation}
\begin{aligned}
\sum_{d_\ell\ge D} W &\le e^{-D}\,C_\bullet,\\
\sum W\,(d_\ell+1)\,e^{d_\ell/3}\,\mathbf{1}[d_\ell\ge D] &\le 3\,e^{-D/2}\,C_\bullet ,
\end{aligned}
\end{equation}
the sums running over the terms of that clause.
\end{subequations}
\end{lemma}

\begin{proof}
\emph{Proof of \textup{(pc-1)}.} By the format hypothesis, the two kinds of boundary terms,
$\mathbf{B}^{(\ell)}$ and $\mathbf{B}'^{(\ell)}$, are both stored at rate at least $\kappa$. This is the
format bound (f2) of Lemma~\ref{4.21:lem-format-class}. For a term $T$ of either kind, we combine:
\begin{itemize}
\item the three rate cases of the analysis at the preparation site;
\item the payment rule for the format class;
\item the label form of the bound on boundary terms at the preparation site.
\end{itemize}
Together they give, on labels,
\begin{equation}\label{4.21:eq-per-cube-bounds-1}
\|T\|\,e^{w(T)}\le A_\flat B_0\,e^{-\sigma_* d_\ell(A_T)} .
\end{equation}
Here $\|T\|$ is the stored size of $T$ and $e^{w(T)}$ its stored decay factor, the two factors from which
$\mathfrak{w}(T)$ is formed in Notation~\ref{4.21:not-weights-incidences}.
Passing from this stored weight to $W(T)\,e^{d_\ell(A_T)}$ costs at most the factor
$e^{\mathsf{e}+\mathsf{f}}e^{d_\ell}$ (Notation~\ref{4.21:not-weights-incidences}). Moreover
$e^{\mathsf{e}+\mathsf{f}}e^{d_\ell}\le e^{3d_\ell}$. Let $c_d$ be the dimensional constant with which
the summation bound \cite[(1.26)]{B-CMP116} is used below. Since $\sigma_*-3\ge c_d+1$ and
$d_\ell(A_T)\ge0$, \eqref{4.21:eq-per-cube-bounds-1} yields
\begin{equation*}
W(T)\,e^{d_\ell(A_T)}\le A_\flat B_0\,e^{-(\sigma_*-3)d_\ell(A_T)}
\le A_\flat B_0\,e^{-(c_d+1)d_\ell(A_T)} .
\end{equation*}
We sum over the labels $A\in\mathbf{D}_\ell$ with $A\ni\square'$, one kind at a time.
By \cite[(1.26)]{B-CMP116},
\begin{equation*}
\sum_{A\in\mathbf{D}_\ell,\ A\ni\square'}e^{-(c_d+1)d_\ell(A)}\le K(d).
\end{equation*}
Each kind therefore contributes at most $K(d)A_\flat B_0$. The two kinds together give
$2K(d)A_\flat B_0$.

\emph{Proof of \textup{(pc-2)}.} We recall the norm of stage $i$ at the preparation site:
\begin{equation}\label{4.21:eq-per-cube-bounds-2}
\|\mathbb{C}\|^{A,\mathbb{C}}_{(i)}
:=\sup_x\sum_{S\ni x}\mathfrak{N}^{A}(c)\,e^{\kappa_A|\mathcal{S}|}\,
e^{a^{(i)}_{m(S)}d_{m(S)}(S)} .
\end{equation}
Here membership is literal membership of the point $x$ in $S$, and every level $m(S)$ enters the one
sum. If $S\ni\square'$, then $S\ni x$ for every point $x$ of $\square'$. Hence a sum over $S\ni\square'$
of the summands of \eqref{4.21:eq-per-cube-bounds-2} is bounded by the norm.

We use the bound $\|\mathbb{C}^{(i)}_k\|^{A,\mathbb{C}}_{(i)}\le C_0^{\sharp}$, and we need it on
labels, wrapped terms included. This bound is supplied by the induction over the earlier stages
$i\le n$ of this $\mathbf{R}_k$, which gives the format bound for the relative chain terms
$\mathbb{C}^{(i)}_k$.
A bound on plain domains only would not serve here, since wrapped terms are summed by their labels.
Let $C''$ be a sibling. Its stages $i\le n(C'')\le n$ are members of the class chain
$\mathbb{C}^{(i)}_k$. A norm taken per component is a sub-sum of the norm of the class chain, so it is
also at most $C_0^{\sharp}$.

The weight $\mathfrak{w}(c)$ is computed at the rate $a^{(n+1)}_m$ of the consuming stage $n+1$.

\emph{Levels $m<k$.} By the rate schedule of the preparation site, the rates below the top level do
not run with the stage: $a^{(n+1)}_m=a_m=a^{(i)}_m$. So the weight of $\mathfrak{w}$ is
the stored one, $\mathsf{e}=\mathsf{f}=0$, and $W=\mathfrak{w}$. The sum over $S\ni\square'$ of one
stage is therefore at most $C_0^{\sharp}$. A term of level $\ell<k$ lies in its own component, by the
statement that a chain term filed below $k$ has its label inside its own class component. Hence only
the chain of $C'$ occurs, and summing over its stages
$i\le n(C')$ gives
\begin{equation*}
\sum_{i\le n(C')}\bigl\|\mathbb{C}^{(i)}_k[C']\bigr\|^{A,\mathbb{C}}_{(i)}
\le n(C')\,C_0^{\sharp}\le N C_0^{\sharp}.
\end{equation*}

\emph{Level $m=k$: the weight.} Let $c$ be a term of level $k$ stored at stage $i\le n$. Here the
rate $a^{(i)}_k$ at which $c$ was stored exceeds the consuming rate $a^{(n+1)}_k$ by the running spare
$(n+1-i)/N\ge 1/N$. Write
$\Pi(c):=\mathfrak{N}^{A}(c)e^{\kappa_A|\mathcal{S}|}e^{a^{(i)}_{k}d_k(S)}$ for the stored summand of
\eqref{4.21:eq-per-cube-bounds-2}. The spare gives $\mathfrak{w}(c)\le\Pi(c)\,e^{-\mathsf{e}(c)}$, and
therefore
\begin{equation*}
W(c)=\mathfrak{w}(c)\,e^{(2/3)\mathsf{e}(c)}\le \Pi(c)\,e^{-\mathsf{e}(c)/3}
=\Pi(c)\,e^{-d_k(S)/(3N)} .
\end{equation*}

\emph{Level $m=k$: the geometry.} Let $c$ be a term of level $k$ of the chain of a component $C''$,
with $S\ni\square'$. Its label $S$ meets its band in $C''$. By
Lemma~\ref{4.21:lem-tree-length-facts}\,(g4),
\begin{equation*}
d_k(S)\ge \operatorname{dist}_k(\square',C'')-1 .
\end{equation*}
The class components are pairwise at least $\check{R}_k$ apart, by \eqref{4.21:eq-strips-geometry-a}.
For $u\ge0$, the components with $\operatorname{dist}_k(\square',C'')<(u+1)\check{R}_k$ all meet a box of
side $2(u+1)\check{R}_k+1$. By the packing statement, Lemma~\ref{4.21:lem-tree-length-facts}\,(g6),
there are at most $(2u+4)^d$ of them. A component with $\operatorname{dist}_k(\square',C'')\ge u\check{R}_k$
carries a factor at most
\begin{equation*}
e^{-(u\check{R}_k-1)/(3N)}=e^{1/(3N)}e^{-u\check{R}_k/(3N)}\le e^{1/3}e^{-u\theta_0}.
\end{equation*}
This uses $N\ge1$ and $\theta_0\le\check{R}_k/(3N)$; the latter follows from the size condition on
$\check{R}_k$ in \eqref{4.21:eq-strips-geometry-a} alone.

\emph{Level $m=k$: the sum.} Group the components $C''$ according to the integer $u\ge0$ with
$u\check{R}_k\le\operatorname{dist}_k(\square',C'')<(u+1)\check{R}_k$. Each component has at most $N$
stages, each stage contributes at most $C_0^{\sharp}$ before the factor above, and we put
\begin{equation*}
\Sigma_{\theta_0}:=\sum_{u\ge0}(2u+4)^d e^{-u\theta_0}.
\end{equation*}
Then
\begin{equation}\label{4.21:eq-per-cube-bounds-3}
\sum_{C''}\sum_i\sum_{S\ni\square'}W(c)\le N C_0^{\sharp}\,e^{1/3}\,\Sigma_{\theta_0}
\le c_{\mathrm{sib}}\,N C_0^{\sharp},
\end{equation}
the last step by the choice of $c_{\mathrm{sib}}$ in \eqref{4.21:eq-weights-incidences-a}.

\emph{The point form.} Let $y$ be a point. The terms of level below $k$ whose label contains $y$ are
those of the local chain, the chain of the class component $C'$ containing $y$. They lie inside
$\sum_i\|\mathbb{C}^{(i)}_k[C']\|^{A,\mathbb{C}}_{(i)}\le NC_0^{\sharp}$, all their levels at once,
because \eqref{4.21:eq-per-cube-bounds-2} sums every level at a point. The terms of level $k$ of all
chains add at most $e^{1/3}\Sigma_{\theta_0}NC_0^{\sharp}$; this is the argument just given, with the
point $y$ in place of $\square'$. The terms of level $k$ of the local chain are thereby counted twice,
which only enlarges the bound. The total is at most $[1+e^{1/3}\Sigma_{\theta_0}]NC_0^{\sharp}$.

It remains to compare this with $c_{\mathrm{sib}}$. The term $u=0$ shows
$\Sigma_{\theta_0}\ge 4^d$, and $(e-e^{1/3})4^d\ge1$. Hence
\begin{equation*}
e\,\Sigma_{\theta_0}=e^{1/3}\Sigma_{\theta_0}+(e-e^{1/3})\Sigma_{\theta_0}
\ge e^{1/3}\Sigma_{\theta_0}+1 .
\end{equation*}
So the total is at most $e\,\Sigma_{\theta_0}NC_0^{\sharp}\le c_{\mathrm{sib}}NC_0^{\sharp}$, by the
same choice of $c_{\mathrm{sib}}$.

\emph{Proof of \textup{(pc-3)}.} We use the third case of the payment rule, together with the rule for
the consumption of left-overs. A left-over of level $\ell$ is consumed at the rate of the output.
\begin{itemize}
\item If $\ell<k_a$, it is stored above $\ell$ by its consuming term.
\item If $\ell=k_a$, it is treated as a boundary term of rate $\kappa$, stored at rate
$(1+3\beta)\kappa$, where $\beta$ is the constant of the rate schedule of the preparation site.
\end{itemize}
In both cases it carries the hull factor $A_\flat$ and the surplus $\sigma_*$, by the three rate cases
of the analysis at the preparation site. Thus
\begin{equation*}
\mathfrak{w}(c)\le A_\flat\bigl[\mathfrak{N}^{A}(c)\,e^{\kappa_A|\mathcal{S}|}\,
e^{a\,d_\ell(S)}\bigr]e^{-\sigma_* d_\ell(S)} ,
\end{equation*}
where $a=a^{(i)}_\ell$ is the rate of stage $i$ of step $k_a$ at level $\ell$, so that the bracket is the
summand of \eqref{4.21:eq-per-cube-bounds-2} for $\mathbb{C}^{(i)}_{k_a}$.
Summed over $S\ni\square'$ within one stage, the bracket is at most
$\|\mathbb{C}^{(i)}_{k_a}\|^{A,\mathbb{C}}_{(i)}\le C_0^{\sharp}$. This is the format bound that the
induction over stages gives for the relative chain terms at the earlier step $k_a$, on labels. Since
$e^{\mathsf{e}+\mathsf{f}+d_\ell}\le e^{3d_\ell}$, the sum of (pc-3) is at most
\begin{equation*}
A_\flat C_0^{\sharp}\sup_{d_\ell\ge D}e^{-(\sigma_*-3)d_\ell}=A_\flat C_0^{\sharp}e^{-(\sigma_*-3)D},
\end{equation*}
because $\sigma_*-3\ge c_d+1>0$.

\emph{Proof of \textup{($\alpha$)}, geometry.} A top term of level $\ell$ belongs to an attempt
$\mathfrak{a}$ at step $\ell$. Its label $S$ meets both its band, which lies in $X_{\mathfrak{a}}$, and
$\square'$. By Lemma~\ref{4.21:lem-tree-length-facts}\,(g4),
\begin{equation*}
d_\ell(S)\ge\operatorname{dist}_\ell(\square',X_{\mathfrak{a}})-1 .
\end{equation*}
The attempts at step $\ell$ lie on distinct components of the class at $\mathbf{R}_\ell$. These are
pairwise at least $\check{R}_\ell$ apart, by \eqref{4.21:eq-strips-geometry-a}. By
Lemma~\ref{4.21:lem-tree-length-facts}\,(g6), those at distance less than $(u+1)\check{R}_\ell$ from
$\square'$ number at most $(2u+4)^d$. Each attempt has at most $N(g_\ell)$ stages.

\emph{Proof of \textup{($\alpha$)}, the sum.} For an attempt at distance at least $u\check{R}_\ell$,
the bound just proved applies with $D\ge(u\check{R}_\ell-1)_+$ and $\sigma_*-3\ge c_d+1$. Hence
\begin{equation*}
\sum_{\mathfrak{a}}\sum_i\sum W(c)\,e^{d_\ell(S)}
\le N(g_\ell)A_\flat C_0^{\sharp}\sum_{u\ge0}(2u+4)^d e^{-(c_d+1)(u\check{R}_\ell-1)_+} .
\end{equation*}
The term $u=0$ equals $4^d$. For $u\ge1$ we have
\begin{equation*}
(2u+4)^d=4^d(1+u/2)^d\le 4^d e^{du/2}.
\end{equation*}
By the condition on the radii in \eqref{4.21:eq-strips-geometry-a}, $\check{R}_\ell\ge38$; together
with $c_d\ge1$, so that $c_d+1\ge2$, this gives for $u\ge1$
\begin{equation*}
(c_d+1)(u\check{R}_\ell-1)\ge 2(38u-1)\ge 74u .
\end{equation*}
At $d=4$ the term of index $u\ge1$ is therefore at most $4^d e^{2u}e^{-74u}=4^d e^{-72u}$, and
$\sum_{u\ge1}e^{-72u}\le1$. The whole series is thus at most $2\cdot 4^d\le 3\cdot 4^d$, and
\begin{equation*}
\sum_{\mathfrak{a}}\sum_i\sum W(c)\,e^{d_\ell(S)}\le 3\cdot 4^d\,N(g_\ell)A_\flat C_0^{\sharp}.
\end{equation*}
This is ($\alpha$) with $c_{\mathrm{att}}:=3\cdot4^d$. The restriction $\ell<k$ is needed here. Chain
terms of the same step carry only the spare $(n+1-i)/N$, not the surplus $\sigma_*$ of the third case
of the payment rule. They are treated in \textup{(pc-2)}.

\emph{Proof of \textup{($\beta$)}.} A confined term of level $\ell$ with $S\ni\square'$ belongs to an
attempt with $X_{\mathfrak{a}}\ni\square'$ and $k_a\ge\ell$. By the nesting property of live arrested
attempts, the live attempts whose regions contain a given point form a chain $(\mathfrak{a}_s)_s$ with
\begin{equation*}
k_{a_{s+1}}\ge k_{a_s}+N_{k_{a_{s+1}}}>k_{a_s}.
\end{equation*}
So there is at most one such attempt per step $t\in[\ell,k]$, that is, at most $k-\ell+1$ attempts.
Each has at most
\begin{equation*}
N(g_t)\le C_N\bigl(\mathsf{T}_k+c_{k-t}\bigr)^{r_{\mathrm{pr}}}
\le C_N\bigl(\mathsf{T}_k+c_{k-\ell}\bigr)^{r_{\mathrm{pr}}}
\end{equation*}
stages. For each attempt and stage, the bound of \textup{(pc-3)} together with $\sigma_*-3>0$ and
$D\ge0$ gives at most $A_\flat C_0^{\sharp}$. Multiplying the three counts gives both inequalities
of ($\beta$).

\emph{Proof of \textup{(pc-4)}.} We apply the bound on the terms $\mathfrak{Q}(\mathfrak{a}')$ of a
live attempt at the rate
$r':=(\text{rate of the hull})+2+1$. The rate of the hull is at most $(8/5)\kappa$, so
$r'\le(8/5)\kappa+3\le 2\kappa+3$. The hull factor is $A_\flat$, as for the boundary terms. At this
rate and with this hull factor, that bound, summed over the terms $F$ with $A_F\ni\square'$ filed at one
level $\ell$, is the bound $A_\flat b_\delta$ of \textup{(pc-4)}, top and confined terms together.

\emph{Proof of the cutoff forms.} As above, $C_\bullet$ is the right-hand side of the bound concerned;
it bounds $\sum W e^{d_\ell}$ over the terms of that bound.
Since $\mathbf{1}[d_\ell\ge D]\le e^{-D}e^{d_\ell}$,
\begin{equation*}
\sum_{d_\ell\ge D}W\le e^{-D}\sum W e^{d_\ell}\le e^{-D}C_\bullet .
\end{equation*}
For the second bound, write
\begin{equation*}
(d_\ell+1)e^{d_\ell/3}\mathbf{1}[d_\ell\ge D]
=e^{d_\ell}\cdot(d_\ell+1)e^{-d_\ell/6}\cdot e^{-d_\ell/2}\mathbf{1}[d_\ell\ge D].
\end{equation*}
The function $t\mapsto(t+1)e^{-t/6}$ on $t\ge0$ is maximal at $t=5$, where it equals
$6e^{-5/6}<3$. Also $e^{-d_\ell/2}\mathbf{1}[d_\ell\ge D]\le e^{-D/2}$. Hence the left side is at most
$3e^{-D/2}e^{d_\ell}$. Summing against $W$ gives the second cutoff bound in
\eqref{4.21:eq-per-cube-bounds-a}.
\end{proof}

\begin{remark}[On the constant $c_{\mathrm{sib}}$]
In \eqref{4.21:eq-per-cube-bounds-2} the norm is that of the class chain. With this reading,
$\sum_{\text{all chains}}\mathfrak{w}\,e^{\mathsf{e}}\le NC_0^{\sharp}$ holds at a point directly, and
$c_{\mathrm{sib}}$ could be replaced by $1$. The constant is kept, at $\theta_0$, so that
Lemma~\ref{4.21:lem-per-cube-bounds} also holds when the norm is read per component.
\end{remark}

\begin{corollary}[The per-cube constant of the format class]\label{4.21:cor-per-cube-constant}
Work at stage $n+1$ of the operation $\mathbf{R}_k$, in the setting and under the hypotheses of
Lemma~\ref{4.21:lem-per-cube-bounds}, with the format class $\mathfrak{F}$, the labels $A_c$ and
the weights $W(c)$ of Notation~\ref{4.21:not-weights-incidences}. Put
\begin{equation}\label{4.21:eq-per-cube-constant-a}
\begin{gathered}
\mathsf{b}^{P}(s) := A_{\flat}\bigl[\,b_0+(s+1)\,b_1\,(\mathsf{T}_k+c_s)^{r_{\mathrm{pr}}}\bigr],
\qquad s\ge 0,\\
b_0 := 2K(d)B_0+b_{\delta}+c_{\mathrm{sib}}N(g_k)C_0^{\sharp},
\qquad
b_1 := (2+c_{\mathrm{att}})\,C_N\,C_0^{\sharp}.
\end{gathered}
\end{equation}
In this corollary and its proof, $b_0$ and $b_1$ denote the two constants just defined, not the
constant $b_0$ of the glossary (Notation~\ref{0.2:not-glossary}). For $c\in\mathfrak{F}$ let
$\iota(c)=1$ if $c$ is a leaf (Notation~\ref{4.21:not-weights-incidences}) or a confined old
term in the classification of terms of Lemma~\ref{4.21:lem-per-cube-bounds}, and $\iota(c)=0$
otherwise. Then for every level $\ell\le k$ and every cube
$\square'\in\pi_\ell$,
\begin{equation}\label{4.21:eq-per-cube-constant-1}
\sum_{\substack{c\in\mathfrak{F}\ \text{of level}\ \ell\\ A_c\ni\square'}}
W(c)\,e^{d_\ell(A_c)\,\iota(c)}
\;\le\; \mathsf{b}^{P}(k-\ell).
\end{equation}
Consequently $\mathsf{b}^{P}$ is a per-cube majorant for $\mathfrak{F}$ in the sense of the
anchoring lemma: it satisfies that lemma's per-cube hypothesis for all probes, and its per-cube
hypothesis for a class of probes, with parameter $\varsigma_0:=1$, for every class of probes
under which the chain terms of this $\mathbf{R}_k$ and the local sibling terms (in the sense of
Lemma~\ref{4.21:lem-per-cube-bounds}) have no cube of their own finer than the probe. Moreover
$\mathsf{b}^{P}$ is the standard per-cube majorant of the remark to the anchoring lemma, with the
constants $c_{\mathrm{att}}$ and $c_{\mathrm{sib}}$ absorbed into $b_1$ and $b_0$, and it obeys the
growth law for per-cube majorants with
\begin{equation}\label{4.21:eq-per-cube-constant-2}
p_{\mathsf{b}} = 1+r_{\mathrm{pr}},\qquad C_{\mathsf{b}} = (1+3\lambda)^{r_{\mathrm{pr}}},
\end{equation}
where $\lambda$ is the constant of that name in the growth law for per-cube majorants.
\end{corollary}

\begin{proof}
Fix $\ell\le k$ and $\square'\in\pi_\ell$. The sum on the left of
\eqref{4.21:eq-per-cube-constant-1} is split according to the kind of the term $c$, and each
part is bounded by the corresponding one of the per-cube bounds (pc-1)--(pc-4) of
Lemma~\ref{4.21:lem-per-cube-bounds}, displayed in \eqref{4.21:eq-per-cube-bounds-a}; these
bounds hold at the given level $\ell$ and the given cube $\square'$ because the corollary is
stated under the hypotheses of that lemma.

\emph{Level $\ell=k$.} Split the sum into three parts: the leaves of level $k$, which are the
terms of $\mathbf{B}^{(k)}\cup\mathbf{B}'^{(k)}$ bounded in (pc-1), with their factor
$e^{d_k(A_c)}$; the terms filed at level $k$ that are covered by the case $\ell=k$ of
(pc-2); and the terms bounded by (pc-4). By (pc-1) the first part is at most
$2K(d)A_{\flat}B_0$; by (pc-2) the second part is at most $c_{\mathrm{sib}}N(g_k)C_0^{\sharp}$;
the third part is at most the bound of (pc-4), which enters through the summand $b_{\delta}$.
These three bounds are included in $A_{\flat}b_0$ through the summands $2K(d)B_0$,
$c_{\mathrm{sib}}N(g_k)C_0^{\sharp}$ and $b_{\delta}$ of $b_0$; hence the sum is at most
$A_{\flat}b_0$. Since $b_1\ge0$ and $(s+1)(\mathsf{T}_k+c_s)^{r_{\mathrm{pr}}}\ge 0$,
\begin{equation*}
A_{\flat}b_0 \;\le\; A_{\flat}\bigl[\,b_0+b_1(\mathsf{T}_k+c_0)^{r_{\mathrm{pr}}}\bigr]
\;=\;\mathsf{b}^{P}(0),
\end{equation*}
which is \eqref{4.21:eq-per-cube-constant-1} for $\ell=k$.

\emph{Level $\ell<k$.} By (pc-2), the confined chain terms of this $\mathbf{R}_k$ of level $\ell$
whose label contains $\square'$ are those of the local chain only, and they contribute at most
$NC_0^{\sharp}$, where $N$ is the stage-count bound of (pc-2) in
Lemma~\ref{4.21:lem-per-cube-bounds} (not the rank of the gauge group). The stage count is
bounded in terms of the depth $k-\ell$:
\begin{equation*}
NC_0^{\sharp}\;\le\; C_N C_0^{\sharp}\,(\mathsf{T}_k+c_{k-\ell})^{r_{\mathrm{pr}}}.
\end{equation*}
The contribution of the local chain, together with the contributions
bounded by (pc-3), in its two parts $(\alpha)$ and $(\beta)$, is included in
$(k-\ell+1)\,b_1\,(\mathsf{T}_k+c_{k-\ell})^{r_{\mathrm{pr}}}$. The leaves of level $\ell$, which
are the terms of $\mathbf{B}^{(\ell)}\cup\mathbf{B}'^{(\ell)}$, with their factor
$e^{d_\ell(A_c)}$, are bounded by (pc-1), and the terms bounded by (pc-4) by that bound; as at
level $k$, these are included in $A_{\flat}b_0$. Hence the sum is at most
\begin{equation*}
A_{\flat}\bigl[\,b_0+(k-\ell+1)\,b_1\,(\mathsf{T}_k+c_{k-\ell})^{r_{\mathrm{pr}}}\bigr]
\;=\;\mathsf{b}^{P}(k-\ell),
\end{equation*}
by \eqref{4.21:eq-per-cube-constant-a} with $s=k-\ell$. This proves
\eqref{4.21:eq-per-cube-constant-1} for every $\ell\le k$.

\emph{The consequences.} Inequality \eqref{4.21:eq-per-cube-constant-1} holds for every cube
$\square'$ of every level $\ell\le k$, with no restriction on the probe; this is the anchoring
lemma's per-cube hypothesis for all probes, with majorant $\mathsf{b}^{P}$. For a class of probes
under which the chain terms of this $\mathbf{R}_k$ and the local sibling terms have no cube of
their own finer than the probe, the same inequality gives that lemma's per-cube hypothesis for
the class, with $\varsigma_0:=1$. Finally, $\mathsf{b}^{P}$ is the standard per-cube majorant of
the remark to the anchoring lemma, with $c_{\mathrm{att}}$ absorbed into $b_1$ and
$c_{\mathrm{sib}}$ into $b_0$; by that remark it obeys the growth law for per-cube majorants with
the exponent $p_{\mathsf{b}}$ and the constant $C_{\mathsf{b}}$ of
\eqref{4.21:eq-per-cube-constant-2}.
\end{proof}

\begin{lemma}[Bounds at a window point and a window block]\label{4.21:lem-window-bounds}
Consider stage $n+1$ of the operation $\mathbf{R}_k$. Assume the format hypothesis under which the
per-cube bounds for the format class (Lemma~\ref{4.21:lem-per-cube-bounds}) are stated, the
region-size condition of \eqref{4.21:eq-strips-geometry-a}, and the level condition
$(\mathrm{c\text{-}lev})'$ of \eqref{4.21:eq-weights-incidences-a}. Here $N$ is the parameter of
Notation~\ref{4.21:not-weights-incidences}, not the rank of the gauge group; in particular the
levels $h,\dots,k$ of the operation number at most $N+1$. Let $\mathfrak{t}^{\mathrm{win}}$
be the constant of \eqref{4.21:eq-weights-incidences-a},
\begin{equation*}
\mathfrak{t}^{\mathrm{win}}=c_{\mathrm{sib}}\,N\,C_0^{\sharp}+(N+1)\,A_{\flat}\,b_{\delta}
+(3+C_{\mathrm{fr}})\,\mathfrak{c},
\end{equation*}
and for a window block $\square_0$ put
\begin{equation*}
\mathfrak{T}^{\mathrm{co}}(\square_0):=\sum\bigl\{W(c):(\square_0,c)\ \text{coarse}\bigr\},
\qquad
\mathfrak{T}(\square_0):=\sum\bigl\{W(c):\square_0\in\mathcal{A}_c\bigr\},
\end{equation*}
both sums running over $c\in\mathfrak{F}$, with the incidences of
Notation~\ref{4.21:not-weights-incidences}. Then the following hold.

For every window point $y$ and every window block $\square_0$ of zone $i_0$, with $j$ the top
zone of the present operation as in \eqref{4.21:eq-weights-incidences-a},
\begin{equation}\label{4.21:eq-window-bounds-a}
\begin{aligned}
&\text{(a)}\quad \sum_{c\in\mathfrak{F}:\,A_c\ni y}W(c)\le\mathfrak{t}^{\mathrm{win}},\\
&\text{(b)}\quad \mathfrak{T}^{\mathrm{co}}(\square_0)\le(5L)^d\,\mathfrak{t}^{\mathrm{win}}
\quad\text{and}\quad
\mathfrak{T}(\square_0)\le\mathfrak{l}_j\,\mathfrak{t}^{\mathrm{win}}.
\end{aligned}
\end{equation}

\begin{itemize}
\item[(c)] (The window's tile number.) The constant $\mathfrak{l}_j$ of
\eqref{4.21:eq-weights-incidences-a} satisfies
\begin{equation*}
\mathfrak{l}_j\le c_{\mathfrak{l}}\,\mathsf{l}_j^2,\qquad c_{\mathfrak{l}}:=3(5L^4)^d,
\end{equation*}
where $d=4$ is used in reading the data-size condition on the constants $\mathsf{l}_j$; and
\begin{equation*}
\mathfrak{l}_j\le 3\bigl(5L^3\,\mathsf{R}^{\uparrow}(\check{\mathsf{T}}_k)\bigr)^d,
\end{equation*}
a bound depending on $L$ and $\check{\mathsf{T}}_k$ only, by the comparison of
$\check{R}^{\uparrow}$ with $\mathsf{R}^{\uparrow}(\check{\mathsf{T}}_k)$ in
\eqref{4.21:eq-strips-geometry-a}; here $L^{N_0}=L\check{R}_{k_0+1}\le L\check{R}^{\uparrow}$.
\end{itemize}
No constant in \textup{(a)}--\textup{(c)} depends on an age, or on the number of pockets, of
windows, of attempts or of components; $N$ enters $\mathfrak{t}^{\mathrm{win}}$ to the first
power, through the norms of the chains, through the number, at most $N+1$, of levels at which
the terms $\mathfrak{Q}(\mathfrak{a})$ of a sibling attempt $\mathfrak{a}$ are filed, and through
$N^{\mathrm{fr}}$.
\end{lemma}

\begin{proof}
Throughout, the levels, bases and floors are those of the lemma on what lies under a block
(Lemma~\ref{4.21:lem-under-block}), and the kinds of terms in $\mathfrak{F}$ are those of
Lemma~\ref{4.21:lem-format-class}. We call a term $c\in\mathfrak{F}$ of level $\ell$ a
\emph{term of boundary type} if it is a member of $\mathbf{B}^{(\ell)}\cup\mathbf{B}'^{(\ell)}$
(item $(\alpha)$ of the format class), a top frozen chain term, or a top term of item $(\delta)$.

\emph{Part (a).} Let $y$ be a window point and let $b:=b(y)$ be its base. By
Lemma~\ref{4.21:lem-window-probe} (what a window probe meets), $b\ge h$. We bound separately the
contributions to $\sum_{c\in\mathfrak{F}:\,A_c\ni y}W(c)$ of four groups of terms.

\emph{Chains of the present operation $\mathbf{R}_k$.} The bound (pc-2) in
\eqref{4.21:eq-per-cube-bounds-a} of Lemma~\ref{4.21:lem-per-cube-bounds}, taken in its point form
at $y$ and applied to all chains and
all filing levels at once, bounds their contribution by $c_{\mathrm{sib}}NC_0^{\sharp}$.

\emph{Confined old terms.} By Remark~\ref{4.21:rem-under-block-remarks}, in the form
\eqref{4.21:eq-under-block-remarks-a}, no confined old term $c$ has $A_c\ni y$. Their contribution
is zero.

\emph{Local sibling terms.} These occur only if $y$ lies in the window region of a sibling
attempt with stage count $n^{+}\le n$. In that case the bound (pc-4) of
Lemma~\ref{4.21:lem-per-cube-bounds}, applied without cutoff at each of the at most $N+1$ levels
$\ell\in[h,k]$, bounds the contribution of each level by $A_{\flat}b_{\delta}$. The total is at most
$(N+1)A_{\flat}b_{\delta}$.

\emph{Terms of boundary type.} Let $\ell$ be a level and let $\square'_\ell(y)$ be the
$\pi_\ell$-cube containing $y$. The terms of boundary type of level $\ell$ with $A_c\ni y$ are
controlled by part (e) of Lemma~\ref{4.21:lem-under-block}, together with
Lemma~\ref{4.21:lem-per-cube-bounds} applied at the cube $\square'_\ell(y)$. We distinguish the
levels according to their position relative to $b$.

Short levels $\ell\in\{b,b+1\}$. Here no cutoff is used. For $\ell\ge b\ge h$ one has
$k-\ell\le N$ and $N(g_\ell)\le N^{\mathrm{fr}}$. The bounds (pc-1), (pc-3)$(\alpha)$ and (pc-4)
of Lemma~\ref{4.21:lem-per-cube-bounds} at $\square'_\ell(y)$ therefore give, at each of the two
levels, at most
\begin{equation*}
2K(d)A_{\flat}B_0+c_{\mathrm{att}}N^{\mathrm{fr}}A_{\flat}C_0^{\sharp}+A_{\flat}b_{\delta}
=\mathfrak{c},
\end{equation*}
the equality being the definition of $\mathfrak{c}$ in \eqref{4.21:eq-weights-incidences-a}. The
two short levels together contribute at most $2\mathfrak{c}$.

Two remarks complete this case. At $\ell=k$ there is no top frozen chain term, and
(pc-3)$(\alpha)$ is stated for levels $\ell<k$ only; the bound at that level holds a fortiori.
When $y$ lies in the window
region of a sibling attempt, the contribution of (pc-4) at these two levels is also contained in
the bound for the local sibling terms. It is thereby counted twice, which only enlarges the bound.

The level $\ell=b-1$. The cutoff below the base, (L$\downarrow$) in
the corresponding display, gives the cutoff $D=\check{R}_k-3$ in the cutoff form
\eqref{4.21:eq-per-cube-bounds-a} of
Lemma~\ref{4.21:lem-per-cube-bounds}. Since $k-(b-1)\le N+1$, we have
$N(g_{b-1})\le\mathfrak{g}(0)N^{\mathrm{fr}}$. The contribution of this level is therefore at most
\begin{equation*}
\mathfrak{g}(0)\,e^{3-\check{R}_k}\,\mathfrak{c}
\le\tfrac14 e^{-1}\,\mathfrak{c},
\end{equation*}
the last inequality by the level condition $(\mathrm{c\text{-}lev})'$.

The levels $\ell=b-1-s$ with $s\ge1$. The cutoff below the base gives $D\ge 38L^s-3$. Since
$k-\ell\le N+1+s$, we have $N(g_\ell)\le\mathfrak{g}(s)N^{\mathrm{fr}}$. The level $\ell=b-1-s$
thus contributes at most $\mathfrak{g}(s)e^{3-38L^s}\mathfrak{c}$, and all these levels together
contribute at most
\begin{equation*}
\mathfrak{c}\sum_{s\ge1}\mathfrak{g}(s)\,e^{3-38L^s}\le C_{\mathrm{fr}}\,\mathfrak{c}.
\end{equation*}
The inequality holds because $L^{d(s+1)}\ge1$ and
$\sum_{s\ge1}L^{d(s+1)}\mathfrak{g}(s)e^{3-38L^s}\le C_{\mathrm{fr}}$ by the definition of
$C_{\mathrm{fr}}$ in \eqref{4.21:eq-weights-incidences-a}.

The levels $\ell\ge b+2$. Since $y$ is a window point, the cutoff above the base,
(L$\uparrow$) in the corresponding display, gives $D\ge\check{R}_k-3$ at each such level.
There are at most $k-b-1\le N$ such levels, and at each of them $N(g_\ell)\le N^{\mathrm{fr}}$.
Their contribution is at most
\begin{equation*}
N\,e^{3-\check{R}_k}\,\mathfrak{c}\le e^{-1}\,\mathfrak{c},
\end{equation*}
again by the level condition $(\mathrm{c\text{-}lev})'$.

Adding the four ranges of levels, the terms of boundary type contribute at most
\begin{equation*}
\bigl[2+\tfrac14e^{-1}+C_{\mathrm{fr}}+e^{-1}\bigr]\mathfrak{c}\le(3+C_{\mathrm{fr}})\,\mathfrak{c},
\end{equation*}
since $\tfrac54e^{-1}<1$. Together with the bounds for the chains and for the local sibling terms,
this gives
\begin{equation*}
\sum_{c\in\mathfrak{F}:\,A_c\ni y}W(c)
\le c_{\mathrm{sib}}NC_0^{\sharp}+(N+1)A_{\flat}b_{\delta}+(3+C_{\mathrm{fr}})\mathfrak{c}
=\mathfrak{t}^{\mathrm{win}},
\end{equation*}
which is (a) of \eqref{4.21:eq-window-bounds-a}.

\emph{Part (b).} Let $\square_0$ be a window block of zone $i_0$, and put
$\mathfrak{f}:=\mathfrak{f}(i_0)$, the floor of part (b) of Lemma~\ref{4.21:lem-under-block}. Let
$C'$ be the class component holding the block's cell, if there is one, and put $J:=j(C')\le j+1$.

The probe $\tilde{\square}_0$ is the block $\square_0$ enlarged by one layer of
$\pi_{i_0}$-cubes; it is a union of $\pi_{i_0}$-cubes, and its side is $4$ such cubes, in
particular at most $5$, in agreement with the general bound on the side of a probe.
The partitions $\pi_\ell$ are nested. Hence every cube of $\pi_\ell$ with $\ell\le i_0$ that meets
$\tilde{\square}_0$ lies in $\tilde{\square}_0$. Each point of such a cube lies in a cell met by
the probe, and is therefore a window point by Lemma~\ref{4.21:lem-window-probe}.

\emph{Coarse incidences.} Let $(\square_0,c)$ be coarse. Then an own cube $Q$ of $c$ with
$\operatorname{lev}(Q)\ge i_0-1$ meets $\tilde{\square}_0$, and $Q$ contains a $\pi_{i_0-1}$-cube
$Q'\subset\tilde{\square}_0$. Every point $y'$ of $Q'$ is a window point with $A_c\ni y'$. The
probe $\tilde{\square}_0$ contains at most $(5L)^d$ cubes of $\pi_{i_0-1}$. Applying (a) at one
point of each of them gives $\mathfrak{T}^{\mathrm{co}}(\square_0)\le(5L)^d\mathfrak{t}^{\mathrm{win}}$.

\emph{All incidences: the kinds of terms.} Since $\square_0$ is a window block, there are no
incidences through strips. A fine incidence passes through an own cube $Q\subset\tilde{\square}_0$
of some level $\ell\le i_0-2$. The kinds of terms behave as follows.
\begin{itemize}
\item For the chain terms of the present $\mathbf{R}_k$ with filing level below $k$, and for the
local sibling terms, parts (d) and (b) of Lemma~\ref{4.21:lem-under-block} give
$\ell\ge b\ge\mathfrak{f}$ at the points of $Q$. Terms with filing level $k$ are never fine.
\item There are no confined old terms.
\item For terms of boundary type the level $\ell$ is free. If $\ell\le\mathfrak{f}-2$, then
$\ell\le b(y')-2$ at every point $y'$ of $\tilde{\square}_0$, and the cutoff below the base gives
\begin{equation*}
d_\ell(A_c)\ge 38L^{\mathfrak{f}-1-\ell}-3.
\end{equation*}
\end{itemize}

\emph{All incidences: positions.} Let $P$ be the set of the $\pi_\ell$-cubes contained in
$\tilde{\square}_0$ with $\mathfrak{f}-1\le\ell\le i_0-1$. Since $L^{-d}\le\tfrac12$,
\begin{equation*}
\#P\le\sum_{\ell=\mathfrak{f}-1}^{i_0-1}\bigl(5L^{i_0-\ell}\bigr)^d
\le 2\bigl(5L^{i_0-\mathfrak{f}+1}\bigr)^d.
\end{equation*}
Every coarse incidence puts a cube of $P$ inside $A_c$, namely the cube of $\pi_{i_0-1}$ found
above. So does every fine incidence with $\ell\ge\mathfrak{f}-1$. Applying (a) at one window
point of each position, these terms contribute at most $\#P\cdot\mathfrak{t}^{\mathrm{win}}$.

There remain the terms of boundary type of level $\ell=\mathfrak{f}-1-s$, $s\ge1$, that meet
$\tilde{\square}_0$. Each of them passes through one of the at most
$(5L^{i_0-\ell})^d=(5L^{i_0-\mathfrak{f}+1})^dL^{ds}$ cubes of $\pi_\ell$ in $\tilde{\square}_0$.
Since $\mathfrak{f}\ge h$, we have $k-\ell\le N+1+s$. By the displayed lower bound on
$d_\ell(A_c)$ and the first cutoff form of Lemma~\ref{4.21:lem-per-cube-bounds}, each such cube
carries at most $\mathfrak{g}(s)\mathfrak{c}e^{3-38L^s}$. These terms therefore contribute at most
\begin{equation*}
\begin{aligned}
\bigl(5L^{i_0-\mathfrak{f}+1}\bigr)^d\,\mathfrak{c}\sum_{s\ge1}L^{ds}\mathfrak{g}(s)e^{3-38L^s}
&\le\bigl(5L^{i_0-\mathfrak{f}+1}\bigr)^d\,C_{\mathrm{fr}}\,\mathfrak{c}\\
&\le\bigl(5L^{i_0-\mathfrak{f}+1}\bigr)^d\,\mathfrak{t}^{\mathrm{win}},
\end{aligned}
\end{equation*}
where the last step uses $C_{\mathrm{fr}}\mathfrak{c}\le(3+C_{\mathrm{fr}})\mathfrak{c}\le
\mathfrak{t}^{\mathrm{win}}$. Altogether,
\begin{equation}\label{4.21:eq-window-bounds-1}
\mathfrak{T}(\square_0)\le 3\bigl(5L^{i_0-\mathfrak{f}+1}\bigr)^d\,\mathfrak{t}^{\mathrm{win}},
\end{equation}
and $i_0-\mathfrak{f}(i_0)+1\le\max\{3,\,i_0-k_0+2\}$ by part (b) of
Lemma~\ref{4.21:lem-under-block}. It remains to compare the exponent with $\mathfrak{l}_j$. By
\eqref{4.21:eq-weights-incidences-a}, $\mathfrak{l}_j=3(5L^4)^d$ when $j\le k_0+1$, and
$\mathfrak{l}_j=3(5L^{N_0+2})^d$ when $j\ge k_0+2$.

\emph{Blocks of the prepared zones of $C'$.} Here $i_0\le J\le\min\{j+1,k\}$, so the exponent
satisfies
\begin{equation*}
i_0-\mathfrak{f}+1\le\max\{3,\,J-k_0+2\}\le\max\{3,\,j-k_0+3\}.
\end{equation*}
If $j\le k_0+1$, the right side is at most $4$. If $j\ge k_0+2$, it equals $j-k_0+3$. Since
$k_0+N_0=k$, with $N_0$ as in the corresponding display, the inequality
$j-k_0+3\le N_0+2$ is equivalent to $j\le k-1$, and the top zone $j$ of the present operation
satisfies $j\le k-1$.
In both cases \eqref{4.21:eq-window-bounds-1} gives
$\mathfrak{T}(\square_0)\le\mathfrak{l}_j\mathfrak{t}^{\mathrm{win}}$.

\emph{Blocks of the zones made by the $\mathbf{T}$-step, $i_0\ge J+1$, and blocks of zone $k$
off $Z$.} By part (b)(iv) of Lemma~\ref{4.21:lem-under-block}, the probes of these blocks meet
only zones $\ge i_0-1\ge J$.

Case $i_0=J+1$. The cells of zone $J$ that the probe meets have
$b\ge\mathfrak{f}':=\max\{h,\min\{J-1,k_0\}\}$ by part (a) of Lemma~\ref{4.21:lem-under-block}.
The same count, with positions of levels in $[\mathfrak{f}'-1,J]$ in place of
$[\mathfrak{f}-1,i_0-1]$, gives
\begin{equation*}
\mathfrak{T}(\square_0)\le 3\bigl(5L^{J+2-\mathfrak{f}'}\bigr)^d\,\mathfrak{t}^{\mathrm{win}}.
\end{equation*}
Moreover,
\begin{equation*}
J+2-\mathfrak{f}'\le\max\{3,\,J+2-k_0\}\le\max\{3,\,j+3-k_0\}.
\end{equation*}
As in the previous case, this gives $\mathfrak{T}(\square_0)\le\mathfrak{l}_j\mathfrak{t}^{\mathrm{win}}$.

Case $i_0\ge J+2$. The cells the probe meets are unraised, so $b=\operatorname{lev}\ge i_0-1$
on them, and no fine incidence comes from a chain term.

The terms of boundary type of level $i_0-2\le b-1$ pass through at most $(5L^2)^d$ positions.
By the cutoff below the base, and since $k-(i_0-2)\le N$, each position carries at most
$\mathfrak{g}(0)\mathfrak{c}e^{3-\check{R}_k}$. Their contribution is therefore at most
\begin{equation*}
5^dL^d\cdot\bigl[L^d\mathfrak{g}(0)e^{3-\check{R}_k}\bigr]\mathfrak{c}
\le\tfrac14(5L)^d\,\mathfrak{c},
\end{equation*}
by the level condition $(\mathrm{c\text{-}lev})'$.

The terms of boundary type of level $i_0-2-s$ with $s\ge1$, which is $\le b-1-s$, pass through at
most $(5L)^dL^{d(s+1)}$ positions. Each position carries at most
$\mathfrak{g}(s)\mathfrak{c}e^{3-38L^s}$, so their contribution is at most
$(5L)^dC_{\mathrm{fr}}\mathfrak{c}$.

Adding the coarse incidences, and using $\mathfrak{c}\le\mathfrak{t}^{\mathrm{win}}$ and
$C_{\mathrm{fr}}\mathfrak{c}\le\mathfrak{t}^{\mathrm{win}}$, we obtain
\begin{equation*}
\mathfrak{T}(\square_0)\le\bigl[(5L)^d+\tfrac14(5L)^d+(5L)^d\bigr]\mathfrak{t}^{\mathrm{win}}
\le 3(5L)^d\,\mathfrak{t}^{\mathrm{win}}\le\mathfrak{l}_j\,\mathfrak{t}^{\mathrm{win}}.
\end{equation*}

In every case $\mathfrak{T}(\square_0)\le\mathfrak{l}_j\mathfrak{t}^{\mathrm{win}}$. This completes
(b) of \eqref{4.21:eq-window-bounds-a}.

\emph{Part (c).} Let $j\ge k_0+2$, so that $j>k_0$ and $\mathfrak{l}_j=3(5L^{N_0+2})^d$. Then
\begin{equation*}
\bigl(5L^{N_0+2}\bigr)^d=(5L^2)^d\,L^{dN_0},\qquad
L^{dN_0}\le 2^{-d}\bigl(2L^{N_0}\bigr)^d\le 2^{-d}\,\mathsf{l}_j^2.
\end{equation*}
The last inequality is the data-size condition on the constants $\mathsf{l}_j$ for $j>k_0$, read
at $d=4$. Hence
\begin{equation*}
\mathfrak{l}_j\le 3(5L^2)^d\,\mathsf{l}_j^2\le c_{\mathfrak{l}}\,\mathsf{l}_j^2.
\end{equation*}
For $j\le k_0+1$ we have, since $\mathsf{l}_j\ge1$,
\begin{equation*}
\mathfrak{l}_j=3(5L^4)^d=c_{\mathfrak{l}}\le c_{\mathfrak{l}}\,\mathsf{l}_j^2.
\end{equation*}

For the second bound, note first that $L^{N_0-1}=\check{R}_{k_0+1}$ by
the corresponding display. By the monotonicity of the region sizes in
\eqref{4.21:eq-strips-geometry-a}, together with $h+1\le k_0+1$, we get
$\check{R}_{k_0+1}\le\check{R}_{h+1}=\check{R}^{\uparrow}$, the last identity by
\eqref{4.21:eq-weights-incidences-a}. Therefore
\begin{equation*}
L^{N_0+2}=L^3L^{N_0-1}=L^3\check{R}_{k_0+1}\le L^3\check{R}_{h+1}=L^3\check{R}^{\uparrow}.
\end{equation*}
Also $L^4\le L^3\check{R}^{\uparrow}$, since $\check{R}^{\uparrow}\ge L$. Both values of
$\mathfrak{l}_j$ are therefore at most $3(5L^3\check{R}^{\uparrow})^d$. Finally, the comparison of
$\check{R}^{\uparrow}$ with $\mathsf{R}^{\uparrow}(\check{\mathsf{T}}_k)$ in
\eqref{4.21:eq-strips-geometry-a} gives
\begin{equation*}
\mathfrak{l}_j\le 3\bigl(5L^3\check{R}^{\uparrow}\bigr)^d
\le 3\bigl(5L^3\,\mathsf{R}^{\uparrow}(\check{\mathsf{T}}_k)\bigr)^d.
\end{equation*}
\end{proof}

\begin{lemma}[Fine incidences pass through grain-fine cubes]\label{4.21:lem-fine-incidences}
We work at stage $n+1$ of $\mathbf{R}_k$, in the notation of Notation~\ref{4.21:not-stage-grain},
Notation~\ref{4.21:not-weights-incidences} and Lemma~\ref{4.21:lem-under-block}. Let
$\square_0$ be a window block of zone $i_0$, with probe $\tilde{\square}_0$, and let $C'$ be the
class component holding the cells of $\square_0$.
\begin{enumerate}
\item[(a)] Let the incidence $(\square_0,c)$ be fine through the own cube $Q$ of $c$, with
$\ell:=\operatorname{lev}(Q)\le i_0-2$. Then $Q$ is contained in a member of the grain
$\pi^{\flat}$ of stage $n+1$. Thus every fine incidence is through an own cube which is fine with
respect to the grain.
\item[(b)] At window blocks, the only chain terms with fine incidences are the terms of filing
level below $k$ of the local chain of $C'$; terms of filing level $k$ are never fine. Let
$c=\mathbb{C}^{(t)}_k[C'](S)$, a chain term of some stage $t$ of the local chain of $C'$, have a
fine own cube at $\square_0$, where
$i_0\le\min\{j(C')+1,k\}$. Then this cube is a $\pi_m$-cube of $S$, where $m=m(S)$ satisfies
\begin{equation*}
m\le i_0-2\le k-2,
\end{equation*}
and $S$ lies in $Z_{m+1}\cap C'$, which for $m+1\in\,]h,k]$ is one parallelepiped consisting of
at most $(100L\check{R}^{\uparrow})^d$ cubes of $\pi_m$. The cubes of $\pi_m$ in
$Z_{m+1}\cap C'$, $h+1\le m\le k-2$, are called the \emph{fine chain positions} of $C'$. The
local sibling terms of a sibling $C''$ of $C'$ in the sense of
Lemma~\ref{4.21:lem-format-class}, which is itself a class component, have their fine own cubes,
of levels $m_0$ with $h\le m_0\le k-2$, in the boxes $Z_{m_0+1}\cap C''$, and these give at most
$(\check{R}_k+1)(100L\check{R}^{\uparrow})^d$ positions of $C''$ in all. At blocks of zone
$i_0\ge j(C')+2$ there are no fine own cubes of chain terms. A chain term of level $m<k$ has no
own cube in a class component other than that of its band. A chain term of level $m=k$ is not
confined and is never fine.
\end{enumerate}
In short, for $Q$ as in (a) and $c$, $S$, $m$ as in (b),
\begin{equation}\label{4.21:eq-fine-incidences-a}
\begin{gathered}
Q\subset G\quad\text{for some } G\in\pi^{\flat},\\
S\subset\Omega^{c}_{m+1}\cap C'\subset Z_{m+1}\cap C',\\
\#\{\text{fine chain positions in } C'\}\le(\check{R}_k-2)(100L\check{R}^{\uparrow})^d
\le(\check{R}_k+1)(100L\check{R}^{\uparrow})^d .
\end{gathered}
\end{equation}
\end{lemma}

\begin{proof}
\emph{Part (a).} Let $\Delta$ be a cell meeting $Q$, and put $i:=\operatorname{lev}(\Delta)$.
By clause (b) of Lemma~\ref{4.21:lem-under-block}, $i\le i_0+1$. We distinguish the possible
locations of $\Delta$.

Suppose first that $\Delta\subset Z\cap\Omega_{j+1}$ or $\Delta\subset\Lambda_k$. Then $\Delta$
is made by $\mathbf{T}$. Its level satisfies $i\ge j+1$ in the first alternative and $i=k$ in
the second. There are two subcases.
\begin{itemize}
\item If $i_0\le j(C')\le j+1$, then $i\ge i_0>\ell$.
\item If $i_0\ge j(C')+1$, or if the block lies off $Z$, then by clause (b)(iv) of
Lemma~\ref{4.21:lem-under-block} the probe meets cells of levels at least $i_0-1$ only. Hence
\begin{equation*}
i\ge i_0-1\ge\ell+1 .
\end{equation*}
\end{itemize}
In either subcase the level $\ell$ of $Q$ is below the level $i$ of $\Delta$, and $Q$ meets
$\Delta$. Hence $Q\subset\Delta$. The cell $\Delta$ is a member of the grain: in the first
alternative by clause $(\flat c)$ of Notation~\ref{4.21:not-stage-grain} (cells in
$Z\cap\Omega_{j+1}$), in the second by clause $(\flat a)$ there ($\pi_k$-cubes off $Z$).

Suppose next that $\Delta\subset Z\cap\Omega^{c}_{j+1}$. By the description of the grain in
Notation~\ref{4.21:not-stage-grain}, the cells there have level at most $j$, or level $j+1$ on a
band of the current stage of the kind described in that Notation. By clause $(\flat b)$ of that
Notation, the members of the grain there are $\pi_{j+1}$-cubes.

The zones of level at least $j+2$ are made by $\mathbf{T}$ and are thick in every class
component, because $j(C')\le j+1$. Hence, by clause (b)(iv) of Lemma~\ref{4.21:lem-under-block},
a block of zone at least $j+3$ has its probe inside $\Omega_{j+2}\subset\Omega_{j+1}$, and that
probe does not meet $\Delta$. Therefore $i_0\le j+2$, and consequently
\begin{equation*}
\ell\le i_0-2\le j<j+1 .
\end{equation*}
So the cube $Q$ has level below $j+1$, and it lies in one of the $\pi_{j+1}$-cubes that are the
members of the grain there.

In every case $Q$ lies in a member of $\pi^{\flat}$. This is the first line of
\eqref{4.21:eq-fine-incidences-a}.

\emph{Part (b).} We first identify the chain terms concerned. By Lemma~\ref{4.21:lem-window-probe},
in the form \eqref{4.21:eq-window-probe-a}, and by \eqref{4.21:eq-under-block-remarks-a}, the
only chain terms with fine incidences at window blocks are the terms of filing level below $k$
of the local chain of the class component $C'$ holding the block's cells. Terms of filing level
$k$ are never fine.

Next we locate the own cubes. Let $c=\mathbb{C}^{(t)}_k[C'](S)$, a chain term of some stage $t$
of the local chain of $C'$, have a fine own cube at a window
block of zone $i_0\le\min\{j(C')+1,k\}$. This cube is a $\pi_m$-cube of $S$, where $m=m(S)$
satisfies
\begin{equation*}
m\le i_0-2\le k-2 .
\end{equation*}
By \cite[(1.63)]{B-CMP122a}, in the unprimed form in which it is read in clause (d) of
Lemma~\ref{4.21:lem-under-block}, we have $S\subset\Omega^{c}_{m+1}$. By the inclusion of the
complements of the small-field domains in the zones, applicable because $m+1\le k$,
\begin{equation}\label{4.21:eq-fine-incidences-1}
\Omega^{c}_{m+1}\subset Z_{m+1}\subset\Lambda_k^{c}.
\end{equation}
The label $S$ is connected and meets its band in $C'$, and $C'$ is a connected component of
$\Lambda_k^{c}$. Since $S\subset\Lambda_k^{c}$ by \eqref{4.21:eq-fine-incidences-1}, it follows
that $S\subset C'$. Hence
\begin{equation*}
S\subset\Omega^{c}_{m+1}\cap C'\subset Z_{m+1}\cap C',
\end{equation*}
which is the second line of \eqref{4.21:eq-fine-incidences-a}.

We now count. For $m+1\in\,]h,k]$, the nesting statement for the zones of a class component says
that $Z_{m+1}\cap C'$ is one parallelepiped of at most $(100\check{R}_{m+1})^d$ cubes of
$\pi_{m+1}$. This statement applies to siblings as well, since a sibling is itself a class
component. Each cube of $\pi_{m+1}$ is the union of $L^d$ cubes of $\pi_m$, and
$\check{R}_{m+1}\le\check{R}^{\uparrow}$. So the parallelepiped consists of at most
\begin{equation*}
(100L\check{R}_{m+1})^d\le(100L\check{R}^{\uparrow})^d
\end{equation*}
cubes of $\pi_m$. Each level $m$ with $h+1\le m\le k-2$ thus contributes at most
$(100L\check{R}^{\uparrow})^d$ fine chain positions, and there are $k-2-h\le\check{R}_k-2$ such
levels. Hence the fine chain positions in a class component number at most
\begin{equation*}
(\check{R}_k-2)(100L\check{R}^{\uparrow})^d\le(\check{R}_k+1)(100L\check{R}^{\uparrow})^d ,
\end{equation*}
which is the third line of \eqref{4.21:eq-fine-incidences-a}.

We turn to the siblings. Let $C''$ be a sibling of $C'$ in the sense of
Lemma~\ref{4.21:lem-format-class}. Let $e$ be a local sibling term of $C''$, of level $m_0$ and
label $A_e$; by that lemma it is a confined term of the class $\mathfrak{Q}$, and
\begin{equation*}
A_e\subset\Omega^{c}_{m_0+1}\cap C''\subset Z_{m_0+1}\cap C''.
\end{equation*}
The fine own cubes of these terms, of levels $h\le m_0\le k-2$, therefore lie in the same boxes,
namely the parallelepipeds $Z_{m_0+1}\cap C''$ of the preceding count. There are
$k-1-h\le\check{R}_k-1$ such levels, each contributing at most $(100L\check{R}^{\uparrow})^d$
positions; this gives at most $(\check{R}_k+1)(100L\check{R}^{\uparrow})^d$ positions of $C''$
in all.

At blocks of zone $i_0\ge j(C')+2$, which lie in zones made by $\mathbf{T}$, there are no fine
own cubes of chain terms, by clause (b)(iv) of Lemma~\ref{4.21:lem-under-block}.

Finally, consider a chain term with an own cube in a class component other than that of its band.
For $m<k$ this is impossible by the second line of \eqref{4.21:eq-fine-incidences-a}, which places
$S$ inside the component $C'$ of its band. For $m=k$ the term is not confined and is never fine.
It is counted at its coarse incidences, and through the position entering its budget
$\mathsf{e}(c)$ of Notation~\ref{4.21:not-weights-incidences}, in the counting statement of this
section for coarse incidences.
\end{proof}

\begin{lemma}[Fine incidences at window blocks counted once]\label{4.21:lem-fine-count}
Work at stage $n+1$ of the operation $\mathbf{R}_k$, with the terms $c$, their weights $W(c)$,
the class components, the window blocks and the fine incidences $(\square_0,c)$ as in
Notation~\ref{4.21:not-weights-incidences} and Lemma~\ref{4.21:lem-fine-incidences}, and with
the constants $\theta_0$, $C_0^{\sharp}$, $A_{\flat}$,
$b_{\delta}$, the region sizes $\check{R}_k$, $\check{R}^{\uparrow}$, the envelope
$\mathsf{R}^{\uparrow}(\check{\mathsf{T}}_k)$, the function $\mathfrak{g}(s)$ and the numbers
$\mathfrak{c}$, $\mathfrak{c}^{\circ}$ as in Notation~\ref{4.21:not-weights-incidences} and
\eqref{4.21:eq-weights-incidences-a}. Here $N:=\check{R}_k$ is the letter declared at step $k$
for the region size of the step; in this section it is a counting letter of the step, not the rank
of the gauge group. Let $\delta_P$, $d^{\wedge}(\square_0,R_n)$ and
$c_{\mathrm{loc}}$ be as in Lemma~\ref{4.21:lem-grain-block-sum} and
\eqref{4.21:eq-grain-block-sum-a}, and $\mathsf{c}_d$ as in part (g1) of
\eqref{4.21:eq-tree-length-facts-a} in Lemma~\ref{4.21:lem-tree-length-facts}. For a point $x$
and a block, cube or class component $\Diamond$ we define the position weight
$\pi_x(\Diamond)$, and the constants $C_{\mathrm{pos}}$, $c_{\star}''$, $C_{\mathrm{far}}$,
$C_{\mathrm{old}}$ and $\mathfrak{F}^{\mathrm{fi}}$, by the definitions in the display
\eqref{4.21:eq-fine-count-a} below.
The position weight $\pi_x(\Diamond)$ is non-increasing in $\operatorname{dist}_k(x,\Diamond)$.
The numbers $C_{\mathrm{pos}}$, $c_{\star}''$, $C_{\mathrm{far}}$, $C_{\mathrm{old}}$ are finite
and are fixed together with $d$, $L$, the parameter $\lambda$ of this chapter and the
stage-count constant $r_{\mathrm{pr}}$. Then, for every point $x$, the
inequalities (a) and (b) in the last lines of \eqref{4.21:eq-fine-count-a} hold; in (a) the sum
runs over all class components $\bar c$.
\begin{equation}\label{4.21:eq-fine-count-a}
\begin{aligned}
\pi_x(\Diamond) &:= \min\bigl\{1,\ e^{-(\operatorname{dist}_k(x,\Diamond)-204\check{R}_k)/(3N)}\bigr\},\\
C_{\mathrm{pos}} &:= e\Bigl[416^d+\sum_{t\ge 1}(2t+420)^d e^{-t\theta_0}\Bigr],\\
c_{\star}'' &:= 4\sum_{u\ge 0}(2u+7)^d\min\bigl\{1,\ e^{-(u-205)\theta_0}\bigr\},\\
C_{\mathrm{far}} &:= \frac{c_{\star}''}{4}\,\sup_{R\ge 38}\ \sum_{s\ge 1}(RL^{s+1})^d\,
\mathfrak{g}(s)\,e^{2-RL^{s+1}},\\
C_{\mathrm{old}} &:= (110L)^d\bigl(1+\mathfrak{g}(0)\bigr)^{d+1}\\
&\qquad+2^{d-1}\sum_{s\ge 1}\bigl[(104L)^dL^{ds}+6^d\mathfrak{g}(s)^d\bigr]\,
\mathfrak{g}(s)\,e^{3/2-19L^{s}},\\
\mathfrak{F}^{\mathrm{fi}} &:= 2C_{\mathrm{pos}}(N+1)(100L\check{R}^{\uparrow})^d
\bigl(NC_0^{\sharp}+A_{\flat}b_{\delta}\bigr)\\
&\qquad+6e^{4/3}\mathsf{c}_dC_{\mathrm{pos}}\bigl(L\mathsf{R}^{\uparrow}(\check{\mathsf{T}}_k)\bigr)^d\\
&\qquad\qquad\times\Bigl\{(N+1)(103L)^d\mathfrak{c}^{\circ}
+\bigl[(103L)^d+C_{\mathrm{old}}\bigr]\mathfrak{c}\Bigr\}\\
&\qquad+2C_{\mathrm{far}}\,\mathfrak{c},\\
&\text{(a)}\quad \sum_{\bar c}\pi_x(\bar c)\ \le\ \frac{C_{\mathrm{pos}}}{e},\\
&\text{(b)}\quad \sum\Bigl\{W(c)\,e^{-\frac12\delta_P d^{\wedge}(\square_0,R_n)}\,\pi_x(\square_0):\
(\square_0,c)\ \text{a fine incidence},\\
&\qquad\qquad\square_0\ \text{a window block}\Bigr\}
\ \le\ c_{\mathrm{loc}}\,\mathfrak{F}^{\mathrm{fi}}.
\end{aligned}
\end{equation}
\end{lemma}

\begin{proof}[Proof of part (a)]
Fix $x$. We sort the class components $\bar c$ by their distance from $x$. Every $\bar c$ either
satisfies $\operatorname{dist}_k(x,\bar c)<206\check{R}_k$, or has
$\operatorname{dist}_k(x,\bar c)$ in exactly one of the intervals
$[(205+t)\check{R}_k,(206+t)\check{R}_k)$ with $t\ge 1$. These cases therefore exhaust all
class components, and we bound the contribution of each case.

First, the class components with $\operatorname{dist}_k(x,\bar c)<206\check{R}_k$. Each of them
meets one and the same box, namely the box of side $412\check{R}_k$ centred at $x$, and
distinct class components are at mutual distance at least $\check{R}_k$. The packing bound, part
(g6) of \eqref{4.21:eq-tree-length-facts-a} in Lemma~\ref{4.21:lem-tree-length-facts}, applied
once to this box, shows that they number
at most $413^d\le 416^d$. Each of them has $\pi_x(\bar c)\le 1$ by the definition of $\pi_x$ in
\eqref{4.21:eq-fine-count-a}. Their contribution to the sum in (a) is therefore at
most $416^d$. The packing bound must be applied to the whole box at once. The box contains
$206$ shells of width $\check{R}_k$ around $x$; applying the bound to each shell separately gives
$206$ bounds, each at most $416^d$, whose sum is not bounded by $416^d$.

Next, fix $t\ge 1$ and consider the class components with
\begin{equation*}
(205+t)\check{R}_k\ \le\ \operatorname{dist}_k(x,\bar c)\ <\ (206+t)\check{R}_k .
\end{equation*}
All of them meet the box of side $2(206+t)\check{R}_k$ centred at $x$, and they are again at
mutual distance at least $\check{R}_k$. By the packing bound (g6) of
\eqref{4.21:eq-tree-length-facts-a} in Lemma~\ref{4.21:lem-tree-length-facts}, applied to this
box, they number at most $(2t+414)^d$, and
$(2t+414)^d\le(2t+420)^d$. For each of them
$\operatorname{dist}_k(x,\bar c)-204\check{R}_k\ge(t+1)\check{R}_k$. Since $\pi_x$ is
non-increasing in the distance, this gives
\begin{equation*}
\pi_x(\bar c)\ \le\ e^{-(t+1)\check{R}_k/(3N)}\ \le\ e^{-t\theta_0},
\end{equation*}
where the last inequality uses the value of $\theta_0$ fixed in
\eqref{4.21:eq-weights-incidences-a}. The contribution of the $t$-th interval is therefore at
most $(2t+420)^de^{-t\theta_0}$.

Adding the first contribution and the contributions of all $t\ge 1$ gives
\begin{equation*}
\sum_{\bar c}\pi_x(\bar c)\ \le\ 416^d+\sum_{t\ge 1}(2t+420)^de^{-t\theta_0}
\ =\ \frac{C_{\mathrm{pos}}}{e},
\end{equation*}
by the definition of $C_{\mathrm{pos}}$ in \eqref{4.21:eq-fine-count-a}.
\end{proof}

The proof of part (b) is given in a separate lemma later in this section.

\begin{remark}[Degrees of the constants in the logarithmic scale variable]
We record, for use in a later section of this chapter, how the constants in
\eqref{4.21:eq-fine-count-a}, and the numbers $N^{\mathrm{fr}}$, $\mathfrak{t}^{\mathrm{win}}$,
$\mathfrak{l}_j$ of \eqref{4.21:eq-weights-incidences-a}, grow with the
logarithmic scale variable $\mathsf{T}$, for $d=4$; the degree of a quantity is the exponent of
its polynomial growth in $\mathsf{T}$, and $\mathsf{R}^{\uparrow}$ stands for
$\mathsf{R}^{\uparrow}(\check{\mathsf{T}}_k)$. Degrees are expressed through the exponent $r$
that appears in the condition $p_0\ge 5r$ on the degree function $p_0(\cdot)$, and we say that a
quantity has degree $a+0^{+}$ if it has degree
at most $a+\varepsilon$ for every $\varepsilon>0$. The factors have the following degrees in
$\mathsf{T}$:
\begin{itemize}
\item $(L\check{R}^{\uparrow})^d$ and $(L\mathsf{R}^{\uparrow})^d$ have degree $4r+0^{+}$;
\item $N$ has degree $r$;
\item the number $N^{\mathrm{fr}}$, which enters $\mathfrak{c}$, gives $\mathfrak{c}$ degree
$r+0^{+}$, and $\mathfrak{c}^{\circ}$ has degree $0$;
\item $\mathfrak{t}^{\mathrm{win}}$ has degree $r+0^{+}$ and $\mathfrak{l}_j$ has degree $4r+0^{+}$;
\item $c_{\mathrm{loc}}$, $C_{\mathrm{pos}}$, $c_{\star}''$, $C_{\mathrm{far}}$ and
$C_{\mathrm{old}}$ have degree $0$.
\end{itemize}
Consequently the chain part $N(N+1)(100L\check{R}^{\uparrow})^d$ of the first summand of
$\mathfrak{F}^{\mathrm{fi}}$ has degree $6r+0^{+}$, and the leaf part, the second summand, has
degree $5r+0^{+}$. Indeed, in the second summand the factor $N^{\mathrm{fr}}$, entering through
$\mathfrak{c}$, multiplies $(103L)^d+C_{\mathrm{old}}$ and never $(N+1)(103L)^d$; the latter is
multiplied only by $\mathfrak{c}^{\circ}$, which has degree $0$. This is the split already made
in the earlier bound from which the second summand is taken.
\end{remark}

\begin{proof}[Proof of Lemma~\ref{4.21:lem-fine-count}]\label{4.21:prf-fine-count-proof}
We write $\Sigma^{\mathrm{fi}}$ for the left-hand side of the bound of
Lemma~\ref{4.21:lem-fine-count}: the sum of
$W(c)\,e^{-\frac12\delta_P d^{\wedge}(\square_0,R_n)}\,\pi_x(\square_0)$ over the fine incidences
$(\square_0,c)$ in which $\square_0$ is a window block.

\emph{Reduction to own cubes.} By part (F1) of Lemma~\ref{4.21:lem-fine-incidences}, every fine
incidence $(\square_0,c)$ passes through a grain-fine own cube $Q\subset\tilde{\square}_0$ of $c$.
For each incidence we fix one such cube $Q$. The map sending an incidence to the triple
$(Q,c,\square_0)$ is injective, and $\tilde{\square}_0$ meets $Q$. Fix a pair $(Q,c)$ that occurs. The
bound for the blocks whose probe meets a grain member, stated in
\eqref{4.21:eq-grain-block-sum-a} as a consequence of Lemma~\ref{4.21:lem-grain-block-sum}, is
applied with $\beta=\frac12\delta_P$; this is admissible because
$\frac12\delta_P\ge\delta_P/16$. It gives
\begin{equation}\label{4.21:eq-fine-count-proof-1}
\sum_{\square_0:\ \tilde{\square}_0\cap Q\neq\emptyset}
e^{-\frac12\delta_P d^{\wedge}(\square_0,R_n)}\le c_{\mathrm{loc}} .
\end{equation}
For every block $\square_0$ occurring in \eqref{4.21:eq-fine-count-proof-1}, both $\square_0$ and $Q$
lie in the probe $\tilde{\square}_0$. This probe is $\square_0$ with one layer of
$\pi_{i_0}$-cubes added. Its side is $4$ in units of its own level $i_0\le k$, hence at most $4$
level-$k$ units. The general bound on probes gives side at most $5$, and the argument below holds
with either bound. Consequently $\operatorname{dist}_k(x,\square_0)\ge\operatorname{dist}_k(x,Q)-4$.
By the definition of the position weight in \eqref{4.21:eq-fine-count-a}, this gives
\begin{equation*}
\pi_x(\square_0)\le e^{4/(3N)}\,\pi_x(Q)\le 2\,\pi_x(Q).
\end{equation*}
The last inequality holds because $N\ge 3$ gives $e^{4/(3N)}\le e^{4/9}\le 2$. With side $5$ one
uses $e^{5/(3N)}\le e^{5/9}\le 2$ instead. Summing over the triples and using
\eqref{4.21:eq-fine-count-proof-1}, we obtain
\begin{equation}\label{4.21:eq-fine-count-proof-2}
\Sigma^{\mathrm{fi}}\le 2c_{\mathrm{loc}}\sum_{(Q,c)}W(c)\,\pi_x(Q),
\end{equation}
where the sum runs over the pairs $(Q,c)$ that occur.

A term $F$ is called \emph{leaf-type} if it is read through its label $A_F$, carries a filing
level $\ell$ and an anchor cube $Q_F$, and is of one of the kinds whose per-cube constants enter
clauses (pc-1), (pc-3)($\alpha$) and (pc-4) of Lemma~\ref{4.21:lem-per-cube-bounds}; these kinds
include the top frozen chain terms treated below. By Lemma~\ref{4.21:lem-window-probe} and
Remark~\ref{4.21:rem-under-block-remarks}, every term $c$ that occurs is of one of three kinds:
\begin{itemize}
\item a term of the local chain filed at a level below $k$;
\item a local sibling term;
\item a leaf-type term.
\end{itemize}
We split the
sum in \eqref{4.21:eq-fine-count-proof-2} into three parts:
\begin{itemize}
\item $\Sigma_{\mathrm{A}}$, the pairs with $c$ of the first two kinds;
\item $\Sigma_{\mathrm{B1}}$ and $\Sigma_{\mathrm{B2}}$, the pairs with $c$ leaf-type, in the two
cases B1 and B2 described below.
\end{itemize}

\emph{Case A: local chain terms and local sibling terms.} By part (F2) of
Lemma~\ref{4.21:lem-fine-incidences}, $Q$ is a $\pi_m$-cube contained in $Z_{m+1}\cap\bar c$, where
$\bar c$ is the class component holding $c$ and $m\in[h,k-2]$. Fix such a position $Q$. For the
chain terms, clause (pc-2) of Lemma~\ref{4.21:lem-per-cube-bounds} applies, with
$W(c)=\mathfrak{w}(c)$. For the local sibling terms, which are the $\mathfrak{Q}$-terms, clause
(pc-4) of the same lemma applies at the level of $Q$, once the exponential factor in the tree
length $d_m$ that this clause carries is dropped. Together these clauses give
\begin{equation*}
\sum_{c:\ S_c\supseteq Q}W(c)\le NC_0^{\sharp}+A_{\flat}b_{\delta}.
\end{equation*}
There are at most $(N+1)(100L\check{R}^{\uparrow})^d$ positions $Q$ in one component. Since
$Q\subset\bar c$, we have $\pi_x(Q)\le\pi_x(\bar c)$. The sum over class components of their
position weights is controlled by the position bound of \eqref{4.21:eq-fine-count-a}, which gives
the factor $C_{\mathrm{pos}}/e$. Therefore
\begin{equation}\label{4.21:eq-fine-count-proof-3}
\Sigma_{\mathrm{A}}\le\frac{C_{\mathrm{pos}}}{e}\,(N+1)(100L\check{R}^{\uparrow})^d
\bigl(NC_0^{\sharp}+A_{\flat}b_{\delta}\bigr).
\end{equation}

\emph{Leaf-type terms.} Let $F$ be a leaf-type term. Its level satisfies
$\ell\le i_0-2\le k-2$, and $Q\in\pi_\ell$ is an own cube of $F$ consisting of window points. Let
$K_t$, for $t\ge\ell+1$, be the components constructed by induction in
\eqref{4.21:eq-under-block-proof-a}. In particular $K_k$ is a current component of
$Z^{\mathbf{T}}_k=\Lambda_k^c$, either a class component or an untreated one, and $Q_F\subset K_k$;
hence $Q_F\not\subset\Lambda_k$. Case B1 consists of the pairs for which $Q$ lies in the
component $K_k$. Case B2 consists of the remaining pairs.

\emph{Case B1: $Q$ lies in a class component $C'$, in $C'\setminus Z_h$, and $K_k=C'$.} We
enumerate $F$ by its anchor.

\emph{Position weight.} By part (g1) of Lemma~\ref{4.21:lem-tree-length-facts}, the number of own
cubes of $F$ is at most
$N_\ell(A_F)\le\mathsf{c}_d\bigl(d_\ell(A_F)+1\bigr)$. Moreover
$\operatorname{dist}_k(Q,Q_F)\le d_\ell(A_F)L^{\ell-k}+1$. Hence
\begin{equation*}
\pi_x(Q)\le e^{(d_\ell(A_F)+1)/(3N)}\pi_x(Q_F)\le e^{1/3}e^{d_\ell(A_F)/3}\pi_x(Q_F).
\end{equation*}
Since $Q_F\subset K_k=C'$, we also have $\pi_x(Q_F)\le\pi_x(C')$.

\emph{The constants $\mathfrak{c}_\ell$ and $D_\ell$.} For each level $\ell$, let
$\mathfrak{c}_\ell$ be the sum of the constants of clauses (pc-1), (pc-3)($\alpha$) and (pc-4) of
Lemma~\ref{4.21:lem-per-cube-bounds}, taken over those kinds of term that occur at level $\ell$. Let
$D_\ell$ be the least value of $d_\ell(A_F)$ among the terms $F$ counted at level $\ell$. The second
of the cutoff forms of the per-cube bounds of Lemma~\ref{4.21:lem-per-cube-bounds} then gives
\begin{equation}\label{4.21:eq-fine-count-proof-4}
\begin{aligned}
\Sigma_{\mathrm{B1}}
&\le e^{1/3}\mathsf{c}_d\sum_{C'}\pi_x(C')\sum_{\ell}\sum_{\square'}
\sum_{F:\ A_F\ni\square'}W(F)\bigl(d_\ell(A_F)+1\bigr)e^{d_\ell(A_F)/3}
\mathbf{1}[d_\ell(A_F)\ge D_\ell]\\
&\le e^{1/3}\mathsf{c}_d\sum_{C'}\pi_x(C')\sum_{\ell}
(\#\text{anchors at level }\ell)\cdot 3\,\mathfrak{c}_\ell\,e^{-D_\ell/2},
\end{aligned}
\end{equation}
where $\square'$ runs over the anchor cubes at level $\ell$.

\emph{Frozen chain terms sit at level at most $h$.} Let $F$ be a top frozen chain term with
$K_k=C'$. Let $\mathfrak{a}$ be the arrested attempt to whose frozen chain $F$ belongs, $k_a$ its
step and $X_{\mathfrak{a}}$ its region. The level of $F$ is $\ell=k_a$, and we show $k_a\le h$. Its
anchor is a band cube
$Q_F\subset X_{\mathfrak{a}}$. The region $X_{\mathfrak{a}}$ is connected, lies in
$Z_{k_a}\subset\Lambda_k^c$, and meets $K_k=C'$; hence $X_{\mathfrak{a}}\subset C'$. An arrest at a
step of $\,]h,k[\,$ inside the class component $C'$ would be new large field in the last $N$ steps.
This is excluded by condition (ii) of \cite[p.~177]{B-CMP122a}. The statement on the levels of
arrested attempts, applied to $\mathfrak{a}$, then bounds its step by the level of the attempt
$\mathfrak{b}$ of the present step that treats $C'$: writing $h_b$ for that level, $k_a\le h_b$, and
here $h_b=h$. Hence no top frozen chain term occurs
at the levels $h<\ell\le k-2$, and there only the constants collected in $\mathfrak{c}^{\circ}$
enter; at $\ell=h$ the constants are bounded by $\mathfrak{c}$; and at the levels $\ell=h-1-s$
below $h$, where $k-\ell=N+1+s$, they are bounded by $\mathfrak{c}$ times the factor
$\mathfrak{g}(s)$ of \eqref{4.21:eq-weights-incidences-a}. That is,
\begin{equation*}
\mathfrak{c}_\ell\le\mathfrak{c}^{\circ}\quad(h<\ell\le k-2),\qquad
\mathfrak{c}_h\le\mathfrak{c},\qquad
\mathfrak{c}_\ell\le\mathfrak{g}(s)\,\mathfrak{c}\quad(\ell=h-1-s,\ s\ge0).
\end{equation*}

\emph{Levels $\ell\ge h$.} The anchor $Q_F$ lies in $(Z_\ell^{\boxplus})^{(1)}$. Hence it is within
$\check{R}_\ell+2<8L\check{R}_{\ell+1}$ of a component $K'$ of $Z_\ell$. Let $Z'^{\sim 8}_\ell$ be
the enlargement of the domain $Z'_\ell$ by $8$ used in the construction of the domain sequences.
The component of $Z'^{\sim 8}_\ell$ containing $K'$ is connected and lies inside
$Z^{\mathbf{T}}_{\ell+1}$, and the sets $K_t$ increase with $t$. Therefore
\begin{equation*}
K'\subset K_{\ell+1}\subset K_k=C', \qquad\text{so}\qquad K'\subset C'\cap Z_\ell .
\end{equation*}
For $\ell\in[h,k]$, the set $C'\cap Z_\ell$ is one box of side at most $100\check{R}_\ell$. For
these levels this follows from the nesting of the zones inside a class component, together with
the second part of condition (ii) on the domain sequences at the levels $\ell\in[h,k]$. At $\ell=h$
it follows from the uniqueness clause invoked in the proof of part (g) of
Lemma~\ref{4.21:lem-live-region-geometry}. The anchors at level $\ell$ lie in this box enlarged by
the collar of width $\check{R}_\ell+2$ and one cube, so their number is at most
\begin{equation*}
(100\check{R}_\ell+2\check{R}_\ell+6)^d\le(103\check{R}_\ell)^d\le(103L\check{R}^{\uparrow})^d .
\end{equation*}
The last step uses $\check{R}_\ell\le\check{R}_h\le L\check{R}_{h+1}$. Here $D_\ell\ge 0$, and the
levels $h<\ell\le k-2$ are at most $N+1$ in number. The levels $\ell\ge h$ therefore contribute to
the sum over $\ell$ in \eqref{4.21:eq-fine-count-proof-4} at most
\begin{equation*}
(103L\check{R}^{\uparrow})^d\cdot 3\bigl[(N+1)\mathfrak{c}^{\circ}+\mathfrak{c}\bigr].
\end{equation*}

\emph{Levels $\ell\le h-1$: localisation.} Put $s:=h-1-\ell\ge 0$. Then
\begin{equation*}
Q_F\subset K_h\subset Z_h\cap C'\subset\mathsf{C}_{C'}.
\end{equation*}
By part (g) of
Lemma~\ref{4.21:lem-live-region-geometry}, $\mathsf{C}_{C'}$ is a box of side at most
$104L\check{R}_{h+1}$ level-$h$ units. This is at most $104L\check{R}^{\uparrow}L^{s+1}$ level-$\ell$
units. Adding the collar of width $\check{R}_\ell+2$, the number of anchors is at most
\begin{equation*}
(104L^2\check{R}^{\uparrow}L^{s}+2\check{R}_\ell+4)^d .
\end{equation*}
By Notation~\ref{4.21:not-weights-incidences},
$\check{R}_\ell\le\mathfrak{g}(s)\mathsf{R}^{\uparrow}(\check{\mathsf{T}}_k)$. In the rest of this
proof we write $\mathsf{R}^{\uparrow}$ for $\mathsf{R}^{\uparrow}(\check{\mathsf{T}}_k)$.

\emph{Levels $\ell\le h-1$: damping.} The cube $Q$ consists of window points. By
Lemma~\ref{4.21:lem-window-probe}, its base satisfies $b\ge h\ge\ell+1$. The downward length bound
among the corresponding displays therefore applies. It gives
$D_\ell\ge\check{R}_k-3$ for $s=0$, and $D_\ell\ge 38L^{s}-3$ for $s\ge1$.

\emph{Levels $\ell\le h-1$: the count for $s=0$.} Using $\check{R}^{\uparrow}\le\mathsf{R}^{\uparrow}$
and $2\mathfrak{g}(0)+4\le 6L^2\mathfrak{g}(0)$, the number of anchors is at most
\begin{equation*}
\begin{aligned}
(104L^2\check{R}^{\uparrow}+2\mathfrak{g}(0)\mathsf{R}^{\uparrow}+4)^d
&\le\bigl(110L^2(1+\mathfrak{g}(0))\mathsf{R}^{\uparrow}\bigr)^d\\
&=(110L)^d(1+\mathfrak{g}(0))^d(L\mathsf{R}^{\uparrow})^d .
\end{aligned}
\end{equation*}
Per anchor, the factor $3\mathfrak{c}_\ell e^{-D_\ell/2}$ is at most
$3\mathfrak{g}(0)\mathfrak{c}\,e^{3/2-\check{R}_k/2}\le 3\mathfrak{g}(0)\mathfrak{c}$. The
contribution of $s=0$ is therefore at most
\begin{equation*}
3\mathfrak{c}\,(L\mathsf{R}^{\uparrow})^d(110L)^d(1+\mathfrak{g}(0))^{d+1}.
\end{equation*}

\emph{Levels $\ell\le h-1$: the count for $s\ge1$.} We use $(a+b)^d\le 2^{d-1}(a^d+b^d)$ and
$2\mathfrak{g}(s)+4\le 6\mathfrak{g}(s)$. The number of anchors is then at most
\begin{equation*}
\begin{aligned}
(104L^2\check{R}^{\uparrow}L^{s}+2\mathfrak{g}(s)\mathsf{R}^{\uparrow}+4)^d
&\le 2^{d-1}\bigl[(104L)^dL^{ds}+6^d\mathfrak{g}(s)^d\bigr](L\mathsf{R}^{\uparrow})^d .
\end{aligned}
\end{equation*}
Per anchor, the factor is at most $3\mathfrak{g}(s)\mathfrak{c}\,e^{3/2-19L^{s}}$. By the definition
of $C_{\mathrm{old}}$ in \eqref{4.21:eq-fine-count-a}, the two counts together give
\begin{equation*}
\sum_{s\ge0}(\#\text{anchors at level }h-1-s)\cdot 3\mathfrak{c}_{h-1-s}e^{-D_{h-1-s}/2}
\le 3\mathfrak{c}\,\bigl(L\mathsf{R}^{\uparrow}(\check{\mathsf{T}}_k)\bigr)^dC_{\mathrm{old}} .
\end{equation*}

\emph{Case B1 collected.} We insert both level ranges into \eqref{4.21:eq-fine-count-proof-4}. We use
the position bound of \eqref{4.21:eq-fine-count-a} for $\sum_{C'}\pi_x(C')$, which gives
$C_{\mathrm{pos}}/e$. We also use
$(103L\check{R}^{\uparrow})^d\le(103L)^d(L\mathsf{R}^{\uparrow})^d$. This gives
\begin{equation}\label{4.21:eq-fine-count-proof-5}
\Sigma_{\mathrm{B1}}\le 3e^{1/3}\mathsf{c}_d\,\frac{C_{\mathrm{pos}}}{e}\,
\bigl(L\mathsf{R}^{\uparrow}(\check{\mathsf{T}}_k)\bigr)^d
\Bigl\{(N+1)(103L)^d\mathfrak{c}^{\circ}+\bigl[(103L)^d+C_{\mathrm{old}}\bigr]\mathfrak{c}\Bigr\}.
\end{equation}

\emph{Case B2: $K_k$ is not the component holding $Q$.} Either $Q\subset\Lambda_k$, or $Q\subset C'$
for a class component $C'$ with $K_k\neq C'$. In the second alternative $K_k$ is another class
component (running or sibling) or an untreated one, near $x$ or far from it. The label $A_F$ is a
connected union of $\pi_\ell$-cubes containing $Q$ and $Q_F\subset K_k$. Hence it crosses the gap
between $K_k$ and $Q$. We estimate this gap in the two alternatives.
\begin{itemize}
\item If $Q\subset C'\neq K_k$, the two sets lie in distinct components of $\Lambda_k^c$. By the
separation of such components stated in \eqref{4.21:eq-strips-geometry-a}, their
$\operatorname{dist}_k$ is at least $\check{R}_k$, that is, at least $\check{R}_kL^{k-\ell}$
level-$\ell$ units.
\item If
\begin{equation*}
Q\subset\Lambda_k=\Lambda^{\mathbf{T}}_k\subset\Omega^{\mathbf{T}}_k,
\end{equation*}
then, since
$\ell+1\le k-1$ and $F$ is live, $Q_F\subset K_{k-1}\subset Z_{k-1}$. By part (d) of
Lemma~\ref{4.21:lem-live-region-geometry},
$\operatorname{dist}_{k-1}(Z_{k-1},\Omega^{\mathbf{T}}_k)\ge 8L\check{R}_k$. This is at least
$8\check{R}_kL^{k-\ell}$ level-$\ell$ units.
\end{itemize}
By part (g4) of Lemma~\ref{4.21:lem-tree-length-facts}, in either alternative
\begin{equation*}
d_\ell(A_F)\ge\check{R}_kL^{k-\ell}-2 .
\end{equation*}

\emph{Case B2: counting through $Q$.} We enumerate these pairs by the cube $Q$ itself. Fix an
integer $u\ge 0$. The $\pi_\ell$-cubes $Q$ with
$\operatorname{dist}_k(x,Q)\in[u\check{R}_k,(u+1)\check{R}_k[$ lie in a box of side
$2(u+1)\check{R}_k+3\le(2u+7)\check{R}_k$ level-$k$ units. Their number is at most
$(2u+7)^d(\check{R}_kL^{k-\ell})^d$. On that shell,
$\pi_x(Q)\le\min\{1,e^{-(u-205)\theta_0}\}$.

Now fix $Q$ and put $s:=k-1-\ell\ge 1$ (a depth measured from $k$, different from the $s$ of case
B1). The per-cube bounds of
Lemma~\ref{4.21:lem-per-cube-bounds} hold at every cube of $\pi_\ell$, so they apply at $Q$. We use
the first of their cutoff forms, at $D=\check{R}_kL^{k-\ell}-2$. The constant at level $\ell$ is at
most $\mathfrak{g}(s)\mathfrak{c}$, because
$N(g_\ell)\le\mathfrak{g}((s-N)_+)N^{\mathrm{fr}}\le\mathfrak{g}(s)N^{\mathrm{fr}}$. Hence
\begin{equation*}
\sum_{F\text{ in case B2}:\ A_F\ni Q}W(F)\le\mathfrak{g}(s)\,\mathfrak{c}\,e^{2-\check{R}_kL^{s+1}}.
\end{equation*}
Note that $k-\ell=s+1$. The sum over $u$ is at most $c_{\star}''/4$, and $c_{\star}''/4$ times the
supremum over $R\ge38$ displayed below is $C_{\mathrm{far}}$; both constants are defined in
\eqref{4.21:eq-fine-count-a}. We also use
$\check{R}_k\ge 38$, stated in \eqref{4.21:eq-strips-geometry-a}. Therefore
\begin{equation}\label{4.21:eq-fine-count-proof-6}
\begin{aligned}
\Sigma_{\mathrm{B2}}
&\le\mathfrak{c}\Bigl[\sum_{u\ge0}(2u+7)^d\min\{1,e^{-(u-205)\theta_0}\}\Bigr]
\sum_{s\ge1}(\check{R}_kL^{s+1})^d\,\mathfrak{g}(s)\,e^{2-\check{R}_kL^{s+1}}\\
&\le\mathfrak{c}\cdot\frac{c_{\star}''}{4}\cdot
\sup_{R\ge38}\sum_{s\ge1}(RL^{s+1})^d\,\mathfrak{g}(s)\,e^{2-RL^{s+1}}
=C_{\mathrm{far}}\,\mathfrak{c}.
\end{aligned}
\end{equation}
This is a pure number times $\mathfrak{c}$. No volume factor enters, and no power of $N$ beyond the
factor $N^{\mathrm{fr}}$ contained in $\mathfrak{c}$.

\emph{Collecting.} By \eqref{4.21:eq-fine-count-proof-2},
$\Sigma^{\mathrm{fi}}\le 2c_{\mathrm{loc}}(\Sigma_{\mathrm{A}}+\Sigma_{\mathrm{B1}}+\Sigma_{\mathrm{B2}})$.
We insert
\eqref{4.21:eq-fine-count-proof-3}, \eqref{4.21:eq-fine-count-proof-5} and
\eqref{4.21:eq-fine-count-proof-6}, and use $C_{\mathrm{pos}}/e\le C_{\mathrm{pos}}$ and
$e^{1/3}\le e^{4/3}$. The three contributions are then bounded by the three terms of
$\mathfrak{F}^{\mathrm{fi}}$ in \eqref{4.21:eq-fine-count-a}, so
$\Sigma^{\mathrm{fi}}\le c_{\mathrm{loc}}\,\mathfrak{F}^{\mathrm{fi}}$.
\end{proof}

\begin{remark}
In case B2 of the proof of Lemma~\ref{4.21:lem-fine-count}, a leaf-type term that is
anchored elsewhere crosses either a collar made by the $\mathbf{T}$-step or a gap between
components. Pairs whose anchor lies in a far component, which would spoil the count of case B1,
fall under case B2 and are enumerated through the cube $Q$ itself. Consequently the bound of case B1
keeps only the volume factor
$(L\mathsf{R}^{\uparrow}(\check{\mathsf{T}}_k))^d$, and the bound of case B2 is a pure number times
$\mathfrak{c}$.
\end{remark}

\begin{lemma}[No finer terms at an old block, and its depth]\label{4.21:lem-old-block-depth}
Work at stage $n+1$ of the operation $\mathbf{R}_k$. We use the old blocks $\mathfrak{O}^{+}$, the
class components (running or sibling) and the untreated components of $\Lambda_k^c$ of
Definition~\ref{4.21:def-old-window-blocks}, the current collar region $R_n$, the scaled distance
$d^{\wedge}$ with the stratum unit $w=M/M_1$, and the per-cube constant $\mathsf{b}^{P}$ of
Corollary~\ref{4.21:cor-per-cube-constant}, see \eqref{4.21:eq-per-cube-constant-a}. Let
$\square_0\in\mathfrak{O}^{+}$, let $n_0:=\operatorname{zone}(\square_0)$ be the zone index of
$\square_0$, and let $\mathsf{Q}:=\tilde{\square}_0$ be its probe, a probe of level $n_0$. Recall
that $k=h+N$, with $h$ and $N$ as fixed for the operation $\mathbf{R}_k$.
We say that $\square_0$ is \emph{of treated kind} if $\mathsf{Q}$ meets a cell of
$Z_h\cap C'$ for a class component $C'$, or the strip of a live pocket of a class component.
\begin{enumerate}
\item[(a)] No chain term of the present operation $\mathbf{R}_k$ and no local sibling term has an
own cube meeting $\mathsf{Q}$ that is finer than $\mathsf{Q}$, that is, of level less than $n_0$.
In the bound of Corollary~\ref{4.21:cor-per-cube-constant}, the leaf-type terms and the confined
old terms, those singled out by the indicator in \eqref{4.21:eq-per-cube-constant-a}, carry the
factor $e^{d_\ell(A_c)}$; that is, the block-incidence hypothesis of
the anchoring lemma for old blocks holds on $\mathfrak{O}^{+}$ with the parameter $\varsigma_0$ of
that lemma equal to $1$.
\item[(b)] The quantities $\varrho(\square_0)$ and $N_*$ are defined in the first line of the
display below. The second line holds if $\square_0$ is of treated kind; the third line holds if
$\mathsf{Q}$ meets an untreated component of $\Lambda_k^c$; the fifth line holds if
$\mathsf{Q}\cap Y^{\sharp}\neq\emptyset$ for a live pocket $(Y,m)$; the fourth and the sixth line
hold in every case:
\begin{equation}\label{4.21:eq-old-block-depth-a}
\begin{aligned}
&\varrho(\square_0):=1+\frac{d^{\wedge}(\square_0,R_n)}{w},\qquad N_{*}:=N+2,\\
&d^{\wedge}(\square_0,R_n)\ge w\Bigl[(h-n_0-2)_{+}+\tfrac52\check{R}_h-5\Bigr],\\
&d^{\wedge}(\square_0,R_n)\ge w\bigl[(k-n_0-2)_{+}+\check{R}_k-5\bigr],\\
&k-n_0+1\le N_{*}+\varrho(\square_0),\\
&\check{R}_m\le 2\varrho(\square_0)+2,\\
&\varrho(\square_0)\ge\check{R}_k-4,\qquad d^{\wedge}(\square_0,R_n)\ge w(\check{R}_k-5).
\end{aligned}
\end{equation}
\end{enumerate}
\end{lemma}

\begin{proof}
\emph{Part} (a). Every chain term of the present operation $\mathbf{R}_k$ and every local sibling
term has all its levels at least $h$. For the chain terms this holds because their filing levels
$j_F$ are at least $h+1$. For a local sibling term, write $\mathfrak{Q}$ for its set of stored
terms, and $\hat q$ and $j_W$ for the cube and the level that the format of local sibling terms
attaches to it. Every element $e$ of $\mathfrak{Q}$ has original level satisfying
\begin{equation*}
m_0(e)\ge\operatorname{lev}(\hat q)\ge j_W\ge h ;
\end{equation*}
the first two inequalities are part of the format of local sibling terms, and the last is the
lower bound $j_W\ge h$ satisfied by every local sibling.
On the other hand, $n_0\le h$ in two cases. If $\mathsf{Q}$ meets a cell of $Z_h\cap C'$, this is
the second clause of Lemma~\ref{4.21:lem-under-block}(b). If $\mathsf{Q}$ meets the strip of a
pocket $(Y,m)$ of a class component, then clauses (i1) and (g) of
Lemma~\ref{4.21:lem-live-region-geometry} give $n_0\le m\le h$. In both cases every own cube of the
terms considered has level at least $h\ge n_0$, and none is finer than $\mathsf{Q}$.

Suppose next that $\mathsf{Q}$ meets an untreated component of $\Lambda_k^c$. Then $\mathsf{Q}$ lies
at distance at least $\check{R}_k-5>0$ from $Z$, and the own cubes of every term of the kinds
considered filed at a level below $k$ lie in $Z$; both facts are read off the display
\eqref{4.21:eq-fine-incidences-a} of fact (F2) in Lemma~\ref{4.21:lem-fine-incidences}. Hence no
such term meets $\mathsf{Q}$. The terms filed at level $k$ have level $k\ge n_0$, so they
are not finer than $\mathsf{Q}$ either.

Finally, in Corollary~\ref{4.21:cor-per-cube-constant} the leaf-type terms and the confined old
terms, those singled out by the indicator in \eqref{4.21:eq-per-cube-constant-a}, enter with the
factor $e^{d_\ell(A_c)}$. This is the block-incidence hypothesis of
the anchoring lemma for old blocks with $\varsigma_0=1$, and part (a) is proved.

\emph{Part} (b). The first line of \eqref{4.21:eq-old-block-depth-a} is the definition of
$\varrho(\square_0)$ and $N_*$.

For the second line, note first that by clause (i3) of Lemma~\ref{4.21:lem-live-region-geometry},
every point of $\square_0$ lies in a zone of index at most $n_0+1$.

Let $\square_0$ be of treated kind. By the argument of part (a), $n_0\le h$. If $\square_0$ is not
contained in $Z_h$, then its own cell, which has level $n_0$, lies in $\Lambda_h\cap C'$. This
forces $n_0=h$ and places $\square_0$ within two level-$h$ units of $Z_h$.

The region $R_n$ lies in zones of index at least $h$, outside $Z''_{h+1}$, in a running
component. Since every point of $\square_0$ lies in a zone of index at most $n_0+1$, a contour
from $\square_0$ to $R_n$ therefore crosses the whole zones $n_0+2,\dots,h-1$, which lie in
$\Omega_h^c\subset Z_h$; their number is $(h-n_0-2)_{+}$.

After its last point in $Z_h$ (or from $\square_0$ itself if $\square_0\not\subset Z_h$), the
contour runs through cells of level at most $h$ inside $Z''_{h+1}$. There it covers
\begin{equation*}
\text{at least}\quad 3L\check{R}_{h+1}-4\ \ge\ 3\check{R}_h-4\quad\text{level-$h$ units.}
\end{equation*}
This length comes from two facts. By clause (h) of Lemma~\ref{4.21:lem-live-region-geometry},
\begin{equation*}
\operatorname{dist}_h\bigl(Z_h\cap C',(Z''_{h+1})^c\bigr)\ge 3\check{R}_h
\end{equation*}
in any class component. Moreover, the strips reach at most $\tfrac12\check{R}_h$ beyond $Z_h$.

The crossing of the zones $n_0+2,\dots,h-1$ and this run through $Z''_{h+1}$ give the second line
of \eqref{4.21:eq-old-block-depth-a} by \eqref{4.21:eq-cell-block-sums-a}; the coefficient
$\tfrac52$ there is $3-\tfrac12$, the distance of clause (h) less the reach of the strips.

For the third line, let $\mathsf{Q}$ meet an untreated component. Then $\square_0$ belongs also to
the set $\mathfrak{O}$ of old blocks of the anchoring lemma, since that set and
$\mathfrak{O}^{+}$ differ only in their second clause
(Definition~\ref{4.21:def-old-window-blocks}); the third line of
\eqref{4.21:eq-old-block-depth-a} is the depth bound of that lemma.

For the fourth line, consider first the treated kind. Since $k=h+N$,
\begin{equation*}
k-n_0+1=N+(h-n_0)+1 .
\end{equation*}
By the second line of \eqref{4.21:eq-old-block-depth-a},
\begin{equation*}
\varrho(\square_0)\ge 1+(h-n_0-2)_{+}+\tfrac52\check{R}_h-5 .
\end{equation*}
Using $(h-n_0-2)_{+}\ge h-n_0-2$ and $\check{R}_h\ge 2$ (only $\tfrac52\check{R}_h\ge 5$ is
needed; it follows from the lower bound $38$ on the radii $\check{R}_i$, fixed with the radii
themselves), we get
\begin{equation*}
N_*+\varrho(\square_0)\ge N+(h-n_0)+\tfrac52\check{R}_h-4\ge N+(h-n_0)+1,
\end{equation*}
and this gives $k-n_0+1\le N_*+\varrho(\square_0)$.

In the untreated case,
\begin{equation*}
k-n_0+1=3+(k-n_0-2)\le 3+(k-n_0-2)_{+} .
\end{equation*}
By the third line of \eqref{4.21:eq-old-block-depth-a},
\begin{equation*}
\varrho(\square_0)\ge 1+(k-n_0-2)_{+}+\check{R}_k-5 ,
\end{equation*}
so that
\begin{equation*}
N_*+\varrho(\square_0)\ge 3+(k-n_0-2)_{+}+N+\check{R}_k-5\ge 3+(k-n_0-2)_{+},
\end{equation*}
the last inequality because $N=k-h\ge 0$ and $\check{R}_k-5>0$ (part (a)). This gives
$3+(k-n_0-2)_{+}\le N_*+\varrho(\square_0)$.

For the fifth line, let $\mathsf{Q}$ meet the strip $Y^{\sharp}$ of a live pocket $(Y,m)$. The
region $R_n$ lies outside the $\tfrac52\check{R}_m$-neighbourhood of $Y^{\sharp}$, while every
point of $\square_0$ lies within $5$ level-$m$ units of $Y^{\sharp}$. Hence, by
\eqref{4.21:eq-cell-block-sums-a} as for the second line,
\begin{equation*}
d^{\wedge}(\square_0,R_n)\ge w\bigl(\tfrac52\check{R}_m-5\bigr),
\qquad\text{so}\qquad
\varrho(\square_0)\ge\tfrac52\check{R}_m-4 .
\end{equation*}
Consequently $\check{R}_m\le\tfrac25\bigl(\varrho(\square_0)+4\bigr)$. Since
$\varrho(\square_0)\ge 1$,
\begin{equation*}
\tfrac25\bigl(\varrho(\square_0)+4\bigr)\le 2\varrho(\square_0)+2 ,
\end{equation*}
and therefore $\check{R}_m\le 2\varrho(\square_0)+2$. This is the computation of the anchoring
lemma for the set $\mathfrak{O}$ of old blocks.

For the last line, the anchoring lemma for old blocks gives, in all cases,
$d^{\wedge}(\square_0,R_n)\ge w(\check{R}_k-5)$. In the untreated case this is also contained in
the third line of \eqref{4.21:eq-old-block-depth-a}. By the definition of $\varrho(\square_0)$,
it gives
\begin{equation*}
\varrho(\square_0)\ge 1+(\check{R}_k-5)=\check{R}_k-4 .
\end{equation*}
\end{proof}

\begin{lemma}[The total weight at an old block]\label{4.21:lem-old-block-total}
Assume the floors on the auxiliary parameters under which the anchoring lemma for stored terms is
proved, and the
hypotheses of Lemma~\ref{4.21:lem-per-cube-bounds}. Let $\mathfrak{F}$ be the format class of
Lemma~\ref{4.21:lem-format-class}. For $c\in\mathfrak{F}$ let $W(c)$ be its weight and
$\mathcal{A}_c$ the set of blocks incident to it, as in Notation~\ref{4.21:not-weights-incidences}.
Let $d^{\wedge}(\square_0,R_n)$ be the scaled distance of a block $\square_0$ to the collar
region $R_n$ of the current stage $n+1$ of $\mathbf{R}_k$, and let $\delta_P$ be the decay rate.
Then for every $\square_0\in\mathfrak{O}^{+}$
\begin{equation}\label{4.21:eq-old-block-total-1}
\mathfrak{S}^{\mathrm{all}}(\square_0)
:=\sum\bigl\{W(c):\ c\in\mathfrak{F},\ \square_0\in\mathcal{A}_c\bigr\}
\le C^{P,\mathrm{all}}_{\mathsf{a}}\,e^{\frac18\delta_P d^{\wedge}(\square_0,R_n)} .
\end{equation}
Here
\begin{equation*}
C^{P,\mathrm{all}}_{\mathsf{a}}:=\sup_{\varrho\ge1}\Pi^{\mathrm{all}}(\varrho)\,e^{2-\varrho}<\infty ,
\end{equation*}
and $\Pi^{\mathrm{all}}$ is the function of the depth defined in
\eqref{4.21:eq-old-block-total-a} in the proof. It is a polynomial times a polylogarithm.
The number $C^{P,\mathrm{all}}_{\mathsf{a}}$ is fixed in the order of choice before $\gamma$ and
depends only on
$d,L,\lambda,c_2,r_{\mathrm{pr}},C_N,C_0^{\sharp},B_0,b_{\delta},A_{\flat},K(d)$.
The bound contains no age, no quantity $N_0(k_a)$ attached to an arrested attempt of step $k_a$
(Lemma~\ref{4.21:lem-under-block}(c)), and no number of arrested attempts or of live
second-class regions.
\end{lemma}

\begin{proof}
We fix the following notation.
\begin{itemize}
\item $\mathsf{Q}:=\tilde{\square}_0$ is the probe of $\square_0$.
\item $n_0$ is the level of $\square_0$.
\item $\varrho:=\varrho(\square_0)$ is the depth of $\square_0$, and $N_*$ is the constant of
\eqref{4.21:eq-old-block-depth-a}.
\item $\operatorname{lev}$, $m_0$ and $b$ are the current level, the original level and the base
of a cell, as in Lemma~\ref{4.21:lem-base-below-level}.
\item $h$, $Z$, $Z_h$ and $Z''_{h+1}$ are the level and the regions of the current stage, as in
Lemma~\ref{4.21:lem-under-block}.
\end{itemize}
By Notation~\ref{4.21:not-weights-incidences}, $\square_0\in\mathcal{A}_c$ holds exactly when
$\mathsf{Q}$ meets
\begin{equation*}
X_c=A_c\cup\bigcup_{(Y,m)\in\mathcal{P}_c}Y^{\sharp}.
\end{equation*}
For a term $c$ of type $(\delta)$ attached to the element $e$ one has $X_c=S^+(e)\subset A_e$.
The cubes of the label $A_c$ are
called the own cubes of $c$.

Accordingly, the incidences $(\square_0,c)$ are of three kinds.
\begin{itemize}
\item[(I)] $\mathsf{Q}$ meets an own cube $Q$ of $c$ with $\ell:=\operatorname{lev}(Q)\ge n_0-1$.
\item[(II)] $\mathsf{Q}$ meets an own cube $Q$ of $c$ with $\operatorname{lev}(Q)\le n_0-2$; such
an own cube is fine at an old probe.
\item[(III)] $\mathsf{Q}$ meets $X_c$ only through a strip $Y^{\sharp}$ with
$(Y,m)\in\mathcal{P}_c$. In this case $c$ is wrapped.
\end{itemize}
We bound the three contributions separately. Their sum bounds $\mathfrak{S}^{\mathrm{all}}(\square_0)$.

\emph{Kind} (I). First let $\ell=n_0-1$. Then $Q$ is one of the at most $(5L)^d$ cubes of
$\pi_{n_0-1}$ meeting $\mathsf{Q}$. Next let $\ell\ge n_0$. Then $Q$ is one of the at most $5^d$
cubes of $\pi_\ell$ meeting $\mathsf{Q}$. Both counts follow from the conventions on blocks, probes
and cubes fixed with the anchoring lemma for stored terms.

Fix a level $\ell$ and a cube $\square'\in\pi_\ell$. By Corollary~\ref{4.21:cor-per-cube-constant}
(display \eqref{4.21:eq-per-cube-constant-a}), the terms of $\mathfrak{F}$ of level $\ell$ with own
cube $\square'$ have total weight at most $\mathsf{b}^P(k-\ell)$. This bound does not depend on the
class of the terms. The exponential factors $e^{d_\ell(A_c)}$ that
\eqref{4.21:eq-per-cube-constant-a} attaches to boundary-type and confined old terms are at
least $1$ and are dropped.

Now sum over the levels $\ell\in[n_0-1,k]$. For these levels $k-\ell\le k-n_0+1$, and $\mathsf{b}^P$
is non-decreasing. The depth relations \eqref{4.21:eq-old-block-depth-a} give
$k-n_0\le N_*+\varrho-1$, hence $k-n_0+2\le N_*+\varrho+1$. Using these facts and the growth law of
the per-cube constant, kind (I) contributes at most
\begin{equation}\label{4.21:eq-old-block-total-2}
\begin{aligned}
\bigl[(5L)^d+5^d(k-n_0+1)\bigr]\mathsf{b}^P(k-n_0+1)
&\le(5L)^d(k-n_0+2)\,\mathsf{b}^P(k-n_0+1)\\
&\le(5L)^d(N_*+\varrho+1)\,C_{\mathsf{b}}(2+\varrho)^{p_{\mathsf{b}}}\,\mathsf{b}^P(N_*)
=:\Pi_{\mathrm{I}}(\varrho).
\end{aligned}
\end{equation}
The first inequality uses $5^d\le(5L)^d$.

\emph{Kind} (II). By Lemma~\ref{4.21:lem-under-block}(b), the cells met by $\mathsf{Q}$ have
level at most $n_0+1$.

Suppose $\mathsf{Q}$ meets a cell of level at most $n_0-2$. That cell is not separated from the cell
of $\square_0$ by a thick zone. Indeed, a zone made by $\mathbf{T}$, or a thick prepared stratum,
stops a reach of two own units; this follows from the widths in
Lemma~\ref{4.21:lem-under-block}(b). Hence $\mathsf{Q}$ crosses only thin strata. By
Lemma~\ref{4.21:lem-under-block}(c) it then lies in the window of a single preparation. The cell
is one of that window's cells: either a raised cell, or a cell of the window's unraised bottom
stratum (raise $0$). What is used below is only the floor on the base $b$.

Four cases exhaust kind (II).

(a) \emph{$\mathsf{Q}$ lies in the window of a live arrested attempt $\mathfrak{a}$ of an earlier
step.} Then, by Lemma~\ref{4.21:lem-under-block}(c), $n_0\le k_a$. Moreover
\begin{equation*}
b\ge\mathfrak{f}_{\mathfrak{a}}:=\min\{n_0-2,\ k_0(k_a)-1\}
\end{equation*}
on everything $\mathsf{Q}$ meets.

(b) \emph{$\mathsf{Q}$ lies in a window of the present $\mathbf{R}_k$,} either a running
preparation's or a sibling's. This is impossible, for the following reason.

The cells raised during this $\mathbf{R}_k$, together with their thin strata, lie in
$Z\setminus Z''_{h+1}$. They are at least $3L\check{R}_{h+1}$ level-$h$ units from $Z_h$ and from
every strip of a component of the class; this is
Lemma~\ref{4.21:lem-live-region-geometry}(h). They also lie off every untreated component. In both
respects they are farther away than any old probe reaches. Indeed, an old block of treated kind
lies within two own units of $Z_h$ or of such a strip, inside $Z''_{h+1}$. There all cells have
level at most $h$ and are not raised by this $\mathbf{R}_k$. An old block of untreated kind lies
within two own units of a component that is not of the class.

(c) \emph{$\mathsf{Q}$ meets a live second-class region $(Y,m)$}, in the sense of
Definition~\ref{4.21:def-strips-geometry}. By \cite[(1.33)]{B-CMP122b} the label on $Y$ is the
single domain $Y^{(m)}$. Hence, by Definition~\ref{4.21:def-strips-geometry},
$\operatorname{lev}=m_0=b=m$ on every cell of $Y$.

Since $\mathsf{Q}$ meets $Y\subset Y^{\sharp}$, Lemma~\ref{4.21:lem-live-region-geometry}(i1)
gives $n_0\le m$. The probe $\mathsf{Q}$ lies within $5$ units of level $n_0$ of $Y$, and so
within at most $5$ units of level $m$. Every cell in that neighbourhood either belongs to $Y$ or
lies beside it in $\Lambda^{\mathbf{T}}_m$, where $\operatorname{lev}=b=m$. Two facts justify this.
\begin{itemize}
\item $\Omega^{\mathbf{T}}_{m+1}$ is at distance at least $8L\check{R}_{m+1}$, by
Lemma~\ref{4.21:lem-live-region-geometry}(d).
\item Consider the later raises made by any attempt $\mathfrak{b}$ on the component of $Y$. For it,
$m\le h_{\mathfrak{b}}$ by Lemma~\ref{4.21:lem-live-region-geometry}(g). These raises lie at least
$3L\check{R}_{h_{\mathfrak{b}}+1}$ away from $Z_{h_{\mathfrak{b}}}\supset Y$, by
Lemma~\ref{4.21:lem-live-region-geometry}(h) and \eqref{4.21:eq-stages-components-a}.
\end{itemize}
Therefore no cell of level at most $n_0-2$ lies under $\mathsf{Q}$. Also
$b=\operatorname{lev}=m\ge n_0$ on everything $\mathsf{Q}$ meets. This is case (d) below.

In fact, no term of $\mathfrak{F}$ has an own cube of level $\ell\le n_0-2$ inside $Y$ at all. We
check this for each type of term.
\begin{itemize}
\item A confined old term, or a term of the present $\mathbf{R}_k$, has own level at least $b=m$
there. This follows from Lemma~\ref{4.21:lem-under-block}(d) and
Lemma~\ref{4.21:lem-old-block-depth}.
\item Take a term of boundary type, in the sense of
Corollary~\ref{4.21:cor-per-cube-constant}, with level $\ell\le n_0-2<m$. It was born at step
$\ell<m$, and its reading set $X_c$ meets $Y$. For terms of type $(\alpha)$ this holds because
$X_c\supseteq A_c$. For terms of type $(\delta)$ it holds because $Y$ is a union of
$\pi_m$-cubes, $\ell<m$, and every own cube of $A_e$ meets $X_c$ by the definition of $A_e$
(see Lemma~\ref{4.21:lem-terms-cease}). Now $Y$ lies in a component carried through the completed
large-field operation $\mathbf{R}_m$. By Lemma~\ref{4.21:lem-terms-cease}, the term ceased there.
\item The arrests inside $Y$ of steps earlier than $m$ were dissolved at $\mathbf{R}_m$, together
with their frozen terms and their terms of type $(\delta)$. This follows from the last sentence of
the definition of arrested attempts, and from the rule that an arrest and its terms are dissolved
together.
\item A term born at a step $m$ or later has level at least $m\ge n_0$. For a term of boundary type
the level is the birth step. A later attempt $\mathfrak{b}$ with component containing $Y$ has
$m\le h_{\mathfrak{b}}$, by Lemma~\ref{4.21:lem-live-region-geometry}(g). Its chain terms are filed
at levels at least $h_{\mathfrak{b}}+1$. Its terms of type $(\delta)$ have
$m_0(e)\ge j_W\ge h_{\mathfrak{b}}$, with $j_W$ the level so denoted in
Lemma~\ref{4.21:lem-old-block-depth}. Such a term is therefore of kind (I).
\end{itemize}
So the count of case (d) on $Y$ is void. Counting it there only over-estimates, which is harmless.

(d) \emph{No cell of level at most $n_0-2$ lies under $\mathsf{Q}$.} This covers unraised thick
zones. It also covers the flat cells of a live second-class region and their surroundings, as
shown in case (c). In this case $b=\operatorname{lev}\ge n_0-1$ on everything $\mathsf{Q}$ meets.
Own cubes of the chains of the present $\mathbf{R}_k$, and of local sibling terms, do not occur in
kind (II), by Lemma~\ref{4.21:lem-old-block-depth}.

\emph{Count in case} (a). Own cubes of confined old terms have level at least $b$ at their
points, hence at least $\mathfrak{f}_{\mathfrak{a}}$; this is Lemma~\ref{4.21:lem-under-block}(d).
Own cubes of terms of boundary type may have any level $\ell$. For
$\ell\le\mathfrak{f}_{\mathfrak{a}}-2$, and so $\ell\le b-2$ at every point met, the lower bound gives
\begin{equation*}
d_\ell\ge 38L^{\mathfrak{f}_{\mathfrak{a}}-1-\ell}-3 .
\end{equation*}
That lower bound holds at any point.

\emph{Positions.} Let $P_{\mathfrak{a}}$ be the set of $\pi_\ell$-cubes meeting $\mathsf{Q}$ with
$\mathfrak{f}_{\mathfrak{a}}-1\le\ell\le n_0-2$. From the definition of
$\mathfrak{f}_{\mathfrak{a}}$,
\begin{equation*}
n_0-\mathfrak{f}_{\mathfrak{a}}+1\le\max\{3,\ n_0-k_0(k_a)+2\}\le N_0(k_a)+2 .
\end{equation*}
Hence
\begin{equation*}
\#P_{\mathfrak{a}}\le 2\bigl(5L^{n_0-\mathfrak{f}_{\mathfrak{a}}+1}\bigr)^d
\le 2(5L^3)^dL^{dN_0(k_a)} .
\end{equation*}
Per position and level, the weight is at most $\mathsf{b}^P(k-\ell)$. Moreover
$k-k_a\le k-n_0\le N_*+\varrho-1$, by Lemma~\ref{4.21:lem-under-block}(c) and
\eqref{4.21:eq-old-block-depth-a}. Therefore
\begin{equation*}
\begin{aligned}
k-\ell\le k-\mathfrak{f}_{\mathfrak{a}}+1
&\le\max\{k-n_0+3,\ (k-k_a)+N_0(k_a)+2\}\\
&\le N_*+\varrho+N_0(k_a)+1 .
\end{aligned}
\end{equation*}
We multiply three bounds: the number of positions, the number
$n_0-\mathfrak{f}_{\mathfrak{a}}\le N_0(k_a)+2$ of levels in the range, and the bound per position
and level. With the growth law of the per-cube constant, the positions contribute at most
\begin{equation}\label{4.21:eq-old-block-total-3}
2(5L^3)^dL^{dN_0(k_a)}\bigl(N_0(k_a)+2\bigr)\,
C_{\mathsf{b}}\bigl(N_0(k_a)+\varrho+2\bigr)^{p_{\mathsf{b}}}\,\mathsf{b}^P(N_*) .
\end{equation}

\emph{Tail.} Consider the terms of boundary type of level
$\ell=\mathfrak{f}_{\mathfrak{a}}-1-s$, with $s\ge1$. Under $\mathsf{Q}$ there are at most
$\bigl(5L^{n_0-\mathfrak{f}_{\mathfrak{a}}+1}\bigr)^dL^{ds}$ cubes of $\pi_\ell$.

By the first cutoff form of Lemma~\ref{4.21:lem-per-cube-bounds}
(display \eqref{4.21:eq-per-cube-bounds-a}), each such cube carries weight at most
$\mathfrak{g}(s'')\,\mathfrak{c}\,e^{3-38L^s}$. Here
\begin{equation*}
s'':=(k-1-\ell-N)_+\le a+s+1,
\qquad a:=(h-n_0)_++N_0(k_a)+1 ,
\end{equation*}
where $N=N(g_k)$ is the number of stages of $\mathbf{R}_k$ (see
\eqref{4.21:eq-strips-geometry-a}); this $N$ is not the rank of the gauge group.
By Notation~\ref{4.21:not-weights-incidences} and \eqref{4.21:eq-weights-incidences-a},
$\mathfrak{g}(s'')\le\mathfrak{g}(a)\,\mathfrak{g}(s)$. Moreover
$\mathfrak{g}(a)\le\mathfrak{g}(\varrho+N_0(k_a)+2)$.

The sum $\sum_{s\ge1}L^{ds}\mathfrak{g}(s)e^{3-38L^s}$ is at most the constant $C_{\mathrm{fr}}$ of
\eqref{4.21:eq-weights-incidences-a}. Hence the tail contributes at most
\begin{equation}\label{4.21:eq-old-block-total-4}
\begin{aligned}
&2(5L^3)^dL^{dN_0(k_a)}\,\mathfrak{g}\bigl(\varrho+N_0(k_a)+2\bigr)\,\mathfrak{c}
\sum_{s\ge1}L^{ds}\mathfrak{g}(s)e^{3-38L^s}\\
&\qquad\le 2(5L^3)^dL^{dN_0(k_a)}\,\mathfrak{g}\bigl(\varrho+N_0(k_a)+2\bigr)\,
C_{\mathrm{fr}}\,\mathfrak{c} .
\end{aligned}
\end{equation}

\emph{Count in case} (d). Here kind (II) reduces to the terms of boundary type of levels
$\ell\le n_0-2$. At every point met, $n_0-2\le b-1$.
\begin{itemize}
\item At level $n_0-2$ there are at most $(5L^2)^d$ positions. Each carries at most
$\mathfrak{g}(\varrho+2)\,\mathfrak{c}\,e^{3-\check{R}_k}\le\mathfrak{g}(\varrho+2)\,\mathfrak{c}$.
\item At level $n_0-2-s$, with $s\ge1$, there are at most $(5L)^dL^{d(s+1)}=(5L^2)^dL^{ds}$
positions. Each carries at most $\mathfrak{g}(\varrho+2)\,\mathfrak{g}(s)\,\mathfrak{c}\,e^{3-38L^s}$.
\end{itemize}
Summing over $s$ with the definition of $C_{\mathrm{fr}}$ as before, case (d) contributes at most
\begin{equation}\label{4.21:eq-old-block-total-5}
(5L^2)^d\,(1+C_{\mathrm{fr}})\,\mathfrak{g}(\varrho+2)\,\mathfrak{c} .
\end{equation}

\emph{Removing $N_0(k_a)$.} We use Lemma~\ref{4.21:lem-under-block}(c), with its bounds on $N_0$
displayed in the corresponding display. We combine them with the relations
$k-k_a\le N_*+\varrho$ and $k-k_0(k_a)-1\le N_*+\varrho+N_0(k_a)$, and with the monotonicity of
$s\mapsto c_s$. This gives
\begin{equation}\label{4.21:eq-old-block-total-6}
\begin{aligned}
L^{dN_0(k_a)}&\le L^{2d}\bigl(\check{\mathsf{T}}_k+1
+c_{N_*+\varrho+\mathfrak{N}_0(\varrho)}\bigr)^{d r_{\mathrm{pr}}},\\
N_0(k_a)&\le\mathfrak{N}_0(\varrho)
:=4+2r_{\mathrm{pr}}\log_L\bigl(\check{\mathsf{T}}_k+3+c_{N_*+\varrho}\bigr).
\end{aligned}
\end{equation}

On an old block $\varrho\ge\check{R}_k-4$, by \eqref{4.21:eq-old-block-depth-a}. The relations in
\eqref{4.21:eq-strips-geometry-a}
between $\check{\mathsf{T}}_k$, $\mathsf{T}_k$, $\check{R}_k$ and the stage count $N$ then give
\begin{equation}\label{4.21:eq-old-block-total-7}
\check{\mathsf{T}}_k^{\,r_{\mathrm{pr}}}\le\check{R}_k\le\varrho+4,
\qquad
\mathsf{T}_k\le\check{\mathsf{T}}_k+1,
\qquad
N\le C_N2^{r_{\mathrm{pr}}}(\varrho+4).
\end{equation}
Insert \eqref{4.21:eq-old-block-total-6} and \eqref{4.21:eq-old-block-total-7} into
\eqref{4.21:eq-old-block-total-3}, \eqref{4.21:eq-old-block-total-4} and
\eqref{4.21:eq-old-block-total-5}. This includes the factor $N^{\mathrm{fr}}$ contained in
$\mathfrak{c}$. Every factor of these bounds then becomes a polynomial in $\varrho$ times a
polylogarithm in $\varrho$. Its coefficients are formed from the constants named in the
statement. Let $\Pi_{\mathrm{II}}(\varrho)$ be the sum of the right-hand sides of
\eqref{4.21:eq-old-block-total-3}, \eqref{4.21:eq-old-block-total-4} and
\eqref{4.21:eq-old-block-total-5} after these substitutions. Whichever of the cases (a) and (d)
occurs, kind (II) contributes at most $\Pi_{\mathrm{II}}(\varrho)$.

\emph{Kind} (III). By Lemma~\ref{4.21:lem-live-region-geometry}(i2), $\mathsf{Q}$ meets at most
$2^d$ strips. By Lemma~\ref{4.21:lem-live-region-geometry}(i1), each of them has $m\ge n_0$.

For one such strip we apply clause (iii) of the anchoring lemma for stored terms. It applies
to the present format class because its proof uses the terms only through their per-cube bound,
which here is $\mathsf{b}^P$ (Corollary~\ref{4.21:cor-per-cube-constant}). With
\eqref{4.21:eq-old-block-depth-a} and the growth law of the per-cube constant it gives
\begin{equation}\label{4.21:eq-old-block-total-8}
\begin{aligned}
\sum_{c:\ \mathcal{P}_c\ni(Y,m)}W(c)
&\le(k-m+1)\bigl(101\check{R}_m+2\bigr)^d\,\mathsf{b}^P(k-m)\\
&\le(N_*+\varrho)\,(202\varrho+204)^d\,C_{\mathsf{b}}(1+\varrho)^{p_{\mathsf{b}}}\,
\mathsf{b}^P(N_*)
=:\Pi_{\mathrm{III}}(\varrho).
\end{aligned}
\end{equation}

\emph{Conclusion.} Adding the three kinds, and multiplying kind (III) by the number $2^d$ of
strips,
\begin{equation*}
\mathfrak{S}^{\mathrm{all}}(\square_0)\le\Pi_{\mathrm{I}}(\varrho)+\Pi_{\mathrm{II}}(\varrho)
+2^d\Pi_{\mathrm{III}}(\varrho).
\end{equation*}
We check that this right-hand side is a polynomial times a polylogarithm in $\varrho\ge1$. The
remaining occurrences of $N$, $\check{\mathsf{T}}_k$ and $\mathsf{T}_k+c_s$ sit inside
\begin{equation*}
\mathsf{b}^P(N_*)=A_{\flat}\bigl[b_0+(N_*+1)\,b_1\,(\mathsf{T}_k+c_{N_*})^{r_{\mathrm{pr}}}\bigr]
\end{equation*}
of \eqref{4.21:eq-per-cube-constant-a}. In addition, $N(g_k)\le C_N\mathsf{T}_k^{r_{\mathrm{pr}}}$
sits inside $b_0$. All of these are bounded as in \eqref{4.21:eq-old-block-total-7}.

Since $e^{2-\varrho'}$ decays faster than any polynomial times polylogarithm grows, the supremum
below is finite. We obtain
\begin{equation*}
\mathfrak{S}^{\mathrm{all}}(\square_0)
\le\Bigl[\sup_{\varrho'\ge1}\Pi^{\mathrm{all}}(\varrho')\,e^{2-\varrho'}\Bigr]e^{\varrho-2}
=C^{P,\mathrm{all}}_{\mathsf{a}}\,e^{\varrho-2},
\end{equation*}
where
\begin{equation}\label{4.21:eq-old-block-total-a}
\begin{gathered}
\Pi^{\mathrm{all}}(\varrho):=\Pi_{\mathrm{I}}(\varrho)+\Pi_{\mathrm{II}}(\varrho)
+2^d\,\Pi_{\mathrm{III}}(\varrho),
\qquad
C^{P,\mathrm{all}}_{\mathsf{a}}:=\sup_{\varrho'\ge1}\Pi^{\mathrm{all}}(\varrho')\,e^{2-\varrho'},\\
e^{\varrho-2}\le e^{2(\varrho-1)}=e^{(2/w)\,d^{\wedge}(\square_0,R_n)}
\le e^{\frac18\delta_P\,d^{\wedge}(\square_0,R_n)} .
\end{gathered}
\end{equation}

We justify the three steps in the second line of \eqref{4.21:eq-old-block-total-a}.
\begin{itemize}
\item The first inequality is $\varrho-2\le2\varrho-2$, valid for $\varrho\ge0$.
\item The equality is the relation $\varrho-1=d^{\wedge}(\square_0,R_n)/w$ of
\eqref{4.21:eq-old-block-depth-a}.
\item The last inequality holds because $\frac18\delta_P w\ge2$, by
\eqref{4.21:eq-cell-block-sums-a}.
\end{itemize}
This proves \eqref{4.21:eq-old-block-total-1}. No age, no $N_0(k_a)$ and no number of arrested
attempts or of live second-class regions occurs in $C^{P,\mathrm{all}}_{\mathsf{a}}$.
\end{proof}

\begin{remark}[Remarks on the old-block total]\label{4.21:rem-old-block-remarks}
We make three observations on Lemma~\ref{4.21:lem-old-block-total} and on the old blocks of
Definition~\ref{4.21:def-old-window-blocks}.
\begin{enumerate}
\item Let $\mathfrak{a}$ be an attempt arrested at step $k_a=h$ inside a class component $C'$,
excluded after all $N(g_h)$ stages. Let $\square$ be a block of zone $h$ in the raised common
domain
\begin{equation*}
\mathsf{D}_{\mathfrak{a}}:=\Omega^{c}_{k_0(h)+1}\setminus Z^{\prime\prime(\mathfrak{a})}_h ,
\end{equation*}
lying over $(4\check{R}_{k_0(h)+1})^d$ terms of the thin stages, each supported on a single cube,
and over the short domain sequences $\mathbb{B}^{(m)}$ with $k_0(h)<m<h$. Then
$\square\in\mathfrak{O}^{+}$,
and its contribution to $\mathfrak{S}^{\mathrm{all}}(\square)$ is bounded by case (II-a) of the
proof of Lemma~\ref{4.21:lem-old-block-total}.
\item The factor $L^{dN_0(k_a)}$ in $\Pi^{\mathrm{all}}$, which counts positions, is paid by the
depth factor, because a term is incident to a block only through the probe of that block.
\item The old-anchor sum $\mathfrak{S}^{\mathfrak{K}}(\square_0)$ of the kept members of live
pockets at an old block $\square_0$ is not part of $\mathfrak{S}^{\mathrm{all}}(\square_0)$.
\end{enumerate}
\end{remark}

\begin{proof}
(1) The probe $\tilde{\square}$ lies in $Z_h\cap C'$. Hence $\tilde{\square}$ meets a cell contained
in $Z_h\cap C'$ for the class component $C'$. By the second clause of
Definition~\ref{4.21:def-old-window-blocks}, $\square$ is therefore an old block.

In the proof of Lemma~\ref{4.21:lem-old-block-total} this block falls under case (II-a). There
$n_0=h=k_a$, and the floor of the window is
\begin{equation*}
\mathfrak{f}_{\mathfrak{a}}=\min\{h-2,\;k_0(h)-1\}.
\end{equation*}
The number of positions is at most $2(5L^3)^dL^{dN_0(h)}$. This number is paid by the factor
$e^{\frac18\delta_P d^{\wedge}}$, since
\begin{equation*}
d^{\wedge}\ge w\bigl(\tfrac52\check{R}_h-5\bigr).
\end{equation*}
Consequently the bound of Lemma~\ref{4.21:lem-old-block-total} holds for this configuration.

In the order of steps of Ba{\l}aban's large-field operation \cite{B-CMP122b} one has
$k_a\le h-1$, so this configuration does not arise there. The proof of
Lemma~\ref{4.21:lem-old-block-total} does not use this fact, and the configuration is covered all
the same.

(2) For a block $\square$, write $\mathsf{s}_\square$ for the decoupling parameter attached to
$\square$. By the probe-locality of the interpolation, a term $c$ with region $X_c$ responds to
$\mathsf{s}_\square$ only if
\begin{equation*}
X_c\cap\tilde{\square}\neq\emptyset .
\end{equation*}
Thus a term is incident to a block only through the probe of that block. It follows that the
count of positions is attached to the probe. Here the probe is an old probe, and the count
$L^{dN_0(k_a)}$ is absorbed by the depth factor $e^{\frac18\delta_P d^{\wedge}(\square_0,R_n)}$ of
Lemma~\ref{4.21:lem-old-block-total}.

(3) The sum over kept members obeys
\begin{equation*}
\mathfrak{S}^{\mathfrak{K}}(\square_0)\le e\,\mathfrak{n}^{\mathrm{kept}}\,
e^{\frac1{16}\delta_P d^{\wedge}}.
\end{equation*}
This is the old-anchor bound for the kept members of live pockets, read with the set of old
blocks taken to be the set $\mathfrak{O}^{+}$ of Definition~\ref{4.21:def-old-window-blocks}. The
sum enters the single and double totals over the kept members, not the sum
$\mathfrak{S}^{\mathrm{all}}(\square_0)$.
\end{proof}

\begin{definition}[Bounds carried from the plain-unit expansion]\label{4.21:not-plain-unit-data}
In this section the setting of the plain-unit expansion is in force, with its standing reading of
the constants. The statements proved in that setting before the block-sum proposition are used as
proved: its base case, the lemma fixing the grain constants $N^{\flat}$, $\hat{b}_t$,
$c_{\flat}$, the strong-set statement, and the five clauses of its lemma at the preparation site. The
letters and bounds below are taken from there; those not listed keep the meaning they have there.

\emph{The grain number.} The number $\mathsf{x}$ of the plain-unit expansion is subject to the
condition
\begin{equation}\label{4.21:eq-plain-unit-data-1}
2e^{2}\,\hat{b}_t^{2}\,N^{\flat}\,\mathsf{x}\le 1 .
\end{equation}

\emph{Weights and incidences.} For a term $c$ the weight $\mathfrak{w}(c)$, the set of incident
blocks $\mathcal{A}_c$, the letters $\mathsf{e}(c)$, $\mathsf{f}(c)$, $\mathsf{p}^{\mathrm{rel}}(c)$
and the constant $c_{\mathsf{p}}:=212^{d}$ are those of
Notation~\ref{4.21:not-weights-incidences}. The set $\mathcal{A}_c$ is the disjoint union
$\mathcal{A}_c=\mathcal{A}^{\mathrm{pl}}_c\sqcup\mathcal{A}^{\mathrm{isl}}_c$ of its plain part
and its island part. The letters $m_{\square}$, $x_{\square}$, the corridor
$\operatorname{Cor}(\square)$ of a block $\square$, the relative domains
$\mathsf{d}^{\mathrm{rel}}(e)$ and $\mathsf{d}^{\mathrm{rel}}(\hat e)$ of an element $e$ and of a
group $\hat e$, and the notion of a confined term are those of the plain-unit expansion. The set
$\mathfrak{O}$ of old blocks used there is replaced throughout by its enlargement
$\mathfrak{O}^{+}$.

\emph{Groups.} A group $\hat e$ of the expansion of a term $c$ carries a set $A$ of blocks, a
set $Q$ of cubes and a family $\hat{\mathsf{y}}$ of grains, and has budget $w(\hat e)$;
$\mathcal{S}_j(c)$ is the strong set of $c$, and $w(c)$ the budget of $c$. The numbers $a_k$,
$c_J$, $c_{\flat}$, $\kappa_A$, $K_Q$, $r_A$, $\varepsilon_S$, $C'_V$, the block weights
$\vartheta(\square)$ and their bound $\vartheta$, and the collar regions $R_n$ and $R(\bar c)$
are those of the plain-unit expansion; $K_A$ is taken in the form given by the covering bound
for $K_A$ in this chapter. We put
\begin{equation}\label{4.21:eq-plain-unit-data-a}
\begin{aligned}
K_2^{\flat}&:=K_A\exp\bigl[(a_k+1+3^{d}c_{\flat})N^{\flat}+1\bigr],\\
V'&:=\sup_{n,\bar c}\sum_{\square}e^{-\frac{1}{16}\delta_P d^{\wedge}(\square,R(\bar c))},\\
y_{\square}&:=K_2^{\flat}\,\vartheta(\square)\,e^{(a_k+1)|\operatorname{Cor}(\square)|},\\
\Sigma_c&:=\bigl(K_2^{\flat}\vartheta+K_Q/r_A\bigr)\,V'\,\mathsf{p}^{\mathrm{rel}}(c),\\
\tilde\nu(\hat e)&\le 2\,\mathfrak{w}(c)\,e^{c_J}\,e^{-(\kappa_A-c_{\flat})|\mathcal{S}_j(c)|}\\
&\qquad\times\prod_{\square\in A}\bigl[e^{-1}y_{\square}\bigr]\prod_{q\in Q}\bigl[K_Q/r_A\bigr]
\,\mathsf{x}^{|\hat{\mathsf{y}}|},\\
y_{\square}&\le K_2^{\flat}\vartheta\,e^{-\frac14\delta_P d^{\wedge}(\square,R_n)}
\qquad(\square\in\mathcal{A}^{\mathrm{isl}}_c),\\
y_{\square}&=K_2^{\flat}\vartheta\,e^{-\frac12\delta_P d^{\wedge}(\square,R_n)}
\qquad(\square\in\mathcal{A}^{\mathrm{pl}}_c),\\
\sum_{\hat e\text{ of }c}\tilde\nu(\hat e)&\le 2\,\mathfrak{w}(c)\,e^{c_J}
\bigl[(e^{\Sigma_c}-1)+\varepsilon_S e^{\Sigma_c}\bigr],\\
\sum_{\substack{\hat e\text{ of }c,\ A\ni\square_0\\ (\square_0,\mathsf{E})}}\tilde\nu(\hat e)
&\le 2\,\mathfrak{w}(c)\,e^{c_J}\,y_{\square_0}\,e^{\Sigma_c}\\
&\qquad\times e^{-(\operatorname{dist}_k(\mathsf{E},\square_0)-103\check{R}_k)_+}.
\end{aligned}
\end{equation}
The supremum defining $V'$ runs over all stages $n$ and all terms $\bar c$, and the sum over all
blocks. In the last line the sum runs over the groups of $c$ whose block set contains
$\square_0$ and which are anchored at the pair $(\square_0,\mathsf{E})$ in the sense of the
plain-unit expansion; the set $\mathsf{E}$ and the distance $\operatorname{dist}_k$ are those of
that expansion. If instead of $\square_0\in A$ one prescribes a cube $q_0\in Q$, the
same sum is at most $2\,\mathfrak{w}(c)\,e^{c_J}(K_Q/r_A)e^{\Sigma_c}$. If one prescribes
$q_0\in\mathcal{S}_j(c)$, it is at most $2\,\mathfrak{w}(c)\,e^{c_J}\varepsilon_S e^{\Sigma_c}$.

\emph{Block sums.} With $\check{R}^{\uparrow}$ the largest zone radius,
\begin{equation}\label{4.21:eq-plain-unit-data-2}
\begin{gathered}
V'\le C'_V\,(104\,L\,\check{R}^{\uparrow})^{d}=:\bar V',\\
\sum_{\square\in\mathcal{A}_c}e^{-\frac12\delta_P d^{\wedge}(\square,R_n)}
e^{(a_k+1)|\operatorname{Cor}(\square)|}\le V'\,\mathsf{p}^{\mathrm{rel}}(c).
\end{gathered}
\end{equation}
The counting device of the block-sum proposition is used in the form stated there. Clause~(d) of
that proposition is not used; in its place we use Lemma~\ref{4.21:lem-old-block-total}.

\emph{Budgets and grain sums.} For every group $\hat e$ of $c$,
\begin{equation}\label{4.21:eq-plain-unit-data-3}
\begin{gathered}
w(\hat e)\le w(c)+(a_k+1)\Bigl[\sum_{\square\in A}\bigl(N^{\flat}+|\operatorname{Cor}(\square)|\bigr)
+|\hat{\mathsf{y}}|\Bigr],\\
\sum_{\hat{\mathsf{y}}}\mathsf{x}^{|\hat{\mathsf{y}}|}\le e^{|A|}.
\end{gathered}
\end{equation}
The bound obtained in the third step of the proof of clause~(e) of the block-sum proposition is
used as stated there.
\end{definition}

\begin{proof}
Condition \eqref{4.21:eq-plain-unit-data-1} is assumed. The first four lines of
\eqref{4.21:eq-plain-unit-data-a} are definitions.
The bound $V'\le\bar V'$ and the weighted block sum in \eqref{4.21:eq-plain-unit-data-2} are the
conclusions retained from clauses~(a), (a$'$), (b) and~(c) of the block-sum proposition of the
plain-unit expansion, and the counting device is the tool stated with that proposition.

The first bound of \eqref{4.21:eq-plain-unit-data-3} is proved in the first step of the proof of
clause~(e) of that proposition. The same step proves the bound on $\tilde\nu(\hat e)$ in
\eqref{4.21:eq-plain-unit-data-a}, together with the bound on $y_{\square}$ for
$\square\in\mathcal{A}^{\mathrm{isl}}_c$ and the identity for
$\square\in\mathcal{A}^{\mathrm{pl}}_c$.

The second bound of \eqref{4.21:eq-plain-unit-data-3} is proved in the second step of that proof.
The same step proves the total bound, the anchored bound at $(\square_0,\mathsf{E})$, and the two
anchored bounds at a prescribed $q_0\in Q$ and at a prescribed $q_0\in\mathcal{S}_j(c)$.

The set $\mathfrak{O}$ is replaced by $\mathfrak{O}^{+}$, a convention adopted in this section, and
clause~(d) is replaced by Lemma~\ref{4.21:lem-old-block-total}, which is proved in this section.
\end{proof}

\begin{definition}[The local data factor at a block]\label{4.21:def-local-data-factor}
Fix the step $k$ and Ba{\l}aban's operation $\mathbf{R}_k$ at that step, a stage $n$ of it with
its current collar region $R_n$, and the index $j$ that labels the pieces of type
$\mathbb{P}^{(j)}$ below.
Let $\square$ be a block of the cover, with enlargements
$\tilde{\square}^{4}\subset\tilde{\square}^{5}$. Consider a term of the expansion of the
correlation theorem whose near cube is the single block, $A=\{\square\}$, and whose bounds are
those of Notation~\ref{4.21:not-plain-unit-data}.
Its near groups, its $\mathbb{P}^{(j)}$-pieces and $W$-pieces, and its far pieces are the three
kinds of pieces into which the decomposition of such a term at its near cube splits it. These pieces
carry the \emph{near-group weight} $\mathfrak{n}(\square)$. The \emph{local data factor}
$\mathsf{l}^{2}(\square)$ is the per-block factor in that weight. Both, together with the
constant $\mathfrak{n}_0$, are given by
\begin{equation}\label{4.21:eq-local-data-factor-a}
\begin{aligned}
\mathfrak{n}(\square) &:= \mathsf{l}^{2}(\square)\,\mathfrak{n}_0\,
\bigl(\vartheta(\square)+\mathfrak{p}_j\,\mathbf{1}[\square\text{ meets }R_n]\bigr),\qquad
\mathfrak{n}_0 := c_\alpha C_R(E_0+1)e^{C_2\kappa_1},\\
\mathsf{l}^{2}(\square) &=: L_{\circ}^{6s(\square)}\,\mathfrak{d}(\square),\qquad
s(\square) := \max\{\mathsf{r}_0(\Delta):\ \Delta \text{ a cell meeting } \tilde{\square}^{5}\}.
\end{aligned}
\end{equation}
Here $\vartheta(\square)$ is the decay factor of the block $\square$ and $\mathfrak{p}_j$ is the
weight of the $\mathbb{P}^{(j)}$- and $W$-pieces, both as they enter the near-group total of
that term.
The constant $\mathfrak{n}_0$ is the constant of the bound on the near groups in the preparation
of the components of $\mathbf{R}_k$; the constants $c_\alpha$, $C_R$, $E_0$, $C_2$ are those of
that bound.
The pieces of type $\mathbb{P}^{(j)}$ and of type $W$ are functions of the
integration variables and sit only at blocks meeting $R_n$. At such blocks the decay factor
takes a constant value, written $\vartheta$: $\vartheta(\square)=\vartheta$. The quantities
entering \eqref{4.21:eq-local-data-factor-a}
are as follows.
\begin{enumerate}
\item[(a)] For a cell $\Delta$, $\operatorname{lev}(\Delta)$ is its current level, and
$\min_{\Delta} m_0$ is the least original level $m_0(x)$ of a point $x\in\Delta$. We put
\begin{equation*}
\mathsf{r}_0(\Delta) := \operatorname{lev}(\Delta)-\min_{\Delta} m_0 .
\end{equation*}
For each cell this is the raise made on it during this operation $\mathbf{R}_k$. The raise
is made either by the running preparation of the cell's own component, or by the
preparation of a sibling, frozen at that sibling's exit (raised cells and arrested attempts
as in Definition~\ref{4.21:def-arrested-attempts}). The near groups are localised in
$\tilde{\square}^{4}$ by the same decomposition, and $\tilde{\square}^{4}\subset\tilde{\square}^{5}$.
Hence every cell that
a near group at $\square$ can read meets $\tilde{\square}^{5}$ and enters the maximum defining
$s(\square)$.
\item[(b)] Let $\mathcal{W}$ be the set of cells that the domain sequence of $\mathbf{R}_k$
inherited, already raised, from the preparation of a live arrested attempt of an earlier step.
The
\emph{inherited-datum factor} $\mathfrak{d}(\square)\ge 1$ is defined as follows. If
$\tilde{\square}^{5}\cap\mathcal{W}=\emptyset$, then $\mathfrak{d}(\square):=1$. Otherwise
$\mathfrak{d}(\square)$ is the factor by which the near-group total at $\square$ exceeds
\begin{equation}\label{4.21:eq-local-data-factor-1}
L_{\circ}^{6s(\square)}\,\mathfrak{n}_0\,
\bigl(\vartheta(\square)+\mathfrak{p}_j\,\mathbf{1}[\square\text{ meets }R_n]\bigr),
\end{equation}
where $\mathfrak{n}_0 = c_\alpha C_R(E_0+1)e^{C_2\kappa_1}$ as in
\eqref{4.21:eq-local-data-factor-a}.
\item[(c)] The fourth of the bounds of Notation~\ref{4.21:not-plain-unit-data}, which compares
two machines $\mathsf{A}$, $\mathsf{B}$ on one
scaffold, uses the same weight. In it the size $\mathfrak{N}(c)$ of a stored term $c$ is replaced
by $\sup|\delta c|$, where $\delta c:=c^{\mathsf{B}}-c^{\mathsf{A}}$. For this bound
$\mathfrak{d}(\square)$ is read as the same factor, now for the difference between the two
machines of the near-group total at $\square$. The comparison quantity is
\begin{equation}\label{4.21:eq-local-data-factor-2}
L_{\circ}^{6s(\square)}\,\mathfrak{n}_0|_1\,
\bigl(\vartheta(\square)+\mathfrak{p}_j\,\mathbf{1}[\square\text{ meets }R_n]\bigr),
\qquad \mathfrak{n}_0|_1 := c_\alpha C_R e^{C_2\kappa_1},
\end{equation}
multiplied by the size constant of the two-machine difference of the near-group total in the
correlation theorem.
Thus $\mathfrak{n}_0|_1$ is $\mathfrak{n}_0$ at unit size, that is, with $E_0+1$ replaced by $1$.
This size constant stands in the place of $E_0+1$; it is a member of the data-size constants
$\mathsf{l}_{n+1}$.
\item[(d)] $L_{\circ}$ is the small blocking constant of Ba{\l}aban's large-field construction
\cite{B-CMP122a}; it is the constant $\mathsf{L}_0$ of the small blocking, chosen with
$\mathsf{L}_0^2\le L/3$ as in \cite{B-CMP122a}, so
that $L_{\circ}^2\le L/3$. Raising this inequality to the third power,
and using $d=4$ and $L\ge 1$, gives
\begin{equation}\label{4.21:eq-local-data-factor-3}
L_{\circ}^{6} \le \Bigl(\frac{L}{3}\Bigr)^{3} = \frac{L^{3}}{27} \le L^{d}.
\end{equation}
\end{enumerate}
\end{definition}

\begin{remark}[What Ba{\l}aban's regularity bound decides]\label{4.21:rem-regularity-scope}
Let
\begin{equation*}
\mathsf{l}^{2}(\square)=L_{\circ}^{6s(\square)}\,\mathfrak{d}(\square),\qquad
s(\square)=\max\{\mathsf{r}_0(\Delta):\ \Delta\ \text{a cell meeting}\ \tilde{\square}^{5}\},
\end{equation*}
be the local data factor at a block $\square$ of
Definition~\ref{4.21:def-local-data-factor}. Here $L_{\circ}$ is the blocking constant of
\cite{B-CMP122a}, which satisfies $2\le L_{\circ}<L/2$; $\mathsf{r}_0(\Delta)$ is the raise
that the operation $\mathbf{R}_k$ of the present step gives to the cell $\Delta$; and
$\mathfrak{d}(\square)$ is the inherited-datum factor. Write $\mathcal{W}$ for the set of cells
that $\mathbb{B}_k$ inherited raised from the preparation of a live arrested attempt of an
earlier step (Definition~\ref{4.21:def-arrested-attempts}).
\begin{enumerate}
\item[(a)] At every cell $\Delta$ of the preparation made by $\mathbf{R}_k$, the regularity
bound \cite[(1.24)]{B-CMP122a} carries the power
\begin{equation*}
L_{\circ}^{2(\mathsf{r}_0(\Delta)-1)_{+}},\qquad (\mathsf{r}_0(\Delta)-1)_{+}\le \mathsf{r}_0(\Delta).
\end{equation*}
The data enter at most cubically, so this preparation contributes at most the factor
$L_{\circ}^{6s(\square)}$ to the data constant at $\square$.
\item[(b)] The bound \cite[(1.24)]{B-CMP122a} says nothing about the cells of $\mathcal{W}$.
Consequently it does not determine $\mathfrak{d}(\square)$ at a block with
$\tilde{\square}^{5}\cap\mathcal{W}\neq\emptyset$.

The estimates of this section assume none of the three candidate bounds for
$\mathfrak{d}(\square)$ displayed below.
They use only the inequality of the inherited-datum hypothesis, and that hypothesis is posed at
every block.
\item[(c)] Let $(Y,m)$ be a live second-class region (Lemma~\ref{4.21:lem-live-region-geometry}).
Then $\mathsf{r}=\mathsf{r}_0=0$ on its cells and $Y\cap\mathcal{W}=\emptyset$, so no cell of
$Y$ contributes to $s(\square)$ or to $\mathfrak{d}(\square)$. A block $\square$ with
$\tilde{\square}^{5}\cap Y\neq\emptyset$ either belongs to $\mathfrak{O}^{+}$
(Definition~\ref{4.21:def-old-window-blocks}) or is exceptional in the sense of clause~(i) of
the lemma on the inherited-datum factor.
\end{enumerate}
\end{remark}

\begin{proof}
\emph{(a) The exponent on the cells of this preparation.}
We use the notation of \cite[(1.24), pp.~181--182]{B-CMP122a}: $k_0$ is the original level in
that display, $j>k_0$ is the level of the integration zone, $\varepsilon_l$ and $\eta$ are the
regularity parameters of that display, and the bounds hold after $n$ integrations. That
display bounds $|U^{(n)}_{k,Z}(\partial p)-1|$ by a numerical factor times a power of
$L_{\circ}$ times $\varepsilon_{l}(L^{k-l}\eta)^{2}$, the power and the index $l$ being as
follows.
\begin{itemize}
\item On the thin strata $Z''_{m_1+1}\setminus Z''_{m_1}$ with $k_0\le m_1\le j-1$, the power is
$L_{\circ}^{2(m_1-k_0-1)_{+}}$, with $l=m_1$.
\item On $\Omega^{c}_{k_0+1}\setminus Z''_{j}$, the power is $L_{\circ}^{2(j-k_0-1)}$, with $l=j$.
\item On the original shells $\Omega_{m_0}\setminus\Omega_{m_0+1}$ with
$k_0+1\le m_0\le j-2$, the power is $L_{\circ}^{2(j-m_0-1)}$, with $l=j$.
\item On $\Omega_{j-1}\setminus\Omega_{j+1}$ there is no power of $L_{\circ}$, and $l=j$.
\end{itemize}
Since the raise of a cell in this preparation is the level to which it is raised minus its
original level (Definition~\ref{4.21:def-arrested-attempts}), the raise that this preparation
gives is as follows: $m_1-k_0$ on a thin stratum; $j-k_0$ on
$\Omega^{c}_{k_0+1}\setminus Z''_{j}$; $j-m_0$ on a shell; and at most one on
$\Omega_{j-1}\setminus\Omega_{j+1}$. Reading each stratum with this raise, half the exponent is
$(\mathsf{r}_0-1)_{+}$ in every case, and $(\mathsf{r}_0-1)_{+}\le\mathsf{r}_0$.

The power of $L_{\circ}$ is thus attached to raised cells only, each through its own raise. For
this reason $s(\square)$ is defined through the raises of the cells that meet
$\tilde{\square}^{5}$. It is not defined through the zone level of the block minus the least
original level $m_0$ of a cell meeting $\tilde{\square}^{5}$ (positive part). Here and below
$C'$ denotes a class component (running or sibling), and the level $h$ and the large-field
domain $Z_h$ are those of Definition~\ref{4.21:def-old-window-blocks}; $h$ is a level, not a
history. That difference would charge
deep unraised cells, for instance the cells of $Z_h\cap C'$ under a block of zone $h$, where
that difference is unbounded. It would also give a factor on unraised
neighbours. With $s(\square)$ as defined, every cell met by $\tilde{\square}^{5}$ is raised by at
most the difference between the level $j(C')$ of Lemma~\ref{4.21:lem-base-below-level} and the
original level $k_0(C')$, the level $k_0$ of \cite[(1.24)]{B-CMP122a} for the preparation of
$C'$; hence $s(\square)\le j(C')-k_0(C')$.

In the proof of the uniform bound on the data constants, the data constant on the raised zone
is $L_{\circ}^{2(j-m_0)}$ times the data constant of level $j$. Here $m_0$ is the original level
of the cell, that is, its level in $\mathbb{B}_k$. The terms are at most cubic in this constant;
with $N_0$ the bound on $j-m_0$ used in that proof, the resulting factor is bounded there by
\begin{equation*}
L_{\circ}^{6N_0}\le L^{4N_0}\le \mathsf{l}_j^{2}.
\end{equation*}
By the halved exponent just computed, the cubic dependence gives at $\square$ at most
$L_{\circ}^{6s(\square)}$.

The uniform bound on the data constants is stated with the block-independent factor
$\mathsf{l}_j^{2}$. Taken at every block of the cover, whose number is comparable to $V'$, this
would give $\mathsf{l}_j^{2}$ to a power of order $V'$, which the estimates do not admit. This is
why the local form $\mathsf{l}^{2}(\square)$, read off from its proof, is used. This proves (a).

\emph{(b) Silence on $\mathcal{W}$.}
The first row of \cite[(1.24)]{B-CMP122a} is stated for $p\in Z''_{h+1}\cap\Omega_h$
\cite[p.~181]{B-CMP122a}, with $\Omega_h$ the domain of level $h$ in the notation of that
display. There Ba{\l}aban also states that ``the only characteristic functions
which remain, and which are connected with these integrations, are
$\chi^{(n)}_k\chi_h(Z_h\cap\Omega_h)$''. Since $\mathcal{W}\subset\Omega_h^{c}$, the cells of
$\mathcal{W}$ lie outside the domain of the display, and the display says nothing about them.

In the uniform bound on the data constants, $m_0$ is a cell's level in $\mathbb{B}_k$; that
bound, including its case $j\le k_0$, concerns the cells of the present preparation and does not
determine the datum on the cells of $\mathcal{W}$.

Ba{\l}aban's construction retains the structures of arrested attempts in later steps
\cite{B-CMP122a}, with the sentence ``\dots and it is easy to generalize the definition of the
spaces in such a case for the next steps''. Their frozen structure
$(\ell_{\mathfrak{a}},w_{\mathfrak{a}},\mathcal{Q}_{\mathfrak{a}})$ carries its
$\mathsf{L}_0$-weight (Definition~\ref{4.21:def-arrested-attempts}).

At a block with $\tilde{\square}^{5}\cap\mathcal{W}\neq\emptyset$ three readings of
$\mathfrak{d}(\square)$ are possible.
\begin{itemize}
\item The first is $\mathfrak{d}\equiv1$.
\item The second is
\begin{equation*}
\mathsf{l}^{2}(\square)=L_{\circ}^{6\hat{s}(\square)},\qquad
\hat{s}(\square)=\max\{\mathsf{r}(\Delta):\ \Delta\ \text{meets}\ \tilde{\square}^{5}\},
\end{equation*}
that is,
\begin{equation*}
\mathfrak{d}(\square)\le\max\{L_{\circ}^{6\mathsf{r}(\Delta)}:\ \Delta\in\mathcal{W}
\ \text{meets}\ \tilde{\square}^{5}\}.
\end{equation*}
\item The third is the bound in clause~(iv) of the lemma on the inherited-datum factor.
\end{itemize}
None of the three is used below. The estimates use the one inequality
they consume, the inherited-datum hypothesis. This proves (b).

\emph{(c) Live second-class regions.}
A live second-class region $(Y,m)$ carries the single domain $Y^{(m)}$ in $\mathbb{B}_k$
\cite[(1.33)]{B-CMP122b}. On its cells $m_0=\operatorname{lev}=b=m$ and
$\mathsf{r}=\mathsf{r}_0=0$, and $Y\cap\mathcal{W}=\emptyset$. Hence a cell of $Y$ contributes
nothing to $s(\square)$ and does not enter $\mathfrak{d}(\square)$.

What persists of the completed preparation of $(Y,m)$ persists as data of the density on $Y$.
These are ``the previous functions $\chi_{k,\Lambda_i}$'' of \cite{B-CMP122b}, together with
$\chi^{c}_m(Y^{\sim-6})$. They enter, as data of the density on $Y$, the near-group total at a
block \eqref{4.21:eq-local-data-factor-a} and the bounds on the kept members of $Y$. The old
functions
$\chi_h(Z_h\cap\Omega_h)\mathbf{T}_h(Z_h)$ of \cite{B-CMP122a} on the unraised cells of
$Z_h\cap C'$ enter in the same way.

Now let $\tilde{\square}^{5}$ meet $Y$. Read the six distance conditions of clause~(i) of the
lemma on the inherited-datum factor with $Y\subset Z_h\cap C'$, respectively $Y\subset U$, in
place of $\mathcal{W}$; the two alternatives are the two cases of that clause, and $U$ is the set
so denoted there (not a bond configuration). They show that $\square$ belongs to
$\mathfrak{O}^{+}$ or is exceptional. For
$m\le h-1$ this applies because, by \eqref{4.21:eq-stages-components-a},
\begin{equation*}
Y\subset Z_{h-1}\subset(\Omega^{\mathbf{T}}_h)^{c}.
\end{equation*}

At zone $h$, for a region created at $m=h$, its own collar $Y\cap\Omega^{\mathbf{T}}_h$ leaves
those distance conditions no margin. There the first clause of the definition of
$\mathfrak{O}^{+}$ is used instead: if $\tilde{\square}$ meets $Y^{\sharp}$, then
$\square\in\mathfrak{O}^{+}$. A window probe misses the strip $Y^{\sharp}$, so the block's
own cell lies outside $Y^{\sharp}$, by the remark on the width of the strip $Y^{\sharp}$ at a
distance of at least
\begin{equation*}
\lfloor\tfrac12(\check{R}_h-1)\rfloor
\end{equation*}
level-$h$ units from $Y$. By Lemma~\ref{4.21:lem-live-region-geometry}(a), with $m=h$,
$\check{R}_h\ge 38$, so this number is at least $18>7$. In every case the inherited-datum
hypothesis is posed at every block.
This proves (c).
\end{proof}

\begin{definition}[The inherited-datum hypothesis]\label{4.21:def-datum-hypothesis}
For a block $\square$ of the cover and for either of the two quantities listed below, let
$\mathfrak{d}(\square)$ denote the inherited-datum factor of that quantity at $\square$. It is
defined in Definition~\ref{4.21:def-local-data-factor} through the factorisation
\eqref{4.21:eq-local-data-factor-a} of the local data factor $\mathsf{l}^{2}(\square)$ as
$L_{\circ}^{6s(\square)}\mathfrak{d}(\square)$. Let $n$ be the stage of the operation
$\mathbf{R}_k$ under consideration, $R_n$ the collar region of that stage,
$d^{\wedge}(\square,R_n)$ the scaled distance from $\square$ to $R_n$, and $\delta_P$ the decay
rate of the exponential profiles in the scaled distance $d^{\wedge}$. The \emph{inherited-datum
hypothesis} is the following statement.

There is a number $C_{\mathrm{dat}}$, fixed before $\gamma$ in the order of choice of the auxiliary
parameters, with
\begin{equation*}
C_{\mathrm{dat}}\ge C_{\mathrm{dat,pr}},
\end{equation*}
where $C_{\mathrm{dat,pr}}$ is the constant of part~(ii) of the lemma on depth-paid shares of the
inherited datum (display~\eqref{4.21:eq-datum-depth-shares-a}), such that the following holds.
The two quantities are the following:
\begin{itemize}
\item the near-group total $v$ at $\square$ of one expansion;
\item the difference $\delta v=v^{\mathsf{B}}-v^{\mathsf{A}}$ of that total between two expansions
$\mathsf{A}$ and $\mathsf{B}$ carried out on one common scaffold.
\end{itemize}
For each of them, at every stage of every $\mathbf{R}_k$ and at every block $\square$ of the cover,
\begin{equation}\label{4.21:eq-datum-hypothesis-a}
\mathfrak{d}(\square)\le 1+C_{\mathrm{dat}}\,e^{\frac18\delta_P\,d^{\wedge}(\square,R_n)}.
\end{equation}

Suppose $\tilde{\square}^{5}$ meets no cell of $\mathcal{W}$, the set of raised cells inherited
from the preparation of a live arrested attempt of an earlier step. Then
Definition~\ref{4.21:def-local-data-factor} gives $\mathfrak{d}(\square)=1$, and the inequality
\eqref{4.21:eq-datum-hypothesis-a} is vacuous at $\square$.
\end{definition}

The bound \eqref{4.21:eq-datum-hypothesis-a} is an input of the correlation theorem, namely of its
bound on the near-group total of the small-field terms. It is carried as an explicit hypothesis of
the main theorem of the present chapter.

Proof outstanding.

\textit{Route.} The lemma on depth-paid shares of the inherited datum does two things toward
\eqref{4.21:eq-datum-hypothesis-a}:
\begin{itemize}
\item Its parts (i)--(iii) work from the geometry of the cover and from the bound on the old
factors stated there. They show that the depth of the block pays every share of
$\mathfrak{d}(\square)$ that comes from a cell $\Delta\in\mathcal{W}$ with
$t(\Delta)\ge n_0-2$. Here $t(\Delta)$ is the step at which $\Delta$ was raised and
$n_0$ is the index of the zone containing $\square$.
\item Its part (iv) derives \eqref{4.21:eq-datum-hypothesis-a} at every block from the additional
hypothesis under which that part is stated.
\end{itemize}

What remains is exactly the residual set of blocks: the blocks whose enlargement
$\tilde{\square}^{5}$, but not whose probe $\tilde{\square}$, reaches a cell $\Delta\in\mathcal{W}$
with $t(\Delta)\le n_0-3$. Such blocks occur in particular in the level-$k$ band left when the
operation $\mathbf{R}_k$ terminates at its last stage; the blocks of that band belong to the window
of the counting lemmas.
At a residual block none of the near-group terms has support meeting the cells of $\mathcal{W}$;
the only object that reaches these cells is the cube $\square_0=\tilde{\square}^{5}$ of the
hypothesis of \cite[Lemma~4]{B-CMP109}.

A proof must therefore establish \eqref{4.21:eq-datum-hypothesis-a} at the residual blocks, the
level-$k$ band above included. It must do so both for the near-group total of one expansion and for
the difference of that total between two expansions on one scaffold.

\begin{corollary}[Near-group block weight under the inherited-datum hypothesis]
\label{4.21:cor-near-group-weight}
Assume the inherited-datum hypothesis of Definition~\ref{4.21:def-datum-hypothesis} (whose proof is
outstanding), in the form
\eqref{4.21:eq-datum-hypothesis-a}, and write $C_{\mathrm{dat}}$ for the constant of that hypothesis.
Then at every stage $n$
of every operation $\mathbf{R}_k$ and at every block $\square$ of the cover,
\begin{equation}\label{4.21:eq-near-group-weight-a}
\mathfrak{n}(\square)\le (1+C_{\mathrm{dat}})\,L_{\circ}^{6s(\square)}\,\mathfrak{n}_0\,
\vartheta\Bigl(1+\frac{\mathfrak{p}_j}{\vartheta}\Bigr)\,
e^{-\frac38\delta_P d^{\wedge}(\square,R_n)} .
\end{equation}
Here $\mathfrak{n}(\square)$, $\mathfrak{n}_0$, $s(\square)$ and $L_{\circ}$ are as in
Definition~\ref{4.21:def-local-data-factor}; $\vartheta$ is the block-weight constant and
$\mathfrak{p}_j$ the
weight of the pieces at the rim; $\delta_P$ is the decay rate of the block weights in the scaled
distance $d^{\wedge}$; and $d^{\wedge}(\square,R_n)$ is the
scaled distance from $\square$ to the collar region $R_n$. The bound holds for the near-group total
of one machine, and alike for the difference of the near-group totals of the two machines
$\mathsf{A}$, $\mathsf{B}$ on one scaffold. In the second case $\mathfrak{n}_0$ is replaced by
$\mathfrak{n}_0|_1$ multiplied by the difference size of the near-group total, which is a member of
the list $\mathsf{l}_{n+1}$ of data-size constants.
\end{corollary}

\begin{proof}
Fix a stage $n$ of an operation $\mathbf{R}_k$ and a block $\square$ of the cover. Write
$d^{\wedge}=d^{\wedge}(\square,R_n)$, and note that $d^{\wedge}\ge 0$.

First we bound the local data factor. By Definition~\ref{4.21:def-local-data-factor}, in particular
\eqref{4.21:eq-local-data-factor-a}, we have
$\mathsf{l}^2(\square)=L_{\circ}^{6s(\square)}\,\mathfrak{d}(\square)$. The inherited-datum hypothesis
of Definition~\ref{4.21:def-datum-hypothesis}, whose proof is outstanding, gives, in the form
\eqref{4.21:eq-datum-hypothesis-a},
$\mathfrak{d}(\square)\le 1+C_{\mathrm{dat}}\,e^{\frac18\delta_P d^{\wedge}}$. It gives this at this
block and this stage, for the near-group total of one machine and for the difference on one scaffold
alike. For every real $C$ and every $a\ge0$ one has
\begin{equation*}
(1+C)e^{a}-(1+Ce^{a})=e^{a}-1\ge0 .
\end{equation*}
We apply this with $C=C_{\mathrm{dat}}$ and $a=\frac18\delta_P d^{\wedge}$, and obtain
\begin{equation}\label{4.21:eq-near-group-weight-1}
\mathsf{l}^2(\square)\le L_{\circ}^{6s(\square)}\,(1+C_{\mathrm{dat}})\,
e^{\frac18\delta_P d^{\wedge}} .
\end{equation}

Next we bound the two kinds of contribution to $\mathfrak{n}(\square)$. Recall that
$\mathfrak{n}(\square)$ and $\mathfrak{n}_0$ are as in Definition~\ref{4.21:def-local-data-factor}
and \eqref{4.21:eq-local-data-factor-a}. The quantity $\mathfrak{n}(\square)$ is
$\mathsf{l}^2(\square)\,\mathfrak{n}_0$ times the sum of the weights of its pieces, and a piece
carries either the weight of the block, $\vartheta(\square)=\vartheta\,e^{-\frac12\delta_P d^{\wedge}}$,
or the rim weight $\mathfrak{p}_j$; the pieces of the second kind occur only at blocks with
$d^{\wedge}=0$. By
\eqref{4.21:eq-near-group-weight-1} the pieces carrying the block weight contribute at most
\begin{equation*}
(1+C_{\mathrm{dat}})\,L_{\circ}^{6s(\square)}\,\mathfrak{n}_0\,\vartheta\,
e^{\frac18\delta_P d^{\wedge}-\frac12\delta_P d^{\wedge}}
=(1+C_{\mathrm{dat}})\,L_{\circ}^{6s(\square)}\,\mathfrak{n}_0\,\vartheta\,
e^{-\frac38\delta_P d^{\wedge}} .
\end{equation*}
At a block with $d^{\wedge}=0$ both exponentials equal $1$. By
\eqref{4.21:eq-near-group-weight-1} the pieces carrying the weight $\mathfrak{p}_j$ therefore
contribute at most
\begin{equation*}
(1+C_{\mathrm{dat}})\,L_{\circ}^{6s(\square)}\,\mathfrak{n}_0\,\mathfrak{p}_j
=(1+C_{\mathrm{dat}})\,L_{\circ}^{6s(\square)}\,\mathfrak{n}_0\,\vartheta\,
\frac{\mathfrak{p}_j}{\vartheta}\,e^{-\frac38\delta_P d^{\wedge}} .
\end{equation*}
At blocks with $d^{\wedge}>0$ they contribute nothing. Adding the two contributions gives
\eqref{4.21:eq-near-group-weight-a}.

For the difference of the two machines' totals on one scaffold, the same derivation applies with
$\mathfrak{n}_0$ replaced throughout by $\mathfrak{n}_0|_1$ times the difference size of the
near-group total, a member of the list $\mathsf{l}_{n+1}$. Inequality
\eqref{4.21:eq-near-group-weight-1} is unchanged,
because the hypothesis \eqref{4.21:eq-datum-hypothesis-a} is assumed for the difference as well.
\end{proof}

\begin{lemma}[The raise read by a block]\label{4.21:lem-block-raise}
Consider the operation $\mathbf{R}_k$, with the class components $C'$ (the running one and its
siblings), the base level $h$, the level $k_0$, the top zones $j(C')$, the level $j$, and the
regions $Z$, $Z_h$, $Z'_h$, $Z''_{h+1}$ and $\Lambda_k$ as in
Definition~\ref{4.21:def-old-window-blocks} and Lemma~\ref{4.21:lem-under-block}. Let
$\mathfrak{O}^{+}$ be the set of old blocks of Definition~\ref{4.21:def-old-window-blocks}. For a
block $\square$ let $s(\square)$ be the raise read by $\square$ in
Definition~\ref{4.21:def-local-data-factor}, that is, the largest of the raises
$\mathsf{r}_0(\Delta)$ made during this $\mathbf{R}_k$ over the cells $\Delta$ meeting
$\tilde{\square}^5$. Write $\operatorname{zone}(\square)$ for the zone of $\square$. Then the
following hold.
\begin{enumerate}
\item[(i)] $s(\square)=0$ for every $\square\in\mathfrak{O}^{+}$, and for every block $\square$ such
that $\tilde{\square}^5$ meets no cell raised during this $\mathbf{R}_k$.
\item[(ii)] $\mathsf{r}_0\le 1$ on raised cells of level at most $k_0+1$, and
$\mathsf{r}_0\le i-k_0$ on raised cells of level $i\ge k_0+1$.
\item[(iii)] $s(\square)\le 1$ if $\tilde{\square}^5$ meets only cells of levels at most $k_0+1$, and
$s(\square)\le 2$ if $\operatorname{zone}(\square)\le k_0+1$.
\item[(iv)] Let $N_0$ be as in \eqref{4.21:eq-strips-geometry-a}. If $\tilde{\square}^5$ meets a
cell raised during this $\mathbf{R}_k$, let $C'$ be the class component containing the raised
cells that $\tilde{\square}^5$ meets. Then
\begin{equation}\label{4.21:eq-block-raise-a}
\begin{gathered}
s(\square)\le j(C')-k_0\le j+1-k_0\le N_0
\quad\text{if $\tilde{\square}^5$ meets a raised cell of $C'$},\\
L_{\circ}^{6s(\square)}\le L^{d}\,\mathsf{l}_j^{2}\quad\text{for every block $\square$}.
\end{gathered}
\end{equation}
\end{enumerate}
\end{lemma}

\begin{proof}
Fix a block $\square$ and put $i_0=\operatorname{zone}(\square)$. We call the units of level
$i_0$ the block's own units, and the cell of level $i_0$ carrying $\square$ its own cell.

\emph{Levels met by $\tilde{\square}^5$.} A point of $\tilde{\square}^5$ is within $7$ own units
of the block's own cell. The own cell lies in $\Omega^c_{i_0+1}$ in every case, by the statement
that fixes the zone of a block. This includes the case of a cell of a live second-class region
$(Y,m)$: by Lemma~\ref{4.21:lem-under-block}(a) such a cell has level $m$, so that $i_0=m$, and
Lemma~\ref{4.21:lem-live-region-geometry}(b) gives $Y^{\sharp}\subset\Omega^c_{m+1}=
\Omega^c_{i_0+1}$ and says that a cell meeting $Y^{\sharp}$ has level at most $m$ and lies in
$Y^{\sharp}$. We apply
condition (A3) of \eqref{4.21:eq-strips-geometry-a} at zone $i_0+1$: it states that
$\Omega_{i_0+2}$ is at distance at least $L^2\ge 9$ level-$i_0$ units from $\Omega^c_{i_0+1}$.
Since the points of $\tilde{\square}^5$ lie within $7<9$ own units of a cell contained in
$\Omega^c_{i_0+1}$, the cells meeting $\tilde{\square}^5$ have level at most
$\operatorname{zone}(\square)+1$.

\emph{At most one class component.} The set $\tilde{\square}^5$ has diameter at most $14$
level-$k$ units. By condition (G0) of \eqref{4.21:eq-strips-geometry-a} and
$\check{R}_k\ge 38$, $\tilde{\square}^5$ meets at most one class component. In particular the
component $C'$ in \eqref{4.21:eq-block-raise-a} is well defined whenever $\tilde{\square}^5$
meets a cell raised during this $\mathbf{R}_k$. The cells raised during this $\mathbf{R}_k$ in
$C'$ have levels in the interval $[h+1,j(C')]$.

\emph{Proof of (i).} By Definition~\ref{4.21:def-old-window-blocks}, a block of
$\mathfrak{O}^{+}$ falls under one of two cases.
\begin{itemize}
\item It has zone at most $h$ and lies within $2$ own units of $Z_h\cap C'$ or of a strip; each
of these sets lies inside $Z''_{h+1}$.
\item It lies within $2$ own units of an untreated component of $\Lambda_k^c$.
\end{itemize}
The cells raised during this $\mathbf{R}_k$, on the other hand, lie inside $Z$ and outside
$Z''_{h+1}$. The set $Z''_{h+1}$ contains the collar of width $3L\check{R}_{h+1}$ about $Z'_h$.
Therefore, in either case, $\tilde{\square}^5$ meets no cell raised during this $\mathbf{R}_k$,
and $s(\square)=0$ for every $\square\in\mathfrak{O}^{+}$.

For a block $\square$ whose $\tilde{\square}^5$ meets no cell raised during this
$\mathbf{R}_k$, every cell $\Delta$ meeting $\tilde{\square}^5$ has $\mathsf{r}_0(\Delta)=0$,
because $\mathsf{r}_0$ is the raise made during this $\mathbf{R}_k$. Hence $s(\square)=0$ by its
definition as a maximum. Statement (i) concerns $s$, that is, the raise made during this
$\mathbf{R}_k$ only; it asserts nothing about the inherited-datum factor
$\mathfrak{d}(\square)$.

\emph{Proof of (ii).} Lemma~\ref{4.21:lem-under-block}(a) states that a cell
$\Delta\subset Z\cap C'$ of current level $i$, with $h+1\le i\le j(C')$, satisfies
$b\ge\min\{i-1,k_0\}$ on $\Delta$. It applies to every class component $C'$ and not only to
the running one, because a sibling is a preparation with the same $h$ and the same $k_0$. Let
$\Delta$ be such a raised cell, of level $i$. The raise made on $\Delta$ during this
$\mathbf{R}_k$ is at most the total raise of $\Delta$, which is its current level $i$ minus its
original level $m_0(\Delta)$; with the inequality $m_0\ge b$ between the original level and the
base this gives
\[
\mathsf{r}_0(\Delta)\le i-m_0(\Delta)\le i-b\le i-\min\{i-1,k_0\}.
\]
Hence $\mathsf{r}_0\le 1$ on raised cells of level at most $k_0+1$, and $\mathsf{r}_0\le i-k_0$ on
raised cells of level $i\ge k_0+1$. The two bounds are the two cases of
\[
i-\min\{i-1,k_0\}=
\begin{cases}
1, & i\le k_0+1,\\
i-k_0, & i\ge k_0+1 .
\end{cases}
\]

\emph{Proof of (iii).} If $\tilde{\square}^5$ meets only cells of levels at most $k_0+1$, then
every raised cell it meets has $\mathsf{r}_0\le 1$ by (ii), so $s(\square)\le 1$.

If $\operatorname{zone}(\square)\le k_0+1$, the first step of the proof shows that every cell
meeting $\tilde{\square}^5$ has level at most $k_0+2$. By (ii), $\mathsf{r}_0$ is at most $1$
on such raised cells of level at most $k_0+1$, and at most $(k_0+2)-k_0=2$ on those of level
$k_0+2$. Hence $s(\square)\le 2$.

\emph{Proof of the first line of \eqref{4.21:eq-block-raise-a}.} Let $\tilde{\square}^5$ meet a
cell raised during
this $\mathbf{R}_k$. The raised cells met by $\tilde{\square}^5$ lie in the one class component
$C'$ and have levels at most $j(C')$. By (ii), $\mathsf{r}_0\le i-k_0\le j(C')-k_0$ on those of
level $i$ with $k_0+1\le i\le j(C')$, and $\mathsf{r}_0\le 1$ on those of level at most $k_0+1$.
The top zone of $C'$ satisfies $j(C')\le j+1$, and $j+1-k_0\le N_0$. Therefore
\[
s(\square)\le j(C')-k_0\le j+1-k_0\le N_0 .
\]

\emph{Proof of the second line of \eqref{4.21:eq-block-raise-a}.} There are two cases.

If $s(\square)\le 1$, then $L_{\circ}^{6s(\square)}\le L_{\circ}^{6}$, since $L_{\circ}\ge 2$.
Moreover $L_{\circ}^{6}\le L^{d}$ and $\mathsf{l}_j\ge 1$, so
$L_{\circ}^{6s(\square)}\le L^{d}\mathsf{l}_j^{2}$.

If $s(\square)\ge 2$, then by (ii) the set $\tilde{\square}^5$ meets a raised cell of level at
least $k_0+2$, since raised cells of level at most $k_0+1$ have $\mathsf{r}_0\le 1$. That cell
lies in $C'$ and has level at most $j(C')$, so $j+1\ge j(C')\ge k_0+2$, and hence $j>k_0$. For
$j>k_0$ the data-size condition that fixes $\mathsf{l}$ and $\mathsf{l}_j$ gives
$\mathsf{l}_j=4\mathsf{l}$, together with $L_{\circ}^{6N_0}\le L^{4N_0}\le\mathsf{l}^{2}$. With
the first line of \eqref{4.21:eq-block-raise-a}, which applies because $\tilde{\square}^5$ meets
a raised cell, and
$L_{\circ}\ge 2$ we obtain
\[
L_{\circ}^{6s(\square)}\le L_{\circ}^{6N_0}\le L^{4N_0}\le\mathsf{l}^{2}
\le\mathsf{l}_j^{2}\le L^{d}\,\mathsf{l}_j^{2}.
\]
This proves the second line of \eqref{4.21:eq-block-raise-a} for every block.
\end{proof}

\begin{lemma}[Depth-paid shares of the inherited datum]\label{4.21:lem-datum-depth-shares}
Work at stage $n+1$ of $\mathbf{R}_k$. Assume the floor on $w$ recorded in
\eqref{4.21:eq-cell-block-sums-a} and the condition on $\check{R}$ recorded in
\eqref{4.21:eq-strips-geometry-a}. For a block $\square$ of the cover put
$n_0:=\operatorname{zone}(\square)$, the zone index of $\square$
(Lemma~\ref{4.21:lem-block-raise}), and
\begin{equation*}
\varrho(\square):=1+\frac{d^{\wedge}(\square,R_n)}{w}.
\end{equation*}
For a cell $\Delta\in\mathcal{W}$ let $t(\Delta):=k_a$ be the step of the arrested attempt
$\mathfrak{a}$ whose preparation raised $\Delta$; this attempt, the raiser of $\Delta$, is a
live arrested attempt. We call $L_{\circ}^{6\mathsf{r}(\Delta)}$ the \emph{share} of $\Delta$.

Part (i) below involves no bound on $\mathfrak{d}$. Parts (ii) and (iii) are stated under
the \emph{block form of the inherited datum} at $\square$, namely the equality
\begin{equation*}
\mathsf{l}^{2}(\square)=\max\Bigl\{L_{\circ}^{6s(\square)},\ \max\bigl\{L_{\circ}^{6\mathsf{r}(\Delta)}:\
\Delta\in\mathcal{W},\ \Delta\cap\tilde{\square}^{5}\neq\emptyset\bigr\}\Bigr\};
\end{equation*}
this is the at-most-cubic growth bound in the proof of the overlap estimate of the preparation,
read, for each cell $\Delta\in\mathcal{W}$, at the step $k_a$, the level $k_0(k_a)$ and the top
zone $j_a$ of the raiser of $\Delta$, with the frozen $\mathsf{L}_0$-weight $w_{\mathfrak{a}}$
of Definition~\ref{4.21:def-arrested-attempts}. Since
$\mathsf{l}^{2}(\square)=L_{\circ}^{6s(\square)}\mathfrak{d}(\square)$, it implies
\begin{equation}\label{4.21:eq-datum-depth-shares-1}
\mathfrak{d}(\square)\le\max\Bigl\{1,\ \max\bigl\{L_{\circ}^{6\mathsf{r}(\Delta)}:\
\Delta\in\mathcal{W},\ \Delta\cap\tilde{\square}^{5}\neq\emptyset\bigr\}\Bigr\}.
\end{equation}
Under \eqref{4.21:eq-datum-depth-shares-1} the bound on $\mathfrak{d}(\square)$ is a maximum,
over the cells $\Delta\in\mathcal{W}$ met by $\tilde{\square}^{5}$, of their shares; parts (ii)
and (iii) are therefore stated cell by cell, and the right side of
\eqref{4.21:eq-datum-depth-shares-1} depends on the scaffold alone, so that it is the same for
the near-group total of one machine and for the difference of the two machines' totals on one
scaffold, the two cases of Definition~\ref{4.21:def-datum-hypothesis}. Part (iv) uses, in place
of \eqref{4.21:eq-datum-depth-shares-1}, its own hypothesis \eqref{4.21:eq-datum-depth-shares-3},
a level-discounted form of the same shares, and of the rest of the lemma only part (i) and the
computation of part (ii).
\begin{enumerate}
\item[(i)] (Scope.) Let $L\ge 8$. If $\tilde{\square}^{5}$ meets a cell of $\mathcal{W}$, then
either $\square\in\mathfrak{O}^{+}$ (Definition~\ref{4.21:def-old-window-blocks}), or
$\square$ is \emph{exceptional}: $n_0=k$, and $\square$ lies over the band of a component
$C''$ exiting at the last stage, whose stage level is $j=k-1$ in the notation of
Lemma~\ref{4.21:lem-base-below-level}, so that $j(C'')=k=j+1$; by
Lemma~\ref{4.21:lem-base-below-level}, the level $j(C'')=j+1$ occurs only for a component
exiting at the current stage in the branch described there. In both cases
\begin{equation}\label{4.21:eq-datum-depth-shares-2}
\varrho(\square)\ge\check{R}_k-4,\qquad k-n_0+1\le N_{*}+\varrho(\square).
\end{equation}
Here $N_{*}$ is the depth constant of Lemma~\ref{4.21:lem-old-block-depth}, fixed in
\eqref{4.21:eq-old-block-depth-a}.
Of part (i), parts (ii)--(iv) use only the two inequalities
\eqref{4.21:eq-datum-depth-shares-2}, and they apply to every block satisfying them, whether
it lies in $\mathfrak{O}^{+}$, is exceptional, or neither.
\item[(ii)] (The depth-paid shares.) Let $\square$ be a block satisfying
\eqref{4.21:eq-datum-depth-shares-2}; this includes every block of part (i). For every
$\Delta\in\mathcal{W}$ met by $\tilde{\square}^{5}$ with $t(\Delta)\ge n_0-2$,
\begin{equation}\label{4.21:eq-datum-depth-shares-a}
\begin{gathered}
L_{\circ}^{6\mathsf{r}(\Delta)}\le C^{\mathtt{a}}_{\mathrm{pr}}\,
e^{\frac18\delta_P d^{\wedge}(\square,R_n)},\\
C^{\mathtt{a}}_{\mathrm{pr}}:=\sup_{\varrho\ge1}L^{12}
\bigl(\varrho+7+c_{\mathsf{n}(\varrho)}\bigr)^{6r_{\mathrm{pr}}}e^{2-\varrho},\qquad
\mathsf{n}(\varrho):=C_N2^{r_{\mathrm{pr}}}(\varrho+4)+\varrho+4,
\end{gathered}
\end{equation}
where $C^{\mathtt{a}}_{\mathrm{pr}}$ is a number determined by $L$, $\lambda$, $c_2$ and
$r_{\mathrm{pr}}$. Hence, under \eqref{4.21:eq-datum-depth-shares-1},
\begin{equation*}
\mathfrak{d}(\square)\le C^{\mathtt{a}}_{\mathrm{pr}}\,e^{\frac18\delta_P d^{\wedge}(\square,R_n)}
\end{equation*}
whenever every $\Delta\in\mathcal{W}$ met by $\tilde{\square}^{5}$ has $t(\Delta)\ge n_0-2$.
The share of a cell $\Delta\in\mathcal{W}$ met by the probe $\tilde{\square}$ is always one of
those bounded here, since contact with the probe gives $n_0\le t(\Delta)$; in particular, if
$\tilde{\square}$ meets the raised cells of an arrested attempt $\mathfrak{a}$, so that
$n_0\le k_a$, their shares are bounded cell by cell. Contact of the probe with the raised cells
of one arrested attempt does not bound the shares of raised cells of older arrested attempts
met by $\tilde{\square}^{5}\setminus\tilde{\square}$: a block may have probe contact with the
raised cells of one arrested attempt and fall under (iii) through those of another.
\item[(iii)] (The residual.) What (ii) leaves is exactly the share of a cell
$\Delta\in\mathcal{W}$ with $t(\Delta)\le n_0-3$ met by $\tilde{\square}^{5}$; such a cell is
never met by the probe $\tilde{\square}$, and it belongs to a structure older than the
block's own zone index. For such a cell
\begin{equation}\label{4.21:eq-datum-depth-shares-4}
L_{\circ}^{6\mathsf{r}(\Delta)}\le L^{12}
\bigl(\check{\mathsf{T}}_k+3+c_{(k-n_0)+(n_0-t(\Delta))}\bigr)^{6r_{\mathrm{pr}}}.
\end{equation}
Of the index, $k-n_0$ is controlled by the depth through the second inequality of
\eqref{4.21:eq-datum-depth-shares-2}; the age $n_0-t(\Delta)$ is bounded above by none of the
conditions in force, the condition of Definition~\ref{4.21:def-arrested-attempts} on the gaps
of a nest bounding them from below only. The right side of \eqref{4.21:eq-datum-depth-shares-4}
is thus a fixed power of $\check{\mathsf{T}}_k+3+c_u$ taken at an index $u$ containing the age
$n_0-t(\Delta)$, for which no upper bound is available.
\item[(iv)] (The level-discounted form implies the inherited-datum hypothesis.) Let $L\ge8$ as
in part (i). Suppose that for a number $a>0$ fixed before $\gamma$, at every block,
\begin{equation}\label{4.21:eq-datum-depth-shares-3}
\mathsf{l}^{2}(\square)\le2\max\Bigl\{1,\ \max\bigl\{L_{\circ}^{6\mathsf{r}(\Delta)}
L^{-a(n_0-\operatorname{lev}\Delta)_{+}}:\ \Delta\text{ a cell meeting }
\tilde{\square}^{5}\bigr\}\Bigr\}.
\end{equation}
Then the inherited-datum hypothesis (Definition~\ref{4.21:def-datum-hypothesis}, a hypothesis
whose proof is outstanding) holds, in the form \eqref{4.21:eq-datum-hypothesis-a}, with
\begin{equation}\label{4.21:eq-datum-depth-shares-5}
C^{\mathtt{a}}:=1+2^{6r_{\mathrm{pr}}+1}(1+\mathfrak{w}_a)C^{\mathtt{a}}_{\mathrm{pr}},\qquad
\mathfrak{w}_a:=\sup_{v\ge0}c_v^{6r_{\mathrm{pr}}}L^{-av}<\infty.
\end{equation}
\end{enumerate}
\end{lemma}

\begin{proof}
Throughout we write $s:=s(\square)$ and use that raises are nonnegative, that $L_{\circ}>1$, and
that the numbers $c_u$, which by their definition are logarithms of $1+3\lambda u+2c_2\ge1$, are
nonnegative and nondecreasing in $u$.

\emph{The block form.} By the equality for $\mathsf{l}^{2}(\square)$ displayed in the statement
and $L_{\circ}^{6s}\ge1$,
\begin{equation*}
\mathsf{l}^{2}(\square)\le L_{\circ}^{6s}\max\Bigl\{1,\ \max\bigl\{L_{\circ}^{6\mathsf{r}(\Delta)}:\
\Delta\in\mathcal{W},\ \Delta\cap\tilde{\square}^{5}\neq\emptyset\bigr\}\Bigr\};
\end{equation*}
dividing by $L_{\circ}^{6s(\square)}$ and using
$\mathsf{l}^{2}(\square)=L_{\circ}^{6s(\square)}\mathfrak{d}(\square)$
(Definition~\ref{4.21:def-local-data-factor}) gives
\eqref{4.21:eq-datum-depth-shares-1}.

\emph{The bound on one share.} Let $\Delta\in\mathcal{W}$ and $t:=t(\Delta)$. By part (a) of
Lemma~\ref{4.21:lem-under-block}, $\mathsf{r}(\Delta)\le N_0(t)$; by part (c) of the same lemma
(the corresponding display),
$N_0(t)\le4+2r_{\mathrm{pr}}\log_L(\check{\mathsf{T}}_k+3+c_{k-t})$; and
$L_{\circ}^{6}\le L^{3}$, by the condition relating $L_{\circ}$ to $L$ where $L_{\circ}$ is
fixed. Therefore
\begin{equation}\label{4.21:eq-datum-depth-shares-6}
L_{\circ}^{6\mathsf{r}(\Delta)}\le L^{3\mathsf{r}(\Delta)}\le L^{3N_0(t)}
\le L^{12+6r_{\mathrm{pr}}\log_L(\check{\mathsf{T}}_k+3+c_{k-t})}
=L^{12}\bigl(\check{\mathsf{T}}_k+3+c_{k-t}\bigr)^{6r_{\mathrm{pr}}}.
\end{equation}

\emph{Part (i).} Let $\square$ be a window block, that is, $\square\notin\mathfrak{O}^{+}$
(Definition~\ref{4.21:def-old-window-blocks}), and let $\Delta_0$ be its own cell. By
Lemma~\ref{4.21:lem-block-raise}, every point of $\tilde{\square}^{5}$ lies within $7$ units of
level $n_0$ of $\Delta_0$. It therefore suffices to show that, apart from one case, the distance
from $\Delta_0$ to the union of the cells of $\mathcal{W}$ exceeds $7$ such units.

We first locate $\mathcal{W}$:
\begin{equation}\label{4.21:eq-datum-depth-shares-7}
\mathcal{W}\subset\bigcup_{C'}\bigl(Z_h\cap(\Omega^{\mathbf{T}}_h)^c\cap C'\bigr)\ \cup\
\bigcup_{U}\bigl((\Omega^{\mathbf{T}}_k)^c\cap U\bigr),
\end{equation}
where $C'$ runs over the class components, running or sibling, and $U$ over the untreated
components. Indeed, a raiser $\mathfrak{a}$ of step $t$ has its raised cells in
\begin{equation*}
W_{\mathfrak{a}}=X_{\mathfrak{a}}\cap(\Omega^{\mathbf{T}}_t)^c\subset X_{\mathfrak{a}}\subset Z_t,
\end{equation*}
as in the proof of Lemma~\ref{4.21:lem-window-probe}. Inside a class
component $C'$ one has $t\le h$ and $X_{\mathfrak{a}}\subset Z_h$; for $t=h$ this gives
$W_{\mathfrak{a}}\subset Z_h\cap(\Omega^{\mathbf{T}}_h)^c$ directly, and for $t\le h-1$ one has
$X_{\mathfrak{a}}\subset Z_{h-1}\subset(\Omega^{\mathbf{T}}_h)^c$ by
\eqref{4.21:eq-stages-components-a}. Inside an untreated component $U$ one has $t\le k-1$ and
$X_{\mathfrak{a}}\subset Z_{k-1}\subset(\Omega^{\mathbf{T}}_k)^c$.

We now bound the distance from $\Delta_0$ to the set on the right of
\eqref{4.21:eq-datum-depth-shares-7}, in units of level $n_0$, according to the position of
$\Delta_0$.
\begin{itemize}
\item $\Delta_0\subset\Lambda_k$. The distance is at least $2\check{R}_k$, by the collar
property of the $\mathbf{T}$-step at step $k$ stated in \cite[(1.10)]{B-CMP122a}, the collar
being $\Omega^{\mathbf{T}}_k\setminus\Lambda_k$.
\item $\Delta_0\subset C'$, in a zone made by the $\mathbf{T}$-step, with
$n_0\ge j(C')+1$. Then $\Delta_0\subset\Omega_{n_0}$, while $\Omega_h^c\subset Z_h\subset
Z_{n_0-1}$; by part (d) of Lemma~\ref{4.21:lem-live-region-geometry} the distance is at least
$8\check{R}_{n_0}$.
\item $\Delta_0\subset C'$, in a prepared zone, so that $\Delta_0\subset Z\setminus
Z''_{n_0}$. By the widths of part (b) of Lemma~\ref{4.21:lem-under-block}, the distance is at
least $3\check{R}_{n_0}$ if $h+1\le n_0\le k_0$; at least $\check{R}_{k_0+1}$ if $n_0=k_0+1$;
and at least $L^{k-n_0}\ge L$ if $k_0+2\le n_0\le k-1$.
\item $n_0=h$, so that $\Delta_0\subset\Lambda_h\cap C'$. If
$\Delta_0\subset\Lambda^{\mathbf{T}}_h$, the distance is at least $2\check{R}_h$, by the collar
property of the $\mathbf{T}$-step at step $h$ stated in \cite[(1.10)]{B-CMP122a}: the collar
$\Omega^{\mathbf{T}}_h\setminus\Lambda^{\mathbf{T}}_h$ consists of two layers of unraised
level-$h$ cells. If $\Delta_0$ lies in a component healed at $\mathbf{R}_h$, every live
$W_{\mathfrak{a}}$ lies in another component $X''$ of $Z^{\mathbf{T}}_h$, at distance at least
$\check{R}_h$ by \eqref{4.21:eq-strips-geometry-a}, and behind the collar
$X''\cap\Omega^{\mathbf{T}}_h$ of $X''$.
\end{itemize}
Each of these lower bounds exceeds $7$, since $\check{R}_m\ge38$ for every $m$ by
\eqref{4.21:eq-strips-geometry-a} and $L\ge8$. In all these cases
$\tilde{\square}^{5}$ therefore meets no cell of $\mathcal{W}$. The one case left is
$n_0=k=j(C'')$, where the distance is $3$ units: this is the exceptional block of part (i).
Hence, if $\tilde{\square}^{5}$ meets a cell of $\mathcal{W}$, then $\square\in\mathfrak{O}^{+}$
or $\square$ is exceptional.

It remains to prove \eqref{4.21:eq-datum-depth-shares-2}. For $\square\in\mathfrak{O}^{+}$ the
two inequalities are the depth bounds of Lemma~\ref{4.21:lem-old-block-depth}, displayed in
\eqref{4.21:eq-old-block-depth-a}. Let $\square$ be exceptional. It lies within $1$ level-$k$
unit of $C''$; the region $R_n$ lies in running components, at distance at least
$\check{R}_k$ from $C''$ by \eqref{4.21:eq-strips-geometry-a}; and all levels are at most $k$.
By \eqref{4.21:eq-cell-block-sums-a}, therefore, $d^{\wedge}(\square,R_n)\ge w(\check{R}_k-2)$,
whence
\begin{equation*}
\varrho(\square)\ge1+\check{R}_k-2\ge\check{R}_k-4 .
\end{equation*}
Moreover $n_0=k$, so $k-n_0+1=1$, and $1\le\varrho(\square)\le N_{*}+\varrho(\square)$, as
$N_{*}\ge0$.

\emph{Part (ii).} Let $\square$ satisfy \eqref{4.21:eq-datum-depth-shares-2}, write
$\varrho:=\varrho(\square)\ge1$, and let $\Delta\in\mathcal{W}$ meet $\tilde{\square}^{5}$ with
$t:=t(\Delta)\ge n_0-2$. By the second inequality of \eqref{4.21:eq-datum-depth-shares-2},
\begin{equation*}
k-t\le k-n_0+2\le N_{*}+\varrho+1=(N_{*}-2)+\varrho+3 .
\end{equation*}
The number $N_{*}-2$ is at most $C_N2^{r_{\mathrm{pr}}}\check{R}_k$, by the bound recorded in
\eqref{4.21:eq-strips-geometry-a} together with the definition of $N_{*}$ in
\eqref{4.21:eq-old-block-depth-a}. With this and with
$\check{R}_k\le\varrho+4$, the first inequality of \eqref{4.21:eq-datum-depth-shares-2}, we
obtain
\begin{equation*}
k-t\le C_N2^{r_{\mathrm{pr}}}(\varrho+4)+\varrho+3\le\mathsf{n}(\varrho).
\end{equation*}
Further $\check{\mathsf{T}}_k\le\check{R}_k^{1/r_{\mathrm{pr}}}\le\varrho+4$, the first
inequality by \eqref{4.21:eq-strips-geometry-a}. Inserting these
into \eqref{4.21:eq-datum-depth-shares-6}, and using that $c_u$ is nondecreasing in $u$,
\begin{equation*}
L_{\circ}^{6\mathsf{r}(\Delta)}\le L^{12}\bigl(\varrho+7+c_{\mathsf{n}(\varrho)}\bigr)^{6r_{\mathrm{pr}}}
\le C^{\mathtt{a}}_{\mathrm{pr}}\,e^{\varrho-2},
\end{equation*}
the last inequality by the definition of $C^{\mathtt{a}}_{\mathrm{pr}}$ as a supremum over
$\varrho\ge1$. Finally $e^{\varrho(\square)-2}\le e^{\frac18\delta_Pd^{\wedge}(\square,R_n)}$ by
the depth inequality recorded in \eqref{4.21:eq-old-block-total-a}. This proves the first line
of \eqref{4.21:eq-datum-depth-shares-a}.

Under \eqref{4.21:eq-datum-depth-shares-1}, $\mathfrak{d}(\square)$ is at most the maximum of $1$
and of the shares of the cells of $\mathcal{W}$ met by $\tilde{\square}^{5}$. When every such
cell has $t(\Delta)\ge n_0-2$, each share obeys the bound just proved; and so does $1$, since
$C^{\mathtt{a}}_{\mathrm{pr}}\ge1$ (the quantity under the supremum at $\varrho=2$ is at least
$L^{12}9^{6r_{\mathrm{pr}}}\ge1$) and $d^{\wedge}\ge0$. Hence
$\mathfrak{d}(\square)\le C^{\mathtt{a}}_{\mathrm{pr}}e^{\frac18\delta_Pd^{\wedge}(\square,R_n)}$.

An exceptional block never falls under this last conclusion: no $\Delta\in\mathcal{W}$ that it
meets has $t(\Delta)\ge n_0-2$, since $t\le h\le k-3=n_0-3$. The statement about cells met by
the probe is proved in part (iii).

\emph{Part (iii).} The cells of $\mathcal{W}$ met by $\tilde{\square}^{5}$ not covered by (ii)
are those with $t(\Delta)\le n_0-3$. For such a cell, \eqref{4.21:eq-datum-depth-shares-6} with
$k-t=(k-n_0)+(n_0-t)$ is \eqref{4.21:eq-datum-depth-shares-4}.

The probe $\tilde{\square}$ never meets such a cell. Indeed, if $\tilde{\square}$ meets a cell
$\Delta\in\mathcal{W}$ raised by $\mathfrak{a}$, then $n_0\le t(\Delta)=k_a$, by the first step
of the proof of part (c) of Lemma~\ref{4.21:lem-under-block}. That computation uses of
$\mathfrak{a}$ only that $X_{\mathfrak{a}}\subset Z_{k_a}$ and that every later attempt on its
component, of level $h_1$ say, has $h_1\ge k_a$, which is one of the conditions of
Definition~\ref{4.21:def-arrested-attempts}; both hold for the raiser of every cell of
$\mathcal{W}$. Thus $t(\Delta)\ge n_0>n_0-3$ at probe contact. This also proves the statement of
part (ii) that the share of a cell of $\mathcal{W}$ met by the probe is among those bounded
there.

\emph{Part (iv).} Let $\Delta$ be a cell meeting $\tilde{\square}^{5}$. If
$\Delta\notin\mathcal{W}$, then $\mathsf{r}(\Delta)=\mathsf{r}_0(\Delta)\le s$ by the definition
of $s(\square)$ (Definition~\ref{4.21:def-local-data-factor}), so that its term in
\eqref{4.21:eq-datum-depth-shares-3} is at most $L_{\circ}^{6s}$.

Let $\Delta\in\mathcal{W}$ and $t:=t(\Delta)$. Then $\operatorname{lev}(\Delta)\le t$, so that
$(n_0-\operatorname{lev}\Delta)_{+}\ge v:=(n_0-t)_{+}$. Put
$A:=\check{\mathsf{T}}_k+3+c_{k-n_0}\ge1$. Since $k-t\le(k-n_0)+v$, and since
$c_{u+v}\le c_u+c_v$, which follows, $c_u$ being a logarithm of $1+3\lambda u+2c_2$, from
\begin{equation*}
1+3\lambda(u+v)+2c_2\le(1+3\lambda u+2c_2)(1+3\lambda v+2c_2),
\end{equation*}
we have $c_{k-t}\le c_{k-n_0}+c_v$. The number $\mathfrak{w}_a$ is finite because $c_v$ grows
at most logarithmically in $v$ while $L^{-av}$ decays geometrically. By
\eqref{4.21:eq-datum-depth-shares-6}, using
$(x+y)^{p}\le2^{p}(x^{p}+y^{p})$ for $x,y,p\ge0$, $L^{-av}\le1$ and $A\ge1$,
\begin{align*}
L_{\circ}^{6\mathsf{r}(\Delta)}L^{-a(n_0-\operatorname{lev}\Delta)_{+}}
&\le L^{12}(A+c_v)^{6r_{\mathrm{pr}}}L^{-av}
\le2^{6r_{\mathrm{pr}}}L^{12}\bigl(A^{6r_{\mathrm{pr}}}+c_v^{6r_{\mathrm{pr}}}L^{-av}\bigr)\\
&\le2^{6r_{\mathrm{pr}}}L^{12}\bigl(A^{6r_{\mathrm{pr}}}+\mathfrak{w}_a\bigr)
\le2^{6r_{\mathrm{pr}}}(1+\mathfrak{w}_a)L^{12}A^{6r_{\mathrm{pr}}}.
\end{align*}
Since $\Delta\in\mathcal{W}$ meets $\tilde{\square}^{5}$, part (i) shows that $\square$
satisfies \eqref{4.21:eq-datum-depth-shares-2}. As in part (ii), $k-n_0\le k-n_0+2\le
\mathsf{n}(\varrho)$ and $\check{\mathsf{T}}_k\le\varrho+4$ with $\varrho=\varrho(\square)$,
so that
\begin{equation*}
L^{12}A^{6r_{\mathrm{pr}}}\le L^{12}\bigl(\varrho+7+c_{\mathsf{n}(\varrho)}\bigr)^{6r_{\mathrm{pr}}}
\le C^{\mathtt{a}}_{\mathrm{pr}}e^{\frac18\delta_Pd^{\wedge}(\square,R_n)}.
\end{equation*}
Hence
\begin{equation*}
L_{\circ}^{6\mathsf{r}(\Delta)}L^{-a(n_0-\operatorname{lev}\Delta)_{+}}
\le\Xi:=2^{6r_{\mathrm{pr}}}(1+\mathfrak{w}_a)C^{\mathtt{a}}_{\mathrm{pr}}\,
e^{\frac18\delta_Pd^{\wedge}(\square,R_n)} .
\end{equation*}
Since $L_{\circ}^{6s}\ge1$, each term of the maximum in \eqref{4.21:eq-datum-depth-shares-3},
the term $1$ included, is at most $L_{\circ}^{6s}(1+\Xi)$. Therefore, as
$e^{\frac18\delta_Pd^{\wedge}(\square,R_n)}\ge1$,
\begin{equation*}
\mathsf{l}^{2}(\square)\le2L_{\circ}^{6s}(1+\Xi)
\le L_{\circ}^{6s}\bigl(1+e^{\frac18\delta_Pd^{\wedge}(\square,R_n)}+2\Xi\bigr)
=L_{\circ}^{6s}\bigl(1+C^{\mathtt{a}}e^{\frac18\delta_Pd^{\wedge}(\square,R_n)}\bigr),
\end{equation*}
with $C^{\mathtt{a}}$ as in \eqref{4.21:eq-datum-depth-shares-5}. Dividing by $L_{\circ}^{6s}$ and using
$\mathsf{l}^{2}(\square)=L_{\circ}^{6s(\square)}\mathfrak{d}(\square)$
(Definition~\ref{4.21:def-local-data-factor}) gives
$\mathfrak{d}(\square)\le1+C^{\mathtt{a}}e^{\frac18\delta_Pd^{\wedge}(\square,R_n)}$ at every block,
which is \eqref{4.21:eq-datum-hypothesis-a}; and $C^{\mathtt{a}}\ge C^{\mathtt{a}}_{\mathrm{pr}}$ since
$2^{6r_{\mathrm{pr}}+1}(1+\mathfrak{w}_a)\ge1$.
\end{proof}

\begin{remark}[The residual configuration and exceptional window blocks]
\label{4.21:rem-residual-configuration}
We use the raised cells, arrested attempts, collars and frozen data of
Definition~\ref{4.21:def-arrested-attempts}, the old blocks, window blocks and window points of
Definition~\ref{4.21:def-old-window-blocks}, and parts~(i)--(iv) of
Lemma~\ref{4.21:lem-datum-depth-shares} on the depth-paid shares of the inherited datum. We
write $\mathcal{W}$ for the set of cells inherited by the current step in raised state from the
preparation of a live arrested attempt of an earlier step.
\begin{enumerate}
\item[(1)] \emph{The residual configuration.} The residual of part~(iii) of
Lemma~\ref{4.21:lem-datum-depth-shares} occurs exactly at the blocks $\square$ of the following
kind.
\begin{itemize}
\item $\square$ belongs to the completed top zone, of level $k_b=n_0$, of a later live arrested
attempt $\mathfrak{b}$ with $j_b=k_b$, arising in one of the branches of
Definition~\ref{4.21:def-arrested-attempts} for which $j_b=k_b$ (for the exit branch, only at
the last stage of the operation); here $n_0$ denotes the level of the zone of the block, as in
part~(i) of Lemma~\ref{4.21:lem-datum-depth-shares}.
\item $\square$ lies at distance $3$ units of level $k_b$ from $Z'_{k_0(k_b)}$, in the sense of
\cite[(1.10)]{B-CMP122a}.
\item The enlargement $\tilde{\square}^5$ of $\square$ covers the raised cells of an older live
arrested attempt $\mathfrak{a}$ nested in the core of $\mathfrak{b}$, so that $k_a\le h_b$.
\end{itemize}
The older structure is always such a live arrested attempt. If the region size parameter
$\check{R}$ is constant along the nest, the gap between $\mathfrak{a}$ and $\mathfrak{b}$ is
smaller than $4+11/(L-1)$ units of level $k_b$, whatever the size of $\check{R}$. The
configuration is therefore generic; it does not require $\check{R}$ to be close to its lower
bound. The same set of blocks is the one singled out by parts~(b) and~(c2) of the localisation
lemma for the inherited datum.
\item[(2)] \emph{Exceptional window blocks.} The exceptional blocks of part~(i) of
Lemma~\ref{4.21:lem-datum-depth-shares} are window blocks of the configuration of item~(1). They
arise only when a sibling component $C''$ of Definition~\ref{4.21:def-arrested-attempts} has
$j(C'')=j+1$ with $j=k-1$, in the notation of that definition, so no zone of level $\le k-1$
contains one. At zone $k$ they do occur. For this reason the inherited-datum hypothesis
\eqref{4.21:eq-datum-hypothesis-a} is posed at every block, and its constant is carried inside
the total $\mathfrak{s}_J$ of the perturbative-site step.
\item[(3)] \emph{Zone $h$.} Beside the components of an operation completed at
$\mathbf{R}_h$, no window block of zone $h$ meets a cell of $\mathcal{W}$. Moreover, let $C'$
be a class component, and let $\square$ be a block of zone $h$ whose enlargement
$\tilde{\square}^5$ reaches $Z_h\cap C'$. Then $\square$ satisfies the
hypothesis of part~(ii) of Lemma~\ref{4.21:lem-datum-depth-shares}, in the depth form
\eqref{4.21:eq-old-block-depth-a}; so nothing beyond the inherited-datum hypothesis is needed
at $\square$, even if it meets inherited raised cells there.
\item[(4)] \emph{A sufficient discount.} Let $i\ge\ell$ be levels, let $s_j$ denote the length
scale of level $j$, so that $s_\ell/s_i=L^{-(i-\ell)}$, and let $\alpha'$ be the exponent of the
near-group bound at a block. The bound on the near-group total at a block carries, as its own
factor on its share of level $\ell$, the factor
\begin{equation}\label{4.21:eq-residual-configuration-1}
L^{4(i-\ell)}\Bigl(\frac{s_\ell}{s_i}\Bigr)^{4+\alpha'}=L^{-\alpha'(i-\ell)} .
\end{equation}
Charged as a discount on the level-$\ell$ share of the inherited datum, it would give the
hypothesis of part~(iv) of Lemma~\ref{4.21:lem-datum-depth-shares} with $a=\alpha'$. We do not
prove that it may be so charged: at a residual block the older structure enters only through
the hypothesis cube $\tilde{\square}^5$ of \cite[Lemma~4]{B-CMP109}, which no responding
near-group term meets, so that the discount, as the bound is stated, would be charged on cells
that no responding near-group term meets.
\end{enumerate}
\end{remark}

\begin{proof}
\emph{Item} (1). By the nesting statement for arrested attempts, the older attempt
$\mathfrak{a}$ nested in the core of $\mathfrak{b}$ satisfies $k_a\le h_b$; by the gap statement
for nests of arrested attempts, when $\check{R}$ is constant along the nest the gap is less than
$4+11/(L-1)$ in units of level $k_b$. That
estimate involves no lower bound on $\check{R}$, so it holds at every size of $\check{R}$.

It remains to see that the older structure is a live arrested attempt, and that no such attempt
lies inside a live pocket. We first show that a completed operation $\mathbf{R}$ leaves no cell
of $\mathcal{W}$. Each
of its components either is healed as a whole \cite{B-CMP122b}, or is thrown back with the
flat label \cite[(1.33)]{B-CMP122b}. In neither case does a raised cell of the operation
survive into a later step.

Second, no live arrested attempt outlives a completed $\mathbf{R}$ on its component: an
arrested attempt is, by Definition~\ref{4.21:def-arrested-attempts}, an attempt of an operation
not completed on that component. Hence no older live arrested attempt lies inside a component
thrown back with the flat label. Consequently the structures older than $\mathfrak{b}$ whose
cells $\tilde{\square}^5$ can cover are live arrested attempts, as stated.

\emph{Item} (2). By Definition~\ref{4.21:def-arrested-attempts}, a sibling component $C''$ has
$j(C'')=j+1$ only when it arises from an exit of the current stage. It then enters the next
stage through its geometry alone; for all other siblings whose frozen terms are read, one has
$j(C'')\le j$. An exceptional block of part~(i) of Lemma~\ref{4.21:lem-datum-depth-shares}
arises only through the case $j(C'')=j+1$ at $j=k-1$ of the analysis of siblings. Hence there
is no exceptional block in a zone of level $\le k-1$.

The case $j=k-1$ does occur in Ba{\l}aban's construction. In \cite[p.~188]{B-CMP122a}, an exit
at a stage, \cite[(1.55)]{B-CMP122a}, is performed before the fluctuation integral of the same
stage, \cite[(1.60)]{B-CMP122a}. A component exiting at that stage is therefore present, with
$j(C'')=j+1$, when the fluctuation integral of the stage is taken; the case $j=k-1$ is not
empty. Exceptional window blocks of
zone $k$ therefore exist, and the inherited-datum hypothesis
\eqref{4.21:eq-datum-hypothesis-a} cannot be restricted to the zones below $k$. It is posed at
every block.

\emph{Item} (3), \emph{first assertion}. We treat in turn the two kinds of component of an
operation completed at $\mathbf{R}_h$: those healed as a whole and those thrown back as pockets.

A component healed as a whole \cite{B-CMP122b} is treated as follows. By
Lemma~\ref{4.21:lem-window-probe}, its points have $m_0=b\ge h$, so they carry no inherited
raise, and they hold no cell of $\mathcal{W}$. Hence, by the rule that sets
$\mathfrak{d}(\square)=1$ at a block meeting no cell of $\mathcal{W}$, one has $\mathfrak{d}=1$
there, and the datum takes its local form.

A component thrown back as a pocket $Y$ carries the flat label $Y^{(h)}$
\cite[(1.33)]{B-CMP122b}. Its strip keeps every window probe at distance at least
$\tfrac12(\check{R}_h-1)-2$ from $Y$. This lies beyond the reach $7$ of $\tilde{\square}^5$,
and the interior of $Y$ contains no inherited raised cell.

Finally, in the case $n_0=h$ of part~(i) of Lemma~\ref{4.21:lem-datum-depth-shares}, every
cell of $\mathcal{W}$ lies at distance at least $2\check{R}_h$ from the cell of a window block
of zone $h$. This follows from the separation property of the construction, used together with
the collars $\mathfrak{c}_{\mathfrak{a}}$ of Definition~\ref{4.21:def-arrested-attempts}. So no
window block of zone $h$ meets $\mathcal{W}$.

\emph{Item} (3), \emph{second assertion}. Let $\square$ be a block of zone $h$ whose
enlargement $\tilde{\square}^5$ reaches $Z_h\cap C'$. Every point of $\square$ is then within
$9$ units of level $h$ of $Z_h\cap C'$. Since $9<\check{R}_h$, this gives
\begin{equation}\label{4.21:eq-residual-configuration-2}
\square\subset Z\cap Z'^{\sim 3}_h\subset Z''_{h+1}.
\end{equation}

The collar region $R_n$ of the current stage lies outside $Z''_{h+1}$. By the geometry of the
domains made at step $h$, one has at level $h$
\begin{equation*}
\operatorname{dist}_h\bigl(Z_h\cap C',(Z''_{h+1})^c\bigr)\ge 3\check{R}_h .
\end{equation*}
Hence a contour from a cell meeting $\square$ to $R_n$ runs through at least
$3\check{R}_h-9$ units of level $h$, through cells of level $\le h$. By the contour
characterisation of the scaled distance $d^{\wedge}$, with $w=M/M_1$,
\begin{equation}\label{4.21:eq-residual-configuration-3}
d^{\wedge}(\square,R_n)\ge w\,(3\check{R}_h-9).
\end{equation}

Moreover, the depth of $\square$ satisfies
\begin{equation}\label{4.21:eq-residual-configuration-4}
\varrho(\square)\ge 3\check{R}_h-8\ge \check{R}_k-4 .
\end{equation}
Here $n_0=h$, and the number of stages of $\mathbf{R}_k$ is set equal to the region size
parameter $\check{R}_k$. Then
\begin{equation}\label{4.21:eq-residual-configuration-5}
k-n_0+1=\check{R}_k+1\le N_*+\varrho(\square).
\end{equation}

The bounds \eqref{4.21:eq-residual-configuration-3}--\eqref{4.21:eq-residual-configuration-5}
are the hypothesis of part~(ii) of Lemma~\ref{4.21:lem-datum-depth-shares}, in the form
\eqref{4.21:eq-old-block-depth-a}. Part~(ii) then applies. For every cell $\Delta$ with
$t(\Delta)\ge h-2$, the share of $\Delta$ is paid for, in the form
\eqref{4.21:eq-datum-depth-shares-a}, by the constant of that display together with the factor
$e^{\frac18\delta_P d^{\wedge}(\square,R_n)}$. A cell $\Delta$ with $t(\Delta)\le h-3$ falls
under the residual of part~(iii). That residual is the inherited-datum hypothesis
\eqref{4.21:eq-datum-hypothesis-a} itself, which is posed at every block by item~(2). Nothing
beyond that hypothesis is used at such a block.

\emph{Item} (4). By the identity \eqref{4.21:eq-residual-configuration-1}, the displayed
factor equals $L^{-\alpha'(i-\ell)}$ on the level-$\ell$ share. This is the level discount
required in the hypothesis of part~(iv) of Lemma~\ref{4.21:lem-datum-depth-shares}, with
$a=\alpha'$.

We now describe what happens at a residual block. There, every responding near-group term has
its region inside $\tilde{\square}^2$, and none of these terms meets the older structure. The
older structure enters only through the cube entering the hypothesis of
\cite[Lemma~4]{B-CMP109}, which here is $\tilde{\square}^5$. The level-share discount
\eqref{4.21:eq-residual-configuration-1} would therefore be applied to cells that no responding
term meets.
\end{proof}

\begin{definition}[The small-field constant at a block]\label{4.21:not-small-field-constant}
We use the following constants.
\begin{itemize}
\item $K_S$ and $E_0$ are the constants of the bound on the far small-field terms of the step of
the correlation theorem: at any block $\square$ of the cover, the total of these terms at $\square$
is at most $K_S(E_0+1)$, the long terms being controlled through their diameter.
\item $A_{\flat}$ is the constant of the hull statement for these terms.
\item $K(d)$ is the stage-count constant.
\item We write $C_{\mathrm{dat}}$ for the constant of the inherited-datum hypothesis
(Definition~\ref{4.21:def-datum-hypothesis}, whose proof is outstanding). It is the number for which
that hypothesis asserts $\mathfrak{d}(\square)\le 1+C_{\mathrm{dat}}\,
e^{\delta_P d^{\wedge}(\square,R_n)/8}$ at every block of the cover.
\item $\mathfrak{n}(\square)$ and $\mathfrak{n}_0$ are the near-group datum and its constant, as in
\eqref{4.21:eq-local-data-factor-a} of the local data factor at a block
(Definition~\ref{4.21:def-local-data-factor}).
\item $K_A=4\mathsf{K}_r e^{C_2\kappa_1}\ge 1$ is the block constant of the covering bound for
$A$-class elements.
\item $\mathfrak{p}_j$ is the rim weight and $\vartheta$ is the block decay factor. Both are
functions of the scale variable $\mathsf{T}$, the logarithm of the inverse coupling, and
$\vartheta(\square)$ denotes the block decay factor at the block $\square$; in
\eqref{4.21:eq-small-field-constant-a} these arguments are suppressed.
\end{itemize}
The \emph{small-field constant at a block} is
\begin{equation}\label{4.21:eq-small-field-constant-a}
\mathfrak{s}_J:=\Bigl[A_{\flat}K_SK(d)(E_0+1)
+(1+C_{\mathrm{dat}})\,\frac{\mathfrak{n}_0}{K_A}\Bigr]
\Bigl(1+\frac{\mathfrak{p}_j}{\vartheta}\Bigr).
\end{equation}
In \eqref{4.21:eq-small-field-constant-a} the number $1+C_{\mathrm{dat}}$ stands inside the bracket
and multiplies the near-group term at every block. The factor $1+\mathfrak{p}_j/\vartheta$, which we
call the scaffold factor, is kept displayed as a separate factor.

The near-group term enters in the form $\mathfrak{n}_0/K_A$ for the following reason. An element of
the $A$-class at the block $\square$ carries the weight $\mathfrak{N}^{\natural}K_A\vartheta(\square)$,
that is, its size $\mathfrak{N}^{\natural}$ multiplied by the block factor $K_A\vartheta(\square)$.
A near group at the block $\square$ has activity
\begin{equation}\label{4.21:eq-small-field-constant-1}
\nu=\mathfrak{n}(\square)\,e^{-(\kappa_1-1)|\hat{\mathsf{y}}|}.
\end{equation}
It is read in the same way as an element of the $A$-class, as a term weight measured against the
same block factor $K_A\vartheta(\square)$; this is the origin of the division by $K_A$ in
\eqref{4.21:eq-small-field-constant-a}. Here $K_A\ge1$.

The rim weight and the block decay factor are subject to the condition
\begin{equation}\label{4.21:eq-small-field-constant-2}
\sup_{\mathsf{T}}\frac{\mathfrak{p}_j(\mathsf{T})}{\vartheta(\mathsf{T})}\le 1 ,
\end{equation}
which is one of the conditions imposed on the weights $\mathfrak{p}_j$ and $\vartheta$. Under it the
scaffold factor in \eqref{4.21:eq-small-field-constant-a} is at most $2$.
\end{definition}

\begin{proof}
At each value of $\mathsf{T}$ the scaffold factor is $1+\mathfrak{p}_j(\mathsf{T})/\vartheta(\mathsf{T})$.
The quotient $\mathfrak{p}_j(\mathsf{T})/\vartheta(\mathsf{T})$ is bounded by its supremum over
$\mathsf{T}$. By \eqref{4.21:eq-small-field-constant-2} that supremum is at most $1$. Hence
\begin{equation*}
1+\frac{\mathfrak{p}_j(\mathsf{T})}{\vartheta(\mathsf{T})}
\le 1+\sup_{\mathsf{T}'}\frac{\mathfrak{p}_j(\mathsf{T}')}{\vartheta(\mathsf{T}')}\le 2 .
\end{equation*}
\end{proof}

\begin{lemma}[Position-weighted sums over blocks and rim cubes]\label{4.21:lem-position-sums}
Let $\pi_x(\Diamond)$ be the position weight of a block, a cube or a component $\Diamond$, and let
$C_{\mathrm{pos}}$, $c_{\star}''$ be the constants fixed in \eqref{4.21:eq-fine-count-a}. Let $V'$ be
the block sum of Notation~\ref{4.21:not-plain-unit-data}. Let $C_N$ and $r_{\mathrm{pr}}$ be the
constants, and $\tau_d(\cdot)$ the tail sum, of \eqref{4.21:eq-strips-geometry-a}. Let $N$ denote
the stage count of \eqref{4.21:eq-strips-geometry-a}; in this lemma $N$ is not the rank of the gauge
group. Assume the conditions on the admissible geometry displayed in
\eqref{4.21:eq-strips-geometry-a}; of these, the bound on the stage count and the tail condition are
used below. Let $R_n$ be the collar region of the current stage, $\bar c$ its components, with
regions $R(\bar c)$ and rim cubes as fixed with the collar region, and $d^{\wedge}$ the scaled
distance. Define $\mathsf{r}'''$ by the first line of \eqref{4.21:eq-position-sums-a}. Then
$\mathsf{r}'''\ge 2+2N/\check{R}_k$, and for every point $x$ and every $\beta\ge\tfrac18\delta_P$ the
two bounds of the second line hold, the first sum running over all blocks $\square_0$ of the cover and
the second over all rim cubes $q$ of the components $\bar c$:
\begin{equation}\label{4.21:eq-position-sums-a}
\begin{gathered}
\mathsf{r}''' := 2+2C_N\cdot 2^{r_{\mathrm{pr}}},\\
\sum_{\square_0}e^{-\beta d^{\wedge}(\square_0,R_n)}\,\pi_x(\square_0)\le c_{\star}''\,V',
\qquad
\sum_{q}\pi_x(q)\le C_{\mathrm{pos}}\,V' .
\end{gathered}
\end{equation}
\end{lemma}

\begin{proof}
\emph{The inequality for $\mathsf{r}'''$.} After subtracting $2$ and dividing by $2$, the inequality
$\mathsf{r}'''\ge 2+2N/\check{R}_k$ is equivalent to
\begin{equation*}
\frac{N}{\check{R}_k}\le C_N\cdot 2^{r_{\mathrm{pr}}} .
\end{equation*}
This follows from the bound on the stage count displayed in \eqref{4.21:eq-strips-geometry-a}.

\emph{The block sum: the shells.} Write $\operatorname{dist}_k$ for the distance measured at level
$k$. Fix $x$ and $\beta\ge\tfrac18\delta_P$. For each integer $u\ge0$, the $u$-shell is the set of
blocks $\square_0$ with
\begin{equation*}
\operatorname{dist}_k(x,\square_0)\in\bigl[u\check{R}_k,(u+1)\check{R}_k\bigr[ .
\end{equation*}
The blocks of the $u$-shell lie in a box of side $2(u+1)\check{R}_k+4$.

\emph{Splitting $R_n$ into components.} Let $\bar c$ run over the components making up $R_n$. Then,
for every block $\square$,
\begin{equation}\label{4.21:eq-position-sums-1}
e^{-\beta d^{\wedge}(\square,R_n)}\le\sum_{\bar c}e^{-\beta d^{\wedge}(\square,R(\bar c))}.
\end{equation}
We insert \eqref{4.21:eq-position-sums-1} into the sum over the blocks of one shell. We then sort the
components $\bar c$ by the distance of $R(\bar c)$ from the box of that shell.

\emph{Components near the box.} We use the comparison of scaled distances with level-$k$ distances
among the bounds carried from the plain-unit expansion in Notation~\ref{4.21:not-plain-unit-data}.
It applies at $\beta\ge\tfrac18\delta_P$, with $w=M/M_1$ the stratum unit, because
$\tfrac1{16}\delta_Pw\ge1$. Together with the packing bound of
Lemma~\ref{4.21:lem-tree-length-facts}, displayed in \eqref{4.21:eq-tree-length-facts-a}, it shows
that at most $(2u+6)^d$ components $\bar c$ have their regions within distance $\check{R}_k$ of the
box. For each such $\bar c$, the sum over blocks of $e^{-\beta d^{\wedge}(\square,R(\bar c))}$ is at
most $V'$. Indeed, $\beta\ge\tfrac18\delta_P\ge\tfrac1{16}\delta_P$, and $V'$ is the supremum over
$n$ and $\bar c$ of the sum over all blocks of $e^{-\frac1{16}\delta_Pd^{\wedge}(\square,R(\bar c))}$.

\emph{Components far from the box.} Let $l\ge1$. By the same comparison and the same packing bound,
the components $\bar c$ whose regions lie at distance in
$\bigl[l\check{R}_k,(l+1)\check{R}_k\bigr[$ from the box number at most $(2u+2l+6)^d$. For
$u\ge0$ and $l\ge1$ we have
\begin{equation*}
(2u+7)(2l+2)=4ul+4u+14l+14\ge 2u+2l+6 ,
\end{equation*}
and therefore
\begin{equation*}
(2u+2l+6)^d\le(2u+7)^d(2l+2)^d .
\end{equation*}
The same comparison shows that each such component carries, besides the factor $V'$, the factor
$e^{-(l\check{R}_k-5)}$.

\emph{The bound per shell.} The near components contribute at most $(2u+6)^dV'\le(2u+7)^dV'$. The
far components contribute at most
\begin{equation*}
(2u+7)^dV'\sum_{l\ge1}(2l+2)^de^{-(l\check{R}_k-5)} .
\end{equation*}
Under the tail condition among \eqref{4.21:eq-strips-geometry-a}, the two contributions together are
at most $2(2u+7)^dV'$ per $u$-shell.

\emph{The position weight.} On the $u$-shell the position weight satisfies
\begin{equation*}
\pi_x(\square_0)\le\min\bigl\{1,e^{-(u-205)\theta_0}\bigr\}
\end{equation*}
by \eqref{4.21:eq-fine-count-a}. Summing over $u\ge0$ gives
\begin{equation*}
\sum_{\square_0}e^{-\beta d^{\wedge}(\square_0,R_n)}\pi_x(\square_0)
\le\sum_{u\ge0}2(2u+7)^d\min\bigl\{1,e^{-(u-205)\theta_0}\bigr\}\,V'
\le\tfrac12c_{\star}''\,V' ,
\end{equation*}
with the constant $c_{\star}''$ of \eqref{4.21:eq-fine-count-a}. This is at most $c_{\star}''V'$,
which is the first bound of the second line of \eqref{4.21:eq-position-sums-a}.

\emph{Rim cubes.} Fix a component $\bar c$. The number of rim cubes of $\bar c$ is at most the number
of cells of the closure of $R(\bar c)$. By part (c) of Lemma~\ref{4.21:lem-cell-block-sums},
\eqref{4.21:eq-cell-block-sums-a}, read at scaled distance $d^{\wedge}=0$, the number of those cells
is at most $V'$. Hence each component has at most $V'$ rim cubes.

A rim cube $q$ of $\bar c$ satisfies $\pi_x(q)\le\pi_x(\bar c)$. Therefore
\begin{equation*}
\sum_q\pi_x(q)\le V'\sum_{\bar c}\pi_x(\bar c)\le C_{\mathrm{pos}}V' ,
\end{equation*}
where the last step is the bound on the summed position weights of the components displayed in
\eqref{4.21:eq-fine-count-a}. This is the second bound of the second line of
\eqref{4.21:eq-position-sums-a}.

No smallness of the tail sum $\tau_d$ is used in this proof beyond what the tail condition of
\eqref{4.21:eq-strips-geometry-a} requires, namely finiteness.
\end{proof}

\begin{lemma}[The near groups' sum over blocks]\label{4.21:lem-near-group-sum}
Let the step $k$, the stage $n$, the class components and the blocks of the cover be as in this
section. The sum below runs over the blocks of every zone, and each block is measured in the
units of its own zone. For a block $\square$:
\begin{itemize}
\item $s(\square)$ is the raise read by $\square$ (Definition~\ref{4.21:def-local-data-factor});
\item $d^{\wedge}(\square,R_n)$ is its scaled distance to the collar region $R_n$;
\item $\pi_x(\square)$ is its position weight at the point $x$.
\end{itemize}
The constants $C_{\mathrm{pos}}$ and $c_{\star}''$ are those of \eqref{4.21:eq-fine-count-a}, and
$V'$ is the volume sum in the bound \eqref{4.21:eq-position-sums-a}. The radius
$\check{R}^{\uparrow}$ is that of
Definition~\ref{4.21:def-strips-geometry}, and $\mathsf{l}_j$ is the data-size constant at the top
level $j$ of the running class component.

Assume the condition on the stratum unit $w$ in \eqref{4.21:eq-cell-block-sums-a}, and the size
condition and the tail condition on the radii $\check{R}_i$ in \eqref{4.21:eq-strips-geometry-a}.
Then for every point $x$ and every rate $\beta\ge\frac18\delta_P$, with $d=4$ wherever a numerical
value of the dimension enters,
\begin{equation}\label{4.21:eq-near-group-sum-a}
\begin{gathered}
\sum_{\square}L_{\circ}^{6s(\square)}e^{-\beta d^{\wedge}(\square,R_n)}\pi_x(\square)
\le L^{2d}c_{\star}''V'
+C_{\mathrm{pos}}\bigl[4^d(100L\check{R}^{\uparrow})^d+C_{\mathrm{nr}}\mathsf{l}_j^2\bigr],\\
C_{\mathrm{nr}}:=64^d+3\cdot 112^dL^4 .
\end{gathered}
\end{equation}
The constant $C_{\mathrm{nr}}$ depends on $d$ and $L$ only. In \eqref{4.21:eq-near-group-sum-a}
the volume sum $V'$ multiplies only the number $L^{2d}c_{\star}''$. The terms
$(100L\check{R}^{\uparrow})^d$ and $C_{\mathrm{nr}}\mathsf{l}_j^2$ are added to it; they are not
multiplied by it.
\end{lemma}

\begin{proof}
Throughout, $d=4$ wherever a numerical value of the dimension is used.

We recall the levels and constants of this section that enter the proof. The levels $k_0$ and $h$
and the stage bound $N_0$ are those fixed with the data-size constants; of $N_0$ we use only
$k-k_0\le N_0$. By Lemma~\ref{4.21:lem-block-raise}, the cells raised during $\mathbf{R}_k$ in a
class component $C'$ have levels in $[h+1,j(C')]$. By the choice of the data-size constants,
$\mathsf{l}_j=4\mathsf{l}$ for $j>k_0$, and $L^{4N_0}\le\mathsf{l}^2$. We write $Z=\Lambda^c$ for
the complement of the current small-field region.

We use two bounds on $L_{\circ}$, the small blocking constant. It is chosen with $2\le L_{\circ}$
and $L_{\circ}^2\le L/3$; in \cite[p.~182]{B-CMP122a} it is subject to $2\le L_{\circ}<\frac12L$.
The first bound is $L_{\circ}^6\ge2$, which follows from $L_{\circ}\ge2$. The second is
$L_{\circ}^6\le L^3/27$, which follows from $L_{\circ}^2\le L/3$ by taking cubes. From the second
bound we get
\begin{equation*}
L_{\circ}^{6}L^{-4}\le\frac{1}{27L}\le 1,
\qquad
L_{\circ}^{12}\le \frac{L^{6}}{729}\le L^{8}=L^{2d}.
\end{equation*}

We split the blocks according to the value of $s(\square)$. We write $\Sigma^{(1)}$,
$\Sigma^{(2)}$ and $\Sigma^{(3)}$ for the parts of the left side of
\eqref{4.21:eq-near-group-sum-a} contributed by the blocks of cases (i), (ii-a) and (ii-b) below.
Every block falls under exactly one of these cases.

\emph{Case} (i): $s(\square)\le2$. By Lemma~\ref{4.21:lem-block-raise} this case contains the
following blocks:
\begin{itemize}
\item every block of $\mathfrak{O}^{+}$;
\item every block off $Z$ whose $\tilde{\square}^5$ meets no raised cell;
\item every block of zone at most $k_0+1$.
\end{itemize}
For a block of this case, $L_{\circ}^{6s(\square)}\le L_{\circ}^{12}\le L^{2d}$. Hence
\begin{equation*}
\Sigma^{(1)}\le L^{2d}\sum_{\square}e^{-\beta d^{\wedge}(\square,R_n)}\pi_x(\square)
\le L^{2d}c_{\star}''V' .
\end{equation*}
The second inequality is the first bound of \eqref{4.21:eq-position-sums-a}
(Lemma~\ref{4.21:lem-position-sums}), which applies because $\beta\ge\frac18\delta_P$.

\emph{Case} (ii): $s(\square)\ge3$. By Definition~\ref{4.21:def-local-data-factor}, $s(\square)$
is the largest raise $\mathsf{r}_0(\Delta)$ in this operation $\mathbf{R}_k$ of a cell $\Delta$
meeting $\tilde{\square}^5$. So $\tilde{\square}^5$ meets a cell $\Delta$ raised during this
$\mathbf{R}_k$ with $\mathsf{r}_0(\Delta)\ge3$.

By Lemma~\ref{4.21:lem-block-raise}, such a cell has current level $i\in[k_0+3,J]$ and lies inside
one class component $C'$, where $J:=j(C')\le j+1$. If $C'$ is the running component then $J=j$; if
$C'$ is a sibling component, $J$ is its frozen top level. Such cells exist only if $J\ge k_0+3$.
Hence $j\ge k_0+2>k_0$, and therefore $\mathsf{l}_j=4\mathsf{l}$ by the choice of the data-size
constants. The raised cells are either cells
of zone $J$ or cells of the thin strata, with $k_0+3\le i\le J-1$.

\emph{The zone of a block in case} (ii). By Lemma~\ref{4.21:lem-block-raise}, the cells meeting
$\tilde{\square}^5$ have level at most the zone of $\square$ plus one. Hence the zone of $\square$
is at least $i-1\ge k_0+2$.

It is also at most $J+1$. Consider a block of a zone $i'\ge J+2$ made by $\mathbf{T}$. Its
$\tilde{\square}^5$ lies within $7$ level-$i'$ units of its own cell. The zone $i'-1\ge J+1$ below
it is made by $\mathbf{T}$. By Lemma~\ref{4.21:lem-live-region-geometry} and the size condition in
\eqref{4.21:eq-strips-geometry-a}, that zone is at least $8\check{R}_{i'}\ge 8\cdot 38$ level-$i'$
units thick. So its $\tilde{\square}^5$ meets no cell of level at most $i'-2$, and in particular
none of the raised cells of $C'$, whose levels are at most $J$.

Therefore the blocks of case (ii) belong to the zones $J$ or $J+1$, or to the thin zones
$m\in[k_0+2,J-1]$ of $C'$. By Lemma~\ref{4.21:lem-block-raise}, $\tilde{\square}^5$ meets at most
one class component, so every block of case (ii) determines its component $C'$.

\emph{The position weight and the exponential in case} (ii). In every case $\square$ lies within
$9$ units of its own level of $C'$. These units are at most level-$k$ units, so $\square$ lies
within $9$ level-$k$ units of $C'$. Hence
\begin{equation*}
\pi_x(\square)\le e^{9/(3N)}\pi_x(C')\le e\,\pi_x(C'),
\end{equation*}
the last inequality because $N\ge3$; here $N$ denotes the length parameter of the position weight
$\pi_x$ of \eqref{4.21:eq-fine-count-a}, not the rank of the gauge group. Of the weight
$e^{-\beta d^{\wedge}(\square,R_n)}$ only the
bound $e^{-\beta d^{\wedge}(\square,R_n)}\le1$ is used.

\emph{Case} (ii-a): \emph{blocks of the zones $J$ and $J+1$.}

\emph{Localisation of the blocks with a given raise.} Let $3\le s\le s(\square)$. Then
$\tilde{\square}^5$ meets, at some point, a cell whose original level $m_0$ is at most its current
level minus $s$, hence at most $J-s$. That is, $\tilde{\square}^5$ meets
$\Omega^{(k)c}_{J-s+1}\cap C'$. By the zone relations in \eqref{4.21:eq-strips-geometry-a},
\begin{equation*}
\Omega^{(k)c}_{J-s+1}\cap C'\subset Z_{J-s+1}\cap C' .
\end{equation*}
This set is one box of side at most $100\check{R}_{J-s+1}$ level-$(J-s+1)$ units, by the nesting of
the zones of a class component, which holds at every level in $[h+1,k]$. It applies here because
$J-s+1\in[k_0+1,k]\subset[h+1,k]$, as $s\le J-k_0$. In level-$J$ units the side is
$100\check{R}_{J-s+1}L^{1-s}$, which is at most $100\check{R}^{\uparrow}L^{1-s}$.

\emph{Counting.} The blocks of zone $J$ are unions of $2^d$ cubes of $\pi_J$. If $\tilde{\square}^5$
meets the box, the cubes of $\square$ lie in the box enlarged by $9$ on each side. So the number of
such blocks is at most
\begin{equation*}
2^d\bigl(100\check{R}^{\uparrow}L^{1-s}+20\bigr)^d
\le 2^{2d-1}\bigl[(100L\check{R}^{\uparrow})^dL^{-ds}+20^d\bigr].
\end{equation*}
Here we used $(a+b)^d\le 2^{d-1}(a^d+b^d)$ for $a,b\ge0$ and
$(100\check{R}^{\uparrow}L^{1-s})^d=(100L\check{R}^{\uparrow})^dL^{-ds}$. The blocks of zone $J+1$
have coarser units, and there are none if $J=k$. Their number is at most the same. So, for each
$C'$ and each $s$, the blocks of zones $J$ and $J+1$ with $s(\square)\ge s$ number at most
$4^d[(100L\check{R}^{\uparrow})^dL^{-ds}+20^d]$.

\emph{Resolving the raise into classes.} For a block of this case,
\begin{equation*}
\mathbf{1}[s(\square)\ge3]\,L_{\circ}^{6s(\square)}\le\sum_{3\le s\le s(\square)}L_{\circ}^{6s}.
\end{equation*}
Indeed, all summands are positive and the last one alone equals the left side. Moreover
$s(\square)\le J-k_0\le N_0$. We exchange the sum over blocks with the sum over $s$, and use
$\pi_x(\square)\le e\,\pi_x(C')$ and the count above. By the position bound of
\eqref{4.21:eq-fine-count-a}, the sum $e\sum_{C'}\pi_x(C')$ is at most $C_{\mathrm{pos}}$. This
gives
\begin{align*}
\Sigma^{(2)}
&\le e\sum_{C'}\pi_x(C')\cdot\sum_{s=3}^{N_0}L_{\circ}^{6s}\,
4^d\bigl[(100L\check{R}^{\uparrow})^dL^{-ds}+20^d\bigr]\\
&\le C_{\mathrm{pos}}\,4^d\Bigl[(100L\check{R}^{\uparrow})^d
\sum_{s\ge1}\bigl(L_{\circ}^{6}L^{-d}\bigr)^s+20^d\cdot 2L_{\circ}^{6N_0}\Bigr]\\
&\le C_{\mathrm{pos}}\bigl[4^d(100L\check{R}^{\uparrow})^d+2\cdot 80^dL^{4N_0}\bigr]\\
&\le C_{\mathrm{pos}}\bigl[4^d(100L\check{R}^{\uparrow})^d+80^d\mathsf{l}_j^2/8\bigr].
\end{align*}

\emph{Justification of the steps.}
\begin{itemize}
\item In the second line we used
$\sum_{s=3}^{N_0}L_{\circ}^{6s}L^{-ds}\le\sum_{s\ge1}(L_{\circ}^6L^{-d})^s$, since all terms are
nonnegative, and $\sum_{s\le N_0}L_{\circ}^{6s}\le 2L_{\circ}^{6N_0}$, because $L_{\circ}^6\ge2$.
\item In the third line, the geometric series has ratio $L_{\circ}^6L^{-4}\le 1/(27L)$, and
$\sum_{s\ge1}(27L)^{-s}<1$. Also
$4^d\cdot 20^d=80^d$, and $L_{\circ}^{6N_0}\le L^{4N_0}$ because $L_{\circ}^6L^{-4}\le1$.
\item The last line uses $L^{4N_0}\le\mathsf{l}^2=\mathsf{l}_j^2/16$, from the choice of the
data-size constants and $\mathsf{l}_j=4\mathsf{l}$.
\end{itemize}

\emph{Case} (ii-b): \emph{blocks of the thin zones $m\in[k_0+2,J-1]$.} No depth of the block, and no
lower bound on $d^{\wedge}(\square,R_n)$, is used in this case.

\emph{Counting.} By Lemma~\ref{4.21:lem-live-region-geometry}, the stratum $m$ of $C'$ is contained
in $Z''_{m+1}\subset Z''_k\subset\Lambda_{C'}$, where $\Lambda_{C'}$ is the box of that lemma
attached to $C'$.
The set $\Lambda_{C'}$ is a box of side at most $110$ level-$k$ units, that is, $110L^{k-m}$
level-$m$ units. So each component has at most $2^d(112L^{k-m})^d$ blocks of zone $m$.

\emph{The raise.} The cells met by $\tilde{\square}^5$ for such a block have level at most $m+1$.
So $s(\square)\le(m+1)-k_0$ by Lemma~\ref{4.21:lem-block-raise}.

\emph{Summing over the thin zones.} Put $t:=m-k_0\ge2$. Since $k-m\le N_0-t$, $d=4$ and
$L_{\circ}^6\le L^3/27$,
\begin{align*}
(112L^{k-m})^dL_{\circ}^{6(m+1-k_0)}
&\le 112^dL^{4(N_0-t)}L^{3(t+1)}27^{-(t+1)}\\
&=112^dL^{4N_0+3}\cdot L^{-t}27^{-(t+1)} .
\end{align*}
Summing the geometric series,
\begin{equation*}
\sum_{t\ge2}L^{-t}27^{-(t+1)}=L^{-2}27^{-3}\Bigl(1-\frac{1}{27L}\Bigr)^{-1}\le L^{-2}.
\end{equation*}
Summing over $m$ therefore gives at most $2^d112^dL^{4N_0+1}$ per component.

\emph{Conclusion of case} (ii-b). With $\pi_x(\square)\le e\,\pi_x(C')$ and the position bound of
\eqref{4.21:eq-fine-count-a},
\begin{align*}
\Sigma^{(3)}
&\le e\sum_{C'}\pi_x(C')\cdot 2^d112^dL^{4N_0+1}
\le C_{\mathrm{pos}}\,2^d112^dL\,\mathsf{l}^2\\
&= C_{\mathrm{pos}}\,112^dL\,\mathsf{l}_j^2
\le C_{\mathrm{pos}}\cdot 3\cdot112^dL^4\mathsf{l}_j^2 .
\end{align*}
The second step uses the bound $L^{4N_0}\le\mathsf{l}^2$ from the choice of the data-size
constants. The equality uses
$\mathsf{l}^2=\mathsf{l}_j^2/16$ and $2^d/16=1$ at $d=4$.

\emph{Total.} Adding the three cases,
\begin{equation*}
\Sigma^{(1)}+\Sigma^{(2)}+\Sigma^{(3)}
\le L^{2d}c_{\star}''V'
+C_{\mathrm{pos}}\Bigl[4^d(100L\check{R}^{\uparrow})^d
+\Bigl(\frac{80^d}{8}+3\cdot112^dL^4\Bigr)\mathsf{l}_j^2\Bigr].
\end{equation*}
At $d=4$ we have $(80/64)^4=(5/4)^4=625/256<8$, so $80^d/8\le 64^d$. The right side is therefore at
most the right side of \eqref{4.21:eq-near-group-sum-a}.
\end{proof}

\begin{remark}
We add four comments on the proof of Lemma~\ref{4.21:lem-near-group-sum}.
\begin{enumerate}
\item The mechanism has three parts. There is one box per raise class, by the nesting of the zones
of a class component. The factor $L_{\circ}^6L^{-d}<1$ sums the classes. Only the top classes cost
$\mathsf{l}_j^2$.

Three features matter. First, $s(\square)$ is the raise of what the block reads, not the
difference $j-m_0$; for a deep block of a thin zone that difference would be an age. Second, each
class component is counted at its own top level $J$: the top zone $J<j$ of a sibling component is
not a thin stratum, and it falls under case (ii-a). Third, the $(112L^{k-m})^d$ blocks of the thin
zones are beaten by $L_{\circ}^6\le L^3/27$ alone.
\item Near the level $h$ the collar fits in a few cubes of $\pi_k$, and the blocks of zone $j$
within $8$ level-$k$ units of $x$ are comparable in number to $V'$. Their share is the term
$L^{2d}c_{\star}''V'$ of case (i), which is $V'$ times a number. The share carrying
$\mathsf{l}_j^2$ comes from case (ii), and it is added to that term.
\item The argument for the thin zones uses no lower bound on $d^{\wedge}(\square,R_n)$. Such a bound
is not available from the block's own cell alone in the form $w(j-m-1)$. The reason is that
$d^{\wedge}(\square,\cdot)$ is a minimum over the cells meeting $\square$, and a block of zone $m$
may meet cells of level $m+1$ (Lemma~\ref{4.21:lem-live-region-geometry}). From the own cell one
obtains only $w(j-m-2)_+$. With that bound and the condition on $w$ in
\eqref{4.21:eq-cell-block-sums-a}, in the form $e^{-\frac12\delta_Pw}\le L^{-4d}$, the sum over
$m$ is at most $3L^{4N_0+4}$. This route leads to the coefficient $3\cdot112^dL^4$ of
$\mathsf{l}_j^2$ in $C_{\mathrm{nr}}$, which also dominates the coefficient $112^dL$ obtained in the
proof.
\item The lemma concerns the raise produced by the present operation $\mathbf{R}_k$ only. The
inherited datum enters the near groups' total through the factor one plus the constant of the
inherited-datum hypothesis. The rate
$\frac38\delta_P$ at which the lemma is applied there satisfies $\frac38\delta_P\ge\frac18\delta_P$.
\end{enumerate}
\end{remark}

\begin{subequations}\label{4.21:eq-group-totals-a}
\begin{proposition}[The group totals in sum form]\label{4.21:prop-group-totals}
In the setting of Notation~\ref{4.21:not-plain-unit-data}, at stage $n+1$ of the operation
$\mathbf{R}_k$, assume the following.
\begin{itemize}
\item The condition $(F_{\kappa_1})^{\flat}$ of Notation~\ref{4.21:not-plain-unit-data}, the
condition $(F_{w})$ of \eqref{4.21:eq-cell-block-sums-a}, the condition $(\mathrm{c}\text{-lev})'$ of
Notation~\ref{4.21:not-weights-incidences}, the floor $(F_{M})$ on the parameter $M$, the conditions
$(F_{\kappa})^{P}$, $(\mathrm{c}\text{-}\check{R})$ and $(\mathrm{c}\text{-tail})$, which are among
the hypotheses of Lemma~\ref{4.21:lem-window-bounds},
Lemma~\ref{4.21:lem-old-block-total} and Lemma~\ref{4.21:lem-near-group-sum}, and the six
supremum conditions (q1)--(q6) on the coupling displayed in \eqref{4.21:eq-coupling-conditions-a}.
\item The inherited-datum hypothesis of Definition~\ref{4.21:def-datum-hypothesis}, whose proof is
outstanding, for both quantities to which it applies: the near-group total of one machine at a
block, and the difference of that total between the two machines $\mathsf{A}$, $\mathsf{B}$ on one
scaffold.
\item The bound for the number of fine small-field leaves stated in the counting lemma for those
leaves, and the proposition on the group totals of the kept members, which supplies the numbers
$\mu^{\mathrm{LEAF}}$ and $\mathsf{M}^{\mathrm{LEAF}}$ displayed in \eqref{4.21:eq-group-totals-4}.
\end{itemize}
Then the group bound \eqref{4.21:eq-plain-unit-data-a}, in its un-anchored form and in its anchored
forms, holds for the groups $\hat{e}$. Put
\begin{equation}\label{4.21:eq-group-totals-1}
\mu^{\flat}:=2e^{c_J+\frac12}\Bigl\{2K_2^{\flat}\vartheta c_{\mathrm{loc}}
\bigl[C^{P,\star}_{\mathsf{a}}+\mathfrak{l}_j\mathfrak{t}^{\mathrm{win}}
+2L^d\mathsf{l}_j^2\mathfrak{s}_J\bigr]
+\Bigl(\frac{K_Q}{r_A}+\varepsilon_S\Bigr)\bigl(2\mathfrak{t}^{\mathrm{win}}+3^d\mathfrak{s}_J\bigr)\Bigr\},
\end{equation}
\begin{equation}\label{4.21:eq-group-totals-2}
\begin{split}
\mathsf{M}^{\flat}:=2e^{c_J+\frac12}\Bigl\{K_2^{\flat}\vartheta\Bigl[
&c_{\star}''V'\bigl(C^{P,\star}_{\mathsf{a}}+(5L)^d\mathfrak{t}^{\mathrm{win}}+\mathfrak{s}_J\bigr)
+c_{\mathrm{loc}}\bigl(\mathfrak{F}^{\mathrm{fi}}+\mathfrak{F}^{\mathrm{fi}}_J\bigr)\\
&+\mathfrak{s}_J\Bigl(L^{2d}c_{\star}''V'
+C_{\mathrm{pos}}\bigl(4^d(100L\check{R}^{\uparrow})^d+C_{\mathrm{nr}}\mathsf{l}_j^2\bigr)\Bigr)\Bigr]\\
&+\Bigl(\frac{K_Q}{r_A}+\varepsilon_S\Bigr)C_{\mathrm{pos}}V'
\bigl(2\mathfrak{t}^{\mathrm{win}}+3^d\mathfrak{s}_J\bigr)\Bigr\},
\end{split}
\end{equation}
where
\begin{equation}\label{4.21:eq-group-totals-3}
C^{P,\star}_{\mathsf{a}}:=C^{P,\mathrm{all}}_{\mathsf{a}}+\mathfrak{s}_J,
\qquad
\mathfrak{F}^{\mathrm{fi}}_J:=C_{\mathrm{pos}}4^d(N_0+2)(104L\check{R}^{\uparrow})^d\mathfrak{s}_J,
\end{equation}
and put
\begin{equation}\label{4.21:eq-group-totals-4}
\mu^{\mathrm{LEAF}}:=4e^{c_J+\frac32}K_2^{\flat}\vartheta\, c_{\mathrm{loc}}\,\mathfrak{n}^{\mathrm{kept}},
\qquad
\mathsf{M}^{\mathrm{LEAF}}:=2e^{c_J+\frac32}K_2^{\flat}\vartheta\, c_{\star}''V'\,\mathfrak{n}^{\mathrm{kept}}.
\end{equation}
Here $c_{\mathrm{loc}}$ is the constant of Lemma~\ref{4.21:lem-grain-block-sum};
$\mathfrak{t}^{\mathrm{win}}$ and $\mathfrak{l}_j$ are those of
Notation~\ref{4.21:not-weights-incidences} and Lemma~\ref{4.21:lem-window-bounds};
$\mathfrak{s}_J$ is the constant of Notation~\ref{4.21:not-small-field-constant};
$C^{P,\mathrm{all}}_{\mathsf{a}}$ is that of Lemma~\ref{4.21:lem-old-block-total};
$\mathfrak{F}^{\mathrm{fi}}$ and $C_{\mathrm{pos}}$ are those of Lemma~\ref{4.21:lem-fine-count};
$c_{\star}''$ is that of Lemma~\ref{4.21:lem-position-sums}; $C_{\mathrm{nr}}$ is that of
Lemma~\ref{4.21:lem-near-group-sum}; and $K_2^{\flat}$, $V'$ and $c_{\mathsf{p}}=212^d$ are those of
Notation~\ref{4.21:not-plain-unit-data}. Further, $\mathsf{l}_j$ is the number through which
Lemma~\ref{4.21:lem-window-bounds} bounds $\mathfrak{l}_j$; $c_J$, $\vartheta$, $K_Q$, $r_A$,
$\varepsilon_S$ and $\tilde{\mathfrak{l}}$ are the constants of the proposition on the summation of
the groups at the preparation site; $\mathfrak{n}^{\mathrm{kept}}$ is the kept-member constant of the
proposition on the group totals of the kept members; and $N_0$ is the constant of the counting lemma
for fine small-field leaves whose value increased by two enters the count of those leaves. The number
$C^{P,\star}_{\mathsf{a}}$ accounts, at an old
block, also for the small-field terms whose constant at a block is $\mathfrak{s}_J$
(Notation~\ref{4.21:not-small-field-constant}): their far leaves, and their
near groups with $s=0$, whose inherited-datum factor $\mathfrak{d}$ is bounded through the
inherited-datum hypothesis assumed above (Definition~\ref{4.21:def-datum-hypothesis}, whose proof
is outstanding); both are contained in $\mathfrak{s}_J$. The number
$\mathfrak{F}^{\mathrm{fi}}_J$ bounds the fine small-field leaves among those terms, each anchored
once: their number is the count of the counting lemma for fine small-field leaves, used as stated
there, and the blocks around each of their own cubes are summed by
Lemma~\ref{4.21:lem-grain-block-sum}. The numbers $\mu^{\mathrm{LEAF}}$ and
$\mathsf{M}^{\mathrm{LEAF}}$ are obtained from the terms of $\mu^{\flat}$ and $\mathsf{M}^{\flat}$
that contain $C^{P,\star}_{\mathsf{a}}$ (the two old-block terms) by putting
$e\,\mathfrak{n}^{\mathrm{kept}}$ in the place of $C^{P,\star}_{\mathsf{a}}$.

With these numbers, the totals $\mu$ and $\mathsf{M}_*$ of the proposition on the summation of the
groups at the preparation site satisfy
\begin{equation}\label{4.21:eq-group-totals-5}
\mu\le\mu^{\flat}+\mu^{\mathrm{LEAF}},\qquad
\mathsf{M}_*\le\mathsf{M}^{\flat}+\mathsf{M}^{\mathrm{LEAF}}.
\end{equation}
Consequently, under (q1) and (q2),
\begin{equation}\label{4.21:eq-group-totals-6}
\mu+\tilde{\mathfrak{l}}\le\tfrac14+\tilde{\mathfrak{l}}\le\tfrac12,
\qquad
4\bigl(\mathsf{M}_*+2^dc_{\mathsf{p}}V'\tilde{\mathfrak{l}}\bigr)\le\tfrac12,
\end{equation}
and that summation proposition, the chain-term bound of the preparation site and its corollary hold
with the class $\mathfrak{E}$ replaced by $\hat{\mathfrak{E}}$ together with the groups $\hat{e}$ of
the kept members, with $\mathsf{d}$ replaced by $\mathsf{d}^{\mathrm{rel}}$, and with $K_c$ of that
summation proposition replaced by the label $S_c$. In particular, for the chain terms of stage $n+1$
of $\mathbf{R}_k$, with the constants $a_k$, $a_m$ and in the norm on labels,
\begin{equation}\label{4.21:eq-group-totals-7}
\|\mathbb{C}^{(n+1)}_k\|^{A,\mathbb{C}}_{(n+1)}\le C_0^{\sharp},
\end{equation}
and the outputs are in the format of stored terms, with label $A:=X$ and with $\mathcal{P}$ formed by
the sets $\mathcal{P}_c$ of the constituents (Notation~\ref{4.21:not-weights-incidences}).

All factors of $\mu^{\flat}$ and $\mathsf{M}^{\flat}$ are displayed in
\eqref{4.21:eq-group-totals-1}--\eqref{4.21:eq-group-totals-3}. The quantities $\mathfrak{l}_j$,
$\mathsf{l}_j^2$, $(100L\check{R}^{\uparrow})^d$ and $\mathfrak{F}^{\mathrm{fi}}$ enter as addends
to the terms carrying $V'$ and never multiply them; $V'$ multiplies only numbers and suprema taken
per point or per coarse block. The sum over fine incidences at window blocks entering
$\mathfrak{F}^{\mathrm{fi}}$ is the one bounded in Lemma~\ref{4.21:lem-fine-count}; no bound for a
sum over terms at a window anchor is taken from another statement. The shape of the further
condition on the coupling displayed beside \eqref{4.21:eq-coupling-conditions-a}, beside (q1)--(q6),
is not used.
\end{proposition}

\begin{proof}[Proof of Proposition~\ref{4.21:prop-group-totals}: the first steps]
The first three steps of the proof are the steps of the plain-unit expansion recalled in
Notation~\ref{4.21:not-plain-unit-data}, together with the bounds stated there and in
\eqref{4.21:eq-plain-unit-data-a}; we use them as stated there.

\emph{The factor $e^{\Sigma_c}$.} Let $c\in\mathfrak{F}$, with label $S_c$ and with the quantities
$\mathsf{e}(c)$, $\mathsf{f}(c)$, $\mathsf{p}^{\mathrm{rel}}(c)$ of
Notation~\ref{4.21:not-weights-incidences}. Write $D_c:=d_k([S_c]_k)$ if $c$ is not confined and
$D_c:=0$ if $c$ is confined, so that $\mathsf{e}(c)=D_c/N$, where $N$ is the parameter of
Notation~\ref{4.21:not-weights-incidences} dividing $d_k([S_c]_k)$ in the definition of
$\mathsf{e}(c)$ (not the rank of the gauge group). By condition (q3) of
\eqref{4.21:eq-coupling-conditions-a}, $\Sigma_c\le\mathsf{p}^{\mathrm{rel}}(c)/(6Nc_{\mathsf{p}})$;
with $\mathsf{p}^{\mathrm{rel}}(c)\le c_{\mathsf{p}}(2+D_c/\check{R}_k)$ (see
Notation~\ref{4.21:not-plain-unit-data}, where $\mathsf{p}^{\mathrm{rel}}$ and $c_{\mathsf{p}}$ are
fixed) this gives
\begin{equation*}
\Sigma_c\le\frac{\mathsf{p}^{\mathrm{rel}}(c)}{6Nc_{\mathsf{p}}}
\le\frac{2+D_c/\check{R}_k}{6N}
=\frac{1}{3N}+\frac{\mathsf{e}(c)}{6\check{R}_k}
\le\frac12+\frac{\mathsf{e}(c)}{12}.
\end{equation*}
The equality is the identity $D_c/(6N\check{R}_k)=\mathsf{e}(c)/(6\check{R}_k)$, and the last
inequality holds because $N\ge1$ gives $1/(3N)\le\frac13\le\frac12$ and $\check{R}_k\ge2$ gives
$\mathsf{e}(c)/(6\check{R}_k)\le\mathsf{e}(c)/12$. Hence
\begin{equation}\label{4.21:eq-group-totals-8}
e^{\Sigma_c}\le e^{\frac12+\mathsf{e}(c)/12}.
\end{equation}
Moreover, by the radius condition $N/\check{R}_k\le C_N2^{r_{\mathrm{pr}}}$ (the condition through
which Lemma~\ref{4.21:lem-position-sums} has $\mathsf{r}'''\ge2+2N/\check{R}_k$) and the elementary
inequality $t\le2e^{t/2}$ for $t\ge0$,
\begin{equation*}
2+\frac{D_c}{\check{R}_k}=2+\frac{N}{\check{R}_k}\,\mathsf{e}(c)
\le2+C_N2^{r_{\mathrm{pr}}}\mathsf{e}(c)
\le\mathsf{r}'''e^{\mathsf{e}(c)/2},
\end{equation*}
where $\mathsf{r}'''=2+2C_N2^{r_{\mathrm{pr}}}$ is the number of
Lemma~\ref{4.21:lem-position-sums}: indeed $2\le2e^{\mathsf{e}(c)/2}$ and
$C_N2^{r_{\mathrm{pr}}}\mathsf{e}(c)\le2C_N2^{r_{\mathrm{pr}}}e^{\mathsf{e}(c)/2}$. For $\Sigma_c$ as
carried in \eqref{4.21:eq-plain-unit-data-a} this gives
\begin{equation*}
\Sigma_c\le\Bigl(K_2^{\flat}\vartheta+\frac{K_Q}{r_A}\Bigr)V'c_{\mathsf{p}}\,
\mathsf{r}'''e^{\mathsf{e}(c)/2}.
\end{equation*}
This last bound is used only in the un-anchored form of the group bound
\eqref{4.21:eq-plain-unit-data-a}, which the fifth and sixth steps of the proof do not use.

\emph{The budget of a non-confined chain term.} Let $c$ be a chain term that is not confined. Its
budget is $e^{\mathsf{e}(c)}=e^{d_k([S_c]_k)/N}$: by the rate display of the proposition on the
summation of the groups at the preparation site, the running decay rate there
exceeds the rate in the weight by at least $1/N$, so that the stored summand bounds
$\mathfrak{w}(c)e^{\mathsf{e}(c)}$. This budget is divided as follows. The part
$e^{\mathsf{e}(c)/12}$ pays the $\mathsf{e}$-dependent factor of $e^{\Sigma_c}$ in
\eqref{4.21:eq-group-totals-8}. The part $e^{\mathsf{e}(c)/3}$ is not included in the weight
$W(c)=\mathfrak{w}(c)e^{\frac23\mathsf{e}(c)}$; it sums the chain terms of this $\mathbf{R}_k$ over
their $\check{R}_k$-separated components, through the constant $c_{\mathrm{sib}}$ of
Notation~\ref{4.21:not-weights-incidences}. A further part $e^{\mathsf{e}(c)/3}$ pays the position
weight, as shown next. Since $\frac1{12}+\frac13+\frac13=\frac34\le1$, a quarter of the exponent is
not used.

\emph{The position weight.} Let $c$ be a non-confined chain term, let $x$ be a point such that
$\Delta_k(x)$ meets $S_c$, and let $\Diamond$ meet $X_c$. The set $X_c$ is contained in the
$(101\check{R}_k+1)$-neighbourhood of $[S_c]_k$, $\Diamond$ lies within distance $3$ of $X_c$, and
$x$ lies within distance $1$ of $[S_c]_k$; hence
\begin{equation*}
\operatorname{dist}_k(x,\Diamond)\le d_k([S_c]_k)+103\check{R}_k,
\end{equation*}
and therefore, for the position weight $\pi_x$ of \eqref{4.21:eq-fine-count-a},
$\pi_x(\Diamond)\ge e^{-\mathsf{e}(c)/3}$.

For every $c\in\mathfrak{F}$, every point $x$ such that $\Delta_k(x)$ meets $S_c$ and every
$\Diamond$ meeting $X_c$, the following holds. If $c$ is a leaf-type term, a confined old term or a
local sibling term, then $\mathsf{f}(c)=d_k([S_c]_k)$, and this gives
\begin{equation*}
e^{-\mathsf{f}(c)}\le e^{-(\operatorname{dist}_k(x,\Diamond)-103\check{R}_k)_+}\le\pi_x(\Diamond).
\end{equation*}
If $c$ is a confined chain term, then $\operatorname{dist}_k(x,\Diamond)\le204\check{R}_k$ when the
hull of the group $\hat{e}$ containing $c$ reaches $x$ through
$S_c\cup\hat{\mathsf{K}}\cup\operatorname{Cor}$: the set $X_c$, together
with its strips, lies inside one $\bar{c}$ of the kind over which the supremum defining $V'$ in
Notation~\ref{4.21:not-plain-unit-data} is taken, of side at most $100\check{R}_k$, and
$\hat{\mathsf{K}}$ and $\operatorname{Cor}$ lie within $103\check{R}_k$ of it; hence
$\pi_x(\Diamond)=1$.

\emph{The weight inequality.} In every case
\begin{equation}\label{4.21:eq-group-totals-9}
\mathfrak{w}(c)e^{\mathsf{e}(c)/12}\le W(c)\,\pi_x(\Diamond).
\end{equation}
For a non-confined chain term this follows from $\pi_x(\Diamond)\ge e^{-\mathsf{e}(c)/3}$:
\begin{equation*}
W(c)\pi_x(\Diamond)\ge\mathfrak{w}(c)e^{\frac23\mathsf{e}(c)-\frac13\mathsf{e}(c)}
=\mathfrak{w}(c)e^{\mathsf{e}(c)/3}\ge\mathfrak{w}(c)e^{\mathsf{e}(c)/12}.
\end{equation*}
For the other terms, the factor $e^{\mathsf{f}(c)}$ of the weight is compensated by the lower bound
of $\pi_x(\Diamond)$ obtained above:
\begin{equation*}
W(c)\pi_x(\Diamond)\ge\mathfrak{w}(c)e^{\frac23\mathsf{e}(c)+\mathsf{f}(c)-\mathsf{f}(c)}
\ge\mathfrak{w}(c)e^{\frac23\mathsf{e}(c)}\ge\mathfrak{w}(c)e^{\mathsf{e}(c)/12}.
\end{equation*}
Finally, when the hull of the group reaches $x$ only through a connected subfamily of
$\hat{\mathsf{y}}$ attached to $\hat{\mathsf{K}}_{\square_0}$, the anchored form of the group bound
\eqref{4.21:eq-plain-unit-data-a}, applied with $\mathsf{E}:=\Delta_k(x)$, gives the factor
$e^{-(\operatorname{dist}_k(x,\square_0)-103\check{R}_k)_+}$, and by clause~(iv) of the lemma on the
plain-unit expansion recalled in Notation~\ref{4.21:not-plain-unit-data} this factor is at most
$\pi_x(\square_0)$.

The inequalities \eqref{4.21:eq-group-totals-8} and \eqref{4.21:eq-group-totals-9} are used in the
remaining steps of the proof, which follow.
\end{proof}
\end{subequations}

\begin{lemma}[Sums at one anchor and the first total]\label{4.21:prf-one-anchor-sums}
Work in the setting and under the hypotheses of
Proposition~\ref{4.21:prop-group-totals}. Among these hypotheses is the inherited-datum
hypothesis \eqref{4.21:eq-datum-hypothesis-a} of
Definition~\ref{4.21:def-datum-hypothesis}, whose proof is outstanding.

For a block $\square_0$ of the cover, write $\Sigma^{\mathrm{sf}}(\square_0)$ for the total
weight at $\square_0$ of the small-field species. These species are the far small-field
leaves and the near groups, and their weights are measured in units of the block factor
$K_A\vartheta(\square_0)$.

A \emph{rim cube} is a cube $q$ of $\pi_{j+1}$ with the following properties: it lies in
the zone $j+1$, made by the operation $\mathbf{T}$, of a running component $C'$, and
$b=j+1$ at its points.

Then the following hold.
\begin{enumerate}
\item[(a)] For an old block $\square_0\in\mathfrak{O}^+$ one has
$\Sigma^{\mathrm{sf}}(\square_0)\le\mathfrak{s}_J\,e^{\frac18\delta_P d^{\wedge}(\square_0,R_n)}$.
\item[(b)] For a window block $\square_0$ one has
\begin{equation*}
\mathfrak{T}^{\mathrm{co}}(\square_0)\le(5L)^d\,\mathfrak{t}^{\mathrm{win}} .
\end{equation*}
Moreover, the part of $\Sigma^{\mathrm{sf}}(\square_0)$ due to the far small-field leaves is
at most $\mathfrak{s}_J$, with no factor $\mathsf{l}^2$.
\item[(c)] For a rim cube $q$, the small-field species anchored at $q$ have total weight at
most $3^d\mathfrak{s}_J$. Only far small-field leaves are anchored at a rim cube; near groups
never are.
\item[(d)] The three bounds
\begin{equation}\label{4.21:eq-one-anchor-sums-a}
\begin{aligned}
&\mathfrak{S}^{\mathrm{all}}(\square_0)+\Sigma^{\mathrm{sf}}(\square_0)
 \le C^{P,\star}_{\mathsf{a}}\,e^{\frac18\delta_P d^{\wedge}(\square_0,R_n)}
 &&(\square_0\in\mathfrak{O}^+),\\
&\mathfrak{T}(\square_0)\le\mathfrak{l}_j\,\mathfrak{t}^{\mathrm{win}},\qquad
 \Sigma^{\mathrm{sf}}(\square_0)
 \le L^d\mathsf{l}_j^2\,\mathfrak{s}_J\,e^{\frac18\delta_P d^{\wedge}(\square_0,R_n)}
 &&(\square_0\ \text{a window block}),\\
&\sum\bigl\{W(c):\ c\in\mathfrak{F},\ S_c\cap q\neq\emptyset\bigr\}
 \le 2\,\mathfrak{t}^{\mathrm{win}}
 &&(q\ \text{a rim cube})
\end{aligned}
\end{equation}
hold.
\item[(e)] Fix a probe cube $q\in\mathsf{Q}$, and consider the groups $\hat e$ with
$q\in\mathsf{b}(\hat e)$. We write $c$ for the term of $\hat e$ and $A$ for its set of
blocks.
\begin{itemize}
\item Those for which $q$ belongs to the set $\mathcal{Q}^{w}_c$ of rim cubes of $c$ or to
its strong set $\mathcal{S}_j(c)$, so that $S_c$ meets the rim cube $q$, contribute to
$\sum\tilde\nu(\hat e)$ at most
\begin{equation}\label{4.21:eq-one-anchor-sums-1}
2e^{c_J+\frac12}\Bigl(\frac{K_Q}{r_A}+\varepsilon_S\Bigr)
\bigl(2\mathfrak{t}^{\mathrm{win}}+3^d\mathfrak{s}_J\bigr).
\end{equation}
\item Those with $q\subset\hat{\mathsf{K}}_A\cup\bigcup\hat{\mathsf{y}}$ contribute at most
\begin{equation}\label{4.21:eq-one-anchor-sums-2}
4e^{c_J+\frac12}K_2^{\flat}\,\vartheta\,c_{\mathrm{loc}}
\bigl[C^{P,\star}_{\mathsf{a}}+\mathfrak{l}_j\mathfrak{t}^{\mathrm{win}}
+L^d\mathsf{l}_j^2\mathfrak{s}_J\bigr].
\end{equation}
This bound contains no factor $V'$, no position factor and no sum over components.
\end{itemize}
Write $\mu^{\mathrm{kept}}$ for the total of the groups of kept members displayed in
\eqref{4.21:eq-group-totals-a}. Together with these groups one has
$\mu\le\mu^{\flat}+\mu^{\mathrm{kept}}$, where $\mu$ and $\mu^{\flat}$ are as in
\eqref{4.21:eq-group-totals-a}.
\end{enumerate}
\end{lemma}

\begin{proof}
Throughout, $d^{\wedge}$ stands for $d^{\wedge}(\square_0,R_n)$.

\emph{Old blocks.}
Let $\square_0\in\mathfrak{O}^+$. By the total weight at an old block
(Lemma~\ref{4.21:lem-old-block-total}, display \eqref{4.21:eq-old-block-total-a}),
\begin{equation*}
\mathfrak{S}^{\mathrm{all}}(\square_0)\le C^{P,\mathrm{all}}_{\mathsf{a}}\,
e^{\frac18\delta_P d^{\wedge}} .
\end{equation*}

We now bound the small-field species at $\square_0$. The far small-field leaves contribute
at most $A_\flat K_S(E_0+1)$. This is the bound on the far leaves stated with the
constant $\mathfrak{s}_J$ in Notation~\ref{4.21:not-small-field-constant} ($K_S$ and $E_0$
as there).

The near groups contribute at most
$[\mathfrak{n}_0/K_A]\,\mathsf{l}^2(\square_0)(1+\mathfrak{p}_j/\vartheta)$. The local data
factor is $\mathsf{l}^2(\square_0)=L_\circ^{6s(\square_0)}\mathfrak{d}(\square_0)$, by
\eqref{4.21:eq-local-data-factor-a}. On $\mathfrak{O}^+$ the raise read by the block is
$s(\square_0)=0$, by the facts on the raise read by a block
(Lemma~\ref{4.21:lem-block-raise}). Hence $\mathsf{l}^2(\square_0)=\mathfrak{d}(\square_0)$.

The inherited-datum hypothesis \eqref{4.21:eq-datum-hypothesis-a} (proof outstanding) gives
$\mathfrak{d}(\square_0)\le1+C_{\mathrm{dat}}e^{\frac18\delta_P d^{\wedge}}$, where
$C_{\mathrm{dat}}$ denotes the constant of that hypothesis; this is the constant that enters
the definition of $\mathfrak{s}_J$ in \eqref{4.21:eq-small-field-constant-a}. Since the
exponential is at least one, it follows that
\begin{equation*}
\mathfrak{d}(\square_0)\le(1+C_{\mathrm{dat}})\,e^{\frac18\delta_P d^{\wedge}} .
\end{equation*}
So the near groups contribute at most
$(1+C_{\mathrm{dat}})[\mathfrak{n}_0/K_A](1+\mathfrak{p}_j/\vartheta)\,
e^{\frac18\delta_P d^{\wedge}}$.

By the definition of $\mathfrak{s}_J$ in \eqref{4.21:eq-small-field-constant-a}, the two
contributions together are at most $\mathfrak{s}_J e^{\frac18\delta_P d^{\wedge}}$. This is
part~(a).

Adding the bound on $\mathfrak{S}^{\mathrm{all}}(\square_0)$ gives
\begin{equation*}
\mathfrak{S}^{\mathrm{all}}(\square_0)+\Sigma^{\mathrm{sf}}(\square_0)
\le\bigl(C^{P,\mathrm{all}}_{\mathsf{a}}+\mathfrak{s}_J\bigr)e^{\frac18\delta_P d^{\wedge}} .
\end{equation*}
This is the first line of \eqref{4.21:eq-one-anchor-sums-a}, since
$C^{P,\mathrm{all}}_{\mathsf{a}}+\mathfrak{s}_J\le C^{P,\star}_{\mathsf{a}}$ by the definition
of $C^{P,\star}_{\mathsf{a}}$ in \eqref{4.21:eq-group-totals-a}.

\emph{Window blocks.}
Let $\square_0$ be a window block. Part~(b) of the bounds at a window point and
a window block (Lemma~\ref{4.21:lem-window-bounds}, display
\eqref{4.21:eq-window-bounds-a}) gives
$\mathfrak{T}^{\mathrm{co}}(\square_0)\le(5L)^d\mathfrak{t}^{\mathrm{win}}$ and
$\mathfrak{T}(\square_0)\le\mathfrak{l}_j\mathfrak{t}^{\mathrm{win}}$.

We turn to the small-field species at $\square_0$. The far leaves contribute at most
$A_\flat K_S(E_0+1)$, which is at most $\mathfrak{s}_J$. This bound holds per block and
carries no factor $\mathsf{l}^2$, because the far leaves are long by diameter (the bound
stated in Notation~\ref{4.21:not-small-field-constant}).

The near groups contribute at most
\begin{equation*}
[\mathfrak{n}_0/K_A](1+\mathfrak{p}_j/\vartheta)\,L_\circ^{6s(\square_0)}\,
\mathfrak{d}(\square_0) .
\end{equation*}
The factor $L_\circ^{6s(\square_0)}$ is bounded by Lemma~\ref{4.21:lem-block-raise}, and the
factor $\mathfrak{d}(\square_0)$ by the inherited-datum hypothesis
\eqref{4.21:eq-datum-hypothesis-a} (proof outstanding). Together these give at most
$L^d\mathsf{l}_j^2\mathfrak{s}_J e^{\frac18\delta_P d^{\wedge}}$.

By part~(i) of the depth-paid shares of the inherited datum
(Lemma~\ref{4.21:lem-datum-depth-shares}), the exponential factor is needed only at the
exceptional blocks of that lemma.

In total,
$\Sigma^{\mathrm{sf}}(\square_0)\le L^d\mathsf{l}_j^2\mathfrak{s}_J
e^{\frac18\delta_P d^{\wedge}}$. Within this total, the far-leaf part is bounded by
$\mathfrak{s}_J$ and is free of $\mathsf{l}^2$. This proves part~(b) and the second line of
\eqref{4.21:eq-one-anchor-sums-a}.

\emph{Rim cubes.}
Let $q$ be a rim cube. By the geometric description of the running components, $q$ is
contained in $C'\setminus Z_h$ ($Z_h$ as in Lemma~\ref{4.21:lem-under-block}) and meets no
strip. Hence every point of $q$ is a window
point. Let $c\in\mathfrak{F}$ with $S_c\cap q\ne\emptyset$. We sort the terms $c$ by the
level of the cubes of $S_c$ that meet $q$.

\emph{Cubes of level at least $j+1$.} Such a cube meeting $q$ contains $q$, because $q$ is
a cube of $\pi_{j+1}$. Hence $A_c$ contains a point of $q$, and that point is a window point.
By part~(a) of Lemma~\ref{4.21:lem-window-bounds}, the total weight of these terms is at
most $\mathfrak{t}^{\mathrm{win}}$.

\emph{Cubes of level at most $j$.} Such a cube lies inside $q$. We consider three kinds of
terms in turn.
\begin{itemize}
\item The following terms have no cubes of this kind: the terms of this $\mathbf{R}_k$'s
chains filed below level $k$, the local sibling terms, and the confined old terms. By
part~(d) of what lies under a block (Lemma~\ref{4.21:lem-under-block}) together with
\eqref{4.21:eq-window-probe-a}, their level is at least $b=j+1$.
\item Terms of leaf type of level $j=b-1$. There are at most $L^d$ positions of $\pi_j$
in $q$. At each position, the corresponding display and the per-cube bounds for the
format class (Lemma~\ref{4.21:lem-per-cube-bounds}) bound the weight by
$\mathfrak{g}(0)\,\mathfrak{c}\,e^{3-\check{R}_k}$. By the second member of the condition
on levels in \eqref{4.21:eq-weights-incidences-a},
\begin{equation*}
L^d\,\mathfrak{g}(0)\,\mathfrak{c}\,e^{3-\check{R}_k}\le\tfrac14\mathfrak{c} .
\end{equation*}
\item Terms of leaf type of level $\ell=j-s$, with $s\ge1$. There are $L^{d(s+1)}$
positions of $\pi_{j-s}$ in $q$. At each, the corresponding display is read with
$b-1-\ell=s$ and $k-\ell\le\check{R}_k+s$. It bounds the weight by
$\mathfrak{g}(s)\,\mathfrak{c}\,e^{3-38L^s}$. By the definition of $C_{\mathrm{fr}}$ in
\eqref{4.21:eq-weights-incidences-a},
\begin{equation*}
\mathfrak{c}\sum_{s\ge1}L^{d(s+1)}\mathfrak{g}(s)\,e^{3-38L^s}=C_{\mathrm{fr}}\,\mathfrak{c}.
\end{equation*}
\end{itemize}

Adding the three contributions,
\begin{equation*}
\sum\bigl\{W(c):\ c\in\mathfrak{F},\ S_c\cap q\ne\emptyset\bigr\}
\le\mathfrak{t}^{\mathrm{win}}+\bigl(\tfrac14+C_{\mathrm{fr}}\bigr)\mathfrak{c}
\le\mathfrak{t}^{\mathrm{win}}+(3+C_{\mathrm{fr}})\mathfrak{c}
\le2\,\mathfrak{t}^{\mathrm{win}} .
\end{equation*}
The last inequality holds because, by the definition of $\mathfrak{t}^{\mathrm{win}}$ in
\eqref{4.21:eq-weights-incidences-a}, $\mathfrak{t}^{\mathrm{win}}$ contains the summand
$(3+C_{\mathrm{fr}})\mathfrak{c}$ besides non-negative summands. This is the third line of
\eqref{4.21:eq-one-anchor-sums-a}.

For the small-field species at a rim cube we use the rim-cube bound of the plain-unit
expansion. It holds at a zone made by $\mathbf{T}$, because there a label $S$ meeting the
rim cube has level $m(S)\ge j+1$, by the analysis of the preparation site. This bound gives
at most $3^d\mathfrak{s}_J$.

Only far small-field leaves enter here; their total at a block is at most
$A_\flat K_S(E_0+1)\le\mathfrak{s}_J$, which is the per-block input of the rim-cube bound.
The near groups have block set $A=\{\square\}\ne\emptyset$ and are never anchored at a rim
cube. They are counted in the second case below and in the near-group total. This is
part~(c).

\emph{The first total.}
Fix $q\in\mathsf{Q}$, and consider the groups $\hat e$ with $q\in\mathsf{b}(\hat e)$.

\emph{First case: $q$ belongs to the set $\mathcal{Q}^{w}_c$ of rim cubes of $c$ or to its
strong set $\mathcal{S}_j(c)$.} Then $S_c$ meets $q$, and $q$ is a rim cube in the sense
defined above. The anchored forms of the bound
\eqref{4.21:eq-plain-unit-data-a} supply the prefactor
$2e^{c_J}(K_Q/r_A+\varepsilon_S)$ of this case and give
$e^{\Sigma_c}\le e^{\frac12+\mathsf{e}/12}$. By \eqref{4.21:eq-weights-incidences-a},
$\mathfrak{w}\,e^{\mathsf{e}/12}\le W$. Inserting these facts into the third line of
\eqref{4.21:eq-one-anchor-sums-a} and into part~(c), the groups of this case contribute at
most the quantity \eqref{4.21:eq-one-anchor-sums-1}.

\emph{Second case: $q\subset\hat{\mathsf{K}}_A\cup\bigcup\hat{\mathsf{y}}$.} Choose
$\square_0\in A$, and choose $G_0$ as follows.
\begin{itemize}
\item If $q\subset\hat{\mathsf{K}}_{\square_0}$, put $G_0:=q$.
\item Otherwise, let $G_0$ be a member, touching $\hat{\mathsf{K}}_{\square_0}$, of the
animal of $\hat{\mathsf{y}}$ through $q$ (animals of $\hat{\mathsf{y}}$ as in
Notation~\ref{4.21:not-plain-unit-data}).
\end{itemize}

Fix the triple $(c,\square_0,G_0)$. The per-block bound of the plain-unit expansion, carried
in \eqref{4.21:eq-plain-unit-data-a}, bounds the contribution of the triple
$(c,\square_0,G_0)$ to $\tilde\nu(\hat e)$ by
\begin{equation*}
2\,\mathfrak{w}(c)\,e^{c_J}\,y_{\square_0}\,e^{\Sigma_c}\,a(q,G_0).
\end{equation*}
Here $a(q,G_0):=1$ in the first choice of $G_0$. In the second choice,
$a(q,G_0):=\sum\{\mathsf{x}^{|\text{animal}|}\}$, summed over the animals containing $q$ and
$G_0$. In this bound, the other animals and the other blocks are summed by part~(iv) of the
resummation statement carried in Notation~\ref{4.21:not-plain-unit-data}. The product of the
factors $e^{-1}$ over the blocks of $A$ against $e^{|A|}$ is at most one.

Summing over $G_0$,
\begin{equation*}
\sum_{G_0}a(q,G_0)\le1+\sum_{n\ge1}n\,(e\hat b_t\mathsf{x})^n\le2 .
\end{equation*}
Here the summand $1$ accounts for the choice $G_0=q$, and the $n$-th summand for animals with
$n$ members. The last inequality is the smallness condition on $\mathsf{x}$ of
Notation~\ref{4.21:not-plain-unit-data}.

Insert $e^{\Sigma_c}\le e^{\frac12+\mathsf{e}/12}$ and interchange the sums over groups and
blocks. The groups of this case contribute at most
\begin{equation*}
4e^{c_J+\frac12}\sum_{\square_0}\Bigl[\mathbf{1}\{\hat{\mathsf{K}}_{\square_0}
\text{ contains or touches }G_0\}\;y_{\square_0}
\sum_{c:\ \square_0\in\mathcal{A}_c}\mathfrak{w}(c)\,e^{\mathsf{e}(c)/12}\Bigr].
\end{equation*}
The inner sum runs over $\mathfrak{F}$ and the small-field species. Moreover
$\mathfrak{w}\,e^{\mathsf{e}/12}\le W$.

\emph{Old anchors.} Let $\square_0$ be old, island blocks included. By
\eqref{4.21:eq-plain-unit-data-a},
$y_{\square_0}\le K_2^{\flat}\vartheta e^{-\frac14\delta_P d^{\wedge}}$. By the first line of
\eqref{4.21:eq-one-anchor-sums-a}, the inner sum is at most
$C^{P,\star}_{\mathsf{a}}e^{\frac18\delta_P d^{\wedge}}$. The product leaves
$e^{-\frac18\delta_P d^{\wedge}}$.

The blocks around one grain member sum to a constant
(Lemma~\ref{4.21:lem-grain-block-sum}, display \eqref{4.21:eq-grain-block-sum-a}).
Applied at $\beta=\frac18\delta_P\ge\delta_P/16$, it bounds the sum of these factors over
the blocks $\square_0$ with $\hat{\mathsf{K}}_{\square_0}$ containing or touching $G_0$ by
$c_{\mathrm{loc}}$.

\emph{Window anchors.} Let $\square_0$ be a window block with
$\operatorname{Cor}(\square_0)=\emptyset$. By
\eqref{4.21:eq-plain-unit-data-a},
$y_{\square_0}=K_2^{\flat}\vartheta e^{-\frac12\delta_P d^{\wedge}}$. By the second line of
\eqref{4.21:eq-one-anchor-sums-a}, the inner sum is at most
$[\mathfrak{l}_j\mathfrak{t}^{\mathrm{win}}+L^d\mathsf{l}_j^2\mathfrak{s}_J]
e^{\frac18\delta_P d^{\wedge}}$. The product leaves $e^{-\frac38\delta_P d^{\wedge}}$; for
these blocks alone, Lemma~\ref{4.21:lem-grain-block-sum} at $\beta=\frac38\delta_P$ gives at
most $c_{\mathrm{loc}}$.

At an island window block, that is one with $\operatorname{Cor}(\square_0)\ne\emptyset$,
the factor $y_{\square_0}$ carries only $e^{-\frac14\delta_P d^{\wedge}}$, by
\eqref{4.21:eq-plain-unit-data-a}. This leaves
$e^{-\frac18\delta_P d^{\wedge}}$, since $\frac14-\frac18=\frac18\ge\frac1{16}$.
The bound $c_{\mathrm{loc}}$ for the window anchors is unchanged. Indeed, since
$e^{-\frac38\delta_P d^{\wedge}}\le e^{-\frac18\delta_P d^{\wedge}}$, every window anchor,
island or not, leaves a factor at most $e^{-\frac18\delta_P d^{\wedge}}$, and a single
application of Lemma~\ref{4.21:lem-grain-block-sum} at $\beta=\frac18\delta_P\ge\delta_P/16$
bounds the sum over all window anchors by $c_{\mathrm{loc}}$, because $c_{\mathrm{loc}}$
serves every $\beta\ge\delta_P/16$.

Each anchor $\square_0$ is either old or a window block, and each of the two classes of
anchors contributes at most $c_{\mathrm{loc}}$ times its constant. The groups of the second
case therefore contribute at most \eqref{4.21:eq-one-anchor-sums-2}.

No volume factor $V'$, no position factor and no sum over components has entered this
bound.

\emph{Conclusion.}
The bounds \eqref{4.21:eq-one-anchor-sums-1} and \eqref{4.21:eq-one-anchor-sums-2}, added,
give the number $\mu^{\flat}$ of \eqref{4.21:eq-group-totals-a}; by the two cases above it
bounds, for every $q\in\mathsf{Q}$, the total of $\tilde\nu(\hat e)$ over the groups $\hat e$
with $q\in\mathsf{b}(\hat e)$ other than the groups of kept members.
It remains to add the groups of kept members, whose total is
$\mu^{\mathrm{kept}}$ of \eqref{4.21:eq-group-totals-a}. At the number used in
Proposition~\ref{4.21:prop-group-totals}, this total is likewise free of $V'$, by
Lemma~\ref{4.21:lem-grain-block-sum} at $\beta=\frac3{16}\delta_P$. Adding these groups
gives $\mu\le\mu^{\flat}+\mu^{\mathrm{kept}}$.
\end{proof}

\begin{proof}[Proof of Proposition~\ref{4.21:prop-group-totals}, continued: the per-point total of the
groups]\label{4.21:prf-per-point-total}
We bound $\mathsf{M}_{*}$ and then conclude. Throughout, $x$ is a fixed point, and for a group
$\hat e$ of a term $c$ we write $A$ for its set of blocks and $\hat{\mathsf{y}}$ for its family of
grains, so that
\begin{equation*}
\mathsf{d}^{\mathrm{rel}}(\hat e)=S_c\cup\bigcup_{\square\in A}\bigl(\hat{\mathsf{K}}_{\square}\cup
\operatorname{Cor}(\square)\bigr)\cup\bigcup\hat{\mathsf{y}} .
\end{equation*}
We sum $\tilde\nu(\hat e)$ over the groups $\hat e$ with $[\mathsf{d}^{\mathrm{rel}}(\hat e)]_k\ni x$,
that is, over the groups for which $\Delta_k(x)$ meets $\mathsf{d}^{\mathrm{rel}}(\hat e)$. This is
the sum that $\mathsf{M}_{*}$ bounds; it counts every group at least as often as the anchoring of
groups in the proposition on the resummed group expansion does, and we keep this over-count. We
write $J$ for the family of small-field terms whose constant at a block is
$\mathfrak{s}_J$ of \eqref{4.21:eq-small-field-constant-a}.

\emph{Anchoring.} Every group with $A\neq\emptyset$ is anchored at one block $\square_0\in A$, chosen as
follows. If $\Delta_k(x)$ is reached through $\hat{\mathsf{K}}_{\square_0}\cup\operatorname{Cor}(\square_0)$
or through an animal on $\hat{\mathsf{K}}_{\square_0}$, then $\square_0$ is the block realising this
contact, and the anchored form of the group bound among the bounds carried from the plain-unit
expansion \eqref{4.21:eq-plain-unit-data-a} supplies a factor at most $\pi_x(\square_0)$. Otherwise the
contact is through $S_c$, and $\square_0$ is the first block of $A$. In both cases the bound for the
groups of a single term displayed with the group totals in \eqref{4.21:eq-group-totals-a} shows that
the groups of $c$ anchored at $\square_0$ weigh at most
\begin{equation}\label{4.21:eq-per-point-total-1}
2e^{c_J+\frac12}\,y_{\square_0}\,W(c)\,\pi_x(\square_0).
\end{equation}
When we later sum over all $\square_0\in\mathcal{A}_c$ rather than over the single anchor, we only
over-count, and this enlarges the bound. A group with $A=\emptyset$ is anchored at a rim cube $q_0$
meeting $S_c$, with $q_0$ in the fluctuation-cube family $\mathsf{Q}$, respectively
$q_0\in\mathcal{S}_j(c)$; by the
same bound of \eqref{4.21:eq-group-totals-a} its weight is at most
\begin{equation}\label{4.21:eq-per-point-total-2}
2e^{c_J+\frac12}\,\Bigl(\frac{K_Q}{r_A},\ \text{respectively}\ \varepsilon_S\Bigr)\,W(c)\,\pi_x(q_0).
\end{equation}
We now sum over the anchors, in six classes (a)--(f).

\emph{(a) Old blocks.} By the old-block line of the sums at one anchor \eqref{4.21:eq-one-anchor-sums-a},
with the weight $y_{\square_0}$ of an island block (a block of $\mathcal{A}^{\mathrm{isl}}_c$) carrying
the rate $\frac14\delta_P$, and by the position-weighted block sum of
Lemma~\ref{4.21:lem-position-sums} at $\beta=\frac18\delta_P$,
\begin{align*}
&\sum_{\square_0\in\mathfrak{O}^{+}}y_{\square_0}\,\pi_x(\square_0)
\Bigl[\mathfrak{S}^{\mathrm{all}}(\square_0)
+\sum\bigl\{W(c):\ c\in J,\ \square_0\in\mathcal{A}_c\bigr\}\Bigr]\\
&\qquad\le K_2^{\flat}\,\vartheta\,C^{P,\star}_{\mathsf{a}}\sum_{\square_0}
e^{-\frac18\delta_P d^{\wedge}(\square_0,R_n)}\pi_x(\square_0)
\le K_2^{\flat}\,\vartheta\,C^{P,\star}_{\mathsf{a}}\,c_{\star}''\,V'.
\end{align*}
Here $\mathfrak{S}^{\mathrm{all}}(\square_0)$ is the total weight at an old block of
Lemma~\ref{4.21:lem-old-block-total}. The near groups at old blocks are counted here and again in
(e); this over-count only enlarges the bound.

\emph{(b) Window blocks: coarse incidences, and the far small-field terms of $J$ at each block.} By
Lemma~\ref{4.21:lem-window-bounds}\,(b) the coarse term weight at a window block is
$\mathfrak{T}^{\mathrm{co}}(\square_0)\le(5L)^d\mathfrak{t}^{\mathrm{win}}$, and the far small-field
terms of $J$ enter with the constant $\mathfrak{s}_J$ alone, without the per-block factor
$\mathsf{l}^2$, by the clause of the statement on the family $J$ that fixes the constant of its far
terms. At a plain window block the weight $y_{\square_0}$ carries
$e^{-\frac12\delta_P d^{\wedge}}$; at an island block it carries $e^{-\frac14\delta_P d^{\wedge}}$, by
the definition of the weight $y_{\square}$ among the bounds carried from the plain-unit expansion
\eqref{4.21:eq-plain-unit-data-a}. Lemma~\ref{4.21:lem-position-sums} applies to every rate
$\beta\ge\frac18\delta_P$, in particular to $\beta=\frac14\delta_P$ and $\beta=\frac12\delta_P$, with
the same constant $c_{\star}''$. Hence
\begin{equation*}
\sum_{\square_0}K_2^{\flat}\vartheta\,e^{-\frac12\delta_P d^{\wedge}(\square_0,R_n)}\pi_x(\square_0)
\bigl[(5L)^d\mathfrak{t}^{\mathrm{win}}+\mathfrak{s}_J\bigr]
\le K_2^{\flat}\vartheta\,c_{\star}''V'\bigl[(5L)^d\mathfrak{t}^{\mathrm{win}}+\mathfrak{s}_J\bigr],
\end{equation*}
and the island blocks obey the same bound.

\emph{(c) Window blocks: fine incidences.} By the count of fine incidences at window blocks,
Lemma~\ref{4.21:lem-fine-count},
\begin{equation*}
\sum_{(\square_0,c)\ \text{fine}}W(c)\,y_{\square_0}\,\pi_x(\square_0)
\le K_2^{\flat}\,\vartheta\,c_{\mathrm{loc}}\,\mathfrak{F}^{\mathrm{fi}} .
\end{equation*}
At island blocks Lemma~\ref{4.21:lem-fine-count} is applied with $\frac14\delta_P$ in place of
$\frac12\delta_P$. Its proof uses the rate only through the bound on the blocks around one grain
member, Lemma~\ref{4.21:lem-grain-block-sum}, in the form \eqref{4.21:eq-grain-block-sum-a}, and
that lemma holds for every $\beta\ge\delta_P/16$. The constant $c_{\mathrm{loc}}\mathfrak{F}^{\mathrm{fi}}$ is
therefore unchanged.

\emph{(d) The fine small-field terms of $J$} (at the levels $k_0,\dots,j-1$, under window blocks). Each
such term is anchored once, at a cube $q$ of its own. The per-cube constant of the fine
small-field terms of $J$, $A_{\flat}K(d)K_S(E_0+1)$,
carries no factor $\mathsf{l}^2$: $\mathsf{l}^2$ is the per-block multiplicity, and the count in which
each term is taken once replaces it; this is the count of the small-field terms of one component,
used here by statement. The blocks around each cube are summed by
Lemma~\ref{4.21:lem-grain-block-sum}, which gives the factor $c_{\mathrm{loc}}$, a $(d,L)$-number, and
$\pi_x(\square_0)\le2\pi_x(q)$. By the same count the anchors in one component number at most
$4^d(N_0+2)(104L\check{R}^{\uparrow})^d$, and the components are summed by the position bound
\eqref{4.21:eq-fine-count-a}. The result is at most
$K_2^{\flat}\vartheta\,c_{\mathrm{loc}}\,\mathfrak{F}^{\mathrm{fi}}_J$, with
$\mathfrak{F}^{\mathrm{fi}}_J$ as in \eqref{4.21:eq-group-totals-a}. The terms of $J$ that are
anchored in another component are accounted for in the bound for the family $J$ itself.

\emph{(e) Near groups, at all blocks.} A near group at a block $\square$ has
$\mathsf{d}=\hat{\mathsf{K}}_{\square}\cup\bigcup\hat{\mathsf{y}}$, and its bond set $\mathsf{b}$ lies
in its members, by the properties of the grain neighbourhood recalled in
Notation~\ref{4.21:not-plain-unit-data}. Hence
\begin{equation*}
\mathfrak{w}(\hat e)\le(a_k+1)(N^{\flat}+|\hat{\mathsf{y}}|),\qquad
|\mathsf{b}(\hat e)|\le N^{\flat}+|\hat{\mathsf{y}}| ,
\end{equation*}
and, by clause (ii) of the relative small-field bounds recalled in
Notation~\ref{4.21:not-plain-unit-data}, applied to groups of the third class,
$\nu(\hat e)\le\mathfrak{n}(\square)e^{-(\kappa_1-1)|\hat{\mathsf{y}}|}$. Consequently
\begin{equation*}
\tilde\nu(\hat e)\le2\,\mathfrak{n}(\square)\,e^{c_J}\,e^{(a_k+1+c_{\flat})N^{\flat}}\,
\mathsf{x}^{|\hat{\mathsf{y}}|}.
\end{equation*}
The condition $[\mathsf{d}]_k\ni x$ forces one of two things. Either
$\operatorname{dist}_k(x,\square)\le7$, because $\Delta_k(x)$ meets $\hat{\mathsf{K}}_{\square}$, which
lies within $6$ of $\square$ by the same properties of the grain neighbourhood; or an animal on
$\hat{\mathsf{K}}_{\square}$ meets $\Delta_k(x)$. Both clauses of clause (iv) of the relative
small-field bounds, with $|A|=1$, therefore give
\begin{equation*}
\sum_{\hat{\mathsf{y}}}\mathsf{x}^{|\hat{\mathsf{y}}|}\,\mathbf{1}[\,[\mathsf{d}]_k\ni x\,]
\le e\min\bigl\{1,e^{-(\operatorname{dist}_k(x,\square)-8)_+}\bigr\}\le e\,\pi_x(\square).
\end{equation*}
For the second inequality we use, with $t=\operatorname{dist}_k(x,\square)$,
\begin{equation*}
\min\{1,e^{-(t-8)_+}\}\le\min\{1,e^{-(t-204\check{R}_k)/(3N)}\}=\pi_x(\square),
\end{equation*}
the equality being the definition of the position weight in \eqref{4.21:eq-fine-count-a}.
For $t\le204\check{R}_k$ the middle member equals $1$. For $t>204\check{R}_k$ one has
$t-8\ge(t-204\check{R}_k)/(3N)$.

By the definition of $K_2^{\flat}$ in \eqref{4.21:eq-plain-unit-data-a},
$e^{(a_k+1+c_{\flat})N^{\flat}}\cdot e\le K_2^{\flat}/K_A$. The near-group block weight under the
inherited-datum hypothesis, Corollary~\ref{4.21:cor-near-group-weight} in the form
\eqref{4.21:eq-near-group-weight-a}, bounds $\mathfrak{n}(\square)$; we write $C^{\mathrm{dat}}$ for the
constant of the inherited-datum hypothesis \eqref{4.21:eq-datum-hypothesis-a}, which enters
\eqref{4.21:eq-near-group-weight-a} through the factor
$1+C^{\mathrm{dat}}$. Inserting \eqref{4.21:eq-near-group-weight-a}, and then, by the definition
\eqref{4.21:eq-small-field-constant-a} of $\mathfrak{s}_J$, the near groups' block sum,
Lemma~\ref{4.21:lem-near-group-sum}, at the rate $\beta=\frac38\delta_P\ge\frac18\delta_P$, we obtain
\begin{align}
\sum_{\text{near groups}\ni x}\tilde\nu
&\le2e^{c_J}K_2^{\flat}\Bigl[\frac{(1+C^{\mathrm{dat}})\,\mathfrak{n}_0\,(1+\mathfrak{p}_j/\vartheta)}
{K_A}\Bigr]\vartheta\sum_{\square}L_{\circ}^{6s(\square)}
e^{-\frac38\delta_P d^{\wedge}(\square,R_n)}\pi_x(\square)\notag\\
&\le2e^{c_J+\frac12}K_2^{\flat}\vartheta\,\mathfrak{s}_J\Bigl[L^{2d}c_{\star}''V'
+C_{\mathrm{pos}}\bigl(4^d(100L\check{R}^{\uparrow})^d+C_{\mathrm{nr}}\mathsf{l}_j^2\bigr)\Bigr].
\label{4.21:eq-per-point-total-3}
\end{align}
This is the near groups' member of $\mathsf{M}^{\flat}$: $V'$ times a number, plus two added volumes.

\emph{(f) Rim anchors.} By the rim line of \eqref{4.21:eq-one-anchor-sums-a} and the rim-cube sum of
Lemma~\ref{4.21:lem-position-sums},
\begin{equation*}
\sum_{q_0}\pi_x(q_0)\bigl(2\mathfrak{t}^{\mathrm{win}}+3^d\mathfrak{s}_J\bigr)
\le C_{\mathrm{pos}}V'\bigl(2\mathfrak{t}^{\mathrm{win}}+3^d\mathfrak{s}_J\bigr).
\end{equation*}

\emph{Collecting.} We add the bounds (a), (b), (c), (d), (e) and (f). The prefactor
$2e^{c_J+\frac12}$ of \eqref{4.21:eq-per-point-total-1} and \eqref{4.21:eq-per-point-total-2} is
applied to (a), (b), (c), (d) and (f); the bound (e) already carries it. The sum gives
$\mathsf{M}_{*}\le\mathsf{M}^{\flat}$ on $\mathfrak{F}$ and on the family $J$, with $\mathsf{M}^{\flat}$ as in
\eqref{4.21:eq-group-totals-a}. We write $\mathsf{M}^{\mathrm{kept}}$ and $\mu^{\mathrm{kept}}$ for the
two totals of the groups of kept members in \eqref{4.21:eq-group-totals-a}. The groups of kept members
contribute the total $\mathsf{M}^{\mathrm{kept}}$, taken by name from the statement on groups of kept
members, with the value displayed in \eqref{4.21:eq-group-totals-a}: the line (a) with
$e\,\mathfrak{n}^{\mathrm{kept}}e^{\frac1{16}\delta_P d^{\wedge}}$ in place of the term sum. Therefore
\begin{equation}\label{4.21:eq-per-point-total-4}
\mathsf{M}_{*}\le\mathsf{M}^{\flat}+\mathsf{M}^{\mathrm{kept}} .
\end{equation}

\emph{Conclusion.} The proof ends as the concluding step of the plain-unit expansion does, with the
totals just obtained. The proposition on the resummed group expansion reads the elements only through
$(\mathsf{b},\mathsf{d},\nu)$. It is therefore applied to $\hat{\mathfrak{E}}$ together with the groups
of kept members: by clause (v) of the relative small-field bounds recalled in
Notation~\ref{4.21:not-plain-unit-data}, and by the statement on groups of kept members. Its uses of
$\mathsf{d}$ are served by clause (a$'$) of the proposition recalled in
Notation~\ref{4.21:not-plain-unit-data}.

Clause (iii) of the proposition on the resummed group expansion and the coupling-comparison
recursion both require $\mu\le\frac14$. The first of the six conditions on the coupling in
\eqref{4.21:eq-coupling-conditions-a} gives
\begin{equation*}
\mu\le\mu^{\flat}+\mu^{\mathrm{kept}}\le\tfrac14 .
\end{equation*}
Hence $\mu+\tilde{\mathfrak{l}}\le\frac12$, under the bound $\tilde{\mathfrak{l}}\le\frac14$ of the
preparation site. The second condition in \eqref{4.21:eq-coupling-conditions-a}, together with
\eqref{4.21:eq-per-point-total-4}, gives
\begin{equation*}
4\bigl(\mathsf{M}_{*}+2^d c_{\mathsf{p}}V'\tilde{\mathfrak{l}}\bigr)\le\tfrac12 .
\end{equation*}
The outputs are read on $X$ together with the strips. On labels, the two lemmas on labels applied in
the concluding step of the plain-unit expansion apply unchanged.

\emph{Remarks on the bounds.}
\begin{enumerate}
\item[(1)] \emph{What is added and what is multiplied.} $V'$ multiplies only per-point or per-coarse-block
suprema and numbers, namely $\mathfrak{t}^{\mathrm{win}}$, $(5L)^d\mathfrak{t}^{\mathrm{win}}$,
$C^{P,\star}_{\mathsf{a}}$, $\mathfrak{s}_J$, $L^{2d}\mathfrak{s}_J$ and
$\mathfrak{n}^{\mathrm{kept}}$. The fine counts, in which each term is taken once, stand beside $V'$ as
summands and not as factors: the count $(N+1)(100L\check{R}^{\uparrow})^dNC_0^{\sharp}$ inside
$\mathfrak{F}^{\mathrm{fi}}$, the total $\mathfrak{F}^{\mathrm{fi}}_J$, and the terms
$(100L\check{R}^{\uparrow})^d$ and $\mathsf{l}_j^2$ of the near groups. The tile number
$\mathfrak{l}_j$, which is at most $c_{\mathfrak{l}}\mathsf{l}_j^2$ by
Lemma~\ref{4.21:lem-window-bounds}\,(c), multiplies only $c_{\mathrm{loc}}$ in $\mu$. The products
$V'\mathsf{l}_j^2$, $V'\mathfrak{l}_j$, $V'(L\check{R}^{\uparrow})^d$ and $V'\cdot V'$ occur nowhere.
\item[(2)] \emph{Positions.} The position factor is paid in five ways.
\begin{itemize}
\item Terminal terms of $\mathfrak{F}$ and confined terms pay through the factor $\mathsf{f}(c)$.
\item Chain terms that are not confined pay through one third of the running rate $1/N$.
\item Animals pay through clause (iv) of the relative small-field bounds.
\item Near groups pay through the properties of the grain neighbourhood together with that clause.
\item Confined terms pay through the size of their component.
\end{itemize}
Lemma~\ref{4.21:lem-position-sums} and the position bound of \eqref{4.21:eq-fine-count-a} need only
that the components are $\check{R}_k$-sparse and that $N/\check{R}_k\le C_N2^{r_{\mathrm{pr}}}$; they
need finiteness, not smallness.
\item[(3)] \emph{Coarse anchoring of chain terms.} A term whose reading set $S$ meets a component
meeting $\Delta_k(x)$ is anchored at a block at which $S$ is coarse. Every chain term has a coarse
incidence at the zone of its own cubes, which are of a single level. The terms with no coarse
incidence are terminal terms of $\mathfrak{F}$ lying wholly in raised cells, and these are counted
in (c).\qedhere
\end{enumerate}
\end{proof}

\begin{definition}[Six supremum conditions on the coupling]\label{4.21:def-coupling-conditions}
The functions below are functions of a variable $\mathsf{T}\ge\mathsf{T}_{\gamma}$; in the estimates
of this chapter they are evaluated at $\mathsf{T}=\check{\mathsf{T}}_k$, with $k$ the level.
The conditions below are the conditions on the coupling that enter the group totals of
Proposition~\ref{4.21:prop-group-totals}. They are one-sided, and they are written in sum form.
Each of the six is a single inequality of the form
\begin{equation*}
\sup_{\mathsf{T}\ge\mathsf{T}_{\gamma}}\Phi(\mathsf{T})\le c,
\end{equation*}
for a function $\Phi$ of $\mathsf{T}$ and a number $c$.

\emph{An inequality that is not assumed.} The inequality
\begin{equation*}
e^{-(\kappa_1-1)}\,\bigl(10\,\check{R}_{k_0+1}/L\bigr)^{d-1}\le\mathrm{const}
\end{equation*}
is not among the conditions of this definition. No statement of this chapter uses it, and
neither does the proof of the correlation theorem.

\emph{Upper functions of the scaffold functions.} The counting quantities of
Notation~\ref{4.21:not-weights-incidences} and Definition~\ref{4.21:def-strips-geometry} (the
scaffold functions) are bounded above by the following functions of $\mathsf{T}$ (their upper
functions):
\begin{equation}\label{4.21:eq-coupling-conditions-1}
\begin{aligned}
N&\le C_N(\mathsf{T}+1)^{r_{\mathrm{pr}}}=:N(\mathsf{T}),
 \qquad N_0\le\mathsf{N}_0(\mathsf{T}),
 \qquad \check{R}^{\uparrow}\le\mathsf{R}^{\uparrow}(\mathsf{T}),\\
V'&\le\bar V'(\mathsf{T}):=C'_V\bigl(104L\,\mathsf{R}^{\uparrow}(\mathsf{T})\bigr)^d,
 \qquad \mathfrak{l}_j\le\mathfrak{l}(\mathsf{T}):=3\bigl(5L^3\mathsf{R}^{\uparrow}(\mathsf{T})\bigr)^d,\\
\mathsf{l}_j^2&\le\mathsf{l}^2(\mathsf{T}):=16\,\mathsf{l}(x)^2\ \text{ at } x=e^{\mathsf{T}},
 \qquad N^{\mathrm{fr}}(\mathsf{T}):=C_N\bigl(\mathsf{T}+2+c_{N(\mathsf{T})+1}\bigr)^{r_{\mathrm{pr}}}.
\end{aligned}
\end{equation}
Here $\mathsf{l}(x)$ is the data-size function behind the constants $\mathsf{l}_j$ of
Lemma~\ref{4.21:lem-window-bounds}\textup{(c)}, and $r$ is the exponent fixed with that function;
$\mathsf{l}^2(\mathsf{T})$ is of degree $4r$ up to an arbitrarily small positive excess.
The bound $\check{R}^{\uparrow}\le\mathsf{R}^{\uparrow}(\mathsf{T})$ is the bound of
\eqref{4.21:eq-strips-geometry-a}. Every factor $(L\check{R}^{\uparrow})^d$ in this chapter is
evaluated at this bound. The quantities $\mathfrak{t}^{\mathrm{win}}(\mathsf{T})$,
$\mathfrak{F}^{\mathrm{fi}}(\mathsf{T})$ and $\mathfrak{F}^{\mathrm{fi}}_J(\mathsf{T})$ are defined
by their displays: \eqref{4.21:eq-weights-incidences-a}, the display of
Lemma~\ref{4.21:lem-fine-count}, and \eqref{4.21:eq-group-totals-a} respectively. In each of them
the scaffold functions are replaced by the upper functions in
\eqref{4.21:eq-coupling-conditions-1}. Finally, the separation radius $\check{R}_k$ of
Definition~\ref{4.21:def-strips-geometry} is read at the lower bound
$\check{R}_k\ge\mathsf{T}^{r_{\mathrm{pr}}}$; this relation is one-sided.

\emph{Paying functions.} The paying functions are
\begin{equation}\label{4.21:eq-coupling-conditions-2}
\vartheta(\mathsf{T})=K_{\vartheta}\,\frac{A_0}{C_*}\Bigl(\tfrac12\mathsf{T}\Bigr)^{p_1-q}
\quad(q\ge q_*),
\qquad \frac{1}{r_A(\mathsf{T})},
\qquad \varepsilon_S(\mathsf{T})=2e^{-(\kappa_A(\mathsf{T})-c_{\flat})},
\end{equation}
together with $\tilde{\mathfrak{l}}(\mathsf{T})$ and $\mathfrak{p}_j(\mathsf{T})$. The function
$\varepsilon_S$ is as fixed with the weight $\kappa_A$ of the strong set, and
$\tilde{\mathfrak{l}}$, $\mathfrak{p}_j$ are as fixed at the preparation site of the step. The
constants $K_{\vartheta}$, $C_*$, $p_1$, $q_*$ in $\vartheta$, and the constant $c_R$ in
\textup{(q4)} below, are the constants with which the block weight $\vartheta$ and its condition
are fixed at the preparation site of the step. The constant $C'_V$ is the constant in the bound
on the volume sum $V'$ carried from the plain-unit expansion, \eqref{4.21:eq-plain-unit-data-a};
the constant $c_{\mathrm{loc}}$ is that of Lemma~\ref{4.21:lem-fine-count}. The totals of the kept
members, repeated here from the statement in which they are defined, are
\begin{equation}\label{4.21:eq-coupling-conditions-3}
\begin{aligned}
\mu^{\mathrm{kept}}(\mathsf{T})&:=4e^{c_J+3/2}K_2^{\flat}\,\vartheta(\mathsf{T})\,
 c_{\mathrm{loc}}\,\mathfrak{n}^{\mathrm{kept}},\\
\mathsf{M}^{\mathrm{kept}}(\mathsf{T})&:=2e^{c_J+3/2}K_2^{\flat}\,\vartheta(\mathsf{T})\,
 c_{\star}''\,\bar V'(\mathsf{T})\,\mathfrak{n}^{\mathrm{kept}}.
\end{aligned}
\end{equation}

\emph{Unit sizes.} For an expression $\Phi$ built from the quantities of this section, we write
$\Phi|_1$ for the same expression in which the size constants listed below are replaced by $1$,
every count, every scaffold function and every numerical factor being kept. The size constants
so replaced are the following:
\begin{itemize}
\item the product $NC_0^{\sharp}$ carried by the chain terms, where it occurs in
$\mathfrak{t}^{\mathrm{win}}$, in the first member of $\mathfrak{F}^{\mathrm{fi}}$, and as
$N(g_k)C_0^{\sharp}$ in $b_0$;
\item $C_0^{\sharp}$ inside the factor $c_{\mathrm{att}}A_{\flat}C_0^{\sharp}N^{\mathrm{fr}}$ of
$\mathfrak{c}$, and inside $b_1$, for the frozen terms;
\item $B_0$ and $b_{\delta}$;
\item the factor $E_0+1$ inside $\mathfrak{s}_J$;
\item $C_{\Sigma}$ inside $\mathfrak{n}^{\mathrm{kept}}$, so that
$\mathfrak{n}^{\mathrm{kept}}|_1=2^d\mathfrak{s}_{\Sigma}$.
\end{itemize}
Write $C_{\mathrm{inh}}$ for the constant of the inherited-datum hypothesis. The factor
$1+C_{\mathrm{inh}}$ is kept inside $\mathfrak{s}_J|_1$: it enters through the inherited-datum
hypothesis, which is among the hypotheses of Proposition~\ref{4.21:prop-group-totals} and is used
there in the bound on the difference of the near groups between the two machines $\mathsf{A}$ and
$\mathsf{B}$. The unit sizes $\mathsf{a}_{\mu}$, $\mathsf{a}_{\mathsf{M}}$,
$\mathsf{a}^{\mathrm{kept}}_{\mu}$, $\mathsf{a}^{\mathrm{kept}}_{\mathsf{M}}$ of the totals are
defined in the first line of \eqref{4.21:eq-coupling-conditions-a} below.

\emph{The six conditions.} The \emph{six supremum conditions on the coupling} are the
inequalities \textup{(q1)}--\textup{(q6)} below. Each is required for all
$\mathsf{T}\ge\mathsf{T}_{\gamma}$; the first line of the display defines the unit sizes that
enter \textup{(q6)}:
\begin{equation}\label{4.21:eq-coupling-conditions-a}
\begin{aligned}
&\mathsf{a}_{\mu}:=\mu^{\flat}|_1,\qquad
\mathsf{a}_{\mathsf{M}}:=\mathsf{M}^{\flat}|_1,\qquad
\mathsf{a}^{\mathrm{kept}}_{\mu}:=\mu^{\mathrm{kept}}|_1,\qquad
\mathsf{a}^{\mathrm{kept}}_{\mathsf{M}}:=\mathsf{M}^{\mathrm{kept}}|_1 ;\\
&\textup{(q1)}\quad \mu^{\flat}(\mathsf{T})+\mu^{\mathrm{kept}}(\mathsf{T})\le\tfrac14;\\
&\textup{(q2)}\quad 4\bigl(\mathsf{M}^{\flat}(\mathsf{T})+\mathsf{M}^{\mathrm{kept}}(\mathsf{T})
 +2^dc_{\mathsf{p}}\bar V'(\mathsf{T})\tilde{\mathfrak{l}}(\mathsf{T})\bigr)\le\tfrac12;\\
&\textup{(q3)}\quad 6N(\mathsf{T})\,c_{\mathsf{p}}\bar V'(\mathsf{T})
 \bigl(K_2^{\flat}\vartheta(\mathsf{T})+K_Q/r_A(\mathsf{T})\bigr)\le1;\\
&\textup{(q4)}\quad \vartheta\le c_R/2,\quad r_A\ge2,\quad
 \kappa_A(\mathsf{T})\ge c_{\flat}+c_J+1,\quad \check{R}\ge2,\quad \tilde{\mathfrak{l}}\le\tfrac14,\\
&\hphantom{\textup{(q4)}\quad} N(\mathsf{T})\,e^{-\sigma_*\mathsf{T}^{r_{\mathrm{pr}}}/L}\le1,\qquad
 \sup_{\mathsf{T}}\,\mathfrak{p}_j(\mathsf{T})/\vartheta(\mathsf{T})\le1;\\
&\textup{(q5)}\quad \text{the level condition in \eqref{4.21:eq-weights-incidences-a}};\\
&\textup{(q6)}\quad \varepsilon^{N}_{\flat,A}(\mathsf{T}):=N(\mathsf{T})\Bigl[
 8\bigl(\mathsf{a}_{\mathsf{M}}+\mathsf{a}^{\mathrm{kept}}_{\mathsf{M}}\bigr)
 +4\bigl(\mathsf{a}_{\mu}+\mathsf{a}^{\mathrm{kept}}_{\mu}\bigr)\Bigr](\mathsf{T})\le\tfrac14 .
\end{aligned}
\end{equation}
The first five inequalities of \textup{(q4)} are the conditions of the preparation site. The
inequality $N(\mathsf{T})e^{-\sigma_*\mathsf{T}^{r_{\mathrm{pr}}}/L}\le1$ is also a condition of the
preparation site, and it is used in the proof of the correlation theorem. The last inequality of
\textup{(q4)} is of the same form as the others, namely a supremum of a function of $\mathsf{T}$:
$\mathfrak{p}_j$ is proportional to $g_j$ times a power of a logarithm, so it tends to zero faster
than $\vartheta$. Condition \textup{(q5)} is the level condition of
Notation~\ref{4.21:not-weights-incidences}. Condition \textup{(q6)} is the hypothesis of the
coupling-comparison recursion of the step. It is written in sum form and includes the groups of
the kept members.
\end{definition}

\begin{definition}[The difference hypothesis for kept members. Proof outstanding]\label{4.21:def-kept-difference}
Let $\mathsf{A}$ and $\mathsf{B}$ be two machines built on one and the same scaffold, in the setting
of the preparation site of this chapter, where machines and scaffolds are defined. For a kept
member $v$ of a live pocket $Y$, write $v^{\mathsf{A}}$ and $v^{\mathsf{B}}$ for its realisations under
$\mathsf{A}$ and $\mathsf{B}$, and put
\begin{equation*}
  \delta v := v^{\mathsf{B}} - v^{\mathsf{A}} .
\end{equation*}
Here $\mathfrak{O}(\delta v)$ denotes the oscillation $\mathfrak{O}^{\mathcal{K}}(\delta v)$ of
$\delta v$ over the paired chart system $\mathcal{K}^{\sharp}(Y^{\sharp})$ of the strip
$Y^{\sharp}$ of $Y$. The constant $\mathsf{l}^{\mathfrak{K}} = \mathsf{l}^{\mathfrak{K}}(k)$ is given
with the hypothesis at step $k$; it is the size constant of the differences of kept members in
\eqref{4.21:eq-kept-difference-a}. Further, $\mathcal{A}_v$ is the anchor set of $v$, the set of
blocks of the cover incident to $v$, and
$\mathfrak{O}^{+}$ is the set of old blocks of
Definition~\ref{4.21:def-old-window-blocks}; the letter $\mathfrak{O}$ with an argument and no
superscript always denotes an oscillation. Finally,
$n$ denotes the current stage of $\mathbf{R}_k$.

The \emph{difference hypothesis for kept members} is the following statement. For every pocket
$Y$ of level $m$, written $(Y,m)$, that is live at $\mathbf{R}_k$, and for every kept
member $v\in\mathfrak{K}(Y)$, the difference $\delta v$ obeys clauses (a)--(d) of the kept-member
bounds with
the oscillation $\mathfrak{O}(v)$ replaced by $\mathfrak{O}(\delta v)$, where $\mathfrak{O}(\delta v)$
satisfies the first bound below; consequently the second bound below holds at every
$\square_0\in\mathfrak{O}^{+}$:
\begin{equation}\label{4.21:eq-kept-difference-a}
\begin{aligned}
  &\mathfrak{O}(\delta v) \le \mathsf{l}^{\mathfrak{K}}\,\check{R}_m^{\,p_\Sigma},\\
  &\sum_{v:\ \square_0\in\mathcal{A}_v} \mathfrak{O}(\delta v)
    \le e\cdot 2^{d}\,\mathfrak{s}_\Sigma\,\mathsf{l}^{\mathfrak{K}}\,
    e^{\frac{1}{16}\,\delta_P\, d^{\wedge}(\square_0,R_n)} .
\end{aligned}
\end{equation}
In the second line, $e$ is Euler's number and $d=4$. The rate $\delta_P$ is the decay rate of the
kept-member bounds. The quantity $d^{\wedge}(\square_0,R_n)$ is the scaled distance from
$\square_0$ to the collar region $R_n$ at the current stage $n$ of $\mathbf{R}_k$. The constants
$\mathfrak{s}_\Sigma$ and $p_\Sigma$ are those of the kept-member bounds, and $\check{R}_m$ is the
region size parameter at level $m$.
\end{definition}

Proof outstanding. What is to be proved is that the kept members produced by the construction of
this chapter satisfy the difference hypothesis.

\noindent\textit{Route.} A proof has to establish the kept-member bounds (a)--(d) for the
difference $\delta v$ of two machines on one scaffold, with the oscillation taken over the paired
chart system $\mathcal{K}^{\sharp}(Y^{\sharp})$. It must also establish the bound in the first line
of \eqref{4.21:eq-kept-difference-a} with the constant $\mathsf{l}^{\mathfrak{K}}(k)$; the second
line then follows. The input from which such a proof would proceed is a version for differences of
the base estimates on which the kept-member bounds rest. The proofs of the base estimates are
linear in the stored terms, so they are to be applied to $v^{\mathsf{B}}-v^{\mathsf{A}}$ in place of $v$.

\begin{lemma}[Differences of two machines on one scaffold]\label{4.21:lem-machine-differences}
Work at stage $n+1$ of $\mathbf{R}_k$. Let $\mathsf{A}$ and $\mathsf{B}$ be two machines on one scaffold, and for every quantity $v$ computed by both write $\delta v:=v^{\mathsf{B}}-v^{\mathsf{A}}$. For every stage $i$ of $\mathbf{R}_k$ put
\[
\mathsf{s}_i:=\|\delta\mathbb{C}^{(i)}\|^{\mathcal{K},A}_{(i)},
\]
the size of the difference of the chain terms of stage $i$. Let $\nu^m$ be the weights of the coupling-comparison recursion: they assign to every term $c$ the number $\sup|\delta c|$ in place of its size $\mathfrak{N}(c)$. Define $\mathsf{l}_{n+1}$ as the largest of the following three quantities:
\begin{itemize}
\item[(a)] the per-cube constants of the differences of the stored terms, one for each kind of stored term entering the displays \eqref{4.21:eq-per-cube-bounds-a} of Lemma~\ref{4.21:lem-per-cube-bounds};
\item[(b)] the difference sizes of the small-field terms of the correlation theorem that stand, inside $\mathfrak{s}_J$, in the places of the factor $E_0+1$, the size factor of those terms in $\mathfrak{s}_J$; these are the difference size of the near groups, which multiplies $\mathfrak{n}_0|_1$, and the difference size of the far terms, which multiplies the far-term constant $K_S$ alone, the size factor $E_0+1$ being read as $1$ there;
\item[(c)] the constant $\mathsf{l}^{\mathfrak{K}}$ of the difference hypothesis for kept members \eqref{4.21:eq-kept-difference-a} (hypothesis (ii) below, whose proof is outstanding).
\end{itemize}
In the coupling-comparison recursion the constant $\mathsf{l}_{n+1}$ is defined through sums over the terms $T$ of the step in which the bound $c_Te^{-\kappa_Td_\ell}$ of a term is replaced by $\sup|\delta T-\delta T(\mathbf{1})|$, $\mathbf{1}$ the all-ones parameter point; those sums are not formed here, and $\mathsf{l}_{n+1}$ as just defined takes their place in the recursion's display.
Here and below $N$ denotes the bound on the number of stages of $\mathbf{R}_k$ that enters \eqref{4.21:eq-coupling-conditions-a}, not the rank of the gauge group.
For every stage $n'$ of $\mathbf{R}_k$ let $\mathsf{l}_{n'}$ denote the quantity so defined at stage $n'$. Assume:
\begin{itemize}
\item[(i)] the hypotheses of Proposition~\ref{4.21:prop-group-totals}, among them the members (q1), (q2) and (q6) of the six supremum conditions on the coupling \eqref{4.21:eq-coupling-conditions-a} (Definition~\ref{4.21:def-coupling-conditions}) and the inherited-datum hypothesis \eqref{4.21:eq-datum-hypothesis-a} at members (3) and (4), the latter read for differences at member (4);
\item[(ii)] the difference hypothesis for kept members \eqref{4.21:eq-kept-difference-a}, stated in Definition~\ref{4.21:def-kept-difference}, whose proof is outstanding;
\item[(iii)] the coupling-comparison recursion of the perturbative-site step, with its weights $\nu^m$: at every stage $n'+1$ of $\mathbf{R}_k$ it bounds $\mathsf{s}_{n'+1}$ by any bound obtained at that stage for the difference $\delta\mathbb{C}_I$ of the terms $\mathbb{C}_I$ formed at that stage, in the norm of the recursion, plus the remainder $\mathsf{z}_{n'+1}$ of the recursion at that stage; and, if the display \eqref{4.21:eq-machine-differences-a} holds at every stage of $\mathbf{R}_k$ and $\varepsilon^N_{\flat,A}\le\frac14$, it yields
\[
\max_n\mathsf{s}_n\le 2\max_n\Bigl(\frac{\varepsilon^N_{\flat,A}}{N}\,\mathsf{l}_n+\mathsf{z}_n\Bigr);
\]
\item[(iv)] the third clause of the bound, by the group totals, of the terms $\mathbb{C}_I$ formed at stage $n+1$: if the group total $\mu$ formed at the sizes $\mathfrak{N}$ satisfies $\mu\le\frac14$, then
\[
\|\delta\mathbb{C}_I\|\le 8\,\mathsf{M}_*[\nu^m]+32\,\mu[\nu^m]\bigl(\mathsf{M}_*+2^dc_{\mathsf{p}}V'\tilde{\mathfrak{l}}\bigr),
\]
where $\mu[\nu^m]$ and $\mathsf{M}_*[\nu^m]$ denote the group totals of \eqref{4.21:eq-group-totals-a} formed at the weights $\nu^m$, the unbracketed totals $\mu$, $\mathsf{M}_*$ being formed at the sizes $\mathfrak{N}$, and the norm is that of the coupling-comparison recursion at stage $n+1$.
\end{itemize}
Then, with $\varepsilon^N_{\flat,A}$ as in \eqref{4.21:eq-coupling-conditions-a},
\begin{equation}\label{4.21:eq-machine-differences-a}
\mathsf{s}_{n+1}\le\frac{\varepsilon^N_{\flat,A}}{N}\Bigl[\sum_{i\le n}\mathsf{s}_i+\mathsf{l}_{n+1}\Bigr]+\mathsf{z}_{n+1},
\end{equation}
and, if hypotheses (i)--(iv) hold at every stage of $\mathbf{R}_k$ and member (q6) of \eqref{4.21:eq-coupling-conditions-a}, that is $\varepsilon^N_{\flat,A}\le\frac14$, holds, then
\begin{equation}\label{4.21:eq-machine-differences-1}
\max_n\mathsf{s}_n\le 2\max_n\Bigl(\frac{\varepsilon^N_{\flat,A}}{N}\,\mathsf{l}_n+\mathsf{z}_n\Bigr).
\end{equation}
\end{lemma}

\begin{proof}
\emph{The input bounds at the weights $\nu^m$.} The steps of the proof of Proposition~\ref{4.21:prop-group-totals} that form the group totals use Lemma~\ref{4.21:lem-per-cube-bounds}, the window bounds of Lemma~\ref{4.21:lem-window-bounds}, the count of fine incidences of Lemma~\ref{4.21:lem-fine-count}, the total weight at an old block of Lemma~\ref{4.21:lem-old-block-total}, the near groups' block sum of Lemma~\ref{4.21:lem-near-group-sum}, the grain block sum of Lemma~\ref{4.21:lem-grain-block-sum}, and the two further bounds those steps invoke. Each of these steps and lemmas bounds a sum of weights of terms, is linear in those weights, and anchors the terms through the scaffold only. The scaffold is the same for $\mathsf{A}$ and $\mathsf{B}$, so every such estimate applies verbatim to the weights $\nu^m(c)=\sup|\delta c|$ in place of the sizes $\mathfrak{N}(c)$, with each size constant replaced by the corresponding difference constant. Each input bound is then read as follows.

The bound for the chain terms, which at the sizes $\mathfrak{N}$ gives at most $c_{\mathrm{sib}}NC_0^{\sharp}$ from \eqref{4.21:eq-per-cube-bounds-a}, gives at the weights $\nu^m$, in place of $c_{\mathrm{sib}}NC_0^{\sharp}$, the quantity
\[
c_{\mathrm{sib}}\sum_{i\le n}\|\delta\mathbb{C}^{(i)}\|^{\mathcal{K},A}_{(i)}=c_{\mathrm{sib}}\sum_{i\le n}\mathsf{s}_i .
\]
The bound for every other kind of stored term gives that kind's difference size.

The linearity just invoked is a property of the lemmas listed above. It is not asserted of the near-group total of the correlation theorem. That bound is read instead through the near-group block weight \eqref{4.21:eq-near-group-weight-a} of Corollary~\ref{4.21:cor-near-group-weight} at member (4). This reading is available under the inherited-datum hypothesis \eqref{4.21:eq-datum-hypothesis-a} at members (3) and (4) of hypothesis (i). At member (4) the constant $\mathfrak{n}_0$ is replaced by $\mathfrak{n}_0|_1$ times the difference size of the near groups, which is one of the quantities in (b).

The bound for the kept members is read from the difference hypothesis for kept members \eqref{4.21:eq-kept-difference-a}, hypothesis (ii), whose proof is outstanding. Its constant is $\mathsf{l}^{\mathfrak{K}}$, the quantity (c).

\emph{The constant $\mathsf{l}_{n+1}$.} This constant is the one defined in the statement in place of the recursion's sums. By definition it is the largest of the quantities (a), (b) and (c). These are exactly the size constants occupying, in the bounds just listed, the places of the constants of the stored terms in the per-cube displays, of the factor $E_0+1$ inside $\mathfrak{s}_J$, and of the kept members. Hence every bound for stored terms holds for the weights $\nu^m$ with its size constant replaced by $\mathsf{l}_{n+1}$. The bound for the chain terms holds with its constant replaced by $\sum_{i\le n}\mathsf{s}_i$.

\emph{The group totals at the weights $\nu^m$.} Run the derivation of the group totals of Proposition~\ref{4.21:prop-group-totals}, in the sum form \eqref{4.21:eq-group-totals-a}, with the weights $\nu^m$. Let $\mathsf{a}_\mu,\mathsf{a}_{\mathsf{M}},\mathsf{a}^{\mathrm{kept}}_\mu,\mathsf{a}^{\mathrm{kept}}_{\mathsf{M}}$ be the coefficients of the group totals \eqref{4.21:eq-group-totals-a}, the superscript marking the element formed by the kept members; through these coefficients $\varepsilon^N_{\flat,A}$ is defined in \eqref{4.21:eq-coupling-conditions-a}. The derivation is linear in the input bounds, and each input constant is now $\sum_{i\le n}\mathsf{s}_i$ or $\mathsf{l}_{n+1}$. Therefore
\begin{align*}
\mu[\nu^m]&\le(\mathsf{a}_\mu+\mathsf{a}^{\mathrm{kept}}_\mu)\max\Bigl\{\sum_{i\le n}\mathsf{s}_i,\ \mathsf{l}_{n+1}\Bigr\}
\le(\mathsf{a}_\mu+\mathsf{a}^{\mathrm{kept}}_\mu)\Bigl(\sum_{i\le n}\mathsf{s}_i+\mathsf{l}_{n+1}\Bigr),\\
\mathsf{M}_*[\nu^m]&\le(\mathsf{a}_{\mathsf{M}}+\mathsf{a}^{\mathrm{kept}}_{\mathsf{M}})\Bigl(\sum_{i\le n}\mathsf{s}_i+\mathsf{l}_{n+1}\Bigr).
\end{align*}
The second inequality in the first line holds because both quantities are non-negative, so their maximum is at most their sum.

\emph{The difference $\delta\mathbb{C}_I$.} Condition (q1) of \eqref{4.21:eq-coupling-conditions-a} gives $\mu\le\frac14$ for the group total at the sizes $\mathfrak{N}$. This is the hypothesis of the bound in hypothesis (iv). That bound therefore bounds the difference of the terms $\mathbb{C}_I$ formed at stage $n+1$, in the norm of the coupling-comparison recursion at that stage, by
\[
\|\delta\mathbb{C}_I\|\le 8\,\mathsf{M}_*[\nu^m]+32\,\mu[\nu^m]\bigl(\mathsf{M}_*+2^dc_{\mathsf{p}}V'\tilde{\mathfrak{l}}\bigr).
\]
Condition (q2) of \eqref{4.21:eq-coupling-conditions-a} gives $\mathsf{M}_*+2^dc_{\mathsf{p}}V'\tilde{\mathfrak{l}}\le\frac18$, so the second term is at most $4\mu[\nu^m]$. Inserting the bounds on $\mu[\nu^m]$ and $\mathsf{M}_*[\nu^m]$ obtained above gives
\begin{align*}
\|\delta\mathbb{C}_I\|
&\le\bigl[8(\mathsf{a}_{\mathsf{M}}+\mathsf{a}^{\mathrm{kept}}_{\mathsf{M}})
+4(\mathsf{a}_\mu+\mathsf{a}^{\mathrm{kept}}_\mu)\bigr]\Bigl(\sum_{i\le n}\mathsf{s}_i+\mathsf{l}_{n+1}\Bigr)\\
&=\frac{\varepsilon^N_{\flat,A}}{N}\Bigl(\sum_{i\le n}\mathsf{s}_i+\mathsf{l}_{n+1}\Bigr).
\end{align*}
The last equality is the definition of $\varepsilon^N_{\flat,A}$ in \eqref{4.21:eq-coupling-conditions-a}.

\emph{Conclusion.} The bound on $\|\delta\mathbb{C}_I\|$ just obtained is the input that the display of the coupling-comparison recursion, hypothesis (iii), requires at stage $n+1$. With this input that display holds in the form \eqref{4.21:eq-machine-differences-a}. If hypotheses (i)--(iv) hold at every stage of $\mathbf{R}_k$, the argument above, carried out at each stage, gives \eqref{4.21:eq-machine-differences-a} at every stage. Under member (q6) of \eqref{4.21:eq-coupling-conditions-a}, that is $\varepsilon^N_{\flat,A}\le\frac14$, hypothesis (iii) then yields \eqref{4.21:eq-machine-differences-1}. Thus member (q6) of \eqref{4.21:eq-coupling-conditions-a} is precisely the smallness hypothesis under which the coupling-comparison recursion yields \eqref{4.21:eq-machine-differences-1}.

\emph{The role of hypothesis (ii).} The difference hypothesis for kept members enters the argument only through the quantity (c) and through the reading of the kept members' bound. Without it the derivation above covers only the elements $\hat{\mathfrak{E}}$ on which the group totals are formed. The elements over which the second stage of the expansion runs are, however, $\hat{\mathfrak{E}}\cup\{\hat e\}$, where $\hat e$ is the element formed by the kept members; hypothesis (ii) supplies the difference size (c) for that element.
\end{proof}

\begin{proposition}[The coupling conditions hold at small coupling]\label{4.21:prop-small-coupling}
Write $\mathsf{T}$ for the variable $\check{\mathsf{T}}_k$, which ranges over
$[\mathsf{T}_\gamma,\infty)$. Write each of the six members of the supremum conditions on
the coupling (Definition~\ref{4.21:def-coupling-conditions}) as one inequality, with the member's
constant $c>0$,
\begin{equation}\label{4.21:eq-small-coupling-1}
\sup_{\mathsf{T}\ge\mathsf{T}_\gamma}\Phi(\mathsf{T})\le c .
\end{equation}
The functions entering $\Phi$ are read at their upper functions of $\mathsf{T}$, as in
\eqref{4.21:eq-coupling-conditions-a}. The volume-type functions in particular are read at the
bound $\check{R}^{\uparrow}\le\mathsf{R}^{\uparrow}(\mathsf{T})$ fixed there.

Let $q_*$ denote the lower bound on Ba{\l}aban's exponent $q$ that is fixed in the order of choice
of Ba{\l}aban's auxiliary parameters $\mathfrak{p}$ (Notation~\ref{2.1:not-prefix}), so that
$q\ge q_*$. Assume the lower bound on $q_*$ fixed in that order of choice,
\begin{equation}\label{4.21:eq-small-coupling-2}
q_*\ge p_0+(1+d+d\beta_0^{\mathrm{pr}})\,r_{\mathrm{pr}}+1 ,
\end{equation}
in which $r_{\mathrm{pr}}$ is the exponent in the stage-count bound
$N\le C_N(\mathsf{T}+1)^{r_{\mathrm{pr}}}$ of \eqref{4.21:eq-coupling-conditions-a}, and
$\beta_0^{\mathrm{pr}}$ is an exponent fixed in the same order of choice before $q_*$. At $d=4$
the bound \eqref{4.21:eq-small-coupling-2} reads $q_*\ge p_0+7r+1$. Assume also $p_1<p_0$, as in
\cite{B-CMP122a}. No other lower bound on $q$ than $q\ge q_*$ is assumed.

The \emph{degree} of a function of $\mathsf{T}$ is the exponent of $\mathsf{T}$ in its upper
function. The degree $a+0^{+}$ stands for the exponent $a$ together with powers of
$\log\mathsf{T}$. Then the following hold.
\begin{enumerate}
\item[(a)] Each $\Phi$ is a finite sum of terms. Each term is the product of a number fixed before
$\gamma$ and a function of $\mathsf{T}$ alone. Moreover $\Phi(\mathsf{T})\to0$ as
$\mathsf{T}\to\infty$.
\item[(b)] Hence the $\gamma$ at which all six members hold form an interval
$\mathopen{]}0,\gamma_P]$ with
$\gamma_P>0$. The six members together therefore require only that the coupling be small enough.
No member bounds $g$, $g_k$ or $\gamma$ from below. No member bounds $L$, $K$, $N$, $N_0$,
$\check{R}$ or an age from above.
\item[(c)] At the reading $\check{R}^{\uparrow}\le\mathsf{R}^{\uparrow}(\mathsf{T})$ the degrees
are as follows.
\begin{itemize}
\item $\bar V'$, $\mathfrak{l}$, $\mathsf{l}^{2}$ and $(L\mathsf{R}^{\uparrow})^d$ have degree
$4r+0^{+}$.
\item $N$, $N^{\mathrm{fr}}$, $\mathfrak{c}$ and $\mathfrak{t}^{\mathrm{win}}$ have degree
$r+0^{+}$.
\item $\mathfrak{c}^{\circ}$ has degree $0$.
\item $\mathfrak{F}^{\mathrm{fi}}$ has degree $6r+0^{+}$.
  \begin{itemize}
  \item Its chain part $N(N+1)(100L\mathsf{R}^{\uparrow})^d$ has this degree.
  \item Its other part, the second summand of \eqref{4.21:eq-fine-count-a}, has degree
  $5r+0^{+}$: there $(L\mathsf{R}^{\uparrow})^d$ is multiplied, apart from constants of degree $0$,
  by
  \begin{equation*}
  (N+1)(103L)^d\mathfrak{c}^{\circ}+[(103L)^d+C_{\mathrm{old}}]\mathfrak{c},
  \end{equation*}
  which has degree $r+0^{+}$ (schematically $[N+N^{\mathrm{fr}}]$), and not by $N(N+1)$.
  \end{itemize}
\item $\mathfrak{F}^{\mathrm{fi}}_J$ has degree $4r+0^{+}$.
\item The following constants have degree $0$: $C^{P,\star}_{\mathsf{a}}$, $\mathfrak{s}_J$,
$c_{\mathrm{loc}}$, $c_{\mathrm{sib}}$, $c_{\star}''$, $C_{\mathrm{pos}}$, $C_{\mathrm{fr}}$,
$C_{\mathrm{far}}$, $C_{\mathrm{old}}$, $C_{\mathrm{nr}}$, the constant $C^{\mathrm{inh}}$ of the
hypothesis on inherited data among the hypotheses of
Proposition~\ref{4.21:prop-group-totals}, $c_{\mathfrak{l}}$
and $\mathfrak{n}^{\mathrm{kept}}$.
\item $\vartheta$ has degree $p_1-q$, and $p_1-q\le-(7r+1)-(p_0-p_1)$.
\item $K_Q/r_A$ has degree at most $-(q_*-p_0)$, and $-(q_*-p_0)\le-(7r+1)$.
\item $\varepsilon_S$, $\tilde{\mathfrak{l}}$ and $\mathfrak{p}_j$ tend to zero faster than is
needed in the members in which they occur.
\end{itemize}
\item[(d)] Let $\mu^{\mathrm{kept}}$ and $\mathsf{M}^{\mathrm{kept}}$ be the single and the double
total of the kept members in \eqref{4.21:eq-group-totals-a}, and $\mathsf{a}_{\mathsf{M}}^{\mathrm{kept}}$
the bracket of the sixth member that corresponds to $\mathsf{M}^{\mathrm{kept}}$ at unit sizes.
The worst products, and their degrees, are
\begin{align*}
&c_{\mathrm{loc}}\mathfrak{F}^{\mathrm{fi}}\vartheta\ \text{(second member)}:
 && 6r+0^{+}-7r-1<0,\\
&c_{\star}''\bar V'\mathfrak{t}^{\mathrm{win}}\vartheta: && 5r+0^{+}-7r-1,\\
&C_{\mathrm{pos}}C_{\mathrm{nr}}\mathsf{l}^{2}\vartheta\ \text{and}\
 (100L\mathsf{R}^{\uparrow})^d\vartheta\ \text{(near groups)}: && 4r+0^{+}-7r-1,\\
&\mathsf{M}^{\mathrm{kept}}: && 4r+0^{+}-7r-1,\\
&\mathfrak{l}\,\mathfrak{t}^{\mathrm{win}}\vartheta\ \text{(first member)}: && 5r+0^{+}-7r-1,\\
&\mu^{\mathrm{kept}}: && -7r-1,\\
&\text{third member}: && r+4r-7r-1 .
\end{align*}
The sixth member is $N$ times the same brackets at unit sizes, marked $|_1$. In these brackets the
number $N$ of chains multiplying the chain part has become $1$. In the sixth member the following
brackets, each taken before the small factor it carries in that member ($\vartheta$ in the first
three, $K_Q/r_A$ in the fourth), have the following degrees:
\begin{align*}
&N c_{\mathrm{loc}}\mathfrak{F}^{\mathrm{fi}}|_1: && r+5r+0^{+},\\
&N c_{\star}''\bar V'(5L)^d\mathfrak{t}^{\mathrm{win}}|_1: && r+4r+r+0^{+},\\
&N c_{\mathrm{loc}}\mathfrak{l}\,\mathfrak{t}^{\mathrm{win}}|_1: && r+4r+r+0^{+},\\
&N C_{\mathrm{pos}}\bar V'\mathfrak{t}^{\mathrm{win}}|_1: && r+4r+r+0^{+},\\
&N\mathsf{a}_{\mathsf{M}}^{\mathrm{kept}}: && r+4r .
\end{align*}
The worst degree over all six members is $6r+0^{+}-7r-1<0$.
\item[(e)] No member contains a product of two volume-type functions. The products
$\bar V'\mathsf{l}^{2}$, $\bar V'\mathfrak{l}$, $\bar V'(L\mathsf{R}^{\uparrow})^d$ and
$\bar V'^{2}$ would have degree $8r+0^{+}-7r-1\ge0$; they are absent.
\end{enumerate}
\end{proposition}

\begin{proof}
\emph{The form of $\Phi$.} In each member of Definition~\ref{4.21:def-coupling-conditions}, as
displayed in \eqref{4.21:eq-coupling-conditions-a}, the function $\Phi$ is a finite sum of
products. Each product consists of a coefficient fixed before $\gamma$, some of the counting
functions of this chapter, and small factors. The counting functions are the factors to be
compensated, and they are read at their upper functions of $\mathsf{T}$. The small factors, which
compensate them, are also read at their upper functions of $\mathsf{T}$. After these replacements
each product is a number fixed before $\gamma$ times a function of $\mathsf{T}$ alone. This is the
first assertion of~(a).

\emph{The degrees.} The degrees in~(c) are read off the following displays:
\begin{itemize}
\item the upper functions in \eqref{4.21:eq-coupling-conditions-a};
\item the displays \eqref{4.21:eq-weights-incidences-a} of Notation~\ref{4.21:not-weights-incidences};
\item \eqref{4.21:eq-fine-count-a} of Lemma~\ref{4.21:lem-fine-count} and
\eqref{4.21:eq-group-totals-a} of Proposition~\ref{4.21:prop-group-totals}.
\end{itemize}
All volume-type functions are taken at $\check{R}^{\uparrow}\le\mathsf{R}^{\uparrow}(\mathsf{T})$.

Two of the entries are sums. In $\mathfrak{F}^{\mathrm{fi}}$ the chain part is
$N(N+1)(100L\mathsf{R}^{\uparrow})^d$. Its factors $N$ and $N+1$ have degree $r+0^{+}$ each, and
$(100L\mathsf{R}^{\uparrow})^d$ has degree $4r+0^{+}$. So the chain part has degree $6r+0^{+}$. In
the other part, the second summand of \eqref{4.21:eq-fine-count-a}, the factor multiplying
$(L\mathsf{R}^{\uparrow})^d$ is, apart from constants of degree $0$, the sum displayed in~(c); it
has degree $r+0^{+}$, since $N+1$ and $\mathfrak{c}$ have degree $r+0^{+}$ while
$\mathfrak{c}^{\circ}$, $(103L)^d$ and $C_{\mathrm{old}}$ have degree $0$. The volume factor has
degree $4r+0^{+}$. So that part has degree $5r+0^{+}$. The degree of
$\mathfrak{F}^{\mathrm{fi}}$ is therefore $6r+0^{+}$.

For the small factors the bounds come from \eqref{4.21:eq-small-coupling-2} and $q\ge q_*$:
\begin{itemize}
\item Since $q\ge q_*\ge p_0+7r+1$, the degree $p_1-q$ of $\vartheta$ satisfies
\begin{equation}\label{4.21:eq-small-coupling-3}
p_1-q\le p_1-p_0-7r-1=-(7r+1)-(p_0-p_1).
\end{equation}
\item The degree of $K_Q/r_A$ is at most $-(q_*-p_0)$. By \eqref{4.21:eq-small-coupling-2} at
$d=4$ this is at most $-(7r+1)$.
\item Since $p_1<p_0$, the right side of \eqref{4.21:eq-small-coupling-3} is at most $-7r-1$.
\end{itemize}

\emph{Degrees of products.} Let $f$ have degree $a+0^{+}$ and $h$ have degree $a'+0^{+}$. Then
their upper functions multiply, and the product of two powers of $\log\mathsf{T}$ is a power of
$\log\mathsf{T}$. So $fh$ has degree $a+a'+0^{+}$. A factor of degree $0$ does not change the
degree. The degrees in~(d) are the sums of the entries of~(c); we treat the products of largest
degree, the others being smaller.
\begin{itemize}
\item For $c_{\mathrm{loc}}\mathfrak{F}^{\mathrm{fi}}\vartheta$ the degree is at most
$0+(6r+0^{+})-7r-1$.
\item In the sixth member the chain part of $\mathfrak{F}^{\mathrm{fi}}|_1$ carries $1$ in place
of $N$. So $\mathfrak{F}^{\mathrm{fi}}|_1$ has degree $5r+0^{+}$, and
$Nc_{\mathrm{loc}}\mathfrak{F}^{\mathrm{fi}}|_1$ has degree $r+5r+0^{+}$. With the factor
$\vartheta$, of degree at most $-7r-1$, the product has degree at most $6r+0^{+}-7r-1$.
\item The next three products in the sixth member have degree $r+4r+r+0^{+}$ before their small
factor $\vartheta$ or $K_Q/r_A$. Each small factor has degree at most $-(7r+1)$, so each product
again has degree at most $6r+0^{+}-7r-1$.
\item The last bracket $N\mathsf{a}_{\mathsf{M}}^{\mathrm{kept}}$ of the sixth member, of degree
$r+4r$, is likewise followed in that member by a small factor of degree at most $-(7r+1)$, so
the product has degree at most $5r-7r-1$.
\end{itemize}
The remaining entries of~(d) have smaller degree.

\emph{Convergence to zero.} The worst degree is $6r+0^{+}-7r-1=-(r+1)+0^{+}<0$. Hence every term
of $\Phi$ not containing $\varepsilon_S$, $\tilde{\mathfrak{l}}$ or $\mathfrak{p}_j$ is bounded by
\begin{equation*}
C\,\mathsf{T}^{-(r+1)}(\log\mathsf{T})^{b},
\end{equation*}
with a number $C$ fixed before $\gamma$ and a fixed exponent $b$. This bound tends to zero as
$\mathsf{T}\to\infty$. The terms containing $\varepsilon_S$, $\tilde{\mathfrak{l}}$ or
$\mathfrak{p}_j$ tend to zero, and faster than any rate needed here, by their upper functions in
\eqref{4.21:eq-coupling-conditions-a}. A finite sum of terms tending to zero tends
to zero, so $\Phi(\mathsf{T})\to0$. This completes~(a).

\emph{Absence of double volume products.} The displays \eqref{4.21:eq-coupling-conditions-a},
\eqref{4.21:eq-weights-incidences-a},
\eqref{4.21:eq-fine-count-a} and \eqref{4.21:eq-group-totals-a} contain no product of two of the
volume-type functions $\bar V'$, $\mathsf{l}^{2}$, $\mathfrak{l}$, $(L\mathsf{R}^{\uparrow})^d$.
Such a product would have degree $8r+0^{+}$ before the small factor, and $8r+0^{+}-7r-1\ge0$
after it. The argument of the preceding paragraph would then fail. This is~(e).

\emph{The interval of admissible $\gamma$.} Fix one member and put
\begin{equation*}
\Sigma(\gamma):=\sup_{\mathsf{T}\ge\mathsf{T}_\gamma}\Phi(\mathsf{T}).
\end{equation*}
Since $\mathsf{T}_\gamma$ increases as $\gamma$ decreases, the set
$\{\mathsf{T}\ge\mathsf{T}_\gamma\}$ shrinks as $\gamma$ decreases. So $\Sigma(\gamma)$ does not
increase as $\gamma$ decreases. Consequently, if $\Sigma(\gamma)\le c$, then
$\Sigma(\gamma')\le c$ for every $\gamma'<\gamma$. Thus $\{\gamma:\Sigma(\gamma)\le c\}$ is a
down-set.

It is non-empty, for the following reason. Since $\Phi(\mathsf{T})\to0$ and $c>0$, there is
$\mathsf{T}_*$ with $\Phi\le c$ on $[\mathsf{T}_*,\infty)$. The supremum is then at most $c$ as
soon as $\mathsf{T}_\gamma\ge\mathsf{T}_*$, and this holds for all small $\gamma$, since
$\mathsf{T}_\gamma$ increases without bound as $\gamma$ decreases to $0$.

The set of $\gamma$ at which all six members hold is the intersection of six non-empty down-sets
of $\mathopen{]}0,\infty\mathclose{[}$, and is therefore again a non-empty down-set. Choosing
$\gamma_P>0$ in it, we obtain $\gamma_P>0$ with every $\gamma\in\mathopen{]}0,\gamma_P]$
admissible, which is~(b).

\emph{What the members bound.} Every member is an inequality of the form
\eqref{4.21:eq-small-coupling-1}. In it the quantities $N$, $N_0$, $\check{R}$ and the other
counting functions, and the ages, enter only through their upper functions of $\mathsf{T}$. The
coefficients are numbers fixed before $\gamma$. Each member is therefore a condition on $\gamma$
through $\mathsf{T}_\gamma$ alone, and, by the previous paragraph, it holds for all $\gamma$ small
enough. So no member bounds $g$, $g_k$ or $\gamma$ from below, and no member bounds $L$, $K$, $N$,
$N_0$, $\check{R}$ or an age from above. This completes~(b).
\end{proof}

\begin{remark}[The radius reading of the volume functions]\label{4.21:rem-radius-reading}
Let $\check{\mathsf{T}}_k$, $\mathsf{T}_\gamma$, $\mathsf{R}^{\uparrow}$, $\bar{V}'$ and the six supremum
conditions \eqref{4.21:eq-coupling-conditions-a} be as in
Definition~\ref{4.21:def-coupling-conditions}. Let $\check{R}^{\uparrow}=\check{R}_{h+1}$ and the level $h$
be as in \eqref{4.21:eq-strips-geometry-a}. Then
\begin{equation}\label{4.21:eq-radius-reading-1}
\check{R}^{\uparrow}\le\mathsf{R}^{\uparrow}(\check{\mathsf{T}}_k),
\end{equation}
where $\mathsf{R}^{\uparrow}$ has degree $r+0^{+}$ in $\mathsf{T}$, degrees in $\mathsf{T}$ and the symbol
$0^{+}$ being as in Proposition~\ref{4.21:prop-small-coupling}. The bound
\eqref{4.21:eq-radius-reading-1} follows from one-sided inequalities only and imposes no cap in the sense
of Definition~\ref{2.3:def-cap}. The second and the sixth of the conditions
\eqref{4.21:eq-coupling-conditions-a} are taken with the volume-type functions read at
\eqref{4.21:eq-radius-reading-1}. Each of them is a single condition
$\sup_{\mathsf{T}\ge\mathsf{T}_\gamma}\Phi(\mathsf{T})\le c$, with $\Phi$ and $c$ the function and the
constant of that condition as in Definition~\ref{4.21:def-coupling-conditions}. Under the floor
$q\ge p_0+7r+1$ of
Proposition~\ref{4.21:prop-small-coupling}, the $\mathsf{T}$-exponent of each is at most $-r-1+0^{+}$.
Hence no new lower bound on $q$ arises.
\end{remark}

\begin{proof}
Write $x_i=g_i^{-2}$. In this proof $N$ is the stage count; $N$, $C_N$ and
$r_{\mathrm{pr}}$ are as in Definition~\ref{4.21:def-coupling-conditions}, where the bound
$N\le C_N(\check{\mathsf{T}}_k+1)^{r_{\mathrm{pr}}}$ is stated. Let $\lambda$ and $c_2$ be the constants of
the comparison of inverse couplings in the correlation theorem.

\emph{The bound on the radius.} By the definition of the separation radii
in \eqref{4.21:eq-strips-geometry-a}, we have
$\check{R}^{\uparrow}=\check{R}_{h+1}=R(\min_{i\le h+1}x_i)$, where $R(\cdot)$ is the radius function of
that display, evaluated at an inverse coupling. By the growth bound for the radius function,
$R(\min_{i\le h+1}x_i)<L\max\{1,(\log x_{h+1})^{r_{\mathrm{pr}}}\}$. By the comparison of inverse couplings
across the stages of one step in the correlation theorem, $x_{h+1}\le x_k(1+3\lambda N+2c_2)$. By
Definition~\ref{4.21:def-coupling-conditions}, the stage count satisfies
$N\le C_N(\check{\mathsf{T}}_k+1)^{r_{\mathrm{pr}}}$. Since $\log$ and $t\mapsto t^{r_{\mathrm{pr}}}$ are
increasing, these three inequalities give
\begin{equation}\label{4.21:eq-radius-reading-2}
\begin{aligned}
\check{R}^{\uparrow}
&=R\bigl(\min\nolimits_{i\le h+1}x_i\bigr)
<L\max\bigl\{1,(\log x_{h+1})^{r_{\mathrm{pr}}}\bigr\}\\
&\le L\max\Bigl\{1,\Bigl(\log\bigl(x_k\bigl(1+3\lambda C_N(\check{\mathsf{T}}_k+1)^{r_{\mathrm{pr}}}
+2c_2\bigr)\bigr)\Bigr)^{r_{\mathrm{pr}}}\Bigr\}.
\end{aligned}
\end{equation}
By the definitions of the logarithmic scale variable $\check{\mathsf{T}}_k$ and of the upper function
$\mathsf{R}^{\uparrow}$ in Definition~\ref{4.21:def-coupling-conditions}, the right side of
\eqref{4.21:eq-radius-reading-2} is at most $\mathsf{R}^{\uparrow}(\check{\mathsf{T}}_k)$; this is
\eqref{4.21:eq-radius-reading-1}. The function
$\mathsf{R}^{\uparrow}$ has degree $r+0^{+}$. Each inequality used above is one-sided, and none bounds
$L$, $N$ or a radius from above by a cap.

\emph{Comparison with the cruder reading.} One may instead read the radius through the cruder bound
$\check{R}\le(L+1)N^{\beta_0}R_k$, which rests on \cite[(2.9)]{B-CMP119}, with $\beta_0$ and $R_k$ as in
\cite[(2.9)]{B-CMP119}. The floor $q\ge p_0+7r+1$ of Proposition~\ref{4.21:prop-small-coupling} is
computed on this reading. Under it, $\bar{V}'$ has degree $6r$. The chain part $N^2\bar{V}'$ of the total
$\mathfrak{F}^{\mathrm{fi}}$ would then have $\mathsf{T}$-exponent $8r-(q-p_1)$, where $p_1<p_0$ is the
exponent so denoted in Proposition~\ref{4.21:prop-small-coupling}; this number is not
negative by its form alone. For this reason the second and sixth conditions of
\eqref{4.21:eq-coupling-conditions-a} are read at \eqref{4.21:eq-radius-reading-1}. Each is its own
supremum condition, of the kind used in \cite[p.~186]{B-CMP122a}: there quantities of this form are made
arbitrarily small by the choice of $\gamma$, with constants that do not depend on any scale.

\emph{Fitting the envelope.} The correlation theorem carries the envelope
\begin{equation}\label{4.21:eq-radius-reading-3}
24K'\,\mathsf{T}^{(1+d+d\beta_0^{\mathrm{pr}})r_{\mathrm{pr}}}
\bigl[\mathsf{T}^{-(q-p_1)}+\cdots\bigr],
\end{equation}
in which the dots stand for its further terms, and where $K'$ and $\beta_0^{\mathrm{pr}}$ are the
constants of \eqref{4.21:eq-radius-reading-3} as stated in the correlation theorem.
By Proposition~\ref{4.21:prop-small-coupling},
$(1+d+d\beta_0^{\mathrm{pr}})r_{\mathrm{pr}}=7r$ at $d=4$. At \eqref{4.21:eq-radius-reading-1},
$\bar{V}'$ has degree
$4r+0^{+}$. Hence $N^2\bar{V}'$ has degree $6r+0^{+}$, which is at most $7r$. Both conditions therefore fit
\eqref{4.21:eq-radius-reading-3} by degree, with $K'$ replaced by a larger constant, although not term by
term. The resulting exponent satisfies
\begin{equation}\label{4.21:eq-radius-reading-4}
6r+0^{+}-(q-p_1)\le 6r+0^{+}-(p_0+7r+1)+p_1=-r-1+0^{+}-(p_0-p_1)
\end{equation}
under $q\ge p_0+7r+1$. This is at most $-r-1+0^{+}$ because $p_1<p_0$. The floor on $q$ is therefore the
one already in force, and no new floor arises.

\emph{What is not assumed.}
The condition $e^{-(\kappa_1-1)}(10\check{R}_{k_0+1}/L)^{d-1}\le\mathrm{const}$, displayed beside
\eqref{4.21:eq-coupling-conditions-a}, is not assumed and is used nowhere. The other condition that
Definition~\ref{4.21:def-coupling-conditions} sets aside is likewise not assumed.
\end{proof}

\begin{theorem}[Invariance of the relative small-field format]\label{4.21:thm-format-invariance}
In this statement $(T)$, $(P)$ and $(L)$ denote the three sites of a step of the induction:
the $\mathbf{T}$-step from step $k-1$ to step $k$; the stages $n=0,\dots,N-1$, where
$N=N_k:=\check{R}_k$, at which the chain terms $\mathbb{C}^{(n+1)}_k$ are prepared; and the
localisation of the operation $\mathbf{R}_k$.

Let the parameters of the correlation theorem satisfy, besides the floors of that theorem, the
six inequalities
\begin{equation*}
\begin{gathered}
\tfrac14\,\beta_s\kappa\ \ge\ \lambda_\theta,\qquad
\sigma_{*}\ \ge\ c_d+4,\qquad
\kappa_1\ \ge\ 1+c_l(2\kappa+4)+\lambda_{*}(d,L),\\
2e^{2}\hat{b}_t^{2}N^{\flat}\mathsf{x}\ \le\ 1,\qquad
\tfrac14\,\delta_P w\ \ge\ 110\,d\,(a_k+1),\qquad
\tfrac1{16}\,\delta_P w\ \ge\ c_{\boxplus}\log\hat{b}_t+1 .
\end{gathered}
\end{equation*}
The first two are implied by the correlation theorem's floor $\kappa_d$. The fourth is imposed
after $L$ and $\kappa$ have been fixed. Let $\gamma$ satisfy the following ceilings:
\begin{itemize}
\item the ceiling on the region size parameter and the tail ceiling, both contained in
\eqref{4.21:eq-strips-geometry-a};
\item the two relative ceilings at the sites $(T)$ and $(L)$. Wherever a kept member occurs in
the sum, these are taken in their re-maximised forms, in which the number
$\mathfrak{n}^{\mathrm{kept}}$ of kept members enters, and the ceiling of the lemma on deep units
is taken in the same re-maximised form;
\item the level ceiling of Notation~\ref{4.21:not-weights-incidences};
\item the six conditions on the coupling of Definition~\ref{4.21:def-coupling-conditions},
\eqref{4.21:eq-coupling-conditions-a}, the second and the sixth of them taken at the upper
function $\check{R}^{\uparrow}\le\mathsf{R}^{\uparrow}(\check{\mathsf{T}}_k)$ displayed
in \eqref{4.21:eq-strips-geometry-a};
\item the two sharp ceilings, taken at the value $2+\mathsf{K}^{0}$, where $\mathsf{K}^{0}$ is
the constant so named in the statement of those two ceilings;
\item the ceilings on the coupling under which the third-factor lemma holds.
\end{itemize}
The relative small-field machine is run with clause (S) of its definition. Assume the following,
each at the place in the order of the induction, lexicographic in step and site, at which it is
used:
\begin{enumerate}
\item[(i)] the third-factor lemma;
\item[(ii)] the input bound on plain labels at the site $(L)$, with its two companion statements on
the decoupling family, under the pointwise form of its base condition;
\item[(iii)] the statements of the correlation theorem taken by statement on labels. No bound on a
sum over terms at a window anchor is taken from that theorem by statement. Among the statements
taken are, in particular, the following.
\begin{itemize}
\item Its bounds on the small-field terms and on the near groups, in their second to fifth
clauses, the fifth in its first alternative.
\item In the third of these clauses, which bounds the near-group total, the data factor is taken
in the local form $L_{\circ}^{6s(\square)}\mathfrak{d}(\square)$. Here $s(\square)$ is the raise
of the cells met by $\tilde{\square}^{5}$ in the current operation $\mathbf{R}_k$, and
$\mathfrak{d}(\square)$ is the inherited-datum factor.
\item The inherited-datum hypothesis, Definition~\ref{4.21:def-datum-hypothesis}, whose proof is
outstanding, with its bound \eqref{4.21:eq-datum-hypothesis-a}. It is assumed both for the
near-group total of one machine and for the difference of the near-group totals of two machines
$\mathsf{A}$, $\mathsf{B}$ on one scaffold.
\item Parts (a) and (c) of the summation lemma for the small-field terms of the correlation
theorem;
\end{itemize}
\item[(iv)] the containment statements at the three sites; the base lemmas of the machine on
localisation and on the data constants; the first part of the scaffold comparison; the
admissible geometry \eqref{4.21:eq-strips-geometry-a}, together with the declaration
$N=N_k:=\check{R}_k$ (one class and one pair $(h,k)$ of levels per step), the nesting statement,
the definitions of arrested attempts and of the stages at which their terms are read, and the
clock lemma for arrested attempts;
\item[(v)] the orbit-data statement for stored terms, with its statements for kept members. Write
$\mu^{\mathrm{kept}}$ and $\mathsf{M}^{\mathrm{kept}}$ for the totals of the latter at the site
$(P)$. These totals enter the first, the second and the sixth condition of
Definition~\ref{4.21:def-coupling-conditions}, with the constants of
Proposition~\ref{4.21:prop-group-totals}. The difference statement for kept members,
Definition~\ref{4.21:def-kept-difference}, whose proof is outstanding, with its bound
\eqref{4.21:eq-kept-difference-a}, is used in the derivation of the sixth condition
(Lemma~\ref{4.21:lem-machine-differences});
\item[(vi)] hence, for the class $\mathfrak{S}$ of wrapped descendants and with the constants of
Proposition~\ref{4.21:prop-group-totals} entering hypothesis~(v), the
source-tube statement. It is assumed in its bound on the source functional over the source tube
and in its two-point oscillation bound, no convergence of the latter to zero being assumed. It
comes with its two size bounds for the probe family $\mathsf{Q}$, in the form derived from the
list of statements displayed with Proposition~\ref{4.21:prop-group-totals};
\item[(vii)] the theorem on the terms created by live arrested attempts, in its clause on the format
and orbit datum of these terms and in its clause bounding their summed weights, with the
definitions it uses, the datum being $\mathfrak{N}^{G}$. No statement on the levels of these terms
is assumed;
\item[(viii)] the reading of the base: $b(x)$ is read off a completed operation $\mathbf{R}$ of
either class of components, the flat label being carried on a thrown-back second-class component,
in the sense of \cite[(1.33)]{B-CMP122b};
\item[(ix)] the volume counts of the correlation theorem on which part (a) of the volume
proposition rests: the count of \cite[(1.26)]{B-CMP116} for the terms $\mathbf{B}$, $\mathbf{B}'$,
and the corresponding counts for the chain terms, which are assumed here and whose proof is
outstanding. Part (a) of the volume proposition discharges the correlation theorem's volume
displays from these counts;
\item[(x)] the statements named in the proof below, as they are stated and proved in this book:
the propositions at the sites $(T)$ and $(L)$, the boundary and volume propositions and the
corollary of the boundary proposition, the lemmas and totalling statements listed in the proof,
and the list of the eighteen displays of the correlation theorem in their relative form, with $X$
replaced by $A_F$; together with the proof, given in this book under the anchored bound on the set
$\mathfrak{O}$ of old blocks and part (e) of the group-total proposition, of the passage of the
relative small-field format from step $k$ to step $k+1$.
\end{enumerate}
Then the following hold. All of them are taken inside the single induction on (step, site) that
carries, jointly, the correlation theorem, the classification of stored terms into the classes
$\mathfrak{A}$ and $\mathfrak{S}$, and the source-tube statement.
\begin{enumerate}
\item[(a)] For every $k<K$, the inductive assertion of the correlation theorem at step $k$,
together with the relative small-field format at step $k$, implies the relative small-field format
at step $k+1$.
\item[(b)] The relative small-field format at a step is the following assertion. Every display of
the correlation theorem, in its parts treating the sites $(T)$, $(P)$, $(L)$ and the boundary and
volume terms, that carries one of $d_j(X)$, $N(X)$, $\bar{X}$, $\bar{S}_v$, $K_v$,
$\mathsf{d}(e)$ of a boundary-type term holds, in the norm $\mathfrak{N}^{A,\mathbb{C}}$, with $X$
replaced by the label $A_F$ of the stored term $F$; here $N(X)$, $\bar{X}$, $\bar{S}_v$ and $K_v$
are the quantities so denoted in the correlation theorem for a boundary-type term. These are the
assertions of the propositions at the sites $(T)$, $(P)$ (with
Proposition~\ref{4.21:prop-group-totals}) and $(L)$, of the boundary proposition, of the volume
proposition and of the corollary of the boundary proposition. The labels, index sets, floors and
ceilings are the same in every norm of the list of norms of the relative format. The radii are
$\varepsilon-\eta^{T}$ at
$(T)$, the radius of the class-chain corollary of the group-total proposition at $(P)$, and
$\varepsilon'=\varepsilon-\eta^{L}-\ell 7$ at $(L)$. Deep units are included, and the orbit data
$\mathsf{e}$ are carried on the frames born at the step. The displays concerned are the eighteen
displays of the correlation theorem in their relative form, with $X$ replaced by $A_F$,
enumerated in the list of these displays.
\item[(c)] The following are not altered: the second group of bounds of
\cite[(1.99)]{B-CMP122b}; the torus bound and the running-shift statement of the correlation
theorem; and that theorem's statements on inputs, on reference pairs, on memory and on real values.
\item[(d)] The third-factor lemma is used twice: once in the lemma on deep units, and a second time
through hypothesis~(ii).
\end{enumerate}
\end{theorem}

\begin{proof}
The proof, given in this book, of the passage of the relative small-field format from step $k$ to
step $k+1$, inside the induction and site by site in the order of the induction, under the
anchored bound on the set $\mathfrak{O}$ of old blocks and part (e) of the group-total
proposition, is referred to below as that proof. It uses the statements listed below, each with
its statement and its proof:
\begin{itemize}
\item the proposition at $(T)$;
\item the proposition at $(L)$, with its lemma on the starred norm and the lemma on deep units;
\item the boundary proposition, the volume proposition and the corollary of the boundary
proposition;
\item the grain lemma;
\item the five parts of the resummation statement at $(P)$;
\item the list of members of the stored sums at $(P)$;
\item the cell and block sums of Lemma~\ref{4.21:lem-cell-block-sums};
\item the summation, anchoring and boundary lemmas and the summation device;
\item parts (a)--(c) of the group-total proposition, with the steps of its proof that establish
them;
\item the totalling statements for majorants, for constants and for families, and the corollary
that totals the constants;
\item the transfer to second-class components;
\item the linearity remark, that these estimates are linear in the terms and anchor by the scaffold
only;
\item every statement named in the order of the induction.
\end{itemize}
These statements, and that proof, are hypothesis~(x); we take them as they stand, statements
together with proofs, and do not re-derive them. The proof of the theorem is that proof with the
two substitutions described next, followed by the check that nothing else in it is affected.

\emph{First substitution.} Wherever that proof uses part (e) of the group-total proposition, we use
Proposition~\ref{4.21:prop-group-totals} instead. Its hypotheses are all among those of the present
theorem:
\begin{itemize}
\item the fourth, second, fifth and sixth of the displayed floors;
\item the ceiling on the region size parameter, the tail ceiling and the level ceiling;
\item the six conditions of Definition~\ref{4.21:def-coupling-conditions};
\item the inherited-datum hypothesis, Definition~\ref{4.21:def-datum-hypothesis}, whose proof is
outstanding, for one machine and for the difference of two machines on one scaffold, assumed in
hypothesis~(iii).
\end{itemize}
Inside the first, second and sixth of those conditions, Proposition~\ref{4.21:prop-group-totals}
uses the totals $\mu^{\mathrm{kept}}$ and $\mathsf{M}^{\mathrm{kept}}$ of hypothesis~(v) with its
own constants.

\emph{Second substitution.} Wherever that proof bounds the total weight at an old block by the
anchored bound on the set $\mathfrak{O}$, we use instead Lemma~\ref{4.21:lem-old-block-total}.
That lemma bounds the total weight $\mathfrak{S}^{\mathrm{all}}(\square_0)$ of all terms of the
format class incident to $\square_0$, at every block $\square_0$ of the enlarged set
$\mathfrak{O}^{+}$. Its hypotheses are the floors under which that proof is given, the floor
$\sigma_{*}\ge c_d+4$ and the ceiling on the region size parameter, all of which are assumed here.

\emph{Nothing else is affected.} No part of that proof, and no statement it uses, other than the two
places just substituted, uses either the steps of the proof of the group-total proposition that
establish its part (e) or the set $\mathfrak{O}$. We verify this item by item.
\begin{enumerate}
\item The proposition at $(L)$ uses the group-total proposition only through the radius of its
class-chain corollary and through the outputs of the site $(P)$ in format.
\item The boundary proposition, the volume proposition and the concluding part of that proof use
the group-total proposition only through the bound
$\|\mathbb{C}^{(n+1)}_k\|^{A,\mathbb{C}}\le C_0^{\sharp}$ and through the format.
\item The orbit-data statements at the site $(P)$ use the groups, which the substitutions leave
unchanged.
\item The set $\mathfrak{O}$ occurs in the statements used only in two places: the part of the
group-total proposition that treats old blocks, together with its proofs, and the first of the
eighteen displays of the correlation theorem in their relative form named in conclusion (b). At
each of these
occurrences $\mathfrak{O}$ is replaced by $\mathfrak{O}^{+}$. The bound used there is then the bound
of Lemma~\ref{4.21:lem-old-block-total}, which is stated for every $\square_0\in\mathfrak{O}^{+}$.
\item Where that proof speaks of the current chain, it uses the norm of the chain of the class. A
sibling $C'$ of the current component enters that chain only through its stages
$i\le n(C')\le n$, $n(C')$ its stage count, which are members of the class chain
(Lemma~\ref{4.21:lem-base-below-level}). Hence the presence of siblings does not alter these
dependences.
\end{enumerate}

With the two substitutions and these checks, that proof gives (a) and (b) for every $k<K$. Neither
substitution touches the statements listed in (c), and the third-factor lemma is used exactly where
that proof uses it: in the lemma on deep units and through hypothesis~(ii). This proves (d).
\end{proof}

\begin{remark}[Gauge covariance at complex sources]\label{4.21:rem-complex-covariance}
In the setting of Theorem~\ref{4.21:thm-format-invariance}, let the source $\zeta$ be complex, not necessarily real. Consider the block-local gauge map carried by a link,
\begin{equation}\label{4.21:eq-complex-covariance-1}
\mathsf{u}(\mathsf{s})=\mathfrak{u}\bigl[e^{i\eta\mathcal{Y}(\mathsf{s})}\mathbf{K};\mathbf{K}\bigr],
\end{equation}
built from the link's layer $\mathcal{Y}(\mathsf{s})$ and the configuration $\mathbf{K}$ at the lower end of the link, where $\eta$ is the parameter, fixed in the setting of Theorem~\ref{4.21:thm-format-invariance}, by which the layer is multiplied in the exponent. Consider also the gauge maps of the frames of that setting. At such $\zeta$ these maps take values in the complexification $G^{\mathbb{C}}$ of the gauge group $G$.

Accordingly, every object obtained by conjugation with these maps is defined by evaluation on the regular representative of the configuration. Among these objects are:
\begin{itemize}
\item the orbit datum $F^{\sharp}(\dots;h\,\mathsf{u}\,g')$ of a stored term, evaluated at the product of $\mathsf{u}$ with gauge transformations $h$ and $g'$;
\item the rotated current rider $\operatorname{Ad}_h\mathcal{J}$;
\item the slots of a stored term rotated by these maps.
\end{itemize}
On the real slice each such object is identified with the object of the same name defined there for $G$-valued maps.

The two identities of the setting of Theorem~\ref{4.21:thm-format-invariance} in which these conjugated objects enter, namely the identity stated for the frames and the invariance identity for stored terms, are read in this sense at complex $\zeta$.
\end{remark}

\begin{proof}
The first assertion is that the map $\mathsf{u}(\mathsf{s})$ of \eqref{4.21:eq-complex-covariance-1} and the gauge maps of the frames are $G^{\mathbb{C}}$-valued at complex $\zeta$. This is the reason why the conjugated objects listed in the remark must be given a definition at complex $\zeta$ at all. The definitions of these objects in force at real $\zeta$ presuppose that the conjugating maps take values in $G$; at complex $\zeta$ they take values in $G^{\mathbb{C}}$, so a definition has to be supplied.

The definition adopted is evaluation on the regular representative. At real $\zeta$ the map $\mathsf{u}(\mathsf{s})$ of \eqref{4.21:eq-complex-covariance-1} and the gauge maps of the frames are $G$-valued, and evaluating the conjugated objects on the regular representative returns the objects as defined for $G$-valued maps; this is the identification asserted in the remark. Hence at real $\zeta$ every conjugated object is the same object as before. Every statement about these objects at real $\zeta$ therefore keeps both its content and its truth.

At complex $\zeta$ these two identities are statements about the objects as just defined. Reading them so changes no statement on the real slice.
\end{proof}

We use the cover of the lattice by blocks $\square$ and its partition of unity $\{\zeta_\square\}$, which are introduced in \cite[p.~25]{B-CMP109} as follows:
\begin{quote}
We construct a cover of the space $T$ by cubes $\square$, which are unions of $2^d$ neighbouring
cubes from $\pi_k$. For this cover we take a partition of unity
$1=\sum_{\square}\zeta_{\square}$ \dots\ if $y$ is a center of the cube $\square$, then
$\zeta_{\square}(x)=\prod_{\mu=1}^{d}\zeta(M^{-1}(x_\mu-y_\mu))$, where
$\zeta\in C_0^{\infty}(\mathbb{R}^1)$, $\zeta(t)=1$ for $|t|\le 1/3$, $\zeta(t)=0$ for
$|t|\ge 2/3$, $\zeta$ has derivatives up to the second order bounded by~5\,[\dots].
\end{quote}

A block $\square$ with centre $y$ has side $2M$, the cubes of $\pi_k$ having side $M$; that is,
$\square=\{x:\ |x_\mu-y_\mu|\le M,\ \mu=1,\dots,d\}$. We put
\begin{equation}\label{4.21:eq-cover-overlap-1}
\square^{\zeta}:=\Bigl\{x:\ |x_\mu-y_\mu|\le \tfrac23 M,\ \mu=1,\dots,d\Bigr\}.
\end{equation}

For the overlap count we use the blocks themselves and not their enlargements
$\tilde{\square}$. The enlargements do not have an overlap number bounded independently of
the number of zones. On the old blocks $\mathfrak{O}$, consider a point $y$ of zone $k_0+1$
lying just below a window. It lies in $\tilde{\square}$ for a block $\square$ of every window
zone $k_0+1,\dots,j$, because the thin strata between these zones may have total thickness
less than two level-$i$ units, $i$ being the current level of the cells of those strata. This
is the geometry of
Lemma~\ref{4.21:lem-under-block}(b). The overlap number of $\{\tilde{\square}\}$ is therefore
controlled only by the number $j-k_0$ of window zones, which is bounded only by one less than
the stage count of the window (of the arrested attempt, at an arrested window). Constants
built from it, such as $K_A$, which enters $K_2^{\flat}$, would then grow with
$\log_L\check{R}_k$, a number of levels, instead of depending only on $(d,L,\kappa,\kappa_1)$.

\begin{lemma}[Bounded overlap of the block cover]\label{4.21:lem-cover-overlap}
Let $\{\square\}$ be the cover by blocks of all zones, each zone's blocks $\square$ and
functions $\zeta_\square$ being those
of \cite[p.~25]{B-CMP109} described above, with the partition of unity $\{\zeta_\square\}$
over all zones formed as in \cite{B-CMP119}, and let $\square^{\zeta}$ be the boxes of
\eqref{4.21:eq-cover-overlap-1}.
\begin{enumerate}
\item[(a)] For every block, $\operatorname{supp}\zeta_\square\subset\square^{\zeta}\subset\square$.
\item[(b)] For every point $x$,
\begin{equation}\label{4.21:eq-cover-overlap-2}
\#\{\square:\ x\in\square\}\le 3\cdot 2^d\le 4^d .
\end{equation}
Consequently the overlap number $\nu_d$ of $\{\square^{\zeta}\}$ satisfies
\begin{equation*}
\nu_d:=\sup_x\#\{\square:\ x\in\square^{\zeta}\}
\le\sup_x\#\{\square:\ x\in\square\}\le N_{\mathrm{cov}}:=4^d .
\end{equation*}
\item[(c)] Let $\delta>0$ satisfy $\delta w\ge c_{\boxplus}\log\hat{b}_t+1$, the hypothesis
of Lemma~\ref{4.21:lem-cell-block-sums}(b) with $\beta=\delta$, and define the block sum
$\mathsf{c}_{\square}$ by the left-hand side of \eqref{4.21:eq-cover-overlap-3}, where
$\operatorname{dist}(y,A)$ denotes the distance from the point $y$ to the set $A$, measured as
in Notation~\ref{4.4:not-units-distances}. Then
\begin{equation}\label{4.21:eq-cover-overlap-3}
\mathsf{c}_{\square}:=\sup_y\sum_{\square}e^{-\delta\operatorname{dist}(y,\square^{\zeta})}
\le C_{\mathrm{blk}}\bigl(1+\hat{b}_t C_{\mathrm{cell}}\bigr),
\end{equation}
and the right-hand side is a number depending only on $d$ and $L$.
\end{enumerate}
\end{lemma}

\begin{proof}
(a) Let $\square$ have centre $y$.

\emph{Support of $\zeta_\square$.} If $x\notin\square^{\zeta}$, then $|x_\mu-y_\mu|>\frac23 M$ for
some $\mu$. Hence $|M^{-1}(x_\mu-y_\mu)|>2/3$, so $\zeta(M^{-1}(x_\mu-y_\mu))=0$, and the
product defining $\zeta_\square(x)$ vanishes. The set $\square^{\zeta}$ is closed, so
$\operatorname{supp}\zeta_\square\subset\square^{\zeta}$.

\emph{Change of zone.} Where the cover passes from one zone to the next, the functions
$\zeta_\square$ of the blocks of each zone already sum to one on that zone. The normalisation
producing the
partition of unity over all zones in \cite{B-CMP119} therefore moves no support, and the
inclusion persists.

\emph{The box lies in the block.} Finally $\frac23M\le M$ gives $\square^{\zeta}\subset\square$.
The support lies strictly inside the block, and in particular not merely in $\tilde{\square}$.

(b) Let $x$ be a point and $\Sigma$ its cell. By Lemma~\ref{4.21:lem-live-region-geometry},
item (i3), the cell $\Sigma$ meets only blocks of three kinds:
\begin{itemize}
\item blocks of its own zone, of which at most $2^d$ contain $x$;
\item blocks of the zone above, of which at most $2^d$ contain $x$;
\item blocks of the zone below, and these only within two units of that zone, of which at most
$2^d$ contain $x$.
\end{itemize}
Summing the three bounds gives $\#\{\square:\ x\in\square\}\le 3\cdot2^d$. Since $d=4$ we have
$3\le 2^d$, so $3\cdot 2^d\le 4^d$, which is \eqref{4.21:eq-cover-overlap-2}.

Since $\square^{\zeta}\subset\square$ by (a), we have
$\{\square:\ x\in\square^{\zeta}\}\subset\{\square:\ x\in\square\}$ for every $x$. Taking the
supremum over $x$ gives the bound on $\nu_d$. The statement on $N_{\mathrm{cov}}$ is
\eqref{4.21:eq-cover-overlap-2}.

(c) Fix a point $y$ and let $\Sigma$ be its cell. We establish the chain
\begin{align}
\sum_{\square}e^{-\delta\operatorname{dist}(y,\square^{\zeta})}
&\le\sum_{\square}e^{-\delta d^{\wedge}(y,\square)}\label{4.21:eq-cover-overlap-4}\\
&\le\sum_{\square}e^{-\delta d^{\wedge}(\square,\Sigma)}\label{4.21:eq-cover-overlap-5}\\
&\le C_{\mathrm{blk}}\bigl(1+\hat{b}_tC_{\mathrm{cell}}\bigr),\label{4.21:eq-cover-overlap-6}
\end{align}
where $\square$ runs over all blocks of the cover.

\emph{First inequality.} By (a), $\square^{\zeta}\subset\square$, hence
$\operatorname{dist}(y,\square)\le \operatorname{dist}(y,\square^{\zeta})$. The scaled distance
does not exceed the distance, by \eqref{4.21:eq-cell-block-sums-a}. Thus
$d^{\wedge}(y,\square)\le \operatorname{dist}(y,\square^{\zeta})$ for every block, and
\eqref{4.21:eq-cover-overlap-4} follows term by term.

\emph{Second inequality.} The point $y$ lies in $\Sigma$, and by the definition of the scaled
distance recalled in \eqref{4.21:eq-cell-block-sums-a} the scaled distance between two
sets does not exceed the scaled distance from any point of the one to the other; hence
$d^{\wedge}(\square,\Sigma)\le d^{\wedge}(y,\square)$ for every block, and
\eqref{4.21:eq-cover-overlap-5} follows term by term.

\emph{Third inequality.} Apply Lemma~\ref{4.21:lem-cell-block-sums}(c) with $\beta=\delta$ and
with $\Sigma$ the cell of $y$; its hypothesis is that of item (b) of the same lemma, which is
the hypothesis of part (c) of the present lemma. It bounds the sum over all blocks in
\eqref{4.21:eq-cover-overlap-5}
by $C_{\mathrm{blk}}(1+\hat{b}_tC_{\mathrm{cell}})\cdot\#\Sigma$, where $\#\Sigma$ is the
number of cells of $\Sigma$, and $\#\Sigma=1$ since $\Sigma$ is
a single cell. This is \eqref{4.21:eq-cover-overlap-6}.

Taking the supremum over $y$ gives \eqref{4.21:eq-cover-overlap-3}. The constants
$C_{\mathrm{blk}}$, $\hat{b}_t$, $C_{\mathrm{cell}}$ of Lemma~\ref{4.21:lem-cell-block-sums}
depend only on $d$ and $L$; no level, age, number or thickness of zones enters them. So does
the bound.
\end{proof}

\begin{remark}[The cover constants are computed from the blocks]
The following quantities are defined in terms of the boxes $\square^{\zeta}$ and the blocks
$\square$, with $\operatorname{supp}\zeta_\square\subset\square^{\zeta}\subset\square$ by
part (a) of Lemma~\ref{4.21:lem-cover-overlap}, and not of the enlargements $\tilde{\square}$:
the overlap number $\nu_d$ of the cover; the cover constant $N_{\mathrm{cov}}$ and the constant
$c_R^0$ built from it; the block sum $\mathsf{c}_{\square}$, the two-block sum
$\mathsf{c}_{\square\square}$, and the block quantity $q_\square$ of the bounds in which these
sums occur. In those bounds the support of $\zeta_\square$ enters through the containment
$\operatorname{supp}\zeta_\square\subset\square^{\zeta}\subset\square$ of part (a) of
Lemma~\ref{4.21:lem-cover-overlap}, which is sharper than
$\operatorname{supp}\zeta_\square\subset\tilde{\square}$. The quantities $\nu_d$,
$N_{\mathrm{cov}}$ and
$\mathsf{c}_{\square}$ are bounded by that lemma, and the two-block sum
$\mathsf{c}_{\square\square}$ is estimated in the same way as $\mathsf{c}_{\square}$ in the
proof of part (c) of that lemma, by a number depending only on $d$ and $L$.

\emph{Triangle inequality for block quantities.} Let $q_\square$ be a quantity attached to a
block that is a supremum over points $y'\in\square^{\zeta}$. For every point $y$ and every
$y'\in\square^{\zeta}$ one has $y'\in\square$ by part (a), hence
$\operatorname{dist}(y,\square)\le \operatorname{dist}(y,y')$. Therefore the triangle inequality
used in the bound on $q_\square$, and the estimate obtained from it, hold with
$\operatorname{dist}(y,\square)$ in place of $\operatorname{dist}(y,\tilde{\square})$.

\emph{Dependence of the constants.} The constant $K_A$, which carries a factor exponential in
$\kappa_1$ and enters $K_2^{\flat}$, the constant $K_2^{\flat}$ itself, and the constants formed
from
them in the same way, such as $K_{\flat}''$, are built, by their definitions, from $\kappa$,
$\kappa_1$, from numbers depending only on $d$ and $L$, and from the geometry of the cover only
through $\nu_d$, $N_{\mathrm{cov}}$, $\mathsf{c}_{\square}$ and $\mathsf{c}_{\square\square}$. By
Lemma~\ref{4.21:lem-cover-overlap} and the estimate of $\mathsf{c}_{\square\square}$ just
described,
these are numbers depending only on $d$ and $L$. Hence $K_A$, $K_2^{\flat}$ and $K_{\flat}''$
depend only on $(d,L,\kappa,\kappa_1)$.
\end{remark}

\begin{definition}[Site geometries, readers and labels]\label{4.21:def-site-geometry}
Fix a level $i$. Cells, their walls and their levels $\operatorname{lev}(\cdot)$, the partitions
$\pi_\ell$ into cubes of level $\ell$, the never-cut blocks, and the distance $\operatorname{dist}_i$ at
level $i$ are those of the geometry of cells and blocks set up earlier in this chapter.

\emph{Site geometry.} A \emph{site geometry} at level $i$ is a triple
$(\pi^{\mathfrak{s}},\mathfrak{H},\mathfrak{F})$ with the following properties.
\begin{enumerate}
\item $\pi^{\mathfrak{s}}$ is a finite family of cells, all of level at most $i$; its elements are called
\emph{members}. Two members that share a wall have levels differing by at most one; this is the level
property of the cells used throughout this chapter.
\item $\mathfrak{H}$ is a finite, possibly empty, family of \emph{hubs} (the never-cut blocks). Each hub is
a box made of $\pi_i$-cubes whose sides are at most
\begin{equation*}
h_{\sharp}=112,
\end{equation*}
the bound on the sides of never-cut blocks supplied by the geometry lemma of this chapter. The interior
of each hub is disjoint from the members of $\pi^{\mathfrak{s}}$, and distinct hubs are at distance at
least $3$ from each other.
\item $\mathfrak{F}\subset\pi^{\mathfrak{s}}$ is the family of \emph{source cells}, all of level $i$. If
$\mathfrak{H}\neq\emptyset$, then $\mathfrak{F}$ is the family of the members sharing a wall with a hub.
\end{enumerate}
Two members are \emph{adjacent} if they share a wall, or if both share a wall with one and the same hub.
The words \emph{connected} and \emph{component}, applied to subfamilies of $\pi^{\mathfrak{s}}$, refer to
this adjacency. We put
\begin{equation*}
\hat{b}:=2dL^{d-1}+2d(h_{\sharp}+2)^{d-1}.
\end{equation*}
A member has at most $\hat{b}$ adjacent members. Indeed, by the level property, across each of its $2d$
walls it has at most $L^{d-1}$ wall-neighbours. Moreover, a member sharing a wall with a hub is a source
cell by the third item, hence of level $i$, that is, a $\pi_i$-cube, and it shares a wall with at most one
hub, since hubs are at least $3$ apart. Finally, the
level-$i$ cells sharing a wall with a box of sides at most $h_{\sharp}$ number at most
$2d(h_{\sharp}+2)^{d-1}$.

For $\Delta\in\pi^{\mathfrak{s}}$, let $\operatorname{pa}(\Delta)$, the \emph{parent} of $\Delta$, be the
$\pi_i$-cube containing
$\Delta$. Let $\rho(\Delta)$ be the least cardinality of a connected subfamily of $\pi^{\mathfrak{s}}$
containing both $\Delta$ and a source cell, with $\rho(\Delta)=\infty$ if there is none. Then
\begin{equation}\label{4.21:eq-site-geometry-b}
\begin{gathered}
\rho(\Delta)\ \ge\ i-\operatorname{lev}(\Delta)+1,\qquad
\rho(\Delta)\ \ge\ \operatorname{dist}_i\Bigl(\Delta,\textstyle\bigcup\mathfrak{F}\Bigr),\\
\#\bigl\{\mathsf{y}\subset\pi^{\mathfrak{s}}:\ \mathsf{y}\text{ connected},\ \#\mathsf{y}=n,\
\Delta\in\mathsf{y}\bigr\}\ \le\ (e\hat{b})^{n}\qquad(n\ge1).
\end{gathered}
\end{equation}

\emph{Readers.} Put $c_{\rho}:=120$, so that $c_{\rho}\ge h_{\sharp}$. A member \emph{sees} a closed set
$\mathcal{R}$ in either of two cases: it meets $\mathcal{R}$; or it is a source cell sharing a wall with a
hub that meets $\mathcal{R}$.

A \emph{reader} is a triple $F=(A_F,\mathcal{R}_F,\|F\|)$ with three components:
\begin{itemize}
\item $A_F\in\mathbf{D}_i$, a domain of Ba{\l}aban's class of localisation domains at level $i$, called
the \emph{label} of $F$;
\item $\mathcal{R}_F$, a closed set, the set read by $F$;
\item $\|F\|\ge0$.
\end{itemize}
These are subject to the condition
\begin{equation}\label{4.21:eq-site-geometry-a}
\text{every }\Delta\in\pi^{\mathfrak{s}}\text{ seeing }\mathcal{R}_F\text{ satisfies }\
\operatorname{dist}_i\bigl(\operatorname{pa}(\Delta),A_F\bigr)\le c_{\rho}\,\rho(\Delta).
\end{equation}

Let $\mathfrak{Y}(F)$ be the set of finite non-empty $\mathsf{y}\subset\pi^{\mathfrak{s}}$ such that every
component of $\mathsf{y}$ contains a source cell and a member seeing $\mathcal{R}_F$. Fix a total order on
$\pi^{\mathfrak{s}}$. For $\mathsf{y}\in\mathfrak{Y}(F)$ with components $\mathsf{y}_a$, we define the
following objects.
\begin{itemize}
\item $\Delta_a$ is the first member of $\mathsf{y}_a$ seeing $\mathcal{R}_F$.
\item $Q_a:=\operatorname{pa}(\Delta_a)$ and $x_a:=\operatorname{dist}_i(Q_a,A_F)$.
\item $\mathrm{Cor}_a$ is the set of intermediate cubes of a chain of $\pi_i$-cubes, of the kind
constructed in the geometry lemma of this chapter, from $Q_a$ to a nearest cube of $A_F$. The chain is
chosen by a fixed rule, and $\mathrm{Cor}_a=\emptyset$ if $Q_a\subset A_F$.
\item Fix a spanning tree of $\mathsf{y}_a$. Consider an edge of this tree joining two members that share
no wall. Both are then source cells at one hub, and their parents are at most $h_{\sharp}$ apart. Take the
intermediate cubes of a chain of the same kind between the two parents. $\mathrm{Hub}_a$ is the union of
these sets over all such edges.
\end{itemize}
With $[\mathsf{y}]_i$ the $\pi_i$-hull of $\bigcup\mathsf{y}$, the \emph{label} of the pair
$(F,\mathsf{y})$ is
\begin{equation}\label{4.21:eq-site-geometry-1}
\operatorname{lab}(F,\mathsf{y}):=A_F\cup[\mathsf{y}]_i\cup\bigcup_a\bigl(\mathrm{Cor}_a\cup
\mathrm{Hub}_a\bigr).
\end{equation}

Let $\mathfrak{L}$ be a family of readers and let $r\ge0$. With $d_i(\cdot)$ as in the glossary
(Notation~\ref{0.2:not-glossary}) and $e^{-\infty}:=0$, we put
\begin{equation}\label{4.21:eq-site-geometry-2}
\begin{aligned}
\mathsf{T}^{*}_{r}[\mathfrak{L}]&:=\sup_{\Delta\in\pi^{\mathfrak{s}}}e^{-\rho(\Delta)}
\sum_{F\in\mathfrak{L}:\ \Delta\text{ sees }\mathcal{R}_F}\|F\|\,e^{r\,d_i(A_F)},\\
\mathsf{Q}_{r}[\mathfrak{L}]&:=\sup_{Q\in\pi_i}\ \sum_{F\in\mathfrak{L}:\ A_F\ni Q}\|F\|\,
e^{r\,d_i(A_F)}.
\end{aligned}
\end{equation}

\emph{Constants.} We put
\begin{equation}\label{4.21:eq-site-geometry-3}
c_{l}:=1+d(h_{\sharp}+1)+d(c_{\rho}+1)=1+234d,
\end{equation}
a number depending on $d$ only. With sums over integers $x\ge0$, we put
\begin{equation}\label{4.21:eq-site-geometry-4}
S_{d}:=\sum_{x\ge0}(2x+3)^{d}e^{-x},\qquad C_{\Sigma}:=2e^{2}S_{d}+e,\qquad
C_{e}:=\sup_{\rho\ge1}(2c_{\rho}\rho+3)^{d}e^{-\rho},
\end{equation}
where in $C_e$ the letter $\rho$ is a free variable. For $\lambda>0$ let
\begin{equation}\label{4.21:eq-site-geometry-5}
\lambda':=\lambda-\log(e\hat{b}),\qquad
\varepsilon_{\lambda}:=4e^{-\frac12\lambda'}\sum_{x\ge0}(2x+3)^{d}e^{-\lambda'x/2c_{\rho}},
\end{equation}
with $\varepsilon_\lambda=+\infty$ when the series diverges. Let $\mathsf{c}_d$ denote the counting
constant, depending on $d$ only, of the position-weight counting estimate of this chapter. Let
$\lambda_{*}(d,L)$ be the least number
such that, for all $\lambda\ge\lambda_{*}(d,L)$, the following four conditions hold:
\begin{equation}\label{4.21:eq-site-geometry-c}
\begin{gathered}
L^{d}e^{-\frac12\lambda'}\le\tfrac12,\qquad \varepsilon_{\lambda}\,\mathsf{c}_{d}\le1,\qquad
\lambda'-c_{\rho}-1\ge d\log L+\log2,\\
\sum_{n\ge1}n\,(2c_{\rho}n+3)^{d}\,(e^{2}\hat{b}L^{d})^{n}e^{-\lambda n}\le1.
\end{gathered}
\end{equation}
Each of these conditions is monotone in $\lambda$, and $\lambda_{*}(d,L)\le C(d)(1+\log L)$ with $C(d)$
depending on $d$ only.
\end{definition}

\begin{proof}[Proof of \eqref{4.21:eq-site-geometry-b} and of the assertions on $\lambda_{*}(d,L)$]
Let $\Delta\in\pi^{\mathfrak{s}}$. If $\rho(\Delta)=\infty$, the two inequalities in the first line of
\eqref{4.21:eq-site-geometry-b} hold trivially. Otherwise, let $\mathsf{y}$ be a connected subfamily with
$\rho(\Delta)$ members containing $\Delta$ and a source cell $\Delta'$. Since $\mathsf{y}$ is connected,
there is a path $\Delta=\Delta_0,\Delta_1,\dots,\Delta_n=\Delta'$ in $\mathsf{y}$ in which consecutive
members are adjacent. Choosing a shortest such path, its members are distinct, so
$n+1\le\rho(\Delta)$.

First inequality. Consider a step from $\Delta_j$ to $\Delta_{j+1}$. If the two members share a wall,
their levels differ by at most one, by the level property. If they do not share a wall, both share a wall
with one hub. Then $\mathfrak{H}\neq\emptyset$, so both are source cells, and both have level $i$. In
either case $|\operatorname{lev}(\Delta_{j+1})-\operatorname{lev}(\Delta_j)|\le1$. Since
$\operatorname{lev}(\Delta_n)=i$, summing along the path gives
\begin{equation*}
i-\operatorname{lev}(\Delta)\le n\le\rho(\Delta)-1.
\end{equation*}

Second inequality. Let $m\le n$ be the least index for which $\Delta_m$ is a source cell. For $j<m$ the
member $\Delta_j$ is not a source cell. If $\mathfrak{H}\neq\emptyset$, this means that $\Delta_j$
shares no wall with a hub. Hence for $j<m$ the adjacency of $\Delta_j$ and $\Delta_{j+1}$ is
wall-sharing, and $\Delta_0,\dots,\Delta_m$ is a chain of members in which consecutive members share a
wall. If $m=0$, then $\operatorname{dist}_i(\Delta,\bigcup\mathfrak{F})=0$ and there is nothing to
prove. If $m\ge1$, choose points $z_j\in\Delta_j\cap\Delta_{j+1}$ for $0\le j\le m-1$. For $1\le j\le
m-1$ both $z_{j-1}$ and $z_j$ lie in $\Delta_j$. Since $\Delta_j$ has level at most $i$, its side is at
most one in level-$i$ units, and hence $\operatorname{dist}_i(z_{j-1},z_j)\le1$. As $z_0\in\Delta$ and
$z_{m-1}\in\Delta_m\subset\bigcup\mathfrak{F}$, the triangle inequality gives
\begin{equation*}
\operatorname{dist}_i\Bigl(\Delta,\textstyle\bigcup\mathfrak{F}\Bigr)\le m-1<\rho(\Delta).
\end{equation*}

Third inequality (the counting bound). We fix a total order on $\pi^{\mathfrak{s}}$. To each member
$\Delta''$ we attach a list of $\hat{b}$ slots. The first slots hold the members adjacent to $\Delta''$
in increasing order, and the remaining slots are blank; this is possible because a member has at most
$\hat{b}$ adjacent members.

Let $\mathsf{y}$ be connected with $\#\mathsf{y}=n$ and $\Delta\in\mathsf{y}$. We explore $\mathsf{y}$
as follows. Initially only $u_1:=\Delta$ is discovered. At step $t$ we process $u_t$, the $t$-th
discovered member. For each slot $s=1,\dots,\hat{b}$ of its list we set $\varepsilon_{t,s}:=1$ if the
slot holds a member of $\mathsf{y}$ not yet discovered, and we then declare that member discovered, the
next in order. Otherwise we set $\varepsilon_{t,s}:=0$.

At the start of step $t\le n$ at least $t$ members are discovered. We argue by induction on $t$; for
$t=1$ this holds since $u_1$ is discovered. Let $2\le t\le n$. By the inductive hypothesis at least $t-1$
members were discovered at the start of step $t-1$, hence at least $t-1$ are discovered at the start of
step $t$. Suppose that exactly $t-1$ are. They are then $u_1,\dots,u_{t-1}$, all of them already
processed. Since $t-1<n$ and $\mathsf{y}$ is connected, some undiscovered member of $\mathsf{y}$ is
adjacent to a processed one, and it would have been discovered when that member was processed, a
contradiction. Hence the exploration runs for $n$ steps and discovers every member of $\mathsf{y}$. It
therefore produces a word $(\varepsilon_{t,s})\in\{0,1\}^{n\hat{b}}$ with exactly $n-1$ entries equal to
$1$.

The word determines $\mathsf{y}$: rerunning the exploration from $\Delta$, and reading membership from
the word instead of from $\mathsf{y}$, reproduces the discovered members in order. Hence the number in
question is at most $\binom{n\hat{b}}{n-1}$. For $n=1$ this is $1\le e\hat{b}$. For $n\ge2$ we use
$\binom{N}{k}\le(eN/k)^k$, $(1+\frac1{n-1})^{n-1}\le e$ and $\hat{b}\ge1$:
\begin{equation*}
\binom{n\hat{b}}{n-1}\le\Bigl(\frac{e\,n\hat{b}}{n-1}\Bigr)^{n-1}
=(e\hat{b})^{n-1}\Bigl(1+\frac1{n-1}\Bigr)^{n-1}\le e\,(e\hat{b})^{n-1}\le(e\hat{b})^{n}.
\end{equation*}

Assertions on $\lambda_{*}(d,L)$. Since $\lambda'$ increases with $\lambda$, the left-hand sides of the
first condition and of the last condition in \eqref{4.21:eq-site-geometry-c} are non-increasing in
$\lambda$. The same holds for $\varepsilon_{\lambda}$: for $\lambda'>0$ it is a sum of terms decreasing in
$\lambda'$, and for $\lambda'\le0$ the series diverges, so $\varepsilon_\lambda=+\infty$. The left-hand
side of the third condition is increasing in $\lambda$. Each condition is closed in $\lambda$. The first
and third are continuous conditions. On $\lambda'>0$ the quantity $\varepsilon_{\lambda}$ is continuous,
because its series converges uniformly on $\lambda'\ge\delta$ for each $\delta>0$, and
$\varepsilon_\lambda\to+\infty$ as $\lambda'\downarrow0$. The last series consists of non-negative
continuous terms, so its sum is lower semicontinuous. Hence the set of $\lambda$ satisfying all four
conditions is a closed half-line, once it is non-empty, and $\lambda_{*}(d,L)$ is its left end.

It remains to show non-emptiness together with the bound. Since $L\ge1$, we have
$\hat{b}\le(2d+2d(h_{\sharp}+2)^{d-1})L^{d-1}$, so $\log(e\hat{b})\le C_1(d)(1+\log L)$. Thus
$\lambda'\ge\lambda-C_1(d)(1+\log L)$. We check the four conditions in turn.
\begin{itemize}
\item The first condition holds as soon as $\lambda'\ge2d\log L+2\log2$.
\item The third condition holds as soon as $\lambda'\ge c_{\rho}+1+d\log L+\log2$.
\item For the second condition, let $\lambda'\ge1$. Then
\begin{equation*}
\varepsilon_\lambda\le4e^{-\frac12\lambda'}\sum_{x\ge0}(2x+3)^de^{-x/2c_\rho}.
\end{equation*}
Since $\mathsf{c}_d$ depends on $d$ only, the second condition holds once $\lambda'$ exceeds a number
depending on $d$ only.
\item For the last condition, write $\lambda=\log(e^2\hat{b}L^d)+\mu$. The sum becomes
$\sum_{n\ge1}n(2c_\rho n+3)^de^{-\mu n}$. For $\mu>0$ this sum is finite and decreases to $0$ as
$\mu\to\infty$, by monotone convergence. Hence it is at most $1$ for $\mu\ge\mu_0(d)$, and
$\log(e^2\hat{b}L^d)\le C_2(d)(1+\log L)$.
\end{itemize}
All four conditions therefore hold for $\lambda\ge C(d)(1+\log L)$, with $C(d)$ depending on $d$ only.
This gives $\lambda_{*}(d,L)\le C(d)(1+\log L)$.
\end{proof}

\begin{lemma}[Summation over source-anchored families]\label{4.21:lem-source-anchored-sums}
Let $(\pi^{\mathfrak{s}},\mathfrak{H},\mathfrak{F})$ be a site geometry of level~$i$, let $F$ be a
reader and let $\mathsf{y}\in\mathfrak{Y}(F)$, all in the sense of
Definition~\ref{4.21:def-site-geometry}. The label $\operatorname{lab}(F,\mathsf{y})$, the functionals
$\mathsf{T}^{*}_{r}$ and $\mathsf{Q}_{r}$, the constants $c_{l}$, $S_{d}$, $C_{\Sigma}$, $C_{e}$,
$\mathsf{c}_{d}$ and $\lambda_{*}=\lambda_{*}(d,L)$, and, for $\lambda>0$, the numbers $\lambda'$ and
$\varepsilon_{\lambda}$ are those of that definition; $|\mathsf{y}|$ is the number of members of
$\mathsf{y}$.
\begin{itemize}
\item[(a)] (Label.) $\operatorname{lab}(F,\mathsf{y})\in\mathbf{D}_i$, and
\begin{equation}\label{4.21:eq-source-anchored-sums-1}
d_i(\operatorname{lab}(F,\mathsf{y}))\le d_i(A_F)+c_{l}|\mathsf{y}|,\qquad
N_i(\operatorname{lab}(F,\mathsf{y}))\le N_i(A_F)+c_{l}|\mathsf{y}|.
\end{equation}
\item[(b)] (Free sum.) Let $\lambda\ge\lambda_{*}$. For $\Delta\in\pi^{\mathfrak{s}}$ let
$\varphi_\lambda(\Delta)$ be the sum of $e^{-\lambda|\mathsf{y}_1|}$ over all connected subfamilies
$\mathsf{y}_1\subset\pi^{\mathfrak{s}}$ that contain $\Delta$ and a source cell, and let $\Phi_F$ be the
sum of $\varphi_\lambda(\Delta)$ over all $\Delta\in\pi^{\mathfrak{s}}$ that see $\mathcal{R}_F$. Then
\begin{align}
\sum_{\mathsf{y}\in\mathfrak{Y}(F)}e^{-\lambda|\mathsf{y}|}&\le e^{\Phi_F}-1,
\label{4.21:eq-source-anchored-sums-2}\\
\Phi_F&\le\varepsilon_{\lambda}N_i(A_F)\le d_i(A_F)+1,
\label{4.21:eq-source-anchored-sums-3}\\
\varphi_\lambda(\Delta)&\le 2e^{-\lambda'\rho(\Delta)}.
\label{4.21:eq-source-anchored-sums-4}
\end{align}
\item[(c)] (Cube-anchored $\ell^1$-sum.) Let $\mathfrak{L}$ be a family of readers, let $r\ge0$, and
assume $\kappa_1-1\ge c_{l}r+\lambda_{*}$. Then for every $Q_*\in\pi_i$
\begin{equation}\label{4.21:eq-source-anchored-sums-5}
\sum_{F\in\mathfrak{L}}\ \sum_{\substack{\mathsf{y}\in\mathfrak{Y}(F):\\
\operatorname{lab}(F,\mathsf{y})\ni Q_*}}
\|F\|\,e^{-(\kappa_1-1)|\mathsf{y}|}\,e^{r\,d_i(\operatorname{lab}(F,\mathsf{y}))}
\le C_{\Sigma}\,\mathsf{T}^{*}_{r+2}[\mathfrak{L}].
\end{equation}
\item[(d)] (Smallness.) If moreover $d_i(\operatorname{lab}(F,\mathsf{y}))\ge D$ for all pairs
$(F,\mathsf{y})$ occurring in the left member of \eqref{4.21:eq-source-anchored-sums-5}, and
$\kappa_1-1\ge c_{l}(r+\varsigma)+\lambda_{*}$, then the left member of
\eqref{4.21:eq-source-anchored-sums-5} is at most
$e^{-\varsigma D}C_{\Sigma}\,\mathsf{T}^{*}_{r+\varsigma+2}[\mathfrak{L}]$.
\item[(e)] (Closure under reception.) Let $\mathfrak{L}'$ be the family, indexed by the pairs
$(F,\mathsf{y})$ with $F\in\mathfrak{L}$ and $\mathsf{y}\in\mathfrak{Y}(F)$, of the triples
$\bigl(\operatorname{lab}(F,\mathsf{y}),\mathcal{R}',\|F\|e^{-(\kappa_1-1)|\mathsf{y}|}\bigr)$, with an
arbitrary closed set $\mathcal{R}'$ as reading set (the functional $\mathsf{Q}_r$ does not depend on
it). Under the hypotheses of (c),
\eqref{4.21:eq-source-anchored-sums-5} reads
\begin{equation}\label{4.21:eq-source-anchored-sums-6}
\mathsf{Q}_{r}[\mathfrak{L}']\le C_{\Sigma}\,\mathsf{T}^{*}_{r+2}[\mathfrak{L}].
\end{equation}
Moreover, for every family $\mathfrak{L}$, in every site geometry of level $i$, and with any reading
sets for which the reader condition \eqref{4.21:eq-site-geometry-a} holds,
\begin{equation}\label{4.21:eq-source-anchored-sums-7}
\mathsf{T}^{*}_{r}[\mathfrak{L}]\le C_{e}\,\mathsf{Q}_{r}[\mathfrak{L}].
\end{equation}
\end{itemize}
The constants $c_{l}$, $S_{d}$, $C_{\Sigma}$, $C_{e}$ depend on $d$ only and $\lambda_{*}$ on $d$ and
$L$ only; none of $c_{l}$, $S_{d}$, $C_{\Sigma}$, $C_{e}$, $\lambda_{*}$ depends on the site geometry
$(\pi^{\mathfrak{s}},\mathfrak{H},\mathfrak{F})$, on the reading sets, or on the family $\mathfrak{L}$.
\end{lemma}

\begin{proof}
Throughout, the $\mathsf{y}_a$ are the components of $\mathsf{y}$, and $\Delta_a$,
$Q_a=\operatorname{pa}(\Delta_a)$, $x_a=\operatorname{dist}_i(Q_a,A_F)$, $\mathrm{Cor}_a$ and
$\mathrm{Hub}_a$ are as in Definition~\ref{4.21:def-site-geometry}, so that
\begin{equation*}
\operatorname{lab}(F,\mathsf{y})=A_F\cup[\mathsf{y}]_i\cup\bigcup_a(\mathrm{Cor}_a\cup\mathrm{Hub}_a),
\qquad [\mathsf{y}]_i=\bigcup_a[\mathsf{y}_a]_i .
\end{equation*}

\emph{Proof of (a).} Fix a component $\mathsf{y}_a$ and its fixed spanning tree. If an edge of the
tree joins two members that share a wall, their parents coincide or share a wall. Every other edge
joins two members that share no wall; by the definition of adjacency both are then source cells
sharing a wall with one hub, their parents are at most $h_{\sharp}$ apart, and $\mathrm{Hub}_a$
contains the intermediate cubes of a chain, as provided by clause (g5) of the geometry lemma, between
the two parents. Hence $[\mathsf{y}_a]_i\cup\mathrm{Hub}_a$ is wall-connected. It contains $Q_a$, the
parent of $\Delta_a\in\mathsf{y}_a$, and $\mathrm{Cor}_a$ consists of the intermediate cubes of a chain
of the same kind from $Q_a$ to a nearest cube of $A_F$ (and is empty if $Q_a\subset A_F$). Thus
$\operatorname{lab}(F,\mathsf{y})$ arises from $A_F$ by adjoining, for each $a$, the wall-connected set
$[\mathsf{y}_a]_i\cup\mathrm{Cor}_a\cup\mathrm{Hub}_a$, which is chained to $A_F$.

We count the adjoined cubes. The family $\mathsf{y}_a$ is connected, contains $\Delta_a$ and, by the
definition of $\mathfrak{Y}(F)$, a source cell; by the definition of $\rho$,
$|\mathsf{y}_a|\ge\rho(\Delta_a)$. Since $\Delta_a$ sees $\mathcal{R}_F$, the reader condition
\eqref{4.21:eq-site-geometry-a} gives $x_a\le c_{\rho}\rho(\Delta_a)\le c_{\rho}|\mathsf{y}_a|$.
The set $[\mathsf{y}_a]_i$ has at most $|\mathsf{y}_a|$ cubes, one parent per member. The spanning tree
has fewer than $|\mathsf{y}_a|$ edges, and by clause (g5) the chain bridging an edge contributes at
most $d(h_{\sharp}+1)$ cubes; by the same clause $\mathrm{Cor}_a$ has at most $d(x_a+1)$ cubes. Using
$|\mathsf{y}_a|\ge1$ and $c_{l}=1+d(h_{\sharp}+1)+d(c_{\rho}+1)$, the cubes adjoined for
$\mathsf{y}_a$ number at most
\begin{equation*}
|\mathsf{y}_a|+|\mathsf{y}_a|\,d(h_{\sharp}+1)+d(x_a+1)
\le |\mathsf{y}_a|+|\mathsf{y}_a|\,d(h_{\sharp}+1)+d(c_{\rho}|\mathsf{y}_a|+1)
\le c_{l}|\mathsf{y}_a| .
\end{equation*}
Summing over $a$, at most $c_{l}|\mathsf{y}|$ cubes are adjoined to $A_F$. Clause (g2) of the
geometry lemma, applied to $A_F\in\mathbf{D}_i$ and to these adjoined sets, gives
$\operatorname{lab}(F,\mathsf{y})\in\mathbf{D}_i$ and \eqref{4.21:eq-source-anchored-sums-1}.

\emph{Proof of (b).} A member $\mathsf{y}$ of $\mathfrak{Y}(F)$ is the union of its components, which
are pairwise distinct connected families, each containing a source cell and a member seeing
$\mathcal{R}_F$, and $|\mathsf{y}|=\sum_a|\mathsf{y}_a|$. Let $\Psi_F$ be the sum of
$e^{-\lambda|\mathsf{y}_1|}$ over all connected $\mathsf{y}_1\subset\pi^{\mathfrak{s}}$ containing a
source cell and a member seeing $\mathcal{R}_F$. A set of $n$ pairwise distinct such families
corresponds to $n!$ ordered $n$-tuples, so
\begin{equation*}
\sum_{\mathsf{y}\in\mathfrak{Y}(F)}e^{-\lambda|\mathsf{y}|}\le\sum_{n\ge1}\frac{\Psi_F^n}{n!}
=e^{\Psi_F}-1 .
\end{equation*}
Every such $\mathsf{y}_1$ contains a member $\Delta$ seeing $\mathcal{R}_F$, and
$e^{-\lambda|\mathsf{y}_1|}$ is then a term of $\varphi_\lambda(\Delta)$; hence $\Psi_F\le\Phi_F$, and
\eqref{4.21:eq-source-anchored-sums-2} follows.

For \eqref{4.21:eq-source-anchored-sums-4}: a connected $\mathsf{y}_1$ containing $\Delta$ and a source
cell has $n\ge\rho(\Delta)$ members, and by the third property in \eqref{4.21:eq-site-geometry-b}
there are at most $(e\hat{b})^n$ connected subfamilies with $n$ members containing $\Delta$. With
$\lambda'=\lambda-\log(e\hat{b})$,
\begin{equation*}
\varphi_\lambda(\Delta)\le\sum_{n\ge\rho(\Delta)}(e\hat{b})^ne^{-\lambda n}
=\frac{e^{-\lambda'\rho(\Delta)}}{1-e^{-\lambda'}}\le 2e^{-\lambda'\rho(\Delta)},
\end{equation*}
because the first condition in \eqref{4.21:eq-site-geometry-c}, $L^de^{-\lambda'/2}\le\frac12$, gives
$e^{-\lambda'/2}\le\frac12$, hence $\lambda'\ge\log2$.

For \eqref{4.21:eq-source-anchored-sums-3} we sort the members $\Delta$ seeing $\mathcal{R}_F$ by
their parent $Q=\operatorname{pa}(\Delta)\in\pi_i$ and by $t:=i-\operatorname{lev}(\Delta)\ge0$;
a cube $Q$ contains at most $L^{dt}$ members $\Delta$ with $\operatorname{lev}(\Delta)=i-t$. Put
$x:=\operatorname{dist}_i(Q,A_F)$. The reader condition \eqref{4.21:eq-site-geometry-a} gives
$x\le c_{\rho}\rho(\Delta)$, and the first property in \eqref{4.21:eq-site-geometry-b} gives
$\rho(\Delta)\ge t+1$; therefore
\begin{equation*}
\rho(\Delta)\ge\max\Bigl\{\frac{x}{c_{\rho}},\,t+1\Bigr\}
\ge\frac{x}{2c_{\rho}}+\frac{t+1}{2}.
\end{equation*}
With \eqref{4.21:eq-source-anchored-sums-4} and the first condition in
\eqref{4.21:eq-site-geometry-c},
\begin{align*}
\sum_{\substack{\Delta\subset Q:\\ \Delta\text{ sees }\mathcal{R}_F}}\varphi_\lambda(\Delta)
&\le 2e^{-\lambda'x/(2c_{\rho})}e^{-\lambda'/2}\sum_{t\ge0}\bigl(L^de^{-\lambda'/2}\bigr)^t
\le 4e^{-\lambda'/2}e^{-\lambda'x/(2c_{\rho})}.
\end{align*}
The cubes of $\pi_i$ at distance $x$ from $A_F$ number at most $N_i(A_F)(2x+3)^d$. Summing over
$x\ge0$,
\begin{equation*}
\Phi_F\le N_i(A_F)\,4e^{-\lambda'/2}\sum_{x\ge0}(2x+3)^de^{-\lambda'x/(2c_{\rho})}
=\varepsilon_{\lambda}N_i(A_F).
\end{equation*}
The second condition in \eqref{4.21:eq-site-geometry-c}, $\varepsilon_{\lambda}\mathsf{c}_d\le1$,
together with clause (g1) of the geometry lemma, gives $\varepsilon_{\lambda}N_i(A_F)\le
d_i(A_F)+1$.

\emph{Proof of (c).} By \eqref{4.21:eq-source-anchored-sums-1} and $r\ge0$,
$e^{r\,d_i(\operatorname{lab}(F,\mathsf{y}))}\le e^{r\,d_i(A_F)}e^{c_{l}r|\mathsf{y}|}$. Put
$\lambda:=\kappa_1-1-c_{l}r$; by hypothesis $\lambda\ge\lambda_{*}$, so (b) applies with this
$\lambda$, and each term of \eqref{4.21:eq-source-anchored-sums-5} is at most
$\|F\|e^{r\,d_i(A_F)}e^{-\lambda|\mathsf{y}|}$. If $\operatorname{lab}(F,\mathsf{y})\ni Q_*$, then
either (i) $Q_*\in[\mathsf{y}_a]_i\cup\mathrm{Cor}_a\cup\mathrm{Hub}_a$ for some $a$, or (ii)
$Q_*\in A_F$. We bound the contributions of the two cases separately.

\emph{Case (ii).} We sum over all $\mathsf{y}\in\mathfrak{Y}(F)$ and over the $F\in\mathfrak{L}$
with $A_F\ni Q_*$. By \eqref{4.21:eq-source-anchored-sums-2}, the inequality $e^u-1\le ue^u$ for
$u\ge0$, and \eqref{4.21:eq-source-anchored-sums-3},
\begin{equation*}
\sum_{\mathsf{y}\in\mathfrak{Y}(F)}e^{-\lambda|\mathsf{y}|}\le\Phi_Fe^{\Phi_F}
\le e\,e^{d_i(A_F)}\sum_{\Delta\text{ sees }\mathcal{R}_F}\varphi_\lambda(\Delta).
\end{equation*}
Interchanging the sums over $F$ and $\Delta$, the contribution of Case (ii) is at most
\begin{equation*}
e\sum_{\Delta\in\pi^{\mathfrak{s}}}\varphi_\lambda(\Delta)
\sum_{\substack{F\in\mathfrak{L}:\ A_F\ni Q_*,\\ \Delta\text{ sees }\mathcal{R}_F}}
\|F\|\,e^{(r+1)d_i(A_F)}.
\end{equation*}
For such $F$, the set $A_F$ contains $Q_*$ and, by the reader condition
\eqref{4.21:eq-site-geometry-a}, a cube $Q'$ with
$\operatorname{dist}_i(Q',\operatorname{pa}(\Delta))\le c_{\rho}\rho(\Delta)$. Clause (g4) of the
geometry lemma, applied to the cubes $Q_*$, $Q'$ of $A_F$, together with the triangle inequality
\begin{equation*}
\operatorname{dist}_i(Q_*,\operatorname{pa}(\Delta))\le\operatorname{dist}_i(Q_*,Q')
+\operatorname{dist}_i(Q',\operatorname{pa}(\Delta)),
\end{equation*}
gives
\begin{equation*}
d_i(A_F)\ge\operatorname{dist}_i(Q_*,\operatorname{pa}(\Delta))-c_{\rho}\rho(\Delta)-1.
\end{equation*}
Hence
\begin{equation*}
e^{(r+1)d_i(A_F)}=e^{(r+2)d_i(A_F)}e^{-d_i(A_F)}
\le e^{(r+2)d_i(A_F)}e^{-\operatorname{dist}_i(Q_*,\operatorname{pa}(\Delta))+c_{\rho}\rho(\Delta)+1}.
\end{equation*}
Dropping the condition $A_F\ni Q_*$, the definition of $\mathsf{T}^{*}_{r+2}$ bounds the sum of
$\|F\|e^{(r+2)d_i(A_F)}$ over the $F$ with $\Delta$ seeing $\mathcal{R}_F$ by
$e^{\rho(\Delta)}\mathsf{T}^{*}_{r+2}[\mathfrak{L}]$. With
\eqref{4.21:eq-source-anchored-sums-4}, grouping the members $\Delta$ by their parent $Q$, the
contribution of Case (ii) is at most
\begin{equation*}
2e^2\,\mathsf{T}^{*}_{r+2}[\mathfrak{L}]\sum_{Q\in\pi_i}e^{-\operatorname{dist}_i(Q_*,Q)}
\sum_{\Delta\subset Q}e^{-(\lambda'-c_{\rho}-1)\rho(\Delta)}.
\end{equation*}
In the inner sum there are at most $L^{dt}$ members of level $i-t$ in $Q$, each with
$\rho(\Delta)\ge t+1$ by the first property in \eqref{4.21:eq-site-geometry-b}, and the third
condition in \eqref{4.21:eq-site-geometry-c} gives $\lambda'-c_{\rho}-1\ge d\log L+\log2$; hence
\begin{equation*}
\sum_{\Delta\subset Q}e^{-(\lambda'-c_{\rho}-1)\rho(\Delta)}
\le\sum_{t\ge0}L^{dt}e^{-(d\log L+\log2)(t+1)}=L^{-d}\sum_{t\ge0}2^{-(t+1)}\le1.
\end{equation*}
At most $(2x+3)^d$ cubes of $\pi_i$ lie at distance $x$ from $Q_*$, so
$\sum_{Q}e^{-\operatorname{dist}_i(Q_*,Q)}\le S_{d}$, and the contribution of Case (ii) is at most
$2e^2S_{d}\,\mathsf{T}^{*}_{r+2}[\mathfrak{L}]$.

\emph{Case (i).} Let $n:=|\mathsf{y}_a|$. The cubes of $[\mathsf{y}_a]_i$ are parents of members of
$\mathsf{y}_a$; by clause (g5) the cubes of $\mathrm{Cor}_a$ lie within $x_a$ of
$Q_a=\operatorname{pa}(\Delta_a)$, and those of $\mathrm{Hub}_a$ within $h_{\sharp}$ of a parent. Since
$x_a\le c_{\rho}n$ (proof of (a)) and $h_{\sharp}\le c_{\rho}$, every cube of
$[\mathsf{y}_a]_i\cup\mathrm{Cor}_a\cup\mathrm{Hub}_a$ is within
$\max\{x_a,h_{\sharp}\}\le c_{\rho}n$ of the parent of a member of $\mathsf{y}_a$. So $\mathsf{y}_a$
has a member $\Delta_1$ with $\operatorname{dist}_i(\operatorname{pa}(\Delta_1),Q_*)\le c_{\rho}n$, and
$\operatorname{pa}(\Delta_1)$ is one of at most $(2c_{\rho}n+3)^d$ cubes. Since $\mathsf{y}_a$ is
connected and contains $\Delta_1$ and a source cell, the first property in
\eqref{4.21:eq-site-geometry-b} gives $n\ge\rho(\Delta_1)\ge i-\operatorname{lev}(\Delta_1)+1$; inside
a given parent the candidates for $\Delta_1$ therefore number at most
$\sum_{t<n}L^{dt}\le L^{dn}$. Given $\Delta_1$, there are at most $(e\hat{b})^n$ connected families
$\mathsf{y}_a$ with $n$ members containing $\Delta_1$, by the third property in
\eqref{4.21:eq-site-geometry-b}.

Given $\mathsf{y}_a$, the reader $F$ has $\mathcal{R}_F$ seen by one of the $n$ members $\Delta$ of
$\mathsf{y}_a$, and each of them has $\rho(\Delta)\le n$, because $\mathsf{y}_a$ is connected, has $n$
members and contains a source cell. The union of the other components of $\mathsf{y}$ is empty or a
member of $\mathfrak{Y}(F)$, so by \eqref{4.21:eq-source-anchored-sums-2} and
\eqref{4.21:eq-source-anchored-sums-3} the sum over them of $e^{-\lambda|\mathsf{y}|}$ is at most
$e^{-\lambda n}e^{\Phi_F}\le e^{-\lambda n}e\,e^{d_i(A_F)}$. The sum of $\|F\|e^{(r+1)d_i(A_F)}$
over the $F$ with a given $\Delta\in\mathsf{y}_a$ seeing $\mathcal{R}_F$ is at most
$e^{\rho(\Delta)}\mathsf{T}^{*}_{r+1}[\mathfrak{L}]\le e^n\mathsf{T}^{*}_{r+1}[\mathfrak{L}]$. Every
pair $(F,\mathsf{y})$ of Case (i) is reached at least once by choosing $n$,
$\operatorname{pa}(\Delta_1)$, $\Delta_1$, $\mathsf{y}_a$, $F$ and the remaining components; all terms
are non-negative, so by the fourth condition in \eqref{4.21:eq-site-geometry-c} the contribution of
Case (i) is at most
\begin{align*}
&e\sum_{n\ge1}(2c_{\rho}n+3)^dL^{dn}(e\hat{b})^ne^{-\lambda n}\,n\,e^{n}\,
\mathsf{T}^{*}_{r+1}[\mathfrak{L}]\\
&\qquad=e\,\mathsf{T}^{*}_{r+1}[\mathfrak{L}]\sum_{n\ge1}n(2c_{\rho}n+3)^d(e^2\hat{b}L^d)^n
e^{-\lambda n}\\
&\qquad\le e\,\mathsf{T}^{*}_{r+1}[\mathfrak{L}].
\end{align*}

Since $d_i(A_F)\ge0$, $\mathsf{T}^{*}_{r+1}[\mathfrak{L}]\le\mathsf{T}^{*}_{r+2}[\mathfrak{L}]$.
Adding the two cases, the left member of \eqref{4.21:eq-source-anchored-sums-5} is at most
$(2e^2S_{d}+e)\mathsf{T}^{*}_{r+2}[\mathfrak{L}]=C_{\Sigma}\mathsf{T}^{*}_{r+2}[\mathfrak{L}]$.

\emph{Proof of (d).} For every pair occurring, $d_i(\operatorname{lab}(F,\mathsf{y}))\ge D$, so
\begin{equation*}
e^{r\,d_i(\operatorname{lab}(F,\mathsf{y}))}
\le e^{-\varsigma D}e^{(r+\varsigma)\,d_i(\operatorname{lab}(F,\mathsf{y}))}.
\end{equation*}
The left member of
\eqref{4.21:eq-source-anchored-sums-5} is therefore at most $e^{-\varsigma D}$ times the same sum with
$r+\varsigma$ in place of $r$; the hypothesis $\kappa_1-1\ge c_{l}(r+\varsigma)+\lambda_{*}$ is the
hypothesis of (c) at $r+\varsigma$, and (c) at $r+\varsigma$ gives the bound
$e^{-\varsigma D}C_{\Sigma}\mathsf{T}^{*}_{r+\varsigma+2}[\mathfrak{L}]$.

\emph{Proof of (e).} The functional $\mathsf{Q}_r$ reads only the label and the weight of each
member of a family. For the member of $\mathfrak{L}'$ indexed by $(F,\mathsf{y})$ these are
$\operatorname{lab}(F,\mathsf{y})$ and $\|F\|e^{-(\kappa_1-1)|\mathsf{y}|}$, so the left member of
\eqref{4.21:eq-source-anchored-sums-5}, with the supremum over $Q_*\in\pi_i$, is
$\mathsf{Q}_r[\mathfrak{L}']$, and (c) is \eqref{4.21:eq-source-anchored-sums-6}.

For \eqref{4.21:eq-source-anchored-sums-7}, let $\Delta\in\pi^{\mathfrak{s}}$. If $\Delta$ sees
$\mathcal{R}_F$, then by the reader condition \eqref{4.21:eq-site-geometry-a} the set $A_F$ contains
one of the at most $(2c_{\rho}\rho(\Delta)+3)^d$ cubes $Q\in\pi_i$ with
$\operatorname{dist}_i(Q,\operatorname{pa}(\Delta))\le c_{\rho}\rho(\Delta)$. Hence
\begin{align*}
e^{-\rho(\Delta)}\sum_{\substack{F\in\mathfrak{L}:\\ \Delta\text{ sees }\mathcal{R}_F}}
\|F\|e^{r\,d_i(A_F)}
&\le e^{-\rho(\Delta)}\sum_{Q}\ \sum_{F\in\mathfrak{L}:\,A_F\ni Q}\|F\|e^{r\,d_i(A_F)}\\
&\le(2c_{\rho}\rho(\Delta)+3)^de^{-\rho(\Delta)}\,\mathsf{Q}_r[\mathfrak{L}],
\end{align*}
the sum over $Q$ running over those cubes. By the first property in \eqref{4.21:eq-site-geometry-b},
$\rho(\Delta)\ge1$, so the right member is at most $C_{e}\mathsf{Q}_r[\mathfrak{L}]$; taking the
supremum over $\Delta$ gives \eqref{4.21:eq-source-anchored-sums-7}.

Finally, $c_{l}=1+d(h_{\sharp}+1)+d(c_{\rho}+1)$, $S_{d}$, $C_{\Sigma}=2e^2S_{d}+e$ and $C_{e}$
involve only $d$ and the absolute numbers $h_{\sharp}$, $c_{\rho}$, and $\lambda_{*}(d,L)$ is
defined by the four conditions in \eqref{4.21:eq-site-geometry-c}, which involve only $d$ and $L$
(through $\hat{b}$, $c_{\rho}$ and $\mathsf{c}_d$); none of them involves the site geometry, the
reading sets or $\mathfrak{L}$.
\end{proof}

\begin{corollary}[The relative length in Ba{\l}aban's localisation passage]
\label{4.21:cor-localisation-length}
Work in the setting and notation of Ba{\l}aban's localisation passage
\cite[(1.66)--(1.67)]{B-CMP122b}: the sets $\Lambda$, $\Lambda_C$, $Z$, $Y_1,\dots,Y_4$ and
$Y'_2,Y'_3,Y'_4$ are those of that passage. Let a term of the first part of the expansion have
plain domain $X$. Put
\begin{equation*}
  Y_0:=[X]_k,\qquad \mathsf{y}:=\bigl(Y'_2\cup Y'_3\cup Y'_4\bigr)\setminus\Lambda ,
\end{equation*}
and let $Y_4$ be taken together with the components of $Z$ that it meets. Then
\begin{equation}\label{4.21:eq-localisation-length-1}
  d_{k,Z}(Y_4)\;\le\; d_k(Y_0)+c_l\sum_{i=2}^{4}M^{-d}\,|Y'_i| .
\end{equation}
\end{corollary}

\begin{proof}
We first place the passage inside a site geometry in the sense of
Definition~\ref{4.21:def-site-geometry}, with the level of that definition equal to $k$ (so that
its $d_i$, $N_i$ are $d_k$, $N_k$). The family $\pi^{\mathfrak{s}}$
consists of the cells disjoint from $\Lambda$, and the family $\mathfrak{H}$ of hubs consists of
the sets $\Lambda_C$. That the $\Lambda_C$ meet the requirements which
Definition~\ref{4.21:def-site-geometry} places on hubs is clause (f) of the geometric lemma of this
chapter.

In this site geometry, the requirement of \cite[(1.66)--(1.67)]{B-CMP122b} that each component of
each added set $Y'_i$ contain at least one component of $\Lambda$ reads as follows: each component of
$\mathsf{y}$, with respect to the adjacency of Definition~\ref{4.21:def-site-geometry}, contains a
source cell. Thus $\mathsf{y}$ is a family of the kind to which
Lemma~\ref{4.21:lem-source-anchored-sums} applies, once a reader has been specified.

At the localisation producing $Y_i$, $i\in\{2,3,4\}$, we take as reader the pair whose label is
the label obtained so far and whose reading set is the boundary layer of $Y_{i-1}$. We verify the
reader condition
\eqref{4.21:eq-site-geometry-a} for this pair. Let $\Delta$ be a cell meeting the boundary layer
of $Y_{i-1}$. Then $\Delta$ touches the label, or it touches the rim of a component $C\subset
Y_{i-1}$ of $Z$. In the second case the label meets $C$. In either case, by clause (f) of the
geometric lemma, $\Delta$ lies at distance at most $100\check{R}_k+1$ from the label. On the other
hand $\rho(\Delta)\ge 9\check{R}_k-1$, by the same clause together with $(\rho2)$
of \eqref{4.21:eq-site-geometry-b}. These two bounds give \eqref{4.21:eq-site-geometry-a}.

For every such component $C$, the corridor from a rim cube of $C$ to a cube of the label is
taken inside the box $C$, as clause (g5) of Definition~\ref{4.21:def-site-geometry} permits.

We now apply part (a) of Lemma~\ref{4.21:lem-source-anchored-sums} three times, once at each of
the three localisations. Each application gives that the new label belongs to $\mathbf{D}_k$.
It also shows that its length exceeds the length of the previous label by at most $c_l$
times the number of cells cut at that localisation. Adding the three inequalities, and writing
$\operatorname{lab}$ for the final label, we obtain
\begin{equation*}
  d_k(\operatorname{lab})\;\le\; d_k(Y_0)+c_l|\mathsf{y}| .
\end{equation*}

Finally, the construction gives
\begin{equation*}
  Y_4\setminus Z\;\subset\;\operatorname{lab}\;\subset\;Y_4 .
\end{equation*}
Consequently a tree of minimal length for $d_k(\operatorname{lab})$ is admissible in the
minimisation defining $d_{k,Z}(Y_4)$. Hence $d_{k,Z}(Y_4)\le d_k(\operatorname{lab})$. The cells
of $\mathsf{y}$ lie in $Y'_2\cup Y'_3\cup Y'_4$, and the proof of
\eqref{4.21:eq-localisation-length-1} is completed by the bound
\begin{equation*}
  |\mathsf{y}|\;\le\;\sum_{i=2}^{4}M^{-d}\,|Y'_i|
\end{equation*}
on their number.

The corollary is stated for the index sets of \cite[(1.66)--(1.67)]{B-CMP122b}. The condition
$\tfrac12(\kappa_1-1)\ge(1+3\beta)\kappa c_l$, with $\beta$, $\kappa$, $\kappa_1$ the constants of
\cite{B-CMP122b}, which \cite{B-CMP122b} imposes on these index sets, is not used in the argument
above; nor is it a hypothesis of the theorem on the relative class chain. Inequality
\eqref{4.21:eq-localisation-length-1} is \cite[(1.67)]{B-CMP122b} without its third term; in
\cite{B-CMP122b} that term is used in the summation over $X$.
\end{proof}

\begin{definition}[Wrapped terms, probes and depth data]\label{4.21:def-wrapped-terms}
Fix a point of the construction at which the current level is $i$, together with the cells and
probes present there. A probe is a cell, a block or a fattened block, lying inside a strip or not.
Each probe $\mathsf{Q}$ is a subset of the lattice and carries a level, written
$\operatorname{lev}(\mathsf{Q})$. Let $\mathfrak{Q}$ be a class of probes.

A family $\mathfrak{L}$ of \emph{wrapped terms} is a family of objects $F$, each of which carries the
following data:
\begin{itemize}
\item a level $j_F\le i$;
\item a label $A^{\circ}_F\in\mathbf{D}_{j_F}$;
\item a finite non-empty set $\mathcal{P}_F$ of live pockets $(Y,m)$ with $m\le j_F$, whose strips
$Y^{\sharp}$ touch $A^{\circ}_F$;
\item a weight $W(F)\ge 0$.
\end{itemize}
The region of $F$ is
\begin{equation}\label{4.21:eq-wrapped-terms-1}
X_F:=A^{\circ}_F\cup\bigcup_{(Y,m)\in\mathcal{P}_F}Y^{\sharp}.
\end{equation}
The pockets $(Y,m)$, their strips $Y^{\sharp}$ and the zone radius $\check{R}_m$ at level $m$ are
those of the definition of live pockets, their strips and the zone radii.

We assume that there are
\begin{itemize}
\item a function $\mathsf{b}\ge 1$ on $[0,\infty[$;
\item a number $\varsigma_0>0$;
\item constants $C_{\mathsf{b}}$ and $p_{\mathsf{b}}$;
\item numbers $\varrho(\mathsf{Q})\ge 1$, one for each $\mathsf{Q}\in\mathfrak{Q}$, and a number
$N_*\ge 0$;
\end{itemize}
such that the following hold. In the following display $\pi_j$ denotes the partition into cubes of
level $j$, and $d_j(X)$ denotes Ba{\l}aban's size at level $j$ of a domain $X$.
\begin{equation}\label{4.21:eq-wrapped-terms-a}
\begin{gathered}
\textup{(a)}\quad\sum_{F\in\mathfrak{L}:\ j_F=j,\ A^{\circ}_F\ni\square'}W(F)\le\mathsf{b}(i-j)
\qquad\text{for all } j\le i,\ \square'\in\pi_j;\\[3pt]
\textup{(b)}\quad\begin{aligned}[t]
&\sum_{F\in\mathfrak{L}:\ j_F=j,\ A^{\circ}_F\ni\square'}W(F)\,e^{\varsigma_0 d_j(A^{\circ}_F)}
\le\mathsf{b}(i-j)\\
&\text{for all } \mathsf{Q}\in\mathfrak{Q},\ j<\operatorname{lev}(\mathsf{Q}),\
\square'\in\pi_j \text{ with } \square'\subset\mathsf{Q};
\end{aligned}\\[3pt]
\textup{(c)}\quad\mathsf{b}(s+t)\le C_{\mathsf{b}}(1+t)^{p_{\mathsf{b}}}\,\mathsf{b}(s)
\qquad\text{for all } s,t\ge 0;\\[3pt]
\textup{(d)}\quad i-\operatorname{lev}(\mathsf{Q})+1\le\varrho(\mathsf{Q})+N_*
\qquad\text{for all } \mathsf{Q}\in\mathfrak{Q};\\[3pt]
\textup{(e)}\quad\check{R}_m\le 2\varrho(\mathsf{Q})+2
\qquad\text{whenever } \mathsf{Q}\in\mathfrak{Q} \text{ meets the strip of a live pocket } (Y,m).
\end{gathered}
\end{equation}
The hypotheses \textup{(a)}--\textup{(e)} of \eqref{4.21:eq-wrapped-terms-a} are read as follows.
\begin{itemize}
\item Hypothesis \textup{(a)}, the one-level bound, bounds the total weight of the terms of one
level $j$ whose label passes through one cube of that level.
\item Hypothesis \textup{(b)}, the spare-rate bound, asks for a spare decay rate $\varsigma_0$ in
the size $d_j(A^{\circ}_F)$ of the label. It asks this only of labels finer than the probe, that
is, of levels $j<\operatorname{lev}(\mathsf{Q})$ through cubes contained in $\mathsf{Q}$.
\item Hypothesis \textup{(c)} is the growth law of $\mathsf{b}$.
\item Hypotheses \textup{(d)} and \textup{(e)} are the two depth conditions. The numbers
$\varrho(\mathsf{Q})$, $\mathsf{Q}\in\mathfrak{Q}$, together with $N_*$, are called a
\emph{depth datum}.
\end{itemize}

For a real number $x$ write $(x)_+:=\max\{x,0\}$. Put
\begin{equation}\label{4.21:eq-wrapped-terms-2}
C^{<}:=C_{\mathsf{b}}\sum_{t\ge 1}(1+t)^{p_{\mathsf{b}}}\,(5L^t)^d\,
e^{-\varsigma_0\left(\frac12 L^t-2\right)_+}.
\end{equation}
For $\beta>0$ put
\begin{equation}\label{4.21:eq-wrapped-terms-3}
C_{\mathsf{a}}(\beta):=C_{\mathsf{b}}\sup_{\varrho\ge 1}
\Bigl[\varrho\bigl(5^d+2^d(202\varrho+204)^d\bigr)+C^{<}\Bigr](1+\varrho)^{p_{\mathsf{b}}}
e^{-\beta\varrho}.
\end{equation}
In \eqref{4.21:eq-wrapped-terms-3} the supremum runs over real numbers $\varrho\ge 1$; this
$\varrho$ is a variable of the supremum and not a value of the depth datum. The numbers $C^{<}$ and
$C_{\mathsf{a}}(\beta)$ are called the \emph{anchoring constants}; they are finite and depend only
on $(d,L,\varsigma_0,p_{\mathsf{b}},C_{\mathsf{b}},\beta)$.
\end{definition}

\begin{proof}
We verify the last assertion of Definition~\ref{4.21:def-wrapped-terms}, that $C^{<}$ and
$C_{\mathsf{a}}(\beta)$ are finite and depend only on the parameters listed.
We first treat $C^{<}$. Since $(\tfrac12 L^t-2)_+\ge\tfrac12 L^t-2$, the $t$-th term of the series in
\eqref{4.21:eq-wrapped-terms-2} is at most
\begin{equation*}
5^d e^{2\varsigma_0}\,(1+t)^{p_{\mathsf{b}}}\,\bigl(L^{td}e^{-\varsigma_0 L^t/4}\bigr)\,
e^{-\varsigma_0 L^t/4}.
\end{equation*}
For $u>0$ one has $u^d e^{-\varsigma_0 u/4}\le(4d/(e\varsigma_0))^d$, because the left side is
maximal at $u=4d/\varsigma_0$; so the middle factor is bounded uniformly in $t$. The blocking factor
satisfies $L>1$, and $L^t\ge 1+t\log L$, so that
\begin{equation*}
e^{-\varsigma_0 L^t/4}\le e^{-(\varsigma_0\log L/4)\,t}.
\end{equation*}
A power of $1+t$ times this geometric factor is summable over $t\ge 1$. Hence the series converges,
and $C^{<}$ is a finite number depending only on $(d,L,\varsigma_0,p_{\mathsf{b}},C_{\mathsf{b}})$.

We now treat $C_{\mathsf{a}}(\beta)$. For $\varrho\ge 1$ the bracket in
\eqref{4.21:eq-wrapped-terms-3} is a polynomial in $\varrho$, so its product with
$(1+\varrho)^{p_{\mathsf{b}}}$ is bounded by a power of $1+\varrho$. For $\beta>0$ the product
$(1+\varrho)^{q}e^{-\beta\varrho}$ is continuous on $[1,\infty[$ for every real $q$ and tends to
$0$ as $\varrho\to\infty$. It is therefore bounded on $[1,\infty[$, and the supremum is finite. Its
value involves only $d$, $p_{\mathsf{b}}$, $\beta$ and $C^{<}$. Together with the factor
$C_{\mathsf{b}}$, this shows that $C_{\mathsf{a}}(\beta)$ depends only on
$(d,L,\varsigma_0,p_{\mathsf{b}},C_{\mathsf{b}},\beta)$.
\end{proof}

\begin{lemma}[The mass read at one probe]\label{4.21:lem-probe-mass}
Work in the setting of Definition~\ref{4.21:def-wrapped-terms}, at current level $i$, under the
hypotheses collected in \eqref{4.21:eq-wrapped-terms-a}, and assume in addition that the function
$\mathsf{b}$ there is non-decreasing. This is no restriction: the non-decreasing hull of $\mathsf{b}$
satisfies the same hypotheses with the same constants, as shown at the beginning of the proof.
Let $\varsigma_0$, $C_{\mathsf{b}}$ and
$p_{\mathsf{b}}$ be the constants appearing there, and put
\begin{equation}\label{4.21:eq-probe-mass-1}
\begin{aligned}
C^{<} &:= C_{\mathsf{b}}\sum_{t\ge 1}(1+t)^{p_{\mathsf{b}}}\,(5L^{t})^{d}\,
e^{-\varsigma_0\left(\frac12 L^{t}-2\right)_{+}},\\
C_{\mathsf{a}}(\beta) &:= C_{\mathsf{b}}\sup_{\varrho\ge 1}
\Bigl[\varrho\bigl(5^{d}+2^{d}(202\varrho+204)^{d}\bigr)+C^{<}\Bigr]
(1+\varrho)^{p_{\mathsf{b}}}\,e^{-\beta\varrho}\qquad(\beta>0).
\end{aligned}
\end{equation}
These are finite numbers depending only on $(d,L,\varsigma_0,p_{\mathsf{b}},C_{\mathsf{b}},\beta)$.
For every probe $\mathsf{Q}\in\mathfrak{Q}$, with $n:=\operatorname{lev}(\mathsf{Q})$, the following hold.
\begin{enumerate}
\item[(i)] (Labels not finer than the probe.)
\begin{equation*}
\sum_{F:\ j_F\ge n,\ A^{\circ}_F\cap\mathsf{Q}\neq\emptyset}W(F)\le 5^{d}(i-n+1)\,\mathsf{b}(i-n).
\end{equation*}
\item[(ii)] (Labels finer than the probe.)
\begin{equation*}
\sum_{F:\ j_F<n,\ A^{\circ}_F\cap\mathsf{Q}\neq\emptyset}W(F)\le C^{<}\,\mathsf{b}(i-n).
\end{equation*}
\item[(iii)] (The fixed strip.) If $\mathsf{Q}$ meets the strip $Y^{\sharp}$ of a live pocket
$(Y,m)$, then, with $s_j(Y)$ the side parameter of the strip at level $j$ fixed in the geometric
lemma on live strips and probes,
\begin{equation*}
\sum_{F:\ (Y,m)\in\mathcal{P}_F}W(F)
\le\sum_{j=m}^{i}\bigl(s_j(Y)+3\bigr)^{d}\,\mathsf{b}(i-j)
\le(i-m+1)\bigl(101\check{R}_m+2\bigr)^{d}\,\mathsf{b}(i-m).
\end{equation*}
\item[(iv)] (Total.) For every $\beta>0$,
\begin{equation*}
\mathsf{S}(\mathsf{Q}):=\sum_{F\in\mathfrak{L}:\ X_F\cap\mathsf{Q}\neq\emptyset}W(F)
\le e^{\beta\varrho(\mathsf{Q})}\,C_{\mathsf{a}}(\beta)\,(1+N_{*})\,\mathsf{b}(N_{*}).
\end{equation*}
\end{enumerate}
In particular, no age (number of steps since the creation of a pocket), no number of pockets and
no face or volume of a pocket enters a constant, and no factor $L^{d\cdot\mathrm{age}}$ occurs in
\textup{(i)}--\textup{(iv)}.
\end{lemma}

\begin{proof}
Here $\operatorname{dist}_j$ denotes distance measured in units of level $j$.

\emph{Finiteness of the constants.} Since $L>1$, the factor
$e^{-\varsigma_0(\frac12L^{t}-2)_+}$ decays faster than any power of $L^{t}$, so the series defining
$C^{<}$ converges. In $C_{\mathsf{a}}(\beta)$ the bracket and $(1+\varrho)^{p_{\mathsf{b}}}$ grow
at most polynomially in $\varrho$ while $e^{-\beta\varrho}$ decays exponentially, so the supremum
is finite. Both numbers are built from $d,L,\varsigma_0,p_{\mathsf{b}},C_{\mathsf{b}},\beta$ alone.

\emph{Reduction.} The monotonicity of $\mathsf{b}$ assumed in the statement is no restriction: if
$\mathsf{b}$ satisfies the hypotheses of \eqref{4.21:eq-wrapped-terms-a}, then its non-decreasing hull
$\bar{\mathsf{b}}(s):=\sup_{0\le u\le s}\mathsf{b}(u)$ satisfies them with the same constants, so the
lemma may be applied with $\bar{\mathsf{b}}$ in place of $\mathsf{b}$. Indeed $\bar{\mathsf{b}}\ge\mathsf{b}\ge1$,
so the single-cube bound for terms of one level,
\begin{equation*}
\sum_{F\in\mathfrak{L}:\ j_F=j,\ A^{\circ}_F\ni\square'}W(F)\le\mathsf{b}(i-j)
\qquad(j\le i,\ \square'\in\pi_j),
\end{equation*}
and, for $\mathsf{Q}\in\mathfrak{Q}$, $j<\operatorname{lev}(\mathsf{Q})$ and $\square'\in\pi_j$ with
$\square'\subset\mathsf{Q}$, the weighted bound for labels finer than the probe,
\begin{equation*}
\sum_{F\in\mathfrak{L}:\ j_F=j,\ A^{\circ}_F\ni\square'}W(F)\,e^{\varsigma_0 d_j(A^{\circ}_F)}
\le\mathsf{b}(i-j),
\end{equation*}
both in \eqref{4.21:eq-wrapped-terms-a}, hold with $\bar{\mathsf{b}}$ in place of $\mathsf{b}$.
For the growth law $\mathsf{b}(s+t)\le C_{\mathsf{b}}(1+t)^{p_{\mathsf{b}}}\mathsf{b}(s)$
($s,t\ge0$) of \eqref{4.21:eq-wrapped-terms-a}, note first that
$C_{\mathsf{b}}\ge1$ (take $t=0$ and use $\mathsf{b}\ge1$), and that $p_{\mathsf{b}}\ge0$: the growth
law at $s=0$ gives $1\le\mathsf{b}(t)\le C_{\mathsf{b}}(1+t)^{p_{\mathsf{b}}}\mathsf{b}(0)$ for all
$t\ge0$, which is impossible for $p_{\mathsf{b}}<0$. Now let $s,t\ge0$ and $0\le u\le s+t$. If
$u\le s$, then
\begin{equation*}
\mathsf{b}(u)\le\bar{\mathsf{b}}(s)\le C_{\mathsf{b}}(1+t)^{p_{\mathsf{b}}}\bar{\mathsf{b}}(s).
\end{equation*}
If $s<u\le s+t$, the growth law with the pair $(s,u-s)$ gives
\begin{equation*}
\mathsf{b}(u)\le C_{\mathsf{b}}(1+u-s)^{p_{\mathsf{b}}}\mathsf{b}(s)
\le C_{\mathsf{b}}(1+t)^{p_{\mathsf{b}}}\bar{\mathsf{b}}(s).
\end{equation*}
Taking the supremum over $u$ gives
$\bar{\mathsf{b}}(s+t)\le C_{\mathsf{b}}(1+t)^{p_{\mathsf{b}}}\bar{\mathsf{b}}(s)$.
In what follows $\mathsf{b}$ is non-decreasing, as assumed in the statement. All weights $W(F)$ are
non-negative, so every sum below
may be bounded by a sum in which each term is counted at least once.

\emph{Proof of \textup{(i)}.} Fix $j$ with $n\le j\le i$ (recall $j_F\le i$ for every $F$). By the
definition of probes, $\mathsf{Q}$ meets at most $5^{d}$ cubes of $\pi_j$. If $j_F=j$ and
$A^{\circ}_F$ meets $\mathsf{Q}$, one of these cubes belongs to $A^{\circ}_F$. Hence, by the
single-cube bound for terms of one level in \eqref{4.21:eq-wrapped-terms-a},
\begin{equation*}
\sum_{F:\ j_F=j,\ A^{\circ}_F\cap\mathsf{Q}\neq\emptyset}W(F)
\le\sum_{\substack{\square'\in\pi_j\\ \square'\cap\mathsf{Q}\neq\emptyset}}\
\sum_{F:\ j_F=j,\ A^{\circ}_F\ni\square'}W(F)\le 5^{d}\,\mathsf{b}(i-j).
\end{equation*}
Summing over $j=n,\dots,i$ and using $\mathsf{b}(i-j)\le\mathsf{b}(i-n)$ for $j\ge n$ gives (i).

\emph{Proof of \textup{(ii)}.} Let $t\ge1$ and $j=n-t$; one unit of level $n$ is $L^{t}$ units of
level $j$. Let $F$ satisfy $j_F=j$ and $A^{\circ}_F\cap\mathsf{Q}\neq\emptyset$. Then $A^{\circ}_F$
contains a cube $\square_2\in\pi_j$ with $\square_2\subset\mathsf{Q}$. Since $\mathcal{P}_F$ is
non-empty and the strips of its pockets touch $A^{\circ}_F$, $A^{\circ}_F$ also contains a cube
$\square_1\in\pi_j$ touching a strip $Y'^{\sharp}$ with $(Y',m')\in\mathcal{P}_F$ and
$m'\le j<n$. Let $\Omega_n$ be the region of level $n$ fixed together with the probes of level $n$.
The clause of the geometric lemma on live strips and probes that separates live strips from
$\Omega_n$, applied at level $n$ to this pocket,
gives $\operatorname{dist}_n(Y'^{\sharp},\Omega_n)\ge\frac52$. The cell from which the probe is
formed lies in $\Omega_n$, and $\mathsf{Q}$ lies within distance $2$ of it in level-$n$ units.
Therefore $\operatorname{dist}_n(\mathsf{Q},Y'^{\sharp})\ge\frac12$, that is, the distance is at
least $\frac12L^{t}$ in level-$j$ units. Since $\square_2\subset\mathsf{Q}$ and every point of
$\square_1$ lies within one level-$j$ unit of $Y'^{\sharp}$,
\begin{equation*}
\operatorname{dist}_j(\square_1,\square_2)\ge\tfrac12L^{t}-1 .
\end{equation*}
Hence, by the lower bound for the size $d_j$ of a domain of $\mathbf{D}_j$ in terms of the distance
between two of its cubes, $d_j(A^{\circ}_F)\ge\frac12L^{t}-2$. As $d_j(A^{\circ}_F)\ge0$ as well,
$e^{-\varsigma_0 d_j(A^{\circ}_F)}\le e^{-\varsigma_0(\frac12L^{t}-2)_+}$.

The probe $\mathsf{Q}$ contains at most $(5L^{t})^{d}$ cubes of $\pi_j$. Assign each such $F$ to its
cube $\square_2$ and apply the weighted bound for labels finer than the probe in
\eqref{4.21:eq-wrapped-terms-a}, which holds for $j<\operatorname{lev}(\mathsf{Q})$ and
$\square'\in\pi_j$, $\square'\subset\mathsf{Q}$. The level $j$ then contributes at most
\begin{equation*}
\sum_{\substack{\square'\in\pi_j\\ \square'\subset\mathsf{Q}}}
e^{-\varsigma_0(\frac12L^{t}-2)_+}\!\!\sum_{F:\ j_F=j,\ A^{\circ}_F\ni\square'}\!\!
W(F)\,e^{\varsigma_0 d_j(A^{\circ}_F)}
\le(5L^{t})^{d}e^{-\varsigma_0(\frac12L^{t}-2)_+}\,\mathsf{b}(i-n+t),
\end{equation*}
since $i-j=i-n+t$. The growth law gives
$\mathsf{b}(i-n+t)\le C_{\mathsf{b}}(1+t)^{p_{\mathsf{b}}}\mathsf{b}(i-n)$. Summing over $t\ge1$ and
comparing with \eqref{4.21:eq-probe-mass-1} yields (ii).

\emph{Proof of \textup{(iii)}.} Let $F$ satisfy $(Y,m)\in\mathcal{P}_F$, and put $j=j_F$, so that
$m\le j\le i$. Since the strip $Y^{\sharp}$ touches $A^{\circ}_F$, the label $A^{\circ}_F$ contains
a cube of $\pi_j$ touching the box $Y^{\sharp}$. Such cubes lie in the box $Y^{\sharp}$ enlarged by
one level-$j$ unit, so there are at most $(s_j(Y)+3)^{d}$ of them. The single-cube bound for terms
of one level in \eqref{4.21:eq-wrapped-terms-a} bounds the weight of the terms of level $j$ through
each of them by $\mathsf{b}(i-j)$. Summing over the cubes and over $j=m,\dots,i$ gives the first
inequality.

For the second, the clause of the geometric lemma on live strips and probes that bounds the side
parameters of a strip gives, for $j\ge m$,
\begin{equation*}
s_j(Y)\le s_m(Y)\le 101\check{R}_m-1 ,
\end{equation*}
so that $(s_j(Y)+3)^{d}\le(101\check{R}_m+2)^{d}$. Moreover $\mathsf{b}(i-j)\le\mathsf{b}(i-m)$,
and there are $i-m+1$ values of $j$.

\emph{Proof of \textup{(iv)}.} By the definition $X_F=A^{\circ}_F\cup\bigcup_{\mathcal{P}_F}Y^{\sharp}$,
the probe $\mathsf{Q}$ meets $X_F$ if and only if it meets $A^{\circ}_F$ or a strip $Y^{\sharp}$
with $(Y,m)\in\mathcal{P}_F$. By the two clauses of the geometric lemma on live strips and probes
that count the strips met by a probe and bound their levels, $\mathsf{Q}$ meets
at most $2^{d}$ strips of live pockets, and each such pocket $(Y,m)$ has $m\ge n$. Bound the first
alternative by (i) and (ii). Bound the second by (iii), summed over the at most $2^{d}$ strips met
by $\mathsf{Q}$, using $i-m+1\le i-n+1$ and $\mathsf{b}(i-m)\le\mathsf{b}(i-n)$ for $m\ge n$. This
gives
\begin{equation*}
\mathsf{S}(\mathsf{Q})\le\Bigl[(i-n+1)\bigl(5^{d}+2^{d}\max_m(101\check{R}_m+2)^{d}\bigr)
+C^{<}\Bigr]\mathsf{b}(i-n).
\end{equation*}
Here the maximum runs over the levels $m$ of the strips met by $\mathsf{Q}$ and is read as $0$ if
there are none.

Write $\varrho=\varrho(\mathsf{Q})\ge1$. By the first depth condition
$i-\operatorname{lev}(\mathsf{Q})+1\le\varrho(\mathsf{Q})+N_{*}$ of
\eqref{4.21:eq-wrapped-terms-a},
\begin{equation*}
i-n+1\le\varrho+N_{*}\le\varrho(1+N_{*}).
\end{equation*}
The same condition gives $i-n\le\varrho+N_{*}$. Monotonicity and the growth law with
$s=N_{*}$, $t=\varrho$ then give
\begin{equation*}
\mathsf{b}(i-n)\le\mathsf{b}(\varrho+N_{*})
\le C_{\mathsf{b}}(1+\varrho)^{p_{\mathsf{b}}}\,\mathsf{b}(N_{*}).
\end{equation*}
By the second depth condition of \eqref{4.21:eq-wrapped-terms-a}, which gives
$\check{R}_m\le2\varrho(\mathsf{Q})+2$ whenever $\mathsf{Q}$ meets the strip of a live pocket
$(Y,m)$, every strip met by $\mathsf{Q}$ has $\check{R}_m\le2\varrho+2$, so
$101\check{R}_m+2\le202\varrho+204$. Inserting these bounds and using $1\le 1+N_{*}$ for the
term $C^{<}$, we obtain
\begin{align*}
\mathsf{S}(\mathsf{Q})
&\le\Bigl[\varrho(1+N_{*})\bigl(5^{d}+2^{d}(202\varrho+204)^{d}\bigr)+C^{<}\Bigr]
C_{\mathsf{b}}(1+\varrho)^{p_{\mathsf{b}}}\,\mathsf{b}(N_{*})\\
&\le e^{\beta\varrho}(1+N_{*})\,\mathsf{b}(N_{*})\cdot C_{\mathsf{b}}
\Bigl[\varrho\bigl(5^{d}+2^{d}(202\varrho+204)^{d}\bigr)+C^{<}\Bigr]
(1+\varrho)^{p_{\mathsf{b}}}e^{-\beta\varrho}\\
&\le e^{\beta\varrho(\mathsf{Q})}\,C_{\mathsf{a}}(\beta)\,(1+N_{*})\,\mathsf{b}(N_{*}).
\end{align*}
The last step uses $\varrho\ge1$ and the definition \eqref{4.21:eq-probe-mass-1}.

The constants entering (i)--(iv) are $5^{d}$, $2^{d}$, $C^{<}$ and $C_{\mathsf{a}}(\beta)$. Each
depends only on $(d,L,\varsigma_0,p_{\mathsf{b}},C_{\mathsf{b}},\beta)$, which gives the final
sentence of the statement.
\end{proof}

\begin{remark}[Remarks on the probe bound; the standard majorant]\label{4.21:rem-standard-majorant}
We comment on Lemma~\ref{4.21:lem-probe-mass} (the mass read at one probe).
\begin{enumerate}
\item In the final bound of clause (iii) of Lemma~\ref{4.21:lem-probe-mass}, the number of levels
$i-m+1$ and the per-cube majorant $\mathsf{b}$ enter as separate factors.
\item In clause (iii), the terms around a pocket are counted in cubes of their own level $j$, not
in cubes of level $i$. No factor $L^{d(i-j)}$ therefore arises, and no weight has to compensate for
one. For this reason the lemma holds also where the sources lie at lower levels and where the core
is read.
\item \emph{The standard majorant.} Let $b_0,b_1\ge 0$ and $\mathsf{T}\ge 1$, and put
\begin{equation}\label{4.21:eq-standard-majorant-1}
\mathsf{b}(s):=b_0+(s+1)\,b_1\,(\mathsf{T}+c_s)^{r_{\mathrm{pr}}},\qquad s\ge 0 .
\end{equation}
Here $c_s$ is the logarithmic scale variable and $\lambda$, $c_2$, $r_{\mathrm{pr}}$ are the
constants fixed together with the logarithmic coupling variables; of them we use only
$\lambda,c_2,r_{\mathrm{pr}}\ge 0$ and $c_s=\log(1+3\lambda s+2c_2)$. In the application
$\mathsf{T}$ is the logarithmic coupling variable; here only $\mathsf{T}\ge 1$ is used.
Then $\mathsf{b}$ obeys the growth law of the per-cube majorant with the constants
\begin{equation}\label{4.21:eq-standard-majorant-2}
p_{\mathsf{b}}:=1+r_{\mathrm{pr}},\qquad C_{\mathsf{b}}:=(1+3\lambda)^{r_{\mathrm{pr}}},
\end{equation}
that is,
\begin{equation}\label{4.21:eq-standard-majorant-3}
\mathsf{b}(s+t)\le C_{\mathsf{b}}\,(1+t)^{p_{\mathsf{b}}}\,\mathsf{b}(s)\qquad (s,t\ge 0).
\end{equation}
\item Clause (i) of Lemma~\ref{4.21:lem-probe-mass} uses no strip $Y^{\sharp}$. It holds for every
family of terms that satisfies the first part of the weight bound under which that lemma is stated,
the plain terms included. Clause (ii) does use a strip. The weight bound is stated in two parts
because, at the preparation site, a chain term filed below the top level spends its whole decay
rate on its weight.
\end{enumerate}
\end{remark}

\begin{proof}[Proof of \eqref{4.21:eq-standard-majorant-3}]
Let $s,t\ge 0$.

The first step is an inequality between the arguments of the logarithms. Expanding the product gives
\begin{equation*}
(1+3\lambda s+2c_2)(1+3\lambda t)
=1+3\lambda(s+t)+2c_2+9\lambda^2 st+6\lambda c_2 t .
\end{equation*}
Since the last two terms are non-negative,
\begin{equation}\label{4.21:eq-standard-majorant-4}
1+3\lambda(s+t)+2c_2\le (1+3\lambda s+2c_2)(1+3\lambda t).
\end{equation}

Taking logarithms in \eqref{4.21:eq-standard-majorant-4} and using $c_s=\log(1+3\lambda s+2c_2)$
gives $c_{s+t}\le c_s+\log(1+3\lambda t)$. Moreover $c_s\ge 0$, because $1+3\lambda s+2c_2\ge 1$;
together with $\mathsf{T}\ge 1$ this gives $\mathsf{T}+c_s\ge 1$, hence
$\log(1+3\lambda t)\le(\mathsf{T}+c_s)\log(1+3\lambda t)$. Adding $\mathsf{T}+c_s$ to both sides,
we obtain
\begin{equation}\label{4.21:eq-standard-majorant-5}
\mathsf{T}+c_{s+t}\le(\mathsf{T}+c_s)\bigl(1+\log(1+3\lambda t)\bigr).
\end{equation}

Next we bound the logarithmic factor. The inequality $\log(1+x)\le x$ for $x\ge 0$ gives the first
step below, and $(1+3\lambda)(1+t)=1+3\lambda t+t+3\lambda$ gives the second:
\begin{equation}\label{4.21:eq-standard-majorant-6}
1+\log(1+3\lambda t)\le 1+3\lambda t\le(1+3\lambda)(1+t).
\end{equation}

Inserting \eqref{4.21:eq-standard-majorant-6} into \eqref{4.21:eq-standard-majorant-5} and raising
both sides to the power $r_{\mathrm{pr}}$ (both sides are at least $1$, and $r_{\mathrm{pr}}\ge 0$),
we get
\begin{equation*}
(\mathsf{T}+c_{s+t})^{r_{\mathrm{pr}}}
\le C_{\mathsf{b}}\,(1+t)^{r_{\mathrm{pr}}}\,(\mathsf{T}+c_s)^{r_{\mathrm{pr}}} .
\end{equation*}
Moreover $s+t+1\le(s+1)(1+t)$, because the difference of the two sides is $st\ge 0$.

Both facts enter the definition \eqref{4.21:eq-standard-majorant-1}, and $b_1\ge 0$. This gives
\begin{align*}
\mathsf{b}(s+t)
&=b_0+(s+t+1)\,b_1\,(\mathsf{T}+c_{s+t})^{r_{\mathrm{pr}}}\\
&\le b_0+C_{\mathsf{b}}\,(1+t)^{1+r_{\mathrm{pr}}}\,(s+1)\,b_1\,(\mathsf{T}+c_s)^{r_{\mathrm{pr}}}.
\end{align*}
Finally, $b_0\ge 0$ and $C_{\mathsf{b}}(1+t)^{p_{\mathsf{b}}}\ge 1$, so
$b_0\le C_{\mathsf{b}}(1+t)^{p_{\mathsf{b}}}b_0$. Substituting this bound into the last line yields
\eqref{4.21:eq-standard-majorant-3}, with the constants \eqref{4.21:eq-standard-majorant-2}.
\end{proof}

\begin{remark}[Ba{\l}aban's decoupling and the one-component form; {\cite[(1.9)--(1.11)]{B-CMP116}}, {\cite{B-CMP119}}]
\label{4.21:rem-one-component-form}
We state the decoupling of \cite{B-CMP116} in the form given there, together
with the choice of the decoupling family made there for the boundary terms $\mathbf{B}_k$, on which
\cite{B-CMP119} relies. The proofs are those of the cited papers; only the outline of the locality
argument is repeated below. All letters not in the glossary are those of the cited papers.

\emph{The decoupling parameters.} In \cite{B-CMP116} the decoupling is built on a regular partition
$\sigma_k$ of the space $T$ on which the random walks live into cubes $\Delta$ of side $R_1M_1$. For
the particular example treated first, the decoupling family $\sigma_0$ is the family of cubes of
side $M$ that are disjoint from the interior of $\tilde{\square}^4$. To each $\Delta\in\sigma_0$ a
variable $s(\Delta)$ is assigned. For a random walk $\omega$ localised in
$X_0\cup X_1\cup\dots\cup X_n$ one takes the family $\{\Delta_1,\dots,\Delta_m\}$ of all cubes of
$\sigma_0$ that meet this localisation domain, and multiplies the term of the walk expansion
\cite[(1.6)]{B-CMP116} corresponding to $\omega$ by $s(\Delta_1)\cdots s(\Delta_m)$. This defines the
$s$-dependent propagators $H(s)$, $G(s)$, $H_0(s)$, which coincide with the original ones at
$s=1$ (all parameters equal to one). The cube $\tilde{\square}^4$ is thus exempted: no parameter is
attached to a cube meeting its interior. Before the parameters are introduced, \cite{B-CMP116}
observes that in the equations considered there the propagators may be replaced by arbitrary
operators having the same regularity properties and satisfying the same bounds; only these bounds
enter.

\emph{The expansion.} Let $E(s)$ denote the functional to be decoupled, evaluated with the
propagators $H(s)$, $G(s)$, $H_0(s)$. As in cluster expansions, the fundamental theorem of
calculus, applied in each variable $s(\Delta)$, $\Delta\in\sigma_0$, gives \cite[(1.9)]{B-CMP116}
\begin{equation}\label{4.21:eq-one-component-form-1}
E(s)\big|_{s=1}
=\sum_{\sigma\subset\sigma_0}\ \prod_{\Delta\in\sigma}\int_0^1 ds(\Delta)\,
\frac{\partial}{\partial s(\Delta)}\,E(s)\big|_{s(\sigma_0\setminus\sigma)=0}.
\end{equation}

\emph{Component-wise locality.} For $\sigma\subset\sigma_0$ let $Y(\sigma)$ be the region determined
by $\sigma$ in \cite{B-CMP116}; its components are taken with respect to connectedness through
common $(d-1)$-dimensional walls, and $Y_0$ denotes the component of $Y(\sigma)$ singled out in
\cite{B-CMP116}, which contains the exempted cube $\tilde{\square}^4$. In the
notation of \cite{B-CMP116}, consider the function
\begin{equation}\label{4.21:eq-one-component-form-2}
\bigl(t\tilde{\zeta}_{\square}+t_{\square}\zeta_{\square}\bigr)\,
H_k\bigl(H(s),G(s),H_0(s)B'\bigr),
\qquad s(\sigma_0\setminus\sigma)=0 .
\end{equation}
In \eqref{4.21:eq-one-component-form-2} and in the three steps below, $t$, $t_{\square}$,
$\zeta_{\square}$, $\tilde{\zeta}_{\square}$, $A'$, $B'$, $D(\cdot)$ and $A_0(\cdot)$ are the letters of
\cite{B-CMP116}; in particular $A_0(\cdot)$ denotes the solution so written there, not the auxiliary
parameter $A_0$ of the glossary, and $t$ is not the variable $t=X-\hat{x}$ of the glossary.
\cite{B-CMP116} proves that this function depends on $B$ and on $s$ only through
their restrictions to $Y_0$. The argument given there runs as follows. First, when
$s(\sigma_0\setminus\sigma)=0$, the kernels of $H(s)$, $G(s)$, $H_0(s)$ vanish unless both arguments
lie in the interior of one component of $Y(\sigma)$, indeed at distance greater than $M_1$ from
the boundary of that component. Second, the solution $D(H(s),A')$ restricted to the interior of a
component depends on $A'$ and on $H(s)$, hence on $s$, only through their restrictions to that
component. Third, the solution $A_0(H(s),G(s),H_0(s)B')$ vanishes on a neighbourhood of the
complement of $Y(\sigma)$, and on each component of $Y(\sigma)$ it depends only on the propagators
and on $B'$ restricted to that component.

\emph{The one-component form.} If $Y(\sigma)$ has components other than $Y_0$, the term of
\eqref{4.21:eq-one-component-form-1} indexed by $\sigma$ contains derivatives with respect to
parameters $s$ restricted to those components; since, by the component-wise locality just
described, the function \eqref{4.21:eq-one-component-form-2} does not depend on these parameters,
these derivatives render the term equal to zero. This is the simplification of the sum
\eqref{4.21:eq-one-component-form-1} stated as \cite[(1.10)]{B-CMP116}: only those $\sigma$
survive for which $Y(\sigma)$ has no component other than $Y_0$. We note what this step uses. The
one-component form is derived from the fact that the function
\eqref{4.21:eq-one-component-form-2} is localised in the exempted cube $\tilde{\square}^4$:
combined with the component-wise locality of the $s$-dependent operators and solutions, this gives
dependence on $s$ through $s\restriction Y_0$ alone, and the vanishing of the remaining derivatives
follows. The expansion \eqref{4.21:eq-one-component-form-1} and the component-wise locality are used
as stated in \cite{B-CMP116}.

\emph{The bound.} In \cite{B-CMP116} the parameters $s(Y_0)$ are complex valued, with
$|s(\Delta)|\le e^{\kappa_1}$, and the resulting estimate \cite[(1.11)]{B-CMP116} is a bound by
$B_0|X|$, with $B_0$ the constant and $X$ the region of \cite[(1.11)]{B-CMP116}, times a supremum
over walks $\omega$ in which $\omega$ contributes a power of $M_1$
with exponent proportional to $|\omega|$, the factor $\exp(-\delta_0 d(\omega))$ and the factor
$e^{m\kappa_1}$.
Here $d(\omega)$ is the length of a shortest tree graph meeting all localisation domains of
$\omega$, and $m$ is the number of the parameters $s$ connected with the walk $\omega$. If $m$ is
large, for example $m>2^4$, then $\delta_0 d(\omega)\ge\delta_1 mM$ for a positive constant
$\delta_1$ depending on $\delta_0$ only, and the factor $e^{m\kappa_1}$ is controlled by the decay
provided $\delta_1M\ge\kappa_1$. If $m\le 2^4$, the factor is bounded by $e^{16\kappa_1}$ only.

\emph{The family for the boundary terms.} For the expression \cite[(3.7)]{B-CMP109}, or rather its
analytic extension \cite[(3.15)]{B-CMP109}, \cite{B-CMP116} notes that it is determined by the
functions $H_j$ and $(\delta/\delta B)H_j$, besides $H_k$, and that it depends on the first two
only through their restrictions to $X$. The parameters $s$ are introduced there in the same way as
above, with a different family: $\sigma_0$ is now the family of all cubes $\Delta$ disjoint from
the interior of $X_0$, where $X_0$ now denotes the smallest localisation domain in $\mathbf{D}_k$
containing $X$. In \cite{B-CMP119} the expression determined by the boundary terms $\mathbf{B}_k$ is
treated by applying \cite[(3.6), (3.7)]{B-CMP109} to its terms and then localising the result; the
localised expansion runs over domains $Y_0\in\mathbf{D}_k$ with $Y_0\cap\Lambda_{k+1}\ne\emptyset$,
and the estimates there use the localisation of the fluctuation field and the exponential decay of
the minimising function $H_k$. Read together, the two passages fix the decoupling family used in
\cite{B-CMP116} and \cite{B-CMP119} for the boundary terms $\mathbf{B}_k$: the cubes disjoint from
the interior of $X_0$, where $X_0\supseteq X$ is the smallest domain of $\mathbf{D}_k$ containing
$X$. The decoupling family adopted later in this chapter for these terms is a different family,
fixed by its own definition there.
\end{remark}

\begin{lemma}[Decoupled operators are blind to the exempted set]\label{4.21:lem-exempted-set}
Let $\pi$ be the partition into cells of the region on which the decoupling of
Remark~\ref{4.21:rem-one-component-form} is performed. Let $\sigma_0\subset\pi$ be an arbitrary
subfamily, and let
\begin{equation*}
\mathsf{K}:=\bigcup\,(\pi\setminus\sigma_0)
\end{equation*}
be the \emph{exempted set}; it may be empty. For a walk $\omega$ write
$\operatorname{dom}\omega$ for its localisation domain. For every operator
$K=\sum_\omega K_\omega$ with a walk expansion to which the localisation lemma for decoupled
operators applies, and for every complex function $\sigma$ on $\sigma_0$, put
\begin{equation}\label{4.21:eq-exempted-set-1}
K(\sigma):=\sum_{\omega}\Bigl(\prod_{\Delta\in\sigma_0,\ \Delta\cap\operatorname{dom}\omega\neq\emptyset}
\sigma(\Delta)\Bigr)K_\omega ,
\end{equation}
and
\begin{equation*}
Y(\sigma):=\mathsf{K}\cup\bigcup\{\Delta\in\sigma_0:\ \sigma(\Delta)\neq 0\}.
\end{equation*}
Assume the following statements, each as stated and proved for the decoupling family of its own
setting:
\begin{itemize}
\item the localisation lemma for decoupled operators, clauses (a)--(c);
\item the field-datum lemma;
\item the layer theorem, clauses (a)--(f), and its corollary;
\item the statement on the current rider of the layer.
\end{itemize}
Then all of these statements hold, with the same constants, for the arbitrary subfamily
$\sigma_0$ above, when $K(\sigma)$ is defined by~\eqref{4.21:eq-exempted-set-1} and $Y(\sigma)$
is as above.

In particular, let $\mathsf{s}$ be a decoupling parameter (a complex function on $\sigma_0$, in
the role of $\sigma$), let $\mathbf{U},\mathbf{J}$ be the configuration and the current of the
layer theorem, let $\mathcal{Y}(\mathsf{s};b)$ be the layer of that theorem at the datum $b$,
and let $N_y$ be the local norm at the point $y$ in which that theorem is stated. Write
$\mathbf{1}$ for the parameter identically equal to $1$. Then:
\begin{enumerate}
\item[(i)] $\mathcal{Y}(\mathbf{1};b)$ is the layer of the undecoupled operators; at real points
it is the layer of Ba{\l}aban's construction, by clause (d) of the layer theorem.
\item[(ii)] If $\mathsf{s}(\Delta)=0$ for every $\Delta\in\sigma_0\setminus\mathsf{y}$, where
$\mathsf{y}\subset\sigma_0$, then on each component of $\mathsf{K}\cup\bigcup\mathsf{y}$ the layer
$\mathcal{Y}$ depends on $(\mathsf{s},b,\mathbf{U},\mathbf{J})$ only through their restrictions to
that component; $\mathcal{Y}$ vanishes off $\mathsf{K}\cup\bigcup\mathsf{y}$; and $\mathcal{Y}$
vanishes identically on every component on which $b=0$. If $\mathsf{K}=\emptyset$, then
$\mathcal{Y}(0;b)=0$.
\item[(iii)] For all $\mathsf{s}$ with $|\mathsf{s}(\Delta)|\le e^{\kappa_1}$ for every
$\Delta\in\sigma_0$, and at every $y$,
\begin{equation}\label{4.21:eq-exempted-set-2}
N_y\bigl(\mathcal{Y}(\mathsf{s};b)\bigr)\le B_{\sharp}\,\sup|b|\;
e^{-\delta_P d^{\wedge}(y,\operatorname{supp} b)} .
\end{equation}
\item[(iii)$'$] If the datum is the field datum $\mathfrak{j}=\tilde{G}_0(\mathsf{s})\mathsf{J}$,
then
\begin{equation}\label{4.21:eq-exempted-set-3}
N_y(\mathcal{Y})\le B_{\sharp}K_{\sharp}\,\sup \ell^3|\mathsf{J}|\cdot
e^{-\delta_P d^{\wedge}(y,\operatorname{supp}\mathsf{J})},
\end{equation}
where $\sup\ell^3|\mathsf{J}|$ is the size of the source $\mathsf{J}$ in the norm of the
field-datum lemma, $\ell$ denoting the scale factor of that norm; and $\mathfrak{j}(\mathsf{s})$
and $\mathcal{Y}$ vanish on every component of $Y(\mathsf{s})$ that
does not meet $\operatorname{supp}\mathsf{J}$.
\end{enumerate}
Here components are taken with respect to connectedness through common $(d-1)$-dimensional walls,
as in \cite{B-CMP116}.
\end{lemma}

\begin{proof}
The family $\sigma_0$ enters the localisation lemma for decoupled operators in two places only:
through the display~\eqref{4.21:eq-exempted-set-1}, and, in the proof of that lemma, through the
one number $m(\omega)$ of factors in the coefficient $\sigma^{\omega}$ of $K_\omega$, where
\begin{gather*}
m(\omega):=\#\{\Delta\in\sigma_0:\ \Delta\cap\operatorname{dom}\omega\neq\emptyset\},\\
\sigma^{\omega}:=\prod_{\Delta\in\sigma_0,\ \Delta\cap\operatorname{dom}\omega\neq\emptyset}
\sigma(\Delta).
\end{gather*}
Let $m^{\pi}(\omega)$ be the same count with $\sigma_0$ replaced
by $\pi$. Since $\sigma_0\subset\pi$,
\begin{equation*}
m(\omega)\le m^{\pi}(\omega).
\end{equation*}
We write $d(\omega)$ for the length of a shortest tree graph meeting all localisation domains of
$\omega$, as in \cite[(1.11)]{B-CMP116}.

\emph{Clause (a).} In the proof of the localisation lemma for decoupled operators the number
$m=m(\omega)$ is used only in the bound on $\sigma^\omega$, in two cases. If $m\le 2^4$, then
$\sigma^\omega$ is a product of at most $2^4$ factors of modulus at most $e^{\kappa_1}$, so
$|\sigma^{\omega}|\le e^{16\kappa_1}$. If $m>2^4$, then also $m^{\pi}>2^4$, and the inequality
stated after \cite[(1.11)]{B-CMP116} for this case, namely $\delta_0 d(\omega)\ge\delta_1 mM$ with
a constant $\delta_1>0$ depending on $\delta_0$ only, is a statement about the partition $\pi$ and
holds with $m^\pi$ in place of $m$. Indeed, $\operatorname{dom}\omega$ is a connected union of sets
of diameter $O(M_1)$, and $M_1=M/w$ is small compared with $M$, $w=M/M_1$ being large; a tree
graph through all localisation domains of $\omega$
therefore passes within distance $O(M_1)$ of each of the $m^{\pi}(\omega)$ cells of $\pi$ that
$\operatorname{dom}\omega$ meets; and more than $2^d=2^4$ cells of the partition $\pi$ are
not met inside one cube of side $M$ measured at the cells' own scale. Hence
\begin{equation*}
\delta_0 d(\omega)\ge\delta_1 m^{\pi}(\omega)M\ge\delta_1 m(\omega)M .
\end{equation*}
With the condition $\delta_1 M\ge16\kappa_1$ of the localisation lemma this gives, as in its
proof,
\begin{equation*}
|\sigma^\omega|\le e^{m\kappa_1}\le e^{\delta_1 mM/16}\le e^{\delta_0 d(\omega)/16}.
\end{equation*}
Both cases therefore give the bounds used in the proof of clause (a), with the same constants,
and clause (a) holds verbatim for the family $\sigma_0$.

\emph{Clause (b).} Let $\omega$ be a walk whose coefficient $\sigma^\omega$ in
\eqref{4.21:eq-exempted-set-1} is non-zero. Then $\sigma(\Delta)\neq0$ for every
$\Delta\in\sigma_0$ that meets $\operatorname{dom}\omega$, while every cell of $\pi$ not in
$\sigma_0$ lies in $\mathsf{K}$. Since $\pi$ is a partition, $\operatorname{dom}\omega$ is contained
in the union of the cells it meets, and hence
\begin{equation*}
\operatorname{dom}\omega\subset\mathsf{K}\cup\bigcup\{\Delta\in\sigma_0:\ \sigma(\Delta)\neq0\}
=Y(\sigma).
\end{equation*}
The domain $\operatorname{dom}\omega$ is connected in the sense of the first sentence of the
argument of \cite{B-CMP116} on connectedness by common $(d-1)$-dimensional walls: consecutive
localisation cubes of a walk overlap in full dimension, so a walk domain does not cross a contact
of cells along a corner. Therefore $\operatorname{dom}\omega$ lies in one component of $Y(\sigma)$.
Since, by the definition of its localisation domain, the kernel of $K_\omega$ has both arguments
in $\operatorname{dom}\omega$ and $K_\omega$ depends on the configuration in
$\operatorname{dom}\omega$ only, the kernel of $K(\sigma)$ vanishes unless both arguments lie in
one component of $Y(\sigma)$, and it reads the configuration there only; this is the content of
clause (b).

\emph{Clause (c).} Neither the statement nor the proof of clause (c) involves $\sigma_0$.

\emph{The remaining statements.} The field-datum lemma, the layer theorem, its corollary and the
statement on the current rider of the layer use the decoupled operators $K(\sigma)$ only through
clauses (a)--(c) of the localisation lemma for decoupled operators: factor by factor through
clause (a), and, for the locality clause (f) of the layer theorem, through clause (b) of the
field-datum lemma together with the component-wise argument of \cite[p.~4]{B-CMP116}. By the three
preceding paragraphs, these statements hold for the family $\sigma_0$ with the same constants.

\emph{Item (i).} By clause (a) of the localisation lemma for decoupled operators,
$K(\mathbf{1})=K$ for each of the operators, so at $\mathsf{s}=\mathbf{1}$ the layer is built from
the undecoupled operators; the identification at real points is clause (d) of the layer theorem.

\emph{Item (ii).} Under its hypothesis $Y(\mathsf{s})\subset\mathsf{K}\cup\bigcup\mathsf{y}$, and
we apply clauses (f) and (a) of the layer theorem on each component: the equations defining the
layer separate over the components, and the solution is unique in the ball of that theorem. On a
component on which $b=0$ the solution is therefore the one at datum zero, which vanishes by
clause (a). If $\mathsf{K}=\emptyset$, that is $\sigma_0=\pi$, then every walk has a non-empty
domain, which meets at least one cell of $\pi=\sigma_0$, so every term of
\eqref{4.21:eq-exempted-set-1} carries a factor $\mathsf{s}(\Delta)$; the case
$\mathsf{y}=\emptyset$ of the preceding statements then gives $\mathcal{Y}(0;b)=0$.

\emph{Item (iii).} Put $p:=e^{-\delta_P d^{\wedge}(\cdot,\operatorname{supp} b)}$. This weight
belongs to the class $\mathfrak{M}_\delta$ of admissible weights of clause (b) of the layer
theorem, by the admissibility of exponential weights in the scaled distance $d^{\wedge}$. Clause
(b) of the layer theorem, applied at this weight $p$ with
the second datum equal to zero, at which the layer vanishes by clause (a), gives a bound on the
layer in the norm weighted by $p$ by $B_{\sharp}$ times the weighted norm of $b$. Since
$d^{\wedge}(y,\operatorname{supp}b)=0$ for $y\in\operatorname{supp}b$, the weight $p$ equals $1$ on
$\operatorname{supp}b$, so the weighted norm of $b$ is $\sup|b|$, and the weighted bound on the
layer, read at the point $y$, is \eqref{4.21:eq-exempted-set-2}.

\emph{Item (iii)$'$.} Put $p:=e^{-\delta_P d^{\wedge}(\cdot,\operatorname{supp}\mathsf{J})}$; as in
item (iii), $p\in\mathfrak{M}_\delta$, and $\mathfrak{M}_\delta\subset\mathfrak{M}_{2\delta}$.
Clause (a) of the field-datum lemma gives
\begin{equation*}
\|\mathfrak{j}\|^{\Delta}_{p}\le K_{\sharp}\|\mathsf{J}\|_{p}=K_{\sharp}\sup\ell^3|\mathsf{J}| ,
\end{equation*}
and clause (b) of the layer theorem at the weight $p$ then gives
\eqref{4.21:eq-exempted-set-3}. The statements on supports follow from clause (b) of the
localisation lemma for decoupled operators.
\end{proof}

We note three points about the use of Lemma~\ref{4.21:lem-exempted-set}. First, the expansion
by the fundamental theorem of calculus, \cite[(1.9)]{B-CMP116}, requires no multi-affinity in the
parameters: the functionals to which it is applied here are products of decoupled factors, and the
expansion holds for every continuously differentiable function on $[0,1]^{\sigma_0}$. Second, it
is the decaying bound \eqref{4.21:eq-exempted-set-2}, and not only the bound by $\sup|b|$, that is
independent of the choice of $\sigma_0$. Third, thin strata of cells do not affect the polydisc
$|\sigma(\Delta)|\le e^{\kappa_1}$: walks step in cubes of side $M_1$, and such a stratum is at
least $w=M/M_1$ of its own units thick, so the number $m(\omega)$ is paid for as in the proof of
clause (a).

\begin{definition}[The decoupling clause]\label{4.21:def-decoupling-clause}
Let $F$ be a stored boundary-type term in the relative format of stored terms. Write $A_F$ for its
label and $\mathcal{P}_F$ for its list of pairs $(Y,m)$; each such pair carries its strip
$Y^{\sharp}$. The term $F$ is called \emph{wrapped} if $\mathcal{P}_F\neq\emptyset$ and
\emph{plain} if $\mathcal{P}_F=\emptyset$. For wrapped $F$ we put
\begin{equation}\label{4.21:eq-decoupling-clause-1}
X_F:=A_F\cup\bigcup_{(Y,m)\in\mathcal{P}_F}Y^{\sharp}.
\end{equation}
The stored terms are expanded at three sites of the step, denoted (T), (L) and (P): (T) is the
site of the operation $\mathbf{T}$, (L) is the site of the localisations inside the operation
$\mathbf{R}_k$, and (P) is the site at which the family of fluctuation cubes is expanded. At
the site (L) the stored terms are read as the \emph{units} $v$ of Ba{\l}aban's first and second
parts, and for a unit $v$ we write $X_v$ for its localisation domain. The first-part and
second-part units, the bracket of a unit, live and dead strips, and a chain of localisations
read by a unit are understood in the sense of the relative format. To expand the bracket of a
term by \cite[(1.9)]{B-CMP116} with a family $\sigma_0$ of cells means the following. A
parameter $\mathsf{s}(\Delta)$ is attached to every cell $\Delta\in\sigma_0$, and the
fundamental theorem of calculus is applied in these parameters as in \cite[(1.9)]{B-CMP116}.
A cell of $\sigma_0$ in whose parameter a derivative is taken is said to be \emph{cut}. In the
notation of Lemma~\ref{4.21:lem-exempted-set}, the union $\mathsf{K}$ of the cells of the
site's partition that do not belong to $\sigma_0$ is the \emph{exempted set}.

The choices made below do not depend on the complex source parameters. The same choices are
used in both of the two expansions that are compared with one another on a common family of
domains. Each prescription below declares a decoupling family under the declaration rule of
this chapter; by Lemma~\ref{4.21:lem-exempted-set} the bounds on the decoupled operators hold
with the same constants whichever cells are exempted.

The \emph{decoupling clause} consists of the following prescriptions (a)--(e) for a wrapped
term $F$. The decoupling families they fix are
\begin{equation}\label{4.21:eq-decoupling-clause-a}
\sigma_0=
\begin{cases}
\pi^{T}, & \text{at the site (T)},\\
\{\Delta:\ \Delta\cap\Lambda^{\sim}=\emptyset\}, & \text{first operation at the site (L)},\\
\{\Delta:\ \Delta\cap\Lambda=\emptyset\}, & \text{three later localisations at the site (L)}.
\end{cases}
\end{equation}
Here $\Delta$ ranges over the cells of the site's partition, $\Lambda$ is the small-field
region of the step, and $\Lambda^{\sim}$ is its enlargement used in Ba{\l}aban's first
operation.

\begin{enumerate}
\item[(a)] \emph{The site (T).} At the site (T) the bracket of $F$ is expanded by
\cite[(1.9)]{B-CMP116} with $\sigma_0:=\pi^{T}$. Here $\pi^{T}$ is the family of all cells of
the cover $\{\square\}$ of \cite{B-CMP119}, so that $\mathsf{K}=\emptyset$. This
expansion is used in place of \cite[(3.6)--(3.7)]{B-CMP109} and in place of the core
$\bar{S}_v$ of the unit $v$.

\emph{Comparison with Ba{\l}aban's construction.} For the boundary terms at this site,
Ba{\l}aban's family $\sigma_0$ is the family of all cubes disjoint with the interior of $X_0$
\cite{B-CMP116}, \cite{B-CMP119}. Here $X$ is the localisation region of the term in the
notation of \cite{B-CMP116}, not a large-field region, and $X_0\supseteq X$ is the smallest
localisation domain from $\mathbf{D}_k$ containing $X$; see
Remark~\ref{4.21:rem-one-component-form}. Prescription (a) uses the family of all cells instead.

\item[(b)] \emph{The site (L): the first operation.} At the site (L) of $\mathbf{R}_k$ the
following is done for every wrapped unit, whether its strips are live or dead. In Ba{\l}aban's
first operation, the bracket is expanded by \cite[(1.9)]{B-CMP116} in $\mathsf{s}$ alone, with
$\sigma_0$ the family of cells disjoint with $\Lambda^{\sim}$. This applies to each of the at
most four brackets of a second-part unit and to the one bracket of a first-part unit. In this
expansion the block $\square_0$ of the cover attached to the unit (or the domain $X$ of the
term), which Ba{\l}aban's first operation exempts, is not exempted. Ba{\l}aban's amplitude
$t_{\square}$ is not taken either; this is
the bracket telescoped with $\zeta_{\square}$ and estimated by Cauchy's formula in
$t_{\square}$, of \cite[(3.4)--(3.7)]{B-CMP109} and \cite[(1.30)]{B-CMP116}.

\emph{Comparison with Ba{\l}aban's construction.} Prescription (b) differs from Ba{\l}aban's
first operation in two respects: the exempted set, and the omitted amplitude $t_{\square}$.
For a unit with $X_v\subset\Lambda$ the family $\sigma_0$ of (b) is Ba{\l}aban's own, and
only the second respect differs. It is this second respect that produces the sub-case, in the
analysis of the pieces, in which the family of cut cells is $\mathsf{y}_1$ alone.

\item[(c)] \emph{The site (L): the three later localisations.} The three later localisations
run over Ba{\l}aban's own families, without change: $\sigma_0$ is the family of cells disjoint
with $\Lambda$. No cell meeting $\Lambda$ is ever cut in a chain of localisations that a unit
reads.

\item[(d)] \emph{The site (L): domains.} This prescription declares domains and nothing more.
Let $\operatorname{lab}_t$ be the label after the $t$-th localisation, and let $Z$ be the
large-field region. Put
\begin{equation}\label{4.21:eq-decoupling-clause-2}
\begin{split}
Y^{\mathrm{tr}}_t:={}&\operatorname{lab}_t\cup\bigcup\{C:\ C\text{ a component of }Z,\
C\cap\operatorname{lab}_t\neq\emptyset\}\\
&\cup\bigcup\{Y^{\sharp}:\ (Y,m)\in\mathcal{P}_F\text{ live after }\mathbf{R}_k\}.
\end{split}
\end{equation}
This set is a union of cells in proper scales. Its boundary layer is taken of width at most
$2M_1$ in the corresponding scale, as in Ba{\l}aban's later localisations. The set
$Y^{\mathrm{tr}}_t$ enters only two things:
\begin{itemize}
\item[($\alpha$)] the domains $Y_4(\mathfrak{p})$ and $Y^{+}$ of a piece $\mathfrak{p}$ of the
expansion;
\item[($\beta$)] the reading sets $\mathcal{R}_{t-1}$.
\end{itemize}
It enters nothing else. No determining set
$\mathbf{B}''_k(Y^{\mathrm{tr}}_t)$, $\mathbf{B}_1(\cdot)$ or $\mathbf{B}_k(\cdot)$ is formed.
No representation \cite[(1.61)--(1.63)]{B-CMP122b} is substituted. The layers used in the
second, third and fourth localisations are the global layers of the carriers, as fixed for the
carriers in this chapter.

\emph{Comparison with Ba{\l}aban's construction.} Ba{\l}aban's domain $Y_1\in\mathbf{D}_k$,
with its localised representations, is not formed at the site (L) in the expansion used in
the proof of the correlation theorem. That difference therefore belongs to that expansion,
not to the present clause. For a unit without live strip, $Y^{\mathrm{tr}}_t$ belongs to
$\mathbf{D}_k$ and coincides, as a set, with Ba{\l}aban's own $Y_t$. The closure
$[Y^{\sharp}]_k$ is not taken: it would violate condition (f1) of the relative format and
would merge strips.

\item[(e)] \emph{The site (P).} At the site (P) the expansion is not changed; it is the
expansion already fixed for that site (the perturbative-site expansion), used unchanged. Only
the domain function $\mathsf{d}$ of the relative format is thinned. The
resummation lemma at the site (P) is a resummation inside a proof and declares nothing.
\end{enumerate}

Whether $F$ is wrapped or not, the clause changes the localisation of boundary terms and never
the exponent; this is part~(i) of the lemma on the source form of the boundary-term expansion.

For the expansion used in the proof of the correlation theorem the following convention is
also fixed here. At the site (L), every interior first-part unit (that is, every first-part
unit with $X_v\subset Z$), plain or wrapped, is expanded by (b), (c) and (d). No
$\zeta_{\square}$ and no interpolation time $\mathsf{t}$ is used for such a unit. The
interpolation time $\mathsf{t}$ of the boundary-term lemma used in that proof is used on
plain first-part units with $K_v\setminus Z\neq\emptyset$, where $K_v$ is the domain that
lemma attaches to the unit $v$, and there only in the fourth localisation. Consequently no
piece carries two auxiliary parameter discs.

The remaining terms with $\mathcal{P}_F=\emptyset$ are expanded as in the proof of the
correlation theorem. Nothing in the present construction is proved about plain terms except
three statements: the resummation lemma at the site (P), the lemma on the constants
$K_2^{\flat}$ and $c_{\flat}$ at that site, and part~(c) of the statement on the thinned domain
function; none of the three distinguishes wrapped terms from plain ones.
\end{definition}

\begin{definition}[Links, readers at amplitude one, and holomorphy]\label{4.21:def-links-readers}
Fix a site, (T) or (L), of the decoupling clause (Definition~\ref{4.21:def-decoupling-clause}).
Fix also its level $i$, its partition $\pi$ into cells and the subfamily $\sigma_0\subset\pi$
chosen there, and put $\mathsf{K}:=\bigcup(\pi\setminus\sigma_0)$, the exempted set of
Lemma~\ref{4.21:lem-exempted-set}. For a parameter $\sigma$ on $\sigma_0$ we write
\begin{equation*}
Y(\sigma):=\mathsf{K}\cup\bigcup\{\Delta:\sigma(\Delta)\neq0\},
\end{equation*}
as in Lemma~\ref{4.21:lem-exempted-set}. There $\mathcal{Y}(\mathsf{s};b)$ denotes the layer at
source datum $b$ and decoupling parameter $\mathsf{s}$, supplied by the theorem on the layer and
carried over by that lemma.

\smallskip\noindent\emph{(a) One link.} A \emph{link} of the site consists of the following data.
\begin{enumerate}
\item A site geometry $(\pi^{\mathfrak{s}},\mathfrak{H},\mathfrak{F})$ in the sense of
Definition~\ref{4.21:def-site-geometry}, with $\pi^{\mathfrak{s}}:=\sigma_0$ and with $\mathfrak{H}$
the family of components of $\mathsf{K}$.
\begin{itemize}
\item At the site (T) one has $\mathsf{K}=\emptyset$, and $\mathfrak{F}$ is the family of level-$i$
cells holding free bonds of the domain $\Omega_{i+1}$ of Ba{\l}aban's step.
\item At the site (L), where the fixed level $i$ is the level $k$ of the step, one has
$\mathsf{K}=\Lambda^{\sim}$, respectively $\mathsf{K}=\Lambda$, as
fixed in the decoupling clause \eqref{4.21:eq-decoupling-clause-a}. By clause (f) of the geometric
lemma, the components of $\mathsf{K}$ are boxes with sides at most $h_{\sharp}$ and at mutual
distance at least $19\check{R}_k$. The members of $\pi^{\mathfrak{s}}$ sharing a wall with them
have level $k$.
\end{itemize}
\item A source datum $b$ vanishing off $\mathsf{K}\cup\bigcup\mathfrak{F}$.
\item The \emph{link map} on the lower end $(\mathbf{K},\mathbf{J}_{\mathbf{K}})$ of the link, a
configuration--current pair:
\begin{equation*}
\mathsf{R}(\mathsf{s}):=\bigl(e^{i\eta\mathcal{Y}(\mathsf{s};b)}\mathbf{K},\;
\mathbf{J}_{\mathbf{K}}+\mathcal{J}[\mathcal{Y}(\mathsf{s};b)]\bigr).
\end{equation*}
The form of this map, with its unit $\eta$, is that of \cite[(3.10)--(3.13)]{B-CMP109}. The map
$\mathcal{J}$ is the current
correction $\mathfrak{s}^{\sharp}_{\mathsf{s}}[\cdot]$ of the current-rider statement. It is a local
function of $(A',X,\mathcal{Y})$, up to the action of the decoupled operators $P_1(\mathsf{s})$ and
$\mathfrak{m}_H(\mathsf{s})$. The supports of these operators obey the locality clause of the
localisation lemma for the decoupled operators (the clause stating that, if $\sigma=0$ on
$\sigma_0\setminus Y$, the kernel vanishes unless both arguments lie in one component of
$Y(\sigma)$, and then reads the configuration there only). Here $A'$ and $X$ are the fields of
the theorem on the layer of which $\mathcal{Y}$ is the layer.
\end{enumerate}
The link map is carried together with the gauge map of \cite{B-CMP122b},
\begin{equation*}
\mathsf{u}(\mathsf{s}):=\mathfrak{u}\bigl[e^{i\eta\mathcal{Y}(\mathsf{s})}\mathbf{K};\mathbf{K}\bigr].
\end{equation*}
By clause (i) of the gauge-map lemma, $\mathsf{u}(\mathsf{s})$ is block-local in $\mathcal{Y}$: on
each block it is an analytic function of $\mathbf{H}''_k$ restricted to that block. It acts on the
pair and, jointly, on the slots of a stored term $F$. The composition $F\circ\mathsf{R}$ means the
interpolated term of $F$ at interpolation time $\mathsf{t}=0$, namely
\begin{equation*}
(F\circ\mathsf{R})(\mathsf{s}):=F^{\sharp}\bigl(\mathsf{R}(\mathsf{s})^{h};\mathbf{A},
\boldsymbol{\Phi},h\,\mathsf{u}(\mathsf{s})\,g'\bigr).
\end{equation*}
Here $h$ is a gauge transformation, $g'$ is the gauge-transformation argument (the gauge slot) of
the orbit datum $F^{\sharp}$, $\mathsf{R}(\mathsf{s})^{h}$ denotes the pair
$\mathsf{R}(\mathsf{s})$ transformed by $h$, and the last argument is the product
$h\,\mathsf{u}(\mathsf{s})\,g'$. At the site (T), $h=1$. At the site (L), $h$ is the gauge
transformation $h^{\mathfrak{b}}_q$ of the frame term $F^{\mathfrak{b}}(m;g')$ of the third-factor
bound, $\mathsf{u}$ is its carried gauge map $\mathfrak{Q}^{\mathfrak{b}_{s-1}}\mathfrak{w}_s$, and
$F\circ\mathsf{R}$ is that frame term. The gauge map inherits the locality of
the layer:
\begin{equation}\label{4.21:eq-links-readers-1}
\mathsf{u}(\mathsf{s})=1\quad\text{wherever }\mathcal{Y}(\mathsf{s})=0 .
\end{equation}

\smallskip\noindent\emph{(b) Readers at amplitude one.} A stored term $F$ is a \emph{reader at
amplitude one} if it reads the $\mathsf{s}$-dependent objects only through the triple
\begin{equation}\label{4.21:eq-links-readers-a}
\mathcal{A}(\mathsf{s}):=\bigl(\mathcal{Y}(\mathsf{s};b),\;\mathcal{J}[\mathcal{Y}(\mathsf{s};b)],
\;\mathsf{u}(\mathsf{s})\bigr)
\end{equation}
restricted to a closed set $\mathcal{R}_F$, its reading set, and at amplitude one. That is, no
response kernel $(\delta/\delta b)\mathcal{Y}$ and no derivative in an amplitude is read. The
triple has the following property. On a component $C$ of $Y(\mathsf{s})$,
$\mathcal{A}(\mathsf{s})\restriction C$ is a function of
$(b\restriction C,\mathsf{s}\restriction C,(\mathbf{U},\mathbf{J})\restriction C)$, and
\begin{equation}\label{4.21:eq-links-readers-2}
\mathcal{A}(\mathsf{s})\restriction C=(0,0,1)\qquad\text{if }b\restriction C=0 .
\end{equation}

\smallskip\noindent\emph{(c) Holomorphy.} A stored term $F$ has the \emph{holomorphy property} if
the map
\begin{equation}\label{4.21:eq-links-readers-b}
(\mathsf{s},b,(\mathbf{U},\mathbf{J}),\mathbf{A},w,g')\longmapsto
F\bigl(\mathsf{R}(\mathsf{s})\restriction X_F\bigr)
\end{equation}
is holomorphic, and bounded in modulus by the size functional $\mathfrak{N}(F)$ of (d) below, on
the product of the following sets:
\begin{itemize}
\item the polydisc $\{|\mathsf{s}(\Delta)|\le e^{\kappa_1}:\Delta\in\pi^{\mathfrak{s}}\}$;
\item the source polydisc of the site;
\item the output tube of the site;
\item the domains of the slots of $F$, of the twist parameter $w$ of the step
(Definition~\ref{4.4:def-twisted-argument}) and of the chart, the chart radius being lowered by
the amount fixed for the site.
\end{itemize}

\smallskip\noindent\emph{(d) Norm convention.} $\mathfrak{N}$ denotes any of the size functionals
of a stored boundary term used in the correlation theorem. Each is the supremum over the whole
parameter domain of the term; among them is $\mathfrak{N}^{\mathbb{C}}_{\varepsilon}$. Together
with $\mathfrak{N}$ we use
\begin{equation}\label{4.21:eq-links-readers-c}
\mathfrak{N}^{A}(F):=\sum_{\mathcal{S}}e^{\kappa_A|\mathcal{S}|}\,
\mathfrak{N}\bigl(F(\cdot,\mathcal{S})\bigr).
\end{equation}
Here the sum runs over the strong sets $\mathcal{S}$, $F(\cdot,\mathcal{S})$ is the part of $F$
with strong set $\mathcal{S}$, and $\kappa_A$ is the weight of the strong sets.
$\mathfrak{N}^{A,\mathbb{C}}$ is the functional \eqref{4.21:eq-links-readers-c} built on
$\mathfrak{N}^{\mathbb{C}}$. The expansion of the stored terms along the links has two stages: the
operations of Stage~A are linear in $F$ at fixed strong set and fixed parameters, and the
operations of Stage~B are majorant series in the suprema $\mathfrak{N}$. Every numerical estimate
of the operations of Stage~A and of Stage~B holds in $\mathfrak{N}$, $\mathfrak{N}^A$,
$\mathfrak{N}^{\mathbb{C}}$ and $\mathfrak{N}^{A,\mathbb{C}}$ alike, given the holomorphy property
\eqref{4.21:eq-links-readers-b} in the same norm.

For a wrapped $F$ (that is, $\mathcal{P}_F\neq\emptyset$, as in the decoupling clause), the
holomorphy property \eqref{4.21:eq-links-readers-b} is available only in $\mathfrak{N}^{\mathbb{C}}$.
Accordingly, wherever that expansion writes $\mathfrak{N}^A$, the correlation theorem it serves is
read in $\mathfrak{N}^{A,\mathbb{C}}$. In that norm, the holomorphy property for wrapped terms is
supplied by the statements on wrapped terms that follow the expansion.
\end{definition}

\begin{proof}
We prove the locality of the gauge map \eqref{4.21:eq-links-readers-1}, the property of the
triple \eqref{4.21:eq-links-readers-a} stated in (b), and the submultiplicativity of
$\mathfrak{N}^A$ on which the convention (d) rests.

\emph{The locality of the gauge map.} Let a block be given on which $\mathcal{Y}(\mathsf{s})$
vanishes. On that block $e^{i\eta\mathcal{Y}(\mathsf{s})}\mathbf{K}=\mathbf{K}$. Since
$\mathsf{u}(\mathsf{s})$ is block-local in $\mathcal{Y}$, its value on that block is the value of
$\mathfrak{u}[\mathbf{K};\mathbf{K}]$ there, which is $1$ by the normalisation of the gauge map in
the gauge-map lemma. This is
\eqref{4.21:eq-links-readers-1}.

\emph{The triple on a component.} Let $C$ be a component of $Y(\mathsf{s})$. Put
$\mathsf{y}:=\{\Delta\in\sigma_0:\mathsf{s}(\Delta)\neq0\}$, so that $\mathsf{s}=0$ off
$\mathsf{y}\subset\sigma_0$ and $Y(\mathsf{s})=\mathsf{K}\cup\bigcup\mathsf{y}$.

First entry. Clause (ii) of Lemma~\ref{4.21:lem-exempted-set} applies. It gives that on $C$ the
layer $\mathcal{Y}(\mathsf{s};b)$ reads $(\mathsf{s},b,\mathbf{U},\mathbf{J})$ on $C$ only. It
also gives that $\mathcal{Y}$ vanishes identically on $C$ if $b\restriction C=0$.

Second entry. The current correction $\mathcal{J}[\mathcal{Y}]$ is a local function of
$(A',X,\mathcal{Y})$, up to the terms $P_1(\mathsf{s})A'$ and $\mathfrak{m}_H(\mathsf{s})X$.
Lemma~\ref{4.21:lem-exempted-set} carries the localisation lemma for the decoupled operators over
to the present $\sigma_0$ with the same constants and with $Y(\sigma)$ as above. Its locality
clause, applied with $\sigma=\mathsf{s}$, shows that the kernels of $P_1(\mathsf{s})$ and
$\mathfrak{m}_H(\mathsf{s})$ keep both legs in $C$. Hence
$(P_1(\mathsf{s})A')\restriction C$ is determined by $A'\restriction C$, and
$(\mathfrak{m}_H(\mathsf{s})X)\restriction C$ is determined by $X\restriction C$. Therefore
$\mathcal{J}[\mathcal{Y}]\restriction C$ is determined by
$(A',X,\mathcal{Y})\restriction C$. By clause (f) of the theorem on the layer, which
Lemma~\ref{4.21:lem-exempted-set} carries over with $Y(\sigma)$ as above, the fields
$A'\restriction C$ and $X\restriction C$ read $(\mathsf{s},b,\mathbf{U},\mathbf{J})$ on $C$ only.
Together with the first entry, $\mathcal{J}[\mathcal{Y}]\restriction C$ is therefore a function of
$(b\restriction C,\mathsf{s}\restriction C,(\mathbf{U},\mathbf{J})\restriction C)$.

Now suppose $b\restriction C=0$. Then $\mathcal{Y}=0$ on $C$, by the first entry. Also
$A'=X=0$ on $C$: since $A'\restriction C$ and $X\restriction C$ read $b$ on $C$ only, their values
at a datum with $b\restriction C=0$ equal their values at the zero datum, and these vanish by
clause (a) of the theorem on the layer (the fields vanish at zero datum), which
Lemma~\ref{4.21:lem-exempted-set} carries over. Since the two decoupled operators keep both
legs in $C$, the terms $P_1(\mathsf{s})A'$ and $\mathfrak{m}_H(\mathsf{s})X$ vanish on $C$. Hence
$\mathcal{J}[\mathcal{Y}]=\mathcal{J}[0]=0$ on $C$.

Third entry. By block-locality, $\mathsf{u}(\mathsf{s})$ on $C$ is determined by $\mathcal{Y}$ and
$\mathbf{K}$ there. Here $\mathbf{K}$ is the configuration entry of the pair
$(\mathbf{U},\mathbf{J})$ at the lower end of the link, so by the first entry
$\mathsf{u}(\mathsf{s})\restriction C$ is a function of
$(b\restriction C,\mathsf{s}\restriction C,(\mathbf{U},\mathbf{J})\restriction C)$. If
$b\restriction C=0$, then $\mathcal{Y}=0$ on $C$, and
\eqref{4.21:eq-links-readers-1} gives $\mathsf{u}(\mathsf{s})=1$ on $C$.

Together these give the functional dependence asserted in (b) and the identity
\eqref{4.21:eq-links-readers-2}. The identity \eqref{4.21:eq-links-readers-2} concerns the triple
only. Its analogue fails for a reader of a bare response kernel $(\delta/\delta b)\mathcal{Y}$;
under the decoupling clause (Definition~\ref{4.21:def-decoupling-clause}) no stored term reads such
a kernel.

\emph{The norm convention.} By their definitions, each operation of Stage~A is linear in $F$ at
fixed $\mathcal{S}$ and fixed parameters, and each operation of Stage~B is a majorant series in the
suprema $\mathfrak{N}$, which on strong sets acts by $\mathcal{S}\mapsto\bigcup_j\mathcal{S}_j$.

Let $F_1,\dots,F_n$ be the terms entering one product of such a series. The part of
$F_1\cdots F_n$ with strong set $\mathcal{S}$ is the sum of the products
$\prod_jF_j(\cdot,\mathcal{S}_j)$ over the tuples $(\mathcal{S}_1,\dots,\mathcal{S}_n)$ with
$\bigcup_j\mathcal{S}_j=\mathcal{S}$. Since $\mathfrak{N}$ is a supremum, it is subadditive, and
the supremum of a product is at most the product of the suprema. Since moreover
$|\bigcup_j\mathcal{S}_j|\le\sum_j|\mathcal{S}_j|$ and $\kappa_A\ge0$,
\begin{align*}
\mathfrak{N}^A(F_1\cdots F_n)
&\le\sum_{\mathcal{S}_1,\dots,\mathcal{S}_n}e^{\kappa_A|\bigcup_j\mathcal{S}_j|}
\prod_{j=1}^{n}\mathfrak{N}\bigl(F_j(\cdot,\mathcal{S}_j)\bigr)\\
&\le\sum_{\mathcal{S}_1,\dots,\mathcal{S}_n}\prod_{j=1}^{n}e^{\kappa_A|\mathcal{S}_j|}\,
\mathfrak{N}\bigl(F_j(\cdot,\mathcal{S}_j)\bigr)
=\prod_{j=1}^{n}\mathfrak{N}^A(F_j),
\end{align*}
so $\mathfrak{N}^A$ is submultiplicative. Consequently the two kinds of operation are controlled
in exactly the same way by each of the four suprema $\mathfrak{N}$, $\mathfrak{N}^A$,
$\mathfrak{N}^{\mathbb{C}}$, $\mathfrak{N}^{A,\mathbb{C}}$. Their estimates therefore hold in each
of these norms, given the holomorphy property \eqref{4.21:eq-links-readers-b} in that norm.

For a wrapped $F$, every member of $\mathcal{P}_F$ is a frame with law sites, by the description
of the pockets of a wrapped term. At points with
$|\mathsf{s}(\Delta)|=e^{\kappa_1}$, the carried gauge map $\mathsf{u}(\mathsf{s})$ is not
$G$-valued. Hence on the full polydisc $F$ is evaluated at arguments in the complexified group
$G^{\mathbb{C}}$, and the only size functional bounding the map \eqref{4.21:eq-links-readers-b}
there is the one taken over the complexified parameter domain, $\mathfrak{N}^{\mathbb{C}}$. This is
why, for wrapped terms, the bounds stated in $\mathfrak{N}^A$ are read in $\mathfrak{N}^{A,\mathbb{C}}$.
\end{proof}

\begin{lemma}[The sourced first-stage expansion at one link]\label{4.21:lem-first-stage-expansion}
Fix a link of a site in the sense of Definition~\ref{4.21:def-links-readers}. It consists of:
\begin{itemize}
\item a site geometry $(\pi^{\mathfrak{s}},\mathfrak{H},\mathfrak{F})$ of Definition~\ref{4.21:def-site-geometry}, whose hubs are the components of the exempted set $\mathsf{K}$;
\item a source datum $b$;
\item the link map $\mathsf{s}\mapsto\mathsf{R}(\mathsf{s})$.
\end{itemize}
Let $F$ be a stored boundary term with region $X_F$, read through the link map as in Definition~\ref{4.21:def-links-readers}. Here $F(\mathsf{R}(\mathsf{s}))$ stands for
$F(\mathsf{R}(\mathsf{s})\restriction X_F)$, as in \eqref{4.21:eq-links-readers-b}. Assume:
\begin{itemize}
\item the holomorphy condition \eqref{4.21:eq-links-readers-b};
\item the reader condition \eqref{4.21:eq-links-readers-a};
\item the conclusions of Lemma~\ref{4.21:lem-exempted-set} for this link.
\end{itemize}
Let $\mathcal{R}$ be the closed set of \eqref{4.21:eq-links-readers-a} on which $F$ reads the link's response. Then the following hold.
\begin{enumerate}
\item[(i)] (Identity.) We have
\begin{equation}\label{4.21:eq-first-stage-expansion-1}
F(\mathsf{R}(\mathbf{1}))-F(\mathsf{R}(0))=\sum_{\mathsf{y}\in\mathfrak{Y}(F)}F_{\mathsf{y}},
\end{equation}
where
\begin{equation}\label{4.21:eq-first-stage-expansion-2}
F_{\mathsf{y}}:=\int_{[0,1]^{\mathsf{y}}}\prod_{\Delta\in\mathsf{y}}d\mathsf{s}(\Delta)\,
\partial_{\mathsf{s}(\Delta)}\,F(\mathsf{R}(\mathsf{s}))\Big|_{\mathsf{s}=0\ \text{off}\ \mathsf{y}} .
\end{equation}
The sum runs over the finite non-empty families $\mathsf{y}\subset\pi^{\mathfrak{s}}$ with the following property: every component of $\mathsf{y}$, for the adjacency of Definition~\ref{4.21:def-site-geometry}, contains a source cell and a member that sees $\mathcal{R}$ in the sense of Definition~\ref{4.21:def-site-geometry}. The set of these families is the set $\mathfrak{Y}(F)$ of admissible families of cells for the reader $F$ of that definition.
\item[(ii)] (Bound.) For every $\mathsf{y}\in\mathfrak{Y}(F)$,
\begin{equation}\label{4.21:eq-first-stage-expansion-3}
|F_{\mathsf{y}}|\le\mathfrak{N}(F)\,e^{-(\kappa_1-1)|\mathsf{y}|},
\end{equation}
with $\mathfrak{N}(F)$ the bound of \eqref{4.21:eq-links-readers-b}. More generally, let $c$ be independent of $\mathsf{s}$. If $|F(\mathsf{R}(\mathsf{s}))-c|\le\mathfrak{n}$ on the polydisc $\{|\mathsf{s}(\Delta)|\le e^{\kappa_1}:\Delta\in\pi^{\mathfrak{s}}\}$ of \eqref{4.21:eq-links-readers-b}, then
$|F_{\mathsf{y}}|\le\mathfrak{n}\,e^{-(\kappa_1-1)|\mathsf{y}|}$. Moreover $F_{\mathsf{y}}$ is linear in $F$. It is holomorphic, uniformly, in every parameter in which $F\circ\mathsf{R}$ is holomorphic.
\item[(iii)] (Locality.) Put
\begin{equation}\label{4.21:eq-first-stage-expansion-4}
\mathsf{K}_{\mathsf{y}}:=\bigcup\bigl\{\mathsf{h}\in\mathfrak{H}:\ \mathsf{h}\ \text{shares a wall with a member of}\
\mathsf{y},\ \text{or}\ \mathsf{h}\cap\mathcal{R}\neq\emptyset\bigr\}.
\end{equation}
Then $F_{\mathsf{y}}$ reads $b$ only on $\mathsf{K}_{\mathsf{y}}\cup\bigcup\mathsf{y}$. It reads $((\mathbf{U},\mathbf{J}),\mathbf{A},w)$ only on $X_F\cup\mathsf{K}_{\mathsf{y}}\cup\bigcup\mathsf{y}$.
\item[(iv)] (Invariance.) $F_{\mathsf{y}}$ has the invariance of $F$ stated in \cite[(2.41)(iii)]{B-CMP119}.
\end{enumerate}
\end{lemma}

\begin{proof}
Put $f(\mathsf{s}):=F(\mathsf{R}(\mathsf{s})\restriction X_F)$. Its arguments are $\mathsf{s}=(\mathsf{s}(\Delta))_{\Delta\in\pi^{\mathfrak{s}}}$, a family indexed by the finite set $\pi^{\mathfrak{s}}$. For $\mathsf{y}\subset\pi^{\mathfrak{s}}$ write $\partial_{\mathsf{y}}:=\prod_{\Delta\in\mathsf{y}}\partial_{\mathsf{s}(\Delta)}$.

\emph{Proof of (i).} By \eqref{4.21:eq-links-readers-b}, $f$ is holomorphic on the polydisc $\{|\mathsf{s}(\Delta)|\le e^{\kappa_1}\}$. This polydisc contains $[0,1]^{\pi^{\mathfrak{s}}}$. The interpolation formula \cite[(1.9)]{B-CMP116} (Lemma~\ref{4.4:lem-vertex-formula}) therefore applies to $f$ on the finite set $\pi^{\mathfrak{s}}$ and gives
\begin{equation}\label{4.21:eq-first-stage-expansion-5}
f(\mathbf{1})=\sum_{\mathsf{y}\subset\pi^{\mathfrak{s}}}\int_{[0,1]^{\mathsf{y}}}
\prod_{\Delta\in\mathsf{y}}d\mathsf{s}(\Delta)\,\partial_{\mathsf{y}}f(\mathsf{s})
\Big|_{\mathsf{s}=0\ \text{off}\ \mathsf{y}} .
\end{equation}
The term $\mathsf{y}=\emptyset$ is $f(0)$. For $\mathsf{y}\neq\emptyset$ the term is $F_{\mathsf{y}}$ of \eqref{4.21:eq-first-stage-expansion-2}. It remains to show that $F_{\mathsf{y}}=0$ unless $\mathsf{y}\in\mathfrak{Y}(F)$.

Fix $\mathsf{y}\neq\emptyset$ and let $\mathsf{y}_a$ be the components of $\mathsf{y}$ for the adjacency of Definition~\ref{4.21:def-site-geometry}. The set $Y(\sigma)$ of Lemma~\ref{4.21:lem-exempted-set}, taken at $\sigma=\mathsf{s}$ with $\mathsf{s}=0$ off $\mathsf{y}$, is $\mathsf{K}\cup\bigcup\mathsf{y}$. The components of $Y(\sigma)$ that contain members of $\mathsf{y}$ are the sets
\begin{equation}\label{4.21:eq-first-stage-expansion-6}
Y_a:=\bigcup\mathsf{y}_a\cup\bigcup\bigl\{\mathsf{h}\in\mathfrak{H}:\ \mathsf{h}\ \text{shares a wall with a member of}\
\mathsf{y}_a\bigr\}.
\end{equation}
Indeed, connectedness in \cite{B-CMP116} is connectedness through walls. Call a \emph{wall-component} of $\mathsf{y}$ a component of $\mathsf{y}$ for the adjacency by shared walls only, hubs being ignored. Then two wall-components of $\mathsf{y}$ lie in one component of $Y(\sigma)$ if and only if a chain of hubs and wall-components of $\mathsf{y}$ joins them. This is exactly the adjacency of Definition~\ref{4.21:def-site-geometry}.

Suppose first that $\mathsf{y}_a$ contains no source cell. If $\mathfrak{H}\neq\emptyset$, the source cells are, by Definition~\ref{4.21:def-site-geometry}, the members sharing a wall with a hub. So no member of $\mathsf{y}_a$ shares a wall with a hub, and $Y_a=\bigcup\mathsf{y}_a$ contains no hub. If $\mathfrak{H}=\emptyset$ there is no hub at all. In both cases $Y_a$ contains no hub and no source cell. Since $b$ vanishes off $\mathsf{K}\cup\bigcup\mathfrak{F}$ (Definition~\ref{4.21:def-links-readers}), $b=0$ on $Y_a$. By the reader condition \eqref{4.21:eq-links-readers-a}, two things follow:
\begin{itemize}
\item the triple $\mathcal{A}(\mathsf{s})=(\mathcal{Y}(\mathsf{s};b),\mathcal{J}[\mathcal{Y}(\mathsf{s};b)],\mathsf{u}(\mathsf{s}))$ of \eqref{4.21:eq-links-readers-a} equals $(0,0,1)$ on $Y_a$, whatever $\mathsf{s}\restriction\mathsf{y}_a$ is;
\item on every other component $C$ of $Y(\sigma)$, $\mathcal{A}(\mathsf{s})\restriction C$ is a function of $(b\restriction C,\mathsf{s}\restriction C,(\mathbf{U},\mathbf{J})\restriction C)$, so it does not read $\mathsf{s}\restriction\mathsf{y}_a$.
\end{itemize}
Off $Y(\sigma)$, the locality clause (ii) of Lemma~\ref{4.21:lem-exempted-set}, at parameters vanishing off $\mathsf{y}$, gives that $\mathcal{A}(\mathsf{s})$ does not read $\mathsf{s}\restriction\mathsf{y}_a$ either. Hence $\mathcal{A}(\mathsf{s})$ off $Y_a$ does not read $\mathsf{s}\restriction\mathsf{y}_a$. Since $F$ reads the $\mathsf{s}$-dependent objects only through $\mathcal{A}(\mathsf{s})$, we get $\partial_{\mathsf{s}(\Delta)}f=0$ for $\Delta\in\mathsf{y}_a$. Hence the integrand $\partial_{\mathsf{y}}f=\partial_{\mathsf{y}\setminus\{\Delta\}}\partial_{\mathsf{s}(\Delta)}f$ vanishes, and $F_{\mathsf{y}}=0$.

Suppose next that no member of $\mathsf{y}_a$ sees $\mathcal{R}$. By the definition of seeing in Definition~\ref{4.21:def-site-geometry}, $\mathcal{R}$ then meets neither $\bigcup\mathsf{y}_a$ nor a hub of $Y_a$, so $\mathcal{R}\cap Y_a=\emptyset$. By \eqref{4.21:eq-links-readers-a}, $F$ reads $\mathcal{A}(\mathsf{s})$ only on $\mathcal{R}$. Off $Y_a$, as shown in the previous paragraph (on the other components of $Y(\sigma)$ by \eqref{4.21:eq-links-readers-a}, and off $Y(\sigma)$ by the locality clause (ii) of Lemma~\ref{4.21:lem-exempted-set}), $\mathcal{A}(\mathsf{s})$ does not read $\mathsf{s}\restriction\mathsf{y}_a$. So $f$ reads nothing on $Y_a$, and again $\partial_{\mathsf{s}(\Delta)}f=0$ for $\Delta\in\mathsf{y}_a$, giving $F_{\mathsf{y}}=0$.

Note that a wall-component of $\mathsf{y}$ may miss $\mathcal{R}$ but share a hub with a wall-component that meets $\mathcal{R}$. Such a wall-component does not give a vanishing term; this is why the adjacency joins members through hubs. Thus only the families in $\mathfrak{Y}(F)$ remain in \eqref{4.21:eq-first-stage-expansion-5}. Subtracting $f(0)$ gives \eqref{4.21:eq-first-stage-expansion-1}.

\emph{Proof of (ii).} Let $\mathsf{y}\neq\emptyset$ and let $c$ be independent of $\mathsf{s}$. Then $\partial_{\mathsf{y}}f=\partial_{\mathsf{y}}(f-c)$. Fix $\mathsf{s}$ with $\mathsf{s}(\Delta)\in[0,1]$ for $\Delta\in\mathsf{y}$ and $\mathsf{s}=0$ off $\mathsf{y}$, and put $r:=e^{\kappa_1}-1$. By Cauchy's formula in the variables $\sigma(\Delta)$, $\Delta\in\mathsf{y}$, with $\sigma=0$ off $\mathsf{y}$,
\begin{equation}\label{4.21:eq-first-stage-expansion-7}
\partial_{\mathsf{y}}f(\mathsf{s})=\oint_{|\sigma(\Delta)-\mathsf{s}(\Delta)|=r,\ \Delta\in\mathsf{y}}
\ \prod_{\Delta\in\mathsf{y}}\frac{d\sigma(\Delta)}{2\pi i\,(\sigma(\Delta)-\mathsf{s}(\Delta))^{2}}
\,\bigl(f(\sigma)-c\bigr).
\end{equation}
The circles lie inside the polydisc, because $|\sigma(\Delta)|\le\mathsf{s}(\Delta)+r\le e^{\kappa_1}$. On them $|f-c|\le\mathfrak{n}$, so $|\partial_{\mathsf{y}}f(\mathsf{s})|\le\mathfrak{n}\,r^{-|\mathsf{y}|}$. Since $r=e^{\kappa_1}-1\ge e^{\kappa_1-1}$, this gives $|\partial_{\mathsf{y}}f(\mathsf{s})|\le\mathfrak{n}\,e^{-(\kappa_1-1)|\mathsf{y}|}$. Integrating over $[0,1]^{\mathsf{y}}$, a set of volume one, gives the general bound. The choice $c=0$, $\mathfrak{n}=\mathfrak{N}(F)$, which is admissible by \eqref{4.21:eq-links-readers-b}, gives \eqref{4.21:eq-first-stage-expansion-3}.

By \eqref{4.21:eq-first-stage-expansion-2} and \eqref{4.21:eq-first-stage-expansion-7}, $F_{\mathsf{y}}$ is an integral, over the compact set $[0,1]^{\mathsf{y}}$ and the circles, of the values of $F\circ\mathsf{R}$ against a kernel that depends on no other parameter. Hence $F_{\mathsf{y}}$ is linear in $F$. It is also holomorphic in every parameter in which $F\circ\mathsf{R}$ is holomorphic, and the bound above is uniform in that parameter.

\emph{Proof of (iii).} The integrand of $F_{\mathsf{y}}$ is evaluated at parameters vanishing off $\mathsf{y}$. The claim follows from the locality clause (ii) of Lemma~\ref{4.21:lem-exempted-set} at such parameters, together with \eqref{4.21:eq-links-readers-a}. By the latter, $F$ reads the $\mathsf{s}$-dependent objects only through $\mathcal{A}(\mathsf{s})\restriction\mathcal{R}$. On each component $C$ of $Y(\sigma)$, $\mathcal{A}(\mathsf{s})\restriction C$ is a function of $(b\restriction C,\mathsf{s}\restriction C,(\mathbf{U},\mathbf{J})\restriction C)$. The components read in this way are the sets $Y_a$ and the hubs met by $\mathcal{R}$. Their union lies in $\mathsf{K}_{\mathsf{y}}\cup\bigcup\mathsf{y}$.

\emph{Proof of (iv).} This follows from clause (c) of the theorem that constructs the layer $\mathcal{Y}(\mathsf{s};b)$, which holds at the decoupled parameters with the same constants by Lemma~\ref{4.21:lem-exempted-set}, together with \cite{B-CMP116}.

No disc in an amplitude parameter enters this proof. The separation between the source cells and $\mathcal{R}$ enters only through the vanishing of the terms outside $\mathfrak{Y}(F)$ in (i).
\end{proof}

\begin{lemma}[Holomorphy of wrapped terms at every link]\label{4.21:lem-wrapped-holomorphy}
Let $F$ be a wrapped stored boundary-type term in the format of the decoupling clause
(Definition~\ref{4.21:def-decoupling-clause}), with region $X_F$. In particular $F$ is analytic,
over $X_F$, on the space on which it is constructed. Take as hypotheses, without proving them here,
the following statements: the containment statements (a)--(c) of the construction of the
correlation theorem, at Ba{\l}aban's decoupling family $\sigma_0$, and the bookkeeping statement (d):
\begin{enumerate}
\item[(a)] at the site (T): the homogeneous containment lemma;
\item[(b)] at the site (P): the localised stage-space corollary, together with the weighted
containment bound at the preparation site read over $X_F$;
\item[(c)] at the site (L): the key containment theorem, together with its two conjuncts, the
inheritance statement and the arrested-term statement;
\item[(d)] the statement on wrapped descendants, together with its rider at interpolation time
$\mathsf{t}=0$, up to the ordering site.
\end{enumerate}
Then the holomorphy property \eqref{4.21:eq-links-readers-b} holds for $F$ at every link
(Definition~\ref{4.21:def-links-readers}) of the decoupling clause
\eqref{4.21:eq-decoupling-clause-a}.
\end{lemma}

\begin{proof}
The hypotheses (a)--(c) assert that the relevant objects lie in their domains for Ba{\l}aban's choice of
the family $\sigma_0$. What has to be shown is that the same memberships, and the holomorphy that
comes with them, persist when $\sigma_0$ is replaced by the families that the decoupling clause
prescribes at each site. We treat the three sites in turn.

\emph{The site \textup{(T)}.} Here we use the proof of the homogeneous containment lemma, not
only its statement. That proof consists of five displays quoted from Ba{\l}aban's papers and three
deductions from them. Its conclusion is that the containment \cite[(3.17)]{B-CMP109} is a bound on
a single number, homogeneous of degree one in $\sup g_k|B|$, where $B$ is Ba{\l}aban's field in
that containment. The homogeneity rests on the following sentence of \cite[\S3]{B-CMP109}:
\begin{quote}
Let us stress that it follows only from the fact that $H_j(B(t))$ satisfies (3.14) with
$\frac{1}{3}\alpha_2$, and $\delta H_j$ satisfies (3.9). Later we will change these functions by a
localization procedure, but new localized functions will satisfy the bounds (3.14), (3.9), hence
(3.17) will be preserved also.
\end{quote}
The family $\sigma_0$ enters this proof at exactly two places. The first is
\cite[(1.21)]{B-CMP116}, which is one of the five quoted displays. The second is one of the three
deductions: under $|\sigma|\le e^{\kappa_1}$ the constant $C_H$ of the lemma becomes
$C_He^{16\kappa_1}$. Both are instances of Lemma~\ref{4.21:lem-exempted-set}, so both hold for every
subfamily $\sigma_0$ of the partition, with the same constants.

At $\mathsf{t}=0$ the hypothesis of the homogeneous containment lemma reads
\begin{equation*}
K_Rg_k\sup|\mathsf{B}|\le\alpha_{*,k},
\end{equation*}
where $\mathsf{B}$ is the field whose supremum, multiplied by $g_k$, this hypothesis controls,
$K_R$ is the constant multiplying $g_k\sup|\mathsf{B}|$ in it, and $\alpha_{*,k}$ is the threshold
at level $k$ that it imposes. This holds in the construction of the
correlation theorem, because the parameter $\theta^{\mathbb{C}}$ fixed in that construction
satisfies $\theta^{\mathbb{C}}\le\frac12$.

Over the strip part of $X_F$ the layer schedule used in the second deduction is the one given by
the layer-schedule statement, which holds on every layer $n\le k$. The tube of $F$ there is the
tube of its ancestor, by the tube statement over strips. The chart at this site is the one supplied
by the chart statement at the site (T).

Hence the containment, and with it the holomorphy, holds at the site (T) for the family
$\sigma_0=\pi^{T}$ of all cells that the decoupling clause prescribes.

\emph{The site \textup{(P)}.} Here we use statements only. The localised stage-space corollary
asserts that
\begin{equation*}
|t_\square|\le r_\square^{\sharp}\quad\text{for all blocks $\square$ of the cover simultaneously,}
\end{equation*}
holomorphically in $(\mathbb{U},\mathbb{J},b,\sigma,t)$; here $t=(t_\square)$ are the coordinates
of the stage space, one for each block $\square$ of the cover, and $r^{\sharp}_\square$ is the radius
that the corollary assigns to $\square$. The blocks of the cover are all blocks,
strip blocks included. The stage space is imposed shell by shell over the region $\hat X$ of that
corollary. The clauses that fix the decoupling family at the preparation site already decouple
every cube, including the cubes inside the region $S$ over which the weighted containment bound at
the preparation site is stated. The chart is the one of the chart statement at the preparation site.

The only requirement going beyond these statements is the weighted containment bound at the
preparation site with $S$ replaced by $X_F$. This is hypothesis (b). In the construction of the
correlation theorem it comes from two sources. At the locations in the set $Z^{\mathrm{on}}$ of that
construction it is the third containment condition of that construction. Off those locations it
comes from the tube statement over strips.

\emph{The site \textup{(L)}: the statement used.} At this site we read the proof of the key
containment theorem for its dependence on $\sigma_0$. In that theorem
$(\hat{\mathcal{U}}^{\mathfrak{b}},\hat{\mathcal{J}}^{\mathfrak{b}})$ denotes the pair presented at
the base point $\mathfrak{b}$ and $h$ its presentation; $(W_s,J_s)$ is the pair at stage $s$ and
$\operatorname{arg}_s$ the argument read at stage $s$, $\varepsilon_\sharp$ is the smallness
threshold for that argument, $\gamma_0$ is the parameter of the regularity class used below,
$\operatorname{supp}\Pi$ is the set over which the real field $\boldsymbol{\Phi}=\varpi$ ranges, and
$t$ are the coordinates of the stage space, as at the site (P). The entries and their thresholds are
those named in the statement of that theorem. The theorem concerns
\begin{itemize}
\item every element $\mathfrak{E}\in\mathcal{C}$;
\item every $\zeta$ in the source polydisc $\mathfrak{S}_{\mathrm{src}}$;
\item every real $\boldsymbol{\Phi}=\varpi\in\operatorname{supp}\Pi$;
\item every $\sigma$ with $|\sigma|\le e^{\kappa_1}$;
\item every base point $\mathfrak{b}\in\{\mathfrak{b}_2,\mathfrak{b}_3,\mathfrak{b}_4\}$, and each of
the pairs $(\mathfrak{b}_2,\mathfrak{b}_3)$ and $(\mathfrak{b}_3,\mathfrak{b}_4)$.
\end{itemize}
It has a containment clause and a holomorphy clause. The containment clause reads
\begin{equation}\label{4.21:eq-wrapped-holomorphy-1}
\bigl((\hat{\mathcal{U}}^{\mathfrak{b}})^{h},\operatorname{Ad}_h\hat{\mathcal{J}}^{\mathfrak{b}}
\bigr)\restriction X\in\mathfrak{D}_{\mathfrak{E}},
\end{equation}
and in addition each entry is at most $\frac1{16}$ of its threshold on
$X\cap(\Lambda^{\natural}\cup\mathfrak{N}^{\mathrm{rec}})$. The holomorphy clause states that
\begin{equation}\label{4.21:eq-wrapped-holomorphy-2}
(\zeta,\sigma,t)\longmapsto F^{\mathfrak{b}}_{\mathfrak{E}}\quad\text{is holomorphic,}
\end{equation}
where $F^{\mathfrak{b}}_{\mathfrak{E}}$ denotes the element $\mathfrak{E}$ read at the base point
$\mathfrak{b}$.

The class $\mathcal{C}$ consists of four kinds of terms:
\begin{itemize}
\item the old terms of Ba{\l}aban's expansion denoted there $\mathbf{E}$, $\mathbf{R}$ and
$\mathbf{B}^{(j)}$, with $j\le k$ and $X\cap Z\neq\emptyset$;
\item the terms of the domains $\mathbb{B}''$;
\item the $Y_0$-terms and the terms of \cite[(1.59)]{B-CMP122b};
\item the frozen terms of arrested components.
\end{itemize}

The pair $(\mathfrak{b}_3,\mathfrak{b}_4)$ is the bracket of the operation $\mathbf{B}$. The stage
$s=4$ corresponds to \cite[(1.51)]{B-CMP122b}, and the stage $s=3$ to \cite[(1.54)]{B-CMP122b}, with
\begin{equation*}
\mathfrak{j}_3:=\tilde G_0(\sigma)\mathsf{J}_3,
\end{equation*}
where $\mathsf{J}_3$ is the right-hand bracket of \cite[(1.55)]{B-CMP122b}. The elements of
$\mathcal{C}$ are of two kinds, elements of a first part and elements of a second part. An element of
the first part runs this one bracket and is otherwise read at $\mathfrak{b}_3$ through $\sigma$. An
element of the second part runs the stages $s=4,\dots,1$ on $X\subset Z^c$, as in the second-part
claim.

The statement concerns the globally presented chain, presented by $h$ and restricted to $X$, which is
the true domain of the term. The pair is restricted to this domain $X$ itself; no auxiliary domain
enters the statement in place of $X$. This is the expansion used at the site (L), and
it is the one through which the carriers are read.

\emph{The site \textup{(L)}: where $\sigma$ enters the proof.} The key containment theorem is proved
by induction on $s=1,\dots,4$. The induction hypothesis at stage $s$ is that the pair
$(W_s,J_s)$ lies in Ba{\l}aban's regularity class $\mathfrak{K}(\mathfrak{B}_s;\gamma_0)$ of
\cite{B-CMP109}, and that
\begin{equation*}
\|\operatorname{arg}_s\|_1\le\varepsilon_\sharp .
\end{equation*}
From this hypothesis the proof derives, in order:
\begin{enumerate}
\item[(1)] the layer theorem and clause (i) of its corollary;
\item[(2)] the localisation clause at stage $s$;
\item[(3)] the entries at $\mathfrak{b}_s$;
\item[(4)] the induction hypothesis at stage $s+1$.
\end{enumerate}
The parameter $\sigma$ is read only through the statements used in this induction: the layer theorem
and its corollary, the companion lemma of the layer theorem, the rider of the localisation lemma,
and the localisation clause. Each of them reads $\sigma$ only through the localisation lemma for
decoupled operators.

\emph{The site \textup{(L)}: independence of the choice of $\sigma_0$.} In every statement named in
the preceding paragraph, $\sigma_0$ is an arbitrary subfamily of the partition. None of them
specifies $\sigma_0$. None of them has the hypothesis that $X$ is contained in the exempted set.

By Lemma~\ref{4.21:lem-exempted-set}, each of these statements holds for every
$\sigma_0\subset\pi$, with the same constants. Hence the containment clause
\eqref{4.21:eq-wrapped-holomorphy-1} and the holomorphy clause \eqref{4.21:eq-wrapped-holomorphy-2}
hold for every $\sigma_0$, with $\mathsf{s}$ free also on the cells of $X$.

The same applies to the increments $P$, $P'$, $P''$ of the second-part claim, which enter the
second part. They are suprema of the layers over the part of $X$ lying in the region where that
claim takes its suprema, and each
is bounded by a constant multiple of $e^{-3\delta_1R_k}$; here $\delta_1$ is the constant of the
standing condition $\delta_1M\ge16\kappa_1$ of the localisation lemma for decoupled operators and
$R_k$ is the radius at level $k$ fixed in the
second-part claim. This bound is obtained from the layer bound, with the constant $B_H$ of the
layer bound read on the polydisc. By clause (iii) of Lemma~\ref{4.21:lem-exempted-set} these bounds
on the increments hold for every $\sigma_0$ as well.

\emph{The site \textup{(L)}: the position of $\sigma_0$.} Seven places in the arguments at this
site read where the members of $\sigma_0$ lie. Between them they impose a single demand.

The first is the proof of the localisation of the second current at the step where the far bound is
obtained. There $G''_kQ''^{*}=0$ is an algebraic identity of the full operator ($\sigma=\mathbf{1}$,
the all-ones parameter point),
and in the bracket only walks touching $\sigma_0$ survive. Ba{\l}aban takes $\sigma_0$ to be
\begin{quote}
the set \ldots\ of cubes $\Delta$ disjoint with $\Lambda$
\end{quote}
\cite{B-CMP122b}. Since $D_2\subset Z''_{h+1}\subset\Lambda$, a walk starting from $z\in Z_h$ must
leave $\Lambda$. This argument supplies the far reading-cost margin $\ell 6$ for a law site in
$Z_h$, which is used in the bookkeeping of the theorem on the region $\Lambda$.

Five further places read the position of $\sigma_0$:
\begin{itemize}
\item the kernel bound for $K(\sigma)-K(\mathbf{1})$, where $K(\sigma)$ is the decoupled operator of
Lemma~\ref{4.21:lem-exempted-set}, between points of $\Lambda$, which is stated under
$\sigma_0\cap\Lambda=\emptyset$;
\item the summation over the shells, which uses $\sigma_0\cap\Lambda=\emptyset$;
\item the statement that, because $\sigma_0\cap\Lambda=\emptyset$ on the shells, a walk carrying
$\sigma$ from a tower point has crossed every shell;
\item the statement that the cubes of $\sigma_0$ lie off $\Lambda$;
\item the statement that $Q_b\subset\Lambda^{a}$ meets no cube of $\sigma_0$.
\end{itemize}
These six places ask for one thing only:
\begin{equation}\label{4.21:eq-wrapped-holomorphy-3}
\Delta\cap\Lambda=\emptyset\qquad\text{for every }\Delta\in\sigma_0 .
\end{equation}
The family prescribed by \eqref{4.21:eq-decoupling-clause-a} at the first step at the site (L)
consists of cells disjoint
with $\Lambda^{\sim}$. The three later families consist of cells disjoint with $\Lambda$. All four
therefore satisfy \eqref{4.21:eq-wrapped-holomorphy-3}.

The seventh place is a locality statement. It says that the other cubes meeting $Q_b$ meet
$\Lambda$, or, for $s=4$, meet $X$ or $\Lambda^{\sim}$. It names Ba{\l}aban's exempt set at the
first operation. Together with the statements that locality holds within the piece domain $Y_4$ and
that Ba{\l}aban's $\sigma_0$-localisations of $\mathbf{H}''_k$, of his functions $H_{B_1}$ and
$H_k$, and of the bookkeeping of
\cite[(1.66)]{B-CMP122b} are unchanged, it amounts to locality within $Y_4$. Under the decoupling
clause this is clause (iii) of Lemma~\ref{4.21:lem-first-stage-expansion}, read with the domains of
the second-stage pieces.

It follows that the bookkeeping of the statement on wrapped descendants, hypothesis (d), stands
without change under the decoupling clause. A value, or a piece of the $\sigma$-localisation,
carries no interpolation time: it is the piece at $\mathsf{t}=0$. So the bound for the pieces at
$\mathsf{t}=0$ in the statement on wrapped descendants applies to it directly, with
$\eta:=\eta_0$, where $\eta$ is the parameter in the link map of
Definition~\ref{4.21:def-links-readers} and $\eta_0$ is the value fixed for it in that bound. It
enters with the reading-cost margins $\ell 5$, $\ell 6$ and $\ell 7$ of hypothesis (d) and with
nothing added, for deep terms as well.

Conversely, if a presented chain had a family with a member inside $\Lambda$, the argument for the
far bound described above would not apply to it. The demand \eqref{4.21:eq-wrapped-holomorphy-3}
binds only the families of a chain that is presented, summed, charted and written as a Cauchy
integral. A field in which cells inside $\Lambda$ are switched off, but which is never presented,
summed, charted or written as a Cauchy integral, and for which the identity $G''Q''^{*}=0$ is not
read, is not subject to it. The families of the presented chain are exactly the families of the
decoupling clause, which satisfy \eqref{4.21:eq-wrapped-holomorphy-3}.

\emph{The site \textup{(L)}: the role of Ba{\l}aban's exemption of $X$.} In Ba{\l}aban's
construction the cells of $X$ are exempted at this step. The exemption serves the one-component form
of Ba{\l}aban's expansion. The decoupling clause uses instead the interpolation formula for the
parameter $\mathsf{s}$ prescribed in Definition~\ref{4.21:def-decoupling-clause}, together with
component-wise vanishing, which is clause (i) of
Lemma~\ref{4.21:lem-first-stage-expansion}.

Ba{\l}aban's construction itself cuts the cells of $X$ at three of the four localisations, and of
the first of these it states \cite{B-CMP122b}:
\begin{quote}
the analyticity domain is restricted on the cube $\square_0$, or on the domain $X$, by the
analyticity properties of the original term. Outside the cube $\square_0$, or the domain $X$, this
analyticity domain is quite large, restricted only by the analyticity properties of the function
$\mathbf{H}''_k$.
\end{quote}
On $X$ the restriction in question is membership in $\mathfrak{D}_{\mathfrak{E}}$, and this is the
containment clause \eqref{4.21:eq-wrapped-holomorphy-1}.

\emph{The site \textup{(L)}: the composed chain.} The pieces of the second-stage expansion are
Cauchy integrals over the product polydisc in
$(\mathsf{s}^{(1)},\dots,\mathsf{s}^{(4)})$. Each link has its own family. The induction in the
proof of the key containment theorem runs through the four links in turn, each at
$|\sigma|\le e^{\kappa_1}$. Therefore the holomorphy clause \eqref{4.21:eq-wrapped-holomorphy-2} is
joint holomorphy in all four parameters $\mathsf{s}^{(1)},\dots,\mathsf{s}^{(4)}$ together, which is
what the argument $\sigma$ stands for.

It remains to identify the domain of $F$ over the strips. Over a live strip it is given by the tube
statement over strips. Over a dead strip, which is contained in the core $\mathsf{C}_C$ of a
component $C$, with $\mathsf{C}_C\subset\Lambda^{\natural}$,
it is given by the containment clause \eqref{4.21:eq-wrapped-holomorphy-1} itself. If $F$ is a chain
term or a frozen term, one uses in addition the inheritance statement and the arrested-term statement
of hypothesis (c).

\emph{Conclusion.} At each of the three sites, the hypotheses give membership and holomorphy at
Ba{\l}aban's family $\sigma_0$. The arguments above carry both over to the families of the decoupling
clause, jointly in the parameters of the product polydisc. This is \eqref{4.21:eq-links-readers-b}
at every link of the decoupling clause.
\end{proof}

\begin{definition}[The minimal form of the third-factor lemma]\label{4.21:def-third-factor-form}
We work at site (L) of the operation $\mathbf{R}_k$, in the setting and notation of the
third-factor lemma. Thus $C$ is a component of the large-field region $Z$, $\mathsf{C}_C$ is its
core, and $d^{\wedge}$ is the scaled distance with respect to $\mathbf{B}''_k$. The decay rate is
$\delta_P=\tfrac18\delta_0$, where $\delta_0$ is the decay rate fixed in the setting of the
third-factor lemma. For $t=1,\dots,4$, link $t$ carries its decoupling family
$\sigma_0^{(t)}$ and its parameters $\mathsf{s}^{(t)}$, and
$\vec{\mathsf{s}}=(\mathsf{s}^{(1)},\dots,\mathsf{s}^{(4)})$. We write
$\operatorname{chain}(\vec{\mathsf{s}})$ for the chain of the four links at the parameters
$\vec{\mathsf{s}}$. The numbers $r^{0}>0$ and $C_{3F}=2/r^{0}$ are the constants of that lemma,
$\mathfrak{G}_{\varepsilon'}(X)$ is its chart of radius $\varepsilon'=\varepsilon-\eta_0-\mu$, and
$\mathfrak{N}^{\mathbb{C}}_{\varepsilon}(v)$ is the size of a unit $v$ at the radius $\varepsilon$
in force, where $\tfrac34\mathsf{e}\le\varepsilon\le\mathsf{e}$ and $\mathsf{e}$ is the radius
function of the setting. In the radius $\varepsilon'$, $\eta_0$ is the supremum over the parameter
polydisc of the gauge size of the carried gauge map, and $\mu$ is the reduction of the radius defined
in that lemma; its value $\mu(s)$ at a law site $s$ is recalled in the law-site bound below.

A unit $v$ with true domain $X_v$ is \emph{deep} if $X_v\subset\mathsf{C}_C$. For a wrapped unit in
the sense of Definition~\ref{4.21:def-decoupling-clause}, $X_v=A_v\cup\bigcup Y^{\sharp}$, the union
running over the list $\mathcal{P}_v$ of that definition. We put
$D_X:=d^{\wedge}(X,\mathsf{C}_C^{c})$.

\begin{enumerate}
\item The \emph{oscillation bound} is clause (a) of the third-factor lemma, read at
$\vec{\mathsf{s}}{}'=\vec{0}$ and $\mathfrak{b}=\mathfrak{b}_4$. Here
$v^{\sharp}(\operatorname{chain}(\vec{\mathsf{s}}))$ stands for the function
$F^{\mathfrak{b}_4}(m;g')$ of that lemma at $m=(\zeta,\varpi,\vec{\mathsf{s}})$. The bound says
that for every deep unit $v$, all $m=(\zeta,\varpi,\vec{\mathsf{s}})$ in the parameter polydisc
$\mathfrak{P}$ of that lemma (so that $|\mathsf{s}^{(t)}(\Delta)|\le e^{\kappa_1}$ for every $t$ and
every cell $\Delta$), all $g'\in\mathfrak{G}_{\varepsilon'}(X_v)$ and all values of the remaining
arguments of $F^{\mathfrak{b}_4}$,
\begin{equation}\label{4.21:eq-third-factor-form-1}
\bigl|v^{\sharp}(\operatorname{chain}(\vec{\mathsf{s}}))
-v^{\sharp}(\operatorname{chain}(\vec{0}))\bigr|
\le C_{3F}\,e^{-\frac12\delta_P\,d^{\wedge}(X_v,\mathsf{C}_C^{c})}\,
\mathfrak{N}^{\mathbb{C}}_{\varepsilon}(v).
\end{equation}
The oscillation is taken over the four families jointly.
The first link is decoupled as the first-link rule of \eqref{4.21:eq-decoupling-clause-a}
prescribes, that is, by the fundamental theorem of calculus in $\mathsf{s}^{(1)}$ alone and with no
amplitude factor.

\item The \emph{piece bound} is clause (b) of the third-factor lemma. Let $v$ be any deep unit,
plain or wrapped, and for either base point $\mathfrak{b}\in\{\mathfrak{b}_3,\mathfrak{b}_4\}$ of
the third-factor lemma let $F^{\mathfrak{b}}$ be the function of that lemma formed from $v$. Let
$(\mathsf{y}_1,\dots,\mathsf{y}_4)$ be any tuple of families $\mathsf{y}_t\subset\sigma_0^{(t)}$,
not all empty; the tuples in which only $\mathsf{y}_1$ is non-empty are included. Define the piece
\begin{equation*}
\mathfrak{p}:=\prod_{t}\prod_{\Delta\in\mathsf{y}_t}\int_0^1 d\mathsf{s}^{(t)}(\Delta)\,
\partial_{\mathsf{s}^{(t)}(\Delta)}F^{\mathfrak{b}}\Big|_{\mathsf{s}^{(t)}=0\text{ off }\mathsf{y}_t}.
\end{equation*}
The piece bound says that
\begin{equation}\label{4.21:eq-third-factor-form-2}
|\mathfrak{p}|\le C_{3F}\,e^{-\frac12\delta_P D_{X_v}}\,\mathfrak{N}^{\mathbb{C}}_{\varepsilon}(v)\,
e^{-(\kappa_1-1)\sum_t|\mathsf{y}_t|}
\end{equation}
holds on the whole chart $\mathfrak{G}_{\varepsilon'}(X_v)$. Since $\mathfrak{p}$ is linear in the
unit, the same bound holds with $\mathfrak{N}^{\mathbb{C}}_{\varepsilon}$ replaced by
$\mathfrak{N}^{A}$ and by $\mathfrak{N}^{A,\mathbb{C}}$.

\item The \emph{law-site bound} is clause (c) of the third-factor lemma, in the form of its
law-site lemma; law sites and frames are those of that law-site lemma. Let $v$ be a deep unit with
true domain $X=X_v$, and let $s$ be a law site of $v$, either a pointwise law site or a law site of
a frame $W$ of any age;
$\Sigma_W$ denotes the set of law sites of $W$, and $h$, $k_W$ are the levels fixed for
$s\in\Sigma_W$ in the law-site lemma. Put
$\mathsf{d}(s):=0$ if $s$ is pointwise, and $\mathsf{d}(s):=h-k_W$ if $s\in\Sigma_W$; this is the
depth of the law site. Let $Z(s)$ denote the quantity at $s$ that the law-site lemma estimates; it
is not the large-field region $Z$. In the bound below, $K_u$ is the bound on the carried gauge map,
$\alpha_{*,k}$ is the threshold radius at level $k$, $w=M/M_1$, $\check{R}_k$ is the region size
parameter at step $k$, $\mathsf{e}(s)$ is the value at $s$ of the radius function $\mathsf{e}$
above, and $\mu_k$,
$\varepsilon_{\star}$ are the constants through which the radius reduction $\mu$ is defined. The
auxiliary disc has radius $\rho_X:=r^{0}e^{\frac12\delta_P D_X}$. The bound says that for
$|\lambda|\le\rho_X$
\begin{equation}\label{4.21:eq-third-factor-form-3}
\begin{aligned}
|\lambda|\,|Z(s)|\le\vartheta^{3F}(s)
&:=K_u\,\alpha_{*,k}\,e^{-\frac12\delta_P w\check{R}_k}\,e^{-2(\mathsf{d}(s)-1)_+}\\
&\le\tfrac12\,\mu_k\,\mathsf{e}(s)/\varepsilon_{\star}.
\end{aligned}
\end{equation}
This is asserted under the hypothesis lettered (vi)$''$, under clause~(ix) of the conditions on
live regions, and under the reference
hypothesis for pairs of level at most $k$. The factor $e^{-2(\mathsf{d}(s)-1)_+}$ is the decay in
the depth of the law site that the reference hypothesis supplies for the frame, and it is kept
explicit in the bound. The right-hand member
is $\tfrac12\mu(s)$, where $\mu(s)=\mu_k\mathsf{e}(s)/\varepsilon_{\star}$ is the reduction of the
chart radius in $\varepsilon'=\varepsilon-\eta_0-\mu$. So the law-site bound uses one half of that
reduction, and no further one.

\item The \emph{minimal form of the third-factor lemma} is the following statement. There are
$C_{3F}$ as above and a rate
\begin{equation*}
\varsigma_3\in[\tfrac18\delta_P,\delta_P]
\end{equation*}
with this property. Let $v$ be any wrapped unit with $X_v\subset\mathsf{C}_C$, and let
$(\mathsf{y}_t)$ be any tuple of families $\mathsf{y}_t\subset\sigma_0^{(t)}$ with
$\bigcup_t\mathsf{y}_t\not\subset Z$. Then the piece $\mathfrak{p}$ of the piece bound formed from
$v$, $(\mathsf{y}_t)$ and either base point $\mathfrak{b}\in\{\mathfrak{b}_3,\mathfrak{b}_4\}$
satisfies
\begin{equation}\label{4.21:eq-third-factor-form-a}
|\mathfrak{p}|\le C_{3F}\,e^{-\varsigma_3\,d^{\wedge}(X_v,\mathsf{C}_C^{c})}\,\mathfrak{N}(v)\,
e^{-(\kappa_1-1)\sum_t|\mathsf{y}_t|},
\end{equation}
where $\mathfrak{N}$ is any one of the norms of \eqref{4.21:eq-links-readers-c}, on the chart
$\mathfrak{G}_{\varepsilon'}(X_v)$.
\end{enumerate}

The interval allowed for $\varsigma_3$ is wider than the value $\tfrac12\delta_P$ that the piece
bound supplies, so that the statement on deep units is proved from
\eqref{4.21:eq-third-factor-form-a} whatever the value of $\varsigma_3$ in this interval, in
particular for $\varsigma_3=\tfrac18\delta_P$. The theorem on the relative format at complex sources
uses the third-factor lemma itself and
not only \eqref{4.21:eq-third-factor-form-a}. The reason is that for deep units which are plain,
that is, not wrapped, it needs the piece bound \eqref{4.21:eq-third-factor-form-2}.

If the piece bound \eqref{4.21:eq-third-factor-form-2} holds, then the minimal form
\eqref{4.21:eq-third-factor-form-a} holds with $\varsigma_3=\tfrac12\delta_P$.
\end{definition}

\begin{proof}[Proof of the last assertion]
Assume \eqref{4.21:eq-third-factor-form-2}, and put $\varsigma_3:=\tfrac12\delta_P$. Since
$\delta_P>0$, we have
\begin{equation*}
\tfrac18\delta_P\le\tfrac12\delta_P\le\delta_P,
\end{equation*}
so $\varsigma_3$ lies in the required interval. The constant is the same $C_{3F}$.

Let $v$ be a wrapped unit with $X_v\subset\mathsf{C}_C$. Then $v$ is deep, and its true domain is
$X_v=A_v\cup\bigcup Y^{\sharp}$. By definition,
$D_{X_v}=d^{\wedge}(X_v,\mathsf{C}_C^{c})$, and hence
\begin{equation*}
e^{-\frac12\delta_P D_{X_v}}=e^{-\varsigma_3\,d^{\wedge}(X_v,\mathsf{C}_C^{c})}.
\end{equation*}

Let $(\mathsf{y}_t)$ be a tuple of families $\mathsf{y}_t\subset\sigma_0^{(t)}$ with
$\bigcup_t\mathsf{y}_t\not\subset Z$. The empty set is contained in $Z$, so
$\bigcup_t\mathsf{y}_t\neq\emptyset$. Hence at least one $\mathsf{y}_t$ is non-empty, and the tuple
is not all empty. The piece bound therefore applies to $v$ and this tuple. After the substitution
just displayed, it is \eqref{4.21:eq-third-factor-form-a} with
$\mathfrak{N}=\mathfrak{N}^{\mathbb{C}}_{\varepsilon}$, valid on the whole chart
$\mathfrak{G}_{\varepsilon'}(X_v)$.

The piece bound also holds with $\mathfrak{N}^{A}$ and with $\mathfrak{N}^{A,\mathbb{C}}$ in place
of $\mathfrak{N}^{\mathbb{C}}_{\varepsilon}$, because the piece is linear in the unit. Together
these give \eqref{4.21:eq-third-factor-form-a} in each of the norms of
\eqref{4.21:eq-links-readers-c}, which are these three norms.
\end{proof}

\begin{remark}[Why the deep class is generic]\label{4.21:rem-deep-class}
Let $C$ be a component of $Z$ in the setting of the third-factor lemma, let $\mathsf{C}_C$ be its
core, and let $h$ be the level at which the core and the radius $\check{R}_h$ are defined. For a
stored term $v$ we use the label $A_v$, the list $\mathcal{P}_v$ and the region
$X_v=A_v\cup\bigcup_{\mathcal{P}_v}Y^{\sharp}$ of Definition~\ref{4.21:def-decoupling-clause}. A
stored term is called \emph{deep} if its true domain is contained in $\mathsf{C}_C$; for a wrapped
term $v$ (that is, one with $\mathcal{P}_v\neq\emptyset$) the true domain is $X_v$. For wrapped
terms this is the hypothesis under which Definition~\ref{4.21:def-third-factor-form} bounds the
pieces of $v$. The symbol $\operatorname{dist}_h$ denotes the level-$h$ distance in which the core
$\mathsf{C}_C$ is defined.

In the notation of \cite{B-CMP122b}, whose letters are used in the next display and in this
paragraph only, the factor
\begin{equation*}
\exp\Bigl(-\delta\, d\bigl((Y_1^0)^{c}\cap Y_1\cap Y'_{i_0},X\bigr)\Bigr)
\end{equation*}
is stated in \cite{B-CMP122b} without proof for $Y'_2\cup Y'_3\cup Y'_4\neq\emptyset$; there one
half, $\tfrac12\delta d$, of this decay is kept to control the resummations. For the first
large-field operation $\mathbf{R}$ the sum over the regions is controlled in a different way:
\cite{B-CMP122b} refers for it to bounds of an earlier paper, and in \cite{B-CMP116} this sum is
controlled by the first exponential in \cite[(1.30)]{B-CMP116}. That exponential is a factor of the
Cauchy bound in the interpolation parameters $t_\square$ of \cite{B-CMP116}, which the decoupling at
the first link in Definition~\ref{4.21:def-decoupling-clause} does not use. The tuples in which only
the family $\mathsf{y}_1$ is non-empty are therefore covered here by clause (b) of the third-factor
lemma, its bound on pieces.

Deep terms cannot be treated together with the terms covered by the locality lemma for the chain.
What is needed for a deep term $v$ is the oscillation, in the decoupling parameters
$\mathsf{s}$, of the chain of the four links at site (L) of $\mathbf{R}_k$ on $X_v\subset\Lambda$,
where $\Lambda$ is the region of the setting of the third-factor lemma and the families
$\sigma_0^{(t)}$ are taken with no cell meeting the interior of $\Lambda$, so that every family
$\sigma_0^{(t)}$ lies off $\Lambda$. The stretch from $X_v$ to $\partial\Lambda$ is
not cut, and the locality lemma for the chain reads $\sigma_0$ only through the datum $m(\omega)$,
so it does not see the distance from $X_v$ to $\partial\Lambda$; and the fourth step of the proof
of the localisation bound used at site (L) of $\mathbf{R}_k$ does not apply to a chain cut inside
$\Lambda$.

The class of deep terms is generic, in the following sense.
\begin{enumerate}
\item[(a)] Let $v=\mathbf{B}'^{(m)}(\operatorname{sk}(X))$ be the term of the witness construction,
with $X=Y^{\sharp}\cup\Delta_1$,
where $\mathcal{P}_v=\{Y\}$, the member $Y$ is of level $m\le h$ and lies in $C$, $Y^{\sharp}$ is
its strip, and $\Delta_1$ is the set adjoined to $Y^{\sharp}$ in the skeleton construction of
$\operatorname{sk}(X)$. Suppose $\check{R}_h>3$. Then $X_v\subset\mathsf{C}_C$.
\item[(b)] Let $v$ be a wrapped term of level $j\le h$. Suppose that every member of
$\mathcal{P}_v$ lies in $C$, and that
\begin{equation}\label{4.21:eq-deep-class-1}
d_j(A_v)<\tfrac12\bigl(\check{R}_h-1\bigr)L^{h-j}-1 .
\end{equation}
Then $v$ is deep.
\end{enumerate}
Consequently, for a wrapped term $v$ that is deep by (a), under the condition $\check{R}_h>3$
stated there, or by (b), the following rest on the third-factor lemma at every step at which one of
the members of $v$'s own list $\mathcal{P}_v$ is healed: clauses (d) and (e) of the proposition on
the healing step into which the third-factor lemma enters, and through them three further statements
at that step, namely the statement obtained from those two clauses, the bound on $\alpha^{\sharp}$,
and the creation of the strips born at that step. The third-factor lemma is the statement whose
minimal form is Definition~\ref{4.21:def-third-factor-form}. It is therefore used at every such
healing; this is the sense in which the deep class is generic.
\end{remark}

\begin{proof}
\emph{Part (a).} For every member $Y$ of a list $\mathcal{P}_v$ that lies in $C$, clause (g) of the
geometric lemma that defines the core $\mathsf{C}_C$ gives
\begin{equation}\label{4.21:eq-deep-class-2}
\operatorname{dist}_h\bigl(Y^{\sharp},\mathsf{C}_C^{c}\bigr)\ \ge\ \tfrac12\bigl(\check{R}_h+1\bigr).
\end{equation}
The member $Y$ of $\mathcal{P}_v$ in (a) lies in $C$, so its strip $Y^{\sharp}$ satisfies
\eqref{4.21:eq-deep-class-2}. For the term
$v=\mathbf{B}'^{(m)}(\operatorname{sk}(X))$ with $X=Y^{\sharp}\cup\Delta_1$, the witness
construction places the label $A_v$ within two cubes of level $m$ of the strip $Y^{\sharp}$, that
is, at level-$h$ distance at most
$2L^{m-h}$ from it. Since $\mathcal{P}_v=\{Y\}$, the region $X_v$ is the union of $A_v$ with the
strip $Y^{\sharp}$. By \eqref{4.21:eq-deep-class-2}, $Y^{\sharp}$ is contained in $\mathsf{C}_C$.
For the label, since $m\le h$ we have $2L^{m-h}\le 2$, and since $\check{R}_h>3$ by the hypothesis
of (a) we have
$\tfrac12(\check{R}_h+1)>2$; hence the level-$h$ length $2L^{m-h}$ of two level-$m$ cubes is
smaller than the lower bound $\tfrac12(\check{R}_h+1)$ in \eqref{4.21:eq-deep-class-2}. For
$a\in A_v$ the triangle inequality then gives
\begin{equation*}
\operatorname{dist}_h\bigl(a,\mathsf{C}_C^{c}\bigr)
\ \ge\ \tfrac12\bigl(\check{R}_h+1\bigr)-2L^{m-h}\ >\ 0 .
\end{equation*}
Hence $A_v\subset\mathsf{C}_C$, and therefore $X_v\subset\mathsf{C}_C$.

\emph{Part (b).} Let $v$ be as in (b). By Definition~\ref{4.21:def-decoupling-clause}, $X_v$ is
the union of $A_v$ with the strips $Y^{\sharp}$ of the members of $\mathcal{P}_v$. We show that
each of these sets lies in $\mathsf{C}_C$.

The strips come first. Every member of $\mathcal{P}_v$ lies in $C$, so by the geometric lemma its
strip obeys \eqref{4.21:eq-deep-class-2}. Its level-$h$ distance from $\mathsf{C}_C^{c}$ is
therefore at least $\tfrac12(\check{R}_h+1)>0$, and hence the strip is contained in
$\mathsf{C}_C$.

Now the label. The inequality \eqref{4.21:eq-deep-class-1} is read through the following property
of the label of a wrapped term of level $j$ and of the size function $d_j$, which we take from the
definition of $d_j$ (Notation~\ref{0.2:not-glossary}) together with the format of a wrapped term in
Definition~\ref{4.21:def-decoupling-clause}: for every point $a$ of
$A_v$ there is a member of $\mathcal{P}_v$ whose strip $Y^{\sharp}$ satisfies
\begin{equation}\label{4.21:eq-deep-class-3}
\operatorname{dist}_j\bigl(a,Y^{\sharp}\bigr)\ \le\ d_j(A_v)+1,
\end{equation}
where $\operatorname{dist}_j$ is the distance in level-$j$ units. One level-$h$ unit
equals $L^{h-j}$ level-$j$ units, so by \eqref{4.21:eq-deep-class-3} the point $a$ lies at
level-$h$ distance at most
$\bigl(d_j(A_v)+1\bigr)L^{j-h}$ from that strip. We rewrite \eqref{4.21:eq-deep-class-1} by
adding $1$ to both sides and multiplying by $L^{j-h}>0$; this gives
\begin{equation}\label{4.21:eq-deep-class-4}
\bigl(d_j(A_v)+1\bigr)L^{j-h}\ <\ \tfrac12\bigl(\check{R}_h-1\bigr).
\end{equation}
It follows that
\begin{equation*}
\operatorname{dist}_h\bigl(a,Y^{\sharp}\bigr)<\tfrac12\bigl(\check{R}_h-1\bigr).
\end{equation*}
By the triangle inequality for $\operatorname{dist}_h$, together with \eqref{4.21:eq-deep-class-2},
we obtain
\begin{align*}
\operatorname{dist}_h\bigl(a,\mathsf{C}_C^{c}\bigr)
&\ge \operatorname{dist}_h\bigl(Y^{\sharp},\mathsf{C}_C^{c}\bigr)
      -\operatorname{dist}_h\bigl(a,Y^{\sharp}\bigr)\\
&> \tfrac12\bigl(\check{R}_h+1\bigr)-\tfrac12\bigl(\check{R}_h-1\bigr)\ =\ 1 .
\end{align*}
Hence $a\in\mathsf{C}_C$. Since $a$ was an arbitrary point of $A_v$, we get $A_v\subset\mathsf{C}_C$.
Together with the inclusion of the strips, this gives $X_v\subset\mathsf{C}_C$. The term $v$ is
wrapped by hypothesis, so it is deep.

\emph{The consequence.} Deep wrapped terms are those to which
Definition~\ref{4.21:def-third-factor-form} applies. By (a), under its condition
$\check{R}_h>3$, and by (b), the term of (a) and every
wrapped term whose label satisfies \eqref{4.21:eq-deep-class-1}, with all members of its list in
$C$, are deep. At every step at which one of the members of such a term $v$'s own list
$\mathcal{P}_v$ is healed, the statements listed before the proof are therefore applied to a deep
term. Hence they use the third-factor lemma at that step.
\end{proof}

\begin{remark}[The four layers are read on the global chain]\label{4.21:rem-global-chain}
At site (L) a unit $v$ reads the global chain in its $h$-presented form, that is, presented
through the gauge transformation $h$ of the statement on the global chain, on its own true
domain $X_v$. In particular, on the domain $D_2\supset\mathsf{C}_C$ it reads
\begin{equation*}
  e^{i\eta\mathbb{H}_4}\,e^{i\eta\mathbb{H}_3}\,U_2
\end{equation*}
on $X_v$ itself, by the frame identity of the third-factor bound. Here $D_2$, $U_2$ and $\eta$
are the domain, the configuration and the parameter in the exponents of that frame identity,
and $\mathsf{C}_C$ is the core of the component $C$. The global layers
$\mathbb{H}_4,\dots,\mathbb{H}_1$, through which the chain is carried from link to link, are
called the \emph{carriers}, and reading links two to four on them is called the
\emph{carriers' reading}. Its features are the following.
\begin{enumerate}
\item The four layers $\mathbb{H}_4,\mathbb{H}_3,\mathbb{H}_2,\mathbb{H}_1$ are the functions
  supplied by the layer holomorphy statement at the determining sets
  $\mathbf{B}''_k,\mathbf{B}''_k,\mathbf{B}_1,\mathbf{B}_k$ of the construction, respectively.
\item Each layer is decoupled over its own family $\sigma_0^{(t)}$.
\item The chain, as a function of $\vec{\mathsf{s}}=(\mathsf{s}^{(t)})_t$, is one function,
  jointly holomorphic on the product of the parameter polydiscs of the links, by the joint
  holomorphy clause for the chain.
\item No localised set occurs. The layer holomorphy statement is used only at a determining set
  of the construction.
\end{enumerate}

The alternative reading of links two to four, which would substitute the representations
\cite[(1.61)--(1.63)]{B-CMP122b} at determining sets
$\mathbf{B}''_k(Y^{\mathrm{tr}}_1)$, $\mathbf{B}_1(Y^{\mathrm{tr}}_2)$ and
$\mathbf{B}_k(Y^{\mathrm{tr}}_3)$ formed on the scaffold domains, is not used in this book. The
scaffold domains
$Y^{\mathrm{tr}}_t$ of \eqref{4.21:eq-decoupling-clause-a} enter only the domains of the
pieces and the reading sets $\mathcal{R}_{t-1}$.

The carriers' reading has the following consequences for the localisation at site (L).
\begin{enumerate}
\item[(a)] The index sets of the localisation step are those of
  Lemma~\ref{4.21:lem-first-stage-expansion}\,(i), taken at the reading sets $\mathcal{R}_{t-1}$.
  Each such set contains everything the carriers read: the true domain $X_v$, the earlier
  families and the hubs. For instance, $\mathbb{H}_4$ reads $\mathbb{H}_3$ on
  $\mathsf{K}\cup\bigcup\mathsf{y}_1$.
  In addition, $\mathcal{R}_{t-1}$ contains the whole layer of $Y^{\mathrm{tr}}_{t-1}$, a set
  larger than what is read through it.
\item[(b)] A family $\mathsf{y}$ that sees $\mathcal{R}_{t-1}$ only through that layer gives
  the piece $F_{\mathsf{y}}\equiv0$.
\item[(c)] Such families are kept in the index sets. As a result, the index sets contain the
  families that meet the boundary layer prescribed in \cite{B-CMP122b}, of width at most
  $2M_1$ in the corresponding scale, and the majorants can only grow.
\end{enumerate}
\end{remark}

\begin{proof}
We first give three reasons for reading links two to four on the global chain.

First, at site (L) holomorphy of a wrapped term $F$ at the link is the conclusion of
Lemma~\ref{4.21:lem-wrapped-holomorphy}. That lemma is proved as an implication from the
containment statements named in it, and its base is a statement on the global chain. In the
reading through \cite[(1.61)--(1.63)]{B-CMP122b}, holomorphy and membership are instead the
content of \cite[(1.65)]{B-CMP122b}. Of that display, \cite{B-CMP122b} says that it ``has been
proved already, in a slightly less general form, in the proof of the equality (1.89) [IV]'', and
adds ``we do not repeat them here''.
The argument of this chapter establishes the corresponding statement only for the carriers'
reading.

Second, in the representations \cite[(1.61)--(1.63)]{B-CMP122b} the chain at link $t$ depends
on the families $(\mathsf{y}_1,\dots,\mathsf{y}_{t-1})$ through $Y^{\mathrm{tr}}_{t-1}$.

Third, the third-factor lemma, whose minimal form is
Definition~\ref{4.21:def-third-factor-form}, is proved for the carriers' reading only.

We turn to the consequences. The reading set $\mathcal{R}_{t-1}$ of link $t$ is by construction
the set holding what is read at that link. Hence the index sets of the localisation step are
those of Lemma~\ref{4.21:lem-first-stage-expansion}\,(i) with $\mathcal{R}=\mathcal{R}_{t-1}$.
This is consequence (a).

The layer of $Y^{\mathrm{tr}}_{t-1}$ is contained in $\mathcal{R}_{t-1}$ in addition to what the
carriers read.
A family that sees $\mathcal{R}_{t-1}$ only through this layer gives the zero piece, which is
consequence (b). Keeping such families therefore adds only zero terms to the identity of
Lemma~\ref{4.21:lem-first-stage-expansion}\,(i), which consequently still holds.

With these families kept, the index sets contain the families that meet the boundary layer
prescribed in \cite{B-CMP122b}, of width at most $2M_1$ in the corresponding scale. The
majorant of a sum is a sum of non-negative bounds over the index set, and
a sum of non-negative terms over a larger index set is no smaller. This gives consequence (c).

Finally, the convention removes nothing that is proved: by the first and third reasons above,
holomorphy and the third-factor bound are established here only for the carriers' reading.
Whether the relative form of the layer holomorphy
statement would also carry containment to localised sets is not decided here.
\end{proof}

\begin{remark}[Order of dependence between the theorem and its inputs]\label{4.21:rem-input-order}
The theorem of this section is the relative bound for Ba{\l}aban's boundary terms $\mathbf{B}'$ in the
norms $\mathfrak{N}^{\mathbb{C}}$, together with its table of constants. We call it \emph{the relative
bound}. It stands in the following order of dependence with two of its inputs, which are taken here by
statement.
\begin{itemize}
\item The tube bound for the two kept members. In the notation in which the tube bound is stated, let
$Y_i$ be a second-class component created at step $k_W$. The two kept members are the two functionals
of $(\mathbf{U},\mathbf{J})$ kept under $\mathbf{T}'_{k_W}(Y_i)$. The tube bound is a bound on these two
functionals.
\item The two-class statement on stored terms. Let $\mathfrak{S}$ be the least class of stored terms
that contains every term ever wrapped and the two kept members, and that is closed under passing to
terms built from its elements. Let $\mathfrak{A}$ be the class of the remaining stored terms. The
statement has one conclusion for $\mathfrak{A}$ and one for $\mathfrak{S}$, each with its own list of
hypotheses. It also supplies a list of auxiliary constructions for all stored terms, whatever their
class. Its last clause is a theorem asserted for all stored terms. That theorem closes with an
inequality $\eta<\varepsilon$, where $\varepsilon$ is the radius of the charts
$\mathfrak{G}_{\varepsilon}$ and $\eta$ is the quantity defined in that theorem which it bounds by
$\varepsilon$.
\end{itemize}
These statements are related as follows.
\begin{enumerate}
\item[(1)] The conclusion of the two-class statement for $\mathfrak{S}$ is asserted under several
hypotheses. By its statement, one of them is the bound on the terms $\mathbf{B}'$ in the norms
$\mathfrak{N}^{\mathbb{C}}$. Taken with its table of constants, this hypothesis is the relative bound.
In turn, the relative bound uses the auxiliary constructions that the two-class statement supplies for
all stored terms.
\item[(2)] Neither the relative bound nor the two-class statement is a hypothesis of the other as a
whole. The relative bound enters only the hypotheses of the conclusion for $\mathfrak{S}$. Only the
auxiliary constructions of the two-class statement enter the relative bound.
\item[(3)] The two kept members arise in operations of step $k_W$, and the tube bound bounds them
there. The tube bound is used only at later steps, and there only by the two-class statement. It enters
through one share in the residual part of the third clause of the bound on live contributions, namely
the share that belongs to outputs using a kept member. Hence the closing inequality $\eta<\varepsilon$,
on the charts $\mathfrak{G}_{\varepsilon}$ used by the relative bound, depends on the tube bound.
\item[(4)] No estimate of this section uses the tube bound. The two kept members do not belong to the
class of stored terms fixed for the estimates of the construction.
\item[(5)] Every instance of the relative bound is taken on a frame, hence on charts
$\mathfrak{G}_{\varepsilon}$ on which the closing inequality $\eta<\varepsilon$ depends on the tube
bound. The relative bound therefore depends on the tube bound through its frames, although, by
item~(4), no estimate of this section uses the tube bound.
\end{enumerate}
\end{remark}

\begin{proof}
Item~(1) comes from the two statements themselves. The bound on $\mathbf{B}'$ in
$\mathfrak{N}^{\mathbb{C}}$ appears by name among the hypotheses for $\mathfrak{S}$. The proof of the
relative bound uses the auxiliary constructions that the two-class statement supplies for all stored
terms.

Item~(2) follows from item~(1) and from the order in which the statements of the construction are
arranged. In that order each statement uses only statements that precede it. The auxiliary constructions
of the two-class statement precede the relative bound. The relative bound precedes the conclusion of the
two-class statement for $\mathfrak{S}$.

For item~(3), the tube bound enters at one point, the residual share in the bound on live
contributions. From there it passes to the closing inequality on the charts used by the relative bound.

Item~(4) comes from the hypotheses of the estimates of this section, among which the tube bound does not
appear, and from the class of stored terms fixed for those estimates, which does not contain the two
kept members.

For item~(5), every instance of the relative bound is taken on a frame and uses the charts
$\mathfrak{G}_{\varepsilon}$ of item~(3). On these charts, by item~(3), the closing inequality
$\eta<\varepsilon$ depends on the tube bound, and the relative bound depends on the tube bound through
them.
\end{proof}

\begin{definition}[The relative format of boundary-type terms]\label{4.21:def-relative-format}
Fix the level $k$. The \emph{boundary-type terms} are the terms of the following four kinds.
\begin{itemize}
\item[($\alpha$)] The terms of $\mathbf{B}^{w}_k$: the terms $\mathbf{B}^{(j)}(\cdot,\mathcal{S})$,
$\mathcal{S}$ a strong set, produced by the operation $\mathbf{T}$, and the terms $\mathbf{B}'^{(j)}$
produced by the operation $\mathbf{R}$.
\item[($\beta$)] The frozen chain terms $\mathbb{C}^{(i)}_{k_1}$ of the pieces of live arrested
attempts, where $k_1<k$ is the step of the arrested attempt, as they stand in the arrest
construction.
\item[($\gamma$)] Inside a step $\mathbf{R}_k$: the chain terms $\mathbb{C}^{(n)}_k(\cdot,\mathcal{S})$.
\item[($\delta$)] The change-of-variables terms of an interrupted preparation, in the form in which
the filing statement for interrupted preparations enters them.
\end{itemize}
In what follows the kind of a boundary-type term is used only through clauses (f0)--(f4) below
and through the bound displayed in the lemma on boundary terms. The class of boundary-type terms
does not contain the following terms. On each region $Y_i$ to which the operation $\mathbf{T}'$
is applied in the tube bound for the source functional, two members are kept whose domain is the
strip $Y_i^{\sharp}$ itself; we write $V^{\sharp}(Y_i^{\sharp})$ for them, and they have the form of
\cite[(1.71)]{B-CMP122b}. Their domain $Y_i^{\sharp}$ lies inside the region $W$ of that bound, and
they are controlled by that bound.

We say that the \emph{relative format holds at step $k$} if, in every term of the expansion
\cite[(2.18)]{B-CMP119} obtained after the operation $\mathbf{R}_k$, every boundary-type term $F$
of level $j\le k$ carries \emph{scaffold data}
\begin{equation}\label{4.21:eq-relative-format-a}
F\ \longmapsto\ (A_F,\mathcal{P}_F),\qquad F\ \text{a boundary-type term of level}\ j\le k,
\end{equation}
consisting of a domain $A_F$, called the \emph{label} of $F$, and a set $\mathcal{P}_F$, such that
the following clauses hold.
\begin{itemize}
\item[(f0)] $A_F\in\mathbf{D}_j$, and $A_F$ satisfies the side conditions imposed on $X$ in
\cite[(2.41)]{B-CMP119}, respectively in \cite[(1.63)]{B-CMP122b} for the chain terms, in the form
\begin{equation*}
A_F\cap\Omega_j\neq\emptyset,\qquad\text{and}\qquad
A_F\cap Z_j^{\boxplus}\neq\emptyset\ \ \text{or}\ \ \mathcal{P}_F\neq\emptyset;
\end{equation*}
the alternative in the second condition is imposed on every boundary-type term. The set
$\mathcal{P}_F$ is a finite set of pairs $(Y,m)$ with $m\le j$, where $Y$ is a region of level $m$,
with strip $Y^{\sharp}$, that is live at level $k$ in the sense of the glossary
(Notation~\ref{0.2:not-glossary}), and for every $(Y,m)\in\mathcal{P}_F$ the strip
$Y^{\sharp}$ touches $A_F$. The boundary-type terms are indexed by the triple (kind, label $A$,
strong set $\mathcal{S}$): for each such triple there is at most one term.
\item[(f1)] Put
\begin{equation}\label{4.21:eq-relative-format-1}
X_F:=A_F\cup\bigcup_{(Y,m)\in\mathcal{P}_F}Y^{\sharp}.
\end{equation}
The term $F$ depends on $(\mathbf{U},\mathbf{J})$ and on $\mathbf{A}$ only through their
restrictions to $X_F$, and on $w$ only through its restriction to $X_F^{+}$. Moreover $F$ has an
orbit datum $F^{\sharp}$, in the sense of \eqref{4.21:eq-orbit-data-a}, of the current radius
$\varepsilon_k$, on
\begin{equation*}
\mathfrak{U}_F(X_F)\times\mathfrak{D}_{\mathcal{S}}\times\operatorname{supp}\boldsymbol{\Phi}
\times\mathfrak{G}_{\varepsilon_k}(X_F),
\end{equation*}
where the factors $\mathfrak{D}_{\mathcal{S}}$, $\operatorname{supp}\boldsymbol{\Phi}$ and
$\mathfrak{G}_{\varepsilon_k}(X_F)$ are those of the domain of an orbit datum in
\eqref{4.21:eq-orbit-data-a}, and $\mathfrak{U}_F$ is the stored domain of
\eqref{4.21:eq-orbit-data-a}, which lies inside the
space of the kind of $F$ over $X_F$. These spaces over $X_F$ are meaningful: the spaces
\cite[(2.34)--(2.39)]{B-CMP119}, the spaces \cite[(1.64)--(1.69)]{B-CMP122b} and the space
$\mathfrak{D}^{\flat}$ are cut out by local clauses imposed shell by shell, and are therefore
defined over any union of cells. Over a strip, the domain of the kind of $F$ lies inside the tube
of the term produced by $\mathbf{R}$ from which $F$ descends, by the tube-containment statement,
and inside its stored domain, by the stored-domain containment statement.
\item[(f2)] $F$ obeys the bound of its kind with $d_j(X)$ replaced by $d_j(A_F)$, in the size
functional $\mathfrak{N}^{A,\mathbb{C}}_{\varepsilon_k}$, written $\mathfrak{N}^{A,\mathbb{C}}$
below. For the terms of kind ($\alpha$) produced by $\mathbf{T}$ this reads
\begin{align*}
\mathfrak{N}^{A,\mathbb{C}}\bigl(\mathbf{B}^{(j)}(A,\cdot)\bigr)
&\le B_0\,e^{-\kappa_{\mathrm{out}}d_j(A)}
&&\text{for fresh terms, with }\kappa_{\mathrm{out}}\ge2\kappa,\\
\mathfrak{N}^{A,\mathbb{C}}\bigl(\mathbf{B}^{(j)}(A,\cdot)\bigr)
&\le B_0\,e^{-\kappa d_j(A)}
&&\text{for stored terms,}
\end{align*}
as in \cite[(2.3), (2.42)]{B-CMP119}; for those produced by $\mathbf{R}$ it reads
\begin{equation}\label{4.21:eq-relative-format-2}
\mathfrak{N}^{A,\mathbb{C}}\bigl(\mathbf{B}'^{(j)}(A,\cdot)\bigr)
\le e^{\lambda_{\theta}}C_{99}\,c_1^{\flat}(g_j)\,e^{-(1+\frac14\beta_s)\kappa d_j(A)}
\le B_0\,e^{-\kappa d_j(A)},
\end{equation}
where $\lambda_{\theta}$, $C_{99}$, $c_1^{\flat}(\cdot)$ and $\beta_s$ are the constants of the bound
for the terms produced by $\mathbf{R}$.
For the chain terms,
\begin{equation*}
\|\mathbb{C}^{(n)}_k\|^{A,\mathbb{C}}_{(n)}\le C_0^{\sharp},
\end{equation*}
where the norm is taken with $a_k=\frac85\kappa$ at the top level and $a_m=\frac45\kappa$ at the
levels $m<k$, and the sum and the weight defining the norm run over labels. For the terms of kind
($\delta$), the bound is the one displayed in the filing statement for interrupted preparations.
\item[(f3)] $A_F$ and $\mathcal{P}_F$ depend on the scaffold only: they do not depend on $z$, and
they are the same for two constructions $\mathsf{A}$, $\mathsf{B}$ carried out on one scaffold,
for differences, and for marked and dial terms.
\item[(f4)] (Anchoring.) $A_F$ has a cube in $(Z_j^{\boxplus})^{(1)}$.
\end{itemize}
When $\mathcal{P}_F=\emptyset$, the relative format is the format of \cite[(2.41)--(2.42)]{B-CMP119}
for the terms of kind ($\alpha$), and of \cite[(1.63), (1.70)]{B-CMP122b} for the chain terms,
together with the orbit datum of
\eqref{4.21:eq-orbit-data-a}. The relative format asserts nothing about $d_j(X_F)$ or about
$d_j([X_F]_j)$.
\end{definition}

\begin{proof}[Proof that clause \textup{(f4)} follows from clause \textup{(f0)}]
Let $F$ be a boundary-type term of level $j\le k$ whose scaffold data $(A_F,\mathcal{P}_F)$ satisfy
(f0). By the alternative in (f0), either $A_F\cap Z_j^{\boxplus}\neq\emptyset$ or
$\mathcal{P}_F\neq\emptyset$.

Suppose first that $\mathcal{P}_F\neq\emptyset$, and pick $(Y,m)\in\mathcal{P}_F$; by (f0),
$m\le j$. If $m<j$, the strip $Y^{\sharp}$, which is a strip of level $m<j$, lies in $Z_j$, by the
strip-inclusion statement; since $Z_j^{\boxplus}$ is the one-layer enlargement of $Z_j$, it
contains $Z_j$, so $Y^{\sharp}\subset Z_j^{\boxplus}$. If $m=j$, the strip $Y^{\sharp}$ lies in
$Y^{\boxplus}$, and $Y^{\boxplus}\subset Z_j^{\boxplus}$ by the same statement; so again
$Y^{\sharp}\subset Z_j^{\boxplus}$. By (f0) the strip $Y^{\sharp}$ touches $A_F$, so some cube of
$A_F$ touches $Y^{\sharp}$ and therefore touches $Z_j^{\boxplus}$. A cube touching
$Z_j^{\boxplus}$ lies in $(Z_j^{\boxplus})^{(1)}$; hence $A_F$ has a cube in
$(Z_j^{\boxplus})^{(1)}$.

Suppose next that $A_F\cap Z_j^{\boxplus}\neq\emptyset$. Then a cube of $A_F$ meets
$Z_j^{\boxplus}$, hence touches it, and by the fact used in the previous case it lies in
$(Z_j^{\boxplus})^{(1)}$. In both cases (f4) holds.
\end{proof}

\begin{theorem}[Propagation of the complex relative format from step to step]\label{4.21:thm-format-propagation}
Let the parameters of the correlation theorem satisfy, besides the floors imposed in that theorem,
\begin{equation}\label{4.21:eq-format-propagation-a}
\begin{aligned}
&\tfrac14\beta_s\kappa\ \ge\ \lambda_{\theta},
&&\sigma_{*}\ \ge\ c_d+4,\\
&\kappa_1\ \ge\ 1+c_{l}(2\kappa+4)+\lambda_{*}(d,L),
&&\tfrac14\delta_P\,w\ \ge\ 110\,d\,(a_k+1),
\end{aligned}
\end{equation}
where $w=M/M_1$.
In the first line, $\sigma_{*}:=\min\{\beta\kappa,\kappa/(4c^{\sharp})\}$ is the surplus of the estimate at the preparation
site, $c^{\sharp}$ being a constant entering that estimate. The first inequality is a floor on $\kappa$ depending on $d$
only, and the second is a floor on $\kappa$
depending on $(d,\beta)$. Both are implied by the floor $\kappa_d$ of the correlation theorem, so they impose no new floor
on $\kappa$. Let the parameters further satisfy:
\begin{itemize}
\item the floor \eqref{4.21:eq-preparation-constants-a} on $\kappa_1-1$, imposed after the choice of $L$ and $\kappa$;
\item the floor on the ratio $w=M/M_1$ that is stated as a condition in its own right, separately from the fourth
inequality in \eqref{4.21:eq-format-propagation-a}.
\end{itemize}
Let $\gamma$ satisfy:
\begin{itemize}
\item the conditions of the correlation theorem on the region size $\check{R}_k$ and on the tail;
\item the ceilings at the fluctuation integral, \eqref{4.21:eq-fluctuation-site-a};
\item the ceilings at the large-field site, \eqref{4.21:eq-large-field-site-a};
\item the condition \eqref{4.21:eq-coupling-interval-a} at the preparation site;
\item two further ceilings on $\gamma$ imposed in the correlation theorem, each read at the value $2+\mathsf{K}^{0}$ of
its argument, where $\mathsf{K}^{0}$ is the quantity so denoted in the correlation theorem;
\item the ceilings on $\gamma$ of the third-factor lemma.
\end{itemize}
The expansions at the sites are carried out with the decoupling clause
(Definition~\ref{4.21:def-decoupling-clause}). Assume the following, each at its place in the fixed order of the
statements on which the induction rests:
\begin{enumerate}
\item[(a)] the third-factor lemma, in two of the forms displayed in \eqref{4.21:eq-third-factor-form-a}: its
oscillation bound, at plain and at wrapped units, and the third of the four bounds displayed there;
\item[(b)] the input statement at plain units, together with its two companion forms in the decoupling
parameter $\sigma$;
\item[(c)] the inputs in the fourth, fifth, seventh and tenth places of the list of inputs on which the induction
rests, each as stated in that list, and the input in the eighth place of that list: the inclusion, over a strip, of
the stored space $\mathfrak{U}_F$ of a term in the tube of its $\mathbf{R}$-born ancestor;
\item[(d)] the hereditary clauses for stored terms, for the class $\mathfrak{A}$ and for the class $\mathfrak{S}$ of
stored terms (the classes defined with those clauses), and, as the clauses for $\mathfrak{S}$ require, the tube bound
for the two members kept by the operation $\mathbf{T}'$ at the step that created a second-class component;
\item[(e)] the assignment to the format of the change-of-variables terms of an interrupted preparation;
\item[(f)] the inherited-datum hypothesis in its third and fourth parts, the oscillation bound for kept members over the
paired chart system, and the threshold condition on $\alpha^{+}_k$.
\end{enumerate}
Read the clause on the terms $\mathbf{B}$ in the large-field bounds of the correlation theorem, and the locality clause of
its proposition on the $\mathbf{R}$-operation, in the relative format of
Definition~\ref{4.21:def-relative-format}.

Then the following holds inside the single induction on pairs (step, site). That induction carries the induction
hypothesis of the correlation theorem, the hereditary clauses of hypothesis~(d), and the tube bound for the two kept
members. For every $k<K$:
if the induction hypothesis of the correlation theorem holds at level $k$ and the relative format
\eqref{4.21:eq-relative-format-a} holds at level $k$, then the relative format holds at level $k+1$.

Moreover, consider any display of the parts of the proof of the correlation theorem that treat the three sites and the
bounds following them, and suppose it carries, for a boundary-type term $F$, the distance $d_j(X)$, the quantities
$N(X)$ and $\bar{X}$ defined from the region $X$ in the correlation theorem, the core $\bar{S}_v$ of a unit $v$
(Definition~\ref{4.21:def-decoupling-clause}), the quantity $K_v$ defined from $v$ in the correlation theorem, or the
domain $\mathsf{d}(e)$ of an element $e$. Then it holds with $X$ replaced by $A_F$, in $\mathfrak{N}^{A,\mathbb{C}}$.
This applies to the
propositions of that theorem at the sites $(T)$, $(P)$, $(L)$, to the two propositions that follow them, and to the
corollary on the boundary terms. It holds with the same labels, index sets, constants, floors and ceilings, in every
norm of \eqref{4.21:eq-links-readers-c}. The radii of the outputs are as follows:
\begin{itemize}
\item $\varepsilon-\eta^{T}_k$ at $(T)$;
\item at $(P)$, the radius of the relative form of the preparation-site corollary;
\item $\varepsilon'=\varepsilon-\eta^{L}_k-\ell 7$ at $(L)$, where $\ell 7$ is the reserved radius margin; deep units
are included.
\end{itemize}
The orbit datum $\mathsf{e}$ sits on the frames created at the step.

The following are not affected: the second group of \cite[(1.99)]{B-CMP122b}, the running-shift statement, the torus
bound, and four further clauses of the correlation theorem, each as stated there.

The third-factor lemma is used at one place in the proof of this theorem, the lemma on deep units at the large-field
site, and a second time through hypothesis~(b).
\end{theorem}

At level zero the relative format is vacuous. The passage from level $k$ to level $k+1$ consists of the operation
$\mathbf{T}$, followed, at level $k+1$, by the preparation site $(P)$, the large-field site $(L)$, and the operations of
\cite[(1.90)--(1.99)]{B-CMP122b}. The proof follows this order and is given at the end of the treatment of the
large-field site further on in this chapter.

\begin{lemma}[Per-cube bound on stored wrapped terms]\label{4.21:lem-wrapped-cube-bound}
Let $i$ be the current level, and consider a moment of the operation $\mathbf{R}_i$ at one of
its three sites at which stored terms are read: the $\mathbf{T}$-step site $(\mathrm{T})$, the
preparation site $(\mathrm{P})$, and the site $(\mathrm{L})$ at which the terms of the current
chain are stored. Recall the floors
\begin{equation*}
\kappa\ge\kappa_d\ge 20(c_d+2),\qquad \tfrac12\delta\kappa\ge c_d+1,
\end{equation*}
fixed for the correlation theorem. Here $\kappa_d$ is the lower bound imposed there on the decay
rate $\kappa$, $c_d$ is the constant for which a spare rate of at least $c_d+1$ allows the
summation estimate \cite[(1.26)]{B-CMP116} to be applied, and $\delta$ is the constant that fixes
the admissible weight rate $(1-\frac12\delta)\kappa$ at $(\mathrm{T})$. Recall also $L\ge 3$.
For $s\ge 0$ put
\begin{equation}\label{4.21:eq-wrapped-cube-bound-1}
\mathsf{b}_i(s):=b_0+(s+1)\,C_N C_0^{\sharp}\,(\mathsf{T}_i+c_s)^{r_{\mathrm{pr}}},
\end{equation}
where $b_0:=2K(d)B_0+b_\delta$ at $(\mathrm{T})$, and $b_0:=2K(d)B_0+b_\delta+N(g_i)C_0^{\sharp}$
at $(\mathrm{P})$ and at $(\mathrm{L})$; this is the standard majorant of item~(3) of
Remark~\ref{4.21:rem-standard-majorant}, with $b_1=C_NC_0^{\sharp}$ and $\mathsf{T}=\mathsf{T}_i$.
Assume that the relative format of boundary-type terms
(Definition~\ref{4.21:def-relative-format}) holds at the moment considered, let $\mathfrak{L}$ be
the family of all stored wrapped terms, and put
$W(F):=\mathfrak{N}^{A}(F)\,e^{r\,d_i([A_F]_i)}$ for $F\in\mathfrak{L}$. Let the rates $r$ and
$\varsigma_0$ be as follows: at $(\mathrm{T})$, $r\le(1-\frac12\delta)\kappa$ and
$\varsigma_0=\frac25\kappa$; at $(\mathrm{P})$ and at $(\mathrm{L})$, $r\le\frac85\kappa$ and
$\varsigma_0=\frac15\kappa$. Then for every $j\le i$ and every $\square'\in\pi_j$
\begin{equation}\label{4.21:eq-wrapped-cube-bound-2}
\sum_{F\in\mathfrak{L}:\ j_F=j,\ A_F\ni\square'} W(F)\,
e^{\varsigma_0 d_j(A_F)\,\mathbf{1}[j<i]}\;\le\;\mathsf{b}_i(i-j).
\end{equation}
In particular the family $\mathfrak{L}$ satisfies the per-cube hypotheses
\eqref{4.21:eq-wrapped-terms-a} of Definition~\ref{4.21:def-wrapped-terms}, in the form without
and in the form with the depth factor, with $\mathsf{b}=\mathsf{b}_i$, for every probe.
\end{lemma}

\begin{proof}
The sum in \eqref{4.21:eq-wrapped-cube-bound-2} runs over the stored terms of one level $j$ whose
labels contain one cube $\square'$ of that level. We split it according to the kind of term: the
boundary terms $\mathbf{B}^{(j)}$ and $\mathbf{B}'^{(j)}$ created at step $j$, the chain terms of
level $j$, and the terms of type~$(\delta)$ of the classification of stored terms; this tag names
a class of terms and is unrelated to the constant $\delta$ of the rate $(1-\frac12\delta)\kappa$.
We treat the cases $j=i$ and $j<i$ separately.

\emph{The case $j=i$.} Here $\mathbf{1}[j<i]=0$, so no depth factor is present and the summand is
$W(F)$.

The terms $\mathbf{B}^{(i)}$. At $(\mathrm{T})$ such a term is stored, by the format bounds
\eqref{4.21:eq-relative-format-a} of the relative format, with rate $\kappa$ and prefactor
$B_0$. The weight $W(F)$ costs
the rate $r\le(1-\frac12\delta)\kappa$, so the rate left over is at least
$\frac12\delta\kappa\ge c_d+1$. At $(\mathrm{L})$ the term is fresh and carries the rate
$2\kappa$, against which the weight costs $r\le\frac85\kappa$; the rate left over is at least
$\frac25\kappa$, and $\frac25\kappa\ge 8(c_d+2)\ge c_d+1$ because $\kappa\ge 20(c_d+2)$. In both
cases the sum over the labels through $\square'$ of the remaining factor $e^{-(c_d+1)d_i(A_F)}$ is
bounded by the summation estimate \cite[(1.26)]{B-CMP116}, and the contribution of the terms
$\mathbf{B}^{(i)}$ is at most $K(d)B_0$.

The terms $\mathbf{B}'^{(i)}$, which occur at $(\mathrm{T})$ only, are estimated in the same way,
with the same rate $\kappa$ against $r\le(1-\frac12\delta)\kappa$, and contribute at most
$K(d)B_0$.

The chain terms at their top level $i$: at $(\mathrm{L})$ these are the terms of the current
chain, and at $(\mathrm{T})$ those of a chain interrupted at $\mathbf{R}_i$. Their rate in the
norm is $\frac85\kappa$, and $\frac85\kappa\ge r$ at both sites. For a stage of a chain the sum
over the
labels through a cube of the weighted norms of its terms is the norm of that stage itself (the
label-sum identity for chain stages), and this norm is at most $C_0^{\sharp}$. Hence each stage
contributes at most $C_0^{\sharp}$, and since there are at most $N(g_i)$ stages, the chain terms
of level $i$ contribute at most $N(g_i)C_0^{\sharp}$.

The terms of type~$(\delta)$. By the dichotomy for terms of type~$(\delta)$, such a term
either is a member of a stage sum of its chain, and is then already counted among the chain terms
above, or satisfies, with the present $W(F)$, $r$ and $\varsigma_0$, the per-cube bound by the
number $b_\delta$. Their contribution is therefore at most $b_\delta$.

Adding up, the left side of \eqref{4.21:eq-wrapped-cube-bound-2} for $j=i$ is at most
$2K(d)B_0+N(g_i)C_0^{\sharp}+b_\delta$. At $(\mathrm{P})$ and $(\mathrm{L})$ this is at most
$b_0$ by the definition of $b_0$, hence at most $\mathsf{b}_i(0)$. At $(\mathrm{T})$ the bound on
the number
of stages of a chain run at step $i$ gives $N(g_i)\le C_N(\mathsf{T}_i+c_0)^{r_{\mathrm{pr}}}$, so
that, by \eqref{4.21:eq-wrapped-cube-bound-1} with $s=0$,
\begin{equation*}
2K(d)B_0+b_\delta+N(g_i)C_0^{\sharp}\le b_0+C_NC_0^{\sharp}(\mathsf{T}_i+c_0)^{r_{\mathrm{pr}}}
=\mathsf{b}_i(0).
\end{equation*}

\emph{The case $j<i$.} The comparison of the level-$i$ size of the $\pi_i$-hull of a label with the
level-$j$ size of the label bounds the weight factor:
$e^{r d_i([A_F]_i)}\le e^{(r/L)d_j(A_F)}$. At $(\mathrm{P})$ and $(\mathrm{L})$ we have
$r\le\frac85\kappa$ and $L\ge 3$, so $r/L\le\frac8{15}\kappa$. At $(\mathrm{T})$ we have
$r\le(1-\frac12\delta)\kappa$, so $r/L\le\frac13(1-\frac12\delta)\kappa\le\frac13\kappa$. The
summand of \eqref{4.21:eq-wrapped-cube-bound-2} now carries in addition the depth factor
$e^{\varsigma_0 d_j(A_F)}$.

The terms $\mathbf{B}^{(j)}$ and $\mathbf{B}'^{(j)}$. They carry the rate $\kappa$. At
$(\mathrm{P})$ and $(\mathrm{L})$ the rate left after the weight and the depth factor is
\begin{equation*}
\kappa-\tfrac8{15}\kappa-\tfrac15\kappa=\tfrac4{15}\kappa ,
\end{equation*}
and at $(\mathrm{T})$ it is
\begin{equation*}
\kappa-\tfrac13\kappa-\tfrac25\kappa=\tfrac4{15}\kappa .
\end{equation*}
Since $\kappa\ge 20(c_d+2)$, we have $\frac4{15}\kappa\ge\frac{16}3(c_d+2)\ge c_d+1$. The
summation estimate \cite[(1.26)]{B-CMP116} over the labels through $\square'$ then bounds the
contribution of the terms $\mathbf{B}^{(j)}$ by $K(d)B_0$, and likewise that of the terms
$\mathbf{B}'^{(j)}$; together they contribute at most $2K(d)B_0$.

The chain terms of level $j$. These are the frozen chain terms of level $j$, among them the terms
of type~$(\delta)$ in the first alternative of the dichotomy, and the terms of the current chain
stored at the level $j<i$. Their rate in the norm is at least $\frac45\kappa$. At $(\mathrm{P})$
and $(\mathrm{L})$,
\begin{equation*}
\tfrac8{15}\kappa+\varsigma_0=\tfrac8{15}\kappa+\tfrac15\kappa=\tfrac{11}{15}\kappa
\le\tfrac45\kappa ,
\end{equation*}
and at $(\mathrm{T})$,
\begin{equation*}
\tfrac13\bigl(1-\tfrac12\delta\bigr)\kappa+\tfrac25\kappa\le\tfrac{11}{15}\kappa\le\tfrac45\kappa .
\end{equation*}
So the weight and the depth factor together cost no more than the rate in the norm, and by the
label-sum identity for chain stages each stage of each chain contributes at most $C_0^{\sharp}$.
The chains that hold terms of level $j$ were run at steps $h$ with $j\le h\le i$, at most one
chain at each such step. A chain run at step $h$ has at most $N(g_h)$ stages, and the bound on the
number of stages gives
\begin{equation*}
N(g_h)\le C_N(\mathsf{T}_i+c_{i-h})^{r_{\mathrm{pr}}}\le C_N(\mathsf{T}_i+c_{i-j})^{r_{\mathrm{pr}}} .
\end{equation*}
Summing over the $i-j+1$ steps, the chain terms of level $j$ contribute at most
\begin{equation*}
(i-j+1)\,C_NC_0^{\sharp}\,(\mathsf{T}_i+c_{i-j})^{r_{\mathrm{pr}}} .
\end{equation*}

The terms of type~$(\delta)$ in the second alternative of the dichotomy contribute at most
$b_\delta$, as in the case $j=i$.

Adding up, and using $b_0\ge 2K(d)B_0+b_\delta$ at each of the three sites, the left side of
\eqref{4.21:eq-wrapped-cube-bound-2} for $j<i$ is at most
\begin{equation*}
2K(d)B_0+b_\delta+(i-j+1)C_NC_0^{\sharp}(\mathsf{T}_i+c_{i-j})^{r_{\mathrm{pr}}}
\le\mathsf{b}_i(i-j).
\end{equation*}

This proves \eqref{4.21:eq-wrapped-cube-bound-2}. Since the depth factor is at least $1$, the
bound without it follows, and since \eqref{4.21:eq-wrapped-cube-bound-2} holds for every cube
$\square'$ of every level $j\le i$, both per-cube hypotheses \eqref{4.21:eq-wrapped-terms-a} hold
with $\mathsf{b}=\mathsf{b}_i$ whatever the probe.

No count of nested arrested attempts is used in this proof. The sums over levels and over the
steps at which the chains were run are carried by the polynomial growth of $\mathsf{b}_i$ in its
argument. They are paid by the depth factor in clause~(iv) of the probe bound.
\end{proof}

\begin{proposition}[The relative machine at the fluctuation integral]
\label{4.21:prop-fluctuation-site}
Fix a level $i$ and consider the fluctuation integral of the step $i\to i+1$. Assume the following
statements, in the form in which they are stated in the proof of the correlation theorem:
\begin{enumerate}
\item[(1)] the complex-potential lemma of that proof: its analyticity clauses (a) and (b), taken
over the analyticity space of a domain and over the tube $\mathcal{W}_i(\cdot)$; its weighted
clause (d); and the step of its proof in which the smallness constant $\varepsilon_V$ of the
potential enters;
\item[(2)] the two summation lemmas of that proof;
\item[(3)] the flat theorem on the fluctuation integral of that proof;
\item[(4)] the two corollaries of that proof which are used for the outputs of the step, the
second being the counterpart of \cite[(2.42)]{B-CMP119}.
\end{enumerate}
Let $\mathfrak{L}_T$ be the family of the stored wrapped terms $F$, each regarded as a reader
with data $([A_F]_i, X_F, \mathfrak{N}^{A}(F))$, and use the site geometry
$(\pi^{T},\emptyset,\mathfrak{F})$ of Definition~\ref{4.21:def-site-geometry}, where $\pi^{T}$
is the cell family at the fluctuation integral of the step $i\to i+1$, there are no hubs, and the
source cells $\mathfrak{F}$ are the level-$i$ cells holding free bonds of $\Omega_{i+1}$. Write
$F_{\mathsf{y}}$, $\mathsf{y}\in\mathfrak{Y}(F)$, for the pieces of
Lemma~\ref{4.21:lem-first-stage-expansion}. For $Y\in\mathbf{D}_i$ put
\begin{equation}\label{4.21:eq-fluctuation-site-1}
\begin{gathered}
V^{\mathrm{rel}}(Y):=\sum_{F\in\mathfrak{L}_T}\ \sum_{\mathsf{y}\in\mathfrak{Y}(F):\,
\operatorname{lab}(F,\mathsf{y})=Y}F_{\mathsf{y}},\\
Y^{+}:=Y\cup\bigcup\bigl\{Y_p^{\sharp}:\ (Y_p,m)\ \text{live},\ Y_p^{\sharp}\ \text{touches}\ Y
\bigr\}.
\end{gathered}
\end{equation}
For an arbitrary domain $X$, the set $X^{+}$ is defined in the same way with $X$ in place of $Y$.
Then the following hold.
\begin{itemize}
\item[(a)] With $\eta$ and the layer $\mathcal{Y}$ of the step, the sum of the brackets is
reproduced exactly:
\begin{equation}\label{4.21:eq-fluctuation-site-2}
\sum_{Y\in\mathbf{D}_i}V^{\mathrm{rel}}(Y)=\sum_{F\in\mathfrak{L}_T}
\bigl[F(e^{i\eta\mathcal{Y}}\mathbf{U},\dots)-F(\mathbf{U},\dots)\bigr].
\end{equation}
\item[(b)] $V^{\mathrm{rel}}(Y)$ has the properties of clauses (a) and (b) of the
complex-potential lemma, with the analyticity space of $Y$ and the tube $\mathcal{W}_i(Y)$ taken
over $Y^{+}$; it reads $\mathsf{B}$ only on the free bonds of $Y$.
\item[(c)] With $\kappa_V:=(1-2\delta)\kappa$,
\begin{equation}\label{4.21:eq-fluctuation-site-3}
|V^{\mathrm{rel}}(Y)|\le\varepsilon^{\mathrm{rel}}_i\,e^{-\kappa_V d_i(Y)},\qquad
\varepsilon^{\mathrm{rel}}_i:=C_{\Sigma}\,C_{\mathsf{a}}(1)\,\mathsf{b}_i(0)\,
e^{-\frac14\delta\kappa(7\check{R}_i-2)}.
\end{equation}
\item[(d)] The weighted clause (d) of the complex-potential lemma holds for $V^{\mathrm{rel}}$.
\item[(e)] Assume
\begin{equation}\label{4.21:eq-fluctuation-site-a}
\sup_{\mathsf{T}\ge\mathsf{T}_{\gamma}}C_{\Sigma}\,C_{\mathsf{a}}(1)
\bigl[b_0+C_N C_0^{\sharp}(\mathsf{T}+1+c_0)^{r_{\mathrm{pr}}}\bigr]
e^{-\frac14\delta\kappa(7\mathsf{T}^{r_{\mathrm{pr}}}-2)}\le\tfrac14\varepsilon_V^{*}.
\end{equation}
Then the complex-potential lemma holds for $V:=V^{\mathrm{pl}}+V^{\mathrm{rel}}$, where
$V^{\mathrm{pl}}$ is the plain potential of the step (the kept members' share of the potential is
added to $V$ elsewhere in this chapter and is not treated in this proposition); the statements
(2), (3) and (4) hold
verbatim, with \emph{local in $X$} read as \emph{reads $((\mathbf{U},\mathbf{J}),\mathbf{A},w)$
on $X^{+}$}; no output with $X\subset\Lambda_{i+1}$ contains a $V^{\mathrm{rel}}$; and the
terms of level $i+1$ born in the $\mathbf{T}$-step are in the relative format
\eqref{4.21:eq-relative-format-a} with $A:=X$ and $\mathcal{P}:=$ the pockets of the
constituents.
\end{itemize}
\end{proposition}

\begin{proof}
\emph{(a)} By Lemma~\ref{4.21:lem-source-anchored-sums}\,(a), every label
$\operatorname{lab}(F,\mathsf{y})$ lies in $\mathbf{D}_i$, so summing
\eqref{4.21:eq-fluctuation-site-1} over $Y\in\mathbf{D}_i$ regroups, for each $F$, the sum over
all $\mathsf{y}\in\mathfrak{Y}(F)$. By Lemma~\ref{4.21:lem-first-stage-expansion}\,(i) this sum
equals $F(\mathsf{R}(\mathbf{1}))-F(\mathsf{R}(0))$. At the parameter $\mathbf{1}$ the link map
gives the configuration $e^{i\eta\mathcal{Y}}\mathbf{U}$; at the parameter $0$ the layer vanishes,
$\mathcal{Y}(0)=0$, by clause (ii) of Lemma~\ref{4.21:lem-exempted-set} applied with empty
exempted set $\mathsf{K}=\emptyset$. Hence $F(\mathsf{R}(0))=F(\mathbf{U},\dots)$
and \eqref{4.21:eq-fluctuation-site-2} follows.

\emph{(b)} By clauses (ii)--(iv) of Lemma~\ref{4.21:lem-first-stage-expansion}, the piece
$F_{\mathsf{y}}$ reads $\mathsf{B}$ on $\mathsf{y}$, and $\mathsf{y}\subset Y$ for
$Y=\operatorname{lab}(F,\mathsf{y})$; it reads the remaining arguments on
$X_F\cup\bigcup\mathsf{y}$, and this set lies in $Y^{+}$, because the strips of the pockets in
$\mathcal{P}_F$ touch $A_F$, and $A_F\subset[A_F]_i\subset Y$. Thus $V^{\mathrm{rel}}(Y)$ is a
finite sum of functions each analytic over a sub-domain of $Y^{+}$, and such a sum is analytic
over $Y^{+}$ by restriction: as in \cite{B-CMP116}, each term is defined and analytic on the
corresponding space, and after restriction to the domain all the terms are defined and analytic on
it. The statement about $\mathsf{B}$ follows from the first sentence of this paragraph.

\emph{(c)} We verify the three hypotheses under which the summation over source-anchored families
applies to $\mathfrak{L}_T$.

\emph{Reader condition \eqref{4.21:eq-site-geometry-a}.} Let $\Delta$ be a cell read by $F$. If
$\Delta$ meets $A_F$, then $\operatorname{pa}(\Delta)\subset[A_F]_i$. If
$\Delta\subset Y_p^{\sharp}$ with $(Y_p,m)\in\mathcal{P}_F$, then a cube of $A_F$ touches
$Y_p^{\sharp}$, and therefore, with $s_i(Y_p)$ the size of the pocket $Y_p$ in level-$i$ units
(the quantity entering Lemma~\ref{4.21:lem-probe-mass}\,(iii)),
\begin{equation*}
\operatorname{dist}_i\bigl(\operatorname{pa}(\Delta),[A_F]_i\bigr)\le s_i(Y_p)+1
\le 101\check{R}_i+1\le c_{\rho}\cdot 7\check{R}_i\le c_{\rho}\,\rho(\Delta),
\end{equation*}
by clauses (a) and (d) of the geometric lemma on pockets and strips and by the second condition of
\eqref{4.21:eq-site-geometry-b}.

\emph{Depth data \eqref{4.21:eq-wrapped-terms-a}.} Take as probes $\mathfrak{Q}$ the cells, with
depth $\varrho:=\rho$ and $N_{*}:=0$. The first depth condition is the first condition of
\eqref{4.21:eq-site-geometry-b}. For the second, let $\Delta\subset Y_p^{\sharp}$ with $(Y_p,m)$ a
live pocket; a wall-chain
from $\Delta$ to a source cell, which lies in $\Omega_{i+1}\subset\Omega_{m+1}$, must leave the
$\tfrac52\check{R}_m$-neighbourhood of $Y_p^{\sharp}$, and in this neighbourhood all cells have
sides at most $1$ in level-$m$ units by clause (c) of the geometric lemma on pockets and strips.
Hence $\rho(\Delta)\ge\tfrac52\check{R}_m$.

\emph{Per-cube bound.} Lemma~\ref{4.21:lem-wrapped-cube-bound} applies at the site of the
fluctuation integral with $r=(1-\tfrac12\delta)\kappa$ and gives the per-cube hypotheses of
\eqref{4.21:eq-wrapped-terms-a}. Since there are no hubs, a cell sees the reading set of a
reader if and only if it meets it. Lemma~\ref{4.21:lem-probe-mass}\,(iv) at $\beta=1$ therefore
bounds the reception functional:
\begin{equation}\label{4.21:eq-fluctuation-site-4}
\mathsf{T}^{*}_{(1-\frac12\delta)\kappa}[\mathfrak{L}_T]\le C_{\mathsf{a}}(1)\,\mathsf{b}_i(0).
\end{equation}

\emph{Distance.} Every label $\operatorname{lab}(F,\mathsf{y})$ holds a source cell, which
contains a point of $\Omega_{i+1}$, and a cube of $[A_F]_i$ with a point on a strip (the pocket
set $\mathcal{P}_F$ is non-empty by Definition~\ref{4.21:def-wrapped-terms}, and its strips touch
$A_F$). By clause (d) of the geometric lemma on pockets and strips, together with the counting
statement for the cut cells lying between a live strip and the region $\Omega_{i+1}$ made at this
step,
\begin{equation}\label{4.21:eq-fluctuation-site-5}
d_i\bigl(\operatorname{lab}(F,\mathsf{y})\bigr)\ge 7\check{R}_i-2 .
\end{equation}

\emph{Summation.} Apply clause (d) of Lemma~\ref{4.21:lem-source-anchored-sums} with
$r=\kappa_V$, $\varsigma=\tfrac14\delta\kappa$, and any $Q_{*}\subset Y$. Its conditions hold:
\begin{equation*}
r+\varsigma+2=(1-\tfrac12\delta)\kappa-\tfrac54\delta\kappa+2\le(1-\tfrac12\delta)\kappa,
\end{equation*}
as $\tfrac54\delta\kappa\ge2$; and $\kappa_1-1\ge c_l(r+\varsigma)+\lambda_{*}$ by the condition
on $\kappa_1$ in \eqref{4.21:eq-format-propagation-a}, since $r+\varsigma\le\kappa\le2\kappa+4$.
The clause gives
\begin{equation*}
|V^{\mathrm{rel}}(Y)|\le C_{\Sigma}\,
\mathsf{T}^{*}_{(1-\frac12\delta)\kappa}[\mathfrak{L}_T]\,
e^{-\varsigma d_i(Y)}\,e^{-\kappa_V d_i(Y)} .
\end{equation*}
If the sum in \eqref{4.21:eq-fluctuation-site-1} is not empty, $Y$ is a label
$\operatorname{lab}(F,\mathsf{y})$, so \eqref{4.21:eq-fluctuation-site-5} gives
$e^{-\varsigma d_i(Y)}\le e^{-\frac14\delta\kappa(7\check{R}_i-2)}$; together with
\eqref{4.21:eq-fluctuation-site-4} this yields
\begin{equation*}
|V^{\mathrm{rel}}(Y)|\le C_{\Sigma}\,C_{\mathsf{a}}(1)\,\mathsf{b}_i(0)\,
e^{-\frac14\delta\kappa(7\check{R}_i-2)}\,e^{-\kappa_V d_i(Y)}
=\varepsilon^{\mathrm{rel}}_i\,e^{-\kappa_V d_i(Y)},
\end{equation*}
which is \eqref{4.21:eq-fluctuation-site-3}.

\emph{(d)} The map $F\mapsto F_{\mathsf{y}}$ is linear. We attach $F_{\mathsf{y}}$ to the strong
set $\mathcal{S}'$ that $F$ carries in its format; the weighted clause then holds for
$V^{\mathrm{rel}}$ by condition (N) of that clause.

\emph{(e)} Evaluate \eqref{4.21:eq-fluctuation-site-a} at $\mathsf{T}:=\check{\mathsf{T}}_i$. By
the comparison of the coupling envelopes $\check{R}_i$, $\mathsf{T}_i$ with
$\check{\mathsf{T}}_i$ in the proof of the correlation theorem,
$\check{R}_i\ge\check{\mathsf{T}}_i^{r_{\mathrm{pr}}}$ and
$\mathsf{T}_i\le\check{\mathsf{T}}_i+1$. The per-cube constant $\mathsf{b}_i(0)$ of
Lemma~\ref{4.21:lem-wrapped-cube-bound} satisfies
\begin{equation*}
\mathsf{b}_i(0)\le b_0+C_N C_0^{\sharp}(\mathsf{T}_i+c_0)^{r_{\mathrm{pr}}}.
\end{equation*}
Since the bracket in \eqref{4.21:eq-fluctuation-site-a} is non-decreasing in $\mathsf{T}$ and the
exponential factor of $\varepsilon^{\mathrm{rel}}_i$ is decreasing in $\check{R}_i$, the
inequalities $\mathsf{T}_i\le\check{\mathsf{T}}_i+1$ and
$\check{R}_i\ge\check{\mathsf{T}}_i^{r_{\mathrm{pr}}}$ give
\begin{align*}
\varepsilon^{\mathrm{rel}}_i
&=C_{\Sigma}\,C_{\mathsf{a}}(1)\,\mathsf{b}_i(0)\,e^{-\frac14\delta\kappa(7\check{R}_i-2)}\\
&\le C_{\Sigma}\,C_{\mathsf{a}}(1)
\bigl[b_0+C_N C_0^{\sharp}(\check{\mathsf{T}}_i+1+c_0)^{r_{\mathrm{pr}}}\bigr]
e^{-\frac14\delta\kappa(7\check{\mathsf{T}}_i^{r_{\mathrm{pr}}}-2)}
\le\tfrac14\varepsilon_V^{*},
\end{align*}
the last inequality being \eqref{4.21:eq-fluctuation-site-a} at
$\mathsf{T}=\check{\mathsf{T}}_i$. In the step
of the proof of the complex-potential lemma in which $\varepsilon_V$ enters, replace
$\varepsilon_V$ by $\varepsilon_V+\varepsilon^{\mathrm{rel}}_i$. By (b), (c) and (d), the
complex-potential lemma then holds for $V=V^{\mathrm{pl}}+V^{\mathrm{rel}}$.

The second stage of the expansion does not see $V$, nor where $V$ reads the background. Indeed,
$V$ enters the fluctuation integral only through Mayer factors. The interpolation bound estimates
$\exp[\sum\tau V]$, with $\tau$ the interpolation parameters attached to the terms $V(Y)$, by
$\exp[\sum|\tau||V|]$, and after that no $V$ occurs. The steps of the proof
of the flat theorem on the fluctuation integral read $V$ only through the products
$\prod_{Y\in D}(e^{V(Y)}-1)$ over families $D$ of domains, through the bound
$|e^{V(Y)}-1|\le2\nu(Y)$ with the majorant $\nu(Y)$ of that proof, through holomorphy,
and through the property that $V(Y)$ reads $\mathsf{B}$ on $Y$. Conditioning, decoupling and
factorisation are performed over the integration variables $(\psi,x)$, and a strip holds none of
them by clause (d) of the geometric lemma on pockets and strips. The data
$(\mathbf{U},\mathbf{J},\mathbf{A}_{<i},w)$
on the strips enter as parameters under uniform bounds: background domains may overlap freely,
because they carry no integration variable. This proves that the statements (2), (3) and (4) hold
verbatim, with \emph{local in $X$} read as \emph{reads $((\mathbf{U},\mathbf{J}),\mathbf{A},w)$
on $X^{+}$}.

An output with domain $X$ reads only $V(Y)$ with $Y\subset X$, and the strips touching such a $Y$
touch $X$. The $\mathbf{T}$-step treats no component, so the pockets stay live. A domain $Y$
carrying $V^{\mathrm{rel}}$ holds a cube with a point on a strip, and the strip lies in
$\Omega_{i+1}^{c}\subset\Lambda_{i+1}^{c}$; hence $X\not\subset\Lambda_{i+1}$ for every output
domain $X\supset Y$. Thus no output with $X\subset\Lambda_{i+1}$ contains a $V^{\mathrm{rel}}$.

For the terms of level $i+1$ born in the $\mathbf{T}$-step, put $A:=X$ and let $\mathcal{P}$ be
the pockets of the constituents. Property (f1) of \eqref{4.21:eq-relative-format-a} over the
strips holds by the reading statement for the terms of the $\mathbf{T}$-step. The old terms
$F(\cdot,U_{i+1})$ stay stored, with their data unchanged.
\end{proof}

\begin{definition}[The preparation site: collar, cover and fluctuation cubes]\label{4.21:not-preparation-site}
Throughout this section we work inside the large-field operation $\mathbf{R}_k$, at its
stage $n+1$, and we write
\begin{equation}\label{4.21:eq-preparation-site-1}
j = h+n
\end{equation}
for the level at which this stage operates; here $h$ is the level of the first stage of
$\mathbf{R}_k$ (the case $n=0$), and the letter is unrelated to the histories $h$ of the glossary.
We call this stage, with the collar, the cover and the fluctuation cubes fixed below, the
\emph{preparation site}.

\emph{The collar and its treated parts.} The collar of stage $n+1$ is
\begin{equation}\label{4.21:eq-preparation-site-2}
R_n = Z_{j+1}\setminus Z''_{j+1}.
\end{equation}
For a component $\bar c$ treated at stage $n+1$ we put $R(\bar c) := R_n\cap\bar c$.

\emph{The cover and the block weights.} We write $\{\square\}$ for the cover of the
preparation site and $\{\zeta_\square\}$ for the partition of unity subordinate to it. To each
block $\square$ of the cover we attach the weight
\begin{equation}\label{4.21:eq-preparation-site-3}
\vartheta(\square) = \vartheta\, e^{-\frac12\delta_P\, d^{\wedge}(\square,R_n)},
\end{equation}
where $\vartheta$ is the block-weight constant of the preparation site (not the
separation-window parameter that the glossary denotes by the same letter), $\delta_P$ is the
decay rate of the
preparation site, and $d^{\wedge}$ is the scaled distance belonging to the domain sequence
$\mathbb{B}^{(n)}_k$, as fixed where that sequence is introduced.

\emph{The fluctuation cubes.} We write $\mathsf{Q}$ for the family of cubes of the partition
$\pi_{j+1}$ that meet the bonds of integration of stage $n+1$, and call its members
\emph{fluctuation cubes}; in this section $\mathsf{Q}$ denotes this family wherever fluctuation
cubes are meant, and it is neither the window set of cubes of the counting lemmas of this
chapter, for which the same letter is used, nor the symbol $\mathsf{Q}$ of the glossary
(Notation~\ref{0.2:not-glossary}). The bonds of integration are the bonds of the set
$(\Omega^c_{j+1}\setminus Z''_{j+1})^{(j)*}$, which carry the integration variables $B_j$,
and the bonds of the rim
\begin{equation}\label{4.21:eq-preparation-site-4}
Z_{j+1}\cap\Omega_{j+1}\cap Z ,
\end{equation}
which carry the integration variables $A_j$.

\emph{Further symbols.} The symbols
\begin{equation*}
\begin{gathered}
\mathcal{A}_c,\ \mathcal{Q}^{w}_c,\ \mathsf{b},\ \mathsf{d},\ \nu,\ \tilde{\nu},\ w(\cdot),\
\mathsf{p},\ V,\ K_A,\ K_Q,\ r_A,\ \varepsilon_S,\\
c_J,\ c_{\flat},\ \mu,\ \mathsf{M}_{*},\ \tilde{\mathfrak{l}},\ \sigma_{*},\ A_{\flat},\
\mathsf{r}
\end{gathered}
\end{equation*}
retain the meanings assigned to them where the perturbative expansion at the preparation site
and its bound are set up, and none of them carries in this section any other meaning that the
same letter has elsewhere in the book. Of these, $\mathcal{A}_c$ and $\mathcal{Q}^{w}_c$ are
attached to an element $c$ of that expansion; $\mathsf{b}$, $\mathsf{d}$, $\nu$ are the three
data of an element through which the proposition in item (p-ii) below reads it; and $w(\cdot)$
is a function on the elements. For an element $c$ of that expansion we put
\begin{equation}\label{4.21:eq-preparation-site-5}
\mathfrak{w}(c) := \mathfrak{N}^{A}(c)\, e^{w(c)} ,
\end{equation}
where $w(c)$ is the value of the function $w(\cdot)$ at the element $c$ (the letter denotes
neither the complex source nor the twist parameter).

\emph{What the estimates of this section take as given.} The estimate at the preparation site
is carried out afresh in this section; it is not quoted as a whole. Three statements are used
without being reproved:
\begin{enumerate}
\item[(p-i)] the lemma giving an identity and two Cauchy estimates for a function holomorphic
on the M\"obius polydisc;
\item[(p-ii)] the proposition on the expansion at the preparation site, which reads the
elements of the expansion only through the triple $(\mathsf{b},\mathsf{d},\nu)$;
\item[(p-iii)] the sums, over terms whose domain meets a given block, a given rim cube or a
given point, that are established in the counting for the expansion at the preparation site
and in the third step of the bound at that site, and these only at window anchors, in the
sense specified below in this section.
\end{enumerate}
The sum over the families $\mathsf{y}$ at fixed $(c,A,Q)$, with $c$ an element of the
expansion, is not among these ingredients; it
is carried out in the summation of the families by grain, later in this section.
\end{definition}

\begin{definition}[The mixed-level grain partition and fattened blocks]\label{4.21:def-grain-partition}
We work in the setting of Notation~\ref{4.21:not-preparation-site}: the operation
$\mathbf{R}_k$ is at stage $n+1$, at level $j=h+n$. We use the following objects, fixed with
the stages of $\mathbf{R}_k$.

\begin{itemize}
\item the label $\mathbb{B}^{(n)}_k$ read at stage $n+1$, and its cells;
\item the large-field domain $Z$, and the domain $\Omega_{j+1}$ made by the step $\mathbf{T}$,
which has not yet been re-labelled at this stage;
\item for each level $\ell$, the partition $\pi_\ell$ of the lattice into cubes of level $\ell$.
\end{itemize}

The \emph{mixed-level grain partition}, or grain of stage $n+1$, is the family $\pi^{\flat}$
consisting of the sets of the following three kinds:
\begin{enumerate}
\item[(a)] the cubes of $\pi_k$ not contained in $Z$;
\item[(b)] the cubes of $\pi_{j+1}$ contained in $Z\cap\Omega^{c}_{j+1}$;
\item[(c)] the cells of $\mathbb{B}^{(n)}_k$ contained in $Z\cap\Omega_{j+1}$.
\end{enumerate}

For $G\in\pi^{\flat}$ the level $\operatorname{lev}(G)$ is defined according to the kind of $G$:
\begin{equation}\label{4.21:eq-grain-partition-1}
\operatorname{lev}(G):=
\begin{cases}
k & \text{if $G$ is of kind (a)},\\
j+1 & \text{if $G$ is of kind (b)},\\
\text{the level of the cell $G$} & \text{if $G$ is of kind (c)}.
\end{cases}
\end{equation}

On $Z$ the family $\pi^{\flat}$ coincides with the decomposition of \cite{B-CMP122a}, which
states, for a domain $X$ decomposed along the zones $\Omega_m\setminus\Omega_{m+1}$, that
\begin{equation*}
X\cap\Omega^{c}_{j+2}\in\mathbf{D}_{j+1},
\qquad
X\cap(\Omega_m\setminus\Omega_{m+1})\in\mathbf{D}_m,
\quad m=j+2,\dots,k,
\end{equation*}
where $\Omega_m$, $m=j+2,\dots,k$, are the small-field domains at the levels $j+2,\dots,k$.
In particular, the cells of $\mathbb{B}^{(n)}_k$ contained in
$Z\cap(\Omega_{j+1}\setminus\Omega_{j+2})$ are also cubes of $\pi_{j+1}$.

Let $\square$ be a block of the cover of Notation~\ref{4.21:not-preparation-site}, and let
$\tilde{\square}^{4}\supset\square$ be the fattened block of $\square$.
The \emph{grain neighbourhood} of $\square$ is
\begin{equation}\label{4.21:eq-grain-partition-2}
\hat{\mathsf{K}}_{\square}
:=\bigcup\bigl\{G\in\pi^{\flat}:\ G\cap\tilde{\square}^{4}\neq\emptyset\bigr\}.
\end{equation}
For a finite set $A$ of blocks we put
\begin{equation}\label{4.21:eq-grain-partition-3}
\hat{\mathsf{K}}_A:=\bigcup_{\square\in A}\hat{\mathsf{K}}_{\square},
\qquad
\hat{\sigma}_0(A):=\bigl\{G\in\pi^{\flat}:\ G\not\subset\hat{\mathsf{K}}_A\bigr\}.
\end{equation}
We recall that $\sigma_0(A)$ denotes the family of cells of $\mathbb{B}^{(n)}_k$ that do not meet
$\bigcup_{\square\in A}\tilde{\square}^{4}$. Finally, we put
\begin{equation}\label{4.21:eq-grain-partition-4}
N^{\flat}:=(10L+2)^{d}.
\end{equation}

The grain $\pi^{\flat}$ is determined by the label $\mathbb{B}^{(n)}_k$, the domains $Z$
and $\Omega_{j+1}$, and the pair $(k,j)$ alone. It does not
depend on the source variables $z$, and its definition imposes no condition.
\end{definition}

\begin{lemma}[The grain partition: levels, neighbours and hulls]\label{4.21:lem-grain-properties}
Let the step $k$, the level $j$, the stage $n$, the large-field region $Z$, the domains $\Omega_\ell$,
the cell partition $\mathbb{B}^{(n)}_k$, the grain partition $\pi^{\flat}$ with its level function
$\operatorname{lev}$, the sets $\hat{\mathsf{K}}_{\square}$ and the number $N^{\flat}=(10L+2)^d$ be as in
Definition~\ref{4.21:def-grain-partition}. The cover and its blocks $\square$, with the zone and the
own cell of a block, the family $\mathsf{Q}$ of fluctuation cubes with its rim and the bond set $B_j$
carried by its cubes, and the regions $Z_{j+1}$ and $\Lambda_k$ are those of the present step, fixed
earlier in this chapter. The elements of $\pi^{\flat}$ are called its members. Throughout this lemma a
closed cube \emph{meets} a set $A$ if an interior point of the cube lies in $A$, and a set meets
$\bigcup\mathsf{Q}$ if some cube of $\mathsf{Q}$ meets it in this sense; it is in this sense that
members meet $\tilde{\square}^{4}$ in the definition of $\hat{\mathsf{K}}_{\square}$, that members
meet $\mathsf{D}$ in clause \textup{(5)}, and that hulls are formed.
\begin{enumerate}
\item[\textup{(1)}] $\pi^{\flat}$ is a partition of the torus into closed cubes, each of which is a
union of cells of $\mathbb{B}^{(n)}_k$.
\item[\textup{(2)}] Every member $G$ satisfies $j+1\le\operatorname{lev}(G)\le k$. Touching members
differ in level by at most one. A member touches at most $\hat{b}_t$ members, where $\hat{b}_t$ is the
constant of this chapter's counting lemma for cells, part (b).
\item[\textup{(3)}] Every cube $q\in\mathsf{Q}$ is a member of $\pi^{\flat}$.
\item[\textup{(4)}] Let $\square$ be any block of the cover, of any zone, whether it lies in $Z$ or
not. Then $\hat{\mathsf{K}}_{\square}\supset\tilde{\square}^{4}$, and $\hat{\mathsf{K}}_{\square}$ is
wall-connected: any two of its members are joined by a chain of its members in which consecutive
members meet in a wall, that is, in a $(d-1)$-dimensional face. Moreover $\hat{\mathsf{K}}_{\square}$
consists of at most $N^{\flat}$ members, and it lies within $6$ units of level $k$ of $\square$.
\item[\textup{(5)}] Let $\mathsf{D}$ be a connected closed set which meets
$\bigcup\mathsf{Q}$, put
\begin{equation*}
m:=m(\mathsf{D})=\min\{m'\ge j+1:\ \mathsf{D}\subset\Omega^c_{m'+1}\},
\end{equation*}
and let $\hat{\mathsf{D}}\supset\mathsf{D}$ be the union of $\mathsf{D}$ and of the members of
$\pi^{\flat}$ meeting $\mathsf{D}$. Then $m(\hat{\mathsf{D}})=m$, and
\begin{equation}\label{4.21:eq-grain-properties-1}
[\hat{\mathsf{D}}]_{m'}=[\mathsf{D}]_{m'}\qquad\text{for every } m'\ge m,
\end{equation}
where $[A]_{\ell}$ denotes the $\pi_\ell$-hull of a set $A$, the union of the cubes of $\pi_\ell$
meeting $A$.
\end{enumerate}
No quantity of this lemma, neither the constants $\hat{b}_t$, $N^{\flat}$ and $6$ nor the hulls of
clause \textup{(5)}, involves the region size parameters $\check{R}_\ell$ (in particular $\check{R}_k$),
the stage-count constant $\mathsf{N}_0$, the age of any large-field region, or the number of windows.
\end{lemma}

\begin{proof}
Following Definition~\ref{4.21:def-grain-partition}, we say that a member of $\pi^{\flat}$ is of the
first kind if it is a cube of $\pi_k$ not contained in $Z$, of the second kind if it is a cube of
$\pi_{j+1}$ contained in $Z\cap\Omega^c_{j+1}$, and of the third kind if it is a cell of
$\mathbb{B}^{(n)}_k$ contained in $Z\cap\Omega_{j+1}$; its level is $k$, $j+1$, or the level of the
cell, respectively. Here $\Omega_{j+1}$ is the domain made by the $\mathbf{T}$-step, not yet
re-labelled at the present stage. For $i\le k$ one unit of level $i$ is at most one unit of level $k$.

\emph{Clause (1).} By the geometric normalisation of the large-field region, $Z$ is a union of cubes
of $\pi_k$. By the first assertion of the structural statement on the $\mathbf{T}$-made domains,
$\Omega_{j+1}$ is a union of cubes of side $L^{-(k-j-1)}M\check{R}_{j+1}$, hence a union of cubes of
$\pi_{j+1}$. Consequently the members of the first kind tile the closure of the complement of $Z$,
those of the second kind tile $Z\cap\Omega^c_{j+1}$, and those of the third kind tile
$Z\cap\Omega_{j+1}$; together they form a partition of the torus into closed cubes. The cells of
$\mathbb{B}^{(n)}_k$ off $Z$ have level at most $k$, the cells in $\Omega^c_{j+1}$ have level at most
$j$, and the partitions $\pi_\ell$ are nested. Hence a cube of $\pi_k$ not contained in $Z$ is a union
of cells, and so is a cube of $\pi_{j+1}$ contained in $\Omega^c_{j+1}$; a member of the third kind is
itself a cell.

\emph{Clause (2).} Inside $Z$ the zones $j+1,\dots,k$, that is the sets $\Omega_\ell\setminus\Omega_{\ell+1}$
for $\ell=j+1,\dots,k$, are made by the $\mathbf{T}$-step: by the statement on arrested attempts, no
earlier arrested attempt re-labelled a level above its own level, and the present preparation has
re-labelled only the levels $\le j$. A member of the third kind lies in one of these zones, and its
level is its zone, so it lies between $j+1$ and $k$; members of the first and second kinds have levels
$k$ and $j+1$. This gives $j+1\le\operatorname{lev}(G)\le k$.

We compare the levels of touching members according to their kinds. Two members of the first kind
both have level $k$, and two members of the second kind both have level $j+1$. Since the zones
$j+1,\dots,k$ are $\mathbf{T}$-made, the adjacency property of the $\mathbf{T}$-made partition holds
among the members of the third kind, so touching members of the third kind differ in level by at
most one. A member of the second kind which touches $\Omega_{j+1}$ touches cells of zone $j+1$, which
have its own level $j+1$; this holds because zone $j+1$ is at least $8\check{R}_{j+2}$ units of level
$j+2$ thick. Let $C$ be a component of $Z$. The cells of $C$ within $2\check{R}_k$ of $\partial C$ lie
in $\Omega_k\setminus\Lambda_k$, by the property of the $\mathbf{T}$-made domains at the boundary of a
component and by the nesting of the zones; hence they have level $k$, the level of the members of the
first kind, and a member of the third kind touching a member of the first kind meets $\partial C$ and
so has level $k$. A member of the second kind does not touch a member of the first kind unless $j+1=k$,
in which case both have level $k$. In every case touching members differ in level by at most one.

The bound $\hat{b}_t$ on the number of members touching a given member is the one given by this
chapter's counting lemma for cells, part (b).

\emph{Clause (3).} The bonds of $B_j$ lie in $\Omega^c_{j+1}\cap Z$, so the cubes of $\mathsf{Q}$
carrying them are members of the second kind. The rim lies in $Z_{j+1}\cap\Omega_{j+1}$, which is
contained in $\Omega_{j+1}\setminus\Omega_{j+2}$ by the nesting of the zones; this is zone $j+1$,
whose cells are cubes of $\pi_{j+1}$, so the cubes of the rim are members of the third kind.

\emph{Clause (4).} Let $i$ be the zone of $\square$ and $\Delta_i$ its own cell. Every point of
$\tilde{\square}^{4}$ lies within $5$ units of level $i$ of $\Delta_i$. Since the members cover the
torus by clause (1), every point of $\tilde{\square}^{4}$ lies in a member which meets
$\tilde{\square}^{4}$; hence $\hat{\mathsf{K}}_{\square}\supset\tilde{\square}^{4}$.

\emph{The case $i\le j+1$ or $\Delta_i\not\subset Z$.} Every member meeting $\tilde{\square}^{4}$ has
side at least one unit of level $\max\{i,j+1\}$; off $Z$ a member has side one unit of level $k$,
which is at least that. If
$\Delta_i\not\subset Z$ and $\tilde{\square}^{4}$ reaches into $Z$, then $\Delta_i$ lies within $5$ of
its own units of $\partial Z$, so it lies in $\Lambda_k$, and $i=k$, hence $\max\{i,j+1\}=k$. Distinct
members meeting
$\tilde{\square}^{4}$ contain distinct cubes of $\pi_{\max\{i,j+1\}}$ meeting $\tilde{\square}^{4}$,
and there are at most $11^d$ such cubes. Thus $\hat{\mathsf{K}}_{\square}$ has at most $11^d\le
N^{\flat}$ members.

\emph{The case $i\ge j+2$ and $\Delta_i\subset Z$.} Zone $i-1\ge j+1$ is $\mathbf{T}$-made, and by
part (d) of the geometric lemma of this chapter it is at least $8\check{R}_i$ units of level $i$
thick, which exceeds $5$ units of level $i$. Hence $\tilde{\square}^{4}\cap Z\subset\Omega_{i-1}$, and
the members met by $\tilde{\square}^{4}$ have level at least $i-1$. There are at most $(10L+1)^d\le
N^{\flat}$ of them.

\emph{Wall-connectedness.} Two interior points of the cube $\tilde{\square}^{4}$ are joined inside it
by a path which misses every $(d-2)$-dimensional face of the partition $\pi^{\flat}$, which is locally
finite. Such a path passes from member to member through walls, and every member it visits meets
$\tilde{\square}^{4}$; hence $\hat{\mathsf{K}}_{\square}$ is wall-connected.

\emph{Distance.} A point of $\hat{\mathsf{K}}_{\square}$ lies in a member meeting
$\tilde{\square}^{4}$; it is therefore at distance at most five units of level $i$ plus the side of
one member from $\square$. Each of these lengths is at most one unit of level $k$ per unit, so
$\hat{\mathsf{K}}_{\square}$ lies within $6$ units of level $k$ of $\square$.

\emph{Clause (5).} Let $G$ be a member of $\pi^{\flat}$ meeting $\mathsf{D}$ and added to it in
$\hat{\mathsf{D}}$.

If $G$ is of the second kind, then $G\subset\Omega^c_{j+1}\subset\Omega^c_{m+1}$, because $m\ge j+1$.

If $G$ is of the third kind, a cell of zone $i'$, then $G$ has an interior point in $\mathsf{D}$,
which lies in $\Omega^c_{m+1}$, and $G$ lies in $\Omega_{i'}$; therefore $i'\le m$. Since $G$ lies in
zone $i'$, we get $G\subset\Omega^c_{i'+1}\subset\Omega^c_{m+1}$.

If $G$ is of the first kind, then $\mathsf{D}$, which is connected, has a point off $Z$ and a point in
a fluctuation cube of $\mathsf{Q}$. By the proof of clause (3), such a cube is a member of the second
or third kind and so is contained in $Z$. Hence $\mathsf{D}$ crosses $\partial Z$. On
both sides of $\partial Z$ the cells have level $k$, by the proof of clause (2), so $\mathsf{D}$ meets
$\Omega_k$. Consequently $m=k$ and $\Omega_{m+1}=\Omega_{k+1}=\emptyset$, and trivially
$G\subset\Omega^c_{m+1}$.

In every case $G\subset\Omega^c_{m+1}$; since also $\mathsf{D}\subset\Omega^c_{m+1}$, we have
$\hat{\mathsf{D}}\subset\Omega^c_{m+1}$ and $m(\hat{\mathsf{D}})\le m$. Conversely, if
$\hat{\mathsf{D}}\subset\Omega^c_{m'+1}$ then $\mathsf{D}\subset\Omega^c_{m'+1}$, because
$\mathsf{D}\subset\hat{\mathsf{D}}$, so $m(\hat{\mathsf{D}})\ge m$. Hence $m(\hat{\mathsf{D}})=m$.

In every case $\operatorname{lev}(G)\le m$: the level is $j+1\le m$ for the second kind, $i'\le m$
for the third kind, and $k=m$ for the first kind. Let $m'\ge m$. Then $\operatorname{lev}(G)\le m'$,
and since the partitions are nested, $G$ lies in one cube of $\pi_{m'}$; that cube meets $\mathsf{D}$,
because $G$ does. If a cube $P$ of $\pi_{m'}$ meets $\hat{\mathsf{D}}$, it meets $\mathsf{D}$ or some
added member $G$; in the latter case $P$ is the one cube of $\pi_{m'}$ containing $G$, which meets
$\mathsf{D}$. Hence $[\hat{\mathsf{D}}]_{m'}\subset[\mathsf{D}]_{m'}$, and the reverse inclusion holds
because $\mathsf{D}\subset\hat{\mathsf{D}}$; this proves \eqref{4.21:eq-grain-properties-1}.
\end{proof}

\begin{definition}[Functions reading the decoupling parameters near a block set]
\label{4.21:def-block-read-functions}
Work in the setting of Notation~\ref{4.21:not-preparation-site}. Let $\pi$ be the partition of the
site into cells, and let $K(\mathsf{s})$ be the decoupled operators, as in
Lemma~\ref{4.21:lem-exempted-set}. For a data point $b$, let $\mathcal{J}(\mathsf{s};b)$ be the current
of the layer of that lemma. Let $A$ be a finite non-empty set of blocks of the cover, and
let $\sigma_0(A)\subset\pi$ be the decoupling family of $A$. Put
\begin{equation}\label{4.21:eq-block-read-functions-1}
\mathfrak{P}_A:=\bigl\{\mathsf{s}:\ |\mathsf{s}(\Delta)|\le e^{\kappa_1}
\ \text{for all}\ \Delta\in\sigma_0(A)\bigr\},
\qquad
\mathsf{K}_A:=\bigcup\bigl(\pi\setminus\sigma_0(A)\bigr).
\end{equation}
The exempted set $\mathsf{K}_A$ contains $\bigcup_{\square\in A}\tilde{\square}^4$.

Let $f$ be a function on $\mathfrak{P}_A$ times the domains of its remaining parameters, and let
$\mathfrak{n}$ be a number. We say that $f$ \emph{reads the decoupling parameters near $A$ with bound
$\mathfrak{n}$} if both of the following hold.
\begin{enumerate}
\item[(i)] $f$ is holomorphic there and $|f|\le\mathfrak{n}$.
\item[(ii)] $f$ reads $\mathsf{s}$ only through objects built from the decoupled operators
$K(\mathsf{s})$, restricted to $\bigcup_{\square\in A}\tilde{\square}^4$, namely objects built from the
responses $\zeta_\square\eta\,\mathbb{H}^{(n)}_k(\mathsf{s};b)$ and
$\zeta_\square\mathcal{J}(\mathsf{s};b)$, and from the operators $D$ and $A_0$ of
\cite[(1.5)]{B-CMP116}.
\end{enumerate}

For a family $\mathsf{y}\subset\sigma_0(A)$ of cells, write $\partial_{\mathsf{y}}$ for
$\prod_{\Delta\in\mathsf{y}}\partial/\partial\mathsf{s}(\Delta)$, and write
$f|_{\mathsf{s}=0\text{ off }\mathsf{y}}$ for $f$ evaluated at $\mathsf{s}(\Delta)=0$ for all
$\Delta\notin\mathsf{y}$. Put
\begin{equation}\label{4.21:eq-block-read-functions-2}
f_{\mathsf{y}}:=\int_{[0,1]^{\mathsf{y}}}\partial_{\mathsf{y}}f\big|_{\mathsf{s}=0\text{ off }\mathsf{y}},
\end{equation}
where the integration runs over $(\mathsf{s}(\Delta))_{\Delta\in\mathsf{y}}\in[0,1]^{\mathsf{y}}$.

Such a function $f$ has the following locality property. Let $\mathsf{s}=0$ off $\mathsf{y}$, and write, in
the notation of Lemma~\ref{4.21:lem-exempted-set} for the subfamily $\sigma_0(A)$,
$Y(\sigma)=\mathsf{K}_A\cup\bigcup\mathsf{y}$. Beyond what it reads independently of $\mathsf{s}$,
$f$ reads $\mathsf{s}$, $b$ and $(\mathbb{U},\mathbb{J})$ only on those components of $Y(\sigma)$
that meet $\bigcup_{\square\in A}\tilde{\square}^4$.

The following functions read the decoupling parameters near $A$ with the bound $\mathfrak{n}$
indicated.
\begin{enumerate}
\item[(1)] Let $c$ be a stored term of the class with size functional $\mathfrak{N}^{A}$ of
Notation~\ref{4.21:not-preparation-site}, plain or wrapped. Let $Q$ be a finite set of cubes such that
$(A,Q)$ indexes a term of the identity for $c$ of the lemma giving an identity and two Cauchy
estimates for a function holomorphic on the M\"obius polydisc, and let $F(t,u;\mathsf{s})$, with
$t\in\mathbb{C}^{A}$ and $u\in\mathbb{C}^{Q}$, be the function of $c$ to which that identity is
applied. Write $\partial_{t_A}=\prod_{\square\in A}\partial/\partial t_\square$ and
$\partial_{u_Q}=\prod_{q\in Q}\partial/\partial u_q$. Take
\begin{equation}\label{4.21:eq-block-read-functions-3}
f:=c_{A,Q}(\mathsf{s})
=\int_{[0,1]^{A\sqcup Q}}\partial_{t_A}\partial_{u_Q}
F(t\mathbf{1}_A,u\mathbf{1}_Q;\mathsf{s}),
\end{equation}
where $t\mathbf{1}_A$ and $u\mathbf{1}_Q$ are the arguments of $F$ as in that lemma, with bound
\begin{equation}\label{4.21:eq-block-read-functions-4}
\mathfrak{n}:=\mathfrak{N}^{\natural}(c)\prod_{\square\in A}K_A\vartheta(\square)
\prod_{Q}(2/r_A).
\end{equation}
\item[(2)] The elements produced, with $Q=\emptyset$, by the identity of (1) from the far terms of
the two expansion families of the preparation site, with the bound
\eqref{4.21:eq-block-read-functions-4} for $Q=\emptyset$.
\item[(3)] Take $A=\{\square\}$. The members are the near groups, the far pieces of a term that has a
near cube, and the $\mathbb{P}$- and $W$-pieces of the preparation site. Their bound is
\begin{equation}\label{4.21:eq-block-read-functions-5}
\mathfrak{n}:=c_\alpha C_R(E_0+1)e^{C_2\kappa_1}\bigl(\vartheta(\square)+\mathfrak{p}_j\bigr),
\end{equation}
where $c_\alpha$, $C_R$, $E_0$, $C_2$ are the constants of the bounds for these pieces at the
preparation site.
\item[(4)] Let two machines $\mathsf{A}$, $\mathsf{B}$ act on one scaffold, and take the differences
$\delta c=c^{\mathsf{B}}-c^{\mathsf{A}}$ of the functions in (1)--(3). The bound is the same as in
(1)--(3), with $\sup|\delta c|$ in place of $\mathfrak{N}(c)$.
\end{enumerate}
\end{definition}

\begin{proof}
We first prove the locality property.

Let $\mathsf{y}\subset\sigma_0(A)$ and $\mathsf{s}=0$ off $\mathsf{y}$. Apply
Lemma~\ref{4.21:lem-exempted-set} with the subfamily $\sigma_0(A)$, so that its exempted set is
$\mathsf{K}_A$ and $Y(\sigma)=\mathsf{K}_A\cup\bigcup\mathsf{y}$. Part (ii) of that lemma then gives
two facts:
\begin{itemize}
\item on each component of $Y(\sigma)$, the objects of clause (ii), built from the decoupled
operators $K(\mathsf{s})$, read $(\mathsf{s},b,\mathbb{U},\mathbb{J})$ on that component only;
\item they vanish off $Y(\sigma)$.
\end{itemize}
By clause (ii) of the definition, $f$ reads $\mathsf{s}$ only through these objects, restricted to
$\bigcup_{\square\in A}\tilde{\square}^4$. This set is contained in $\mathsf{K}_A$, hence in
$Y(\sigma)$. The restriction therefore involves a component of $Y(\sigma)$ only if that component
meets $\bigcup_{\square\in A}\tilde{\square}^4$. This gives the locality property.

Next we treat the members of the class. For (1) and (2), the function $F$ of the stored term $c$
depends on $\mathsf{s}$ only through the responses and the operators $D$, $A_0$ of clause (ii),
restricted to $\bigcup_{\square\in A}\tilde{\square}^4$, by the construction of the stored terms in
the setting of Notation~\ref{4.21:not-preparation-site}; differentiation and integration in $(t,u)$
do not involve $\mathsf{s}$, so clause (ii) holds for $c_{A,Q}$. For (3), the same holds for the near
groups, the far pieces and the $\mathbb{P}$- and $W$-pieces by their construction at the preparation
site.

For (1), for every $\mathsf{s}\in\mathfrak{P}_A$ the function $F(\cdot,\cdot;\mathsf{s})$ is
holomorphic on the product polydisc in $(t,u)$, that is, on the product of the discs in the
variables $t_\square$, $\square\in A$, and $u_q$, $q\in Q$, of the lemma named in (1); moreover $F$
is holomorphic jointly in $(t,u,\mathsf{s})$ on that product polydisc times $\mathfrak{P}_A$. Both
facts are shown in the proof of part (ii) of the lemma named in (1).
\begin{itemize}
\item By the joint holomorphy, $\partial_{t_A}\partial_{u_Q}F$ is holomorphic in $\mathsf{s}$, and
integration over the compact set $[0,1]^{A\sqcup Q}$ preserves holomorphy in $\mathsf{s}$. Hence
$c_{A,Q}$ is holomorphic on $\mathfrak{P}_A$.
\item At each fixed $\mathsf{s}\in\mathfrak{P}_A$, the Cauchy estimates in $(t,u)$ of that proof bound
$c_{A,Q}$ by \eqref{4.21:eq-block-read-functions-4}. Since the holomorphy on the product polydisc holds
for every $\mathsf{s}\in\mathfrak{P}_A$, the bound is uniform on $\mathfrak{P}_A$.
\end{itemize}

For (2), the same applies with $Q=\emptyset$.

For (3), these pieces are holomorphic on $\mathfrak{P}_{\{\square\}}$, and they obey
\eqref{4.21:eq-block-read-functions-5} by the bounds for the near groups, the far pieces and the
$\mathbb{P}$- and $W$-pieces at the preparation site.

For (4), the difference of two holomorphic functions is holomorphic, and it reads $\mathsf{s}$ only
through objects read by the two functions whose difference it is, so clause (ii) holds for it. The
bound is the one of (1)--(3) for the two machines on one scaffold, by the estimates for the two
machines at the preparation site, with $\sup|\delta c|$ in place of $\mathfrak{N}(c)$.
\end{proof}

\begin{lemma}[Resummation of the decoupling expansion by grains]\label{4.21:lem-grain-resummation}
Let $A$ be a finite non-empty set of blocks, and let $f$ be a function of the class of
Definition~\ref{4.21:def-block-read-functions} with bound $\mathfrak{n}$ at $A$. In particular, $f$ is
defined on the polydisc $\mathfrak{P}_A$ in the variables $\mathsf{s}(\Delta)$,
$\Delta\in\sigma_0(A)$. Let $\pi^{\flat}$, $\hat{\mathsf{K}}_A$, $\hat{\mathsf{K}}_\square$ (for a
block $\square$), $N^{\flat}$, $\sigma_0(A)$ and $\hat\sigma_0(A)$ be as in
Definition~\ref{4.21:def-grain-partition}, and let $\hat b_t$ be the bound, in
Lemma~\ref{4.21:lem-grain-properties}, on the number of members of $\pi^{\flat}$ touched by one
member. Wall-connectedness and wall-components of families of members of $\pi^{\flat}$ are meant as
in Lemma~\ref{4.21:lem-grain-properties}.

For $s\in\mathbb{C}^{\hat\sigma_0(A)}$ define $\mathsf{s}_s\in\mathbb{C}^{\sigma_0(A)}$ by
\begin{equation}\label{4.21:eq-grain-resummation-1}
\mathsf{s}_s(\Delta):=
\begin{cases}
s(G) & \text{if } \Delta\subset G\in\hat\sigma_0(A),\\
1 & \text{if } \Delta\subset\hat{\mathsf{K}}_A,
\end{cases}
\qquad \Delta\in\sigma_0(A),
\end{equation}
and put $\hat f(s):=f(\mathsf{s}_s)$. For $\mathsf{y}\subset\sigma_0(A)$ let, as in
Definition~\ref{4.21:def-block-read-functions},
\[
f_{\mathsf{y}}:=\int_{[0,1]^{\mathsf{y}}}\partial_{\mathsf{y}}f\Big|_{\mathsf{s}=0\text{ off }\mathsf{y}} .
\]
For $\hat{\mathsf{y}}\subset\hat\sigma_0(A)$ put
\begin{equation}\label{4.21:eq-grain-resummation-2}
f_{\hat{\mathsf{y}}}:=\sum\Big\{f_{\mathsf{y}}:\ \mathsf{y}\subset\sigma_0(A),\
\{G\in\hat\sigma_0(A): G\text{ contains a cell of }\mathsf{y}\}=\hat{\mathsf{y}}\Big\}.
\end{equation}
The set of families $\mathsf{y}$ occurring in \eqref{4.21:eq-grain-resummation-2} is called the
\emph{fibre} of $\hat{\mathsf{y}}$. The cells of $\mathsf{y}$ inside $\hat{\mathsf{K}}_A$ are
unrestricted in it. Then the following hold.
\begin{enumerate}
\item[(i)] (Identity.) For every $\hat{\mathsf{y}}\subset\hat\sigma_0(A)$,
\begin{equation}\label{4.21:eq-grain-resummation-3}
f_{\hat{\mathsf{y}}}=\int_{[0,1]^{\hat{\mathsf{y}}}}\prod_{G\in\hat{\mathsf{y}}}ds(G)\,
\partial_{s(G)}\hat f(s)\Big|_{s=0\text{ off }\hat{\mathsf{y}}} .
\end{equation}
Moreover,
\begin{equation}\label{4.21:eq-grain-resummation-4}
\sum_{\hat{\mathsf{y}}}f_{\hat{\mathsf{y}}}=\sum_{\mathsf{y}}f_{\mathsf{y}}=f(\mathbf{1}).
\end{equation}
\item[(ii)] (Bound.) We have
\begin{equation}\label{4.21:eq-grain-resummation-5}
|f_{\hat{\mathsf{y}}}|\le\mathfrak{n}\,e^{-(\kappa_1-1)|\hat{\mathsf{y}}|}.
\end{equation}
Here $|\hat{\mathsf{y}}|$ is the number of members of $\pi^{\flat}$ in $\hat{\mathsf{y}}$, however
many cells they contain. Further, $f_{\hat{\mathsf{y}}}$ is linear in $f$, and it is holomorphic,
uniformly, in every parameter in which $f$ is holomorphic.
\item[(iii)] (Locality and vanishing.) Beyond what $f$ reads independently of $\mathsf{s}$,
$f_{\hat{\mathsf{y}}}$ reads the datum $b$ and the pair $(\mathbb{U},\mathbb{J})$ only on
$\hat{\mathsf{K}}_A\cup\bigcup\hat{\mathsf{y}}$. It vanishes unless every wall-component of
$\hat{\mathsf{y}}$ shares a wall with $\hat{\mathsf{K}}_A$.
\item[(iv)] (Entropy.) Let $\mathsf{x}>0$ satisfy
$2e^2\hat b_t^{\,2}N^{\flat}\mathsf{x}\le 1$. Call $\hat{\mathsf{y}}$ \emph{admissible} if every
wall-component of $\hat{\mathsf{y}}$ shares a wall with $\hat{\mathsf{K}}_A$. Then
\begin{equation}\label{4.21:eq-grain-resummation-6}
\sum_{\hat{\mathsf{y}}\ \text{admissible}}(e\mathsf{x})^{|\hat{\mathsf{y}}|}\le e^{|A|}.
\end{equation}
Let $\mathsf{E}$ be a closed set, let $\square_0\in A$, and write $\operatorname{dist}_k$ for
distance measured in level-$k$ units. Restrict the sum to those admissible $\hat{\mathsf{y}}$ one
of whose wall-components shares a wall with $\hat{\mathsf{K}}_{\square_0}$ and meets or touches
$\mathsf{E}$, and take the weight $\mathsf{x}^{|\hat{\mathsf{y}}|}$ in place of
$(e\mathsf{x})^{|\hat{\mathsf{y}}|}$. The restricted sum is at most
\[
e^{-(\operatorname{dist}_k(\mathsf{E},\hat{\mathsf{K}}_{\square_0})-1)_+}\,e^{|A|}.
\]
\item[(v)] (The same functions.) Let $\mathfrak{E}$ be the set of elements of the polymer expansion
of the preparation step. An element has the form $e=(\cdot,A,\cdot,\mathsf{y})$, with $A$ a finite
set of blocks and $\mathsf{y}\subset\sigma_0(A)$. The symbols $c$, $Q$ and $K_c$ below denote the
term, the cube set and the region attached to $e$ in the preparation step; they are attached to
the entries of $e$ other than $\mathsf{y}$, which the fibre map below leaves unchanged. Let $v_e$ be
the activity of $e$, and let $\mathsf{b}(e)$ and $\mathsf{d}(e)$ be its cube set and its domain.

Let $e\mapsto\hat e:=(\cdot,A,\cdot,\hat{\mathsf{y}})$ be the fibre map, where $\hat{\mathsf{y}}$ is
the set of members of $\hat\sigma_0(A)$ containing a cell of $\mathsf{y}$. Let
$\hat{\mathfrak{E}}$ be its image, and put $v_{\hat e}:=\sum_{e\mapsto\hat e}v_e$ and
\begin{align*}
\mathsf{b}(\hat e)&:=\big\{q\in\mathsf{Q}:\ q\subset\hat{\mathsf{K}}_A\cup\textstyle\bigcup
\hat{\mathsf{y}}\big\}\cup Q\cup\mathcal{S}_j(c),\\
\mathsf{d}(\hat e)&:=K_c\cup\hat{\mathsf{K}}_A\cup\textstyle\bigcup\hat{\mathsf{y}} .
\end{align*}
Then for every $e\mapsto\hat e$ with $\mathsf{b}(e)\neq\emptyset$,
\begin{equation}\label{4.21:eq-grain-resummation-7}
\mathsf{b}(e)=\mathsf{b}(\hat e),\qquad m(\mathsf{d}(e))=m(\mathsf{d}(\hat e)),\qquad
[\mathsf{d}(e)]_{m'}=[\mathsf{d}(\hat e)]_{m'}\quad(m'\ge m(\mathsf{d}(e))),
\end{equation}
with $m(\cdot)$ and $[\cdot]_{m'}$ as in Lemma~\ref{4.21:lem-grain-properties}. The exponent is
unchanged:
\[
\sum_{e}v_e=\sum_{\hat e}v_{\hat e}.
\]
For every $n\ge1$ and every hull $X$, let $a_n(X)$ be the finite sum of all cluster terms with $n$
polymers and hull $X$. It is the same number whether the polymers are built on $\mathfrak{E}$ or
on $\hat{\mathfrak{E}}$. Hence $\mathbb{C}_I(X)=\sum_n a_n(X)$ is one and the same function in the
two constructions, and the convergence proposition for the polymer expansion of the preparation
step, run on $\hat{\mathfrak{E}}$, bounds it.
\end{enumerate}
\end{lemma}

\begin{proof}
Throughout, write $T:=\sigma_0(A)$; it is a finite set.

\emph{Decomposition of $T$.} By part~($\flat1$) of Lemma~\ref{4.21:lem-grain-properties},
$\pi^{\flat}$ is a partition of the torus into closed cubes, each a union of cells. So every cell
$\Delta\in T$ lies in exactly one member $G$ of $\pi^{\flat}$. Either
$G\subset\hat{\mathsf{K}}_A$, or $G\in\hat\sigma_0(A)$.

Moreover, every cell of a member $G\in\hat\sigma_0(A)$ belongs to $T$. Indeed, a cell of $G$
meeting $\tilde\square^4$ for some $\square\in A$ would make $G$ meet $\tilde\square^4$. Then
$G\subset\hat{\mathsf{K}}_\square\subset\hat{\mathsf{K}}_A$, contrary to $G\in\hat\sigma_0(A)$.

Hence
\[
T=T_0\sqcup\bigsqcup_{G\in\hat\sigma_0(A)}T_G ,
\]
where $T_0$ is the set of cells of $T$ contained in $\hat{\mathsf{K}}_A$, and $T_G$ is the
(non-empty) set of all cells of $G$. In particular \eqref{4.21:eq-grain-resummation-1} defines
$\mathsf{s}_s(\Delta)$ for every $\Delta\in T$, and without ambiguity.

\emph{(i).} Since $e^{\kappa_1}\ge1$, the cube $[0,1]^T$ lies in $\mathfrak{P}_A$, so $f$ is smooth
on it. On functions of $\mathsf{s}\in[0,1]^T$, let $\varepsilon^1_\Delta$ and
$\varepsilon^0_\Delta$ denote evaluation at $\mathsf{s}(\Delta)=1$ and at $\mathsf{s}(\Delta)=0$
respectively. Evaluations in distinct variables commute. By the fundamental theorem of calculus in
the variable $\mathsf{s}(\Delta)$, for every function $\psi$ of $\mathsf{s}$ smooth on $[0,1]^T$,
\[
\int_0^1\partial_{\mathsf{s}(\Delta)}\psi\,d\mathsf{s}(\Delta)
=(\varepsilon^1_\Delta-\varepsilon^0_\Delta)\psi .
\]
Applying this in each variable of $\mathsf{y}$ gives
\begin{equation}\label{4.21:eq-grain-resummation-8}
f_{\mathsf{y}}=\prod_{\Delta\in\mathsf{y}}(\varepsilon^1_\Delta-\varepsilon^0_\Delta)
\prod_{\Delta\in T\setminus\mathsf{y}}\varepsilon^0_\Delta\,f .
\end{equation}

By the decomposition of $T$, a family $\mathsf{y}\subset T$ lies in the fibre of $\hat{\mathsf{y}}$
if and only if
\[
\mathsf{y}=\mathsf{y}_0\sqcup\bigsqcup_{G\in\hat{\mathsf{y}}}\mathsf{y}_G ,
\]
where $\mathsf{y}_0\subset T_0$ is arbitrary, $\mathsf{y}_G\subset T_G$ is non-empty for
$G\in\hat{\mathsf{y}}$, and $\mathsf{y}\cap T_G=\emptyset$ for $G\notin\hat{\mathsf{y}}$. The
operator in \eqref{4.21:eq-grain-resummation-8} is a product of commuting factors, one group for
$T_0$ and one for each $T_G$. So the sum over the fibre factorises.

Over all subsets of $T_0$ we obtain
\[
\prod_{\Delta\in T_0}\big[(\varepsilon^1_\Delta-\varepsilon^0_\Delta)+\varepsilon^0_\Delta\big]
=\prod_{\Delta\in T_0}\varepsilon^1_\Delta .
\]
For $G\in\hat{\mathsf{y}}$, over the non-empty subsets of $T_G$ we obtain
\[
\prod_{\Delta\in T_G}\big[(\varepsilon^1_\Delta-\varepsilon^0_\Delta)+\varepsilon^0_\Delta\big]
-\prod_{\Delta\in T_G}\varepsilon^0_\Delta=\varepsilon^{1}_G-\varepsilon^0_G .
\]
Here $\varepsilon^{1}_G:=\prod_{\Delta\in T_G}\varepsilon^1_\Delta$ sets all cells of $G$
to $1$, and $\varepsilon^0_G:=\prod_{\Delta\in T_G}\varepsilon^0_\Delta$ sets all of them to $0$.
For $G\notin\hat{\mathsf{y}}$ the factor is $\varepsilon^0_G$. Therefore
\[
f_{\hat{\mathsf{y}}}=\prod_{G\in\hat{\mathsf{y}}}(\varepsilon^{1}_G-\varepsilon^0_G)
\prod_{G\in\hat\sigma_0(A)\setminus\hat{\mathsf{y}}}\varepsilon^0_G
\prod_{\Delta\in T_0}\varepsilon^1_\Delta\,f .
\]

For $a\in\{0,1\}^{\hat\sigma_0(A)}$, the operator
$\prod_{T_0}\varepsilon^1_\Delta\prod_G\varepsilon^{a(G)}_G$ sends $f$ to $f(\mathsf{s}_a)$, that
is, to $\hat f(a)$, by \eqref{4.21:eq-grain-resummation-1}. Hence the last display is the same
expression for $\hat f$ on $[0,1]^{\hat\sigma_0(A)}$: with $\hat\varepsilon^{1}_G$ and
$\hat\varepsilon^{0}_G$ the evaluations at $s(G)=1$ and at $s(G)=0$,
\[
f_{\hat{\mathsf{y}}}=\prod_{G\in\hat{\mathsf{y}}}(\hat\varepsilon^{1}_G-\hat\varepsilon^{0}_G)
\prod_{G\notin\hat{\mathsf{y}}}\hat\varepsilon^{0}_G\,\hat f .
\]
The fundamental theorem of calculus in each variable $s(G)$, $G\in\hat{\mathsf{y}}$, turns this
into \eqref{4.21:eq-grain-resummation-3}.

The fibres of the subsets $\hat{\mathsf{y}}\subset\hat\sigma_0(A)$ partition the set of all
subsets of $T$, because each $\mathsf{y}$ determines its $\hat{\mathsf{y}}$. Hence
$\sum_{\hat{\mathsf{y}}}f_{\hat{\mathsf{y}}}=\sum_{\mathsf{y}}f_{\mathsf{y}}$. By
\eqref{4.21:eq-grain-resummation-8},
\[
\sum_{\mathsf{y}}f_{\mathsf{y}}
=\prod_{\Delta\in T}\big[(\varepsilon^1_\Delta-\varepsilon^0_\Delta)+\varepsilon^0_\Delta\big]f
=\prod_{\Delta\in T}\varepsilon^1_\Delta f=f(\mathbf{1}),
\]
which is \eqref{4.21:eq-grain-resummation-4}. No term is dropped or added. The sums run over all
subsets, including the families $\mathsf{y}$ whose contribution $f_{\mathsf{y}}$ vanishes.

\emph{(ii).} Each coordinate of $\mathsf{s}_s$ is either a coordinate of $s$ or the constant $1$.
So $s\mapsto\mathsf{s}_s$ is affine. Since $1\le e^{\kappa_1}$, it maps the polydisc
$\{|s(G)|\le e^{\kappa_1}: G\in\hat\sigma_0(A)\}$ into $\mathfrak{P}_A$. Therefore $\hat f$ is
holomorphic on that polydisc and satisfies $|\hat f|\le\mathfrak{n}$ there. This is the
restriction of the first condition of Definition~\ref{4.21:def-block-read-functions} (holomorphy on
$\mathfrak{P}_A$ with $|f|\le\mathfrak{n}$).

No new estimate enters. The localised response bounds and the analyticity bounds on
$\mathfrak{P}_A$, from which the first condition is obtained for the members of the class, are
used as they stand, on a diagonal of their own polydisc. On that diagonal the weight of a walk
becomes a power of $s(G)$, and neither Cauchy's inequality nor the fundamental theorem of calculus
depends on the degree. The bound on the polydisc is uniform in the number of cells. This is
because the number $m(\omega)$ of cells met by a walk $\omega$ is paid for by its length
$d(\omega)$, by \cite[(1.11)]{B-CMP116} and the locality bounds of the walk expansions, which hold
with the same constants by Lemma~\ref{4.21:lem-exempted-set}.

Let $s$ satisfy $s(G)\in[0,1]$ for $G\in\hat{\mathsf{y}}$ and $s=0$ off $\hat{\mathsf{y}}$. For
$G\in\hat{\mathsf{y}}$ the circle of radius $e^{\kappa_1}-1$ about $s(G)$ lies in the disc of
radius $e^{\kappa_1}$ about $0$, since $|s(G)|\le1$. The inequality
$e^{\kappa_1}-1\ge e^{\kappa_1-1}$ used in the second step of the next display is equivalent to
$e^{\kappa_1}(1-e^{-1})\ge1$, that is, to $\kappa_1\ge\log\frac{e}{e-1}$. Cauchy's inequality for
the mixed first derivative in the $|\hat{\mathsf{y}}|$ variables $s(G)$, $G\in\hat{\mathsf{y}}$, on
the polydisc of these radii around $s$ gives
\[
\Big|\prod_{G\in\hat{\mathsf{y}}}\partial_{s(G)}\hat f(s)\Big|
\le\mathfrak{n}\,(e^{\kappa_1}-1)^{-|\hat{\mathsf{y}}|}
\le\mathfrak{n}\,e^{-(\kappa_1-1)|\hat{\mathsf{y}}|} .
\]
Integrating over $[0,1]^{\hat{\mathsf{y}}}$, a set of measure one, and using
\eqref{4.21:eq-grain-resummation-3} gives \eqref{4.21:eq-grain-resummation-5}. Only the number of
variables $s(G)$, that is, of members of $\pi^{\flat}$ in $\hat{\mathsf{y}}$, enters.

The map $f\mapsto\hat f$ and the operations in \eqref{4.21:eq-grain-resummation-3} are linear, so
$f_{\hat{\mathsf{y}}}$ is linear in $f$. If $f$ is holomorphic in a further parameter, so is
$\hat f$. So are the derivatives of $\hat f$ used above, which by Cauchy's formula are integrals of
$\hat f$ over products of circles. The integrand of \eqref{4.21:eq-grain-resummation-3} is bounded
by $\mathfrak{n}e^{-(\kappa_1-1)|\hat{\mathsf{y}}|}$ uniformly in that parameter, since only
$|f|\le\mathfrak{n}$ was used. Hence the integral is holomorphic in it, with the same bound.

\emph{(iii).} Let $s=0$ off $\hat{\mathsf{y}}$. By \eqref{4.21:eq-grain-resummation-1},
$\mathsf{s}_s=0$ on every cell of $T$ not contained in $\hat{\mathsf{K}}_A\cup\bigcup\hat{\mathsf{y}}$.
Write $\mathsf{K}_A$ for the union of the cells meeting $\bigcup_A\tilde\square^4$. Then
$\mathsf{K}_A\subset\hat{\mathsf{K}}_A$, since a cell meeting $\tilde\square^4$ lies in a member of
$\pi^{\flat}$ meeting $\tilde\square^4$. With $Y(\cdot)$ as in Lemma~\ref{4.21:lem-exempted-set},
this gives
\[
Y(\mathsf{s}_s)\subset\hat{\mathsf{K}}_A\cup\textstyle\bigcup\hat{\mathsf{y}} .
\]
By the second condition of Definition~\ref{4.21:def-block-read-functions} ($f$ reads $\mathsf{s}$
only through the decoupled operators restricted to $\bigcup_A\tilde\square^4$) and part~(ii) of
Lemma~\ref{4.21:lem-exempted-set}, $f$ at $\mathsf{s}_s$ reads $\mathsf{s}$, $b$ and
$(\mathbb{U},\mathbb{J})$, beyond its $\mathsf{s}$-independent reading, only on the components of
$Y(\mathsf{s}_s)$ that meet $\bigcup_A\tilde\square^4$. The derivatives in
\eqref{4.21:eq-grain-resummation-3} are taken in variables $s(G)$ with $G\in\hat{\mathsf{y}}$, and
so stay in the slice $\{s=0\text{ off }\hat{\mathsf{y}}\}$. The locality statement follows.

For the vanishing, consider a wall-component of $\hat{\mathsf{K}}_A\cup\bigcup\hat{\mathsf{y}}$. By
part~($\flat4$) of Lemma~\ref{4.21:lem-grain-properties}, each $\hat{\mathsf{K}}_\square$,
$\square\in A$, is wall-connected and contains $\tilde\square^4$. So the component either contains
some $\hat{\mathsf{K}}_\square$, or it is a wall-component $\hat{\mathsf{y}}_a$ of
$\hat{\mathsf{y}}$ sharing no wall with $\hat{\mathsf{K}}_A$.

In the second case, $\hat{\mathsf{y}}_a$ does not meet $\bigcup_A\tilde\square^4$. Indeed, a member
of $\hat\sigma_0(A)$ meeting some $\tilde\square^4$ would be contained in $\hat{\mathsf{K}}_A$.
Hence, on the slice $\{s=0\text{ off }\hat{\mathsf{y}}\}$, $\hat f$ does not depend on the
restriction $s\restriction\hat{\mathsf{y}}_a$. So the derivative in
\eqref{4.21:eq-grain-resummation-3}, which contains $\partial_{s(G)}$ for every
$G\in\hat{\mathsf{y}}_a$, vanishes. This is the mechanism of \cite{B-CMP116}: derivatives with
respect to the parameters restricted to such components make the term equal to zero.

\emph{(iv).} An admissible $\hat{\mathsf{y}}$ is the disjoint union of its wall-components. These
form a set of pairwise distinct wall-connected families of members of $\pi^{\flat}$, which we call
animals. Each animal $\mathsf{h}$ contains a member $G_0$ sharing a wall with $\hat{\mathsf{K}}_A$,
and $|\hat{\mathsf{y}}|=\sum|\mathsf{h}|$. The map from $\hat{\mathsf{y}}$ to its set of animals
is injective. Hence
\[
\sum_{\hat{\mathsf{y}}\ \text{admissible}}(e\mathsf{x})^{|\hat{\mathsf{y}}|}
\le\prod_{\mathsf{h}}\big(1+(e\mathsf{x})^{|\mathsf{h}|}\big)
\le\exp\Big[\sum_{\mathsf{h}}(e\mathsf{x})^{|\mathsf{h}|}\Big],
\]
the product and the sum running over the animals that contain a member sharing a wall with
$\hat{\mathsf{K}}_A$.

By part~($\flat4$) of Lemma~\ref{4.21:lem-grain-properties}, $\hat{\mathsf{K}}_A$ consists of at
most $N^{\flat}|A|$ members. By part~($\flat2$), each of them touches at most $\hat b_t$ members. A
member $G_0$ sharing a wall with $\hat{\mathsf{K}}_A$ touches a member of $\hat{\mathsf{K}}_A$, so
there are at most $\hat b_tN^{\flat}|A|$ choices of $G_0$. By part~(b) of the cell-counting lemma,
the wall-connected families of $n$ members of $\pi^{\flat}$ containing a given member $G_0$ number
at most $(e\hat b_t)^n$, a member touching at most $\hat b_t$ members by part~($\flat2$). Counting
each animal through one of its members of this kind,
\[
\sum_{\mathsf{h}}(e\mathsf{x})^{|\mathsf{h}|}
\le\hat b_tN^{\flat}|A|\sum_{n\ge1}(e^2\hat b_t\mathsf{x})^n .
\]
Put $y:=e^2\hat b_t\mathsf{x}$. The hypothesis gives $y\le(2\hat b_tN^{\flat})^{-1}\le\tfrac12$,
since $\hat b_t\ge1$ and $N^{\flat}=(10L+2)^d\ge1$. Hence $\sum_{n\ge1}y^n=y/(1-y)\le 2y$, and
\[
\sum_{\hat{\mathsf{y}}\ \text{admissible}}(e\mathsf{x})^{|\hat{\mathsf{y}}|}
\le\exp\big[2e^2\hat b_t^{\,2}N^{\flat}\mathsf{x}\,|A|\big]\le e^{|A|},
\]
which is \eqref{4.21:eq-grain-resummation-6}.

For the second clause, let $\hat{\mathsf{y}}$ be one of the restricted families. One of its
wall-components is a wall-connected family joining $\hat{\mathsf{K}}_{\square_0}$ to $\mathsf{E}$.
Each of its members is of side at most one level-$k$ unit, so it has at least
$\operatorname{dist}_k(\mathsf{E},\hat{\mathsf{K}}_{\square_0})$ members. Hence
\[
|\hat{\mathsf{y}}|\ge N_{\mathsf{E}}
:=(\operatorname{dist}_k(\mathsf{E},\hat{\mathsf{K}}_{\square_0})-1)_+ ,
\]
and
\[
\mathsf{x}^{|\hat{\mathsf{y}}|}=e^{-|\hat{\mathsf{y}}|}(e\mathsf{x})^{|\hat{\mathsf{y}}|}
\le e^{-N_{\mathsf{E}}}(e\mathsf{x})^{|\hat{\mathsf{y}}|} .
\]
Summing over the restricted admissible families and applying \eqref{4.21:eq-grain-resummation-6}
gives the bound $e^{-N_{\mathsf{E}}}e^{|A|}$.

\emph{(v), cube sets and domains.} By part~($\flat3$) of Lemma~\ref{4.21:lem-grain-properties},
every $q\in\mathsf{Q}$ is a member of $\pi^{\flat}$. So $q$ meets $\bigcup_A\tilde\square^4$ if and
only if $q\subset\hat{\mathsf{K}}_A$. Also, $q$ meets $\bigcup\mathsf{y}$ outside
$\hat{\mathsf{K}}_A$ if and only if $q\in\hat{\mathsf{y}}$. Comparing with the cube set
$\mathsf{b}(e)$ of the element in the preparation step gives $\mathsf{b}(e)=\mathsf{b}(\hat e)$.
The bonds within distance $M_1$ of the boundary of $Y(\mathsf{s})$ do not enter, since they are not
read \cite{B-CMP116}.

The set $\mathsf{d}(\hat e)$ is obtained from $\mathsf{d}(e)$ by adjoining members of $\pi^{\flat}$
meeting $\mathsf{d}(e)$. Let $\hat{\mathsf{D}}$ be the union of $\mathsf{d}(e)$ and all members of
$\pi^{\flat}$ meeting it. Then
\[
\mathsf{d}(e)\subset\mathsf{d}(\hat e)\subset\hat{\mathsf{D}} .
\]
Both $m(\cdot)$ and $[\cdot]_{m'}$ are monotone under inclusion. Part~($\flat5$) of
Lemma~\ref{4.21:lem-grain-properties}, applied with $\mathsf{D}=\mathsf{d}(e)$, gives
$m(\hat{\mathsf{D}})=m(\mathsf{d}(e))$ and $[\hat{\mathsf{D}}]_{m'}=[\mathsf{d}(e)]_{m'}$ for
$m'\ge m(\mathsf{d}(e))$. The sandwich then gives the remaining two identities in
\eqref{4.21:eq-grain-resummation-7}.

\emph{(v), the finite identity.} The lattice is finite, so $\mathfrak{E}$, the bond sets $P$ and
the sets $Z$ entering the expansion are finitely many, and every sum below is finite. The fibres
of $e\mapsto\hat e$ partition $\mathfrak{E}$, so $\sum_e v_e=\sum_{\hat e}v_{\hat e}$. Moreover,
\[
\exp(v_{\hat e})-1=\prod_{e\mapsto\hat e}\big(1+(\exp(v_e)-1)\big)-1
=\sum_{\emptyset\neq D\subset\{e:\,e\mapsto\hat e\}}\ \prod_{e\in D}(\exp(v_e)-1).
\]
For a finite $\hat D\subset\hat{\mathfrak{E}}$, write $D\mapsto\hat D$ when the image of
$D\subset\mathfrak{E}$ under the fibre map is exactly $\hat D$. Multiplying the last identity over
$\hat e\in\hat D$ gives
\begin{equation}\label{4.21:eq-grain-resummation-9}
\prod_{\hat e\in\hat D}(\exp(v_{\hat e})-1)=\sum_{D\mapsto\hat D}\ \prod_{e\in D}(\exp(v_e)-1).
\end{equation}

We use the notation of the proof of the convergence proposition for the polymer expansion of the
preparation step. The step of that proof that forms the quantity $\mathcal{T}(D,P,Z)$ builds it
from a finite set $D$ of elements, a bond set $P$ and sets $Z_0$, $Z$. The operations involved read
only $(Y_0,P,Z_0,Z)$, and they are linear in
\[
\mathcal{G}:=(-1)^{|P|}\chi_{Y_0}\chi^c_P\prod_{e\in D}(\exp(v_e)-1),
\]
where $\chi_{Y_0}$ and $\chi^c_P$ are the indicator factors of that step. Here $Y_0$ is the set of
bonds of $\bigcup_{e\in D}\mathsf{b}(e)$. By \eqref{4.21:eq-grain-resummation-7}, $Y_0$ is the
same for $D$ and for $\hat D$ whenever $D\mapsto\hat D$. By linearity and
\eqref{4.21:eq-grain-resummation-9},
\[
\mathcal{T}(\hat D,P,Z)=\sum_{D\mapsto\hat D}\mathcal{T}(D,P,Z).
\]

The step of that proof that joins the components of $Z$ through the cube sets $\mathsf{b}(e)$
forms polymers $\gamma$ with activity $z(\gamma)$ and support $Z_\gamma$. These cube sets are the
same for $D$ and for $\hat D$. So, writing $\gamma\mapsto\hat\gamma$ for the replacement of the
elements by their images,
\[
z(\hat\gamma)=\sum_{\gamma\mapsto\hat\gamma}z(\gamma),\qquad Z_{\hat\gamma}=Z_\gamma .
\]
Compatibility of polymers and the Ursell coefficient $\rho^T$ read the supports $Z_\gamma$ only;
the superscript $T$ of $\rho^T$ stands for truncated and is unrelated to the cell set $T$.

The hull of a cluster is $[\bigcup_o\mathsf{d}(o)]_m$ with $m=\max_o m(\mathsf{d}(o))$, the union
and the maximum running over the elements $o$ of its polymers. The $\pi_m$-hull of a union is the
union of the $\pi_m$-hulls, and $m\ge m(\mathsf{d}(o))$ for every $o$. So by
\eqref{4.21:eq-grain-resummation-7} the hull is the same for the elements $o$ and for their images
$\hat o$.

So is the strong set assigned to a polymer in the preparation step,
$\mathcal{S}(\gamma)=\bigcup_{e\in\gamma}\mathcal{S}(c(e))$, since $c$ is constant on each fibre.
Accordingly, the identity below holds for the sums $a_n(X,\mathcal{S})$ with prescribed hull and
strong set, and the argument is read with these throughout.

Group the ordered $n$-tuples of polymers on $\mathfrak{E}$ according to their image tuples. Then
\[
\frac1{n!}\sum_{(\hat\gamma_i):\ \text{hull }X}\rho^T(\hat\gamma_1,\dots,\hat\gamma_n)
\prod_{i=1}^n z(\hat\gamma_i)
=\frac1{n!}\sum_{(\gamma_i):\ \text{hull }X}\rho^T(\gamma_1,\dots,\gamma_n)
\prod_{i=1}^n z(\gamma_i)=:a_n(X).
\]
Both sides are finite sums. They are equal term-group by term-group: $\rho^T$ and the hull depend
on the image tuple only, and $\prod_iz(\hat\gamma_i)$ is the sum of $\prod_iz(\gamma_i)$ over the
tuples $(\gamma_i)$ with image $(\hat\gamma_i)$.

Hence $\mathbb{C}_I(X)=\sum_na_n(X)$ is one and the same function in the two constructions. The
remaining steps of the proof of the convergence proposition, other than the two used above, run on
$\hat{\mathfrak{E}}$, give $\sum_n|a_n(X)|<\infty$ together with its norm bound.
\end{proof}

\begin{remark}[What the resummation by grains leaves unchanged]\label{4.21:rem-resummation-remarks}
We record five facts about the resummation of the decoupling expansion by grains
(Lemma~\ref{4.21:lem-grain-resummation}).
\begin{enumerate}
\item The resummation is performed inside a proof, and it changes none of the following:
the expansion to which the lemma is applied, its elements, its exponent, the hull rule, and
the terms $\mathbb{C}_I(X)$. What changes is only the order in which a convergent
series is majorised. In particular the decoupling clause \eqref{4.21:eq-decoupling-clause-a}
holds at the preparation site exactly as stated.
\item In the published construction the resummation is the step performed at the preparation
site in \cite{B-CMP122a}, and at the site of the operation $\mathbf{T}$ it is the step stated in
\cite[p.~276]{B-CMP119} in the words ``Then we resum them''. At the site of $\mathbf{T}$ no
thin layer occurs, since the layers of \cite{B-CMP119} at that site are those made by the
operation $\mathbf{T}$. The sentence cited at the preparation site concerns the resummation
only and not the count of the resummed families.
\item In the estimate at the preparation site, the fine cells that occur lie either in the
grain neighbourhood $\hat{\mathsf{K}}_{\square}$ of a block $\square$ or in at most
$\hat{b}_t N^{\flat}$ neighbouring grains $G$. They are summed exactly: a cell in
$\hat{\mathsf{K}}_{\square}$ carries the parameter value $\mathsf{s}=1$, and a cell in a grain
$G$ carries the single variable $s(G)$. The windows of earlier arrested attempts, which lie at
levels at most $k_a-1\le h-1$ by the arrest condition, the thin zones of the present
preparation, which lie at levels at most $j-1$, and every other cell below level $j+1$ are
absorbed by clause~(b) of the lemma bounding sums over fine cells. Cells off the large-field
region $Z$ are absorbed by clause~(a) of the same lemma.
\item Branching between adjacent levels costs only a lower bound on $\kappa_1$ that depends on
$d$ and $L$ alone. Unbounded relative fineness of cells is not controlled by a lower bound of
this kind.
\item The same geometry does not reach the sealed polydisc. This is the content of the last
bracketed clause of the lemma on the decoupling family $\sigma_0$.
\end{enumerate}
\end{remark}

\begin{proof}
We use the notation of Lemma~\ref{4.21:lem-grain-resummation} in its application at the
preparation site, with $A$ as in that lemma. Item~2 describes \cite[p.~276]{B-CMP119} and the
sentence of \cite{B-CMP122a} cited at the preparation site; item~5 restates the last bracketed
clause of the lemma on the decoupling family $\sigma_0$; item~4 records the lower bound on
$\kappa_1$ imposed by branching between adjacent levels. We prove items~1 and~3.

Let $\mathsf{y}\subset\sigma_0(A)$ be a family of cut cells. Consider the members of
$\hat{\sigma}_0(A)$ that hold a cell of $\mathsf{y}$. These form a set $\hat{\mathsf{y}}$ that
is determined by $\mathsf{y}$, so $f_{\mathsf{y}}$ enters the fibre $f_{\hat{\mathsf{y}}}$ of
exactly one family of grains.

Part~(i) of Lemma~\ref{4.21:lem-grain-resummation} gives
\begin{equation*}
\sum_{\hat{\mathsf{y}}} f_{\hat{\mathsf{y}}}
=\sum_{\mathsf{y}} f_{\mathsf{y}}
=f(\mathbf{1}).
\end{equation*}
Thus the grain sums regroup the terms $f_{\mathsf{y}}$ of the original expansion. The
functional $f(\mathbf{1})$ and each term $f_{\mathsf{y}}$ are the same before and after the
regrouping. The regrouping does not alter the expansion, its elements, its exponent, the hull
rule or the terms $\mathbb{C}_I(X)$, none of which is redefined by it; only the order of
majorisation differs. This is the first item.

For the third item, recall the definition of $\mathsf{s}_s$ in
Lemma~\ref{4.21:lem-grain-resummation}: a cell $\Delta\in\sigma_0(A)$ with
$\Delta\subset G\in\hat{\sigma}_0(A)$ receives $\mathsf{s}_s(\Delta)=s(G)$, and a cell with
$\Delta\subset\hat{\mathsf{K}}_A$ receives $\mathsf{s}_s(\Delta)=1$. Moreover the cells of
$\mathsf{y}$ inside $\hat{\mathsf{K}}_A$ are free in the fibre. Therefore the sum over these
cells is carried out exactly, at the value $1$, inside $f_{\hat{\mathsf{y}}}$. Likewise the sum
over the cells of one grain $G$ is carried out exactly through the one variable $s(G)$ in the
integral of part~(i).
\end{proof}

\begin{definition}[Constants of the preparation step after resummation]\label{4.21:not-preparation-constants}
Let $L$, $\kappa$, $\kappa_1$, $M$, $q$ be as in the quantifier prefix (Notation~\ref{2.1:not-prefix}),
chosen there in their fixed order; $d=4$ throughout. Let
$a_k$, $\hat{b}_t$, $N^{\flat}$, $c_{\flat}$, $c_J$, $\sigma_{*}$, $K_A$,
$K_2$, $K_{\flat}$ be the constants of the preparation step. Here $a_k$ is the rate in the format
bounds, $\hat{b}_t$ bounds the number of touching members, and $N^{\flat}$ is the member count of the
grain partition $\pi^{\flat}$. The constant $c_{\flat}$ is the one fixed in the flat bound of the
preparation step, where $c_{\flat}\le 9\kappa-8$. The constant $K_A$ is the one fixed with the
estimates on the Möbius polydisc of the preparation step. The constant $K_2$ is the one of the flat
bound of the preparation step that enters $K_{\flat}$ linearly and contains the factor $e^{c_J}$. We put
\begin{equation}\label{4.21:eq-preparation-constants-a}
\begin{aligned}
\mathsf{x}&:=\exp\bigl(-(\kappa_1-1)+(a_k+1)+3^d c_{\flat}\bigr),\\
&\phantom{:=}\ \ 2e^2\,\hat{b}_t^{\,2}\,N^{\flat}\,\mathsf{x}\le 1,\\
K_2^{\flat}&:=K_A\exp\bigl[(a_k+1+3^d c_{\flat})N^{\flat}+1\bigr],\\
c_{\bigstar}&:=214^d,\\
K_{\flat}''&:=e^{c_J+\sigma_{*}+1}\,c_{\bigstar}\,\frac{K_2^{\flat}}{K_2}\,K_{\flat}\,
\bigl(1+C^{P}_{\mathsf{a}}\bigr),
\end{aligned}
\end{equation}
and we require the inequality in the second line of \eqref{4.21:eq-preparation-constants-a}, the
resummation floor on $\kappa_1$.

That inequality is a lower bound on $\kappa_1$, imposed after $L$ and $\kappa$ have been fixed. With
$\mathsf{x}$ as in the first line, it is the hypothesis
\begin{equation*}
\mathsf{x}>0,\qquad 2e^2\,\hat{b}_t^{\,2}\,N^{\flat}\,\mathsf{x}\le 1
\end{equation*}
of clause~(iv) of Lemma~\ref{4.21:lem-grain-resummation}. It does not replace the per-cube
condition of the preparation step,
\begin{equation*}
e^{-(\kappa_1-1)+c_J+c_{\flat}}\le e^{-4},
\end{equation*}
which remains the lower bound on $\kappa_1$ imposed there. That condition is used for the decoupled
cubes in the fifth step of the flat bound of the preparation step. These cubes belong to the family
of fluctuation cubes of the step and meet no window cube.

The constant $K_2^{\flat}$ depends on $(d,L,\kappa,\kappa_1)$ only.

Recall from the flat bound of the preparation step that the constant $K_{\flat}$ is
\begin{equation}\label{4.21:eq-preparation-constants-1}
K_{\flat}=2^{3d+6}(5L)^d\,K_S\,K(d)\,A_{\flat}\,K_2\,M^{q}\,e^{C_2\kappa_1}.
\end{equation}
Here $K_S$, $K(d)$, $A_{\flat}$ and $C_2$ are constants fixed in that bound; none of them involves
$K_2$. The right
side of \eqref{4.21:eq-preparation-constants-1} is linear in $K_2$. Hence the factor
$K_2^{\flat}/K_2$ in the last line of \eqref{4.21:eq-preparation-constants-a} replaces $K_2$ by
$K_2^{\flat}$ in $K_{\flat}$. In \eqref{4.21:eq-preparation-constants-1} the factor $e^{c_J}$ is
contained in $K_2$. It is not contained in $K_2^{\flat}$, and it is restored as the factor $e^{c_J}$
in front in $K_{\flat}''$. The constant $C^{P}_{\mathsf{a}}$ is the anchoring constant of clause~(d)
of the proposition on the sums of term weights over old anchors below. The constants in
\eqref{4.21:eq-preparation-constants-a} are chosen generously and are not
optimised. All of them are fixed before $\gamma$.
\end{definition}

\begin{lemma}[The admissible couplings at the preparation site form an interval]\label{4.21:lem-coupling-interval}
Let $\mathsf{x}$, $K_2^{\flat}$, $c_{\bigstar}$ and $K_{\flat}''$ be the constants of the preparation step after
resummation (Notation~\ref{4.21:not-preparation-constants}). Let $V'$ be the block sum and $c_{\mathsf{p}}$ the
multiplicity constant of the preparation site, both defined later in this chapter. Write $\mathsf{T}$ for the
logarithm of the inverse coupling, and $\mathsf{T}_\gamma$ for its value at the coupling $\gamma$, so that
$\mathsf{T}_\gamma$ increases to $+\infty$ as $\gamma$ decreases to $0$. The number $\gamma$ bounds the running
couplings: at every step considered $g_k\le\gamma$, so that the value of $\mathsf{T}$ at $g_k$ is at least
$\mathsf{T}_\gamma$.

We say that the \emph{coupling condition of the preparation site} holds at $\gamma$ if, for every $K$, every
$k\le K$ and every coupling flow with $g_k\le\gamma$, the following two requirements are met.
\begin{itemize}
\item[(a)] The inequality $2e^2\hat{b}_t^2N^{\flat}\mathsf{x}\le 1$ of \eqref{4.21:eq-preparation-constants-a} holds.
It involves $(d,L,\kappa)$ and no coupling.
\item[(b)] The smallness hypotheses of the group-total bound of the preparation site hold, and the remaining
smallness hypotheses of the per-cube bound of the preparation site hold when read with the triple
$(V,K_2,K_{\flat})$ replaced by $(c_{\mathsf{p}}V',K_2^{\flat},K_{\flat}'')$. Together these are the inequalities
\begin{equation}\label{4.21:eq-coupling-interval-a}
\begin{gathered}
K_{\flat}''\,\mathsf{r}\,\check{R}_k\,c_{\mathsf{p}}V'\,\bigl(C_0^{\sharp}+E_0+B_0+2\bigr)
\bigl(\vartheta+K_Q/r_A+\varepsilon_S\bigr)\le\tfrac14,\\
K_{\flat}''\,\mathsf{r}\,c_{\mathsf{p}}V'\,\bigl(\mathfrak{p}_j+\tilde{\mathfrak{l}}\bigr)\le\tfrac14,
\qquad
6\check{R}_k\,c_{\mathsf{p}}V'\,\bigl(K_2^{\flat}\vartheta+K_Q/r_A\bigr)\le 1,\\
\vartheta\le c_R/2,\qquad r_A\ge 2,\qquad \kappa_A\ge c_{\flat}+c_J+1,\\
\check{R}_k\,e^{-\sigma_*R_*/L}\le 1,\qquad R_k\ge 2 .
\end{gathered}
\end{equation}
\end{itemize}
Here $\check{R}_k$ is the region size parameter at step $k$ and $\mathsf{r}$ is the raise read by a block of the
cover; $\vartheta$ is the per-block weight and $\mathfrak{p}_j$ the rim weight; $\varepsilon_S$ and
$\tilde{\mathfrak{l}}$ are small quantities of the group totals of the preparation site, $r_A$ is its large
parameter, which enters through $1/r_A$, and $K_Q$ is the factor multiplying $1/r_A$; the quantities $\vartheta$,
$1/r_A$, $\varepsilon_S$, $\mathfrak{p}_j$ and $\tilde{\mathfrak{l}}$ depend on the coupling, and their smallness
makes the first three members hold;
$\kappa_A$ is the weight of the strong set; $c_J$ and $c_{\flat}$ are the two constants in the exponent of the
per-cube floor $e^{-(\kappa_1-1)+c_J+c_{\flat}}\le e^{-4}$ on $\kappa_1$ in the per-cube bound of the preparation
site; $C_0^{\sharp}$ and $B_0$ are
constants of the format bounds of stored terms. The letters $E_0$, $R_*$ and $R_k$ denote quantities of the
group-total bound and of the per-cube bound of the preparation site.
The number $\sigma_*$ is the constant of the preparation site that also enters the exponent $c_J+\sigma_*+1$ in the
definition of $K_{\flat}''$ (Notation~\ref{4.21:not-preparation-constants}), and it is fixed before $\gamma$; $c_R$
is related to the constant $K_A$ of
Notation~\ref{4.21:not-preparation-constants} by $K_A=2/c_R$.

Assume that requirement (a) holds. Then the following hold.
\begin{enumerate}
\item[(i)] Each member of \eqref{4.21:eq-coupling-interval-a} is implied, for every $k\le K$ and every $K$, by one
inequality of the form
\begin{equation}\label{4.21:eq-coupling-interval-1}
\sup_{\mathsf{T}\ge\mathsf{T}_\gamma}\Phi(\mathsf{T})\le c .
\end{equation}
Here $\Phi$ is a function of $\mathsf{T}$ alone, its coefficients are fixed before $\gamma$, and
$\Phi(\mathsf{T})\to0$ as $\mathsf{T}\to\infty$; $c>0$ is a number fixed before $\gamma$. For the first three
members $\Phi$ is the left-hand side of \eqref{4.21:eq-coupling-interval-2} and $c=1$, that is, the inequality is
\eqref{4.21:eq-coupling-interval-2} for every $\mathsf{T}\ge\mathsf{T}_\gamma$:
\begin{equation}\label{4.21:eq-coupling-interval-2}
K'\,\mathsf{T}^{(1+d+d\beta_0^{\mathrm{pr}})r}
\Bigl[\mathsf{T}^{-(q_*-p_1)}+\mathsf{T}^{-(q_*-p_0)}+e^{-\kappa_A^{\mathrm{pr}}+c_{\flat}}\Bigr]\le 1 .
\end{equation}
Here $K'$ is a number fixed before $\gamma$; $p_0$, $p_1$, $q_*$ and $r$ are among Ba{\l}aban's free integer
parameters; $\beta_0^{\mathrm{pr}}$ is
a further exponent, and $\kappa_A^{\mathrm{pr}}$ is a weight distinct from the strong-set weight $\kappa_A$ of the
sixth member of \eqref{4.21:eq-coupling-interval-a}; these are the exponents and the weight with which the
smallness hypotheses of the group-total bound of the preparation site are stated.
Inequality \eqref{4.21:eq-coupling-interval-2} is used under the floor on Ba{\l}aban's free integer parameters that
is among the hypotheses of the group-total bound of the preparation site.
\item[(ii)] Call $\gamma$ \emph{admissible} if, for each member of \eqref{4.21:eq-coupling-interval-a}, the
inequality \eqref{4.21:eq-coupling-interval-1} of (i) that implies it for every $K$ and every $k\le K$ holds at
$\gamma$. The set of admissible $\gamma$ is an
interval $(0,\gamma_P]$. The number $\gamma_P>0$ is fixed by the constants chosen before $\gamma$. If $\gamma$ is
admissible, every $\gamma'<\gamma$ is admissible, and at every admissible $\gamma$ the coupling condition of the
preparation site holds. Moreover, if the coupling condition of the preparation site holds at $\gamma$, it holds
at every $\gamma'<\gamma$.
\item[(iii)] Admissibility only requires the coupling to be small. No member of the coupling condition bounds $g$,
$g_k$ or $\gamma$ from below.
\end{enumerate}
\end{lemma}

\begin{proof}
We first prove (i). In each member of \eqref{4.21:eq-coupling-interval-a} we distinguish its cost side, the
quantities whose growth makes the member harder to satisfy, from its paying side, the small quantities that make it
hold; in the first three members the left-hand side is a product of quantities of both kinds. We read every quantity
of the cost side at its upper function of $\mathsf{T}$:
\begin{itemize}
\item the region size parameter obeys $\check{R}_k\le C_N\mathsf{T}^{r_{\mathrm{pr}}}$;
\item the block sum obeys $V'\le C'_V\bigl(104L\,\mathsf{R}^{\uparrow}(\mathsf{T})\bigr)^d$, by the volume bound
proved together with the definition of $V'$, with the constant $C'_V$ of that bound, fixed before $\gamma$;
\item the raise obeys $\mathsf{r}\le e\,(2+6C_N2^{r_{\mathrm{pr}}})$, by the bound that compares the stage count with
the region size.
\end{itemize}
We read the paying side, the small factors $\vartheta$, $1/r_A$, $\varepsilon_S$, $\mathfrak{p}_j$ and
$\tilde{\mathfrak{l}}$, at their upper functions of $\mathsf{T}$ as well. The remaining five members compare a single
quantity, or the product of $\check{R}_k$ with an exponential factor, with a number, and they are read in the same
way. The remaining letters of \eqref{4.21:eq-coupling-interval-a} are read in the same way: $c_R$, $c_{\flat}$, $c_J$
and $\sigma_*$ are numbers fixed before $\gamma$ (Notation~\ref{4.21:not-preparation-constants}), and $C_0^{\sharp}$,
$E_0$, $B_0$, $K_Q$, $R_*$ and $R_k$ are taken as functions of $\mathsf{T}$, a number fixed before $\gamma$ being a
constant such function. After these replacements all quantities are functions of $\mathsf{T}$.

The remaining constants $K_2^{\flat}$ and $K_{\flat}''$, and $c_{\bigstar}$, which enters through $K_{\flat}''$, are
numbers fixed before $\gamma$ (Notation~\ref{4.21:not-preparation-constants}); $c_{\mathsf{p}}$ is a number as well,
by its definition later in this chapter. None of them depends on $\mathsf{T}$. Hence each member
is bounded by a function $\Phi$ of $\mathsf{T}$ alone, with coefficients fixed before $\gamma$. The degree of $\Phi$ in
$\mathsf{T}$ is that of the smallness hypotheses of the group-total bound of the preparation site.

For the first three members this function is the left-hand side of \eqref{4.21:eq-coupling-interval-2}. That
inequality is read under the floor of (i). This proves (i).

We now prove (ii). Fix one of the functions $\Phi$ of (i). As $\gamma$ decreases, $\mathsf{T}_\gamma$ increases,
because $\mathsf{T}$ is the logarithm of the inverse coupling. A supremum of $\Phi$ over
$\{\mathsf{T}\ge\mathsf{T}_\gamma\}$
does not increase when $\mathsf{T}_\gamma$ increases. So if \eqref{4.21:eq-coupling-interval-1} holds at $\gamma$, it
holds at every $\gamma'<\gamma$.

Hence if $\gamma$ is admissible, every $\gamma'<\gamma$ is admissible.
At an admissible $\gamma$ requirement (a) holds by assumption, and every member of
\eqref{4.21:eq-coupling-interval-a} holds, since by (i) it is implied by
an inequality of the form \eqref{4.21:eq-coupling-interval-1} that holds there; so the coupling condition of the
preparation site holds at $\gamma$.

If the coupling condition of the preparation site holds at $\gamma$ and $\gamma'<\gamma$, then every coupling flow
with $g_k\le\gamma'$ also has $g_k\le\gamma$, so requirement (b) holds for it, and requirement (a) involves no
coupling; hence the coupling condition holds at $\gamma'$.

The set is also non-empty. Since $\Phi(\mathsf{T})\to0$ as $\mathsf{T}\to\infty$, and $\mathsf{T}_\gamma$ grows
without bound as $\gamma$ decreases to $0$, the supremum in \eqref{4.21:eq-coupling-interval-1} is at most $c$ once
$\gamma$ is small enough. This holds for each of the finitely many members. The threshold so obtained depends only on
the functions $\Phi$, whose coefficients are fixed before $\gamma$. This gives the interval $(0,\gamma_P]$ with
$\gamma_P>0$ fixed by the constants chosen before $\gamma$, which is (ii).

Finally we prove (iii). By (i), each member is implied by a bound of the form
\eqref{4.21:eq-coupling-interval-1}, and admissibility asks for exactly these bounds. By the
monotonicity just shown, such a bound is preserved when $\gamma$ is decreased. So no member imposes a lower bound on
$g$, on $g_k$ or on $\gamma$.
\end{proof}

\begin{remark}
The following inequality is a hypothesis of no statement in this book:
\begin{equation*}
e^{-(\kappa_1-1)}\bigl(10\,\check{R}_{k_0+1}/L\bigr)^{d-1}\le\mathrm{const}.
\end{equation*}
Here $\check{R}_{k_0+1}$ is the region size parameter at step $k_0+1$.
It is what the per-cube bound of the preparation site would need beside a window, that is, beside a probe cube,
if read term by term. Two facts
make it unsuitable. First, $\kappa_1$ is fixed before $\gamma$. Second, $\check{R}_{k_0+1}$ increases to infinity as
$g$ decreases to $0$. Together these mean the inequality holds only when $g_k$ lies above a threshold. It therefore
bounds the running coupling $g_k$ from below, which makes it a cap (compare Definition~\ref{2.3:def-cap}), and for
this reason no member of the coupling condition of the
preparation site has this form. In Ba{\l}aban's construction \cite{B-CMP122a} the coupling is restricted in
the same one-sided way: the relevant quantities are made arbitrarily small by choosing $\gamma$ sufficiently small,
with bounds that do not depend on any scale, and $g_k$ is assumed sufficiently small.
\end{remark}

\begin{remark}\label{4.21:prop-preparation-machine}
The statement of this proposition is not written out here; parts of its proof are printed below.
\end{remark}

\begin{lemma}[The block sum in clause (c) of the relative machine at the preparation site]
\label{4.21:prf-block-sum-proof}
Work in the setting and notation of Proposition~\ref{4.21:prop-preparation-machine}. Thus $c$ is a
wrapped member of the class considered there, with core $S_c$, reading set $X_c$, pocket list
$\mathcal{P}_c$, the block set $\mathcal{A}_c=\mathcal{A}_c^{\mathrm{pl}}\sqcup\mathcal{A}_c^{\mathrm{isl}}$
and the corridors $\operatorname{Cor}(\square)$, empty on $\mathcal{A}_c^{\mathrm{pl}}$. The regions
$R(\bar c)$ are those attached to the cells $\bar c$ of stage $n$, and $R_n$ is the collar region at
stage $n$; $d^{\wedge}$ is the scaled distance, $w=M/M_1$, $\delta_P$ the decay rate, and
\begin{equation*}
V'=\sup_{n,\bar c}\ \sum_{\square}e^{-\frac{1}{16}\delta_P d^{\wedge}(\square,R(\bar c))},
\end{equation*}
the sum running over all blocks. Further $c_{\mathsf{p}}=212^d$, $\mathsf{p}^{\mathrm{rel}}(c)=2$ if $c$ is
confined and $\mathsf{p}^{\mathrm{rel}}(c)=c_{\mathsf{p}}(2+d_k([S_c]_k)/\check{R}_k)$ otherwise, and
$\check{R}^{\uparrow}=\max_{h<i\le k}\check{R}_i=\check{R}_{h+1}$, $h$ being the level fixed in the setting of
Proposition~\ref{4.21:prop-preparation-machine}. For a family $\mathfrak{B}$ of blocks
and a region $R$ put $d^{\wedge}(\mathfrak{B},R):=\min_{\square\in\mathfrak{B}}d^{\wedge}(\square,R)$.
Assume:
\begin{enumerate}
\item clause (b) of Proposition~\ref{4.21:prop-preparation-machine}: for
$\square\in\mathcal{A}_c^{\mathrm{isl}}$ one has $(a_k+1)d(x_\square+1)\le\frac14\delta_P
d^{\wedge}(\square,R_n)$, while $\operatorname{Cor}(\square)$ consists of fewer than $d(x_\square+1)$
cubes;
\item part (c) of the cell lemma, which at $\beta=\delta_P/16$ bounds the block sum against a region
made of a family $\Sigma$ of cells by $C'_V\,\#\Sigma$, where $C'_V:=C_{\mathrm{blk}}(1+\hat b_t
C_{\mathrm{cell}})$;
\item the size bound on the regions: each $R(\bar c)$ lies in a box of side at most
$100\check{R}_{j+1}$ level-$(j+1)$ units and consists of cells of two consecutive levels $j$ and
$j+1$, for some $j=j(\bar c)$ with $h<j+1\le k$;
\item the counting bounds for cells: distinct cells of the stage lie at level-$k$ distance at least
$\check{R}_k$ from one another; for every $l\ge1$, the cells $\bar c$ with
$\operatorname{dist}_k(\bar c_0,\bar c)\in[l\check{R}_k,(l+1)\check{R}_k)$ from a given cell $\bar c_0$
number at most $(2l+104)^d$; the cells within $\check{R}_k$ of a box of side $208\check{R}_k$ number at
most $211^d$, and those at level-$k$ distance in $[l\check{R}_k,(l+1)\check{R}_k)$ from it at most
$(2l+211)^d$;
\item the geometric statement on strips and probes: a strip of a pocket in $\mathcal{P}_c$ touches
$S_c$ and has sides at most $101\check{R}_k$, the probe $\tilde\square$ of a block of $\mathcal{A}_c$
has side at most $5$ level-$k$ units and meets $S_c$ or such a strip, and every cube of $[S_c]_k$ meets
an optimal tree $\mathcal{T}$ of $[S_c]_k$; from these facts that statement concludes that, whenever
points $z_0,\dots,z_r$ of $\mathcal{T}$ are such that every point of $\mathcal{T}$ is within
$\check{R}_k$ of one of them, every block of $\mathcal{A}_c$ lies in one of the boxes
$\{y:\ |y-z_a|_\infty\le104\check{R}_k\}$, $a=0,\dots,r$;
\item the tail sum $\tau_d(\check{R}_k)$ majorises the two series
\begin{equation*}
\sum_{l\ge1}(2l+104)^d\,e^{-\frac3{16}\delta_P w(l\check{R}_k-5)}
\qquad\text{and}\qquad
\sum_{l\ge1}(2l+211)^d\,e^{-3l\check{R}_k},
\end{equation*}
and the tail condition under which clause (c) of
Proposition~\ref{4.21:prop-preparation-machine} is stated gives $\tau_d(\check{R}_k)\le1$;
\item $\frac{3}{16}\delta_P w\ge3$;
\item the collar region $R_n$ is contained in the union of the regions $R(\bar c)$, $\bar c$ running
over the cells of stage $n$;
\item for every block $\square$ and every cell $\bar c$ of the stage,
$d^{\wedge}(\square,R(\bar c))\ge w\operatorname{dist}_k(\square,\bar c)$, where $\operatorname{dist}_k$
is the distance in level-$k$ units.
\end{enumerate}
Then the following hold.
\begin{enumerate}
\item[(A)] $V'\le C'_V(104L\check{R}^{\uparrow})^d$.
\item[(B)] For every family $\mathfrak{B}$ of blocks, every $\bar c$ and every $\beta\ge\delta_P/16$,
\begin{equation}\label{4.21:eq-block-sum-proof-a}
\sum_{\square\in\mathfrak{B}}e^{-\beta d^{\wedge}(\square,R(\bar c))}
\le e^{-(\beta-\delta_P/16)\,d^{\wedge}(\mathfrak{B},R(\bar c))}\,V',
\end{equation}
and, for every block $\square$,
\begin{equation}\label{4.21:eq-block-sum-proof-1}
e^{-\beta d^{\wedge}(\square,R_n)}\le\sum_{\bar c}e^{-\beta d^{\wedge}(\square,R(\bar c))}.
\end{equation}
\item[(C)] The block sum of clause (c) of Proposition~\ref{4.21:prop-preparation-machine} obeys
\begin{equation}\label{4.21:eq-block-sum-proof-2}
\sum_{\square\in\mathcal{A}_c}e^{-\frac12\delta_P d^{\wedge}(\square,R_n)}\,
e^{(a_k+1)|\operatorname{Cor}(\square)|}\le V'\,\mathsf{p}^{\mathrm{rel}}(c).
\end{equation}
\end{enumerate}
\end{lemma}

\begin{proof}
\emph{The bound on $V'$.} Fix $n$ and $\bar c$, and let $\Sigma$ be the family of cells of which
$R(\bar c)$ consists. By assumption (3), $R(\bar c)$ lies in a box of side at most $100\check{R}_{j+1}$
level-$(j+1)$ units and $\Sigma$ consists of cells of levels $j$ and $j+1$; counting the cells of these
two levels in such a box gives $\#\Sigma\le(104L\check{R}_{j+1})^d$. Part (c) of the cell lemma,
assumption (2), at $\beta=\delta_P/16$ bounds the block sum against $R(\bar c)$ by $C'_V\#\Sigma$, hence
by $C'_V(104L\check{R}_{j+1})^d\le C'_V(104L\check{R}^{\uparrow})^d$, since $h<j+1\le k$. Taking the
supremum over $n$ and $\bar c$ gives (A).

\emph{The tool.} Let $\beta\ge\delta_P/16$ and $\square\in\mathfrak{B}$. Since
$d^{\wedge}(\square,R(\bar c))\ge d^{\wedge}(\mathfrak{B},R(\bar c))$ and $\beta-\delta_P/16\ge0$,
\begin{equation*}
e^{-\beta d^{\wedge}(\square,R(\bar c))}
=e^{-(\beta-\delta_P/16)d^{\wedge}(\square,R(\bar c))}\,e^{-\frac1{16}\delta_P d^{\wedge}(\square,R(\bar c))}
\le e^{-(\beta-\delta_P/16)d^{\wedge}(\mathfrak{B},R(\bar c))}\,
e^{-\frac1{16}\delta_P d^{\wedge}(\square,R(\bar c))}.
\end{equation*}
Summing over $\square\in\mathfrak{B}$ and bounding the sum of the last factors over $\mathfrak{B}$ by
the sum over all blocks, which is at most $V'$ by the definition of $V'$, gives
\eqref{4.21:eq-block-sum-proof-a}. For \eqref{4.21:eq-block-sum-proof-1}, fix a block $\square$ and
$\beta\ge\delta_P/16$. By assumption (8), $R_n$ lies in the union of the regions $R(\bar c)$ of the
stage, which are finitely many, the torus being finite. Since $d^{\wedge}(\square,\cdot)$ is a scaled
distance from $\square$ to a set, it does not increase when the set is enlarged, and for a finite union
it equals the minimum over the members of the union; hence
$d^{\wedge}(\square,R_n)\ge d^{\wedge}(\square,R(\bar c_*))$ for a cell $\bar c_*$ attaining that minimum.
As $\beta\ge\delta_P/16>0$, the left side of \eqref{4.21:eq-block-sum-proof-1} is at most
$e^{-\beta d^{\wedge}(\square,R(\bar c_*))}$, which is one of the non-negative terms of the sum on the
right.

\emph{Reduction of the summands.} Every summand on the left of \eqref{4.21:eq-block-sum-proof-2} is at
most $e^{-\frac14\delta_P d^{\wedge}(\square,R_n)}$. On $\mathcal{A}_c^{\mathrm{pl}}$ this holds because
$\operatorname{Cor}(\square)=\emptyset$, so the summand is
$e^{-\frac12\delta_P d^{\wedge}(\square,R_n)}\le e^{-\frac14\delta_P d^{\wedge}(\square,R_n)}$. On
$\mathcal{A}_c^{\mathrm{isl}}$, assumption (1) gives
\begin{equation*}
(a_k+1)|\operatorname{Cor}(\square)|\le(a_k+1)d(x_\square+1)\le\tfrac14\delta_P\,d^{\wedge}(\square,R_n),
\end{equation*}
so the second factor is at most $e^{\frac14\delta_P d^{\wedge}(\square,R_n)}$. By
\eqref{4.21:eq-block-sum-proof-1} at $\beta=\frac14\delta_P$ it therefore suffices to bound
\begin{equation}\label{4.21:eq-block-sum-proof-3}
\sum_{\bar c}\ \sum_{\square\in\mathcal{A}_c}e^{-\frac14\delta_P d^{\wedge}(\square,R(\bar c))},
\end{equation}
and every application of \eqref{4.21:eq-block-sum-proof-a} below is at $\beta=\frac14\delta_P$, where
$\beta-\delta_P/16=\frac3{16}\delta_P$. We use assumption (7), $\frac3{16}\delta_P w\ge3$, throughout.

\emph{The case in which $c$ is confined.} By the definition of confinement in
Proposition~\ref{4.21:prop-preparation-machine}, $X_c$ lies in a single treated component; we write
$\bar c_0$ for the cell of stage $n$ in which $c$ is then confined. Every block of $\mathcal{A}_c$ has
its probe $\tilde\square$, of side at most $5$ level-$k$ units, meeting $X_c$, by the definition of
$\mathcal{A}_c$; so the blocks of $\mathcal{A}_c$ lie within $5$ level-$k$ units of $\bar c_0$. Apply
\eqref{4.21:eq-block-sum-proof-a} with $\mathfrak{B}=\mathcal{A}_c$. Against $\bar c_0$ itself the
exponential prefactor is at most $1$, and the contribution is at most $V'$. Against $\bar c\ne\bar c_0$
one has $d^{\wedge}(\mathcal{A}_c,R(\bar c))\ge w(\operatorname{dist}_k(\bar c_0,\bar c)-5)$. By the
counting bounds, assumption (4), distinct cells lie at level-$k$ distance at least $\check{R}_k$, so every
$\bar c\ne\bar c_0$ has $\operatorname{dist}_k(\bar c_0,\bar c)\in[l\check{R}_k,(l+1)\check{R}_k)$ for
exactly one $l\ge1$, and for each $l\ge1$ these cells number at most $(2l+104)^d$; for such a cell the
prefactor is at most $e^{-\frac3{16}\delta_P w(l\check{R}_k-5)}$. Summing over $l\ge1$, the
contribution of all $\bar c\ne\bar c_0$ is at most
\begin{equation*}
V'\sum_{l\ge1}(2l+104)^d\,e^{-\frac3{16}\delta_P w(l\check{R}_k-5)}\le\tau_d(\check{R}_k)V',
\end{equation*}
the inequality being assumption (6). By the tail condition, also assumption (6), the total
\eqref{4.21:eq-block-sum-proof-3} is at most
\begin{equation*}
(1+\tau_d(\check{R}_k))V'\le2V'=V'\,\mathsf{p}^{\mathrm{rel}}(c).
\end{equation*}

\emph{The case in which $c$ is not confined.} Let $\mathcal{T}$ be an optimal tree of $[S_c]_k$ and
write $x$ for its length, measured in level-$k$ units, so that $x=d_k([S_c]_k)$. We choose points
$z_0,\dots,z_r\in\mathcal{T}$, with $r\le x/\check{R}_k+1$, such that every point of $\mathcal{T}$ is
within $\check{R}_k$ of some $z_a$. To do so, traverse $\mathcal{T}$ by its Euler tour, the closed walk
which runs along every edge of $\mathcal{T}$ once in each direction and so has length $2x$; cut this
tour into at most $x/\check{R}_k+1$ consecutive arcs of length at most $2\check{R}_k$ (one arc if
$x=0$, and $\lceil x/\check{R}_k\rceil$ arcs if $x>0$), and let $z_0,\dots,z_r$ be their midpoints, so
that $r+1\le x/\check{R}_k+1$; in particular $r\le x/\check{R}_k+1$ and $r+1\le2+x/\check{R}_k$.
Every point of $\mathcal{T}$ lies on some arc, hence within arc length $\check{R}_k$, and so within
distance $\check{R}_k$, of the midpoint of that arc.

Every block of $\mathcal{A}_c$ lies in one of the boxes
\begin{equation*}
\mathsf{B}_a:=\{y:\ |y-z_a|_\infty\le104\check{R}_k\},\qquad a=0,\dots,r.
\end{equation*}
Indeed, every cube of $[S_c]_k$ meets $\mathcal{T}$; by assumption (5) a strip of $\mathcal{P}_c$
touches $S_c$ and has sides at most $101\check{R}_k$, and a block of $\mathcal{A}_c$ has its probe
$\tilde\square$, of side at most $5$ level-$k$ units, meeting $S_c$ or such a strip. From these facts the
geometric statement, assumption (5), applied to the points $z_0,\dots,z_r$ just chosen, places every
block of $\mathcal{A}_c$ in some $\mathsf{B}_a$.

Fix $a$ and apply \eqref{4.21:eq-block-sum-proof-a} to the family of blocks of $\mathcal{A}_c$ lying in
$\mathsf{B}_a$. By the counting bounds, assumption (4), the cells $\bar c$ within $\check{R}_k$ of
$\mathsf{B}_a$ number at most $211^d$, and each carries a prefactor at most $1$; every other cell lies at
level-$k$ distance in $[l\check{R}_k,(l+1)\check{R}_k)$ from $\mathsf{B}_a$ for exactly one $l\ge1$, and
for each $l\ge1$ these cells number at most $(2l+211)^d$. For such a cell $\bar c$ and a block $\square$
in $\mathsf{B}_a$ one has $\operatorname{dist}_k(\square,\bar c)\ge l\check{R}_k$, hence
$d^{\wedge}(\square,R(\bar c))\ge wl\check{R}_k$ by assumption (9), and the prefactor is at most
$e^{-\frac3{16}\delta_P wl\check{R}_k}\le e^{-3l\check{R}_k}$ by assumption (7). Hence the blocks of
$\mathcal{A}_c$ in $\mathsf{B}_a$ contribute to \eqref{4.21:eq-block-sum-proof-3} at most
\begin{equation*}
\Bigl(211^d+\sum_{l\ge1}(2l+211)^d\,e^{-3l\check{R}_k}\Bigr)V'
\le(211^d+\tau_d(\check{R}_k))V'\le c_{\mathsf{p}}V',
\end{equation*}
the last two steps by assumption (6) and $211^d+1\le212^d=c_{\mathsf{p}}$.

There are $r+1\le2+x/\check{R}_k$ boxes. Summing the bound over them counts each block of
$\mathcal{A}_c$ once for every box in which it lies; as all terms are non-negative, this over-counts.
Therefore \eqref{4.21:eq-block-sum-proof-3} is at most
\begin{equation*}
c_{\mathsf{p}}\Bigl(2+\frac{d_k([S_c]_k)}{\check{R}_k}\Bigr)V'=V'\,\mathsf{p}^{\mathrm{rel}}(c),
\end{equation*}
which together with the confined case proves \eqref{4.21:eq-block-sum-proof-2}.
\end{proof}

\begin{remark}[The multiplicity of a core that is not confined]
In the proof of Lemma~\ref{4.21:prf-block-sum-proof} no sum over the strips of $\mathcal{P}_c$ is taken:
\eqref{4.21:eq-block-sum-proof-a} sums all blocks of a box at once. This matters because many old
strips may touch one coarse cube. A long core $S_c$ in a large untreated component can have of order
$d_k([S_c]_k)/\check{R}_k$ treated neighbours. The block-sum bound at the plain preparation site carries
the multiplicity $2$ for filing levels $m<k$, and this multiplicity falls short by that count; the
shortfall is harmless there as well, since the terms concerned carry a surplus that covers it. Here the
count is carried explicitly, by the term $d_k([S_c]_k)/\check{R}_k$ in the multiplicity
$\mathsf{p}^{\mathrm{rel}}(c)$ of a member $c$ that is not confined.
\end{remark}

\begin{lemma}[The relative machine at the preparation site: old anchors]\label{4.21:prf-old-anchor-proof}
This is clause~(d) of Proposition~\ref{4.21:prop-preparation-machine}. The setting and the
notation are those of that proposition and of Definition~\ref{4.21:def-wrapped-terms}: the
operation $\mathbf{R}_k$ and its stages, the level $h$ and the integer $N=k-h$ of that
proposition (here $N$ is not the rank of the gauge group), the bound $N(g_k)$ on the number of
stages of $\mathbf{R}_k$, the collar $R_n$ of the current stage, the set
$\mathfrak{O}$ of old blocks, the anchor sets $\mathcal{A}_c$, the weights $\mathfrak{w}(c)$
and the budgets $\mathsf{e}(c)$, $\mathsf{f}(c)$ of the wrapped members $c$, the scaled
distance $d^{\wedge}$ and $w = M/M_1$. Assume the following three statements.
\begin{itemize}
\item[(i)] The bound on boundary terms at the preparation site, evaluated with labels in place
of domains: for
$\ell\le k$, every boundary term $T$ of level $\ell$, of $\mathbf{B}^{(\ell)}$ or of
$\mathbf{B}'^{(\ell)}$, with label $A_T$, satisfies
$\|T\|e^{w(T)}\le A_{\flat}c_T e^{-\sigma_{*}d_\ell(A_T)}$, where $w(T)$ is the weight
exponent of the stored term $T$ (not the unit $w=M/M_1$) and $c_T$ is its constant in that
bound; the weight of $T$ satisfies $\mathfrak{w}(T)\le\|T\|e^{w(T)}$.
\item[(ii)] The bound on the chain terms of the current operation: for $\ell\le k$ and
$\square'\in\pi_\ell$, the chain terms of each stage of $\mathbf{R}_k$ filed at $\ell$ whose
label $S$ contains $\square'$ have total $\|\cdot\|^{A}$-norm at most $C_0^{\sharp}$, where at
$\ell=k$ the norm $\|\cdot\|^{A}$ includes the factor $e^{d_k(S)/N}$; each such term $c$
satisfies $\mathfrak{w}(c)\,e^{\mathbf{1}[\ell=k]\,d_k(S)/N}\le\|c\|^{A}$; the number $n$ of
stages performed satisfies $n\le N(g_k)$.
\item[(iii)] The per-cube bound at the preparation site for the class of wrapped members
whose per-cube constant is $b_{\delta}$: for $\ell\le k$ and $\square'\in\pi_\ell$,
\begin{equation*}
\sum_{c}W(c)\,e^{d_\ell(A_c)}\ \le\ A_{\flat}b_{\delta},
\end{equation*}
the sum running over the members $c$ of that class of level $\ell$ with $A_c\ni\square'$,
and $W(c)$ being the weight of (a) below.
\end{itemize}
Then the following hold.
\begin{itemize}
\item[(a)] With $i := k$, the class of probes
$\mathfrak{Q} := \{\tilde{\square}_0 : \square_0\in\mathfrak{O}\}$, the wrapped members $c$
(level $j_c$, label $A_c$, pockets $\mathcal{P}_c$, region $X_c$) and the weights
$W(c) := \mathfrak{w}(c)e^{\mathsf{e}(c)+\mathsf{f}(c)}$, the hypotheses
\eqref{4.21:eq-wrapped-terms-a} (the per-cube bound on terms of one level, the per-cube
bound with a spare rate, the growth law, and the depth datum) hold with the data
\begin{equation}\label{4.21:eq-old-anchor-proof-a}
\mathsf{b}^{P} := A_{\flat}\mathsf{b}_k,\qquad \varsigma_0 := 1,\qquad
\varrho(\tilde{\square}_0) := 1+\frac{d^{\wedge}(\square_0,R_n)}{w},\qquad N_{*} := N+2 ,
\end{equation}
where $\mathsf{b}_k$ is the per-cube function of Lemma~\ref{4.21:lem-wrapped-cube-bound} with
$i=k$.
\item[(b)] Every old block $\square_0\in\mathfrak{O}$ satisfies
\begin{equation}\label{4.21:eq-old-anchor-proof-1}
d^{\wedge}(\square_0,R_n)\ \ge\ w\,(\check{R}_k-5).
\end{equation}
\item[(c)] For every $\square_0\in\mathfrak{O}$,
\begin{equation*}
\mathfrak{S}(\square_0)
:= \sum\bigl\{\mathfrak{w}(c)e^{\mathsf{e}(c)+\mathsf{f}(c)} :
c \text{ wrapped},\ \square_0\in\mathcal{A}_c\bigr\}
\ \le\ C^{P}_{\mathsf{a}}\,e^{\frac18\delta_P d^{\wedge}(\square_0,R_n)},
\end{equation*}
with the constant $C^{P}_{\mathsf{a}}$ of \eqref{4.21:eq-old-anchor-proof-2} below; it depends
on no age and on none of $\check{R}_k$, $N$, $k$, $g_k$.
\end{itemize}
\end{lemma}

\begin{proof}
Let $\square_0\in\mathfrak{O}$, and write $n_0$ for the zone of $\square_0$; it is the level
of the probe $\tilde{\square}_0$ in Lemma~\ref{4.21:lem-probe-mass}. By the definition of
$\mathcal{A}_c$ in Proposition~\ref{4.21:prop-preparation-machine}, $\square_0\in\mathcal{A}_c$
if and only if the probe $\tilde{\square}_0$ meets $X_c$. Hence, with the weights $W(c)$ of
(a),
\begin{equation}\label{4.21:eq-old-anchor-proof-3}
\mathfrak{S}(\square_0)=\sum_{c:\ \tilde{\square}_0\text{ meets }X_c}W(c)
=\mathsf{S}(\tilde{\square}_0),
\end{equation}
the total of part~(iv) of the lemma on the mass read at one probe
(Lemma~\ref{4.21:lem-probe-mass}), applied with $i:=k$ and the class $\mathfrak{Q}$ of (a).
The first three paragraphs below verify the hypotheses of that lemma with the data
\eqref{4.21:eq-old-anchor-proof-a}; the last one draws the conclusion.

\smallskip\noindent\emph{The per-cube bound on terms of one level.} Let $\ell\le k$ and
$\square'\in\pi_\ell$. Consider first the boundary terms $T$ of level $\ell$ with
$A_T\ni\square'$. Hypothesis~(i) gives
$\|T\|e^{w(T)}\le A_{\flat}c_Te^{-\sigma_{*}d_\ell(A_T)}$; the bound on labels from which it
comes distinguishes three cases, which compare decay rates through the bound on the size of a
stored term in the stored-term format and through its filing level. By the definitions
of $\mathsf{e}$ and $\mathsf{f}$ in Proposition~\ref{4.21:prop-preparation-machine}, and since
$N\ge1$, one has $\mathsf{e}(T)+\mathsf{f}(T)\le2d_k([S_T]_k)$, where $S_T=[A_T]_m$ is the set
of that proposition; and $d_k([S_T]_k)\le d_\ell(A_T)$ by the comparison of the size of a
label with the sizes of its hulls at coarser levels. Hence
$e^{\mathsf{e}(T)+\mathsf{f}(T)}\le e^{2d_\ell(A_T)}$. Since
$\mathfrak{w}(T)\le\|T\|e^{w(T)}$ by hypothesis~(i), the factors combine to
\begin{equation*}
W(T)\,e^{d_\ell(A_T)}\ \le\ A_{\flat}c_T\,e^{-(\sigma_{*}-3)d_\ell(A_T)},
\end{equation*}
and $\sigma_{*}-3\ge c_d+1$ by the condition on $\kappa$ at the preparation site among the
conditions \eqref{4.21:eq-format-propagation-a}. Summing over the labels through $\square'$,
\begin{equation}\label{4.21:eq-old-anchor-proof-4}
\sum_{T:\ A_T\ni\square'}W(T)\,e^{d_\ell(A_T)}\ \le\ K(d)\,A_{\flat}B_0 ,
\end{equation}
both for the terms of $\mathbf{B}^{(\ell)}$ and for those of $\mathbf{B}'^{(\ell)}$.

For the frozen chain terms the estimate leading to \eqref{4.21:eq-old-anchor-proof-4}, with
the factor $e^{d_\ell(A)}$ of the label $A$, holds in the norm of \cite[(1.26)]{B-CMP116},
with $A_{\flat}C_0^{\sharp}$ per stage and per chain in place of $K(d)A_{\flat}B_0$, the stages
and the chains being counted as in Lemma~\ref{4.21:lem-wrapped-cube-bound}. Indeed, in the
accounting of decay rates for frozen chain terms at the preparation site, with $k_1$ the level
attached to the frozen chain in that accounting, a frozen term considered at a level $\ell<k_1$
is filed above $\ell$, by the lemma on the filing levels of frozen terms, and at $\ell=k_1$ it is
estimated as a boundary term of rate $\kappa$; in both cases the surplus of the decay rate is
at least $\sigma_{*}$. The members of the class covered by hypothesis~(iii) contribute
$A_{\flat}b_{\delta}$.

Finally consider the chain terms of the current operation $\mathbf{R}_k$ filed at $\ell$,
with label $S\ni\square'$. By hypothesis~(ii), in the norm $\|\cdot\|^{A}$, each of the $n$
stages performed contributes at most $C_0^{\sharp}$, so that their contribution is at most
$nC_0^{\sharp}\le N(g_k)C_0^{\sharp}$. At $\ell=k$ this bound includes, by hypothesis~(ii),
the factor $e^{d_k(S)/N}$, which is the factor $e^{\mathsf{e}(c)}$ of such a member. At
$\ell<k$ one has $\mathsf{e}(c)=0$: such an $S$ lies in one treated component, so that the
member is confined. And $\mathsf{f}(c)=0$ for every term of the current chain, by the
definition of $\mathsf{f}$ in Proposition~\ref{4.21:prop-preparation-machine}.

Collecting these contributions, the per-cube bound on terms of one level in
\eqref{4.21:eq-wrapped-terms-a} holds with $\mathsf{b}^{P}=A_{\flat}\mathsf{b}_k$, where
$\mathsf{b}_k$ is the per-cube function of Lemma~\ref{4.21:lem-wrapped-cube-bound} at level
$k$. The growth law in \eqref{4.21:eq-wrapped-terms-a} for $\mathsf{b}^{P}$ is the growth law
of $\mathsf{b}_k$, fixed together with $\mathsf{b}_k$ and its constants $C_{\mathsf{b}}$,
$p_{\mathsf{b}}$ in the format of the stored terms, multiplied on both sides by $A_{\flat}$;
its constants are therefore those of $\mathsf{b}_k$.

\smallskip\noindent\emph{The per-cube bound with a spare rate, $\varsigma_0=1$.} Let
$\ell<n_0$. For boundary terms and frozen terms the required bound is
\eqref{4.21:eq-old-anchor-proof-4}, whose left-hand side carries the factor
$e^{\varsigma_0 d_\ell(A_T)}$ with $\varsigma_0=1$; for the members of the class of
hypothesis~(iii) it is that hypothesis, whose left-hand side carries the same factor. No term
of the current chain is finer than an old probe. Indeed, the filing levels of such terms are
at least $h+1$. On the other hand $n_0\le h$ if
$\tilde{\square}_0$ meets a strip of a treated component, by parts (i1) and (g) of the lemma
on the geometry of zones, strips and probes, and $n_0\le h$ if $\tilde{\square}_0$ meets a
cell of level less than $h$ of a treated component, by part (i4) of that lemma. If
$\tilde{\square}_0$ meets an untreated component, then $\tilde{\square}_0$ lies at distance at
least $\check{R}_k-5$ from $Z$, and $Z$ contains every $S$ filed below $k$; hence no such label
$S$ contains a cube of $\pi_\ell$ inside $\tilde{\square}_0$, and these terms do not enter the
bound with a spare rate.

\smallskip\noindent\emph{The depth datum.} Put $\varrho(\tilde{\square}_0)$ and $N_{*}$ as in
\eqref{4.21:eq-old-anchor-proof-a}; then $\varrho(\tilde{\square}_0)\ge1$, and we write
$\varrho$ for $\varrho(\tilde{\square}_0)$. By part (i3) of the geometry lemma, every point of
$\square_0$ lies in a zone at most $n_0+1$.

Treated case (the first two kinds of old blocks in the definition of $\mathfrak{O}$). The
collar $R_n$ lies in zones at least $h$ and outside $Z''_{h+1}$. A path from $\square_0$ to
$R_n$ crosses the whole zones $n_0+2,\dots,h-1$, which lie in
$\Omega_h^{c}\subset Z_h$, and then, after its last point in $Z_h$, the stretch to
$(Z''_{h+1})^{c}$. By part (h) of the geometry lemma that stretch is
$3L\check{R}_{h+1}\ge3\check{R}_h$ level-$h$ units long and runs through cells of level at
most $h$. If $\square_0$ itself is not contained in $Z_h$, then $n_0=h$, and
$\tilde{\square}_0$, whose side is at most $5$, meets $Z_h$ or a strip reaching at most
$\frac12\check{R}_h$ beyond $Z_h$. Hence
\begin{equation}\label{4.21:eq-old-anchor-proof-5}
d^{\wedge}(\square_0,R_n)\ \ge\ w\Bigl[(h-n_0-2)_{+}+\tfrac52\check{R}_h-5\Bigr].
\end{equation}
Untreated case. The zones $n_0+2,\dots,k-1$ lie in $\Lambda_k^{c}$, and the gap between
$\tilde{\square}_0$ and $Z$, of width at least $\check{R}_k-5$ by the preceding paragraph, lies
in $\Lambda_k$, so that
\begin{equation}\label{4.21:eq-old-anchor-proof-6}
d^{\wedge}(\square_0,R_n)\ \ge\ w\bigl[(k-n_0-2)_{+}+\check{R}_k-5\bigr].
\end{equation}

First condition of the depth datum. In the treated case
\begin{equation*}
k-n_0+1=N+(h-n_0)+1\ \le\ N+3+(h-n_0-2)_{+}\ \le\ N_{*}+\varrho ,
\end{equation*}
the first inequality because $h-n_0-2\le(h-n_0-2)_{+}$, the second by
\eqref{4.21:eq-old-anchor-proof-5} and the definition of $\varrho$. In the untreated case
$k-n_0+1\le 3+(k-n_0-2)_{+}$, and \eqref{4.21:eq-old-anchor-proof-6} bounds the right-hand
side by $N_{*}+\varrho$ in the same way.

Second condition of the depth datum. Let $\tilde{\square}_0$ meet the strip $Y^{\sharp}$ of a
live pocket $(Y,m)$. Then $R_n$ lies outside the $\frac52\check{R}_m$-neighbourhood of
$Y^{\sharp}$: for $m<h$ because the cells of $R_n$ have level at least $h>m$, and by part (c)
of the geometry lemma cells of level above $m$ lie outside that neighbourhood; for $m=h$ by
part (h); in the untreated case because that
neighbourhood lies in the component of $Y$, by part (c). Every point of $\square_0$ lies
within $5$ level-$m$ units of $Y^{\sharp}$, by part (i1). Therefore
$d^{\wedge}(\square_0,R_n)\ge w(\frac52\check{R}_m-5)$, that is,
$\frac52\check{R}_m\le\varrho+4$, and
\begin{equation*}
\check{R}_m\ \le\ \tfrac25(\varrho+4)\ \le\ 2\varrho+2
\qquad(\varrho\ge1).
\end{equation*}

In all cases $\varrho\ge\check{R}_k-4$, that is, \eqref{4.21:eq-old-anchor-proof-1} holds for
every old block. This proves (a) and (b).

\smallskip\noindent\emph{Conclusion.} By \eqref{4.21:eq-old-anchor-proof-3} and part~(iv) of
Lemma~\ref{4.21:lem-probe-mass} with $\beta=1$, applied with the data
\eqref{4.21:eq-old-anchor-proof-a},
\begin{equation*}
\mathfrak{S}(\square_0)\ \le\ e^{\varrho}\,C^{P}(1)\,(N+3)\,\mathsf{b}^{P}(N+2),
\end{equation*}
where $C^{P}(1)$ denotes the constant $C_{\mathsf{a}}(1)$ of that lemma evaluated at
$(\varsigma_0,\mathsf{b})=(1,\mathsf{b}^{P})$. By the comparison of the stage count and of the
logarithmic scale variable with the region size, together with $\varrho\ge\check{R}_k-4$,
\begin{equation*}
\begin{gathered}
N\le C_N2^{r_{\mathrm{pr}}}(\varrho+4),\\
\mathsf{T}_k^{r_{\mathrm{pr}}}\le2^{r_{\mathrm{pr}}}(\varrho+4),\\
(\mathsf{T}_k+c_s)^{r_{\mathrm{pr}}}\le\mathsf{T}_k^{r_{\mathrm{pr}}}(1+c_s)^{r_{\mathrm{pr}}}.
\end{gathered}
\end{equation*}
Consequently $(N+3)\mathsf{b}^{P}(N+2)\le\Pi_0(\varrho)$, where $\Pi_0$ is a function on
$[1,\infty[$, a polynomial in $\varrho$ times a polylogarithmic factor, whose coefficients
depend only on
\begin{equation*}
d,\ L,\ \lambda,\ c_2,\ r_{\mathrm{pr}},\ C_N,\ C_0^{\sharp},\ B_0,\ b_{\delta},\
A_{\flat},\ K(d).
\end{equation*}
Put
\begin{equation}\label{4.21:eq-old-anchor-proof-2}
C^{P}_{\mathsf{a}}\ :=\ C^{P}(1)\cdot\sup_{\varrho\ge1}\Pi_0(\varrho)\,e^{2-\varrho};
\end{equation}
it is finite, since $e^{-\varrho}$ decays faster than $\Pi_0$ grows. Writing
$e^{\varrho}=e^{2-\varrho}e^{2(\varrho-1)}$ we obtain
$\mathfrak{S}(\square_0)\le C^{P}_{\mathsf{a}}e^{2(\varrho-1)}$. By the definition of
$\varrho$,
\begin{equation*}
e^{2(\varrho-1)}=e^{(2/w)\,d^{\wedge}(\square_0,R_n)}
\ \le\ e^{\frac18\delta_P d^{\wedge}(\square_0,R_n)},
\end{equation*}
since $\frac18\delta_P w\ge2$ by the condition on $w$ among the format conditions of the
preparation site, and $d^{\wedge}\ge0$. This is (c).
\end{proof}

\begin{proof}[Proof of Proposition~\ref{4.21:prop-preparation-machine}: groups, anchored totals and block sums]
\label{4.21:prf-group-totals-proof}
In this part of the proof we establish the totals of clause~(e) of
Proposition~\ref{4.21:prop-preparation-machine} for the groups $\hat e$. We then prove two sums over
blocks, which are used when the anchored totals are summed over their anchors. The hypotheses
are those of the proposition. Of the conditions on the constants, three are used below:
\begin{itemize}
\item the smallness condition $2e^2\hat b_t^2N^{\flat}\mathsf{x}\le 1$ of
\eqref{4.21:eq-preparation-constants-a};
\item the inequality
$6N c_{\mathsf{p}}V'(K_2^{\flat}\vartheta+K_Q/r_A)\le 1$, which is among the conditions listed in
Lemma~\ref{4.21:lem-coupling-interval}; here $N$ is the number that enters $\mathsf{e}(c)$ in
Proposition~\ref{4.21:prop-preparation-machine}, not the rank of the gauge group;
\item the tail condition on $\tau_d(\check{R}_k)$ under which clause~(c) of
Proposition~\ref{4.21:prop-preparation-machine} is stated; it is used below in the form
$4\tau_d(\check{R}_k)\le1$.
\end{itemize}

We fix a wrapped member $c$ in format. Its label $A_c$, its core $S_c$ (the label $A_c$ for a chain
term, the hull $[A_T]_m$ at the filing level $m$ for a leaf $T$), its reading set $X_c$, its incident
blocks
$\mathcal{A}_c=\mathcal{A}_c^{\mathrm{pl}}\sqcup\mathcal{A}_c^{\mathrm{isl}}$, the corridors
$\operatorname{Cor}(\square)$, the hulls $\mathsf{d}^{\mathrm{rel}}(\hat e)$, the block sum $V'$, the
old blocks $\mathfrak{O}$ and the numbers $c_{\mathsf{p}}$, $\mathsf{p}^{\mathrm{rel}}(c)$,
$\mathsf{e}(c)$ are those of the statement of the proposition. The groups
$\hat e=(c,A,Q,\hat{\mathsf{y}})$ are those of its clause~(a$'$). There
\begin{equation*}
\nu(\hat e)=\mathfrak{N}^{\natural}(c)\prod_{\square\in A}K_A\vartheta(\square)
\prod_{q\in Q}\frac{2}{r_A}\,e^{-(\kappa_1-1)|\hat{\mathsf{y}}|},
\qquad
|\mathsf{b}(\hat e)|\le N^{\flat}|A|+|\hat{\mathsf{y}}|+|Q|+|\mathcal{S}_j(c)|,
\end{equation*}
where $\mathsf{b}(\hat e)$ is the bond set of the group and $\mathcal{S}_j(c)$ is the strong set of
level $j$ of $S_c$, with $j=j_c$ the level of the label $A_c$. We also write $\mathcal{Q}^w_c$ for
the rim cubes of $S_c$. Following clause~(e), we put
\begin{equation*}
y_{\square}:=K_2^{\flat}\vartheta(\square)e^{(a_k+1)|\operatorname{Cor}(\square)|},
\qquad
\Sigma_c:=\bigl(K_2^{\flat}\vartheta+K_Q/r_A\bigr)V'\mathsf{p}^{\mathrm{rel}}(c).
\end{equation*}
Here the per-block weight is $\vartheta(\square)=\vartheta\,e^{-\frac12\delta_Pd^{\wedge}(\square,R_n)}$,
so that
\begin{equation*}
y_{\square}=K_2^{\flat}\vartheta\,e^{-\frac12\delta_Pd^{\wedge}(\square,R_n)}
e^{(a_k+1)|\operatorname{Cor}(\square)|}.
\end{equation*}
We write $w(\cdot)$ for the weight of a term in its format, and $w(\hat e)$ for the weight of the
group filed at the level of its hull; the weight is always written with its argument, and is not to
be confused with the number $w=M/M_1$ that enters the scaled distance $d^{\wedge}$ and appears in
Step~3 only as a factor in explicit products.

\emph{Step 1: the weight of one group.}
Let $m:=m(\mathsf{d}^{\mathrm{rel}}(\hat e))$ be the filing level of the hull. Then
$[\mathsf{d}^{\mathrm{rel}}(\hat e)]_m\supset[S_c]_m$, and both sets belong to $\mathbf{D}_m$. We
claim that the $\pi_m$-cubes of $[\mathsf{d}^{\mathrm{rel}}(\hat e)]_m$ that are not in
$[S_c]_m$ number at most
\begin{equation*}
\sum_{\square\in A}\bigl(N^{\flat}+|\operatorname{Cor}(\square)|\bigr)+|\hat{\mathsf{y}}|.
\end{equation*}
Indeed, $\mathsf{d}^{\mathrm{rel}}(\hat e)$ is the union of $S_c$, of the sets
$\hat{\mathsf{K}}_{\square}\cup\operatorname{Cor}(\square)$ for $\square\in A$, and of the members of
$\hat{\mathsf{y}}$. We treat the three kinds of contribution in turn.
\begin{itemize}
\item By the hull statement (\(\flat 5\)) of Lemma~\ref{4.21:lem-grain-properties}, every member of
$\pi^{\flat}$ contained in $\mathsf{d}^{\mathrm{rel}}(\hat e)$ has level at most $m$, and so lies in a
single $\pi_m$-cube. By (\(\flat 4\)) of the same lemma, $\hat{\mathsf{K}}_{\square}$ has at most
$N^{\flat}$ members. The family $\hat{\mathsf{y}}$ has $|\hat{\mathsf{y}}|$ members.
\item By clause~(b) of Proposition~\ref{4.21:prop-preparation-machine}, the corridor
$\operatorname{Cor}(\square)$ adds at most $|\operatorname{Cor}(\square)|$ cubes.
\item Finally, by the hull property of the geometry of the labels, adjoining to $S_c$ a set that lies
in a single $\pi_m$-cube adds at most one cube to the $\pi_m$-hull; applied to the members of the
sets $\hat{\mathsf{K}}_{\square}$ and of $\hat{\mathsf{y}}$, this completes the count.
\end{itemize}
By the comparison of the format rates across levels, $a^{(n+1)}_m\le a_k+1$. By the weight hypothesis
on labels, re-filing $S_c$ at the coarser level $m$ does not raise its weight. Together with the count
of added cubes, these two facts give
\begin{equation}\label{4.21:eq-group-totals-proof-1}
w(\hat e)\le w(c)+(a_k+1)\Bigl[\sum_{\square\in A}\bigl(N^{\flat}+|\operatorname{Cor}(\square)|\bigr)
+|\hat{\mathsf{y}}|\Bigr].
\end{equation}
Put
\begin{equation*}
\tilde\nu(\hat e):=2\nu(\hat e)e^{w(\hat e)+c_J+c_{\flat}|\mathsf{b}(\hat e)|}.
\end{equation*}
We insert the expression for $\nu(\hat e)$, the bound on $|\mathsf{b}(\hat e)|$ and
\eqref{4.21:eq-group-totals-proof-1}, and collect the factors block by block, cube by cube and grain
by grain. This gives
\begin{align*}
\tilde\nu(\hat e)\le{}&2\,\mathfrak{N}^{\natural}(c)e^{w(c)}e^{c_J}e^{c_{\flat}|\mathcal{S}_j(c)|}
\prod_{\square\in A}\Bigl[K_A\vartheta(\square)e^{(a_k+1)(N^{\flat}+|\operatorname{Cor}(\square)|)}
e^{c_{\flat}N^{\flat}}\Bigr]\\
&\times\prod_{q\in Q}\Bigl[\frac{2}{r_A}e^{c_{\flat}}\Bigr]\cdot
\Bigl[e^{-(\kappa_1-1)}e^{a_k+1}e^{c_{\flat}}\Bigr]^{|\hat{\mathsf{y}}|}.
\end{align*}
Four comparisons now reduce this bound.
\begin{itemize}
\item The constants of Notation~\ref{4.21:not-preparation-constants} satisfy
$K_Ae^{(a_k+1)N^{\flat}}e^{c_{\flat}N^{\flat}}\le e^{-1}K_2^{\flat}$. Hence each factor of the first
product is at most $e^{-1}y_{\square}$.
\item We have $(2/r_A)e^{c_{\flat}}=K_Q/r_A$.
\item By the definition of $\mathsf{x}$ in \eqref{4.21:eq-preparation-constants-a}, and since
$3^dc_{\flat}\ge c_{\flat}$, we have $e^{-(\kappa_1-1)}e^{a_k+1}e^{c_{\flat}}\le\mathsf{x}$.
\item The weight $\mathfrak{w}(c)$ contains the strong weight $e^{\kappa_A|\mathcal{S}_j(c)|}$ of
level $j$, which pays the factor $e^{c_{\flat}|\mathcal{S}_j(c)|}$:
\begin{equation*}
\mathfrak{N}^{\natural}(c)e^{w(c)}e^{c_{\flat}|\mathcal{S}_j(c)|}
\le\mathfrak{w}(c)\,e^{-(\kappa_A-c_{\flat})|\mathcal{S}_j(c)|}.
\end{equation*}
\end{itemize}
Therefore
\begin{equation}\label{4.21:eq-group-totals-proof-2}
\tilde\nu(\hat e)\le 2\mathfrak{w}(c)e^{c_J}\,e^{-(\kappa_A-c_{\flat})|\mathcal{S}_j(c)|}
\prod_{\square\in A}\bigl[e^{-1}y_{\square}\bigr]\prod_{q\in Q}\bigl[K_Q/r_A\bigr]\,
\mathsf{x}^{|\hat{\mathsf{y}}|}.
\end{equation}
On $\mathcal{A}_c^{\mathrm{pl}}$, where $\operatorname{Cor}(\square)=\emptyset$, we have
$y_{\square}=K_2^{\flat}\vartheta e^{-\frac12\delta_Pd^{\wedge}(\square,R_n)}$. On the islands, the
corridor has fewer than $d(x_{\square}+1)$ cubes, and clause~(b) of
Proposition~\ref{4.21:prop-preparation-machine} gives
\begin{equation*}
(a_k+1)|\operatorname{Cor}(\square)|\le(a_k+1)d(x_{\square}+1)
\le\tfrac14\delta_Pd^{\wedge}(\square,R_n),
\end{equation*}
so that $y_{\square}\le K_2^{\flat}\vartheta e^{-\frac14\delta_Pd^{\wedge}(\square,R_n)}$.

\emph{Step 2: summation over the families, then over $(A,Q)$.}
Under the smallness condition on $\mathsf{x}$ in \eqref{4.21:eq-preparation-constants-a},
clause~(iv) of Lemma~\ref{4.21:lem-grain-resummation} gives
$\sum_{\hat{\mathsf{y}}}\mathsf{x}^{|\hat{\mathsf{y}}|}\le e^{|A|}$. The factor $e^{|A|}$ cancels the
factors $e^{-1}$ in \eqref{4.21:eq-group-totals-proof-2}. So at fixed $(A,Q)$ the groups sum to at
most
\begin{equation}\label{4.21:eq-group-totals-proof-3}
2\mathfrak{w}(c)e^{c_J}e^{-(\kappa_A-c_{\flat})|\mathcal{S}_j(c)|}
\prod_{\square\in A}y_{\square}\prod_{q\in Q}\bigl[K_Q/r_A\bigr].
\end{equation}
Next we sum over $A\subset\mathcal{A}_c$ and $Q\subset\mathcal{Q}^w_c$. By the product formula,
\begin{equation*}
\sum_{A\subset\mathcal{A}_c}\prod_{\square\in A}y_{\square}
=\prod_{\square\in\mathcal{A}_c}(1+y_{\square})
\le\exp\Bigl(\sum_{\square\in\mathcal{A}_c}y_{\square}\Bigr),
\end{equation*}
and similarly
\begin{equation*}
\sum_{Q\subset\mathcal{Q}^w_c}(K_Q/r_A)^{|Q|}\le\exp\bigl(|\mathcal{Q}^w_c|K_Q/r_A\bigr).
\end{equation*}
Two bounds control the exponents.
\begin{itemize}
\item By the expression for $y_{\square}$ and clause~(c) of
Proposition~\ref{4.21:prop-preparation-machine},
\begin{equation*}
\sum_{\square\in\mathcal{A}_c}y_{\square}
=K_2^{\flat}\vartheta\sum_{\square\in\mathcal{A}_c}e^{-\frac12\delta_Pd^{\wedge}(\square,R_n)}
e^{(a_k+1)|\operatorname{Cor}(\square)|}
\le K_2^{\flat}\vartheta V'\mathsf{p}^{\mathrm{rel}}(c).
\end{equation*}
\item A rim cube has $d^{\wedge}=0$, as in the group-total bound at the preparation site. Hence
$|\mathcal{Q}^w_c|\le V'\mathsf{p}^{\mathrm{rel}}(c)$.
\end{itemize}
The sum of the two exponents is therefore at most $\Sigma_c$. With these three inputs (the product
formula, the bound on $\sum_{\mathcal{A}_c}y_{\square}$ from clause~(c), and the bound on
$|\mathcal{Q}^w_c|$), the group-total bound at the preparation site holds for the groups $\hat e$,
with the exponent $\Sigma_c$:
\begin{equation}\label{4.21:eq-group-totals-proof-4}
\sum_{\hat e\text{ of }c}\tilde\nu(\hat e)\le 2\mathfrak{w}(c)e^{c_J}
\Bigl[\bigl(e^{\Sigma_c}-1\bigr)+\varepsilon_Se^{\Sigma_c}\Bigr].
\end{equation}
Assume $6Nc_{\mathsf{p}}V'(K_2^{\flat}\vartheta+K_Q/r_A)\le1$. Then
$\Sigma_c\le\mathsf{p}^{\mathrm{rel}}(c)/(6Nc_{\mathsf{p}})$. Moreover
$\mathsf{p}^{\mathrm{rel}}(c)/(6Nc_{\mathsf{p}})\le\frac12+\mathsf{e}(c)$. If $c$ is confined, the
left side is $2/(6Nc_{\mathsf{p}})\le\frac12=\frac12+\mathsf{e}(c)$. Otherwise we use $N\ge1$ and
$\check{R}_k\ge1$, and obtain
\begin{equation*}
\frac{\mathsf{p}^{\mathrm{rel}}(c)}{6Nc_{\mathsf{p}}}
=\frac{2+d_k([S_c]_k)/\check{R}_k}{6N}
\le\frac13+\frac{d_k([S_c]_k)}{6N\check{R}_k}
\le\frac12+\frac{d_k([S_c]_k)}{N}=\frac12+\mathsf{e}(c).
\end{equation*}
Thus $\Sigma_c\le\frac12+\mathsf{e}(c)$.

\emph{Anchored totals, with position.}
Let $\square_0\in A$ and let $\mathsf{E}$ be a closed set. We say that $(\square_0,\mathsf{E})$
\emph{holds for} $\hat e$ if one of two conditions is met:
\begin{itemize}
\item $\mathsf{E}$ meets $\hat{\mathsf{K}}_{\square_0}\cup\operatorname{Cor}(\square_0)$; or
\item a connected component of the family $\hat{\mathsf{y}}$ (an animal of $\hat{\mathsf{y}}$, in the
sense of clause~(iv) of Lemma~\ref{4.21:lem-grain-resummation}) that shares a wall with
$\hat{\mathsf{K}}_{\square_0}$ meets or touches $\mathsf{E}$.
\end{itemize}
Then
\begin{equation}\label{4.21:eq-group-totals-proof-a}
\sum\bigl\{\tilde\nu(\hat e):\ \hat e\text{ of }c,\ A\ni\square_0,\ (\square_0,\mathsf{E})
\text{ holds for }\hat e\bigr\}
\le 2\mathfrak{w}(c)e^{c_J}\,y_{\square_0}\,e^{\Sigma_c}\,
e^{-(\operatorname{dist}_k(\mathsf{E},\square_0)-103\check{R}_k)_+}.
\end{equation}
This is the anchored form of the group-total bound at the preparation site, for the groups $\hat e$;
the position factor arises from the two alternatives in the definition of $(\square_0,\mathsf{E})$,
as follows.

Suppose first that $\mathsf{E}$ meets $\hat{\mathsf{K}}_{\square_0}\cup\operatorname{Cor}(\square_0)$.
By the corridor construction in the geometry of the labels, $\operatorname{Cor}(\square_0)$ lies in
the box spanned by $\tilde{\square}_0$ and a cube touching the strip. By (\(\flat 4\)) of
Lemma~\ref{4.21:lem-grain-properties}, $\hat{\mathsf{K}}_{\square_0}$ lies within $6$ level-$k$
units of $\square_0$. Since the strip met by $\tilde{\square}_0$ has sides at most
$101\check{R}_k$ level-$k$ units and $\tilde{\square}_0$ has side at most $5$, it follows that
$\operatorname{dist}_k(\mathsf{E},\square_0)\le103\check{R}_k$, and the
last factor of \eqref{4.21:eq-group-totals-proof-a} equals $1$.

Suppose next that an animal of $\hat{\mathsf{y}}$ sharing a wall with $\hat{\mathsf{K}}_{\square_0}$
meets or touches $\mathsf{E}$. The second statement of clause~(iv) of
Lemma~\ref{4.21:lem-grain-resummation} provides the factor
$e^{-(\operatorname{dist}_k(\mathsf{E},\hat{\mathsf{K}}_{\square_0})-1)_+}$ in the sum over
$\hat{\mathsf{y}}$. Since $\hat{\mathsf{K}}_{\square_0}$ lies within $6$ of $\square_0$, we have
\begin{equation*}
e^{-(\operatorname{dist}_k(\mathsf{E},\hat{\mathsf{K}}_{\square_0})-1)_+}
\le e^{-(\operatorname{dist}_k(\mathsf{E},\square_0)-7)_+}
\le e^{-(\operatorname{dist}_k(\mathsf{E},\square_0)-103\check{R}_k)_+},
\end{equation*}
the last step because $7\le103\check{R}_k$.

The two anchored forms of the group-total bound at the preparation site with a cube $q_0\in Q$,
respectively $q_0\in\mathcal{S}_j(c)$, in place of the block $\square_0$ hold for the groups $\hat e$
word for word, with the same proof.

\emph{Step 3: two sums over blocks.}
For every point or cube $x$,
\begin{equation}\label{4.21:eq-group-totals-proof-b}
\begin{gathered}
\sum_{\square_0\in\mathfrak{O}}e^{-\frac18\delta_Pd^{\wedge}(\square_0,R_n)}
e^{-(\operatorname{dist}_k(x,\square_0)-103\check{R}_k)_+}\le4\tau_d(\check{R}_k)V'\le V',\\
\sum_{\square_0}e^{-\frac18\delta_Pd^{\wedge}(\square_0,R_n)}
e^{-(\operatorname{dist}_k(x,\square_0)-103\check{R}_k)_+}\le c_{\bigstar}V',
\end{gathered}
\end{equation}
where the second sum runs over all blocks. For $u\ge0$, let $\mathfrak{B}_u$ be the set of blocks
$\square_0$ of the class summed (old blocks in the first inequality, all blocks in the second) with
$\operatorname{dist}_k(x,\square_0)<(u+1)\check{R}_k$. These blocks lie in a box of side at most
$2(u+1)\check{R}_k+4$.

Let $\bar c$ range over the cells indexing the pieces $R(\bar c)$ of $R_n$, as in the block-sum part
of the proof. We apply the bound
$e^{-\beta d^{\wedge}(\square,R_n)}\le\sum_{\bar c}e^{-\beta d^{\wedge}(\square,R(\bar c))}$ and then
\eqref{4.21:eq-block-sum-proof-a}, both at $\beta=\frac18\delta_P$. This gives
\begin{equation*}
\sum_{\square_0\in\mathfrak{B}_u}e^{-\frac18\delta_Pd^{\wedge}(\square_0,R_n)}
\le\sum_{\bar c}\sum_{\square_0\in\mathfrak{B}_u}e^{-\frac18\delta_Pd^{\wedge}(\square_0,R(\bar c))}
\le V'\sum_{\bar c}e^{-\frac1{16}\delta_Pd^{\wedge}(\mathfrak{B}_u,R(\bar c))}.
\end{equation*}
We count the cells $\bar c$ by the counting property of cells near a box in the geometry of the
labels, and we use $\frac1{16}\delta_P\,w\ge1$, with $w=M/M_1$; this follows from the condition
$\frac18\delta_P\,w\ge2$ used at the end of the old-anchor part of the proof of
Proposition~\ref{4.21:prop-preparation-machine}. The cells $\bar c$ with
$\operatorname{dist}_k(\mathfrak{B}_u,\bar c)<\check{R}_k$ number at most $(2u+6)^d$. For $l\ge1$,
those with
$\operatorname{dist}_k(\mathfrak{B}_u,\bar c)\in[l\check{R}_k,(l+1)\check{R}_k)$ number at most
$(2l+2u+6)^d$, and each carries a factor at most $e^{-l\check{R}_k}$.

\emph{The first inequality.} By the depth bound proved in the old-anchor part of the proof,
\eqref{4.21:eq-old-anchor-proof-a}, every old block satisfies
$d^{\wedge}(\square_0,R_n)\ge(\check{R}_k-5)\,w$, with $w=M/M_1$. So the near cells $\bar c$ carry a
factor at most
$e^{-(\check{R}_k-5)}$. Their contribution $(2u+6)^de^{-(\check{R}_k-5)}V'$ is at most the term $l=1$
of the series
\begin{equation*}
V'\sum_{l\ge1}(2l+2u+6)^de^{-(l\check{R}_k-5)}.
\end{equation*}
The contribution of the far cells is at most this series as well. Hence
\begin{equation*}
\sum_{\square_0\in\mathfrak{B}_u}e^{-\frac18\delta_Pd^{\wedge}(\square_0,R_n)}
\le 2V'\sum_{l\ge1}(2l+2u+6)^de^{-(l\check{R}_k-5)}.
\end{equation*}
At $u=103$, by the definition of the tail sum $\tau_d(\cdot)$, this is at most
$2\tau_d(\check{R}_k)V'$. On $\mathfrak{B}_{103}$ the second factor of the summand is at most $1$.

Now let $u\ge104$ and put $t:=u-103$. Then $2l+2u+6\le(2l+215)(2t+1)$, since the difference of the
two sides is $4lt+428t+3>0$. On $\mathfrak{B}_u\setminus\mathfrak{B}_{u-1}$ the second factor of the
summand is at most $e^{-t\check{R}_k}$. By the definition of $\tau_d(\cdot)$,
$\sum_{l\ge1}(2l+215)^de^{-(l\check{R}_k-5)}\le\tau_d(\check{R}_k)$ and
\begin{equation*}
\sum_{t\ge1}(2t+1)^de^{-t\check{R}_k}\le\tau_d(\check{R}_k)\le1,
\end{equation*}
the last inequality by the tail condition $4\tau_d(\check{R}_k)\le1$.
Hence the blocks outside $\mathfrak{B}_{103}$ contribute at most
\begin{equation*}
2V'\sum_{t\ge1}(2t+1)^de^{-t\check{R}_k}\sum_{l\ge1}(2l+215)^de^{-(l\check{R}_k-5)}
\le2\tau_d(\check{R}_k)^2V'.
\end{equation*}
In total the left side of the first inequality is at most
\begin{equation*}
2\tau_d(\check{R}_k)V'+2\tau_d(\check{R}_k)^2V'\le4\tau_d(\check{R}_k)V',
\end{equation*}
since $\tau_d(\check{R}_k)\le1$. Finally, $4\tau_d(\check{R}_k)V'\le V'$ by the tail condition on
$\tau_d(\check{R}_k)$ under which clause~(c) of Proposition~\ref{4.21:prop-preparation-machine} is
stated.

\emph{The second inequality.} Now the near cells carry a factor at most $1$, so
\begin{equation*}
\sum_{\square_0\in\mathfrak{B}_u}e^{-\frac18\delta_Pd^{\wedge}(\square_0,R_n)}
\le V'\Bigl[(2u+6)^d+\sum_{l\ge1}(2l+2u+6)^de^{-l\check{R}_k}\Bigr].
\end{equation*}
At $u=103$, by the definition of $\tau_d(\cdot)$, this is at most
$\bigl(212^d+\tau_d(\check{R}_k)\bigr)V'$.

For $u=103+t$ with $t\ge1$, the blocks of $\mathfrak{B}_u\setminus\mathfrak{B}_{u-1}$ carry the factor
$e^{-t\check{R}_k}$. By the definition of $\tau_d(\cdot)$,
$\sum_{t\ge1}(2t+212)^de^{-t\check{R}_k}\le\tau_d(\check{R}_k)$. As before, using
$2l+2t+212\le(2l+215)(2t+1)$, the double sum
\begin{equation*}
\sum_{t\ge1}\sum_{l\ge1}(2l+2t+212)^de^{-l\check{R}_k}e^{-t\check{R}_k}
\end{equation*}
is at most $\tau_d(\check{R}_k)^2$. The total is therefore at most
\begin{equation*}
\bigl(212^d+2\tau_d(\check{R}_k)+\tau_d(\check{R}_k)^2\bigr)V'\le(212^d+2)V'\le214^dV'.
\end{equation*}
Here the first inequality holds because $4\tau_d(\check{R}_k)\le1$ by the tail condition. The
second holds because $214^d-212^d\ge2$ for $d\ge1$. This is the second inequality of
\eqref{4.21:eq-group-totals-proof-b}, with the constant $c_{\bigstar}$ of
\eqref{4.21:eq-preparation-constants-a}.
\end{proof}

\begin{lemma}[Proof of the preparation proposition: sums over terms and hull totals]
\label{4.21:prf-term-sums-proof}
We use the notation of Proposition~\ref{4.21:prop-preparation-machine}. For a wrapped member $c$
of the $A$-class this means its label $A_c$ of level $j_c$, the set $S_c$, the region $X_c$, the
incident blocks $\mathcal{A}_c=\mathcal{A}_c^{\mathrm{pl}}\sqcup\mathcal{A}_c^{\mathrm{isl}}$, the
corridors $\operatorname{Cor}(\square)$, the domains $\mathsf{d}^{\mathrm{rel}}$, and the volume sum
$V'$. It also means the set $\mathfrak{O}$ of old blocks (all other blocks are \emph{window blocks})
and the quantities $c_{\mathsf{p}}$, $\mathsf{p}^{\mathrm{rel}}(c)$, $\mathsf{e}(c)$, $\mathsf{f}(c)$
and $\mathfrak{S}(\square_0)$.

Let $\hat e=(c,A,Q,\hat{\mathsf{y}})$ run through the groups of clause (a$'$) of
Proposition~\ref{4.21:prop-preparation-machine}. Let $\tilde\nu(\hat e)$, $y_\square$, $\Sigma_c$,
and the relation ``$(\square_0,\mathsf{E})$ holds for $\hat e$'' be as in
Lemma~\ref{4.21:prf-group-totals-proof}.

By $\mu$ and $\mathsf{M}_{*}$ we mean the two totals of the Stage~B proposition of the perturbative
step. In $\mu$ the groups are collected by a cube $q\in\mathsf{Q}$ of their bond set $\mathsf{b}(\hat e)$.
In $\mathsf{M}_{*}$ they are collected by a point $x$ of the $\pi_k$-hull $[\mathsf{d}^{\mathrm{rel}}(\hat e)]_k$
of their domain. Assume:
\begin{enumerate}
\item[(i)] the conditions $(F_{\kappa_1})^{\flat}$ of \eqref{4.21:eq-preparation-constants-a} and
$(Q^{+})^{A,\flat}$ of \eqref{4.21:eq-coupling-interval-a};
\item[(ii)] the window total bound, assumed for window blocks, which are blocks of zones at least
$h$: with
\begin{equation}\label{4.21:eq-term-sums-proof-1}
\mathfrak{T}^{\mathrm{win}}:=\sup_{\square_0\ \text{window}}
\sum\bigl\{\mathfrak{w}(c)e^{\mathsf{e}(c)}:\ \square_0\in\mathcal{A}_c\bigr\},
\end{equation}
one has
\begin{equation}\label{4.21:eq-term-sums-proof-2}
\mathfrak{T}^{\mathrm{win}}\le e^{\sigma_{*}}\bigl[K_S K(d)A_{\flat}(E_0+B_0+1)
+(5L)^d N C_0^{\sharp}\bigr];
\end{equation}
here $K_S$ denotes the constant of the window-anchor bound for stored terms of the perturbative
step;
\item[(iii)] the following statements, in the form in which they are proved for the perturbative
step and for the $A$-class:
  \begin{itemize}
  \item the rooted form of the window-anchor bound of the perturbative step;
  \item the anchored totals for the $A$-class at window anchors, in their two anchored forms;
  \item the Stage~B proposition, in the steps where it uses the domains of its elements (the
  weight, the connectedness of the union of domains, the rooting of polymers) and its clause (i)
  on the domain of the output;
  \item the two statements of the perturbative step on the domains of outputs.
  \end{itemize}
\end{enumerate}
Then the following hold.
\begin{enumerate}
\item[(I)] Let $q\in\mathsf{Q}$. The groups with $q\in\mathsf{b}(\hat e)$ and $q\in Q\cup\mathcal{S}_j(c)$
are bounded by the anchored totals for the $A$-class at window anchors, read with
$(V,\mathsf{p},K_2)$ replaced by $(c_{\bigstar}V',\mathsf{p}^{\mathrm{rel}},K_2^{\flat})$. The groups
with $q\in\mathsf{b}(\hat e)$ and $q\subset\hat{\mathsf{K}}_A\cup\bigcup\hat{\mathsf{y}}$ satisfy
\begin{equation}\label{4.21:eq-term-sums-proof-3}
\sum\tilde\nu(\hat e)\le 2e^{c_J+\frac12}K_2^{\flat}\vartheta V'\bigl[C^{P}_{\mathsf{a}}
+c_{\bigstar}\mathfrak{T}^{\mathrm{win}}\bigr].
\end{equation}
\item[(II)] Let $x$ be a point. The groups with $x\in[\mathsf{d}^{\mathrm{rel}}(\hat e)]_k$ for which
$(\square_0,\Delta_k(x))$ holds for some $\square_0\in A$ obey \eqref{4.21:eq-term-sums-proof-3}.
Consider the groups with $x\in[\mathsf{d}^{\mathrm{rel}}(\hat e)]_k$ and $\Delta_k(x)\cap S_c\ne\emptyset$.
Their contribution through window blocks, rim cubes and strong sets is bounded by the window-anchor
totals read at $(V'\mathsf{p}^{\mathrm{rel}},K_2^{\flat})$. Their contribution through old blocks
is, for chain terms filed at level $k$, a window-anchor contribution bounded in the same way, and
for the other members at most $2e^{c_J+\frac12}K_2^{\flat}\vartheta C^{P}_{\mathsf{a}}V'$.
\item[(III)] $\mu+\tilde{\mathfrak{l}}\le\frac12$ and
$4(\mathsf{M}_{*}+2^dV\tilde{\mathfrak{l}})\le\frac12$. The Stage~B proposition holds for the family
$\hat{\mathfrak{E}}$ of groups, with $\mathsf{d}$ replaced by $\mathsf{d}^{\mathrm{rel}}$. Its outputs
read the fields on $X$ together with strips that touch $X$.
\end{enumerate}
Together with Lemma~\ref{4.21:prf-group-totals-proof}, these statements complete the proof of
clause (e) of Proposition~\ref{4.21:prop-preparation-machine}.
\end{lemma}

\begin{proof}
\emph{Sums over terms at one block.} Let $\square_0$ be a block.

If $\square_0\in\mathfrak{O}$, the sum over the wrapped members incident to $\square_0$ is bounded by
clause (d) of Proposition~\ref{4.21:prop-preparation-machine}:
\begin{equation*}
\mathfrak{S}(\square_0)=\sum\bigl\{\mathfrak{w}(c)e^{\mathsf{e}(c)+\mathsf{f}(c)}:\ \square_0\in
\mathcal{A}_c\bigr\}\le C^{P}_{\mathsf{a}}e^{\frac18\delta_P d^{\wedge}(\square_0,R_n)}.
\end{equation*}

Now let $\square_0$ be a window block. Its probe $\tilde\square_0$ meets no strip $Y^{\sharp}$ of a
member of $\mathcal{P}_c$, since a block whose probe meets such a strip is old by the first clause
of the definition of $\mathfrak{O}$. Since $X_c=A_c\cup\bigcup Y^{\sharp}$, the condition
$\square_0\in\mathcal{A}_c$ therefore means $\tilde\square_0\cap A_c\ne\emptyset$. As
$A_c\subset S_c$, it follows that $\square_0\in\mathcal{A}_c^{\mathrm{pl}}$.

Every cell met by $\tilde\square_0$ has level at least $h$, because a cell of level below $h$ of a
treated component makes the block old by the second clause of that definition. Every such cell lies
in a treated component or in $\Lambda_k$, because an untreated component of $\Lambda_k^c$ makes the
block old by the third clause. In these regions the zones $h+1,\dots,k$ are made by the operation
$\mathbf{T}$ or by the present preparation; this is the zone statement for arrested attempts. The
thin strata of an older window lie in zones at most $h-1$, and a block whose probe meets one of them
is old by the second clause of the definition of $\mathfrak{O}$.

So window blocks are blocks of zones at least $h$. This is the setting in which the window-anchor
bound of the perturbative step and the summation of the relative Stage~B theorem are stated.

In these results the anchoring conditions $S\cap\tilde\square_0\ne\emptyset$, $S\cap q\ne\emptyset$
and $S\ni x$ are read as conditions on $S_c$. Their proofs read a term only through three things:
its domain in $\mathbf{D}_\ell$ or $\mathbb{D}_m$, its bound, and the side conditions on that domain.
These are supplied by clauses (f0), (f2) and (f4) of the relative format
\eqref{4.21:eq-relative-format-a}.

Clause (f4) places the anchor in $(Z_\ell^{\boxplus})^{(1)}$ instead of $Z_\ell^{\boxplus}$. Hence
each lower bound on $d_\ell$ used in the window-anchor bound decreases by at most $1$. This costs a
factor at most $e^{\sigma_{*}}$, which is included in the constant $K_{\flat}''$ of
\eqref{4.21:eq-preparation-constants-a}.

The quantity $\mathfrak{T}^{\mathrm{win}}$ of \eqref{4.21:eq-term-sums-proof-1} is then bounded by
\eqref{4.21:eq-term-sums-proof-2}, which is hypothesis (ii). This bound is taken by statement from
three results: the window-anchor bound for stored terms at a block of the perturbative step; the
summation over stored terms in the relative Stage~B theorem; and the first clause of the norm of
chain terms of the perturbative step. In it the factor $e^{\mathsf{e}(c)}$
enters as follows.
\begin{itemize}
\item For a member $c=T$ with $S_c=[A_T]_m$, $m$ its filing level, it is absorbed by lowering the
decay rate $\sigma_{*}$ by $1$. The
remaining rate $\sigma_{*}-1\ge c_d+1$ is the rate at which the summation over stored terms of the
relative Stage~B theorem is carried out.
\item For a chain term, $e^{\mathsf{e}(c)}$ is its factor $e^{d_k(S)/N}$ in the first clause of the
norm of chain terms.
\end{itemize}

\emph{The total $\mu$.} Fix $q\in\mathsf{Q}$ and consider the groups $\hat e$ with
$q\in\mathsf{b}(\hat e)$.

(\(\mu\)-a) Suppose $q\in Q\cup\mathcal{S}_j(c)$. These are window-anchor cases, with the anchoring
condition $S_c\cap q\ne\emptyset$ read on $S_c$. The groups are bounded by the two anchored forms
obtained in Lemma~\ref{4.21:prf-group-totals-proof} at a cube $q_0\in Q$, resp.
$q_0\in\mathcal{S}_j(c)$. Their sum over $c$ is bounded by the anchored totals for the $A$-class
at window anchors, with $e^{\Sigma_c}$ bounded as in that lemma.

(\(\mu\)-b) Suppose $q\subset\hat{\mathsf{K}}_A\cup\bigcup\hat{\mathsf{y}}$. Then
$(\square_0,q)$ holds for $\hat e$ for some $\square_0\in A$: either the block whose
$\hat{\mathsf{K}}_{\square_0}$ contains $q$, or the block on whose $\hat{\mathsf{K}}_{\square_0}$ the
animal of $\hat{\mathsf{y}}$ through $q$ hangs, that is, shares a wall with it. The blocks of $A$
belong to $\mathcal{A}_c$.

We apply the anchored bound \eqref{4.21:eq-group-totals-proof-a} with $\mathsf{E}:=q$, and bound
$e^{\Sigma_c}\le e^{\frac12+\mathsf{e}(c)}$ by $\Sigma_c\le\frac12+\mathsf{e}(c)$, as in
Lemma~\ref{4.21:prf-group-totals-proof}. This gives
\begin{equation*}
\sum_{(\mu\text{-b})}\tilde\nu(\hat e)\le\sum_{\square_0}\ \sum_{c:\,\square_0\in\mathcal{A}_c}
2e^{c_J+\frac12}\,\mathfrak{w}(c)e^{\mathsf{e}(c)}\,y_{\square_0}\,
e^{-(\operatorname{dist}_k(q,\square_0)-103\check{R}_k)_+}.
\end{equation*}
In Lemma~\ref{4.21:prf-group-totals-proof}:
\begin{itemize}
\item on islands, $y_{\square_0}\le K_2^{\flat}\vartheta e^{-\frac14\delta_P d^{\wedge}(\square_0,R_n)}$;
\item on $\mathcal{A}^{\mathrm{pl}}$, $y_{\square_0}=K_2^{\flat}\vartheta e^{-\frac12\delta_P d^{\wedge}(\square_0,R_n)}$.
\end{itemize}

\emph{Old blocks.} Let $\square_0\in\mathfrak{O}$. Since $\mathsf{f}\ge0$, the sum over $c$ of
$\mathfrak{w}(c)e^{\mathsf{e}(c)}y_{\square_0}$ is at most $\mathfrak{S}(\square_0)$ times the bound on
$y_{\square_0}$. By clause (d),
\begin{equation*}
y_{\square_0}\,\mathfrak{S}(\square_0)\le K_2^{\flat}\vartheta C^{P}_{\mathsf{a}}
e^{-\frac18\delta_P d^{\wedge}(\square_0,R_n)}.
\end{equation*}
On an island the exponents combine to $\frac14-\frac18=\frac18$; on a plain old block to
$\frac12-\frac18\ge\frac18$. The first inequality of \eqref{4.21:eq-group-totals-proof-b}, with
$x:=q$, then bounds the sum over $\square_0\in\mathfrak{O}$ by
$2e^{c_J+\frac12}K_2^{\flat}\vartheta C^{P}_{\mathsf{a}}V'$.

\emph{Window blocks.} By \eqref{4.21:eq-term-sums-proof-1}, the sum over $c$ of
$\mathfrak{w}(c)e^{\mathsf{e}(c)}$ is at most $\mathfrak{T}^{\mathrm{win}}$. Moreover
$y_{\square_0}\le K_2^{\flat}\vartheta e^{-\frac18\delta_P d^{\wedge}(\square_0,R_n)}$. The second
inequality of \eqref{4.21:eq-group-totals-proof-b}, a sum over all blocks, bounds this part by
$2e^{c_J+\frac12}K_2^{\flat}\vartheta c_{\bigstar}V'\mathfrak{T}^{\mathrm{win}}$.

Adding the two parts gives \eqref{4.21:eq-term-sums-proof-3}.

The position of $\square_0$ is essential here. Without the factor
$e^{-(\operatorname{dist}_k(q,\square_0)-103\check{R}_k)_+}$, the sum over all blocks of
$y_{\square_0}$ times the total term weight at $\square_0$ would run over all treated components. Their
number is a volume. The factor localises the sum; it is obtained from the spare factor $e$ in the
condition $(F_{\kappa_1})^{\flat}$. No transfer from a point to a block through a bounded number of
blocks per cube is used.

The share of $\mu$ coming from blocks carries the factor $V'$, which the anchored totals for the
$A$-class do not carry. The first inequality of $(Q^{+})^{A,\flat}$ in
\eqref{4.21:eq-coupling-interval-a} already contains the product of $N$, $c_{\mathsf{p}}V'$ and
$\vartheta$; hence it gives $\mu\le\frac14$. The Stage~B proposition requires only
$\mu+\tilde{\mathfrak{l}}\le\frac12$.

\emph{The total $\mathsf{M}_{*}$.} Let $x\in[\mathsf{d}^{\mathrm{rel}}(\hat e)]_k$.

(x1) Suppose $(\square_0,\Delta_k(x))$ holds for some $\square_0\in A$. The argument of
(\(\mu\)-b), with $\mathsf{E}:=\Delta_k(x)$ in place of $q$, gives the same bound.

(x2) Suppose $\Delta_k(x)$ meets $S_c$. Then every group of $c$ counts. In the bound of
Lemma~\ref{4.21:prf-group-totals-proof} for the sum over the groups of $c$, we use
$e^{\Sigma_c}-1\le\Sigma_c e^{\Sigma_c}$. The part of $\Sigma_c$ coming from blocks is
$\sum_{\square_0\in\mathcal{A}_c}y_{\square_0}$, and we interchange the sums over $c$ and over
$\square_0$.

\emph{Window blocks, rim cubes and strong sets.} These parts are window anchors. They are bounded by
the rooted form of the window-anchor bound and by the anchored totals for the $A$-class,
applied with $(V'\mathsf{p}^{\mathrm{rel}},K_2^{\flat})$.

\emph{Old blocks $\square_0$.} We distinguish three kinds of member.
\begin{itemize}
\item A chain term of level $k$. Here $S_c\in\mathbb{D}_k$ meets $\Delta_k(x)$, so $x\in S_c$. This is
a window anchor, bounded through the second clause of the norm of chain terms. Its factor
$\Sigma_c e^{\Sigma_c}$ is controlled by clause (c) of
Proposition~\ref{4.21:prop-preparation-machine}. From $t\le e^t$ with $t=d_k(S_c)/2N$ one has
$d_k(S_c)\le 2Ne^{d_k(S_c)/2N}$, hence
\begin{equation*}
\mathsf{p}^{\mathrm{rel}}(c)\le c_{\mathsf{p}}\bigl(2+2N/\check{R}_k\bigr)e^{d_k(S_c)/2N},
\end{equation*}
which is the factor $\mathsf{r}$ of the perturbative step up to a numerical factor. For these terms
$\mathsf{f}(c)=0$, and no further decay is needed.
\item A chain term of level below $k$. It is confined, so it lies, together with $\square_0$, in the
treated component meeting $\Delta_k(x)$. Hence
$\operatorname{dist}_k(x,\square_0)\le 101\check{R}_k$, and the factor
$e^{-(\operatorname{dist}_k(x,\square_0)-103\check{R}_k)_+}$ equals $1$.
\item A member $c=T$ with $S_c=[A_T]_m$, $m$ its filing level, or a frozen term. The hull
$[S_c]_k$ is connected, contains $\Delta_k(x)$, and
contains a cube within $101\check{R}_k+2$ of $\square_0$. By clause (g4) of the geometry lemma,
$\mathsf{f}(c)=d_k([S_c]_k)\ge\operatorname{dist}_k(x,\square_0)-103\check{R}_k$. Since also
$\mathsf{f}(c)\ge0$, we get
\begin{equation*}
\mathfrak{w}(c)e^{\mathsf{e}(c)}\le\mathfrak{w}(c)e^{\mathsf{e}(c)+\mathsf{f}(c)}
e^{-(\operatorname{dist}_k(x,\square_0)-103\check{R}_k)_+},
\end{equation*}
and so
\begin{equation*}
\sum\bigl\{\mathfrak{w}(c)e^{\mathsf{e}(c)}:\ \square_0\in\mathcal{A}_c,\ [S_c]_k\ni x\bigr\}
\le e^{-(\operatorname{dist}_k(x,\square_0)-103\check{R}_k)_+}\,\mathfrak{S}(\square_0).
\end{equation*}
This is the principle of the rooted window-anchor bound, that the surplus length of a long term sums
the distant components; here the decay that pays for it is $\mathsf{f}(c)$.
\end{itemize}
In the first case the contribution is a window-anchor contribution, bounded, as stated there,
through the second clause of the norm of chain terms with the factor of clause (c). In the second
and third cases, $e^{\Sigma_c}\le e^{\frac12+\mathsf{e}(c)}$ and clause (d) bound the sum over $c$
by the per-block factor
\begin{equation*}
2e^{c_J+\frac12}K_2^{\flat}\vartheta C^{P}_{\mathsf{a}}
e^{-\frac18\delta_P d^{\wedge}(\square_0,R_n)}
e^{-(\operatorname{dist}_k(x,\square_0)-103\check{R}_k)_+}.
\end{equation*}
The first inequality of \eqref{4.21:eq-group-totals-proof-b} bounds the sum of the last two
factors over old blocks by $V'$. The
contribution is therefore at most $2e^{c_J+\frac12}K_2^{\flat}\vartheta C^{P}_{\mathsf{a}}V'$.
This proves (I) and (II).

\emph{Stage B.} The Stage~B proposition uses the elements it sums only through the triple
$(\mathsf{b},\mathsf{d},\nu)$. We apply it to the family $\hat{\mathfrak{E}}$ of groups, as clause
(v) of Lemma~\ref{4.21:lem-grain-resummation} allows.

In its proof the domain $\mathsf{d}$ enters in four places:
\begin{itemize}
\item in the weight $w(\hat e)$;
\item in the connectedness of the union of the domains $\mathsf{d}(o)$ of the elements $o$ of a
polymer;
\item in the rooting of polymers at a point;
\item in the domain of the output in its clause (i).
\end{itemize}
In each place the domain $\mathsf{d}^{\mathrm{rel}}(\hat e)$ serves, by clause (a$'$) of
Proposition~\ref{4.21:prop-preparation-machine}: it is connected, contains $\mathsf{b}(\hat e)$, and
is common to the fibre. The factorisation used in that proof is over integration variables and does
not involve domains.

By the relation among the constants of the preparation step
(Notation~\ref{4.21:not-preparation-constants}), and by \eqref{4.21:eq-term-sums-proof-2} and
\eqref{4.21:eq-term-sums-proof-3}, the totals of the two preceding paragraphs are at most
\begin{equation*}
\tfrac14K_{\flat}''V'\bigl(\vartheta+K_Q/r_A+\varepsilon_S\bigr)
\bigl(E_0+B_0+2+NC_0^{\sharp}\bigr).
\end{equation*}
The conditions $(Q^{+})^{A,\flat}$ then give
\begin{equation*}
\mu+\tilde{\mathfrak{l}}\le\tfrac12,\qquad 4\bigl(\mathsf{M}_{*}+2^dV\tilde{\mathfrak{l}}\bigr)
\le\tfrac12.
\end{equation*}
By clauses (a) and (a$'$) of Proposition~\ref{4.21:prop-preparation-machine}, the outputs read the
fields on $X$ together with strips $Y^{\sharp}$, and these strips touch $X$. The two statements of
the perturbative step on the domains of outputs are statements about domains, and they hold with the
domains read on labels. This proves (III).
\end{proof}

\begin{remark}[Constants, polydiscs, and the sums over terms at a block]
\begin{enumerate}
\item[(1)] The degree in $\log g^{-2}$ of the conditions $(Q^{+})^{A,\flat}$ is that of
$(Q^{+})^{A}$. Indeed, by clause (c) of Proposition~\ref{4.21:prop-preparation-machine},
$V'\le C(d,L)(\check{R}^{\uparrow})^d$, as for $V$. The constants $K_2^{\flat}$, $K_{\flat}''$,
$c_{\mathsf{p}}$ and $c_{\bigstar}$ are numbers fixed before $\gamma$, also beside a window; this is
what Lemma~\ref{4.21:lem-grain-resummation} provides. The number $N$ enters to the first power, as
in the anchored totals for the $A$-class (at window anchors), and not at all at old anchors.
\item[(2)] The parameter polydiscs at the preparation site are the sealed half-rate polydisc fixed
earlier, which is not shrunk here. Island blocks and old blocks use it as all
blocks do. Lemma~\ref{4.21:lem-grain-resummation} uses a diagonal of its $\mathsf{s}$-part.
\item[(3)] Lemma~\ref{4.21:lem-grain-resummation} does not treat the sums over terms at one block.
For wrapped terms at old blocks these sums are bounded by clause (d) of
Proposition~\ref{4.21:prop-preparation-machine}. At window blocks they are hypothesis (ii) of
Lemma~\ref{4.21:prf-term-sums-proof}, the window total bound \eqref{4.21:eq-term-sums-proof-2}.
For stored terms $T$ that are not wrapped, at blocks of old zones, a count of at most $N$ levels in
the window-anchor bound does not suffice; clause (c) of
Proposition~\ref{4.21:prop-preparation-machine} gives the
volume sum, since with $\operatorname{Cor}=\emptyset$ it holds for such terms too; the decay is paid
by $\vartheta(\square_0)$; and the count of levels for these terms is the one made in the proof of
the correlation theorem, and is not supplied here.
\end{enumerate}
\end{remark}

\begin{definition}[Classes of wrapped units at the large-field site and the starred norm]
\label{4.21:def-wrapped-classes}
We first describe the scheme in which the classes below are used. The expansion at the
large-field site of $\mathbf{R}_k$ follows \cite{B-CMP122b}, and is read on the global chain as in
Remark~\ref{4.21:rem-global-chain}. It is sequential in its index sets. The layer of link $t$
is read where the links before it have carried the reading, so the family $\mathsf{y}_t$ of cut
cells is anchored on the reading set $\mathcal{R}_{t-1}$. As a function, however, the chain is a
single function of the decoupling parameters
$(\mathsf{s}^{(1)},\mathsf{s}^{(2)},\mathsf{s}^{(3)},\mathsf{s}^{(4)})$ of the four links.

The bracket of a unit of the second part (defined below) has a lower end that is not a free
background. Its piece belonging to the family $\mathsf{y}_1$ re-enters, as a unit of the first
part, in the reading of links two to four. No cell of a live strip is left uncut: each family
runs either over all cells outside $\Lambda^{\sim}$ or over all cells outside $\Lambda$, where
$\Lambda$ and $\Lambda^{\sim}$ are the two domains of the large-field operation of
\cite{B-CMP122b} on whose complements the families of the scheme are indexed; in
either case the cells of $X_v$ are included.

Let $v$ be a wrapped unit at the large-field site, of level $j\le k$. That is, $v$ is a stored
boundary-type term in the relative format of Definition~\ref{4.21:def-relative-format}, and its
set $\mathcal{P}_v$ of pairs $(Y,m)$ is non-empty (Definition~\ref{4.21:def-decoupling-clause}).
Its region is $X_v$. We split
\begin{equation*}
  \mathcal{P}_v=\mathcal{P}_v^{\dagger}\sqcup\mathcal{P}_v^{\circ},
\end{equation*}
where the two parts are as follows.
\begin{itemize}
  \item $\mathcal{P}_v^{\dagger}$ consists of the pairs with $Y\subset Z$. These do not survive
    the present step.
  \item $\mathcal{P}_v^{\circ}$ consists of the pairs whose $Y$ lies in a component of the
    large-field region that $\mathbf{R}_k$ does not treat. These are live after $\mathbf{R}_k$.
\end{itemize}
We put $K_v:=[A_v]_k$, the $\pi_k$-hull of the label $A_v$, and $a:=\tfrac{31}{20}\,\kappa$.

The units with $X_v\cap Z=\emptyset$ are the units of the \emph{second part}; we call the
remaining units, those with $X_v\cap Z\neq\emptyset$, the units of the \emph{first part}. We
define four classes.
\begin{itemize}
  \item $\mathfrak{L}^{\mathrm{II}}$ consists of the units with $X_v\cap Z=\emptyset$, that is,
    of the units of the second part. For these, $\mathcal{P}_v=\mathcal{P}_v^{\circ}$.
  \item $\mathfrak{L}^{\circ}$ consists of the units of the first part with
    $\mathcal{P}_v^{\circ}\neq\emptyset$.
  \item $\mathfrak{L}^{\mathrm{sh}}$ consists of the units of the first part with
    $\mathcal{P}_v^{\circ}=\emptyset$ and with $X_v$ contained in no core $\mathsf{C}_C$ of a
    component $C$ of $Z$, the cores being those of the core geometry of this chapter. Its
    units are called \emph{shallow}.
  \item $\mathfrak{L}^{\mathrm{dp}}$ consists of the units with
    $X_v\subset\mathsf{C}_C$ for a component $C$ of $Z$. Its units are called \emph{deep}.
\end{itemize}
The deep class is the class that occurs generically at a step at which the region $Y$ of one of
the unit's own pairs $(Y,m)\in\mathcal{P}_v$ is healed, that is, at which its large-field factor
is absorbed back into the small-field format and the region is re-filed into the regular
expansion.

We put
\begin{equation*}
  \mathfrak{L}^{\mathrm{I}}:=\mathfrak{L}^{\circ}\sqcup\mathfrak{L}^{\mathrm{sh}}
  \sqcup\mathfrak{L}^{\mathrm{dp}},
  \qquad
  \mathfrak{L}_L:=\mathfrak{L}^{\mathrm{I}}\sqcup\mathfrak{L}^{\mathrm{II}}.
\end{equation*}
The set $X_v=A_v\cup\bigcup_{(Y,m)\in\mathcal{P}_v}Y^{\sharp}$ of
Definition~\ref{4.21:def-relative-format} is connected, and, by the core geometry of this
chapter, the core $\mathsf{C}_C$ of a component $C$ of $Z$ lies in $C$,
so that the cores of distinct components of $Z$ are disjoint. Since $X_v\supset A_v$ is
non-empty, a deep unit therefore determines the component $C$ with $X_v\subset\mathsf{C}_C$.

\emph{The starred norm.} Let $C_{3F}$ and $\varsigma_3$ be the constant and the rate of the
minimal form of the third-factor lemma, \eqref{4.21:eq-third-factor-form-a}. By the third-factor
lemma, \eqref{4.21:eq-third-factor-form-a} holds with $\varsigma_3=\tfrac12\delta_P$. We may
assume, without loss, that $C_{3F}\ge1$. For a wrapped unit $v$ at the large-field site we put
\begin{equation}\label{4.21:eq-wrapped-classes-1}
  \|v\|_{\star}:=
  \begin{cases}
    \mathfrak{N}^{A}(v), & v\notin\mathfrak{L}^{\mathrm{dp}},\\[2pt]
    C_{3F}\,e^{-\varsigma_3\,d^{\wedge}(X_v,\,\mathsf{C}_C^{c})}\,\mathfrak{N}^{A}(v),
      & v\in\mathfrak{L}^{\mathrm{dp}},\ X_v\subset\mathsf{C}_C.
  \end{cases}
\end{equation}
Here $\mathfrak{N}^{A}$ is the size functional of stored terms in the relative format of
Definition~\ref{4.21:def-relative-format}, and $d^{\wedge}(X_v,\mathsf{C}_C^{c})$ is the scaled
distance from $X_v$ to the complement $\mathsf{C}_C^{c}$ of the core $\mathsf{C}_C$.

\emph{The fattened domain and exterior pieces.} For a union $Y$ of $\pi_k$-cubes we put
\begin{equation}\label{4.21:eq-wrapped-classes-2}
  Y^{+}:=Y\cup\bigcup\bigl\{Y_p^{\sharp}:\ (Y_p,m)\ \text{live after}\ \mathbf{R}_k,\ m<k,\
  Y_p^{\sharp}\ \text{touches}\ Y\bigr\}.
\end{equation}
A piece $\mathfrak{p}$ of the decoupling expansion is \emph{exterior} if its label
$\operatorname{lab}_4$, the label after the fourth link, has a cube not contained in $Z$. Otherwise
it is \emph{interior}.
\end{definition}

\begin{lemma}[The sum of deep units at a core]\label{4.21:lem-deep-unit-sum}
Work at the large-field site in the setting of Definition~\ref{4.21:def-wrapped-classes}, at level
$k$. Let $Z$, the components $C$ of $Z$, the cores $\mathsf{C}_C$, the class $\mathfrak{L}^{\mathrm{dp}}$
of deep units and the hulls $K_v=[A_v]_k$ be as there. Let $\mathsf{b}_k$ be the per-cube majorant of
Lemma~\ref{4.21:lem-wrapped-cube-bound}, and let $\check{R}_k$ be the region size parameter at
level $k$. Let
\[
\varsigma_3\in\bigl[\tfrac18\delta_P,\delta_P\bigr].
\]
Then for every component $C$ of $Z$, with $h$ the level of the cells of $\mathsf{C}_C$ that touch
$\mathsf{C}_C^{c}$ (part (g) of the geometry lemma for zones, strips and cores), and every
$r\le\frac85\kappa$,
\begin{equation}\label{4.21:eq-deep-unit-sum-1}
\begin{split}
&\sum_{v\in\mathfrak{L}^{\mathrm{dp}},\,X_v\subset\mathsf{C}_C}
\mathfrak{N}^{A}(v)\,e^{r\,d_k(K_v)}\,e^{-\varsigma_3 d^{\wedge}(X_v,\mathsf{C}_C^{c})}\\
&\qquad\le e\,C_{\mathsf{a}}(1)\,C_{\mathrm{cell}}\cdot 2d\,\bigl(104L\check{R}_{h+1}+2\bigr)^{d-1}
\cdot(2+\check{R}_k)\,\mathsf{b}_k(\check{R}_k+1).
\end{split}
\end{equation}
Here $e$ is Euler's number, $C_{\mathsf{a}}(1)$ is the anchoring constant of
Definition~\ref{4.21:def-wrapped-terms} at $\beta=1$, and $C_{\mathrm{cell}}$ is the constant of
part (b) of the cell-summation lemma.

Consequently, by the minimal form of the third-factor lemma
(Definition~\ref{4.21:def-third-factor-form}), which follows from the third-factor lemma, the
exterior pieces of the deep units of every kind enter the functional $\mathsf{Q}^{\star}_r$ with the
bound \eqref{4.21:eq-deep-unit-sum-1} multiplied by $C_{3F}$. This covers the pieces with families
$\mathsf{y}_t$ at the links $t=2,3,4$ and the pieces with a family at the link $t=1$ alone.
\end{lemma}

\begin{proof}
The argument uses no property of the fields. Throughout, $w=M/M_1$, $d^{\wedge}$ is the scaled
distance, $\operatorname{lev}(\Delta)$ is the level of a cell $\Delta$, and $\operatorname{dist}_h$
is the distance measured in units of level $h$.

\emph{Reduction to cells.} Let $v$ be a unit in the sum. Since $X_v\subset\mathsf{C}_C$, there is a
point $x\in X_v$ with $d^{\wedge}(x,\mathsf{C}_C^{c})=d^{\wedge}(X_v,\mathsf{C}_C^{c})$. Let $\Delta_v$
be a cell of $\mathsf{C}_C$ containing such a point. The distance of a set is at most that of each
of its points, so
\[
d^{\wedge}(\Delta_v,\mathsf{C}_C^{c})\le d^{\wedge}(X_v,\mathsf{C}_C^{c}),
\]
and $\Delta_v$ meets $X_v$.

Now take, in Definition~\ref{4.21:def-wrapped-terms}, the following data:
\begin{itemize}
\item $\mathfrak{L}$ is the family of all wrapped units;
\item $\mathfrak{Q}$ is the class of cells of $\mathsf{C}_C$;
\item the weight is $W(v):=\mathfrak{N}^{A}(v)e^{r d_k(K_v)}$;
\item the current level is $i=k$.
\end{itemize}
Let $\mathsf{S}$ be the total mass of Lemma~\ref{4.21:lem-probe-mass}. Each term of the left member
of \eqref{4.21:eq-deep-unit-sum-1} is non-negative. Each deep unit is a wrapped unit, and it is
counted in $\mathsf{S}(\Delta_v)$. Its exponential factor is at most
$e^{-\varsigma_3 d^{\wedge}(\Delta_v,\mathsf{C}_C^{c})}$. The left member is therefore at most
\begin{equation}\label{4.21:eq-deep-unit-sum-2}
\sum_{\Delta\subset\mathsf{C}_C}e^{-\varsigma_3 d^{\wedge}(\Delta,\mathsf{C}_C^{c})}\,\mathsf{S}(\Delta).
\end{equation}
For this weight, the two per-cube bounds among the hypotheses \eqref{4.21:eq-wrapped-terms-a} hold
with $\mathsf{b}=\mathsf{b}_k$. This is Lemma~\ref{4.21:lem-wrapped-cube-bound} at the large-field
site, where $r\le\frac85\kappa$ is the range allowed there.

\emph{The depth datum.} We put
\[
\varrho(\Delta):=1+\frac{d^{\wedge}(\Delta,\mathsf{C}_C^{c})}{w},\qquad N_{*}:=\check{R}_k+1 .
\]
Then $\varrho(\Delta)\ge1$ and $N_*\ge0$. We verify the two conditions of the depth datum in
\eqref{4.21:eq-wrapped-terms-a}.

The first condition concerns levels. By part (g) of the geometry lemma for zones, strips and cores,
every cell of $\mathsf{C}_C$ that touches $\mathsf{C}_C^{c}$ has level $h$. A path from
$\mathsf{C}_C^{c}$ to $\Delta$ therefore crosses the zones of levels
$h-1,\dots,\operatorname{lev}(\Delta)+1$. Hence
\[
d^{\wedge}(\Delta,\mathsf{C}_C^{c})\ge w\,\bigl(h-\operatorname{lev}(\Delta)-1\bigr)_+ ,
\]
that is, $(h-\operatorname{lev}(\Delta)-1)_+\le\varrho(\Delta)-1$, so that
$\varrho(\Delta)\ge h-\operatorname{lev}(\Delta)$. Since moreover $k-h\le\check{R}_k$, it follows that
\[
k-\operatorname{lev}(\Delta)+1\le\check{R}_k+1+\varrho(\Delta)=\varrho(\Delta)+N_* .
\]

The second condition concerns strips. Let $\Delta\subset Y^{\sharp}$ for a live pocket $(Y,m)$.
By part (g) of the geometry lemma the levels in $\mathsf{C}_C$ are at most $h$, so $m\le h$ and two
cases arise.

If $m=h$, then by part (g) of the geometry lemma,
\[
d^{\wedge}(\Delta,\mathsf{C}_C^{c})\ge w\,\operatorname{dist}_h(Y^{\sharp},\mathsf{C}_C^{c})
\ge\tfrac12 w\,(\check{R}_h+1).
\]
This gives $\varrho(\Delta)\ge1+\frac12(\check{R}_h+1)>\frac12\check{R}_h$.

If $m<h$, consider the collar of $Y^{\sharp}$ given by part (c) of the geometry lemma. It lies in
\[
\Omega_{m+1}^{c}\cap C\subset Z_h\subset\mathsf{C}_C ,
\]
so it separates $\Delta$ from $\mathsf{C}_C^{c}$, and
$d^{\wedge}(\Delta,\mathsf{C}_C^{c})\ge\frac52 w\,\check{R}_m$. This gives
$\varrho(\Delta)\ge1+\frac52\check{R}_m>\frac12\check{R}_m$.

In both cases $\check{R}_m\le2\varrho(\Delta)\le2\varrho(\Delta)+2$.

\emph{Splitting the rate.} Write
\[
e^{-\varsigma_3 d^{\wedge}}=e^{-\frac12\varsigma_3 d^{\wedge}}\cdot e^{-\frac12\varsigma_3 d^{\wedge}} .
\]
Since $\varsigma_3\ge\frac18\delta_P$, the lower bound imposed on $w$ in the choice of parameters
gives $\frac12\varsigma_3 w\ge\frac1{16}\delta_P w\ge1$. Therefore
\[
e^{-\frac12\varsigma_3 d^{\wedge}(\Delta,\mathsf{C}_C^{c})}
\le e^{-d^{\wedge}(\Delta,\mathsf{C}_C^{c})/w}=e^{1-\varrho(\Delta)} .
\]
Lemma~\ref{4.21:lem-probe-mass}, part (iv), at $\beta=1$ gives
\[
\mathsf{S}(\Delta)\le e^{\varrho(\Delta)}\,C_{\mathsf{a}}(1)\,(1+N_*)\,\mathsf{b}_k(N_*).
\]
Its hypotheses are the per-cube bounds supplied by Lemma~\ref{4.21:lem-wrapped-cube-bound}, the
growth law in \eqref{4.21:eq-wrapped-terms-a} for $\mathsf{b}_k$, and the depth datum verified
above. Multiplying the two bounds,
\begin{equation}\label{4.21:eq-deep-unit-sum-3}
e^{-\frac12\varsigma_3 d^{\wedge}(\Delta,\mathsf{C}_C^{c})}\,\mathsf{S}(\Delta)
\le e\,C_{\mathsf{a}}(1)\,(2+\check{R}_k)\,\mathsf{b}_k(\check{R}_k+1).
\end{equation}

\emph{Summing over cells.} Let $\Xi$ be the closure of $\mathsf{C}_C^{c}$. By part (g) of the
geometry lemma, $\Xi$ is a union of cells. Apply part (b) of the cell-summation lemma with
$\Xi$ and the rate $\beta=\frac12\varsigma_3\ge\frac1{16}\delta_P$. This gives
\begin{equation}\label{4.21:eq-deep-unit-sum-4}
\sum_{\Delta\subset\mathsf{C}_C}e^{-\frac12\varsigma_3 d^{\wedge}(\Delta,\mathsf{C}_C^{c})}
\le C_{\mathrm{cell}}\,\#\partial^{+}\Xi .
\end{equation}
Here $\partial^{+}\Xi$ is the outermost layer of level-$h$ cells of the box $\mathsf{C}_C$, so
\[
\#\partial^{+}\Xi\le2d\,(104L\check{R}_{h+1}+2)^{d-1}.
\]
In \eqref{4.21:eq-deep-unit-sum-2}, bound each summand by
\eqref{4.21:eq-deep-unit-sum-3} times $e^{-\frac12\varsigma_3 d^{\wedge}(\Delta,\mathsf{C}_C^{c})}$,
and sum with \eqref{4.21:eq-deep-unit-sum-4}. This gives \eqref{4.21:eq-deep-unit-sum-1}.

\emph{The final assertion.} The minimal form of the third-factor lemma bounds each exterior piece
of a deep unit $v$ with the factor $C_{3F}\,e^{-\varsigma_3 d^{\wedge}(X_v,\mathsf{C}_C^{c})}$ in
front of the size of $v$. Combined with \eqref{4.21:eq-deep-unit-sum-1}, this is the final
assertion of the lemma.
\end{proof}

\begin{remark}
The bound does not depend on the history of the core: its volume, the number, sizes and ages of
its pockets, and its thin strata do not enter. The cells are summed by the cell-summation lemma at
a rate that beats their multiplication per crossed cell, and the terms at one cell by
Lemma~\ref{4.21:lem-probe-mass}.

At the rate $\varsigma_3=\frac12\delta_P$ of the third-factor lemma, the rate $\frac14\delta_P$ goes
to Lemma~\ref{4.21:lem-probe-mass} and $\frac14\delta_P$ to the cell summation.

For plain deep units the same two lemmas are applied, with the total weight of the stored terms
meeting a cell in place of $\mathsf{S}$, in the estimate for plain deep units. Part (i) of
Lemma~\ref{4.21:lem-probe-mass} serves every family satisfying the per-cube bound of one level in
\eqref{4.21:eq-wrapped-terms-a}; parts (ii) and (iii) are particular to wrapped terms. That estimate
is not carried out in the present lemma.
\end{remark}

\begin{lemma}[Wrapped mass over one top-level cube]\label{4.21:lem-wrapped-mass}
Work at the large-field site at level $k$, with the wrapped units $v$, their labels $A_v$, regions
$X_v$, pocket sets $\mathcal{P}_v$, hulls $K_v=[A_v]_k$, the classes $\mathfrak{L}^{\circ}$,
$\mathfrak{L}^{\mathrm{sh}}$, $\mathfrak{L}^{\mathrm{dp}}$, $\mathfrak{L}^{\mathrm{I}}$, the starred norm
$\|\cdot\|_{\star}$ and the cores $\mathsf{C}_C$ of the components $C$ of $Z$ of
Definition~\ref{4.21:def-wrapped-classes}. Let $h$ be the level of the cells of a core that touch its
complement (clause (g) of the lemma on the geometry of the cores). Let $\mathsf{Q}^{\star}_r$ be the
cube-anchored functional $\mathsf{Q}_r[\mathfrak{L}^{\mathrm{I}}]$ computed with $\|\cdot\|_{\star}$ in
place of $\mathfrak{N}^A$, and let $\mathsf{b}_k$ be the per-cube majorant of
Lemma~\ref{4.21:lem-wrapped-cube-bound} at current level $k$. For $r\le\frac85\kappa$,
\begin{equation}\label{4.21:eq-wrapped-mass-1}
\mathsf{Q}^{\star}_r\le C_Q\bigl[1+(\check{R}_{h+1})^d+N(g_k)\bigr]\bigl(2+N(g_k)\bigr)\,
\mathsf{b}_k\bigl(N(g_k)+1\bigr),
\end{equation}
where $C_Q$ is a number fixed before $\gamma$ which contains the factor $C_{3F}$. When the
right-hand side of \eqref{4.21:eq-wrapped-mass-1} is evaluated at $\mathsf{T}:=\check{\mathsf{T}}_k$ by
the bounds of this chapter on the stage count $N(\cdot)$ and on the zone radii, it is at most a
polynomial in $\mathsf{T}$ with coefficients fixed before $\gamma$.
\end{lemma}

\begin{proof}
Fix a cube $Q\in\pi_k$. We bound the mass over $Q$, that is, the sum of $\|v\|_{\star}e^{r d_k(K_v)}$
over the units $v\in\mathfrak{L}^{\mathrm{I}}$ with $Q\subset K_v$, class by class. We put
$W(v):=\mathfrak{N}^A(v)\,e^{r d_k(K_v)}$ and $N:=N(g_k)$ (in this proof $N$ abbreviates the stage
count $N(g_k)$, not the rank of the gauge group). For $v\notin\mathfrak{L}^{\mathrm{dp}}$ the
starred norm is $\mathfrak{N}^A(v)$, so the contribution of such $v$ is the sum of their weights $W(v)$.
Since $r\le\frac85\kappa$, the per-cube bound on stored wrapped terms
(Lemma~\ref{4.21:lem-wrapped-cube-bound}), applied at the large-field site with current level $k$,
holds with $\varsigma_0=\frac15\kappa$: for every $j\le k$ and every $\square'\in\pi_j$,
\begin{equation}\label{4.21:eq-wrapped-mass-2}
\sum_{v:\ j_v=j,\ A_v\ni\square'}W(v)\,e^{\varsigma_0 d_j(A_v)\mathbf{1}[j<k]}\le\mathsf{b}_k(k-j).
\end{equation}
Since $\pi_j$ refines $\pi_k$ for $j\le k$ and $A_v$ is a union of $\pi_j$-cubes, a unit $v$ of level
$j$ with $Q\subset K_v$ has a label containing one of the $L^{d(k-j)}$ cubes of $\pi_j$ contained
in $Q$.

\emph{The class $\mathfrak{L}^{\circ}$.} Let $v\in\mathfrak{L}^{\circ}$ have level $j$. Its label $A_v$
joins a cube containing a point of a live strip, which lies at distance at least $\check{R}_k$ from $Z$,
to a cube meeting $Z$; hence $d_j(A_v)\ge(\check{R}_k-2)L^{k-j}$. We anchor $v$ at the
$L^{d(k-j)}$ cubes of $\pi_j$ in $Q$ and apply \eqref{4.21:eq-wrapped-mass-2} at each of them. For
$j<k$ the factor $e^{\varsigma_0 d_j(A_v)}$ is present in \eqref{4.21:eq-wrapped-mass-2}, so
$W(v)\le e^{-\varsigma_0(\check{R}_k-2)L^{k-j}}\,W(v)e^{\varsigma_0 d_j(A_v)}$; writing $s=k-j$, the
levels $j<k$ contribute at most
\begin{equation*}
\sum_{s\ge1}L^{ds}e^{-\varsigma_0(\check{R}_k-2)L^s}\,\mathsf{b}_k(s)\le C\,\mathsf{b}_k(0),
\end{equation*}
by the growth law of the per-cube majorant $\mathsf{b}$ and the lower bound $\check{R}\ge 38$ on the
zone radii. Here and below $C$ denotes a constant depending only on $d$, $L$, $\varsigma_0$ and the
constants of the growth law of $\mathsf{b}$, hence fixed before $\gamma$; its value may change from one
occurrence to the next. The level $j=k$ contributes at most $\mathsf{b}_k(0)$, by
\eqref{4.21:eq-wrapped-mass-2} at $\square'=Q$.

\emph{The class $\mathfrak{L}^{\mathrm{sh}}$.} Let $v\in\mathfrak{L}^{\mathrm{sh}}$, of level $j$. The
pocket set $\mathcal{P}_v$ is non-empty (Definition~\ref{4.21:def-wrapped-terms}) and contains no pocket
that is live after $\mathbf{R}_k$; so $v$ has a pocket $(Y,m)$ with $Y$ in $Z$, which dies at this step,
and $Y$ lies in a component $C$ of $Z$. Then $A_v$ contains a $\pi_j$-cube $Q_1$ touching
$Y^{\sharp}$, and $Y^{\sharp}\subset\mathsf{C}_C$, because the strips of the pockets of $C$ lie in
$\mathsf{C}_C$. Moreover $A_v$ contains a point off $\mathsf{C}_C$: indeed
$X_v\not\subset\mathsf{C}_C$, since $v$ is shallow; if $A_v\subset\mathsf{C}_C$, some strip in $X_v$
would lie outside $\mathsf{C}_C$, hence would be the strip of a pocket in another component, and such
a pocket forces $A_v$ out of $C$. We bound the whole class attached to $C$,
\begin{equation*}
S^{\mathrm{sh}}(C):=\sum\bigl\{W(v):\ v\in\mathfrak{L}^{\mathrm{sh}},\ \mathcal{P}_v\ \text{contains a
pocket of}\ C\bigr\}.
\end{equation*}
In both ranges of levels below we use that the number $k-h$ of levels above $h$ is at most $N$ and
that $\mathsf{b}_k$ is non-decreasing, so that $\mathsf{b}_k(k-j)\le\mathsf{b}_k(N+h-j)$ for $j\le h$
and $\mathsf{b}_k(k-j)\le\mathsf{b}_k(N)$ for $h<j\le k$.

Levels $j\le h$. Put $u:=h-j$. The candidates for $Q_1$ are the $\pi_j$-cubes touching
$\mathsf{C}_C$ or contained in it; by clause (g) of the lemma on the geometry of the cores there are
at most $(104L\check{R}_{h+1}L^u+3)^d$ of them. Since $A_v$ contains $Q_1$ and a point off
$\mathsf{C}_C$, and $Q_1$ has side at most one in units of level $h$, we have, with
$\operatorname{dist}_j$ the distance measured in units of level $j$,
\begin{equation*}
d_j(A_v)\ge\operatorname{dist}_j(Q_1,\mathsf{C}_C^c)-1\ge\bigl(\tfrac12(\check{R}_h+1)-1\bigr)L^u-1
=\tfrac12(\check{R}_h-1)L^u-1.
\end{equation*}
We anchor at the candidates, apply \eqref{4.21:eq-wrapped-mass-2} at each, which gives at most the
factor $\mathsf{b}_k(k-j)\le\mathsf{b}_k(N+u)$, and use $L\check{R}_{h+1}\le L\check{R}_h$. The levels
$j\le h$ contribute at most
\begin{equation*}
\sum_{u\ge0}\bigl(104L\check{R}_hL^u+3\bigr)^d e^{-\varsigma_0(\frac12(\check{R}_h-1)L^u-1)}\,
\mathsf{b}_k(N+u)\le C\,\mathsf{b}_k(N).
\end{equation*}

Levels $h<j\le k$. At level $j$ the candidates for $Q_1$ number at most
$(104L\check{R}_{h+1}L^{-(j-h)}+3)^d$. By $(a+b)^d\le2^{d-1}(a^d+b^d)$, summation of the geometric
series in $L^{-d(j-h)}$, and the fact that there are $k-h\le N$ of these levels, these counts, summed
over $j$, are at most $2^d[2(104\check{R}_{h+1})^d+3^dN]$. Each is multiplied by
$\mathsf{b}_k(k-j)\le\mathsf{b}_k(N)$, from \eqref{4.21:eq-wrapped-mass-2}. No decay is used in this
range: the bound is a count of candidate cubes; the decay for these terms is supplied by the factor
$e^{-(\check{R}_k-2)}$ under which this count stands in clause (d) of the proposition at the
large-field site that follows. Hence
\begin{equation*}
S^{\mathrm{sh}}(C)\le\Bigl(C+2^d\bigl[2(104\check{R}_{h+1})^d+3^dN\bigr]\Bigr)\mathsf{b}_k(N).
\end{equation*}

The cube $Q$ lies in at most one component $C$ of $Z$. The units $v\in\mathfrak{L}^{\mathrm{sh}}$ with
$Q\subset K_v$ that are attached to that $C$ are among those summed in $S^{\mathrm{sh}}(C)$. The others
are attached only to components $C'$ with $Q\not\subset C'$, and their labels join $\mathsf{C}_{C'}$ to
$Q$; hence $d_j(A_v)\ge(9\check{R}_k-1)L^{k-j}$. They are anchored at the $L^{d(k-j)}$ cubes of $\pi_j$
in $Q$, as for the class $\mathfrak{L}^{\circ}$, and no sum over $C'$ is taken. Since
$9\check{R}_k-1\ge\check{R}_k-2$, the computation made for $\mathfrak{L}^{\circ}$ bounds their
contribution by $C\,\mathsf{b}_k(0)+\mathsf{b}_k(0)$.

\emph{The class $\mathfrak{L}^{\mathrm{dp}}$.} For $v\in\mathfrak{L}^{\mathrm{dp}}$ with
$X_v\subset\mathsf{C}_C$ the starred norm is
$C_{3F}e^{-\varsigma_3 d^{\wedge}(X_v,\mathsf{C}_C^c)}\mathfrak{N}^A(v)$ with
$\varsigma_3=\frac12\delta_P\in[\frac18\delta_P,\delta_P]$. So the deep-unit lemma
(Lemma~\ref{4.21:lem-deep-unit-sum} below) applies with $r\le\frac85\kappa$, and its left member
carries the weight $e^{r d_k(K_v)}$. By the consequence stated in that lemma, the deep units enter
$\mathsf{Q}^{\star}_r$ with its right-hand side times $C_{3F}$, that is, with at most
\begin{equation*}
C_{3F}\,e\,C_{\mathsf{a}}(1)\,C_{\mathrm{cell}}\,2d\bigl(104L\check{R}_{h+1}+2\bigr)^{d-1}
(2+N)\,\mathsf{b}_k(N+1).
\end{equation*}

Adding the contributions of the three classes, the mass over $Q$ is at most
\begin{equation}\label{4.21:eq-wrapped-mass-3}
\begin{split}
&2(C+1)\,\mathsf{b}_k(0)
+\Bigl(C+2^d\bigl[2(104\check{R}_{h+1})^d+3^dN\bigr]\Bigr)\mathsf{b}_k(N)\\
&\quad+C_{3F}\,e\,C_{\mathsf{a}}(1)\,C_{\mathrm{cell}}\,2d\bigl(104L\check{R}_{h+1}+2\bigr)^{d-1}
(2+N)\,\mathsf{b}_k(N+1).
\end{split}
\end{equation}
Since $\mathsf{b}_k$ is non-decreasing, $\mathsf{b}_k(0)\le\mathsf{b}_k(N)\le\mathsf{b}_k(N+1)$, and
since $\check{R}_{h+1}\ge1$,
\begin{equation*}
\bigl(104L\check{R}_{h+1}+2\bigr)^{d-1}\le(106L)^{d-1}\bigl(1+\check{R}_{h+1}^{\,d}\bigr).
\end{equation*}
Hence \eqref{4.21:eq-wrapped-mass-3} is of the form \eqref{4.21:eq-wrapped-mass-1}, with $C_Q$
depending only on $d$, $L$, $C$, $e\,C_{\mathsf{a}}(1)C_{\mathrm{cell}}$ and $C_{3F}$, all fixed
before $\gamma$; in particular $C_Q$ contains $C_{3F}$.
\end{proof}

\begin{proposition}[The relative machine at the large-field site]\label{4.21:prop-large-field-site}
We work at the large-field site of the operation $\mathbf{R}_k$. The wrapped units $v$ of level
$j\le k$, their labels $A_v$, their regions $X_v$, the hulls $K_v:=[A_v]_k$, the pocket lists
$\mathcal{P}_v=\mathcal{P}_v^{\dagger}\sqcup\mathcal{P}_v^{\circ}$, the region $Z$ of the step with
its components $C$ and their cores $\mathsf{C}_C$, the classes $\mathfrak{L}^{\mathrm{II}}$,
$\mathfrak{L}^{\mathrm{I}}$ and $\mathfrak{L}^{\mathrm{dp}}$ (the deep units), the starred norm
$\|v\|_{\star}$, the number $a=\tfrac{31}{20}\kappa$, the enlargement $Y^{+}$ of a union $Y$ of cubes
of $\pi_k$, and the exterior pieces are as in Definition~\ref{4.21:def-wrapped-classes}. The
decoupling clause (Definition~\ref{4.21:def-decoupling-clause}) prescribes its families of cells at
this site relative to two domains $\Lambda\subset\Lambda^{\sim}$ of the step, the second an
enlargement of the first; and to each component $C$ of $Z$ there is attached a domain $\Lambda_C$
with $\mathsf{C}_C\subset\Lambda_C\subset C$, together with its enlargement
$\Lambda_C^{\sim}\supset\Lambda_C$.
Assume the following.
\begin{itemize}
\item[(I)] The third-factor lemma, in its minimal form (Definition~\ref{4.21:def-third-factor-form}).
\item[(II)] From the proof of the correlation theorem: the value-domain lemma; clauses (i), (i$'$)
and (ii) of the supply statement; and clause (c) of the boundary-family proposition; each as stated
there.
\item[(III)] The strip theorem.
\end{itemize}
Then the following hold.
\begin{itemize}
\item[(a)] (Parts.) $X_v$ meets $Z$ if and only if $K_v$ meets $Z$. A strip of a pocket in
$\mathcal{P}_v^{\dagger}$ lies in a core $\mathsf{C}_C\subset\Lambda_C$ and is touched by a cube of
$A_v$ lying in $\Lambda_C$. If $X_v$ meets $\Lambda_C^{\sim}$, so does $K_v$.
\item[(b)] (Values.) The value-domain lemma holds with the hull $K_v=[A_v]_k$, with
\begin{equation*}
Y_1(v):=K_v\cup\bigcup\{C: C \text{ a component of } Z \text{ meeting } K_v\},
\end{equation*}
and with
$\mathrm{val}_v$, the piece of part (c) in which all four families $\mathsf{y}_t$ are empty,
analytic over $Y_1(v)^{+}$. Clauses
(i), (i$'$) and (ii) of the supply statement hold as well, their geometric hypothesis being true of
labels: for every unit $v$ of the supply statement's exterior family $\mathcal{V}^{\mathrm{ext}}$
whose label is rooted in $\Omega_k^{c}$,
\begin{equation}\label{4.21:eq-large-field-site-1}
d_j(A_v)\ \ge\ (2\check{R}_k-1)\,L^{k-j}.
\end{equation}
The bound on the terms $\mathbf{B}'$ holds with $d_j(X)$
replaced by $d_j(A_v)$; no length of $X_v$ occurs. A deep $v$ has $K_v\subset Z$; it belongs to no
exterior family $\mathcal{V}^{\mathrm{ext}}$.
\item[(c)] (Pieces.) Under the prescriptions of the decoupling clause at the large-field site
(Definition~\ref{4.21:def-decoupling-clause}, display~\eqref{4.21:eq-decoupling-clause-a}), what
is not a value resolves into pieces
$\mathfrak{p}=(v;s;\mathsf{y}_1,\mathsf{y}_2,\mathsf{y}_3,\mathsf{y}_4)$, where $s\le 4$ is the link
of the first operation, $\mathsf{y}_1$ is a family of cells disjoint from $\Lambda^{\sim}$,
$\mathsf{y}_2,\mathsf{y}_3,\mathsf{y}_4$ are families of cells disjoint from $\Lambda$, not all four
families are empty, and $\mathsf{y}_1\neq\emptyset$ if $v\in\mathfrak{L}^{\mathrm{II}}$. The labels
are $\operatorname{lab}_0:=K_v$ and, for $t=1,\dots,4$,
\begin{equation*}
\operatorname{lab}_t:=
\begin{cases}
\operatorname{lab}\bigl((\operatorname{lab}_{t-1},\mathcal{R}_{t-1}),\mathsf{y}_t\bigr)
& \text{if } \mathsf{y}_t\neq\emptyset,\\
\operatorname{lab}_{t-1} & \text{if } \mathsf{y}_t=\emptyset,
\end{cases}
\end{equation*}
where $\mathcal{R}_{t-1}$ is the set on which link $t$ reads (the closed set $\mathcal{R}$ of
Lemma~\ref{4.21:lem-first-stage-expansion} for that link) and $\operatorname{lab}(\cdot,\mathsf{y})$
is the label of a piece cut along the family $\mathsf{y}$;
and the domain of $\mathfrak{p}$ is
\begin{equation}\label{4.21:eq-large-field-site-2}
Y_4(\mathfrak{p}):=\operatorname{lab}_4\cup\bigcup\{C: C \text{ a component of } Z
\text{ meeting } \operatorname{lab}_4\}\ \in\ \mathbf{D}_k^{Z},
\end{equation}
with $d_{k,Z}(Y_4(\mathfrak{p}))\le d_k(\operatorname{lab}_4)$. The piece $\mathfrak{p}$ is analytic
over $Y_4(\mathfrak{p})^{+}$ and reads the integration variables of the components of $Z$ contained
in $Y_4(\mathfrak{p})$ only. Always
\begin{equation*}
|\mathfrak{p}|\ \le\ \mathfrak{N}^{A}(v)\,e^{-(\kappa_1-1)\sum_t|\mathsf{y}_t|}.
\end{equation*}
An exterior piece of a deep $v$ has $\bigcup_t\mathsf{y}_t\not\subset Z$, and for every exterior
piece of every $v$, those with $\mathsf{y}_2=\mathsf{y}_3=\mathsf{y}_4=\emptyset$ included,
\begin{equation*}
|\mathfrak{p}|\ \le\ \|v\|_{\star}\,e^{-(\kappa_1-1)\sum_t|\mathsf{y}_t|}.
\end{equation*}
\item[(d)] (Exterior sum.) Put
$\mathfrak{v}^{\mathrm{rel}}(Y_4):=\sum_{\mathfrak{p}:\,Y_4(\mathfrak{p})=Y_4}\mathfrak{p}$. Then
\begin{equation}\label{4.21:eq-large-field-site-3}
\begin{aligned}
\|\mathfrak{v}^{\mathrm{rel}}\|_{a,1}
&:=\sup_{\Delta\not\subset Z}\ \sum_{Y_4\ni\Delta}e^{a\,d_{k,Z}(Y_4)}
\,|\mathfrak{v}^{\mathrm{rel}}(Y_4)|
\ \le\ \varepsilon^{\mathrm{rel}}_L(k),\\
\varepsilon^{\mathrm{rel}}_L(k)
&:=e^{-(\check{R}_k-2)}(1+C_{\Sigma}C_e)^3
\bigl[(1+4h_{\sharp}^{d}C_{\Sigma})\,\mathsf{Q}^{\star}_{a+9}
+4C_{\Sigma}C_{\mathsf{a}}(1)\,\mathsf{b}_k(0)\bigr],
\end{aligned}
\end{equation}
the supremum being over the cubes $\Delta$ of $\pi_k$ not contained in $Z$ and the sum over the
domains $Y_4\in\mathbf{D}_k^{Z}$ containing $\Delta$,
where $\mathsf{Q}^{\star}_r:=\mathsf{Q}_r[\mathfrak{L}^{\mathrm{I}}]$ is the cube-anchored functional
of the family $\mathfrak{L}^{\mathrm{I}}$ computed in the norms $\|\cdot\|_{\star}$, bounded in
Lemma~\ref{4.21:lem-wrapped-mass}.
\item[(e)] Assume in addition the input condition of clause (c) of the boundary-family proposition
for the plain family $\mathfrak{v}$ of the correlation theorem at this step, with constant
$\varepsilon_{\mathrm{br}}(k)$. Let $\mathsf{P}_L$ be the polynomial majorant, given by
Lemma~\ref{4.21:lem-wrapped-mass}, of the bracket in \eqref{4.21:eq-large-field-site-3}. If
\begin{equation}\label{4.21:eq-large-field-site-a}
\sup_{\mathsf{T}\ge\mathsf{T}_{\gamma}}\ \mathsf{P}_L(\mathsf{T})\,
e^{-(\mathsf{T}^{r_{\mathrm{pr}}}-2)}\ \le\ 1,
\end{equation}
then that input condition holds with $\mathfrak{v}$ replaced by
$\mathfrak{v}+\mathfrak{v}^{\mathrm{rel}}$ and $\varepsilon_{\mathrm{br}}(k)$ replaced by
$\varepsilon_{\mathrm{br}}(k)+1$, each member $(\mathfrak{v}+\mathfrak{v}^{\mathrm{rel}})(Y_4)$ of the
new family being analytic
over $Y_4^{+}$; the volume clause of the input condition is supplied by clause (a) of the volume
proposition.
\end{itemize}
\end{proposition}

\begin{proof}[Proof of parts (a) and (b)]
(a) By clause (f1) of the relative format (Definition~\ref{4.21:def-relative-format},
display~\eqref{4.21:eq-relative-format-a}) the region of $v$ is
$X_v=A_v\cup\bigcup_{(Y,m)\in\mathcal{P}_v}Y^{\sharp}$, and by clause (f0) each strip $Y^{\sharp}$ of
a pocket of $v$ touches $A_v$. Let $(Y,m)\in\mathcal{P}_v^{\dagger}$, so that $Y\subset Z$, and let
$C$ be the component of $Z$ containing $Y$. By clause (g) of the geometric lemma of the large-field
site, $Y^{\sharp}\subset\mathsf{C}_C\subset\Lambda_C$; and a cube of $A_v$ touching $Y^{\sharp}$,
which exists by clause (f0), lies in $\Lambda_C\subset Z$. This is the second assertion of (a).

By clause (e) of the geometric lemma of the large-field site, the strips of the pockets in
$\mathcal{P}_v^{\circ}$, which are
live after $\mathbf{R}_k$, lie at distance at least $\check{R}_k$ from $Z$. Therefore $X_v$ meets $Z$
(respectively $\Lambda_C^{\sim}$) only through $A_v$ or through a dead strip, that is, a strip of a
pocket in $\mathcal{P}_v^{\dagger}$; and in the second case, by the assertion just proved, $A_v$ has a
cube in $\Lambda_C$. Since $A_v\subset[A_v]_k=K_v$, in either case $K_v$ meets $Z$ (respectively
$\Lambda_C^{\sim}$). Conversely, $K_v=[A_v]_k$ is the $\pi_k$-hull of $A_v$, each cube of which
contains points of $A_v$, and $Z$ is a union of cubes of $\pi_k$; hence a cube of $K_v$ meeting $Z$
lies in $Z$ and contains points of $A_v\subset X_v$, so that $X_v$ meets $Z$. This gives the first
and the third assertions of (a).

(b) The proof of the value-domain lemma reads the pair $(K_v,Y_1(v))$ through the following
properties: $Y_1(v)\in\mathbf{D}_k^{Z}$; $Y_1(v)\setminus Z=K_v\setminus Z$; and an optimal tree of
$K_v$, in the sense of the value-domain lemma, lies in $Y_1$ and meets every cube of
$K_v\supseteq Y_1\setminus Z$. These are true of
$K_v=[A_v]_k$; the second is immediate from the definition of $Y_1(v)$, since the sets adjoined to
$K_v$ there are components of $Z$.

The region on which $v$ reads, $X_v$ by clause (f1), satisfies $X_v\subset Y_1(v)^{+}$. Indeed,
$A_v\subset K_v\subset Y_1(v)$. The dead strips lie, by part (a), in cores
$\mathsf{C}_C\subset\Lambda_C$ of components $C$ of $Z$ in which $A_v$ has a cube; such a $C$ meets
$K_v$, so it is contained in $Y_1(v)$. The live strips touch $A_v\subset K_v\subset Y_1(v)$, and
are therefore among the strips adjoined to $Y_1(v)$ in forming $Y_1(v)^{+}$
(Definition~\ref{4.21:def-wrapped-classes}).

The geometric hypothesis of the supply statement concerns a connected domain rooted in
$\Omega_k^{c}$: such a domain must leave its component of $\Lambda_k^{c}$, crossing the collar of that
component. It is thus a statement about a connected union of cubes having a cube in $\Omega_k^{c}$.
The label $A_v$ is such a union: it meets $Z_j^{\boxplus}$ or has a point on a strip of level at
most $j$. For $j<k$ both of these sets lie in
$\Omega_{j+1}^{c}\subset\Omega_k^{c}$. Hence the geometric hypothesis applies to $A_v$ and yields
\eqref{4.21:eq-large-field-site-1}.

Finally, the sums of the supply statement are sums over domains of $\mathbf{D}_j$ containing a given
cube, taken under pointwise bounds of the form $e^{-c\,d_j(X)}$ with a fixed rate $c>0$. For the
units $v$ these
pointwise bounds are furnished by clause (f2) of the relative format
(Definition~\ref{4.21:def-relative-format}).
\end{proof}

\begin{lemma}[The large-field-site proposition: the pieces]\label{4.21:prf-pieces-proof}
This lemma is the part of clause~(c) of Proposition~\ref{4.21:prop-large-field-site} that concerns
the links, the bound on the pieces, the reader condition, the domain of a piece and the exterior
pieces of deep units. Assume the following statements:
\begin{enumerate}
\item the decoupling clause at the large-field site, in both its large-field parts
(Definition~\ref{4.21:def-decoupling-clause}, display~\eqref{4.21:eq-decoupling-clause-a});
\item holomorphy of wrapped terms at every link, at the large-field site
(Lemma~\ref{4.21:lem-wrapped-holomorphy});
\item the blindness of the decoupled operators to the exempted set
(Lemma~\ref{4.21:lem-exempted-set});
\item the sourced first-stage expansion at one link (Lemma~\ref{4.21:lem-first-stage-expansion}),
for readers of the class defined in \eqref{4.21:eq-links-readers-a};
\item the minimal form of the third-factor lemma (Definition~\ref{4.21:def-third-factor-form},
display~\eqref{4.21:eq-third-factor-form-a}).
\end{enumerate}
Let $v\in\mathfrak{L}_L$ be a wrapped unit at the large-field site, with the classes of
Definition~\ref{4.21:def-wrapped-classes}, and $K_v=[A_v]_k$. Let
$\mathfrak{p}=(v;s;\mathsf{y}_1,\mathsf{y}_2,\mathsf{y}_3,\mathsf{y}_4)$ be a piece of $v$, with the
labels $\operatorname{lab}_0=K_v$, $\operatorname{lab}_1,\dots,\operatorname{lab}_4$ and the domain
$Y_4(\mathfrak{p})$ of Proposition~\ref{4.21:prop-large-field-site}(c), and write
$Y^{\mathrm{tr}}_t$ for the scaffold domains. Put $\mathcal{R}_0:=X_v$ and, for
$t=2,3,4$,
\begin{equation}\label{4.21:eq-pieces-proof-1}
\begin{split}
\mathcal{R}_{t-1}:={}&X_v\cup\bigcup_{r<t}\bigcup\mathsf{y}_r\\
&\cup\bigl(\text{the hubs touched or met by }X_v,\mathsf{y}_1,\dots,\mathsf{y}_{t-1}\bigr)\\
&\cup\bigl(\text{the boundary layer of }Y^{\mathrm{tr}}_{t-1}\bigr).
\end{split}
\end{equation}
Then:
\begin{enumerate}
\item[(a)] for $v\in\mathfrak{L}^{\mathrm{I}}$, the index set of the $t$-th link is the set
$\mathfrak{Y}$ of Lemma~\ref{4.21:lem-first-stage-expansion}(i) at $\mathcal{R}_{t-1}$, and
\begin{equation}\label{4.21:eq-pieces-proof-2}
\sum_{(\mathsf{y}_t)}\mathfrak{p}=v^{\sharp}\bigl(\operatorname{chain}(\vec{\mathbf{1}})\bigr),
\end{equation}
the all-empty tuple contributing
$\mathrm{val}_v=v^{\sharp}(\operatorname{chain}(\vec{0}))$; for $v\in\mathfrak{L}^{\mathrm{II}}$
each of the four brackets is resolved in the same way;
\item[(b)] every piece satisfies
$|\mathfrak{p}|\le\mathfrak{N}^A(v)\,e^{-(\kappa_1-1)\sum_t|\mathsf{y}_t|}$;
\item[(c)] for every $t$ the reader $(\operatorname{lab}_{t-1},\mathcal{R}_{t-1})$ of the $t$-th link
satisfies the reader condition \eqref{4.21:eq-site-geometry-a};
\item[(d)] $Y_4(\mathfrak{p})\setminus Z\subset\operatorname{lab}_4\subset Y_4(\mathfrak{p})$, an
optimal tree of $\operatorname{lab}_4$ is admissible for $d_{k,Z}(Y_4(\mathfrak{p}))$, the piece
reads sources only in components of $Z$ meeting $\operatorname{lab}_4$, and backgrounds only on
the set on the left-hand side of the inclusion
\begin{equation}\label{4.21:eq-pieces-proof-4}
X_v\cup\bigcup_t\mathsf{y}_t\cup(\text{the hubs touched or met})\subset Y_4(\mathfrak{p})^{+};
\end{equation}
\item[(e)] if $v\in\mathfrak{L}^{\mathrm{dp}}$, $X_v\subset\mathsf{C}_C$ and
$\bigcup_t\mathsf{y}_t\subset Z$, then $\operatorname{lab}_t\subset C$ for all $t$; consequently every
exterior piece of a deep unit has $\bigcup_t\mathsf{y}_t\not\subset Z$ and satisfies
$|\mathfrak{p}|\le\|v\|_{\star}\,e^{-(\kappa_1-1)\sum_t|\mathsf{y}_t|}$.
\end{enumerate}
\end{lemma}

\begin{proof}
From Proposition~\ref{4.21:prop-large-field-site} only clause~(a) and the definitions of the
pieces, of $\operatorname{lab}_t$ and of $Y_4(\mathfrak{p})$ in clause~(c) are used below. The four
layers are read on the global chain, as in Remark~\ref{4.21:rem-global-chain}.

\emph{Links.} For $v\in\mathfrak{L}^{\mathrm{II}}$ the operations to be resolved are the four
brackets of \cite{B-CMP122b}; for $v\in\mathfrak{L}^{\mathrm{I}}$ it is the first operation of
Ba{\l}aban's construction \cite{B-CMP122b}. Each of them is a link in the sense of
Definition~\ref{4.21:def-links-readers} with exempted set $\mathsf{K}=\Lambda^{\sim}$. Indeed, in
\cite{B-CMP122b} the first operation makes changes only inside the domain $\Lambda$, and
$B$, in the notation of \cite{B-CMP122b}, is taken on $\mathfrak{B}_0\subset\Lambda_0$; the
arguments of the layers
$\mathbb{H}_3,\mathbb{H}_2,\mathbb{H}_1$ are supported in $\Lambda^{\sim}$ by the support statements
for these arguments in the relative small-field machine. At $s=3$ the source datum is not a function
supported in $\mathsf{K}\cup\bigcup\mathfrak{F}$
but the field $\mathfrak{j}_3=\tilde{G}_0(\mathsf{s})\mathsf{J}_3$, by the construction of the third
layer in the relative small-field machine. Its primitive source $\mathsf{J}_3$, the right-hand
bracket of \cite[(1.55)]{B-CMP122b}, is supported in $\Lambda$ by the clause of
Lemma~\ref{4.21:lem-exempted-set} that concerns this primitive source. Hence the reader class
defined in \eqref{4.21:eq-links-readers-a} and
Lemma~\ref{4.21:lem-first-stage-expansion}(i) apply, with the condition that the source datum
vanish on a component read as the condition $\mathsf{J}_3=0$ there.

By the first large-field part of the decoupling clause \eqref{4.21:eq-decoupling-clause-a}, by
Lemma~\ref{4.21:lem-wrapped-holomorphy} at the large-field site and by
Lemma~\ref{4.21:lem-first-stage-expansion}, the bracket equals
$\sum_{\mathsf{y}_1\in\mathfrak{Y}}v_{\mathsf{y}_1}$, for the reader
$(K_v,\mathcal{R}_0:=X_v)$. For $v\in\mathfrak{L}^{\mathrm{II}}$ the value $v(\mathsf{R}(0))$ is the
lower end of the link: $\mathcal{Y}(0)$ vanishes off $\Lambda^{\sim}$, and $X_v$ does not meet $Z$.

The piece reads the lower layers $\mathbb{H}_3,\mathbb{H}_2,\mathbb{H}_1$ on $X_v$ and wherever
the layers above them read their background, that is, on $\bigcup_{r<t}\bigcup\mathsf{y}_r$ and on
the hubs touched or met. The same holds, for $v\in\mathfrak{L}^{\mathrm{I}}$, for the term
$\mathsf{y}_1=\emptyset$, which carries the response of $B$ confined to the hubs. These
lower layers are links with $\mathsf{K}=\Lambda$. We take $\mathcal{R}_{t-1}$ as in
\eqref{4.21:eq-pieces-proof-1}; its last member, the boundary layer of $Y^{\mathrm{tr}}_{t-1}$, is
an over-cover in the sense of Remark~\ref{4.21:rem-global-chain}.

Fix $t$. Consider the function of $\mathsf{s}^{(t)}$ obtained from the function already
differentiated and integrated in $\mathsf{s}^{(1)},\dots,\mathsf{s}^{(t-1)}$, the later parameters
being left free. It is a reader of the class defined in \eqref{4.21:eq-links-readers-a} with reading
set $\mathcal{R}:=\mathcal{R}_{t-1}$. Lemma~\ref{4.21:lem-first-stage-expansion}(i), applied to it,
gives the index set $\mathfrak{Y}$ at $\mathcal{R}_{t-1}$. This index set contains the index set of
\cite{B-CMP122b}.

Lemma~\ref{4.21:lem-first-stage-expansion}(i), applied at each of the four links, is a
fundamental-theorem-of-calculus identity in the parameters of that link. Summing the four identities
gives, for $v\in\mathfrak{L}^{\mathrm{I}}$, the identity \eqref{4.21:eq-pieces-proof-2}. In it the
all-empty tuple contributes $\mathrm{val}_v=v^{\sharp}(\operatorname{chain}(\vec{0}))$. For
$v\in\mathfrak{L}^{\mathrm{II}}$ the four brackets telescope, as in \cite{B-CMP122b}. Each bracket
is resolved in the same way, its later links being the layers below $s$. This proves~(a).

\emph{Bound.} By Lemma~\ref{4.21:lem-wrapped-holomorphy}, in its clause on the joint holomorphy of
the chain, the chain is jointly
holomorphic on the product polydisc. Cauchy's estimate in all the parameters $\mathsf{s}$ on that
polydisc, in the form of Lemma~\ref{4.21:lem-first-stage-expansion}(ii), gives
\begin{equation}\label{4.21:eq-pieces-proof-3}
|\mathfrak{p}|\le\mathfrak{N}^A(v)\,e^{-(\kappa_1-1)\sum_t|\mathsf{y}_t|}.
\end{equation}
This proves~(b).

\emph{The reader condition.} We verify \eqref{4.21:eq-site-geometry-a} link by link, with
$c_{\rho}$ the constant of that condition.

Case $t=1$; the reader is $(K_v,X_v)$. We go through the members $\Delta$ that see
$\mathcal{R}_0=X_v$.
\begin{itemize}
\item A member meeting $A_v$ is at distance $0$ from $K_v$.
\item For a member $\Delta$ in the strip $Y^{\sharp}$ of a live large-field region $Y$ we use clauses
(a), (e) and (f) of the lemma on the geometry of the large-field site, with $s_k(Y)$ as in that lemma,
together with
the lower bound on $\rho(\Delta)$ in \eqref{4.21:eq-site-geometry-b}. They give
\begin{equation*}
\operatorname{dist}_k(\operatorname{pa}(\Delta),K_v)\le s_k(Y)+1\le101\check{R}_k+1
\le c_{\rho}\cdot10\check{R}_k\le c_{\rho}\,\rho(\Delta).
\end{equation*}
\item No member meets a dead strip, since such a strip lies in $\Lambda$.
\item Let $\Delta$ be a source cell at a hub met by $X_v$. By
Proposition~\ref{4.21:prop-large-field-site}(a), $K_v$ meets that hub. Hence the distance is at most
$h_{\sharp}\le c_{\rho}$.
\end{itemize}

Case $t\ge2$. The label is $\operatorname{lab}_{t-1}$. A member seeing $\mathcal{R}_{t-1}$ falls
into one of the following cases.
\begin{itemize}
\item It sees $X_v$. Then it is treated as at $t=1$, since $K_v\subset\operatorname{lab}_{t-1}$.
\item It belongs to an earlier family $\mathsf{y}_r$, $r<t$. Its parent lies in
$\operatorname{lab}_{t-1}$, so the distance is $0$.
\item It is a source cell at a hub touched or met before. A cube of $\operatorname{lab}_{t-1}$
shares a wall with that hub or lies in it, so the distance is at most $h_{\sharp}\le c_{\rho}$.
\item It touches the rim of $\operatorname{lab}_{t-1}$; the distance is at most $1$.
\item It touches the rim of a live strip; the estimate of the case $t=1$ holds with $1$ added.
\item It touches the rim of a component $C\subset Y^{\mathrm{tr}}_{t-1}$ of $Z$, which
$\operatorname{lab}_{t-1}$ meets. By clause~(f) of the geometry lemma, the nearest hub is at
distance at least $9\check{R}_k$ from any rim of $Z$. Hence a layer cell inside $C$ along
$\partial C$ has $\rho(\Delta)\ge9\check{R}_k-2$, and
\begin{equation*}
\operatorname{dist}_k(\operatorname{pa}(\Delta),\operatorname{lab}_{t-1})\le100\check{R}_k+1
\le c_{\rho}(9\check{R}_k-2)\le c_{\rho}\,\rho(\Delta).
\end{equation*}
\end{itemize}
The boundary layer of $Y^{\mathrm{tr}}_{t-1}$ meets no hub, so no further case arises from it.
This proves~(c).

\emph{Domain.} Let $Y_4=Y_4(\mathfrak{p})$ be as in
Proposition~\ref{4.21:prop-large-field-site}(c). Then $Y_4\setminus Z\subset\operatorname{lab}_4$,
and $\operatorname{lab}_4\subset Y_4$ by the definition of $Y_4$. An optimal tree of
$\operatorname{lab}_4$ is admissible for $d_{k,Z}(Y_4)$, by the argument of the optimal-tree
comparison for relative distances. By Lemma~\ref{4.21:lem-first-stage-expansion}(iii), sources are
read on the hubs touched or met, that is, in components of $Z$ meeting $\operatorname{lab}_4$.
Backgrounds are read on the set on the left-hand side of \eqref{4.21:eq-pieces-proof-4}, which lies
in $Y_4^{+}$. This
proves~(d).

\emph{Exterior pieces of a deep unit leave $Z$ by a family.} Let $v\in\mathfrak{L}^{\mathrm{dp}}$
with $X_v\subset\mathsf{C}_C$, and suppose $\bigcup_t\mathsf{y}_t\subset Z$. We claim
$\operatorname{lab}_t\subset C$ for all $t$.

First, $\operatorname{lab}_0=K_v\subset\Lambda_C$. This follows from clause~(g) of the geometry
lemma: a cube of level at most $k$ meeting $\mathsf{C}_C\subset Z''_k$ lies in
$(Z''_k)^{(1)}$.

Now suppose $\operatorname{lab}_{t-1}\subset C$. Then $Y^{\mathrm{tr}}_{t-1}=C$, since a deep unit
has $\mathcal{P}_v^{\circ}=\emptyset$, and therefore $\mathcal{R}_{t-1}\subset C$. Let
$\mathsf{y}_a$ be a component of $\mathsf{y}_t$. It is connected and lies in $Z$. Hence it lies in
one component $C'$ of $Z$: two components of $Z$ are at distance at least $\check{R}_k$ from each
other, and the hub adjacency stays within one component. The component $\mathsf{y}_a$ sees
$\mathcal{R}_{t-1}\subset C$, so $C'=C$. Its parents therefore lie in $C$. The corridor cubes
$\mathrm{Cor}_a$ and the hub-bridge cubes $\mathrm{Hub}_a$ that the label adds lie in the smallest
box containing their two ends, by the box property of corridors and hub bridges, and $C$ is a box.
Hence $\operatorname{lab}_t\subset C$.

By induction, $\operatorname{lab}_4\subset C\subset Z$, so the piece is interior: its label
$\operatorname{lab}_4$ lies in $Z$. Consequently
every exterior piece of a deep unit has $\bigcup_t\mathsf{y}_t\not\subset Z$. The minimal form of
the third-factor lemma, \eqref{4.21:eq-third-factor-form-a}, therefore applies to every exterior
piece of a deep unit, whichever of its families are non-empty. With the starred norm of
Definition~\ref{4.21:def-wrapped-classes}, it gives
\begin{equation*}
|\mathfrak{p}|\le\|v\|_{\star}\,e^{-(\kappa_1-1)\sum_t|\mathsf{y}_t|}.
\end{equation*}
This proves~(e).
\end{proof}

\begin{proof}[Proof of Proposition~\ref{4.21:prop-large-field-site}, parts (d) and (e)]
\label{4.21:prf-exterior-sum-proof}
\emph{Part (d).} Throughout, $\mathfrak{v}^{\mathrm{rel}}(Y_4)$ is the sum of the pieces
$\mathfrak{p}$ of part (c) with $Y_4(\mathfrak{p})=Y_4$. Fix a cube $\Delta\not\subset Z$. We bound
\begin{equation*}
\sum_{Y_4\ni\Delta}e^{a\,d_{k,Z}(Y_4)}\,\bigl|\mathfrak{v}^{\mathrm{rel}}(Y_4)\bigr|
\end{equation*}
by $\varepsilon^{\mathrm{rel}}_L(k)$, uniformly in $\Delta$.

\emph{Reduction to exterior pieces.} By part (c), $Y_4(\mathfrak{p})$ is the union of
$\operatorname{lab}_4$ and of the components of $Z$ meeting $\operatorname{lab}_4$. Hence if
$\Delta\in Y_4(\mathfrak{p})$ and $\Delta\not\subset Z$, then $\Delta\in\operatorname{lab}_4$: only
exterior pieces occur. Bounding $|\mathfrak{v}^{\mathrm{rel}}(Y_4)|$ by the sum of the moduli of its
pieces and using $d_{k,Z}(Y_4(\mathfrak{p}))\le d_k(\operatorname{lab}_4)$ from part (c), the left
member is at most $\mathsf{Q}_a[\mathfrak{L}^4]$ restricted to the exterior pieces, where
\begin{align*}
\mathfrak{L}^1&:=\mathfrak{L}^{\mathrm{I}}\cup\bigsqcup_{s=1}^{4}(\mathfrak{L}_L)'_{(s)},\\
\mathfrak{L}^t&:=\mathfrak{L}^{t-1}\cup(\mathfrak{L}^{t-1})'\qquad(t=2,3,4),
\end{align*}
in which $(\mathfrak{L}_L)'_{(1)},\dots,(\mathfrak{L}_L)'_{(4)}$ are four disjoint copies of
$(\mathfrak{L}_L)'$,
the prime being that of Lemma~\ref{4.21:lem-source-anchored-sums}, and all readers are taken in
the norms $\|\cdot\|_{\star}$. This is a majorant book-keeping. The intermediate readers of a deep
$v$ carry the weight $\|v\|_{\star}$, although the families of cut cells chosen for them so far may
be empty or may lie inside $Z$. Only tuples whose final label leaves $Z$ are summed, and for each of
these part (c) gives exactly the bound
$|\mathfrak{p}|\le\|v\|_{\star}e^{-(\kappa_1-1)\sum_t|\mathsf{y}_t|}$ that
$\mathsf{Q}_a[\mathfrak{L}^4]$ assigns to it. Dropping the restriction to exterior pieces only adds
non-negative terms.

\emph{Smallness.} Let $\operatorname{lab}_4$ hold a cube $\Delta\not\subset Z$.

If $v\in\mathfrak{L}^{\mathrm{II}}\cup\mathfrak{L}^{\circ}$, then $\operatorname{lab}_4$ holds a cube
containing a point of a live strip and a cube containing a point of $Z$. For
$v\in\mathfrak{L}^{\mathrm{II}}$ the second cube comes from $\mathsf{y}_1\ne\emptyset$: this family
holds a source cell, and source cells lie in $Z$ by part (f) of the geometry lemma of the
large-field site. By part (e) of that lemma the two points are at distance at least
$\check{R}_k$. Hence $d_k(\operatorname{lab}_4)\ge\check{R}_k-2$, by the lower bound on the
level-$k$ size of a label in terms of the distance between two of its points.

If $v\in\mathfrak{L}^{\mathrm{sh}}\cup\mathfrak{L}^{\mathrm{dp}}$, then $\operatorname{lab}_4$ holds a
cube in some $\Lambda_C$, by part (a) of Proposition~\ref{4.21:prop-large-field-site}. It also holds
$\Delta\not\subset Z$. Part (f) of the geometry lemma then gives
$d_k(\operatorname{lab}_4)\ge 9\check{R}_k-1$.

Put $D:=\check{R}_k-2$. In both cases $d_k(\operatorname{lab}_4)\ge D$, since
$9\check{R}_k-1\ge\check{R}_k-2$. Consequently
\begin{equation*}
e^{a\,d_k(\operatorname{lab}_4)}\le e^{-D}\,e^{(a+1)\,d_k(\operatorname{lab}_4)},
\end{equation*}
and we obtain $\mathsf{Q}_a\le e^{-D}\mathsf{Q}_{a+1}$ on the pieces kept. Therefore
\begin{equation}\label{4.21:eq-exterior-sum-proof-1}
\sum_{Y_4\ni\Delta}e^{a\,d_{k,Z}(Y_4)}\bigl|\mathfrak{v}^{\mathrm{rel}}(Y_4)\bigr|
\le e^{-(\check{R}_k-2)}\,\mathsf{Q}_{a+1}[\mathfrak{L}^4].
\end{equation}

\emph{The chain of links.} Apply parts (c) and (e) of Lemma~\ref{4.21:lem-source-anchored-sums}.
For $t=2,3,4$ they give
\begin{equation}\label{4.21:eq-exterior-sum-proof-2}
\mathsf{Q}_r[\mathfrak{L}^t]
\le\mathsf{Q}_r[\mathfrak{L}^{t-1}]+C_{\Sigma}\,\mathsf{T}^{*}_{r+2}[\mathfrak{L}^{t-1}]
\le(1+C_{\Sigma}C_e)\,\mathsf{Q}_{r+2}[\mathfrak{L}^{t-1}].
\end{equation}
For the first family they give
\begin{equation}\label{4.21:eq-exterior-sum-proof-3}
\mathsf{Q}_r[\mathfrak{L}^1]\le\mathsf{Q}^{\star}_r+4C_{\Sigma}\,\mathsf{T}^{*}_{r+2}[\mathfrak{L}_L].
\end{equation}

Part (c) of that lemma requires $\kappa_1-1\ge c_l r+\lambda_*$. In the chain below parts (c)
and (e) are used at indices up to $r+2=a+9$, so it suffices that
\begin{equation*}
\kappa_1-1\ge c_l(a+9)+\lambda_* .
\end{equation*}
This holds by the condition on $\kappa_1$ among \eqref{4.21:eq-format-propagation-a}. That
condition is available because $a+9\le\tfrac85\kappa$, which follows from $\kappa\ge180$, a
condition contained in the lower bound $\kappa\ge\kappa_d$ in the hypotheses of
Lemma~\ref{4.21:lem-wrapped-cube-bound}.

Apply \eqref{4.21:eq-exterior-sum-proof-2} with $r=a+1$, $a+3$, $a+5$, and then
\eqref{4.21:eq-exterior-sum-proof-3} with $r=a+7$. This gives
\begin{equation}\label{4.21:eq-exterior-sum-proof-4}
\begin{split}
\mathsf{Q}_{a+1}[\mathfrak{L}^4]
&\le(1+C_{\Sigma}C_e)^3\,\mathsf{Q}_{a+7}[\mathfrak{L}^1]\\
&\le(1+C_{\Sigma}C_e)^3\bigl(\mathsf{Q}^{\star}_{a+7}
+4C_{\Sigma}\,\mathsf{T}^{*}_{a+9}[\mathfrak{L}_L]\bigr).
\end{split}
\end{equation}

\emph{The reception functional $\mathsf{T}^{*}_{a+9}[\mathfrak{L}_L]$.} Let $\Delta$ be a member of
$\pi^{\mathfrak{s}}$, that is, a cell, disjoint from $\Lambda^{\sim}$. Consider the wrapped terms $F$
with $X_F$ meeting $\Delta$. None
of them is deep, since $X_F\not\subset\Lambda$. We apply part (iv) of
Lemma~\ref{4.21:lem-probe-mass}, with $\beta=1$ and with the cells as probes. The depth data of
\eqref{4.21:eq-wrapped-terms-a} are taken with $\varrho:=\rho$ and $N_*:=0$. We check the two depth
conditions there.
\begin{itemize}
\item The first depth condition is the first condition of \eqref{4.21:eq-site-geometry-b}.
\item For the second: $\Delta$ lies in a live strip. A chain of wall-neighbouring cells from
$\Delta$ to a source cell (such a cell has level $k$ and lies at a hub) leaves the collar of width
$\tfrac52\check{R}_m$ of part (c) of the geometry lemma.
\end{itemize}
Hence the weighted mass of these terms at $\Delta$ is at most
\begin{equation*}
e^{\rho(\Delta)}\,C_{\mathsf{a}}(1)\,\mathsf{b}_k(0).
\end{equation*}

Suppose moreover that $\Delta$ is a source cell at a hub $H$. By part (a) of
Proposition~\ref{4.21:prop-large-field-site}, the $F$ with $X_F$ meeting $H$ are terms of
$\mathfrak{L}^{\mathrm{I}}$ whose $K_F$ meets $H$. The hub $H$ is a box of cubes of $\pi_k$ with
sides at most $h_{\sharp}$, so it consists of at most $h_{\sharp}^d$ such cubes. The terms anchored
at each of them contribute at most $\mathsf{Q}^{\star}_{a+9}$. Their total is therefore at most
$h_{\sharp}^d\,\mathsf{Q}^{\star}_{a+9}$.

This is the one place at the large-field site where the sum over all admissible
localization domains $X$ in Ba{\l}aban's bound is performed. Its terms are the classes into which
the proof of
Lemma~\ref{4.21:lem-wrapped-mass} divides the wrapped terms. Each class is paid for by one named
source of smallness: the decay of a cut cell, the weight of the term's own label, a count absorbed
by the factor $e^{-(\check{R}_k-2)}$ of \eqref{4.21:eq-exterior-sum-proof-1}, or the bound of the
third-factor lemma.

The two bounds just obtained are read in the reception functional, by its definition in
Lemma~\ref{4.21:lem-source-anchored-sums}; they give
\begin{equation*}
\mathsf{T}^{*}_{a+9}[\mathfrak{L}_L]\le h_{\sharp}^d\,\mathsf{Q}^{\star}_{a+9}
+C_{\mathsf{a}}(1)\,\mathsf{b}_k(0).
\end{equation*}
Insert this into \eqref{4.21:eq-exterior-sum-proof-4} and use
$\mathsf{Q}^{\star}_{a+7}\le\mathsf{Q}^{\star}_{a+9}$. Then insert the result into
\eqref{4.21:eq-exterior-sum-proof-1}. The left member is bounded by
\begin{equation*}
e^{-(\check{R}_k-2)}(1+C_{\Sigma}C_e)^3\Bigl[(1+4h_{\sharp}^dC_{\Sigma})\,\mathsf{Q}^{\star}_{a+9}
+4C_{\Sigma}C_{\mathsf{a}}(1)\,\mathsf{b}_k(0)\Bigr]=\varepsilon^{\mathrm{rel}}_L(k),
\end{equation*}
uniformly in $\Delta\not\subset Z$. This is part (d).

\emph{Part (e).} The family bound of part (e) is stated in the norm $\|\cdot\|_a$, and
$\|\cdot\|_a\le\|\cdot\|_{a,1}$. By part (d), therefore,
$\|\mathfrak{v}^{\mathrm{rel}}\|_a\le\varepsilon^{\mathrm{rel}}_L(k)$.

Now let $\mathsf{P}_L$ be, as in part (e), the polynomial majorant provided by
Lemma~\ref{4.21:lem-wrapped-mass} for the bracket in the definition of
$\varepsilon^{\mathrm{rel}}_L(k)$ in part (d), the constant factor $(1+C_{\Sigma}C_e)^3$ in front of
the bracket included, evaluated at $\mathsf{T}=\check{\mathsf{T}}_k$ by the comparison of
$\check{R}_k$ with $\check{\mathsf{T}}_k$ used in that lemma; thus
$e^{\check{R}_k-2}\,\varepsilon^{\mathrm{rel}}_L(k)\le\mathsf{P}_L(\check{\mathsf{T}}_k)$. Then
\begin{equation*}
\begin{split}
\varepsilon^{\mathrm{rel}}_L(k)
&\le e^{-(\check{R}_k-2)}\,\mathsf{P}_L(\check{\mathsf{T}}_k)\\
&\le\mathsf{P}_L(\check{\mathsf{T}}_k)\,e^{-(\check{\mathsf{T}}_k^{\,r_{\mathrm{pr}}}-2)}\\
&\le1 .
\end{split}
\end{equation*}
The first inequality is Lemma~\ref{4.21:lem-wrapped-mass} so evaluated; the second is the
comparison $\check{R}_k\ge\check{\mathsf{T}}_k^{\,r_{\mathrm{pr}}}$ of $\check{R}_k$ with
$\check{\mathsf{T}}_k$; the third is the condition on $\mathsf{P}_L$ in part (e), a supremum over
$\mathsf{T}\ge\mathsf{T}_{\gamma}$, applied at $\mathsf{T}=\check{\mathsf{T}}_k\ge\mathsf{T}_{\gamma}$.

By the hypothesis of part (e), the plain family $\mathfrak{v}$ alone satisfies the family bound at
$\varepsilon_{\mathrm{br}}(k)$, that is $\|\mathfrak{v}\|_a\le\varepsilon_{\mathrm{br}}(k)$. Hence
\begin{equation*}
\|\mathfrak{v}+\mathfrak{v}^{\mathrm{rel}}\|_a
\le\varepsilon_{\mathrm{br}}(k)+\varepsilon^{\mathrm{rel}}_L(k)
\le\varepsilon_{\mathrm{br}}(k)+1 .
\end{equation*}
The strip summation theorem accepts any family of numbers, or of functions in the sup norm; it
therefore applies to $\mathfrak{v}+\mathfrak{v}^{\mathrm{rel}}$ at $\varepsilon_{\mathrm{br}}(k)+1$.
Each $\mathfrak{v}^{\mathrm{rel}}(Y_4)$
is analytic over $Y_4^{+}$, because each of its pieces is, by part (c).
\end{proof}

\begin{remark}
\begin{enumerate}
\item The constant
\begin{equation*}
\alpha^{\sharp}=C_{\mathrm{fib}}(\varepsilon_{\mathrm{br}}+\mathsf{K}_{\mathrm{val}})
e^{-\frac12\beta_{\sharp}\kappa(m_{\sharp}-1)}
\end{equation*}
requires $\varepsilon_{\mathrm{br}}$ to be bounded, not small. Here $\alpha^{\sharp}$ is the small
factor of the correlation theorem, $\varepsilon_{\mathrm{br}}$ is the smallness constant of the
plain boundary family and $m_{\sharp}$ is the strip width; $C_{\mathrm{fib}}$,
$\mathsf{K}_{\mathrm{val}}$, $\beta_{\sharp}$, $\mathsf{w}_L$, $\mathsf{K}^0$ and $m_{\mathrm{lo}}$
are quantities fixed in the correlation theorem. This is why the smallness condition
of the correlation theorem that involves $\varepsilon_{\mathrm{br}}$ can do without the factor
$\mathsf{w}_L^{-1}$ that it would otherwise carry. In the two conditions of the
correlation theorem on its constants carrying the index $\sharp$, and in its quantity
$m_{\mathrm{lo}}$, the factor $1+\mathsf{K}^0$ that occurs in them is taken as $2+\mathsf{K}^0$.
\item The healing of a term's own pocket. An own pocket receives no strip block, no corridor and
no amplitude: its strip lies in the never-cut set $\Lambda_C$, and the term is a reader of the hub.
Terms are anchored by the core $\mathsf{C}_C$, level by level, at the cubes of their own level.
With $h$ the level that enters Lemma~\ref{4.21:lem-wrapped-mass} through $\check{R}_{h+1}$, the
shallow ones number at most
a constant multiple of $\check{R}_{h+1}^d+N(g_k)$ per unit of $\mathsf{b}$ and need no decay. The deep
ones
are bounded by the deep-term lemma, which is derived from the first input hypothesis of the step.
The partition-of-unity anchor lemma, with its telescoping over the $\zeta_{\square}$ and the
supremum $\mathsf{W}$ over anchors that it carries, is used on no wrapped term.
\item A mixed term, with one dead and one live pocket, belongs to $\mathfrak{L}^{\circ}$.
\item A term in the ring $\Lambda_C\setminus\mathsf{C}_C$, a straddling term, and a term of level
$j>h$ are all shallow; no depth is claimed for them.
\end{enumerate}
\end{remark}

\begin{proposition}[The relative machine at the birth of new boundary terms]
\label{4.21:prop-birth-site}
Fix the step $\mathbf{R}_k$. Let $Y_i$ be the large-field regions created at $\mathbf{R}_k$ and
$Y_i^{\sharp}$ their strips. Put
\begin{equation*}
\mathsf{I}:=\{Y_i^{\sharp}\}\quad\text{(the strips born at }\mathbf{R}_k\text{)},\qquad
W:=\bigcup\mathsf{I}.
\end{equation*}
A large-field region created at a step earlier than $k$ and live after $\mathbf{R}_k$ (live in the
sense of the glossary, Notation~\ref{0.2:not-glossary}: still carrying its large-field factor at
the level read) is called an \emph{older live region}, and its strip an \emph{older live strip}.
Assume the following hypotheses.
\begin{enumerate}
\item[(1)] The input statement of the step $\mathbf{R}_k$, which supplies the constituents
$\mathfrak{v}(Y_4,\cdot)$ of the step, $Y_4$ being their piece domains, with their bounds in
$\mathfrak{N}^{A}$.
\item[(2)] The strip theorem, the corollary on the strip potentials $V^{\sharp}$, and the
combinatorial lemma for strips with its clauses (u1)--(u7), each as stated for the step
$\mathbf{R}_k$.
\item[(3)] Parts (a)--(c) of the skeleton lemma, which attaches to every
$X\in\mathbf{D}_k^{\mathsf{I}}$ its skeleton $\operatorname{sk}(X)$, and clause (d) of the distance
inequalities of the skeleton construction.
\item[(4)] The parameters of the step obey the floor $\frac14\beta_s\kappa\ge\lambda_\theta$ of
\eqref{4.21:eq-format-propagation-a}.
\end{enumerate}
Then the following hold.
\begin{itemize}
\item[(a)] The strip theorem, the corollary on the strip potentials and clauses (u1)--(u7) of the
combinatorial lemma hold without change of wording, with $V^{\sharp}(Y)$, the activities and
$\mathbf{R}'^{(k)}(X;w)$ analytic over $Y^{+}$, $X'^{+}$ and $X^{+}$ respectively; and they hold in
the weighted norm $\mathfrak{N}^{A}$ (and in $\mathfrak{N}^{A,\mathbb{C}}$). Older live strips enter
no sum, no length and no compatibility relation.
\item[(b)] Second group of the boundary terms created at $\mathbf{R}_k$, those of
\cite[(1.100)]{B-CMP122b}: $X^{+}=X$. The displays \cite[(1.100)]{B-CMP122b} and
\cite[(2.30)--(2.32)]{B-CMP119}, and the inclusion statement, hold unchanged.
\item[(c)] First group. For $X\in\mathbf{D}_k^{\mathsf{I}}$ put $A:=\operatorname{sk}(X)$; for $X$
meeting no $Y_i^{\sharp}$ put $A:=X$ (such a term is plain, and the constant $c_1$ entering its
bound is $e^{-p_0(g_k)}$, $p_0(\cdot)$ being the degree function). Write
$A(X)$ for the region so attached to $X$, and define
\begin{gather*}
\mathbf{B}'^{(k)}(A,\mathcal{S}):=\sum_{X:\,A(X)=A}\mathbf{B}'^{(k)}_{0}(X,\mathcal{S}),
\\
\begin{split}
\mathcal{P}_A:={}&\{(Y_i,k):Y_i^{\sharp}\text{ touches }A\}\\
&\cup\{(Y,m):Y\text{ created at step }m<k,\ Y\text{ live after }\mathbf{R}_k,\
Y^{\sharp}\text{ touches }A\},
\end{split}
\end{gather*}
where $\mathbf{B}'^{(k)}_{0}(X,\mathcal{S})$ are the boundary terms of the first group
created at $\mathbf{R}_k$ in \cite[(1.90)--(1.99)]{B-CMP122b}. Then, with the data
$(A,\mathcal{P}_A)$, the terms
$\mathbf{B}'^{(k)}(A,\cdot)$ satisfy the clauses \eqref{4.21:eq-relative-format-a} of the relative
format (Definition~\ref{4.21:def-relative-format}), and
\begin{equation}\label{4.21:eq-birth-site-1}
\mathfrak{N}^{A}\bigl(\mathbf{B}'^{(k)}(A,\cdot)\bigr)
\le e^{\lambda_\theta}C_{99}\,c_1^{\flat}(g_k)\,e^{-(1+\frac14\beta_s)\kappa d_k(A)}.
\end{equation}
\end{itemize}
\end{proposition}

\begin{proof}
\emph{Part (a).} The strip theorem and clauses (u3)--(u7) of the combinatorial lemma are statements
about numbers indexed by the domain classes $\mathbf{D}_k^{Z}$ and $\mathbf{D}_k^{\sharp}$. We feed
them with the numbers $\mathfrak{N}^{A}(\mathfrak{v}(Y_4,\cdot))$ of the constituents
$\mathfrak{v}(Y_4,\cdot)$ supplied at the input, $Y_4$ being their piece domains. They then return
the norms $\mathfrak{N}^{A}$ of the outputs. Indeed, every operation performed between the input
hypothesis and \cite[(1.99)]{B-CMP122b}, that is, in \cite[(1.90)--(1.99)]{B-CMP122b}, is of one of
three kinds:
\begin{itemize}
\item a sum;
\item a product, under which the strong sets $\mathcal{S}$ of the factors go to their union at the
surviving coordinates; the norm $\mathfrak{N}^{A}$ is submultiplicative under such products, by
\eqref{4.21:eq-links-readers-c};
\item an integral $\mathbf{T}'_k(X_j)$ over variables localised in $X_j$, under which the slots
exterior to $X_j$ are parameters, so that, strong-set component by strong-set component,
\begin{equation*}
\bigl|\mathbf{T}'_k(X_j)F\bigr|\le \mathbf{T}'_k(X_j)1\cdot\sup|F|.
\end{equation*}
\end{itemize}
That the terms have this structure is the content of the corollary on the preparation site and of
part (b) of the corollary on complex branches: every term created at the large-field site of the
step is a finite sum of stored terms $(c,S,\mathcal{S})$, where $\mathcal{S}$ is the union of the
strong sets of its constituents at the surviving coordinates. Hence the weighted sum over
$\mathcal{S}$ defining $\mathfrak{N}^{A}$ is bounded term by term, and the resulting bounds are
indexed by the labels $A$ rather than by the regions $X$. The same reasoning,
read in $\mathfrak{N}^{A,\mathbb{C}}$, gives the statements in that norm.

Clause (u1). A piece domain $Y_4\subset C^{\sharp}$, where $C$ is a component of $\Lambda_k^{c}$
and $C^{\sharp}$ is the enlargement of $C$ appearing in the statements of hypothesis~(2), touches
no older live
strip: the older live
strips lie in other components of $\Lambda_k^{c}$, at distance at least $\check{R}_k-m_{\sharp}$, by
the separation property of the strips.

Clause (u2). The integral $\mathbf{T}'_k(X_j)$ integrates variables localised in $X_j$. The
variables on an older live strip belong to the operation $\mathbf{T}_k(Z_k\cap Z^{c})$
\cite{B-CMP122b}; under $\mathbf{T}'_k(X_j)$ they are therefore parameters. Polymers are compatible
if and only if their parts off $\bigcup Y_i^{\sharp}$ neither meet nor touch; this is the
compatibility relation of \cite{B-CMP122b}, which reads
\begin{quote}
we take components of $X'_0\setminus\bigcup Y_i$, and we consider two such components as connected,
if they are contained in one localization domain $Y$ of a term in the product. This means that each
component has a correct tree graph decay factor
\end{quote}
together with clause (u2) of the combinatorial lemma. Two polymers that share only the background on
an older live strip factorise. Thus older live strips enter no sum, no length and no compatibility
relation, and part (a) is proved.

\emph{Part (b).} An older live strip lies in the part of $Z_k$ not treated at the step $\mathbf{R}_k$,
by parts (b) and (e) of the geometry lemma. A cube touching it therefore lies in $Z_k^{\boxplus}$.
A region $X$ of the second group is disjoint from $(Z_k\cup\bigcup Y_i)^{\boxplus}$. Hence no older
live strip touches $X$, so $X^{+}=X$, and \cite[(1.100)]{B-CMP122b},
\cite[(2.30)--(2.32)]{B-CMP119} and the inclusion statement are untouched.

\emph{Part (c), clause (f2): the bound.} For $X$ of the first group, $X\cap W$ is a union of
strips, by clause
(u2) of the combinatorial lemma. Moreover $X\not\subset W$: a term lying inside a strip is a member
of $V^{\sharp}(Y_i^{\sharp})$, kept under $\mathbf{T}'$, and belongs to the class of terms carried in
the tube over the strip, not to the first group. By clause (d) of the distance inequalities of the
skeleton construction and parts (b) and (c) of the skeleton lemma, summing over the $X$ with
$A(X)=A$ gives, with $C_{99}$ the constant of \cite[(1.99)]{B-CMP122b},
\begin{equation}\label{4.21:eq-birth-site-2}
\mathfrak{N}^{A}\bigl(\mathbf{B}'^{(k)}(A,\cdot)\bigr)
\le e^{\lambda_\theta(d_k(A)+1)}C_{99}\,c_1^{\flat}(g_k)\,e^{-(1+\frac12\beta_s)\kappa d_k(A)}.
\end{equation}
The floor $\frac14\beta_s\kappa\ge\lambda_\theta$ of hypothesis~(4) gives
$\lambda_\theta d_k(A)\le\frac14\beta_s\kappa d_k(A)$, since $d_k(A)\ge0$. Hence
\begin{align*}
e^{\lambda_\theta(d_k(A)+1)}e^{-(1+\frac12\beta_s)\kappa d_k(A)}
&=e^{\lambda_\theta}\,e^{\lambda_\theta d_k(A)-(1+\frac12\beta_s)\kappa d_k(A)}\\
&\le e^{\lambda_\theta}\,e^{-(1+\frac14\beta_s)\kappa d_k(A)},
\end{align*}
and inserting this into \eqref{4.21:eq-birth-site-2} gives \eqref{4.21:eq-birth-site-1}.

\emph{Part (c), clause (f1): the reading region.} By part (a) of the skeleton lemma,
\begin{equation*}
X\subset A\cup\bigcup\{Y_i^{\sharp}:Y_i^{\sharp}\text{ touches }A\}.
\end{equation*}
An older live strip read by a constituent $\mathfrak{v}(Y_4,\cdot)$ touches a cube of $Y_4$ lying at
distance at least $1$ outside the set $Z^{\sharp}$ appearing in the statements of
hypothesis~(2); this cube lies in
$X\setminus W\subset A$, so the strip
touches $A$ and its region belongs to $\mathcal{P}_A$. Hence $\mathbf{B}'^{(k)}(A,\cdot)$ reads its
arguments only on $A\cup\bigcup_{(Y,m)\in\mathcal{P}_A}Y^{\sharp}$. For analyticity,
\cite{B-CMP122b} states that the terms are ``defined on the space
$\tilde{U}^c_k(X,\tilde{\alpha}_0,\tilde{\alpha}_1)$'', and this is said of every $X$ before the
division into groups.

\emph{Part (c), clause (f0): the label.} By part (a) of the skeleton lemma, a cube of
$X\setminus W\subset A$
touches a strip $Y_i^{\sharp}$, so that $(Y_i,k)\in\mathcal{P}_A$. By the separation property of the
strips this cube is at distance at least $\check{R}_k-m_{\sharp}-1>0$ from every other component of
$\Lambda_k^{c}$, and it is off $Y_i$; hence it lies in $\Lambda_k\subset\Omega_k$, and
$A\cap\Omega_k\neq\emptyset$.

\emph{Part (c), clause (f3): dependence on the scaffold only.} The skeleton $\operatorname{sk}$
reads only
$(\pi_k,\mathsf{I},X)$; hence $A$ and $\mathcal{P}_A$ read the scaffold only.

\emph{Part (c), clause (f4).} Clause (f4) follows from clause (f0), as recorded with
Definition~\ref{4.21:def-relative-format}.
\end{proof}

\begin{definition}[The complexified gauge group and the polar size of an element]
\label{4.22:not-complex-group}
Let $G=SU(N)$ with $N\ge 2$ fixed, and let $\operatorname{tr}$ be the trace normalised by
$\operatorname{tr}\mathbb{1}=1$.
\begin{enumerate}
\item $\mathfrak{g}$ is the real vector space of Hermitian traceless $N\times N$ matrices. This is
the convention of \cite{B-CMP119}, where a field is written $U'=\exp(i\xi A')$, in the notation
of that paper, with $A'$ taking values in the algebra $\mathfrak{g}^{\mathbb{C}}$.
\item The complexification is
\begin{equation*}
\mathfrak{g}^{\mathbb{C}}=\mathfrak{g}\oplus i\,\mathfrak{g}=\mathfrak{sl}(N,\mathbb{C}),
\end{equation*}
and $G^{\mathbb{C}}=SL(N,\mathbb{C})$. A field is called \emph{real} if it takes values in
$\mathfrak{g}$, respectively in $G$.
\item For $g\in G^{\mathbb{C}}$ and a matrix $A$ we write $\operatorname{Ad}_gA:=gAg^{-1}$ (in
this definition $g$ denotes an element of $G^{\mathbb{C}}$, not a coupling); in
Ba{\l}aban's notation this conjugation is written $R(g)A$, where the letter $R$ has nothing to do
with the radius $R$ of the glossary.
\item For $A,B\in\mathfrak{g}^{\mathbb{C}}$ we put $\langle A,B\rangle:=\operatorname{tr}AB$.
This pairing is $\mathbb{C}$-bilinear, invariant under $\operatorname{Ad}_g$ for every
$g\in G^{\mathbb{C}}$, and positive on $\mathfrak{g}$.
\item $\|\cdot\|$ is the operator norm of matrices acting on $\mathbb{C}^N$ with its Hermitian
norm. In his estimates Ba{\l}aban uses the operator norm for matrices \cite[(19)]{B-CMP98};
accordingly $|\cdot|$ denotes the operator norm on all $N\times N$ matrices, $|\cdot|=\|\cdot\|$,
so that a bound of the form $|Z|\le z$ and a bound of the form $\|Y\|<\varepsilon$ measure one
and the same size. The two letters are kept because $|\cdot|$ enters the estimates only
through the properties
\begin{equation}\label{4.22:eq-complex-group-a}
|XAY|\le\|X\|\,|A|\,\|Y\|,\qquad \|A\|\le|A|,\qquad |A^{*}|=|A|,
\end{equation}
for all $N\times N$ matrices $X,Y$ and all $A\in\mathfrak{g}^{\mathbb{C}}$. These properties
hold for the operator norm and for the Hilbert--Schmidt norm
$|A|_{\mathrm{HS}}:=\bigl(\sum_{i,j}|A_{ij}|^2\bigr)^{1/2}$, and they imply that $|\cdot|$ is
$\operatorname{Ad}_G$-invariant. If $|\cdot|$ is taken
to be the Hilbert--Schmidt norm instead, the frame constant $C_{\mathrm{fr}}$ acquires a factor
$\sqrt{N}$, a number depending on $N$ only and therefore fixed before the coupling ceiling
$\gamma$ in the order of choice of the auxiliary parameters.
\item For $g\in G^{\mathbb{C}}$ the \emph{polar size} of $g$ is
\begin{equation*}
\operatorname{pol}(g):=\log\max\{\|g\|,\|g^{-1}\|\},
\end{equation*}
and for $\varepsilon>0$ we put
$\mathfrak{G}_{\varepsilon}:=\{g\in G^{\mathbb{C}}:\operatorname{pol}(g)<\varepsilon\}$.
\end{enumerate}
\end{definition}

\begin{proof}
We verify the assertions made above: the direct-sum decomposition of $\mathfrak{g}^{\mathbb{C}}$,
the properties of the pairing, the properties \eqref{4.22:eq-complex-group-a} for the two norms,
and the $\operatorname{Ad}_G$-invariance of $|\cdot|$.

\emph{The decomposition
$\mathfrak{g}^{\mathbb{C}}=\mathfrak{g}\oplus i\,\mathfrak{g}=\mathfrak{sl}(N,\mathbb{C})$.}
Let $A$ be traceless. Put $H_1:=(A+A^{*})/2$ and
$H_2:=(A-A^{*})/(2i)$. Both are Hermitian, and both are traceless because
$\operatorname{tr}A^{*}=\overline{\operatorname{tr}A}=0$; and $A=H_1+iH_2$. If $H_1+iH_2=0$
with $H_1,H_2$ Hermitian, taking adjoints gives $H_1-iH_2=0$, hence $H_1=H_2=0$. Thus the
traceless matrices are the direct sum $\mathfrak{g}\oplus i\,\mathfrak{g}$.

\emph{The pairing.} Bilinearity over $\mathbb{C}$ is that of the trace and of matrix
multiplication. For $g\in G^{\mathbb{C}}$, cyclicity of the trace gives
\begin{equation*}
\operatorname{tr}\bigl(gAg^{-1}\,gBg^{-1}\bigr)=\operatorname{tr}\bigl(gABg^{-1}\bigr)
=\operatorname{tr}AB .
\end{equation*}
For $A\in\mathfrak{g}$ we have $A=A^{*}$, so
$\langle A,A\rangle=\operatorname{tr}AA^{*}=N^{-1}\sum_{i,j}|A_{ij}|^2$, which is positive
unless $A=0$.

\emph{The properties \eqref{4.22:eq-complex-group-a} for the operator norm.} The first is
submultiplicativity of $\|\cdot\|$, the second is an equality, and the third holds because
$\|A^{*}\|=\|A\|$.

\emph{The properties \eqref{4.22:eq-complex-group-a} for the Hilbert--Schmidt norm.} Let
$e_1,\dots,e_N$ be the standard basis. Then
\begin{equation*}
|XA|_{\mathrm{HS}}^2=\sum_j\|XAe_j\|^2\le\|X\|^2\sum_j\|Ae_j\|^2=\|X\|^2|A|_{\mathrm{HS}}^2 .
\end{equation*}
The identity $|A^{*}|_{\mathrm{HS}}=|A|_{\mathrm{HS}}$ holds because both sides are the square
root of $\sum_{i,j}|A_{ij}|^2$. Hence, applying the preceding bound to $Y^{*}A^{*}$ and then
this identity,
\begin{equation*}
|AY|_{\mathrm{HS}}=|Y^{*}A^{*}|_{\mathrm{HS}}\le\|Y^{*}\|\,|A^{*}|_{\mathrm{HS}}
=\|Y\|\,|A|_{\mathrm{HS}},
\end{equation*}
and combining the two one-sided bounds gives the first property. For a vector $u$, writing
$Au=\sum_j u_jAe_j$, the triangle inequality and then the Cauchy--Schwarz inequality give
\begin{equation*}
\|Au\|\le\sum_j|u_j|\,\|Ae_j\|\le\|u\|\,|A|_{\mathrm{HS}},
\end{equation*}
which is the second property.

\emph{Invariance.} Let $g\in G$. Since $g^{*}g=\mathbb{1}$, $\|g\|=\|g^{-1}\|=1$. The first
property in \eqref{4.22:eq-complex-group-a} with $X=g$, $Y=g^{-1}$ gives
$|\operatorname{Ad}_gA|\le|A|$, and with $X=g^{-1}$, $Y=g$ applied to $\operatorname{Ad}_gA$ it
gives $|A|\le|\operatorname{Ad}_gA|$. Hence $|\operatorname{Ad}_gA|=|A|$.
\end{proof}

\begin{lemma}[Polar decomposition, identity principle and convexity in the complexified group]
\label{4.22:lem-polar-decomposition}
Let $\mathfrak{g}$, $\mathfrak{g}^{\mathbb{C}}=\mathfrak{sl}(N,\mathbb{C})$, $G=SU(N)$,
$G^{\mathbb{C}}=SL(N,\mathbb{C})$, $\operatorname{Ad}$, the invariant form
$\langle\cdot,\cdot\rangle$, the norms $\|\cdot\|$ and $|\cdot|$, the polar size
$\operatorname{pol}$ and the sets $\mathfrak{G}_{\varepsilon}$ be as in
Definition~\ref{4.22:not-complex-group}. In particular, $\mathfrak{g}$ is the space of
Hermitian traceless $N\times N$ matrices. Then the following hold.
\begin{enumerate}
\item[(g1)] Every $g\in G^{\mathbb{C}}$ can be written in exactly one way as $g=we^{Y}$
with $w\in G$ and $Y\in\mathfrak{g}$. It can also be written in exactly one way as
$g=e^{Y'}w$ with $w\in G$ and $Y'\in\mathfrak{g}$; here $w$ is the same and
$Y'=\operatorname{Ad}_wY$. Moreover $\operatorname{pol}(g)=0$ if and only if $g\in G$.
For $g=we^{Y}$ written in this way, for $v,v'\in G$, for $g_1,g_2\in G^{\mathbb{C}}$ and
for $Z\in\mathfrak{g}^{\mathbb{C}}$,
\begin{equation}\label{4.22:eq-polar-decomposition-a}
\begin{gathered}
g=we^{Y}=e^{\operatorname{Ad}_wY}w,\qquad \operatorname{pol}(g)=\|Y\|,\qquad
\operatorname{pol}(vgv')=\operatorname{pol}(g),\\
\operatorname{pol}(g_1g_2)\le\operatorname{pol}(g_1)+\operatorname{pol}(g_2),\qquad
\operatorname{pol}(e^{Z})\le\|Z\|\le|Z|.
\end{gathered}
\end{equation}
In particular, $|Z|\le\varepsilon$ implies $\operatorname{pol}(e^{Z})\le\varepsilon$.
\item[(g2)] For $A\in\mathfrak{g}^{\mathbb{C}}$ and $g=we^{Y}$ as in (g1),
\begin{equation}\label{4.22:eq-polar-decomposition-b}
|\operatorname{Ad}_gA|\le e^{2\operatorname{pol}(g)}|A|,\qquad
|\operatorname{Ad}_gA-\operatorname{Ad}_wA|\le\bigl(e^{2\operatorname{pol}(g)}-1\bigr)|A|.
\end{equation}
\item[(g3)] (Identity principle.) For $\varepsilon>0$ the set $\mathfrak{G}_{\varepsilon}$ is
open and connected, satisfies $v\mathfrak{G}_{\varepsilon}v'=\mathfrak{G}_{\varepsilon}$ for
$v,v'\in G$, and contains $G$. Let $\mathfrak{s}$ be a finite set, and let
$\varepsilon_s>0$ for $s\in\mathfrak{s}$. If $f$ is holomorphic on
$\prod_{s\in\mathfrak{s}}\mathfrak{G}_{\varepsilon_s}$, then
\begin{equation}\label{4.22:eq-polar-decomposition-c}
f\restriction G^{\mathfrak{s}}=0\quad\Longrightarrow\quad f\equiv0\ \text{ on }\
\prod_{s\in\mathfrak{s}}\mathfrak{G}_{\varepsilon_s}.
\end{equation}
\item[(g4)] (Real points of a complex orbit.) Let $Y\in\mathfrak{g}$. For $A\in\mathfrak{g}$
and for $V\in G$,
\begin{equation}\label{4.22:eq-polar-decomposition-d}
e^{Y}Ae^{-Y}\in\mathfrak{g}\ \Longrightarrow\ YA=AY,\qquad
e^{Y}Ve^{-Y}\in G\ \Longrightarrow\ YV=VY.
\end{equation}
\item[(g5)] (Equivariant maps.) Let $\varphi\colon\mathfrak{g}\to\mathfrak{g}$ satisfy
$\varphi(\operatorname{Ad}_wA)=\operatorname{Ad}_w\varphi(A)$ for all $w\in G$ and
$A\in\mathfrak{g}$. If $Y\in\mathfrak{g}$ and $A\in\mathfrak{g}$, then
\begin{equation}\label{4.22:eq-polar-decomposition-e}
YA=AY\quad\Longrightarrow\quad Y\varphi(A)=\varphi(A)Y.
\end{equation}
The gradient with respect to $\langle\cdot,\cdot\rangle$ of an
$\operatorname{Ad}_G$-invariant $C^1$ function on $\mathfrak{g}$ is such a map $\varphi$.
\item[(g6)] (Convexity.) Let $g\in G^{\mathbb{C}}$ and $Y\in\mathfrak{g}$. Then the map
$t\mapsto\operatorname{pol}(ge^{tY})$ is convex on $\mathbb{R}$, and
\begin{equation}\label{4.22:eq-polar-decomposition-f}
\operatorname{pol}\bigl(ge^{(t+is)Y}\bigr)=\operatorname{pol}\bigl(ge^{tY}\bigr)
\qquad(t,s\in\mathbb{R}).
\end{equation}
\end{enumerate}
\end{lemma}

\begin{proof}
\emph{(g1).} Let $g\in G^{\mathbb{C}}$.

\emph{Existence.} The matrix $g^{*}g$ is Hermitian and positive definite. Its determinant
is $|\det g|^2=1$. By the spectral theorem it has exactly one Hermitian logarithm; call it
$2Y$. The eigenvalues of $2Y$ are the logarithms of those of $g^{*}g$, so they sum to
$\log\det(g^{*}g)=0$. Hence $Y$ is traceless, $Y\in\mathfrak{g}$, and $e^{2Y}=g^{*}g$.
Put $w:=ge^{-Y}$. Then $w^{*}w=e^{-Y}g^{*}ge^{-Y}=1$ and $\det w=\det g\,\det e^{-Y}=1$.
Thus $w\in SU(N)=G$ and $g=we^{Y}$.

\emph{Uniqueness.} Suppose $g=w_1e^{Y_1}$ with $w_1\in G$ and $Y_1\in\mathfrak{g}$. Since
$w_1$ is unitary and $Y_1$ is Hermitian,
\begin{equation*}
g^{*}g=e^{Y_1}w_1^{*}w_1e^{Y_1}=e^{2Y_1},
\end{equation*}
so $2Y_1$ is a Hermitian logarithm of $g^{*}g$. Therefore $Y_1=Y$ and $w_1=ge^{-Y}=w$.

\emph{The second form.} We have $we^{Y}=(we^{Y}w^{-1})w=e^{\operatorname{Ad}_wY}w$, and
$\operatorname{Ad}_wY\in\mathfrak{g}$. Conversely, suppose $g=e^{Y'}w'$ with $w'\in G$ and
$Y'\in\mathfrak{g}$. Then $g=w'e^{\operatorname{Ad}_{w'^{-1}}Y'}$, and uniqueness of the
first form gives $w'=w$ and $Y'=\operatorname{Ad}_wY$.

\emph{The polar size.} Let $\lambda_{\max}(Y)$ and $\lambda_{\min}(Y)$ be the largest and
smallest eigenvalues of $Y$. Since $w$ is unitary,
\begin{equation*}
\|g\|=\|e^{Y}\|=e^{\lambda_{\max}(Y)},\qquad
\|g^{-1}\|=\|e^{-Y}w^{-1}\|=\|e^{-Y}\|=e^{-\lambda_{\min}(Y)}.
\end{equation*}
As $Y$ is Hermitian, $\|Y\|=\max\{\lambda_{\max}(Y),-\lambda_{\min}(Y)\}$, and so
$\operatorname{pol}(g)=\|Y\|$. Since $Y$ is traceless, $\lambda_{\min}(Y)\le0\le\lambda_{\max}(Y)$.
It follows that $\operatorname{pol}(g)=0$ if and only if $Y=0$, that is, if and only if
$g=w\in G$.

\emph{Invariance.} The operator norm is invariant under multiplication by unitary matrices.
Hence $\|vgv'\|=\|g\|$ and
\begin{equation*}
\|(vgv')^{-1}\|=\|v'^{-1}g^{-1}v^{-1}\|=\|g^{-1}\|,
\end{equation*}
which gives $\operatorname{pol}(vgv')=\operatorname{pol}(g)$.

\emph{Sub-additivity.} By sub-multiplicativity,
\begin{equation*}
\|g_1g_2\|\le\|g_1\|\,\|g_2\|\le e^{\operatorname{pol}(g_1)+\operatorname{pol}(g_2)}.
\end{equation*}
The same bound holds for $\|(g_1g_2)^{-1}\|=\|g_2^{-1}g_1^{-1}\|$. Taking the maximum and
the logarithm gives the sub-additivity.

\emph{The exponential.} Let $Z\in\mathfrak{g}^{\mathbb{C}}$. Since $Z$ is traceless,
$\det e^{Z}=1$, so $e^{Z}\in G^{\mathbb{C}}$. The exponential series gives
$\|e^{\pm Z}\|\le\sum_{n\ge0}\|Z\|^n/n!=e^{\|Z\|}$, hence $\operatorname{pol}(e^{Z})\le\|Z\|$.
The inequality $\|Z\|\le|Z|$ is part of \eqref{4.22:eq-complex-group-a}.

\emph{(g2).} By \eqref{4.22:eq-complex-group-a},
\begin{equation*}
|gAg^{-1}|\le\|g\|\,|A|\,\|g^{-1}\|\le e^{2\operatorname{pol}(g)}|A|,
\end{equation*}
which is the first bound. For the second, write
\begin{equation*}
\operatorname{Ad}_gA-\operatorname{Ad}_wA=\operatorname{Ad}_w\bigl(e^{Y}Ae^{-Y}-A\bigr)
=\operatorname{Ad}_w\int_0^1e^{tY}[Y,A]e^{-tY}\,dt .
\end{equation*}
The integrand is the $t$-derivative of $e^{tY}Ae^{-tY}$. By \eqref{4.22:eq-complex-group-a}:
\begin{itemize}
\item $\operatorname{Ad}_w$ does not increase $|\cdot|$, because $\|w\|=\|w^{-1}\|=1$;
\item $|[Y,A]|\le|YA|+|AY|\le2\|Y\|\,|A|$;
\item $|e^{tY}Be^{-tY}|\le\|e^{tY}\|\,\|e^{-tY}\|\,|B|\le e^{2t\|Y\|}|B|$ for every $B$.
\end{itemize}
Estimating the norm of the integral by the integral of the norm, and using
$\|Y\|=\operatorname{pol}(g)$ from (g1), we obtain
\begin{equation*}
|\operatorname{Ad}_gA-\operatorname{Ad}_wA|\le2\|Y\|\,|A|\int_0^1e^{2t\|Y\|}\,dt
=\bigl(e^{2\operatorname{pol}(g)}-1\bigr)|A|.
\end{equation*}

\emph{(g3).} The set $\mathfrak{G}_{\varepsilon}$ is open, because $\operatorname{pol}$ is
continuous on $G^{\mathbb{C}}$. It is invariant under $g\mapsto vgv'$ for $v,v'\in G$ by
(g1). It contains $G$, since $\operatorname{pol}=0<\varepsilon$ on $G$.

\emph{Connectedness.} Consider the map
\begin{equation*}
(w,Y)\mapsto we^{Y},\qquad G\times\{Y\in\mathfrak{g}:\|Y\|<\varepsilon\}\to\mathfrak{G}_{\varepsilon}.
\end{equation*}
It is continuous. By (g1) it is a bijection onto $\mathfrak{G}_{\varepsilon}$, because
$\operatorname{pol}(we^{Y})=\|Y\|$. Its inverse is
$g\mapsto(ge^{-Y(g)},Y(g))$ with $Y(g)=\tfrac12\log(g^{*}g)$, which is continuous since the
Hermitian logarithm is continuous on positive definite matrices. So the map is a
homeomorphism. Its domain is connected, because $SU(N)$ is connected and the ball is
convex. Hence $\mathfrak{G}_{\varepsilon}$ is connected, and so is the finite product
$\prod_{s}\mathfrak{G}_{\varepsilon_s}$.

\emph{Identity principle.} Let $n$ be the real dimension of $\mathfrak{g}^{\mathfrak{s}}$,
and fix $w_0\in G^{\mathfrak{s}}$. Consider the chart
\begin{equation*}
Z\mapsto\bigl(w_0(s)e^{iZ_s}\bigr)_{s\in\mathfrak{s}},\qquad
Z\in(\mathfrak{g}^{\mathbb{C}})^{\mathfrak{s}}\cong\mathbb{C}^n .
\end{equation*}
It is holomorphic and, near $Z=0$, a biholomorphism onto a neighbourhood of $w_0$. It maps
real points $Z\in\mathfrak{g}^{\mathfrak{s}}\cong\mathbb{R}^n$ into $G^{\mathfrak{s}}$,
since $iZ_s$ is then anti-Hermitian and traceless. Hence $f$ composed with the chart is
holomorphic on a polydisc about $0$ and vanishes on its real points. Its complex partial
derivatives at $0$ coincide with the real ones, which vanish. So its power series is zero,
and $f$ vanishes on a neighbourhood of $w_0$.

Now let $O$ be the set of points of $\prod_{s}\mathfrak{G}_{\varepsilon_s}$ near which $f$
vanishes identically. The set $O$ is open, and it is non-empty because it contains
$G^{\mathfrak{s}}$. It is also closed in the product. Indeed, let $g_0$ be a point of the
product that is a limit of points of $O$. The chart $Z\mapsto\bigl(g_0(s)e^{iZ_s}\bigr)_{s}$
about $g_0$ is holomorphic and, for the same reason as above, a biholomorphism near $Z=0$
onto a neighbourhood of $g_0$. In this chart all derivatives of $f$ at $Z=0$ vanish, being
limits of derivatives at points of $O$, where they vanish; so the power series of $f$ about
$g_0$ is zero, and $g_0\in O$. By connectedness $O$ is the whole product, and $f\equiv0$.

\emph{(g4).} First let $A\in\mathfrak{g}$, and put $B:=e^{Y}Ae^{-Y}$. Since $Y$ and $A$ are
Hermitian, $B^{*}=e^{-Y}Ae^{Y}$. The hypothesis $B\in\mathfrak{g}$ gives $B=B^{*}$, that is,
$e^{Y}Ae^{-Y}=e^{-Y}Ae^{Y}$. Multiplying on the left and on the right by $e^{Y}$ gives
$e^{2Y}A=Ae^{2Y}$.

Next let $V\in G$, and put $B:=e^{Y}Ve^{-Y}$. The hypothesis $B\in G$ gives $B^{*}B=1$, that is,
\begin{equation*}
e^{-Y}V^{*}e^{2Y}Ve^{-Y}=1 .
\end{equation*}
Hence $V^{*}e^{2Y}V=e^{2Y}$, and since $V$ is unitary, $e^{2Y}V=Ve^{2Y}$.

In both cases it remains to pass from $e^{2Y}$ to $Y$. The matrix $e^{2Y}$ is positive, with
eigenvalues $e^{2\lambda}$ for the eigenvalues $\lambda$ of $Y$, and distinct $\lambda$ give
distinct $e^{2\lambda}$. A Lagrange interpolation polynomial $p$ with $p(e^{2\lambda})=\lambda$
therefore satisfies $p(e^{2Y})=Y$. So $Y$ is a polynomial in $e^{2Y}$, and it commutes with
$A$, respectively $V$.

The hypothesis of (g4) is a genuine restriction. Let $N=2$, and let
$\sigma_1,\sigma_2,\sigma_3$ be the Pauli matrices
\begin{equation*}
\sigma_1=\begin{pmatrix}0&1\\1&0\end{pmatrix},\quad
\sigma_2=\begin{pmatrix}0&-i\\i&0\end{pmatrix},\quad
\sigma_3=\begin{pmatrix}1&0\\0&-1\end{pmatrix}.
\end{equation*}
Take $A=\sigma_3$ and $Y=t\sigma_1$ with real $t\ne0$. Since $\sigma_3\sigma_1=-\sigma_1\sigma_3$
and $\sigma_1\sigma_3=-i\sigma_2$,
\begin{equation*}
e^{Y}Ae^{-Y}=e^{2t\sigma_1}\sigma_3=\sigma_3\cosh2t-i\sigma_2\sinh2t .
\end{equation*}
This matrix is not Hermitian, so it does not lie in $\mathfrak{g}$.

\emph{(g5).} Since $Y\in\mathfrak{g}$ is Hermitian and traceless, $e^{itY}\in G$ for
$t\in\mathbb{R}$. Since $Y$ commutes with $A$, $\operatorname{Ad}_{e^{itY}}A=A$. Equivariance
gives
\begin{equation*}
\varphi(A)=\varphi\bigl(\operatorname{Ad}_{e^{itY}}A\bigr)=\operatorname{Ad}_{e^{itY}}\varphi(A)
\qquad\text{for all }t\in\mathbb{R}.
\end{equation*}
Differentiating at $t=0$ gives $0=i[Y,\varphi(A)]$.

For the gradient, let $\psi$ be an $\operatorname{Ad}_G$-invariant $C^1$ function on
$\mathfrak{g}$. Define $\nabla\psi(A)\in\mathfrak{g}$ by
$\langle\nabla\psi(A),B\rangle=d\psi(A)[B]$ for $B\in\mathfrak{g}$; this is possible because
$\langle\cdot,\cdot\rangle$ is positive, hence non-degenerate, on $\mathfrak{g}$.
Differentiating $\psi(\operatorname{Ad}_wA)=\psi(A)$ in $A$ in the direction $B$, and using
the invariance of $\langle\cdot,\cdot\rangle$, we get
\begin{equation*}
\langle\nabla\psi(\operatorname{Ad}_wA),\operatorname{Ad}_wB\rangle=\langle\nabla\psi(A),B\rangle
=\langle\operatorname{Ad}_w\nabla\psi(A),\operatorname{Ad}_wB\rangle .
\end{equation*}
As $B$ ranges over $\mathfrak{g}$, so does $\operatorname{Ad}_wB$. Non-degeneracy therefore gives
$\nabla\psi(\operatorname{Ad}_wA)=\operatorname{Ad}_w\nabla\psi(A)$.

\emph{(g6).} For unit vectors $u,v\in\mathbb{C}^N$ the function
$\tau\mapsto u^{*}ge^{\tau Y}v$ is entire, so $\log|u^{*}ge^{\tau Y}v|$ is subharmonic. Put
\begin{equation*}
\Psi(\tau):=\log\|ge^{\tau Y}\|=\sup_{u,v}\log|u^{*}ge^{\tau Y}v| .
\end{equation*}
The function $\Psi$ is continuous, since $ge^{\tau Y}$ is invertible and the norm is continuous. As a
continuous supremum of subharmonic functions, $\Psi$ is subharmonic. For real $s$ the matrix
$e^{isY}$ is unitary, so $\Psi(t+is)=\Psi(t)$. Thus $\Psi$ depends on $\operatorname{Re}\tau$
only. A continuous subharmonic function that depends on $\operatorname{Re}\tau$ only is convex
in $\operatorname{Re}\tau$: its distributional Laplacian is its second derivative in
$\operatorname{Re}\tau$, which is therefore non-negative.

The same argument applies to the inverse. Since $Y$ is Hermitian,
\begin{equation*}
\|(ge^{\tau Y})^{-1}\|=\|e^{-\tau Y}g^{-1}\|=\|(g^{-1})^{*}e^{-\bar\tau Y}\|,
\end{equation*}
and $(g^{-1})^{*}\in G^{\mathbb{C}}$. By the preceding paragraph, applied to $(g^{-1})^{*}$ in
place of $g$ and to the variable $-\bar\tau$, whose real part is $-\operatorname{Re}\tau$, the map
$t\mapsto\log\|(ge^{tY})^{-1}\|$ is convex.

Finally, $\operatorname{pol}(ge^{tY})$ is the maximum of the two convex functions
$t\mapsto\log\|ge^{tY}\|$ and $t\mapsto\log\|(ge^{tY})^{-1}\|$, and a maximum of convex
functions is convex. The identity \eqref{4.22:eq-polar-decomposition-f} follows from
$ge^{(t+is)Y}=ge^{tY}e^{isY}$, from $e^{isY}\in G$, and from the invariance part of (g1).
\end{proof}

\begin{definition}[Steps, their three sites and the scales of the flow]\label{4.22:not-steps-sites}
\emph{Steps.} Following the convention of \cite{B-CMP122a}, in which the $k$th density is
considered instead of the $(k+1)$st, \emph{step $k$} is indexed by the $k$th density
$\rho_k$. It has three sites.
\begin{itemize}
\item[$(\mathrm{T})$] The \emph{fluctuation integral}: the analysis of \cite[\S3]{B-CMP119},
taken at the index passage $k-1\to k$. It integrates the fluctuation field $A_{k-1}$ of that
analysis on the small-field region $\Lambda_k$. It creates the new rim slots of
level $k-1$, which are carried by the rim $\Omega_k\setminus\Lambda_k$, where
$\Omega_k\supset\Lambda_k$ is the region so denoted in the analysis of \cite[\S3]{B-CMP119}.
\item[$(\mathrm{P})$] The \emph{large-field stages}: the $N_k$ stages of
\cite[(1.25)--(1.63)]{B-CMP122a}, with $N_k$ fixed below.
\item[$(\mathrm{L})$] The \emph{presentation}: the construction of
\cite[(1.49)--(1.70)]{B-CMP122b}.
\end{itemize}
The sites $(\mathrm{P})$ and $(\mathrm{L})$, together with the operations
$\mathbf{T}'_k(X)$ and the exponentiation \cite[(1.98)]{B-CMP122b}, form the
\emph{$\mathbf{R}$-site} of step $k$.

\emph{The flow.} Write $x_k:=g_k^{-2}$. Let $c_5$ be the constant of
Definition~\ref{1.2:def-flow-constants}, and let $\lambda_\sharp$, $B_\sharp$ be the constants
of the coupling flow in clause~(0) of Theorem~0 (Theorem~\ref{4.1:thm-clause-zero}). We use
the two-sided bound on the running inverse couplings asserted in
Theorem~\ref{4.1:thm-clause-zero}:
\begin{equation}\label{4.22:eq-steps-sites-a}
\lambda_\sharp\,(j-i)-c_5\;\le\;x_i-x_j\;\le\;B_\sharp\,(j-i),
\qquad 0\le i\le j\le K .
\end{equation}

\emph{The scales.} Ba{\l}aban's sequence of scales $R_j$ of \cite{B-CMP122b} is taken to be
$R_j:=\check{R}_j$, the scales fixed by the choice of scales made earlier in this chapter:
with the exponent $r\ge 2$ of that choice, they are non-increasing in $j$ and satisfy
\begin{equation}\label{4.22:eq-steps-sites-b}
\bigl(\log(x_j-c_5)\bigr)^{r}\;\le\;\check{R}_j\;<\;L\,(\log x_j)^{r}.
\end{equation}

\emph{Number of stages.} The number of large-field stages of step $k$ is
\begin{equation*}
N_k:=\check{R}_k .
\end{equation*}
This is admissible: \cite{B-CMP122b} states that the choice $N=R_j$ satisfies the
conditions (i), (ii) imposed there on the number of stages.

One index letter is local to this section: we put
\begin{equation*}
h:=k-N_k .
\end{equation*}

Of \eqref{4.22:eq-steps-sites-b}, in the form it takes for $N_k$, the upper bound
\begin{equation*}
N_k\le L(\log x_k)^{r}
\end{equation*}
is the one used in general; the lower bound
\begin{equation*}
N_k\ge\bigl(\log(x_k-c_5)\bigr)^{r}
\end{equation*}
is used only where this is said explicitly.
Put
\begin{equation*}
R_{\min}:=\min_{j\le K}\check{R}_j ;
\end{equation*}
since $j\mapsto\check{R}_j$ is non-increasing, $R_{\min}=\check{R}_K$.

\emph{Large-field regions.} Put
\begin{equation*}
Z_j:=\Lambda_j^{c}
\end{equation*}
as in \cite[(2.3)]{B-CMP119}; the sets $Z_j$ increase with $j$. For a class $Z$ of
large-field regions that the large-field operation has already treated, the intersections
$Z_j\cap Z$ are again written $Z_j$, following \cite{B-CMP122a}.
\end{definition}

\begin{definition}[Pointwise coordinates, frames and their formal rotations]
\label{4.22:def-coordinates-frames}
Throughout, $\mathfrak{g}$ is the Lie algebra of Hermitian traceless $N\times N$ matrices,
$\mathfrak{g}^{\mathbb{C}}$ its complexification, and $G^{\mathbb{C}}$, the polar size
$\operatorname{pol}$ and the sets $\mathfrak{G}_{\varepsilon}$ are as in
Definition~\ref{4.22:not-complex-group}. For $g\in G^{\mathbb{C}}$ and
$X\in\mathfrak{g}^{\mathbb{C}}$ we write $\operatorname{Ad}_gX=gXg^{-1}$. The three sites of a
step are those of Definition~\ref{4.22:not-steps-sites}.

\emph{Pointwise coordinates.} A \emph{pointwise coordinate} is a triple of one of the two
following kinds.
\begin{enumerate}
\item $a=(A,j',b)$, standing for the rim fluctuation variable $A_{j'}(b)$, where $b$ is a bond
of the lattice of level $j'$ contained in $Z_{j'+1}\cap\Omega_{j'+1}$, as in \cite{B-CMP119}.
\item $a=(B,j',b)$, standing for the collar variable $B_{j'}(b)$ of a component arrested in the
sense of \cite[(1.55)]{B-CMP122a}, where
$b\in(\Omega^c_{j'+1}\setminus Z''_{j'+1})^{(j')*}$, as in \cite{B-CMP122a}.
\end{enumerate}
In both cases the variable is $\mathfrak{g}$-valued. We write $j'(a)$ for the level $j'$ of
$a$. The \emph{law site} of $a$ is
\begin{equation*}
z(a):=b_-.
\end{equation*}
Thus a fluctuation-type field is rotated at its own site of the block lattice, as in
\cite{B-CMP119}.

\emph{Frames.} A \emph{frame} is a component $W$ whose outcome at the $\mathbf{R}$-site of some
step (Definition~\ref{4.22:not-steps-sites}) is the second class of the large-field operation,
as in \cite{B-CMP122b}; we write $k_W$ for that step, and we take $W$ in the representation
\cite[(1.102)]{B-CMP122b}. Let $\Lambda$ denote the region of Ba{\l}aban's construction in
\cite{B-CMP122a} on whose intersection with $W$ the axial gauge of the frame is imposed. We put
$\Lambda_W:=\Lambda\cap W$, and we let
$y_0^W\in\partial^+\Lambda_W$ be the root of the axial gauge on $\Lambda_W$, as in
\cite{B-CMP122a}. The \emph{data} $\Phi_W$ of the frame are all the integration variables of
$\mathbf{T}'_{k_W}(W)$ in \cite[(1.71)]{B-CMP122b}, namely: the variables $B$ on
$\mathfrak{B}_0\cap W$; the variables $V_h$ on $(\Omega''^{\,\sim2}_{h+1})^c\cap W$; and the
inner variables of $\mathbf{T}_h(Z_h\cap W)$, that is, the $G$-valued $V_j$ and the
$\mathfrak{g}$-valued $A_j$ with $j<h$, together with the data of older frames. The
\emph{law sites} of $W$ form the set
\begin{equation*}
\Sigma_W:=\{y_0^W\}\cup\{b_-:\ b\ \text{an inward $\partial\Lambda_W$-crossing bond of }
\mathfrak{B}_0\}\subset\partial^+\Lambda_W,
\end{equation*}
whose elements are sites of the lattice of level $k_W$.

For a $G$-valued map $v$ on $\Sigma_W$ the \emph{rotated data} $\Phi_W^{[v]}$ are obtained from
$\Phi_W$ by the following substitutions:
\begin{equation*}
\begin{aligned}
X&\mapsto\operatorname{Ad}_{v(y_0^W)}X &&\text{on the interior $\mathfrak{g}$-valued members,}\\
V&\mapsto v(y_0^W)\,V\,v(y_0^W)^{-1} &&\text{on the $G$-valued members,}\\
B(b)&\mapsto\operatorname{Ad}_{v(b_-)}B(b) &&\text{on the crossing bonds.}
\end{aligned}
\end{equation*}
So the rotation is constant inside the region $\Lambda_W$ and acts by conjugation with the
value at $b_-$ on the inward crossing bonds of $\mathfrak{B}_0$. This is the rotation induced at
the root $y_0^W$ on variables declared in a frame rooted at $y_0^W$, as in \cite{B-CMP119}.
These are exactly the substitutions by which the frame data enter the presentation's rotated
unit $F^{\mathfrak{b}}_{\mathfrak{E}}$.

\emph{Terms with a domain.} Let a term with domain $S$ be carried in the expansion. We write
$\mathbf{A}$ for its pointwise coordinates, $\mathfrak{F}(S)$ for the set of its frames, and
$\boldsymbol{\Phi}:=(\Phi_W)_{W\in\mathfrak{F}(S)}$. We put
\begin{equation*}
s(a):=(\text{kind of }a,\ j'(a),\ z(a)),
\end{equation*}
where the kind is $A$ or $B$. The law sites of the term are the sets
\begin{align*}
\mathfrak{s}^{\mathrm{pt}}(S)&:=\{s(a):\ a\ \text{a pointwise coordinate of the term}\},\\
\mathfrak{s}^{\mathrm{fr}}(S)&:=\textstyle\bigsqcup_{W\in\mathfrak{F}(S)}\Sigma_W,\\
\mathfrak{s}(S)&:=\mathfrak{s}^{\mathrm{pt}}(S)\sqcup\mathfrak{s}^{\mathrm{fr}}(S).
\end{align*}

\emph{Rotation of pointwise coordinates.} For $g:\mathfrak{s}(S)\to G^{\mathbb{C}}$ we put
\begin{equation}\label{4.22:eq-coordinates-frames-1}
\mathbf{A}^{[g]}(a):=\operatorname{Ad}_{g(s(a))}\mathbf{A}(a).
\end{equation}
For two such maps $g_1,g_2$ the product $g_1g_2$ is taken pointwise,
$(g_1g_2)(s):=g_1(s)g_2(s)$.

\emph{Rotation of frames.} For $G$-valued $g$ we put
$\boldsymbol{\Phi}^{[g]}:=(\Phi_W^{[g\restriction\Sigma_W]})_{W\in\mathfrak{F}(S)}$. For
$G^{\mathbb{C}}$-valued $g$ this formula does not define $\boldsymbol{\Phi}^{[g]}$; at complex
$g$ the symbol is a name, whose meaning is that of the formally rotated frame data
$\boldsymbol{\Phi}^{\langle g\rangle}$ introduced in the definition of formally rotated frame
data at complex rotations.

\emph{Site maps.} To every law site $s$ there belongs a site of the lattice: $z$ if
$s=(\text{kind},j',z)\in\mathfrak{s}^{\mathrm{pt}}(S)$, and $s$ itself if
$s\in\mathfrak{s}^{\mathrm{fr}}(S)$. A map $u$ defined on sites acts on maps
$g:\mathfrak{s}(S)\to G^{\mathbb{C}}$ by $(ug)(s):=u(\text{site of }s)\,g(s)$. The rotations used
in the presentation are of the form $g(s):=u(\text{site of }s)$ for a map $u$ defined on sites.
Indexing law sites by kind and
level, and not by site alone, allows a stage to treat the level it integrates separately.

\emph{Rotation sets.} For $\varepsilon:\mathfrak{s}(S)\to\,]0,\infty[$ we put
\begin{equation*}
\mathfrak{G}_{\varepsilon}(S):=\prod_{s\in\mathfrak{s}(S)}\mathfrak{G}_{\varepsilon(s)}.
\end{equation*}

\emph{Background pairs.} Write $P:=(\mathbb{U},\mathbb{J})$. For a gauge transformation $u$
put $P^u:=(\mathbb{U}^u,\operatorname{Ad}_u\mathbb{J})$, as in \cite{B-CMP109}.

Finally, for all $g_1,g_2:\mathfrak{s}(S)\to G^{\mathbb{C}}$,
\begin{equation}\label{4.22:eq-coordinates-frames-2}
\bigl(\mathbf{A}^{[g_2]}\bigr)^{[g_1]}=\mathbf{A}^{[g_1g_2]}.
\end{equation}
\end{definition}

\begin{proof}
We prove \eqref{4.22:eq-coordinates-frames-2}. Fix a pointwise coordinate $a$ and write
$s=s(a)$. Apply \eqref{4.22:eq-coordinates-frames-1} first to $g_1$ and the coordinates
$\mathbf{A}^{[g_2]}$, then to $g_2$ and $\mathbf{A}$. This gives
\begin{equation*}
\bigl(\mathbf{A}^{[g_2]}\bigr)^{[g_1]}(a)
=\operatorname{Ad}_{g_1(s)}\operatorname{Ad}_{g_2(s)}\mathbf{A}(a)
=g_1(s)g_2(s)\,\mathbf{A}(a)\,g_2(s)^{-1}g_1(s)^{-1}.
\end{equation*}
Since $g_2(s)^{-1}g_1(s)^{-1}=(g_1(s)g_2(s))^{-1}$ and $g_1(s)g_2(s)=(g_1g_2)(s)$, the
right-hand side equals $\operatorname{Ad}_{(g_1g_2)(s)}\mathbf{A}(a)$. By
\eqref{4.22:eq-coordinates-frames-1} applied to $g_1g_2$, this is $\mathbf{A}^{[g_1g_2]}(a)$.
\end{proof}

\begin{definition}[Constants inherited from the correlation theorem]\label{4.22:not-inherited-constants}
The letters of the first four groups below keep the meaning they have in the correlation theorem
and in the chain estimates stated with it; those of the last group keep the meaning they have in
the presentation of stored terms. They are used here unchanged. The correlation theorem treats the step from level $k$ to level $k+1$; its letters are nevertheless written here at the level at which each is defined. The three sites $(\mathrm{T})$, $(\mathrm{P})$, $(\mathrm{L})$ of a step are those of Definition~\ref{4.22:not-steps-sites}. Only the letters fixed by a formula are displayed; the others are named and used as in their home statements.

\emph{The threshold and the new slot.} The small-field threshold at level $k$ is
\begin{equation}\label{4.22:eq-inherited-constants-1}
\mathfrak{b}=\mathfrak{b}_k=A_1\bigl(\log g_k^{-2}\bigr)^{p_0}.
\end{equation}
The set $\Lambda$ is the output domain of the flat form of the fluctuation-integral statement at
the site $(\mathrm{T})$. The letters $\bar{\Lambda}$ and $\mathfrak{r}$ are those of the
correlation theorem. The new slot is $x:=A\restriction\mathfrak{r}$, the restriction to
$\mathfrak{r}$ of the field argument $A$ of the correlation theorem. The letter $x$ without
subscript denotes this slot here and the argument of the local layer norm in
\eqref{4.22:eq-inherited-constants-8}; with a subscript, $x_k=g_k^{-2}$ as in the glossary.
Further letters of the correlation theorem used here are
\begin{equation*}
T(\sigma),\ T_0,\ \mathsf{m},\ \Theta_1,\ \Theta_{\flat},\ \tau_{\flat},\ \varpi,\ \varpi_{\gamma},\
\mathcal{S}(P),\ \mathsf{N}(b),\ N_A,\ N_Z,\ \gamma_2,\ c_0'.
\end{equation*}
Here $\mathsf{m}$ is the mean field of the correlation theorem.

\emph{Weights of the strong sets.} The constants $c_{\eta}$ and $\mathsf{p}$ are
\begin{equation}\label{4.22:eq-inherited-constants-2}
c_{\eta}=\Bigl[\frac{\gamma_2}{16\,(c_0'+1)}\Bigr]^{1/2}\le\frac{1}{16},
\qquad
\mathsf{p}=\frac{\gamma_2\,\mathfrak{b}^2}{64}.
\end{equation}
With $\mathsf{p}(g)$ the value of $\mathsf{p}$ at the coupling $g$, we set
\begin{equation}\label{4.22:eq-inherited-constants-3}
\kappa_A^{\mathrm{pr}}(g)=\frac{\mathsf{p}(g)}{N_A(g)},
\qquad
\kappa_A(k):=\min_{i\le k}\kappa_A^{\mathrm{pr}}(g_i).
\end{equation}
The sequence $k\mapsto\kappa_A(k)$ is non-increasing. Indeed, $\kappa_A(k+1)$ is a minimum over the
indices $i\le k+1$. This set contains the indices $i\le k$ over which $\kappa_A(k)$ is the minimum;
hence $\kappa_A(k+1)\le\kappa_A(k)$.

\emph{Radii read from the coupling.} The radii at level $k$ are
\begin{align}
\alpha_{*,k}&:=\min\{\alpha_{0,k},\alpha_{1,k}\}=g_k\,C_*\,(\log x_k)^{q_*},
\label{4.22:eq-inherited-constants-4}\\
\alpha^+_k&:=\max\{\tilde\alpha_{0,k},\tilde\alpha_{1,k},\delta'_k\},
\label{4.22:eq-inherited-constants-5}
\end{align}
where $C_*$ and $q_*$ are the constants of the correlation theorem's radii.
With $\alpha_*$ read at the coupling $g$, we set
\begin{equation}\label{4.22:eq-inherited-constants-6}
\varrho=\frac{c_A\,\alpha_*}{g},
\qquad
r_A=\frac{\varrho}{\mathfrak{b}},
\qquad
T^{\mathbb{C}}=\mathfrak{b}+(\Theta_{\flat}+1)\,\varrho ,
\end{equation}
where $c_A$ is the constant of the correlation theorem in the radius $\varrho$. In this section
$\varrho$ denotes this radius.
The correlation theorem also supplies the following objects:

\begin{itemize}
\item the letters $\mathfrak{b}^{\sharp}$ and $C_{\mathrm{far}}$, and the real datum set $\mathfrak{X}^{\sharp}$;
\item the flat class $A^{\flat}$ with its parameter domain $\mathfrak{D}_{\mathcal{S}}$ and its sup norm $\mathfrak{N}$;
\item its floor conditions, with the constant $\gamma_V$;
\item the constant $K_R$ and its homogeneity condition.
\end{itemize}

\noindent All of these are used here as stated there.

\emph{Chain domains.} The chain estimates stated with the correlation theorem fix the following letters:

\begin{itemize}
\item the domains $\mathfrak{D}^{\flat}_n(X)$ and the norms $\|\cdot\|^{A,\flat}_{(n)}$;
\item the constants $C_0^{\sharp}$ and $a_m$;
\item the background map $\Phi_{t,\sigma}(b)$;
\item the constant $\varkappa_{\mathsf{w}}$, the radius $R^{\natural}_{\square}$ and the weighted
contour distance $d^{\wedge}$.
\end{itemize}

\noindent In these estimates $\mathsf{w}=M/M_1$, and $\delta$ denotes the decay rate fixed there.
The kernel condition fixed there is used at the parameter values $(\mathsf{w}_0,\tfrac12)$, and
the two further conditions stated there beside it are used as stated. Finally, the constant
$\beta$ of these estimates (not to be confused with the coupling quantity $\beta_{k+1}$ of the
glossary) and the layer number $\theta_S$ satisfy
\begin{equation}\label{4.22:eq-inherited-constants-7}
\beta=\frac15,
\qquad
\theta_S\in\Bigl[\frac12,\frac89\Bigr].
\end{equation}

\emph{Letters of the presentation.} The presentation of stored terms fixes the following objects:

\begin{itemize}
\item the sourced set $\mathfrak{S}_{\mathrm{src}}$;
\item a member $\mathfrak{b}$ of the chain $\mathfrak{ch}=(\mathfrak{b}_0,\dots,\mathfrak{b}_4)$;
here $\mathfrak{b}$ denotes a member of the chain and not the threshold of
\eqref{4.22:eq-inherited-constants-1}, and in what follows it occurs only as a superscript;
\item the pairs $(\mathcal{U}_q^{\mathfrak{b}},\mathcal{J}_q^{\mathfrak{b}})$ and the map $h^{\mathfrak{b}}=h_q^{\mathfrak{b}}\,\mathfrak{Q}^{\mathfrak{b}}$;
\item the maps $\mathfrak{u}_s$, $\mathfrak{w}_s$ and the layers $\mathbb{H}_s$;
\item the constants $K_u$, $\varepsilon_{\mathfrak{Q}}$ and the product $K''a_1$;
\item the set $D_2$, the strip, and the sets $\Lambda^{\natural}$ and $\Lambda^{\sim}$.
\end{itemize}

\noindent The local layer norm, in the scaling of the tube spaces of the presentation, is
\begin{equation}\label{4.22:eq-inherited-constants-8}
\Theta_0^{\mathrm{loc}}(\mathbb{H})(x):=L^{m(x)}\,\eta\cdot\sup_{B^{\mathrm{top}}(x)}|\mathbb{H}|.
\end{equation}
Here the exponent $m(x)$, the factor $\eta$ and the set $B^{\mathrm{top}}(x)$ over which the
supremum runs are those fixed in the presentation for the argument $x$ of the norm.
The rotation size used in this section is
\begin{equation}\label{4.22:eq-inherited-constants-9}
\varepsilon_{\mathrm{live}}(k):=K''a_1+\varepsilon_{\mathfrak{Q}} .
\end{equation}
By the lemma of the presentation of stored terms in which $\varepsilon_{\mathrm{live}}(k)$ is
introduced (its third assertion), it is of size $g_k$ times a power of $\log g_k^{-2}$.
\end{definition}

\begin{definition}[The rotation radius, the frame unit and the margin]
\label{4.22:def-rotation-radius}
The steps of the flow, their three sites $(\mathrm{T})$, $(\mathrm{P})$, $(\mathrm{L})$ and the
law sites are those of Definition~\ref{4.22:not-steps-sites}. For a stored term with domain $S$,
$\mathfrak{s}^{\mathrm{pt}}$ denotes its pointwise law sites and $\Sigma_W$ the law sites of a
frame $W$. The following numbers and functions are fixed.

\begin{enumerate}
\item[(a)] \emph{Numbers chosen after $M$ and before $A_0$, $A_1$, $C_{\mathrm{far}}$, $c_A$ and
$\gamma$.} The \emph{rotation radius} at pointwise sites is
\begin{equation}\label{4.22:eq-rotation-radius-a}
\begin{gathered}
\varepsilon_{\star}
:=\frac12\log\Bigl(1+\frac{c_{\eta}}{5(\Theta_{\flat}+1)}\Bigr),
\\
\text{equivalently}\qquad
\bigl(e^{2\varepsilon_{\star}}-1\bigr)(\Theta_{\flat}+1)=\frac{c_{\eta}}{5}.
\end{gathered}
\end{equation}
Here $c_{\eta}$ and $\Theta_{\flat}$ are the constants of
Definition~\ref{4.22:not-inherited-constants}, and $e^{2\varepsilon_{\star}}\le 1+1/160$. Further,
\begin{equation}\label{4.22:eq-rotation-radius-1}
\varpi_{\mathrm{fr}}:=\tfrac18(1-\theta_S)\beta,
\qquad
C_{\mathrm{fr}}:=3^d(C_{\chi}+4),
\end{equation}
where $\beta$ is the schedule constant of the layered analyticity spaces
(Definition~\ref{4.4:def-post-step-space}), $\theta_S$ is the layer number and $C_{\chi}$ is the
bump constant. The rate $\delta_{\natural}$ and the constant
$C_{\natural}$ of the far debits are numbers depending only on $d$, $L$ and $M$.

\item[(b)] \emph{Functions of the coupling.} The \emph{margin} and the \emph{frame radius} are
\begin{equation}\label{4.22:eq-rotation-radius-2}
\mu(x):=\varepsilon_{\star}(x-c_5)^{-2},\qquad \mu_k:=\mu(x_k),\qquad
\varepsilon^{\mathrm{fr}}(W):=\frac{\varpi_{\mathrm{fr}}\,\alpha_{*,k_W}}{4C_{\mathrm{fr}}}.
\end{equation}

\item[(c)] \emph{The unit of a law site} (the frame unit at frame sites) is
\begin{equation}\label{4.22:eq-rotation-radius-3}
\mathsf{e}(s):=
\begin{cases}
\varepsilon_{\star}, & s\in\mathfrak{s}^{\mathrm{pt}},\\[2pt]
\varepsilon^{\mathrm{fr}}(W), & s\in\Sigma_W.
\end{cases}
\end{equation}

\item[(d)] \emph{Numbers chosen after $\kappa_1$ and before $\gamma$.}
\begin{itemize}
\item[--] $C_b\ge 1$ is the constant of the composition, with itself, of the residual bound at
the site $(\mathrm{T})$.
\item[--] $C_A^{(1/2)}$ is the resummation constant, at half the decay rate, of the resummation
lemma in the proof of the correlation theorem.
\item[--] Let $B_H$ be the constant of the family bound, and let $K_R$ be the further number
entering the combined constant below. Let $B_{\mathbb{H}}$ be the constant of clause (b) of the
bound on the Landau layer $\mathbb{H}$, and $K_{\mathbb{H}}$ the constant of clause (a) of the
lemma on orbit data.
The family bound is used in what follows with the combined constant
\begin{equation}\label{4.22:eq-rotation-radius-4}
B_H^S:=\max\bigl\{B_H,\ B_{\mathbb{H}}C_1/K_R,\ B_{\mathbb{H}}(1+K_{\mathbb{H}})\bigr\},
\end{equation}
where $C_1$ is the constant of \cite[(1.19)]{B-CMP116}.
\end{itemize}

\item[(e)] \emph{The reserve} at a law site $s$ and level $k$ is
\begin{equation}\label{4.22:eq-rotation-radius-5}
\eta^{\mathrm{inc}}_k(s):=\frac{\mu_k\,\mathsf{e}(s)}{\varepsilon_{\star}}.
\end{equation}

\item[(f)] The scaffold assigns to every stored term $F$ a length $\mathsf{d}_{\ell}(F)$.

\item[(g)] \emph{Disc and decay letters.} For a piece $p$, let $\mathsf{w}_{\mathrm{site}}$ be
the disc width of a site and $\eta_p$ the debit attached to $p$. The disc radius is
\begin{equation*}
\rho_p:=\mathsf{w}_{\mathrm{site}}\,(C_b\eta_p)^{-1/2}.
\end{equation*}
Let $r^0$ be the base radius and $\delta_0$ the base decay rate of the three-field bound entering
hypothesis (H6) (Definition~\ref{2.2:not-hypotheses}).
The decay constants are $C_{3F}=2/r^0$ and $\delta_P=\tfrac18\delta_0$.
For a core $\mathsf{C}_C$ and a region $X$, the depth and the disc radius attached to $X$
are
\begin{equation}\label{4.22:eq-rotation-radius-6}
D_X:=d^{\wedge}\bigl(X,\mathsf{C}_C^{c}\bigr),\qquad
\rho_X:=r^0e^{\frac12\delta_P D_X}.
\end{equation}

\item[(h)] \emph{Letters of hypothesis (H6)} (Definition~\ref{2.2:not-hypotheses}).
\begin{itemize}
\item[--] $\Sigma^{\mathfrak{q}}(x)$ is the function of the coupling bounding the kept quadratic
form. It is non-increasing in $x$ and tends to $0$ as $x\to\infty$.
\item[--] $\mathfrak{q}_{\max}(j)$ is the size of the kept quadratic form at step $j$.
\item[--] $\Sigma^V(j)$ is the two-point bound of the kept member $V$.
\item[--] The total is
\begin{equation}\label{4.22:eq-rotation-radius-7}
\begin{gathered}
\Sigma^{\mathrm{tot}}(j):=\Sigma^V(j)+2\,\mathfrak{q}_{\max}(j)\le C_{\Sigma}\,\check{R}_j^{p_{\Sigma}},
\\
p_{\Sigma}:=2d+6+r_{\mathrm{pr}}+\bigl\lceil 2p_1/r_{\mathrm{pr}}\bigr\rceil .
\end{gathered}
\end{equation}
\end{itemize}
The exponent $p_{\Sigma}$ is formed from $d$ and from the numbers $r_{\mathrm{pr}}$ and $p_1$,
which are fixed before it.
Here $C_{\Sigma}$ is a number chosen after $\kappa_1$ and $M$, and before
$\mathfrak{n}^{\mathrm{kept}}$ and $\gamma$.

\item[(i)] \emph{Numbers chosen after $C_{\Sigma}$ and before $\gamma$.} For a rate variable
$\beta>0$ (distinct from the schedule constant $\beta$ of (a), and evaluated at $\beta=1$ below) set
\begin{equation}\label{4.22:eq-rotation-radius-8}
\begin{gathered}
\mathfrak{s}_{\Sigma}(\beta):=\sup_{\varrho\ge1}(2\varrho+2)^{p_{\Sigma}}e^{-\beta\varrho},
\\
\mathfrak{n}^{\mathrm{kept}}(\beta):=2^dC_{\Sigma}\,\mathfrak{s}_{\Sigma}(\beta),
\\
\mathfrak{n}^{\mathrm{kept}}:=\mathfrak{n}^{\mathrm{kept}}(1).
\end{gathered}
\end{equation}
$\mathfrak{O}(v)$ denotes the oscillation seminorm of a kept member $v$ over the tube and the
chart.
\end{enumerate}
\end{definition}

\begin{proof}
\emph{The two forms of \eqref{4.22:eq-rotation-radius-a}.} Doubling the definition of
$\varepsilon_{\star}$ and exponentiating gives
$e^{2\varepsilon_{\star}}=1+c_{\eta}/(5(\Theta_{\flat}+1))$. Subtracting $1$ and multiplying by
$\Theta_{\flat}+1$ gives the second form. Conversely, the second form, divided by
$\Theta_{\flat}+1$, yields
\begin{equation*}
e^{2\varepsilon_{\star}}=1+\frac{c_{\eta}}{5(\Theta_{\flat}+1)}.
\end{equation*}
Taking $\tfrac12\log$ of this recovers the first form. The same identity shows that the bound
$e^{2\varepsilon_{\star}}\le 1+1/160$ is the inequality
\begin{equation*}
\frac{c_{\eta}}{5(\Theta_{\flat}+1)}\le\frac1{160},
\qquad\text{that is}\qquad
c_{\eta}\le\frac{\Theta_{\flat}+1}{32}.
\end{equation*}
So the bound $e^{2\varepsilon_{\star}}\le 1+1/160$ stated in (a) holds exactly when $c_{\eta}$ and
$\Theta_{\flat}$ satisfy this inequality.

\emph{The reserve.} Inserting \eqref{4.22:eq-rotation-radius-2} into
\eqref{4.22:eq-rotation-radius-5} gives
\begin{equation*}
\eta^{\mathrm{inc}}_k(s)=\mathsf{e}(s)\,(x_k-c_5)^{-2}.
\end{equation*}
By \eqref{4.22:eq-rotation-radius-3}, this equals $\mu_k$ when $s\in\mathfrak{s}^{\mathrm{pt}}$,
and $\mu_k\,\varepsilon^{\mathrm{fr}}(W)/\varepsilon_{\star}$ when $s\in\Sigma_W$.

\emph{Finiteness of $\mathfrak{s}_{\Sigma}(\beta)$.} Let $\beta>0$ and let $n$ be an integer with
$n\ge p_{\Sigma}+1$. For $\varrho\ge1$ we have $2\varrho+2\le 4\varrho$ and
$e^{\beta\varrho}\ge(\beta\varrho)^n/n!$. Hence
\begin{equation*}
(2\varrho+2)^{p_{\Sigma}}e^{-\beta\varrho}
\le\frac{4^{p_{\Sigma}}\,n!}{\beta^{n}}\,\varrho^{\,p_{\Sigma}-n}
\le\frac{4^{p_{\Sigma}}\,n!}{\beta^{n}},
\end{equation*}
since $p_{\Sigma}-n\le-1<0$ and $\varrho\ge1$. So the supremum in
\eqref{4.22:eq-rotation-radius-8} is finite. It follows that $\mathfrak{n}^{\mathrm{kept}}(\beta)$ and
$\mathfrak{n}^{\mathrm{kept}}$ are finite numbers.
\end{proof}

\begin{remark}[Why a margin per step and an age factor at frames]\label{4.22:rem-margin-age-factor}
The radii of Definition~\ref{4.22:def-rotation-radius} are diminished at every step by the
debits at the three sites $(\mathrm{T})$, $(\mathrm{P})$, $(\mathrm{L})$. Three observations show
that these debits alone do not close, and each is met by one further device:
\begin{equation}\label{4.22:eq-margin-age-factor-a}
\begin{gathered}
\mu_k=\mu(x_k)=\varepsilon_\star\,(x_k-c_5)^{-2}\ \text{per step},\qquad
\eta_p\ \longmapsto\ \eta_p^{1/2},\\
e^{-2(k-k_W)}\ \text{against}\quad
\frac{\alpha_{*,k}}{\alpha_{*,k_W}}\le 2\Bigl(1+\frac{B_\sharp\,(k-k_W)}{x_k}\Bigr)^{1/2}.
\end{gathered}
\end{equation}
The first is the margin $\mu_k$ charged at each step for the bracket increment; the second is the
price of reading the response on its Cauchy disc with site-wise decay; the third is the age factor
carried by every debit charged to a frame $W$ created at level $k_W$, at age $k-k_W$, until the
frame is closed.
\end{remark}

\begin{proof}
\emph{The margin.} Every bound at a site that reads a presented unit $\mathfrak{E}$ only through
$\sup|\mathfrak{E}|$ holds with the supremum taken over the whole chart. This does not cover the
bracket increment. The complex bracket bound is a Cauchy estimate in the rotation, an element of
the chart, at the identity rotation, with room $\tfrac12\varepsilon_\star$. The output of the step is
stored on the chart $\mathfrak{G}_{\varepsilon_k-\eta_k}$. On its boundary the room left for the
Cauchy estimate is $0$, so the smallness of the new terms is not uniform on the chart on which they
are stored. One further debit cures this: the margin $\mu_k$ of
\eqref{4.22:eq-margin-age-factor-a} is subtracted at each step. Then
\begin{equation*}
\sum_k\mu_k=\varepsilon_\star\sum_k(x_k-c_5)^{-2}
\end{equation*}
is finite by the summability of $(x_k-c_5)^{-2}$ along the flow, \eqref{4.22:eq-steps-sites-a},
the same bound that sums the other debits. With this margin, clause~(b) of the complex bracket
bound holds on the whole chart, at the loss factor
\begin{equation*}
\frac{4.2\,K_u}{\mu_k}=\frac{4.2\,K_u\,(x_k-c_5)^2}{\varepsilon_\star}.
\end{equation*}
Since $x_k=g_k^{-2}$, this factor is bounded by a constant times a fixed power of $g_k^{-2}$.
Every live law site is $R$-far, so the terms read there are smaller than every power of $g_k$.
After multiplication by the loss factor they therefore remain below every power of $g_k$.

\emph{Dilation on the Cauchy disc.} The response is small at $R$-far law sites only in its
undilated form. In the correlation theorem the response is read on a Cauchy disc
$|\mathsf{t}|\le\mathsf{w}_{\mathrm{site}}/\eta_p$, where $\eta_p$ is the parameter that, together
with the disc width $\mathsf{w}_{\mathrm{site}}$, fixes the radius of the disc at the piece $p$;
this disc dilates the response by exactly its decay. On the disc, therefore, the rotation of the
response at a law site is comparable to
$K_u\alpha_{*}$ and is not small in $R$. If it were debited at that size, the sum
$\sum_k\alpha_{*,k}$ would diverge with $K$, and a cap in the sense of
Definition~\ref{2.3:def-cap} would be needed. The cure is the disc-rotation lemma, applied with the
site-wise rotation bound $\vartheta_p(s)$ at the fluctuation-integral site. The dilation of the
disc is calibrated at the near end of the set $S_v$ of the site-wise bound,
\cite[(3.15)]{B-CMP109}. At depth the rotation retains the factor
$e^{-\frac12\delta_1\mathfrak{d}}$, where $\delta_1$ is the decay constant of the site-wise rotation
bound $\vartheta_p(s)$ and $\mathfrak{d}$ is the depth of the law site below that near end.
Hence there is no loss factor. The only price is the replacement $\eta_p\mapsto\eta_p^{1/2}$
of \eqref{4.22:eq-margin-age-factor-a}, as stated in the disc-rotation lemma. At the site
$(\mathrm{P})$ nothing needs to be cured: the inheritance theorem, applied with its parameter
$\theta=\tfrac12$, retains the factor $e^{-\frac14\delta d^{\wedge}}$, where $\delta$ is the decay
constant of the inheritance theorem.

\emph{Frames.} The twist of a frame $W$ is a relative orientation of its data $\Phi_W$ against the
background. By the definition of the domains, its radius
$\varepsilon^{\mathrm{fr}}(W)=\varpi_{\mathrm{fr}}\alpha_{*,k_W}/(4C_{\mathrm{fr}})$ is proportional to
$\alpha_{*,k_W}$, that is, to the coupling at the creation of $W$. The debits at later levels
$k>k_W$, however, read the later and larger couplings through $\alpha_{*,k}$. The sum of debits as
first set up carries no decay in the age in the class of terms touched by a second-class outcome,
so it does not close against this credit. Two facts cure this. First, by clause~(v) of the frame
lemma, the account of a frame
is closed at the next completed processing of its component. Second, until then every debit
charged to $W$ carries the factor $e^{-2(k-k_W)}$ from the shells crossed, by clause~(ii) of the
shell lemma. This factor beats the growth of $\alpha_{*,k}/\alpha_{*,k_W}$ bounded in
\eqref{4.22:eq-margin-age-factor-a}.
\end{proof}

\begin{proposition}[The variables the presentation rotates]\label{4.22:prop-rotated-variables}
Let $\mathcal{C}$ be the class of presented terms, consisting of:
\begin{itemize}
\item the old $\mathbf{E}^{(j)}$-, $\mathbf{R}^{(j)}$- and $\mathbf{B}^{(j)}$-terms with $j\le k$ whose domain
$X$ meets a large-field region $Z$;
\item the terms of the structure $\mathbb{B}''$;
\item the $Y_0$-terms, that is, the terms indexed by $Y_0$ in \cite[\S1]{B-CMP122b}, and the
terms of \cite[(1.59)]{B-CMP122a};
\item the frozen terms of arrested components.
\end{itemize}
For a pair $(\mathbf{U},\mathbf{J})$ and a $G$-valued map $u$ we write
$(\mathbf{U}^u,\operatorname{Ad}_u\mathbf{J})$ for its gauge transform.
For $\mathfrak{E}\in\mathcal{C}$ with domain $X$, let $\zeta$ be the source variable, let $h$ be the
$G$-valued gauge map of the presentation and $(\hat{\mathcal{U}}^{\mathfrak{b}},\hat{\mathcal{J}}^{\mathfrak{b}})$
its background pair, and let $\Phi$ denote the integration variables of the current step
(kind~(d) below). The rotated presentation
of $\mathfrak{E}$ is
\begin{equation}\label{4.22:eq-rotated-variables-1}
F^{\mathfrak{b}}_{\mathfrak{E}}(\zeta;\Phi)
:=\mathfrak{E}\bigl(X;(\hat{\mathcal{U}}^{\mathfrak{b}})^{h},
\operatorname{Ad}_h\hat{\mathcal{J}}^{\mathfrak{b}};\mathbf{A}^{[h]}\bigr).
\end{equation}
Here $\mathbf{A}^{[h]}$ stands for all further arguments of $\mathfrak{E}$ other than the pair,
namely its pointwise coordinates and its frame data in the sense of
Definition~\ref{4.22:def-coordinates-frames}, each rotated at its law site by the site map
$g(s):=h(\text{site of }s)$; for a frame $W$ the rotation is $\Phi_W^{[v]}$ with $v$ the
restriction of $h$ to $\Sigma_W$.
Then $F^{\mathfrak{b}}_{\mathfrak{E}}$ pre-rotates exactly those arguments of $\mathfrak{E}$,
other than the pair $(\mathbf{U},\mathbf{J})$, on which the joint $G$-invariance of
$\mathfrak{E}$ does not act as the identity. Every argument of $\mathfrak{E}$ other than
$(\mathbf{U},\mathbf{J})$ is of one of the following five kinds.
\begin{enumerate}
\item[(a)] \emph{Rim slots.} These are the rim fluctuation variables $A_{j'}(b)$ of the following terms:
the $\mathbf{B}^{(j)}$-terms; the terms of $\mathbb{B}''=\mathbb{C}^{(n)}$; the terms
$\mathbf{B}'^{(j)}$; the arrested terms $\mathbb{C}^{(i)}$; and the $Y_0$-terms that arise
from the $\mathbf{B}$-terms of the second group in \cite[\S1]{B-CMP122b}, which are of
$\mathbf{B}$-type. Their law site is $b_-$. They enter $\mathfrak{E}$ through its slot for
the fluctuation variable $A$, and this slot may be weak, strong or silent.
\item[(b)] \emph{Arrested collar variables.} These are the variables $B_{j'}(b)$ of a component
arrested in the branch \cite[(1.55)]{B-CMP122a}. Their law site is $b_-$. They enter
$\mathfrak{E}$ analytically through the background configuration
\begin{equation*}
U^{(n)}_k\bigl(\exp i[g_jCB_j-\dots]\,V^{(j)}\bigr),
\end{equation*}
with $C$ and $V^{(j)}$ as in the branch \cite[(1.55)]{B-CMP122a}.
They also enter through the cut-off function of Ba{\l}aban's change of variables.
Finally, they enter as a polynomial of degree two through the
kept Gaussian of \cite[(1.60)]{B-CMP122a},
\begin{equation}\label{4.22:eq-rotated-variables-2}
-\tfrac12\bigl\langle B_j,\,C^{*}\Delta^{(j)}C\,B_j\bigr\rangle ,
\end{equation}
with $C$ and $\Delta^{(j)}$ as there.
\item[(c)] \emph{Frame data.} These are the data $\Phi_W$ of an older component $W$ of the second class,
in the representation \cite[(1.102)]{B-CMP122b} of Definition~\ref{4.22:def-coordinates-frames}.
They comprise the $\mathfrak{g}$-valued variables $B$; the $G$-valued variables $V_h$, the
restriction of $V''$ to the core, and the inner variables $V_j$; the re-rooted inner
variables $A_j$; and the data of older frames. Their law sites are the points of
$\Sigma_W$. They enter through the inner background configurations $U''_{k_W}(B)$ and
$U_2(V'')$ and through the slots of the inner terms.
\item[(d)] \emph{The integration variables of the current step.} These are the variables
$\Phi$ of the current step, namely $B$, the restriction $V''\restriction D_2$ of $V''$ to the
region $D_2$ on which it is integrated at this step, $V_h$ and the inner tree variables.
They have no law site. They enter through the pair
$(\mathcal{U}_q^{\mathfrak{b}},\mathcal{J}_q^{\mathfrak{b}})$. They are not rotated:
the presentation leaves $\Phi$ untouched and makes no change of integration
variables, so that $B$ is unrotated and real.
\item[(e)] \emph{Gauge-field data.} These are the variables $V_j$ on determining sets, including those
of the arrest branches other than that of \cite[(1.55)]{B-CMP122a}. They have no law site. They are
read through $\mathbf{U}$: the term depends on the gauge-field variable $V$ through
the background field $U_k(V)$ \cite[\S2]{B-CMP119}. They have no slot of their own. In the
extension of the term to complex rotations of its arguments they are replaced by a
function of $\mathbf{U}$.
\end{enumerate}
The arguments so pre-rotated are exactly those of kinds (a), (b) and (c), each rotated at its
law site; the arguments of kinds (d) and (e) are not rotated.
\end{proposition}

\begin{proof}
\emph{Exactness.} The rotated presentation \eqref{4.22:eq-rotated-variables-1} is defined
by inserting into $\mathfrak{E}$ the transformed pair
$\bigl((\hat{\mathcal{U}}^{\mathfrak{b}})^{h},\operatorname{Ad}_h\hat{\mathcal{J}}^{\mathfrak{b}}\bigr)$
together with the rotated further arguments $\mathbf{A}^{[h]}$. It therefore applies to
$\mathfrak{E}$ a single $G$-valued map $h$, acting on the pair by gauge transformation and on
the further arguments by the site law, and the joint $G$-invariance of $\mathfrak{E}$ used
here is invariance under this one map.

The site law is the action of the gauge group on the stored variables.
For fluctuation-type variables it is the rotation at the variable's own block-lattice site
$b_-$. For variables declared in a frame rooted at $y_0^W$ it is the induced rotation at
the root, as in the rotation $\Phi_W^{[v]}$ of Definition~\ref{4.22:def-coordinates-frames}.

Consequently the arguments of $\mathfrak{E}$ other than $(\mathbf{U},\mathbf{J})$ that
\eqref{4.22:eq-rotated-variables-1} pre-rotates are precisely those on which the site law
does not act as the identity. It remains to show that every argument of a term of $\mathcal{C}$
belongs to one of the five kinds listed. We go through the sources of arguments in turn.

\emph{Terms of \cite{B-CMP119}.} By \cite[(2.23)]{B-CMP119}, the action is a function of $U_k$ and
of the system of fluctuation fields $\{A_{j-1}\}$. The fluctuation fields give kind~(a);
the dependence on $U_k$ gives kind~(e).

\emph{Terms added by \cite{B-CMP122a}.} We go through the arrest branches one at a time.
\begin{itemize}
\item In the branches in which the integral \cite[(1.25)]{B-CMP122a} is kept in the
variables $V_j$ and $A_j$, nothing new is added.
\item In the branch of \cite[(1.55)]{B-CMP122a}, the collar variables $B'_j$ are added.
This is kind~(b).
\item In the branches producing the structure $\mathbb{B}''_k$, that structure is added
together with its data $V''$. This is kind~(e).
\end{itemize}
The splitting $V''=V'V_0$ of \cite[(1.81)]{B-CMP122a} adds nothing. It occurs only inside
characteristic functions and $\delta$-functions. These belong to the
$\mathbf{T}$-operations and not to the terms.

\emph{Terms added by \cite{B-CMP122b}.} A term there depends on the new background field
$U_k$ restricted to $Y$, and also on the fields $A$, $B$ and $V''$ localised in $Y$
\cite[\S1]{B-CMP122b}.
\begin{itemize}
\item For components of the first class, these fields are integrated out by
$\mathbf{T}'_k$.
\item For components of the second class, they remain under $\mathbf{T}'_k(Y_i)$ in the
representation \cite[(1.102)]{B-CMP122b}. They are then frame data of an older component,
which is kind~(c).
\end{itemize}

\emph{$\mathbf{E}$- and $\mathbf{R}$-terms.} These carry no further argument, by
\cite[(2.27)(i), (2.30)]{B-CMP119}.

\emph{The two representations.} The proposition is stated in the representation
\cite[(1.102)]{B-CMP122b}. In the representation \cite[(1.104)]{B-CMP122b}, there are no
integral operations connected with the terms of second-class components. These terms then
depend only on the new field variables $V_k$ restricted to $Y_i$, so kind~(c) is void.

The entry of the arrested variable $B_j$ through the kept Gaussian
\eqref{4.22:eq-rotated-variables-2} is of kind~(b), as recorded there.

This exhausts the sources of arguments of the terms of $\mathcal{C}$. This proves the
proposition.
\end{proof}

\begin{definition}[Orbit data and the rotated sup norm]\label{4.22:def-orbit-datum}
Let $(c,S,\mathcal{S})$ be a stored term with domain $S$ and strong sets $\mathcal{S}$. Let
$\mathbf{A}$ be its pointwise coordinates, $\boldsymbol{\Phi}$ its frames and $\mathfrak{s}(S)$ its
law sites, as in Definition~\ref{4.22:def-coordinates-frames}. Let $P=(\mathbb{U},\mathbb{J})$ denote
the background pair. Thus $c=c(S;P,\mathbf{A},\boldsymbol{\Phi})$ is defined for
\begin{equation}\label{4.22:eq-orbit-datum-1}
(P,\mathbf{A},\boldsymbol{\Phi})\in
\mathfrak{U}(S)\times\mathfrak{D}_{\mathcal{S}}(S)\times\operatorname{supp}(\boldsymbol{\Phi}).
\end{equation}
The following properties of the stored term are used.
\begin{itemize}
\item The datum $A^{\flat}$ attached to the stored term is among the pointwise coordinates
$\mathbf{A}$. The coordinates in which it enters $\mathbf{A}$ are weak, at the radius
$e^{-2\varepsilon_{\star}}\varrho^{\mathsf{B}}$ fixed for the weak collar variables in the lemma on
weak collar variables. The exception is its coordinates on the bonds of the set
$\mathfrak{S}^{22}$ of Ba{\l}aban's full change of variables, which are silent.
\item $\boldsymbol{\Phi}$ is real and ranges in the supports of its measures, and these supports
are stable under $\operatorname{Ad}_G$.
\item $\mathfrak{U}(S)$ is open and $\mathcal{G}$-invariant.
\item $c$ is jointly invariant: for every $G$-valued site map $v$,
\begin{equation}\label{4.22:eq-orbit-datum-2}
c\bigl(S;P^{v},\mathbf{A}^{[v]},\boldsymbol{\Phi}^{[v]}\bigr)=c(S;P,\mathbf{A},\boldsymbol{\Phi}).
\end{equation}
This is the invariance of \cite[(2.41)(iii)]{B-CMP119}, extended to stored terms by clause (ii)
of the gauge-invariance theorem of this chapter.
\end{itemize}
Moreover, $\mathfrak{D}_{\mathcal{S}}(S)$ is stable under $\mathbf{A}\mapsto\mathbf{A}^{[g]}$ for every
$G$-valued map $g$ on the law sites $\mathfrak{s}(S)$. The reason is that the boxes and polydiscs
of which it is built are balls for the norm $|\cdot|$.

Let $\varepsilon\colon\mathfrak{s}(S)\to\,]0,\infty[$ satisfy $\varepsilon(s)\le\mathsf{e}(s)$ for
every $s\in\mathfrak{s}(S)$, where $\mathsf{e}(s)$ is the unit of the law site $s$
(Definition~\ref{4.22:def-rotation-radius}). Let $\mathfrak{G}_{\varepsilon}(S)$ be the set of
rotations of Definition~\ref{4.22:def-coordinates-frames}, that is, of maps
$g\colon\mathfrak{s}(S)\to G^{\mathbb{C}}$ with $\operatorname{pol}(g(s))<\varepsilon(s)$
at every law site. An \emph{orbit datum of radius $\varepsilon$} for $c$ is a continuous function
\begin{equation}\label{4.22:eq-orbit-datum-3}
c^{\sharp}\colon
\mathfrak{U}(S)\times\mathfrak{D}_{\mathcal{S}}(S)\times\operatorname{supp}(\boldsymbol{\Phi})
\times\mathfrak{G}_{\varepsilon}(S)\longrightarrow\mathbb{C}
\end{equation}
with the following two properties. Property \textup{(O1)} is that
\begin{equation}\label{4.22:eq-orbit-datum-a}
\begin{gathered}
\text{for fixed real values of the coordinates of $\mathbf{A}$ that are kept real,}\\
\text{namely the strong and the silent ones, and for fixed $\boldsymbol{\Phi}$,}\\
c^{\sharp}(P;\mathbf{A},\boldsymbol{\Phi},g)\ \text{is holomorphic jointly in $P$,}\\
\text{in the weak coordinates of $\mathbf{A}$, and in $g$;}
\end{gathered}
\end{equation}
and property \textup{(O2)} is that, for every $G$-valued $g\in\mathfrak{G}_{\varepsilon}(S)$,
\begin{equation}\label{4.22:eq-orbit-datum-b}
c^{\sharp}(P;\mathbf{A},\boldsymbol{\Phi},g)=c\bigl(S;P,\mathbf{A}^{[g]},\boldsymbol{\Phi}^{[g]}\bigr).
\end{equation}
At $G$-valued $g$ the right member of \eqref{4.22:eq-orbit-datum-b} is defined for two reasons.
First, $\mathfrak{D}_{\mathcal{S}}(S)$ is stable under $G$-valued rotations. Second, the supports of
the measures of $\boldsymbol{\Phi}$ are $\operatorname{Ad}_G$-stable.

The \emph{rotated sup norm} of $c$ at radius $\varepsilon$ is
\begin{equation}\label{4.22:eq-orbit-datum-4}
\mathfrak{N}^{\mathbb{C}}_{\varepsilon}(c)
:=\sup\bigl|c^{\sharp}(P;\mathbf{A},\boldsymbol{\Phi},g)\bigr|,
\end{equation}
the supremum taken over the whole domain in \eqref{4.22:eq-orbit-datum-3}. Take for $g$ the
constant map $g\equiv 1$ in \eqref{4.22:eq-orbit-datum-b}. Since $\operatorname{Ad}_{1}$ is the
identity, $\mathbf{A}^{[1]}=\mathbf{A}$ and $\boldsymbol{\Phi}^{[1]}=\boldsymbol{\Phi}$, so
$c^{\sharp}(\cdot\,,1)=c$. Hence
$\mathfrak{N}^{\mathbb{C}}_{\varepsilon}(c)\ge\mathfrak{N}(c)$.

Suppose $c^{\sharp}$ is built from the holomorphic lift of the full change of variables, with the
cut-off data $(\mathsf{s},\mathsf{v})$ frozen. Then the supremum in
\eqref{4.22:eq-orbit-datum-4} is taken also over the admissible frozen values of
$(\mathsf{s},\mathsf{v})$.

The norms $\|\cdot\|^{A,\mathbb{C}}$ and $\|\cdot\|^{A,\flat,\mathbb{C}}_{(n)}$ are the norms
$\|\cdot\|^{A}$ and $\|\cdot\|^{A,\flat}_{(n)}$ of stored terms defined earlier in this chapter, with
$\mathfrak{N}$ replaced by $\mathfrak{N}^{\mathbb{C}}_{\varepsilon}$ throughout. The strong and silent
coordinates and the frame data stay real. The strong sets $\mathcal{S}$, the weights and the weak
radii are unchanged.

The quantity $\mathfrak{N}^{\mathbb{C}}_{\varepsilon}(c)$ carries no length. In a pointwise bound,
as in the pointwise lemma for the rotated sup norm, the length stands in the right member. In the
norms $\|\cdot\|^{A,\flat,\mathbb{C}}_{(n)}$ it stands in the left member, through the relation
between the anchored sum and the label weights: for terms of the class $\mathfrak{S}$ (the terms
touched by a second-class outcome), the anchored sum of these norms is taken on labels, and so is
their weight $e^{a_m\mathsf{d}_{\ell}(F)}$, where $\mathsf{d}_{\ell}(F)$ is the length of the term
$F$ on labels and $a_m$ is the rate of this weight in the norms $\|\cdot\|^{A,\flat}_{(n)}$.
\end{definition}

\begin{definition}[The radius budget at a site: decrements, margins and reserve]\label{4.22:def-radius-budget}
Let $\varepsilon_{\star}$, the unit $\mathsf{e}(s)$ of a law site $s$ and $\mu_k=\mu(x_k)$ be as in
Definition~\ref{4.22:def-rotation-radius}. For a complex rotation $g\in G^{\mathbb{C}}$ write
$\operatorname{pol}(g)=\log\max\{\|g\|,\|g^{-1}\|\}$ for its polar size.

Each step $k$ acts on a term at three sites:
\begin{itemize}
\item $(\mathrm{T})$, the fluctuation integral;
\item $(\mathrm{P})$, the large-field stages;
\item $(\mathrm{L})$, the presentation.
\end{itemize}
For a law site $s$ and a step $k$ we define the following quantities.

\begin{enumerate}
\item[(a)] The \emph{decrement at the fluctuation-integral site}, $\eta^T_k(s)$, is the supremum of
$\operatorname{pol}(\mathfrak{u}_p(0,\sigma)(s))$ over the units read at the site $(\mathrm{T})$ of step $k$
whose chart contains $s$. Here $\mathfrak{u}_p(\mathsf{t},\sigma)$ is the $G^{\mathbb{C}}$-valued response
rotation of a unit given by the response lemma, and $(s)$ denotes its value at $s$. The supremum is taken
under the following constraints:
\begin{itemize}
\item the Cauchy variable is $\mathsf{t}=0$;
\item the interpolation parameters satisfy $|\sigma(\Delta)|\le e^{\kappa_1}$ for every cell $\Delta$;
\item the configuration lies in the stored tube of the unit;
\item the lifted variables satisfy $|\mathsf{B}|<T^{\mathbb{C}}$.
\end{itemize}
The cells carrying $\sigma$ are those of the index set $\sigma_0$ as the step runs it. For units that are
not touched by a second-class outcome this is the index set of the correlation theorem. For
units that are touched by one, it is the index set of the sourced interpolation at the
fluctuation-integral site, and all of its cells are included.

\item[(b)] The \emph{decrement at the large-field stages} is
\begin{equation}\label{4.22:eq-radius-budget-1}
\eta^P_k(s)=\sum_{l\le N_k}\sup\operatorname{pol}\bigl(\mathfrak{u}^{(l)}_{t,\sigma}(b)(s)\bigr).
\end{equation}
For each $l$ the supremum runs over the rotation families of the $l$-th stage that the response lemma
supplies for the proofs of the stages; $\mathfrak{u}^{(l)}_{t,\sigma}(b)$ denotes the
$G^{\mathbb{C}}$-valued rotation of such a family at the parameters $t,\sigma$ and the background $b$, and
$(s)$ its value at $s$. The constraints are $|t_{\square}|\le R^{\natural}_{\square}$,
$|\sigma|\le e^{\kappa_1}$, and $b$ complex in its ball.

\item[(c)] The \emph{decrement at the presentation}, $\eta^L_k(s)$, is a sum of suprema of
$\operatorname{pol}(h^{\mathfrak{b}}(\zeta;\Phi)(s))$, where $h^{\mathfrak{b}}(\zeta;\Phi)$ is the
$G^{\mathbb{C}}$-valued rotation given by the presenter map of the pair $(\mathfrak{E},\mathfrak{b})$ and
$(s)$ its value at $s$. The sum runs over the at most two successive
presentations that a term undergoes at step $k$: first in the second part, then, for its $Y_0$-term (the
term attached to the region $Y_0$ of the first part of Ba{\l}aban's operation), in
the first part \cite{B-CMP122b}. For each presentation the supremum runs over the following ranges:
\begin{itemize}
\item the pairs $(\mathfrak{E},\mathfrak{b})$ presented with $s$ in the chart of $\mathfrak{E}$;
\item $\zeta\in\mathfrak{S}_{\mathrm{src}}$;
\item $\Phi$ real in the supports;
\item $|\sigma|\le e^{\kappa_1}$;
\item $t$ in the range over which the presentation runs it.
\end{itemize}
The cells carrying $\sigma$ are again those of $\sigma_0$ as the step runs it. For units not touched by a
second-class outcome, $\sigma_0$ is the index set of \cite{B-CMP122b}. On units touched by a second-class
outcome, and on every interior unit of the first part, the index sets used are those of the two sourced
presentations, the first and the second. For a bracket, $t$ ranges over $[0,1]$.
\end{enumerate}
Each of $\eta^T_k(s)$, $\eta^P_k(s)$, $\eta^L_k(s)$ is set equal to $0$ if there is no such reading at
step $k$.

The \emph{margins} at step $k$ total $3\mu_k\mathsf{e}(s)/\varepsilon_{\star}$. One margin
$\mu_k\mathsf{e}(s)/\varepsilon_{\star}$ is spent at $(\mathrm{T})$, by the rotation bound at the
fluctuation-integral site. One is spent at each of the at most two presentations, by item (b) of the
rotation bound for presented units, and a presentation's margin is a single allowance shared by the
bounds of all its pieces, in the sense of the shared-room lemma. The \emph{reserve} at step $k$ is
\begin{equation*}
\eta^{\mathrm{inc}}_k(s)=\mu_k\mathsf{e}(s)/\varepsilon_{\star}.
\end{equation*}
It is fixed by this choice and is not drawn on in the correlation theorem.

The \emph{total debit} and the \emph{current radius} at $s$ are
\begin{equation}\label{4.22:eq-radius-budget-2}
\begin{aligned}
\eta_k(s)&=\eta^T_k(s)+\eta^P_k(s)+\eta^L_k(s)+3\mu_k\mathsf{e}(s)/\varepsilon_{\star},\\
\varepsilon_k(s)&=\mathsf{e}(s)-\sum_{k'<k}\bigl[\eta_{k'}(s)+\eta^{\mathrm{inc}}_{k'}(s)\bigr].
\end{aligned}
\end{equation}
The sum runs over the steps $k'$ after the creation of $s$ and before its account is closed.

The site $s$ is created in one of three ways:
\begin{itemize}
\item at the fluctuation-integral site that creates the rim slot containing it;
\item at an arrest;
\item at a second-class outcome.
\end{itemize}
Its account is closed in one of two ways:
\begin{itemize}
\item when the variable at $s$ is integrated;
\item when its component is processed and $s$ joins a new frame, in the sense of item (v) of the frame
lemma.
\end{itemize}

Within step $k$ the radius is lowered site by site. After the site $(\mathrm{T})$ it is
\begin{equation*}
\varepsilon_k(s)-\eta^T_k(s)-\mu_k\mathsf{e}(s)/\varepsilon_{\star}.
\end{equation*}
After the $l$-th large-field stage it is lowered further by the $l$-th summand of
\eqref{4.22:eq-radius-budget-1}, and so on through the presentations. In what follows ``$\varepsilon_k$''
denotes the radius current at the site in question.

The \emph{budget condition} is
\begin{equation}\label{4.22:eq-radius-budget-a}
\sup_s\ \frac{1}{\mathsf{e}(s)}\sum_{k'}\bigl[\eta_{k'}(s)+\eta^{\mathrm{inc}}_{k'}(s)\bigr]\le\frac14,
\qquad
\sup_k\ \varepsilon_{\mathrm{live}}(k)\le\frac14\,\varepsilon_{\star},
\end{equation}
where $\varepsilon_{\mathrm{live}}(k)$ is the live rotation size of step $k$. The first inequality is
supplied by the theorem bounding the total debit, under the conditions on the live radius. The second is
the fourth of those conditions. Under \eqref{4.22:eq-radius-budget-a}, for every law site $s$ and every
step $k$,
\begin{equation}\label{4.22:eq-radius-budget-3}
\tfrac34\,\mathsf{e}(s)\le\varepsilon_k(s)\le\mathsf{e}(s).
\end{equation}
\end{definition}

\begin{proof}
We prove \eqref{4.22:eq-radius-budget-3}.

First, every summand in \eqref{4.22:eq-radius-budget-2} is non-negative. For $g\in G^{\mathbb{C}}$ we
have $1=\|gg^{-1}\|\le\|g\|\,\|g^{-1}\|$, so $\max\{\|g\|,\|g^{-1}\|\}\ge1$ and hence
$\operatorname{pol}(g)\ge0$. Each of $\eta^T_k(s)$, $\eta^P_k(s)$, $\eta^L_k(s)$ is a supremum, or a
finite sum of suprema, of such quantities, or is $0$; so each is non-negative. The margin is
non-negative because
\begin{equation*}
\mu_k=\varepsilon_{\star}(x_k-c_5)^{-2}\ge0,
\end{equation*}
$\varepsilon_{\star}>0$, and $\mathsf{e}(s)>0$. The last fact holds by
Definition~\ref{4.22:def-rotation-radius}, since $\mathsf{e}(s)$ is either $\varepsilon_{\star}$ or
$\varepsilon^{\mathrm{fr}}(W)$, both positive. The same two facts give $\eta^{\mathrm{inc}}_k(s)\ge0$.

The radius current at a site of step $k$ is $\varepsilon_k(s)$ diminished by those terms of
$\eta_k(s)$ that have already been spent at the earlier sites of step $k$. It is therefore of the form
\begin{equation*}
\mathsf{e}(s)-\sum_{k'<k}\bigl[\eta_{k'}(s)+\eta^{\mathrm{inc}}_{k'}(s)\bigr]-\delta,
\qquad 0\le\delta\le\eta_k(s),
\end{equation*}
where $\delta$ is a partial sum of the non-negative summands of $\eta_k(s)$.

Since every subtracted quantity is non-negative, the current radius is at most $\mathsf{e}(s)$. This is
the upper bound in \eqref{4.22:eq-radius-budget-3}.

For the lower bound, note that the subtracted total does not exceed the full sum
$\sum_{k'}[\eta_{k'}(s)+\eta^{\mathrm{inc}}_{k'}(s)]$ over all steps of the account of $s$, because the
terms omitted are non-negative. By the first inequality of \eqref{4.22:eq-radius-budget-a}, multiplied by
$\mathsf{e}(s)>0$, this full sum is at most $\tfrac14\mathsf{e}(s)$. Hence the current radius is at least
\begin{equation*}
\mathsf{e}(s)-\tfrac14\mathsf{e}(s)=\tfrac34\,\mathsf{e}(s),
\end{equation*}
which is the lower bound in \eqref{4.22:eq-radius-budget-3}.
\end{proof}

Recall from Definition~\ref{4.22:def-radius-budget} that the current radius $\varepsilon_k(s)$ is a
function of the law site $s$ alone. The decrements $\eta^T_k(s)$, $\eta^P_k(s)$, $\eta^L_k(s)$ entering it
are suprema over all units whose chart contains $s$. Hence a term of the plain class and a
wrapped term whose charts both contain $s$ carry the same radius at $s$. An example in which a plain term
and a wrapped term share such a site is given further on.

Let $\mathfrak{A}$ be the plain class of Definition~\ref{4.22:def-cores-wrapped-class}. Let
$\eta^{T,\mathfrak{A}}_k(s)$, $\eta^{P,\mathfrak{A}}_k(s)$, $\eta^{L,\mathfrak{A}}_k(s)$ be the same suprema,
restricted to the units of $\mathfrak{A}$ and to the families on which those units are run:
\begin{itemize}
\item at the fluctuation-integral site $(\mathrm{T})$, the families of the correlation theorem;
\item at the large-field stages $(\mathrm{P})$, the families of the statement on which the large-field
stages are carried out;
\item at the presentation site $(\mathrm{L})$, Ba{\l}aban's families, the brackets of the correlation
theorem at $\mathsf{t}=0$, and, on interior plain units, the families of the two presentation-site
statements for the wrapped class.
\end{itemize}
Put
\begin{equation}\label{4.22:eq-lineage-1}
\begin{aligned}
\eta^{\mathfrak{A}}_k(s)
&:=\eta^{T,\mathfrak{A}}_k(s)+\eta^{P,\mathfrak{A}}_k(s)+\eta^{L,\mathfrak{A}}_k(s)
+3\mu_k\,\mathsf{e}(s)/\varepsilon_{\star},\\
\varepsilon^{\mathfrak{A}}_k(s)
&:=\mathsf{e}(s)-\sum_{k'<k}\bigl[\eta^{\mathfrak{A}}_{k'}(s)+\eta^{\mathrm{inc}}_{k'}(s)\bigr].
\end{aligned}
\end{equation}
Each restricted supremum is taken over a subcollection of the units and families of the unrestricted one.
By the convention of Definition~\ref{4.22:def-radius-budget}, a decrement is zero when there is no reading.
Hence $\eta^{T,\mathfrak{A}}_{k'}(s)\le\eta^{T}_{k'}(s)$, and likewise for the decrements at $(\mathrm{P})$
and $(\mathrm{L})$. The margin and the reserve
terms are the same in both radii. Comparing \eqref{4.22:eq-lineage-1} with the definition of
$\varepsilon_k(s)$ term by term therefore gives
\begin{equation}\label{4.22:eq-lineage-2}
\varepsilon^{\mathfrak{A}}_k(s)\ \ge\ \varepsilon_k(s).
\end{equation}

\begin{lemma}[The radius is lowered only along a term's lineage]\label{4.22:lem-lineage}
Let $F$ be a stored term and $s$ a law site of its domain.
\begin{enumerate}
\item The radius at $s$ of the orbit datum $F^{\sharp}$ is lowered only at readings of a term of the
lineage of $F$, that is, of $F$ itself or of a constituent of $F$ at any depth.
\item At each such reading, the radius is lowered by the rotation applied at that reading.
\item Consequently every term of $\mathfrak{A}$ has an orbit datum of radius $\varepsilon^{\mathfrak{A}}_k$
and, by restriction, one of radius $\varepsilon_k$.
\end{enumerate}
\end{lemma}

\begin{proof}
The radius of an orbit datum is changed only in the following four statements:
\begin{itemize}
\item the complexified corollary for the effective density;
\item the complexified corollary for the large-field stages;
\item part~(a) of the presentation proposition;
\item part~(iii) of the transfer lemma.
\end{itemize}
Each of these statements acts on the unit that is read at the site in question. It lowers the
radius of the output built from that unit, and of nothing else, by the decrement of that one reading
increased by its margin.

The suprema over all units whose chart contains $s$ were taken in Definition~\ref{4.22:def-radius-budget}
only to have one letter for the decrement at $s$. They play no role in the amount by which a given output
is lowered. Hence the radius of $F^{\sharp}$ at $s$ changes only when the unit being read is $F$ or a term
from which $F$ is built, i.e.\ a constituent of $F$ at some depth. At such a reading it changes by the
rotation applied there. This proves the first two assertions.

Now let $F\in\mathfrak{A}$. By Definition~\ref{4.22:def-cores-wrapped-class}, the class $\mathfrak{A}$ is
closed under passage to constituents. Hence every term of the lineage of $F$ lies in $\mathfrak{A}$. By
the first assertion, every lowering of the radius of $F^{\sharp}$ at $s$ during step $k'$ occurs at a
reading of a unit of $\mathfrak{A}$, on one of the families on which such units are run.

At each of the three sites of step $k'$, this lowering is therefore bounded by the corresponding
restricted supremum $\eta^{T,\mathfrak{A}}_{k'}(s)$, $\eta^{P,\mathfrak{A}}_{k'}(s)$ or
$\eta^{L,\mathfrak{A}}_{k'}(s)$. The margins of step $k'$ contribute $3\mu_{k'}\mathsf{e}(s)/\varepsilon_{\star}$,
and the reserve contributes $\eta^{\mathrm{inc}}_{k'}(s)$. Starting from $\mathsf{e}(s)$ and summing over
$k'<k$, the radius of $F^{\sharp}$ at $s$ is at least $\varepsilon^{\mathfrak{A}}_k(s)$, as defined in
\eqref{4.22:eq-lineage-1}. That is, $F^{\sharp}$ is an orbit datum of radius $\varepsilon^{\mathfrak{A}}_k$.

Finally, an orbit datum of radius $\varepsilon$ restricts to an orbit datum of every radius
$\varepsilon'\le\varepsilon$, and
\begin{equation*}
\mathfrak{N}^{\mathbb{C}}_{\varepsilon'}(F^{\sharp})\le\mathfrak{N}^{\mathbb{C}}_{\varepsilon}(F^{\sharp}).
\end{equation*}
By \eqref{4.22:eq-lineage-2}, $\varepsilon_k\le\varepsilon^{\mathfrak{A}}_k$. Restriction therefore gives an
orbit datum of radius $\varepsilon_k$, with chart supremum no larger than at radius
$\varepsilon^{\mathfrak{A}}_k$.
\end{proof}

This lemma is used as follows. Conjunct~(B) of the theorem of this section is proved with the radius
$\varepsilon^{\mathfrak{A}}_k$ and is read with the radius $\varepsilon_k$. Its proof uses nothing of the
wrapped terms, nor of the families on which the fluctuation-integral statement for those terms is run.

The two radii $\varepsilon_k$ and $\varepsilon^{\mathfrak{A}}_k$ obey one and the same majorant. It is the
majorant of part~(iii) of the decrement lemma, and of that lemma's version for the wrapped class.
Accordingly, the theorem bounding the accumulated decrements relative to $\mathsf{e}(s)$ is a single
theorem for both radii.

\begin{definition}[Second-class components, cores and the wrapped class of stored terms]
\label{4.22:def-cores-wrapped-class}
Steps, their sites $(\mathrm{T})$, $(\mathrm{P})$, $(\mathrm{L})$ and the $\mathbf{R}$-site, the scales
$\check{R}_j$ of \eqref{4.22:eq-steps-sites-b}, $N_k=\check{R}_k$ and $h=k-N_k$ are those of
Notation~\ref{4.22:not-steps-sites}. Frames $W$, their levels $k_W$ and their data are those of
Definition~\ref{4.22:def-coordinates-frames}.

\emph{Second-class components.} A \emph{second-class component} is a component $Y$ that took the
second class at the $\mathbf{R}$-site of some step $m$, in the representation
\cite[(1.102)]{B-CMP122b}; we record it together with its step as the pair $(Y,m)$. It is
\emph{live} at a later moment if no $\mathbf{R}$-site of a later
step has carried the component of the large-field region containing $Y$ through the presentation
$(\mathrm{L})$. Its \emph{strip} $Y^{\sharp}$ is $Y$ together with $m_{\sharp}$ layers of $M$-cubes
of level $m$ around it, where
\begin{equation}\label{4.22:eq-cores-wrapped-class-1}
m_{\sharp}:=\bigl\lfloor \tfrac12(\check{R}_m-1)\bigr\rfloor ;
\end{equation}
it is one and the same set at all later levels, and it is the set written $Y^{\sim}$ in
\cite{B-CMP122b}, with the width $m_{\sharp}$ used for $Y^{\sim}$ in the correlation theorem. In the
representation \cite[(1.102)]{B-CMP122b} a second-class component is a frame, with $W=Y$ and
$k_W=m$.

\emph{The format.} A stored boundary-type term $F$ of the correlation theorem carries a pair
$(A_F,\mathcal{P}_F)$, consisting of a \emph{label} $A_F\in\mathbf{D}_\ell$, where $\ell$ is the
level at which $F$ is stored, and a set
$\mathcal{P}_F$ of live second-class components whose strips touch $A_F$. Its \emph{true domain}
is
\begin{equation}\label{4.22:eq-cores-wrapped-class-2}
X_F:=A_F\cup\bigcup_{(Y,m)\in\mathcal{P}_F}Y^{\sharp}.
\end{equation}
The term $F$ is \emph{wrapped} at a given moment if $\mathcal{P}_F\neq\emptyset$. This is not the
condition that the true domain of $F$ contain a live strip: a term with $\mathcal{P}_F=\emptyset$
whose label covers a strip cube by cube is not wrapped; it is bounded, like every term with
$\mathcal{P}_F=\emptyset$, by the estimates of the correlation theorem.

\emph{Cores.} Let $C$ be a component of the class $Z$ treated at step $k$. Its \emph{core} is
\begin{equation}\label{4.22:eq-cores-wrapped-class-3}
\mathsf{C}_C:=\bigl(\Omega''^{\,\sim2}_{h+1}\bigr)^{c}\cap C=Z'^{\,\sim}_{h}\cap C ,
\end{equation}
in the notation of \cite{B-CMP122a}: the part of $C$ that the presentation $(\mathrm{L})$ never
integrates. A first-part unit $v$, that is, one of the stored terms consumed by the presentation
$(\mathrm{L})$ of the step, with true domain $X_v$ is \emph{interior} if $X_v\subset Z$,
\emph{deep} if $X_v\subset\mathsf{C}_C$ for some component $C$, and \emph{shallow} if it is
interior and not deep.

\emph{The classes.} For a second-class component $Y_i$, let $\mathfrak{K}(Y_i)$ be the set of the
two members depending on $(\mathbf{U},\mathbf{J})$ that are kept under $\mathbf{T}'_{k_W}(Y_i)$ in
\cite[(1.102)]{B-CMP122b} (here $W=Y_i$): the term $V(Y_i^{\sim})$, and the quadratic form of
\cite[(1.71)]{B-CMP122b} (see \cite[p.~378]{B-CMP122b}),
\begin{equation}\label{4.22:eq-cores-wrapped-class-4}
\mathsf{Q}_i:=-\tfrac12\bigl\langle DH''_{1,k_W,Y_i}B,\;
\zeta\, DH''_{1,k_W,Y_i}B\bigr\rangle ,
\end{equation}
with $\zeta$ as in \cite[(1.71)]{B-CMP122b} (Ba{\l}aban's operator of that display, not the
complex-source shift $\zeta$ of the glossary). Put $\mathfrak{K}:=\bigcup_i\mathfrak{K}(Y_i)$; its
elements are the \emph{kept members}. The class $\mathfrak{S}$ is the least class of stored terms
(of all steps and all labels) such that
\begin{enumerate}
\item[(i)] every term that is wrapped at some moment of its life belongs to $\mathfrak{S}$;
\item[(ii)] $\mathfrak{K}\subset\mathfrak{S}$;
\item[(iii)] a term created at a site from constituents one of which belongs to $\mathfrak{S}$
belongs to $\mathfrak{S}$.
\end{enumerate}
In (iii) the constituents are: at $(\mathrm{T})$, the terms $V(Y)$, $Y\subset X$, on which the
output depends, where $X$ is the domain of the output, and the units whose pieces are assigned to
that site;
at $(\mathrm{P})$, the
elements of the cluster expansion in the output's cluster and their kept members; at
$(\mathrm{L})$, the terms presented into $V(Y)$, $\mathbf{R}'^{(k)}$ and $\mathbf{B}'^{(k)}$, and
the kept members. A stored term being a sum of products, (iii) means that some summand contains a
constituent belonging to $\mathfrak{S}$. The class $\mathfrak{A}$ consists of all other stored
terms.

\emph{The assigned length.} For $F\in\mathfrak{A}$, $\mathsf{d}_\ell(F):=d_\ell(X_F)$, Ba{\l}aban's
length. For $F\in\mathfrak{S}$ in the format, $\mathsf{d}_\ell(F):=d_\ell(A_F)$, where the label
$A_F$ is as follows: at the birth of a term $\mathbf{B}'^{(k)}(X)$ of the first group, $A$ is the
skeleton $\operatorname{sk}X$ of its domain $X$ (by the skeleton statement and clause (c) of the
proposition on the format of boundary-type terms); for descendants born at $(\mathrm{T})$, at a
chain of the large-field stages $(\mathrm{P})$ or at
$(\mathrm{L})$, $A$ is the label built at that site; for a healed descendant
($\mathcal{P}_F=\emptyset$), $A=X$ and $\mathsf{d}=d$. For $v\in\mathfrak{K}(Y_i)$,
\begin{equation}\label{4.22:eq-cores-wrapped-class-5}
A_v:=\emptyset,\qquad \mathcal{P}_v:=\{(Y_i,k_W)\},\qquad \mathsf{d}(v):=0 ;
\end{equation}
a piece of $v$ has the label built at its site, $\operatorname{lab}(\mathsf{y})$, which touches
$Y_i^{\sharp}$, and as $\mathsf{d}$ its plain length. The kept members are not in the format: the
format of a first-group term asks $X\not\subset W$, and the relative format requires the label to
meet $\Omega_j$, that is, $A_F\cap\Omega_j\neq\emptyset$. That the estimates accept them is clause
(leaf-3) of the
source-tube hypothesis \eqref{4.22:eq-source-tube-a}, which is proved in the section on the source
tube. When a kept member enters the estimates at a site of level $l$, it is given the one-cube label
$Q_{Y_i}(l)$, a single cube of level $l$, with $d_l=0$; this is clause (leaf-1) of
\eqref{4.22:eq-source-tube-a}. At
$(\mathrm{P})$ the level is $l:=m_v$, namely $h+1$ at the member's own chain and $k$ at a foreign
chain, and never the step's $k$ at the own chain. The
length $\mathsf{d}$ is a datum of the term, read off from its construction; it is not a function of
its true domain.

\emph{Properties.}
\begin{enumerate}
\item The least class $\mathfrak{S}$ exists. The relation \emph{is built from} is well-founded with
respect to the pair (step, site). A term of $\mathfrak{A}$ has no constituent in $\mathfrak{S}$ at
any depth, so $\mathfrak{A}$ is closed downward under \emph{constituent of}.
\item The class $\mathfrak{S}$ contains the set of terms wrapped at a given moment, in general
strictly: if a second-class component ceases to be live at step $k$, the term
$\mathbf{B}'^{(k)}(X)$ over the healed component $X$, and every term built from it, lies in
$\mathfrak{S}$ although it is wrapped round nothing.
\item Membership of $\mathfrak{S}$ is a datum of the construction, independent of $z$, and is fixed
at the birth of the term.
\item $\mathfrak{S}$ may be empty while the cores are full: if no second-class outcome occurs at any
step, then $\mathfrak{S}=\emptyset$; in a label with no second-class outcome every stored term of
that label lies in $\mathfrak{A}$, while the old terms $\mathbf{B}^{(j)}(X)$ with
$X\subset\mathsf{C}_C$ are present in as great a number as the past of the core allows. In
particular the plain deep units lie in $\mathfrak{A}$, and the estimates for deep units are needed
for terms of either class.
\item For every stored term $F$,
\begin{equation}\label{4.22:eq-cores-wrapped-class-6}
d_{\ell,\bigcup_{(Y,m)\in\mathcal{P}_F}Y^{\sharp}}(X_F)\le\mathsf{d}_\ell(F).
\end{equation}
\item At the birth of a first-group term $\mathbf{B}'^{(k)}(X)$, $\mathsf{d}_k=d_k(\operatorname{sk}X)$
is at most $d_{k,\cup Y_i}(X)$ plus a count of boundary cubes. This
excess is absorbed into the rate and the constant of the bound for the term concerned:
$(1+\tfrac12\beta)$ is replaced by $(1+\tfrac14\beta)$, and $C_{99}$ by $e^{\lambda_\theta}C_{99}$.
\end{enumerate}
The length $\mathsf{d}$ admits no lower bound by $d(X_F)$: a term whose true domain is a strip
together with a set $\Delta$, $X=Y^{\sharp}\cup\Delta$, can carry $\mathsf{d}=0$ while $d(X)$ is
unbounded, and \cite[(1.99)]{B-CMP122b} gives exactly the factor $e^{0}$.
\end{definition}

\begin{proof}
(1) The class of all stored terms satisfies (i)--(iii). Let a family of classes each satisfy
(i)--(iii). Conditions (i) and (ii) assert that fixed sets of terms are contained in a class, so
they hold for the intersection of the family. If a term has a constituent in the intersection, that
constituent lies in every member of the family; by (iii) for each member, the term lies in every
member, hence in the intersection. Thus the intersection of all classes satisfying (i)--(iii)
satisfies them, and it is the least such class.

Now let $F\in\mathfrak{A}$, and suppose $F$ had a constituent in $\mathfrak{S}$. Then (iii) would
put $F$ in $\mathfrak{S}$, which is impossible. Hence every constituent of $F$ lies in
$\mathfrak{A}$. Since \emph{is built from} is well-founded in (step, site), induction along it shows
that no constituent at any depth lies in $\mathfrak{S}$.

(2) Every term wrapped at a given moment lies in $\mathfrak{S}$ by (i). Suppose a second-class
component ceases to be live at step $k$. By (i), its wrapped terms lie
in $\mathfrak{S}$, and at that $(\mathrm{L})$ they are first-part units. Built from them is the term
$\mathbf{B}'^{(k)}(X)$ over the healed component, with $c_1=e^{-p_0(g_k)}$ and
$d_{k,\cup Y_i}(X)=d_k(X)$ by \cite{B-CMP122b}. By (iii) this term lies in $\mathfrak{S}$, although
it is wrapped round nothing. By (iii) again, so is every term built from it.

(3) Conditions (i)--(iii) refer only to wrappings, kept members and constituents, which are data of
the construction and do not involve $z$. Membership is fixed at birth because an old term that meets
a new second-class component $Y$ becomes a first-part unit and ceases.

(4) If no second-class outcome occurs at any step, no term is ever wrapped and
$\mathfrak{K}=\emptyset$. The empty class therefore satisfies (i)--(iii), so
$\mathfrak{S}=\emptyset$ and every stored term lies in $\mathfrak{A}$. Nevertheless the old terms
$\mathbf{B}^{(j)}(X)$ with
$X\subset\mathsf{C}_C$ are present in as great a number as the past of the core allows.

(5) For $F\in\mathfrak{A}$, $F$ is never wrapped, so $\mathcal{P}_F=\emptyset$, and both sides of
\eqref{4.22:eq-cores-wrapped-class-6} equal $d_\ell(X_F)$. For $v\in\mathfrak{K}(Y_i)$,
\eqref{4.22:eq-cores-wrapped-class-2} and \eqref{4.22:eq-cores-wrapped-class-5} give
$X_v=Y_i^{\sharp}$, so the left side is $0=\mathsf{d}(v)$. Now let $F$ be in the format. By the
corollary on the format of boundary-type terms, an optimal tree
of $A_F$ is admissible for the length of \cite[(1.67)]{B-CMP122b} relative to
$\bigcup_{\mathcal{P}_F}Y^{\sharp}$, because
\begin{equation*}
A_F\supseteq X_F\setminus\bigcup_{(Y,m)\in\mathcal{P}_F}Y^{\sharp}
\end{equation*}
by \eqref{4.22:eq-cores-wrapped-class-2}. The relative length is at most the length of that tree,
which is $d_\ell(A_F)=\mathsf{d}_\ell(F)$.

(6) This follows from clauses (b) and (c) of the skeleton statement, and from its
$\theta$-variant.
\end{proof}

\begin{definition}[The hypotheses of the holomorphic extension]\label{4.22:def-extension-hypotheses}
Fix a level $k$, let $\mathbf{R}_k$ be the large-field operation of step $k$, $Z$ its large-field
region and $\Lambda$ the small-field region of the step, and write $(\mathrm{T})$, $(\mathrm{P})$,
$(\mathrm{L})$ for the three sites of the step: the
fluctuation integral, the large-field stages and the presentation. Plain and wrapped units, their
cores and the wrapped class $\mathfrak{S}$ are as in
Definition~\ref{4.22:def-cores-wrapped-class}. The rotated sup norm
$\mathfrak{N}^{\mathbb{C}}_{\varepsilon}$ is that of Definition~\ref{4.22:def-orbit-datum}. The
rotation radius $\varepsilon_{\star}$, the frame unit and the margin $\mu_k$ are those of
Definition~\ref{4.22:def-rotation-radius}. The unit $\mathsf{e}(s)$ of a law site $s$ and the
reserve are those of the radius budget, Definition~\ref{4.22:def-radius-budget}. For a component $C$
of $Z$ we write $\mathsf{C}_C$ for its core and $\mathcal{C}$ for the class of presented units.
In this definition $\mathfrak{p}$ denotes a piece of the expansion of a unit, $\mathsf{d}$ the length
attached to a stored term (written $\mathsf{d}_{\ell}(F)$ for a term $F$; we write $\mathsf{d}_k$
for this length at level $k$ and $\mathsf{d}_i(Y)$ for that of a domain $Y$ at level $i$), $N$ and
$N_0$ the two parameters so named in the correlation theorem, the first of which is the subscript of
the stored domains $\mathfrak{U}_N$, $k_W$ the level attached to a frame $W$, and, in the sum of
condition (r1) below, $\mathcal{S}$ runs over the sets indexing Ba{\l}aban's boundary terms
$\mathbf{B}'^{(k)}(A,\mathcal{S})$, where $A$ is the label of the term. Further, $m$ and $m'$ denote
parameter tuples of the presentation (not the torus index), $\sigma$ the second argument of the
response $\mathsf{u}_p$, $V^{\mathrm{pl}}$ and $V^{\mathrm{rel}}$ the plain and relative parts of
the potential $V$ at $(\mathrm{T})$, and $V^{\sharp}$ the potential of the step into which terms are
resummed.

The hypotheses refer to the organisation of the correlation theorem's expansion at step $k$ and use
its terms by name, as follows. At the site $(\mathrm{L})$ the expansion passes through four
successive \emph{links}, the first of which is Ba{\l}aban's first operation there. The correlation
theorem proceeds through a numbered list of \emph{estimates}; those that take terms at a site as
inputs are its \emph{receiving estimates}, and their inputs are called \emph{input terms}. The units
presented at $(\mathrm{L})$ are divided by the correlation theorem into \emph{first-part} and
\emph{second-part} units; a first-part unit $v$ is \emph{interior} if its domain $X_v$ satisfies
$X_v\subset Z$, and units are divided into \emph{deep} and \emph{shallow} ones according to the
correlation theorem's division of units by depth. A term is \emph{in format} when it is presented in
the form in which the receiving estimates take their input terms; on the wrapped class this form is
read with a label in place of the core and the length $\mathsf{d}$ in place of $d$, as in condition
(r3) below. The remaining terms of that organisation used below, namely the seal
and the objects called sealed, the brackets of second-part units, the pile bounds, the iterations of
the correlation theorem, the carrier of a unit, the led rotations and the led-rotation condition on
the parameters of the step, and the strips and the corridors through them, are those of the
correlation theorem and are used here only by name.

The \emph{hypotheses of the holomorphic extension} are the five statements
$(\mathrm{HE}1)$--$(\mathrm{HE}5)$ below.

\medskip
\noindent\textup{(HE1)} \emph{The presented-difference bound.} Consider the site $(\mathrm{L})$ of
$\mathbf{R}_k$, and let the following be given:
\begin{itemize}
\item a component $C$ of $Z$;
\item a unit $\mathfrak{E}\in\mathcal{C}$, wrapped or plain, whose true domain $X$ satisfies
$X\subset\mathsf{C}_C$;
\item an index $\mathfrak{b}\in\{\mathfrak{b}_3,\mathfrak{b}_4\}$ of the presentation's chain;
\item families $\sigma_0^{(t)}$, $t=1,\dots,4$, of cells disjoint with $\Lambda$, that is,
$\sigma_0\cap\Lambda=\emptyset$; these are the cubes $\Delta$ disjoint with $\Lambda$ of
\cite{B-CMP122b};
\item parameter tuples $m=(\zeta,\varpi,\vec{\mathsf{s}})$ and $m'=(\zeta,\varpi,\vec{\mathsf{s}}')$
with $|\mathsf{s}^{(t)}(\Delta)|\le e^{\kappa_1}$;
\item an element $g'$ of the chart $\mathfrak{G}_{\varepsilon'}(X)$, with
$\varepsilon'=\varepsilon-\eta_0-\mu$, on which the chart proposition of the presentation delivers
its output; here $\eta_0$ is the radius decrement of that proposition, $\varepsilon$ is the radius at
which the rotated sup norm $\mathfrak{N}^{\mathbb{C}}_{\varepsilon}(\mathfrak{E})$ below is taken, and
$\mu=\mu_k$ the margin of Definition~\ref{4.22:def-rotation-radius}.
\end{itemize}
For all such data, writing $F^{\mathfrak{b}}=F^{\mathfrak{b}}_{\mathfrak{E}}$ for the presentation's
rotated unit of $\mathfrak{E}$ at the chain index $\mathfrak{b}$,
\begin{equation}\label{4.22:eq-extension-hypotheses-1}
\bigl|F^{\mathfrak{b}}(m;g')-F^{\mathfrak{b}}(m';g')\bigr|
\le C_{3F}\,e^{-\frac12\delta_P D_X}\,\mathfrak{N}^{\mathbb{C}}_{\varepsilon}(\mathfrak{E}).
\end{equation}
Consequently, on the whole chart, every piece $\mathfrak{p}$ whose families
$(\mathsf{y}_t)$ are not all empty obeys
\begin{equation}\label{4.22:eq-extension-hypotheses-2}
|\mathfrak{p}|\le C_{3F}\,e^{-\frac12\delta_P D_X}\,\mathfrak{N}^{\mathbb{C}}_{\varepsilon}(\mathfrak{E})\,
e^{-(\kappa_1-1)\sum_t|\mathsf{y}_t|}.
\end{equation}
Here $C_{3F}$ is a number fixed before $\gamma$. It depends on none of the following: $K$, $k$, the
age, the parameters $N$ and $N_0$, $\check{R}$, the older large-field regions,
$|X|$, the volume.

The bound is asserted for the unit read on its carrier, namely as
$e^{i\eta\mathbb{H}_4}e^{i\eta\mathbb{H}_3}U_2$ on $X$ itself, with $U_2$, $\eta$, $\mathbb{H}_3$ and
$\mathbb{H}_4$ as in the presentation. Ba{\l}aban's estimate in \cite{B-CMP122b}, of which
\eqref{4.22:eq-extension-hypotheses-1} is the counterpart, is not used in what follows.

The class $\mathcal{C}$ contains no kept member. Let $v\in\mathfrak{K}(Y)$, where $Y\subset C$ is an
older region and $X_v=Y^{\sharp}\subset\mathsf{C}_C$. For such $v$, the two displays
\eqref{4.22:eq-extension-hypotheses-1} and \eqref{4.22:eq-extension-hypotheses-2}, with
$\mathfrak{N}^{\mathbb{C}}_{\varepsilon}(\mathfrak{E})$ replaced by $\tfrac12\mathfrak{O}(v)$, are not
part of $(\mathrm{HE}1)$. They are supplied by clause (iv) of the kept-members lemma, which is a
consequence of clause (c) of the holomorphy lemma for kept members.

Two conventions are part of this hypothesis. First, at complex $\zeta$ the block-local map
$\mathfrak{u}$ is $G^{\mathbb{C}}$-valued. Objects conjugated by it are defined on the regular
representative and identified with the real objects on the real slice, and the two statements in
which such conjugated objects occur are read in this sense. Second, on the first three estimates of
the correlation theorem the norm $\mathfrak{N}$ does not depend on $t$, and the product bound
used there is the one of the pointwise proposition of the presented-difference bound.

\medskip
\noindent\textup{(HE2)} \emph{The led containment of the presented difference.} Consider the disc
$|\lambda|\le\rho_X$ of the one complex interpolation parameter $\lambda$ through which the
presented-difference bound $(\mathrm{HE}1)$ is proved. Along it only the difference of the presented
triples at $m$ and $m'$ is dilated. We write $P_\lambda$ for the interpolated presented datum at
$\lambda$ and $Z^{\mathrm{dil}}(s)$ for the generator at the law site $s$ of the dilated difference;
both are the objects of the containment and radius lemmas attached to the presented-difference
bound, and the superscript distinguishes the latter from the large-field region $Z$.

On this disc, $P_\lambda$ lies in the stored domain of $\mathfrak{E}$, the frame condition of that
domain included: each entry is at most $3/32+\frac12$ of its threshold, and
$3/32+\frac12<1-\varpi_{\mathrm{fr}}$. Moreover, at every law site $s\in\mathfrak{s}(X)$, pointwise
or belonging to a frame of any age, and under the led-rotation condition on the parameters of the
step,
\begin{equation}\label{4.22:eq-extension-hypotheses-3}
\begin{split}
|\lambda|\,|Z^{\mathrm{dil}}(s)|
&\le\vartheta^{3F}(s)
:=K_u\,\alpha_{*,k}\,e^{-\frac12\delta_P\mathsf{w}\check{R}_k}\,
e^{-2(\mathsf{a}_k(s)-1)_+}\\
&\le\tfrac12\,\mu_k\,\mathsf{e}(s)/\varepsilon_{\star},
\end{split}
\end{equation}
where $\mathsf{w}$ is the width factor of the radius lemma attached to the presented-difference
bound. Here $\mathsf{a}_k(s)=(h-k_W)_+$ at frame sites, with $h$ and $k_W$ the two levels of the
correlation theorem's age count at the frame $W$ containing $s$, and the age factor is $1$ at
pointwise sites. This hypothesis is an input of both conclusions of the holomorphic extension.

Let $v\in\mathfrak{K}(Y)$ with $Y\subset C$ an older region. For such $v$, the containment is
supplied by clause (c) of the holomorphy lemma for kept members. The site bound
\eqref{4.22:eq-extension-hypotheses-3} at the law sites $s\in\Sigma_Y$ of $Y$ is the one of the
radius lemma attached to the presented-difference bound, whose proof does not involve $\mathfrak{E}$.

\medskip
\noindent\textup{(HE3)} \emph{The led containment of the correlation theorem's response.} This
hypothesis concerns the disc
\begin{equation*}
|\mathsf{t}|\le\rho_p=\mathsf{w}_{\mathrm{site}}\,(C_b\eta_p)^{-1/2}
\end{equation*}
of the response, at $(\mathrm{L})$ and at $(\mathrm{T})$, where $\eta_p$ is the constant of the
response disc that the correlation theorem fixes. It applies to every unit $v$ to which the
fourth estimate of the correlation theorem applies and on which that theorem runs this disc, namely
to
\begin{itemize}
\item the plain second-part units (wrapped second-part units are treated by condition (r2) of
$(\mathrm{HE}5)$ and by clause (iii) of the correlation theorem's lemma treating them);
\item the plain first-part units with $K_v\setminus Z\ne\emptyset$, where $K_v$ is the set that the
correlation theorem attaches to $v$, in the fourth iteration only.
\end{itemize}
For every such $v$, under the led-rotation condition on the parameters of the step taken at the full
margin $\mu_k$, the following hold.
\begin{itemize}
\item The response set $\mathsf{R}_p^h$ is contained in the stored domain of $v$. This is the
content of the half-rate membership lemma and, outside the correlation theorem's window set, of the
correlation theorem's half-rate membership and arrest conditions together with the two conditions
supplied by the seal.
\item At every law site $s$, with $D(s)$ the depth of $s$ and $\mathsf{a}_{\mathrm{site}}$ the site
factor that the correlation theorem attaches to $s$,
\begin{equation}\label{4.22:eq-extension-hypotheses-4}
\begin{split}
\bigl|\log\bigl[\mathsf{u}_p(0,\sigma)^{-1}\mathsf{u}_p(\mathsf{t},\sigma)\bigr](s)\bigr|
&\le 2.1\,K_u\,\mathsf{a}_{\mathrm{site}}\,e^{-\frac12\delta_{\natural}(D(s)-1)}\\
&\le\mu_k\,\mathsf{e}(s)/\varepsilon_{\star}.
\end{split}
\end{equation}
\end{itemize}
This hypothesis is an input of both conclusions of the holomorphic extension. It is used at the
units leaving $Z$ and at the brackets of the plain second-part units.

\medskip
\noindent\textup{(HE4)} \emph{The inputs and floors of the correlation theorem.} These are the
following five groups of statements.
\begin{enumerate}
\item[(i)] The supply statement for the fourth estimate.
\item[(ii)] The expansion rule at $(\mathrm{L})$. Every interior first-part unit, plain or wrapped,
is expanded by the two expansion identities of the correlation
theorem at $(\mathrm{L})$, the first of which is the identity of the first link, with neither the
partition of unity $\zeta_{\square}$ nor the disc variable
$\mathsf{t}$. The disc variable $\mathsf{t}$ of the correlation theorem (in the fourth iteration
only) is used on the plain first-part units with $K_v\setminus Z\ne\emptyset$.

In its sourced form the rule applies also to one further family of units of the correlation
theorem, with the estimate that the correlation theorem attaches to that family, and it is taken in
its pointwise form together with a floor condition.

The domain list of the sourced rule and of one further statement of the correlation theorem includes
the kept members
$\mathfrak{K}(Y)$ of live regions. They enter as bearers of clauses in the index family
$\mathfrak{I}_n$ of the correlation theorem, and not as terms of the first-stage expansion. The
sourced clause is thereby declared for them.
\item[(iii)] The containment statements at the families so run:
\begin{itemize}
\item at $(\mathrm{T})$, the homogeneity containment;
\item at $(\mathrm{P})$, two containment statements of the large-field stages, one of them in
localised form;
\item at $(\mathrm{L})$, the key sourced containment together with the inheritance and arrest
statements, and the sourced containment lemma at $(\mathrm{L})$, in which $\sigma_0$ is a free
variable ranging over families of cells of the sealed base with $\sigma_0\cap\Lambda=\emptyset$.
\end{itemize}
\item[(iv)] The sums of the large-field stages, with the window-block bound at the window blocks.
\item[(v)] Clauses (c), (d) and (e) of the resummation proposition of the correlation theorem for the
plain family, together with the following:
\begin{itemize}
\item the bracket containment;
\item the per-piece bound, which for each piece is supplied by $(\mathrm{HE}1)$;
\item the summed bound over the plain interior units. For deep units it holds by the proof of the
deep-unit lemma together with the correlation theorem's pile bound. For shallow units it holds by
counting the cells of $C\setminus\mathsf{C}_C$. It is subject to its own ceiling.
\end{itemize}
\end{enumerate}

\medskip
\noindent\textup{(HE5)} \emph{The complex relative bound for the boundary terms.} This is the
relative bound for Ba{\l}aban's boundary terms $\mathbf{B}'^{(k)}$ of \cite{B-CMP122b}, in the norm
$\mathfrak{N}^{\mathbb{C}}$, with every unit read on its carrier as in $(\mathrm{HE}1)$. Its
hypotheses include the following:
\begin{itemize}
\item the per-piece bound of $(\mathrm{HE}4)$(v) at plain units;
\item a first companion bound;
\item the second companion bound, at $(\mathrm{P})$, in the weight of the stage's boundary-term
bounds.
\end{itemize}
It is required on the wrapped class $\mathfrak{S}$, for terms in format, in the form of the package
of conditions
\begin{equation}\label{4.22:eq-extension-hypotheses-a}
\text{(r0)--(r4) and (fam), with (M2) contained in (r2) and (maj) in (r3),}
\end{equation}
which read as follows.
\begin{enumerate}
\item[(r0)] \emph{Parametric form.} Every estimate that the statement supplying $(\mathrm{HE}5)$ runs
on a unit consists of two parts:
\begin{itemize}
\item an identity valid at each value of the site's parameters and of the chart variable $g'$;
\item a majorant that reads the unit through one supremum of one holomorphic function over the
product of datum, parameter polydisc and chart.
\end{itemize}
Hence every bound of the package holds with $\mathfrak{N}$ replaced by
$\mathfrak{N}^{\mathbb{C}}_{\varepsilon}$ and $\mathfrak{N}^{A}$ by $\mathfrak{N}^{A,\mathbb{C}}$,
where $\mathfrak{N}^{A}$ is the companion norm of the correlation theorem and
$\mathfrak{N}^{A,\mathbb{C}}$ its complex counterpart.
\item[(r1)] \emph{Production.} Take \cite[(1.99)]{B-CMP122b} in the form of the correlation theorem,
that is, as clause (e) of the resummation proposition of the correlation theorem, which follows from
the combinatorial resummation lemma.
It implies, at the step at which the term is created, with the weights on labels of clauses (a) and
(c) of the proposition that fixes them,
\begin{equation}\label{4.22:eq-extension-hypotheses-5}
\begin{split}
\sum_{\mathcal{S}}e^{\kappa_A(k)|\mathcal{S}|}\,
\mathfrak{N}^{\mathbb{C}}\bigl(\mathbf{B}'^{(k)}(A,\mathcal{S})\bigr)
&\le e^{\lambda_\theta}\,C_{99}\,c_1^{\flat}(g_k)\,
e^{-(1+\frac14\beta_s)\kappa\mathsf{d}_k}\\
&\le B_0\,e^{-\kappa\mathsf{d}_k},
\end{split}
\end{equation}
where $c_1^{\flat}(\cdot)$ is the function of the coupling and $\lambda_\theta$, $C_{99}$, $\beta_s$
and $B_0$ are the constants of that form.
\item[(r2)] \emph{Sites: identity $(\mathrm{M}1)$ and containment $(\mathrm{M}2)$.} On a wrapped
unit each site runs its first-stage identity at amplitude one, as follows.
\begin{itemize}
\item At $(\mathrm{T})$: the fundamental theorem of calculus in $\mathsf{s}$, on $\sigma_0$ equal to
the set of all cells.
\item At the first link of $(\mathrm{L})$: the fundamental theorem of calculus in $\mathsf{s}$ alone,
on the cells disjoint with the enlargement $\Lambda^{\sim}$ of $\Lambda$.
\item At links two to four: Ba{\l}aban's own $\sigma_0$, that is, the cells disjoint with $\Lambda$.
\item At the links of $(\mathrm{L})$: $|\mathsf{s}(\Delta)|\le e^{\kappa_1}$, and the disc variable
$\mathsf{t}$ of the correlation theorem is not used.
\item At $(\mathrm{P})$: the stage's identity is taken unchanged, on the sealed half-rate polydisc.
\item The pieces of a deep unit are bounded on the disc $|\lambda|\le\rho_X$ of $(\mathrm{HE}2)$, at
links one to four.
\end{itemize}
Moreover:

$(\mathrm{M}2)$: the response maps the whole product of parameter polydisc and tube into the
stored domain $\mathfrak{U}_N(X_F)$, its frame condition included. Every later membership statement
must be proved for this smaller set.

The containment $(\mathrm{M}2)$ follows from the correlation theorem's estimates at Ba{\l}aban's
$\sigma_0$, the cell lemma, and one further bound of the correlation theorem's list, taken at
amplitude one.
\item[(r3)] \emph{Receiving estimates, closure and majorants.} The receiving estimates at
$(\mathrm{T})$,
$(\mathrm{P})$ and $(\mathrm{L})$ accept input terms in format, with $d$ replaced by $\mathsf{d}$ and
the core by the label. They return every output of $\mathfrak{S}$ bounded in $\mathsf{d}$ at the
estimate's rate and constant, and the anchored totals of each estimate are included and taken in
$\mathsf{d}$. These totals are:
\begin{itemize}
\item the total, over the domains $Y$ containing a bond $b$, of the bounds $\nu(Y)$ of the
corresponding estimate, $\sum_{Y\ni b}\nu(Y)\le C_\kappa\varepsilon_V$, with the constants
$C_\kappa$ and $\varepsilon_V$ of that estimate;
\item its analogue at $(\mathrm{P})$;
\item the pile bound of the weak-coordinate lemma;
\item the pile bound of the estimate at $(\mathrm{P})$;
\item the dichotomy on which the complex-potential corollary rests.
\end{itemize}
The decay of the factor $S_v$ of a unit $v$ and the walk factors are never halved. The corridors
through strips are paid by the $\kappa_1$-weight of the term, with constant $c_l$.

$(\mathrm{maj})$:
\begin{itemize}
\item The first-stage pieces are linear in the input term and in the parameters passed with it.
\item The index sets $\mathfrak{Y}(F)$ of a term $F$, the labels, the cores and the counts read the
scaffold, and the place where the input term reads the response, never the response's values.
\item The second stage is a majorant series in those suprema.
\end{itemize}
At $(\mathrm{T})$, $V:=V^{\mathrm{pl}}+V^{\mathrm{rel}}$, where, under the relative ceiling at
$(\mathrm{T})$, with $\varepsilon^{\mathrm{rel}}_i$ and $\kappa_V$ the constants of that ceiling, for
a domain $Y$ at level $i$,
\begin{equation}\label{4.22:eq-extension-hypotheses-6}
|V^{\mathrm{rel}}(Y)|\le\varepsilon^{\mathrm{rel}}_i\,e^{-\kappa_V\mathsf{d}_i(Y)}.
\end{equation}
\item[(r4)] \emph{End of life.} A term never outlives the large-field regions it is attached to. At
their treatment it is a first-part unit at the site $(\mathrm{L})$ of the step that treats those
regions: a shallow one by the
shallow-unit lemma, a deep one by the deep-unit lemma, which rests on $(\mathrm{HE}1)$. It is then
resummed into $V^{\sharp}$.
\item[(fam)] On the families on which the sourced clause runs, the base-point profile of the
profile lemma holds at every link when expressed in the units $\mathsf{e}(s)$, the margin $\mu_k$
and the reserve of Definition~\ref{4.22:def-radius-budget}. This condition is proved in this
chapter as a separate statement, whose inputs are clauses (iii) and (iii)$'$ of the cell lemma and
one further condition. It does not distinguish between the
plain and the wrapped class, so it applies to the interior plain units of $(\mathrm{HE}4)$ as
well.
\end{enumerate}
\end{definition}

\begin{definition}[The source-tube hypothesis on the kept quadratic forms]\label{4.22:def-source-tube}
Let $Y_i$ be a second-class component in the sense of
Definition~\ref{4.22:def-cores-wrapped-class}. Let $k_W$ be the step that created it, and write
$C:=Y_i$. Let $Y_i^{\sharp}$ be its strip. Let $\mathfrak{K}(Y_i)$ be the pair of kept members of
that definition: the term $V(Y_i^{\sim})$ and the quadratic form $\mathsf{Q}_i$, both kept under
$\mathbf{T}'_{k_W}(Y_i)$ in the representation \cite[(1.102)]{B-CMP122b}. In that representation
$Y_i$ is a frame with creating step $k_W$. We write $\Sigma_{Y_i}$ for its law sites. We write
$\mathfrak{G}_{\varepsilon}(\Sigma_{Y_i})$ for the chart of $G^{\mathbb{C}}$-valued maps $g$ with
$\operatorname{pol}(g)<\varepsilon$ at the sites of $\Sigma_{Y_i}$.

\emph{The tube and the source tube.} The tube on which a kept member $v\in\mathfrak{K}(Y_i)$ is
holomorphic is
\begin{equation*}
\mathfrak{U}^{\mathrm{tr}}_v=\tilde{U}^{c}_{k_W}(Y_i^{\sharp},\tilde\alpha_0,\tilde\alpha_1).
\end{equation*}
This is the complex configuration space of level $k_W$ with the radii
$\tilde\alpha_0,\tilde\alpha_1$, with its clauses restricted to the strip $Y_i^{\sharp}$.
The \emph{source tube} $\mathfrak{S}_{\mathrm{src}}:=\mathfrak{S}_{\mathrm{src}}(Y_i^{\sharp})$ is
the sourced set of pairs $(\mathbf{U},\mathbf{J})$ over $Y_i^{\sharp}$. It contains
$\mathfrak{U}^{\mathrm{tr}}_v$. Conditions (F), (Q) and (V) below are required on this set. They
are required uniformly in the further argument $\boldsymbol{\Phi}$, which takes real values on
the supports.

The \emph{source-tube conditions} for $Y_i$ are the following conditions (F), (Q), (V),
(Z) and (U).
\begin{itemize}
\item[(F)] Write $\mathfrak{q}_C:=\mathsf{Q}_i$. This form is the localisation to $C$, in
Ba{\l}aban's sense, of the quadratic form of \cite[(1.71)]{B-CMP122b}. It is built from the
$C$-localised minimiser $U^{0}_{k_W,C}$ on $\mathbf{B}''_{k_W}(C)$, taken under the hypothesis on
the treated large-field region that defines that minimiser. It reads
$(\mathbf{U},\mathbf{J})$ only through the $k_W$-fold block average of $\mathbf{U}$
(Definition~\ref{1.5:def-block-average}) restricted to $C\setminus\Lambda$, where $\Lambda$
denotes the small-field region of step $k_W$, and it does
not depend on $\mathbf{J}$. The first-order term of the $C$-localised minimiser, its derivative at
zero, is $H''_{1,k_W,C}$, which is Ba{\l}aban's $H_1$ of
\cite[(3.129)]{B-CMP99b}, \cite[(2.7)--(2.8)]{B-CMP109} and \cite[(1.20)]{B-CMP122b}; it enters
$\mathfrak{q}_C$ through the expression displayed in (Q).
Two bounds are required:
\begin{itemize}
\item[(F1)] the deviation of $\mathfrak{q}_C$ from its value at the base point $(U,0)$ of an
admissible decomposition (see (U)(a)) is bounded,
\begin{equation*}
\sup_{(\mathbf{U},\mathbf{J})\in\mathfrak{S}_{\mathrm{src}}}
\bigl|\mathfrak{q}_C(\mathbf{U},\mathbf{J})-\mathfrak{q}_C(U,0)\bigr|
\le\Sigma^{\mathfrak{q}}(x_{k_W}),
\end{equation*}
where $\Sigma^{\mathfrak{q}}$ is one non-increasing function of the inverse coupling $x$, the
same for all components and steps, with $\Sigma^{\mathfrak{q}}(x)\to0$ as $x\to\infty$; it is
evaluated at $x_{k_W}=g_{k_W}^{-2}$;
\item[(F2)] $\sup_{\mathfrak{S}_{\mathrm{src}}}|\mathfrak{q}_C|\le\mathfrak{q}_{\max}(k_W)$.
\end{itemize}
\item[(Q)] The form $\mathfrak{q}_C(\mathbf{U},\mathbf{J};\varpi)$, with $\varpi$ the parameter of
the correlation theorem, has the following properties:
\begin{itemize}
\item it does not depend on the splitting $\mathbf{U}=U'U$;
\item it is holomorphic on $\mathfrak{S}_{\mathrm{src}}(Y_i^{\sharp})$;
\item it reads $(\mathbf{U},\mathbf{J})$ only through the $k_W$-fold block average of
$\mathbf{U}$ restricted to $C\setminus\Lambda$;
\item it is invariant under gauge transformations acting jointly on its arguments, in the same
sense in which the other kept member is gauge invariant;
\item at real arguments it equals the form of \cite[(1.71)]{B-CMP122b},
\begin{equation*}
-\tfrac12\bigl\langle DH''_{1,k_W,C}B,\;\hat\chi\, DH''_{1,k_W,C}B\bigr\rangle ,
\end{equation*}
with $B$ and $\hat\chi$ as there.
\end{itemize}
\item[(V)] Write $V^{\sharp}(C^{\sharp};\cdot)$ for the orbit datum of $V(Y_i^{\sim})$ on the
source tube, with $C^{\sharp}=Y_i^{\sharp}$. Then, uniformly in $\boldsymbol{\Phi}$,
\begin{align*}
\sup_{(\mathbf{U},\mathbf{J}),(\mathbf{U}',\mathbf{J}')\in\mathfrak{S}_{\mathrm{src}}}
\bigl|V^{\sharp}(C^{\sharp};\mathbf{U},\mathbf{J})-V^{\sharp}(C^{\sharp};\mathbf{U}',\mathbf{J}')\bigr|
&\le\Sigma^{V}(k_W),\\
\Sigma^{\mathrm{tot}}(k_W):=\Sigma^{V}(k_W)+2\mathfrak{q}_{\max}(k_W)
&\le C_\Sigma\,\check{R}_{k_W}^{\,p_\Sigma}.
\end{align*}
The
first bound is a two-point oscillation over the whole source tube. The number $C_\Sigma$ is fixed
before $\gamma$ in the order of choice, and it does not depend on $\mathfrak{n}^{\mathrm{kept}}$.
No right-hand side depends on any of the following:
\begin{itemize}
\item the quantity $\sum_j|\Gamma^0_j\cap Y_i|$, where $\Gamma^0_j$ is the large-field set of
level $j$ in Ba{\l}aban's construction;
\item the volume of the core;
\item the nesting depth;
\item an age;
\item $K$ or $k$;
\item the number, the sizes or the histories of the second-class components and of the arrests
inside $Y_i$.
\end{itemize}
Condition (V) does not require that $\Sigma^V$ tend to zero. Its bound involves
$\check{R}_{k_W}$, the scale of the step $k_W$. At every later site this bound is paid through the
collar of the second-class component $Y_i$ itself, counted in cells of its own scale $k_W$ and in
the distance $d^{\wedge}$, as prescribed in (K4), and never through the reference estimate of the
construction.
\item[(Z)] The three-family lemma in the $(\mathbf{U},\mathbf{J})$-directions, with all five of its
clauses, holds at the creating step $k_W$.
\item[(U)] (a) The supremum in (F1) runs over every admissible decomposition
\begin{equation*}
\mathbf{U}=U'U=e^{i\eta A'}U
\end{equation*}
that the clauses of the space \cite[(2.34)--(2.39)]{B-CMP119} admit, with $\eta$ and $A'$ as
there. In such a decomposition $U$
has values in the group $G$, and $A'$ is $\mathfrak{g}^{\mathbb{C}}$-valued, real values
included. The point $(U,0)$ exists only relative to a decomposition, so this is the one
well-posed reading of (F1). (b), (c): these are the two clauses of (K3) below.
\end{itemize}

\emph{Kept-member clauses.} For each $v\in\mathfrak{K}(Y_i)$ we impose the following.
\begin{itemize}
\item[(K1)] The member $v$ is stored with region equal to its true domain,
$X_v=Y_i^{\sharp}$, and with norm $\mathfrak{O}(v)$. At a site of level $l$ it is read with the
one-cube label of level $l$ attached to $Y_i$, whose length $d_l$, in the sense of
Definition~\ref{4.22:def-cores-wrapped-class}, is $0$. The level
$l$ is fixed as follows at the stage of index $h$ of a chain of large-field operations carried
out during step $k$:
\begin{itemize}
\item $l=h+1$ on the chain belonging to the component of $Y_i$;
\item $l=k$ on a chain of a foreign component.
\end{itemize}
In particular, $l$ is never the level $k$ of the step on the chain belonging to the component of
$Y_i$.
\item[(K2)] Every later site that moves the background on $Y_i^{\sharp}$ maps the product of its
own parameter polydisc and its own configuration tube holomorphically into
$\mathfrak{U}_v(Y_i^{\sharp})\times\mathfrak{G}_{\varepsilon}(\Sigma_{Y_i})$, with $\varepsilon$
the chart radius in force at that site. Here
$\mathfrak{U}_v$ is the stored domain of $v$, with its reserve. The sites concerned are:
\begin{itemize}
\item the site $(\mathrm{T})$ of each step;
\item each stage $(\mathrm{P})$ of a chain, on the component of $Y_i$ or on a foreign one,
running or interrupted;
\item each link of $(\mathrm{L})$, foreign or the treatment of the component of $Y_i$ itself.
\end{itemize}
\item[(K3)] This clause has two parts, lettered (b) and (c) so as to continue (U)(a).
(b) The membership in $\mathfrak{U}_v(Y_i^{\sharp})$ asserted in (K2) is witnessed by
the decomposition $(U_{\mathbf{K}},\{u_{\square}\})$ of the base $\mathbf{K}$ of the family at the
site. This decomposition does not depend on the family parameter $\tau$, and the reserve
$\varpi_{\mathrm{fr}}r_a$, with $r_a$ the radius attached to the base decomposition at the site, is
counted in it.
(c) At the healing of the component of $Y_i$ itself, the membership holds for the pairs
$P_\lambda$ of the three-family lemma with $|\lambda|\le\rho_X$, on the $\lambda$-disc of that
lemma.
\item[(K4)] The estimates at the sites $(\mathrm{T})$, $(\mathrm{P})$ and $(\mathrm{L})$ accept
the datum $(\emptyset,\{(Y_i,k_W)\})$ of $v$, with the oscillation norm $\mathfrak{O}(v)$ in place
of the sup norm $\mathfrak{N}$. They label the pieces of $v$ by $\operatorname{lab}(\mathsf{y})$.
For each second-class component they carry two additional units in each of the following
estimates, totals included:
\begin{itemize}
\item the summation estimate;
\item the two anchoring estimates, at the steps and at the large-field stages;
\item the estimate marked with a star among the site estimates;
\item the estimate for deep units.
\end{itemize}
In these estimates the share of the far factors is made explicit in cells and in the distance
$d^{\wedge}$.
\end{itemize}
The \emph{kept-member clauses} are the conjunction
\begin{equation}\label{4.22:eq-source-tube-a}
(\mathrm{K}1)\wedge(\mathrm{K}2)\wedge(\mathrm{K}3)\wedge(\mathrm{K}4).
\end{equation}
The \emph{source-tube hypothesis} for $Y_i$ is the conjunction of the source-tube conditions
(F), (Q), (V), (Z), (U) with the kept-member clauses \eqref{4.22:eq-source-tube-a}.
\end{definition}

For real configurations, \cite[Proposition~1]{B-CMP122b} is the first of the statements from
which conditions (F) and (Q) are derived. The estimate
\cite[(1.41)]{B-CMP122b} is used with its variable $B'$ restricted to $|B'|_\infty$ below a radius
that is fixed after $M$ and before $\gamma$ in the order of choice.

The conditions enter the construction through two implications, both parts of the oscillation
lemma for kept members.

First, (V), (Q), (F2) and (K2) together imply the following bound for both members
$v\in\mathfrak{K}(Y_i)$. Writing $v^{\sharp}$ for the orbit datum of $v$ and
$\operatorname{osc}_{\mathfrak{P}}(v^{\sharp})$ for its oscillation over the site's parameter
polydisc $\mathfrak{P}$, we have, on the orbit chart $\mathfrak{G}_{\varepsilon}(\Sigma_{Y_i})$ of
(K2), at every site later than $k_W$,
\begin{equation*}
\operatorname{osc}_{\mathfrak{P}}(v^{\sharp})\le\mathfrak{O}(v)\le C_\Sigma\,\check{R}_{k_W}^{\,p_\Sigma}.
\end{equation*}

Second, (F), clauses (a) and (b) of (U) and the bound on the logarithm used in the proof of that
lemma together imply the following bound at a site of step $k$:
\begin{equation*}
\operatorname{osc}_{\mathfrak{P}}(\mathsf{Q}_i^{\sharp})\le2\Sigma^{\mathfrak{q}}(x_{k_W})
\le2\Sigma^{\mathfrak{q}}(x_k-c_5).
\end{equation*}

The kept-member clauses (K1)--(K4) are established in the proof of the holomorphic extension. They
are derived from the containment statements of the fourth and the lemmas, taken by statement, of
the fifth of the hypotheses of Definition~\ref{4.22:def-extension-hypotheses}, from the first two
of those hypotheses through their proofs, from (V), (Q) and (F2), and from one further declaration
in the list of domains of those hypotheses. They
therefore lengthen the list formed by those hypotheses and the source-tube conditions by that one
declaration only. In substance, the last hypothesis of the holomorphic extension consists of the
source-tube conditions.

\begin{remark}[The decay lemma covers all four families]\label{4.22:rem-four-families}
At the presentation site $(\mathrm{L})$ of the large-field operation $\mathbf{R}_k$, every interior
unit of the first part of the expansion at this site is presented with the index
$\mathfrak{b}=\mathfrak{b}_4$ and four families
$(\mathsf{y}_1,\mathsf{y}_2,\mathsf{y}_3,\mathsf{y}_4)$. Here $\mathfrak{b}_4$ is a member of the
chain $\mathfrak{b}_0,\dots,\mathfrak{b}_4$ of indices of the presentations $F^{\mathfrak{b}}$. It
is not the small-field threshold $\mathfrak{b}_k$. The presentation proceeds in four links, the
first of which corresponds to Ba{\l}aban's first operation, and the $i$-th link contributes the
family $\mathsf{y}_i$. This holds whether the unit is plain or wrapped, a wrapped unit being one
touched by a second-class outcome.

The decay lemma, hypothesis \textup{(H1)} of Definition~\ref{4.22:def-extension-hypotheses},
applies to every piece whose tuple of families is not all empty, that is, to the first, second,
third and fourth families alike. In particular it applies to the pieces with
\begin{equation*}
\mathsf{y}_1\neq\emptyset,\qquad \mathsf{y}_2=\mathsf{y}_3=\mathsf{y}_4=\emptyset .
\end{equation*}
Such pieces occur. Ba{\l}aban's bound on the pieces of his expansion in \cite{B-CMP122b}, the
bound in whose place the decay lemma is used and which is not itself used in this book, addresses
the pieces to which at least one of the second, third and fourth links contributes; the pieces
displayed above are not among them. The scope is not restricted to the second, third and fourth
families by hypothesis
\textup{(H1)} or hypothesis \textup{(H2)} of Definition~\ref{4.22:def-extension-hypotheses}. Nor is
it restricted by the site identities for wrapped units, or by the displayed gloss on these
hypotheses. In this section a disc in the variable $\mathsf{t}$ always means the disc of
the correlation theorem. Presenting a unit with four families, the first included, changes neither
the number of presentations of the unit nor the margins $\mu_k$.

None of the following is claimed or used:
\begin{enumerate}
\item[(i)] any bound expressed through Ba{\l}aban's size function of the true domain $X_F$ of a
wrapped term $F$;
\item[(ii)] any estimate of the kind established in the correlation theorem, on Ba{\l}aban's
family $\sigma_0$ of cells disjoint with $\Lambda$, other than those established there;
\item[(iii)] a $\mathsf{t}$-disc on a wrapped unit or on an interior unit;
\item[(iv)] a parameter on a cell meeting the interior of the small-field region $\Lambda$ of the
step in a presented chain.
\end{enumerate}
\end{remark}

\begin{proof}
At the first link of the presentation site, Ba{\l}aban's amplitude $t_{\square}$ at the first
operation, \cite{B-CMP122b}, is not taken. Consequently the expansion produces pieces whose first
family $\mathsf{y}_1$ is non-empty while $\mathsf{y}_2$, $\mathsf{y}_3$ and $\mathsf{y}_4$ are
empty. Write $Y'_2$, $Y'_3$, $Y'_4$ for the sets in Ba{\l}aban's notation for the bound recalled
above that correspond to the second, third and fourth links. In that notation these pieces have
\begin{equation*}
Y'_2\cup Y'_3\cup Y'_4=\emptyset ,
\end{equation*}
whereas that bound addresses the pieces for which this union is non-empty.

The conclusion of the decay lemma bounds every piece whose tuple
$(\mathsf{y}_1,\dots,\mathsf{y}_4)$ is not all empty, for plain and wrapped units alike. No index
$i$ is singled out. A tuple with only $\mathsf{y}_1$ non-empty, like every other non-empty tuple,
is not all empty; hence the pieces of the first paragraph are among those the decay lemma bounds,
and the lemma covers the first, second, third and fourth families.

Item (iv) is not needed, for the following reason. In the proof of the decay lemma the comparator
is the auxiliary zero field. This field is never presented, never summed, never placed on a rotation
chart, and never made the variable of a Cauchy integral. Hence no parameter is attached to a cell
meeting the interior of $\Lambda$ in a presented chain. Items (i)--(iii) do not enter the argument
above.
\end{proof}

\begin{proposition}[Holomorphic extension of stored terms in the rotation variables]
\label{4.22:prop-holomorphic-extension}
Let the following standing statements be given, each taken as the statement it is.
\begin{itemize}
\item The statements and floors collected in the closing section of the chapter of the
correlation theorem which the auxiliary lemmas of this section invoke by name, together with
the statements that the correlation theorem itself takes by statement.
\item The hypotheses of the inheritance theorem at the parameter pair
$(\mathsf{w}_0,\tfrac12)$, with layer number $\theta_S\in[\tfrac12,\tfrac89]$, including, by
name, its near statements for frames and for cores.
\item The hypotheses and floors of the presentation statement at the site $(\mathrm{L})$, whose
clauses are read on the background pairs $P=(\mathbb{U},\mathbb{J})$ on which the consumers of
its presented units read them.
\item The arrest statement, together with its initial case.
\item The two declared choices: the invariant smooth cut-off and its holomorphic lift
\eqref{4.22:eq-invariant-cutoff-a}, and the choice of the stored frame domains.
\item The conditions on the live radii, namely conditions (i)--(x), with condition (iii) in
its reserve form read separately for each class of terms, condition (vi) with the constant
$8e^{2\varepsilon_\star}$, and the primed forms (vi)$'$, (viii)$'$, (ix)$'$ together with (x);
the condition on the constants of the source tube; the condition on the frame constants; the
condition \eqref{4.22:eq-weak-collar-a} on the arrested collar variable; and the regularity
ceiling. Each of these is a single inequality in one real variable $x$ (an inverse squared
coupling), required for all $x\ge\gamma_H^{-2}$, and refers to no object of a particular
renormalisation run.
\end{itemize}
For each of the two classes of stored terms of Definition~\ref{4.22:def-cores-wrapped-class},
the plain class $\mathfrak{A}$ and the wrapped lineage $\mathfrak{S}$, the \emph{extension
statement for that class} consists of clauses (a), (b), (c) and (g) below, together with
clause (iv) of the anchored-sum lemma at the site $(\mathrm{T})$ and clause (iv) of its
counterpart at the site $(\mathrm{P})$, all asserted for the terms of that class. In clause (b)
for a class, the norm inequality is asserted for the sub-family of the terms of that class,
with the anchored totals entering the norm restricted to that class; such a restricted total is
a sub-sum of a sum of non-negative terms, and a cluster of the plain class contains elements of
the plain class only.

Then, for every $K$ and every $k\le K$, the following hold.
\begin{itemize}
\item[(A)] From the standing statements alone: clauses (d)--(f), for every term that has an
orbit datum; clause (h), in the following reading: the data supplied by the weak-collar lemma,
by clause (ii) of the frame lemma and by clause (i) of the entry statement follow from the
standing statements, with two exceptions. The first exception is clause (ii) of the frame lemma
for the kept quadratic form $\mathsf{Q}_i$: the holomorphy, the bound and the joint invariance
of $\mathsf{Q}_i$ on the tube $\mathfrak{U}^{\mathrm{tr}}_v$ are the two properties of
$\mathsf{Q}_i$ required in the source-tube hypothesis
(Definition~\ref{4.22:def-source-tube}), and they rest on a list of statements headed by the
real form of a proposition printed with its proof in \cite{B-CMP122b}. The second exception is
the bounds of the oscillation lemma for kept members and the bounds of the clauses
\eqref{4.22:eq-source-tube-a}: the latter are clauses of the source-tube hypothesis
(Definition~\ref{4.22:def-source-tube}), and the former rest on that hypothesis.
Also from the standing statements alone follow clause (i) below; each of the following
auxiliary statements, as stated: the polar decomposition lemma
(Lemma~\ref{4.22:lem-polar-decomposition}), the uniqueness, equivariance and Cauchy
properties of orbit data \eqref{4.22:eq-orbit-datum-properties-a}, the complexified lemma at
the site $(\mathrm{P})$, the margin lemma at the site $(\mathrm{T})$, the Cauchy-disc lemma, the
further variant of the margin lemma, the transfer lemma from the sup
norm to the rotated sup norm with its form for the oscillation norm, the shared-margin lemma,
the base-point lemma, the weak-collar lemma, clauses (i)--(iii) of the anchored-sum lemma at
the site $(\mathrm{T})$ and of its counterpart at the site $(\mathrm{P})$, the frame lemma, the
oscillation lemma for kept members, the holomorphy lemma for kept members, the entry
statement, the orbit lemma, the Cauchy corollary for presentations, clause (a) of the
output-chart proposition, and the law-site rule; and clause (j), which on the families of the
relative bound for the boundary terms rests on the cell-family lemma of that bound.
\item[(B)] Assume
\begin{itemize}
\item the containment and led-rotation bound on the discs of the correlation theorem: on each
such disc, for every unit to which the correlation theorem is applied on that disc, the response
lies in the stored domain of the unit, and at every law site $s$ the rotation led by the disc
is at most $\mu_k\mathsf{e}(s)/\varepsilon_\star$;
\item the containment and led-rotation bound on the disc of the difference lemma: on the disc of
the parameter $\lambda$ of the difference lemma at the site $(\mathrm{L})$, the displaced pair
lies in the stored domain of the unit, the choice of the stored frame domains included, and at
every law site $s$, pointwise or of a frame of any age, the rotation led is at most
$\tfrac12\mu_k\mathsf{e}(s)/\varepsilon_\star$;
\item the difference lemma at the site $(\mathrm{L})$, which bounds the change of a presented
unit of $\mathcal{C}$ whose true domain lies in a core, under a change of the parameters of that
site, by a factor decaying exponentially in the depth $D_X$ times the rotated sup norm of the
unit; and
\item the inputs and floors of the correlation theorem on the families it runs: the supply
statement for the units of its fourth family, the expansion of the interior first-part units,
the containment statements at
the sites $(\mathrm{T})$, $(\mathrm{P})$, $(\mathrm{L})$, the sums over the stages with their
bound at window blocks, and clauses (c)--(e) of the proposition of the correlation theorem on
the operation $\mathbf{R}$ with the bound per piece and its sum over the plain interior units.
\end{itemize}
Then the extension statement for the plain class $\mathfrak{A}$ holds.
\item[(C)] Assume the four statements listed in (B), and in addition
\begin{itemize}
\item the relative bound for the boundary terms $\mathbf{B}'^{(k)}$ in the rotated sup norm
$\mathfrak{N}^{\mathbb{C}}$, in the form \eqref{4.22:eq-extension-hypotheses-a} of
Definition~\ref{4.22:def-extension-hypotheses}, and
\item the source-tube hypothesis of Definition~\ref{4.22:def-source-tube}, with its clauses
\eqref{4.22:eq-source-tube-a}.
\end{itemize}
Then the extension statement for the wrapped lineage $\mathfrak{S}$ holds.
\end{itemize}
Both (B) and (C) are implications. Neither class is free of the difference lemma: the terms of
the plain class lying inside a core depend on it exactly as the wrapped terms do (see the remark
accompanying Definition~\ref{4.22:def-cores-wrapped-class} on plain terms lying inside a core).
The only hypotheses of (C) absent from (B) are the relative bound for the boundary terms and the
source-tube hypothesis.

The clauses (a)--(j) are the following. Throughout, orbit data, their radii and the rotated sup
norm $\mathfrak{N}^{\mathbb{C}}$ are those of Definition~\ref{4.22:def-orbit-datum}, and
$\varepsilon_k(s)$ denotes the radius current at the law site $s$ at the point of the step in
question, as in Definition~\ref{4.22:def-radius-budget}.
\begin{itemize}
\item[(a)] (Creation at the site $(\mathrm{T})$.) Every term $\mathbf{B}^{(k)}(X,\mathcal{S})$
created at the site $(\mathrm{T})$ of step $k$ under the bound of the correlation theorem at
that site has an orbit datum, of radius $\varepsilon_\star$ at the coordinates of its new slot,
which belong to level $k-1$, and of the current radius $\varepsilon_k(s)$ at the older law
sites $s$; and the bound of the correlation theorem at the site $(\mathrm{T})$ holds with
$\mathfrak{N}^{\mathbb{C}}$ in the attached length $\mathsf{d}_k$ of the term at level $k$ (for a
stored term $F$ at level $\ell$ we write $\mathsf{d}_{\ell}(F)$ for its attached length):
\begin{equation}\label{4.22:eq-holomorphic-extension-1}
\begin{split}
&\sum_{\mathcal{S}}\exp\Bigl[\kappa_A^{\mathrm{pr}}(g_{k-1})\,|\mathcal{S}_{k-1}|
+\kappa_A(k-1)\sum_{j'<k-1}|\mathcal{S}_{j'}|\Bigr]\,
\mathfrak{N}^{\mathbb{C}}\bigl(\mathbf{B}^{(k)}(X,\mathcal{S})\bigr)\\
&\qquad\le B_0\,e^{-(1+4\beta_s)\kappa\,\mathsf{d}_k}.
\end{split}
\end{equation}
Here $B_0$ and $\beta_s$ are the constant and the rate constant of that bound of the
correlation theorem (the subscript of $\beta_s$ belongs to its name and does not denote a law
site). For terms of $\mathfrak{S}$ the right member of \eqref{4.22:eq-holomorphic-extension-1}
--- its rate, its constant and its strong-set weights --- is the one supplied by clause (r3) of
\eqref{4.22:eq-extension-hypotheses-a} at the site $(\mathrm{T})$; the replacement of
$\mathfrak{N}$ by $\mathfrak{N}^{\mathbb{C}}$ is the contribution of the present statement,
through the transfer lemma.
\item[(b)] (Chain.) For every $n\le\bar n$, where $\bar n$ denotes the number of stages of the
chain representation of the large-field stages of step $k$, each chain term
$\mathbb{C}^{(n)}_k(X,\mathcal{S})$ of that representation, read together with the uniformity
property of that representation, has an orbit datum on
\begin{equation*}
\mathfrak{D}^{\flat}_n(X)\times\mathfrak{D}_{\mathcal{S}}(X)\times
\operatorname{supp}\boldsymbol{\Phi}\times\mathfrak{G}_{\varepsilon_k}(X),
\end{equation*}
and $\|\mathbb{C}^{(n)}_k\|^{A,\flat,\mathbb{C}}_{(n)}\le C_0^{\sharp}$ at the rates of that
representation, where $C_0^{\sharp}$ is the constant of the chain representation; for chain
terms of $\mathfrak{S}$ this norm is read in the attached length, as
in clause (r3) of \eqref{4.22:eq-extension-hypotheses-a} at the site $(\mathrm{P})$.
\item[(c)] (The holomorphic extension.) For every real $\mathbf{A}\in\mathfrak{X}^{\sharp}$,
silent parameters included, and every real $\boldsymbol{\Phi}$, put
\begin{equation*}
\mathbb{C}^{(n),\sharp}_k(X;\cdot):=\sum_{\mathcal{S}}\mathbb{C}^{(n)}_k(X,\mathcal{S})^{\sharp}.
\end{equation*}
Then the map
$(P,g)\mapsto\mathbb{C}^{(n),\sharp}_k(X;P;\mathbf{A},\boldsymbol{\Phi},g)$ is holomorphic on
$\mathfrak{D}^{\flat}_n(X)\times\mathfrak{G}_{\varepsilon_k}(X)$ and local in $X$; for
$G$-valued $g$
\begin{equation*}
\mathbb{C}^{(n),\sharp}_k(X;P;\mathbf{A},\boldsymbol{\Phi},g)
=\mathbb{C}^{(n)}_k\bigl(X;P,\mathbf{A}^{[g]},\boldsymbol{\Phi}^{[g]}\bigr);
\end{equation*}
and, uniformly in $(\mathbf{A},\boldsymbol{\Phi})$,
\begin{equation}\label{4.22:eq-holomorphic-extension-2}
\bigl|\mathbb{C}^{(n),\sharp}_k(X;P;\mathbf{A},\boldsymbol{\Phi},g)\bigr|
\le C_0^{\sharp}\,e^{-a_m\mathsf{d}_m},
\end{equation}
where $e^{-a_m\mathsf{d}_m}$ is the decay factor of the chain representation at its rates (the
subscript $m$ is that representation's index and is not the exponent of the torus
$\mathbb{T}_m$).
At pointwise law sites
$\mathfrak{G}_{\varepsilon_k}\supset\{|\log g|\le\varepsilon_{\mathrm{live}}(k)\}$, with a room
of at least $\tfrac12\varepsilon_\star$. At the law sites $s$ of a frame $W$, the rotation $g$
applied at a reading of the term satisfies, by the radius budget
\eqref{4.22:eq-radius-budget-a},
\begin{equation*}
\operatorname{pol}(g)\le\eta^L_k(s)+(\text{the margins and the reserve of the step})
\le\tfrac14\,\varepsilon^{\mathrm{fr}}(W).
\end{equation*}
The same holds for every unit $\mathfrak{E}\in\mathcal{C}$ of the type of $\mathbf{B}$, with
the right member of the bound of its own site; for $\mathfrak{E}\in\mathfrak{S}$ this is the
right member of clauses (r1) and (r3) of \eqref{4.22:eq-extension-hypotheses-a}.
\item[(d)] (A function of the rotated variables.) For every term $\mathfrak{E}$ with orbit
datum $\mathfrak{E}^{\sharp}$, the value of $\mathfrak{E}^{\sharp}$ depends on
$(\mathbf{A},\boldsymbol{\Phi},g)$ only through the formally rotated variables
$(\mathbf{A}^{[g]},\boldsymbol{\Phi}^{\langle g\rangle})$. In particular the expression
$\mathfrak{E}(X;P;A^{[h]})$ used by the presentation statement at complex $h$, where $A$ denotes
the coordinate argument of the unit in the notation of that statement, is well defined.
\item[(e)] (Joint invariance under $\mathcal{G}^{\mathbb{C}}$.) For every $G^{\mathbb{C}}$-valued
site map $u$ joined to $1$ by a path $u_t$ such that $(P^{u_t},u_tg)$ lies in the domain of
$\mathfrak{E}^{\sharp}$ for all $t$,
\begin{equation*}
\mathfrak{E}^{\sharp}(P^u;\mathbf{A},\boldsymbol{\Phi},ug)
=\mathfrak{E}^{\sharp}(P;\mathbf{A},\boldsymbol{\Phi},g).
\end{equation*}
\item[(f)] (Cauchy bounds in the rotation.) The Cauchy bounds of
\eqref{4.22:eq-orbit-datum-properties-a} hold, and the Cauchy corollary for presentations
holds, in its clause (a) at the point of reading and in its clause (b) on the whole chart.
\item[(g)] (Closure.) The terms created by the large-field operation $\mathbf{R}$ of step $k$
--- the terms $V(Y)$, $\mathbf{R}'^{(k)}$, $\mathbf{B}'^{(k)}$, and, for every second-class
component $Y_i$, the two kept members $\mathfrak{K}(Y_i)$ --- are finite sums of stored terms
$(c,S,\mathcal{S})$ with orbit data of radius $\varepsilon_{k+1}$ on their surviving law sites
and of radius $\mathsf{e}(s)$ on the law sites $s$ of the new frames, and with the right
members of the site in $\mathcal{S}$-weighted form, by the output-chart proposition. These
right members are: for $\mathfrak{A}$, those of clauses (c)--(e) of the proposition of the
correlation theorem on the operation $\mathbf{R}$, with the bound per piece and its sum over
the plain interior units; for $\mathfrak{S}$, those of clauses (r1) and (r3) of
\eqref{4.22:eq-extension-hypotheses-a}, in the attached length; for $\mathbf{B}'^{(k)}$, the
relative bound \cite[(1.99)]{B-CMP122b}, as printed; for $\mathfrak{K}(Y_i)$, the bound
\cite[(1.69)]{B-CMP122b} at creation only, every later reading being by the source-tube
hypothesis with its clauses \eqref{4.22:eq-source-tube-a}.
\item[(h)] The weak coordinates of stored terms and their frame data, including those of the
kept members, are supplied: by the weak-collar lemma, at the radius $6.6\,C_1\delta'_j$, where
$\delta'_j$ is the parameter of level $j$ in the radii of the weak collar variables, and
under \eqref{4.22:eq-weak-collar-a}; by the counterpart at the site $(\mathrm{P})$ of the
anchored-sum lemma; by the entry statement; by the frame lemma; by the oscillation lemma for
kept members; and by the clauses \eqref{4.22:eq-source-tube-a}.
\item[(i)] The uniformity property of the chain representation is carried through the sites
$(\mathrm{T})$ and $(\mathrm{P})$, by the base-point lemma with its clause (iv) at the site
$(\mathrm{T})$ in reserve form.
\item[(j)] The radius budget \eqref{4.22:eq-radius-budget-a} holds. By the radius-budget
theorem, at each step $k'$ the debit satisfies
$\eta_{k'}(s)\le4\mu_{k'}\mathsf{e}(s)/\varepsilon_\star$, with $\nu(s)$ added once; with the
reserve included, the step costs at most $5\mu_{k'}\mathsf{e}(s)/\varepsilon_\star$; and for
every $K$ the total is at most $\tfrac{7}{32}\mathsf{e}(s)<\tfrac14\mathsf{e}(s)$. This holds
unconditionally: the margins and the reserve are chosen, and the two led-rotation bounds are
obligations of the users of the Cauchy-disc lemma, hypotheses of the displays that use bounds
taken on those discs, and not hypotheses of this clause. Moreover the profile lemma holds on
Ba{\l}aban's families and on those of the correlation theorem by the standing statements, and
on the families of the relative bound for the boundary terms its clause (iii), in the form for
those families, rests on the cell-family lemma.
\end{itemize}
No constant in (a)--(j) depends on $k$, $K$, $n$, $\bar n$, $N_0$, the scales $\check{R}$,
ages, the number of attempts, or the volume. The constants $\varepsilon_\star$,
$\varpi_{\mathrm{fr}}$, $C_{\mathrm{fr}}$, $\delta_{\natural}$, $C_{\natural}$,
$C_A^{(1/2)}$, $C_b$, $B_H^S$ depend on $(d,L,M)$, $\varpi_{\mathrm{fr}}$ also on $\theta_S$;
the constants $B_H$, $B_H^S$, $C_A^{(1/2)}$, $C_b$ are fixed after $\kappa_1$. The margin
$\mu_k$ depends on $g_k$, and $\varepsilon^{\mathrm{fr}}(W)$ on $g_{k_W}$; the number $N_0$
enters through condition (x) on the live radii only. For $\mathfrak{S}$ the constants and rates
of the right members are those of the relative bound for the boundary terms. The majorant
$C_{\Sigma}\check{R}_{k_W}^{p_{\Sigma}}$ of $\mathfrak{O}(v)$ depends on $\check{R}_{k_W}$, the
scale of the creating step $k_W$; it is paid through the source-tube hypothesis with its clauses
\eqref{4.22:eq-source-tube-a} and is never converted into a quantity at a later scale.
\end{proposition}

Clauses (a)--(c) and (g)--(j) are proved in this chapter by a single induction over the pairs
formed by a step and one of its sites, taken in the order of the steps and of their sites fixed
for this chapter; clauses (d)--(f) are consequences of the properties
\eqref{4.22:eq-orbit-datum-a} and \eqref{4.22:eq-orbit-datum-b} of orbit data alone.

For $G^{\mathbb{C}}$-valued maps $g$ on $\mathfrak{s}(S)$ we write $\boldsymbol{\Phi}^{\langle g\rangle}$ for the
formally rotated frame data, that is, the rotation formulas of
Definition~\ref{4.22:def-coordinates-frames} read in $\mathfrak{g}^{\mathbb{C}}$ and $G^{\mathbb{C}}$.

\begin{lemma}[Uniqueness, equivariance and Cauchy bounds for orbit data]
\label{4.22:lem-orbit-datum-properties}
Let $(c,S,\mathcal{S})$ be a stored term with an orbit datum $c^{\sharp}$ of radius $\varepsilon$ in the
sense of Definition~\ref{4.22:def-orbit-datum}. The arguments $(P;\mathbf{A},\boldsymbol{\Phi},g)$ below
range over the domain of $c^{\sharp}$. Then the following hold.
\begin{itemize}
\item[(O3)] The orbit datum $c^{\sharp}$ is unique.
\item[(O4)] For every $G$-valued map $v\colon\mathfrak{s}(S)\to G$,
\begin{equation*}
c^{\sharp}(P;\mathbf{A}^{[v]},\boldsymbol{\Phi}^{[v]},g)=c^{\sharp}(P;\mathbf{A},\boldsymbol{\Phi},gv).
\end{equation*}
\item[(O5)] For every $G$-valued site map $v$, with $vg$ the pointwise product on $\mathfrak{s}(S)$,
\begin{equation*}
c^{\sharp}(P^{v};\mathbf{A},\boldsymbol{\Phi},vg)=c^{\sharp}(P;\mathbf{A},\boldsymbol{\Phi},g).
\end{equation*}
\item[(O6)] Let $\mathbf{A},\mathbf{A}'$ be real in all coordinates, let $\boldsymbol{\Phi},\boldsymbol{\Phi}'$ be
real, let $g,g'\in\mathfrak{G}_{\varepsilon}(S)$, and suppose
$(\mathbf{A}^{[g]},\boldsymbol{\Phi}^{\langle g\rangle})
=(\mathbf{A}'^{[g']},\boldsymbol{\Phi}'^{\langle g'\rangle})$. Then
\begin{equation}\label{4.22:eq-orbit-datum-properties-b}
c^{\sharp}(P;\mathbf{A},\boldsymbol{\Phi},g)=c^{\sharp}(P;\mathbf{A}',\boldsymbol{\Phi}',g').
\end{equation}
\item[(O7)] Let $u_t$, $t\in[0,1]$, be a piecewise $C^1$ path of $G^{\mathbb{C}}$-valued site maps with
$u_0=1$ and $(P^{u_t},u_tg)\in\mathfrak{U}(S)\times\mathfrak{G}_{\varepsilon}(S)$ for all $t$. Then
\begin{equation}\label{4.22:eq-orbit-datum-properties-1}
c^{\sharp}(P^{u_1};\mathbf{A},\boldsymbol{\Phi},u_1g)=c^{\sharp}(P;\mathbf{A},\boldsymbol{\Phi},g).
\end{equation}
\item[(O8)] Let $g_0\in\mathfrak{G}_{\varepsilon}(S)$ and $Z\colon\mathfrak{s}(S)\to\mathfrak{g}^{\mathbb{C}}$,
and suppose that for every $s\in\mathfrak{s}(S)$
\begin{equation*}
\rho_s:=\frac{\varepsilon(s)-\operatorname{pol}(g_0(s))}{\|Z_s\|}\ \ge\ 2
\qquad(\rho_s:=\infty\ \text{if}\ Z_s=0).
\end{equation*}
Then
\begin{equation}\label{4.22:eq-orbit-datum-properties-a}
\bigl|c^{\sharp}(P;\mathbf{A},\boldsymbol{\Phi},g_0e^{Z})-c^{\sharp}(P;\mathbf{A},\boldsymbol{\Phi},g_0)\bigr|
\le\mathfrak{N}^{\mathbb{C}}_{\varepsilon}(c)\sum_{s\in\mathfrak{s}(S)}
\frac{2\|Z_s\|}{\varepsilon(s)-\operatorname{pol}(g_0(s))}.
\end{equation}
\item[(O9)] If every variable with law site $s$ is at its neutral value ($0$ for
$\mathfrak{g}$-valued, $1$ for $G$-valued variables), then $c^{\sharp}$ does not depend on $g(s)$:
\begin{equation}\label{4.22:eq-orbit-datum-properties-c}
c^{\sharp}(P;\mathbf{A},\boldsymbol{\Phi},g)=c^{\sharp}(P;\mathbf{A},\boldsymbol{\Phi},\tilde g)
\end{equation}
whenever $g,\tilde g\in\mathfrak{G}_{\varepsilon}(S)$ agree at every law site other than $s$.
\end{itemize}
\end{lemma}

\begin{proof}
Throughout, $\mathfrak{G}_{\varepsilon}(S)$ is the product over $s\in\mathfrak{s}(S)$ of the sets
$\mathfrak{G}_{\varepsilon(s)}$, and we use the clauses \eqref{4.22:eq-polar-decomposition-a},
\eqref{4.22:eq-polar-decomposition-c}, \eqref{4.22:eq-polar-decomposition-d} and
\eqref{4.22:eq-polar-decomposition-f} of Lemma~\ref{4.22:lem-polar-decomposition}. The clauses are
proved in the order (O3), (O4), (O5), (O9), (O8), (O7), (O6).

\emph{(O3).} Let $c^{\sharp}_1,c^{\sharp}_2$ be two orbit data of radius $\varepsilon$ for $(c,S,\mathcal{S})$,
and fix $(P,\mathbf{A},\boldsymbol{\Phi})$. By \eqref{4.22:eq-orbit-datum-a} the difference
$g\mapsto c^{\sharp}_1(P;\mathbf{A},\boldsymbol{\Phi},g)-c^{\sharp}_2(P;\mathbf{A},\boldsymbol{\Phi},g)$ is
holomorphic on $\mathfrak{G}_{\varepsilon}(S)$, and by \eqref{4.22:eq-orbit-datum-b} both terms equal
$c(S;P,\mathbf{A}^{[g]},\boldsymbol{\Phi}^{[g]})$ when $g$ is $G$-valued, so the difference vanishes on
$G^{\mathfrak{s}(S)}$. By the identity principle \eqref{4.22:eq-polar-decomposition-c} it vanishes
identically. As $(P,\mathbf{A},\boldsymbol{\Phi})$ was arbitrary, $c^{\sharp}_1=c^{\sharp}_2$.

\emph{(O4).} The pair $(\mathbf{A}^{[v]},\boldsymbol{\Phi}^{[v]})$ lies in the domain, since
$\mathfrak{D}_{\mathcal{S}}(S)$ is stable under rotation by $G$-valued maps and the supports of the
measures of $\boldsymbol{\Phi}$ are $\operatorname{Ad}_G$-stable (Definition~\ref{4.22:def-orbit-datum}).
By \eqref{4.22:eq-polar-decomposition-a}, $\operatorname{pol}(g(s)v(s))=\operatorname{pol}(g(s))$, so $gv\in\mathfrak{G}_{\varepsilon}(S)$
for $g\in\mathfrak{G}_{\varepsilon}(S)$. Hence by \eqref{4.22:eq-orbit-datum-a} both sides of (O4) are
holomorphic functions of $g$ on $\mathfrak{G}_{\varepsilon}(S)$. For $G$-valued $g$, applying
\eqref{4.22:eq-orbit-datum-b} to the left side at the $G$-valued map $g$ and to the right side at the
$G$-valued map $gv$, and using that rotating first by $v$ and then by $g$ is rotating by $gv$, both
sides equal
\begin{equation*}
c\bigl(S;P,(\mathbf{A}^{[v]})^{[g]},(\boldsymbol{\Phi}^{[v]})^{[g]}\bigr)
=c\bigl(S;P,\mathbf{A}^{[gv]},\boldsymbol{\Phi}^{[gv]}\bigr).
\end{equation*}
Their difference is therefore holomorphic on $\mathfrak{G}_{\varepsilon}(S)$ and vanishes on
$G^{\mathfrak{s}(S)}$, and \eqref{4.22:eq-polar-decomposition-c} gives (O4).

\emph{(O5).} Since $\mathfrak{U}(S)$ is $\mathcal{G}$-invariant, $P^{v}\in\mathfrak{U}(S)$; since
$\operatorname{pol}(v(s)g(s))=\operatorname{pol}(g(s))$ by \eqref{4.22:eq-polar-decomposition-a},
$vg\in\mathfrak{G}_{\varepsilon}(S)$. By \eqref{4.22:eq-orbit-datum-a} both sides of (O5) are holomorphic
in $g$ on $\mathfrak{G}_{\varepsilon}(S)$. For $G$-valued $g$, \eqref{4.22:eq-orbit-datum-b} at $vg$, the
rule that rotating by $g$ and then by $v$ is rotating by $vg$, the joint invariance of $c$
(Definition~\ref{4.22:def-orbit-datum}) and \eqref{4.22:eq-orbit-datum-b} at $g$ give
\begin{equation*}
c^{\sharp}(P^{v};\mathbf{A},\boldsymbol{\Phi},vg)
=c\bigl(S;P^{v},(\mathbf{A}^{[g]})^{[v]},(\boldsymbol{\Phi}^{[g]})^{[v]}\bigr)
=c\bigl(S;P,\mathbf{A}^{[g]},\boldsymbol{\Phi}^{[g]}\bigr)
=c^{\sharp}(P;\mathbf{A},\boldsymbol{\Phi},g).
\end{equation*}
The identity principle \eqref{4.22:eq-polar-decomposition-c} gives (O5) for all $g\in\mathfrak{G}_{\varepsilon}(S)$.

\emph{(O9).} For $g\in\mathfrak{G}_{\varepsilon}(S)$ let $g_{s\to1}$ agree with $g$ off $s$ and equal $1$
at $s$; it lies in $\mathfrak{G}_{\varepsilon}(S)$, as $\operatorname{pol}(1)=0<\varepsilon(s)$. The function
$g\mapsto c^{\sharp}(P;\mathbf{A},\boldsymbol{\Phi},g)-c^{\sharp}(P;\mathbf{A},\boldsymbol{\Phi},g_{s\to1})$ is
holomorphic on $\mathfrak{G}_{\varepsilon}(S)$ by \eqref{4.22:eq-orbit-datum-a}. For $G$-valued $g$ both
terms are given by \eqref{4.22:eq-orbit-datum-b}; the variables with law site $s$ are rotated through
$g(s)$ only, and their neutral values are fixed by every rotation
($\operatorname{Ad}_{g(s)}0=0$, $g(s)1g(s)^{-1}=1$), so the rotated data at $g$ and at $g_{s\to1}$
coincide and the difference vanishes on $G^{\mathfrak{s}(S)}$. By
\eqref{4.22:eq-polar-decomposition-c} it vanishes identically, which is
\eqref{4.22:eq-orbit-datum-properties-c}, applied to $g$ and to $\tilde g$.

\emph{(O8).} For $\tau=(\tau_s)_{s\in\mathfrak{s}(S)}$ put
$f(\tau):=c^{\sharp}\bigl(P;\mathbf{A},\boldsymbol{\Phi},(g_0(s)e^{\tau_sZ_s})_s\bigr)$. By the subadditivity of
$\operatorname{pol}$ and the bound $\operatorname{pol}(e^{Z})\le\|Z\|$ of \eqref{4.22:eq-polar-decomposition-a},
\begin{equation*}
\operatorname{pol}\bigl(g_0(s)e^{\tau_sZ_s}\bigr)\le\operatorname{pol}(g_0(s))+|\tau_s|\,\|Z_s\|,
\end{equation*}
which is $<\varepsilon(s)$ when $|\tau_s|<\rho_s$. Hence, by \eqref{4.22:eq-orbit-datum-a}, $f$ is
holomorphic on the polydisc $\prod_s\{|\tau_s|<\rho_s\}$, and $|f|\le\mathfrak{N}^{\mathbb{C}}_{\varepsilon}(c)$
there, the rotated sup norm being the supremum of $|c^{\sharp}|$ over its domain. Write $\mathbb{1}$ for
the point with all coordinates $1$; since $\rho_s\ge2$, the segment $\{t\mathbb{1}:0\le t\le1\}$ lies in
the polydisc, and $f(\mathbb{1})$, $f(0)$ are the two values in \eqref{4.22:eq-orbit-datum-properties-a}.
By the chain rule,
\begin{equation*}
f(\mathbb{1})-f(0)=\int_0^1\sum_{s}\partial_sf(t\mathbb{1})\,dt .
\end{equation*}
Fix $t\in[0,1]$ and $s$ with $Z_s\neq0$ (if $Z_s=0$, $f$ does not depend on $\tau_s$ and
$\partial_sf=0$). As a function of $\tau_s$ alone, the other coordinates held at $t$, $f$ is
holomorphic on $|\tau_s|<\rho_s$, and for every $r<\rho_s-t$ the circle $|\tau_s-t|=r$ lies in this
disc. Cauchy's inequality on it gives $|\partial_sf(t\mathbb{1})|\le\mathfrak{N}^{\mathbb{C}}_{\varepsilon}(c)/r$;
letting $r$ increase to $\rho_s-t\ge\rho_s-1\ge\tfrac12\rho_s$,
\begin{equation*}
|\partial_sf(t\mathbb{1})|\le\frac{2\,\mathfrak{N}^{\mathbb{C}}_{\varepsilon}(c)}{\rho_s}
=\mathfrak{N}^{\mathbb{C}}_{\varepsilon}(c)\,\frac{2\|Z_s\|}{\varepsilon(s)-\operatorname{pol}(g_0(s))}.
\end{equation*}
Summing over $s$ and integrating over $t\in[0,1]$ gives \eqref{4.22:eq-orbit-datum-properties-a}.

\emph{(O7).} The argument is carried out on the complex manifold
$\mathfrak{U}(S)\times\mathfrak{G}_{\varepsilon}(S)$ with $(\mathbf{A},\boldsymbol{\Phi})$ held fixed. Write
$Dc^{\sharp}_{(P,g)}$ for the differential at $(P,g)$ of $(P,g)\mapsto c^{\sharp}(P;\mathbf{A},\boldsymbol{\Phi},g)$.
Let $Z$ be a $\mathfrak{g}$-valued site map. For real $\sigma$, $e^{i\sigma Z}$ is a $G$-valued site map,
so (O5) with $v=e^{i\sigma Z}$ gives
\begin{equation*}
c^{\sharp}(P^{e^{i\sigma Z}};\mathbf{A},\boldsymbol{\Phi},e^{i\sigma Z}g)
=c^{\sharp}(P;\mathbf{A},\boldsymbol{\Phi},g);
\end{equation*}
differentiating at $\sigma=0$,
\begin{equation*}
Dc^{\sharp}_{(P,g)}[\delta_Z]=0,\qquad
\delta_Z:=\frac{d}{d\sigma}\Big|_{\sigma=0}\bigl(P^{e^{i\sigma Z}},e^{i\sigma Z}g\bigr).
\end{equation*}
The action $(u,(P,g))\mapsto(P^{u},ug)$ of $G^{\mathbb{C}}$-valued site maps is holomorphic in $u$, so
the same formula defines $\delta_Z$ for $\mathfrak{g}^{\mathbb{C}}$-valued $Z$: it is the complex
differential at $u=1$ of $u\mapsto(P^{u},ug)$ applied to $iZ$, hence $\mathbb{C}$-linear in $Z$. By
\eqref{4.22:eq-orbit-datum-a}, $c^{\sharp}$ is holomorphic in $(P,g)$, so $Dc^{\sharp}_{(P,g)}$ is
$\mathbb{C}$-linear. Writing a $\mathfrak{g}^{\mathbb{C}}$-valued $Z$ as $Z_1+iZ_2$ with $Z_1,Z_2$
$\mathfrak{g}$-valued, we get
\begin{equation*}
Dc^{\sharp}_{(P,g)}[\delta_Z]=Dc^{\sharp}_{(P,g)}[\delta_{Z_1}]+i\,Dc^{\sharp}_{(P,g)}[\delta_{Z_2}]=0
\end{equation*}
for all $\mathfrak{g}^{\mathbb{C}}$-valued $Z$ and at every point
$(P,g)$ of $\mathfrak{U}(S)\times\mathfrak{G}_{\varepsilon}(S)$.

Now take the path $u_t$. On each interval where it is $C^1$ put $Z_t:=-i\,\dot u_tu_t^{-1}$, so
$iZ_t=\dot u_tu_t^{-1}$; since $u_t$ is $G^{\mathbb{C}}$-valued, $\det u_t\equiv1$, so by Jacobi's
formula the matrix $\dot u_tu_t^{-1}$ is traceless, and $Z_t$ is $\mathfrak{g}^{\mathbb{C}}$-valued.
The map $h\mapsto u_{t+h}u_t^{-1}$ equals $1$ at $h=0$ and has derivative $iZ_t$ there, and, the map
$(u,(P,g))\mapsto(P^{u},ug)$ being an action,
\begin{equation*}
(P^{u_{t+h}},u_{t+h}g)
=\bigl((P^{u_t})^{u_{t+h}u_t^{-1}},\,u_{t+h}u_t^{-1}\,u_tg\bigr).
\end{equation*}
Hence
\begin{equation*}
\frac{d}{dt}\,c^{\sharp}(P^{u_t};\mathbf{A},\boldsymbol{\Phi},u_tg)
=Dc^{\sharp}_{(P^{u_t},u_tg)}[\delta_{Z_t}]=0,
\end{equation*}
the last equality by the previous paragraph at the point $(P^{u_t},u_tg)$, which lies in
$\mathfrak{U}(S)\times\mathfrak{G}_{\varepsilon}(S)$ by hypothesis. No condition on the values of $u_t$ at
the roots of the frames enters: on the chart $\mathfrak{G}_{\varepsilon}(S)$ the law sites of the frames
are holomorphic coordinates like all the others. The function
$t\mapsto c^{\sharp}(P^{u_t};\mathbf{A},\boldsymbol{\Phi},u_tg)$ is continuous and piecewise $C^1$ with
vanishing derivative, hence constant; its values at $t=1$ and at $t=0$ (where $u_0=1$) give
\eqref{4.22:eq-orbit-datum-properties-1}.

\emph{(O6).} Put $u:=g'^{-1}g$ sitewise and, by \eqref{4.22:eq-polar-decomposition-a}, write
$u(s)=v(s)e^{Y(s)}$ with $v(s)\in G$, $Y(s)\in\mathfrak{g}$. The hypothesis
$(\mathbf{A}^{[g]},\boldsymbol{\Phi}^{\langle g\rangle})=(\mathbf{A}'^{[g']},\boldsymbol{\Phi}'^{\langle g'\rangle})$
says, variable by variable, that for every variable with law site $s$,
\begin{equation*}
\operatorname{Ad}_{u(s)}\mathbf{A}(a)=\mathbf{A}'(a)\in\mathfrak{g}
\qquad(\text{resp. } u(s)Vu(s)^{-1}=V'\in G),
\end{equation*}
so that $\operatorname{Ad}_{e^{Y(s)}}\mathbf{A}(a)=\operatorname{Ad}_{v(s)^{-1}}\mathbf{A}'(a)\in\mathfrak{g}$
(resp. $e^{Y(s)}Ve^{-Y(s)}=v(s)^{-1}V'v(s)\in G$), while $\mathbf{A}(a)\in\mathfrak{g}$ (resp. $V\in G$)
by reality. By \eqref{4.22:eq-polar-decomposition-d}, $Y(s)$ commutes with every variable with law
site $s$. Then $\operatorname{Ad}_{e^{Y(s)}}$ fixes each of them, and
$(\mathbf{A}',\boldsymbol{\Phi}')=(\mathbf{A}^{[v]},\boldsymbol{\Phi}^{[v]})$. By (O4),
\begin{equation*}
c^{\sharp}(P;\mathbf{A}',\boldsymbol{\Phi}',g')=c^{\sharp}(P;\mathbf{A},\boldsymbol{\Phi},g''),
\qquad g'':=g'v,\qquad g=g'u=g''e^{Y}.
\end{equation*}
Let $\varphi(\tau):=c^{\sharp}(P;\mathbf{A},\boldsymbol{\Phi},g''e^{-i\tau Y})$ on the set $\mathcal{D}$ of
$\tau\in\mathbb{C}$ with $g''e^{-i\tau Y}\in\mathfrak{G}_{\varepsilon}(S)$; $\mathcal{D}$ is open, since
$\mathfrak{G}_{\varepsilon}(S)$ is open by \eqref{4.22:eq-polar-decomposition-c}, and $\varphi$ is
holomorphic on it by \eqref{4.22:eq-orbit-datum-a}. For $\tau=\sigma+it$ with $\sigma$ real and
$0\le t\le1$ we have $-i\tau=t-i\sigma$, and \eqref{4.22:eq-polar-decomposition-f} gives, at each $s$,
\begin{align*}
\operatorname{pol}\bigl(g''(s)e^{-i\tau Y(s)}\bigr)
&=\operatorname{pol}\bigl(g''(s)e^{tY(s)}\bigr)
\le\max\bigl\{\operatorname{pol}(g''(s)),\operatorname{pol}(g''(s)e^{Y(s)})\bigr\}\\
&=\max\bigl\{\operatorname{pol}(g'(s)),\operatorname{pol}(g(s))\bigr\}<\varepsilon(s),
\end{align*}
using convexity on $[0,1]$ and $\operatorname{pol}(g'(s)v(s))=\operatorname{pol}(g'(s))$. Thus the closed
strip $\{0\le\operatorname{Im}\tau\le1\}$ lies in the open set $\mathcal{D}$. For real $\tau$,
$e^{-i\tau Y}$ is $G$-valued and fixes $(\mathbf{A},\boldsymbol{\Phi})$, because $Y(s)$ commutes with every
variable at $s$; so (O4) gives $\varphi(\tau)=c^{\sharp}(P;\mathbf{A},\boldsymbol{\Phi},g'')=\varphi(0)$. The
strip is connected, so it lies in one connected component of $\mathcal{D}$; on that component
$\varphi-\varphi(0)$ is holomorphic and vanishes on $\mathbb{R}$, hence vanishes identically. In
particular $\varphi(i)=\varphi(0)$, that is,
\begin{equation*}
c^{\sharp}(P;\mathbf{A},\boldsymbol{\Phi},g''e^{Y})=c^{\sharp}(P;\mathbf{A},\boldsymbol{\Phi},g''),
\end{equation*}
which with $g''e^{Y}=g$ and the display above is \eqref{4.22:eq-orbit-datum-properties-b}.
\end{proof}

\begin{remark}
We record three observations on the scope of Lemma~\ref{4.22:lem-orbit-datum-properties}.
\begin{itemize}
\item[(1)] Clause (O6) of Lemma~\ref{4.22:lem-orbit-datum-properties} holds for every orbit datum,
with no proviso on the kind of coordinate involved. In particular it covers the coordinates attached
to the bonds of the full change of variables and the frame data $\boldsymbol{\Phi}$, as well as all
other coordinates of $\mathbf{A}$.
\item[(2)] Clause (O9) is the counterpart, on the chart $\mathfrak{G}_{\varepsilon}(S)$, of the first
clause of the lemma on the slots of a completed chain, which states that the own slots of a completed
chain sit at their neutral value $0$ and are never rotated.
\item[(3)] At a weak coordinate $a$ the term $c$ is already holomorphic in the complex slot.
By the second clause of the lemma on the slots of a completed chain (the clause on weak
coordinates), $c^{\sharp}(P;\mathbf{A},\boldsymbol{\Phi},g)$ is obtained from $c$ by rotating the weak
coordinate $\mathbf{A}(a)$ through $g$, that is, by evaluating $c$ at $\mathbf{A}^{[g]}$ in place of
$\mathbf{A}$, wherever
\begin{equation*}
e^{2\operatorname{pol}(g)}|\mathbf{A}(a)|<\varrho .
\end{equation*}
The chart $\mathfrak{G}_{\varepsilon}(S)$ keeps the two conditions $|\mathbf{A}(a)|<\varrho$ and
$\operatorname{pol}(g)<\varepsilon$ at the same time. A tube in the slot variable alone cannot keep both.
\end{itemize}
\end{remark}

\begin{lemma}[The cluster expansion with a complex-manifold parameter]\label{4.22:lem-cluster-manifold}
Assume the following two statements:
\begin{itemize}
\item the lemma on the cluster expansion with parameters;
\item the proposition on the parameter-free majorants of the cluster expansion, clauses (ii)
and (iii).
\end{itemize}
In the lemma on the cluster expansion with parameters, let the parameter range of the element $e$
be
\begin{equation*}
D_e := (\text{real box coordinates})\times(\text{complex polydisc coordinates})\times\mathfrak{M}_e,
\end{equation*}
where $\mathfrak{M}_e$ is a finite-dimensional complex manifold. Let $v_e$ be continuous on
$D_e$, holomorphic in the polydisc coordinates and in the coordinates of $\mathfrak{M}_e$ for
fixed box coordinates, and such that $|v_e|\le\nu(e)$ on $D_e$. Suppose that all the other objects
entering that lemma do not depend on the parameter: the pieces of the expansion listed there, the
large-field bonds, the decoupled cubes, the Gaussian measure with its decoupling parameters
$\mathsf{s}$, the products $\chi^{(j)}\chi'^{(j)}$ of cut-off functions of that lemma, the
partitions and the contours. Let $D$ be the fibre product of the ranges $D_e$ over the coordinates
shared between elements, as in that lemma. Then:
\begin{enumerate}
\item[(a)] each term $\mathbb{C}_I(X)$ of the expansion, indexed as in that lemma by $I$ and $X$,
is defined and continuous on $D$;
\item[(b)] $\mathbb{C}_I(X)$ is holomorphic in every coordinate that is a polydisc coordinate or a
coordinate of $\mathfrak{M}_e$ for every element $e$ in whose range $D_e$ it occurs;
\item[(c)] clauses (ii) and (iii) of the proposition on the parameter-free majorants of the cluster
expansion hold with the supremum over $D$ taken inside the norm.
\end{enumerate}
\end{lemma}

\begin{proof}
The argument is that of the lemma on the cluster expansion with parameters, which runs through the
steps of the proof of the proposition on the parameter-free majorants of the cluster expansion at
each point of $D$; we follow those steps. The frame data $\boldsymbol{\Phi}$ enter as real box
coordinates, so they are covered by the box factor of $D_e$.

The first step of the proof of the proposition on the parameter-free majorants of the cluster
expansion is an identity that holds at each fixed parameter value; since every object other than
the functions $v_e$ is parameter-free, it holds at every point of $D$.

The second to sixth steps of that proof majorise every term of the Mayer--polymer--Ursell series by a
majorant built from the numbers $\nu(e)$ alone. Since $|v_e|\le\nu(e)$ holds on all of $D_e$, and
the pieces, bonds, cubes, Gaussian measure, products of cut-off functions, partitions and contours
do not depend on the parameter, this majorant is independent of the point of $D$. Hence the series
converges absolutely and uniformly on $D$, and the bounds of clauses (ii) and (iii) of that
proposition hold for the supremum over $D$ of each term, that is, with $\sup_D$ inside the norm.
This is conclusion (c).

Each term of the series is a finite sum of products of integrals, over parameter-free contours,
of finitely many of the functions $v_e$. Such a term is continuous on $D$, because the $v_e$ are
continuous and the contours do not move; and it is holomorphic in every coordinate that is a
polydisc or manifold coordinate for every element $e$ whose range $D_e$ contains that coordinate,
because for fixed box coordinates
the integrand is holomorphic in that coordinate and the contour of integration is fixed. By the
uniform convergence just established, $\mathbb{C}_I(X)$ is a uniform limit of continuous
functions on $D$, hence continuous, which gives (a). Uniform limits of holomorphic functions on a
complex manifold are holomorphic: in a chart this is Weierstrass's theorem. Applied in charts of
the polydisc and of $\mathfrak{M}_e$, at fixed box coordinates, this gives (b).
\end{proof}

In the next lemma the letters $\rho,\rho',\Lambda,\varepsilon,\eta_0,\vartheta$ and $m$ are used locally;
$\Lambda$ is a number, not a small-field region, and $m$ is a point of $\mathfrak{M}$.

\begin{lemma}[Cauchy estimate with a riding rotation]\label{4.22:lem-cauchy-riding}
Let $\mathfrak{M}$ be a complex manifold, let $S$ be a lattice domain with a finite set of law
sites $\mathfrak{s}(S)$, and let $\varepsilon,\eta_0,\vartheta,\mu$ be functions on $\mathfrak{s}(S)$.
Write $\mathfrak{G}_{\varepsilon}(S)$ for the set of maps $g\colon\mathfrak{s}(S)\to G^{\mathbb{C}}$
with $\operatorname{pol}(g(s))<\varepsilon(s)$ for every $s$, and $u g'$ for the sitewise product
$(ug')(s)=u(s)g'(s)$. Let $\rho>0$ and let $f(\mathsf{t},m,g)$ be holomorphic on
$\{|\mathsf{t}|<\rho\}\times\mathfrak{M}\times\mathfrak{G}_{\varepsilon}(S)$ with
$|f|\le\mathfrak{N}$. Let
\begin{equation*}
u\colon\{|\mathsf{t}|<\rho\}\times\mathfrak{M}\to(G^{\mathbb{C}})^{\mathfrak{s}(S)}
\end{equation*}
be holomorphic, and assume that for all $s\in\mathfrak{s}(S)$, $|\mathsf{t}|<\rho$ and $m\in\mathfrak{M}$
\begin{equation*}
\operatorname{pol}\bigl(u(0,m)(s)\bigr)\le\eta_0(s),\qquad
\bigl|\log\bigl[u(0,m)(s)^{-1}u(\mathsf{t},m)(s)\bigr]\bigr|\le\vartheta(s)\le 10^{-2}.
\end{equation*}
Let $\mu(s)>0$ for every $s$, and assume that $\varepsilon':=\varepsilon-\eta_0-\mu>0$ on
$\mathfrak{s}(S)$. Put
\begin{equation*}
\rho':=\rho\cdot\min\Bigl\{1,\min_{s}\frac{\mu(s)}{\vartheta(s)}\Bigr\},\qquad
\Lambda:=\max\Bigl\{1,\max_{s}\frac{\vartheta(s)}{\mu(s)}\Bigr\},
\end{equation*}
where a quotient $\mu(s)/\vartheta(s)$ with $\vartheta(s)=0$ is read as $+\infty$. Then
\begin{equation*}
F(\mathsf{t},m,g'):=f\bigl(\mathsf{t},m,u(\mathsf{t},m)g'\bigr)
\end{equation*}
is holomorphic on $\{|\mathsf{t}|<\rho'\}\times\mathfrak{M}\times\mathfrak{G}_{\varepsilon'}(S)$,
satisfies $|F|\le\mathfrak{N}$ there, and
\begin{equation}\label{4.22:eq-cauchy-riding-1}
\bigl|\partial_{\mathsf{t}}F(0,m,g')\bigr|\le\frac{\mathfrak{N}}{\rho}\,\Lambda
\qquad\bigl(m\in\mathfrak{M},\ g'\in\mathfrak{G}_{\varepsilon'}(S)\bigr).
\end{equation}
\end{lemma}

\begin{proof}
Fix $m\in\mathfrak{M}$ and $s\in\mathfrak{s}(S)$, and put
\begin{equation*}
Z_s(\mathsf{t}):=\log\bigl[u(0,m)(s)^{-1}u(\mathsf{t},m)(s)\bigr],\qquad|\mathsf{t}|<\rho .
\end{equation*}
The map $\mathsf{t}\mapsto u(0,m)(s)^{-1}u(\mathsf{t},m)(s)$ is holomorphic, and the logarithm is
holomorphic on the neighbourhood of the identity in which it is taken by hypothesis; hence
$Z_s$ is holomorphic with values in $\mathfrak{g}^{\mathbb{C}}$. Moreover $Z_s(0)=0$, since the
argument of the logarithm is the identity at $\mathsf{t}=0$, and
$|Z_s|\le\vartheta(s)$ by hypothesis.

We claim that
\begin{equation}\label{4.22:eq-cauchy-riding-2}
|Z_s(\mathsf{t})|\le\vartheta(s)\,\frac{|\mathsf{t}|}{\rho}\qquad(|\mathsf{t}|<\rho).
\end{equation}
If $\vartheta(s)=0$ then $Z_s\equiv0$ and there is nothing to prove. Otherwise let $\ell$ be a
linear functional on the matrix space whose dual norm with respect to $|\cdot|$ is one. The function
$\varphi(\zeta):=\ell(Z_s(\rho\zeta))/\vartheta(s)$ is holomorphic on the unit disc, bounded by
one there, and vanishes at $\zeta=0$; by Schwarz's lemma $|\varphi(\zeta)|\le|\zeta|$, that is,
$|\ell(Z_s(\mathsf{t}))|\le\vartheta(s)|\mathsf{t}|/\rho$. Since $|Z|$ is the supremum of
$|\ell(Z)|$ over such functionals, \eqref{4.22:eq-cauchy-riding-2} follows.

Since the exponential inverts the logarithm, $u(\mathsf{t},m)(s)=u(0,m)(s)\,e^{Z_s(\mathsf{t})}$.
Let $g'\in\mathfrak{G}_{\varepsilon'}(S)$ and $|\mathsf{t}|<\rho'$. By the subadditivity
$\operatorname{pol}(g_1g_2)\le\operatorname{pol}(g_1)+\operatorname{pol}(g_2)$, applied twice, and the bound
$\operatorname{pol}(e^{Z})\le\|Z\|\le|Z|$ for $Z\in\mathfrak{g}^{\mathbb{C}}$, both from
\eqref{4.22:eq-polar-decomposition-a} in Lemma~\ref{4.22:lem-polar-decomposition}, together with
\eqref{4.22:eq-cauchy-riding-2},
\begin{align*}
\operatorname{pol}\bigl(u(\mathsf{t},m)(s)g'(s)\bigr)
&\le\operatorname{pol}\bigl(u(0,m)(s)\bigr)+\operatorname{pol}\bigl(e^{Z_s(\mathsf{t})}\bigr)
+\operatorname{pol}\bigl(g'(s)\bigr)\\
&\le\eta_0(s)+\vartheta(s)\frac{|\mathsf{t}|}{\rho}+\operatorname{pol}\bigl(g'(s)\bigr).
\end{align*}
If $\vartheta(s)>0$, then $|\mathsf{t}|<\rho'\le\rho\,\mu(s)/\vartheta(s)$ gives
$\vartheta(s)|\mathsf{t}|/\rho<\mu(s)$; if $\vartheta(s)=0$, the middle term is
$0<\mu(s)$. As $\operatorname{pol}(g'(s))<\varepsilon'(s)$,
\begin{equation*}
\operatorname{pol}\bigl(u(\mathsf{t},m)(s)g'(s)\bigr)<\eta_0(s)+\mu(s)+\varepsilon'(s)=\varepsilon(s).
\end{equation*}
This holds at every law site, so $u(\mathsf{t},m)g'\in\mathfrak{G}_{\varepsilon}(S)$; also
$\rho'\le\rho$. Hence $(\mathsf{t},m,g')\mapsto(\mathsf{t},m,u(\mathsf{t},m)g')$ maps
$\{|\mathsf{t}|<\rho'\}\times\mathfrak{M}\times\mathfrak{G}_{\varepsilon'}(S)$ holomorphically
(sitewise matrix multiplication of holomorphic factors) into the domain of $f$. As a
composition of holomorphic maps, $F$ is holomorphic there, and $|F|\le\mathfrak{N}$ because $F$
takes values of $f$.

Fix $m$ and $g'$. For $0<\rho''<\rho'$ the Cauchy formula on the circle $|\mathsf{t}|=\rho''$
gives
\begin{equation*}
\bigl|\partial_{\mathsf{t}}F(0,m,g')\bigr|
=\Bigl|\frac{1}{2\pi i}\oint_{|\mathsf{t}|=\rho''}\frac{F(\mathsf{t},m,g')}{\mathsf{t}^2}\,d\mathsf{t}\Bigr|
\le\frac{\mathfrak{N}}{\rho''}.
\end{equation*}
Letting $\rho''\uparrow\rho'$ yields $|\partial_{\mathsf{t}}F(0,m,g')|\le\mathfrak{N}/\rho'$.
Finally, for positive numbers $a_s$ (some possibly $+\infty$) one has
$1/\min\{1,\min_s a_s\}=\max\{1,\max_s 1/a_s\}$; with $a_s=\mu(s)/\vartheta(s)$ this reads
$\rho/\rho'=\Lambda$, which is \eqref{4.22:eq-cauchy-riding-1}.
\end{proof}

In this lemma the radius of the chart on which $F$ lives is reduced by $\eta_0+\mu$ only, that
is, by the size of the rotation at $\mathsf{t}=0$ and by a margin $\mu$ that may be chosen freely;
the variation $\vartheta$ of the rotation in $\mathsf{t}$ is not charged to the chart radius
$\varepsilon'$; it shrinks only the disc in $\mathsf{t}$, and so appears in the bound
\eqref{4.22:eq-cauchy-riding-1} through the factor $\Lambda$.

\begin{lemma}[No shrinking of the chart on a polydisc]\label{4.22:lem-no-shrink}
Let finitely many blocks $\square$ be given, with radii $\rho_{\square}>0$, and let
\begin{equation*}
\mathfrak{P}:=\prod_{\square}\{t_{\square}\in\mathbb{C}:|t_{\square}|<\rho_{\square}\}
\end{equation*}
be the corresponding polydisc. Let $\mathfrak{M}$ be a complex manifold, let $S$ be a domain, with
$\mathfrak{s}(S)$ its set of law sites, and let $\varepsilon$ be a positive function on $\mathfrak{s}(S)$. Let $\eta_0$ and
$\mu$ be functions on $\mathfrak{s}(S)$ with values in $[0,\infty)$. Let $f$ be
holomorphic on $\mathfrak{P}\times\mathfrak{M}\times\mathfrak{G}_{\varepsilon}(S)$ with
$|f|\le\mathfrak{N}$. Let
\begin{equation*}
u:\mathfrak{P}\times\mathfrak{M}\to(G^{\mathbb{C}})^{\mathfrak{s}(S)}
\end{equation*}
be holomorphic. Assume that for every $s\in\mathfrak{s}(S)$ and every $m\in\mathfrak{M}$ one has
$\operatorname{pol}(u(0,m)(s))\le\eta_0(s)$, and that
\begin{equation}\label{4.22:eq-no-shrink-1}
\sup_{(t,m)\in\mathfrak{P}\times\mathfrak{M}}
\bigl|\log\bigl[u(0,m)(s)^{-1}u(t,m)(s)\bigr]\bigr|\le\mu(s)
\qquad(s\in\mathfrak{s}(S)),
\end{equation}
where $\log h\in\mathfrak{g}^{\mathbb{C}}$ denotes the logarithm of $h\in G^{\mathbb{C}}$, with
$e^{\log h}=h$. Assume also that
$\varepsilon'(s):=\varepsilon(s)-\eta_0(s)-\mu(s)>0$ for every $s\in\mathfrak{s}(S)$. Define
\begin{equation}\label{4.22:eq-no-shrink-2}
F(t,m,g'):=f\bigl(t,m,u(t,m)g'\bigr),
\end{equation}
where the product is taken site by site, $(u(t,m)g')(s)=u(t,m)(s)\,g'(s)$. Then $F$ is holomorphic
on the whole of $\mathfrak{P}\times\mathfrak{M}\times\mathfrak{G}_{\varepsilon'}(S)$, and
$|F|\le\mathfrak{N}$ there. Consequently every Cauchy, M\"obius, fundamental-theorem-of-calculus or
Schwarz estimate which the site of application takes on $\mathfrak{P}$ holds for $F$ with that site's
constants and with the same bound $\mathfrak{N}$.
\end{lemma}

\begin{proof}
Fix $(t,m)\in\mathfrak{P}\times\mathfrak{M}$ and $s\in\mathfrak{s}(S)$, and write $Z_s(t,m)$ for the
logarithm in \eqref{4.22:eq-no-shrink-1}. It is an element of $\mathfrak{g}^{\mathbb{C}}$ with
\begin{equation*}
u(0,m)(s)^{-1}u(t,m)(s)=e^{Z_s(t,m)},
\end{equation*}
and by hypothesis \eqref{4.22:eq-no-shrink-1} its norm satisfies $|Z_s(t,m)|\le\mu(s)$. By the last
assertion of \eqref{4.22:eq-polar-decomposition-a} in
Lemma~\ref{4.22:lem-polar-decomposition}, namely
$\operatorname{pol}(e^{Z})\le\|Z\|\le|Z|$ for $Z\in\mathfrak{g}^{\mathbb{C}}$, we obtain
\begin{equation*}
\operatorname{pol}\bigl(e^{Z_s(t,m)}\bigr)\le\|Z_s(t,m)\|\le|Z_s(t,m)|\le\mu(s).
\end{equation*}

Now let $g'\in\mathfrak{G}_{\varepsilon'}(S)$, so that
$\operatorname{pol}(g'(s))<\varepsilon'(s)$ for every $s$. At the site $s$ we have
$u(t,m)(s)g'(s)=u(0,m)(s)\,e^{Z_s(t,m)}\,g'(s)$. The sub-additivity
$\operatorname{pol}(g_1g_2)\le\operatorname{pol}(g_1)+\operatorname{pol}(g_2)$ of
\eqref{4.22:eq-polar-decomposition-a}, applied twice, gives
\begin{align*}
\operatorname{pol}\bigl(u(t,m)(s)g'(s)\bigr)
&\le\operatorname{pol}\bigl(u(0,m)(s)\bigr)+\operatorname{pol}\bigl(e^{Z_s(t,m)}\bigr)
+\operatorname{pol}\bigl(g'(s)\bigr)\\
&<\eta_0(s)+\mu(s)+\varepsilon'(s)=\varepsilon(s).
\end{align*}
This holds at every site $s\in\mathfrak{s}(S)$ and for every
$(t,m)\in\mathfrak{P}\times\mathfrak{M}$. Hence the map
\begin{equation*}
\Psi:(t,m,g')\longmapsto\bigl(t,m,u(t,m)g'\bigr)
\end{equation*}
sends $\mathfrak{P}\times\mathfrak{M}\times\mathfrak{G}_{\varepsilon'}(S)$ into
$\mathfrak{P}\times\mathfrak{M}\times\mathfrak{G}_{\varepsilon}(S)$, the domain of $f$. No
restriction on $t$ beyond $t\in\mathfrak{P}$ is used.

The map $\Psi$ is holomorphic. Indeed, $u$ is holomorphic by hypothesis, and the site-wise matrix
product $(u,g')\mapsto ug'$ is polynomial in the entries. Moreover $\mathfrak{G}_{\varepsilon'}(S)$
is open in $(G^{\mathbb{C}})^{\mathfrak{s}(S)}$ by Lemma~\ref{4.22:lem-polar-decomposition}.
Therefore $F=f\circ\Psi$ is holomorphic on
$\mathfrak{P}\times\mathfrak{M}\times\mathfrak{G}_{\varepsilon'}(S)$, and
$|F|\le\sup|f|\le\mathfrak{N}$.

Finally fix $(m,g')\in\mathfrak{M}\times\mathfrak{G}_{\varepsilon'}(S)$. The function
$t\mapsto F(t,m,g')$ is holomorphic on the same polydisc $\mathfrak{P}$ on which $f(\cdot,m,g)$ is
holomorphic for every $g\in\mathfrak{G}_{\varepsilon}(S)$, and it is bounded by the same number
$\mathfrak{N}$. An estimate of Cauchy, M\"obius,
fundamental-theorem-of-calculus or Schwarz type on $\mathfrak{P}$ uses only these two facts and the
radii $\rho_{\square}$. It therefore holds for $F$ with the same constants and with $\mathfrak{N}$.
\end{proof}

\begin{remark}
Lemma~\ref{4.22:lem-cauchy-riding} (one disc, shrunk) and Lemma~\ref{4.22:lem-no-shrink} (a
polydisc, unshrunk) are used as follows. The shrunk one-block form is
Lemma~\ref{4.22:lem-cauchy-riding}
(Cauchy estimate with a riding rotation). It is used only at the fluctuation-integral site of a
step, with the site-wise rotation bound $\vartheta_p(s)$ furnished there by clause (iv) of the lemma
on the response, in its form at that site. At that site the Cauchy disc is
$|\mathsf{t}|<\mathsf{w}/\eta_p$, in the Cauchy variable $\mathsf{t}$ of the correlation theorem,
where $\mathsf{w}$ is the disc width and $\eta_p$ is the decay factor, so that the disc of width
$\mathsf{w}$ is dilated by the factor $\eta_p^{-1}$.
Under the condition which makes the constant $\Lambda^{\mathrm{res}}$ of the lemma on the response at
most one, one has only
$\vartheta_p(s)/\mu(s)\le\eta_p^{-1/2}\Lambda^{\mathrm{res}}$, and $\Lambda^{\mathrm{res}}$ contains
the factor $C_b^{1/2}$. Lemma~\ref{4.22:lem-cauchy-riding} then shrinks the disc of radius
$\mathsf{w}/\eta_p$ to a disc of radius $\rho'\ge\mathsf{w}\eta_p^{-1/2}$. This shrinking is the price
paid at the fluctuation-integral site, namely $\eta_p^{1/2}\max\{1,\Lambda^{\mathrm{res}}\}$. On the
full disc the
rotation at the near end of the domain $S_v$ of the kept member $v$ is of the order of
$\vartheta^{T}$, which is much larger than $\mu$; see \eqref{4.22:eq-margin-age-factor-a}.

On every other disc or polydisc, the statement that supplies the rotation exhibits the site-wise
supremum in \eqref{4.22:eq-no-shrink-1} as bounded by its share of $\mu(s)$. This is the case for
the rotations led by the correlation theorem, for the rotations with the site-wise profile
$\vartheta^{3F}(s)$, which the lemma leading them exhibits site by site, and, by clause (iii) of the
corollary to the one-run
expansions on labels, for every family running over all blocks with rate at most $1/2$.
There Lemma~\ref{4.22:lem-no-shrink} applies with no shrinking, that is, with factor one. No
alternative is available there: a shrinking factor $\min_s\mu(s)/\vartheta(s)$ as in
Lemma~\ref{4.22:lem-cauchy-riding}, with $\vartheta(s)$ the rotation bound of that lemma, at the law
sites $s\in\Sigma_W$ of a frame $W$ would read the coupling $g_{k_W}$, once per block.
\end{remark}

\begin{lemma}[The joint polydisc with shares of the margin]\label{4.22:lem-joint-polydisc}
Let $\mathfrak{P}=\mathfrak{P}_1\times\cdots\times\mathfrak{P}_n$, each $\mathfrak{P}_i$ a finite
polydisc, and let $\mathfrak{M}$, $f$, $\mathfrak{N}$ and $\varepsilon$ be as in
Lemma~\ref{4.22:lem-no-shrink}. Let
\begin{equation*}
u\colon \mathfrak{P}\times\mathfrak{M}\to\bigl(G^{\mathbb{C}}\bigr)^{\mathfrak{s}(S)}
\end{equation*}
be holomorphic, with $\operatorname{pol}(u(0,m)(s))\le\eta_0(s)$ for all $m\in\mathfrak{M}$ and
$s\in\mathfrak{s}(S)$. Let $\theta_1,\dots,\theta_n\ge0$ be numbers and put
$\Theta:=\sum_{i=1}^n\theta_i$. For $\tau=(\tau_1,\dots,\tau_n)\in\mathfrak{P}$ and
$0\le i\le n$ write $\tau^{(i)}:=(\tau_1,\dots,\tau_i,0,\dots,0)$, and assume that for every
$i\in\{1,\dots,n\}$, every $(\tau,m)\in\mathfrak{P}\times\mathfrak{M}$ and every
$s\in\mathfrak{s}(S)$
\begin{equation}\label{4.22:eq-joint-polydisc-1}
\bigl|\log\bigl[u(\tau^{(i-1)},m)(s)^{-1}\,u(\tau^{(i)},m)(s)\bigr]\bigr|\le\theta_i\,\mu(s).
\end{equation}
If $\varepsilon':=\varepsilon-\eta_0-\Theta\mu>0$, these quantities being read at each law site
$s$ as in Lemma~\ref{4.22:lem-no-shrink}, then the conclusion of
Lemma~\ref{4.22:lem-no-shrink} holds on the whole of
$\mathfrak{P}\times\mathfrak{M}\times\mathfrak{G}_{\varepsilon'}(S)$: the function
$F(\tau,m,g'):=f(\tau,m,u(\tau,m)g')$ is holomorphic there, $|F|\le\mathfrak{N}$, and every
Cauchy, M\"obius, fundamental-theorem-of-calculus or Schwarz estimate taken on $\mathfrak{P}$
at the site of the step holds for $F$ with that site's constants and with $\mathfrak{N}$.
\end{lemma}

\begin{proof}
Fix $(\tau,m)\in\mathfrak{P}\times\mathfrak{M}$ and a law site $s\in\mathfrak{s}(S)$, and
abbreviate $u(\sigma):=u(\sigma,m)(s)$. Each $\mathfrak{P}_j$ is a polydisc about the origin,
so every $\tau^{(i)}$ lies in $\mathfrak{P}$; moreover $\tau^{(0)}=0$ and $\tau^{(n)}=\tau$.
Hence the product
\begin{equation}\label{4.22:eq-joint-polydisc-2}
u(\tau)=u(0)\cdot\prod_{i=1}^{n}\bigl[u(\tau^{(i-1)})^{-1}u(\tau^{(i)})\bigr],
\end{equation}
with the factors taken in the order $i=1,\dots,n$, telescopes. Let $Z_i\in\mathfrak{g}^{\mathbb{C}}$
be the logarithm in \eqref{4.22:eq-joint-polydisc-1}, so that the $i$-th factor in
\eqref{4.22:eq-joint-polydisc-2} equals $e^{Z_i}$. By the properties
\eqref{4.22:eq-polar-decomposition-a} of Lemma~\ref{4.22:lem-polar-decomposition} and by
\eqref{4.22:eq-joint-polydisc-1},
\begin{equation*}
\operatorname{pol}\bigl(e^{Z_i}\bigr)\le\|Z_i\|\le|Z_i|\le\theta_i\,\mu(s).
\end{equation*}
Let $g'\in\mathfrak{G}_{\varepsilon'}(S)$. The sub-additivity of $\operatorname{pol}$ in
\eqref{4.22:eq-polar-decomposition-a}, applied to the product
\eqref{4.22:eq-joint-polydisc-2} followed by $g'(s)$, together with
$\operatorname{pol}(u(0))\le\eta_0(s)$, gives
\begin{align*}
\operatorname{pol}\bigl(u(\tau)g'(s)\bigr)
&\le\operatorname{pol}(u(0))+\sum_{i=1}^n\operatorname{pol}\bigl(e^{Z_i}\bigr)
+\operatorname{pol}(g'(s))\\
&\le\eta_0(s)+\Theta\,\mu(s)+\operatorname{pol}(g'(s))<\varepsilon,
\end{align*}
the last inequality because $\operatorname{pol}(g'(s))<\varepsilon'=\varepsilon-\eta_0-\Theta\mu$
at $s$. As $s$ was arbitrary, $u(\tau,m)g'\in\mathfrak{G}_{\varepsilon}(S)$.

Thus the map $(\tau,m,g')\mapsto(\tau,m,u(\tau,m)g')$ sends
$\mathfrak{P}\times\mathfrak{M}\times\mathfrak{G}_{\varepsilon'}(S)$ into
$\mathfrak{P}\times\mathfrak{M}\times\mathfrak{G}_{\varepsilon}(S)$, the domain of $f$; it is
holomorphic, since $u$ is holomorphic and the product $u(\tau,m)(s)g'(s)$ is polynomial in
the matrix entries. Composing with $f$, the function $F$ is holomorphic on the whole of
$\mathfrak{P}\times\mathfrak{M}\times\mathfrak{G}_{\varepsilon'}(S)$ and $|F|\le\mathfrak{N}$
there. The remaining part of the conclusion of Lemma~\ref{4.22:lem-no-shrink} follows from
these two facts exactly as it does in that lemma.

No Baker--Campbell--Hausdorff expansion enters: $\operatorname{pol}$ is sub-additive, and the
logarithm of each factor of \eqref{4.22:eq-joint-polydisc-2} is controlled by hypothesis.
We note that the hypothesis \eqref{4.22:eq-joint-polydisc-1} on the $i$-th factor is required
at every value of the earlier coordinates $\tau_1,\dots,\tau_{i-1}$, that is, on the whole
product $\mathfrak{P}$.
\end{proof}

In what follows the indexing is that of the step from level $k$ to level $k+1$ in \cite[\S3]{B-CMP119}. We consider one output term of the fluctuation integral, in the flat form used by the correlation theorem. The new slot is $x=A_k\restriction\mathfrak{r}$, and every coupling-dependent letter is evaluated at $g_k$. The letters $\mathfrak{b}_k$, $\varrho$, $c_\eta$, $\Theta_\flat$, $\tau_\flat$, $\varpi_\gamma$, $T(\sigma)$, $\mathcal{S}(P)$, $\mathsf{N}(b)$, $\mathfrak{X}^{\sharp}$, the regions $\Lambda$, $\bar\Lambda$, $Z_0$, $Z$ and the floors (F1), (F2) are those listed among the constants inherited from the correlation theorem (Definition~\ref{4.22:not-inherited-constants}).

We first describe how $x$ enters. The integral over $\Lambda$ in \cite[(3.25)]{B-CMP119} is a factor independent of $x$ times
\begin{equation*}
\mathsf{I}(x)=\int d\mu_{C_\Lambda}(\psi)\prod_{b\in\bar\Lambda}
\chi\bigl(|\psi_b+\mathsf{m}_b(x)|<\mathfrak{b}_k\bigr)\,
\exp\sum_{Y}V\bigl(Y;x\oplus(\psi+\mathsf{m}(x))\bigr).
\end{equation*}
Here $\mu_{C_\Lambda}$ is the Gaussian measure with covariance $C_\Lambda$, $\chi(\cdot)$ is the indicator of the condition displayed in its argument, and $V(Y;\cdot)$ are the localised potentials.

Now consider the one-term formula $\mathcal{T}$ for this integral. The following ingredients of $\mathcal{T}$ do not depend on $x$: the Gaussian weight, the determinant, the factors denoted $\mathsf{F}$ and $\Gamma$ in that formula, the factors $h(|\psi_b|)$, the set $\mathcal{S}(P)$ and the numbers $\mathsf{N}(b)$. The slot $x$ occurs only in the four places
\begin{equation}\label{4.22:eq-rotated-mean-field-a}
\begin{aligned}
&\text{(e1)}\quad \mathsf{m}=T(\sigma)x\restriction Z_0,
&&\text{(e2)}\quad x\ \text{in the argument of}\ V(Y;x\oplus(\psi+\mathsf{m})),\\
&\text{(e3)}\quad V(Y;x),\ \ Y\cap\bar\Lambda=\emptyset,
&&\text{(e4)}\quad \langle x_c,S_{cc'}x_{c'}\rangle .
\end{aligned}
\end{equation}
In (e1), $x$ enters through the cut-off factors $\mathsf{L}_b$, $\mathsf{L}^c_b$ evaluated at $\mathsf{m}_b$, and through the argument $\psi+\mathsf{m}$ of $V$. Places (e3) and (e4) lie outside the integral: (e3) consists of the potentials that do not depend on $\psi$, and (e4) of the remnants, which are polynomials in $x$. The rim factors, in the flat form of \cite[(3.21)]{B-CMP119}, belong to $\mathbf{T}^{(k+1)}$ and not to a stored term.

Let $g$ be a map from the law sites of $\mathfrak{r}$ to $G^{\mathbb{C}}$. We put
\begin{equation*}
x^g:=x^{[g]},\qquad \mathsf{m}^g:=T(\sigma)x^g\restriction Z_0 .
\end{equation*}
For $A\in\mathfrak{g}^{\mathbb{C}}$ we write $\operatorname{Re}A$, $\operatorname{Im}A$ for the components of $A$ in the decomposition $\mathfrak{g}^{\mathbb{C}}=\mathfrak{g}\oplus i\mathfrak{g}$, so that $A=\operatorname{Re}A+i\operatorname{Im}A$. Since $|A^*|=|A|$, both components have norm at most $|A|$.

\begin{lemma}[The rotated mean field stays in the widened box]\label{4.22:lem-rotated-mean-field}
Assume the floors \textup{(F1)}, \textup{(F2)}. Let $P\subset Z_0\cap\bar\Lambda$. Let $x$ be as in clause \textup{(ii)} of the flat mean-field lemma of the correlation theorem, that is:
\begin{itemize}
\item $x$ is real, with values in $\mathfrak{X}^{\sharp}$, on the cubes of $\mathcal{S}(P)$ (the strong components of $x$);
\item $x$ is complex with $|x|<\varrho$ elsewhere (the weak components of $x$);
\item $x$ is real on the silent bonds (in the sense of the flat mean-field lemma).
\end{itemize}
Let $\operatorname{pol}(g)<\varepsilon_\star$ at every law site of $\mathfrak{r}$. Then $|x^g|<e^{2\varepsilon_\star}\varrho$, and $\mathsf{m}^g$ lies in the set
\begin{equation}\label{4.22:eq-rotated-mean-field-b}
\begin{aligned}
\mathfrak{M}_P^{+}:=\Bigl\{\mathsf{m}:\ &|\mathsf{m}_b|\le e^{2\varepsilon_\star}\Theta_\flat\varrho
\ \text{for all } b;\\
&|\operatorname{Re}\mathsf{m}_b|\le\tfrac{7}{16}\mathfrak{b}_k,\ \
|\operatorname{Im}\mathsf{m}_b|\le\tfrac{6}{5}c_\eta\mathfrak{b}_k\ \text{for } b\in P\Bigr\}.
\end{aligned}
\end{equation}
Moreover, $\mathsf{m}^g$ is holomorphic in the weak components of $x$ and in $g$, $\sigma$, $\mathbf{U}$, $\mathbf{J}$. If $\sigma=0$ off $Z$, then each component $\mathsf{m}^g_b$ depends only on the restrictions of $x$, $g$, $\mathbf{U}$, $\mathbf{J}$ to $Z$.
\end{lemma}

\begin{proof}
\emph{Splitting off the unitary part.} By the polar decomposition \eqref{4.22:eq-polar-decomposition-a}, we write $g=we^{Y_g}$ at each law site. Here $w$ is $G$-valued, $Y_g$ is $\mathfrak{g}$-valued (it is the exponent written $Y$ in that lemma; we write $Y_g$ because $Y$ denotes polymers here), and $\operatorname{pol}(g)=\|Y_g\|$.

Put $x'':=x^{[w]}$. Then $x''$ again satisfies the hypotheses of clause (ii) of the flat mean-field lemma, for two reasons: the balls $\{|z|<\varrho\}$ are invariant under $\operatorname{Ad}_G$, and the rotation by the $G$-valued $w$ takes real values to real values. That lemma therefore gives $\mathsf{m}'':=T(\sigma)x''\in\mathfrak{M}_P$, where $\mathfrak{M}_P$ is the mean-field box of the flat mean-field lemma, whose conditions at $b\in P$ include $|\operatorname{Im}\mathsf{m}_b|\le c_\eta\mathfrak{b}_k$. Its last assertion gives, for $b\in P$,
\begin{equation*}
|\operatorname{Re}\mathsf{m}''_b|\le\bigl(\tfrac5{16}+c_\eta\bigr)\mathfrak{b}_k .
\end{equation*}

\emph{The non-unitary part.} Apply \eqref{4.22:eq-polar-decomposition-b} at each bond $b'$, with the values of $g$ and $w$ at its law site. This gives $x^g=x''+\delta$ with
\begin{equation*}
|\delta(b')|\le\bigl(e^{2\operatorname{pol}(g)}-1\bigr)|x(b')|
\le\bigl(e^{2\varepsilon_\star}-1\bigr)|x(b')| .
\end{equation*}
Since $T(\sigma)$ is linear, $\mathsf{m}^g=\mathsf{m}''+T(\sigma)\delta\restriction Z_0$.

\emph{Estimate of $T(\sigma)\delta$ at $b\in P$.} Fix $b\in P$ and split the bonds $b'$ into two classes.
\begin{itemize}
\item Bonds of the cubes of $\mathcal{S}(P)$ (the strong cubes). On these $|x|\le\mathfrak{b}_k$, since none of them is far in the sense of the flat mean-field lemma. The row sum of $T(\sigma)$ is at most $\Theta_\flat$. Their contribution is therefore at most $(e^{2\varepsilon_\star}-1)\Theta_\flat\mathfrak{b}_k$.
\item All other bonds of the rim $\mathfrak{r}$. These lie at distance beyond $\varpi_\gamma$ from $b$, and on them $|x|\le\varrho$. Their contribution is at most
\begin{equation*}
\bigl(e^{2\varepsilon_\star}-1\bigr)\tau_\flat(\varpi_\gamma)\varrho
\le\bigl(e^{2\varepsilon_\star}-1\bigr)\cdot\tfrac12c_\eta\mathfrak{b}_k .
\end{equation*}
\end{itemize}
Adding the two contributions and using $\tfrac12c_\eta\le1$, we get
\begin{equation*}
|T(\sigma)\delta|_b\le\bigl(e^{2\varepsilon_\star}-1\bigr)(\Theta_\flat+1)\mathfrak{b}_k
=\tfrac15c_\eta\mathfrak{b}_k .
\end{equation*}
The equality is the definition \eqref{4.22:eq-rotation-radius-a} of the rotation radius, which says that $(e^{2\varepsilon_\star}-1)(\Theta_\flat+1)=c_\eta/5$.

\emph{Imaginary and real parts.} Membership of $\mathsf{m}''$ in $\mathfrak{M}_P$ gives $|\operatorname{Im}\mathsf{m}''_b|\le c_\eta\mathfrak{b}_k$ for $b\in P$. Using $|\operatorname{Im}A|,|\operatorname{Re}A|\le|A|$ for $A=(T(\sigma)\delta)_b$, we obtain $|\operatorname{Im}\mathsf{m}^g_b|\le(1+\tfrac15)c_\eta\mathfrak{b}_k$. For the real part,
\begin{equation*}
|\operatorname{Re}\mathsf{m}^g_b|\le\bigl(\tfrac5{16}+\tfrac65c_\eta\bigr)\mathfrak{b}_k
\le\tfrac7{16}\mathfrak{b}_k .
\end{equation*}
The last inequality holds because $c_\eta\le\tfrac1{16}$ (Definition~\ref{4.22:not-inherited-constants}) gives $\tfrac65c_\eta\le\tfrac{6/5}{16}<\tfrac2{16}$.

\emph{First clause.} By \eqref{4.22:eq-polar-decomposition-b}, $|x^g(b')|\le e^{2\operatorname{pol}(g)}|x(b')|$. On the weak components, $|x|<\varrho$ and $\operatorname{pol}(g)<\varepsilon_\star$ give $|x^g|<e^{2\varepsilon_\star}\varrho$. Multiplying the supremum of $|x^g|$ by the row-sum bound $\Theta_\flat$ of $T(\sigma)$ gives $|\mathsf{m}^g_b|\le e^{2\varepsilon_\star}\Theta_\flat\varrho$ for every $b$. Hence $\mathsf{m}^g\in\mathfrak{M}_P^{+}$.

\emph{Holomorphy.} The operator $T(\sigma;\mathbf{U},\mathbf{J})$ is holomorphic in $(\sigma,\mathbf{U},\mathbf{J})$ by its random-walk representation at zeroth order. Moreover, $x^g$ is linear in $x$ and polynomial in $g^{\pm1}$. Hence $\mathsf{m}^g$ is holomorphic in the weak components of $x$ and in $g$, $\sigma$, $\mathbf{U}$, $\mathbf{J}$.

\emph{Locality.} This is clause (iii) of the flat mean-field lemma.
\end{proof}

No constant is changed in this lemma. The floor (F1) and the choice of $\varpi_\gamma$ keep the split $\tfrac12c_\eta+\tfrac12c_\eta$, and the rotation is absorbed by the widening $\mathfrak{M}_P\subset\mathfrak{M}_P^{+}$. The set of admitted data is a product of three factors: the weak components of $x$ range over balls of radius $\varrho$, the strong components range over real values, and $g$ ranges over rotations of polar size less than $\varepsilon_\star$. In particular, the mixed real and complex point used in the subsequent analytic-continuation argument lies in this product.

\begin{lemma}[The one-term bound on the widened box]\label{4.22:lem-one-term-bound}
Assume the hypotheses of the one-term lemma, including its floors \textup{(F1)}--\textup{(F3)},
and assume moreover the floor \textup{(F2)}$^{+}$:
\begin{equation}\label{4.22:eq-one-term-bound-a}
\mathfrak{b}_k\ge 64 .
\end{equation}
Let $P\subset Z_0\cap\bar{\Lambda}$ be as in Lemma~\ref{4.22:lem-rotated-mean-field}, and let
$\mathsf{m}\in\mathfrak{M}_P^{+}$, the widened mean-field box of
\eqref{4.22:eq-rotated-mean-field-b}. Let $\mathsf{X}$, the argument $x$ of the potentials $V$, be
complex, with $|\mathsf{X}|<e^{2\varepsilon_{\star}}\varrho$ on the bonds outside the set
$\mathfrak{S}^{22}$ of bonds of the change of variables \cite[(1.22)]{B-CMP119}. Let $\mathcal{T}$
be the term given by the one-term formula of the one-term lemma. Then the following hold; the
clauses carry the letters of the corresponding clauses of the one-term lemma.
\begin{itemize}
\item[(a)] $\mathcal{T}$ is holomorphic in
$(\mathsf{m},\mathsf{X},\sigma,\mathbf{U},\mathbf{J},w)$ and in the parameters.
\item[(c)] With $D$, $\nu(Y)$, $c_F$ and $Z$ as in the one-term lemma,
\begin{equation}\label{4.22:eq-one-term-bound-1}
|\mathcal{T}|\le \prod_{Y\in D}2\nu(Y)\cdot e^{-3\mathsf{p}|P|}\cdot e^{2c_F|Z|}.
\end{equation}
This is clause (c) of the one-term lemma, unchanged.
\end{itemize}
\end{lemma}

\begin{proof}
The proof of the one-term lemma uses the mean field $\mathsf{m}$ at three places. We show
that each of them remains valid for $\mathsf{m}\in\mathfrak{M}_P^{+}$ and $\mathsf{X}$ as in the
statement.

\emph{First place.} The proof of the one-term lemma applies the auxiliary shift lemma at
$\mu=\mathsf{m}$. The hypothesis of that lemma is the bound $|\mu_r|\le 7\mathfrak{b}_k/16$ on the
real part $\mu_r$ of the shift. At $\mu=\mathsf{m}$ it is the bound
$|\operatorname{Re}\mathsf{m}_b|\le 7\mathfrak{b}_k/16$ for $b\in P$, which is part of the
definition of $\mathfrak{M}_P^{+}$ in \eqref{4.22:eq-rotated-mean-field-b}.

\emph{Second place.} The arguments of $V$ must lie in the polydisc of radius $T^{\mathbb{C}}$ on
which the holomorphic extension $V^{\mathbb{C}}$ of $V$ is controlled. For
$\mathsf{m}\in\mathfrak{M}_P^{+}$ and $\mathsf{X}$ as in the statement, the moduli of these
arguments are bounded by
\begin{equation*}
\tfrac12\mathfrak{b}_k+e^{2\varepsilon_{\star}}\Theta_{\flat}\varrho,\qquad
\bigl(1+\tfrac65 c_{\eta}\bigr)\mathfrak{b}_k,\qquad
e^{2\varepsilon_{\star}}\varrho .
\end{equation*}
The distances of these bounds from $T^{\mathbb{C}}$ satisfy
\begin{equation}\label{4.22:eq-one-term-bound-b}
\begin{aligned}
T^{\mathbb{C}}-\bigl(\tfrac12\mathfrak{b}_k+e^{2\varepsilon_{\star}}\Theta_{\flat}\varrho\bigr)
&\ge \tfrac12\mathfrak{b}_k+\varrho\bigl(1-\tfrac{c_{\eta}}{5}\bigr),\\
T^{\mathbb{C}}-\bigl(1+\tfrac65 c_{\eta}\bigr)\mathfrak{b}_k&\ge 2\varrho-\mathfrak{b}_k,\\
T^{\mathbb{C}}-e^{2\varepsilon_{\star}}\varrho&\ge \mathfrak{b}_k .
\end{aligned}
\end{equation}
These inequalities use three facts:
\begin{equation*}
(e^{2\varepsilon_{\star}}-1)\Theta_{\flat}\le \frac{c_{\eta}}{5},\qquad \Theta_{\flat}\ge1,\qquad
e^{2\varepsilon_{\star}}<2 .
\end{equation*}
The first follows from the definition \eqref{4.22:eq-rotation-radius-a} of $\varepsilon_{\star}$.
That definition gives $e^{2\varepsilon_{\star}}-1=c_{\eta}/(5(\Theta_{\flat}+1))$, and
$\Theta_{\flat}/(\Theta_{\flat}+1)\le1$. Each of the three right-hand sides in
\eqref{4.22:eq-one-term-bound-b} is at least $2$ (recall $r_A\ge2$). Hence the bound
$|e^{V}-1|\le 2\nu(Y)$ of the proof of the one-term lemma holds on the present domain.

\emph{Third place.} The shift $\upsilon$ of the proof of the one-term lemma is supported on $P$. By
the definition \eqref{4.22:eq-rotated-mean-field-b} of $\mathfrak{M}_P^{+}$ it satisfies
\begin{equation*}
|\upsilon_b|\le|\operatorname{Im}\mathsf{m}_b|\le\tfrac65 c_{\eta}\mathfrak{b}_k\qquad(b\in P).
\end{equation*}
Let $c_0'$ and $\gamma_2$ be the constants of the Gaussian factor and of the weights, and $\xi_b$
the fluctuation variables, of the proof of the one-term lemma, and let
$\varsigma:=\frac12\mathfrak{b}_k-1$ be the weight threshold of that proof.
This shift enters the modulus of the Gaussian factor as the factor
$e^{(\frac12 c_0'+\frac12)|\upsilon|^2}$. It enters the weights as a factor at most
$(1+|\upsilon_b|)\mathbb{1}(|\xi_b|\ge \varsigma)$. On the support of the weight, the lower bound
on the fluctuation variables established in the proof of the one-term lemma gives
$|\xi_b|\ge 9\mathfrak{b}_k/16$, and $9\mathfrak{b}_k/16>\frac12\mathfrak{b}_k>\varsigma$. The
auxiliary estimate of the proof of the one-term lemma that acts on the factors not containing
$\mathsf{m}$ is applied to objects that do not depend on $\mathsf{m}$, and is applied unchanged.
Collecting the contributions of the Gaussian factor and of the weights, each bond of $P$
contributes to the bound a factor at most
\begin{equation}\label{4.22:eq-one-term-bound-2}
2\bigl(1+\tfrac65 c_{\eta}\mathfrak{b}_k\bigr)\,
\exp\Bigl[-\tfrac12\gamma_2\varsigma^2+\bigl(\tfrac12c_0'+\tfrac12\bigr)\tfrac{36}{25}c_{\eta}^2
\mathfrak{b}_k^2\Bigr].
\end{equation}
We bound the three ingredients of \eqref{4.22:eq-one-term-bound-2} in turn.

First, $\varsigma=\frac12\mathfrak{b}_k-1\ge\frac{31}{64}\mathfrak{b}_k$, since this inequality is
equivalent to $\mathfrak{b}_k\ge64$, which is \eqref{4.22:eq-one-term-bound-a}. Therefore
\begin{equation*}
\tfrac12\gamma_2\varsigma^2\ge\tfrac12\cdot\tfrac{961}{4096}\,\gamma_2\mathfrak{b}_k^2
=\tfrac{961}{8192}\,\gamma_2\mathfrak{b}_k^2\ge 7.5\,\mathsf{p}.
\end{equation*}

Second, $(\frac12c_0'+\frac12)c_{\eta}^2=\gamma_2/32$. Hence the second term in the exponent is
\begin{equation*}
\bigl(\tfrac12c_0'+\tfrac12\bigr)\tfrac{36}{25}c_{\eta}^2\mathfrak{b}_k^2
=\tfrac{9}{200}\,\gamma_2\mathfrak{b}_k^2=2.88\,\mathsf{p}.
\end{equation*}

Third, the prefactor satisfies $2(1+\frac65c_{\eta}\mathfrak{b}_k)\le\frac65e^{\mathsf{p}}$ by
the floor (F2).

Inserting these three bounds into \eqref{4.22:eq-one-term-bound-2}, the factor per bond of $P$ is
at most
\begin{equation*}
\tfrac65\,e^{\mathsf{p}}\,e^{-7.5\mathsf{p}+2.88\mathsf{p}}=\tfrac65\,e^{-3.62\mathsf{p}}
\le e^{-3\mathsf{p}} .
\end{equation*}
The last inequality holds because $\mathsf{p}\ge5$ by the floor (F3), so that
$e^{0.62\mathsf{p}}\ge e^{3.1}>\frac65$. The product over the bonds of $P$ is therefore at most
$e^{-3\mathsf{p}|P|}$. With the first two places treated as above and every other step of the
proof of clause (c) of the one-term lemma unchanged, this gives \eqref{4.22:eq-one-term-bound-1}.

\emph{Clause (a).} The integrand of the one-term formula is holomorphic in
$(\mathsf{m},\mathsf{X},\sigma,\mathbf{U},\mathbf{J},w)$ and the parameters, and is bounded,
uniformly in them on the domain of the statement, by the integrable majorant assembled in the
proof of (c). By dominated convergence, $\mathcal{T}$ is therefore continuous there. Fix all
variables but one, and integrate over a closed triangle in the remaining complex variable. By
Fubini's theorem, applicable because of the uniform majorant, the order of integration may be
exchanged, and the inner integral vanishes by Cauchy's theorem. Morera's theorem
then shows that $\mathcal{T}$ is holomorphic in each variable separately, hence holomorphic.

The integrals over the cut-off region enter with $\mu$ a component of $\mathsf{m}$. They have the
form
\begin{equation*}
\mathsf{T}f(\mu):=\int_{|A|<\mathfrak{b}_k}f(A-\mu)\,dA .
\end{equation*}
For holomorphic $f$ such an integral is holomorphic in $\mu$: the domain of integration does not
depend on $\mu$, and the integrand is holomorphic in $\mu$.
\end{proof}

\begin{remark}
The proof of the one-term lemma used $\varsigma^2\ge3\mathfrak{b}_k^2/16$, valid for
$\mathfrak{b}_k\ge16$, which gives $6\mathsf{p}$ in the exponent. The actual value
$\frac12\gamma_2\varsigma^2=8\mathsf{p}(1-O(1/\mathfrak{b}_k))$ pays for the widening of the
mean-field box to $\mathfrak{M}_P^{+}$. The margin used is the gap between $6\mathsf{p}$ and
$8\mathsf{p}$; the spare third $\mathsf{p}$ kept in reserve in the proof of the flat form of the
one-term lemma is not drawn on here, and is spent later in the argument.

Clause (a) of the one-term lemma asserts only continuity of $\mathcal{T}$ in the strong components
of $x$. This understates what holds: the strong components of $x$ enter $\mathcal{T}$ only through
$(\mathsf{m},\mathsf{X})$, in which $\mathcal{T}$ is holomorphic by clause (a) of
Lemma~\ref{4.22:lem-one-term-bound}.
\end{remark}

\begin{theorem}[Orbit data for terms created by the fluctuation integral]
\label{4.22:thm-fluctuation-orbit-data}
We work at step $k+1$ of Ba{\l}aban's construction, with one output term of the fluctuation
integral, the new slot $x=A_k\restriction\mathfrak{r}$ and all couplings read at $g_k$. For a
set $X$ and a family $\mathcal{S}$ of strong sets of the new slot, let $\mathfrak{L}(X,\mathcal{S})$
denote the polymer activity of the fluctuation-integral polymer theorem, and write
$\mathfrak{L}(\cdot\,;x)$ for its dependence on the slot. Assume the following.
\begin{itemize}
\item The fluctuation-integral polymer theorem, with its hypotheses and with the identities and
estimates of the five steps of its proof; the interpolated-term bound established in the third of
those steps; and the summation lemma of the fluctuation step.
\item Hypothesis \eqref{4.22:eq-one-term-bound-a}, the invariant cut-off \eqref{4.22:eq-invariant-cutoff-a}
and the properties \eqref{4.22:eq-lifted-cutoff-a} of the lifted cut-off terms.
\item The stored terms of the older slots have orbit data
(Definition~\ref{4.22:def-orbit-datum}) of radius $\varepsilon$.
\end{itemize}
Then every $\mathfrak{L}(X,\mathcal{S})$ has an orbit datum $\mathfrak{L}(X,\mathcal{S})^{\sharp}$, of
radius $\varepsilon_{\star}$ at the coordinates of $x$ and of radius
$\varepsilon-\eta^{T}-\mu\,\mathsf{e}/\varepsilon_{\star}$ at the older law sites, which is local in
$X$ and satisfies
\begin{equation}\label{4.22:eq-fluctuation-orbit-data-1}
\sum_{\mathcal{S}} e^{\kappa_A^{\mathrm{pr}}(g_k)|\mathcal{S}|}\,
\mathfrak{N}^{\mathbb{C}}\bigl(\mathfrak{L}(X,\mathcal{S})\bigr)
\le \tfrac18 B_0\, e^{-(1+4\beta_s)\kappa\,\mathsf{d}_{k+1}} .
\end{equation}
Here the sum runs over the families $\mathcal{S}$ of strong sets of the new slot; the right member is
that of the fluctuation-integral polymer theorem in the present letters, with $B_0$, $\beta_s$,
$\kappa$ and the length $\mathsf{d}_{k+1}$ of the term measured at level $k+1$ as in its third
clause, whose output rate
satisfies $\kappa_{\mathrm{out}}\ge(1+4\beta_s)\kappa$; the weights of the strong sets of older slots
are carried by the bounds $\nu$ at the rate $\kappa_A(k)$. The same holds, with the sizes given by
the summation lemma of the fluctuation step, for
$\mathbf{B}^{\mathrm{sl}}:=\mathfrak{L}(\cdot\,;x)-\mathfrak{L}(\cdot\,;0)$, for the terms free of
the fluctuation field $\psi$, and for the term $\mathbf{B}^{(k+1)}$ of the third item of that lemma.
\end{theorem}

\begin{proof}
Write $x^{g}$ and $\mathsf{m}^{g}$ for the rotated slot and the rotated mean field of
Lemma~\ref{4.22:lem-rotated-mean-field}. For $g$ of polar size below $\varepsilon_{\star}$ at the law
sites of $\mathfrak{r}$ and below the stated radius at the older law sites (this set of rotations is
the chart of Definition~\ref{4.22:def-orbit-datum}), define
$\mathfrak{L}(X,\mathcal{S})^{\sharp}(\cdot\,;x,g)$ by the series of the fluctuation-integral polymer
theorem for $\mathfrak{L}(X,\mathcal{S})$, in which every one-term expression $\mathcal{T}$ is
evaluated at $(\mathsf{m},\mathsf{X}):=(\mathsf{m}^{g},x^{g})$ and every potential $V$ is replaced by
its orbit datum $V^{\sharp}$: at the older law sites $V^{\sharp}$ is supplied by the corollary giving
orbit data for the older potentials, and at the bonds of $\mathfrak{S}^{22}$, those of Ba{\l}aban's
full change of variables, by \eqref{4.22:eq-invariant-cutoff-a} and \eqref{4.22:eq-lifted-cutoff-a}.

\emph{The first two steps.} The first two steps of the proof of the fluctuation-integral polymer
theorem consist of the Mayer expansion; the splittings $\chi_b=h+\rho_b$ and $\chi_b=1-\chi^{c}_b$
of the cut-off factor $\chi_b$ at each bond $b$; the grouping of the terms according to their
region $Z_0$; the conditioning;
the interpolation of the operators and of the mean field $\mathsf{m}$ together; and the
factorisation at interpolation parameter $s=0$ away from the region $Z$ carried by the term. These
are identities between integrals
over real fields. They are used here only for $G$-valued $g$. For such $g$ the rotated slot $x^{g}$ is
an admissible real datum on the strong cubes, and the identities are those of the fluctuation-integral
polymer theorem applied at $x^{g}$; this gives the relation \eqref{4.22:eq-orbit-datum-b}, term by
term. At complex $g$ only the formula of each term and its bound are used.

\emph{The third step.} Let $g$ be complex. By Lemma~\ref{4.22:lem-rotated-mean-field},
$|x^{g}|<e^{2\varepsilon_{\star}}\varrho$ and $\mathsf{m}^{g}$ lies in the widened box
$\mathfrak{M}_P^{+}$ of \eqref{4.22:eq-rotated-mean-field-b}, where $P\subset Z_0\cap\bar{\Lambda}$ is
the set of bonds of the term as in that lemma. Hence the one-term bound on the widened
box (Lemma~\ref{4.22:lem-one-term-bound}) applies, under \eqref{4.22:eq-one-term-bound-a}, with
$(\mathsf{m},\mathsf{X})=(\mathsf{m}^{g},x^{g})$, and its clause (c) gives the one-term bound used in
the third step of that proof, unchanged. The Cauchy estimate on the circles
$|\sigma(\Delta)|=e^{\kappa_1}$ in the interpolation parameters then yields the interpolated-term bound
of the third step word for word.

\emph{The fourth and fifth steps.} These steps see a term only through the interpolated-term bound,
and the sets $\mathcal{S}(P)$, the counting numbers $N_A$ and $N_Z$ of those steps, and the weights
do not depend on $(x,g)$.
They therefore run unchanged on the chart: the polymer series defining $\mathfrak{L}(X,\mathcal{S})^{\sharp}$
converge uniformly on the chart under the majorants of the fluctuation-integral polymer theorem, and
the bound of that theorem holds for the sup over the chart, which is
\eqref{4.22:eq-fluctuation-orbit-data-1}. Each term is holomorphic by clause (a) of
Lemma~\ref{4.22:lem-one-term-bound}, composed with the holomorphic dependence of $\mathsf{m}^{g}$ on
the weak coordinates of $x$ and on $(g,\sigma,\mathbf{U},\mathbf{J})$ given by
Lemma~\ref{4.22:lem-rotated-mean-field}, with the dependence through $\mathsf{X}=x^{g}$, which is
linear in $x$ and polynomial in $g^{\pm1}$, and with the dependence through the orbit data
$V^{\sharp}$, which are holomorphic on their charts; a uniform limit
of holomorphic functions is holomorphic by the Weierstrass theorem. This gives
\eqref{4.22:eq-orbit-datum-a}.

\emph{The slot difference.} Since the rotation of the zero slot is zero, $0^{[g]}=0$, the datum
constructed at $x$ minus the datum constructed at the slot $0$ is an orbit datum for
$\mathbf{B}^{\mathrm{sl}}$.

\emph{The terms free of $\psi$.} Among the four ways in which $x$ enters, listed in
\eqref{4.22:eq-rotated-mean-field-a}, two lie outside the integral. The first consists of the potentials
$V(Y;x)$ with $Y\cap\bar{\Lambda}=\emptyset$: their datum is $V(Y;x^{g})$, which is defined since
$|x^{g}|<T^{\mathbb{C}}$, and it is bounded by $\nu(Y)$. The second consists of the remnants
$\langle x_c,S_{cc'}x_{c'}\rangle$, with $S_{cc'}$ the remnant kernel between the cubes $c$, $c'$: their
datum $\langle x^{g}_c,S_{cc'}x^{g}_{c'}\rangle$ is a polynomial in $(x,g^{\pm1})$, and, with $C$
and $e^{-\delta_1 M R_{k+1}}$ the constant and the decay factor of the remnant bound in the
fluctuation-integral polymer theorem,
\begin{equation*}
\bigl|\langle x^{g}_c,S_{cc'}x^{g}_{c'}\rangle\bigr|
\le C\,(e^{2\varepsilon_{\star}}\varrho)^{2}\,e^{-\delta_1 M R_{k+1}} ,
\end{equation*}
which is at most four times the bound of the remnant in the fluctuation-integral polymer theorem; by
the third clause of the ceiling on the rotation sizes fixed in the order of choice, this stays inside
the same allowance $\tfrac14 B_0$.

\emph{The term of the summation lemma.} The first item of the summation lemma of the fluctuation step,
the run with $A_k=0$, reads only the older law sites. There the orbit data are supplied by the corollary
giving orbit data for the older potentials and by the summation lemma of the correlation theorem, which
gives analyticity in all parameters.
\end{proof}

Ba{\l}aban's change of variables is made in \cite[(1.22)--(1.24)]{B-CMP119}, in the following words:
\begin{quotation}
Now we make a change of variables for each bond variable $A(b)$, $b\in S_1^{*}$. There are two
forms \dots\ For $b\in(R_1^{\sim2}\cap S_1)^{*}$, we take
\begin{equation*}
\begin{split}
A(b)&=\frac{1}{ig_0}\log \exp ig_0A'(b)\exp\bigl(-ig(A'(b))\mathbf{H}_{1,Ax}(b)\bigr)\\
&=A'(b)-\frac{1}{g_0}g(A'(b))\mathbf{H}_{1,Ax}(b)
-\frac{1}{g_0}\mathbf{F}\bigl(g_0A'(b),-g(A'(b))\mathbf{H}_{1,Ax}(b)\bigr).
\end{split}
\tag*{(1.22)}
\end{equation*}
Here $g$ is a $C^{\infty}$-function defined on the Lie algebra $\mathfrak{g}$, $0\le g(A')\le 1$,
$g(A')=0$ on $\{A'\in\mathfrak{g}:|A'|\le\frac{4}{3}g_0^{-1}\delta_0\}$, $g(A')=1$ on
$\{A'\in\mathfrak{g}:|\exp ig_0A'-1|\ge\frac{5}{3}\delta_0\}$ \dots\ The function $\mathbf{F}$ is
defined as a sum of terms of orders higher than 1 in the Baker-Campbell-Haussdorf formula \dots\
For $b\in S_1^{*}\setminus(R_1^{\sim2})^{*}$ we make a simpler change of variables; we take
(1.22) with the function $g(A'(b))$ replaced by 1 \dots\ The Jacobians of the transformations
(1.22) can be exponentiated \dots\ (1.24) Expanding the terms of the effective action \dots\ with
respect to the last two perturbative terms in (1.22) up to the first order, we obtain a new
expression $V^{(0)}(S_1,A,\mathbf{H}_{1,Ax})$ \dots\ a sum over bonds $b\in S_1$ of the
expressions which depend on $A$, $\mathbf{H}_{1,Ax}$ almost locally. The dependence on
$\mathbf{H}_{1,Ax}$ is analytic, and the expressions are small, i.e., they can be bounded by any
positive power of $g_0$.
\end{quotation}
Of $\mathbf{H}_{1,Ax}$, \cite{B-CMP119} says that it
\begin{quotation}
depends on its argument restricted to a small neighborhood of $\partial\Lambda_1$ \dots\ The
dependence is analytic
\end{quotation}
and it does not involve $A'$. At the step from level $k$ to level $k+1$,
\cite{B-CMP119} makes the same change:
\begin{quotation}
$H^{(k)}$ is an analytic function of the background field $U_{k+1}$ restricted to
$\Lambda_{k+1}^{c}\cap\Lambda_k$, bounded on the set $S_{k+1}$ by $O(1)\varepsilon_k\exp(-R_k)$
\dots\ For bonds in $(R^{\sim2}_{k+1}\cap S_{k+1})^{(k)*}$ we take the change of variables given by
the formula (1.22), with $H^{(k)}(b)$ \dots\ For the remaining bonds we take the simpler change of
variables \dots\ The other characteristic functions are unchanged. The action is changed by terms
corresponding to (1.24), and by terms obtained by the first order expansion of the action with
respect to the perturbative parts of the changes of variables. We denote this new part of the
action by $V^{(k)}(S_{k+1},A_k,H^{(k)})$. All terms of the previous action contribute to it
\end{quotation}
At the large-field stages \cite{B-CMP122a} repeats it:
\begin{quotation}
Now we make the change of variables (1.22) [III], with $H^{(j)}_Z$ instead of $H_{1,Ax}$, for the
bonds belonging to $(\Omega^{c}_{j+1}\setminus Z''_{j+1})^{(j)*}$, or rather to the components of
this set with the large field function $1-\chi'^{(j)}$ \dots\ The change of variables gives also
the usual new terms in the action
\end{quotation}
The reference [III] in that sentence is to \cite{B-CMP119}, whose change of variables is quoted
above. In \cite{B-CMP122a}, $H^{(j)}_Z$ is an analytic function of the configuration restricted to
$Z$, and on $(\Omega^{c}_{j+1}\setminus Z''_j)^{(j)}$ it is bounded by
$O(1)\exp(-R_j)\varepsilon_j$, with $\varepsilon_j$ as in \cite[(1.57)]{B-CMP122a}. Only the
consequence that this bound lies below every power of $g_j$ is used.

We write $\mathsf{g}$ for Ba{\l}aban's cut-off $g$, since the letter $g$ is reserved for the
coupling. We also write
\begin{equation*}
S^{*}:=S^{(k)*}_{k+1},\qquad
\mathfrak{S}^{22}:=(R^{\sim2}_{k+1}\cap S_{k+1})^{(k)*},\qquad
H:=H^{(k)}.
\end{equation*}
We set $H_{\max}:=\sup_{S_{k+1}}|H|$. On the tube of background fields of the step,
$H_{\max}\le O(1)\varepsilon_k e^{-R_k}$, with $\varepsilon_k$ and $R_k$ the quantities so denoted
in \cite{B-CMP119}; this is the bound just quoted from \cite{B-CMP119}, in
the form in which it is part of the analyticity statement for the effective densities at complex
arguments. Throughout,
$G=SU(N)$ acts on $\mathfrak{g}$ and $G^{\mathbb{C}}$ acts on $\mathfrak{g}^{\mathbb{C}}$ by
$\operatorname{Ad}_u X=uXu^{-1}$. The symbol $du$ denotes normalised Haar measure on $G$.

\begin{definition}[An invariant smooth cut-off and its holomorphic lift]
\label{4.22:def-invariant-cutoff}
\emph{The cut-off.} Let $\mathfrak{b}_k$ be the small-field threshold of level $k$. At level $k$,
the parameter $\delta_0$ of \cite[(1.22)]{B-CMP119} is the number $g_k\mathfrak{b}_k$. Set the
inner plateau and the outer plateau
\begin{equation*}
\Pi^{(0)}_k:=\{A'\in\mathfrak{g}:|A'|\le\tfrac43\mathfrak{b}_k\},\qquad
\Pi^{(1)}_k:=\{A'\in\mathfrak{g}:|\exp(ig_kA')-\mathbb{1}|\ge\tfrac53g_k\mathfrak{b}_k\}.
\end{equation*}
Of the cut-off, \cite{B-CMP119} prescribes only that it is $C^{\infty}$, its range and its two
plateaux. We fix the cut-off
$\mathsf{g}=\mathsf{g}_k$ of the change of variables at step $k+1$ to be a $C^{\infty}$
function on $\mathfrak{g}$ with the following properties, for a constant $G_1=G_1(N)$ depending on
$N$ only, fixed in the proof below; the invariance and the gradient bound are additional to what
\cite{B-CMP119} prescribes:
\begin{equation}\label{4.22:eq-invariant-cutoff-a}
\begin{gathered}
\mathsf{g}_k(\operatorname{Ad}_uA')=\mathsf{g}_k(A')\quad(u\in G,\ A'\in\mathfrak{g}),\qquad
0\le\mathsf{g}_k\le1,\\
\mathsf{g}_k=0\ \text{on}\ \Pi^{(0)}_k,\qquad \mathsf{g}_k=1\ \text{on}\ \Pi^{(1)}_k,\qquad
\sup|\nabla\mathsf{g}_k|\le G_1/\mathfrak{b}_k .
\end{gathered}
\end{equation}
Here $\sup|\nabla\mathsf{g}_k|\le G_1/\mathfrak{b}_k$ means that
$|D\mathsf{g}_k(A')[X]|\le(G_1/\mathfrak{b}_k)|X|$ for all $A',X\in\mathfrak{g}$. At the
large-field stages, the cut-off $\mathsf{g}'_j$ of the change of variables of \cite{B-CMP122a} is
chosen in the same way, with $\mathfrak{b}_j$ replaced by $\delta'_j/g_j$ in both plateaux and in
the gradient bound, so that the number $g_j\mathfrak{b}_j$ in the place of $\delta_0$ becomes
$\delta'_j$; here $\delta'_j$ is the corresponding parameter of the large-field stages in
\cite{B-CMP122a}.

\emph{The lift.} Let $c_B>0$ and $C_B$ be the constants fixed in the proof below from the
Baker--Campbell--Hausdorff series; $c_B$ is a radius for that series, shrunk so that
$C_Bc_B\le1$. For $(a,\eta)$ in
\begin{equation*}
\mathcal{D}_B:=\{(a,\eta)\in\mathfrak{g}^{\mathbb{C}}\times\mathfrak{g}^{\mathbb{C}}:
g_k|a|+|\eta|<c_B\},
\end{equation*}
put
\begin{equation}\label{4.22:eq-invariant-cutoff-1}
\mathfrak{c}(a,\eta):=\frac{1}{ig_k}\log\bigl[\exp(ig_ka)\exp(-i\eta)\bigr].
\end{equation}
Then $\mathfrak{c}$ is holomorphic on $\mathcal{D}_B$ and $\mathfrak{c}(a,0)=a$. With this constant
$C_B$, the following bounds hold on $\mathcal{D}_B$:
\begin{equation}\label{4.22:eq-invariant-cutoff-2}
|\mathfrak{c}(a,\eta)-a+\eta/g_k|\le C_B|a||\eta|,\qquad
\|\partial_a\mathfrak{c}-1\|\le C_B|\eta|,\qquad
\|\partial_\eta\mathfrak{c}\|\le 2/g_k .
\end{equation}
Moreover, for $u\in G^{\mathbb{C}}$, whenever both $(a,\eta)$ and
$(\operatorname{Ad}_ua,\operatorname{Ad}_u\eta)$ lie in $\mathcal{D}_B$ (for $u\in G$ the second
condition follows from the first),
\begin{equation}\label{4.22:eq-invariant-cutoff-3}
\operatorname{Ad}_u\mathfrak{c}(a,\eta)=\mathfrak{c}(\operatorname{Ad}_ua,\operatorname{Ad}_u\eta).
\end{equation}

Let $\mathcal{A}ct(\mathsf{B})$ be the exponent of the fluctuation integral of the step, in the form
in which it enters the proof of the correlation theorem, taken before the change of variables. It
reads
\begin{equation*}
\mathcal{A}ct(\mathsf{B})=-\tfrac12\langle\mathsf{B},\mathcal{M}\mathsf{B}\rangle
+\sum_YV^{\mathrm{pre}}(Y;\mathsf{B}),
\end{equation*}
where $\mathcal{M}$ is the operator of its quadratic part and the $V^{\mathrm{pre}}(Y;\mathsf{B})$
are its remaining terms, indexed by their domains $Y$.
Let $\mathsf{B}$ be complex, and let $\mathsf{s}=(\mathsf{s}_b)$ and $\mathsf{v}=(\mathsf{v}_b)$,
$b\in\mathfrak{S}^{22}$, be auxiliary variables, with $\mathsf{s}_b$ complex and
$\mathsf{v}_b\in\mathfrak{g}^{\mathbb{C}}$. On $S^{*}\setminus\mathfrak{S}^{22}$ put
$\mathsf{s}_b:=1$ and $\mathsf{v}_b:=0$. For $(\mathsf{B},\mathsf{s},\mathsf{v})$ such that
$(\mathsf{B}(b),\mathsf{s}_bH(b))\in\mathcal{D}_B$ for every $b\in S^{*}$, and for $b\in S^{*}$,
define
\begin{equation}\label{4.22:eq-invariant-cutoff-4}
\begin{split}
\delta_b&:=\mathfrak{c}(\mathsf{B}(b),\mathsf{s}_bH(b))-\mathsf{B}(b),\\
\mathsf{J}_b&:=\partial_a\mathfrak{c}+\bigl(\partial_\eta\mathfrak{c}[H(b)]\bigr)
\langle\mathsf{v}_b,\cdot\,\rangle\ \in\operatorname{End}\mathfrak{g}^{\mathbb{C}},
\end{split}
\end{equation}
with the derivatives taken at $(\mathsf{B}(b),\mathsf{s}_bH(b))$. Finally set
\begin{equation}\label{4.22:eq-invariant-cutoff-5}
f_b^{\uparrow}(\mathsf{B},\mathsf{s},\mathsf{v})
:=\int_0^1dt\,\bigl\langle\partial_{\mathsf{B}(b)}\mathcal{A}ct(\mathsf{B}+t\delta),
\delta_b\bigr\rangle+\operatorname{Tr}\log\mathsf{J}_b,
\qquad
V^{(k)\uparrow}:=\sum_{b\in S^{*}}f_b^{\uparrow}.
\end{equation}
Here $\log\mathsf{J}_b$ is the principal logarithm, defined where $\|\mathsf{J}_b-1\|<1$;
$\delta=(\delta_b)_{b\in S^{*}}$, and $\mathsf{B}+t\delta$ is the configuration $\mathsf{B}$
with $\mathsf{B}(b)$ replaced by $\mathsf{B}(b)+t\delta_b$ for every $b\in S^{*}$.
\end{definition}

\begin{proof}
\emph{A cut-off with the properties \eqref{4.22:eq-invariant-cutoff-a} exists.}

First, the two sets are $\operatorname{Ad}_G$-invariant. For $u\in G$ the map $\operatorname{Ad}_u$
preserves the operator norm, since $u$ is unitary. Hence $\Pi^{(0)}_k$ is invariant. From the power
series,
\begin{equation*}
\exp(ig_k\operatorname{Ad}_uA')=u\exp(ig_kA')u^{-1}.
\end{equation*}
It follows that
\begin{equation*}
|\exp(ig_k\operatorname{Ad}_uA')-\mathbb{1}|=|u(\exp(ig_kA')-\mathbb{1})u^{-1}|
=|\exp(ig_kA')-\mathbb{1}|,
\end{equation*}
so $\Pi^{(1)}_k$ is invariant as well.

Next, the two sets are at distance at least $\mathfrak{b}_k/3$. For $X\in\mathfrak{g}$ the matrix
$X$ is Hermitian, so $e^{iX}-\mathbb{1}$ is normal with eigenvalues $e^{i\theta}-1$, where $\theta$
runs over the real eigenvalues of $X$. Since $|e^{i\theta}-1|\le|\theta|$, we get
$|e^{iX}-\mathbb{1}|\le|X|$. Take $X=g_kA'$ with $A'\in \Pi^{(1)}_k$. Then
$g_k|A'|\ge\frac53g_k\mathfrak{b}_k$, that is,
\begin{equation*}
\Pi^{(1)}_k\subset\{A'\in\mathfrak{g}:|A'|\ge\tfrac53\mathfrak{b}_k\}.
\end{equation*}
Let $A'_0\in \Pi^{(0)}_k$ and $A'_1\in \Pi^{(1)}_k$. By the triangle inequality,
\begin{equation*}
|A'_1-A'_0|\ge\tfrac53\mathfrak{b}_k-\tfrac43\mathfrak{b}_k=\tfrac13\mathfrak{b}_k .
\end{equation*}

Now take a $C^{\infty}$ function $\varphi$ on $\mathfrak{g}$ with values in $[0,1]$ that separates
the two sets: $\varphi=0$ on $\Pi^{(0)}_k$, $\varphi=1$ on $\Pi^{(1)}_k$, and
$|D\varphi(A')[X]|\le(G_1/\mathfrak{b}_k)|X|$. Such a function exists with $G_1$ depending on $N$
only. Indeed, fix for each $N$ a $C^{\infty}$ function $\psi\ge0$ on
$\mathfrak{g}\cong\mathbb{R}^{N^2-1}$, supported in $\{|X'|\le1\}$, with integral one for Lebesgue
measure, and put $c_\psi:=\sup_{|X|\le1}\int|D\psi(X')[X]|\,dX'$, a number depending on $N$ only.
Let
$\widehat{\Pi}^{(1)}_k$ be the set of points at distance at most $\mathfrak{b}_k/12$ from
$\Pi^{(1)}_k$, and let $\varphi$ be the convolution of the indicator of $\widehat{\Pi}^{(1)}_k$
with $(12/\mathfrak{b}_k)^{N^2-1}\psi(12\,\cdot/\mathfrak{b}_k)$. Then $\varphi$ is $C^{\infty}$
with values in $[0,1]$. It equals $1$ on $\Pi^{(1)}_k$, because the ball of radius
$\mathfrak{b}_k/12$ about a point of $\Pi^{(1)}_k$ lies in $\widehat{\Pi}^{(1)}_k$. It equals $0$
at every point at distance more than $\mathfrak{b}_k/6$ from $\Pi^{(1)}_k$, hence on
$\Pi^{(0)}_k$, which is at distance at least $\mathfrak{b}_k/3$. Differentiating the mollifier,
$|D\varphi(A')[X]|\le(12c_\psi/\mathfrak{b}_k)|X|$, so that $G_1:=12c_\psi$ will do. Average
$\varphi$ over $\operatorname{Ad}_G$:
\begin{equation*}
\mathsf{g}_k(A'):=\int_G\varphi(\operatorname{Ad}_uA')\,du .
\end{equation*}
We check each property in turn.
\begin{itemize}
\item Range. The values lie in $[0,1]$, because the integrand does and $du$ is a probability measure.
\item Plateaux. If $A'\in \Pi^{(0)}_k$, then $\operatorname{Ad}_uA'\in \Pi^{(0)}_k$ for every $u$, so
the integrand vanishes identically and $\mathsf{g}_k(A')=0$. In the same way $\mathsf{g}_k=1$ on
$\Pi^{(1)}_k$.
\item Smoothness. The group $G$ is compact, $\varphi$ is $C^{\infty}$, and
$(u,A')\mapsto\operatorname{Ad}_uA'$ is smooth. Differentiation under the integral is therefore
legitimate, and $\mathsf{g}_k$ is $C^{\infty}$.
\item Gradient. We have
$D(\varphi\circ\operatorname{Ad}_u)(A')[X]=D\varphi(\operatorname{Ad}_uA')[\operatorname{Ad}_uX]$.
This has modulus at most $(G_1/\mathfrak{b}_k)|\operatorname{Ad}_uX|=(G_1/\mathfrak{b}_k)|X|$.
Integrating over $u$ gives the bound in \eqref{4.22:eq-invariant-cutoff-a}.
\item Invariance. For $v\in G$,
\begin{equation*}
\mathsf{g}_k(\operatorname{Ad}_vA')=\int_G\varphi(\operatorname{Ad}_{uv}A')\,du=\mathsf{g}_k(A'),
\end{equation*}
by right invariance of Haar measure.
\end{itemize}
The same argument, with $\mathfrak{b}_j$ replaced by $\delta'_j/g_j$, gives $\mathsf{g}'_j$.

\emph{Properties of the lift.} By the Baker--Campbell--Hausdorff theorem there are numbers $c>0$
and $C\ge1$ with the following property: for $X_1,X_2\in\mathfrak{g}^{\mathbb{C}}$ with
$|X_1|+|X_2|<c$, the logarithm of $e^{X_1}e^{X_2}$, given by the holomorphic functional calculus,
satisfies
\begin{equation*}
\log(e^{X_1}e^{X_2})=X_1+X_2+F_{\mathrm{BCH}}(X_1,X_2),
\end{equation*}
where $F_{\mathrm{BCH}}$, the sum of the terms of order higher than one of the series, is
holomorphic on this set and obeys
\begin{equation*}
|F_{\mathrm{BCH}}(X_1,X_2)|\le C|X_1||X_2|,\qquad
\|\partial_{X_1}F_{\mathrm{BCH}}\|\le C|X_2|,\qquad
\|\partial_{X_2}F_{\mathrm{BCH}}\|\le C|X_1|.
\end{equation*}
Fix $c_B:=\min\{c,1/C\}$ and $C_B:=C$. For $(a,\eta)\in\mathcal{D}_B$ put $X_1=ig_ka$ and
$X_2=-i\eta$; then $|X_1|+|X_2|=g_k|a|+|\eta|<c_B\le c$, and
\begin{equation*}
\mathfrak{c}(a,\eta)=a-\frac{\eta}{g_k}+\frac{1}{ig_k}F_{\mathrm{BCH}}(ig_ka,-i\eta).
\end{equation*}
Hence $\mathfrak{c}$ is holomorphic on $\mathcal{D}_B$. The bound on $F_{\mathrm{BCH}}$ gives
$F_{\mathrm{BCH}}(X_1,0)=0$, so $\mathfrak{c}(a,0)=a$. Differentiating,
$\partial_a\mathfrak{c}=1+\partial_{X_1}F_{\mathrm{BCH}}$ and
$\partial_\eta\mathfrak{c}=-g_k^{-1}(1+\partial_{X_2}F_{\mathrm{BCH}})$, with the derivatives of
$F_{\mathrm{BCH}}$ taken at $(ig_ka,-i\eta)$. The bounds on $F_{\mathrm{BCH}}$ then give
\begin{align*}
|\mathfrak{c}(a,\eta)-a+\eta/g_k|&\le g_k^{-1}Cg_k|a||\eta|=C_B|a||\eta|,\\
\|\partial_a\mathfrak{c}-1\|&\le C|\eta|=C_B|\eta|,\\
\|\partial_\eta\mathfrak{c}\|&\le g_k^{-1}(1+Cg_k|a|)\le g_k^{-1}(1+Cc_B)\le2/g_k .
\end{align*}
This proves \eqref{4.22:eq-invariant-cutoff-2}.

\emph{Equivariance of the lift.} Fix $u\in G^{\mathbb{C}}$, and suppose that both $(a,\eta)$ and
$(\operatorname{Ad}_ua,\operatorname{Ad}_u\eta)$ lie in $\mathcal{D}_B$. From the power series,
\begin{equation*}
\exp(ig_k\operatorname{Ad}_ua)\exp(-i\operatorname{Ad}_u\eta)=u\exp(ig_ka)\exp(-i\eta)u^{-1}.
\end{equation*}
The logarithm in \eqref{4.22:eq-invariant-cutoff-1} is given by the holomorphic functional
calculus, which commutes with conjugation by an invertible matrix; applied to the product above,
\begin{equation*}
\log\bigl(u\exp(ig_ka)\exp(-i\eta)u^{-1}\bigr)=u\log\bigl(\exp(ig_ka)\exp(-i\eta)\bigr)u^{-1}.
\end{equation*}
Dividing by $ig_k$ gives \eqref{4.22:eq-invariant-cutoff-3}. For $u\in G$ the condition on
$(\operatorname{Ad}_ua,\operatorname{Ad}_u\eta)$ follows from that on $(a,\eta)$, since
$\operatorname{Ad}_u$ preserves the operator norm.

\emph{The terms of $\mathcal{A}ct$ are defined at all free bonds.} Each term $V^{\mathrm{pre}}(Y)$
is of one of the kinds of term treated by the analyticity statement for the effective densities at
complex arguments, namely:
\begin{itemize}
\item the terms of the effective action that this statement lists individually;
\item the terms denoted $\mathbf{P}^{(k)}$ in \cite{B-CMP119};
\item the terms carrying the complex source $\zeta_k$;
\item the stored units of the effective action.
\end{itemize}
By that analyticity statement, each $V^{\mathrm{pre}}(Y)$ is holomorphic, with modulus at most the
number $\nu(Y)$ that statement attaches to $Y$, on
$|\mathsf{B}|<T^{\mathbb{C}}$ at all free bonds of $Y$. No cut-off has been applied to
$\mathcal{A}ct$, so that, apart from its quadratic part, the first term of
\eqref{4.22:eq-invariant-cutoff-5} involves only these terms, evaluated at the shifted
configurations $\mathsf{B}+t\delta$.
\end{proof}

The invariance in \eqref{4.22:eq-invariant-cutoff-a} is not optional. The action
$V^{(k)}(S_{k+1},A_k,H^{(k)})$ of \cite{B-CMP119} contains the terms
$\mathsf{g}(A(b))H(b)$. These are boundary terms \cite{B-CMP119}, and \cite[(2.41)(iii)]{B-CMP119}
is asserted for them. The same invariance is used, without the cut-off being named, in two places:
the analyticity statement for the effective densities at complex arguments, and one step of the
proof of the flat form of the correlation theorem. In both places the norm $|\cdot|$, the
thresholds and the sets $\mathcal{S}(\cdot)$ entering those statements are
$\operatorname{Ad}$-invariant objects.

The choice changes no density. For every cut-off with Ba{\l}aban's range and plateaux,
\cite[(1.22)]{B-CMP119} is a change of the integration variable. The only properties of the
cut-off used in the display \cite[(1.23)]{B-CMP119}, which follows the change of variables, are
its plateaux.

\begin{lemma}[Holomorphy, invariance and smallness of the lifted cut-off terms]
\label{4.22:lem-lifted-cutoff}
The notation is that of Definition~\ref{4.22:def-invariant-cutoff}. Let $S^{*}:=S^{(k)*}_{k+1}$ be the set of
bonds at which Ba{\l}aban's change of variables \cite[(1.22)]{B-CMP119} is made at step $k+1$,
and let $\mathfrak{S}^{22}:=(R^{\sim 2}_{k+1}\cap S_{k+1})^{(k)*}$ be the subset at which this change
carries the cut-off. Write $H:=H^{(k)}$ for Ba{\l}aban's function of the background entering the
change, and $H_{\max}$ for the supremum of $|H|$ over $S_{k+1}$ on the tube of complex background
pairs $(\mathbf{U},\mathbf{J})$; by the holomorphy statement for the stored terms in the apparatus of the
correlation theorem, $H_{\max}\le O(1)\,\varepsilon_ke^{-R_k}$ on this tube, where $\varepsilon_k$ and $R_k$ are
Ba{\l}aban's parameters at level $k$ in the bound on $H^{(k)}$ in \cite{B-CMP119}. Let
$\mathfrak{c}$, $\delta_b$, $\mathsf{J}_b$, $f_b^{\uparrow}$ and $V^{(k)\uparrow}$ be the holomorphic
lift of the change of variables \cite[(1.22)]{B-CMP119} of
Definition~\ref{4.22:def-invariant-cutoff}, with $\mathsf{s}_b=1$ and $\mathsf{v}_b=0$
at $b\in S^{*}\setminus\mathfrak{S}^{22}$. Let $c_B$ and $C_B$ be the radius and the constant of the
Baker--Campbell--Hausdorff bounds for the map $\mathfrak{c}$. Let the cut-off
$\mathsf{g}=\mathsf{g}_k$ satisfy \eqref{4.22:eq-invariant-cutoff-a}, with the constant $G_1=G_1(N)$
of its gradient bound, and write $\delta_k:=g_k\mathfrak{b}_k$, Ba{\l}aban's threshold $\delta$ of
\cite[(1.22)]{B-CMP119} at level $k$, expressed through the small-field threshold $\mathfrak{b}_k$.
Assume
\begin{equation}\label{4.22:eq-lifted-cutoff-a}
\begin{gathered}
g_kT^{\mathbb{C}}+2H_{\max}<c_B,\qquad C_Bg_kT^{\mathbb{C}}\le 1,\\
\frac{4H_{\max}}{g_k}\le 1,\qquad 2C_BH_{\max}+\frac{4G_1H_{\max}}{\delta_k}\le\frac12 .
\end{gathered}
\end{equation}
\begin{enumerate}
\item[(i)] Let $\mathsf{B}$ be real at the bonds of $\mathfrak{S}^{22}$, and put
$\mathsf{s}_b=\mathsf{g}(\mathsf{B}(b))$ and $\mathsf{v}_b=\nabla\mathsf{g}(\mathsf{B}(b))$ for
$b\in\mathfrak{S}^{22}$. Then $V^{(k)\uparrow}$ coincides, bond by bond, with Ba{\l}aban's term
$V^{(k)}(S_{k+1},A_k,H^{(k)})$ of \cite{B-CMP119}.
\item[(ii)] For every $b\in S^{*}$ the function $f_b^{\uparrow}$ is holomorphic on the product of
the tube of background pairs $(\mathbf{U},\mathbf{J})$ with the set
\begin{equation*}
\{|\mathsf{B}|<T^{\mathbb{C}}-2\}\times\prod_{b'\in\mathfrak{S}^{22}}
\bigl(\{|\mathsf{s}_{b'}|<2\}\times\{|\mathsf{v}_{b'}|<2G_1/\mathfrak{b}_k\}\bigr),
\end{equation*}
where the bound on $\mathsf{B}$ is imposed at every bond, the bonds of $\mathfrak{S}^{22}$ included.
\item[(iii)] For every $G$-valued gauge transformation $u$, the function $f_b^{\uparrow}$ is invariant
under
\begin{equation*}
(\mathbf{U},\mathbf{J},\mathsf{B},\mathsf{s},\mathsf{v})\longmapsto
(\mathbf{U}^{u},\operatorname{Ad}_u\mathbf{J},\operatorname{Ad}_u\mathsf{B},\mathsf{s},
\operatorname{Ad}_u\mathsf{v}).
\end{equation*}
\item[(iv)] On the domain of \textup{(ii)},
\begin{equation}\label{4.22:eq-lifted-cutoff-1}
|f_b^{\uparrow}|\le C_{22}\,(1+G_1)\,\frac{A_0}{A_1}\,(T^{\mathbb{C}})^2e^{-R_k},
\end{equation}
where $C_{22}$ depends only on $d$, $L$, $M$, $\kappa_1$ and $N$. The right-hand side is smaller
than every power of $g_k$: $(\log(x_k-c_5))^r\le\check{R}_k$ with $r\ge 2$
(see \eqref{4.22:eq-steps-sites-b}), and $T^{\mathbb{C}}$ is polylogarithmic in the coupling.
\end{enumerate}
\end{lemma}

\begin{proof}
We use the following properties of
$\mathfrak{c}(a,\eta)=(1/ig_k)\log[\exp(ig_ka)\exp(-i\eta)]$, fixed with the holomorphic lift of the
change of variables in Definition~\ref{4.22:def-invariant-cutoff}: it is holomorphic on
$\{g_k|a|+|\eta|<c_B\}\subset\mathfrak{g}^{\mathbb{C}}\times\mathfrak{g}^{\mathbb{C}}$,
$\mathfrak{c}(a,0)=a$, and there
\begin{equation}\label{4.22:eq-lifted-cutoff-2}
|\mathfrak{c}(a,\eta)-a+\eta/g_k|\le C_B|a||\eta|,\qquad
\|\partial_a\mathfrak{c}-1\|\le C_B|\eta|,\qquad \|\partial_\eta\mathfrak{c}\|\le 2/g_k ;
\end{equation}
moreover $\operatorname{Ad}_u\mathfrak{c}(a,\eta)=\mathfrak{c}(\operatorname{Ad}_ua,\operatorname{Ad}_u\eta)$
for all $u\in G^{\mathbb{C}}$. The action before the change of variables, that is the exponent of
\cite[(3.15)]{B-CMP119} before the change, is
\begin{equation*}
\mathcal{A}ct(\mathsf{B})=-\frac12\langle\mathsf{B},\mathcal{A}\mathsf{B}\rangle
+\sum_YV^{\mathrm{pre}}(Y;\mathsf{B}),
\end{equation*}
where $\mathcal{A}$ here denotes the operator of the quadratic part of the fluctuation action, as in
the apparatus of the correlation theorem; the sum runs over the polymers $Y$ of the expansion before the
change, and by the holomorphy statement for the stored terms each term $V^{\mathrm{pre}}(Y)$ is
holomorphic on $|\mathsf{B}|<T^{\mathbb{C}}$ at all free bonds of $Y$, no cut-off having yet been
applied, and bounded there by its activity $\nu(Y)$.

\emph{(i).} The real variable $\mathsf{B}$ plays the role of the new integration variable $A'$ of
\cite[(1.22)]{B-CMP119}. The function $H$ depends on the background alone, not on $A'$, and the change
of variables is made bond by bond: $A_k(b)=\mathfrak{c}(A'(b),\mathsf{g}(A'(b))H(b))$ for
$b\in\mathfrak{S}^{22}$, and the simpler change, with $\mathsf{g}$ replaced by $1$, at
$b\in S^{*}\setminus\mathfrak{S}^{22}$. Hence, after the change, the integrand is the integrand of
\cite[(3.15)]{B-CMP119} before the change evaluated at $A_k$, multiplied by
$\prod_b|\det\partial A_k(b)/\partial A'(b)|$. Since $H$ does not depend on $A'$, the chain rule gives
\begin{equation*}
\frac{\partial A_k(b)}{\partial A'(b)}=\partial_a\mathfrak{c}
+\partial_\eta\mathfrak{c}[H(b)]\,\langle\nabla\mathsf{g}(A'(b)),\cdot\,\rangle ,
\end{equation*}
evaluated at $(A'(b),\mathsf{g}(A'(b))H(b))$; with $\mathsf{s}_b,\mathsf{v}_b$ as in (i) this is
$\mathsf{J}_b$, and the same holds on $S^{*}\setminus\mathfrak{S}^{22}$ with $\mathsf{s}_b=1$,
$\mathsf{v}_b=0$. These are the Jacobians exponentiated in \cite[(1.24)]{B-CMP119}. By the bound on
$\|\mathsf{J}_b-1\|$ proved in (ii) below, which does not use (i), for real $\mathsf{B}$ in the domain
of (ii) $\log\mathsf{J}_b$ is given by its
convergent series; at the real point $\mathsf{J}_b$ maps $\mathfrak{g}$ to itself, being the Jacobian of
a real change of variables, so $\operatorname{Tr}\log\mathsf{J}_b$ is real and
$|\det\mathsf{J}_b|=\exp\operatorname{Tr}\log\mathsf{J}_b$. Next,
$\delta_b=A_k(b)-A'(b)$; writing $\delta_{S^{*}}:=(\delta_b)_{b\in S^{*}}$, the fundamental theorem
of calculus along $t\mapsto A'+t\delta_{S^{*}}$ gives
\begin{equation*}
\mathcal{A}ct(A'+\delta_{S^{*}})-\mathcal{A}ct(A')=\sum_{b\in S^{*}}\int_0^1
\langle\partial_{\mathsf{B}(b)}\mathcal{A}ct(A'+t\delta_{S^{*}}),\delta_b\rangle\,dt .
\end{equation*}
This is the first-order expansion of the action with respect to the perturbative parts of the change of
variables, in the words of \cite{B-CMP119} at step $k+1$, taken with its exact integral remainder; no
first-order surrogate is used.
Adding $\sum_b\operatorname{Tr}\log\mathsf{J}_b$, the exponent after the change is
$\mathcal{A}ct(A')+\sum_{b\in S^{*}}f_b^{\uparrow}=\mathcal{A}ct(A')+V^{(k)\uparrow}$, and the
term at $b$ is the term Ba{\l}aban collects at $b$ in $V^{(k)}(S_{k+1},A_k,H^{(k)})$.

\emph{(ii).} Let $|\mathsf{B}|<T^{\mathbb{C}}-2$ at every bond, and let $|\mathsf{s}_{b'}|<2$,
$|\mathsf{v}_{b'}|<2G_1/\mathfrak{b}_k$ for all $b'\in\mathfrak{S}^{22}$ (at the bonds of
$S^{*}\setminus\mathfrak{S}^{22}$ one has $\mathsf{s}_b=1$, $\mathsf{v}_b=0$). For every $b\in S^{*}$,
at $(a,\eta)=(\mathsf{B}(b),\mathsf{s}_bH(b))$ we have
$g_k|a|+|\eta|<g_kT^{\mathbb{C}}+2H_{\max}<c_B$ by the first member of
\eqref{4.22:eq-lifted-cutoff-a}, so $\mathfrak{c}$ is holomorphic there. By the first bound of
\eqref{4.22:eq-lifted-cutoff-2} and the second and third members of
\eqref{4.22:eq-lifted-cutoff-a},
\begin{equation*}
|\delta_b|\le\frac{|\mathsf{s}_b||H|}{g_k}\bigl(1+C_Bg_k|\mathsf{B}(b)|\bigr)
\le\frac{4H_{\max}}{g_k}\le 1,
\end{equation*}
at every $b\in S^{*}$, so $|\mathsf{B}+t\delta_{S^{*}}|<T^{\mathbb{C}}-1$ for $t\in[0,1]$. Thus
$\partial_{\mathsf{B}(b)}\mathcal{A}ct(\mathsf{B}+t\delta_{S^{*}})$ and $\delta_b$ are compositions of
holomorphic maps, continuous in $t$ and bounded, and the $t$-integral is holomorphic. By the second
and third bounds of \eqref{4.22:eq-lifted-cutoff-2} and the fourth member of
\eqref{4.22:eq-lifted-cutoff-a}, with $g_k\mathfrak{b}_k=\delta_k$,
\begin{equation*}
\|\mathsf{J}_b-1\|\le C_B|\mathsf{s}_b||H|+\frac{2}{g_k}|H||\mathsf{v}_b|
\le 2C_BH_{\max}+\frac{4G_1H_{\max}}{\delta_k}\le\frac12 ,
\end{equation*}
so the series of $\log\mathsf{J}_b$ converges and $\operatorname{Tr}\log\mathsf{J}_b$ is holomorphic.

\emph{(iii).} The function $H$ satisfies $H(\mathbf{U}^u)=\operatorname{Ad}_uH(\mathbf{U})$, since $H(b)$
is the logarithm of a ratio of averages at the point $b_-$. By the equivariance of $\mathfrak{c}$, the
map sends $\delta_b$ to $\operatorname{Ad}_u\delta_b$; differentiating the equivariance, it sends
$\partial_a\mathfrak{c}$ to $\operatorname{Ad}_u\,\partial_a\mathfrak{c}\,\operatorname{Ad}_u^{-1}$ and
the vector $\partial_\eta\mathfrak{c}[H]$ to $\operatorname{Ad}_u\partial_\eta\mathfrak{c}[H]$,
and, $\langle\cdot,\cdot\rangle$ being $\operatorname{Ad}_u$-invariant,
$\langle\operatorname{Ad}_u\mathsf{v}_b,\cdot\rangle=\langle\mathsf{v}_b,\operatorname{Ad}_u^{-1}\cdot\rangle$.
Hence $\mathsf{J}_b\mapsto\operatorname{Ad}_u\mathsf{J}_b\operatorname{Ad}_u^{-1}$, and
$\operatorname{Tr}\log\mathsf{J}_b$ is unchanged. The action $\mathcal{A}ct$ is invariant by the invariance
part of the holomorphy statement for the stored terms, so its gradient at $b$ is mapped by
$\operatorname{Ad}_u$, and the pairing with $\operatorname{Ad}_u\delta_b$ is unchanged.

\emph{(iv).} At $|\mathsf{B}+t\delta_{S^{*}}|<T^{\mathbb{C}}-1$, Cauchy's estimate at radius $1$ gives
$|\partial_{\mathsf{B}(b)}V^{\mathrm{pre}}(Y)|\le(\dim\mathfrak{g})\,\nu(Y)$, and, by the summed
activity bound of the correlation theorem's apparatus,
$\sum_{Y\ni b}\nu(Y)\le C_\kappa\varepsilon_V$, with that bound's constant $C_\kappa$ and activity
size $\varepsilon_V$. The quadratic part contributes $(\mathcal{A}\mathsf{B})(b)$, and by the
random-walk bound at zeroth order of the same apparatus
$|(\mathcal{A}\mathsf{B})(b)|\le B_WC_\rho(\delta_W)T^{\mathbb{C}}$, where $B_W$ and
$C_\rho(\delta_W)$ are the constants of that bound. With
$H_{\max}\le O(1)\varepsilon_ke^{-R_k}$, recalled before \eqref{4.22:eq-lifted-cutoff-a}, and
$\varepsilon_k/g_k=(A_0/A_1)\mathfrak{b}_k$, the bound of (ii) gives
$|\delta_b|\le 4\,O(1)(A_0/A_1)\mathfrak{b}_ke^{-R_k}$. Since $\|\mathsf{J}_b-1\|\le\frac12$, the
series of $\log$ gives $\|\log\mathsf{J}_b\|\le2\|\mathsf{J}_b-1\|$, hence
\begin{equation*}
|\operatorname{Tr}\log\mathsf{J}_b|\le 2(\dim\mathfrak{g})\|\mathsf{J}_b-1\|
\le 2(\dim\mathfrak{g})\Bigl[2C_B\varepsilon_k+4G_1\frac{A_0}{A_1}\Bigr]O(1)e^{-R_k}.
\end{equation*}
Adding the two contributions,
\begin{align*}
|f_b^{\uparrow}|\le{}&\bigl(B_WC_\rho(\delta_W)T^{\mathbb{C}}+(\dim\mathfrak{g})C_\kappa\varepsilon_V\bigr)
\,4\,O(1)\frac{A_0}{A_1}\mathfrak{b}_ke^{-R_k}\\
&+2(\dim\mathfrak{g})\Bigl[2C_B\varepsilon_k+4G_1\frac{A_0}{A_1}\Bigr]O(1)e^{-R_k},
\end{align*}
from which \eqref{4.22:eq-lifted-cutoff-1} follows once $\mathfrak{b}_k$, $\varepsilon_k$ and the
constants are absorbed into $C_{22}(1+G_1)(A_0/A_1)(T^{\mathbb{C}})^2$.
\end{proof}

\begin{corollary}[The orbit datum at a cut-off coordinate]\label{4.22:cor-cutoff-coordinate}
Assume the hypotheses of Lemma~\ref{4.22:lem-lifted-cutoff}. In clause~(a) of the
complex-extension statement for the densities $V$, the exception of that clause, under which the
bonds of the set $(R^{\sim 2}\cap S)^{*}$ singled out there are silent, remains in force for the
slot. On the chart, let $x$ be a
real field on $\mathfrak{S}^{22}$, the set of bonds of Ba{\l}aban's full change of variables,
with $|x|\le C_{\mathrm{far}}\mathfrak{b}_k$, where $C_{\mathrm{far}}$ denotes the far constant.
Let $g$ be a $G^{\mathbb{C}}$-valued site map with
$\operatorname{pol}(g)<\varepsilon_{\star}$ at every site. For every bond $b$ of
$\mathfrak{S}^{22}$, with law site $b_-$, put
\begin{equation}\label{4.22:eq-cutoff-coordinate-1}
\mathsf{B}(b):=\operatorname{Ad}_{g(b_-)}x(b),\qquad
\mathsf{s}_b:=\mathsf{g}\bigl(x(b)\bigr),\qquad
\mathsf{v}_b:=\operatorname{Ad}_{g(b_-)}\nabla\mathsf{g}\bigl(x(b)\bigr).
\end{equation}
Then the following hold:
\begin{itemize}
\item the $V^{(k)}$-terms, and hence the terms $V(Y)$ that contain them, are holomorphic in $g$
on $\mathfrak{G}_{\varepsilon_{\star}}$ and continuous in $x$;
\item they satisfy \eqref{4.22:eq-orbit-datum-b};
\item clause~(c) of the complex-extension statement holds on the chart.
\end{itemize}
\end{corollary}

\begin{proof}
Fix a bond $b$ of $\mathfrak{S}^{22}$ and write $h:=g(b_-)$, so that
$\operatorname{pol}(h)<\varepsilon_{\star}$.

\emph{The lifted variables lie in the domain of holomorphy.} By clause~(g2) of
Lemma~\ref{4.22:lem-polar-decomposition}, $|\operatorname{Ad}_hA|\le e^{2\operatorname{pol}(h)}|A|$
for every $A\in\mathfrak{g}^{\mathbb{C}}$. Applied to $A=x(b)$, this gives
\begin{equation*}
|\mathsf{B}(b)|\le e^{2\varepsilon_{\star}}|x(b)|
\le e^{2\varepsilon_{\star}}C_{\mathrm{far}}\mathfrak{b}_k
\le e^{2\varepsilon_{\star}}\varrho< T^{\mathbb{C}}-2 .
\end{equation*}
Here the third inequality uses $C_{\mathrm{far}}\mathfrak{b}_k\le\varrho$, and the last is
\eqref{4.22:eq-one-term-bound-b}, by which
$T^{\mathbb{C}}-e^{2\varepsilon_{\star}}\varrho\ge\mathfrak{b}_k$. Since $x(b)$ is real and the
cut-off $\mathsf{g}$ of Definition~\ref{4.22:def-invariant-cutoff} takes its values in $[0,1]$, we
have $\mathsf{s}_b\in[0,1]$, so that $|\mathsf{s}_b|<2$.

By the gradient bound $\sup|\nabla\mathsf{g}|\le G_1/\mathfrak{b}_k$ of
Definition~\ref{4.22:def-invariant-cutoff} and, once more, clause~(g2) of
Lemma~\ref{4.22:lem-polar-decomposition}, applied to $A=\nabla\mathsf{g}(x(b))$, we obtain
\begin{equation*}
|\mathsf{v}_b|\le e^{2\varepsilon_{\star}}\,\frac{G_1}{\mathfrak{b}_k}<\frac{2G_1}{\mathfrak{b}_k}.
\end{equation*}
Here $e^{2\varepsilon_{\star}}=1+c_{\eta}/(5(\Theta_{\flat}+1))<2$ is used, as in the proof of
Lemma~\ref{4.22:lem-one-term-bound}.

Hence $(\mathsf{B},\mathsf{s},\mathsf{v})$ lies in the domain of clause~(ii) of
Lemma~\ref{4.22:lem-lifted-cutoff}, in which the bonds of $\mathfrak{S}^{22}$ are included. On this
domain the lifted bond term $f_b^{\uparrow}$ is holomorphic.

\emph{Holomorphy in $g$ and continuity in $x$.} Since $\operatorname{Ad}_hA=hAh^{-1}$, the
variables $\mathsf{B}(b)$ and $\mathsf{v}_b$ are polynomial in the entries of $h$ and of $h^{-1}$,
whereas $\mathsf{s}_b$ does not depend on the rotation at all. Both $h$ and $h^{-1}$ depend
holomorphically on $h\in G^{\mathbb{C}}$. The composition of $f_b^{\uparrow}$ with these maps is
therefore holomorphic in $g$ on $\mathfrak{G}_{\varepsilon_{\star}}$. The maps
$x\mapsto(\mathsf{B},\mathsf{s},\mathsf{v})$ are continuous, because $\mathsf{g}$ is smooth, and
hence the composition is continuous in $x$. This proves the first assertion for the
$V^{(k)}$-terms, and so for the terms $V(Y)$ that contain them.

\emph{The identity \eqref{4.22:eq-orbit-datum-b}.} Let $g$ be $G$-valued, and recall the rotated
slot $x^g$, given on the bonds of $\mathfrak{S}^{22}$ by $x^g(b)=\operatorname{Ad}_{g(b_-)}x(b)$.
This is a real field on $\mathfrak{S}^{22}$, and
$\mathsf{B}=x^g$ by \eqref{4.22:eq-cutoff-coordinate-1}. The cut-off $\mathsf{g}$ is
$\operatorname{Ad}_G$-invariant by Definition~\ref{4.22:def-invariant-cutoff}, and this gives
$\mathsf{s}_b=\mathsf{g}(x(b))=\mathsf{g}(x^g(b))$. By \eqref{4.22:eq-polar-decomposition-e},
$\nabla\mathsf{g}(x^g(b))=\operatorname{Ad}_{g(b_-)}\nabla\mathsf{g}(x(b))$, that is,
$\mathsf{v}_b=\nabla\mathsf{g}(x^g(b))$.

Thus $\mathsf{B}$ is real and
$\mathsf{s}_b=\mathsf{g}(\mathsf{B}(b))$, $\mathsf{v}_b=\nabla\mathsf{g}(\mathsf{B}(b))$, which are
the hypotheses of clause~(i) of Lemma~\ref{4.22:lem-lifted-cutoff} at the real datum $x^g$. That
clause gives, bond by bond, $V^{(k)\uparrow}=V^{(k)}(S_{k+1},A_k,H^{(k)})$ at $x^g$, the latter
being the term of Ba{\l}aban's construction. This is \eqref{4.22:eq-orbit-datum-b} for these terms.

\emph{Localisation in the background.} The localisation of these terms in the background
variables, through the walk expansions of $H^{(k)}$ and of the second background quantity on which
these terms depend, is the one obtained in the localisation step of the proof of the
complex-extension statement; we invoke it without repeating the argument. It
does not involve $(\mathsf{s},\mathsf{v},g)$, and the parameter lemma of the correlation theorem
carries the parameters through it.

\emph{Size.} By clause~(iv) of Lemma~\ref{4.22:lem-lifted-cutoff}, the bound on $f_b^{\uparrow}$
lies below every power of $g_k$. On the lifted polydisc, this is the smallness by an arbitrary
positive power of $g_k$ that the localisation step of the proof of the complex-extension statement
requires; it is
ensured by the third item of the ceiling condition on the live radius. Hence clause~(c) of the
complex-extension statement holds on the chart.
\end{proof}

\begin{lemma}[The cut-off coordinates at a large-field stage]\label{4.22:lem-cutoff-large-field}
We work at the site $(\mathrm{P})$ of the large-field stages (see
\eqref{4.22:eq-steps-sites-b}). Let a frame $W_{\mathfrak{a}}$ be arrested at stage $n+1$ of the
large-field operation, that is, at level $j=h+n$, by the large-field factor $1-\chi'^{(j)}$ of
\cite{B-CMP122a}, in the sense of the definition of arrests at the large-field stages. The objects
of Lemma~\ref{4.22:lem-lifted-cutoff} are transported to the stage by the substitutions
\begin{equation}\label{4.22:eq-cutoff-large-field-1}
\begin{gathered}
\bigl(k,\ A_k,\ \mathfrak{b}_k,\ \delta_k,\ \mathsf{g}_k,\ H^{(k)},\ S^{*},\ T^{\mathbb{C}}\bigr)\\
\longmapsto\ \bigl(j,\ B_j,\ \mathfrak{b}'_j,\ \delta'_j,\ \mathsf{g}'_j,\ H^{(j)}_Z,\ S^{P},\
\varrho^B_j\bigr).
\end{gathered}
\end{equation}
Here $\mathfrak{b}'_j:=\delta'_j/g_j$, the function $\mathsf{g}'_j$ is the cut-off function of
Definition~\ref{4.22:def-invariant-cutoff} with $\delta_j$ replaced by $\delta'_j$, the set $S^{P}$
is given by
\begin{equation*}
S^{P}:=\bigl(\Omega^{c}_{j+1}\setminus Z''_{j+1}\bigr)^{(j)*}\cap W_{\mathfrak{a}} ,
\end{equation*}
and $\varrho^B_j$ is the complex radius of the arrested collar variable of
Lemma~\ref{4.22:lem-weak-collar}. Let
\begin{equation}\label{4.22:eq-cutoff-large-field-2}
\mathcal{A}ct^{P}(\mathsf{B}):=-\tfrac12\bigl\langle\mathsf{B},C^{*}\Delta^{(j)}C\mathsf{B}\bigr\rangle
+\mathbb{P}^{(j)}(g_j,A_j,\mathsf{B})+\sum_{e}v_e(\mathsf{B})
\end{equation}
be the exponent of the integrand of the stage before the change of variables (in the words of
\cite{B-CMP122a}, ``an expression analogous to the expression on the right-hand side of
(3.15) [III]''; in \cite{B-CMP122a} the reference [III] is \cite{B-CMP119}, so that the expression
meant is the right-hand side of \cite[(3.15)]{B-CMP119}), written in the localised form of the
stage's action. In \eqref{4.22:eq-cutoff-large-field-2},
$\langle\mathsf{B},C^{*}\Delta^{(j)}C\mathsf{B}\rangle$ is the quadratic form of the stage, and
$B_j$, the image of $A_k$ under \eqref{4.22:eq-cutoff-large-field-1}, is the variable of the stage
in the role that $A_k$ plays in Lemma~\ref{4.22:lem-lifted-cutoff}, as in \cite{B-CMP122a}. Here
$e$ runs over the
elements of the localised expansion of the stage's action, read in the frame convention of the
complex presentation of stored terms. By clauses (i) and (iv) of the lemma on the localised form of
the stage's action and by Lemma~\ref{4.22:lem-weak-collar}, each $v_e$ is holomorphic in the complex
collar variable $\mathsf{B}$ on $|\mathsf{B}|<\varrho^B_j$, depends on $\mathsf{B}$ only through its
restriction to the set $\mathsf{b}(e)$, and satisfies $|v_e|\le\nu^{\natural}(e)$ there.

Assume Definition~\ref{4.22:def-invariant-cutoff}, the conditions
\eqref{4.22:eq-lifted-cutoff-a} read through \eqref{4.22:eq-cutoff-large-field-1}, and the
smallness condition of the localised expansion of the stage's action. Let $f_b^{\uparrow,P}$ be the
lifted bond term formed from $\mathcal{A}ct^{P}$ through \eqref{4.22:eq-cutoff-large-field-1} in the
way in which $f_b^{\uparrow}$ is formed in Lemma~\ref{4.22:lem-lifted-cutoff}. Then:
\begin{enumerate}
\item[(a)] clauses (i)--(iii) of Lemma~\ref{4.22:lem-lifted-cutoff} hold word for word through
\eqref{4.22:eq-cutoff-large-field-1};
\item[(b)] clause (iv) of that lemma is replaced by the following: for every bond $b$,
\begin{equation}\label{4.22:eq-cutoff-large-field-3}
\bigl|f_b^{\uparrow,P}\bigr|\le C_{22}^{P}\,(1+G_1)
\Bigl[\frac{\varepsilon_j}{g_j}\,\varrho^B_j+\frac{\varepsilon_j}{\delta'_j}\Bigr]e^{-R_j},
\end{equation}
where $\varepsilon_j$ is the parameter of \cite[(1.57)]{B-CMP122a} at level $j$, which satisfies
$\varepsilon_j/g_j=(A_0/A_1)\mathfrak{b}_j$ with $\mathfrak{b}_j$ the small-field threshold at
level $j$; and the right-hand side is smaller than every positive power of $g_j$.
\end{enumerate}
Consequently the ``usual new terms in the action'' produced by the change of variables at the stage
\cite{B-CMP122a} have orbit data of radius $\varepsilon_{\star}$ at the cut-off coordinates of the
bonds of $S^{P}$, given by the three substitutions of Corollary~\ref{4.22:cor-cutoff-coordinate}.
\end{lemma}

\begin{proof}
\emph{Clauses (i)--(iii).} Of its setting, the proof of Lemma~\ref{4.22:lem-lifted-cutoff} uses
only the following properties: the function $H^{(k)}$ does not depend on the variables being
changed and is holomorphic in the background; the change of variables is made bond by bond; the
action before the change is holomorphic on a polydisc which exceeds the real range of the changed
variables by at least $2$; and that action is invariant under $\operatorname{Ad}$. Each of them
holds at the stage. By \cite{B-CMP122a}, $H^{(j)}_Z$ is an analytic function of the configuration
restricted to $Z$, free of the changed variable, and the change of variables at the stage is the one
of Lemma~\ref{4.22:lem-lifted-cutoff} with $H^{(j)}_Z$ in place of $H^{(k)}$, made bond by bond on
the bonds of $S^{P}$. The exponent $\mathcal{A}ct^{P}$ is holomorphic on the polydisc
$|\mathsf{B}|<\varrho^B_j$, which under \eqref{4.22:eq-cutoff-large-field-1} takes the place of the
polydisc $|\mathsf{B}|<T^{\mathbb{C}}$ and exceeds the real range by at least $2$. It is
$\operatorname{Ad}$-invariant: for the terms $v_e$ this is the gauge-invariance clause of the
invariance and counting lemma of the localised expansion, and for $\mathbb{P}^{(j)}$ it is
\cite[Lemma~2]{B-CMP116}. With these four properties, the proofs of clauses (i)--(iii) of
Lemma~\ref{4.22:lem-lifted-cutoff} apply without change through
\eqref{4.22:eq-cutoff-large-field-1}.

\emph{Clause (iv).} The first-order part of $f_b^{\uparrow,P}$ is the integral over $t\in[0,1]$ of
$\langle\partial_{\mathsf{B}(b)}\mathcal{A}ct^{P}(\mathsf{B}+t\delta),\delta_b\rangle$, as in the
proof of Lemma~\ref{4.22:lem-lifted-cutoff}(i). By that proof read through
\eqref{4.22:eq-cutoff-large-field-1}, the points $\mathsf{B}+t\delta$ satisfy
$|\mathsf{B}+t\delta|<\varrho^B_j-1$. Hence the Cauchy estimate at radius $1$ in the coordinate
$\mathsf{B}(b)$, applied to $v_e$, which is holomorphic and bounded by $\nu^{\natural}(e)$ on
$|\mathsf{B}|<\varrho^B_j$, gives
\begin{equation*}
\bigl|\partial_{\mathsf{B}(b)}v_e\bigr|\le(\dim\mathfrak{g})\,\nu^{\natural}(e),
\end{equation*}
and $\partial_{\mathsf{B}(b)}v_e=0$ unless $b\in\mathsf{b}(e)$, because $v_e$ reads $\mathsf{B}$ only
on $\mathsf{b}(e)$. Summing over $e$,
\begin{equation*}
\Bigl|\sum_{e}\partial_{\mathsf{B}(b)}v_e\Bigr|
\le(\dim\mathfrak{g})\sum_{e:\,\mathsf{b}(e)\ni b}\nu^{\natural}(e)
\le\tfrac12\dim\mathfrak{g},
\end{equation*}
where the bound $\sum_{e:\,\mathsf{b}(e)\ni b}\nu^{\natural}(e)\le\frac12$ is the anchored-total
estimate of the invariance and counting lemma of the localised expansion. That estimate holds under
the counting
hypothesis of that expansion and covers all the chains of the expansion. The contributions of the
quadratic form and of $\mathbb{P}^{(j)}$ are estimated exactly as in the proof of
Lemma~\ref{4.22:lem-lifted-cutoff}(iv) at the site $(\mathrm{T})$. Finally, as in the proof of
Lemma~\ref{4.22:lem-lifted-cutoff}(ii),
\begin{equation*}
|\delta_b|\le\frac{|\mathsf{s}_b|\,|H^{(j)}_Z(b)|}{g_j}\bigl(1+C_B\,g_j|\mathsf{B}(b)|\bigr).
\end{equation*}
We have $|\mathsf{s}_b|<2$. By the second condition of \eqref{4.22:eq-lifted-cutoff-a} read through
\eqref{4.22:eq-cutoff-large-field-1}, $1+C_Bg_j|\mathsf{B}(b)|\le2$. Together with the bound
$|H^{(j)}_Z(b)|\le O(1)\varepsilon_je^{-R_j}$, which \cite{B-CMP122a} gives on the set
$(\Omega^{c}_{j+1}\setminus Z''_{j})^{(j)}$ and which is applied here at the bonds $b$ of $S^{P}$,
this gives
\begin{equation*}
|\delta_b|\le4\,O(1)\,\frac{\varepsilon_j}{g_j}\,e^{-R_j}.
\end{equation*}
The Jacobian term is estimated as in the proof of Lemma~\ref{4.22:lem-lifted-cutoff}(iv). There the
quantity $H_{\max}/\delta_k$ becomes $H_{\max}/\delta'_j$ through
\eqref{4.22:eq-cutoff-large-field-1}, with $H_{\max}\le O(1)\varepsilon_je^{-R_j}$. This produces
the second term in the bracket of \eqref{4.22:eq-cutoff-large-field-3}. The factor $1+G_1$ comes
from the gradient bound of the cut-off. These estimates together give
\eqref{4.22:eq-cutoff-large-field-3}.

\emph{Form and size of the bound.} The substitutions \eqref{4.22:eq-cutoff-large-field-1} replace
$\mathfrak{b}_k$ by $\mathfrak{b}'_j$, but the quantity $\varepsilon_j$ is not transformed. In the
bound for $|\delta_b|$ one has $\varepsilon_j/g_j=(A_0/A_1)\mathfrak{b}_j$, with $\mathfrak{b}_j$ the
small-field threshold at level $j$, and not $(A_0/A_1)\mathfrak{b}'_j$. Moreover
$\mathfrak{b}_j/\varrho^B_j$ is a constant multiple of $(\log g_j^{-2})^{p_0-p_1}$, which is
unbounded as $g_j\to0$ because $p_1<p_0$ \cite{B-CMP122a}. Therefore the bound cannot be written in
the form of
Lemma~\ref{4.22:lem-lifted-cutoff}(iv) with $T^{\mathbb{C}}$ replaced by $\varrho^B_j$, and it is
kept in the form \eqref{4.22:eq-cutoff-large-field-3}. The same remark applies to
$H_{\max}/\delta'_j$ in the Jacobian term. As for the size, the first term of the bracket is
$(A_0/A_1)\mathfrak{b}_j\varrho^B_j$, with $\mathfrak{b}_j=A_1(\log g_j^{-2})^{p_0}$; since
$\mathfrak{b}_j/\varrho^B_j$ is a constant multiple of $(\log g_j^{-2})^{p_0-p_1}$, this term is a
constant multiple of $(\log g_j^{-2})^{p_0+p_1}$. The second term, $\varepsilon_j/\delta'_j$, is the
Jacobian term of the proof of Lemma~\ref{4.22:lem-lifted-cutoff}(iv) read through
\eqref{4.22:eq-cutoff-large-field-1}. Hence, with the factor $e^{-R_j}$ and for the same reason as
in Lemma~\ref{4.22:lem-lifted-cutoff}(iv), the right-hand side is smaller than every positive power
of $g_j$.

\emph{The new terms.} Clauses (i)--(iii) and the bound \eqref{4.22:eq-cutoff-large-field-3} take at
the stage the place of the properties of Lemma~\ref{4.22:lem-lifted-cutoff} used in the proof of
Corollary~\ref{4.22:cor-cutoff-coordinate}. Applying the three substitutions of that corollary at
the bonds of $S^{P}$ therefore gives orbit data of radius $\varepsilon_{\star}$ for the new terms of
the stage.
\end{proof}

\begin{remark}[Why the invariant cut-off needs no holomorphic tube]\label{4.22:rem-cutoff-remarks}
Let $\mathsf{g}=\mathsf{g}_k$ be the invariant cut-off of
Definition~\ref{4.22:def-invariant-cutoff}, and let $\mathsf{B}$, $\mathsf{s}$, $\mathsf{v}$ and
$f_b^{\uparrow}$ be the lifted variables and the lifted bond term of
Lemma~\ref{4.22:lem-lifted-cutoff}. For a real slot $x$ on $\mathfrak{S}^{22}$ and a
$G^{\mathbb{C}}$-valued map $g$ on the law sites $\mathfrak{s}(S)$ of the stored term's domain
$S$ (Lemma~\ref{4.22:lem-orbit-datum-properties}), write $x^{[g]}$ for the rotated slot,
$x^{[g]}(b)=\operatorname{Ad}_{g(b_-)}x(b)$. This is the first of the three substitutions of
Corollary~\ref{4.22:cor-cutoff-coordinate}.
\begin{enumerate}
\item[(1)] The function $\mathsf{g}$ has no holomorphic extension to a tube around the real
slots. One might think that one of two things is then needed: either no live bond of
$\mathfrak{S}^{22}$ lies at a site where $g$ is
not $G$-valued, or
the presentation of the rotated terms leaves the slots on such bonds unrotated. Neither is
needed.
\item[(2)] The value defined by the substitutions of
Corollary~\ref{4.22:cor-cutoff-coordinate} depends on the complex slot $x^{[g]}$ only. No
choice is hidden in it.
\item[(3)] The same three substitutions serve every dependence on the slot of the form
$\Psi(x,\varphi_1(x),\dots)$, where $\Psi$ is holomorphic and each $\varphi_i$ is continuous
and $\operatorname{Ad}_G$-equivariant, that is,
$\varphi_i(\operatorname{Ad}_u y)=\operatorname{Ad}_u\varphi_i(y)$ for $G$-valued $u$, with
$\operatorname{Ad}$ acting by conjugation on the range of $\varphi_i$, a space of $N\times N$
matrices or of tuples of them, an invariant scalar value being read as a multiple of the
identity.
\item[(4)] No derivative of $\mathsf{g}$ of order higher than one ever occurs.
\item[(5)] The anchored totals, that is, the sums of the right-hand sides of the term bounds
over all domains containing a fixed bond, used in the proofs of clause~(iv) of
Lemma~\ref{4.22:lem-lifted-cutoff} and of Lemma~\ref{4.22:lem-cutoff-large-field} are to be
read class by class. For the plain class $\mathfrak{A}$ they are the totals of the
correlation theorem. For the class $\mathfrak{S}$ of terms touched by a second-class outcome
they are the totals of the bound (r3) among the hypotheses
\eqref{4.22:eq-extension-hypotheses-a} of the holomorphic extension.
\end{enumerate}
\end{remark}

\begin{proof}
(1) The function $\mathsf{g}$ takes two different constant values on its two plateaux, the
two closed sets of Definition~\ref{4.22:def-invariant-cutoff}, each of which has non-empty
interior. A real-analytic function on the connected space $\mathfrak{g}$ that is constant on a
non-empty open set is constant. Hence $\mathsf{g}$ is not real-analytic and has no holomorphic
extension to any tube around the real slots. The presentation, however, does not ask for
$\mathsf{g}$ at a complex argument. For a
unit $\mathfrak{E}$ it asks for the map
$\mathfrak{m}\mapsto\mathfrak{E}(\dots;x^{[g(\mathfrak{m})]})$ of the complex parameter
$\mathfrak{m}$ of the presentation, ranging over its parameter domain $\mathfrak{M}$,
at a fixed real slot $x$. Along a $G^{\mathbb{C}}$-orbit the invariants of $x$ do not move, so
$\mathsf{g}$ is read at the real slot $x$. Accordingly, in the notation of
Corollary~\ref{4.22:cor-cutoff-coordinate}, one has
$\mathsf{s}_b=\mathsf{g}(x(b))$, which does
not depend on the rotation at all. Moreover $\mathsf{B}(b)=\operatorname{Ad}_{g(b_-)}x(b)$ and
$\mathsf{v}_b=\operatorname{Ad}_{g(b_-)}\nabla\mathsf{g}(x(b))$ are polynomial in $g^{\pm1}$.

(2) Let $S$ be the domain of the stored term and $\mathfrak{s}(S)$ its law sites, as in
Lemma~\ref{4.22:lem-orbit-datum-properties}. By Corollary~\ref{4.22:cor-cutoff-coordinate}
the value defined by the substitutions is an orbit datum of radius $\varepsilon_{\star}$ on
$\mathfrak{G}_{\varepsilon_{\star}}(S)$, so Lemma~\ref{4.22:lem-orbit-datum-properties}
applies with $\varepsilon=\varepsilon_{\star}$. Let $\mathbf{A},\mathbf{A}'$ be real pointwise
coordinates, the slot $x$, respectively $x'$, being one of the coordinates of $\mathbf{A}$,
respectively $\mathbf{A}'$, let $\boldsymbol{\Phi},\boldsymbol{\Phi}'$ be real frames, and let
$g,g'\in\mathfrak{G}_{\varepsilon_{\star}}(S)$. Suppose that the rotated data coincide:
\begin{equation*}
(\mathbf{A}^{[g]},\boldsymbol{\Phi}^{\langle g\rangle})
=(\mathbf{A}'^{[g']},\boldsymbol{\Phi}'^{\langle g'\rangle}).
\end{equation*}
Property \eqref{4.22:eq-orbit-datum-properties-b} of
Lemma~\ref{4.22:lem-orbit-datum-properties}, applied to the two data, gives equal values of
the orbit datum $c^{\sharp}$. Hence the value depends on the complex slot alone, and not on
the real slot and the rotation chosen to represent it.

(3) The substitutions of Corollary~\ref{4.22:cor-cutoff-coordinate} evaluate $\mathsf{g}$ and
$\nabla\mathsf{g}$ at the real slot $x(b)$. They rotate the slot by
$\operatorname{Ad}_{g(b_-)}$, rotate the equivariant value $\nabla\mathsf{g}(x(b))$ in the same
way, and leave the invariant value $\mathsf{g}(x(b))$ unchanged.

Put $\varphi_i$ in place of $\mathsf{g}$ and $\nabla\mathsf{g}$, with each value
$\varphi_i(x(b))$ rotated by the action of $\operatorname{Ad}_{g(b_-)}$ on the range of
$\varphi_i$. Here, as in item~(3), $\operatorname{Ad}$ acts by conjugation on the range of
$\varphi_i$; for
$G^{\mathbb{C}}$-valued $g$ the rotation $\operatorname{Ad}_{g(b_-)}$ is again conjugation,
polynomial in $g^{\pm1}$. The resulting value has three properties:
\begin{itemize}
\item it is holomorphic in $g$, because $\Psi$ is holomorphic and the rotated arguments are
polynomial in $g^{\pm1}$;
\item it is continuous in $x$, because the $\varphi_i$ are continuous;
\item for $G$-valued $g$ it equals $\Psi$ at the real datum $x^{[g]}$, by the equivariance of
the $\varphi_i$.
\end{itemize}

(4) In the proof of Lemma~\ref{4.22:lem-lifted-cutoff}, the bracket of
\cite[(1.24)]{B-CMP119} is the Jacobian
\begin{equation*}
\mathsf{J}_b
=\partial_a\mathfrak{c}+\partial_\eta\mathfrak{c}[H]
\langle\nabla\mathsf{g}(A'(b)),\cdot\rangle ,
\end{equation*}
where $\mathfrak{c}(a,\eta)$ is the logarithm map of the lift, $H=H^{(k)}(b)$, and $A'(b)$ is
the new bond variable of the change of variables in Lemma~\ref{4.22:lem-lifted-cutoff}.
It contains $\nabla\mathsf{g}$ once and no higher derivative. After the lift,
$\nabla\mathsf{g}$ enters only through the variable $\mathsf{v}_b$. Every later estimate
treats $\mathsf{v}_b$ as an independent variable, so no further derivative of $\mathsf{g}$
arises.

(5) The two anchored sums of right-hand sides over the domains through one bond $b$ are
\begin{align}
\sum_{Y\ni b}\nu(Y)&\le C_{\kappa}\varepsilon_V ,\label{4.22:eq-cutoff-remarks-1}\\
\sum_{e:\,\mathsf{b}(e)\ni b}\nu^{\natural}(e)&\le\tfrac12 .\label{4.22:eq-cutoff-remarks-2}
\end{align}
In the first, $\nu(Y)$ is the right-hand side of the bound for the term on the domain $Y$,
$C_{\kappa}$ is the constant and $\varepsilon_V$ the smallness parameter of the term bounds of
the correlation theorem, and the sum is the anchored total of these right-hand sides in that
theorem. In the second, $\nu^{\natural}(e)$ bounds the cluster element $e$ on its support
$\mathsf{b}(e)$; the sum is the anchored total of the cluster-element majorants of the
localised form at a large-field stage, valid under the smallness condition of that form, and
it is quoted in the proof of clause~(iv) of Lemma~\ref{4.22:lem-cutoff-large-field}.
The bound \eqref{4.22:eq-cutoff-remarks-2} also enters the storage clause of the lemma fixing
the complex radius $\varrho^B_j$ of the weak collar variable, $j$ being the level of the stage
in Lemma~\ref{4.22:lem-cutoff-large-field}, and it recurs in a later section of this chapter.

For $\mathfrak{A}$, the term $f_b^{\uparrow}$ is a single sum over $Y\ni b$ of terms of both
classes. The assignment of each summand to its domain $Y$ in
Corollary~\ref{4.22:cor-cutoff-coordinate} sorts this sum by class. The $\mathfrak{A}$-part is
a sub-sum of a sum of non-negative majorants, and is therefore bounded by the total of the
correlation theorem.

For $\mathfrak{S}$, the relative term $V^{\mathrm{rel}}(Y)$ of the class $\mathfrak{S}$ is
indexed by its label $Y$, a domain in the class $\mathbf{D}_i$ of domains at level $i$. It
obeys the bound $\nu(Y)$ at the plain distance
$d_i(Y)$ of such $Y$. The sum over $Y\ni b$ is therefore the sum of \cite[(1.26)]{B-CMP119},
and no count of fibres enters. The fibres of the labelling map, that is, the sets of terms of
the class $\mathfrak{S}$ that share a label, are controlled by the fibre bound imposed at the
step that created these terms.

The constant $C_{\kappa}$ and the constant of \eqref{4.22:eq-cutoff-remarks-2} are adjusted
to the class of the term. Clause (ii) of the condition on live rotations is used in the same
form for both classes. These
sums, the bound \eqref{4.22:eq-cutoff-remarks-2} in its later use, and clause (ii) of the
corollary on the complexified terms are the only places where the proofs of the holomorphic
extension quote a sum or a comparison of lengths.
\end{proof}

Throughout this subsection $j$ is a step of the large-field stages of \cite{B-CMP122a}. The indices $n$ and $k$, the coupling $g_j$ and the parameter $\delta'_j$ are those of that paper. Ba{\l}aban introduces the collar variable $B'_j$ as follows (\cite[(1.53)]{B-CMP122a}):
\begin{quote}
$V'_j=V_j(V^{(j)})^{-1}$, $V^{(j)}=M_j(U^{(n+1)}_k)$, $B'_j=(1/i)\log V'_j$ \dots\ Thus
$B'_j$ is small, $|B'_j|<6\delta'_j$, and we expand all the expressions in the integral with
respect to $B'_j$, we linearize the expressions in the $\delta$-functions, we remove the
$\delta$-functions using the operator $C$, and finally we perform the scaling transformation
$B'_j=g_jB_j$
\end{quote}
He then splits each component of $Z$ (\cite[(1.55)]{B-CMP122a}):
\begin{quote}
$1=\chi'^{(j)}+(1-\chi'^{(j)})$ in each component of $Z$ \dots\ We exclude from $Z$ the
components with the large field function $1-\chi'^{(j)}$
\end{quote}
We call \emph{arrested} a component of $Z$ treated in the branch of \cite[(1.55)]{B-CMP122a} whose variables are $B'_j$, with $|B'_j|<6\delta'_j$ as in \cite[(1.54)]{B-CMP122a}, together with $V_{j+1}$. We write $\mathfrak{a}$ for such an arrest.

By \cite[(1.60)]{B-CMP122a} the terms of the stage's action on $\mathfrak{a}$ are evaluated at
\begin{equation*}
U^{(n)}_k\bigl(\exp i[g_jCB_j-h\tilde{D}(g_jCB_j)]\,V^{(j)}\bigr).
\end{equation*}
Here $C$ is the operator of \cite[(1.53)]{B-CMP122a} with which the $\delta$-functions are removed, and $h$ and $\tilde D$ are as in \cite[(1.60)]{B-CMP122a}. Put
\begin{equation}\label{4.22:eq-weak-collar-1}
b(\mathsf{B}):=g_jC\mathsf{B}-h\tilde{D}(g_jC\mathsf{B}).
\end{equation}
Every constituent $c$ of the stage's action on $\mathfrak{a}$ is stored as its \emph{composite at $\mathfrak{a}$}:
\begin{equation}\label{4.22:eq-weak-collar-2}
(P,\mathsf{B})\longmapsto c\bigl(S;\Phi_{\mathbb{1},\sigma}(b(\mathsf{B}))[P];\,\cdot\,\bigr).
\end{equation}
Here $\Phi_{t,\sigma}(b)[P]$ is the background map of the inheritance theorem at the pair $P$ with amplitude $b$.

The old terms at an arrest enter through a first-order expansion in a shift. They are carried on the domain $\mathfrak{S}^{(k_a)}_{n+1}$, which consists of a shift together with a Cauchy margin. We write $\mathfrak{S}^{22}$ for the set of bonds of the full change of variables; on these bonds the collar variables are silent coordinates (Definition~\ref{4.22:def-orbit-datum}).

For a map $g$ from law sites to $G^{\mathbb{C}}$ and a family $\mathsf{B}=(\mathsf{B}(b))$ indexed by bonds, write
\begin{equation*}
\mathsf{B}^{[g]}(b):=\operatorname{Ad}_{g(b_-)}\mathsf{B}(b).
\end{equation*}

\begin{lemma}[the arrested collar variable is a weak coordinate]\label{4.22:lem-weak-collar}
\begin{enumerate}
\item[(i)] Let $c$ be a stored term in the sense of Definition~\ref{4.22:def-orbit-datum}.
Assume that $c$ depends, besides $(P,\mathbf{A},\boldsymbol{\Phi})$, on $\mathfrak{g}$-valued variables $\mathsf{B}=(\mathsf{B}(b))$, and that the law site of $\mathsf{B}(b)$ is $b_-$. Assume further:
\begin{itemize}
\item $c$ extends holomorphically to $|\mathsf{B}(b)|<\varrho^{\mathsf{B}}$, jointly with $P$ and the weak coordinates of $\mathbf{A}$, and is bounded there in modulus by $\mathfrak{N}$;
\item $c$ is jointly $G$-invariant with $\mathsf{B}\mapsto\mathsf{B}^{[v]}$, that is, $c$ is unchanged under $(P,\mathbf{A},\mathsf{B},\boldsymbol{\Phi})\mapsto (P^v,\mathbf{A}^{[v]},\mathsf{B}^{[v]},\boldsymbol{\Phi}^{[v]})$ for every $G$-valued site map $v$;
\item the real range of the variables is $|\mathsf{B}(b)|\le\mathfrak{b}^{\mathsf{B}}$, where
\begin{equation}\label{4.22:eq-weak-collar-3}
e^{2\varepsilon_\star}\,\mathfrak{b}^{\mathsf{B}}\le\varrho^{\mathsf{B}}.
\end{equation}
\end{itemize}
Then the $\mathsf{B}(b)$ are weak coordinates in the sense of Definition~\ref{4.22:def-orbit-datum} at the radius $e^{-2\varepsilon_\star}\varrho^{\mathsf{B}}$. More precisely, the function
\begin{equation}\label{4.22:eq-weak-collar-4}
c^{\sharp}(P;(\mathbf{A},\mathsf{B}),\boldsymbol{\Phi},g)
:=c^{\sharp}\bigl(P;(\mathbf{A},\mathsf{B}^{[g]}),\boldsymbol{\Phi},g\restriction\mathrm{rest}\bigr)
\end{equation}
is an orbit datum of radius $\varepsilon_\star$ at the law sites of the $\mathsf{B}(b)$, and $\mathfrak{N}^{\mathbb{C}}\le\mathfrak{N}$. On the right of \eqref{4.22:eq-weak-collar-4}, $c^\sharp$ is the orbit datum of $c$ in its remaining coordinates, and $g\restriction\mathrm{rest}$ is the restriction of $g$ to their law sites.
\item[(ii)] Let $\mathfrak{a}$ be an arrest at step $j$, and let $C_1$ be the constant of \cite[(1.19)]{B-CMP116}, so that $|b(\mathsf{B})|\le C_1g_j\sup|\mathsf{B}|$. Assume
\begin{equation}\label{4.22:eq-weak-collar-a}
\begin{gathered}
\text{the ceilings on }\gamma\text{ and the amplitude window of the inheritance theorem}\\
\text{hold with its amplitude }1.1\,C_1\delta'_j\text{ replaced by }6.6\,C_1\delta'_j,\\
\text{and }\frac{0.6\,\delta'_j}{g_j}\ge 3 .
\end{gathered}
\end{equation}
Then part (i) applies, off the bonds of $\mathfrak{S}^{22}$, to the composite \eqref{4.22:eq-weak-collar-2} of every constituent of the stage's action at $\mathfrak{a}$, with
\begin{equation}\label{4.22:eq-weak-collar-5}
\mathfrak{b}^B_j:=\frac{6\,\delta'_j}{g_j},\qquad \varrho^B_j:=\frac{6.6\,\delta'_j}{g_j}.
\end{equation}
The composite lies on the chain domain $\mathfrak{D}^{\flat}_{n+1}(\cdot)$, respectively (for the old terms at the arrest) on $\mathfrak{S}^{(k_a)}_{n+1}$, with the bound of its parent term.

If the parent is a kept member $v\in\mathfrak{K}(Y)$, the composite is the sum of the term of its expansion at $\mathsf{B}=0$ and of the grouped terms $\sum_{\hat e}v_{\hat e}$. It is estimated through part (a) of the proposition on the expansion of kept members, at the amplitude $6.6\,C_1\delta'_j$ of \eqref{4.22:eq-weak-collar-a}. This is legitimate because the bounds of the lemma on the background map, part (i), are homogeneous of degree one in $\sup|b|$. It is not estimated through $\mathfrak{N}^{\mathbb{C}}(v)$, nor through Ba{\l}aban's bound on the kept member.
\end{enumerate}
\end{lemma}

\begin{proof}
(i) Let $g\in\mathfrak{G}_{\varepsilon_\star}$, so that $\operatorname{pol}(g(b_-))<\varepsilon_\star$, and let $|\mathsf{B}(b)|<e^{-2\varepsilon_\star}\varrho^{\mathsf{B}}$. By \eqref{4.22:eq-polar-decomposition-b},
\begin{equation*}
|\mathsf{B}^{[g]}(b)|\le e^{2\operatorname{pol}(g(b_-))}|\mathsf{B}(b)|
<e^{2\varepsilon_\star}e^{-2\varepsilon_\star}\varrho^{\mathsf{B}}=\varrho^{\mathsf{B}}.
\end{equation*}
Hence the right side of \eqref{4.22:eq-weak-collar-4} is evaluated inside the ball on which $c$ extends holomorphically and is bounded by $\mathfrak{N}$, so $\mathfrak{N}^{\mathbb{C}}\le\mathfrak{N}$. By \eqref{4.22:eq-weak-collar-3}, the real range $|\mathsf{B}(b)|\le\mathfrak{b}^{\mathsf{B}}$ lies in the closed ball of radius $e^{-2\varepsilon_\star}\varrho^{\mathsf{B}}$.

The map $\mathsf{B}^{[g]}(b)=g(b_-)\mathsf{B}(b)g(b_-)^{-1}$ is polynomial in the entries of $g(b_-)$ and $g(b_-)^{-1}$, and linear in $\mathsf{B}(b)$. Composing it with the holomorphic $c^\sharp$ in the remaining coordinates therefore gives property \eqref{4.22:eq-orbit-datum-a}. For $G$-valued $g$, property \eqref{4.22:eq-orbit-datum-b} is the definition \eqref{4.22:eq-weak-collar-4} combined with property \eqref{4.22:eq-orbit-datum-b} of the orbit datum in the remaining coordinates.

(ii) We first establish holomorphy. By part (i) of the lemma on the background map of the inheritance theorem, the map
\begin{equation*}
(\mathbb{U},\mathbb{J},b,\sigma,t)\mapsto\Phi_{t,\sigma}(b)
\end{equation*}
is holomorphic for $b$ complex in its ball. Moreover, every bound of the first three steps of its construction is homogeneous of degree one in $\sup|b|$: the constant $\varkappa_{\mathsf{w}}$, the bounds on the increments of the entries of the background map, and the remaining bounds of those steps are all linear in the amplitudes.

The first member of \eqref{4.22:eq-weak-collar-a} is a re-reading of the ceilings on $\gamma$ at the larger amplitude. Under it, part (a) of the inductive hypothesis on the domains at stage $n+1$ places $\Phi_{\mathbb{1},\sigma}(b)\restriction S$ in the domain of $c$ for all complex $b$ with $|b|\le 6.6\,C_1\delta'_j$. Now let $|\mathsf{B}|<\varrho^B_j$. Then
\begin{equation*}
|b(\mathsf{B})|\le C_1g_j\sup|\mathsf{B}|<C_1g_j\cdot\frac{6.6\,\delta'_j}{g_j}=6.6\,C_1\delta'_j .
\end{equation*}
So the composite \eqref{4.22:eq-weak-collar-2} is holomorphic for $|\mathsf{B}|<\varrho^B_j$. Its values are values of $c$, so it carries the parent's bound.

Next, the real range. On the real range, after the $\delta$-functions are removed, $B'_j=g_jCB_j$, and by \cite[(1.53)]{B-CMP122a} $|B'_j|<6\delta'_j$ on all bonds. On the bonds where $B_j$ is an independent variable we have $(CB_j)(b)=B_j(b)$, so there
\begin{equation*}
|B_j(b)|<\frac{6\,\delta'_j}{g_j}=\mathfrak{b}^B_j .
\end{equation*}
With $\varepsilon_\star$ as in \eqref{4.22:eq-rotation-radius-a}, $e^{2\varepsilon_\star}\le 1+1/160$, and therefore
\begin{equation*}
6\,e^{2\varepsilon_\star}\le 6\Bigl(1+\frac{1}{160}\Bigr)<6.6 .
\end{equation*}
This gives $e^{2\varepsilon_\star}\mathfrak{b}^B_j<\varrho^B_j$, which is \eqref{4.22:eq-weak-collar-3}.

The amplitude $1.1\,C_1\delta'_j$ of the inheritance theorem corresponds to the case $|B_j|<\delta'_j/g_j$ of the cut-off $\chi'^{(j)}$. At an arrest the factor $6$ of \cite[(1.54)]{B-CMP122a} enters instead.

Finally, invariance. The Landau layer $\mathbb{H}$ and the current datum $\mathcal{J}^{\sharp}[1]$ from which the background map is built are $\operatorname{Ad}$-covariant, by part (ii) of the recursive hypothesis of the presentation. The operators $C$ and $\tilde D$ are built from covariant operators at the background \cite[(3.33)--(3.34)]{B-CMP99b}. Hence
\begin{equation*}
b(\mathsf{B}^{[v]})[P^v]=\operatorname{Ad}_v\,b(\mathsf{B})[P].
\end{equation*}
\end{proof}

At the later sites of the step the composites are kept in localised form. Let $v_e$ be the elements of the localised expansion of the stage. Each $v_e$ has domain $\mathsf{d}(e)$, depends on $\mathsf{B}$ only through its values on $\mathsf{b}(e)$, and satisfies $|v_e|\le\nu^{\natural}(e)$ on the complex ball. Their anchored sums are at most $\tfrac12$, by the anchored-sum bound of the localised expansion. With these elements, $\sum_e v_e$ is the composite minus its value at $\mathsf{B}=0$ and at vanishing rim slots, $A_j=0$.

\begin{definition}[Tent functions and the near-identity gauge map on a frame]
\label{4.22:not-tent-gauge-map}
We recall the place in Ba{\l}aban's construction to which the following objects are attached.
In the large-field step of \cite{B-CMP122b} the components carrying the second large-field
function, that is, the large-field characteristic functions $\chi^c_k(Y_i^{\sim 6})$ together with
the previous characteristic functions $\chi_{k,\Lambda_i}$, $\chi_i$ and the $\delta$-functions
$\delta_{G_i}(V'_k)$, are included again into the large-field domain of the new action, the
integral expression in \cite[(1.70)]{B-CMP122b} being transformed further. The operation
$\mathbf{T}'_k(X)$ acts on a component $X$ of either class, the sum in its definition
\cite[(1.71)]{B-CMP122b} acts also on the last exponential in \cite[(1.72)]{B-CMP122b}, the new
large-field region is $Z_k\cup\bigcup_i Y_i$, and the operation $\mathbf{T}_k$ on a component has
the form \cite[(1.102)]{B-CMP122b}. Consequently, in \cite[(1.102)]{B-CMP122b} the new boundary
terms $\mathbf{B}'^{(k)}(X)$ with $X\supset Y_i$ remain under $\mathbf{T}'_k(Y_i)$, and they are
functions of $U_k$, of the frame data $\Phi_{Y_i}$ and of the surviving exterior variables.

A \emph{frame} $W$ is a component of the second class; $\Phi_W$ denotes its data and $\Sigma_W$ its
law sites, a finite set of distinct sites of the unit lattice of level $k_W$, so that they are
pairwise separated by at least the spacing of that lattice; here $k_W$ is the step at which
$W$ is created. Let $N$ be a term created at the large-field operation $\mathbf{R}$ of step $k_W$ whose
domain $X$ contains $W$. This includes the two members of the kept lineage
$\mathfrak{K}(Y_i)$ under $\mathbf{T}'_{k_W}(Y_i)$ in \cite[(1.102)]{B-CMP122b}: the term
$V(Y_i^{\sim})$, and the quadratic form $\mathsf{Q}_i$ of \cite[(1.71)]{B-CMP122b}, which is taken in
the localised form $\mathsf{Q}_i:=\mathfrak{q}_C(\zeta;\varpi)$ with $C=Y_i$ of the definition of the
kept quadratic form; at real $\zeta$ this is the form
\begin{equation*}
-\tfrac12\bigl\langle DH''_{1,k_W,C}B,\ \hat\chi\, DH''_{1,k_W,C}B\bigr\rangle
\end{equation*}
of \cite[(1.71)]{B-CMP122b}, with the $C$-localised operators of that definition; here $B$ is the
variable, $D$ the derivative and $\hat\chi$ the characteristic factor as they appear in
\cite[(1.71)]{B-CMP122b}. By
\cite{B-CMP122b}, only the quadratic form in $B$ and the term $V(X^{\sim})$ depend on
$(\mathbf{U},\mathbf{J})$. The domain $X$ may contain one or several frames; they are treated
independently, their sets $\mathsf{B}_W$ defined below being disjoint.

Let $\ell$ be the spacing of the unit lattice of level $k_W$, that is, $L^{k_W}$ times the spacing
of the finest lattice. Let $\mathfrak{U}^{\mathrm{tr}}_N(X)$ be the space on which
$N$ is holomorphic and bounded, as established at the large-field operation of step $k_W$, for every
real value of the
frame data $\boldsymbol{\Phi}$ in the supports of its characteristic functions; by
\cite[(1.68), (1.99)]{B-CMP122b} this is the tube
$\tilde{U}^{c}_{k_W}(X,\tilde\alpha_0,\tilde\alpha_1)$. It is described by eight clauses,
\begin{equation}\label{4.22:eq-tent-gauge-map-1}
\mathfrak{U}^{\mathrm{tr}}_N(X)
=\bigl\{\,Q_a(y)<\theta_a(y)\ \text{for } a=1,\dots,8 \text{ and all clause locations } y\,\bigr\},
\end{equation}
where $Q_a(y)$ is the quantity of clause $a$ at the location $y$, $\theta_a(y)$ its threshold and
$r_a(y)$ its radius; clause $5$ bounds $|U(\partial p)-1|$ for the unitary part $U$, clauses $6$ and
$7$ bound $\ell|A'|$ and $\ell^2|\nabla_U A'|$ for the $\mathfrak{g}^{\mathbb{C}}$-valued field $A'$,
and clause $8$ is the cube clause. After the large-field operation of step $k_W$ the frame $W$ lies in
the new domain of level $k_W$ of Ba{\l}aban's sequence of domains in \cite{B-CMP122b}, because the
determining set $\mathbf{B}_k$ restricted to the large-field region is the unit lattice of level $k$
\cite[(1.33)]{B-CMP122b}; hence on $W$ the layer is $n=k_W$ and $\theta_a=r_a$.

The term $N$, each of the two kept members included, is moreover jointly invariant under every
$G$-valued gauge transformation $g$: its value does not change when $g$ acts simultaneously on its
field argument $\zeta$, giving $\zeta^{g}$, and, through the induced transformation $\tilde g$, on the
frame data. For the new boundary terms this is clause~(ii) of the gauge-invariance theorem for the
boundary terms,
\begin{equation*}
\mathbf{B}'^{(k)}(X;\zeta^{g};\boldsymbol{\Phi}^{\tilde g})
=\mathbf{B}'^{(k)}(X;\zeta;\boldsymbol{\Phi})
\qquad\text{for every $G$-valued } g,
\end{equation*}
and for $\mathsf{Q}_i$ the holomorphy, the bound and the invariance are supplied by the statements on
the kept quadratic form. Of $N$ only these three properties, holomorphy and a bound on
$\mathfrak{U}^{\mathrm{tr}}_N(X)$ and joint invariance under every $G$-valued $g$, are used in what
follows.

\emph{Tent functions.} For $s\in\Sigma_W$ let $\chi_s$ be the smoothed tent functions of the unit
lattice of level $k_W$: real-valued functions on the lattice sites such that
\begin{itemize}
\item $\chi_s(s')=\delta_{ss'}$ for all unit sites $s'$, and $\sum_{s}\chi_s\le 1$;
\item $\operatorname{supp}\chi_s\subset\{x:\ \operatorname{dist}(x,s)<\ell\}$;
\item $\ell|\nabla\chi_s|\le C_\chi$ and $\ell^2|\nabla\nabla\chi_s|\le C_\chi$;
\item at most $3^d$ of the functions $\chi_s$ overlap at any point.
\end{itemize}
Put
\begin{equation}\label{4.22:eq-tent-gauge-map-2}
\mathsf{B}_W:=\bigcup_{s\in\Sigma_W}\operatorname{supp}\chi_s .
\end{equation}

\emph{Decompositions.} A \emph{decomposition} $\mathsf{r}=(U,\{u_{\square}\})$ of a pair
$P=(\mathbb{U},\mathbb{J})$ consists of a factorisation $\mathbb{U}=U'U$, with $U$ the unitary part,
together with the cube gauges $u_{\square}$ of $U$ in the sense of \cite[(2.38)]{B-CMP119}: for every
cube $\square$ of size $C^{119}ML^n\xi$ the $G$-valued gauge transformation $u_\square$ on $\square$
satisfies $U^{u_\square}=\exp(i\xi A)$ with
\begin{equation*}
L^n\xi\,|A|<B^{119}C^{119}M\alpha_{0,n},\qquad (L^n\xi)^2\,|\nabla A|<B^{119}C^{119}M\alpha_{0,n},
\end{equation*}
where $\xi$ is the lattice spacing and $B^{119}$, $C^{119}$ are the constants of that display. For each
$s\in\Sigma_W$ fix one such cube $\square_s$ with $\square_s\supset\operatorname{supp}\chi_s$.

\emph{The gauge map.} For $Z=(Z_s)_{s\in\Sigma_W}$ with $Z_s\in\mathfrak{g}^{\mathbb{C}}$ define, at
every lattice site $x$,
\begin{equation}\label{4.22:eq-tent-gauge-map-3}
\mathsf{Y}_{\mathsf{r}}[Z](x)
:=\sum_{s\in\Sigma_W}\chi_s(x)\,
\operatorname{Ad}_{u_{\square_s}(x)^{-1}u_{\square_s}(s)}Z_s,
\qquad
\rho_{\mathsf{r}}[Z]:=\exp\bigl(-\mathsf{Y}_{\mathsf{r}}[Z]\bigr),
\end{equation}
the summand with index $s$ being read as $0$ at sites $x\notin\operatorname{supp}\chi_s$ (at the
remaining sites $x\in\operatorname{supp}\chi_s\subset\square_s$, so $u_{\square_s}(x)$ is defined).
Then, writing $\rho=\rho_{\mathsf{r}}[Z]$:
\begin{enumerate}
\item[(a)] $\rho(s)=e^{-Z_s}$ for every $s\in\Sigma_W$;
\item[(b)] $\rho(x)=1$ for every site $x\notin\mathsf{B}_W$;
\item[(c)] if $Z_s\in i\mathfrak{g}$ for every $s\in\Sigma_W$, then $\rho$ is $G$-valued;
\item[(d)] at frozen $\mathsf{r}$, $\rho(x)$ is a holomorphic function of
$Z\in(\mathfrak{g}^{\mathbb{C}})^{\Sigma_W}$ for every site $x$.
\end{enumerate}
\end{definition}

\begin{proof}
Throughout, $\mathsf{r}$ is fixed and we abbreviate $h_s(x):=u_{\square_s}(x)^{-1}u_{\square_s}(s)$
for $x\in\operatorname{supp}\chi_s$. Since $\chi_s(s)=1$, the site $s$ lies in
$\operatorname{supp}\chi_s\subset\square_s$, so $u_{\square_s}(s)$ is defined and $h_s(x)$ makes sense.
Recall that $\operatorname{Ad}_h Z=hZh^{-1}$.

(a) Let $s\in\Sigma_W$. Every $s'\in\Sigma_W$ is a site of the unit lattice of level $k_W$, and $s$ is
one as well, so $\chi_{s'}(s)=\delta_{s's}$ by the first property of the tent functions. Hence in the
sum \eqref{4.22:eq-tent-gauge-map-3} evaluated at $x=s$ only the summand $s'=s$ survives, with
coefficient $\chi_s(s)=1$. Its group element is
$h_s(s)=u_{\square_s}(s)^{-1}u_{\square_s}(s)=\mathbb{1}$, and $\operatorname{Ad}_{\mathbb{1}}$ is the
identity. Therefore $\mathsf{Y}_{\mathsf{r}}[Z](s)=Z_s$, and $\rho(s)=\exp(-Z_s)$.

(b) Let $x\notin\mathsf{B}_W$. By \eqref{4.22:eq-tent-gauge-map-2}, $x\notin\operatorname{supp}\chi_s$
for every $s\in\Sigma_W$, so every summand in \eqref{4.22:eq-tent-gauge-map-3} is $0$. Hence
$\mathsf{Y}_{\mathsf{r}}[Z](x)=0$ and $\rho(x)=\exp(0)=\mathbb{1}$.

(c) Assume $Z_s\in i\mathfrak{g}$ for all $s$, so each $Z_s$ is anti-Hermitian and traceless. Fix $x$
and $s$ with $x\in\operatorname{supp}\chi_s$. The gauge transformation $u_{\square_s}$ is $G$-valued, so
$h_s(x)\in G$; since the elements of $G$ are unitary, $h_s(x)^{-1}=h_s(x)^*$. Then
\begin{equation*}
\bigl(h_s(x)Z_sh_s(x)^{*}\bigr)^{*}=h_s(x)Z_s^{*}h_s(x)^{*}=-\,h_s(x)Z_sh_s(x)^{*},
\end{equation*}
and $\operatorname{tr}\bigl(h_s(x)Z_sh_s(x)^{-1}\bigr)=\operatorname{tr}Z_s=0$, so
$\operatorname{Ad}_{h_s(x)}Z_s\in i\mathfrak{g}$. The coefficients $\chi_s(x)$ are real and
$i\mathfrak{g}$ is a real vector space; hence $Y:=\mathsf{Y}_{\mathsf{r}}[Z](x)\in i\mathfrak{g}$. For
anti-Hermitian $Y$ one has
\begin{equation*}
(e^{-Y})^{*}=e^{-Y^{*}}=e^{Y}=(e^{-Y})^{-1},
\end{equation*}
so $e^{-Y}$ is unitary; and $\det e^{-Y}=1$ because $Y$ is traceless. Thus $\rho(x)=e^{-Y}$ is unitary
of determinant one, that is, $\rho(x)\in G$.

(d) For fixed $x$ and $s$, the map $Z_s\mapsto\operatorname{Ad}_{h_s(x)}Z_s=h_s(x)Z_sh_s(x)^{-1}$ is
$\mathbb{C}$-linear on $\mathfrak{g}^{\mathbb{C}}$ (it is matrix multiplication by fixed matrices),
and $\chi_s(x)$ is a fixed scalar. Hence $Z\mapsto\mathsf{Y}_{\mathsf{r}}[Z](x)$ is a
$\mathbb{C}$-linear map from the finite-dimensional space $(\mathfrak{g}^{\mathbb{C}})^{\Sigma_W}$ to
the algebra $\mathfrak{A}$ of all complex matrices of the size of the elements of $G$, in particular
holomorphic. The exponential is given on $\mathfrak{A}$ by a
power series converging everywhere, so it is entire, and $\rho(x)=\exp(-\mathsf{Y}_{\mathsf{r}}[Z](x))$
is holomorphic in $Z$ as a composition of holomorphic maps.
\end{proof}

\begin{remark}
The purpose of $\rho_{\mathsf{r}}[Z]$ is the following. Rotating the frame data themselves is not
available: at complex values of the source the gauge element which is applied to them later is not
$G$-valued, and for $G$-valued data $V$ far from the identity the perturbation $cVc^{-1}V^{-1}$
produced by a rotation $c$ has size comparable to $\operatorname{pol}(c)$, which would have to be held
against analyticity radii proportional to $\alpha$. A frame, however, is rotated only at the finitely
many separated law sites of $\Sigma_W$, and by (a) and (b) a prescribed value at these sites extends
to a gauge map equal to $\mathbb{1}$ off $\mathsf{B}_W$; that this map is regular when the prescribed
values are near the identity is the content of the frame lemma which follows. The frame data therefore
remain real and unrotated, and the rotation is carried by the background of the term $N$ itself. In
\cite[(3.25)]{B-CMP109} gauge transformations that are explicitly given analytic functions of $U$ are
applied in the same way to the further variables $J$, $B$.
\end{remark}

\begin{definition}[The stored domain of a frame-carrying term]\label{4.22:def-frame-domain}
Let $N$ be a term created by the large-field operation $\mathbf{R}$ at a step, numbered $k_W$, with
each of its frames $W$ contained in the large-field region $X$. Let $P=(\mathbb{U},\mathbb{J})$
denote the background pair. Let $\mathfrak{U}^{\mathrm{tr}}_N(X)$ be the tube on which $N$ is
holomorphic and bounded,
written as the set of pairs $P$ with
\begin{equation*}
Q_a(y)<\theta_a(y),\qquad a=1,\dots,8,
\end{equation*}
at every location $y$ at which the $a$-th condition is imposed. Here $Q_a(y)$ is the $a$-th
regularity quantity at $y$, $\theta_a(y)$ is its threshold and $r_a(y)$ its radius. The entry
$a=5$ bounds the deviation from the identity of the plaquette variables of the unitary part, and
$a=8$ is the condition on cubes.

For each frame $W$, let $\mathsf{B}_W$ be the union of the supports of the tent functions about the
sites of $W$ (Notation~\ref{4.22:not-tent-gauge-map}), and let $\ell$ be the length bounding the
supports of these tent functions: the support of the tent function about a site $s$ lies in the
set of points at distance less than $\ell$ from $s$. When $N$ carries several frames, their sets
$\mathsf{B}_W$ are pairwise disjoint and the frames are treated independently. Let
$\theta_S$ be the layer number, and let $\beta$ be the schedule constant of the layered analyticity
spaces (Definition~\ref{4.4:def-post-step-space}). Put
\begin{equation}\label{4.22:eq-frame-domain-1}
\varpi_{\mathrm{fr}}:=\tfrac18\,(1-\theta_S)\,\beta .
\end{equation}
The \emph{stored domain} of $N$ is the set $\mathfrak{U}_N(X)$ of those
$P\in\mathfrak{U}^{\mathrm{tr}}_N(X)$ that satisfy
\begin{equation}\label{4.22:eq-frame-domain-2}
Q_a(y)<\theta_a(y)-\varpi_{\mathrm{fr}}\,r_a(y)
\end{equation}
for every $a\in\{1,\dots,8\}\setminus\{5,8\}$ and every location $y$ of the $a$-th condition
whose stencil, the set of lattice points on which $Q_a(y)$ depends, meets $\mathsf{B}_W$ for some
frame $W$ of $N$.

These locations are understood to include the blocks of side at most $\ell$ centred in
$\mathsf{B}_W$ on which the derived entries $a=2$ and $a=4$ are conjugated by the block-constant
gauge map $\tilde\rho$.

The term $N$ is regarded as defined on $\mathfrak{U}_N(X)$. This is a choice of domain only: the
term $N$ and the densities are not changed by it. The set $\mathfrak{U}_N(X)$ is an open subset of
$\mathfrak{U}^{\mathrm{tr}}_N(X)$ and is invariant under the action $P\mapsto P^u$ of $\mathcal{G}$.
\end{definition}

\begin{lemma}[The frame lemma: increments, orbit datum and closing]\label{4.22:lem-frame-lemma}
Let the term $N$ created at the $\mathbf{R}$-site of step $k_W$, its domain $X$, the frame
$W\subset X$ with law sites $\Sigma_W$ and data $\Phi_W$, the set $\mathsf{B}_W$, the tube
$\mathfrak{U}^{\mathrm{tr}}_N(X)$ with its clause quantities $Q_a$, thresholds $\theta_a$ and
radii $r_a$, the decompositions $\mathsf{r}=(U,\{u_{\square}\})$ of a pair $P$, and the gauge maps
$\rho_{\mathsf{r}}[Z]$ be as in Definition~\ref{4.22:not-tent-gauge-map}. Let
$\mathfrak{U}_N(X)$ be the stored domain of Definition~\ref{4.22:def-frame-domain}, with its
reserve fraction $\varpi_{\mathrm{fr}}$. Assume the following.
\begin{itemize}
\item[(F1)] The second member of \eqref{4.22:eq-coupling-ceilings-a}, read at the layer
$j=k_W$ of the frame:
\begin{equation*}
2.1\,(2C_\chi+1)\,BCM\,\bar\alpha(x_{k_W})\le 1 ,
\end{equation*}
where $B$, $C$, $M$ are the constants of the cube-gauge condition
\cite[(2.38)]{B-CMP119}. Since
\begin{equation*}
\bar\alpha(x_{k_W})=\alpha^+_{k_W}\ge\tilde\alpha_{0,k_W}\ge\alpha_{0,k_W},
\end{equation*}
where the equality is the value of $\bar\alpha$ at $x_{k_W}$, the first inequality holds because
$\alpha^+_{k}=\max\{\tilde\alpha_{0,k},\tilde\alpha_{1,k},\delta'_k\}$, and the second is the
comparison $\tilde\alpha_{0,k}\ge\alpha_{0,k}$ of the cube-gauge radius of the tube with the
radius $\alpha_{0,k}$ of \cite[(2.38)]{B-CMP119}, taken by statement,
this bound implies the first member of \eqref{4.22:eq-coupling-ceilings-a} at
$x_{k_W}$, and it covers the coefficient of the cube-gauge clause
\cite[(2.38)]{B-CMP119} of the tube
$\tilde{U}^{c}_{k_W}(\cdot,\tilde\alpha_0,\tilde\alpha_1)$, which is the clause read in (i).
\item[(F2)] The ceiling \eqref{4.22:eq-coupling-ceilings-b}, read at $j=k_W$:
$\mathfrak{c}\,\bar\alpha(x_{k_W})\le 1$, with
$\mathfrak{c}=\max\{20\bar r_6,\mathfrak{c}_6,\mathfrak{c}_7\}$ and $\bar r_6$ as in (F3) below.
\item[(F3)] The clause list of the stored tube: $\mathfrak{U}^{\mathrm{tr}}_N(X)$ is the set of
pairs $P$ with $Q_a(y)<\theta_a(y)$ for $a=1,\dots,8$ at every clause location $y$; here
entry~$5$ is the clause on $|\partial U-1|$ for the unitary part, entry~$8$ is the cube-gauge
clause of \cite[(2.38)]{B-CMP119}, entries~$6$ and~$7$ are the clauses on $\ell|A'|$ and
$\ell^2|\nabla_U A'|$, where $\ell$ is the unit length of the level-$k_W$ lattice and $A'$ is
the $\mathfrak{g}^{\mathbb{C}}$-valued field of the factor $U'$ in the decomposition
$\mathbb{U}=U'U$ of the pair, with the thresholds of \cite[(2.39)]{B-CMP119}, and entries~$1$--$4$
are the remaining four clauses of the list. On $W$ one has $\theta_a=r_a$, and
$\bar r_a:=r_a/\alpha_{1,k_W}\ge 1$ for $a=6,7$.
\item[(F4)] Part~(ii) of the orbit-datum theorem at pointwise law sites, which supplies the
function $N^{\sharp}(P;\mathbf{A},\Phi,g)$ of the pointwise rotations $g$ appearing on the
right of \eqref{4.22:eq-frame-lemma-1} below, and the polar decomposition of complex
rotations, both taken by statement.
\item[(F5)] For (iv) only: the following five room statements of this chapter for a stored
term at the locations of its frames, each taken by statement: part~(ii) of the room statement
for the regular data; the homogeneity statement; clause~(d2) of the room statement of the
variation class; the fifth room statement of the step; and the room proposition accompanying
the sixth-family lemma.
\item[(F6)] For (v) only: the classification of the outcomes of a processed component into a
first and a second class, the convention by which rotated data are read at the presentation
of a step, the arrest lemma for the large-field stages, and the statements on Haar
integration at complex rotations, all taken by statement.
\end{itemize}
Hypothesis (F1) is used for (i) only; hypothesis (F2) is used for (i) and, through the
inequality $\mathfrak{c}\ge\mathfrak{c}_6\ge 200\,\bar p_1$ with $\bar p_1=3^d(C_\chi+1)$, in the
proof of (ii). Then the following hold.
\begin{enumerate}
\item[(i)] \emph{Increments.} Let $P\in\mathfrak{U}^{\mathrm{tr}}_N(X)$, let
$z\le 10^{-2}$, and let $Z\colon\Sigma_W\to\mathfrak{g}^{\mathbb{C}}$ satisfy $|Z_s|\le z$ for
every $s\in\Sigma_W$; we say that $\rho_{\mathsf{r}}[Z]$ is a map of size $z$. With the same
decomposition $\mathsf{r}=(U,\{u_{\square}\})$, the pair $P^{\rho_{\mathsf{r}}[Z]}$ has the
following properties: entries~$5$ and~$8$ are unchanged; entries~$1$--$4$ are multiplied by
at most $e^{2z}$; no clause quantity changes at a clause location whose stencil does not meet
$\mathsf{B}_W$ in the sense of Definition~\ref{4.22:def-frame-domain}. If moreover
\begin{equation}\label{4.22:eq-frame-lemma-2}
z\le 4\,\varepsilon^{\mathrm{fr}}(W)
=\frac{\varpi_{\mathrm{fr}}\,\alpha_{*,k_W}}{C_{\mathrm{fr}}},
\end{equation}
then entry~$6$ increases by at most
$[3^d(C_\chi+1)+1]\,z\,r_6/\alpha_{1,k_W}$ and entry~$7$ by at most
$[3^d(C_\chi+2)+1]\,z\,r_7/\alpha_{1,k_W}$, and each of these is at most
$C_{\mathrm{fr}}\,z\,r_a/\alpha_{1,k_W}$ for the respective $a\in\{6,7\}$. Consequently, on the
range \eqref{4.22:eq-frame-lemma-2} a map of size $z$ increases the quantity of every clause
by at most
\begin{equation*}
\frac{z}{\varepsilon^{\mathrm{fr}}(W)}\cdot\frac{1}{4}\,\varpi_{\mathrm{fr}}\,r_a ;
\end{equation*}
in particular a map of size $z\le\varepsilon^{\mathrm{fr}}(W)$ increases every clause by at
most $\frac14\varpi_{\mathrm{fr}}r_a$, and the increments of successive maps taken at the one
decomposition $\mathsf{r}$ fixed with the stored domain add.
\item[(ii)] \emph{The orbit datum.} For
$c\in\mathfrak{G}_{\varepsilon^{\mathrm{fr}}(W)}(\Sigma_W)$ write $c(s)=e^{Y_s}\omega_s$, $s\in\Sigma_W$,
for its polar decomposition, and put
\begin{equation}\label{4.22:eq-frame-lemma-1}
N^{\sharp}\bigl(P;\mathbf{A},\Phi_W,(g,c)\bigr)
:=N^{\sharp}\bigl(P^{\rho_{\mathsf{r}}[Y]};\mathbf{A},\Phi_W^{[\omega]},g\bigr).
\end{equation}
This function does not depend on the decomposition $\mathsf{r}$. It is defined on the
product of $\mathfrak{U}_N(X)$, the polydiscs of the weak collar variables of $N$, the domain
of $g$, and
$\mathfrak{G}_{\varepsilon^{\mathrm{fr}}(W)}(\Sigma_W)$; it is holomorphic in $P$, in the weak
collar variables, in $g$ and in $c$; at $G$-valued $(g,c)$ it equals
$N(P;\mathbf{A}^{[g]},\Phi_W^{[c]})$; and
\begin{equation*}
\bigl|N^{\sharp}\bigl(P;\mathbf{A},\Phi_W,(g,c)\bigr)\bigr|
\le\sup_{\mathfrak{U}^{\mathrm{tr}}_N(X)}|N| ,
\end{equation*}
which is the right member of the bound proved for $N$ at the site that creates it. It is the
orbit datum of $N$ at $\Sigma_W$, of radius $\mathsf{e}=\varepsilon^{\mathrm{fr}}(W)$.
\item[(iii)] \emph{No spill.} No pointwise law site of $N$ and no older live frame of $N$
meets $\mathsf{B}_W$.
\item[(iv)] \emph{Rooms.} On the stored domain $\mathfrak{U}_N(X)$ every room statement of
this chapter holds at the locations of $\mathsf{B}_W$, in the form
\eqref{4.22:eq-frame-proof-datum-b}.
\item[(v)] \emph{Accounts.} Suppose that the component holding $W$ is next processed to the
end, that is, its chain of large-field stages completes and the presentation site
$(\mathrm{L})$ and the operation $\mathbf{T}'$ are reached (an arrested attempt only adds far
debits of the first kind $(f_1)$), and let $k_2$ be the step at which this happens; then
necessarily $k_W\le h_2:=k_2-\check{R}_{k_2}$, where $\check{R}_k$ is the range number of
step $k$. If the outcome is of the first class, $\Phi_W$ is integrated by the measure
$d\nu_i$ of the operation $\mathbf{T}'$ of step $k_2$ at real values. If it is of the second
class, $\Phi_W$, together with every interior
pointwise variable and every other interior frame, becomes a member of the real data
$\Phi_{W_2}$ of the new frame $W_2$, whose orbit datum is the one given by (ii) for the new
terms. In both cases $N^{\sharp}$ has been read at the rotation
$c=(h\mathfrak{q})\restriction\Sigma_W$ supplied by the presentation of
step $k_2$ and is never read at $\Sigma_W$ again: the accounts of $\Sigma_W$ and of the
interior pointwise law sites are closed at $k_2$.
\item[(vi)] Under \cite[(1.104)]{B-CMP122b} the class of terms considered here is void.
\end{enumerate}
\end{lemma}

The proof of Lemma~\ref{4.22:lem-frame-lemma}, from the hypotheses (F1)--(F6), occupies the
remainder of this section.

The proof of clause~(i) of the frame lemma (Lemma~\ref{4.22:lem-frame-lemma}) begins here. In this first part we fix the range of the clause, the conventions of the proof and the bounds on the cube gauges. We then prove the elementary exponential estimates, and the bounds on the first and second covariant differences of the tent field $\mathsf{Y}_{\mathsf{r}}[Z]$ of Notation~\ref{4.22:not-tent-gauge-map}.

Throughout, $\|\cdot\|$ is the operator norm on $\mathbb{C}^N$; it is the matrix norm $|\cdot|$ that appears in the cube-gauge and threshold conditions. For $N\times N$ matrices $\Pi$, $\Pi'$ we use $\|[\Pi,\Pi']\|\le 2\|\Pi\|\,\|\Pi'\|$, and $e^{i\eta A}\in U(N)$ for $A\in\mathfrak{g}$. No constant in what follows depends on the rank $N$ of the gauge group.

\emph{Range.} For entries 1--4, 5 and 8 clause~(i) is proved on the whole range $\|Z_s\|\le z\le 10^{-2}$ of its statement. For entries 6 and 7 it is proved on the range
\begin{equation}\label{4.22:eq-frame-proof-tent-1}
z\le z_{\max}:=4\varepsilon^{\mathrm{fr}}(W)
=\frac{\varpi_{\mathrm{fr}}\,\alpha_{*,k_W}}{C_{\mathrm{fr}}}.
\end{equation}
One has $z_{\max}<10^{-3}$. Indeed, since $\theta_S\ge\tfrac12$ and $\beta=\tfrac15$,
\begin{equation*}
\varpi_{\mathrm{fr}}=\tfrac18(1-\theta_S)\beta\le\tfrac18\cdot\tfrac12\cdot\tfrac15=\tfrac1{80};
\end{equation*}
moreover $\alpha_{*,k_W}\le 1$, and $C_{\mathrm{fr}}=3^d(C_\chi+4)\ge 15$, since $C_\chi\ge1$ (item~3 of Lemma~\ref{4.22:prf-frame-proof-tent} below). Hence
\begin{equation*}
z_{\max}\le\frac{1/80}{15}=\frac1{1200}<10^{-3}.
\end{equation*}
Every application of the clause on entries 6 and 7 in the remaining parts of the proof of the frame lemma is to a map, or a succession of maps, of total size less than $z_{\max}$. For $z_{\max}<z\le10^{-2}$ nothing is asserted about entries 6 and 7.

\emph{Conventions.} Let $T_\eta$ denote the fine lattice, of spacing $\eta$, on which the pair $P=(\mathbb{U},\mathbb{J})$ lives. A bond is $b=(x,\mu):=\langle x,x+\eta\mu\rangle$, with $b_-=x$ and $b_+=x+\eta\mu$. Put $\ell:=L^{k_W}\eta$, where $k_W\ge1$. Since the blocking factor $L$ is an integer larger than one,
\begin{equation*}
\eta/\ell=L^{-k_W}\le\tfrac12 .
\end{equation*}
This is used in the overlap count $\mathfrak{n}_\chi$ below; elsewhere only $\eta/\ell\le1$ is used. Put $\mathfrak{A}:=\mathfrak{gl}(N,\mathbb{C})$.

The gauge action is $P^g:=(\mathbb{U}^g,\operatorname{Ad}_g\mathbb{J})$, with $\mathbb{U}^g(b):=g(b_-)\,\mathbb{U}(b)\,g(b_+)^{-1}$. With the opposite convention one replaces $\rho_{\mathsf{r}}[Z]$ by its inverse $\rho_{\mathsf{r}}[-Z]$. Every bound below is even in $Z$, so nothing changes.

Let $\mathsf{r}=(U,\{u_\square\})$ be a decomposition of $P$ as in Notation~\ref{4.22:not-tent-gauge-map}. Thus $\mathbb{U}=U'U$ with $U'(b)=e^{i\eta A'_\mu(x)}$ for $b=(x,\mu)$, where $U$ is the unitary part, taken with its cube gauges. For an $\mathfrak{A}$-valued site function $f$ the covariant difference is
\begin{equation*}
(\nabla^U_\mu f)(x):=\eta^{-1}\bigl[\operatorname{Ad}_{U(x,\mu)}f(x+\eta\mu)-f(x)\bigr].
\end{equation*}
A bond field, regarded at fixed $\mu$ as a function of the site $b_-$, is differentiated by the same operator. The plain differences are $(\partial_\nu a)(x):=\eta^{-1}(a(x+\eta\nu)-a(x))$, for scalar or $\mathfrak{A}$-valued site functions $a$ and for the cube-gauge bond field $A_\mu$ below. For an operator-valued function $x\mapsto E(x)$ on $\mathfrak{A}$, $\partial_\nu E$ is defined in the same way.

\emph{Thresholds of entries 6 and 7.} Among the eight entries $Q_a$ whose bounds $Q_a(y)<\theta_a(y)$ at the clause locations $y$ define $\mathfrak{U}^{\mathrm{tr}}_N(X)$, entry 6 is $\sup\ell\|A'_\mu(x)\|$ and entry 7 is $\sup_{\mu,\nu}\ell^2\|(\nabla^U_\nu A'_\mu)(x)\|$. In both cases the supremum runs over the stencils of the clause locations. On $W$ the layer is $k_W$ and $\theta_a=r_a$. The thresholds are those of \cite[(2.39)]{B-CMP119}, proportional to the radius $\alpha_1:=\alpha_{1,k_W}$:
\begin{equation*}
r_6=\bar r_6\,\alpha_1,\qquad r_7=\bar r_7\,\alpha_1 .
\end{equation*}
Here $\bar r_6$, $\bar r_7$ are the threshold numbers of \cite[(2.39)]{B-CMP119} in units of $\alpha_{1,k_W}$. We use the normalisation $\bar r_6,\bar r_7\ge1$; if it fails, $C_{\mathrm{fr}}$ is multiplied by $\max\{1,1/\bar r_6,1/\bar r_7\}$ wherever it occurs in the parameter conditions, and nothing else changes. Thus $P\in\mathfrak{U}^{\mathrm{tr}}_N(X)$ gives
\begin{equation*}
\ell\|A'_\mu\|<r_6,\qquad \ell^2\|\nabla^U_\nu A'_\mu\|<r_7
\end{equation*}
on those stencils.

\emph{Cube gauge.} For each $s\in\Sigma_W$ we read \cite[(2.38)]{B-CMP119} as providing a cube $\square_s$ of its cube gauge that contains the $2\eta$-neighbourhood of $\operatorname{supp}\chi_s$. On $\square_s$ there are a $G$-valued $u:=u_{\square_s}$ and a $\mathfrak{g}$-valued bond field $A$, with
\begin{equation*}
U(x,\mu)=u(x)^{-1}e^{i\eta A_\mu(x)}u(x+\eta\mu).
\end{equation*}
The index $s$ of $u$ and $A$ is suppressed. They satisfy
\begin{equation}\label{4.22:eq-frame-proof-tent-2}
\ell\sup\|A\|<BCM\tilde\alpha_{0,k_W},\qquad
\ell^2\sup\max_{\mu,\nu}\|\partial_\nu A_\mu\|<BCM\tilde\alpha_{0,k_W}.
\end{equation}
Here $B$ and $C$ are the constants of \cite[(2.38)]{B-CMP119} and $M$ is the cube-size parameter of the glossary, the cubes of the cover having side $CM\ell$. Two remarks explain \eqref{4.22:eq-frame-proof-tent-2}.
\begin{itemize}
\item The derivative bound of \cite[(2.38)]{B-CMP119} is used for all difference quotients $\partial_\nu A_\mu$.
\item $\mathfrak{U}^{\mathrm{tr}}_N(X)$ is the tube $\tilde U^c_{k_W}(X,\tilde\alpha_0,\tilde\alpha_1)$, on which the cube-gauge clause holds with coefficient at most $\tilde\alpha_{0,k_W}$. Were the coefficient the radius $\alpha_{0,k_W}\le\tilde\alpha_{0,k_W}$, the bound would only improve.
\end{itemize}
Moreover $\tilde\alpha_{0,k_W}\le\alpha^+_{k_W}$, since $\alpha^+_k$ is the maximum of a list of radii at level $k$ that contains $\tilde\alpha_{0,k}$; and $\alpha^+_{k_W}=\bar\alpha(x_{k_W})$ by the first clause of the theorem that fixes the function $\bar\alpha$ entering the coupling ceilings \eqref{4.22:eq-coupling-ceilings-a}.

Put
\begin{equation*}
\lambda:=2\ell\sup\|A\|,\qquad \lambda':=2\ell^2\sup\max_{\mu,\nu}\|\partial_\nu A_\mu\|.
\end{equation*}
The second member of the coupling ceiling \eqref{4.22:eq-coupling-ceilings-a}, $2.1(2C_\chi+1)BCM\bar\alpha(x)\le1$, is a hypothesis of Lemma~\ref{4.22:lem-frame-lemma} at $x=x_{k_W}$. With \eqref{4.22:eq-frame-proof-tent-2} it gives
\begin{equation}\label{4.22:eq-frame-proof-tent-3}
\lambda,\ \lambda'<2BCM\tilde\alpha_{0,k_W}\le 2BCM\bar\alpha(x_{k_W})
\le\frac{2}{2.1(2C_\chi+1)}=:\mathsf{t}=\frac{20/21}{2C_\chi+1}\le\frac{20}{63}<\frac13 ,
\end{equation}
where the inequality $\mathsf{t}\le 20/63$ uses $C_\chi\ge1$ (item~3 below).

\emph{The tent field.} Let $E_\mu(x):=\operatorname{Ad}_{e^{i\eta A_\mu(x)}}$ and $D_\mu(x):=\eta^{-1}(E_\mu(x)-1)$. Let $Z\colon\Sigma_W\to\mathfrak{g}^{\mathbb{C}}$ with $\|Z_s\|\le z$, and put
\begin{equation*}
\tilde Z_s:=\operatorname{Ad}_{u(s)}Z_s,\qquad
f_s(x):=\chi_s(x)\operatorname{Ad}_{u(x)^{-1}}\tilde Z_s=\operatorname{Ad}_{u(x)^{-1}}\bigl[\chi_s(x)\tilde Z_s\bigr].
\end{equation*}
Here $f_s$ is defined on $\square_s$ and extended by zero off $\operatorname{supp}\chi_s$. By Notation~\ref{4.22:not-tent-gauge-map}, $\mathsf{Y}:=\sum_s f_s=\mathsf{Y}_{\mathsf{r}}[Z]$ and $\rho:=e^{-\mathsf{Y}}=\rho_{\mathsf{r}}[Z]$.

We use the following properties of the tents $\chi_s$, $s\in\Sigma_W$:
\begin{itemize}
\item $\chi_s\ge0$ and $\sum_s\chi_s\le1$;
\item $\chi_s(s')=\delta_{ss'}$ on the unit sites of level $k_W$;
\item $\operatorname{supp}\chi_s\subset\{\operatorname{dist}(\cdot,s)<\ell\}$;
\item $\ell|\partial_\mu\chi_s|\le C_\chi$ and $\ell^2|\partial_\nu\partial_\mu\chi_s|\le C_\chi$ for all $\mu,\nu$.
\end{itemize}
Put $m_1:=\max_s\sup\ell|\partial_\mu\chi_s|$ and $m_2:=\max_s\sup\ell^2|\partial_\nu\partial_\mu\chi_s|$, both at most $C_\chi$. Put also
\begin{equation*}
\mathfrak{n}_\chi:=\max_x\#\bigl\{s:\operatorname{supp}\chi_s\cap(x+[0,2\eta]^d)\ne\emptyset\bigr\}\le 3^d .
\end{equation*}
To prove the bound on $\mathfrak{n}_\chi$: $\operatorname{supp}\chi_s$ lies in the open sup-norm cube of half-side $\ell$ about $s$, and $s$ runs over the level-$k_W$ unit lattice of spacing $\ell$. If this cube meets $x+[0,2\eta]^d$, each coordinate of $s$ lies in an open interval of length $2\ell+2\eta\le3\ell$ (as $\eta/\ell\le\frac12$). Such an interval contains at most three points of spacing $\ell$.

\emph{Transport identity.} From $U(x,\nu)u(x+\eta\nu)^{-1}=u(x)^{-1}e^{i\eta A_\nu(x)}$ we get, for $x,x+\eta\nu\in\square_s$,
\begin{equation}\label{4.22:eq-frame-proof-tent-4}
\operatorname{Ad}_{U(x,\nu)}\operatorname{Ad}_{u(x+\eta\nu)^{-1}}=\operatorname{Ad}_{u(x)^{-1}}E_\nu(x).
\end{equation}

\begin{lemma}[Proof of the frame lemma: exponential estimates and tent differences]
\label{4.22:prf-frame-proof-tent}
In the setting just described the following hold.
\begin{enumerate}
\item \emph{Exponential estimates.} Let $A,A_0,A_1\in\mathfrak{g}$ and $\Gamma,\Gamma_0,\Gamma_1,\Pi\in\mathfrak{A}$. Let $E:=\operatorname{Ad}_{e^{i\eta A}}$ (so that $E_\mu(x)$ is $E$ at $A=A_\mu(x)$), let $t\mapsto\Gamma_t$ be a $C^1$ path in $\mathfrak{A}$, and put $\mathfrak{b}_1:=2.05$. Then
\begin{equation}\label{4.22:eq-frame-proof-tent-a}
\begin{aligned}
&\text{(e1)}\quad \|E\Pi\|=\|\Pi\|;\\
&\text{(e2)}\quad \|\eta^{-1}(E-1)\Pi\|\le\|[A,\Pi]\|\le2\|A\|\,\|\Pi\|;\\
&\text{(e3)}\quad \|\operatorname{Ad}_{e^{i\eta A_1}}\Pi-\operatorname{Ad}_{e^{i\eta A_0}}\Pi\|
\le2\eta\|A_1-A_0\|\,\|\Pi\|;\\
&\text{(e4)}\quad \|\operatorname{Ad}_{e^{-\Gamma}}\Pi\|\le e^{2\|\Gamma\|}\|\Pi\|,\quad
\|\operatorname{Ad}_{e^{-\Gamma}}\Pi-\Pi\|\le(e^{2\|\Gamma\|}-1)\|\Pi\|;\\
&\text{(e5)}\quad \bigl\|\tfrac{d}{dt}e^{\Gamma_t}\bigr\|\le\|\dot\Gamma_t\|e^{\|\Gamma_t\|},\quad
\|e^{\Gamma_1}-e^{\Gamma_0}\|\le\|\Gamma_1-\Gamma_0\|e^{\max(\|\Gamma_0\|,\|\Gamma_1\|)},\\
&\phantom{\text{(e5)}\quad}\|e^{\Gamma}-1\|\le e^{\|\Gamma\|}-1;\\
&\text{(e6)}\quad e^{2y}-1\le\mathfrak{b}_1y\ \text{ and }\ 2e^{2y}\le\mathfrak{b}_1
\ \text{ for }0\le y\le10^{-2};\\
&\phantom{\text{(e6)}\quad}e^{2a}-1-2a\le3a^2\ \text{ for }0\le a\le\tfrac12 .
\end{aligned}
\end{equation}
For the cube gauge, whenever $x,x+\eta\nu\in\square_s$,
\begin{equation*}
\|(\partial_\nu E_\mu)(x)\Pi\|\le2\eta\|(\partial_\nu A_\mu)(x)\|\,\|\Pi\|,\qquad
\|(\partial_\nu D_\mu)(x)\Pi\|\le2\|(\partial_\nu A_\mu)(x)\|\,\|\Pi\|.
\end{equation*}
\item \emph{Leibniz rules.} Let $a$ be a scalar-valued site function and $F$ an $\mathfrak{A}$-valued one. Let $F$ also denote an $\mathfrak{A}$-valued function on $\square_s$, and put $f:=\operatorname{Ad}_{u^{-1}}F$ there. Then (L1) holds at every site $x$, and (L2) holds at every $x$ with $x,x+\eta\nu\in\square_s$:
\begin{equation}\label{4.22:eq-frame-proof-tent-b}
\begin{aligned}
&\text{(L1)}\quad \partial_\nu(aF)(x)=(\partial_\nu a)(x)F(x+\eta\nu)+a(x)(\partial_\nu F)(x);\\
&\text{(L2)}\quad (\nabla^U_\nu f)(x)
=\operatorname{Ad}_{u(x)^{-1}}\bigl[E_\nu(x)(\partial_\nu F)(x)+D_\nu(x)F(x)\bigr]\\
&\phantom{\text{(L2)}\quad (\nabla^U_\nu f)(x)}
=\operatorname{Ad}_{u(x)^{-1}}\bigl[(\partial_\nu F)(x)+D_\nu(x)F(x+\eta\nu)\bigr].
\end{aligned}
\end{equation}
\item $m_1\ge1$; in particular $C_\chi\ge1$.
\item \emph{First covariant differences.} For every $s$, $x$, $\mu$ such that $\{x,x+\eta\mu\}$ meets $\operatorname{supp}\chi_s$,
\begin{equation*}
(\nabla^U_\mu f_s)(x)=\operatorname{Ad}_{u(x)^{-1}}G^s_\mu(x),\qquad
G^s_\mu(x):=(\partial_\mu\chi_s)(x)E_\mu(x)\tilde Z_s+\chi_s(x)D_\mu(x)\tilde Z_s .
\end{equation*}
Otherwise $(\nabla^U_\mu f_s)(x)=0$. In all cases
\begin{equation*}
\ell\|(\nabla^U_\mu f_s)(x)\|\le(m_1+\lambda\chi_s(x))z\le(m_1+\lambda)z<(C_\chi+1)z .
\end{equation*}
\item \emph{Second covariant differences.} Let $\mathfrak{s}(x;\mu,\nu):=\{x,x+\eta\mu,x+\eta\nu,x+\eta\mu+\eta\nu\}$. For every $s$, $x$, $\mu$, $\nu$ such that $\mathfrak{s}(x;\mu,\nu)$ meets $\operatorname{supp}\chi_s$, one has $\mathfrak{s}(x;\mu,\nu)\subset\square_s$. Writing $\chi:=\chi_s$ and $\tilde Z:=\tilde Z_s$, we have $(\nabla^U_\nu\nabla^U_\mu f_s)(x)=\operatorname{Ad}_{u(x)^{-1}}[T_1+\dots+T_6]$, where
\begin{align*}
T_1&=(\partial_\nu\partial_\mu\chi)(x)E_\mu(x+\eta\nu)\tilde Z, &
T_2&=(\partial_\mu\chi)(x)(\partial_\nu E_\mu)(x)\tilde Z,\\
T_3&=(\partial_\nu\chi)(x)D_\mu(x+\eta\nu)\tilde Z, &
T_4&=\chi(x)(\partial_\nu D_\mu)(x)\tilde Z,\\
T_5&=(\partial_\mu\chi)(x+\eta\nu)D_\nu(x)E_\mu(x+\eta\nu)\tilde Z, &
T_6&=\chi(x+\eta\nu)D_\nu(x)D_\mu(x+\eta\nu)\tilde Z .
\end{align*}
These satisfy the bounds
\begin{align*}
\ell^2\|T_1\|&\le m_2z, & \ell^2\|T_2\|&\le m_1\lambda'(\eta/\ell)z, &
\ell^2\|T_3\|&\le m_1\lambda z,\\
\ell^2\|T_4\|&\le\lambda'\chi(x)z, & \ell^2\|T_5\|&\le m_1\lambda z, &
\ell^2\|T_6\|&\le\lambda^2\chi(x+\eta\nu)z .
\end{align*}
If $\mathfrak{s}(x;\mu,\nu)$ misses $\operatorname{supp}\chi_s$, the second difference vanishes. Hence, in all cases, $\ell^2\|(\nabla^U_\nu\nabla^U_\mu f_s)(x)\|\le K_\chi z$, where
\begin{equation*}
K_\chi:=m_2+m_1\bigl(2\lambda+\lambda'\eta/\ell\bigr)+\lambda'+\lambda^2<C_\chi+2 ,
\end{equation*}
and a fortiori $\ell^2\|(\nabla^U_\nu\nabla^U_\mu f_s)(x)\|\le(C_\chi+3)z$.
\end{enumerate}
In particular, with $p_1:=\mathfrak{n}_\chi(m_1+\lambda)$, $p_2:=\mathfrak{n}_\chi K_\chi$, $\bar p_1:=3^d(C_\chi+1)$ and $\bar p_2:=3^d(4C_\chi+2)$,
\begin{equation}\label{4.22:eq-frame-proof-tent-c}
\begin{gathered}
\sup_x\|\mathsf{Y}(x)\|\le z,\qquad
\sup_{x,\mu}\ell\|(\nabla^U_\mu\mathsf{Y})(x)\|\le p_1z,\qquad p_1<\bar p_1,\\
\sup_{x,\mu,\nu}\ell^2\|(\nabla^U_\nu\nabla^U_\mu\mathsf{Y})(x)\|\le p_2z,\qquad
p_2<3^d(C_\chi+2)\le\bar p_2 .
\end{gathered}
\end{equation}
\end{lemma}

\begin{proof}
\emph{Item 1.}

(e1) Since $A\in\mathfrak{g}$ is Hermitian and $\eta$ is real, $e^{i\eta A}$ is unitary. Conjugation by a unitary matrix preserves the operator norm.

(e2) For $t\in[0,1]$,
\begin{equation*}
\frac{d}{dt}\bigl(e^{it\eta A}\Pi e^{-it\eta A}\bigr)=e^{it\eta A}[i\eta A,\Pi]e^{-it\eta A}.
\end{equation*}
Integrating over $t\in[0,1]$ gives
\begin{equation*}
E\Pi-\Pi=\int_0^1e^{it\eta A}[i\eta A,\Pi]e^{-it\eta A}\,dt .
\end{equation*}
The outer factors are unitary, so $\|E\Pi-\Pi\|\le\eta\|[A,\Pi]\|$. Divide by $\eta$ and use $\|[A,\Pi]\|\le2\|A\|\,\|\Pi\|$.

(e3) Put $A_t:=A_0+t(A_1-A_0)\in\mathfrak{g}$ and $V_t:=e^{i\eta A_t}$. Duhamel's formula gives
\begin{equation*}
\dot V_t=\int_0^1e^{i\sigma\eta A_t}\,i\eta(A_1-A_0)\,e^{i(1-\sigma)\eta A_t}\,d\sigma .
\end{equation*}
The outer factors are unitary, so $\|\dot V_t\|\le\eta\|A_1-A_0\|$. Since $V_t^*=V_t^{-1}$, we have $V_t\Pi V_t^*=\operatorname{Ad}_{e^{i\eta A_t}}\Pi$ and
\begin{equation*}
\frac{d}{dt}(V_t\Pi V_t^*)=\dot V_t\Pi V_t^*+V_t\Pi\dot V_t^* .
\end{equation*}
This has norm at most $2\eta\|A_1-A_0\|\,\|\Pi\|$. Integrating over $t\in[0,1]$ gives (e3).

For the cube gauge,
\begin{equation*}
(\partial_\nu E_\mu)(x)=\eta^{-1}\bigl[\operatorname{Ad}_{e^{i\eta A_\mu(x+\eta\nu)}}-\operatorname{Ad}_{e^{i\eta A_\mu(x)}}\bigr],
\end{equation*}
and $\|A_\mu(x+\eta\nu)-A_\mu(x)\|=\eta\|(\partial_\nu A_\mu)(x)\|$. So (e3) gives $\|(\partial_\nu E_\mu)(x)\Pi\|\le2\eta\|(\partial_\nu A_\mu)(x)\|\,\|\Pi\|$. Finally $\partial_\nu D_\mu=\eta^{-1}\partial_\nu E_\mu$.

(e4) From the exponential series, $\|e^{\mp\Gamma}\|\le e^{\|\Gamma\|}$, which gives the first bound. For the second,
\begin{equation*}
\operatorname{Ad}_{e^{-\Gamma}}\Pi-\Pi=\int_0^1\frac{d}{dt}\bigl(e^{-t\Gamma}\Pi e^{t\Gamma}\bigr)dt
=\int_0^1e^{-t\Gamma}[\Pi,\Gamma]e^{t\Gamma}\,dt .
\end{equation*}
Its norm is at most
\begin{equation*}
2\|\Gamma\|\,\|\Pi\|\int_0^1e^{2t\|\Gamma\|}dt=(e^{2\|\Gamma\|}-1)\|\Pi\| .
\end{equation*}
In what follows (e4) is applied with $\Gamma$ a value of the tent field $\mathsf{Y}$.

(e5) By the product rule, $\frac{d}{dt}\Gamma_t^n$ is a sum of $n$ terms, so it has norm at most $n\|\Gamma_t\|^{n-1}\|\dot\Gamma_t\|$. The series of these derivatives is therefore dominated by $\|\dot\Gamma_t\|\sum_{n\ge1}\|\Gamma_t\|^{n-1}/(n-1)!$. It converges uniformly on compact $t$-intervals, so the exponential series may be differentiated termwise, and this gives the first bound.

For the second bound, apply the first along $\Gamma_t:=\Gamma_0+t(\Gamma_1-\Gamma_0)$, $t\in[0,1]$. There $\dot\Gamma_t=\Gamma_1-\Gamma_0$, and by convexity of the norm $\|\Gamma_t\|\le\max(\|\Gamma_0\|,\|\Gamma_1\|)$; integrate over $t$.

For the third, apply the first along $\Gamma_t:=t\Gamma$, keeping the exact exponent:
\begin{equation*}
\|e^{\Gamma}-1\|\le\int_0^1\|\Gamma\|e^{t\|\Gamma\|}dt=e^{\|\Gamma\|}-1 .
\end{equation*}

(e6) For $0\le y\le10^{-2}$,
\begin{equation*}
e^{2y}-1=\int_0^y2e^{2t}dt\le2ye^{2y}\le2\cdot1.0203\cdot y,
\end{equation*}
because $e^{0.02}<1.0203$. Also $2e^{2y}\le2\cdot1.0203<2.05=\mathfrak{b}_1$. For $0\le a\le\frac12$ we have $2a\le1$, hence
\begin{equation*}
e^{2a}-1-2a=4a^2\sum_{n\ge2}\frac{(2a)^{n-2}}{n!}\le4a^2(e-2)<3a^2 .
\end{equation*}

\emph{Item 2.} (L1) is the identity
\begin{equation*}
a(x+\eta\nu)F(x+\eta\nu)-a(x)F(x)
=\bigl(a(x+\eta\nu)-a(x)\bigr)F(x+\eta\nu)+a(x)\bigl(F(x+\eta\nu)-F(x)\bigr),
\end{equation*}
divided by $\eta$.

For (L2), put $x':=x+\eta\nu$. By the definition of $\nabla^U$ and the transport identity \eqref{4.22:eq-frame-proof-tent-4},
\begin{equation*}
(\nabla^U_\nu f)(x)=\eta^{-1}\bigl[\operatorname{Ad}_{U(x,\nu)}\operatorname{Ad}_{u(x')^{-1}}F(x')
-\operatorname{Ad}_{u(x)^{-1}}F(x)\bigr]
=\operatorname{Ad}_{u(x)^{-1}}\eta^{-1}\bigl[E_\nu(x)F(x')-F(x)\bigr].
\end{equation*}
Now
\begin{equation*}
E_\nu F(x')-F(x)=E_\nu\bigl(F(x')-F(x)\bigr)+(E_\nu-1)F(x)=\bigl(F(x')-F(x)\bigr)+(E_\nu-1)F(x').
\end{equation*}
Dividing by $\eta$, the first decomposition gives the first form of (L2) and the second gives the second form.

\emph{Item 3.} We have $\chi_s(s)=1$, and $\chi_s(s+\ell e_1)=0$ because $\operatorname{dist}(s+\ell e_1,s)=\ell$ and $\operatorname{supp}\chi_s\subset\{\operatorname{dist}(\cdot,s)<\ell\}$. Telescoping along the $\ell/\eta=L^{k_W}$ bonds from $s$ to $s+\ell e_1$ gives
\begin{equation*}
1=\chi_s(s)-\chi_s(s+\ell e_1)\le\sum_{j=0}^{L^{k_W}-1}\eta\,\bigl|(\partial_1\chi_s)(s+j\eta e_1)\bigr|
\le\ell\max|\partial_1\chi_s|\le m_1 .
\end{equation*}
Hence $C_\chi\ge m_1\ge1$.

\emph{Item 4.} If neither $x$ nor $x+\eta\mu$ lies in $\operatorname{supp}\chi_s$, then $f_s$ vanishes at both points and $(\nabla^U_\mu f_s)(x)=0$. Otherwise both points lie in the $2\eta$-neighbourhood of $\operatorname{supp}\chi_s$, hence in $\square_s$. Apply the first form of (L2) with $F:=\chi_s\tilde Z_s$. Since $\tilde Z_s$ is constant, $(\partial_\mu F)(x)=(\partial_\mu\chi_s)(x)\tilde Z_s$, and (L2) yields the displayed formula with $G^s_\mu$.

For the bound we collect four facts:
\begin{itemize}
\item $\operatorname{Ad}_{u(x)^{-1}}$ is an isometry, because $u(x)\in G\subset U(N)$;
\item $\ell|(\partial_\mu\chi_s)(x)|\le m_1$;
\item $\|E_\mu(x)\tilde Z_s\|=\|\tilde Z_s\|$ by (e1), and $\|\tilde Z_s\|=\|Z_s\|$ since $u(s)\in G$, so this is at most $z$;
\item $\ell\|D_\mu(x)\tilde Z_s\|\le2\ell\|A_\mu(x)\|z\le\lambda z$, by (e2).
\end{itemize}
Hence $\ell\|(\nabla^U_\mu f_s)(x)\|\le(m_1+\lambda\chi_s(x))z$. Then $\chi_s\le1$ gives the bound $(m_1+\lambda)z$, and $m_1\le C_\chi$ together with $\lambda<\mathsf{t}<1$ (by \eqref{4.22:eq-frame-proof-tent-3}) gives the bound $(C_\chi+1)z$.

\emph{Item 5.} If $\mathfrak{s}(x;\mu,\nu)$ meets $\operatorname{supp}\chi_s$, its points lie within distance $2\eta$ of that support, hence in $\square_s$. If it misses $\operatorname{supp}\chi_s$, then $f_s$ vanishes on $\mathfrak{s}(x;\mu,\nu)$. The second difference at $x$ involves only the values of $f_s$ at the four points of $\mathfrak{s}(x;\mu,\nu)$, so it vanishes.

In the first case, item~4 applies at $x$ and at $x+\eta\nu$: the pairs $\{x,x+\eta\mu\}$ and $\{x+\eta\nu,x+\eta\nu+\eta\mu\}$ both lie in $\square_s$. Where one of these pairs misses $\operatorname{supp}\chi_s$, both sides of the formula of item~4 vanish. So $\nabla^U_\mu f_s=\operatorname{Ad}_{u^{-1}}G^s_\mu$ at $x$ and at $x+\eta\nu$. The second form of (L2), applied with $F:=G^s_\mu$, gives
\begin{equation*}
(\nabla^U_\nu\nabla^U_\mu f_s)(x)
=\operatorname{Ad}_{u(x)^{-1}}\bigl[(\partial_\nu G^s_\mu)(x)+D_\nu(x)G^s_\mu(x+\eta\nu)\bigr].
\end{equation*}
We split the right-hand side into the six terms.
\begin{itemize}
\item (L1) with $a:=\partial_\mu\chi$ and $F:=E_\mu\tilde Z$ gives $\partial_\nu[(\partial_\mu\chi)E_\mu\tilde Z](x)=T_1+T_2$.
\item (L1) with $a:=\chi$ and $F:=D_\mu\tilde Z$ gives $\partial_\nu[\chi D_\mu\tilde Z](x)=T_3+T_4$.
\item Directly from the definition of $G^s_\mu$, $D_\nu(x)G^s_\mu(x+\eta\nu)=T_5+T_6$.
\end{itemize}

The bounds on the six terms use the following facts:
\begin{itemize}
\item $\ell^2|\partial_\nu\partial_\mu\chi|\le m_2$, $\ell|\partial\chi|\le m_1$ and $0\le\chi\le1$;
\item $\|E_\mu\tilde Z\|=\|\tilde Z\|\le z$, by (e1);
\item $\ell\|D_\mu\Pi\|\le2\ell\|A_\mu\|\,\|\Pi\|\le\lambda\|\Pi\|$, by (e2);
\item from item~1 (the consequence of (e3)),
\begin{equation*}
\ell^2\cdot\ell^{-1}m_1\|(\partial_\nu E_\mu)\tilde Z\|\le m_1\ell\cdot2\eta\|\partial_\nu A_\mu\|z
=m_1(\eta/\ell)\bigl(2\ell^2\|\partial_\nu A_\mu\|\bigr)z\le m_1\lambda'(\eta/\ell)z ;
\end{equation*}
\item also from item~1, $\ell^2\|(\partial_\nu D_\mu)\tilde Z\|\le2\ell^2\|\partial_\nu A_\mu\|z\le\lambda'z$;
\item $\ell^2\|D_\nu D_\mu\tilde Z\|\le\lambda\cdot\lambda z$, by (e2) applied twice and (e1).
\end{itemize}
Applied term by term, these give the six displayed bounds. For $T_5$, the norm $\ell\|D_\nu(x)E_\mu(x+\eta\nu)\tilde Z\|$ is at most $\lambda\|E_\mu(x+\eta\nu)\tilde Z\|$ by (e2), hence at most $\lambda z$ by (e1). Summing, and using $\chi\le1$, $\operatorname{Ad}_{u(x)^{-1}}$ isometric, gives the bound $K_\chi z$.

It remains to show $K_\chi<C_\chi+2$. We use $m_1,m_2\le C_\chi$, $\eta/\ell\le1$ and $\lambda,\lambda'<\mathsf{t}$. The estimates $m_1\cdot2\lambda<2C_\chi\mathsf{t}$, $m_1\lambda'\eta/\ell+\lambda'<(C_\chi+1)\mathsf{t}$ and $\lambda^2<\mathsf{t}^2$ give
\begin{equation*}
K_\chi-C_\chi<2C_\chi\mathsf{t}+(C_\chi+1)\mathsf{t}+\mathsf{t}^2 .
\end{equation*}
We bound the three terms in turn; the second bound uses $(C_\chi+1)/(2C_\chi+1)\le\tfrac23$, which is equivalent to $C_\chi\ge1$ (item~3):
\begin{align*}
2C_\chi\mathsf{t}&=\frac{20}{21}\cdot\frac{2C_\chi}{2C_\chi+1}<\frac{60}{63},\\
(C_\chi+1)\mathsf{t}&=\frac{20}{21}\cdot\frac{C_\chi+1}{2C_\chi+1}\le\frac{20}{21}\cdot\frac23=\frac{40}{63},\\
\mathsf{t}^2&\le\Bigl(\frac{20}{63}\Bigr)^2<\frac{7}{63}.
\end{align*}
The sum is less than $107/63<2$, so $K_\chi<C_\chi+2\le C_\chi+3$.

\emph{The display \eqref{4.22:eq-frame-proof-tent-c}.} Since $\chi_s\ge0$, $\sum_s\chi_s\le1$ and the transports $\operatorname{Ad}_{u(x)^{-1}u(s)}$ are by $G$-valued elements,
\begin{equation*}
\|\mathsf{Y}(x)\|\le\sum_s\chi_s(x)\|Z_s\|\le z .
\end{equation*}
The operators $\nabla^U_\mu$ and $\nabla^U_\nu\nabla^U_\mu$ are linear, so $\nabla^U_\mu\mathsf{Y}=\sum_s\nabla^U_\mu f_s$, and likewise for the second difference. At a given $x$, only tents whose support meets $\{x,x+\eta\mu\}$, respectively $\mathfrak{s}(x;\mu,\nu)$, contribute. Both sets lie in $x+[0,2\eta]^d$, so there are at most $\mathfrak{n}_\chi$ such tents.

Items~4 and~5 then give the two difference bounds of \eqref{4.22:eq-frame-proof-tent-c}, with $p_1=\mathfrak{n}_\chi(m_1+\lambda)$ and $p_2=\mathfrak{n}_\chi K_\chi$. Finally, $\mathfrak{n}_\chi\le3^d$, $m_1+\lambda<C_\chi+1$ and $K_\chi<C_\chi+2$ give $p_1<\bar p_1$ and
\begin{equation*}
p_2<3^d(C_\chi+2)\le3^d(4C_\chi+2)=\bar p_2 .
\end{equation*}
\end{proof}

The six terms of item~5 have a continuum counterpart. For $u=1$ and $\eta\to0$, with $D^A_\mu F:=\partial_\mu F+i[A_\mu,F]$ for an $\mathfrak{A}$-valued function $F$, and $Z\in\mathfrak{A}$ constant, one has
\begin{equation*}
\begin{aligned}
D^A_\nu D^A_\mu(\chi Z)={}&(\partial_\nu\partial_\mu\chi)Z+i(\partial_\nu\chi)[A_\mu,Z]
+i(\partial_\mu\chi)[A_\nu,Z]\\
&+i\chi[\partial_\nu A_\mu,Z]-\chi[A_\nu,[A_\mu,Z]] .
\end{aligned}
\end{equation*}
The five terms on the right correspond to $T_1$, $T_3$, $T_5$, $T_4$, $T_6$ respectively. The term $T_2$, of relative size $O(\eta/\ell)$, is an effect of the lattice. The two cross terms $T_3$ and $T_5$ are distinct and do not combine, which is why each contributes its own term $m_1\lambda$ to $K_\chi$.

This completes the exponential estimates and the bounds on the first and second covariant differences of the tent field $\mathsf{Y}_{\mathsf{r}}[Z]$. The estimates of the increments of the entries in clause~(i) of the frame lemma (Lemma~\ref{4.22:lem-frame-lemma}) use these bounds.

\begin{lemma}[Proof of the frame lemma: the logarithm of a triple product]
\label{4.22:prf-frame-proof-logarithm}
Let $\mathfrak{A}:=\mathfrak{gl}(N,\mathbb{C})$ carry the operator norm $\|\cdot\|$ on
$\mathbb{C}^N$. Put $\rho_1:=1/10$,
\begin{equation*}
\mathcal{P}:=\bigl\{(Y,X,W)\in\mathfrak{A}^3:\ \|Y\|,\|X\|,\|W\|<\rho_1\bigr\},\qquad
\mathcal{P}':=\bigl\{(Y,X,W)\in\mathfrak{A}^3:\ \|Y\|+\|X\|<\rho_1,\ \|W\|<\rho_1\bigr\},
\end{equation*}
so that $\mathcal{P}'\subset\mathcal{P}$, and let
$\log(1+M):=\sum_{n\ge1}(-1)^{n+1}M^n/n$ for $\|M\|<1$. On $\mathcal{P}$ define
\begin{equation}\label{4.22:eq-frame-proof-logarithm-1}
F(Y,X,W):=\log\bigl(e^{-Y}e^{X}e^{Y+W}\bigr),\qquad r:=F-X-W,\qquad
r_0(Y,X):=r(Y,X,0),\qquad R_W:=r-r_0,
\end{equation}
and put $\mathfrak{b}_2:=2\rho_1^{-2}=200$, $\mathfrak{b}_3:=4\rho_1^{-2}=400$. Then the
following hold.

The clauses \textup{(R0)}--\textup{(R4)} are
\begin{equation}\label{4.22:eq-frame-proof-logarithm-a}
\begin{aligned}
&(\mathrm{R}0)\quad \bigl\|e^{-Y}e^{X}e^{Y+W}-1\bigr\|<\tfrac12,\qquad
e^{F}=e^{-Y}e^{X}e^{Y+W},\\
&\qquad\ \ \|F\|<\log 2,\qquad \|r\|<1,\qquad \|R_W\|<2,\\
&\qquad\ \ r(\operatorname{Ad}_gY,\operatorname{Ad}_gX,\operatorname{Ad}_gW)
=\operatorname{Ad}_g\,r(Y,X,W),\\
&\qquad\ \ r(Y,0,0)=r(0,X,0)=r(0,0,W)=0;\\
&(\mathrm{R}1)\quad r_0(Y,X)=\operatorname{Ad}_{e^{-Y}}X-X,\qquad
\|r_0(Y,X)\|\le\bigl(e^{2\|Y\|}-1\bigr)\|X\|,\\
&\qquad\ \ \|r_0(Y_1,X_1)-r_0(Y_0,X_0)\|\le\bigl(e^{2\bar y}-1\bigr)\|X_1-X_0\|
+2e^{2\bar y}\|X_0\|\,\|Y_1-Y_0\|;\\
&(\mathrm{R}2)\quad \|R_W\|\le 2\|W\|/\rho_1\ \text{ on }\mathcal{P},\qquad
\|R_W\|\le\mathfrak{b}_2\bigl(\|Y\|+\|X\|\bigr)\|W\|\ \text{ on }\mathcal{P}';\\
&(\mathrm{R}3)\quad \|R_W(P_1)-R_W(P_0)\|\\
&\qquad\ \ \le\mathfrak{b}_3\Bigl[\|W_1\|\bigl(\|Y_1-Y_0\|+\|X_1-X_0\|\bigr)
+\bigl(\|Y_0\|+\|X_0\|\bigr)\|W_1-W_0\|\Bigr];\\
&(\mathrm{R}4)\quad \|r\|\le\bigl(e^{2\|Y\|}-1\bigr)\|X\|
+\mathfrak{b}_2\bigl(\|Y\|+\|X\|\bigr)\|W\|\ \text{ on }\mathcal{P}'.
\end{aligned}
\end{equation}
Here \textup{(R0)} holds on $\mathcal{P}$, where $F$ and $r$ are holomorphic; its
equivariance holds for every $g\in U(N)$, and indeed for every invertible $g$ such that both
$(Y,X,W)$ and $(\operatorname{Ad}_gY,\operatorname{Ad}_gX,\operatorname{Ad}_gW)$ lie in
$\mathcal{P}$; and $\operatorname{tr}F=0$ when $Y,X,W$ are traceless.
\textup{(R1)} holds on $\mathcal{P}$ (for the difference bound: at two points
$(Y_i,X_i,0)\in\mathcal{P}$, $i=0,1$, with $\bar y:=\max_i\|Y_i\|$); in \textup{(R3)},
$P_i=(Y_i,X_i,W_i)$, $i=0,1$, are triples all six of whose coordinates have norm at most
$\tfrac12\rho_1$, and $\|Y_0\|<\tfrac12\rho_1$. In words, \textup{(R4)} says that every term
of $r$ carries a factor $W$, or a factor $Y$ together with a factor $X$; at $W=0$,
$r=\operatorname{Ad}_{e^{-Y}}X-X$, whose leading term $-[Y,X]$ is linear in $\|Y\|$ at fixed
$X$, and $r(\cdot,\cdot,0)$ vanishes identically for commuting data: this is a non-abelian
fact. All bounds are independent of $N$.
\end{lemma}

\begin{proof}
Throughout, $\varphi$ ranges over the linear functionals on $\mathfrak{A}$ of norm one, so that
$\|M\|=\sup_\varphi|\varphi(M)|$ by the Hahn--Banach theorem; $\mathfrak{A}$ is a
finite-dimensional complex vector space, and holomorphy is meant in the matrix entries. We use
the clauses (e4), (e5), (e6) of \eqref{4.22:eq-frame-proof-tent-a}, together with the
elementary bound $\|e^{B}\|\le e^{\|B\|}$, which follows from the exponential series.

\emph{(R0).} Write $e^{-Y}=1+\Delta_1$, $e^{X}=1+\Delta_2$, $e^{Y+W}=1+\Delta_3$. By the third
inequality of (e5), $\|\Delta_1\|\le e^{\|Y\|}-1$, $\|\Delta_2\|\le e^{\|X\|}-1$ and
$\|\Delta_3\|\le e^{\|Y+W\|}-1\le e^{\|Y\|+\|W\|}-1$. Expanding
$(1+\Delta_1)(1+\Delta_2)(1+\Delta_3)-1$ as the sum of the non-empty ordered sub-products of
the $\Delta_i$ and estimating each sub-product by the product of the norms,
\begin{equation*}
\bigl\|(1+\Delta_1)(1+\Delta_2)(1+\Delta_3)-1\bigr\|
\le\prod_{i=1}^3\bigl(1+\|\Delta_i\|\bigr)-1\le e^{2\|Y\|+\|X\|+\|W\|}-1 .
\end{equation*}
On $\mathcal{P}$ the exponent is $<4\rho_1=0.4$, and $e^{0.4}<1.4919$, so the right-hand side
is $<\tfrac12$. The logarithmic series converges locally uniformly on $\{\|M\|<1\}$, hence $F$,
the composition of this series with the entire map $(Y,X,W)\mapsto e^{-Y}e^Xe^{Y+W}-1$, is
holomorphic on $\mathcal{P}$, and so is $r=F-X-W$. With $M:=e^{-Y}e^Xe^{Y+W}-1$,
\begin{equation*}
\|F\|\le\sum_{n\ge1}\frac{\|M\|^n}{n}=-\log\bigl(1-\|M\|\bigr)<-\log\bigl(1-\tfrac12\bigr)=\log2 .
\end{equation*}
Hence
\begin{equation*}
\|r\|\le\|F\|+\|X\|+\|W\|<\log2+2\rho_1<0.9<1 .
\end{equation*}
Since $(Y,X,0)\in\mathcal{P}$ whenever
$(Y,X,W)\in\mathcal{P}$, also $\|r_0(Y,X)\|<1$, and $\|R_W\|\le\|r\|+\|r_0\|<2$.

Two identities of series are used. First, $\log(e^B)=B$ for $\|B\|<\log2$: indeed
$\|e^B-1\|\le e^{\|B\|}-1<1$ by (e5), and
$\sum_n\frac1n\bigl(e^{\|B\|}-1\bigr)^n=-\log\bigl(2-e^{\|B\|}\bigr)<\infty$, so the composed
double series converges absolutely. Second, $e^{\log(1+M)}=1+M$ for $\|M\|<1$: here
\begin{equation*}
\sum_{k\ge0}\frac1{k!}\Bigl(\sum_{n\ge1}\frac{\|M\|^n}{n}\Bigr)^{k}
=\exp\bigl(-\log(1-\|M\|)\bigr)=\bigl(1-\|M\|\bigr)^{-1}<\infty .
\end{equation*}
In both cases the composed series converge absolutely, so they may be rearranged, and the
formal identities $\log\circ\exp=\mathrm{id}$, $\exp\circ\log=\mathrm{id}$ of power series in
one commuting variable yield the matrix identities. In particular
$e^{F}=e^{-Y}e^Xe^{Y+W}$ on $\mathcal{P}$.

Equivariance: for invertible $g$, the map $M\mapsto gMg^{-1}$ commutes with the exponential
series, with products and with the logarithmic series; hence
$F(\operatorname{Ad}_gY,\operatorname{Ad}_gX,\operatorname{Ad}_gW)=\operatorname{Ad}_gF(Y,X,W)$
whenever both triples lie in $\mathcal{P}$, and since $X+W$ enters $r$ linearly, the same holds
for $r$. For $g\in U(N)$ the map $\operatorname{Ad}_g$ preserves the operator norm, so it maps
$\mathcal{P}$ onto itself.

Vanishing: $r(Y,0,0)=\log(e^{-Y}e^{Y})=\log1=0$; $r(0,X,0)=\log(e^X)-X=0$ and
$r(0,0,W)=\log(e^W)-W=0$, by the first identity above, since $\|X\|,\|W\|<\rho_1<\log2$.

Trace: let $Y,X,W$ be traceless. The set $\mathcal{P}$ is star-shaped about $0$, so
$\tau\mapsto\operatorname{tr}F(\tau Y,\tau X,\tau W)$ is continuous on $[0,1]$. By
$e^{F}=e^{-Y}e^Xe^{Y+W}$ at the point $(\tau Y,\tau X,\tau W)$ and $\det e^{B}=e^{\operatorname{tr}B}$,
\begin{equation*}
e^{\operatorname{tr}F(\tau Y,\tau X,\tau W)}
=\det\bigl(e^{-\tau Y}e^{\tau X}e^{\tau(Y+W)}\bigr)
=e^{\tau(-\operatorname{tr}Y+\operatorname{tr}X+\operatorname{tr}Y+\operatorname{tr}W)}=1 ,
\end{equation*}
so the function takes values in $2\pi i\mathbb{Z}$; it vanishes at $\tau=0$ (where
$F=\log1=0$), hence identically, and $\operatorname{tr}F(Y,X,W)=0$.

\emph{(R1).} Conjugation commutes with the exponential series, so
$e^{-Y}e^{X}e^{Y}=\exp\bigl(\operatorname{Ad}_{e^{-Y}}X\bigr)$. By the first inequality of (e4),
on $\mathcal{P}$,
\begin{equation*}
\bigl\|\operatorname{Ad}_{e^{-Y}}X\bigr\|\le e^{2\|Y\|}\|X\|<e^{0.2}/10<0.123<\log2 ,
\end{equation*}
so by $\log\circ\exp=\mathrm{id}$ on $\{\|B\|<\log2\}$ we get
$F(Y,X,0)=\operatorname{Ad}_{e^{-Y}}X$, that is $r_0(Y,X)=\operatorname{Ad}_{e^{-Y}}X-X$. The
bound $\|r_0(Y,X)\|\le(e^{2\|Y\|}-1)\|X\|$ is the second inequality of (e4). For the
differences, a direct expansion gives
\begin{equation*}
r_0(Y_1,X_1)-r_0(Y_0,X_0)
=\bigl(\operatorname{Ad}_{e^{-Y_1}}-1\bigr)(X_1-X_0)
+\bigl(\operatorname{Ad}_{e^{-Y_1}}-\operatorname{Ad}_{e^{-Y_0}}\bigr)X_0 ,
\end{equation*}
and
\begin{equation*}
\bigl(\operatorname{Ad}_{e^{-Y_1}}-\operatorname{Ad}_{e^{-Y_0}}\bigr)X_0
=\bigl(e^{-Y_1}-e^{-Y_0}\bigr)X_0e^{Y_1}+e^{-Y_0}X_0\bigl(e^{Y_1}-e^{Y_0}\bigr).
\end{equation*}
By the second inequality of (e4) the first term has norm at most
$(e^{2\|Y_1\|}-1)\|X_1-X_0\|\le(e^{2\bar y}-1)\|X_1-X_0\|$. By the second inequality of (e5),
$\|e^{\mp Y_1}-e^{\mp Y_0}\|\le\|Y_1-Y_0\|e^{\bar y}$, and $\|e^{Y_1}\|,\|e^{-Y_0}\|\le e^{\bar y}$;
so each of the last two terms has norm at most $e^{2\bar y}\|X_0\|\,\|Y_1-Y_0\|$. Adding gives
the difference bound of (R1).

\emph{(R2).} Let $(Y,X,W)\in\mathcal{P}$ and put $a:=\|Y\|+\|X\|$, $\omega:=\|W\|$. If
$\omega=0$, then $R_W(Y,X,0)=r(Y,X,0)-r_0(Y,X)=0$ by definition; if $a=0$, then
$R_W(0,0,W)=r(0,0,W)-r(0,0,0)=0$ by (R0). In these cases there is nothing to prove. Let
$\omega>0$. The function $\tau\mapsto\varphi\bigl(R_W(Y,X,\tau W/\omega)\bigr)$ is holomorphic
on $|\tau|<\rho_1$, since there $\|\tau W/\omega\|=|\tau|<\rho_1$ and the point lies in
$\mathcal{P}$; it vanishes at $\tau=0$ and has modulus $<2$ by (R0). Schwarz's lemma on the
disc of radius $\rho_1$ gives modulus at most $2|\tau|/\rho_1$. Taking $\tau:=\omega<\rho_1$
and the supremum over $\varphi$ yields $\|R_W\|\le2\|W\|/\rho_1$.

Now let $(Y,X,W)\in\mathcal{P}'$ with $a,\omega>0$; thus $a<\rho_1$. Put
\begin{equation*}
\psi(\sigma,\tau):=\varphi\bigl(R_W(\sigma Y/a,\ \sigma X/a,\ \tau W/\omega)\bigr),
\qquad |\sigma|,|\tau|<\rho_1 .
\end{equation*}
Since $\|\sigma Y/a\|+\|\sigma X/a\|=|\sigma|$, each of the first two arguments has norm
$<\rho_1$, and the point lies in $\mathcal{P}$; so $\psi$ is holomorphic on the bidisc and
$|\psi|<2$ by (R0). It vanishes on $\{\tau=0\}$ by definition of $R_W$, and on
$\{\sigma=0\}$ because $R_W(0,0,\cdot)=0$. Hence $\psi=\tau\psi_1$ with
$\psi_1(\sigma,\tau):=\int_0^1\partial_\tau\psi(\sigma,s\tau)\,ds$ holomorphic; for
$\tau\ne0$, $\psi_1(0,\tau)=\psi(0,\tau)/\tau=0$, and by continuity also at $\tau=0$; so in the
same way $\psi_1=\sigma\tilde\psi$ with $\tilde\psi$ holomorphic on the bidisc, and
$\psi=\sigma\tau\tilde\psi$. Fix $\varsigma<\rho_1$ and a point with $|\sigma|,|\tau|\le\varsigma$.
The maximum principle in $\sigma$ at fixed $\tau$, followed by the maximum principle in $\tau$
at fixed $\sigma'$, bounds $|\tilde\psi(\sigma,\tau)|$ by the maximum of $|\tilde\psi|$ over
$|\sigma'|=|\tau'|=\varsigma$, where $|\tilde\psi|=|\psi|/\varsigma^2<2/\varsigma^2$. Letting
$\varsigma\uparrow\rho_1$ gives $|\tilde\psi|\le2/\rho_1^2$ on the bidisc, hence
$|\psi(\sigma,\tau)|\le2|\sigma||\tau|/\rho_1^2$. At $(\sigma,\tau):=(a,\omega)$ the argument of
$R_W$ is $(Y,X,W)$; taking the supremum over $\varphi$,
$\|R_W(Y,X,W)\|\le2a\omega/\rho_1^2=\mathfrak{b}_2(\|Y\|+\|X\|)\|W\|$.

\emph{(R3).} We first record a Cauchy--Lipschitz estimate. Let $f$ be holomorphic, with values
in $\mathfrak{A}$ and $\|f\|\le S_f$, on the open ball of radius $\varsigma_0$ about $0$ in a
finite-dimensional complex normed space, and let $p_0,p_1$ have norm at most
$\varsigma_1<\varsigma_0$. Then
\begin{equation}\label{4.22:eq-frame-proof-logarithm-2}
\|f(p_1)-f(p_0)\|\le\frac{S_f\,\|p_1-p_0\|}{\varsigma_0-\varsigma_1}.
\end{equation}
Indeed, if $p_1\ne p_0$, put $p_t:=p_0+t(p_1-p_0)$, so $\|p_t\|\le\varsigma_1$ by convexity,
and $\hat v:=(p_1-p_0)/\|p_1-p_0\|$. For $|\zeta|<\varsigma_0-\varsigma_1$ the point
$p_t+\zeta\hat v$ has norm $<\varsigma_0$, so $\zeta\mapsto f(p_t+\zeta\hat v)$ is holomorphic
on that disc and
bounded by $S_f$. Cauchy's estimate, applied to $\varphi$ of this function and followed by the
supremum over $\varphi$, bounds its derivative at $\zeta=0$ by $S_f/(\varsigma_0-\varsigma_1)$.
Since $\frac{d}{dt}f(p_t)$ is $\|p_1-p_0\|$ times this derivative, integrating over
$t\in[0,1]$ gives \eqref{4.22:eq-frame-proof-logarithm-2}.

Let $P_0,P_1$ be as in (R3). Since $\|Y_0\|<\tfrac12\rho_1$ and $\|X_0\|\le\tfrac12\rho_1$, we
have $\|Y_0\|+\|X_0\|<\rho_1$.

\emph{Step 1: $W$ moves.} The map $W\mapsto R_W(Y_0,X_0,W)$ is holomorphic on the ball
$\|W\|<\rho_1$; there $(Y_0,X_0,W)\in\mathcal{P}'$, so by the second bound of (R2) it is bounded
by $\mathfrak{b}_2(\|Y_0\|+\|X_0\|)\rho_1$. Applying \eqref{4.22:eq-frame-proof-logarithm-2}
with $\varsigma_0=\rho_1$ and $\varsigma_1=\tfrac12\rho_1$,
\begin{equation*}
\|R_W(Y_0,X_0,W_1)-R_W(Y_0,X_0,W_0)\|
\le2\mathfrak{b}_2\bigl(\|Y_0\|+\|X_0\|\bigr)\|W_1-W_0\|
=\mathfrak{b}_3\bigl(\|Y_0\|+\|X_0\|\bigr)\|W_1-W_0\| .
\end{equation*}

\emph{Step 2: $(Y,X)$ moves at $W_1$.} On $\mathfrak{A}^2$ use the norm
$\max(\|Y\|,\|X\|)$. The map $(Y,X)\mapsto R_W(Y,X,W_1)$ is holomorphic on the ball
$\max(\|Y\|,\|X\|)<\rho_1$; there $(Y,X,W_1)\in\mathcal{P}$ because
$\|W_1\|\le\tfrac12\rho_1$, so by the first bound of (R2) it is bounded by $2\|W_1\|/\rho_1$.
The points $(Y_i,X_i)$ have norm at most $\tfrac12\rho_1$, and
\eqref{4.22:eq-frame-proof-logarithm-2} with $\varsigma_0=\rho_1$, $\varsigma_1=\tfrac12\rho_1$
gives
\begin{equation*}
\|R_W(Y_1,X_1,W_1)-R_W(Y_0,X_0,W_1)\|
\le\frac{4}{\rho_1^2}\|W_1\|\max\bigl(\|Y_1-Y_0\|,\|X_1-X_0\|\bigr)
\le\mathfrak{b}_3\|W_1\|\bigl(\|Y_1-Y_0\|+\|X_1-X_0\|\bigr).
\end{equation*}
Adding the two steps, by way of
\begin{equation*}
R_W(P_1)-R_W(P_0)=\bigl[R_W(Y_1,X_1,W_1)-R_W(Y_0,X_0,W_1)\bigr]
+\bigl[R_W(Y_0,X_0,W_1)-R_W(Y_0,X_0,W_0)\bigr],
\end{equation*}
gives (R3).

\emph{(R4).} On $\mathcal{P}'\subset\mathcal{P}$ write $r=r_0+R_W$ and bound $r_0$ by (R1) and
$R_W$ by the second bound of (R2). As to the verbal reading: $r_0$ vanishes at $Y=0$ and at
$X=0$ by (R0), so every bihomogeneous term of its Taylor expansion in $(Y,X)$ contains both
$Y$ and $X$, and $R_W$ vanishes at $W=0$, so every term of it contains $W$. At $W=0$,
$r=r_0=\operatorname{Ad}_{e^{-Y}}X-X$ by (R1), which vanishes when $[Y,X]=0$.
\end{proof}

\begin{remark}
In the decomposition of the gauge-transformed bond variable carried out below, the three
exponents of the triple
product are $-Y=-\mathsf{Y}(x)$, $X=i\eta A'_\mu(x)$ and
$Y+W=\mathsf{Y}(x)+\eta(\nabla^U_\mu\mathsf{Y})(x)$, at a bond $b=(x,\mu)$. The remainder $r$
is not quadratic in the two exponents $-Y$ and $Y+W$ alone: there is no constant $c$ such that
$\|r(Y,X,W)\|\le c\,(\|Y\|+\|Y+W\|)^2$ near $0$ in $\mathcal{P}$. Indeed, let $N\ge2$, let
$E_{ij}$ be the matrix units, and put $Y_0:=E_{12}+E_{21}$, $X_0:=i\epsilon(E_{11}-E_{22})$
with $0<\epsilon<\rho_1$; both are traceless, $\|Y_0\|=1$, $\|X_0\|=\epsilon$, and
$[Y_0,X_0]=2i\epsilon(E_{21}-E_{12})$ has norm $2\epsilon\ne0$. For $0<t<\rho_1$ the point
$(tY_0,X_0,0)$ lies in $\mathcal{P}$, and by (R1)
\begin{equation*}
r(tY_0,X_0,0)=\operatorname{Ad}_{e^{-tY_0}}X_0-X_0
=-t[Y_0,X_0]+\sum_{n\ge2}\frac{(-t)^n}{n!}\operatorname{ad}_{Y_0}^{\,n}X_0 ,
\end{equation*}
by the identity $\operatorname{Ad}_{e^{-tY_0}}=\exp(-t\operatorname{ad}_{Y_0})$ of absolutely
convergent series, where $\operatorname{ad}_{Y_0}M:=[Y_0,M]$ for $M\in\mathfrak{A}$. Since
$\|\operatorname{ad}_{Y_0}M\|\le2\|Y_0\|\|M\|$, the sum over $n\ge2$
has norm at most $(e^{2a}-1-2a)\|X_0\|$ with $a:=t\|Y_0\|$, and for $a\le\tfrac12$ the last
clause of (e6) in \eqref{4.22:eq-frame-proof-tent-a} bounds this by $3a^2\|X_0\|$. Hence
\begin{equation*}
\|r(tY_0,X_0,0)\|\ge t\,\|[Y_0,X_0]\|-3t^2\|Y_0\|^2\|X_0\|=2\epsilon t-3\epsilon t^2 ,
\end{equation*}
while $(\|Y\|+\|Y+W\|)^2=4t^2$ at these points; the left side exceeds $4ct^2$ for all small
$t>0$. For commuting data $r_0\equiv0$, so the obstruction is non-abelian. Majorants in which
$X$ appears as a factor, such as (R4) in \eqref{4.22:eq-frame-proof-logarithm-a}, are
consistent with this example; (R4) is the bound used below, and the contribution of $r$ to the
increments is absorbed there by means of the coupling ceiling
\eqref{4.22:eq-coupling-ceilings-b}.
\end{remark}

\begin{proof}[Proof of the frame lemma, continued: increments of the eight entries]
\label{4.22:prf-frame-proof-increments}
We continue the proof of clause (i). Everything fixed above stays in force:
\begin{itemize}
\item the fine lattice, the covariant differences $\nabla^U$ and the cube gauges $u=u_{\square_s}$;
\item the tent field $\mathsf{Y}=\mathsf{Y}_{\mathsf{r}}[Z]$ and the map $\rho=e^{-\mathsf{Y}}=\rho_{\mathsf{r}}[Z]$;
\item the numbers $\lambda,\lambda',m_1,\mathfrak{n}_\chi,K,p_1,p_2$;
\item the estimates (e1)--(e6) of \eqref{4.22:eq-frame-proof-tent-a} and the bounds
\eqref{4.22:eq-frame-proof-tent-c} on $\mathsf{Y}$ and its first and second covariant differences;
\item the function $r=F-X-W$ on the polydiscs $\mathcal{P}'\subset\mathcal{P}$, its parts $r_0$ and $R_W$, the
radius $\rho_1=1/10$ and the numbers $\mathfrak{b}_2=200$, $\mathfrak{b}_3=400$ of
\eqref{4.22:eq-frame-proof-logarithm-a}.
\end{itemize}
Entries $1$--$4$, $5$ and $8$ are treated for $\max_s\|Z_s\|\le z\le 10^{-2}$. Entries $6$ and $7$ are treated
only for
\begin{equation*}
z\le z_{\max}:=4\varepsilon^{\mathrm{fr}}(W)=\varpi_{\mathrm{fr}}\alpha_{*,k_W}/C_{\mathrm{fr}},
\end{equation*}
and we recall that $z_{\max}<10^{-3}$. We write $\alpha_1:=\alpha_{1,k_W}$ and
$\alpha_*:=\alpha_{*,k_W}=\min\{\alpha_{0,k_W},\alpha_{1,k_W}\}$, so that $\alpha_*\le\alpha_1$. The thresholds of
entries $6$ and $7$ are $r_6=\bar r_6\alpha_1$ and $r_7=\bar r_7\alpha_1$, with $\bar r_6,\bar r_7\ge1$. For an
entry $a$ we write $Q_a(P)$ for the quantity it bounds; $\theta_a=r_a$ is its threshold on the tube
$\mathfrak{U}^{\mathrm{tr}}_N(X)$.

\emph{Decomposition.} Let $b=(x,\mu)$ be a bond. By the decomposition $\mathsf{r}=(U,\{u_\square\})$ we have
$\mathbb{U}(b)=U'(b)U(b)$ with $U'(b)=e^{i\eta A'_\mu(x)}$. With the convention
$\mathbb{U}^g(b)=g(b_-)\mathbb{U}(b)g(b_+)^{-1}$, inserting $U(b)^{-1}U(b)$ gives
\begin{equation*}
\mathbb{U}^\rho(b)=\rho(x)U'(b)U(b)\rho(x+\eta\mu)^{-1}
=\bigl[\rho(x)U'(b)\cdot U(b)\rho(x+\eta\mu)^{-1}U(b)^{-1}\bigr]\cdot U(b).
\end{equation*}
Since $\rho(x+\eta\mu)^{-1}=e^{\mathsf{Y}(x+\eta\mu)}$, the definition of $\nabla^U$ gives
\begin{equation*}
U(b)\rho(x+\eta\mu)^{-1}U(b)^{-1}=\exp\bigl(\operatorname{Ad}_{U(x,\mu)}\mathsf{Y}(x+\eta\mu)\bigr)
=\exp\bigl(\mathsf{Y}(x)+\eta(\nabla^U_\mu\mathsf{Y})(x)\bigr).
\end{equation*}
Put
\begin{equation*}
Y:=\mathsf{Y}(x),\qquad X:=i\eta A'_\mu(x),\qquad W:=\eta(\nabla^U_\mu\mathsf{Y})(x).
\end{equation*}
With $\rho(x)=e^{-Y}$ this yields
\begin{equation}\label{4.22:eq-frame-proof-increments-1}
\mathbb{U}^\rho=U'_\rho U,\qquad U'_\rho(b)=e^{-Y}e^{X}e^{Y+W},\qquad \mathbb{J}^\rho=\operatorname{Ad}_\rho\mathbb{J},
\end{equation}
with the same unitary part $U$ and the same cube gauges $\{u_\square\}$.

\emph{The coupling ceiling at $x=x_{k_W}$.} The ceiling \eqref{4.22:eq-coupling-ceilings-b} states that
$\mathfrak{c}\bar\alpha(x)\le1$ for all $x\ge\gamma_H^{-2}$, where
\begin{equation*}
\begin{gathered}
\mathfrak{c}:=\max\{20\bar r_6,\mathfrak{c}_6,\mathfrak{c}_7\},\qquad
\mathfrak{c}_6:=\mathfrak{b}_1+\mathfrak{b}_2(\bar p_1+\varpi_{\mathrm{fr}}),\\
\mathfrak{c}_7:=\mathfrak{b}_1\Bigl(1+\frac{\bar p_1\bar r_6}{\bar r_7}\Bigr)
+\mathfrak{b}_3\Bigl[\bar p_1(1+\varpi_{\mathrm{fr}})+4\varpi_{\mathrm{fr}}+\frac{\bar p_2\bar r_6}{\bar r_7}\Bigr],
\end{gathered}
\end{equation*}
and $\bar p_1=3^d(C_\chi+1)$, $\bar p_2=3^d(4C_\chi+2)$.

The radii at $k_W$ are controlled as follows. Since $\alpha\le\tilde\alpha$ for the coefficients of the tube, and
$\alpha^+_{k}$ is by definition the maximum of $\tilde\alpha_{0,k}$, $\tilde\alpha_{1,k}$ and the radius
$\delta'_{k}$ that fixes the range $6\delta'_k/g_k$ of the weak collar variables at level $k$, we have
\begin{equation*}
\alpha_1=\alpha_{1,k_W}\le\tilde\alpha_{1,k_W}\le\alpha^+_{k_W}=\bar\alpha(x_{k_W}).
\end{equation*}
The last equality is the identification $\bar\alpha(x_k)=\alpha^+_k$ made by the first clause of the theorem on the
functions of the coupling. Moreover $x_{k_W}>\gamma_H^{-2}$ (clause (0) of Theorem~0,
Theorem~\ref{4.1:thm-clause-zero}). Hence
\begin{equation}\label{4.22:eq-frame-proof-increments-2}
20\bar r_6\alpha_1\le1,\qquad \mathfrak{c}_6\alpha_1\le1,\qquad \mathfrak{c}_7\alpha_1\le1 .
\end{equation}
In particular $r_6=\bar r_6\alpha_1\le 1/20=\tfrac12\rho_1$.

\emph{Sizes (for entries $6$, $7$, so $z\le z_{\max}$).} We estimate the three exponents in turn.
\begin{itemize}
\item By \eqref{4.22:eq-frame-proof-tent-c}, $\|Y\|=\|\mathsf{Y}(x)\|\le z\le z_{\max}<10^{-3}$.
\item On the stencils, the tube gives $\ell\|A'_\mu\|<r_6$. Hence, by
\eqref{4.22:eq-frame-proof-increments-2},
\begin{equation*}
\|X\|=\eta\|A'_\mu(x)\|<(\eta/\ell)r_6\le r_6\le\tfrac12\rho_1 .
\end{equation*}
\item By \eqref{4.22:eq-frame-proof-tent-c}, using $\eta\le\ell$, $p_1<\bar p_1=3^d(C_\chi+1)\le C_{\mathrm{fr}}$ and
$\alpha_*\le1$,
\begin{equation*}
\|W\|\le(\eta/\ell)p_1z\le\bar p_1z_{\max}=\bar p_1\varpi_{\mathrm{fr}}\alpha_*/C_{\mathrm{fr}}
\le\varpi_{\mathrm{fr}}\le1/80 .
\end{equation*}
\end{itemize}
Since $\|Y\|+\|X\|<10^{-3}+1/20<1/10$ and $\|W\|\le1/80<1/10$, the triple $(Y,X,W)$ lies in $\mathcal{P}'$. Each of
its coordinates has norm at most $\tfrac12\rho_1$, and $\|Y\|<\tfrac12\rho_1$. The same holds at every site. It
therefore holds also for the transports of these field values by unitaries, since $\operatorname{Ad}_g$ is an
isometry for $g\in U(N)$.

The fields $\mathsf{Y}$ and $A'$ are $\mathfrak{g}^{\mathbb{C}}=\mathfrak{sl}(N,\mathbb{C})$-valued, and $\nabla^U$
preserves tracelessness, so $Y$, $X$, $W$ are traceless. By (R0) of \eqref{4.22:eq-frame-proof-logarithm-a}:
\begin{itemize}
\item $\|U'_\rho(b)-1\|<\tfrac12$;
\item $F(Y,X,W)$ is a logarithm of $U'_\rho(b)$, that is, $e^{F}=U'_\rho(b)$;
\item $\|F(Y,X,W)\|<\log2$ and $F(Y,X,W)$ is traceless.
\end{itemize}
Hence
\begin{equation*}
A'_{\rho,\mu}(x):=(i\eta)^{-1}F(Y,X,W)
\end{equation*}
is a $\mathfrak{g}^{\mathbb{C}}$-valued field with $e^{i\eta A'_\rho}=U'_\rho$. The tube \cite[(2.39)]{B-CMP119}
requires such a field $A'$ to exist, and $A'_\rho$ is the one used for $P^\rho$. Since $F=X+W+r$, we have
$i\eta A'_{\rho,\mu}=i\eta A'_\mu+\eta\nabla^U_\mu\mathsf{Y}+r$, that is,
\begin{equation}\label{4.22:eq-frame-proof-increments-3}
\delta A'_\mu(x):=A'_{\rho,\mu}(x)-A'_\mu(x)=-i\bigl[(\nabla^U_\mu\mathsf{Y})(x)+\eta^{-1}r_\mu(x)\bigr],
\qquad r_\mu(x):=r(Y,X,W).
\end{equation}

Suppose the stencil of a clause location misses $\mathsf{B}_W=\bigcup_s\operatorname{supp}\chi_s$. Then $\mathsf{Y}=0$
on that stencil, so $Y=W=0$, and $r(0,X,0)=0$ by (R0). Also $\rho=1$ there. Nothing moves at such locations, and
no reserve is kept there.

\emph{Entries $5$ and $8$.} These depend only on $U$ and its cube gauges $\{u_\square\}$. By
\eqref{4.22:eq-frame-proof-increments-1} both are unchanged, for $z\le10^{-2}$.

\emph{Entries $1$--$4$ ($z\le10^{-2}$).} The curvature-type entries $1$--$4$ and the current $\mathbb{J}$ are
conjugated by $\rho$. The entries derived from them are conjugated likewise, by the covariance identities for the
averaging operation and for the background configuration under complexified gauge transformations:
\cite[(11)]{B-CMP98}, valid for $G^{\mathbb{C}}$-valued transformations, and the complex form of
\cite[(181)]{B-CMP102b}. The block-constant transformation $\tilde\rho$ that enters there is taken, as in the
description of the frame domain, with $\|\log\tilde\rho\|\le z$.

Two bounds control the conjugations. First, (e4) gives $\|\operatorname{Ad}_\rho M\|\le e^{2\|\mathsf{Y}\|}\|M\|$.
Second, the transports $\operatorname{Ad}_{u(x)^{-1}u(s)}$ are $G$-valued, so by
\eqref{4.22:eq-frame-proof-tent-c}
\begin{equation*}
\|\mathsf{Y}(x)\|\le\sum_s\chi_s(x)\|Z_s\|\le z .
\end{equation*}
The same bound $e^{2z}$ holds for $\operatorname{Ad}_{\tilde\rho}$. Hence $Q_a(P^\rho)\le e^{2z}Q_a(P)$. Since
$Q_a(P)<\theta_a=r_a$ on the tube, (e6) gives
\begin{equation*}
Q_a(P^\rho)-Q_a(P)\le(e^{2z}-1)\theta_a\le\mathfrak{b}_1zr_a<2.1\,zr_a,\qquad a=1,\dots,4 .
\end{equation*}

\emph{Entry $6$ ($z\le z_{\max}$).} Take a bond $(x,\mu)$. We use \eqref{4.22:eq-frame-proof-increments-3}, the first
difference bound $\ell\|(\nabla^U_\mu\mathsf{Y})(x)\|\le p_1z$ of \eqref{4.22:eq-frame-proof-tent-c}, and (R4) on
$\mathcal{P}'$. This gives
\begin{equation*}
\begin{split}
\ell\|\delta A'_\mu(x)\|&\le\ell\|(\nabla^U_\mu\mathsf{Y})(x)\|+(\ell/\eta)\|r_\mu(x)\|\\
&\le p_1z+(\ell/\eta)\bigl[(e^{2\|Y\|}-1)\|X\|+\mathfrak{b}_2(\|Y\|+\|X\|)\|W\|\bigr].
\end{split}
\end{equation*}
We insert three bounds: $e^{2\|Y\|}-1\le\mathfrak{b}_1z$ (by (e6), since $\|Y\|\le z\le10^{-2}$);
$\|X\|<(\eta/\ell)r_6$; and $\|W\|\le(\eta/\ell)p_1z$. This gives
\begin{equation*}
\ell\|\delta A'_\mu(x)\|\le p_1z+(\ell/\eta)\bigl[\mathfrak{b}_1z(\eta/\ell)r_6
+\mathfrak{b}_2\bigl(z+(\eta/\ell)r_6\bigr)(\eta/\ell)p_1z\bigr].
\end{equation*}
Using $\eta\le\ell$, this is at most the number $\Delta_6$ of \eqref{4.22:eq-frame-proof-increments-a} below.
Entry $6$ therefore increases by at most $\sup\ell\|\delta A'_\mu\|\le\Delta_6$.

\emph{Entry $7$ ($z\le z_{\max}$).} We first split the second difference. Since $\nabla^U_\nu$ is linear,
\begin{equation*}
\ell^2\|(\nabla^U_\nu\delta A'_\mu)(x)\|\le\ell^2\|(\nabla^U_\nu\nabla^U_\mu\mathsf{Y})(x)\|
+(\ell^2/\eta)\|(\nabla^U_\nu r_\mu)(x)\|.
\end{equation*}
The first term is at most $p_2z$ by \eqref{4.22:eq-frame-proof-tent-c}.

For the second term, $(\nabla^U_\nu r_\mu)(x)=\eta^{-1}[\operatorname{Ad}_{U(x,\nu)}r_\mu(x+\eta\nu)-r_\mu(x)]$. By the
equivariance in (R0) under the $G$-valued $U(x,\nu)$, $\operatorname{Ad}_{U(x,\nu)}r_\mu(x+\eta\nu)=r(P_1)$, where
\begin{equation*}
\begin{split}
P_1&:=\bigl(\operatorname{Ad}_{U(x,\nu)}\mathsf{Y}(x+\eta\nu),\,
\operatorname{Ad}_{U(x,\nu)}i\eta A'_\mu(x+\eta\nu),\,
\operatorname{Ad}_{U(x,\nu)}\eta(\nabla^U_\mu\mathsf{Y})(x+\eta\nu)\bigr)\\
&\phantom{:}=P_0+(\delta Y,\delta X,\delta W),\qquad P_0:=(Y,X,W).
\end{split}
\end{equation*}
The definition of $\nabla^U$, applied to the site functions $\mathsf{Y}$, $X_\mu:=i\eta A'_\mu$ and
$W_\mu:=\eta\nabla^U_\mu\mathsf{Y}$, gives the three increments:
\begin{equation*}
\delta Y=\eta(\nabla^U_\nu\mathsf{Y})(x),\qquad
\delta X=i\eta^2(\nabla^U_\nu A'_\mu)(x),\qquad
\delta W=\eta^2(\nabla^U_\nu\nabla^U_\mu\mathsf{Y})(x).
\end{equation*}
By \eqref{4.22:eq-frame-proof-tent-c} and by the bound $\ell^2\|\nabla^U_\nu A'_\mu\|<r_7$ that the tube gives on
the stencils,
\begin{equation*}
\|\delta Y\|\le(\eta/\ell)p_1z,\qquad
\|\delta X\|<(\eta/\ell)^2r_7,\qquad
\|\delta W\|\le(\eta/\ell)^2p_2z .
\end{equation*}

The three coordinates of $P_1$ are transports by unitaries of field values at $x+\eta\nu$. By the paragraph on
sizes, $P_0$ and $P_1$ therefore both lie in $\mathcal{P}'$, all six coordinates have norm at most
$\tfrac12\rho_1$, and $\|Y_0\|<\tfrac12\rho_1$. No smallness of the increments is used for the domain. Writing
$r=r_0+R_W$, we get
\begin{equation*}
\frac{\ell^2}{\eta}\|(\nabla^U_\nu r_\mu)(x)\|=\Bigl(\frac{\ell}{\eta}\Bigr)^2\|r(P_1)-r(P_0)\|
\le\Bigl(\frac{\ell}{\eta}\Bigr)^2\bigl[\|r_0(Y_1,X_1)-r_0(Y_0,X_0)\|+\|R_W(P_1)-R_W(P_0)\|\bigr].
\end{equation*}

\emph{The $r_0$ part.} Apply (R1), with $\bar y:=\max\|Y_i\|\le z\le10^{-2}$. By (e6),
$e^{2\bar y}-1\le\mathfrak{b}_1z$ and $2e^{2\bar y}\le\mathfrak{b}_1$. With $\|X_0\|<(\eta/\ell)r_6$ this gives
\begin{equation*}
\|r_0(Y_1,X_1)-r_0(Y_0,X_0)\|\le\mathfrak{b}_1z(\eta/\ell)^2r_7+\mathfrak{b}_1(\eta/\ell)r_6\,(\eta/\ell)p_1z .
\end{equation*}

\emph{The $R_W$ part.} Apply (R3), with
$\|W_1\|=\|\eta(\nabla^U_\mu\mathsf{Y})(x+\eta\nu)\|\le(\eta/\ell)p_1z$. This gives
\begin{equation*}
\|R_W(P_1)-R_W(P_0)\|\le\mathfrak{b}_3\Bigl[\frac{\eta}{\ell}p_1z\Bigl(\frac{\eta}{\ell}p_1z
+\Bigl(\frac{\eta}{\ell}\Bigr)^2r_7\Bigr)+\Bigl(z+\frac{\eta}{\ell}r_6\Bigr)\Bigl(\frac{\eta}{\ell}\Bigr)^2p_2z\Bigr].
\end{equation*}

Multiplying both parts by $(\ell/\eta)^2$ and using $\eta\le\ell$ yields
\begin{equation*}
\frac{\ell^2}{\eta}\|(\nabla^U_\nu r_\mu)(x)\|\le\mathfrak{b}_1zr_7+\mathfrak{b}_1p_1zr_6
+\mathfrak{b}_3p_1z(p_1z+r_7)+\mathfrak{b}_3p_2z(z+r_6).
\end{equation*}
Entry $7$ therefore increases by at most $\Delta_7$, where, together with entry $6$,
\begin{equation}\label{4.22:eq-frame-proof-increments-a}
\begin{gathered}
\Delta_6:=p_1z+\mathfrak{b}_1zr_6+\mathfrak{b}_2p_1z(z+r_6),\\
\Delta_7:=p_2z+\mathfrak{b}_1z(r_7+p_1r_6)+\mathfrak{b}_3\bigl[p_1z(p_1z+r_7)+p_2z(z+r_6)\bigr].
\end{gathered}
\end{equation}

\emph{Conversion of $\Delta_6$ into units of $zr_6/\alpha_1$.} Divide $\Delta_6$ term by term by
$zr_6/\alpha_1=z\bar r_6$.
\begin{itemize}
\item First term: $p_1/\bar r_6\le p_1<\bar p_1$, since $\bar r_6\ge1$.
\item Second term: $\mathfrak{b}_1\alpha_1$.
\item Third term, with $r_6$: $\mathfrak{b}_2p_1\alpha_1\le\mathfrak{b}_2\bar p_1\alpha_1$.
\item Third term, with $z$: we have $z_{\max}=\varpi_{\mathrm{fr}}\alpha_*/C_{\mathrm{fr}}\le\varpi_{\mathrm{fr}}\alpha_1/C_{\mathrm{fr}}$
(as $\alpha_*\le\alpha_1$) and $C_{\mathrm{fr}}\bar r_6\ge C_{\mathrm{fr}}\ge\bar p_1$. Hence
\begin{equation*}
\frac{\mathfrak{b}_2p_1z}{\bar r_6}\le\frac{\mathfrak{b}_2\bar p_1z_{\max}}{\bar r_6}
\le\frac{\mathfrak{b}_2\bar p_1\varpi_{\mathrm{fr}}\alpha_1}{C_{\mathrm{fr}}\bar r_6}\le\mathfrak{b}_2\varpi_{\mathrm{fr}}\alpha_1 .
\end{equation*}
\end{itemize}
The share of $r$ is thus at most
\begin{equation*}
[\mathfrak{b}_1+\mathfrak{b}_2(\bar p_1+\varpi_{\mathrm{fr}})]\alpha_1\cdot z\bar r_6
=\mathfrak{c}_6\alpha_1\cdot z\,r_6/\alpha_1 .
\end{equation*}
By \eqref{4.22:eq-frame-proof-increments-2} this is at most $zr_6/\alpha_1$. Therefore
\begin{equation*}
\Delta_6\le(\bar p_1+1)\,z\,\frac{r_6}{\alpha_1}=\bigl(3^d(C_\chi+1)+1\bigr)z\,\frac{r_6}{\alpha_1}
\le C_{\mathrm{fr}}\,z\,\frac{r_6}{\alpha_{1,k_W}},
\end{equation*}
because $3^d(C_\chi+1)+1\le3^d(C_\chi+4)$.

\emph{Conversion of $\Delta_7$ into units of $zr_7/\alpha_1$.} Divide $\Delta_7$ term by term by
$zr_7/\alpha_1=z\bar r_7$.
\begin{itemize}
\item First term: $p_2/\bar r_7\le p_2<3^d(C_\chi+2)$, since $\bar r_7\ge1$ and by \eqref{4.22:eq-frame-proof-tent-c}.
\item $\mathfrak{b}_1zr_7$ gives $\mathfrak{b}_1\alpha_1$.
\item $\mathfrak{b}_1p_1zr_6$ gives $\mathfrak{b}_1p_1\alpha_1\bar r_6/\bar r_7\le\mathfrak{b}_1\bar p_1\alpha_1\bar r_6/\bar r_7$.
\item $\mathfrak{b}_3p_1^2z^2$ gives, since $C_{\mathrm{fr}}\bar r_7\ge C_{\mathrm{fr}}\ge\bar p_1$,
\begin{equation*}
\frac{\mathfrak{b}_3p_1^2z}{\bar r_7}\le\frac{\mathfrak{b}_3\bar p_1\cdot\bar p_1\varpi_{\mathrm{fr}}\alpha_1}{C_{\mathrm{fr}}\bar r_7}
\le\mathfrak{b}_3\bar p_1\varpi_{\mathrm{fr}}\alpha_1 .
\end{equation*}
\item $\mathfrak{b}_3p_1zr_7$ gives $\mathfrak{b}_3p_1\alpha_1\le\mathfrak{b}_3\bar p_1\alpha_1$.
\item $\mathfrak{b}_3p_2z^2$ gives, since $p_2\le\bar p_2$ and
$\bar p_2/C_{\mathrm{fr}}=(4C_\chi+2)/(C_\chi+4)\le4$,
\begin{equation*}
\frac{\mathfrak{b}_3p_2z}{\bar r_7}\le\frac{\mathfrak{b}_3\bar p_2\varpi_{\mathrm{fr}}\alpha_1}{C_{\mathrm{fr}}\bar r_7}
\le4\mathfrak{b}_3\varpi_{\mathrm{fr}}\alpha_1 .
\end{equation*}
With the sharper $p_2<3^d(C_\chi+2)\le C_{\mathrm{fr}}$ one even gets $\mathfrak{b}_3\varpi_{\mathrm{fr}}\alpha_1$.
\item $\mathfrak{b}_3p_2zr_6$ gives $\mathfrak{b}_3p_2\alpha_1\bar r_6/\bar r_7\le\mathfrak{b}_3\bar p_2\alpha_1\bar r_6/\bar r_7$.
\end{itemize}
The share of $r$ is thus at most $\mathfrak{c}_7\alpha_1\cdot zr_7/\alpha_1$, where $\mathfrak{c}_7$ is the number
\begin{equation*}
\mathfrak{c}_7=\mathfrak{b}_1\Bigl(1+\frac{\bar p_1\bar r_6}{\bar r_7}\Bigr)
+\mathfrak{b}_3\Bigl[\bar p_1(1+\varpi_{\mathrm{fr}})+4\varpi_{\mathrm{fr}}+\frac{\bar p_2\bar r_6}{\bar r_7}\Bigr].
\end{equation*}
By \eqref{4.22:eq-frame-proof-increments-2} this share is at most $zr_7/\alpha_1$. Therefore
\begin{equation*}
\Delta_7\le\bigl(3^d(C_\chi+2)+1\bigr)z\,\frac{r_7}{\alpha_1}\le C_{\mathrm{fr}}\,z\,\frac{r_7}{\alpha_{1,k_W}},
\end{equation*}
because $3^d(C_\chi+2)+1\le3^d(C_\chi+4)$.

\emph{The constant $C_{\mathrm{fr}}$.} Three inputs combine here:
\begin{itemize}
\item at most $3^d$ tents overlap, since $\mathfrak{n}_\chi\le3^d$; this uses $\eta/\ell=L^{-k_W}\le\tfrac12$, which
holds for every blocking factor because $k_W\ge1$;
\item the per-tent counts are $C_\chi+1$ and $K<C_\chi+2\le C_\chi+3$, both under
\eqref{4.22:eq-coupling-ceilings-a};
\item the share of $r$ is at most $z\,r_a/\alpha_{1,k_W}$ under \eqref{4.22:eq-coupling-ceilings-b}, since
$\bar r_a\ge1$.
\end{itemize}
Hence, for $z\le4\varepsilon^{\mathrm{fr}}(W)$,
\begin{equation}\label{4.22:eq-frame-proof-increments-b}
\Delta_a\le C_{\mathrm{fr}}\,z\,\frac{r_a}{\alpha_{1,k_W}}\quad(a=6,7),\qquad C_{\mathrm{fr}}=3^d(C_\chi+4).
\end{equation}

\emph{The last sentence of clause (i).} Recall $\varepsilon^{\mathrm{fr}}(W)=\varpi_{\mathrm{fr}}\alpha_*/(4C_{\mathrm{fr}})$.
Consider first a map of size $z\le\varepsilon^{\mathrm{fr}}(W)$.
\begin{itemize}
\item Entries $6$, $7$: since $\alpha_*\le\alpha_{1,k_W}$,
\begin{equation*}
C_{\mathrm{fr}}z/\alpha_{1,k_W}\le C_{\mathrm{fr}}\varepsilon^{\mathrm{fr}}(W)/\alpha_*=\tfrac14\varpi_{\mathrm{fr}} .
\end{equation*}
\item Entries $1$--$4$: since $\alpha_*\le1$ and $C_{\mathrm{fr}}\ge5$, we have $2.1\alpha_*\le C_{\mathrm{fr}}$, and
\begin{equation*}
(e^{2z}-1)\theta_a\le2.1\,zr_a\le\tfrac14\varpi_{\mathrm{fr}}r_a .
\end{equation*}
\end{itemize}
More generally, a map of size $z\le4\varepsilon^{\mathrm{fr}}(W)$ increases every entry by at most
$(z/\varepsilon^{\mathrm{fr}}(W))\cdot\tfrac14\varpi_{\mathrm{fr}}r_a$.
\begin{itemize}
\item For entries $6$, $7$,
\begin{equation*}
C_{\mathrm{fr}}zr_a/\alpha_1=\frac{z}{\varepsilon^{\mathrm{fr}}(W)}\cdot\frac{\varpi_{\mathrm{fr}}}{4}\cdot
\frac{\alpha_*}{\alpha_1}\,r_a\le\frac{z}{\varepsilon^{\mathrm{fr}}(W)}\cdot\frac{\varpi_{\mathrm{fr}}}{4}\,r_a .
\end{equation*}
\item For entries $1$--$4$, since $2.1\alpha_*\le C_{\mathrm{fr}}$,
\begin{equation*}
2.1\,zr_a\le C_{\mathrm{fr}}zr_a/\alpha_*=\frac{z}{\varepsilon^{\mathrm{fr}}(W)}\cdot\frac{\varpi_{\mathrm{fr}}}{4}\,r_a .
\end{equation*}
\end{itemize}

\emph{Successive maps at one decomposition.} Let $P\in\mathfrak{U}_N(X)$, the stored domain, which keeps at every
entry $a$ the reserve $\varpi_{\mathrm{fr}}r_a$ below the threshold of the tube. Let
$\rho^{(1)},\dots,\rho^{(m)}$ be given by $\rho^{(i)}=\rho_{\mathsf{r}}[Z_i]$ at one and the same decomposition
$\mathsf{r}$, with sizes $z_i=\max_s\|(Z_i)_s\|$ and $\sum_iz_i<4\varepsilon^{\mathrm{fr}}(W)$. Put $P^{(0)}:=P$ and
$P^{(j)}:=(P^{(j-1)})^{\rho^{(j)}}$.

Then every $P^{(j)}$ lies in $\mathfrak{U}^{\mathrm{tr}}_N(X)$ and is decomposed with the same
$(U,\{u_\square\})$. Moreover, for $a\in\{1,\dots,4,6,7\}$,
\begin{equation}\label{4.22:eq-frame-proof-increments-c}
Q_a(P^{(j)})-Q_a(P)\le\frac{\sum_{i\le j}z_i}{\varepsilon^{\mathrm{fr}}(W)}\cdot\frac{\varpi_{\mathrm{fr}}}{4}\,r_a
<\varpi_{\mathrm{fr}}r_a ,
\end{equation}
while entries $5$ and $8$ do not move.

We argue by induction on $j$. The case $j=0$ is trivial. Assume the claim for $j-1$: $P^{(j-1)}$ lies in
$\mathfrak{U}^{\mathrm{tr}}_N(X)$, is decomposed with the same $U$ and cube gauges, and each entry has increased so
far by at most
\begin{equation*}
\frac{\sum_{i<j}z_i}{\varepsilon^{\mathrm{fr}}(W)}\cdot\frac{\varpi_{\mathrm{fr}}}{4}\,r_a<\varpi_{\mathrm{fr}}r_a .
\end{equation*}
The estimates above therefore apply to $P^{(j-1)}$: its own field $A'$ obeys the thresholds,
$z_j\le\sum_iz_i<z_{\max}<10^{-2}$, and $\rho^{(j)}$ is built from the same $U$ and $\{u_\square\}$.

The increments of entries $6$ and $7$ add, giving the bound \eqref{4.22:eq-frame-proof-increments-c}.

Entries $1$--$4$ compose multiplicatively:
$Q_a(P^{(j)})\le e^{2z_j}Q_a(P^{(j-1)})\le e^{2\sum_{i\le j}z_i}Q_a(P)$. Indeed, the increment after two
successive maps of sizes $z'$ and $z''$ is at most
\begin{equation*}
(e^{2z'}-1)\theta_a+e^{2z'}(e^{2z''}-1)\theta_a=(e^{2(z'+z'')}-1)\theta_a ,
\end{equation*}
the increment of one map of size $z'+z''$.
The total increment is therefore at most
\begin{equation*}
\bigl(e^{2\sum_{i\le j}z_i}-1\bigr)\theta_a\le2.1\Bigl(\sum_{i\le j}z_i\Bigr)r_a
\le\frac{\sum_{i\le j}z_i}{\varepsilon^{\mathrm{fr}}(W)}\cdot\frac{\varpi_{\mathrm{fr}}}{4}\,r_a ,
\end{equation*}
by (e6) and $2.1\alpha_*\le C_{\mathrm{fr}}$.

Entries $5$ and $8$ never move. The reserve $\varpi_{\mathrm{fr}}r_a$ of $P\in\mathfrak{U}_N(X)$ absorbs the total,
so $P^{(j)}\in\mathfrak{U}^{\mathrm{tr}}_N(X)$. This completes the induction, and with it the estimates of clause (i)
on the increments of the eight entries; the proof of the frame lemma continues with the profile class and
clauses (ii)--(vi).
\end{proof}

\begin{corollary}[The profile class and the cost of a gauge map in it]\label{4.22:cor-profile-class}
Assume the hypotheses of the frame lemma, in particular the two regularity ceilings on the
coupling read at $x_{k_W}$ (the one bounding the cube-gauge potential $A$, on which the
tent-difference bounds rest, and $\mathfrak{c}\,\bar\alpha(x_{k_W})\le1$), and
$\bar r_6,\bar r_7\ge1$. We use the notation of the proof of the frame lemma
(\ref{4.22:prf-frame-proof-tent} and \ref{4.22:prf-frame-proof-increments}):
\begin{itemize}
\item $T_\eta$ is the fine lattice of spacing $\eta$, and $\ell=L^{k_W}\eta$;
\item $\mathsf{r}=(U,\{u_{\square}\})$ is a decomposition of a pair
$P=(\mathbb{U},\mathbb{J})$, with $\mathbb{U}(b)=e^{i\eta A'_\mu(x)}U(b)$ for $b=(x,\mu)$;
\item for an $\mathfrak{A}$-valued site function $f$, where
$\mathfrak{A}:=\mathfrak{gl}(N,\mathbb{C})$, the covariant difference is
$(\nabla^U_\mu f)(x):=\eta^{-1}[\operatorname{Ad}_{U(x,\mu)}f(x+\eta\mu)-f(x)]$;
\item $\mathsf{B}_W=\bigcup_{s\in\Sigma_W}\operatorname{supp}\chi_s$;
\item $Q_a$, $a=1,\dots,8$, are the eight entries defining the tube
$\mathfrak{U}^{\mathrm{tr}}_N(X)$, with thresholds $\theta_a=r_a$; here
$r_6=\bar r_6\alpha_{1,k_W}$ and $r_7=\bar r_7\alpha_{1,k_W}$.
\end{itemize}
Put $\bar p_1:=3^d(C_\chi+1)$ and $\hat p_2:=3^d(C_\chi+2)$. Then
$\hat p_2\le\bar p_2:=3^d(4C_\chi+2)$ and $\hat p_2+1\le C_{\mathrm{fr}}$.

For $z'\ge0$ let $\mathfrak{P}(z')$ be the set of site functions
$\mathsf{Y}\colon T_\eta\to\mathfrak{g}^{\mathbb{C}}=\mathfrak{sl}(N,\mathbb{C})$ such that
$\mathsf{Y}=0$ off $\mathsf{B}_W$ and
\begin{equation}\label{4.22:eq-profile-class-a}
\begin{gathered}
\sup_x\|\mathsf{Y}(x)\|\le z',\qquad
\sup_{x,\mu}\ell\,\|(\nabla^U_\mu\mathsf{Y})(x)\|\le\bar p_1z',\\
\sup_{x,\mu,\nu}\ell^2\,\|(\nabla^U_\nu\nabla^U_\mu\mathsf{Y})(x)\|\le\hat p_2z'.
\end{gathered}
\end{equation}
The class depends on $U$ and $W$; we say that the profile is measured with the $U$ of $\mathsf{r}$.

The class has the following elementary properties:
\begin{itemize}
\item $\mathfrak{P}(z')=-\mathfrak{P}(z')$;
\item $\mathfrak{P}(z')\subset\mathfrak{P}(z'')$ for $z'\le z''$;
\item $t\mathsf{Y}\in\mathfrak{P}(tz')$ for $\mathsf{Y}\in\mathfrak{P}(z')$ and $t\ge0$;
\item $\mathsf{Y}_{\mathsf{r}}[Z]\in\mathfrak{P}(z)$ whenever $\max_s\|Z_s\|\le z$.
\end{itemize}
Write $z_{\max}:=4\varepsilon^{\mathrm{fr}}(W)=\varpi_{\mathrm{fr}}\alpha_{*,k_W}/C_{\mathrm{fr}}$.
Note that $z_{\max}<10^{-3}$, since $\varpi_{\mathrm{fr}}\le1/80$, $\alpha_{*,k_W}\le1$ and
$C_{\mathrm{fr}}\ge15$ (\ref{4.22:prf-frame-proof-tent}).
\begin{itemize}
\item[(a)] Let $\mathsf{Y}\in\mathfrak{P}(z')$ with $z'\le z_{\max}$, and let
$P\in\mathfrak{U}^{\mathrm{tr}}_N(X)$ carry the decomposition $\mathsf{r}$. In particular
$\mathbb{U}=e^{i\eta A'}U$ with $\ell\|A'_\mu\|<r_6$ and $\ell^2\|\nabla^U_\nu A'_\mu\|<r_7$ on
the stencils of the locations at which entries $6$ and $7$ are read, and $U$ comes with its
cube gauges $\{u_{\square}\}$. Then
$P^{e^{-\mathsf{Y}}}=(U'_{e^{-\mathsf{Y}}}U,\operatorname{Ad}_{e^{-\mathsf{Y}}}\mathbb{J})$
is decomposed with the same $(U,\{u_{\square}\})$, and:
\begin{itemize}
\item entries $5$ and $8$ are unchanged;
\item entries $1$--$4$ are multiplied by at most $e^{2z'}$, so that their increments are at
most $\mathfrak{b}_1z'r_a<2.1z'r_a$;
\item entries $6$ and $7$ increase by at most $(\bar p_1+1)z'r_6/\alpha_{1,k_W}$ and
$(\hat p_2+1)z'r_7/\alpha_{1,k_W}$ respectively, both at most
$C_{\mathrm{fr}}z'r_a/\alpha_{1,k_W}$;
\item nothing moves at a location whose stencil misses $\mathsf{B}_W$.
\end{itemize}
Hence the map costs at most
$(z'/\varepsilon^{\mathrm{fr}}(W))\cdot\tfrac14\varpi_{\mathrm{fr}}r_a$ at every entry $a$.
\item[(b)] \textup{(Successive maps.)} Let $P\in\mathfrak{U}_N(X)$ carry a decomposition
$\mathsf{r}=(U,\{u_{\square}\})$, and let
$\mathsf{Y}_j\in\mathfrak{P}(z_j)$, $j=1,\dots,m$, with profiles measured with this $U$, with
$\sum_jz_j<4\varepsilon^{\mathrm{fr}}(W)$. Then every successive image
$P^{(j)}:=(P^{(j-1)})^{e^{-\mathsf{Y}_j}}$, $P^{(0)}:=P$, lies in
$\mathfrak{U}^{\mathrm{tr}}_N(X)$ and is decomposed with the same $(U,\{u_{\square}\})$.
The total cost at every entry $a$ is at most
\begin{equation}\label{4.22:eq-profile-class-1}
\frac{\sum_{j=1}^m z_j}{\varepsilon^{\mathrm{fr}}(W)}\cdot\frac{\varpi_{\mathrm{fr}}}{4}\,r_a
<\varpi_{\mathrm{fr}}\,r_a .
\end{equation}
\end{itemize}
\end{corollary}

\begin{proof}
\emph{Elementary properties.} For every $\mu$ the covariant difference $\nabla^U_\mu$ is
$\mathbb{C}$-linear, because $\operatorname{Ad}_{U(x,\mu)}$ is linear. Hence, for
$\mathsf{Y}\in\mathfrak{P}(z')$ and $t\ge0$, the three suprema in
\eqref{4.22:eq-profile-class-a} for $-\mathsf{Y}$ equal those for $\mathsf{Y}$, and those for
$t\mathsf{Y}$ are $t$ times those for $\mathsf{Y}$. Moreover $-\mathsf{Y}$ and $t\mathsf{Y}$ are
$\mathfrak{sl}(N,\mathbb{C})$-valued and vanish off $\mathsf{B}_W$. This gives the first and
third properties. The second is immediate from the definition.

Let $\max_s\|Z_s\|\le z$. We have $\mathsf{Y}_{\mathsf{r}}[Z]=\sum_sf_s$ with
$f_s=\chi_s\operatorname{Ad}_{u^{-1}}\tilde Z_s$. Each $f_s$ vanishes off
$\operatorname{supp}\chi_s\subset\mathsf{B}_W$, and it is
$\mathfrak{sl}(N,\mathbb{C})$-valued since $Z_s$ is and $\operatorname{Ad}$ preserves
tracelessness. By the first and second tent-difference bounds
\eqref{4.22:eq-frame-proof-tent-c}:
\begin{itemize}
\item $\|\mathsf{Y}_{\mathsf{r}}[Z](x)\|\le z$;
\item $\ell\|\nabla^U_\mu\mathsf{Y}_{\mathsf{r}}[Z]\|\le p_1z$ with $p_1<\bar p_1$;
\item $\ell^2\|\nabla^U_\nu\nabla^U_\mu\mathsf{Y}_{\mathsf{r}}[Z]\|\le p_2z$ with
$p_2<3^d(C_\chi+2)=\hat p_2$.
\end{itemize}
Thus $\mathsf{Y}_{\mathsf{r}}[Z]\in\mathfrak{P}(z)$.

\emph{Proof of (a): the decomposition.} Put $\rho:=e^{-\mathsf{Y}}$ and fix
$b=(x,\mu)$. With $\mathbb{U}(b)=U'(b)U(b)$ and $U'(b)=e^{i\eta A'_\mu(x)}$,
\begin{align*}
\mathbb{U}^\rho(b)&=\rho(x)U'(b)U(b)\rho(x+\eta\mu)^{-1}\\
&=\bigl[\rho(x)U'(b)\cdot U(b)\rho(x+\eta\mu)^{-1}U(b)^{-1}\bigr]U(b).
\end{align*}
Since $\rho^{-1}=e^{\mathsf{Y}}$ and $Ve^{M}V^{-1}=e^{\operatorname{Ad}_VM}$ for invertible $V$,
and by the definition of $\nabla^U_\mu$ in the last step,
\begin{align*}
U(b)\rho(x+\eta\mu)^{-1}U(b)^{-1}
&=U(b)e^{\mathsf{Y}(x+\eta\mu)}U(b)^{-1}\\
&=\exp\bigl(\operatorname{Ad}_{U(x,\mu)}\mathsf{Y}(x+\eta\mu)\bigr)\\
&=\exp\bigl(\mathsf{Y}(x)+\eta(\nabla^U_\mu\mathsf{Y})(x)\bigr).
\end{align*}
Hence $\mathbb{U}^\rho=U'_\rho U$ with the same $U$, and therefore the same cube gauges, where
\begin{equation}\label{4.22:eq-profile-class-2}
U'_\rho(b)=e^{-Y}e^{X}e^{Y+W},\qquad
Y:=\mathsf{Y}(x),\quad X:=i\eta A'_\mu(x),\quad W:=\eta(\nabla^U_\mu\mathsf{Y})(x).
\end{equation}
Also $\mathbb{J}^\rho=\operatorname{Ad}_\rho\mathbb{J}$. This identity holds for every site
function $\mathsf{Y}$; in particular it is the one used for $\mathsf{Y}_{\mathsf{r}}[Z]$ in
\ref{4.22:prf-frame-proof-increments}.

\emph{Sizes.} Let $\alpha_1:=\alpha_{1,k_W}$ and $\alpha_*:=\alpha_{*,k_W}\le\alpha_1$,
$\alpha_*\le 1$. The regularity ceiling on the coupling, $\mathfrak{c}\,\bar\alpha(x)\le1$
for $x\ge\gamma_H^{-2}$ with $\mathfrak{c}=\max\{20\bar r_6,\mathfrak{c}_6,\mathfrak{c}_7\}$,
is taken at $x_{k_W}$ in \ref{4.22:prf-frame-proof-increments}. There it gives
$20\bar r_6\alpha_1\le1$, $\mathfrak{c}_6\alpha_1\le1$ and $\mathfrak{c}_7\alpha_1\le1$, where
\begin{align*}
\mathfrak{c}_6&=\mathfrak{b}_1+\mathfrak{b}_2(\bar p_1+\varpi_{\mathrm{fr}}),\\
\mathfrak{c}_7&=\mathfrak{b}_1\bigl(1+\bar p_1\bar r_6/\bar r_7\bigr)
+\mathfrak{b}_3\bigl[\bar p_1(1+\varpi_{\mathrm{fr}})+4\varpi_{\mathrm{fr}}
+\bar p_2\bar r_6/\bar r_7\bigr].
\end{align*}
In particular $r_6\le1/20$. From \eqref{4.22:eq-profile-class-a} and $z'\le z_{\max}<10^{-3}$
we get $\|Y\|\le z'<10^{-3}$. Also
\begin{equation*}
\|X\|<(\eta/\ell)r_6\le r_6\le1/20 .
\end{equation*}
Finally, using
$\bar p_1\le C_{\mathrm{fr}}$, $\alpha_*\le1$ and $\varpi_{\mathrm{fr}}\le 1/80$,
\begin{equation*}
\|W\|\le(\eta/\ell)\bar p_1z'\le\bar p_1z_{\max}
=\bar p_1\varpi_{\mathrm{fr}}\alpha_*/C_{\mathrm{fr}}\le\varpi_{\mathrm{fr}}\le 1/80 .
\end{equation*}
With $\rho_1=1/10$, the triple $(Y,X,W)$ therefore lies in the domain on which
$\|Y\|+\|X\|<\rho_1$ and $\|W\|<\rho_1$. Each coordinate has norm at most $\tfrac12\rho_1$,
and $\|Y\|<\tfrac12\rho_1$. The same holds at every site. Since $U$ is $G$-valued and
$\operatorname{Ad}_g$ is an isometry for $g\in U(N)$, it also holds for the unitary transports
of these field values.

These are exactly the hypotheses of the structure lemma for the logarithm $F(Y,X,W)$ of
$e^{-Y}e^Xe^{Y+W}$, stated in the proof of the frame lemma. By that lemma
$\|U'_\rho(b)-1\|<\tfrac12$, and $F$ is the logarithm of $U'_\rho(b)$, with norm $<\log2$.
The matrices $Y$ and $W$ are traceless since $\mathsf{Y}$ is and $\operatorname{Ad}$ preserves
the trace, and $X$ is traceless since $A'$ is $\mathfrak{g}^{\mathbb{C}}$-valued; by the
structure lemma $F$ is then traceless, and
$A'_{\rho,\mu}(x):=(i\eta)^{-1}F(Y,X,W)$ is $\mathfrak{g}^{\mathbb{C}}$-valued with
$e^{i\eta A'_\rho}=U'_\rho$. Writing $r=F-X-W$ and $r_\mu(x):=r(Y,X,W)$, we obtain
\begin{equation*}
\delta A'_\mu(x):=A'_{\rho,\mu}(x)-A'_\mu(x)
=-i\bigl[(\nabla^U_\mu\mathsf{Y})(x)+\eta^{-1}r_\mu(x)\bigr].
\end{equation*}

\emph{Entries $5$ and $8$.} They depend only on $U$ and its cube gauges, which do not change.

\emph{Entries $1$--$4$.} By the covariance of entries $1$--$4$ under conjugation recalled in
\ref{4.22:prf-frame-proof-increments}, these entries are conjugated either by $\rho$ or by a
block-constant value $\tilde\rho$ of $\rho$. Since $\tilde\rho$ is a value of $\rho$ on its
block, $\|\log\tilde\rho\|\le\sup\|\mathsf{Y}\|\le z'$. By the conjugation estimate in
\eqref{4.22:eq-frame-proof-tent-a}, $\|\operatorname{Ad}_\rho\|\le e^{2\|\mathsf{Y}\|}$, so
$Q_a(P^\rho)\le e^{2z'}Q_a(P)$. The increment is therefore at most
$(e^{2z'}-1)\theta_a\le\mathfrak{b}_1z'r_a<2.1z'r_a$, by the elementary bound
$e^{2y}-1\le\mathfrak{b}_1y$ for $0\le y\le10^{-2}$ in \eqref{4.22:eq-frame-proof-tent-a}.

\emph{Entry $6$.} The structure lemma gives, on the domain containing $(Y,X,W)$ found above,
$\|r\|\le(e^{2\|Y\|}-1)\|X\|+\mathfrak{b}_2(\|Y\|+\|X\|)\|W\|$. Together with the size bounds
on $Y$, $X$, $W$ established above, the bound $e^{2y}-1\le\mathfrak{b}_1y$ and $\eta\le\ell$,
this yields
\begin{align*}
\ell\|\delta A'_\mu(x)\|
&\le\ell\|(\nabla^U_\mu\mathsf{Y})(x)\|+(\ell/\eta)\|r_\mu(x)\|\\
&\le\bar p_1z'+(\ell/\eta)\bigl[\mathfrak{b}_1z'(\eta/\ell)r_6
+\mathfrak{b}_2\bigl(z'+(\eta/\ell)r_6\bigr)(\eta/\ell)\bar p_1z'\bigr]\\
&\le\bar p_1z'+\mathfrak{b}_1z'r_6+\mathfrak{b}_2\bar p_1z'(z'+r_6).
\end{align*}
This is the bound $\Delta_6$ of \eqref{4.22:eq-frame-proof-increments-a} with $p_1$ replaced
by $\bar p_1$.

Divide by $z'r_6/\alpha_1=z'\bar r_6$. The first term gives
$\bar p_1/\bar r_6\le\bar p_1$, since $\bar r_6\ge1$. The remaining terms give
$\mathfrak{b}_1\alpha_1$, then $\mathfrak{b}_2\bar p_1\alpha_1$, and
\begin{equation*}
\mathfrak{b}_2\bar p_1z'/\bar r_6
\le\mathfrak{b}_2\bar p_1\varpi_{\mathrm{fr}}\alpha_1/(C_{\mathrm{fr}}\bar r_6)
\le\mathfrak{b}_2\varpi_{\mathrm{fr}}\alpha_1 .
\end{equation*}
Here we used $z'\le z_{\max}\le\varpi_{\mathrm{fr}}\alpha_1/C_{\mathrm{fr}}$ and
$C_{\mathrm{fr}}\bar r_6\ge C_{\mathrm{fr}}\ge\bar p_1$. These three bounds sum to
$[\mathfrak{b}_1+\mathfrak{b}_2(\bar p_1+\varpi_{\mathrm{fr}})]\alpha_1=\mathfrak{c}_6\alpha_1$,
and $\mathfrak{c}_6\alpha_1\le1$. Hence entry $6$ increases by at most
\begin{equation*}
(\bar p_1+1)z'r_6/\alpha_1\le C_{\mathrm{fr}}z'r_6/\alpha_{1,k_W},
\end{equation*}
because $3^d(C_\chi+1)+1\le3^d(C_\chi+4)$.

\emph{Entry $7$.} We have
\begin{equation*}
\ell^2\|(\nabla^U_\nu\delta A'_\mu)(x)\|
\le\ell^2\|(\nabla^U_\nu\nabla^U_\mu\mathsf{Y})(x)\|
+(\ell^2/\eta)\|(\nabla^U_\nu r_\mu)(x)\|,
\end{equation*}
and the first term is at most $\hat p_2z'$. For the second, the structure lemma's
equivariance under the $G$-valued $U(x,\nu)$ gives
$\operatorname{Ad}_{U(x,\nu)}r_\mu(x+\eta\nu)=r(P_1)$, where
\begin{equation*}
P_1:=\bigl(\operatorname{Ad}_{U(x,\nu)}\mathsf{Y}(x+\eta\nu),\,
\operatorname{Ad}_{U(x,\nu)}i\eta A'_\mu(x+\eta\nu),\,
\operatorname{Ad}_{U(x,\nu)}\eta(\nabla^U_\mu\mathsf{Y})(x+\eta\nu)\bigr).
\end{equation*}
By the definition of $\nabla^U$, $P_1=P_0+(\delta Y,\delta X,\delta W)$ with $P_0:=(Y,X,W)$,
\begin{align*}
\delta Y&=\eta(\nabla^U_\nu\mathsf{Y})(x), & \|\delta Y\|&\le(\eta/\ell)\bar p_1z',\\
\delta X&=i\eta^2(\nabla^U_\nu A'_\mu)(x), & \|\delta X\|&<(\eta/\ell)^2r_7,\\
\delta W&=\eta^2(\nabla^U_\nu\nabla^U_\mu\mathsf{Y})(x), & \|\delta W\|&\le(\eta/\ell)^2\hat p_2z'.
\end{align*}
The coordinates of $P_1$ are unitary transports of field values at $x+\eta\nu$. By the size
bounds established above, $P_0$ and $P_1$ both lie in the domain of the structure lemma, with
all six coordinates of norm at most $\tfrac12\rho_1$ and $\|Y\|<\tfrac12\rho_1$.

Write $P_i=(Y_i,X_i,W_i)$, $i=0,1$, and $\bar y:=\max\{\|Y_0\|,\|Y_1\|\}$. By the size bounds,
$\bar y\le z'\le10^{-2}$, so $e^{2\bar y}-1\le\mathfrak{b}_1z'$ and
$2e^{2\bar y}\le\mathfrak{b}_1$. Since $r=r_0+R_W$ with $r_0=r(\cdot,\cdot,0)$,
\begin{equation*}
\|r(P_1)-r(P_0)\|\le\|r_0(Y_1,X_1)-r_0(Y_0,X_0)\|+\|R_W(P_1)-R_W(P_0)\|.
\end{equation*}
The structure lemma's Lipschitz bound for $r_0$, with $\|X_1-X_0\|=\|\delta X\|$,
$\|X_0\|<(\eta/\ell)r_6$ and $\|Y_1-Y_0\|=\|\delta Y\|$, gives
\begin{equation*}
\|r_0(Y_1,X_1)-r_0(Y_0,X_0)\|
\le\mathfrak{b}_1z'(\eta/\ell)^2r_7+\mathfrak{b}_1(\eta/\ell)r_6\cdot(\eta/\ell)\bar p_1z'.
\end{equation*}
Its Lipschitz bound for $R_W$, with
$\|W_1\|=\|\eta(\nabla^U_\mu\mathsf{Y})(x+\eta\nu)\|\le(\eta/\ell)\bar p_1z'$,
$\|Y_0\|+\|X_0\|\le z'+(\eta/\ell)r_6$ and $\|W_1-W_0\|=\|\delta W\|$, gives
\begin{align*}
\|R_W(P_1)-R_W(P_0)\|
&\le\mathfrak{b}_3\Bigl[(\eta/\ell)\bar p_1z'\bigl((\eta/\ell)\bar p_1z'+(\eta/\ell)^2r_7\bigr)\\
&\qquad+\bigl(z'+(\eta/\ell)r_6\bigr)(\eta/\ell)^2\hat p_2z'\Bigr].
\end{align*}
Multiplying by $(\ell/\eta)^2$ and using $\eta\le\ell$,
\begin{align*}
(\ell^2/\eta)\|\nabla^U_\nu r_\mu\|
&=(\ell/\eta)^2\|r(P_1)-r(P_0)\|\\
&\le\mathfrak{b}_1z'(r_7+\bar p_1r_6)
+\mathfrak{b}_3\bigl[\bar p_1z'(\bar p_1z'+r_7)+\hat p_2z'(z'+r_6)\bigr].
\end{align*}
So entry $7$ increases by at most the bound $\Delta_7$ of
\eqref{4.22:eq-frame-proof-increments-a}, with $p_1$ and $p_2$ replaced by $\bar p_1$ and
$\hat p_2$.

Divide by $z'r_7/\alpha_1=z'\bar r_7$. The first term gives $\hat p_2/\bar r_7\le\hat p_2$.
The remaining terms give:
\begin{itemize}
\item $\mathfrak{b}_1\alpha_1$ and $\mathfrak{b}_1\bar p_1\alpha_1\bar r_6/\bar r_7$;
\item $\mathfrak{b}_3\bar p_1^2z'/\bar r_7\le\mathfrak{b}_3\bar p_1\varpi_{\mathrm{fr}}\alpha_1$,
since $C_{\mathrm{fr}}\bar r_7\ge C_{\mathrm{fr}}\ge\bar p_1$;
\item $\mathfrak{b}_3\bar p_1\alpha_1$;
\item $\mathfrak{b}_3\hat p_2z'/\bar r_7\le4\mathfrak{b}_3\varpi_{\mathrm{fr}}\alpha_1$,
since $\hat p_2\le\bar p_2$, $\bar p_2/C_{\mathrm{fr}}=(4C_\chi+2)/(C_\chi+4)$ and
$(4C_\chi+2)/(C_\chi+4)\le4$;
\item $\mathfrak{b}_3\hat p_2\alpha_1\bar r_6/\bar r_7\le
\mathfrak{b}_3\bar p_2\alpha_1\bar r_6/\bar r_7$.
\end{itemize}
These sum to at most $\mathfrak{c}_7\alpha_1\le1$. Hence entry $7$ increases by at most
\begin{equation*}
(\hat p_2+1)z'r_7/\alpha_1\le C_{\mathrm{fr}}z'r_7/\alpha_{1,k_W}.
\end{equation*}

\emph{Locations off $\mathsf{B}_W$.} If the stencil of a location misses $\mathsf{B}_W$, then
$\mathsf{Y}=0$ on it. Hence $Y=W=0$ there, $r(0,X,0)=0$ and $\rho=1$, and nothing moves.

\emph{Cost.} Since $\varepsilon^{\mathrm{fr}}(W)=\varpi_{\mathrm{fr}}\alpha_*/(4C_{\mathrm{fr}})$,
we have $C_{\mathrm{fr}}/\alpha_*=\varpi_{\mathrm{fr}}/(4\varepsilon^{\mathrm{fr}}(W))$. At
entries $6$ and $7$, using $\alpha_*\le\alpha_1$,
\begin{equation*}
C_{\mathrm{fr}}z'r_a/\alpha_1
=\frac{z'}{\varepsilon^{\mathrm{fr}}(W)}\cdot\frac{\varpi_{\mathrm{fr}}}{4}
\cdot\frac{\alpha_*}{\alpha_1}\,r_a
\le\frac{z'}{\varepsilon^{\mathrm{fr}}(W)}\cdot\frac{\varpi_{\mathrm{fr}}}{4}\,r_a .
\end{equation*}
At entries $1$--$4$ we have $2.1\alpha_*\le C_{\mathrm{fr}}$, because $\alpha_*\le1$ and
$C_{\mathrm{fr}}\ge5$. Hence
\begin{equation*}
2.1z'r_a\le C_{\mathrm{fr}}z'r_a/\alpha_*
=(z'/\varepsilon^{\mathrm{fr}}(W))\cdot\tfrac14\varpi_{\mathrm{fr}}r_a .
\end{equation*}
This proves (a).

\emph{Proof of (b).} We argue by induction on $j$. For $j=1$ the induction hypothesis below
holds: $P^{(0)}=P$ is decomposed with $U$, its cost so far is zero, and it lies in
$\mathfrak{U}^{\mathrm{tr}}_N(X)$ because the reserve $\varpi_{\mathrm{fr}}r_a$ of the stored
domain $\mathfrak{U}_N(X)$ absorbs a cost zero. Suppose that
$P^{(j-1)}\in\mathfrak{U}^{\mathrm{tr}}_N(X)$ is decomposed with the same $U$, and that its
cost so far is at most
$(\sum_{i<j}z_i/\varepsilon^{\mathrm{fr}}(W))\tfrac14\varpi_{\mathrm{fr}}r_a<\varpi_{\mathrm{fr}}r_a$.
Its own $A'$ obeys the thresholds $r_6$ and $r_7$, because $P^{(j-1)}$ lies in the tube. The
profile of $\mathsf{Y}_j$ is measured with this same $U$, and
$z_j\le\sum_iz_i<z_{\max}$. So (a) applies to $P^{(j-1)}$ and $\mathsf{Y}_j$.

The increments of entries $6$ and $7$ add. Entries $1$--$4$ compose multiplicatively: by
$(e^{2\sigma}-1)+e^{2\sigma}(e^{2\sigma'}-1)=e^{2(\sigma+\sigma')}-1$ for
$\sigma,\sigma'\ge0$, their total increment is at most
$(e^{2\sum z_i}-1)\theta_a$. Since $\sum_iz_i<10^{-3}$, this is at most
\begin{equation*}
2.1\Bigl(\sum_i z_i\Bigr)r_a
\le\Bigl(\sum_iz_i/\varepsilon^{\mathrm{fr}}(W)\Bigr)\tfrac14\varpi_{\mathrm{fr}}r_a .
\end{equation*}
Entries $5$ and $8$ never move. This gives \eqref{4.22:eq-profile-class-1}. By the definition
of the stored domain $\mathfrak{U}_N(X)$, whose reserve at every entry is
$\varpi_{\mathrm{fr}}r_a$, the total is absorbed. Hence $P^{(j)}\in\mathfrak{U}^{\mathrm{tr}}_N(X)$,
decomposed with the same $(U,\{u_{\square}\})$.

Since $\mathsf{Y}_{\mathsf{r}}[Z_j]\in\mathfrak{P}(\max_s\|(Z_j)_s\|)$, the multiple-map statement
\eqref{4.22:eq-frame-proof-increments-c} is the case $\mathsf{Y}_j=\mathsf{Y}_{\mathsf{r}}[Z_j]$
of (b).
\end{proof}

\begin{lemma}[Profile of the logarithm of a product of exponentials]\label{4.22:lem-profile-logarithm}
The conventions of the proof of the frame lemma (\ref{4.22:prf-frame-proof-tent}) are in force.
\begin{itemize}
\item $T_\eta$ is the fine lattice of spacing $\eta$, and $k_W\ge 1$ is the level of the frame $W$.
\item $\ell=L^{k_W}\eta$, so that $\eta\le\ell$.
\item $\mathfrak{A}:=\mathfrak{gl}(N,\mathbb{C})$ carries the operator norm $\|\cdot\|$.
\end{itemize}
Let $\mathsf{r}=(U,\{u_{\square}\})$ be a decomposition; its unitary part $U$ is $G$-valued. For an
$\mathfrak{A}$-valued site function $f$ on $T_\eta$ the covariant difference is
\begin{equation*}
(\nabla^U_\mu f)(x):=\eta^{-1}\bigl[\operatorname{Ad}_{U(x,\mu)}f(x+\eta\mu)-f(x)\bigr].
\end{equation*}
Let $\rho_1:=1/10$, the same radius as in the treatment of the logarithm of a triple product in the
proof of the frame lemma (\ref{4.22:prf-frame-proof-logarithm}). For $\vec Y=(Y_1,\dots,Y_n)\in\mathfrak{A}^n$
write $\|\vec Y\|_1:=\sum_i\|Y_i\|$. For $n\ge1$ put
\begin{equation}\label{4.22:eq-profile-logarithm-1}
\mathcal{D}_n:=\bigl\{\vec Y\in\mathfrak{A}^n:\ \|\vec Y\|_1<\rho_1\bigr\},\qquad
\Xi(\vec Y):=\log\Bigl(\prod_{i=1}^{n}e^{Y_i}\Bigr),
\end{equation}
where the product is ordered with $i$ increasing from left to right, and
$\log(1+M):=\sum_{k\ge1}(-1)^{k+1}M^k/k$. Put further
\begin{equation*}
\Upsilon(\vec Y):=\Xi(\vec Y)-\sum_{i=1}^nY_i,\qquad
M_\Upsilon:=\sup_{\mathcal{D}_n}\|\Upsilon\|,\qquad \mathfrak{m}:=e^{1/10}-1 .
\end{equation*}
Thus $\Upsilon$ is the non-linear part of $\Xi$. The following four numerical bounds hold, the last
two, for the quantities $z_+$ and $\delta$ of (d), under the hypotheses of (e):
\begin{equation}\label{4.22:eq-profile-logarithm-a}
\mathfrak{m}<0.1052,\qquad M_\Upsilon<0.212,\qquad
\bar p_1z_+<\frac{1}{30000},\qquad \delta<0.03 .
\end{equation}
Then the following hold.
\begin{enumerate}
\item[(a)] On $\mathcal{D}_n$ we have $\|\prod_ie^{Y_i}-1\|\le e^{\|\vec Y\|_1}-1<\mathfrak{m}$, and
$\mathfrak{m}$ obeys the first inequality of \eqref{4.22:eq-profile-logarithm-a}. The function $\Xi$
is holomorphic on $\mathcal{D}_n$ and $e^{\Xi(\vec Y)}=\prod_ie^{Y_i}$. Moreover
$\|\Xi\|<0.1114$, and $M_\Upsilon$ obeys the second inequality of
\eqref{4.22:eq-profile-logarithm-a}. For $g\in U(N)$,
$\Xi((\operatorname{Ad}_gY_i)_i)=\operatorname{Ad}_g\Xi(\vec Y)$, and the same holds for $\Upsilon$.
If every $Y_i$ is traceless, then $\operatorname{tr}\Xi(\vec Y)=0$. If $Y_j=0$ for all $j\ne i$, then
$\Xi(\vec Y)=Y_i$ and $\Upsilon(\vec Y)=0$; in particular $\Upsilon(0)=0$. For every
$\hat{\vec Y}\in\mathfrak{A}^n$ the function $\zeta\mapsto\Upsilon(\zeta\hat{\vec Y})$ vanishes to
second order at $\zeta=0$; in particular $\Upsilon(0)=0$ and $D\Upsilon(0)=0$.
\item[(b)] Let $\vec Y,\vec W,\vec W'\in\mathfrak{A}^n$, and put $a:=\|\vec Y\|_1$ and
$b:=\|\vec W\|_1$.
\begin{itemize}
\item[(b1)] If $a<\rho_1$, then $\|\Upsilon(\vec Y)\|\le 22a^2$.
\item[(b2)] If $a,b<\tfrac12\rho_1$, then
$\|\Upsilon(\vec Y+\vec W)-\Upsilon(\vec Y)\|\le 170\,b\max\{a,b\}$.
\item[(b3)] Suppose $a\le10^{-3}$ and $\|\vec W\|_1=\|\vec W'\|_1=1$. Let $\sigma,\tau\in\mathbb{C}$
with $|\sigma|,|\tau|<\tfrac12(\rho_1-a)$. Then
\begin{equation*}
\bigl\|\Upsilon(\vec Y+\sigma\vec W+\tau\vec W')-\Upsilon(\vec Y+\sigma\vec W)
-\Upsilon(\vec Y+\tau\vec W')+\Upsilon(\vec Y)\bigr\|\le 350|\sigma||\tau| .
\end{equation*}
\end{itemize}
\item[(c)] Let $\mathsf{Y}_1,\dots,\mathsf{Y}_n$ be $\mathfrak{A}$-valued site functions with
$\sum_i\sup_x\|\mathsf{Y}_i(x)\|<\rho_1$, and let $\mathcal{L}(x):=\Xi((\mathsf{Y}_i(x))_i)$.
Fix $x,\mu,\nu$ and put
\begin{equation*}
\vec Y:=(\mathsf{Y}_i(x))_i,\quad \vec W_\mu:=((\nabla^U_\mu\mathsf{Y}_i)(x))_i,\quad
\vec W_\nu:=((\nabla^U_\nu\mathsf{Y}_i)(x))_i,\quad
\vec V:=((\nabla^U_\nu\nabla^U_\mu\mathsf{Y}_i)(x))_i .
\end{equation*}
By the definition of $\nabla^U$ these satisfy
\begin{equation*}
\operatorname{Ad}_{U(x,\mu)}\mathsf{Y}_i(x+\eta\mu)=Y_i+\eta W_{\mu,i},\qquad
\operatorname{Ad}_{U(x,\nu)}(\nabla^U_\mu\mathsf{Y}_i)(x+\eta\nu)=W_{\mu,i}+\eta V_i .
\end{equation*}
Then
\begin{equation}\label{4.22:eq-profile-logarithm-2}
(\nabla^U_\mu\mathcal{L})(x)=\sum_iW_{\mu,i}
+\eta^{-1}\bigl[\Upsilon(\vec Y+\eta\vec W_\mu)-\Upsilon(\vec Y)\bigr]
\end{equation}
and
\begin{equation}\label{4.22:eq-profile-logarithm-3}
\begin{split}
(\nabla^U_\nu\nabla^U_\mu\mathcal{L})(x)=\sum_iV_i
+\eta^{-2}\bigl[&\Upsilon(\vec Y+\eta\vec W_\nu+\eta\vec W_\mu+\eta^2\vec V)
-\Upsilon(\vec Y+\eta\vec W_\nu)\\
&-\Upsilon(\vec Y+\eta\vec W_\mu)+\Upsilon(\vec Y)\bigr].
\end{split}
\end{equation}
The four arguments of $\Upsilon$ in \eqref{4.22:eq-profile-logarithm-3} are images, under
$\operatorname{Ad}$ of unitaries, of the tuples of values of the $\mathsf{Y}_i$ at
$x+\eta\mu+\eta\nu$, $x+\eta\nu$, $x+\eta\mu$ and $x$. Hence they lie in $\mathcal{D}_n$. Put
$\tilde{\vec Y}:=\vec Y+\eta\vec W_\nu+\eta\vec W_\mu$. If moreover
$\tilde{\vec Y}\in\mathcal{D}_n$, the bracket in \eqref{4.22:eq-profile-logarithm-3} equals
\begin{equation}\label{4.22:eq-profile-logarithm-4}
\bigl[\Upsilon(\tilde{\vec Y}+\eta^2\vec V)-\Upsilon(\tilde{\vec Y})\bigr]
+\bigl[\Upsilon(\tilde{\vec Y})-\Upsilon(\vec Y+\eta\vec W_\nu)
-\Upsilon(\vec Y+\eta\vec W_\mu)+\Upsilon(\vec Y)\bigr].
\end{equation}
\item[(d)] Put $\bar p_1:=3^d(C_\chi+1)$ and $\hat p_2:=3^d(C_\chi+2)$. Let
$\mathsf{Y}_i\in\mathfrak{P}(z_i)$ for $i=1,\dots,n$. Here $\mathfrak{P}(z')$ is the profile class
of Corollary~\ref{4.22:cor-profile-class} (see \eqref{4.22:eq-profile-class-a}), and all the
profiles are measured with the same $U$. Let $z_+:=\sum_iz_i$ with $\bar p_1z_+\le10^{-3}$. Then:
\begin{itemize}
\item the sitewise function $\mathcal{L}:=\Xi((\mathsf{Y}_i)_i)$ is defined on $T_\eta$ and is
$\mathfrak{sl}(N,\mathbb{C})$-valued;
\item $e^{\mathcal{L}}=\prod_ie^{\mathsf{Y}_i}$ sitewise;
\item $\mathcal{L}(x)=0$ wherever every $\mathsf{Y}_i(x)=0$, so in particular off $\mathsf{B}_W$;
\item $\mathcal{L}(x)=0$ wherever $\prod_ie^{\mathsf{Y}_i(x)}=1$.
\end{itemize}
Moreover $\mathcal{L}\in\mathfrak{P}((1+\delta)z_+)$ with $\delta:=860\,\bar p_1z_+$. Precisely,
\begin{equation}\label{4.22:eq-profile-logarithm-5}
\begin{aligned}
\sup_x\|\mathcal{L}(x)\|&\le z_+(1+22z_+),\\
\sup_{x,\mu}\ell\|(\nabla^U_\mu\mathcal{L})(x)\|&\le \bar p_1z_+(1+170\,\bar p_1z_+),\\
\sup_{x,\mu,\nu}\ell^2\|(\nabla^U_\nu\nabla^U_\mu\mathcal{L})(x)\|
&\le \hat p_2z_+(1+860\,\bar p_1z_+).
\end{aligned}
\end{equation}
\item[(e)] Assume the ceiling \eqref{4.22:eq-coupling-ceilings-b} for all $x\ge\gamma_H^{-2}$,
where $x_{k_W}>\gamma_H^{-2}$, and assume $z_+\le 2.1\,\varepsilon^{\mathrm{fr}}(W)$. Then
$\bar p_1z_+<1/30000$. Hence (d) applies with
$\delta<0.03$, and $\mathcal{L}\in\mathfrak{P}(1.1\,z_+)$. This is the factor $1.1$ entering
\eqref{4.22:eq-frame-proof-datum-a}.
\end{enumerate}
The numbers $\rho_1$, $\mathfrak{m}$, $M_\Upsilon$,
$22$, $170$, $350$ and $860$ are absolute, and every bound of the lemma is independent of $N$.
\end{lemma}

\begin{proof}
\emph{(a).} Write $1+\mathsf{a}_i:=e^{Y_i}$.

\emph{The bound on the product.} By the estimate $\|e^B-1\|\le e^{\|B\|}-1$ of the exponential
estimates in the proof of the frame lemma (\ref{4.22:prf-frame-proof-tent}),
$\|\mathsf{a}_i\|\le e^{\|Y_i\|}-1$. Expand
$\prod_i(1+\mathsf{a}_i)-1$ as the sum, over the non-empty subsets of $\{1,\dots,n\}$, of the
ordered products of the corresponding $\mathsf{a}_i$. This gives
\begin{equation*}
\Bigl\|\prod_i(1+\mathsf{a}_i)-1\Bigr\|\le\prod_i(1+\|\mathsf{a}_i\|)-1
\le e^{\|\vec Y\|_1}-1<e^{\rho_1}-1=\mathfrak{m}.
\end{equation*}
For the first inequality of \eqref{4.22:eq-profile-logarithm-a}, write
\begin{equation*}
\mathfrak{m}-\tfrac1{10}=\sum_{k\ge2}\frac{10^{-k}}{k!}
=\frac1{200}+\frac1{6000}+\sum_{k\ge4}\frac{10^{-k}}{k!}.
\end{equation*}
Here $1/6000<0.00017$. Moreover
$\sum_{k\ge4}10^{-k}/k!\le(10^{-4}/24)\sum_{j\ge0}10^{-j}<10^{-5}$. Hence
\begin{equation*}
\mathfrak{m}<0.1+0.005+0.00017+0.00001<0.1052<1 .
\end{equation*}

\emph{Holomorphy.} Put $M:=\prod_ie^{Y_i}-1$. It is an entire function of the matrix entries of
the $Y_i$. On a compact subset of $\mathcal{D}_n$ we have $\|\vec Y\|_1\le r$ for some $r<\rho_1$.
There $\|M\|\le e^{r}-1$, a number less than $1$. So the logarithm series converges uniformly on
that subset. Therefore $\Xi$, a uniform limit of holomorphic functions, is holomorphic on
$\mathcal{D}_n$.

\emph{The bound on $\Xi$.} We have $\|\Xi(\vec Y)\|\le\sum_k\|M\|^k/k=-\log(1-\|M\|)$. Since
$\sum_{k\ge2}\mathfrak{m}^k/k\le\frac12\sum_{k\ge2}\mathfrak{m}^k$, this gives
\begin{equation*}
\|\Xi(\vec Y)\|\le\mathfrak{m}+\frac{\mathfrak{m}^2}{2(1-\mathfrak{m})}
<0.1052+\frac{0.01107}{1.7896}<0.1052+0.0062=0.1114 .
\end{equation*}
Here we used $0.1052^2<0.01107$ and $1.7896\cdot0.0062>0.01107$. Hence
\begin{equation*}
M_\Upsilon\le\sup_{\mathcal{D}_n}\bigl(\|\Xi\|+\|\vec Y\|_1\bigr)\le0.1114+0.1<0.212,
\end{equation*}
which is the second inequality of \eqref{4.22:eq-profile-logarithm-a}.

\emph{The identity $e^{\Xi}=\prod_ie^{Y_i}$.} We have
$\sum_k(1/k!)\bigl(-\log(1-\|M\|)\bigr)^k=(1-\|M\|)^{-1}<\infty$. So the composed series
$\exp(\log(1+M))$ converges absolutely. The identity $\exp\circ\log=\mathrm{id}$ of formal power
series in one commuting variable may therefore be rearranged, as in the treatment of the triple
product (\ref{4.22:prf-frame-proof-logarithm}). This gives $e^{\Xi}=1+M=\prod_ie^{Y_i}$.

\emph{Equivariance.} The exponential, finite products and the logarithm series all commute with
$M\mapsto gMg^{-1}$. For $g\in U(N)$ we have $\|\operatorname{Ad}_gY_i\|=\|Y_i\|$, so
$\operatorname{Ad}_g$ maps $\mathcal{D}_n$ onto itself. This proves the identity for $\Xi$. Since
$\operatorname{Ad}_g$ is linear, the identity for $\Upsilon$ follows.

\emph{Trace.} The set $\mathcal{D}_n$ is star-shaped, because $\|\tau\vec Y\|_1=\tau\|\vec Y\|_1$ for
$\tau\in[0,1]$. The function $\tau\mapsto\operatorname{tr}\Xi(\tau\vec Y)$ is therefore continuous
on $[0,1]$. When every $Y_i$ is traceless,
\begin{equation*}
e^{\operatorname{tr}\Xi(\tau\vec Y)}=\det e^{\Xi(\tau\vec Y)}=\prod_i\det e^{\tau Y_i}
=\prod_ie^{\tau\operatorname{tr}Y_i}=1 .
\end{equation*}
So $\operatorname{tr}\Xi(\tau\vec Y)\in2\pi i\mathbb{Z}$. At $\tau=0$ we have
$\Xi(0)=\log1=0$, so the function vanishes there. Being continuous, it vanishes identically, and
$\operatorname{tr}\Xi(\vec Y)=0$. Continuity decides here, not smallness: the bound
$|\operatorname{tr}\Xi|\le N\|\Xi\|$ may exceed $2\pi$.

\emph{One factor.} Suppose $Y_j=0$ for $j\ne i$. Then $\Xi(\vec Y)=\log(e^{Y_i})$, and
$\|Y_i\|<\rho_1<\log2$. For $\|B\|<\log2$ we have $\|e^B-1\|\le e^{\|B\|}-1<1$ by the same
estimate, and
$\sum_m(1/m)(e^{\|B\|}-1)^m<\infty$. So $\log(e^{B})$ converges absolutely, and the formal identity
$\log\circ\exp=\mathrm{id}$ may be rearranged. This gives $\log(e^B)=B$. Hence $\Xi(\vec Y)=Y_i$ and
$\Upsilon(\vec Y)=0$.

\emph{Second order.} Fix $\hat{\vec Y}$ and put $S:=\sum_i\hat Y_i$. Each factor can be written
\begin{equation*}
e^{\zeta\hat Y_i}=1+\zeta\hat Y_i+\zeta^2(\text{a function holomorphic in }\zeta).
\end{equation*}
Multiplying out, for $\zeta$ near $0$ we get
\begin{equation*}
\prod_ie^{\zeta\hat Y_i}=1+\zeta S+\zeta^2B(\zeta),
\end{equation*}
with $B$ holomorphic. Thus $M=\zeta(S+\zeta B(\zeta))$. In the series
$\sum_{k\ge1}(-1)^{k+1}M^k/k$, every term with $k\ge2$ carries the factor $\zeta^2$. Hence
$\Xi(\zeta\hat{\vec Y})=\zeta S+\zeta^2(\text{holomorphic})$, and
\begin{equation*}
\Upsilon(\zeta\hat{\vec Y})=\Xi(\zeta\hat{\vec Y})-\zeta S=\zeta^2(\text{holomorphic}).
\end{equation*}
Therefore $\Upsilon(0)=0$, and $D\Upsilon(0)[\hat{\vec Y}]$, the derivative of
$\Upsilon(\zeta\hat{\vec Y})$ at $\zeta=0$, vanishes for every direction $\hat{\vec Y}$.

\emph{(b).} Throughout, $\varphi$ ranges over the norm-one linear functionals on $\mathfrak{A}$, and
$\|M\|=\sup_\varphi|\varphi(M)|$.

\emph{(b1).} If $a=0$, then $\vec Y=0$ and $\Upsilon(0)=0$. Otherwise put
$\psi_1(\zeta):=\varphi(\Upsilon(\zeta\vec Y/a))$.
\begin{itemize}
\item $\psi_1$ is holomorphic on $|\zeta|<\rho_1$, because $\|\zeta\vec Y/a\|_1=|\zeta|$.
\item $|\psi_1|\le M_\Upsilon$.
\item $\psi_1(0)=\psi_1'(0)=0$ by (a).
\end{itemize}
Hence $\psi_1=\zeta^2\tilde\psi_1$ with $\tilde\psi_1$ holomorphic. The maximum principle on the
circles $|\zeta|=\rho'$, letting $\rho'\uparrow\rho_1$, gives $|\tilde\psi_1|\le M_\Upsilon/\rho_1^2$
on the disc. At $\zeta=a$,
\begin{equation*}
|\varphi(\Upsilon(\vec Y))|\le M_\Upsilon a^2/\rho_1^2\le0.212\cdot100\cdot a^2<22a^2 .
\end{equation*}
Take the supremum over $\varphi$.

\emph{(b2).} If $b=0$ there is nothing to prove. If $a=0$, then (b1), applicable since
$b<\rho_1$, gives
\begin{equation*}
\|\Upsilon(\vec W)-\Upsilon(0)\|=\|\Upsilon(\vec W)\|\le22b^2\le170\,b\max\{a,b\}.
\end{equation*}
Otherwise put
\begin{equation*}
\psi_2(\sigma,\tau):=\varphi\bigl(\Upsilon(\sigma\vec Y/a+\tau\vec W/b)\bigr)
-\varphi\bigl(\Upsilon(\sigma\vec Y/a)\bigr)
\end{equation*}
on the bidisc $|\sigma|,|\tau|<\tfrac12\rho_1$. There
$\|\sigma\vec Y/a+\tau\vec W/b\|_1\le|\sigma|+|\tau|<\rho_1$ and $\|\sigma\vec Y/a\|_1<\rho_1$, so
both arguments lie in $\mathcal{D}_n$. Thus $\psi_2$ is holomorphic, $|\psi_2|\le2M_\Upsilon$, and
$\psi_2(\sigma,0)=0$.

Hence $\psi_2=\tau\tilde\psi_2$, where
$\tilde\psi_2(\sigma,\tau)=\int_0^1\partial_\tau\psi_2(\sigma,s\tau)\,ds$ is holomorphic on the
bidisc. The maximum principle in $\tau$ at fixed $\sigma$, on the circles $|\tau|=\rho'$ with
$\rho'\uparrow\frac12\rho_1$, gives
\begin{equation*}
|\tilde\psi_2|\le2M_\Upsilon/(\tfrac12\rho_1)=4M_\Upsilon/\rho_1 .
\end{equation*}
Moreover $\tilde\psi_2(0,0)=\lim_{\tau\to0}\varphi(\Upsilon(\tau\vec W/b))/\tau=0$ by (a).

Put $\hat m:=\max\{a,b\}<\frac12\rho_1$. The function
$\zeta\mapsto\tilde\psi_2(\zeta a/\hat m,\zeta b/\hat m)$ is holomorphic on $|\zeta|<\frac12\rho_1$,
because $|\zeta a/\hat m|,|\zeta b/\hat m|\le|\zeta|$. It is bounded by $4M_\Upsilon/\rho_1$ and
vanishes at $\zeta=0$. Schwarz's lemma gives
\begin{equation*}
|\tilde\psi_2(\zeta a/\hat m,\zeta b/\hat m)|\le\frac{4M_\Upsilon}{\rho_1}
\cdot\frac{|\zeta|}{\frac12\rho_1}=\frac{8M_\Upsilon|\zeta|}{\rho_1^2}.
\end{equation*}
At $\zeta=\hat m$, which is real and less than $\frac12\rho_1$,
\begin{equation*}
\bigl|\varphi\bigl(\Upsilon(\vec Y+\vec W)-\Upsilon(\vec Y)\bigr)\bigr|=|\psi_2(a,b)|
=b|\tilde\psi_2(a,b)|\le\frac{8M_\Upsilon b\hat m}{\rho_1^2}
\le8\cdot0.212\cdot100\cdot b\hat m<170\,b\max\{a,b\}.
\end{equation*}
Take the supremum over $\varphi$.

\emph{(b3).} Put
\begin{equation*}
\psi_3(\sigma,\tau):=\varphi\bigl(\Upsilon(\vec Y+\sigma\vec W+\tau\vec W')
-\Upsilon(\vec Y+\sigma\vec W)-\Upsilon(\vec Y+\tau\vec W')+\Upsilon(\vec Y)\bigr).
\end{equation*}
On the bidisc $|\sigma|,|\tau|<\frac12(\rho_1-a)$, every argument has
$\|\cdot\|_1\le a+|\sigma|+|\tau|<\rho_1$. So $\psi_3$ is holomorphic there,
$|\psi_3|\le4M_\Upsilon$, and $\psi_3(\sigma,0)=0=\psi_3(0,\tau)$ identically.

First factor out $\tau$. We get $\psi_3=\tau\tilde\psi_3$ with $\tilde\psi_3$ holomorphic. The
maximum principle in $\tau$ at fixed $\sigma$ gives
$|\tilde\psi_3|\le4M_\Upsilon/(\frac12(\rho_1-a))$. For $\tau\ne0$,
$\tilde\psi_3(0,\tau)=\psi_3(0,\tau)/\tau=0$, and at $\tau=0$ this holds by continuity.

Next factor out $\sigma$. We get $\tilde\psi_3=\sigma\hat\psi_3$ with $\hat\psi_3$ holomorphic. The
maximum principle in $\sigma$ at fixed $\tau$ gives
$|\hat\psi_3|\le4M_\Upsilon/(\frac12(\rho_1-a))^2$. Since $a\le10^{-3}$, we have
$\rho_1-a\ge0.099$, and so
\begin{equation*}
|\psi_3(\sigma,\tau)|\le\frac{16M_\Upsilon|\sigma||\tau|}{(\rho_1-a)^2}
\le\frac{16\cdot0.212}{(0.099)^2}|\sigma||\tau|<350|\sigma||\tau| .
\end{equation*}
The last step uses $16\cdot0.212=3.392<3.43=350\cdot0.0098$ and $0.099^2=0.009801>0.0098$. Take
the supremum over $\varphi$.

\emph{(c).} Every tuple of values of the $\mathsf{Y}_i$ at one site has $\|\cdot\|_1<\rho_1$. The
values $U(x,\mu)$ lie in $G\subset U(N)$, and $\operatorname{Ad}$ of a unitary preserves the norm.
So transported tuples lie in $\mathcal{D}_n$ as well.

\emph{First identity.} By the definition of $\nabla^U$, the equivariance in (a) under $U(x,\mu)$,
and the first relation of (c),
\begin{equation*}
\begin{split}
(\nabla^U_\mu\mathcal{L})(x)
&=\eta^{-1}\bigl[\operatorname{Ad}_{U(x,\mu)}\Xi((\mathsf{Y}_i(x+\eta\mu))_i)-\Xi(\vec Y)\bigr]\\
&=\eta^{-1}\bigl[\Xi((\operatorname{Ad}_{U(x,\mu)}\mathsf{Y}_i(x+\eta\mu))_i)-\Xi(\vec Y)\bigr]
=\eta^{-1}\bigl[\Xi(\vec Y+\eta\vec W_\mu)-\Xi(\vec Y)\bigr].
\end{split}
\end{equation*}
Now write $\Xi=\sum_i(\cdot)+\Upsilon$. The sum is linear, and its contribution is
$\eta^{-1}\sum_i\eta W_{\mu,i}$. This gives \eqref{4.22:eq-profile-logarithm-2}.

\emph{Second identity.} Apply $\nabla^U_\nu$ to the site function $\nabla^U_\mu\mathcal{L}$:
\begin{equation*}
(\nabla^U_\nu\nabla^U_\mu\mathcal{L})(x)
=\eta^{-1}\bigl[\operatorname{Ad}_{U(x,\nu)}(\nabla^U_\mu\mathcal{L})(x+\eta\nu)
-(\nabla^U_\mu\mathcal{L})(x)\bigr].
\end{equation*}
Use the first identity at the site $x+\eta\nu$, in its form with $\Xi$, and then apply
$\operatorname{Ad}_{U(x,\nu)}$, using equivariance and linearity again. The quantity
$\operatorname{Ad}_{U(x,\nu)}(\nabla^U_\mu\mathcal{L})(x+\eta\nu)$ equals
\begin{equation*}
\eta^{-1}\Bigl[\Xi\bigl((\operatorname{Ad}_{U(x,\nu)}
[\mathsf{Y}_i(x+\eta\nu)+\eta(\nabla^U_\mu\mathsf{Y}_i)(x+\eta\nu)])_i\bigr)
-\Xi\bigl((\operatorname{Ad}_{U(x,\nu)}\mathsf{Y}_i(x+\eta\nu))_i\bigr)\Bigr].
\end{equation*}
By the two relations of (c) this is the right side of the first line below. Subtracting
$\eta^{-1}[\Xi(\vec Y+\eta\vec W_\mu)-\Xi(\vec Y)]$ and dividing by $\eta$ gives the second line:
\begin{align*}
\operatorname{Ad}_{U(x,\nu)}(\nabla^U_\mu\mathcal{L})(x+\eta\nu)
&=\eta^{-1}\bigl[\Xi(\vec Y+\eta\vec W_\nu+\eta\vec W_\mu+\eta^2\vec V)
-\Xi(\vec Y+\eta\vec W_\nu)\bigr],\\
(\nabla^U_\nu\nabla^U_\mu\mathcal{L})(x)
&=\eta^{-2}\bigl[\Xi(\vec Y+\eta\vec W_\nu+\eta\vec W_\mu+\eta^2\vec V)
-\Xi(\vec Y+\eta\vec W_\nu)\\
&\qquad\qquad-\Xi(\vec Y+\eta\vec W_\mu)+\Xi(\vec Y)\bigr].
\end{align*}
The linear parts contribute $\eta^{-2}\sum_i\eta^2V_i=\sum_iV_i$. The rest is the bracket of
\eqref{4.22:eq-profile-logarithm-3}.

\emph{The arguments.} The four arguments are images under $\operatorname{Ad}$ of unitaries of the
tuples of values at $x+\eta\mu+\eta\nu$, $x+\eta\nu$, $x+\eta\mu$ and $x$. For instance,
\begin{equation*}
\mathsf{Y}_i(x+\eta\nu)+\eta(\nabla^U_\mu\mathsf{Y}_i)(x+\eta\nu)
=\operatorname{Ad}_{U(x+\eta\nu,\mu)}\mathsf{Y}_i(x+\eta\nu+\eta\mu).
\end{equation*}
Therefore
\begin{equation*}
\vec Y+\eta\vec W_\nu+\eta\vec W_\mu+\eta^2\vec V
=\bigl(\operatorname{Ad}_{U(x,\nu)}\operatorname{Ad}_{U(x+\eta\nu,\mu)}
\mathsf{Y}_i(x+\eta\nu+\eta\mu)\bigr)_i .
\end{equation*}
So each argument has $\|\cdot\|_1<\rho_1$. When $\tilde{\vec Y}\in\mathcal{D}_n$, adding and
subtracting $\Upsilon(\tilde{\vec Y})$ gives \eqref{4.22:eq-profile-logarithm-4}.

\emph{(d).} Fix a site $x$ and use the notation of (c). Since $\mathsf{Y}_i\in\mathfrak{P}(z_i)$
and $\eta\le\ell$,
\begin{equation*}
a:=\|\vec Y\|_1\le z_+,\qquad
\eta\|\vec W_\mu\|_1\le(\eta/\ell)\sum_i\bar p_1z_i\le\bar p_1z_+ .
\end{equation*}
Likewise $\eta\|\vec W_\nu\|_1\le\bar p_1z_+$. Since $C_\chi+2\le2(C_\chi+1)$,
\begin{equation*}
\eta^2\|\vec V\|_1\le(\eta/\ell)^2\hat p_2z_+\le\hat p_2z_+\le2\bar p_1z_+ .
\end{equation*}
Since $\bar p_1\ge1$, it follows that
\begin{equation*}
\|\tilde{\vec Y}\|_1\le z_++2\bar p_1z_+\le3\bar p_1z_+\le3\cdot10^{-3}<\tfrac12\rho_1 .
\end{equation*}
So $\tilde{\vec Y}\in\mathcal{D}_n$, and (c) applies in the regrouped form
\eqref{4.22:eq-profile-logarithm-4}.

\emph{Basic properties of $\mathcal{L}$.} We have
$\sum_i\sup\|\mathsf{Y}_i\|\le z_+$, and $z_+\le10^{-3}<\rho_1$, so $\mathcal{L}$ is defined on
$T_\eta$.
Each $\mathsf{Y}_i$ is $\mathfrak{sl}(N,\mathbb{C})$-valued, so $\mathcal{L}$ is traceless by (a).
By (a), $e^{\mathcal{L}}=\prod_ie^{\mathsf{Y}_i}$ sitewise. Where every $\mathsf{Y}_i(x)=0$,
$\mathcal{L}(x)=\Xi(0)=0$. Members of $\mathfrak{P}$ vanish off $\mathsf{B}_W$, so this includes
every site off $\mathsf{B}_W$. Where $\prod_ie^{\mathsf{Y}_i(x)}=1$, we have $M=0$, and the
logarithm series gives $\mathcal{L}(x)=0$.

\emph{Sup bound.} By (b1), $\|\mathcal{L}(x)\|\le\|\vec Y\|_1+\|\Upsilon(\vec Y)\|\le z_++22z_+^2$.

\emph{First difference.} Apply \eqref{4.22:eq-profile-logarithm-2}, together with (b2) at the base
point $\vec Y$ with increment $\eta\vec W_\mu$. Here $a\le z_+<\frac12\rho_1$ and
$b:=\eta\|\vec W_\mu\|_1\le(\eta/\ell)\bar p_1z_+$, so that $b<\frac12\rho_1$. Further,
$(\ell/\eta)b=\ell\|\vec W_\mu\|_1\le\bar p_1z_+$ and $\max\{a,b\}\le\bar p_1z_+$. Hence
\begin{equation*}
\ell\|(\nabla^U_\mu\mathcal{L})(x)\|\le\ell\|\vec W_\mu\|_1+(\ell/\eta)\cdot170\,b\max\{a,b\}
\le\bar p_1z_++170\,\bar p_1z_+\cdot\bar p_1z_+ .
\end{equation*}

\emph{Second difference.} By \eqref{4.22:eq-profile-logarithm-3} and
\eqref{4.22:eq-profile-logarithm-4}, $\ell^2\|(\nabla^U_\nu\nabla^U_\mu\mathcal{L})(x)\|$ is at
most
\begin{equation*}
\ell^2\|\vec V\|_1+(\ell/\eta)^2\bigl\|\Upsilon(\tilde{\vec Y}+\eta^2\vec V)
-\Upsilon(\tilde{\vec Y})\bigr\|
+(\ell/\eta)^2\bigl\|\Upsilon(\tilde{\vec Y})-\Upsilon(\vec Y+\eta\vec W_\nu)
-\Upsilon(\vec Y+\eta\vec W_\mu)+\Upsilon(\vec Y)\bigr\| .
\end{equation*}
We bound the three terms in turn.
\begin{itemize}
\item The first term is at most $\hat p_2z_+$.
\item For the second term apply (b2) at the base point $\tilde{\vec Y}$ with increment
$\eta^2\vec V$, and put $\tilde a:=\|\tilde{\vec Y}\|_1$ and $\tilde b:=\eta^2\|\vec V\|_1$. Here
\begin{equation*}
\tilde a\le3\bar p_1z_+<\tfrac12\rho_1,\qquad \tilde b\le2\bar p_1z_+<\tfrac12\rho_1,\qquad
(\ell/\eta)^2\tilde b=\ell^2\|\vec V\|_1\le\hat p_2z_+ ,
\end{equation*}
so $\max\{\tilde a,\tilde b\}\le3\bar p_1z_+$, and the second term is at most
\begin{equation*}
(\ell/\eta)^2\cdot170\,\tilde b\cdot3\bar p_1z_+\le510\,\bar p_1z_+\cdot\hat p_2z_+ .
\end{equation*}
\item The third term vanishes if $\vec W_\nu=0$ or $\vec W_\mu=0$. Otherwise apply (b3) at the base
point $\vec Y$, where $a\le z_+\le10^{-3}$. Take the unit directions
$\vec W:=\vec W_\nu/\|\vec W_\nu\|_1$ and $\vec W':=\vec W_\mu/\|\vec W_\mu\|_1$. Take the real
numbers $\sigma:=\eta\|\vec W_\nu\|_1$ and $\tau:=\eta\|\vec W_\mu\|_1$. Both are at most
$\bar p_1z_+\le10^{-3}<\frac12(\rho_1-a)$. Then $\vec Y+\sigma\vec W+\tau\vec W'=\tilde{\vec Y}$.
Using $\bar p_1\le\hat p_2$, the third term is at most
\begin{equation*}
(\ell/\eta)^2\cdot350\,\sigma\tau=350\,\ell\|\vec W_\nu\|_1\,\ell\|\vec W_\mu\|_1
\le350(\bar p_1z_+)^2\le350\,\bar p_1z_+\cdot\hat p_2z_+ .
\end{equation*}
\end{itemize}
Summing the three terms gives at most $\hat p_2z_+(1+860\,\bar p_1z_+)$. This proves
\eqref{4.22:eq-profile-logarithm-5}.

\emph{Membership in the class.} Since $\bar p_1\ge1$, we have $22z_+\le170\,\bar p_1z_+\le860\,\bar p_1z_+$.
So the three bounds of \eqref{4.22:eq-profile-logarithm-5} are those defining
$\mathfrak{P}((1+860\,\bar p_1z_+)z_+)$. Also $\mathcal{L}$ is $\mathfrak{sl}(N,\mathbb{C})$-valued
and vanishes off $\mathsf{B}_W$. Hence $\mathcal{L}\in\mathfrak{P}((1+\delta)z_+)$.

\emph{(e).} By hypothesis the ceiling \eqref{4.22:eq-coupling-ceilings-b} holds at
$x=x_{k_W}$, since $x_{k_W}>\gamma_H^{-2}$; that is, $\mathfrak{c}\bar\alpha(x_{k_W})\le1$. The
paragraph on the ceiling in the proof of the increments
(\ref{4.22:prf-frame-proof-increments}) gives
\begin{equation*}
\alpha_{*,k_W}\le\alpha_{1,k_W}\le\tilde\alpha_{1,k_W}\le\bar\alpha(x_{k_W}).
\end{equation*}
The member $\mathfrak{c}_6$ of the ceiling satisfies
\begin{equation*}
\mathfrak{c}\ge\mathfrak{c}_6=\mathfrak{b}_1+\mathfrak{b}_2(\bar p_1+\varpi_{\mathrm{fr}})
\ge\mathfrak{b}_2\bar p_1=200\,\bar p_1 .
\end{equation*}
Hence $200\,\bar p_1\alpha_{*,k_W}\le\mathfrak{c}\bar\alpha(x_{k_W})\le1$, that is,
$\bar p_1\alpha_{*,k_W}\le1/200$.

Next,
\begin{equation*}
2.1\,\varepsilon^{\mathrm{fr}}(W)=\frac{2.1\,\varpi_{\mathrm{fr}}\alpha_{*,k_W}}{4C_{\mathrm{fr}}}
=\frac{21}{40}\,\frac{\varpi_{\mathrm{fr}}\alpha_{*,k_W}}{C_{\mathrm{fr}}}.
\end{equation*}
The range of the frame lemma's clause
(i) (\ref{4.22:prf-frame-proof-tent}) gives $\varpi_{\mathrm{fr}}\le1/80$, and
$C_{\mathrm{fr}}=3^d(C_\chi+4)\ge1$. Therefore
\begin{equation*}
\bar p_1z_+\le\frac{21}{40}\cdot\frac1{80}\cdot\frac1{200}=\frac{21}{640000}
<\frac1{30000}<10^{-3},
\end{equation*}
because $21\cdot3=63<64$. This is the third inequality of \eqref{4.22:eq-profile-logarithm-a}.
Hence (d) applies, with
\begin{equation*}
\delta=860\,\bar p_1z_+<\frac{860}{30000}<0.03,
\end{equation*}
because $86<90$. This is the fourth inequality of \eqref{4.22:eq-profile-logarithm-a}. Finally,
$\mathfrak{P}(z')\subset\mathfrak{P}(z'')$ for $z'\le z''$, so
$\mathcal{L}\in\mathfrak{P}(1.03\,z_+)\subset\mathfrak{P}(1.1\,z_+)$.
\end{proof}

\begin{remark}
The lemma does not sum a Dynkin series. The leading term $\frac12[\mathsf{Y}_1,\mathsf{Y}_2]$ of
that series shows why a coupling-dependent smallness is needed.

By the Leibniz rule for $\nabla^U$, the second differences of this term contain
$[\nabla^U_\nu\mathsf{Y}_1,\nabla^U_\mu\mathsf{Y}_2]$. Its size is at most
$2\bar p_1^2z_1z_2/\ell^2$, to be compared with the allowance $\hat p_2z_+/\ell^2$ of the class.
The relative size is of the order $\bar p_1^2z_+/\hat p_2\le\bar p_1z_+$, which is not small by a
number alone.

In the proof this contribution lies inside the term $350\,\sigma\tau$ of (b3). It is paid for by
$\bar p_1z_+<1/30000$, that is, by $\mathfrak{c}_6\ge200\,\bar p_1$ in
\eqref{4.22:eq-coupling-ceilings-b}, on the range $z_+\le2.1\,\varepsilon^{\mathrm{fr}}(W)$,
which is proportional to $\alpha_{*,k_W}$.

All profiles in the lemma are measured with one and the same $U$. No comparison of $\nabla^U$ with
the covariant difference of a second decomposition is needed here. Independence of the
decomposition is established separately, in the proof of clause (ii) of the frame lemma.
\end{remark}

\begin{remark}
Clause (e) uses only the following:
\begin{itemize}
\item the member $\mathfrak{c}\ge\mathfrak{c}_6\ge200\,\bar p_1$ of the ceiling
\eqref{4.22:eq-coupling-ceilings-b};
\item the chain $\alpha_{*,k_W}\le\alpha_{1,k_W}\le\tilde\alpha_{1,k_W}\le\bar\alpha(x_{k_W})$;
\item $\varpi_{\mathrm{fr}}\le1/80$;
\item $C_{\mathrm{fr}}\ge1$.
\end{itemize}
Clause (e) imposes no condition on the parameters beyond the ceiling, with the constant
$\mathfrak{c}$ as defined there.

The ceiling is of supremum form over $x\ge\gamma_H^{-2}$. Its left member $\mathfrak{c}\bar\alpha(x)$
is of the form (coupling)$\times$(polylogarithm) and tends to $0$. It is therefore a bound on the
coupling from above only, met by the choice of $\gamma$, which comes last in the order of choice.
It is not a cap in the sense of Definition~\ref{2.3:def-cap}.

The numbers $\mathfrak{b}_1=2.05$, $\mathfrak{b}_2=200$ and $\mathfrak{b}_3=400$ are absolute. The
number $\mathfrak{c}$ depends on $d$, $C_\chi$, $\bar r_6$, $\bar r_7$ and $\varpi_{\mathrm{fr}}$
only, where $\bar r_6$ and $\bar r_7$ are the threshold numbers of entries six and seven in units of
$\alpha_{1,k_W}$. It is fixed together with $C_{\mathrm{fr}}$, after the constant $K_R$ of the
order of choice and before $A_0$, $A_1$ and $\gamma$. Clause (e) imposes no condition on the
constants $C_{\mathrm{fr}}=3^d(C_\chi+4)$, $\varepsilon^{\mathrm{fr}}(W)$ and
$\varpi_{\mathrm{fr}}$ beyond their definitions.

Nothing is asked of $L$ beyond $L\ge2$, which gives $\eta/\ell=L^{-k_W}\le\frac12$. Nothing is asked
of $M$, $M_1$, $\kappa_1$, $K$, $k$ or $N$, nor of $N_0$, of the scale $\check R_k$, or of the
ages of large-field regions, and no condition is imposed on $L_0$. Every
bound is independent of $N$: the norm is the operator norm,
$\|\operatorname{Ad}_gM\|=\|M\|$ for $g\in U(N)$, and $\det=e^{\operatorname{tr}}$.
\end{remark}

\begin{proof}[Proof of Lemma~\ref{4.22:lem-frame-lemma}: clauses (iii), (ii), (iv), (v) and (vi)]
\label{4.22:prf-frame-proof-datum}
Clause (i) of the frame lemma has been proved above: the increments of the eight entries under one
gauge map are bounded in the proof of the increments,
\eqref{4.22:eq-frame-proof-increments-c}, and successive maps in the profile class are controlled
by Corollary~\ref{4.22:cor-profile-class}. Clause (iii) is proved next, because the proof of
clause (ii) uses it. Throughout, $N^{\sharp}$ is the continuation of the term $N$ of the frame
lemma, with arguments $(P;\mathbf{A},\Psi,g)$: the pair $P$, the pointwise coordinates
$\mathbf{A}$, the frame data $\Psi$, and the rotations $g$ at the law sites of $N$ other than
$\Sigma_W$. When a rotation $c$ at the roots $\Sigma_W$ is also displayed, the arguments are
$(P;\mathbf{A},\Phi_W,(g,c))$.

\emph{Clause (iii).} The surviving pointwise law sites of $N$ are of class $(\alpha)$ at step
$k_W$: for every surviving pointwise coordinate $a$ its law site $z(a)$ lies outside the
large-field region, by the first clause of the lemma on pointwise law sites. The set
$W\setminus\Lambda_W$ contains at least $10$ layers of
$MR_{k_W}$-cubes \cite{B-CMP122a}, and $\mathsf{B}_W$ lies within distance $\ell$ of
$\partial^{+}\Lambda_W$. The interior variables are members of $\Phi_W$.

\emph{Clause (ii), Step 0: invariance under maps equal to $1$ on $\Sigma_W$.} Let
$Q\in\mathfrak{U}^{\mathrm{tr}}_N(X)$.
Let $t\mapsto u_t$, $t\in[0,1]$, be a $C^1$ path of $G^{\mathbb{C}}$-valued site maps with
$u_0=1$, $u_t\restriction\Sigma_W=1$ and $Q^{u_t}\in\mathfrak{U}^{\mathrm{tr}}_N(X)$ for every $t$.
We claim that for real $\Psi$
\begin{equation}\label{4.22:eq-frame-proof-datum-1}
N^{\sharp}(Q^{u_1};\mathbf{A},\Psi,u_1g)=N^{\sharp}(Q;\mathbf{A},\Psi,g).
\end{equation}
For a $G$-valued map $u$ with $u\restriction\Sigma_W=1$ this identity is clause (ii) of the
invariance theorem for the boundary terms under real gauge transformations, since then
$\Psi^{[u]}=\Psi$, combined with property (O5) of the definition of $N^{\sharp}$. Consequently the
derivative of $N^{\sharp}$ in the direction $\delta_Z$ of the infinitesimal gauge transformation
generated by a $\mathfrak{g}$-valued site function $Z$ vanishing on $\Sigma_W$ is zero. Since
$N^{\sharp}$ is holomorphic, this derivative is $\mathbb{C}$-linear in $Z$. By the argument of
property (O7) of the same definition it therefore vanishes for every
$\mathfrak{g}^{\mathbb{C}}$-valued $Z$ vanishing on $\Sigma_W$.

Now $Q^{u_{t'}}=(Q^{u_t})^{u_{t'}u_t^{-1}}$ and $u_{t'}g=(u_{t'}u_t^{-1})u_tg$. Hence the
derivative of $t\mapsto N^{\sharp}(Q^{u_t};\mathbf{A},\Psi,u_tg)$ at $t$ is the derivative of
$N^{\sharp}$ at $(Q^{u_t};\mathbf{A},\Psi,u_tg)$ in the direction $\delta_Z$ with
$Z=\dot u_tu_t^{-1}$. This $Z$ vanishes on $\Sigma_W$, because $u_t=1$ there. So the function is
constant on $[0,1]$, which is \eqref{4.22:eq-frame-proof-datum-1}.

\emph{Step 1: only the restriction to $\Sigma_W$ matters, at one decomposition.} For a site map $r$
that is a product of maps $\rho_{\mathsf{r}}[Z]$, write $z(r)$ for the sum, over the factors
$\rho_{\mathsf{r}}[Z]$ of $r$, of $\max_s\|Z_s\|$.

Let $\mathsf{r}=(U,\{u_{\square}\})$ be a decomposition of $P\in\mathfrak{U}_N(X)$. Let $r_1,r_2$ be
site maps, each a product of at most two maps $\rho_{\mathsf{r}}[\cdot]$ at this $\mathsf{r}$, with
\[
z(r_1),\,z(r_2)\le 1.05\,\varepsilon^{\mathrm{fr}}(W),\qquad
r_1\restriction\Sigma_W=r_2\restriction\Sigma_W .
\]
Write $r_1=\rho_{\mathsf{r}}[Z_{(2)}]\rho_{\mathsf{r}}[Z_{(1)}]$ and
$r_2=\rho_{\mathsf{r}}[Z'_{(2)}]\rho_{\mathsf{r}}[Z'_{(1)}]$. An absent factor is $1$, that is,
its $Z$ is $0$. Then $u:=r_2r_1^{-1}=e^{\mathsf{Y}_1}e^{\mathsf{Y}_2}e^{\mathsf{Y}_3}e^{\mathsf{Y}_4}$
with
\[
(\mathsf{Y}_1,\mathsf{Y}_2,\mathsf{Y}_3,\mathsf{Y}_4)
:=\bigl(-\mathsf{Y}_{\mathsf{r}}[Z'_{(2)}],\,-\mathsf{Y}_{\mathsf{r}}[Z'_{(1)}],\,
\mathsf{Y}_{\mathsf{r}}[Z_{(1)}],\,\mathsf{Y}_{\mathsf{r}}[Z_{(2)}]\bigr).
\]
By the first and second covariant-difference bounds \eqref{4.22:eq-frame-proof-tent-c}, as restated
in Corollary~\ref{4.22:cor-profile-class}, and since the profile class is even, each
$\mathsf{Y}_i\in\mathfrak{P}(z_i)$, where $z_i$ is the maximum over $s$ of the norms of the
corresponding $Z_s$. Put
\[
z_+:=\sum_iz_i=z(r_1)+z(r_2)\le 2.1\,\varepsilon^{\mathrm{fr}}(W).
\]
By clauses (d) and (e) of the lemma on the profile of the logarithm of a product of exponentials
(Lemma~\ref{4.22:lem-profile-logarithm}, read under the ceiling \eqref{4.22:eq-coupling-ceilings-b}
at $x_{k_W}$), the function
\[
\log u:=\mathcal{L}=\Xi(\mathsf{Y}_1,\mathsf{Y}_2,\mathsf{Y}_3,\mathsf{Y}_4)
\in\mathfrak{P}(1.1\,z_+)\subset\mathfrak{P}\bigl(2.31\,\varepsilon^{\mathrm{fr}}(W)\bigr)
\]
has the following properties:
\begin{itemize}
\item it is $\mathfrak{sl}(N,\mathbb{C})$-valued and $e^{\mathcal{L}}=u$;
\item $\mathcal{L}=0$ off $\mathsf{B}_W$;
\item $\mathcal{L}=0$ on $\Sigma_W$, because there $u(s)=r_2(s)r_1(s)^{-1}=1$, and the logarithm
vanishes wherever the product of the exponentials is $1$.
\end{itemize}

Put $u_t:=e^{t\mathcal{L}}$, $t\in[0,1]$. This is a $C^1$ path of site maps, which are
$G^{\mathbb{C}}=SL(N,\mathbb{C})$-valued because
$\det e^{t\mathcal{L}}=e^{t\operatorname{tr}\mathcal{L}}=1$. It satisfies $u_0=1$, $u_1=u$,
$u_t\restriction\Sigma_W=1$ and $u_t=1$ off $\mathsf{B}_W$. Moreover
$u_t=e^{-(-t\mathcal{L})}$ with $-t\mathcal{L}\in\mathfrak{P}(2.31\,\varepsilon^{\mathrm{fr}}(W))$,
since the class is even, $t\mathcal{L}\in\mathfrak{P}(t\cdot2.31\,\varepsilon^{\mathrm{fr}}(W))$,
and the classes increase with their index.

Apply the successive-maps clause (b) of Corollary~\ref{4.22:cor-profile-class},
\eqref{4.22:eq-profile-class-a}, to $P$ with the three successive maps
$\rho_{\mathsf{r}}[Z_{(1)}]$, $\rho_{\mathsf{r}}[Z_{(2)}]$ and $e^{-(-t\mathcal{L})}$. Their sizes
sum to at most $(1.05+2.31)\,\varepsilon^{\mathrm{fr}}(W)<4\varepsilon^{\mathrm{fr}}(W)$. Since
successive gauge maps compose as $(P^{r})^{u}=P^{ur}$, it follows that
\[
P^{u_tr_1}=(P^{r_1})^{u_t}\in\mathfrak{U}^{\mathrm{tr}}_N(X)\qquad\text{for every }t\in[0,1],
\]
decomposed with the same $(U,\{u_{\square}\})$. At every entry of the tube, with threshold $r_a$,
the total cost is at most
\begin{equation}\label{4.22:eq-frame-proof-datum-a}
\Bigl[1.05+1.1\cdot 2.1\Bigr]\cdot\tfrac14\varpi_{\mathrm{fr}}\,r_a
=0.84\,\varpi_{\mathrm{fr}}\,r_a<\varpi_{\mathrm{fr}}\,r_a ,
\end{equation}
that is, $3.36$ of the $4$ available units $\tfrac14\varpi_{\mathrm{fr}}\,r_a$.

So Step 0 applies to $Q:=P^{r_1}$ and the path $u_t$. Since $ur_1=r_2$, it gives for real $\Psi$
\[
N^{\sharp}(P^{r_2};\mathbf{A},\Psi,ug)=N^{\sharp}((P^{r_1})^{u};\mathbf{A},\Psi,ug)
=N^{\sharp}(P^{r_1};\mathbf{A},\Psi,g).
\]
Moreover $ug=g$. Indeed $g$ lives on the law sites of $N$ other than $\Sigma_W$; by clause (iii)
none of these meets $\mathsf{B}_W$, and off $\mathsf{B}_W$ we have $u=1$. Therefore
\begin{equation}\label{4.22:eq-frame-proof-datum-2}
N^{\sharp}(P^{r_2};\mathbf{A},\Psi,g)=N^{\sharp}(P^{r_1};\mathbf{A},\Psi,g).
\end{equation}

\emph{Step 2: holomorphy.} Every $c\in G^{\mathbb{C}}$ is written in polar form
$c=e^{Y}\mathring c$,
with $\mathring c\in G$ and $\|Y\|=\operatorname{pol}(c)$, as in the lemma on the polar charts
$\mathfrak{G}_\varepsilon$.

Fix $(P_0,c_0)$ with $c_0=e^{Y_0}\mathring c_0$, and fix a decomposition $\mathsf{r}$ of $P_0$. This
decomposition serves for every $P$ near $P_0$, with the same $U$ and with
$U':=\mathbb{U}U^{-1}$, because the conditions defining $\mathfrak{U}^{\mathrm{tr}}_N(X)$ are open,
by the first clause of the openness lemma for the tube. Put
\[
f(P,c):=N^{\sharp}\bigl(P^{\rho_{\mathsf{r}}[\log(c\mathring c_0^{-1})]};\mathbf{A},
\Phi_W^{[\mathring c_0]},g\bigr).
\]
It is holomorphic near $(P_0,c_0)$, as a composition of holomorphic maps. The norm of
$\log(c\mathring c_0^{-1})$ is near $\|Y_0\|<\varepsilon^{\mathrm{fr}}(W)$.

For $c=e^{Y}\mathring c$ near $c_0$ put $v:=\mathring c\,\mathring c_0^{-1}$. This map is
$G$-valued and near $1$.
Put $\hat v:=\rho_{\mathsf{r}}[-\log v]$. Since $\log v(s)\in i\mathfrak{g}$ for $v(s)\in G$ near
$1$, the map $\hat v$ is $G$-valued. Since $\mathsf{Y}_{\mathsf{r}}$ is linear in $Z$, we have
$\hat v^{-1}=\rho_{\mathsf{r}}[\log v]$, which is $G$-valued as well; moreover
$\hat v\restriction\Sigma_W=v$.
Apply clause (ii) of the invariance theorem for the boundary terms under $\hat v^{-1}$:
\[
N^{\sharp}\bigl(P^{\rho_{\mathsf{r}}[Y]};\mathbf{A},(\Phi_W^{[\mathring c_0]})^{[v]},g\bigr)
=N^{\sharp}\bigl(P^{\hat v^{-1}\rho_{\mathsf{r}}[Y]};\mathbf{A},\Phi_W^{[\mathring c_0]},g\bigr).
\]
The rotation slot is here $\hat v^{-1}g$, which equals $g$ off $\Sigma_W$ by clause (iii).
Moreover $(\Phi_W^{[\mathring c_0]})^{[v]}=\Phi_W^{[\mathring c]}$, since $v\,\mathring c_0=\mathring c$
and the rotation by $v$ applied after the rotation by $\mathring c_0$ is the rotation by
$v\,\mathring c_0$. At the roots,
\[
(\hat v^{-1}\rho_{\mathsf{r}}[Y])(s)=v(s)^{-1}e^{-Y_s}=\bigl(c(s)\mathring c_0(s)^{-1}\bigr)^{-1}
=\rho_{\mathsf{r}}[\log(c\mathring c_0^{-1})](s).
\]

Now apply Step 1 at this one $\mathsf{r}$ with the following two maps.
\begin{itemize}
\item $r_1:=\hat v^{-1}\rho_{\mathsf{r}}[Y]=\rho_{\mathsf{r}}[\log v]\rho_{\mathsf{r}}[Y]$, by the
identity $\hat v^{-1}=\rho_{\mathsf{r}}[\log v]$ just recalled.
This map has two factors. For $c$ close to $c_0$ we have
$\max_s\|\log v(s)\|\le 0.05\,\varepsilon^{\mathrm{fr}}(W)$, since $v\to1$ as $c\to c_0$. Also
$\max_s\|Y_s\|=\max_s\operatorname{pol}(c(s))<\varepsilon^{\mathrm{fr}}(W)$. Hence
$z(r_1)\le 1.05\,\varepsilon^{\mathrm{fr}}(W)$.
\item $r_2:=\rho_{\mathsf{r}}[\log(c\mathring c_0^{-1})]$. This map has one factor. We have
$\max_s\|\log(c(s)\mathring c_0(s)^{-1})\|<\varepsilon^{\mathrm{fr}}(W)$ for $c$ close to $c_0$, its
value at $c_0$ being $\|Y_0\|<\varepsilon^{\mathrm{fr}}(W)$. Hence
$z(r_2)\le 1.05\,\varepsilon^{\mathrm{fr}}(W)$.
\end{itemize}
The two maps agree on $\Sigma_W$, as just computed, and $P\in\mathfrak{U}_N(X)$ near $P_0$.
Step 1 now identifies
$N^{\sharp}(P^{\hat v^{-1}\rho_{\mathsf{r}}[Y]};\mathbf{A},\Phi_W^{[\mathring c_0]},g)$ with $f(P,c)$.
Hence, for $(P,c)$ near $(P_0,c_0)$,
\[
N^{\sharp}\bigl(P^{\rho_{\mathsf{r}}[Y]};\mathbf{A},\Phi_W^{[\mathring c]},g\bigr)=f(P,c),
\]
a holomorphic function of $(P,c)$.

Property (O2) of the definition of $N^{\sharp}$ holds: for $G$-valued $c$ one has $Y=0$, and the
map is $\rho_{\mathsf{r}}[0]=1$. The bound asserted in clause (ii) of the frame lemma holds because
$P^{\rho}\in\mathfrak{U}^{\mathrm{tr}}_N(X)$, by \eqref{4.22:eq-frame-proof-increments-c}.

\emph{Independence of the decomposition.} Let $\mathsf{r},\mathsf{r}'$ be two decompositions of
$P\in\mathfrak{U}_N(X)$. Fix $\mathbf{A}$, $g$ and the real $\Phi_W$. For
\[
c\in\mathfrak{G}_{\varepsilon^{\mathrm{fr}}}(\Sigma_W)
=\prod_{s\in\Sigma_W}\mathfrak{G}_{\varepsilon^{\mathrm{fr}}(W)},
\qquad c(s)=e^{Y_s}\mathring c_s ,
\]
put $\Gamma_{\mathsf{r}}(c):=N^{\sharp}(P^{\rho_{\mathsf{r}}[Y]};\mathbf{A},\Phi_W^{[\mathring c]},g)$,
and define $\Gamma_{\mathsf{r}'}(c)$ in the same way with $\rho_{\mathsf{r}'}$. Both are defined, by
the bound of Step 2 at $\mathsf{r}$ and at $\mathsf{r}'$ respectively.

Run Step 2 at $\mathsf{r}$ about an arbitrary $c_0\in\mathfrak{G}_{\varepsilon^{\mathrm{fr}}}(\Sigma_W)$,
with $P_0:=P$. It exhibits $\Gamma_{\mathsf{r}}$ near $c_0$ as the holomorphic function
$f(P,\cdot)$. So $\Gamma_{\mathsf{r}}$ is holomorphic on
$\mathfrak{G}_{\varepsilon^{\mathrm{fr}}}(\Sigma_W)$, and the
same holds for $\Gamma_{\mathsf{r}'}$. At $G$-valued $c$ we have $Y=0$ and
$\rho_{\mathsf{r}}[0]=1=\rho_{\mathsf{r}'}[0]$, so both functions equal
$N^{\sharp}(P;\mathbf{A},\Phi_W^{[c]},g)$. The domain
$\mathfrak{G}_{\varepsilon^{\mathrm{fr}}}(\Sigma_W)$ is a finite product of charts
$\mathfrak{G}_\varepsilon$. By clause (g3) of the lemma on the polar charts, two holomorphic
functions on such a product that agree at all $G$-valued points coincide. Hence
$\Gamma_{\mathsf{r}}-\Gamma_{\mathsf{r}'}\equiv0$.

So $N^{\sharp}(P;\mathbf{A},\Phi_W,(g,c))$ does not depend on $\mathsf{r}$. Step 2's function $f$
is built on a decomposition of $P_0$ that is a decomposition of every $P$ near $P_0$, so $f$ is the
restriction of $N^{\sharp}$ near $(P_0,c_0)$. Now $f$ is the inner $N^{\sharp}$, which is
holomorphic jointly in all its arguments by property (O1), composed with maps holomorphic in
$(P,c)$. Therefore $N^{\sharp}$ is holomorphic jointly in $P$, the weak coordinates, $g$ and $c$.

\emph{Clause (iv).} Consider first the first fluctuation-integral site $(\mathrm{T})$ after the
birth of $W$: the step $k_W+1$ of the debit lemma of the fluctuation-integral step,
at layer $n=k_W$. The factor $1$ of the stored domain exceeds the layer factor
\[
F_n=f_{k_W+1,k_W}+\tfrac12\theta_S\beta=1-\tfrac12\beta+\tfrac12\theta_S\beta
\]
of layer $n$ by $\tfrac12(1-\theta_S)\beta$; here $f_{k_W+1,k_W}$ is the rate factor of the debit
lemma from level $k_W$ to level $k_W+1$. By clause (ii) of the debit lemma at $n=k_W$, where there is
no attenuation, and by the homogeneity condition of that step, the debit $\pi_n$ of layer $n$,
measured in units of the radius $\alpha_{0,n}$, obeys
\begin{equation}\label{4.22:eq-frame-proof-datum-3}
\begin{aligned}
\frac{\pi_n}{\alpha_{0,n}}
&\le\frac{C_H(1+\beta_0)^2\|C\|\,(1+\eta_p|\mathsf{t}|)\,g_{k_W}\sup|\mathsf{B}|}{\alpha_{0,k_W}}\\
&\le\frac{C_H(1+\beta_0)^2\|C\|}{K_R}\le\tfrac14(1-\theta_S)\beta .
\end{aligned}
\end{equation}
Here $C_H$, $\beta_0$, $\|C\|$ and $\eta_p$ are the constants of the debit lemma, $\mathsf{t}$ is the
disc parameter and $\sup|\mathsf{B}|$ the supremum of the field variable $\mathsf{B}$ of its
estimate, and $g_{k_W}$ is the running coupling at the step. The last inequality is the lower bound
on the constant $K_R$ fixed in the order of constants,
\[
K_R\ge\frac{4(1+\beta_0)^2\|C\|\,C_He^{16\kappa_1}}{(1-\theta_S)\beta},
\]
in which the constant $C_H$ of the debit lemma is replaced by $C_He^{16\kappa_1}$ under the
restriction $|\sigma|\le e^{\kappa_1}$ on the interpolation variable $\sigma$ of that lemma. The
share $\tfrac14(1-\theta_S)\beta$ and the reserve
$\varpi_{\mathrm{fr}}=\tfrac18(1-\theta_S)\beta$ together fit into the excess
$\tfrac12(1-\theta_S)\beta$, since $\tfrac14+\tfrac18\le\tfrac12$. At the later
fluctuation-integral sites, $q\ge2$ steps after the birth of $W$, the excess is
$\beta[1-(1+\theta_S)2^{-q}]\ge\tfrac12\beta$.

At the large-field site $(\mathrm{P})$ and the presentation site $(\mathrm{L})$ of a step $k>k_W$,
write $s^{(k)}_{k_W}$ for the layer factor reached at step $k$ by a frame born at step $k_W$, and
$r_T$ for the radius share left free at the site once the reserve $\varpi_{\mathrm{fr}}$ is set
aside. Then
\begin{equation}\label{4.22:eq-frame-proof-datum-4}
\begin{aligned}
r_T&=(1-\varpi_{\mathrm{fr}})-s^{(k)}_{k_W}
=\beta\bigl(1-2^{-(k-k_W)}\bigr)-\varpi_{\mathrm{fr}}\\
&\ge\tfrac12\beta-\tfrac18(1-\theta_S)\beta\ \ge\ \tfrac12\theta_S\beta
\ \ge\ \beta\theta_S2^{-(k-k_W)} .
\end{aligned}
\end{equation}
The middle inequality is equivalent to $\tfrac38\ge\tfrac38\theta_S$. The last holds because
$k-k_W\ge1$. Hence the radius condition $r_{\min}(y)\ge\beta e_*(j)$ of the statement on the minimal
radius at a site, in which $r_{\min}(y)$ is the minimal radius required at a site $y$ and
$e_*(j)=\theta_S2^{-(k-j)}$ for a frame born at level $j$, and the rate
$\min(\theta_S,\tfrac12)\beta$ of the proposition on layer rates stand. In summary, the reserve fits
at every site:
\begin{equation}\label{4.22:eq-frame-proof-datum-b}
\begin{aligned}
&\tfrac14(1-\theta_S)\beta+\varpi_{\mathrm{fr}}\le\tfrac12(1-\theta_S)\beta
&&\text{(first fluctuation-integral site)},\\
&\beta\bigl[1-(1+\theta_S)2^{-q}\bigr]\ge\tfrac12\beta\quad(q\ge2)
&&\text{(later fluctuation-integral sites)},\\
&r_T\ge\beta\theta_S2^{-(k-k_W)}
&&\text{(sites $(\mathrm{P})$, $(\mathrm{L})$ of a step $k>k_W$)}.
\end{aligned}
\end{equation}

\emph{Clause (v).} A second-class outcome is a creation \cite[\S1]{B-CMP122b}. Write $k$ for the
current step, $h$ for the last step up to $k$ in whose preparatory operations new large fields are
introduced, $h_2$ for the corresponding level of the second domain $\Lambda_2$ of the term, and
$k_a$ for the step at which a pointwise coordinate $a$ was created.
Ba{\l}aban's condition (ii) on the large-field characteristic functions, together with the condition
on the creation steps of the domain, gives $k_W\le h_2$ by the argument of the lemma on the last
creation step. This is the
use Ba{\l}aban himself makes of condition (ii). Of the last operations of the sequence he writes in
\cite{B-CMP122a}: ``by the condition (ii) no large field characteristic functions are included in
them''; and, where new large fields are introduced, ``we have introduced new large fields,
therefore they do not satisfy the condition (ii)''. The
clock is reset in \cite[\S1]{B-CMP122b}: ``a new large field was introduced in the preparatory
operations \ldots\ $K=R_{j+1}$ for $Z$''. So $k_W$ and $k_a$ do not lie in $(h,k]$, that is,
$k_W,k_a\le h$; this is the inequality $h_b\ge k_a$ of the lemma on the last creation step, $h_b$
being the last creation step of that lemma. We do not claim the sharper bound $k_W,k_a\le h-1$.

Furthermore $W\subset Z_{k_W}\subset Z_{h_2}\subset\Lambda_2$. The data $\Phi_W$ are among the inner
variables of $\mathbf{T}_h(Z_h\cap X)$, by clause (f3) of the Haar-invariance statement for the
inner variables. They are either integrated against the measure $d\nu_i$ of a component $Y_i$ of
the large-field operation, by the integration statement for the inner variables, or kept under
the large-field operation $\mathbf{T}'_{k_2}(Y_i)$ on that component at the level $k_2$. The new
terms are functions of the real data, jointly invariant under the
simultaneous constant rotation $c$ of every inner variable, by clauses (f2), (f3) of the same
statement and clause (ii) of the invariance theorem for the boundary terms. Clause (ii) of the
frame lemma therefore applies to them afresh.

\emph{Clause (vi).} This clause follows from Ba{\l}aban's construction in the passage of
\cite{B-CMP122b} that comes after his description of the new large-field region and of the form
\cite[(1.102)]{B-CMP122b} of the operation $\mathbf{T}_k$ on a component; it is used here as
stated there.
\end{proof}

\begin{remark}
(1) By property (O3) of the definition of $N^{\sharp}$, the function $N^{\sharp}$ is the
holomorphic continuation in $c$ of $c\mapsto N(P;\Phi_W^{[c]})$. By property (O6) it is a function
of the formally rotated data, so the notation $\boldsymbol{\Phi}^{\langle g\rangle}$ of the
presentation is to be read literally for frames too.

(2) \emph{Composites.} A later term built from $N$ carries $\Phi_W$ through the chart variable of
$N^{\sharp}$, by the propositions on composite terms. Its own $(\mathbf{U},\mathbf{J})$-domain is
not touched: the twist accumulates in the chart of $N$. Hence the twist is accounted for at
$\Sigma_W$, in the unit $\varepsilon^{\mathrm{fr}}(W)$.

(3) \emph{Why the reserve is not a credit depending on $(d,L,M)$.} The reserve is a fraction of a
radius proportional to $\alpha_{1,k_W}$; no larger domain is established here. A stronger
statement --- that the walk-type dependence on the twist is entire, with a radius depending on
$(d,L)$ only, each crossing of $\partial\Lambda_W$ by a walk costing a factor
$e^{2\operatorname{pol}(c)}$ against the decay of the kernel --- is neither needed nor asserted.

(4) The three arrest cases of the large-field operation that split $V'$ carry no frame datum in the
chapter's list of terms: that splitting lives in the characteristic functions $\chi$ and the
$\delta$-functions.
\end{remark}

\begin{remark}[The kept terms at creation: what the printed step asserts]
\label{4.22:rem-kept-terms-creation}
Let $W$ be a frame contained in a component $Y_i$, let $k_W$ be the step at which the
large-field operation $\mathbf{T}'_{k_W}(Y_i)$ of \cite[(1.71)]{B-CMP122b} is formed on $Y_i$,
and let $Y_i^{\sim}$ be the enlarged
domain attached to $Y_i$ in \cite{B-CMP122b}; the symbol $Y_i^{\sim}$ is the one written
$X^{\sim}$ there, with $X=Y_i$. By \cite[(1.102)]{B-CMP122b} two terms are kept inside
$\mathbf{T}'_{k_W}(Y_i)$:
\begin{itemize}
\item the potential $V(Y_i^{\sim})$;
\item the quadratic form in $B$ of \cite[(1.71)]{B-CMP122b}, denoted $\mathsf{Q}_i$.
\end{itemize}
They form the kept lineage $\mathfrak{K}(Y_i)$. In \cite[pp.~378--379]{B-CMP122b} the following
is asserted about $\mathfrak{K}(Y_i)$.
\begin{enumerate}
\item[(a)] Smallness is asserted only for the first-order term of the two members in the
variables $A'$, $J$ in which \cite[(1.71)]{B-CMP122b} is expanded, that is, for a change. It
is asserted for a single purpose: the comparison, for each
decomposition $\mathsf{r}=(U,\{u_{\square}\})$ of the background pair $(\mathbf{U},\mathbf{J})$
(Definition~\ref{4.22:not-tent-gauge-map}), of
$\mathbf{T}'_{k_W}(Y_i,(\mathbf{U},\mathbf{J}))$ with $\mathbf{T}'_{k_W}(Y_i,(U,0))$. This is a
comparison at the creating step's own operation $\mathbf{T}'_{k_W}$.
\item[(b)] For the size of $V(Y_i^{\sim})$, \cite[pp.~378--379]{B-CMP122b} keeps the count of
old cells of \cite[(1.69)]{B-CMP122b}: every large-field domain $Z_j$ contributes a constant
\begin{equation*}
O(1)\log g_j^{-2}|Z_j|\le O(1)M^dR_j^{d+1}d'_j(Z_j),
\end{equation*}
with $d'_j$ as in \cite[(1.69)]{B-CMP122b}.
\item[(c)] The cited pages do not refer to the sites of later steps. The resummation they
describe is given there by reference to the way discussed for the other terms $V(Y)$.
\end{enumerate}
The uses of the kept lineage made below do not rest on \cite[pp.~378--379]{B-CMP122b}. At the
creating step $k_W$ the kept lineage is used in two places only:
\begin{enumerate}
\item[(1)] $V(Y_i^{\sim})$ is bounded on the tube $\mathfrak{U}^{\mathrm{tr}}(Y_i)$ by the
weighted form of \cite[(1.69)]{B-CMP122b}; here $\mathfrak{U}^{\mathrm{tr}}(Y_i)$ is the space
of complex configurations on which the kept members are proved holomorphic and bounded at
step $k_W$. The route through the small term $\sigma$ of
\cite[(1.74)--(1.75)]{B-CMP122b} is not used.
\item[(2)] The complex part of $\mathsf{Q}_i$ does not depend on the complex source parameter
$z$, and it is small. This is the only use of $\mathsf{Q}_i$ at the creating step $k_W$.
\end{enumerate}
Both are inputs of the proof of the correlation theorem.

Every later site that moves the background on $Y_i^{\sim}$ evaluates the two members of
$\mathfrak{K}(Y_i)$. There are three kinds of such site:
\begin{itemize}
\item a site $(\mathrm{T})$;
\item each stage of a chain at a site $(\mathrm{P})$, because the layer moved at a stage has
tails;
\item each link of a site $(\mathrm{L})$. This includes the link at the healing of $Y_i$
itself, that is, at the later step at which its large-field factor is absorbed back into the
small-field format; there $Y_i^{\sim}$ lies in the core of the link.
\end{itemize}
What enters such a site is a difference of values of a member at two backgrounds.

If a difference of this kind were estimated by $\sup|v|$ for a member $v$, that is, by
\cite[(1.69)]{B-CMP122b}, the estimate would carry the nesting depth
\begin{equation*}
D=k_W-j(W),
\end{equation*}
where $j(W)\le k_W$ is the step from which the nesting depth of $W$ is counted.
This depth is unbounded at fixed $x_{k_W}$. The member $V(Y_i^{\sim})$ is therefore estimated
instead by the two-point bound over the whole tube, which is proved as an implication from the
base bound for $V$; its right member is
$C_{\Sigma}\check{R}_{k_W}^{p_\Sigma}$. That right member is carried unchanged into the bounds
of the later sites, where the collar of $Y_i$ pays it.
\end{remark}

\begin{proof}
\emph{The quoted text.} On \cite[pp.~378--379]{B-CMP122b}, in the notation of that paper (its
step $k$ and component $X$ are, here, $k_W$ and $Y_i$), one reads:
\begin{quote}
The sum in the last exponential does not include terms with localization domains equal to
one of the components of $Z^{\sim}$; these terms are included into the operations
$\mathbf{T}'_k$.
\end{quote}
and
\begin{quote}
In the definition (1.71) only the quadratic form in $B$ and the term $V(X^{\sim})$ depend on
$(\mathbf{U},\mathbf{J})$. We expand them up to the first order in $A'$, $J$, the first order
terms are already small. This is clear for the expansion of the quadratic form, for which the
first order term can be bounded by
\begin{align*}
&O(1)B_3^2A_1^2p_1^2(g_k)|\mathbf{B}_0|(\alpha_{0,k}+\alpha_{1,k})\\
&\qquad< O(1)B_3^2A_1^2C_0M^dR_k^{d+2}p_1^2(g_k)q_0(g_k)g_k,
\end{align*}
and this bound is small. The expression $V(X^{\sim})$ is a sum of a big number of terms, mainly
boundary terms from all the renormalization steps, depending on the configuration
(1.61)--(1.63) localized in $X^{\sim}$. For all these terms the expansion in $A'$, $J$ creates
a connection between their localization domains and the set
$\mathbf{B}_0\cup(\Gamma''_k\cap X^{\sim})$, a connection with a proper exponential decay.
This allows us to resum all these terms, in the way discussed in connection with the other
terms $V(Y)$. The resummed expression has the same bound as above, with a different constant
$O(1)$. Thus the operation $\mathbf{T}'_k(X)$ can be represented as the operation for a
regular $G$-valued configuration $U$, with the additional small, and possible complex valued,
term in the exponential. All the expressions in the definition (1.71), for the configuration
$U$, are real, and the exponential density in the integral is positive. This implies the
inequalities [\dots] (1.73) [\dots] (1.74), where $\sigma$ is the small term of the first
order in $A'$, $J$, and $F$ is a function of the integration variables in the integral (1.71).
\end{quote}
Further on in \cite[pp.~378--379]{B-CMP122b}:
\begin{quote}
The expression $V(X^{\sim})$ is bounded in (1.69) [\dots] every large field domain $Z_j$
contributes the constant
\begin{equation*}
O(1)\log g_j^{-2}|Z_j|\le[\dots]\le O(1)M^dR_j^{d+1}d'_j(Z_j)
\end{equation*}
\end{quote}

\emph{The kept members.} By the first quoted sentence, the terms whose localisation domains
are components of $Z^{\sim}$ are not summed in the exponential. They are included in the
operations $\mathbf{T}'_k$. By the first sentence of the second quotation, the only parts of
\cite[(1.71)]{B-CMP122b} that depend on $(\mathbf{U},\mathbf{J})$ are the quadratic form in
$B$ and $V(X^{\sim})$. With $k=k_W$ and $X=Y_i$ these are $\mathsf{Q}_i$ and
$V(Y_i^{\sim})$, the two members of $\mathfrak{K}(Y_i)$ kept in \cite[(1.102)]{B-CMP122b}.

\emph{Clause (a).} Both members are expanded to first order in $A'$ and $J$. The sentence
``the first order terms are already small'' attaches smallness to the first-order terms. These
are changes of the members between the configuration $(\mathbf{U},\mathbf{J})$ and the
configuration at $A'=0$, $J=0$. The displayed bound is a bound on the first-order term of the
quadratic form. For $V(X^{\sim})$, the resummed first-order expression is asserted to have
the same bound with a different constant $O(1)$.

The sentence beginning ``Thus'' says what the smallness is for. It is used to represent
$\mathbf{T}'_k(X)$ as the operation for the regular $G$-valued configuration $U$, with an
additional small term in the exponential. Reality at $U$ and positivity of the density then
give \cite[(1.73), (1.74)]{B-CMP122b}, with $\sigma$ the small first-order term. This is a
comparison of $\mathbf{T}'_{k_W}(Y_i,(\mathbf{U},\mathbf{J}))$ with
$\mathbf{T}'_{k_W}(Y_i,(U,0))$. It is made for the configuration $U$ of a decomposition, and
inside the operation of step $k_W$ itself.

\emph{Clause (b).} The third quotation bounds the size of $V(X^{\sim})$ through
\cite[(1.69)]{B-CMP122b}. It does so by counting, for each large-field domain $Z_j$, the
constant $O(1)\log g_j^{-2}|Z_j|$, which is bounded by $O(1)M^dR_j^{d+1}d'_j(Z_j)$. This is a
count of old cells.

\emph{Clause (c).} No sentence quoted above refers to a step later than $k$. The resummation
of the terms of $V(X^{\sim})$ is described by reference to ``the way discussed in connection
with the other terms $V(Y)$'', and the argument for it is given there by that reference.

\emph{The two uses.} Use (1) bounds $V(Y_i^{\sim})$ by size, through the weighted form of
\cite[(1.69)]{B-CMP122b} on the tube. It therefore needs neither $\sigma$ nor
\cite[(1.74)--(1.75)]{B-CMP122b}. Use (2) concerns $\mathsf{Q}_i$ at step $k_W$ only. It rests
on the holomorphy and the bounds of the kept quadratic form on the tube. It does not rest on
the first-order smallness asserted in
\cite[pp.~378--379]{B-CMP122b}. Hence neither use rests on
\cite[pp.~378--379]{B-CMP122b}.

\emph{Later sites.} By \cite[(1.102)]{B-CMP122b} the two members stand inside
$\mathbf{T}'_{k_W}(Y_i)$. A later site whose background moves on $Y_i^{\sim}$ therefore
evaluates them at moved backgrounds. This covers a fluctuation-integral site $(\mathrm{T})$,
every stage of a chain at $(\mathrm{P})$ whose layer has a tail on $Y_i^{\sim}$, and every
link of $(\mathrm{L})$, including the healing of $Y_i$. What enters each such site is the
difference of the member's values at the two backgrounds.

Estimating this difference by $\sup|v|$, that is, by \cite[(1.69)]{B-CMP122b}, brings in the
count of old cells of clause (b). That count runs over the nested large-field domains, and so
carries the nesting depth $D=k_W-j(W)$. By \cite[conditions (i), (ii)]{B-CMP122a}, $D$ is
unbounded at fixed $x_{k_W}$. An age (the age of a region at a step being the number of steps
since the step that created it) would thus appear explicitly at each of the three kinds
of site.

\emph{The two-point bound.} A one-point bound of the form
\begin{equation*}
|v(\mathbf{U},\mathbf{J})-v(U,0)|\le\Sigma(x_{k_W}),
\end{equation*}
for some non-increasing function $\Sigma$ of the inverse coupling with $\Sigma(x)\to0$ as
$x\to\infty$, for both members would need a unitary reference configuration $U$ independent
of the parameter of the family of backgrounds along which the members are compared; this need
is the reason why the comparison of clause (a) is read for each decomposition. It would be
free of age, by the coupling flow of Theorem~\ref{4.1:thm-clause-zero}, $\Sigma$ being
monotone. No bound of this form is used for $V(Y_i^{\sim})$.

What is used is the two-point bound for $V(Y_i^{\sim})$, proved as an implication from the base
bound for $V$: the difference of the values of $V(Y_i^{\sim})$ at two arbitrary
points of the tube is bounded by $C_{\Sigma}\check{R}_{k_W}^{p_\Sigma}$. It requires no
reference point.
Its right member, however, belongs to the scale of step $k_W$, and it grows as
$g_{k_W}\downarrow0$. Expressed through the coupling of a later step $k'>k_W$ at which the
member is evaluated, it would leave a power of $\log(1+B_\sharp(k'-k_W))$, in which the age
$k'-k_W$ of $Y_i$ at step $k'$ appears explicitly. For this reason it is not converted.
It is carried as $C_{\Sigma}\check{R}_{k_W}^{p_\Sigma}$ into the bounds of the later sites,
where the collar of $Y_i$ pays it.
\end{proof}

\begin{lemma}[Oscillation of the kept terms on the orbit chart]\label{4.22:lem-kept-oscillation}
\emph{Setting.} Let $Y_i$ be a component carrying a large-field operation of the form
\cite[(1.71)]{B-CMP122b}, and let its kept lineage $\mathfrak{K}(Y_i)$ consist of the two members
that the operation keeps, the quadratic form $\mathsf{Q}_i$ and the term $V(Y_i^{\sim})$. Put
$W:=Y_i$, let $k_W$ be the level of the frame $W$, and let $v\in\mathfrak{K}(Y_i)$. Let
$v^{\sharp}$ be the orbit datum of $v$ at the roots $\Sigma_W$ given by clause (ii) of the frame
lemma (Lemma~\ref{4.22:lem-frame-lemma}), with radius $\varepsilon^{\mathrm{fr}}(W)$, on the
stored domain
\begin{equation*}
\mathfrak{U}_v:=\mathfrak{U}_v(Y_i^{\sharp})\subset\mathfrak{U}^{\mathrm{tr}}_v
:=\tilde{U}^{c}_{k_W}(Y_i^{\sharp},\tilde\alpha_0,\tilde\alpha_1)
\subset\mathfrak{S}_{\mathrm{src}}:=\mathfrak{S}_{\mathrm{src}}(Y_i^{\sharp}),
\end{equation*}
the clauses of the tube $\mathfrak{U}^{\mathrm{tr}}_v$ being restricted to the strip
$Y_i^{\sharp}$. The member $v$ has no law site other than those of $\Sigma_W$: its further
arguments are members of $\Phi_W$. Put
\begin{equation}\label{4.22:eq-kept-oscillation-a}
\mathfrak{O}(v):=\sup\bigl|v^{\sharp}(P;\Phi,c)-v^{\sharp}(P';\Phi,c'')\bigr|,
\end{equation}
the supremum being taken over real $\Phi$ in the supports and over all pairs
$(P,c),(P',c'')\in\mathfrak{U}_v\times\mathfrak{G}_{\varepsilon^{\mathrm{fr}}(W)}(\Sigma_W)$;
$\mathfrak{O}$ is a seminorm, it vanishes on constants, and
$\mathfrak{O}(v)\le 2\mathfrak{N}^{\mathbb{C}}(v)$, where $\mathfrak{N}^{\mathbb{C}}(v)$ is the
supremum of $|v^{\sharp}|$ over the same set. Put further
$\mathfrak{o}(V(Y_i^{\sim})):=\Sigma^V(k_W)$ and $\mathfrak{o}(\mathsf{Q}_i):=2\mathfrak{q}_{\max}(k_W)$.

\emph{Hypotheses.} The clauses below use the following inputs, each where it is named.
\begin{itemize}
\item[(KH1)] The two-point bound for the kept potential term on the sourced set: for
$Q,Q'\in\mathfrak{S}_{\mathrm{src}}$ and real further arguments $\Psi$,
$|V(Y_i^{\sim})(Q;\Psi)-V(Y_i^{\sim})(Q';\Psi)|\le\Sigma^V(k_W)$.
\item[(KH1$'$)] The bound of the total share of the kept members:
\begin{equation*}
\Sigma^V(k_W)+2\mathfrak{q}_{\max}(k_W)\le\Sigma^{\mathrm{tot}}(k_W)
\le C_\Sigma\check{R}_{k_W}^{p_\Sigma}.
\end{equation*}
\item[(KH2)] The kept quadratic form $\mathsf{Q}_i$ is holomorphic on $\mathfrak{U}^{\mathrm{tr}}_v$
and jointly invariant under $G$-valued gauge transformations acting on the background and on
the further arguments.
\item[(KH3)] The bound of the kept quadratic form on the sourced set, uniformly in real
further arguments: $|\mathsf{Q}_i|\le\sup|\mathfrak{q}_C|\le\mathfrak{q}_{\max}(k_W)$ there.
\item[(KH4)] For the family $\mathsf{R}(\tau)$, $\tau\in\mathfrak{P}$, and the riding map
$\mathfrak{u}(\tau)$ described before clause (ii$_2$):
$(\mathsf{R}(\tau),\mathfrak{u}(\tau)g')\in\mathfrak{U}_v\times
\mathfrak{G}_{\varepsilon^{\mathrm{fr}}(W)}(\Sigma_W)$ for all $\tau\in\mathfrak{P}$.
\item[(KH5)] $\mathsf{R}(\tau)\in\mathfrak{U}_v$ with the decomposition
$(U_{\mathbf{K}},\{u_{\square}\})$ of the base (the fifth and eighth clauses of the tube do not
move along the family).
\item[(KH6)] In the setting of clause (iv), for the pairs $P_\lambda$ and the radius $\rho_X$
defined there: $P_\lambda\in\mathfrak{U}_v$ for $|\lambda|\le\rho_X$.
\item[(KH7)] At every $s\in\Sigma_W$, for a number $\varsigma(s)$,
\begin{equation}\label{4.22:eq-kept-oscillation-b}
\begin{gathered}
\sup_{\mathfrak{P}}\|\log\mathfrak{u}(\tau)(s)\|\le\varsigma(s)\le\varepsilon^{\mathrm{fr}}(W)/16,\\
\operatorname{pol}(g'(s))<\varepsilon^{\mathrm{fr}}(W),\qquad
\operatorname{pol}(\mathfrak{u}(\tau)(s)g'(s))<\varepsilon^{\mathrm{fr}}(W).
\end{gathered}
\end{equation}
\item[(KH8)] The one-point comparison of the kept form with its unitary reference: for
$(\mathbf{U},\mathbf{J})\in\mathfrak{U}^{\mathrm{tr}}_v$ written $\mathbf{U}=e^{i\eta A'}U$ in a
decomposition with unitary part $U$ and $\mathfrak{g}^{\mathbb{C}}$-valued $A'$, and for real
further arguments $\Psi$ in the supports,
$|\mathsf{Q}_i(\mathbf{U},\mathbf{J};\Psi)-\mathsf{Q}_i(U,0;\Psi)|\le\Sigma^{\mathfrak{q}}(x_{k_W})$.
\item[(KH9)] The bounds of the led rotations at $\Sigma_W$ in the setting of clause (iv), for the
rotations $\mathsf{u}_\lambda=\mathsf{u}'e^{\lambda\mathcal{Z}}$ defined there:
$\operatorname{pol}(\mathsf{u}_0(s))\le\eta_0(s)$ and
$|\lambda|\,\|\mathcal{Z}(s)\|\le\vartheta^{3F}(s)\le\tfrac12\mu(s)$ for $|\lambda|\le\rho_X$.
\item[(KH10)] Clause (ii) of the gauge invariance theorem for $V(Y_i^{\sim})$, for all
$G$-valued gauge transformations: for every $G$-valued gauge transformation $g$ on
$Y_i^{\sharp}$, every $Q\in\mathfrak{U}^{\mathrm{tr}}_v$ and every real $\Phi$,
$V(Y_i^{\sim})(Q^{g};\Phi^{[g]})=V(Y_i^{\sim})(Q;\Phi)$.
\item[(KH11)] The domain statement of the lemma on the orbit datum along led rotations, for the
data of clause (iv): whenever \textup{(KH6)} and \textup{(KH9)} hold, then for $|\lambda|<\rho_X$
the pair $(P_\lambda,\mathsf{u}_\lambda g')$ lies in
$\mathfrak{U}_v\times\mathfrak{G}_{\varepsilon}(\Sigma_W)$ with
$\varepsilon\le\varepsilon^{\mathrm{fr}}(W)$, and
$\lambda\mapsto v^{\sharp}(P_\lambda;\Phi_W,\mathsf{u}_\lambda g')$ is holomorphic on
$|\lambda|<\rho_X$.
\item[(KH12)] The invariance of Ba{\l}aban's spaces of regular configurations under $G$-valued
gauge transformations \cite{B-CMP119}: ``The spaces defined above are invariant with respect to
$G$-valued gauge transformations.''
\end{itemize}

\emph{(o)} Assume \textup{(KH1)}, \textup{(KH1$'$)}, \textup{(KH2)}, \textup{(KH3)},
\textup{(KH10)}, \textup{(KH12)}. Then
\begin{equation}\label{4.22:eq-kept-oscillation-1}
\begin{gathered}
\mathfrak{O}(v)\le\sup_{\Psi\ \mathrm{real}}\ \sup_{Q,Q'\in\mathfrak{U}^{\mathrm{tr}}_v}
|v(Q;\Psi)-v(Q';\Psi)|\le\mathfrak{o}(v),\\
\mathfrak{O}(V(Y_i^{\sim}))+\mathfrak{O}(\mathsf{Q}_i)\le\Sigma^{\mathrm{tot}}(k_W)
\le C_\Sigma\check{R}_{k_W}^{p_\Sigma}.
\end{gathered}
\end{equation}

At a later site of a step $k$, let $\mathsf{R}(\tau)$, $\tau\in\mathfrak{P}$, be the family of
background pairs that the consuming term evaluates on $Y_i^{\sharp}$, with first component
$e^{i\eta\mathcal{K}(\tau)}\mathbf{K}$, where $\mathbf{K}=e^{i\eta A'_{\mathbf{K}}}U_{\mathbf{K}}$ is
its base with the decomposition $\mathsf{r}=(U_{\mathbf{K}},\{u_{\square}\})$, independent of
$\tau$; let $\mathfrak{u}(\tau)$ be the riding map (at the presentation site, the product
$h_q\mathfrak{Q}^{\mathfrak{b}}$ of the presenter maps), and $g'$ a point of the output chart.

\emph{(ii$_2$)} Assume the hypotheses of \textup{(o)} and \textup{(KH4)}. Then
\begin{equation}\label{4.22:eq-kept-oscillation-2}
\begin{aligned}
\operatorname{osc}_{\mathfrak{P}}(v^{\sharp})
&:=\sup_{g'}\sup_{\tau,\tau'\in\mathfrak{P}}
\bigl|v^{\sharp}(\mathsf{R}(\tau);\Phi_W,\mathfrak{u}(\tau)g')
-v^{\sharp}(\mathsf{R}(\tau');\Phi_W,\mathfrak{u}(\tau')g')\bigr|\\
&\le\mathfrak{O}(v)\le\mathfrak{o}(v)\le C_\Sigma\check{R}_{k_W}^{p_\Sigma},
\end{aligned}
\end{equation}
the variation of the riding map included. The bound contains no nesting depth, no $|Y_i|$, no
measure of the earlier large-field regions meeting $Y_i$, no age, and no number, size or
history of an inner large-field region; it uses membership in the set only, and neither a
decomposition nor \eqref{4.22:eq-kept-oscillation-b} nor a one-point comparison. Its right member
depends on $\check{R}_{k_W}$: it is not a function of the current coupling, and it is not
converted into one.

\emph{(i)} Assume \textup{(KH2)}, \textup{(KH5)}, \textup{(KH7)}, \textup{(KH10)} and
\textup{(KH12)}. Write, by polar
decomposition, $g'(s)=e^{Y'_s}\omega'_s$ and $\mathfrak{u}(\tau)(s)=\omega_u(s)e^{Y_u(s)}$. Then
\begin{equation}\label{4.22:eq-kept-oscillation-3}
\begin{gathered}
v^{\sharp}(\mathsf{R}(\tau);\Phi_W,\mathfrak{u}(\tau)g')
=v\bigl(Q(\tau);\Phi_W^{[\omega']}\bigr),\\
Q(\tau):=\mathsf{R}(\tau)^{\rho_{\mathsf{r}}[Y']\rho_{\mathsf{r}}[Y_u]\rho_{\mathsf{r}}[\log\omega_u]}
\in\mathfrak{U}^{\mathrm{tr}}_v(Y_i^{\sharp}),
\end{gathered}
\end{equation}
and $Q(\tau)$ has a decomposition with unitary part $U_{\mathbf{K}}$.

\emph{(ii$_{\mathsf{Q}}$)} Assume the hypotheses of (i) and \textup{(KH8)}. Then
\begin{equation}\label{4.22:eq-kept-oscillation-4}
\bigl|\mathsf{Q}_i^{\sharp}(\mathsf{R}(\tau);\Phi_W,\mathfrak{u}(\tau)g')
-\mathsf{Q}_i(U_{\mathbf{K}},0;\Phi_W^{[\omega']})\bigr|\le\Sigma^{\mathfrak{q}}(x_{k_W});
\end{equation}
the subtracted number depends neither on $\tau$ nor on $\mathfrak{u}$, so that
\begin{equation*}
\operatorname{osc}_{\mathfrak{P}}(\mathsf{Q}_i^{\sharp})\le2\Sigma^{\mathfrak{q}}(x_{k_W})
\le2\Sigma^{\mathfrak{q}}(x_k-c_5),
\end{equation*}
a function of the current coupling alone, in supremum form, decreasing to zero.

\emph{(iii)} In the transfer of majorant bounds to the orbit chart, the primed form of the third
majorant bound among \eqref{4.22:eq-majorant-transfer-a} holds for $v$: the piece derivatives
$D_\pi v^{\sharp}$ and the constants $c_\pi$ of that display satisfy
$|D_\pi v^{\sharp}|\le c_\pi\mathfrak{O}(v)$; for $\mathsf{Q}_i$ also
$|D_\pi\mathsf{Q}_i^{\sharp}|\le 2c_\pi\Sigma^{\mathfrak{q}}(x_k-c_5)$.

\emph{(iv)} Let $C\ni Y_i$ be treated by the operation $\mathbf{R}$ at step $k$, and
$X:=Y_i^{\sharp}\subset\mathsf{C}_C$; let $h$, with $k_W\le h$, be the level of the large-field
operation $\mathbf{T}_h$ among whose inner variables $\Phi_W$ lies, its region $Z_h$ being
contained in $\mathsf{C}_C$. For the two stages $\mathfrak{b}\in\{\mathfrak{b}_3,\mathfrak{b}_4\}$
of the chain of stages of the presentation (stage letters, not the numbers $\mathfrak{b}_3=400$
of the proof of Lemma~\ref{4.22:lem-frame-lemma}) and
points $\mathfrak{z}=(\zeta,\varpi,\vec{\mathsf{s}})$ of the polydisc $\mathfrak{P}$ of the
three-factor construction put, for $g'\in\mathfrak{G}_{\varepsilon'}(\Sigma_W)$ with
$\varepsilon'(s)=\varepsilon(s)-\eta_0(s)-\mu(s)$, where $\varepsilon(s)\le\varepsilon^{\mathrm{fr}}(W)$
is the chart radius at $s$ of \textup{(KH11)}, $\eta_0(s)$ the bound of \textup{(KH9)} on the
polar size of the led rotation, and $\mu(s)$ the margin at $s$,
\begin{equation*}
F_v^{\mathfrak{b}}(\mathfrak{z};g'):=v^{\sharp}\Bigl(
\bigl((\mathcal{U}_q^{\mathfrak{b}})^{h_q},\operatorname{Ad}_{h_q}\mathcal{J}_q^{\mathfrak{b}}\bigr)
(\mathfrak{z})\restriction X;\ \Phi_W,\ h_q\mathfrak{Q}^{\mathfrak{b}}(\mathfrak{z})g'\Bigr).
\end{equation*}
For $\mathfrak{z}=(\zeta,\varpi,\vec{\mathsf{s}})$ and
$\mathfrak{z}'=(\zeta',\varpi,\vec{\mathsf{s}}\,')$ in $\mathfrak{P}$ with the same $\varpi$,
whether or not $\zeta'=\zeta$, let
$(\mathbb{W},\mathbb{J},\mathsf{u}';\mathcal{K}',\mathfrak{j},\mathcal{Z})$ be the data that the
three-factor construction attaches to the two presented triples at $\mathfrak{z}'$ and
$\mathfrak{z}$, evaluated on $X=Y_i^{\sharp}$, and put
\begin{equation*}
P_\lambda:=\bigl(e^{i\eta\lambda\mathcal{K}'}\mathbb{W},\ \mathbb{J}+\lambda\mathfrak{j}\bigr)
\restriction X,\qquad \mathsf{u}_\lambda:=\mathsf{u}'e^{\lambda\mathcal{Z}},
\end{equation*}
entire in $\lambda$; $(P_0,\mathsf{u}_0)$ and $(P_1,\mathsf{u}_1)$ are the presented triples
at $\mathfrak{z}'$ and $\mathfrak{z}$, the gauge $h_q$ being $1$ on them by the normalisation of
the presentation. Let $\rho_X=r^{0}e^{\frac12\delta_PD_X}$, with $D_X$ as below and
$1/r^{0}=\tfrac12C_{3F}$, be the radius of the $\lambda$-disc of the construction.
Assume the hypotheses of (o), \textup{(KH4)} at this presentation site, \textup{(KH6)},
\textup{(KH9)} and \textup{(KH11)}. Put
$\mathfrak{n}_X:=\tfrac12C_{3F}e^{-\frac12\delta_PD_X}\mathfrak{O}(v)$. Then for all such
pairs $\mathfrak{z}$, $\mathfrak{z}'$ with the same $\varpi$, whether or not $\zeta'=\zeta$,
\begin{equation}\label{4.22:eq-kept-oscillation-5}
|F_v^{\mathfrak{b}}(\mathfrak{z};g')-F_v^{\mathfrak{b}}(\mathfrak{z}';g')|\le\mathfrak{n}_X,
\qquad
|\mathfrak{p}|\le\mathfrak{n}_X\,e^{-(\kappa_1-1)\sum_t|\mathsf{y}_t|}
\end{equation}
for every piece $\mathfrak{p}$ of the expansion of the three-factor bounds (not the parameter
tuple $\mathfrak{p}$ of Theorem~0) whose families $(\mathsf{y}_t)$ are not all empty, and
$|\mathfrak{p}(\zeta)-\mathfrak{p}(\zeta')|$ is bounded by the same right member for every
tuple, the empty tuple included. Thus the three-factor bounds, clauses (a) and (b), and their
$\zeta$-difference form, clauses (a) and (b), hold for $v$ with the rotated supremum replaced by
$\tfrac12\mathfrak{O}(v)$, in the currency of the weighted contour distance $d^{\wedge}$. Here
\begin{equation}\label{4.22:eq-kept-oscillation-6}
D_X=d^{\wedge}(Y_i^{\sharp},\mathsf{C}_C^{c})\ \ge\
\begin{cases}
\tfrac12\mathsf{w}_{\wedge}(\check{R}_{k_W}+1), & k_W=h,\\[2pt]
\tfrac52\mathsf{w}_{\wedge}\check{R}_{k_W}, & k_W<h,
\end{cases}
\end{equation}
where $\mathsf{w}_{\wedge}$ is the lower bound, per unit of level $n$, of the contour distance
$d^{\wedge}$ on a set whose cells have level at most $n$.
\end{lemma}

\begin{proof}
\emph{Clause (o).} Let $(P,c)$ be a point of
$\mathfrak{U}_v\times\mathfrak{G}_{\varepsilon^{\mathrm{fr}}(W)}(\Sigma_W)$ and $\Phi$ real.
Write $c=e^{Y}\omega$ by the polar decomposition at each $s\in\Sigma_W$, with $Y(s)$ Hermitian,
$\|Y(s)\|=\operatorname{pol}(c(s))<\varepsilon^{\mathrm{fr}}(W)$ and $\omega$ $G$-valued. By
clause (ii) of the frame lemma, $v^{\sharp}(P;\Phi,c)=v(P^{\rho_{\mathsf{r}}[Y]};\Phi^{[\omega]})$
for a decomposition $\mathsf{r}$ of $P$. By clause (i) of the frame lemma, in the form of its
corollary \eqref{4.22:eq-frame-proof-increments-c} applied to the single map
$\rho_{\mathsf{r}}[Y]$ of size $\|Y\|<\varepsilon^{\mathrm{fr}}(W)$, this map costs at most
$\tfrac14\varpi_{\mathrm{fr}}r_a$ at every clause $a$ of the tube, $r_a$ being that clause's
threshold; since $P$ lies in the stored domain, $P^{\rho_{\mathsf{r}}[Y]}\in
\mathfrak{U}^{\mathrm{tr}}_v$.

Let $\hat\omega$ be any $G$-valued site map on $Y_i^{\sharp}$ with
$\hat\omega\restriction\Sigma_W=\omega$ (for instance $\hat\omega:=1$ off $\Sigma_W$). For
$V(Y_i^{\sim})$, clause (ii) of the gauge invariance theorem \textup{(KH10)} holds for all
$G$-valued gauge transformations; for $\mathsf{Q}_i$, the joint invariance and holomorphy on
$\mathfrak{U}^{\mathrm{tr}}_v$ are \textup{(KH2)}. The rotated further arguments
$\Phi^{[\hat\omega]}$ depend on $\hat\omega$ only through its restriction to $\Sigma_W$, so
$\Phi^{[\hat\omega]}=\Phi^{[\omega]}$. Hence, for $Q\in\mathfrak{U}^{\mathrm{tr}}_v$, applying
the invariance with the gauge transformation $\hat\omega$ to the pair $Q^{\hat\omega^{-1}}$,
\begin{equation*}
v(Q;\Phi^{[\omega]})=v\bigl((Q^{\hat\omega^{-1}})^{\hat\omega};\Phi^{[\hat\omega]}\bigr)
=v(Q^{\hat\omega^{-1}};\Phi).
\end{equation*}
The tube $\mathfrak{U}^{\mathrm{tr}}_v$ is invariant under $G$-valued gauge transformations
\textup{(KH12)}; no regularity of $\hat\omega$ is required there. Applying this with
$Q=P^{\rho_{\mathsf{r}}[Y]}$, and likewise at $(P',c'')$, we obtain
$v^{\sharp}(P;\Phi,c)=v(Q_1;\Phi)$ and $v^{\sharp}(P';\Phi,c'')=v(Q_1';\Phi)$ with
$Q_1,Q_1'\in\mathfrak{U}^{\mathrm{tr}}_v$ and the same real $\Phi$. Taking the supremum gives
the first inequality of \eqref{4.22:eq-kept-oscillation-1}.

For the second, $\mathfrak{U}^{\mathrm{tr}}_v\subset\mathfrak{S}_{\mathrm{src}}(Y_i^{\sharp})$.
For $v=V(Y_i^{\sim})$, \textup{(KH1)} bounds the two-point difference by
$\Sigma^V(k_W)=\mathfrak{o}(v)$. For $v=\mathsf{Q}_i$, \textup{(KH3)} gives
\begin{equation*}
|\mathsf{Q}_i(Q)-\mathsf{Q}_i(Q')|\le|\mathsf{Q}_i(Q)|+|\mathsf{Q}_i(Q')|
\le2\sup|\mathfrak{q}_C|\le2\mathfrak{q}_{\max}(k_W)=\mathfrak{o}(v).
\end{equation*}
Adding, and using the bound \textup{(KH1$'$)} of the total share of the kept members,
$\Sigma^V(k_W)+2\mathfrak{q}_{\max}(k_W)\le\Sigma^{\mathrm{tot}}(k_W)\le
C_\Sigma\check{R}_{k_W}^{p_\Sigma}$, gives the second line of
\eqref{4.22:eq-kept-oscillation-1}. Since $\Sigma^V(k_W)$ and $2\mathfrak{q}_{\max}(k_W)$ bound
moduli, both are non-negative, so each of $\mathfrak{o}(V(Y_i^{\sim}))$ and
$\mathfrak{o}(\mathsf{Q}_i)$ is at most $C_\Sigma\check{R}_{k_W}^{p_\Sigma}$.

\emph{Clause (ii$_2$).} By \textup{(KH4)}, for $\tau,\tau'\in\mathfrak{P}$ both points
$(\mathsf{R}(\tau),\mathfrak{u}(\tau)g')$ and $(\mathsf{R}(\tau'),\mathfrak{u}(\tau')g')$ lie in
$\mathfrak{U}_v\times\mathfrak{G}_{\varepsilon^{\mathrm{fr}}(W)}(\Sigma_W)$, and both are
evaluated at the one real datum $\Phi_W$. By the definition
\eqref{4.22:eq-kept-oscillation-a} the difference is at most $\mathfrak{O}(v)$; the remaining
inequalities of \eqref{4.22:eq-kept-oscillation-2} are clause (o) and its last remark.

\emph{Clause (i). Sizes.} By \eqref{4.22:eq-kept-oscillation-b},
$\|Y_u\|=\operatorname{pol}(\mathfrak{u})\le\varsigma$. Since $\omega_u=\mathfrak{u}e^{-Y_u}$ and
$\mathfrak{u}=e^{\log\mathfrak{u}}$ with $\|\log\mathfrak{u}\|\le\varsigma$,
\begin{equation*}
\|\omega_u-1\|\le e^{\|\log\mathfrak{u}\|+\|Y_u\|}-1\le e^{2\varsigma}-1\le2.03\,\varsigma .
\end{equation*}
A unitary $\omega$ with $\|\omega-1\|=r\le0.03$ has $\|\log\omega\|\le2\arcsin(r/2)\le1.01r$;
therefore $\|\log\omega_u\|\le2.1\varsigma$. Further $\|Y'\|=\operatorname{pol}(g')<
\varepsilon^{\mathrm{fr}}(W)$. All maps used below are products of at most three maps
$\rho_{\mathsf{r}}[\cdot]$ at the one decomposition $\mathsf{r}$, of total size at most
$\varepsilon^{\mathrm{fr}}(W)+(1.5+2.1)\varsigma\le1.23\,\varepsilon^{\mathrm{fr}}(W)$ (the norm
$|\cdot|$ being the operator norm), where the coefficient $1.5$ bounds $|\zeta|$ on the
rectangle of Step B below. By clause (i) of the frame lemma, in the form of its corollary
\eqref{4.22:eq-frame-proof-increments-c}, applied once per factor with the same
$(U_{\mathbf{K}},\{u_{\square}\})$, these maps cost at most
$1.23\cdot\tfrac14\varpi_{\mathrm{fr}}r_a<\varpi_{\mathrm{fr}}r_a$ at every clause. Hence,
starting from the stored domain, one stays in $\mathfrak{U}^{\mathrm{tr}}_v$, and the unitary
part is $U_{\mathbf{K}}$ throughout. The computations of clause (i) of the frame lemma are
algebra valid for $\mathfrak{g}^{\mathbb{C}}$-valued arguments $Z$, $i\mathfrak{g}$-valued ones
included: a $G$-valued regular regauge near $1$ is read as a real increment of $A'$.

\emph{Step A: the unitary part of the riding map.} Put
$\hat\omega:=\rho_{\mathsf{r}}[\log\omega_u]^{-1}$. Since $\log\omega_u$ is
$i\mathfrak{g}$-valued, $\hat\omega$ is a $G$-valued site map; since
$\rho_{\mathsf{r}}[Z](s)=e^{-Z_s}$ at $s\in\Sigma_W$, it equals $\omega_u$ on $\Sigma_W$, and it
equals $1$ off $\mathsf{B}_W$. By the covariance of the orbit datum under $G$-valued gauge
transformations, together with clause (iii) of the frame lemma (no law site of $v$ other than
those of $\Sigma_W$ meets $\mathsf{B}_W$),
\begin{equation*}
v^{\sharp}(\mathsf{R};\Phi_W,\omega_u\tilde c)=v^{\sharp}(\tilde P;\Phi_W,\tilde c),\qquad
\tilde P:=\mathsf{R}^{\rho_{\mathsf{r}}[\log\omega_u]},\quad
\tilde c:=e^{Y_u}e^{Y'}\omega',
\end{equation*}
and $\tilde P\in\mathfrak{U}_v$ because $\mathfrak{U}_v$ is invariant under $G$-valued gauge
transformations. Note that $\mathfrak{u}g'=\omega_ue^{Y_u}e^{Y'}\omega'=\omega_u\tilde c$.

\emph{Step B: the Hermitian part.} Let $\delta_0\in(0,1/10]$ be the number fixed below, and put
$\mathsf{Rec}:=\{\zeta:-\delta_0<\operatorname{Re}\zeta<1+\delta_0,\ |\operatorname{Im}\zeta|<1\}$
and
\begin{equation*}
\varphi_1(\zeta):=v^{\sharp}(\tilde P;\Phi_W,e^{\zeta Y_u}e^{Y'}\omega'),\qquad
\varphi_2(\zeta):=v\Bigl(\bigl(\tilde P^{\rho_{\mathsf{r}}[\zeta Y_u]}\bigr)^{\rho_{\mathsf{r}}[Y']};
\Phi_W^{[\omega']}\Bigr).
\end{equation*}

\emph{$\varphi_1$ is holomorphic.} Write $\zeta=t+i\sigma$. Since $Y_u$ is Hermitian,
$e^{\zeta Y_u}=e^{i\sigma Y_u}e^{tY_u}$ with $e^{i\sigma Y_u}$ unitary; unitary factors on the
left or on the right do not change $\operatorname{pol}$, so
$\operatorname{pol}(e^{\zeta Y_u}e^{Y'}\omega')=\operatorname{pol}(e^{tY_u}e^{Y'})$. Since
$\operatorname{pol}(g^{*})=\operatorname{pol}(g)$ and $Y_u$, $Y'$ are Hermitian, this equals
$\operatorname{pol}(e^{Y'}e^{tY_u})$, which is a convex function of $t$ by the convexity of
$\operatorname{pol}$ along such products. At $t=0$ it equals $\operatorname{pol}(g')$, and at
$t=1$ it equals $\operatorname{pol}(e^{Y_u}e^{Y'})=\operatorname{pol}(\mathfrak{u}g')$; both are
$<\varepsilon^{\mathrm{fr}}(W)$ by \eqref{4.22:eq-kept-oscillation-b}. By convexity it is
$<\varepsilon^{\mathrm{fr}}(W)$ on $[0,1]$, and, the condition being open, on
$]-\delta_0,1+\delta_0[$ for some $\delta_0\in(0,1/10]$; this fixes $\delta_0$, and
$|\zeta|<1.5$ on $\mathsf{Rec}$. By the holomorphy of the orbit datum on the chart,
$\varphi_1$ is holomorphic on $\mathsf{Rec}$.

\emph{$\varphi_2$ is holomorphic.} At the frozen decomposition $\mathsf{r}$,
$\rho_{\mathsf{r}}[\zeta Y_u]$ is holomorphic in $\zeta$; its size is
$|\zeta|\|Y_u\|\le1.5\varsigma$, so by the Sizes paragraph the point stays in
$\mathfrak{U}^{\mathrm{tr}}_v$, where $v$ is holomorphic.

\emph{They agree on the segment $\zeta=it$, $|t|<1$.} There $e^{itY_u}\in G$ and
$\rho_{\mathsf{r}}[itY_u]$ is $G$-valued. By the covariance of the orbit datum under
$G$-valued gauge transformations,
$\varphi_1(it)=v^{\sharp}(\tilde P^{\rho_{\mathsf{r}}[itY_u]};\Phi_W,e^{Y'}\omega')$, and the
definition of the orbit datum in clause (ii) of the frame lemma, at the polar form
$e^{Y'}\omega'$ and with the decomposition $\mathsf{r}$ (the definition does not depend on the
decomposition), gives $\varphi_2(it)$. Since $\mathsf{Rec}$ is connected and the segment has
accumulation points in it, $\varphi_1\equiv\varphi_2$ on $\mathsf{Rec}$. At $\zeta=1$, using
$(V^a)^b=V^{ba}$,
\begin{equation*}
\varphi_2(1)=v\bigl(\tilde P^{\rho_{\mathsf{r}}[Y']\rho_{\mathsf{r}}[Y_u]};\Phi_W^{[\omega']}\bigr)
=v(Q(\tau);\Phi_W^{[\omega']}),\qquad
\varphi_1(1)=v^{\sharp}(\tilde P;\Phi_W,\tilde c),
\end{equation*}
and Step A identifies $\varphi_1(1)$ with $v^{\sharp}(\mathsf{R}(\tau);\Phi_W,\mathfrak{u}(\tau)g')$.
This is \eqref{4.22:eq-kept-oscillation-3}; the membership of $Q(\tau)$ in
$\mathfrak{U}^{\mathrm{tr}}_v$ and its decomposition with unitary part $U_{\mathbf{K}}$ are the
Sizes paragraph.

\emph{Clause (ii$_{\mathsf{Q}}$).} Apply \textup{(KH8)} at $(\mathbf{U},\mathbf{J}):=Q(\tau)$, in the
decomposition with $U=U_{\mathbf{K}}$ furnished by clause (i); the increments of Steps A and B are
real and complex parts of $A'$ over the same unitary part. The datum $\Phi_W^{[\omega']}$ is
real, the supports being $\operatorname{Ad}_G$-stable. Together with
\eqref{4.22:eq-kept-oscillation-3} this gives \eqref{4.22:eq-kept-oscillation-4}. The rotation
$\omega'$ is read off $g'$ and $U_{\mathbf{K}}$ off the base, so the subtracted number depends
neither on $\tau$ nor on $\mathfrak{u}$; at fixed $g'$ the triangle inequality gives
$\operatorname{osc}_{\mathfrak{P}}(\mathsf{Q}_i^{\sharp})\le2\Sigma^{\mathfrak{q}}(x_{k_W})$. By the
lower bound of the coupling flow (Theorem~\ref{4.1:thm-clause-zero}), for $k\ge k_W$,
\begin{equation*}
x_{k_W}\ge x_k+\lambda_\sharp(k-k_W)-c_5\ge x_k-c_5,
\end{equation*}
and $\Sigma^{\mathfrak{q}}$ is non-increasing, whence the last inequality.

\emph{Clause (iii).} Insert \eqref{4.22:eq-kept-oscillation-2}, respectively the bound on
$\operatorname{osc}_{\mathfrak{P}}(\mathsf{Q}_i^{\sharp})$ of clause (ii$_{\mathsf{Q}}$), into the
primed form of the third majorant bound among \eqref{4.22:eq-majorant-transfer-a}.

\emph{Clause (iv).} Let the data $(\mathbb{W},\mathbb{J},\mathsf{u}';\mathcal{K}',\mathfrak{j},\mathcal{Z})$,
the pairs $P_\lambda$, the rotations $\mathsf{u}_\lambda$ and the radius $\rho_X$ be those of the
statement of clause (iv), for a pair $\mathfrak{z}$, $\mathfrak{z}'$ with the same $\varpi$.

If $\rho_X\le1$: by \textup{(KH4)} both points lie in
$\mathfrak{U}_v\times\mathfrak{G}_{\varepsilon^{\mathrm{fr}}(W)}(\Sigma_W)$, so
$|F_v^{\mathfrak{b}}(\mathfrak{z})-F_v^{\mathfrak{b}}(\mathfrak{z}')|\le\mathfrak{O}(v)\le
\mathfrak{O}(v)/\rho_X$.

Let $\rho_X>1$. By \textup{(KH6)}, $P_\lambda\in\mathfrak{U}_v$ for $|\lambda|\le\rho_X$; by
\textup{(KH9)}, $\operatorname{pol}(\mathsf{u}_0(s))\le\eta_0(s)$ and
$|\lambda|\|\mathcal{Z}(s)\|\le\vartheta^{3F}(s)\le\tfrac12\mu(s)$. By the domain statement
\textup{(KH11)}, the function
$\Gamma(\lambda):=v^{\sharp}(P_\lambda;\Phi_W,\mathsf{u}_\lambda g')$ is holomorphic on
$|\lambda|<\rho_X$, and $(P_\lambda,\mathsf{u}_\lambda g')\in\mathfrak{U}_v\times
\mathfrak{G}_{\varepsilon}(\Sigma_W)$ with $\varepsilon\le\varepsilon^{\mathrm{fr}}(W)$. By the
definition \eqref{4.22:eq-kept-oscillation-a}, $|\Gamma(\lambda)-\Gamma(0)|\le\mathfrak{O}(v)$ on
the disc. Schwarz's lemma applied to $\Gamma-\Gamma(0)$ on the disc of radius $\rho_X$, at
$\lambda=1$, gives
\begin{equation*}
|\Gamma(1)-\Gamma(0)|\le\frac{\mathfrak{O}(v)}{\rho_X}
=\frac{1}{r^{0}}e^{-\frac12\delta_PD_X}\mathfrak{O}(v)
=\tfrac12C_{3F}e^{-\frac12\delta_PD_X}\mathfrak{O}(v)=\mathfrak{n}_X,
\end{equation*}
and $\Gamma(1)-\Gamma(0)=F_v^{\mathfrak{b}}(\mathfrak{z};g')-F_v^{\mathfrak{b}}(\mathfrak{z}';g')$. In
both cases the first bound of \eqref{4.22:eq-kept-oscillation-5} holds.

\emph{Pieces.} Apply the piece expansion of the three-factor bounds without change, with
constant term $\mathsf{c}:=F_v^{\mathfrak{b}}((\zeta,\varpi,\vec{0});g')$ and bound
$\mathfrak{n}_X$; the holomorphy on the product polydisc is \textup{(KH4)} at this presentation
site. This gives the bound on $|\mathfrak{p}|$ in \eqref{4.22:eq-kept-oscillation-5}. For the
$\zeta$-difference take the function
\begin{equation*}
\Delta(\vec{\mathsf{s}}):=F_v^{\mathfrak{b}}((\zeta,\varpi,\vec{\mathsf{s}});g')
-F_v^{\mathfrak{b}}((\zeta',\varpi,\vec{\mathsf{s}});g'),
\end{equation*}
which satisfies $|\Delta|\le\mathfrak{n}_X$ on the polydisc by the first bound; Cauchy's estimate
in each cell gives the bound on $|\mathfrak{p}(\zeta)-\mathfrak{p}(\zeta')|$, the empty tuple
giving $\Delta(\vec{0})$.

\emph{The depth $D_X$.} The two lower bounds in \eqref{4.22:eq-kept-oscillation-6} are the two
cases of the depth statement for $C$; they rest on the following. If $k_W=h$, then
$\operatorname{dist}_h(Y_i^{\sharp},\mathsf{C}_C^{c})\ge\tfrac12(\check{R}_h+1)$, and all cells of
$\mathsf{C}_C$ have level at most $h$, by the geometric lemma on the cells of $\mathsf{C}_C$. If
$k_W<h$, the own collar of $Y_i$, given by the collar statement of the geometric lemma, lies in
the complement, inside $C$, of the small-field region of level $k_W+1$, hence in
$Z_h\subset\mathsf{C}_C$, and its cells have level at most $k_W$. On a set whose cells have level
at most $n$, a contour pays at least $\mathsf{w}_{\wedge}$ per unit of level $n$ in the distance
$d^{\wedge}$; this gives \eqref{4.22:eq-kept-oscillation-6}.
\end{proof}

\begin{remark}[Remarks on the rotation bound and the two-point oscillation bound]\label{4.22:rem-two-point-remarks}
The notation is that of Lemma~\ref{4.22:lem-kept-oscillation} (oscillation of the kept terms on the
orbit chart): $v\in\mathfrak{K}(Y_i)$, $W:=Y_i$, $\Sigma_W$ the law sites of the frame $W$,
$\varepsilon^{\mathrm{fr}}(W)$ the frame radius, $\mathfrak{u}$ the riding map, and $\varsigma(s)$ the
number bounding $|\log\mathfrak{u}(\tau)(s)|$ in the rotation bound
\eqref{4.22:eq-kept-oscillation-b}.
\begin{enumerate}
\item[(1)] The rotation bound \eqref{4.22:eq-kept-oscillation-b} is a statement about the riding map
$\mathfrak{u}$, not about the kept member $v$. It is used only in the representation clause~(i) of
Lemma~\ref{4.22:lem-kept-oscillation} and, through it, in the form clause~(ii)$_{\mathsf{Q}}$. It
holds at every site of a step at which that lemma is applied, with
\begin{equation}\label{4.22:eq-two-point-remarks-1}
\varsigma(s)\le 2(x_k-c_5)^{-2}\,\varepsilon^{\mathrm{fr}}(W)\le \tfrac{2}{32}\,
\varepsilon^{\mathrm{fr}}(W).
\end{equation}
The factor needed to absorb it is already present in the weighted contour distance $d^{\wedge}$, and
the bound does not accumulate from one step to the next.
\item[(2)] Under the hypotheses \textup{[V]}, \textup{(q0)}, \textup{(q1)} and $(\mathrm{hol})_v$ of
Lemma~\ref{4.22:lem-kept-oscillation}, the oscillation clause of the strengthened sup bound for the
kept members holds on the chart, at each of the three kinds of site $(\mathrm{T})$, $(\mathrm{P})$,
$(\mathrm{L})$ of a step, with the variation of the riding map included. The far factor of a piece
is the constant $c_\pi$ of the family in question. The number $\mathfrak{n}^{\mathrm{kept}}$ depends
neither on $\check R_{k_W}$, nor on the depth $D$, nor on the age $k-k_W$. The source-tube bound
itself remains a hypothesis.
\item[(3)] The bound \cite[(1.69)]{B-CMP122b} is what clause (g) of the statement on the kept lineage
states of $\mathfrak{K}(Y_i)$ at the step $k_W$ at which the region is created. It is used there by
the modification statement, and by no later step.
\item[(4)] Under \cite[(1.104)]{B-CMP122b} the kept lineage $\mathfrak{K}(Y_i)$ is empty.
\item[(5)] The decomposition clause is used, at the real regauges in the proof of clause~(i) of
Lemma~\ref{4.22:lem-kept-oscillation}, only for the bound on the kept quadratic form. The two-point
clause~(ii)$_2$ of that lemma does not use it.
\item[(6)] The bound $C_\Sigma\check R_{k_W}^{p_\Sigma}$ of Lemma~\ref{4.22:lem-kept-oscillation} is not
rewritten in terms of the current coupling $x_k$.
\end{enumerate}
\end{remark}

\begin{proof}
\emph{Item~(1).} The display \eqref{4.22:eq-kept-oscillation-b} bounds $|\log\mathfrak{u}(\tau)(s)|$
and the polar sizes $\operatorname{pol}(g'(s))$ and $\operatorname{pol}(\mathfrak{u}(\tau)(s)g'(s))$.
Each of these quantities is attached to the riding map and to the chart point $g'$, and none of them
involves $v$. In Lemma~\ref{4.22:lem-kept-oscillation} the display enters only as a hypothesis of
clause~(i), and through clause~(i) as a hypothesis of clause~(ii)$_{\mathsf{Q}}$. Clause~(ii)$_2$
uses only the membership of the two points in
$\mathfrak{U}_v\times\mathfrak{G}_{\varepsilon^{\mathrm{fr}}(W)}(\Sigma_W)$.

We now check that the display holds at every kind of site.

The proofs of clause~(iii) of the presentation-site lemma and of its version on the second-class
families bound $|\log\mathfrak{u}|$ and not only $\operatorname{pol}(\mathfrak{u})$.
Clause~(i) of the riding-map lemma gives, at a law site $z$,
$|\log\mathfrak{u}(z)|\le K_u\Theta_0^{\mathrm{loc}}$. For a product $\prod_s\mathfrak{w}_s$ of
presenter maps, the Baker--Campbell--Hausdorff formula at sizes at most $10^{-2}$ costs at most a
factor $1.01$, and this factor is absorbed in the factor $e^{2}$ contained in $C_\natural$.

By the ceiling \eqref{4.22:eq-coupling-ceilings-k}, one reading contributes at most
$\mu_k\mathsf{e}(s)/\varepsilon_\star$ to $|\log\mathfrak{u}(\tau)(s)|$. By the led-rotation bounds for
the increment discs, the increment disc adds at most $\mu_k\mathsf{e}(s)/\varepsilon_\star$ more. Since
the margin is $\mu_k=\varepsilon_\star(x_k-c_5)^{-2}$, this gives
\begin{equation*}
\varsigma(s)\le \frac{2\mu_k\mathsf{e}(s)}{\varepsilon_\star}=2(x_k-c_5)^{-2}\,\mathsf{e}(s).
\end{equation*}
By the definition of the unit $\mathsf{e}(s)$ at the law sites of a frame, $\mathsf{e}(s)$ is at most
the frame radius $\varepsilon^{\mathrm{fr}}(W)$ at the law sites of $W$, so that
$\varsigma(s)\le 2(x_k-c_5)^{-2}\varepsilon^{\mathrm{fr}}(W)$, and the ceiling
\eqref{4.22:eq-coupling-ceilings-c} gives
$2(x_k-c_5)^{-2}\varepsilon^{\mathrm{fr}}(W)\le\frac{2}{32}\varepsilon^{\mathrm{fr}}(W)$. This is
\eqref{4.22:eq-two-point-remarks-1}. Since $\frac{2}{32}=\frac1{16}$, it is the bound
$\varsigma(s)\le\varepsilon^{\mathrm{fr}}(W)/16$ required in \eqref{4.22:eq-kept-oscillation-b}.

The factor that absorbs the display is already present in $d^{\wedge}$. Write $B^{\mathrm{top}}(s)$
for the top block attached to the law site $s$ in the version of clause~(iii) of the
presentation-site lemma on the second-class families, and $\Delta_{\mathrm{src}}$ for the lattice set
that supports the source term. On these families the displays of that clause carry the factor
\begin{equation*}
e^{-\delta_P\, d^{\wedge}(B^{\mathrm{top}}(s),\,\Delta_{\mathrm{src}})}.
\end{equation*}
Each zone crossed between $B^{\mathrm{top}}(s)$ and $\Delta_{\mathrm{src}}$ contributes
$\mathsf{w}_{\wedge}$ to $d^{\wedge}$, where $\mathsf{w}_{\wedge}$ is the amount that $d^{\wedge}$
charges per zone. This decay is set against the factor $(1+\mathsf{a}_k(s))^{1/2}$ of the age
weight.

Finally, the rotation bound is imposed site by site. The term that survives a step is $v$ at the new
base, and there $\mathfrak{u}=1$. Hence the bound \eqref{4.22:eq-kept-oscillation-b} is applied afresh
at each step and never accumulates.

\emph{Item~(2).} Under the hypotheses \textup{[V]}, \textup{(q0)}, \textup{(q1)} and
$(\mathrm{hol})_v$ of Lemma~\ref{4.22:lem-kept-oscillation}, clause~(ii)$_2$ of
Lemma~\ref{4.22:lem-kept-oscillation} bounds the oscillation $\operatorname{osc}_{\mathfrak{P}}(v^\sharp)$
over the family by
$\mathfrak{O}(v)\le\mathfrak{o}(v)\le C_\Sigma\check R_{k_W}^{p_\Sigma}$; see
\eqref{4.22:eq-kept-oscillation-a}. The bound holds on the chart and at each of the sites
$(\mathrm{T})$, $(\mathrm{P})$, $(\mathrm{L})$.

The variation of the riding map is included in this bound. By the definition of the frame datum
$v^\sharp$ at $\Sigma_W$, an increment at $\Sigma_W$ is an increment of $v$ in its pair argument $P$,
and $\mathfrak{O}(v)$ takes the supremum over all such pairs.

By clause~(iii) of the same lemma, a piece is bounded by $c_\pi\,\mathfrak{O}(v)$. So the far factor
of a piece is the constant $c_\pi$ of the family in question. The corollary on the kept members' share
of the far factor shows, through a displayed share of that far factor, that $\mathfrak{n}^{\mathrm{kept}}$
is a number that depends neither on $\check R_{k_W}$, nor on the depth $D$, nor on $k-k_W$.

None of these arguments proves the source-tube bound. It remains a hypothesis.

\emph{Item~(3).} Clause (g) of the statement on the kept lineage states \cite[(1.69)]{B-CMP122b} for
$\mathfrak{K}(Y_i)$ at the step $k_W$ at which the region is created. The modification statement uses
it at that step, and no later step uses it. The modification statement carries the bound at creation
forward.

At the healing step, the fully decoupled value of the member is a constituent of the new term
$V(C^{\sim})$, where $C\ni Y_i$ is the component treated by the large-field operation at the healing
step. The size of $V(C^{\sim})$ is again given by the count of the sourced form of
\cite[(1.69)]{B-CMP122b}.

\emph{Item~(4).} Under \cite[(1.104)]{B-CMP122b} there is no kept lineage: $\mathfrak{K}(Y_i)$ is
empty.

\emph{Item~(5).} In the proof of clause~(i) of Lemma~\ref{4.22:lem-kept-oscillation}, the two steps
treating the unitary and the Hermitian parts of the riding map perform regauges, and the increments
these regauges produce are real and complex parts of the $\mathfrak{g}^{\mathbb{C}}$-valued part $A'$
of the decomposition of the base (Lemma~\ref{4.22:lem-kept-oscillation}) over the same unitary part.
The decomposition clause is what allows the bound on the kept quadratic form to be applied after these
regauges, with their increments interpreted in this way.

Suppose the form bound were stated only for purely imaginary $A'$. Then, for a $G$-valued site map
$\hat\omega$, the difference
\begin{equation*}
\mathsf{Q}(U^{\hat\omega^{-1}},0;\Phi)-\mathsf{Q}(U,0;\Phi)
\end{equation*}
would be left without a bound.

For clause~(ii)$_2$ nothing of this kind is needed. In clause~(o) of
Lemma~\ref{4.22:lem-kept-oscillation}, the whole twist, unitary part included, is moved into the pair
by an arbitrary $G$-valued site map $\hat\omega$. Moreover, the two-point bound does not ask where in
the tube the two points lie.

\emph{Item~(6).} One could combine $\check R_{k_W}<L(\log x_{k_W})^r$ with
$x_{k_W}\le x_k+B_\sharp(k-k_W)$ to bound $C_\Sigma\check R_{k_W}^{p_\Sigma}$. The result would be
\begin{equation*}
C_\Sigma L^{p_\Sigma}\bigl(\log(x_k+B_\sharp(k-k_W))\bigr)^{rp_\Sigma}.
\end{equation*}
This majorant depends explicitly on the age $k-k_W$, which no right-hand side of the estimates of
this section is allowed to do; the substitution is therefore not made.
\end{proof}

The pieces of the Gaussian kept under the large-field operation $\mathbf{T}'_{k_W}(W)$ at an
arrested chain are estimated through the oscillation of their kernel when the background pair
moves, whereas the slot is read at an argument rotated by a block-constant rotation, which is not
a small change of the pair. The following definition therefore fixes the slot, the pieces, their
sizes, and the bound on their oscillation at a real unrotated slot.

\begin{definition}[The kept Gaussian form of an arrested chain]\label{4.22:not-kept-gaussian}
Let $\mathfrak{a}=(W,k_W,n+1)$ be an arrested chain of the fourth kind with frame $W$. Let
$j:=h+n$ be the level
of $\mathfrak{a}$. Here $h$ and $n$ are the integers fixed with the arrested chain $\mathfrak{a}$,
and $h$ does not denote a history.

\begin{enumerate}
\item[(a)] \emph{The slot.} The slot is
\begin{equation*}
\phi:=B_j\restriction S^P_W ,
\end{equation*}
the field $B_j$ of \cite[(1.60)]{B-CMP122a} restricted to the set $S^P_W$ attached to the frame
$W$ at the large-field stages $(\mathrm{P})$ of the step. It is written in the coordinates fixed
for the slots of the large-field stages, and its value at a bond $b$ has law site $b_-$. Its real
range is
\begin{equation}\label{4.22:eq-kept-gaussian-1}
|\phi(b)|\le \mathfrak{b}_W:=\mathfrak{b}^B_j=\frac{6\delta'_j}{g_j},
\end{equation}
where $\mathfrak{b}^B_j$ is the real range of a weak collar variable at level $j$, and
$\delta'_j$ is the value at level $j$ of Ba{\l}aban's parameter $\delta'$ of the large-field
construction \cite{B-CMP122a}.

\item[(b)] \emph{The kernel.} For a complex parameter $\sigma$ and a pair
$P=(\mathbb{U},\mathbb{J})$, put
\begin{equation*}
\mathcal{A}(\sigma;\mathbb{U},\mathbb{J}):=C(\mathbb{U})^{*}\,
\Delta^{(j)}(\sigma;\mathbb{U},\mathbb{J})\,C(\mathbb{U}),
\end{equation*}
where $C(\mathbb{U})$ and $\Delta^{(j)}(\sigma;\mathbb{U},\mathbb{J})$ are the operators of
\cite[(1.60)]{B-CMP122a}.
This is the kernel of the quadratic form
$-\tfrac12\langle B_j,C^{*}\Delta^{(j)}CB_j\rangle$ of \cite[(1.60)]{B-CMP122a}. For arrested
chains of the fourth kind the kept form is this one, and not the form of
\cite[(1.71)]{B-CMP122b}.

\item[(c)] \emph{The pieces.} A piece is a triple $\mathsf{p}=(q_1,q_2,\mathsf{y})$ of the
expansion of $\mathcal{A}$ into families. Here $q_1,q_2$ are the sets of bonds over which the
piece sums, and $\mathsf{y}$ is its family, of length $|\mathsf{y}|$. We write
$\mathcal{A}_{\mathsf{y}}$ for the corresponding summand of the kernel, $S_{\mathsf{p}}$ for the
support of the piece and $n_q$ for the number entering its size. The piece of the kept
Gaussian and its size are
\begin{align*}
\mathsf{q}_{\mathsf{p}}(\phi;P)&:=-\frac12\sum_{b\in q_1,\ b'\in q_2}
\bigl\langle\phi(b),\mathcal{A}_{\mathsf{y}}[P](b,b')\,\phi(b')\bigr\rangle,\\
\bar{\mathsf{q}}_{\mathsf{p}}&:=\frac12\,\mathfrak{b}_W^{2}\,n_q\,B_{\mathcal{A}}\,
e^{-(\kappa_1-1)|\mathsf{y}|},
\end{align*}
where $B_{\mathcal{A}}$ is the constant of the kernel bound in (d).

\item[(d)] \emph{Hypothesis on the kernel.} The following properties are taken as stated in
clauses (i), (iii) and (iv) of the lemma on the random-walk expansion of the kept kernel at the
large-field stages, with constants
$B_{\mathcal{A}}$ and $C'_{\mathcal{A}}$. On the domain $\mathfrak{K}_P$ of background pairs $P$
fixed with the kept members, times the disc $\{|\sigma|\le e^{\kappa_1}\}$, the kernel
$\mathcal{A}_{\mathsf{y}}$
\begin{itemize}
\item is holomorphic;
\item defines a symmetric bilinear form;
\item is $\operatorname{Ad}$-covariant, that is, for every gauge transformation $v$,
\begin{equation}\label{4.22:eq-kept-gaussian-2}
\mathcal{A}[\mathbb{U}^v,\operatorname{Ad}_v\mathbb{J}](b,b')
=\operatorname{Ad}_{v(b_-)}\,\mathcal{A}[\mathbb{U},\mathbb{J}](b,b')\,
\operatorname{Ad}_{v(b'_-)}^{-1};
\end{equation}
\item satisfies the bound
\begin{equation*}
|\mathcal{A}_{\mathsf{y}}(b,b')|\le B_{\mathcal{A}}\,e^{-(\kappa_1-1)|\mathsf{y}|}.
\end{equation*}
\end{itemize}
Moreover, let $P,P'$ be two pairs presented, in the sense of the presentation of the step, by an
operation $\mathfrak{O}$ on the presentation set $\widetilde{S}_{\mathsf{p}}$ of the piece (a set
distinct from its support $S_{\mathsf{p}}$).
Let $\Delta_{\mathfrak{O}}(S_{\mathsf{p}})$ be the increment of $\mathfrak{O}$ on
$S_{\mathsf{p}}$, measured in threshold units. Let $\bar\alpha$ be evaluated at the inverse
coupling $x_\ell=g_\ell^{-2}$ of the level $\ell=\ell(\mathsf{y})$ of the family itself, not at
the level $j$ of the chain. Then
\begin{equation}\label{4.22:eq-kept-gaussian-3}
\bigl|[\mathcal{A}_{\mathsf{y}}[P]-\mathcal{A}_{\mathsf{y}}[P']](b,b')\bigr|
\le C'_{\mathcal{A}}\,\bar\alpha(x_\ell)\,\Delta_{\mathfrak{O}}(S_{\mathsf{p}})\,
e^{-(\kappa_1-1)|\mathsf{y}|}.
\end{equation}

\item[(e)] \emph{The induced bound.} Put
\begin{equation}\label{4.22:eq-kept-gaussian-4}
\mathfrak{N}_{\mathfrak{O}}(\mathsf{q}_{\mathsf{p}}):=\frac12\,n_q\,C'_{\mathcal{A}}\,
\mathfrak{b}_W^{2}\,\bar\alpha(x_\ell)\,\Delta_{\mathfrak{O}}(S_{\mathsf{p}})\,
e^{-(\kappa_1-1)|\mathsf{y}|}.
\end{equation}
Then $\mathfrak{N}_{\mathfrak{O}}(\mathsf{q}_{\mathsf{p}})$ bounds the oscillation of
$\mathsf{q}_{\mathsf{p}}$ between $P$ and $P'$ at a real unrotated slot $\phi$ obeying
\eqref{4.22:eq-kept-gaussian-1}:
\begin{equation*}
|\mathsf{q}_{\mathsf{p}}(\phi;P)-\mathsf{q}_{\mathsf{p}}(\phi;P')|
\le\mathfrak{N}_{\mathfrak{O}}(\mathsf{q}_{\mathsf{p}}).
\end{equation*}
This is clause (ii) of the oscillation lemma for kept members.

The shape of \eqref{4.22:eq-kept-gaussian-4} is read off from the definitions. The slot $\phi$ is
the same in both terms and $\mathsf{q}_{\mathsf{p}}$ is linear in the kernel, so
\begin{equation*}
\mathsf{q}_{\mathsf{p}}(\phi;P)-\mathsf{q}_{\mathsf{p}}(\phi;P')
=-\frac12\sum_{b\in q_1,\ b'\in q_2}
\bigl\langle\phi(b),[\mathcal{A}_{\mathsf{y}}[P]-\mathcal{A}_{\mathsf{y}}[P']](b,b')\,
\phi(b')\bigr\rangle .
\end{equation*}
In \eqref{4.22:eq-kept-gaussian-4} the factor $\mathfrak{b}_W^{2}$ corresponds to the two slot
values $\phi(b)$, $\phi(b')$, each bounded by \eqref{4.22:eq-kept-gaussian-1}; the factors
$C'_{\mathcal{A}}\,\bar\alpha(x_\ell)\,\Delta_{\mathfrak{O}}(S_{\mathsf{p}})\,
e^{-(\kappa_1-1)|\mathsf{y}|}$ are those of \eqref{4.22:eq-kept-gaussian-3}; and the factor
$\tfrac12 n_q$ corresponds to the prefactor $-\tfrac12$ and the sum over
$(b,b')\in q_1\times q_2$. Since the real range \eqref{4.22:eq-kept-gaussian-1} of the slot
enters, the bound is stated for a real unrotated slot.
\end{enumerate}

The letter $\mathfrak{O}$ here denotes an operation and occurs only as a subscript. It is
distinct from the oscillation seminorm $\mathfrak{O}(v)$ of a kept member, which always carries
an argument.
\end{definition}

\begin{proposition}[Oscillation of the kept Gaussian form under rotation]
\label{4.22:prop-gaussian-oscillation}
Let $\mathfrak{a}$ be an arrest of level $j$, and let the slot $\phi$, the law sites $b_-$ of its
coordinates, the slot bound $\mathfrak{b}_W$, the kernel $\mathcal{A}$ with its decoupling parameter
$\sigma$, the pieces
$\mathsf{p}=(q_1,q_2,\mathsf{y})$ with their sets $S_{\mathsf{p}}$, kernels $\mathcal{A}_{\mathsf{y}}$
and counts $n_q$, the form $\mathsf{q}_{\mathsf{p}}(\phi;P)$ and its size
$\bar{\mathsf{q}}_{\mathsf{p}}$ be as in Notation~\ref{4.22:not-kept-gaussian}. Let
$\widetilde{S}_{\mathsf{p}}$ be the enlarged set attached to $S_{\mathsf{p}}$ and $\mathfrak{K}_P$ the
domain of background pairs $P=(\mathbb{U},\mathbb{J})$ on which the kept kernel is controlled. Assume
the following statements.
\begin{itemize}
\item[(a)] \emph{The bounds on the kept kernel.} On
$\mathfrak{K}_P\times\{|\sigma|\le e^{\kappa_1}\}$ the kernel $\mathcal{A}_{\mathsf{y}}$ is
holomorphic and symmetric bilinear; it is covariant,
\begin{equation}\label{4.22:eq-gaussian-oscillation-1}
\mathcal{A}[\mathbb{U}^v,\operatorname{Ad}_v\mathbb{J}](b,b')
=\operatorname{Ad}_{v(b_-)}\mathcal{A}[\mathbb{U},\mathbb{J}](b,b')\operatorname{Ad}_{v(b'_-)}^{-1},
\end{equation}
and $|\mathcal{A}_{\mathsf{y}}(b,b')|\le B_{\mathcal{A}}e^{-(\kappa_1-1)|\mathsf{y}|}$. For two pairs
$P,P'$ which an operation $\mathfrak{O}$ presents on $\widetilde{S}_{\mathsf{p}}$,
\begin{equation}\label{4.22:eq-gaussian-oscillation-2}
\bigl|[\mathcal{A}_{\mathsf{y}}[P]-\mathcal{A}_{\mathsf{y}}[P']](b,b')\bigr|
\le C_{\mathcal{A}}'\,\bar\alpha(x_\ell)\,\Delta_{\mathfrak{O}}(S_{\mathsf{p}})\,
e^{-(\kappa_1-1)|\mathsf{y}|},
\end{equation}
where $\Delta_{\mathfrak{O}}(S_{\mathsf{p}})$ is the increment of the operation in threshold units and
$\bar\alpha$ is taken at the own-level coupling $x_\ell$, $\ell=\ell(\mathsf{y})$.
\item[(b)] \emph{The lemma on the kept Gaussian form, clauses (i) and (ii).} The slot $\phi$ is
stored as a real variable in its box $|\phi(b)|\le\mathfrak{b}_W$ and is not itself complexified;
and for $P,P'$ as in (a) the oscillation of $\mathsf{q}_{\mathsf{p}}$ at a real unrotated slot is
bounded by
\begin{equation}\label{4.22:eq-gaussian-oscillation-3}
\mathfrak{N}_{\mathfrak{O}}(\mathsf{q}_{\mathsf{p}})
:=\tfrac12\,n_qC_{\mathcal{A}}'\,\mathfrak{b}_W^2\,\bar\alpha(x_\ell)\,
\Delta_{\mathfrak{O}}(S_{\mathsf{p}})\,e^{-(\kappa_1-1)|\mathsf{y}|}.
\end{equation}
\item[(c)] \emph{The holomorphy lemma for the riding map, clauses (i) and (ii).} At a law site $s$
the value $\mathfrak{u}(s)$ of the riding map is a holomorphic function of the layer and of the
complex part of the base, both restricted to the block $B^{\mathrm{top}}(s)$, on the class of
configurations of \cite[(158)]{B-CMP98}; the radius of this class in the complex part of the base
is the number $\alpha_1^{98}$ of \cite[Proposition~7]{B-CMP98}, which depends on $(d,L)$ only and
not on the coupling. This number is the small constant of \cite[(158)]{B-CMP98}; it is not the
radius $\alpha_1$, $\alpha_{1,k}$ or $\tilde\alpha_1$ of \cite{B-CMP119}, \cite{B-CMP122a} and
\cite{B-CMP122b}, and with it the number $\varrho_0$ of (f) does not depend on the coupling.
\item[(d)] \emph{Clause (iv) of the bound lemma for the riding map: the one-layer logarithm bound.}
For a layer $\mathcal{K}$ of a site, $\tau,\tau'$ in
the parameter polydisc $\mathfrak{P}$ and a law site $s$, put
$a_0:=\Theta_0^{\mathrm{loc}}(\mathcal{K}(\tau'))(s)$ and
$\xi:=\Theta_0^{\mathrm{loc}}(\mathcal{K}(\tau)-\mathcal{K}(\tau'))(s)$. If $a_0\le\varrho_0/3$,
$\xi\le\varrho_0/3$ and $a_0+\xi\le\varrho_0$, then
$|\log[\mathfrak{u}(\tau')^{-1}\mathfrak{u}(\tau)](s)|\le 2.1K_u\,\xi$.
\item[(d$'$)] \emph{The bound lemma for the riding map, clause (i), with the Baker--Campbell--Hausdorff
formula.} For a law site $s$, a layer $X_0$ with $\Theta_0^{\mathrm{loc}}(X_0)(s)\le a_0$ and two
bases $W$, $W'$ in the class of configurations of (c),
\begin{equation*}
\bigl|\log\bigl[\mathfrak{u}[e^{i\eta X_0}W;W]^{-1}\,\mathfrak{u}[e^{i\eta X_0}W';W']\bigr](s)\bigr|
\le 1.05\cdot 2K_u\,a_0;
\end{equation*}
and $\operatorname{pol}(\mathfrak{u}(\tau')(s))\le 10^{-2}$ at the law sites of the coordinates of
$\phi$ and for every $\tau'\in\mathfrak{P}$.
\item[(e)] \emph{The far-site lemma, clauses (i)--(iv).} The law sites of the coordinates of $\phi$
are never of the near kind $(\mathrm{n})$; their kinds are among $(\mathrm{T})$,
$(\mathrm{f}_1)$, $(\mathrm{f}_2)$, $(\alpha)$, $(\beta)$ (the class letters of the far-site
lemma); they lie at distance at least $\check{R}_k-1$ from every source.
At such a law site $s=b_-$ one has $h_q=1$, and the partial products of the links of the riding map
have polar size at most $\eta^L_k(s)$.
\item[(e$'$)] \emph{The layer bound at far law sites.} At every law site $s=b_-$ of the coordinates
of $\phi$, every layer $\mathcal{K}$ of the site and every $\tau\in\mathfrak{P}$ satisfy
\begin{equation}\label{4.22:eq-gaussian-oscillation-4}
\Theta_0^{\mathrm{loc}}(\mathcal{K}(\tau))(s)\le e^2B_H^S\,\bar\alpha(x_k)\,E(x_k)^{\delta_\natural}.
\end{equation}
\item[(f)] \emph{The coupling ceilings.} The ceilings \eqref{4.22:eq-coupling-ceilings-d} and
\eqref{4.22:eq-coupling-ceilings-f} hold, with $B_H$ in
\eqref{4.22:eq-coupling-ceilings-f} read as $B_H^S$ on the families of the wrapped class
$\mathfrak{S}$; moreover
$\eta^L_k(s)\le\mu_k$ and $\mu_k\le\varepsilon_\star/32$. In particular
$B_H\bar\alpha(x)\le\varrho_0$ and $B_H^S\bar\alpha(x)\le\varrho_0$, where
$\varrho_0:=\min\{10^{-2}/K_u,\,c_4,\,\alpha_1^{98}\}$, and
$E(x)^{\delta_\natural}\le(x-c_5)^{-4}\le 2^{-10}$. Further, the constant $B_H$ contains the factor
$e^{16\kappa_1}$, so that $B_H\ge 6$.
\item[(g)] \emph{The threshold bound for presented layers, first inequality.} For a layer (or an end
of a link) presented by $\mathfrak{O}$ on $\widetilde{S}_{\mathsf{p}}$ and a law site $s$ whose
$M$-cube lies in $\widetilde{S}_{\mathsf{p}}$, the increment between $\tau$ and $\tau'$ satisfies
\begin{equation}\label{4.22:eq-gaussian-oscillation-10}
\Theta_0^{\mathrm{loc}}(\cdot)(s)\le\sup_{y\in B^{\mathrm{top}}(s)}N_y(\cdot)
\le\bar\alpha(x_\ell)\Delta_{\mathfrak{O}}(S_{\mathsf{p}}),
\end{equation}
where, for a point $y$, $N_y(\cdot)$ denotes the pointwise norm at $y$ in which the threshold units
of the operation are read.
\end{itemize}
Then the following hold.
\begin{itemize}
\item[(i)] \emph{The datum.} The function
$\mathsf{q}_{\mathsf{p}}^{\sharp}(P;\phi,g):=\mathsf{q}_{\mathsf{p}}(\phi^{[g]};P)$ is the orbit
datum of $\mathsf{q}_{\mathsf{p}}$, of radius $\varepsilon_\star$ at the law sites of the bonds of
$q_1\cup q_2$, on $\mathfrak{K}_P\restriction\widetilde{S}_{\mathsf{p}}$, hence on every tube and for
every further argument and every source admitted there; and its chart supremum at radius
$\varepsilon_\star$ satisfies
\begin{equation}\label{4.22:eq-gaussian-oscillation-5}
\mathfrak{N}^{\mathbb{C}}_{\varepsilon_\star}(\mathsf{q}_{\mathsf{p}})
\le e^{4\varepsilon_\star}\bar{\mathsf{q}}_{\mathsf{p}}\le 1.013\,\bar{\mathsf{q}}_{\mathsf{p}}.
\end{equation}
\item[(ii)] \emph{Oscillation with the riding map.} Let a site of the step present
$\mathsf{R}(\tau)\in\mathfrak{K}_P$ and the riding map $\mathfrak{u}(\tau)$, $\tau\in\mathfrak{P}$.
Let $g'$ satisfy $\operatorname{pol}(\mathfrak{u}(\tau)(s)g'(s))<\varepsilon_\star$ at every law site
$s$ and every $\tau\in\mathfrak{P}$, and suppose that for $\tau,\tau'\in\mathfrak{P}$ and every law
site $s$ the element $\mathfrak{u}(\tau)(s)\mathfrak{u}(\tau')(s)^{-1}$ is a product of at most
eight factors $e^{\tilde Z_i}$ with $\sum_i|\tilde Z_i|\le\mathsf{z}(s)\le 10^{-2}$, the number of
factors being one at the sites $(\mathrm{T})$ and $(\mathrm{P})$. Then on the output chart
\begin{equation}\label{4.22:eq-gaussian-oscillation-6}
\operatorname{osc}_{\mathfrak{P}}(\mathsf{q}_{\mathsf{p}}^{\sharp})
\le 1.013\,\mathfrak{N}_{\mathfrak{O}}(\mathsf{q}_{\mathsf{p}})
+4.12\,\sup_s\mathsf{z}(s)\,\bar{\mathsf{q}}_{\mathsf{p}}.
\end{equation}
\item[(iii)] \emph{The size of the riding map.} At the sites $(\mathrm{T})$ and $(\mathrm{P})$, where
there is one layer $\mathcal{K}$,
\begin{equation}\label{4.22:eq-gaussian-oscillation-7}
\mathsf{z}(s)\le 2.15K_u\,\Theta_0^{\mathrm{loc}}(\mathcal{K}(\tau)-\mathcal{K}(\tau'))(s)
\le 2.15K_u\,\bar\alpha(x_\ell)\,\Delta_{\mathfrak{O}}(S_{\mathsf{p}}).
\end{equation}
At the site $(\mathrm{L})$ the riding map is the product $\prod\mathfrak{w}_{s'}$ over at most four
links $s'$, not a single $\mathfrak{u}[\mathcal{K}]$, and the net increment of the presented pair may
cancel between links where the increment of the riding map does not. There
$\Delta_{\mathfrak{O}}(S_{\mathsf{p}})$ is taken per link and layer: it is the largest, over the at
most four links $s'$, of the increments between $\tau$ and $\tau'$, in threshold units on
$\widetilde{S}_{\mathsf{p}}$, of the layer $\mathbb{H}_{s'}$ of the link and of its two ends
$W_{s'}$ and $e^{i\eta\mathbb{H}_{s'}}W_{s'}$. Since the presented pair is such an end, the bound
\eqref{4.22:eq-gaussian-oscillation-2} holds with $\Delta_{\mathfrak{O}}$ so taken, and
$\mathsf{z}(s)\le 8\cdot 2.15K_u\bar\alpha(x_\ell)\Delta_{\mathfrak{O}}(S_{\mathsf{p}})$. Hence at
every site of the step
\begin{equation}\label{4.22:eq-gaussian-oscillation-a}
\begin{gathered}
\operatorname{osc}_{\mathfrak{P}}(\mathsf{q}_{\mathsf{p}}^{\sharp})
\le\tfrac12\,n_q\,C_{\mathcal{A}}^{\sharp}\,\mathfrak{b}_W^2\,\bar\alpha(x_\ell)\,
\Delta_{\mathfrak{O}}(S_{\mathsf{p}})\,e^{-(\kappa_1-1)|\mathsf{y}|},\\
C_{\mathcal{A}}^{\sharp}:=1.013\,C_{\mathcal{A}}'+71.2\,K_uB_{\mathcal{A}}.
\end{gathered}
\end{equation}
This is the bound \eqref{4.22:eq-gaussian-oscillation-3} of clause (ii) of the lemma on the kept
Gaussian form with $C_{\mathcal{A}}'$ replaced by $C_{\mathcal{A}}^{\sharp}$; it has the same shape,
the coupling of the operation entering through the power $\bar\alpha(x_\ell)\propto x_\ell^{-1/2}$.
The ceiling on the kept constants is taken with the larger constant $C_{\mathcal{A}}^{\sharp}$, a
number fixed before $\gamma$ in the order of constants, in supremum form; its power of $x$ is
unchanged.
\item[(iv)] \emph{The radius account.} The law sites $b_-$ are those of the coordinates of $\phi$:
they receive the credit $\mathrm{c}1$ at the arrest; their kinds are those listed in (e), never
$(\mathrm{n})$; the family $\mathfrak{Q}(\mathfrak{a})$ is carried by each law site's own families
and parameter discs; the debits $\ell1$, $\ell4$, $\ell5$, $\ell6$ and the margins are those of
every coordinate of $\phi$, as in \eqref{4.22:eq-radius-accounting-a}. No entry and no amount is
added: the datum enters only as a summand of the bound. The account closes together with its
component, by clause (v) of the frame lemma and through the entries $\mathrm{x}1$, $\mathrm{x}2$ of
\eqref{4.22:eq-radius-accounting-a}.
\end{itemize}
\end{proposition}

\begin{proof}
\emph{Clause (i).} The form $\mathsf{q}_{\mathsf{p}}(\phi;P)$ is a polynomial of degree two in
$\phi$. By the covariance \eqref{4.22:eq-gaussian-oscillation-1} of hypothesis (a) it is jointly
invariant under the simultaneous rotation of the slot and gauge transformation of the pair. The
lemma on orbit data of jointly invariant polynomials, clause (i), applies at any radius; applied at
radius $\varepsilon_\star$ at the law sites of the bonds of $q_1\cup q_2$, it shows that
$\mathsf{q}_{\mathsf{p}}^{\sharp}$ is the orbit datum of $\mathsf{q}_{\mathsf{p}}$ on
$\mathfrak{K}_P\restriction\widetilde{S}_{\mathsf{p}}$. For the bound, recall that
$\operatorname{pol}(g)=\log\max\{\|g\|,\|g^{-1}\|\}$, so that
$|\operatorname{Ad}_gB|\le\|g\|\,\|g^{-1}\|\,|B|\le e^{2\operatorname{pol}(g)}|B|$; this is the
elementary bound on rotated coordinates, and it gives
$|\phi^{[g]}(b)|\le e^{2\operatorname{pol}(g)}\mathfrak{b}_W\le e^{2\varepsilon_\star}\mathfrak{b}_W$
on the chart of radius $\varepsilon_\star$, since $\phi$ is a real slot in its box by hypothesis
(b). The modulus of $\mathsf{q}_{\mathsf{p}}$ is at most the kernel bound of (a), times the number
$n_q$ of terms, times $\frac12$, times the square of the supremum of the slot:
\begin{equation*}
|\mathsf{q}_{\mathsf{p}}(\phi^{[g]};P)|
\le\tfrac12\,n_qB_{\mathcal{A}}e^{-(\kappa_1-1)|\mathsf{y}|}\,e^{4\varepsilon_\star}\mathfrak{b}_W^2
=e^{4\varepsilon_\star}\bar{\mathsf{q}}_{\mathsf{p}}.
\end{equation*}
The number $\varepsilon_\star$ satisfies $e^{2\varepsilon_\star}\le 1+1/160$, so
$e^{4\varepsilon_\star}\le(1+1/160)^2\le 1.013$, which is
\eqref{4.22:eq-gaussian-oscillation-5}. Were $\phi$ stored instead as a weak collar variable of
complex radius $e^{-2\varepsilon_\star}\varrho^B_j$, the rotated slot would be bounded by
$\varrho^B_j$ and, since $\varrho^B_j/\mathfrak{b}^B_j=6.6/6$, the bound would grow by the constant
factor $(6.6/6)^2$.

\emph{Clause (ii).} Fix $\tau,\tau'\in\mathfrak{P}$ and put
$\phi_\tau:=\phi^{[\mathfrak{u}(\tau)g']}$, $\phi_{\tau'}:=\phi^{[\mathfrak{u}(\tau')g']}$; write
$q(\psi;P)$ for the $\mathbb{C}$-bilinear form
$-\frac12\sum_{b\in q_1,b'\in q_2}\langle\psi(b),\mathcal{A}_{\mathsf{y}}[P](b,b')\psi(b')\rangle$.
The oscillation $\operatorname{osc}_{\mathfrak{P}}(\mathsf{q}_{\mathsf{p}}^{\sharp})$ is the supremum
of the modulus of $q(\phi_\tau;\mathsf{R}(\tau))-q(\phi_{\tau'};\mathsf{R}(\tau'))$, and
\begin{align*}
q(\phi_\tau;\mathsf{R}(\tau))-q(\phi_{\tau'};\mathsf{R}(\tau'))
&=\bigl[q(\phi_\tau;\mathsf{R}(\tau))-q(\phi_\tau;\mathsf{R}(\tau'))\bigr]\\
&\quad+\bigl[q(\phi_\tau;\mathsf{R}(\tau'))-q(\phi_{\tau'};\mathsf{R}(\tau'))\bigr].
\end{align*}
The first bracket is the kernel difference of (a) evaluated at the slot $\phi_\tau$. By the
hypothesis $\operatorname{pol}(\mathfrak{u}(\tau)(s)g'(s))<\varepsilon_\star$ and the elementary
bound on rotated coordinates, $|\phi_\tau|\le e^{2\varepsilon_\star}\mathfrak{b}_W$; comparing with
\eqref{4.22:eq-gaussian-oscillation-3}, in which the slot enters through $\mathfrak{b}_W^2$, the
first bracket is at most $e^{4\varepsilon_\star}\mathfrak{N}_{\mathfrak{O}}(\mathsf{q}_{\mathsf{p}})$.
By bilinearity and symmetry of the kernel the second bracket equals
\begin{multline*}
-\tfrac12\sum_{b\in q_1,\,b'\in q_2}\Bigl[\bigl\langle(\phi_\tau-\phi_{\tau'})(b),
\mathcal{A}_{\mathsf{y}}(b,b')\phi_\tau(b')\bigr\rangle\\
+\bigl\langle\phi_{\tau'}(b),\mathcal{A}_{\mathsf{y}}(b,b')(\phi_\tau-\phi_{\tau'})(b')\bigr\rangle\Bigr],
\end{multline*}
with $\mathcal{A}_{\mathsf{y}}=\mathcal{A}_{\mathsf{y}}[\mathsf{R}(\tau')]$. Using the kernel bound of
(a), $|\phi_\tau|,|\phi_{\tau'}|\le e^{2\varepsilon_\star}\mathfrak{b}_W$ and the definition of
$\bar{\mathsf{q}}_{\mathsf{p}}$, its modulus is at most
$\bar{\mathsf{q}}_{\mathsf{p}}\cdot(\sup|\phi_\tau-\phi_{\tau'}|/\mathfrak{b}_W)
\cdot 2e^{2\varepsilon_\star}$.

It remains to bound $\phi_\tau-\phi_{\tau'}$. Write
$\mathfrak{u}(\tau)\mathfrak{u}(\tau')^{-1}=\prod_ie^{\tilde Z_i}$ at the law site concerned; then
$\phi_\tau=\operatorname{Ad}_{\prod_ie^{\tilde Z_i}}\phi_{\tau'}$. For one factor,
\begin{equation*}
(\operatorname{Ad}_{e^{\tilde Z}}-1)B=\int_0^1\operatorname{Ad}_{e^{t\tilde Z}}[\tilde Z,B]\,dt,
\end{equation*}
and $|\operatorname{Ad}_{e^{t\tilde Z}}|\le e^{2t|\tilde Z|}$, $|[\tilde Z,B]|\le 2|\tilde Z|\,|B|$,
so $|(\operatorname{Ad}_{e^{\tilde Z}}-1)B|\le(e^{2|\tilde Z|}-1)|B|$; this is the integrated form
of the elementary bound on rotated coordinates. The identity
$\operatorname{Ad}_{g_1g_2}-1=\operatorname{Ad}_{g_1}(\operatorname{Ad}_{g_2}-1)
+(\operatorname{Ad}_{g_1}-1)$, applied with $g_1$ the product of the first factors, of total size
$\varsigma_1$, and $g_2$ the next factor, of size $\varsigma_2$, together with
$e^{2\varsigma_1}(e^{2\varsigma_2}-1)+(e^{2\varsigma_1}-1)=e^{2(\varsigma_1+\varsigma_2)}-1$, gives
by induction on the number of factors
$|(\operatorname{Ad}_{\prod_ie^{\tilde Z_i}}-1)B|\le(e^{2\sum_i|\tilde Z_i|}-1)|B|$; no
Baker--Campbell--Hausdorff formula is used. Hence, with $|\phi_{\tau'}|\le
e^{2\varepsilon_\star}\mathfrak{b}_W$ and $\mathsf{z}\le 10^{-2}$,
\begin{equation*}
|\phi_\tau-\phi_{\tau'}|\le(e^{2\mathsf{z}}-1)e^{2\varepsilon_\star}\mathfrak{b}_W
\le 2.03\,e^{2\varepsilon_\star}\mathsf{z}\,\mathfrak{b}_W ,
\end{equation*}
since $t\mapsto(e^{2t}-1)/t$ is increasing on $t>0$ and $100(e^{0.02}-1)\le 2.03$. The second
bracket is therefore at most $2\cdot 2.03\,e^{4\varepsilon_\star}\mathsf{z}\,\bar{\mathsf{q}}_{\mathsf{p}}$;
with $e^{4\varepsilon_\star}\le 1.013$ and $2\cdot 1.013\cdot 2.03\le 4.12$ this is at most
$4.12\,\mathsf{z}\,\bar{\mathsf{q}}_{\mathsf{p}}$. With $e^{4\varepsilon_\star}\le 1.013$ for the
first bracket this gives \eqref{4.22:eq-gaussian-oscillation-6}.

\emph{Clause (iii): the premise.} By hypothesis (e) the law sites of the coordinates of $\phi$ are
never of kind $(\mathrm{n})$ and lie at distance at least $\check R_k-1$ from every source, and, by
hypothesis (e$'$), \eqref{4.22:eq-gaussian-oscillation-4} bounds every layer at $s=b_-$. By
hypothesis (f),
$B_H^S\bar\alpha(x_k)\le\varrho_0$ and $E(x_k)^{\delta_\natural}\le 2^{-10}$; since
$e^2\cdot 2^{-10}\le 7.4/1024$, which is below $\frac16$ and below $0.0073$,
\begin{equation}\label{4.22:eq-gaussian-oscillation-8}
\Theta_0^{\mathrm{loc}}(\mathcal{K}(\tau))(s)\le 0.0073\,\varrho_0\le\varrho_0/6
\end{equation}
for every $\tau\in\mathfrak{P}$ and every layer of the site. Thus
$a_0:=\Theta_0^{\mathrm{loc}}(\mathcal{K}(\tau'))(s)\le\varrho_0/6\le\varrho_0/3$; since
$\Theta_0^{\mathrm{loc}}$ is a seminorm,
\begin{equation*}
\xi:=\Theta_0^{\mathrm{loc}}(\mathcal{K}(\tau)-\mathcal{K}(\tau'))(s)
\le 2\cdot 0.0073\,\varrho_0\le 0.015\,\varrho_0\le\varrho_0/3;
\end{equation*}
and $a_0+\xi\le\varrho_0$. These are the hypotheses of (d).

\emph{Clause (iii): the smallness hypothesis of (ii).} At the site $(\mathrm{L})$ the riding map is
a product of $r\le 4$ link factors, indexed by $i$ (see \eqref{4.22:eq-gaussian-oscillation-9}
below). Let $A''_i$ be the increment of the lower end
of the $i$-th link, built from the differences of at most three lower layers. By
\eqref{4.22:eq-gaussian-oscillation-8} and the seminorm property,
\begin{equation*}
\xi':=\Theta_0^{\mathrm{loc}}(A''_i)(s)\le 3\cdot 2\cdot 0.0073\,\varrho_0\le 0.044\,\varrho_0,
\end{equation*}
since
$6\cdot 0.0073=0.0438$. Each of the at most $2r\le 8$ logarithms entering the factorisation
obtained below is at most $2.15K_u\xi$ or $2.15K_u\xi'$, four of each kind, so, with
$K_u\varrho_0\le 10^{-2}$ from the definition of $\varrho_0$,
\begin{equation*}
\mathsf{z}\le 4\cdot 2.15\,(0.015+0.044)K_u\varrho_0\le 8.6\cdot 0.059\cdot 10^{-2}
<0.51\cdot 10^{-2}<10^{-2},
\end{equation*}
using $4\cdot 2.15=8.6$ and $8.6\cdot 0.059=0.5074$. At $(\mathrm{T})$ and $(\mathrm{P})$ there is one
factor and $\mathsf{z}\le 2.15\cdot 0.015\cdot 10^{-2}$. This second bound on $\mathsf{z}$ comes from
the distance of the law site from the sources; the display
\eqref{4.22:eq-gaussian-oscillation-7} keeps the first bound, through
$\Delta_{\mathfrak{O}}$. No smallness of $\Delta_{\mathfrak{O}}$ is required; what is required is the
bound \eqref{4.22:eq-gaussian-oscillation-8} on the base layer at $b_-$.

\emph{Clause (iii): one layer.} By the premise, the one-layer logarithm bound (d) gives
$|\log[\mathfrak{u}(\tau')^{-1}\mathfrak{u}(\tau)](s)|\le 2.1K_u\xi$. Since
$\mathfrak{u}(\tau)\mathfrak{u}(\tau')^{-1}=\mathfrak{u}(\tau')\bigl(\mathfrak{u}(\tau')^{-1}
\mathfrak{u}(\tau)\bigr)\mathfrak{u}(\tau')^{-1}$, the logarithm of
$\mathfrak{u}(\tau)\mathfrak{u}(\tau')^{-1}$ is the conjugate by $\mathfrak{u}(\tau')$ of that of
$\mathfrak{u}(\tau')^{-1}\mathfrak{u}(\tau)$. By the second part of hypothesis (d$'$),
$\operatorname{pol}(\mathfrak{u}(\tau')(s))\le 10^{-2}$, so this conjugation costs at most
$e^{0.02}\le 1.021$; and $2.1\cdot 1.021\le 2.15$. The value $\mathfrak{u}(s)$ depends on the layer
only through its restriction to $B^{\mathrm{top}}(s)$, which lies in the $M$-cube of $s$, contained
in $\widetilde{S}_{\mathsf{p}}$; so the bound \eqref{4.22:eq-gaussian-oscillation-10} of hypothesis
(g) gives
\begin{equation*}
\xi\le\sup_{y\in B^{\mathrm{top}}(s)}N_y(\mathcal{K}(\tau)-\mathcal{K}(\tau'))
\le\bar\alpha(x_\ell)\Delta_{\mathfrak{O}}(S_{\mathsf{p}}).
\end{equation*}
This proves
\eqref{4.22:eq-gaussian-oscillation-7}. Inserted in \eqref{4.22:eq-gaussian-oscillation-6}, it
contributes $4.12\cdot 2.15\le 8.9$ times
$K_u\bar\alpha(x_\ell)\Delta_{\mathfrak{O}}\bar{\mathsf{q}}_{\mathsf{p}}$.

\emph{Clause (iii): several links.} At a law site of the coordinates of $\phi$ one has $h_q=1$ by
(e), so $\mathfrak{u}=\mathfrak{Q}^{\mathfrak{b}}=\Gamma_1\cdots\Gamma_r$ with $r\le 4$ and
$\Gamma_i:=\mathfrak{w}_{s_i}=\mathfrak{u}[e^{i\eta\mathbb{H}_{s_i}}W_{s_i};W_{s_i}]$. Put
$p_i:=\Gamma_1\cdots\Gamma_i$, $p_0:=1$ and $e_i:=\Gamma_i(\tau)\Gamma_i(\tau')^{-1}$. Then
\begin{equation*}
p_r(\tau)p_r(\tau')^{-1}
=\bigl[p_{r-1}(\tau)e_rp_{r-1}(\tau)^{-1}\bigr]\cdot p_{r-1}(\tau)p_{r-1}(\tau')^{-1},
\end{equation*}
as one sees by inserting $p_r=p_{r-1}\Gamma_r$; by induction on $r$,
\begin{equation}\label{4.22:eq-gaussian-oscillation-9}
\mathfrak{u}(\tau)\mathfrak{u}(\tau')^{-1}=\prod_{i=r}^{1}p_{i-1}(\tau)\,e_i\,p_{i-1}(\tau)^{-1},
\end{equation}
the factor with $i=r$ standing on the left. Each $e_i$ factors as
$e_i=(\Gamma_i(\tau)\hat\Gamma_i^{-1})(\hat\Gamma_i\Gamma_i(\tau')^{-1})$, with the intermediate
element
$\hat\Gamma_i:=\mathfrak{u}[e^{i\eta\mathbb{H}_{s_i}(\tau')}W_{s_i}(\tau);W_{s_i}(\tau)]$. In the
first factor the layer moves at the fixed lower end $W_{s_i}(\tau)$; this is the one-layer step
above. In the second factor the lower end moves, $W_{s_i}(\tau)=e^{i\eta A''_i}W_{s_i}(\tau')$, at
the fixed layer $X_0:=\mathbb{H}_{s_i}(\tau')$. With $W:=W_{s_i}(\tau')$ and
$W_\lambda:=e^{i\eta\lambda A''_i}W$ put
\begin{equation*}
Z(\lambda):=\log\bigl[\mathfrak{u}[e^{i\eta X_0}W;W]^{-1}\,
\mathfrak{u}[e^{i\eta X_0}W_\lambda;W_\lambda]\bigr](s).
\end{equation*}
Let $b_0$ be $\Theta_0^{\mathrm{loc}}$ of the complex part of $W$. The complex part $A'$ of the
tube is at most $\bar\alpha(x)\le\varrho_0/B_H\le\varrho_0/6$ by (f), and the at most three lower
layers contribute at most $3\cdot 0.0073\varrho_0$ by \eqref{4.22:eq-gaussian-oscillation-8}, so
$b_0\le\varrho_0/3$. The complex part of the base of $W_\lambda$ is at most $b_0+|\lambda|\xi'$;
by hypothesis (c), whose radius $\alpha_1^{98}$ is at least $\varrho_0$, the function $Z$ is
holomorphic on the disc $|\lambda|\xi'<\varrho_0-b_0$. It satisfies $Z(0)=0$, and
$|Z|\le 1.05\cdot 2K_ua_0$ there by the first part of hypothesis (d$'$), applied with the layer
$X_0$, for which $\Theta_0^{\mathrm{loc}}(X_0)(s)$ is the quantity $a_0$ of the premise for the
layer $\mathcal{K}=\mathbb{H}_{s_i}$, and with the bases $W$ and $W_\lambda$. Since
\begin{equation*}
\xi'\le\varrho_0/3<2\varrho_0/3\le\varrho_0-b_0,
\end{equation*}
the
point $\lambda=1$ lies inside the disc, and the Schwarz lemma gives
\begin{equation*}
|Z(1)|\le\frac{2.1K_u\,a_0\,\xi'}{\varrho_0-b_0}
\le 2.1K_u\,\xi'\cdot\frac{1/6}{2/3}\le 2.1K_u\,\xi'.
\end{equation*}
Thus $\mathfrak{u}(\tau)\mathfrak{u}(\tau')^{-1}$ is, by \eqref{4.22:eq-gaussian-oscillation-9}, a
product of at most $2r\le 8$ factors $e^{\tilde Z_i}$. Since $\Delta_{\mathfrak{O}}(S_{\mathsf{p}})$
is taken per link and layer, the increments of the layer and of the lower end of each link are
bounded in threshold units by $\Delta_{\mathfrak{O}}(S_{\mathsf{p}})$, and the bound
\eqref{4.22:eq-gaussian-oscillation-10} of hypothesis (g) gives
$\xi,\xi'\le\bar\alpha(x_\ell)\Delta_{\mathfrak{O}}(S_{\mathsf{p}})$; so each logarithm is at most
$2.1K_u\bar\alpha(x_\ell)\Delta_{\mathfrak{O}}(S_{\mathsf{p}})$ before conjugation. Passing from the
quotient $\mathfrak{u}(\tau')^{-1}\mathfrak{u}(\tau)$ to $\mathfrak{u}(\tau)\mathfrak{u}(\tau')^{-1}$
costs at most $1.021$, as in the one-layer step. Conjugation by
$p_{i-1}(\tau)$ costs at most $e^{2\operatorname{pol}(p_{i-1})}$; at a law site of the coordinates
of $\phi$, by (e) and (f),
\begin{equation*}
\operatorname{pol}(p_{i-1})\le\eta^L_k(s)\le\mu_k\le\varepsilon_\star/32,
\end{equation*}
and $\varepsilon_\star/32\le 10^{-4}$ because $2\varepsilon_\star\le\log(1+1/160)\le 1/160$; so this
cost is at most $e^{2\cdot 10^{-4}}\le 1.0003$. As $2.1\cdot 1.021\cdot 1.0003\le 2.15$, each
$|\tilde Z_i|\le 2.15K_u\bar\alpha(x_\ell)\Delta_{\mathfrak{O}}(S_{\mathsf{p}})$, whence
$\mathsf{z}(s)\le 8\cdot 2.15K_u\bar\alpha(x_\ell)\Delta_{\mathfrak{O}}(S_{\mathsf{p}})$. The
smallness hypothesis of (ii) was verified above, so \eqref{4.22:eq-gaussian-oscillation-6} applies
at every site; with \eqref{4.22:eq-gaussian-oscillation-3} and the definition of
$\bar{\mathsf{q}}_{\mathsf{p}}$, and since $8\cdot 4.12\cdot 2.15\le 70.9\le 71.2=8\cdot 8.9$, it gives
\eqref{4.22:eq-gaussian-oscillation-a}. At $(\mathrm{T})$ and $(\mathrm{P})$ the factor is
$4.12\cdot 2.15\le 8.9\le 71.2$, so \eqref{4.22:eq-gaussian-oscillation-a} holds there as well.

\emph{Clause (iv).} Clauses (i) and (iv) of the far-site lemma, hypothesis (e), identify the law
sites and their kinds. Nothing in the proofs of clauses (i)--(iii) changes a radius, so no credit
or debit of \eqref{4.22:eq-radius-accounting-a} is added.
\end{proof}

\begin{remark}[Derivatives of the datum along the parameter]
In the notation of \eqref{4.22:eq-majorant-transfer-a}, write $D_\pi$ for the derivative along the
parameter $\pi$ and $c_\pi$ for its constant. By clause $(\mathrm{M}3)'$ of the transfer of
majorant bounds to the orbit chart, \eqref{4.22:eq-majorant-transfer-a}, the derivatives of the
datum along the parameter satisfy
$|D_\pi\mathsf{q}_{\mathsf{p}}^{\sharp}|\le c_\pi\operatorname{osc}_{\mathfrak{P}}
(\mathsf{q}_{\mathsf{p}}^{\sharp})$, so that \eqref{4.22:eq-gaussian-oscillation-a} bounds them.
\end{remark}

\begin{remark}[The base class of the riding map read in the small-field class]
If the base of the riding map in hypothesis (c) is taken instead in the configuration class of
\cite{B-CMP109} of radius $\gamma_0^{(k)}=5w_\ast\sup\alpha$, where
$w_\ast=\mathsf{L}_0^{2N_0}\ge 3^4$ since $\mathsf{L}_0\ge 3$ and $N_0\ge 2$, so that
$\gamma_0^{(k)}\ge 405\bar\alpha(x_k)$, the Schwarz step of the proof closes as well. Under
\eqref{4.22:eq-coupling-ceilings-m} every layer is at most $\frac34\bar\alpha$; hence
$a_0\le\frac34\bar\alpha$, $b_0\le\bar\alpha+3\cdot\frac34\bar\alpha=\frac{13}4\bar\alpha$ and
$\xi'\le 3\cdot 2\cdot\frac34\bar\alpha=\frac92\bar\alpha$. Since
$405-\frac{13}{4}>401>\frac92$, one has $\xi'<\gamma_0^{(k)}-b_0$, so that $\lambda=1$ lies inside
the disc, and $a_0/(\gamma_0^{(k)}-b_0)\le\frac34/401<\frac14$. With this choice of base class the
bound $2.1K_u\xi'$ and the constant $71.2$ of \eqref{4.22:eq-gaussian-oscillation-a} are unchanged.
\end{remark}

\begin{definition}[Geometry and counting of live large-field regions. Proof outstanding]
\label{4.22:not-live-region-geometry}
Let $Y$ be a large-field region created at step $m$ (in this section $m$ denotes the creation
step, not the torus index); the pair is written $(Y,m)$, and $Y$ is \emph{live} at a level if it
still carries its large-field factor at that level. The region is treated as a frame
$W=Y$ of level $k_W=m$, with data $\Phi_Y:=\Phi_W$ and law sites $\Sigma_Y:=\Sigma_W$. Its
strip is written $Y^{\sharp}$.

The statements below involve the following further objects:
\begin{itemize}
\item for each level $\ell$, a map $\pi_\ell$;
\item the distance $\operatorname{dist}_\ell$, which is the supremum-norm distance measured in
units of level $\ell$;
\item two relations between sets, written \emph{touch} and \emph{meet};
\item sets $[S]_i$, attached to a set $S$ and a level $i$, and the numbers
$N_\ell:=\check{R}_\ell$;
\item \emph{cells} $\Delta$, each with a level $\operatorname{lev}(\Delta)$, and \emph{zones};
\item \emph{blocks} $\square$, with the enlarged blocks $\widetilde{\square}$ and
$\widetilde{\square}^{4}$;
\item \emph{probes} $\wp$;
\item a set $\operatorname{pa}(\Delta)$ attached to each cell $\Delta$, and numbers $s_i(Y)$
attached to $Y$ and a level $i$.
\end{itemize}
The following statements are used in this chapter as hypotheses.

\emph{Geometric hypotheses.} Let $(Y,m)$ be as above.
\begin{itemize}
\item[(a)] The number $s_i(Y)$ satisfies
\begin{equation}\label{4.22:eq-live-region-geometry-1}
s_i(Y)\le (101\,\check{R}_m-1)\,L^{-(i-m)}\le 101\,\check{R}_i .
\end{equation}
\item[(b)] A cell meeting $Y^{\sharp}$ has level at most $m$ and lies in $Y^{\sharp}$. The
strips of distinct live large-field regions are disjoint, and a cell lies in at most one of
them.
\item[(c)] For the current small-field region $\Omega_{m+1}$,
$\operatorname{dist}_m(Y^{\sharp},\Omega_{m+1})\ge \tfrac52\check{R}_m$. Every cell at
$\operatorname{dist}_m$-distance less than $\tfrac52\check{R}_m$ from $Y^{\sharp}$ has level at
most $m$, and, for every $i>m$, this neighbourhood of $Y^{\sharp}$ lies in the component of
$\Lambda_i^{c}$ containing $Y$. This is the region's own collar. The collar
made by the operation $\mathbf{T}$ is wider; after an interrupted preparation whose index
$h_1$ equals $m$ only the width $\tfrac52\check{R}_m$ is left, and clause (c) claims no more
than this width.
\item[(d)] At the site $(\mathrm{T})$ of the step from level $i$ to level $i+1$,
$\operatorname{dist}_i(Y^{\sharp},\Omega_{i+1})\ge 7\check{R}_i$.
\item[(e)] At the site $(\mathrm{L})$ or $(\mathrm{P})$ of $\mathbf{R}_k$, let $C''$ be the
component of $\Lambda_k^{c}$ containing $Y$. If $C''$ is not treated, then
$Y^{\sharp}\subset C''$ and $\operatorname{dist}_k(Y^{\sharp},Z)\ge\check{R}_k$.
\item[(f)] Let $C$ be the component of $\Lambda_k^{c}$ containing $Y$ (the component written
$C''$ in (e)), and let $\Lambda_C$ and $\widetilde{\Lambda}_C$ be the two sets attached to $C$
by the large-field operation $\mathbf{R}_k$. The
set $\widetilde{\Lambda}_C$ is a box whose sides are at most $h_\sharp=112$,
measured at level $k$, and $\operatorname{dist}_k(\widetilde{\Lambda}_C,Z^{c})\ge 9\check{R}_k$.
\item[(g)] Let the component $C$ containing $Y$ be treated at $\mathbf{R}_k$, and let $h$
denote the level attached to this treatment. Then $m\le h$.
The set $\mathsf{C}_C$ is a box, a union of cells of levels at most $h$, with sides of at most
$104\,L\,\check{R}_{h+1}$ units of level $h$. Moreover $Y^{\sharp}\subset\mathsf{C}_C$ and
\begin{equation}\label{4.22:eq-live-region-geometry-2}
\operatorname{dist}_h\bigl(Y^{\sharp},\mathsf{C}_C^{c}\bigr)\ge\tfrac12\bigl(\check{R}_h+1\bigr).
\end{equation}
A cube of any level at most $k$ touching $Y^{\sharp}$ lies in $\Lambda_C$.
\item[(h)] At the site $(\mathrm{P})$ of $\mathbf{R}_k$: if the component containing $Y$ is
treated, then $\operatorname{dist}_{h+1}(Y^{\sharp},R_\nu)\ge\tfrac52\check{R}_{h+1}$ for every
collar $R_\nu$; if it is not treated, then
$\operatorname{dist}_k(Y^{\sharp},R_\nu)\ge\check{R}_k$ for every collar $R_\nu$. No integration
bond of $(\mathrm{P})$ lies on a strip.
\item[(i1)] A probe meeting $Y^{\sharp}$ has level at most $m$.
\item[(i2)] A probe meets the strips of at most $2^d$ live large-field regions.
\end{itemize}

\emph{Contour hypothesis.} Let $\mathsf{c}_{\wedge}>0$ be the zone weight of the
weighted contour distance $d^{\wedge}$. Then $d^{\wedge}$ obeys the triangle inequality. A
contour pays at least $\mathsf{c}_{\wedge}$ for each zone it crosses. On a set all of whose
cells have level at most $n$, a contour pays at least $\mathsf{c}_{\wedge}$ for each level-$n$
unit of its Euclidean length. Payments on disjoint arcs add.

\emph{Cell hypothesis.} The weight $\mathsf{c}_{\wedge}$ satisfies
\begin{equation}\label{4.22:eq-live-region-geometry-3}
\tfrac{1}{16}\,\delta_P\,\mathsf{c}_{\wedge}\ \ge\ c_{\boxplus}\log\hat{b}_t+1\ \ge\ 1 ,
\end{equation}
where $c_{\boxplus}$ and $\hat{b}_t$ are fixed constants. Clauses (b) and (c) of the cell
estimate of this chapter hold for every value, not smaller than $\delta_P/16$, of the parameter
in which that estimate is stated.

\emph{Summation hypothesis.} Let a site geometry
$(\pi^{\mathfrak{s}},\mathfrak{H},\mathfrak{F})$ of level $i$ be given: a family
$\pi^{\mathfrak{s}}$ among whose members are distinguished members called source cells and
hubs, together with two further data $\mathfrak{H}$ and $\mathfrak{F}$ (the latter distinct
from the frame family $\mathfrak{F}(S)$). For a member $\Delta$
of $\pi^{\mathfrak{s}}$, $\rho(\Delta)$ is the
least cardinality of a connected subfamily of $\pi^{\mathfrak{s}}$ that contains $\Delta$ and a
source cell. A member \emph{sees} a set $\mathcal{E}$ if it meets $\mathcal{E}$, or if it is a
source cell sharing a wall with a hub that meets $\mathcal{E}$. A \emph{labelled term} is a
triple $F=(A_F,\mathcal{E}_F,\|F\|)$ with $A_F\in\mathbf{D}_i$ satisfying
\begin{itemize}
\item[(R)] $\operatorname{dist}_i(\operatorname{pa}(\Delta),A_F)\le c_\rho\,\rho(\Delta)$ for
every member $\Delta$ seeing $\mathcal{E}_F$, where $c_\rho=120$.
\end{itemize}
For a labelled term $F$, $\mathfrak{Y}(F)$ is the set of families $\mathsf{y}$ every component of
which contains a source cell and a member seeing $\mathcal{E}_F$, and
$\operatorname{lab}(F,\mathsf{y})$ is the label of the pair $(F,\mathsf{y})$. For a family
$\mathfrak{L}$ of labelled terms and $r\ge0$,
\begin{equation}\label{4.22:eq-live-region-geometry-4}
\mathsf{T}^{*}_{r}[\mathfrak{L}]
:=\sup_{\Delta}\ e^{-\rho(\Delta)}
\sum_{\substack{F\in\mathfrak{L}:\\ \Delta\text{ sees }\mathcal{E}_F}}
\|F\|\,e^{r\,d_i(A_F)} .
\end{equation}
Put
$c_l=1+234d$, let $\lambda_*=\lambda_*(d,L)$ be a number depending on $d$ and $L$, let
$C^{\mathrm{sum}}_{\Sigma}$ be a constant, and fix an element $Q_*$; the inner sums below run
over the families $\mathsf{y}$ whose label contains $Q_*$ as a member. Then:
\begin{itemize}
\item[(1)] $d_i(\operatorname{lab}(F,\mathsf{y}))\le d_i(A_F)+c_l|\mathsf{y}|$.
\item[(2)] If $\kappa_1-1\ge c_l r+\lambda_*$, then
\begin{equation}\label{4.22:eq-live-region-geometry-5}
\sum_{F\in\mathfrak{L}}\ \sum_{\mathsf{y}:\ \operatorname{lab}(F,\mathsf{y})\ni Q_*}
\|F\|\,e^{-(\kappa_1-1)|\mathsf{y}|}\,e^{r\,d_i(\operatorname{lab}(F,\mathsf{y}))}
\ \le\ C^{\mathrm{sum}}_{\Sigma}\,\mathsf{T}^{*}_{r+2}[\mathfrak{L}] .
\end{equation}
\item[(3)] Let $D$ and $\varsigma$ be numbers. If every label has
$d_i(\operatorname{lab}(F,\mathsf{y}))\ge D$ and
$\kappa_1-1\ge c_l(r+\varsigma)+\lambda_*$, then the left side of
\eqref{4.22:eq-live-region-geometry-5} is at most
\begin{equation}\label{4.22:eq-live-region-geometry-6}
e^{-\varsigma D}\,C^{\mathrm{sum}}_{\Sigma}\,\mathsf{T}^{*}_{r+\varsigma+2}[\mathfrak{L}] .
\end{equation}
\end{itemize}

\emph{Depth hypothesis.} The depth datum $\varrho$, a function on probes, satisfies
$\varrho(\wp)\ge1$ for every probe $\wp$; if $\wp$ meets the strip of a live
large-field region $(Y,m)$, then $\check{R}_m\le 2\varrho(\wp)+2$.

\emph{Condition on $\kappa_1$.} The parameter $\kappa_1$ satisfies
\begin{equation}\label{4.22:eq-live-region-geometry-7}
\kappa_1\ \ge\ 1+c_l(2\kappa+4)+\lambda_* .
\end{equation}

\emph{Tree graphs.} A tree graph meeting two cubes has length at least their distance.
\end{definition}

\medskip\noindent
Route. The geometric hypotheses (a)--(i2) are to be derived from the property
\cite[(ii)]{B-CMP122a} of the large-field regions together with the nesting of the large-field
regions created at successive steps. The remaining hypotheses are taken by statement.

\begin{definition}[A kept term as a reader: label, norm and law sites]\label{4.22:def-kept-term-reader}
We use the notation of Definition~\ref{4.22:not-live-region-geometry}, whose geometric clauses are
stated there with proof outstanding.
In that notation $(Y,m)$ is a pocket and $Y^{\sharp}$ its strip. For each level $\ell$, $\pi_\ell$ is
the partition into cubes of level $\ell$. For a set in $\mathbf{D}_\ell$, $d_\ell(\cdot)$ is
Ba{\l}aban's size function in level-$\ell$ units, as in the glossary
(Definition~\ref{0.2:not-glossary}), and $N_\ell(\cdot)$ is the number of cubes of $\pi_\ell$
contained in the set. Finally, $\Sigma_Y$ is the set of law sites of the
frame of $Y$. Let
$v\in\mathfrak{K}(Y)$ be a kept member of the pocket $(Y,m)$.

\emph{Format.} The region read by $v$ and its true domain are the strip:
\begin{equation*}
\mathcal{R}_v=X_v=Y^{\sharp},
\end{equation*}
and its set of pockets is $\mathcal{P}_v=\{(Y,m)\}$. The term $v$ is wrapped, in the sense of the
definition of wrapped terms, and at every step it is carried by the clause governing wrapped terms
that is stated with that definition. That clause acts through its part at the
fluctuation integral $(\mathrm{T})$, its two parts at the presentation $(\mathrm{L})$ and its part at
the large-field stages $(\mathrm{P})$. At no step is a Cauchy disc in the correlation theorem's
contour variable $\mathsf{t}$ opened for $v$, and at no step is a piece debit $\eta_p$ attached to
it. The stored label of $v$ is
empty, and its stored length is $\mathsf{d}(v)=0$.

\emph{Label.} At a site of level $i$, $v$ is read with the label
\begin{equation*}
A_v:=Q_Y(i),
\end{equation*}
the cube of $\pi_i$ that contains the first vertex of the box $Y^{\sharp}$. This is a fixed rule and
depends on the scaffold only. Thus $A_v$ consists of one cube of $\pi_i$, and it contains a point of
$Y^{\sharp}$. Since $A_v$ is a single cube of $\pi_i$, we have
\begin{equation*}
d_i(A_v)=0,\qquad N_i(A_v)=1.
\end{equation*}
At the large-field stages $(\mathrm{P})$ the level is $i:=m_v$, where $m_v$ is the level assigned by
the proposition on kept terms at the large-field stages. On the chain of stages of the component that
contains $Y$ this level is $m_v=h+1$, where $h$ is the level of clause~(g) of
Definition~\ref{4.22:not-live-region-geometry} (proof outstanding): the core $\mathsf{C}_C$ of the
component $C\ni Y$ treated at $\mathbf{R}_k$ is a union
of cells of levels at most $h$. On any other chain it is $m_v=k$. In particular, on the chain of
the component containing $Y$, the level $i$ is never the level $k$ of the step.

\emph{Norm.} The norm of $v$ is the oscillation norm of \eqref{4.22:eq-kept-oscillation-a}:
\begin{equation*}
\|v\|:=\mathfrak{O}(v).
\end{equation*}
Thus at a site of level $i$ the term $v$ enters as a reader with data $(A_v,\mathcal{R}_v,\|v\|)$,
in the sense of Definition~\ref{4.22:not-live-region-geometry}, its label $A_v$ lying in
$\mathbf{D}_i$.

\emph{Law sites.} The term $v$ has no strong set, no pointwise law site, and no integration variable
of a later site. Its law sites are the frame sites $\Sigma_Y\subset Y$ together with, for each such
site $s$, the set $B^{\mathrm{top}}(s)$ attached to $s$, which lies in $Y^{\sharp}=\mathcal{R}_v$.

\emph{At the step treating $Y$.} Let $C\ni Y$ be the component of $\Lambda_k^c$ treated by
$\mathbf{R}_k$. At that step, $K_v$ is defined to be one cube of $\pi_k$ that meets $Y^{\sharp}$.
Clause~(g) of Definition~\ref{4.22:not-live-region-geometry} (proof outstanding) states that a cube
of any level at most $k$ touching $Y^{\sharp}$ lies in
the set $\Lambda_C$ of that clause. Hence $K_v\subset\Lambda_C$.
\end{definition}

\begin{definition}[Kept terms in the domain list of the preparatory chain]\label{4.22:def-kept-domain-list}
Fix a level $k$, and let $Y$ be a large-field region created at step $j_Y$ and
live at the large-field operation $\mathbf{R}$ of step $k$. Write $Y^{\sharp}$ for the strip of
$Y$, $\Sigma_Y\subset Y$ for its roots, $\mathsf{B}_Y$ for the union of the supports of the
tents round $\Sigma_Y$, and $\mathfrak{K}(Y)$ for the set of kept members of $Y$
(Definition~\ref{4.22:def-kept-term-reader}). A chain is either the own chain of $Y$ or a
foreign one, as in Definition~\ref{4.22:def-kept-term-reader}, and $\Lambda^{\mathrm{ch}}$
denotes the region on which the domain clauses of the chain are imposed. For an entry
$a\in\{1,\dots,8\}$ of the regularity clauses and a clause location $y$, $r_a(y)$ is the
threshold of the $a$-th entry at $y$, and $Q_a(y)$ is the value at $y$ of the $a$-th clause
quantity of the background. Let $\varpi_{\mathrm{fr}}=\tfrac18(1-\theta_S)\beta$ be the frame
reserve fraction, where $\beta$ is the schedule constant of the layered analyticity spaces
(Definition~\ref{4.4:def-post-step-space}) and $\theta_S$ is the layer number of the inheritance
theorem; recall that $\theta_S\ge\tfrac12$.

Every $v\in\mathfrak{K}(Y)$ enters the domain list of the preparatory chain, that is, the
domain clauses $\mathfrak{I}_n$ and the clauses $(\mathrm{H}^{\natural})_l(a)$, as a bearer of
clauses. It enters in the same way as the clause-bearing terms of the class
$\mathcal{T}^{\mathbb{B}}$ of chain terms of the large-field stages already in the list, and in
the same way as the terms $\mathcal{C}_{\mathrm{arr}}$ of the arrests, which were entered in the
list as clause-bearers of that class.
Its data in the list are its domain $X_v$, its level $\ell_v$, its stored domain
$\mathfrak{U}_v$, its clause set $\operatorname{Cl}(v)$ and its reference thresholds
$\mathbf{r}^v_a$:
\begin{gather*}
X_v:=Y^{\sharp},\qquad \ell_v:=j_Y,\qquad
\mathfrak{U}_v:=\mathfrak{U}_v(Y^{\sharp}),\\
\operatorname{Cl}(v):=\text{the clauses of }v,\qquad
\mathbf{r}^v_a:=r_a .
\end{gather*}
Here $\mathfrak{U}_v(Y^{\sharp})$ is the stored domain with reserve of $v$ on $Y^{\sharp}$.
The clause thresholds of $v$ are
\begin{equation}\label{4.22:eq-kept-domain-list-1}
\theta_{v,a}(y):=
\begin{cases}
r_a(y)\bigl(1-\varpi_{\mathrm{fr}}\,\mathbb{1}[\text{the stencil of }(y,a)
\text{ meets }\mathsf{B}_Y]\bigr), & a\notin\{5,8\},\\[2pt]
r_a(y), & a\in\{5,8\}.
\end{cases}
\end{equation}

The term $v$ is not a term of the flat sum of the inheritance bound. It carries no size factor
$c^{\mathrm{sz}}e^{-\kappa\,\mathsf{d}_{\ell}(v)}$, where $\mathsf{d}_{\ell}(v)$ is the length
assigned to a stored term, and no constant $B_0$, and it has no entry of its own
in the chapter's table of term data. Its pieces are the units of the proposition on the pieces
of kept members at the large-field stages of the step. Each piece has size
$\mathfrak{O}(v)\prod_A K_A\vartheta^{\natural}(\square)$, by clause (iv) of the increment lemma
for the background map; the product over $A$ and the size factors $K_A$ and
$\vartheta^{\natural}(\square)$ are those of that clause. The total of the pieces is carried by
the anchor bound at the large-field stages of the step: its constant
$C^{\mathrm{P},\star}_{\mathsf{a}}$ is increased by
$e\,\mathfrak{n}^{\mathrm{kept}}$. This is the increase that the complex induction hypothesis at
level $k$ provides for kept members in its exception.

Clause (S) of Definition~\ref{4.22:def-kept-term-reader}, by which wrapped terms are run, is
imposed on these units, which means that the kept members are added to the set of
terms run by clause (S). The list so extended is read from the scaffold, and no constant is
chosen in this definition.

Assume the following statements of this chapter:
\begin{itemize}
\item the domain list of the preparatory chain, with its clauses $\mathfrak{I}_n$ and
$(\mathrm{H}^{\natural})_l(a)$;
\item clauses (i) and (iv) of the increment lemma for the background map
$\Phi_{t,\sigma}(b)$;
\item the inheritance theorem, with steps (2) and (4)$(\alpha)$ of its proof and its clause (c);
\item the corollary on the minimal room $r_{\min}$;
\item the geometric lemma on the preparatory chain, which gives, at a foreign chain,
\begin{equation*}
Y\subset C''\subset Z^{c}\subset\Lambda^{\mathrm{ch}},
\end{equation*}
where $C''$ and $Z^{c}$ are the nested sets it attaches to that chain, and, at the own chain,
\begin{equation*}
Y^{\sharp}\subset\mathsf{C}_C\subset Z^{\mathrm{on}},
\end{equation*}
where $\mathsf{C}_C$ is the core of the chain and $Z^{\mathrm{on}}$ the region of the own chain
on which the domain clause is read;
\item the datum \eqref{4.22:eq-frame-proof-datum-b} in the proof of the frame lemma
(Lemma~\ref{4.22:prf-frame-proof-datum}).
\end{itemize}
Then the following hold for every $v\in\mathfrak{K}(Y)$; they are proved below.
\begin{enumerate}
\item[(i)] Let the chain be foreign, that is, not the own chain of $Y$, and let
$Y'\supseteq Y^{\sharp}$. Then, for every layer $n$ of the chain,
$\mathfrak{I}_n(Y'\cap\Lambda^{\mathrm{ch}})$ imposes on
$\operatorname{Cl}(v)$ the conditions
\begin{equation}\label{4.22:eq-kept-domain-list-2}
Q_a(y)<\theta_{v,a}(y)-\sum_{l\le n}\iota^a_l(y)\,r_a(y),
\end{equation}
where $\iota^a_l(y)\ge0$ is the allowance that the domain list spends at layer $l$ on entry $a$
at $y$.
\item[(ii)] Let the chain be foreign, and let the background be an input of the chain domain
at layer $n$, so that \eqref{4.22:eq-kept-domain-list-2} holds at the clause locations of $v$.
Then, simultaneously for all blocks, for every $y\in Y$, all
$|t_{\square}|\le R^{\natural}_{\square}$, all $|\sigma|\le e^{\kappa_1}$ and every $b$ complex
in its ball, where $t_{\square}$, $\sigma$ and $b$ are the arguments of the background map
$\Phi_{t,\sigma}(b)$ of the inheritance theorem,
\begin{equation*}
\Phi_{t,\sigma}(b)\restriction Y^{\sharp}\in\mathfrak{U}_v(Y^{\sharp}).
\end{equation*}
The decomposition is the one with the unitary part $U$ and the cube gauges $\{u_{\square}\}$ of
the background before the shift, which the shift keeps, and the reserve is included.
\item[(iii)] The increment bounds used in (ii) are homogeneous of degree one in $\sup|b|$, and
therefore hold, scaled accordingly, on the ball of radius $6.6C_1\delta'_j$ on which clause (ii)
of the lemma of this chapter that treats a parent $v\in\mathfrak{K}(Y)$ reads $b$; here
$\delta'_j$ is the primed radius at the level $j$ of that clause.
\item[(iv)] If $k>j_Y$, then for $a\notin\{5,8\}$ the chain's input pair obeys the clause of $v$
with room at least $\tfrac{17}{48}\beta r_a(y)>0$. The entries $a\in\{5,8\}$ are never
incremented.
\item[(v)] The inheritance theorem holds with $v$ in the list, and its proof is unchanged. The
minimal room $r_{\min}$ of the corollary on the minimal room is unchanged, and no cap in the
sense of Definition~\ref{2.3:def-cap} is introduced.
\end{enumerate}
\end{definition}

\begin{proof}
(i) By the geometric lemma on the preparatory chain, assumed above, at a foreign chain
\begin{equation*}
Y\subset C''\subset Z^{c}\subset\Lambda^{\mathrm{ch}} .
\end{equation*}
This inclusion places $Y$ inside $\Lambda^{\mathrm{ch}}$. Take $Y'\supseteq Y^{\sharp}$. The
clause $\mathfrak{I}_n(Y'\cap\Lambda^{\mathrm{ch}})$ is defined in the domain list and ranges
over every clause-bearer of the list. The term $v$ is such a bearer, with domain
$X_v=Y^{\sharp}\subseteq Y'$, so
this clause imposes \eqref{4.22:eq-kept-domain-list-2} on $\operatorname{Cl}(v)$. Hence the
domain of the chain at layer $n$ contains no background that violates these thresholds at the
clause locations of $v$.

(ii) Run step (2) of the proof of the inheritance theorem at the term $T=v$. This step invokes
clause (i) of the increment lemma, which holds for all blocks at once, for every $y\in Y$, for
all $|t_{\square}|\le R^{\natural}_{\square}$ and $|\sigma|\le e^{\kappa_1}$, and for $b$
complex in its ball. For entries $5$ and $8$ that clause states that the decomposition keeps
$U$ and the cube gauges, whatever the shift. The increments $\Delta_a$ of the clause quantities
under $\Phi_{t,\sigma}(b)$ therefore satisfy
\begin{equation*}
\Delta_5=\Delta_8=0,\qquad \Delta_a\le\iota^a_n(y)\,r_a(y)\quad(a\notin\{5,8\}).
\end{equation*}
For $a\in\{1,3,6,7\}$ the bound is clause (i) as stated. For $a\in\{2,4\}$ it comes from the
last clause of clause (i), which applies, since by
\eqref{4.22:eq-kept-domain-list-2} and \eqref{4.22:eq-kept-domain-list-1}
\begin{equation*}
Q_a<\theta_{v,a}\le r_a\le 2\mathbf{r}^v_a .
\end{equation*}
Clause (i) is read at the level of $v$. This is the correct level, because
$\ell_v=j_Y=j(y)$ on $Y^{\sharp}$, where $j(y)$ is the level at which the domain list reads the
location $y$.

Now take $a\notin\{5,8\}$. By the definition of $\Delta_a$ and by
\eqref{4.22:eq-kept-domain-list-2},
\begin{align*}
Q_a\bigl(\Phi_{t,\sigma}(b)\bigr)(y)
&\le Q_a(y)+\Delta_a(y)\\
&<\theta_{v,a}(y)-\sum_{l\le n}\iota^a_l(y)\,r_a(y)+\iota^a_n(y)\,r_a(y)\\
&=\theta_{v,a}(y)-\sum_{l\le n-1}\iota^a_l(y)\,r_a(y).
\end{align*}
For $a\in\{5,8\}$ nothing moves. Since the allowances are non-negative, it follows that
$\Phi_{t,\sigma}(b)\in\mathfrak{I}_{n-1}$ and
$Q_a(\Phi_{t,\sigma}(b)\restriction Y^{\sharp})<\theta_{v,a}$ for every entry $a$.

Step (4)$(\alpha)$ of the proof of the inheritance theorem then gives
$\Phi_{t,\sigma}(b)\restriction Y^{\sharp}\in\mathfrak{U}_v(Y^{\sharp})$, in the decomposition
with the unitary part $U$ and the cube gauges $\{u_{\square}\}$ of the background before the
shift. The reserve is included, because it is part of
$\theta_{v,a}$. At the clause locations whose stencils meet $\mathsf{B}_Y$ we have
$\theta_{v,a}=(1-\varpi_{\mathrm{fr}})r_a$, so the amount $\varpi_{\mathrm{fr}}r_a$ is kept.
The argument uses no input beyond those listed in the hypotheses, and it concerns foreign
chains only: at every foreign chain the hypothesis that the restriction
of the background to $Y^{\sharp}$ lies in $\mathfrak{U}_v(Y^{\sharp})$ is discharged by this
argument. At the own chain of
$Y$, where the geometric lemma on the preparatory chain gives
\begin{equation*}
Y^{\sharp}\subset\mathsf{C}_C\subset Z^{\mathrm{on}},
\end{equation*}
the membership is not derived here; it remains a hypothesis,
supplied there by the domain clause restricted to $Z^{\mathrm{on}}$.

(iii) The bounds of clause (i) of the increment lemma are homogeneous of degree one in
$\sup|b|$. This homogeneity is also used in the proof of clause (ii) of the lemma of this
chapter that treats a parent $v\in\mathfrak{K}(Y)$, described in item (iii) of the statement.
Consequently
these bounds hold, scaled accordingly, at the radius $6.6C_1\delta'_j$. This is the radius at
which that clause, for a parent $v\in\mathfrak{K}(Y)$, reads $b$.

(iv) Assume $k>j_Y$. Let $\theta_{\mathrm{src},a}$ be the threshold that the chain's input pair
obeys at entry
$a$, and put $r_v:=(\theta_{v,a}-\theta_{\mathrm{src},a})/r_a$. By the third line of
\eqref{4.22:eq-frame-proof-datum-b}, read with the level $j_Y$ of $v$ in place of $k_W$,
\begin{equation*}
r_v=(1-\varpi_{\mathrm{fr}})-s^{(k)}_{j_Y}
=\beta\bigl(1-2^{-(k-j_Y)}\bigr)-\varpi_{\mathrm{fr}} ,
\end{equation*}
where $s^{(k)}_{j_Y}$ is the schedule fraction, at step $k$, of a term created at step $j_Y$,
as it appears in that line.
Since $k>j_Y$, $\varpi_{\mathrm{fr}}=\tfrac18(1-\theta_S)\beta$ and $\theta_S\ge\tfrac12$,
\begin{equation*}
r_v\ge\frac{\beta}{2}-\frac{(1-\theta_S)\beta}{8}\ge\frac{7\beta}{16}\ge\frac{3\beta}{8}.
\end{equation*}
By clause (c) of the inheritance theorem, $\sum_l\iota^a_l(y)\le\tfrac12\mathsf{w}_0\beta$,
where $\mathsf{w}_0$ is the constant of the domain list which bounds the total allowance. The
domain list fixes $1/\mathsf{w}_0=12/\min\{\tfrac12,\theta_S\}$, and since $\theta_S\ge\tfrac12$,
\begin{equation*}
\frac{1}{\mathsf{w}_0}=\frac{12}{\min\{\tfrac12,\theta_S\}}=24 .
\end{equation*}
Hence
\begin{equation*}
\sum_l\iota^a_l(y)\le\tfrac12\mathsf{w}_0\beta\le\frac{\beta}{48}.
\end{equation*}
Subtracting the allowances from the room $r_v r_a$ gives a room of at least
\begin{equation*}
\Bigl(\frac{18}{48}-\frac{1}{48}\Bigr)\beta\,r_a(y)=\frac{17}{48}\beta\,r_a(y)>0 .
\end{equation*}
By (ii), the entries $a\in\{5,8\}$ are never incremented. The first inequality above, read
together with $\theta_S\le1$,
\begin{equation*}
r_v\ge\frac{\beta}{2}-\frac{(1-\theta_S)\beta}{8}\ge\frac{\theta_S\beta}{2},
\end{equation*}
is the inequality from which \eqref{4.22:eq-frame-proof-datum-b} derives
$r_{\min}(y)\ge\beta e_\ast(j_Y)$ with $e_\ast(j_Y)=\theta_S2^{-(k-j_Y)}$.

(v) In clause (iv) of the increment lemma the exponent parameter $\theta$ of the profile
$\vartheta^{\natural}(\square)$ equals $\tfrac12$, and at this value the exponent of
$\vartheta^{\natural}(\square)$ is that of the profile $\vartheta(\square)$. The allocation of
shares among
the terms therefore does not change. The number $1/\mathsf{w}_0=24$ obtained in (iv) enters the
constant $K_2^{\flat}\vartheta^{\natural}$ only through the ceilings of sup form on the coupling
in the definition of the stage constants, so no cap is introduced.

The inheritance theorem and its proof apply with $v$ in the list. The clause $\mathfrak{I}_n$
quantifies over a list, and one more clause-bearer only shrinks the chain domain
$\mathfrak{D}^{\flat}_n$. The steps of the proof use, for a clause-bearer $T$, only the
inequalities
\begin{equation*}
\theta_{T,a}\le 2\mathbf{r}^T_a,\qquad r_a\le\mathbf{r}^T_a .
\end{equation*}
For $T=v$ both hold, because $\theta_{v,a}\le r_a=\mathbf{r}^v_a$ by
\eqref{4.22:eq-kept-domain-list-1}. By (iv), when $k>j_Y$, the chain's input pair obeys the
clause of $v$ with room. Hence the minimal room $r_{\min}$ of the corollary on the minimal room
is unchanged.
\end{proof}

\begin{lemma}[The shifted family stays in a kept term's domain at a foreign region]
\label{4.22:lem-foreign-domain-membership}
Fix the step $k$ and a stage index $l\ge 1$ of the inheritance theorem, not exceeding the number
of its stages. Let
$(Y,m)$ be a live large-field region created at step $m<k$, with $Y$ foreign in the sense of
Definition~\ref{4.22:not-live-region-geometry} (geometry and counting of live large-field
regions); let
$v\in\mathfrak{K}(Y)$ be a kept member of $Y$; and let $X\supseteq Y^{\sharp}$ be a union of
cells.

The following notation is used:
\begin{itemize}
\item for a location $y$, $m(y)$ is its shell, the level attached to $y$ in the inheritance
theorem;
\item $\mathfrak{T}$ and $Z^{\mathrm{on}}$ are the two regions of locations distinguished in the
inheritance theorem, the second being called the on-region;
\item $\tilde{U}^{(l)c,\Sigma\flat}_k(X)$ is the stage-$l$ domain of the inheritance theorem on
$X$, and $\mathfrak{N}^{(l)}(X)$ is its stage-$l$ input class on $X$;
\item $Q_a(\mathbb{U},\mathbb{J})(y)$, $1\le a\le 8$, are the eight entries of the clauses of a
stored background domain, evaluated at a pair $(\mathbb{U},\mathbb{J})$ and a location $y$, and
$r_a(y)$ are their thresholds; entries $5$ and $8$ are the ones read
from the unitary part $U$ and the cube gauges $\{u_{\square}\}$ of a decomposition;
\item $Q_a(\Phi_{t,\sigma}(b))(y)$ are the same entries of the pair $\Phi_{t,\sigma}(b)$
introduced below;
\item $\theta_{v,a}$ and $\mathbf{r}^v_a=r_a$ are the thresholds of the domain
$\mathfrak{U}_v(Y^{\sharp})$ of $v$ fixed in Definition~\ref{4.22:def-kept-domain-list}.
\end{itemize}

Assume:
\begin{enumerate}
\item[(a)] (the unweighted-clause hypothesis) at every location of $\mathfrak{T}$ of shell less
than $k$, the clause of the
stage-$l$ domain $\tilde{U}^{(l)c,\Sigma\flat}_k$ carries no weight, that is, both of
Ba{\l}aban's weight factors in it are equal to one;
\item[(b)] (the structure-inclusion hypothesis) for $a\in\{2,4\}$, the structures
$X'\subset Y^{\sharp}$, that is, the sub-regions on which the entries $a\in\{2,4\}$ of the
clauses of $v$ are evaluated,
are among the structures over $X$ at which the clause of the stage-$l$ domain is read;
\item[(c)] the threshold bound of the inheritance theorem;
\item[(d)] clause (i) of the increment lemma of the inheritance theorem, which is the localised
response corollary;
\item[(e)] the relation \eqref{4.22:eq-frame-proof-datum-b} of the proof of the frame lemma
(\ref{4.22:prf-frame-proof-datum}).
\end{enumerate}

Let
\[
(\mathbb{U},\mathbb{J})\in\tilde{U}^{(l)c,\Sigma\flat}_k(X)\cap\mathfrak{N}^{(l)}(X),
\]
and let $\Phi_{t,\sigma}(b)$ be the image of $(\mathbb{U},\mathbb{J})$ under the shifted
background map of the inheritance theorem, with block parameters $t=(t_{\square})$, a complex
parameter $\sigma$ and a complex variable $b$ ranging over its ball.
Then for all block parameters $|t_{\square}|\le r^{\sharp}_{\square}$, with
$r^{\sharp}_{\square}$ the block radii of the inheritance theorem, all
$|\sigma|\le e^{\kappa_1}$ and every complex $b$ in its ball,
\begin{equation}\label{4.22:eq-foreign-domain-membership-1}
\Phi_{t,\sigma}(b)\restriction Y^{\sharp}\in\mathfrak{U}_v(Y^{\sharp}),
\end{equation}
holomorphically, in the sense of the last sentence of clause (i) of the increment lemma
(hypothesis (d)). The membership holds in the decomposition with the same
unitary part $U$ and cube gauges $\{u_{\square}\}$ as the decomposition of the base pair
$(\mathbb{U},\mathbb{J})$, and with the reserve $\varpi_{\mathrm{fr}}r_a$ included in the
thresholds $\theta_{v,a}$. It holds with room at least $\tfrac{7}{16}\beta r_a$, with $\beta$
the constant of \eqref{4.22:eq-frame-proof-datum-b}: for
$a\notin\{5,8\}$ and $y\in Y^{\sharp}$,
\begin{equation}\label{4.22:eq-foreign-domain-membership-2}
Q_a(\Phi_{t,\sigma}(b))(y)<\theta_{v,a}(y)-\tfrac{7}{16}\beta r_a(y).
\end{equation}
\end{lemma}

\begin{proof}
Since $Y$ is foreign, clause (e) of Definition~\ref{4.22:not-live-region-geometry}, on the
geometry and counting of live large-field regions, whose proof is outstanding, together
with the inheritance theorem, gives, for the regions $C''$ and $Z^{c}$ of that statement,
\begin{equation}\label{4.22:eq-foreign-domain-membership-3}
Y^{\sharp}\subset C''\subset Z^{c}\subset\mathfrak{T},
\end{equation}
and $m(y)=m$ for every $y\in Y^{\sharp}$. Thus every location of $Y^{\sharp}$ is a location of
$\mathfrak{T}$ of shell $m<k$, and hypothesis (a) applies there.

By the inheritance theorem, under hypothesis (a) the stage-$l$ domain is contained in
$\mathfrak{N}^{(l)}(X)$ at every location of $X$: at locations of $\mathfrak{T}$ because the
threshold $\theta^{(l)}_a$ of the $a$-th clause of the stage-$l$ domain satisfies
$\theta^{(l)}_a\le r_a$ there, and at locations of the on-region $Z^{\mathrm{on}}$ because the
on-region threshold $3\varsigma_*$ lies below the class threshold $5\varsigma_*$, $\varsigma_*$
being the unit fixed in the inheritance theorem. Consequently the intersection with
$\mathfrak{N}^{(l)}(X)$ in the hypothesis imposes, beyond the clauses of the stage-$l$ domain,
only the ambient clause on the enlargement $X^{+}$ (the inheritance theorem).

\emph{Step (1): the input.} The hypothesis of the threshold bound (c) is that the stage-$l$
clause carries no weight, and this is hypothesis (a). For $y\in Y^{\sharp}$ the bound gives,
in terms of its source threshold $\theta_{\mathrm{src},a}$, its source factor $s^{(k)}_m$ and
its stage rates $\rho^{(l')}_m$,
\begin{align}
Q_a(\mathbb{U},\mathbb{J})(y)&<\theta^{(l)}_a(y)
=\theta_{\mathrm{src},a}(y)-\sum_{l'\le l}\rho^{(l')}_m r_a(y),
\label{4.22:eq-foreign-domain-membership-4}\\
\theta_{\mathrm{src},a}(y)&=s^{(k)}_m r_a(y)\qquad(a\le 7),\notag
\end{align}
and $Q_8<r_8$. For $a\in\{2,4\}$ the bound is needed at the structures $X'\subset Y^{\sharp}$
at which the clauses of $v$ are read. By hypothesis (b) these are among the structures over
$X$ at which the stage-$l$ clause is read, so
\eqref{4.22:eq-foreign-domain-membership-4} holds there as well. Since
$\theta^{(l)}_a\le r_a$ on $\mathfrak{T}$, in particular $Q_a(\mathbb{U},\mathbb{J})<r_a$ on
$Y^{\sharp}$ for every $a$.

\emph{Step (2): the move.} We apply (d), that is, clause (i) of the increment lemma. In terms of
its width $\mathsf{w}$, its layer fraction $\theta$ and its constants $K_{\min}$ and $\vartheta$,
its admissible range is
\[
K_{\min}\vartheta\le\mathsf{w}\le 1,\qquad \theta\in[\tfrac14,\tfrac12],
\]
and this range contains $(\mathsf{w},\theta)=(1,\tfrac12)$. Its block radii are
$R^{\natural}_{\square}=r^{\sharp}_{\square}$, so it covers all $|t_{\square}|\le
r^{\sharp}_{\square}$. The lemma coincides with the localised response corollary. Its input
class $\mathfrak{N}(X)$ contains $\mathfrak{N}^{(l)}(X)$, so $(\mathbb{U},\mathbb{J})$ is an
admissible input. Write
\[
\Delta_a(y):=Q_a(\Phi_{t,\sigma}(b))(y)-Q_a(\mathbb{U},\mathbb{J})(y)
\]
for the increment of entry $a$. The lemma gives three things.
\begin{itemize}
\item The decomposition of the image keeps the unitary part $U$ and the cube gauges
$\{u_{\square}\}$ for any shift, so $\Delta_5=\Delta_8=0$.
\item For $a\in\{1,3,6,7\}$, with $\iota_l(y)$ the per-stage increment fraction of the lemma,
\begin{equation}\label{4.22:eq-foreign-domain-membership-5}
\Delta_a(y)\le\iota_l(y)r_a(y)\le\tfrac16\rho^{(l)}_m r_a(y).
\end{equation}
\item For $a\in\{2,4\}$ the same bound holds by the last clause of the lemma, which applies
because, by Step (1),
\[
Q_a(\mathbb{U},\mathbb{J})<r_a\le 2\mathbf{r}^v_a,\qquad \mathbf{r}^v_a:=r_a .
\]
\end{itemize}

\emph{Step (3): the thresholds.} Since $l\ge1$, the term $\rho^{(l)}_m r_a(y)$ is one of the
summands subtracted in \eqref{4.22:eq-foreign-domain-membership-4}; the other summands are
nonnegative. Hence, for $a\in\{1,2,3,4,6,7\}$ and $y\in Y^{\sharp}$,
\begin{align*}
Q_a(\Phi_{t,\sigma}(b))(y)&=Q_a(\mathbb{U},\mathbb{J})(y)+\Delta_a(y)\\
&<\theta_{\mathrm{src},a}(y)-\sum_{l'\le l}\rho^{(l')}_m r_a(y)
+\tfrac16\rho^{(l)}_m r_a(y)\\
&\le\theta_{\mathrm{src},a}(y)=s^{(k)}_m r_a(y).
\end{align*}

For the region $W=Y$, whose creation step is $k_W=m$, the relation
\eqref{4.22:eq-frame-proof-datum-b} reads
\[
(1-\varpi_{\mathrm{fr}})-s^{(k)}_m=\beta\bigl(1-2^{-(k-m)}\bigr)-\varpi_{\mathrm{fr}}.
\]
Since $m<k$, we have $1-2^{-(k-m)}\ge\tfrac12$. Since
$\varpi_{\mathrm{fr}}=\tfrac18(1-\theta_S)\beta$ by the definition of the frame reserve
fraction, with $\beta$ the constant of \eqref{4.22:eq-frame-proof-datum-b}, and
$\theta_S\ge\tfrac12$, as in the proof of the frame lemma (\ref{4.22:prf-frame-proof-datum}),
we have
$\varpi_{\mathrm{fr}}\le\tfrac1{16}\beta$. Therefore
\[
(1-\varpi_{\mathrm{fr}})-s^{(k)}_m\ge\tfrac{\beta}{2}-\tfrac{\beta}{16}=\tfrac{7}{16}\beta.
\]
By Definition~\ref{4.22:def-kept-domain-list}, for $a\notin\{5,8\}$
\[
\theta_{v,a}(y)=r_a(y)\bigl(1-\varpi_{\mathrm{fr}}\mathbb{1}[\text{the stencil of }(y,a)
\text{ meets }\mathsf{B}_Y]\bigr)\ge(1-\varpi_{\mathrm{fr}})r_a(y),
\]
so that
\[
\theta_{\mathrm{src},a}(y)=s^{(k)}_m r_a(y)\le\theta_{v,a}(y)-\tfrac{7}{16}\beta r_a(y).
\]
Together with the previous chain this is \eqref{4.22:eq-foreign-domain-membership-2}.

For the remaining entries, $Q_5(\Phi_{t,\sigma}(b))=Q_5<r_5$ and
$Q_8(\Phi_{t,\sigma}(b))=Q_8<r_8$, by Steps (1) and (2). Here
$r_5=\theta_{v,5}$ and $r_8=\theta_{v,8}$.

Thus every clause of $\mathfrak{U}_v(Y^{\sharp})$ holds for
$\Phi_{t,\sigma}(b)\restriction Y^{\sharp}$, with thresholds $\theta_{v,a}$ that contain the
reserve $\varpi_{\mathrm{fr}}r_a$. By Step (2) this holds in the decomposition with the base's
$(U,\{u_{\square}\})$. This is \eqref{4.22:eq-foreign-domain-membership-1}.

\emph{Step (4): holomorphy.} Holomorphy in the parameters is the last sentence of clause (i)
of the increment lemma, hypothesis (d).

We note the following. Hypothesis (a) is used in Step (1). Suppose that at some
$y\in Y^{\sharp}$ the stage-$l$ clause
carried a weight, with one weight factor greater than one or the other equal to three. Then
its threshold $\theta^{(l)}_a(y)$ could exceed $\theta_{v,a}(y)$, which is at most $r_a(y)$,
and the conclusion \eqref{4.22:eq-foreign-domain-membership-1} would fail.

At an arrested step $j$, under the arrest condition, the increment bound of Step (2) holds with
the same $\iota_l$ at the radius $6.6C_1\delta'_j$ of that step. This follows from the clause of
the arrest lemma for a parent kept member $v\in\mathfrak{K}(Y)$.

The lemma supplies the part of the holomorphy requirement on the kept member $v$ that concerns
the field variables. The complementary part, holomorphy in the chart variables at width
$\mathsf{w}=1$, is a separate statement: clause (iv) of the holomorphy requirement, read at the
site of a step at which the large-field stages are performed, which is stated for all
$\mathsf{w}\le1$.
\end{proof}

\begin{lemma}[Kept terms add one summand to the probe functional]\label{4.22:lem-probe-summand}
Work in the setting of the anchor estimate of Definition~\ref{4.22:not-live-region-geometry}, with
probes $\wp$ of its class, strips $Y^{\sharp}$ of live large-field regions $(Y,m)$, and any depth
datum $\varrho(\cdot)$ satisfying condition (D2) of that definition. Denote by $\mathsf{S}$ the
functional bounded in clause (iv) of the anchor estimate, and by
$C_{\mathsf{a}}(\beta)(1+N_*)\mathsf{b}(N_*)$ the constant of its right member. For every $\beta>0$,
$\mathfrak{s}_{\Sigma}(\beta)$ and $\mathfrak{n}^{\mathrm{kept}}(\beta)$ are the numbers given by the
chapter's definition of the kept members' share; $\mathfrak{n}^{\mathrm{kept}}(\beta)$ is the
number of \eqref{4.22:eq-fixed-number-entry-a}. Assume the following.
\begin{itemize}
\item For every live $(Y,m)$, hypotheses $[\mathrm{V}]$, $(\mathrm{q}0)$, $(\mathrm{q}1)$ of
Lemma~\ref{4.22:lem-kept-oscillation} hold on $\mathfrak{S}_{\mathrm{src}}(Y^{\sharp})$.
\item Clauses (i2) and (b) of the geometric part of
Definition~\ref{4.22:not-live-region-geometry}: a probe meets the
strips of at most $2^d$ live large-field regions; a cell meeting $Y^{\sharp}$ has level at most $m$
and lies in $Y^{\sharp}$, strips of distinct live regions are disjoint, and a cell lies in at most
one of them.
\item Clause (iv) of the anchor estimate, stated in
Definition~\ref{4.22:not-live-region-geometry} with its proof outstanding.
\item Condition (D2): if $\wp$ meets the strip of a live $(Y,m)$, then
$\check{R}_m\le 2\varrho(\wp)+2$.
\end{itemize}
Then for every probe $\wp$ of the class and every $\beta>0$,
\begin{equation}\label{4.22:eq-probe-summand-1}
\begin{aligned}
\mathsf{S}^{\mathfrak{K}}(\wp)
&:=\sum\bigl\{\mathfrak{O}(v):\ v\in\mathfrak{K}(Y),\ (Y,m)\ \text{live},
\ \wp\ \text{meets}\ Y^{\sharp}\bigr\}\\
&\le 2^dC_{\Sigma}\bigl(2\varrho(\wp)+2\bigr)^{p_{\Sigma}}
\le \mathfrak{n}^{\mathrm{kept}}(\beta)\,e^{\beta\varrho(\wp)},
\end{aligned}
\end{equation}
and if $\wp$ is a cell the factor $2^d$ may be replaced by $1$. Consequently clause (iv) of the anchor
estimate holds for $\mathsf{S}+\mathsf{S}^{\mathfrak{K}}$, the constant of its right member being
replaced as follows:
\begin{equation}\label{4.22:eq-probe-summand-2}
C_{\mathsf{a}}(\beta)(1+N_*)\mathsf{b}(N_*)\ \longmapsto\
C_{\mathsf{a}}(\beta)(1+N_*)\mathsf{b}(N_*)+\mathfrak{n}^{\mathrm{kept}}(\beta).
\end{equation}
Thus the kept members are added to the bound of clause (iv), never multiplied into it.
\end{lemma}

\begin{proof}
We follow four steps.

\emph{Counting the strips.} Fix a probe $\wp$. By clause (i2), stated in
Definition~\ref{4.22:not-live-region-geometry} with its proof outstanding, $\wp$ meets the strips
$Y^{\sharp}$ of at most $2^d$ live large-field regions $(Y,m)$. If $\wp$ is a cell, clause (b) of
the geometric part of the same definition, stated there with its proof outstanding, applies
instead. A cell meeting $Y^{\sharp}$ lies in
$Y^{\sharp}$, the strips of distinct live regions are disjoint, and a cell lies in at most one strip.
Hence a cell meets the strip of at most one live region. The sum defining
$\mathsf{S}^{\mathfrak{K}}(\wp)$ therefore runs over the kept members of at most $2^d$ live
regions, and of at most one if $\wp$ is a cell.

\emph{One region.} Fix a live $(Y,m)$ whose strip $\wp$ meets. Apply clause (o) of
Lemma~\ref{4.22:lem-kept-oscillation} with $Y_i=Y$ and $W=Y$, so that $k_W=m$; its hypotheses
$[\mathrm{V}]$, $(\mathrm{q}0)$, $(\mathrm{q}1)$ hold on $\mathfrak{S}_{\mathrm{src}}(Y^{\sharp})$ by
assumption. The kept lineage $\mathfrak{K}(Y)$ consists of the two members $V(Y^{\sim})$ and the
kept quadratic form $\mathsf{Q}$ of $Y$. By the last bound in \eqref{4.22:eq-kept-oscillation-a},
these two members, counted together, satisfy
\begin{equation*}
\mathfrak{O}(V(Y^{\sim}))+\mathfrak{O}(\mathsf{Q})\le C_{\Sigma}\check{R}_m^{p_{\Sigma}}.
\end{equation*}

\emph{The depth datum.} Since $\wp$ meets $Y^{\sharp}$, condition (D2) gives
$\check{R}_m\le 2\varrho(\wp)+2$. Because $p_{\Sigma}>0$, this yields
$\check{R}_m^{p_{\Sigma}}\le(2\varrho(\wp)+2)^{p_{\Sigma}}$. Each region met thus contributes at
most $C_{\Sigma}(2\varrho(\wp)+2)^{p_{\Sigma}}$. Summing over the at most $2^d$ regions of the first
step gives the first inequality of \eqref{4.22:eq-probe-summand-1}, and over the at most one region
when $\wp$ is a cell, its variant with $2^d$ replaced by $1$.

\emph{The constant.} The second inequality of \eqref{4.22:eq-probe-summand-1} holds for every
$\beta>0$ by the definition of $\mathfrak{s}_{\Sigma}(\beta)$, the kept members' share, with
$\mathfrak{n}^{\mathrm{kept}}(\beta)$ as in \eqref{4.22:eq-fixed-number-entry-a}.

For the consequence, add \eqref{4.22:eq-probe-summand-1} to clause (iv) of the anchor estimate,
stated in Definition~\ref{4.22:not-live-region-geometry} with its proof outstanding, which bounds
$\mathsf{S}$. The kept members enter only through the additive term
$\mathfrak{n}^{\mathrm{kept}}(\beta)$ in the constant of the right member, as recorded in
\eqref{4.22:eq-probe-summand-2}.

No sum over levels occurs: a kept member has one level and one label.
\end{proof}

\begin{lemma}[Holomorphy of the kept terms along all later sites]\label{4.22:lem-kept-holomorphy}
Let $Y=W$ be a frame, that is, a second-class component, with $k_W=m$, and consider the pair
$(Y,m)$; write
$\Phi_Y:=\Phi_W$, $\Sigma_Y:=\Sigma_W$ and $\mathsf{B}_Y:=\mathsf{B}_W$. Let
$v\in\mathfrak{K}(Y)$. Write $k$ for the level of the step at which a site is read. Assume the
following statements.
\begin{itemize}
\item[(I)] The containment statements of the correlation theorem. Among them are the
following six; item $(j)$ of this list is referred to as (I.$j$).
\begin{itemize}
\item[(1)] the membership statement for stored terms: for every stored term $F$, at every
link of the site clause by which wrapped terms are run, on the whole of its parameter
polydisc, source polydisc and output tube,
\begin{equation*}
(\mathsf{R}(\tau)^{h^{\mathfrak{b}}})\restriction X_F\in\mathfrak{U}_F(X_F),
\end{equation*}
the stored domain of $F$, where $\mathsf{R}(\tau)$ and $h^{\mathfrak{b}}$ are the family and
the cut-off of the link (see below), the cases being the first and the later fluctuation
integrals $(\mathrm{T})$, the stages $(\mathrm{P})$, and the foreign and the own links of the
presentation;
\item[(2)] the homogeneity bound, together with the separate statement fixing the layer
schedule at $(\mathrm{T})$, on every layer $n\le k$;
\item[(3)] the one-sixteenth bound, the quantitative form of the containment statement at the
presentation: each entry is at most $\frac1{16}$ of its threshold on
$X_F\cap(\Lambda^{\natural}\cup\mathfrak{N}^{\mathrm{rec}})$, where $X_F$ is the domain of the
stored term $F$ at which the bound is applied; here $\Lambda^{\natural}$ is the
region of the presentation that contains the core $\mathsf{C}_C$ of the treated component,
and $\mathfrak{N}^{\mathrm{rec}}$ is the further region on which the bound is stated with it;
\item[(4)] at the presentation $(\mathrm{L})$, the increments of the containment statement
whose quantitative form is (I.3), and the increments of clause $(\beta)$ of the inheritance
claim;
\item[(5)] the containment statement of the correlation theorem on the own region
$Z^{\mathrm{on}}$, the presentation region of the own component;
\item[(6)] the source-polydisc containment clause of the site clause.
\end{itemize}
\item[(II)] The lemma stating that the containment increments at $(\mathrm{L})$ hold for
every set $\sigma_0$ disjoint from the region $\Lambda$ named in that lemma.
\item[(III)] The three inequalities \eqref{4.22:eq-frame-proof-datum-b} established in the
proof of the frame lemma, in the part showing that nothing moves off $\mathsf{B}_W$ and fixing
the orbit datum at the roots.
\item[(IV)] Clause (iii) of the law-site lemma and its display
\eqref{4.22:eq-law-site-b}, together with the theorem that the radius budget closes
uniformly in all scales, in the form \eqref{4.22:eq-budget-closes-a},
\eqref{4.22:eq-budget-closes-b}.
\item[(V)] For conclusion (c):
\begin{itemize}
\item the amplitude bound and the core clause of the analysis of presented triples;
\item the lemma on the two presented triples;
\item the membership lemma for presented triples, in the form established by its proof: its
conclusion $P_\lambda\in\mathfrak{D}_{\mathfrak{E}}$ holds for every unit $\mathfrak{E}$ with
domain $\mathfrak{D}_{\mathfrak{E}}$ such that, at every location $y$ of the domain and for
every entry $a$, the following four conditions hold:
\begin{itemize}
\item the threshold inequality $\theta_a(y)\ge\varphi_{\flat}\mathbf{r}_a(n(y))$, where
$\theta_a(y)$ is the threshold that the domain $\mathfrak{D}_{\mathfrak{E}}$ imposes on the
entry $a$ at $y$, $n(y)$ the clause index at $y$, $\varphi_{\flat}\le1$ the threshold fraction
of that proof and $\mathbf{r}_a(n)$ the threshold of the entry $a$ at clause index $n$;
\item the base inequality $Q_a(P_0)(y)\le\frac1{16}\theta_a(y)$, where $P_0$ is the base pair
of the membership lemma for presented triples;
\item the clause-index inequality $n(y)\le\pi(y)$, where $\pi(y)$ is the level assigned to
$y$ by Ba{\l}aban's determining set $\mathbf{B}_k$;
\item the age inequality of the display of the analysis of presented triples, in its form at a
location where the depth factor equals one,
\begin{equation*}
e^{-\frac12\delta_Pd^{\wedge}(y,\mathsf{C}_C^{\,c})}\,\frac{\alpha_{*,k}}{\alpha_{*,n(y)}}
\le C_{\mathrm{age}},
\end{equation*}
where $\delta_P$ is the decay constant of that analysis, $d^{\wedge}$ the weighted contour
distance, and $C_{\mathrm{age}}$ the age constant of that display.
\end{itemize}
\end{itemize}
\item[(VI)] At the stages $(\mathrm{P})$:
\begin{itemize}
\item the declaration of the kept terms in the domain list
(Definition~\ref{4.22:def-kept-domain-list});
\item clauses (i) and (iv) of the lemma of the inheritance analysis that bounds the entries
under a simultaneous shift of all blocks;
\item from the proof of the inheritance statement, the step bounding the increment of a
stage, the step concluding membership, and the bound on the sum of the stage increments.
\end{itemize}
\item[(VII)] On the run of the preparatory chain on which the stage link is taken in
Ba{\l}aban's form, the weight convention under which Ba{\l}aban's clause is unweighted, with
unit weights at every location of shell $m<k$, together with the locality of the tube
clauses ((X) below), through Lemma~\ref{4.22:lem-foreign-domain-membership}.
\item[(VIII)] If $v=\mathsf{Q}_i$: the two bounds on the kept quadratic form $\mathsf{Q}_i$
that are hypotheses of Lemma~\ref{4.22:lem-kept-oscillation}.
\item[(IX)] The linear form of the containment proofs. Each containment statement in (I)
proves membership as
\begin{itemize}
\item the entries of the base, with room,
\item plus increments of the entries $1$--$4$, $6$, $7$ that are linear in a majorant of the
layer and are taken at fixed unitary part and fixed cube gauges;
\end{itemize}
the entries $5$, $8$ are those of the base.
\item[(X)] The locality of the clauses of the tube. The stored tubes
$\tilde{U}^{c}_m(\cdot,\tilde\alpha_0,\tilde\alpha_1)$ are cut, as the spaces of
\cite[(2.34)--(2.39)]{B-CMP119} are, by clauses at locations, each reading the pair on a
stencil in the domain. This property is here adopted also for the stored domains
$\mathfrak{U}_v(Y^{\sharp})$ with their reserve.
\end{itemize}
Here $Q_a$, $a=1,\dots,8$, are the eight entries of these clauses and $r_a$ their
thresholds.

Consider every site after
$\mathbf{R}_m$ up to and including the presentation
$(\mathrm{L})$ that treats the component of $Y$. These sites are:
\begin{itemize}
\item the fluctuation integral $(\mathrm{T})$ of each step, run under the
fluctuation-integral part of the site clause;
\item each stage of each chain $(\mathrm{P})$, on the component of $Y$ or on a foreign one;
\item each link of $(\mathrm{L})$, run under the two presentation parts of the site clause,
either at a foreign component (written $(\mathrm{L}^{\circ})$) or at the own component
(written $(\mathrm{L}^{\dagger})$).
\end{itemize}
At each such site, write $\mathsf{R}(\tau)$, $\tau\in\mathfrak{P}$, for the family of the
site, $h^{\mathfrak{b}}$ for its cut-off (written $h^{\mathfrak{b}}_q$ at the own component)
and $\mathfrak{u}(\tau)$ for its riding map.
The parameters collected in $\tau$ are the site's variable $\mathsf{s}$ at $(\mathrm{T})$, the
triple $(t,\sigma,b)$ of the stage at $(\mathrm{P})$, and the tuple $\vec{\mathsf{s}}$ of the
link at $(\mathrm{L})$, together with $\varpi$ at $(\mathrm{L}^{\dagger})$. Let $\eta_0(s)$ be
the debit at the law site $s$
in the sense of the law-site lemma (not to be confused with the fine-lattice spacing $\eta$
of (b)) and $\mu(s)$ the
margin at the law site $s$, and let $\varepsilon(s)=\varepsilon_k(s)$ be the current radius
at $s$. Then the following hold on the whole of the
site's parameter polydisc $\mathfrak{P}$, source polydisc and output tube.
\begin{itemize}
\item[(a)] The family satisfies
\begin{equation}\label{4.22:eq-kept-holomorphy-1}
(\mathsf{R}(\tau)^{h^{\mathfrak{b}}})\restriction Y^{\sharp}\in\mathfrak{U}_v(Y^{\sharp}),
\end{equation}
the stored domain of $v$, and $h^{\mathfrak{b}}=1$ on $Y^{\sharp}$. The map
\begin{equation*}
(\tau,\dots,g')\longmapsto
v^{\sharp}\bigl(\mathsf{R}(\tau)\restriction Y^{\sharp};\Phi_Y,\mathfrak{u}(\tau)g'\bigr),
\end{equation*}
the dots standing for the source and tube variables, is holomorphic on
\begin{equation*}
\mathfrak{P}\times(\text{sources})\times(\text{tube})
\times\mathfrak{G}_{\varepsilon-\eta_0-\mu}(\Sigma_Y).
\end{equation*}
Moreover
\begin{equation*}
\bigl(\mathsf{R}(\tau)\restriction Y^{\sharp},\mathfrak{u}(\tau)g'\bigr)\in
\mathfrak{U}_v\times\mathfrak{G}_{\varepsilon}(\Sigma_Y),
\end{equation*}
and $\varepsilon\le\varepsilon^{\mathrm{fr}}(Y)$.
\item[(b)] The membership in (a) is witnessed in the decomposition
\begin{equation*}
\mathsf{R}(\tau)=e^{i\eta A'(\tau)}U_{\mathbf{K}},
\end{equation*}
where $\eta$ denotes the spacing of the fine lattice on which the configurations
$\mathsf{R}(\tau)$ live and $A'(\tau)$ is $\mathfrak{g}^{\mathbb{C}}$-valued,
with the $\tau$-free decomposition $\mathsf{r}=(U_{\mathbf{K}},\{u_{\square}\})$ of the base
$\mathbf{K}$:
\begin{itemize}
\item the entries $5$, $8$ are those of the base;
\item the clauses on the entries $1$--$4$, $6$, $7$ hold in that decomposition, with the
reserve $\varpi_{\mathrm{fr}}r_a$ at the locations of $\mathsf{B}_Y$.
\end{itemize}
Here $U_{\mathbf{K}}$ is the $G$-valued part of:
\begin{itemize}
\item the point $\mathbf{U}$ of the output tube, at $(\mathrm{T})$;
\item $\mathbb{U}$, at $(\mathrm{P})$;
\item the field $\mathbf{U}$ carried by the unit presented at the foreign link, at
$(\mathrm{L}^{\circ})$;
\end{itemize}
and, at $(\mathrm{L}^{\dagger})$, $U_{\mathbf{K}}:=U_2(\varpi)$, the base configuration of
the own link in the analysis of presented triples.
\item[(c)] At $(\mathrm{L}^{\dagger})$, put $X:=Y^{\sharp}\subset\mathsf{C}_C$. Let
$\mathfrak{z},\mathfrak{z}'\in\mathfrak{P}$ have the same $\varpi$, and let $|\lambda|\le\rho_X$,
where $\rho_X$, the radius of the $\lambda$-disc, is
\begin{equation*}
\rho_X=r^0e^{\frac12\delta_PD_X},
\end{equation*}
with $\delta_P$ the decay constant and $D_X$ the depth of $X$ in the analysis of presented
triples, and $r^0$ the constant factor of that radius.
Let $P_\lambda$ be the pair of the membership lemma for presented triples formed at
$\mathfrak{z}$, $\mathfrak{z}'$. Then $P_\lambda\in\mathfrak{U}_v(Y^{\sharp})$, with unitary
part $\bar U=U_2$. At every $y\in Y^{\sharp}$ each entry is at most $\tfrac3{32}+\tfrac12$
times its threshold, and $\tfrac3{32}+\tfrac12<1-\varpi_{\mathrm{fr}}$. No age enters this
bound.
\end{itemize}
\end{lemma}

\begin{proof}
\emph{Step~0: the stored domain is local.} By (X), the tube
\begin{equation*}
\mathfrak{U}^{\mathrm{tr}}_v(Y^{\sharp})=\tilde{U}^{c}_m(Y^{\sharp},\tilde\alpha_0,\tilde\alpha_1)
\end{equation*}
and the stored domain $\mathfrak{U}_v(Y^{\sharp})\subset\mathfrak{U}^{\mathrm{tr}}_v(Y^{\sharp})$
are cut by clauses at locations, as the spaces of \cite[(2.34)--(2.39)]{B-CMP119} are. There
the clauses are imposed on the sets $X\cap(\Omega_n\setminus\Omega_{n+1})$, each clause
reading the pair on a stencil in the domain, through gauge transformations defined on
$\square\cap X$.
The composite entries $2$ and $4$ are imposed at every admissible triple $X'$ contained in
the domain. Thus the stored domains are cut by local, shell-wise clauses, and they are
defined over any union of cells.

Let $X\supset Y^{\sharp}$ be a union of cells. Every clause of
$\tilde{U}^{c}_m(Y^{\sharp})$ is a clause of $\tilde{U}^{c}_m(X)$, read on a smaller or equal
set. A decomposition and its cube gauges restrict to $Y^{\sharp}$. The reserve in the
definition of the stored domains stands at the same locations, since $\mathsf{B}_Y\subset Y$.
Hence, for every term $N$ created by $\mathbf{R}_m$ with $Y\subset X$,
\begin{equation}\label{4.22:eq-kept-holomorphy-2}
P\in\mathfrak{U}_N(X)\ \Longrightarrow\ P\restriction Y^{\sharp}\in\mathfrak{U}_v(Y^{\sharp}).
\end{equation}
Restriction in this sense is used in the same way in \cite{B-CMP116}.

\emph{(a), the field half.} Let $\Delta_1$ be a cube of the level-$m$ partition $\pi_m$
sharing a wall with $Y^{\sharp}$, and put
\begin{equation*}
F_0:=\mathbf{B}'^{(m)}\bigl(\operatorname{sk}(Y^{\sharp}\cup\Delta_1),\cdot\bigr),
\end{equation*}
where $\operatorname{sk}(Y^{\sharp}\cup\Delta_1)$ is the set attached to $Y^{\sharp}\cup\Delta_1$
in the analysis of the $\mathbf{R}$-born boundary terms, at which the boundary term
$\mathbf{B}'^{(m)}$ is evaluated.
The properties of $F_0$ are as follows.
\begin{itemize}
\item $F_0$ is a term of the format class (the class of stored terms carrying the data listed,
for the kept terms, in Definition~\ref{4.22:def-kept-term-reader}: label, true domain, the set
$\mathcal{P}_{F}$ of frames $(Y,m)$ to which the term is attached, and stored domain), with
$(Y,m)\in\mathcal{P}_{F_0}$, by clause (c) of
the birth-format proposition of the analysis of the $\mathbf{R}$-born boundary terms.
\item It carries these data from its creation: by the same clause there is exactly one such
term for each kind of boundary term, each label and each setting, whether or not it vanishes.
\item It is stored unchanged as long as $(Y,m)$ is live, by the last line of clause (e) of the
storage proposition and by clause (b) of the proposition on the stored terms of live frames
of the same analysis; and it is a wrapped unit at every site.
\item Its stored domain is, by the definition of the stored domains,
\begin{equation*}
\mathfrak{U}_{F_0}(X_{F_0})\subset\tilde{U}^{c}_m(X_{F_0},\tilde\alpha_0,\tilde\alpha_1),
\qquad X_{F_0}\supset Y^{\sharp}.
\end{equation*}
\end{itemize}
Apply the membership statement (I.1) to the stored term $F_0$. At every link of the site
clause, on the whole of its parameter polydisc, source polydisc and output tube,
\begin{equation*}
(\mathsf{R}(\tau)^{h^{\mathfrak{b}}})\restriction X_{F_0}\in\mathfrak{U}_{F_0}(X_{F_0}).
\end{equation*}
The cases of that statement are the first and the later fluctuation integrals, the stages,
and the links $(\mathrm{L}^{\circ})$ and $(\mathrm{L}^{\dagger})$. The term $F_0$ is created
by $\mathbf{R}_m$ and $X_{F_0}\supset Y^{\sharp}$; the implication
\eqref{4.22:eq-kept-holomorphy-2} is applied with $N=F_0$, $X=X_{F_0}$ and
$P=(\mathsf{R}(\tau)^{h^{\mathfrak{b}}})\restriction X_{F_0}$, and gives
\eqref{4.22:eq-kept-holomorphy-1}.

The same membership is also obtained clause by clause; in this route the term $F_0$ enters
only at the own stage. The proofs of
the membership statement (I.1) and of the source-polydisc containment clause (I.6) of the
site clause verify the clauses one at a time at the locations $y$ of $X_F$, and of a stored
term $F$ they use only the information which clause is imposed at $y$. The clauses are
obtained as follows.
\begin{itemize}
\item Off $\mathsf{B}_Y$, at $(\mathrm{T})$: the homogeneity bound (I.2) with its schedule on
every layer $n\le k$.
\item At a foreign stage $(\mathrm{P})$. Here
\begin{equation*}
Y\subset C''\subset Z^{c}\subset\Omega_k\cup Z^{c},
\end{equation*}
by the geometry of the scaffold (the untreated component $C''$ of $\Lambda_k^{c}$ holding
$Y$ contains $Y^{\sharp}$); the last set is, by definition, the region on which the
inheritance analysis is carried. The conclusion
\begin{equation*}
\Phi_{t,\sigma}(b)\restriction Y^{\sharp}\in\mathfrak{U}_v(Y^{\sharp}),
\end{equation*}
witnessed by the decomposition $(U,\{u_{\square}\})$ of the base, arrests included, follows
from the declaration of Definition~\ref{4.22:def-kept-domain-list}, by the remarks following
it, which carry out the steps of (VI) for the term $v$. This uses neither (I.5), nor the
localisation item of the inheritance analysis on which the localised membership corollary
below rests, nor the term
$F_0$. So at every foreign stage the appeal to (I.5) is not needed:
\begin{itemize}
\item on the run on which the kept terms are declared in the domain list, by
Definition~\ref{4.22:def-kept-domain-list};
\item on the run on which the stage link is taken in Ba{\l}aban's form, by
Lemma~\ref{4.22:lem-foreign-domain-membership} under (VII);
\item alternatively, on that run, by appeal to the localised membership corollary and
that localisation item, combined with Step~0; this alternative is
not used in what follows.
\end{itemize}
\item At the own stage $(\mathrm{P})$. Here $Y^{\sharp}\subset\mathsf{C}_C\subset
Z^{\mathrm{on}}$. The domain family $\mathfrak{L}_n$ is cut by local clauses on
$Z^{\mathrm{on}}$, and the membership of $F_0$, restricted to $Y^{\sharp}$ by Step~0,
supplies the clauses of $\mathfrak{L}_n$ there. The containment on $Z^{\mathrm{on}}$ at
$F_0$ is hypothesis (I.5), applied to $F_0$. On
either run, this is the only place where (I.5) is used.
\item At $(\mathrm{L})$: the containment increments (I.4), each for every $\sigma_0$
disjoint from the region $\Lambda$ named in (II), by (II).
\item At $\mathsf{B}_Y$: the three lines \eqref{4.22:eq-frame-proof-datum-b} at amplitude one.
\end{itemize}
The locations of $Y^{\sharp}$ are among these locations.

\emph{The cut-off.} Consider first a foreign component. By the definition of the cut-off of
the presentation and the geometry of the scaffold,
\begin{equation*}
h^{\mathfrak{b}}=1\ \text{ on }\ \{\operatorname{dist}(\cdot,\Lambda^{\natural})\ge\tfrac12\}
\supset Y^{\sharp}.
\end{equation*}
At the own component, $h^{\mathfrak{b}}_q=1$ on $D_2\supset\mathsf{C}_C$ for
$\mathfrak{b}\ge\mathfrak{b}_2$, where $D_2$ is the second of the nested domains of the
presentation at the own component and $\mathfrak{b}_2$ is the second of the chain thresholds
of the presentation. In either case the cut-off equals $1$ on $Y^{\sharp}$.

\emph{(a), the chart half.} By clause (i) of the riding-map lemma, $\mathfrak{u}(\tau)(s)$
depends only on the layer on $B^{\mathrm{top}}(s)\subset Y$, the top block of the law site
$s$. The bound
\begin{equation*}
\sup_{\mathfrak{P}}\operatorname{pol}\bigl(\mathfrak{u}(\tau)(s)\bigr)\le\eta_0(s)+\mu(s)
\end{equation*}
holds at the fluctuation integrals and at the foreign links by clause (iii) of the law-site
lemma and \eqref{4.22:eq-law-site-b}, which are stated at every law site in the chart of a
unit run by the site clause, whatever its class; at the stages $(\mathrm{P})$ by clause (iv)
of the riding-map lemma; and at the own link $(\mathrm{L}^{\dagger})$ by the profile bound of
the led rotations.
Let $g'\in\mathfrak{G}_{\varepsilon-\eta_0-\mu}(\Sigma_Y)$. For invertible matrices
$\Gamma_1,\Gamma_2$ one has $\|\Gamma_1\Gamma_2\|\le\|\Gamma_1\|\|\Gamma_2\|$ and
$\|(\Gamma_1\Gamma_2)^{-1}\|\le\|\Gamma_2^{-1}\|\|\Gamma_1^{-1}\|$, so the polar size is
subadditive under products.
Hence, at every $s\in\Sigma_Y$,
\begin{align*}
\operatorname{pol}\bigl((\mathfrak{u}(\tau)g')(s)\bigr)
&\le\operatorname{pol}(\mathfrak{u}(\tau)(s))+\operatorname{pol}(g'(s))\\
&<\eta_0(s)+\mu(s)+\varepsilon(s)-\eta_0(s)-\mu(s)=\varepsilon(s).
\end{align*}
By the theorem that the radius budget closes uniformly in all scales,
\eqref{4.22:eq-budget-closes-a}, one has $\varepsilon\le\mathsf{e}=\varepsilon^{\mathrm{fr}}(Y)$.
Thus $\mathfrak{u}(\tau)g'\in\mathfrak{G}_{\varepsilon}(\Sigma_Y)$.

\emph{Holomorphy.} By the holomorphy clause of part (ii) of the frame lemma, the datum
$v^{\sharp}$ is holomorphic in its arguments
jointly on $\mathfrak{U}_v\times\mathfrak{G}_{\varepsilon^{\mathrm{fr}}(Y)}(\Sigma_Y)$, which
contains $\mathfrak{U}_v\times\mathfrak{G}_{\varepsilon}(\Sigma_Y)$ since
$\varepsilon\le\varepsilon^{\mathrm{fr}}(Y)$. For $v=\mathsf{Q}_i$ this
rests on the two bounds of (VIII). The map of (a) is the composition of $v^{\sharp}$ with the
maps $(\tau,\text{sources},\text{tube})\mapsto\mathsf{R}(\tau)\restriction Y^{\sharp}$ and
$(\tau,g')\mapsto\mathfrak{u}(\tau)g'$, holomorphic in all these variables. By the two
preceding paragraphs these maps take values
in $\mathfrak{U}_v\times\mathfrak{G}_{\varepsilon}(\Sigma_Y)$. A composition of holomorphic
maps is holomorphic, as in the holomorphy proofs for stored terms at the sites and at the
stages. This proves (a).

\emph{(b).} By the linear form (IX), each containment statement used in (a) proves membership
as the entries of the base, with room, plus increments of the entries $1$--$4$, $6$, $7$.
These increments are linear in a majorant of the layer and are taken at fixed unitary part
and fixed cube gauges. The entries $5$ and $8$ are those of the base, because the unitary part
and the cube gauges are kept. The arithmetic of \eqref{4.22:eq-frame-proof-datum-b} is
arithmetic of these increments, and it keeps the reserve, through the three inequalities
\begin{equation*}
\tfrac14+\tfrac18\le\tfrac12,\qquad r_T\ge\tfrac12\theta_S\beta,\qquad
\tfrac1{16}<1-\varpi_{\mathrm{fr}};
\end{equation*}
here $r_T$ denotes the room left, in threshold units, at the stages and the links, $\beta$ is
the schedule constant of Definition~\ref{4.4:def-post-step-space} and $\theta_S$ the layer
number.
So the clauses on the entries $1$--$4$, $6$, $7$ hold in
the decomposition with the base's unitary part and cube gauges, with the reserve
$\varpi_{\mathrm{fr}}r_a$ at the locations of $\mathsf{B}_Y$.

It remains to see that the base does not depend on $\tau$.
\begin{itemize}
\item At $(\mathrm{T})$, the point of the output tube is a parameter beside the site's
variable $\mathsf{s}$.
\item At $(\mathrm{P})$, $\mathbb{U}$ is a parameter beside $(t,\sigma,b)$.
\item At $(\mathrm{L}^{\circ})$, off the box $\Lambda^{\sim}$ of the geometry of the scaffold,
the chain of the link is
\begin{equation*}
\operatorname{chain}(\vec{\mathsf{s}})
=e^{i\eta\mathbb{H}_4}e^{i\eta\mathbb{H}_3}e^{i\eta\mathbb{H}_2}e^{i\eta\mathbb{H}_1}\mathbf{U}
=:e^{i\eta\mathcal{K}(\vec{\mathsf{s}})}\mathbf{U},
\end{equation*}
with $\mathbf{U}$ the field carried by the unit presented at the foreign link.
\item At $(\mathrm{L}^{\dagger})$, on $D_2$, the chain is
$e^{i\eta\mathbb{H}_4}e^{i\eta\mathbb{H}_3}U_2(\varpi)$. By the bounds of the analysis of
presented triples, $U_2$ is real and independent of the presented unit and of
$\vec{\mathsf{s}}$.
\end{itemize}
This proves (b).

\emph{(c).} Let $h$ denote the level of the component $C$ treated at $\mathbf{R}_k$, as in the
geometry of the scaffold, so that $m\le h$; from here on the letter $h$ is used only in this
sense. Apply the membership lemma for presented triples in the form (V), with the domain
$\mathfrak{D}_{\mathfrak{E}}:=\mathfrak{U}_v$; it suffices to verify the four conditions of
(V) at every $y\in Y^{\sharp}$. Of the unit $\mathfrak{E}$ that form reads only the threshold
inequality, the base inequality and the clause-index inequality; the age inequality concerns
the location $y$ and is verified below from the position of $Y^{\sharp}$ in $\mathsf{C}_C$.
The remaining input of that proof is supplied by the lemma on the two presented triples,
stated on the enlargement $X^{+}\subset D_2$ of $X$, and concerns only those two triples.
The linear form (IX), the
bound \eqref{4.22:eq-budget-closes-b} and the amplitude bound read nothing of $\mathfrak{E}$.

\emph{The indices.} On $Y^{\sharp}$ both indices equal $m$: $n(y)=m$, $\theta_a=r_a$ (factor
one), and $\pi(y)=m$. For the last, by definition \cite[(1.33)]{B-CMP122b}, $\mathbf{B}_k$
equals $\mathbf{B}''_k$ on $Z^{c}$ and $\{Z^{(k)}\}$ on $Z$. Take $k=m$; this holds on every
component of $Z$, second-class components included, and the new background $U_m$ is built on
$\mathbf{B}_m$. The
following facts then give $\pi(y)=m$:
\begin{itemize}
\item the strip lies in $\Lambda_m\setminus\Omega_{m+1}$, that is, in zone $m$;
\item nothing re-blocks $Y^{\sharp}$ before the presentation of its own component
\cite{B-CMP122a}, and the geometry of the scaffold gives the own collar;
\item on the core, the modified small-field regions $\Omega''_j$ of \cite{B-CMP122a} coincide
with $\Omega_j$ for $j\le h$.
\end{itemize}
The geometric statement that a cell meeting $Y^{\sharp}$ has level at most $m$ is the safe
form of this.

\emph{Threshold inequality.} The threshold $r_a$ is $\tilde\alpha_{\cdot,m}$ times the
threshold number of the entry $a$ in \cite[(2.39)]{B-CMP119}, and
$r_a\ge\varphi_{\flat}\mathbf{r}_a(m)$, since
$\varphi_{\flat}\le1$. In the notation of the inheritance analysis,
$r_a(y):=\mathbf{r}_{a,m(y)}$, and $n(y)=m$ on $Y^{\sharp}$. So the threshold inequality
$\theta_a=r_a\ge\varphi_{\flat}\mathbf{r}_a(m)$ is an identity followed by an inequality with
room. In
\cite[(2.34)--(2.39)]{B-CMP119} slightly larger spaces are introduced for the terms
$E^{(k)}$, $R^{(k)}$, $B^{(k)}$ born in $\mathbf{T}$, after the operation $\mathbf{T}$. The
domain of $v$ belongs to terms born in $\mathbf{R}$ \cite[(1.68), (1.99)]{B-CMP122b}, so that
remark is cited here by analogy only.

\emph{Age inequality.} Take the age pair $(a',b'):=(h-m,0)$. Then
$k-n(y)=N_{\mathrm{prep}}+a'$, where $N_{\mathrm{prep}}$ is the number of preparatory
operations. A contour from $y\in Y^{\sharp}$ to $\mathsf{C}_C^{\,c}$ crosses the zones
$m+1,\dots,h-1$. Since a contour pays at least $\mathsf{w}_{\mathrm{zone}}$ per crossed zone,
\begin{equation*}
d^{\wedge}(y,\mathsf{C}_C^{\,c})\ge\mathsf{w}_{\mathrm{zone}}(a'-1)_+ .
\end{equation*}
With the depth factor at $y$ equal to one, the display of the analysis of presented triples
gives
\begin{equation*}
e^{-\frac12\delta_Pd^{\wedge}(y,\mathsf{C}_C^{\,c})}\,\frac{\alpha_{*,k}}{\alpha_{*,m}}
\le2(1+B_\sharp)^{1/2}(1+a')^{1/2}e^{-2(a'-1)_+}\le C_{\mathrm{age}} .
\end{equation*}
Since $n(y)=m$, this is the age inequality of (V).

\emph{Base inequality.} Apply the one-sixteenth bound (I.3) to
$\mathfrak{E}:=F_0\in\mathcal{C}$. This is legitimate
because $F_0$ is a term $\mathbf{B}'$ born in $\mathbf{R}$. Each entry is then at most
$\frac1{16}$ of its threshold on $X_{F_0}\cap(\Lambda^{\natural}\cup\mathfrak{N}^{\mathrm{rec}})$,
a set which contains $Y^{\sharp}$.
On $Y^{\sharp}\subset\mathsf{C}_C\subset\Lambda^{\natural}$ the thresholds of $F_0$ are those of
$v$, by Step~0.

\emph{Conclusion.} The membership lemma for presented triples, in the form (V), now bounds
every additive increment by $\frac12$ in threshold units. Hence, at every $y\in Y^{\sharp}$:
\begin{itemize}
\item the entries $1$, $3$, $6$, $7$ are at most $\frac1{16}+\frac12$ of their thresholds;
\item the entries $2$, $4$ are at most $\frac1{16}\cdot\frac32+\frac12=\frac3{32}+\frac12$;
\item the entries $5$, $8$ are at most $\frac1{16}$, with $\bar U=U_2$.
\end{itemize}
By the definition of the stored domains, $\varpi_{\mathrm{fr}}\le\frac1{80}$, and therefore
\begin{equation*}
\tfrac3{32}+\tfrac12<1-\tfrac1{80}\le1-\varpi_{\mathrm{fr}} .
\end{equation*}
So $P_\lambda\in\mathfrak{U}_v(Y^{\sharp})$, each entry is at most $\frac3{32}+\frac12$ of its
threshold at every $y\in Y^{\sharp}$, and no age of the frame is read in this bound. This
proves (c).
\end{proof}

\begin{proposition}[Kept terms at a fluctuation site]\label{4.22:prop-fluctuation-site}
Consider the fluctuation site $(\mathrm{T})$ of the step $i\to i+1$. Throughout, a large-field region
is called \emph{live} if it is live at the moment of this site. Assume the following.
\begin{itemize}
\item The two-point bound on the kept members and the two hypotheses on the kept quadratic form, on
the sourced set $\mathfrak{S}_{\mathrm{src}}(Y^{\sharp})$ of every live region $(Y,m)$; and the
membership statement for the stored domains of the kept members at the large-field stages.
\item The conclusions of Lemma~\ref{4.22:lem-kept-holomorphy} (holomorphy of the kept terms along all
later sites), and the bounds of Lemma~\ref{4.22:lem-kept-oscillation} (oscillation of the kept terms
on the orbit chart) in the form \eqref{4.22:eq-kept-oscillation-a}.
\item The proposition on the fluctuation site for the plain and relative terms (without kept
members), parts (a)--(e); the summation lemma for readers at a
site, parts (a), (c), (d); the anchor lemma, part (iv); the base-point lemma, part (ii); the
source-interpolation statement, parts (i)--(iv); the transfer lemma for orbit data, part (iii).
\item The statements (a)--(d) of the conventions~\ref{4.22:not-live-region-geometry} on the geometry
and counting of live large-field regions, whose proof is outstanding.
\item The floor
\begin{equation}\label{4.22:eq-fluctuation-site-1}
\kappa_1\ \ge\ 1+c_l(2\kappa+4)+\lambda_*,\qquad c_l=1+234d,
\end{equation}
with $\lambda_*=\lambda_*(d,L)$ the constant of the summation lemma; the floor on the zone weight;
the further floors on $M$ and on $\kappa_1$ and the ceilings on the auxiliary parameters under
which the present section is set; and the inequality $(1-2\delta)\kappa\ge c_d+1$, with $c_d$ and
$K(d)$ the constants of \cite[(1.26)]{B-CMP116}.
\end{itemize}
The site geometry is that of the proposition on the fluctuation site: the members form the family
$\pi^{\mathfrak{s}}:=\pi^T$ of all cells; the set $\mathfrak{H}$ of hubs is empty; the source cells
form the family $\mathfrak{F}$ of level-$i$ cells holding free bonds of $\Omega_{i+1}$ (in this
proposition $\mathfrak{F}$ denotes this family of source cells, not a family of frames). The
interpolation parameter is $\mathsf{s}=(\mathsf{s}(\Delta))_{\Delta\in\pi^T}$; the interpolated
configuration $\mathsf{R}(\mathsf{s})=(e^{i\eta\mathcal{Y}(\mathsf{s};b)}\mathbf{U},\dots)$ and the
riding map $\mathsf{u}(\mathsf{s})$ are those of that proposition. We write $\mathbb{1}$ and $0$ for
the constant families $\mathsf{s}(\Delta)=1$, resp.\ $\mathsf{s}(\Delta)=0$, for all $\Delta$. For a
member $\Delta$, $\rho(\Delta)$
is the least cardinality of a connected subfamily of $\pi^T$ holding $\Delta$ and a source cell; a
member \emph{sees} a set $\mathcal{R}$ if it meets $\mathcal{R}$, or is a source cell sharing a wall
with a hub that meets $\mathcal{R}$.

For a live region $(Y,m)$ with $m\le i$ and $v\in\mathfrak{K}(Y)$, the kept member $v$ is a reader with
label $Q_Y(i)$, read region $Y^{\sharp}$ and norm $\mathfrak{O}(v)$
(Definition~\ref{4.22:def-kept-term-reader}). Put
\begin{equation*}
\mathfrak{L}^{\mathfrak{K}}_T:=\bigl\{(Q_Y(i),Y^{\sharp},\mathfrak{O}(v)):\ v\in\mathfrak{K}(Y),\
(Y,m)\ \text{live},\ m\le i\bigr\},
\end{equation*}
and, for a family $\mathfrak{L}$ of readers $F=(A_F,\mathcal{R}_F,\|F\|)$ and $r\ge0$, let
$\mathsf{T}^{*}_{r}[\mathfrak{L}]$ be the anchor functional of the summation lemma,
\begin{equation*}
\mathsf{T}^{*}_{r}[\mathfrak{L}]:=\sup_{\Delta}\,e^{-\rho(\Delta)}
\sum_{F\in\mathfrak{L}:\ \Delta\ \text{sees}\ \mathcal{R}_F}\|F\|\,e^{r d_i(A_F)}.
\end{equation*}
For $g'$ in the chart $\mathfrak{G}_{\varepsilon-\eta^T}(\Sigma_Y)$, where $\varepsilon$ is the
current radius at the law sites $\Sigma_Y$ and $\eta^T$ the debit of the fluctuation site, set
\begin{equation*}
f_v(\mathsf{s};g'):=v^{\sharp}\bigl(\mathsf{R}(\mathsf{s})\restriction Y^{\sharp};\Phi_Y,
\mathsf{u}(\mathsf{s})g'\bigr),
\qquad
v_{\mathsf{y}}:=\int_{[0,1]^{\mathsf{y}}}\prod_{\Delta\in\mathsf{y}}d\mathsf{s}(\Delta)\,
\partial_{\mathsf{s}(\Delta)}f_v\Big|_{\mathsf{s}=0\ \text{off}\ \mathsf{y}},
\end{equation*}
let $\mathfrak{Y}(v)$ be the set of families $\mathsf{y}\subset\pi^T$ every component of which holds
a source cell and a member seeing $Y^{\sharp}$, let $\operatorname{lab}(v,\mathsf{y})\in\mathbf{D}_i$
be the label the summation lemma assigns to $(v,\mathsf{y})$, and, for $Y'\in\mathbf{D}_i$,
\begin{equation*}
V^{\mathrm{kept}}(Y'):=\sum_{v}\ \sum_{\mathsf{y}\in\mathfrak{Y}(v):\
\operatorname{lab}(v,\mathsf{y})=Y'}v_{\mathsf{y}},
\end{equation*}
the outer sum running over $v\in\mathfrak{K}(Y)$, $(Y,m)$ live, $m\le i$. Let
$\kappa_V:=(1-2\delta)\kappa$, and consider
\begin{equation}\label{4.22:eq-fluctuation-site-a}
\begin{aligned}
&\varepsilon^{\mathrm{kept}}_i:=C_\Sigma^{\mathrm{sum}}\,\mathfrak{n}^{\mathrm{kept}}\,
e^{-\frac14\delta\kappa(7\check R_i-2)},\\
&\sup_{\mathsf{T}\ge\mathsf{T}_\gamma}C_\Sigma^{\mathrm{sum}}\Bigl\{C_{\mathsf{a}}(1)
\bigl[b_0+C_NC_0^{\sharp}(\mathsf{T}+1+c_0)^{r_{\mathrm{pr}}}\bigr]
+\mathfrak{n}^{\mathrm{kept}}\Bigr\}\\
&\qquad\times e^{-\frac14\delta\kappa(7\mathsf{T}^{r_{\mathrm{pr}}}-2)}\le\tfrac14\varepsilon_V^{*}.
\end{aligned}
\end{equation}
Here $C_\Sigma^{\mathrm{sum}}$ is the constant of part (c) of the summation lemma, and
$\mathsf{T}_\gamma$, $C_{\mathsf{a}}(1)$, $b_0$, $C_N$, $C_0^{\sharp}$, $c_0$, $r_{\mathrm{pr}}$ and
$\varepsilon_V^{*}$ are the constants of the same names entering the proposition on the fluctuation
site; the variable $\mathsf{T}$ of the supremum is unrelated to the functional
$\mathsf{T}^{*}_{r}$. In (e), $V^{\mathrm{pl}}$, $V^{\mathrm{rel}}$, $\mathsf{b}_i(0)$ and
$\varepsilon^{\mathrm{rel}}_i$ are the activities and quantities of that proposition, and $A$, $X$,
$\mathcal{P}$ denote an output's label, true domain and set of large-field regions.
Then the following hold.
\begin{enumerate}
\item[(a)] For every kept member $v$,
\begin{equation*}
f_v(\mathbb{1})-f_v(0)=\sum_{\mathsf{y}\in\mathfrak{Y}(v)}v_{\mathsf{y}},
\qquad\text{hence}\qquad
\sum_{Y'}V^{\mathrm{kept}}(Y')=\sum_v\bigl[f_v(\mathbb{1})-f_v(0)\bigr]
\end{equation*}
exactly; and $f_v(0)=v^{\sharp}(\mathbf{U}\restriction Y^{\sharp};\Phi_Y,g')$. Thus the kept member
survives the site: it is the member at the new base with riding map equal to $1$, and it stays under
$\mathbf{T}'_m(Y)$. This is the telescoping identity listed in
\eqref{4.22:eq-majorant-transfer-a}.
\item[(b)] For every $g'\in\mathfrak{G}_{\varepsilon-\eta^T}(\Sigma_Y)$,
$|v_{\mathsf{y}}|\le\mathfrak{O}(v)e^{-(\kappa_1-1)|\mathsf{y}|}$. Every
$\mathsf{y}\in\mathfrak{Y}(v)$ satisfies $|\mathsf{y}|\ge\rho(\Delta)\ge\frac52\check R_m$ for its
member $\Delta\subset Y^{\sharp}$. The piece $v_{\mathsf{y}}$ reads $\mathsf{B}$ on
$\bigcup\mathsf{y}$ only, $((\mathbf{U},\mathbf{J}),w)$ on $Y^{\sharp}\cup\bigcup\mathsf{y}$,
and $\Phi_Y$; it is an orbit datum of radius $\varepsilon-\eta^T$ at $\Sigma_Y$, with the invariance
of $v$. The reader condition holds for $v$, where $\operatorname{pa}(\Delta)$ denotes the set
attached to the member $\Delta$ in the conventions~\ref{4.22:not-live-region-geometry}:
\begin{equation}\label{4.22:eq-fluctuation-site-2}
\operatorname{dist}_i(\operatorname{pa}(\Delta),Q_Y(i))\le c_\rho\,\rho(\Delta)
\quad\text{for every }\Delta\text{ seeing }Y^{\sharp},\qquad c_\rho=120.
\end{equation}
These are the piece bound and the reader condition listed in \eqref{4.22:eq-majorant-transfer-a}.
\item[(c)] For every $r\ge0$,
\begin{equation*}
\mathsf{T}^{*}_{r}[\mathfrak{L}^{\mathfrak{K}}_T]\le C_\Sigma\mathfrak{s}_\Sigma(1)
\le\mathfrak{n}^{\mathrm{kept}}.
\end{equation*}
\item[(d)] For every $Y'\in\mathbf{D}_i$ and every bond $b$,
\begin{equation*}
|V^{\mathrm{kept}}(Y')|\le\varepsilon^{\mathrm{kept}}_ie^{-\kappa_Vd_i(Y')},
\qquad
\sum_{Y'\ni b}|V^{\mathrm{kept}}(Y')|\le K(d)\,\varepsilon^{\mathrm{kept}}_i,
\end{equation*}
with $\varepsilon^{\mathrm{kept}}_i$ as defined in \eqref{4.22:eq-fluctuation-site-a}.
\item[(e)] Parts (b)--(e) of the proposition on the fluctuation site hold for
$V:=V^{\mathrm{pl}}+V^{\mathrm{rel}}+V^{\mathrm{kept}}$, with $C_{\mathsf{a}}(1)\mathsf{b}_i(0)$
replaced by $C_{\mathsf{a}}(1)\mathsf{b}_i(0)+\mathfrak{n}^{\mathrm{kept}}$ in
$\varepsilon^{\mathrm{rel}}_i$, under the supremum condition in
\eqref{4.22:eq-fluctuation-site-a}. An output that reads some $V^{\mathrm{kept}}(Y')$ is in format,
with $A:=X$, with $(Y,m)$ in its set $\mathcal{P}$ of large-field regions, and in the class
$\mathfrak{S}$.
\end{enumerate}
\end{proposition}

\begin{proof}
\emph{(a).} By part (a) of Lemma~\ref{4.22:lem-kept-holomorphy}, $f_v$ is holomorphic on the
polydisc $\{|\mathsf{s}(\Delta)|\le e^{\kappa_1}\}$ in the interpolation parameters, times the
remaining parameter, source and chart domains of the site. It reads the objects that depend on
$\mathsf{s}$ only through the restriction
$(\mathcal{Y},\mathcal{J}[\mathcal{Y}],\mathsf{u})(\mathsf{s})\restriction Y^{\sharp}$, at amplitude
one, where $\mathcal{J}[\mathcal{Y}]$ is the object of the proposition on the fluctuation site
determined by $\mathcal{Y}$; for the riding map this holds because, by the property of the riding
map of that proposition, its value $\mathsf{u}(\mathsf{s})$ at a law site $s$ reads $\mathcal{Y}$
only on a set $B^{\mathrm{top}}(s)\subset Y$. Hence $f_v$ lies in the class of sourced readers of the
source-interpolation statement, with read region $\mathcal{R}=Y^{\sharp}$. Part (i) of that statement
gives $f_v(\mathbb{1})-f_v(0)=\sum_{\mathsf{y}\in\mathfrak{Y}(v)}v_{\mathsf{y}}$. Each
$\mathsf{y}\in\mathfrak{Y}(v)$ carries exactly one label $\operatorname{lab}(v,\mathsf{y})\in
\mathbf{D}_i$, so summing over $v$ and regrouping the pieces by their labels yields
$\sum_{Y'}V^{\mathrm{kept}}(Y')=\sum_v[f_v(\mathbb{1})-f_v(0)]$, with no term added or lost. By part
(ii) of the base-point lemma, applied with empty exceptional set $\mathsf{K}=\emptyset$, one has
$\mathcal{Y}(0)=0$ and $\mathsf{u}(0)=1$; so the first entry of $\mathsf{R}(0)$ is $\mathbf{U}$ and
$f_v(0)=v^{\sharp}(\mathbf{U}\restriction Y^{\sharp};\Phi_Y,g')$, which is the member at the new
base with riding map equal to $1$.

\emph{(b).} The number $c:=f_v(0)$ does not depend on $\mathsf{s}$, and
\begin{equation*}
|f_v(\mathsf{s})-c|\le\operatorname{osc}_{\mathfrak{P}}(v^{\sharp})\le\mathfrak{O}(v)
\end{equation*}
by the oscillation bound \eqref{4.22:eq-kept-oscillation-a} of
Lemma~\ref{4.22:lem-kept-oscillation}. Part (ii) of the source-interpolation statement, in its second
form, applied to $f_v$ with this constant $c$, gives
$|v_{\mathsf{y}}|\le\mathfrak{O}(v)e^{-(\kappa_1-1)|\mathsf{y}|}$ for every $g'$ in the chart
$\mathfrak{G}_{\varepsilon-\eta^T}(\Sigma_Y)$.

Since there are no hubs, a member seeing $Y^{\sharp}$ meets it; by the geometry statement (b)
of~\ref{4.22:not-live-region-geometry}, whose proof is outstanding, it then lies in $Y^{\sharp}$.
Let $\Delta\subset Y^{\sharp}$ and consider a chain of cells, consecutive ones sharing a wall, from
$\Delta$ to a source cell. The source cell lies in $\Omega_{i+1}\subset\Omega_{m+1}$. Since
$\operatorname{dist}_m(Y^{\sharp},\Omega_{m+1})\ge\frac52\check R_m$ and every cell at level-$m$
distance less than $\frac52\check R_m$ from $Y^{\sharp}$ has level at most $m$, hence sides of at
most one level-$m$ unit (the geometry statement (c) of~\ref{4.22:not-live-region-geometry}, whose
proof is outstanding), the chain leaves the set of points at level-$m$ distance less than
$\frac52\check R_m$ from $Y^{\sharp}$ only after passing through cells of that size.
Hence the chain has at least
$\frac52\check R_m$ members, and $\rho(\Delta)\ge\frac52\check R_m$. Since every
$\mathsf{y}\in\mathfrak{Y}(v)$ contains a connected subfamily holding a member $\Delta$ that sees,
hence lies in, $Y^{\sharp}$, together with a source cell, $|\mathsf{y}|\ge\rho(\Delta)$. This is the
argument by which the depth condition is established in the proof of part (c) of the proposition on
the fluctuation site.

That $v_{\mathsf{y}}$ reads $\mathsf{B}$ on $\bigcup\mathsf{y}$ only, the pair and the source on
$Y^{\sharp}\cup\bigcup\mathsf{y}$, and otherwise only $\Phi_Y$, is part (iii) of the
source-interpolation statement; that it has the invariance of $v$ is part (iv). That it is an orbit
datum of radius $\varepsilon-\eta^T$ at $\Sigma_Y$ is obtained exactly as in the proof of part (iii)
of the transfer lemma for orbit data, with $f_v$ in place of the datum there.

For \eqref{4.22:eq-fluctuation-site-2}, let $\Delta$ see $Y^{\sharp}$, so $\Delta\subset Y^{\sharp}$.
Let $s_i(Y)$ be the number attached to $Y$ at level $i$ in the
conventions~\ref{4.22:not-live-region-geometry}.
By the geometry statements (a) and (d) of~\ref{4.22:not-live-region-geometry}, whose proof is
outstanding, and since $\rho(\Delta)$ is at least the level-$i$ distance from $\Delta$ to the source
cells, which lie in $\Omega_{i+1}$,
\begin{equation*}
\operatorname{dist}_i(\operatorname{pa}(\Delta),Q_Y(i))\le s_i(Y)+1\le101\check R_i+1
\le c_\rho\cdot7\check R_i\le c_\rho\,\rho(\Delta);
\end{equation*}
the first inequality uses $\Delta\subset Y^{\sharp}$ and that $Q_Y(i)$ holds a vertex of $Y^{\sharp}$,
the second is the geometry statement (a), and the last uses
$\operatorname{dist}_i(Y^{\sharp},\Omega_{i+1})\ge7\check R_i$, the geometry statement (d).

\emph{(c).} Since there are no hubs, a cell $\Delta$ sees $Y^{\sharp}$ only if it meets it, hence
only if $\Delta\subset Y^{\sharp}$, and by the geometry statement (b)
of~\ref{4.22:not-live-region-geometry}, whose proof is outstanding, this happens for at most one live
region $(Y,m)$. Moreover $d_i(Q_Y(i))=0$ by Definition~\ref{4.22:def-kept-term-reader}, so the factor
$e^{rd_i(A_F)}$ equals one and $\mathsf{T}^{*}_{r}[\mathfrak{L}^{\mathfrak{K}}_T]$ does not depend on
$r$. Take as depth datum $\varrho:=\rho$ on cells. It satisfies the depth condition
$\check R_m\le2\varrho+2$ whenever the cell meets the strip $Y^{\sharp}$ of a live $(Y,m)$, because
(b) gives $\check R_m\le\frac25\rho(\Delta)$. Lemma~\ref{4.22:lem-probe-summand}, in its form for
cells (where the factor $2^d$ is replaced by $1$), therefore gives, for the unique $(Y,m)$ with
$\Delta\subset Y^{\sharp}$,
\begin{equation*}
e^{-\rho(\Delta)}\sum_{v\in\mathfrak{K}(Y)}\mathfrak{O}(v)
\le C_\Sigma\bigl(2\rho(\Delta)+2\bigr)^{p_\Sigma}e^{-\rho(\Delta)}.
\end{equation*}
Taking the supremum over $\Delta$ and using the definition of $\mathfrak{s}_\Sigma(\beta)$ at
$\beta=1$ bounds the right member by $C_\Sigma\mathfrak{s}_\Sigma(1)$, and the definition of
$\mathfrak{n}^{\mathrm{kept}}$ in \eqref{4.22:eq-fixed-number-entry-a} gives
$C_\Sigma\mathfrak{s}_\Sigma(1)\le\mathfrak{n}^{\mathrm{kept}}$.

\emph{(d).} Every label $\operatorname{lab}(v,\mathsf{y})$ holds a source cell, which contains a
point of $\Omega_{i+1}$, and the cube $Q_Y(i)$, which contains a point of $Y^{\sharp}$. Since
$\operatorname{dist}_i(Y^{\sharp},\Omega_{i+1})\ge7\check R_i$ by the geometry statement (d)
of~\ref{4.22:not-live-region-geometry}, whose proof is outstanding, and a tree graph meeting two
cubes has length at least their distance, every label satisfies
$d_i(\operatorname{lab}(v,\mathsf{y}))\ge7\check R_i-2$.

Apply part (d) of the summation lemma to the family $\mathfrak{L}^{\mathfrak{K}}_T$, with
$r=\kappa_V=(1-2\delta)\kappa$, $\varsigma=\frac14\delta\kappa$, $D=7\check R_i-2$, and a cube
$Q_*\subset Y'$. Its hypothesis $\kappa_1-1\ge c_l(r+\varsigma)+\lambda_*$ follows from the floor
\eqref{4.22:eq-fluctuation-site-1}, because $r+\varsigma=(1-\frac74\delta)\kappa\le2\kappa+4$. The
readers are $v$ with $\|v\|=\mathfrak{O}(v)$ and $A_v=Q_Y(i)$, and they satisfy the reader condition
\eqref{4.22:eq-fluctuation-site-2} by (b). Every pair $(v,\mathsf{y})$ with
$\operatorname{lab}(v,\mathsf{y})=Y'$ has $Q_*\in\operatorname{lab}(v,\mathsf{y})$. Hence, by the
piece bound of (b), by the summation lemma (d), and by (c) at $r+\varsigma+2$,
\begin{align*}
e^{\kappa_Vd_i(Y')}|V^{\mathrm{kept}}(Y')|
&\le\sum_v\sum_{\mathsf{y}:\ \operatorname{lab}(v,\mathsf{y})\ni Q_*}
\mathfrak{O}(v)\,e^{-(\kappa_1-1)|\mathsf{y}|}\,e^{\kappa_Vd_i(\operatorname{lab}(v,\mathsf{y}))}\\
&\le e^{-\frac14\delta\kappa(7\check R_i-2)}\,C_\Sigma^{\mathrm{sum}}\,
\mathsf{T}^{*}_{\kappa_V+\varsigma+2}[\mathfrak{L}^{\mathfrak{K}}_T]
\le\varepsilon^{\mathrm{kept}}_i.
\end{align*}
For the anchored total, the bound just proved gives
\begin{equation*}
\sum_{Y'\ni b}|V^{\mathrm{kept}}(Y')|\le\varepsilon^{\mathrm{kept}}_i\sum_{Y'\ni b}
e^{-\kappa_Vd_i(Y')}\le K(d)\,\varepsilon^{\mathrm{kept}}_i,
\end{equation*}
where the last sum is bounded, with the constant $K(d)$ of \cite[(1.26)]{B-CMP116}, by that
estimate, used at $\kappa_V\ge c_d+1$.

\emph{(e).} The supremum of a sum is at most the sum of the suprema, so, with $\mathfrak{L}_T$ the
family of readers of the proposition on the fluctuation site,
\begin{equation*}
\mathsf{T}^{*}_{r}[\mathfrak{L}_T\cup\mathfrak{L}^{\mathfrak{K}}_T]
\le\mathsf{T}^{*}_{r}[\mathfrak{L}_T]+\mathsf{T}^{*}_{r}[\mathfrak{L}^{\mathfrak{K}}_T],
\end{equation*}
and the second member is at most $\mathfrak{n}^{\mathrm{kept}}$ by (c); this accounts for the
replacement of $C_{\mathsf{a}}(1)\mathsf{b}_i(0)$ by
$C_{\mathsf{a}}(1)\mathsf{b}_i(0)+\mathfrak{n}^{\mathrm{kept}}$ in $\varepsilon^{\mathrm{rel}}_i$,
and the supremum condition in \eqref{4.22:eq-fluctuation-site-a} is assumed. The rest of
the proof of the proposition on the fluctuation site reads the activity $V$ only through a bound
$|V(Y')|\le\nu(Y')$, through its holomorphy, and through the property that $V(Y')$ reads
$\mathsf{B}$ on $Y'$; the stage of that proof in which the integration variables are treated does
not involve $V$ at all. The data $\Phi_Y$ enter $V^{\mathrm{kept}}$ as a parameter under bounds
uniform in it. Finally, an output that reads some $V^{\mathrm{kept}}(Y')$ is in format, with $A:=X$,
with $(Y,m)\in\mathcal{P}$, and in $\mathfrak{S}$, by the third clause of the format convention for
wrapped terms.
\end{proof}

\begin{proposition}[Kept terms at a preparatory stage]\label{4.22:prop-preparatory-stage}
Consider the site $(\mathrm{P})$ of the operation $\mathbf{R}_k$ at stage $n+1$, at level $j=h+n$.
The chain may be the chain of the component under treatment or a chain of another component,
and the preparation may be running or interrupted. Let $d^{\wedge}$ be the weighted contour
distance of the chain term $\mathbb{B}^{(n)}_k$, and let $w_{\wedge}>0$ be the amount that
$d^{\wedge}$ charges a contour per crossed zone.

We use the notation of the proposition on the large-field stages of the step: the collar $R_n$;
the cover $\{\square\}$ by blocks, with fattenings $\widetilde{\square}$; the factor
\begin{equation}\label{4.22:eq-preparatory-stage-1}
\vartheta(\square)=\vartheta\, e^{-\frac12\delta_P d^{\wedge}(\square,R_n)};
\end{equation}
the constants $K_A$, $K_2^{\flat}$, $V'$, $c_{\mathsf{p}}$; the relative count
$\mathsf{p}^{\mathrm{rel}}$; the quantities $x_{\square}$, $y_{\square}$, $a_k$ and the sets
$\operatorname{Cor}(\square)$; the families $\hat{\mathsf{y}}$; the functional $\mathfrak{S}$
and the constant $C^P_{\mathsf{a}}$ of its clause (d); and its class $\mathfrak{Old}$ of old
blocks. This class contains every block $\square$ whose fattening $\widetilde{\square}$ meets
the strip of a live large-field region. For use in the statement that follows this proposition
we also take from it its two totals $\mu$ and $\mathsf{M}_{*}$, its block quantities
$\hat{\mathsf{K}}_{\square}$, its weights, and the anchoring carried out in the last steps of
its proof.

Assume the following.
\begin{itemize}
\item The proposition on the large-field stages of the step, including its clauses (b), (c), (d)
and their proofs, and the verification made there that the depth datum of the stage satisfies
condition (D2) on the depth datum (stated in the geometry
statement~\ref{4.22:not-live-region-geometry} and assumed in
Lemma~\ref{4.22:lem-probe-summand}).
\item Clauses (a), (c), (h), (i1) and the payment rule $(g^{\wedge})$ of the geometry of live
large-field regions~\ref{4.22:not-live-region-geometry}, whose proof is outstanding.
\item The lower bound on the scales $\check{R}_i$ imposed among the conditions on the
parameters, which gives
\begin{equation}\label{4.22:eq-preparatory-stage-2}
101\check{R}_{h+1}+2\le 110\bigl(\tfrac52\check{R}_{h+1}-1\bigr),
\qquad
101\check{R}_k+2\le 110\bigl(\check{R}_k-1\bigr).
\end{equation}
\item A condition in the order of choice of the parameters under which, for every block
$\square$ of the cover,
\begin{equation*}
x_{\square}+1\le \frac{110\,d^{\wedge}(\square,R_n)}{w_{\wedge}}
\ \Longrightarrow\
(a_k+1)|\operatorname{Cor}(\square)|\le\tfrac14\delta_Pd^{\wedge}(\square,R_n).
\end{equation*}
\item The condition on the zone charge, which includes $\frac1{16}\delta_P w_{\wedge}\ge 1$.
\item Clause (ii) of the interpolation lemma of the stage (its Cauchy bounds).
\item The first clause of the locality statement for the background map of the stage.
\item Clause (a) of Lemma~\ref{4.22:lem-kept-holomorphy}.
\item Clause $(\mathrm{ii})_2$ of Lemma~\ref{4.22:lem-kept-oscillation}, in the notation of
\eqref{4.22:eq-kept-oscillation-a}.
\item Lemma~\ref{4.22:lem-probe-summand}.
\item Clauses (i)--(iii) and (v) of the representation rules of the stage.
\end{itemize}

Let $(Y,m)$ be a live large-field region created at step $m<k$, with strip $Y^{\sharp}$, and
let $v\in\mathfrak{K}(Y)$, with orbit datum $v^{\sharp}$. Put $m_v:=h+1$ if the component of
$\Lambda_k^c$ containing $Y$ is treated at $\mathbf{R}_k$, and $m_v:=k$ if it is not. Let the
core of $v$ be $S_v:=Q_Y(m_v)$, where $Q_Y(\cdot)$ is the core assignment of the proposition on
the large-field stages; here the core is a single cube. Put
$\mathcal{A}_v:=\{\square:\widetilde{\square}\cap Y^{\sharp}\neq\emptyset\}$. Then
$\mathcal{A}_v\subset\mathfrak{Old}$, by the first clause of the definition of $\mathfrak{Old}$.

In the notation of the proposition on the large-field stages, the two sets which that
proposition attaches to a term, $\mathcal{Q}^{w}_v$ and $\mathcal{S}_j(v)$, are empty for $v$.
The weight of $v$ is $\mathfrak{w}(v):=\mathfrak{O}(v)$, since a core consisting of one cube has
linear size $0$. The two numbers which that proposition attaches to a term, $\mathsf{e}(v)$ and
$\mathsf{f}(v)$, vanish, and $\mathsf{p}^{\mathrm{rel}}(v)\le 2c_{\mathsf{p}}$. With
$t=(t_{\square})$ the interpolation
parameters and $\mathsf{s}$ the auxiliary parameter of the stage, put
\begin{equation}\label{4.22:eq-preparatory-stage-3}
F_v(t;\mathsf{s}):=v^{\sharp}\bigl(\Phi_{t,\mathsf{s}}(b)\restriction Y^{\sharp};\,
\Phi_Y,\,\mathfrak{u}^{(l)}_{t,\mathsf{s}}(b)\,g'\bigr).
\end{equation}
Then the following hold.
\begin{enumerate}
\item[(a)] \emph{Elements.} With $\hat e=(v,A,\emptyset,\hat{\mathsf{y}})$,
\begin{equation}\label{4.22:eq-preparatory-stage-4}
\begin{gathered}
F_v(\mathbb{1})-F_v(0)=\sum_{\emptyset\neq A\subset\mathcal{A}_v}\ \sum_{\hat{\mathsf{y}}}
v_{\hat e},\\
|v_{\hat e}|\le \nu(\hat e):=\mathfrak{O}(v)\prod_{\square\in A}K_A\vartheta(\square)\cdot
e^{-(\kappa_1-1)|\hat{\mathsf{y}}|}.
\end{gathered}
\end{equation}
The element $v_{\hat e}$ depends on the pair $(\mathbb{U},\mathbb{J})$ only on
$\mathsf{d}^{\mathrm{rel}}(\hat e)\cup Y^{\sharp}$. It depends on the integration variables
only on $\mathsf{b}(\hat e)$, and it depends on $\Phi_Y$. Moreover
$F_v(0)=v^{\sharp}(\mathbb{U}\restriction Y^{\sharp};\Phi_Y,g')$: the kept term survives at
the base of the stage.
\item[(b)] For $\square\in\mathcal{A}_v$,
\begin{equation}\label{4.22:eq-preparatory-stage-5}
x_{\square}+1\le \frac{110\,d^{\wedge}(\square,R_n)}{w_{\wedge}} .
\end{equation}
Hence, by that parameter condition,
$(a_k+1)|\operatorname{Cor}(\square)|\le\frac14\delta_Pd^{\wedge}(\square,R_n)$. Also
\begin{equation}\label{4.22:eq-preparatory-stage-6}
y_{\square}\le K_2^{\flat}\vartheta\, e^{-\frac14\delta_Pd^{\wedge}(\square,R_n)} .
\end{equation}
\item[(c)] The sum over the blocks of $\mathcal{A}_v$ satisfies
\begin{equation}\label{4.22:eq-preparatory-stage-7}
\sum_{\square\in\mathcal{A}_v}e^{-\frac12\delta_Pd^{\wedge}(\square,R_n)}
e^{(a_k+1)|\operatorname{Cor}(\square)|}\le V'\mathsf{p}^{\mathrm{rel}}(v).
\end{equation}
\item[(d)] \emph{The collar, in $d^{\wedge}$.} Let $\square_0\in\mathfrak{Old}$ be such that
$\widetilde{\square}_0$ meets the strip of $(Y,m)$. Whether $R_n$ belongs to the component of
$Y$ or to another component,
\begin{equation}\label{4.22:eq-preparatory-stage-8}
d^{\wedge}(\square_0,R_n)\ge w_{\wedge}\bigl(\tfrac52\check{R}_m-5\bigr).
\end{equation}
Moreover
\begin{equation}\label{4.22:eq-preparatory-stage-9}
\mathfrak{S}^{\mathfrak{K}}(\square_0):=\sum\{\mathfrak{O}(v):\square_0\in\mathcal{A}_v\}
\le e\,\mathfrak{n}^{\mathrm{kept}}\,e^{\frac1{16}\delta_Pd^{\wedge}(\square_0,R_n)} .
\end{equation}
Consequently clause (d) of the proposition on the large-field stages holds for
$\mathfrak{S}+\mathfrak{S}^{\mathfrak{K}}$, with $C^P_{\mathsf{a}}$ replaced by
$C^{P,+}_{\mathsf{a}}:=C^P_{\mathsf{a}}+e\,\mathfrak{n}^{\mathrm{kept}}$.
\end{enumerate}

The stage is carried out in one of two settings: on the chain domains $\mathfrak{D}^{\flat}_n$,
or in the lettering of the inheritance theorem itself. In (a) and throughout,
$\vartheta(\square)$ denotes the factor of the setting in force. In the first setting, the
$t_{\square}$-radius is $R^{\natural}_{\square}$. It is determined by
\begin{equation}\label{4.22:eq-preparatory-stage-10}
\frac{2}{R^{\natural}_{\square}}=K_A\vartheta^{\natural}(\square),\qquad
\vartheta^{\natural}(\square)=\vartheta^{\natural}e^{-\frac12\delta_Pd^{-}_{\square}},\qquad
\vartheta^{\natural}:=\frac{\vartheta}{\mathsf{w}_0},
\end{equation}
where $d^{-}_{\square}$ is the distance attached to the block $\square$ in the inheritance
theorem and $\mathsf{w}_0$ is the constant of that theorem fixed below.
This is the radius of the inheritance theorem at $\theta=\frac12$, read through the dictionary
of its companion lemma, which replaces $\vartheta(\square)$ by $\vartheta^{\natural}(\square)$
and $\vartheta$ by $\vartheta^{\natural}$. For $\theta_S\in[\frac12,\frac89]$ one has
\begin{equation*}
\frac1{\mathsf{w}_0}=\frac{12}{\min\{\frac12,\theta_S\}}=24 .
\end{equation*}
In the second setting, the $t_{\square}$-radius is $r^{\sharp}_{\square}$, with
$2/r^{\sharp}_{\square}=K_A\vartheta(\square)$, as in the inheritance theorem. The exponent
$\frac12\delta_Pd^{\wedge}$ is the same in both settings. Hence clauses (b)--(d), the totals for
the kept terms stated after this proposition and the corollary on the kept terms are the same in
both settings. The factor $24$ enters only through $K_2^{\flat}\vartheta^{\natural}$, in
smallness conditions of supremum form imposed at the value of $\vartheta$ of the setting in
force, and it creates no cap in the sense of Definition~\ref{2.3:def-cap}.
\end{proposition}

\begin{proof}
\emph{(a).} The interpolation of the stage is used unchanged.

\emph{Dependence on $t$.} By the first clause of the locality statement for the background
map, the restrictions $\Phi_{t,\mathsf{s}}(b)\restriction Y^{\sharp}$ and
$\mathfrak{u}^{(l)}_{t,\mathsf{s}}(b)\restriction\Sigma_Y$ depend on $t$ only through the
$t_{\square}$ with $\square\in\mathcal{A}_v$. Hence, by \eqref{4.22:eq-preparatory-stage-3},
$F_v$ depends on $t$ only through $(t_{\square})_{\square\in\mathcal{A}_v}$.

\emph{Decomposition.} Apply M\"obius inversion on the subsets of $\mathcal{A}_v$ to the set
function $A'\mapsto F_v(\mathbb{1}_{A'})$. This gives
$F_v(\mathbb{1}_{A'};\mathsf{s})=\sum_{A\subset A'}v_A(\mathsf{s})$, where
\begin{equation}\label{4.22:eq-preparatory-stage-11}
v_A(\mathsf{s})=\sum_{B\subset A}(-1)^{|A\setminus B|}F_v(\mathbb{1}_B;\mathsf{s})
=\int_{[0,1]^A}\partial_{t_A}F_v(t\mathbb{1}_A;\mathsf{s})\,dt .
\end{equation}
The second equality is the fundamental theorem of calculus applied once in each variable
$t_{\square}$, $\square\in A$. Since $v_{\emptyset}=F_v(0)$, taking $A'=\mathcal{A}_v$ gives
\begin{equation*}
F_v(\mathbb{1})-F_v(0)=\sum_{\emptyset\neq A\subset\mathcal{A}_v}v_A .
\end{equation*}

\emph{Holomorphy.} By clause (a) of Lemma~\ref{4.22:lem-kept-holomorphy}, $F_v$ is holomorphic
on the product polydisc
\begin{equation*}
\{|t_{\square}|\le R^{\natural}_{\square}\}\times\{|\mathsf{s}|\le e^{\kappa_1}\}
\end{equation*}
at the site $(\mathrm{P})$. In the second setting, $R^{\natural}_{\square}$ is replaced by
$r^{\sharp}_{\square}$. In each setting, this clause of Lemma~\ref{4.22:lem-kept-holomorphy}
rests on the inputs listed in that lemma for that setting.

\emph{The constant term.} Put $c_0:=F_v(0;\cdot)$. At $t=0$ the layer is switched off and
$\mathfrak{u}=1$, so $c_0$ depends neither on $t$ nor on $\mathsf{s}$. It equals
$v^{\sharp}(\mathbb{U}\restriction Y^{\sharp};\Phi_Y,g')$, which is the last assertion of (a).
For $A\neq\emptyset$ the derivative $\partial_{t_A}$ annihilates $c_0$, so
$\partial_{t_A}F_v=\partial_{t_A}(F_v-c_0)$. By clause $(\mathrm{ii})_2$ of
Lemma~\ref{4.22:lem-kept-oscillation} (see \eqref{4.22:eq-kept-oscillation-a}),
$\sup|F_v-c_0|\le\mathfrak{O}(v)$ on the polydisc.

\emph{Cauchy bound.} The Cauchy bounds of clause (ii) of the interpolation lemma of the stage,
applied to $F_v-c_0$ on this polydisc, use $2/R^{\natural}_{\square}=K_A\vartheta(\square)$ in
the notation of the setting in force. They give, uniformly for $\mathsf{s}$ in the parameter
polydisc $\mathfrak{P}_A$,
\begin{equation}\label{4.22:eq-preparatory-stage-12}
|v_A(\mathsf{s})|\le\mathfrak{O}(v)\prod_{\square\in A}K_A\vartheta(\square).
\end{equation}

\emph{Representation.} Thus $v_A$ belongs to the first member of the class of small
functionals of the stage, with the number $\mathfrak{n}$ of that class equal to the right
member of \eqref{4.22:eq-preparatory-stage-12}. Clauses (i)--(iii) and (v) of the
representation rules of the stage then decompose each $v_A$ into elements $v_{\hat e}$,
$\hat e=(v,A,\emptyset,\hat{\mathsf{y}})$. This yields the sum
\eqref{4.22:eq-preparatory-stage-4}, with the bound $\nu(\hat e)$, and with the stated
dependence of $v_{\hat e}$ on $(\mathbb{U},\mathbb{J})$, on the integration variables and on
$\Phi_Y$. By clause (h) of the geometry statement~\ref{4.22:not-live-region-geometry}, whose
proof is outstanding, no integration bond of the site $(\mathrm{P})$ lies on a strip. Hence
$v$ itself has no slot for the integration variables.

\emph{(b).} The proof of clause (b) of the proposition on the large-field stages uses only two
facts: that $\widetilde{\square}$ is a probe meeting $Y^{\sharp}$, and that the core $S_v$ has a
point within distance $s_{m_v}(Y)+1$ of it. Both the bound on $x_{\square}$ and the lower bound
on $d^{\wedge}(\square,R_n)$ below are obtained from that proof. The lower bound on $d^{\wedge}$
comes from the collar distances of clause (h) of the geometry
statement~\ref{4.22:not-live-region-geometry}, whose proof is outstanding.

If the component of $Y$ is treated, then $m_v=h+1$. By clause (a) of the geometry statement,
whose proof is outstanding, $s_{h+1}(Y)\le101\check{R}_{h+1}$. Hence
\begin{equation*}
x_{\square}+1\le s_{h+1}(Y)+2\le 101\check{R}_{h+1}+2,\qquad
d^{\wedge}(\square,R_n)\ge w_{\wedge}\bigl(\tfrac52\check{R}_{h+1}-1\bigr).
\end{equation*}
If it is not treated, then $m_v=k$, and
\begin{equation*}
x_{\square}+1\le 101\check{R}_k+2,\qquad
d^{\wedge}(\square,R_n)\ge w_{\wedge}\bigl(\check{R}_k-1\bigr).
\end{equation*}
In both cases \eqref{4.22:eq-preparatory-stage-2} bounds the quotient of the right members by
$110$. This is \eqref{4.22:eq-preparatory-stage-5}.

That parameter condition converts \eqref{4.22:eq-preparatory-stage-5} into
$(a_k+1)|\operatorname{Cor}(\square)|\le\frac14\delta_Pd^{\wedge}(\square,R_n)$. From this
point the proof of clause (b) of the proposition on the large-field stages gives
\eqref{4.22:eq-preparatory-stage-6} unchanged.

\emph{(c).} The proof of clause (c) of the proposition on the large-field stages sees the strips
only through the set of blocks to which it is applied, here $\mathcal{A}_v$. It therefore
applies to
$\mathcal{A}_v$. If the core is confined, in the sense of that proposition, the proof gives the
bound $2V'$. If it is not confined, then $[S_v]_k$ is one cube and its tree has length $0$. The
sum is then bounded by at most two boxes $\mathsf{B}_a$ of that proof, each contributing at most
$c_{\mathsf{p}}V'$. This gives \eqref{4.22:eq-preparatory-stage-7}.

\emph{(d), the collar.} We follow the verification of (D2) in the proposition on the large-field
stages.

First, $R_n$ lies outside the $\frac52\check{R}_m$-neighbourhood of $Y^{\sharp}$. Three cases
arise.
\begin{itemize}
\item $R_n$ belongs to the component of $Y$, and $m<h$. The cells of $R_n$ have level at least
$h>m$. By clause (c) of the geometry statement~\ref{4.22:not-live-region-geometry}, whose proof
is outstanding, the cells of the neighbourhood have level at most $m$.
\item $R_n$ belongs to the component of $Y$, and $m=h$. By clause (h) of the same statement,
whose proof is outstanding, the distance from $Y^{\sharp}$ to $R_n$ is at least
$\frac52\check{R}_{h+1}$ level-$(h+1)$ units. This is at least $\frac52\check{R}_h$ level-$h$
units.
\item $R_n$ belongs to another component. By clause (c) of that statement, whose proof is
outstanding, the neighbourhood lies in the
component of $\Lambda_k^c$ containing $Y$, which is not treated. On the other hand
$R_n\subset Z$, the region treated at $\mathbf{R}_k$. Hence $R_n$ does not meet the
neighbourhood.
\end{itemize}

Second, the probe $\widetilde{\square}_0$ meets $Y^{\sharp}$. By clause (i1) of that statement,
whose proof is outstanding, it has level at
most $m$, so every point of $\square_0$ lies within $5$ level-$m$ units of $Y^{\sharp}$.

Third, all cells of the neighbourhood have level at most $m$. A contour from $\square_0$ to
$R_n$ therefore covers a Euclidean length of at least $\frac52\check{R}_m-5$ level-$m$ units
inside a set all of whose cells have level at most $m$. By the payment rule $(g^{\wedge})$ of
that statement, whose proof is outstanding, it
pays at least $w_{\wedge}$ per such unit. This is \eqref{4.22:eq-preparatory-stage-8}.

\emph{(d), the depth condition.} Put $\varrho(\widetilde{\square}_0):=1+d^{\wedge}(\square_0,R_n)/w_{\wedge}$.
Then $\varrho(\widetilde{\square}_0)\ge1$. By \eqref{4.22:eq-preparatory-stage-8},
$\varrho(\widetilde{\square}_0)\ge\frac52\check{R}_m-4$. Hence
\begin{equation*}
\check{R}_m\le\tfrac25\bigl(\varrho(\widetilde{\square}_0)+4\bigr)
\le2\varrho(\widetilde{\square}_0)+2 .
\end{equation*}
The argument applies to every live region whose strip $\widetilde{\square}_0$ meets. So this
datum satisfies condition (D2) on the depth datum (stated in the geometry
statement~\ref{4.22:not-live-region-geometry} and assumed in
Lemma~\ref{4.22:lem-probe-summand}).

\emph{(d), the sum over kept terms.} By definition of $\mathcal{A}_v$, the condition
$\square_0\in\mathcal{A}_v$ means that the probe $\widetilde{\square}_0$ meets $Y^{\sharp}$.
Hence $\mathfrak{S}^{\mathfrak{K}}(\square_0)$ is the sum $\mathsf{S}^{\mathfrak{K}}$ of
Lemma~\ref{4.22:lem-probe-summand} at the probe $\widetilde{\square}_0$. That lemma at
$\beta=1$, with $\mathfrak{n}^{\mathrm{kept}}=\mathfrak{n}^{\mathrm{kept}}(1)$, gives
\begin{equation*}
\mathfrak{S}^{\mathfrak{K}}(\square_0)\le\mathfrak{n}^{\mathrm{kept}}
e^{\varrho(\widetilde{\square}_0)}
=e\,\mathfrak{n}^{\mathrm{kept}}e^{d^{\wedge}(\square_0,R_n)/w_{\wedge}}
\le e\,\mathfrak{n}^{\mathrm{kept}}e^{\frac1{16}\delta_Pd^{\wedge}(\square_0,R_n)} .
\end{equation*}
The last inequality holds because $\frac1{16}\delta_Pw_{\wedge}\ge1$, so
$1/w_{\wedge}\le\frac1{16}\delta_P$, and $d^{\wedge}\ge0$. This is
\eqref{4.22:eq-preparatory-stage-9}. The last assertion of (d) is this bound added to clause (d)
of the proposition on the large-field stages for $\mathfrak{S}$, the two constants being added.
\end{proof}

\begin{proposition}[Contribution of the kept terms to the stage totals.
Stated; proof incomplete at the step marked $(\ast)$]\label{4.22:prop-stage-totals}
Work in the setting of Proposition~\ref{4.22:prop-preparatory-stage}: the site $(\mathrm{P})$ of
$\mathbf{R}_k$, stage $n+1$, level $j=h+n$, with clauses (a)--(d) of that proposition in force.
On the run on the chain domains $\mathfrak{D}^{\flat}_n$ the factor $\vartheta$ is read as
$\vartheta^{\natural}=\vartheta/\mathsf{w}_0$ throughout, as in clause~(a) of
Proposition~\ref{4.22:prop-preparatory-stage}; on the run with cube radii $r^{\sharp}_{\square}$,
$2/r^{\sharp}_{\square}=K_A\vartheta(\square)$, it is $\vartheta$ itself.

For a kept member $v\in\mathfrak{K}(Y)$, with $(Y,m)$ live and $m<k$, let
$\hat e=(v,A,\emptyset,\hat{\mathsf{y}})$ be its groups, as in clause~(a) of
Proposition~\ref{4.22:prop-preparatory-stage}.

Let $\mu$ and $\mathsf{M}_{*}$ be the two totals of the flat cluster-expansion proposition at this
stage, taken over the family of its elements enlarged by the groups $\hat e$ of the kept members.
They are formed with the enlarged element bound
\begin{equation*}
\tilde\nu(e):=2\nu(e)\,e^{\varsigma(e)+c_J+c_{\flat}|\mathsf{b}(e)|},
\end{equation*}
where $\varsigma(e)$ is the exponent attached to the element $e$, $c_J$ and $c_{\flat}$ are the two
constants of the flat cluster-expansion proposition, and $\mathsf{b}(e)$ is the set of integration
bonds read by $e$; this exponent is non-negative. The total $\mu$ bounds, for every $q$ in the set
$\mathcal{I}^{\mathrm{fl}}$ of integration bonds of the stage, the sum of $\tilde\nu(e)$ over the
elements $e$ with $q\in\mathsf{b}(e)$. The total $\mathsf{M}_{*}$ bounds, for every point $x$, the sum
of $\tilde\nu(e)$ over the elements $e$ with $x\in[\mathsf{d}^{\mathrm{rel}}(e)]_k$, where
$[\mathsf{d}^{\mathrm{rel}}(e)]_k$ is the relative support of $e$ read at level $k$.

Let $\mu^{\flat}$ and $\mathsf{M}^{\flat}$ be the totals of the stage bound at the site $(\mathrm{P})$
for the terms of the stage other than the kept members. Let $C^{P,\star}_{\mathsf{a}}$ be the constant
of the totals clause of the stage bound that stands in that clause's two old-block members, namely in
the member of $\mu^{\flat}$ that carries $c_{\mathrm{loc}}$ and in the member of $\mathsf{M}^{\flat}$
that carries $c_{\star\star}''V'$. Put
$\mu^{\mathrm{leaf}}$ and $\mathsf{M}^{\mathrm{leaf}}$ as in the first two lines of the display below,
where $\mathfrak{n}^{\mathrm{kept}}$ is the fixed number of \eqref{4.22:eq-fixed-number-entry-a}.
Assume the six smallness conditions of the stage bound at the site $(\mathrm{P})$, imposed on each of
the two runs at its own factor ($\vartheta$ on the run with radii $r^{\sharp}_{\square}$,
$\vartheta^{\natural}$ on the chain-domain run), in supremum form, with $\mu^{\flat}$ and
$\mathsf{M}^{\flat}$ replaced by $\mu^{\flat}+\mu^{\mathrm{leaf}}$ and
$\mathsf{M}^{\flat}+\mathsf{M}^{\mathrm{leaf}}$ (equivalently, with $C^{P,\star}_{\mathsf{a}}$
increased by $e\,\mathfrak{n}^{\mathrm{kept}}$) in the first and second conditions and in the unmarked
factor of the sixth condition; if that factor is already bounded by $\tfrac18$ through the second
condition, no replacement is needed there. The first two conditions then read as in the last two
lines of the display, in which $\bar V'$ and $\tilde{\mathfrak{l}}$ are the constants of the second
condition:
\begin{equation}\label{4.22:eq-stage-totals-a}
\begin{aligned}
\mu^{\mathrm{leaf}}&:=4e^{c_J+3/2}K_2^{\flat}\vartheta\,c_{\mathrm{loc}}\,\mathfrak{n}^{\mathrm{kept}},\\
\mathsf{M}^{\mathrm{leaf}}&:=2e^{c_J+3/2}K_2^{\flat}\vartheta\,c_{\star\star}''V'\,
\mathfrak{n}^{\mathrm{kept}},\\
\mu^{\flat}+\mu^{\mathrm{leaf}}&\le\tfrac14,\\
4\bigl(\mathsf{M}^{\flat}+\mathsf{M}^{\mathrm{leaf}}
+2^dc_{\mathsf{p}}\bar V'\tilde{\mathfrak{l}}\bigr)&\le\tfrac12 .
\end{aligned}
\end{equation}
Then the following hold.
\begin{enumerate}
\item[(1)] The groups of the kept members add to the totals of the flat cluster-expansion proposition at
most $\mu^{\mathrm{leaf}}$ to $\mu$ and at most $\mathsf{M}^{\mathrm{leaf}}$ to $\mathsf{M}_{*}$. The
quantity $\mu^{\mathrm{leaf}}$ contains no factor $V'$.

Equivalently, the totals clause of the stage bound at the site $(\mathrm{P})$ holds for the family of
the stage's terms enlarged by the kept members, with
$C^{P,\star}_{\mathsf{a}}$ replaced by $C^{P,\star}_{\mathsf{a}}+e\,\mathfrak{n}^{\mathrm{kept}}$ in its
two old-block members. That is,
$\mu\le\mu^{\flat}+\mu^{\mathrm{leaf}}$ and $\mathsf{M}_{*}\le\mathsf{M}^{\flat}+\mathsf{M}^{\mathrm{leaf}}$.
\item[(2)] No group of a kept member is anchored at a window block or at a rim cube. The window data and
the rim data of the stage enter none of the bounds in~(1).
\item[(3)] The third condition is used in the proof below, through the bound $\Sigma_v\le 1/(3N)$ on the
exponent $\Sigma_v$ of $v$ in the element estimate of the stage bound. On the chain-domain
run it is imposed at $\vartheta^{\natural}$, with the constants $K_2^{\flat}$ and $c_{\mathsf{p}}V'$.

The products $K_2^{\flat}\vartheta$ and $K_2^{\flat}\vartheta\bar V'$ already occur in $\mu^{\flat}$
and $\mathsf{M}^{\flat}$ with numerical coefficients. The replacements therefore introduce no new kind
of term and move no degree. In the degree bookkeeping used for caps
(Definition~\ref{2.3:def-cap}), the degree of $\mu^{\mathrm{leaf}}$ is at most $-(7r+1)$ and that of
$\mathsf{M}^{\mathrm{leaf}}$ is at most $4r-(7r+1)+\epsilon$ for every $\epsilon>0$; no cap arises.
\item[(4)] Under the hypotheses above, the flat cluster-expansion proposition applies to its family of
elements enlarged by the groups $\hat e$ of the kept members. Its premise $\nu(e)\le\tfrac12$ holds for
these groups in the form
\begin{equation*}
\nu(\hat e)\le\tfrac12\tilde\nu(\hat e)\le\tfrac12\mu\le\tfrac18 .
\end{equation*}
This bound is supplied by the first smallness condition, which controls the supremum over the collar
multiplied by $\vartheta$, and never by any smallness of $\mathfrak{O}(v)$.

The outputs of the expansion that contain such a group are in the format of the stage. The set
$\mathcal{P}$ attached to such an output contains $(Y,m)$, and the output belongs to the class
$\mathfrak{S}$.
\end{enumerate}
\end{proposition}

\begin{proof}
Throughout, $v$ is a kept member, $v\in\mathfrak{K}(Y)$ with $(Y,m)$ live and $m<k$. On the
chain-domain run every $\vartheta$ below is $\vartheta^{\natural}$.

\emph{Format data of the kept members.} By the definitions fixed in
Proposition~\ref{4.22:prop-preparatory-stage}:
\begin{itemize}
\item $\mathcal{A}_v\subset\mathfrak{Old}$, where $\mathfrak{Old}$ is the class of old blocks of
Proposition~\ref{4.22:prop-preparatory-stage};
\item $\mathcal{Q}_v=\mathcal{S}_j(v)=\emptyset$;
\item $\mathfrak{w}(v)=\mathfrak{O}(v)$ and $\mathsf{e}(v)=\mathsf{f}(v)=0$.
\end{itemize}
Hence the weight of the stage bound at the site $(\mathrm{P})$ is
\begin{equation*}
W(v):=\mathfrak{w}(v)e^{\mathsf{e}(v)+\mathsf{f}(v)}=\mathfrak{O}(v).
\end{equation*}

Every group has $A\neq\emptyset$, because the sum in clause~(a) of
Proposition~\ref{4.22:prop-preparatory-stage} runs over the non-empty subsets $A$ of $\mathcal{A}_v$.

By the third smallness condition, and because $\mathsf{p}^{\mathrm{rel}}(v)\le 2c_{\mathsf{p}}$,
\begin{equation*}
\Sigma_v\le\frac{\mathsf{p}^{\mathrm{rel}}(v)}{6Nc_{\mathsf{p}}}\le\frac{1}{3N}\le\frac12,
\qquad\text{so}\qquad e^{\Sigma_v}\le e^{1/2}.
\end{equation*}

Let $\square_0$ be a block. If $\square_0$ lies in no $\mathcal{A}_v$, then
$\mathfrak{S}^{\mathfrak{K}}(\square_0)=0$. Otherwise $\square_0\in\mathcal{A}_v$ for some kept member
$v$, so $\widetilde{\square}_0$ meets the strip $Y^{\sharp}$ and $\square_0\in\mathfrak{Old}$. Then
clauses (b) and (d) of Proposition~\ref{4.22:prop-preparatory-stage} give
\begin{equation*}
y_{\square_0}\le K_2^{\flat}\vartheta\,e^{-\frac14\delta_Pd^{\wedge}(\square_0,R_n)},
\end{equation*}
and
\begin{equation*}
\mathfrak{S}^{\mathfrak{K}}(\square_0)=\sum_{v:\ \square_0\in\mathcal{A}_v}\mathfrak{O}(v)
\le e\,\mathfrak{n}^{\mathrm{kept}}\,e^{\frac1{16}\delta_Pd^{\wedge}(\square_0,R_n)}.
\end{equation*}
Multiplying the two bounds and using $\tfrac14-\tfrac1{16}=\tfrac3{16}$, we obtain, for every block
$\square_0$,
\begin{equation}\label{4.22:eq-stage-totals-b}
y_{\square_0}\,\mathfrak{S}^{\mathfrak{K}}(\square_0)\le e\,K_2^{\flat}\vartheta\,
\mathfrak{n}^{\mathrm{kept}}\,e^{-\frac3{16}\delta_Pd^{\wedge}(\square_0,R_n)}.
\end{equation}

$(\ast)$ \emph{The statements taken as given.} The rest of the proof uses, as stated, the following
parts of the stage bound at the site $(\mathrm{P})$:
\begin{itemize}
\item its old-block summation for $\mu$;
\item of its anchoring of elements for $\mathsf{M}_{*}$, only the anchoring itself, the contact
condition and the position factor $\pi_x$ (its treatment of the near groups is not used);
\item the element estimate at a fixed anchor;
\item the local anchor-sum lemma;
\item the weighted anchor-sum bound;
\item the anchoring bound;
\item properties (iv) and (v) of its class of elements;
\item the format of the outputs, in the form asserted in~(4): an output of the expansion that contains
a group $\hat e$ is in the format of the stage, the set $\mathcal{P}$ attached to it contains $(Y,m)$,
and it belongs to the class $\mathfrak{S}$.
\end{itemize}
These statements enter the argument below as hypotheses, in the form listed above; the present proof
does not establish them, and this is the point at which it is incomplete.
If those statements are read with an enlarged class of old blocks containing $\mathfrak{Old}$, nothing
below changes, since the anchors used here lie in $\mathcal{A}_v\subset\mathfrak{Old}$ by the first
clause of the definition of $\mathfrak{Old}$.

\emph{Bound on $\mu$.} Fix $q\in\mathcal{I}^{\mathrm{fl}}$ and consider the groups $\hat e$ with
$q\in\mathsf{b}(\hat e)$. The part of the $\mu$-summation of the stage bound that concerns
$\mathcal{Q}_v$ and $\mathcal{S}_j(v)$ is void, since both sets are empty.

The old-block summation is applied without change. Take $\square_0\in A$. Let $\Gamma_0:=q$ if
$q\subset\hat{\mathsf{K}}_{\square_0}$; otherwise let $\Gamma_0$ be the member of the animal of
$\hat{\mathsf{y}}$ through $q$ that touches $\hat{\mathsf{K}}_{\square_0}$.

At fixed $(v,\square_0,\Gamma_0)$, the element estimate at a fixed anchor bounds the sum of
$\tilde\nu(\hat e)$ over the corresponding groups by
$2\mathfrak{w}(v)e^{c_J}y_{\square_0}e^{\Sigma_v}\mathrm{a}(q,\Gamma_0)$, where $\mathrm{a}(q,\Gamma_0)$
is the animal weight of the element estimate. Here the remaining blocks of $A$
and the animals are summed by property (iv) of the element class, together with the inequality
$\prod e^{-1}\cdot e^{|A|}\le 1$. Moreover $\sum_{\Gamma_0}\mathrm{a}(q,\Gamma_0)\le 2$ by the
animal-summation condition on $\kappa_1$.

Use $\mathfrak{w}(v)=\mathfrak{O}(v)$ and $e^{\Sigma_v}\le e^{1/2}$. Interchange the sum over $v$ with
the sum over $\square_0$; the inner sum over $v$ with $\square_0\in\mathcal{A}_v$ is then
$\mathfrak{S}^{\mathfrak{K}}(\square_0)$. This gives
\begin{align*}
\sum_{\hat e:\ q\in\mathsf{b}(\hat e)}\tilde\nu(\hat e)
&\le 2\cdot 2e^{c_J+\frac12}\sideset{}{'}\sum_{\square_0}
y_{\square_0}\,\mathfrak{S}^{\mathfrak{K}}(\square_0)\\
&\le 4e^{c_J+\frac32}K_2^{\flat}\vartheta\,\mathfrak{n}^{\mathrm{kept}}
\sideset{}{'}\sum_{\square_0}e^{-\frac3{16}\delta_Pd^{\wedge}(\square_0,R_n)}\\
&\le 4e^{c_J+\frac32}K_2^{\flat}\vartheta\,\mathfrak{n}^{\mathrm{kept}}\,c_{\mathrm{loc}}
=\mu^{\mathrm{leaf}},
\end{align*}
where the prime on the sum restricts it to the blocks $\square_0$ such that
$\hat{\mathsf{K}}_{\square_0}$ contains $\Gamma_0$ or a member of $\hat{\mathsf{K}}_{\square_0}$
touches $\Gamma_0$.
The second inequality is \eqref{4.22:eq-stage-totals-b}, with $4e^{c_J+1/2}\cdot e=4e^{c_J+3/2}$. The
third is the local anchor-sum lemma at $\beta=\tfrac3{16}\delta_P$, which applies because
$\tfrac3{16}\delta_P\ge\tfrac1{16}\delta_P$ and the lemma holds at $\beta\ge\tfrac1{16}\delta_P$ by the
width condition.

No factor $V'$ and no position factor enter, and there is no sum over components, live pairs or strips.
The kept members of the different live pairs at $\square_0$ are summed inside
$\mathfrak{S}^{\mathfrak{K}}(\square_0)$. At most $2^d$ strips meet a probe, and this count is contained
in $\mathfrak{n}^{\mathrm{kept}}$.

\emph{Bound on $\mathsf{M}_{*}$.} Fix $x$ and sum $\tilde\nu(\hat e)$ over the groups of the kept
members with $x\in[\mathsf{d}^{\mathrm{rel}}(\hat e)]_k$. The anchoring of the stage bound is applied
without change. Let $\Delta_k(x)$ be the level-$k$ cube at $x$. Each group (recall $A\neq\emptyset$) is
anchored at one block $\square_0\in A\subset\mathcal{A}_v$, chosen as follows.
\begin{itemize}
\item If the contact $(\square_0,\Delta_k(x))$ is realised through
$\hat{\mathsf{K}}_{\square_0}\cup\operatorname{Cor}(\square_0)$ or through an animal on
$\hat{\mathsf{K}}_{\square_0}$, the anchor is the block realising it. The factor is at most
$\pi_x(\square_0)$, by the anchoring bound and property (iv) of the element class.
\item Otherwise the contact is through $S_v$, and the anchor is the first block of $A$, subject to the
contact condition.
\end{itemize}
The contact condition holds for $v$ trivially. Indeed, $\Delta_k(x)$ meets $S_v$, which is a cube of
level at most $k$ with a point on $Y^{\sharp}$, and $\widetilde{\square}_0$ meets $Y^{\sharp}$. By
part~(a) of the geometric lemma of the stage, with $s_k(Y)$ the size of $Y$ in level-$k$ units as in
that lemma, it follows that
\begin{equation*}
\operatorname{dist}_k(x,\square_0)\le 2+s_k(Y)+3\le 101\check{R}_k+5<204\check{R}_k,
\end{equation*}
and $\pi_x(\square_0)=1$.

The groups of $v$ anchored at $\square_0$ weigh at most
$2e^{c_J+\frac12}y_{\square_0}W(v)\pi_x(\square_0)$; the factor $e^{1/2}$ bounds $e^{\Sigma_v}$. Summing
this over all $\square_0\in\mathcal{A}_v$, rather than over the anchors only, over-counts, which is
harmless for an upper bound.

The anchors are old blocks only. Indeed $\mathcal{A}_v\subset\mathfrak{Old}$; the fattening of a window
block meets no strip; and no rim cube occurs, since $A\neq\emptyset$. This also proves~(2).

Since $W(v)=\mathfrak{O}(v)$, summing over $v$ gives $\mathfrak{S}^{\mathfrak{K}}(\square_0)$ at each
$\square_0$. By \eqref{4.22:eq-stage-totals-b} we then have
\begin{align*}
\sum_{\square_0\in\mathfrak{Old}}y_{\square_0}\pi_x(\square_0)\mathfrak{S}^{\mathfrak{K}}(\square_0)
&\le e\,K_2^{\flat}\vartheta\,\mathfrak{n}^{\mathrm{kept}}\sum_{\square_0}
e^{-\frac3{16}\delta_Pd^{\wedge}(\square_0,R_n)}\pi_x(\square_0)\\
&\le e\,K_2^{\flat}\vartheta\,\mathfrak{n}^{\mathrm{kept}}\,c_{\star\star}''V'.
\end{align*}
The last step is the weighted anchor-sum bound, which is stated at $\beta=\tfrac18\delta_P$. It applies
because $e^{-\frac3{16}\delta_Pd^{\wedge}}\le e^{-\frac18\delta_Pd^{\wedge}}$.

Multiplying by $2e^{c_J+1/2}$, the total is at most
\begin{equation*}
2e^{c_J+\frac32}K_2^{\flat}\vartheta\,c_{\star\star}''V'\,\mathfrak{n}^{\mathrm{kept}}
=\mathsf{M}^{\mathrm{leaf}}.
\end{equation*}
Together with the bound on $\mu$, this proves the first sentence of~(1).

\emph{Equivalence.} In the old-block line of the summation for $\mu$, the sum of the terms at
$\square_0$ is at most $C^{P,\star}_{\mathsf{a}}e^{\frac18\delta_Pd^{\wedge}(\square_0,R_n)}$. The kept
members add
\begin{equation*}
\mathfrak{S}^{\mathfrak{K}}(\square_0)\le e\,\mathfrak{n}^{\mathrm{kept}}
e^{\frac1{16}\delta_Pd^{\wedge}(\square_0,R_n)}
\le e\,\mathfrak{n}^{\mathrm{kept}}e^{\frac18\delta_Pd^{\wedge}(\square_0,R_n)}.
\end{equation*}
Hence the old-block lines of the old-block summation and of the anchoring run with
$C^{P,\star}_{\mathsf{a}}+e\,\mathfrak{n}^{\mathrm{kept}}$ in place of $C^{P,\star}_{\mathsf{a}}$.

Their members are $4e^{c_J+\frac12}K_2^{\flat}\vartheta c_{\mathrm{loc}}C^{P,\star}_{\mathsf{a}}$ in
$\mu^{\flat}$ and $2e^{c_J+\frac12}K_2^{\flat}\vartheta c_{\star\star}''V'C^{P,\star}_{\mathsf{a}}$ in
$\mathsf{M}^{\flat}$. Under the replacement they grow by exactly $\mu^{\mathrm{leaf}}$ and
$\mathsf{M}^{\mathrm{leaf}}$, since $4e^{c_J+1/2}\cdot e=4e^{c_J+3/2}$ and
$2e^{c_J+1/2}\cdot e=2e^{c_J+3/2}$.

The concluding step of the stage bound depends on the elements only through
$(\mathsf{b},\mathsf{d},\nu)$, by property (v) of the element class, so it applies to the enlarged
family unchanged. This proves the second sentence of~(1).

\emph{Allocation of the exponent.} Of the exponent $\tfrac12\delta_Pd^{\wedge}(\square_0,R_n)$ carried
by $\vartheta(\square_0)$:
\begin{itemize}
\item the share $\tfrac14$ pays for the corridor, by clause (b) of
Proposition~\ref{4.22:prop-preparatory-stage} and the corridor condition;
\item the share $\tfrac1{16}$ pays for the polynomial of the kept members; it lies inside the share
$\tfrac18$ of the old-block term sum, by clause (d) of
Proposition~\ref{4.22:prop-preparatory-stage};
\item of the remaining $\tfrac3{16}$, the local anchor-sum lemma needs $\tfrac1{16}$ (for $\mu$) and
the weighted anchor-sum bound needs $\tfrac18$ (for $\mathsf{M}_{*}$).
\end{itemize}
Nothing is spent twice: in each of the two bounds one factor serves one sum at one block. The share
$\tfrac1{16}\delta_P$ that pays for the factor $V'$ of clause (c) of
Proposition~\ref{4.22:prop-preparatory-stage} is used inside the weighted anchor-sum bound, for
$\mathsf{M}_{*}$. It is idle for $\mu$, because the local anchor-sum lemma carries no factor $V'$.

\emph{The premise $\nu(\hat e)\le\frac12\mu$.} Since the exponent in the definition of $\tilde\nu$ is
non-negative, $\nu(\hat e)\le\tfrac12\tilde\nu(\hat e)$. Two cases remain.
\begin{itemize}
\item Suppose $\mathsf{b}(\hat e)=\emptyset$. Then $v_{\hat e}$ depends on no integration variable, so
it equals its value at $b=0$. There the Landau layer vanishes, $\mathbb{H}(\sigma;0)=0$ for each value
of its first argument $\sigma$, and $\Phi_{t,\mathsf{s}}(0)$ is the base
configuration for every $t$. Since $A\neq\emptyset$, the derivative $\partial_{t_A}F_v$, from which
$v_{\hat e}$ is formed in clause~(a) of Proposition~\ref{4.22:prop-preparatory-stage}, vanishes. Hence
$v_{\hat e}=0$.
\item Otherwise pick $q\in\mathsf{b}(\hat e)$. Then
\begin{equation*}
\tilde\nu(\hat e)\le\sum_{e:\ q\in\mathsf{b}(e)}\tilde\nu(e)\le\mu
\le\mu^{\flat}+\mu^{\mathrm{leaf}}\le\tfrac14,
\end{equation*}
by~(1) and the first smallness condition. Hence $\nu(\hat e)\le\tfrac12\mu\le\tfrac18$.
\end{itemize}
This is not circular: the proof of $\mu\le\mu^{\flat}+\mu^{\mathrm{leaf}}$ above sums $\tilde\nu$ and
nowhere uses the premise $\nu\le\tfrac12$. This proves the displayed bound in~(4). The remaining
assertions of~(4) are among the statements taken as given at $(\ast)$.
\end{proof}

\begin{proposition}[Kept terms at a large-field operation on a foreign component]
\label{4.22:prop-foreign-operation}
Work at the presentation site $(\mathrm{L})$ of the large-field operation $\mathbf{R}_k$. Let $(Y,m)$
be a live large-field region lying in a component of $\Lambda_k^c$ that is not treated at
$\mathbf{R}_k$, and let $Z$ be the large-field region of $\mathbf{R}_k$. Let
$v\in\mathfrak{K}(Y)$ be a kept member of $(Y,m)$, with domain $X_v$, and put
$K_v:=Q_Y(k)$, the level-$k$ set attached to $Y$ that every label below contains. Assume the
following:
\begin{itemize}
\item clauses (c) and (d) of the presentation-site proposition, together with the first and
second conditions of the presentation site;
\item clause (a) of Lemma~\ref{4.22:lem-kept-holomorphy};
\item clause (ii)$_2$ of Lemma~\ref{4.22:lem-kept-oscillation}, display
\eqref{4.22:eq-kept-oscillation-a};
\item clauses (a), (c), (e), (f) of the geometry of live large-field regions
in Definition~\ref{4.22:not-live-region-geometry}, whose proof is outstanding;
\item Lemma~\ref{4.22:lem-probe-summand};
\item clauses (c) and (e) of the summation statement in
Definition~\ref{4.22:not-live-region-geometry}, whose proof is outstanding, and the floor on
$\kappa_1$ recorded there.
\end{itemize}
Then $v$ is a second-part unit, that is $X_v\cap Z=\emptyset$, and belongs to the second-part
class $\mathfrak{L}^{\mathrm{II}}$. Its four brackets resolve into pieces
$\mathfrak{p}=(v;s;\mathsf{y}_1,\dots,\mathsf{y}_4)$, with $s$ a law site of the presentation,
$\mathsf{y}_1\neq\emptyset$ and
labels $\operatorname{lab}_t$ as in clause (c) of the presentation-site proposition, all of
them exterior, and
\begin{equation}\label{4.22:eq-foreign-operation-1}
|\mathfrak{p}|\le \mathfrak{O}(v)\,e^{-(\kappa_1-1)\sum_t|\mathsf{y}_t|},
\qquad
\sum_t|\mathsf{y}_t|\ge \tfrac52\,\check{R}_m .
\end{equation}
The distance condition of the summation statement in
Definition~\ref{4.22:not-live-region-geometry} (whose proof is outstanding) --- the level-$k$
distance from $\operatorname{pa}(\Delta)$ to the term's label is at most
$c_\rho\,\rho(\Delta)$ for every member $\Delta$ seeing the term's region --- holds at each link
$t$ of the chain of clause (c) of the presentation-site proposition, with the label $K_v$ at
$t=1$ and $\operatorname{lab}_{t-1}$ at $t\ge2$, and the kept member survives at the chain's
lower end. For every member
$\Delta$ disjoint from the box $\Lambda^{\sim}$ of clause (f) of
Definition~\ref{4.22:not-live-region-geometry} (proof outstanding),
\begin{equation}\label{4.22:eq-foreign-operation-2}
e^{-\rho(\Delta)}\,\mathsf{S}^{\mathfrak{K}}(\Delta)\le C_\Sigma\,\mathfrak{s}_\Sigma(1).
\end{equation}
Consequently clause (d) of the presentation-site proposition holds, in its own notation, for
the sum $\mathfrak{v}^{\mathrm{rel}}+\mathfrak{v}^{\mathrm{kept}}$ of the relative and the kept
contributions, with the constant
\begin{equation}\label{4.22:eq-foreign-operation-3}
4C_\Sigma^{\mathrm{sum}}C_{\mathsf{a}}(1)\,\mathsf{b}_k(0)
\quad\text{replaced by}\quad
4C_\Sigma^{\mathrm{sum}}\bigl[C_{\mathsf{a}}(1)\,\mathsf{b}_k(0)
+\mathfrak{n}^{\mathrm{kept}}\bigr]
\end{equation}
inside the quantity $\varepsilon^{\mathrm{rel}}_L(k)$ of that clause, where
$\mathfrak{n}^{\mathrm{kept}}$ is the number of \eqref{4.22:eq-fixed-number-entry-a}; the factor
$e^{-(\check{R}_k-2)}$ in $\varepsilon^{\mathrm{rel}}_L(k)$ stands unchanged.
\end{proposition}

\begin{proof}
That $X_v\cap Z=\emptyset$ is clause (e) of the geometry of live large-field regions (whose
proof is outstanding), applied to the untreated component of $\Lambda_k^c$ holding $Y$; thus
$v$ lies in the second-part class $\mathfrak{L}^{\mathrm{II}}$.

\emph{The pieces.} We follow the proof of clauses (c) and (d) of the presentation-site
proposition for the class $\mathfrak{L}^{\mathrm{II}}$, with the norm $\mathfrak{N}^A(v)$
there replaced throughout by $\mathfrak{O}(v)$. The function to which that argument is applied
is the map $\operatorname{chain}(\vec{\mathsf{s}})\mapsto v^{\sharp}$, sending the chain of
variables $\vec{\mathsf{s}}$ to the orbit datum of $v$; by clause (a) of
Lemma~\ref{4.22:lem-kept-holomorphy} it is holomorphic on the product polydisc. The lower end
of each bracket does not depend on the first variable $\mathsf{s}^{(1)}$. Writing $f$ for the
function and $c$ for the lower end of a bracket, Cauchy's estimate on the product polydisc is
applied to $f-c$, and $|f-c|\le\mathfrak{O}(v)$ by clause (ii)$_2$ of
Lemma~\ref{4.22:lem-kept-oscillation}. This gives the pieces, their labels and the first bound
of \eqref{4.22:eq-foreign-operation-1}; the second bound is obtained below, from the own
collar of $Y$.

\emph{The distance condition.} At $t=1$, for a member $\Delta$ in the live strip
$Y^{\sharp}$, by clauses (a), (e), (f) of the geometry of live
large-field regions (whose proof is outstanding),
\begin{equation*}
\operatorname{dist}_k(\operatorname{pa}(\Delta),K_v)\le s_k(Y)+1\le 101\,\check{R}_k+1
\le c_\rho\cdot 10\,\check{R}_k\le c_\rho\,\rho(\Delta).
\end{equation*}
Here the second inequality is clause (a), the third is arithmetic with $c_\rho=120$, and the
last, $10\,\check{R}_k\le\rho(\Delta)$, is read from clauses (e) and (f), which separate the
live strip $Y^{\sharp}$ from $Z$ by at least $\check{R}_k$ and $Z^c$ from $\Lambda^{\sim}$ by at
least $9\,\check{R}_k$ level-$k$ units.
At $t\ge2$ the distance condition is the one in the list of clause (c) of the presentation-site
proposition, and it holds because $K_v\subset\operatorname{lab}_{t-1}$, the distance of
$\operatorname{pa}(\Delta)$ to $\operatorname{lab}_{t-1}$ being at most its distance to
$K_v$.

\emph{The depth of a member.} Let $\Delta\subset Y^{\sharp}$, and consider a chain of members,
consecutive ones sharing a wall, from $\Delta$ to a source cell; a source cell has level $k$ and
sits at a hub in $Z$. By clause (c) of the geometry of live large-field regions (whose proof is
outstanding), the own collar of $Y$, namely the cells within level-$m$ distance less than
$\tfrac52\,\check{R}_m$ of $Y^{\sharp}$, all of level at most $m$, lies in the component holding
$Y$, hence not in $Z$; so the chain leaves the own collar, and
$\rho(\Delta)\ge\tfrac52\,\check{R}_m$.
The same argument applies to every connected family of members holding a member of
$Y^{\sharp}$ and a source cell. Each component of the non-empty family $\mathsf{y}_1$ is such a
family, since it holds a source cell and a member seeing the region of $v$, which at $t=1$ is a
member of $Y^{\sharp}$; hence $\sum_t|\mathsf{y}_t|\ge|\mathsf{y}_1|\ge\tfrac52\,\check{R}_m$,
the second inequality of \eqref{4.22:eq-foreign-operation-1}.
In particular $\check{R}_m\le 2\rho(\Delta)+2$, which is the depth condition of
Lemma~\ref{4.22:lem-probe-summand} for the depth datum $\rho$. That lemma, in its form for
cells (the factor $2^d$ replaced by $1$), then gives \eqref{4.22:eq-foreign-operation-2}.

\emph{The size of the last label.} The label $\operatorname{lab}_4$ contains $Q_Y(k)$, which is
not contained in $Z$ and has a point on a live strip, and it contains a source cell in $Z$.
Hence, by clause (e) of the geometry of live large-field regions (whose proof is outstanding),
which separates the live strip from $Z$ by at least $\check{R}_k$ level-$k$ units,
\begin{equation*}
d_k(\operatorname{lab}_4)\ge\check{R}_k-2 .
\end{equation*}

\emph{Summation.} Clauses (c) and (e) of the summation statement (whose proof is outstanding)
are applied as in the proof of clause (d) of the presentation-site proposition, with the member
bound \eqref{4.22:eq-foreign-operation-2} for the kept summand and the lower bound on
$d_k(\operatorname{lab}_4)$ just obtained; this yields clause (d) for
$\mathfrak{v}^{\mathrm{rel}}+\mathfrak{v}^{\mathrm{kept}}$ with the constant
\eqref{4.22:eq-foreign-operation-3} and the factor $e^{-(\check{R}_k-2)}$ standing. The
hypothesis on $\kappa_1$ that these clauses require, namely
\begin{equation*}
\kappa_1-1\ge c_l\,(a+9)+\lambda_*,
\end{equation*}
with $a$ the number carried by their application in the proof of clause (d) of the
presentation-site proposition, follows from the floor on $\kappa_1$ in
Definition~\ref{4.22:not-live-region-geometry}, whose proof is outstanding.
\end{proof}

\begin{proposition}[Kept terms at the healing of their own component]\label{4.22:prop-own-healing}
Let $(Y,m)$ be a live large-field region, let $C\ni Y$ be the component of $\Lambda_k^{c}$ treated
at the operation $\mathbf{R}_k$, let $\mathsf{C}_C$ be the core of $C$, and let $h\ge m$ be the level
attached to $C$ in clause (g) of the geometry of live regions, a statement listed in
Notation~\ref{4.22:not-live-region-geometry} whose proof is outstanding: $\mathsf{C}_C$ is a box that
is a union of cells of levels at most $h$, of sides at most
$104L\check{R}_{h+1}$ level-$h$ units. Let
$d^{\wedge}$ be the weighted contour distance of the boundary terms $\mathbf{B}''_k$, with payment
$\mathsf{w}_{\wedge}$ per crossed zone. Assume the following statements.
\begin{itemize}
\item The conclusions of clause (iv) of Lemma~\ref{4.22:lem-kept-oscillation} (oscillation of the
kept terms on the orbit chart).
\item The conclusions of the lemma on deep units of the first part, with its depth datum, and of the
lemma bounding the exterior functional $\mathsf{Q}^{\star}_r$, two statements listed in
Notation~\ref{4.22:not-live-region-geometry} whose proofs are outstanding; and the relative
inequality at the presentation site $(\mathrm{L})$, with its constant $\mathsf{P}_L$.
\item Clauses (b) and (g) of the geometry of live regions and clause (b) of the cell-sum bound,
statements listed in Notation~\ref{4.22:not-live-region-geometry} (geometry and counting of live
large-field regions) whose proofs are outstanding; and the zone-width condition recorded there.
\item Clause (c) of the proposition on the presentation site $(\mathrm{L})$, a statement listed in
Notation~\ref{4.22:not-live-region-geometry} whose proof is outstanding.
\item The conclusions of Lemma~\ref{4.22:lem-probe-summand} (kept terms add one summand to the probe
functional).
\end{itemize}
Suppose that, for every live $(Y',m')$ with $Y'\subset C$ and every $v'\in\mathfrak{K}(Y')$, the
following hold: $v'$ is a deep unit of the first part whose domain is
$X_{v'}=Y'^{\sharp}\subset\mathsf{C}_C$; no member of a family meets $Y'^{\sharp}\subset\Lambda$; and a
member sees $Y'^{\sharp}$ only as a source cell at the hub, so that the distance condition of the
summation over families holds with distance at most $h_\sharp\le c_\rho$. Here $h_\sharp=112$ is the
side bound of clause (f) of the geometry of live regions, a statement listed in
Notation~\ref{4.22:not-live-region-geometry} whose proof is outstanding, and $c_\rho=120$ is the
constant of the distance condition. Fix $v\in\mathfrak{K}(Y)$
and put
\begin{equation}\label{4.22:eq-own-healing-1}
\|v\|_{\star}:=\tfrac12 C_{3F}\,
e^{-\frac12\delta_P d^{\wedge}(Y^{\sharp},\mathsf{C}_C^{c})}\,\mathfrak{O}(v).
\end{equation}
Then the following hold.
\begin{enumerate}
\item[(a)] Every piece $\mathfrak{p}$ of $v$ with a non-empty family $(\mathsf{y}_t)$ satisfies
$|\mathfrak{p}|\le\|v\|_{\star}\,e^{-(\kappa_1-1)\sum_t|\mathsf{y}_t|}$.
\item[(b)] The kept members inside $C$ satisfy
\begin{equation}\label{4.22:eq-own-healing-a}
\sum_{(Y',m'):\,Y'\subset C}\ \sum_{v'\in\mathfrak{K}(Y')}\mathfrak{O}(v')\,
e^{-\frac12\delta_P d^{\wedge}(Y'^{\sharp},\mathsf{C}_C^{c})}
\le e\,\mathfrak{n}^{\mathrm{kept}}\,C_{\mathrm{cell}}\cdot 2d\,
\bigl(104L\check{R}_{h+1}+2\bigr)^{d-1},
\end{equation}
where $C_{\mathrm{cell}}$ is the constant of clause (b) of the cell-sum bound, a statement listed in
Notation~\ref{4.22:not-live-region-geometry} whose proof is outstanding.
The right member does not depend on the past of the core: neither the number of live regions
inside $C$, nor their sizes, ages or pasts enter it.
\item[(c)] The exterior pieces of the kept members enter $\mathsf{Q}^{\star}_r$ for every $r$, since
$d_k(K_{v'})=0$ for the set $K_{v'}$ carried by each such member $v'$ in the functional
$\mathsf{Q}^{\star}_r$, with
$\tfrac12 C_{3F}$ times the right member of \eqref{4.22:eq-own-healing-a}.
Accordingly, the lemma bounding $\mathsf{Q}^{\star}_r$ and the relative inequality at
$(\mathrm{L})$ hold with the maximum defining $\mathsf{P}_L$ taken also over this summand.
\item[(d)] The value of $v$ and its interior pieces are constituents of the term
$V^{\sharp}(C^{\sharp})$ of the component $C$. Their oscillation obeys the oscillation bound for the
constituents of $V^{\sharp}(C^{\sharp})$, and their size obeys the modified form of
\cite[(1.69)]{B-CMP122b} used in this chapter; they enter no other bound of the step.
\end{enumerate}
\end{proposition}

\begin{proof}
\emph{Clause (a).} Apply clause (iv) of Lemma~\ref{4.22:lem-kept-oscillation} with $Y_i=Y$ and
$X=Y^{\sharp}\subset\mathsf{C}_C$; see \eqref{4.22:eq-kept-oscillation-a}. Its second inequality
holds for every piece whose family is not empty. It gives
\begin{equation*}
|\mathfrak{p}|\le\tfrac12 C_{3F}e^{-\frac12\delta_P D_X}\mathfrak{O}(v)\,
e^{-(\kappa_1-1)\sum_t|\mathsf{y}_t|},
\end{equation*}
where $D_X=d^{\wedge}(Y^{\sharp},\mathsf{C}_C^{c})$. By \eqref{4.22:eq-own-healing-1} this is
$\|v\|_{\star}e^{-(\kappa_1-1)\sum_t|\mathsf{y}_t|}$.

\emph{Clause (b), reduction to cells.} Let $(Y',m')$ be live with $Y'\subset C$. By clause (g) of
the geometry of live regions, a statement listed in Notation~\ref{4.22:not-live-region-geometry}
whose proof is outstanding, $\mathsf{C}_C$ is a union
of cells and $Y'^{\sharp}\subset\mathsf{C}_C$. Let $\Delta_{Y'}$ be a cell of $Y'^{\sharp}$ that holds
a point of $Y'^{\sharp}$ nearest to $\mathsf{C}_C^{c}$ in $d^{\wedge}$. Since
$\Delta_{Y'}\subset Y'^{\sharp}$ and $\Delta_{Y'}$ contains a nearest point, we have
$d^{\wedge}(\Delta_{Y'},\mathsf{C}_C^{c})=d^{\wedge}(Y'^{\sharp},\mathsf{C}_C^{c})$.

Clause (b) of the geometry of live regions, a statement listed in
Notation~\ref{4.22:not-live-region-geometry} whose proof is outstanding, says two things: strips of
distinct live regions are disjoint, and a cell lies in at most one strip. Hence $Y'\mapsto\Delta_{Y'}$ is
injective. The cell $\Delta_{Y'}$ meets $Y'^{\sharp}$ and $(Y',m')$ is live, so the summands
$\mathfrak{O}(v')$, $v'\in\mathfrak{K}(Y')$, all occur in the sum defining
$\mathsf{S}^{\mathfrak{K}}(\Delta_{Y'})$ in Lemma~\ref{4.22:lem-probe-summand}. All terms are
non-negative and the map is injective, so the left member of \eqref{4.22:eq-own-healing-a} is at most
\begin{equation}\label{4.22:eq-own-healing-2}
\sum_{\Delta\subset\mathsf{C}_C}
e^{-\frac12\delta_P d^{\wedge}(\Delta,\mathsf{C}_C^{c})}\,\mathsf{S}^{\mathfrak{K}}(\Delta).
\end{equation}

\emph{The depth datum.} The depth datum of the lemma on deep units of the first part, a statement
listed in Notation~\ref{4.22:not-live-region-geometry} whose proof is outstanding, is
$\varrho(\Delta):=1+d^{\wedge}(\Delta,\mathsf{C}_C^{c})/\mathsf{w}_{\wedge}$. Since $d^{\wedge}\ge0$,
we have $\varrho(\Delta)\ge1$. We check the condition imposed on a depth datum in
Lemma~\ref{4.22:lem-probe-summand}, namely $\check{R}_{m'}\le2\varrho(\Delta)+2$ whenever $\Delta$
meets the strip of a live $(Y',m')$.

Such a cell lies in $Y'^{\sharp}$ by clause (b) of the geometry of live regions, whose proof is
outstanding. Hence
$d^{\wedge}(\Delta,\mathsf{C}_C^{c})\ge D_{Y'}$, where $D_{Y'}:=d^{\wedge}(Y'^{\sharp},\mathsf{C}_C^{c})$.
Since $m'\le h$ by clause (g) of the geometry of live regions, whose proof is outstanding, clause (iv) of
Lemma~\ref{4.22:lem-kept-oscillation}, with $k_W=m'$, gives one of two lower bounds for $D_{Y'}$.
\begin{itemize}
\item If $m'=h$, then $D_{Y'}\ge\tfrac12\mathsf{w}_{\wedge}(\check{R}_{m'}+1)$. Hence
$\varrho(\Delta)\ge1+\tfrac12(\check{R}_{m'}+1)$, that is, $\check{R}_{m'}\le2\varrho(\Delta)-3$.
\item If $m'<h$, then $D_{Y'}\ge\tfrac52\mathsf{w}_{\wedge}\check{R}_{m'}$. Hence
$\check{R}_{m'}\le\tfrac25(\varrho(\Delta)-1)$.
\end{itemize}
In both cases $\check{R}_{m'}\le2\varrho(\Delta)\le2\varrho(\Delta)+2$.

\emph{Splitting the decay.} Write $\tfrac12\delta_P=\tfrac14\delta_P+\tfrac14\delta_P$. The
zone-width condition recorded in Notation~\ref{4.22:not-live-region-geometry}
gives $\tfrac1{16}\delta_P\mathsf{w}_{\wedge}\ge1$, hence $\tfrac14\delta_P\mathsf{w}_{\wedge}\ge1$.
Therefore
\begin{equation*}
\tfrac14\delta_P\, d^{\wedge}(\Delta,\mathsf{C}_C^{c})
\ge d^{\wedge}(\Delta,\mathsf{C}_C^{c})/\mathsf{w}_{\wedge}=\varrho(\Delta)-1,
\end{equation*}
that is, $e^{-\frac14\delta_P d^{\wedge}(\Delta,\mathsf{C}_C^{c})}\le e^{1-\varrho(\Delta)}$.

Now apply Lemma~\ref{4.22:lem-probe-summand} to cells, at $\beta=1$, with the depth datum just
checked; for a cell it gives
$\mathsf{S}^{\mathfrak{K}}(\Delta)\le C_\Sigma\bigl(2\varrho(\Delta)+2\bigr)^{p_\Sigma}$. With the
preceding bound this gives
\begin{equation}\label{4.22:eq-own-healing-3}
\begin{split}
e^{-\frac14\delta_P d^{\wedge}(\Delta,\mathsf{C}_C^{c})}\,\mathsf{S}^{\mathfrak{K}}(\Delta)
&\le e^{1-\varrho(\Delta)}\,C_\Sigma\bigl(2\varrho(\Delta)+2\bigr)^{p_\Sigma}\\
&\le e\,C_{\Sigma}\,\mathfrak{s}_{\Sigma}(1)\le e\,\mathfrak{n}^{\mathrm{kept}},
\end{split}
\end{equation}
where $\mathfrak{s}_\Sigma(1):=\sup_{\varrho\ge1}(2\varrho+2)^{p_\Sigma}e^{-\varrho}$, the supremum
being over $\varrho\ge1$ because $\varrho(\Delta)\ge1$. The last inequality holds because
$\mathfrak{n}^{\mathrm{kept}}$, the number of \eqref{4.22:eq-fixed-number-entry-a}, is $2^dC_\Sigma$
times the supremum of $(2\varrho+2)^{p_\Sigma}e^{-\varrho}$ over the depth data, hence at least
$C_\Sigma\mathfrak{s}_\Sigma(1)$.

\emph{The cell sum.} The remaining factor is summed by clause (b) of the cell-sum bound, a statement
listed in Notation~\ref{4.22:not-live-region-geometry} whose proof is outstanding. That bound holds at
every $\beta\ge\delta_P/16$; we take $\beta=\tfrac14\delta_P$ and let $S_C$ be the closure of
$\mathsf{C}_C^{c}$; the bound is stated through the cardinality $\#\partial^{+}S_C$ of the boundary
set $\partial^{+}S_C$ of $S_C$. By clause (g) of the geometry of live regions, whose proof is
outstanding,
$\mathsf{C}_C$ is a box of cells of levels at most $h$ with sides at most $104L\check{R}_{h+1}$
level-$h$ units. Hence
\begin{equation}\label{4.22:eq-own-healing-4}
\sum_{\Delta\subset\mathsf{C}_C}e^{-\frac14\delta_P d^{\wedge}(\Delta,\mathsf{C}_C^{c})}
\le C_{\mathrm{cell}}\,\#\partial^{+}S_C
\le C_{\mathrm{cell}}\cdot2d\,\bigl(104L\check{R}_{h+1}+2\bigr)^{d-1}.
\end{equation}

Inserting \eqref{4.22:eq-own-healing-3} termwise into \eqref{4.22:eq-own-healing-2} bounds
\eqref{4.22:eq-own-healing-2} by $e\,\mathfrak{n}^{\mathrm{kept}}$ times the left member of
\eqref{4.22:eq-own-healing-4}. With \eqref{4.22:eq-own-healing-4} this is
\eqref{4.22:eq-own-healing-a}. None of the constants used depends on the number, sizes, ages or
pasts of the live regions in $C$.

\emph{Clause (c).} By clause (c) of the proposition on the presentation site $(\mathrm{L})$, a
statement listed in Notation~\ref{4.22:not-live-region-geometry} whose proof is outstanding,
an exterior piece of a deep unit leaves $Z$ through a family. Its family is therefore non-empty, so
clause (iv) of Lemma~\ref{4.22:lem-kept-oscillation} applies in the form of clause (a).

Define $\|v'\|_{\star}$ by \eqref{4.22:eq-own-healing-1} with $Y'$, $v'$ in place of $Y$, $v$. Summing
$\|v'\|_{\star}$ over $(Y',m')$ with $Y'\subset C$ and over $v'\in\mathfrak{K}(Y')$ gives
$\tfrac12 C_{3F}$ times the left member of
\eqref{4.22:eq-own-healing-a}. This is at most $\tfrac12 C_{3F}$ times its right member. The weight
$e^{r d_k(K_{v'})}$ equals one because $d_k(K_{v'})=0$, so the entry is the same for every $r$. The
lemma bounding $\mathsf{Q}^{\star}_r$, a statement listed in
Notation~\ref{4.22:not-live-region-geometry} whose proof is outstanding, and the relative inequality at
$(\mathrm{L})$ then hold with $\mathsf{P}_L$ re-maximised over this summand.

\emph{Clause (d).} The oscillation bound for the constituents of $V^{\sharp}(C^{\sharp})$ applies
to the value and the interior pieces of $v$. This is because clause (iv) of
Lemma~\ref{4.22:lem-kept-oscillation} bounds the $\zeta$-difference for every tuple of families,
the empty tuple included. Their size is bounded by the modified form of \cite[(1.69)]{B-CMP122b}
used in this chapter.
\end{proof}

\begin{remark}[How each far factor is shared among the payments]\label{4.22:rem-far-factor-shares}
Let $(Y,m)$ be a live large-field region and $v\in\mathfrak{K}(Y)$ a kept member. Four sites
expand $v$ into pieces:
\begin{itemize}
\item the fluctuation integral $(\mathrm{T})$ of the step from level $i$ to level $i+1$
(Proposition~\ref{4.22:prop-fluctuation-site});
\item the large-field operation $\mathbf{R}_k$ when $Y$ lies in a component that $\mathbf{R}_k$
does not treat, written $(\mathrm{L}^{\circ})$ (Proposition~\ref{4.22:prop-foreign-operation});
\item a preparatory stage $(\mathrm{P})$ of $\mathbf{R}_k$, on its own chain or on a foreign one
(Proposition~\ref{4.22:prop-preparatory-stage}, and Proposition~\ref{4.22:prop-stage-totals},
which is stated with its proof incomplete at one marked step);
\item the healing $(\mathrm{L}^{\dagger})$ of the component containing $Y$ at $\mathbf{R}_k$
(Proposition~\ref{4.22:prop-own-healing}).
\end{itemize}
At each site every piece carries a far factor, exponential in a distance. The rate of that factor
is $\kappa_1-1$ per cell at $(\mathrm{T})$ and $(\mathrm{L}^{\circ})$, and $\frac12\delta_P$ per
unit of $d^{\wedge}$ at $(\mathrm{P})$ and $(\mathrm{L}^{\dagger})$.

The table below divides this rate into slices and lists what each slice pays. Each slice is
applied to the whole distance. The slices of one site add up to no more than the rate, so no
part of a far factor is spent twice. The slice called the anchor slice is the one in which the
factor $\check{R}_m^{p_\Sigma}$ carried by the kept members is paid. A slice called idle is not
used.

In the table:
\begin{itemize}
\item $\ell_{\wedge}$ is the unit of length of the weighted distance $d^{\wedge}$; the depth datum
is $\varrho=1+d^{\wedge}/\ell_{\wedge}$;
\item $a_{\mathrm{L}}:=\frac{31}{20}\kappa$ is the label rate of the foreign operation
(Proposition~\ref{4.22:prop-foreign-operation});
\item $c_l$, $\lambda_*$, $c_\rho$ and $\hat b$ are the constants of the summation lemma for
families of cells and of its counting bounds;
\item $\kappa_V$ and $\delta$ are the decay rate and the fraction of
Proposition~\ref{4.22:prop-fluctuation-site};
\item $m$ is the level of the large-field operation $\mathbf{T}'_m(Y)$ attached to the live
region $(Y,m)$, and $h$ is the level for which stage $n+1$ of $\mathbf{R}_k$ takes place at
level $h+n$;
\item the collar $R_n$, the block set $A$ of a group, the corridor $\operatorname{Cor}(\square)$,
the block weight $y_{\square}$, the constants $K_A$, $K_2^{\flat}$, $V'$ and the stage totals
$\mu$, $\mathsf{M}_*$ are those of Proposition~\ref{4.22:prop-preparatory-stage};
\item $\square_0$ is an old block whose fattening $\widetilde{\square}_0$ meets the strip
$Y^{\sharp}$ of $(Y,m)$;
\item the cells $\pi^{\mathrm{T}}$ and the core $Q_Y(i)$ are those of
Proposition~\ref{4.22:prop-fluctuation-site}; $K_v$ is the core of the unit $v$, with
$K_v=Q_Y(k)$ in Proposition~\ref{4.22:prop-foreign-operation}; $\Lambda$ and $\Lambda^{\sim}$ are
the regions named in Propositions~\ref{4.22:prop-foreign-operation}
and~\ref{4.22:prop-own-healing};
\item $\mathsf{C}_C$ is the core of Proposition~\ref{4.22:prop-own-healing}, and $D_X$ is the depth
that enters the decay lemma through that proposition; $\#\partial^{-}\mathsf{C}_C$ is the number
of boundary cells of the core $\mathsf{C}_C$;
\item $\beta$ denotes the rate argument of the cell-sum lemma and of the anchor bounds.
\end{itemize}

\begin{center}
\begin{tabular}{p{3.2cm}p{8.3cm}}
\hline
\multicolumn{2}{l}{Site $(\mathrm{T})$, step $i\to i+1$}\\
\hline
far factor & cells of $\pi^{\mathrm{T}}$: $e^{-(\kappa_1-1)|\mathsf{y}|}$; floor on $\kappa_1$:
$\kappa_1-1\ge c_l(2\kappa+4)+\lambda_*$\\
label, corridor & $c_l\kappa_V$ per cell: $e^{\kappa_V d_i(\operatorname{lab})}$
(summation lemma, clause (a))\\
counting & $\lambda_*-1$ per cell: $\log(e\hat b)$ for the families; $c_\rho$ and
$d\log L+\log 2$ for parents and finer cells (the four counting bounds)\\
anchor slice & $1$ per cell of a shortest family: $e^{-\rho(\Delta)}$, a weight already inside
$\mathsf{T}^{*}$ (the $-1$ of the third counting bound)\\
left for smallness & $c_l\cdot\frac14\delta\kappa$ per cell: $e^{-\frac14\delta\kappa(7\check{R}_i-2)}$
(summation lemma, clause (d)); idle: at least $c_l\kappa$ per cell\\
own collar, as crossed & $\rho(\Delta)\ge\frac52\check{R}_m$ cells of level at most $m$
(collar geometry)\\
\hline
\multicolumn{2}{l}{Site $(\mathrm{L}^{\circ})$, foreign}\\
\hline
far factor & cells off $\Lambda^{\sim}$, $\Lambda$: $e^{-(\kappa_1-1)\sum_t|\mathsf{y}_t|}$;
$\kappa_1-1\ge c_l(a_{\mathrm{L}}+9)+\lambda_*$\\
label, corridor & $c_la_{\mathrm{L}}$ per cell: $e^{a_{\mathrm{L}}d_k(\operatorname{lab}_4)}$;
$c_l\cdot 8$: the four additive constants $2$ in the label distances invoked in
Proposition~\ref{4.22:prop-foreign-operation}\\
counting & $\lambda_*-1$ per cell\\
anchor slice & $1$ per cell: $e^{-\rho(\Delta)}$\\
left for smallness & $c_l\cdot 1$ per cell: $e^{-(\check{R}_k-2)}$\\
own collar, as crossed & $\rho(\Delta)\ge\frac52\check{R}_m$ (collar geometry: the collar lies in
$Y$'s own, untreated component)\\
\hline
\multicolumn{2}{l}{Site $(\mathrm{P})$, own or foreign chain}\\
\hline
far factor & $d^{\wedge}$ of $\mathbb{B}^{(n)}_k$:
$\vartheta(\square)=\vartheta e^{-\frac12\delta_Pd^{\wedge}(\square,R_n)}$ per block of $A$; the
factor $e^{-(\kappa_1-1)|\hat{\mathsf{y}}|}$ of the families is not drawn on\\
label, corridor & $\frac14\delta_P$: $\operatorname{Cor}(\square)$, by the floor on $M$\\
counting & $\frac1{16}\delta_P$: the cell-sum lemma, clause (c), which produces $V'$, inside the
anchored block sum (for $\mathsf{M}_*$); idle for $\mu$ (the local block sum does not contain
$V'$)\\
anchor slice & $\frac1{16}\delta_P$ of the $\frac18\delta_P$ that the anchoring bound in clause (d)
of the proposition on the terms of the preparatory stage without kept members reserves for the
anchor; the kept members use it through $e^{\varrho}\le e\cdot e^{\frac1{16}\delta_Pd^{\wedge}}$,
Proposition~\ref{4.22:prop-preparatory-stage}, clause (d)\\
left for smallness & $\frac1{16}\delta_P$: the far components in the anchored block sum, where only
finiteness is used and no smallness; the smallness of the output comes from $\vartheta$ alone;
idle: $\frac1{16}\delta_P$ in $\mathsf{M}_*$, $\frac18\delta_P$ in $\mu$\\
own collar, as crossed & $d^{\wedge}(\square_0,R_n)\ge\ell_{\wedge}(\frac52\check{R}_m-5)$
(depth inequality at the stage)\\
\hline
\multicolumn{2}{l}{Site $(\mathrm{L}^{\dagger})$, own}\\
\hline
far factor & $d^{\wedge}$ of $\mathbf{B}''_k$: $e^{-\frac12\delta_PD_X}$ of the decay lemma
with constant $C_{3F}$ (clause (iv) of the kept-members lemma);
$e^{-(\kappa_1-1)\sum_t|\mathsf{y}_t|}$ goes whole to the summation lemma\\
label, corridor & none ($K_v$ is one cube)\\
counting & $\frac14\delta_P$: cell-sum lemma, clause (b): $\#\partial^{-}\mathsf{C}_C$\\
anchor slice & $\frac14\delta_P$: $e^{1-\varrho(\Delta)}$ (anchor bound for kept members at
$\beta=1$)\\
left for smallness & $e^{-(\check{R}_k-2)}$ from the families
($d_k(\operatorname{lab}_4)\ge 9\check{R}_k-1$)\\
own collar, as crossed & $D_X\ge\frac12\ell_{\wedge}(\check{R}_m+1)$ if $m=h$;
$D_X\ge\frac52\ell_{\wedge}\check{R}_m$ if $m<h$\\
\hline
\end{tabular}
\end{center}
\end{remark}

\begin{proof}
We take the four sites in turn. For each we say where every slice is used, and we check that the
slices do not exceed the rate.

\emph{Site $(\mathrm{T})$.} By Proposition~\ref{4.22:prop-fluctuation-site}, clause (b), each piece
satisfies $|v_{\mathsf{y}}|\le\mathfrak{O}(v)e^{-(\kappa_1-1)|\mathsf{y}|}$, the family
$\mathsf{y}$ consists of cells of $\pi^{\mathrm{T}}$, and
$|\mathsf{y}|\ge\rho(\Delta)\ge\frac52\check{R}_m$ for its member $\Delta\subset Y^{\sharp}$. This
gives the first and the last entry.

Clause (d) of that proposition is obtained from the summation lemma, clause (d), with label rate
$\kappa_V=(1-2\delta)\kappa$ and output rate $\frac14\delta\kappa$. Its hypothesis
$\kappa_1-1\ge c_l(\kappa_V+\frac14\delta\kappa)+\lambda_*$ follows from the floor on $\kappa_1$.
The slice $c_l\kappa_V$ pays the label weight $e^{\kappa_Vd_i(\operatorname{lab})}$, as in clause
(a) of the summation lemma. The slice $c_l\cdot\frac14\delta\kappa$ gives the factor
$e^{-\frac14\delta\kappa(7\check{R}_i-2)}$ of clause (d) of
Proposition~\ref{4.22:prop-fluctuation-site}. The slice $\lambda_*-1$ is spent on the counting
bounds.

The remaining unit per cell, taken along the cells of a shortest family, is the weight
$e^{-\rho(\Delta)}$. This weight already enters the anchor functional $\mathsf{T}^{*}$, as the $-1$
of the third counting bound, and it is the weight in which clause (c) of
Proposition~\ref{4.22:prop-fluctuation-site} pays the factor $\check{R}_m^{p_\Sigma}$.

It remains to check that these slices fit inside the rate. Since
$\kappa_V+\frac14\delta\kappa=(1-\frac74\delta)\kappa$, the floor on $\kappa_1$ gives
\begin{align*}
(\kappa_1-1)-\Bigl[c_l\kappa_V+(\lambda_*-1)+1+\tfrac14c_l\delta\kappa\Bigr]
&\ge c_l(2\kappa+4)-c_l\bigl(1-\tfrac74\delta\bigr)\kappa\\
&=c_l(\kappa+4)+\tfrac74c_l\delta\kappa\ \ge\ c_l\kappa .
\end{align*}
The last inequality uses that $c_l$, $\delta$ and $\kappa$ are positive. The slices therefore fit
inside the rate, and at least $c_l\kappa$ per cell is left idle.

\emph{Site $(\mathrm{L}^{\circ})$.} By Proposition~\ref{4.22:prop-foreign-operation}, each piece
satisfies
\begin{equation*}
|\mathfrak{p}|\le\mathfrak{O}(v)e^{-(\kappa_1-1)\sum_t|\mathsf{y}_t|},\qquad
\sum_t|\mathsf{y}_t|\ge\frac52\check{R}_m ,
\end{equation*}
and all its labels are exterior; this gives the first and the last entry.

The slice $c_l\cdot1$ gives the factor $e^{-(\check{R}_k-2)}$ that stands in the conclusion of
Proposition~\ref{4.22:prop-foreign-operation}. The slice $c_la_{\mathrm{L}}$ pays
$e^{a_{\mathrm{L}}d_k(\operatorname{lab}_4)}$, and $c_l\cdot8$ pays the four additive constants $2$
in the label distances. The unit slice is the weight $e^{-\rho(\Delta)}$ of the anchor bound for
cells, which pays $\check{R}_m^{p_\Sigma}$ as at $(\mathrm{T})$. By the floor on $\kappa_1$, the
slices add up to
\begin{equation*}
c_la_{\mathrm{L}}+8c_l+(\lambda_*-1)+1+c_l=c_l(a_{\mathrm{L}}+9)+\lambda_*\le\kappa_1-1 .
\end{equation*}

\emph{Site $(\mathrm{P})$.} By Proposition~\ref{4.22:prop-preparatory-stage}, clause (a), each
group satisfies
\begin{equation*}
|v_{\hat e}|\le\mathfrak{O}(v)\prod_{\square\in A}K_A\vartheta(\square)\cdot
e^{-(\kappa_1-1)|\hat{\mathsf{y}}|} .
\end{equation*}
Only the factors $\vartheta(\square)$ are drawn on.

By clause (b), the floor on $M$ bounds the weight of the corridor $\operatorname{Cor}(\square)$ by
$e^{\frac14\delta_Pd^{\wedge}(\square,R_n)}$, and the block weight satisfies
$y_{\square}\le K_2^{\flat}\vartheta e^{-\frac14\delta_Pd^{\wedge}(\square,R_n)}$.

By clause (d), every old block $\square_0$ whose fattening $\widetilde{\square}_0$ meets the strip
of $(Y,m)$ satisfies $d^{\wedge}(\square_0,R_n)\ge\ell_{\wedge}(\frac52\check{R}_m-5)$, whether
$R_n$ belongs to $Y$'s own component or to a foreign one, and
\begin{equation*}
\mathfrak{S}^{\mathfrak{K}}(\square_0)
\le e\,\mathfrak{n}^{\mathrm{kept}}e^{\frac1{16}\delta_Pd^{\wedge}(\square_0,R_n)} .
\end{equation*}
Here $\mathfrak{S}^{\mathfrak{K}}(\square_0)$ is the sum of $\mathfrak{O}(v)$ over the kept members
$v$ whose strip meets $\widetilde{\square}_0$. This is the anchor slice, namely $\frac1{16}\delta_P$
of the $\frac18\delta_P$ that the anchoring bound in clause (d) of the proposition on the terms of
the preparatory stage without kept members reserves for the anchor.

By these two bounds and $\frac14-\frac1{16}=\frac3{16}$, the product of the block weight
$y_{\square_0}$ and $\mathfrak{S}^{\mathfrak{K}}(\square_0)$ is at most
$eK_2^{\flat}\vartheta\,\mathfrak{n}^{\mathrm{kept}}e^{-\frac3{16}\delta_Pd^{\wedge}(\square_0,R_n)}$.

The totals are formed in the proof of Proposition~\ref{4.22:prop-stage-totals}, whose proof is
incomplete at one step. The total $\mu$ uses the local block sum at $\beta=\frac3{16}\delta_P$, and
$\frac3{16}\delta_P\ge\frac1{16}\delta_P$. That sum needs $\frac1{16}\delta_P$ and contains no
$V'$; the remaining $\frac18\delta_P$ is idle. The total $\mathsf{M}_*$ uses the anchored block sum,
stated at $\beta=\frac18\delta_P$, through
$e^{-\frac3{16}\delta_Pd^{\wedge}}\le e^{-\frac18\delta_Pd^{\wedge}}$. That $\frac18\delta_P$
consists of $\frac1{16}\delta_P$ producing $V'$ through the cell-sum lemma, clause (c), and
$\frac1{16}\delta_P$ for the far components, of which only finiteness is used; the remaining
$\frac1{16}\delta_P$ is idle. Thus $\frac12=\frac14+\frac1{16}+\frac3{16}$.

Each group is anchored at one block $\square_0\in A$, and $\vartheta(\square_0)$ enters a single sum
at that block. The totals contributed are proportional to $\vartheta$, so their smallness comes
from $\vartheta$ and not from any smallness of $\mathfrak{O}(v)$.

\emph{Site $(\mathrm{L}^{\dagger})$.} By Proposition~\ref{4.22:prop-own-healing}, every piece with a
non-empty family satisfies $|\mathfrak{p}|\le\|v\|_{\star}e^{-(\kappa_1-1)\sum_t|\mathsf{y}_t|}$.
The second factor goes whole to the summation lemma, and $\|v\|_\star$ carries the factor
$e^{-\frac12\delta_PD_X}$ with $D_X=d^{\wedge}(Y^{\sharp},\mathsf{C}_C^{c})$
(Proposition~\ref{4.22:prop-own-healing}), supplied by the decay lemma through clause (iv) of the
kept-members lemma. There is no label slice, since $K_v$ is one cube and $d_k(K_v)=0$.

In the proof of the bound on the sum over $(Y,m)$ with $Y\subset C$, the rate is split as
$\frac12\delta_P=\frac14\delta_P+\frac14\delta_P$. Since $\frac14\delta_P\ell_{\wedge}\ge1$ by the
floor on $\ell_{\wedge}$, one has $e^{-\frac14\delta_Pd^{\wedge}}\le e^{1-\varrho(\Delta)}$. The
anchor bound for kept members at $\beta=1$ then gives, for every cell $\Delta\subset\mathsf{C}_C$,
\begin{equation*}
e^{-\frac14\delta_Pd^{\wedge}(\Delta,\mathsf{C}_C^{c})}\,\mathsf{S}^{\mathfrak{K}}(\Delta)
\le e\,C_\Sigma\mathfrak{s}_\Sigma(1)\le e\,\mathfrak{n}^{\mathrm{kept}},
\end{equation*}
where $\mathsf{S}^{\mathfrak{K}}(\Delta)$ is the sum of the kept members' summands at the cell
$\Delta$; this is the anchor slice.
Here $\varrho(\Delta)$ obeys $\check{R}_m\le2\varrho$ in the two cases $m=h$ and $m<h$ of the last
entry. The other $\frac14\delta_P$ is used in the cell-sum lemma, clause (b), at
$\beta=\frac14\delta_P\ge\frac1{16}\delta_P$, which counts the boundary cells of $\mathsf{C}_C$.
The factor $e^{-(\check{R}_k-2)}$ comes from the families.

The two slices exhaust the rate $\frac12\delta_P$. At none of the four sites does the sum of the
slices exceed the rate of the far factor.
\end{proof}

\begin{corollary}[Kept terms enter every later bound through one fixed number. Stated; the proof
below is incomplete at the step marked $(\ast)$]
\label{4.22:cor-fixed-number-entry}
Assume the following:
\begin{itemize}
\item the containment conditions named in hypothesis (H4) (Definition~\ref{2.2:not-hypotheses});
\item the lemmas on the preparatory stages of the large-field operation that
Proposition~\ref{4.22:prop-stage-totals} takes by statement;
\item hypotheses (H1)--(H2);
\item the two-point bound on the kept potential $V$, with constants $C_\Sigma$ and $p_\Sigma$;
\item the two bounds on the kept quadratic forms $\mathsf{Q}_i$;
\item the standing condition on the kept terms at the preparatory stages under which
Lemma~\ref{4.22:lem-kept-holomorphy} is stated.
\end{itemize}
Then the following hold:
\begin{itemize}
\item the opening kept-term statement of this section;
\item the holomorphy of the kept terms along all later sites together with its companion bounds,
which make up parts (b) and (c) of the three-part statement on the kept terms
(Lemma~\ref{4.22:lem-kept-holomorphy});
\item the closing kept-term statement of this section.
\end{itemize}
Moreover, clauses (i)--(iv) below hold. The sites of a step are: a fluctuation site
$(\mathrm{T})$; a preparatory stage $(\mathrm{P})$ of a large-field operation; a large-field
operation on a foreign component $(\mathrm{L}^{\circ})$; and the healing of a large-field
region's own component $(\mathrm{L}^{\dagger})$.
\begin{enumerate}
\item[(i)] Let $(Y,m)$ be a live large-field region created at step $m$. At every later site,
the two kept members $v\in\mathfrak{K}(Y)$, the kept potential $V$ and the kept quadratic form,
enter the bound at that site only through their
oscillation norms $\mathfrak{O}(v)$ of \eqref{4.22:eq-kept-oscillation-a}. Consider their joint
contribution to the anchor functional of that bound; this functional is $\mathsf{T}^{*}$ at
$(\mathrm{T})$ and $(\mathrm{L}^{\circ})$, it is
$\mathfrak{S}(\square_0)e^{-\frac18\delta_P d^{\wedge}}$ at $(\mathrm{P})$, and it is
$e^{-\frac14\delta_P d^{\wedge}}\mathsf{S}(\Delta)$ at $(\mathrm{L}^{\dagger})$. At
$(\mathrm{T})$ and $(\mathrm{L}^{\circ})$ this joint contribution is at most
\begin{equation}\label{4.22:eq-fixed-number-entry-a}
\mathfrak{n}^{\mathrm{kept}}
=2^dC_\Sigma\sup_{\varrho\ge1}\,(2\varrho+2)^{p_\Sigma}e^{-\varrho},
\end{equation}
and at $(\mathrm{P})$ and $(\mathrm{L}^{\dagger})$ it is at most $e\,\mathfrak{n}^{\mathrm{kept}}$.
The number $\mathfrak{n}^{\mathrm{kept}}$ is fixed in the order of choice before $\gamma$. It
depends on none of the following: $\check R_m$; the depth $D$; the age $i-m$ at the current
step $i$; the coupling; the number, the sizes or the pasts of the large-field regions.
\item[(ii)] The share of the kept terms in the bound at each site is
$\mathfrak{n}^{\mathrm{kept}}$, times constants of that bound, times the small factor of that
bound. Specifically, the share is:
\begin{itemize}
\item the number $\varepsilon^{\mathrm{kept}}_i$ of Proposition~\ref{4.22:prop-fluctuation-site}
at $(\mathrm{T})$;
\item at $(\mathrm{L}^{\circ})$, with the constants $C_\Sigma^{\mathrm{sum}}$ and $C_e$ of the
chain of clause (d) of the chapter's general proposition on the treatment of a component at a
large-field operation (site $(\mathrm{L})$), the number
\begin{equation*}
4\bigl(1+C_\Sigma^{\mathrm{sum}}C_e\bigr)^3C_\Sigma^{\mathrm{sum}}\,
\mathfrak{n}^{\mathrm{kept}}\,e^{-(\check R_k-2)};
\end{equation*}
\item at $(\mathrm{P})$, with the stage constant $c_J$ of the preparatory stages, the numbers
through which the kept terms enter the two stage totals,
\begin{equation}\label{4.22:eq-fixed-number-entry-1}
\mu^{\mathrm{leaf}}=4e^{c_J+3/2}K_2^{\flat}\,\vartheta\,c_{\mathrm{loc}}\,
\mathfrak{n}^{\mathrm{kept}},
\qquad
\mathsf{M}^{\mathrm{leaf}}=2e^{c_J+3/2}K_2^{\flat}\,\vartheta\,c_{\star\star}''\,V'\,
\mathfrak{n}^{\mathrm{kept}}.
\end{equation}
The first enters the stage total $\mu$ and does not involve $V'$; the second enters the stage
total $\mathsf{M}_*$. Thus $e\,\mathfrak{n}^{\mathrm{kept}}$ stands in place of the anchoring
constant $C^{P,\star}_{\mathsf{a}}$ of the stage totals in the members
of the stage totals that come from old blocks, and
\begin{equation*}
\mu\le\mu^\flat+\mu^{\mathrm{leaf}},\qquad
\mathsf{M}_*\le\mathsf{M}^\flat+\mathsf{M}^{\mathrm{leaf}},
\end{equation*}
as in Proposition~\ref{4.22:prop-stage-totals}, whose proof is incomplete at one step (see
\eqref{4.22:eq-stage-totals-a});
\item at $(\mathrm{L}^{\dagger})$, where $C$ is the component containing $Y$ that is treated at
the operation $\mathbf{R}_k$, $\mathsf{C}_C$ is its core and $h$ is the level of
Proposition~\ref{4.22:prop-own-healing}, with the cell-count constant $C_{\mathrm{cell}}$ of
\eqref{4.22:eq-own-healing-a}, the number
\begin{equation*}
\tfrac12C_{3F}\,e\,\mathfrak{n}^{\mathrm{kept}}\,C_{\mathrm{cell}}\,2d\,
\bigl(104L\check R_{h+1}+2\bigr)^{d-1}
\end{equation*}
in the quantity $\mathsf{Q}^{\star}$ of Proposition~\ref{4.22:prop-own-healing}, under the
factor $e^{-(\check R_k-2)}$.
\end{itemize}
Each of these shares tends to zero together with the bound at its site, in supremum form.
\item[(iii)] None of the lower bounds imposed on the decay rates changes. At $(\mathrm{T})$
and $(\mathrm{L}^{\circ})$, the part of the rate on which the anchor is paid is the unit per
cell already contained in the constant $\lambda_*$ of the summation lemma for $\mathsf{T}^{*}$.
At $(\mathrm{P})$ and $(\mathrm{L}^{\dagger})$, it lies inside the parts of the rate that
clause (d) of the chapter's general proposition on the preparatory stages and the lemma on deep
first-part units already reserve under their rate condition. In particular, no additional unit
of rate is spent on the kept terms.
Three ceilings on $\gamma$ are taken again as maxima that include the number
$\mathfrak{n}^{\mathrm{kept}}$, in supremum form, with the powers of the variable $\mathsf{T}$
in which the three ceilings are written unchanged:
\begin{itemize}
\item the relative ceiling at $(\mathrm{T})$, into which $\mathfrak{n}^{\mathrm{kept}}$ enters
through the share $\varepsilon^{\mathrm{kept}}_i$ of Proposition~\ref{4.22:prop-fluctuation-site};
\item the relative ceiling at $(\mathrm{L})$ of \eqref{4.22:eq-own-healing-a};
\item the list of six ceilings of the preparatory stages in \eqref{4.22:eq-stage-totals-a}
(Proposition~\ref{4.22:prop-stage-totals}, whose proof is incomplete at one step).
\end{itemize}
The list of six is imposed at both values $\vartheta$ and
$\vartheta^{\natural}=\vartheta/\mathsf{w}_0$ of the parameter at which the chapter applies the
preparatory stages, one for each of the two applications, $\mathsf{w}_0$ being the fixed divisor
relating the second application's parameter to the first's. At each occurrence of $\mu^\flat$
and $\mathsf{M}^\flat$ in the list that Proposition~\ref{4.22:prop-stage-totals} designates for
the replacement, $\mu^\flat$ is replaced by
$\mu^\flat+\mu^{\mathrm{leaf}}$ and
$\mathsf{M}^\flat$ by $\mathsf{M}^\flat+\mathsf{M}^{\mathrm{leaf}}$. These sums stand in
place of the constants $\mu_J$ and $\mathsf{M}_J$, which do not occur in the list
(Proposition~\ref{4.22:prop-stage-totals}, whose proof is incomplete at one step).
\item[(iv)] In this section the reference comparison of couplings and the two-sided flow
comparison are nowhere used to pay for, convert or compare the factor
$\check R_m^{p_\Sigma}$. They enter only in three places:
\begin{itemize}
\item the age clause of part (c) of Lemma~\ref{4.22:lem-kept-holomorphy};
\item the display of the three-factor bound, as a ratio of couplings paid for by the zones
crossed;
\item the bound supplied for the orbit chart on which the oscillation norms of
\eqref{4.22:eq-kept-oscillation-a} are taken.
\end{itemize}
No implication producing the two-point bound on $V$ is used. The following are not used by the
closing kept-term statement: clauses (ii) for $\mathsf{Q}$ and (i) of
\eqref{4.22:eq-kept-oscillation-a}; the logarithmic radius bound; and parts (a) and (b) of
the three-part statement on the kept terms.
\end{enumerate}
\end{corollary}

\begin{proof}
Parts (b) and (c) of the three-part statement on the kept terms, among them the holomorphy of
the kept terms along all later sites, are Lemma~\ref{4.22:lem-kept-holomorphy}. Of that lemma's
hypotheses, the containment conditions of (H4), the standing condition at the preparatory stages
and, for $v=\mathsf{Q}_i$, the two bounds on the kept quadratic forms are hypotheses of the
corollary. The opening and closing kept-term statements, and clauses
(i)--(iv), follow from Lemma~\ref{4.22:lem-probe-summand}, Lemma~\ref{4.22:lem-kept-holomorphy}
and Propositions~\ref{4.22:prop-fluctuation-site}, \ref{4.22:prop-preparatory-stage},
\ref{4.22:prop-foreign-operation} and~\ref{4.22:prop-own-healing}, as follows.

\emph{Clause (i).} Fix a probe $\wp$ of the class of Lemma~\ref{4.22:lem-probe-summand}. By
Lemma~\ref{4.22:lem-probe-summand},
applied under the two-point bound on $V$ and the two bounds on the kept quadratic forms, the kept
members of the live regions met by $\wp$ add one summand $\mathsf{S}^{\mathfrak{K}}(\wp)$ to
the probe functional. This summand satisfies
\begin{equation*}
\mathsf{S}^{\mathfrak{K}}(\wp)\le 2^dC_\Sigma\bigl(2\varrho(\wp)+2\bigr)^{p_\Sigma}
\le\mathfrak{n}^{\mathrm{kept}}(\beta)\,e^{\beta\varrho(\wp)}.
\end{equation*}
We take $\beta=1$, so that $\mathfrak{n}^{\mathrm{kept}}(1)$ is the number
\eqref{4.22:eq-fixed-number-entry-a}. The members enter only through $\mathfrak{O}(v)$,
because this is the only quantity of $v$ on which that bound depends.

The summand is added to the anchor functional and is never multiplied into it. At
$(\mathrm{T})$ and $(\mathrm{L}^{\circ})$ the anchor functional $\mathsf{T}^{*}$ carries the
weight $e^{-\rho(\Delta)}$. Propositions~\ref{4.22:prop-fluctuation-site}
and~\ref{4.22:prop-foreign-operation} apply the above bound at each member $\Delta$ against
this weight, and this gives the bound $\mathfrak{n}^{\mathrm{kept}}$ at these two sites. At
$(\mathrm{P})$ and $(\mathrm{L}^{\dagger})$, Propositions~\ref{4.22:prop-preparatory-stage}
and~\ref{4.22:prop-own-healing} apply the same bound against the weights
$e^{-\frac18\delta_Pd^{\wedge}}$ and $e^{-\frac14\delta_Pd^{\wedge}}$ respectively. At
$(\mathrm{P})$ the depth datum satisfies $e^{\varrho}\le e\,e^{\frac1{16}\delta_Pd^{\wedge}}$,
so that the factor $e^{\varrho}$ is paid, up to the factor $e$, by a sixteenth of the rate taken
from the eighth in the weight $e^{-\frac18\delta_Pd^{\wedge}}$; at $(\mathrm{L}^{\dagger})$ the
choice $\beta=1$ leaves $e^{1-\varrho(\Delta)}$ against $e^{-\frac14\delta_Pd^{\wedge}}$. In
both cases the factor $e$ remains, whence $e\,\mathfrak{n}^{\mathrm{kept}}$.

The right side of \eqref{4.22:eq-fixed-number-entry-a} contains only $d$, $C_\Sigma$ and
$p_\Sigma$. It therefore depends on none of the quantities listed in (i), and it is fixed with
$C_\Sigma$ and $p_\Sigma$, before $\gamma$.

\emph{Clause (ii).} We take the sites in turn.
\begin{itemize}
\item At $(\mathrm{T})$ the share is the number $\varepsilon^{\mathrm{kept}}_i$ of
Proposition~\ref{4.22:prop-fluctuation-site}.
\item At $(\mathrm{L}^{\circ})$, Proposition~\ref{4.22:prop-foreign-operation} shows that
clause (d) of the chapter's general proposition on the treatment of a component at a
large-field operation (site $(\mathrm{L})$) holds with
the kept summand included. Through the chain of that clause, a contribution to
$\mathsf{T}^{*}$ is carried into the bound with the factor
$\bigl(1+C_\Sigma^{\mathrm{sum}}C_e\bigr)^3\cdot4C_\Sigma^{\mathrm{sum}}$. By (i), the part of
$\mathsf{T}^{*}$ due to the kept terms is at most $\mathfrak{n}^{\mathrm{kept}}$. Under the
small factor $e^{-(\check R_k-2)}$ of that bound, this gives the displayed share.
\item At $(\mathrm{P})$: $(\ast)$ the bounds $\mu\le\mu^\flat+\mu^{\mathrm{leaf}}$ and
$\mathsf{M}_*\le\mathsf{M}^\flat+\mathsf{M}^{\mathrm{leaf}}$, with the numbers
\eqref{4.22:eq-fixed-number-entry-1}, are the stage totals of
Proposition~\ref{4.22:prop-stage-totals}, whose proof is incomplete at one step. They are
obtained there by inserting the bound $e\,\mathfrak{n}^{\mathrm{kept}}$ from (i) in the place
of $C^{P,\star}_{\mathsf{a}}$ in the members that come from old blocks.
\item At $(\mathrm{L}^{\dagger})$, Proposition~\ref{4.22:prop-own-healing} sets, for each
deep kept member,
\begin{equation*}
\|v\|_\star:=\tfrac12C_{3F}\,e^{-\frac12\delta_Pd^{\wedge}(Y^{\sharp},\mathsf{C}_C^{c})}\,
\mathfrak{O}(v).
\end{equation*}
Summing over the regions $(Y,m)$ with $Y\subset C$ and over $v\in\mathfrak{K}(Y)$, the bound
\eqref{4.22:eq-own-healing-a} gives
\begin{equation*}
\sum_{(Y,m):\,Y\subset C}\ \sum_{v\in\mathfrak{K}(Y)}\|v\|_\star
\le\tfrac12C_{3F}\,e\,\mathfrak{n}^{\mathrm{kept}}\,C_{\mathrm{cell}}\,2d\,
\bigl(104L\check R_{h+1}+2\bigr)^{d-1}.
\end{equation*}
Proposition~\ref{4.22:prop-own-healing} places this sum in $\mathsf{Q}^{\star}$, under the
factor $e^{-(\check R_k-2)}$.
\end{itemize}
Each share is $\mathfrak{n}^{\mathrm{kept}}$ times constants of the bound at its site times the
small factor of that bound: at $(\mathrm{T})$ the share is the number
$\varepsilon^{\mathrm{kept}}_i$ of Proposition~\ref{4.22:prop-fluctuation-site}; at
$(\mathrm{L}^{\circ})$ the small factor is $e^{-(\check R_k-2)}$; at $(\mathrm{P})$ it is
$\vartheta$; at $(\mathrm{L}^{\dagger})$ the factor $(104L\check R_{h+1}+2)^{d-1}$ stands
against $e^{-(\check R_k-2)}$. Each share therefore tends to zero with the bound at its site, in
supremum form.

\emph{Clause (iii).} This is read off from the four propositions. At $(\mathrm{T})$ and
$(\mathrm{L}^{\circ})$, Propositions~\ref{4.22:prop-fluctuation-site}
and~\ref{4.22:prop-foreign-operation} pay the anchor with the weight $e^{-\rho(\Delta)}$,
which is already inside $\mathsf{T}^{*}$; this is the unit per cell contained in $\lambda_*$.
At $(\mathrm{P})$ and $(\mathrm{L}^{\dagger})$, Propositions~\ref{4.22:prop-preparatory-stage}
and~\ref{4.22:prop-own-healing} pay it from the part of the rate already reserved for the
anchor by clause (d) of the chapter's general proposition on the preparatory stages and by the
lemma on deep first-part units. The
three ceilings then change only through the numbers of (ii). These enter, at $(\mathrm{T})$,
through $\varepsilon^{\mathrm{kept}}_i$ in the relative ceiling of
Proposition~\ref{4.22:prop-fluctuation-site}, and otherwise as stated in
\eqref{4.22:eq-own-healing-a} and \eqref{4.22:eq-stage-totals-a}, the latter with the
replacements of $\mu^\flat$ and $\mathsf{M}^\flat$ made in
Proposition~\ref{4.22:prop-stage-totals}, whose proof is incomplete at one step.

\emph{Clause (iv).} Clause (iv) concerns the proofs of this section as a whole. The reference
comparison of couplings and the two-sided flow comparison occur only inside the three statements
listed in (iv), and there not to pay for, convert or compare $\check R_m^{p_\Sigma}$; the proofs
of (i)--(iii) above do not use them. These proofs use the two-point bound on $V$ only as
assumed. The closing kept-term statement draws on none of the clauses listed at the end
of (iv).
\end{proof}

\begin{remark}[Remarks on the own collar, foreign stages and additivity]\label{4.22:rem-collar-remarks}
Throughout this remark $Y^{\sharp}$ is a strip carrying kept members; $m$ is the level at which $Y^{\sharp}$ was created, so that at a site of level $i$ its age is $i-m$, and the cells of its own collar have level at most $m$ (in this section $m$ denotes a level, not the torus index); $k$ is the current level, and $(\mathrm{T})$, $(\mathrm{L}^{\circ})$, $(\mathrm{P})$, $(\mathrm{L}^{\dagger})$ are the sites of the share table in Remark~\ref{4.22:rem-far-factor-shares}. By $d^{\wedge}$ we mean the weighted contour distance. We write $\mathsf{w}_{\wedge}$ for the constant of the lower bounds for $d^{\wedge}$ in terms of crossed cubes.

\begin{enumerate}
\item[(1)] \emph{The collar pays at a coarse level.} Let a term be read at a site of level $i$ much larger than $m$. Measured in units of level $i$, the collar has width
\begin{equation*}
\tfrac52\,\check{R}_m\,L^{-(i-m)},
\end{equation*}
which is negligible. The payment does not come from this width. The cell clause cuts every cell in its proper scale, and $d^{\wedge}$ counts $M_1$-units of each cell's own level. The cells of the collar have level at most $m$. Hence a family of cells, or a contour, that reaches $Y^{\sharp}$ counts at least $\tfrac52\check{R}_m$ such cells, whatever $i$ is. This is why the first kept-member clause of Corollary~\ref{4.22:cor-fixed-number-entry}, whose proof is incomplete at one step, places the kept members under the cell clause. The decay lemma of the correlation theorem, whose decay $\eta_p$ is measured in units of the level at which the term is read, is never applied to them.

\item[(2)] \emph{The stage site over a foreign strip.} Let $\square_0$ be a block at the site $(\mathrm{P})$ lying over a foreign strip. Three payments are made at $\square_0$: for the corridor, for the counting of blocks, and for the anchor, as in the share table. All three draw on one distance $d^{\wedge}(\square_0,R_n)$ to the region $R_n$ of the $n$-th stage, as in the entry for the site $(\mathrm{P})$ in Remark~\ref{4.22:rem-far-factor-shares}. This distance satisfies both
\begin{equation*}
d^{\wedge}(\square_0,R_n)\ \ge\ \mathsf{w}_{\wedge}\,(\check{R}_k-5)
\quad\text{and}\quad
d^{\wedge}(\square_0,R_n)\ \ge\ \mathsf{w}_{\wedge}\bigl(\tfrac52\check{R}_m-5\bigr).
\end{equation*}
The first lower bound is at the current scale. It pays the corridor, which is measured in level-$k$ cubes. The corridor's count $\mathsf{x}_{\square}$ of level-$k$ cubes satisfies
\begin{equation*}
\mathsf{x}_{\square}+1\le 101\,\check{R}_k+2 .
\end{equation*}
The second lower bound is at the old scale. It pays the factor $\check{R}_m^{p_{\Sigma}}$.

The payments draw on disjoint slices of the rate $\tfrac12\delta_P$: the corridor on $\tfrac14\delta_P$, and the counting and the anchor on slices $\tfrac1{16}\delta_P$ each. Each slice multiplies the whole distance $d=d^{\wedge}(\square_0,R_n)$, so that
\begin{equation*}
e^{-(\frac14+\frac1{16}+\frac1{16})\delta_P d}
=e^{-\frac14\delta_P d}\,e^{-\frac1{16}\delta_P d}\,e^{-\frac1{16}\delta_P d}.
\end{equation*}
The first factor on the right, belonging to the corridor, is estimated with the first lower bound on $d$, and the third, belonging to the anchor, with the second. Neither lower bound is therefore used up by the other.

The number of blocks over $Y^{\sharp}$ is never counted: the block sum of the cell-counting statement sums all blocks at once at the rate $\tfrac1{16}\delta_P$. The number of regions carrying kept members is never counted either: a probe $\wp$ meets at most $2^d$ strips, and the blocks over them are summed by the local block sum, the doubled block sum, the diamond sum, the weight of the first display of the block-by-block estimate of the large-field stages, and the anchoring steps of that estimate, all taken by statement.

\item[(3)] \emph{An interrupted own preparation.} Suppose that the strip's own preparation is interrupted at the stage $h_1=m$. What is left is exactly the collar width $\tfrac52\check{R}_m$ of the third clause of the collar geometry statement. The share table of Remark~\ref{4.22:rem-far-factor-shares} claims no more than this.

\item[(4)] \emph{Add, never multiply.} The lemma on sums at window cubes concerns the sum of the terms at one cover cube of a window stage $j>k_0$. It states that this sum carries the factor $\mathsf{l}_j^{2}$, with $k_0$ and $\mathsf{l}_j$ as in that lemma, and that this factor cannot be dropped. That sum therefore may not be multiplied into the stage constant $V'$. Accordingly, in each of the propositions on which Corollary~\ref{4.22:cor-fixed-number-entry} (whose proof is incomplete at one step) rests, the kept members enter as a further summand of an anchor functional. No product of two counts is formed. The kept-member proposition at the stage site $(\mathrm{P})$ meets no window cube at all.

\item[(5)] \emph{What is not asserted.} The following statements are not asserted or proved in this section:
\begin{itemize}
\item the collar geometry, cell-counting, summation and anchoring statements on which the share table draws, and the stage-site relative statement;
\item the local block sum, the doubled block sum, the diamond sum, the weight of the first display of the block-by-block estimate and its anchoring steps, named in item (2) above, together with the anchored form of the stage bound and the second multiplicity estimate of the block-by-block estimate;
\item the fifth clause of the stage-site proposition at window anchors;
\item the relative comparison at complex arguments;
\item the two-point bound on the kept members, with constants $C_{\Sigma}$ and $p_{\Sigma}$, and the two accompanying conditions stated over the sourced set $\mathfrak{S}_{\mathrm{src}}(Y^{\sharp})$;
\item the three further members of the list entering the two-point bound.
\end{itemize}
Those three members change only the constant $C_{\Sigma}$ of $\Sigma^{V}$. They change no line of the argument of this section.

\item[(6)] \emph{The same bound counted in cells.} Let $\operatorname{dep}(Y)$ be the number of cells one must cross to leave the core of $Y$. Put
\begin{equation*}
\kappa_{\boxplus}=\min\bigl\{\kappa_1-1,\ \delta_P\,\mathsf{w}_{\wedge}/c_{\boxplus}\bigr\},
\end{equation*}
with $c_{\boxplus}$ the cell-conversion constant of the pointwise kept-member bound. Counted in cells, the entry of the kept members in the bound at the own presentation site $(\mathrm{L}^{\dagger})$ reads
\begin{equation*}
\tfrac12\,C^{\zeta}\,e^{-\frac12\kappa_{\boxplus}\operatorname{dep}(Y^{\sharp})}\,\mathfrak{O}(v),
\end{equation*}
where $C^{\zeta}$ is the constant of the decay bound used at $(\mathrm{L}^{\dagger})$, with
\begin{equation*}
\operatorname{dep}(Y^{\sharp})\ \ge\ \tfrac12\,\check{R}_m
\end{equation*}
by the pointwise form of the kept-member bound in the cell count. In this form cells are used throughout. Either way of measuring serves. The decay bound used at $(\mathrm{L}^{\dagger})$, which supplies the factor $e^{-\frac12\delta_P D_X}$, is stated in $d^{\wedge}$.
\end{enumerate}
\end{remark}

\begin{remark}[The Landau-gauge and axial-gauge representatives of one orbit]
\label{4.22:rem-gauge-representatives}
At the site $(\mathrm{T})$ (the fluctuation integral) and at the site $(\mathrm{P})$ (the large-field
stages) of a step, the response maps of this chapter are in Landau form. These are
the response map $\mathsf{R}_p$ of the fluctuation integral at $(\mathrm{T})$ and, at $(\mathrm{P})$,
\begin{equation}\label{4.22:eq-gauge-representatives-1}
\Phi_{t,\sigma}(b)=\bigl(e^{i\eta\chi_t\mathbb{H}}\,\mathbb{U},\ \dots\bigr).
\end{equation}
Each of them moves the background by its Landau layer and leaves all other variables, in
particular the further arguments $\mathbf{A},\boldsymbol{\Phi}$ of a stored term, unrotated.

The background of Ba{\l}aban's construction is the axial representative of the same gauge orbit.
It is obtained by a gauge transformation $u$. Of $u$ only the averages $\bar{R}_0u$ are normalised,
$\bar{R}_0$ being the averaging operation on gauge transformations of \cite{B-CMP99a};
$u$ is not normalised to $1$ at individual points. Consequently $u\neq 1$ at live law sites in
general, and the following holds.
\begin{itemize}
\item[(a)] For a stored term of type $\mathbf{E}$ or $\mathbf{R}$, gauge invariance removes $u$.
\item[(b)] For a stored term of type $\mathbf{B}$, joint invariance turns $u$ into a rotation of
every further argument.
\end{itemize}
\end{remark}

\begin{proof}
\emph{Landau form.} At $(\mathrm{P})$ the map \eqref{4.22:eq-gauge-representatives-1} acts on the
background pair only, by the factor $e^{i\eta\chi_t\mathbb{H}}$. Here $\chi_t\mathbb{H}$ is the Landau
layer of the stage. The remaining entries of \eqref{4.22:eq-gauge-representatives-1} are not
rotated. The same holds for $\mathsf{R}_p$ at $(\mathrm{T})$, which multiplies the background by the
exponential of its layer. Thus neither map rotates the exterior variables.

\emph{Ba{\l}aban's background is the axial representative.} In \cite{B-CMP102b} the background at
level $k$ is defined as follows: ``we apply to it a gauge transformation $u$ satisfying the
conditions $R_0u=1$ on $\Lambda_j$, and such that the gauge transformed configuration
$(U_1U_0)^u$ satisfies the axial gauge conditions $\mathrm{Ax}_k(\mathfrak{B}_k,U_0)$. Let us
define $U_k=(U_1U_0)^u$''. Here $R_0$ is the averaging operation of \cite{B-CMP102b}, and
$U_1U_0$ is the configuration to which $u$ is applied; it and $U_k=(U_1U_0)^u$ lie in one gauge
orbit. By \cite[Proposition~9]{B-CMP102b} the Landau representative and the axial representative
of that orbit are related by the gauge transformation satisfying the axial conditions and the
normalisation, which is unique by \cite[Theorem~2]{B-CMP99a}. The background is therefore the
gauge transform by $u$ of a configuration in the orbit, and $u$ is determined by the two quoted
conditions.

The same structure carries over to the large-field stages. In the notation of
\cite{B-CMP122a}, whose letters are used only inside this quotation, \cite[(1.35)]{B-CMP122a} reads
\begin{equation*}
\text{``}V^{(j)}_{\dots}=\bigl(M_j(\exp i\xi H)M_j(U^{(n+1)})\bigr)^{u_j}
=\exp iH^{(j)}V^{(j)}\text{''},
\end{equation*}
with the bounds taken from \cite[(106)--(108), (159)--(163)]{B-CMP98}. The gauge transformation
$u_{j+1,\square'}$ is carried in \cite[(1.40)]{B-CMP122a}.

In \cite{B-CMP122b} Ba{\l}aban states: ``we have to keep the gauge transformation $u_k$ \dots\ We
cannot remove it \dots\ $u_k(x)$ is an analytic function of $H_k$ restricted to this block of the
set $\mathbf{B}_k$, to which the point $x$ belongs'', the formula in question being
\cite[(106)]{B-CMP98}. In
every case Ba{\l}aban's background is the axial representative of the orbit of which the Landau
pair is another representative.

\emph{Only means of $u$ are normalised.} By the normalisation convention for gauge transformations
adopted in this chapter, which follows \cite{B-CMP99a}, only the averages $\bar{R}_0u$ of $u$ are
normalised. No pointwise condition $u(x)=1$ is imposed at any site; in particular none is imposed
at the live law sites. Hence $u\neq 1$ at live law sites in general.
As the passage of \cite{B-CMP122b} quoted above says, $u_k$ cannot be removed.

\emph{Consequence for terms of type $\mathbf{E}$ and $\mathbf{R}$.} Such a term depends on its
background through a gauge-invariant expression. Its value at the axial representative therefore
equals its value at the Landau representative, so $u$ drops out. This is (a).

\emph{Consequence for terms of type $\mathbf{B}$.} For a term $c(S;\,\cdot\,;\mathbf{A},
\boldsymbol{\Phi})$ of type $\mathbf{B}$ with domain $S$, invariance holds only jointly. For a
gauge transformation $v$, write $P^v$ for the gauge transform of a background pair $P$, and
$\mathbf{A}^{[v]},\boldsymbol{\Phi}^{[v]}$ for the rotated further arguments. Joint invariance is
the identity
\begin{equation}\label{4.22:eq-gauge-representatives-2}
c\bigl(S;P^{v};\mathbf{A}^{[v]},\boldsymbol{\Phi}^{[v]}\bigr)
=c\bigl(S;P;\mathbf{A},\boldsymbol{\Phi}\bigr).
\end{equation}
Rotations compose as gauge transformations do:
$(\mathbf{A}^{[a]})^{[b]}=\mathbf{A}^{[ba]}$, $(\boldsymbol{\Phi}^{[a]})^{[b]}=\boldsymbol{\Phi}^{[ba]}$
and $(P^a)^b=P^{ba}$.

Let $P$ be the Landau pair and $P^u$ the axial one. Apply \eqref{4.22:eq-gauge-representatives-2}
with $v=u$ to the further arguments $\mathbf{A}^{[u^{-1}]},\boldsymbol{\Phi}^{[u^{-1}]}$. Since
$(\mathbf{A}^{[u^{-1}]})^{[u]}=\mathbf{A}$ and likewise for $\boldsymbol{\Phi}$, this gives
\begin{equation*}
c\bigl(S;P^{u};\mathbf{A},\boldsymbol{\Phi}\bigr)
=c\bigl(S;P;\mathbf{A}^{[u^{-1}]},\boldsymbol{\Phi}^{[u^{-1}]}\bigr).
\end{equation*}
The left side is the term evaluated at Ba{\l}aban's axial background; the right side is the same
term evaluated at the Landau pair, which is the pair the response maps of $(\mathrm{T})$ and $(\mathrm{P})$
produce. On the right the gauge transformation has become a rotation of every further argument.
Since $u\neq 1$ at live law sites, this rotation is not the identity there. This is (b).
\end{proof}

\begin{lemma}[The axialising gauge map at a site]\label{4.22:lem-axialising-map}
Consider one of two sites of a renormalisation step: the fluctuation integral $(\mathrm{T})$, or
stage $l$ of the large-field operations $(\mathrm{P})$.

\emph{Data.}
\begin{itemize}
\item The \emph{base background} $\mathbf{K}$ is the background $\mathbf{U}$ of the output tube
at $(\mathrm{T})$, and a configuration $\mathbb{U}$ in the class $\mathfrak{N}(Y)$ of
configurations of the stage on the region $Y$ at $(\mathrm{P})$.
\item The site's \emph{Landau layer} $\mathcal{K}$ is $\chi_t\mathbb{H}(\sigma;b)$ at
$(\mathrm{P})$, stage $l$: the layer $\mathbb{H}(\sigma;b)$ of the stage multiplied by the
factor $\chi_t$ of the stage, which depends on parameters $t=(t_\square)_\square$ indexed by the
cubes $\square$. At $(\mathrm{T})$ it is $\mathcal{K}_p(\mathsf{t},\sigma)$, the layer
of the response $\mathsf{R}_p(\mathsf{t},\sigma)$, whose first entry is
$e^{i\eta\mathcal{K}_p(\mathsf{t},\sigma)}\mathbf{U}$ and whose second entry is the current
partner of the first.
\item $\mathbf{B}$ is the determining set of the configuration expanded: $\mathbb{B}^{(l-1)}_k$ at
$(\mathrm{P})$, and the determining set $\mathbf{B}$ of the step's input at $(\mathrm{T})$.
\end{itemize}
Put
\begin{equation}\label{4.22:eq-axialising-map-1}
\mathfrak{u}:=\mathfrak{u}\bigl[e^{i\eta\mathcal{K}}\mathbf{K};\mathbf{K};\mathbf{B}\bigr],
\end{equation}
the axialising map of $e^{i\eta\mathcal{K}}\mathbf{K}$ relative to the base $\mathbf{K}$ and the
determining set $\mathbf{B}$, in the notation of the axial-map lemma (the first hypothesis
below). Write
$P:=(e^{i\eta\mathcal{K}}\mathbf{K},\mathbb{J})$, where $\mathbb{J}$ is the current partner of
$e^{i\eta\mathcal{K}}\mathbf{K}$. Put
\begin{equation}\label{4.22:eq-axialising-map-2}
\varrho_0:=\min\bigl\{10^{-2}/K_u,\ c_4(d,L),\ \alpha_1^{98}\bigr\},
\end{equation}
where $\alpha_1^{98}$ is the coupling-free number of \cite[(158)]{B-CMP98}.

\emph{Hypotheses.} Assume the following.
\begin{itemize}
\item \emph{The lemma on the axial gauge map of a complex configuration} (the axial-map lemma).
It applies to a base that is the product of a $G$-valued configuration, regular in the sense of
\cite{B-CMP98}, and
$e^{i\eta A'}$ with $A'$ complex and small, and to a layer $\mathcal{K}$ with
$\Theta_0(\mathcal{K})\le c_4(d,L)$, where $\Theta_0$ is the layer norm of that lemma.
\begin{itemize}
\item[--] Its clauses (i), (ii) and (iv) give the properties of $\mathfrak{u}$ listed in (i) and
(iii) below. In particular the value $\mathfrak{u}(z)$ depends on the layer only through its
restriction to $B^{\mathrm{top}}(z)$, the top block containing $z$ of the lemma's partition into
cubes of side $M=L^{m'}$. It is holomorphic in the layer on the ball
$\{\Theta_0^{\mathrm{loc}}(\mathcal{K})(z)<\varrho_0\}$, and there
$|\log\mathfrak{u}(z)|\le K_u\Theta_0^{\mathrm{loc}}(\mathcal{K})(z)$.
\item[--] Its clause (iii) treats real data. There the gauge transformation carrying the
Landau representative of a critical orbit to its axial representative, subject to the axial
gauge conditions and the normalisation conditions, is the map $\mathfrak{u}$.
\end{itemize}
\item \emph{The base conditions.} At $(\mathrm{P})$, the inequality
$600\,\mathsf{c}_*\,\alpha_{\sup}\le1$ holds, where $\mathsf{c}_*$ is the constant of the base
condition of the stage and $\alpha_{\sup}$ is the supremum of the radii $\alpha$; under it the
class $\mathfrak{N}(Y)$ of configurations of the stage is contained in the class of bases
admitted by the axial-map lemma. At $(\mathrm{T})$, the
layer-wise tube of the step consists
of bases of the form required by the axial-map lemma.
\item \cite[Proposition~9]{B-CMP102b}, identifying the Landau and the axial representatives of
one critical orbit; \cite[Theorem~2]{B-CMP99a}, the uniqueness of the gauge transformation
with the axial and normalisation conditions; and the transformation of the currents,
\cite[(3.11), (1.10)]{B-CMP109}.
\item The joint gauge invariance of the coefficient function $c$ of every term of type
$\mathbf{B}$. For a gauge transformation $a$,
$c(S;P^{a};\mathbf{A}^{[a]},\boldsymbol{\Phi}^{[a]})=c(S;P;\mathbf{A},\boldsymbol{\Phi})$.
Gauge transformations compose as $(V^a)^b=V^{ba}$.
\item The coupling ceiling \eqref{4.22:eq-coupling-ceilings-f}, which gives
$B_H\alpha_*\le\varrho_0$.
\item \emph{The layer bound at stage $l$}: its clause (i), and its decay inequality together with
its side conditions $11\,\hat B\,\mathsf{c}_\square/\mathsf{K}_r\le1$ and $\mathsf{w}\le1$.
Here $\hat B$, $\mathsf{c}_\square$ and $\mathsf{K}_r$ are constants of that lemma, and
$\mathsf{w}$ is the disc width of the site.
\item \emph{The homogeneity bound of the unit shift.} With its constant $K_R$, on the disc of (c)
below, $(1+\eta_p|\mathsf{t}|)K_Rg_k\sup|\mathsf{B}|\le\alpha_*$.
\item \emph{The decay bound of the layer's profile}, resting on \cite[(190)]{B-CMP102b} for the
domains $\{\Omega_j\}$. Here $\Omega$ is the step's fluctuation domain, and $\delta_1$ is the
decay rate of this bound in the scaled distance $d_{\mathrm{sc}}$.
\item \emph{The layer lemma of the correlation theorem}, which introduces the pieces $\delta H_p$,
indexed by pairs $p=(\square,Y_0)$ of a cube $\square$ of the cover and a region $Y_0$, and
their decay factors $\eta_p$. It is used together with \cite[(3.6)--(3.9),
(3.14)]{B-CMP109}.
\item For (a)--(c) below only:
\begin{itemize}
\item[--] the ceiling \eqref{4.22:eq-coupling-ceilings-g};
\item[--] the inequalities $\delta_1\ge\delta_\natural$ and
$\check R_k\ge R_{\min}\ge7$, where $R_{\min}$ is the minimal range number;
\item[--] clause (i), and clause (ii) at $(\mathrm{T})$, of \emph{the lemma on the geometry of
the law sites}; these are geometric statements and do not use the present lemma.
\end{itemize}
\end{itemize}

\emph{Conclusions.} Under these hypotheses the following hold.
\begin{enumerate}
\item[(i)] $\mathfrak{u}$ exists and is unique. It is holomorphic in all the site's parameters
and $G$-valued at real data. It is identically $1$ on top blocks where $\mathcal{K}$ vanishes
identically. The value $\mathfrak{u}(z)$ depends on $(\mathcal{K},\mathbf{K})$ only through
their restrictions to $B^{\mathrm{top}}(z)$, and
\begin{equation}\label{4.22:eq-axialising-map-3}
|\log\mathfrak{u}(z)|\le K_u\,\Theta_0^{\mathrm{loc}}(\mathcal{K})(z).
\end{equation}
$B^{\mathrm{top}}(z)$ lies in the $M$-cube of the local scale of $z$. If $z$ is a law site of a
term with domain $S$, then $B^{\mathrm{top}}(z)$ lies in $S$.
\item[(ii)] (Real point.) Take real data and the full layer: $t=\mathbb{1}$ at $(\mathrm{P})$,
the whole shift at $(\mathrm{T})$, and $\sigma=\mathbb{1}$. Then
$(e^{i\eta\mathcal{K}}\mathbf{K})^{\mathfrak{u}^{-1}}$ is Ba{\l}aban's background
$\mathbf{K}^{\mathrm{pr}}$, jointly with its current, that is,
$P^{\mathfrak{u}^{-1}}=(\mathbf{K}^{\mathrm{pr}},\mathcal{J}(\mathbf{K}^{\mathrm{pr}}))$, where
$\mathcal{J}(\mathbf{K})$ denotes the current determined by a background $\mathbf{K}$. Here
$\mathbf{K}^{\mathrm{pr}}$ is $U^{(l-1)}_k$ of \cite[(1.18), (1.44)]{B-CMP122a} at
$(\mathrm{P})$. At $(\mathrm{T})$ it is the background
$U_{k-1}(\exp(i[\,\dots])V^{(k-1)})$ of \cite[(3.15)]{B-CMP119}, with the bracket as there.
Hence for every term of type $\mathbf{B}$ with domain $S$,
\begin{equation}\label{4.22:eq-axialising-map-4}
c\bigl(S;\mathbf{K}^{\mathrm{pr}},\mathcal{J}(\mathbf{K}^{\mathrm{pr}});
\mathbf{A},\boldsymbol{\Phi}\bigr)
=c\bigl(S;e^{i\eta\mathcal{K}}\mathbf{K},\mathbb{J};
\mathbf{A}^{[\mathfrak{u}]},\boldsymbol{\Phi}^{[\mathfrak{u}]}\bigr).
\end{equation}
\item[(iii)] For a $G$-valued gauge transformation $a$,
$\mathfrak{u}\bigl[(e^{i\eta\mathcal{K}}\mathbf{K})^a;\mathbf{K}^a;\mathbf{B}\bigr]
=a\,\mathfrak{u}\,a^{-1}$.
\item[(iv)] (Sizes.) At $(\mathrm{P})$, on the whole family $\mathfrak{W}^{(l),\natural}$ of
stage $l$ used in the proof of the layer bound at stage $l$, with parameters
$|t_\square|\le R^\natural_\square$,
\begin{equation}\label{4.22:eq-axialising-map-5}
\operatorname{pol}\bigl(\mathfrak{u}^{(l)}_{t,\sigma}(b)(z)\bigr)
\le K_u\sup_{y\in B^{\mathrm{top}}(z)}\varkappa^{(l)}_{\mathsf{w}}(y)
\le K_u\alpha^{+}\rho^{(l)}_{m}\,e^{-\frac14\delta\, d^{\wedge}(\tilde z,\Gamma^{(l)})}.
\end{equation}
Here $\alpha^{+}$ is the radius $\alpha^{+}_k$ read from the coupling, $\rho^{(l)}_m$ is the
amplitude of the layer bound at stage $l$, the index $m$ being that lemma's scale index (not the
torus exponent), $\Gamma^{(l)}$ is the contour of stage $l$,
$\tilde z$ is the point attached to $z$ in the weighted contour distance $d^{\wedge}$, and
$\delta$ is the decay rate of the layer bound at stage $l$.
At $(\mathrm{T})$, at $\mathsf{t}=0$ and for $|\sigma|\le e^{\kappa_1}$,
\begin{equation}\label{4.22:eq-axialising-map-6}
\operatorname{pol}\bigl(\mathfrak{u}_p(0,\sigma)(z)\bigr)
\le K_uB_H\alpha_*\,e^{-\delta_1(d_{\mathrm{sc}}(z,\Omega)-1)}.
\end{equation}
Moreover, at $(\mathrm{T})$, let $p=(\square,Y_0)$ and let $\delta H_p$ be the unit piece indexed
by $p$ of the layer in the layer lemma of the correlation theorem. The layer is then
$\mathcal{K}_p(\mathsf{t},\sigma)=\mathcal{K}_p(0,\sigma)+\mathsf{t}\,\delta H_p$, which is
linear in $\mathsf{t}$, as in the passage from \cite[(3.6)]{B-CMP109} to
\cite[(3.7)]{B-CMP109}, where $A$ is replaced by
$H_j(B(t))+t_{\square}\langle\,\cdots\rangle$. Let $\mathsf{U}$, the unit shift of the
homogeneity bound, and $e_p$, the profile of the piece, be defined by the first line of
\eqref{4.22:eq-axialising-map-a}: $\mathsf{U}$ is $B_HK_Rg_k\sup|\mathsf{B}|$ taken at
$\sup|\mathsf{B}|=T^{\mathbb{C}}$, one normalisation for (a) and (b), and the supremum defining
$e_p$ is over the parameter manifold
$\mathfrak{M}\ni(\sigma,\mathbf{U},\mathbf{J},\mathsf{B},w)$, where $w$ is the parameter of that
name on which the pieces depend, so that $\vartheta_p(s)$
below does not depend on the point of $\mathfrak{M}$. Let $v$ be a kept member, with domain
$S_v$. Then statements (a), (b) (for every $y$) and (c) of
\eqref{4.22:eq-axialising-map-a} hold, (c) on the closed disc
$|\mathsf{t}|\le\mathsf{w}/\eta_p$, $|\sigma|\le e^{\kappa_1}$ and at every law site $s$ of
$v$:
\begin{equation}\label{4.22:eq-axialising-map-a}
\begin{gathered}
\mathsf{U}:=B_HK_Rg_k\,T^{\mathbb{C}},\qquad
e_p(y):=\frac{1}{\mathsf{U}}\,\sup_{\mathfrak{M}}\Theta_0^{\mathrm{loc}}(\delta H_p)(y),\\
\text{(a)}\quad e_p\le\eta_p\ \text{ on } S_v,\qquad
\text{(b)}\quad e_p(y)\le C_b\,e^{-\delta_1(d_{\mathrm{sc}}(y,\Omega)-1)},\\
\text{(c)}\quad
\bigl|\log[\mathfrak{u}_p(0,\sigma)^{-1}\mathfrak{u}_p(\mathsf{t},\sigma)](s)\bigr|
\le\vartheta_p(s):=\vartheta^{T}\,\frac{e_p(s)}{\eta_p}\le\vartheta^{T},\\
\vartheta^{T}:=2.1\,K_uB_H\alpha_*.
\end{gathered}
\end{equation}
In (b), $C_b\ge1$ is the constant of the composition of the decay bound with itself. The sum
over $x\in\square$ is estimated pointwise, so that no factor of the diameter of $\square$ is
lost. The constant $B_H$, the one entering $\varepsilon_{\mathfrak{Q}}$, is not altered.
\end{enumerate}
\end{lemma}

\begin{proof}
\emph{Clauses (i) and (iii).} These are clauses (i), (ii) and (iv) of the axial-map lemma,
applied to the base $\mathbf{K}$, the layer $\mathcal{K}$ and the determining set $\mathbf{B}$.
We check its hypotheses.
\begin{itemize}
\item The base must be the product of a $G$-valued regular configuration and $e^{i\eta A'}$,
with $A'$ complex and small. At $(\mathrm{P})$ this holds because
$600\,\mathsf{c}_*\alpha_{\sup}\le1$ is assumed and, under it, $\mathfrak{N}(Y)$ is contained
in the class of bases admitted by the axial-map lemma.
At $(\mathrm{T})$ it holds because the
output tube is the layer-wise tube of the step.
\item The layer must satisfy $\Theta_0(\mathcal{K})\le c_4(d,L)$. This follows from the bounds
on the layer established in the proof of (iv) below and, on the disc of (c), from
\eqref{4.22:eq-axialising-map-13}, together with \eqref{4.22:eq-coupling-ceilings-f} and
$\varrho_0\le c_4(d,L)$. Those bounds concern $\mathcal{K}$ itself, not $\mathfrak{u}$.
\end{itemize}
The blocks of the axial-map lemma are cells of the partitions into cubes of side $M=L^{m'}$.
These partitions are compatible with all the other partitions of the step
\cite[\S3]{B-CMP119}. Hence
the top block $B^{\mathrm{top}}(z)$ lies in the $M$-cube of the local scale of $z$. If $z$ is a law
site of a term with domain $S$, it lies in $S$.

\emph{Clause (ii).} This is clause (iii) of the axial-map lemma, combined with the published
inputs as follows.
\begin{itemize}
\item At real data and full layer, $e^{i\eta\mathcal{H}}U(\mathbf{B},V_0)$ and
$U(\mathbf{B},V'V_0)$, in the notation of \cite[Proposition~9]{B-CMP102b}, are the Landau and
the axial representative of one critical orbit, by that proposition.
\item They are related by the gauge transformation satisfying the axial gauge conditions and the
normalisation conditions. This transformation is unique by \cite[Theorem~2]{B-CMP99a}, and by
clause (iii) of the axial-map lemma it is $\mathfrak{u}$.
\item At $(\mathrm{P})$ the configuration $U^{(l-1)}_k$ of \cite[(1.44)]{B-CMP122a} has the form
of \cite[(1.30), (1.34)--(1.36)]{B-CMP122a}; it is in this form that $\mathbf{K}^{\mathrm{pr}}$
at $(\mathrm{P})$ is read in the first two items.
\item The currents transform accordingly by \cite[(3.11), (1.10)]{B-CMP109}.
\end{itemize}
This gives
$P^{\mathfrak{u}^{-1}}=(\mathbf{K}^{\mathrm{pr}},\mathcal{J}(\mathbf{K}^{\mathrm{pr}}))$.

Now apply the joint invariance of $c$ to the pair $P^{\mathfrak{u}^{-1}}$ and the gauge
transformation $\mathfrak{u}$. This gives
\begin{equation}\label{4.22:eq-axialising-map-7}
c\bigl(S;P^{\mathfrak{u}^{-1}};\mathbf{A},\boldsymbol{\Phi}\bigr)
=c\bigl(S;(P^{\mathfrak{u}^{-1}})^{\mathfrak{u}};
\mathbf{A}^{[\mathfrak{u}]},\boldsymbol{\Phi}^{[\mathfrak{u}]}\bigr).
\end{equation}
By $(V^a)^b=V^{ba}$ we have
$(P^{\mathfrak{u}^{-1}})^{\mathfrak{u}}=P^{\mathfrak{u}\mathfrak{u}^{-1}}=P$. Inserting
$P^{\mathfrak{u}^{-1}}=(\mathbf{K}^{\mathrm{pr}},\mathcal{J}(\mathbf{K}^{\mathrm{pr}}))$ and
$P=(e^{i\eta\mathcal{K}}\mathbf{K},\mathbb{J})$ into \eqref{4.22:eq-axialising-map-7} gives
\eqref{4.22:eq-axialising-map-4}.

\emph{A remark used in (iv).} If $g=e^{Y}$, then $\|g\|\le e^{\|Y\|}$ and
$\|g^{-1}\|\le e^{\|Y\|}$. Hence $\operatorname{pol}(g)\le|\log g|$, and
\eqref{4.22:eq-axialising-map-3} bounds $\operatorname{pol}(\mathfrak{u}(z))$ by
$K_u\Theta_0^{\mathrm{loc}}(\mathcal{K})(z)$.

\emph{Clause (iv) at $(\mathrm{P})$.}
\begin{itemize}
\item The local norm $N_y$ at $y$ of the layer bound at stage $l$ carries the factor
$L^m\eta|\cdot|$, where $L^m$ is the scale factor of that norm; here $m$ is a scale exponent of
stage $l$, not the torus exponent. Since
$\Theta_0^{\mathrm{loc}}(\mathcal{K})(z)$ depends on the layer only through its restriction to
$B^{\mathrm{top}}(z)$,
\[
\Theta_0^{\mathrm{loc}}(\mathcal{K})(z)\le\sup_{y\in B^{\mathrm{top}}(z)}N_y(\mathcal{K}).
\]
\item By clause (i) of the layer bound at stage $l$, $N_y(\chi_t\mathbb{H})\le
\varkappa_{\mathsf{w}}(y)$ at every $y$ and for all $|t_\square|\le R^\natural_\square$. Hence
this holds on the whole family $\mathfrak{W}^{(l),\natural}$.
\item Together with the remark above this gives the first inequality of
\eqref{4.22:eq-axialising-map-5}.
\item The second inequality is the decay inequality of the layer bound at stage $l$. It is used
with its side conditions $11\hat B\mathsf{c}_\square/\mathsf{K}_r\le1$ and $\mathsf{w}\le1$.
The decay rate is
split as $\delta=\frac12\delta+\frac14\delta+\frac14\delta$: the first half is spent on the
amplitude, one quarter on the absolute summation, and the last quarter is kept. The kept quarter
is the factor $e^{-\frac14\delta d^{\wedge}(\tilde z,\Gamma^{(l)})}$.
\end{itemize}

\emph{Clause (iv) at $(\mathrm{T})$, $\mathsf{t}=0$.}
\begin{itemize}
\item By the homogeneity bound, $(1+\eta_p|\mathsf{t}|)K_Rg_k\sup|\mathsf{B}|\le\alpha_*$. This
bounds the argument of the layer.
\item The decay bound of the layer's profile bounds its profile.
\item At $\mathsf{t}=0$ the layer is the undilated remainder, the part carried by
$\zeta'_\square$, and its source lies in $\Omega$. Here $\zeta_\square$ is the partition-of-unity
function of the cube $\square$ of the cover, and $\zeta'_\square$ its companion function, in the
localisation of \cite[(3.7)]{B-CMP109}.
\item Together these give, at $\mathsf{t}=0$ and $|\sigma|\le e^{\kappa_1}$,
\begin{equation}\label{4.22:eq-axialising-map-8}
\Theta_0^{\mathrm{loc}}(\mathcal{K}_p(0,\sigma))(z)
\le\mathsf{U}\,e^{-\delta_1(d_{\mathrm{sc}}(z,\Omega)-1)}
\le B_H\alpha_*\,e^{-\delta_1(d_{\mathrm{sc}}(z,\Omega)-1)}.
\end{equation}
The second inequality is the homogeneity bound at $\mathsf{t}=0$, multiplied by $B_H$.
\item By the remark above, \eqref{4.22:eq-axialising-map-8} gives
\eqref{4.22:eq-axialising-map-6}.
\end{itemize}

\emph{Bound (a).} By its definition in the layer lemma of the correlation theorem, $\eta_p$ is
the decay factor of the piece $\delta H_p$ between $\square$ and $S_v$. It is taken in the same
normalisation by $\mathsf{U}$ as $e_p$, so $e_p\le\eta_p$ on $S_v$.

\emph{Bound (b).} By \cite[(3.7)--(3.8)]{B-CMP109},
\begin{equation}\label{4.22:eq-axialising-map-9}
\delta H_p=\Bigl\langle\frac{\delta}{\delta B}H_j,\;Q\bigl(\zeta_\square H_k(B')\bigr)\Bigr\rangle,
\end{equation}
where the pairing is a sum over $x\in\square$. We estimate the two factors.
\begin{itemize}
\item By \cite[\S3]{B-CMP109}, which refers for this to \cite[Proposition~9, (190)]{B-CMP102b},
the
derivative $(\delta/\delta B)H$ decays exponentially, with decay rate $\delta_0$ on the scale of
the lattice spacing. In the scaled distance this is a factor
$e^{-\delta_1 d_{\mathrm{sc}}(y,\square)}$.
\item The datum $\zeta_\square H_k(B')$ has $B'$ living on $\Omega$. By the decay bound of the
layer's profile, in the normalisation of $\mathsf{U}$, it carries the factor
$e^{-\delta_1 d_{\mathrm{sc}}(\square,\Omega)}$. This is the same decay bound as was used at
$\mathsf{t}=0$.
\end{itemize}
The sum over $x\in\square$ is estimated pointwise: both bounds are applied at the summation point
$x\in\square$, where they give the factors $e^{-\delta_1 d_{\mathrm{sc}}(y,x)}$ and
$e^{-\delta_1 d_{\mathrm{sc}}(x,\Omega)}$. The triangle inequality gives
$d_{\mathrm{sc}}(y,\Omega)\le d_{\mathrm{sc}}(y,x)+d_{\mathrm{sc}}(x,\Omega)$ for each
$x\in\square$. So every summand carries $e^{-\delta_1 d_{\mathrm{sc}}(y,\Omega)}$, and no factor
of the diameter of $\square$ enters. The sum is absorbed by the constant $C_b$ of the
composition of the decay bound with itself.

The norm $\Theta_0^{\mathrm{loc}}(\cdot)(y)$ depends on its argument only through its restriction
to $B^{\mathrm{top}}(y)$. Passing
from $y$ to this block costs the $-1$ in the exponent. The bound is uniform on
$|\sigma|\le e^{\kappa_1}$, since the new localised functions satisfy the bounds
\cite[(3.14), (3.9)]{B-CMP109}. This proves (b).

For a term of type $\mathbf{B}$ whose piece is the direct $\zeta_\square\mathbb{H}$, of the type
of \cite[(1.60)]{B-CMP122b}, (b) is the decay bound of the layer's profile at $y$ itself.

\emph{Bound (c).} We use the Schwarz lemma on the top block.

\emph{Setting up.} Fix $(\sigma,\mathbf{U},\mathbf{J},\mathsf{B},w)\in\mathfrak{M}$ and a law
site $s$ of $v$. Put $\mathcal{K}_0:=\mathcal{K}_p(0,\sigma)$,
$a_0:=\Theta_0^{\mathrm{loc}}(\mathcal{K}_0)(s)$ and
$\hat a:=\Theta_0^{\mathrm{loc}}(\delta H_p)(s)$. By the definition of $e_p$ as a supremum over
$\mathfrak{M}$, $\hat a\le\mathsf{U}e_p(s)$. Let $\varrho_0$ be as in
\eqref{4.22:eq-axialising-map-2}.

By clauses (i) and (ii) of the axial-map lemma, $\mathfrak{u}(s)$ depends on the layer only
through its restriction to $B^{\mathrm{top}}(s)$. It is holomorphic in the layer on the ball
$\{\Theta_0^{\mathrm{loc}}(\cdot)(s)<\varrho_0\}$, with
$|\log\mathfrak{u}(s)|\le K_u\Theta_0^{\mathrm{loc}}(\text{layer})(s)$. For complex $\lambda$ put
\begin{equation}\label{4.22:eq-axialising-map-10}
\Psi(\lambda):=\log\bigl[\mathfrak{u}[\mathcal{K}_0]^{-1}\,
\mathfrak{u}[\mathcal{K}_0+\lambda\,\delta H_p]\bigr](s).
\end{equation}
By the triangle inequality for $\Theta_0^{\mathrm{loc}}(\cdot)(s)$,
$\Theta_0^{\mathrm{loc}}(\mathcal{K}_0+\lambda\delta H_p)(s)\le a_0+|\lambda|\hat a$.

\emph{Holomorphy and the a priori bound.} On $|\lambda|\hat a<\varrho_0-a_0$ both factors in
\eqref{4.22:eq-axialising-map-10} are holomorphic, and $\Psi(0)=0$. Their logarithms have norms
at most $K_ua_0$ and $K_u(a_0+|\lambda|\hat a)$. Both are at most
$K_u\varrho_0\le10^{-2}$. The Baker--Campbell--Hausdorff series for two matrices of norm at most
$10^{-2}$ bounds the logarithm of the product by $1.05$ times the sum of the two norms. Since
\[
K_ua_0+K_u(a_0+|\lambda|\hat a)<K_u(\varrho_0+a_0),
\]
we obtain $|\Psi|\le1.05K_u(\varrho_0+a_0)$ there.

\emph{Schwarz.} By the Schwarz lemma on the disc $|\lambda|<(\varrho_0-a_0)/\hat a$,
\begin{equation}\label{4.22:eq-axialising-map-11}
|\Psi(\lambda)|\le1.05K_u\,\frac{\varrho_0+a_0}{\varrho_0-a_0}\,|\lambda|\hat a
\le2.1K_u|\lambda|\hat a.
\end{equation}
If $\hat a=0$, $\Psi$ is entire and bounded with $\Psi(0)=0$, hence zero, and
\eqref{4.22:eq-axialising-map-11} holds trivially.

\emph{The second inequality of \eqref{4.22:eq-axialising-map-11}.} It uses $a_0\le\varrho_0/3$,
which holds for three reasons.
\begin{itemize}
\item By clause (ii) at $(\mathrm{T})$ of the lemma on the geometry of the law sites,
$a_0\le B_H\alpha_*e^{-\delta_1(8\check R_k-1)}$.
\item By \eqref{4.22:eq-coupling-ceilings-f}, $B_H\alpha_*\le\varrho_0$.
\item $e^{-\delta_1(8\check R_k-1)}<\frac13$. Indeed, \eqref{4.22:eq-coupling-ceilings-g} gives
$5\delta_\natural R_{\min}\ge2$. Since $\check R_k\ge1$ we have $8\check R_k-1\ge7\check R_k$.
With $\delta_1\ge\delta_\natural$ and $\check R_k\ge R_{\min}$ this gives
\[
\delta_1(8\check R_k-1)\ge7\delta_\natural\check R_k\ge7\delta_\natural R_{\min}
\ge\tfrac{14}{5}>\log3.
\]
\end{itemize}
Since $(\varrho_0+a_0)/(\varrho_0-a_0)$ increases with $a_0/\varrho_0$,
$a_0\le\varrho_0/3$ gives
$(\varrho_0+a_0)/(\varrho_0-a_0)\le(1+\frac13)/(1-\frac13)=2$. Moreover $1.05\cdot2=2.1$.

\emph{The value $\lambda=\mathsf{t}$ is in range.} Let $\mathsf{t}$ lie in the disc
$|\mathsf{t}|\le\mathsf{w}/\eta_p$. The homogeneity bound, multiplied by $B_H$, gives
$(1+\eta_p|\mathsf{t}|)\mathsf{U}\le B_H\alpha_*$. In particular
$\eta_p|\mathsf{t}|\,\mathsf{U}\le B_H\alpha_*$, so
\begin{equation}\label{4.22:eq-axialising-map-12}
|\mathsf{t}|\hat a\le|\mathsf{t}|\,\mathsf{U}\,e_p(s)\le B_H\alpha_*\,\frac{e_p(s)}{\eta_p}.
\end{equation}
By the decay bound, as in \eqref{4.22:eq-axialising-map-8},
$a_0\le\mathsf{U}e^{-\delta_1(d_{\mathrm{sc}}(s,\Omega)-1)}\le\mathsf{U}$. By (a),
$e_p(s)\le\eta_p$. Hence
\begin{equation}\label{4.22:eq-axialising-map-13}
a_0+|\mathsf{t}|\hat a\le\mathsf{U}+|\mathsf{t}|\,\mathsf{U}\eta_p
=\mathsf{U}(1+\eta_p|\mathsf{t}|)\le B_H\alpha_*\le\varrho_0,
\end{equation}
the last step by \eqref{4.22:eq-coupling-ceilings-f}. So $\lambda=\mathsf{t}$ is in the range of
\eqref{4.22:eq-axialising-map-11}.

\emph{Conclusion.} By the linearity of the layer in $\mathsf{t}$,
$\mathcal{K}_0+\mathsf{t}\,\delta H_p=\mathcal{K}_p(\mathsf{t},\sigma)$, so
$\Psi(\mathsf{t})=\log[\mathfrak{u}_p(0,\sigma)^{-1}\mathfrak{u}_p(\mathsf{t},\sigma)](s)$.
By \eqref{4.22:eq-axialising-map-11} and \eqref{4.22:eq-axialising-map-12},
\[
|\Psi(\mathsf{t})|\le2.1K_uB_H\alpha_*\,\frac{e_p(s)}{\eta_p}=\vartheta_p(s).
\]
Finally $\vartheta_p(s)\le\vartheta^{T}$ by (a). The bound is uniform in the fixed point of
$\mathfrak{M}$, which proves (c).
\end{proof}

\begin{remark}
The dilation statement for the pieces, that on the disc $|\mathsf{t}|\le\mathsf{w}/\eta_p$ the
piece seen on $S_v$ has size at most the unit shift $\mathsf{U}$, stands as a statement about the
near end of $S_v$: there $e_p$ is comparable to $\eta_p$, and no decay is left. Bound (b) of
Lemma~\ref{4.22:lem-axialising-map} is what that statement leaves at depth.
\end{remark}

\begin{definition}[Reading of block-type terms at fluctuation and preparatory sites]
\label{4.22:def-block-term-reading}
Let the site be $(\mathrm{T})$, the fluctuation integral of a step, or $(\mathrm{P})$, a
large-field stage. Let $\mathbf{K}$ be the base background of the site, $\mathcal{K}$ its
Landau layer and $\mathbf{B}$ the determining set of the configuration expanded at the site,
and let
\begin{equation*}
\mathfrak{u}=\mathfrak{u}\bigl[e^{i\eta\mathcal{K}}\mathbf{K};\mathbf{K};\mathbf{B}\bigr]
\end{equation*}
be the axialising gauge map of the site, all as in Lemma~\ref{4.22:lem-axialising-map}. The
\emph{Landau pair} of the site is the background $e^{i\eta\mathcal{K}}\mathbf{K}$ together with
its partner current, which we denote $\mathcal{J}$. At $(\mathrm{T})$ and at $(\mathrm{P})$,
every term of type $\mathbf{B}$ (Ba{\l}aban's term species, not the determining set $\mathbf{B}$
above) with domain $S$, pointwise coordinates $\mathbf{A}$ and frames $\boldsymbol{\Phi}$ is read
as the orbit datum
\begin{equation*}
c^{\sharp}\bigl(S;\,e^{i\eta\mathcal{K}}\mathbf{K},\mathcal{J};\,\mathbf{A},\boldsymbol{\Phi},
\mathfrak{u}\cdot g\bigr),
\end{equation*}
that is, evaluated at the Landau pair, with $\mathfrak{u}\cdot g$ in the rotation argument
in place of $g$; here $g$ denotes the rotation argument of the orbit datum, a complex rotation
at the law sites of $S$, and not the reference coupling.
\end{definition}

\begin{corollary}[One preparatory stage on the orbit chart]\label{4.22:cor-preparatory-chart}
Assume the following statements, each of which is used by its statement only:
\begin{itemize}
\item the one-stage corollary of the stage expansion at real data, the stage expansion
identity, and clause (iv) of the stage expansion lemma;
\item clauses (iii) and (iv) of the first stage lemma, and the second stage lemma;
\item clauses (i)--(iii) of Lemma~\ref{4.22:lem-axialising-map} (the axialising gauge map
at a site);
\item the complex-parameter stage lemma;
\item properties (O1), (O2), (O4) and (O5) of orbit data;
\item clause (a) of the domain hypothesis at stage $l$.
\end{itemize}
Fix a step $k$ and a stage $l$ of its large-field part, at which the level $j=\ell_l$ is
integrated. Let $\mathfrak{u}^{(l)}$ be the axialising gauge map of
Lemma~\ref{4.22:lem-axialising-map} at stage $l$, and for a law site $s$ put
$\eta^{(l)}(s):=\sup\operatorname{pol}(\mathfrak{u}^{(l)}(s))$. Let every constituent $c$ of
the plain class have an orbit datum $c^{\sharp}$ of radius $\varepsilon$, where
$\varepsilon(s)>\eta^{(l)}(s)$ at every law site $s$. For such a $c$ with domain $S$ write
$\mathcal{A}_c$ and $\mathcal{Q}^w_c$ for the two index sets of its stage expansion, as fixed
in the stage expansion identity; the first carries the cube parameters $t_\square$, the second
the dilation variables $u_q$ of the level-$j$ coordinates. In the stage
expansion identity and in the stage expansion lemma make, in accordance with the reading of
Definition~\ref{4.22:def-block-term-reading}, the replacement
\begin{equation}\label{4.22:eq-preparatory-chart-1}
c\bigl(S;\Phi_{t,\sigma}(b),A_j^u,\lambda\bigr)\ \longmapsto\
c^{\sharp}\bigl(S;\Phi_{t,\sigma}(b);(A_j^u,\lambda),\boldsymbol{\Phi},
\mathfrak{u}^{(l)}_{t,\sigma}(b)\cdot g\bigr),
\end{equation}
where $g\in\mathfrak{G}_{\varepsilon-\eta^{(l)}}(S)$, and $g:=1$ at the law sites of the
integrated rim $\mathsf{R}_j$. Then:
\begin{enumerate}
\item[(1)] at real data and $G$-valued $g$, the left member of the stage expansion identity is
the bracket of \cite[(1.60)]{B-CMP122a} at the exterior variables
$(\lambda^{[g]},\boldsymbol{\Phi}^{[g]})$;
\item[(2)] every statement of the real-data analysis of the stage on which the one-stage
corollary at real data rests, together with clauses (iii), (iv) of the first stage lemma and
the second stage lemma, holds under the dictionary
$\lambda\mapsto(\lambda,\boldsymbol{\Phi},g)$,
$\mathfrak{N}\mapsto\mathfrak{N}^{\mathbb{C}}$, with the same index sets
$\mathcal{A}_c$, $\mathcal{Q}^w_c$ and the same supports $\mathsf{b}(e)$, $\mathsf{d}(e)$;
\item[(3)] the outputs $\mathbb{C}^{(l)}_k(X,\mathcal{S})^{\sharp}$ are orbit data of radius
$\varepsilon-\eta^{(l)}$ on the surviving law sites.
\end{enumerate}
\end{corollary}

\begin{proof}
\textit{Which arguments are read.} By clause (iii) of the first stage lemma, read against the
bracket of \cite[(1.60)]{B-CMP122a}, the background map $\Phi_{t,\sigma}(b)$, the Landau layer
$\mathbb{H}$, the factor $\mathcal{J}^{\sharp}[1]$ of that bracket, the functions
$\zeta_\square$, the Gaussian integral, the product $\chi^{(j)\sharp}\chi'^{(j)}$, the factor
$\mathbb{P}^{(j)}$, and the three remaining factors of that bracket read neither the slot
arguments of the constituents nor
$g$. By Lemma~\ref{4.22:lem-axialising-map}(i), the map $\mathfrak{u}^{(l)}$ reads
$(t,\sigma,b,\mathbb{U})$ on $S$, and does not read the slot arguments. Consequently, in
\eqref{4.22:eq-preparatory-chart-1} the rotation enters only through the last argument of
$c^{\sharp}$, and every other factor of the stage is the same as at real data.

\textit{(i) The identity.} For $A'\subset\mathcal{A}_c$ and $Q'\subset\mathcal{Q}^w_c$ let
$\Phi_{\mathbb{1}_{A'}}$ and $\mathfrak{u}_{\mathbb{1}_{A'}}$ denote $\Phi_{t,\sigma}(b)$ and
$\mathfrak{u}^{(l)}_{t,\sigma}(b)$ at $t=\mathbb{1}_{A'}$, the indicator function of $A'$, and
let $A_j^{\mathbb{1}_{Q'}}$ denote the level-$j$ coordinates at the dilation
$u=\mathbb{1}_{Q'}$. Put
\begin{equation}\label{4.22:eq-preparatory-chart-2}
f(A',Q'):=c^{\sharp}\bigl(S;\Phi_{\mathbb{1}_{A'}};(A_j^{\mathbb{1}_{Q'}},\lambda),
\boldsymbol{\Phi},\mathfrak{u}_{\mathbb{1}_{A'}}g\bigr).
\end{equation}
The stage expansion identity is the M\"obius inversion formula on the Boolean lattice
$2^{\mathcal{A}_c}\times2^{\mathcal{Q}^w_c}$ applied to this $f$. M\"obius inversion is a
finite combinatorial identity valid for every
function on $2^{\mathcal{A}_c}\times2^{\mathcal{Q}^w_c}$; it therefore holds at every value of
the parameters $(\lambda,\boldsymbol{\Phi},g)$, complex ones included. The left member of the
identity is $f(\mathcal{A}_c,\mathcal{Q}^w_c)$. At real data and $G$-valued $g$, this is the
constituent $c(S;U^{(l-1)}_k;A_j,\lambda^{[g]},\boldsymbol{\Phi}^{[g]})$ of
\cite[(1.60)]{B-CMP122a}, by Lemma~\ref{4.22:lem-axialising-map}(ii) and properties (O2),
(O4) of orbit data; property (O4) applies because the restriction $\Phi\restriction S$ and the
restriction $\mathfrak{u}\restriction\mathfrak{s}(S)$ read $t_\square$ only for
$\square\in\mathcal{A}_c$. This is (1). At $A'=\emptyset$ the Landau layer is switched off and
$\mathfrak{u}=1$. Hence the coefficient $c_{\emptyset,\emptyset}$, the term $c_{\mathrm{str}}$
and the subtracted term $c(S;\Phi_0,0)$ of the stage expansion identity are those of the
one-stage corollary at real data.

\textit{(ii) The bounds.} Put
\begin{equation}\label{4.22:eq-preparatory-chart-3}
F(t,u):=c^{\sharp}\bigl(S;\Phi_{t,\sigma}(b)\restriction S;(A_j^u,\lambda),
\boldsymbol{\Phi},\mathfrak{u}_{t,\sigma}(b)g\bigr).
\end{equation}
By property (O1) of orbit data, $F$ is holomorphic on
\begin{equation*}
\prod_{\square\in\mathcal{A}_c}\{|t_\square|\le R^{\natural}_\square\}\times
\prod_{q\in\mathcal{Q}^w_c}\{|u_q|\le r_A\}\times
\mathfrak{G}_{\varepsilon-\eta^{(l)}}(S).
\end{equation*}
The hypotheses of (O1) are met for three reasons. First, the level-$j$ coordinates dilated by
$u$ are weak coordinates. Second, clause (a) of the domain hypothesis at stage $l$ puts
$\Phi\restriction S$ in the domain of $c$. Third, $\operatorname{pol}(\mathfrak{u}g)<\varepsilon$
at every law site $s$ of $S$. Indeed, for $y_1,y_2\in G^{\mathbb{C}}$ the inequalities
$\|y_1y_2\|\le\|y_1\|\,\|y_2\|$ and
$\|(y_1y_2)^{-1}\|\le\|y_2^{-1}\|\,\|y_1^{-1}\|$ give $\operatorname{pol}(y_1y_2)\le
\operatorname{pol}(y_1)+\operatorname{pol}(y_2)$, whence
\begin{equation*}
\operatorname{pol}\bigl(\mathfrak{u}(s)g(s)\bigr)\le\eta^{(l)}(s)+\operatorname{pol}(g(s))
<\eta^{(l)}(s)+\varepsilon(s)-\eta^{(l)}(s)=\varepsilon(s).
\end{equation*}
For the same reason $|F|\le\mathfrak{N}^{\mathbb{C}}_{\varepsilon}(c)$ on that domain. The
Cauchy estimates of clause (iv) of the first stage lemma, applied to $F$ with this bound, give,
with the constant $K_A$ and the cube profile $\vartheta^{\natural}(\square)$ of those Cauchy
estimates,
\begin{equation}\label{4.22:eq-preparatory-chart-4}
|c_{A,Q}|\le\mathfrak{N}^{\mathbb{C}}_{\varepsilon}(c)\prod_{\square\in A}K_A
\vartheta^{\natural}(\square)\prod_{q\in Q}\frac{2}{r_A},\qquad
|c_{\mathrm{str}}|\le2\,\mathfrak{N}^{\mathbb{C}}_{\varepsilon}(c).
\end{equation}
No loss factor enters. The dilation of the family keeps the decay factor
$e^{-\frac14\delta d^{\wedge}}$ in the weighted contour distance $d^{\wedge}$, so that
$\eta^{(l)}$ is already the whole debit of the stage.

\textit{(iii) Clause (iv) of the stage expansion lemma.} The substitution $b\mapsto0$ made
there kills the Landau layer $\mathbb{H}$. Hence $\mathfrak{u}=1$, and $F$ does not depend
on $t$.

\textit{(iv) The expansion.} Apply the complex-parameter stage lemma with parameter manifolds
\begin{equation*}
\mathfrak{M}_e:=\mathfrak{G}_{\varepsilon-\eta^{(l)}}(S(e))
\end{equation*}
for the elements $e$ of the expansion, $S(e)$ denoting the domain of the element $e$.

\textit{(v) Property (O2) for the output.} For $G$-valued $g$, each output term $v_e^{\sharp}$
is $v_e$ at the exterior variables $(\lambda^{[g]},\boldsymbol{\Phi}^{[g]})$. This follows term
by term from (i), applied to each constituent entering $v_e$.

\textit{(vi) Invariance.} For a gauge transformation $v$,
\begin{equation*}
c^{\sharp}\bigl(S;\Phi^{v};\cdot,(v\mathfrak{u}v^{-1})(vg)\bigr)
=c^{\sharp}\bigl(S;\Phi;\cdot,\mathfrak{u}g\bigr),
\end{equation*}
by Lemma~\ref{4.22:lem-axialising-map}(iii) and property (O5) of orbit data. The assignment
of the outputs to their domains, the strong sets and the weights do not involve the
rotation $g$.

Step (i) gives (1); steps (i)--(iv) give (2); steps (ii) and (iv)--(vi) show that the outputs
are holomorphic and bounded on the chart of radius $\varepsilon-\eta^{(l)}$, reduce to the
real outputs at $G$-valued $g$ and are gauge invariant, which is (3).
\end{proof}

\begin{lemma}[Transfer of majorant bounds to the orbit chart]
\label{4.22:lem-majorant-transfer}
Fix a step and one of its sites $(\mathrm{T})$, $(\mathrm{P})$, $(\mathrm{L})$. Let $v$ be a stored
term with domain $X_v$, stored domain
$\mathfrak{U}_v\times\mathfrak{D}_{\mathcal{S}}\times\operatorname{supp}\boldsymbol{\Phi}$ and orbit
datum $f_v^{\sharp}$ of radius $\varepsilon=(\varepsilon(s))_s$, a function on the law sites $s$ of
$v$. Write $\mathfrak{N}^{\mathbb{C}}_{\varepsilon}(v)$ for the
supremum of $|f_v^{\sharp}|$ over its chart. Let $v$ be expanded along a family
$\tau\mapsto\mathsf{R}(\tau)$ of configurations by an expansion with majorant bounds that
satisfies the following four conditions.
\begin{itemize}
\item[(M1)] (Identity.) Let $\mathfrak{P}=\mathfrak{P}^{\mathrm{base}}\times\mathfrak{P}^{\mathrm{inc}}$
be the site's parameter polydisc. For $\tau\in\mathfrak{P}$ put
$f(\tau):=f_v(\mathsf{R}(\tau)\restriction X_v)$. Then
\begin{equation}\label{4.22:eq-majorant-transfer-a}
f(\mathrm{end})-f(\mathrm{start})=\sum_{\pi\in\Pi(v)}D_{\pi}f ,
\end{equation}
where $\mathrm{end}$ and $\mathrm{start}$ are the end-point and the starting point of the
expansion in $\tau$. Here $\mathfrak{P}^{\mathrm{base}}$ collects the parameter families that are
debited by their supremum: the $\mathsf{s}$-families, at amplitude one; the parameters
$(t,\sigma,b,u)$ of a preparatory stage at $(\mathrm{P})$, whose debit is the $\eta^{(l)}$ of
Corollary~\ref{4.22:cor-preparatory-chart}; and a parameter $t\in[0,1]$ of a bracket. The factor
$\mathfrak{P}^{\mathrm{inc}}$ consists of at most one increment
disc: at $(\mathrm{T})$ the disc of the variable $\mathsf{t}$ of the correlation theorem, and at
$(\mathrm{L})$ the disc of radius $\rho_p$. For a piece of an $\mathsf{s}$-family there is no
increment disc. The complex parameter $\lambda$ used to bound such pieces belongs to the bound,
and it is not a coordinate of the identity~\eqref{4.22:eq-majorant-transfer-a}.

Each $D_\pi$ is a finite composition of evaluations, integrations over $[0,1]$, derivatives and
M\"obius differences in $\tau$. It contains at least one derivative, one integral over $[0,1]$ of
a derivative, or one M\"obius difference, so that $D_\pi c=0$ for every $c$ independent of
$\tau$. For the source expansion this is the non-emptiness of the family $\mathsf{y}$. For the
preparatory expansion of Corollary~\ref{4.22:cor-preparatory-chart} it is the non-emptiness of
$(A,Q)$. For the correlation theorem at $(\mathrm{T})$ it is the derivative $\partial_{\mathsf{t}}$.

The index set $\Pi(v)$ is fixed by two things: the combinatorial data of the expansion, and a
vanishing clause. The vanishing clause uses only the locality of $\mathsf{R}$ and a closed set
$\mathcal{R}_v\subset X_v$ outside of which $f_v$ does not depend on the response. The set
$\mathcal{R}_v$ contains the $M$-cube, at the scale of $s$, of every law site $s$ of $v$. This is
consistent because law sites lie in the term's domain, and because, by
Lemma~\ref{4.22:lem-axialising-map}\,(i), the top block $B^{\mathrm{top}}(s)$ lies in the $M$-cube of
the scale of $s$.
\item[(M2)] (Containment.) For every configuration of the site's tube and every
$\tau\in\mathfrak{P}$, the configuration $\mathsf{R}(\tau)$ lies in the stored domain:
\begin{equation}\label{4.22:eq-majorant-transfer-1}
\mathsf{R}(\mathfrak{P}\times\text{tube})\subset\mathfrak{U}_v .
\end{equation}
\item[(M3)] For every $\phi$ holomorphic on $\mathfrak{P}$,
\begin{equation}\label{4.22:eq-majorant-transfer-2}
|D_\pi\phi|\le c_\pi\sup_{\mathfrak{P}}|\phi| ,
\end{equation}
where the $c_\pi$ are numbers fixed by the combinatorial data of the expansion.
\item[(M4)] (Count.) The passage from the numbers $c_\pi$ and the right-hand sides of the bounds on
the terms expanded, as functions of $\mathsf{d}$, to the right-hand sides of the bounds on the
output terms is numerical. It consists of anchored sums over labels and a majorant series, and it
reads no value of $f$.
\end{itemize}
Let $\mathfrak{u}(\tau):=\mathfrak{u}[\mathsf{R}(\tau)]$ be the axialising gauge map of
Lemma~\ref{4.22:lem-axialising-map} at the configuration $\mathsf{R}(\tau)$. At $(\mathrm{L})$ it is
instead the presenter's map $h_q$, given by clause (a) of the presentation
proposition. Assume that at every law site $s$
\begin{equation}\label{4.22:eq-majorant-transfer-3}
\sup_{\mathfrak{P}}\operatorname{pol}\bigl(\mathfrak{u}(\tau)(s)\bigr)\le\eta_0(s)+\mu(s)<\varepsilon(s).
\tag{R}
\end{equation}
Here $\eta_0(s)$ is the supremum of $\operatorname{pol}(\mathfrak{u}(\tau)(s))$ over $\tau$ ranging
over $\mathfrak{P}^{\mathrm{base}}$ at the origin of $\mathfrak{P}^{\mathrm{inc}}$. The excess
contributed by the increment disc is at most the margin $\mu(s)$, by the hypothesis of whichever
of the two increment lemmas applies at the site. Put
\begin{equation}\label{4.22:eq-majorant-transfer-4}
f^{\sharp}(\tau,g'):=f_v^{\sharp}\bigl(\mathsf{R}(\tau);\mathbf{A},\boldsymbol{\Phi},
\mathfrak{u}(\tau)g'\bigr).
\end{equation}
Then the following hold.
\begin{itemize}
\item[(i)] The identity~\eqref{4.22:eq-majorant-transfer-a} holds for $f^{\sharp}(\cdot,g')$ with the
same index set $\Pi(v)$. At $G$-valued $g'$ its left-hand side is the corresponding bracket of
Ba{\l}aban's construction at the rotated arguments
$(\mathbf{A}^{[g']},\boldsymbol{\Phi}^{[g']})$.
\item[(ii)] For every $\pi\in\Pi(v)$,
\begin{equation}\label{4.22:eq-majorant-transfer-5}
|D_\pi f^{\sharp}|\le c_\pi\,\mathfrak{N}^{\mathbb{C}}_{\varepsilon}(v)
\quad\text{on }\mathfrak{G}_{\varepsilon-\eta_0-\mu}(X_v).
\end{equation}
\item[(ii)$'$] For every $\pi\in\Pi(v)$,
\begin{equation}\label{4.22:eq-majorant-transfer-6}
|D_\pi f^{\sharp}|\le c_\pi\operatorname{osc}_{\mathfrak{P}}(f^{\sharp}),
\tag{M$3'$}
\end{equation}
where
\begin{equation*}
\operatorname{osc}_{\mathfrak{P}}(f^{\sharp})
:=\sup_{g'}\,\sup_{\tau,\tau'\in\mathfrak{P}}\bigl|f^{\sharp}(\tau,g')-f^{\sharp}(\tau',g')\bigr|
\le2\,\mathfrak{N}^{\mathbb{C}}_{\varepsilon}(v),
\end{equation*}
and the supremum is over $g'\in\mathfrak{G}_{\varepsilon-\eta_0-\mu}(X_v)$.
\item[(iii)] The output terms are orbit data of radius $\varepsilon-\eta_0-\mu$. They obey the
right-hand sides of the expansion's bounds, as functions of $\mathsf{d}$, with $\mathfrak{N}$
replaced by $\mathfrak{N}^{\mathbb{C}}$.
\end{itemize}
\end{lemma}

\begin{proof}
\emph{(i).} The identity~\eqref{4.22:eq-majorant-transfer-a}, with the fixed index set $\Pi(v)$, holds
for every function of $\tau$ that has the locality of $f$. We show that $f^{\sharp}(\cdot,g')$ has
this locality.

By Lemma~\ref{4.22:lem-axialising-map}\,(i), $\mathfrak{u}(\tau)(s)$ depends on $\mathsf{R}(\tau)$
only through its values on $B^{\mathrm{top}}(s)$. This block is contained in the $M$-cube of $s$ at
its scale, and that cube is contained in $\mathcal{R}_v$ by (M1). Hence $f^{\sharp}$ depends on
$\mathsf{R}$ exactly where $f$ does.

The vanishing clause of \cite{B-CMP116} reads: ``If there are other components, then the derivatives
with respect to $s$ restricted to these components render the term equal to $0$.'' (In this
sentence the letter $s$ denotes the parameters of the cluster expansion of \cite{B-CMP116}, not a
law site.) It refers only
to the places where the term depends on the configuration. Therefore the vanishing clause and the
index set $\Pi(v)$ are unchanged when $f$ is replaced by $f^{\sharp}(\cdot,g')$.

On top blocks on which the site's Landau layer $\mathcal{K}$ vanishes identically,
$\mathfrak{u}(\tau)=1$. This follows from Lemma~\ref{4.22:lem-axialising-map}\,(i)
and from the second clause of the regauge lemma: if the regauge generator $\mathsf{Y}$ vanishes
identically, then the carried gauge transformation $\mathsf{u}$ is identically $1$.

At the end-points, three facts identify the left-hand side of~\eqref{4.22:eq-majorant-transfer-a}
for $f^{\sharp}$ at $G$-valued $g'$ with Ba{\l}aban's bracket at
$(\mathbf{A}^{[g']},\boldsymbol{\Phi}^{[g']})$. They are Lemma~\ref{4.22:lem-axialising-map}\,(ii)
and properties (O2) and (O4) of orbit data.

\emph{(ii).} For $a,b\in G^{\mathbb{C}}$ we have $\|ab\|\le\|a\|\,\|b\|$ and
$\|(ab)^{-1}\|\le\|b^{-1}\|\,\|a^{-1}\|$. Hence $\operatorname{pol}(ab)\le\operatorname{pol}(a)+\operatorname{pol}(b)$.
Let $g'\in\mathfrak{G}_{\varepsilon-\eta_0-\mu}(X_v)$ and $\tau\in\mathfrak{P}$. At every law site
$s$ this subadditivity and~\eqref{4.22:eq-majorant-transfer-3} give
\begin{align*}
\operatorname{pol}\bigl(\mathfrak{u}(\tau)(s)g'(s)\bigr)
&\le\operatorname{pol}\bigl(\mathfrak{u}(\tau)(s)\bigr)+\operatorname{pol}\bigl(g'(s)\bigr)\\
&\le\eta_0(s)+\mu(s)+\operatorname{pol}\bigl(g'(s)\bigr)<\varepsilon(s).
\end{align*}
Together with the containment (M2), this shows that the map
$(\tau,g')\mapsto(\mathsf{R}(\tau),\mathfrak{u}(\tau)g')$ sends
$\mathfrak{P}\times\mathfrak{G}_{\varepsilon-\eta_0-\mu}(X_v)$ into
$\mathfrak{U}_v\times\mathfrak{G}_{\varepsilon}(X_v)$. The map is holomorphic there; for
$\mathfrak{u}$ this is the holomorphy of Lemma~\ref{4.22:lem-axialising-map}\,(i).

By property (O1) of the orbit datum $f_v^{\sharp}$, the function $f^{\sharp}(\cdot,g')$ is therefore
holomorphic on $\mathfrak{P}$, and
$\sup_{\mathfrak{P}}|f^{\sharp}(\cdot,g')|\le\mathfrak{N}^{\mathbb{C}}_{\varepsilon}(v)$. Applying
(M3) with $\phi:=f^{\sharp}(\cdot,g')$ gives~\eqref{4.22:eq-majorant-transfer-5}.

\emph{(ii)$'$.} Fix $g'\in\mathfrak{G}_{\varepsilon-\eta_0-\mu}(X_v)$ and $\tau_0\in\mathfrak{P}$, and
put $c:=f^{\sharp}(\tau_0,g')$. The number $c$ does not depend on $\tau$, so $D_\pi c=0$ by (M1).
Since $D_\pi$ is a composition of linear operations, $D_\pi f^{\sharp}=D_\pi(f^{\sharp}-c)$.

The function $\phi:=f^{\sharp}(\cdot,g')-c$ is holomorphic on $\mathfrak{P}$ by the proof of (ii).
By (M3),
\begin{equation*}
|D_\pi f^{\sharp}|\le c_\pi\sup_{\mathfrak{P}}|f^{\sharp}(\cdot,g')-c|
\le c_\pi\operatorname{osc}_{\mathfrak{P}}(f^{\sharp}),
\end{equation*}
where the second inequality holds because $\tau_0\in\mathfrak{P}$. This
is~\eqref{4.22:eq-majorant-transfer-6}.

For the bound on the oscillation, let $\tau,\tau'\in\mathfrak{P}$. By the triangle inequality and
the bound on $\sup_{\mathfrak{P}}|f^{\sharp}(\cdot,g')|$ established in (ii),
\begin{equation*}
|f^{\sharp}(\tau,g')-f^{\sharp}(\tau',g')|\le|f^{\sharp}(\tau,g')|+|f^{\sharp}(\tau',g')|
\le2\,\mathfrak{N}^{\mathbb{C}}_{\varepsilon}(v).
\end{equation*}

\emph{(iii).} Condition (M4) reads numbers only: the $c_\pi$ and the right-hand sides of the bounds
on the terms expanded. It therefore applies unchanged when the bounds on the $D_\pi f$ are replaced
by the bounds~\eqref{4.22:eq-majorant-transfer-5}, in which $\mathfrak{N}$ is replaced by
$\mathfrak{N}^{\mathbb{C}}_{\varepsilon}$.

Property (O2) holds for the output terms termwise, by (i). The output terms are holomorphic on the
chart of radius $\varepsilon-\eta_0-\mu$, and they satisfy their bounds. Both facts follow from the
uniform convergence of the anchored sums and of the majorant series under the majorants, by the
argument of the lemma on the complexified preparatory stage.
\end{proof}

\begin{remark}
We add three observations on the scope of Lemma~\ref{4.22:lem-majorant-transfer}.

(a) The lemma does not assert that any particular expansion satisfies (M1)--(M4). For each
expansion to which it is applied, these conditions are established from that expansion's own
identities and bounds.

(b) The inequality~\eqref{4.22:eq-majorant-transfer-6} is the single form in which the oscillation
bounds of this chapter enter. Three bounds take this form.
\begin{itemize}
\item For a deep term $v$, clause (a) of the oscillation lemma for the four $\mathsf{s}$-families
bounds $\operatorname{osc}_{\mathfrak{P}}(f^{\sharp})$ over the product polydisc of these families by
$C_{3F}e^{-\frac12\delta_PD_X}\mathfrak{N}^{\mathbb{C}}_{\varepsilon}(v)$; this is a bound on the
right-hand side of~\eqref{4.22:eq-majorant-transfer-6}. Clause (b) of the same lemma is
\eqref{4.22:eq-majorant-transfer-6} with $c_\pi=e^{-(\kappa_1-1)\sum_t|\mathsf{y}_t|}$.
\item Clause (ii)$_2$ of the kept-lineage lemma, the oscillation bound for the kept lineage
$\mathfrak{K}(Y_i)$, is an instance.
\item Clause (iii) of the arrest entry, the oscillation bound for $\mathfrak{Q}(\mathfrak{a})$
attached to an arrest $\mathfrak{a}$, is an instance.
\end{itemize}
In each case the variation of the map $\mathfrak{u}(\tau)$ is included in
$\operatorname{osc}_{\mathfrak{P}}(f^{\sharp})$.

(c) A bound on the supremum over the datum implies nothing for the supremum over the chart, since
$\mathfrak{N}^{\mathbb{C}}_{\varepsilon}\ge\mathfrak{N}$. An expansion that asserts a bound without
satisfying (M1)--(M4) transfers nothing to the chart; the hypothesis of the lemma is (M1)--(M4).
\end{remark}

\begin{corollary}[Older law sites through a fluctuation site]\label{4.22:cor-older-law-sites}
Indices are those of the rotated bound on the localised terms, that is, of the form of the bound on
the localised terms $V(Y,\mathcal{S}')$ in which the arguments are rotated by complex gauge maps. The
price paid in $(\alpha)$ below, namely the replacement of $\eta_p$ by $\eta_p^{1/2}$ together with the
bound $\Lambda^{\mathrm{res}}\le1$, is computed in the step index, in which step $k$ produces the
$k$-th density.
At the fluctuation site $(\mathrm{T})$ every block-type unit is read as in
Definition~\ref{4.22:def-block-term-reading}. The \emph{older stored units} at step $k$ are the
following units, created at earlier steps, which enter the fluctuation integral of step $k$:
\begin{itemize}
\item Ba{\l}aban's boundary terms $\mathbf{B}^{(j)}$;
\item the complex values left by chains of the large-field stages arrested at a level $k_1<k$
\cite{B-CMP122a};
\item the families $\mathfrak{Q}(\mathfrak{a})$ attached to the live arrests $\mathfrak{a}$ of the
block-type chains;
\item the terms $\mathbf{B}'^{(j)}$ created by the large-field operation $\mathbf{R}$, which are
wrapped and are indexed by their labels;
\item the terms $V^{(k)}(S)$;
\item what \cite[(1.102)]{B-CMP122b} keeps under $\mathbf{T}'(W)$ and enters a later fluctuation
integral,
namely the kept lineage $\mathfrak{K}(Y_i)$; by the tube statement for kept members these are the
terms $V^{\mathrm{kept}}$ of Proposition~\ref{4.22:prop-fluctuation-site}.
\end{itemize}
A unit is \emph{wrapped} if it belongs to the class $\mathfrak{S}$ of terms touched by a second-class
outcome, and \emph{plain} if it belongs to the class $\mathfrak{A}$. Assume the following.
\begin{itemize}
\item The rotated bound on the localised terms, its clauses (a)--(d), and the dichotomy of the second
step of its proof.
\item The response-expansion lemma of the correlation theorem, with the radii $\eta_p$ in their
positional form.
\item The loss lemma for holomorphic families carrying a riding gauge map, in its stated form.
\item Lemma~\ref{4.22:lem-axialising-map}, clause (ii), and its clause (iv) at $(\mathrm{T})$ in the
residual form \eqref{4.22:eq-axialising-map-a}.
\item The coupling ceilings \eqref{4.22:eq-coupling-ceilings-l}, (v),
\eqref{4.22:eq-coupling-ceilings-e}, \eqref{4.22:eq-coupling-ceilings-f},
\eqref{4.22:eq-coupling-ceilings-g}, \eqref{4.22:eq-coupling-ceilings-h} and (x) of the list of
coupling ceilings.
\item The one-step bounds on the flow of the inverse couplings.
\item Clause (ii) at $(\mathrm{T})$ of the law-site lemma, and the bound
\eqref{4.22:eq-law-site-a}.
\item The radius budget \eqref{4.22:eq-budget-closes-b}.
\item The floors of the order of constants $\kappa\ge10\log L$, $\kappa\ge C_\delta(d)/\delta$,
$\kappa_1-1\ge4L\kappa$ and $c\delta_0M\ge(1+4\beta)\kappa L$, where $C_\delta(d)$ is a constant
depending on $d$ only, and the floor $L\ge L_0$ on the
blocking factor.
\item The separation of the large-field set $Z_k^{\sim}$ from the region $\Omega_{k+1}$ on which the
fluctuation data of the step live by at least $3LR_{k+1}$ $M$-cubes of scale $k$, where $R_{k+1}$ is
Ba{\l}aban's scale number at level $k+1$.
\item Lemma~\ref{4.22:lem-majorant-transfer}.
\end{itemize}
Put $\mu(s):=\mu_k\mathsf{e}(s)/\varepsilon_\star$ at every law site $s$. Let the older stored units
have orbit data of radius $\varepsilon$, where $\varepsilon(s)>\eta^T+\mu(s)$ at every law site $s$,
and put
\begin{equation*}
\varepsilon'(s):=\varepsilon(s)-\eta^T-\mu(s).
\end{equation*}
Then every $V(Y,\mathcal{S}')$ has an orbit datum $V(Y,\mathcal{S}')^{\sharp}$, holomorphic on the
product of the domain of clause (a) of the rotated bound with $\mathfrak{G}_{\varepsilon'}(Y)$. Under
the coupling ceilings listed above, clause (c) of the rotated bound and the bound (d) of
\eqref{4.22:eq-older-law-sites-a} hold with the supremum taken also over the chart. In (d),
$\nu(Y)$ is the right member of clause (d) of the rotated bound, a function of $d_k(Y)$ for $Y$ in
the plain class $\mathfrak{A}$ and of $\mathsf{d}_k(Y)$ for $Y$ in the wrapped class $\mathfrak{S}$.
In the relative formulation it is replaced by
\begin{equation*}
\nu(Y)+(\varepsilon^{\mathrm{rel}}_k+\varepsilon^{\mathrm{kept}}_k)e^{-\kappa_V\mathsf{d}_k(Y)},
\end{equation*}
as in \eqref{4.22:eq-fluctuation-site-a}; here $\kappa_V:=(1-2\delta)\kappa$, and
$\varepsilon^{\mathrm{rel}}_k$ and $\varepsilon^{\mathrm{kept}}_k$ are the constants of the relative
formulation fixed in \eqref{4.22:eq-fluctuation-site-a}. Moreover, for every older stored unit $v$
and every datum:
\begin{itemize}
\item[$(\alpha)$] if $v$ is not wrapped, then on the $\mathsf{t}$-disc of the response-expansion
lemma its piece indexed by Ba{\l}aban's datum $p$,
\begin{equation}\label{4.22:eq-older-law-sites-1}
v_p^{\sharp}:=\prod_{\Delta}\int_0^1d\mathsf{s}(\Delta)\,\partial_{\mathsf{s}(\Delta)}\,
\partial_{\mathsf{t}}\big|_{\mathsf{t}=0}\,
f_v^{\sharp}\bigl(\mathsf{R}_p(\mathsf{t},\mathsf{s});\mathbf{A}_{<k},\boldsymbol{\Phi},
\mathfrak{u}_p(\mathsf{t},\mathsf{s})g'\bigr),
\end{equation}
satisfies line $(\alpha)$ of \eqref{4.22:eq-older-law-sites-a}, and $\Lambda^{\mathrm{res}}\le1$;
\item[$(\beta)$] if $v$ is wrapped and the $\mathsf{s}$-expansion hypothesis at the fluctuation site
holds, then every piece $F^{\sharp}_{\mathsf{y}}$ of its expansion in the
parameters $\mathsf{s}$ satisfies line $(\beta)$ of
\eqref{4.22:eq-older-law-sites-a}. There is no further factor, and the radius is
$\varepsilon-\eta^T$. Here $\eta^T$ is taken over the family of that expansion, as in clause (iii)
of the law-site lemma for that family, and the margin $\ell2$ of
\eqref{4.22:eq-radius-accounting-a} is not drawn;
\item[$(\gamma)$] in $(\alpha)$ and $(\beta)$ (line $(\gamma)$ of
\eqref{4.22:eq-older-law-sites-a}), $\mathfrak{N}^{\mathbb{C}}_{\varepsilon}(v)$ may be
replaced by $\operatorname{osc}_{\mathfrak{P}}(v^{\sharp})$, as in $(\mathrm{M3})'$ of
\eqref{4.22:eq-majorant-transfer-a}, and then further as follows.
\begin{itemize}
\item For $v\in\mathfrak{K}(Y_i)$ (such a $v$ always falls under $(\beta)$), it may be replaced by
$\mathfrak{O}(v)$. Here $\mathfrak{O}(v)\le C_\Sigma\check{R}_{k_W}^{p_\Sigma}$ by
$(\mathrm{ii})_2$ of \eqref{4.22:eq-kept-oscillation-a}, and this factor is paid by
$|\mathsf{y}|\ge\frac52\check{R}_{k_W}$, by Proposition~\ref{4.22:prop-fluctuation-site}.
\item For $v=\mathsf{Q}_i$, it may also be replaced by $2\Sigma^{\mathfrak{q}}(x_k-c_5)$, by
$(\mathrm{ii})_{\mathsf{Q}}$ of \eqref{4.22:eq-kept-oscillation-a}. This alternative is admissible
under the rule for the arrow of the kept quadratic form, and no bound below needs it.
\item For $v=\mathsf{q}_{\mathsf{p}}$, it may be replaced by the right member of
\eqref{4.22:eq-gaussian-oscillation-a}.
\end{itemize}
\end{itemize}
The displays are
\begin{equation}\label{4.22:eq-older-law-sites-a}
\begin{aligned}
&(\mathrm{d}) &&\sum_{\mathcal{S}'}\exp\Bigl[\kappa_A(k)\sum_{j'<k}|\mathcal{S}'_{j'}|\Bigr]
\sup|V(Y,\mathcal{S}')^{\sharp}|\le\nu(Y),\\
&(\alpha) &&|v_p^{\sharp}|\le2\mathsf{w}_{\mathrm{site}}^{-1}e^{-(\kappa_1-1)m(p)}
\mathfrak{N}^{\mathbb{C}}_{\varepsilon}(v)\,\eta_p^{1/2}\max\{1,\Lambda^{\mathrm{res}}\},\\
&(\beta) &&|F^{\sharp}_{\mathsf{y}}|\le\mathfrak{N}^{\mathbb{C}}_{\varepsilon}(v)
e^{-(\kappa_1-1)|\mathsf{y}|},\\
&(\gamma) &&\mathfrak{N}^{\mathbb{C}}_{\varepsilon}(v)\ \text{replaced by}\
\operatorname{osc}_{\mathfrak{P}}(v^{\sharp});\quad
\mathfrak{O}(v)\le C_\Sigma\check{R}_{k_W}^{p_\Sigma}\ \ (v\in\mathfrak{K}(Y_i));\\
& &&2\Sigma^{\mathfrak{q}}(x_k-c_5)\ \ (v=\mathsf{Q}_i);\quad
\text{the right member of \eqref{4.22:eq-gaussian-oscillation-a}}\ \ (v=\mathsf{q}_{\mathsf{p}}),\\
& &&\Lambda^{\mathrm{res}}:=\max_s\Bigl\{\frac{\vartheta^T}{\mu(s)}\,C_b^{1/2}
e^{-\frac12\delta_1(d_{\mathrm{sc}}(s,\Omega)-1)}\Bigr\}.
\end{aligned}
\end{equation}
Here $\mathsf{w}_{\mathrm{site}}$ is the width of the $\mathsf{t}$-disc of the response-expansion
lemma, $m(p)$ is the length attached to Ba{\l}aban's datum $p$ in that lemma, so
that $e^{-(\kappa_1-1)m(p)}$ is its walk factor, and $\delta_1$ is the decay rate of the loss lemma,
which satisfies $\delta_1\ge\delta_\natural$. Throughout, $\Omega$ denotes the region $\Omega_k$ of
step $k$ in the step index fixed above.
\end{corollary}

\begin{proof}
\emph{Proof of $(\alpha)$.} Let $v$ be an older stored unit that is not wrapped, and let
$v_p^{\sharp}$ be its piece \eqref{4.22:eq-older-law-sites-1} indexed by Ba{\l}aban's datum $p$.
This is the piece of the response-expansion lemma, with $v$ read as in
Definition~\ref{4.22:def-block-term-reading}. The derivatives $\partial_{\mathsf{t}}$ and
$\partial_{\mathsf{s}}$ act both on the background $\mathsf{R}_p(\mathsf{t},\mathsf{s})$ and on the
gauge map $\mathfrak{u}_p$.

Summed over Ba{\l}aban's data $p$ (\cite[(3.4), (3.7)]{B-CMP109}, \cite[(1.10)]{B-CMP116}), these
pieces are increments of the single family
$f_v^{\sharp}(\mathsf{R};\,\cdot\,,\mathfrak{u}[\mathsf{R}]g')$, and therefore they reproduce the
difference of its values at the two end-points: the value of $f_v$ at Ba{\l}aban's axial background,
which is the end-point described by clause (ii) of Lemma~\ref{4.22:lem-axialising-map}, minus the
value of $f_v$ at the background $\mathbf{U}$, where $\mathfrak{u}=\mathbb{1}$.

We apply the loss lemma in its stated form; its function $\vartheta(s)$ is defined site by site. The
data of the application are $\rho:=\mathsf{w}_{\mathrm{site}}/\eta_p$, the parameter manifold
$\mathfrak{M}$ of the variables $(\sigma,\mathbf{U},\mathbf{J},\mathsf{B},w)$, the riding gauge map
$u:=\mathfrak{u}_p$, $\eta_0:=\eta^T$, the function $\vartheta(s):=\vartheta_p(s)$ of clause (iv) at
$(\mathrm{T})$ in the residual form \eqref{4.22:eq-axialising-map-a}, and
$\mu(s)=\mu_k\mathsf{e}(s)/\varepsilon_\star$. With $\mathfrak{N}$ the supremum norm and $\Lambda$
the loss factor of the loss lemma, these data give
\begin{equation*}
\frac{\mathfrak{N}}{\rho}\Lambda=\frac{\mathfrak{N}}{\mathsf{w}_{\mathrm{site}}}
\max\Bigl\{\eta_p,\ \max_s\frac{\vartheta^Te_p(s)}{\mu(s)}\Bigr\}.
\end{equation*}
By clauses (a) and (b) of \eqref{4.22:eq-axialising-map-a} and the inequality
$\min\{a,b\}\le(ab)^{1/2}$ for $a,b\ge0$,
\begin{equation}\label{4.22:eq-older-law-sites-2}
e_p(s)\le\eta_p^{1/2}C_b^{1/2}e^{-\frac12\delta_1(d_{\mathrm{sc}}(s,\Omega)-1)}.
\end{equation}
Hence, on $\mathfrak{G}_{\varepsilon'}$, line $(\alpha)$ of \eqref{4.22:eq-older-law-sites-a} holds
with $\Lambda^{\mathrm{res}}$ as defined there, whatever the index of $v$ and the age of its law
sites; the derivatives in $\sigma$ are estimated by the Cauchy inequality on the circles
$|\sigma(\Delta)|=e^{\kappa_1}$, as in the response-expansion lemma. It is the product
$(\mathfrak{N}/\mathsf{w}_{\mathrm{site}})\vartheta^Te_p(s)/\mu(s)$ of the loss lemma that carries
clause (b) of \eqref{4.22:eq-axialising-map-a}, also for an old unit with $s$ at the near end of
its domain $S_v$, where $\vartheta_p(s)$ is of the order of $\vartheta^T$.

\emph{The bound $\Lambda^{\mathrm{res}}\le1$.} We work in the step index: step $k$ produces the
$k$-th density, $\Omega=\Omega_k$, and $k_W\le k-1$. The coupling of the fluctuation site is
$g_{k-1}$, whereas $\mu_k$ is evaluated at $x_k$. By the one-step bounds on the flow of the inverse
couplings, ceiling (v) and \eqref{4.22:eq-coupling-ceilings-h},
\begin{equation*}
\frac{\alpha_{*,k-1}}{\alpha_{*,k}}\le\sqrt2\Bigl(\frac{\log(x_k+B_\sharp)}{\log x_k}\Bigr)^{q_*}
\le2\sqrt2,
\end{equation*}
where $q_*$ is the exponent appearing in \eqref{4.22:eq-coupling-ceilings-h},
and this factor is taken into the maximum of \eqref{4.22:eq-coupling-ceilings-e}.

If $s$ is a pointwise law site, then
\begin{equation*}
\frac{\vartheta^T}{\mu(s)}=\frac{2.1K_uB_H\alpha_{*,k}(x_k-c_5)^2}{\varepsilon_\star},\qquad
d_{\mathrm{sc}}(s,\Omega)\ge8\check{R}_k,
\end{equation*}
the second by clause (ii) at $(\mathrm{T})$ of the law-site lemma.

If $s\in\Sigma_W$ is a law site of a frame $W$, put $\mathsf{a}':=k-k_W$. In a step the fluctuation
site precedes the site of the large-field operation $\mathbf{R}$, so $k_W\le k-1$ and
$\mathsf{a}'\ge1$. By \eqref{4.22:eq-budget-closes-b} at $m=k_W$,
\begin{align*}
\frac{\vartheta^T}{\mu(s)}
&=2.1K_uB_H\alpha_{*,k}(x_k-c_5)^2\,\frac{4C_{\mathrm{fr}}}{\varpi_{\mathrm{fr}}\alpha_{*,k_W}}\\
&\le2.1K_uB_H\cdot8C_{\mathrm{fr}}\varpi_{\mathrm{fr}}^{-1}(x_k-c_5)^2
\Bigl(1+\frac{B_\sharp\mathsf{a}'}{x_k}\Bigr)^{1/2},
\end{align*}
and, by clause (ii) at $(\mathrm{T})$ of the law-site lemma with $m(s):=k_W$ and by
\eqref{4.22:eq-law-site-a},
\begin{equation*}
d_{\mathrm{sc}}(s,\Omega)\ge8\check{R}_k+5R_{\min}(\mathsf{a}'-1),
\end{equation*}
where $R_{\min}$ is the scale constant of \eqref{4.22:eq-law-site-a} and of
\eqref{4.22:eq-coupling-ceilings-g}.
Now $\frac12\delta_\natural\cdot5R_{\min}\ge1$ by \eqref{4.22:eq-coupling-ceilings-g}, and
$\delta_1\ge\delta_\natural$. Further,
\begin{equation}\label{4.22:eq-older-law-sites-3}
\Bigl(1+\frac{B_\sharp\mathsf{a}'}{x_k}\Bigr)^{1/2}e^{-(\mathsf{a}'-1)}\le e(1+B_\sharp)^{1/2},
\end{equation}
which follows from $1+B_\sharp\mathsf{a}'/x_k\le(1+B_\sharp)(1+\mathsf{a}')$, valid since
$x_k\ge1$ (indeed $\log(x_k-c_5)\ge1$ by \eqref{4.22:eq-coupling-ceilings-h} and $c_5\ge0$), and
from $(1+a)^{1/2}e^{-a}\le1$ for $a\ge0$; and, since $\check{R}\ge7$ and by the
statement relating $\check{R}_k$ to $E(x_k)$,
\begin{equation*}
e^{-\frac12\delta_\natural(8\check{R}_k-1)}\le e^{-\delta_\natural\check{R}_k}
\le E(x_k)^{\delta_\natural}.
\end{equation*}
Inserting these bounds, together with the factor $2\sqrt2$ above, into the definition of
$\Lambda^{\mathrm{res}}$ gives
\begin{align*}
\Lambda^{\mathrm{res}}
&\le(x_k-c_5)^2E(x_k)^{\delta_\natural}\,2\sqrt2\,C_b^{1/2}K_uB_H\\
&\qquad\times\max\Bigl\{\frac{17e\,C_{\mathrm{fr}}(1+B_\sharp)^{1/2}}{\varpi_{\mathrm{fr}}},\
\frac{2.1\alpha_*(x_k)}{\varepsilon_\star}\Bigr\}\le1
\end{align*}
by \eqref{4.22:eq-coupling-ceilings-e}, using $2.1\cdot8\le17$. Thus the loss factor of the loss
lemma disappears, and the price of $(\alpha)$ is the replacement of $\eta_p$ by $\eta_p^{1/2}$.

\emph{The price is affordable, and no other budget is changed.}

(i) The resummation over $p$ in the response-expansion lemma produces a constant $C_A$, a sum of
factors of the type of $\eta_p$ taken at the full response rate. For the older stored units the same
resummation runs at half rate:
\begin{equation*}
\sum_p\eta_p^{1/2}e^{-(\kappa_1-1)m(p)}\le C_A^{(1/2)}\cdot(\text{the other factors of }C_A),
\end{equation*}
where $C_A^{(1/2)}$ depends on $d$, $L$, $M$, $\kappa_1$ and $\delta_0M$ only, is fixed before
$\gamma$ in the order of constants, and enters only the share of the older stored units, of order
$o_\gamma(1)$, in clause (iii) of the condition on live terms.

(ii) The dichotomy of the second step of the proof of the rotated bound states that either
$d_j(S_v)$ or the length of the response and of the walks is at least $cR_{k+1}$, half of that decay
giving $e^{-2R_{k+1}}$ and the other half $\kappa_Vd_k(Y)$, by the argument that splits the decay
between these two factors. This dichotomy survives with the response factor at half rate, each half
now being a quarter of the exponent of that argument. The share $e^{-2R_{k+1}}$ survives because
$\min\{\kappa,c\delta_0M\}\cdot3LR_{k+1}/4\ge2R_{k+1}$, which follows from
$\min\{\kappa,c\delta_0M\}\cdot3L\ge8$, a consequence of the floors $\kappa\ge10\log L$ and
$c\delta_0M\ge(1+4\beta)\kappa L$ together with the floors on the blocking factor contained in the
condition $L\ge L_0$ of the order of constants. The share $\kappa_V\mathsf{d}_k(Y)$ survives
because the floor $c\delta_0M\ge(1+4\beta)\kappa L$ on the walk rate leaves, at quarter rate,
\begin{equation*}
\frac{c\delta_0M}{4}\ge\frac{(1+4\beta)\kappa L}{4}\ge(1+4\beta)\kappa\ge\kappa\ge
\kappa_V=(1-2\delta)\kappa
\end{equation*}
per $M$-cube of scale $k$ (clause (c) of the rotated bound), the second inequality by $L/4\ge1$,
which is among the floors on the blocking factor contained in the condition $L\ge L_0$ of the order
of constants; the bound $d_{k+1}(\mathrm{hull})\le d_k/L$ of the same clause
is used as before. The branch $d_j(S_v)$, the walk factors $e^{-(\kappa_1-1)m(p)}$ and the factor
$e^{-\kappa\mathsf{d}_j(S_v)}$ contained in $\mathfrak{N}(v)$ are untouched: the decay on $S_v$
and the walk factors are never halved, since the response-expansion lemma keeps the rate
$(1-\frac12\delta)\kappa$ on $\bar S_v$. Let $n_1'$ be the number of $M$-cubes of scale $k$ of
$S_v$ in the gap, and let $n_2$, $n_3$ be the numbers of the response's and of the walks' own
$M$-cubes of scale $k$. The surplus satisfies
\begin{equation*}
\tfrac12\delta\kappa n_1'+\bigl(\tfrac12c\delta_0M-\kappa_V\bigr)n_2
+(\kappa_1-1-\kappa_V)n_3\ge2R_{k+1},
\end{equation*}
because $3\delta\kappa L\ge4$ (from $\kappa\ge C_\delta(d)/\delta$), $\kappa_1-1\ge4L\kappa$ and
$c\delta_0M\ge(1+4\beta)\kappa L$ are floors of the order of constants. This argument concerns a
unit $v$ that is not wrapped. For a wrapped $v$ the share $\kappa_V\mathsf{d}_k(Y)$ requires the
corridors through strips and is supplied by the corridor bound for wrapped terms, at rate
$\kappa_1$ with its constant $c_l$; in the relative formulation a wrapped $v$ meets no
$\mathsf{t}$-disc at the fluctuation site, so (ii) does not apply to it.

The gap lies off every live strip. Indeed, a unit meets $Z_k^{\sim}$, while the data live on
$\Omega_{k+1}$, and these are at least $3LR_{k+1}$ $M$-cubes of scale $k$ apart, by the separation
hypothesis and \cite[(3.2)--(3.5)]{B-CMP119}. A live strip, belonging to a live pair $(Y,m)$ as in
Proposition~\ref{4.22:prop-fluctuation-site}, exceeds its region $Y$ by its width
\begin{equation*}
m_\sharp\le\tfrac12\check{R}_mL^{m-k}\le LR_{k+1}
\end{equation*}
$M$-cubes of scale $k$ (\cite{B-CMP122a}; the ratios of the numbers $\check{R}$ by the statement
relating $\check{R}$ to $E$, with a factor at most $2$ from \eqref{4.22:eq-coupling-ceilings-h}).
So a live strip lies under one of the eight layers of \cite[(3.2)--(3.5)]{B-CMP119}, whereas the gap
is three layers wide. Moreover $\mathsf{d}$ counts these $M$-cubes: by clause (a) of the strip
lemma, the part $X\setminus\bigcup Y^{\sharp}$ of the component $X$ lying outside the strips is
contained in the set $A$ of that lemma, and by the tree-length bound of the cluster expansion a tree
meeting two cubes has length at least their distance.

(iii) By the share of the older stored units we mean the total contribution of the units listed in
the statement; it is bounded separately for the plain class and for the wrapped class. Line
$(\alpha)$ of \eqref{4.22:eq-older-law-sites-a} multiplies that share by $C_A^{(1/2)}/C_A$ and by
$\max\{1,\Lambda^{\mathrm{res}}\}=1$, so the accounting rule that charges this share spends no
factor $e^{-R}$ of the rotated bound.

\emph{Proof of $(\beta)$.} Apply Lemma~\ref{4.22:lem-majorant-transfer} with empty increment
polydisc. Its hypothesis $(\mathrm{M1})$ is the identity of clause (i) of the source expansion, and
its hypothesis $(\mathrm{M3})$ is the Cauchy inequality on the circles
$|\mathsf{s}(\Delta)-\mathsf{s}_0|=e^{\kappa_1}-1$, whose radius is at least $e^{\kappa_1-1}$.

\emph{Proof of $(\gamma)$.} This is $(\mathrm{M3})'$ of \eqref{4.22:eq-majorant-transfer-a}: the
derivative $\partial_{\mathsf{t}}$ in $(\alpha)$, respectively the derivatives in $\mathsf{s}(\Delta)$,
$\Delta\in\mathsf{y}$, with $\mathsf{y}\ne\emptyset$ in $(\beta)$, annihilate every constant that
does not depend on the parameters.

Finally, the response-expansion lemma is linear in $f_v$ and carries the parameters through; the
resummations over $p$, over the blocks $\square$ and over the sets $Y_0$ of that resummation do not
involve the parameters; each
term is assigned to the strong set $\mathcal{S}'$ of its unit, with the weight
$e^{\kappa_A(k)|\mathcal{S}|}$ of the complex inductive bound; and the property that at $G$-valued
rotations the orbit datum is the term evaluated at the rotated arguments holds term by term.
\end{proof}

\begin{remark}[Remarks on the axial reading and the role of the chart]\label{4.22:rem-axial-reading-remarks}
\begin{enumerate}
\item[(1)] \emph{What the reading rule changes.} Take the data real: real parameters and
$G$-valued rotations. At such data two objects differ from the corresponding bracket of
Ba{\l}aban's construction (for the preparatory stages, the bracket identified in clause (1) of
Corollary~\ref{4.22:cor-preparatory-chart}): the M\"obius-inversion identity of the preparatory
stages, read
according to Definition~\ref{4.22:def-block-term-reading}, and the pieces of the correlation
theorem's resummation over the data $p$ at the fluctuation site $(\mathrm{T})$. In both
cases the difference is the rotation of the exterior variables
\begin{equation*}
\lambda\longmapsto\lambda^{[\mathfrak{u}]},\qquad \mathfrak{u}=\mathfrak{u}(b,\mathbb{U}),
\end{equation*}
where $\mathfrak{u}$ is the rotation of Definition~\ref{4.22:def-block-term-reading},
$G$-valued at real data, a function of the argument $b$ of the background map and of
the configuration $\mathbb{U}$. The function $\mathbb{C}^{(l)}$ of a preparatory stage
therefore changes, but no bound changes. The parameter domain
$\mathfrak{D}_{\mathcal{S}}$, the boxes and the supports are stable under
$\operatorname{Ad}_G$, and the majorants do not see the rotation. From here on, the function
of a preparatory stage is the one of Corollary~\ref{4.22:cor-preparatory-chart}, namely the
output $\mathbb{C}^{(l)}_k(X,\mathcal{S})^{\sharp}$, which is an orbit datum on the surviving
law sites.
\item[(2)] The tube form of the stored units cannot carry the reading of
Definition~\ref{4.22:def-block-term-reading}, on the Cauchy polydiscs, at those slots of a
stored unit that are of the two kinds called strong and silent. The orbit chart
$\mathfrak{G}_{\varepsilon}(S)$ carries it, and it
requires no idea beyond the lemma bounding a stored unit read through a complex rotation that
moves with the parameters of the site, the lemma applied in part $(\alpha)$ of the
proof of Corollary~\ref{4.22:cor-older-law-sites}.
\item[(3)] The lemma of the correlation theorem that treats the fluctuation site
$(\mathrm{T})$, and the estimates of the preparatory stages at $(\mathrm{P})$
(Corollary~\ref{4.22:cor-preparatory-chart}), both act at sites
of the correlation theorem itself. These are the sites at which that theorem reads its
stored block-type units, the units treated in
Corollary~\ref{4.22:cor-older-law-sites}. The Cauchy discs at these two sites contain
complex values of $(\mathsf{t},\sigma,t,\mathbf{U})$. Here $\mathsf{t}$ is the Cauchy
variable of the correlation theorem, $\sigma$ and $t$ are the parameters of the background
map $\Phi_{t,\sigma}(b)$, and $\mathbf{U}$ is the background gauge field. At such values the
rotation $\mathfrak{u}$ of Definition~\ref{4.22:def-block-term-reading} is not $G$-valued.
Consequently the correlation theorem reads, at complex rotations, those slots of these units
that are of the kind called live, and it does so through
Corollaries~\ref{4.22:cor-older-law-sites}
and~\ref{4.22:cor-preparatory-chart}. At $(\mathrm{T})$ and at $(\mathrm{P})$ this reading is
independent of the reading at the presentation site $(\mathrm{L})$. Beyond these two
corollaries the correlation theorem needs nothing new: its constants are unchanged except
that, for these units, the positional factor $\eta_p$ of
Corollary~\ref{4.22:cor-older-law-sites} enters the resummation over the data $p$ as
$\eta_p^{1/2}$ and the constant $C_A$ of that resummation becomes $C_A^{(1/2)}$; and the
ceilings listed among the hypotheses of
Corollary~\ref{4.22:cor-older-law-sites} are in force.
\end{enumerate}
\end{remark}

\begin{definition}[The induction hypothesis on orbit data at level $k$]\label{4.22:def-induction-hypothesis}
Let $0\le k\le K$. Consider the effective density $\rho_k$. A \emph{stored term} of its action is
one of the terms listed below, taken at the point where it enters the large-field operation
$\mathbf{R}$ of the step. They are among the terms carried, with their clauses, in the
domain list of the preparatory chain (Definition~\ref{4.22:def-kept-domain-list}), namely:
\begin{itemize}
\item[(1)] Ba{\l}aban's boundary terms $\mathbf{B}^{(\ell)}$, $\ell\le k$;
\item[(2)] the terms $\mathbf{B}'$ produced by the earlier applications of $\mathbf{R}$,
each on its stored background domain in the sense of the frame lemma
(Lemma~\ref{4.22:lem-frame-lemma}), those of the class $\mathfrak{S}$ of terms touched by a
second-class outcome included;
\item[(3)] the arrested chain terms $\mathbb{C}^{(i)}$, the composite terms formed at the
arrests considered in~(b) below, and the change-of-variables terms.
\end{itemize}
The kept members of~(c) below, also carried in that list, are not stored terms.
For every law site $s$, $\varepsilon_k(s)$ denotes the current radius at $s$ and
$\eta_{k'}(s)$, $\eta^{\mathrm{inc}}_{k'}(s)$ denote the total debit and the reserve
incurred at $s$ in the step from level $k'$. The \emph{induction hypothesis at level $k$} is
the conjunction of the following four requirements.
\begin{itemize}
\item[(a)] Every stored term $T$ has an orbit datum $T^{\sharp}$ of the current radius
$\varepsilon_k(s)$ at each of its law sites $s$. Moreover $T$ satisfies the flat,
strong-set-weighted size bound of the inheritance theorem, with the rotated sup norm
$\mathfrak{N}^{\mathbb{C}}$ and the length $\mathsf{d}_{\ell}(T)$ assigned to $T$:
\begin{equation}\label{4.22:eq-induction-hypothesis-a}
\sum_{\mathcal{S}} e^{\kappa_A(k)|\mathcal{S}|}\,
\mathfrak{N}^{\mathbb{C}}\bigl(T(X_T,\mathcal{S})\bigr)
\le c^{\mathrm{sz}}_T\,e^{-\kappa_T\,\mathsf{d}_{\ell}(T)}.
\end{equation}
Here $X_T$ is the domain of $T$, and the sum runs over the strong sets $\mathcal{S}$ of the
strong-set weights of the inheritance theorem, $\kappa_A(k)$ being the weight constant at
level $k$; in this definition the letter $\mathcal{S}$ denotes only this summation variable.
For a term of the class $\mathfrak{S}$, \eqref{4.22:eq-induction-hypothesis-a} is the
creation bound of the lineage statement when the term has just been created, and its
descendant bound otherwise; in both cases $c^{\mathrm{sz}}_T$, $\kappa_T$ and the core of
$T$, a set attached to $T$ together with these two constants, are those fixed in the
statement on the large-field stages of the step. For the other stored terms these constants
and the core are those fixed where the term is produced.
\item[(b)] Let $\mathfrak{a}$ be a composite arrest that is live, that is, still carrying its
large-field factor at level $k$. Its families $\mathfrak{Q}(\mathfrak{a})$ of pieces
$\mathsf{q}_{\mathsf{p}}$ of the kept Gaussian have as orbit datum the datum of clause~(i) of
Proposition~\ref{4.22:prop-gaussian-oscillation}, and they are estimated only by means of
its clause~(iii), the oscillation bound \eqref{4.22:eq-gaussian-oscillation-a}. They are
never estimated by the sizes $\bar{\mathsf{q}}_{\mathsf{p}}$ of their pieces.
\item[(c)] Let $W=Y_i$ be a frame, that is, a second-class component, and let $k_W$ denote
the level of its large-field operation $\mathbf{T}'_{k_W}(Y_i)$, which is live at level $k$.
The two members of the kept lineage $\mathfrak{K}(Y_i)$, one of which is the kept quadratic
form $\mathsf{Q}_i$, are kept under it; they sit inside the operation
$\mathbf{T}'_{k_W}(Y_i)$ of \cite[(1.71)]{B-CMP122b}. Each such member $v$ has as
orbit datum the datum of clause~(ii) of the
frame lemma (Lemma~\ref{4.22:lem-frame-lemma}) on $\mathfrak{U}_v(Y_i^{\sharp})$. For the
kept quadratic form $\mathsf{Q}_i$ this datum follows from the two bounds on $\mathsf{Q}_i$
assumed on the sourced set $\mathfrak{S}_{\mathrm{src}}(Y_i^{\sharp})$ in
Lemma~\ref{4.22:lem-kept-oscillation}. The member $v$ is estimated by means of the
strengthened source-tube bound, which yields, by clause~(o) of
Lemma~\ref{4.22:lem-kept-oscillation}, displayed in \eqref{4.22:eq-kept-oscillation-a},
\begin{equation*}
\mathfrak{O}(v)\le C_{\Sigma}\,\check{R}_{k_W}^{p_{\Sigma}}.
\end{equation*}
The cost of this bound is absorbed in the leaf majorant $\mu^{\mathrm{leaf}}$ and
the leaf total $\mathsf{M}^{\mathrm{leaf}}$ of the counting ceilings for the kept members.
The member $v$ is never estimated by a size. The kept members are not terms to which
\eqref{4.22:eq-induction-hypothesis-a} is applied; they enter only the domain list, as
declared in Definition~\ref{4.22:def-kept-domain-list}.
\item[(d)] For every law site $s$,
\begin{equation}\label{4.22:eq-induction-hypothesis-1}
\sum_{k'<k}\bigl[\eta_{k'}(s)+\eta^{\mathrm{inc}}_{k'}(s)\bigr]\le\tfrac14\,\mathsf{e}(s).
\end{equation}
\end{itemize}
For $k\ge1$, at every law site $s$ of a coordinate of level $k-1$ one has
$\varepsilon_k(s)=\varepsilon_{\star}$, since no part of the radius has yet been consumed there.

The families in~(b) are kept out of \eqref{4.22:eq-induction-hypothesis-a} because the
rotated sup norm of a piece $\mathsf{q}_{\mathsf{p}}$ is of the order of
$\mathfrak{b}_W^2$, a positive power of the logarithm of an inverse coupling of an earlier
level, so that a constant $c^{\mathrm{sz}}_T$ dominating it would depend on the age of the
arrest.
\end{definition}

\begin{lemma}[The order of sites within a step is well founded]\label{4.22:lem-site-order}
In this lemma the \emph{run} is the inductive construction, step by step, of the effective densities.
Four statements are proved together in one induction on the step along the run: the closure
statement of the induction, the treated-component lemma, the relative statement on the orbit chart,
and the source-tube bounds. An item \emph{reads} another if its statement or its proof uses that
other item as an input. A statement is \emph{run-free} if it reads no object of the run.
For a step $k$ we write $(\mathrm{g})_k$ for the output of the $\mathbf{R}$-site of step $k$: the
terms of that site with orbit data of radius $\varepsilon_{k+1}$, the frames created there with data
at $\mathsf{e}$, and the source-tube bounds for every large-field region $W$ created at a step
$k_W\le k$. We write $(\mathrm{a})_k$, $(\mathrm{b})_k$, $(\mathrm{c})_k$ for the three lettered
clauses at step $k$ of the correlation theorem's induction in chart form; clauses $(\mathrm{b})_k$ and
$(\mathrm{c})_k$ enter, together with the induction hypothesis at level $k$, the chart form of its
induction hypothesis at step $k$ (input (i1) below). The sites of step $k$ are traversed in the order
\begin{equation}\label{4.22:eq-site-order-1}
\begin{gathered}
(\mathrm{g})_{k-1}\;\to\;(\mathrm{T})\;\to\;
\text{[the induction hypothesis at level }k\text{, Definition~\ref{4.22:def-induction-hypothesis}]}
\;\to\;(\mathrm{P})\\
\to\;(\mathrm{c})_k\;\to\;\text{[the treated-component lemma]}\;\to\;(\mathrm{L})\;\to\;
(\mathrm{g})_k .
\end{gathered}
\end{equation}
Here the treated-component lemma reads the radius accounting as it stands at that point of the chain.
Inside the passage from $(\mathrm{L})$ to $(\mathrm{g})_k$ the order is: the presentations and their
lists; then the source-tube bounds at step $k$; then the modification statement and
Ba{\l}aban's formulas \cite[(1.90)--(1.99)]{B-CMP122b}. Accordingly the items established at
step $k$ are arranged in rows, and pairs (step; row) are ordered lexicographically:
\begin{equation}\label{4.22:eq-site-order-a}
\begin{gathered}
\mathrm{O}_0<\mathrm{O}_1<\mathrm{O}_2<\mathrm{O}_3<\mathrm{O}_4<\mathrm{O}_5<\mathrm{O}_6
<\mathrm{O}_{7a}<\mathrm{O}_{7b}<\mathrm{O}_{7c}<\mathrm{O}_8,\\
(k;\mathrm{O})<(k';\mathrm{O}')\iff k<k'\ \text{ or }\ \bigl(k=k'\ \text{and}\ \mathrm{O}<\mathrm{O}'\bigr),
\end{gathered}
\end{equation}
where the rows $\mathrm{O}_0$--$\mathrm{O}_8$ of \eqref{4.22:eq-site-order-a}, with what each
establishes at step $k$, by what it is established, and what it reads, are the following.
\begin{itemize}
\item[$\mathrm{O}_0$] \emph{Item:} the two-sided flow bounds for the running couplings at indices at
most $k$; the label of the step: the regions $Z_j$, $\Omega_j$, $\Lambda_j$ ($j\le k$), the treated
class $Z$, the large-field regions live or not, the cores $\mathsf{C}_C$, the formats
$(A_F,\mathcal{P}_F)$, membership of the wrapped lineage $\mathfrak{S}$, and the families $\sigma_0$.
\emph{Established by:} the coupling-fixing statement (its clause (0) and its ordering clause); the
stage of the correlation theorem's run that establishes the flow bounds, through the level-independence
statement of the increment; the declarations of the source-free scaffold.
\emph{Reads:} data of levels at most $k-1$. From this row on, the scaffold lemma (clauses
(i)--(v) and its clause (iii) for the wrapped lineage), the localisation statement on the scaffold
maps, and the radius theorem
at step $k$ are available: they read row $\mathrm{O}_0$ and the live ceiling only.
\item[$\mathrm{O}_1$] \emph{Item:} $(\mathrm{g})_{k-1}$: the terms of the $\mathbf{R}$-site of step
$k-1$, with data of radius $\varepsilon_k$; the frames created there, with data at $\mathsf{e}$; the
source-tube bounds for every large-field region $W$ with $k_W\le k-1$. \emph{Established by:} row
$\mathrm{O}_8$ of step $k-1$. \emph{Reads:} nothing further.
\item[$\mathrm{O}_2$] \emph{Item:} the fluctuation site $(\mathrm{T})$: the orbit data
$V(Y,\mathcal{S}')^{\sharp}$, the units of $V^{(k-1)}$, the orbit data $\mathfrak{K}(X,\mathcal{S})^{\sharp}$
of the kept members, and $V^{\mathrm{kept}}$, each with its right member. \emph{Established by:}
Corollary~\ref{4.22:cor-older-law-sites}, clauses $(\alpha)$--$(\gamma)$; the change-of-variables
lemma and its corollary; the chart statement of the mean field, the one-term formula in its extended
form, and the remaining chart statement of the fluctuation site;
Lemma~\ref{4.22:lem-majorant-transfer}; Lemma~\ref{4.22:lem-kept-holomorphy} at $(\mathrm{T})$;
Proposition~\ref{4.22:prop-fluctuation-site}. \emph{Reads:} row $\mathrm{O}_1$ (data and right
members); Lemma~\ref{4.22:lem-axialising-map} at this site; row $\mathrm{O}_0$ (clauses (i),
(ii) at $(\mathrm{T})$ and (iii) of the scaffold lemma, its clause (iii) for the wrapped lineage,
and the geometric, summation and anchoring counts of the label); the correlation theorem's
statements at the fluctuation site, namely its homogeneity statement, its chart statement on the
terms $V$, its two flat bounds and its supply statement; for the wrapped lineage, clause (r2), its
fluctuation-site statement with condition (M2) of Lemma~\ref{4.22:lem-majorant-transfer}, and the
fluctuation-site part of its clause (r3); the source-tube bounds at steps at most $k-1$, through
clauses (o) and $(\mathrm{ii})_2$ of Lemma~\ref{4.22:lem-kept-oscillation};
Proposition~\ref{4.22:prop-gaussian-oscillation}.
\item[$\mathrm{O}_3$] \emph{Item:} clause $(\mathrm{a})_k$ and the induction hypothesis at level $k$
(Definition~\ref{4.22:def-induction-hypothesis}). \emph{Established by:} row $\mathrm{O}_2$ for the
new terms; row $\mathrm{O}_1$, with the radius lowered by the debits of the fluctuation site, for the
older leaves. \emph{Reads:} rows $\mathrm{O}_1$, $\mathrm{O}_2$.
\item[$\mathrm{O}_4$] \emph{Item:} the preparatory stages $(\mathrm{P})$, $l=1,\dots,N_k$: the stage
condition at stage $l$; the orbit data $\mathbb{C}^{(l)}_k(X,\mathcal{S})^{\sharp}$; at an arrest,
the left-over chain terms $\mathbb{C}^{(i)}$, the composites, the families
$\mathfrak{Q}(\mathfrak{a})$ and the usual new terms. \emph{Established by:} the inheritance theorem;
Corollary~\ref{4.22:cor-preparatory-chart}, the chart proposition of the stage and the second lemma
of the inheritance group; the
arrest lemma; the change-of-variables lemma at a stage;
Proposition~\ref{4.22:prop-gaussian-oscillation}(i); Lemma~\ref{4.22:lem-kept-holomorphy} at
$(\mathrm{P})$; Proposition~\ref{4.22:prop-preparatory-stage}. \emph{Reads:} row $\mathrm{O}_3$; the
stages before $l$; clause (iv) of Lemma~\ref{4.22:lem-axialising-map} at $(\mathrm{P})$; row
$\mathrm{O}_0$; the domain-list statement and the two further preparatory-stage statements, with the
stage restriction; for the wrapped lineage, its preparatory-stage statement, the preparatory-stage
part of clause (r3) (its anchoring) and condition (M2) of Lemma~\ref{4.22:lem-majorant-transfer} at
$(\mathrm{P})$; the source-tube bounds at steps at most $k-1$, through clause $(\mathrm{ii})_2$ of
Lemma~\ref{4.22:lem-kept-oscillation}, at old anchors only.
\item[$\mathrm{O}_5$] \emph{Item:} clauses $(\mathrm{b})_k$, $(\mathrm{c})_k$ and the chart form of the
corollary establishing them; clauses (i)--(v) of the scaffold chain at step $k$ on its pairs; the key
containment of the source-presentation statements; the source containment at $(\mathrm{L})$;
condition (M2) of Lemma~\ref{4.22:lem-majorant-transfer} at $(\mathrm{L})$ at amplitude one;
clauses (a), (b) of Lemma~\ref{4.22:lem-kept-holomorphy} at $(\mathrm{L})$. \emph{Established by:}
the two lines of that corollary; the orbit lemma; the two theorems of the scaffold chain; the
source-presentation statements; the lemma on the families $\sigma_0$. \emph{Reads:} row
$\mathrm{O}_4$; row $\mathrm{O}_0$ (the part of the scaffold chain in $(\mathbf{U},\mathbf{J})$ does
not involve the slot variable).
\item[$\mathrm{O}_6$] \emph{Item:} the treated-component lemma and its form for the led rotations
(the rotations with profiles $\vartheta^{3F}(s)$) at step $k$; its
source-blind form at step $k$, for every component $C$ of $Z$ (the split of classes comes later);
clause (c) of Lemma~\ref{4.22:lem-kept-holomorphy} and clause (iv) of
Lemma~\ref{4.22:lem-kept-oscillation} for the members of the older large-field regions $Y\subset Z$.
\emph{Established by:} the proof of the treated-component lemma and of its source-blind form, whose
inputs are exactly (i1)--(i7) below; clause (c) of Lemma~\ref{4.22:lem-kept-holomorphy} and clause
(iv) of Lemma~\ref{4.22:lem-kept-oscillation} by their proofs. \emph{Reads:} the inputs (i1)--(i7),
which lie in rows
$\mathrm{O}_0$, $\mathrm{O}_3$, $\mathrm{O}_5$ or are run-free; of the radius accounting, the radius
current at $(\mathrm{L})$,
\begin{equation}\label{4.22:eq-site-order-2}
\varepsilon^{(\mathrm{L})}:=\varepsilon_k-\eta^T_k-\mu_k\mathsf{e}/\varepsilon_{\star}-\eta^P_k
\in\bigl[\tfrac34\mathsf{e},\,\mathsf{e}\bigr],
\end{equation}
and the number $\eta_0$ of the presentation and its further number, which we denote
$\mu^{(\mathrm{L})}$; for clause (iv) of
Lemma~\ref{4.22:lem-kept-oscillation}, in addition, for a level $h<k$ as in that clause, the norms
$\mathfrak{O}(v)$ of members kept by
regions with $k_W\le h$, that is, the source-tube bounds at steps at most $h$, which are row
$\mathrm{O}_1$ of earlier steps.
\item[$\mathrm{O}_{7a}$] \emph{Item:} the presentation site $(\mathrm{L})$: the presentations and
their pieces; $V(Y)$; for every component $C$ of $Z$ the list of the $V^{\sharp}(C^{\sharp})$ and
the form $\mathfrak{q}_C$. \emph{Established by:} clause (a) and the first three parts of clause (b)
of the presentation proposition; the chart corollary of the presentation;
Lemma~\ref{4.22:lem-majorant-transfer} with its companion clause; the lemma on the presentation's
number $\mu^{(\mathrm{L})}$; the law-site rule; the two propositions on kept members at the
presentation site.
\emph{Reads:} rows $\mathrm{O}_5$, $\mathrm{O}_6$; the correlation theorem's statement for the led
rotations and its
statements at the presentation site, namely clauses (c)--(e) of its proposition on that site, a
containment statement, and its two statements on the families $\sigma$; for the wrapped lineage,
its two presentation-site statements,
the presentation-site part of clause (r3) (with its starred bound and the depth statement) and clause
(r4); for members
kept by older regions, the source-tube bounds at steps before $k$ and clauses $(\mathrm{ii})_2$,
(iv) of Lemma~\ref{4.22:lem-kept-oscillation}.
\item[$\mathrm{O}_{7b}$] \emph{Item:} the source-tube bounds at step $k$: the form clause, (q0),
(q1) and the two-point bound for every component $C$ of $Z$ at $\mathbf{R}_k$. \emph{Established
by:} the statement of the form clause with its rider list; the base statement of the two-point bound.
\emph{Reads:} for the form clause, (q0), (q1): the fixed list and rider list of those clauses
together with the geometry of step $k$, which are run-free; for the two-point bound: the lists of row
$\mathrm{O}_{7a}$, the units' $\mathfrak{N}$, $\mathfrak{O}$ and labels; the source-blind form of
the treated-component lemma (row $\mathrm{O}_6$); the source-tube bounds at steps $m'<k$, through
clauses (o), (iv) of Lemma~\ref{4.22:lem-kept-oscillation}.
\item[$\mathrm{O}_{7c}$] \emph{Item:} the modification statement (its second clause reads the form
clause at $k_W=k$, in this step); the formulas \cite[(1.90)--(1.99)]{B-CMP122b}: the terms
$\mathbf{R}'^{(k)}$, $\mathbf{B}'^{(k)}$; first and second class; $\mathfrak{K}(Y_i)$; the new
frames; the closings. \emph{Established by:} the last two parts of clause (b) of the presentation
proposition; clauses (ii), (v) of the frame proposition; clause (r1) of the wrapped lineage.
\emph{Reads:} rows $\mathrm{O}_{7a}$, $\mathrm{O}_{7b}$.
\item[$\mathrm{O}_8$] \emph{Item:} $(\mathrm{g})_k$. \emph{Established by:} row
$\mathrm{O}_{7c}$. \emph{Reads:} rows $\mathrm{O}_{7a}$--$\mathrm{O}_{7c}$.
\end{itemize}
The inputs of the treated-component lemma at step $k$ are the following.
\begin{itemize}
\item[(i1)] The induction hypothesis of the correlation theorem: the two-sided flow bounds for pairs
at most $k$ (row $\mathrm{O}_0$); its chart form at step $k$, which is the induction hypothesis at
level $k$ together with clauses $(\mathrm{b})_k$ and $(\mathrm{c})_k$ (rows $\mathrm{O}_3$,
$\mathrm{O}_5$): every $\mathfrak{E}\in\mathcal{C}$ has an orbit datum of the current radius
$\varepsilon$, $\frac34\mathsf{e}\le\varepsilon\le\mathsf{e}$, and is read through the one number
$\mathfrak{N}^{\mathbb{C}}_{\varepsilon}(\mathfrak{E})$. Clause $(\mathrm{c})_k$ is a true superset
established earlier: no chain term of step $k$ is deep; the lemma reads $(\mathrm{c})_{<k}$
through the induction hypothesis at level $k$.
\item[(i2)] The scaffold chain and the source-free scaffold maps (rows $\mathrm{O}_0$,
$\mathrm{O}_5$): $(\mathcal{U}_q^{\mathfrak{b}},\mathcal{J}_q^{\mathfrak{b}})$,
$\mathfrak{Q}^{\mathfrak{b}}$, $h_q^{\mathfrak{b}}$, with the seven frozen facts of the chain.
\item[(i3)] Containment at $(\mathrm{L})$ (row $\mathrm{O}_5$): the key containment, clause (i) of the
scaffold chain (each entry at most $1/16$ of its threshold) and its clause (ii) with the estimates
of its proof, together with the inheritance and arrest containments; at families with
$\sigma_0\cap\Lambda=\emptyset$, where $\Lambda$ is the region so denoted in the containment
statements listed here: the lemma on the families $\sigma_0$, the source admissibility
statement (ii), the source containment at $(\mathrm{L})$ with its six clauses; the geometric
clause (g) and the depth clause of the depth statement.
\item[(i4)] The base of the presentation, which is run-free: its five clauses and its
localisation statement in the source variable; its sharp lemma; the layer theorem for $\mathbb{H}$
(clauses (a), (b), (f) and the six steps of its proof); the rider statement; clauses (a), (c) of
the assignment statement and its bending statement; clause (b) of the second localisation
statement; the kernel lemma, the sharpened fourth lemma and the thirteenth law-site
statement of the scaffold; and, of the source-presentation statements, one lemma of that group, its
margin condition, its two bounds, its three clauses on the families $\sigma_0$, its condition at the
bump and its bump lemma, its restricted-source statement and its second ceiling clause.
\item[(i5)] Statements of this chapter: $\operatorname{pol}$ and its first property; the definition
of orbit data; the chart lemma; clause (a) of the presentation proposition (a device which reads one
unit's datum); item $(\beta)$ of clause (i) of the scaffold lemma at $(\mathrm{L})$; the second of
the two containment
facts of this chapter on which Lemma~\ref{4.22:lem-kept-holomorphy} rests; the numbers $\eta_0$,
$\mu^{(\mathrm{L})}$ of the presentation; the radius
theorem at steps at most $k$ (row $\mathrm{O}_0$), that is, the radius accounting up to this point.
\item[(i6)] From the increment statements: the linearisation statement; the Schwarz step of clause
(b) of the sharpened datum lemma, with the hypothesis on which that step rests; the statement on
values and
$\sigma$-pieces; the companion theorem of that group; and the ceiling of the increment statements
which serves in (i7) as the ceiling (c3) of the treated-component lemma.
\item[(i7)] Ceilings, run-free and in sup form: the ceiling of the treated-component lemma, read as a
ceiling on the supremum over the parameter polydisc $\mathfrak{P}$ of $\|J_3\|_1$ plus
$4K_{\sharp}\mathsf{a}_{\star}$, where $J_3$, $K_{\sharp}$ and $\mathsf{a}_{\star}$ are the
quantities so denoted in the treated-component lemma; the ceilings (c0)--(c2) of the
treated-component lemma and its ceiling
(c3), which is the ceiling of the increment statements named at the end of (i6); clause (ix) of the
live
ceiling; the floors on $\kappa_1$, on the
level, the base floor and the floor on $M_1$. No new condition enters.
\end{itemize}
The treated-component lemma reads no length, no label, no accumulated family of stored terms,
no sum; no right member; no output of
$(\mathrm{L})$ at step $k$; nothing of $(\mathrm{g})_k$; no row of the fifth hypothesis in the list
of hypotheses of the induction; no
source-tube bound; and no part of \cite[(1.90)--(1.99)]{B-CMP122b}.

Then, in the induction described above, every item reads only items strictly earlier in the order
\eqref{4.22:eq-site-order-a}, and run-free statements. The order is well founded: there is no circle.
\end{lemma}

\begin{proof}
We go through the rows and show, for each, that what could point forward does not.

\emph{Row $\mathrm{O}_0$.} The couplings $g_0,\dots,g_k$ are fixed before the $\mathbf{R}$-site of
step $k$ is performed; this is the ordering clause of the coupling-fixing statement. In the run of
the correlation theorem, the two-sided flow bounds for pairs at most $k$ are established at the stage
that precedes the fluctuation step. They are obtained from a derivative, at the trivial
configuration, of the frozen whole-lattice family, and this derivative needs only data of levels at
most $k-1$: this is the level-independence statement of the increment, by which the increment of
index $k+1$ needs only data of levels at most $k$, so that the increment of index $k$ needs only data
of levels at most $k-1$. That stage comes before the
fluctuation step with domains and before the $\mathbf{R}$-step. The scaffold lemma, its clause (iii)
for the wrapped lineage and the localisation statement on the scaffold maps are statements about the
maps
$\mathfrak{u}_p$, $\mathfrak{u}^{(l)}$, $\mathfrak{Q}$, $\mathfrak{w}_s$, $h_q$, which are
functions of $((\mathbf{U},\mathbf{J}),\sigma,\mathsf{s},t,\mathsf{B},b)$ built from the
source-free scaffold, and about distances in the label; no value, orbit datum or bound of a stored
term occurs in them. The radius theorem reads these statements, the flow bounds at indices at most
$k$, and the live ceiling. Hence the radii of step $k$ are known before any term of step $k$ is
bounded. The list declaration of the leaves is a declaration about the scaffold, and so belongs to
row $\mathrm{O}_0$.

\emph{Row $\mathrm{O}_2$.} The restricted form of clause (iv) of
Lemma~\ref{4.22:lem-axialising-map} at $(\mathrm{T})$ cites clauses (i) and (ii) at $(\mathrm{T})$
of the scaffold lemma; these are geometry, hence row $\mathrm{O}_0$. The fluctuation-site part of
clause (r3) of the wrapped lineage reads, of the closure statement, the datum of the leaves, which
is row $\mathrm{O}_1$; it reads Lemma~\ref{4.22:lem-axialising-map} at $\mathsf{t}=0$, which is
among the inputs of this row, and the radius accounting up to this site, which is row
$\mathrm{O}_0$.
Lemma~\ref{4.22:lem-kept-oscillation} reads the source-tube bounds at $k_W\le k-1$; since the
fluctuation site of a step precedes its $\mathbf{R}$-site, these are row $\mathrm{O}_1$.
Lemma~\ref{4.22:lem-kept-holomorphy} at $(\mathrm{T})$ reads the correlation theorem's homogeneity
statement, clause (ii) of the domain statement, the first of those two containment facts, the lemma
on the
families $\sigma_0$ and clause (iii) of the scaffold lemma for the wrapped lineage; all of these are
statements about $(\mathbf{U},\mathbf{J})$ and the scaffold maps, in which the slot variable does not
occur. Proposition~\ref{4.22:prop-fluctuation-site} reads these statements, the counts of the label
(row $\mathrm{O}_0$) and the numbers $\mathfrak{O}(v)$ (row $\mathrm{O}_1$).

\emph{Row $\mathrm{O}_4$.} The verification of the stage condition at stage $l$ concerns
$(\mathbb{U},\mathbb{J})$ and the families $\Phi_{t,\sigma}(b)$ (the inheritance theorem's
background map) and
$\mathcal{W}_T(t,\tau)$ alone, in which the slot variable does not occur. Stage $l$ reads the norms
$\|\mathbb{C}^{(i)}\|$ of the chain terms of the stages $i<l$, which precede it within the same row.
Lemma~\ref{4.22:lem-kept-holomorphy} at $(\mathrm{P})$ and
Proposition~\ref{4.22:prop-preparatory-stage} are treated as at row $\mathrm{O}_2$: they read
statements about the background pair and the scaffold maps, in which the slot variable does not
occur, the counts of the label (row $\mathrm{O}_0$) and the numbers $\mathfrak{O}(v)$ (row
$\mathrm{O}_1$); in addition they read the localisation corollary of the majorant statements, the
third flat bound of the inheritance group (the first is the bound displayed in
Definition~\ref{4.22:def-induction-hypothesis}), and the stage restriction.

\emph{Row $\mathrm{O}_5$.} The membership statements of the scaffold chain do not involve the slot
variable; its clauses on live arguments read clause $(\mathrm{c})_k$. Condition (M2) of
Lemma~\ref{4.22:lem-majorant-transfer} at $(\mathrm{L})$ and clauses (a), (b) of
Lemma~\ref{4.22:lem-kept-holomorphy} at $(\mathrm{L})$ read the key containment of the
source-presentation statements, the lemma on the families $\sigma_0$, the first of the two
containment facts, and the linearisation statement; none of them reads a term of the presentation
site $(\mathrm{L})$ of step $k$.

\emph{Row $\mathrm{O}_6$.} Each of the inputs (i1)--(i7) lies in row $\mathrm{O}_0$, $\mathrm{O}_3$
or $\mathrm{O}_5$, or is run-free. The rotation on the disc of radius $\rho_X$ in the parameter
$\lambda$ is compared with $\mu^{(\mathrm{L})}$, which is
a function of $(g_k;\mathsf{e}(s))$, and with $\eta_0$, which is a supremum of scaffold maps; both
are numbers of row $\mathrm{O}_0$. By its list of inputs, the treated-component lemma reads
no length, no label and no sum; therefore it can be cited in the passage from the relative statement
at step $k$ to the relative statement at step $k+1$, and in the closure of the induction at step $k$,
without a loop. The source-blind form of the treated-component lemma is the chain of implications
of that lemma's proof read without the source variable $\zeta$, in eight steps; it has the same
inputs. Clause (c) of Lemma~\ref{4.22:lem-kept-holomorphy} reads clause (i) of the scaffold chain
(row $\mathrm{O}_5$), a lemma and the amplitude clause of the treated-component lemma, both listed
among the hypotheses of Lemma~\ref{4.22:lem-kept-holomorphy} for its clause (c), and the second of
the two containment facts. Clause (iv) of
Lemma~\ref{4.22:lem-kept-oscillation} reads, moreover, the bounds
$\mathfrak{O}(v)\le\mathfrak{o}(v)$ for large-field regions with $k_W\le h<k$, that is, the
source-tube bounds of earlier steps.

\emph{Row $\mathrm{O}_{7a}$.} The depth statement and the two statements of the correlation theorem
at the presentation site on the families $\sigma$ read the treated-component lemma piece by piece
(row $\mathrm{O}_6$) and counts of the label (row $\mathrm{O}_0$). The correlation theorem's
statement for the led rotations reads clause (ii) of the scaffold lemma, a geometric statement on the
label, and the ceiling of the increment statements named in (i6), all of row $\mathrm{O}_0$. The
two propositions on kept
members at the presentation site read clauses $(\mathrm{ii})_2$ and (iv) of
Lemma~\ref{4.22:lem-kept-oscillation} at large-field regions with $k_W<k$.

\emph{Row $\mathrm{O}_{7b}$.} The form clause, (q0) and (q1) of the source-tube bounds read the
fixed list of those clauses, with its rider list, together with the geometry of step $k$, through the
run-free clause of the statement of the form clause; they read no object of the run. The two-point
bound at step $k$ reads the lists of row $\mathrm{O}_{7a}$ and the units' $\mathfrak{N}$,
$\mathfrak{O}$ and labels, which is why the source-tube bounds of step $k$ come after the
presentation site; it reads the source-blind form of the treated-component lemma (row
$\mathrm{O}_6$), and the source-tube bounds at steps $m'<k$, through clauses (o) and (iv) of
Lemma~\ref{4.22:lem-kept-oscillation}. Row $\mathrm{O}_{7a}$ reads the source-tube bounds at steps
$m'<k$ only. Hence there is no circle between rows $\mathrm{O}_{7a}$ and $\mathrm{O}_{7b}$.

\emph{Row $\mathrm{O}_{7c}$.} The second clause of the modification statement reads the form clause
at $k_W=k$, which is established in row $\mathrm{O}_{7b}$ of the same step. That the source-tube
bounds of step $k$ are used only at steps after $k$ holds for the two-point bound and for what
Lemma~\ref{4.22:lem-kept-oscillation} and the propositions on kept members read, but not for the
kept Gaussian $\mathsf{Q}_i$ on the correlation theorem's own row; this is harmless, since the form
clause is run-free. A frame created at a step is first read at a later step, because the law-site
rule speaks of the chart variable; so Lemma~\ref{4.22:lem-kept-oscillation} and the propositions on
kept members use the source-tube bounds of step $k$ at steps after $k$ only. The closings of clause
(v) of the frame proposition read the class of the component, which is decided in row
$\mathrm{O}_{7c}$.

\emph{Row $\mathrm{O}_8$.} Row $\mathrm{O}_8$ of step $k$ is row $\mathrm{O}_1$ of step $k+1$.

\emph{Ceilings.} The live ceiling, the ceiling of the increment statements named in (i6) and one
further ceiling on $\gamma$ of the same kind, the half
ceilings on members and on arrests, the relative ceilings at $(\mathrm{T})$ and at $(\mathrm{L})$,
the flat ceiling of the inheritance group, the tail, interior and kept ceilings, the ceiling of the
treated-component lemma and the ceilings (c0)--(c2) are inequalities on $\gamma$; the function
$\Sigma^{\mathfrak{q}}$ is a function of the coupling; and $C_\Sigma$, $p_\Sigma$,
$\mathfrak{n}^{\mathrm{kept}}$ are numbers fixed before $\gamma$. All of them are in sup form: they
read no object of the run, and may be cited at any row.
\end{proof}

\begin{remark}
The order \eqref{4.22:eq-site-order-a} is also the order in which the right member of a term is
copied: a display at row $\mathrm{O}_i$ of step $k$ is computed by a majorant from displays at rows
before $\mathrm{O}_i$ or at steps before $k$.
\end{remark}

\begin{proof}[Proof of the orbit-chart theorem, by induction on the level]
\label{4.22:prf-orbit-chart-proof}
The hypotheses are those of the orbit-chart theorem in the form \eqref{4.22:eq-two-implications-a}:
the clauses of the correlation theorem; its led-rotation bound for brackets on the disc of its
Cauchy variable $\mathsf{t}$; the deep-unit lemma and its led-rotation bound. For the class
$\mathfrak{S}$ they include, in addition, the relation bound for the $\mathbf{B}'$-terms, taken in the
norm $\mathfrak{N}^{\mathbb{C}}$, and the strengthened tube statement for the kept members. We argue by
induction on the level $k$. Inside a step, the assertions are established in the order of sites of
Lemma~\ref{4.22:lem-site-order} (displayed in \eqref{4.22:eq-site-order-a}). In that order every
assertion uses only assertions that precede it strictly, together with statements and constants
that involve no term bounded in the course of the induction.

\emph{The level $k=1$.} The step at this level is the one of \cite[\S1]{B-CMP119}. It is the
fluctuation site at index $0$. The orbit-chart form of the fluctuation-integral theorem applies to it
in the case in which no older law site is present.

\emph{From the induction hypothesis at level $k$ to clause \textup{(b)}.} In the order of sites this
part runs from the induction hypothesis at level $k$
(Definition~\ref{4.22:def-induction-hypothesis}) through the large-field stages to the site of clauses
(b) and (c), where it yields clause (b). We argue by induction on the stage $n$, in the way the
inheritance theorem does.

The proof that the inheritance hypothesis holds at stage $n$ consists of four steps. These steps
concern only the pair $(\mathbb{U},\mathbb{J})\in\mathfrak{D}^{\flat}_n(Y)$ and the families
$\Phi_{t,\sigma}(b)$ and $\mathcal{W}_T(t,\tau)$ of the inheritance theorem. Neither a slot variable,
nor the axialising map $\mathfrak{u}$, nor a rotation $g$ enters any of them.

Two kinds of input to the domain list $\mathfrak{I}_n$ of stage $n$ need a word.
\begin{itemize}
\item The stored terms that carry a frame enter $\mathfrak{I}_n$ on their stored background domain.
This is clause (iv) of the frame lemma (Lemma~\ref{4.22:lem-frame-lemma}).
\item For every live large-field region $Y$, the kept members $\mathfrak{K}(Y)$ enter
$\mathfrak{I}_n$ as carriers of clauses (Definition~\ref{4.22:def-kept-domain-list}).
\end{itemize}
Consequently the second step, and part $(\alpha)$ of the fourth step, of the proof of the inheritance
hypothesis at stage $n$ deliver, for each kept member $v$, the half of its holomorphy statement that
concerns the field variables. They deliver it at every strip other than $v$'s own.

With the inheritance hypothesis at stage $n$ in hand, Corollary~\ref{4.22:cor-preparatory-chart}
carries the stage lemma of the large-field chain over to the orbit chart. By it, each
$\mathbb{C}^{(n)}_k(X,\mathcal{S})^{\sharp}$ is an orbit datum on
\begin{equation*}
\mathfrak{D}^{\flat}_n(X)\times\mathfrak{D}_{\mathcal{S}}(X)
\times\operatorname{supp}\boldsymbol{\Phi}\times\mathfrak{G}_{\varepsilon}(X).
\end{equation*}
The total bounds of the stage take their inputs through two quantities. The first is the weighted
bound \eqref{4.22:eq-induction-hypothesis-a}. The second consists of the norms
$\|\mathbb{C}^{(i)}\|$, $i<n$,
now taken in the norm $\mathfrak{N}^{\mathbb{C}}$. The closing inequality of these totals is
$3c_G+\vartheta_{\mathrm{tot}}+\tfrac12\le C_0^{\sharp}$, in which $c_G$ and
$\vartheta_{\mathrm{tot}}$ are constants entering the total bounds and $C_0^{\sharp}$ is the constant
that bounds those totals. By induction on $n$ this gives clause (b).

Two remarks complete this description.
\begin{itemize}
\item The two terms of the chain denoted $\mathbb{C}_G$ and $\Delta W^{\mathrm{ch}}$ take no stored
term as input. For them $\mathcal{S}=\emptyset$, and no rotation enters.
\item The total bounds are the anchored sums of the large-field stages themselves, in their own
notation. For the plain class $\mathfrak{A}$ they are the anchored sums of the preparatory stages,
with their restricted form at the stage. For the class $\mathfrak{S}$ they are the large-field-stage
part of the descendant bound of the wrapped lineage, taken on labels. The kept members $\mathfrak{K}$,
and the families $\mathfrak{Q}(\mathfrak{a})$ attached to the live arrests $\mathfrak{a}$ of the
chain, enter these sums through their oscillation, by the oscillation form of the stage bound: the
kept members through Lemma~\ref{4.22:lem-kept-oscillation}, the families
$\mathfrak{Q}(\mathfrak{a})$ through clause (iii) of Proposition~\ref{4.22:prop-gaussian-oscillation}.
\end{itemize}

A chain arrested at stage $n$ leaves its terms $\mathbb{C}^{(i)}$, $i\le n$, with these orbit data
\cite{B-CMP122a}. If the arrest is by \cite[(1.55)]{B-CMP122a}, the composite terms of the arrest and
its new terms also have orbit data, by clause (ii) of the arrest-composite lemma and the new-terms
lemma at the large-field stages; the families $\mathfrak{Q}(\mathfrak{a})$ have theirs by clause (i)
of Proposition~\ref{4.22:prop-gaussian-oscillation}. This proves clause (h) for the terms produced at
the large-field stages.

\emph{From clause \textup{(b)} to clause \textup{(c)}.} In the order of sites this takes place at
the position of clauses (b) and (c). It rests on the two lines of the terminal corollary of the
large-field stages.

The first line is the inclusion $\mathfrak{X}^{\sharp}\subset\mathfrak{D}_{\mathcal{S}}$ for every
$\mathcal{S}$, the domain inclusion of the large-field stages. The second line is the chain
\begin{equation}\label{4.22:eq-orbit-chart-proof-1}
\begin{split}
\Bigl|\sum_{\mathcal{S}}\mathbb{C}(X,\mathcal{S})^{\sharp}\Bigr|
&\le\sum_{\mathcal{S}}e^{\kappa_A|\mathcal{S}|}\,
\mathfrak{N}^{\mathbb{C}}\bigl(\mathbb{C}(X,\mathcal{S})\bigr)\\
&\le\|\mathbb{C}^{(n)}\|^{A,\flat,\mathbb{C}}_{(n)}\,e^{-a_{\ell}\mathsf{d}_{\ell}}.
\end{split}
\end{equation}
Here $\kappa_A$ is the weight $\kappa_A(k)$ of the induction hypothesis, $\mathsf{d}_{\ell}$ is the
assigned length of the chain term at its level $\ell$, and $a_{\ell}$ is the decay rate attached to that
length. A finite sum of orbit data is again an orbit datum.

Two facts about the radii are used here.
\begin{itemize}
\item For a complex rotation $g\in G^{\mathbb{C}}$, the bound
$|\log g|\le\varepsilon_{\mathrm{live}}$ implies
$\operatorname{pol}(g)\le\varepsilon_{\mathrm{live}}$. Indeed $g=\exp(\log g)$ gives
$\|g\|\le e^{|\log g|}$ and $\|g^{-1}\|\le e^{|\log g|}$, so that
\begin{equation*}
\operatorname{pol}(g)=\log\max\{\|g\|,\|g^{-1}\|\}\le|\log g|.
\end{equation*}
\item The difference satisfies $\varepsilon_k-\varepsilon_{\mathrm{live}}\ge\tfrac12\varepsilon_{\star}$,
by the budget inequalities \eqref{4.22:eq-budget-closes-a}.
\end{itemize}

\emph{From clause \textup{(c)} to clauses \textup{(g)} and \textup{(h)}.} In the order of sites this
runs from clauses (b) and (c), through the site of the deep-unit lemma, through the presentation and
the tube statements of the step, to the output of the step. At the presentation site and at the
output of the step, clause (c) gives clauses (g) and (h). This follows from the presentation
proposition and the frame lemma (Lemma~\ref{4.22:lem-frame-lemma}).

The pieces of deep units are supplied by the deep-unit lemma, at its own site. The pieces of brackets
on the disc of the correlation theorem are supplied by the led-rotation bound of that theorem. Both
statements are among the hypotheses listed above.

\emph{Clause \textup{(j)}.} This clause belongs to the first position in the order of sites. The
radius-budget theorem at step $k$ (see \eqref{4.22:eq-budget-closes-a}) takes only two inputs. The
first is the running-shift statement, in both its forms, at indices at most $k$. The second is the
geometry of the steps at most $k$.

\emph{From clause \textup{(g)} to clause \textup{(a)} and the induction hypothesis at level $k+1$.}
In the order of sites this runs from the output of step $k$ to the first sites of step $k+1$. These
are the terms of the large-field operation of step $k$, which are the output of that step, then the
fluctuation integral of step $k+1$, and clause (a) together with the induction hypothesis.

At the next fluctuation site three things are used.
\begin{itemize}
\item The orbit-chart form of the fluctuation-integral theorem is applied there.
\item The stored terms of older steps are carried through that site by
Corollary~\ref{4.22:cor-older-law-sites}.
\item The new slot enters at the radius $\varepsilon_{\star}$.
\end{itemize}
The right members of the bounds are copied unchanged from the bounds for the same terms at real
configurations. For
the class $\mathfrak{S}$ they are taken from the birth bound and the descendant bound of the wrapped
lineage.

\emph{The weight.} The fluctuation-integral theorem gives the weight $\kappa_A^{\mathrm{pr}}(g_k)$ on
the sets of the new slot. Clause (d) of Corollary~\ref{4.22:cor-older-law-sites} carries the older
sets at the weight $\kappa_A(k)$. Together these give the display of clause (a) at level $k+1$. The
sizes enter as the fluctuation-site summation statement of the correlation theorem allots them:
\begin{equation*}
\tfrac14+2\cdot\tfrac18+\tfrac14\le1,
\end{equation*}
in units of that statement's constant $B_0$. This is exactly the way in which the weighted
corollary of the fluctuation step at
real data delivers that display. Every later use takes the weight $e^{\kappa_A(k+1)|\mathcal{S}|}$.
This weight is dominated by the weight just obtained, because $\kappa_A$ is non-increasing.

\emph{Clause \textup{(i)}.} The axialising map lemma (Lemma~\ref{4.22:lem-axialising-map}) is applied
at each of the sites above.
\end{proof}

\begin{remark}[The orbit-chart theorem as two implications, one per class]\label{4.22:rem-two-implications}
The terms carried through the steps fall into two classes:
\begin{itemize}
\item the \emph{plain class} $\mathfrak{A}$;
\item the class $\mathfrak{S}$ of terms whose lineage has been touched by a second-class outcome, the
\emph{wrapped} terms.
\end{itemize}
For a term of $\mathfrak{S}$, its \emph{first wrapped ancestor} is the first ancestor along its
lineage that lies in $\mathfrak{S}$. Consider the following six statements.
\begin{itemize}
\item[(h1)] The rotation bound for the bracket pieces that the correlation theorem takes on its disc.
\item[(h2)] The rotation bound for the pieces of the plain deep units.
\item[(h3)] The third-factor lemma. This is the bound, piece by piece, of one half of the third
factor of the exponential that remains after \cite[(1.67)]{B-CMP122b} has been extracted.
\item[(h4)] The rows of the correlation theorem.
\item[(h5)] The relative bound for Ba{\l}aban's terms $\mathbf{B}'$ of \cite{B-CMP122b}, read in the
rotated sup norm $\mathfrak{N}^{\mathbb{C}}$.
\item[(h6)] The strengthened source-tube bounds. Among them are the two-point bound and the two
bounds on the kept quadratic form $\mathsf{Q}_i$.
\end{itemize}
Then the statement proved as the orbit-chart theorem is the conjunction of two implications, one
for each class:
\begin{equation}\label{4.22:eq-two-implications-a}
\begin{aligned}
&\bigl[(\mathrm{h}1)\wedge(\mathrm{h}2)\wedge(\mathrm{h}3)\wedge(\mathrm{h}4)\\
&\qquad\;\Longrightarrow\;\text{the theorem's conclusion for the terms of }
\mathfrak{A}\bigr]\\
&\wedge\;\bigl[(\mathrm{h}1)\wedge(\mathrm{h}2)\wedge(\mathrm{h}3)\wedge(\mathrm{h}4)
\wedge(\mathrm{h}5)\wedge(\mathrm{h}6)\\
&\qquad\;\Longrightarrow\;\text{the theorem's conclusion for the terms of }
\mathfrak{S}\bigr].
\end{aligned}
\end{equation}
\end{remark}

\begin{proof}
\emph{The split is exact.}
Consider the induction on the level by which the orbit-chart theorem is proved. Two facts govern
the bounds it produces.
\begin{itemize}
\item The right member of the bound of a term is copied from the right members of its constituents,
through the majorant of the site at which the term is formed. This copying is supplied by the fourth
majorant relation of the sites, by the fourth and fifth steps of the proof of the flat form of the
fluctuation theorem, by the totals of the large-field stages, and by the third item of part~(b) of
the proposition used in passing from clause~(c) to clause~(g).
\item The radius of a term is lowered along its own lineage, by the radius recursion of the
induction.
\end{itemize}
As recorded in the remark on the two classes, a term of $\mathfrak{A}$ has no constituent in
$\mathfrak{S}$.
So at every site of the order \eqref{4.22:eq-site-order-a}, the display of a term of $\mathfrak{A}$
reads only two kinds of input:
\begin{itemize}
\item the displays of terms of $\mathfrak{A}$, taken on the plain radius $\varepsilon^{\mathfrak{A}}_k$;
\item the rows of the correlation theorem, that is, (h4).
\end{itemize}
A term of $\mathfrak{S}$ reads three further kinds of input, each at a definite place along its
lineage:
\begin{itemize}
\item the two relative rows from which the right members of wrapped terms are copied, together with
the relative bound (h5), at its first wrapped ancestor;
\item the strengthened source-tube bounds (h6), at the first ancestor that lies among the kept
members $\mathfrak{K}$;
\item the transfer statement for later sites, at every ancestor after that.
\end{itemize}
Hence the plain class is closed under the induction, with inputs (h1)--(h4) only. The wrapped class
needs, in addition, exactly the inputs read at these places. Those inputs are identified below.

\emph{The inputs of the first implication, and where each is used.}
None of (h1)--(h4) is idle.

(h2) and (h3) are used together at the site $\mathrm{O}_7$ of the order
\eqref{4.22:eq-site-order-a}, for the plain deep units. In \cite{B-CMP122b}, the quantity that sums
the domains $X$ inside a core into the constant $\alpha$ of \cite[(1.68)]{B-CMP122b} is the
exponential kept after \cite[(1.67)]{B-CMP122b} has been extracted, namely, in the notation of
\cite[p.~376]{B-CMP122b},
\begin{equation*}
\exp\Bigl(-\beta\kappa d_j(X)
-\tfrac12\delta\, d\bigl((Y_1^0)^c\cap Y_1\cap Y'_{i_0},X\bigr)
-\sum_i\tfrac12(\kappa_1-1)M^{-d}|Y'_i|\Bigr).
\end{equation*}
Half of its third factor performs this summation, precisely in the
following way.
\begin{itemize}
\item It sums the core units whose chain exits, that is, those for which some $Y'_i$ is non-empty.
\item The core units for which every $Y'_i$ is empty and whose domain lies in a component are not
summed by this factor. They enter \cite[(1.69)]{B-CMP122b} by count.
\end{itemize}
In this book that factor is bounded piece by piece by the third-factor lemma (h3). The sum over the
pieces is the summation of the plain pieces over a core.

This bound is established at clause~(g) of the orbit-chart theorem for the plain class. From there
it passes, at every later step, to:
\begin{itemize}
\item clause~(a);
\item the induction hypothesis \eqref{4.22:eq-induction-hypothesis-a};
\item clause~(b);
\item clause~(c) for the plain class.
\end{itemize}

(h1) is also used at the site $\mathrm{O}_7$. It bounds every bracket piece that the correlation
theorem takes on its disc, namely:
\begin{itemize}
\item the plain units of the second part;
\item the plain units of the first part leaving $Z$.
\end{itemize}
Here the first and second parts are those into which the units entering the correlation theorem's
brackets are divided.

(h4), the rows of the correlation theorem, is used at every site.

\emph{What the second implication uses beyond the first.}
The relative bound (h5) is used at the first wrapped ancestor.

Of the strengthened source-tube bounds (h6), the two-point bound and the two bounds on the kept
quadratic form are used at every later site. They enter in two ways:
\begin{itemize}
\item the oscillation norm $\mathfrak{O}(v)$ of the kept members $v$;
\item the datum of $\mathsf{Q}_i$.
\end{itemize}
At those sites they are read through clauses (o), $(\mathrm{ii})_2$ and (iv) of
\eqref{4.22:eq-kept-oscillation-a}. They are also read through the statement that kept terms enter
every later bound through one fixed number, Corollary~\ref{4.22:cor-fixed-number-entry}, whose proof
is incomplete at one step. Consequently the lists of hypotheses under which the source-tube bounds
are proved stand under every later row that reads a kept quadratic form $\mathsf{Q}_i$. These lists
are headed by the reality statement for the kept data.

Two further statements are used, at three places.
\begin{itemize}
\item The implication contained in the form statement is used at one place only. That place is the
second modification of the creating step, a sub-site of the site $\mathrm{O}_7$. This is a row of the
correlation theorem, not of the orbit-chart theorem.
\item A variant of the third-factor lemma is used inside the proof of the two-point bound, at a
further sub-site of the site $\mathrm{O}_7$.
\item For the kept members of the inner pockets among the pockets $\mathcal{P}_F$ of a stored term
$F$, that variant is used as clause~(iv) of \eqref{4.22:eq-kept-oscillation-a}.
\end{itemize}
The implication contained in the kept source-term statement, a further statement of this chapter,
is used nowhere.

\emph{The coupling ceilings \eqref{4.22:eq-coupling-ceilings-l} read per class.}
Consider a fluctuation site $(\mathrm{T})$. The share of the older stored units in clause~(iii) of
Corollary~\ref{4.22:cor-older-law-sites} is a sum over the units whose pieces enter the term $V(Y)$
that an output reads. That share enters the ceilings \eqref{4.22:eq-coupling-ceilings-l}.

\emph{Outputs of $\mathfrak{A}$.} For an output of $\mathfrak{A}$ the sum ranges over units of
$\mathfrak{A}$ only, since otherwise the output would lie in $\mathfrak{S}$. For the plain units, the
share is supplied by clause~(ii) of the supply statement of the correlation theorem, together with
its filing statement. The presented units $\mathfrak{Q}(\mathfrak{a})$ at an arrest $\mathfrak{a}$
are included. The cost is two inequalities on $\gamma$. The row of
$\mathfrak{A}$ in \eqref{4.22:eq-coupling-ceilings-l} is the sum of the second step of the
proof of Corollary~\ref{4.22:cor-older-law-sites}, with the constants $C_A^{(1/2)}$ and $C_A$ there.

\emph{Outputs of $\mathfrak{S}$.} For an output of $\mathfrak{S}$ the sum contains two further
families of units.
\begin{itemize}
\item[$(\alpha)$] The wrapped units. For these, the relative bound (h5) supplies the relative term
$V^{\mathrm{rel}}$ and its radius $\varepsilon^{\mathrm{rel}}_k$, under the relative ceiling at the
fluctuation site. That ceiling is part of (h5).
\item[$(\beta)$] The members of $\mathfrak{K}(Y_i)$. For these, the proposition on kept terms at a
fluctuation site, Proposition~\ref{4.22:prop-fluctuation-site}, supplies the term $V^{\mathrm{kept}}$
with
\begin{equation}\label{4.22:eq-two-implications-1}
\varepsilon^{\mathrm{kept}}_k
=C_\Sigma^{\mathrm{sum}}\,\mathfrak{n}^{\mathrm{kept}}\,
e^{-\frac14\delta\kappa(7\check R_k-2)},
\end{equation}
under the strengthened relative ceiling of \eqref{4.22:eq-fluctuation-site-a}; here
$C_\Sigma^{\mathrm{sum}}$, $\delta$ and $\kappa$ are the constants of that proposition.
\end{itemize}
The right-hand side of \eqref{4.22:eq-two-implications-1} is a fixed number times the row's own far
factor at the current scale $\check R_k$. It therefore tends to zero with the row. The two units
belonging to one pocket are added at the anchor, by the anchoring statement for kept members.
Neither the implication of the kept source-term statement nor the reference statement, a further
statement of this chapter, is used at this site. Finally, the per-unit displays of
Corollary~\ref{4.22:cor-older-law-sites}
hold whatever the list of units over which the share is summed.
\end{proof}

\begin{corollary}[Chain terms as holomorphic functions of the complex slot]
\label{4.22:cor-chain-terms-holomorphic}
Assume the following three statements:
\begin{itemize}
\item clause (c) of the orbit-chart theorem (\ref{4.22:prf-orbit-chart-proof}), proved by
induction on the level;
\item the statement that an orbit datum defines a function of the rotated slot;
\item the statement fixing the joint $\mathcal{G}^{\mathbb{C}}$-invariance of orbit data.
\end{itemize}
Let $k$, $n$ and $X$ be as in the orbit-chart theorem, and let $\Phi_A$ be a real further
argument of the chain term, lying in the supports on which the orbit datum of clause (c) is
defined. Let $\mathfrak{R}_k(X)$ denote the set of $G^{\mathbb{C}}$-valued maps $g$ on
the law sites subject to two conditions. At the pointwise law sites,
$|\log g|\le\varepsilon_{\mathrm{live}}(k)$. At the law sites of a frame $W$,
$\operatorname{pol}(g)$ is at most $\eta^{L}_k$ plus the step's margins and reserve; by the
frame lemma (Lemma~\ref{4.22:lem-frame-lemma}), this quantity is at most
$\tfrac14\varepsilon^{\mathrm{fr}}(W)$. For such $g$, $\Phi_A^{[g]}$ denotes the rotation of
$\Phi_A$ by $g$.

Then the map
\begin{equation}\label{4.22:eq-chain-terms-holomorphic-1}
\bigl((\mathbb{U},\mathbb{J}),g\bigr)\longmapsto
\mathbb{C}^{(n)}_k\bigl(X;(\mathbb{U},\mathbb{J}),\Phi_A^{[g]}\bigr)
:=\mathbb{C}^{(n),\sharp}_k\bigl(X;(\mathbb{U},\mathbb{J});\Phi_A,g\bigr)
\end{equation}
is well defined as a function of the complex slot. It is holomorphic on
$\mathfrak{D}^{\flat}_n(X)\times\mathfrak{R}_k(X)$, and it obeys
\begin{equation}\label{4.22:eq-chain-terms-holomorphic-2}
\bigl|\mathbb{C}^{(n)}_k\bigl(X;(\mathbb{U},\mathbb{J}),\Phi_A^{[g]}\bigr)\bigr|
\le C_0^{\sharp}\,e^{-a_m\mathsf{d}_m}.
\end{equation}
Here $C_0^{\sharp}$ is the constant of the orbit-chart theorem bounding the norms of the
chain terms, and $e^{-a_m\mathsf{d}_m}$ is the weight attached to $X$ in the norm of
clause (c), with decay rate $a_m$ and length $\mathsf{d}_m=\mathsf{d}_m(X)$ of $X$.
The map is jointly $\mathcal{G}^{\mathbb{C}}$-invariant in the
sense of the third statement assumed above. Finally, the set $\mathfrak{R}_k(X)$ is connected
and contains the identity.
\end{corollary}

\begin{proof}
The right member of \eqref{4.22:eq-chain-terms-holomorphic-1} is the orbit datum
$\mathbb{C}^{(n),\sharp}_k(X;\cdot)$ evaluated at the real argument $\Phi_A$ and the rotation
$g$. The left member of \eqref{4.22:eq-chain-terms-holomorphic-1} is defined to be this
value.

Clause (c) of the orbit-chart theorem provides $\mathbb{C}^{(n),\sharp}_k(X;\cdot)$ as an
orbit datum on $\mathfrak{D}^{\flat}_n(X)$ at the rotations admitted in $\mathfrak{R}_k(X)$.
It also provides holomorphy there and the bound \eqref{4.22:eq-chain-terms-holomorphic-2}.
At the frame sites, the admissible size $\tfrac14\varepsilon^{\mathrm{fr}}(W)$ is the one
supplied by Lemma~\ref{4.22:lem-frame-lemma}. For chain terms of the class $\mathfrak{S}$,
the right member of \eqref{4.22:eq-chain-terms-holomorphic-2} is the one supplied by the
large-field stages' bound on labels for that class.

The second and the third statements assumed above, applied to the orbit datum
$\mathbb{C}^{(n),\sharp}_k(X;\cdot)$, show that its value at $(\Phi_A,g)$ is a well-defined
function of the complex slot $\Phi_A^{[g]}$, which justifies the notation on the left of
\eqref{4.22:eq-chain-terms-holomorphic-1}, and that the map
\eqref{4.22:eq-chain-terms-holomorphic-1} is jointly $\mathcal{G}^{\mathbb{C}}$-invariant in
the sense of the third of them.

The identity lies in $\mathfrak{R}_k(X)$, since $\log\mathbb{1}=0$ and
$\operatorname{pol}(\mathbb{1})=0$, and both right members in the definition of
$\mathfrak{R}_k(X)$ are non-negative.

The conditions defining $\mathfrak{R}_k(X)$ are imposed site by site, so $\mathfrak{R}_k(X)$
is the product over the law sites of the sets of admissible values at each site, and it
suffices to join every admissible value at a site to $\mathbb{1}$ by a path of admissible
values. At a pointwise law site, the hypothesis $|\log g|\le\varepsilon_{\mathrm{live}}(k)$
is read with the logarithm inverse to $\exp$ on the ball of radius
$\varepsilon_{\mathrm{live}}(k)$ in $\mathfrak{g}^{\mathbb{C}}$, the hypothesis itself
presupposing that $\exp$ is one-to-one on this ball; put $Y=\log g$, so that
$g=\exp Y$ and $|Y|\le\varepsilon_{\mathrm{live}}(k)$. For $0\le t\le1$ the element
$\exp(tY)$ lies in $G^{\mathbb{C}}$, because $\operatorname{tr}(tY)=0$, and
\begin{equation*}
|\log\exp(tY)|=t|Y|\le\varepsilon_{\mathrm{live}}(k),
\end{equation*}
so $t\mapsto\exp(tY)$ is a path of admissible values from $\mathbb{1}$ to $g$. At a law site
of a frame, let $g=u\,e^{H}$ be the polar decomposition: $e^{H}=(g^*g)^{1/2}$ with $H$
Hermitian and $\operatorname{tr}H=\log|\det g|=0$, and $u=g\,e^{-H}$ is unitary with
$\det u=\det g/\det e^{H}=1$, so $u\in SU(N)$. For $0\le t\le1$, since $u$ is unitary,
\begin{equation*}
\|u\,e^{tH}\|=\|e^{tH}\|,\qquad \|(u\,e^{tH})^{-1}\|=\|e^{-tH}\|,
\end{equation*}
and the larger of the two is $e^{t\|H\|}$, because the eigenvalues of $tH$ are real and the
largest modulus among them is $t\|H\|$. Hence
\begin{equation*}
\operatorname{pol}(u\,e^{tH})=t\|H\|=t\operatorname{pol}(g)\le\operatorname{pol}(g),
\end{equation*}
so $t\mapsto u\,e^{tH}$ is a path of admissible values from $u$ to $g$. Finally $u$ is joined
to $\mathbb{1}$ by a path in the connected group $SU(N)$, along which $\operatorname{pol}$
vanishes. Thus $\mathfrak{R}_k(X)$ is connected.
\end{proof}

\begin{lemma}[The orbit lemma on a connected parameter manifold]\label{4.22:lem-orbit-lemma}
Let $\mathcal{N}\subset\mathbb{C}^{n_{\mathcal{N}}}$ be a connected open set of complex parameters
whose real slice $\mathcal{N}_{\mathbb{R}}:=\mathcal{N}\cap\mathbb{R}^{n_{\mathcal{N}}}$ has non-empty
interior in $\mathbb{R}^{n_{\mathcal{N}}}$. Let
\begin{equation*}
(\mathcal{P},\mathfrak{q})\colon \mathcal{N}\longrightarrow
\operatorname{Conf}^{\mathbb{C}}(X)\times (G^{\mathbb{C}})^{\mathfrak{s}(X)}
\end{equation*}
be holomorphic and real on $\mathcal{N}_{\mathbb{R}}$. Here $\operatorname{Conf}^{\mathbb{C}}(X)$ is the space of
complexified background pairs $(\mathbf{U},\mathbf{J})$ on $X$. Reality means that on $\mathcal{N}_{\mathbb{R}}$ the pair
$\mathcal{P}$ is real and $\mathfrak{q}$ is $G$-valued. Let $\mathfrak{D}\subset\operatorname{Conf}^{\mathbb{C}}(X)$
be open and $\mathcal{G}$-invariant. Let $\mathfrak{E}$ be a term with an orbit datum
$\mathfrak{E}^{\sharp}$ of radius $\varepsilon$ on $\mathfrak{D}$, and fix the further arguments
$(\mathbf{A},\boldsymbol{\Phi})$. Let $h\colon\mathcal{N}\to\mathcal{G}^{\mathbb{C}}(X)$ be continuous and
$G$-valued on $\mathcal{N}_{\mathbb{R}}$, and assume that for every $\zeta\in\mathcal{N}$
\begin{equation*}
\bigl(\mathcal{P}(\zeta)^{h(\zeta)},\,h(\zeta)\mathfrak{q}(\zeta)\bigr)\in
\mathfrak{D}\times\mathfrak{G}_{\varepsilon}(X).
\end{equation*}
Then the function
\begin{equation}\label{4.22:eq-orbit-lemma-1}
F(\zeta):=\mathfrak{E}^{\sharp}\bigl(\mathcal{P}(\zeta)^{h(\zeta)};\mathbf{A},\boldsymbol{\Phi},
h(\zeta)\mathfrak{q}(\zeta)\bigr)
\end{equation}
has the following four properties.
\begin{itemize}
\item It is holomorphic on $\mathcal{N}$.
\item It satisfies $|F|\le\mathfrak{N}^{\mathbb{C}}_{\varepsilon}(\mathfrak{E})$.
\item It does not depend on the choice of $h$ among maps with the properties above.
\item For $\zeta\in\mathcal{N}_{\mathbb{R}}\cap\mathcal{P}^{-1}(\mathfrak{D})$ it satisfies
\begin{equation}\label{4.22:eq-orbit-lemma-2}
F(\zeta)=\mathfrak{E}\bigl(X;\mathcal{P}(\zeta),\mathbf{A}^{[\mathfrak{q}(\zeta)]},
\boldsymbol{\Phi}^{[\mathfrak{q}(\zeta)]}\bigr).
\end{equation}
\end{itemize}
\end{lemma}

\begin{proof}
\emph{Local freezing of the gauge transformation.} Fix $\zeta_0\in\mathcal{N}$. For $\zeta$ near
$\zeta_0$ put $\varsigma:=h(\zeta)h(\zeta_0)^{-1}$. By continuity of $h$, $\varsigma$ is close to the
identity at every site of $X$. Hence $\log\varsigma$ is defined sitewise, and so is
$e^{t\log\varsigma}$ for $t\in[0,1]$. Consider the path
\begin{equation*}
t\longmapsto\Bigl(\bigl(\mathcal{P}(\zeta)^{h(\zeta_0)}\bigr)^{e^{t\log\varsigma}},\;
e^{t\log\varsigma}\,h(\zeta_0)\,\mathfrak{q}(\zeta)\Bigr),\qquad t\in[0,1].
\end{equation*}
At $t=0$ its value is
$\bigl(\mathcal{P}(\zeta)^{h(\zeta_0)},h(\zeta_0)\mathfrak{q}(\zeta)\bigr)$. At $t=1$ its value is
$\bigl(\mathcal{P}(\zeta)^{h(\zeta)},h(\zeta)\mathfrak{q}(\zeta)\bigr)$, because
$e^{\log\varsigma}h(\zeta_0)=h(\zeta)$.

The set $\mathfrak{D}$ is open by assumption. The set $\mathfrak{G}_{\varepsilon}(X)$ is open, being
defined by the strict inequalities $\operatorname{pol}<\varepsilon$ at the finitely many sites of
$X$. Hence $\mathfrak{D}\times\mathfrak{G}_{\varepsilon}(X)$ is open. At $\zeta=\zeta_0$ we have
$\varsigma=\mathbb{1}$, so the path is constant and equal to
$\bigl(\mathcal{P}(\zeta_0)^{h(\zeta_0)},h(\zeta_0)\mathfrak{q}(\zeta_0)\bigr)$, which lies in this
open set by hypothesis. The point of the path depends continuously on $(\zeta,t)$, since
$\mathcal{P}$ and $\mathfrak{q}$ are holomorphic and $h$ is continuous. As $[0,1]$ is compact, the
whole path stays in $\mathfrak{D}\times\mathfrak{G}_{\varepsilon}(X)$ for all $\zeta$ close enough to
$\zeta_0$.

For such $\zeta$ we apply the path-invariance property in the definition of an orbit datum: the
value of $\mathfrak{E}^{\sharp}$ is unchanged along a compact path of this form that stays in
$\mathfrak{D}\times\mathfrak{G}_{\varepsilon}(X)$. It gives
\begin{equation}\label{4.22:eq-orbit-lemma-3}
F(\zeta)=\mathfrak{E}^{\sharp}\bigl(\mathcal{P}(\zeta)^{h(\zeta_0)};\mathbf{A},\boldsymbol{\Phi},
h(\zeta_0)\mathfrak{q}(\zeta)\bigr).
\end{equation}
In \eqref{4.22:eq-orbit-lemma-3} the gauge transformation $h(\zeta_0)$ is fixed. The map
$\zeta\mapsto\mathcal{P}(\zeta)^{h(\zeta_0)}$ is the holomorphic map $\mathcal{P}$ followed by the
action of a fixed element of $\mathcal{G}^{\mathbb{C}}(X)$. The map
$\zeta\mapsto h(\zeta_0)\mathfrak{q}(\zeta)$ is holomorphic. The function $\mathfrak{E}^{\sharp}$ is
holomorphic on $\mathfrak{D}\times\mathfrak{G}_{\varepsilon}(X)$, and the argument lies there (it is
the point $t=0$ of the path). So the right side of \eqref{4.22:eq-orbit-lemma-3} is a composition of
holomorphic maps. Thus $F$ is holomorphic near $\zeta_0$, and since $\zeta_0$ was arbitrary, on all
of $\mathcal{N}$.

\emph{The bound.} For every $\zeta$ the argument of $\mathfrak{E}^{\sharp}$ in
\eqref{4.22:eq-orbit-lemma-1} lies in $\mathfrak{D}\times\mathfrak{G}_{\varepsilon}(X)$ by
hypothesis. The number $\mathfrak{N}^{\mathbb{C}}_{\varepsilon}(\mathfrak{E})$ is the supremum of
$|\mathfrak{E}^{\sharp}|$ over this chart, so
$|F(\zeta)|\le\mathfrak{N}^{\mathbb{C}}_{\varepsilon}(\mathfrak{E})$.

\emph{The real slice.} Let $\zeta\in\mathcal{N}_{\mathbb{R}}$. Then $\mathfrak{q}(\zeta)$ and
$h(\zeta)$ are $G$-valued, and hence so is $h(\zeta)\mathfrak{q}(\zeta)$. If moreover
$\mathcal{P}(\zeta)\in\mathfrak{D}$, two facts combine:
\begin{itemize}
\item the property of the orbit datum that identifies $\mathfrak{E}^{\sharp}$ at $G$-valued
rotations with $\mathfrak{E}$ at the rotated further arguments;
\item the joint invariance of $\mathfrak{E}$ under $G$-valued gauge transformations acting on the
configuration and on the further arguments together.
\end{itemize}
Together they give
\begin{align*}
F(\zeta)&=\mathfrak{E}^{\sharp}\bigl(\mathcal{P}(\zeta)^{h(\zeta)};\mathbf{A},\boldsymbol{\Phi},
h(\zeta)\mathfrak{q}(\zeta)\bigr)\\
&=\mathfrak{E}\bigl(X;\mathcal{P}(\zeta)^{h(\zeta)},\mathbf{A}^{[h(\zeta)\mathfrak{q}(\zeta)]},
\boldsymbol{\Phi}^{[h(\zeta)\mathfrak{q}(\zeta)]}\bigr)\\
&=\mathfrak{E}\bigl(X;\mathcal{P}(\zeta),\mathbf{A}^{[\mathfrak{q}(\zeta)]},
\boldsymbol{\Phi}^{[\mathfrak{q}(\zeta)]}\bigr),
\end{align*}
the second equality by the first fact, applied at the $G$-valued rotation
$h(\zeta)\mathfrak{q}(\zeta)$, and the third by the second fact, applied to the $G$-valued gauge
transformation $h(\zeta)$. This is \eqref{4.22:eq-orbit-lemma-2}.

\emph{Independence of $h$.} Let $h_1,h_2$ both satisfy the hypotheses on $h$, and let $F_1,F_2$ be
the corresponding functions \eqref{4.22:eq-orbit-lemma-1}. By the first step, both are holomorphic
on the connected set $\mathcal{N}$.

On $\mathcal{N}_{\mathbb{R}}$ the quotient $h_2h_1^{-1}$ is $G$-valued. We apply the invariance of
the orbit datum under a $G$-valued gauge transformation acting on both the configuration slot and
the group slot, to $h_2h_1^{-1}$ at the point
$\bigl(\mathcal{P}(\zeta)^{h_1(\zeta)},h_1(\zeta)\mathfrak{q}(\zeta)\bigr)$. This gives
$F_1=F_2$ on $\mathcal{N}_{\mathbb{R}}$.

The difference $F_1-F_2$ is therefore holomorphic on $\mathcal{N}$ and vanishes on the totally real
slice $\mathcal{N}_{\mathbb{R}}$, which has non-empty interior. At an interior point of
$\mathcal{N}_{\mathbb{R}}$, every complex partial derivative of $F_1-F_2$ equals the corresponding
real partial derivative along $\mathcal{N}_{\mathbb{R}}$, and these vanish. Hence the Taylor
expansion of $F_1-F_2$ at that point is zero, and $F_1-F_2$ vanishes on a neighbourhood of it. Since
$\mathcal{N}$ is connected, the identity theorem gives $F_1=F_2$ on $\mathcal{N}$.
\end{proof}

The lemma is applied with
$(\mathcal{P},\mathfrak{q},h):=\bigl((\mathcal{U}_q^{\mathfrak{b}},\mathcal{J}_q^{\mathfrak{b}}),
\mathfrak{Q}^{\mathfrak{b}},h_q^{\mathfrak{b}}\bigr)$, the background pair, the rotation and the
presenter map of the presentation. With this choice $F$ is the rotated term
$F^{\mathfrak{b}}_{\mathfrak{E}}(\zeta;\Phi)$ of the presentation.

In that application no invariance of $\mathfrak{E}$ under the rough rotation
$\mathfrak{Q}^{\mathfrak{b}}$ is used at complex $\zeta$. The rotation $\mathfrak{Q}^{\mathfrak{b}}$
enters only by composition, through the group slot of $\mathfrak{E}^{\sharp}$.

The rule of this chapter on regauging fine configuration slots is also respected. No fine
configuration slot is regauged by $\mathfrak{Q}^{\mathfrak{b}}$. The rule is motivated by the jumps
across faces measured against \cite[(2.39)]{B-CMP119}. That reason does not apply to the gauge maps
of the frames inside the orbit datum to which part~(i) of the frame lemma
(Lemma~\ref{4.22:lem-frame-lemma}) refers, since there they are regular by that part.

\begin{corollary}[Increment of a bracket along the orbit]\label{4.22:cor-bracket-increment}
Let $\mathfrak{E}$ be a unit with domain $X$, let $\mathfrak{E}^{\sharp}$ be its orbit datum, in the
sense in which orbit data enter Lemma~\ref{4.22:lem-orbit-lemma}, and let
$\mathfrak{N}^{\mathbb{C}}(\mathfrak{E})$ be the rotated sup norm of $\mathfrak{E}^{\sharp}$ on its chart.
Fix the further arguments $(\mathbf{A},\boldsymbol{\Phi})$ and a parameter $\zeta$ of the manifold
$\mathcal{N}$ of Lemma~\ref{4.22:lem-orbit-lemma}, on which $P_s$ and the presenter maps depend. Let $s$
be a step of the presentation at level $k$. At this step the presenter maps are updated by
\begin{equation*}
\mathfrak{Q}^{\mathfrak{b}_s}=\mathfrak{Q}^{\mathfrak{b}_{s-1}}\,\mathfrak{w}_s ,
\end{equation*}
with gauge transformation $h_s$ and background pair
$(\mathcal{U}_q^{\mathfrak{b}_s},\mathcal{J}_q^{\mathfrak{b}_s})$. Put
\begin{equation*}
P_s:=\bigl((\mathcal{U}_q^{\mathfrak{b}_s})^{h_s},\ \operatorname{Ad}_{h_s}\mathcal{J}_q^{\mathfrak{b}_s}\bigr).
\end{equation*}
For $g'$ in the chart of the output, define
\begin{equation}\label{4.22:eq-bracket-increment-1}
\Delta_s^{A}(\zeta;g')
:=\mathfrak{E}^{\sharp}\bigl(P_s;\mathbf{A},\boldsymbol{\Phi},h_s\mathfrak{Q}^{\mathfrak{b}_s}g'\bigr)
-\mathfrak{E}^{\sharp}\bigl(P_s;\mathbf{A},\boldsymbol{\Phi},h_s\mathfrak{Q}^{\mathfrak{b}_{s-1}}g'\bigr).
\end{equation}
For a law site $s'$, write $\varepsilon(s')$ for the current radius at $s'$, $\eta^{L}(s')$ for the
debit at the presentation site, $\mathsf{e}(s')$ for the credit at $s'$ and $\tilde s'$ for the
point at which the local layer norm $\Theta_0^{\mathrm{loc}}(\mathbb{H}_s)$ is read for the site $s'$.
The constants $K_u$, $\mu_k$, $\varepsilon_{\star}$ are those of the presentation.

Assume the following statements.
\begin{itemize}
\item \emph{The first-order bound of orbit data in the rotation slot}, in the following form, which is
the one used below. Let $g_0$ be a rotation and let $Z$ be a $\mathfrak{g}^{\mathbb{C}}$-valued generator
on the law sites, and suppose that the ratio which that bound attaches to $(g_0,Z)$ is at least $2$; we
call this the \emph{ratio condition}. Then
\begin{equation}\label{4.22:eq-bracket-increment-2}
\bigl|\mathfrak{E}^{\sharp}(P;\mathbf{A},\boldsymbol{\Phi},g_0e^{Z})
-\mathfrak{E}^{\sharp}(P;\mathbf{A},\boldsymbol{\Phi},g_0)\bigr|
\le\mathfrak{N}^{\mathbb{C}}(\mathfrak{E})\sum_{s'}
\frac{2\,\|Z(s')\|}{\varepsilon(s')-\operatorname{pol}(g_0(s'))},
\end{equation}
where the sum runs over the live law sites $s'$ of $X$. We call the denominator the \emph{room} at $s'$.
\item \emph{The bracket bound.} By the first display of the bracket lemma together with part~(ii) of
the lemma on the background domain $\mathfrak{U}$, every law site $z$ satisfies
\begin{equation}\label{4.22:eq-bracket-increment-3}
|\log\mathfrak{w}_s(z)|\le K_u\,\Theta_0^{\mathrm{loc}}(\mathbb{H}_s)(\tilde z).
\end{equation}
\item \emph{The coupling ceilings.} These are the ceiling
$\varepsilon_{\mathrm{live}}\le\tfrac14\varepsilon_{\star}$ on the rotation size and the ceilings
\eqref{4.22:eq-coupling-ceilings-d}.
\item \emph{The adjoint bound}, in the form used below: for every $g'$ in the chart of the output,
every law site $s'$ and every $Y\in\mathfrak{g}^{\mathbb{C}}$, one has
$\|\operatorname{Ad}_{g'(s')^{-1}}Y\|\le e^{2\varepsilon_{\star}}\|Y\|$.
\item \emph{The far-site lemma.} This is the lemma on the distance of the live law sites.
\end{itemize}
Then the following hold.
\begin{itemize}
\item[(a)] At $g'=1$,
\begin{equation}\label{4.22:eq-bracket-increment-4}
|\Delta_s^{A}(\zeta;1)|\le\mathfrak{N}^{\mathbb{C}}(\mathfrak{E})
\sum_{s'\ \text{live law site of}\ X}\frac{4K_u}{\mathsf{e}(s')}\,
\Theta_0^{\mathrm{loc}}(\mathbb{H}_s)(\tilde s').
\end{equation}
\item[(b)] For every $g'$ in the chart of the output, that is, with
\begin{equation*}
\operatorname{pol}(g'(s'))<\varepsilon(s')-\eta^{L}(s')-\mu_k\,\mathsf{e}(s')/\varepsilon_{\star}
\end{equation*}
at every law site $s'$, one has
\begin{equation}\label{4.22:eq-bracket-increment-5}
|\Delta_s^{A}(\zeta;g')|\le\mathfrak{N}^{\mathbb{C}}(\mathfrak{E})\sum_{s'}
\frac{4.2\,K_u\,\varepsilon_{\star}}{\mu_k\,\mathsf{e}(s')}\,
\Theta_0^{\mathrm{loc}}(\mathbb{H}_s)(\tilde s').
\end{equation}
\end{itemize}
\end{corollary}

\begin{proof}
The proof rests on the bracket bound \eqref{4.22:eq-bracket-increment-3}.

(a) Since $\mathfrak{Q}^{\mathfrak{b}_s}=\mathfrak{Q}^{\mathfrak{b}_{s-1}}\mathfrak{w}_s$, the definition
\eqref{4.22:eq-bracket-increment-1} at $g'=1$ reads
\begin{equation*}
\Delta_s^{A}(\zeta;1)
=\mathfrak{E}^{\sharp}(P_s;\mathbf{A},\boldsymbol{\Phi},g_0e^{Z})
-\mathfrak{E}^{\sharp}(P_s;\mathbf{A},\boldsymbol{\Phi},g_0).
\end{equation*}
Here
\begin{equation*}
g_0:=h_s\mathfrak{Q}^{\mathfrak{b}_{s-1}},\qquad Z:=\log\mathfrak{w}_s .
\end{equation*}
This is because $e^{Z}=\mathfrak{w}_s$ sitewise. We apply the first-order bound
\eqref{4.22:eq-bracket-increment-2} to these $g_0$ and $Z$, and check its hypotheses.

At every live law site $s'$ one has $\operatorname{pol}(g_0(s'))\le\tfrac14\mathsf{e}(s')$. This is where
the ceiling $\varepsilon_{\mathrm{live}}\le\tfrac14\varepsilon_{\star}$ enters. Hence the room satisfies
$\varepsilon(s')-\operatorname{pol}(g_0(s'))\ge\tfrac12\mathsf{e}(s')$. The ratio condition
holds by the ceilings \eqref{4.22:eq-coupling-ceilings-d}.

The bound \eqref{4.22:eq-bracket-increment-2} therefore applies. At each live law site $s'$ we insert
$\|Z(s')\|\le K_u\Theta_0^{\mathrm{loc}}(\mathbb{H}_s)(\tilde s')$ from
\eqref{4.22:eq-bracket-increment-3} and the lower bound $\tfrac12\mathsf{e}(s')$ on the room. The
summand at $s'$ is then at most
\begin{equation*}
\frac{2K_u\Theta_0^{\mathrm{loc}}(\mathbb{H}_s)(\tilde s')}{\tfrac12\mathsf{e}(s')}
=\frac{4K_u}{\mathsf{e}(s')}\Theta_0^{\mathrm{loc}}(\mathbb{H}_s)(\tilde s').
\end{equation*}
This is \eqref{4.22:eq-bracket-increment-4}.

(b) Let $g'$ be as in (b). At each site,
$\mathfrak{w}_sg'=g'\,(g'^{-1}\mathfrak{w}_sg')$. With $\operatorname{Ad}_gA=gAg^{-1}$ one has
$g'^{-1}e^{\log\mathfrak{w}_s}g'=\exp(\operatorname{Ad}_{g'^{-1}}\log\mathfrak{w}_s)$. Hence
\begin{equation*}
h_s\mathfrak{Q}^{\mathfrak{b}_{s-1}}\mathfrak{w}_sg'=g_0e^{Z'},\qquad
g_0:=h_s\mathfrak{Q}^{\mathfrak{b}_{s-1}}g',\qquad
Z':=\operatorname{Ad}_{g'^{-1}}\log\mathfrak{w}_s .
\end{equation*}
Thus $\Delta_s^{A}(\zeta;g')$ is the difference of $\mathfrak{E}^{\sharp}(P_s;\mathbf{A},\boldsymbol{\Phi},\cdot)$
at $g_0e^{Z'}$ and at $g_0$.

By the adjoint bound and \eqref{4.22:eq-bracket-increment-3},
\begin{equation*}
\|Z'(s')\|\le e^{2\varepsilon_{\star}}K_u\Theta_0^{\mathrm{loc}}(\mathbb{H}_s)(\tilde s').
\end{equation*}
Polar size is subadditive, since $\|ab\|\le\|a\|\|b\|$ and $\|(ab)^{-1}\|\le\|b^{-1}\|\|a^{-1}\|$. The
presenter rotation $h_s\mathfrak{Q}^{\mathfrak{b}_{s-1}}$ contributes at most the debit $\eta^{L}(s')$.
Together with the condition on $g'$ this gives
\begin{equation*}
\operatorname{pol}(g_0(s'))\le\eta^{L}(s')+\operatorname{pol}(g'(s'))
<\varepsilon(s')-\mu_k\mathsf{e}(s')/\varepsilon_{\star}.
\end{equation*}
So the room at $s'$ is at least $\mu_k\mathsf{e}(s')/\varepsilon_{\star}$. The ratio condition
holds again by the ceilings \eqref{4.22:eq-coupling-ceilings-d}.

We now apply \eqref{4.22:eq-bracket-increment-2} with $g_0$ and $Z'$. The summand at $s'$ is at most
\begin{equation*}
\frac{2e^{2\varepsilon_{\star}}K_u\,\varepsilon_{\star}}{\mu_k\,\mathsf{e}(s')}
\Theta_0^{\mathrm{loc}}(\mathbb{H}_s)(\tilde s')
\le\frac{4.2\,K_u\,\varepsilon_{\star}}{\mu_k\,\mathsf{e}(s')}
\Theta_0^{\mathrm{loc}}(\mathbb{H}_s)(\tilde s').
\end{equation*}
The last step uses $2e^{2\varepsilon_{\star}}\le4.2$. This is \eqref{4.22:eq-bracket-increment-5}.
\end{proof}

\begin{remark}
Both bounds consist of one term per live law site. Each term is of first order in
$\Theta_0^{\mathrm{loc}}(\mathbb{H}_s)$, and is supported and decaying where $\mathbb{H}_s$ is.

At pointwise sites the coefficient in (a) is $4K_u/\varepsilon_{\star}$. This is a constant depending on
$(d,L,M)$. Part~(b) is the form in which the bounds of the new terms are read on their chart.

Every live law site is far in the sense of the far-site lemma. Its second part gives the scaled distance
$d_{\mathrm{sc}}(s')\ge\check R_k-2$. Hence the ceilings \eqref{4.22:eq-coupling-ceilings-d} give
\begin{equation*}
(\varepsilon_{\star}/\mu_k)\,e^{-\frac12\delta_{\natural}(\check R_k-2)}\le1,
\end{equation*}
together with the analogous inequality at frame sites. Consequently the right member of
\eqref{4.22:eq-bracket-increment-5} is at most
\begin{equation*}
\mathfrak{N}^{\mathbb{C}}(\mathfrak{E})\sum_{s'}\frac{4.2\,K_uC_{\natural}}{\varepsilon_{\star}}\,
\alpha^{+}_k\,e^{-\frac12\delta_{\natural}d_{\mathrm{sc}}(s')} .
\end{equation*}
This carries the factor $\alpha^{+}_k$ and is therefore small under the coupling ceilings
\eqref{4.22:eq-coupling-ceilings-d}, and it decays exponentially at the rate
$\tfrac12\delta_{\natural}$, which depends on $(d,L,M)$.

The localisation factors $e^{-(\kappa_1-1)M^{-d}|Y'_i|}$, in the notation of
\cite[(1.66)]{B-CMP122b}, do not come from this rate. They come from the $\sigma$-expansion of
$\mathbb{H}_s$, which \cite{B-CMP122b} carries into the function $u_{k\Lambda}$ of that paper.

Suppose the radius were proportional to $\varepsilon_{\mathrm{live}}$ instead of $\varepsilon_{\star}$.
Then the quotient in (a) would only be bounded by a coefficient of order one times
$e^{-\delta\,d_{\mathrm{sc}}(s')}$, with $\delta>0$ a decay rate, and the smallness of the data in
the coupling ceilings would be lost. The credit with constants depending on $(d,L,M)$ is therefore
necessary.
\end{remark}

\begin{remark}
The analyticity of the terms of the class $\mathcal{C}$ in the rotation variables is asserted, at every step, on a chart of one fixed radius. The class $\mathcal{C}$ at a step $k_2$ contains terms that were built at an earlier step $k_1<k_2$ from a function of the form
\begin{equation*}
F=\mathfrak{E}_1\bigl(\dots;\mathbf{A}^{[h_1(\zeta_1)]}\bigr).
\end{equation*}
Here $\mathfrak{E}_1$ is a term presented at step $k_1$, and $\zeta_i$ and $h_i$ denote the parameter and the presenting gauge map of the presentation at step $k_i$, in the notation of the orbit lemma (Lemma~\ref{4.22:lem-orbit-lemma}). At step $k_2$ the parameter $\zeta_1$ has become a function of $\zeta_2$, and when such a term is presented again at step $k_2$, it is read at
\begin{equation*}
\mathfrak{E}_1\bigl(\dots;\mathbf{A}^{[h_1(\zeta_1(\zeta_2))\,h_2(\zeta_2)]}\bigr).
\end{equation*}
Closedness of $\mathcal{C}$ at a fixed radius is therefore not automatic. In this remark $g,g'$ denote rotations, that is elements of $G^{\mathbb{C}}$, and not couplings. The polar size is sub-additive: for $g,g'\in G^{\mathbb{C}}$, from $\|gg'\|\le\|g\|\,\|g'\|$ and $\|(gg')^{-1}\|\le\|g'^{-1}\|\,\|g^{-1}\|$ one gets $\operatorname{pol}(gg')\le\operatorname{pol}(g)+\operatorname{pol}(g')$. Sub-additivity alone, however, bounds the composed rotation after presentations at steps $k_1<k_2<\cdots$ only by $\sum_i\varepsilon_{\mathrm{live}}(k_i)$.

This sum is of the order of $\sum_k g_k$ times a polylogarithmic factor in $g_k$, and it diverges with $K$. Indeed, the upper bound $x_k\le x_K+B_\sharp(K-k)$ of the coupling flow (Theorem~\ref{4.1:thm-clause-zero}) gives
\begin{equation*}
g_k=x_k^{-1/2}\ge\bigl(x_K+B_\sharp(K-k)\bigr)^{-1/2},
\end{equation*}
so that
\begin{equation*}
\sum_{k=0}^{K}g_k\ge\sum_{j=0}^{K}\bigl(x_K+B_\sharp j\bigr)^{-1/2},
\end{equation*}
which grows like $K^{1/2}$, while the polylogarithmic factor is bounded below by a positive constant.

The restriction on the total rotation is genuine. If a strong coordinate is rotated by $g\in G^{\mathbb{C}}$, then the imaginary part of the mean field there satisfies
\begin{equation*}
\operatorname{Im}\mathsf{m}\asymp\operatorname{pol}(g)\,\Theta\,A_1(\log g_k^{-2})^{p_0},
\end{equation*}
where $A_1(\log g_k^{-2})^{p_0}$ is the small-field threshold at level $k$ and $\Theta$ is the constant of this bound. The modulus bound $e^{c(\operatorname{Im}\mathsf{m})^2}$ for the integrand at a complex mean field, given by the mean-field modulus lemma, then exceeds the decay
\begin{equation*}
\exp\bigl(-c'A_1^2(\log g_k^{-2})^{2p_0}\bigr)
\end{equation*}
of the small-field cut-off factor at the same coordinate as soon as $\operatorname{pol}(g)\gtrsim c_\eta/\Theta_\flat$; here $c,c'>0$ are the constants of these two factors. Hence the radius $\varepsilon_\star$ is sharp in order, and the rotation used up along a lineage of presentations has to be counted rather than estimated by sub-additivity. The lemma below fixes, as the first step of that count, at which law site each argument is rotated at a presentation.
\end{remark}

\begin{lemma}[Each argument is rotated at its own law site]\label{4.22:lem-law-site-rule}
Let $\mathfrak{E}\in\mathcal{C}$ be a term with domain $X$, presented at a step by the presenting map $h$. Assume the following two clauses of the construction of presentations.
\begin{enumerate}
\item[(i)] The presentation identity. For the presentation with presenting map $h_q$,
\begin{equation}\label{4.22:eq-law-site-rule-1}
\begin{split}
F^{\mathfrak{b}}(U_k;\Phi)
&=\mathfrak{E}\bigl(X;(\mathfrak{b}^{\mathrm{pr}},\mathcal{J}(\mathfrak{b}^{\mathrm{pr}}),\mathbf{A})^{q}\bigr)\\
&=\mathfrak{E}\bigl(X;\mathfrak{b}^{\mathrm{pr}},\mathcal{J}(\mathfrak{b}^{\mathrm{pr}});\mathbf{A}\bigr).
\end{split}
\end{equation}
Here $U_k$ is the block field at level $k$, $\Phi$ the data carried at the step of the presentation (distinct from the further arguments $\boldsymbol{\Phi}$ of $\mathfrak{E}$), $\mathfrak{b}^{\mathrm{pr}}$ the presented background and $\mathcal{J}(\mathfrak{b}^{\mathrm{pr}})$ its current. The superscript $q$ denotes the triple transformed by the presenting map $h_q$; the component of this action on the further arguments is the rotation law~(ii).
\item[(ii)] The rotation law of further arguments. Each further argument of $\mathfrak{E}$ has a law site in $\mathfrak{s}(X)$. The term $\mathfrak{E}$ is invariant under the joint action of a gauge map $g$ on background, current and further arguments. The component of this action on the further arguments is $\mathbf{A}\mapsto\mathbf{A}^{[g]}$ and $\boldsymbol{\Phi}\mapsto\boldsymbol{\Phi}^{[g]}$, each further argument being rotated by the value of $g$ at its law site.
\end{enumerate}
Then:
\begin{enumerate}
\item[(a)] at the presentation, each further argument of $\mathfrak{E}$ is rotated by $h$ at the law site that it has as an argument of $\mathfrak{E}$;
\item[(b)] the root of a frame which the step may be about to create is a law site of the new terms only, and it is first read at a later step.
\end{enumerate}
\end{lemma}

\begin{proof}
(a) Apply the presentation identity~\eqref{4.22:eq-law-site-rule-1} with presenting map $h_q=h$. The presented function is $\mathfrak{E}$ evaluated at the triple transformed by $h$. The right-hand side of~\eqref{4.22:eq-law-site-rule-1} is a value of $\mathfrak{E}$ itself, and its second equality is the invariance of $\mathfrak{E}$ under the joint action of~(ii), so the only invariance that enters is that of $\mathfrak{E}$. No other transformation law is applied to the further arguments.

Consequently the presenter acts on the further arguments only through the law of $\mathfrak{E}$ itself, namely by $\mathbf{A}\mapsto\mathbf{A}^{[h]}$ in the sense of the rotation law~(ii). By~(ii), this rotates each further argument by the value of $h$ at the law site that the argument has as an argument of $\mathfrak{E}$. This proves (a).

The construction also contains a rule by which the arguments carried at sites interior to the region treated by the step are transformed. That rule is not an action on the further arguments of $\mathfrak{E}$: it is the covariance law of the presented function $F$ under a gauge map $g$ of its own argument $\zeta$, and it assigns to no further argument of $\mathfrak{E}$ a law site other than the one given by~(ii). It therefore does not alter (a).

(b) In the notation of the frame lemma (Lemma~\ref{4.22:lem-frame-lemma}), let $W$ be a frame which the step may create, with root $y_0^W$ and data $\Phi_W$. At the presentation of $\mathfrak{E}$ the frame $W$ is not yet a frame of $\mathfrak{E}$, so the data $\Phi_W$ are not among the further arguments of $\mathfrak{E}$. By~(a), the presentation rotates the arguments of $\mathfrak{E}$ only at the law sites they have as arguments of $\mathfrak{E}$. Hence no value of the presenter at the root $y_0^W$, taken as the root of $W$, rotates any argument of $\mathfrak{E}$. In particular, no rotation at a root is used to present arguments that were never framed.

The root $y_0^W$ is thus a law site only of the data $\Phi_W$, and $\Phi_W$ are arguments only of the terms that the step creates. Those terms enter the class $\mathcal{C}$ at a later step, and they are presented from then on. Applying~(a) to them, the data $\Phi_W$ are rotated at their law sites as arguments of the new terms, and this happens first at a later step.
\end{proof}

\begin{proposition}[A presentation lowers only the chart radius; closure under the large-field operation]\label{4.22:prop-presentation-radius}
\begin{enumerate}
\item[(a)] \emph{One presentation.} Assume the orbit lemma (Lemma~\ref{4.22:lem-orbit-lemma}) and
work in its setting, with its connected parameter manifold $\mathcal{N}$, its
holomorphic pair $(\mathcal{P},\mathfrak{q})$, the term $\mathfrak{E}$ with its orbit datum
$\mathfrak{E}^{\sharp}$ of radius $\varepsilon$, and the presenting map $h$. Assume in addition
that $(\mathcal{P},\mathfrak{q},h)$ does not depend on the variables $(\mathbf{A},\boldsymbol{\Phi})$
of $\mathfrak{E}$, that at every law site $s$
\begin{equation}\label{4.22:eq-presentation-radius-1}
\sup_{\zeta\in\mathcal{N}}\operatorname{pol}\bigl((h\mathfrak{q})(\zeta)(s)\bigr)\le\eta(s),
\end{equation}
and that $\mu(s)\ge0$ and $\eta+\mu<\varepsilon$. Then
\begin{equation}\label{4.22:eq-presentation-radius-2}
F^{\sharp}(\zeta;\mathbf{A},\boldsymbol{\Phi},g')
:=\mathfrak{E}^{\sharp}\bigl(\mathcal{P}(\zeta)^{h(\zeta)};\mathbf{A},\boldsymbol{\Phi},
h(\zeta)\mathfrak{q}(\zeta)g'\bigr)
\end{equation}
is an orbit datum of radius $\varepsilon-\eta-\mu$ of
$F(\zeta;\mathbf{A},\boldsymbol{\Phi}):=F^{\sharp}(\zeta;\mathbf{A},\boldsymbol{\Phi},1)$ on
$\mathcal{N}\times\mathfrak{D}_{\mathcal{S}}$: it is holomorphic in $\zeta$, in the weak
coordinates and in $g'$, and
$|F^{\sharp}|\le\mathfrak{N}^{\mathbb{C}}_{\varepsilon}(\mathfrak{E})$. The parameter domain
$\mathfrak{D}_{\mathcal{S}}$ (weak radii, strong boxes, strong sets) is unchanged; only the chart
radius is lowered.

\item[(b)] \emph{Closure.} In the setting of part (a), assume the following statements.
\begin{itemize}
\item The structure of Ba{\l}aban's large-field operation at the $\mathbf{R}$-site, as printed in
\cite[(1.68)--(1.72), (1.98), pp.~378--379]{B-CMP122b}, where
\cite[(1.98)]{B-CMP122b} is \cite[(2.12)--(2.13)]{B-CMP116}, used as printed and not re-derived:
\begin{equation}\label{4.22:eq-presentation-radius-b}
\begin{aligned}
\mathbf{T}'_k(X_i)&=\int d\nu_i(\Phi_i)\,I_i(\Phi_i;\mathbf{U},\mathbf{J})\,(\,\cdot\,),\\
I_i&=\exp\bigl[\mathsf{Q}_i+V(X_i^{\sim};\mathbf{U},\mathbf{J};\Phi_i)\bigr],\\
\mathbf{R}'^{(k)}(X)&=\sum_{n}\frac{1}{n!}\sum\rho_T(Z_1,\dots,Z_n)\prod_{j=1}^{n}H(Z_j).
\end{aligned}
\end{equation}
In the first line, by \cite[(1.71)]{B-CMP122b}, only the quadratic form in $B$ and the term
$V(X_i^{\sim})$ depend on $(\mathbf{U},\mathbf{J})$. In the second, $\mathsf{Q}_i$ is the
quadratic form in $B$ of \cite[(1.71)]{B-CMP122b}, namely minus one half of the pairing of
$DH''_{1,k,X}B$ with itself, $D$ being the operator so denoted in \cite[(1.71)]{B-CMP122b},
weighted by the factor printed there between the two entries, and
$V(Y;\cdot)$ is
the resummation \cite[\S1]{B-CMP122b} of localised consumer terms, each of them a presented term
$F^{\mathfrak{b}}_{\mathfrak{E}}(\zeta;\Phi)$. In the third, $\rho_T$ is combinatorial in the
domains, and the activities $H(Z)$ are built from $\mathbf{T}'_k$- and $V$-objects interpolated
in the cluster-expansion parameters of \cite[(2.12)--(2.13)]{B-CMP116}.
\item The decomposition of each term into pieces over strong sets on one common parameter domain
(the inheritance theorem), and the joint invariance of each piece (clause (b) of the flat-bound
statement and clause (b) of the second analyticity lemma).
\item The bracket lemma, and clause (b) of Corollary~\ref{4.22:cor-bracket-increment}.
\item The bracket-piece bound accompanying the correlation theorem, with
Lemma~\ref{4.22:lem-majorant-transfer} and its constant $C'$.
\item The two structural conditions on $\mathsf{s}$-pieces; the deep-core lemma, clauses (a)
and (b); and the margin bound of the deep-core lemma.
\item Clause (iii) of the exclusion lemma: no piece carries two of the three devices used below
(the layer bracket, the disc bracket and the interpolated $\mathsf{s}$-piece).
\item The induction hypothesis on orbit data at level $k$
(Definition~\ref{4.22:def-induction-hypothesis}), and the bounds of the step at this site, with
the right members made precise in the proof.
\item Clause (ii) of the joint-invariance theorem; the invariance property of the kept quadratic
form $\mathsf{Q}_i$; and the identity principle on finite products of charts
$\mathfrak{G}_{\varepsilon}$.
\item Clauses (ii) and (v) of the frame lemma (Lemma~\ref{4.22:lem-frame-lemma}).
\end{itemize}
Let $\mathcal{C}$ be the class of terms presented at step $k$. Assume that every piece
$\mathfrak{E}(X,\mathcal{S})$ of every $\mathfrak{E}\in\mathcal{C}$ has an orbit datum of the
current radius $\varepsilon$, where
\begin{equation}\label{4.22:eq-presentation-radius-3}
\varepsilon>\eta^{L}_k+2\mu_k\mathsf{e}/\varepsilon_{\star}+\eta^{\mathrm{inc}}_k .
\end{equation}
Here the middle term is one margin per presentation, there being at most two presentations, and
$\eta^{\mathrm{inc}}_k$ is the reserve.

Then the terms created at the $\mathbf{R}$-site are $V(Y)$, $\mathbf{R}'^{(k)}(X)$ and
$\mathbf{B}'^{(k)}(X)$ of \cite[(1.68)--(1.72), (1.98)--(1.100)]{B-CMP122b} and, for every
component $Y_i$ of the second class, the two members of $\mathfrak{K}(Y_i)$. Each of them is a
finite sum of stored terms $(c,S,\mathcal{S})$, in which $\mathcal{S}$ is the union of the strong
sets of its constituents at the surviving coordinates. Each stored term has an orbit datum of
radius at least $\varepsilon_{k+1}$ at the surviving law sites, where
\begin{equation}\label{4.22:eq-presentation-radius-4}
\varepsilon_{k+1}=\varepsilon-\eta^{L}_k-2\mu_k\mathsf{e}/\varepsilon_{\star}
-\eta^{\mathrm{inc}}_k ,
\end{equation}
and of radius $\mathsf{e}$ at the law sites of the new frames, on the
stored domain of the frame lemma. The bounds of the site hold in the form
\begin{equation}\label{4.22:eq-presentation-radius-5}
\sum_{\mathcal{S}}e^{\kappa_A(k)|\mathcal{S}|}\,
\mathfrak{N}^{\mathbb{C}}\bigl(c(S,\mathcal{S})\bigr)\le
\bigl(\text{the right member of the bound of the step at this site}\bigr),
\end{equation}
the right member being the one assigned to the term by the correlation theorem or, for wrapped
consumers, by the two propositions on wrapped consumers, as made precise in the proof.
\end{enumerate}
\end{proposition}

\begin{proof}
\emph{Part (a).} Let $g'\in\mathfrak{G}_{\varepsilon-\eta-\mu}(X)$. For $a,b\in G^{\mathbb{C}}$
one has $\|ab\|\le\|a\|\,\|b\|$ and $\|(ab)^{-1}\|\le\|b^{-1}\|\,\|a^{-1}\|$. Hence
$\max\{\|ab\|,\|(ab)^{-1}\|\}$ is at most the product of $\max\{\|a\|,\|a^{-1}\|\}$ and
$\max\{\|b\|,\|b^{-1}\|\}$, that is,
$\operatorname{pol}(ab)\le\operatorname{pol}(a)+\operatorname{pol}(b)$. With
\eqref{4.22:eq-presentation-radius-1} this gives, at every law site,
\begin{equation*}
\operatorname{pol}\bigl(h(\zeta)\mathfrak{q}(\zeta)g'\bigr)
\le\eta+\operatorname{pol}(g')<\eta+(\varepsilon-\eta-\mu)\le\varepsilon ,
\end{equation*}
since $\mu\ge0$. Consider
$\mathcal{N}':=\mathcal{N}\times\mathfrak{G}_{\varepsilon-\eta-\mu}(X)$. It is connected, and its
real slice $\mathcal{N}_{\mathbb{R}}\times G^{\mathfrak{s}(X)}$ has non-empty interior in it.

We apply the orbit lemma on $\mathcal{N}'$ to the triple $(\mathcal{P},\mathfrak{q}g',h)$. Its
hypotheses hold for the following reasons.
\begin{itemize}
\item The map $(\zeta,g')\mapsto(\mathcal{P}(\zeta),\mathfrak{q}(\zeta)g')$ is holomorphic.
\item It is real on the real slice: $\mathcal{P}$ is real on $\mathcal{N}_{\mathbb{R}}$ by the
setting of the orbit lemma, and $\mathfrak{q}(\zeta)g'$ is $G$-valued there because
$\mathfrak{q}$ and $g'$ are.
\item The map $h$ is continuous and $G$-valued on the real slice.
\item $\mathcal{P}(\zeta)^{h(\zeta)}$ lies in $\mathfrak{D}$, as in the setting of the orbit
lemma.
\item $h(\zeta)\mathfrak{q}(\zeta)g'$ lies in $\mathfrak{G}_{\varepsilon}(X)$ by the inequality
just proved.
\end{itemize}
The orbit lemma gives holomorphy of $F^{\sharp}$ in $(\zeta,g')$ and the bound
$|F^{\sharp}|\le\mathfrak{N}^{\mathbb{C}}_{\varepsilon}(\mathfrak{E})$. The weak coordinates
are holomorphic parameters throughout.

It remains to show that $F^{\sharp}$ reproduces $F$ at real rotations. Let $g'$ be $G$-valued.
An orbit datum absorbs a $G$-valued right factor of its rotation into the rotated arguments
$\mathbf{A}^{[g']}$, $\boldsymbol{\Phi}^{[g']}$. Using this property of $\mathfrak{E}^{\sharp}$,
\begin{equation*}
\begin{aligned}
\mathfrak{E}^{\sharp}\bigl(\mathcal{P}(\zeta)^{h(\zeta)};\mathbf{A},\boldsymbol{\Phi},
h(\zeta)\mathfrak{q}(\zeta)g'\bigr)
&=\mathfrak{E}^{\sharp}\bigl(\mathcal{P}(\zeta)^{h(\zeta)};\mathbf{A}^{[g']},
\boldsymbol{\Phi}^{[g']},h(\zeta)\mathfrak{q}(\zeta)\bigr)\\
&=F\bigl(\zeta;\mathbf{A}^{[g']},\boldsymbol{\Phi}^{[g']}\bigr).
\end{aligned}
\end{equation*}
The second equality holds because $(\mathcal{P},\mathfrak{q},h)$ does not depend on
$(\mathbf{A},\boldsymbol{\Phi})$. This is the defining identity of an orbit datum of $F$ at
$G$-valued rotations. Since $\mathfrak{D}_{\mathcal{S}}$ has entered only as the domain of the
fixed parameters, it is unchanged.

\emph{Part (b): termwise presentation.} The first hypothesis of part (a) holds for the
presentations of the step: their backgrounds, layers and carried gauges are functions of $\zeta$,
of the data and of the present-step field $\Phi$, never of any other argument of a consumer term
(\cite{B-CMP119}; the corollary on presenters), and an older frame's data $\Phi_W$ are read by
stored terms only (Lemma~\ref{4.22:lem-frame-lemma}, clause (v)).
Write $\mathfrak{E}=\sum_{\mathcal{S}}\mathfrak{E}(X,\mathcal{S})$.
By the decomposition of the inheritance theorem, all pieces live on one parameter domain
$\mathfrak{D}_{\mathfrak{E}}$. Each piece is jointly invariant, by clause (b) of the flat-bound
statement and clause (b) of the second analyticity lemma. The strong sets $\mathcal{S}$ do not
depend on $(\zeta,g')$. We present each piece: it
receives the datum $F^{\sharp}_{\mathcal{S}}$ of part (a), with margin
$\mu:=\mu_k\mathsf{e}/\varepsilon_{\star}$.

\emph{Part (b): the four kinds of consumer terms.} A localised consumer term is of one of four
kinds. In the bounds below, $\mathfrak{B}^{\mathrm{rec}}_s$ is the first term of the
decomposition of a bracket given by the bracket lemma, and $v^{\sharp}_p$ is the orbit datum of
the piece $p$ of the unit $v$; $\mathsf{w}_L$, $\eta_p$, $m(p)$ and $C_b$ are the numbers, and
$K_v$ is the domain, attached to $p$ and $v$ in the bracket-piece bound accompanying the
correlation theorem. The bounds of the four kinds are collected in
\begin{equation}\label{4.22:eq-presentation-radius-a}
\begin{aligned}
&\text{values, $\sigma$-pieces:}\quad F^{\sharp}_{\mathcal{S}}\ \text{as in part (a)};\\
&\text{layer brackets:}\quad \mathfrak{B}^{\mathrm{rec}}_s+\Delta^{A}_s;\\
&\text{disc brackets:}\quad |v^{\sharp}_p|\le2\mathsf{w}_L^{-1}(C_b\eta_p)^{1/2}
e^{-(\kappa_1-1)m(p)}\,\mathfrak{N}^{\mathbb{C}}_{\varepsilon}(v);\\
&\text{$\mathsf{s}$-pieces:}\quad |\mathfrak{p}|\le\mathfrak{N}^{\mathbb{C}}_{\varepsilon}(v)\,
e^{-(\kappa_1-1)\sum_t|\mathsf{y}_t|},\\
&\phantom{\text{$\mathsf{s}$-pieces:}}\quad
|\mathfrak{p}|\le C_{3F}\,e^{-\frac12\delta_PD_X}\,\mathfrak{N}^{\mathbb{C}}_{\varepsilon}(v)\,
e^{-(\kappa_1-1)\sum_t|\mathsf{y}_t|}\quad\text{if }X_v\subset\mathsf{C}_C .
\end{aligned}
\end{equation}
We treat the four kinds in turn.

\emph{Values and $\sigma$-pieces.} A value or a $\sigma$-piece $F_{\mathcal{S}}$ does not involve
the variable $\mathsf{t}$. For it part (a) applies as it stands; the margin is reserved and not
used.

\emph{Layer brackets.} A bracket in the arrangement of the presentation has, by the bracket lemma,
the form $\mathfrak{B}^{\mathrm{rec}}_s+\Delta^{A}_s$. In $\mathfrak{B}^{\mathrm{rec}}_s$ both
members carry the common rotation $\mathfrak{Q}^{\mathfrak{b}_{s-1}}g'$. Its size, a Cauchy
estimate in the increment of the presented layer, therefore reads $\mathfrak{E}^{\sharp}$ through
$\mathfrak{N}^{\mathbb{C}}$, with the rotation a fixed parameter. The increment $\Delta^{A}_s$ is
bounded on the whole chart by clause (b) of Corollary~\ref{4.22:cor-bracket-increment}.

\emph{Disc brackets.} Consider a bracket piece of the correlation theorem's expansion, on the disc
$|\mathsf{t}|\le\rho_p$ of the Cauchy variable $\mathsf{t}$ given by the disc-rider statement,
under the six disc conditions of the disc lemma. The pieces of this
kind are:
\begin{itemize}
\item the plain second-part terms, at $s=4,\dots,1$ (wrapped ones are of the fourth kind, by the
filing rule for wrapped second-part terms);
\item the plain first-part terms with $K_v\setminus Z\neq\emptyset$, which occur at $s=4$ only.
\end{itemize}
For these we use Lemma~\ref{4.22:lem-majorant-transfer}, with its constant $C'$ taken at the full
share of the margin $\mu$. Its input is the bracket-piece bound accompanying the correlation
theorem, which is the third line of \eqref{4.22:eq-presentation-radius-a} and holds on the whole
chart.

\emph{$\mathsf{s}$-pieces.} An $\mathsf{s}$-piece is
$\mathfrak{p}=D_{\mathsf{y}_4}D_{\mathsf{y}_3}D_{\mathsf{y}_2}D_{\mathsf{y}_1}F^{\mathfrak{b}}$,
under the two structural conditions on $\mathsf{s}$-pieces. The pieces of this kind are:
\begin{itemize}
\item every interior first-part term ($X_v\subset Z$), plain or wrapped;
\item every wrapped term beside a live second-class component, live meaning that it still
carries its large-field factor at the level read;
\item the members $\mathfrak{K}(Y)$ of older second-class components, the class $\mathcal{C}$
being understood to contain them.
\end{itemize}
By the fundamental theorem of calculus in
$(\mathsf{s}^{(1)},\dots,\mathsf{s}^{(4)})$, with no partition-of-unity factor $\zeta_{\square}$
of \cite[p.~276]{B-CMP119} and no variable
$\mathsf{t}$, Lemma~\ref{4.22:lem-majorant-transfer} applies with $\mathfrak{P}^{\mathrm{inc}}=\emptyset$.
It gives the fourth line of \eqref{4.22:eq-presentation-radius-a}.

If moreover $X_v\subset\mathsf{C}_C$ (the deep case), clause (b) of the deep-core lemma, through
its first four links, gives the
last line. This is the oscillation clause $(\mathrm{M}3)'$ of
\eqref{4.22:eq-majorant-transfer-a}, with clause (a) of the deep-core lemma supplying the
oscillation. The disc in $\lambda$ used by the deep-core lemma is an instrument of the bound
only. By the margin bound of the deep-core lemma it is accommodated inside the same margin, and
uses at most $\frac12\mu$.

For $v\in\mathfrak{K}(Y)$ two replacements are made:
\begin{itemize}
\item in the first $\mathsf{s}$-piece bound, $\mathfrak{N}^{\mathbb{C}}_{\varepsilon}(v)$ is
replaced by $\mathfrak{O}(v)$, by clause $(\mathrm{ii})_2$ of
\eqref{4.22:eq-kept-oscillation-a} and the first proposition on the kept members'
contributions;
\item in the second, $C_{3F}\mathfrak{N}^{\mathbb{C}}_{\varepsilon}(v)$ is replaced by
$\frac12C_{3F}\mathfrak{O}(v)$, by clause (iv) of \eqref{4.22:eq-kept-oscillation-a} (which rests
on clause (c) of the holomorphy lemma) and the second proposition on the kept members'
contributions.
\end{itemize}

By clause (iii) of the exclusion lemma, no piece carries two of the devices of the last three
kinds. Hence every bound of the site reads a piece through
$\mathfrak{N}^{\mathbb{C}}(\mathfrak{E}(X,\mathcal{S}))$, or through the oscillation of
$(\mathrm{M}3)'$ in \eqref{4.22:eq-majorant-transfer-a}, on the whole chart.

\emph{Part (b): assignment of strong sets.} A product of pieces is assigned to the union of their
strong sets, restricted to the surviving coordinates. The weights are sub-multiplicative. Cubes at
coordinates that $d\nu_i$ integrates are dropped; since $e^{\kappa_A(k)|\cdot|}\ge1$, this errs
on the
safe side. The bounds \cite[(1.66)--(1.69), (1.91)--(1.97)]{B-CMP122b} are used in the form
\begin{equation*}
\sum_{\mathcal{S}}e^{\kappa_A(k)|\mathcal{S}|}\,
\mathfrak{N}^{\mathbb{C}}\bigl(\mathfrak{E}(X,\mathcal{S})\bigr)
\le\bigl(\text{the right member for }\mathfrak{E}\bigr)
\end{equation*}
wherever they read $\sup|\mathfrak{E}|$. The right member is the one provided by the induction
hypothesis on orbit data at level $k$ (Definition~\ref{4.22:def-induction-hypothesis}) and by the
hypothesis of part (b). These right members are read as follows.
\begin{enumerate}
\item[(i)] The majorants at this site are those of the statement that runs the consumer. For
consumers that are not
wrapped they are those of the remainder proposition of the correlation theorem, clauses (c)--(e),
with its bracket bound, its $\sigma$-bound and the sum of the $\sigma$-bounds. For wrapped
consumers they are the label, summation and anchor bounds, the starred bound and the deep
bound of the two propositions on wrapped
consumers. For wrapped consumers the plain quantities $d_j(X)$, $d_k(Y_0)$ of
\cite[(1.66)]{B-CMP122b} are never used. Healed terms of the wrapped class $\mathfrak{S}$, healed
meaning re-filed into the regular expansion by a later step, have
$\mathcal{P}_F=\emptyset$, so that their length $\mathsf{d}_{\ell}$ is the plain distance $d$,
and they are bounded by the estimates of the
correlation theorem.
\item[(ii)] The third factor of \cite[(1.66)]{B-CMP122b} is supplied by the deep-core lemma, for
every first-part consumer whose true domain lies in a core, plain or wrapped. The corresponding
estimate of \cite[\S1]{B-CMP122b} is not used for any class.
\item[(iii)] The bound \cite[(1.69)]{B-CMP122b} is the bound at birth of the members of
$\mathfrak{K}(Y_i)$, a count of old cells. It is read at step $k$ by the corresponding clause of
the correlation theorem, which carries the bound at birth on, and by no bound at a later
step. Every later bound on these members is the tube bound, with $W:=Y_i$,
$\mathfrak{O}(v)\le C_{\Sigma}\check{R}_{k_W}^{p_{\Sigma}}$ of clauses $(\mathrm{ii})_2$ and
(iv) in \eqref{4.22:eq-kept-oscillation-a}. It is paid through the component's own collar, in
cells and in the distance $d^{\wedge}$.
\end{enumerate}

\emph{Part (b): the operations.} The surviving variables and $g'$ enter through the functions $F$
only: the measures $d\nu_i$, the operator of the quadratic form $\mathsf{Q}_i$, the factors
$\rho_T$ and the interpolation in the cluster-expansion parameters do not depend on them. Put
$g':=1$ at the law sites
of the variables that $d\nu_i$ integrates.

The remaining operations are finite sums, integrals over real $\Phi_i$ against a finite measure
free of the parameters, and series converging under the majorants of the preceding paragraph. On
the correlation theorem's route through \cite[(1.102)]{B-CMP122b} no quotient is formed. The
identities
\cite[(1.73)--(1.75)]{B-CMP122b}, which serve that paper for \cite[(1.103)--(1.104)]{B-CMP122b},
are used by no bound and no device of this argument, nor by the correlation theorem; nor does
the lifted bond term $V^{(k)\uparrow}$ enter anywhere, not even as notation. These
operations preserve, term by term:
\begin{itemize}
\item continuity;
\item holomorphy in $\zeta$, the weak coordinates and $g'$, by the argument of the
complex-parameter lemma;
\item the defining identity of an orbit datum at $G$-valued rotations.
\end{itemize}

\emph{Part (b): the second class.} For a component $Y_i$ of the second class, the integration
$d\nu_i$ is not performed in \cite[(1.102)]{B-CMP122b}. All of $\Phi_i$, the interior pointwise
variables (the class $(\beta)$ of \eqref{4.22:eq-older-law-sites-a}) and the interior older
frames included, become the real data $\Phi_{Y_i}$.

The joint invariance of the new terms is clause (ii) of the joint-invariance theorem. It applies
to the terms $V$ and $\mathbf{B}'$; for $\mathsf{Q}_i$ it is the invariance property of the kept
quadratic form assumed above. For the chart
variables it is supplemented by the identity principle on finite
products of charts $\mathfrak{G}_{\varepsilon}$. Clause (ii) of the frame lemma
(Lemma~\ref{4.22:lem-frame-lemma}) gives the datum at $\Sigma_{Y_i}$, and clause (v) supplies the
orbit data of the interior variables. The two factors of $I_i$ in
\eqref{4.22:eq-presentation-radius-b}
remain under $\mathbf{T}'_k(Y_i)$ as the members of $\mathfrak{K}(Y_i)$, with the datum of clause
(ii) of the frame lemma.

\emph{Part (b): the radius.} Each presentation lowers $\varepsilon$ by its debit $\eta$ and its
margin $\mu$, by part (a) and clause (iii) of Lemma~\ref{4.22:lem-majorant-transfer}. Along a
lineage there are at most two presentations:
\begin{itemize}
\item the four brackets of a second-part term \cite[\S1]{B-CMP122b} are four separate outputs,
each presented once, so their costs enter as a maximum and not as a sum;
\item the $Y_0$-term of a bracket is presented once more \cite[\S1]{B-CMP122b};
\item a first-part term is presented once \cite[\S1]{B-CMP122b}.
\end{itemize}
The reserve $\eta^{\mathrm{inc}}_k$ is then removed by choice. This gives the radius in
\eqref{4.22:eq-presentation-radius-4}.
\end{proof}

\begin{remark}
\begin{enumerate}
\item[(1)] Restricting the rotation $g'$ of part (a) to the diagonal gives, by the third of the
defining properties of orbit data, the same function as the datum of clause (ii) of the frame
lemma wherever both exist.
Its radius, however, would have to be measured in units of $\varepsilon_{\star}$, and $G$-valued
members carry no orbit datum at that scale. Clause (v) of the frame lemma is used instead.
\item[(2)] At the fluctuation integral $(\mathrm{T})$ and at the large-field stages
$(\mathrm{P})$ the radii are kept, apart from the radius adjustment made at those sites (the
statement on the radii at the fluctuation and large-field sites).
At the presentation $(\mathrm{L})$ the parameter domain is kept and the chart radius is lowered,
as in part (a).
\item[(3)] The tube form of the bounds cannot carry this argument. At a weak coordinate it would
shrink the weak radius by the factor $e^{-2\sum\operatorname{pol}}$, and at a strong coordinate it
has nothing to shrink. It is first converted into an orbit datum.
\end{enumerate}
\end{remark}

\begin{definition}[Credits and debits of chart radius over a term's life]
\label{4.22:def-radius-accounting}
Fix a stored term with domain $S$, one of its stored variables, and a law site $s\in\mathfrak{s}(S)$ of
that variable. A step has three sites: the fluctuation integral $(\mathrm{T})$, the stages
$(\mathrm{P})$ of the large-field operation, and the presentation $(\mathrm{L})$. At level $k$ the
variable is carried with an orbit datum whose chart at $s$ consists of the complex rotations $g'$
with $\operatorname{pol}(g'(s))<\varepsilon_k(s)$; $\varepsilon_k(s)$ is the \emph{current radius} at
$s$. Here $\mathsf{e}(s)$ is the unit of the law site $s$, $\varepsilon_\star$ is the rotation radius
at pointwise sites, and $\mu_k$ is the margin of step $k$; these are the quantities in which the
radius budget \eqref{4.22:eq-budget-closes-a} is stated.

The \emph{accounting of chart radius at $s$} is the following list. A \emph{credit} opens the
account and fixes its initial radius; a \emph{debit} lowers the current radius; the
\emph{reserve} is radius kept back by choice; a \emph{closing} ends the account. The amounts are:
\begin{equation}\label{4.22:eq-radius-accounting-a}
\begin{array}{lll}
(\mathrm{c}1) & \text{credit} & \varepsilon_\star,\\[2pt]
(\mathrm{c}2) & \text{credit} &
  \varepsilon^{\mathrm{fr}}(W)=\varpi_{\mathrm{fr}}\,\alpha_{*,k_W}/(4C_{\mathrm{fr}}),\\[2pt]
(\ell 1) & \text{debit} & \eta^T_k(s),\\[2pt]
(\ell 2) & \text{debit} & \mu_k\mathsf{e}(s)/\varepsilon_\star,\\[2pt]
(\ell 3) & \text{debit} & \nu(s)\le\varepsilon_\star/16,\ \text{incurred once},\\[2pt]
(\ell 4) & \text{debit} & \text{part of }\eta^P_k(s),\\[2pt]
(\ell 5) & \text{debit} & \text{part of }\eta^L_k(s),\\[2pt]
(\ell 6) & \text{debit} & \text{part of }\eta^L_k(s),\\[2pt]
(\ell 7) & \text{debit} & \mu_k\mathsf{e}(s)/\varepsilon_\star\ \text{for each presentation},\\[2pt]
(\ell 8) & \text{reserve} & \eta^{\mathrm{inc}}_k(s)=\mu_k\mathsf{e}(s)/\varepsilon_\star\
  \text{per step},\\[2pt]
(\mathrm{x}1) & \text{closing} & \text{the account is closed},\\[2pt]
(\mathrm{x}2) & \text{closing} & \text{the account is closed; (c2) opens at }\Sigma_{W_2}.
\end{array}
\end{equation}
For each entry we state when it is made, which map, if any, rotates the variable at $s$, the
nature of the bound on the amount, and the statements that supply it. An amount called
\emph{far} is a far debit, bounded by the far rates of the lemma bounding the debits at the three
sites (below, \emph{the debit lemma}); an amount called \emph{chosen} is fixed by the choice of
margins.
\begin{itemize}
\item[(c1)] Made at the creation of a pointwise coordinate: at the fluctuation site that creates
the new slot on the rim, or at the arrest that creates such a coordinate. No map rotates the
variable. The amount
$\varepsilon_\star$ is a number depending on $(d,L,M)$ only. It is supplied, at a fluctuation site,
by the theorem on orbit data at a fluctuation site and its corollary, and, at an arrest, by the
lemma on orbit data at an arrest and its companion lemma; these give the new pointwise
coordinates their orbit data.
\item[(c2)] Made at a second-class outcome at level $k_W$, for the frame $W$. No map rotates the
variable. The amount $\varepsilon^{\mathrm{fr}}(W)$ of \eqref{4.22:eq-radius-accounting-a} depends on
the coupling $g_{k_W}$ through $\alpha_{*,k_W}$. It is supplied by clause (ii) of the frame lemma
(Lemma~\ref{4.22:lem-frame-lemma}) together with the definition of the frame radius.
\item[$(\ell 1)$] Made at the fluctuation site $(\mathrm{T})$ of every later step $k$. The variable
is rotated by $\mathfrak{u}_p(0,\sigma)$, the map from Landau to axial gauge of the undilated
response. The amount $\eta^T_k(s)$ is far; it is supplied by clause (iv) of the axialising gauge
map lemma (Lemma~\ref{4.22:lem-axialising-map}) and clause (iii) of the debit lemma.
\item[$(\ell 2)$] Made at the same site. No map rotates the variable: the entry is the margin for
the disc in the Cauchy variable $\mathsf{t}$. The amount $\mu_k\mathsf{e}(s)/\varepsilon_\star$ is
chosen; it is supplied by the margin lemma at the fluctuation site.
\item[$(\ell 3)$] Made at the stages of the one chain of large-field stages in whose range
$[h,k-1]$ of levels the level of the slot lies. The variable is rotated by
$\mathfrak{u}^{(l)}_{t,\sigma}(b)$
on the polydisc of all blocks, near the collars. The amount is $\nu(s)\le\varepsilon_\star/16$,
incurred once; it is supplied by clause (iv) of the debit lemma.
\item[$(\ell 4)$] Made at every stage of every other chain. The variable is rotated by the same map,
now far from the collars. The amount is part of $\eta^P_k(s)$ and is far; it is supplied by
Lemma~\ref{4.22:lem-axialising-map}(iv) and clause (iii) of the debit lemma.
\item[$(\ell 5)$] Made at a presentation, when $s\notin Z$. The variable is rotated by the product
$\mathfrak{Q}^{\mathfrak{b}}=\prod_i\mathfrak{u}_i$ of the presenter maps along the presentation
chain of background pairs, whose members are written $\mathfrak{b}_0,\dots,\mathfrak{b}_4$; these
members are not the small-field thresholds $\mathfrak{b}_k$. The carried gauge is
$h_q=1$. The amount is
part of $\eta^L_k(s)$ and is far; it is supplied by clause (iii)($\alpha$) of the debit lemma.
\item[$(\ell 6)$] Made at a presentation, when $s\in Z_h$, where $Z_h$ is a large-field region
contained in the region $D_2$ fixed with the presentation. The variable is rotated by
$\mathfrak{u}_3\mathfrak{u}_4$, with $h_q=1$ and $\mathfrak{w}_1=\mathfrak{w}_2=1$. The amount is part
of $\eta^L_k(s)$ and is far; it is supplied by clause (iii)($\beta$) of the debit lemma and the
localisation of the current beyond the cut, \eqref{4.22:eq-current-localisation-a}.
\item[$(\ell 7)$] Made at each of the at most two presentations of a step. No map rotates the
variable: the entry is a margin, one room per presentation, shared by the bounds of all its
pieces, namely the Cauchy estimate in $g'$ on the chart, the $\mathsf{t}$-disc of the correlation
theorem at radius $\rho_p$, where $\rho_p$ is the radius of the correlation theorem's
$\mathsf{t}$-disc at the site in question (its lemma's at $(\mathrm{T})$, its rider's at
$(\mathrm{L})$), given for each site in the list of users below, and the $\lambda$-disc of
the lemma on deep first-part units (below,
\emph{the deep-unit lemma}). The amount $\mu_k\mathsf{e}(s)/\varepsilon_\star$ for each presentation
is chosen. It is supplied by Corollary~\ref{4.22:cor-bracket-increment}(b), by the margin lemma at
a presentation together with the conditions that keep the correlation theorem's rider and the
deep-unit lemma within the margin, and by the lemma on one margin for many bounds.
\item[$(\ell 8)$] Made at the presentation $(\mathrm{L})$ of every later step. No map rotates the
variable: the entry is room kept for a second device on one piece; it is idle in the
correlation theorem's treatment of bracket pieces. The amount
$\eta^{\mathrm{inc}}_k(s)=\mu_k\mathsf{e}(s)/\varepsilon_\star$ per step is chosen. It is supplied by
the lemma on several devices on one piece and by clause (iv) of the lemma on one margin for many
bounds.
\item[(x1)] Made when the variable is integrated, at stage $l$ or by the measure $d\nu_i$: the
account is closed. This is supplied by Corollary~\ref{4.22:cor-preparatory-chart} and by the step on
operations in the proof of Proposition~\ref{4.22:prop-presentation-radius}(b).
\item[(x2)] Made when the component of the variable takes the second class. The variable joins the
data $\Phi_{W_2}$ of the new frame $W_2$ and its account is closed; the credit (c2) opens at the
law sites $\Sigma_{W_2}$. This is supplied by clause (v) of the frame lemma
(Lemma~\ref{4.22:lem-frame-lemma}).
\end{itemize}

\emph{Total debit and reserve.} Let $\eta_k(s)$ be the total debit of step $k$ at $s$, as in the
radius budget \eqref{4.22:eq-budget-closes-a}. The reserve $(\ell 8)$ carries its own letter
$\eta^{\mathrm{inc}}_k(s)$ and is not part of $\eta_k(s)$. Consequently the bounds of the radius
budget on the total debit,
\begin{equation}\label{4.22:eq-radius-accounting-1}
\eta_k(s)\le\frac{4\mu_k\mathsf{e}(s)}{\varepsilon_\star},\qquad
\sum_{k'\le k}\eta_{k'}(s)\le\frac{3}{16}\,\mathsf{e}(s),
\end{equation}
are unaffected by the reserve. What the current radius pays at step $k$ is
$\eta_k(s)+\eta^{\mathrm{inc}}_k(s)$. Since $\eta^{\mathrm{inc}}_k(s)=\mu_k\mathsf{e}(s)/\varepsilon_\star$,
the first bound of \eqref{4.22:eq-radius-accounting-1} gives
\begin{equation}\label{4.22:eq-radius-accounting-2}
\eta_k(s)+\eta^{\mathrm{inc}}_k(s)\le\frac{5\mu_k\mathsf{e}(s)}{\varepsilon_\star},
\end{equation}
and, by the radius budget \eqref{4.22:eq-budget-closes-a}, the total paid over all steps is at most
$\tfrac{7}{32}\,\mathsf{e}(s)$.

\emph{Users of the entries.} The entries are drawn on as follows. This list adds no entry and no
amount to \eqref{4.22:eq-radius-accounting-a}. Throughout the list, and in the proof below,
$\mu:=\mu_k\mathsf{e}(s)/\varepsilon_\star$ denotes the margin amount of one presentation. For each
user we name the pieces it bounds, the map
that rides along at $s$, the disc on which it works, the entries it draws, and its obligation.
\begin{itemize}
\item The correlation theorem, with its lemma at $(\mathrm{T})$, on units that are not wrapped. It
bounds the $\mathsf{t}$-pieces $v_p$. The riding map is $\mathfrak{u}_p(\mathsf{t},\sigma)$. The disc
is the disc of radius $\rho'\ge\mathsf{w}\,\eta_p^{-1/2}$ supplied by the margin lemma at the
fluctuation site, inside $|\mathsf{t}|\le\mathsf{w}/\eta_p$, that is, the disc of radius $\rho_p$
at $(\mathrm{T})$;
here $\mathsf{w}$ is the disc width at $(\mathrm{T})$ and $\eta_p$ the positional quantity attached
to the piece $p$ by the correlation theorem. It draws
$(\ell 1)+(\ell 2)$. Its obligation is clause (iv) of
Lemma~\ref{4.22:lem-axialising-map} in its restricted form at $(\mathrm{T})$, together with the
primed clause (vi)$'$ of the liveness conditions.
\item The clause governing $\mathsf{s}$-pieces, written (S) below, at $(\mathrm{T})$ and at
$(\mathrm{L})$, on wrapped units and, at $(\mathrm{L})$,
also on every interior first-part unit. It bounds the $\mathsf{s}$-pieces of the last case of
\eqref{4.22:eq-presentation-radius-a} and the values. The riding map is $\mathsf{u}(\mathsf{s})$,
respectively $h\mathfrak{Q}^{\mathfrak{b}}(\vec{\mathsf{s}})$, with $h$ the carried gauge. It works
at amplitude one, on
$|\mathsf{s}|\le e^{\kappa_1}$, with no $\mathsf{t}$. It draws $(\ell 1)$ (not $(\ell 2)$), and
$(\ell 5)$ or $(\ell 6)$ with $(\ell 7)$, the latter idle. Its obligation is clause (iii) of the
debit lemma in its form for clause (S).
\item The majorant estimate at $(\mathrm{P})$ under which
Corollary~\ref{4.22:cor-preparatory-chart} is stated, on every unit. It bounds the constituents
$c_{A,Q}$. The
riding map is $\mathfrak{u}^{(l)}_{t,\sigma}(b)$. The disc is $|t_{\square}|\le R^{\natural}_{\square}$,
already at half rate. It draws $(\ell 3)$ or $(\ell 4)$ and no margin. Its obligation is clause
(iv) of Lemma~\ref{4.22:lem-axialising-map} at $(\mathrm{P})$.
\item The presentation's treatment of brackets. It bounds the increment $\Delta^A$ of
Corollary~\ref{4.22:cor-bracket-increment} of a bracket in the presentation's own filing, the
second case of \eqref{4.22:eq-presentation-radius-a}. The riding map is
$h\mathfrak{Q}^{\mathfrak{b}}e^{\tau Z'}g'$, with $h$ the carried gauge at the current member of
the chain and $\mathfrak{Q}^{\mathfrak{b}}$ the product of presenter maps up to the previous member,
as in that corollary, and $Z'$ the direction of the layer increment there. It works by the Cauchy
estimate in $g'$, in the
room $\mu$. It draws $(\ell 5)$ or $(\ell 6)$, with $(\ell 7)$. Its obligation is clause (vi) of
the liveness conditions.
\item The rider of the correlation theorem at $(\mathrm{L})$. It bounds the bracket pieces in the
correlation theorem's filing, the third case of
\eqref{4.22:eq-presentation-radius-a}: plain second-part units, and plain first-part units $v$
with $K_v\setminus Z\ne\emptyset$, these last only in the round of the fourth presenter map
$\mathfrak{w}_4$,
where $K_v$ is the smallest domain in $\mathbf{D}_k$ containing
$X_v$. The riding map is
$h_q\mathfrak{Q}^{\mathfrak{b}}\mathfrak{w}\bigl[(t+\mathsf{t}\hat\chi_{\square})\mathbb{H}\bigr]$,
with $\mathfrak{Q}^{\mathfrak{b}}$ at the previous member of the chain and $\mathfrak{w}$,
$\mathbb{H}$ at the current one. Here $\hat\chi_{\square}$ is the cut-off of the block $\square$
used by the rider. The disc is $|\mathsf{t}|\le\rho_p$, where
$\rho_p=\mathsf{w}_L(C_b\eta_p)^{-1/2}$, with $\mathsf{w}_L$ the disc width at $(\mathrm{L})$. It
draws $(\ell 5)$ or $(\ell 6)$, with $(\ell 7)$ up to
the full $\mu$. Its obligation is the condition that keeps this rider within the margin.
\item The deep-unit lemma. It bounds the $\mathsf{s}$-pieces and the value of a deep first-part
unit, with $X\subset\mathsf{C}_C$, plain or wrapped, through each of the four links numbered in
that lemma, at the members
$\mathfrak{b}_3$ or $\mathfrak{b}_4$ of the presentation chain. The riding map is
$\mathfrak{Q}^{\mathfrak{b}}(m')\exp[\lambda Z(m,m')]\,g'$, with $Z(m,m')$ the generator of the
link from $m$ to $m'$ in the deep-unit lemma,
$\mathfrak{Q}^{\mathfrak{b}}=\mathfrak{u}_3$ or $\mathfrak{u}_4\mathfrak{u}_3$, $h_q=1$,
$\mathfrak{w}_1=\mathfrak{w}_2=1$, in class ($\beta$), that of interior pointwise variables. The
disc is $|\lambda|<\rho_X$, where
$\rho_X=r^0e^{\frac12\delta_PD_X}$ is the square-root radius, with $r^0$ the base radius of the
$\lambda$-disc, and $\vec{\mathsf{s}}$ ranges over
$|\mathsf{s}|\le e^{\kappa_1}$, with no $\mathsf{t}$. It draws $(\ell 6)+(\ell 7)$, the latter up
to $\tfrac12\mu$. Its obligation is the condition that keeps this lemma within the margin. The
distances it needs, measured in $d^{\wedge}$ to the complement of the core $\mathsf{C}_C$, are given
by clause (ii)($\beta$) of the debit lemma combined with the shell-distance estimate in
$d^{\wedge}$; they are also given, independently, by the lemma on distances to the complement of a
core, a self-contained lemma used for the $\lambda$-disc only.
\item The kept Gaussian carried by a live arrest $\mathfrak{a}$. Its pieces
$\mathsf{q}_{\mathsf{p}}$ ride
with the site's own map at $b_-$, on the site's own disc, and draw the site's own entries. They
enter as a summand of the site's bound.
\item The two kept members $\mathfrak{K}(Y_i)$. They are treated under clause (S) as terminal
members of their lineage; when they are healed deep inside a core they are treated by the item on
the deep-unit lemma above, at $X=Y_i^{\sharp}$. They ride with the site's own map at $\Sigma_W$, on
the site's own disc,
never on a $\mathsf{t}$-disc, and draw the site's own entries in the unit
$\varepsilon^{\mathrm{fr}}(W)$. Their obligations are clauses (ii)$_2$ and (iv) of the kept-member
lemma. These clauses rest on the two-point bound on $V$, on the two bounds on the kept
quadratic form $\mathsf{Q}_i$, and on the holomorphy lemma. The bound on $\mathfrak{O}(v)$ is paid by
the
estimate of the terminal members of a lineage through their own collar, not by the entries of
\eqref{4.22:eq-radius-accounting-a}.
\end{itemize}
\end{definition}

\begin{remark}[Completeness of the accounting]
No operation other than those of the debits $(\ell 1)$, $(\ell 3)$--$(\ell 6)$ rotates a stored
variable on the background pairs at which the entries of \eqref{4.22:eq-radius-accounting-a} are
made; at the members $\mathfrak{b}_0,\mathfrak{b}_1$ of the presentation chain this rests on the
statement on rotations of stored variables at the first two members of the presentation chain.
\end{remark}

\begin{proof}[Proof of the completeness assertion]
Let a stored variable be given, with law site $s$. A rotation at $s$ can come from a parameter site,
from an operation of the large-field step, from a real gauge rotation, from the presenter maps of a
presentation, or from the discs used by
the lemmas that provide bounds at a presentation. We take these in turn. At the members
$\mathfrak{b}_0,\mathfrak{b}_1$ of the presentation chain, it is the statement on rotations of
stored variables at the first two members of the presentation chain that places every operation
acting on those pairs among these cases.

First, the parameter sites. These are the fluctuation site $(\mathrm{T})$ and the stages
$(\mathrm{P})$. They rotate stored variables only through the rule under which
Corollaries~\ref{4.22:cor-preparatory-chart} and~\ref{4.22:cor-older-law-sites} are stated: the
site's axialising map of Lemma~\ref{4.22:lem-axialising-map} rides along the stored variables. At
$(\mathrm{T})$ this rotation is the entry $(\ell 1)$. At the stages of the chain whose range
contains the level of the slot it is $(\ell 3)$, and at the stages of every other chain it is
$(\ell 4)$.

Second, the operations of the large-field step do not rotate a stored variable. These are the
operation $\mathbf{T}'_k$, the terms of \cite[(1.98)]{B-CMP122b}, the freezing at an arrest, and the
re-filing of terms. By the step on operations in the proof of
Proposition~\ref{4.22:prop-presentation-radius}(b), the
surviving variables and the chart variable $g'$ enter these operations only through the presented
functions. The measure $d\nu_i$, the operator of the kept quadratic form, the combinatorial
coefficients and the interpolation are free of them.

Third, rotations by $G$-valued maps cost no radius. These are the rotations of clause (iii) of the
lemma on orbit data at an arrest and those used in the joint-invariance theorem. If
$g'(s)\in G$, then $g'(s)$ is
unitary, so
$\|g'(s)\|=\|g'(s)^{-1}\|=1$, and
\begin{equation*}
\operatorname{pol}(g'(s))=\log\max\{\|g'(s)\|,\|g'(s)^{-1}\|\}=0.
\end{equation*}

Fourth, at the presentation $(\mathrm{L})$ the rotation of the presented unit by its presenter maps
is the entry $(\ell 5)$ when $s\notin Z$ and the entry $(\ell 6)$ when $s\in Z_h$. The increment discs
of the lemmas that provide these bounds lie inside the margin $(\ell 7)$. These discs are the Cauchy
disc in $g'$ of
Corollary~\ref{4.22:cor-bracket-increment}(b), the $\mathsf{t}$-disc of the correlation theorem's
rider, and the $\lambda$-disc of the deep-unit lemma. By the lemma on one margin for many bounds,
each of these estimates bounds one piece on the whole chart of radius
$\varepsilon-\eta_0-\mu$. Here $\eta_0$ is the member of $\eta^L_k(s)$ belonging to the
presentation and $\mu$ is the margin amount fixed before the list of users. Each estimate uses of
the radius
$\varepsilon$ only that its auxiliary points stay in the chart of radius $\varepsilon$. It therefore
uses the room between the charts of radii $\varepsilon-\eta_0-\mu$ and $\varepsilon-\eta_0$,
changes no chart, and so appears in the list of users above, not as a further entry.

Finally, once an account is closed by (x1) or (x2), the variable is no longer a stored variable at
$s$: it is either integrated, or it belongs to the data $\Phi_{W_2}$ of a frame, whose account is
the new one opened by (c2). Hence the only rotations of a stored variable at $s$ are those of
$(\ell 1)$, $(\ell 3)$--$(\ell 6)$.
\end{proof}

\begin{lemma}[One margin per presentation serves every piece]\label{4.22:lem-shared-margin}
Assume the following statements:
\begin{itemize}
\item clauses (a) and (b) of Proposition~\ref{4.22:prop-presentation-radius};
\item clauses (i) and (iii) of the transfer lemma (Lemma~\ref{4.22:lem-majorant-transfer});
\item clause (b) of Corollary~\ref{4.22:cor-bracket-increment};
\item the margin lemma for disc variables, by which a piece may be bounded on the disc of a
  variable $\mathsf{t}$, or $\lambda$, that draws the share $\theta\mu$ of the margin $\mu$ of its
  presentation, and the margin lemma for several devices;
\item the device-assignment rule, which files the pieces of every wrapped unit and of every
  first-part unit with $X_v\subset Z$ as $\mathsf{s}$-pieces, and the bracket pieces of plain
  units, second-part or first-part with $K_v\setminus Z\neq\emptyset$, on the disc of
  $\mathsf{t}$ (both kinds of piece in the sense of \eqref{4.22:eq-presentation-radius-a}), where
  $K_v$ is the smallest domain of $\mathbf{D}_k$ containing $X_v$; in particular neither the
  partition-of-unity factors of the cover nor the variable $\mathsf{t}$ occur on a unit with
  $X_v\subset Z$ or on a wrapped unit;
\item the deep-unit decay lemma, whose bound is obtained by the Schwarz lemma in $\lambda$ and
  Cauchy estimates in $\mathsf{s}$, and which is applied to $\mathsf{t}$-free pieces only;
\item the bracket-piece estimate on the disc of $\mathsf{t}$, which takes no Cauchy estimate in the
  rotation variable $g'$.
\end{itemize}
Let $v$ be a unit with an orbit datum $f_v^{\sharp}$ of radius $\varepsilon$, presented once at
step $k$, with gauge $h$ and riding rotation $\mathsf{u}$ of the presentation. At each law site $s$
let $\eta_0(s)$ be the member of $\eta^L_k(s)$ belonging to this
presentation, $\mu(s):=\mu_k\mathsf{e}(s)/\varepsilon_{\star}$, and
$\varepsilon':=\varepsilon-\eta_0-\mu$. The three devices are: the Cauchy estimate in $g'$ on the
chart, the disc of $\mathsf{t}$, and the disc of $\lambda$. Then the following hold.
\begin{enumerate}
\item[(i)] Every piece of this presentation is a function on the product of its parameters and
  $\mathfrak{G}_{\varepsilon'}(X_v)$, obtained by the operators $D_\pi$ of the identity
  \eqref{4.22:eq-majorant-transfer-a} from
  \begin{equation}\label{4.22:eq-shared-margin-1}
  F^{\sharp}(\tau_b;g')=f_v^{\sharp}\bigl(\mathsf{R}(\tau_b)^h;\mathbf{A},\boldsymbol{\Phi},
  h\,\mathsf{u}(\tau_b)g'\bigr),\qquad \tau_b\in\mathfrak{P}^{\mathrm{base}};
  \end{equation}
  it is stored with the chart $\varepsilon'$, whichever device bounds it.
\item[(ii)] Each of the following estimates of \eqref{4.22:eq-presentation-radius-a} bounds one
  piece on the whole of $\mathfrak{G}_{\varepsilon'}$: the bound of a bracket increment by
  Corollary~\ref{4.22:cor-bracket-increment}(b); the bound of a bracket piece on the disc of
  $\mathsf{t}$, by the margin lemma for disc variables at share $\theta=1$; the bound of an
  $\mathsf{s}$-piece of a deep first-part unit, by the same lemma at share $\theta=1/2$. Of
  $\varepsilon$ each of them uses only this: for $g'\in\mathfrak{G}_{\varepsilon'}$ the auxiliary
  points
  \begin{equation}\label{4.22:eq-shared-margin-2}
  g'_{\circ}e^{\tau\mathsf{Z}'},\qquad h\,\mathsf{u}(\mathsf{t})g',\qquad \mathsf{u}_\lambda g'
  \end{equation}
  stay in $\mathfrak{G}_{\varepsilon}$. Thus the margin is a room between the charts $\varepsilon'$
  and $\varepsilon-\eta_0$, not a consumable: any number of pieces may each use it.
\item[(iii)] No piece carries two of the three devices.
\item[(iv)] If $n$ devices, drawing the shares $\theta_1,\dots,\theta_n$ of $\mu$, were run on one
  piece in one estimate, the margin lemma for several devices would require
  $\sum_{i}\theta_i\le(\text{room})/\mu$; on $\mathfrak{G}_{\varepsilon_{k+1}}$ the room at a
  presentation is at least $\mu+\eta^{\mathrm{inc}}_k=2\mu$.
\end{enumerate}
\end{lemma}

In \eqref{4.22:eq-shared-margin-2}, $g'_{\circ}$ denotes the rotated point of
Corollary~\ref{4.22:cor-bracket-increment} and $e^{\tau\mathsf{Z}'}$ the one-parameter deformation
in which that corollary takes its Cauchy estimate, with complex parameter $\tau$ (distinct from the
parameter $\tau_b$ of \eqref{4.22:eq-shared-margin-1}) and direction $\mathsf{Z}'$;
$\mathsf{u}(\mathsf{t})$ and
$\mathsf{u}_\lambda$ are the rotations carried on the discs of $\mathsf{t}$ and of $\lambda$.

\begin{proof}
(i) Apply Proposition~\ref{4.22:prop-presentation-radius}(a) to the presentation of $v$, with
$\eta_0$ as the bound on the presentation's rotation and $\mu$ as the margin. It gives that
$F^{\sharp}$ of \eqref{4.22:eq-shared-margin-1} is an orbit datum of radius
$\varepsilon-\eta_0-\mu=\varepsilon'$. Clause (i) of Lemma~\ref{4.22:lem-majorant-transfer},
applied to the expansion that runs $v$ with this orbit datum, gives every piece as $D_\pi$ applied
to $F^{\sharp}(\tau_b;\cdot)$, $\tau_b\in\mathfrak{P}^{\mathrm{base}}$, hence as a function on the
product of its parameters and $\mathfrak{G}_{\varepsilon'}(X_v)$. Clause (iii) of the same lemma
gives that the piece is stored with the chart of the orbit datum, $\varepsilon'$. This chart is
fixed by $\varepsilon$, $\eta_0$ and $\mu$ alone, none of which refers to the device used to
bound the piece.

(ii) The proof of Corollary~\ref{4.22:cor-bracket-increment}(b) reads, for $g'$ in its chart,
\begin{equation*}
\operatorname{pol}(g'_{\circ})\le\eta^L+\operatorname{pol}(g')<\varepsilon-\mu_k\mathsf{e}/\varepsilon_{\star},
\end{equation*}
so that the room left for the Cauchy deformation $e^{\tau\mathsf{Z}'}$ is at least
$\mu_k\mathsf{e}/\varepsilon_{\star}=\mu$. The proof of the margin lemma for disc variables reads,
for $g'\in\mathfrak{G}_{\varepsilon'}$, that is for $\operatorname{pol}(g')<\varepsilon-\eta_0-\mu$,
\begin{equation*}
\operatorname{pol}\bigl(h\,\mathsf{u}(\mathsf{t})g'\bigr)\le\eta_0+\mu+\operatorname{pol}(g')<\varepsilon,
\end{equation*}
and this is the form in which the bracket pieces (share $1$) and the deep $\mathsf{s}$-pieces
(share $1/2$, with $\lambda$ in place of $\mathsf{t}$) use it. Each of these is a statement about
one function, the piece being bounded, on the whole of $\mathfrak{G}_{\varepsilon'}$; its only
use of $\varepsilon$ is that the auxiliary points \eqref{4.22:eq-shared-margin-2} lie in
$\mathfrak{G}_{\varepsilon}$, where $f_v^{\sharp}$ is defined; and none of them changes a chart.
Since
\begin{equation*}
(\varepsilon-\eta_0)-\varepsilon'=\mu,
\end{equation*}
the margin is the gap between the charts $\varepsilon'$ and $\varepsilon-\eta_0$. An estimate that
reads it leaves the chart $\varepsilon'$ of every other piece unchanged, so any number of pieces
may each use it.

(iii) Consider first an $\mathsf{s}$-piece; it has the form
$D_{\mathsf{y}_4}D_{\mathsf{y}_3}D_{\mathsf{y}_2}D_{\mathsf{y}_1}F^{\mathfrak{b}}$. By the
device-assignment rule it belongs to a wrapped unit or to a first-part unit with $X_v\subset Z$,
and on such a unit neither the partition-of-unity factors of the cover nor $\mathsf{t}$ occur; so
it holds no $\mathsf{t}$. Its bound takes no Cauchy estimate in $g'$: it is bounded by
Lemma~\ref{4.22:lem-majorant-transfer} with $\mathfrak{P}^{\mathrm{inc}}=\emptyset$, through the
fundamental theorem of calculus in the $\mathsf{s}$-variables, and, when its unit is deep, by the
deep-unit decay lemma, which uses the Schwarz lemma in $\lambda$ and Cauchy estimates in
$\mathsf{s}$ only. Hence it carries at most one device, the disc
of $\lambda$, and that only when its unit is deep.

Next consider a bracket piece on the disc of $\mathsf{t}$. By the device-assignment rule it belongs
to a plain unit that is either second-part or first-part with $K_v\setminus Z\neq\emptyset$, and
it belongs to a unit with $X_v\not\subset Z$. Now $Z$ is a union of level-$k$ $M$-cubes: in
\cite{B-CMP119}, $Z_k=\Lambda_k^c$ with $\Lambda_k$ a union of cubes of side $MR_k$, where $R_k$
is as in Lemma~\ref{4.4:lem-collar-avoidance}, and \cite{B-CMP122a} states that $Z''_k$ and
$Z''_k\setminus Z''_{k-1}$ are unions of $M$-cubes of the lattice of spacing $\eta$ of that paper
(there $\eta$ is Ba{\l}aban's lattice spacing, not a debit). Since $K_v$ is the smallest domain of
$\mathbf{D}_k$ containing $X_v$, it follows that $X_v\subset Z$ if and only if $K_v\subset Z$; hence
$K_v\setminus Z\neq\emptyset$ gives $X_v\not\subset Z$, and the two cases of the
device-assignment rule are complementary. So $X_v$ lies
in no core $\mathsf{C}_C\subset Z$, and the deep-unit decay lemma is not applied to this piece; in
any case that lemma is applied to $\mathsf{t}$-free pieces only. The bracket-piece estimate takes
no Cauchy estimate in $g'$. Hence the only device of such a piece is the disc of $\mathsf{t}$.

Finally, a bracket increment $\Delta^A_s$ holds neither the disc of $\mathsf{t}$ nor $\lambda$.
In the part $\mathfrak{B}^{\mathrm{rec}}_s$ of the bracket, in which both members carry the common
rotation $\mathfrak{Q}^{\mathfrak{b}_{s-1}}g'$, this rotation is a fixed parameter, so neither disc
is run there. That $\Delta^A_s$ carries neither $\mathsf{t}$ nor $\lambda$ is the property of orbit
data that governs the increment $\Delta^A_s$. Its only device is therefore the Cauchy estimate in
$g'$ of Corollary~\ref{4.22:cor-bracket-increment}(b).

(iv) The first assertion is the margin lemma for several devices. By
Proposition~\ref{4.22:prop-presentation-radius}(b),
\begin{equation*}
\varepsilon_{k+1}=\varepsilon-\eta^L_k-2\mu-\eta^{\mathrm{inc}}_k .
\end{equation*}
Hence the room at the presentation on $\mathfrak{G}_{\varepsilon_{k+1}}$ is
\begin{equation}\label{4.22:eq-shared-margin-3}
\begin{split}
(\varepsilon-\eta_0)-\varepsilon_{k+1}&=(\eta^L_k-\eta_0)+2\mu+\eta^{\mathrm{inc}}_k\\
&\ge2\mu+\eta^{\mathrm{inc}}_k\ge\mu+\eta^{\mathrm{inc}}_k .
\end{split}
\end{equation}
Here the first inequality holds because $\eta_0$ is one of the non-negative members of
$\eta^L_k$, so $\eta^L_k-\eta_0\ge0$, and the second because $\mu\ge0$. By the definition of the
reserve in \eqref{4.22:eq-radius-accounting-a},
$\eta^{\mathrm{inc}}_k(s)=\mu_k\mathsf{e}(s)/\varepsilon_{\star}=\mu(s)$. The last member of
\eqref{4.22:eq-shared-margin-3} is therefore $2\mu$.
\end{proof}

\begin{definition}[Depth of a law site and constants of the far rotations]\label{4.22:not-law-site-depth}
The following letters are fixed for the lemma on live law sites and far debits, which follows, and for the statements that depend on it.

\begin{enumerate}
\item \emph{Depth of a law site.} Let $s$ be the law site of a pointwise coordinate. Its depth is
\begin{equation}\label{4.22:eq-law-site-depth-1}
m(s):=\min\{\,j:\ s\in Z_j\,\},
\end{equation}
where $Z_j$ denotes the region of level $j$ in the family of regions $(Z_j)_j$ used in this chapter. Let $s\in\Sigma_W$ be a law site of a frame $W$. Its depth is
\begin{equation}\label{4.22:eq-law-site-depth-2}
m(s):=k_W,
\end{equation}
where $k_W$ is the top level of $W$.

The choice \eqref{4.22:eq-law-site-depth-2} is forced by the way a processed frame is re-blocked. By \cite[(1.33)]{B-CMP122b}, the determining set $\mathbf{B}_k$ equals $\mathbf{B}''_k$ on $Z^c$ and equals $\{Z^{(k)}\}$ on $Z$. A processed frame $W$ is therefore re-blocked at its top level, and the shells of $W$ formed before the operation $\mathbf{R}$ are erased. The scaled distance $d_{\mathrm{sc}}$ defined below is measured relative to the current determining set \cite{B-CMP122a}, so it counts none of these shells. The distance from $\Sigma_W$ to the complement $W^c$ is nevertheless still at least $10\check{R}_{k_W}$. This follows from \cite{B-CMP122a} and from clause (iii) of the frame lemma, Lemma~\ref{4.22:lem-frame-lemma}. This separation has not been used in \eqref{4.22:eq-law-site-depth-2} and remains available.

\item \emph{Scaled distance.} The symbol $d_{\mathrm{sc}}$ denotes the scaled distance. It is measured in $M$-cubes of the local scale, relative to the current determining set. In \cite{B-CMP122a} the corresponding distances are defined in terms of $M_1$-cubes on the corresponding scales. The weighted contour distance $d^{\wedge}$ and the scaled distance are compared by
\begin{equation}\label{4.22:eq-law-site-depth-3}
d^{\wedge}\ \ge\ \mathsf{w}\, d_{\mathrm{sc}},
\end{equation}
where $\mathsf{w}>0$ is the width that $d^{\wedge}$ charges for each crossed zone, as in the comparison estimate between the two distances.

\item \emph{The layer rate.} The letter $\delta_1$ denotes the decay rate, per unit of $d_{\mathrm{sc}}$, of the layer estimate for the large-field layers. Elsewhere in the book the same letter denotes the decay rate of the kernel of the operator $T$. That is a different number.

\item \emph{The far rate.} Let $\delta_L$ be the decay rate of the localisation estimate for the presentation maps; $\delta_L$ depends only on $d$ and $L$. Let $C_L$ be the constant of the same estimate. Put
\begin{equation}\label{4.22:eq-law-site-depth-4}
\delta_{\natural}:=\min\{\delta_1,\ \delta_L,\ 1\}.
\end{equation}

\item \emph{The constant of the presentation radius.} A \emph{$(d,L,M)$-number} is a number depending only on $d$, $L$ and $M$. Since $\kappa$ and $\kappa_1$ are fixed before $M$, such a number may also depend on $\kappa$ and $\kappa_1$. We write $K_{\mathfrak{Q}}$ for the $(d,L,M)$-number such that
\begin{equation}\label{4.22:eq-law-site-depth-5}
\varepsilon_{\mathfrak{Q}}(k)\ \le\ K_{\mathfrak{Q}}\,\alpha^{+}_k .
\end{equation}
The bound on $\varepsilon_{\mathfrak{Q}}(k)$ from which \eqref{4.22:eq-law-site-depth-5} is read is a linear expression in $\tilde\alpha_0+\tilde\alpha_1$, in $\delta'_k$ and in $\alpha_0+\alpha_1$.

\item \emph{The constant of the far debits.} Put
\begin{equation}\label{4.22:eq-law-site-depth-6}
C_{\natural}:=2e^{2}\cdot\max\{\,C_L K_{\mathfrak{Q}},\ K_u B_H,\ K_u\,\}.
\end{equation}
Here $B_H$ is read on the parameters $\sigma$ with $|\sigma|\le e^{\kappa_1}$. It then contains the factor $e^{16\kappa_1}$ of \cite[(1.21)]{B-CMP116}. Since $\kappa_1$ is fixed before $M$, a factor $e^{16\kappa_1}$ is admitted in a $(d,L,M)$-number in the sense fixed above.

\emph{Convention for the families of presentations.} Inside $C_{\natural}$, and only there, two replacements are made. The letter $B_H$ is read as $B_H^S$, and $K_{\mathfrak{Q}}$ is read as $K_{\mathfrak{Q}}B_H^S/B_H$, by the definition of $B_H^S$. With these replacements $C_{\natural}$ serves for the bounds on the families of presentations in the lemma that follows. Everywhere else the letters $B_H$, $\varepsilon_{\mathfrak{Q}}$ and $\varepsilon_{\mathrm{live}}$ keep their meaning, in particular in the floor conditions of the parameter order. The replacements are absorbed by one re-maximisation of the constants. That re-maximisation is made before $\gamma$ is fixed, inside the three ceilings on the coupling that are stated in supremum form: clause (vi) of the conditions on live law sites, the ceiling attached to the two-point bound of the kept members, and the ceiling attached to the frames.

\item \emph{Age of a law site.} For a level $k$ and a law site $s$ put
\begin{equation}\label{4.22:eq-law-site-depth-7}
\mathsf{a}_k(s):=\bigl(k-N_k-m(s)\bigr)_{+},
\end{equation}
where $(x)_+=\max\{x,0\}$ and $m(s)$ is given by \eqref{4.22:eq-law-site-depth-1} or \eqref{4.22:eq-law-site-depth-2}. At a law site $s\in\Sigma_W$ of a frame this reads
\begin{equation*}
\mathsf{a}_k(s)=(k-N_k-k_W)_{+}.
\end{equation*}
\end{enumerate}
\end{definition}

\begin{lemma}[The law-site lemma: positions, distances and far debits]\label{4.22:lem-law-site}
Let $k$ be a step. The depth $m(s)$ of a law site $s$, the scaled distance $d_{\mathrm{sc}}$ (in
$M$-cubes of the local scale), the rates and constants $\delta_{\natural}$, $\delta_L$, $C_{\natural}$,
$K_{\mathfrak{Q}}$ and the age $\mathsf{a}_k(s)=(k-N_k-m(s))_+$ are those of
Notation~\ref{4.22:not-law-site-depth}; recall $N_k=\check R_k$, and put $h:=k-N_k$. The
letter $Z$ denotes the large-field region on whose class the chain of large-field stages runs
at step $k$, and $Z_i$ the large-field region of level $i$. At a frame $W$ (a second-class
component) created at step $k_W$, with law sites $\Sigma_W$, one has $m(s)=k_W$ for
$s\in\Sigma_W$. For each level $j$, $\Omega_j$ denotes Ba{\l}aban's small-field region of level
$j$ and $\Omega''_j$ the region of \cite[(1.11)]{B-CMP122a}, with $\Omega''_j=\Omega_j$ for
$j\le h$; for a region $X$ of level $j$ and $n\ge1$, $X^{\sim n}$ is its enlargement by $n$
layers of cubes of the scale of level $j$, and $X^{\sim}:=X^{\sim1}$. The letter $\Lambda$ denotes
the region of the presentation at step $k$, in whose enlargement $\Lambda^{\sim}$ by one layer the
arguments of the presented terms are supported; $\Lambda_j$ denotes the corresponding region of
level $j$, and $\partial^+\Lambda_j$ its outer boundary. Finally $R_{\min}$ denotes a lower bound
for the range numbers $R_j$ of Lemma~\ref{4.4:lem-collar-avoidance} over the levels $j\le k$, such
that each shell of the
nest of large-field regions has $d_{\mathrm{sc}}$-thickness at least $10R_{\min}$ where no
large-field operation has re-blocked it \cite{B-CMP122a}. Assume the following; the coupling
ceilings and the coupling-flow bound of the sixth item are assumed at every index $j\le k$:
\begin{itemize}
\item the floors of the order of constants of this chapter;
\item the decay of the Landau layers at rate $\delta_1$ (the rate of
Notation~\ref{4.22:not-law-site-depth}) per unit of $d_{\mathrm{sc}}$, the
separation of far arguments, and the clause of the separation of far arguments giving
$d_{\mathrm{sc}}(s,\Lambda^{\sim})\ge 5\check R_k$ at live law sites outside $Z$;
\item the level floor, which gives $\delta_P\mathsf{w}\ge 4$, together
with the comparison $d^{\wedge}\ge\mathsf{w}\,d_{\mathrm{sc}}$ of the weighted contour distance
$d^{\wedge}$ with the scaled distance, $\mathsf{w}$ being the zone width of $d^{\wedge}$; recall
that $\delta_{\natural}\le1$ by Notation~\ref{4.22:not-law-site-depth};
\item $r\ge 2$; the definition of the scaffold's scale $\check R_j$, which is non-increasing in
$j$; the rule that the large-field operation acts on a component at the first step at which it
qualifies, with the nesting of the large-field regions;
\item the arrest-timing lemma, and the bookkeeping rules for attempts of the large-field
operation that are arrested or excluded on a component;
\item the coupling ceilings \eqref{4.22:eq-coupling-ceilings-g}, \eqref{4.22:eq-coupling-ceilings-h}
and \eqref{4.22:eq-coupling-ceilings-i}, and the coupling-flow bound;
\item the rim-removal statement for completed chains of large-field stages, with its companion
on the stage inputs;
\item the rim separation of the inheritance theorem;
\item the field-datum clause of the subfamily lemma;
\item the identity $U_{1,2}=U_2$ on the region $D_2$ defined in (i) below
\cite[(1.53)]{B-CMP122b};
\item the presentation statement at the indices $\mathfrak{b}\in\{\mathfrak{b}_2,\mathfrak{b}_3,
\mathfrak{b}_4\}$ of the presentation chain and at the pairs
$(\mathfrak{b}_2,\mathfrak{b}_3)$, $(\mathfrak{b}_3,\mathfrak{b}_4)$.
\end{itemize}
A \emph{reading} of a variable is one of the rotations that the three operations of the step,
the fluctuation integral $(\mathrm{T})$, the large-field stages $(\mathrm{P})$ and the
presentation $(\mathrm{L})$, apply at its law site. Then at step $k$ the following hold.
\begin{enumerate}
\item[(i)] (Which law sites are live, and where.) At $(\mathrm{T})$, every law site of a term of
the summation of boundary terms at finer zones lies in $Z_{k-1}$. At $(\mathrm{P})$, along a
chain on the class of $Z$, a law site $s$ of a constituent is of one of the following kinds:
\begin{itemize}
\item[(n)] $s$ is the site of a rim slot of $Z$ of level $j'\in[h,k-1]$;
\item[($\mathrm{f}_1$)] $s\in Z_h$;
\item[($\mathrm{f}_2$)] $s\notin Z$.
\end{itemize}
At $(\mathrm{L})$, a live law site $s$ of a presented term $\mathfrak{E}$ satisfies either
\begin{itemize}
\item[($\alpha$)] $s\notin Z$, or
\item[($\beta$)] $s\in Z_h\subset D_2$, where $D_2:=Z\setminus(\Omega''_{h+1})^{\sim}$; and then
$\mathfrak{E}$ is a first-part term, presented at
$\mathfrak{b}\in\{\mathfrak{b}_2,\mathfrak{b}_3,\mathfrak{b}_4\}$ only.
\end{itemize}
The law sites of frames and of arrested variables are never of kind (n).
\item[(ii)] (Distances.) At every law site $s$ of the respective kind the lower bounds before the
age increments are
\begin{align*}
(\mathrm{T})&:\ d_{\mathrm{sc}}(s,\Omega_k)\ge 8\check R_k,
&(\mathrm{f}_1)&:\ d_{\mathrm{sc}}(s,\mathcal{C}^{(l)})\ge\check R_k,\\
(\mathrm{f}_2)&:\ d_{\mathrm{sc}}(s,\mathcal{C}^{(l)})\ge 2\check R_k,
&(\alpha)&:\ d_{\mathrm{sc}}(s,\Lambda^{\sim})\ge 5\check R_k,\\
(\beta)&:\ d_{\mathrm{sc}}(s,\mathsf{S})\ge\check R_k-1,
\end{align*}
where $\mathcal{C}^{(l)}$ is the region on which the constituents of the $l$-th large-field stage
of the chain are localised (a union over the components of the class), $\Lambda^{\sim}$ the
enlargement of $\Lambda$ by one layer, and $\mathsf{S}$ the source set of the localisation of the
current beyond the cut, \eqref{4.22:eq-current-localisation-a}. In $(\mathrm{f}_2)$ the rim of $Z$
alone gives the first member when $m(s)=k$; when $m(s)\le k-1$ the rim of $s$ itself adds
$2\check R_k$. The
age increments are asserted at frame sites $s\in\Sigma_W$ only, and rest on
\eqref{4.22:eq-law-site-a}: the right-hand sides of $(\mathrm{T})$, $(\mathrm{f}_2)$ and $(\alpha)$
increase by $5R_{\min}(k-1-k_W)_+$, those of $(\mathrm{f}_1)$ and $(\beta)$ by
$5R_{\min}(h-k_W)_+$. At pointwise law sites the increments hold with $10R_{\min}$ in place of
$5R_{\min}$ inside a nest of shells that no large-field operation has re-blocked
\cite{B-CMP122a}, but they are not asserted across regions re-blocked by the large-field
operation \cite[(1.33)]{B-CMP122b}. No estimate uses them at pointwise sites except the sum over
law sites in the boundary-term corollary at complex pairs, and that sum is controlled by
$d_{\mathrm{sc}}$ itself; the pointwise estimates carry no age factor.
\item[(iii)] (Far debits.) In each of the cases $(\mathrm{T})$, $(\mathrm{f}_1)$, $(\mathrm{f}_2)$,
$(\alpha)$, $(\beta)$ one reading at $s$ costs at most
\begin{equation*}
\tfrac12 C_{\natural}\,\alpha^+_k\,e^{-\delta_{\natural}(\check R_k-2)}\,e^{-2\mathsf{a}_k(s)},
\end{equation*}
and at $(\mathrm{L})$ one has $h_q^{\mathfrak{b}}(s)=1$. Consequently, at every law site $s$ not of
kind (n),
\begin{equation*}
\eta^T_k(s)+\eta^P_k(s)+\eta^L_k(s)
\le (N_k+3)\,C_{\natural}\,\alpha^+_k\,e^{-\delta_{\natural}(\check R_k-2)}\,e^{-2\mathsf{a}_k(s)}.
\end{equation*}
The member $\check R_k-2$, obtained after passage from $s$ to the site $\tilde s$ at which the
local layer norm is evaluated, is common to all law sites. The factor $e^{-2\mathsf{a}_k(s)}$ is
used at frame sites, where $\mathsf{a}_k(s)=(k-N_k-k_W)_+$, and is replaced by $1$ at pointwise
sites. The bounds are proved for the modulus of the logarithm of each rotation, not only for its
polar size $\operatorname{pol}$.
\item[(iii)$^{S}$] (The same on the cell families, whatever the class of the
term.) On the cell families on which the statement running the step on families of cells
expands the step --- all cells at $(\mathrm{T})$, the cells disjoint from $\Lambda^{\sim}$ at
the first link of the presentation chain, the cells disjoint from $\Lambda$ at its second to
fourth links --- let the complex decoupling parameters $\mathsf{s}(\Delta)$, $\Delta$ running over
the cells of the family, satisfy $|\mathsf{s}(\Delta)|\le e^{\kappa_1}$. Then at every law site $s$
in the chart of a term
expanded on such a family (a term of the wrapped class, or an interior first-part term of the
plain or of the wrapped class) the bounds \eqref{4.22:eq-law-site-b} below hold. Here
$B_{\sharp}$ and $K_{\sharp}$ are the constants of the profile bound and of the field-datum bound
of the subfamily lemma ($B_{\sharp}$ is distinct from the flow constant of the glossary written
with the same letter); $b$ is the datum of the Landau layer at $(\mathrm{T})$, and $\Omega$ the
region in which $b$ is supported; $B^{\mathrm{top}}(s)$ is the top block of that layer at $s$,
which lies in the $M$-cube of $s$; $\mathrm{arg}_{s'}$ are the arguments of the $s'$-th factor of
the presenter map $\mathfrak{Q}^{\mathfrak{b}}$, and $h_q^{\mathfrak{b}}$ is the cut-off function of
the presentation at the index $\mathfrak{b}$, by which $\mathfrak{Q}^{\mathfrak{b}}$ is multiplied;
$\mathsf{J}_3$ is the current from which the argument that is a field datum is obtained, and
$\sup\ell^3|\mathsf{J}_3|$ is its size as measured in the field-datum bound of the subfamily lemma,
$\ell^3$ being the weight with which that bound measures the current; and $C_L$ is the constant
of \eqref{4.22:eq-current-localisation-a}.
\begin{equation}\label{4.22:eq-law-site-b}
\begin{aligned}
&(\mathrm{T}):\quad |\log\mathsf{u}(\mathsf{s})(s)|
\le K_uB_{\sharp}\sup|b|\;e^{-\delta_P d^{\wedge}(B^{\mathrm{top}}(s),\Omega)}\\
&\qquad\le K_uB_H^S\alpha_*\;e^{-\delta_P\mathsf{w}(d_{\mathrm{sc}}(s,\Omega)-1)},\\
&(\mathrm{L})(\alpha),\ s\notin Z:\quad
|\log(h_q^{\mathfrak{b}}\,\mathfrak{Q}^{\mathfrak{b}}(\vec{\mathsf{s}}))(s)|\\
&\qquad\le 1.01K_uB_{\sharp}\Bigl[\sum_{s'}\sup|\mathrm{arg}_{s'}|
+K_{\sharp}\sup\ell^3|\mathsf{J}_3|\Bigr]\,e^{-\delta_P\mathsf{w}(d_{\mathrm{sc}}(s,\Lambda^{\sim})-1)},\\
&(\mathrm{L})(\beta),\ s\in Z_h:\quad
K_u\Theta_0^{\mathrm{loc}}(\mathbb{H}_{s'})(\tilde s)
\le C_L\varepsilon_{\mathfrak{Q}}(k)\,e^{-\delta_L(d_{\mathrm{sc}}(s,\mathsf{S})-1)},
\qquad s'=3,4;\\
&\text{conversion:}\quad
e^{-\delta_P d^{\wedge}}\le e^{-\delta_P\mathsf{w}\,d_{\mathrm{sc}}}
\le e^{-\delta_{\natural}d_{\mathrm{sc}}};
\end{aligned}
\end{equation}
at $(\mathrm{P})$, no family being altered there, clause (iv) at $(\mathrm{P})$ of
Lemma~\ref{4.22:lem-axialising-map} holds as stated; and the bound at $(\mathrm{L})(\beta)$ is the
display \eqref{4.22:eq-current-localisation-a} as stated, up to the factor $B_H^S/B_H$ inside
$\varepsilon_{\mathfrak{Q}}$, which is paid inside $C_{\natural}$ as fixed in
Notation~\ref{4.22:not-law-site-depth}. The conversion in the last line of
\eqref{4.22:eq-law-site-b} holds since $d^{\wedge}\ge\mathsf{w}\,d_{\mathrm{sc}}$
and $\delta_P\mathsf{w}\ge4\ge1\ge\delta_{\natural}$, and each reading is at most
$\tfrac12 C_{\natural}\alpha^+_ke^{-\delta_{\natural}(\check R_k-2)}e^{-2\mathsf{a}_k(s)}$, with
$C_{\natural}$ as fixed for the families clause: this is the majorant of (iii) itself, with the
same rate $\delta_{\natural}$. Ages at frame sites: at $(\mathrm{T})$ and $(\mathrm{L})(\alpha)$ each
crossed zone contributes $\mathsf{w}$ to $d^{\wedge}$, by the clause on crossed zones of the
geometry lemma for the weighted contour distance; splitting the exponent in halves, with
$\tfrac12\delta_P\mathsf{w}\ge2$,
\begin{equation*}
e^{-\delta_P d^{\wedge}}\le e^{-2d_{\mathrm{sc}}}\,e^{-2\cdot\#\mathrm{zones}},
\qquad \#\mathrm{zones}\ge\mathsf{a}_k(s),
\end{equation*}
and no bound on the zone width beyond $\delta_P\mathsf{w}\ge4$ is used; at $(\mathrm{L})(\beta)$,
where the rate is $\delta_L$ per $M$-cube, the age factor comes from \eqref{4.22:eq-law-site-a}.
\item[(iv)] (Near debits, once.) A pointwise coordinate is of kind (n) in at most one chain of its
life, and in that chain
\begin{equation*}
\nu(s):=\sum_l\sup\operatorname{pol}(\mathfrak{u}^{(l)}(s))\le 3\beta K_u\sup_{j\le k}\alpha^+_j,
\end{equation*}
where $\beta$ is the schedule constant of Definition~\ref{4.4:def-post-step-space}.
\item[(v)] (Mergers.) Let $W_1\subset W_2\subset\cdots$ be a chain of enclosures, the $W_i$ being
second-class outcomes at steps $k_1<k_2<\cdots$, with components joined in between arbitrarily.
At step $k_{i+1}$ the law sites $\Sigma_{W_i}$ of the consumer terms are of kind $(\beta)$; the
account of chart radius at $\Sigma_{W_i}$ is closed there; and the new root, which lies on
$\partial^+\Lambda_{k_{i+1}}$, where $h_q^{\mathfrak{b}}\ne1$, is not read at step $k_{i+1}$. No
reading of any
variable costs the full rotation size $\varepsilon_{\mathrm{live}}(k)$ on the pairs at which that
size is fixed.
\end{enumerate}
\end{lemma}

\begin{proof}
We prove clauses (i) and (ii); clauses (iii), (iii)$^{S}$, (iv) and (v) are proved in the proof
that follows the present one in this section.

\emph{Clause} (i), \emph{at} $(\mathrm{T})$. A rim slot of level $j'\le k-2$ lies in $Z_{j'+1}$,
and $Z_{j'+1}\subset Z_{k-1}$ by the nesting of the large-field regions, since $j'+1\le k-1$.
Arrested pieces of steps $\le k-1$ lie in $\Omega^c_{k_a}\cap Z\subset Z_{k-1}$, $k_a$ being the
step of the arrest; frames of steps $\le k-1$ lie in $Z_{k_W}$, which is contained in $Z_{k-1}$ by
the same nesting, since $k_W\le k-1$.

\emph{Clause} (i), \emph{at} $(\mathrm{P})$ \emph{and} $(\mathrm{L})$. In a completed chain the rim
slots of $Z$ of levels $j'\in[h,k-1]$ are set to zero: this is the rim-removal statement for
completed chains, together with its companion on the stage inputs, in the form of
\cite[(1.60)]{B-CMP122a}, ``$A^{(n)}_k(U^{(n+1)}_k,A_j = 0)$''. Before completion these slots are
of kind (n). Next, let a live variable have its law site in $Z$ and level $j'<h$. Its bond lies
in $Z_{j'+1}\cap\Omega_{j'+1}\subset Z_{j'+1}\subset Z_h$, the last inclusion by nesting, since
$j'+1\le h$.

Consider a variable arrested at step $k_a$ inside $Z$. The arrest ``introduced new large fields''
\cite{B-CMP122a}, so that $k_a\le h$, by Ba{\l}aban's condition on the sequence of operations that
no large-field characteristic functions are included in the last $N_k$ of them, and by the
arrest-timing lemma. Ba{\l}aban uses this condition in exactly this way, (ii) being his numbering
of the conditions on the sequence of operations and not a clause of the present lemma: of the
last $N_k$
operations he writes ``by the condition (ii) no large field characteristic functions are included
in them'', and, after an arrest, ``we have introduced new large fields, therefore they do not
satisfy the condition (ii)'' \cite{B-CMP122a}. The clock is reset accordingly: ``a new large field
was introduced in the preparatory operations \ldots\ $K = R_{j+1}$ for $Z$'' \cite{B-CMP122b}.
Hence $k_a$ does not lie in $(h,k]$, that is, $k_a\le h$; the sharper bound $h-1$ is not
asserted. Therefore the region $W_{\mathfrak{a}}$ of the arrest $\mathfrak{a}$
satisfies $W_{\mathfrak{a}}\subset\Omega^c_{k_a}\cap Z\subset Z_h$.

Consider next a frame created at step $k_W$ inside $Z$. A second-class outcome is a creation
\cite{B-CMP122b}. By the rule that the large-field operation acts at the first step at which a
component qualifies, and by the same argument as for arrests, $k_W$ does not lie in $(h,k]$, that
is, $k_W\le h$; again the sharper bound $h-1$ is not asserted. Hence
\begin{equation*}
\Sigma_W\subset W\subset Z_{k_W}\subset Z_h .
\end{equation*}
Thus the law sites of frames and of arrested variables lie in $Z_h$: this is case
$(\mathrm{f}_1)$ at $(\mathrm{P})$ and case $(\beta)$ at $(\mathrm{L})$. In particular they are
never of kind (n).

Collecting these facts: at $(\mathrm{P})$, a law site of a constituent lying in $Z$ is either the
site of a rim slot of level in $[h,k-1]$ before completion, of kind (n), or lies in $Z_h$, of
kind $(\mathrm{f}_1)$; a law site outside $Z$ is of kind $(\mathrm{f}_2)$. At $(\mathrm{L})$ the
chain is completed, the rim slots of levels in $[h,k-1]$ are zero, and so a live law site lying
in $Z$ lies in $Z_h$, while a live law site outside $Z$ is of kind $(\alpha)$.

\emph{The inclusion} $Z_h\subset D_2$. The construction of \cite[(3.2)--(3.5)]{B-CMP119}
surrounds the region $Z'_h$ of that construction, which contains $Z_h$, successively by four
layers, one layer, one layer and two layers of $LMR_{h+1}$-cubes of the scale of level $h$, that
is by $4+1+1+2=8$ layers;
hence
\begin{equation*}
\Omega_{h+1}\subset\bigl((Z'_h)^{\sim8}\bigr)^c .
\end{equation*}
By \cite[(1.11)]{B-CMP122a}, where Ba{\l}aban writes ``$Z''_j=(\Omega_j^{\sim5})^c\cap Z$ for
$j=k-N_0,\dots,h+1$'' for the regions $Z''_j$ of $Z$ off $\Omega''_j$, with $N_0$ the number of
top levels of an attempt that are treated at the thin scales of \cite[(1.10)--(1.11)]{B-CMP122a},
one has $\Omega''_{h+1}\cap Z=\Omega^{\sim5}_{h+1}\cap Z$; and $\Omega''_j=\Omega_j$
for $j\le h$ \cite{B-CMP122a}. Counting layers, $Z_h$ is therefore separated from
$(\Omega''_{h+1})^{\sim2}$ by at least $8-5-2=1$ layer: $\Omega_{h+1}$ lies at least eight
layers away from $Z'_h\supseteq Z_h$, the passage to $\Omega''_{h+1}\cap Z=\Omega^{\sim5}_{h+1}\cap Z$
uses five of them, and the enlargement by two layers two more. The count is exact, and it is
Ba{\l}aban's own arithmetic in one family of cubes: on \cite[p.~195]{B-CMP122a} he writes
\begin{equation*}
\text{``}\Omega''_{h+1}=\Omega^{\sim5}_{h+1}=Z'^{\,c\sim3}_h,\ \text{hence}\
\Omega''^{\sim2}_{h+1}=\Omega^{\sim7}_{h+1}=(Z'^{\sim}_h)^c\text{''},
\end{equation*}
and on \cite[p.~179]{B-CMP122a} that ``$\Omega_j^{\sim n}$ are unions of $L^{-(k-j)}MR_j$-cubes'',
with, in \cite[(1.10)]{B-CMP122a},
\begin{equation*}
\text{``}Z''_{k-N_0+1}=(Z'_{k-N_0})^{\sim}=(\Omega^{\sim7}_{k-N_0+1})^c\cap Z\text{''}.
\end{equation*}
Thus one full layer of $MR_{h+1}$-cubes of the scale of level $h+1$ --- the same cubes as the
$LMR_{h+1}$-cubes of the scale of level $h$ above, a unit of level $h+1$ being $L$ units of level
$h$ by the quotation from p.~179 --- separates $Z'_h\supseteq Z_h$ from
$(\Omega''_{h+1})^{\sim2}$; and, again by \cite[p.~195]{B-CMP122a}, ``bonds intersecting
$\partial\Omega''^{\sim2}_{h+1}$ do not belong to $\mathbb{B}_0$'', $\mathbb{B}_0$ being the set of
bonds so denoted there. Since $(\Omega''_{h+1})^{\sim}\subset(\Omega''_{h+1})^{\sim2}$, it follows
that $Z_h\subset D_2$, the part of $Z$ off $(\Omega''_{h+1})^{\sim}$, on which $U_{1,2}=U_2$
\cite[(1.53)]{B-CMP122b}.

\emph{First-part terms.} Let $s\in Z_h$ be a live law site of a presented term $\mathfrak{E}$.
The term $\mathfrak{E}$ is a first-part term in Ba{\l}aban's division of the terms into two parts:
``the first part includes \ldots\ these terms \ldots\
for which their localization domains intersect the region $Z$, the second part includes the
remaining terms \ldots\ we expand the terms of the second part around their values at the
configuration $U_k$'' \cite{B-CMP122b}. A first-part term
undergoes the operations \cite[(1.59)--(1.60)]{B-CMP122b} and then the representations
\cite[(1.61)--(1.63)]{B-CMP122b}, which replace data of the boundary layer only; it remains a
first-part term in \cite[(1.70)--(1.71)]{B-CMP122b}. The presentation statement, an input of this
lemma, is stated for $\mathfrak{b}\in\{\mathfrak{b}_2,\mathfrak{b}_3,\mathfrak{b}_4\}$ and for the
pairs $(\mathfrak{b}_2,\mathfrak{b}_3)$, $(\mathfrak{b}_3,\mathfrak{b}_4)$; this is the range of
presentation asserted in case $(\beta)$. Finally, the terms attached to the region $Y_0$ (the
region so denoted in \cite{B-CMP122b}) in the $\mathbf{B}$-terms of the second group of that paper
carry the slots of their parent terms, and their domain
$X$ satisfies $X\cap Z=\emptyset$; their live law sites are therefore of kind $(\alpha)$. This
proves (i).

\emph{Clause} (ii): \emph{the shell bound at frame sites.} Let $s\in\Sigma_W$, let its account of
chart radius be open at step $k$, and let $k_W<m'\le k$. In a nest of shells that no large-field
operation has re-blocked, each level crossed by a path from $W$ to $Z_{m'}^c$ costs
$d_{\mathrm{sc}}\ge10R_{\min}$: ``Each renormalization step adds at least ten layers of
$MR_k$-cubes'' \cite{B-CMP122a}, and the shells are nested; the rims of
\cite[(3.2)--(3.5)]{B-CMP119} are not used here. The shells of levels $k_W+1,\dots,m'$ round $W$,
$m'-k_W$ of them, are of that kind, except those inside the window $[h_2,k_2-1]$, of $N_{k_2}$
levels, of an attempt made at step $k_2$ and arrested or excluded on the component of $W$ (the
arrest-timing lemma and the bookkeeping rules for arrested and excluded attempts); a completed
attempt closes the account, by clause (v) of the frame lemma
(Lemma~\ref{4.22:lem-frame-lemma}). Inside such a window the top levels, at most $N_0$ of them,
are the thin strata ``separated by one layer of $M$-cubes \ldots\ one layer of $L^{-1}M$-cubes,
and so on'' \cite{B-CMP122a}, each of $d_{\mathrm{sc}}$-thickness at least $1$, while the levels
$\le k_0$, $k_0$ being the highest level of the window below these strata, keep their own scales,
each of thickness at least $10R_{\min}$; at $i=k_0$ the kept annulus
$\Gamma''_{k_0}=\Omega^{\sim5}_{k_0}\setminus\Omega^{\sim7}_{k_0+1}$ has thickness at least
$(L+7)R_{\min}\ge10R_{\min}$, since the blocking factor $L$ is an odd integer greater than one.
Windows of distinct attempts are disjoint:
an arrest is a creation, so that $k_2\le h_3$ for the next attempt $[h_3,k_3-1]$ on the
component, by the arrest-timing lemma and the argument in the proof of (i); moreover $k_W\le h_2$.
A window crossed only in part at the end of the path therefore contributes its thick levels
first; this holds also when $k_2=k$, for a component excluded earlier in the same step. Since
$2N_0\le N_{k_2}$ at every step by \eqref{4.22:eq-coupling-ceilings-i}, at least half of the
$m'-k_W$ levels crossed are thick, each contributing at least $10R_{\min}$, and therefore
\begin{equation}\label{4.22:eq-law-site-a}
d_{\mathrm{sc}}(W,Z_{m'}^c)\ge 5R_{\min}(m'-k_W).
\end{equation}

\emph{The five cases.} $(\mathrm{T})$: by (i), $s$ lies in $Z_{k-1}$, and $\Omega_k$ is separated
from $Z_{k-1}$ by the eight layers of \cite[(3.2)--(3.5)]{B-CMP119} at level $k-1$, that is by
$8LR_k$ $M$-cubes of that scale; this gives the first member. At frame sites the age increment
follows from \eqref{4.22:eq-law-site-a} with $m'=k-1$ (there is nothing to add when
$k-1\le k_W$).

$(\mathrm{f}_1)$: here $s\in Z_h$ and $\mathcal{C}^{(l)}\subset\Omega^c_{j+1}\setminus Z''_{j+1}$
with $j\ge h$. The regions $Z''_j$ increase \cite{B-CMP122a}, so $Z''_{j+1}\supset Z''_{h+1}$, and
$Z''_{h+1}=(\Omega^{\sim5}_{h+1})^c\cap Z$ contains $Z_h$ together with $8-5=3$ layers round it.
Hence $d_{\mathrm{sc}}(s,\mathcal{C}^{(l)})\ge LR_{h+1}\ge\check R_k$, the last inequality because
$\check R_j$ is non-increasing in $j$. At frame sites the age increment follows from
\eqref{4.22:eq-law-site-a} with $m'=h$.

$(\mathrm{f}_2)$: here $s\notin Z$ and $\mathcal{C}^{(l)}\subset\Omega_k^c\cap Z$. A path from
$\mathcal{C}^{(l)}$ to $s$ leaves $Z$ through $Z\cap\Omega_k$ and crosses the rim of $Z$, at a
cost of at least $2\check R_k$; when $m(s)\le k-1$ it crosses in addition the rim of $s$ itself,
at a cost of $2\check R_k$ more, by the rim separation of the inheritance theorem. At frame sites
the age increment follows from \eqref{4.22:eq-law-site-a} with $m'=k-1$.

$(\alpha)$: the bound is the separation of far arguments, through the clause of that statement
giving $d_{\mathrm{sc}}(s,\Lambda^{\sim})\ge5\check R_k$, the supports of the arguments lying in
$\Lambda^{\sim}$; for the argument that is a field datum, obtained from the current
$\mathsf{J}_3$, this is the field-datum clause of the subfamily lemma.

$(\beta)$: the layer found in the proof of (i) between $Z_h$ and $(\Omega''_{h+1})^{\sim2}$,
diminished by the $2M_1+3$ blocks of $\mathsf{S}$, which amount to less than one $M$-cube, gives
the first member $\check R_k-1$; at frame sites the age increment follows from
\eqref{4.22:eq-law-site-a} with $m'=h$. This proves (ii).
\end{proof}

\begin{lemma}[Proof of the law-site lemma: the far debits]\label{4.22:prf-law-site-debits}
Let $k$ be a step at which the hypotheses of the law-site lemma (Lemma~\ref{4.22:lem-law-site})
hold and at which its clauses (i) and (ii) hold. Throughout, $\varsigma$ denotes the zone width
of the weighted contour distance, so that
$d^{\wedge}\ge\varsigma\,d_{\mathrm{sc}}$. Assume in addition the following statements.
\begin{itemize}
\item Clauses (i) and (iv) of the axialising gauge map lemma
(Lemma~\ref{4.22:lem-axialising-map}); $\delta$ denotes the decay rate of its clause (iv), and
$\rho^{(l)}$ the stage weights of that clause, indexed by the stage $l$. Together
with these, the bound of the local layer norm of a layer $\mathcal{K}$ at a site $s$ by the layer
norms $N_y$ (described below) at the points $y$ of the top block $B^{\mathrm{top}}(s)$ of the site:
\begin{equation*}
\Theta_0^{\mathrm{loc}}(\mathcal{K})(s)\le\sup_{y\in B^{\mathrm{top}}(s)}N_y(\mathcal{K}) .
\end{equation*}
\item The subfamily lemma, taken by statement. For a subfamily $\sigma_0$ of the family of
cells on which the expansion is run, we write $\sigma=(\mathsf{s}(\Delta))_{\Delta}$ for the
family of interpolation parameters $\mathsf{s}(\Delta)$ of the cells $\Delta$ of $\sigma_0$. For
every such subfamily $\sigma_0$, the
localisation statements for $\mathbf{G}''(\sigma)$ and the layer bounds on which the presentation
rests hold with the same constants; the subfamily enters their proofs only through the single
count of the cluster expansion of \cite{B-CMP116}, and the decaying bound below, not only the
supremum bound, is independent of $\sigma_0$. In it, $N_y$ is the layer norm at the point $y$;
$\mathcal{Y}(\mathsf{s};\varphi)$ is the profile built from a datum $\varphi$ at the interpolation
parameters
$\mathsf{s}(\Delta)$ of the cells $\Delta$; $\tilde G_0(\mathsf{s})$ is the propagator through
which a field datum is built from a datum $\mathsf{J}$; $\ell$ is the unit length of the level;
and $B_{\mathrm{pr}}$, $K_\sharp$ are constants of that lemma. Its profile clause: for all
$|\mathsf{s}(\Delta)|\le e^{\kappa_1}$ and every $y$,
\begin{equation*}
N_y(\mathcal{Y}(\mathsf{s};\varphi))\le B_{\mathrm{pr}}\sup|\varphi|\,
e^{-\delta_P d^{\wedge}(y,\operatorname{supp}\varphi)} .
\end{equation*}
Its field-datum clause: if the datum is $\mathfrak{j}=\tilde G_0(\mathsf{s})\mathsf{J}$, then
\begin{equation*}
N_y(\mathcal{Y})\le B_{\mathrm{pr}}K_\sharp\sup\ell^3|\mathsf{J}|\,
e^{-\delta_P d^{\wedge}(y,\operatorname{supp}\mathsf{J})} .
\end{equation*}
Its source condition: on the families run by the presenter maps $s'=2,3,4$, on which
\eqref{4.22:eq-law-site-b} is stated, no cell of the subfamily meets $\Lambda$.
\item The localisation of the current beyond the cut (Lemma~\ref{4.22:lem-current-localisation}),
that is, the bound \eqref{4.22:eq-current-localisation-a}, whose derivation uses $\sigma$ only
through the bound of a layer by its arguments, the kernel rate of $\mathbf{G}''(\sigma)$ and the
condition that the subfamily of cells does not meet $\Lambda$; $C_L$, $\delta_L$ and the source
set $\mathsf{S}$ are those of that lemma.
\item The level-width condition: with $\delta_P$ the profile rate of the subfamily lemma,
\begin{equation*}
\tfrac14\delta\varsigma\ge 4\ln L+1,\qquad \delta_P\varsigma\ge 4\ge 1\ge\delta_\natural .
\end{equation*}
\item The coupling ceilings \eqref{4.22:eq-coupling-ceilings-g}, through their consequence
$e^{-5\delta_\natural R_{\min}}\le e^{-2}$, where $R_{\min}$ is the least range entering the
frame-site distances of clause (ii) of Lemma~\ref{4.22:lem-law-site}.
\item The ordering of constants, in which
$B_H^S\ge B_H$ and $B_H^S\ge B_{\mathrm{pr}}\max\{C_1/K_R,\,1+K_\sharp\}$; here $B_H$ is the
constant of the bound of a layer by its arguments and $K_R$ the constant of the homogeneity bound,
both stated below, and $B_H^S$ is the constant of \eqref{4.22:eq-law-site-b}.
\item The properties of the presenter maps $h^{\mathfrak{b}}=h^{\mathfrak{b}}_q$ (the index $q$
is suppressed when immaterial): $h^{\mathfrak{b}}=1$ on
$\{\operatorname{dist}(\cdot,\Lambda^{\natural})\ge\tfrac12\}$, $\Lambda^{\natural}$ denoting the
region of the presentation that is fixed together with the presenter maps; for
$\mathfrak{b}\ge\mathfrak{b}_2$, $h^{\mathfrak{b}}=1$ on $D_2$ and on the strip; for $s'=1,2$,
$\mathfrak{w}_{s'}\equiv1$ on $\Lambda^{\sim}$; the recursion
$\mathfrak{Q}^{\mathfrak{b}_{s'}}=\mathfrak{Q}^{\mathfrak{b}_{s'-1}}\mathfrak{w}_{s'}$, started at
$\mathfrak{Q}^{\mathfrak{b}_0}=\mathbb{1}$; and
$|\log\mathfrak{w}_{s'}|\le K_u\Theta_0^{\mathrm{loc}}(\mathbb{H}_{s'})$.
\item The bound of a layer by its arguments, which is \cite[(1.45)]{B-CMP122a} in the present
notation: for every layer $\mathbb{H}$ of the presentation and every block $\square$, the sum
running over the arguments $\psi$ of the layer,
\begin{equation*}
\Theta_0^{\mathrm{loc}}(\mathbb{H})(\square)\le B_H\sum_{\psi}\sup|\psi|\,
e^{-\delta_1 d_{\mathrm{sc}}(\square,\operatorname{supp}\psi)} ,
\end{equation*}
with the constant $B_H$ and the rate $\delta_1$ of that bound. The amplitude
$\varepsilon_{\mathfrak{Q}}=\varepsilon_{\mathfrak{Q}}(k)$ of the presentation is defined so that
$K_uB_H\sum_\psi\sup|\psi|\le\varepsilon_{\mathfrak{Q}}(k)$, the sum running over all arguments of
the layers $\mathbb{H}_{s'}$.
\item The homogeneity bound at $\mathsf{t}=0$: for the variable $\mathsf{B}$ of
\cite[(1.19)]{B-CMP116}, $K_R\,g_k\sup|\mathsf{B}|\le\alpha_*$, where $K_R$ is the constant of
that bound.
\item The timing of arrests: an attempt arrested at step $k$ is a creation, and the base level
$\tilde h$ of every later attempt on the same component, made at a step $\tilde k$, satisfies
$k\le\tilde h$.
\item The stage-weight bound of the inheritance theorem, $\sum_l\rho^{(l)}\le3\beta$, in which
$\beta$ is the constant fixed in that theorem.
\item The law-site rule (Lemma~\ref{4.22:lem-law-site-rule}) and clause (v) of the frame lemma.
\end{itemize}
Then clauses (iii), (iii)$^S$, (iv) and (v) of Lemma~\ref{4.22:lem-law-site} hold at step $k$.
\end{lemma}

\begin{proof}
\emph{Clause (iii): the presenter maps at} $(\mathrm{L})$. In case $(\alpha)$ the value
$h^{\mathfrak{b}}_q(s)=1$ is the first property of the presenter maps. In case $(\beta)$,
clause (i) gives $s\in Z_h\subset D_2$ and $\mathfrak{b}\in\{\mathfrak{b}_2,\mathfrak{b}_3,
\mathfrak{b}_4\}$, so $\mathfrak{b}\ge\mathfrak{b}_2$ and the second property gives
$h^{\mathfrak{b}}_q(s)=1$.

\emph{Case $(\alpha)$.} For $g,g'\in G^{\mathbb{C}}$ one has $\|gg'\|\le\|g\|\,\|g'\|$ and
$\|(gg')^{-1}\|\le\|g^{-1}\|\,\|g'^{-1}\|$, hence
$\operatorname{pol}(gg')\le\operatorname{pol}(g)+\operatorname{pol}(g')$; and, since
$\|\exp(\pm\log\mathfrak{w}_{s'}(s))\|\le e^{|\log\mathfrak{w}_{s'}(s)|}$, one has
$\operatorname{pol}(\mathfrak{w}_{s'}(s))\le|\log\mathfrak{w}_{s'}(s)|$. By the recursion
$\mathfrak{Q}^{\mathfrak{b}_{s'}}=\mathfrak{Q}^{\mathfrak{b}_{s'-1}}\mathfrak{w}_{s'}$,
$\mathfrak{Q}^{\mathfrak{b}}(s)$ is the ordered product of the $\mathfrak{w}_{s'}(s)$. With the
bound on $|\log\mathfrak{w}_{s'}|$ this yields
\begin{equation}\label{4.22:eq-law-site-debits-1}
\operatorname{pol}(\mathfrak{Q}^{\mathfrak{b}}(s))
\le\sum_{s'}|\log\mathfrak{w}_{s'}(s)|
\le K_u\sum_{s'}\Theta_0^{\mathrm{loc}}(\mathbb{H}_{s'})(\tilde s),
\end{equation}
where $\tilde s$ is the block at which the local layer norm is read for the site $s$; it lies at
scaled distance at most one from $s$. Apply the bound of a layer by its arguments
(\cite[(1.45)]{B-CMP122a} in the present notation) at
$\square=\tilde s$ to each $\mathbb{H}_{s'}$. Replacing $\tilde s$ by $s$ costs at most one unit
of $d_{\mathrm{sc}}$, since
$d_{\mathrm{sc}}(\tilde s,\operatorname{supp}\psi)\ge
d_{\mathrm{sc}}(s,\operatorname{supp}\psi)-1$. Moreover
$K_uB_H\sum\sup|\psi|\le\varepsilon_{\mathfrak{Q}}$ by the definition of
$\varepsilon_{\mathfrak{Q}}$. Hence the right side of \eqref{4.22:eq-law-site-debits-1} is at
most $\varepsilon_{\mathfrak{Q}}$ times the largest of the factors
$e^{-\delta_1(d_{\mathrm{sc}}(s,\operatorname{supp}\psi)-1)}$, the supports of the arguments
lying in $\Lambda^{\sim}$, so that
$d_{\mathrm{sc}}(s,\operatorname{supp}\psi)\ge d_{\mathrm{sc}}(s,\Lambda^{\sim})$.

\emph{Case $(\beta)$.} For $s'=1,2$ the factors $\mathfrak{w}_{s'}$ are identically one on
$\Lambda^{\sim}$ and contribute nothing. For $s'=3,4$,
\eqref{4.22:eq-current-localisation-a} gives
\begin{equation*}
K_u\Theta_0^{\mathrm{loc}}(\mathbb{H}_{s'})(\tilde s)
\le C_L\varepsilon_{\mathfrak{Q}}(k)\,e^{-\delta_L(d_{\mathrm{sc}}(s,\mathsf{S})-1)} .
\end{equation*}

\emph{Cases $(\mathrm{T})$, $(f_1)$, $(f_2)$.} The bound is that of clause (iv) of
Lemma~\ref{4.22:lem-axialising-map}. Since $\tfrac14\delta\varsigma\ge4\ln L+1\ge1$ by the
level-width condition and $d_{\mathrm{sc}}\ge0$, its exponential factor satisfies
\begin{equation*}
e^{-\frac14\delta d^{\wedge}}\le e^{-(\frac14\delta\varsigma)d_{\mathrm{sc}}}\le e^{-d_{\mathrm{sc}}},
\end{equation*}
and its stage weights satisfy $\rho^{(l)}\le1$.

\emph{The exponent.} In each of the five cases, by the distances of clause (ii), the exponent is
at most $-\delta_\natural(\check R_k-2)$.

\emph{Frame sites.} At a frame site $s\in\Sigma_W$ the exponent is moreover decreased by
$5\delta_\natural R_{\min}$ for each level crossed. The age increments of clause (ii), which rest
on \eqref{4.22:eq-law-site-a}, are $5R_{\min}(k-1-k_W)_+$ in the cases $(\mathrm{T})$, $(f_2)$,
$(\alpha)$ and $5R_{\min}(h-k_W)_+$ in the cases $(f_1)$, $(\beta)$. Since $k-1\ge h=k-N_k$,
both counts of levels are at least
\begin{equation*}
\mathsf{a}_k(s)=(k-N_k-k_W)_+ .
\end{equation*}
By \eqref{4.22:eq-coupling-ceilings-g} each level crossed contributes a factor
$e^{-5\delta_\natural R_{\min}}\le e^{-2}$, which gives the factor $e^{-2\mathsf{a}_k(s)}$. At
pointwise sites the factor $e^{-2\mathsf{a}_k(s)}$ of clause (iii) is read as $1$.

\emph{Count.} There is one reading at $(\mathrm{T})$, one at each of the $N_k$ stages and at most
two at presentations, that is, at most $N_k+3$ readings. Each is bounded as in clause (iii), and
summing gives the displayed bound on $\eta^T_k(s)+\eta^P_k(s)+\eta^L_k(s)$. All components of
the class are processed simultaneously by a single family. The sets $\mathcal{C}^{(l)}$ and
$\Lambda^{\sim}$ are unions over these components, and the sums over the components are contained
in the block constant $\mathsf{c}_{\square}$ of the stage bounds and in $B_H$.

\emph{Clause (iii)$^S$.} We prove the bounds \eqref{4.22:eq-law-site-b} on the families on which
the sourced clause runs.

At $(\mathrm{T})$, the site's Landau layer $\mathcal{K}$ is the profile
$\mathcal{Y}(\mathsf{s};\varphi)$
of a datum $\varphi$ with $\operatorname{supp}\varphi\subset\Omega_k$. By clause (i) of
Lemma~\ref{4.22:lem-axialising-map}, by $\Theta_0^{\mathrm{loc}}(\mathcal{K})(s)\le
\sup_{y\in B^{\mathrm{top}}(s)}N_y(\mathcal{K})$, and by the profile clause of the subfamily
lemma,
\begin{align*}
|\log\mathsf{u}(\mathsf{s})(s)|
&\le K_u\Theta_0^{\mathrm{loc}}(\mathcal{K})(s)
\le K_u\sup_{y\in B^{\mathrm{top}}(s)}N_y(\mathcal{K})\\
&\le K_uB_{\mathrm{pr}}\sup|\varphi|\,e^{-\delta_P d^{\wedge}(B^{\mathrm{top}}(s),\Omega_k)} .
\end{align*}
By \cite[(1.19)]{B-CMP116} and the homogeneity bound at $\mathsf{t}=0$,
$\sup|\varphi|\le C_1g_k\sup|\mathsf{B}|\le(C_1/K_R)\alpha_*$. The block
$B^{\mathrm{top}}(s)$ lies in
the $M$-cube of $s$, which costs one unit of $d_{\mathrm{sc}}$, and $d^{\wedge}\ge\varsigma
d_{\mathrm{sc}}$. Hence
\begin{equation*}
|\log\mathsf{u}(\mathsf{s})(s)|\le K_uB_{\mathrm{pr}}(C_1/K_R)\,\alpha_*\,
e^{-\delta_P\varsigma(d_{\mathrm{sc}}(s,\Omega_k)-1)} .
\end{equation*}

At $(\mathrm{P})$ the bound is clause (iv)$(\mathrm{P})$ of Lemma~\ref{4.22:lem-axialising-map}
word for word: no family changes there.

At $(\mathrm{L})(\alpha)$, $h^{\mathfrak{b}}_q(s)=1$ by clause (iii). The recursion
$\mathfrak{Q}^{\mathfrak{b}_{s'}}=\mathfrak{Q}^{\mathfrak{b}_{s'-1}}\mathfrak{w}_{s'}$ and
$|\log\mathfrak{w}_{s'}|\le K_u\Theta_0^{\mathrm{loc}}(\mathbb{H}_{s'})$ apply. The
Baker--Campbell--Hausdorff formula at sizes at most $10^{-2}$ bounds the norm of the logarithm of
the product by $1.01$ times the sum of the norms of the logarithms of its factors. The profile
clause then applies to the
arguments supported in $\Lambda^{\sim}$, and the field-datum clause to the field datum
$\mathfrak{j}_3=\tilde G_0(\mathsf{s})\mathsf{J}_3$ among the arguments of the layers, at the
points $y\in B^{\mathrm{top}}(\tilde s)$ through which
$\Theta_0^{\mathrm{loc}}(\mathbb{H}_{s'})(\tilde s)$ is bounded. The supports entering these
weights lie in $\Lambda^{\sim}$, and the block $\tilde s$ and its top block
$B^{\mathrm{top}}(\tilde s)$ are read in the $M$-cube of $s$, so that passing from $s$ to any point
$y$ of $B^{\mathrm{top}}(\tilde s)$ costs at most one unit of $d_{\mathrm{sc}}$, as at
$(\mathrm{T})$. Hence for every such $y$ and every such support
\begin{equation*}
d^{\wedge}(y,\operatorname{supp}\psi)\ge\varsigma\,(d_{\mathrm{sc}}(s,\Lambda^{\sim})-1) .
\end{equation*}
This gives the bound of \eqref{4.22:eq-law-site-b} at
$(\mathrm{L})(\alpha)$, with the bracket
$\sum_{s'}\sup|\psi_{s'}|+K_\sharp\sup\ell^3|\mathsf{J}_3|$, where $\psi_{s'}$ runs over the
arguments of $\mathbb{H}_{s'}$.

At $(\mathrm{L})(\beta)$, the proof of Lemma~\ref{4.22:lem-current-localisation} uses $\sigma$
at three places:
\begin{itemize}
\item the bound of a layer by its arguments (\cite[(1.45)]{B-CMP122a} in the present notation)
on the polydisc, which holds by the profile clause;
\item the kernel rate of $\mathbf{G}''(\sigma)$, which is independent of the subfamily;
\item the step that uses that the subfamily does not meet $\Lambda$, which holds there by the
source condition.
\end{itemize}
All three hold on the families of \eqref{4.22:eq-law-site-b}. These families differ from those of
\cite{B-CMP122b} only by cells of the large-field region $X$ off $\Lambda^{\sim}$ at the first
operation, and no cell
meeting $\Lambda$ is ever cut. Hence \eqref{4.22:eq-current-localisation-a} holds unchanged, up
to the factor $B_H^S/B_H$ inside $\varepsilon_{\mathfrak{Q}}$.

\emph{Conversion.} By the ordering of constants,
$K_uB_{\mathrm{pr}}C_1/K_R\le K_uB_H^S$ and $K_uB_{\mathrm{pr}}(1+K_\sharp)\le K_uB_H^S$. The
bracket at $(\mathrm{L})(\alpha)$ satisfies
\begin{equation*}
\sum_{s'}\sup|\psi_{s'}|+K_\sharp\sup\ell^3|\mathsf{J}_3|
\le(1+K_\sharp)\max\Bigl\{\sum_{s'}\sup|\psi_{s'}|,\ \sup\ell^3|\mathsf{J}_3|\Bigr\}.
\end{equation*}
Each member of the maximum is at most the quantity bounded in the definition of
$\varepsilon_{\mathfrak{Q}}$, up to the factor $B_H^S/B_H$. This factor, like the losses $-1$ in
the exponents and the factor $1.01$, is absorbed in the constant $C_\natural$ of clause (iii),
which contains a factor $e^{2}$ and, in \eqref{4.22:eq-law-site-b}, the ratio $B_H^S/B_H$. The
distances are those of clause (ii). Since $d^{\wedge}\ge\varsigma d_{\mathrm{sc}}$ and, by the
level-width condition,
\begin{equation*}
\delta_P\varsigma\ge1\ge\delta_\natural ,
\end{equation*}
one has $e^{-\delta_P d^{\wedge}}\le e^{-\delta_\natural d_{\mathrm{sc}}}$, which brings each
reading to the majorant of clause (iii); this conversion is part of
\eqref{4.22:eq-law-site-b}.

\emph{Clause (iv).} Let $a=(A,j',b)$ be of kind (n) at step $k$, so $h\le j'\le k-1$ and $b$
lies in the processed class. If the chain completes, $a$ is integrated. If it does not, $a$ is
arrested at step $k$, and the arrest is a creation. The next attempt on its component, made at a
step $\tilde k$, then has base level $\tilde h\ge k$ by the timing of arrests. Thus
$j'\le k-1<\tilde h$, and $a$ is of kind $(f_1)$ from then on.

Before step $k$ no chain ran on the component of $a$ with $j'$ in range. Indeed, $\mathbf{R}$
acts on a component at the first step at which the first-time condition among the hypotheses of
Lemma~\ref{4.22:lem-law-site} is met. An
earlier arrested attempt at a step $\hat k<k$ would give $\hat k\le h\le j'$ by the timing of
arrests, whereas a chain at step $\hat k$ reads levels at most $\hat k-1<j'$ only.

Hence $a$ is of kind (n) in at most one chain of its life. The bound on $\nu(s)$ follows from
clause (iv) of Lemma~\ref{4.22:lem-axialising-map} at each
stage and from $\sum_l\rho^{(l)}\le3\beta$. The coordinates of the arrested variables are
created by the arrest itself, so they do not exist for the chain at step $k$; at every later
attempt their law sites are law sites of arrested variables, which by the last sentence of
clause (i) are never of kind (n).

\emph{Clause (v).} By the law-site rule (Lemma~\ref{4.22:lem-law-site-rule}), each argument is
rotated at its own law site. By clause (i), at step $k_{i+1}$ the consumers' law sites
$\Sigma_{W_i}$ are of case $(\beta)$. The account of
$\Sigma_{W_i}$ is closed there by clause (v) of the frame lemma. By the second sentence of the
law-site rule, the new root is the
law site of the new terms only, and it is first read at a later step. Joining components, as in
\cite{B-CMP122b}, changes no law site.
\end{proof}

In \cite[(1.54)--(1.55)]{B-CMP122b} the operator $G''_k$ is applied to the current $J_{1,2}$ of
the configuration $U_{1,2}$, and Ba{\l}aban explains the localisation of this composition as follows:
\begin{quotation}
We have written explicitly the form of the operator dependence of the function $H''_k$ on the
fields $J_{1,2}$, $Q(\eta A_0)$. The second field is obviously localized in the domain
$Z''_k\cap\Omega''^{\sim2}_{h+1}$ \dots{} The field $J_{1,2}$ is not localized, but in the
composition $G''_kJ_{1,2}$ we may localize it,
because of the properties of the kernel of the operator $G''_k$. The property important here is
the equality $G''_kQ''^{*} = 0$, see the definition (3.148) [13] and the identities
(3.149)--(3.153) [13]. The field $J_{1,2}$ restricted to $\Omega''^{\sim}_{h+1}$ is equal to
$J_1$, and it satisfies the criticality condition $\langle\delta A,J_1\rangle = 0$ for $\delta A$
such, that $Q_{\mathbf{B}_1}\delta A = 0$. This condition for $\delta A$, and the condition for
$G''_k$ coincide on the domain $\Omega''^{\sim}_{h+1}$, except the boundary layer of the width
$2M_1$ \dots
\end{quotation}
The reference [13] there is \cite{B-CMP99b}. In \cite[(1.55)]{B-CMP122b} this localisation is
carried out on $\Omega''^{\sim}_{h+1}$, by means of the criticality condition of $J_1$ just
quoted. The result is
\begin{equation*}
G''_kJ_{1,2}=G''_k(1-\theta)Q^{*}\theta_1H^{*}J_1+G''_k\theta J_{1,2},
\end{equation*}
with $\theta_1$ and $H^{*}$ as in that paper.
The term $G''_k\theta J_{1,2}$, with $\theta=1$ on $(\Omega''^{\sim}_{h+1})^c$ and on the boundary
layer of width $2M_1$ at the cut, is kept whole; inside the set $Z$ of that paper nothing more is
needed. The following lemma localises the composition on the other side of the cut, on the domain
$D_2$.

\begin{lemma}[Localisation of the current beyond the cut]\label{4.22:lem-current-localisation}
We use the notation of \cite[\S1]{B-CMP122b}, in particular of (1.50)--(1.55) there, for the
following objects:
\begin{itemize}
\item the level $k$ and the stage index $h$;
\item the domains $\Omega''^{\sim}_{h+1}$ and $\Omega''^{\sim2}_{h+1}$, and the cut
$\partial\Omega''^{\sim}_{h+1}$;
\item the region $\Lambda$ and the sets $Z_h$, $Z''_k$, $Z''_{h+1}$;
\item the determining set $\mathbf{B}''_k$, its averaging operator $Q''$ with adjoint $Q''^{*}$,
and the operator $G''_k$.
\end{itemize}
Write $\mathbf{G}''(\sigma)$ for $G''_k$ carrying the decoupling parameters
$\sigma=\{\sigma_\Delta\}$ of its random-walk expansion. As in \cite[\S1]{B-CMP122b}, these
parameters are introduced only on the cubes of the set $\sigma_0$ of cubes $\Delta$ disjoint with
$\Lambda$, so that $\mathbf{G}''(\mathbb{1})=G''_k$.

Set $\mathbf{B}_2:=\mathbf{B}''_k\cup\mathbf{B}_h((\Omega''^{\sim}_{h+1})^c)$ as in
\cite[(1.52)]{B-CMP122b}, where $\mathbf{B}_h(\cdot)$ denotes the determining set of order $h$ on
the indicated domain, and $Q:=Q_{\mathbf{B}_2}$. Let $U_2=U_{\mathbf{B}_2}(\cdot)$ be the
minimiser of the problem \cite[(1.53)]{B-CMP122b}. Write
$\mathcal{J}(U)$ for the current of a configuration
$U$, and set $J_2:=\mathcal{J}(U_2)$ and $\lambda_2:=(QQ^{*})^{-1}QJ_2$. Let
\begin{equation*}
C_Q:=1+\|Q^{*}\|\,\|(QQ^{*})^{-1}Q\|,
\end{equation*}
the norms being operator norms between supremum norms; by \cite{B-CMP98}, $C_Q$ depends on $d$ and
$L$ only.

Consider the presentation of index $\mathfrak{b}_2$ with source $\zeta\in\mathfrak{S}_{\mathrm{src}}$.
Let $D_2\subset Z''_{h+1}\subset\Lambda$ be the domain of this presentation lying beyond the cut
$\partial\Omega''^{\sim}_{h+1}$.
Write $\hat R^{\mathfrak{b}_2}$ for its configuration, $\hat{\mathcal{J}}^{\mathfrak{b}_2}$ for its
current and $\hat{\mathcal{U}}^{\mathfrak{b}_2}$ for its background. Among its layers we consider
$\mathbb{H}_3$ and $\mathbb{H}_4$, with local layer norm $\Theta_0^{\mathrm{loc}}$, evaluated at
the point $\tilde z$ which the presentation associates with a point $z$, and constants $K_u$,
$B_H$, $\varepsilon_{\mathfrak{Q}}(k)$. Let $d_{\mathrm{sc}}$ be the scaled distance in $M$-cubes
of the local scale. Assume the following.
\begin{enumerate}
\item[(a)] \emph{Euler--Lagrange identity of a $\mathbf{B}$-minimiser.} If $U_{\mathbf{B}}$ is
the minimiser of the constrained problem with determining set $\mathbf{B}$ (the constrained minimum
of \cite[(2.12)]{B-CMP119}), then $\langle\delta A,\mathcal{J}(U_{\mathbf{B}})\rangle=0$ for every
$\delta A$ with $Q_{\mathbf{B}}\delta A=0$.
\item[(b)] \emph{The deep-current identity of the presentation.} On every stencil of the current
meeting $D_2\cup B_S$, where $B_S$ is the set attached to the presentation on which this identity
is stated, $\hat{\mathcal{J}}^{\mathfrak{b}_2}=\mathcal{J}(\hat R^{\mathfrak{b}_2})$.
\item[(c)] \emph{The configuration of the presentation beyond the cut.}
$\hat R^{\mathfrak{b}_2}\restriction D_2=U_2$, and $U_2$ is real, of order $h$ and
independent of $\mathbb{U}$.
\item[(d)] \emph{The layer bound for $\mathbb{H}_s$} (cf.\ \cite[(1.45)]{B-CMP122a}). It is
stated for an argument localised in
one $M$-cube $y$. It bounds $\Theta_0^{\mathrm{loc}}(\mathbb{H}_s)(\tilde z)$ by $B_H$ times the
supremum of the argument times $e^{-\delta_1 d_{\mathrm{sc}}(z,y)}$, where $\delta_1$ is its rate
per unit of $d_{\mathrm{sc}}$.
\item[(d$'$)] \emph{The budget of the arguments.} The arguments of $\mathbb{H}_3$ and
$\mathbb{H}_4$ satisfy $K_uB_H\sum\sup|\cdot|\le\varepsilon_{\mathfrak{Q}}(k)$.
\item[(e)] The statements of \cite[Theorem~3.13, Theorem~3.14, (3.148)--(3.154)]{B-CMP99b} for
$G''_k$ at the backgrounds of its class, namely:
\begin{itemize}
\item the identity $G''_kQ''^{*}=0$;
\item the kernel decay of $G''_k$, with rate $\delta_G$ per $M$-cube, obtained from the walks of
the propagator $\mathfrak{G}$ of \cite[Theorem~3.13]{B-CMP99b};
\item the walk representation \cite[Theorem~3.14, (3.154)]{B-CMP99b} of the dependence of
$\mathbf{G}''(\sigma)$ on $\sigma$.
\end{itemize}
\end{enumerate}
Let $D_2^{\circ}$ be the union of the $\mathbf{B}''_k$-blocks of $D_2$ at block-distance at least
$2$ from the boundary layer of width $2M_1$ at the cut $\partial\Omega''^{\sim}_{h+1}$. Let
$\mathsf{S}$ be the $(2M_1+3\text{ blocks})$-neighbourhood of $\Omega''^{\sim2}_{h+1}$ in
$\Lambda$. Let $\zeta\in\mathfrak{S}_{\mathrm{src}}$ and $|\sigma|\le e^{\kappa_1}$. Then there is a
field $r$, supported in the blocks adjacent to $\partial D_2^{\circ}$ and satisfying
$\sup|r|\le C_Q\sup|J_2|$, such that
\begin{equation}\label{4.22:eq-current-localisation-a}
\begin{split}
\mathbf{G}''(\sigma)\hat{\mathcal{J}}^{\mathfrak{b}_2}
&=\mathbf{G}''(\sigma)\bigl[\mathbb{1}_{(D_2^{\circ})^c}\hat{\mathcal{J}}^{\mathfrak{b}_2}+r\bigr]\\
&\quad+\bigl[\mathbf{G}''(\sigma)-\mathbf{G}''(\mathbb{1})\bigr]
Q''^{*}(\lambda_2\restriction D_2^{\circ}).
\end{split}
\end{equation}
Moreover, for $z\in Z_h$ and $s\in\{3,4\}$,
\begin{equation}\label{4.22:eq-current-localisation-1}
K_u\Theta_0^{\mathrm{loc}}(\mathbb{H}_s)(\tilde z)
\le C_L\,\varepsilon_{\mathfrak{Q}}(k)\,e^{-\delta_L(d_{\mathrm{sc}}(z,\mathsf{S})-1)},
\qquad \delta_L:=\tfrac12\min\{\delta_1,\delta_G\},
\end{equation}
where $C_L=C(d,L,M)\ge C_Q$.
\end{lemma}

\begin{proof}
\emph{The current on $D_2$.} By hypothesis (b), on every stencil meeting $D_2\cup B_S$ we have
$\hat{\mathcal{J}}^{\mathfrak{b}_2}=\mathcal{J}(\hat R^{\mathfrak{b}_2})$. By hypothesis (c),
$\hat R^{\mathfrak{b}_2}\restriction D_2=U_2$, and $U_2$ is real, of order $h$ and independent of
$\mathbb{U}$. Hence $J_2=\mathcal{J}(U_2)$ is real and independent of $\zeta$, and on every bond
whose stencil is contained in $D_2$, in particular on the blocks of $D_2^{\circ}$, we have
$\hat{\mathcal{J}}^{\mathfrak{b}_2}=J_2$. At complex $\zeta$ the
identity proved below is therefore needed only on the side of $D_2$. The sources on the side of
$\Omega''^{\sim}_{h+1}$, localised or not, lie in $\mathsf{S}$.

\emph{Criticality of $J_2$.} The configuration $U_2=U_{\mathbf{B}_2}(\cdot)$ minimises its own
problem \cite[(1.53)]{B-CMP122b}. Ba{\l}aban states the corresponding property for $J_1$. Here we
apply hypothesis (a), the Euler--Lagrange identity of a $\mathbf{B}$-minimiser, with
$\mathbf{B}=\mathbf{B}_2$. It gives $\langle\delta A,J_2\rangle=0$ whenever
$Q_{\mathbf{B}_2}\delta A=0$, that is,
\begin{equation*}
J_2\in(\ker Q_{\mathbf{B}_2})^{\perp}=\operatorname{ran}Q^{*}_{\mathbf{B}_2}.
\end{equation*}
Write $J_2=Q^{*}_{\mathbf{B}_2}\lambda$. Applying $Q=Q_{\mathbf{B}_2}$ gives $QJ_2=QQ^{*}\lambda$,
hence $\lambda=(QQ^{*})^{-1}QJ_2=\lambda_2$. So
\begin{equation*}
J_2=Q^{*}_{\mathbf{B}_2}\lambda_2 .
\end{equation*}
The bounds on $Q^{*}$ and $(QQ^{*})^{-1}Q$ entering $C_Q$ are those of the averaging operators of
\cite{B-CMP98}, and $C_Q$ depends on $d$ and $L$ only.

\emph{Comparison of the determining sets.} By \cite[(1.52)]{B-CMP122b},
$\mathbf{B}_2=\mathbf{B}''_k\cup\mathbf{B}_h((\Omega''^{\sim}_{h+1})^c)$. On $D_2$ this set has the
same blocks as $\mathbf{B}''_k$, except within the graded layer, of width at most $2M_1$, at the
cut $\partial\Omega''^{\sim}_{h+1}$; there the blocks of $\mathbf{B}_h(\cdot)$ are finer. This is
read off from \cite[(2.13), (2.14)]{B-CMP119}; the letters $\Omega$, $\Omega_n$, $\Omega_j$,
$\xi$, $\mathbf{B}$ and $\mathbf{B}_j(\Omega)$ in the next sentences are those of that paper.
By \cite[(2.13)]{B-CMP119}, $\Omega_n$ is a union of
$L^n\xi M_1$-cubes and $\Omega^{\sim-2}\subset\Omega_j$. By \cite[(2.14)]{B-CMP119},
$\partial\Omega\subset\Omega_j\setminus\Omega_{j+1}$, and the same formula bounds a distance in
this configuration from below by $2M_1$. The
display of \cite[(2.14)]{B-CMP119} reads
\begin{equation*}
\mathbf{B}\cup\mathbf{B}_j(\Omega)
=\bigl(\mathbf{B}\cap\Omega^{\sim-2}\bigr)\cup
\bigl(\mathbf{B}_j(\Omega)\cap(\Omega\setminus\Omega^{\sim-2})\bigr).
\end{equation*}
The symmetric statement for $\mathbf{B}_1$ is made in \cite[\S1]{B-CMP122b}. At the cut,
$\mathbf{B}''_k$ is of order $h$ by \cite[(1.13)]{B-CMP122a}.

For a bond $b$, the value $(Q^{*}\lambda)(b)$ reads $\lambda$ only on the block bonds touching the
block of $b$. On $D_2^{\circ}$, which lies at least two blocks away from that layer, these are
common blocks of $\mathbf{B}_2$ and $\mathbf{B}''_k$. Hence $Q''^{*}(\lambda_2\restriction
D_2^{\circ})=Q^{*}_{\mathbf{B}_2}(\lambda_2\restriction D_2^{\circ})$. Moreover this field agrees
with $\mathbb{1}_{D_2^{\circ}}Q^{*}_{\mathbf{B}_2}\lambda_2=\mathbb{1}_{D_2^{\circ}}J_2$ except on
the blocks adjacent to $\partial D_2^{\circ}$. Therefore
\begin{equation}\label{4.22:eq-current-localisation-2}
\mathbb{1}_{D_2^{\circ}}J_2=Q''^{*}(\lambda_2\restriction D_2^{\circ})+r,
\qquad r:=\mathbb{1}_{D_2^{\circ}}J_2-Q''^{*}(\lambda_2\restriction D_2^{\circ}),
\end{equation}
with $r$ supported in the blocks adjacent to $\partial D_2^{\circ}$. From the expression for $r$,
\begin{equation*}
\sup|r|\le\sup|J_2|+\|Q^{*}\|\sup|\lambda_2|
\le\bigl(1+\|Q^{*}\|\,\|(QQ^{*})^{-1}Q\|\bigr)\sup|J_2|=C_Q\sup|J_2| .
\end{equation*}

\emph{The display.} Write
$\hat{\mathcal{J}}^{\mathfrak{b}_2}=\mathbb{1}_{(D_2^{\circ})^c}\hat{\mathcal{J}}^{\mathfrak{b}_2}
+\mathbb{1}_{D_2^{\circ}}\hat{\mathcal{J}}^{\mathfrak{b}_2}$. By the first paragraph of this proof,
$\mathbb{1}_{D_2^{\circ}}\hat{\mathcal{J}}^{\mathfrak{b}_2}=\mathbb{1}_{D_2^{\circ}}J_2$, and we
insert \eqref{4.22:eq-current-localisation-2}. The identity $G''_kQ''^{*}=0$ of hypothesis (e) is
an algebraic identity of the full operator, that is, at $\sigma=\mathbb{1}$, and holds at any
background of its class. We apply it at the background $\hat{\mathcal{U}}^{\mathfrak{b}_2}$, which
equals $U_2$ on the blocks of $D_2^{\circ}$. This gives
$\mathbf{G}''(\mathbb{1})Q''^{*}(\lambda_2\restriction D_2^{\circ})=0$. Subtracting this vanishing
term from $\mathbf{G}''(\sigma)Q''^{*}(\lambda_2\restriction D_2^{\circ})$ yields
\eqref{4.22:eq-current-localisation-a}.

\emph{The bracket.} In $\mathbf{G}''(\sigma)-\mathbf{G}''(\mathbb{1})$ only the walks touching
$\sigma_0$ survive. By \cite[\S1]{B-CMP122b}, $\sigma_0$ is the set of cubes $\Delta$ disjoint with
$\Lambda$, while $D_2\subset Z''_{h+1}\subset\Lambda$. Hence a walk from $z\in Z_h$ that
contributes to the bracket must leave $\Lambda$. Such walks are of the type of
\cite[Theorem~3.14, (3.154)]{B-CMP99b}, which is used by statement through hypothesis (e). They
carry the factor
\begin{equation*}
e^{-\delta_G d_{\mathrm{sc}}(z,\Lambda^c)}\le e^{-\delta_G d_{\mathrm{sc}}(z,\mathsf{S})},
\end{equation*}
because, as seen from $Z_h$, $\Lambda^c$ lies beyond $\mathsf{S}$.

\emph{The layer bound.} Consider the first argument of $\mathbb{H}_3$, namely the right side of
\eqref{4.22:eq-current-localisation-a}. The sources $\mathbb{1}_{(D_2^{\circ})^c}
\hat{\mathcal{J}}^{\mathfrak{b}_2}$ and $r$ are placed as described in the first paragraph and in
\eqref{4.22:eq-current-localisation-2}. The kernel of $\mathbf{G}''(\sigma)$ decays at rate
$\delta_G$ by hypothesis (e), and the bracket obeys the preceding estimate. Together these make the
first argument of $\mathbb{H}_3$ a field $f_1$ with
\begin{equation*}
|f_1(y)|\le C\sup|\hat{\mathcal{J}}^{\mathfrak{b}_2}|\,e^{-\delta_G d_{\mathrm{sc}}(y,\mathsf{S})}
\end{equation*}
in scaled units, with $C=C(d,L,M)$. This field is not localised in $\mathsf{S}$; it only decays
away from $\mathsf{S}$. By contrast, the second argument of $\mathbb{H}_3$ and the argument $g_kB$
of $\mathbb{H}_4$, with $B$ the field of \cite[(1.50)]{B-CMP122b}, are supported in $\mathsf{S}$.
Indeed, by \cite[(1.50), (1.54)]{B-CMP122b} these arguments are supported in the sets
$\Lambda\cap\Omega''^{\sim2}_{h+1}$ and, for the tower terms, $Z''_k\cap\Omega''^{\sim2}_{h+1}$,
both of which lie in $\mathsf{S}$.

The layer bound of hypothesis (d) is stated for a localised argument. We cut $f_1$ into its
restrictions to the $M$-cubes $y$, apply the bound to each piece and sum. Since
$\min\{\delta_1,\delta_G\}=2\delta_L$, we have
\begin{equation*}
\delta_1d_{\mathrm{sc}}(z,y)+\delta_Gd_{\mathrm{sc}}(y,\mathsf{S})
\ge2\delta_L\bigl(d_{\mathrm{sc}}(z,y)+d_{\mathrm{sc}}(y,\mathsf{S})\bigr).
\end{equation*}
The triangle inequality gives $d_{\mathrm{sc}}(z,y)+d_{\mathrm{sc}}(y,\mathsf{S})\ge
d_{\mathrm{sc}}(z,\mathsf{S})$, and nonnegativity gives
$d_{\mathrm{sc}}(z,y)+d_{\mathrm{sc}}(y,\mathsf{S})\ge d_{\mathrm{sc}}(z,y)$. Hence
\begin{equation*}
d_{\mathrm{sc}}(z,y)+d_{\mathrm{sc}}(y,\mathsf{S})
\ge\max\{d_{\mathrm{sc}}(z,\mathsf{S}),d_{\mathrm{sc}}(z,y)\}.
\end{equation*}
Since twice a maximum dominates the sum of the two numbers,
\begin{equation*}
\delta_1d_{\mathrm{sc}}(z,y)+\delta_Gd_{\mathrm{sc}}(y,\mathsf{S})
\ge\delta_Ld_{\mathrm{sc}}(z,\mathsf{S})+\delta_Ld_{\mathrm{sc}}(z,y).
\end{equation*}
Put $C_\rho(\delta):=\sup_z\sum_y e^{-\delta d_{\mathrm{sc}}(z,y)}$, the sum running over
$M$-cubes. Then
\begin{equation*}
\sum_y e^{-\delta_1d_{\mathrm{sc}}(z,y)-\delta_Gd_{\mathrm{sc}}(y,\mathsf{S})}
\le e^{-\delta_Ld_{\mathrm{sc}}(z,\mathsf{S})}\sum_y e^{-\delta_Ld_{\mathrm{sc}}(z,y)}
\le C_\rho(\delta_L)\,e^{-\delta_Ld_{\mathrm{sc}}(z,\mathsf{S})}.
\end{equation*}
The sizes of the arguments enter through $K_uB_H\sum\sup|\cdot|\le\varepsilon_{\mathfrak{Q}}(k)$,
by hypothesis (d$'$). The field $r$ costs the factor $C_Q$, since $\sup|r|\le C_Q\sup|J_2|$.
Set $C_L:=C_QC_\rho(\delta_L)C$; we take $C$ and $C_\rho(\delta_L)$ at least $1$, so that
$C_L=C(d,L,M)\ge C_Q$. Multiplying the layer bound by $K_u$ and collecting these factors gives
\begin{equation*}
K_u\Theta_0^{\mathrm{loc}}(\mathbb{H}_s)(\tilde z)\le C_L\,\varepsilon_{\mathfrak{Q}}(k)\,
e^{-\delta_Ld_{\mathrm{sc}}(z,\mathsf{S})}.
\end{equation*}
Since the exponential on the right is at most
$e^{-\delta_L(d_{\mathrm{sc}}(z,\mathsf{S})-1)}$, this is
\eqref{4.22:eq-current-localisation-1}.
\end{proof}

\begin{remark}
The estimate of the bracket in the proof of Lemma~\ref{4.22:lem-current-localisation} is the only
place in the construction of the presented chain where it is used that no decoupling parameter is
placed on a cube meeting $\Lambda$. It is used only at the families of the presented chain. The
lemma giving the profiles $\vartheta^{3F}(s)$ of the led rotations uses neither the
identity $G''_kQ''^{*}=0$ nor Lemma~\ref{4.22:lem-current-localisation}; its profile for the class
of terms for which the side of $D_2$ has to be read is supplied by the profile lemma for that
class of terms.
\end{remark}

\begin{theorem}[The radius budget closes uniformly in all scales]\label{4.22:thm-budget-closes}
Assume the following:
\begin{itemize}
\item the hypotheses of the law-site lemma (Lemma~\ref{4.22:lem-law-site});
\item the upper and the lower flow bounds for the running inverse couplings $x_j=g_j^{-2}$,
at all indices $\le k$: for all $i\le j\le k$,
\begin{equation*}
x_j+\lambda_\sharp(j-i)-c_5\le x_i\le x_j+B_\sharp(j-i);
\end{equation*}
\item the running inverse couplings lie above the ceiling of the coupling window:
$x_j>\gamma_H^{-2}$ for all $j\le k$;
\item the ceilings (v)--(x) of the statement on coupling ceilings and the order of constants,
among them \eqref{4.22:eq-coupling-ceilings-c} and \eqref{4.22:eq-coupling-ceilings-h},
together with the ceilings \eqref{4.22:eq-coupling-ceilings-j} and
\eqref{4.22:eq-coupling-ceilings-k};
\item the radius $\alpha_{*,j}=\alpha_*(x_j)$ is a constant multiple of
$x_j^{-1/2}(\log x_j)^{q_*}$, with an exponent $q_*\ge0$, the power one half being that of
\cite[(2.6)]{B-CMP119} (the same power as in \cite{B-CMP122a}).
\end{itemize}
Put $y:=x_k-2c_5$ and $y_\flat:=\gamma_H^{-2}-2c_5$. Then for every $K$, every $k\le K$ and
every law site $s$ whose account is open at step $k$ in the sense of
Definition~\ref{4.22:def-radius-accounting}, the following hold, clauses \textup{(1)} and
\textup{(1$^{\mathrm{res}}$)} for each step $k'\le k$:
\begin{equation}\label{4.22:eq-budget-closes-a}
\begin{aligned}
&\textup{(1)} &&\eta_{k'}(s)\le 4\mu_{k'}\,\mathsf{e}(s)/\varepsilon_{\star}
  \;\bigl(+\,\nu(s)\bigr),\\
&\textup{(1$^{\mathrm{res}}$)} &&\eta_{k'}(s)+\eta^{\mathrm{inc}}_{k'}(s)
  \le 5\mu_{k'}\,\mathsf{e}(s)/\varepsilon_{\star}\;\bigl(+\,\nu(s)\bigr),\\
&\textup{(2)} &&\sum_{k'\le k}\eta_{k'}(s)
  \le\mathsf{e}(s)\Bigl\{4\bigl[y^{-2}+(\lambda_\sharp y)^{-1}\bigr]+\tfrac{1}{16}\Bigr\}\\
& &&\qquad\le\bigl(\tfrac18+\tfrac1{16}\bigr)\mathsf{e}(s)<\tfrac14\,\mathsf{e}(s),\\
&\textup{(2$^{\mathrm{res}}$)} &&\sum_{k'\le k}\bigl[\eta_{k'}(s)
  +\eta^{\mathrm{inc}}_{k'}(s)\bigr]\\
& &&\qquad\le\mathsf{e}(s)\Bigl\{5\bigl[y^{-2}+(\lambda_\sharp y)^{-1}\bigr]
  +\tfrac{1}{16}\Bigr\}\\
& &&\qquad\le\bigl(\tfrac5{32}+\tfrac2{32}\bigr)\mathsf{e}(s)=\tfrac7{32}\,\mathsf{e}(s)
  <\tfrac14\,\mathsf{e}(s),\\
&\textup{(3)} &&\tfrac34\,\mathsf{e}(s)\le\varepsilon_k(s)\le\mathsf{e}(s),
\end{aligned}
\end{equation}
where the summand $\nu(s)$ in parentheses in \textup{(1)} and \textup{(1$^{\mathrm{res}}$)}
is present only if $s$ is pointwise, and then at one step $k'$ only.

The right members of \textup{(2)} and \textup{(2$^{\mathrm{res}}$)} depend on none of $K$, $k$,
the age of $s$, the number of attempts on the component of $s$, the number $N_k$ of large-field
stages of the step, Ba{\l}aban's integer $N_0$ of \cite{B-CMP122a} (the number of top levels of an
attempt at which the large-field strata are thin), or the volume;
their quotients by $\mathsf{e}(s)$ depend on the step only through $y$ and are bounded using
$y\ge y_\flat$ alone, and, with $\nu(s)/\varepsilon_{\star}$ in place of $1/16$, they tend to
$0$ as $\gamma_H\to0$. No hypothesis on the estimates that produce the debits enters beyond those
listed: the margins and the reserve are chosen quantities. The conclusions hold for the radius
$\varepsilon_k(s)$ and equally for the radius $\varepsilon^{\mathfrak{A}}_k(s)$ of the plain
class, since the far-debit majorant of clause (iii) of Lemma~\ref{4.22:lem-law-site} and that of
its form on families, \eqref{4.22:eq-law-site-b}, are one and the same.
\end{theorem}

\begin{proof}
Fix $K$, $k\le K$ and a law site $s$ whose account is open at step $k$. We use the entries
\eqref{4.22:eq-radius-accounting-a} of Definition~\ref{4.22:def-radius-accounting}. The debit
$\eta_{k'}(s)$ charged against $s$ at step $k'$ consists of the far debits
$\eta^{T}_{k'}(s)+\eta^{P}_{k'}(s)+\eta^{L}_{k'}(s)$; of the margin
$\mu_{k'}\mathsf{e}(s)/\varepsilon_{\star}$ of the fluctuation integral; of one margin
$\mu_{k'}\mathsf{e}(s)/\varepsilon_{\star}$ for each of the at most two presentations of the
step; and, at a step at which $s$ is the site of a rim slot in a running chain, of the near
debit $\nu(s)$. The reserve is $\eta^{\mathrm{inc}}_{k'}(s)=\mu_{k'}\mathsf{e}(s)/
\varepsilon_{\star}$. The credit is $\mathsf{e}(s)=\varepsilon_{\star}$ at a pointwise site, and
$\mathsf{e}(s)=\varepsilon^{\mathrm{fr}}(W)=\varpi_{\mathrm{fr}}\alpha_{*,k_W}/(4C_{\mathrm{fr}})$
at a frame site $s\in\Sigma_W$. The margin is $\mu_{k'}=\mu(x_{k'})$ with
$\mu(x)=\varepsilon_{\star}(x-c_5)^{-2}$. The hypotheses of the law-site lemma are assumed at
indices $\le k$, so that lemma applies at every step $k'\le k$. By hypothesis
$x_{k'}>\gamma_H^{-2}$ for every $k'\le k$, so the ceilings \eqref{4.22:eq-coupling-ceilings-h},
\eqref{4.22:eq-coupling-ceilings-j} and \eqref{4.22:eq-coupling-ceilings-k}, which are stated for
$x\ge\gamma_H^{-2}$, apply at $x=x_{k'}$.

\emph{Clause \textup{(1)} at pointwise sites.} Let $s$ be pointwise. Clause (iii) of
Lemma~\ref{4.22:lem-law-site}, applied at step $k'$, bounds the far debits, apart from the near
debit, by $(N_{k'}+3)C_{\natural}\alpha^+_{k'}e^{-\delta_{\natural}(\check R_{k'}-2)}
e^{-2\mathsf{a}_{k'}(s)}$; at a pointwise site the factor $e^{-2\mathsf{a}_{k'}(s)}$ is read as
$1$, and we drop it in any case, as it is at most $1$. Here $N_{k'}=\check R_{k'}$ is the range
number of step $k'$, and we use the bounds
\begin{equation*}
(\log(x_{k'}-c_5))^r\le\check R_{k'}=N_{k'}\le L(\log x_{k'})^r
\end{equation*}
that accompany its definition (the upper one as in \cite{B-CMP122a});
since $\delta_{\natural}>0$, the lower one gives
$e^{-\delta_{\natural}(\check R_{k'}-2)}\le\exp(-\delta_{\natural}[(\log(x_{k'}-c_5))^r-2])$.
Let $\bar\alpha(x)$ be the function with $\alpha^+_j=\bar\alpha(x_j)$. Then the far debits of
step $k'$ are at most $\mathsf{F}(x_{k'})$, where $\mathsf{F}$ is the function of
\eqref{4.22:eq-coupling-ceilings-j}, namely
\begin{equation}\label{4.22:eq-budget-closes-1}
\mathsf{F}(x):=\bigl(L(\log x)^r+3\bigr)C_{\natural}\,\bar\alpha(x)\,
\exp\bigl(-\delta_{\natural}\bigl[(\log(x-c_5))^r-2\bigr]\bigr).
\end{equation}
The ceiling \eqref{4.22:eq-coupling-ceilings-j} states $\mathsf{F}(x)\le\mu(x)$ for all
$x\ge\gamma_H^{-2}$. Applied at $x=x_{k'}$, it bounds the far debits by
$\mu_{k'}=\mu_{k'}\mathsf{e}(s)/\varepsilon_{\star}$, since $\mathsf{e}(s)=\varepsilon_{\star}$.
Adding the three margins gives $\eta_{k'}(s)\le4\mu_{k'}\mathsf{e}(s)/\varepsilon_{\star}$, plus
$\nu(s)$ at a step at which $s$ is of kind (n). By clause (iv) of
Lemma~\ref{4.22:lem-law-site}, this happens in at most one chain of the life of $s$, and $\nu(s)$
is the total near debit over that chain.

\emph{Clause \textup{(1)} at frame sites.} Let $s\in\Sigma_W$, so that its depth is
$m(s):=k_W$ (Notation~\ref{4.22:not-law-site-depth}). The account of $s$ opens at step $k_W$
with the credit $\varepsilon^{\mathrm{fr}}(W)$, so no debit is charged against $s$ at steps
$k'<k_W$, and we may take $k_W\le k'\le k$. Clause (iii) of
Lemma~\ref{4.22:lem-law-site} now carries the factor $e^{-2\mathsf{a}_{k'}(s)}$ with
$\mathsf{a}_{k'}(s)=(k'-N_{k'}-k_W)_+$. With the same two bounds on the range number as above,
the far debits are at most $\mathsf{F}(x_{k'})e^{-2\mathsf{a}_{k'}(s)}$. Dividing by
$\mathsf{e}(s)=\varpi_{\mathrm{fr}}\alpha_{*,k_W}/(4C_{\mathrm{fr}})$, their quotient by the
credit is at most
\begin{equation*}
\mathsf{F}(x_{k'})\,e^{-2\mathsf{a}_{k'}(s)}\cdot
\frac{4C_{\mathrm{fr}}}{\varpi_{\mathrm{fr}}\alpha_{*,k_W}}.
\end{equation*}
To compare $\alpha_{*,k_W}$ with $\alpha_{*,k'}=\alpha_*(x_{k'})$, let $m\le k'$. By the
hypothesis on the form of $\alpha_*$ (the power one half of \cite[(2.6)]{B-CMP119} and the power
$q_*$ of the logarithm),
\begin{equation}\label{4.22:eq-budget-closes-b}
\frac{\alpha_{*,k'}}{\alpha_{*,m}}
=\Bigl(\frac{x_m}{x_{k'}}\Bigr)^{1/2}\Bigl(\frac{\log x_{k'}}{\log x_m}\Bigr)^{q_*}
\le2\Bigl(1+\frac{B_\sharp(k'-m)}{x_{k'}}\Bigr)^{1/2}.
\end{equation}
The inequality is obtained factor by factor. For the first factor we use the upper flow bound
$x_m\le x_{k'}+B_\sharp(k'-m)$. For the second, the lower flow bound gives
$x_m\ge x_{k'}-c_5$, hence
\begin{equation*}
\frac{\log x_{k'}}{\log x_m}\le\frac{\log x_{k'}}{\log(x_{k'}-c_5)},
\end{equation*}
the logarithms being positive by the member $\log(x-c_5)\ge1$ of
\eqref{4.22:eq-coupling-ceilings-h}. Since $q_*\ge0$, the member
$(\log x/\log(x-c_5))^{q_*}\le2$ of \eqref{4.22:eq-coupling-ceilings-h} then bounds the second
factor by $2$.

Take $m=k_W$. By the definition of $\mathsf{a}_{k'}(s)$ with $m(s)=k_W$,
\begin{equation*}
k'-k_W\le N_{k'}+\mathsf{a}_{k'}(s).
\end{equation*}
The members $L(\log x)^r\le x$ and $\log(x-c_5)\ge1$ of \eqref{4.22:eq-coupling-ceilings-h},
together with $N_{k'}\le L(\log x_{k'})^r$, give $N_{k'}\le x_{k'}$ and $x_{k'}\ge1$. Hence
\begin{equation*}
1+\frac{B_\sharp(k'-k_W)}{x_{k'}}
\le1+B_\sharp+B_\sharp\mathsf{a}_{k'}(s)
\le(1+B_\sharp)\bigl(1+\mathsf{a}_{k'}(s)\bigr).
\end{equation*}
By \eqref{4.22:eq-budget-closes-b}, the ratio $\alpha_{*,k'}/\alpha_{*,k_W}$ is therefore at most
$2(1+B_\sharp)^{1/2}(1+\mathsf{a}_{k'}(s))^{1/2}$. Moreover
$(1+\mathsf{a})^{1/2}e^{-2\mathsf{a}}\le1$ for $\mathsf{a}\ge0$, since $1+\mathsf{a}\le
e^{\mathsf{a}}\le e^{4\mathsf{a}}$. Writing $1/\alpha_{*,k_W}=(\alpha_{*,k'}/\alpha_{*,k_W})/
\alpha_*(x_{k'})$, the quotient of the far debits by $\mathsf{e}(s)$ is at most
\begin{equation*}
\mathsf{F}(x_{k'})\cdot
\frac{8C_{\mathrm{fr}}(1+B_\sharp)^{1/2}}{\varpi_{\mathrm{fr}}\,\alpha_*(x_{k'})}
\le(x_{k'}-c_5)^{-2}=\mu_{k'}/\varepsilon_{\star},
\end{equation*}
by the ceiling \eqref{4.22:eq-coupling-ceilings-k}. Thus the far debits are at most
$\mu_{k'}\mathsf{e}(s)/\varepsilon_{\star}$, and with the three margins
$\eta_{k'}(s)\le4\mu_{k'}\mathsf{e}(s)/\varepsilon_{\star}$. Frame sites carry no near debit:
by clause (i) of Lemma~\ref{4.22:lem-law-site}, the law sites of frames are never of kind (n).

\emph{Clause \textup{(1$^{\mathrm{res}}$)}.} By definition
$\eta^{\mathrm{inc}}_{k'}(s)=\mu_{k'}\mathsf{e}(s)/\varepsilon_{\star}$. Adding this to
\textup{(1)} gives \textup{(1$^{\mathrm{res}}$)}.

\emph{Clauses \textup{(2)} and \textup{(2$^{\mathrm{res}}$)}.} For $k'\le k$, the lower flow
bound gives $x_{k'}\ge x_k+\lambda_\sharp(k-k')-c_5$, that is,
\begin{equation*}
x_{k'}-c_5\ge x_k+\lambda_\sharp(k-k')-2c_5=y+\lambda_\sharp(k-k').
\end{equation*}
Since $x_k>\gamma_H^{-2}$, we have $y\ge y_\flat$, and $y_\flat>0$ by
\eqref{4.22:eq-coupling-ceilings-c}. Hence
\begin{equation*}
\mu_{k'}/\varepsilon_{\star}=(x_{k'}-c_5)^{-2}\le\bigl(y+\lambda_\sharp(k-k')\bigr)^{-2}.
\end{equation*}
As $k'$ runs over the steps $\le k$, the integer $n=k-k'\ge0$ takes each value at most once.
For $n\ge1$, the function $t\mapsto(y+\lambda_\sharp t)^{-2}$ is decreasing on $[n-1,n]$, so
$(y+\lambda_\sharp n)^{-2}\le\int_{n-1}^{n}(y+\lambda_\sharp t)^{-2}\,dt$. The ceiling
\eqref{4.22:eq-coupling-ceilings-c} states $y_\flat^{-2}+(\lambda_\sharp y_\flat)^{-1}\le1/32$.
The left member is decreasing in $y>0$, so the same bound holds at $y\ge y_\flat$. Therefore,
and recording the arithmetic used below,
\begin{equation}\label{4.22:eq-budget-closes-c}
\begin{gathered}
\sum_{n\ge0}(y+\lambda_\sharp n)^{-2}\le y^{-2}+\int_0^\infty(y+\lambda_\sharp t)^{-2}\,dt
=y^{-2}+(\lambda_\sharp y)^{-1}\le\tfrac1{32},\\
\tfrac5{32}+\tfrac2{32}=\tfrac7{32}<\tfrac8{32}=\tfrac14.
\end{gathered}
\end{equation}
Consequently
\begin{align*}
\sum_{k'\le k}4\mu_{k'}/\varepsilon_{\star}&\le4\bigl[y^{-2}+(\lambda_\sharp y)^{-1}\bigr]
\le\tfrac18,\\
\sum_{k'\le k}5\mu_{k'}/\varepsilon_{\star}&\le5\bigl[y^{-2}+(\lambda_\sharp y)^{-1}\bigr]
\le\tfrac5{32}.
\end{align*}
At a pointwise site, clause (iv) of Lemma~\ref{4.22:lem-law-site} gives
\begin{equation*}
\nu(s)\le3\beta K_u\sup_{j\le k}\alpha^+_j=3\beta K_u\sup_{j\le k}\bar\alpha(x_j).
\end{equation*}
The member $6\beta K_u\bar\alpha(x)\le\varepsilon_{\star}/16$ of the ceiling (vii) of the
statement on coupling ceilings and the order of constants, applied at $x=x_j>\gamma_H^{-2}$,
gives $3\beta K_u\bar\alpha(x_j)\le\varepsilon_{\star}/32$, hence
$\nu(s)\le\varepsilon_{\star}/16=\mathsf{e}(s)/16$. This summand occurs once, and at frame sites
it is absent. Summing \textup{(1)} over $k'\le k$ gives
\begin{equation*}
\sum_{k'\le k}\eta_{k'}(s)\le
\mathsf{e}(s)\Bigl\{4\bigl[y^{-2}+(\lambda_\sharp y)^{-1}\bigr]+\tfrac1{16}\Bigr\},
\end{equation*}
which is at most $(\tfrac18+\tfrac1{16})\mathsf{e}(s)<\tfrac14\mathsf{e}(s)$. This is
\textup{(2)}. The same argument with \textup{(1$^{\mathrm{res}}$)}, that is, with $5$ in place
of $4$, gives the first inequality of \textup{(2$^{\mathrm{res}}$)}. The bound
$(\tfrac5{32}+\tfrac2{32})\mathsf{e}(s)<\tfrac14\mathsf{e}(s)$ is the second line of
\eqref{4.22:eq-budget-closes-c}.

\emph{Clause \textup{(3)} and the uniformity.} The radius $\varepsilon_k(s)$ is the credit
$\mathsf{e}(s)$ diminished by the debits and reserves charged against it up to step $k$. These
are non-negative, and by \textup{(2$^{\mathrm{res}}$)} they total less than
$\tfrac14\mathsf{e}(s)$; this gives \textup{(3)}. The right members of \textup{(2)} and
\textup{(2$^{\mathrm{res}}$)} are functions of $\mathsf{e}(s)$, $y$ and $\lambda_\sharp$ alone,
and the numerical bounds used $y\ge y_\flat$ alone. As $\gamma_H\to0$, one has
$y_\flat=\gamma_H^{-2}-2c_5\to\infty$, and $y\ge y_\flat$ gives
\begin{equation*}
y^{-2}+(\lambda_\sharp y)^{-1}\le y_\flat^{-2}+(\lambda_\sharp y_\flat)^{-1}\longrightarrow0;
\end{equation*}
at a pointwise site the bound above gives
$\nu(s)/\varepsilon_{\star}\le3\beta K_u\sup_{x>\gamma_H^{-2}}\bar\alpha(x)/\varepsilon_{\star}$.
\end{proof}

\begin{remark}[The coupling flow within the level induction]\label{4.22:rem-flow-induction}
The radius budget (Theorem~\ref{4.22:thm-budget-closes}) assumes the upper and lower bounds of the
coupling flow at indices at most $k$. These are the two bounds of clause~(0) of Theorem~0
(Theorem~\ref{4.1:thm-clause-zero}). The appeal to them is not circular: at step $k$ the
radius budget uses only bounds that the level induction has already supplied at that stage.

Indeed, the lower and upper flow bounds make up clause~(0) of Theorem~0
(Theorem~\ref{4.1:thm-clause-zero}). Clause~(0) is established at the same place as clause~(i):
clauses~(0) and~(i) are proved by one simultaneous induction on the level $k$.
Theorem~\ref{4.22:thm-budget-closes}, applied at step $k$, reads the flow bounds only at indices
$k'\le k$. That is, it reads them only for the couplings $g_0,\dots,g_k$, equivalently for the
inverse couplings $x_{k'}=g_{k'}^{-2}$ with $k'\le k$. By the order of operations within a step,
the couplings $g_0,\dots,g_k$ are fixed before Ba{\l}aban's large-field operation $\mathbf{R}$ of
step $k$ is performed. Hence every quantity about which the flow hypothesis of
Theorem~\ref{4.22:thm-budget-closes} speaks at step $k$ is determined before that step's
large-field site is reached. The flow hypothesis at step $k$ is therefore read inside the single
induction on $k$ by which clauses~(0) and~(i) are proved, and no circularity arises.

We add three comments on the radius budget.

First, consider sites of untouched layered regions. Here a \emph{layered region} is a large-field
region together with its nested layers, and it is \emph{untouched} if none of its layers has been
touched by a second-class outcome of the large-field stages. For these sites a stronger factor is
available than the one used in the proof of Theorem~\ref{4.22:thm-budget-closes}, where the age
factor enters only through $e^{-2\mathsf{a}}\le 1$ at pointwise sites and through
$(1+\mathsf{a})^{1/2}e^{-2\mathsf{a}}\le 1$ at frame sites. By the decay statement for untouched
layered regions, in an untouched layered region created at step $m$ the factor is
\begin{equation*}
e^{-5\delta_1 R_k}\,e^{-2(k-1-m)},
\end{equation*}
which is geometric in the age; here $\delta_1$ and $R_k$ are the decay rate and the radius at
level $k$ of that statement, and $m$ denotes this step, not the torus exponent. The lower flow
bound is needed in three places: for the sites of regions that are not untouched; for the first
fluctuation bound and for the fluctuation-integral site $(\mathrm{T})$ at age $\mathsf{a}=0$; and
for the margins $\mu_k$.

Second, without the layer count $\theta_S$, or an equivalent localisation statement for the
current in its place, a further condition on the step would have to be carried as a named
hypothesis. The class of untouched layered regions, which is the class the inheritance theorem
works with, closes without that hypothesis.

Third, the condition $r\ge 2$ on the exponent in the bounds $N_{k'}\le L(\log x_{k'})^r$ and
$(\log(x_{k'}-c_5))^r\le\check{R}_{k'}$ is used in two places only: the ceilings on the far
debits at pointwise sites and at frame sites in the proof of
Theorem~\ref{4.22:thm-budget-closes}. What it supplies there is the super-polynomial smallness of
$e^{-\delta_\natural\check{R}_{k'}}$. At $r=1$ the factor
$\exp(-\delta_\natural[(\log(x-c_5))^r-2])$ of the far-debit majorant becomes
\begin{equation*}
\exp\bigl(-\delta_\natural[\log(x-c_5)-2]\bigr)=e^{2\delta_\natural}(x-c_5)^{-\delta_\natural},
\end{equation*}
and it is multiplied by the prefactor $(L\log x+3)C_\natural\bar\alpha(x)$, which contains the
growing factor $L\log x+3$. The ceilings compare this product with the margin
$\mu(x)=\varepsilon_\star(x-c_5)^{-2}$, so at $r=1$ they would require $\delta_\natural$ to be
strictly larger than $2$, with room to spare. That is a lower bound on a rate depending on
$(d,L,M)$, and no such bound is available. The condition $r\ge 2$ belongs to the setting of
Theorem~0.
\end{remark}

\begin{definition}[Coupling ceilings and the order of constants]\label{4.22:def-coupling-ceilings}
Throughout, $x=g^{-2}$ denotes an inverse squared coupling and $x_k=g_k^{-2}$. The two ceilings
$\gamma$ and $\gamma_H$ of the glossary (Notation~\ref{0.2:not-glossary}) are chosen last (see
the order of choice below), and, by Theorem~\ref{4.1:thm-clause-zero}, $x_k>\gamma_H^{-2}$ at every
level $k\le K$.

\emph{Constants and functions.} The following are fixed before $\gamma$: the rotation radius
$\varepsilon_{\star}$ at pointwise law sites, which depends on $(\Theta_{\flat},c_{\eta})$ only; the
frame reserve fraction $\varpi_{\mathrm{fr}}=\tfrac18(1-\theta_S)\beta$, where $\beta$ is the
schedule constant of Definition~\ref{4.4:def-post-step-space}; the frame constant
$C_{\mathrm{fr}}=3^d(C_{\chi}+4)$; and the rate $\delta_{\natural}$ and constant $C_{\natural}$ of
the far debits. The margin per step is the function $\mu(x)=\varepsilon_{\star}(x-c_5)^{-2}$, with
$\mu_k:=\mu(x_k)$; a frame $W$ created at step $k_W$ has the radius $\varepsilon^{\mathrm{fr}}(W)$,
where $4\varepsilon^{\mathrm{fr}}(W)=\varpi_{\mathrm{fr}}\alpha_{*,k_W}/C_{\mathrm{fr}}$. Two readings
are declared with these constants: Ba{\l}aban's cut-off function of the change of variables is
written $\mathsf{g}$, and the data of a stored term at a frame are kept on the stored domain of
Lemma~\ref{4.22:lem-frame-lemma}, which carries at every clause the reserve $\varpi_{\mathrm{fr}}$
times the threshold of that clause. We write
$\bar\alpha(x)$, $\alpha_{*}(x)$, $\mathfrak{b}(x)$, $T^{\mathbb{C}}(x)$,
$\varepsilon_{\mathrm{live}}(x)$ and $\varepsilon(x)$ for the explicit functions of $x$ of the
construction (the last being the function through which the lift lemma of item (ii) below bounds
$H_{\max}$), so that
$\bar\alpha(x_k)=\alpha^+_k$, $\alpha_*(x_k)=\alpha_{*,k}$, $\mathfrak{b}(x_k)=\mathfrak{b}_k$ and
$\varepsilon_{\mathrm{live}}(x_k)=\varepsilon_{\mathrm{live}}(k)$, and we put
\begin{equation}\label{4.22:eq-coupling-ceilings-1}
E(x):=\exp\bigl(-\bigl[(\log(x-c_5))^r-2\bigr]\bigr).
\end{equation}
Also fixed before $\gamma$ are the constants introduced later in this chapter: the composition
constant $C_b\ge1$ of Lemma~\ref{4.22:lem-axialising-map}, the half-rate resummation constant
$C_A^{(1/2)}$ of Corollary~\ref{4.22:cor-older-law-sites}, the bound $B_H^S$ that replaces $B_H$
on the cell families of the family clause of Lemma~\ref{4.22:lem-law-site} (item (viii) below),
the constants $C_\Sigma$, $p_\Sigma$ of the bound
$\mathfrak{O}(v)\le C_\Sigma\check R_{k_W}^{p_\Sigma}$ on the oscillation seminorm of a kept
member $v$, the constant $\mathfrak{s}_\Sigma$ of the leaf terms, and the number
$\mathfrak{n}^{\mathrm{kept}}$ of \eqref{4.22:eq-fixed-number-entry-a}; and $C_{\natural}$ is taken
in the normalisation of the family clause of Lemma~\ref{4.22:lem-law-site}, so that one majorant
serves both that clause and clause (iii) of the lemma. The margins of the radius
budget of Theorem~\ref{4.22:thm-budget-closes} and its reserve
$\eta^{\mathrm{inc}}_k(s):=\mu_k\mathsf{e}(s)/\varepsilon_{\star}$ are chosen, not derived.

\emph{Convention on the flow constants.} Inside every ceiling below, $B_\sharp$ is read as
$\bar B=C_\beta+1$ and, where stated, $\lambda_\sharp$ as $\underline{\lambda}=c_0/4$
(Definition~\ref{1.2:def-flow-constants}): the ceilings are evaluated at the triple
$(c_0/4,\,c_5,\,C_\beta+1)$. This is necessary because $B_\sharp$ itself exceeds $C_\beta$ by a
quantity depending on $\gamma$, so that reading $B_\sharp$ literally would place $\gamma$ inside a
condition imposed on $\gamma$.

\emph{Form of the ceilings.} Except for \eqref{4.22:eq-coupling-ceilings-c} and
\eqref{4.22:eq-coupling-ceilings-g}, which are read at the single point $\gamma_H^{-2}$, each ceiling
is one inequality required for every $x\ge\gamma_H^{-2}$. None involves a quantity depending on the
run: the scale $\check R_k$ and the number $N_k$ enter only through the bounds
\begin{equation}\label{4.22:eq-coupling-ceilings-2}
(\log(x_k-c_5))^r\le\check R_k<L(\log x_k)^r,\qquad N_k\le L(\log x_k)^r ,
\end{equation}
which hold because the scale $R(x)$ of Ba{\l}aban's construction \cite[\S1]{B-CMP122a} satisfies
$(\log x)^r\le R(x)<L(\log x)^r$ and is non-decreasing in $x$, because the monotone radius is
$\check R_k=R(\min_{i\le k}x_i)$, where
\begin{equation*}
x_k-c_5\le\min_{i\le k}x_i\le x_k
\end{equation*}
by the lower bound of the coupling flow (Theorem~\ref{4.1:thm-clause-zero}), and because
$N_k:=\check R_k$.
Every ceiling holds once $\gamma_H$ is small, every constant chosen before $\gamma$ being held fixed:
the left members of \eqref{4.22:eq-coupling-ceilings-c}, \eqref{4.22:eq-coupling-ceilings-d},
\eqref{4.22:eq-coupling-ceilings-e}, \eqref{4.22:eq-coupling-ceilings-f},
\eqref{4.22:eq-coupling-ceilings-a}, \eqref{4.22:eq-coupling-ceilings-b} and of the remaining upper
bounds tend to zero; in \eqref{4.22:eq-coupling-ceilings-j} and \eqref{4.22:eq-coupling-ceilings-k}
the quotient of the left member by the right member tends to zero; the left members of
\eqref{4.22:eq-coupling-ceilings-3}, \eqref{4.22:eq-coupling-ceilings-g} and of $\log(x-c_5)\ge1$
tend to infinity; the two ratios in
\eqref{4.22:eq-coupling-ceilings-h} tend to $1$; and in $L(\log x)^r\le x$ and in
\eqref{4.22:eq-coupling-ceilings-i} the right member dominates the left as $x\to\infty$. The
ceilings are imposed last. Each bounds
the coupling from above only; none is a cap in the sense of Definition~\ref{2.3:def-cap}.

\emph{The ceilings.} For all $x\ge\gamma_H^{-2}$, with $g=x^{-1/2}$:
\begin{itemize}
\item[(i)] The small-field threshold is large:
\begin{equation}\label{4.22:eq-coupling-ceilings-3}
\mathfrak{b}(x)\ge 64 .
\end{equation}
\item[(ii)] Let $c_B$, $C_B$, $\delta_{\mathrm{lift}}$, $G_1(N)$ be the constants of the lemma on the
holomorphic lift of Ba{\l}aban's change of variables, and let $H_{\max}$ be the supremum appearing
there. Then
\begin{equation}\label{4.22:eq-coupling-ceilings-4}
\begin{gathered}
gT^{\mathbb{C}}(x)+2H_{\max}<c_B,\qquad C_BgT^{\mathbb{C}}(x)\le1,\qquad 4H_{\max}/g\le1,\\
2C_BH_{\max}+4G_1(N)H_{\max}/\delta_{\mathrm{lift}}\le\tfrac12 .
\end{gathered}
\end{equation}
Here $H_{\max}$ is bounded by a constant multiple of $\varepsilon(x)E(x)$, because
$e^{-\check R_k}\le E(x_k)$. The same four conditions are imposed with the constants of the version
of that lemma at the large-field stages.
\item[(iii)] \emph{The reserved shares.} Let $C_A$ be the resummation constant, at the full response
rate, of the piece lemma of the correlation theorem, and $C_A^{(1/2)}$ its half-rate form introduced
in Corollary~\ref{4.22:cor-older-law-sites}; let
$\Sigma_4$ be the share of the stored terms treated in that corollary; let
$\mathfrak{m}_{\mathrm{lift}}$ be the number of lifted terms of the change of variables counted in
the dichotomy step of its proof; and let $C_{\mathrm{rn}}\varrho^2e^{-\delta_1MR_{\mathrm{rn}}}$ be
the bound on the remainder terms entering these shares, $C_{\mathrm{rn}}$ being its constant and
$R_{\mathrm{rn}}$ its length scale. The shares reserved in
Corollary~\ref{4.22:cor-older-law-sites} and in the flat total of the large-field stages, which tend
to zero with $\gamma$, contain the following quantities:
\begin{equation}\label{4.22:eq-coupling-ceilings-l}
\begin{gathered}C_{22}(N)\bigl(1+G_1(N)\bigr)\frac{A_0}{A_1}\,T^{\mathbb{C}}(x)^2E(x)\cdot\mathfrak{m}_{\mathrm{lift}},
\qquad 4\,C_{\mathrm{rn}}\varrho^2e^{-\delta_1MR_{\mathrm{rn}}},\\
\Sigma_{4}\cdot C_A^{(1/2)}/C_A .
\end{gathered}
\end{equation}
The first quantity is the lifted term of the change of variables, multiplied by
$\mathfrak{m}_{\mathrm{lift}}$; it enters together with its counterpart at the large-field
stages, the quantity bounded in clause (iv) of the stage version of the lift lemma. The second is
four times the bound on the remainder terms. The third is the share $\Sigma_4$ multiplied by
$C_A^{(1/2)}/C_A$; the price of the margin is paid by \eqref{4.22:eq-coupling-ceilings-e} and is not
a factor here. This share is taken separately for the plain class $\mathfrak{A}$ and the wrapped
class $\mathfrak{S}$. For $\mathfrak{S}$, the share can be made to tend to zero with $\gamma$ because
the ceiling \eqref{4.22:eq-fluctuation-site-a} places the number $\mathfrak{n}^{\mathrm{kept}}$
under the factor $e^{-\frac14\delta\kappa(7\check R_k-2)}$ carried by these terms, $\delta$ being the
small parameter in the exponential weights of these terms; no decay factor of the terms
$V(Y,\mathcal{S}')$ of that corollary is drawn on.
\item[(iv)] The rotation size used at the presentation satisfies
\begin{equation}\label{4.22:eq-coupling-ceilings-5}
\varepsilon_{\mathrm{live}}(x)=K''a_1+\varepsilon_{\mathfrak{Q}}\le\tfrac14\varepsilon_{\star},
\end{equation}
where $\varepsilon_{\mathfrak{Q}}$ is the size of the presenter maps $\mathfrak{Q}^{\mathfrak{b}}$
and $K''a_1$, a constant $K''$ times an amplitude $a_1$, is the other summand of the rotation size.
\item[(v)] With $\lambda_\sharp$ read as $c_0/4$,
\begin{equation}\label{4.22:eq-coupling-ceilings-c}
y_\flat^{-2}+(\lambda_\sharp y_\flat)^{-1}\le\frac1{32}\qquad\text{at}\qquad
y_\flat:=\gamma_H^{-2}-2c_5>0 .
\end{equation}
\item[(vi)] Both of the following hold:
\begin{equation}\label{4.22:eq-coupling-ceilings-d}
\begin{gathered}
(x-c_5)^2E(x)^{\delta_\natural/2}\le1,\\
(x-c_5)^2\cdot 8e^{2\varepsilon_\star}\,C_{\mathrm{fr}}(1+B_\sharp)^{1/2}C_\natural\,
\frac{\bar\alpha(x)\,E(x)^{\delta_\natural/2}}{\varpi_{\mathrm{fr}}\,\alpha_*(x)}\le1 .
\end{gathered}
\end{equation}
The factor $8e^{2\varepsilon_\star}$ of the second member, which satisfies
$8e^{2\varepsilon_\star}\le8.05$, is the one required at frames in the half-rate form of
Corollary~\ref{4.22:cor-bracket-increment}\,(b). That form requires at frames that a parameter
$\rho_{\mathrm{fr}}$ of it be at least $2$; this holds by (xii) and (xiii) below, one reading costing
at most $\mathsf{F}(x)/6$ (see the proof below).
\item[(vii)] The following holds:
\begin{equation}\label{4.22:eq-coupling-ceilings-e}
(x-c_5)^2E(x)^{\delta_\natural}\cdot2\sqrt2\,C_b^{1/2}K_uB_H\cdot
\max\Bigl\{\frac{17e\,C_{\mathrm{fr}}(1+B_\sharp)^{1/2}}{\varpi_{\mathrm{fr}}},\,
\frac{2.1\,\alpha_*(x)}{\varepsilon_\star}\Bigr\}\le1 .
\end{equation}
Its left member tends to zero as $x\to\infty$.
\item[(viii)] With $\varrho_0:=\min\{10^{-2}/K_u,\;c_4(d,L),\;\alpha_1^{98}\}$, where
$\alpha_1^{98}$ is the number of \cite[(158)]{B-CMP98} (it depends on $d$ and $L$ only),
\begin{equation}\label{4.22:eq-coupling-ceilings-f}
B_H\,\bar\alpha(x)\le\varrho_0,\qquad 2.1K_uB_H\alpha_*(x)\le10^{-2},\qquad
6\beta K_u\bar\alpha(x)\le\varepsilon_\star/16 .
\end{equation}
The middle member is the bound $\vartheta^T\le10^{-2}$ on the quantity $\vartheta^T$ of
\eqref{4.22:eq-axialising-map-a}. By the first member, the supremum of
$\Theta_0^{\mathrm{loc}}(\mathcal{K})$, for the layers $\mathcal{K}$ of
Lemma~\ref{4.22:lem-axialising-map}, lies in the range of the existence statement for the
axialising gauge map on which clause (i) of Lemma~\ref{4.22:lem-axialising-map} rests. On the cell
families treated in the family clause of Lemma~\ref{4.22:lem-law-site}, $B_H$ is read as $B_H^S$ in
\eqref{4.22:eq-coupling-ceilings-f}; this is the premise of clause (iii) of
Proposition~\ref{4.22:prop-gaussian-oscillation}.
\item[(ix)] The following holds:
\begin{equation}\label{4.22:eq-coupling-ceilings-g}
5\delta_\natural\bigl(\log(\gamma_H^{-2}-c_5)\bigr)^r\ge2 .
\end{equation}
\item[(x)] All of the following hold:
\begin{equation}\label{4.22:eq-coupling-ceilings-h}
\begin{gathered}
\Bigl(\frac{\log x}{\log(x-c_5)}\Bigr)^{q_*}\le2,\qquad L(\log x)^r\le x,\\
\log(x-c_5)\ge1,\qquad \Bigl(\frac{\log(x+B_\sharp)}{\log x}\Bigr)^{q_*}\le2 .
\end{gathered}
\end{equation}
Here $q_*$ is the exponent of $\log x$ in the radius $\alpha_*(x)$.
\item[(xi)] The following holds:
\begin{equation}\label{4.22:eq-coupling-ceilings-i}
4+2r\log_L\log\bigl((1+B_\sharp)x\bigr)\le\bigl(\log(x-c_5)\bigr)^r .
\end{equation}
\item[(xii)] \emph{The sum ceiling.} The following holds:
\begin{equation}\label{4.22:eq-coupling-ceilings-j}
\mathsf{F}(x):=\bigl(L(\log x)^r+3\bigr)C_\natural\,\bar\alpha(x)\,E(x)^{\delta_\natural}
\le\mu(x).
\end{equation}
\item[(xiii)] \emph{The frame ceiling.} The following holds:
\begin{equation}\label{4.22:eq-coupling-ceilings-k}
\mathsf{F}(x)\cdot\frac{8C_{\mathrm{fr}}(1+B_\sharp)^{1/2}}{\varpi_{\mathrm{fr}}\,\alpha_*(x)}
\le(x-c_5)^{-2}.
\end{equation}
\item[(xiv)] \emph{The arrest ceilings.} Three conditions are imposed at the amplitude
$6.6\,C_1\delta'(x)$, where $\delta'(x)$ is the explicit function of $x$ with
$\delta'(x_k)=\delta'_k$, the third of the radii whose maximum is $\alpha^+_k$: the absorption
condition, the condition on the weak collar variables, and the
window condition of the inheritance theorem. This amplitude is six times the amplitude
$1.1\,C_1\delta'(x)$ at which they are imposed for sealed terms, and each condition is homogeneous of
degree one in the amplitude. In addition, $0.6\,\delta'(x)/g\ge3$.
\item[(xv)] \emph{The regularity ceilings.} With $B$, $C$ the constants of the cube-gauge clause
\cite[(2.38)]{B-CMP119},
\begin{equation}\label{4.22:eq-coupling-ceilings-a}
2.1(2C_\chi+1)BCM\,\alpha_0(x)\le1\qquad\text{and}\qquad
2.1(2C_\chi+1)BCM\,\bar\alpha(x)\le1 ,
\end{equation}
and
\begin{equation}\label{4.22:eq-coupling-ceilings-b}
\mathfrak{c}\,\bar\alpha(x)\le1,\qquad \mathfrak{c}:=\max\{20\bar r_6,\mathfrak{c}_6,\mathfrak{c}_7\},
\end{equation}
where
\begin{equation*}
\begin{aligned}
\mathfrak{c}_6&:=\mathfrak{b}_1+\mathfrak{b}_2(\bar p_1+\varpi_{\mathrm{fr}}),\\
\mathfrak{c}_7&:=\mathfrak{b}_1\Bigl(1+\frac{\bar p_1\bar r_6}{\bar r_7}\Bigr)
+\mathfrak{b}_3\Bigl[\bar p_1(1+\varpi_{\mathrm{fr}})+4\varpi_{\mathrm{fr}}
+\frac{\bar p_2\bar r_6}{\bar r_7}\Bigr],
\end{aligned}
\end{equation*}
with the absolute numbers $\mathfrak{b}_1=2.05$, $\mathfrak{b}_2=200$, $\mathfrak{b}_3=400$, and
$\bar p_1=3^d(C_\chi+1)$, $\bar p_2=3^d(4C_\chi+2)$. Here the thresholds of the sixth and seventh
entries of \cite[(2.39)]{B-CMP119} at the level $k_W$ of a frame $W$ are
$r_a=\bar r_a\alpha_{1,k_W}$ ($a=6,7$), with pure numbers $\bar r_6,\bar r_7$, taken in the
normalisation of these thresholds on $W$ used in the proof of Lemma~\ref{4.22:lem-frame-lemma}\,(i),
in which $\bar r_6,\bar r_7\ge1$. Thus
$\mathfrak{c}$ is a number determined by $(d,C_\chi,\bar r_6,\bar r_7,\varpi_{\mathrm{fr}})$. When
the regularity ceiling is invoked without further specification, both members of
\eqref{4.22:eq-coupling-ceilings-a} together with \eqref{4.22:eq-coupling-ceilings-b} are meant.
\end{itemize}

\emph{Where they are used.}
\begin{itemize}
\item \eqref{4.22:eq-coupling-ceilings-3} is used in the strengthened bound of the fluctuation-site
analysis that reads the small-field threshold.
\item \eqref{4.22:eq-coupling-ceilings-4} is used in the lift lemma and its stage version.
\item \eqref{4.22:eq-coupling-ceilings-l} is used in the corollary of the lift lemma, in the stage
bound for the lifted terms, and in Corollary~\ref{4.22:cor-older-law-sites}.
\item \eqref{4.22:eq-coupling-ceilings-5} is used in Corollary~\ref{4.22:cor-bracket-increment}\,(a).
\item \eqref{4.22:eq-coupling-ceilings-c} is used in Theorem~\ref{4.22:thm-budget-closes}\,(2).
\item \eqref{4.22:eq-coupling-ceilings-d} is used in
Corollary~\ref{4.22:cor-bracket-increment}\,(b).
\item \eqref{4.22:eq-coupling-ceilings-e} is used in Corollary~\ref{4.22:cor-older-law-sites}, where
it gives $\Lambda^{\mathrm{res}}\le1$ in \eqref{4.22:eq-older-law-sites-a}.
\item \eqref{4.22:eq-coupling-ceilings-f} is used in Lemma~\ref{4.22:lem-axialising-map}, in
Lemma~\ref{4.22:lem-law-site}\,(iv), in
Proposition~\ref{4.22:prop-gaussian-oscillation}\,(iii), and in the lemma through which the pieces
are bounded in the proof of Corollary~\ref{4.22:cor-older-law-sites}\,$(\alpha)$.
\item \eqref{4.22:eq-coupling-ceilings-g} is used in Lemma~\ref{4.22:lem-law-site}\,(iii), in
Corollary~\ref{4.22:cor-older-law-sites} and in clause (iv) of
Lemma~\ref{4.22:lem-axialising-map}.
\item \eqref{4.22:eq-coupling-ceilings-h} is used in the proof of
Theorem~\ref{4.22:thm-budget-closes}, in \eqref{4.22:eq-coupling-ceilings-j} and in the price
estimate of Corollary~\ref{4.22:cor-older-law-sites}.
\item \eqref{4.22:eq-coupling-ceilings-i} is used in the shell estimate at frame sites in the proof of
Lemma~\ref{4.22:lem-law-site}.
\item \eqref{4.22:eq-coupling-ceilings-j} is used in Theorem~\ref{4.22:thm-budget-closes}\,(1), and
\eqref{4.22:eq-coupling-ceilings-k} in the same clause at frames.
\item The arrest ceilings are used in the arrest lemma and in the stage version of the lift lemma.
\item The second member of \eqref{4.22:eq-coupling-ceilings-a} is used in
Lemma~\ref{4.22:lem-frame-lemma}\,(i) only.
\item \eqref{4.22:eq-coupling-ceilings-b} is used in Lemma~\ref{4.22:lem-frame-lemma}\,(i), for the
domain of the triple-product logarithm ($r_6\le1/20$) and for the share of its non-linear part in the
sixth and seventh entries for gauge maps of size $z\le4\varepsilon^{\mathrm{fr}}(W)$. Through the
member $\mathfrak{c}_6$, it is also used in Lemma~\ref{4.22:lem-profile-logarithm}\,(e), within the
first step of Lemma~\ref{4.22:lem-frame-lemma}\,(ii).
\end{itemize}

\emph{Further ceilings imposed by their users.} The following are imposed in the same one-sided form,
for all $x$ in the stated range, and are discharged by the same last choice of $\gamma$ and
$\gamma_H$.
\begin{itemize}
\item[(a)] For all $x\ge\gamma_H^{-2}$, with the amplitudes $\mathsf{a}_T(x)$, $\mathsf{a}_L(x)$ of
the led rotations at the
fluctuation site and at the presentation (not to be confused with the age $\mathsf{a}_k(s)$ of a
frame site),
\begin{equation}\label{4.22:eq-coupling-ceilings-6}
\begin{split}
&2.1K_u\max\{\mathsf{a}_T(x),\mathsf{a}_L(x)\}\,(x-c_5)^2E(x)^{\delta_\natural/2}\\
&\qquad\cdot\max\Bigl\{\frac1{\varepsilon_\star},\,
\frac{8e\,C_{\mathrm{fr}}(1+B_\sharp)^{1/2}}{\varpi_{\mathrm{fr}}\,\alpha_*(x)}\Bigr\}\le1 .
\end{split}
\end{equation}
Here
\begin{equation*}
\begin{gathered}
K_u\mathsf{a}_T(x)=K_uB_H\alpha_*(x),\\
K_u\mathsf{a}_L(x)\le C_{\theta'}^{-1}(\log x+2)^{q-p_1}\cdot\tfrac12e^{-2}C_\natural\bar\alpha(x).
\end{gathered}
\end{equation*}
Here $q$ and $p_1$ are the exponents of the polylogarithmic factor of the amplitude at the
presentation (this $p_1$ is not the constant $p_1$ of Lemma~\ref{4.22:lem-frame-lemma}), and
$\theta'$ is the disc parameter of the correlation theorem, with constant $C_{\theta'}$.
Both the polylogarithmic factor $(\log x+2)^{q-p_1}$ and the factor $E(x)^{\delta_\natural/2}$,
which is smaller than every power of $x$, are kept, and the disc width at the presentation is
$\mathsf{w}_L:=1/\theta'$.
\item[(b)] The following holds:
\begin{equation*}
5\delta_\natural\bigl(\log(\gamma^{-2}-c_5)\bigr)^r\ge4 .
\end{equation*}
\item[(c)] The relative ceilings \eqref{4.22:eq-fluctuation-site-a} at the fluctuation site and
\eqref{4.22:eq-own-healing-a} at the healing of the own component, and the list of flat stage-total
conditions \eqref{4.22:eq-stage-totals-a}, are imposed. They belong to
Propositions~\ref{4.22:prop-fluctuation-site} and~\ref{4.22:prop-own-healing} and to
Proposition~\ref{4.22:prop-stage-totals}, whose proof is incomplete at one step. Each is the form,
printed in the proposition to which it belongs, obtained from the corresponding relative ceiling of
the analysis of the relative terms, on which the kept-member propositions build, by multiplying its
left member by the number
$\mathfrak{n}^{\mathrm{kept}}$, all other factors being unchanged; each is required as one bound
over all $x$ in its range. The last is the whole list of flat stage-total conditions at the value of
$\vartheta$ in force, imposed both as it stands and with
$\mu^\flat$ replaced by $\mu^\flat+\mu^{\mathrm{leaf}}$ and $\mathsf{M}^\flat$ by
$\mathsf{M}^\flat+\mathsf{M}^{\mathrm{leaf}}$, $\mu^{\mathrm{leaf}}$ and $\mathsf{M}^{\mathrm{leaf}}$
being the leaf majorant and the leaf total of \eqref{4.22:eq-stage-totals-a}.
\item[(d)] With $x_\gamma:=\gamma^{-2}$,
\begin{equation}\label{4.22:eq-coupling-ceilings-m}
\sup_{x\ge x_\gamma}\tfrac{4e^2}{3}\,B_H^S\,E(x)^{\delta_\natural}\le1 .
\end{equation}
For $x\ge\gamma_H^{-2}$, the inequality under the supremum is implied by the first member of
\eqref{4.22:eq-coupling-ceilings-d} together with
$(x-c_5)^4\ge\frac{4e^2}3B_H^S$. The constant $B_H^S$ is chosen after $\kappa_1$ and before
$\gamma$. Under the reading through \cite[(158)]{B-CMP98} of the existence statement for the
axialising gauge map on which clause (i) of Lemma~\ref{4.22:lem-axialising-map} rests, stated in
clause (iii) of Proposition~\ref{4.22:prop-gaussian-oscillation}, this ceiling is not consumed; it
would be consumed only under the alternative reading of that statement through the estimates of
\cite{B-CMP109}, and imposing it costs nothing.
\item[(e)] The remaining ceilings imposed by their users are: the membership and arrest ceilings in
their half forms; the ceiling on $\gamma$ of the three-field analysis, together with its three
accompanying conditions; the tail ceiling, the ceiling on the scale $\check R_k$, a further ceiling
of the kept-member analysis, and the ceiling on the kept terms, the last taken with the constant
$C_{\mathcal{A}}^\sharp$ of the kept terms at its value $C_{\mathcal{A}}^\sharp=71.2$; the ceilings
of the source-tube bounds on the bounds $\mathfrak{q}$ of the
kept quadratic form, on nesting and on the window, together with one further ceiling of those bounds;
and one further ceiling of the same one-sided form. They are stated, each in the form in which it is
imposed, in Proposition~\ref{4.22:prop-gaussian-oscillation}, in the kept-member
Propositions~\ref{4.22:prop-fluctuation-site}, \ref{4.22:prop-stage-totals} (whose proof is
incomplete at one step) and~\ref{4.22:prop-own-healing}, in the led-rotation accounts of the
correlation theorem and of the three-field analysis, and in the hypotheses under which those accounts
are stated.
\end{itemize}

\emph{Non-voidness.} For $r\ge2$ and $\log(x-c_5)\ge1$, and every $\varsigma>0$,
\begin{equation}\label{4.22:eq-coupling-ceilings-7}
E(x)^{\varsigma}\le e^{2\varsigma}(x-c_5)^{-\varsigma\log(x-c_5)} .
\end{equation}
Hence $E(x)^\varsigma$ decays faster than every power of $x$. Moreover $\bar\alpha(x)/\alpha_*(x)$ is
bounded by a power of $\log x$, and is a constant when the two exponents $q_0$, $q_1$ of the radii
coincide.

\emph{Order of choice.} The constants are chosen in the order
\begin{equation*}
\begin{gathered}
L\to\bigl(K_0,\ \delta_1,\ \Theta_\flat,\ c_0',\ \delta_W,\ \delta_G,\ K_u,\ C_Q,\ C_\chi,\ c_4,\
\alpha_1^{98}\bigr)\to\kappa\to\kappa_1\to\bigl(B_H,\ B^{\mathrm{prof}}_{\sharp},\ K_\sharp\bigr)
\to M\\
\to(\gamma_2,c_\eta)\to\varepsilon_\star\to\theta_S\to K_R\to
\bigl(\varpi_{\mathrm{fr}},\ C_{\mathrm{fr}},\ \mathfrak{c},\ \delta_\natural,\ B_H^S,\ C_\natural,\
C_A^{(1/2)},\ C_b\bigr)\\
\to A_0/A_1\to C_{\mathrm{far}}\to c_A\to(C_\Sigma,\ p_\Sigma)\to\mathfrak{n}^{\mathrm{kept}}
\to\gamma\to\gamma_H .
\end{gathered}
\end{equation*}
The constants in the first bracket after $L$ are numbers depending on $(d,L)$; the various decay
rates denoted $\delta_1$ are among these $(d,L)$-numbers. Among them, $K_u$ is the constant of the
bound $|\log\mathfrak{u}(z)|\le K_u\Theta_0^{\mathrm{loc}}(\mathcal{K})(z)$ of
Lemma~\ref{4.22:lem-axialising-map}\,(i); $c_4$ is the bound on $\Theta_0(\mathcal{K})$ in the
hypothesis of the existence statement for the axialising gauge map; $\delta_G$ is the decay rate,
per $M$-cube, of the kernels of the propagators of \cite[Theorem~3.13]{B-CMP99b}; and
$C_Q=1+\|Q^*\|\,\|(QQ^*)^{-1}Q\|$ is the constant of the averaging operators.
$B_H$ is the bound on the polydisc $|\sigma|\le e^{\kappa_1}$ of
the decoupling parameters. $B^{\mathrm{prof}}_{\sharp}$ and $K_\sharp$ are the constants of the
profile bound of the Landau layer; the former is so written to distinguish it from the flow constant
$B_\sharp$. The constant $\gamma_2$ is chosen together with the constant $c_\eta$ of the strong-set
weights; $K_R$ is the constant in the unit shift $\mathsf{U}$ of
\eqref{4.22:eq-axialising-map-a}; and $C_{\mathrm{far}}$ and $c_A$ are constants fixed at the
positions shown.

The choices obey the following dependencies.
\begin{itemize}
\item $\varepsilon_\star$ depends on $(\Theta_\flat,c_\eta)$ only, and $\varpi_{\mathrm{fr}}$ on
$(\theta_S,\beta)$ only.
\item $\mathfrak{c}$ depends on $(d,C_\chi,\bar r_6,\bar r_7,\varpi_{\mathrm{fr}})$ and is chosen
together with $C_{\mathrm{fr}}$. The number $2C_\chi+1$ of \eqref{4.22:eq-coupling-ceilings-a} is
chosen together with $C_\chi$, among the $(d,L)$-numbers.
\item No constant chosen before these depends on them, and no choice returns to $M$ or to $L$.
\item $C_\Sigma$ is chosen after the constant $\mathsf{K}^0$ and before
$\mathfrak{n}^{\mathrm{kept}}$ and $\gamma$, and does not depend on $\mathfrak{n}^{\mathrm{kept}}$.
In turn, $\mathfrak{n}^{\mathrm{kept}}$ enters the ceilings on $\gamma$ only, so the order has no
cycle.
\end{itemize}
The order is compatible with the order
$L\to\mathfrak{M}\ (R_1\ge L^2)\to\kappa\to\kappa_1\to M_1\to M\to\gamma$ of the three-field
analysis, which is acyclic; here $\mathfrak{M}$ is the parameter of that analysis subject to
$R_1\ge L^2$, $R_1$ being its first length scale. The constants $C_{3F}$ and
$C_{\mathcal{A}}^\sharp$ are chosen after $M$
and before $\gamma$. The conditions of the correlation theorem on its parameters are read as they
stand, in particular
\begin{equation*}
\kappa_1\ge\max\{16\kappa+1,\ 4L\kappa+1,\ 1+c_l(2\kappa+4)+\lambda_*,\ \kappa_1^{\flat}\},
\end{equation*}
where $\kappa_1^{\flat}$ denotes the lower bound on $\kappa_1$ of the flat form of the correlation
theorem's condition on $\kappa_1$, together with its conditions on $M/M_1$ and on $C_{\theta'}$,
and $\mathsf{L}_0\ge3$; here $c_l$, $\lambda_*$, $C_{\theta'}$ and $\mathsf{L}_0$ are constants of
the correlation theorem. The third and fourth members of this maximum, and the conditions on
$M/M_1$, are used as they stand in the kept-member propositions.
\end{definition}

\begin{proof}
We verify the assertions made in the definition.

\emph{Reading of the flow constants.} In each ceiling in which $B_\sharp$ occurs, the left member
is non-decreasing in $B_\sharp$. These ceilings are the second member of
\eqref{4.22:eq-coupling-ceilings-d}, \eqref{4.22:eq-coupling-ceilings-e}, the last member of
\eqref{4.22:eq-coupling-ceilings-h}, \eqref{4.22:eq-coupling-ceilings-i},
\eqref{4.22:eq-coupling-ceilings-k} and \eqref{4.22:eq-coupling-ceilings-6}; the check is made on
the factors $(1+B_\sharp)^{1/2}$, $\log(x+B_\sharp)$ and $\log((1+B_\sharp)x)$. Since
$B_\sharp\le C_\beta+1$, each of these ceilings, read at $C_\beta+1$, implies the same inequality at
the actual $B_\sharp$. The left member of \eqref{4.22:eq-coupling-ceilings-c} is decreasing in
$\lambda_\sharp$, and $\lambda_\sharp\ge c_0/4$; so \eqref{4.22:eq-coupling-ceilings-c} read at
$c_0/4$ implies it at $\lambda_\sharp$.

\emph{Monotonicity in the ceiling.} A condition required for all $x\ge\gamma_H^{-2}$ remains true
when $\gamma_H$ is decreased. The quantity $y_\flat=\gamma_H^{-2}-2c_5$ increases as $\gamma_H$
decreases, so the left member of \eqref{4.22:eq-coupling-ceilings-c} decreases. Likewise the left
member of \eqref{4.22:eq-coupling-ceilings-g} increases as $\gamma_H$ decreases. The same holds in
$\gamma$ for (b) and \eqref{4.22:eq-coupling-ceilings-m}. Hence the set of admissible $\gamma_H$
(resp.\ $\gamma$) is closed under decrease.

\emph{The bound $e^{-\check R_k}\le E(x_k)$.} By \eqref{4.22:eq-coupling-ceilings-2},
$\check R_k\ge(\log(x_k-c_5))^r$. Hence
\begin{equation*}
e^{-\check R_k}\le\exp\bigl(-(\log(x_k-c_5))^r\bigr)=e^{-2}E(x_k)\le E(x_k).
\end{equation*}

\emph{The cost of one reading.} By Lemma~\ref{4.22:lem-law-site}\,(iii), one reading at step $k$
costs at most $\frac12C_\natural\alpha^+_ke^{-\delta_\natural(\check R_k-2)}$, the factor
$e^{-2\mathsf{a}_k(s)}\le1$ being dropped. By \eqref{4.22:eq-coupling-ceilings-2},
$\check R_k-2\ge(\log(x_k-c_5))^r-2$, so that
$e^{-\delta_\natural(\check R_k-2)}\le E(x_k)^{\delta_\natural}$. Since
$\alpha^+_k=\bar\alpha(x_k)$ and $L(\log x_k)^r+3\ge3$,
\begin{equation*}
\tfrac12C_\natural\alpha^+_ke^{-\delta_\natural(\check R_k-2)}
\le\tfrac12C_\natural\bar\alpha(x_k)E(x_k)^{\delta_\natural}\le\tfrac16\mathsf{F}(x_k).
\end{equation*}
From this bound, \eqref{4.22:eq-coupling-ceilings-j} and \eqref{4.22:eq-coupling-ceilings-k} give
$\rho_{\mathrm{fr}}\ge2$ at frames in Corollary~\ref{4.22:cor-bracket-increment}\,(b).

\emph{Non-voidness.} Put $t:=\log(x-c_5)\ge1$. Since $r\ge2$ we have $t^r\ge t^2$. Therefore
\begin{equation*}
E(x)^\varsigma=e^{2\varsigma}e^{-\varsigma t^r}\le e^{2\varsigma}e^{-\varsigma t^2}
=e^{2\varsigma}(x-c_5)^{-\varsigma\log(x-c_5)},
\end{equation*}
which is \eqref{4.22:eq-coupling-ceilings-7}. Its right member is eventually below every power of
$x$. The ratio $\bar\alpha/\alpha_*$ is bounded by a power of $\log x$. Consequently, in
\eqref{4.22:eq-coupling-ceilings-d}, \eqref{4.22:eq-coupling-ceilings-e},
\eqref{4.22:eq-coupling-ceilings-6} and \eqref{4.22:eq-coupling-ceilings-m}, and in
\eqref{4.22:eq-coupling-ceilings-j} and \eqref{4.22:eq-coupling-ceilings-k} after both members have
been multiplied by $(x-c_5)^2$, the power
$E(x)^{\delta_\natural}$ or $E(x)^{\delta_\natural/2}$ beats the factor $(x-c_5)^2$, the
polylogarithmic factors (among them $\bar\alpha/\alpha_*$), the factors $\bar\alpha(x)$ and
$\alpha_*(x)$ that stand in numerators, which tend to zero, and the constants chosen before
$\gamma$. Each of these left members, so normalised, therefore tends to zero as $x\to\infty$, while
the right members of the normalised \eqref{4.22:eq-coupling-ceilings-j} and
\eqref{4.22:eq-coupling-ceilings-k} are the constants $\varepsilon_\star$ and $1$; in particular this
holds for \eqref{4.22:eq-coupling-ceilings-e}.

\emph{The two regularity members.} The second member of \eqref{4.22:eq-coupling-ceilings-a} implies
the first, because $\bar\alpha\ge\alpha_0$ as functions of $x$. Indeed,
\begin{equation*}
\bar\alpha(x_k)=\alpha^+_k=\max\{\tilde\alpha_{0,k},\tilde\alpha_{1,k},\delta'_k\}
\ge\tilde\alpha_{0,k}\ge\alpha_{0,k}.
\end{equation*}
The second member gives
\begin{equation*}
2BCM\bar\alpha(x)\le\frac{2}{2.1(2C_\chi+1)}=\frac{20/21}{2C_\chi+1},
\end{equation*}
which is the form in which Lemma~\ref{4.22:lem-frame-lemma}\,(i) uses it. Since
$\mathfrak{c}\ge20\bar r_6$, \eqref{4.22:eq-coupling-ceilings-b} gives
$20\bar r_6\bar\alpha(x)\le1$. Since moreover
\begin{equation*}
\mathfrak{c}\ge\mathfrak{c}_6=\mathfrak{b}_1+\mathfrak{b}_2(\bar p_1+\varpi_{\mathrm{fr}})\ge
\mathfrak{b}_2\bar p_1=200\,\bar p_1,
\end{equation*}
it also gives $\bar p_1\bar\alpha(x)\le1/200$.

\emph{Consequences of \eqref{4.22:eq-coupling-ceilings-g}.} Let $R_{\min}$ be the smallest of the
scales $\check R_k$ of a run. By \eqref{4.22:eq-coupling-ceilings-2} and $x_k>\gamma_H^{-2}$,
\begin{equation*}
\check R_k\ge(\log(x_k-c_5))^r\ge(\log(\gamma_H^{-2}-c_5))^r,
\end{equation*}
so that $R_{\min}\ge(\log(\gamma_H^{-2}-c_5))^r$, and \eqref{4.22:eq-coupling-ceilings-g} gives
$5\delta_\natural R_{\min}\ge2$. Hence $e^{-5\delta_\natural R_{\min}}\le e^{-2}$ per shell, as used
in Lemma~\ref{4.22:lem-law-site}\,(iii), and $\frac12\delta_\natural\cdot5R_{\min}\ge1$, as used in
Corollary~\ref{4.22:cor-older-law-sites}.

\emph{Consequence of \eqref{4.22:eq-coupling-ceilings-i}: $2N_0\le N_k$ at every step.} Let $N_0$
be the number of the levels $k_0+1,\dots,k$, with $k_0:=k-N_0$, that Ba{\l}aban's construction
\cite[\S1]{B-CMP122a} separates by single layers of cubes of decreasing size; in that construction
$L^{N_0-1}=\check R_{k_0+1}$. By
\eqref{4.22:eq-coupling-ceilings-2} at level $k_0+1$,
\begin{equation*}
L^{N_0-1}=\check R_{k_0+1}<L(\log x_{k_0+1})^r .
\end{equation*}
Taking $\log_L$ gives $N_0-1<1+r\log_L\log x_{k_0+1}$, that is,
\begin{equation*}
N_0<2+r\log_L\log x_{k_0+1}.
\end{equation*}
By the upper bound of the coupling flow (Theorem~\ref{4.1:thm-clause-zero}),
$x_{k_0+1}\le x_k+B_\sharp N_0$.

Suppose first that $N_0>x_k$. Then $N_0>\gamma_H^{-2}$, so \eqref{4.22:eq-coupling-ceilings-i} and
\eqref{4.22:eq-coupling-ceilings-h} may be applied at $x:=N_0$. Moreover
$x_{k_0+1}\le x_k+B_\sharp N_0<(1+B_\sharp)N_0$. Since $\log_L\log$ is increasing, we obtain the
chain
\begin{equation*}
N_0<2+r\log_L\log\bigl((1+B_\sharp)N_0\bigr)\le\tfrac12\bigl(\log(N_0-c_5)\bigr)^r
\le\tfrac12(\log N_0)^r\le\frac{N_0}{2L}.
\end{equation*}
The second inequality is \eqref{4.22:eq-coupling-ceilings-i} at $x=N_0$, divided by two. The last is
the second member of \eqref{4.22:eq-coupling-ceilings-h} at $x=N_0$. The chain is absurd, so
$N_0\le x_k$.

Hence $x_{k_0+1}\le x_k+B_\sharp N_0\le(1+B_\sharp)x_k$, and
$N_0<2+r\log_L\log((1+B_\sharp)x_k)$. Combining with \eqref{4.22:eq-coupling-ceilings-i} at $x=x_k$
and with \eqref{4.22:eq-coupling-ceilings-2},
\begin{equation*}
2N_0<4+2r\log_L\log\bigl((1+B_\sharp)x_k\bigr)\le\bigl(\log(x_k-c_5)\bigr)^r\le\check R_k=N_k .
\end{equation*}

Finally, the left member of \eqref{4.22:eq-coupling-ceilings-i} is decreasing in $L$ at fixed
$B_\sharp$. Indeed, $\log(x-c_5)\ge1$ by \eqref{4.22:eq-coupling-ceilings-h} gives $x>e$, so that
$\log\log((1+B_\sharp)x)>0$, and $\log_L=\log/\log L$. Thus
\eqref{4.22:eq-coupling-ceilings-i} imposes no upper bound on $L$ and is not a cap.

\emph{The ceiling \eqref{4.22:eq-coupling-ceilings-m}.} For $x\ge\gamma_H^{-2}$, the first member of
\eqref{4.22:eq-coupling-ceilings-d} gives $E(x)^{\delta_\natural}\le(x-c_5)^{-4}$. Hence
\begin{equation*}
\tfrac{4e^2}{3}B_H^SE(x)^{\delta_\natural}\le\tfrac{4e^2}{3}B_H^S(x-c_5)^{-4},
\end{equation*}
and the right member is at most $1$ as soon as $(x-c_5)^4\ge\frac{4e^2}3B_H^S$.
\end{proof}

\begin{remark}[Dependence of the constants on $N$]\label{4.22:rem-n-dependence}
Throughout this chapter $G=SU(N)$ with $N\ge 2$ fixed. This $N$ is unrelated to the stage count
$N_k$ of the glossary (Definition~\ref{0.2:not-glossary}).

Every step of the lemmas of this construction, from the polar-decomposition lemma in
$SL(N,\mathbb{C})$ through the frame lemma and its sub-lemmas to the last lemma of the
construction, is independent of $N$ except for finitely many numbers. They are fixed numbers,
entering the order of choice of the parameters before the coupling ceiling $\gamma$:
\begin{itemize}
\item the dimension $\dim\mathfrak{g}=N^2-1$;
\item the constants $G_1(N)$, $C_{22}(N)$ and $C_{22}^{P}$ of the lemma and corollaries in which
they are fixed; the dimension of $\mathfrak{g}$ and the norm factor of a rank-one term enter
$G_1(N)$ and $C_{22}(N)$;
\item a factor $\sqrt{N}$ inside the frame constant $C_{\mathrm{fr}}$, present only if matrices
are measured in the norm $(\operatorname{Tr}M^{*}M)^{1/2}$, where
$\operatorname{Tr}=N\operatorname{tr}$ is the unnormalised trace, rather than in the operator norm.
\end{itemize}
\end{remark}

\begin{proof}
The arguments of these lemmas use the following facts: the polar decomposition in
$G^{\mathbb{C}}=SL(N,\mathbb{C})$; the connectedness of $SU(N)$; Haar averages over
$\operatorname{Ad}_{SU(N)}$; the operator norm $\|\cdot\|$ on matrices acting on $\mathbb{C}^N$, as
in \cite[(19)]{B-CMP98}; the inclusion $U(x,\mu)\in G\subset U(N)$ for every bond variable; and the
identity $\det e^{A}=e^{\operatorname{Tr}A}$. The polar decomposition and the connectedness are
used in the polar-decomposition lemma; the last three facts, and the three norm relations below,
are used in the frame lemma, its four sub-lemmas and the product-logarithm lemma.

Together with these facts, three norm relations are used.
First, if $g\in U(N)$, then $\|\operatorname{Ad}_g M\|=\|gMg^{-1}\|=\|M\|$. This holds because
multiplication by a unitary matrix, on either side, preserves the operator norm.
Second, for $A\in\mathfrak{g}$ and real $\eta$ the matrix $e^{i\eta A}$ is unitary. Indeed, $A$ is
Hermitian, so $(e^{i\eta A})^{*}=e^{-i\eta A}=(e^{i\eta A})^{-1}$.
Third, for all matrices $A$ and $B$,
\begin{equation*}
\|[A,B]\|\le\|AB\|+\|BA\|\le 2\|A\|\|B\|.
\end{equation*}

None of these statements carries a constant depending on $N$. Hence $N$ enters the estimates of
these lemmas only through $\dim\mathfrak{g}=N^2-1$, through $G_1(N)$ and $C_{22}(N)$, into which
this dimension enters, through $C_{22}^{P}$, and, for the norm of the unnormalised trace, through
the factor $\sqrt{N}$ in $C_{\mathrm{fr}}$. Indeed, if the operator norm is replaced by
$(\operatorname{Tr}M^{*}M)^{1/2}$, the comparison
$(\operatorname{Tr}M^{*}M)^{1/2}\le\sqrt{N}\,\|M\|$ introduces a factor $\sqrt{N}$, which is
absorbed into $C_{\mathrm{fr}}$. All of these numbers are fixed before $\gamma$ in the order of
choice.
\end{proof}

\begin{definition}[The source tube and the oscillation seminorm]\label{4.23:not-oscillation-seminorm}
The following conventions are in force throughout this section.
\begin{enumerate}
\item[(a)] \emph{Level and site.} We put $k:=k_W$. We work at the site of the large-field operation
$\mathbf{R}_k$ fixed in the ordering of the operations of this chapter, before the classes are split.
In \cite{B-CMP122b} the classes are formed only after \cite[(1.70)]{B-CMP122b}, and a component $C$
becomes one of the regions $X_i$ or $Y_i$ only afterwards. We write $N:=N(g_k)$ for the window (a
number of levels), $h:=k-N$ for its bottom level, and $k_0:=k-N_0(g_k)$, where the function $N_0$ is
the one fixed together with the window in this chapter. In this section the unadorned letter $N$
denotes the window only; the rank of the gauge group does not occur, and $N_\ell(\cdot)$ in item (b)
is the cube-counting function of the geometric conventions. Further, $\eta:=L^{-k}$ and
$\ell_j:=L^j\eta$, so that $\ell_k=1$; $x:=g^{-2}$ and $x_k:=g_k^{-2}$; $\mathsf{T}_k:=\log x_k$, and
$\check{\mathsf{T}}_k$ is the variant of $\mathsf{T}_k$ fixed in the conventions of this chapter. In
counts we use $d=4$, and $\check{R}_k$ denotes the window radius; where no confusion can arise we
write $R$ for $\check{R}_k$, and the ball radius $R$ of the glossary does not occur in this section.
\item[(b)] \emph{Geometry.} $Z$ denotes the class treated at this site, and $C$ any component of $Z$.
The geometric notation is that of this chapter: $\pi_\ell$, $\operatorname{dist}_\ell$, the relations
``touch'' and ``meet'', $[S]_i$, $N_\ell$, $\mathbf{D}_\ell$, $d_\ell$, $d_{k,Z}$; the cells, their
levels and the zones of $\mathfrak{B}:=\mathbf{B}''_k$; $\Lambda$, $\Lambda_C:=\Lambda\cap C$ and
$\Lambda_C^{\sim}$; the regions $\Lambda_k$, $\Omega_k$ of the fixed level $k$; and the
set $D_2$ fixed with the geometric conventions. The \emph{core} of $C$ is
\begin{equation*}
\mathsf{C}:=\mathsf{C}_C:=\bigl(\Omega''^{\sim 2}_{h+1}\bigr)^c\cap C .
\end{equation*}
$C^{\sharp}$ is the set $C$ enlarged by
$m_\sharp(k):=\lfloor \tfrac12(\check{R}_k-1)\rfloor$ layers of $\pi_k$-cubes. It belongs to
$\mathbf{D}_k^Z$, the class of domains $Y_4\in\mathbf{D}_k$ for which $Y_4\cap Z$ is a non-empty union of
components of $Z$. We write $d^{\wedge}:=d^{\wedge}_{\mathfrak{B}}$ for the own-scale contour distance
of the conventions of this chapter, $\delta_P:=\tfrac18\delta_0$, where
$\delta_0$ is the constant so denoted in the conventions of this chapter, and, for
a cell $\Delta$,
\begin{equation}\label{4.23:eq-oscillation-seminorm-1}
\varrho(\Delta):=1+\frac{M_1}{M}\,d^{\wedge}(\Delta,\mathsf{C}^c)\quad\text{if }\Delta\subset\mathsf{C},
\qquad \varrho(\Delta):=1\quad\text{otherwise}.
\end{equation}
In this section a bare $\varrho$ denotes a depth variable, $\varrho\ge 1$, standing for a value of
$\varrho(\cdot)$; the auxiliary parameter $\varrho$ of the glossary (Notation~\ref{0.2:not-glossary})
is never written bare in this section. The letters of the interpolation chain are
those fixed with the chain in this chapter. For a point $\mathsf{m}$ of the presented polydisc
$\mathfrak{P}$ of item (g) below, $\operatorname{chain}(\mathsf{m})$ is the carrier pair of the chain at
its layer $\mathfrak{b}_4$ at the present site, and $\mathcal{U}(\mathsf{m})$ is its configuration
component.
\item[(c)] \emph{Expansion at the site.} The step is run under the standing clause of this chapter that
governs the expansion at the site. At the
present site every interior unit $v$ of the first part of the expansion at this site
(interior meaning $X_v\subset Z$), plain or wrapped, is expanded by the two small-field expansion
rules of this chapter. Neither a partition-of-unity factor $\zeta_{\square}$ of \cite{B-CMP119} nor the
interpolation parameter $\mathsf{t}$ enters this expansion; the interpolation variables
$\vec{\mathsf{s}}$ of the chain (item (g)) are a different object and are not affected. The carrier
pairs are read as in this chapter, and the pointed
description of the initial data is used together with its compatibility with the restriction of
carriers.
\item[(d)] \emph{The real variables $\varpi$ and the norm.} We write $\varpi$ for the real variables
of \cite[(1.71)]{B-CMP122b}; they range over the set $\operatorname{supp}\Pi$, the support of the
function $\Pi$ fixed with them in this chapter, and we write $\varpi\in\operatorname{supp}\Pi$.
These are the variable $B$ on $\mathbf{B}_0\cap C$,
subject to $|B|<\mathfrak{b}'_k$, where
\begin{equation}\label{4.23:eq-oscillation-seminorm-2}
\mathfrak{b}'_k:=\frac{\delta'_k}{g_k}=A_1\,p_1(g_k),\qquad p_1(g):=(\log x)^{p_1}
\end{equation}
(the same letter denoting the function and its fixed exponent $p_1$, as for the degree function
$p_0(\cdot)$; the quantity $\delta'_k$ is that of \cite{B-CMP122a}), together with the $G$-valued
variable $V''$
restricted to $\mathsf{C}$, where $G$ denotes the gauge group, as in
Definition~\ref{4.4:def-regular-backgrounds}. On $\mathfrak{g}^{\mathbb{C}}$, the complexification of
the Lie algebra $\mathfrak{g}$ of $G$, we fix an $\operatorname{Ad}_G$-invariant
norm $|\cdot|$ with $|[X,Y]|\le 2|X||Y|$, and we put
\begin{equation*}
c_{\mathrm{tr}}:=\sup_{X,Y\neq 0}\frac{|\operatorname{tr}XY|}{|X||Y|}.
\end{equation*}
\item[(e)] \emph{The sets $\mathfrak{L}_C$, $\mathfrak{B}^0$, $\mathfrak{B}^C$.} With $\sim$ taken for
cubes of the fine block size $\mathfrak{M}$ fixed in this chapter,
the inner rim of $C$ is
$\mathfrak{L}_C:=C\setminus C^{\sim-2}$. Next, $\mathfrak{B}^0:=\mathbb{B}_k(C)$ is the minimal
determining set of \cite{B-CMP119} with support $C$. Finally, following \cite[(2.14)]{B-CMP119},
\begin{equation}\label{4.23:eq-oscillation-seminorm-3}
\mathfrak{B}^C:=\mathbf{B}''_k(C):=\bigl(\mathbf{B}''_k\cap C^{\sim-2}\bigr)\cup
\bigl(\mathbb{B}_k(C)\cap\mathfrak{L}_C\bigr).
\end{equation}
\item[(f)] \emph{The source tube.} The source tube $\mathfrak{S}_{\mathrm{src}}(Y_4)=
\tilde{U}^{c}_{k}(Y_4,\tilde\alpha_0,\tilde\alpha_1)$ is defined separately for each domain $Y_4$ on
which a term is filed. Written without argument, $\mathfrak{S}_{\mathrm{src}}$ denotes the source tube
over the filing domain of the term under consideration, whichever that domain is; in particular
$\mathfrak{S}_{\mathrm{src}}=\mathfrak{S}$ when that filing domain is $C^{\sharp}$. We read it at
$Y_4:=C^{\sharp}$:
\begin{equation}\label{4.23:eq-oscillation-seminorm-4}
\mathfrak{S}:=\mathfrak{S}_{\mathrm{src}}(C^{\sharp}):=
\tilde{U}^{c}_{k}(C^{\sharp},\tilde\alpha_0,\tilde\alpha_1).
\end{equation}
In this section $\mathfrak{S}$ denotes this tube only, not the class of admissible background pairs of
Definition~\ref{4.5:def-admissible-pairs}.
It is the set of pairs $\zeta=(\mathbf{U},\mathbf{J})$ on $C^{\sharp}$ satisfying
\cite[(2.34)--(2.39)]{B-CMP119} at the new set $\mathbf{B}_k$, which is of level $k$ on all of
$C^{\sharp}$. Restriction of configurations maps $\mathfrak{S}$ into
$\tilde{U}^c_k(Y_4,\tilde\alpha_0,\tilde\alpha_1)$ for every $Y_4\subset C^{\sharp}$.
The estimates of this chapter on tube points are stated for every point of the source tube over every
filing domain; on $\mathfrak{S}$ they are used by restriction to the filing domain $C^{\sharp}$ only.
The unadorned letter $\zeta$ always denotes a point of the tube; Ba{\l}aban's cut-off
and partition-of-unity functions (among them $\zeta_0$, $\zeta_1$, $\zeta_{\square}$) are never
denoted by an unadorned $\zeta$ in this section, and neither the complex-source letter $\zeta$ nor the
running shifts $\zeta_k$, $\zeta'_k$ of the glossary (Notation~\ref{0.2:not-glossary}) are meant by an
unadorned $\zeta$ here.
\item[(g)] \emph{The slice of the presented polydisc.} Let $\vec{\mathsf{s}}=(\mathsf{s}^{(t)}(\Delta))$
be the interpolation variables of the families of the chain, and let $\sigma_0^{(t)}$ be the set of
cells of the $t$-th family. The presented polydisc is
\begin{equation*}
\mathfrak{P}:=\mathfrak{S}_{\mathrm{src}}\times\operatorname{supp}\Pi\times\prod_{t}
\bigl\{|\mathsf{s}^{(t)}(\Delta)|\le e^{\kappa_1}:\ \Delta\in\sigma_0^{(t)}\bigr\}.
\end{equation*}
Its slice over $C^{\sharp}$ is
\begin{equation}\label{4.23:eq-oscillation-seminorm-5}
\mathfrak{P}^{\sharp}:=\mathfrak{S}\times\operatorname{supp}\Pi\times\prod_{t}
\Bigl\{|\mathsf{s}^{(t)}(\Delta)|\le e^{\kappa_1}:\ \Delta\in\sigma_0^{(t)},\ \Delta\subset C^{\sharp}
\Bigr\}.
\end{equation}
It is regarded as a subset of $\mathfrak{P}$ by putting $\mathsf{s}^{(t)}(\Delta):=0$ for
$\Delta\not\subset C^{\sharp}$. On $\mathfrak{P}^{\sharp}$ every layer of the chain, restricted to
$C^{\sharp}$, depends on the variables $(\vec{\mathsf{s}},\varpi,\zeta)$ on $C^{\sharp}$ only.
Every piece filed in $C^{\sharp}$ is evaluated on
$\mathfrak{P}^{\sharp}$, and every input stated ``for every point of $\mathfrak{P}$'' is read on
$\mathfrak{P}^{\sharp}$.
\item[(h)] \emph{Admissible splittings.} An \emph{admissible splitting} is a pair $(\zeta,U)$ with
$\mathbf{U}=U'U$, where the following hold.
\begin{itemize}
\item $U$ is $G$-valued and satisfies the conditions among \cite[(2.34)--(2.39)]{B-CMP119} that are
imposed on the group-valued factor of a tube point.
\item $U'=e^{i\eta A'}$ with $A'$ a $\mathfrak{g}^{\mathbb{C}}$-valued field satisfying the third
clause of the tube, which is the smallness condition $\ell_k|A'|<\tilde\alpha_{1,k}$ on $C$, where
$\tilde\alpha_{1,k}$ is the radius $\tilde\alpha_1$ of the tube read at level $k$; $A'$ may
be real.
\item $(U,0)\in\mathfrak{S}_{\mathrm{src}}$.
\end{itemize}
We let $\mathfrak{V}_{\mathbb{R}}$ be the set of the restrictions to $C\setminus\Lambda$ of the
$k$-fold block averages (Definition~\ref{1.5:def-block-average}) of all such $U$. By the
linearisation statement for the block-averaging map of this chapter, the $k$-fold block average of
$\mathbf{U}$ equals $e^{iB'}$ times the $k$-fold block average of $U$, with, for the constant $c_Q$ of
that statement,
\begin{equation}\label{4.23:eq-oscillation-seminorm-6}
|B'|_\infty\le c_Q\sup\ell_k|A'|<c_Q\,\tilde\alpha_{1,k}
\end{equation}
on the whole of $\mathfrak{S}_{\mathrm{src}}$, and $\tilde\alpha_{1,k}\le\bar\alpha_1(x_k)$, where
$\bar\alpha_1(\cdot)$ is the radius ceiling fixed with the tubes in this chapter.
\item[(i)] \emph{The real variables $\Phi_C$.} $\Phi_C$ denotes the collection of all real variables
on which the localised expression $V^{\sharp}(C^{\sharp})$, the sum of the terms filed in the domains
$Y_4\subset C^{\sharp}$, depends besides $\zeta$. These are:
\begin{itemize}
\item $\varpi$;
\item the variables $(\mathbf{A},\boldsymbol{\Phi})$ in the live slots of the boundary-type units
with domain in $C$;
\item the data $\Phi_{Y'}$ attached to the nested regions $(Y',\mathsf{a}')$ with $Y'\subset C$;
\item the collar variables of the arrested regions nested in $C$.
\end{itemize}
Each of these ranges over its support. The complex source $w$, ranging over the space
$\mathcal{W}_k(C^{\sharp})$ of complex sources over $C^{\sharp}$, and the weak coordinates (the further
holomorphic parameters of the response proposition) enter as
holomorphic parameters, under majorants that do not depend on $w$, by clauses (w1)--(w3) of part (a)
of the response proposition of this chapter.
\item[(j)] \emph{The oscillation.} For a complex-valued function $f$ on
$\mathfrak{S}\times\{\Phi_C\}\times\mathcal{W}_k$ we put
\begin{equation}\label{4.23:eq-oscillation-seminorm-7}
\operatorname{osc}(f):=\sup_{\Phi_C\ \mathrm{real},\ w}\ \sup_{\zeta,\zeta'\in\mathfrak{S}}
\bigl|f(\zeta;\Phi_C,w)-f(\zeta';\Phi_C,w)\bigr| .
\end{equation}
Thus two arbitrary points of the tube are compared at one and the same $\Phi_C$ and $w$; two
different values of $\Phi_C$ are never compared. The functional $\operatorname{osc}$ is a seminorm;
$\operatorname{osc}(f+f_0)=\operatorname{osc}(f)$ whenever $f_0$ does not depend on $\zeta$; and
$\operatorname{osc}(f)\le 2\sup|f|$, where $\sup|f|$ is the supremum of $|f(\zeta;\Phi_C,w)|$ over
$\zeta\in\mathfrak{S}$, real $\Phi_C$ and $w$. It is the quantity $\operatorname{osc}(f)$ that the
oscillation bound of this chapter uses. Each statement of this chapter whose conclusion is used here
is formulated ``for all
$(\mathbf{A},
\boldsymbol{\Phi})$ and $w$'' and ``for every real value of the variables of
\cite[(1.71)]{B-CMP122b}'', written $\Phi$ in those statements and $\varpi\in\operatorname{supp}\Pi$
here, so the
uniformity in these variables is provided by the inputs themselves.
\item[(k)] \emph{Format data.}
\begin{itemize}
\item $\mathfrak{N}^{A,\mathbb{C}}(v)$ is the format datum of a stored boundary-type unit, or of a
wrapped unit, $v$ on the complex presented polydisc. It sums over the strong sets (the
sets so designated in the format of this chapter) and includes the supremum over the family
$\mathfrak{G}_{\varepsilon'}(X_v)$ attached to $X_v$. It is identified with
the datum
$\mathfrak{N}^{A}(v)$ of the same unit as it enters the bound of this chapter on deep units and two
further bounds of this chapter in which that datum appears; the
superscript $\mathbb{C}$ only records that the tube is complex.
\item $\mathfrak{N}^{\flat}(v)$ is the pointed datum, with $\mathfrak{N}^{\flat}:=\mathfrak{N}$ on
the set $\mathfrak{L}^B$ of this chapter, where $\mathfrak{N}$ denotes the format datum.
\item The per-cube datum $\mathfrak{N}^{G}(e;\mathfrak{O})$ inside the bound $b^{G}_{\delta}$ of the
Gaussian units (here $\mathfrak{O}$ is an oscillation datum of the arrested region at the cube $e$)
and the datum $\mathfrak{N}^{G,L}(e)$ are identified; their common value, the \emph{oscillation
datum} $\mathfrak{N}^G$, is what enters the total bound of this chapter.
\item For a unit $v_e$ that is not deep, the datum is $2\nu^{\natural}(e)$.
\item In $\mathfrak{N}^{G}(e;\mathfrak{O})$ the datum $\mathfrak{O}$ is read as the arrested region's
own oscillation datum, involving neither the bound $\mathfrak{O}\le C_\Sigma\check{R}_k^{\,p_\Sigma}$
satisfied by the datum of a nested region, with its constants $C_\Sigma$ and $p_\Sigma$, nor the
kept-term constant $\mathfrak{n}^{\mathrm{kept}}$. This reading is not derived in this section; it
enters as a hypothesis of the lemma of this section that fixes the order of the constants.
\end{itemize}
\end{enumerate}
\end{definition}

\begin{proof}
We justify the assertions contained in items (b), (f), (g), (h) and (j), in that order.

\emph{Item (b).} That $C^{\sharp}$ belongs to $\mathbf{D}_k^Z$ is part of the lemma of this chapter on
the enlarged component $C^{\sharp}$. That lemma also gives
$C^{\sharp}\setminus C\subset\Lambda_k$. It further gives that $C^{\sharp}$ holds exactly one
component of $\Lambda_k^c$, and that the components of $\Lambda_k^c$ other than the one contained in
$C^{\sharp}$ lie at distance at least
$\check{R}_k-m_\sharp(k)>0$ from $C^{\sharp}$.

\emph{Item (f).}

(1) \emph{Level of $\mathbf{B}_k$.} In \cite[(1.33)]{B-CMP122b} the set $\mathbf{B}_k$ is taken equal
to $\mathbf{B}''_k$ on $Z^c$ and to $\{Z^{(k)}\}$ on $Z$; in particular, on $C$, which is contained in
$Z$, it is $\{Z^{(k)}\}$. Moreover $C^{\sharp}\setminus C\subset
\Lambda_k$ by item (b), and $\Lambda_k\subset\Omega_k$ by the construction of the regions $\Lambda_k$
and $\Omega_k$. Hence $\mathbf{B}_k$ is of level $k$ on all of $C^{\sharp}$, and
the conditions \cite[(2.34)--(2.39)]{B-CMP119} are imposed at this set.

(2) \emph{The case $C=Y_i$.} Suppose $C$ is one of the regions $Y_i$. Then
\eqref{4.23:eq-oscillation-seminorm-4} reads
$\mathfrak{S}=\tilde{U}^c_{k}(Y_i^{\sharp},\tilde\alpha_0,\tilde\alpha_1)$ (recall $k=k_W$). This is
exactly the tube
$\mathfrak{U}^{\mathrm{tr}}_v$ on which the pieces of the corresponding unit $v$ are presented, so the
two tubes are equal. The hypothesis of this chapter on the presented pieces of the unit $v$ is
formulated over a set containing
$\tilde{U}^c_{k}(Y_i^{\sim},\tilde\alpha_0,\tilde\alpha_1)$; it is applied on the single tube
$\mathfrak{S}$, which is a subset of the set over which the hypothesis is formulated. Since a supremum
over a subset does not exceed the supremum over the whole set, every bound of the hypothesis remains
valid on $\mathfrak{S}$.

(3) \emph{Restriction.} The estimates of this chapter on tube points are stated for every $\zeta$ in
the source tube over every filing domain. In particular they hold for every $\zeta$ in the source
tube over $C^{\sharp}$, that is, on $\mathfrak{S}$, by restriction to this filing domain, and only
this restriction is used in what follows.

(4) \emph{Sub-domains.} By the restriction property of the tubes, restriction of configurations maps
$\mathfrak{S}$ into the tube $\tilde{U}^c_k(Y_4,\tilde\alpha_0,\tilde\alpha_1)$ over every
$Y_4\subset C^{\sharp}$.

\emph{Item (g).}

(1) \emph{$\mathfrak{P}^{\sharp}\subset\mathfrak{P}$.} Setting $\mathsf{s}^{(t)}(\Delta):=0$ for
$\Delta\not\subset C^{\sharp}$ satisfies $|\mathsf{s}^{(t)}(\Delta)|\le e^{\kappa_1}$. Together with
$\mathfrak{S}=\mathfrak{S}_{\mathrm{src}}(C^{\sharp})$, the source tube over the filing domain
$C^{\sharp}$ of the pieces evaluated on $\mathfrak{P}^{\sharp}$, this places $\mathfrak{P}^{\sharp}$
inside $\mathfrak{P}$.

(2) \emph{Locality on the slice.} We claim: on $\mathfrak{P}^{\sharp}$, every layer of the chain,
restricted to $C^{\sharp}$, depends on the variables $(\vec{\mathsf{s}},\varpi,\zeta)$ on $C^{\sharp}$
only. This follows from three facts. First, by the locality clause of the lemma on the index sets
$\sigma_0^{(t)}$, for every term of a layer of the chain, with index $\sigma$ and support
$Y(\sigma)$, one has $Y(\sigma)\subset\Lambda^{\sim}\cup C^{\sharp}$, where $\Lambda^{\sim}$ is the
enlargement of
$\Lambda$ in the sense of the geometric conventions of item (b). Second, by item (b), the components
of $\Lambda_k^c$ other than the one contained in $C^{\sharp}$ lie at distance at
least $\check{R}_k-m_\sharp(k)>0$ from $C^{\sharp}$. Third, every cell $\Delta$ between $C^{\sharp}$
and those components is switched off on $\mathfrak{P}^{\sharp}$, that is, $\mathsf{s}^{(t)}(\Delta)=0$
there. Combining the three facts: the support of every term of a layer lies in
$\Lambda^{\sim}\cup C^{\sharp}$; the components of $\Lambda_k^c$ outside $C^{\sharp}$ are separated
from $C^{\sharp}$ by a positive distance; and every cell in between carries the value zero of its
interpolation variable on $\mathfrak{P}^{\sharp}$. Hence no variable outside $C^{\sharp}$ enters a
layer restricted to $C^{\sharp}$, which is the claim. It is for
this reason that pieces filed in $C^{\sharp}$, and inputs stated on $\mathfrak{P}$, may be evaluated
on $\mathfrak{P}^{\sharp}$.

\emph{Item (h).} Here $\mathbf{U}=U'U$ with $U'=e^{i\eta A'}$. The linearisation statement for the
block-averaging map of this chapter, valid on the whole of
$\mathfrak{S}_{\mathrm{src}}$, gives that the $k$-fold block average of $\mathbf{U}$ equals $e^{iB'}$
times the $k$-fold block average of $U$, with the constant $c_Q$ of that statement,
$|B'|_\infty\le c_Q\sup\ell_k|A'|$. The smallness
condition $\ell_k|A'|<\tilde\alpha_{1,k}$ then gives the
strict inequality in \eqref{4.23:eq-oscillation-seminorm-6}. The bound
$\tilde\alpha_{1,k}\le\bar\alpha_1(x_k)$ is the bound on the radii fixed with the tubes in this
chapter.

\emph{Item (j).} Throughout, $\operatorname{osc}$ is considered on the functions $f$ for which
\eqref{4.23:eq-oscillation-seminorm-7} is finite; by step (3) below, this includes all bounded $f$.

(1) \emph{$\operatorname{osc}$ is a seminorm.} Clearly $\operatorname{osc}(f)\ge 0$. Let $f_1,f_2$ be
two such functions and $c\in\mathbb{C}$. For fixed real $\Phi_C$, fixed $w$ and
$\zeta,\zeta'\in\mathfrak{S}$, write $f_j(\zeta)$ for $f_j(\zeta;\Phi_C,w)$. Then
\begin{align*}
\bigl|(f_1+f_2)(\zeta)-(f_1+f_2)(\zeta')\bigr|
&\le|f_1(\zeta)-f_1(\zeta')|+|f_2(\zeta)-f_2(\zeta')|\\
&\le\operatorname{osc}(f_1)+\operatorname{osc}(f_2).
\end{align*}
Taking the suprema over $\zeta,\zeta'$, over real $\Phi_C$ and over $w$ gives the triangle inequality
$\operatorname{osc}(f_1+f_2)\le\operatorname{osc}(f_1)+\operatorname{osc}(f_2)$. Also
$|cf_1(\zeta)-cf_1(\zeta')|=|c|\,|f_1(\zeta)-f_1(\zeta')|$; taking the same suprema gives
$\operatorname{osc}(cf_1)=|c|\operatorname{osc}(f_1)$.

(2) \emph{$\operatorname{osc}$ annihilates terms independent of $\zeta$.} Suppose $f_0$ does not
depend on $\zeta$. Then $f_0(\zeta;\Phi_C,w)=f_0(\zeta';\Phi_C,w)$ for all $\zeta,\zeta'$, so
$\operatorname{osc}(f_0)=0$. By step (1), for every $f$,
\begin{equation*}
\operatorname{osc}(f+f_0)\le\operatorname{osc}(f)+\operatorname{osc}(f_0)=\operatorname{osc}(f).
\end{equation*}
Applying the same inequality to $f+f_0$ and $-f_0$ gives
$\operatorname{osc}(f)\le\operatorname{osc}(f+f_0)$. Hence $\operatorname{osc}(f+f_0)=
\operatorname{osc}(f)$, so terms independent of $\zeta$ may be dropped. In particular
$\operatorname{osc}$ is not a norm.

(3) \emph{The crude bound.} With $\sup|f|$ as in item (j), for every pair $\zeta,\zeta'$,
\begin{equation*}
|f(\zeta;\Phi_C,w)-f(\zeta';\Phi_C,w)|\le|f(\zeta;\Phi_C,w)|+|f(\zeta';\Phi_C,w)|\le 2\sup|f|.
\end{equation*}
Taking suprema gives $\operatorname{osc}(f)\le 2\sup|f|$.

(4) \emph{Uniformity.} In \eqref{4.23:eq-oscillation-seminorm-7} the variables $\Phi_C$ and $w$ are
the same in both evaluations, and the supremum over them is taken outside. A bound on
$\operatorname{osc}(f)$ therefore requires, for each fixed real $\Phi_C$ and each $w$, only a bound on
the variation of $f$ in $\zeta$, uniform in $(\Phi_C,w)$. The inputs used with
$\operatorname{osc}$ are stated for all $(\mathbf{A},\boldsymbol{\Phi})$, all $w$ and every real value
of the variables of \cite[(1.71)]{B-CMP122b}, written $\Phi$ in those inputs and
$\varpi\in\operatorname{supp}\Pi$ here, and so provide exactly this uniformity.
\end{proof}

\begin{theorem}[The tube statement for a treated large-field component]\label{4.23:thm-tube-statement}
Assume the standing hypotheses \textup{(G1)}--\textup{(G26)} of the large-field analysis of this
chapter. Assume further the following statements, each in the form in which it is stated at its
own place:
\begin{itemize}
\item the seam identity;
\item the bound on the total constant of the large-field analysis, in the form in which the third
item of that bound uses hypothesis \textup{(G15)};
\item clauses \textup{(i)} and \textup{(ii)} of the lemma on the subtracted copies;
\item the closed form of the quantity $\mathsf{b}_k$ of the large-field analysis, and the cut
condition on the auxiliary parameters;
\item the inequality $s_{\mathrm{box}}(3)\ge 1$ for the function $s_{\mathrm{box}}$ of the
large-field analysis of this chapter;
\item the table of the large-field analysis of this chapter;
\item hypotheses \textup{(a)} and \textup{(b)} of the lemma on the order of the constants of the
large-field analysis.
\end{itemize}

Then for every $K$, every $k\le K$, every label of the large-field expansion and every component
$C$ of the large-field region
$Z$ at the operation $\mathbf{R}_k$, the tube statement at level $k$ holds for $C$. We use the
following notation.
\begin{itemize}
\item $C^{\sharp}$ is the component $C$ together with
$m_\sharp(k)=\lfloor(\check{R}_k-1)/2\rfloor$ layers of cubes of the level-$k$ partition
$\pi_k$ around it, where $\check{R}_k$ is the window radius function.
\item $\mathfrak{S}=\tilde{U}^{c}_{k}(C^{\sharp},\tilde\alpha_0,\tilde\alpha_1)$ is the tube
over $C^{\sharp}$. It is Ba{\l}aban's space of complex configurations over $C^{\sharp}$
\cite[(2.34)--(2.39)]{B-CMP119}, and $\tilde\alpha_0,\tilde\alpha_1$ are its radii.
\item $\mathfrak{S}_{\mathrm{src}}$ is the source tube: the same space
$\tilde{U}^{c}_{k}(\cdot,\tilde\alpha_0,\tilde\alpha_1)$ of complex configurations, taken over the
filing domain of $C$.
\item $\zeta=(\mathbf{U},\mathbf{J})$ is a point of the tube; in \textup{(F)} and \textup{(q0)} it
ranges over $\mathfrak{S}_{\mathrm{src}}$.
\item $\varpi$ denotes the real variables of \cite[(1.71)]{B-CMP122b}; it ranges over the set
$\operatorname{supp}\Pi$, and every clause below is asserted for every such $\varpi$.
\item $\mathfrak{q}_C$ is the second kept member, that is, the quadratic form of the fluctuation
field. It is the quantity of the definition of the second kept member, which precedes the form
theorem (bound, variation and sign of the second kept member); the splitting referred to below is
an admissible splitting as in the form theorem.
\item $\mathsf{T}^{\mathsf{L}}$ is the linearised localised minimiser map, and $K_{\mathsf{T}}$
denotes its norm bound, fixed in the form theorem; at $\mathsf{T}=\mathsf{T}^{\mathsf{L}}$ this
bound is $K_{\mathsf{T}}=B_\sharp$, the constant that enters clause \textup{(t)} of
\eqref{4.23:eq-oscillation-theorem-a}.
\item $c_{\mathrm{tr}}$ is the level-independent constant and $\mathsf{n}_{\mathrm{sh}}(k)$ the
level-dependent factor of the bound in clause \textup{(X1)} of \eqref{4.23:eq-form-theorem-a}.
\item For a gauge transformation $\lambda$ with values in the gauge group $G$ of
Definition~\ref{4.4:def-regular-backgrounds}, $\zeta^{\lambda}$ and $\varpi^{\lambda}$
are the transforms of $\zeta$ and $\varpi$ by $\lambda$ under the joint action of gauge
transformations on configurations, currents and these real variables that is fixed in the
large-field analysis of this chapter.
\end{itemize}
The tube statement at level $k$ is the conjunction of four clauses \textup{(F)}, \textup{(q0)},
\textup{(q1)} and \textup{(V)}. The
first three are
\begin{equation}\label{4.23:eq-tube-statement-a}
\begin{aligned}
&\text{\textup{(F)}}\quad \text{clauses \textup{(X2)} and \textup{(X3)} of
  \eqref{4.23:eq-form-theorem-a} hold for }
  \mathfrak{q}_C;\\
&\text{\textup{(q0)}}\quad \text{at } \mathsf{T}=\mathsf{T}^{\mathsf{L}}:\ \mathfrak{q}_C
  \text{ does not depend on the splitting,}\\
&\qquad\quad \text{is holomorphic on } \mathfrak{S}_{\mathrm{src}},\\
&\qquad\quad \text{is at real points the form of \cite[(1.71)]{B-CMP122b}, and}\\
&\qquad\quad \mathfrak{q}_C(\zeta^{\lambda};\varpi^{\lambda})=\mathfrak{q}_C(\zeta;\varpi)
  \quad\text{for every $G$-valued $\lambda$};\\
&\text{\textup{(q1)}}\quad \sup_{\mathfrak{S}}|\mathfrak{q}_C|\le \mathfrak{q}_{\max}(k)\\
&\qquad\quad\;\;
  =15\,c_{\mathrm{tr}}\,B_{\sharp}^{2}\,(\mathfrak{b}'_k)^{2}\,M^{d}\,\mathsf{n}_{\mathrm{sh}}(k).
\end{aligned}
\end{equation}
Clause \textup{(q0)} is clause \textup{(X0)} of \eqref{4.23:eq-form-theorem-a} at
$\mathsf{T}=\mathsf{T}^{\mathsf{L}}$ together with the joint $G$-invariance in its last line; of
the conclusions of \textup{(X0)} the display lists three: independence of the splitting,
holomorphy on $\mathfrak{S}_{\mathrm{src}}$, and the form at real points. Clause \textup{(q1)} is
clause \textup{(X1)} of
\eqref{4.23:eq-form-theorem-a} at
$\mathsf{T}=\mathsf{T}^{\mathsf{L}}$, restricted to the tube $\mathfrak{S}$. Clause \textup{(V)} is
the conjunction of the clauses $(\sharp)$, \textup{(v)}, \textup{(t)}, \textup{(b)}, \textup{(d)},
\textup{(g)}, \textup{(o)}, \textup{(n)} of
\eqref{4.23:eq-oscillation-theorem-a}.
\end{theorem}

By clause (X5) of \eqref{4.23:eq-form-theorem-a}, the bound $\mathfrak{q}_{\max}(k)$ in (q1)
depends on the level through the number of levels $N=N(g_k)$ of the window, the window radius
$\check{R}_{h+1}$ at the bottom level $h=k-N$, and $\mathfrak{b}'_k$.

The proof of Theorem~\ref{4.23:thm-tube-statement} combines the form theorem (bound, variation and
sign of the second kept member), together with its clause (X0), with the oscillation theorem for
the resummed density. It is given in the assembly of the tube statement, which follows the two
theorems in this chapter. There the hypotheses required by the two theorems are derived from the
hypotheses listed in Theorem~\ref{4.23:thm-tube-statement}; for the oscillation theorem they
include the hypotheses of the form theorem at every step $\le k$. The conclusions of the two
theorems then give the clauses (F), (q0), (q1) and (V). In clause (t) of
\eqref{4.23:eq-oscillation-theorem-a} the bound
$\sup_{\mathfrak{S}}|\mathfrak{q}_C|\le\mathfrak{q}_{\max}(k)$ enters as an input, supplied by
clause (q1), and not as a conclusion of the oscillation theorem.

\begin{definition}[Analytic solvability of the equation for the $\Lambda$-datum; printed without
proof in {\cite[p.~367, (1.40)--(1.41)]{B-CMP122b}}]
\label{4.23:def-lambda-datum}
Let $\Lambda$ be the region of Ba{\l}aban's equation \cite[(1.40)]{B-CMP122b} on which the solution
of that equation is a field, and let $C$ be the component of
Notation~\ref{4.23:not-oscillation-seminorm}. The solution of the equation, its
bound and its analyticity are taken from the following sentence, which is stated in
\cite[p.~367, (1.40)--(1.41)]{B-CMP122b} without proof; cited as printed:
\begin{quote}
We can prove again that Eq.~(1.40) has exactly one solution $B'_{\Lambda}(\tilde B')$, which is
an analytic function of the $\mathfrak{g}^c$-valued small field $\tilde B'$, satisfying the bound
\begin{equation*}
|B'_{\Lambda}(\tilde B')| \leq 4d(100M)^{5}\gamma_0^{-1}B_3^{2}|\tilde B'|. \tag*{(1.41)}
\end{equation*}
\end{quote}
The sentence states no radius: $\tilde B'$ is only required to be sufficiently small
\cite[p.~367]{B-CMP122b}.

The field denoted $\tilde B'$ in \cite{B-CMP122b} is written $B'$ here; the solution map keeps
its name $B'_{\Lambda}$, and its value, a field on $\Lambda$, is written $b_\Lambda$. Write
$\gamma_0^{\mathrm{pos}}$ for the positivity constant of \cite[(1.7), (1.9)]{B-CMP122b}; it is
this constant that appears as $\gamma_0$ in the quoted bound. It is neither the parameter
$\gamma_0$ bounding the current in the class $\mathfrak{K}_{109}(\mathfrak{B};\gamma_0)$ of
background pairs nor the constant $\gamma_0^{(k)}$ used elsewhere in this chapter. Set
\begin{equation}\label{4.23:eq-lambda-datum-1}
K_\Lambda := 4d(100M)^{5}(\gamma_0^{\mathrm{pos}})^{-1}B_3^{2},
\end{equation}
where $B_3$ is the constant so denoted in \cite[(1.41)]{B-CMP122b}, and $M$ and $d$ are as in
the glossary (Notation~\ref{0.2:not-glossary}).

The quoted sentence is used only for fields $B'$ with $|B'|_\infty < c_Q\tilde\alpha_{1,k}$.
This is the bound on the block-field twist $B'$ over the whole source tube
$\mathfrak{S}_{\mathrm{src}}$, defined domain by domain, of which the tube of
Notation~\ref{4.23:not-oscillation-seminorm} is the instance over $C^{\sharp}$.
Let $\mathfrak{V}_{\mathbb{R}}$ be the class of real block fields on $C\setminus\Lambda$ obtained
by restricting to $C\setminus\Lambda$ the level-$k$ block fields of the $G$-valued
configurations $U$, $G$ the gauge group, with $(U,0)\in\mathfrak{S}_{\mathrm{src}}$, and let
$V\in\mathfrak{V}_{\mathbb{R}}$ be arbitrary. Let $B'$ be any $\mathfrak{g}^c$-valued field
with $|B'|_\infty < c_Q\tilde\alpha_{1,k}$. The sentence is then taken to assert the following
three clauses.
\begin{itemize}
\item[(a)] Consider equation \cite[(1.40)]{B-CMP122b} at the background configuration
$U_0 = U_{k,C}(V,V_\Lambda(V))$ of the step on the component $C$, where $V_\Lambda$ is the real
$\Lambda$-datum, the function of clause~(c) at real arguments. It has the solution
$b_\Lambda = B'_\Lambda(B';V)$, which satisfies
\begin{equation}\label{4.23:eq-lambda-datum-2}
|b_\Lambda|_\infty \le K_\Lambda\,|B'|_\infty .
\end{equation}
In particular $B'_\Lambda(0;V) = 0$.
\item[(b)] The map $B'\mapsto B'_\Lambda(B';V)$ is analytic on the open ball
$\{|B'|_\infty < c_Q\tilde\alpha_{1,k}\}$.
\item[(c)] The function $V_\Lambda$ is extended to complex arguments by
\begin{equation}\label{4.23:eq-lambda-datum-3}
V_\Lambda(e^{iB'}V) := e^{ib_\Lambda}\,V_\Lambda(V).
\end{equation}
As in \cite[(1.38)]{B-CMP122b}, this identity is the definition of the function.
The right-hand side depends on $B'$ and $V$ only through $\mathsf{V} = e^{iB'}V$.
At real $B'$ it coincides with the real function $V_\Lambda$ of
\cite[Proposition~1]{B-CMP122a},
that is, the function $V_\Lambda$ used at the radius $c_Q\tilde\alpha_{1,k}$ in this chapter.
\end{itemize}

Only the bound \eqref{4.23:eq-lambda-datum-2} of clause (a) enters the estimate of the
$\Lambda$-datum in this chapter; clauses (b) and (c) serve there as side conditions only.
\end{definition}

\begin{definition}[boundary terms: absolute bound and gauge invariance; printed without proof in
{\cite[(1.23)--(1.48), (1.11), p.~369]{B-CMP122b}}]\label{4.23:def-boundary-terms}
Fix the level $k$ and the large-field region $Z$, and let $C$, its collared enlargement
$C^{\sharp}$, the tube variable $\zeta=(\mathbf{U},\mathbf{J})$ and the real variables $\varpi$
be as in Notation~\ref{4.23:not-oscillation-seminorm}. Let $\mathbf{C}_k$ be Ba{\l}aban's sum
of boundary terms, introduced in \cite[p.~369]{B-CMP122b} by the sentence quoted below, which
includes the term $I(U_0)$ coming from the denominator. Write $\mathbf{D}^{Z}_k$ for the class of
localisation domains at level $k$ used in \cite[\S1]{B-CMP122b} relative to $Z$.

The members of $\mathbf{C}_k$ that do not depend on the coefficients $z$ of the complex source
and are localised in $C^{\sharp}$ form a family
\begin{equation*}
  \mathfrak{v}^{\mathbf{C}}(Y_4), \qquad Y_4\in\mathbf{D}^{Z}_k,\quad Y_4\subset C^{\sharp},
\end{equation*}
with the following properties.
\begin{enumerate}
\item[(1)] Each $\mathfrak{v}^{\mathbf{C}}(Y_4)$ is analytic on the tube
  $\tilde{U}^{c}_{k}(Y_4,\tilde\alpha_0,\tilde\alpha_1)$ over $Y_4$, the space of complex
  configurations of \cite[(2.34)--(2.39)]{B-CMP119}.
\item[(2)] Each $\mathfrak{v}^{\mathbf{C}}(Y_4)$ is invariant under gauge transformations acting
  jointly on $\zeta$ and on $\varpi$: for every $G$-valued gauge transformation $\lambda$,
  \begin{equation*}
    \mathfrak{v}^{\mathbf{C}}(Y_4)(\zeta^{\lambda};\varpi^{\lambda})
    =\mathfrak{v}^{\mathbf{C}}(Y_4)(\zeta;\varpi),
  \end{equation*}
  where $\zeta^{\lambda}$ and $\varpi^{\lambda}$ denote the transforms of $\zeta$ and $\varpi$
  under $\lambda$ in the law by which gauge transformations act jointly on the variables of the
  tube and on the real variables.
\item[(3)] The family is bounded in absolute value, summed over localisation domains:
  \begin{equation}\label{4.23:eq-boundary-terms-1}
    \sum_{Y_4\subset C^{\sharp}}\sup\bigl|\mathfrak{v}^{\mathbf{C}}(Y_4)\bigr|
    \le \mathsf{c}_{\mathbf{C}}:=1+c_I M^5 ,
  \end{equation}
  where $c_I$ is the constant written $O(1)$ in \cite[(1.11)]{B-CMP122b}, namely in
  \begin{equation*}
    |I(U_0)|<O(1)\log M|\Lambda|<O(1)M^5 ,
  \end{equation*}
  with $I(U_0)$, the term from the denominator, and the region $\Lambda$ as there; the supremum
  is taken over the domain of $\mathfrak{v}^{\mathbf{C}}(Y_4)$, that is, over the tube over
  $Y_4$ and over the real variables $\varpi$. The smallness of the sum asserted in
  \cite{B-CMP122b} is read as the bound \eqref{4.23:eq-boundary-terms-1} on the sum of the
  suprema of the absolute values.
\end{enumerate}
\end{definition}

The following sentences, on which properties (1) and (3) rest, are printed without proof in
\cite[(1.23)--(1.48), (1.11), p.~369]{B-CMP122b}; cited as printed. Property (2) is a hypothesis
clause of the statement and is not among them. In \cite[p.~369]{B-CMP122b}:
\begin{quotation}
they are all small, in fact their sum is also small, although we need boundedness only, and they
depend on background fields restricted to a neighborhood of the region $Z$, therefore they are
boundary terms \ldots{} We denote the sum of all these terms, including the term $I(U_0)$ from
the denominator, by $\mathbf{C}_k$
\end{quotation}
In \cite[\S1]{B-CMP122b}:
\begin{quotation}
the bounds are preserved if we replace this field by the corresponding $G^c$-valued field
\end{quotation}
\begin{quotation}
see the estimate (1.47), so we omit details now
\end{quotation}
Of the estimates behind these sentences, \cite{B-CMP122b} displays one chain, that of
\cite[(1.47), p.~368]{B-CMP122b}, and calls the others similar. Let $N_{\mathrm{av}}$ denote the
number of averaging steps in \cite[(1.23)--(1.47)]{B-CMP122b}; it is distinct from the window
$N=N(g_k)$ of this chapter. With $\beta_0$, $B_3$, $B_5$, $\nu$, $\varepsilon_k$, $R_k$ and the
function $q_1(\cdot)$ of the coupling as in \cite[(1.23)--(1.48)]{B-CMP122b} (there $R_k$ is
neither the ball radius $R$ nor the window radius function $\check{R}_k$), the displayed chain
bounds the left-hand side of \cite[(1.47)]{B-CMP122b} first by
\begin{equation*}
  O(1)\,g_k^{-2}\,(\alpha_{0,k}+\alpha_{1,k})\,\varepsilon_k\,L^{-N_{\mathrm{av}}}\,
  N_{\mathrm{av}}^{1+\beta_0}\,B_3^3B_5M^6\sum\sum\cdots
\end{equation*}
and then by
\begin{equation*}
  O(1)\,A_0C_1B_3^3B_5M^{10}\,p_0(g_k)\,q_1(g_k)\,R_k^7\,L^{-N_{\mathrm{av}}},
\end{equation*}
where $\sum\sum\cdots$ is the double sum displayed there. The powers of $g_k$ cancel, and the only
small factor is $L^{-N_{\mathrm{av}}}$. The condition on $N_{\mathrm{av}}$ used in
\cite[\S1]{B-CMP122b}, which determines that number there, is the lower bound
\begin{equation*}
  N_{\mathrm{av}}\ge \frac{\nu}{\log L}\,\log\log g_k^{-2};
\end{equation*}
it is a lower bound, not a cap in the sense of Definition~\ref{2.3:def-cap}, not a smallness
condition on the coupling, and not the restriction $N_{\mathrm{av}}\le R_k$. It is met by taking
$N_{\mathrm{av}}=\check{R}_k$, where $\check{R}_k\ge(\log(x_k-c_5))^{r_{\mathrm{pr}}}$ with
$r_{\mathrm{pr}}$ the exponent in the lower bound of the window radius function. With this choice
the bound \eqref{4.23:eq-boundary-terms-1} is used under the single condition, in supremum form,
\begin{equation}\label{4.23:eq-boundary-terms-2}
  L^{-\check{R}_k}\,p_0(g_k)\,q_1(g_k)\,R_k^7\le 1 .
\end{equation}
The choice $N_{\mathrm{av}}=\check{R}_k$ and the condition \eqref{4.23:eq-boundary-terms-2} are
made in the statement that fixes the window radius function, not in this definition.

The sizes given in \cite{B-CMP122b} for these terms assume that the number written $N$ there is
at most $R_k$, whereas the number of levels $N=N(g_k)$ of the window of this chapter satisfies
$N\le c_N\check{R}_k$ and $N\le L(\log x_k)^{r}$, with $c_N$ as fixed with the window and $r$ the
exponent of the upper bound on the window. Passing from one to the other
changes the values of constants only, by a power of $c_N$, a number fixed after $L$ and
$r_{\mathrm{pr}}$; the constant $\mathsf{c}_{\mathbf{C}}$ may have to be enlarged accordingly,
and in either reading it is of degree zero in $\log x_k$.

The family $\mathfrak{v}^{\mathbf{C}}$ is not among the terms estimated by the hypotheses of the
correlation theorem: the $z$-free part of $\mathbf{C}_k$ stands in none of them. Properties
(1)--(3) therefore enter the correlation theorem as an additional hypothesis, properties (1) and
(3) being taken by statement from \cite{B-CMP122b}.

\begin{definition}[Conventions for extensions, restriction and real bases]
\label{4.23:def-extension-conventions}
Fix the level $k$, a component $C$ with its collared domain $C^{\sharp}$, and the tube
\begin{equation*}
\mathfrak{S}=\tilde{U}^{c}_{k}(C^{\sharp},\tilde\alpha_0,\tilde\alpha_1)
\end{equation*}
of Notation~\ref{4.23:not-oscillation-seminorm}: Ba{\l}aban's space of complex configurations
\cite[(2.34)--(2.39)]{B-CMP119} over $C^{\sharp}$, with the second and the fourth of its
composite conditions imposed, in this chapter, on every admissible triple $X'$ contained in
$C^{\sharp}$. Its points are written
$\zeta=(\mathbf{U},\mathbf{J})$. The core of the component $C$ is denoted by
$\mathsf{C}=\mathsf{C}_C$, as fixed in this chapter. The nucleus of a term of the expansion, that
is, the unit of the expansion to which the term belongs, is denoted by $v$; it is a different
object from the core $\mathsf{C}$. A nucleus
is classed as smooth or rough, as interior or not, and as far or not, according to the
classification of nuclei fixed in this chapter. For a smooth nucleus $v$ at level $j$ we write $p_v$
and $Q_v$ for the level and the box that the box-selection rule of this chapter assigns to it;
the rule depends on a pair of parameters $(n,\mathsf{W})$. We write $M^{p}(\mathbf{U})$ for the
averages of $\mathbf{U}$ at level $p$ under Ba{\l}aban's block-averaging map
(Definition~\ref{1.5:def-block-average}), and $M^{\cdot}(\mathbf{U})$ for the family of these
averages. Conventions (a)--(e) below are in force throughout this section; convention (f) is in
force only where it is invoked by name.

\begin{enumerate}
\item[(a)] \emph{Extension of far nuclei.} Let $v$ be a smooth interior far nucleus. A term with nucleus
$v$ is evaluated at the box minimiser on $Q_v$ of the level-$p_v$ averages of its configuration
argument. Every other term is evaluated on the chain of the expansion itself, at the point where
the chain is evaluated. This includes every term whose nucleus is rough. One weight
$\bar{\mathsf{w}}_v$ is attached to each nucleus $v$, and it is evaluated at the fine weight
$\mathsf{L}_{\mathrm{f}}^{2e}$, where the base $\mathsf{L}_{\mathrm{f}}$ and the exponent $e$ are
fixed in this chapter. This agrees with \cite[p.~375]{B-CMP122b},
where the configuration
$\mathbf{U}$ enters through its averages $M^{\cdot}(\mathbf{U})$.

\item[(b)] \emph{Extension of the subtracted copies.} Convention (a) also governs the subtracted
copies of the add-and-subtract step, with the point $\zeta$ in place of the point at which the
chain is evaluated. For these copies the box-selection rule is used with
parameters $(n,\mathsf{W})=(k,3)$, and the configuration argument is the component $\mathbf{U}$
of $\zeta$. For a smooth far nucleus $v$ of such a copy we write
\begin{equation*}
\operatorname{ext}_v(\zeta):=(U_{p_v},J_{p_v})\bigl(Q_v,M^{p_v}(\mathbf{U}\restriction Q_v)\bigr)
\end{equation*}
for the pair, configuration and current, of box minimisers on $Q_v$ whose data are the
level-$p_v$ averages of the restriction of $\mathbf{U}$ to $Q_v$.
  \begin{itemize}
  \item A subtracted copy whose nucleus $v$ is a smooth far nucleus is evaluated at
  $\operatorname{ext}_v(\zeta)$.
  \item Every other subtracted copy, in particular every copy whose nucleus is rough, is
  evaluated at $\zeta$ itself.
  \end{itemize}
The subtracted copies are those of \cite{B-CMP122b}; there, as at \cite[p.~375]{B-CMP122b},
the configuration $\mathbf{U}$ enters through its averages $M^{\cdot}(\mathbf{U})$.

\item[(c)] \emph{Sums over localisation domains.} Where \cite{B-CMP122b} states that the sum of
the localised terms over their localisation domains is also small, the assertion is understood
to refer to the absolute sum of the terms over the localisation domains, weighted as in the
bounds for these terms. Each of the bounds \cite[(1.23), (1.27), (1.31), (1.47)]{B-CMP122b}
displayed for these terms is a weighted absolute sum, and smallness refers to the same
quantity.

\item[(d)] \emph{Restriction.} Every $\zeta\in\mathfrak{S}$ restricts to a point of
$\tilde{U}^{c}_{k}(Y_4,\tilde\alpha_0,\tilde\alpha_1)$ for every $Y_4\in\mathbf{D}_k$ with
$Y_4\subset C^{\sharp}$. This holds because, in the definition of the tube, the second and the
fourth of the composite conditions are imposed on every admissible triple $X'$ contained in
$C^{\sharp}$ (the admissible triples of the definition of the tube in
Notation~\ref{4.23:not-oscillation-seminorm}) and not only on $C^{\sharp}$ itself; it is not a
consequence of the conditions of \cite[p.~261]{B-CMP119}. Condition~(i) of
\cite[p.~261]{B-CMP119} bounds the minimisers
\begin{equation*}
U_{p,X}(M^{\cdot}(\mathbf{U})),\qquad J_{p,X}(M^{\cdot}(\mathbf{U}))
\end{equation*}
of the domain $X$, as in \cite[(2.35), (2.37)]{B-CMP119}, and does not by itself give the
corresponding condition for a smaller domain $Y_4\subset X$. An analogous restriction, with
different constants, occurs in
\cite{B-CMP116}; it is a precedent only and does not imply the property stated here.

\item[(e)] \emph{The real base.} Let $\mathfrak{V}_{\mathbb{R}}$ denote the class of real data
$V$ of this section, let $\Pi$ be the function of the real variables $\varpi$ of
Notation~\ref{4.23:not-oscillation-seminorm} that is fixed together with them, and let
$\operatorname{supp}\Pi$ denote its support, over which the variables $\varpi$ range. Let
$V\in\mathfrak{V}_{\mathbb{R}}$ and $\varpi\in\operatorname{supp}\Pi$. Consider Ba{\l}aban's
real objects of \cite{B-CMP122b}
\begin{equation*}
Y:=V_{\Lambda}(V),\qquad U_0:=U_{k,C}(V,Y)\ \text{on } \mathfrak{B}^0,\qquad
U^0:=U^0_{k,C}\ \text{on } \mathfrak{B}^C ,
\end{equation*}
where $\mathfrak{B}^0=\mathbb{B}_k(C)$ is the minimal determining set and
$\mathfrak{B}^C=\mathbf{B}''_k(C)$ the two-level localised set. These objects are required to
satisfy three conditions.
  \begin{itemize}
  \item The two pairs lie in the classes
  \begin{equation*}
  (U_0,\mathcal{J}(U_0))\in\mathfrak{K}_{109}(\mathfrak{B}^0;\tfrac12\gamma^{\flat}),\qquad
  (U^0,\mathcal{J}(U^0))\in\mathfrak{K}_{109}(\mathfrak{B}^C;\tfrac12\gamma^{\flat}),
  \end{equation*}
  with $A'=0$, under the condition \cite[(14)]{B-CMP102b}. Here $\mathcal{J}(U)$ denotes the
  current of the configuration $U$; $\mathfrak{K}_{109}(\mathfrak{B};\gamma_0)$ is the class of
  pairs over $\mathfrak{B}$ fixed with the localised small solution of this chapter, $A'$ is the
  potential entering its definition, and the class is taken at $\gamma_0=\tfrac12\gamma^{\flat}$,
  where $\gamma^{\flat}$ is the size at which the classes $\mathfrak{K}_{109}$ are taken in this
  chapter.
  \item Both currents $\mathcal{J}(U_0)$ and $\mathcal{J}(U^0)$ satisfy the gauge condition
  $\tilde{G}_0\mathcal{J}=0$ of Definition~\ref{4.4:def-local-data-classes}.
  \item The quantity $\mathsf{p}$ is bounded:
  \begin{equation}\label{4.23:eq-extension-conventions-1}
  \mathsf{p}:=\sup\,\ell^{2}\eta^{-2}\,\bigl|\partial U^0-1\bigr|\le 1 .
  \end{equation}
  \end{itemize}
Condition \eqref{4.23:eq-extension-conventions-1} is the condition
\cite[(3.35)]{B-CMP99b} at that paper's radius $\alpha_0$. It coincides with
\cite[(2), (14)]{B-CMP102b} after $\varepsilon_0$ is replaced by
$O(1)B_3^2B_5M^5\varepsilon_k$, as in \cite{B-CMP122b}; here $\ell$, $\eta$ and $\partial$ are
as in \cite[(3.35)]{B-CMP99b}, $B_3$, $B_5$ and $\varepsilon_k$ are as in \cite{B-CMP122b}, and
$M=L^{m'}$ is Ba{\l}aban's auxiliary parameter of that name.

Every threshold in (e) is an absolute number, as in \cite{B-CMP122b}. Every size it bounds is
small with $\gamma$. In this sense (e) is a condition that $\gamma$ be small.

On the core $\mathsf{C}$, (e) is imposed here and is not a statement of \cite{B-CMP122b}. In
\cite[p.~361]{B-CMP122b} the regularity of $U^0$ is given only on
$\Omega^{c}_{k}\cap\Omega''^{\sim}_{h+1}$, that is, on the support of $\hat\chi$; these sets
and the function $\hat\chi$ are those of \cite[p.~361]{B-CMP122b}. Nothing is stated there
about $U^0$ on $\mathsf{C}$.

\item[(f)] \emph{The real base at perturbed data.} The conditions of (e) are required also at
every real datum $(e^{iB}V,e^{ib}Y)$ whose real fields $B$ and $b$ satisfy
\begin{equation*}
|B|_\infty,\ |b|_\infty\le r_{\mathbb{R}}:=3.3\,(1+K_{\Lambda})\,c_Q\,\tilde\alpha_{1,k},
\end{equation*}
where $K_{\Lambda}$ and $c_Q$ are constants of this chapter and $\tilde\alpha_{1,k}$ is the
radius $\tilde\alpha_1$ at level $k$; at such a datum the objects of (e) are formed with
$(V,Y)$ replaced by $(e^{iB}V,e^{ib}Y)$. This requirement is additional to (e); like the
requirement of (e) on the core, it is imposed here and is not stated in \cite{B-CMP122b}. Apart
from the values of its thresholds, which are left open, requirement (f) is implied by (e): at
such a datum the perturbed configurations are real and lie in the class $\mathfrak{K}_{109}$, by
the statement on real pairs that accompanies the localised small solution of this chapter; the
condition \cite[(2)]{B-CMP102b} is gauge invariant, as stated in \cite{B-CMP102b}; and the
increments of the quantities bounded in (e) under the perturbation are
proportional to $r_{\mathbb{R}}$.
As in (e), the thresholds are absolute numbers and the sizes they bound are small with
$\gamma$.
\end{enumerate}
\end{definition}

\begin{definition}[Coupling-free radii for the second kept member]\label{4.23:def-coupling-free-radii}
Let $K_\sharp$ and $K_\natural$ be the constants of the theorem on the localised small solution, so that
$K_\natural\ge 6dK_\sharp^2$, and let $B_{\mathbb{H}}$ be the constant in the bounds of that
theorem. That theorem is stated under the smallness condition
$K_\natural(\gamma_0+\varepsilon_3)\le 1$ on its two free parameters $\gamma_0$ and $\varepsilon_3$, and
its data radius is $\varepsilon_\sharp=\varepsilon_3/(4K_\sharp)$. Let $c_Q$, $K_u$ and $c_4$ be the constants
fixed earlier in this chapter, and let $K_J$ be the constant of the current bound.
Let $\varepsilon_\alpha$ denote the radius $\alpha_1$ of \cite[(3.37)]{B-CMP99b}; it does not depend on the
coupling. Set
\begin{equation}\label{4.23:eq-coupling-free-radii-1}
\begin{gathered}
\gamma^\flat:=\varepsilon_3^\flat:=\frac{1}{2K_\natural},\qquad
\varepsilon_\sharp^\flat:=\frac{\varepsilon_3^\flat}{4K_\sharp},\\
K_{\mathfrak{s}}:=3+18d+2.2\,K_J,\qquad
K_1^0:=(c_Q+1.01\,K_u)\,B_{\mathbb{H}}.
\end{gathered}
\end{equation}
The theorem on the localised small solution is used at $(\gamma_0,\varepsilon_3)=(\gamma^\flat,\varepsilon_3^\flat)$;
since these are free parameters of that theorem, the choice places no condition on $L$ or on
the auxiliary parameters $\mathfrak{p}$. Define
\begin{equation}\label{4.23:eq-coupling-free-radii-a}
\begin{gathered}
\rho_{\mathrm{abs}}:=\min\Bigl\{\varepsilon_\sharp^\flat,\ \frac{\varepsilon_\sharp^\flat}{K_1^0},\
\frac{\min\{c_4,\,10^{-2}/K_u,\,1/20\}}{B_{\mathbb{H}}(1+K_1^0)},\\
\frac{\varepsilon_\alpha}{2B_{\mathbb{H}} K_1^0},\
\frac{\gamma^\flat}{2K_{\mathfrak{s}}B_{\mathbb{H}} K_1^0}\Bigr\},\\
\mathbb{D}(\rho):=\bigl\{(B',b):\ \max\{|B'|_\infty,|b|_\infty\}<\rho\bigr\}.
\end{gathered}
\end{equation}
The number $\rho_{\mathrm{abs}}$ depends only on $d$, $L$ and $\kappa_1$: it involves neither a coupling,
nor the constant $K_\Lambda$ attached to the small-field region $\Lambda$, nor any of the constants that
Ba{\l}aban's construction fixes by a condition of the form ``sufficiently small''. In the definition of
$\mathbb{D}(\rho)$, for a
component $C$ at level $k$, $B'$ is a $\mathfrak{g}^c$-valued field on the level-$k$ bonds of $C$
that do not lie in $\Lambda$, $b$ is a $\mathfrak{g}^c$-valued field on the bonds of the region
$\Lambda_C$ associated with $C$, and $|B'|_\infty$, $|b|_\infty$ denote the suprema over these
respective bonds of the norm in $\mathfrak{g}^c$ of the value of the field.
\end{definition}

\begin{proof}
We check that $(\gamma^\flat,\varepsilon_3^\flat)$ is an admissible choice of the two free parameters of the
theorem on the localised small solution, and that $\varepsilon_\sharp^\flat$ is the corresponding data radius.
By \eqref{4.23:eq-coupling-free-radii-1}, $\gamma^\flat+\varepsilon_3^\flat=2/(2K_\natural)=1/K_\natural$.
Hence $K_\natural(\gamma^\flat+\varepsilon_3^\flat)=1$, and the smallness condition
$K_\natural(\gamma_0+\varepsilon_3)\le1$ holds, with equality. Moreover,
\begin{equation*}
6dK_\sharp^2\,\gamma^\flat=\frac{6dK_\sharp^2}{2K_\natural}\le\frac12,
\end{equation*}
because $K_\natural\ge 6dK_\sharp^2$. Finally, substituting $\varepsilon_3=\varepsilon_3^\flat$ into
$\varepsilon_\sharp=\varepsilon_3/(4K_\sharp)$ gives $\varepsilon_\sharp^\flat$ as defined in
\eqref{4.23:eq-coupling-free-radii-1}. Thus the theorem applies at $(\gamma^\flat,\varepsilon_3^\flat)$ with
data radius $\varepsilon_\sharp^\flat$; since $\gamma_0$ and $\varepsilon_3$ are free parameters of that
theorem, no condition on $L$ or on $\mathfrak{p}$ is imposed.
\end{proof}

\begin{definition}[Construction of the second kept member]\label{4.23:def-kept-member-construction}
Fix a real configuration $V\in\mathfrak{V}_{\mathbb{R}}$, the class of real configurations over
which the equation for the $\Lambda$-datum is solved, and the real variables $\varpi$ (the
block field $B$ on $\mathbf{B}_0\cap C$ and the field $V''$ on the core $\mathsf{C}=\mathsf{C}_C$).
Let $Y:=V_\Lambda(V)$, and let $U_0:=U_{k,C}(V,Y)$ on $\mathfrak{B}^0=\mathbb{B}_k(C)$ and
$U^0:=U^0_{k,C}$ on $\mathfrak{B}^C=\mathbf{B}''_k(C)$ be the real configurations of
Ba{\l}aban's construction fixed in the base-point hypothesis. Let $\eta$
be Ba{\l}aban's lattice-spacing parameter, the factor of the generator in the exponential
parametrisation $e^{i\eta\mathsf{H}}U$ of configurations used below, $\mathcal{J}(\cdot)$ the
current of a configuration,
$\mathbb{H}[\mathfrak{B};(W,J);\sigma;\mathsf{d}]$ the localised small solution on
$\mathfrak{B}$ at the pair $(W,J)$, with third argument $\sigma$ (in this construction always
$\mathbf{1}$) and datum $\mathsf{d}$,
$\mathfrak{s}^{\sharp}_{\sigma}[\,\cdot\,]$ its source term, and
$\mathfrak{u}[W;U;\mathfrak{B}]$ the group-valued field attached on $\mathfrak{B}$ to a
configuration $W$ relative to $U$. Let $\rho_{\mathrm{abs}}$ and $\gamma^{\flat}$ be the
coupling-free radii of Definition~\ref{4.23:def-coupling-free-radii}. The first step of the
construction is the solution $B'_\Lambda(B';V)$ of the equation for the $\Lambda$-datum
(Definition~\ref{4.23:def-lambda-datum}, printed without proof in
\cite[p.~367]{B-CMP122b}; cited as printed). For a pair $(B',b)$ in the domain
$\mathbb{D}(\rho_{\mathrm{abs}})$ attached to
the radius $\rho_{\mathrm{abs}}$, in which
$b$ is the variable that stands for the $\Lambda$-datum, the remaining steps are the
following.
\begin{enumerate}
\item[(a)] On bonds of level $k$ let $\mathfrak{d}^0$ be equal to $B'$ off the set $\Lambda$ of
the equation for the $\Lambda$-datum and to $b$
on $\Lambda_C$, the set on which the $\Lambda$-datum $b$ is given. Lift $\mathfrak{d}^0$ to the
bonds of the finer lattices by $Q^{s*}$: on a bond of a finer lattice
joining two blocks of level $k$ the lift takes the value of $\mathfrak{d}^0$ on the level-$k$
bond joining these blocks, and it is $0$ on bonds inside a block. Put
\begin{equation*}
\mathsf{H}^0:=\mathbb{H}\bigl[\mathfrak{B}^0;(U_0,\mathcal{J}(U_0));\mathbf{1};(0,\mathfrak{d}^0)\bigr],
\qquad
\mathfrak{u}^0:=\mathfrak{u}\bigl[e^{i\eta\mathsf{H}^0}U_0;U_0;\mathfrak{B}^0\bigr].
\end{equation*}
\item[(b)] The block field $\mathfrak{f}$ on $\mathfrak{B}^C$ is defined as follows. On the
bonds $c$, of level $j$, that carry data of $U_0$, $\mathfrak{f}(c):=\mathfrak{b}_j(c)$, the
level-$j$ block field of the block-field statement evaluated at the pair $(\mathsf{H}^0,U_0)$;
these fields have the form of the expressions
\cite[(1.44)--(1.45)]{B-CMP122b}, and only that form, and none of the estimates of that paper,
is used here. Inside $\mathsf{C}$, $\mathfrak{f}(c):=0$. On the bonds $c$ crossing
$\partial\mathsf{C}$, which carry $V''$,
\begin{equation*}
\mathfrak{f}(c):=\frac1i\log\bigl[V''(c)\,\mathfrak{u}^0(c_+)^{-1}\,V''(c)^{-1}\bigr],
\quad\text{respectively}\quad
\mathfrak{f}(c):=\frac1i\log\mathfrak{u}^0(c_-).
\end{equation*}
\item[(c)] Put
\begin{gather*}
\mathsf{H}^C:=\mathbb{H}\bigl[\mathfrak{B}^C;(U^0,\mathcal{J}(U^0));\mathbf{1};(0,\mathfrak{f})\bigr],
\qquad
W:=e^{i\eta\mathsf{H}^C}U^0,\\
\mathfrak{j}^C:=\mathfrak{s}^{\sharp}_{\mathbf{1}}[\mathsf{H}^C],\qquad
J:=\mathcal{J}(U^0)+\mathfrak{j}^C,\qquad
\mathfrak{u}^C:=\mathfrak{u}[W;U^0;\mathfrak{B}^C].
\end{gather*}
Let $\tilde v$ be the block-constant
extension of the field equal to $\mathfrak{u}^0$ at base points off $\mathsf{C}$ and to $1$ on
$\mathsf{C}$, and put $\mathsf{v}:=\mathfrak{u}^C\cdot\tilde v$.
\item[(d)] For $(W',J')\in\mathfrak{K}_{109}(\mathfrak{B}^C;\gamma^{\flat})$, let
$\mathsf{T}[W',J']$ be the linearised localised minimiser map at $(W',J')$, whose definition
follows the present one,
and let $a,a'$ be block fields on $\mathbf{B}_0\cap C$. With the objects of steps (a)--(c), put
\begin{equation}\label{4.23:eq-kept-member-construction-a}
\begin{aligned}
\mathsf{Q}[W',J'](a,a')&:=-\frac12\sum_p\eta^d\,\hat\chi(x(p))\\
&\qquad\times\operatorname{tr}\Bigl[\bigl(D_{W'}\mathsf{T}[W',J']a\bigr)(p)\cdot
\bigl(D_{W'}\mathsf{T}[W',J']a'\bigr)(p)\Bigr],\\
\tilde{\mathfrak{q}}[V](B',b)&:=\mathsf{Q}[W,J](\hat B,\hat B),\qquad
\hat B(c):=\operatorname{Ad}_{\mathsf{v}(c_-)}B(c),\\
\mathfrak{q}[V](B')&:=\tilde{\mathfrak{q}}[V]\bigl(B',B'_\Lambda(B';V)\bigr),\\
\mathfrak{q}_C(\zeta;\varpi)&:=\mathfrak{q}[M^k(U)](B'),\qquad
e^{iB'}:=M^k(\mathbf{U})\,M^k(U)^{-1}.
\end{aligned}
\end{equation}
The form $\mathsf{Q}[W',J']$ is bilinear and symmetric, and no complex conjugation is taken in
it; $\hat\chi$ is the cut-off function of the localised construction. Here $D_W$ is the
operator of \cite[(3.4)]{B-CMP99b}: for the plaquette $p$ spanned at $x$ by the directions
$\mu,\nu$,
\begin{equation*}
(D_WA)(p)=\bigl(\nabla_{W,\mu}A_\nu-\nabla_{W,\nu}A_\mu\bigr)(x).
\end{equation*}
\end{enumerate}
In \eqref{4.23:eq-kept-member-construction-a}, $\zeta=(\mathbf{U},\mathbf{J})$ is a point of the
source tube $\mathfrak{S}$ of Notation~\ref{4.23:not-oscillation-seminorm}, $U$ is the real
configuration of the tube about which $\mathbf{U}$ is taken, and $M^k$ is Ba{\l}aban's
block-averaging map to level $k$ (Definition~\ref{1.5:def-block-average}).
The form $\mathfrak{q}_C$ is the \emph{second kept member}. The variable $\mathbf{J}$ does not
enter it. At complex points the objects $W$, $J$, $\mathsf{v}$ and $\hat B$ are defined by
steps (a)--(d) applied to the regular representative of the complex point; no gauge covariance
is used off the real slice. Finally, at the base point $\zeta=(U,0)$, with $V=M^k(U)$, one has
$B'=0$ and $B'_\Lambda(0;V)=0$, and all objects of steps (a)--(c) are trivial at
$(B',b)=(0,0)$: $\mathsf{H}^0=\mathsf{H}^C=0$, $\mathfrak{u}^0=1$, $\mathfrak{f}=0$,
$\mathfrak{j}^C=0$ and $\mathsf{v}=1$, so that $W=U^0$, $J=\mathcal{J}(U^0)$ and $\hat B=B$; and
\begin{equation}\label{4.23:eq-kept-member-construction-b}
\mathfrak{q}_C((U,0);\varpi)=\tilde{\mathfrak{q}}[V](0,0)
=\mathsf{Q}\bigl[U^0,\mathcal{J}(U^0)\bigr](B,B).
\end{equation}
\end{definition}

\begin{proof}[Proof of \eqref{4.23:eq-kept-member-construction-b}]
Take $(B',b)=(0,0)$. Then $\mathfrak{d}^0$ vanishes on every level-$k$ bond, and so does its
lift by $Q^{s*}$, so the datum $(0,\mathfrak{d}^0)$ is zero. By clause (a) of the theorem on
the localised small solution, the solution $\mathbb{H}[\,\cdot\,;\cdot\,;\cdot\,;\mathsf{d}]$
vanishes at $\mathsf{d}=0$; hence $\mathsf{H}^0=0$. Consequently
$e^{i\eta\mathsf{H}^0}U_0=U_0$, and, by the definition of the field
$\mathfrak{u}[\cdot;\cdot;\cdot]$, $\mathfrak{u}^0=\mathfrak{u}[U_0;U_0;\mathfrak{B}^0]=1$.
The first argument of the block fields $\mathfrak{b}_j$ is $\mathsf{H}^0=0$, and at the pair
$(0,U_0)$ the block-field statement gives $\mathfrak{b}_j=0$; so
$\mathfrak{f}=0$ on the bonds carrying data of $U_0$; $\mathfrak{f}=0$ inside $\mathsf{C}$ by
definition; and on a bond $c$ crossing $\partial\mathsf{C}$, since $\mathfrak{u}^0=1$,
\begin{equation*}
\frac1i\log\bigl[V''(c)\,\mathfrak{u}^0(c_+)^{-1}V''(c)^{-1}\bigr]=\frac1i\log 1=0,
\qquad
\frac1i\log\mathfrak{u}^0(c_-)=\frac1i\log1=0.
\end{equation*}
Hence $\mathfrak{f}=0$ on $\mathfrak{B}^C$, and the datum $(0,\mathfrak{f})$ of step (c) is
zero. By clause (a) of the theorem on the localised small solution again, $\mathsf{H}^C=0$,
so $W=U^0$. Hence the pair $(W,\mathcal{J}(U^0))$ is the pair $(U^0,\mathcal{J}(U^0))$, which
lies in the class $\mathfrak{K}_{109}(\mathfrak{B}^C;\tfrac12\gamma^{\flat})$ with component
$A'=0$ by the base-point hypothesis. The same theorem gives
$\mathfrak{s}^{\sharp}_{\mathbf{1}}[0]=0$ (the functional $F$ entering the source term
satisfies $F(W,0)=0$); so $\mathfrak{j}^C=\mathfrak{s}^{\sharp}_{\mathbf{1}}[0]=0$ and
$J=\mathcal{J}(U^0)$.
Further $\mathfrak{u}^C=\mathfrak{u}[U^0;U^0;\mathfrak{B}^C]=1$, and $\tilde v=1$ because
$\mathfrak{u}^0=1$ off $\mathsf{C}$ and the extended field is $1$ on $\mathsf{C}$; thus
$\mathsf{v}=1$ and $\hat B(c)=\operatorname{Ad}_{1}B(c)=B(c)$. By
\eqref{4.23:eq-kept-member-construction-a},
\begin{equation*}
\tilde{\mathfrak{q}}[V](0,0)=\mathsf{Q}\bigl[U^0,\mathcal{J}(U^0)\bigr](B,B).
\end{equation*}

It remains to reach $\tilde{\mathfrak{q}}[V](0,0)$ from $\mathfrak{q}_C((U,0);\varpi)$. At
$\zeta=(U,0)$ one has $\mathbf{U}=U$, so
$e^{iB'}=M^k(U)M^k(U)^{-1}=1$ and $B'=0$; here $V=M^k(U)$ by
\eqref{4.23:eq-kept-member-construction-a}. The solution of the equation for the
$\Lambda$-datum satisfies the bound quoted in Definition~\ref{4.23:def-lambda-datum} (printed
without proof in \cite[p.~367]{B-CMP122b}; cited as printed): $|B'_\Lambda|$ is at most a
constant times
the modulus of its argument. At the argument $0$ this gives $B'_\Lambda(0;V)=0$. Therefore
\begin{equation*}
\mathfrak{q}_C((U,0);\varpi)=\mathfrak{q}[V](0)
=\tilde{\mathfrak{q}}\bigl[V\bigr]\bigl(0,B'_\Lambda(0;V)\bigr)
=\tilde{\mathfrak{q}}[V](0,0),
\end{equation*}
and together with the previous display this is
\eqref{4.23:eq-kept-member-construction-b}.
\end{proof}

\begin{lemma}[The linearised localised minimiser map]\label{4.23:lem-linearised-minimiser}
Let $\gamma^{\flat}$ be the coupling-free radius of
Definition~\ref{4.23:def-coupling-free-radii}. Let $C$ be the component under consideration and
$\mathfrak{B}^C=\mathbf{B}''_k(C)$ its two-level localised set. Write
$\mathfrak{K}_{109}(\mathfrak{B}^C;\gamma^{\flat})$ for the class of pairs $(W',J')$ over
$\mathfrak{B}^C$ whose configuration satisfies \cite[(3.35)]{B-CMP99b}, whose potential
satisfies \cite[(3.37)]{B-CMP99b}, and
whose current satisfies $\sup \ell^3|J'|\le\gamma^{\flat}$. Let $\mathbf{1}$ denote the function
identically equal to one. For $(W',J')\in\mathfrak{K}_{109}(\mathfrak{B}^C;\gamma^{\flat})$
define two linear maps on the bounded fields $a$ that the theorem on the localised small solution
admits as the second component of its data:
\begin{align*}
\mathsf{T}^{\sharp}[W',J'] &:= H(\mathbf{1};W',J'),\\
\mathsf{T}^{\mathsf{L}}[W',J']\,a &:= \frac{d}{dt}\,
  \mathbb{H}\bigl[\mathfrak{B}^C;(W',J');\mathbf{1};(0,ta)\bigr]\Big|_{t=0}.
\end{align*}
Here $H(\sigma;W,J)$ is the operator of the statement on Ba{\l}aban's linearising operator,
with first argument $\sigma$; this operator is distinct from the Hamiltonian $H$ of
Definition~\ref{0.2:not-glossary}. The
map $\mathbb{H}[\mathfrak{B};(W,J);\sigma;\mathsf{d}]$ is the
localised small solution with data $\mathsf{d}$. The data $(0,ta)$ has vanishing first
component and second component $ta$. The map $\mathsf{T}^{\sharp}[W',J']$ is an operator by
definition. The map $\mathsf{T}^{\mathsf{L}}[W',J']$ is linear in $a$.

Let $\mathsf{T}$ denote either of $\mathsf{T}^{\sharp}$, $\mathsf{T}^{\mathsf{L}}$. Put
$K_{\mathsf{T}}:=2K_\sharp$ for $\mathsf{T}=\mathsf{T}^{\sharp}$ and
$K_{\mathsf{T}}:=B_{\mathbb{H}}$ for $\mathsf{T}=\mathsf{T}^{\mathsf{L}}$. Here $K_\sharp$ is the
constant of the statement on the linearising operator, and $B_{\mathbb{H}}$ is the constant of
the theorem on the localised small solution. We write $N^{\Delta}_y$ for the norm, indexed by
the point $y$, in which the theorem on the localised small solution bounds its solution, and
$\|\cdot\|_1$ for the norm on the data in that theorem. Then the following hold:
\begin{equation}\label{4.23:eq-linearised-minimiser-a}
\begin{gathered}
(W',J')\longmapsto\mathsf{T}[W',J']\ \text{ is holomorphic on }
  \mathfrak{K}_{109}(\mathfrak{B}^C;\gamma^{\flat}),\\
\sup_y N^{\Delta}_y\bigl(\mathsf{T}[W',J']\,a\bigr)\le K_{\mathsf{T}}\,|a|_\infty .
\end{gathered}
\end{equation}
\end{lemma}

\begin{proof}
\emph{Linearity of $\mathsf{T}^{\mathsf{L}}$.} By clause~(a) of the theorem on the localised
small solution the map
$\mathsf{d}\mapsto\mathbb{H}[\mathfrak{B}^C;(W',J');\mathbf{1};\mathsf{d}]$ is holomorphic, and
by the chain rule $\mathsf{T}^{\mathsf{L}}[W',J']\,a$ is its differential at $\mathsf{d}=0$
evaluated at $(0,a)$. A differential is linear, so $\mathsf{T}^{\mathsf{L}}[W',J']$ is linear
in $a$.

\emph{Holomorphy of $\mathsf{T}^{\sharp}$.} Clause~(b) of the statement on the linearising
operator $H(\sigma;W,J)$ asserts that this operator is holomorphic in $(W,J)$ on the class
$\mathfrak{K}_{109}$. We apply it with $\sigma=\mathbf{1}$ on the localised set
$\mathfrak{B}^C$ at radius $\gamma^{\flat}$, which is admissible by
Definition~\ref{4.23:def-coupling-free-radii}. This is the first line of
\eqref{4.23:eq-linearised-minimiser-a} for $\mathsf{T}^{\sharp}$.

\emph{Holomorphy of $\mathsf{T}^{\mathsf{L}}$.} Clause~(a) of the theorem on the localised
small solution, applied at the radius $\gamma^{\flat}$, which is admissible for it by the same
definition, asserts that $\mathbb{H}$ is holomorphic jointly in $(W,J,\mathsf{d})$. Fix
$a$, and fix $r>0$ such that $(0,ta)$ lies in the domain of the data for $|t|\le r$. The
map $t\mapsto\mathbb{H}[\mathfrak{B}^C;(W',J');\mathbf{1};(0,ta)]$ is then holomorphic on
$|t|\le r$. Cauchy's formula for the derivative at $t=0$ gives
\begin{equation*}
\mathsf{T}^{\mathsf{L}}[W',J']\,a=\frac{1}{2\pi i}\oint_{|t|=r}
  t^{-2}\,\mathbb{H}\bigl[\mathfrak{B}^C;(W',J');\mathbf{1};(0,ta)\bigr]\,dt .
\end{equation*}
By joint holomorphy, the integrand is continuous in $(t,W',J')$ and holomorphic in
$(W',J')$ for each $t$ on the circle. Hence the integral is holomorphic in $(W',J')$ on
$\mathfrak{K}_{109}(\mathfrak{B}^C;\gamma^{\flat})$. This is the first line of
\eqref{4.23:eq-linearised-minimiser-a} for $\mathsf{T}^{\mathsf{L}}$.

\emph{The bound for $\mathsf{T}^{\sharp}$.} Clause~(a) of the statement on the linearising
operator bounds $\sup_y N^{\Delta}_y$ of the image of $a$. The bound is a constant times
$|a|_\infty$, and the constant equals $2K_\sharp$ when the weight function of that clause is
identically one, which is the case of the operator $H(\mathbf{1};W',J')$. Since
$\mathsf{T}^{\sharp}=H(\mathbf{1};W',J')$, this gives the second line
of \eqref{4.23:eq-linearised-minimiser-a} with $K_{\mathsf{T}}=2K_\sharp$.

\emph{The bound for $\mathsf{T}^{\mathsf{L}}$.} Clause~(b) of the theorem on the localised
small solution bounds $\sup_y N^{\Delta}_y$ of
$\mathbb{H}[\mathfrak{B}^C;(W',J');\mathbf{1};\mathsf{d}]$ by $B_{\mathbb{H}}\|\mathsf{d}\|_1$.
We take $\mathsf{d}=(0,ta)$, whose first component vanishes, so that
$\|(0,ta)\|_1\le|t|\,|a|_\infty$. This gives
\begin{equation*}
\sup_y N^{\Delta}_y\Bigl(\mathbb{H}\bigl[\mathfrak{B}^C;(W',J');\mathbf{1};(0,ta)\bigr]\Bigr)
 \le B_{\mathbb{H}}\,\|(0,ta)\|_1\le B_{\mathbb{H}}\,|t|\,|a|_\infty .
\end{equation*}
At $t=0$ the right-hand side vanishes, so every $N^{\Delta}_y$ of the solution at the zero
datum is zero. By the triangle inequality for $N^{\Delta}_y$, the difference quotient
\begin{equation*}
t^{-1}\Bigl(\mathbb{H}\bigl[\mathfrak{B}^C;(W',J');\mathbf{1};(0,ta)\bigr]
 -\mathbb{H}\bigl[\mathfrak{B}^C;(W',J');\mathbf{1};(0,0)\bigr]\Bigr)
\end{equation*}
therefore has $\sup_y N^{\Delta}_y$ at most $B_{\mathbb{H}}|a|_\infty$ for every $t\neq0$.
Letting $t\to0$ gives the second line of \eqref{4.23:eq-linearised-minimiser-a} with
$K_{\mathsf{T}}=B_{\mathbb{H}}$.
\end{proof}

\begin{remark}[The choice of $\mathsf{T}$ and Ba{\l}aban's operator $H''_{1,k,C}$]
The estimates of this section into which the map $\mathsf{T}$ enters use it only through the
two properties \eqref{4.23:eq-linearised-minimiser-a}. They therefore hold for either choice
of $\mathsf{T}$. The choice made in the construction of the second kept member is the
following.

By clause~(d) of the statement on the linearising operator, $\mathsf{T}^{\sharp}$ is the
operator $H$ of \cite[(3.126)]{B-CMP99b}; in Ba{\l}aban's notation there,
$HB=GQ^{*}(QGQ^{*})^{-1}B$. This is the
operator $H$ entering the map $\mathcal{H}=A'_1-HD(A'_1)$ in the notation of
\cite{B-CMP99b}.

The operator $H''_{1,k,Z}$ of \cite[(1.20)]{B-CMP122b} is a different operator. We write
$\varsigma$ for the cut-off function that Ba{\l}aban uses there. The display
\cite[(1.20), p.~361]{B-CMP122b} reads
\begin{align*}
A(\varsigma,U''_{k,Z})
 &=A\bigl(\varsigma,\exp\bigl(i\eta\mathbf{H}''_{k,Z}(g_kB)\bigr)U^0_{k,Z}\bigr)\\
 &=A(\varsigma,U^0_{k,Z})
  +g_k\bigl\langle DH''_{1,k,Z}B,\;\varsigma\eta^{-2}\operatorname{Im}\partial U^0_{k,Z}\bigr\rangle\\
 &\quad+\tfrac12 g_k^2\bigl\langle H''_{1,k,Z}B,\;\Delta_1(\varsigma)H''_{1,k,Z}B\bigr\rangle
  +V\bigl(\varsigma,g_kH''_{1,k,Z}B\bigr).
\end{align*}
It is obtained there using \cite[(2.8)]{B-CMP109}. Display \cite[(2.6)]{B-CMP109} contains
both operators $H_{1,k}$ and $H_k$, with reference to their definitions
\cite[(3.127), (3.128)]{B-CMP99b}. In \cite{B-CMP102b} the expansion of $\mathcal{H}$ begins
with the first-order term $H_1B$. The operator $H_1$ of \cite[(3.129)]{B-CMP99b}, in the same
notation $H_1B=G_1Q^{*}(QG_1Q^{*})^{-1}B$, carries the new quadratic terms of the
linearisation. These
terms are connected with the linear term and are small because the current $J$ is small. The
operator $H$ of \cite[(3.126)]{B-CMP99b} does not carry them.

Consequently, with $\mathcal{H}$ the map named above and taking $Z=C$, the component fixed in
Definition~\ref{4.23:def-kept-member-construction}, clause~(d) of the theorem on the localised
small solution gives, at real critical bases in the sense of that clause,
\begin{equation}\label{4.23:eq-linearised-minimiser-b}
H''_{1,k,C}=\partial_B\mathcal{H}\big|_{B=0}=\mathsf{T}^{\mathsf{L}} .
\end{equation}
The difference $\mathsf{T}^{\mathsf{L}}-\mathsf{T}^{\sharp}$, by contrast, is of relative size
$O(\ell^3|J|)$ and does not vanish identically. It comes from the current-dependent terms of the
linearisation.

Accordingly, in the construction of the second kept member
(Definition~\ref{4.23:def-kept-member-construction}) the map $\mathsf{T}$ is taken to be
$\mathsf{T}^{\mathsf{L}}$. With this choice the constant in the bound
$\mathfrak{q}_{\max}(k)$ on the second
kept member is $C_{\mathfrak{q}}=15c_{\mathrm{tr}}B_{\mathbb{H}}^2$, where $c_{\mathrm{tr}}$ is
the trace constant of that construction. With $\mathsf{T}^{\sharp}$ one would obtain instead
$15c_{\mathrm{tr}}(2K_\sharp)^2=60c_{\mathrm{tr}}K_\sharp^2$; the two values differ by the
factor $B_{\mathbb{H}}^2/(4K_\sharp^2)$, the ratio
of the squares of the two values of $K_{\mathsf{T}}$. The degrees of the bound
$\mathfrak{q}_{\max}(k)$ are the same for either choice, and the constant $C_\Sigma$ of
\eqref{4.23:eq-degree-constant-a}, which depends on $C_{\mathfrak{q}}$, is taken with this
value of $C_{\mathfrak{q}}$.

The identification of the second kept member, at real points, with the quadratic form of
\cite[(1.71)]{B-CMP122b} holds for $\mathsf{T}^{\mathsf{L}}$ only.
\end{remark}

\begin{theorem}[The form theorem: bound, variation, sign]\label{4.23:thm-form-theorem}
Let $\mathsf{T}\in\{\mathsf{T}^{\mathsf{L}},\mathsf{T}^{\sharp}\}$ be one of the two operators of the
linearised localised minimiser map (Lemma~\ref{4.23:lem-linearised-minimiser}), and let
$K_{\mathsf{T}}$ be its norm bound from that lemma. Assume the following.
\begin{itemize}
\item The four floor conditions (one-sided lower bounds) of the identification statement for the
localised solution: those on $\kappa_1$, on $M_1$ and on the fine block size $\mathfrak{M}$,
and the level floor of that statement.
\item (F1)(a): clause~(a) of the analytic solvability of the equation for the
$\Lambda$-datum (Definition~\ref{4.23:def-lambda-datum}, printed without proof in
\cite[(1.41), p.~367]{B-CMP122b}; cited as printed); $K_\Lambda$ denotes the constant of that
clause.
\item (F2): the localisation statement for the localised small solution
$\mathbb{H}[\mathfrak{B};(W,J);\sigma;\mathsf{d}]$ relative to the large-field region $Z$,
imposed at both sets $\mathfrak{B}^0=\mathbb{B}_k(C)$ and $\mathfrak{B}^C=\mathbf{B}''_k(C)$,
together with the three scope clauses stated with it in this section.
\item (F3): the radius hypothesis on $\mathbb{H}$, read at the coupling-free radii
$(\gamma^\flat,\varepsilon_3^\flat)$ of Definition~\ref{4.23:def-coupling-free-radii}.
\item (F4): the block-field hypothesis, which supplies the block fields
$\mathfrak{b}_j(c)$.
\item (F5): the second clause of the hypothesis $\mathfrak{U}$ on the map
$\mathfrak{u}[\cdot;\cdot;\cdot]$ that enters the construction of the second kept member
(Definition~\ref{4.23:def-kept-member-construction}).
\item (F6): the real-base hypothesis, which supplies the real base
$U^0=U^0_{k,C}$ on $\mathfrak{B}^C$, with its clause on the parameter $\mathsf{p}$ and its
size conditions.
\item (F7): the bound on the count $\mathsf{n}_{\mathrm{sh}}(k)$ supplied by the nesting
statement, which enters the definition of $\mathfrak{q}_{\max}(k)$ below.
\item The five conditions used in the proof of (X3) below: the radius condition that
accompanies the construction of the second kept member, the window condition of the
localisation statement (F2), the reference condition, the definition of the window radius
$\check{R}_k$, and the fifth amplitude clause.
\item The ceiling
\begin{equation}\label{4.23:eq-form-theorem-b}
14\,(1+K_\Lambda)\,c_Q\,\bar{\alpha}_1(x)\le\rho_{\mathrm{abs}}
\qquad\text{for all } x\ge\gamma^{-2},
\end{equation}
with $\rho_{\mathrm{abs}}$ as in \eqref{4.23:eq-coupling-free-radii-a}. Here
$\bar{\alpha}_1(\cdot)$ and, in (X3) below, $\bar{\mathfrak{q}}_{\max}(\cdot)$ denote the
functions of the inverse coupling $x$ that correspond to $\tilde{\alpha}_{1,k}$ and to
$\mathfrak{q}_{\max}(k)$; they are fixed where (X3) is proved.
\end{itemize}
For the rider (X0) assume in addition clauses (b) and (c) of
Definition~\ref{4.23:def-lambda-datum}, printed without proof in
\cite[(1.38), (1.41), p.~367]{B-CMP122b}; cited as printed (clause (c) comprising the definition
\cite[(1.38)]{B-CMP122b} and the identification stated with it); and (F8): the identification
statement for the localised solution at
$\mathfrak{B}^0$ and at $\mathfrak{B}^C$, the identity \cite[(181)]{B-CMP102b}, and the
strengthened real-base hypothesis. Let $\vartheta_X$ be the number defined in
\eqref{4.23:eq-data-ray-a} below,
\begin{equation}\label{4.23:eq-form-theorem-1}
\vartheta_X=(1+K_\Lambda)\max\{\vartheta^0,\,2/\rho_{\mathrm{abs}}\},
\end{equation}
where $\vartheta^0$ is the number of \eqref{4.23:eq-kernel-disc-a}, and put
\begin{equation}\label{4.23:eq-form-theorem-2}
C^{\mathsf{T}}_{\mathfrak{q}}:=15\,c_{\mathrm{tr}}K_{\mathsf{T}}^2,\qquad
\mathfrak{q}_{\max}(k):=C^{\mathsf{T}}_{\mathfrak{q}}\,(\mathfrak{b}'_k)^2\,M^d\,
\mathsf{n}_{\mathrm{sh}}(k),
\end{equation}
where $c_{\mathrm{tr}}$ is the trace constant.
Then for every $K$, at the level $k=k_W\le K$, for every label of the
expansion term considered,
every component $C$ of $Z$ at
$\mathbf{R}_k$, every $\varpi$ in the support of $\Pi$ (the measure over the real variables
$\varpi$ of \cite[(1.71)]{B-CMP122b}) and every admissible splitting
$(\zeta,U)$ with $\zeta=(U'U,\mathbf{J})\in\mathfrak{S}_{\mathrm{src}}$, whether or not the
field $A'$ is real (the field on which the second of the three conditions defining the class
$\mathfrak{K}_{109}(\mathfrak{B};\gamma_0)$ that enters
Lemma~\ref{4.23:lem-linearised-minimiser} is imposed), the following hold, where $B'$ is the
$\mathfrak{g}^c$-valued field, $\mathfrak{g}^c$ the complexified Lie algebra of the gauge
group $G$, with
$U'=e^{iB'}$:
\begin{equation}\label{4.23:eq-form-theorem-a}
\begin{gathered}
\text{clauses (X1)--(X5) below, and, under the additional}\\
\text{hypotheses of the rider, clause (X0).}
\end{gathered}
\end{equation}
\begin{itemize}
\item[(X1)] For every real configuration $V$ of
Definition~\ref{4.23:def-kept-member-construction}, on the polydisc
$\mathbb{D}(\rho_{\mathrm{abs}})$ of \eqref{4.23:eq-coupling-free-radii-a},
\begin{equation*}
|\tilde{\mathfrak{q}}[V]|\le\mathfrak{q}_{\max}(k);
\end{equation*}
consequently $|\mathfrak{q}_C|\le\mathfrak{q}_{\max}(k)$ on the source tube $\mathfrak{S}$
(Notation~\ref{4.23:not-oscillation-seminorm}).
\item[(X2)] The variation of the form is bounded by
\begin{equation}\label{4.23:eq-form-theorem-3}
\begin{split}
|\mathfrak{q}_C(\zeta;\varpi)-\mathfrak{q}_C((U,0);\varpi)|
&\le\vartheta_X\,|B'|_\infty\,\mathfrak{q}_{\max}(k)\\
&\le\vartheta_X\,c_Q\,\tilde{\alpha}_{1,k}\,\mathfrak{q}_{\max}(k).
\end{split}
\end{equation}
\item[(X3)] The right member of \eqref{4.23:eq-form-theorem-3} is at most
$\Sigma^{\mathfrak{q}}(x_k)$, with $x_k=x_{k_W}$, where
\begin{equation}\label{4.23:eq-form-theorem-4}
\Sigma^{\mathfrak{q}}(x):=\sup_{x'\ge x}\vartheta_X\,c_Q\,\bar{\alpha}_1(x')\,
\bar{\mathfrak{q}}_{\max}(x').
\end{equation}
The function $\Sigma^{\mathfrak{q}}$ is finite, non-increasing, and tends to $0$ as its
argument tends to infinity; at any later step $i$,
\begin{equation}\label{4.23:eq-form-theorem-5}
\Sigma^{\mathfrak{q}}(x_{k_W})\le\Sigma^{\mathfrak{q}}(x_i-2c_2).
\end{equation}
\item[(X4)] $\mathfrak{q}_C((U,0);\varpi)$ is real and $\le0$; hence
$\operatorname{Re}\mathfrak{q}_C\le\Sigma^{\mathfrak{q}}(x_k)$ and
$|\operatorname{Im}\mathfrak{q}_C|\le\Sigma^{\mathfrak{q}}(x_k)$.
\item[(X5)] The quantities $\vartheta_X$, $C^{\mathsf{T}}_{\mathfrak{q}}$,
$\rho_{\mathrm{abs}}$ and the ceiling \eqref{4.23:eq-form-theorem-b} depend neither on $K$,
$k$, the age of any large-field region (the number of steps since the step that created it),
the number $N$ of levels of the window, the parameters $N_0$ and $D$ (the letter $D$ here
being a parameter, not the covariant derivative of (X0)),
the window radius $\check{R}$,
the size $|\mathsf{C}|$ of the core $\mathsf{C}$ of $C$, the number, size or past of the
nested large-field regions $(Y',m')$ with $Y'\subset C$ or of the arrested regions nested in
$C$, nor on the volume; the core enters through the real base $U^0$ alone;
$\mathfrak{q}_{\max}$ depends on $N$, $\check{R}_{h+1}$ and $\mathfrak{b}'_k$.
\item[(X0)] (Rider.) Under the additional hypotheses above and for
$\mathsf{T}=\mathsf{T}^{\mathsf{L}}$: $\mathfrak{q}_C(\zeta;\varpi)$ does not depend on the
splitting, is holomorphic on $\mathfrak{S}_{\mathrm{src}}$, depends on $\zeta$ only through
the restriction $\mathsf{V}$ to $C\setminus\Lambda$ of the level-$k$ configuration
$M^k(\mathbf{U})$ determined by $\mathbf{U}$ (here $M^k(\cdot)$ denotes an operation on
configurations, not a power of the block size $M$), and at real $\zeta$ equals the form of
\cite[(1.71)]{B-CMP122b},
\begin{equation}\label{4.23:eq-form-theorem-6}
-\tfrac12\bigl\langle DH''_{1,k,C}B,\;\hat{\chi}\,DH''_{1,k,C}B\bigr\rangle ,
\end{equation}
in which $D$ is the lattice covariant derivative, $\hat{\chi}$, $H''_{1,k,C}$ are as in
\cite[(1.71)]{B-CMP122b}, and $B$ is the component of $\varpi$ on $\mathbf{B}_0\cap C$;
moreover $\mathfrak{q}_C(\zeta^g;\varpi^g)=\mathfrak{q}_C(\zeta;\varpi)$
for every $G$-valued gauge transformation $g$, acting jointly on $\zeta$ and $\varpi$.
\end{itemize}
The convergence to $0$ in (X3) is that of the function $\Sigma^{\mathfrak{q}}$ of the inverse
coupling and holds at each fixed splitting. Nothing is asserted about two points of the tube with
different unitary parts.
\end{theorem}

\begin{proof}
Each of the clauses (X1)--(X5) is proved from the hypotheses listed before the rider, and
the rider (X0) from these together with the additional hypotheses of the rider, in the
results that follow in this section: clause (X1) in
the bound on the second kept member on the polydisc $\mathbb{D}(\rho_{\mathrm{abs}})$; clause
(X2) in the variation bound for the form, which rests on the kernel disc
\eqref{4.23:eq-kernel-disc-a} and the data ray \eqref{4.23:eq-data-ray-a}; clauses (X3) and
(X5) in the statement on the majorant $\Sigma^{\mathfrak{q}}$ and on the dependence of the
constants; clause (X4) and the rider (X0) in the sign and identification statement for the
form; together these give the theorem.
\end{proof}

\begin{lemma}[The localised minimiser lemma at both sets]\label{4.23:lem-minimiser-scope}
Let $C$ be a component of the construction of this chapter; by that construction $C$ is a union
of cubes of side $M\check{R}_k$, and the fine block size $\mathfrak{M}$ divides $M\check{R}_k$.
Let $\mathfrak{B}^0=\mathbb{B}_k(C)$ be the
minimal determining set and $\mathfrak{B}^C=\mathbf{B}''_k(C)$ the two-level localised set. Then:
\begin{enumerate}
\item[(1)] the provisos of Ba{\l}aban's localisation \cite[(2.14)]{B-CMP119}, read with $M_1$
replaced by the fine block size $\mathfrak{M}$, hold at $\mathfrak{B}^0$ and at $\mathfrak{B}^C$;
\item[(2)] the three scope conditions (c1), (c2), (c3) of the localised small-solution hypothesis,
which enters Theorem~\ref{4.23:thm-form-theorem} as one of its hypotheses, hold at
$\mathfrak{B}^0$ and at $\mathfrak{B}^C$.
\end{enumerate}
Consequently the scope of that hypothesis is met at $\mathfrak{B}\in\{\mathfrak{B}^0,\mathfrak{B}^C\}$
with $\sigma\equiv\mathbf{1}$. Hence, under the floors on the auxiliary parameters that accompany
that hypothesis, its conclusions hold at both sets, with $\sigma\equiv\mathbf{1}$, for the
localised small solution $\mathbb{H}[\mathfrak{B};(W,J);\mathbf{1};\mathsf{d}]$ and its source
term $\mathfrak{s}^{\sharp}_{\sigma}[\mathbb{H}]$, whenever $(W,J)$ lies in the class on which
the hypothesis is formulated and $\|\mathsf{d}\|_1\le\varepsilon^{\flat}_{\sharp}$.
\end{lemma}

\begin{proof}
\emph{The cube structure of $C$.} Since the fine block size $\mathfrak{M}$ divides
$M\check{R}_k$, every cube of side $M\check{R}_k$ is a union of cubes of side $\mathfrak{M}$, and
so is $C$, which is a union of such cubes.

\emph{The provisos of \cite[(2.14)]{B-CMP119}.} Let $Z''_k\subset\Lambda_C$ and
$\Lambda_C^{\sim}$ be the sets attached to $C$ in the construction of $\mathbf{B}''_k(C)$, and let
$\operatorname{dist}_k$ denote the distance at level $k$. Off $Z''_k$ the parent
set $\mathbf{B}''_k$ is of level $k$ in $C$. By the statement of the geometric lemma on the
separation of $\Lambda_C^{\sim}$ from the complement of $C$,
\begin{equation*}
\operatorname{dist}_k\bigl(\Lambda_C^{\sim},C^{c}\bigr)\ \ge\ 9\check{R}_k .
\end{equation*}
The provisos of \cite[(2.14)]{B-CMP119}, read with $\mathfrak{M}$ in place of $M_1$, are stated in
Ba{\l}aban's notation, in which $\Omega$ is the set on which one localises and
$\Omega_j\supset\Omega_{j+1}$ are consecutive strata of the parent; these letters are used in this
paragraph only. Here $\Omega=C$, with $j=k$, so that $\Omega_k\setminus\Omega_{k+1}$ is the
stratum of level $k$ of the parent in $C$. The provisos require three things. The boundary
$\partial\Omega$ must lie in $\Omega_j\setminus\Omega_{j+1}$, that is, $\partial\Omega$ meets one
stratum of the parent only. The set $\Omega$ must be a union of $\mathfrak{M}$-cubes. The third
proviso is the separation condition: $\partial\Omega$ must lie at distance at least
$2\mathfrak{M}$ from the other strata of the parent. The first requirement is met on the
strength of the fact noted above, that off $Z''_k$ the parent set is of level $k$ in $C$. The
second requirement was established in the preceding paragraph. For the distance requirement, the
distance in question is controlled by the separation of $\Lambda_C^{\sim}$ from $C^{c}$; the
level-$k$ distance being counted in units of $M$ sites, the separation bound gives a room
of $9\check{R}_kM$ sites. Since $\mathfrak{M}$ divides $M\check{R}_k$, we have
$\mathfrak{M}\le M\check{R}_k$, and therefore
\begin{equation}\label{4.23:eq-minimiser-scope-1}
9\check{R}_kM\ \ge\ 9\mathfrak{M}\ \ge\ 2\mathfrak{M}.
\end{equation}
At $\mathfrak{B}^0$ the parent is the one-level set $\mathbf{B}_k\restriction C$ (see the
verification of (c1) below); it has a single level, so the first proviso is immediate there, and
the second and third are the properties of $C$ just established.
This proves (1).

\emph{Condition (c1).} Condition (c1) asks for three properties: nested levels; sets that are
unions of $\mathfrak{M}$-cubes of their own scale; and every stratum of width at least
$\mathfrak{M}$ sites of its own scale.

At either set the following holds. By the layer property of the parent set in the construction
of this chapter, each stratum of the parent contains at least one layer of $M$-cubes of its own
scale. By \cite{B-CMP119}, each grade of the localisation contains at least $L\mathfrak{M}$ sites
of its own scale. Accordingly, boundaries of distinct strata of the parent are at least one layer
of own-scale $M$-cubes apart, and boundaries of distinct grades are at least $L\mathfrak{M}$
own-scale sites apart. The requirement that the sets be unions of $\mathfrak{M}$-cubes rests on
the cube structure of $C$ established in the first paragraph.

At $\mathfrak{B}^0$, moreover, the parent is the one-level set $\mathbf{B}_k\restriction C$, the
restriction to $C$ of the set $\mathbf{B}_k$ of \cite[(1.33)]{B-CMP122b}. The nesting condition
is therefore met trivially, and the stratum and grade bounds of the preceding paragraph apply to
its single level $k$. Thus (c1) holds at both sets.

\emph{Condition (c2).} Condition (c2) asks that the localised set refine the set from which it is
built and coincide with it off the localisation layer. Both $\mathfrak{B}^0$ and $\mathfrak{B}^C$
are constructed with this property, so (c2) holds at both sets.

\emph{Condition (c3).} Condition (c3) constrains the set $\sigma_0$ attached to the interpolation
parameter $\sigma$ to lie in the host cubes of the parent. In the present application only
$\sigma\equiv\mathbf{1}$ occurs, and (c3) is satisfied for this value. This proves (2), and the
scope of the localised small-solution hypothesis is met at both sets.

\emph{The conclusions used.} At $\mathfrak{B}\in\{\mathfrak{B}^0,\mathfrak{B}^C\}$, with
$\sigma\equiv\mathbf{1}$ and $\|\mathsf{d}\|_1\le\varepsilon^{\flat}_{\sharp}$, the conclusions
of that hypothesis invoked in this section are of two kinds. The main ones are clauses (a) and
(b) of the theorem on the localised small solution $\mathbb{H}$, together with clauses (a) and (b)
of the lemma on its source term $\mathfrak{s}^{\sharp}_{\sigma}[\mathbb{H}]$. The remaining ones,
namely clauses (c), (d) and (e) of that theorem and clause (i) of the parametrised form of that
theorem, enter only as side clauses accompanying these.
\end{proof}

\begin{lemma}[Auxiliary sizes on the coupling-free bidisc]\label{4.23:lem-bidisc-sizes}
Assume the standing hypotheses of this section numbered (F2)--(F6), under which the second kept
member is constructed in
Definition~\ref{4.23:def-kept-member-construction}. Let $\rho_{\mathrm{abs}}$ and the bidisc
$\mathbb{D}(\rho_{\mathrm{abs}})$ be those of Definition~\ref{4.23:def-coupling-free-radii}
(see \eqref{4.23:eq-coupling-free-radii-a}). For $(B',b)\in\mathbb{D}(\rho_{\mathrm{abs}})$ put,
in this lemma and its proof only, $s:=\max\{|B'|_\infty,|b|_\infty\}$. Then the objects
$\mathfrak{d}^0$, $\mathsf{H}^0$,
$\mathfrak{u}^0$, $\mathfrak{f}$, $\mathsf{H}^C$, $W$, $\mathfrak{j}^C$, $J$, $\mathfrak{u}^C$,
$\tilde{\mathsf{v}}$ and $\mathsf{v}$ of the construction
\eqref{4.23:eq-kept-member-construction-a} exist on the whole of
$\mathbb{D}(\rho_{\mathrm{abs}})$, they are holomorphic in $(B',b)$, and
\begin{equation}\label{4.23:eq-bidisc-sizes-a}
\begin{aligned}
&\text{(A1)}\quad && \sup N^{\Delta}(\mathsf{H}^0)\le B_\sharp s;\\
&\text{(A2)} && |\mathfrak{b}_j(c)|\le c_Q B_\sharp s,\qquad
|\log \mathfrak{u}^0|\le K_u B_\sharp s;\\
&\text{(A3)} && |\mathfrak{f}|_\infty\le K_1^0 s;\\
&\text{(A4)} && \sup N^{\Delta}(\mathsf{H}^C)\le B_\sharp K_1^0 s\le \tfrac12\varepsilon_\alpha;\\
&\text{(A5)} && \sup \ell^3|\mathfrak{j}^C|\le K_{\mathfrak{s}}B_\sharp K_1^0 s
\le \tfrac12\gamma^\flat;\\
&\text{(A6)} && |\log \mathsf{v}|_\infty\le 1.05\,K_u B_\sharp(1+K_1^0)s\le 1.05\cdot 10^{-2};\\
&\text{(A7)} && (W,J)\in\mathfrak{K}_{109}(\mathfrak{B}^C;\gamma^\flat),
\end{aligned}
\end{equation}
where in (A4) the norm $N^{\Delta}$ is taken at the kernel $U^0$. Here $N^{\Delta}$ denotes the
weighted norm in which clause (b) of the theorem on the localised small solution $\mathbb{H}$ is
stated, computed with the covariant derivatives of the kernel at which it is taken, and $\mathcal{N}_y$ and
$\mathcal{N}$ denote the associated norm at a point $y$ and the global norm; $\ell$ and
$\ell_j$ are the scale weights entering these norms (at a bond of level $j$ the weight is
$\ell_j$), $\ell_{\mathfrak{B}}(y)$ is the weight of the set $\mathfrak{B}$ at $y$, and
$\|\cdot\|_1$ is the norm of the datum of $\mathbb{H}$, which is the supremum of a block field.
The proof uses only the following relations, which hold by the definitions of these norms: for a
field $A$ and a kernel $U$,
\begin{equation*}
\max\{\ell|A|,\ \ell^2|\nabla_UA|,\ \ell^3|\Delta_UA|\}\le N^{\Delta}(A),\qquad
\mathcal{N}_y(A)\le\sup N^{\Delta}(A),\qquad \mathcal{N}(A)\le\sup N^{\Delta}(A),
\end{equation*}
where $\nabla_U$ is the covariant derivative at $U$ and $\Delta_U$ is the second-order operator
of the identity for the source term used in the proof of (A5).
\end{lemma}

\begin{proof}
Since $(B',b)\in\mathbb{D}(\rho_{\mathrm{abs}})$ we have $s<\rho_{\mathrm{abs}}$, and
$\rho_{\mathrm{abs}}$ is the minimum of the five numbers listed in
\eqref{4.23:eq-coupling-free-radii-a}. Hence $s$ lies below each of them, that is,
\begin{equation}\label{4.23:eq-bidisc-sizes-1}
\begin{aligned}
& s<\varepsilon^\flat_\sharp, \qquad K_1^0 s<\varepsilon^\flat_\sharp,\qquad
B_\sharp(1+K_1^0)s<\min\{c_4,\,10^{-2}/K_u,\,1/20\},\\
& 2B_\sharp K_1^0 s<\varepsilon_\alpha,\qquad 2K_{\mathfrak{s}}B_\sharp K_1^0 s<\gamma^\flat .
\end{aligned}
\end{equation}
We refer to these five inequalities, in the order written, as the first to the fifth
inequality of \eqref{4.23:eq-bidisc-sizes-1}. The clauses are proved in order.

\emph{Proof of} (A1). By the construction \eqref{4.23:eq-kept-member-construction-a}, the
field $\mathfrak{d}^0$ on the level-$k$ bonds takes the values of $B'$ or of $b$, and its lift to
the finer grades either repeats such a value (on a grade bond joining two level-$k$ blocks) or
puts $0$ (inside a block). The norm $\|\cdot\|_1$ of a block field is its supremum. Therefore,
by the first inequality of \eqref{4.23:eq-bidisc-sizes-1},
\begin{equation*}
\|(0,\mathfrak{d}^0)\|_1\le s<\varepsilon^\flat_\sharp .
\end{equation*}
By (F6), $(U_0,\mathcal{J}(U_0))\in\mathfrak{K}_{109}(\mathfrak{B}^0;\tfrac12\gamma^\flat)$,
where $\mathcal{J}(U)$ is the current attached to a configuration $U$ in the construction,
and this class is contained in $\mathfrak{K}_{109}(\mathfrak{B}^0;\gamma^\flat)$, because the
only condition in its definition that depends on the last argument is that $\gamma_0$ bound the
supremum of $\ell^3$ times the modulus of the current component. The theorem on the localised
small solution $\mathbb{H}$ applies at
$\mathfrak{B}^0$ by Lemma~\ref{4.23:lem-minimiser-scope}. Its clause (a) gives the existence of
$\mathsf{H}^0$ and its holomorphy in the datum, and its clause (b) gives
\begin{equation*}
N^{\Delta}(\mathsf{H}^0)\le B_\sharp\|(0,\mathfrak{d}^0)\|_1\le B_\sharp s .
\end{equation*}

\emph{Proof of} (A2). By the third inequality of \eqref{4.23:eq-bidisc-sizes-1},
\begin{equation*}
B_\sharp s\le B_\sharp(1+K_1^0)s<\min\{c_4,\,10^{-2}/K_u\}.
\end{equation*}
We apply the statement that defines the block fields $\mathfrak{b}_j$ of a small field relative
to a configuration, to the small field $\mathsf{H}^0$ and the configuration $U_0$. It gives
$|\mathfrak{b}_j(c)|\le c_Q\sup\ell_j|\mathsf{H}^0|$. The set $\mathfrak{B}^C$ refines
$\mathfrak{B}^0$. The two sets coincide on the rim $\mathfrak{L}_C$. On the subset
$C^{\sim-2}$ of $C$ the bonds
of $\mathfrak{B}^C$ have levels at most $k$, whereas those of $\mathfrak{B}^0$ have level
$k$. Consequently, at a bond of level $j$ at the point $y$ one has
$\ell_j\le\ell_{\mathfrak{B}^0}(y)$, and therefore, by the definition of $\mathcal{N}_y$ and by (A1),
\begin{equation*}
\ell_j|\mathsf{H}^0|\le \mathcal{N}_y(\mathsf{H}^0)\le B_\sharp s .
\end{equation*}
This proves $|\mathfrak{b}_j(c)|\le c_QB_\sharp s$.

For $\mathfrak{u}^0=\mathfrak{u}[e^{i\eta\mathsf{H}^0}U_0;U_0;\mathfrak{B}^0]$, with $\eta$ as in
the construction, we use clause (ii)
of the statement on the gauge map $\mathfrak{u}[\cdot;\cdot;\cdot]$ at the configuration $U_0$. Its
hypothesis holds, since, by (A1), $\mathcal{N}(\mathsf{H}^0)\le B_\sharp s<c_4$. It gives
\begin{equation*}
|\log\mathfrak{u}^0|\le K_u\sup \mathcal{N}_y(\mathsf{H}^0)\le K_uB_\sharp s .
\end{equation*}

\emph{Proof of} (A3). The block field $\mathfrak{f}$ on $\mathfrak{B}^C$ is defined separately
on three kinds of bonds.
\begin{itemize}
\item On the bonds $c$ of level $j$ that carry data of $U_0$, $\mathfrak{f}(c)=\mathfrak{b}_j(c)$.
By (A2), $|\mathfrak{f}(c)|\le c_QB_\sharp s$ there.
\item On the bonds crossing $\partial\mathsf{C}$, which carry $V''$, the field is given by
$\mathfrak{f}(c)=(1/i)\log[V''(c)\mathfrak{u}^0(c_+)^{-1}V''(c)^{-1}]$, respectively by
$\mathfrak{f}(c)=(1/i)\log\mathfrak{u}^0(c_-)$. Since $V''(c)\in G$, the first expression
equals $-(1/i)\operatorname{Ad}_{V''(c)}\log\mathfrak{u}^0(c_+)$, and $|\cdot|$ is
$\operatorname{Ad}_G$-invariant. Hence $|\mathfrak{f}(c)|=|\log\mathfrak{u}^0(c_\pm)|$, and by
(A2) this is at most $K_uB_\sharp s$.
\item On the bonds inside the core $\mathsf{C}$, $\mathfrak{f}(c)=0$.
\end{itemize}
Since $K_1^0=(c_Q+1.01K_u)B_\sharp$, it follows that
\begin{equation*}
|\mathfrak{f}|_\infty\le\max\{c_Q,K_u\}B_\sharp s\le K_1^0 s .
\end{equation*}

\emph{Proof of} (A4). By (A3) and the second inequality of \eqref{4.23:eq-bidisc-sizes-1},
$|\mathfrak{f}|_\infty\le K_1^0s<\varepsilon^\flat_\sharp$. By (F6),
$(U^0,\mathcal{J}(U^0))\in\mathfrak{K}_{109}(\mathfrak{B}^C;\gamma^\flat)$. The theorem on the
localised small solution applies at $\mathfrak{B}^C$ by Lemma~\ref{4.23:lem-minimiser-scope}.
Its clause (a) gives the existence and holomorphy of
\begin{equation*}
\mathsf{H}^C=\mathbb{H}[\mathfrak{B}^C;(U^0,\mathcal{J}(U^0));\mathbf{1};(0,\mathfrak{f})].
\end{equation*}
Its clause (b), with the norm at the kernel $U^0$, gives
\begin{equation*}
N^{\Delta}(\mathsf{H}^C)\le B_\sharp|\mathfrak{f}|_\infty\le B_\sharp K_1^0s
\le\tfrac12\varepsilon_\alpha,
\end{equation*}
where the last step is the fourth inequality of \eqref{4.23:eq-bidisc-sizes-1}.

\emph{Proof of} (A5). We write the identity for the source term $\mathfrak{s}^\sharp_\sigma$
that accompanies the theorem on the localised small solution, for the base pair
$(U^0,\mathcal{J}(U^0))$ and $\sigma\equiv\mathbf{1}$, that is, for the problem whose solution is
$\mathsf{H}^C$, and bound it term by term. It reads
\begin{equation}\label{4.23:eq-bidisc-sizes-2}
\mathfrak{j}^C=(\Delta_{U^0}-P_1+\mathsf{K})A'-\mathfrak{m}_H(X)+F(U^0,\mathsf{H}^C),
\end{equation}
where $A'=\mathsf{H}^C$ is the solution and $X$ the auxiliary field of the variational problem
whose solution is $\mathsf{H}^C$; $\Delta_{U^0}$, $P_1$ and $\mathsf{K}$ are the operators of the
identity acting on $A'$, $\mathfrak{m}_H(X)$ is its term in the auxiliary field, and
$F(U^0,\mathsf{H}^C)$ is its remainder term.
The terms are estimated as follows.
\begin{itemize}
\item By the definition of the norm, $\ell^3|\Delta_{U^0}A'|\le N^{\Delta}(A')$. This
quantity, and likewise $\ell^3|P_1A'|$ and $\ell^3|\mathfrak{m}_H(X)|$, is at most
$B_\sharp|\mathfrak{f}|_\infty$ by clause (b) of the theorem on the localised small solution.
\item By the bound on the operator $\mathsf{K}$, $\ell^3|\mathsf{K}A'|\le 18d\,\mathsf{p}\,\ell|A'|$,
where $\mathsf{p}$ is the parameter of the base configuration that enters the hypotheses of
the remainder lemma.
Hypothesis (F6) gives $\mathsf{p}\le1$, and $\ell|A'|\le N^{\Delta}(A')$ by the definition of
the norm. Hence
$\ell^3|\mathsf{K}A'|\le 18dB_\sharp|\mathfrak{f}|_\infty$.
\item By the sharpened remainder bound, valid under the hypotheses of the remainder lemma,
\begin{equation*}
\ell^3|F(U^0,\mathsf{H}^C)|\le K_J(\mathsf{p}+\Theta_0+\Theta_1)(\Theta_0+\Theta_1),
\end{equation*}
where $\Theta_0,\Theta_1$ are the two size parameters of the field in that lemma.
These hypotheses hold at $U^0$: the base $U^0$ is near-unitary,
$\mathsf{p}\le1$, and, by (A4) and the third inequality of \eqref{4.23:eq-bidisc-sizes-1},
\begin{equation*}
\Theta_0+\Theta_1\le 2\mathcal{N}(\mathsf{H}^C)\le 2B_\sharp K_1^0 s
\le 2\cdot\tfrac1{20}\cdot\frac{K_1^0}{1+K_1^0}<\tfrac1{10}.
\end{equation*}
Thus $\mathsf{p}+\Theta_0+\Theta_1\le1.1$. Since also
$\Theta_0+\Theta_1\le 2\mathcal{N}(\mathsf{H}^C)\le 2B_\sharp|\mathfrak{f}|_\infty$, we obtain
$\ell^3|F(U^0,\mathsf{H}^C)|\le 1.1K_J\cdot 2B_\sharp|\mathfrak{f}|_\infty$.
\end{itemize}
The bound on $\mathsf{K}$ and the remainder bound are local statements in the configuration
and the field. They do not refer to the set, so they apply on $\mathfrak{B}^C$ as they stand.
Summing the five contributions in \eqref{4.23:eq-bidisc-sizes-2} and using
$K_{\mathfrak{s}}=3+18d+2.2K_J$, (A3) and the fifth inequality of
\eqref{4.23:eq-bidisc-sizes-1}, we get
\begin{equation*}
\begin{aligned}
\ell^3|\mathfrak{j}^C|
&\le(1+1+1+18d+2.2K_J)B_\sharp|\mathfrak{f}|_\infty
=K_{\mathfrak{s}}B_\sharp|\mathfrak{f}|_\infty\\
&\le K_{\mathfrak{s}}B_\sharp K_1^0s\le\tfrac12\gamma^\flat .
\end{aligned}
\end{equation*}

\emph{Proof of} (A6). We use clause (ii) of the statement on the gauge map at
$(\mathsf{H}^C,U^0)$, for $\mathfrak{u}^C=\mathfrak{u}[W;U^0;\mathfrak{B}^C]$. Its hypothesis
holds, since by (A4) and the third inequality of \eqref{4.23:eq-bidisc-sizes-1}
\begin{equation*}
\mathcal{N}(\mathsf{H}^C)\le B_\sharp K_1^0s\le B_\sharp(1+K_1^0)s<c_4 .
\end{equation*}
It gives $|\log\mathfrak{u}^C|\le K_uB_\sharp K_1^0s$. The field $\tilde{\mathsf{v}}$ is the
block-constant extension of $\mathfrak{u}^0$ at the base points off $\mathsf{C}$ and of $1$ on
$\mathsf{C}$, so (A2) gives $|\log\tilde{\mathsf{v}}|\le K_uB_\sharp s$.

Now $\mathsf{v}=\mathfrak{u}^C\cdot\tilde{\mathsf{v}}=e^{E_1}e^{E_2}$ with
$E_1=\log\mathfrak{u}^C$ and
$E_2=\log\tilde{\mathsf{v}}$. By the third inequality of \eqref{4.23:eq-bidisc-sizes-1},
\begin{equation*}
|E_1|+|E_2|\le K_uB_\sharp(1+K_1^0)s<10^{-2}.
\end{equation*}
At this size $|\log(e^{E_1}e^{E_2})|\le1.05(|E_1|+|E_2|)$. Indeed, the first correction in the
Baker--Campbell--Hausdorff series satisfies $\tfrac12|[E_1,E_2]|\le|E_1||E_2|$, and
\begin{equation*}
|E_1||E_2|\le\tfrac14(|E_1|+|E_2|)^2\le0.0025\,(|E_1|+|E_2|).
\end{equation*}
This correction and the tail of the series together stay below $0.05(|E_1|+|E_2|)$ when
$|E_1|+|E_2|<10^{-2}$. Therefore
\begin{equation*}
|\log\mathsf{v}|_\infty\le1.05\,K_uB_\sharp(1+K_1^0)s\le1.05\cdot10^{-2}.
\end{equation*}

\emph{Proof of} (A7). Since $W=e^{i\eta\mathsf{H}^C}U^0$, membership of $(W,J)$ in
$\mathfrak{K}_{109}(\mathfrak{B}^C;\gamma^\flat)$ asks that $U^0$ satisfy the condition
\cite[(3.35)]{B-CMP99b}, which holds by (F6), and that $A'=\mathsf{H}^C$ obey the condition
\cite[(3.37)]{B-CMP99b}, in which the derivative is
the covariant derivative $\nabla$ at $U=U^0$. The latter holds: by (A4),
\begin{equation*}
\max\{\ell|\mathsf{H}^C|,\ \ell^2|\nabla_{U^0}\mathsf{H}^C|\}\le\tfrac12\varepsilon_\alpha
<\varepsilon_\alpha.
\end{equation*}
For the current $J=\mathcal{J}(U^0)+\mathfrak{j}^C$ we use (A5) for the second term:
\begin{equation*}
\sup\ell^3|J|\le\sup\ell^3|\mathcal{J}(U^0)|+\sup\ell^3|\mathfrak{j}^C|
\le\tfrac12\gamma^\flat+\tfrac12\gamma^\flat=\gamma^\flat .
\end{equation*}
The class $\mathfrak{K}_{109}(\mathfrak{B}^C;\gamma^\flat)$ imposes no other condition on $J$.
Hence $(W,J)\in\mathfrak{K}_{109}(\mathfrak{B}^C;\gamma^\flat)$.

\emph{Existence and holomorphy.} Each object of the construction is obtained from $(B',b)$ by
composing the following maps:
\begin{itemize}
\item the localised small solution, holomorphic by clause (a) of its theorem;
\item the block fields $\mathfrak{b}_j$;
\item the gauge map $\mathfrak{u}$;
\item the logarithm near $1$, the exponential, and $\operatorname{Ad}$.
\end{itemize}
The bounds (A1)--(A6), obtained above, keep the argument of each of these maps inside the range
in which it was applied. On that range each map exists and is holomorphic. Hence the objects
exist on the whole of $\mathbb{D}(\rho_{\mathrm{abs}})$ and are holomorphic in $(B',b)$.
\end{proof}

\begin{lemma}[The quadratic form and the shallow-shell count]\label{4.23:lem-shell-count-bound}
Let $\mathsf{T}$ be either of the maps $\mathsf{T}^{\sharp}$, $\mathsf{T}^{\mathsf{L}}$ of
Lemma~\ref{4.23:lem-linearised-minimiser}, and let $K_{\mathsf{T}}$ be its constant in
clause (t2) of \eqref{4.23:eq-linearised-minimiser-a}. For
$(W',J')\in\mathfrak{K}_{109}(\mathfrak{B}^C;\gamma^{\flat})$ let $\mathsf{Q}[W',J']$ be the
bilinear form
\begin{equation*}
\mathsf{Q}[W',J'](a,a')
=-\tfrac12\sum_{p}\eta^{d}\,\hat\chi(p)\,
\operatorname{tr}\bigl[(D_{W'}\mathsf{T}a)(p)\,(D_{W'}\mathsf{T}a')(p)\bigr],
\end{equation*}
with $\mathsf{T}=\mathsf{T}[W',J']$, and with $\eta$, $\hat\chi$ and the covariant plaquette
derivative $D_{W'}$ as in \eqref{4.23:eq-kept-member-construction-a}. Then $\mathsf{Q}[W',J']$ is
holomorphic in $(W',J')$ on $\mathfrak{K}_{109}(\mathfrak{B}^C;\gamma^{\flat})$, and for all
$a,a'$ on which $\mathsf{T}$ acts
\begin{equation}\label{4.23:eq-shell-count-bound-a}
\bigl|\mathsf{Q}[W',J'](a,a')\bigr|
\le 12\,c_{\mathrm{tr}}\,K_{\mathsf{T}}^{2}\,|a|_\infty\,|a'|_\infty\,M^{d}\,
\mathsf{n}_{\mathrm{sh}}(k).
\end{equation}
Here $c_{\mathrm{tr}}$ is the trace constant of this chapter, and $\mathsf{n}_{\mathrm{sh}}(k)$
is the number fixed in the shallow-shell count used in the proof below.
\end{lemma}

\begin{proof}
Fix a plaquette $p$ and let $y$ be the block with $x(p)\in\Delta(y)$; write $\ell=\ell(y)$
for its side. The covariant plaquette derivative is bounded by twice the covariant
gradient,
\begin{equation*}
|(D_{W'}A)(p)|\le 2\sup|\nabla_{W'}A|\le 2\,\ell^{-2}\,N^{\Delta}_y(A).
\end{equation*}
The second inequality holds because the gradient entering the block norm $N^{\Delta}_y$ is
the covariant gradient of the kernel hypotheses, and because the norm is taken over the
enlarged block $\tilde\Delta(y)$, which covers the plaquettes that straddle the boundary of
$\Delta(y)$. The supremum is therefore controlled by $N^{\Delta}_y(A)$ for every $p$ with
$x(p)\in\Delta(y)$.

Next we count plaquettes. In dimension $d=4$ there are six plaquette orientations at each
site, and $\Delta(y)$ contains $\ell^{4}\eta^{-4}$ sites. Hence
\begin{equation}\label{4.23:eq-shell-count-bound-1}
\sum_{p:\,x(p)\in\Delta(y)}\eta^{4}=6\,\ell(y)^{4}.
\end{equation}
By clause (t2) of \eqref{4.23:eq-linearised-minimiser-a},
$N^{\Delta}_y(\mathsf{T}a)\le K_{\mathsf{T}}|a|_\infty$ for every $y$, and likewise for
$a'$.

Take $A=\mathsf{T}a$ and $A=\mathsf{T}a'$ in the first display. Use the trace constant,
$0\le\hat\chi\le1$ and \eqref{4.23:eq-shell-count-bound-1}. The plaquettes with
barycentre in one block $\Delta(y)$ then contribute to $|\mathsf{Q}[W',J'](a,a')|$ at most
\begin{align*}
&\tfrac12\,c_{\mathrm{tr}}\cdot 6\ell^{4}\cdot
\bigl(2\ell^{-2}K_{\mathsf{T}}|a|_\infty\bigr)\bigl(2\ell^{-2}K_{\mathsf{T}}|a'|_\infty\bigr)\\
&\qquad=12\,c_{\mathrm{tr}}\,K_{\mathsf{T}}^{2}\,|a|_\infty\,|a'|_\infty .
\end{align*}
Only blocks meeting $\operatorname{supp}\hat\chi$ contribute, since $\hat\chi$ vanishes
elsewhere. By the shallow-shell count, the geometric statement on the support of
$\hat\chi$ in force in this section, the support of $\hat\chi$ meets at most
$M^{d}\mathsf{n}_{\mathrm{sh}}(k)$ blocks $\Delta(y)$, all of them blocks of
$\mathfrak{B}^C$. Summing the per-block bound over these blocks gives
\eqref{4.23:eq-shell-count-bound-a}.

It remains to prove holomorphy. The operator $D_{W'}$ is a polynomial in the entries of
$\operatorname{Ad}_{W'}$, and $\operatorname{Ad}_{W'}X=W'X{W'}^{-1}$ has entries that are
polynomials in the entries of $W'$ and ${W'}^{-1}$. By clause (t1) of
\eqref{4.23:eq-linearised-minimiser-a}, $\mathsf{T}[W',J']$ is holomorphic in $(W',J')$ on
$\mathfrak{K}_{109}(\mathfrak{B}^C;\gamma^{\flat})$. The weights $\eta^{d}\hat\chi(p)$ do not
depend on $(W',J')$. Hence $\mathsf{Q}[W',J'](a,a')$ is a finite sum of products of
holomorphic functions of $(W',J')$, and is therefore holomorphic.
\end{proof}

\begin{proof}[Proof of clause \textup{(X1)} of \eqref{4.23:eq-form-theorem-a}]
We work under the hypotheses of Lemma~\ref{4.23:lem-bidisc-sizes}. By clause (A6) of
\eqref{4.23:eq-bidisc-sizes-a}, $|\log\mathsf{v}|_\infty\le 1.05\cdot10^{-2}$, so
\begin{equation*}
\|\operatorname{Ad}_{\mathsf{v}}\|\le e^{2|\log\mathsf{v}|_\infty}\le e^{0.021}.
\end{equation*}
Since $e^{0.021}\le 1.03$, this gives $|\hat B|_\infty\le 1.03\,\mathfrak{b}'_k$ for the
field $\hat B$ of the construction of the second kept member
(Definition~\ref{4.23:def-kept-member-construction}). Moreover
\begin{equation*}
12\cdot(1.03)^{2}=12.7308\le15.
\end{equation*}
By clause (A7) of \eqref{4.23:eq-bidisc-sizes-a},
$(W,J)\in\mathfrak{K}_{109}(\mathfrak{B}^C;\gamma^{\flat})$ at every point of
$\mathbb{D}(\rho_{\mathrm{abs}})$ (see \eqref{4.23:eq-coupling-free-radii-a}), so
Lemma~\ref{4.23:lem-shell-count-bound} applies at $(W',J')=(W,J)$. The value
$\tilde{\mathfrak{q}}[V](B',b)$ is the form $\mathsf{Q}[W,J]$ evaluated at $a=a'=\hat B$.
Inserting the bound on $|\hat B|_\infty$ into \eqref{4.23:eq-shell-count-bound-a} gives, on
$\mathbb{D}(\rho_{\mathrm{abs}})$,
\begin{equation}\label{4.23:eq-shell-count-bound-2}
\bigl|\tilde{\mathfrak{q}}[V](B',b)\bigr|
\le 15\,c_{\mathrm{tr}}\,K_{\mathsf{T}}^{2}\,(\mathfrak{b}'_k)^{2}\,M^{d}\,
\mathsf{n}_{\mathrm{sh}}(k)=\mathfrak{q}_{\max}(k).
\end{equation}
The construction of the second kept member evaluates $\tilde{\mathfrak{q}}[V]$ at the points
$(B',B'_{\Lambda}(B'))$ of \eqref{4.23:eq-kept-member-construction-a}. By the first step of
the proof of the lemma immediately following this one in this section, these points lie in
$\mathbb{D}(\rho_{\mathrm{abs}})$. By \eqref{4.23:eq-shell-count-bound-2}, therefore,
$|\mathfrak{q}_C|\le\mathfrak{q}_{\max}(k)$ on the tube.
\end{proof}

\begin{remark}
The bound (q1) of the tube statement, \eqref{4.23:eq-tube-statement-a}, is by definition
clause (X1) of \eqref{4.23:eq-form-theorem-a} at $\mathsf{T}=\mathsf{T}^{\mathsf{L}}$, for
which $K_{\mathsf{T}}=B_{\sharp}$, restricted to the tube $\mathfrak{S}$.
\end{remark}

\begin{lemma}[Variation of the form: the kernel disc]\label{4.23:prf-kernel-disc}
The following is the variation bound in \eqref{4.23:eq-form-theorem-a} of
Theorem~\ref{4.23:thm-form-theorem}; it fixes the number $\vartheta^0$ entering the constant
$\vartheta_X$ there. Assume the hypotheses of
Theorem~\ref{4.23:thm-form-theorem}, among them the ceiling \eqref{4.23:eq-form-theorem-b}, and let
$\mathsf{T}\in\{\mathsf{T}^{\mathsf{L}},\mathsf{T}^{\sharp}\}$. Fix $K$, a level $k=k_W\le K$, a
label, a component $C$ of $Z$ at $\mathbf{R}_k$ for that label, and $\varpi\in\operatorname{supp}\Pi$.
Let $(\zeta,U)$ be an admissible splitting, $\zeta=(U'U,\mathbf{J})\in\mathfrak{S}_{\mathrm{src}}$,
with $A'$ real or not, whose field $B'$ is not zero. Let $V$ be the $k$-fold block average of $U$
(Definition~\ref{1.5:def-block-average}), and let $b_\Lambda:=B'_\Lambda(B';V)$ be the solution of the
equation for the $\Lambda$-datum (Definition~\ref{4.23:def-lambda-datum}, which is stated in
\cite[(1.41), p.~367]{B-CMP122b} without proof). Put $s_0:=\max\{|B'|_\infty,|b_\Lambda|_\infty\}$.
Then
\begin{equation}\label{4.23:eq-kernel-disc-a}
\begin{gathered}
\bigl|\mathfrak{q}_C(\zeta;\varpi)-\mathfrak{q}_C((U,0);\varpi)\bigr|
\le\mathfrak{q}_{\max}(k)\,s_0\,\vartheta^0,\\
\vartheta^0:=4B_\sharp K_1^0\max\Bigl\{\frac{1}{\varepsilon_{\alpha}},\frac{K_{\mathfrak{s}}}{\gamma^\flat}\Bigr\}
+3.7K_uB_\sharp(1+K_1^0),
\end{gathered}
\end{equation}
and $s_0\le(1+K_\Lambda)|B'|_\infty$.
\end{lemma}

\begin{proof}
If $B'=0$, both members of the variation bound in \eqref{4.23:eq-form-theorem-a} vanish, the left
one by \eqref{4.23:eq-kept-member-construction-b}. We therefore assume $B'\neq0$, as in the
statement, and proceed in six steps.

\emph{Step 1: the point $(B',b_\Lambda)$.} By clause (a) of Definition~\ref{4.23:def-lambda-datum},
which is stated in \cite[(1.41), p.~367]{B-CMP122b} without proof, the solution
$b_\Lambda=B'_\Lambda(B';V)$ exists, and its bound gives the first inequality in
\begin{equation}\label{4.23:eq-kernel-disc-2}
s_0\le(1+K_\Lambda)|B'|_\infty<(1+K_\Lambda)c_Q\tilde\alpha_{1,k}
\le(1+K_\Lambda)c_Q\bar\alpha_1(x_k)\le\frac{\rho_{\mathrm{abs}}}{14}.
\end{equation}
The second and third inequalities are the bounds on the radii of the source tube
$\mathfrak{S}_{\mathrm{src}}$: the field of a point of
$\mathfrak{S}_{\mathrm{src}}$ satisfies $|B'|_\infty<c_Q\tilde\alpha_{1,k}$, and
$\tilde\alpha_{1,k}\le\bar\alpha_1(x_k)$. The last inequality is the ceiling
\eqref{4.23:eq-form-theorem-b} at $x=x_k$, which is admissible because $x_k\ge\gamma^{-2}$ along the
run. In particular $s_0\le(1+K_\Lambda)|B'|_\infty$, the last assertion of the lemma, and
$0<s_0<\rho_{\mathrm{abs}}$, so that $(B',b_\Lambda)$ lies in the bidisc
$\mathbb{D}(\rho_{\mathrm{abs}})$ of \eqref{4.23:eq-coupling-free-radii-a}.

\emph{Step 2: the sizes at the one point.} Since $(B',b_\Lambda)\in\mathbb{D}(\rho_{\mathrm{abs}})$,
Lemma~\ref{4.23:lem-bidisc-sizes} applies at the single point $(B',b_\Lambda)$. It provides the
objects of the construction of Definition~\ref{4.23:def-kept-member-construction}, in particular
$\mathsf{H}^C$, $\mathfrak{j}^C$, $\mathsf{v}$ and $\hat B$, and the bounds (A4)--(A6) of
\eqref{4.23:eq-bidisc-sizes-a} at $s=s_0$. Write $\mathcal{N}(\mathsf{H}^C)$ for
$\sup_\Delta\mathcal{N}^\Delta(\mathsf{H}^C)$, the norm taken at the kernel $U^0$. Then
\begin{align*}
\mathcal{N}(\mathsf{H}^C)&\le B_\sharp K_1^0s_0\le\tfrac12\varepsilon_{\alpha},\\
\sup\ell^3|\mathfrak{j}^C|&\le K_{\mathfrak{s}}B_\sharp K_1^0s_0\le\tfrac12\gamma^\flat,\\
|\log\mathsf{v}|_\infty&\le1.05K_uB_\sharp(1+K_1^0)s_0\le1.05\cdot10^{-2}.
\end{align*}
The variable $B$ satisfies $|B|_\infty\le\mathfrak{b}'_k$, and $\hat B$ arises from $B$ by the
rotation by $\mathsf{v}$ of the construction. Hence
\begin{equation*}
|\hat B-B|_\infty\le\bigl(e^{2|\log\mathsf{v}|_\infty}-1\bigr)\mathfrak{b}'_k
\le2.1|\log\mathsf{v}|_\infty\,\mathfrak{b}'_k .
\end{equation*}
Here we used $e^{2u}-1\le2.1u$ for $0\le u\le0.0105$, which follows from $e^{2u}-1\le2ue^{2u}$ and
$e^{0.021}<1.05$, together with $|\log\mathsf{v}|_\infty\le0.0105$ from the third bound above.
Consequently $|\hat B-B|_\infty\le2.1\cdot0.0105\,\mathfrak{b}'_k$, and therefore
\begin{equation*}
|\hat B|_\infty\le1.03\,\mathfrak{b}'_k,\qquad
|\hat B+B|_\infty\le|\hat B-B|_\infty+2|B|_\infty\le2.03\,\mathfrak{b}'_k.
\end{equation*}

\emph{Step 3: the kernel disc.} For $\tau\in\mathbb{C}$ put
\begin{equation*}
W(\tau):=e^{i\eta\tau\mathsf{H}^C}U^0,\qquad J(\tau):=\mathcal{J}(U^0)+\tau\mathfrak{j}^C,
\qquad
T:=\frac{\min\{\tfrac12\varepsilon_{\alpha},\gamma^\flat/(2K_{\mathfrak{s}})\}}{B_\sharp K_1^0s_0}.
\end{equation*}
By the definition of $\rho_{\mathrm{abs}}$ in \eqref{4.23:eq-coupling-free-radii-a},
\begin{equation*}
\rho_{\mathrm{abs}}\le\min\Bigl\{\frac{\varepsilon_{\alpha}}{2B_\sharp K_1^0},
\frac{\gamma^\flat}{2K_{\mathfrak{s}}B_\sharp K_1^0}\Bigr\},
\end{equation*}
and the right-hand side equals $Ts_0$. Hence, by \eqref{4.23:eq-kernel-disc-2},
$T\ge\rho_{\mathrm{abs}}/s_0>14$. Let $|\tau|\le T$. The configuration $W(\tau)$ is
$e^{i\eta A'}U^0$ with $A'=\tau\mathsf{H}^C$, and by Step~2
\begin{equation*}
\max\bigl\{\ell|\tau\mathsf{H}^C|,\ \ell^2|\nabla_{U^0}\tau\mathsf{H}^C|\bigr\}
\le T\mathcal{N}(\mathsf{H}^C)\le TB_\sharp K_1^0s_0\le\tfrac12\varepsilon_{\alpha}<\varepsilon_{\alpha},
\end{equation*}
where $TB_\sharp K_1^0s_0=\min\{\frac12\varepsilon_{\alpha},\gamma^\flat/(2K_{\mathfrak{s}})\}$ by the
definition of $T$. For the current, the estimate below uses the bound
$\sup\ell^3|\mathcal{J}(U^0)|\le\frac12\gamma^\flat$ on the current of the base, which is its first
summand; by Step~2 and $TK_{\mathfrak{s}}B_\sharp K_1^0s_0\le\frac12\gamma^\flat$,
\begin{equation*}
\sup\ell^3|J(\tau)|\le\tfrac12\gamma^\flat+TK_{\mathfrak{s}}B_\sharp K_1^0s_0\le\gamma^\flat.
\end{equation*}
The disc is built directly around the base $U^0$, at which $A'=0$, so no recombination of
exponentials enters. The class $\mathfrak{K}_{109}(\mathfrak{B}^C;\gamma^\flat)$ imposes no
criticality condition linking the current to the configuration: the current is arbitrary, subject
only to the bound on $\sup\ell^3|J|$. Consequently
\begin{equation}\label{4.23:eq-kernel-disc-3}
(W(\tau),J(\tau))\in\mathfrak{K}_{109}(\mathfrak{B}^C;\gamma^\flat)\qquad(|\tau|\le T).
\end{equation}
Its ends are $(U^0,\mathcal{J}(U^0))$ at $\tau=0$ and the pair $(W,J)$ of the construction at $\tau=1$.

\emph{Step 4: the variation along the disc.} Keep $\hat B$ fixed and put
$\Phi(\tau):=\mathsf{Q}[W(\tau),J(\tau)](\hat B,\hat B)$. The variable $\tau$ enters $A'=\tau\mathsf{H}^C$
linearly and $J(\tau)$ affinely. By \eqref{4.23:eq-kernel-disc-3} and the holomorphy of
$\mathsf{Q}[W',J']$ in $(W',J')$ on the class (Lemma~\ref{4.23:lem-shell-count-bound}), $\Phi$ is
holomorphic on $|\tau|<T$. By the bound \eqref{4.23:eq-shell-count-bound-a} of the same lemma,
\begin{equation*}
|\Phi(\tau)|\le12c_{\mathrm{tr}}K_{\mathsf{T}}^2|\hat B|_\infty^2M^d\mathsf{n}_{\mathrm{sh}}(k)
\le\mathfrak{q}_{\max}(k)\qquad(|\tau|<T).
\end{equation*}
The last inequality holds because $|\hat B|_\infty\le1.03\,\mathfrak{b}'_k$ (Step~2),
$12\cdot1.03^2<15$, and, by the definition of $\mathfrak{q}_{\max}(k)$,
\begin{equation*}
\mathfrak{q}_{\max}(k)\ge15c_{\mathrm{tr}}K_{\mathsf{T}}^2\mathfrak{b}'^2_kM^d\mathsf{n}_{\mathrm{sh}}(k).
\end{equation*}

The function $\Phi-\Phi(0)$ is therefore holomorphic on $|\tau|<T$, vanishes at $0$, and has
modulus at most $2\mathfrak{q}_{\max}(k)$ there. The Schwarz lemma, applied to
$\tau'\mapsto(\Phi(T\tau')-\Phi(0))/(2\mathfrak{q}_{\max}(k))$ on the unit disc, gives
$|\Phi(\tau)-\Phi(0)|\le2\mathfrak{q}_{\max}(k)|\tau|/T$ for $|\tau|<T$. Since $1<14<T$, we may take
$\tau=1$, and by the definition of $T$
\begin{equation}\label{4.23:eq-kernel-disc-4}
|\Phi(1)-\Phi(0)|\le\frac{2\mathfrak{q}_{\max}(k)}{T}
=4B_\sharp K_1^0\max\Bigl\{\frac{1}{\varepsilon_{\alpha}},\frac{K_{\mathfrak{s}}}{\gamma^\flat}\Bigr\}
s_0\,\mathfrak{q}_{\max}(k).
\end{equation}

\emph{Step 5: the rotation.} The form $\mathsf{Q}[U^0,\mathcal{J}(U^0)]$ is bilinear and symmetric, so
\begin{equation*}
\Phi(0)-\mathsf{Q}[U^0,\mathcal{J}(U^0)](B,B)=\mathsf{Q}[U^0,\mathcal{J}(U^0)](\hat B-B,\hat B+B).
\end{equation*}
By Step~2, $|\hat B-B|_\infty\le2.1|\log\mathsf{v}|_\infty\,\mathfrak{b}'_k$ and
$|\hat B+B|_\infty\le2.03\,\mathfrak{b}'_k$.
The pair $(U^0,\mathcal{J}(U^0))$ lies in $\mathfrak{K}_{109}(\mathfrak{B}^C;\gamma^\flat)$; it is the
value of \eqref{4.23:eq-kernel-disc-3} at $\tau=0$. Lemma~\ref{4.23:lem-shell-count-bound} applies at
this pair and gives
\begin{align*}
\bigl|\mathsf{Q}[U^0,\mathcal{J}(U^0)](\hat B-B,\hat B+B)\bigr|
&\le12\cdot2.1\cdot2.03\,c_{\mathrm{tr}}K_{\mathsf{T}}^2|\log\mathsf{v}|_\infty\,
\mathfrak{b}'^2_kM^d\mathsf{n}_{\mathrm{sh}}(k)\\
&\le3.5|\log\mathsf{v}|_\infty\,\mathfrak{q}_{\max}(k)\\
&\le3.5\cdot1.05K_uB_\sharp(1+K_1^0)s_0\,\mathfrak{q}_{\max}(k)\\
&\le3.7K_uB_\sharp(1+K_1^0)s_0\,\mathfrak{q}_{\max}(k).
\end{align*}
The second line uses $12\cdot2.1\cdot2.03=51.156\le3.5\cdot15$ and the lower bound on
$\mathfrak{q}_{\max}(k)$ recalled in Step~4. The third line is (A6) of
\eqref{4.23:eq-bidisc-sizes-a} at $s=s_0$, and the fourth is $3.5\cdot1.05=3.675\le3.7$.

\emph{Step 6: the sum.} By the construction \eqref{4.23:eq-kept-member-construction-a},
\begin{equation*}
\Phi(1)=\mathsf{Q}[W,J](\hat B,\hat B)=\tilde{\mathfrak{q}}[V](B',b_\Lambda)=\mathfrak{q}_C(\zeta;\varpi),
\end{equation*}
and by \eqref{4.23:eq-kept-member-construction-b},
$\mathfrak{q}_C((U,0);\varpi)=\mathsf{Q}[U^0,\mathcal{J}(U^0)](B,B)$. Write the difference
$\mathfrak{q}_C(\zeta;\varpi)-\mathfrak{q}_C((U,0);\varpi)$ as the sum of $\Phi(1)-\Phi(0)$ and
$\Phi(0)-\mathsf{Q}[U^0,\mathcal{J}(U^0)](B,B)$. The triangle inequality, together with
\eqref{4.23:eq-kernel-disc-4} and Step~5, gives \eqref{4.23:eq-kernel-disc-a} with $\vartheta^0$ as
defined there. The bound $s_0\le(1+K_\Lambda)|B'|_\infty$ was obtained in Step~1.
\end{proof}

\begin{remark}
The bound \eqref{4.23:eq-kernel-disc-a} has the shape of Ba{\l}aban's own bound
$O(1)B_3^2A_1^2p_1^2(g_k)|\mathbf{B}_0|(\alpha_{0,k}+\alpha_{1,k})$ in \cite{B-CMP122b}. Under the
comparison, the factor $O(1)B_3^2$ corresponds to $\vartheta_XC_{\mathfrak{q}}c_Q$, where
$C_{\mathfrak{q}}$ is the constant entering $\mathfrak{q}_{\max}(k)$;
$A_1^2p_1^2(g_k)$ corresponds to $\mathfrak{b}'^2_k$, and $|\mathbf{B}_0|$ corresponds to
$M^d\mathsf{n}_{\mathrm{sh}}(k)$. This comparison is not used in the proof.
\end{remark}

\begin{lemma}[Variation of the form: the data ray]\label{4.23:prf-data-ray}
In the setting of the form theorem, and under its hypotheses, let $b_\Lambda$ be the vector
supplied by hypothesis \textup{(H1)(a)} of the form theorem at the true
value of $B'$, let $s_0$ be as in Lemma~\ref{4.23:prf-kernel-disc}, and let $K_\Lambda$ be the
constant of the bound $s_0\le(1+K_\Lambda)|B'|_\infty$ established there. Then
\begin{equation}\label{4.23:eq-data-ray-1}
\begin{split}
\bigl|\mathfrak{q}_C(\zeta;\varpi)-\mathfrak{q}_C((U,0);\varpi)\bigr|
&\le \frac{2\,\mathfrak{q}_{\max}(k)\,s_0}{\rho_{\mathrm{abs}}}\\
&\le \frac{2(1+K_\Lambda)}{\rho_{\mathrm{abs}}}\,|B'|_\infty\,\mathfrak{q}_{\max}(k).
\end{split}
\end{equation}
Put
\begin{equation}\label{4.23:eq-data-ray-a}
\vartheta_X:=(1+K_\Lambda)\max\Bigl\{\vartheta^0,\ \frac{2}{\rho_{\mathrm{abs}}}\Bigr\},
\end{equation}
with $\vartheta^0$ the constant of \eqref{4.23:eq-kernel-disc-a}. With this constant, the bound
\eqref{4.23:eq-kernel-disc-a} of Lemma~\ref{4.23:prf-kernel-disc} and the bound
\eqref{4.23:eq-data-ray-1} each deliver the variation bound \textup{(X2)} of
\eqref{4.23:eq-form-theorem-a}.
\end{lemma}

\begin{proof}
Of the argument of Lemma~\ref{4.23:prf-kernel-disc} this proof uses only the quantity $s_0$ and
the inequality $s_0\le(1+K_\Lambda)|B'|_\infty$ established there.

By the auxiliary sizes on the coupling-free bidisc (Lemma~\ref{4.23:lem-bidisc-sizes}), by the
bound for the quadratic form and the shallow-shell count (Lemma~\ref{4.23:lem-shell-count-bound}),
and by the bound \textup{(X1)} of \eqref{4.23:eq-form-theorem-a}, the function
$\tilde{\mathfrak{q}}[V]$ is holomorphic on the coupling-free bidisc
$\mathbb{D}(\rho_{\mathrm{abs}})$ of \eqref{4.23:eq-coupling-free-radii-a}, and it satisfies
$|\tilde{\mathfrak{q}}[V]|\le\mathfrak{q}_{\max}(k)$ there.

Define, for a complex number $z$,
\begin{equation*}
\varphi(z):=\tilde{\mathfrak{q}}[V](zB',\,zb_\Lambda).
\end{equation*}
For $|z|<\rho_{\mathrm{abs}}/s_0$ the point $(zB',zb_\Lambda)$ lies in
$\mathbb{D}(\rho_{\mathrm{abs}})$. Since $z\mapsto(zB',zb_\Lambda)$ is linear, $\varphi$ is
holomorphic on the disc $\{|z|<\rho_{\mathrm{abs}}/s_0\}$, whose radius exceeds $14$; in
particular the disc contains $z=1$. On this disc the triangle inequality and the bound on
$\tilde{\mathfrak{q}}[V]$ give
\begin{equation*}
|\varphi(z)-\varphi(0)|\le|\varphi(z)|+|\varphi(0)|\le 2\,\mathfrak{q}_{\max}(k).
\end{equation*}
At $z=1$ the point is $(B',b_\Lambda)$, the true datum, and at $z=0$ it is the origin; by
\eqref{4.23:eq-kept-member-construction-b} these give $\varphi(1)=\mathfrak{q}_C(\zeta;\varpi)$
and $\varphi(0)=\mathfrak{q}_C((U,0);\varpi)$.

We now apply the Schwarz lemma. Its form here is the following: if $f$ is holomorphic on
$\{|z|<R\}$ and $|f-f(0)|\le A$ there, then $|f(z)-f(0)|\le A|z|/R$. Take $f=\varphi$,
$R=\rho_{\mathrm{abs}}/s_0$, $A=2\mathfrak{q}_{\max}(k)$ and $z=1$. This gives
\begin{equation*}
\bigl|\mathfrak{q}_C(\zeta;\varpi)-\mathfrak{q}_C((U,0);\varpi)\bigr|
\le\frac{2\,\mathfrak{q}_{\max}(k)\,s_0}{\rho_{\mathrm{abs}}},
\end{equation*}
which is the first inequality of \eqref{4.23:eq-data-ray-1}. The second follows from
$s_0\le(1+K_\Lambda)|B'|_\infty$.

It remains to show that each of the two bounds delivers \textup{(X2)} with the constant
\eqref{4.23:eq-data-ray-a}. By the definition of the maximum,
\begin{equation*}
s_0\vartheta^0\le(1+K_\Lambda)\vartheta^0|B'|_\infty\le\vartheta_X|B'|_\infty,
\qquad
\frac{2(1+K_\Lambda)}{\rho_{\mathrm{abs}}}\le\vartheta_X .
\end{equation*}
The first inequality, inserted in \eqref{4.23:eq-kernel-disc-a}, gives \textup{(X2)}. The second,
inserted in \eqref{4.23:eq-data-ray-1}, gives \textup{(X2)} as well. It is not asserted that the
two proofs agree on the constant.
\end{proof}

Along the ray, the datum on $\Lambda$ is $zb_\Lambda$, a fixed vector multiplied by a scalar. The
quantity $B'_\Lambda(zB')$ is never formed. For $z\ne1$ the point $(zB',zb_\Lambda)$ lies off the
critical graph $b=B'_\Lambda(B')$, on which the datum on $\Lambda$ is the value $B'_\Lambda(B')$
determined by $B'$. There $\tilde{\mathfrak{q}}[V]$ is an auxiliary function only: it
is not an object of Ba{\l}aban's construction.

Neither Lemma~\ref{4.23:prf-kernel-disc} nor the proof above evaluates \cite[(1.41)]{B-CMP122b}
away from the true value of $B'$, and no coupling-free reading of \cite[(1.41)]{B-CMP122b} is used.

The constants $\vartheta^0$ and $\rho_{\mathrm{abs}}$ involve none of the quantities of hypothesis
\textup{(H1)} of the form theorem. They are numbers fixed after $\kappa_1$ in the order of choice;
$\vartheta_X$ is fixed after $M$.

\begin{lemma}[Monotone hull of the variation bound]\label{4.23:prf-monotone-hull}
This lemma gives clauses \textup{(X3)} and \textup{(X5)} of the form theorem
(Theorem~\ref{4.23:thm-form-theorem}; see \eqref{4.23:eq-form-theorem-a}): items (a)--(c) below
constitute clause \textup{(X3)}, and item (d) constitutes clause \textup{(X5)}. Let
$\mathfrak{q}_{\max}(k)$ be
the bound on the quadratic form fixed in that theorem, and let $C_{\mathfrak{q}}$ be the constant
occurring in it. Let $r$ be the exponent in the defining bound $\check{R}_k<L(\log x_k)^r$ of the
window radius function $\check{R}_k$. Let $h=k-N(g_k)$ be the bottom level of the window, and let
$c_R$ be the constant in the radius-ratio bound $\check{R}_{h+1}\le c_R\check{R}_k$ of the window.
Let $A_1$ and $p_1$ be the amplitude and the exponent in
$\mathfrak{b}'_k=\delta'_k/g_k=A_1(\log x_k)^{p_1}$, where $\delta'_k$ is as defined in
\cite{B-CMP122a}. Let $\vartheta_X$ be the constant of \eqref{4.23:eq-data-ray-a}. Let $c_Q$ be
as in Theorem~\ref{4.23:thm-form-theorem}, and let $\bar\alpha_1$ be the amplitude ceiling of
the amplitude bound for the source tube, the function of the inverse coupling recalled in the
proof of (b) below; it bounds the radius $\tilde\alpha_{1,k}$ of the source tube by
$\tilde\alpha_{1,k}\le\bar\alpha_1(x_k)$. Put
\begin{equation}\label{4.23:eq-monotone-hull-a}
\begin{gathered}
\bar{\mathfrak{q}}_{\max}(x):=C_{\mathfrak{q}}A_1^2M^d(\log x)^{2p_1}
\bigl(L(\log x)^r+1\bigr)\bigl(110L^2c_R(\log x)^r\bigr)^d,\qquad x>1,\\
\Sigma^{\mathfrak{q}}(x):=\sup_{x'\ge\max\{x,e\}}\,\vartheta_X\,c_Q\,\bar\alpha_1(x')\,
\bar{\mathfrak{q}}_{\max}(x'),\qquad x\in\mathbb{R}.
\end{gathered}
\end{equation}
Then the following hold.
\begin{enumerate}
\item[(a)] For every level $k$, $\mathfrak{q}_{\max}(k)\le\bar{\mathfrak{q}}_{\max}(x_k)$.
\item[(b)] $\Sigma^{\mathfrak{q}}$ is finite and non-increasing on $\mathbb{R}$, and
$\Sigma^{\mathfrak{q}}(x)\to0$ as $x\to\infty$.
\item[(c)] For every level $i\ge k_W$ we have $x_{k_W}\ge x_i-2c_2$, and hence
$\Sigma^{\mathfrak{q}}(x_{k_W})\le\Sigma^{\mathfrak{q}}(x_i-2c_2)$.
\item[(d)] Each of the constants $\rho_{\mathrm{abs}}$ (the radius of the coupling-free bidisc
of clause \textup{(X1)} of Theorem~\ref{4.23:thm-form-theorem}), $\vartheta^0$ (the constant of
that name in Theorem~\ref{4.23:thm-form-theorem}), $\vartheta_X$ and
$C_{\mathfrak{q}}$, and the ceiling of the form theorem, is a number determined by
$(d,L,\kappa_1,M)$. The bound $\mathfrak{q}_{\max}(k)$ depends on the level only through
the window length $N=N(g_k)$, the radius $\check{R}_{h+1}$, and $\mathfrak{b}'_k$, each taken as
its own function of the
coupling. The core $\mathsf{C}$ enters only through two conditions: the base configuration $U^0$
of Theorem~\ref{4.23:thm-form-theorem} lies in the class
$\mathfrak{K}_{109}(\mathfrak{B};\gamma_0)$ of pairs $(W,J)$ of that theorem's base-configuration
hypothesis, at the set $\mathfrak{B}$ and the bound $\gamma_0$ fixed there, and
$V''(c)\in G$ for every point $c$ of the core $\mathsf{C}$. The support of the localising weight
$\hat\chi$ of the form $\mathsf{Q}$ is disjoint from $\mathsf{C}$.
\end{enumerate}
\end{lemma}

\begin{proof}
\emph{Proof of (a).} By Theorem~\ref{4.23:thm-form-theorem} (see
\eqref{4.23:eq-form-theorem-a}), the bound satisfies
\begin{equation*}
\mathfrak{q}_{\max}(k)\le C_{\mathfrak{q}}(\mathfrak{b}'_k)^2M^d\bigl(N(g_k)+1\bigr)
\bigl(110L\,\check{R}_{h+1}\bigr)^d .
\end{equation*}
The right-hand side depends on the level through three quantities only, and each factor is
non-decreasing in the non-negative quantity it contains. We bound each of them in turn.

First, by the definition of $\delta'_k$ in \cite{B-CMP122a},
\begin{equation*}
\mathfrak{b}'_k=\delta'_k/g_k=A_1(\log x_k)^{p_1},
\end{equation*}
and hence $(\mathfrak{b}'_k)^2=A_1^2(\log x_k)^{2p_1}$.

Second, the bound on the window length gives $N(g_k)\le L(\log x_k)^r$, and hence
$N(g_k)+1\le L(\log x_k)^r+1$.

Third, the radius-ratio bound of the window gives
$\check{R}_{h+1}\le c_R\check{R}_k$. The defining bound of the window radius function gives
$\check{R}_k<L(\log x_k)^r$. Together,
\begin{equation*}
110L\,\check{R}_{h+1}<110L^2c_R(\log x_k)^r .
\end{equation*}
We replace the three level-dependent quantities in the above bound of
$\mathfrak{q}_{\max}(k)$ by these bounds. Since each factor is non-decreasing in the quantity
replaced, the result is a bound of $\mathfrak{q}_{\max}(k)$ by the right-hand side of the first
line of \eqref{4.23:eq-monotone-hull-a} at $x=x_k$. The three bounds just used are statements
about powers of $\log x_k$, taken in the regime $x_k>1$, which is the domain of
$\bar{\mathfrak{q}}_{\max}$ in \eqref{4.23:eq-monotone-hull-a}. This proves (a).

\emph{Proof of (b).} Let $f(x)=\vartheta_Xc_Q\bar\alpha_1(x)\bar{\mathfrak{q}}_{\max}(x)$ on
$[e,\infty)$. The amplitude bound for the source tube gives the amplitude ceiling as
\begin{equation*}
\bar\alpha_1(x)=(1+2\beta)(1+\varrho)C_0\,x^{-1/2}(\log x)^q ,
\end{equation*}
with $\beta$ the schedule constant of Definition~\ref{4.4:def-post-step-space}.
Put $a=q+2p_1+rd$ and
\begin{equation*}
c_f=\vartheta_Xc_Q(1+2\beta)(1+\varrho)C_0\,C_{\mathfrak{q}}A_1^2M^d(110L^2c_R)^d .
\end{equation*}
Multiplying out \eqref{4.23:eq-monotone-hull-a} gives
\begin{equation*}
f(x)=c_f\bigl(L\,x^{-1/2}(\log x)^{a+r}+x^{-1/2}(\log x)^{a}\bigr).
\end{equation*}
On $[e,\infty)$ each summand is a product of continuous functions of $x$, so $f$ is continuous
there. For every real $b$, the substitution $x=e^t$ gives
$x^{-1/2}(\log x)^b=e^{-t/2}t^b$, and this tends to $0$ as $t\to\infty$. Hence
$f(x)\to0$ as $x\to\infty$.

Now fix $x\in\mathbb{R}$. Choose $y_1\ge\max\{x,e\}$ with $|f|\le1$ on $[y_1,\infty)$. By
continuity, $f$ is bounded on the compact interval $[\max\{x,e\},y_1]$. So $f$ is bounded on
$[\max\{x,e\},\infty)$, and $\Sigma^{\mathfrak{q}}(x)$ is finite. If $x\le y$, the supremum
defining $\Sigma^{\mathfrak{q}}(y)$ is taken over a subset of the set defining
$\Sigma^{\mathfrak{q}}(x)$; hence $\Sigma^{\mathfrak{q}}(y)\le\Sigma^{\mathfrak{q}}(x)$.

Finally, let $\varepsilon>0$ and choose $y_\varepsilon\ge e$ with $|f|<\varepsilon$ on
$[y_\varepsilon,\infty)$. Then $\Sigma^{\mathfrak{q}}(x)\le\varepsilon$ for
$x\ge y_\varepsilon$, so $\Sigma^{\mathfrak{q}}(x)\to0$.

This convergence concerns the bound on the quadratic form alone. It is not
a convergence statement for the functional $V$ whose second kept member is $\mathfrak{q}[V]$.

\emph{Proof of (c).} By the coupling flow (Theorem~\ref{4.1:thm-clause-zero}), the inverse
couplings satisfy $\underline{\lambda}(j-i)-2c_2\le x_i-x_j$ for all levels $j\ge i$. Apply this
with $k_W$ in place of $i$ and $i$ in place of $j$. Since $i\ge k_W$, this gives
\begin{equation*}
x_{k_W}\ge x_i+\underline{\lambda}(i-k_W)-2c_2\ge x_i-2c_2 .
\end{equation*}
Since $\Sigma^{\mathfrak{q}}$ is non-increasing on $\mathbb{R}$ by (b), we get
$\Sigma^{\mathfrak{q}}(x_{k_W})\le\Sigma^{\mathfrak{q}}(x_i-2c_2)$.

\emph{Proof of (d).} We take the constants one at a time.
\begin{itemize}
\item Inspection of \eqref{4.23:eq-form-theorem-a}, for $\rho_{\mathrm{abs}}$, $\vartheta^0$,
$C_{\mathfrak{q}}$ and the ceiling of the form theorem, and of \eqref{4.23:eq-data-ray-a}, for
$\vartheta_X$, shows that their definitions involve only $d$, $L$, $\kappa_1$ and $M$.
\item In $\mathfrak{q}_{\max}(k)$ the quantities $N$, $\check{R}_{h+1}$ and
$\mathfrak{b}'_k$ appear as their own functions of the coupling, as in the proof of (a).
\item By the list of hypotheses of Theorem~\ref{4.23:thm-form-theorem}, the core $\mathsf{C}$
enters the hypotheses of the form theorem only through the
condition of its base-configuration hypothesis that $U^0$ lies in
$\mathfrak{K}_{109}(\mathfrak{B};\gamma_0)$, and through the condition $V''(c)\in G$ for every
point $c$ of $\mathsf{C}$.
\item The weight $\hat\chi$ of the form $\mathsf{Q}$ of
Lemma~\ref{4.23:lem-shell-count-bound} is supported off $\mathsf{C}$.
\end{itemize}
This proves (d).
\end{proof}

\begin{lemma}[The form at the trivial splitting]\label{4.23:prf-trivial-splitting}
The following are clauses (X4) and (X0) of the form theorem (Theorem~\ref{4.23:thm-form-theorem},
display \eqref{4.23:eq-form-theorem-a}), each proved under the hypotheses named in it.
For every $V\in\mathfrak{V}_{\mathbb{R}}$ let $Y=V_\Lambda(V)$, $U_0$, $U^0$, $\mathsf{H}^0$,
$\mathfrak{u}^0$, $\mathfrak{f}$, $\mathsf{v}$ and $(W,J)$
be the objects of the construction of the second kept member
(Definition~\ref{4.23:def-kept-member-construction}, displays
\eqref{4.23:eq-kept-member-construction-a} and \eqref{4.23:eq-kept-member-construction-b}); the
hypotheses below are assumed for every such $V$.
Let $\mathsf{T}^{\mathsf{L}}$ and $\mathsf{T}^{\sharp}$ be the operators of the linearised localised
minimiser map (Lemma~\ref{4.23:lem-linearised-minimiser}). Let
\begin{equation*}
\mathsf{Q}[W',J'](a,a')=-\tfrac12\sum_p\eta^d\hat\chi\,\operatorname{tr}[(D\mathsf{T}a)(D\mathsf{T}a')],
\end{equation*}
where $\mathsf{T}$ stands for either $\mathsf{T}^{\mathsf{L}}$ or $\mathsf{T}^{\sharp}$, and let
$\mathfrak{q}_C$, $\tilde{\mathfrak{q}}[V]$, $\mathfrak{q}[V]$ be as in the form theorem
(Theorem~\ref{4.23:thm-form-theorem}); for $\zeta$ presented by an admissible splitting as
$(U'U,\mathbf{J})$ with $U$ real, $(U,0)$ denotes, as there, the real point of the tube with
configuration $U$ and current $0$. Put
\begin{equation}\label{4.23:eq-trivial-splitting-1}
r_{\mathbb{R}}:=3.3\,(1+K_\Lambda)\,c_Q\,\tilde\alpha_{1,k}.
\end{equation}
Under the ceiling \eqref{4.23:eq-form-theorem-b} of the form theorem,
$(1+K_\Lambda)c_Q\tilde\alpha_{1,k}\le\rho_{\mathrm{abs}}/14$, hence
$r_{\mathbb{R}}\le3.3\,\rho_{\mathrm{abs}}/14<\rho_{\mathrm{abs}}$, so real data of size at most
$r_{\mathbb{R}}$ lie in the coupling-free bidisc $\mathbb{D}(\rho_{\mathrm{abs}})$ of
\eqref{4.23:eq-coupling-free-radii-a}.
\begin{enumerate}
\item[(X4)] Assume hypotheses (H2)--(H6) of Theorem~\ref{4.23:thm-form-theorem}. Then at the
origin $(B',b)=(0,0)$ of the bidisc the pair $(W,J)$ is $(U^0,\mathcal{J}(U^0))$; it is real under the
reality condition of the construction of the second kept member, and the current correction
$\mathfrak{j}=\tilde G_0\mathcal{J}$ vanishes. For $\mathsf{T}\in\{\mathsf{T}^{\mathsf{L}},\mathsf{T}^{\sharp}\}$
and every real $\mathfrak{g}$-valued $a$ the field $\mathsf{T}a$ at $(U^0,\mathcal{J}(U^0))$ is
$\mathfrak{g}$-valued; for every real $\mathfrak{g}$-valued $B$ the number
$\mathsf{Q}[U^0,\mathcal{J}(U^0)](B,B)$ is real and $\le0$. If, moreover, all hypotheses of
Theorem~\ref{4.23:thm-form-theorem} hold, so that its clauses (X2) and (X3) are available, then
\begin{equation}\label{4.23:eq-trivial-splitting-2}
\bigl|\operatorname{Re}\mathfrak{q}_C(\zeta;\varpi)-\mathfrak{q}_C((U,0);\varpi)\bigr|\le\Sigma^{\mathfrak{q}},
\qquad
\bigl|\operatorname{Im}\mathfrak{q}_C(\zeta;\varpi)\bigr|\le\Sigma^{\mathfrak{q}},
\end{equation}
with $\Sigma^{\mathfrak{q}}$ as in \eqref{4.23:eq-monotone-hull-a}.
\item[(X0)] Let $\mathsf{T}=\mathsf{T}^{\mathsf{L}}$. Assume the hypotheses of
Theorem~\ref{4.23:thm-form-theorem}, in particular (H1)(a) and (H2)--(H6); assume further clauses
(b) and (c), written (H1)(b) and (H1)(c), of the analytic solvability of the equation for the
$\Lambda$-datum (Definition~\ref{4.23:def-lambda-datum}), of which hypothesis (H1)(a) of that
theorem is clause (a); that statement is printed without proof in
\cite[p.~367, (1.41)]{B-CMP122b}; cited as printed. Assume also hypothesis (H8), consisting of:
\begin{itemize}
\item[(H8)(a)] for $\mathfrak{B}\in\{\mathfrak{B}^0,\mathfrak{B}^C\}$, the real base
$\mathbb{U}\in\{U_0,U^0\}$ of $\mathfrak{B}$, $\mathsf{D}_0:=M^{\mathfrak{B}}(\mathbb{U})$ (here and below
$M^{\mathfrak{B}}$, $M_j$, $M$ and $M'$ denote block-averaging operations of Ba{\l}aban's
construction, not the block size $M$) and every real
$\mathfrak{f}'$ with $|\mathfrak{f}'|_\infty\le K^0_1r_{\mathbb{R}}$, the configuration
\begin{equation*}
e^{i\eta\mathbb{H}[\mathfrak{B};(\mathbb{U},\mathcal{J}(\mathbb{U}));\mathbf{1};(0,\mathfrak{f}')]}\mathbb{U}
\end{equation*}
is the Landau representative of the minimal orbit of $e^{i\mathfrak{f}'}\mathsf{D}_0$, and its
transform by $\mathfrak{u}^{-1}$, where $\mathfrak{u}$ is $\mathfrak{u}[\,\cdot\,;\mathbb{U};\mathfrak{B}]$
of the construction evaluated at this configuration, is Ba{\l}aban's axial minimiser for the data
$e^{i\mathfrak{f}'}\mathsf{D}_0$;
\item[(H8)(b)] the relation \cite[(181)]{B-CMP102b}, which enters below in the identification of $W$
in (i) and in the covariance of $U_{k,C}$ in (v);
\item[(H8)(c)] hypothesis (H6) with $(V,Y)$ replaced by every real datum $(e^{iB'}V,e^{ib}Y)$ with
$|B'|_\infty,|b|_\infty\le r_{\mathbb{R}}$.
\end{itemize}
Assume finally that $\rho_{\mathrm{abs}}\le10^{-2}$, with $\rho_{\mathrm{abs}}$ the radius of
\eqref{4.23:eq-coupling-free-radii-a}.
Then:
\begin{itemize}
\item[(i)] for real $(B',b)$ with $|B'|_\infty,|b|_\infty\le r_{\mathbb{R}}$ and
$\mathsf{D}:=(e^{iB'}V,e^{ib}Y)$,
\begin{equation}\label{4.23:eq-trivial-splitting-3}
\tilde{\mathfrak{q}}[V](B',b)=\mathsf{f}^{\mathrm{pr}}(\mathsf{D};\varpi)
:=-\tfrac12\sum_p\eta^d\hat\chi(x(p))\,\bigl|(DH''_{1,k,C}[\mathsf{D}]B)(p)\bigr|^2 ,
\end{equation}
a function of $\mathsf{D}$ and $\varpi$ alone; here $B$ is the fluctuation variable among the real
variables $\varpi$, and $H''_{1,k,C}[\mathsf{D}]$ is the operator of \cite[(1.71)]{B-CMP122b} at the
data $\mathsf{D}$;
\item[(ii)] $\mathfrak{q}_C(\zeta;\varpi)$ does not depend on the admissible splitting of $\zeta$;
\item[(iii)] $\mathfrak{q}_C(\cdot;\varpi)$ is holomorphic on $\mathfrak{S}_{\mathrm{src}}$, depends on
$\zeta$ only through $\mathsf{V}=e^{iB'}V$, the configuration presented by every admissible
splitting of $\zeta$, and not on $\mathbf{J}$;
\item[(iv)] at real $\zeta$, $\mathfrak{q}_C(\zeta;\varpi)$ is the form of \cite[(1.71)]{B-CMP122b},
namely $\mathsf{f}^{\mathrm{pr}}((V,V_\Lambda(V));\varpi)$;
\item[(v)] for every $G$-valued gauge transformation $g$,
$\mathfrak{q}_C(\zeta^g;\varpi^g)=\mathfrak{q}_C(\zeta;\varpi)$.
\end{itemize}
\end{enumerate}
\end{lemma}

\begin{proof}
Write $\mathfrak{g}^c$ for the complexification of $\mathfrak{g}$, as in
Notation~\ref{4.4:not-kernel-symmetrisation}, and use the convention in which
$\mathfrak{g}$ consists of Hermitian matrices and group elements are written $e^{iX}$.

\emph{Clause (X4).} At $(B',b)=(0,0)$ the construction gives the pair $(U^0,\mathcal{J}(U^0))$; it is
real under the reality condition of the construction of the second kept member, and by hypothesis
(H6) the current correction $\mathfrak{j}=\tilde G_0\mathcal{J}$ vanishes.

Let $a$ be real and $\mathfrak{g}$-valued, and let $t$ be real. At the real pair
$(U^0,\mathcal{J}(U^0))$ the datum $(0,ta)$ is of the kind treated in clause (d) of the theorem on
the localised small solution, which holds at $\mathfrak{B}^C$ by hypothesis (H2). We do not use any
identification of this solution with an object of Ba{\l}aban's papers. The conjugation
$X\mapsto\bar X$ of $\mathfrak{g}^{c}$, the real structure whose fixed set is $\mathfrak{g}$,
maps the small solution of the problem of clause (a) of that theorem, at the real pair
$(U^0,\mathcal{J}(U^0))$ and at a real datum, to a small solution of the same problem. By the
uniqueness asserted in clause (a) it maps that solution to itself. Hence
$\mathbb{H}[\mathfrak{B}^C;(U^0,\mathcal{J}(U^0));\mathbf{1};(0,ta)]$ is $\mathfrak{g}$-valued for every
real $t$, and so is its derivative at $t=0$, which is $\mathsf{T}^{\mathsf{L}}a$ by the definition in
Lemma~\ref{4.23:lem-linearised-minimiser}. For $\mathsf{T}^{\sharp}$ the same conclusion holds
because $\mathsf{T}^{\sharp}$ is a Neumann series of operators that are real at a real pair.

Since $\operatorname{Ad}_{U^0}$ preserves $\mathfrak{g}$, for real $\mathfrak{g}$-valued $B$ the field
$A:=D_{U^0}\mathsf{T}B$ is $\mathfrak{g}$-valued. In the Hermitian convention
$\operatorname{tr}A(p)^2=|A(p)|^2\ge0$, and $\hat\chi\ge0$. Each summand of
$\mathsf{Q}[U^0,\mathcal{J}(U^0)](B,B)$ is therefore real and $\le0$, and so is the sum. This agrees
with \cite[\S1]{B-CMP122b}, where the corresponding term is bounded by $0$; that statement is not
used here.

By \eqref{4.23:eq-kept-member-construction-b} the value $\mathfrak{q}_C((U,0);\varpi)$ is the value of
$\tilde{\mathfrak{q}}[V]$ at $(0,0)$, that is, the form built on $\mathsf{Q}$ at the real pair
$(U^0,\mathcal{J}(U^0))$. By what was just shown it is real. Hence
\begin{align*}
\bigl|\operatorname{Re}\mathfrak{q}_C(\zeta;\varpi)-\mathfrak{q}_C((U,0);\varpi)\bigr|
&\le\bigl|\mathfrak{q}_C(\zeta;\varpi)-\mathfrak{q}_C((U,0);\varpi)\bigr|,\\
\bigl|\operatorname{Im}\mathfrak{q}_C(\zeta;\varpi)\bigr|
=\bigl|\operatorname{Im}\bigl(\mathfrak{q}_C(\zeta;\varpi)-\mathfrak{q}_C((U,0);\varpi)\bigr)\bigr|
&\le\bigl|\mathfrak{q}_C(\zeta;\varpi)-\mathfrak{q}_C((U,0);\varpi)\bigr|,
\end{align*}
and the right-hand side is at most $\Sigma^{\mathfrak{q}}$ by clauses (X2) and (X3) of
\eqref{4.23:eq-form-theorem-a}. This is \eqref{4.23:eq-trivial-splitting-2}.

\emph{Clause (X0)(i): real points.} Let $(B',b)$ be real with $|B'|_\infty,|b|_\infty\le r_{\mathbb{R}}$,
and $\mathsf{D}:=(e^{iB'}V,e^{ib}Y)$. The $\Lambda$-component $b$ is free here. No step of (i) uses
the equation that determines the $\Lambda$-datum.

The datum $\mathfrak{d}^0$ of the construction satisfies
$|\mathfrak{d}^0|_\infty\le r_{\mathbb{R}}\le K^0_1r_{\mathbb{R}}$. So (H8)(a) at $\mathfrak{B}^0$ applies
to $\mathfrak{f}'=\mathfrak{d}^0$. Here $e^{i\mathfrak{d}^0}\mathsf{D}_0=\mathsf{D}$, since
$U_0=U_{k,C}(V,Y)$ has block averages $\mathsf{D}_0=M^{\mathfrak{B}^0}(U_0)=(V,Y)$ and $\mathfrak{d}^0$
carries $B'$ off $\Lambda$ and $b$ on $\Lambda_C$. Hence (H8)(a) gives
\begin{equation*}
\bigl(e^{i\eta\mathsf{H}^0}U_0\bigr)^{(\mathfrak{u}^0)^{-1}}=U_{k,C}(\mathsf{D}),
\end{equation*}
Ba{\l}aban's axial minimiser. By \cite[(11)]{B-CMP98} the averages of $U_{k,C}(\mathsf{D})$ are the
$(\mathfrak{u}^0)^{-1}$-transform of $M_j(e^{i\eta\mathsf{H}^0}U_0)$.

Let $\mathfrak{v}$ equal $\mathfrak{u}^0$ off the core $\mathsf{C}$ and $1$ on $\mathsf{C}$. Let
$\mathsf{D}^{\mathrm{pr}}$ equal $M'(U_{k,C}(\mathsf{D}))$ on $\Omega''^{\sim2}_{h+1}$ and $V''$ on
$\mathsf{C}$; these are Ba{\l}aban's data of \cite[(1.19)]{B-CMP122b}. The definition of
$\mathfrak{f}$ in \eqref{4.23:eq-kept-member-construction-a} then gives, on every bond $c$ of
$\mathfrak{B}^C$ with initial point $c_-$ and final point $c_+$,
\begin{equation*}
e^{i\mathfrak{f}(c)}M(U^0)(c)=\mathfrak{v}(c_-)\,\mathsf{D}^{\mathrm{pr}}(c)\,\mathfrak{v}(c_+)^{-1}.
\end{equation*}
Three kinds of bond occur. On bonds carried by $U_0$ the twist cancels against the block field
$\mathfrak{b}_j$ of hypothesis (H4). On the core both sides reduce to $V''$. On a bond crossing
into the core a single one-sided factor of $\mathfrak{u}^0$ remains, and the identity is an algebraic
identity at $V''$, which is $G$-valued, being one of the real variables $\varpi$ of
\cite[(1.71)]{B-CMP122b}.

By (A3) of \eqref{4.23:eq-bidisc-sizes-a}, $|\mathfrak{f}|_\infty\le K^0_1r_{\mathbb{R}}$. So (H8)(a) at
$\mathfrak{B}^C$ applies, and together with (H8)(b) it gives
$W=(U^0_{k,C}[\mathsf{D}])^{\mathsf{v}}$. Here $\mathsf{v}$ is $G$-valued by clause (iii), read at real
data, of the statement on the gauge factor $\mathsf{v}$ of the construction, whose clause (ii) is
hypothesis (H5).

Further, $J=\mathcal{J}(W)$. This follows from clause (i) of the free-current statement for the
localised small solution, since the operator of \cite[(3.11)]{B-CMP109} annihilates $\mathsf{H}^C$
at the real critical point, $\mathsf{H}^C$ being the localised small solution on $\mathfrak{B}^C$ of
the construction \eqref{4.23:eq-kept-member-construction-a}.

We now apply two clauses of the theorem on the localised small solution. Clause (c) is applied
at the $G$-valued $\mathsf{v}$. Clause (d) is applied at the base $U^0_{k,C}[\mathsf{D}]$, which lies
in the admissible class by (H8)(c). Combined with \eqref{4.23:eq-linearised-minimiser-b}, they give
\begin{equation*}
\mathsf{T}^{\mathsf{L}}[W,J]\operatorname{Ad}_{\mathsf{v}}
=\operatorname{Ad}_{\mathsf{v}}H''_{1,k,C}[\mathsf{D}].
\end{equation*}
The covariant derivative is covariant under $\operatorname{Ad}_{\mathsf{v}}$, and the trace is
invariant. So the form $\tilde{\mathfrak{q}}[V](B',b)$ equals the form built with
$H''_{1,k,C}[\mathsf{D}]$ and the covariant derivative at $U^0_{k,C}[\mathsf{D}]$. In the Hermitian
convention $\operatorname{tr}X^2=|X|^2$ for $\mathfrak{g}$-valued $X$; this gives
\eqref{4.23:eq-trivial-splitting-3}. The right-hand side depends on $\mathsf{D}$ and $\varpi$ only.

\emph{Clause (X0)(ii): independence of the splitting.} Let $(\zeta,U_1)$ and $(\zeta,U_2)$ be
admissible splittings, with data $V_i$, $Y_i$, $B'_i$ and $b_i:=B'_\Lambda(B'_i;V_i)$; recall
$\mathfrak{q}[V_i](B'_i)=\tilde{\mathfrak{q}}[V_i](B'_i,b_i)$. Both splittings present the same
configuration $\mathsf{V}$:
\begin{equation*}
e^{iB'_1}V_1=\mathsf{V}=e^{iB'_2}V_2 .
\end{equation*}
By (H1)(c), a clause of the analytic solvability of the equation for the $\Lambda$-datum
(printed without proof in \cite[p.~367, (1.41)]{B-CMP122b}; cited as printed),
\begin{equation*}
e^{ib_1}Y_1=V_\Lambda(\mathsf{V})=e^{ib_2}Y_2 .
\end{equation*}
Hence $V_2V_1^{-1}=:e^{i\Xi}$ and $Y_2Y_1^{-1}=:e^{i\Xi_\Lambda}$ are $G$-valued. Here $\Xi$ and
$\Xi_\Lambda$ are real, being logarithms of $G$-valued elements near the identity.

Each of $|B'_i|_\infty,|b_i|_\infty$ is at most $(1+K_\Lambda)c_Q\tilde\alpha_{1,k}$: for $B'_i$ by
the admissibility of the splitting, and for $b_i$ by the bound $|b_i|_\infty\le K_\Lambda|B'_i|_\infty$
of (H1)(a), printed without proof in \cite[p.~367, (1.41)]{B-CMP122b}; cited as printed. This
number is at most $\rho_{\mathrm{abs}}/14$ by the ceiling \eqref{4.23:eq-form-theorem-b} of the form
theorem. The estimate near the identity on which (A6) of Lemma~\ref{4.23:lem-bidisc-sizes} rests,
namely that $\frac1i\log(e^{iX}e^{-iX'})$ has size at most $1.05(|X|_\infty+|X'|_\infty)$ when $X$
and $X'$ have size below $10^{-2}$, then gives
\begin{equation*}
|\Xi|_\infty\le1.05(|B'_1|_\infty+|B'_2|_\infty)\le2.1(1+K_\Lambda)c_Q\tilde\alpha_{1,k}
\le\frac{2.1\,\rho_{\mathrm{abs}}}{14}\le\frac{\rho_{\mathrm{abs}}}{6},
\end{equation*}
and likewise for $\Xi_\Lambda$, with $|b_1|_\infty+|b_2|_\infty$ in place of
$|B'_1|_\infty+|B'_2|_\infty$.

In the display above and in what follows that estimate is applied only to factors of size
less than $\rho_{\mathrm{abs}}$, hence, by the assumption $\rho_{\mathrm{abs}}\le10^{-2}$ of (X0), to
factors of size below its threshold $10^{-2}$.

On the bidisc $\mathbb{D}(\rho_{\mathrm{abs}}/4)$ define
\begin{align*}
f_1(B',b)&:=\tilde{\mathfrak{q}}[V_1](B',b),\\
f_2(B',b)&:=\tilde{\mathfrak{q}}[V_2]\Bigl(\tfrac1i\log\bigl(e^{iB'}e^{-i\Xi}\bigr),\,
\tfrac1i\log\bigl(e^{ib}e^{-i\Xi_\Lambda}\bigr)\Bigr).
\end{align*}
Both are holomorphic by the lemma on auxiliary sizes on the coupling-free bidisc
(Lemma~\ref{4.23:lem-bidisc-sizes}). For $f_2$ this uses the size of its arguments:
\begin{equation*}
1.05\bigl(\tfrac14+\tfrac16\bigr)\rho_{\mathrm{abs}}<\rho_{\mathrm{abs}} .
\end{equation*}

Consider the real ball of radius $(1+K_\Lambda)c_Q\tilde\alpha_{1,k}$, which is at most
$\rho_{\mathrm{abs}}/14<\rho_{\mathrm{abs}}/4$. On it both functions equal
$\mathsf{f}^{\mathrm{pr}}((e^{iB'}V_1,e^{ib}Y_1);\varpi)$ by (i). For $f_1$ this is (i) at $V_1$.
For $f_2$, the arguments are real, of size at most
\begin{equation*}
1.05\cdot3.1\,(1+K_\Lambda)c_Q\tilde\alpha_{1,k}\le r_{\mathbb{R}}
=3.3\,(1+K_\Lambda)c_Q\tilde\alpha_{1,k},
\end{equation*}
and, since $e^{-i\Xi}V_2=V_1$ and $e^{-i\Xi_\Lambda}Y_2=Y_1$,
\begin{equation*}
\Bigl(e^{i\frac1i\log(e^{iB'}e^{-i\Xi})}V_2,\ e^{i\frac1i\log(e^{ib}e^{-i\Xi_\Lambda})}Y_2\Bigr)
=(e^{iB'}V_1,\,e^{ib}Y_1).
\end{equation*}

A holomorphic function on a convex domain that vanishes on an open piece of the real slice is
identically zero. Applied to $f_1-f_2$, this gives $f_1\equiv f_2$ on $\mathbb{D}(\rho_{\mathrm{abs}}/4)$.
This bidisc contains $(B'_1,b_1)$. From $e^{iB'_1}V_1=e^{iB'_2}V_2$ we get
$e^{iB'_1}e^{-i\Xi}=e^{iB'_1}V_1V_2^{-1}=e^{iB'_2}$, so
$\frac1i\log(e^{iB'_1}e^{-i\Xi})=B'_2$; likewise the second argument is $b_2$. Evaluating at
$(B'_1,b_1)$ therefore gives
\begin{equation*}
\mathfrak{q}[V_1](B'_1)=\tilde{\mathfrak{q}}[V_2](B'_2,b_2)=\mathfrak{q}[V_2](B'_2).
\end{equation*}
The margin used is taken in the coupling-free bidisc, not in the ball of \cite[(1.41)]{B-CMP122b}.

\emph{Clause (X0)(iii): holomorphy.} Near a point $\zeta_0=(U'_0U_0,\mathbf{J}_0)$ every $\zeta$
splits over the same $U_0$, because the clauses defining the tube are open. In such a splitting:
\begin{itemize}
\item $B'$ is holomorphic in $\mathbf{U}$ by hypothesis (H4);
\item $b_\Lambda$ is holomorphic in $B'$, on the open ball on which (H1)(b) holds, by (H1)(b)
(printed without proof in \cite[p.~367, (1.41)]{B-CMP122b}; cited as printed);
\item $\tilde{\mathfrak{q}}$ is holomorphic in both by Lemma~\ref{4.23:lem-bidisc-sizes}.
\end{itemize}
Hence $\mathfrak{q}_C$ is holomorphic on $\mathfrak{S}_{\mathrm{src}}$. By construction it depends on
$\zeta$ only through $\mathsf{V}$, and not on $\mathbf{J}$.

\emph{Clause (X0)(iv): real $\zeta$.} Split a real $\zeta$ with $U'=1$. By
\eqref{4.23:eq-kept-member-construction-b}, $\mathfrak{q}_C(\zeta;\varpi)=\tilde{\mathfrak{q}}[V](0,0)$.
Clause (i) at $(0,0)$ with $Y=V_\Lambda(V)$ gives $\mathsf{f}^{\mathrm{pr}}((V,V_\Lambda(V));\varpi)$,
the form of \cite[(1.71)]{B-CMP122b}.

\emph{Clause (X0)(v): joint $G$-invariance.} Let $g$ be a $G$-valued gauge transformation. Write
$\zeta^g$ for the transform of $\zeta$, $\varpi^g$ for the transform of the real variables $\varpi$,
and $\tilde g$ for the induced transformation of fine configurations, all by the transformation law
that accompanies the source tube $\mathfrak{S}_{\mathrm{src}}$.
Consider the function
\begin{equation*}
\zeta\mapsto\mathfrak{q}_C(\zeta^g;\varpi^g)-\mathfrak{q}_C(\zeta;\varpi).
\end{equation*}

First, it is holomorphic on the tube. This follows from (iii), since the tube is invariant under
gauge transformations \cite[\S2]{B-CMP119}.

Second, it vanishes at real $\zeta$. By (iv), at real $\zeta$ both terms are the form of
\cite[(1.71)]{B-CMP122b}, at the data $\mathsf{D}$ and $\mathsf{D}^g$ respectively. That form is
gauge invariant, for the following reasons:
\begin{itemize}
\item $V_\Lambda$ is gauge covariant, by clause (b) of the covariance statement for the map
$V_\Lambda$;
\item \cite[(181)]{B-CMP102b} and \cite[(11)]{B-CMP98} give
$U_{k,C}(\mathsf{D}^{g})=U_{k,C}(\mathsf{D})^{\tilde g}$;
\item $H''$ is covariant by clause (c) of the theorem on the localised small solution;
\item $|\cdot|^2$ is $\operatorname{Ad}$-invariant;
\item $\hat\chi$ does not depend on the gauge.
\end{itemize}

Third, consider a fibre
\begin{equation*}
\bigl\{\zeta=(e^{i\eta A'}U,\mathbf{J}):\ U\ \text{fixed and $G$-valued},\
\ell_k|A'|<\tilde\alpha_{1,k}\bigr\}.
\end{equation*}
It is convex in $A'$. It contains the open real piece $\{A'\ \text{real, small}\}$, every point of
which is a real $\zeta$. By the identity principle on the fibre, the function vanishes on the whole
fibre. Every $\zeta\in\mathfrak{S}_{\mathrm{src}}$ lies on such a fibre. Hence
$\mathfrak{q}_C(\zeta^g;\varpi^g)=\mathfrak{q}_C(\zeta;\varpi)$.
\end{proof}

\begin{remark}
If $\mathsf{T}^{\sharp}$ replaces $\mathsf{T}^{\mathsf{L}}$, clauses (i)--(v) of (X0) in
Lemma~\ref{4.23:prf-trivial-splitting} and their proofs go through unchanged with
two substitutions. Clauses (b) and (d) of the lemma on the operator $H(\mathbf{1})$, which holds at
$\mathfrak{B}^C$ by (H2), replace clauses (c) and (d) of the theorem on the localised small
solution; and the condition \cite[(3.36)]{B-CMP109}, which stands between the conditions
\cite[(3.35), (3.37)]{B-CMP109} defining the class $\mathfrak{K}_{109}(\mathfrak{B}^C;\gamma^\flat)$,
is adjoined to (H8)(c).
Then $\mathfrak{q}_C$ is again well defined, holomorphic and invariant. At real points, however,
it is the form built on the operator identified in \eqref{4.23:eq-linearised-minimiser-b}, not
that of \cite[(1.71)]{B-CMP122b}.
\end{remark}

\begin{remark}[The printed size bound counts old cells]\label{4.23:rem-printed-size-bound}
Throughout this remark, the set $Z^{\sim}$ of the quoted text below is read as the large-field
region $Z^{\sharp}$ that replaces it in the domain class fixed earlier in this chapter. The last
resummation of \cite[(1.68), (1.69)]{B-CMP122b} is described there as follows:
\begin{quotation}
the summation is over domains $Y \in \mathbf{D}_k$ satisfying the property that the
intersection $Y\cap Z^{\sim}$ is a union of components of $Z^{\sim}$, containing at least one
component. The expression corresponding to such a domain $Y$ depends on the new background
field $U_k$ restricted to $Y$, and also on the fields $A$, $B$, $V''$ localized in $Y$. We denote
this expression by $V(Y,U_k)$. It has an analytic extension $V(Y,(\mathbf{U},\mathbf{J}))$
defined on the space $\tilde{U}^{c}_{k}(Y,\alpha_0,\alpha_1)$, and satisfying the bounds:
\begin{equation*}
|V(Y,(\mathbf{U},\mathbf{J}))| \le \alpha \exp(-(1+2\beta)\kappa d_{k,Z}(Y))
\tag*{\cite[(1.68)]{B-CMP122b}}
\end{equation*}
for $Y\setminus Z^{\sim} \neq \emptyset$ \dots{} and
\begin{equation*}
|V(Y,(\mathbf{U},\mathbf{J}))| \le O(1)\sum_j|\Gamma^0_j\cap Y|,
\tag*{\cite[(1.69)]{B-CMP122b}}
\end{equation*}
if $Y$ is a component of $Z^{\sim}$. Here $\{\Gamma^0_j\}$ denotes the old determining set
[\dots].
\end{quotation}
In the quoted bounds $\alpha$ and $\beta$ are the constants of \cite[\S1]{B-CMP122b}, and $U_k$ is
the new background field of that section.
The subtracted copies of these expressions carry the same letter,
\cite[\S1]{B-CMP122b}:
\begin{quotation}
For simplicity we denote them by $V(Y)$ also
\end{quotation}
For the resummed expression, \cite[\S1]{B-CMP122b} states:
\begin{quotation}
The expression $V(X^{\sim})$ is a sum of a big number of terms \dots{} The resummed expression
has the same bound as above, with a different constant $O(1)$
\end{quotation}
This statement is not used as a hypothesis anywhere in this book.

The assertion of this remark is the following. The bound \cite[(1.69)]{B-CMP122b} is a bound on
the size of $V(Y,\zeta)$, $\zeta=(\mathbf{U},\mathbf{J})$, for $Y$ a component of $Z^{\sharp}$, and not
on a difference of two such expressions: its right-hand side counts every old cell in $Y$, and it
therefore depends on the nesting depth $D$ of the old large-field regions in $Y$. The bounds
\cite[(1.68), (1.69)]{B-CMP122b} enter this book only as bounds on size; what the later
estimates use is a difference of such expressions, and it is that difference which the estimates of the
following sections bound.
\end{remark}

\begin{proof}
Let $Y$ be a component of $Z^{\sharp}$, so that the second of the two printed bounds, and not the
first, applies to $V(Y,\zeta)$; the first applies only when $Y\setminus Z^{\sharp}\neq\emptyset$,
and its right-hand side decays exponentially in $d_{k,Z}(Y)$. The right-hand side of the second,
\begin{equation}\label{4.23:eq-printed-size-bound-1}
O(1)\sum_j|\Gamma^0_j\cap Y| ,
\end{equation}
is a sum over the members $\Gamma^0_j$ of the old determining set whose $j$-th term is the number
$|\Gamma^0_j\cap Y|\ge 0$ of cells of $\Gamma^0_j$ lying in $Y$. It carries no decay factor, no
term is discounted, and it bounds $|V(Y,\zeta)|$ at a single point $\zeta$ of
$\tilde{U}^{c}_{k}(Y,\alpha_0,\alpha_1)$, without reference to a second expression or a second
point. It is therefore a bound on the size of $V(Y,\zeta)$ by a count that runs over all old
cells of $Y$, at whatever depth of nesting they were created; this is the sense in which the
bound reads the nesting depth $D$.

The subtracted copies are also denoted $V(Y)$; what later estimates use is a difference of such
expressions, not the expression bounded by \eqref{4.23:eq-printed-size-bound-1}, and that
difference is estimated separately in the following sections.
\end{proof}

\begin{definition}[The resummed density of a treated component]\label{4.23:def-resummed-density}
Let $C$ be a treated component at level $k$, let $C^{\sharp}$ be the box obtained from $C$ by
adjoining its collar layers as in the strip lemma, let $Z$ be the large-field region at level
$k$, and let
$\mathfrak{S}=\tilde{U}^{c}_{k}(C^{\sharp},\tilde\alpha_0,\tilde\alpha_1)$ be the source
tube over $C^{\sharp}$, with points $\zeta=(\mathbf{U},\mathbf{J})$. Let $\Phi_C$ be the
family of real variables of Definition~\ref{4.23:not-oscillation-seminorm}, and let $w$ range
over the complex sources admitted there. Write $\mathbf{D}^{Z}_{k}$ for the class of domains
$Y_4\in\mathbf{D}_k$ such that $Y_4\cap Z$ is a non-empty union of components of $Z$. For
$Y_4\in\mathbf{D}^{Z}_{k}$ put
\begin{equation}\label{4.23:eq-resummed-density-1}
\mathfrak{v}
:=\mathfrak{v}^{\mathrm{br}}+\mathfrak{v}^{\mathrm{val}}+\mathfrak{v}^{\mathrm{rel}}
+\mathfrak{v}^{\mathrm{kept}}+\mathfrak{v}^{\mathbf{C}},
\end{equation}
evaluated at $(Y_4;\zeta;\Phi_C,w)$, the last member $\mathfrak{v}^{\mathbf{C}}(Y_4)$ being
given only for $Y_4\subset C^{\sharp}$; and define the \emph{resummed density} of $C$ by
\begin{equation}\label{4.23:eq-resummed-density-2}
V^{\sharp}(C^{\sharp};\zeta;\Phi_C,w)
:=\sum_{\substack{Y_4\in\mathbf{D}^{Z}_{k}\\ Y_4\subset C^{\sharp}}}
\mathfrak{v}(Y_4;\zeta;\Phi_C,w).
\end{equation}
The localisation identity on which the correlation theorem rests writes the first part of
\cite[(1.49)]{B-CMP122b} less its first two terms, certain terms of
\cite[(1.59), (1.60)]{B-CMP122b}, and the subtracted copies entering that identity, as a sum
of localised members over $Y_4\in\mathbf{D}^{Z}_{k}$.
The five members are the following.
\begin{enumerate}
\item[(a)] $\mathfrak{v}^{\mathrm{br}}$ and $\mathfrak{v}^{\mathrm{val}}$ are the plain bracket
members and the value members of the localisation identity on which the correlation theorem
rests. In that identity the relative members are contained in the bracket member, where they
account for one summand of its bound; here they are written apart as
$\mathfrak{v}^{\mathrm{rel}}$, and the total is unchanged.
\item[(b)] $\mathfrak{v}^{\mathrm{rel}}$ collects the pieces of all wrapped units, in the sense
of the localisation proposition for wrapped units, as given by clause~(d) of that proposition.
\item[(c)] $\mathfrak{v}^{\mathrm{kept}}$ collects the pieces of the members kept by the older
regions $(Y',m')$, $Y'\subset C$, whose data enter $\Phi_C$, as given by the two kept-member
lemmas.
\item[(d)] $\mathfrak{v}^{\mathbf{C}}$ is the family of localised members of $\mathbf{C}_k$
that do not depend on the source coordinates, of Definition~\ref{4.23:def-boundary-terms}
(printed without proof in \cite{B-CMP122b}; cited as printed). The first part of the effective
action \cite[(1.49)]{B-CMP122b} includes its first three terms; on \cite[p.~369]{B-CMP122b}
the right-hand side of that equation begins
\begin{equation*}
-\tfrac12\langle DH''B,\hat\chi DH''B\rangle
-A\bigl(g''^{-2}\hat\chi_1,U''_{h+2}\bigr)
+\mathbf{C}_k+\mathbf{E}+\mathbf{R}+\mathbf{B}+\mathbf{B}''+\cdots,
\end{equation*}
in the notation of that paper. The first two terms, the quadratic form and the Wilson
term, are kept apart in \cite[(1.71)]{B-CMP122b}; the third term is $\mathbf{C}_k$, and its
localised members independent of the source coordinates form $\mathfrak{v}^{\mathbf{C}}$.
\end{enumerate}
At real points of $\mathfrak{S}$ each member is Ba{\l}aban's member of \cite{B-CMP122b}, the
averages entering it being taken as in the conventions of
Definition~\ref{4.23:def-extension-conventions}. Away from the real points each member is
defined through its carrier, the analytic function whose values at real points are
Ba{\l}aban's member, by the clause fixing the members away from the real points and by the
conventions of Definition~\ref{4.23:def-extension-conventions}, applied to the subtracted
copies as well.

The sum~\eqref{4.23:eq-resummed-density-2} has the form of clause~(c) of the strip lemma, which
holds for any family indexed by $\mathbf{D}^{Z}_{k}$. For every real value of $\Phi_C$ and
every admitted $w$, the function $\zeta\mapsto V^{\sharp}(C^{\sharp};\zeta;\Phi_C,w)$ is
holomorphic on $\mathfrak{S}$.
\end{definition}

\begin{proof}
Only the last assertion requires proof. Fix a real value of $\Phi_C$ and an admitted $w$.
Since $C^{\sharp}$ is a box of a finite lattice, only finitely many domains
$Y_4\in\mathbf{D}^{Z}_{k}$ satisfy $Y_4\subset C^{\sharp}$. Hence the sum
\eqref{4.23:eq-resummed-density-2} is finite. Each summand
$\mathfrak{v}(Y_4;\cdot\,;\Phi_C,w)$ is analytic in $\zeta$ on the tube over the sub-domain
$Y_4$ of $C^{\sharp}$: for $\mathfrak{v}^{\mathbf{C}}$ by
Definition~\ref{4.23:def-boundary-terms}, which is printed without proof in
\cite{B-CMP122b} and cited as printed; for the other four members by their definition through
their carriers and the conventions of Definition~\ref{4.23:def-extension-conventions}.

We now apply the argument of the proof of the corollary to the strip lemma on finite sums over
sub-domains, namely the corollary asserting that a finite sum of functions analytic on tubes
over sub-domains of $C^{\sharp}$ is analytic on $\mathfrak{S}$. The argument runs as follows.
For $Y_4\subset C^{\sharp}$, the
map taking a point of $\mathfrak{S}$ to its restriction to $Y_4$ is the coordinate projection
onto the variables attached to $Y_4$. It is therefore holomorphic, and, by the proof of the
corollary on finite sums over sub-domains, it maps $\mathfrak{S}$
into the tube over $Y_4$. Composing a function analytic on the tube over $Y_4$ with this
projection gives a function holomorphic on $\mathfrak{S}$. A finite sum of holomorphic
functions is holomorphic. Consequently $V^{\sharp}(C^{\sharp};\cdot\,;\Phi_C,w)$ is
holomorphic on $\mathfrak{S}$.
\end{proof}

\begin{lemma}[Filing domains: the component and its rim]\label{4.23:lem-filing-domains}
Let the level $k$, the large-field region $Z$, its component $C$, the set $C^{\sharp}$ (the
set $C$ together with $m_{\sharp}(k)=\lfloor\frac12(\check{R}_k-1)\rfloor$ collar layers of
$\pi_k$-cubes), the class $\mathbf{D}_k^Z$ of filing domains, each member of which meets $Z$
in a non-empty union of components of $Z$, the totals $\mathfrak{v}(Y_4)$
and the resummed density $V^{\sharp}(C^{\sharp})$ be as in
Definition~\ref{4.23:def-resummed-density}. All of these are taken at one and the same
argument, which is suppressed. For a unit $v$ of the expansion of the correlation theorem we
write $K_v$ for its base label set, $X_v$ for its localisation domain, $\mathcal{P}_v$ for its
set of nested regions and $\mathcal{P}_v^{\circ}\subset\mathcal{P}_v$ for the set of its live
open nested regions. Here a \emph{nested region} $(Y',m')$ of $v$ is a large-field region $Y'$
created at an earlier step $m'$, carried by $v$ together with its strip $Y'^{\sharp}$, a set of
the corresponding number $m_{\sharp}$ of layers of cubes at its own level that is the same set
at every later level; a
strip, or its nested region, is \emph{live}
after $\mathbf{R}_k$ if it still carries its large-field factor after the operation
$\mathbf{R}_k$. The following hold.
\begin{enumerate}
\item[(i)] Every $Y_4\in\mathbf{D}_k^Z$ with $Y_4\subset C^{\sharp}$ contains $C$ and meets no
other component of $Z$. Either $Y_4=C$, or $Y_4\setminus Z\neq\emptyset$, and in the latter
case $Y_4$ contains a cube of $C^{\sharp}\setminus C$.
\item[(ii)] With the rim class $\mathcal{Y}_{\mathrm{rim}}$ defined by the first line of the
following display, the second line holds:
\begin{equation}\label{4.23:eq-filing-domains-a}
\begin{gathered}
\mathcal{Y}_{\mathrm{rim}}:=\{Y_4\in\mathbf{D}_k^Z:\ C\subsetneq Y_4\subset C^{\sharp}\},\\
V^{\sharp}(C^{\sharp})=\mathfrak{v}(C)+\sum_{Y_4\in\mathcal{Y}_{\mathrm{rim}}}\mathfrak{v}(Y_4).
\end{gathered}
\end{equation}
\item[(iii)] A piece, bracket or value filed at $Y_4=C$ has a parent unit $v$ with
$K_v\subset C$, $X_v\subset C$ and $\mathcal{P}_v^{\circ}=\emptyset$. No unit of the second
part of the expansion of the correlation theorem (every unit of that second part has
$X_v\cap Z=\emptyset$) is filed at $C$. Assume in addition the following \emph{filing
hypothesis}: under the classification of the units of that expansion, every term of
\cite[(1.59), (1.60)]{B-CMP122b} indexed by a domain $Y_0$ (a $Y_0$-term) belongs to a parent
unit $v$ with $K_v\setminus Z\neq\emptyset$ or $X_v\cap Z=\emptyset$. Then no $Y_0$-term is
filed at $C$.
\end{enumerate}
\end{lemma}

\begin{proof}
(i) Let $Y_4\in\mathbf{D}_k^Z$ with $Y_4\subset C^{\sharp}$. By the defining property of the
class $\mathbf{D}_k^Z$ recalled in the statement, the set $Y_4\cap Z$ is a non-empty union of
components of $Z$. By the
separation statement for the collared components, $C^{\sharp}$ contains exactly one component
of $\Lambda_k^c$, the complement of the small-field region $\Lambda_k$ at level $k$, namely $C$.
The set $Z$ is a union of components of $\Lambda_k^c$, so every
component of $Z$ contained in $Y_4\subset C^{\sharp}$ is a component of $\Lambda_k^c$ contained
in $C^{\sharp}$, and is therefore $C$. Hence $Y_4\cap Z=C$. In particular $Y_4\supset C$, and
$Y_4$ meets no other component of $Z$. If $Y_4\neq C$, then $Y_4\setminus Z=Y_4\setminus C$
is non-empty. It is a union of cubes of $C^{\sharp}\setminus C$, and so $Y_4$ contains such a
cube.

(ii) By Definition~\ref{4.23:def-resummed-density}, $V^{\sharp}(C^{\sharp})$ is the sum of
$\mathfrak{v}(Y_4)$ over all $Y_4\in\mathbf{D}_k^Z$ with $Y_4\subset C^{\sharp}$. By (i), each
such $Y_4$ satisfies either $Y_4=C$ or $C\subsetneq Y_4\subset C^{\sharp}$, and these two
cases are exclusive. The sum therefore splits into the summand indexed by $Y_4=C$ and the sum
over $\mathcal{Y}_{\mathrm{rim}}$, which is \eqref{4.23:eq-filing-domains-a}.

(iii) First let a piece or bracket $\mathfrak{l}$ be given, and write
$\mathrm{lab}_0(\mathfrak{l})$ and $\mathrm{lab}_4(\mathfrak{l})$ for its label sets at stages
zero and four of its construction. Its filing domain and its label
sets satisfy
\begin{align*}
Y_4(\mathfrak{l})&:=\mathrm{lab}_4(\mathfrak{l})\cup\{\text{components of $Z$ meeting
$\mathrm{lab}_4(\mathfrak{l})$}\}\\
&\supset\mathrm{lab}_4(\mathfrak{l})\supset\mathrm{lab}_0(\mathfrak{l})=K_v ,
\end{align*}
where $v$ is the parent unit, by clause (c) of the proposition on the labels of the wrapped
units together with the locality clause (a) of the statement on the plain and pointed boundary
terms.
For a value of the unit $v$, the filing domain is the set $Y_1(v)$ of the filing rule for values
of the correlation theorem, and that rule places $K_v$ inside $Y_1(v)$. In each case the filing
domain $Y_4=C$ therefore forces $K_v\subset C$.

Next, the localisation domain of the parent is
\begin{equation*}
X_v=A_v\cup\bigcup_{\mathcal{P}_v}Y'^{\sharp},\qquad A_v\subset K_v ,
\end{equation*}
where the union runs over the strips $Y'^{\sharp}$ of the nested regions in $\mathcal{P}_v$.
The strip of each
nested region in $\mathcal{P}_v$ touches $A_v$. Since $A_v\subset K_v\subset C$ and
$C\subset Z$, each such strip touches $Z$. By clause (e) of the geometric lemma on the
large-field regions and their strips, a strip that is live after $\mathbf{R}_k$ lies at
distance at least $\check{R}_k$ from $Z$. Hence no strip of a nested region in $\mathcal{P}_v$
is live after $\mathbf{R}_k$, so no nested region of $v$ is live, in particular none is a live
open nested region, that is,
$\mathcal{P}_v^{\circ}=\emptyset$, and hence $\mathcal{P}_v=\mathcal{P}_v^{\dagger}$, where
$\mathcal{P}_v^{\dagger}:=\mathcal{P}_v\setminus\mathcal{P}_v^{\circ}$.
By clause (g) of the same lemma, the strips of the nested regions in $\mathcal{P}_v$ lie in the
core $\mathsf{C}=\mathsf{C}_C\subset C$. Together with $A_v\subset C$, this gives $X_v\subset C$.

A unit of the second part has $X_v\cap Z=\emptyset$. The parent of anything filed at $C$ has
$X_v\subset C\subset Z$. Therefore no unit of the second part is filed at $C$.

Finally, assume the filing hypothesis and consider the $Y_0$-terms. By the filing hypothesis,
each $Y_0$-term belongs to a
unit with $K_v\setminus Z\neq\emptyset$ or with $X_v\cap Z=\emptyset$. In the first case
$K_v\not\subset C$, because $C\subset Z$. In the second case $X_v\cap Z=\emptyset$, which is
excluded for the parent of anything filed at $C$, exactly as for the units of the second part.
In either case the term fails the conclusion $K_v\subset C$, $X_v\subset C$ established above,
so it is not filed at $C$.
\end{proof}

\begin{remark}
The filing hypothesis of Lemma~\ref{4.23:lem-filing-domains}\,(iii) reflects the following
feature of these terms in \cite[(1.59), (1.60)]{B-CMP122b}: there the domain $Y_0$ contains the
cube $\square_0$ of its indexing point $z$, and $z$ lies in a region of \cite[\S1]{B-CMP122b}
that is contained in $Z^c$. Thus $z$ and its cube lie off $Z$, and the localisation domain of
the parent contains a point of $Z^c$.

The lemma is used to sort the terms of $V^{\sharp}(C^{\sharp})$ by filing domain. The totals
$\mathfrak{v}(Y_4)$, $Y_4\in\mathcal{Y}_{\mathrm{rim}}$, are bounded as totals, by the bound
that the correlation theorem states for them. Only the single
total $\mathfrak{v}(C)$ is opened, into the constituents that define it in
Definition~\ref{4.23:def-resummed-density}. This opening is a definition and not a bound, so no
estimate of a sub-sum of a total is needed.
\end{remark}

\begin{definition}[Constituents filed at the component]\label{4.23:def-component-constituents}
Let $C$ be a component of $Z$ at level $k$, let $C^{\sharp}\supset C$ be $C$ enlarged by
$m_{\sharp}(k)$ collar layers of $\pi_k$-cubes, and let $\mathsf{C}=\mathsf{C}_C\subset C$ be the
core of $C$.
Let $N=N(g_k)$ and $h=k-N$. Let $V^{\sharp}(C^{\sharp})$ be the resummed
density of Definition~\ref{4.23:def-resummed-density}, and let $\mathfrak{v}(C)$ be its term filed
at $Y_4=C$. By Lemma~\ref{4.23:lem-filing-domains}, clause (ii), this term and the terms filed at
the rim domains $Y_4\in\mathcal{Y}_{\mathrm{rim}}$ make up $V^{\sharp}(C^{\sharp})$; see
\eqref{4.23:eq-filing-domains-a}.

We use the following three classes of units:
\begin{equation}\label{4.23:eq-component-constituents-a}
\begin{aligned}
\mathfrak{L}^{\mathrm{pl}}(C)&=\mathfrak{L}^{B}\sqcup\mathfrak{L}^{\mathrm{pt}},\\
\mathfrak{L}^{\mathrm{wr}}(C)&:=\{v\in\mathfrak{L}^{\mathrm{sh}}\cup\mathfrak{L}^{\mathrm{dp}}:\
K_v\subset C\},\\
\mathfrak{L}^{\mathrm{cp}}(C)&:=\{v:\ v\ \text{a subtracted copy with attached domain}\
Y_1(v)=C\}.
\end{aligned}
\end{equation}
Here $\mathfrak{L}^{B}$ and $\mathfrak{L}^{\mathrm{pt}}$ are described in (m1), the shallow and
deep classes $\mathfrak{L}^{\mathrm{sh}}$ and $\mathfrak{L}^{\mathrm{dp}}$ of wrapped units of the
first part in (m2), and the attached domain $Y_1(v)$ of a copy in (m4).
The decomposition in the first line is the one of the plain-format lemma.
By clause (iii) of Lemma~\ref{4.23:lem-filing-domains}, which restricts the parent units of the
terms filed at $Y_4=C$, and by the table of plain brackets and values on which
Definition~\ref{4.23:def-resummed-density} draws, the constituents of $\mathfrak{v}(C)$ are the
following five items.
\begin{itemize}
\item[(m1)] \emph{Plain constituents.} The pieces
$\mathfrak{e}=(v;\mathsf{y}_1,\dots,\mathsf{y}_4)$ of the units $v\in\mathfrak{L}^{\mathrm{pl}}(C)$
whose label set satisfies $\mathrm{lab}_4(\mathfrak{e})\subset C$; a piece consists of a unit $v$
and the four label data $\mathsf{y}_1,\dots,\mathsf{y}_4$ attached to it. The empty piece
$\vec{\emptyset}$ is included; it is the value of $v$.

Here $\mathfrak{L}^{B}$ is the class of plain stored boundary-type terms. It consists of the terms
$\mathbf{B}^{(j)}$ and $\mathbf{B}'^{(j)}$, the chain terms $\mathbb{C}^{(n)}_k$ of the present
chain, the frozen chain terms $\mathbb{C}^{(i)}_{k_a}$, and the arrest terms in their plain-format
filing. Among the arrest terms
are the families $\mathfrak{Q}(\mathfrak{a})$ of the live arrests $\mathfrak{a}$ nested in $C$;
these families form the third member of the decomposition whose second member is the term
of (m5).

$\mathfrak{L}^{\mathrm{pt}}$ is the class of pointed terms of the plain format. It consists of the
first three entries, the fifth entry and the boundary entry of the table of plain brackets and
values named above; their smooth-core data are those supplied by the pointed-format datum
statement.
\item[(m2)] \emph{Wrapped constituents.} The same pieces, with $\mathrm{lab}_4(\mathfrak{e})\subset C$
and values included, of the units of $\mathfrak{L}^{\mathrm{wr}}(C)$. These are the wrapped units
of the first part with $K_v\subset C$; in the classification of the wrapped-unit format lemmas
they fall into the shallow class $\mathfrak{L}^{\mathrm{sh}}$ and the deep class
$\mathfrak{L}^{\mathrm{dp}}$. The cut between these two classes is the one made in those lemmas.
\item[(m3)] \emph{Inner expiring kept regions.} The same pieces of the two members
$v'\in\mathfrak{K}(Y')$ of every kept region $(Y',m')$ in the sense of
Definition~\ref{4.23:def-resummed-density} (a region $Y'$ together with the level $m'$ at which
it expires) with $Y'\subset C$ that expires at the
present step, that is, with $m'\le h$. For each such member one has, by the statement on the
members of kept regions, $X_{v'}=Y'^{\sharp}\subset\mathsf{C}$, and
$K_{v'}$ is a single $\pi_k$-cube.
\item[(m4)] \emph{Copies.} The values, taken at the point $\zeta$ itself and on the new set of the
present step fixed in the statement on subtracted copies, of the units of
$\mathfrak{L}^{\mathrm{cp}}(C)$. These units are the copies subtracted in Ba{\l}aban's
construction \cite{B-CMP122b} whose attached domain $Y_1(v)$ equals $C$.
\item[(m5)] \emph{Localised $z$-free members.} The member $\mathfrak{v}^{\mathbf{C}}(C)$ of the family
of Definition~\ref{4.23:def-boundary-terms} at $Y_4=C$, whose bound rests on
\cite[(1.11)]{B-CMP122b}, stated there without proof and cited as printed; it enters
$\mathfrak{v}(C)$ as the summand
$\mathfrak{v}^{\mathbf{C}}$ of Definition~\ref{4.23:def-resummed-density}. This member is the
second member of the decomposition named in (m1).
\end{itemize}
\end{definition}

\begin{remark}[Robustness of the list]
The theorem bounding $V^{\sharp}(C^{\sharp})$ is robust with respect to the list of
Definition~\ref{4.23:def-component-constituents}. Suppose a further constituent filed at $C$ is of
one of the following kinds:
\begin{itemize}
\item[($\alpha$)] a piece or value of an interior unit, presented in the form $F^{\mathfrak{b}}$ or
as $f_v$ evaluated along the chain, such that the presentation obeys the two bounds of the
totalisation lemma and carries a pile under the pile bound or under the anchoring bound;
\item[($\beta$)] a constituent independent of $\zeta$;
\item[($\gamma$)] a constituent anchored at cells outside $\mathsf{C}$, whose supremum on each cell is
bounded by a polynomial in $\check{R}_k$;
\item[($\delta$)] a $Y_0$-term, were such a term filed at $C$.
\end{itemize}
A term of kind ($\delta$) enters through the entry of the unit that owns it. Its parent is either a
first-part unit under the interior condition, hence of kind ($\alpha$), or the total filed at
another filing domain in the table of plain brackets and values.

In each case the constituent enters the right member of that theorem's bound through its pile.
The constant $C_{\Sigma}$ of that theorem's bound is then taken again as a maximum, now over a
family that includes the new
constituent. All other constants and bounds of that theorem are unchanged.
\end{remark}

\begin{theorem}[The oscillation theorem for the resummed density]\label{4.23:thm-oscillation-theorem}
Let the parameters of the correlation theorem satisfy the floors of that theorem, together with
the floor on $\kappa_1$, the floor on the ratio $M/M_1$, and the further floor fixed with them
in the order of choice of the parameters. Let Ba{\l}aban's steps $1,\dots,k$ be run in the format
fixed at the beginning of this section, with
the conventions for extensions, for restriction and for real bases of
Definition~\ref{4.23:def-extension-conventions} in force, those for
extensions being read also on the subtracted copies. Let $\mathfrak{q}_C$ be the second kept member,
constructed as in \eqref{4.23:eq-kept-member-construction-a} with the linearised localised
minimiser map $\mathsf{T}^{\mathsf{L}}$ in its final step; it depends on
$\zeta=(\mathbf{U},\mathbf{J})$ only through the configuration $\mathsf{V}=M^{k}(\mathbf{U})$
restricted to $C\setminus\Lambda$, $M^{k}(\cdot)$ denoting the operation of Ba{\l}aban's
construction through which $\mathbf{U}$ is read in \cite[p.~375]{B-CMP122b}, and $\Lambda$
denoting the small-field region of the step.
Assume the following hypotheses, each at its place in the order of choice of the parameters.
\begin{itemize}
\item The hypotheses and ceilings of the bound for the wrapped units whose constant multiplies the
first of the two sums over wrapped units in \eqref{4.23:eq-right-member-a}, so that this bound,
its version on the tube, and the lemma established with it under the same hypotheses hold.
\item The nesting condition.
\item The complex relative statement at step $k$.
\item The nine standing properties of the expansion of the correlation theorem, together with
their three accompanying conditions.
\item The input statements on the plain brackets and on the values, at step $k$.
\item Clauses (d) and (e) of the proposition filing the pieces of the wrapped units (the pieces
summed in the resummed density of Definition~\ref{4.23:def-resummed-density}), together with
the statement on the members kept by older nested regions, under its ceiling.
\item Clauses (a) and (c) of the holomorphy lemma and the accompanying bound, at the nested
regions $(Y',m')$, $Y'\subset C$ (in the notation of
Notation~\ref{4.23:not-oscillation-seminorm}).
\item The hypotheses of the form theorem (Theorem~\ref{4.23:thm-form-theorem}), with its ceiling,
at every step $\le k$, and hence the holomorphy-and-invariance clause and the bound for
$\mathfrak{q}_C$ stated in \eqref{4.23:eq-tube-statement-a}, which enter here as hypotheses;
among them the bound
\begin{equation*}
\sup_{\mathfrak{S}}|\mathfrak{q}_C|\le\mathfrak{q}_{\max}(k),
\qquad
\mathfrak{q}_{\max}(k):=15\,c_{\mathrm{tr}}\,B_\sharp^{2}\,(\mathfrak{b}'_k)^{2}\,M^{d}\,
\mathsf{n}_{\mathrm{sh}}(k),
\end{equation*}
which is the bound of the form theorem for the choice $\mathsf{T}^{\mathsf{L}}$ of the minimiser
map (whose norm bound $K_{\mathsf{T}}$ is the constant written $B_\sharp$ in
\eqref{4.23:eq-form-theorem-a}), restricted to the tube $\mathfrak{S}$; $c_{\mathrm{tr}}$ is the
constant and $\mathsf{n}_{\mathrm{sh}}(k)$ the count entering that bound.
\item Clause $(\beta)$ of the filing statement for the terms of the arrests, with its ceiling.
\item The statement on boundary terms of Definition~\ref{4.23:def-boundary-terms}, printed
without proof in \cite[(1.11), (1.23)--(1.48), p.~369]{B-CMP122b} and cited as printed,
with the constant $\mathsf{c}_{\mathbf{C}}=1+c_IM^{5}$.
\item The ceiling on $\check{R}_k$ and the two ceilings of \eqref{4.23:eq-degree-constant-a}.
\item The following identities and statements, used as stated: the identity by which the
constituents of $V^\sharp(C^\sharp)$ are totalled; the statement used in the third line of that
total; clauses (i) and (ii) of
the statement on the subtracted copies; the closed form of the quantity $\mathsf{b}_k$, whose
constants are fixed in \eqref{4.23:eq-degree-constant-a}, and the cut of the proposition filing
the pieces of the wrapped units; the inequality $\iota(3)\ge1$, where $\iota$ is the function
written $s$ in its defining statement, renamed here to distinguish it from the depth $s$ below the
cutoff and from the interpolation variables $\mathsf{s}^{(t)}(\Delta)$; the table of
filings underlying the list \eqref{4.23:eq-component-constituents-a}; and the third and fifth of
the conditions under which a constituent is filed at $C$.
\end{itemize}
Then there are a number $C_\Sigma$ and the integer
\begin{equation*}
p_\Sigma=2d+6+r_{\mathrm{pr}}+\lceil 2p_1/r_{\mathrm{pr}}\rceil,
\end{equation*}
where $p_1$ is the exponent in $\mathfrak{b}'_k=A_1(\log x_k)^{p_1}$ and $r_{\mathrm{pr}}$ is the
integer fixed, with the table of degrees, in \eqref{4.23:eq-degree-constant-a},
such that for every $K$, every $k\le K$, every label of the expansion of the correlation theorem
and every component $C$ of the large-field
region $Z$ at the operation $\mathbf{R}_k$ the following hold, where $V^\sharp(C^\sharp)$ is the
resummed density of Definition~\ref{4.23:def-resummed-density}, $\Phi_C$, the complex source
$w\in\mathcal{W}_k(C^\sharp)$ and $\operatorname{osc}$ are as in
Notation~\ref{4.23:not-oscillation-seminorm}, and $\Sigma^V(k)$ is the explicit right member
\eqref{4.23:eq-right-member-a}, in which the member coming from the boundary terms is
$2\mathsf{c}_{\mathbf{C}}=2(1+c_IM^{5})$:
\begin{equation}\label{4.23:eq-oscillation-theorem-a}
\begin{gathered}
\operatorname{osc}\bigl(V^\sharp(C^\sharp)\bigr)\le\Sigma^V(k),\\
\Sigma^{\mathrm{tot}}(k):=\Sigma^V(k)+2\,\mathfrak{q}_{\max}(k)\le C_\Sigma\,\check{R}_k^{\,p_\Sigma}.
\end{gathered}
\end{equation}
More precisely:
\begin{itemize}
\item[$(\sharp)$] $V^\sharp(C^\sharp;\cdot\,;\Phi_C,w)$ is holomorphic on
$\mathfrak{S}=\tilde{U}^{c}_{k}(C^\sharp,\tilde\alpha_0,\tilde\alpha_1)$, and it depends on
$\zeta$ only through the restriction of $\zeta$ to $C^\sharp$.
\item[(v)] The first inequality of \eqref{4.23:eq-oscillation-theorem-a} holds.
\item[(t)] The second inequality of \eqref{4.23:eq-oscillation-theorem-a} holds, with
$\mathfrak{q}_{\max}(k)$ the quantity of the bound on $\mathfrak{q}_C$ assumed above.
\item[(b)] For every $\varsigma\ge0$, the constituents (m1), (m2) and (m3) of
Definition~\ref{4.23:def-component-constituents}, taken over the units $v$ with
$d^{\wedge}(\hat{X}_v,\mathsf{C}^{c})\ge\varsigma$, $\hat{X}_v$ denoting the domain attached there
to the unit $v$, contribute to the left-hand side of the first
inequality of \eqref{4.23:eq-oscillation-theorem-a} at most
$e^{-\delta_P\varsigma/8}\,\Sigma^V(k)$.
\item[(d)] No right-hand side above contains the quantity $\sum_j|\Gamma^{0}_j\cap C|$ of the
large-field construction inside $C$, the size $|\mathsf{C}|$ of the core $\mathsf{C}$, the
quantity $D$ or the integer $N_0$ of that construction, the
age of any large-field region (the number of steps since the step that created it), $K$, $k$,
or the number, the sizes or the pasts of the nested regions and of the
arrests inside $C$. The window $N(g_k)$ (the number of levels), $\check{R}_{h+1}$, $\check{R}_k$
and $\mathsf{T}_k$, with bottom level $h=k-N(g_k)$, do appear in $\Sigma^V(k)$.
\item[(g)] $V^\sharp(C^\sharp)$ has the properties that clause (ii) of the statement in which the
kept members are consumed requires of a kept member: for every real
$\Phi_C$ it is holomorphic and bounded on $\mathfrak{S}$, and it is jointly invariant under
gauge transformations with values in the gauge group $SU(N)$ acting on $\zeta$ and on the real
variables together, in the law
under which $\mathfrak{q}_C$ is invariant in \eqref{4.23:eq-tube-statement-a}. The conventions of
Definition~\ref{4.23:def-extension-conventions}, read on the configurations and on the
subtracted copies, do not break this invariance.
\item[(o)] The constant $C_\Sigma$ is fixed, in the order of choice of the parameters of the
correlation theorem, after the last parameter on which it depends, namely the constant
$\mathsf{K}^{0}$ entering the first member of \eqref{4.23:eq-right-member-a}, and
before the constant $\mathfrak{n}^{\mathrm{kept}}$ of the kept members and $\gamma$; it does not
depend on $\mathfrak{n}^{\mathrm{kept}}$.
\item[(n)] Nothing is asserted about $\Sigma^V(k)\to 0$: the quantity $\Sigma^V(k)$ increases as
$g_k$ decreases to $0$.
\end{itemize}
\end{theorem}

The proof of Theorem~\ref{4.23:thm-oscillation-theorem} establishes nothing about
$\mathfrak{q}_C$. Of the hypotheses on $\mathfrak{q}_C$ the oscillation estimate uses only
\begin{equation*}
\operatorname{osc}(\mathfrak{q}_C)\le 2\,\mathfrak{q}_{\max}(k),
\end{equation*}
which follows from the bound on $\sup_{\mathfrak{S}}|\mathfrak{q}_C|$ by the triangle inequality
applied to $\mathfrak{q}_C(\zeta)-\mathfrak{q}_C(\zeta')$ for $\zeta,\zeta'\in\mathfrak{S}$. The
two-point oscillation of $\mathfrak{q}_C$ over the tube $\mathfrak{S}$ is not small, and the
variation and sign statements of the form theorem, which hold for each splitting separately, do
not enter.

The proof of Theorem~\ref{4.23:thm-oscillation-theorem} is given in the statements that follow it
in this section: after the auxiliary lemmas, clauses $(\sharp)$ and (v) are proved together, then
clause (g), then clause (t), and then clauses (b), (d) and (n); clause (o) is proved in the lemma
fixing the place of $C_\Sigma$ in the order of choice of the parameters.

\begin{lemma}[The rim lemma]\label{4.23:lem-rim-lemma}
Work in the setting of Lemma~\ref{4.23:lem-filing-domains}. Thus $C$ is a component of the
large-field region $Z$, and $C^{\sharp}$ is $C$ with
$m_\sharp(k)=\lfloor\tfrac12(\check{R}_k-1)\rfloor$ layers of cubes of the level-$k$ partition
adjoined. The family $\mathcal{Y}_{\mathrm{rim}}$ is that of \eqref{4.23:eq-filing-domains-a},
and $\mathfrak{v}$ is the resummed density of Definition~\ref{4.23:def-resummed-density}, with
localised $z$-free member $\mathfrak{v}^{\mathbf{C}}$. Then
\begin{equation}\label{4.23:eq-rim-lemma-a}
\begin{gathered}
\mathsf{rim}(k):=\sum_{Y_4\in\mathcal{Y}_{\mathrm{rim}}}
\sup\bigl|(\mathfrak{v}-\mathfrak{v}^{\mathbf{C}})(Y_4)\bigr|
\le \mathsf{n}_\sharp(k)\,C_{\mathrm{fib}}\,\mathsf{N},\\
\mathsf{N}:=3+\mathsf{K}^0+N(g_k)\,C_0^{\sharp},\qquad
\mathsf{n}_\sharp(k):=(101\,\check{R}_k)^d .
\end{gathered}
\end{equation}
Each supremum is taken over the points $\zeta\in\mathfrak{S}$ of the source tube, at every real
value of the fields localised in $Y_4$. The number $C_{\mathrm{fib}}$ is the constant of the
fibre count. The number $\mathsf{K}^0$ is the $g$-free number, fixed after $M$, such that the
value part $\mathfrak{v}^{\mathrm{val}}$ of the resummed density is bounded, in the weighted
norm of the proof below, by $\mathsf{K}^0+N(g_k)C_0^{\sharp}$, where $N(g_k)$ is the number of
levels of the window. Finally $C_0^{\sharp}$ is the constant fixed first among the
constants of the correlation theorem.
\end{lemma}

\begin{proof}
Throughout, $a=\tfrac{31}{20}\kappa$. Recall the weighted norm $\|\cdot\|_a$ of the bracket
construction. For a density $\mathfrak{f}$ indexed by $\mathbf{D}_k^Z$ (below, $\mathfrak{f}$ is
one of $\mathfrak{v}^{\mathrm{br}}$, $\mathfrak{v}^{\mathrm{val}}$ and
$\mathfrak{v}^{\mathrm{rel}}+\mathfrak{v}^{\mathrm{kept}}$) it is
\begin{equation}\label{4.23:eq-rim-lemma-1}
\|\mathfrak{f}\|_a:=\sup_{Y_4\in\mathbf{D}_k^Z,\ Y_4\setminus Z\neq\emptyset}
e^{a\,d_{k,Z}(Y_4)}\,|\mathfrak{f}(Y_4)| .
\end{equation}
Here each $|\cdot|$ is a supremum on the tube over $Y_4$, at every real value of the fields
localised in $Y_4$. The supremum in \eqref{4.23:eq-rim-lemma-1} runs over all
$Y_4\in\mathbf{D}_k^Z$ with $Y_4\setminus Z\neq\emptyset$; in particular it covers every
$Y_4\in\mathcal{Y}_{\mathrm{rim}}$.

\emph{The pointwise bound.}
Let $Y_4\in\mathcal{Y}_{\mathrm{rim}}$. Then $C\subsetneq Y_4\subset C^{\sharp}$. By clause~(i)
of Lemma~\ref{4.23:lem-filing-domains}, either $Y_4=C$ or $Y_4\setminus Z\neq\emptyset$. The
first is excluded, so $Y_4\setminus Z\neq\emptyset$.

By Definition~\ref{4.23:def-resummed-density},
\begin{equation*}
\mathfrak{v}-\mathfrak{v}^{\mathbf{C}}
=\mathfrak{v}^{\mathrm{br}}+\mathfrak{v}^{\mathrm{val}}
+\bigl(\mathfrak{v}^{\mathrm{rel}}+\mathfrak{v}^{\mathrm{kept}}\bigr).
\end{equation*}
We bound the three parts as follows.
\begin{itemize}
\item The bracket bound, with the pieces of the wrapped units
written apart, gives
$\|\mathfrak{v}^{\mathrm{br}}\|_a\le\varepsilon_{\mathrm{br}}(k)-\varepsilon^{\mathrm{rel}}(k)$;
this is explained below.
\item The value bound of the bracket construction gives
$\|\mathfrak{v}^{\mathrm{val}}\|_a\le\mathsf{K}_{\mathrm{val}}(k)$,
with $\mathsf{K}_{\mathrm{val}}(k)\le\mathsf{K}^0+N(g_k)C_0^{\sharp}$.
\item By clauses (d) and (e) of the wrapping proposition, together with the two statements on
the pieces of the innermost terminal units of the wrapping that are kept by older regions,
\begin{equation*}
\|\mathfrak{v}^{\mathrm{rel}}+\mathfrak{v}^{\mathrm{kept}}\|_a
\le\|\mathfrak{v}^{\mathrm{rel}}+\mathfrak{v}^{\mathrm{kept}}\|_{a,1}
\le\varepsilon^{\mathrm{rel}}_L(k),
\end{equation*}
where $\|\cdot\|_{a,1}$ is the stronger weighted norm of the wrapping proposition, which
dominates $\|\cdot\|_a$. Under the smallness condition on the relative pieces that accompanies
these two statements, $\varepsilon^{\mathrm{rel}}_L(k)\le1$.
\end{itemize}

It remains to explain and bound the bracket term. The number $\varepsilon_{\mathrm{br}}(k)$ of
the bracket bound is the sum of four parts: the members carrying the factor $\theta'(g_k)$, a
function of the coupling $g_k$ that multiplies the first members of the bracket bound; the
member $\varepsilon^{\mathrm{rel}}(k)$; the interior member $\varepsilon^{\mathrm{int}}(k)$; and
the exterior member $\varepsilon^{\mathrm{ext}}(k)$. By the condition on the bracket constants
in the correlation theorem, the members carrying $\theta'(g_k)$ add up to at most $1$. By the
interior smallness condition of the localisation step,
$\varepsilon^{\mathrm{int}}(k)\le\tfrac12$. By the exterior bound of the bracket construction,
$\varepsilon^{\mathrm{ext}}(k)\le\tfrac12$.

By Definition~\ref{4.23:def-resummed-density}, the bracket part of the correlation theorem holds
the pieces of the wrapped units as the member $\varepsilon^{\mathrm{rel}}(k)$ of the bracket
bound; the definition writes these pieces apart, as $\mathfrak{v}^{\mathrm{rel}}$, and the total
is unchanged. They are therefore counted in $\varepsilon^{\mathrm{rel}}_L(k)$, and the bracket
bound bounds $\mathfrak{v}^{\mathrm{br}}$, as written there, by its remaining members:
\begin{equation*}
\|\mathfrak{v}^{\mathrm{br}}\|_a
\le\varepsilon_{\mathrm{br}}(k)-\varepsilon^{\mathrm{rel}}(k)
\le 1+\tfrac12+\tfrac12=2 .
\end{equation*}

Now take the supremum in \eqref{4.23:eq-rim-lemma-a}. By the restriction convention of
Definition~\ref{4.23:def-extension-conventions}, a point $\zeta\in\mathfrak{S}$ restricts into the
tube over $Y_4$. That supremum is therefore a supremum over a subset of the set on which the
$|\cdot|$ of \eqref{4.23:eq-rim-lemma-1} is taken. The triangle inequality now gives
\begin{equation}\label{4.23:eq-rim-lemma-2}
\begin{aligned}
\sup\bigl|(\mathfrak{v}-\mathfrak{v}^{\mathbf{C}})(Y_4)\bigr|
&\le\bigl[\varepsilon_{\mathrm{br}}(k)-\varepsilon^{\mathrm{rel}}(k)
+\varepsilon^{\mathrm{rel}}_L(k)
+\mathsf{K}_{\mathrm{val}}(k)\bigr]e^{-a\,d_{k,Z}(Y_4)}\\
&\le\bigl[2+1+\mathsf{K}^0+N(g_k)C_0^{\sharp}\bigr]e^{-a\,d_{k,Z}(Y_4)}
=\mathsf{N}\,e^{-a\,d_{k,Z}(Y_4)} .
\end{aligned}
\end{equation}
The constant $3$ in $\mathsf{N}$ is not optimal; it is kept as an upper bound.

\emph{A cube of the rim in each domain.}
By clause~(i) of Lemma~\ref{4.23:lem-filing-domains}, the domain $Y_4$ contains $C$ and meets no
other component of $Z$, so $Y_4\cap Z=C$. Since $Y_4\setminus Z\neq\emptyset$, the same clause
gives a cube $\Delta_*\subset Y_4$ with $\Delta_*\subset C^{\sharp}\setminus C$.

This cube is disjoint from $C$. It lies in $Y_4$, which meets no other component of $Z$. Hence
$\Delta_*\cap Z=\emptyset$ and $\Delta_*\subset Y_4\setminus Z$.

Every $Y_4\in\mathcal{Y}_{\mathrm{rim}}$ lies in $\mathbf{D}_k^Z$ and contains at least one such
cube. The terms are non-negative, so \eqref{4.23:eq-rim-lemma-2} gives
\begin{equation}\label{4.23:eq-rim-lemma-3}
\mathsf{rim}(k)\le\mathsf{N}
\sum_{\substack{\Delta_*\subset C^{\sharp}\setminus C\\ \Delta_*\cap Z=\emptyset}}\
\sum_{\substack{Y_4\in\mathbf{D}_k^Z\\ Y_4\supset\Delta_*}}e^{-a\,d_{k,Z}(Y_4)} .
\end{equation}

\emph{The inner sum.}
We apply the fibre count with the following data.
\begin{itemize}
\item The islands are the components of $Z$.
\item The set of islands has union $W:=Z$.
\item The exponent is $\lambda_0:=a$.
\end{itemize}
Write $d_W$ for the distance function of the fibre count relative to the island set $W$; for
$W=Z$ it is $d_{k,Z}$, by the identification of distances.

The components of $Z$ are pairwise at distance at least $\check{R}_k$, by the separation
property of the components of the large-field region. In particular they are pairwise
non-touching.

The exponent is admissible. The auxiliary parameter $\kappa$ obeys the floor
$\kappa/20\ge2(\mathsf{k}_d+1)$, where $\mathsf{k}_d$ is the dimensional constant of the fibre
count: for each $n$, the number of domains $Y$ counted there with $d_W(Y)<n$ is at most
$e^{\mathsf{k}_d(n+1)}$. Therefore
\begin{equation*}
a=\tfrac{31}{20}\kappa\ge\tfrac{1}{20}\kappa\ge2(\mathsf{k}_d+1)\ge\mathsf{k}_d+1 .
\end{equation*}

Each $Y_4\in\mathbf{D}_k^Z$ lies in $\mathbf{D}_k$. By the definition of $\mathbf{D}_k^Z$,
$Y_4\cap W$ is a non-empty union of components of $Z$, that is, a union of islands. For a cube
$\Delta_*$ with $\Delta_*\cap Z=\emptyset$ and $\Delta_*\subset Y_4$, we have
$\Delta_*\subset Y_4\setminus W$.

By the identification of distances, $d_{k,Z}(Y_4)=d_W(Y_4)$ with $W=Z$.
Hence the inner sum in \eqref{4.23:eq-rim-lemma-3} is a sub-sum, with non-negative terms, of
the sum of $e^{-\lambda_0 d_W(Y)}$ over all domains $Y$ of the kind counted by the fibre count.
That sum is at most $C_{\mathrm{fib}}$, hence so is the inner sum.

\emph{The number of cubes.}
By the collar-box lemma, $N_k(C^{\sharp}\setminus C)\le(101\check{R}_k-1)^d$.
The outer sum in \eqref{4.23:eq-rim-lemma-3} has at most that many terms, and
$(101\check{R}_k-1)^d\le\mathsf{n}_\sharp(k)$. Therefore
$\mathsf{rim}(k)\le\mathsf{n}_\sharp(k)\,C_{\mathrm{fib}}\,\mathsf{N}$, which is
\eqref{4.23:eq-rim-lemma-a}.
\end{proof}

\begin{remark}
The sums over $\mathcal{Y}_{\mathrm{rim}}$ contain the exterior pieces of the buried units and of
the innermost terminal units of the wrapping whose pieces make up $\mathfrak{v}^{\mathrm{rel}}$
and $\mathfrak{v}^{\mathrm{kept}}$ in Definition~\ref{4.23:def-resummed-density}. These pieces
are estimated within the bounds of the
units to which they belong. What enters \eqref{4.23:eq-rim-lemma-2} from each of the bounds used
is its right-hand side, which is a number or a polynomial in $\log x_k$.

A bound of the same form as \eqref{4.23:eq-rim-lemma-a}, with another constant, accompanies the
bound on the brackets. It is not used here.
\end{remark}

The terms of the pointed families need the curvature and regularity inputs at pairs of points of the source tube that differ in their first component.

The three-factor bound in the tube direction is available in three clauses. Its two deep clauses carry the depth factor $e^{-\delta_P D_X/2}$, with $\delta_P$ as below and $D_X$ the own-scale contour distance from the true domain $X$ of the term to the complement of the core. Its interior clause carries the factor $2$. For the terms of the first three pointed families, however, the format norm in that bound is the supremum of the modulus of the term over its domain of definition, without any power of the interpolation variables. A cell priced through it carries a constant times the number of levels of the window, as the norm count for the pointed families records, so the estimate grows with the number of steps. The pointed three-factor bound avoids this: it gives the product of the powers of the interpolation variables with the depth factor. It is stated only for pairs $m,m'\in\mathfrak{P}$ with the same $(\zeta,\varpi)$, that is, in the interpolation directions $\vec{\mathsf{s}}$ alone. The oscillation seminorm of Definition~\ref{4.23:not-oscillation-seminorm} compares values at two points $\zeta,\zeta'$ of the source tube $\mathfrak{S}_{\mathrm{src}}$. It therefore needs the curvature and regularity inputs at pairs with $\zeta'\neq\zeta$, and the following lemma supplies them.

Throughout, points of the presented polydisc are written $m=(\zeta,\varpi,\vec{\mathsf{s}})$, with $\zeta=(\mathbf{U},\mathbf{J})\in\mathfrak{S}_{\mathrm{src}}$, $\varpi\in\operatorname{supp}\Pi$ and $\vec{\mathsf{s}}$ the interpolation variables. Let $\mathsf{C}$ be the core of the component under consideration. For a pair $m,m'\in\mathfrak{P}$ we use the notation of the curvature bound along the interpolation and of the regularity statement. Specifically:
\begin{itemize}
\item $P_\lambda$, for complex $\lambda$, is the family formed there from the presented data at $m'$ and at $m$;
\item $\mathsf{r}_S$ is the admissible radius of that family on a set $S$, and $\mathsf{r}^{0}$ is its base radius; they are related by $\mathsf{r}_S=\mathsf{r}^{0}e^{\delta_P d^{\wedge}(S,\mathsf{C}^{c})/2}$, where $d^{\wedge}$ is the own-scale contour distance and $\delta_P$ the constant of the depth factor;
\item $U_2(\varpi)$ and $J_2$ are the real base and the real current of the second link of the chain;
\item $\mathcal{Q}_a$, for $a\in\{1,3,6,7\}$, are the quantities bounded there against the thresholds $\theta_a$;
\item $\mathcal{K}_{\mathrm{pt}}$ and $\mathfrak{j}$ are the kernel and the current that the point lemma bounds at a point $y$, $N_y(\mathcal{K}_{\mathrm{pt}})$ is the norm at $y$ in which it bounds the kernel, and $\mathfrak{k}(y)$ is the common bound at $y$;
\item $\varphi_\flat$, $\mathbf{r}_a$, the projection $\pi$, the length unit $\ell$, the spacing $\eta$, the cube $\Delta_y$, the constant $c_{\mathrm{reg}}$, the cube weight $\mathsf{w}^{\vee}(\cdot)$ and the radius $\alpha_{*,n}$ at the level $n$ are as in those two statements.
\end{itemize}

\begin{lemma}[Curvature and regularity at paired tube points]\label{4.23:lem-paired-tube-points}
Assume the following statements.
\begin{enumerate}
\item[(0)] The curvature bound along the interpolation and the regularity statement, with the derivations given for them, together with the gauge-cube condition that accompanies the pointwise curvature bound.
\item[(1)] The pointwise curvature bound on the presented polydisc, which holds at every $m\in\mathfrak{P}$.
\item[(2)] The point lemma, whose hypotheses are that $m,m'\in\mathfrak{P}$ have the same $\varpi$. Together with the depth-weight bound, it gives
\begin{equation*}
|\lambda|\,N_y(\mathcal{K}_{\mathrm{pt}})\le\mathfrak{k}(y),\qquad
|\lambda|\,\ell^3|\mathfrak{j}|\le\mathfrak{k}(y).
\end{equation*}
We also assume the first and second comparator claims on which its proof rests, and the Baker--Campbell--Hausdorff bound and the covariant-to-flat comparison, which its proof also uses.
\item[(3)] The identity of the real base, with its companion for the current, and the chain identity, each at every $m\in\mathfrak{P}$.
\item[(4)] The elementary half of the linearisation statement, the age bound at the value $b=0$ of its parameter, the amplitude bounds, and the choice of the base radius $\mathsf{r}^{0}$.
\item[(5)] The independence statement established in the proof of the three-factor bound in the tube direction: each link of the proof of the point lemma depends on the pair $m,m'$ only through their common $\varpi$ and through each point's own data separately, and at no link uses the condition $\zeta=\zeta'$.
\end{enumerate}
Let $m=(\zeta,\varpi,\vec{\mathsf{s}})$ and $m'=(\zeta',\varpi,\vec{\mathsf{s}}')$ be points of $\mathfrak{P}$ with the same $\varpi$, whether or not $\zeta'=\zeta$. Then the curvature bound along the interpolation and the second clause of the regularity statement hold for the pair $(m,m')$. To state them, consider the two bounds
\begin{equation}\label{4.23:eq-paired-tube-points-a}
\begin{gathered}
\mathcal{Q}_a(P_\lambda)(y)\le\tfrac{1}{16}\,\theta_a(y)
+\tfrac12\,\varphi_\flat\,\mathbf{r}_a(\pi(y))<\theta_a(y),\\
\ell|\mathfrak{A}^{0}|,\ \ell^2|\nabla\mathfrak{A}^{0}|,\ \ell^3|\mathbf{J}^{0}|
\le\mathfrak{a}':=3c_{\mathrm{reg}}\,\mathsf{w}^{\vee}(\Delta_y)\,\alpha_{*,n},
\end{gathered}
\end{equation}
the second of which concerns the data $\mathfrak{A}^{0},\mathbf{J}^{0}$ of the third item below. With these, for every $S\subset\mathsf{C}$ and every $\lambda\in\mathbb{C}$ with $|\lambda|\le\mathsf{r}_S$, the following hold.
\begin{itemize}
\item The $G$-part of $P_\lambda\restriction S$ is $U_2(\varpi)$.
\item The first line of \eqref{4.23:eq-paired-tube-points-a} holds for $a\in\{1,3,6,7\}$ and $y\in S$.
\item If $S$ is one of the boxes of the gauge-cube condition that accompanies the pointwise curvature bound, then $P_\lambda\restriction S$ is represented, in a single gauge of the complexified group $G^{\mathbb{C}}$, as $(e^{i\eta\mathfrak{A}^{0}},\mathbf{J}^{0})$, where $\mathfrak{A}^{0}$ and $\mathbf{J}^{0}$ satisfy the second line of \eqref{4.23:eq-paired-tube-points-a}.
\end{itemize}
\end{lemma}

\begin{proof}
The derivation of the curvature bound along the interpolation uses the pair $(m,m')$ only through four inputs. For each input we determine whether the two first components $\zeta,\zeta'$ are required to coincide.

\emph{First input.} The base of the derivation is the pointwise curvature bound at $m'$. By hypothesis~(1) this bound holds at every point of $\mathfrak{P}$. It therefore holds at $m'$ whatever $\zeta'$ is.

\emph{Second input.} The derivation uses the point lemma together with the depth-weight bound. These give, at each point $y$,
\begin{equation*}
|\lambda|\,N_y(\mathcal{K}_{\mathrm{pt}})\le\mathfrak{k}(y),\qquad
|\lambda|\,\ell^3|\mathfrak{j}|\le\mathfrak{k}(y).
\end{equation*}
The hypotheses of the point lemma are that $m,m'\in\mathfrak{P}$ have the same $\varpi$; they impose no condition on $\zeta,\zeta'$. Its proof passes through the following links.
\begin{itemize}
\item The first part of the first comparator claim: the comparators depend on $m$ through $\varpi$ only. This holds because the identity of the real base and its companion for the current make $U_2$ and $J_2$ real and free of $\zeta$.
\item The second part of the first comparator claim: it is applied once per point, at that point's own kernel and data.
\item The third part of the first comparator claim: it requires only that the two points have the same $\varpi$.
\item The second comparator claim: it is applied once per point.
\end{itemize}
No link of this chain requires $\zeta=\zeta'$, for the following reasons.
\begin{itemize}
\item The comparator of the first claim depends only on the restriction $(U_2,J_2)\restriction\mathsf{F}$ to the set $\mathsf{F}$ on which it is formed.
\item The proof of the point lemma uses the chain identity at $m$ and, separately, at $m'$. It also uses the Baker--Campbell--Hausdorff bound and the covariant-to-flat comparison.
\item Finally, it uses that the group-valued datum $\mathbf{K}$ entering the proof of the point lemma equals $U_2$. By the identity of the real base, $\mathbf{K}$ is the same at $m$ and at $m'$.
\end{itemize}
That none of these links requires $\zeta=\zeta'$ is exactly the independence statement of hypothesis~(5). Hence the conclusion of the point lemma holds, and with it the two displayed bounds, for every pair with the same $\varpi$, in particular when $\zeta'\neq\zeta$.

\emph{Third input.} The derivation uses that the $G$-part of $P_\lambda\restriction S$ is the same $U_2$ throughout. By hypothesis~(3), the identity of the real base and the chain identity hold at any $m\in\mathfrak{P}$. The first makes $U_2=U_2(\varpi)$ real and free of $\zeta$. The second holds at $m$ and at $m'$ separately. Therefore the $G$-part is $U_2(\varpi)$ at both points, for any $\zeta,\zeta'$.

\emph{Fourth input.} The derivation also uses the four statements of hypothesis~(4): the elementary half of the linearisation statement, the age bound at $b=0$, the amplitude bounds, and the choice of $\mathsf{r}^{0}$. None of them depends on a point of $\mathfrak{P}$, so they apply unchanged to the pair $(m,m')$.

Since the pair enters the derivation only through these four inputs, and none of them requires $\zeta'=\zeta$, the derivation applies verbatim to every pair with the same $\varpi$. This gives the $G$-part assertion and the first line of \eqref{4.23:eq-paired-tube-points-a} for all $S\subset\mathsf{C}$, $|\lambda|\le\mathsf{r}_S$, $a\in\{1,3,6,7\}$ and $y\in S$.

The second clause of the regularity statement uses two facts: the curvature bound along the interpolation, just established for the pair $(m,m')$, and the independence of its gauge transformation $u$ from $\lambda$. The transformation $u$ is the one that gauges $U_2(\varpi)$. Since $U_2(\varpi)$ depends on $\varpi$ alone, $u$ does not depend on $\lambda$, and it is the same for every pair with this $\varpi$, whether or not $\zeta'=\zeta$. Hence, for $S$ a box of the gauge-cube condition, $P_\lambda\restriction S$ is represented in the single $G^{\mathbb{C}}$-gauge as $(e^{i\eta\mathfrak{A}^{0}},\mathbf{J}^{0})$ with the bounds of the second line of \eqref{4.23:eq-paired-tube-points-a}.
\end{proof}

\begin{proposition}[Depth decay of deep pointed units]\label{4.23:prop-pointed-depth-decay}
Fix the component $C$, with core $\mathsf{C}=\mathsf{C}_C$ and presented polydisc
$\mathfrak{P}$, whose points are written $m=(\zeta,\varpi,\vec{\mathsf{s}})$ with
$\zeta=(\mathbf{U},\mathbf{J})$ a point of the source tube, $\varpi$ the real variables and
$\vec{\mathsf{s}}$ the interpolation variables. Let $\mathfrak{L}^{\mathrm{pt}}$ be the set of
pointed units of $C$. For $v\in\mathfrak{L}^{\mathrm{pt}}$ let $f_v$ be the functional of $v$ on
$\mathfrak{P}$, $\hat{X}_v$ its box, $\hat{D}_v\ge 0$ its depth, and $\mathfrak{N}^{\flat}(v)$ its
datum as fixed in the table of format data of pointed units. Call $v$ \emph{deep} if its support
$X_v$ lies in the core $\mathsf{C}$; by convention $\hat{D}_v=0$ whenever
$\hat{X}_v\not\subset\mathsf{C}$. Let $\mathsf{r}^{0}>0$ be the constant in the
analyticity radius $\mathsf{r}_{X}$ attached to boxes $X$, let $C_{\mathrm{osc}}$ be the constant of
the oscillation bound for long single terms and tails, and put
\begin{equation}\label{4.23:eq-pointed-depth-decay-1}
C^{\mathrm{pt}}=\max\{C_{\mathrm{osc}},\,2/\mathsf{r}^{0},\,2\}.
\end{equation}
Then for every deep $v\in\mathfrak{L}^{\mathrm{pt}}$ and all
\begin{equation*}
m=(\zeta,\varpi,\vec{\mathsf{s}}),\qquad m'=(\zeta',\varpi,\vec{\mathsf{s}}{}')\in\mathfrak{P}
\end{equation*}
with the same $\varpi$,
\begin{equation}\label{4.23:eq-pointed-depth-decay-2}
|f_v(m)-f_v(m')|\le C^{\mathrm{pt}}\,e^{-\frac12\delta_P\hat{D}_v}\,\mathfrak{N}^{\flat}(v).
\end{equation}
Neither the constant nor the right-hand side involves the age of $v$, the window numbers $N$ and
$N_0$, the window radius $\check{R}$, or the slot of $v$.
\end{proposition}

\begin{proof}
Put
\begin{equation}\label{4.23:eq-pointed-depth-decay-3}
\mathsf{r}_{\hat{X}_v}=\mathsf{r}^{0}e^{\frac12\delta_P\hat{D}_v}.
\end{equation}
Throughout, $\mathfrak{N}^{\flat}$ stands for $\mathfrak{N}^{\flat}(v)$.

\emph{The case $\hat{D}_v=0$ or $\mathsf{r}_{\hat{X}_v}\le 1$.} By the holomorphy and format
statement for pointed
units, for every $v\in\mathfrak{L}^{\mathrm{pt}}$ the map $m\mapsto f_v(m)$ is holomorphic on
$\mathfrak{P}$ and
\begin{equation*}
|f_v(m)-f_v(1)|\le\mathfrak{N}^{\flat}(v)\qquad(m\in\mathfrak{P}),
\end{equation*}
where $1\in\mathfrak{P}$ denotes the base point of the presented polydisc at which the format
bound is centred.
Since $\zeta$ is a coordinate of $m$, this applies at $m$ and at $m'$ whatever $\zeta,\zeta'$ are,
and the triangle inequality gives $|f_v(m)-f_v(m')|\le 2\mathfrak{N}^{\flat}$. If $\hat{D}_v=0$,
then $2\mathfrak{N}^{\flat}\le C^{\mathrm{pt}}e^{0}\mathfrak{N}^{\flat}$ because
$C^{\mathrm{pt}}\ge 2$ by \eqref{4.23:eq-pointed-depth-decay-1}. If $\mathsf{r}_{\hat{X}_v}\le 1$,
then, by \eqref{4.23:eq-pointed-depth-decay-3} and $2/\mathsf{r}^{0}\le C^{\mathrm{pt}}$,
\begin{equation*}
2\mathfrak{N}^{\flat}\le\frac{2\mathfrak{N}^{\flat}}{\mathsf{r}_{\hat{X}_v}}
=\frac{2}{\mathsf{r}^{0}}\,e^{-\frac12\delta_P\hat{D}_v}\,\mathfrak{N}^{\flat}
\le C^{\mathrm{pt}}e^{-\frac12\delta_P\hat{D}_v}\,\mathfrak{N}^{\flat}.
\end{equation*}

\emph{The remaining case.} Otherwise $\hat{D}_v\neq 0$ and $\mathsf{r}_{\hat{X}_v}>1$; since the
depth is set to zero
when $\hat{X}_v\not\subset\mathsf{C}$, we have $\hat{X}_v\subset\mathsf{C}$. We treat the kinds of
pointed unit in turn. For $|\lambda|<\mathsf{r}_{\hat{X}_v}$ let $P_\lambda$ be the first component
of the presented
triple interpolating between $m'$ and $m$, so that, by the presentation lemma for three-field terms,
$P_0$ and $P_1$ are the first components of the presented triples at $m'$ and $m$ respectively.

\emph{Long single terms and tails.} Here \eqref{4.23:eq-pointed-depth-decay-2} follows from
clause~(a) of the oscillation bound for long single terms and tails, applied as stated with its
datum taken to be $\tfrac12\mathfrak{N}^{\flat}(v)$ and its depth $D_X$ taken to be $\hat{D}_v$,
since $C_{\mathrm{osc}}\le C^{\mathrm{pt}}$ by \eqref{4.23:eq-pointed-depth-decay-1}.

\emph{Short single terms.} Let $X=X_v$ be the support of the term, with functional $F(X,\cdot)$,
and set
\begin{equation*}
G(\lambda):=F(X,P_\lambda)-F(X,1),
\end{equation*}
where $F(X,1)$ is the value of the functional at the base point $1$ of the presented polydisc.
By the analyticity statement for three-field terms, in its form for pairs of points with the same
$\varpi$, $G$ is holomorphic on $|\lambda|<\mathsf{r}_{\hat{X}_v}$, and
$\mathsf{r}_{\hat{X}_v}\le\mathsf{r}_X$. On this disc
$|G|\le\mathfrak{N}^{\flat}(v)$: the regularity bounds in \eqref{4.23:eq-paired-tube-points-a}
(Lemma~\ref{4.23:lem-paired-tube-points}) hold for $P_\lambda$ restricted to the boxes
$S\subset\mathsf{C}$ for which they are stated in that lemma, whether or not $\zeta'=\zeta$,
and the single-term bound on regular pairs, a statement about any pair $(\mathbf{U},\mathbf{J})$
restricted to $S$ that satisfies the three regularity bounds, then gives $|G|\le\mathfrak{N}^{\flat}$.
By the presentation lemma, $G(0)$ and $G(1)$ are the values at $m'$ and $m$, so that
\begin{equation*}
G(1)-G(0)=f_v(m)-f_v(m').
\end{equation*}
The function $h(\lambda):=G(\lambda)-G(0)$ is holomorphic on $|\lambda|<\mathsf{r}_{\hat{X}_v}$,
vanishes at
$\lambda=0$, and satisfies $|h(\lambda)|\le|G(\lambda)|+|G(0)|\le 2\mathfrak{N}^{\flat}$ there.
Schwarz's lemma
gives $|h(\lambda)|\le 2\mathfrak{N}^{\flat}|\lambda|/\mathsf{r}_{\hat{X}_v}$ for
$|\lambda|<\mathsf{r}_{\hat{X}_v}$; since $\mathsf{r}_{\hat{X}_v}>1$ we
may take $\lambda=1$, and by \eqref{4.23:eq-pointed-depth-decay-3}
\begin{equation*}
|G(1)-G(0)|\le\frac{2\mathfrak{N}^{\flat}}{\mathsf{r}_{\hat{X}_v}}
=\frac{2}{\mathsf{r}^{0}}\,e^{-\frac12\delta_P\hat{D}_v}\,\mathfrak{N}^{\flat}
\le C^{\mathrm{pt}}e^{-\frac12\delta_P\hat{D}_v}\,\mathfrak{N}^{\flat}.
\end{equation*}

\emph{The Wilson piece.} The same argument applies, with the first entry of the curvature bounds in
\eqref{4.23:eq-paired-tube-points-a} (Lemma~\ref{4.23:lem-paired-tube-points}) as its input.

\emph{Smooth cores.} Let $j$ be the level at which the core is born, $p=p_v$ the level at which it is
evaluated, $Q_v$ its box at level $p$, and $\mathfrak{i}^{(j)}_y$ its functional at the point $y$.
Write $\mathcal{U}(P_\lambda)$ for the group-valued component of $P_\lambda$,
$\mathcal{U}(P_\lambda)|_{Q_v}$ for its restriction to $Q_v$, and $\operatorname{ext}_p$ for the
extension at level $p$ of Definition~\ref{4.23:def-extension-conventions}, and set
\begin{equation*}
G(\lambda):=\mathfrak{i}^{(j)}_y\bigl(\operatorname{ext}_p(\mathcal{U}(P_\lambda)|_{Q_v})\bigr).
\end{equation*}
By the regularity bounds in \eqref{4.23:eq-paired-tube-points-a}, the companion bound of the
regularity statement and the smooth-core lemma, one has $|G|\le\mathfrak{N}^{\flat}(v)$ on the whole
disc $|\lambda|<\mathsf{r}_{\hat{X}_v}$. The smooth-core lemma is a statement about any
configuration on a box $Q$ that
reads $e^{i\eta\mathcal{A}}$, where $\eta$ is the spacing of the lattice carrying the configuration
and $\mathcal{A}$ is a potential with
\begin{equation*}
\max\{\ell_p|\mathcal{A}|,\,\ell_p^{2}|\nabla\mathcal{A}|\}\le\epsilon\,\mathfrak{r}_p ,
\end{equation*}
$\ell_p$ being the unit of length at level $p$ and $\epsilon\,\mathfrak{r}_p$ the radius fixed in
that lemma at level $p$. The smooth-core lemma, which rests on the admissibility lemma at scale $p$,
delivers the datum $\mathfrak{N}^{\flat}$ of the core, with its one factor $L^{-\hat{\alpha}(p-j)}$,
$\hat{\alpha}$ being the decay exponent of the core, between the birth level $j$ and the evaluation
level $p$. This factor enters once, at level $p$, through the regularity ball of level $p$ given by
the regularity-ball lemma; it is used here by statement, and is not derived, iterated or improved
here.

The ends of $G$ are $G(0)=f_v(m')$ and $G(1)=f_v(m)$. This follows from the convention of
Definition~\ref{4.23:def-extension-conventions} for the evaluation of cores, because by the
presentation lemma the group-valued components of $P_0$ and $P_1$ on $Q_v\subset\mathsf{C}$ are
$\mathcal{U}(m')$ and $\mathcal{U}(m)$, the group-valued configurations presented at the points
$m'$ and $m$; no condition at the real point is used. Schwarz's lemma,
applied to $G-G(0)$ exactly as for short single terms, gives
\eqref{4.23:eq-pointed-depth-decay-2}.

\emph{Rough cores.} A rough core is made of smooth constituents carried by boxes
$\hat{X}\subset\hat{X}_v$. The bound is obtained termwise: the smooth-core argument is applied to
each smooth constituent on its box $\hat{X}\subset\hat{X}_v$.

Thus a deep pointed unit depends on the point $\zeta=(\mathbf{U},\mathbf{J})$ of the tube no more
strongly than on the interpolation variables $\vec{\mathsf{s}}$, and its datum remains
$\mathfrak{N}^{\flat}(v)$.
\end{proof}

\begin{lemma}[Total oscillation of one unit. Stated; the proof below is incomplete at the step marked $(\ast)$]\label{4.23:lem-unit-oscillation}
Let $v$ be a unit of one of the classes \textup{(m1)}--\textup{(m3)} of the list
\eqref{4.23:eq-component-constituents-a} of constituents filed at the component $C$
(Definition~\ref{4.23:def-component-constituents}). Its index sets, its tuples
$\vec{\mathsf{y}}=(\mathsf{y}_1,\dots,\mathsf{y}_4)$, its labels
$\mathrm{lab}_0=K_v,\mathrm{lab}_1,\dots,\mathrm{lab}_4$ and its pieces
$\mathfrak{p}(v;\vec{\mathsf{y}})$ are those of that definition. Put
\begin{equation*}
\operatorname{tot}(v):=\sum_{\vec{\mathsf{y}}:\ \mathrm{lab}_4\subset C^{\sharp}}
\operatorname{osc}\bigl(\mathfrak{p}(v;\vec{\mathsf{y}})\bigr).
\end{equation*}
The sum runs over all tuples of the index sets of $v$ with $\mathrm{lab}_4\subset C^{\sharp}$,
and the empty tuple $\vec{\emptyset}$ is included. Here $\operatorname{osc}$ is the oscillation
seminorm over the tube $\mathfrak{S}$ of Definition~\ref{4.23:not-oscillation-seminorm}. For the
stored units that carry live slots in the sense of that definition, the supremum over the couplings
$g'\in\mathfrak{G}_{\varepsilon'}(X_v)$ at which these units are read is included in
$\operatorname{osc}$, and the sum defining $\mathfrak{N}^{A,\mathbb{C}}(v)$ runs over the strong
sets $\mathcal{S}$ of the unit.

To each $v$ the table in \eqref{4.23:eq-unit-oscillation-a} assigns a region $\hat X_v$. For $v$
in the pointed class $\mathfrak{L}^{\mathrm{pt}}$, $\hat X_v$ is the region of
Proposition~\ref{4.23:prop-pointed-depth-decay}. For a Gaussian family $v_e$, $\hat X_v=S^{+}(e)$
is the region on which $v_e$ reads the chain of localised terms. The last column of the table,
headed ``if $\hat X_v\subset\mathsf{C}$'', is
used for $v$ when $\hat X_v\subset\mathsf{C}$, and in that case we put
\begin{equation*}
\hat D_v:=d^{\wedge}(\hat X_v,\mathsf{C}^{c}).
\end{equation*}
Define $\mathfrak{n}(v)$ as follows. If $\hat X_v\not\subset\mathsf{C}$, $\mathfrak{n}(v)$ is the
entry in the column headed ``always''. If $\hat X_v\subset\mathsf{C}$, $\mathfrak{n}(v)$ is the
minimum of the entries of its line in the columns headed ``always'' and
``if $\hat X_v\subset\mathsf{C}$''. For the units $v'\in\mathfrak{K}(Y')$ one has
$\hat X_{v'}=Y'^{\sharp}\subset\mathsf{C}$, so the last column always applies to them, and
$\mathfrak{n}(v')$ is the entry of that column, with $\mathfrak{O}(v')$ the oscillation datum of $v'$.
The number $\mathfrak{O}(v')$ is defined further on in this chapter; here it enters only as the
datum of the bound in the last line of the table, and no property of it other than that bound is used.
Then, with $\mathfrak{n}(v)$ read from the table, the
inequality in the last line of the following display holds:
\begin{equation}\label{4.23:eq-unit-oscillation-a}
\begin{gathered}
\begin{array}{l|c|c|c}
\text{class of }v & \hat X_v & \text{always}
 & \text{if }\hat X_v\subset\mathsf{C}\\ \hline
\mathfrak{L}^{B}\setminus\bigcup\mathfrak{Q}(\mathfrak{a}),\ \ \mathfrak{L}^{\mathrm{wr}}
 & X_v & 2\,\mathfrak{N}^{A,\mathbb{C}}(v)
 & C_{3F}\,e^{-\frac12\delta_P\hat D_v}\,\mathfrak{N}^{A,\mathbb{C}}(v)\\
\mathfrak{L}^{\mathrm{pt}} & \text{as in Proposition~\ref{4.23:prop-pointed-depth-decay}}
 & 2\,\mathfrak{N}^{\flat}(v)
 & C^{\mathrm{pt}}e^{-\frac12\delta_P\hat D_v}\,\mathfrak{N}^{\flat}(v)\\
v=v_e\in\mathfrak{Q}(\mathfrak{a}) & S^{+}(e) & 2\cdot2\,\nu^{\natural}(e)
 & C_{3F}\,e^{-\frac12\delta_P\hat D_v}\,\mathfrak{N}^{G,L}(e)\\
v=v'\in\mathfrak{K}(Y') & Y'^{\sharp} & \text{---}
 & \tfrac12\,C_{3F}\,e^{-\frac12\delta_P\hat D_v}\,\mathfrak{O}(v')
\end{array}\\[2ex]
\operatorname{tot}(v)\le\mathfrak{n}(v)\,e^{4(d_k(K_v)+1)}.
\end{gathered}
\end{equation}
\end{lemma}

\begin{proof}
\emph{The function $G$.} Fix real variables $\Phi_C$, a source parameter in
$\mathcal{W}_k(C^{\sharp})$ and the weak coordinates, as in
Definition~\ref{4.23:not-oscillation-seminorm}. For units with live slots, fix also
$g'\in\mathfrak{G}_{\varepsilon'}(X_v)$. Fix finally $\zeta,\zeta'\in\mathfrak{S}$. Let $f_v$
denote the function of the unit:
\begin{itemize}
\item for $v$ in the first row of \eqref{4.23:eq-unit-oscillation-a}, the format function
$F^{\mathfrak{b}}_v$ of the stored unit $v$, read as $F^{\mathfrak{b}}_v(\cdot\,;g')$ when $v$
carries live slots;
\item for $v\in\mathfrak{L}^{\mathrm{pt}}$, the function $f_v$ of
Proposition~\ref{4.23:prop-pointed-depth-decay};
\item for $v=v_e$, the format function $F^{\mathfrak{b}}_{v_e}$ of the Gaussian family;
\item for $v=v'\in\mathfrak{K}(Y')$, its format function $F^{\mathfrak{b}}_{v'}$.
\end{itemize}
Put
\begin{equation}\label{4.23:eq-unit-oscillation-1}
G(\vec{\mathsf{s}}):=f_v\bigl((\zeta,\varpi,\vec{\mathsf{s}})\bigr)
-f_v\bigl((\zeta',\varpi,\vec{\mathsf{s}})\bigr).
\end{equation}
The function $G$ is holomorphic on the product polydisc
\begin{equation}\label{4.23:eq-unit-oscillation-2}
\prod_{t=1}^{4}\bigl\{\,|\mathsf{s}^{(t)}(\Delta)|\le e^{\kappa_1}:\ \Delta\subset C^{\sharp}\,\bigr\}.
\end{equation}
The reason depends on the class of $v$:
\begin{itemize}
\item for stored terms, it is the format bound for stored terms on the tube, read together with
the two clauses of the statement that fixes the tube constant $C^{\mathrm{src}}(L)$ which give
holomorphy of the stored terms in the interpolation variables on the polydisc
\eqref{4.23:eq-unit-oscillation-2};
\item for pointed terms, it is the format bound for the pointed class;
\item for the Gaussian families, it is the holomorphy of each $v_e$ on $\mathfrak{P}$;
\item for the units of $\mathfrak{K}(Y')$, it is the holomorphy of $F^{\mathfrak{b}}_{v'}$ on
$\mathfrak{P}$.
\end{itemize}

\emph{The bound $|G|\le\mathfrak{n}(v)$.} The slice $\mathfrak{P}^{\sharp}$ is a product of the
tube, the range of the real variables $\varpi$ and the polydisc of the interpolation variables.
Hence, for every
$\vec{\mathsf{s}}$ in \eqref{4.23:eq-unit-oscillation-2}, the points
$m=(\zeta,\varpi,\vec{\mathsf{s}})$ and $m'=(\zeta',\varpi,\vec{\mathsf{s}})$ both lie in
$\mathfrak{P}$. They have the same $\varpi$, and the same $\vec{\mathsf{s}}$. We bound $|G|$ on the
whole polydisc by each entry of the table that applies to $v$.

The column headed ``always''. For $v$ in the first row of \eqref{4.23:eq-unit-oscillation-a}, the
format bound gives
$|F^{\mathfrak{b}}_v(m)-F^{\mathfrak{b}}_v(1)|\le\mathfrak{N}^{A,\mathbb{C}}(v)$, where $1$ denotes
the reference point of the format bound, at which the presented pair is the identity. Applying it
at $m$ and at
$m'$ and using the triangle inequality gives $|G|\le2\,\mathfrak{N}^{A,\mathbb{C}}(v)$. The interior
clause, with factor two, of the depth-decay lemma for deep constituents is what this bound yields
after the Cauchy estimate below. For $v\in\mathfrak{L}^{\mathrm{pt}}$, the format bound for the
pointed class, applied at $m$ and at $m'$, gives $|G|\le2\,\mathfrak{N}^{\flat}(v)$ in the same way.
For $v=v_e\in\mathfrak{Q}(\mathfrak{a})$, the entry of that column comes in the same way from the
format bound at both points, with the datum $2\nu^{\natural}(e)$ of the norm bound for the non-deep
pieces of the Gaussian families: applied at $m$ and at $m'$ it gives $|G|\le2\cdot2\,\nu^{\natural}(e)$.
After the Cauchy estimate below this is the content of that norm bound, which carries the factor
two and no depth factor. The units of $\mathfrak{K}(Y')$ need no entry here.

The last column. Let $\hat X_v\subset\mathsf{C}$.
\begin{itemize}
\item For $v$ in the first row of \eqref{4.23:eq-unit-oscillation-a} we use the bound for the
function in the depth-decay lemma for deep constituents. That lemma is stated for every deep term
and all
$m=(\zeta,\varpi,\vec{\mathsf{s}})$, $m'=(\zeta',\varpi,\vec{\mathsf{s}}')\in\mathfrak{P}$ with the
same $\varpi$, whether $\zeta'=\zeta$ or not. The unit $v$ is a deep term in the sense of that
lemma, since $X_v=\hat X_v\subset\mathsf{C}$. Here we take
$\vec{\mathsf{s}}'=\vec{\mathsf{s}}$, the datum $\mathfrak{N}^{A,\mathbb{C}}(v)$ and the depth
$\hat D_v=d^{\wedge}(X_v,\mathsf{C}^c)$. It gives
$|G|\le C_{3F}e^{-\frac12\delta_P\hat D_v}\mathfrak{N}^{A,\mathbb{C}}(v)$.
\item For $v\in\mathfrak{L}^{\mathrm{pt}}$, Proposition~\ref{4.23:prop-pointed-depth-decay} with
$\vec{\mathsf{s}}'=\vec{\mathsf{s}}$ gives
$|G|\le C^{\mathrm{pt}}e^{-\frac12\delta_P\hat D_v}\mathfrak{N}^{\flat}(v)$.
\item For $v'\in\mathfrak{K}(Y')$ we use the oscillation bound for the units of $\mathfrak{K}(Y')$.
It states that, for all $m=(\zeta,\varpi,\vec{\mathsf{s}})$,
$m'=(\zeta',\varpi,\vec{\mathsf{s}}')\in\mathfrak{P}$ with the same $\varpi$, whether
$\zeta'=\zeta$ or not, $|F^{\mathfrak{b}}_{v'}(m)-F^{\mathfrak{b}}_{v'}(m')|$ is
at most $\tfrac12 C_{3F}e^{-\frac12\delta_P\hat D_{v'}}\mathfrak{O}(v')$, with
$\hat D_{v'}=d^{\wedge}(Y'^{\sharp},\mathsf{C}^{c})$. Applied with
$\vec{\mathsf{s}}'=\vec{\mathsf{s}}$ it gives the entry of the last line of the table.
\end{itemize}

$(\ast)$ For $v=v_e\in\mathfrak{Q}(\mathfrak{a})$ with $\hat X_v\subset\mathsf{C}$, the entry of the
last column requires the
depth-carrying bounds for deep constituents, transferred to the Gaussian
family $v_e$ with $\mathfrak{N}$ replaced by $\mathfrak{N}^{G,L}(e)$, and read, as in the
depth-decay lemma for deep constituents, with $\zeta'\ne\zeta$
allowed. Precisely, it requires the two bounds
\begin{equation}\label{4.23:eq-unit-oscillation-b}
\begin{aligned}
\bigl|F^{\mathfrak{b}}_{v_e}(m)-F^{\mathfrak{b}}_{v_e}(m')\bigr|
&\le C_{3F}\,e^{-\frac12\delta_P\hat D_v}\,\mathfrak{N}^{G,L}(e),\\
\bigl|\mathfrak{p}(v_e;\vec{\mathsf{y}})(\zeta)-\mathfrak{p}(v_e;\vec{\mathsf{y}})(\zeta')\bigr|
&\le C_{3F}\,e^{-\frac12\delta_P\hat D_v}\,\mathfrak{N}^{G,L}(e)\,
e^{-(\kappa_1-1)\sum_t|\mathsf{y}_t|},
\end{aligned}
\end{equation}
the first for all $m=(\zeta,\varpi,\vec{\mathsf{s}})$ and
$m'=(\zeta',\varpi,\vec{\mathsf{s}}')\in\mathfrak{P}$ with the same $\varpi$, the second for every
tuple $\vec{\mathsf{y}}$, $\vec{\emptyset}$ included, and all $\zeta,\zeta'$ in the source tube
$\mathfrak{S}_{\mathrm{src}}$ over the filing domain of $v_e$, that is, the space
$\tilde{U}^{c}_{k}(\cdot,\tilde\alpha_0,\tilde\alpha_1)$ over that domain, as $\mathfrak{S}$ is the
one over $C^{\sharp}$.
The first line is used for $G$ here; the second is its counterpart for the pieces. The interior
clause with factor two is not the one required: that clause carries no depth factor.

The transfer to the Gaussian families gives the depth bounds for deep constituents at pairs with
the same $(\zeta,\varpi)$. Their extension \eqref{4.23:eq-unit-oscillation-b} to pairs with the
same $\varpi$ only is what the entry requires, and it is not proved in this book. Every later
estimate of this chapter that uses the entry of the last column for the Gaussian families
depends on it. Granting
\eqref{4.23:eq-unit-oscillation-b}, its first line at $\vec{\mathsf{s}}'=\vec{\mathsf{s}}$ gives
$|G|\le C_{3F}e^{-\frac12\delta_P\hat D_v}\mathfrak{N}^{G,L}(e)$.

When $\hat X_v\subset\mathsf{C}$ and the line of $v$ has two entries, both bound $|G|$ on the whole
polydisc, hence so does their minimum; for $v'\in\mathfrak{K}(Y')$ the entry of the last column
is the only one. In all cases
\begin{equation}\label{4.23:eq-unit-oscillation-3}
|G(\vec{\mathsf{s}})|\le\mathfrak{n}(v)\qquad\text{for all }\vec{\mathsf{s}}
\text{ in \eqref{4.23:eq-unit-oscillation-2}.}
\end{equation}

\emph{From $G$ to the pieces.} The piece operator is linear in the function it acts on. Hence
\begin{equation}\label{4.23:eq-unit-oscillation-4}
\mathfrak{p}(v;\vec{\mathsf{y}})(\zeta)-\mathfrak{p}(v;\vec{\mathsf{y}})(\zeta')
=\prod_{t=1}^{4}\prod_{\Delta\in\mathsf{y}_t}\int_0^1 d\mathsf{s}^{(t)}(\Delta)\,
\partial_{\mathsf{s}^{(t)}(\Delta)}\,G\Big|_{\mathsf{s}=0\text{ off }\vec{\mathsf{y}}},
\end{equation}
where the variables of cells not in $\mathsf{y}_t$ are set to $0$.

At $\vec{\mathsf{y}}=\vec{\emptyset}$ the right side is $G(\vec 0)$. By
\eqref{4.23:eq-unit-oscillation-3} its modulus is at most $\mathfrak{n}(v)$, so the values of the
unit are covered: $G$ itself is small.

Otherwise, let $n=\sum_t|\mathsf{y}_t|$. There is one variable $\mathsf{s}(\Delta)$ for each cell
of $\bigcup_t\mathsf{y}_t$, counted with its index $t$, so there are $n$ variables. The mixed
derivative in \eqref{4.23:eq-unit-oscillation-4} is evaluated at points with
$|\mathsf{s}(\Delta)|\le1$. Apply Cauchy's formula in each of these variables on the circle
\begin{equation*}
|\sigma(\Delta)-\mathsf{s}(\Delta)|=e^{\kappa_1}-1 .
\end{equation*}
Since $|\mathsf{s}(\Delta)|\le1$, every point of these circles satisfies
$|\sigma(\Delta)|\le e^{\kappa_1}$, so the product of circles lies in the closed polydisc
\eqref{4.23:eq-unit-oscillation-2}. If a holomorphy statement used above gives holomorphy on the
open polydisc only, run the circles at radius $(1-\varepsilon)(e^{\kappa_1}-1)$ and let
$\varepsilon\downarrow0$. With \eqref{4.23:eq-unit-oscillation-3} this gives
\begin{equation*}
\Bigl|\prod\partial\,G\Bigr|\le\mathfrak{n}(v)\,(e^{\kappa_1}-1)^{-n}.
\end{equation*}
The $n$ integrals over $[0,1]$ cost at most $1$.

Moreover,
\begin{equation*}
e^{\kappa_1}-1\ge e^{\kappa_1-1}
\iff e^{\kappa_1}\Bigl(1-\frac1e\Bigr)\ge1
\iff \kappa_1\ge\log\frac{e}{e-1}.
\end{equation*}
Here $\log(e/(e-1))<1$, because $e/(e-1)<e$ is equivalent to $e>2$; indeed
$\log(e/(e-1))<\tfrac12$. The lower bound $\kappa_1-1\ge3c_l+\lambda_*>0$, already imposed on
$\kappa_1$ in the order of choice, where $c_l$ and $\lambda_*$ are the growth constant and the rate
floor of the summation lemma for the index families used below, gives $\kappa_1>1$. So the
condition $\kappa_1\ge\log(e/(e-1))$
is implied by it, and no new condition on $\kappa_1$ arises. Taking the supremum over real
$\Phi_C$, over the source parameter and the weak coordinates, over $\zeta,\zeta'\in\mathfrak{S}$,
and over $g'$ for units with live slots,
we obtain, for every tuple with $\mathrm{lab}_4\subset C^{\sharp}$, $\vec{\emptyset}$ included,
\begin{equation}\label{4.23:eq-unit-oscillation-5}
\operatorname{osc}\bigl(\mathfrak{p}(v;\vec{\mathsf{y}})\bigr)
\le\mathfrak{n}(v)\,e^{-(\kappa_1-1)\sum_t|\mathsf{y}_t|}.
\end{equation}

\emph{The four free sums.} Each $\mathsf{y}_t$ ranges over the index family that
Definition~\ref{4.23:def-component-constituents} attaches at stage $t$ to the label
$\mathrm{lab}_{t-1}$ and its accompanying datum, together with $\emptyset$. The hypothesis (R)
that the summation lemma places on its index families holds for these families: it is part of the
statements that construct
the index sets of the three classes. Put $\lambda:=\kappa_1-1$. We use two clauses of the
summation lemma. Its clause (b) gives, for every $t\in\{1,\dots,4\}$ and every rate
$\lambda'\ge\lambda_*$,
\begin{equation*}
1+\sum_{\mathsf{y}_t\ne\emptyset}e^{-\lambda'|\mathsf{y}_t|}\le e^{d_k(\mathrm{lab}_{t-1})+1}.
\end{equation*}
Its clause (a)
gives
\begin{equation*}
d_k(\mathrm{lab}_t)\le d_k(\mathrm{lab}_{t-1})+c_l|\mathsf{y}_t|.
\end{equation*}
For $t=0,\dots,3$ let $\Sigma_t$ be the sum of
$\exp\bigl(-\lambda\sum_{u>t}|\mathsf{y}_u|\bigr)$ over $\mathsf{y}_{t+1},\dots,\mathsf{y}_4$. It
is a function of $\mathrm{lab}_t$. Clause (b) at $t=4$ and the rate $\lambda$ gives
$\Sigma_3\le e^{d_k(\mathrm{lab}_3)+1}$. We
claim $\Sigma_{t}\le e^{(4-t)(d_k(\mathrm{lab}_t)+1)}$, and prove it by descending induction on
$t$.

For $t=3$ the claim is the bound just stated. Suppose it holds at $t$, where $1\le t\le3$. By
clause (a), and then clause (b) at the rate $\lambda-(4-t)c_l$,
\begin{align*}
\Sigma_{t-1}
&=\sum_{\mathsf{y}_t}e^{-\lambda|\mathsf{y}_t|}\,\Sigma_t
\le\sum_{\mathsf{y}_t}e^{-\lambda|\mathsf{y}_t|}\,
e^{(4-t)(d_k(\mathrm{lab}_{t-1})+c_l|\mathsf{y}_t|+1)}\\
&=e^{(4-t)(d_k(\mathrm{lab}_{t-1})+1)}
\Bigl[1+\sum_{\mathsf{y}_t\ne\emptyset}e^{-(\lambda-(4-t)c_l)|\mathsf{y}_t|}\Bigr]
\le e^{(5-t)(d_k(\mathrm{lab}_{t-1})+1)}.
\end{align*}
For $t=3$ the right side is
\begin{equation*}
e^{d_k(\mathrm{lab}_2)+1}\Bigl[1+\sum_{\mathsf{y}_3\ne\emptyset}
e^{-(\lambda-c_l)|\mathsf{y}_3|}\Bigr]\le e^{2(d_k(\mathrm{lab}_2)+1)},
\end{equation*}
and the steps $t=2,1$ are the same computation at the rates $\lambda-2c_l$ and $\lambda-3c_l$.

The rates used are $\lambda,\lambda-c_l,\lambda-2c_l,\lambda-3c_l$. The lower bound
$\kappa_1-1\ge3c_l+\lambda_*$ gives $\lambda-3c_l\ge\lambda_*$, and since $c_l\ge0$ (the growth
constant of the summation lemma), every one of these rates is at least $\lambda_*$. Hence
\begin{equation*}
\Sigma_0\le e^{4(d_k(\mathrm{lab}_0)+1)}=e^{4(d_k(K_v)+1)}.
\end{equation*}
Finally, all terms are non-negative, so restricting to tuples with $\mathrm{lab}_4\subset C^{\sharp}$
only decreases the sum. Summing \eqref{4.23:eq-unit-oscillation-5} over these tuples gives
\begin{equation*}
\operatorname{tot}(v)\le\mathfrak{n}(v)\,\Sigma_0\le\mathfrak{n}(v)\,e^{4(d_k(K_v)+1)},
\end{equation*}
which is the last line of \eqref{4.23:eq-unit-oscillation-a}.
\end{proof}

\begin{remark}
Only tuples with $\mathrm{lab}_4\subset C^{\sharp}$ enter $\operatorname{tot}(v)$. The reason is
that the pieces of the remaining tuples are not functions on $\mathfrak{P}^{\sharp}$. The later
summation over the component needs only $\mathrm{lab}_4\subset C$, so $\operatorname{tot}(v)$ is a
majorant for what it uses. Tuples with $\mathrm{lab}_4\subset C^{\sharp}$ but
$\mathrm{lab}_4\not\subset C$ are counted here and again in the estimate of the rim
$\mathfrak{L}_C$. All terms are non-negative, so this double counting only enlarges the majorant.
\end{remark}

\begin{remark}
The following alternative for the step $(\ast)$ rests on a hypothesis that is itself not
established in this book. The Gaussian family $v_e$ has the format of a stored term, with region
$A_e:=[S^{+}(e)]_{m_0(e)}$, the region attached to $S^{+}(e)$ at the index $m_0(e)$ that the
construction of the Gaussian families assigns to $e$, and with the set $\mathcal{P}$ of that format
empty. It is holomorphic on
$\mathfrak{P}$, and it depends on the chain of localised terms only through its restriction to
$S^{+}(e)\subset\mathsf{C}$.

Suppose that the format statement for $v_e$ gives the following: $F^{\mathfrak{b}}_{v_e}$ is a
holomorphic function of the presented pair on $S^{+}(e)$, and, less its value at $1$, it is
bounded by $\mathfrak{N}^{G,L}(e)$ at every pair satisfying the three bounds of
\eqref{4.23:eq-paired-tube-points-a}. Put
\begin{equation*}
G_e(\lambda):=F^{\mathfrak{b}}_{v_e}(X,P_\lambda)-F^{\mathfrak{b}}_{v_e}(X,1),
\end{equation*}
the first argument $X$ being the region argument of the format function of $v_e$, held fixed,
and $\lambda$ being the parameter of the pairs $P_\lambda$ of
Lemma~\ref{4.23:lem-paired-tube-points}, not the rate of the proof above.
By Lemma~\ref{4.23:lem-paired-tube-points}, $G_e$ is then holomorphic, with
$|G_e|\le\mathfrak{N}^{G,L}(e)$, on the disc $|\lambda|<\rho_{S^{+}(e)}$, where
\begin{equation*}
\rho_{S^{+}(e)}=r^0e^{\frac12\delta_P\hat D_v}.
\end{equation*}
Schwarz's lemma, applied exactly as in the proof of
Proposition~\ref{4.23:prop-pointed-depth-decay}, then gives
\begin{equation*}
|G_e(1)-G_e(0)|\le\frac{2}{r^0}\,e^{-\frac12\delta_P\hat D_v}\,\mathfrak{N}^{G,L}(e)
\le C^{\mathrm{pt}}e^{-\frac12\delta_P\hat D_v}\,\mathfrak{N}^{G,L}(e),
\end{equation*}
since $C^{\mathrm{pt}}\ge2/r^0$ by its definition in
Proposition~\ref{4.23:prop-pointed-depth-decay}. This is the entry of the last column for $v_e$
with $C^{\mathrm{pt}}$ in place of $C_{3F}$. The hypothesis made on the format statement for $v_e$ is
not proved here, so the step $(\ast)$ is not proved on this route either.
\end{remark}

\begin{lemma}[The cell sum for plain units]\label{4.23:lem-plain-cell-sum}
Let $C$ be a component at level $k$, with core $\mathsf{C}=\mathsf{C}_C$, window $N=N(g_k)$,
bottom level $h=k-N$ and window radius $\check{R}_{h+1}$. Put $r:=4$. Since
$\kappa\ge 180$, one has $r\le\frac{8}{5}\kappa$, which is the range in which the deep
and the shallow plain-unit sum bounds used below are stated.

In what follows, $b_\delta$ is the per-cube constant with which the stored boundary-type terms
are filed; $\mathsf{K}^{\flat}$ is the constant of the collar bound for plain units (that bound
also carries a rim count, which does not enter below); $C_{\uparrow}$ is the constant of the
window-multiplicity bound recalled in the
remark below; $C_{3F}$ and $C^{\mathrm{pt}}$ are the constants entering the deep column of the
table of total oscillation of Lemma~\ref{4.23:lem-unit-oscillation}; and $\mathsf{b}_k(\cdot)$,
$C^{\mathrm{pl}}_{\mathsf{a}}$, $C_{\mathrm{cell}}$, $C_{\circledast}$, $C^{<}$, $C_{\mathsf{b}}$,
$p_{\mathsf{b}}$ are the function and the constants of the deep and shallow plain-unit sum bounds.
For a plain unit $v\in\mathfrak{L}^{\mathrm{pl}}(C)=\mathfrak{L}^{B}\sqcup\mathfrak{L}^{\mathrm{pt}}$
(Definition~\ref{4.23:def-component-constituents}) and for $s\ge 0$ put
\begin{equation*}
W(v):=\mathfrak{N}^{\flat}(v)\,e^{r\,d_k(K_v)},\qquad
\mathsf{b}(s):=\mathsf{b}_k(s)+(s+1)\,b_\delta ,
\end{equation*}
where $K_v$ and $d_k(K_v)$ are as in the total-oscillation bound
\eqref{4.23:eq-unit-oscillation-a} of Lemma~\ref{4.23:lem-unit-oscillation}: $d_k(K_v)$ is the
quantity attached there to the unit $v$ that enters the exponent of that bound.
Define the counts $\mathsf{n}_\partial(k)$, $\mathsf{n}_{\mathrm{sh}}(k)$ and the constants
$\mathsf{S}^{\mathrm{dp}}_{\mathrm{pl}}(k)$, $\mathsf{S}^{\mathrm{sh}}_{\mathrm{pl}}(k)$, $C^+$ by
\begin{equation}\label{4.23:eq-plain-cell-sum-a}
\begin{split}
\mathsf{n}_\partial(k)&:=2d\,(104L\check{R}_{h+1}+2)^{d-1},\qquad
\mathsf{n}_{\mathrm{sh}}(k):=(N+1)(110L\check{R}_{h+1})^{d},\\
\mathsf{S}^{\mathrm{dp}}_{\mathrm{pl}}(k)&:=e\cdot C^{\mathrm{pl}}_{\mathsf{a}}C_{\mathrm{cell}}
\cdot 2d\,(104L\check{R}_{h+1}+2)^{d-1}(2+N)\,\mathsf{b}(N+1),\\
\mathsf{S}^{\mathrm{sh}}_{\mathrm{pl}}(k)&:=(N+1)(110L\check{R}_{h+1})^{d}
\Bigl[(N+1)\bigl(1+C_{\circledast}\mathsf{K}^{\flat}(2LC_{\uparrow}\check{R}_{h+1})^{3}\bigr)
+C^{<}\\
&\qquad\qquad
+2C_{\mathsf{b}}(1+N)^{p_{\mathsf{b}}}(L\check{R}_{h+1})^{d}\Bigr]\,\mathsf{b}(N),\\
C^{+}&:=\max\{C_{3F},C^{\mathrm{pt}}\}.
\end{split}
\end{equation}
One has $C^{+}\ge 2$. Thus $\mathsf{S}^{\mathrm{dp}}_{\mathrm{pl}}(k)$ is
$e\,C^{\mathrm{pl}}_{\mathsf{a}}C_{\mathrm{cell}}\,\mathsf{n}_\partial(k)$ times
$(2+N)\,\mathsf{b}(N+1)$, and $\mathsf{S}^{\mathrm{sh}}_{\mathrm{pl}}(k)$ is
$\mathsf{n}_{\mathrm{sh}}(k)$ times the bracket times $\mathsf{b}(N)$.

Call $v\in\mathfrak{L}^{\mathrm{pl}}(C)$ \emph{deep} if $X_v\subset\mathsf{C}$ and \emph{shallow}
otherwise; here $X_v$ is the domain of the unit $v$, and $\hat{X}_v$ and $\hat{D}_v\ge 0$ are the
second domain and the depth parameter of $v$ that enter the table of
\eqref{4.23:eq-unit-oscillation-a} (the latter is distinct from the cell depth
$\varrho(\Delta)$). When $\hat{X}_v\not\subset\mathsf{C}$ we set
$\hat{D}_v:=0$. Put
\begin{equation*}
S^{\mathrm{dp}}(C):=\sum_{v\ \mathrm{deep}}W(v)\,e^{-\frac14\delta_P\hat{D}_v},\qquad
S^{\mathrm{sh}}(C):=\sum_{v\ \mathrm{shallow}}W(v),
\end{equation*}
both sums running over the units of $\mathfrak{L}^{\mathrm{pl}}(C)$. Then:
\begin{enumerate}
\item[(a)] $S^{\mathrm{dp}}(C)\le\mathsf{S}^{\mathrm{dp}}_{\mathrm{pl}}(k)$ and
$S^{\mathrm{sh}}(C)\le\mathsf{S}^{\mathrm{sh}}_{\mathrm{pl}}(k)$;
\item[(b)] with $\operatorname{tot}(v)$ the total oscillation of the unit $v$
(Lemma~\ref{4.23:lem-unit-oscillation}, with the oscillation seminorm of
Definition~\ref{4.23:not-oscillation-seminorm}),
\begin{equation}\label{4.23:eq-plain-cell-sum-1}
\sum_{v\in\mathfrak{L}^{\mathrm{pl}}(C)}\operatorname{tot}(v)
\ \le\ e^{4}\bigl[C^{+}S^{\mathrm{dp}}(C)+2S^{\mathrm{sh}}(C)\bigr]
\ \le\ e^{4}\bigl[C^{+}\mathsf{S}^{\mathrm{dp}}_{\mathrm{pl}}(k)
+2\,\mathsf{S}^{\mathrm{sh}}_{\mathrm{pl}}(k)\bigr].
\end{equation}
\end{enumerate}
\end{lemma}

\begin{proof}
\emph{Clause} (a). The first inequality is the deep plain-unit sum bound and the second is the
shallow plain-unit sum bound, each taken by statement, with $W(v)$, $\mathsf{b}(s)$, the notion
of a deep unit and the convention $\hat{D}_v:=0$ exactly as defined above; both bounds are stated
for $r\le\frac{8}{5}\kappa$, which holds for $r=4$ because $\kappa\ge 180$. Their right-hand
sides are, by definition, the constants $\mathsf{S}^{\mathrm{dp}}_{\mathrm{pl}}(k)$ and
$\mathsf{S}^{\mathrm{sh}}_{\mathrm{pl}}(k)$ of \eqref{4.23:eq-plain-cell-sum-a}.

\emph{Clause} (b). Let $v\in\mathfrak{L}^{\mathrm{pl}}(C)$. By the total-oscillation bound
\eqref{4.23:eq-unit-oscillation-a} of Lemma~\ref{4.23:lem-unit-oscillation}, whose proof is
incomplete at one step,
\begin{equation}\label{4.23:eq-plain-cell-sum-2}
\operatorname{tot}(v)\ \le\ \mathfrak{n}(v)\,e^{4(d_k(K_v)+1)}
\ =\ e^{4}\,\mathfrak{n}(v)\,e^{r\,d_k(K_v)},
\end{equation}
where the equality uses $r=4$. It remains to bound $\mathfrak{n}(v)$ by a multiple of
$\mathfrak{N}^{\flat}(v)$, separately for deep and for shallow units.

Let $v$ be deep, so that $X_v\subset\mathsf{C}$. In the table of
\eqref{4.23:eq-unit-oscillation-a}, $\mathfrak{n}(v)$ is a minimum over the columns.

If $\hat{X}_v\subset\mathsf{C}$, then $\mathfrak{n}(v)$ is at most the entry of the
deep column. That entry is a constant, $C_{3F}$ or $C^{\mathrm{pt}}$, times
$e^{-\frac12\delta_P\hat{D}_v}$ times the format datum of $v$, and in every case that datum is
$\mathfrak{N}^{\flat}(v)$:
\begin{itemize}
\item For $v\in\mathfrak{L}^{\mathrm{pt}}$ the format datum is $\mathfrak{N}^{\flat}(v)$ itself.
\item For $v\in\mathfrak{L}^{B}$ other than the distinguished unit $v_e$ described in the next
item, the format datum is $\mathfrak{N}^{A,\mathbb{C}}(v)$, and on
$\mathfrak{L}^{B}$ one has $\mathfrak{N}^{A,\mathbb{C}}(v)=\mathfrak{N}^{\flat}(v)$.
\item For $v_e$, the unit of $\mathfrak{L}^{B}$ whose entry in the table carries as its format
datum the datum with which the arrest terms of $\mathfrak{L}^{B}$ are filed, that datum is
$\mathfrak{N}^{\flat}(v_e)$ itself.
\end{itemize}
Each constant is at most $C^{+}=\max\{C_{3F},C^{\mathrm{pt}}\}$. Since $\delta_P\hat{D}_v\ge 0$,
one has $e^{-\frac12\delta_P\hat{D}_v}\le e^{-\frac14\delta_P\hat{D}_v}$, and hence
\begin{equation}\label{4.23:eq-plain-cell-sum-3}
\mathfrak{n}(v)\ \le\ C^{+}e^{-\frac14\delta_P\hat{D}_v}\,\mathfrak{N}^{\flat}(v).
\end{equation}

If $\hat{X}_v\not\subset\mathsf{C}$, then $\hat{D}_v=0$ by convention, and the table
gives $\mathfrak{n}(v)=2\mathfrak{N}^{\flat}(v)$. Because $C^{+}\ge 2$, this is at most
$C^{+}\mathfrak{N}^{\flat}(v)=C^{+}e^{-\frac14\delta_P\hat{D}_v}\mathfrak{N}^{\flat}(v)$, so
\eqref{4.23:eq-plain-cell-sum-3} holds in this case as well. Inserting
\eqref{4.23:eq-plain-cell-sum-3} into \eqref{4.23:eq-plain-cell-sum-2} and using the definition
of $W(v)$ gives, for every deep $v$,
\begin{equation*}
\operatorname{tot}(v)\ \le\ e^{4}C^{+}\,W(v)\,e^{-\frac14\delta_P\hat{D}_v}.
\end{equation*}

Let $v$ be shallow. The table gives $\mathfrak{n}(v)\le 2\mathfrak{N}^{\flat}(v)$, and
\eqref{4.23:eq-plain-cell-sum-2} yields $\operatorname{tot}(v)\le 2e^{4}\,W(v)$.

Every $v\in\mathfrak{L}^{\mathrm{pl}}(C)$ is either deep or shallow, and not both. Summing the
two bounds over the deep and over the shallow units of $\mathfrak{L}^{\mathrm{pl}}(C)$ gives
\begin{equation*}
\sum_{v\in\mathfrak{L}^{\mathrm{pl}}(C)}\operatorname{tot}(v)
\ \le\ e^{4}C^{+}\sum_{v\ \mathrm{deep}}W(v)\,e^{-\frac14\delta_P\hat{D}_v}
+2e^{4}\sum_{v\ \mathrm{shallow}}W(v)
\ =\ e^{4}\bigl[C^{+}S^{\mathrm{dp}}(C)+2S^{\mathrm{sh}}(C)\bigr].
\end{equation*}
This is the first inequality of \eqref{4.23:eq-plain-cell-sum-1}. The second follows from
clause (a), since $e^{4}C^{+}>0$ and $2e^{4}>0$.
\end{proof}

\begin{remark}
The constants of \eqref{4.23:eq-plain-cell-sum-a} account for the following. The sum at one
cell collects the contributions of all levels $k-n+1$, with $n$ ranging over the level indices of
the plain-unit sum bounds; this sum over $n$ is controlled by the first level-sum bound. It also
carries the window multiplicity
$\mathsf{w}(\Delta)^{d}$ of a cell $\Delta$ in the thin strata of an old arrest and the
weights $\mathsf{w}^{3}$ of the pointed units. These are absorbed by the window bound and its
companion bound, through $\mathsf{w}(\Delta)\le LC_{\uparrow}\varrho(\Delta)$. The cells of the
core are summed by part (b) of the cell-summation lemma. After that summation only the
boundary cells survive, and there are at most $\mathsf{n}_\partial(k)$ of them.
\end{remark}

\begin{lemma}[The cell sum for wrapped units]\label{4.23:lem-wrapped-cell-sum}
Let $C$ be a component of the large-field region $Z$ at level $k$, let
$\mathsf{C}=\mathsf{C}_C$ be its core, let $N=N(g_k)$, assume $N\ge1$, put $h=k-N$, and put
$\varsigma_0=\frac15\kappa$. For a unit $v$ write $K_v$ for its label set, $X_v$ for its
region, $A_v$ for its support, $j_v$ for its level, $\mathfrak{N}^{A}(v)$ for its stored
amplitude (the format datum of the unit), $\mathcal{P}_v$ for its family of pockets, and
$\mathcal{P}_v^{\circ}$ for the subfamily of $\mathcal{P}_v$ that enters clause (iii) of
Lemma~\ref{4.23:lem-filing-domains}. For a pocket $(Y,m)$ write $Y^{\sharp}$ for its strip. For a level $j$, a
\emph{$\pi_j$-cube} is a cube of the block partition $\pi_j$ at level $j$, and
$\operatorname{dist}_h$ is distance in level-$h$ units. Assume the following.
\begin{enumerate}
\item[(a)] \emph{The cut of the wrapped first part into deep and shallow units.} Relative to
the component $C$, the wrapped units $v$ of the first part with $K_v\subset C$ split into
the deep class $\mathfrak{L}^{\mathrm{dp}}$, those with $X_v\subset\mathsf{C}$, and the
shallow class $\mathfrak{L}^{\mathrm{sh}}$, the rest. Every $v\in\mathfrak{L}^{\mathrm{sh}}$
satisfies $X_v\not\subset\mathsf{C}$, and $\mathcal{P}_v$ contains a pocket of $C$.
\item[(b)] \emph{The deep-sum bound with decay rate $\frac14\delta_P$ in the own-scale
contour distance.} With the own-scale contour distance $d^{\wedge}$, the constant
$C_{\mathsf{a}}(\cdot)$ of this bound evaluated at the argument $1$, the cell-sum constant
$C_{\mathrm{cell}}$ and the function
$\mathsf{b}_k(\cdot)$ of the format,
\begin{equation*}
\begin{split}
&\sum_{v\in\mathfrak{L}^{\mathrm{dp}},\,X_v\subset\mathsf{C}}\mathfrak{N}^{A}(v)\,
e^{4d_k(K_v)}\,e^{-\frac14\delta_P d^{\wedge}(X_v,\mathsf{C}^c)}\\
&\qquad\le e\,C_{\mathsf{a}}(1)\,C_{\mathrm{cell}}\,2d\,
(104L\check{R}_{h+1}+2)^{d-1}(2+N)\,\mathsf{b}_k(N+1).
\end{split}
\end{equation*}
\item[(c)] \emph{The anchored weight bound with exponent $4$ on $d_k(K_v)$.} For every
level $j\le k$ and every
$\pi_j$-cube $Q$, the terms $F$ of the format with level $j_F$, support $A_F$ and weight
$W(F)$ satisfy
\begin{equation*}
\sum_{F:\,j_F=j,\,A_F\ni Q}W(F)\,e^{\varsigma_0 d_j(A_F)\mathbf{1}[j<k]}
\le\mathsf{b}_k(k-j),
\end{equation*}
a unit $v$ entering as the term $F=v$, with $j_F=j_v$, $A_F=A_v$ and weight
$W(v)=\mathfrak{N}^{A}(v)e^{4d_k(K_v)}$.
\item[(d)] \emph{The geometry of the core.} The strip $Y^{\sharp}$ of every pocket $(Y,m)$
of $C$ lies in $\mathsf{C}$, and
$\operatorname{dist}_h(Y^{\sharp},\mathsf{C}^c)\ge\frac12(\check{R}_h+1)$. For every level
$j\le k$, with $u=h-j$ (so that $L^{u}$ converts units between the levels $j$ and $h$),
at most $(104L\check{R}_{h+1}L^{u}+3)^d$ cubes of $\pi_j$ lie in or touch $\mathsf{C}$.
\item[(e)] \emph{The distance estimate at a pocket strip.} Let $j\le h$ and $u=h-j$. If a
set $A$ contains a $\pi_j$-cube, of side at most one level-$h$ unit, touching the strip
$Y^{\sharp}$ of a pocket of $C$, and also contains a point off $\mathsf{C}$, then
$d_j(A)\ge\bigl(\frac12(\check{R}_h+1)-1\bigr)L^{u}-1$.
\item[(f)] \emph{The growth bound.} For every $u\ge0$,
$\mathsf{b}_k(N+u)\le C_{\mathsf{b}}(1+u)^{p_{\mathsf{b}}}\mathsf{b}_k(N)$; and
$\mathsf{b}_k(s)\le\mathsf{b}_k(N)$ for $0\le s<N$.
\item[(g)] \emph{The window radius and the blocking factor.} $\check{R}_h\ge38$ and
$\check{R}_{h+1}\le\check{R}_h$; moreover, since $L^d\ge2$,
$\sum_{i\ge1}L^{d(1-i)}=(1-L^{-d})^{-1}\le2$.
\item[(h)] \emph{The depth factor of a deep wrapped unit.} For every
$v\in\mathfrak{L}^{\mathrm{dp}}$, the depth $\hat{D}=\hat{D}(v)$ entering the table of
Lemma~\ref{4.23:lem-unit-oscillation} satisfies
\begin{equation*}
e^{-\frac12\delta_P\hat{D}(v)}\le e^{-\frac14\delta_P d^{\wedge}(X_v,\mathsf{C}^c)}.
\end{equation*}
\end{enumerate}
Define
\begin{equation}\label{4.23:eq-wrapped-cell-sum-a}
\begin{aligned}
\mathsf{S}^{\mathrm{dp}}_{\mathrm{wr}}(k)&:=e\,C_{\mathsf{a}}(1)\,C_{\mathrm{cell}}\,2d\,
(104L\check{R}_{h+1}+2)^{d-1}(2+N)\,\mathsf{b}_k(N+1),\\
\mathsf{S}^{\mathrm{sh}}_{\mathrm{wr}}(k)&:=\bigl[C^{\mathrm{sh}}
+2^d\bigl(2(104\check{R}_{h+1})^d+3^dN\bigr)\bigr]\,\mathsf{b}_k(N),\\
C^{\mathrm{sh}}&:=C_{\mathsf{b}}\sup_{R\ge38}\sum_{u\ge0}(104LRL^{u}+3)^d
(1+u)^{p_{\mathsf{b}}}\,e^{-\varsigma_0(\frac12(R-1)L^{u}-1)}.
\end{aligned}
\end{equation}
Then $C^{\mathrm{sh}}<\infty$ is a number depending on $(d,L,\kappa,p_{\mathsf{b}},
C_{\mathsf{b}})$ only (it is not required to be small), and:
\begin{enumerate}
\item[(1)] the sum in \textup{(b)} is at most $\mathsf{S}^{\mathrm{dp}}_{\mathrm{wr}}(k)$;
\item[(2)] the shallow sum satisfies
\begin{equation}\label{4.23:eq-wrapped-cell-sum-1}
\sum\Bigl\{\mathfrak{N}^{A}(v)e^{4d_k(K_v)}:\ v\ \text{as in (S)}\Bigr\}
\le\mathsf{S}^{\mathrm{sh}}_{\mathrm{wr}}(k),
\end{equation}
where (S) means that $v$ is a wrapped unit of the first part with
$\mathcal{P}_v^{\circ}=\emptyset$, $\mathcal{P}_v$ contains a pocket of $C$, and
$X_v\not\subset\mathsf{C}$;
\item[(3)] with $\mathfrak{L}^{\mathrm{wr}}(C)$ the set of wrapped units among the
constituents filed at $C$ (Definition~\ref{4.23:def-component-constituents}), with
$\operatorname{tot}(v)$ the total oscillation of Lemma~\ref{4.23:lem-unit-oscillation},
and with $C_{3F}$ the constant of the deep wrapped entry of that lemma's table,
\begin{equation}\label{4.23:eq-wrapped-cell-sum-2}
\sum_{v\in\mathfrak{L}^{\mathrm{wr}}(C)}\operatorname{tot}(v)
\le e^{4}\bigl[C_{3F}\,\mathsf{S}^{\mathrm{dp}}_{\mathrm{wr}}(k)
+2\,\mathsf{S}^{\mathrm{sh}}_{\mathrm{wr}}(k)\bigr].
\end{equation}
\end{enumerate}
\end{lemma}

\begin{proof}
\emph{Conclusion \textup{(1)}.} The right side of the inequality in hypothesis (b) is, by
definition \eqref{4.23:eq-wrapped-cell-sum-a}, the number
$\mathsf{S}^{\mathrm{dp}}_{\mathrm{wr}}(k)$.

\emph{Finiteness of $C^{\mathrm{sh}}$.} The series in $C^{\mathrm{sh}}$ converges for every
$R\ge38$, and it is decreasing in $R$ term by term once $\frac12\varsigma_0(R-1)$ exceeds
$d\log(104LR)$. Hence the supremum is finite, and $C^{\mathrm{sh}}$ is a number depending on
$(d,L,\kappa,p_{\mathsf{b}},C_{\mathsf{b}})$ only.

\emph{Conclusion \textup{(2)}: the choice of an anchoring cube.} Let $v$ be as in (S). Fix a
pocket $(Y,m)\in\mathcal{P}_v$ of $C$, and put $j:=j_v$. By hypothesis (d) the strip
$Y^{\sharp}$ of this pocket lies in $\mathsf{C}$. Since $X_v\not\subset\mathsf{C}$, the
support $A_v$ contains a point off $\mathsf{C}$. Moreover $A_v$ contains a $\pi_j$-cube
touching $Y^{\sharp}$; choose one and call it $Q_1=Q_1(v)$. As
$Y^{\sharp}\subset\mathsf{C}$, the cube $Q_1$ lies in or touches $\mathsf{C}$. Grouping the
units as in (S) according to their level $j_v$ and their chosen cube $Q_1(v)$, each of them
is counted exactly once in
\begin{equation*}
\sum_{j\le k}\ \sum_{Q}\ \sum_{v:\,j_v=j,\,Q_1(v)=Q}\mathfrak{N}^{A}(v)e^{4d_k(K_v)},
\end{equation*}
where $Q$ runs over the $\pi_j$-cubes in or touching $\mathsf{C}$, and the inner sum is
restricted to units as in (S). This triple sum is therefore the left side of
\eqref{4.23:eq-wrapped-cell-sum-1}. We split the levels into $j\le h$ and
$h<j\le k$.

\emph{The range $j\le h$.} Put $u:=h-j\ge0$. By hypothesis (d), the cube $Q$ is one of at
most $(104L\check{R}_{h+1}L^{u}+3)^d$ cubes. Fix a unit $v$ of the inner sum, so that
$Q=Q_1(v)$.
The cube $Q_1$ has side at most one level-$h$ unit, because $j\le h$. It touches the strip
$Y^{\sharp}$, and $\operatorname{dist}_h(Y^{\sharp},\mathsf{C}^c)\ge\frac12(\check{R}_h+1)$
by hypothesis (d). The set $A_v$ also contains a point off $\mathsf{C}$. Hypothesis (e)
therefore gives
\begin{equation*}
d_j(A_v)\ge\Bigl(\tfrac12(\check{R}_h+1)-1\Bigr)L^{u}-1
=\tfrac12(\check{R}_h-1)L^{u}-1.
\end{equation*}
On this range $j\le h<k$, so the indicator in hypothesis (c) equals one. With
$W(v)=\mathfrak{N}^{A}(v)e^{4d_k(K_v)}$ as in (c), we obtain
\begin{equation*}
\mathfrak{N}^{A}(v)e^{4d_k(K_v)}
\le W(v)\,e^{\varsigma_0 d_j(A_v)}\,e^{-\varsigma_0(\frac12(\check{R}_h-1)L^{u}-1)}.
\end{equation*}
Since $A_v\ni Q_1(v)=Q$ and $j_v=j$, every unit of the inner sum is a term $F=v$ of the
anchored sum of hypothesis (c) at the cube $Q$. Sum the last inequality over the units of
the inner sum, enlarge the sum to all terms $F$ with $j_F=j$ and $A_F\ni Q$ (all terms are
non-negative), and apply hypothesis (c). Since
$k-j=N+u$, the contribution of one cube $Q$ is at most
\begin{equation*}
e^{-\varsigma_0(\frac12(\check{R}_h-1)L^{u}-1)}\,\mathsf{b}_k(N+u)
\le C_{\mathsf{b}}(1+u)^{p_{\mathsf{b}}}\,
e^{-\varsigma_0(\frac12(\check{R}_h-1)L^{u}-1)}\,\mathsf{b}_k(N),
\end{equation*}
by the growth bound (f). Multiply by the number of cubes. Hypothesis (g) gives
$L\check{R}_{h+1}\le L\check{R}_h$, so the number of cubes is at most
$(104L\check{R}_hL^{u}+3)^d$. Sum over $j\le h$, that is, over
$u\in\{0,\dots,h\}\subset\{u\ge0\}$. The contribution of the range $j\le h$ is then at most
\begin{equation*}
C_{\mathsf{b}}\sum_{u\ge0}(104L\check{R}_hL^{u}+3)^d(1+u)^{p_{\mathsf{b}}}
e^{-\varsigma_0(\frac12(\check{R}_h-1)L^{u}-1)}\,\mathsf{b}_k(N)
\le C^{\mathrm{sh}}\,\mathsf{b}_k(N).
\end{equation*}
For the last inequality, note that $\check{R}_h\ge38$ by hypothesis (g). So the series is
the one in the definition of $C^{\mathrm{sh}}$ at $R=\check{R}_h$, and it is bounded by the
supremum over $R\ge38$.

\emph{The range $h<j\le k$.} Here the level-$h$ box $\mathsf{C}$ is counted in
$\pi_j$-cubes with $j>h$. By hypothesis (d) with $u=h-j=-(j-h)$, there are at most
$(x_j+3)^d$ candidates $Q$, where
\begin{equation*}
x_j:=104L\check{R}_{h+1}L^{-(j-h)}\le104\check{R}_{h+1},
\end{equation*}
since $j-h\ge1$. For a candidate $Q$, every unit $v$ of the inner sum has
$A_v\ni Q_1(v)=Q$ and $j_v=j$, so it is a term $F=v$ of hypothesis (c) at $Q$. The factor
$e^{\varsigma_0 d_j(A_F)\mathbf{1}[j<k]}$ in hypothesis (c) is
at least one. Each candidate therefore contributes at most
$\mathsf{b}_k(k-j)$, and since $0\le k-j<N$ on this range, $\mathsf{b}_k(k-j)\le\mathsf{b}_k(N)$
by hypothesis (f). No decay in $j$ is used on this range; the bound is
a count of cubes. By convexity of $x\mapsto x^d$, $(x_j+3)^d\le2^{d-1}(x_j^d+3^d)$. Next,
\begin{equation*}
\sum_{j>h}x_j^d=(104\check{R}_{h+1})^d\sum_{i\ge1}L^{d(1-i)}
\le2(104\check{R}_{h+1})^d,
\end{equation*}
because $\sum_{i\ge1}L^{d(1-i)}=(1-L^{-d})^{-1}\le2$ by hypothesis (g). Finally,
the range $h<j\le k$ has exactly $N$ levels, so $\sum_{h<j\le k}3^d=3^dN$. Hence
\begin{align*}
\sum_{h<j\le k}(x_j+3)^d
&\le2^{d-1}\sum_{h<j\le k}(x_j^d+3^d)
\le2^{d-1}\bigl[2(104\check{R}_{h+1})^d+3^dN\bigr]\\
&\le2^d\bigl[(104\check{R}_{h+1})^d+3^dN\bigr]
\le2^d\bigl(2(104\check{R}_{h+1})^d+3^dN\bigr).
\end{align*}
The contribution of the range $h<j\le k$ is therefore at most
$2^d\bigl(2(104\check{R}_{h+1})^d+3^dN\bigr)\mathsf{b}_k(N)$.

Adding the two ranges gives
\begin{equation*}
\bigl[C^{\mathrm{sh}}+2^d\bigl(2(104\check{R}_{h+1})^d+3^dN\bigr)\bigr]\mathsf{b}_k(N)
=\mathsf{S}^{\mathrm{sh}}_{\mathrm{wr}}(k),
\end{equation*}
by \eqref{4.23:eq-wrapped-cell-sum-a}, which proves \eqref{4.23:eq-wrapped-cell-sum-1}.

\emph{Conclusion \textup{(3)}.} Let $v\in\mathfrak{L}^{\mathrm{wr}}(C)$. We use clause (iii)
of Lemma~\ref{4.23:lem-filing-domains}, which is proved as an implication from the
statements named there. By that clause, a piece filed at $C$ has a parent unit $v$ with
$K_v\subset C$, $X_v\subset C$ and $\mathcal{P}_v^{\circ}=\emptyset$, and no second-part
unit is filed at $C$. So $v$ is a wrapped unit of the first part with
$\mathcal{P}_v^{\circ}=\emptyset$. By the cut (a), which applies since $K_v\subset C$,
either $v\in\mathfrak{L}^{\mathrm{dp}}$ with $X_v\subset\mathsf{C}$, or $v$ is as in (S).

By Lemma~\ref{4.23:lem-unit-oscillation}, whose proof is incomplete at one step, we have
$\operatorname{tot}(v)\le\mathfrak{n}(v)e^{4(d_k(K_v)+1)}$
(see \eqref{4.23:eq-unit-oscillation-a}), with $\mathfrak{n}(v)$ read from that lemma's
table as a minimum. Since $\mathfrak{n}(v)$ is a minimum over the entries of that table, it
is at most each entry that applies to $v$: for a deep wrapped unit
$\mathfrak{n}(v)\le C_{3F}e^{-\frac12\delta_P\hat{D}}\mathfrak{N}^{A,\mathbb{C}}(v)$, and
for a shallow wrapped unit $\mathfrak{n}(v)\le2\mathfrak{N}^{A,\mathbb{C}}(v)$. Here
$\hat{D}$ is the depth of $v$ entering that table, and
$\mathfrak{N}^{A,\mathbb{C}}(v)=\mathfrak{N}^{A}(v)$ by the identification of the complex
amplitude $\mathfrak{N}^{A,\mathbb{C}}$ with the stored amplitude $\mathfrak{N}^{A}$ made
with the format.
Consequently
\begin{align*}
\operatorname{tot}(v)&\le e^{4}\,C_{3F}\,\mathfrak{N}^{A}(v)e^{4d_k(K_v)}
e^{-\frac12\delta_P\hat{D}}&&(v\in\mathfrak{L}^{\mathrm{dp}}),\\
\operatorname{tot}(v)&\le2e^{4}\,\mathfrak{N}^{A}(v)e^{4d_k(K_v)}
&&(v\ \text{as in (S)}).
\end{align*}
For $v\in\mathfrak{L}^{\mathrm{dp}}$, hypothesis (h) replaces the factor
$e^{-\frac12\delta_P\hat{D}}$ by the larger factor
$e^{-\frac14\delta_P d^{\wedge}(X_v,\mathsf{C}^c)}$. Summing over the deep wrapped units,
all of which satisfy $X_v\subset\mathsf{C}$, and applying hypothesis (b) and conclusion (1),
we bound the sum over the deep wrapped units by
$e^{4}C_{3F}\mathsf{S}^{\mathrm{dp}}_{\mathrm{wr}}(k)$. Conclusion (2) bounds the sum over
the shallow ones by $2e^{4}\mathsf{S}^{\mathrm{sh}}_{\mathrm{wr}}(k)$. Adding the two gives
\eqref{4.23:eq-wrapped-cell-sum-2}.
\end{proof}

\begin{remark}[Gaussian families of nested arrests, already counted]\label{4.23:rem-nested-gaussians}
Let $C$ be a component, and let $\mathfrak{a}$ run over the live arrests nested in $C$. The
Gaussian family $\mathfrak{Q}(\mathfrak{a})$ of such an arrest
is contained in the class $\mathfrak{L}^{B}$ of plain stored boundary-type terms of
Definition~\ref{4.23:def-component-constituents}. Its members are the Gaussian pieces $v_e$,
$e\in\mathfrak{E}^{G}_{W}$, of the decomposition \eqref{4.23:eq-nested-gaussians-2} below, $W$
being the region of $\mathfrak{a}$; each piece carries the datum, an oscillation datum on deep
pieces and a norm datum on the others,
\begin{equation}\label{4.23:eq-nested-gaussians-1}
\mathfrak{N}(v_e):=\mathfrak{N}^{G,L}(e)\quad(e\ \text{deep}),\qquad
\mathfrak{N}(v_e):=2\nu^{\natural}(e)\quad(e\ \text{not deep}),
\end{equation}
and the family is counted in the cell sum for plain units (Lemma~\ref{4.23:lem-plain-cell-sum},
with the constants \eqref{4.23:eq-plain-cell-sum-a}). Nothing is added to the final bound on the
sum over the component. For the oscillations of its pieces, the only input beyond the filing
lemma for the terms of nested arrests is, for the deep pieces, the transfer read at pairs of
points with the same real variables $\varpi$, at the seam \eqref{4.23:eq-unit-oscillation-b} of
Lemma~\ref{4.23:lem-unit-oscillation}, whose proof is incomplete at one step, and, for the other
pieces, the transfer by the norm. The polylogarithmic factor of
the arrest never stands alone, and no factor depending on the depth of nesting occurs.
\end{remark}

\begin{proof}
Let $W$ be the region of the arrest $\mathfrak{a}$, and let $\mathsf{q}_W$ be its Gaussian member.
Let $\mathfrak{E}^{G}_{W}$ be the index set of the Gaussian pieces of $W$ given by the Gaussian
decomposition theorem.
We use the filing lemma for the terms of nested arrests, together with clauses (a)
and (f) of the Gaussian decomposition theorem. By these,
\begin{equation}\label{4.23:eq-nested-gaussians-2}
\mathsf{q}_W=\sum_{e\in\mathfrak{E}^{G}_{W}} v_e
\end{equation}
holds exactly. Each $v_e$ is in the format (f0)--(f4), with support
$A_e:=[S^{+}(e)]_{m_0(e)}$ and with the distinguished set $\mathcal{P}$ of the format empty,
$\mathcal{P}=\emptyset$. Here $S^{+}(e)$ is the region attached to the piece $e$, $m_0(e)$ is the
level at which the piece is read, and $[\,\cdot\,]_{m}$ is the bracket operation of the Gaussian
decomposition theorem, which produces a set of cubes of level $m$. The filing lemma is taken with
the datum \eqref{4.23:eq-nested-gaussians-1}, an oscillation datum $\mathfrak{N}^{G,L}(e)$ on the
deep pieces and a norm datum $2\nu^{\natural}(e)$ on the pieces that are not deep.

By Definition~\ref{4.23:def-component-constituents}, the class $\mathfrak{L}^{B}$ in the splitting
$\mathfrak{L}^{\mathrm{pl}}(C)=\mathfrak{L}^{B}\sqcup\mathfrak{L}^{\mathrm{pt}}$ contains the terms
of nested arrests as the filing lemma files them; applied to the pieces $v_e$ of
\eqref{4.23:eq-nested-gaussians-2}, this gives $\mathfrak{Q}(\mathfrak{a})\subset\mathfrak{L}^{B}$,
with the datum \eqref{4.23:eq-nested-gaussians-1}.

Clause (f) of the Gaussian decomposition theorem is a bound per cube. Its constant $b^{G}_{\delta}$
is a contribution to $b_\delta$. The constant $b_\delta$ enters $\mathsf{b}(s)=\mathsf{b}_k(s)+(s+1)b_\delta$,
which appears in the constants $\mathsf{S}^{\mathrm{dp}}_{\mathrm{pl}}(k)$ and
$\mathsf{S}^{\mathrm{sh}}_{\mathrm{pl}}(k)$ of \eqref{4.23:eq-plain-cell-sum-a}. Hence the families
$\mathfrak{Q}(\mathfrak{a})$ are already counted in Lemma~\ref{4.23:lem-plain-cell-sum}, and nothing
is added to the final bound on the component sum.

Beyond the filing lemma for the terms of nested arrests, the oscillation estimates need two
further inputs.
\begin{itemize}
\item For the deep pieces, they need the transfer of clauses (a) and (b), the depth-carrying
clauses, of the oscillation bound for deep pieces to $v_e$, with $\mathfrak{N}$ replaced by
$\mathfrak{N}^{G,L}(e)$. This transfer is read at
pairs of points of the polydisc $\mathfrak{P}$ with the same real variables $\varpi$, that is, at the
seam \eqref{4.23:eq-unit-oscillation-b} of Lemma~\ref{4.23:lem-unit-oscillation}, whose proof is
incomplete at one step.
\item For the pieces that are not deep, they need the transfer for such pieces, taken by the
norm, with datum $2\nu^{\natural}(e)$.
\end{itemize}

Let $j_W$ be the level of the arrest. Its polylogarithmic factor is
$P_{\mathfrak{b}}(x_{j_W})$, a polylogarithm in the inverse coupling $x_{j_W}$, and it never stands
alone. In the transfer for deep pieces and in the transfer for pieces that are not deep
respectively it enters only through the two bounds below. In them $\Theta^{L}(e)$ is the weight
attached to the deep piece $e$ in the transfer for deep pieces, $\exp(-|Y'|)$ is the exponential
weight in the size $|Y'|$ of the region $Y'$ entering the two transfer statements, and
$\Pi^{L}_{G}$, $\Pi^{\mathrm{sh}}$ are the functions of the inverse coupling furnished by the two
transfers:
\begin{align}
P_{\mathfrak{b}}(x_{j_W})\,\Theta^{L}(e)\,\exp(-|Y'|)&\le\Pi^{L}_{G}(x_k),
\label{4.23:eq-nested-gaussians-3}\\
P_{\mathfrak{b}}(x_{j_W})\,\exp(-|Y'|)&\le\Pi^{\mathrm{sh}}(x_k),\label{4.23:eq-nested-gaussians-4}
\end{align}
where $\Pi^{L}_{G}(x)$ decreases to $0$ as $x\to\infty$. Under the coupling condition of the two
transfers, which includes $\Pi^{\mathrm{sh}}(x_k)\le1$ and $\Pi^{L}_{G}(x_k)\le1$, both right-hand
sides of \eqref{4.23:eq-nested-gaussians-3} and \eqref{4.23:eq-nested-gaussians-4} are at most $1$.

Finally, clause (c) of the counting statement of the Gaussian decomposition theorem is a count
over the roots of the nesting alone, and it carries no factor depending on the depth of nesting.
Consequently, no such factor arises for the arrests nested in $C$.
\end{proof}

\begin{definition}[Reading of the subtracted copies]\label{4.23:def-copies-reading}
Fix the level $k$ and a component $C$ as in Definition~\ref{4.23:def-component-constituents}. Let
$Z$ be a large-field region, $\Lambda_k$ the small-field region at level $k$, $\check{R}_k$ the
window radius function, $\mathfrak{S}_{\mathrm{src}}$ the source tube over $C^{\sharp}$, and
$\zeta=(\mathbf{U},\mathbf{J})$ a point of this tube. Use the extension conventions of
Definition~\ref{4.23:def-extension-conventions}.

\begin{enumerate}
\item[(a)] \emph{The copies.} The \emph{subtracted copies} of \cite{B-CMP122b} are indexed by pairs
$(X,y)$ of a domain $X$ of a term of the expansion (here $X$ denotes a domain of a term, not a
large-field region) with $X\cap Z\neq\emptyset$ and a
point $y$, taken at the single background $U_k$. Farness of a point $y$ is understood in the sense
of the correlation theorem. In this item the admissibility of the prepared large-field data is
used; it is printed without proof in \cite[pp.~178--179]{B-CMP122a}; cited as printed.
\begin{itemize}
\item For a far point $y\in Z$ the copy is the whole group of terms attached to $(X,y)$.
\item For a far point $y\in Z^{c}$ the copy is only the tail of that group.
\end{itemize}
The values of the copies are the differences $F(X,\mathfrak{b},y)-F(X,1,y)$, where
$F(X,\mathfrak{b},y)$ denotes the term of the correlation theorem attached to $(X,y)$ as a function
of its parameter $\mathfrak{b}$, and $F(X,1,y)$ is its value at $\mathfrak{b}=1$; they are taken at
one admissible value $\mathfrak{b}$ of this parameter, in the sense of the correlation theorem's
remark on the copies. By the third clause of the correlation theorem's statement on the filing of
terms by label, they belong to another term of the expansion, the one in which $Z$ is healed, and
not to the term in which $Z$ is live. On the tube the background of a copy is $\zeta$ itself. Its
determining set is the new one: the zone index $n(y)$ equals $k$ on all of $C$, however old the core.

\item[(b)] \emph{The copies attached to $C$.} These are the copies $v$ whose localisation set
$\mathfrak{K}_v$ satisfies $\mathfrak{K}_v\subset C$. They are of two kinds:
\begin{itemize}
\item the copies of the pointed terms of the first three rows of the pointed list in
\eqref{4.23:eq-component-constituents-a};
\item the copies of the counterterm pieces.
\end{itemize}
Both kinds have levels $j\le k$, with $j=k$ included. No copy of boundary type is attached to $C$.
The copies carry no family of the chain, hence no interpolation variables $\vec{\mathsf{s}}$ and
no tuple sum.

\item[(c)] \emph{The reading.} Let $v$ be a copy core term of level $j$ at a far point
$y\in C$, and let $\mathfrak{i}^{(j)}_y$ be the level-$j$ density at $y$ of which $v$ is the copy.
Take the nested boxes $Q^{(p)}$ of the extension construction with outer level $n=k$, the value on
$C$ of the zone index of item (a), and width $\mathsf{W}_v=3$. When the set of levels
\begin{equation*}
\bigl\{p:\ j\le p\le k-2-s_{\mathrm{box}}(3),\ (Q^{(p)})^{+1}\subset Z\bigr\}
\end{equation*}
is non-empty, put
\begin{equation}\label{4.23:eq-copies-reading-1}
\begin{aligned}
p_v&:=\max\bigl\{p:\ j\le p\le k-2-s_{\mathrm{box}}(3),\ (Q^{(p)})^{+1}\subset Z\bigr\},
\qquad Q_v:=Q^{(p_v)},\\
\operatorname{ext}_v(\zeta)&:=(U_{p_v},J_{p_v})\bigl(Q_v,\,M^{p_v}(\mathbf{U}\restriction Q_v)\bigr),
\end{aligned}
\end{equation}
where $(U_p,J_p)(Q,\cdot)$ is the box minimiser on $Q$ of
Definition~\ref{4.23:def-extension-conventions}, which reads $\mathbf{U}$ on the outer boundary
$\partial^{+}Q$, and $(Q)^{+1}$ is $Q$ enlarged by one more layer of level-$p$ cubes.
Here $s_{\mathrm{box}}(3)$ is the name used in this definition for the value at width $3$ of the
function of the box construction of Definition~\ref{4.23:def-extension-conventions} that bounds the
admissible levels $p$ from above, as in \eqref{4.23:eq-copies-reading-1}; the subscript keeps it
apart from the depth $s$ of the glossary, and $M^{p_v}(\mathbf{U}\restriction Q_v)$
denotes the $p_v$-fold block average of $\mathbf{U}\restriction Q_v$, in the notation of
\cite{B-CMP122b}; the letter $M$ with a superscript and an argument always denotes this average,
while without an argument, as in $M^{4}$ below, it is the block size $M$ of the glossary. If the
set displayed above is non-empty, the core $v$ is called
\emph{smooth}, and the value $f_v(\zeta)$ of the copy $v$ at $\zeta$ is, by definition,
$\mathfrak{i}^{(j)}_y$ evaluated at $\operatorname{ext}_v(\zeta)$. Since the box minimiser reads
$\mathbf{U}$ on $\partial^{+}Q_v$, a smooth copy core reads $\zeta$ on $Q_v^{+1}$, and
$Q_v^{+1}\subset C$.

If this set is empty, the core $v$ is called \emph{rough}. Rough cores, and every other copy, are
read at $\zeta$ itself.
\end{enumerate}
\end{definition}

\begin{proof}
We verify the assertions made in (a) and (b), explain why the reading (c) is needed, and show that
it changes no real value and keeps the add-and-subtract exact.

\emph{Item (a).} The differences $F(X,\mathfrak{b},y)-F(X,1,y)$ are taken at one admissible value
$\mathfrak{b}$ of the parameter, as in the correlation theorem's remark on the copies.
By the correlation theorem, these values are bounded by
the same sizes as the terms they are copied from. For the whole groups at far $y\in Z$, the lemma of
the correlation theorem that bounds such groups applies as it stands.

\emph{Item (b): the level $j=k$.} The old action contains the level-$k$ localised term
$E^{(k)}(X)$ only for $X\subset\Lambda_k$. Moreover $\Lambda_k\cap Z=\emptyset$. Since every such
$X$ lies in $\Lambda_k$, no level-$k$ localised term of the old action meets $Z$. Thus at level $k$
the localised terms $E^{(k)}(X)$ of the old action furnish no copies, and the level-$k$ copies
attached to $C$ are those of the pointed and counterterm pieces only; in this sense $j=k$ is
included.

\emph{Item (b): no boundary-type copy.} In the expansion term in which $Z$ is healed, the anchor of a
boundary-type copy $v$ lies, by the anchor property of boundary-type terms, in a component of
$\Lambda_k^{c}$ that is not treated at this step. By the separation property of untreated
components, that component is at distance at least $\check{R}_k$ from $C$. Therefore
$\mathfrak{K}_v\not\subset C$, and $v$ is not attached to $C$.

\emph{Why smooth cores are not read at $\zeta$.} At a rough point of the tube a core has the marginal
bound $\nu^{2}\mathsf{t}^{4}$ and no better one; here $\nu$ and $\mathsf{t}$ are the index and the
variable of the statement on the marginal bound at rough points, the variable being written
$\mathsf{t}$ to keep it apart from the letter $t$ of the glossary. A level-$k$ cell would then
carry the factor $M^{4}E_0\sum_{\nu}\nu^{2}$, where $E_0$ is the constant of that name in the
order in which the constants of the correlation theorem are chosen (not to be confused with the
localised term $E^{(k)}(X)$ above), and the statement on the marginal bound at rough points leaves
this factor uncontrolled. That statement is in force here because, by the statement on rough
configurations of the source tube, such rough points lie in
$\mathfrak{S}_{\mathrm{src}}$. This is why the smooth copy cores are read at the extension
$\operatorname{ext}_v(\zeta)$ rather than at $\zeta$.

\emph{Real values are unchanged.} Consider the real point $\zeta=(U_k,\mathcal{J}(U_k))$, where
$\mathcal{J}(U_k)$ denotes the current that accompanies $U_k$ at this real point of the tube. Apply
the real-slice agreement clause of the statement on nested box extensions of the background (the
statement that also supplies the boxes $Q^{(p)}$ of item (c)), under the third of that statement's
hypotheses, with the data $(j,X,\square_0)$ replaced by $(p_v,[\square_0]_{p_v},Q_v)$.
Here $\square_0$ is the cube at which the density is read, and $[\square_0]_{p}$ is the level-$p$ cube
containing it. The property gives that $\operatorname{ext}_v(\zeta)$ coincides with $U_k$ modulo a
gauge transformation on $[\square_0]_{p_v}$.

Consequently, at real points the value $f_v$ assigned by \eqref{4.23:eq-copies-reading-1} is the
value at $\zeta$ itself. For rough cores and all other copies the reading is at $\zeta$ by
definition. So the reading changes no real value.

\emph{Exactness of the add-and-subtract.} The subtracted copy enters the expansion together with an
added copy of the same term. The two copies agree where the density is evaluated, namely at the real
minimiser. This was shown in the preceding paragraph.

Off the real slice, the extension of each term is a choice made at each use. This is the convention
of Definition~\ref{4.23:def-extension-conventions}, and it is also the convention of the correlation
theorem that a localisation is chosen per use. Hence reading the subtracted copy at
$\operatorname{ext}_v(\zeta)$ does not affect the exactness of the add-and-subtract.

\emph{The added copy.} The added copy is a term of the new action, in the form fixed by the
inductive hypothesis of the correlation theorem at level $k$. It is left untouched by the reading
above.

\emph{Agreement with Ba{\l}aban's construction.} The construction of \cite{B-CMP122b} proceeds in the
same way. In \cite[p.~375]{B-CMP122b} Ba{\l}aban writes that one is to ``replace $M^{\cdot}(U_k)$,
$V''$ by the corresponding $M^{\cdot}(\mathbf{U})$''. Thus $\mathbf{U}$ enters through the averages
$M^{\cdot}(\mathbf{U})$, exactly as in \eqref{4.23:eq-copies-reading-1}.
\end{proof}

\begin{lemma}[Regularity of tube points on coarse boxes]\label{4.23:lem-tube-regularity}
Let $\mathfrak{S}$ be the source tube over $C^{\sharp}$ of
Definition~\ref{4.23:not-oscillation-seminorm}. Let $S\subset C^{\sharp}$ be a box whose sides do
not exceed the side of one cube of the partition $\pi_{k-1}$, and put $\ell:=\ell_{k-1}$. Then for
every $\zeta=(\mathbf{U},\mathbf{J})\in\mathfrak{S}$ there is a $G^{\mathbb{C}}$-valued gauge
transformation on $S$ after which the restriction $\zeta\restriction S$ of $\zeta$ to $S$ reads
$(e^{i\eta\mathcal{A}^0},\mathbf{J}^0)$, where $\mathcal{A}^0$ is a bond function on $S$ with values
in the complexified Lie algebra of $G$, with
\begin{equation}\label{4.23:eq-tube-regularity-a}
\begin{gathered}
\ell\,|\mathcal{A}^0|,\quad \ell^2\,|\nabla\mathcal{A}^0|,\quad \ell^3\,|\mathbf{J}^0|
\;\le\; \mathfrak{a}'_{\mathrm{tb}} := 3\,c_{\mathrm{reg}}\,\alpha_{*,k},\\
|\partial\mathbf{U}-1| \;\le\; \mathfrak{a}'_{\mathrm{tb}}\,(\eta/\ell)^2 ,
\end{gathered}
\end{equation}
where
\begin{equation*}
c_{\mathrm{reg}} = 2L^2\,c_{\Psi}\,c_{\alpha}\,(1+2BM)
\end{equation*}
is the constant of the regularity lemma at free current. Under the smallness condition on
$c_{\mathrm{reg}}\alpha_{*,k}$ of that lemma one has
$\mathfrak{a}'_{\mathrm{tb}}\le 1$.
\end{lemma}

Here $\eta$ is the lattice spacing of the configurations in the tube, $\ell_j$ is the length scale of
level $j$, $\nabla$ is the flat lattice difference, $\partial\mathbf{U}$ stands for the plaquette
products $\mathbf{U}(\partial p)$ of plaquettes $p$ in $S$, and $B$, $M$ are the constants of
\cite[(2.38)]{B-CMP119}. Further, $c_{\alpha}$ is the constant of the amplitude bound
$\tilde\alpha_{\cdot,k}\le c_{\alpha}\alpha_{*,k}$ used in the proof below, $\alpha_{*,k}$ is the
radius of level $k$ with which that bound compares the radii $\tilde\alpha_{0,k}$ and
$\tilde\alpha_{1,k}$, and $c_{\Psi}$ is the constant so denoted in the regularity lemma at free
current.

\begin{proof}
The cubes of $\pi_{k-1}$ are no larger than those of $\pi_k$, since $\ell_{k-1}=\ell_k/L\le\ell_k$;
hence the box $S$ lies in a cube $\square$ made of $2^d$ cubes of $\pi_k$. The whole of $C^{\sharp}$
is of level $k$ in Ba{\l}aban's domain $\mathbf{B}_k$. This follows from \cite[(1.33)]{B-CMP122b}
together with the statement of this chapter that places the collared component $C^{\sharp}$ at
level $k$.

The point $\zeta$ lies in Ba{\l}aban's space of complex configurations
$\tilde{U}^{c}_{k}(C^{\sharp},\tilde\alpha_0,\tilde\alpha_1)$ of \cite[(2.34)--(2.39)]{B-CMP119}.
By the first clause of the description of that space, the configuration factors on
$\square\cap C^{\sharp}$ as $\mathbf{U}=U'U$, with $U$ a $G$-valued configuration and $U'$ the
complex factor.

We apply the second clause of the description of the space of complex configurations in
\cite[p.~261]{B-CMP119} with $n=j=k$ and with the cube size constant there equal to $2$. This is
admissible because $\square$ is a cube of the size allowed there. It gives a gauge transformation
$u$, defined on $\square\cap C^{\sharp}$, such that the real configuration $U$ satisfies
\begin{equation*}
U^{u}=\exp(i\eta A),\qquad
\ell_k\,|A|,\ \ \ell_k^2\,|\nabla_{\eta}A| \;<\; 2BM\,\alpha_{0,k}
\end{equation*}
on $\square\cap C^{\sharp}$. This is \cite[(2.38)]{B-CMP119}, in which the lattice spacing is
written $\xi$; here it is written $\eta$, and $L^k$ times the lattice spacing is written $\ell_k$.

The third clause of the same description writes the complex factor as $U'=\exp(i\eta A')$. It
gives the bounds \cite[(2.39)]{B-CMP119}, namely
\begin{equation*}
\ell_k|A'|<\tilde\alpha_{1,k},
\end{equation*}
together with the remaining bounds listed there. Since $S\subset\square\cap C^{\sharp}$, the
restriction of $u$ to $S$ is the gauge transformation of the statement.

Write $\operatorname{Ad}_u$ for conjugation by $u$ at the initial point of a bond:
$(\operatorname{Ad}_uA')(b)=u(x)A'(b)u(x)^{-1}$ for a bond $b$ from $x$ to $y$. For such a bond,
\begin{equation*}
\mathbf{U}^{u}(b)=u(x)\,U'(b)\,U(b)\,u(y)^{-1}
=\bigl(u(x)\,U'(b)\,u(x)^{-1}\bigr)\,U^{u}(b),
\end{equation*}
and
\begin{equation*}
u(x)\exp\bigl(i\eta A'(b)\bigr)u(x)^{-1}=\exp\bigl(i\eta\,(\operatorname{Ad}_uA')(b)\bigr).
\end{equation*}
Hence in the gauge $u$ the configuration becomes
\begin{equation}\label{4.23:eq-tube-regularity-1}
\mathbf{U}^{u}=e^{i\eta\operatorname{Ad}_uA'}\,e^{i\eta A}.
\end{equation}
The covariant difference of the complex part transforms covariantly:
\begin{equation}\label{4.23:eq-tube-regularity-2}
\nabla_{e^{i\eta A}}\bigl(\operatorname{Ad}_uA'\bigr)=\operatorname{Ad}_u\,\nabla_UA'.
\end{equation}
Hence the bounds of \cite[(2.39)]{B-CMP119} on the covariant differences of $A'$ along $U$ become
bounds on the covariant differences of $\operatorname{Ad}_uA'$ along $e^{i\eta A}$.

The covariant-to-flat comparison of this chapter then replaces these covariant differences by the
flat difference $\nabla$, at the cost of a factor $9/8$. The Baker--Campbell--Hausdorff bound of
this chapter, applied at the base point $V=1$, merges the two exponentials in
\eqref{4.23:eq-tube-regularity-1} into a single exponential $e^{i\eta\mathcal{A}^0}$.

The amplitude bound of this chapter,
\begin{equation*}
\tilde\alpha_{\cdot,k}\le c_{\alpha}\,\alpha_{*,k},
\end{equation*}
expresses the radii $\tilde\alpha_{0,k}$ and $\tilde\alpha_{1,k}$ through $\alpha_{*,k}$. The
bounds on $\ell_k|\mathcal{A}^0|$ and $\ell_k^2|\nabla\mathcal{A}^0|$ at level $k$ are assembled
from three contributions: the bound on $A$ given by \cite[(2.38)]{B-CMP119}; the bounds on
$\operatorname{Ad}_uA'$ given by \cite[(2.39)]{B-CMP119}, with the factor $9/8$ where a covariant
difference was replaced by the flat one; and the correction by which the Baker--Campbell--Hausdorff
bound separates $\mathcal{A}^0$ from $A+\operatorname{Ad}_uA'$. These are added with the arithmetic
of the proof of the first clause of the regularity lemma at free current, which fixes
$c_{\mathrm{reg}}$, and this gives the constant $\mathfrak{a}'_{\mathrm{tb}}$ of
\eqref{4.23:eq-tube-regularity-a}.

It remains to pass from level $k$ to the length scale $\ell=\ell_{k-1}=\ell_k/L\le\ell_k$ of the
statement. Reading the three quantities $\ell|\mathcal{A}^0|$, $\ell^2|\nabla\mathcal{A}^0|$ and
$\ell^3|\mathbf{J}^0|$ at the scale of level $k-1$ only decreases them. Indeed, $\mathcal{A}^0$,
$\nabla\mathcal{A}^0$ and $\mathbf{J}^0$ carry the homogeneities $1$, $2$ and $3$ respectively
between the levels $k-1$ and $k$. The passage is therefore an identity of units, and the regularity
established at level $k$ is kept.

The bound in the second line of \eqref{4.23:eq-tube-regularity-a}, on
$|\partial\mathbf{U}-1|$, holds as it stands among the clauses defining the tube. So does
the bound on $\ell^3|\mathbf{J}^0|$, which is one of the tube's clauses on the current.
\end{proof}

\begin{remark}
The argument above is the proof of the first clause of the regularity lemma at free current, with
the curvature hypothesis of that lemma replaced by the clauses defining the tube, and with the
weight function of that lemma identically equal to $1$. The remark on the curvature hypothesis
made after that lemma is not used here.
\end{remark}

\begin{lemma}[The copies lemma]\label{4.23:lem-copies-lemma}
Let $C$ be a component at level $k$, and let the subtracted copies be read as in
Definition~\ref{4.23:def-copies-reading}. Write $\mathfrak{L}^{\mathrm{cp}}(C)$ for the family of
subtracted copies filed at $C$. For $v\in\mathfrak{L}^{\mathrm{cp}}(C)$ write:
\begin{itemize}
\item $f_v$ for the copy, a function of the point $\zeta=(\mathbf{U},\mathbf{J})$ of the tube
$\mathfrak{S}$;
\item $y_v$ for its point;
\item $X_v$ and $\hat X_v$ for its domains;
\item $Q_v$ and $Q_v^{+1}$ for its box and the box enlarged by one layer, when $v$ is a smooth core.
\end{itemize}
Let $\mathfrak{N}^{\flat}$ be the format table of the pointed constituents with its level parameter
set to $k$ and its second parameter set to $3$, its smooth-core row read at box side
$2(k-p_v)+13$, where $p_v$ is the extraction level of Definition~\ref{4.23:def-copies-reading}.
Assume the following:
\begin{itemize}
\item the analyticity proposition for the pointed constituents, in the form that its conclusion
holds for every background obeying the bounds \eqref{4.23:eq-tube-regularity-a} in place of
clause (i) of its regularity statement.
\item the flat cell-count bound for the pointed constituents, with its constant
$\mathsf{K}^{\flat}$, in the form given by its proof.
\item the inequality $s(3)\ge 1$ for the depth offset $s(\cdot)$ of the pointed-term
construction, as that construction defines it, and $L\ge L_0$.
\item the short-term lemma.
\item the smooth-core lemma, as implied by the admissibility of the level-$p_v$ averages on boxes.
\item the regularity condition on cores of the pointed-term construction, and the
background-regularity lemma at level $p_v$.
\item clause (c) of the inductive hypothesis at level $k$.
\item the cube count of a component: the number $\#_k(C)$ of cubes of $\pi_k$ contained in $C$
satisfies $\#_k(C)\le(100\check{R}_k)^d$.
\end{itemize}
Then the following hold.
\begin{enumerate}
\item[(i)] Every copy $f_v$, $v\in\mathfrak{L}^{\mathrm{cp}}(C)$, is holomorphic on $\mathfrak{S}$.
It depends on $\zeta$ only through its restriction to $\hat X_v\cup X_v$, or to $Q_v^{+1}$ when $v$
is a smooth core; this set is contained in $C$. Moreover $|f_v(\zeta)-f_v(1)|\le\mathfrak{N}^{\flat}(v)$
for every $\zeta\in\mathfrak{S}$, where $f_v(1)$ denotes the value of the copy at the argument $1$,
in the notation of the analyticity proposition for the pointed constituents; it does not depend
on $\zeta$.
\item[(ii)] For every cube $\Delta'$ of $\pi_k$ with $\Delta'\subset C$,
\begin{equation}\label{4.23:eq-copies-lemma-a}
\begin{gathered}
\sum_{v\in\mathfrak{L}^{\mathrm{cp}}(C):\ y_v\in\Delta'}\mathfrak{N}^{\flat}(v)
\le\mathsf{K}^{\mathrm{cp}},\\
\mathsf{K}^{\mathrm{cp}}:=\mathsf{K}^{\flat}
+2M^d\bigl[K(d)\bigl((1+C_B)E_0+1\bigr)+6L^4(b_++r_J)\bigr].
\end{gathered}
\end{equation}
Here $K(d)$ is the constant of \cite[(1.26)]{B-CMP116}; $C_B$, $E_0$, $b_+$ are constants of
the pointed-term construction; $r_J$ is the source radius of the correlation theorem; and
$L^4=L^d$ is the number of cubes of $\pi_{k-1}$ in one cube of $\pi_k$, the sides of the cubes of
$\pi_{k-1}$ being $L^{-1}$ times those of $\pi_k$. The number
$\mathsf{K}^{\mathrm{cp}}$ does not depend on the couplings and is fixed once $M$ is fixed.
\item[(iii)] The oscillation of the sum of the copies satisfies
\begin{equation}\label{4.23:eq-copies-lemma-1}
\operatorname{osc}\Bigl(\sum_{v\in\mathfrak{L}^{\mathrm{cp}}(C)}f_v\Bigr)
\le 2(100\check{R}_k)^d\,\mathsf{K}^{\mathrm{cp}}.
\end{equation}
\end{enumerate}
\end{lemma}

\begin{proof}
\emph{Clause (i).} We follow the proof of the analyticity proposition for the pointed constituents,
class by class. That proof uses the background only through the three bounds of the regularity
statement of the pointed-term construction and through the domain of the constituent itself. We
replace clause (i) of that regularity statement by Lemma~\ref{4.23:lem-tube-regularity}. Its bounds
\eqref{4.23:eq-tube-regularity-a} hold at every $\zeta\in\mathfrak{S}$ on every box
$S\subset C^{\sharp}$ with sides at most one cube of $\pi_{k-1}$.

\emph{Membership.} The copies are terms of the new action, evaluated on the tube of the new action.
For each constituent with domain $X$, the restriction of $\zeta$ to $X$ lies in the domain of that
constituent. This follows from clause (c) of the inductive hypothesis at level $k$, which rests on
\cite[(2.27)(ii)]{B-CMP119}: the spaces descend.

\emph{Short single terms.} These are bounded by the short-term lemma, applied with $S:=\hat X_v$.

\emph{Long terms and tail terms.} Such a term is bounded in absolute value, on the domain of its
constituent, by the amplitude attached to $v$ in the pointed-term construction.

\emph{Counterterm pieces.} These are bounded by the first entry of \eqref{4.23:eq-tube-regularity-a}.

\emph{Smooth cores.} We apply Lemma~\ref{4.23:lem-tube-regularity} with $S:=Q_v^{+1}$, which lies in
$C\subset C^{\sharp}$. We first check that $Q_v^{+1}$ is still a box with sides at most one cube of
$\pi_{k-1}$. Put $p:=p_v$ and $a:=k-p$. The box $Q^{(p)}$ has side $2(k-p)+11$ cubes of $\pi_p$.
The enlargement adds one layer of cubes of $\pi_p$ on each side, so $Q_v^{+1}$ has side
$2(k-p)+13=2a+13$ cubes of $\pi_p$. One cube of $\pi_{k-1}$ has side $L^{k-p-1}=L^{a-1}$ cubes of
$\pi_p$. Hence it suffices to show
\begin{equation*}
2a+13\le L^{a-1}\qquad\text{for } a\ge 3.
\end{equation*}
At $a=3$ this inequality reads $2\cdot 3+13\le L^2$; it is one of the lower bounds on $L$ contained
in the condition $L\ge L_0$. Passing from $a$ to $a+1$, the right member is multiplied by $L$, while
the left member grows by $2$. Indeed, if $2a+13\le L^{a-1}$, then
\begin{gather*}
L^2\ge 2\cdot 3+13>1,\qquad\text{so}\qquad L-1\ge 1,\\
(L-1)L^{a-1}\ge L^{a-1}\ge 2a+13\ge 2,\\
2(a+1)+13\le L^{a-1}+2\le L^{a-1}+(L-1)L^{a-1}=L^{a}.
\end{gather*}
The condition $a\ge 3$ holds because $p\le k-2-s(3)$ by the definition of $p_v$ in
Definition~\ref{4.23:def-copies-reading}, so that
\begin{equation*}
k-p\ge 2+s(3)\ge 3.
\end{equation*}
This uses $s(3)\ge1$, which is the hypothesis on the depth offset $s(\cdot)$ and is not derived
here.

Next, the regularity condition on cores of the pointed-term construction and the definition of
$s(3)$ give the following: the restriction of $\mathbf{U}$ to $Q_v^{+1}$ stays, modulo gauge, in the
ball of the background-regularity lemma at level $p$, with the radius parameter $\epsilon$ of that
ball at most one. The smooth-core lemma then applies. It holds as a consequence of the admissibility
of the level-$p$ averages on boxes; the constant through which that admissibility enters the
smooth-core lemma is a number fixed before $\gamma$. It yields the bound $\mathfrak{N}^{\flat}(v)$
for the core, carrying its one factor $L^{-\hat\alpha(p-j)}$, with $\hat\alpha$ the decay exponent
of the smooth-core lemma. Here $j$ is the level at which the term was born and $p$ the level of
evaluation; the factor is read once, at $p$.

The core lies entirely in the small-field region $\Lambda_j$ of the label of the healed region.
There the region $Z$ of Definition~\ref{4.23:def-copies-reading}, in which the boxes $Q^{(p)}$ are
taken, is small-field at every level, and the other components are at distance at least
$\check{R}_k$.

\emph{Rough cores.} By Definition~\ref{4.23:def-copies-reading} these are read at $\zeta$ itself.
They are bounded termwise.

In each class, therefore, $f_v$ is holomorphic on $\mathfrak{S}$, depends on $\zeta$ only on the sets
named in (i), and satisfies $|f_v(\zeta)-f_v(1)|\le\mathfrak{N}^{\flat}(v)$.

\emph{Clause (ii).} The proof of the flat cell-count bound for the pointed constituents contains no
field. It sees the family only through the following inputs:
\begin{itemize}
\item the format table of the pointed constituents;
\item the number of points of that proof's auxiliary set $T^{(j)}$ in each cell, written there
$M^4t^{-4}$, with $t$ its cell-side parameter;
\item the bound on long terms;
\item the fact that a tail term has $d_j(X)$ at least that proof's tail threshold $\nu$;
\item clause (iii) of the bound on boundary pieces;
\item the count of the rim $\mathfrak{L}_C$.
\end{itemize}
Take the cell to be the cube $\Delta'$ of level $k$. Then the cell weight $\mathsf{w}^{\vee}$ of that
bound equals $1$. Drop the weights $e^{rd_k(K_v)}e^{\varsigma_1\mathsf{r}(v)}$ that the bound
carries, in its notation. The sums of that proof over the levels $j<k$ then stand verbatim, and
contribute at most $\mathsf{K}^{\flat}$.

The level $j=k$ adds, for each point $y$, a sum over domains $X\ni y$; here $\delta$ and
$\kappa_0$ are constants of the flat cell-count bound. Using $g_k^{\kappa_0}\le1$, this sum
satisfies
\begin{equation*}
\sum_{X\ni y}\bigl[2E_0+2C_BE_0+2g_k^{\kappa_0}\bigr]e^{-(1-\frac14\delta)\kappa d_k(X)}
\le 2K(d)\bigl((1+C_B)E_0+1\bigr).
\end{equation*}
The last inequality rests on the count
\begin{equation*}
\sum_{X\ni y}e^{-(1-\frac14\delta)\kappa d_k(X)}\le K(d),
\end{equation*}
which is \cite[(1.26)]{B-CMP116} in the form in which the flat cell-count bound uses it.

The level $j=k$ also adds, for each point $y$, one counterterm piece. Since
$\mathfrak{a}'_{\mathrm{tb}}\le1$ under the condition on the regularity constants stated in
Lemma~\ref{4.23:lem-tube-regularity}, this piece is bounded by
\begin{equation*}
6L^4(b_++r_J)\,\mathfrak{a}'^{\,2}_{\mathrm{tb}}\le 6L^4(b_++r_J).
\end{equation*}
A cube of $\pi_k$ contains $M^d$ points. The contribution of level $k$ to the left member of
\eqref{4.23:eq-copies-lemma-a} is therefore at most
\begin{equation*}
M^d\bigl[2K(d)\bigl((1+C_B)E_0+1\bigr)+6L^4(b_++r_J)\bigr],
\end{equation*}
and $\mathsf{K}^{\mathrm{cp}}-\mathsf{K}^{\flat}$, the bracketed term of
\eqref{4.23:eq-copies-lemma-a}, is an upper bound for it. Adding the contribution
$\mathsf{K}^{\flat}$ of the levels $j<k$ gives \eqref{4.23:eq-copies-lemma-a}.

\emph{Clause (iii).} Put
\begin{equation*}
F:=\sum_{v\in\mathfrak{L}^{\mathrm{cp}}(C)}f_v,\qquad
c:=\sum_{v\in\mathfrak{L}^{\mathrm{cp}}(C)}f_v(1).
\end{equation*}
The oscillation $\operatorname{osc}(F)$ is the supremum of $|F(\zeta)-F(\zeta')|$
over $\zeta,\zeta'\in\mathfrak{S}$, at one fixed value of the real variables $\Phi_C$. The value
$f_v(1)$ does not depend on $\zeta$, so neither does $c$.
Then $|F(\zeta)-F(\zeta')|\le|F(\zeta)-c|+|F(\zeta')-c|$. By the triangle
inequality and clause (i), this gives the first two inequalities below. Grouping the copies according
to the cube $\Delta'\subset C$ of $\pi_k$ that contains $y_v$, and applying
\eqref{4.23:eq-copies-lemma-a} to each cube, gives the third:
\begin{align*}
\operatorname{osc}\Bigl(\sum_{v\in\mathfrak{L}^{\mathrm{cp}}(C)}f_v\Bigr)
&\le 2\sup_{\zeta\in\mathfrak{S}}
\Bigl|\sum_{v\in\mathfrak{L}^{\mathrm{cp}}(C)}\bigl(f_v(\zeta)-f_v(1)\bigr)\Bigr|\\
&\le 2\sum_{\Delta'\subset C}\ \sum_{v:\ y_v\in\Delta'}\mathfrak{N}^{\flat}(v)\\
&\le 2\,\#_k(C)\,\mathsf{K}^{\mathrm{cp}}.
\end{align*}
The cube count of a component gives $\#_k(C)\le(100\check{R}_k)^d$, and
\eqref{4.23:eq-copies-lemma-1} follows.
\end{proof}

However old the core, its copies are counted in cells of level $k$. The lemma bounds the oscillation
of the sum of the copies by this count. It does not assert that the copies vary little on the tube,
and no such conclusion is drawn from it.

\begin{lemma}[Oscillation of the localised boundary terms]\label{4.23:lem-boundary-oscillation}
Let $C$ be a treated component at level $k$, and let $C^{\sharp}$ be $C$ with its collar layers.
The $z$-free part of Ba{\l}aban's sum $\mathbf{C}_k$ of boundary terms \cite[\S1]{B-CMP122b}
localised in $C^{\sharp}$ is the constituent
$\mathfrak{v}^{\mathbf{C}}$ of the resummed density of Definition~\ref{4.23:def-resummed-density}.
Its members $\mathfrak{v}^{\mathbf{C}}(Y_4)$, $Y_4\in\mathbf{D}_k^Z$, $Y_4\subset C^{\sharp}$, with
$\mathbf{D}_k^Z$ the localisation domains of Definition~\ref{4.23:def-resummed-density}, are
those of Definition~\ref{4.23:def-boundary-terms}, printed without proof in
\cite[(1.11), (1.23)--(1.48)]{B-CMP122b}; cited as printed. Here
$\mathsf{c}_{\mathbf{C}}=1+c_IM^5$ is the constant fixed there, $c_I$ being the constant written
$O(1)$ in \cite[(1.11)]{B-CMP122b}. With the oscillation seminorm
$\operatorname{osc}$ of Notation~\ref{4.23:not-oscillation-seminorm}, one has
\begin{equation}\label{4.23:eq-boundary-oscillation-a}
\operatorname{osc}\Bigl(\sum_{Y_4\subset C^{\sharp}}\mathfrak{v}^{\mathbf{C}}(Y_4)\Bigr)
\le 2\sum_{Y_4\subset C^{\sharp}}\sup\bigl|\mathfrak{v}^{\mathbf{C}}(Y_4)\bigr|
\le 2\,\mathsf{c}_{\mathbf{C}}=2\bigl(1+c_IM^5\bigr).
\end{equation}
\end{lemma}

\begin{proof}
We first identify the summand. In Ba{\l}aban's construction the sum $\mathbf{C}_k$ of the boundary
terms, including the term $I(U_0)$ of \cite[(1.11)]{B-CMP122b} coming from the denominator, is one
of the terms of the effective
density that the step carries forward \cite[\S1]{B-CMP122b}. It is therefore part of the resummed
density $V^{\sharp}(C^{\sharp})$ of Definition~\ref{4.23:def-resummed-density}. The part of
$\mathbf{C}_k$ that depends on the complex source $w$ is accounted for in the correlation theorem,
where it is adjoined to one of the terms of that theorem. The $z$-free part of
$\mathbf{C}_k$ occurs in none of the bracket terms and none of the value terms of the correlation
theorem. Hence it is the constituent $\mathfrak{v}^{\mathbf{C}}$ of the resummed density of
Definition~\ref{4.23:def-resummed-density}. Its localised members
$\mathfrak{v}^{\mathbf{C}}(Y_4)$ are those of Definition~\ref{4.23:def-boundary-terms}, which is
printed without proof in \cite[(1.11), (1.23)--(1.48)]{B-CMP122b}; cited as printed. By that
definition each
$\mathfrak{v}^{\mathbf{C}}(Y_4)$ is analytic on the tube over $Y_4$, and the family is invariant
under simultaneous $G$-valued gauge transformations of $(\zeta;\varpi)$, $G$ the gauge group.
The bounds displayed for these terms in \cite[(1.23), (1.27), (1.31), (1.47)]{B-CMP122b} are
weighted absolute sums, and Definition~\ref{4.23:def-boundary-terms} gives the bound on the
absolute sum over the localisation domains:
\begin{equation}\label{4.23:eq-boundary-oscillation-1}
\sum_{Y_4\subset C^{\sharp}}\sup\bigl|\mathfrak{v}^{\mathbf{C}}(Y_4)\bigr|\le\mathsf{c}_{\mathbf{C}}
=1+c_IM^5 .
\end{equation}

We now prove \eqref{4.23:eq-boundary-oscillation-a}. Write
$f=\sum_{Y_4\subset C^{\sharp}}\mathfrak{v}^{\mathbf{C}}(Y_4)$. Fix real $\Phi_C$ and fix $w$, and
let $\zeta,\zeta'\in\mathfrak{S}$. By the triangle inequality,
\begin{equation*}
\bigl|f(\zeta;\Phi_C,w)-f(\zeta';\Phi_C,w)\bigr|
\le\bigl|f(\zeta;\Phi_C,w)\bigr|+\bigl|f(\zeta';\Phi_C,w)\bigr| .
\end{equation*}
The members $\mathfrak{v}^{\mathbf{C}}(Y_4)$ do not depend on $w$. For every
$\zeta''\in\mathfrak{S}$, a second application of the triangle inequality gives
\begin{equation*}
\bigl|f(\zeta'';\Phi_C,w)\bigr|
\le\sum_{Y_4\subset C^{\sharp}}\bigl|\mathfrak{v}^{\mathbf{C}}(Y_4)(\zeta'';\Phi_C)\bigr|
\le\sum_{Y_4\subset C^{\sharp}}\sup\bigl|\mathfrak{v}^{\mathbf{C}}(Y_4)\bigr| ,
\end{equation*}
where each supremum is taken over the tube over $Y_4$ and over the real variables; the restriction
of $\zeta''$ to $Y_4$ lies in the tube $\tilde{U}^{c}_{k}(Y_4,\tilde\alpha_0,\tilde\alpha_1)$ over
$Y_4$, since $\mathfrak{S}=\tilde{U}^{c}_{k}(C^{\sharp},\tilde\alpha_0,\tilde\alpha_1)$,
$Y_4\subset C^{\sharp}$, and the conditions of \cite[(2.34)--(2.39)]{B-CMP119} defining these
spaces are local.
We apply this bound with $\zeta''=\zeta$ and with $\zeta''=\zeta'$. Taking the supremum over
$\zeta,\zeta'\in\mathfrak{S}$, over real $\Phi_C$ and over $w$, as in the definition of
$\operatorname{osc}$, gives the first inequality of \eqref{4.23:eq-boundary-oscillation-a}. The
second inequality is \eqref{4.23:eq-boundary-oscillation-1} multiplied by two.
\end{proof}

We first identify which members of the boundary terms $\mathbf{C}_k$ depend on the point $\zeta$ of the source tube $\mathfrak{S}$ on the core $\mathsf{C}=\mathsf{C}_C$ of a component $C$. These boundary terms are those of Definition~\ref{4.23:def-boundary-terms}, printed without proof in \cite[\S1]{B-CMP122b} and cited as printed.

The identity expressing the boundary terms, which enters the correlation theorem, is exact and linear in the weight of the Wilson action. We take for the weight
\begin{equation*}
x''(p):=x_{\pi(x(p))},
\end{equation*}
which is $1/(g''_k(\cdot))^2$ in the notation of \cite[\S1]{B-CMP122b}. With this weight the third bracket of the identity is
\begin{equation*}
A(x''\hat\chi_1,U''_k)-A(x''\hat\chi_1,U''_{h+2}),
\end{equation*}
where $A(\cdot,\cdot)$ is the weighted Wilson action of Definition~\ref{1.1:def-wilson-action}. This bracket is the telescoped sum, over Ba{\l}aban's intermediate localisations, of the two families of \cite[\S1]{B-CMP122b}.

Points of the presented polydisc are written $m=(\zeta,\varpi,\vec{\mathsf{s}})\in\mathfrak{P}$, and $\mathcal{U}(m)$ denotes the configuration of the chain at the point $m$.

The configuration $U''_{h+2}=U''_{h+2}(e^{ig_kB},V'')$ does not depend on $\zeta$ \cite[\S1]{B-CMP122b}. The configuration $U''_k$ is the top of the chain: it is $\mathcal{U}(m)$ evaluated at the chain's top stage $\mathfrak{b}_4$.

Every other member of $\mathbf{C}_k$ is of one of two kinds.
\begin{itemize}
\item It may depend on the chain only on the support of $1-\hat\chi_1$, which does not meet $\mathsf{C}$. Indeed, Ba{\l}aban's cut-off function $\hat\chi_1$ and region $Z''_{h+1}$ of \cite[\S1]{B-CMP122b} satisfy $\hat\chi_1=1$ on $Z''_{h+1}$ except on a layer of width $2$ \cite[\S1]{B-CMP122b}, and $\mathsf{C}$ lies two layers of width $LM\check{R}_{h+1}$ inside $Z''_{h+1}$.
\item It may depend on $\zeta$ itself, respectively on the variable $\mathsf{V}$, through sets of level $k$. These are a part of the new Wilson action on $Z$, and the members carrying $U_0$ and $I(U_0)$ \cite[\S1]{B-CMP122b}. A member of this kind may cover the region of the core. It is nevertheless counted in cells of level $k$, and by its position its bound carries no depth $D$; in this it is like the subtracted copies.
\end{itemize}

In what follows $N_c$ is the number of colours, the gauge group being $SU(N_c)$ for general $N_c$. We write $N_c$ rather than $N$ because here $N=N(g_k)$ denotes the window and $h=k-N$ the bottom level. We write $\operatorname{Tr}$ for the ordinary trace of $N_c\times N_c$ matrices and $|\cdot|$ for the operator norm. For an invertible complex $N_c\times N_c$ matrix $W$ put
\begin{equation*}
\mathsf{a}(W):=1-\frac{1}{2N_c}\operatorname{Tr}\bigl(W+W^{-1}\bigr).
\end{equation*}
For $W\in SU(N_c)$ this is $1-\operatorname{Re}\operatorname{tr}W$.

Recall that the radius function is
\begin{equation*}
\alpha_*(x)=C_*\,x^{-1/2}(\log x)^{q},
\end{equation*}
and that $\alpha_{*,n}=\alpha_*(x_n)$, so that $x_n\alpha_{*,n}^2=C_*^2(\log x_n)^{2q}$. Throughout, $\ell_n$ is the spacing of level $n$ and $\eta$ that of the fine lattice on which the plaquettes $p$ live; $\eta\le\ell_n$.

\begin{lemma}[Plaquette units of the boundary terms]\label{4.23:lem-plaquette-units}
Let $c$ be a weight on plaquettes with $|c(p)|\le x''(p)$. For a cell $\Delta\subset C$ let $n:=\operatorname{lev}\Delta$ be its level, let $\mathsf{w}(\Delta)$ denote the width of $\Delta$ (the quantity controlled by the window bound on cell widths), and put
\begin{equation}\label{4.23:eq-plaquette-units-1}
f^W_\Delta(m):=\sum_{p:\,x(p)\in\Delta}c(p)\,\mathsf{a}\bigl(\mathcal{U}(m)(\partial p)\bigr),
\qquad
\hat\theta(\Delta):=3c_\alpha\,\mathsf{w}(\Delta)\,\alpha_{*,n},
\qquad
p_\Delta:=\sup_{y\in\Delta}p(y),
\end{equation}
where $c_\alpha$ is the constant of the curvature bound in (a) below, and $p(\cdot)$ is the weight function of the displacement bound below (not to be confused with a plaquette $p$). Assume the following.
\begin{enumerate}
\item[(a)] For every cell $\Delta\subset C$ the curvature bound for plain units on the presented polydisc holds: for every plaquette $p$ with $x(p)\in\Delta$ and every $m\in\mathfrak{P}$,
\begin{equation*}
|\mathcal{U}(m)(\partial p)-1|<\hat\theta(\Delta)(\eta/\ell_n)^2 .
\end{equation*}
\item[(a$'$)] For every cell $\Delta\subset C$ and every plaquette $p$ with $x(p)\in\Delta$, the map $m\mapsto\mathcal{U}(m)(\partial p)$ is holomorphic on $\mathfrak{P}$.
\item[(b)] The displacement bound holds for every cell $\Delta\subset\mathsf{C}$, with a weight function $p(\cdot)\le1$ and a constant $\mathsf{U}\le c_{\mathsf{U}}C_{\mathsf{a}}\alpha_{*,k}$, both independent of $\Delta$, where $c_{\mathsf{U}}$ and $C_{\mathsf{a}}$ are constants of the displacement lemma: if $m,m'\in\mathfrak{P}$ have the same $\varpi$ and $y\in\Delta$, then the displacement field $\mathcal{K}$ of the exponential representation of $\mathcal{U}(m)$ over $\mathcal{U}(m')$ (recalled in the proof) satisfies $\mathsf{N}_y(\mathcal{K})\le\mathsf{U}\,p(y)$. Here $\mathsf{N}_y(\cdot)$ is the local norm at $y$ in which the displacement lemma bounds the field; it majorises $\ell_n|\mathcal{K}(b)|$ on the bonds $b$ of the enlarged cell $\tilde\Delta(y)$ of that representation, and $\ell_n^2$ times the moduli of the covariant differences of $\mathcal{K}$ with respect to the bond variables of $\mathcal{U}(m')$ at the sites of $\tilde\Delta(y)$ (written $\nabla_{\mathbb{W},\mu}\mathcal{K}_\nu$ in the proof, with $\mathbb{W}=\mathcal{U}(m')$).
\item[(b$'$)] For $m'$ and $y$ as in (b), the adjoint actions of the bond variables of $\mathcal{U}(m')$ on $\tilde\Delta(y)$, and of their products along the boundary of a plaquette, have operator norm at most $1.1$.
\item[(c)] The regularity ceiling $\hat\theta(\Delta)\le 1/10$ holds for every cell $\Delta\subset C$.
\item[(d)] The ceiling \textup{(c-W)} of \eqref{4.23:eq-plaquette-units-a} holds.
\end{enumerate}
Define
\begin{equation}\label{4.23:eq-plaquette-units-2}
\begin{aligned}
\mathfrak{N}^W(\Delta)&:=6M^4x_n\hat\theta(\Delta)^2
=54\,c_\alpha^2C_*^2M^4\,\mathsf{w}(\Delta)^2(\log x_n)^{2q},\\
\mathfrak{O}^W(\Delta)&:=6M^4x_n\bigl[4\hat\theta(\Delta)\mathsf{U}p_\Delta
+7\mathsf{U}^2p_\Delta^2\bigr].
\end{aligned}
\end{equation}
Then for every cell $\Delta\subset C$ the function $f^W_\Delta$ is holomorphic on $\mathfrak{P}$, and the following hold:
\begin{itemize}
\item the bound \textup{(W1)} for every cell $\Delta\subset C$;
\item the bound \textup{(W2)} for every cell $\Delta\subset\mathsf{C}$ and all $m,m'\in\mathfrak{P}$ with the same $\varpi$;
\item the two sums \textup{(W3)}, over the cells of $C$.
\end{itemize}
Here the ceiling of (d) and the bounds of the three items are
\begin{equation}\label{4.23:eq-plaquette-units-a}
\begin{aligned}
&\textup{(c-W)}\quad
\sup_{x\ge\gamma^{-2}}40\,c_{\mathsf{U}}C_{\mathsf{a}}\,\alpha_*(x)\le1,\\
&\textup{(W1)}\quad
\sup_{m\in\mathfrak{P}}|f^W_\Delta(m)|\le\mathfrak{N}^W(\Delta),\\
&\textup{(W2)}\quad
|f^W_\Delta(m)-f^W_\Delta(m')|\le\mathfrak{O}^W(\Delta),\\
&\textup{(W3)}\quad
\sum_{\Delta\subset\mathsf{C}}\mathfrak{O}^W(\Delta)
\le C_W(1+N)\check{\mathsf{T}}_k^{2q}\,\mathsf{n}_\partial(k),\\
&\phantom{\textup{(W3)}}\quad
\sum_{\Delta\subset C,\ \Delta\not\subset\mathsf{C}}2\mathfrak{N}^W(\Delta)
\le C_W\,\mathsf{n}_{\mathrm{sh}}(k)\,\check{R}_{h+1}^2\,\check{\mathsf{T}}_k^{2q}.
\end{aligned}
\end{equation}
Here $\mathsf{n}_\partial(k)$ and $\mathsf{n}_{\mathrm{sh}}(k)$ are the counts of \eqref{4.23:eq-plain-cell-sum-a}. The constant $C_W$ is a number fixed before $\gamma$ and after $c_{\mathsf{U}}$, $C_{\mathsf{a}}$, $C_*$, $C_{\mathrm{cell}}$ (the constant of the cell-sum bound). The bounds \textup{(W3)} contain no depth $D$, no volume $|\mathsf{C}|$ and no age.
\end{lemma}

The ceiling (c-W) is imposed for this lemma only. It is a supremum condition on a function of degree $-\tfrac12$ in $x$. Since $\mathsf{U}\le c_{\mathsf{U}}C_{\mathsf{a}}\alpha_{*,k}=c_{\mathsf{U}}C_{\mathsf{a}}\alpha_*(x_k)$ by hypothesis (b), (c-W) applied at $x=x_k\ge\gamma^{-2}$ gives $\mathsf{U}\le1/40$.

\begin{proof}
For $x(p)\in\Delta$ the weight is $x''(p)=x_n$, so $|c(p)|\le x_n$ for every plaquette in the sum \eqref{4.23:eq-plaquette-units-1}.

\emph{Holomorphy.} The function $\mathsf{a}$ is a polynomial in the entries of $W$ and $W^{-1}$, hence holomorphic on invertible matrices. The matrices $\mathcal{U}(m)(\partial p)$, $x(p)\in\Delta$, are invertible, being within $1/10$ of the identity by (a) and (c). The sum \eqref{4.23:eq-plaquette-units-1} is finite. The holomorphy of $f^W_\Delta$ on $\mathfrak{P}$ therefore follows from that of the plaquette variables $m\mapsto\mathcal{U}(m)(\partial p)$, $x(p)\in\Delta$, which is hypothesis (a$'$).

\emph{Proof of (W1).} We first rewrite $\mathsf{a}$. Since $\operatorname{Tr}1=N_c$,
\begin{equation*}
\mathsf{a}(W)=-\frac{1}{2N_c}\operatorname{Tr}\bigl(W+W^{-1}-2\bigr).
\end{equation*}
Moreover $(W-1)^2W^{-1}=W-2+W^{-1}$, so that
\begin{equation*}
\operatorname{Tr}\bigl(W+W^{-1}-2\bigr)=\operatorname{Tr}\bigl((W-1)^2W^{-1}\bigr).
\end{equation*}
If $|W-1|\le\tfrac12$, then $|W^{-1}|\le(1-|W-1|)^{-1}\le2$. Since $|\operatorname{Tr}X|\le N_c|X|$, this gives
\begin{equation}\label{4.23:eq-plaquette-units-3}
|\mathsf{a}(W)|\le\tfrac12|W-1|^2|W^{-1}|\le|W-1|^2
\qquad\text{for }|W-1|\le\tfrac12 .
\end{equation}

By hypothesis (a), applied pointwise, $|\mathcal{U}(m)(\partial p)-1|<\hat\theta(\Delta)(\eta/\ell_n)^2$ for every $m\in\mathfrak{P}$ and every $p$ with $x(p)\in\Delta$. By (c) this is at most $1/10\le\tfrac12$. Hence \eqref{4.23:eq-plaquette-units-3} gives
\begin{equation*}
|\mathsf{a}(\mathcal{U}(m)(\partial p))|\le\hat\theta(\Delta)^2(\eta/\ell_n)^4 .
\end{equation*}

Next we count plaquettes. The cell $\Delta$ holds $6M^4(\ell_n/\eta)^4$ plaquettes: six orientations, $M^4$ sites of level $n$, and $(\ell_n/\eta)^4$ fine sites per site of level $n$. Using $|c(p)|\le x_n$,
\begin{equation*}
|f^W_\Delta(m)|\le6M^4(\ell_n/\eta)^4\,x_n\,\hat\theta(\Delta)^2(\eta/\ell_n)^4
=6M^4x_n\hat\theta(\Delta)^2 .
\end{equation*}
Finally $\hat\theta(\Delta)^2=9c_\alpha^2\mathsf{w}(\Delta)^2\alpha_{*,n}^2$ and $x_n\alpha_{*,n}^2=C_*^2(\log x_n)^{2q}$. This gives the second expression for $\mathfrak{N}^W(\Delta)$ in \eqref{4.23:eq-plaquette-units-2}, and (W1) is proved.

\emph{Proof of (W2).} Let $\Delta\subset\mathsf{C}$, let $m,m'\in\mathfrak{P}$ have the same $\varpi$, and put $\mathbb{W}:=\mathcal{U}(m')$. By the exponential representation of two points of $\mathfrak{P}$ with the same real variables, for each $y\in\Delta$ there is a set $D_2$ containing the enlarged cell $\tilde\Delta(y)$ of that representation, on which, bond by bond,
\begin{equation*}
\mathcal{U}(m)(b)=e^{i\eta\mathcal{K}(b)}\,\mathbb{W}(b).
\end{equation*}
Put $\mathsf{N}:=\mathsf{N}_y(\mathcal{K})$. By hypothesis (b) and (c-W), $\mathsf{N}\le\mathsf{U}p(y)\le\mathsf{U}\le1/40$.

Let $p=(x;\mu,\nu)$ be a plaquette with $x(p)\in\Delta$, and let the point $y\in\Delta$ above be chosen with $p\subset\tilde\Delta(y)$, so that the bonds of $p$ lie in $D_2$. Its bonds are
\begin{equation*}
b_1=\langle x,x+\mu\rangle,\quad b_2=\langle x+\mu,x+\mu+\nu\rangle,\quad
b_3=\langle x+\nu,x+\mu+\nu\rangle,\quad b_4=\langle x,x+\nu\rangle,
\end{equation*}
and we write $\mathcal{K}_i:=\mathcal{K}(b_i)$, $\mathbb{W}_i:=\mathbb{W}(b_i)$ and
\begin{equation*}
W:=\mathbb{W}(\partial p)=\mathbb{W}_1\mathbb{W}_2\mathbb{W}_3^{-1}\mathbb{W}_4^{-1}.
\end{equation*}
Then
\begin{equation*}
\mathcal{U}(m)(\partial p)=e^{i\eta\mathcal{K}_1}\mathbb{W}_1\,e^{i\eta\mathcal{K}_2}\mathbb{W}_2\,
\mathbb{W}_3^{-1}e^{-i\eta\mathcal{K}_3}\,\mathbb{W}_4^{-1}e^{-i\eta\mathcal{K}_4}.
\end{equation*}

We move the four exponential factors to the left, using $Ve^{X}=e^{\operatorname{Ad}_VX}V$ each time:
\begin{itemize}
\item $\mathbb{W}_1e^{i\eta\mathcal{K}_2}=e^{i\eta\operatorname{Ad}_{\mathbb{W}_1}\mathcal{K}_2}\mathbb{W}_1$;
\item since $\mathbb{W}_1\mathbb{W}_2\mathbb{W}_3^{-1}=W\mathbb{W}_4$, the third factor becomes $e^{-i\eta\operatorname{Ad}_W\operatorname{Ad}_{\mathbb{W}_4}\mathcal{K}_3}$;
\item $We^{-i\eta\mathcal{K}_4}=e^{-i\eta\operatorname{Ad}_W\mathcal{K}_4}W$.
\end{itemize}
Therefore
\begin{equation}\label{4.23:eq-plaquette-units-4}
\mathcal{U}(m)(\partial p)W^{-1}
=e^{i\eta\mathcal{K}_1}\,e^{i\eta\operatorname{Ad}_{\mathbb{W}_1}\mathcal{K}_2}\,
e^{-i\eta\operatorname{Ad}_W\operatorname{Ad}_{\mathbb{W}_4}\mathcal{K}_3}\,
e^{-i\eta\operatorname{Ad}_W\mathcal{K}_4}=:e^{\Xi}.
\end{equation}

We now estimate $\Xi$. Write
\begin{equation*}
(\nabla_{\mathbb{W},\mu}\mathcal{K}_\nu)(x):=
\eta^{-1}\bigl[\operatorname{Ad}_{\mathbb{W}(x,x+\mu)}\mathcal{K}_\nu(x+\mu)-\mathcal{K}_\nu(x)\bigr],
\end{equation*}
so that $\eta\nabla_{\mathbb{W},\mu}\mathcal{K}_\nu(x)=\operatorname{Ad}_{\mathbb{W}_1}\mathcal{K}_2-\mathcal{K}_4$ and $\eta\nabla_{\mathbb{W},\nu}\mathcal{K}_\mu(x)=\operatorname{Ad}_{\mathbb{W}_4}\mathcal{K}_3-\mathcal{K}_1$. Adding and subtracting $\operatorname{Ad}_{\mathbb{W}_4}\mathcal{K}_3+\mathcal{K}_4$, the four exponents in \eqref{4.23:eq-plaquette-units-4} sum to
\begin{equation*}
i\eta^2\bigl[(\nabla_{\mathbb{W},\mu}\mathcal{K}_\nu)-(\nabla_{\mathbb{W},\nu}\mathcal{K}_\mu)\bigr](x)
+i\eta(1-\operatorname{Ad}_W)\bigl(\operatorname{Ad}_{\mathbb{W}_4}\mathcal{K}_3+\mathcal{K}_4\bigr).
\end{equation*}
By hypothesis (b), $\mathsf{N}$ majorises $\ell_n|\mathcal{K}_i|$, $i=1,\dots,4$, and $\ell_n^2$ times the moduli of the two covariant derivatives above, and by hypothesis (b$'$) the adjoint actions occurring in \eqref{4.23:eq-plaquette-units-4} have norm at most $1.1$. This controls the two terms as follows.
\begin{itemize}
\item The first term has modulus at most $2(\eta/\ell_n)^2\mathsf{N}$.
\item For the second, $|\eta(\operatorname{Ad}_{\mathbb{W}_4}\mathcal{K}_3+\mathcal{K}_4)|\le2.2(\eta/\ell_n)\mathsf{N}$. Moreover $|1-\operatorname{Ad}_W|\le2.2|W-1|$ at $|W-1|\le\hat\theta(\Delta)\le1/10$. By hypothesis (a) at $m'$, $|W-1|<\hat\theta(\Delta)(\eta/\ell_n)^2$. The second term therefore has modulus at most $2.2(\eta/\ell_n)\mathsf{N}\cdot2.2\hat\theta(\Delta)(\eta/\ell_n)^2$.
\end{itemize}
Each of the four exponents has modulus at most $1.1(\eta/\ell_n)\mathsf{N}$. The Baker--Campbell--Hausdorff expansion of \eqref{4.23:eq-plaquette-units-4} therefore adds at most $10(\eta/\ell_n)^2\mathsf{N}^2$ to the sum of the exponents. Using $\eta\le\ell_n$, then $\hat\theta(\Delta)\le1/10$ by (c) and $\mathsf{N}\le1/40$,
\begin{equation}\label{4.23:eq-plaquette-units-5}
|\Xi|\le(\eta/\ell_n)^2\mathsf{N}\bigl[2+4.84\,\hat\theta(\Delta)+10\mathsf{N}\bigr]
\le3(\eta/\ell_n)^2\mathsf{N}.
\end{equation}

Next we expand $\mathsf{a}$ around $W$. By \eqref{4.23:eq-plaquette-units-4}, $\mathcal{U}(m)(\partial p)=e^{\Xi}W$, and its inverse is $W^{-1}e^{-\Xi}$. Write $e^{-\Xi}-1=-\Xi+(e^{-\Xi}-1+\Xi)$ and use cyclicity of the trace. This gives
\begin{equation*}
\begin{split}
\mathsf{a}(e^{\Xi}W)-\mathsf{a}(W)
&=-\frac{1}{2N_c}\operatorname{Tr}\bigl[\Xi(W-W^{-1})\bigr]\\
&\quad-\frac{1}{2N_c}\operatorname{Tr}\bigl[(e^{\Xi}-1-\Xi)W
+W^{-1}(e^{-\Xi}-1+\Xi)\bigr].
\end{split}
\end{equation*}
For $|W-1|\le1/10$ we have $|W|\le1.1$ and $|W^{-1}|\le10/9$. Also
\begin{equation*}
|W-W^{-1}|\le|W-1|\bigl(1+|W^{-1}|\bigr)\le2.2|W-1|,
\end{equation*}
and $|e^{\pm\Xi}-1\mp\Xi|\le\tfrac12|\Xi|^2e^{|\Xi|}$ with $|\Xi|\le3/40$. Hence
\begin{equation*}
\begin{split}
|\mathsf{a}(e^{\Xi}W)-\mathsf{a}(W)|
&\le1.1|\Xi||W-1|+0.7|\Xi|^2\\
&\le(\eta/\ell_n)^4\bigl[3.3\,\hat\theta(\Delta)\mathsf{N}+6.3\,\mathsf{N}^2\bigr]\\
&\le(\eta/\ell_n)^4\bigl[4\hat\theta(\Delta)\mathsf{U}p_\Delta+7\mathsf{U}^2p_\Delta^2\bigr].
\end{split}
\end{equation*}
The second line uses \eqref{4.23:eq-plaquette-units-5} and $|W-1|<\hat\theta(\Delta)(\eta/\ell_n)^2$. The last line uses $\mathsf{N}\le\mathsf{U}p(y)\le\mathsf{U}p_\Delta$.

Summing over the $6M^4(\ell_n/\eta)^4$ plaquettes of $\Delta$, each treated with its own point $y$, for which in every case $\mathsf{N}\le\mathsf{U}p_\Delta$, and with $|c(p)|\le x_n$, gives $|f^W_\Delta(m)-f^W_\Delta(m')|\le\mathfrak{O}^W(\Delta)$. This is (W2).

\emph{Proof of (W3): preliminaries.} By hypothesis (b), $\mathsf{U}\le c_{\mathsf{U}}C_{\mathsf{a}}\alpha_{*,k}$. The flow inequality in ratio form, with its flow constant $\lambda$, the constant $c_2$, and its level-offset terms $\varsigma_{k-n}$ of the logarithmic form, gives
\begin{equation*}
x_n\le x_k\bigl(1+3\lambda(k-n)+2c_2\bigr),\qquad
\log x_n\le\mathsf{T}_k+\varsigma_{k-n}.
\end{equation*}
Write $P:=1+3\lambda(k-n)+2c_2$. By the formula for $\alpha_*$,
\begin{equation*}
\begin{aligned}
x_n\alpha_{*,n}\alpha_{*,k}
&=C_*^2(x_n/x_k)^{1/2}(\log x_n)^q\mathsf{T}_k^q
\le C_*^2P^{1/2}(\mathsf{T}_k+\varsigma_{k-n})^q\mathsf{T}_k^q,\\
x_n\alpha_{*,k}^2
&=C_*^2(x_n/x_k)\mathsf{T}_k^{2q}\le C_*^2P\,\mathsf{T}_k^{2q}.
\end{aligned}
\end{equation*}
Inserting these and $\hat\theta(\Delta)=3c_\alpha\mathsf{w}(\Delta)\alpha_{*,n}$ into \eqref{4.23:eq-plaquette-units-2} gives
\begin{equation*}
\begin{split}
\mathfrak{O}^W(\Delta)
&\le72\,c_\alpha c_{\mathsf{U}}C_{\mathsf{a}}C_*^2M^4\,\mathsf{w}(\Delta)P^{1/2}
(\mathsf{T}_k+\varsigma_{k-n})^q\mathsf{T}_k^q\,p_\Delta\\
&\quad+42\,c_{\mathsf{U}}^2C_{\mathsf{a}}^2C_*^2M^4\,P\,\mathsf{T}_k^{2q}\,p_\Delta^2 .
\end{split}
\end{equation*}

\emph{Proof of (W3): cells of the core.} Let $\Delta\subset\mathsf{C}$ and write $\varrho:=\varrho(\Delta)\ge1$. We collect the inputs.
\begin{itemize}
\item By the depth bound on the core, $k-n\le N+\varrho$.
\item By the window bound on cell widths, with its width constant $C_{\uparrow}$, $\mathsf{w}(\Delta)\le LC_{\uparrow}\varrho\le LC_{\uparrow}(1+\varrho)$.
\item Since $P\ge1$ and $k-n\le N+\varrho$, $P^{1/2}\le P\le(1+2c_2+3\lambda)(1+N)(1+\varrho)$.
\item With $k-n$ replaced by its upper bound $N+\varrho$, the anchoring estimate and the anchoring cap $\mathsf{T}_k+\varsigma_N\le3\check{\mathsf{T}}_k$ give
\begin{equation*}
\mathsf{T}_k+\varsigma_{N+\varrho}\le(\mathsf{T}_k+\varsigma_N)(1+3\lambda)(1+\varrho)
\le3\check{\mathsf{T}}_k(1+3\lambda)(1+\varrho).
\end{equation*}
\item By the window-size inequality $\mathsf{T}_k\le2\check{\mathsf{T}}_k$, the factor $\mathsf{T}_k^q$ is at most $(2\check{\mathsf{T}}_k)^q$.
\item By the decay of the displacement weight, with its zone width $\mathsf{w}_{\boxplus}$,
\begin{equation*}
p_\Delta\le e^{\delta_P\mathsf{w}_{\boxplus}}e^{-\delta_P(M/M_1)(\varrho-1)} .
\end{equation*}
Since $\varrho\ge1$, the same bound holds for $p_\Delta^2$ with $e^{2\delta_P\mathsf{w}_{\boxplus}}$ in place of $e^{\delta_P\mathsf{w}_{\boxplus}}$.
\end{itemize}
The first term of the preliminary bound then carries $(1+\varrho)^{1+1+q}$ and the second $(1+\varrho)^{1}$. Hence there is a constant $C$, depending only on $c_\alpha,c_{\mathsf{U}},C_{\mathsf{a}},C_*,M,L,C_{\uparrow},\lambda,c_2,q$ and $\delta_P\mathsf{w}_{\boxplus}$, such that
\begin{equation*}
\mathfrak{O}^W(\Delta)\le C(1+N)\check{\mathsf{T}}_k^{2q}(1+\varrho)^{q+2}e^{-\delta_P(M/M_1)(\varrho-1)} .
\end{equation*}

By the lower bound on the contour weight, $\tfrac12\delta_PM/M_1\ge8$. Therefore
\begin{equation*}
(1+\varrho)^{q+2}e^{-\frac12\delta_P(M/M_1)(\varrho-1)}
\le\sup_{t\ge1}(1+t)^{q+2}e^{-8(t-1)},
\end{equation*}
which is a number depending on $q$ only; this uses half of the rate. The other half is $e^{-\frac12\delta_P(M/M_1)(\varrho-1)}=e^{-\frac12\delta_Pd^{\wedge}(\Delta,\mathsf{C}^c)}$, because $\varrho(\Delta)=1+(M_1/M)\,d^{\wedge}(\Delta,\mathsf{C}^c)$. The cell-sum bound states that, for $\beta\ge\delta_P/16$ and a set $\Sigma$, the sum of $e^{-\beta d^{\wedge}(\Delta,\Sigma)}$ over the cells $\Delta\subset\mathsf{C}$ is at most $C_{\mathrm{cell}}$ times the outer boundary count of $\Sigma$. It is applied with $\beta=\tfrac12\delta_P\ge\delta_P/16$ and $\Sigma=\mathsf{C}^c$. For cells $\Delta\subset\mathsf{C}$ the distance $d^{\wedge}(\Delta,\mathsf{C}^c)$ is attained where $\mathsf{C}$ meets its complement, so the boundary count that enters is that of the core, $\#\partial^{+}\mathsf{C}$; by the boundary count of the core in the geometry of components, $\#\partial^{+}\mathsf{C}\le\mathsf{n}_\partial(k)$. This gives
\begin{equation*}
\sum_{\Delta\subset\mathsf{C}}e^{-\frac12\delta_Pd^{\wedge}(\Delta,\mathsf{C}^c)}\le C_{\mathrm{cell}}\,\mathsf{n}_\partial(k).
\end{equation*}
This proves the first inequality of (W3).

\emph{Proof of (W3): cells off the core.} Let $\Delta\subset C$ with $\Delta\not\subset\mathsf{C}$. By the shallow-shell bound for plain units, one of the bounds behind the cell sum \eqref{4.23:eq-plain-cell-sum-a}, such a cell has level $n\ge h$, so $k-n\le N$, it has width $\mathsf{w}(\Delta)\le L\check{R}_{h+1}$, and there are at most $\mathsf{n}_{\mathrm{sh}}(k)$ such cells. Since $k-n\le N$, the flow inequality gives $\log x_n\le\mathsf{T}_k+\varsigma_N$, and this is at most $3\check{\mathsf{T}}_k$ by the anchoring cap used above. Hence by \eqref{4.23:eq-plaquette-units-2}
\begin{equation*}
2\mathfrak{N}^W(\Delta)\le108\,c_\alpha^2C_*^2M^4L^2\,3^{2q}\,\check{R}_{h+1}^2\check{\mathsf{T}}_k^{2q}
\end{equation*}
for each such cell, and summing over at most $\mathsf{n}_{\mathrm{sh}}(k)$ cells gives the second inequality of (W3).

Both constants obtained are free of $D$, of $|\mathsf{C}|$ and of the age.
\end{proof}

Treated as interior units of the first part of the localised family, in the setting of Lemma~\ref{4.23:lem-unit-oscillation}, with $K_v:=[\Delta]_k$, the level-$k$ cube containing $\Delta$, so that $d_k(K_v)=0$ by the rule for interior units of that setting, the Wilson units $f^W_\Delta$ satisfy the total-oscillation bound \eqref{4.23:eq-unit-oscillation-a} of Lemma~\ref{4.23:lem-unit-oscillation}, whose proof is incomplete at one step. The bound is applied with
\begin{itemize}
\item $\mathfrak{n}(v):=\mathfrak{O}^W(\Delta)$ for $\Delta\subset\mathsf{C}$, by (W2);
\item $\mathfrak{n}(v):=2\mathfrak{N}^W(\Delta)$ otherwise, by (W1), since the oscillation of a function is at most twice its supremum.
\end{itemize}
Thus $\operatorname{tot}(v)\le e^{4}\mathfrak{n}(v)$. The sum of $\operatorname{tot}$ over these units is at most $e^4$ times the right members of (W3), so they belong to class (rob)($\alpha$) of \eqref{4.23:eq-component-constituents-a}.

\begin{remark}[Reading of the right member without the boundary-term bound on the core]
The lemma is not used in the proof of the oscillation theorem. That proof uses the bound of Definition~\ref{4.23:def-boundary-terms} (printed without proof in \cite[\S1]{B-CMP122b}; cited as printed) on the whole of $C^{\sharp}$. Suppose instead that this bound is used only for the members localised off the core. Then in the explicit right member \eqref{4.23:eq-right-member-a} the term $2\mathsf{c}_{\mathbf{C}}$ is to be read as $2\mathsf{c}^{\mathrm{off}}_{\mathbf{C}}$, the same sum restricted to the members localised off the core, plus $e^4$ times the right members of (W3), and the exponent $p_\Sigma$ is increased by $\lceil 2q/r_{\mathrm{pr}}\rceil$, where $r_{\mathrm{pr}}$ is the exponent of the window-radius lower bound $\check{R}_j\ge\check{\mathsf{T}}_j^{r_{\mathrm{pr}}}$; nothing else changes.

The flow inequality enters the proof only as a ratio of couplings over the $k-n$ levels between $n$ and $k$. This ratio is paid for by the zones crossed at level $k$, and no power $\check{R}_j^{p_\Sigma}$ of a window radius is spent on it.

The lemma does not treat the term $I(U_0)$, the members depending on $U_0$, or the identification of Ba{\l}aban's intermediate localisations beyond linearity in the weight. These belong to the content of Definition~\ref{4.23:def-boundary-terms}, the statement printed without proof in \cite[\S1]{B-CMP122b} that is cited as printed.
\end{remark}

In what follows the level $k$ is fixed. Let $C$ be a component treated at step $k$, let $\mathsf{C}=\mathsf{C}_C$ be its core, and let $h=k-N(g_k)$ be the bottom level of the window.

A \emph{nested region} of $C$ is a pair $(Y',m')$, where $Y'\subset C$ is a large-field region created at step $m'<k$. For a nested region:
\begin{itemize}
\item $Y'^{\sharp}$ is its strip;
\item $\mathfrak{K}(Y')$ is the family of units born with $Y'$ at step $m'$;
\item $V^{\sharp}(Y'^{\sharp})$ and $\mathfrak{q}_{Y'}$ are the density and the quadratic form born with $Y'$;
\item $\mathfrak{O}(v)$ is the oscillation of a unit $v$, and $\operatorname{tot}(v)$ is the total oscillation of
Lemma~\ref{4.23:lem-unit-oscillation}.
\end{itemize}
The distance $d^{\wedge}$ is the own-scale contour distance and $w=M/M_1$. The rate attached to $d^{\wedge}$ is
$\delta_P=\delta_0/8$, where $\delta_0$ is the decay rate fixed for this chapter. The depth of a cell $\Delta$ is
$\varrho(\Delta)=1+d^{\wedge}(\Delta,\mathsf{C}^c)/w$.

Let $p_\Sigma$ and $C_\Sigma$ be the degree and the constant of clause (t) of
Theorem~\ref{4.23:thm-oscillation-theorem}; they are fixed in \eqref{4.23:eq-degree-constant-a}. Let
$r_{\mathrm{pr}}$ be the exponent in the lower bound on the window radius. Let $\mathsf{T}_\gamma$ be the number, fixed by
the ceilings on $\gamma$, that enters the second inequality of that bound. The bound reads
\begin{equation*}
\check{R}_j\ge\check{\mathsf{T}}_j^{\,r_{\mathrm{pr}}}\ge\mathsf{T}_\gamma^{\,r_{\mathrm{pr}}}\qquad(j\le K).
\end{equation*}
We put
\begin{equation}\label{4.23:eq-nested-regions-a}
\begin{gathered}
(\mathrm{V})_k:=\text{clauses }(\sharp),\ (\mathrm{v}),\ (\mathrm{g})
\text{ of Theorem~\ref{4.23:thm-oscillation-theorem} at step }k,\\
R_\gamma:=\mathsf{T}_\gamma^{\,r_{\mathrm{pr}}},\qquad
\varepsilon_\gamma:=\sup_{\varrho\ge\frac12 R_\gamma}(2\varrho)^{p_\Sigma}e^{-2(\varrho-1)}.
\end{gathered}
\end{equation}
Clause (t) of the oscillation theorem is not part of $(\mathrm{V})_k$. It is an inequality between explicit functions and
is established directly for every $k$, without induction, with the degree and the constant of
\eqref{4.23:eq-degree-constant-a}; the induction concerns clause (v). By the lower bound on the window radius,
$R_\gamma\le\check{R}_j$ for all $j\le K$. The number $\varepsilon_\gamma$ has degree $-\infty$ in $\mathsf{T}_\gamma$,
in the sense of the degree table of \eqref{4.23:eq-degree-constant-a}.

\begin{lemma}[The nested-region lemma]\label{4.23:lem-nested-regions}
Assume the following:
\begin{itemize}
\item for every $m'<k$, the statement $(\mathrm{V})_{m'}$ of \eqref{4.23:eq-nested-regions-a} holds, together with
\textup{(q0)} and \textup{(q1)} of \eqref{4.23:eq-tube-statement-a} at step $m'$;
\item clause \textup{(t)} of Theorem~\ref{4.23:thm-oscillation-theorem} holds at every step $m'<k$.
\end{itemize}
Then
\begin{equation}\label{4.23:eq-nested-regions-b}
\begin{aligned}
\operatorname{nest}(k)&:=\sum_{(Y',m'):\,Y'\subset C}\ \sum_{v'\in\mathfrak{K}(Y')}\operatorname{tot}(v')\\
&\le\operatorname{nest}^*(k):=\tfrac12 C_{3F}e^{4}\cdot\varepsilon_\gamma\cdot
C_{\mathrm{cell}}\,\mathsf{n}_\partial(k)\cdot C_\Sigma .
\end{aligned}
\end{equation}
Here $C_{3F}$ is the constant of the three-family bound. The constant $C_{\mathrm{cell}}$ is that of the cell-counting
bound recalled in the fourth step of the proof below. The number $\mathsf{n}_\partial(k)$ is the boundary count of
\eqref{4.23:eq-plain-cell-sum-a}.

Moreover, for every real number $s$, used here as a distance threshold, the nested regions with
$d^{\wedge}(Y'^{\sharp},\mathsf{C}^c)\ge s$ contribute at most $e^{-\frac18\delta_P s}\operatorname{nest}^*(k)$ to
$\operatorname{nest}(k)$.

The bound does not depend on the number, the sizes, the ages, the histories or the nesting of the nested regions. Nor
does it depend on the sizes of their members, that is, on the bound \cite[(1.69)]{B-CMP122b} at $Y'$.
\end{lemma}

\begin{proof}
We proceed in four steps.

\smallskip
\noindent\textit{First step: oscillation at birth.}
Fix a nested region $(Y',m')$ of $C$. We apply the oscillation bound at birth at step $m'$ and to the component $Y'$.
We use it in the form in which condition (q0) is added to its hypotheses.

Let $v$ be a unit born with a treated component, with tube $\mathfrak{U}^{\mathrm{tr}}_v$. As stated, that bound has
two hypotheses. The first is an oscillation bound for its datum on $\mathfrak{U}^{\mathrm{tr}}_v$, uniform in the real
variables $\Psi$; the second is (q1). Under them it asserts
\begin{equation*}
\mathfrak{O}(v)\le\sup_{\Psi\ \mathrm{real}}\ \sup_{\zeta,\zeta'\in\mathfrak{U}^{\mathrm{tr}}_v}
|v(\zeta;\Psi)-v(\zeta';\Psi)|.
\end{equation*}
It also asserts that the oscillation of the born density plus that of the born quadratic form is at most
$\Sigma^{\mathrm{tot}}$ at the birth step.

Its hypotheses hold here, as follows.
\begin{itemize}
\item The tube is $\mathfrak{U}^{\mathrm{tr}}_{v'}=\tilde{U}^{c}_{m'}(Y'^{\sharp},\tilde\alpha_0,\tilde\alpha_1)$.
\item The oscillation hypothesis on this tube, uniform in real $\Psi$, is clause (v) of $(\mathrm{V})_{m'}$. The tube of
that clause is this very set, and the oscillation it bounds is this very supremum
(Definition~\ref{4.23:not-oscillation-seminorm}).
\item The premises on the datum are holomorphy, a bound and joint invariance of $V^{\sharp}(Y'^{\sharp})$. These are
clause (g) of Theorem~\ref{4.23:thm-oscillation-theorem} at step $m'$.
\item Condition (q0) gives splitting-freeness, holomorphy and joint invariance of $\mathfrak{q}_{Y'}$. The oscillation
bound at birth needs the holomorphy and the joint invariance of the born quadratic form. Splitting-freeness makes
$\mathfrak{q}_{Y'}$ a function on the tube at all.
\item Condition (q1) holds on the tube.
\end{itemize}
The last two items are the form theorem's (X0) together with its joint-invariance clause, and its (X1), in
\eqref{4.23:eq-form-theorem-a}, at step $m'$ and restricted to the tube. This is the content of (q0) and (q1) in
\eqref{4.23:eq-tube-statement-a}, assumed at $m'$. Hence
\begin{equation}\label{4.23:eq-nested-regions-1}
\mathfrak{O}(V^{\sharp}(Y'^{\sharp}))+\mathfrak{O}(\mathfrak{q}_{Y'})
\le\Sigma^{V}(m')+2\mathfrak{q}_{\max}(m')=\Sigma^{\mathrm{tot}}(m')
\le C_\Sigma\check{R}_{m'}^{\,p_\Sigma}.
\end{equation}
The last inequality is clause (t) at step $m'$. Here $\Sigma^{V}(m')$ is the bound of clause (v) at step $m'$. The
quantity $\Sigma^{\mathrm{tot}}(m')$, defined by the equality, is the one bounded in clause (t); both are in
\eqref{4.23:eq-oscillation-theorem-a}.

\smallskip
\noindent\textit{Second step: one nested region.}
Lemma~\ref{4.23:lem-unit-oscillation}, whose proof is incomplete at one step, gives
$\operatorname{tot}(v)\le\mathfrak{n}(v)e^{4(d_k(K_v)+1)}$, with $\mathfrak{n}(v)$ given by the table
\eqref{4.23:eq-unit-oscillation-a}. The members $v'$ of $\mathfrak{K}(Y')$ fall under the fourth entry of that table.

That entry is supplied by the variation bound for members of nested regions. This bound holds for $v'$ because of
three inputs:
\begin{itemize}
\item the holomorphy of $v'$ at the present step, given by clauses (a) and (c) of the holomorphy lemma for members of
nested regions;
\item clause (c) of the second lemma on members of nested regions, in its strengthened form;
\item the three-family bound, taken at the set $\Sigma_{Y'}$ attached to $Y'$ that enters the definition of the
oscillation $\mathfrak{O}$ of its members.
\end{itemize}
The same inputs settle a requirement of that entry: $v'$ does not belong to the class of units written $\mathcal{C}$.
This class is distinct from the component $C$ and from its core $\mathsf{C}$.

The set $K_{v'}$ is a single cube, so $d_k(K_{v'})=0$ and the exponential factor is $e^{4}$. Put
$D':=d^{\wedge}(Y'^{\sharp},\mathsf{C}^c)$. By \eqref{4.23:eq-nested-regions-1},
\begin{equation}\label{4.23:eq-nested-regions-2}
\begin{aligned}
\sum_{v'\in\mathfrak{K}(Y')}\operatorname{tot}(v')
&\le e^{4}\cdot\tfrac12 C_{3F}e^{-\frac12\delta_P D'}
\bigl[\mathfrak{O}(V^{\sharp}(Y'^{\sharp}))+\mathfrak{O}(\mathfrak{q}_{Y'})\bigr]\\
&\le\tfrac12 C_{3F}e^{4}e^{-\frac12\delta_P D'}C_\Sigma\check{R}_{m'}^{\,p_\Sigma}.
\end{aligned}
\end{equation}

\smallskip
\noindent\textit{Third step: a cell for each nested region, and its depth.}
The variation bound for members of nested regions carries the two cases of the depth statement:
\begin{equation*}
D'\ge\tfrac12 w(\check{R}_{m'}+1)\ \text{ if } m'=h,\qquad
D'\ge\tfrac52 w\check{R}_{m'}\ \text{ if } m'<h .
\end{equation*}

Choose a cell $\Delta_{Y'}$ containing a point of $Y'^{\sharp}$ nearest to $\mathsf{C}^c$ in $d^{\wedge}$.
\begin{itemize}
\item By clause (b) of the geometry statement for nested regions (a cell that meets a strip lies in it),
$\Delta_{Y'}\subset Y'^{\sharp}$.
\item Consequently $d^{\wedge}(\Delta_{Y'},\mathsf{C}^c)=D'$. The distance is at least $D'$ because
$\Delta_{Y'}\subset Y'^{\sharp}$, and at most $D'$ because $\Delta_{Y'}$ contains a nearest point.
\item By the anchoring statement for nested regions, the strips of distinct live nested regions are disjoint and a cell
lies in at most one of them. Since $\Delta_{Y'}\subset Y'^{\sharp}$, the map $Y'\mapsto\Delta_{Y'}$ is injective.
\end{itemize}

Now put $\varrho:=\varrho(\Delta_{Y'})=1+D'/w$. We claim $\varrho\ge1+\tfrac12(\check{R}_{m'}+1)$ in both cases.
In the first case this is the displayed lower bound divided by $w$. In the second case it follows from
\begin{equation*}
\tfrac52\check{R}_{m'}-\tfrac12(\check{R}_{m'}+1)=2\check{R}_{m'}-\tfrac12\ge0 .
\end{equation*}
Two consequences follow. First, $\check{R}_{m'}\le2\varrho$. Second, $\varrho\ge\frac12\check{R}_{m'}\ge\frac12R_\gamma$,
by the lower bound on the window radius recorded after \eqref{4.23:eq-nested-regions-a}, applied at $j=m'<k\le K$.

\smallskip
\noindent\textit{Fourth step: three slices of the rate.}
Split the rate as
\begin{equation*}
\tfrac12\delta_P=\tfrac18\delta_P+\tfrac18\delta_P+\tfrac14\delta_P .
\end{equation*}

The first slice absorbs the power of the window radius. The floor $\frac1{16}\delta_P w\ge1$ gives
$\frac18\delta_P w\ge2$. We have $D'=w(\varrho-1)$ and $\varrho-1\ge0$, and $\check{R}_{m'}\le2\varrho$ by the third
step. Hence
\begin{equation*}
e^{-\frac18\delta_P D'}\check{R}_{m'}^{\,p_\Sigma}
=e^{-\frac18\delta_P w(\varrho-1)}\check{R}_{m'}^{\,p_\Sigma}
\le(2\varrho)^{p_\Sigma}e^{-2(\varrho-1)}\le\varepsilon_\gamma .
\end{equation*}
The last inequality is the definition \eqref{4.23:eq-nested-regions-a} of $\varepsilon_\gamma$, because
$\varrho\ge\frac12R_\gamma$.

The second slice is summed over the nested regions with the cell-counting bound. For every rate
$\beta\ge\delta_P/16$ and every set $S$, that bound states
\begin{equation*}
\sum_{\Delta\subset\mathsf{C}}e^{-\beta d^{\wedge}(\Delta,S)}\le C_{\mathrm{cell}}\,\#\partial^{+}S ,
\end{equation*}
where $\#\partial^{+}S$ denotes the number of boundary cells of $S$. We take
$\beta=\frac14\delta_P\ge\delta_P/16$ and $S=\overline{\mathsf{C}^c}$, the closure of $\mathsf{C}^c$. We use three
facts:
\begin{itemize}
\item $d^{\wedge}(\Delta_{Y'},\mathsf{C}^c)=D'$, by the choice of $\Delta_{Y'}$;
\item $Y'\mapsto\Delta_{Y'}$ is injective;
\item $\Delta_{Y'}\subset Y'^{\sharp}\subset\mathsf{C}$, since the strips of the nested regions of $C$ lie in
$\mathsf{C}$ by clause (g) of the geometry statement.
\end{itemize}
With these,
\begin{equation*}
\begin{aligned}
\sum_{(Y',m')}e^{-\frac14\delta_P D'}
&=\sum_{(Y',m')}e^{-\frac14\delta_P d^{\wedge}(\Delta_{Y'},\mathsf{C}^c)}
\le\sum_{\Delta\subset\mathsf{C}}e^{-\frac14\delta_P d^{\wedge}(\Delta,\mathsf{C}^c)}\\
&\le C_{\mathrm{cell}}\,\#\partial^{+}\overline{\mathsf{C}^c}\le C_{\mathrm{cell}}\,\mathsf{n}_\partial(k),
\end{aligned}
\end{equation*}
where the last inequality is again clause (g) of the geometry statement.

The third slice gives the second assertion of the lemma. It satisfies $e^{-\frac18\delta_P D'}\le1$ for every nested
region, and $e^{-\frac18\delta_P D'}\le e^{-\frac18\delta_P s}$ for the nested regions with $D'\ge s$.

\smallskip
\noindent\textit{Conclusion.}
Sum \eqref{4.23:eq-nested-regions-2} over the nested regions of $C$, inserting the three slices:
\begin{equation*}
\begin{aligned}
\operatorname{nest}(k)
&\le\tfrac12 C_{3F}e^{4}C_\Sigma\sum_{(Y',m')}
\Bigl(e^{-\frac18\delta_P D'}\check{R}_{m'}^{\,p_\Sigma}\Bigr)
e^{-\frac18\delta_P D'}e^{-\frac14\delta_P D'}\\
&\le\tfrac12 C_{3F}e^{4}C_\Sigma\cdot\varepsilon_\gamma\cdot C_{\mathrm{cell}}\,\mathsf{n}_\partial(k)
=\operatorname{nest}^*(k).
\end{aligned}
\end{equation*}
Restrict the sum to the nested regions with $D'\ge s$ and use the second bound on the third slice. Their contribution
is then at most $e^{-\frac18\delta_P s}\operatorname{nest}^*(k)$.

No quantity in these bounds depends on the number, sizes, ages, histories or nesting of the nested regions, or on the
sizes of their members.
\end{proof}

\begin{proposition}[Oscillation theorem: the explicit right member]\label{4.23:prf-right-member}
Work in the setting and under the hypotheses of the oscillation theorem for the resummed density
(Theorem~\ref{4.23:thm-oscillation-theorem}). For a step $k\le K$ put
\begin{equation}\label{4.23:eq-right-member-a}
\begin{split}
\Sigma^{V}(k) :={}& 2\,\mathsf{n}_{\sharp}(k)\,C_{\mathrm{fib}}
\bigl[3+\mathsf{K}^{0}+N(g_k)\,C_0^{\sharp}\bigr]\\
&+ e^{4}\bigl[C^{+}\,\mathsf{S}^{\mathrm{dp}}_{\mathrm{pl}}(k)
+2\,\mathsf{S}^{\mathrm{sh}}_{\mathrm{pl}}(k)\bigr]
+ e^{4}\bigl[C_{3\mathrm{F}}\,\mathsf{S}^{\mathrm{dp}}_{\mathrm{wr}}(k)
+2\,\mathsf{S}^{\mathrm{sh}}_{\mathrm{wr}}(k)\bigr]\\
&+ 2\,(100\,\check{R}_k)^{d}\,\mathsf{K}^{\mathrm{cp}}
+ 2\,\bigl(1+c_{\mathrm{I}}M^{5}\bigr) + \operatorname{nest}^{*}(k).
\end{split}
\end{equation}
Here $\mathsf{n}_{\sharp}(k)$, $C_{\mathrm{fib}}$, $\mathsf{K}^{0}$ and $C_0^{\sharp}$ are those of the
rim lemma (Lemma~\ref{4.23:lem-rim-lemma}, \eqref{4.23:eq-rim-lemma-a}); $C^{+}$,
$\mathsf{S}^{\mathrm{dp}}_{\mathrm{pl}}(k)$ and
$\mathsf{S}^{\mathrm{sh}}_{\mathrm{pl}}(k)$ those of \eqref{4.23:eq-plain-cell-sum-a};
$C_{3\mathrm{F}}$ is the constant entering $C^{+}$ in \eqref{4.23:eq-plain-cell-sum-a}, which is also
the constant of that name in \eqref{4.23:eq-nested-regions-b};
$\mathsf{S}^{\mathrm{dp}}_{\mathrm{wr}}(k)$ and $\mathsf{S}^{\mathrm{sh}}_{\mathrm{wr}}(k)$ those of
\eqref{4.23:eq-wrapped-cell-sum-a}; $\mathsf{K}^{\mathrm{cp}}$ that of \eqref{4.23:eq-copies-lemma-a};
$c_{\mathrm{I}}$ and $M$ those of \eqref{4.23:eq-boundary-oscillation-a}; and $\operatorname{nest}^{*}(k)$ that
of \eqref{4.23:eq-nested-regions-b}.

Then, for every label and every component $C$ of $Z$ at the large-field operation $\mathbf{R}$ at
level $k$, clauses $(\sharp)$ and $(\mathrm{v})$ of Theorem~\ref{4.23:thm-oscillation-theorem} hold:
\begin{itemize}
\item[$(\sharp)$] $V^{\sharp}(C^{\sharp};\cdot\,;\Phi_C,w)$ is holomorphic on
$\mathfrak{S}=\tilde{U}^{c}_{k}(C^{\sharp},\tilde\alpha_0,\tilde\alpha_1)$ and reads $\zeta$ on
$C^{\sharp}$ only;
\item[$(\mathrm{v})$] $\operatorname{osc}\bigl(V^{\sharp}(C^{\sharp})\bigr)\le\Sigma^{V}(k)$.
\end{itemize}
\end{proposition}

\begin{proof}
The six members of \eqref{4.23:eq-right-member-a} account, in this order, for the rim of the component;
the plain units, the Gaussian terms filed among them included; the wrapped units; the copies; the
localised $z$-free part of the boundary terms $\mathbf{C}_k$; and the older large-field regions nested
inside the component.

Throughout, $\operatorname{osc}$ is the two-point oscillation over $\mathfrak{S}$ at one fixed real
value of $\Phi_C$ of Theorem~\ref{4.23:thm-oscillation-theorem}. It is a seminorm:
$\operatorname{osc}(f+f')\le\operatorname{osc}(f)+
\operatorname{osc}(f')$ and $\operatorname{osc}(cf)=|c|\operatorname{osc}(f)$. Moreover
$\operatorname{osc}(f)\le 2\sup_{\mathfrak{S}}|f|$, and $\operatorname{osc}(f)=0$ when $f$ does not
depend on $\zeta$. All sums over $Y_4\subset C^{\sharp}$ below run over the filing domains
$Y_4\in\mathbf{D}_k^{Z}$ of Definition~\ref{4.23:def-resummed-density} with $Y_4\subset C^{\sharp}$.

\emph{Clause $(\sharp)$.} By Definition~\ref{4.23:def-resummed-density}, $V^{\sharp}(C^{\sharp})$ is a
finite sum of the densities $\mathfrak{v}(Y_4)$, and each of these is a finite sum of functions, each
analytic on a tube over a sub-domain of $C^{\sharp}$. The argument of the holomorphy corollary for the
resummed sums, with the floor condition attached to that corollary in force, therefore shows that
$V^{\sharp}(C^{\sharp};\cdot\,;\Phi_C,w)$ is holomorphic on $\mathfrak{S}$. For the locality, we use
four inputs. First, by the third clause of the source-locality statement, a term filed at $Y_4$ reads
$\zeta$ on the enlargement $Y_4^{+}$ of $Y_4$; by the strip condition on the filing domains, a domain
$Y_4\subset C^{\sharp}$ touches no strip of an older large-field region that is still live, so that
$Y_4^{+}=Y_4$. Second, we invoke the first clause of the locality statement for the units of the
chain. Third, by part~(i) of the copies lemma (Lemma~\ref{4.23:lem-copies-lemma}) each copy $f_v$
reads $\zeta$ on a set contained in $C$. Fourth, the localised $z$-free members
$\mathfrak{v}^{\mathbf{C}}$ are localised in $C^{\sharp}$, which is part of the input on the localised
boundary terms used in Lemma~\ref{4.23:lem-boundary-oscillation}. Together these give the locality
part of $(\sharp)$: $V^{\sharp}(C^{\sharp})$ reads $\zeta$ on $C^{\sharp}$ only.

\emph{Clause $(\mathrm{v})$.} We argue by induction on $k$. Assume the statement
$(\mathrm{V})_{m'}$ of \eqref{4.23:eq-nested-regions-a} for every $m'<k$. At the first step at which
the operation $\mathbf{R}$ acts, no nested older region exists and $\operatorname{nest}(k)=0$, so this
assumption is void there. Clause $(\mathrm{t})$ of Theorem~\ref{4.23:thm-oscillation-theorem} is
proved in the part of the proof of that theorem which follows the present one, without appeal to
clause $(\mathrm{v})$, so it is available at every step.

By part~(ii) of the filing-domains lemma (Lemma~\ref{4.23:lem-filing-domains},
\eqref{4.23:eq-filing-domains-a}),
\begin{equation*}
V^{\sharp}(C^{\sharp})=\mathfrak{v}(C)+\sum_{Y_4\in\mathcal{Y}_{\mathrm{rim}}}\mathfrak{v}(Y_4).
\end{equation*}
By part~(i) of the same lemma every $Y_4\in\mathbf{D}_k^{Z}$ with $Y_4\subset C^{\sharp}$ contains $C$,
so these domains are exactly $C$ and the elements of
$\mathcal{Y}_{\mathrm{rim}}$. Writing $\mathfrak{v}=(\mathfrak{v}-\mathfrak{v}^{\mathbf{C}})
+\mathfrak{v}^{\mathbf{C}}$ in each term, collecting the $\mathfrak{v}^{\mathbf{C}}$-parts, and using
the seminorm property together with $\operatorname{osc}\le2\sup$ on each rim term, we obtain
\begin{equation}\label{4.23:eq-right-member-1}
\begin{split}
\operatorname{osc}\bigl(V^{\sharp}(C^{\sharp})\bigr)\le{}&
\operatorname{osc}\bigl((\mathfrak{v}-\mathfrak{v}^{\mathbf{C}})(C)\bigr)
+\sum_{Y_4\in\mathcal{Y}_{\mathrm{rim}}}2\sup_{\mathfrak{S}}
\bigl|(\mathfrak{v}-\mathfrak{v}^{\mathbf{C}})(Y_4)\bigr|\\
&+\operatorname{osc}\Bigl(\sum_{Y_4\subset C^{\sharp}}\mathfrak{v}^{\mathbf{C}}(Y_4)\Bigr).
\end{split}
\end{equation}

The second member of \eqref{4.23:eq-right-member-1} is $2\,\mathsf{rim}(k)$, with $\mathsf{rim}(k)$
as in \eqref{4.23:eq-rim-lemma-a}. By the rim lemma
(Lemma~\ref{4.23:lem-rim-lemma}) it is at most
$2\,\mathsf{n}_{\sharp}(k)\,C_{\mathrm{fib}}\,\mathsf{N}$ with
$\mathsf{N}=3+\mathsf{K}^{0}+N(g_k)C_0^{\sharp}$, which is the first member of
\eqref{4.23:eq-right-member-a}.

The third member is at most $2(1+c_{\mathrm{I}}M^{5})$ by \eqref{4.23:eq-boundary-oscillation-a}
(Lemma~\ref{4.23:lem-boundary-oscillation}).

For the first member, the list of constituents filed at the component
(Definition~\ref{4.23:def-component-constituents}, \eqref{4.23:eq-component-constituents-a}) gives
\begin{equation*}
(\mathfrak{v}-\mathfrak{v}^{\mathbf{C}})(C)
=\sum_{v}\ \sum_{\vec{\mathsf{y}}:\ \operatorname{lab}_4\subset C}\mathfrak{p}(v;\vec{\mathsf{y}})
-\sum_{v\in\mathfrak{L}^{\mathrm{cp}}(C)}f_v+(\text{$\zeta$-free terms}),
\end{equation*}
where $v$ runs over the units of (m1)--(m3); here $\mathfrak{p}(v;\vec{\mathsf{y}})$ are the pieces of
the unit $v$ indexed by the tuples $\vec{\mathsf{y}}$ with label set $\operatorname{lab}_4$, as in
Definition~\ref{4.23:def-component-constituents}, and $f_v$, $v\in\mathfrak{L}^{\mathrm{cp}}(C)$, are
the copies of Lemma~\ref{4.23:lem-copies-lemma}. The $\zeta$-free terms have oscillation zero. By the
seminorm property,
\begin{equation*}
\operatorname{osc}\bigl((\mathfrak{v}-\mathfrak{v}^{\mathbf{C}})(C)\bigr)
\le\sum_{v}\operatorname{tot}(v)
+\operatorname{osc}\Bigl(\sum_{v\in\mathfrak{L}^{\mathrm{cp}}(C)}f_v\Bigr).
\end{equation*}
Here $\operatorname{tot}(v)$ is the total oscillation of one unit defined in
Lemma~\ref{4.23:lem-unit-oscillation}, whose proof is incomplete at one step; only its definition
\eqref{4.23:eq-unit-oscillation-a} is used at this point. The inequality holds because
$\operatorname{tot}(v)$ sums the oscillations of $\mathfrak{p}(v;\vec{\mathsf{y}})$ over more tuples,
namely all those with $\operatorname{lab}_4\subset C^{\sharp}$, and the added terms are non-negative.

The sum of $\operatorname{tot}(v)$ over the plain units is at most
$e^{4}[C^{+}\mathsf{S}^{\mathrm{dp}}_{\mathrm{pl}}(k)+2\mathsf{S}^{\mathrm{sh}}_{\mathrm{pl}}(k)]$ by the
cell sum for plain units (Lemma~\ref{4.23:lem-plain-cell-sum}). Over the wrapped units it is at most
$e^{4}[C_{3\mathrm{F}}\mathsf{S}^{\mathrm{dp}}_{\mathrm{wr}}(k)+2\mathsf{S}^{\mathrm{sh}}_{\mathrm{wr}}(k)]$
by the cell sum for wrapped units (Lemma~\ref{4.23:lem-wrapped-cell-sum}). Over the units filed at the
older nested regions it is at most $\operatorname{nest}(k)\le\operatorname{nest}^{*}(k)$ by the
nested-region lemma (Lemma~\ref{4.23:lem-nested-regions}). The hypotheses of that lemma,
$(\mathrm{V})_{m'}$ for every $m'<k$ and clause $(\mathrm{t})$, hold by the induction hypothesis and by
the remark above. Finally, part~(iii) of the copies lemma (Lemma~\ref{4.23:lem-copies-lemma}) bounds
$\operatorname{osc}(\sum_{v\in\mathfrak{L}^{\mathrm{cp}}(C)}f_v)$ by
$2(100\check{R}_k)^{d}\mathsf{K}^{\mathrm{cp}}$.

Adding the bounds for the three members of \eqref{4.23:eq-right-member-1} gives
$\operatorname{osc}(V^{\sharp}(C^{\sharp}))\le\Sigma^{V}(k)$, which is clause $(\mathrm{v})$.
\end{proof}

\begin{lemma}[Oscillation theorem: holomorphy, bound and gauge invariance]\label{4.23:prf-density-invariance}
This is clause (g) of Theorem~\ref{4.23:thm-oscillation-theorem}. Write $\mathcal{N}$ for the resummed density of that theorem, a function of a point $\zeta=(\mathbf{U},\mathbf{J})$ of the tube and of the real variables $\Phi_C$ on which it depends beside $\zeta$. Then $\mathcal{N}$ is holomorphic in $\zeta$, its resummed datum obeys the size bound of clause (c) of the resummation proposition, and for every $G$-valued gauge transformation $g$,
\begin{equation}\label{4.23:eq-density-invariance-1}
\mathcal{N}(\zeta^{g};\Phi_C^{\tilde g})=\mathcal{N}(\zeta;\Phi_C).
\end{equation}
Here $\zeta\mapsto\zeta^{g}$ and $\Phi_C\mapsto\Phi_C^{\tilde g}$ denote the action of $g$ in the gauge-transformation law of the per-filing-domain source statement, with $\tilde g$ the transformation that this law attaches to $g$ on the real variables.
\end{lemma}

\begin{proof}
The argument in which clause (g) is used needs only three properties of the density: its holomorphy, a bound on its resummed datum, and its joint invariance \eqref{4.23:eq-density-invariance-1} for every $G$-valued $g$. We supply the three properties in turn.

\emph{Holomorphy.} This is clause $(\sharp)$ of Theorem~\ref{4.23:thm-oscillation-theorem}, display~\eqref{4.23:eq-oscillation-theorem-a}.

\emph{Bound.} The bound is the size bound on the resummed datum in clause (c) of the resummation proposition. It enters only to guarantee that the datum exists. No inequality of the present proof uses it.

\emph{Joint invariance.} The identity \eqref{4.23:eq-density-invariance-1} is clause (ii) of the gauge-invariance theorem for the density. It holds term by term. The invariance of each term is an input, taken by statement as follows.
\begin{itemize}
\item For the sourced pieces, clause (iv) of the source statement asserts that each sourced piece $F_{\mathsf{y}}$ of Theorem~\ref{4.23:thm-oscillation-theorem} has the invariance of $F$ recorded in \cite[(2.41)(iii)]{B-CMP119}. This is used together with the invariance clause of the per-filing-domain source statement.
\item The pieces of the first, second, third and fifth families and the boundary pieces are invariant under complex gauge transformations, that is, under $G^{c}$-valued maps. This is the opening statement of the presentation lemma. Each of these pieces is therefore a function of the orbit of its argument. Consequently neither the choice of presentation nor the auxiliary data attached to it enter them.
\item The pieces $\mathfrak{Q}(\mathfrak{a})$ of Theorem~\ref{4.23:thm-oscillation-theorem} are jointly invariant under the gauge group. This is clause (i) of the first statement of the block-factor lemma.
\item For the pieces $\mathfrak{K}(Y')$ of Theorem~\ref{4.23:thm-oscillation-theorem} we use two inputs, each applied at the earlier index $\mu'$ of the induction of Theorem~\ref{4.23:thm-oscillation-theorem}: clause (g) at $\mu'$, and the invariance clause of the second kept member at $\mu'$. The latter is the joint $G$-invariance
\begin{equation*}
\mathfrak{q}_C(\zeta^{g};\varpi^{g})=\mathfrak{q}_C(\zeta;\varpi)
\end{equation*}
established in the construction of the second kept member $\mathfrak{q}_C$.
\item For the $z$-free boundary members $\mathfrak{v}^{\mathbf{C}}(Y_4)$, $Y_4\subset C^{\sharp}$, the invariance is the hypothesis clause of Definition~\ref{4.23:def-boundary-terms}, printed without proof in \cite[(1.11), (1.23)--(1.48)]{B-CMP122b}; cited as printed. In particular, the invariance is not derived from the remark that these members are Wilson actions and functions of orbits. Among them is the member $I(U_0)$, an integral at a fixed gauge; of it, \cite[(1.9)]{B-CMP122b} asserts only that the relation displayed there holds for $U_0$ in an arbitrary gauge.
\end{itemize}

\emph{The extension conventions preserve the invariance.} It remains to show that the two declared extensions keep \eqref{4.23:eq-density-invariance-1}. These are the pointwise extension convention of Definition~\ref{4.23:def-extension-conventions} and its reading on the subtracted copies, Definition~\ref{4.23:def-copies-reading}.

Let $v$ be a smooth core, at level $j$, and let $p=p_v$ and $Q_v$ be as fixed in Definition~\ref{4.23:def-copies-reading}. Write $\operatorname{ext}_p(\mathcal{U})$ for the extension at level $p$ over $Q_v$ of a configuration $\mathcal{U}$, formed there from the level-$p$ averages of $\mathcal{U}$ on $Q_v$. Write $\check g$ for the block-constant extension of $g$ from the base points of the level-$p$ blocks of $Q_v$. We claim that for every $G$-valued $g$
\begin{equation}\label{4.23:eq-density-invariance-2}
\operatorname{ext}_p(\mathcal{U}^{g})=\operatorname{ext}_p(\mathcal{U})^{\check g}.
\end{equation}

Two facts give this.
\begin{itemize}
\item The averaging operations preserve gauge transformations, by \cite[(11)]{B-CMP98}. Hence the level-$p$ averages of $\mathcal{U}^{g}$ on $Q_v$ are those of $\mathcal{U}$, transformed by $g$ at the base points.
\item Write $U(\mathfrak{B};\mathsf{D})$ for Ba{\l}aban's minimiser on the set $\mathfrak{B}$ at the data $\mathsf{D}$. For every $G$-valued gauge transformation $\lambda$, \cite[(181)]{B-CMP102b} gives
\begin{equation*}
U(\mathfrak{B};\mathsf{D}^{[\lambda]})=U(\mathfrak{B};\mathsf{D})^{\tilde\lambda},
\end{equation*}
where $\mathsf{D}^{[\lambda]}$ denotes the data transformed by $\lambda$ and $\tilde\lambda$ the induced extension of $\lambda$. Both sides are analytic in the data, so the identity extends to the complex data at which the extension is evaluated.
\end{itemize}
Applying the second fact to the averages supplied by the first yields \eqref{4.23:eq-density-invariance-2}. The same identity is the first step of the proof of the invariance statement for the presented pieces.

Next, the quantity $\mathfrak{i}^{(j)}_y$ on which a smooth core depends is invariant under $G^{c}$-valued gauge transformations, by the opening statement of the presentation lemma used above. For a copy core, $f_v(\zeta)$ is $\mathfrak{i}^{(j)}_y$ evaluated at the extension over $Q_v$. Combining this invariance with \eqref{4.23:eq-density-invariance-2} gives
\begin{equation*}
f_v(\zeta^{g})=f_v(\zeta).
\end{equation*}
Thus $f_v$ is still a function of the orbit of its slot on $Q_v$.

For an interior core, the gauge covariance of the chain acts on $Q_v$ as it acts on the core's own region $X_v$, so the interior core is again a function of the orbit of its slot on $Q_v$.

Hence neither the pointwise extension convention nor its reading on the subtracted copies breaks the invariance. Together with the term-by-term invariance above, this proves \eqref{4.23:eq-density-invariance-1}.
\end{proof}

\begin{lemma}[Oscillation theorem: degree and constant]\label{4.23:prf-degree-constant}
For a level $k$, write $N=N(g_k)$ for the number of levels of the window and $h=k-N$ for its
bottom level, and put $R:=\check{R}_k$. Here $\mathsf{T}_k=\log x_k$ and
$\check{\mathsf{T}}_k$ is its checked form, $r_{\mathrm{pr}}$ is the exponent of the
window radius, and $\varsigma_s$ ($s\ge 0$) are the anchor offsets. The remaining constants are those
of the bounds collected in \eqref{4.23:eq-form-theorem-a}, \eqref{4.23:eq-rim-lemma-a},
\eqref{4.23:eq-plain-cell-sum-a}, \eqref{4.23:eq-wrapped-cell-sum-a},
\eqref{4.23:eq-copies-lemma-a}, \eqref{4.23:eq-nested-regions-a} and
\eqref{4.23:eq-right-member-a}, and of the shallow-shell count of
Lemma~\ref{4.23:lem-shell-count-bound}; among them $c_{\mathrm{tr}}$ and $B_\sharp$ are the
constants of the bound \eqref{4.23:eq-form-theorem-a}. In addition, $C_N$ denotes the
window-count constant, $B_0$ the amplitude constant of the right members of the format, and
$K(d)$ a constant depending on $d$ only; these three enter through hypotheses (b) and (d) below.
Define the constants, the degrees and the quantity $\mathsf{A}(k)$ by
\begin{align}
&c_N:=C_N2^{r_{\mathrm{pr}}},\qquad c_R:=3^{r_{\mathrm{pr}}}L,\notag\\
&c_{\mathsf{b}}:=2K(d)B_0+b_\delta+c_NC_0^{\sharp}+(c_N+2)C_NC_0^{\sharp}3^{r_{\mathrm{pr}}},
\qquad c_{\mathsf{b}}':=c_{\mathsf{b}}+(c_N+2)b_\delta,\notag\\
&c^{\Sigma}_1:=2\cdot 101^{d}C_{\mathrm{fib}}\bigl(3+\mathsf{K}^0+c_NC_0^{\sharp}\bigr),\notag\\
&c^{\Sigma}_2:=e^5C^{+}C^{\mathrm{pl}}_{\mathsf{a}}C_{\mathrm{cell}}\,2d\,(104Lc_R+2)^{d-1}
(2+c_N)\,c_{\mathsf{b}}',\notag\\
&c^{\Sigma}_3:=2e^4(c_N+1)(110Lc_R)^d\Bigl[(c_N+1)\bigl(1+C_{\circledast}\mathsf{K}^{\flat}
(2LC_{\uparrow}c_R)^3\bigr)\notag\\
&\hphantom{c^{\Sigma}_3:=2e^4(c_N+1)(110Lc_R)^d\Bigl[}
+C^{<}+2C_{\mathsf{b}}(1+c_N)^{p_{\mathsf{b}}}(Lc_R)^d\Bigr]c_{\mathsf{b}}',\notag\\
&c^{\Sigma}_4:=e^5C_{3F}C_{\mathsf{a}}(1)C_{\mathrm{cell}}\,2d\,(104Lc_R+2)^{d-1}
(2+c_N)\,c_{\mathsf{b}},\notag\\
&c^{\Sigma}_5:=2e^4\bigl[C^{\mathrm{sh}}+2^d\bigl(2(104c_R)^d+3^dc_N\bigr)\bigr]c_{\mathsf{b}},
\qquad c^{\Sigma}_6:=2\cdot 100^d\mathsf{K}^{\mathrm{cp}},\qquad
c^{\Sigma}_7:=2(1+c_IM^5),\notag\\
&c^{\Sigma}_8:=30c_{\mathrm{tr}}B_\sharp^2A_1^2\cdot 4^{p_1}M^d(c_N+1)(110Lc_R)^d,\notag\\
&p^{\Sigma}_1:=d+1,\quad p^{\Sigma}_2:=d+2,\quad p^{\Sigma}_3:=2d+4+r_{\mathrm{pr}},\quad
p^{\Sigma}_4:=p^{\Sigma}_5:=d+2,\notag\\
&p^{\Sigma}_6:=d,\quad p^{\Sigma}_7:=0,\quad
p^{\Sigma}_8:=d+1+2p_1/r_{\mathrm{pr}},\notag\\
&p_\Sigma:=2d+6+r_{\mathrm{pr}}+\lceil 2p_1/r_{\mathrm{pr}}\rceil,\qquad
C_\Sigma:=2\sum_{i=1}^{8}c^{\Sigma}_i,\notag\\
&\mathsf{A}(k):=\Sigma^{V}(k)-\mathrm{nest}^*(k)+2\mathfrak{q}_{\max}(k),\notag\\
\intertext{and consider the two conditions on the parameters}
&\varsigma_{N+1}\le 2\check{\mathsf{T}}_k-1,\qquad
C_{3F}e^4\,C_{\mathrm{cell}}\,2d\,(104Lc_R+2)^{d-1}\,\varepsilon_\gamma\le 1 .
\label{4.23:eq-degree-constant-a}
\end{align}
Assume that, at every level $k$, the following hold.
\begin{enumerate}
\item[(a)] $\check{R}_k\ge 38$ (the floor on the window radius).
\item[(b)] $N\le c_N\check{R}_k$, $\check{R}_k\ge\check{\mathsf{T}}_k^{\,r_{\mathrm{pr}}}$ and
$\mathsf{T}_k\le 2\check{\mathsf{T}}_k$; in particular
\begin{equation*}
\mathsf{T}_k\le 2\check{\mathsf{T}}_k\le 2\check{R}_k^{1/r_{\mathrm{pr}}}
\end{equation*}
(the bound on the window length and the lower bound on the window radius by the checked
logarithm, from the geometry and counting statements of the window).
\item[(c)] The anchor offsets obey the first of the two conditions recorded at the end of
\eqref{4.23:eq-degree-constant-a}, and $\check{R}_{h+1}\le c_R\check{R}_k$ (the conclusion of
the radius comparison on the scaffold, obtained under that condition).
\item[(d)] The closed form of the sequence $\mathsf{b}_k$ of the format statements: for
$s\le N+1$,
\begin{equation*}
\mathsf{b}_k(s)\le c_{\mathsf{b}}\check{R}_k^{2},
\end{equation*}
with $c_{\mathsf{b}}$ as in \eqref{4.23:eq-degree-constant-a}. This is the closed form of
$\mathsf{b}_k$ obtained from the constant $\mathsf{b}^{0}=2K(d)B_0+b_\delta+NC_0^{\sharp}$
entering the definition of $\mathsf{b}_k$ (the base constant of the weight sequence of the
format statements), in which $b_\delta$ is a numerical constant of the format statements, and
from the growth bound
\begin{equation*}
\mathsf{b}_k(N+u)\le C_{\mathsf{b}}(1+u)^{p_{\mathsf{b}}}\mathsf{b}_k(N),\qquad
p_{\mathsf{b}}=1+r_{\mathrm{pr}}.
\end{equation*}
\item[(e)] The second of the two conditions recorded at the end of
\eqref{4.23:eq-degree-constant-a} holds.
\end{enumerate}
Then, at every level $k$,
\begin{equation*}
\mathsf{A}(k)\le\tfrac12C_\Sigma\check{R}_k^{\,p_\Sigma},\qquad
\mathrm{nest}^*(k)\le\tfrac12C_\Sigma\check{R}_k^{\,p_\Sigma},
\end{equation*}
and therefore the quantity $\Sigma^{\mathrm{tot}}(k)=\mathsf{A}(k)+\mathrm{nest}^*(k)$ bounded in
\textup{(t)} of \eqref{4.23:eq-oscillation-theorem-a} obeys
\begin{equation}\label{4.23:eq-degree-constant-1}
\Sigma^{\mathrm{tot}}(k)\le C_\Sigma\check{R}_k^{\,p_\Sigma}
\end{equation}
at every level $k$. The constant $C_\Sigma$ is a number depending only on the parameters listed
in clause (i) of \eqref{4.23:eq-constant-order-a}, among them $c_I$, which is chosen after $M$;
it involves neither the bound $\mathfrak{n}^{\mathrm{kept}}$ on the kept members of the
expansion nor $\mathrm{nest}^*(k)$.
\end{lemma}

\begin{proof}
The argument is direct: it gives \eqref{4.23:eq-degree-constant-1} at each $k$ separately,
with no induction in $k$, no property of the field variables, and no use of \textup{(v)} of
\eqref{4.23:eq-oscillation-theorem-a}.

\emph{Step 1: the conversions.} By (a), $R=\check{R}_k\ge 38\ge 1$; this is used throughout
to raise a power of $R$ to any larger power. By (b), $N\le c_NR$ with $c_N=C_N2^{r_{\mathrm{pr}}}$,
and $\mathsf{T}_k\le 2R^{1/r_{\mathrm{pr}}}$. By (c), $\check{R}_{h+1}\le c_RR$ with
$c_R=3^{r_{\mathrm{pr}}}L$. We recall how the constant $c_R$ arises in that comparison; the
recollection is not used below. The condition
$\varsigma_{N+1}\le 2\check{\mathsf{T}}_k-1$ gives, for $s\le N+1$,
\begin{equation*}
(\mathsf{T}_k+\varsigma_s)^{r_{\mathrm{pr}}}\le(3\check{\mathsf{T}}_k)^{r_{\mathrm{pr}}}
\le 3^{r_{\mathrm{pr}}}R,
\end{equation*}
the last inequality by $\check{R}_k\ge\check{\mathsf{T}}_k^{\,r_{\mathrm{pr}}}$; and the upper
radius function $\mathsf{R}^{\uparrow}$ satisfies
\begin{equation*}
\mathsf{R}^{\uparrow}(\check{\mathsf{T}}_k)\le L(3\check{\mathsf{T}}_k)^{r_{\mathrm{pr}}}
\le c_R\check{R}_k .
\end{equation*}

By (d), $\mathsf{b}_k(s)\le c_{\mathsf{b}}R^2$ for $s\le N+1$. The sequence
$\mathsf{b}(s)=\mathsf{b}_k(s)+(s+1)b_\delta$ of \eqref{4.23:eq-plain-cell-sum-a} then obeys,
for $s\le N+1$,
\begin{equation*}
\mathsf{b}(s)\le c_{\mathsf{b}}R^2+(N+2)b_\delta\le c_{\mathsf{b}}'R^2,
\end{equation*}
because
\begin{equation*}
N+2\le c_NR+2\le(c_N+2)R\le(c_N+2)R^2
\end{equation*}
by $R\ge1$, and
$c_{\mathsf{b}}'=c_{\mathsf{b}}+(c_N+2)b_\delta$. In particular $\mathsf{b}(N)$ and
$\mathsf{b}(N+1)$ are at most $c_{\mathsf{b}}'R^2$, and $\mathsf{b}_k(N)$, $\mathsf{b}_k(N+1)$ are
at most $c_{\mathsf{b}}R^2$. Finally, since $\mathfrak{b}'_k=A_1\mathsf{T}_k^{p_1}$,
\begin{equation*}
\mathfrak{b}_k'^{\,2}=A_1^2\mathsf{T}_k^{2p_1}\le A_1^2\cdot 4^{p_1}R^{2p_1/r_{\mathrm{pr}}}.
\end{equation*}

\emph{Step 2: the eight members.} By its definition \eqref{4.23:eq-right-member-a}, the
right member $\Sigma^V(k)$ is the sum of seven members and of $\mathrm{nest}^*(k)$:
the rim member, the deep and the shallow plain members, the deep and the shallow wrapped
members, the copies member and the member $2(1+c_IM^5)$ of the boundary terms
$\mathbf{C}_k$. Hence $\mathsf{A}(k)$ is the sum of these seven members and of
$2\mathfrak{q}_{\max}(k)$. In each of the eight expressions we replace $N$,
$\check{R}_{h+1}$, $\mathsf{T}_k$ by the bounds $c_NR$, $c_RR$, $2R^{1/r_{\mathrm{pr}}}$ of
Step~1 and $\mathsf{b}$, $\mathsf{b}_k$ by the bounds of Step~1, and then raise every sub-term
to the highest power of $R$ that occurs in the expression, using $R\ge 1$. This gives the
following eight bounds, the exponent of $R$ in the $i$-th bound being the degree
$p^{\Sigma}_i$ of the member recorded in \eqref{4.23:eq-degree-constant-a}.
\begin{align*}
2\mathsf{n}_\sharp(k)C_{\mathrm{fib}}\bigl[3+\mathsf{K}^0+NC_0^{\sharp}\bigr]
&\le c^{\Sigma}_1R^{d+1},\\
e^4C^{+}\mathsf{S}^{\mathrm{dp}}_{\mathrm{pl}}(k)&\le c^{\Sigma}_2R^{d+2},\\
2e^4\mathsf{S}^{\mathrm{sh}}_{\mathrm{pl}}(k)&\le c^{\Sigma}_3R^{2d+4+r_{\mathrm{pr}}},\\
e^4C_{3F}\mathsf{S}^{\mathrm{dp}}_{\mathrm{wr}}(k)&\le c^{\Sigma}_4R^{d+2},\\
2e^4\mathsf{S}^{\mathrm{sh}}_{\mathrm{wr}}(k)&\le c^{\Sigma}_5R^{d+2},\\
2(100\check{R}_k)^d\mathsf{K}^{\mathrm{cp}}&=c^{\Sigma}_6R^{d},\\
2(1+c_IM^5)&=c^{\Sigma}_7R^{0},\\
2\mathfrak{q}_{\max}(k)&\le c^{\Sigma}_8R^{d+1+2p_1/r_{\mathrm{pr}}}.
\end{align*}
We go through them in order.

The rim member: by \eqref{4.23:eq-rim-lemma-a}, $\mathsf{n}_\sharp(k)=(101R)^d$, and
$3+\mathsf{K}^0+NC_0^{\sharp}\le(3+\mathsf{K}^0+c_NC_0^{\sharp})R$; the degree is $d+1$.

The deep plain member: in $\mathsf{S}^{\mathrm{dp}}_{\mathrm{pl}}(k)$ of
\eqref{4.23:eq-plain-cell-sum-a} the factors are bounded by
$(104L\check{R}_{h+1}+2)^{d-1}\le(104Lc_R+2)^{d-1}R^{d-1}$, $2+N\le(2+c_N)R$ and
$\mathsf{b}(N+1)\le c_{\mathsf{b}}'R^2$; together with the factor $e$ of
$\mathsf{S}^{\mathrm{dp}}_{\mathrm{pl}}(k)$ and the prefactor $e^4C^{+}$ this gives
$c^{\Sigma}_2$ and degree $(d-1)+1+2=d+2$.

The shallow plain member: in $\mathsf{S}^{\mathrm{sh}}_{\mathrm{pl}}(k)$ of
\eqref{4.23:eq-plain-cell-sum-a}, $N+1\le(c_N+1)R$ and
$(110L\check{R}_{h+1})^d\le(110Lc_R)^dR^d$. Inside the bracket the first sub-term is at most
$(c_N+1)(1+C_{\circledast}\mathsf{K}^{\flat}(2LC_{\uparrow}c_R)^3)R^4$, the second is the
constant $C^{<}$, and the third is at most
$2C_{\mathsf{b}}(1+c_N)^{p_{\mathsf{b}}}(Lc_R)^dR^{p_{\mathsf{b}}+d}$, the factor
$(1+N)^{p_{\mathsf{b}}}$ being the one in the shallow plain bound as it is stated in
\eqref{4.23:eq-plain-cell-sum-a}. Since $d=4$ and $p_{\mathsf{b}}\ge 1$, the highest of
these powers is $R^{p_{\mathsf{b}}+d}$, and the bracket is at most the bracket of
$c^{\Sigma}_3$ times $R^{p_{\mathsf{b}}+d}$. Finally $\mathsf{b}(N)\le c_{\mathsf{b}}'R^2$.
With the prefactor $2e^4$ the degree is
\begin{equation*}
1+d+(p_{\mathsf{b}}+d)+2=2d+4+r_{\mathrm{pr}},\qquad p_{\mathsf{b}}=1+r_{\mathrm{pr}}.
\end{equation*}

The deep wrapped member: $\mathsf{S}^{\mathrm{dp}}_{\mathrm{wr}}(k)$ of
\eqref{4.23:eq-wrapped-cell-sum-a} has the same geometric factors as the deep plain member,
with $C_{\mathsf{a}}(1)$ in place of $C^{\mathrm{pl}}_{\mathsf{a}}$ and
$\mathsf{b}_k(N+1)\le c_{\mathsf{b}}R^2$ in place of $\mathsf{b}(N+1)$; with the prefactor
$e^4C_{3F}$ this gives $c^{\Sigma}_4$ and degree $d+2$.

The shallow wrapped member: in $\mathsf{S}^{\mathrm{sh}}_{\mathrm{wr}}(k)$ of
\eqref{4.23:eq-wrapped-cell-sum-a},
\begin{equation*}
C^{\mathrm{sh}}+2^d\bigl(2(104\check{R}_{h+1})^d+3^dN\bigr)
\le\bigl[C^{\mathrm{sh}}+2^d\bigl(2(104c_R)^d+3^dc_N\bigr)\bigr]R^d,
\end{equation*}
and $\mathsf{b}_k(N)\le c_{\mathsf{b}}R^2$; with the prefactor $2e^4$ this gives
$c^{\Sigma}_5$ and degree $d+2$.

The copies member is $2(100R)^d\mathsf{K}^{\mathrm{cp}}$, with $\mathsf{K}^{\mathrm{cp}}$ of
\eqref{4.23:eq-copies-lemma-a}; its degree is $d$. The member of $\mathbf{C}_k$ is the
constant $2(1+c_IM^5)$, of degree $0$.

The member $2\mathfrak{q}_{\max}(k)$: the bound $\sup_{\mathfrak{S}}|\mathfrak{q}_C|\le
\mathfrak{q}_{\max}(k)$ on the second kept member is an input of the oscillation theorem
(Theorem~\ref{4.23:thm-oscillation-theorem}); it is the bound of
\eqref{4.23:eq-form-theorem-a} at the linearised localised minimiser map
$\mathsf{T}^{\mathsf{L}}$, restricted to the tube $\mathfrak{S}$, with the shallow-shell count
of Lemma~\ref{4.23:lem-shell-count-bound}. In the expression of $2\mathfrak{q}_{\max}(k)$
so obtained from \eqref{4.23:eq-form-theorem-a} and Lemma~\ref{4.23:lem-shell-count-bound},
the same substitutions as above, $N+1\le(c_N+1)R$ and
$(110L\check{R}_{h+1})^d\le(110Lc_R)^dR^d$, together with the bound on
$\mathfrak{b}_k'^{\,2}$ of Step~1, give $c^{\Sigma}_8$ and degree
$d+1+2p_1/r_{\mathrm{pr}}$.

\emph{Step 3: the common degree.} The eight degrees $p^{\Sigma}_1,\dots,p^{\Sigma}_8$ of
\eqref{4.23:eq-degree-constant-a} are $d+1$, $d+2$,
$2d+4+r_{\mathrm{pr}}$, $d+2$, $d+2$, $d$, $0$ and $d+1+2p_1/r_{\mathrm{pr}}$. Each is at
most
\begin{equation*}
p_\Sigma=2d+6+r_{\mathrm{pr}}+\lceil 2p_1/r_{\mathrm{pr}}\rceil ,
\end{equation*}
the last one because $d+1\le 2d+6+r_{\mathrm{pr}}$ and
$2p_1/r_{\mathrm{pr}}\le\lceil 2p_1/r_{\mathrm{pr}}\rceil$. The exponent $p_\Sigma$ is used
later only inside a supremum, over the depth, of a polynomial of degree $p_\Sigma$ against a
decaying exponential, which is finite for every value of $p_\Sigma$; the size of $p_\Sigma$
therefore imposes no condition on any other parameter.

\emph{Step 4: the bound on $\mathsf{A}(k)$.} Since $R\ge1$, each of the eight bounds of Step~2
is at most $c^{\Sigma}_iR^{p_\Sigma}$, and adding them,
\begin{equation*}
\mathsf{A}(k)\le\Bigl(\sum_{i=1}^{8}c^{\Sigma}_i\Bigr)R^{p_\Sigma}
=\tfrac12C_\Sigma R^{p_\Sigma}.
\end{equation*}
The coefficients $c^{\Sigma}_i$, and hence $C_\Sigma$, are numbers formed from the parameters
listed in clause (i) of \eqref{4.23:eq-constant-order-a}, among them $c_I$, which is chosen
after $M$; none of them contains the bound $\mathfrak{n}^{\mathrm{kept}}$ on the kept members
or $\mathrm{nest}^*(k)$.

\emph{Step 5: the bound on $\mathrm{nest}^*(k)$.} By \eqref{4.23:eq-nested-regions-b} of the
nested-region lemma (Lemma~\ref{4.23:lem-nested-regions}),
\begin{equation*}
\mathrm{nest}^*(k)=\tfrac12C_{3F}e^4\varepsilon_\gamma C_{\mathrm{cell}}\,
\mathsf{n}_\partial(k)\,C_\Sigma ,
\end{equation*}
an explicit quantity. By the substitution $\check{R}_{h+1}\le c_RR$ and $R\ge 1$, as in
Step~2, the boundary count $\mathsf{n}_\partial(k)$ of \eqref{4.23:eq-plain-cell-sum-a} obeys
$\mathsf{n}_\partial(k)\le 2d(104Lc_R+2)^{d-1}R^{d-1}$. Hence, by hypothesis (e),
\begin{equation*}
\mathrm{nest}^*(k)\le\tfrac12C_\Sigma\bigl[C_{3F}e^4C_{\mathrm{cell}}\,2d\,
(104Lc_R+2)^{d-1}\varepsilon_\gamma\bigr]R^{d-1}
\le\tfrac12C_\Sigma R^{d-1}\le\tfrac12C_\Sigma R^{p_\Sigma},
\end{equation*}
the last step by $d-1\le p_\Sigma$ and $R\ge 1$. The constant $C_\Sigma$ stands on both sides
of the bound $\mathrm{nest}^*(k)\le\frac12C_\Sigma R^{d-1}$ and cancels; this is why the
condition of hypothesis (e) contains no $C_\Sigma$, and why fixing $C_\Sigma$ by
\eqref{4.23:eq-degree-constant-a} before $\mathrm{nest}^*(k)$ is not circular.

\emph{Step 6: conclusion.} By Steps~4 and~5,
\begin{equation*}
\Sigma^{\mathrm{tot}}(k)=\mathsf{A}(k)+\mathrm{nest}^*(k)
\le\tfrac12C_\Sigma\check{R}_k^{\,p_\Sigma}+\tfrac12C_\Sigma\check{R}_k^{\,p_\Sigma}
=C_\Sigma\check{R}_k^{\,p_\Sigma},
\end{equation*}
which is \eqref{4.23:eq-degree-constant-1}, that is, \textup{(t)} of
\eqref{4.23:eq-oscillation-theorem-a}, at the level $k$. Since $k$ was arbitrary, it holds at
every level.
\end{proof}

\begin{lemma}[Oscillation theorem: depth decay and independence]\label{4.23:prf-depth-decay}
Under the hypotheses of Theorem~\ref{4.23:thm-oscillation-theorem}, clauses (b), (d) and (n)
of the oscillation theorem for the resummed density, as displayed in
\eqref{4.23:eq-oscillation-theorem-a}, hold.
\end{lemma}

\begin{proof}
Throughout, $\hat D_v$ denotes the depth of a unit $v$ as it is defined in
\eqref{4.23:eq-unit-oscillation-a} and enters the deep column of the table of the
total-oscillation lemma
(Lemma~\ref{4.23:lem-unit-oscillation}, whose proof is incomplete at one step), and
$\delta_P=\delta_0/8$. Let $\varsigma\ge0$ be the depth parameter of clause (b) of
\eqref{4.23:eq-oscillation-theorem-a}.

\emph{Clause (b).} Let $v$ be a deep unit among the constituents (m1), (m2) of
\eqref{4.23:eq-component-constituents-a}, with $\hat D_v\ge\varsigma$. The deep column of the
table in \eqref{4.23:eq-unit-oscillation-a} carries the factor
$e^{-\frac12\delta_P\hat D_v}$.
The cell sum for plain units (Lemma~\ref{4.23:lem-plain-cell-sum}) and the cell sum for
wrapped units (Lemma~\ref{4.23:lem-wrapped-cell-sum}) use only the factor
$e^{-\frac14\delta_P\hat D_v}$ of it. Since $\hat D_v\ge\varsigma$, we have
$\tfrac12\hat D_v\ge\tfrac14\varsigma+\tfrac14\hat D_v$. Moreover
$e^{-\frac14\delta_P\varsigma}\le e^{-\frac18\delta_P\varsigma}$ for $\varsigma\ge0$.
Together these give
\begin{equation*}
e^{-\frac12\delta_P\hat D_v}
\le e^{-\frac14\delta_P\varsigma}\,e^{-\frac14\delta_P\hat D_v}
\le e^{-\frac18\delta_P\varsigma}\,e^{-\frac14\delta_P\hat D_v}.
\end{equation*}
The factor $e^{-\frac14\delta_P\hat D_v}$ on the right is the one the two cell sums
consume. We now run the summations of Lemma~\ref{4.23:lem-plain-cell-sum} and
Lemma~\ref{4.23:lem-wrapped-cell-sum} over the deep units with $\hat D_v\ge\varsigma$ only.
The contribution of these units to the plain member
$e^4[C^+\mathsf{S}^{dp}_{pl}(k)+2\mathsf{S}^{sh}_{pl}(k)]$ is then at most
$e^{-\frac18\delta_P\varsigma}$ times that member. The same holds for their contribution to
the wrapped member $e^4[C_{3F}\mathsf{S}^{dp}_{wr}(k)+2\mathsf{S}^{sh}_{wr}(k)]$ of
$\Sigma^V(k)$ in \eqref{4.23:eq-right-member-a}.

For the Gaussian families carried inside the plain units, the entry of the deep column
that is used is the deep entry of the third row of the table in
\eqref{4.23:eq-unit-oscillation-a}, which carries the factor
$e^{-\frac12\delta_P\hat D_v}$. That entry is reached through the display
\eqref{4.23:eq-unit-oscillation-b} of Lemma~\ref{4.23:lem-unit-oscillation}, and so depends
on the step of the total-oscillation
lemma at which its proof is incomplete. With this entry the
computation above applies verbatim, and the share of these families in the plain member is
again at most $e^{-\frac18\delta_P\varsigma}$ times that member.

For the constituents (m3), the units filed at the nested regions of $C$
(Lemma~\ref{4.23:lem-nested-regions}), the decay in $\varsigma$ does not come from the table;
it comes from the third clause of that lemma,
which bounds the contribution of the nested regions $(Y',m')$ with
$d^{\wedge}(Y'^{\sharp},\mathsf{C}^c)\ge\varsigma$; this is clause (b) for the constituents
(m3).

\emph{Clause (d).} We read off the right member $\Sigma^V(k)$ of
\eqref{4.23:eq-right-member-a} together with the cell-sum constants
\eqref{4.23:eq-plain-cell-sum-a} and \eqref{4.23:eq-wrapped-cell-sum-a}. It is built from
the following quantities only:
\begin{itemize}
\item $N=N(g_k)$, $\check{R}_{h+1}$ and $\check{R}_k$, where $h=k-N$;
\item $\mathsf{T}_k$, which enters through the function $\mathsf{b}(\cdot)$ of
\eqref{4.23:eq-plain-cell-sum-a} and through $\mathfrak{b}'_k$;
\item constants depending on the fixed parameters only (those displayed in
\eqref{4.23:eq-right-member-a}, \eqref{4.23:eq-plain-cell-sum-a} and
\eqref{4.23:eq-wrapped-cell-sum-a}).
\end{itemize}
The core enters only through the count $\#\partial^+$ of
\eqref{4.23:eq-plain-cell-sum-a}, and only via
the bound $\#\partial^+\le\mathsf{n}_\partial(k)$.

No quantity other than those just listed occurs in $\Sigma^V(k)$. In particular, none of the
following quantities, each as written in clause (d) of
\eqref{4.23:eq-oscillation-theorem-a}, occurs:
\begin{itemize}
\item the sum $\sum_j|\Gamma^0_j\cap C|$;
\item the size $|\mathsf{C}|$ of the core;
\item the quantity $D$;
\item the age of any large-field region;
\item $N_0$;
\item any datum of a nested region;
\item any datum of the arrest.
\end{itemize}
The cutoff index $K$ does not occur in $\Sigma^V(k)$, and the level $k$ enters only through
the quantities of the first list above.

\emph{Clause (n).} Consider three kinds of member of $\Sigma^V(k)$:
\begin{itemize}
\item the rim member, which is the first summand;
\item the shallow members, which are the terms carrying $\mathsf{S}^{sh}_{pl}(k)$ and
$\mathsf{S}^{sh}_{wr}(k)$;
\item the copy member, which is the summand carrying $(100\check{R}_k)^d$.
\end{itemize}
Each of these is a product of counts with the remaining factors displayed in
\eqref{4.23:eq-right-member-a}, \eqref{4.23:eq-plain-cell-sum-a} and
\eqref{4.23:eq-wrapped-cell-sum-a}, among them $N$, $\mathsf{b}(N)$ and $\mathsf{b}_k(N)$.
The counts grow with $\check{R}_k$, and so do these members; hence they
grow as $g_k$ decreases to $0$. Clause (n) asserts nothing about
$\Sigma^V(k)$ tending to zero, and the proof gives nothing of the kind.

Clause (o) of \eqref{4.23:eq-oscillation-theorem-a} is not part of this lemma; it is
proved as a separate lemma in the proof of
Theorem~\ref{4.23:thm-oscillation-theorem}.
\end{proof}

\begin{remark}[On the mechanism of the bound]
The estimate
\cite[(1.69)]{B-CMP122b} is a bound on sizes and therefore sums over every old cell. An
oscillation, by contrast, is paid for cell by cell, at a cost exponentially small in the
own-scale contour distance $d^{\wedge}$ of the cell from the boundary of the core.
The counting clause of the cell lemma controls the number of cells met for each cell
crossed. The depth bounds on piles of anchored units control the growth of the piles
with depth. The nested regions enter only through their oscillations. These are bounded
afresh at each level by \eqref{4.23:eq-right-member-a}, behind the nested regions' own
collars, so the constant does not compound from level to level.
\end{remark}

A kept member $v'\in\mathfrak{K}(Y')$ of a nested region $Y'$ enters the bounds below through one number. The \emph{leaf count} is the number fixed in the corollary on kept members,
\begin{equation}\label{4.23:eq-constant-order-1}
\mathfrak{n}^{\mathrm{kept}}
=2^{d}\,C_{\Sigma}\sup_{\varrho\ge1}(2\varrho+2)^{p_{\Sigma}}e^{-\varrho}.
\end{equation}
Here $C_{\Sigma}$ and $p_{\Sigma}$ are the constant and the degree of \eqref{4.23:eq-degree-constant-a}.

The leaf count is derived from $C_{\Sigma}$. Later steps carry the kept members' share in the form $\mathfrak{n}^{\mathrm{kept}}$ times a far factor. The terms those steps create are stored, lie in later regions $C$, and are summed in $V^{\sharp}(C^{\sharp})$. If a constant of the explicit right member $\Sigma^{V}(k)$ of \eqref{4.23:eq-right-member-a} read $\mathfrak{n}^{\mathrm{kept}}$, then $C_{\Sigma}$ would be defined through itself. The following lemma excludes this, under two explicit hypotheses.

\begin{lemma}[the size-bound constant is fixed before the leaf count]\label{4.23:lem-constant-order}
Assume the following two statements.
\begin{enumerate}
\item[(a)] (The share clause of the corollary on kept members, in the statements that return the format.) Consider a term created from a piece of a kept member in one of two ways:
\begin{itemize}
\item at an earlier step of one of the kinds written (T) and (P) in the statements on kept members, or at a step of the kind written (L) there, taken for a foreign component;
\item at an earlier healing step $\mathbf{R}_{k_1}$ of the component itself, $k_1<k$, from an exterior piece of a kept member that ended there.
\end{itemize}
The statements whose conclusion returns the format after such terms are added return it with unchanged right members. The kept members' share enters only their closing conditions, in one of two forms. In the first form, a kept share $\varepsilon^{\mathrm{kept}}_i$ stands beside $\varepsilon^{\mathrm{rel}}_i$ in a closing condition whose right member is $\frac14\varepsilon_V^*$. In the second form, the amplitude constants $\mu$ and $\mathsf{M}_*$ of such a statement are bounded by the constants $\mu^{\flat}$, $\mathsf{M}^{\flat}$ of the statement before that share is added, plus the kept members' shares $\mu^{\mathrm{kept}}$, $\mathsf{M}^{\mathrm{kept}}$:
\begin{equation*}
\mu\le\mu^{\flat}+\mu^{\mathrm{kept}},\qquad
\mathsf{M}_*\le\mathsf{M}^{\flat}+\mathsf{M}^{\mathrm{kept}}.
\end{equation*}
In each statement of the second form the kept share equals $\mathfrak{n}^{\mathrm{kept}}$ times that statement's far factor, which is the factor by which the statement weights distant pieces.
\item[(b)] (The per-cube bound of the Gaussian step of the arrest theorem.) The per-cube datum $\mathfrak{N}^{G}(e;\mathfrak{O})$ of a Gaussian unit $e$ is the oscillation datum, over its own window, of the Gaussian step of the arrest theorem. This datum enters the per-cube line whose right member is $b_{\delta}^{G}$, through summands
\begin{equation*}
\mathfrak{N}^{G}(e;\mathfrak{O})\,e^{\mathsf{w}(e)}\,e^{r'd_{\ell}(A_e)}.
\end{equation*}
Here $\mathsf{w}(e)$ is the weight of $e$, $A_e$ is its localisation domain, and $r'$ is the decay rate of the per-cube line. The datum $\mathfrak{N}^{G}(e;\mathfrak{O})$ reads none of the following: the oscillation $\mathfrak{O}(v')\le C_{\Sigma}\check{R}^{p_{\Sigma}}$ of a kept member of any nested region; the quantity $\mathsf{S}^{\mathfrak{K}}$ of \eqref{4.23:eq-collar-payment-a}; the leaf count $\mathfrak{n}^{\mathrm{kept}}$.
\end{enumerate}
Then the following hold.
\begin{enumerate}
\item[(i)] $C_{\Sigma}$ is a function of the constants listed below and of no others. Each is chosen before $\gamma$, either in the order of choice of the auxiliary parameters of the correlation theorem or in the order of the statement introducing it. In parentheses after each group stand the constants through which it enters. These are constants of the degree table of \eqref{4.23:eq-degree-constant-a} and of the statements on which its entries depend:
\begin{itemize}
\item $d,L,r_{\mathrm{pr}},p_1,\lambda,c_2$ (through $c_N$, $c_R$, $C_{\uparrow}$, $C_{\mathsf{b}}$, $p_{\mathsf{b}}$);
\item $\mathsf{k}_d$ (through $C_{\mathrm{fib}}$ of the rim lemma);
\item $\kappa,\delta$ (through $C^{<}$, $C^{\mathrm{sh}}$, $C_{\circledast}$, $K(d)$);
\item $\kappa_1$ (through $B_{\sharp}$);
\item $M,M_1$ (through $C_{\mathrm{cell}}$ of the nested-region lemma, through $w=M/M_1$, and through $c_I M^5$, where $c_I$ is the numerical constant written $O(1)$ in Ba{\l}aban's estimate carrying the factor $M^5$, chosen after $M$);
\item $r^{0}$ (through $C_{3\mathrm{F}}$ of the nested-region lemma, $C^{\mathrm{pt}}$, $C^{+}$);
\item $E_0,C_B,C_I,b_{+},r_J$ and the constants of the statement introducing $\mathsf{K}^{\flat}$, $\mathsf{K}^{\mathrm{cp}}$, $\mathsf{K}^{0}$ (through these three);
\item $B_0,C_0^{\sharp},b_{\delta}$ (through $c_{\mathsf{b}}$, $c'_{\mathsf{b}}$, $C_{\mathsf{a}}(1)$, $C^{\mathrm{pl}}_{\mathsf{a}}$);
\item $A_1$ and $c_{\mathrm{tr}}$.
\end{itemize}
\item[(ii)] None of the constants in (i) reads $\mathfrak{n}^{\mathrm{kept}}$, $C_{\Sigma}$, $\gamma$, or a function of the level such as $N(g_k)$, $\check{R}_k$, $\mathsf{T}_k$.
\item[(iii)] The two closing conditions of \eqref{4.23:eq-degree-constant-a} read neither $C_{\Sigma}$ nor $\mathfrak{n}^{\mathrm{kept}}$. The first is the condition attached to the window constant $c_R$. The second is the nested-region smallness condition, which we restate here as
\begin{equation}\label{4.23:eq-constant-order-2}
C_{3\mathrm{F}}e^{4}\,C_{\mathrm{cell}}\,2d\,(104Lc_R+2)^{d-1}\,\varepsilon_{\gamma}\le1 .
\end{equation}
\item[(iv)] Consequently the order of choice
\begin{equation}\label{4.23:eq-constant-order-a}
\begin{gathered}
\cdots\to A_0/A_1\to\cdots\to E_0\to\cdots\to B_0\to\cdots\to c_A\\
\to\mathsf{K}^{0},\ \mathsf{K}^{0\prime},\ \mathsf{K}^{0}_{\mathrm{cell}},\ \alpha_*
\to C_{\Sigma},\ p_{\Sigma}\to\mathfrak{n}^{\mathrm{kept}}\to\gamma
\end{gathered}
\end{equation}
has no cycle. In it the correlation theorem places $E_0$, $B_0$ and $\mathsf{K}^{0}$ after $A_1$, and $C_{\Sigma}$ comes after all of them. The order remains free of cycles when the oscillation theorem (Theorem~\ref{4.23:thm-oscillation-theorem}) is inserted at its own row of the ordered list of conditions at whose rows its hypotheses are assumed.
\end{enumerate}
\end{lemma}

\begin{proof}
\emph{Proof of \textup{(i)}.} By \eqref{4.23:eq-degree-constant-a}, $C_{\Sigma}$ is twice the sum of the eight coefficients of the degree table. By the construction of that table, each of the eight coefficients is an expression in $c_N$, $c_R$, $r_{\mathrm{pr}}$, $A_1$, $p_1$ and in the constants standing in the entries of $\Sigma^{V}(k)$ in \eqref{4.23:eq-right-member-a}. Read against the statements that introduce those constants, the table shows that they are exactly the constants listed in (i). The constant $c_I$ is among them, and it is chosen after $M$.

\emph{Proof of \textup{(iii)}.} The same reading applies to the two closing conditions of \eqref{4.23:eq-degree-constant-a}. The condition \eqref{4.23:eq-constant-order-2} reads $C_{3\mathrm{F}}$, $C_{\mathrm{cell}}$, $d$, $L$ and $c_R$. Through $\varepsilon_{\gamma}$ of \eqref{4.23:eq-nested-regions-a}, it also reads $p_{\Sigma}$ and $\mathsf{T}_{\gamma}$.

It does not read $C_{\Sigma}$, for the following reason. By the nested-region lemma (Lemma~\ref{4.23:lem-nested-regions}),
\begin{equation*}
\mathrm{nest}^*(k)=\tfrac12C_{3\mathrm{F}}e^{4}\varepsilon_{\gamma}C_{\mathrm{cell}}\mathsf{n}_{\partial}(k)\,C_{\Sigma},
\qquad
\mathsf{n}_{\partial}(k)\le2d(104Lc_R+2)^{d-1}R^{d-1}.
\end{equation*}
Here $R\ge1$ is the variable of the degree table of \eqref{4.23:eq-degree-constant-a}.
The inequality $\mathrm{nest}^*(k)\le\frac12C_{\Sigma}R^{d-1}$ has the factor $C_{\Sigma}$ on both sides. Divide it out and use the bound on $\mathsf{n}_{\partial}(k)$ just displayed. The inequality then follows from \eqref{4.23:eq-constant-order-2}, in which $C_{\Sigma}$ does not appear.

Neither condition reads $\mathfrak{n}^{\mathrm{kept}}$. The second condition contains only the constants just listed, and $\mathfrak{n}^{\mathrm{kept}}$ is not among them. Inspection of the table of \eqref{4.23:eq-degree-constant-a} shows that the first condition, the one attached to $c_R$, contains neither $C_{\Sigma}$ nor $\mathfrak{n}^{\mathrm{kept}}$.

\emph{Proof of \textup{(ii)}.} The leaf count $\mathfrak{n}^{\mathrm{kept}}$ enters later steps only through the kept members' share. So a constant of the list in (i) could read $\mathfrak{n}^{\mathrm{kept}}$, $C_{\Sigma}$ or $\gamma$ only through a contribution of a kept member of a nested region to $\Sigma^{V}(k)$. Apart from the per-cube Gaussian datum treated in (5), no line of the proof of the oscillation theorem (Theorem~\ref{4.23:thm-oscillation-theorem}) uses the anchoring statement for kept members, the quantity $\mathsf{S}^{\mathfrak{K}}$ or the count $\mathfrak{n}^{\mathrm{kept}}(\beta)$ of \eqref{4.23:eq-collar-payment-a}. Such a contribution can therefore enter in exactly five ways; we treat them in turn.

(1) \emph{The kept member itself.} These are the value and the interior pieces of $v'\in\mathfrak{K}(Y')$ with $Y'\subset C$.
\begin{itemize}
\item They enter $\mathrm{nest}(k)$ of \eqref{4.23:eq-nested-regions-b}, which is bounded by $\mathrm{nest}^*(k)$ in the nested-region lemma (Lemma~\ref{4.23:lem-nested-regions}).
\item That proof uses two inputs, both taken from the lemma on the oscillation of kept members, together with the induction hypothesis. The first is the bound of the oscillation $\mathfrak{O}(v)$ by the two-point oscillation over the tube. The second is the bound by $\frac12C_{3\mathrm{F}}e^{-\frac12\delta_PD_X}\mathfrak{O}(v')$ on differences of the kept member's functional at two points of $\mathfrak{P}$ with the same $\varpi$; here $D_X$ is the separation parameter of that bound.
\item The bound $\mathrm{nest}^*(k)$ is the number $\frac12C_{3\mathrm{F}}e^{4}$ times $\varepsilon_{\gamma}\,C_{\mathrm{cell}}\,\mathsf{n}_{\partial}(k)\,C_{\Sigma}$.
\item $C_{\Sigma}$ is defined through $\mathsf{A}(k)$ of \eqref{4.23:eq-degree-constant-a}, before the induction. $\mathsf{A}(k)$ contains no term of the nested regions, because $\mathrm{nest}^*(k)$ is subtracted in it.
\item By (iii), \eqref{4.23:eq-constant-order-2} reads neither $C_{\Sigma}$ nor $\mathfrak{n}^{\mathrm{kept}}$.
\end{itemize}
There is also a bound of the deep nested regions in terms of $\mathfrak{n}^{\mathrm{kept}}$, of the form $e\,\mathfrak{n}^{\mathrm{kept}}C_{\mathrm{cell}}2d(\cdots)^{d-1}$. It is true, but it is not used, because it would place $C_{\Sigma}$ after $\mathfrak{n}^{\mathrm{kept}}$. The two suprema involved are distinct objects. The supremum in $\varepsilon_{\gamma}$ runs over $\varrho\ge\frac12R_{\gamma}$, with $R_{\gamma}$ as in \eqref{4.23:eq-nested-regions-a}, and the one in \eqref{4.23:eq-constant-order-1} runs over $\varrho\ge1$.

(2) \emph{Exterior pieces of kept members reaching the strip.} These are the pieces $Y_4\in\mathcal{Y}_{\mathrm{rim}}$.
\begin{itemize}
\item They lie inside the totals $\mathfrak{v}^{\mathrm{rel}}+\mathfrak{v}^{\mathrm{kept}}$. In the rim lemma (Lemma~\ref{4.23:lem-rim-lemma}) these totals are bounded by $1$.
\item The bound $1$ is the closing condition of the exterior-piece bound, in which $\|\cdot\|_{a,1}$ is the weighted norm of the totals:
\begin{equation*}
\|\mathfrak{v}^{\mathrm{rel}}+\mathfrak{v}^{\mathrm{kept}}\|_{a,1}\le\varepsilon^{\mathrm{rel}}_L(k)\le1 .
\end{equation*}
\item The ceiling in that condition reads $\mathfrak{n}^{\mathrm{kept}}$; the bound $1$ does not.
\item Likewise, the condition on the interior parameter $\varepsilon^{\mathrm{int}}$ among the hypotheses of the oscillation theorem (Theorem~\ref{4.23:thm-oscillation-theorem}) reads $C_{\Sigma}$, but its bound $\frac12$ does not.
\item The ceiling is met by choosing $\gamma$ after $\mathfrak{n}^{\mathrm{kept}}$, in accordance with \eqref{4.23:eq-constant-order-a}.
\end{itemize}
Hence the bound of the rim lemma reads neither $\mathfrak{n}^{\mathrm{kept}}$ nor $C_{\Sigma}$.

(3) \emph{Descendants.} Two kinds of term arise.
\begin{itemize}
\item A term created from a piece of a kept member at an earlier step of one of the kinds written (T) and (P) in the statements on kept members, or at a step of the kind written (L) there, taken for a foreign component. By the statements on kept members at such steps, it is in the format, with support $A:=X$, with $(Y,m)$ in its set of pairs $\mathcal{P}$, and with its arguments in the tube $\mathfrak{S}$.
\item A term created at an earlier healing step $\mathbf{R}_{k_1}$, $k_1<k$, from an exterior piece of a kept member that ended there. It is an output $\mathbf{R}'^{(k_1)}$, $\mathbf{B}'^{(k_1)}$ of \cite[(1.90)--(1.99)]{B-CMP122b}, and it is in the format by the format proposition for these outputs. Its share is carried by the closing condition of the exterior-piece bound at step $k_1$.
\end{itemize}
Each is a stored unit. It is counted in the format's sums of stored units through two kinds of constant:
\begin{itemize}
\item the right members $B_0e^{-\kappa d_j(A)}$, $C_0^{\sharp}$, $e^{\lambda_{\theta}}C_{99}c_1^{\flat}$ of the format's bound;
\item the constants $\mathsf{b}_k$ of the format's bookkeeping statement, together with its base constant $2K(d)B_0+b_{\delta}+NC_0^{\sharp}$.
\end{itemize}
These are the format's constants: $C_0^{\sharp}$ is fixed first, and $B_0$ is fixed before $\varepsilon_V^*$, the format's smallness parameter.

By hypothesis (a), the kept members' share enters only the closing conditions of the statements that return the format. In one such condition, whose right member is $\frac14\varepsilon_V^*$, the kept share $\varepsilon^{\mathrm{kept}}_i$ stands beside $\varepsilon^{\mathrm{rel}}_i$. In the others it enters through $\mu\le\mu^{\flat}+\mu^{\mathrm{kept}}$ and $\mathsf{M}_*\le\mathsf{M}^{\flat}+\mathsf{M}^{\mathrm{kept}}$, with kept share $\mathfrak{n}^{\mathrm{kept}}$ times the far factor. The ceilings on $\gamma$ are re-maximised by the number $\mathfrak{n}^{\mathrm{kept}}$, which is fixed before $\gamma$ in \eqref{4.23:eq-constant-order-a}.

The right members of the format do not change. Had a statement returned a right member of the form $B_0+c\,\mathfrak{n}^{\mathrm{kept}}$, the definition of $C_{\Sigma}$ would pass through $\mathfrak{n}^{\mathrm{kept}}$. Hypothesis (a) excludes this.

(4) \emph{Kept members of other components.} These are members kept by nested regions of other components. By the clause on kept members of other components, all their pieces are exterior. Their label set contains the labelling cube $Q_Y(k)$ of the piece. By the strip-separation statement, this cube has a point on a strip at distance at least $\check{R}_k-m_{\sharp}(k)$ from $C^{\sharp}$. Since $m_{\sharp}(k)=\lfloor\frac12(\check{R}_k-1)\rfloor$, this distance is positive. Such pieces are therefore never filed in $C^{\sharp}$.

(5) \emph{The Gaussian datum inside $b_{\delta}$.} The constant $b_{\delta}$ enters $c_{\mathsf{b}}$ and $c'_{\mathsf{b}}$, hence $C_{\Sigma}$. It contains $b_{\delta}^{G}$, the right member of the per-cube line whose summands are $\mathfrak{N}^{G}(e;\mathfrak{O})e^{\mathsf{w}(e)}e^{r'd_{\ell}(A_e)}$. By hypothesis (b), the datum $\mathfrak{N}^{G}(e;\mathfrak{O})$ reads none of $\mathfrak{O}(v')\le C_{\Sigma}\check{R}^{p_{\Sigma}}$, $\mathsf{S}^{\mathfrak{K}}$ or $\mathfrak{n}^{\mathrm{kept}}$. So $b_{\delta}^{G}$, and with it $b_{\delta}$, reads neither $C_{\Sigma}$ nor $\mathfrak{n}^{\mathrm{kept}}$.

The five cases exhaust the ways in which a kept member can reach $\Sigma^{V}(k)$, and (ii) follows.

\emph{Proof of \textup{(iv)}.} By (i) and (ii), every constant on which $C_{\Sigma}$ depends is chosen before $C_{\Sigma}$ and reads none of $C_{\Sigma}$, $\mathfrak{n}^{\mathrm{kept}}$, $\gamma$. By \eqref{4.23:eq-constant-order-1}, $\mathfrak{n}^{\mathrm{kept}}$ reads only $d$, $C_{\Sigma}$ and $p_{\Sigma}$. By (iii), the conditions closing the bound of $C_{\Sigma}$ read neither $C_{\Sigma}$ nor $\mathfrak{n}^{\mathrm{kept}}$. The ceilings on $\gamma$ read $\mathfrak{n}^{\mathrm{kept}}$, and $\gamma$ is chosen last. Hence no arrow of \eqref{4.23:eq-constant-order-a} points backwards, and the order has no cycle.
\end{proof}

In the corollary below and its proof the following notation is fixed.

Let $(Y,m)$ be a large-field region $Y$ created at step $m$. Its \emph{collar strip} $Y^{\sharp}$ consists of $Y$ together with $m_{\sharp}(m)$ collar layers of $\pi_m$-cubes, and it is the same set at every later level.

The two \emph{kept members} of $Y$ are the resummed density $V^{\sharp}(Y^{\sharp})$ on the strip and the second kept member $\mathfrak{q}_Y$. We write $\mathfrak{K}(Y)=\{V^{\sharp}(Y^{\sharp}),\mathfrak{q}_Y\}$.

For $v\in\mathfrak{K}(Y)$ put
\begin{equation*}
\mathfrak{U}^{\mathrm{tr}}_v=\tilde{U}^{c}_{m}(Y^{\sharp},\tilde\alpha_0,\tilde\alpha_1),
\end{equation*}
which is Ba{\l}aban's space of complex configurations $(\mathbf{U},\mathbf{J})$ over $Y^{\sharp}$ at level $m$ \cite[(2.34)--(2.39)]{B-CMP119}. It is the source tube $\mathfrak{S}$ of Definition~\ref{4.23:not-oscillation-seminorm}, taken over $Y^{\sharp}$ at level $m$.

The \emph{datum} $v^{\sharp}$ of $v$ is the function furnished by clause (ii) of the re-presentation lemma for stored data. It depends on three arguments:
\begin{itemize}
\item a point $P$ of a set $\mathfrak{U}_v\subset\mathfrak{U}^{\mathrm{tr}}_v$;
\item the real variables $\Phi$ read by $v$ (the variables $\Phi_C$ of Definition~\ref{4.23:not-oscillation-seminorm} with $Y^{\sharp}$ in place of $C^{\sharp}$; for $v=\mathfrak{q}_Y$ they are the variables $\varpi$);
\item a complex group-valued parameter $c$ on a set $\Sigma_Y$ of sites, ranging over a polydisc $\mathfrak{G}_{\varepsilon}(\Sigma_Y)$.
\end{itemize}

We write $\Phi^{[g]}$ for the variables $\Phi$ transformed by a $G$-valued map $g$. These transformed variables depend on $g$ only through its restriction $g|_{\Sigma_Y}$.

The number $\mathfrak{O}(v)$ is the supremum, over real $\Phi$ and over two points $(P,c),(P',c')$ of $\mathfrak{U}_v\times\mathfrak{G}_{\varepsilon}(\Sigma_Y)$, of
\begin{equation*}
|v^{\sharp}(P;\Phi,c)-v^{\sharp}(P';\Phi,c')|.
\end{equation*}
Its definition asks of the two points only membership in that set.

A kept member is read at sites of four kinds:
\begin{itemize}
\item (T), the step $i\to i+1$ of the expansion, with $i\ge m$;
\item (L$^{\circ}$), the large-field operation $\mathbf{R}_k$ at which the component of $Z$ holding $Y$ is untreated;
\item (P), the stage $n+1$ of $\mathbf{R}_k$, whose region at stage $n$ is $R_n$. Here the stage is taken on the component holding $Y$ or on another (foreign) component, and the stage may be running or interrupted;
\item (L$^{\dagger}$), the operation $\mathbf{R}_k$ at which the component $C\supset Y$ is treated, with core $\mathsf{C}_C$ and bottom level $h=k-N$.
\end{itemize}

At a site of kind (T) or (L$^{\circ}$), $\mathfrak{Y}(v)$ denotes the set of families $\mathsf{y}$ of cells in the expansion of $v$. We write $|\mathsf{y}|$ for the number of members of $\mathsf{y}$ and $v_{\mathsf{y}}$ for its coefficient. At a site of kind (P), $v_{\hat e}$ denotes the coefficient of a contour term $\hat e$ with family $\hat{\mathsf{y}}$ and block set $\Gamma$, and $\mathsf{K}_{\Gamma}\theta_{\mathsf{K}}(\square)$ are the per-block factors in which that expansion is written. At such a site, $\tilde\square_0$ denotes the enlargement that the geometry of large-field strips attaches to a block $\square_0$ of the stage. Clause (i1) of that geometry places every point of $\square_0$ within $5$ level-$m$ units of $Y^{\sharp}$ whenever $\tilde\square_0$ meets $Y^{\sharp}$.

By $(\mathrm{hol})_v$ at a site we mean that the datum presented by the site is holomorphic on the site's presented polydisc $\mathfrak{P}$.

Distances $d^{\wedge}$ are own-scale contour distances, and $w=M/M_1$. For a cell $\Delta$ the depth is
\begin{equation*}
\varrho(\Delta)=1+d^{\wedge}(\Delta,\mathsf{C}^c)/w,
\end{equation*}
so that $\varrho(\Delta)\ge1$. The anchor functional is read at \emph{probes}; we write $\varsigma$ for a probe and $\varrho(\varsigma)\ge1$ for its depth. The depth lemma provides, for the depth datum $\varrho$ concerned, the depth condition $\check{R}_m\le 2\varrho+2$; in particular this holds for $\varrho=\varrho(\varsigma)$ when the probe $\varsigma$ meets $Y^{\sharp}$. The same lemma provides the two-case distance bound stated in (iii) below. We write $\mathsf{S}(\varsigma)$ for the sum already carried at the probe $\varsigma$ by the anchor functional. A region is \emph{live} at the level read if it still carries its large-field factor there.

\begin{corollary}[The depth polynomial paid additively by the region's own collar]
\label{4.23:cor-collar-payment}
Let $(Y,m)$ be a large-field region created at step $m$, and let $v\in\mathfrak{K}(Y)$. Let $v^{\sharp}$ be its datum on $\mathfrak{U}_v\subset\mathfrak{U}^{\mathrm{tr}}_v$.

Assume the following:
\begin{itemize}
\item the conclusions of the oscillation theorem (Theorem~\ref{4.23:thm-oscillation-theorem}) hold at step $m$;
\item the conditions (q0) and (q1) of the tube statement (Theorem~\ref{4.23:thm-tube-statement}, display \eqref{4.23:eq-tube-statement-a}) hold at step $m$;
\item $(\mathrm{hol})_v$ holds at the site considered.
\end{itemize}

For $\beta>0$ and a probe $\varsigma$ set
\begin{equation}\label{4.23:eq-collar-payment-a}
\begin{gathered}
\mathfrak{n}^{\mathrm{kept}}(\beta):=2^{d}C_{\Sigma}\sup_{\varrho\ge1}(2\varrho+2)^{p_{\Sigma}}
e^{-\beta\varrho},\\
\mathsf{S}^{\mathfrak{K}}(\varsigma):=\sum\bigl\{\mathfrak{O}(v'):\ v'\in\mathfrak{K}(Y'),\
(Y',m')\ \text{live},\ \varsigma\ \text{meets}\ Y'^{\sharp}\bigr\}.
\end{gathered}
\end{equation}
Then the following hold.
\begin{enumerate}
\item[(i)] (The datum.) The oscillations of the two kept members are bounded by the depth polynomial:
\begin{equation*}
\mathfrak{O}(V^{\sharp}(Y^{\sharp}))+\mathfrak{O}(\mathfrak{q}_Y)\le\Sigma^{\mathrm{tot}}(m)
\le C_{\Sigma}\check{R}_m^{p_{\Sigma}},
\end{equation*}
where $\Sigma^{\mathrm{tot}}(m)$ is the total oscillation bound at step $m$ of Theorem~\ref{4.23:thm-oscillation-theorem}.
\item[(ii)] (A site reads a difference.) Let $f(\tau)$ be the value of $v^{\sharp}$ at the triple presented by the site at the point $\tau$ of its polydisc. Then
\begin{equation*}
|f(\tau)-f(\tau')|\le\mathfrak{O}(v)
\end{equation*}
for all $\tau,\tau'$ in the whole polydisc of the site. Consequently, under the proper-scale clause,
\begin{equation*}
|v_{\mathsf{y}}|\le\mathfrak{O}(v)\,e^{-(\kappa_1-1)|\mathsf{y}|},
\end{equation*}
and at a site of kind (P),
\begin{equation*}
|v_{\hat e}|\le\mathfrak{O}(v)\prod_{\square\in\Gamma}\mathsf{K}_{\Gamma}\theta_{\mathsf{K}}(\square)\,
e^{-(\kappa_1-1)|\hat{\mathsf{y}}|}.
\end{equation*}
The quantity $D$ of Ba{\l}aban's inductive description of the effective densities enters neither bound.
\item[(iii)] (The own collar is crossed, in own-scale units.)
\begin{itemize}
\item At a site of kind (T) and at a site of kind (L$^{\circ}$), every $\mathsf{y}\in\mathfrak{Y}(v)$ has $|\mathsf{y}|\ge\frac52\check{R}_m$.
\item At a site of kind (P), on the component holding $Y$ or on a foreign one, running or interrupted, every block $\square_0$ with $\tilde\square_0\cap Y^{\sharp}\neq\emptyset$ satisfies
\begin{equation*}
d^{\wedge}(\square_0,R_n)\ge w\bigl(\tfrac52\check{R}_m-5\bigr).
\end{equation*}
\item At a site of kind (L$^{\dagger}$), with $Y\subset C$ and $C$ treated at $\mathbf{R}_k$,
\begin{equation*}
d^{\wedge}(Y^{\sharp},\mathsf{C}_C^c)\ge
\begin{cases}
\tfrac12 w(\check{R}_m+1) & \text{if } m=h,\\
\tfrac52 w\check{R}_m & \text{if } m<h.
\end{cases}
\end{equation*}
\end{itemize}
\item[(iv)] (A number.) For $\beta>0$ one has $\mathfrak{n}^{\mathrm{kept}}(\beta)<\infty$. Whenever a depth datum $\varrho\ge1$ obeys the depth condition $\check{R}_m\le2\varrho+2$, then
\begin{equation*}
C_{\Sigma}\check{R}_m^{p_{\Sigma}}\le2^{-d}\,\mathfrak{n}^{\mathrm{kept}}(\beta)\,e^{\beta\varrho}.
\end{equation*}
The number $\mathfrak{n}^{\mathrm{kept}}(\beta)$ depends on none of $\check{R}_m$, $D$, the age $i-m$, and the coupling.
\item[(v)] (Additivity at a probe.) At every probe $\varsigma$,
\begin{equation*}
\mathsf{S}^{\mathfrak{K}}(\varsigma)\le\mathfrak{n}^{\mathrm{kept}}(\beta)\,e^{\beta\varrho(\varsigma)}.
\end{equation*}
The anchor functional at a probe becomes that of $\mathsf{S}+\mathsf{S}^{\mathfrak{K}}$. Since $\sup(a+b)\le\sup a+\sup b$, the kept members are one more summand of the sum at a probe, and no product of two counts is formed.
\item[(vi)] (No flow comparison.) Nowhere in (i)--(v) is $\check{R}_m^{p_{\Sigma}}$ converted, compared or paid by the coupling-flow comparison or the flow inequalities of the chapter, which enter here by name only. No relation between the density $V$ at two levels is used.
\end{enumerate}
\end{corollary}

\begin{proof}
\emph{(i).}

\emph{Reduction to one real $\Phi$.} Write the parameter in polar form, $c=e^{\eta}u$, with $u$ a $G$-valued map on $\Sigma_Y$ and $e^{\eta}$ its positive part. Let $P^{\rho_{\mathsf{r}}[\eta]}$ denote the radial re-presentation of $P$ by the positive part $e^{\eta}$ of $c$, furnished by the re-presentation lemma. By clause (ii) of that lemma,
\begin{equation*}
v^{\sharp}(P;\Phi,c)=v\bigl(P^{\rho_{\mathsf{r}}[\eta]};\Phi^{[u]}\bigr).
\end{equation*}
Moreover $P^{\rho_{\mathsf{r}}[\eta]}\in\mathfrak{U}^{\mathrm{tr}}_v$. This follows from clause (i) of the same lemma: starting from the stored presentation, a map of size less than the re-presentation threshold $\varepsilon^{\mathrm{fr}}$ costs at most $\frac14\omega_{\mathrm{fr}}r_a$, where $\omega_{\mathrm{fr}}r_a$ is the margin of that clause.

Let $\hat u$ be any $G$-valued site map on $Y^{\sharp}$ with $\hat u|_{\Sigma_Y}=u$. For $V^{\sharp}(Y^{\sharp})$ we use clause (g) of the oscillation theorem, and for $\mathfrak{q}_Y$ the joint-invariance clause of (q0). Each gives joint invariance $v(\zeta^{g};\Phi^{[g]})=v(\zeta;\Phi)$ for every $G$-valued $g$. Since $\Phi^{[g]}$ depends only on $g|_{\Sigma_Y}$, we have $\Phi^{[\hat u]}=\Phi^{[u]}$. Applying invariance with $g=\hat u$ at $\zeta=Q^{\hat u^{-1}}$ gives, for every $Q\in\mathfrak{U}^{\mathrm{tr}}_v$,
\begin{equation*}
v\bigl(Q;\Phi^{[u]}\bigr)=v\bigl((Q^{\hat u^{-1}})^{\hat u};\Phi^{[\hat u]}\bigr)
=v\bigl(Q^{\hat u^{-1}};\Phi\bigr).
\end{equation*}
The space $\mathfrak{U}^{\mathrm{tr}}_v$ is invariant under every $G$-valued gauge transformation on $Y^{\sharp}$; no regularity of the transformation is required \cite[\S2]{B-CMP119}. Hence
\begin{equation*}
Q(P,c):=\bigl(P^{\rho_{\mathsf{r}}[\eta]}\bigr)^{\hat u^{-1}}\in\mathfrak{U}^{\mathrm{tr}}_v,
\qquad v^{\sharp}(P;\Phi,c)=v\bigl(Q(P,c);\Phi\bigr).
\end{equation*}

\emph{Bound by the oscillation.} Take two points $(P,c),(P',c')$ at one real $\Phi$. The two values of the datum are then values of $v$ at the single real $\Phi$ and at two points of $\mathfrak{U}^{\mathrm{tr}}_v=\mathfrak{S}$. Their difference is therefore at most $\operatorname{osc}(v)$ of Definition~\ref{4.23:not-oscillation-seminorm}. Taking the supremum gives $\mathfrak{O}(v)\le\operatorname{osc}(v)$.

\emph{Summing the two members.} By clauses (v) and (t) of the oscillation theorem (display \eqref{4.23:eq-oscillation-theorem-a}), and by (q1) of \eqref{4.23:eq-tube-statement-a}, all at step $m$, the sum of the oscillations of $V^{\sharp}(Y^{\sharp})$ and $\mathfrak{q}_Y$ is at most $\Sigma^{\mathrm{tot}}(m)$. The inequality $\Sigma^{\mathrm{tot}}(m)\le C_{\Sigma}\check{R}_m^{p_{\Sigma}}$ is the degree-and-constant statement of the oscillation theorem, Lemma~\ref{4.23:prf-degree-constant}, with $p_{\Sigma}$ and $C_{\Sigma}$ as in \eqref{4.23:eq-degree-constant-a}.

\emph{(ii).}

\emph{The difference bound.} Let $\tau,\tau'$ lie in the polydisc of the site. Both presented triples consist of a point of $\mathfrak{U}_v$ and a parameter in $\mathfrak{G}_{\varepsilon}(\Sigma_Y)$, at one and the same real $\Phi$. The definition of $\mathfrak{O}(v)$ asks of the two points only membership in $\mathfrak{U}_v\times\mathfrak{G}_{\varepsilon}(\Sigma_Y)$. Hence $|f(\tau)-f(\tau')|\le\mathfrak{O}(v)$.

\emph{Under the proper-scale clause.} The function $f$ is holomorphic on the polydisc of the site by the third hypothesis of the corollary, the holomorphy condition $(\mathrm{hol})_v$ defined before its statement. We apply clause (ii) of the source-cell expansion lemma in its second form. That clause states: if a function $F$ differs, on the whole polydisc, by at most $\mathfrak{O}(v)$ from a constant independent of the interpolation variables $\vec{\mathsf{s}}$, then $|F_{\mathsf{y}}|\le\mathfrak{O}(v)e^{-(\kappa_1-1)|\mathsf{y}|}$. We take $F=f$ and, as the $\vec{\mathsf{s}}$-free constant, the value $f(0)$. The hypothesis holds by the difference bound just proved, with $\tau'=0$. This gives the bound on $v_{\mathsf{y}}$.

\emph{At a site of kind (P).} The Cauchy estimates along lines of the contour expansion are applied to $F_v-F_v(0)$, where $F_v=f$ is the presented function. This is done as in clause (a) of the bound for kept members at stage sites. It yields the product bound on $v_{\hat e}$.

In neither bound does $D$ occur.

\emph{(iii).}

\emph{Geometric input.} We use clause (c) of the geometry of large-field strips. It states three things. First, the level-$m$ distance from $Y^{\sharp}$ to the current $\Omega_{m+1}$ is at least $\frac52\check{R}_m$. Second, every cell at level-$m$ distance less than $\frac52\check{R}_m$ from $Y^{\sharp}$ has level at most $m$. Third, this neighbourhood lies, for every $i>m$, in the component of $\Lambda_i^c$ that holds $Y$.

\emph{Kind (T).} At a site of kind (T) there are no hub cells. Let $\mathsf{y}\in\mathfrak{Y}(v)$. By clause (b) of the same geometry, a member of $\mathsf{y}$ that sees $Y^{\sharp}$ meets it, and hence lies in it. The component of $\mathsf{y}$ containing this member is wall-connected and holds a source cell. That source cell holds free bonds of $\Omega_{i+1}$, and $\Omega_{i+1}\subset\Omega_{m+1}$. The wall-chain from the member in $Y^{\sharp}$ to this source cell therefore leaves the neighbourhood. Inside the neighbourhood every cell has level at most $m$, hence side at most one level-$m$ unit. The chain must cross a width of $\frac52\check{R}_m$ such units. Hence it has at least $\frac52\check{R}_m$ members, and $|\mathsf{y}|\ge\frac52\check{R}_m$.

\emph{Kind (L$^{\circ}$).} The hub cells and the source cells lie in $Z$. The neighbourhood lies in the untreated component of $Y$. Up to its first source cell, a connected family is a wall-chain, by the convention under which connected families are read in the geometry of large-field strips. The count made for kind (T) therefore applies unchanged and gives $|\mathsf{y}|\ge\frac52\check{R}_m$.

\emph{Kind (P).} The region $R_n$ lies outside the neighbourhood, in each of the three possible cases:
\begin{itemize}
\item on the own component with $m<h$, the cells of $R_n$ have level at least $h>m$, whereas cells of the neighbourhood have level at most $m$;
\item on the own component with $m=h$, this is clause (h) of the geometry;
\item on a foreign component, $R_n\subset Z$.
\end{itemize}

Since $\tilde\square_0$ meets $Y^{\sharp}$, clause (i1) of the geometry gives that every point of $\square_0$ lies within $5$ level-$m$ units of $Y^{\sharp}$. A contour from $\square_0$ to $R_n$ therefore runs through at least $\frac52\check{R}_m-5$ level-$m$ units of cells of level at most $m$. By the own-scale contour bound, on such cells a contour pays at least $w$ per level-$m$ unit. Hence $d^{\wedge}(\square_0,R_n)\ge w(\frac52\check{R}_m-5)$. This is the depth bound at stage sites.

\emph{Kind (L$^{\dagger}$).} The two inequalities are the two cases, $m=h$ and $m<h$, of the two-case distance bound of the depth lemma.

\emph{(iv).} A polynomial is dominated by an exponential: $(2\varrho+2)^{p_{\Sigma}}e^{-\beta\varrho}\to0$ as $\varrho\to\infty$. Since this function is continuous on $[1,\infty)$, its supremum there is finite, so $\mathfrak{n}^{\mathrm{kept}}(\beta)<\infty$.

If $\varrho\ge1$ and $\check{R}_m\le2\varrho+2$, then, because $p_{\Sigma}>0$,
\begin{equation*}
C_{\Sigma}\check{R}_m^{p_{\Sigma}}
\le C_{\Sigma}(2\varrho+2)^{p_{\Sigma}}e^{-\beta\varrho}\,e^{\beta\varrho}
\le2^{-d}\,\mathfrak{n}^{\mathrm{kept}}(\beta)\,e^{\beta\varrho}.
\end{equation*}
By \eqref{4.23:eq-collar-payment-a}, $\mathfrak{n}^{\mathrm{kept}}(\beta)$ is built from $d$, $C_{\Sigma}$, $p_{\Sigma}$ and $\beta$ alone. It therefore depends on none of $\check{R}_m$, $D$, the age and the coupling.

\emph{(v).} By clause (i2) of the geometry, a probe meets at most $2^{d}$ strips $Y'^{\sharp}$ of live regions. Fix such a region $(Y',m')$ and bound its contribution in three steps:
\begin{itemize}
\item by (i), its contribution to $\mathsf{S}^{\mathfrak{K}}(\varsigma)$ is at most $C_{\Sigma}\check{R}_{m'}^{p_{\Sigma}}$;
\item by the depth condition, $\check{R}_{m'}\le2\varrho(\varsigma)+2$;
\item by (iv), the contribution is therefore at most $2^{-d}\mathfrak{n}^{\mathrm{kept}}(\beta)e^{\beta\varrho(\varsigma)}$.
\end{itemize}

Summing over at most $2^{d}$ strips gives the bound. This is the anchor bound for kept members. The remaining assertions of (v) follow from $\sup(a+b)\le\sup a+\sup b$, applied at each probe to $a=\mathsf{S}$ and $b=\mathsf{S}^{\mathfrak{K}}$.

\emph{(vi).} In this section, comparisons along the coupling flow occur only at the following places:
\begin{itemize}
\item in the window comparison $N\le c_N\check{R}_k$ and the companion radius comparison of \eqref{4.23:eq-degree-constant-a}, where $\check{R}_{h+1}$ is compared with $\check{R}_k$ over the own window $N$ of step $k$;
\item in the age comparisons inside the lemmas whose hypotheses enter the oscillation theorem;
\item in the rider statement of this section on the Wilson action;
\item in the hull bound among the form clauses of \eqref{4.23:eq-tube-statement-a}.
\end{itemize}

Each of these is either a ratio of couplings paid by zones crossed at the step's own site, or a comparison of a number with its own value at a later level. None of them is a payment of $\check{R}_m^{p_{\Sigma}}$.
\end{proof}

\begin{remark}
The bounds of (iii) are stated in the region's own units, for the following reason.

Consider a site of level $i\gg m$. There the collar of $Y$ is $\frac52\check{R}_mL^{-(i-m)}$ level-$i$ units wide; this is an identity of units between levels $m$ and $i$. Measured in units of the level read, the collar is therefore negligible.

The proper-scale clause, however, cuts all cells in their proper scales, and $d^{\wedge}$ counts own $M_1$-units. Hence, whatever $i$ is, a family or contour that reaches $Y^{\sharp}$ counts at least $\frac52\check{R}_m$ cells of level at most $m$.

For this reason the expansion of a kept member runs under the proper-scale clause. The decay lemma used in the proof of the correlation theorem, whose decay factor decays in the units of the level read, is never applied to it.
\end{remark}

\begin{remark}[Multiplicity at a cover cube forces no lower bound on the coupling]
\label{4.23:rem-cover-multiplicity}
We recall the cover-cube multiplicity lemma. Its bounds concern the sum of terms at one cover cube
$\square'$ of stage $n+1$ of the cover, $n$ denoting here the stage index of the cover and not a
number of points; they are stated in terms of a level $j$, a threshold level $k_0$,
a factor $\mathsf{l}_j$ and constants $B_0$, $K_d$, $N_0$, $C_0$ of that lemma. A per-cube
estimate splits this sum into
three contributions, labelled $(\alpha)$, $(\beta)$, $(\gamma)$; we write $S_{\beta}(\square')$
and $S_{\gamma}(\square')$ for the contributions $(\beta)$ and $(\gamma)$. For every such
$\square'$:
\begin{itemize}
\item the bound on the contribution $(\alpha)$ is the one of the per-cube estimate from which the
three contributions are taken, unchanged;
\item the contributions $(\beta)$ and $(\gamma)$ satisfy the bounds displayed below;
\item the total of the companion per-cube estimate, as bounded through the first three terms of
the recursion that controls it, acquires the factor $\mathsf{l}_j^2$.
\end{itemize}
These bounds are
\begin{equation}\label{4.23:eq-cover-multiplicity-1}
S_{\beta}(\square')\le \mathsf{l}_j^2\cdot 5B_0K_dL^{2d},
\qquad
S_{\gamma}(\square')\le \mathsf{l}_j^2\cdot 3\cdot 4^d(N_0+3)C_0 .
\end{equation}
For $j>k_0$ the factor $\mathsf{l}_j^2$ cannot be dropped from these bounds. It is bounded above
in terms of $N_0$; this is a one-sided bound and not a cap in the sense of
Definition~\ref{2.3:def-cap}.

The factor cannot be dropped because of the following configurations. Take a level $m$ with
$k_0<m<j$ (this $m$ is used only in this paragraph and is not the torus exponent), and the set
$T=\mathbb{B}^{(m)}(X)$ of level-$m$ units attached to a region $X$, with $d_m(X)$
bounded by a constant. A single cover cube of level $j$ can contain as many as $(2L^{j-m})^d$
level-$m$ units of $T$.

Hence the sum of terms at one cover cube of a window stage carries the factor $\mathsf{l}_j^2$,
which is of order $(\log x)^{4r}$, with $x$ the inverse square coupling of the window stage and
$r$ the exponent for which $\mathsf{l}_j$ is of order $(\log x)^{2r}$. A bound in which this
factor multiplied the number $V'$ of
blocks would put a lower bound on the coupling.

In the estimates of this section the multiplicity never multiplies $V'$. It appears undropped,
but only in two places:
\begin{itemize}
\item within a single cell, before the cells are summed;
\item as a factor inside one of the cost terms of the right member $\Sigma^V(k)$ of the
oscillation theorem.
\end{itemize}
No lower bound on the coupling is forced by it.
\end{remark}

\begin{proof}
We follow the three places at which the multiplicity of the cover-cube multiplicity lemma could
enter the estimates of this section.

\emph{First.} No sum over sites of the per-cube total bounded in
\eqref{4.23:eq-cover-multiplicity-1} is formed in these estimates.

\emph{Second.} At the step of these estimates at which the fine terms inside one cell are
bounded, the same multiplicity does appear,
and it is not dropped. It takes the form of the factor $\mathsf{w}(\Delta)^d$ in the bound on the
fine terms inside one cell.
Here
$\mathsf{w}(\Delta)$ is the width, in own cubes, of a fine term inside the cell $\Delta$, so that
$\mathsf{w}(\Delta)^d$ is the number of the fine term's own cubes inside $\Delta$; it
is multiplied by the bound $\mathsf{b}$ per own cube. Two cases arise, according to whether
$\Delta$ lies in the core $\mathsf{C}$ or not.

Suppose first that $\Delta$ lies in the core $\mathsf{C}$. There the factor
$\mathsf{w}(\Delta)^d$ meets the depth of that same cell. This happens in the same supremum, and
before the cells are summed. The window form of the depth inequality gives, with $C_{\uparrow}$
the constant of that inequality,
\begin{equation}\label{4.23:eq-cover-multiplicity-2}
\mathsf{w}(\Delta)\le L\,C_{\uparrow}\,\varrho(\Delta),
\end{equation}
where $\varrho(\Delta)$ is the depth $1+d^{\wedge}(\Delta,\mathsf{C}^c)/(M/M_1)$ of the cell. So
the multiplicity inside one cell of the core is bounded by a power of that cell's own depth. The
subsequent sum over cells, clause~(b) of the lemma bounding sums over cells, sums exponential
factors that decay in
the own-scale contour distance $d^{\wedge}$. It does not sum counts.

Suppose next that $\Delta$ lies off the core. There the width satisfies
\begin{equation}\label{4.23:eq-cover-multiplicity-3}
\mathsf{w}(\Delta)\le L\,\check{R}_{h+1},
\end{equation}
and this product of $\mathsf{w}(\Delta)^d$ with the per-own-cube bound $\mathsf{b}$ is summed over
at most $\mathsf{n}_{\mathrm{sh}}(k)$ cells, $\mathsf{n}_{\mathrm{sh}}(k)$
being the bound on the number of cells off the core. The result is a product
of two polynomials in $\log x_k$. This is the form the cover-cube multiplicity lemma says such a
sum must have. The product stands on the cost side, as one of the terms of the right member
$\Sigma^V(k)$ of the oscillation theorem, where it is required to be neither small nor bounded.
Hence no lower bound on the coupling can arise from it.

\emph{Third.} The depth polynomial $\check{R}_m^{\,p_\Sigma}$, with $p_\Sigma$ its exponent and
$m$ the level of the region to which Corollary~\ref{4.23:cor-collar-payment} applies, in the
notation of that corollary, enters the later estimates through
clause~(v) of the collar payment (Corollary~\ref{4.23:cor-collar-payment};
see~\eqref{4.23:eq-collar-payment-a}). It enters there as a summand at a probe of the collar
payment. When such an
estimate sums over its probes or over its blocks, the quantity that multiplies the number $V'$ of
blocks is the kept number $\mathfrak{n}^{\mathrm{kept}}$ of~\eqref{4.23:eq-collar-payment-a}. The
product
$V'\cdot\mathfrak{n}^{\mathrm{kept}}$ is a count times a number. What multiplies $V'$ is never a
per-block sum containing the factor $\mathsf{l}_j^2$.

The reason lies in the terms $v$ to which the collar payment of
Corollary~\ref{4.23:cor-collar-payment} applies. For every such term $v$, by the statement on the
member families of the terms to which the collar payment applies, the family
$\mathcal{A}_v$ attached to $v$ lies in the class $\mathfrak{O}\mathsf{ld}$ of old members, and
the members of $\mathcal{A}_v$ meet no anchor of the window. In the cover-cube multiplicity lemma
the factor $\mathsf{l}_j^2$ is attached to the sum of terms at a cover cube of a window stage;
since the members of $\mathcal{A}_v$ meet no anchor of the window, no per-block sum containing
$\mathsf{l}_j^2$ multiplies $V'$.

In all three places the multiplicity either is never formed or appears in one of two ways. It may
be absorbed cell by cell into the cell's own depth by~\eqref{4.23:eq-cover-multiplicity-2}. Or it
may stand inside a cost term of $\Sigma^V(k)$ by~\eqref{4.23:eq-cover-multiplicity-3}, where no
smallness or boundedness is required. It never multiplies $V'$. This proves the remark.
\end{proof}

\begin{remark}[The place of the tube statement within one step]\label{4.23:rem-tube-placement}
Fix a step $k$, let $N=N(g_k)$ be the number of levels of the window, and put $h=k-N$. One step of the construction runs through the following consecutive stages, beginning with the closing stage of step $k-1$.
\begin{enumerate}
\item The closing stage of step $k-1$.
\item Four intermediate stages of step $k$, followed by the base stage of step $k$. At the base stage the three-factor bound and its source-dependent form become available. Both use the construction only up to that point.
\item The presentation stage. It supplies the presentations, the pieces and the lists of the step; the proposition of the presentation stage; the branch condition and the value condition; and the two concluding lemmas of the stage.
\item The tube stage. It consists of the tube statement at step $k$ (Theorem~\ref{4.23:thm-tube-statement}). That statement comprises the form conditions, the clause (q0) and the bound (q1) of~\eqref{4.23:eq-tube-statement-a}, together with the oscillation theorem for the resummed density at step $k$ (Theorem~\ref{4.23:thm-oscillation-theorem}), whose clauses are displayed in~\eqref{4.23:eq-oscillation-theorem-a}.
\item The modification stage. It contains the modification statements of the step, the displayed identities of that stage, the domain classes, and the families $\mathfrak{K}(Y_i)$ of units attached to the new regions $Y_i$.
\item The closing stage of step $k$, which opens step $k+1$.
\end{enumerate}
Order the pairs $(\text{step},\text{stage})$ lexicographically.

The inputs of the oscillation theorem at step $k$, and the stage of each, are as follows.
\begin{itemize}
\item[(i)] From the base stage of step $k$ (or earlier): the three-factor bound in its source-dependent form and its companion lemma, and the further bounds of that stage that appear among the hypotheses of Theorem~\ref{4.23:thm-oscillation-theorem}.
\item[(ii)] From the presentation stage of step $k$ or earlier stages of step $k$:
\begin{itemize}
\item the list of units of the treated large-field component $C$ of Theorem~\ref{4.23:thm-tube-statement}, with their labels and the format data $\mathfrak{N}(v)$, $\mathfrak{N}^{\flat}(v)$ of each unit $v$;
\item the complex relative hypothesis at step $k$;
\item the construction of the presentations, in the parts listed among the hypotheses of Theorem~\ref{4.23:thm-oscillation-theorem};
\item the deep-region statement and the boundary statement;
\item clauses (a) and (c) of the proposition of the presentation stage.
\end{itemize}
\item[(iii)] From the presentation stage of step $k$: the branch and value conditions, clauses (d) and (e) of the proposition of that stage, and the lemma of that stage bounding the relative and kept members together.
\item[(iv)] From the tube stage of step $k$ itself: the clauses (q0) and (q1) at step $k$.
\item[(v)] From the tube stages of the steps $m'<k$: the oscillation theorem and the clauses (q0), (q1) at $m'$. These enter through the oscillation bound for units. At the present presentation stage, they are complemented by the difference bound for the units of the family $\mathfrak{K}(Y')$ of a region $Y'$ created at a step $m'<k$ and by clauses (a) and (c) of the holomorphy statement for those units on the presented polydisc $\mathfrak{P}$.
\end{itemize}

The oscillation theorem at step $k$ is an input of the modification stage of step $k$, where $\mathfrak{K}(Y_i)$ is created with oscillation at most $\Sigma^{\mathrm{tot}}(k)$, the total oscillation bound at step $k$. It is also available as an input at every stage of every step later than $k$. No statement of the presentation stage of step $k$ uses it.

Consequently, the placement of the tube stage is consistent with the lexicographic order of the pairs $(\text{step},\text{stage})$: every input of the oscillation theorem at step $k$ lies at a strictly earlier pair or is one of the clauses (q0), (q1) at the same pair, whose derivation does not use the oscillation theorem, and every use of it lies at a strictly later pair.
\end{remark}

\begin{proof}
Every input of a statement at step $k$ must lie at a strictly earlier pair. The only exception is an input at the same pair that does not itself use the statement in question. We check this for the inputs and for the users of the oscillation theorem at step $k$, which sits at the pair $(k,\text{tube stage})$.

Inputs (i), (ii) and (iii) lie at the base stage, the presentation stage or earlier stages of step $k$. All of these precede the tube stage of step $k$. In particular the oscillation theorem uses the presentation stage of step $k$ and is established after it.

Input (iv) lies at the same pair. The derivation of the clauses (q0) and (q1) at step $k$, given with Theorem~\ref{4.23:thm-tube-statement}, uses none of the step's standing hypotheses, no tube statement at an earlier step and no item of the stage's auxiliary apparatus; being carried out outside the inductive construction, it does not use the oscillation theorem at step $k$. Within the tube stage the dependence therefore runs only from the oscillation theorem to (q0) and (q1), never back.

Input (v) consists of statements at the pairs $(m',\text{tube stage})$, $m'<k$, which precede $(k,\text{tube stage})$; it also includes the difference bound and the holomorphy clauses, which belong to the presentation stage of step $k$ and so precede the tube stage. The nested-region lemma (Lemma~\ref{4.23:lem-nested-regions}) is where these inputs at the steps $m'<k$ are used.

Now the users. The modification stage of step $k$ follows the tube stage of step $k$. Every stage of a step $k'>k$ follows every stage of step $k$. So each use of the oscillation theorem at step $k$ occurs at a strictly later pair.

It remains to exclude a use at the presentation stage of step $k$, which precedes the tube stage. The lemma of that stage bounding the relative and kept members together uses the tube statement (Theorem~\ref{4.23:thm-tube-statement}) only at the steps $m'\le h$. Since $h<k$, these are tube stages of strictly earlier steps. The first of the modification statements of the step uses the size of the resummed density through its weighted size bound. It belongs to the modification stage, not the presentation stage. No other statement of the presentation stage uses the oscillation theorem.

Every dependence edge incident to the oscillation theorem at step $k$ therefore runs between strictly ordered pairs, except the one edge to (q0), (q1) at the same pair, which is not reversed. This is the asserted consistency of the placement of the tube stage with the lexicographic order; it concerns the statements that the oscillation theorem at step $k$ uses and the statements that use it.
\end{proof}

\begin{definition}[One-sided conditions on the parameters of the tube statement]
\label{4.23:def-parameter-conditions}
The statements of this section about points $\zeta$ of $\mathfrak{S}_{\mathrm{src}}$, in particular the
form theorem (Theorem~\ref{4.23:thm-form-theorem}) and the degree-and-constant
statement~\ref{4.23:prf-degree-constant} of the oscillation theorem, use conditions on the parameters of
two kinds.

A \emph{floor} is a lower bound on a parameter that is fixed before the smallness parameter $\gamma$.

A \emph{$\gamma$-ceiling} is an inequality required in the supremum form, either for all
$x\ge\gamma^{-2}$, where $x$ ranges over values of the inverse squared coupling, or for all
$\mathsf{T}\ge\mathsf{T}_\gamma$, where $\mathsf{T}=\log x$ and $\mathsf{T}_\gamma=\log\gamma^{-2}$ is
the threshold of $\mathsf{T}$ corresponding to $x=\gamma^{-2}$. Here $\gamma$ is
chosen after every other parameter.

Every condition below is one-sided.

\smallskip
\noindent\emph{Floors.} Each of the following is a lower bound on a parameter fixed before $\gamma$.
None of them is new to this section.
\begin{enumerate}
\item[(F1)] $\kappa_1-1\ge 3c_l+\lambda_*$. This is a floor of the earlier statement of this chapter
that also imposes the condition $\check{R}_k\ge 38$ on the window radius, listed below among the
$\gamma$-ceilings of other statements; we call it the window-radius statement, and $c_l$ and
$\lambda_*$ are constants of that statement. Since $\log\frac{e}{e-1}<1$, this floor also covers the
condition $\kappa_1\ge\log\frac{e}{e-1}$ on the radius of the Cauchy estimate used in the bound of
this chapter on the total sum over the units of the expansion (the total-sum bound). That condition
therefore adds no floor.
\item[(F2)] $\tfrac1{16}\,\delta_P\,M/M_1\ge 1$, a further floor of the window-radius statement.
Multiplying by $2$ and by $8$ gives
\begin{equation}\label{4.23:eq-parameter-conditions-1}
\tfrac18\,\delta_P\,M/M_1\ge 2,\qquad \tfrac12\,\delta_P\,M/M_1\ge 8 .
\end{equation}
\item[(F3)] $\kappa\ge\kappa_d$, where $\kappa_d$ is a threshold on $\kappa$ with $\kappa_d\ge 180$.
In particular $\tfrac85\kappa\ge 288$, so the value $r=4$ taken in this section by the parameter $r$
satisfies $r\le\tfrac85\kappa$.
\item[(F4)] $\tfrac1{20}\kappa\ge 2(\mathsf{k}_d+1)$, a floor of the earlier statement of this
chapter in which the constant $\mathsf{k}_d$ occurs; we call it the $\mathsf{k}_d$-statement. It
follows that the quantity
$a:=\tfrac{31}{20}\kappa$ satisfies $a\ge\mathsf{k}_d+1$. Indeed, (F4) and $\kappa>0$, which holds by
(F3), give
\begin{equation*}
\mathsf{k}_d+1\le\tfrac1{40}\,\kappa\le\tfrac{31}{20}\,\kappa .
\end{equation*}
\item[(F5)] $L\ge 8$, a floor of the statement that supplies the current bound and the count of cells
used in the proof of Lemma~\ref{4.23:lem-plaquette-units}. This floor is used in a box-containment
check and in the geometric sum of a shell bound, both in this chapter.
\item[(F6)] The four floors assumed in Theorem~\ref{4.23:thm-form-theorem}: a further floor on
$\kappa_1$ (distinct from (F1)), a
floor on the levels, a floor on $M_1$, and a floor on the fine block
size $\mathfrak{M}$. They are the floors of the statement on the localised small solution
$\mathbb{H}$ at the large-field region that is assumed in that theorem at the sets
$\mathfrak{B}^0=\mathbb{B}_k(C)$ and
$\mathfrak{B}^C=\mathbf{B}''_k(C)$.
\item[(F7)] The floors of the correlation theorem. Among them are $L\ge 27$, $M\ge 2L^2$ and
$\delta_1M\ge\kappa_1$, where $\delta_1$ is the constant of that theorem so denoted, together with the
condition \cite[(2.59)]{B-CMP109} with the constant $\tfrac18$, read on the quantity $R_1M_1$ of that
theorem.
\item[(F8)] The remaining floors of the window-radius statement and of the
$\mathsf{k}_d$-statement.
\end{enumerate}
The relation between the free radii $\gamma^{\flat}$ and $\varepsilon_3^{\flat}$ of which
Definition~\ref{4.23:def-coupling-free-radii} gives an admissible instance is met by that choice. It
is therefore not a condition on the parameters.

\smallskip
\noindent\emph{$\gamma$-ceilings imposed in this section.}
\begin{enumerate}
\item[(C1)] The following condition, in which $\lambda$ and $c_2$ are the constants of the ratio form
$x_n\le x_k(1+3\lambda(k-n)+2c_2)$, $n\le k$,
of the comparison of inverse squared couplings used in the proof of
Lemma~\ref{4.23:lem-plaquette-units}, $C_N$ is a constant fixed before $\gamma$, and
$r_{\mathrm{pr}}$ is the exponent entering the degree $p_\Sigma$ of
\eqref{4.23:eq-degree-constant-a}:
\begin{equation}\label{4.23:eq-parameter-conditions-2}
1+\log\Bigl(1+3\lambda\bigl(C_N(\mathsf{T}+1)^{r_{\mathrm{pr}}}+2\bigr)+2c_2\Bigr)\le 2\mathsf{T}
\qquad\text{for all }\mathsf{T}\ge\mathsf{T}_\gamma .
\end{equation}
Put $\mathsf{c}(u):=\log(1+3\lambda u+2c_2)$ for $u\ge 0$. Taking logarithms in the ratio form
gives its logarithmic form $\log x_n\le\mathsf{T}_k+\mathsf{c}(k-n)$, where $\mathsf{T}_k=\log x_k$,
and in this notation \eqref{4.23:eq-parameter-conditions-2} reads
\begin{equation*}
1+\mathsf{c}\bigl(C_N(\mathsf{T}+1)^{r_{\mathrm{pr}}}+2\bigr)\le 2\mathsf{T}
\qquad\text{for all }\mathsf{T}\ge\mathsf{T}_\gamma .
\end{equation*}
This is the condition that \eqref{4.23:eq-degree-constant-a} carries beside the condition of (C2).
In the proof of Lemma~\ref{4.23:lem-plaquette-units} it is invoked to bound
$\mathsf{T}_k+\mathsf{c}(N)$, with $N=N(g_k)$ the window, by $3\check{\mathsf{T}}_k$.
\item[(C2)] The following condition, which \eqref{4.23:eq-degree-constant-a} carries beside the
condition of (C1), with the constants of that display:
\begin{equation}\label{4.23:eq-parameter-conditions-3}
C_{3F}\,e^4\,C_{\mathrm{cell}}\cdot 2d\,(104Lc_R+2)^{d-1}\cdot
\sup_{\varrho\ge\frac12\mathsf{T}_\gamma^{r_{\mathrm{pr}}}}(2\varrho)^{p_\Sigma}e^{-2(\varrho-1)}
\le 1 .
\end{equation}
The supremum has degree $-\infty$ in the sense of the degree table of
\eqref{4.23:eq-degree-constant-a}; in particular it tends to zero faster than every power of
$\mathsf{T}_\gamma$ as $\mathsf{T}_\gamma\to\infty$. In \eqref{4.23:eq-degree-constant-a} this
condition is written with the quantity $\varepsilon_\gamma$ of \eqref{4.23:eq-nested-regions-a} in the
place of the supremum. The constant in front of the supremum in
\eqref{4.23:eq-parameter-conditions-3} contains neither the constant $C_\Sigma$ of
the degree-and-constant statement~\ref{4.23:prf-degree-constant} nor the quantity
$\mathfrak{n}^{\mathrm{kept}}$ of the total-sum bound.
\item[(C3)] The condition \eqref{4.23:eq-form-theorem-b} of the form theorem, with the constants of
that display:
\begin{equation}\label{4.23:eq-parameter-conditions-4}
14\,(1+K_\Lambda)\,c_Q\,\bar\alpha_1(x)\le\rho_{\mathrm{abs}}
\qquad\text{for all }x\ge\gamma^{-2},
\end{equation}
with $\rho_{\mathrm{abs}}$ as in Definition~\ref{4.23:def-coupling-free-radii}. The left side has
degree $-\tfrac12$ in the sense of the degree table of \eqref{4.23:eq-degree-constant-a}.
\item[(C4)] The size conditions of the sixth hypothesis of Theorem~\ref{4.23:thm-form-theorem}.
These require that the real base configurations lie in the class
$\mathfrak{K}_{109}(\,\cdot\,;\tfrac12\gamma^{\flat})$: the sizes of Ba{\l}aban's construction that
bound these configurations, which are small when $\gamma$ is small, must not exceed the bounds
defining that class, for all $x\ge\gamma^{-2}$.
\item[(C5)] Used only for Lemma~\ref{4.23:lem-plaquette-units}, on which no other statement of this
section depends, is the condition \eqref{4.23:eq-plaquette-units-a}, with the constants of that
display:
\begin{equation}\label{4.23:eq-parameter-conditions-5}
\sup_{x\ge\gamma^{-2}}\,40\,c_{\mathsf{U}}\,C_{\mathsf{a}}\,\alpha_*(x)\le 1 .
\end{equation}
\end{enumerate}

\smallskip
\noindent\emph{$\gamma$-ceilings of other statements.} Each of these is imposed in the same
supremum form where the statement concerned is stated, and is used here unchanged. They are:
\begin{itemize}
\item the five ceilings of the statement that supplies the bound on the size of the fluctuation used
in the proof of Lemma~\ref{4.23:lem-plaquette-units};
\item the two ceilings of the statement that supplies the current bound and the count of cells used
in the proof of Lemma~\ref{4.23:lem-plaquette-units}, namely the regularity condition under which
that lemma is proved and one further ceiling;
\item the three ceilings of the window-radius statement, among them the
condition $\check{R}_k\ge 38$ on the window radius;
\item three ceilings of the earlier statement of this chapter in which the degree $p_\Sigma$ of
\eqref{4.23:eq-degree-constant-a} enters only through suprema over $\varrho$ of a power of degree
$p_\Sigma$ of $2\varrho+2$ times an exponential decaying in $\varrho$, suprema which are finite for
every value of $p_\Sigma$;
\item one ceiling of a further earlier statement of this chapter, distinct from the statements named
in the other items of this list;
\item one ceiling of the earlier statement of this chapter in which the two product weights
$\Pi^{\mathrm{sh}}$ and $\Pi^{L}_{G}$ occur, together with the bounds $\Pi^{\mathrm{sh}}\le 1$ and
$\Pi^{L}_{G}\le 1$;
\item one ceiling of the earlier statement of this chapter whose second size condition on the second
kept member is replaced here by (C3), its first size condition being (C4);
\item one ceiling of the statement on the localised small solution $\mathbb{H}$ whose floors are
listed in (F6);
\item the standing smallness hypothesis $R_\gamma\le\check{R}_j$ of the correlation theorem, at the
levels $j$ of the steps treated by that theorem, with $R_\gamma$ as in
\eqref{4.23:eq-nested-regions-a}.
\end{itemize}

A further condition, a smallness condition involving the window radius $\check{R}_k$, belongs to the
statement on the boundary terms on which the oscillation theorem of this section relies in place of
Lemma~\ref{4.23:lem-plaquette-units}; it is a condition of that statement and not of this section.
It arises by reading Ba{\l}aban's lower bound on the number of levels of his window with that number
set equal to $\check{R}_k$. It is not a cap in the sense of Definition~\ref{2.3:def-cap}.

\smallskip
\noindent\emph{Properties of the list.}
\begin{itemize}
\item This section introduces no new floor.
\item It introduces no $\gamma$-ceiling beyond (C1)--(C5).
\item If the remark on a logarithmic threshold, made in the proof of
Theorem~\ref{4.23:thm-form-theorem} for the additional clause that this theorem proves under
additional hypotheses, is invoked, only the numerical constant in
\eqref{4.23:eq-parameter-conditions-4} is enlarged; the condition keeps its one-sided supremum form,
and the enlarged constant is needed only for that clause.
\item No condition bounds from above any of $L$, $M$, $K$, $k$, the window $N=N(g_k)$, the
parameter $N_0$, $\check{R}$, the ages of large-field regions (the numbers of steps since the steps
that created them), the quantity $D$, or the number of steps.
\item No condition bounds the coupling from below. As $g\downarrow 0$, the quantity $\mathsf{T}$ and
the ratio $\rho_{\mathrm{abs}}/s_0$ of the coupling-free radius $\rho_{\mathrm{abs}}$ to the quantity
$s_0$ grow, and so does, for every $\gamma$-ceiling, the margin by which its inequality holds at the
point $x=g^{-2}$ (respectively $\mathsf{T}=\log g^{-2}$), that is, the right side minus the left side
evaluated there.
\item No condition is two-sided.
\item No value of the constant $L_0$ is fixed or used.
\end{itemize}
\end{definition}

\begin{lemma}[The two halves of the tube statement: inputs and merged order]
\label{4.23:prf-halves-merged-order}
Fix a step $k$. In the tube statement for a treated large-field component $C$
(Theorem~\ref{4.23:thm-tube-statement}), the \emph{form half} is the form theorem
(Theorem~\ref{4.23:thm-form-theorem}) together with clauses (q0) and (q1) of
\eqref{4.23:eq-tube-statement-a}. The \emph{oscillation half} is the oscillation theorem for the
resummed density (Theorem~\ref{4.23:thm-oscillation-theorem}). Write
\begin{equation}\label{4.23:eq-halves-merged-order-1}
\mathsf{n}_{\mathrm{sh}}(k) = (N+1)\,(110\,L\,\check{R}_{h+1})^{d}
\end{equation}
for the shallow-shell count, where $N=N(g_k)$ is the number of levels of the window at step $k$
(not the integer $N$ of the gauge group $\mathrm{SU}(N)$) and $h=k-N$ is its bottom level, and write
$\mathfrak{b}'_k = A_1(\log x_k)^{p_1}$. Then the following hold.
\begin{enumerate}
\item[(a)] Clause (q1) enters the oscillation half in three places:
  \begin{itemize}
  \item inside clause (t) of \eqref{4.23:eq-oscillation-theorem-a}, through the total bound of
    that clause;
  \item in the first step of the nested-region lemma (Lemma~\ref{4.23:lem-nested-regions}),
    at the earlier steps $m'<k$;
  \item in clause (i) of the collar-payment corollary
    (Corollary~\ref{4.23:cor-collar-payment}), at the step at which the region was created.
  \end{itemize}
  Clause (q0) enters in the last two of these places.
\item[(b)] The form half uses no conclusion of the oscillation half, at any step.
\item[(c)] The two halves are proved without circularity. The dependences among their statements
  are well-founded in the pair (step, position of the statement within the step).
\item[(d)] The inputs common to the two halves carry one and the same meaning in both. In particular
  this holds for the count \eqref{4.23:eq-halves-merged-order-1} and for $\mathfrak{b}'_k$.
\item[(e)] The constants are chosen in the following order:
\begin{equation}\label{4.23:eq-halves-merged-order-2}
\begin{aligned}
&L \to \mathfrak{M} \to \kappa \to \kappa_1 \to M_1 \to M \to \dots \to A_0, A_1 \to \dots
  \to E_0 \to \dots \to B_0 \to \dots \\
&\qquad \to c_A \to \mathsf{K}^0 \to C_\Sigma,\ p_\Sigma \to \mathfrak{n}^{\mathrm{kept}}(\beta)
  \to \gamma .
\end{aligned}
\end{equation}
  Each constant listed in \eqref{4.23:eq-halves-merged-order-2} and below is fixed at the place
  indicated either in \eqref{4.23:eq-constant-order-a} or in the order of choice of the form
  theorem (Theorem~\ref{4.23:thm-form-theorem}).
  Here the following constants are fixed together with $\kappa_1$, among them the norm
  bound $K_{\mathsf{T}}$ of $\mathsf{T}^{\mathsf{L}}$ in clause (q1) of
  \eqref{4.23:eq-tube-statement-a}:
  \begin{equation*}
  K_\sharp,\ K_H,\ K_{\mathsf{T}},\ K_\natural,\ \gamma^\flat,\ \varepsilon_\sharp^\flat,\ 
  K_{\mathfrak{s}},\ K_1^0,\ \rho_{\mathrm{abs}},\ \vartheta^0,\ C_{\mathfrak{q}} .
  \end{equation*}
  The following constants are already fixed as soon as $d$ and $L$ are:
  \begin{equation*}
  c_Q,\ K_u,\ c_4,\ \varepsilon_{\alpha},\ K_J,\ c_{\mathrm{tr}},\ c_N,\ c_R .
  \end{equation*}
  The following are fixed together with $M$:
  \begin{equation*}
  K_\Lambda,\ \vartheta_X,\ r_{\mathbb{R}},\ c_I M^5,\ C_{\mathrm{cell}},\ M/M_1 .
  \end{equation*}
  The constants $\mathsf{K}^\flat$, $\mathsf{K}^{\mathrm{cp}}$, $C^{\mathrm{sh}}$,
  $c_{\mathsf{b}}$, $c_{\mathsf{b}}'$ are fixed together with $\mathsf{K}^0$. No choice returns to
  $L,\mathfrak{M},\kappa,\kappa_1,M_1,M$. The constant $C_\Sigma$ does not depend on
  $\mathfrak{n}^{\mathrm{kept}}(\beta)$, and $\gamma$ is chosen last, against all ceilings.
\end{enumerate}
\end{lemma}

\begin{proof}
For a kept member $\mathfrak{q}$ we write $\mathfrak{O}(\mathfrak{q})$ for the quantity, taken over
the tube, by which the nested-region lemma and the collar-payment corollary measure $\mathfrak{q}$.
We write $\Sigma^{\mathrm{tot}}(k)$ for the total bound of clause (t) of the oscillation theorem,
\eqref{4.23:eq-oscillation-theorem-a}.

\emph{(a).} Clause (q1) is the bound of the form theorem, evaluated at the linearised localised
minimiser map $\mathsf{T}^{\mathsf{L}}$ and restricted from $\mathfrak{S}_{\mathrm{src}}$ to the
tube $\mathfrak{S}\subset\mathfrak{S}_{\mathrm{src}}$. It enters the oscillation half in three
places.

First, it enters $\Sigma^{\mathrm{tot}}(k)$ as the member $2\mathfrak{q}_{\max}(k)$ of clause (t).
This member is the last entry of the table \eqref{4.23:eq-degree-constant-a}, where it is not
estimated but taken as an input.

Second, it enters the first step of the nested-region lemma. That step applies the bound on the
second kept member at each earlier step $m'<k$. For each large-field region $Y'\subset C$ created at
step $m'$, the kept member $\mathfrak{q}_{Y'}$ born with it satisfies
$\mathfrak{O}(\mathfrak{q}_{Y'}) \le 2\mathfrak{q}_{\max}(m')$, and this bound is obtained from
clause (q1) at step $m'$.

Third, it enters clause (i) of the collar-payment corollary, applied at the step $j$ at which the
region was created.

Clause (q0), fixed in \eqref{4.23:eq-tube-statement-a}, supplies in particular the
splitting-freeness, the holomorphy on the tube and the joint gauge invariance of the second kept
member. It is used in the first step of the nested-region lemma and in
clause (i) of the collar-payment corollary. The reason is that one step in the proof of the bound
$\mathfrak{O}(\mathfrak{q}_{Y'}) \le 2\mathfrak{q}_{\max}(m')$ requires the joint gauge invariance
of $\mathfrak{q}_{Y'}$, and its holomorphy, on the tube,
which is invariant under gauge transformations. Splitting-freeness is needed for a more basic
reason: without it, $\mathfrak{q}_{Y'}$ would not be a function on the tube at all.

\emph{(b).} At step $k$ the form theorem is proved, at its place in the step
(Remark~\ref{4.23:rem-tube-placement}), from three kinds of input: the base lemmas established
before the induction, one further standing input of the form theorem, and the geometry of step $k$.
It uses none of the following: a statement of the expansion whose parameters and floors are assumed
in the oscillation theorem (Theorem~\ref{4.23:thm-oscillation-theorem}), or one of the nine
auxiliary statements assumed there as a numbered list; the tube statement at an earlier step; or an
item of the containment argument.

The leaf count, from which $\mathfrak{n}^{\mathrm{kept}}(\beta)$ is formed
(Lemma~\ref{4.23:lem-constant-order}), uses only the bound
$\mathfrak{O}(\mathfrak{q}_C)\le 2\mathfrak{q}_{\max}$. It uses that bound through
$\Sigma^{\mathrm{tot}}$ of the oscillation theorem, whose constant $C_\Sigma$ does not depend on
$\mathfrak{n}^{\mathrm{kept}}(\beta)$ (Lemma~\ref{4.23:lem-constant-order}). The dependence
therefore runs from the form half to the oscillation half, and not conversely.

\emph{(c).} At step $k$ the form theorem uses no quantity produced by the induction over steps. Its
inputs are the base lemmas, the further standing input just named, the hypothesis listed sixth in
the statement of the form theorem, and the shallow-shell count of step $k$, which is the hypothesis
listed seventh there.

Clause (v) of the oscillation theorem, in \eqref{4.23:eq-oscillation-theorem-a}, at step $k$
uses two kinds of input. It uses (q0) and (q1)
at steps $\le k$. It uses the conclusion of the oscillation theorem, together with its
clause (g) in \eqref{4.23:eq-oscillation-theorem-a}, at steps $m'$ with $m'\le h<k$ only. Indeed,
the bound $\mathfrak{O}(\mathfrak{q}_{Y'}) \le 2\mathfrak{q}_{\max}(m')$ of part (a) uses the
oscillation theorem at its own step $m'$ only, where $m'<k$. Conversely, the oscillation theorem
at step $k$ uses that bound only for $m'\le h<k$.

Clause (t) uses no induction hypothesis. The constant $C_\Sigma$ is defined before the induction,
without the nested-region lemma, and it does not depend on $\mathfrak{n}^{\mathrm{kept}}(\beta)$;
this is Lemma~\ref{4.23:lem-constant-order}.

Clause (g) at step $k$ is proved from three inputs: the term-by-term gauge invariance of the
expansion, clause (q0) at earlier steps, and one further clause among the standing inputs of the
oscillation theorem. Hence every statement of either half uses only statements that precede it.
The scheme is therefore well-founded in the pair (step, position within the step); this is the
induction by which the oscillation theorem is proved.

\emph{(d).} The two halves use the following inputs with one meaning.
\begin{itemize}
\item The nesting geometry and the shallow-shell count, with the single count
  \eqref{4.23:eq-halves-merged-order-1}.
\item The size $\mathfrak{b}'_k$.
\item The third of the five inputs that the form theorem cites by name among its hypotheses. The
  form half uses it as the monotonicity of the hull. The oscillation half uses it in ratio form,
  inside the lemmas from which it takes its inputs. Neither half pays through it the factor
  $\check{R}_j^{p_\Sigma}$ attached to a region created at step $j$.
\item The family of base estimates fixed before the induction, among them the hypothesis listed
  fourth in the statement of the form theorem.
\item The tube, fixed once for each domain. The form half works on $\mathfrak{S}_{\mathrm{src}}$;
  the oscillation half works on its restriction $\mathfrak{S}$.
\item The bound $\mathfrak{q}_{\max}$, identical in every symbol.
\item The convention, common to both halves, by which the fifth step of the construction of the
  second kept member is taken to be the map $\mathsf{T}^{\mathsf{L}}$; clause (q1) is evaluated at
  this map in both halves.
\end{itemize}

\emph{(e).} The order of choice of the oscillation half, recorded in
\eqref{4.23:eq-constant-order-a}, is merged with that of the form half. The result is
\eqref{4.23:eq-halves-merged-order-2}.

The following are functions, not entries of the order; they are formed after $\gamma$ is chosen:
\begin{itemize}
\item the functions $\mathsf{n}_\sharp$ and $\mathsf{n}_\partial$, the latter entering
  \eqref{4.23:eq-nested-regions-b}, and the shallow-shell count
  \eqref{4.23:eq-halves-merged-order-1};
\item the collar sums $\mathsf{S}^{\mathfrak{K}}(k)$ of \eqref{4.23:eq-collar-payment-a},
  together with the other sums of the same family;
\item the bound $\mathfrak{q}_{\max}$ of \eqref{4.23:eq-tube-statement-a}, and the quantities
  $\bar\alpha_1$, $\bar{\mathfrak{q}}_{\max}$ and $\Sigma^{\mathfrak{q}}$;
\item the bound $\mathrm{nest}^*$ of \eqref{4.23:eq-nested-regions-b};
\item the sums $\Sigma^{V}$ and $\Sigma^{\mathrm{tot}}$, the latter the total bound of clause (t);
\item the number $\varepsilon_\gamma$ of \eqref{4.23:eq-nested-regions-a}.
\end{itemize}
Here $\varepsilon_\gamma$ of \eqref{4.23:eq-nested-regions-a} is a number once the quantity
$\mathsf{T}_\gamma$, through which $\varepsilon_\gamma$ depends on $\gamma$, is fixed by the ceilings
on $\gamma$. It enters only the nesting ceiling of
\eqref{4.23:eq-degree-constant-a} and the bound $\mathrm{nest}^*$ of
\eqref{4.23:eq-nested-regions-b}.

No choice returns to $L,\mathfrak{M},\kappa,\kappa_1,M_1,M$. By (c), $C_\Sigma$ does not depend on
$\mathfrak{n}^{\mathrm{kept}}(\beta)$. Finally, $\gamma$ is chosen last, against all ceilings.
\end{proof}

\begin{proof}[Proof of Theorem~\ref{4.23:thm-tube-statement}]
Assume the hypotheses of Theorem~\ref{4.23:thm-tube-statement}; they comprise those of the form
theorem (Theorem~\ref{4.23:thm-form-theorem}) and those of the oscillation theorem
(Theorem~\ref{4.23:thm-oscillation-theorem}), each at its place in the step.
By its definition in \eqref{4.23:eq-tube-statement-a}, the tube statement at step $k$ is the
conjunction of four clauses: the form clause, which consists of clauses (X2) and (X3) of the form
theorem in \eqref{4.23:eq-form-theorem-a}; clause (q0), which is clause (X0) of the form theorem
together with the joint gauge-invariance clause of the form half; clause (q1), which is clause (X1)
of the form theorem at $\mathsf{T}=\mathsf{T}^{\mathsf{L}}$, restricted from
$\mathfrak{S}_{\mathrm{src}}$ to the tube $\mathfrak{S}\subset\mathfrak{S}_{\mathrm{src}}$; and the
oscillation clause, which consists of the conclusions ($\sharp$), (v), (t), (b), (d), (g), (o),
(n) of the oscillation theorem in \eqref{4.23:eq-oscillation-theorem-a}.

We proceed by the induction, well-founded in the pair (step, position of the statement within the
step), of part (c) of Lemma~\ref{4.23:prf-halves-merged-order}. Let $k$ be a step, and assume the
four clauses at every step $m'<k$. By parts (b) and (c) of
Lemma~\ref{4.23:prf-halves-merged-order}, the form theorem at step $k$ uses no conclusion of the
oscillation half and no quantity produced by the induction over steps;
Theorem~\ref{4.23:thm-form-theorem} therefore applies at step $k$ without any induction
hypothesis, and gives clauses (X0)--(X3) at step $k$. Together with the joint gauge-invariance
clause of the form half, this yields the form clause, clause (q0) and clause (q1) at step $k$. By
parts (a) and (c) of the same lemma, the oscillation theorem at step $k$ uses clauses (q0) and (q1)
at steps $\le k$, which are now available; its own conclusion and its clause (g) at steps $m'$ with
$m'\le h<k$ only, which are available by the induction assumption; and clause (t), which uses no
induction hypothesis. Theorem~\ref{4.23:thm-oscillation-theorem} therefore applies at step $k$,
and gives the oscillation clause at step $k$. By part
(d) of Lemma~\ref{4.23:prf-halves-merged-order}, the inputs common to the two halves carry one
meaning in both, so the four clauses concern the same objects; by part (e), the constants in them
are chosen in one order in which no choice returns. Hence the tube statement holds at step $k$, and
by induction at every step.
\end{proof}

\begin{definition}[Colour fibres, norms and the adjoint trace]\label{4.24:not-colour-norms}
The following notation is fixed for the rest of this chapter.

\emph{Groups and algebras.} Let $N\ge 2$, let $G:=SU(N)\subset M_N(\mathbb{C})$ and
$G^{\mathbb{C}}:=SL(N,\mathbb{C})$, and let $\mathfrak{g}:=\mathfrak{su}(N)$ and
$\mathfrak{g}^{\mathbb{C}}:=\mathfrak{sl}(N,\mathbb{C})$. We write $n_{\mathfrak{g}}:=N^2-1$,
which is the complex dimension of $\mathfrak{g}^{\mathbb{C}}$. For $g\in G^{\mathbb{C}}$ and
$X\in\mathfrak{g}^{\mathbb{C}}$ we put $\operatorname{Ad}_gX:=gXg^{-1}$.

\emph{Norms on matrices.} The symbol $\|\cdot\|$ denotes the operator norm of a matrix acting on
$\mathbb{C}^N$ with its standard Hermitian norm; it is sub-multiplicative. For
$v\in\mathbb{C}^N$, $|v|$ is its Hermitian length, which is the Hilbert--Schmidt norm below of
$v$ read as an $N\times1$ matrix. We write
$\operatorname{Tr}$ for the ordinary trace on $M_N(\mathbb{C})$, so that
$\operatorname{Tr}\mathbf{1}=N$ and $\operatorname{Tr}=N\operatorname{tr}$. On
$M_N(\mathbb{C})$ we use the Hilbert--Schmidt norm
$|X|:=(\operatorname{Tr}X^*X)^{1/2}$. It satisfies, for all $X,A,Y\in M_N(\mathbb{C})$,
\begin{equation}\label{4.24:eq-colour-norms-a}
|XAY|\le\|X\|\,|A|\,\|Y\|,\qquad \|A\|\le|A|,
\end{equation}
and it is invariant under conjugation by $G$: $|gXg^{-1}|=|X|$ for $g\in G$. Every
smallness quantity introduced from here on that measures a matrix is defined with this norm
$|\cdot|$, so that nothing below depends on the choice of $|\cdot|$ beyond
\eqref{4.24:eq-colour-norms-a} and its invariance under conjugation by $G$. For
$g\in G^{\mathbb{C}}$ we put
\begin{equation*}
\operatorname{pol}(g):=\log\max\{\|g\|,\|g^{-1}\|\};
\end{equation*}
then $\operatorname{pol}\equiv0$ on $G$.

\emph{Inner products on the fibre.} On $\mathfrak{g}$ let $(X,Y):=-\operatorname{Tr}(XY)$.
This is a positive, $\operatorname{Ad}$-invariant inner product, and we extend it
$\mathbb{C}$-bilinearly to $\mathfrak{g}^{\mathbb{C}}$. Since
$\mathfrak{g}^{\mathbb{C}}=\mathfrak{g}\oplus i\mathfrak{g}$, every
$X\in\mathfrak{g}^{\mathbb{C}}$ is $X=A+iB$ with $A,B\in\mathfrak{g}$, and we write
$\bar X:=A-iB$ for the conjugation of $\mathfrak{g}^{\mathbb{C}}$ with respect to
$\mathfrak{g}$. We put $\langle X,Y\rangle:=(\bar X,Y)$; this is a positive Hermitian form,
conjugate-linear in the first argument, and $\langle X,X\rangle=|X|^2$.

\emph{Endomorphisms of the fibre.} For $T\in\operatorname{End}(\mathfrak{g}^{\mathbb{C}})$,
$\|T\|_{e}$ is the operator norm of $T$ on $(\mathfrak{g}^{\mathbb{C}},|\cdot|)$, which is
its operator norm with respect to $\langle\cdot,\cdot\rangle$. We write
$\operatorname{tr}_{\mathfrak{g}}$ for the trace on $\operatorname{End}(\mathfrak{g}^{\mathbb{C}})$. Then
\begin{equation}\label{4.24:eq-colour-norms-1}
|\operatorname{tr}_{\mathfrak{g}}T|\le n_{\mathfrak{g}}\,\|T\|_{e}.
\end{equation}

\emph{The bond fibre space and kernels.} Let $\mathfrak{Bo}$ be the set of bonds of the
lattice on which the step is taken, and for a bond $b$ let $b_{-}$, $b_{+}$ be its initial
and final endpoints. We put
$\mathcal{H}:=\bigoplus_{b\in\mathfrak{Bo}}\mathfrak{g}^{\mathbb{C}}_b$, each
$\mathfrak{g}^{\mathbb{C}}_b$ a copy of $\mathfrak{g}^{\mathbb{C}}$, with the inner
product
\begin{equation*}
\langle\varphi,\psi\rangle:=\sum_{b}\langle\varphi(b),\psi(b)\rangle .
\end{equation*}
A kernel $K=(K(b,b'))_{b,b'\in\mathfrak{Bo}}$ with
$K(b,b')\in\operatorname{End}(\mathfrak{g}^{\mathbb{C}})$ acts by
$(K\varphi)(b):=\sum_{b'}K(b,b')\varphi(b')$. For $E,E'\subset\mathfrak{Bo}$, $P_E$ is the
orthogonal projection of $\mathcal{H}$ onto $\bigoplus_{b\in E}\mathfrak{g}^{\mathbb{C}}_b$,
that is, multiplication by the indicator of $E$, and $K_{EE'}:=P_EKP_{E'}$. A kernel $K$ is an
operator on $\mathcal{H}$, and the letter $K$ stands for a kernel only where it is so declared;
elsewhere $K$ is the number of renormalisation steps. For a kernel
$K$, $\|K\|$ is its operator norm on $(\mathcal{H},\langle\cdot,\cdot\rangle)$; the argument
of $\|\cdot\|$ (a matrix or a kernel) determines which operator norm is meant.

\emph{The gauge action.} For a $G$-valued function $u$ on the sites put
$\mathcal{U}_u:=\bigoplus_{b}\operatorname{Ad}_{u(b_{-})}$, that is,
$(\mathcal{U}_u\varphi)(b)=\operatorname{Ad}_{u(b_{-})}\varphi(b)$. Each
$\operatorname{Ad}_{u(x)}$ is $(\cdot,\cdot)$-orthogonal and commutes with conjugation,
hence is $\langle\cdot,\cdot\rangle$-unitary; $\mathcal{U}_u$ is unitary and
bond-diagonal, and $P_E\mathcal{U}_u=\mathcal{U}_uP_E$ for every $E\subset\mathfrak{Bo}$.

\emph{Real points.} For a pair $(\mathbb{U},\mathbb{J})$ of a bond field $\mathbb{U}$ and a
source form $\mathbb{J}$, the pair is a \emph{real point} if $\mathbb{U}$ is $G$-valued
and $\mathbb{J}$ is $\mathfrak{g}$-valued; $(\mathbf{1},0)$ is the \emph{trivial pair}.

Nothing in what follows is specialised to a particular value of $N$, and the properties of
the invariant form and the $\operatorname{Ad}$-orthogonality used below hold for every compact
group with an $\operatorname{Ad}$-invariant inner product.
\end{definition}

\begin{proof}
\emph{The inequalities \eqref{4.24:eq-colour-norms-a}.} Write $a_1,\dots,a_N$ for the
columns of $A$. The columns of $XA$ are $Xa_j$, so
\begin{equation*}
|XA|^2=\sum_j|Xa_j|^2\le\|X\|^2\sum_j|a_j|^2=\|X\|^2|A|^2 .
\end{equation*}
Since
$|Z^*|=|Z|$ and $\|Z^*\|=\|Z\|$, applying this to $(AY)^*=Y^*A^*$ gives
$|AY|\le\|Y\|\,|A|$; combining the two bounds gives $|XAY|\le\|X\|\,|A|\,\|Y\|$. For a unit
vector $v$, the Cauchy--Schwarz inequality applied to each component
$(Av)_i=\sum_jA_{ij}v_j$ gives
\begin{equation*}
|Av|^2=\sum_i\Bigl|\sum_jA_{ij}v_j\Bigr|^2\le\sum_i\sum_j|A_{ij}|^2\,|v|^2=|A|^2,
\end{equation*}
so $\|A\|\le|A|$.

\emph{Invariance and $\operatorname{pol}$ on $G$.} If $g\in G$ then $g$ is unitary, so
$\|g\|=\|g^{-1}\|=1$. Hence $\operatorname{pol}(g)=\log1=0$, and
\eqref{4.24:eq-colour-norms-a} gives $|gXg^{-1}|\le|X|$ and
$|X|=|g^{-1}(gXg^{-1})g|\le|gXg^{-1}|$.

\emph{The inner products.} For $g\in G^{\mathbb{C}}$,
$\operatorname{Tr}(gXg^{-1})=\operatorname{Tr}X$, so $\operatorname{Ad}_g$ maps
$\mathfrak{g}^{\mathbb{C}}$ into itself. For $X\in\mathfrak{g}$ we have $X^*=-X$, so
$(X,X)=\operatorname{Tr}(X^*X)=|X|^2$, which is positive for $X\ne0$; and
$(\operatorname{Ad}_gX,\operatorname{Ad}_gY)=-\operatorname{Tr}(gXYg^{-1})=(X,Y)$. If $g\in G$
and $X\in\mathfrak{g}$, then $(gXg^{-1})^*=gX^*g^{-1}=-gXg^{-1}$, so $\operatorname{Ad}_g$
preserves $\mathfrak{g}$. Every traceless $X$ is $A+iB$ with $A=(X-X^*)/2$ and
$B=(X+X^*)/(2i)$, both in $\mathfrak{g}$, and $\mathfrak{g}\cap i\mathfrak{g}=\{0\}$, so
$\mathfrak{g}^{\mathbb{C}}=\mathfrak{g}\oplus i\mathfrak{g}$ and $\bar X$ is well defined.
For $X=A+iB$ one has $X^*=-A+iB=-\bar X$, so by bilinearity
\begin{equation*}
\langle X,Y\rangle=-\operatorname{Tr}(\bar XY)=\operatorname{Tr}(X^*Y),
\end{equation*}
which is Hermitian and positive, with $\langle X,X\rangle=|X|^2$. In particular the
operator norm on $(\mathfrak{g}^{\mathbb{C}},|\cdot|)$ is the operator norm with respect
to $\langle\cdot,\cdot\rangle$.

\emph{The gauge action.} Let $g\in G$. As shown, $\operatorname{Ad}_g$ is
$(\cdot,\cdot)$-orthogonal and preserves $\mathfrak{g}$; being $\mathbb{C}$-linear, it maps
$A+iB$ to $\operatorname{Ad}_gA+i\operatorname{Ad}_gB$ with both terms in $\mathfrak{g}$,
hence $\overline{\operatorname{Ad}_gX}=\operatorname{Ad}_g\bar X$. Therefore
$\langle\operatorname{Ad}_gX,\operatorname{Ad}_gY\rangle
=(\operatorname{Ad}_g\bar X,\operatorname{Ad}_gY)=\langle X,Y\rangle$, and
$\operatorname{Ad}_g$ is invertible with inverse $\operatorname{Ad}_{g^{-1}}$, so it is
unitary. Consequently, for $\varphi,\psi\in\mathcal{H}$,
\begin{align*}
\langle\mathcal{U}_u\varphi,\mathcal{U}_u\psi\rangle
&=\sum_b\langle\operatorname{Ad}_{u(b_{-})}\varphi(b),\operatorname{Ad}_{u(b_{-})}\psi(b)\rangle\\
&=\langle\varphi,\psi\rangle,
\end{align*}
and $\mathcal{U}_u$ is invertible with inverse $\mathcal{U}_{u^{-1}}$; thus $\mathcal{U}_u$
is unitary. Since $(\mathcal{U}_u\varphi)(b)$ depends on $\varphi(b)$ only, $\mathcal{U}_u$
is bond-diagonal, and as $P_E$ is multiplication by the indicator of $E$, both
$P_E\mathcal{U}_u\varphi$ and $\mathcal{U}_uP_E\varphi$ equal
$\operatorname{Ad}_{u(b_{-})}\varphi(b)$ at $b\in E$ and $0$ at $b\notin E$.

\emph{The trace bound \eqref{4.24:eq-colour-norms-1}.} Let $e_1,\dots,e_{n_{\mathfrak{g}}}$ be
an orthonormal basis of $(\mathfrak{g}^{\mathbb{C}},\langle\cdot,\cdot\rangle)$. Then
$\operatorname{tr}_{\mathfrak{g}}T=\sum_i\langle e_i,Te_i\rangle$, and by the
Cauchy--Schwarz inequality $|\langle e_i,Te_i\rangle|\le|e_i|\,|Te_i|\le\|T\|_{e}$ for each
$i$. Summing over the $n_{\mathfrak{g}}$ indices gives \eqref{4.24:eq-colour-norms-1}.
\end{proof}

\begin{lemma}[Norm of the adjoint action]\label{4.24:lem-adjoint-norm}
Let $G$, $G^{\mathbb{C}}$, $\mathfrak{g}^{\mathbb{C}}$, the norms $\|\cdot\|$, $|\cdot|$,
$\|\cdot\|_{e}$, the function $\operatorname{pol}$, the map
$\operatorname{Ad}_gX=gXg^{-1}$ and the inner product $\langle\cdot,\cdot\rangle$ on
$\mathfrak{g}^{\mathbb{C}}$, whose norm is $|\cdot|$, be as in
Definition~\ref{4.24:not-colour-norms}.
\begin{enumerate}
\item[(i)] For $g\in G$ and every $X\in M_N(\mathbb{C})$ one has
$|\operatorname{Ad}_gX|=|X|$; on $\mathfrak{g}^{\mathbb{C}}$ the map $\operatorname{Ad}_g$ is
unitary for $\langle\cdot,\cdot\rangle$; and $\|\operatorname{Ad}_g\|_{e}=1$.
\item[(ii)] For $g\in G^{\mathbb{C}}$ one has
$\|\operatorname{Ad}_g\|_{e}\le e^{2\operatorname{pol} g}$.
\end{enumerate}
\end{lemma}

\begin{proof}
For every $g\in G^{\mathbb{C}}$ the map $\operatorname{Ad}_g$ is linear and maps
$\mathfrak{g}^{\mathbb{C}}$ into itself, since
$\operatorname{tr}(gXg^{-1})=\operatorname{tr}X$; it is invertible with inverse
$\operatorname{Ad}_{g^{-1}}$.

(i) Let $g\in G$. The norm $|\cdot|$ is invariant under conjugation by $G$
(Definition~\ref{4.24:not-colour-norms}), so $|\operatorname{Ad}_gX|=|gXg^{-1}|=|X|$
for every $X\in M_N(\mathbb{C})$.
On $\mathfrak{g}^{\mathbb{C}}$ the norm $|\cdot|$ is the norm of the inner product
$\langle\cdot,\cdot\rangle$ (Definition~\ref{4.24:not-colour-norms}). By the polarisation
identity, $\langle X,Y\rangle$ is a fixed
linear combination of the four numbers $|X+i^{q}Y|^2$, $q=0,1,2,3$. The map
$\operatorname{Ad}_g$ is linear and preserves $|\cdot|$, so it preserves each of these
four numbers, and therefore
$\langle\operatorname{Ad}_gX,\operatorname{Ad}_gY\rangle=\langle X,Y\rangle$. Being also
invertible, $\operatorname{Ad}_g$ is unitary. Since $\mathfrak{g}^{\mathbb{C}}\neq\{0\}$, the
operator norm of an isometry of $(\mathfrak{g}^{\mathbb{C}},|\cdot|)$ equals $1$, that is,
$\|\operatorname{Ad}_g\|_{e}=1$.

(ii) Let $g\in G^{\mathbb{C}}$ and $X\in\mathfrak{g}^{\mathbb{C}}$. By the definition of
$\operatorname{pol}$,
\begin{equation*}
\max\{\|g\|,\|g^{-1}\|\}=e^{\operatorname{pol} g}.
\end{equation*}
Hence the norm property \eqref{4.24:eq-colour-norms-a}, applied with $g$ on the left and
$g^{-1}$ on the right, gives
\begin{equation*}
|\operatorname{Ad}_gX|=|gXg^{-1}|\le\|g\|\,|X|\,\|g^{-1}\|
\le e^{2\operatorname{pol} g}\,|X|.
\end{equation*}
Taking the supremum over $X\in\mathfrak{g}^{\mathbb{C}}$ with $|X|\le1$ yields
$\|\operatorname{Ad}_g\|_{e}\le e^{2\operatorname{pol} g}$.
\end{proof}

\begin{lemma}[A subadditive distance to the identity in the group]\label{4.24:lem-group-distance}
Let $G=SU(N)$, $G^{\mathbb{C}}=SL(N,\mathbb{C})$, $\mathfrak{g}^{\mathbb{C}}=\mathfrak{sl}(N,\mathbb{C})$, the norms $\|\cdot\|$, $|\cdot|$, $\|\cdot\|_{e}$ and the function $\operatorname{pol}$ be as in Definition~\ref{4.24:not-colour-norms}. For $h\in G^{\mathbb{C}}$ put
\begin{equation*}
\operatorname{dis}(h):=\inf\Bigl\{\sum_{i=1}^{r}|Z_i|\;:\;h=e^{Z_1}\cdots e^{Z_r},\ Z_i\in\mathfrak{g}^{\mathbb{C}},\ r\ge 1\Bigr\}.
\end{equation*}
Then for all $h,h_1,h_2\in G^{\mathbb{C}}$ and $g\in G^{\mathbb{C}}$, writing $1$ both for the identity matrix and for the identity of $\operatorname{End}(\mathfrak{g}^{\mathbb{C}})$,
\begin{equation}\label{4.24:eq-group-distance-a}
\begin{aligned}
&\text{(d1)}\quad \|h-1\|\le e^{\operatorname{dis}h}-1,\ \text{ hence }\ \operatorname{dis}h=0\iff h=1;\\
&\text{(d2)}\quad \operatorname{dis}(h_1h_2)\le\operatorname{dis}h_1+\operatorname{dis}h_2,\qquad
\operatorname{dis}(h^{-1})=\operatorname{dis}h,\\
&\phantom{\text{(d2)}}\quad \operatorname{dis}(ghg^{-1})\le e^{2\operatorname{pol}g}\operatorname{dis}h,\ \text{ with }\
\operatorname{dis}(ghg^{-1})=\operatorname{dis}h\ \text{ if } g\in G,\\
&\phantom{\text{(d2)}}\quad \operatorname{dis}h\le|\log h|\ \text{ for every logarithm of } h;\\
&\text{(d3)}\quad \operatorname{pol}h\le\operatorname{dis}h;\\
&\text{(d4)}\quad \|\operatorname{Ad}_h-1\|_{e}\le e^{2\operatorname{dis}h}-1;\\
&\text{(d5)}\quad \text{if } h\in G:\ \ \|\operatorname{Ad}_h-1\|_{e}\le\min\{2,\,e^{2\operatorname{dis}h}-1\}\le 6\operatorname{dis}h.
\end{aligned}
\end{equation}
In (d2) a logarithm of $h$ is an element $\log h\in\mathfrak{g}^{\mathbb{C}}$ with $e^{\log h}=h$.
\end{lemma}

\begin{proof}
Throughout, a \emph{presentation} of $h$ is a finite sequence $Z_1,\dots,Z_r\in\mathfrak{g}^{\mathbb{C}}$, $r\ge1$, with $h=e^{Z_1}\cdots e^{Z_r}$; its length is $\sum_i|Z_i|$, and $\operatorname{dis}h$ is the infimum of the lengths of the presentations of $h$. By \eqref{4.24:eq-colour-norms-a} we have $\|A\|\le|A|$, and taking $Y=1$, resp.\ $X=1$, there, $|XA|\le\|X\|\,|A|$ and $|AY|\le|A|\,\|Y\|$ for all matrices.

\emph{(d1).} For $Z\in\mathfrak{g}^{\mathbb{C}}$ the exponential series and the sub-multiplicativity of $\|\cdot\|$ give
\begin{equation*}
\|e^{Z}-1\|=\Bigl\|\sum_{n\ge1}\frac{Z^n}{n!}\Bigr\|\le\sum_{n\ge1}\frac{\|Z\|^n}{n!}=e^{\|Z\|}-1\le e^{|Z|}-1.
\end{equation*}
For matrices $A,B$ one has $AB-1=(A-1)(B-1)+(A-1)+(B-1)$, hence
\begin{equation*}
\|AB-1\|\le\|A-1\|\,\|B-1\|+\|A-1\|+\|B-1\|=(1+\|A-1\|)(1+\|B-1\|)-1.
\end{equation*}
Let $Z_1,\dots,Z_r$ be a presentation of $h$. Applying the last inequality with $A=e^{Z_1}\cdots e^{Z_{j}}$ and $B=e^{Z_{j+1}}$, and inducting on $j$, we get
\begin{equation*}
\|h-1\|\le\prod_{i=1}^{r}\bigl(1+\|e^{Z_i}-1\|\bigr)-1\le\prod_{i=1}^{r}e^{|Z_i|}-1=e^{\sum_i|Z_i|}-1.
\end{equation*}
Since $x\mapsto e^{x}-1$ is continuous and increasing, the infimum over presentations gives $\|h-1\|\le e^{\operatorname{dis}h}-1$. If $\operatorname{dis}h=0$, then $\|h-1\|\le0$, so $h=1$; conversely $1=e^{0}$ is a presentation of length $0$, so $\operatorname{dis}1=0$.

\emph{(d2).} If $Z_1,\dots,Z_r$ presents $h_1$ and $W_1,\dots,W_s$ presents $h_2$, the concatenation $Z_1,\dots,Z_r,W_1,\dots,W_s$ presents $h_1h_2$, with length the sum of the two lengths; taking the infimum over both presentations gives $\operatorname{dis}(h_1h_2)\le\operatorname{dis}h_1+\operatorname{dis}h_2$. If $Z_1,\dots,Z_r$ presents $h$, then $h^{-1}=e^{-Z_r}\cdots e^{-Z_1}$, and $|-Z_i|=|Z_i|$; hence $\operatorname{dis}(h^{-1})\le\operatorname{dis}h$, and the same inequality applied to $h^{-1}$ gives the reverse one. For $g\in G^{\mathbb{C}}$ and $Z\in\mathfrak{g}^{\mathbb{C}}$, $ge^{Z}g^{-1}=e^{gZg^{-1}}=e^{\operatorname{Ad}_gZ}$, and $\operatorname{Ad}_gZ$ is traceless, so lies in $\mathfrak{g}^{\mathbb{C}}$. Thus a presentation $Z_1,\dots,Z_r$ of $h$ yields the presentation $\operatorname{Ad}_gZ_1,\dots,\operatorname{Ad}_gZ_r$ of $ghg^{-1}$, and by Lemma~\ref{4.24:lem-adjoint-norm}(ii), $|\operatorname{Ad}_gZ_i|\le\|\operatorname{Ad}_g\|_{e}|Z_i|\le e^{2\operatorname{pol}g}|Z_i|$. Taking the infimum, $\operatorname{dis}(ghg^{-1})\le e^{2\operatorname{pol}g}\operatorname{dis}h$. If $g\in G$, Lemma~\ref{4.24:lem-adjoint-norm}(i) gives $|\operatorname{Ad}_gZ_i|=|Z_i|$, so $\operatorname{dis}(ghg^{-1})\le\operatorname{dis}h$; applying this with $g^{-1}\in G$ in place of $g$ and $ghg^{-1}$ in place of $h$ gives $\operatorname{dis}h\le\operatorname{dis}(ghg^{-1})$, hence equality. Finally, if $\log h\in\mathfrak{g}^{\mathbb{C}}$ satisfies $e^{\log h}=h$, then the one-term sequence $\log h$ is a presentation of $h$, so $\operatorname{dis}h\le|\log h|$.

\emph{(d3).} For $Z\in\mathfrak{g}^{\mathbb{C}}$, $\|e^{\pm Z}\|\le e^{\|Z\|}\le e^{|Z|}$, so $\operatorname{pol}(e^{Z})=\log\max\{\|e^{Z}\|,\|e^{-Z}\|\}\le|Z|$. The function $\operatorname{pol}$ is subadditive: for $A,B\in G^{\mathbb{C}}$,
\begin{equation*}
\|AB\|\le\|A\|\,\|B\|\le e^{\operatorname{pol}A+\operatorname{pol}B},\qquad
\|(AB)^{-1}\|=\|B^{-1}A^{-1}\|\le e^{\operatorname{pol}A+\operatorname{pol}B},
\end{equation*}
so $\operatorname{pol}(AB)\le\operatorname{pol}A+\operatorname{pol}B$. For a presentation $Z_1,\dots,Z_r$ of $h$ it follows that $\operatorname{pol}h\le\sum_i\operatorname{pol}(e^{Z_i})\le\sum_i|Z_i|$, and the infimum gives $\operatorname{pol}h\le\operatorname{dis}h$.

\emph{(d4).} Let $Z,B\in\mathfrak{g}^{\mathbb{C}}$ and $f(t)=e^{tZ}Be^{-tZ}$. Then $f'(t)=e^{tZ}[Z,B]e^{-tZ}$, so
\begin{equation*}
(\operatorname{Ad}_{e^{Z}}-1)B=f(1)-f(0)=\int_0^1e^{tZ}[Z,B]e^{-tZ}\,dt.
\end{equation*}
By \eqref{4.24:eq-colour-norms-a},
\begin{gather*}
|[Z,B]|\le|ZB|+|BZ|\le2\|Z\|\,|B|\le2|Z|\,|B|,\\
|e^{tZ}[Z,B]e^{-tZ}|\le\|e^{tZ}\|\,\|e^{-tZ}\|\,|[Z,B]|\le e^{2t|Z|}\,2|Z|\,|B|.
\end{gather*}
Integrating,
\begin{equation*}
|(\operatorname{Ad}_{e^{Z}}-1)B|\le|B|\int_0^1 2|Z|\,e^{2t|Z|}\,dt=(e^{2|Z|}-1)\,|B|,
\end{equation*}
that is, $\|\operatorname{Ad}_{e^{Z}}-1\|_{e}\le e^{2|Z|}-1$, and therefore $\|\operatorname{Ad}_{e^{Z}}\|_{e}\le e^{2|Z|}$. For $g_1,g_2\in G^{\mathbb{C}}$,
\begin{equation*}
\operatorname{Ad}_{g_1g_2}-1=\operatorname{Ad}_{g_1}(\operatorname{Ad}_{g_2}-1)+(\operatorname{Ad}_{g_1}-1).
\end{equation*}
If $\|\operatorname{Ad}_{g_1}-1\|_{e}\le e^{2a}-1$ (so $\|\operatorname{Ad}_{g_1}\|_{e}\le e^{2a}$) and $\|\operatorname{Ad}_{g_2}-1\|_{e}\le e^{2b}-1$, this identity gives
\begin{equation*}
\|\operatorname{Ad}_{g_1g_2}-1\|_{e}\le e^{2a}(e^{2b}-1)+(e^{2a}-1)=e^{2(a+b)}-1.
\end{equation*}
For a presentation $Z_1,\dots,Z_r$ of $h$, induction on $j$ with $g_1=e^{Z_1}\cdots e^{Z_j}$, $a=\sum_{i\le j}|Z_i|$, $g_2=e^{Z_{j+1}}$, $b=|Z_{j+1}|$ yields $\|\operatorname{Ad}_h-1\|_{e}\le e^{2\sum_i|Z_i|}-1$, and the infimum over presentations gives (d4).

\emph{(d5).} Let $h\in G$. By Lemma~\ref{4.24:lem-adjoint-norm}(i), $\|\operatorname{Ad}_h\|_{e}=1$, so $\|\operatorname{Ad}_h-1\|_{e}\le2$; together with (d4) this is the first inequality. For the second, suppose first $\operatorname{dis}h\le\frac13$. The function $x\mapsto e^{2x}-1-3x$ is convex, vanishes at $x=0$, and at $x=\frac13$ equals $e^{2/3}-2<0$, because $e^{2}<8$; hence it is $\le0$ on $[0,\frac13]$, and
\begin{equation*}
e^{2\operatorname{dis}h}-1\le3\operatorname{dis}h\le6\operatorname{dis}h.
\end{equation*}
If $\operatorname{dis}h>\frac13$, then $2<6\operatorname{dis}h$. In both cases $\min\{2,e^{2\operatorname{dis}h}-1\}\le6\operatorname{dis}h$.
\end{proof}

\begin{lemma}[The trace of an adjoint rotation]\label{4.24:lem-rotation-trace}
Let $G$, $\mathfrak{g}^{\mathbb{C}}$, $n_{\mathfrak{g}}$, the norms $|\cdot|$ and
$\|\cdot\|_{e}$ and the trace $\operatorname{tr}_{\mathfrak{g}}$ on endomorphisms of
$\mathfrak{g}^{\mathbb{C}}$ be as in Definition~\ref{4.24:not-colour-norms}, and let
$\operatorname{dis}$ be the distance to the identity of
Lemma~\ref{4.24:lem-group-distance}. Put
\begin{equation}\label{4.24:eq-rotation-trace-a}
\varphi(x):=e^{2x}-1-2x ,\qquad x\ge 0 .
\end{equation}
Then for every $\mathfrak{h}\in G$
\begin{equation}\label{4.24:eq-rotation-trace-1}
\bigl|\operatorname{tr}_{\mathfrak{g}}(\operatorname{Ad}_{\mathfrak{h}}-1)\bigr|
\le n_{\mathfrak{g}}\,\min\{2,\varphi(\operatorname{dis}\mathfrak{h})\}.
\end{equation}
Moreover $\varphi(x)\le e^{2x}-1$ for $x\ge0$, and for $x\in[0,\tfrac13]$
\begin{equation}\label{4.24:eq-rotation-trace-2}
\varphi(x)\le 2x^2e^{2x}\le 4x^2 .
\end{equation}
Finally, writing $\operatorname{tr}_{\mathbb{C}^N}$ for the (unnormalised) trace of
matrices acting on $\mathbb{C}^N$ and $\chi_{\mathrm{adj}}$ for the character of the
adjoint representation,
\begin{equation}\label{4.24:eq-rotation-trace-3}
\operatorname{tr}_{\mathfrak{g}}\operatorname{Ad}_{\mathfrak{h}}
=\chi_{\mathrm{adj}}(\mathfrak{h})
=|\operatorname{tr}_{\mathbb{C}^N}\mathfrak{h}|^2-1 .
\end{equation}
\end{lemma}

\begin{proof}
Let $\mathfrak{h}=e^{Z_1}\cdots e^{Z_r}$ with $Z_i\in\mathfrak{g}^{\mathbb{C}}$ be any
presentation of the kind entering the definition of $\operatorname{dis}$. For
$Z\in\mathfrak{g}^{\mathbb{C}}$ write $\operatorname{ad}_Z B=[Z,B]$. Then
$\operatorname{Ad}_{e^{Z}}=e^{\operatorname{ad}_Z}$, and by \eqref{4.24:eq-colour-norms-a}
together with $\|Z\|\le|Z|$ we have
$|[Z,B]|\le|ZB|+|BZ|\le 2|Z|\,|B|$, that is $\|\operatorname{ad}_Z\|_{e}\le 2|Z|$. Hence
\begin{equation*}
\operatorname{Ad}_{e^{Z_i}}=1+\operatorname{ad}_{Z_i}+R_i,\qquad
\|R_i\|_{e}\le\sum_{n\ge2}\frac{(2|Z_i|)^n}{n!}=e^{2|Z_i|}-1-2|Z_i| .
\end{equation*}
Expand the ordered product
$\operatorname{Ad}_{\mathfrak{h}}=\prod_{i=1}^{r}(1+\operatorname{ad}_{Z_i}+R_i)$ by
choosing one of the three summands in each factor. The choice of $1$ everywhere gives
$1$; the choices with exactly one $\operatorname{ad}_{Z_i}$ and $1$ elsewhere give
$\sum_i\operatorname{ad}_{Z_i}$; call $R$ the sum of all remaining terms, so that
$\operatorname{Ad}_{\mathfrak{h}}-1=\sum_i\operatorname{ad}_{Z_i}+R$. By the triangle
inequality and sub-multiplicativity of $\|\cdot\|_{e}$, $\|R\|_{e}$ is at most the sum of
the corresponding products of norms, which is the polynomial
$\prod_i(1+a_i+\rho_i)-1-\sum_i a_i$ in $a_i=\|\operatorname{ad}_{Z_i}\|_{e}$,
$\rho_i=\|R_i\|_{e}$; it has non-negative coefficients and is therefore non-decreasing in
each $a_i$. Using $a_i\le2|Z_i|$ and $1+2|Z_i|+\|R_i\|_{e}\le e^{2|Z_i|}$,
\begin{equation*}
\|R\|_{e}\le\prod_{i}\bigl(1+2|Z_i|+\|R_i\|_{e}\bigr)-1-\sum_i2|Z_i|
\le e^{2\sum_i|Z_i|}-1-2\sum_i|Z_i| .
\end{equation*}

Next, $\operatorname{tr}_{\mathfrak{g}}\operatorname{ad}_Z=0$ for every
$Z\in\mathfrak{g}^{\mathbb{C}}$. Indeed $\mathfrak{g}^{\mathbb{C}}=\mathfrak{sl}(N,\mathbb{C})$
is simple ($N\ge2$), hence equal to its derived algebra, so $Z=\sum_j[X_j,Y_j]$ with
finitely many $X_j,Y_j\in\mathfrak{g}^{\mathbb{C}}$; by the Jacobi identity
$\operatorname{ad}_{[X_j,Y_j]}=[\operatorname{ad}_{X_j},\operatorname{ad}_{Y_j}]$, and the
trace of a commutator vanishes. Consequently
$\operatorname{tr}_{\mathfrak{g}}(\operatorname{Ad}_{\mathfrak{h}}-1)
=\operatorname{tr}_{\mathfrak{g}}R$. The absolute value of the trace of an endomorphism of
the $n_{\mathfrak{g}}$-dimensional space $\mathfrak{g}^{\mathbb{C}}$ is at most
$n_{\mathfrak{g}}$ times its spectral radius, hence at most $n_{\mathfrak{g}}$ times its
operator norm; therefore
\begin{equation*}
\bigl|\operatorname{tr}_{\mathfrak{g}}(\operatorname{Ad}_{\mathfrak{h}}-1)\bigr|
\le n_{\mathfrak{g}}\,\varphi\Bigl(\sum_i|Z_i|\Bigr).
\end{equation*}
The function $\varphi$ is continuous and non-decreasing on $[0,\infty)$, since
$\varphi'(x)=2e^{2x}-2\ge0$. Given $\eta>0$, choose a presentation with
$\sum_i|Z_i|\le\operatorname{dis}\mathfrak{h}+\eta$; then the left side is at most
$n_{\mathfrak{g}}\varphi(\operatorname{dis}\mathfrak{h}+\eta)$, and letting
$\eta\to0$ (the infimum over presentations) gives the bound
$n_{\mathfrak{g}}\varphi(\operatorname{dis}\mathfrak{h})$. The bound $2n_{\mathfrak{g}}$
follows in the same way from $\|\operatorname{Ad}_{\mathfrak{h}}-1\|_{e}\le2$ for
$\mathfrak{h}\in G$, which is clause (d5) of \eqref{4.24:eq-group-distance-a}. This proves
\eqref{4.24:eq-rotation-trace-1}.

For the elementary inequalities: $\varphi(x)\le e^{2x}-1$ because $2x\ge0$. Further
\begin{equation*}
\varphi(x)=\sum_{n\ge2}\frac{(2x)^n}{n!}
=4x^2\sum_{m\ge0}\frac{(2x)^m}{(m+2)!}
\le 2x^2\sum_{m\ge0}\frac{(2x)^m}{m!}=2x^2e^{2x},
\end{equation*}
since $(m+2)!\ge2\,m!$. For $x\le\tfrac13$ we have $e^{2x}\le e^{2/3}<2$, because
$e^{2/3}<2$ is equivalent to $e<2\sqrt2$; hence $2x^2e^{2x}\le4x^2$, which is
\eqref{4.24:eq-rotation-trace-2}.

Finally, $\operatorname{tr}_{\mathfrak{g}}\operatorname{Ad}_{\mathfrak{h}}$ is by definition
the character of the adjoint representation at $\mathfrak{h}$. The action
$X\mapsto\mathfrak{h}X\mathfrak{h}^{-1}=\mathfrak{h}X\mathfrak{h}^{*}$ on
$M_N(\mathbb{C})\cong\mathbb{C}^N\otimes\overline{\mathbb{C}^N}$ has trace
$|\operatorname{tr}_{\mathbb{C}^N}\mathfrak{h}|^2$, since the coefficient of the matrix
unit $E_{ab}$ in $\mathfrak{h}E_{ab}\mathfrak{h}^{*}$ is
$\mathfrak{h}_{aa}\overline{\mathfrak{h}_{bb}}$, and summing over $a,b$ gives
$\sum_{a,b}\mathfrak{h}_{aa}\overline{\mathfrak{h}_{bb}}
=|\operatorname{tr}_{\mathbb{C}^N}\mathfrak{h}|^2$. The decomposition
$M_N(\mathbb{C})=\mathfrak{g}^{\mathbb{C}}\oplus\mathbb{C}\mathbf{1}$ is invariant, and
$\mathfrak{h}$ acts trivially on $\mathbb{C}\mathbf{1}$; subtracting the trace $1$ of the
trivial summand gives \eqref{4.24:eq-rotation-trace-3}.
\end{proof}

\begin{definition}[Gauges and parallel transport along lattice paths]\label{4.24:not-gauge-transport}
The group $G=\mathrm{SU}(N)$, its complexification $G^{\mathbb{C}}$, the algebras
$\mathfrak{g}\subset\mathfrak{g}^{\mathbb{C}}$ and the adjoint action
$\operatorname{Ad}_g X=gXg^{-1}$ are those of
Definition~\ref{4.24:not-colour-norms}. The following conventions hold throughout this
section.
\begin{enumerate}
\item Every bond $b$ has an initial endpoint $b_{-}$ and a final endpoint $b_{+}$. For a set
$Y$ of bonds, $\operatorname{st}(Y)$ is the set of endpoints of the bonds of $Y$.
\item A number $\xi\ge0$, the shift in the resolvents $(\mathcal{A}+\xi)^{-1}$ and
$(\mathcal{A}_{ss}+\xi)^{-1}$ below, is fixed throughout and is mostly suppressed from the
notation. It is not the lattice spacing, which in this section is written $\eta$.
\item A \emph{gauge} is a map $u$ from the sites to $G$. Let $\mathbb{U}$ be a function on
bonds with values in $G^{\mathbb{C}}$ (in the cases below, $G$-valued or of the form
$e^{i\eta A'}U$), and $\mathbb{J}$ a function on bonds with values in
$\mathfrak{g}^{\mathbb{C}}$. A gauge $u$ acts on the pair $(\mathbb{U},\mathbb{J})$ by
\begin{equation}\label{4.24:eq-gauge-transport-1}
\mathbb{U}^{u}(b):=u(b_{-})\,\mathbb{U}(b)\,u(b_{+})^{-1},\qquad
(\operatorname{Ad}_u\mathbb{J})(b):=\operatorname{Ad}_{u(b_{-})}\mathbb{J}(b).
\end{equation}
On the bond fibre space $\mathcal{H}=\bigoplus_b\mathfrak{g}^{\mathbb{C}}_b$ it acts by
$\mathcal{U}_u:=\bigoplus_b\operatorname{Ad}_{u(b_{-})}$. The trivial pair is $(1,0)$, that is,
$\mathbb{U}\equiv1$ and $\mathbb{J}\equiv0$.
\item For the bond fields $\mathbb{U}=e^{i\eta A'}U$, the gauge acts by $U\mapsto U^{u}$ and
$A'\mapsto\operatorname{Ad}_uA'$, with $\operatorname{Ad}_u$ acting on the
$\mathfrak{g}^{\mathbb{C}}$-valued bond function $A'$ as in
\eqref{4.24:eq-gauge-transport-1}. Here $U$ is $G$-valued and satisfies
\cite[(3.35)]{B-CMP99b}, $A'$ is as in \cite[(3.37)]{B-CMP99b}, and $\eta$ is the lattice
spacing of those conditions. This action is consistent with
\eqref{4.24:eq-gauge-transport-1}:
\begin{equation}\label{4.24:eq-gauge-transport-2}
\mathbb{U}^{u}(b)=e^{i\eta\operatorname{Ad}_{u(b_{-})}A'(b)}\,U^{u}(b).
\end{equation}
\item A constant gauge $u\equiv c$ fixes the trivial pair: $(1,0)^{c}=(1,0)$.
\item Let $\gamma=(b_1,\dots,b_n)$ be a lattice path from a site $x$ to a site $y$. Write
$x=x_0,x_1,\dots,x_n=y$ for its successive sites, so that $b_j$ joins $x_{j-1}$ and $x_j$.
Put $\epsilon_j=+1$ if $x_{j-1}=(b_j)_{-}$ and $x_j=(b_j)_{+}$, and $\epsilon_j=-1$ if $b_j$
is traversed from $(b_j)_{+}$ to $(b_j)_{-}$. Then
\begin{equation}\label{4.24:eq-gauge-transport-3}
U(\gamma):=U(b_1)^{\epsilon_1}U(b_2)^{\epsilon_2}\cdots U(b_n)^{\epsilon_n},
\qquad
U^{u}(\gamma)=u(x)\,U(\gamma)\,u(y)^{-1}.
\end{equation}
\item For the rest of this section we fix a bond $b'$ and a walk $\omega=(X_0,\dots,X_p)$,
with $|\omega|:=p$, as follows.
  \begin{itemize}
  \item At the block stage, $b'$ lies in the bond set $s$ onto which the quadratic form
  $\mathcal{A}$ is compressed, $\mathcal{A}_{ss}=P_s\mathcal{A}P_s$. The walk $\omega$ is a
  walk of the generalised random walk expansion of $(\mathcal{A}_{ss}+\xi)^{-1}$ given by the
  walk expansion at the block.
  \item At the full stage, $b'$ lies in the bond set of the stage. The walk $\omega$ is a
  walk of the expansion of $(\mathcal{A}+\xi)^{-1}$ \cite[Theorem~3.10, (3.107)]{B-CMP99b}.
  \end{itemize}
\end{enumerate}
\end{definition}

\begin{proof}
\emph{The identity \eqref{4.24:eq-gauge-transport-2}.} For every $g\in G$ and every square
matrix $X$ we have $g e^{X}g^{-1}=e^{gXg^{-1}}$. Indeed, $gX^kg^{-1}=(gXg^{-1})^k$ for each
$k$, and the exponential series converges absolutely. Insert $u(b_{-})^{-1}u(b_{-})=1$ after
the exponential in the first definition of \eqref{4.24:eq-gauge-transport-1}. This gives
\begin{align*}
\mathbb{U}^{u}(b)
&=u(b_{-})e^{i\eta A'(b)}u(b_{-})^{-1}\cdot u(b_{-})U(b)u(b_{+})^{-1}\\
&=e^{i\eta\operatorname{Ad}_{u(b_{-})}A'(b)}\,U^{u}(b).
\end{align*}
This is \eqref{4.24:eq-gauge-transport-2}. It also shows that the action on $\mathbb{U}$ is
the action $U\mapsto U^{u}$, $A'\mapsto\operatorname{Ad}_uA'$.

\emph{Constant gauges.} If $u\equiv c$, then \eqref{4.24:eq-gauge-transport-1} gives
$1^{c}(b)=c\,1\,c^{-1}=1$. It also gives $(\operatorname{Ad}_c0)(b)=c\,0\,c^{-1}=0$.

\emph{Transport along a path.} Fix $j$. First let $\epsilon_j=+1$. Then $(b_j)_{-}=x_{j-1}$
and $(b_j)_{+}=x_j$, so by definition
\[
U^{u}(b_j)^{\epsilon_j}=u(x_{j-1})U(b_j)^{\epsilon_j}u(x_j)^{-1}.
\]
Now let $\epsilon_j=-1$. Then $(b_j)_{+}=x_{j-1}$ and $(b_j)_{-}=x_j$, and inverting the
product gives
\[
U^{u}(b_j)^{-1}=u(x_{j-1})U(b_j)^{-1}u(x_j)^{-1}.
\]
So in both cases
$U^{u}(b_j)^{\epsilon_j}=u(x_{j-1})U(b_j)^{\epsilon_j}u(x_j)^{-1}$. Take the ordered product
over $j=1,\dots,n$. The inner factors $u(x_j)^{-1}u(x_j)=1$, for $1\le j\le n-1$, cancel. The
result is
\[
U^{u}(\gamma)=u(x_0)U(\gamma)u(x_n)^{-1}=u(x)U(\gamma)u(y)^{-1}.
\]
This is the second identity of \eqref{4.24:eq-gauge-transport-3}.
\end{proof}

\begin{definition}[Structural hypotheses on a walk term]\label{4.24:def-walk-hypotheses}
Fix $\xi\ge0$. Let $\Lambda'$ be a region of the lattice and let $\mathcal{A}(\Lambda')$ be the
compression of the quadratic form $\mathcal{A}$ to its bonds; the block compression
$\mathcal{A}_{ss}=\mathcal{A}(\Lambda'_s)$ is the case $\Lambda'=\Lambda'_s$. By
\cite[Theorem~3.10, (3.107)]{B-CMP99b}, taken by statement under the hypotheses of that theorem,
the shifted compression has the generalised random walk expansion
\begin{equation*}
(\mathcal{A}(\Lambda')+\xi)^{-1}=\sum_{\omega}t^{\xi}_{\omega},
\end{equation*}
over walks $\omega=(X_0,\dots,X_p)$ of localisation domains with consecutive domains
intersecting. Let $\omega$ be such a walk and let $t^{\xi}_{\omega}$ be the term indexed by
$\omega$. Write $|\omega|:=p$, write
$X_r^{\sim}$ for the enlarged localisation domains, and put
$D(\omega):=X_0\cup\bigcup_{r\ge1}X_r^{\sim}$, the domain read by the walk term. Let $b'$ be a
bond.

Pairs $(\mathbb{U},\mathbb{J})$ are those of the class fixed in
Definition~\ref{4.24:not-gauge-transport}, with $\mathbb{U}=e^{i\eta A'}U$ and $U$ $G$-valued;
$(\mathbf{1},0)$ is the trivial pair. Gauges $u$ act on pairs and on $\mathcal{H}$ as there, by
$(\mathbb{U},\mathbb{J})\mapsto(\mathbb{U}^{u},\operatorname{Ad}_u\mathbb{J})$ and by
$\mathcal{U}_u$.

A kernel $F$ on $\mathcal{H}$ has entries $F(b,c)\in\operatorname{End}(\mathfrak{g}^{\mathbb{C}})$.
These entries are measured in $\|\cdot\|_{e}$. Products of kernels are compositions,
$(FF')(b,c)=\sum_{c'}F(b,c')F'(c',c)$.

Let $a,\beta'\ge0$ be numbers, whose sizes are discussed after this definition, and let
$m^{\xi}(\omega)\ge0$ be the majorant of the term $t^{\xi}_{\omega}$ given by
\cite[(3.108)]{B-CMP99b}, displayed after this definition. We say that $t^{\xi}_{\omega}$
\emph{satisfies the structural hypotheses at
$b'$} if there are an integer $q=q(\omega)\ge0$, domains $Y_0,\dots,Y_q$ and kernels
$F_0,\dots,F_q$ for which \textup{(a)}--\textup{(f)} below hold.
\begin{enumerate}
\item[(a)] \emph{Factorisation, locality, covariance.} For every pair,
\begin{equation}\label{4.24:eq-walk-hypotheses-a}
\begin{gathered}
t^{\xi}_{\omega}=F_0F_1\cdots F_q,\qquad q=q(\omega)\le2|\omega|,\\
F_i[\mathbb{U}^{u},\operatorname{Ad}_u\mathbb{J}]
=\mathcal{U}_u\,F_i[\mathbb{U},\mathbb{J}]\,\mathcal{U}_u^{-1}
\quad\text{for every gauge }u\text{ and }0\le i\le q .
\end{gathered}
\end{equation}
Each factor $F_i=F_i[\mathbb{U},\mathbb{J}]$ is a kernel on $\mathcal{H}$ with the following three
properties.
\begin{enumerate}
\item[(i)] It is local: each $Y_i$ is one of the $X_r$ or one of the $X_r^{\sim}$, and
$\bigcup_iY_i=D(\omega)$. Moreover $F_i(b,c)=0$ unless $b,c\in\operatorname{bd}(Y_i)$, and
$F_i$ depends only on the restriction of $(\mathbb{U},\mathbb{J})$ to $\operatorname{bd}(Y_i)$.
\item[(ii)] It is gauge covariant, as displayed in the second line of
\eqref{4.24:eq-walk-hypotheses-a}.
\item[(iii)] It is defined on the whole class of pairs, in particular at the trivial pair. We put
$G_i:=F_i[\mathbf{1},0]$.
\end{enumerate}
Finally, $b'\in\operatorname{bd}(Y_0)\cap\operatorname{bd}(Y_q)$.
\item[(b)] \emph{Majorants.} For each $i$ there is a kernel $\mu_i=\mu_i^{\xi}\ge0$ on
$\operatorname{bd}(Y_i)^2$. For every pair of the class restricted to $Y_i$, the trivial pair
included, it satisfies
$\|F_i[\mathbb{U},\mathbb{J}](b,c)\|_{e}\le\mu_i(b,c)$. The majorants obey the chain bound
\begin{equation}\label{4.24:eq-walk-hypotheses-b}
(\mu_0\cdots\mu_q)(b',b')
:=\sum_{b_1,\dots,b_q}\mu_0(b',b_1)\,\mu_1(b_1,b_2)\cdots\mu_q(b_q,b')
\le m^{\xi}(\omega).
\end{equation}
\end{enumerate}
Items \textup{(c)}--\textup{(f)} are required for every pair $(\mathbb{U},\mathbb{J})$ of the
class, with $\mathbb{U}=e^{i\eta A'}U$; the gauges $u_i$ and the potentials $B_i$ below may depend
on the pair, while the numbers $a$, $\varepsilon$ and $\beta'$ do not.
\begin{enumerate}
\item[(c)] \emph{Cube gauges.} For each $i$ there are a gauge
$u_i\colon\operatorname{st}(Y_i)\to G$, extended arbitrarily to all sites, and a $\mathfrak{g}$-valued
$B_i$ on $\operatorname{bd}(Y_i)$ such that
\begin{equation}\label{4.24:eq-walk-hypotheses-c}
U^{u_i}(b)=e^{i\eta B_i(b)}\quad(b\in\operatorname{bd}(Y_i)),\qquad
a_i:=\sup_{b\in\operatorname{bd}(Y_i)}|i\eta B_i(b)|\le a .
\end{equation}
Unwinding the definition of the gauge action, for $b\in\operatorname{bd}(Y_i)$ the pair in the
gauge $u_i$ reads
\begin{equation*}
\bigl(\mathbb{U}^{u_i},\operatorname{Ad}_{u_i}\mathbb{J}\bigr)(b)
=\bigl(e^{i\eta\operatorname{Ad}_{u_i(b_-)}A'(b)}\,e^{i\eta B_i(b)},\;
\operatorname{Ad}_{u_i(b_-)}\mathbb{J}(b)\bigr).
\end{equation*}
This pair is said to be in \emph{cube form} on $Y_i$.

Cube form is preserved by constant gauges $c\in G$. Indeed,
\begin{equation*}
\bigl(e^{i\eta B_i(b)}\bigr)^{c}=c\,e^{i\eta B_i(b)}c^{-1}=e^{i\eta\operatorname{Ad}_cB_i(b)} .
\end{equation*}
Hence the
gauge $cu_i$ has the properties required of $u_i$, with $B_i$ replaced by
$\operatorname{Ad}_cB_i$, $A'$ replaced by $\operatorname{Ad}_cA'$ and $\mathbb{J}$ replaced by
$\operatorname{Ad}_c\mathbb{J}$. The number $a_i$ is unchanged, because $|\operatorname{Ad}_cX|=|X|$
by the conjugation invariance of $|\cdot|$ (Definition~\ref{4.24:not-colour-norms}).
\item[(d)] \emph{Smallness in cube form.} For every $i$ and every pair $\mathsf{P}$ in cube form
on $Y_i$,
\begin{equation}\label{4.24:eq-walk-hypotheses-d}
\|(F_i[\mathsf{P}]-G_i)(b,c)\|_{e}\le\varepsilon\,\mu_i(b,c),
\end{equation}
where $\varepsilon$ is the per-cube field-dependence number,
\begin{equation*}
\varepsilon=O(1)(BCM\tilde\alpha_0+\tilde\alpha_1),
\end{equation*}
respectively, for the weighted class of pairs, in which the cube-clause bounds carry the weight
$w_{*}$,
\begin{equation*}
\varepsilon=O(1)\bigl[5w_{*}(BCM\tilde\alpha_0+\tilde\alpha_1)+K_{\Delta}\gamma_{\Delta}\bigr].
\end{equation*}
Here $B$, $C$, $\tilde\alpha_0$, $\tilde\alpha_1$, $K_{\Delta}$ and $\gamma_{\Delta}$ denote the
constants of the per-cube field-dependence bound for the respective class of pairs.
\item[(e)] \emph{Overlaps.} Put $Y_{q+1}:=Y_0$ and $B_{q+1}:=B_0$. For $0\le i\le q$ the overlap
$O_i:=\operatorname{st}(Y_i)\cap\operatorname{st}(Y_{i+1})$ is non-empty. The graph on $O_i$ whose
edges are the bonds of $\operatorname{bd}(Y_i)\cap\operatorname{bd}(Y_{i+1})$ is connected.

Fix base sites $s_i^{*}\in O_i$, with $s_q^{*}:=b'_{-}$. This is admissible because
$b'\in\operatorname{bd}(Y_q)\cap\operatorname{bd}(Y_0)$, so $b'_{-}\in O_q$. Put
\begin{equation}\label{4.24:eq-walk-hypotheses-e}
\beta_i:=\sup_{s\in O_i}\;\min_{\gamma\colon s_i^{*}\to s}\;\sum_{b\in\gamma}
\bigl(|i\eta B_i(b)|+|i\eta B_{i+1}(b)|\bigr),
\qquad \bar\beta:=\max_{0\le i\le q}\beta_i .
\end{equation}
Here the minimum runs over paths $\gamma$ from $s_i^{*}$ to $s$ in the graph on $O_i$.
\item[(f)] \emph{Paths.} Each $\operatorname{bd}(Y_i)$ is connected.

For $i=1,\dots,q$, let $\gamma_i$ be a path in $\operatorname{bd}(Y_i)$ from $s_{i-1}^{*}$ to
$s_i^{*}$. Both endpoints lie in $\operatorname{st}(Y_i)$, since $s_{i-1}^{*}\in O_{i-1}$ and
$s_i^{*}\in O_i$.

Let $\gamma_0$ be a path in $\operatorname{bd}(Y_0)$ from $s_q^{*}=b'_{-}$ to $s_0^{*}$. Both
endpoints lie in $\operatorname{st}(Y_0)$, since $b'\in\operatorname{bd}(Y_0)$ and
$s_0^{*}\in O_0$.

These paths satisfy
\begin{equation}\label{4.24:eq-walk-hypotheses-f}
\mathfrak{l}_i:=\sum_{b\in\gamma_i}|i\eta B_i(b)|\le\beta'\quad(0\le i\le q),\qquad
\Gamma_{\omega}:=\gamma_1\gamma_2\cdots\gamma_q\gamma_0 .
\end{equation}
\end{enumerate}
The loop $\Gamma_{\omega}$ is closed at $s_0^{*}$. Indeed, $\gamma_1$ starts at $s_0^{*}$,
each $\gamma_i$ ends where $\gamma_{i+1}$ starts, $\gamma_q$ ends at $s_q^{*}=b'_{-}$, and
$\gamma_0$ runs from $b'_{-}$ back to $s_0^{*}$. The loop lies in $D(\omega)$, because each
$\gamma_i$ lies in $\operatorname{bd}(Y_i)$ and $\bigcup_iY_i=D(\omega)$. It threads the domains
of the walk in their order.
\end{definition}

\begin{remark}
We comment on the hypotheses in turn. None of the following comments is an additional condition.

\emph{On \textup{(a)}.} The count $q\le2|\omega|$ is the count for the alternating form
$R(X_0)W(X_1)R(X_1)\cdots W(X_p)R(X_p)$ of the terms of \cite[(3.107)]{B-CMP99b}, in the notation
of that paper. That form has $2|\omega|+1$ factors. In this chapter only a
bound $q\le c_f|\omega|+c_f'$ with absolute constants $c_f,c_f'$ is used.

Locality and covariance are required of each factor separately. By \cite[Theorem~3.10]{B-CMP99b}
the whole term depends on the configuration restricted to $D(\omega)$. For the individual factors
of a general walk term, locality and covariance are taken by statement from the proof of
\cite[Theorem~3.10]{B-CMP99b}, whose terms are products of covariant local factors; the factors
themselves are not reproduced here. For the walk terms of the block compression
$\mathcal{A}_{ss}=\mathcal{A}(\Lambda'_s)$ of the quadratic form
$\mathcal{A}$, that is, of the compression of $\mathcal{A}$ by the projection $P_s$ onto the bonds
of the block region $\Lambda'_s$, properties \textup{(i)}--\textup{(iii)} of \textup{(a)} are part
of the statement of its walk expansion, taken by statement: the terms are products of covariant
factors, the factors being those of the walk expansion on $\Lambda'_s$. Their
covariance comes from the gauge covariance of $\mathcal{A}$, also taken by statement, since the
projection $P_s$ commutes with $\mathcal{U}_u$: both act
bond by bond, $P_s$ by restriction to a bond set and $\mathcal{U}_u$ fibrewise.

\emph{On \textup{(b)}.} The bound \eqref{4.24:eq-walk-hypotheses-b} is the product form of the
majorant
\begin{equation*}
m(\omega)=O(1)\bigl(O(1)M^{-1/2}\bigr)^{p}e^{-\delta_0d(\omega)}
\end{equation*}
of \cite[(3.108)]{B-CMP99b}, with $d(\omega)$, $\delta_0$ and the constants $O(1)$ as there; the
number $m^{\xi}(\omega)$ in \eqref{4.24:eq-walk-hypotheses-b} is this $m(\omega)$, and the same
majorant serves for every $\xi\ge0$. That
paper prints this majorant as a bound on the whole term $t^{\xi}_{\omega}$. Hypothesis \textup{(b)}
asks for it factor by factor, in the form \eqref{4.24:eq-walk-hypotheses-b}. The telescoping
estimates of this section use the majorants $\mu_i$ of the individual factors, not only a bound
on the whole term.

\emph{On \textup{(c)}.} Let $\square$ be a cube of the class of \cite[(3.35)]{B-CMP99b}, let $j$ be
the index attached to it there, and let $\ell=L^{j}\eta$ be its unit; in these comments $\ell$
always denotes the unit of a cube or of a domain. Then $a$ is the number of that cube clause in
the cube's own unit, that is, the bound it gives for $\ell|B|$. In the
form of \cite[(3.35)]{B-CMP99b} adopted for the class of pairs, one has
\begin{equation*}
\eta|B|\le O(1)M\tilde\alpha_0\cdot\frac{\eta}{\ell}\le O(1)M\tilde\alpha_0 .
\end{equation*}
For the weighted class of pairs, in which the cube-clause bounds carry the weight $w_{*}$, this
number is degraded by the factor $5w_{*}$.

The passage to the gauge $u_i$ leaves unchanged the bounds of \cite[(3.37)]{B-CMP99b} on $A'$ and
the bound on $\mathbb{J}$. This is because $|\operatorname{Ad}_uA'|=|A'|$ and
$|\operatorname{Ad}_u\mathbb{J}|=|\mathbb{J}|$ by conjugation invariance of $|\cdot|$.

The clause \cite[(3.35)]{B-CMP99b} is stated cube by cube. Hypothesis \textup{(c)} asks for one
gauge on each domain $Y_i$, and $Y_i$ may be an enlarged domain $X_r^{\sim}$. It is a hypothesis
in that form.

\emph{On \textup{(d)}.} The number $\varepsilon$ is the per-cube bound on the field dependence of
the local operators, in the gauge of the cube clause. Hypothesis \textup{(d)} is the per-cube
field-dependence bound, taken by statement from the field-dependence bound for the local
operators of the step,
\begin{equation*}
\bigl\|\Delta^{(j)}(\mathbb{U})-\Delta^{(j)}(\mathbf{1})\bigr\|
\le O(1)(BCM\tilde\alpha_0+\tilde\alpha_1)
\end{equation*}
for the local operators $\Delta^{(j)}$ on a cube, in the gauge of the cube clause, read factor by
factor on the domains $Y_i$. The factors of the expansion \cite[(3.107)]{B-CMP99b} are not
reproduced here, and \textup{(d)} is a hypothesis in that form.
Hypothesis \textup{(d)} concerns one factor, on one domain, in one gauge, the gauge $u_i$ of
\textup{(c)}. It is not a bound on the global difference of the quadratic form, and no such bound
is assumed.

Suppose a factor is the resolvent $R[\mathsf{P}]=\Delta[\mathsf{P}]^{-1}$ of a local operator
$\Delta[\mathsf{P}]$ acting on the bonds of $\operatorname{bd}(Y_i)$. Then the second resolvent
identity gives
\begin{equation*}
R[\mathsf{P}]-R[\mathbf{1},0]
=-R[\mathsf{P}]\bigl(\Delta[\mathsf{P}]-\Delta[\mathbf{1},0]\bigr)
\!\upharpoonright_{\operatorname{bd}(Y_i)}R[\mathbf{1},0],
\end{equation*}
where $\upharpoonright_{\operatorname{bd}(Y_i)}$ denotes restriction to $\operatorname{bd}(Y_i)$.
It expresses the difference in
\eqref{4.24:eq-walk-hypotheses-d} through the difference of the local operators in the gauge $u_i$.

\emph{On \textup{(e)} and \textup{(f)}.} The consecutive domains of a walk intersect
\cite[Theorem~3.10]{B-CMP99b}. Beyond this intersection, items \textup{(e)} and \textup{(f)} assume
the following:
\begin{itemize}
\item the connectedness of $\operatorname{bd}(Y_i)$ and of the graph of common bonds on $O_i$;
\item the bound $\mathfrak{l}_i\le\beta'$.
\end{itemize}
These properties belong to the geometry of Ba{\l}aban's domains. They do not follow from
intersection alone.

The intended size of $\bar\beta$ is one-sided and scale-free, as follows; it is not a further
condition. Suppose any two sites of
$O_i$ are joined in its graph by a path of at most $\check n\,\ell/\eta$ bonds. Suppose also that
both potentials obey $|i\eta B_i(b)|,|i\eta B_{i+1}(b)|\le a\,\eta/\ell$ on these bonds, where
$\ell$ is a unit $L^{j}\eta$, the same in both bounds. Then summing along such a path gives
\begin{equation*}
\beta_i\le\check n\,\frac{\ell}{\eta}\cdot2a\,\frac{\eta}{\ell}=2\check n\,a .
\end{equation*}
The right side involves neither $\eta$ nor the levels of $Y_i$ and $Y_{i+1}$. The number $\beta'$
in \eqref{4.24:eq-walk-hypotheses-f} is understood in the same way: it involves neither $\eta$
nor the levels of the domains.
\end{remark}

\begin{lemma}[Trivial-pair factors are colour scalars]\label{4.24:lem-colour-scalars}
Let $N\ge 2$. Let $t^{\xi}_{\omega}=F_0F_1\cdots F_q$ satisfy the structural hypotheses
\eqref{4.24:eq-walk-hypotheses-a} of Definition~\ref{4.24:def-walk-hypotheses}, and let
$G_i=F_i[\mathbf{1},0]$ and the bond $b'$ be as there. Then for every $i$ and all bonds $b,c$
\begin{equation*}
G_i(b,c)=g_i(b,c)\,\mathbf{1}_{\mathfrak{g}^{\mathbb{C}}},\qquad g_i(b,c)\in\mathbb{C}.
\end{equation*}
Consequently
\begin{equation}\label{4.24:eq-colour-scalars-1}
S^{0}:=(G_0G_1\cdots G_q)(b',b')=s^{0}\,\mathbf{1}_{\mathfrak{g}^{\mathbb{C}}},
\end{equation}
where
\begin{equation}\label{4.24:eq-colour-scalars-2}
s^{0}=s^{0}_{\omega}:=\sum_{b_1,\dots,b_q}g_0(b',b_1)\,g_1(b_1,b_2)\cdots g_q(b_q,b'),
\end{equation}
and $\tau^{0}=\tau^{0}_{\omega}:=\operatorname{tr}_{\mathfrak{g}}S^{0}=n_{\mathfrak{g}}\,s^{0}$.
Under the majorant hypothesis \eqref{4.24:eq-walk-hypotheses-b} moreover
\begin{equation}\label{4.24:eq-colour-scalars-3}
|s^{0}|=\|S^{0}\|_{e}\le m^{\xi}(\omega).
\end{equation}
\end{lemma}

\begin{proof}
Fix $\bar u\in SU(N)$ and let $u$ be the constant gauge, $u(x)=\bar u$ for every site $x$.
In the notation of Convention~\ref{4.24:not-gauge-transport}, the trivial pair is fixed by $u$:
for every bond $b$ we have $\mathbf{1}^{u}(b)=\bar u\,\mathbf{1}\,\bar u^{-1}=\mathbf{1}$ and
$(\operatorname{Ad}_u0)(b)=\operatorname{Ad}_{\bar u}0=0$. By the domain clause of
Definition~\ref{4.24:def-walk-hypotheses} each $F_i$ is defined at
the trivial pair, so the covariance clause of that definition, applied at the trivial pair, gives
$G_i=\mathcal{U}_uG_i\,\mathcal{U}_u^{-1}$. Since
$\mathcal{U}_u=\bigoplus_b\operatorname{Ad}_{u(b_{-})}=\bigoplus_b\operatorname{Ad}_{\bar u}$ acts
fibrewise, this identity of kernels reads, for all bonds $b,c$,
\begin{equation*}
G_i(b,c)=\operatorname{Ad}_{\bar u}\,G_i(b,c)\,\operatorname{Ad}_{\bar u}^{-1}.
\end{equation*}
Thus every $G_i(b,c)\in\operatorname{End}(\mathfrak{g}^{\mathbb{C}})$ commutes with
$\operatorname{Ad}_{\bar u}$ for all $\bar u\in SU(N)$.

We show that the representation $\bar u\mapsto\operatorname{Ad}_{\bar u}$ of $SU(N)$ on
$\mathfrak{g}^{\mathbb{C}}=\mathfrak{sl}(N,\mathbb{C})$ is irreducible. Let
$V\subset\mathfrak{sl}(N,\mathbb{C})$ be a complex subspace invariant under
$\operatorname{Ad}(SU(N))$. For $X\in\mathfrak{su}(N)$ and $Y\in V$ the curve
$t\mapsto\operatorname{Ad}_{\exp(tX)}Y$ lies in the finite-dimensional, hence closed, subspace $V$.
Its derivative at $t=0$, which is $\operatorname{ad}_XY$, therefore lies in $V$ as well. So $V$ is
invariant under $\operatorname{ad}(\mathfrak{su}(N))$. Since $V$ is a complex subspace and
$X\mapsto\operatorname{ad}_X$ is complex linear, $V$ is invariant under the $\operatorname{ad}$ of
the complex span of $\mathfrak{su}(N)$, which is $\mathfrak{sl}(N,\mathbb{C})$. Hence $V$ is an
ideal of $\mathfrak{sl}(N,\mathbb{C})$. As $\mathfrak{sl}(N,\mathbb{C})$ is simple for $N\ge2$, we
get $V=0$ or $V=\mathfrak{sl}(N,\mathbb{C})$. Schur's lemma, applied to this irreducible
finite-dimensional complex representation and to the intertwiner $G_i(b,c)$, gives
$G_i(b,c)\in\mathbb{C}\cdot\mathbf{1}_{\mathfrak{g}^{\mathbb{C}}}$. We write $g_i(b,c)$ for the
scalar.

The rest is arithmetic. The kernel product at $(b',b')$ is
$\sum_{b_1,\dots,b_q}G_0(b',b_1)G_1(b_1,b_2)\cdots G_q(b_q,b')$. Each factor is a scalar multiple
of $\mathbf{1}_{\mathfrak{g}^{\mathbb{C}}}$, so each summand equals
$g_0(b',b_1)\cdots g_q(b_q,b')\,\mathbf{1}_{\mathfrak{g}^{\mathbb{C}}}$. Summing gives
\eqref{4.24:eq-colour-scalars-1} with $s^{0}$ as in \eqref{4.24:eq-colour-scalars-2}. Taking the
trace,
\begin{equation*}
\tau^{0}=s^{0}\operatorname{tr}_{\mathfrak{g}}\mathbf{1}_{\mathfrak{g}^{\mathbb{C}}}
=s^{0}\dim_{\mathbb{C}}\mathfrak{sl}(N,\mathbb{C})=n_{\mathfrak{g}}\,s^{0}.
\end{equation*}
The operator norm of the identity of the non-zero space $\mathfrak{g}^{\mathbb{C}}$ is $1$, so
$\|S^{0}\|_{e}=|s^{0}|$.

For the bound, read \eqref{4.24:eq-walk-hypotheses-b} at the trivial pair. For $b,c$ in
$\operatorname{bd}(Y_i)$ it gives
$|g_i(b,c)|=\|G_i(b,c)\|_{e}\le\mu_i(b,c)$. By the locality clause of
Definition~\ref{4.24:def-walk-hypotheses}, $G_i(b,c)=0$ unless
$b,c\in\operatorname{bd}(Y_i)$. Hence in \eqref{4.24:eq-colour-scalars-2} only terms with
$(b_i,b_{i+1})\in\operatorname{bd}(Y_i)^2$ for every $i$ contribute, where $b_0=b_{q+1}=b'$. For
these terms the triangle inequality gives
\begin{equation*}
|s^{0}|\le\sum_{b_1,\dots,b_q}\mu_0(b',b_1)\,\mu_1(b_1,b_2)\cdots\mu_q(b_q,b')
=(\mu_0\cdots\mu_q)(b',b')\le m^{\xi}(\omega).
\end{equation*}
The last inequality is the chain bound in \eqref{4.24:eq-walk-hypotheses-b}. This proves
\eqref{4.24:eq-colour-scalars-3}.

The only property of the gauge group used is the simplicity of $\mathfrak{sl}(N,\mathbb{C})$,
and no step is specific to $N=2$.
\end{proof}

\begin{proposition}[Gauge invariance of walk-term densities]\label{4.24:prop-density-invariance}
Assume the structural hypotheses on a walk term
(Definition~\ref{4.24:def-walk-hypotheses}, display~\eqref{4.24:eq-walk-hypotheses-a}).
Then for every walk $\omega$, every bond $b'\in\operatorname{bd}(Y_0)\cap\operatorname{bd}(Y_q)$
as in Definition~\ref{4.24:def-walk-hypotheses}, every $\xi\ge 0$ and every gauge $u$,
\begin{equation*}
\operatorname{tr}_{\mathfrak{g}}\,\hat t^{\xi}_{\omega}(b',b')\bigl[\mathbb{U}^{u},
\operatorname{Ad}_u\mathbb{J}\bigr]
=\operatorname{tr}_{\mathfrak{g}}\,\hat t^{\xi}_{\omega}(b',b')\bigl[\mathbb{U},\mathbb{J}\bigr].
\end{equation*}
Here, as in the definition of the family of block densities $\mathbb{C}^{s}_{G}(X)$ given
below in this chapter, $\hat t^{\xi}_{\omega}=t^{\xi}_{\omega}$ for $|\omega|\ge1$, and for
$|\omega|=0$ the kernel $\hat t^{\xi}_{\omega}$ is $t^{\xi}_{\omega}$ minus
$(\lambda_0+\xi)^{-1}$ times the share of the identity operator assigned to the trivial walk
$\omega$ in that definition, $\lambda_0$ being the constant of this subtraction fixed there.
\end{proposition}

\begin{proof}
The gauge $u$ acts on pairs by $\mathbb{U}^{u}(b)=u(b_{-})\mathbb{U}(b)u(b_{+})^{-1}$ and
$(\operatorname{Ad}_u\mathbb{J})(b)=\operatorname{Ad}_{u(b_{-})}\mathbb{J}(b)$, and on
$\mathcal{H}=\bigoplus_b\mathfrak{g}^{\mathbb{C}}_b$ by
$\mathcal{U}_u=\bigoplus_b\operatorname{Ad}_{u(b_{-})}$
(Definition~\ref{4.24:not-gauge-transport}). Each $\operatorname{Ad}_{u(b_{-})}$ is invertible
on $\mathfrak{g}^{\mathbb{C}}$ with inverse $\operatorname{Ad}_{u(b_{-})^{-1}}$, so
$\mathcal{U}_u$ is invertible and
$\mathcal{U}_u^{-1}=\bigoplus_b\operatorname{Ad}_{u(b_{-})}^{-1}$. We proceed in five steps.

\emph{Step 1: the factors.} By the structural hypotheses,
$t^{\xi}_{\omega}=F_0F_1\cdots F_q$ with $q=q(\omega)$, and by their covariance clause every
factor satisfies
\begin{equation}\label{4.24:eq-density-invariance-1}
F_i\bigl[\mathbb{U}^{u},\operatorname{Ad}_u\mathbb{J}\bigr]
=\mathcal{U}_u\,F_i[\mathbb{U},\mathbb{J}]\,\mathcal{U}_u^{-1},
\qquad 0\le i\le q .
\end{equation}
For a factor built from the quadratic form $\mathcal{A}$ of the step,
\eqref{4.24:eq-density-invariance-1} is the gauge covariance of that form: under
$(\mathbb{U},\mathbb{J})\mapsto(\mathbb{U}^{u},\operatorname{Ad}_u\mathbb{J})$
one has $\mathcal{A}\mapsto\mathcal{U}_u\mathcal{A}\,\mathcal{U}_u^{-1}$, the operators of which
$\mathcal{A}$ is composed being covariant in the same way. Of this, only the displayed
identity~\eqref{4.24:eq-density-invariance-1}, which is the covariance clause of
Definition~\ref{4.24:def-walk-hypotheses}, is used below, and not the internal
structure of the operators of which $\mathcal{A}$ is composed.

\emph{Step 2: the walk term.} Multiply the identities~\eqref{4.24:eq-density-invariance-1} in
the order $i=0,1,\dots,q$. Between consecutive factors the product
$\mathcal{U}_u^{-1}\mathcal{U}_u$ equals the identity of $\mathcal{H}$, so the product
telescopes:
\begin{equation}\label{4.24:eq-density-invariance-2}
\begin{split}
t^{\xi}_{\omega}\bigl[\mathbb{U}^{u},\operatorname{Ad}_u\mathbb{J}\bigr]
&=\prod_{i=0}^{q}\mathcal{U}_u\,F_i[\mathbb{U},\mathbb{J}]\,\mathcal{U}_u^{-1}
=\mathcal{U}_u\,F_0[\mathbb{U},\mathbb{J}]\cdots F_q[\mathbb{U},\mathbb{J}]\,
\mathcal{U}_u^{-1}\\
&=\mathcal{U}_u\,t^{\xi}_{\omega}[\mathbb{U},\mathbb{J}]\,\mathcal{U}_u^{-1}.
\end{split}
\end{equation}

\emph{Step 3: the diagonal block.} The operator $\mathcal{U}_u$ is bond-diagonal: its kernel
is $\mathcal{U}_u(b,c)=\delta_{bc}\operatorname{Ad}_{u(b_{-})}$. Hence for any kernel $T$ on
$\mathcal{H}$ and any bonds $b,c$,
\begin{equation*}
\bigl(\mathcal{U}_u T\,\mathcal{U}_u^{-1}\bigr)(b,c)
=\operatorname{Ad}_{u(b_{-})}\,T(b,c)\,\operatorname{Ad}_{u(c_{-})}^{-1}.
\end{equation*}
Applied to~\eqref{4.24:eq-density-invariance-2} at $b=c=b'$ this gives
\begin{equation}\label{4.24:eq-density-invariance-3}
t^{\xi}_{\omega}(b',b')\bigl[\mathbb{U}^{u},\operatorname{Ad}_u\mathbb{J}\bigr]
=\operatorname{Ad}_{u(b'_{-})}\,t^{\xi}_{\omega}(b',b')[\mathbb{U},\mathbb{J}]\,
\operatorname{Ad}_{u(b'_{-})}^{-1}.
\end{equation}

\emph{Step 4: the subtracted term.} For $|\omega|\ge1$ one has
$\hat t^{\xi}_{\omega}=t^{\xi}_{\omega}$, and~\eqref{4.24:eq-density-invariance-3} is already
the transformation law of $\hat t^{\xi}_{\omega}(b',b')$. For $|\omega|=0$, by the definition
recalled after the statement, $\hat t^{\xi}_{\omega}$ is $t^{\xi}_{\omega}$ minus
$(\lambda_0+\xi)^{-1}$ times the share of the identity operator assigned to the trivial walk
$\omega$. This subtracted member does not depend on
$(\mathbb{U},\mathbb{J})$, and its block at $(b',b')$ is a scalar multiple
$c\,\mathbf{1}_{\mathfrak{g}^{\mathbb{C}}}$ of the identity of $\mathfrak{g}^{\mathbb{C}}$, with
$c$ independent of the pair. A multiple of the identity is fixed by conjugation:
$\operatorname{Ad}_{u(b'_{-})}\,c\,\mathbf{1}_{\mathfrak{g}^{\mathbb{C}}}\,
\operatorname{Ad}_{u(b'_{-})}^{-1}=c\,\mathbf{1}_{\mathfrak{g}^{\mathbb{C}}}$. Subtracting it
from both sides of~\eqref{4.24:eq-density-invariance-3} therefore gives
\begin{equation}\label{4.24:eq-density-invariance-4}
\begin{split}
\hat t^{\xi}_{\omega}(b',b')\bigl[\mathbb{U}^{u},\operatorname{Ad}_u\mathbb{J}\bigr]
&=\operatorname{Ad}_{u(b'_{-})}\bigl(t^{\xi}_{\omega}(b',b')[\mathbb{U},\mathbb{J}]
-c\,\mathbf{1}_{\mathfrak{g}^{\mathbb{C}}}\bigr)\operatorname{Ad}_{u(b'_{-})}^{-1}\\
&=\operatorname{Ad}_{u(b'_{-})}\,\hat t^{\xi}_{\omega}(b',b')[\mathbb{U},\mathbb{J}]\,
\operatorname{Ad}_{u(b'_{-})}^{-1}.
\end{split}
\end{equation}
So in every case $\hat t^{\xi}_{\omega}(b',b')$ transforms by conjugation with
$\operatorname{Ad}_{u(b'_{-})}$.

\emph{Step 5: the trace.} The trace $\operatorname{tr}_{\mathfrak{g}}$ on
$\operatorname{End}(\mathfrak{g}^{\mathbb{C}})$ is invariant under conjugation by invertible
elements:
$\operatorname{tr}_{\mathfrak{g}}(\Phi\Psi\Phi^{-1})=\operatorname{tr}_{\mathfrak{g}}(\Psi)$ by
cyclicity of the trace. Taking $\operatorname{tr}_{\mathfrak{g}}$ of both sides
of~\eqref{4.24:eq-density-invariance-4}, with $\Phi=\operatorname{Ad}_{u(b'_{-})}$, gives the
assertion.
\end{proof}

\begin{corollary}[Computation of a walk-term density in any gauge]
Assume the hypotheses of Proposition~\ref{4.24:prop-density-invariance}, and for a walk
$\omega$, a bond $b'$ as in Proposition~\ref{4.24:prop-density-invariance} and a pair
$(\mathbb{U},\mathbb{J})$ put
\begin{equation*}
\tau_{\omega}[\mathbb{U},\mathbb{J}]
:=\operatorname{tr}_{\mathfrak{g}}\,\hat t^{\xi}_{\omega}(b',b')[\mathbb{U},\mathbb{J}],
\end{equation*}
with $\xi\ge0$ fixed and suppressed from the notation. Then the value $\tau_{\omega}$ may be
computed in any global gauge: for every gauge $u$,
\begin{equation*}
\tau_{\omega}\bigl[\mathbb{U}^{u},\operatorname{Ad}_u\mathbb{J}\bigr]
=\tau_{\omega}[\mathbb{U},\mathbb{J}].
\end{equation*}
In particular, for every gauge $v$ the value at the trivial pair equals the value at the pure
gauge pair $(1^{v},0)$, where $1^{v}(b)=v(b_{-})v(b_{+})^{-1}$:
\begin{equation*}
\tau_{\omega}[1,0]=\tau_{\omega}\bigl[1^{v},0\bigr].
\end{equation*}
\end{corollary}

\begin{proof}
The first identity is the assertion of Proposition~\ref{4.24:prop-density-invariance},
written with the notation $\tau_{\omega}$. For the second, recall that by the domain clause of
the structural hypotheses (Definition~\ref{4.24:def-walk-hypotheses}) every factor $F_i$, and
hence $\hat t^{\xi}_{\omega}$ and $\tau_{\omega}$, is defined at the trivial pair
$(1,0)$. Apply the first identity with $u=v$ to the trivial pair. By the action of
gauges on pairs (Definition~\ref{4.24:not-gauge-transport}), the transform under $v$ of the
constant bond field $1$ is the field $b\mapsto v(b_{-})\,1\,v(b_{+})^{-1}$, and the transform
of the source $0$ is $\operatorname{Ad}_v 0$; for every bond $b$,
\begin{equation*}
v(b_{-})\,1\,v(b_{+})^{-1}=v(b_{-})v(b_{+})^{-1}=1^{v}(b),
\qquad
(\operatorname{Ad}_v 0)(b)=\operatorname{Ad}_{v(b_{-})}\,0=0,
\end{equation*}
so the transformed pair is the pure-gauge pair $(1^{v},0)$, and the first identity reads
$\tau_{\omega}[1^{v},0]=\tau_{\omega}[1,0]$.
\end{proof}

\begin{remark}
Proposition~\ref{4.24:prop-density-invariance} and the corollary above allow each density
$\tau_{\omega}$ to be evaluated in a gauge chosen for that density alone. They do not give a
single global gauge in which every factor domain $Y_i$ is simultaneously in cube form; in
general no such gauge exists, as is shown below in this chapter.
\end{remark}

\begin{proposition}[Factorwise kernel bounds in cube gauges]\label{4.24:prop-factorwise-bounds}
Let $t^{\xi}_{\omega}=F_0F_1\cdots F_q$ be a walk term satisfying the structural hypotheses
\eqref{4.24:eq-walk-hypotheses-a} of Definition~\ref{4.24:def-walk-hypotheses}. Then for every
factor $F_i$, every pair $(\mathbb{U},\mathbb{J})$ and every gauge $u$,
\begin{equation}\label{4.24:eq-factorwise-bounds-1}
\bigl\|F_i[\mathbb{U}^{u},\operatorname{Ad}_u\mathbb{J}](b,c)\bigr\|_{e}
=\bigl\|F_i[\mathbb{U},\mathbb{J}](b,c)\bigr\|_{e}
\qquad\text{for all bonds } b,c .
\end{equation}
Consequently:
\begin{enumerate}
\item[(i)] if $u_i$ is a cube gauge on the factor domain $Y_i$, that is, a gauge in which the pair
$(\mathbb{U}^{u_i},\operatorname{Ad}_{u_i}\mathbb{J})$ is in cube form on $Y_i$ (its cube-form
representative on $Y_i$), then every kernel majorant $\mu_i$ valid at
$(\mathbb{U}^{u_i},\operatorname{Ad}_{u_i}\mathbb{J})$ on $Y_i$ holds at $(\mathbb{U},\mathbb{J})$ itself;
\item[(ii)] if $\mu_0,\dots,\mu_q$ are kernel majorants as in \eqref{4.24:eq-walk-hypotheses-b}, each
$\mu_i$ majorising $F_i$ at the cube-form representative on $Y_i$ and the chain
$(\mu_0\cdots\mu_q)(b',b')$ being at most $m^{\xi}(\omega)$, then for the base bond
$b'\in\operatorname{bd}(Y_0)\cap\operatorname{bd}(Y_q)$ of \eqref{4.24:eq-walk-hypotheses-a} the chain
bound
\begin{equation}\label{4.24:eq-factorwise-bounds-2}
\bigl\|t^{\xi}_{\omega}(b',b')\bigr\|_{e}\le(\mu_0\cdots\mu_q)(b',b')\le m^{\xi}(\omega)
\end{equation}
is obtained factor by factor, each factor $F_i$ being bounded in its own cube gauge $u_i$. No gauge
transformation relating $u_i$ to $u_{i+1}$, no holonomy of $\mathbb{U}$, and no dependence on
$|D(\omega)|$ other than that already present in \eqref{4.24:eq-walk-hypotheses-b} enters.
\end{enumerate}
\end{proposition}

\begin{proof}
Write $\mathcal{U}_u=\bigoplus_b\operatorname{Ad}_{u(b_{-})}$ for the gauge action on $\mathcal{H}$. It is
block diagonal, with block $\operatorname{Ad}_{u(b_{-})}$ at the bond $b$. The covariance clause of
\eqref{4.24:eq-walk-hypotheses-a} gives
$F_i[\mathbb{U}^{u},\operatorname{Ad}_u\mathbb{J}]=\mathcal{U}_uF_i[\mathbb{U},\mathbb{J}]\mathcal{U}_u^{-1}$.
Reading off the $(b,c)$ entry of this product of kernels therefore gives
\begin{equation}\label{4.24:eq-factorwise-bounds-3}
F_i[\mathbb{U}^{u},\operatorname{Ad}_u\mathbb{J}](b,c)
=\operatorname{Ad}_{u(b_{-})}\,F_i[\mathbb{U},\mathbb{J}](b,c)\,\operatorname{Ad}_{u(c_{-})}^{-1}.
\end{equation}
The gauge $u$ takes values in $G$, and $\operatorname{Ad}_{u(c_{-})}^{-1}=\operatorname{Ad}_{u(c_{-})^{-1}}$
with $u(c_{-})^{-1}\in G$. By part~(i) of the lemma on the norm of the adjoint action
(Lemma~\ref{4.24:lem-adjoint-norm}), both operators flanking $F_i[\mathbb{U},\mathbb{J}](b,c)$ in
\eqref{4.24:eq-factorwise-bounds-3} are isometries of $(\mathfrak{g}^{\mathbb{C}},|\cdot|)$ of
$\|\cdot\|_{e}$-norm $1$.

By submultiplicativity of $\|\cdot\|_{e}$, the norm of the right-hand side of
\eqref{4.24:eq-factorwise-bounds-3} is therefore at most $\|F_i[\mathbb{U},\mathbb{J}](b,c)\|_{e}$.
Now solve \eqref{4.24:eq-factorwise-bounds-3} for $F_i[\mathbb{U},\mathbb{J}](b,c)$. It equals
$\operatorname{Ad}_{u(b_{-})^{-1}}F_i[\mathbb{U}^{u},\operatorname{Ad}_u\mathbb{J}](b,c)\operatorname{Ad}_{u(c_{-})}$,
again flanked by operators of norm $1$, so the reverse inequality holds too. This proves
\eqref{4.24:eq-factorwise-bounds-1}.

(i) Take $u:=u_i$ in \eqref{4.24:eq-factorwise-bounds-1}. If
$\|F_i[\mathbb{U}^{u_i},\operatorname{Ad}_{u_i}\mathbb{J}](b,c)\|_{e}\le\mu_i(b,c)$, then
\begin{equation*}
\|F_i[\mathbb{U},\mathbb{J}](b,c)\|_{e}
=\|F_i[\mathbb{U}^{u_i},\operatorname{Ad}_{u_i}\mathbb{J}](b,c)\|_{e}\le\mu_i(b,c).
\end{equation*}

(ii) Put $b_0=b_{q+1}=b'$. By the definition of the product of kernels on $\mathcal{H}$,
\begin{equation*}
t^{\xi}_{\omega}(b',b')=(F_0\cdots F_q)(b',b')
=\sum_{b_1,\dots,b_q}F_0(b_0,b_1)F_1(b_1,b_2)\cdots F_q(b_q,b_{q+1}).
\end{equation*}
The triangle inequality and submultiplicativity of $\|\cdot\|_{e}$ then give the first inequality
below. For each $i$ separately, part~(i) gives
$\|F_i(b_i,b_{i+1})\|_{e}\le\mu_i(b_i,b_{i+1})$, which is the second inequality:
\begin{equation*}
\bigl\|(F_0\cdots F_q)(b',b')\bigr\|_{e}
\le\sum_{b_1,\dots,b_q}\prod_{i=0}^{q}\|F_i(b_i,b_{i+1})\|_{e}
\le\sum_{b_1,\dots,b_q}\prod_{i=0}^{q}\mu_i(b_i,b_{i+1}).
\end{equation*}
The last sum is $(\mu_0\cdots\mu_q)(b',b')$. The second part of
\eqref{4.24:eq-walk-hypotheses-b} bounds it by $m^{\xi}(\omega)$, which gives
\eqref{4.24:eq-factorwise-bounds-2}.

In this argument each factor enters only through its own estimate in its own gauge $u_i$, and the
only sum performed is the one defining the chain. Hence no transition between the gauges of
consecutive factors, no holonomy, and no dependence on $|D(\omega)|$ beyond that of
\eqref{4.24:eq-walk-hypotheses-b} is involved.
\end{proof}

In the small-field bounds one passes, on each cube, to the gauge of that cube. The proposition
isolates the part of that step which holds term by term: absolute bounds are cube-wise. It says
nothing about the difference between $t^{\xi}_{\omega}$ and its value at the trivial pair, which is
treated separately.

Throughout this proposition $\xi\ge 0$, the bond $b'$ and the walk $\omega$ are fixed as in
Definition~\ref{4.24:def-walk-hypotheses}. The walk term factorises as
$t^{\xi}_{\omega}=F_0F_1\cdots F_q$ by \eqref{4.24:eq-walk-hypotheses-a}, with factor domains
$Y_0,\dots,Y_q$.

Kernels on $\mathcal{H}$ are multiplied by
\begin{equation*}
(KK')(b,c)=\sum_{e\in\mathfrak{Bo}}K(b,e)K'(e,c).
\end{equation*}
For a gauge $u$ the gauge action $\mathcal{U}_u=\bigoplus_b\operatorname{Ad}_{u(b_{-})}$ of
Notation~\ref{4.24:not-gauge-transport} acts on kernels by
\begin{equation}\label{4.24:eq-product-form-1}
(\mathcal{U}_uK\mathcal{U}_{u}^{-1})(b,c)
=\operatorname{Ad}_{u(b_{-})}K(b,c)\operatorname{Ad}_{u(c_{-})}^{-1}.
\end{equation}
An element $C\in G$ is also read as the constant gauge with value $C$. Then
$\mathcal{U}_C=\bigoplus_b\operatorname{Ad}_C$, and
$(\mathcal{U}_CK\mathcal{U}_C^{-1})(b,c)=\operatorname{Ad}_CK(b,c)\operatorname{Ad}_C^{-1}$.

We use the following data.
\begin{itemize}
\item The cube gauges $u_0,\dots,u_q$ of \eqref{4.24:eq-walk-hypotheses-c}, each extended to
all sites, and $u_{q+1}:=u_0$.
\item The overlaps $O_i=\operatorname{st}(Y_i)\cap\operatorname{st}(Y_{i+1})$, with
$Y_{q+1}:=Y_0$, and the base sites $s_i^{*}\in O_i$ of
Definition~\ref{4.24:def-walk-hypotheses}, with $s_q^{*}=b'_{-}$. Since
$b'\in\operatorname{bd}(Y_q)\cap\operatorname{bd}(Y_0)$, indeed $b'_{-}\in O_q$.
\end{itemize}
For $0\le i\le q$ put
\begin{equation}\label{4.24:eq-product-form-2}
\tilde F_i:=F_i[\mathbb{U}^{u_i},\operatorname{Ad}_{u_i}\mathbb{J}],\qquad
h_i:=u_iu_{i+1}^{-1},\qquad c_i:=h_i(s_i^{*})\in G,\qquad k_i:=h_ic_i^{-1}.
\end{equation}
Thus $\tilde F_i$ is the factor at its cube-form representative. The map $h_i$ is $G$-valued
on sites, and only its values on $O_i$ will enter. Finally $k_i(s_i^{*})=1$. Put further
\begin{equation}\label{4.24:eq-product-form-3}
C_0:=1,\qquad C_{i+1}:=C_ic_i,\qquad
\mathfrak{h}=\mathfrak{h}_{\omega}:=C_{q+1}=c_0c_1\cdots c_q\in G,
\end{equation}
and define the rotated objects
\begin{equation}\label{4.24:eq-product-form-4}
\tilde F_i':=\mathcal{U}_{C_i}\tilde F_i\,\mathcal{U}_{C_i}^{-1},\qquad
k_i':=C_ik_iC_i^{-1}.
\end{equation}
Let $\hat t^{\xi}_{\omega}$ be the walk term with the subtraction made at $|\omega|=0$; thus
$\hat t^{\xi}_{\omega}=t^{\xi}_{\omega}$ for $|\omega|\ge 1$. The density of the walk term is
\begin{equation}\label{4.24:eq-product-form-10}
\tau_\omega:=\operatorname{tr}_{\mathfrak{g}}\,\hat t^{\xi}_{\omega}(b',b').
\end{equation}
The element $\mathfrak{h}_\omega$ is called the residual element of the walk.

\begin{proposition}[Product form of a walk term with transition defects]
\label{4.24:prop-product-form}
Assume the structural hypotheses \eqref{4.24:eq-walk-hypotheses-a} and
\eqref{4.24:eq-walk-hypotheses-c} of Definition~\ref{4.24:def-walk-hypotheses}, with the
notation fixed above.
\begin{enumerate}
\item[(i)] \emph{(Operator form.)} One has
\begin{equation}\label{4.24:eq-product-form-5}
t^{\xi}_{\omega}=\mathcal{U}_{u_0}^{-1}\,\tilde F_0\,\mathcal{U}_{h_0}\,\tilde F_1\,
\mathcal{U}_{h_1}\cdots\mathcal{U}_{h_{q-1}}\,\tilde F_q\,\mathcal{U}_{u_q}.
\end{equation}
In the $i$-th junction $\tilde F_i\,\mathcal{U}_{h_i}\,\tilde F_{i+1}$ only the values of
$h_i$ on $O_i$ enter.
\item[(ii)] \emph{(Trace form.)} For every walk put $S_\omega$ as below; if $|\omega|\ge 1$,
then $\tau_\omega$ is given by
\begin{equation}\label{4.24:eq-product-form-a}
\begin{gathered}
S_{\omega}:=\sum_{b_1,\dots,b_q}\tilde F_0'(b',b_1)\operatorname{Ad}_{k_0'((b_1)_{-})}
\tilde F_1'(b_1,b_2)\operatorname{Ad}_{k_1'((b_2)_{-})}\cdots\\
\operatorname{Ad}_{k_{q-1}'((b_q)_{-})}\tilde F_q'(b_q,b'),\\
\tau_{\omega}=\operatorname{tr}_{\mathfrak{g}}\bigl[S_{\omega}\operatorname{Ad}_{\mathfrak{h}_\omega}\bigr]
\qquad\text{if }|\omega|\ge 1.
\end{gathered}
\end{equation}
Let $|\omega|=0$ and suppose that the walk term has a single factor, $q(\omega)=0$ (as for the
alternating form of Definition~\ref{4.24:def-walk-hypotheses}). Then the term
subtracted in $\hat t^{\xi}_{\omega}$ is $(\lambda_0+\xi)^{-1}$ times a kernel whose block at
$(b',b')$ is a scalar multiple $\varsigma_\omega\mathbf{1}$ of the identity, and the identity
for $\tau_\omega$ in \eqref{4.24:eq-product-form-a} holds with $S_\omega$ replaced by
$S_\omega-(\lambda_0+\xi)^{-1}\varsigma_\omega\mathbf{1}$.
\item[(iii)] \emph{(What the pieces are.)} For each $i$,
\begin{equation}\label{4.24:eq-product-form-6}
\tilde F_i'=F_i[P_i'],\qquad
P_i':=\bigl(\mathbb{U}^{C_iu_i},\operatorname{Ad}_{C_iu_i}\mathbb{J}\bigr).
\end{equation}
The pair $P_i'$ is in cube form on $Y_i$: it is the constant rotation by $C_i$ of the
cube-form pair $(\mathbb{U}^{u_i},\operatorname{Ad}_{u_i}\mathbb{J})$. Each
$\operatorname{Ad}_{k_i'(s)}$, $s\in O_i$, is $\langle\cdot,\cdot\rangle$-unitary. The
rotations leave $G_i$ untouched:
\begin{equation*}
\mathcal{U}_{C}G_i\,\mathcal{U}_C^{-1}=G_i\qquad\text{for every }C\in G.
\end{equation*}
\end{enumerate}
\end{proposition}

\begin{proof}
\emph{(i).} Fix $i$ and apply the covariance hypothesis (cov) of
\eqref{4.24:eq-walk-hypotheses-a} with the gauge $u_i$, extended to all sites. This gives
$\tilde F_i=\mathcal{U}_{u_i}F_i\,\mathcal{U}_{u_i}^{-1}$, that is,
\begin{equation}\label{4.24:eq-product-form-7}
F_i=\mathcal{U}_{u_i}^{-1}\,\tilde F_i\,\mathcal{U}_{u_i}.
\end{equation}
The extension of $u_i$ off $\operatorname{st}(Y_i)$ plays no role, for two reasons. First, by
the locality hypothesis (loc) of \eqref{4.24:eq-walk-hypotheses-a}, $\tilde F_i$ reads the pair
$(\mathbb{U}^{u_i},\operatorname{Ad}_{u_i}\mathbb{J})$ only on $\operatorname{bd}(Y_i)$. For
$b\in\operatorname{bd}(Y_i)$ both endpoints $b_{\pm}$ lie in $\operatorname{st}(Y_i)$. Hence
$\mathbb{U}^{u_i}(b)=u_i(b_-)\mathbb{U}(b)u_i(b_+)^{-1}$ and
$(\operatorname{Ad}_{u_i}\mathbb{J})(b)=\operatorname{Ad}_{u_i(b_-)}\mathbb{J}(b)$ are determined
by the restriction of $u_i$ to $\operatorname{st}(Y_i)$. Second, the kernel
$F_i(b,c)$ vanishes unless $b,c\in\operatorname{bd}(Y_i)$. By
\eqref{4.24:eq-product-form-1}, the conjugations in \eqref{4.24:eq-product-form-7} then
involve only $u_i(b_-)$ and $u_i(c_-)$ with $b_-,c_-\in\operatorname{st}(Y_i)$.

Multiplying \eqref{4.24:eq-product-form-7} over $i=0,\dots,q$ gives
\begin{equation*}
t^{\xi}_{\omega}=\mathcal{U}_{u_0}^{-1}\tilde F_0\,\bigl(\mathcal{U}_{u_0}\mathcal{U}_{u_1}^{-1}\bigr)
\,\tilde F_1\,\bigl(\mathcal{U}_{u_1}\mathcal{U}_{u_2}^{-1}\bigr)\cdots
\bigl(\mathcal{U}_{u_{q-1}}\mathcal{U}_{u_q}^{-1}\bigr)\,\tilde F_q\,\mathcal{U}_{u_q}.
\end{equation*}
Each consecutive product is a gauge action, since $\operatorname{Ad}$ is a homomorphism:
\begin{equation*}
\mathcal{U}_{u_i}\mathcal{U}_{u_{i+1}}^{-1}
=\bigoplus_b\operatorname{Ad}_{u_i(b_-)}\operatorname{Ad}_{u_{i+1}(b_-)}^{-1}
=\bigoplus_b\operatorname{Ad}_{u_i(b_-)u_{i+1}(b_-)^{-1}}=\mathcal{U}_{h_i}.
\end{equation*}
This is \eqref{4.24:eq-product-form-5}.

The kernel of the $i$-th junction is
\begin{equation*}
(\tilde F_i\,\mathcal{U}_{h_i}\,\tilde F_{i+1})(b,e)
=\sum_{c}\tilde F_i(b,c)\operatorname{Ad}_{h_i(c_-)}\tilde F_{i+1}(c,e).
\end{equation*}
By (loc), $\tilde F_i(b,c)=0$ unless $c\in\operatorname{bd}(Y_i)$, and
$\tilde F_{i+1}(c,e)=0$ unless $c\in\operatorname{bd}(Y_{i+1})$. Only bonds
$c\in\operatorname{bd}(Y_i)\cap\operatorname{bd}(Y_{i+1})$ therefore contribute, and for them
$c_-\in\operatorname{st}(Y_i)\cap\operatorname{st}(Y_{i+1})=O_i$. Hence only the values of
$h_i$ on $O_i$ enter.

\emph{(ii).} The gauge actions are block diagonal. Taking the kernel of
\eqref{4.24:eq-product-form-5} at $(b',b')$ therefore gives
\begin{equation*}
t^{\xi}_{\omega}(b',b')=\operatorname{Ad}_{u_0(b'_-)}^{-1}\,\Pi\,
\operatorname{Ad}_{u_q(b'_-)},
\end{equation*}
where
\begin{multline*}
\Pi:=\sum_{b_1,\dots,b_q}\tilde F_0(b',b_1)\operatorname{Ad}_{h_0((b_1)_-)}\tilde F_1(b_1,b_2)\\
\cdots\operatorname{Ad}_{h_{q-1}((b_q)_-)}\tilde F_q(b_q,b').
\end{multline*}
Let $|\omega|\ge 1$, so that $\tau_\omega=\operatorname{tr}_{\mathfrak{g}}t^{\xi}_{\omega}(b',b')$
by \eqref{4.24:eq-product-form-10}. By cyclicity of $\operatorname{tr}_{\mathfrak{g}}$, and
since $u_{q+1}=u_0$ gives $h_q=u_qu_0^{-1}$,
\begin{equation*}
\tau_\omega=\operatorname{tr}_{\mathfrak{g}}\bigl(\Pi\operatorname{Ad}_{u_q(b'_-)u_0(b'_-)^{-1}}\bigr)
=\operatorname{tr}_{\mathfrak{g}}\bigl(\Pi\operatorname{Ad}_{h_q(b'_-)}\bigr)
=\operatorname{tr}_{\mathfrak{g}}\bigl(\Pi\operatorname{Ad}_{c_q}\bigr).
\end{equation*}
The last equality holds because $b'_-=s_q^{*}$, so that $k_q(b'_-)=1$ and
$h_q(b'_-)=c_q$.

Next, $h_i=k_ic_i$ by the definition of $k_i$, so
$\operatorname{Ad}_{h_i(s)}=\operatorname{Ad}_{k_i(s)}\operatorname{Ad}_{c_i}$ for every site $s$.
We push the constants to the right. Put $b_0:=b'$ and $b_{q+1}:=b'$. For $0\le j\le q$ let
$\Xi_j$ denote the raw string of the first $j$ factors and junctions,
\begin{multline*}
\Xi_j:=\tilde F_0(b_0,b_1)\operatorname{Ad}_{k_0((b_1)_-)}\operatorname{Ad}_{c_0}\\
\cdots\tilde F_{j-1}(b_{j-1},b_j)\operatorname{Ad}_{k_{j-1}((b_j)_-)}\operatorname{Ad}_{c_{j-1}},
\end{multline*}
and let $\Xi_j'$ denote the corresponding rotated string,
\begin{equation*}
\Xi_j':=\tilde F_0'(b_0,b_1)\operatorname{Ad}_{k_0'((b_1)_-)}\cdots
\tilde F_{j-1}'(b_{j-1},b_j)\operatorname{Ad}_{k_{j-1}'((b_j)_-)}.
\end{equation*}
Both are empty products, equal to $\mathbf{1}$, for $j=0$. We show by induction that
\begin{equation}\label{4.24:eq-product-form-8}
\Xi_j=\Xi_j'\operatorname{Ad}_{C_j}\qquad(0\le j\le q).
\end{equation}
For $j=0$ this is $\mathbf{1}=\operatorname{Ad}_{C_0}$, since $C_0=1$. For the inductive step
we use three facts:
\begin{itemize}
\item $\tilde F_i'(b,c)=\operatorname{Ad}_{C_i}\tilde F_i(b,c)\operatorname{Ad}_{C_i}^{-1}$, which
is \eqref{4.24:eq-product-form-4} read through the kernel formula for a constant gauge;
\item $\operatorname{Ad}_{C_i}\operatorname{Ad}_{k_i(s)}\operatorname{Ad}_{C_i}^{-1}
=\operatorname{Ad}_{C_ik_i(s)C_i^{-1}}=\operatorname{Ad}_{k_i'(s)}$;
\item $\operatorname{Ad}_{C_i}\operatorname{Ad}_{c_i}=\operatorname{Ad}_{C_ic_i}
=\operatorname{Ad}_{C_{i+1}}$.
\end{itemize}
Together they give, for all bonds $b,c$,
\begin{equation}\label{4.24:eq-product-form-9}
\begin{aligned}
\operatorname{Ad}_{C_i}\tilde F_i(b,c)\operatorname{Ad}_{k_i(c_-)}\operatorname{Ad}_{c_i}
&=\bigl[\operatorname{Ad}_{C_i}\tilde F_i(b,c)\operatorname{Ad}_{C_i}^{-1}\bigr]
\bigl[\operatorname{Ad}_{C_ik_i(c_-)C_i^{-1}}\bigr]\operatorname{Ad}_{C_ic_i}\\
&=\tilde F_i'(b,c)\operatorname{Ad}_{k_i'(c_-)}\operatorname{Ad}_{C_{i+1}}.
\end{aligned}
\end{equation}
If \eqref{4.24:eq-product-form-8} holds for some $j<q$, then
\begin{equation*}
\Xi_{j+1}=\Xi_j\,\tilde F_j(b_j,b_{j+1})\operatorname{Ad}_{k_j((b_{j+1})_-)}\operatorname{Ad}_{c_j}
=\Xi_j'\,\operatorname{Ad}_{C_j}\tilde F_j(b_j,b_{j+1})\operatorname{Ad}_{k_j((b_{j+1})_-)}
\operatorname{Ad}_{c_j}.
\end{equation*}
By \eqref{4.24:eq-product-form-9} with $i=j$, this equals
\begin{equation*}
\Xi_j'\,\tilde F_j'(b_j,b_{j+1})\operatorname{Ad}_{k_j'((b_{j+1})_-)}\operatorname{Ad}_{C_{j+1}}
=\Xi_{j+1}'\operatorname{Ad}_{C_{j+1}}.
\end{equation*}
This proves \eqref{4.24:eq-product-form-8}.

The summand of $\Pi\operatorname{Ad}_{c_q}$ for a given $(b_1,\dots,b_q)$ is
$\Xi_q\,\tilde F_q(b_q,b')\operatorname{Ad}_{c_q}$. By \eqref{4.24:eq-product-form-8} with
$j=q$, by \eqref{4.24:eq-product-form-4} for $i=q$, and since
$\operatorname{Ad}_{C_q}\operatorname{Ad}_{c_q}=\operatorname{Ad}_{C_{q+1}}$, it equals
\begin{equation*}
\Xi_q'\,\operatorname{Ad}_{C_q}\tilde F_q(b_q,b')\operatorname{Ad}_{c_q}
=\Xi_q'\,\tilde F_q'(b_q,b')\operatorname{Ad}_{C_{q+1}}.
\end{equation*}
Summing over $b_1,\dots,b_q$ gives $\Pi\operatorname{Ad}_{c_q}=S_\omega\operatorname{Ad}_{C_{q+1}}
=S_\omega\operatorname{Ad}_{\mathfrak{h}_\omega}$. Hence, for $|\omega|\ge 1$,
$\tau_\omega=\operatorname{tr}_{\mathfrak{g}}[S_\omega\operatorname{Ad}_{\mathfrak{h}_\omega}]$,
which is \eqref{4.24:eq-product-form-a}.

Now let $|\omega|=0$; then $\tau_\omega=\operatorname{tr}_{\mathfrak{g}}\hat t^{\xi}_{\omega}(b',b')$
by \eqref{4.24:eq-product-form-10}. The subtracted term contributes at
$(b',b')$ the multiple $(\lambda_0+\xi)^{-1}\varsigma_\omega\mathbf{1}$ of $\mathbf{1}$, and a
multiple of $\mathbf{1}$ is constant under every conjugation
$\operatorname{Ad}_g(\cdot)\operatorname{Ad}_g^{-1}$. By the clause $q(\omega)=0$ of (ii),
$u_1=u_{q+1}=u_0$, and hence
\begin{gather*}
h_0=u_0u_1^{-1}=1,\qquad c_0=h_0(s_0^{*})=1,\\
\mathfrak{h}_\omega=C_{q+1}=C_0c_0=1;
\end{gather*}
moreover, since $C_0=1$,
\begin{equation*}
S_\omega=\tilde F_0'(b',b')=\tilde F_0(b',b')=\Pi.
\end{equation*}
Because
$(\lambda_0+\xi)^{-1}\varsigma_\omega\mathbf{1}$ commutes with
$\operatorname{Ad}_{u_0(b'_-)}$,
\begin{equation*}
\hat t^{\xi}_{\omega}(b',b')
=\operatorname{Ad}_{u_0(b'_-)}^{-1}\bigl(\Pi-(\lambda_0+\xi)^{-1}\varsigma_\omega\mathbf{1}\bigr)
\operatorname{Ad}_{u_0(b'_-)}.
\end{equation*}
Cyclicity of the trace then gives
\begin{align*}
\tau_\omega&=\operatorname{tr}_{\mathfrak{g}}\hat t^{\xi}_{\omega}(b',b')\\
&=\operatorname{tr}_{\mathfrak{g}}\bigl[(S_\omega-(\lambda_0+\xi)^{-1}\varsigma_\omega\mathbf{1})
\operatorname{Ad}_{\mathfrak{h}_\omega}\bigr],
\end{align*}
which is the subtracted form of \eqref{4.24:eq-product-form-a}.

\emph{(iii).} Apply the covariance hypothesis (cov) to the pair
$(\mathbb{U}^{u_i},\operatorname{Ad}_{u_i}\mathbb{J})$ with the constant gauge $C_i$:
\begin{equation*}
\mathcal{U}_{C_i}\tilde F_i\,\mathcal{U}_{C_i}^{-1}
=F_i\bigl[(\mathbb{U}^{u_i})^{C_i},\operatorname{Ad}_{C_i}\operatorname{Ad}_{u_i}\mathbb{J}\bigr].
\end{equation*}
For every bond $b$ we have
\begin{equation*}
(\mathbb{U}^{u_i})^{C_i}(b)=C_iu_i(b_-)\mathbb{U}(b)u_i(b_+)^{-1}C_i^{-1}=\mathbb{U}^{C_iu_i}(b),
\end{equation*}
and
$\operatorname{Ad}_{C_i}\operatorname{Ad}_{u_i(b_-)}\mathbb{J}(b)
=\operatorname{Ad}_{C_iu_i(b_-)}\mathbb{J}(b)$. This proves
\eqref{4.24:eq-product-form-6}.

The pair $(\mathbb{U}^{u_i},\operatorname{Ad}_{u_i}\mathbb{J})$ is in cube form on $Y_i$ by
\eqref{4.24:eq-walk-hypotheses-c}. Cube form is preserved by constant gauges
(Definition~\ref{4.24:def-walk-hypotheses}). Hence $P_i'$ is in cube form on $Y_i$.

For $s\in O_i$ we have $h_i(s)=u_i(s)u_{i+1}(s)^{-1}\in G$ and $c_i\in G$, so
$k_i(s)\in G$. Also $C_i\in G$ as a product of the $c_j$, so $k_i'(s)=C_ik_i(s)C_i^{-1}\in G$.
By Lemma~\ref{4.24:lem-adjoint-norm}(i), $\operatorname{Ad}_{k_i'(s)}$ is
$\langle\cdot,\cdot\rangle$-unitary.

Finally, let $C\in G$. By Lemma~\ref{4.24:lem-colour-scalars},
$G_i(b,c)=g_i(b,c)\mathbf{1}$ with $g_i(b,c)\in\mathbb{C}$, so
\begin{equation*}
(\mathcal{U}_CG_i\,\mathcal{U}_C^{-1})(b,c)
=g_i(b,c)\operatorname{Ad}_C\operatorname{Ad}_C^{-1}=G_i(b,c).
\end{equation*}
Alternatively, (cov) with the constant gauge $C$ gives
\begin{equation*}
\mathcal{U}_CG_i\,\mathcal{U}_C^{-1}=\mathcal{U}_CF_i[1,0]\,\mathcal{U}_C^{-1}=F_i[(1,0)^C].
\end{equation*}
A constant gauge fixes the trivial pair, $(1,0)^C=(1,0)$, because $C\,1\,C^{-1}=1$ and
$\operatorname{Ad}_C0=0$. Hence $F_i[(1,0)^C]=F_i[1,0]=G_i$.
\end{proof}

\begin{remark}
The constants $C_i$ form a connecting gauge for the chain of domains. Take $s\in O_i$. Since
$u_i=h_iu_{i+1}=k_ic_iu_{i+1}$, we get
\begin{equation*}
C_iu_i(s)=\bigl(C_ik_i(s)C_i^{-1}\bigr)\,C_ic_i\,u_{i+1}(s)=k_i'(s)\,C_{i+1}u_{i+1}(s).
\end{equation*}
So the gauge $C_iu_i$ on $Y_i$ differs from the gauge $C_{i+1}u_{i+1}$ on $Y_{i+1}$, on the
overlap $O_i$, only by the transition defect $k_i'$. This defect is small; it is estimated in
the lemma of this section on the transition defects $k_i'$.

Consider the unfolded chain $Y_0,Y_1,\dots,Y_q$, in which the domains are taken as separate
copies in the order of the walk. On it, the gauges $C_iu_i$ on $Y_i$ together form one gauge
that realises every cube form up to $\bar\beta$. Closing the chain returns to $Y_0$ with the
gauge $C_{q+1}u_{q+1}=\mathfrak{h}_\omega u_0$ in place of $C_0u_0=u_0$. The gauge on the
unfolded chain therefore descends to the set $D(\omega)$ of the lattice if and only if
$\mathfrak{h}_\omega=1$.

Accordingly, an argument that brings all factors into cube form with a single gauge
reproduces, on the unfolded walk, the product form of
Proposition~\ref{4.24:prop-product-form}. On $D(\omega)$ itself such an argument fails by
exactly the residual element $\mathfrak{h}_\omega$.
\end{remark}

\begin{lemma}[Smallness of transition defects on overlaps]\label{4.24:lem-transition-defects}
Assume the cube-gauge hypothesis \eqref{4.24:eq-walk-hypotheses-c} and the overlap hypothesis
\eqref{4.24:eq-walk-hypotheses-e}. Use the notation of
Proposition~\ref{4.24:prop-product-form}:
\begin{itemize}
\item $u_0,\dots,u_q$ are the cube gauges of the cube-gauge hypothesis, and $u_{q+1}:=u_0$;
\item $O_i$ are the overlaps and $s_i^{*}$ the base sites of the overlap hypothesis;
\item $h_i:=u_iu_{i+1}^{-1}$, $c_i:=h_i(s_i^{*})\in G$ and $k_i:=h_ic_i^{-1}$;
\item $C_0:=1$, $C_{i+1}:=C_ic_i$ and $k_i':=C_ik_iC_i^{-1}$.
\end{itemize}
Then for every $i$ with $0\le i\le q$ and every $s\in O_i$ we have
\begin{equation*}
\operatorname{dis}k_i(s)\le\beta_i\le\bar\beta,
\qquad
\operatorname{dis}k_i'(s)=\operatorname{dis}k_i(s).
\end{equation*}
Consequently
\begin{equation*}
\|\operatorname{Ad}_{k_i'(s)}-\mathbf{1}\|_{e}\le 6\bar\beta,
\qquad
\|\operatorname{Ad}_{k_i'(s)}\|_{e}=1 .
\end{equation*}
\end{lemma}

\begin{proof}
Fix $i$. Since every $u_j$ takes values in $G$, so do $h_i$ and $k_i$, and $c_i\in G$; hence every
$C_j$, and every $k_i'(s)$, lies in $G$.

Let $b$ be a bond of $\operatorname{bd}(Y_i)\cap\operatorname{bd}(Y_{i+1})$, that is, an edge of
the graph on $O_i$ of the overlap hypothesis, so that $b_{-},b_{+}\in O_i$. By the cube-gauge
hypothesis the cube forms
on $Y_i$ and on $Y_{i+1}$ both hold on $b$. With $X_j:=i\eta B_j(b)$ this reads
$U^{u_i}(b)=e^{X_i}$ and $U^{u_{i+1}}(b)=e^{X_{i+1}}$. Since $u_i=h_iu_{i+1}$ pointwise and
$h_i(b_{+})^{-1}=u_{i+1}(b_{+})u_i(b_{+})^{-1}$,
\begin{align*}
U^{u_i}(b)&=u_i(b_{-})U(b)u_i(b_{+})^{-1}\\
&=h_i(b_{-})\bigl[u_{i+1}(b_{-})U(b)u_{i+1}(b_{+})^{-1}\bigr]u_{i+1}(b_{+})u_i(b_{+})^{-1}\\
&=h_i(b_{-})\,e^{X_{i+1}}\,h_i(b_{+})^{-1}.
\end{align*}
Comparing with $U^{u_i}(b)=e^{X_i}$ gives $h_i(b_{+})=e^{-X_i}h_i(b_{-})e^{X_{i+1}}$. Therefore
\begin{equation*}
h_i(b_{+})h_i(b_{-})^{-1}
=e^{-X_i}\,h_i(b_{-})e^{X_{i+1}}h_i(b_{-})^{-1}
=e^{-X_i}\,e^{\operatorname{Ad}_{h_i(b_{-})}X_{i+1}} .
\end{equation*}
By (d2) of Lemma~\ref{4.24:lem-group-distance}, $\operatorname{dis}$ is subadditive and satisfies
$\operatorname{dis}e^{Z}\le|Z|$, because $Z$ is a logarithm of $e^{Z}$. Part~(i) of
Lemma~\ref{4.24:lem-adjoint-norm} applies because $h_i(b_{-})\in G$, and it gives
$|\operatorname{Ad}_{h_i(b_{-})}X_{i+1}|=|X_{i+1}|$. Together these yield
\begin{equation}\label{4.24:eq-transition-defects-1}
\operatorname{dis}\bigl(h_i(b_{+})h_i(b_{-})^{-1}\bigr)\le|X_i|+|X_{i+1}|
=|i\eta B_i(b)|+|i\eta B_{i+1}(b)|.
\end{equation}
The element $h_i(b_{-})h_i(b_{+})^{-1}$ associated with the reversed bond is the inverse of the
element just estimated. Since $\operatorname{dis}(h^{-1})=\operatorname{dis}h$ by (d2), it obeys
the same bound.

Let $s\in O_i$. By the overlap hypothesis the graph on $O_i$ is connected, so there is a path
$s_i^{*}=z_0,z_1,\dots,z_n=s$ in it. Each pair $z_{m-1},z_m$ is joined by a bond of
$\operatorname{bd}(Y_i)\cap\operatorname{bd}(Y_{i+1})$, traversed in one of its two orientations.
Because $c_i=h_i(s_i^{*})$, the product telescopes:
\begin{equation*}
k_i(s)=h_i(s)h_i(s_i^{*})^{-1}
=\prod_{m=n}^{1}h_i(z_m)h_i(z_{m-1})^{-1},
\end{equation*}
with the factor for $m=n$ on the left. For $n=0$ the product is empty, $k_i(s)=1$, and
$\operatorname{dis}1=0$ by (d1). Each factor $h_i(z_m)h_i(z_{m-1})^{-1}$ is one of the two
elements estimated above for the bond joining $z_{m-1}$ and $z_m$. Subadditivity (d2) and
\eqref{4.24:eq-transition-defects-1} therefore give
\begin{equation*}
\operatorname{dis}k_i(s)\le\sum_{b\text{ on the path}}\bigl(|i\eta B_i(b)|+|i\eta B_{i+1}(b)|\bigr).
\end{equation*}
The left side does not depend on the path, so we may minimise over paths from $s_i^{*}$ to $s$.
By the definition of $\beta_i$ in the overlap hypothesis as a supremum over $s\in O_i$ of these
minima,
and of $\bar\beta$ as the maximum of the $\beta_i$, we obtain
$\operatorname{dis}k_i(s)\le\beta_i\le\bar\beta$.

Next, $k_i'(s)=C_ik_i(s)C_i^{-1}$ with $C_i\in G$. The equality case of (d2) for conjugation by
elements of $G$ gives $\operatorname{dis}k_i'(s)=\operatorname{dis}k_i(s)$.

Finally, $k_i'(s)\in G$, so (d5) applies:
\begin{equation*}
\|\operatorname{Ad}_{k_i'(s)}-\mathbf{1}\|_{e}\le6\operatorname{dis}k_i'(s)
=6\operatorname{dis}k_i(s)\le6\bar\beta .
\end{equation*}
Part~(i) of Lemma~\ref{4.24:lem-adjoint-norm} gives $\|\operatorname{Ad}_{k_i'(s)}\|_{e}=1$.
\end{proof}

No smallness of $h_i$ itself is used in the proof, nor does it hold. On $O_i$ we have
$h_i=k_ic_i$, so up to the small defect $k_i$ the function $h_i$ is an arbitrary constant
$c_i\in G$.

\begin{lemma}[The residual element as a Wilson loop]\label{4.24:lem-residual-loop}
Let $\omega$ be a walk whose term satisfies the cube-gauge hypothesis
\eqref{4.24:eq-walk-hypotheses-c}, the overlap hypothesis \eqref{4.24:eq-walk-hypotheses-e} and the
path hypothesis \eqref{4.24:eq-walk-hypotheses-f} of Definition~\ref{4.24:def-walk-hypotheses}.
These hypotheses supply the following data:
\begin{itemize}
\item the cube gauges $u_0,\dots,u_q$ of the cube-gauge hypothesis, with cube potentials
$B_0,\dots,B_q$;
\item the base sites $s_0^{*},\dots,s_q^{*}$ of the overlap hypothesis, where $s_q^{*}=b'_{-}$;
\item the paths $\gamma_0,\dots,\gamma_q$, the transport budgets $\mathfrak{l}_0,\dots,\mathfrak{l}_q$
and the threading loop $\Gamma_{\omega}=\gamma_1\gamma_2\cdots\gamma_q\gamma_0$ of the path hypothesis.
\end{itemize}
Let
\begin{equation*}
\mathfrak{h}_{\omega}=c_0c_1\cdots c_q,\qquad c_i=u_i(s_i^{*})\,u_{i+1}(s_i^{*})^{-1},\qquad
u_{q+1}:=u_0,
\end{equation*}
be the residual element of Proposition~\ref{4.24:prop-product-form}. Put
\begin{equation*}
W_{\omega}:=\operatorname{Hol}_U(\Gamma_{\omega})
:=U(\gamma_1)U(\gamma_2)\cdots U(\gamma_q)U(\gamma_0)\in G .
\end{equation*}
This is the Wilson loop of $\Gamma_{\omega}$ for the $G$-valued factor $U$ of $\mathbb{U}$, with parallel
transport along lattice paths as in Definition~\ref{4.24:not-gauge-transport}. Then there is an
element $E_{\omega}\in G$ with
\begin{equation*}
\operatorname{dis}E_{\omega}\le\sum_{i=0}^{q}\mathfrak{l}_i
\end{equation*}
such that
\begin{equation}\label{4.24:eq-residual-loop-1}
\mathfrak{h}_{\omega}=u_0(s_0^{*})\,\bigl[W_{\omega}E_{\omega}\bigr]\,u_0(s_0^{*})^{-1}.
\end{equation}
Consequently
\begin{equation}\label{4.24:eq-residual-loop-a}
\bigl|\operatorname{dis}\mathfrak{h}_{\omega}-\operatorname{dis}W_{\omega}\bigr|
\le\sum_{i=0}^{q}\mathfrak{l}_i .
\end{equation}
Changing the base sites or the paths $\gamma_i$ changes $W_{\omega}$ by $G$-conjugation and by factors
whose $\operatorname{dis}$ is at most the corresponding budgets $\mathfrak{l}$. In this sense the
statement is insensitive to the choices made within the budget of the path hypothesis.
\end{lemma}

\begin{proof}
Throughout, $\operatorname{dis}$ is the distance to the identity of
Lemma~\ref{4.24:lem-group-distance}. Its properties (d1)--(d5) are displayed in
\eqref{4.24:eq-group-distance-a}. We use property (d2) in four ways:
\begin{itemize}
\item subadditivity, $\operatorname{dis}(h_1h_2)\le\operatorname{dis}h_1+\operatorname{dis}h_2$;
\item invariance under inversion, $\operatorname{dis}(h^{-1})=\operatorname{dis}h$;
\item the equality $\operatorname{dis}(ghg^{-1})=\operatorname{dis}h$ for $g\in G$;
\item the bound $\operatorname{dis}h\le|Z|$ for every logarithm $Z$ of $h$.
\end{itemize}

\emph{Regrouping.} We write out $\mathfrak{h}_{\omega}=c_0c_1\cdots c_q$ with
$c_i=u_i(s_i^{*})u_{i+1}(s_i^{*})^{-1}$ and $u_{q+1}=u_0$. Grouping the factor
$u_i(s_{i-1}^{*})^{-1}$ ending $c_{i-1}$ with the factor $u_i(s_i^{*})$ opening $c_i$, we obtain
\begin{align*}
\mathfrak{h}_{\omega}
&=u_0(s_0^{*})u_1(s_0^{*})^{-1}\,u_1(s_1^{*})u_2(s_1^{*})^{-1}\cdots
u_q(s_q^{*})u_0(s_q^{*})^{-1}\\
&=u_0(s_0^{*})\prod_{i=1}^{q}\Bigl[u_i(s_{i-1}^{*})^{-1}u_i(s_i^{*})\Bigr]\,
u_0(s_q^{*})^{-1}.
\end{align*}

\emph{The gauge rule along a path.} A gauge $u$ acts by $U^{u}(b)=u(b_{-})U(b)u(b_{+})^{-1}$. Let
$\gamma$ be a path from a site $s$ to a site $s'$. The ordered product along $\gamma$ telescopes, so
\begin{equation*}
U^{u}(\gamma)=u(s)\,U(\gamma)\,u(s')^{-1}.
\end{equation*}

\emph{The transports for $i=1,\dots,q$.} Fix such an $i$. The path $\gamma_i$ lies in
$\operatorname{bd}(Y_i)$ and joins $s_{i-1}^{*}$ to $s_i^{*}$. In the cube gauge $u_i$ the transport
along $\gamma_i$ is
\begin{equation*}
e_i:=U^{u_i}(\gamma_i)=\prod_{b\in\gamma_i}e^{\pm i\eta B_i(b)}\in G,
\end{equation*}
where the sign records the orientation in which $\gamma_i$ traverses $b$. It lies in $G$ because $U$
and $u_i$ are $G$-valued. Each factor $e^{\pm i\eta B_i(b)}$ has the logarithm $\pm i\eta B_i(b)$, so
(d2) gives $\operatorname{dis}e^{\pm i\eta B_i(b)}\le|i\eta B_i(b)|$. Subadditivity over the bonds
of $\gamma_i$ then gives
\begin{equation*}
\operatorname{dis}e_i\le\sum_{b\in\gamma_i}|i\eta B_i(b)|=\mathfrak{l}_i .
\end{equation*}
The gauge rule gives $e_i=u_i(s_{i-1}^{*})U(\gamma_i)u_i(s_i^{*})^{-1}$. Solving,
\begin{equation*}
u_i(s_{i-1}^{*})^{-1}u_i(s_i^{*})=U(\gamma_i)\,\tilde e_i^{\,-1},\qquad
\tilde e_i:=u_i(s_i^{*})^{-1}e_i\,u_i(s_i^{*}).
\end{equation*}
To check this, substitute $U(\gamma_i)=u_i(s_{i-1}^{*})^{-1}e_iu_i(s_i^{*})$ on the right. Since
$u_i(s_i^{*})\in G$, the equality case of (d2) gives $\operatorname{dis}\tilde e_i=\operatorname{dis}e_i$.

\emph{The closing transport.} The path $\gamma_0$ lies in $\operatorname{bd}(Y_0)$ and runs from
$s_q^{*}$ to $s_0^{*}$. As above, put $e_0:=U^{u_0}(\gamma_0)\in G$. Then
$\operatorname{dis}e_0\le\mathfrak{l}_0$ and $e_0=u_0(s_q^{*})U(\gamma_0)u_0(s_0^{*})^{-1}$, so
\begin{equation*}
u_0(s_q^{*})^{-1}=U(\gamma_0)\,u_0(s_0^{*})^{-1}e_0^{-1}
=U(\gamma_0)\,\tilde e_0'^{\,-1}\,u_0(s_0^{*})^{-1},
\end{equation*}
where $\tilde e_0':=u_0(s_0^{*})^{-1}e_0\,u_0(s_0^{*})$. By (d2) again,
$\operatorname{dis}\tilde e_0'=\operatorname{dis}e_0$.

Substituting both identities into the regrouped product gives
\begin{equation*}
\mathfrak{h}_{\omega}=u_0(s_0^{*})\Bigl[U(\gamma_1)\tilde e_1^{\,-1}U(\gamma_2)\tilde e_2^{\,-1}
\cdots U(\gamma_q)\tilde e_q^{\,-1}U(\gamma_0)\tilde e_0'^{\,-1}\Bigr]u_0(s_0^{*})^{-1}.
\end{equation*}

\emph{Moving the defects to the right.} For $1\le i\le q$ let
\begin{equation*}
\Pi_i:=U(\gamma_{i+1})\cdots U(\gamma_q)U(\gamma_0)\in G,
\end{equation*}
so that $\Pi_q=U(\gamma_0)$. Thus $\Pi_i$ is the product of the $G$-valued transports that follow
$\tilde e_i^{\,-1}$ in the bracket. Define
\begin{equation*}
E_{\omega}:=\bigl(\Pi_1^{-1}\tilde e_1^{\,-1}\Pi_1\bigr)\bigl(\Pi_2^{-1}\tilde e_2^{\,-1}\Pi_2\bigr)
\cdots\bigl(\Pi_q^{-1}\tilde e_q^{\,-1}\Pi_q\bigr)\,\tilde e_0'^{\,-1}\in G .
\end{equation*}
We have $W_{\omega}\Pi_1^{-1}=U(\gamma_1)$, $\Pi_i\Pi_{i+1}^{-1}=U(\gamma_{i+1})$ for
$1\le i\le q-1$, and $\Pi_q=U(\gamma_0)$. Multiplying out $W_{\omega}E_{\omega}$ with these three
identities reproduces the bracket exactly. This proves \eqref{4.24:eq-residual-loop-1}.

The element $E_{\omega}$ is a product of $q+1$ $G$-conjugates of the elements $\tilde e_i^{\,-1}$ and
$\tilde e_0'^{\,-1}$. By subadditivity, invariance under inversion and the equality case of (d2),
\begin{equation*}
\operatorname{dis}E_{\omega}\le\sum_{i=1}^{q}\operatorname{dis}\tilde e_i+\operatorname{dis}\tilde e_0'
=\sum_{i=0}^{q}\operatorname{dis}e_i\le\sum_{i=0}^{q}\mathfrak{l}_i .
\end{equation*}

\emph{The comparison of distances.} Put
$\mathfrak{h}'_{\omega}:=u_0(s_0^{*})^{-1}\mathfrak{h}_{\omega}u_0(s_0^{*})$. By
\eqref{4.24:eq-residual-loop-1}, $\mathfrak{h}'_{\omega}=W_{\omega}E_{\omega}$, and by (d2),
$\operatorname{dis}\mathfrak{h}'_{\omega}=\operatorname{dis}\mathfrak{h}_{\omega}$. Hence
\begin{align*}
\operatorname{dis}\mathfrak{h}_{\omega}&=\operatorname{dis}(W_{\omega}E_{\omega})
\le\operatorname{dis}W_{\omega}+\operatorname{dis}E_{\omega},\\
\operatorname{dis}W_{\omega}&=\operatorname{dis}(\mathfrak{h}'_{\omega}E_{\omega}^{-1})
\le\operatorname{dis}\mathfrak{h}_{\omega}+\operatorname{dis}E_{\omega}.
\end{align*}
Together with the bound on $\operatorname{dis}E_{\omega}$, these give
\eqref{4.24:eq-residual-loop-a}.

\emph{Dependence on the choices.} Let $\gamma_i'$ be another path in $\operatorname{bd}(Y_i)$ with the
same endpoints as $\gamma_i$. Let its budget be
$\mathfrak{l}_i':=\sum_{b\in\gamma_i'}|i\eta B_i(b)|$, and put $e_i':=U^{u_i}(\gamma_i')$. For every
$i$ the final site of $\gamma_i$ is $s_i^{*}$: for $i\ge1$ by the description of $\gamma_i$, and for
$i=0$ because $\gamma_0$ ends at $s_0^{*}$. Applying the gauge rule to both paths therefore gives
\begin{equation*}
U(\gamma_i)^{-1}U(\gamma_i')=f_i,\qquad f_i:=u_i(s_i^{*})^{-1}e_i^{-1}e_i'\,u_i(s_i^{*}).
\end{equation*}
By the argument used for $e_i$, $\operatorname{dis}f_i\le\mathfrak{l}_i+\mathfrak{l}_i'$.

Replacing $\gamma_i$ by $\gamma_i'$ replaces $U(\gamma_i)$ by $U(\gamma_i)f_i$ in $W_{\omega}$. Moving
$f_i$ past the $G$-valued transports that follow it, exactly as above, multiplies $W_{\omega}$ on the
right by a $G$-conjugate of $f_i$, whose $\operatorname{dis}$ equals $\operatorname{dis}f_i$.

Re-basing, that is, moving a base site within its overlap, likewise multiplies the transports by
transports inside one domain $Y_i$. In the cube gauge $u_i$ these are products of elements
$e^{\pm i\eta B_i(b)}$, so by (d2) their $\operatorname{dis}$ is at most the corresponding budget
$\mathfrak{l}$. The base point of the loop enters only through $G$-conjugation, which by (d2) leaves
$\operatorname{dis}$ unchanged.
\end{proof}

\begin{remark}
The transitions $c_i$ and the loop $\Gamma_{\omega}$ are read with the $G$-valued factor $U$ of
$\mathbb{U}=e^{i\eta A'}U$ (Definition~\ref{4.24:not-gauge-transport}), and not with $\mathbb{U}$.
The complex potential $A'$ is globally defined and satisfies
$|\operatorname{Ad}_uA'|=|A'|$ for every gauge $u$. It is carried inside the cube form of each
factor, as in the cube-smallness hypothesis on the walk term. Reading the loop with
$\mathbb{U}$ would require moving the defects past the complex transports $\mathbb{U}(\gamma)$. For
such conjugations (d2) gives only the bound
$\operatorname{dis}(ghg^{-1})\le e^{2\operatorname{pol}g}\operatorname{dis}h$, so each move would
cost a factor $e^{2\operatorname{pol}\mathbb{U}(\gamma)}$. Since $\operatorname{pol}\mathbb{U}(\gamma)$
grows with the length of $\gamma$, the loop is read with $U$.
\end{remark}

\begin{proposition}[Difference of a walk-term density from its trivial value]
\label{4.24:prop-density-difference}
Assume the hypotheses of Definition~\ref{4.24:def-walk-hypotheses}:
\begin{itemize}
\item the covariance and locality hypothesis \eqref{4.24:eq-walk-hypotheses-a};
\item the majorant hypothesis \eqref{4.24:eq-walk-hypotheses-b};
\item the cube-gauge hypothesis \eqref{4.24:eq-walk-hypotheses-c};
\item the cube smallness hypothesis \eqref{4.24:eq-walk-hypotheses-d};
\item the overlap hypothesis \eqref{4.24:eq-walk-hypotheses-e};
\item the path hypothesis \eqref{4.24:eq-walk-hypotheses-f}.
\end{itemize}
Fix a walk $\omega$, a number $\xi\ge 0$ and a bond $b'$ as in \eqref{4.24:eq-walk-hypotheses-a}, and let
the following objects be as introduced in the indicated statements:
\begin{itemize}
\item the density $\tau_{\omega}$, the chain $S_{\omega}$ and the residual element $\mathfrak{h}_{\omega}\in G$, as in
Proposition~\ref{4.24:prop-product-form};
\item $s^{0}_{\omega}$, $S^{0}_{\omega}$ and $\tau^{0}_{\omega}$, as in Lemma~\ref{4.24:lem-colour-scalars};
\item the Wilson loop $W_{\omega}$, as in Lemma~\ref{4.24:lem-residual-loop};
\item the numbers $\varepsilon$, $\bar\beta$ and $\mathfrak{l}_0,\dots,\mathfrak{l}_q$, as in the cube smallness, overlap and
path hypotheses;
\item the factor count $q=q(\omega)$ and the walk majorant $m^{\xi}(\omega)$, as in the covariance and locality
hypothesis and the majorant hypothesis;
\item the function $\varphi(x)=e^{2x}-1-2x$ of Lemma~\ref{4.24:lem-rotation-trace}.
\end{itemize}
Then the following hold.
\begin{enumerate}
\item[(i)] (Exact split.) We have
\begin{equation}\label{4.24:eq-density-difference-a}
\tau_{\omega}-\tau^{0}_{\omega}
=\operatorname{tr}_{\mathfrak{g}}\bigl[(S_{\omega}-S^{0}_{\omega})\operatorname{Ad}_{\mathfrak{h}_{\omega}}\bigr]
+s^{0}_{\omega}\operatorname{tr}_{\mathfrak{g}}\bigl(\operatorname{Ad}_{\mathfrak{h}_{\omega}}-1\bigr),
\end{equation}
where $S^{0}_{\omega}=s^{0}_{\omega}\cdot 1=(G_0\cdots G_q)(b',b')$. For $|\omega|=0$ the walk term is read in
its subtracted form $\hat t^{\xi}_{\omega}$. Then the subtracted multiples of the identity are the same at the
given pair and at the trivial pair, and they cancel in the difference. Moreover $q=0$ and
$\mathfrak{h}_{\omega}=1$. Thus single-factor terms carry no holonomy, and for them the bound of (iii) reduces to
$n_{\mathfrak{g}}\,m^{\xi}(\omega)\,\varepsilon$.
\item[(ii)] (The chain part.) We have
\begin{equation}\label{4.24:eq-density-difference-1}
\|S_{\omega}-S^{0}_{\omega}\|_{e}\le\bigl[(q+1)\varepsilon+6q\bar\beta\bigr]\,m^{\xi}(\omega).
\end{equation}
\item[(iii)] (The bound.) We have
\begin{equation}\label{4.24:eq-density-difference-b}
\begin{aligned}
|\tau_{\omega}-\tau^{0}_{\omega}|
&\le n_{\mathfrak{g}}\,m^{\xi}(\omega)\Bigl[(q+1)\varepsilon+6q\bar\beta
+\min\bigl\{2,\varphi(\operatorname{dis}\mathfrak{h}_{\omega})\bigr\}\Bigr]\\
&\le n_{\mathfrak{g}}\,m^{\xi}(\omega)\Bigl[(q+1)\varepsilon+6q\bar\beta
+6\Bigl(\operatorname{dis}W_{\omega}+\sum_{i=0}^{q}\mathfrak{l}_i\Bigr)\Bigr].
\end{aligned}
\end{equation}
Moreover, where $\operatorname{dis}\mathfrak{h}_{\omega}\le\frac13$, the holonomy member
$s^{0}_{\omega}\operatorname{tr}_{\mathfrak{g}}(\operatorname{Ad}_{\mathfrak{h}_{\omega}}-1)$ is bounded in absolute
value by $n_{\mathfrak{g}}\,m^{\xi}(\omega)\cdot 4(\operatorname{dis}\mathfrak{h}_{\omega})^{2}$. It is of second
order.
\end{enumerate}
\end{proposition}

\begin{proof}
Throughout, $\omega$, $\xi$ and $b'$ are fixed. We write $\mathfrak{h}=\mathfrak{h}_{\omega}$, $S=S_{\omega}$,
$S^{0}=S^{0}_{\omega}$ and $s^{0}=s^{0}_{\omega}$.

\emph{(i).} By the trace form in Proposition~\ref{4.24:prop-product-form}(ii),
$\tau_{\omega}=\operatorname{tr}_{\mathfrak{g}}[S\operatorname{Ad}_{\mathfrak{h}}]$. By
Lemma~\ref{4.24:lem-colour-scalars}, $S^{0}=s^{0}\cdot 1$ and
$\tau^{0}_{\omega}=\operatorname{tr}_{\mathfrak{g}}S^{0}$. Since $S^{0}$ is a multiple of the identity,
\begin{equation*}
\operatorname{tr}_{\mathfrak{g}}\bigl[S^{0}\operatorname{Ad}_{\mathfrak{h}}\bigr]
=s^{0}\operatorname{tr}_{\mathfrak{g}}\operatorname{Ad}_{\mathfrak{h}}
=s^{0}\operatorname{tr}_{\mathfrak{g}}\bigl(\operatorname{Ad}_{\mathfrak{h}}-1\bigr)
+\operatorname{tr}_{\mathfrak{g}}S^{0}.
\end{equation*}
Hence
$\tau^{0}_{\omega}=\operatorname{tr}_{\mathfrak{g}}[S^{0}\operatorname{Ad}_{\mathfrak{h}}]
-s^{0}\operatorname{tr}_{\mathfrak{g}}(\operatorname{Ad}_{\mathfrak{h}}-1)$. Subtracting this from the
expression for $\tau_{\omega}$ gives \eqref{4.24:eq-density-difference-a}.

Now let $|\omega|=0$. By \eqref{4.24:eq-walk-hypotheses-a}, $q\le 2|\omega|$, so $q=0$. The residual element is
$\mathfrak{h}=c_0=h_0(s_0^{*})$, with $h_0=u_0u_1^{-1}$. In the setting of
Proposition~\ref{4.24:prop-product-form} one has $u_{q+1}=u_0$ and $s_q^{*}=b'_{-}$, so $u_1=u_0$ and
$s_0^{*}=b'_{-}$. Therefore
\begin{equation*}
\mathfrak{h}=c_0=u_0(b'_{-})\,u_0(b'_{-})^{-1}=1 .
\end{equation*}
In the subtracted form of Proposition~\ref{4.24:prop-product-form}(ii), $S$ is replaced by $S$ minus the
multiple of the identity subtracted in the $|\omega|=0$ form of that statement. This multiple is the same at the
given pair and at the trivial pair. It therefore cancels in $S-S^{0}$, and the identity
\eqref{4.24:eq-density-difference-a} holds with $\mathfrak{h}=1$. Its second member then vanishes.

\emph{(ii).} Put $b_0:=b_{q+1}:=b'$. By Proposition~\ref{4.24:prop-product-form}(ii) and
Lemma~\ref{4.24:lem-colour-scalars},
\begin{equation*}
S=\sum_{b_1,\dots,b_q}\Pi_1\Pi_2\cdots \Pi_{2q+1},\qquad
S^{0}=\sum_{b_1,\dots,b_q}\Pi^{0}_1\Pi^{0}_2\cdots \Pi^{0}_{2q+1}.
\end{equation*}
For fixed $b_1,\dots,b_q$ the slots are defined as follows.
\begin{itemize}
\item The factor slots are $\Pi_{2i+1}:=\tilde F_i'(b_i,b_{i+1})$ and $\Pi^{0}_{2i+1}:=G_i(b_i,b_{i+1})$, for
$0\le i\le q$.
\item The transition slots are $\Pi_{2i}:=\operatorname{Ad}_{k'_{i-1}((b_i)_{-})}$ and $\Pi^{0}_{2i}:=1$, for
$1\le i\le q$.
\end{itemize}
Thus $S$ and $S^{0}$ are the same chain with $2q+1$ slots. In it, the factor slot $\tilde F_i'(b_i,b_{i+1})$
stands against $G_i(b_i,b_{i+1})$, and the transition slot $\operatorname{Ad}_{k_i'}$ stands against $1$. In the
$i$-th junction only the values of the transition on the overlap $O_i$ enter
(Proposition~\ref{4.24:prop-product-form}(i)).

We telescope slot by slot, replacing from the left. For elements of $\operatorname{End}(\mathfrak{g}^{\mathbb{C}})$,
\begin{equation}\label{4.24:eq-density-difference-2}
\Pi_1\cdots \Pi_{2q+1}-\Pi^{0}_1\cdots \Pi^{0}_{2q+1}
=\sum_{j=1}^{2q+1}\Pi^{0}_1\cdots \Pi^{0}_{j-1}\,(\Pi_j-\Pi^{0}_j)\,\Pi_{j+1}\cdots \Pi_{2q+1}.
\end{equation}
Indeed, the $j$-th summand equals
\begin{equation*}
\Pi^{0}_1\cdots \Pi^{0}_{j-1}\Pi_j\cdots \Pi_{2q+1}-\Pi^{0}_1\cdots \Pi^{0}_{j}\Pi_{j+1}\cdots \Pi_{2q+1},
\end{equation*}
and the sum over $j$ collapses. We bound each of the $2q+1$ summands by submultiplicativity of $\|\cdot\|_{e}$, using
the following four estimates.

\begin{itemize}
\item \emph{An undifferenced factor slot.} By Proposition~\ref{4.24:prop-product-form}(iii),
$\tilde F_i'=F_i[P_i']$, where $P_i'$ is the given pair transformed by the gauge $C_iu_i$. By the factorwise
bounds (Proposition~\ref{4.24:prop-factorwise-bounds}) and the majorant hypothesis
\eqref{4.24:eq-walk-hypotheses-b},
\begin{equation*}
\|\tilde F_i'(b,c)\|_{e}=\|F_i(b,c)\|_{e}\le\mu_i(b,c).
\end{equation*}
The majorant hypothesis read at the trivial pair gives $\|G_i(b,c)\|_{e}\le\mu_i(b,c)$.
\item \emph{An undifferenced transition slot.} Such a slot is either $1$, or it is
$\operatorname{Ad}_{k_i'(s)}$ with $k_i'(s)\in G$. In the second case it has norm $1$ by
Lemma~\ref{4.24:lem-adjoint-norm}(i).
\item \emph{The differenced factor slot.} By Proposition~\ref{4.24:prop-product-form}(iii),
$\tilde F_i'=F_i[P_i']$ with $P_i'$ a pair in cube form on $Y_i$. The cube smallness hypothesis
\eqref{4.24:eq-walk-hypotheses-d} therefore applies and gives
\begin{equation*}
\|(\tilde F_i'-G_i)(b,c)\|_{e}\le\varepsilon\,\mu_i(b,c).
\end{equation*}
\item \emph{The differenced transition slot.} By the smallness of transition defects
(Lemma~\ref{4.24:lem-transition-defects}), for $s\in O_i$,
\begin{equation*}
\|\operatorname{Ad}_{k_i'(s)}-1\|_{e}\le 6\bar\beta .
\end{equation*}
\end{itemize}

Consequently, in \eqref{4.24:eq-density-difference-2} the $j$-th summand has norm at most
$\mu_0(b',b_1)\mu_1(b_1,b_2)\cdots\mu_q(b_q,b')$ times a factor. That factor is $\varepsilon$ if $j$ is a factor
slot, and $6\bar\beta$ if $j$ is a transition slot. There are $q+1$ factor slots and $q$ transition slots.
Summing over $j$ and then over $b_1,\dots,b_q$ compares the chains with the majorant chain, and gives
\begin{equation*}
\|S-S^{0}\|_{e}
\le\bigl[(q+1)\varepsilon+q\cdot 6\bar\beta\bigr]\,(\mu_0\cdots\mu_q)(b',b')
\le\bigl[(q+1)\varepsilon+6q\bar\beta\bigr]\,m^{\xi}(\omega).
\end{equation*}
The last inequality is the product form of the majorant hypothesis \eqref{4.24:eq-walk-hypotheses-b}. This
proves \eqref{4.24:eq-density-difference-1}.

\emph{(iii).} For $\Psi\in\operatorname{End}(\mathfrak{g}^{\mathbb{C}})$ we have
$|\operatorname{tr}_{\mathfrak{g}}\Psi|\le n_{\mathfrak{g}}\|\Psi\|_{e}$. This holds because $\mathfrak{g}^{\mathbb{C}}$
has dimension $n_{\mathfrak{g}}$ and every eigenvalue of $\Psi$ has modulus at most $\|\Psi\|_{e}$. Since
$\mathfrak{h}\in G$, Lemma~\ref{4.24:lem-adjoint-norm}(i) gives $\|\operatorname{Ad}_{\mathfrak{h}}\|_{e}=1$.
Hence
\begin{equation*}
\bigl|\operatorname{tr}_{\mathfrak{g}}[(S-S^{0})\operatorname{Ad}_{\mathfrak{h}}]\bigr|
\le n_{\mathfrak{g}}\|S-S^{0}\|_{e}\|\operatorname{Ad}_{\mathfrak{h}}\|_{e}
=n_{\mathfrak{g}}\|S-S^{0}\|_{e}.
\end{equation*}
By Lemma~\ref{4.24:lem-colour-scalars}, $|s^{0}|\le m^{\xi}(\omega)$. By the trace of an adjoint rotation
(Lemma~\ref{4.24:lem-rotation-trace}),
\begin{equation*}
|\operatorname{tr}_{\mathfrak{g}}(\operatorname{Ad}_{\mathfrak{h}}-1)|
\le n_{\mathfrak{g}}\min\{2,\varphi(\operatorname{dis}\mathfrak{h})\}.
\end{equation*}
Apply the triangle inequality to \eqref{4.24:eq-density-difference-a}, and insert these two bounds together with
\eqref{4.24:eq-density-difference-1}. This gives the first inequality of
\eqref{4.24:eq-density-difference-b}.

For the second inequality, write $x=\operatorname{dis}\mathfrak{h}\ge 0$. By
Lemma~\ref{4.24:lem-rotation-trace}, $\varphi(x)\le e^{2x}-1$. Next,
\begin{equation*}
\min\{2,\varphi(x)\}\le\min\{2,e^{2x}-1\}\le 6x,
\end{equation*}
where the last step is the elementary inequality of item (d5) of Lemma~\ref{4.24:lem-group-distance}. It is
checked as follows. For $x\ge\frac13$ one has $6x\ge 2$. For $0\le x\le\frac13$, the mean value theorem gives
\begin{equation*}
e^{2x}-1\le 2xe^{2x}\le 2e^{2/3}x\le 6x .
\end{equation*}
By the residual-loop lemma (Lemma~\ref{4.24:lem-residual-loop}),
$x\le\operatorname{dis}W_{\omega}+\sum_{i=0}^{q}\mathfrak{l}_i$. Since $6x$ is increasing in $x$, it follows that
\begin{equation*}
\min\{2,\varphi(x)\}\le 6\Bigl(\operatorname{dis}W_{\omega}+\sum_{i=0}^{q}\mathfrak{l}_i\Bigr).
\end{equation*}
This is the second inequality of \eqref{4.24:eq-density-difference-b}.

Finally, let $x\le\frac13$. Lemma~\ref{4.24:lem-rotation-trace} gives $\varphi(x)\le 4x^{2}$ on
$[0,\frac13]$. Hence
\begin{equation*}
|s^{0}\operatorname{tr}_{\mathfrak{g}}(\operatorname{Ad}_{\mathfrak{h}}-1)|
\le m^{\xi}(\omega)\,n_{\mathfrak{g}}\,\varphi(x)
\le n_{\mathfrak{g}}\,m^{\xi}(\omega)\cdot 4(\operatorname{dis}\mathfrak{h})^{2}.
\end{equation*}
\end{proof}

\begin{remark}
Three questions about a walk term are answered separately.
\begin{itemize}
\item[(a)] \emph{Value.} The walk term may be evaluated in any gauge.
\item[(b)] \emph{Bounds.} Bounds are obtained factor by factor, each factor on its own cube. They compose through
isometric transitions, and nothing is required of the transitions
(Proposition~\ref{4.24:prop-factorwise-bounds}).
\item[(c)] \emph{Difference.} The difference from the trivial pair composes through the transitions at the
cost of their defects $k_i$ only, never of the transitions $h_i$ themselves; this is the content of
Proposition~\ref{4.24:prop-density-difference}. The defects satisfy $\operatorname{dis}k_i(s)\le\bar\beta$ for
all $s\in O_i$, a bound derived from the cube-gauge hypothesis on the overlaps
(Lemma~\ref{4.24:lem-transition-defects}).
\end{itemize}
There is one exception to (c): the residual element $\mathfrak{h}_{\omega}$. The trace over the closed walk does
not remove it. The net conjugation around $b'\to\cdots\to b'$ is $\operatorname{Ad}_{\mathfrak{h}_{\omega}}$, and
$\operatorname{tr}_{\mathfrak{g}}[S_{\omega}\operatorname{Ad}_{\mathfrak{h}_{\omega}}]
\neq\operatorname{tr}_{\mathfrak{g}}S_{\omega}$ unless
$S_{\omega}$ is a colour scalar or $\mathfrak{h}_{\omega}$ is central. Here $S_{\omega}$ is not a colour scalar,
since it is the field-dependent chain. The precise statement that replaces ``the trace removes the net
conjugation'' is \eqref{4.24:eq-density-difference-a}. The trace removes the net conjugation up to the term
$s^{0}_{\omega}\operatorname{tr}_{\mathfrak{g}}(\operatorname{Ad}_{\mathfrak{h}_{\omega}}-1)$, plus second-order
cross terms that are already contained in
$\operatorname{tr}_{\mathfrak{g}}[(S_{\omega}-S^{0}_{\omega})\operatorname{Ad}_{\mathfrak{h}_{\omega}}]$.
\end{remark}

\begin{definition}[The holonomy hypothesis on threading loops]\label{4.24:def-holonomy-hypothesis}
Consider the two walk expansions of this section: the expansion at $\Lambda'$ (the full stage) and
the expansion at $\Lambda'_s$ (the block stage).
For a walk $\omega$ of either expansion, let $\Gamma_{\omega}=\gamma_1\gamma_2\cdots\gamma_q\gamma_0$ be the threading loop
of $\omega$, with $q=q(\omega)$ as in Definition~\ref{4.24:def-walk-hypotheses}, so that the walk
term is the product $F_0F_1\cdots F_q$. It is the closed lattice loop at the base site $s^{*}_0$ inside $D(\omega)$ built from the paths $\gamma_i$ of the structural hypotheses on a walk term (Definition~\ref{4.24:def-walk-hypotheses}, display~\eqref{4.24:eq-walk-hypotheses-f}), for any choice of the base sites and of the paths permitted there. Following Lemma~\ref{4.24:lem-residual-loop}, write
\begin{equation*}
\operatorname{Hol}_U(\Gamma_{\omega})
 := U(\gamma_1)\,U(\gamma_2)\cdots U(\gamma_q)\,U(\gamma_0)\in G
\end{equation*}
for the holonomy of $U$ around $\Gamma_{\omega}$. Here $U(\gamma)$ is the parallel transport along the path $\gamma$ (Definition~\ref{4.24:not-gauge-transport}).

The \emph{holonomy hypothesis on threading loops} is the following requirement. There exist an
absolute number $C_{\mathfrak{w}}$ and a number $\mathfrak{w}_0$, and for every walk $\omega$ of
either expansion a number $\mathfrak{w}(\omega)$, such that, with
$\Gamma_{\omega}$ given by any admissible choice of base sites and paths
in~\eqref{4.24:eq-walk-hypotheses-f} (by Lemma~\ref{4.24:lem-residual-loop}, two such choices
give holonomies whose values of $\operatorname{dis}$ differ by at most the corresponding
transport budgets $\mathfrak{l}_i$),
\begin{equation}\label{4.24:eq-holonomy-hypothesis-a}
\operatorname{dis}\operatorname{Hol}_U(\Gamma_{\omega})
 \le \mathfrak{w}(\omega)
 \le \mathfrak{w}_0\, C_{\mathfrak{w}}^{|\omega|}.
\end{equation}
In addition, $\mathfrak{w}_0$ must lie within the budget of the smallness parameter
$\tilde\alpha_0$ against which the closure (holonomy) term of a density spanning more than one of
the cubes carrying one gauge each is measured; here $w_{*}$, $B$, $C$ and $M$ denote the fixed
factors multiplying $\tilde\alpha_0$ in that budget. The requirement reads
\begin{equation*}
\mathfrak{w}_0 \le O(1)\cdot 5\, w_{*} B C M \tilde\alpha_0 .
\end{equation*}

A per-walk bound of the form $\mathfrak{w}_0\,P(|\omega|)$, with $P$ a polynomial that is
non-negative on the positive integers, is of the form of the right-hand side
of~\eqref{4.24:eq-holonomy-hypothesis-a} for the walks with $|\omega|\ge 1$. One may take
\begin{equation*}
C_{\mathfrak{w}} := \sup_{p} P(p)^{1/p},
\end{equation*}
the supremum being taken over the positive integers $p$; it is finite, since $P(p)\le c\,p^{n}$
for some $c\ge 1$ and $n\ge 0$, and $(c\,p^{n})^{1/p}\le c\,e^{n/e}$. Then $P(p)\le C_{\mathfrak{w}}^{p}$ for every such $p$, hence $\mathfrak{w}_0\,P(|\omega|)\le\mathfrak{w}_0\,C_{\mathfrak{w}}^{|\omega|}$
whenever $|\omega|\ge 1$.
\end{definition}

Display~\eqref{4.24:eq-holonomy-hypothesis-a} is a hypothesis of this section and is not proved here. Under it, clause~(iii) of the proposition bounding the residual elements of the walk terms is the required termwise bound on the walk terms, with the enlarged per-walk constant given later in this section. The situation without~\eqref{4.24:eq-holonomy-hypothesis-a} is treated at the end of this section.

\begin{proposition}[Walk-term field dependence under the holonomy hypothesis]
\label{4.24:prop-field-dependence}
Fix the block $\Lambda'_s$; write $\operatorname{bd}(\Lambda'_s)$ for its bond set and $P_s$ for the
projection of $\mathcal{H}$ onto the fibres over $\operatorname{bd}(\Lambda'_s)$, so that
$(\cdot)_{ss}=P_s(\cdot)P_s$ is the block compression. The following
objects and properties of the block densities are hypotheses of this proposition; they are fixed in
the walk expansion at the block and in the definition and count of the block densities, and are not
constructed here.
\begin{itemize}
\item[(o1)] The block compressions of the operators of the step have the walk expansion of
\cite[Theorem~3.10]{B-CMP99b} at the block, with walk terms $t^{\xi}_{\omega}$, $\xi\ge0$, and walk
majorant $\bar m(\omega)$ as in (vii) below. Put $\hat t^{\xi}_{\omega}:=t^{\xi}_{\omega}$ for
$|\omega|\ge1$; for the class $\omega_0$ of walks with $|\omega_0|=0$, let
$\hat t^{\xi}_{\omega_0}$ be $t^{\xi}_{\omega_0}$ minus $(\lambda_0+\xi)^{-1}$ times the share of
$\mathbf{1}$ carried by $\omega_0$, with the constant $\lambda_0$ of the walk expansion.
\item[(o2)] $\mathbb{C}^{s}_{G}=(\mathbb{C}^{s}_{G}(X))_{X}$ is the family of block densities at the
block $\Lambda'_s$. Apart from its single-cube members, each member $\mathbb{C}^{s}_{G}(X)$ is the sum
over the bonds $b'\in\operatorname{bd}(\Lambda'_s)$ and over the walks $\omega$ whose localisation
domain is $X$ of
\begin{equation*}
\tfrac12\int_0^\infty d\xi\,\operatorname{tr}_{\mathfrak{g}}\hat t^{\xi}_{\omega}(b',b').
\end{equation*}
Each single-cube member involves one block pair on one cube.
\item[(o3)] In the count of block densities, each cell of the partition $\pi_{j+1}$ of the count
carries $c_{G}'$ densities, with $\pi_{j+1}$ and $c_{G}'$ as in the definition of the normalisation
constants below; $\|\cdot\|_{\mathrm{wt}}$ is the weighted norm in which the count
measures such families, with weights at rate at least $8\kappa L$. The universal series of the
count is a series over walks in which each walk $\omega$ enters through its majorant
$\bar m(\omega)$. For every per-step constant $c'>0$, the series formed in the same way with
$\bar m(\omega)$ replaced by $O(1)(c'M^{-1/2})^{|\omega|}e^{-\delta_0d(\omega)}$, with $\delta_0$
and $d(\omega)$ as in (vii) below, converges whenever $M$ satisfies the standing lower bounds of the
count at per-step constant $c'$; among these bounds, only the hypothesis ``$M$ sufficiently large''
of \cite[Theorem~3.10]{B-CMP99b} depends on the per-step constant. $M$ satisfies the standing lower
bounds of the count at the per-step constant $c$ of (vii) below. The count bounds
the weighted norm as follows. Let a family be formed as in (o2), its members other than the
single-cube ones being sums over $b'\in\operatorname{bd}(\Lambda'_s)$ and over walks $\omega$ with
localisation domain $X$ of
$\tfrac12\int_0^\infty d\xi\,\sigma^{\xi}_{\omega}(b')$, with complex numbers
$\sigma^{\xi}_{\omega}(b')$. Let $\varpi\ge0$ and $\bar m'(\omega)\ge0$ be numbers such that, for every
walk $\omega$,
\begin{equation*}
\sum_{b'\in\operatorname{bd}(\Lambda'_s)}\tfrac12\int_0^\infty d\xi\,|\sigma^{\xi}_{\omega}(b')|
\le\varpi\,\bar m'(\omega).
\end{equation*}
Then the weighted norm of the family, single-cube members excluded, is at most $c_{G}'\,\varpi$ times
the universal series of the count with $\bar m(\omega)$ replaced by $\bar m'(\omega)$.
\item[(o4)] The walk terms are integrable in $\xi$ over $[0,\infty)$ termwise. The majorants
$m^{\xi}(\omega)$ of hypothesis (ii) below are integrable in $\xi$ over $[0,\infty)$, and for every
walk $\omega$
\begin{equation*}
\sum_{b'\in\operatorname{bd}(\Lambda'_s)}\tfrac12\int_0^\infty m^{\xi}(\omega)\,d\xi\le\bar m(\omega),
\end{equation*}
the term through which $\omega$ enters the universal series of the count.
\item[(o5)] The field-dependence estimate for the block densities on one cube holds: it bounds the
difference of each single-cube member from its value at the trivial pair.
\end{itemize}
Write $\mathbb{C}^{s}_{G}\big|_{(1,0)}$ for the same family at the trivial pair
$(\mathbb{U},\mathbb{J})=(1,0)$. Assume moreover the following.
\begin{itemize}
\item[(i)] The walk terms $t^{\xi}_{\omega}$ of the walk expansion of \cite[Theorem~3.10]{B-CMP99b}
compressed to the block are products of the factors $F_i$ of \cite[(3.107)]{B-CMP99b},
$0\le i\le q(\omega)$, on factor domains $Y_i$, with factor count $q(\omega)\le2|\omega|$. These
factors satisfy the termwise gauge
covariance and the locality conditions among the structural hypotheses on a walk term.
\item[(ii)] The factors have majorants $\mu_i$ in product form, valid at the cube-form
representative $(\mathbb{U}^{u_i},\operatorname{Ad}_{u_i}\mathbb{J})$ on $Y_i$, where $u_i$ is
the cube gauge of hypothesis (v) on $Y_i$, and at the trivial pair, with
$(\mu_0\cdots\mu_{q(\omega)})(b',b')\le m^{\xi}(\omega)$.
\item[(iii)] Each factor satisfies the per-cube smallness condition
\eqref{4.24:eq-walk-hypotheses-d} with the number $\varepsilon$.
\item[(iv)] The overlap condition \eqref{4.24:eq-walk-hypotheses-e} holds with transition budgets
$\beta_i\le\bar\beta$. The path condition \eqref{4.24:eq-walk-hypotheses-f} holds with transport
budgets $\mathfrak{l}_i\le\beta'$.
\item[(v)] The cube-gauge condition among the structural hypotheses on a walk term holds on each
factor domain $Y_i$.
\item[(vi)] The holonomy hypothesis of Definition~\ref{4.24:def-holonomy-hypothesis} holds, that
is, \eqref{4.24:eq-holonomy-hypothesis-a} holds with budget $\mathfrak{w}(\omega)$ and constants
$\mathfrak{w}_0$, $C_{\mathfrak{w}}$, where $C_{\mathfrak{w}}\ge1$; there is no loss in this, since
enlarging $C_{\mathfrak{w}}$ preserves \eqref{4.24:eq-holonomy-hypothesis-a}.
\item[(vii)] Write the walk majorant $\bar m(\omega)$ of \cite[Theorem~3.10]{B-CMP99b} in the form
$O(1)(cM^{-1/2})^{|\omega|}e^{-\delta_0 d(\omega)}$, with per-step constant $c$, where
$d(\omega)$ is the walk quantity in the exponent of that majorant, not to be confused with the
dimension $d$. Assume that $M$ satisfies the hypothesis ``$M$ sufficiently large'' of that
theorem with the per-step constant $3C_{\mathfrak{w}}c$ in place of $c$; this holds when
$M\ge(3C_{\mathfrak{w}})^2M_0(\kappa,L,d)$, where $M_0(\kappa,L,d)$ is the threshold at
per-step constant $c$. It is a one-sided lower bound on $M$, imposed after $d$, $L$ and $\kappa$,
on which $M_0(\kappa,L,d)$ depends.
\end{itemize}
Then the following hold.
\begin{itemize}
\item[(a)] For every walk $\omega$, every $\xi\ge0$ and every $b'$, the density
$\operatorname{tr}_{\mathfrak{g}}\hat t^{\xi}_{\omega}(b',b')$ is a gauge-invariant function of
$(\mathbb{U},\mathbb{J})$ and may be computed in any gauge. Each factor's absolute bound may be
established in that factor's own cube gauge $u_i$ on $Y_i$, and these bounds compose to the chain
bound by $m^{\xi}(\omega)$. No single gauge realising all cube gauges at once is asserted to exist.
\item[(b)] For every walk $\omega$, every $\xi\ge0$ and every $b'$,
\begin{equation}\label{4.24:eq-field-dependence-1}
\begin{split}
&\Bigl|\operatorname{tr}_{\mathfrak{g}}\hat t^{\xi}_{\omega}(b',b')
-\bigl(\operatorname{tr}_{\mathfrak{g}}\hat t^{\xi}_{\omega}(b',b')\bigr)\big|_{(1,0)}\Bigr|\\
&\qquad\le n_{\mathfrak{g}}\Bigl[(q(\omega)+1)\varepsilon+6q(\omega)\bar\beta
+6\sum_{i=0}^{q(\omega)}\mathfrak{l}_i+6\mathfrak{w}(\omega)\Bigr]\,m^{\xi}(\omega).
\end{split}
\end{equation}
\item[(c)] Let $\mathfrak{S}_{\partial}$ be the universal series of the count of block densities
with $\bar m(\omega)$ replaced by $C_{\mathfrak{w}}^{|\omega|}3^{|\omega|}\bar m(\omega)$. Then
$\mathfrak{S}_{\partial}$ is finite, by (o3) with $c'=3C_{\mathfrak{w}}c$, whose standing lower
bounds at that constant hold by (o3) and (vii). Let $c_{G}'$ be the number of densities per cell of
the partition $\pi_{j+1}$, and let $\Delta_{\mathrm{sc}}\ge0$ denote the contribution to the weighted
norm of the differences of the single-cube members from their values at the trivial pair. Then
\begin{equation}\label{4.24:eq-field-dependence-a}
\bigl\|\mathbb{C}^{s}_{G}-\mathbb{C}^{s}_{G}\big|_{(1,0)}\bigr\|_{\mathrm{wt}}
\le c_{G}'\,n_{\mathfrak{g}}\,\mathfrak{S}_{\partial}
\bigl[\varepsilon+6\bar\beta+6\beta'+6\mathfrak{w}_0\bigr]+\Delta_{\mathrm{sc}}.
\end{equation}
This contribution $\Delta_{\mathrm{sc}}$ is bounded directly by the field-dependence estimate for
the block densities on one cube.
\end{itemize}
Assume in addition the following hypothesis.
\begin{itemize}
\item[(viii)] The right member of \eqref{4.24:eq-field-dependence-a} obeys
\begin{equation}\label{4.24:eq-field-dependence-2}
\begin{split}
&c_{G}'\,n_{\mathfrak{g}}\,\mathfrak{S}_{\partial}
\bigl[\varepsilon+6\bar\beta+6\beta'+6\mathfrak{w}_0\bigr]+\Delta_{\mathrm{sc}}\\
&\qquad\le c_{G}'\,O(1)\bigl[5w_{*}(BCM\tilde\alpha_0+\tilde\alpha_1)+K_{\Delta}\gamma_0\bigr]
\le\vartheta.
\end{split}
\end{equation}
This is the bound asserted by the field clause of the field-dependence estimate for the block
densities, whose bracket is the one displayed, with its constants $w_{*}$, $\tilde\alpha_0$,
$\tilde\alpha_1$, $K_{\Delta}$, $\gamma_0$ and the number $\vartheta$ as in the definition of the
normalisation constants below, and with $B$ and $C$ the two constants of that bracket
which, together with $M$, multiply $\tilde\alpha_0$; in that estimate the holonomy
hypothesis \eqref{4.24:eq-holonomy-hypothesis-a} enters as a hypothesis; that estimate is assumed,
not derived, in this proposition.
\end{itemize}
Then, under \textup{(o1)}--\textup{(o5)} and \textup{(i)}--\textup{(viii)},
$\bigl\|\mathbb{C}^{s}_{G}-\mathbb{C}^{s}_{G}\big|_{(1,0)}\bigr\|_{\mathrm{wt}}\le\vartheta$.
\end{proposition}

\begin{proof}
\emph{(a) Gauge.} By (i), the walk terms are termwise gauge covariant and their factors are local.
These are the hypotheses of the gauge invariance of walk-term densities
(Proposition~\ref{4.24:prop-density-invariance}). That proposition gives, for every gauge $u$,
\begin{equation*}
\operatorname{tr}_{\mathfrak{g}}\hat t^{\xi}_{\omega}(b',b')[\mathbb{U}^{u},\operatorname{Ad}_u\mathbb{J}]
=\operatorname{tr}_{\mathfrak{g}}\hat t^{\xi}_{\omega}(b',b')[\mathbb{U},\mathbb{J}].
\end{equation*}
So the density is gauge-invariant and may be computed in any gauge.

The same two hypotheses are those of the factorwise kernel bounds in cube gauges
(Proposition~\ref{4.24:prop-factorwise-bounds}). By that proposition, for every gauge $u$,
\begin{equation*}
\|F_i[\mathbb{U}^{u},\operatorname{Ad}_u\mathbb{J}](b,c)\|_{e}=\|F_i[\mathbb{U},\mathbb{J}](b,c)\|_{e}.
\end{equation*}
Hence a majorant $\mu_i$ valid at the cube-form representative
$(\mathbb{U}^{u_i},\operatorname{Ad}_{u_i}\mathbb{J})$ on $Y_i$ holds at $(\mathbb{U},\mathbb{J})$
itself. Part (ii) of that proposition, together with hypothesis (ii), then composes these bounds
factor by factor into
\begin{equation*}
\|t^{\xi}_{\omega}(b',b')\|_{e}\le(\mu_0\cdots\mu_{q(\omega)})(b',b')\le m^{\xi}(\omega).
\end{equation*}
Each factor is used in its own cube gauge $u_i$ only, so no gauge realising all the $u_i$ at once
is needed.

\emph{(b) Difference.} Write $\tau_\omega$ for the density
$\operatorname{tr}_{\mathfrak{g}}\hat t^{\xi}_{\omega}(b',b')$ and $\tau^{0}_\omega$ for its value
at the trivial pair.

By (i) and (v), the hypotheses of the product form of a walk term with transition defects
(Proposition~\ref{4.24:prop-product-form}) hold. By its trace form
\eqref{4.24:eq-product-form-a}, $\tau_\omega=\operatorname{tr}_{\mathfrak{g}}[S_\omega\cdot
\operatorname{Ad}_{\mathfrak{h}_\omega}]$, with the modification stated there for $|\omega|=0$.
Here $S_\omega$ is the chain of cube-form factors and transition defects, and
$\mathfrak{h}_\omega$ is the residual element.

By (v) and (iv), the smallness of transition defects on overlaps
(Lemma~\ref{4.24:lem-transition-defects}) gives, for the transition defects $k_i'$ of
Proposition~\ref{4.24:prop-product-form} on the overlaps $O_i$,
$\|\operatorname{Ad}_{k_i'(x)}-1\|_{e}\le6\bar\beta$ for every site $x\in O_i$.

By (v) and (iv), the residual element as a Wilson loop (Lemma~\ref{4.24:lem-residual-loop}) gives,
with $W_\omega=\operatorname{Hol}_U(\Gamma_\omega)$ as in \eqref{4.24:eq-residual-loop-a},
\begin{equation*}
|\operatorname{dis}\mathfrak{h}_\omega-\operatorname{dis}W_\omega|
\le\sum_{i=0}^{q(\omega)}\mathfrak{l}_i .
\end{equation*}

Hypotheses (i)--(v) together are those of the difference of a walk-term density from its trivial
value (Proposition~\ref{4.24:prop-density-difference}). Its bound \eqref{4.24:eq-density-difference-b}
gives
\begin{equation*}
|\tau_\omega-\tau^{0}_\omega|
\le n_{\mathfrak{g}}\,m^{\xi}(\omega)\Bigl[(q(\omega)+1)\varepsilon+6q(\omega)\bar\beta
+6\Bigl(\operatorname{dis}W_\omega+\sum_{i=0}^{q(\omega)}\mathfrak{l}_i\Bigr)\Bigr].
\end{equation*}
By (vi), $\operatorname{dis}W_\omega=\operatorname{dis}\operatorname{Hol}_U(\Gamma_\omega)
\le\mathfrak{w}(\omega)$ for the threading loop $\Gamma_\omega$, whatever choice is made within the
path condition. Inserting this yields \eqref{4.24:eq-field-dependence-1}.

The argument uses the per-cube condition \eqref{4.24:eq-walk-hypotheses-d} factor by factor and
nothing else of the field-dependence estimate. In particular it uses no bound on the global
difference $\delta\mathcal{A}:=\mathcal{A}-\mathcal{A}|_{(1,0)}$ of the quadratic form. The
inequality
\begin{equation*}
\|(\delta\mathcal{A})_{ss}\|\le\|\delta\mathcal{A}\|\le O(1)(BCM\tilde\alpha_0+\tilde\alpha_1)
\end{equation*}
is false in a general gauge and is not used.

\emph{(c) Absorption of the multiplicities.} Put $p=|\omega|$. We use three facts.
\begin{itemize}
\item By (i), $q(\omega)\le2p$.
\item By (iv), $\sum_{i=0}^{q(\omega)}\mathfrak{l}_i\le(q(\omega)+1)\beta'$.
\item By (vi), $\mathfrak{w}(\omega)\le\mathfrak{w}_0C_{\mathfrak{w}}^{p}$.
\end{itemize}
Hence the bracket in \eqref{4.24:eq-field-dependence-1} is at most
\begin{equation*}
(q+1)\varepsilon+6q\bar\beta+6(q+1)\beta'+6\mathfrak{w}_0C_{\mathfrak{w}}^{p},
\qquad q=q(\omega).
\end{equation*}
Each of $q+1$ and $q$ is at most $2p+1$. Moreover $2p+1\le3^{p}$ for $p\ge1$. For $p=0$ we have
$q=0$ by Proposition~\ref{4.24:prop-density-difference}(i), so the coefficients equal $1$ and $0$.
Since $C_{\mathfrak{w}}\ge1$ by (vi), we have
$3^{p}\le(3C_{\mathfrak{w}})^{p}$ and $C_{\mathfrak{w}}^{p}\le(3C_{\mathfrak{w}})^{p}$. Therefore
\begin{equation}\label{4.24:eq-field-dependence-3}
\begin{split}
n_{\mathfrak{g}}\bigl[(q+1)\varepsilon+6q\bar\beta+6(q+1)\beta'
+6\mathfrak{w}_0C_{\mathfrak{w}}^{p}\bigr]
&\le n_{\mathfrak{g}}\bigl[\varepsilon+6\bar\beta+6\beta'+6\mathfrak{w}_0\bigr]
(3C_{\mathfrak{w}})^{p}.
\end{split}
\end{equation}
A majorant of the form $O(1)(cM^{-1/2})^{p}e^{-\delta_0d(\omega)}$, multiplied by
$(3C_{\mathfrak{w}})^{p}$, becomes
$O(1)(3C_{\mathfrak{w}}c\,M^{-1/2})^{p}e^{-\delta_0d(\omega)}$. So the differenced majorant has the
form of the walk majorant of \cite[Theorem~3.10]{B-CMP99b}. Its per-step constant is multiplied by
the absolute number $3C_{\mathfrak{w}}$, and its rate $\delta_0$ is unchanged.

Two consequences follow.
\begin{itemize}
\item[($\alpha$)] The convergences obtained from that form, namely the sums over lattice animals and
hulls and the weighted norm at rate at least $8\kappa L$, all of which enter the universal series of
the count, hold with $c$ replaced by $3C_{\mathfrak{w}}c$: this is (o3) applied with
$c'=3C_{\mathfrak{w}}c$, provided $M$ satisfies the standing lower bounds of the count at that
per-step constant.
\item[($\beta$)] The one hypothesis of that form which involves the per-step constant is
``$M$ sufficiently large'' of \cite[Theorem~3.10]{B-CMP99b}. Since that hypothesis is a smallness
condition on the per-step quantity $cM^{-1/2}$, imposing it with $3C_{\mathfrak{w}}c$ in place of
$c$ enlarges the threshold on $M$ by the factor $(3C_{\mathfrak{w}})^2$ at most. This is provided
by (vii).
It is a lower bound on $M$ only. It involves neither the step, nor the volume, nor the coupling. In
particular it is not a cap in the sense of Definition~\ref{2.3:def-cap}, and it is of the same
kind as the floor $M\ge M_0(\kappa,L,d)$ among the standing lower bounds of the count.
\end{itemize}

\emph{Count and sum.} Each cell of $\pi_{j+1}$ carries $c_{G}'$ densities, as in (o3). By
\eqref{4.24:eq-field-dependence-1} and \eqref{4.24:eq-field-dependence-3}, for every walk $\omega$,
every $\xi\ge0$ and every $b'\in\operatorname{bd}(\Lambda'_s)$,
\begin{equation*}
|\tau_\omega-\tau^{0}_\omega|\le\varpi\,(3C_{\mathfrak{w}})^{|\omega|}\,m^{\xi}(\omega),
\qquad \varpi:=n_{\mathfrak{g}}\bigl[\varepsilon+6\bar\beta+6\beta'+6\mathfrak{w}_0\bigr].
\end{equation*}
Integrate over $\xi$ termwise by (o4), which gives the $\xi$-integrability of the walk terms and of
their majorants $m^{\xi}(\omega)$ and is unaffected by the factor $(3C_{\mathfrak{w}})^{|\omega|}$,
and sum over $b'\in\operatorname{bd}(\Lambda'_s)$. By the inequality of (o4),
\begin{equation*}
\sum_{b'\in\operatorname{bd}(\Lambda'_s)}\tfrac12\int_0^\infty d\xi\,|\tau_\omega-\tau^{0}_\omega|
\le\varpi\,(3C_{\mathfrak{w}})^{|\omega|}\,\bar m(\omega).
\end{equation*}
By (o2), the members of $\mathbb{C}^{s}_{G}-\mathbb{C}^{s}_{G}\big|_{(1,0)}$ other than the
single-cube ones are the sums over $b'\in\operatorname{bd}(\Lambda'_s)$ and over the walks with
localisation domain $X$ of
$\tfrac12\int_0^\infty d\xi\,(\tau_\omega-\tau^{0}_\omega)$. Apply the bound of the count in (o3) to
this family, with $\sigma^{\xi}_{\omega}(b')=\tau_\omega-\tau^{0}_\omega$, with the number $\varpi$
above and with $\bar m'(\omega):=(3C_{\mathfrak{w}})^{|\omega|}\bar m(\omega)$. The universal series
of the count with $\bar m(\omega)$ replaced by $(3C_{\mathfrak{w}})^{|\omega|}\bar m(\omega)$ is by
definition $\mathfrak{S}_{\partial}$, and by ($\alpha$) and ($\beta$) it is finite. So the weighted
norm of this part of the family is at most
$c_{G}'\,n_{\mathfrak{g}}\,\mathfrak{S}_{\partial}[\varepsilon+6\bar\beta+6\beta'+6\mathfrak{w}_0]$.

The walk terms with $|\omega|=0$ have $q=0$ and $\mathfrak{h}_{\omega}=1$ by
Proposition~\ref{4.24:prop-density-difference}(i), so they carry no holonomy; they are already
contained in the series. The single-cube members are separate additive terms of the family. Each
involves one block pair on one cube, hence no transition and no threading loop, and their
contribution $\Delta_{\mathrm{sc}}$ is bounded by the field-dependence estimate on one cube
directly. Adding this contribution gives \eqref{4.24:eq-field-dependence-a}.

Under (viii), the right member of \eqref{4.24:eq-field-dependence-a} is at most $\vartheta$ by
\eqref{4.24:eq-field-dependence-2}, which is the final assertion.
\end{proof}

\begin{remark}
An alternative absorption of the multiplicities leads to a condition of the same kind.

The factor $2p+1$ can be absorbed into $e^{-\frac12\delta_0d(\omega)}$, since $d(\omega)$
grows at least proportionally to $p=|\omega|$. This enlarges the rate condition
$(M/M_1)\min\{\delta_0,\delta\}\ge8\kappa\sqrt d$ of the standing floors on $M$ (those which also
carry $M\ge M_0(\kappa,L,d)$), in which $\delta$ is the second decay rate entering that clause
alongside $\delta_0$, to
$(M/M_1)\min\{\delta_0,\delta\}\ge16\kappa\sqrt d$, again a lower bound on $M$, chosen after
$\kappa$. This route is not followed above.

Separately, the argument that expands the resolvent difference
\begin{equation*}
-(\mathcal{A}+\xi)^{-1}\,\delta\mathcal{A}\,(\mathcal{A}|_{(1,0)}+\xi)^{-1},
\end{equation*}
which is not used in the proof of (b), meets the same multiplicity: expanding both resolvents in
walks gives a sum over walks with one marked factor, with $q+1$ markings per walk.
\end{remark}

\begin{remark}[The global difference of the form is not gauge covariant]\label{4.24:rem-noncovariant-difference}
Let $\mathcal{A}=\mathcal{A}[\mathbb{U},\mathbb{J}]$ be the quadratic form of
Definition~\ref{4.4:def-quadratic-form}, a kernel on
$\mathcal{H}=\bigoplus_{b}\mathfrak{g}^{\mathbb{C}}_b$, and let
$\mathcal{U}_u=\bigoplus_b\operatorname{Ad}_{u(b_{-})}$ be the action of a gauge $u$ on $\mathcal{H}$.
Put
\begin{equation}\label{4.24:eq-noncovariant-difference-1}
\delta\mathcal{A}[\mathbb{U},\mathbb{J}]:=\mathcal{A}[\mathbb{U},\mathbb{J}]-\mathcal{A}[1,0].
\end{equation}
\begin{enumerate}
\item[(a)] For every gauge $u$,
\begin{equation}\label{4.24:eq-noncovariant-difference-2}
\delta\mathcal{A}[\mathbb{U}^{u},\operatorname{Ad}_u\mathbb{J}]
=\mathcal{U}_u\mathcal{A}[\mathbb{U},\mathbb{J}]\mathcal{U}_u^{-1}-\mathcal{A}[1,0],
\end{equation}
and this equals $\mathcal{U}_u\,\delta\mathcal{A}[\mathbb{U},\mathbb{J}]\,\mathcal{U}_u^{-1}$ if and only
if $\mathcal{U}_u$ commutes with $\mathcal{A}[1,0]$. The latter holds if and only if
$u(b_{-})u(c_{-})^{-1}$ is central in $G$ for every pair of bonds $(b,c)$ with
$\mathcal{A}[1,0](b,c)\neq0$, which fails for a generic non-constant $u$.
\item[(b)] Consequently $\|\delta\mathcal{A}\|$ depends on the gauge. The gauge-invariant object is
the \emph{orbit distance}
\begin{equation}\label{4.24:eq-noncovariant-difference-3}
\eth(\mathbb{U},\mathbb{J}):=\inf_{u}\bigl\|\mathcal{A}[\mathbb{U}^{u},\operatorname{Ad}_u\mathbb{J}]
-\mathcal{A}[1,0]\bigr\|.
\end{equation}
\item[(c)] Let $v$ be a gauge such that $v(b_{-})v(c_{-})^{-1}$ is non-central for some $(b,c)$ with
$\mathcal{A}[1,0](b,c)\neq0$, and let $1^{v}(b)=v(b_{-})v(b_{+})^{-1}$. The pair $(1^{v},0)$ is
gauge-equivalent to the trivial pair, so every walk-term density has at $(1^{v},0)$ the same value as
at the trivial pair, and every cube is in cube form with vanishing cube potentials in the gauge $v^{-1}$;
yet $\delta\mathcal{A}[1^{v},0]\neq0$, with norm up to $2\|\mathcal{A}[1,0]\|$ and no small factor.
\end{enumerate}
In particular the inequality $\|(\delta\mathcal{A})_{ss}\|\le\|\delta\mathcal{A}\|$ for the block
compression $(\delta\mathcal{A})_{ss}=P_s\,\delta\mathcal{A}\,P_s$ holds, but $\|\delta\mathcal{A}\|$
admits no bound by the local smallness factor with a gauge-independent constant. The field dependence
of walk terms is
controlled instead by \eqref{4.24:eq-field-dependence-a}, in which no global difference
$\delta\mathcal{A}$ occurs.
\end{remark}

\begin{proof}
\emph{(a).} By the gauge covariance of the quadratic form of Definition~\ref{4.4:def-quadratic-form},
$\mathcal{A}[\mathbb{U}^{u},\operatorname{Ad}_u\mathbb{J}]=\mathcal{U}_u\mathcal{A}[\mathbb{U},\mathbb{J}]
\mathcal{U}_u^{-1}$, that is, in kernels,
\begin{equation*}
\mathcal{A}[\mathbb{U}^{u},\operatorname{Ad}_u\mathbb{J}](b,c)
=\operatorname{Ad}_{u(b_{-})}\,\mathcal{A}[\mathbb{U},\mathbb{J}](b,c)\,\operatorname{Ad}_{u(c_{-})}^{-1};
\end{equation*}
the same holds for the compression $\mathcal{A}_{ss}$ since $P_s\mathcal{U}_u=\mathcal{U}_uP_s$.
Subtracting the $u$-independent operator $\mathcal{A}[1,0]$ gives
\eqref{4.24:eq-noncovariant-difference-2}. On the other hand
\begin{equation*}
\mathcal{U}_u\,\delta\mathcal{A}[\mathbb{U},\mathbb{J}]\,\mathcal{U}_u^{-1}
=\mathcal{U}_u\mathcal{A}[\mathbb{U},\mathbb{J}]\mathcal{U}_u^{-1}
-\mathcal{U}_u\mathcal{A}[1,0]\mathcal{U}_u^{-1},
\end{equation*}
so the two expressions differ by $\mathcal{U}_u\mathcal{A}[1,0]\mathcal{U}_u^{-1}-\mathcal{A}[1,0]$,
which vanishes exactly when $\mathcal{U}_u$ commutes with $\mathcal{A}[1,0]$.

The kernel $\mathcal{A}$ satisfies the covariance condition of \eqref{4.24:eq-walk-hypotheses-a}
and is defined at the trivial pair, so the argument of Lemma~\ref{4.24:lem-colour-scalars}
(trivial-pair factors are colour scalars) applies to $\mathcal{A}$ itself: for all bonds $b,c$,
$\mathcal{A}[1,0](b,c)=\mathfrak{a}(b,c)\,\mathbf{1}$ with $\mathfrak{a}(b,c)\in\mathbb{C}$. Hence
\begin{equation}\label{4.24:eq-noncovariant-difference-4}
\bigl(\mathcal{U}_u\mathcal{A}[1,0]\mathcal{U}_u^{-1}-\mathcal{A}[1,0]\bigr)(b,c)
=\mathfrak{a}(b,c)\bigl(\operatorname{Ad}_{u(b_{-})u(c_{-})^{-1}}-\mathbf{1}\bigr).
\end{equation}
This vanishes for all $(b,c)$ if and only if
$\operatorname{Ad}_{u(b_{-})u(c_{-})^{-1}}=\mathbf{1}$ on $\mathfrak{g}^{\mathbb{C}}$ whenever
$\mathfrak{a}(b,c)\neq0$. An element $z\in G$ with $\operatorname{Ad}_z=\mathbf{1}$ commutes with
$\mathfrak{g}$, hence with its complex span $\mathfrak{g}^{\mathbb{C}}$ and with the identity matrix,
hence with all $N\times N$ matrices, so $z$ is a scalar matrix, that is, central; conversely a
central $z$ has $\operatorname{Ad}_z=\mathbf{1}$. This proves the stated equivalence. The centre of
$G$ consists of finitely many scalar matrices, so the condition fails for a generic non-constant $u$.

\emph{(b).} By \eqref{4.24:eq-noncovariant-difference-2}, the norm of the difference at the transformed
pair is $\|\mathcal{U}_u\mathcal{A}[\mathbb{U},\mathbb{J}]\mathcal{U}_u^{-1}-\mathcal{A}[1,0]\|$, which by
(a) is not in general $\|\delta\mathcal{A}[\mathbb{U},\mathbb{J}]\|$; item (c) exhibits pairs in one
orbit with $\|\delta\mathcal{A}\|=0$ and $\|\delta\mathcal{A}\|\neq0$. Since gauge transformations form a
group acting on pairs, the infimum \eqref{4.24:eq-noncovariant-difference-3} runs over the whole orbit
of $(\mathbb{U},\mathbb{J})$ and is therefore constant on orbits. A bound of
$\eth(\mathbb{U},\mathbb{J})$ by the local smallness factor, with an absolute constant, would require a
single gauge putting every cube of the whole region simultaneously
in the cube form of \eqref{4.24:eq-walk-hypotheses-d}. This fails whenever $\mathbb{U}$ has a
non-central holonomy around a loop along which $\mathcal{A}[1,0]$ propagates (the holonomy example of
this section), and, even when the region has no holes, it degrades with the diameter of the region
(the estimate chaining cube gauges across the region, in this section).

\emph{(c).} Since $\operatorname{Ad}_v0=0$ and $1^{v}$ is the transform of the trivial bond field by
$v$, the pair $(1^{v},0)$ is the transform $(1,0)^{v}$ of the trivial pair. By the gauge invariance of
walk-term densities (Proposition~\ref{4.24:prop-density-invariance}), every density
$\operatorname{tr}_{\mathfrak{g}}\hat t^{\xi}_{\omega}(b',b')$ takes at $(1^{v},0)$ its value at the
trivial pair. Transforming by $v^{-1}$,
\begin{equation*}
(1^{v})^{v^{-1}}(b)=v(b_{-})^{-1}\,v(b_{-})v(b_{+})^{-1}\,v(b_{+})=1,
\end{equation*}
so on every cube the restriction of $v^{-1}$ is a cube gauge bringing the pair to the trivial pair,
and \eqref{4.24:eq-walk-hypotheses-d} holds with $B_i\equiv0$. Nevertheless, by
\eqref{4.24:eq-noncovariant-difference-2} with $(\mathbb{U},\mathbb{J})=(1,0)$ and $u=v$,
\begin{equation*}
\delta\mathcal{A}[1^{v},0]=\mathcal{U}_v\mathcal{A}[1,0]\mathcal{U}_v^{-1}-\mathcal{A}[1,0],
\end{equation*}
whose $(b,c)$ kernel entry is, by \eqref{4.24:eq-noncovariant-difference-4},
$\mathfrak{a}(b,c)(\operatorname{Ad}_{v(b_{-})v(c_{-})^{-1}}-\mathbf{1})\neq0$ for the chosen $(b,c)$,
because $\mathfrak{a}(b,c)\neq0$ and $v(b_{-})v(c_{-})^{-1}$ is non-central. Its norm is up to
$2\|\mathcal{A}[1,0]\|$, and it carries no small factor whatever.

\emph{The consequence.} The resolvent identity
\begin{equation*}
(\mathcal{A}_{ss}+\xi)^{-1}-(\mathcal{A}_{ss}[1,0]+\xi)^{-1}
=-(\mathcal{A}_{ss}+\xi)^{-1}(\delta\mathcal{A})_{ss}(\mathcal{A}_{ss}[1,0]+\xi)^{-1},
\end{equation*}
with $(\delta\mathcal{A})_{ss}=P_s\,\delta\mathcal{A}\,P_s=\mathcal{A}_{ss}-\mathcal{A}_{ss}[1,0]$, is the
second resolvent identity applied to $\mathcal{A}_{ss}+\xi$ and $\mathcal{A}_{ss}[1,0]+\xi$, and is
correct. The inequality $\|(\delta\mathcal{A})_{ss}\|\le\|\delta\mathcal{A}\|$ also holds; however
$\|\delta\mathcal{A}\|$ admits no bound by the local smallness factor with a gauge-independent constant:
by (c), such a bound would assert a gauge-dependent smallness that is false in a general gauge, even at
a pair gauge-equivalent to the trivial one. The field dependence is controlled instead by
\eqref{4.24:eq-field-dependence-a}, in which no global $\delta\mathcal{A}$ occurs; the smallness that
is available is factorwise, namely the differences of \eqref{4.24:eq-walk-hypotheses-d}, each taken on
one cube in one gauge, joined across cubes by the transition defects of the transition
lemma of this section.
\end{proof}

\begin{proposition}[A flat background with non-central holonomy]\label{4.24:prop-flat-holonomy}
Let $\omega$ be a walk, $\xi\ge0$ and $b'$ a bond. Let the factors $F_0,\dots,F_q$ of the walk term $t^{\xi}_{\omega}$ satisfy the structural hypotheses of Definition~\ref{4.24:def-walk-hypotheses}: covariance and locality \eqref{4.24:eq-walk-hypotheses-a}, the majorant hypothesis \eqref{4.24:eq-walk-hypotheses-b}, the cube-form hypothesis \eqref{4.24:eq-walk-hypotheses-c}, the per-cube smallness hypothesis \eqref{4.24:eq-walk-hypotheses-d}, the overlap hypothesis \eqref{4.24:eq-walk-hypotheses-e} and the path hypothesis \eqref{4.24:eq-walk-hypotheses-f}.

Suppose that the $G$-valued configuration $U$ is flat on $\bigcup_i\operatorname{bd}(Y_i)$ in the following sense. For each $i$, $Y_i$ is in cube form with $B_i\equiv0$, that is, $U^{u_i}\equiv1$ on $\operatorname{bd}(Y_i)$, and $A'\equiv0$ and $\mathbb{J}\equiv0$ there. Here $A'$ is the potential of \eqref{4.24:eq-walk-hypotheses-c}. Suppose further that $W_{\omega}=\operatorname{Hol}_U(\Gamma_{\omega})$ is not central in $G$. Then
\begin{equation}\label{4.24:eq-flat-holonomy-1}
\begin{aligned}
\tau_{\omega}-\tau^{0}_{\omega}
&=s^{0}_{\omega}\,\operatorname{tr}_{\mathfrak{g}}\bigl(\operatorname{Ad}_{W_{\omega}}-1\bigr)
=s^{0}_{\omega}\,\bigl(\chi_{\mathrm{adj}}(W_{\omega})-n_{\mathfrak{g}}\bigr)\\
&=s^{0}_{\omega}\,\bigl(|\operatorname{tr}_{\mathbb{C}^N}W_{\omega}|^2-N^2\bigr).
\end{aligned}
\end{equation}
Here $\operatorname{tr}_{\mathbb{C}^N}$ is the trace on $\mathbb{C}^N$ normalised by $\operatorname{tr}_{\mathbb{C}^N}1=N$. The quantity \eqref{4.24:eq-flat-holonomy-1} vanishes if and only if $s^{0}_{\omega}=0$ or $W_{\omega}$ is central.

Consequently, whatever the rank $N\ge2$, and for every family of factors with $s^{0}_{\omega}\ne0$, no inequality of the form
\begin{equation*}
|\tau_{\omega}-\tau^{0}_{\omega}|\le
\bigl(\text{per-cube smallness on the cubes the term reads}\bigr)\times m(\omega)
\end{equation*}
follows, as an implication, from the hypotheses \eqref{4.24:eq-walk-hypotheses-a}--\eqref{4.24:eq-walk-hypotheses-f}. Here the per-cube smallness is a quantity formed from the numbers $a_i$, $\bar\beta$, $\mathfrak{l}_i$ of the configuration at hand and from the left side of \eqref{4.24:eq-walk-hypotheses-d} at its pair, which vanishes when these all vanish.
\end{proposition}

\begin{proof}
We first compute the per-cube numbers in the situation of the hypothesis.

\emph{The cube budgets.} Since $B_i\equiv0$ on $\operatorname{bd}(Y_i)$, with $\eta$ as in \eqref{4.24:eq-walk-hypotheses-c},
\begin{equation*}
a_i=\sup_{b\in\operatorname{bd}(Y_i)}|i\eta B_i(b)|=0
\end{equation*}
for every $i$.

\emph{The factors in cube form.} In the gauge $u_i$ the pair is
\begin{equation*}
\bigl(e^{i\eta\operatorname{Ad}_{u_i}A'}e^{i\eta B_i},\operatorname{Ad}_{u_i}\mathbb{J}\bigr).
\end{equation*}
On $\operatorname{bd}(Y_i)$ we have $A'\equiv0$, $B_i\equiv0$ and $\mathbb{J}\equiv0$, so this pair is the trivial pair $(1,0)$ there. Cube form is preserved by constant gauges $c\in G$, through $B\mapsto\operatorname{Ad}_cB$, $A'\mapsto\operatorname{Ad}_cA'$ and $\mathbb{J}\mapsto\operatorname{Ad}_c\mathbb{J}$, and each of these maps sends $0$ to $0$. By Proposition~\ref{4.24:prop-product-form}(iii), the pair $P_i'$ at which $\tilde F_i'=F_i[P_i']$ is evaluated is the pair in cube form on $Y_i$ transformed by a constant gauge; it therefore also restricts to $(1,0)$ on $\operatorname{bd}(Y_i)$. By the locality clause of \eqref{4.24:eq-walk-hypotheses-a}, $F_i$ depends only on the restriction of the pair to $\operatorname{bd}(Y_i)$. Therefore
\begin{equation*}
\tilde F_i'=F_i[1,0]=G_i
\end{equation*}
exactly. The difference $F_i[P]-G_i$ entering \eqref{4.24:eq-walk-hypotheses-d} therefore vanishes at the pair of the configuration, and $\varepsilon$ may be taken $0$ for these factors.

\emph{Overlaps and paths.} In the definition of $\beta_i$ in \eqref{4.24:eq-walk-hypotheses-e}, every path sum is a sum of terms $|i\eta B_i(b)|+|i\eta B_{i+1}(b)|$ over bonds $b$ of $\operatorname{bd}(Y_i)\cap\operatorname{bd}(Y_{i+1})$, and each such term is $0$. Hence $\beta_i=0$ for all $i$ and $\bar\beta=0$. Likewise, in \eqref{4.24:eq-walk-hypotheses-f},
\begin{equation*}
\mathfrak{l}_i=\sum_{b\in\gamma_i}|i\eta B_i(b)|=0
\qquad\text{for } i=0,\dots,q.
\end{equation*}

\emph{Transition defects.} By the lemma on the smallness of transition defects on overlaps (Lemma~\ref{4.24:lem-transition-defects}), for every $i$ and every $s\in O_i$,
\begin{equation*}
\operatorname{dis}k_i(s)\le\beta_i\le\bar\beta=0,
\qquad
\operatorname{dis}k_i'(s)=\operatorname{dis}k_i(s),
\end{equation*}
so $k_i\equiv1$. The same lemma gives
\begin{equation*}
\|\operatorname{Ad}_{k_i'(s)}-1\|_{e}\le6\bar\beta=0,
\end{equation*}
so $\operatorname{Ad}_{k_i'(s)}=1$ on every overlap.

\emph{The residual element.} By Lemma~\ref{4.24:lem-residual-loop},
\begin{equation*}
\mathfrak{h}_{\omega}=u_0(s_0^{*})\,\bigl[W_{\omega}E_{\omega}\bigr]\,u_0(s_0^{*})^{-1}
\end{equation*}
with $\operatorname{dis}E_{\omega}\le\sum_{i=0}^{q}\mathfrak{l}_i=0$. Hence $E_{\omega}=1$, and $\mathfrak{h}_{\omega}=u_0(s_0^{*})\,W_{\omega}\,u_0(s_0^{*})^{-1}$ is conjugate to $W_{\omega}$.

\emph{The chain.} In the trace form of Proposition~\ref{4.24:prop-product-form}(ii), insert $\tilde F_i'=G_i$ and $\operatorname{Ad}_{k_i'}=1$. This gives
\begin{equation*}
S_{\omega}=\sum_{b_1,\dots,b_q}G_0(b',b_1)G_1(b_1,b_2)\cdots G_q(b_q,b')
=(G_0G_1\cdots G_q)(b',b')=S^{0}_{\omega}.
\end{equation*}
By the lemma that trivial-pair factors are colour scalars (Lemma~\ref{4.24:lem-colour-scalars}), $S^{0}_{\omega}=s^{0}_{\omega}\cdot1$.

\emph{The exact split.} All hypotheses of Proposition~\ref{4.24:prop-density-difference} are among those assumed, so its exact split \eqref{4.24:eq-density-difference-a} applies. Since $S_{\omega}-S^{0}_{\omega}=0$, its first member vanishes, and it reads
\begin{equation*}
\tau_{\omega}-\tau^{0}_{\omega}=s^{0}_{\omega}\,\operatorname{tr}_{\mathfrak{g}}\bigl(\operatorname{Ad}_{\mathfrak{h}_{\omega}}-1\bigr).
\end{equation*}
With $c:=u_0(s_0^{*})$ we have $\operatorname{Ad}_{\mathfrak{h}_{\omega}}=\operatorname{Ad}_c\operatorname{Ad}_{W_{\omega}}\operatorname{Ad}_c^{-1}$. The trace on $\operatorname{End}(\mathfrak{g}^{\mathbb{C}})$ is invariant under conjugation, so $\operatorname{tr}_{\mathfrak{g}}(\operatorname{Ad}_{\mathfrak{h}_{\omega}}-1)=\operatorname{tr}_{\mathfrak{g}}(\operatorname{Ad}_{W_{\omega}}-1)$. This is the first equality in \eqref{4.24:eq-flat-holonomy-1}.

For the second, $\operatorname{tr}_{\mathfrak{g}}\operatorname{Ad}_{W_{\omega}}=\chi_{\mathrm{adj}}(W_{\omega})$ by definition of the adjoint character, and $\operatorname{tr}_{\mathfrak{g}}1=\dim\mathfrak{g}=n_{\mathfrak{g}}$.

For the third, Lemma~\ref{4.24:lem-rotation-trace} (the trace of an adjoint rotation) gives $\chi_{\mathrm{adj}}(W_{\omega})=|\operatorname{tr}_{\mathbb{C}^N}W_{\omega}|^2-1$. Together with $n_{\mathfrak{g}}=N^2-1$ this yields
\begin{equation*}
\chi_{\mathrm{adj}}(W_{\omega})-n_{\mathfrak{g}}=|\operatorname{tr}_{\mathbb{C}^N}W_{\omega}|^2-N^2.
\end{equation*}

\emph{Vanishing.} The product \eqref{4.24:eq-flat-holonomy-1} vanishes if and only if $s^{0}_{\omega}=0$ or $|\operatorname{tr}_{\mathbb{C}^N}W_{\omega}|=N$. Write the eigenvalues of the unitary matrix $W_{\omega}$ as $e^{i\theta_1},\dots,e^{i\theta_N}$. Then
\begin{equation*}
|\operatorname{tr}_{\mathbb{C}^N}W_{\omega}|=\Bigl|\sum_{j=1}^{N}e^{i\theta_j}\Bigr|\le N,
\end{equation*}
with equality if and only if all $e^{i\theta_j}$ are equal. For a unitary, hence diagonalisable, matrix this means that $W_{\omega}$ is a scalar matrix. The scalar matrices in $G=SU(N)$ form its centre. Hence the quantity vanishes if and only if $s^{0}_{\omega}=0$ or $W_{\omega}$ is central.

\emph{The consequence.} In the situation of the hypothesis, all of \eqref{4.24:eq-walk-hypotheses-a}--\eqref{4.24:eq-walk-hypotheses-f} hold; the numbers of this configuration are $a_i=0$, $\bar\beta=0$ and $\mathfrak{l}_i=0$, and the left side of \eqref{4.24:eq-walk-hypotheses-d} vanishes at the pair of this configuration, so that $\varepsilon$ may be taken $0$ for these factors. Any per-cube smallness number formed from these data is therefore $0$, and the right-hand side of an inequality of the stated form is $0$. On the other hand, if $s^{0}_{\omega}\ne0$ and $W_{\omega}$ is not central, the left-hand side is nonzero by the vanishing criterion just proved. This holds for every $N\ge2$ and every family of factors with $s^{0}_{\omega}\ne0$. So no such inequality is implied by \eqref{4.24:eq-walk-hypotheses-a}--\eqref{4.24:eq-walk-hypotheses-f}.
\end{proof}

\begin{remark}
Configurations $U$ as in Proposition~\ref{4.24:prop-flat-holonomy} exist as soon as $\Gamma_{\omega}$ is not null-homotopic through plaquettes on which $U$ is constrained. Two examples are:
\begin{itemize}
\item $D(\omega)$ is an annulus of domains around a region that the term does not read;
\item $D(\omega)$ is an annulus of domains around a cycle of the torus.
\end{itemize}
In either case one takes $U$ to be a pure gauge on a simply connected cover and lets the deck transformation act by a non-central element of $G$.

Whether the actual factors of the walk expansion have $s^{0}_{\omega}\ne0$ for the walks in question is immaterial to the logical point of the proposition. For the class of walks with $|\omega|=0$, $s^{0}_{\omega_0}$ is the diagonal entry of the trivial resolvent and is nonzero.

What can restore a termwise bound is information about $U$ off the cubes the term reads, namely the flux through $\Gamma_{\omega}$. This is the holonomy hypothesis on threading loops, \eqref{4.24:eq-holonomy-hypothesis-a} of Definition~\ref{4.24:def-holonomy-hypothesis}.

Physically, the proposition expresses the dependence of the diagonal resolvent on holonomies. The sum
\begin{equation*}
\sum_{\omega}\tau_{\omega}=\operatorname{tr}_{\mathfrak{g}}(\mathcal{A}+\xi)^{-1}(b',b')
\end{equation*}
depends on the Wilson loops of $U$ around every cycle that the walks can thread, weighted by $m(\omega)$. Near a hole of the region, this dependence is not proportional to the local curvature entering clause~(ii) of Theorem~\ref{3.1:thm-main}. Examples of such holes are an excluded large-field component, a finer-scale core, and a cycle of the torus.
\end{remark}

\begin{remark}[A single gauge degrades with the diameter]\label{4.24:rem-single-gauge}
An alternative to the product form of Proposition~\ref{4.24:prop-product-form} would be to
bring $U$ into one gauge on the whole domain $D(\omega)$ read by a walk term. Write $\eta$ for the
scale factor with which a cube potential enters the cube form, $U=e^{i\eta B}$; here $B$ denotes
a single potential on the whole union considered, as $B_i$ denotes the potential on the cube
$Y_i$. Consider a contractible union of $\mathsf{q}$ cubes of one scale, each carrying the budget
$a$. On it there is a single gauge in which
\begin{equation*}
U=e^{i\eta B}\quad\text{on every bond of the union,}\qquad \eta|B|\lesssim \mathsf{q}\,a .
\end{equation*}
Such a gauge is obtained axially from a base point. The potential at a bond then integrates the
curvature along a path of length at most a constant multiple of $\mathsf{q}$ cube sides, and this
gives the displayed bound.

The smallness of this gauge therefore degrades linearly in $\mathsf{q}$, where
\begin{equation*}
\mathsf{q}\le|D(\omega)| .
\end{equation*}
The resulting constant depends on this length, that is, on the number $\mathsf{q}$ of cubes in the
union. Such a constant can be absorbed by the decay of the walk only in the regime
\begin{equation*}
\mathsf{q}\,a\ll 1 ,
\end{equation*}
and outside that regime one has only the trivial bound by $2$, which there is itself at most a
constant multiple of $\mathsf{q}\,a$.

Moreover, such a single gauge does not exist at all in two situations. The first is when
$D(\omega)$ is not contractible. The second is when the union surrounds holes across which the
holonomy of $U$ is not central. An instance of the second situation is the flat background with
non-central holonomy of Proposition~\ref{4.24:prop-flat-holonomy}: there every cube potential
vanishes, and yet no gauge on $D(\omega)$ brings $U$ to the identity.

The product form of Proposition~\ref{4.24:prop-product-form} isolates exactly this obstruction in
the residual element $\mathfrak{h}_{\omega}$. It imposes no further condition on $U$.
\end{remark}

\begin{lemma}[Loop holonomy bounded by plaquette holonomies]\label{4.24:lem-loop-holonomy}
Let $U$ be a configuration of bond variables with values in $G$, and for a lattice path $\gamma$
let $U(\gamma)$ be the parallel transport along $\gamma$ of
Definition~\ref{4.24:not-gauge-transport}. For a closed lattice loop $\Gamma$ based at a site
$x$ write $\operatorname{Hol}_U(\Gamma)$ for $U(\Gamma)$, the transport along $\Gamma$ from $x$
back to $x$. Let $\Gamma$ be a closed lattice loop, and suppose that $\Gamma$ is carried to a
constant loop by a finite sequence of moves, each of which is either
\begin{itemize}
\item[(m1)] the insertion or deletion of a backtrack $b b^{-1}$, or
\item[(m2)] the replacement of a sub-path $\pi$ by a path $\pi'$ with the same endpoints such that
$\pi^{-1}\pi'$ (or $\pi'\pi^{-1}$) is the boundary $\partial P$ of one plaquette $P$ traversed once.
\end{itemize}
Then
\begin{equation}\label{4.24:eq-loop-holonomy-a}
\operatorname{dis}\operatorname{Hol}_U(\Gamma)\;\le\;\sum_{\text{(m2)-moves}}\operatorname{dis}U(\partial P),
\end{equation}
where $U(\partial P)$ is the plaquette holonomy based at either endpoint of the replaced sub-path.
The right-hand side does not depend on this choice of base point, by property (d2) of
Lemma~\ref{4.24:lem-group-distance}.
\end{lemma}

\begin{proof}
We first justify the last sentence of the statement. Two plaquette holonomies of $P$ based at two
corners of $\partial P$ are conjugate by the transport along the part of $\partial P$ joining
those corners. This transport is an element of $G$, because $U$ is $G$-valued. By (d2) in
\eqref{4.24:eq-group-distance-a}, $\operatorname{dis}$ is invariant under conjugation by elements
of $G$ and satisfies $\operatorname{dis}(h^{-1})=\operatorname{dis}h$. Hence
$\operatorname{dis}U(\partial P)$ depends neither on the corner at which the holonomy is based nor
on the orientation in which $\partial P$ is traversed.

Write $\Gamma=\Gamma_0,\Gamma_1,\dots,\Gamma_n$ for the loops of the given sequence, so that
$\Gamma_{j+1}$ arises from $\Gamma_j$ by one move and $\Gamma_n$ is the constant loop. Every move
replaces a contiguous segment of the bond sequence by a segment with the same endpoints. The base
point $x$ is therefore the same for all $\Gamma_j$. We run the sequence backwards from $\Gamma_n$
and prove, by downward induction on $j$, that
\begin{equation}\label{4.24:eq-loop-holonomy-1}
\operatorname{dis}\operatorname{Hol}_U(\Gamma_j)\;\le\;
\sum_{\substack{j<i\le n\\ \text{the $i$-th move is of type (m2)}}}\operatorname{dis}U(\partial P_i),
\end{equation}
where $P_i$ is the plaquette of the $i$-th move. For $j=n$ we have
$\operatorname{Hol}_U(\Gamma_n)=1$, and $\operatorname{dis}1=0$ by (d1). The case $j=0$ is
\eqref{4.24:eq-loop-holonomy-a}. It therefore suffices to show that one move changes
$\operatorname{dis}\operatorname{Hol}_U$ by at most $0$ in case (m1), and by at most
$\operatorname{dis}U(\partial P)$ in case (m2). Both kinds of move are symmetric: the reverse of
an (m1)-move is an (m1)-move, and the reverse of an (m2)-move with plaquette $P$ is an (m2)-move
with the same plaquette, since $\pi'^{-1}\pi=(\pi^{-1}\pi')^{-1}$ and
$\pi\pi'^{-1}=(\pi'\pi^{-1})^{-1}$ are again $\partial P$ traversed once. It is thus enough to
compare a loop $\Gamma$ with the loop $\Gamma'$ obtained from it by one move.

\emph{A move of type (m1)} does not change the holonomy. The backtrack $bb^{-1}$ contributes the
factor $U(b)U(b)^{-1}=1$ to the ordered product.

\emph{A move of type (m2).} Write $\Gamma=\alpha\pi\beta$, with $\alpha$ running from $x$ to the
initial point of $\pi$ and $\beta$ from the endpoint $z$ of $\pi$ back to $x$. Either of these
paths may be empty, in which case its transport is $1$. Put $A=U(\alpha)$ and $B=U(\beta)$, both
in $G$. Transport is multiplicative under concatenation and inverted under reversal of a path. Hence
\begin{equation*}
\operatorname{Hol}_U(\Gamma)=A\,U(\pi)\,B,\qquad
\operatorname{Hol}_U(\Gamma')=A\,U(\pi')\,B=A\,U(\pi)\,\bigl[U(\pi)^{-1}U(\pi')\bigr]\,B .
\end{equation*}
Suppose first that $\pi^{-1}\pi'=\partial P$. Then $U(\pi)^{-1}U(\pi')=U(\pi^{-1}\pi')$ is the
transport around the closed loop $\pi^{-1}\pi'$ around one plaquette, based at $z$. That is,
$U(\pi)^{-1}U(\pi')=U(\partial P)^{\pm1}$ with $U(\partial P)$ based at $z$. Consequently
\begin{equation}\label{4.24:eq-loop-holonomy-2}
\operatorname{Hol}_U(\Gamma')=\operatorname{Hol}_U(\Gamma)\cdot B^{-1}U(\partial P)^{\pm1}B .
\end{equation}
By the subadditivity in (d2),
\begin{align*}
\operatorname{dis}\operatorname{Hol}_U(\Gamma')
&\le\operatorname{dis}\operatorname{Hol}_U(\Gamma)
+\operatorname{dis}\bigl(B^{-1}U(\partial P)^{\pm1}B\bigr)\\
&=\operatorname{dis}\operatorname{Hol}_U(\Gamma)+\operatorname{dis}U(\partial P).
\end{align*}
The equality uses the two remaining parts of (d2): conjugation by $B\in G$ leaves
$\operatorname{dis}$ unchanged, and $\operatorname{dis}$ of an inverse equals $\operatorname{dis}$
of the element.

Suppose next that $\pi'\pi^{-1}=\partial P$. Then $U(\pi')U(\pi)^{-1}=U(\partial P)^{\pm1}$, based
at the initial point of $\pi$. This gives $U(\pi')=U(\partial P)^{\pm1}U(\pi)$, and hence
\begin{equation*}
\operatorname{Hol}_U(\Gamma')=A\,U(\partial P)^{\pm1}A^{-1}\cdot\operatorname{Hol}_U(\Gamma).
\end{equation*}
Since $A\in G$, the same three parts of (d2) give the same bound.

Combining the two cases with the case of (m1) proves the induction step, and with it
\eqref{4.24:eq-loop-holonomy-1} for every $j$. For $j=0$ this is
\eqref{4.24:eq-loop-holonomy-a}. No commutativity of the plaquette holonomies is used. In
\eqref{4.24:eq-loop-holonomy-2} they enter conjugated and in the order of the moves, and
$\operatorname{dis}$ discards both the conjugation and the order. Nothing more than this is
claimed.
\end{proof}

\begin{lemma}[The holonomy of one plaquette in a cube gauge]\label{4.24:lem-plaquette-holonomy}
Let $P$ be a plaquette, let $X_1,X_2,X_3,X_4\in\mathfrak{g}^{\mathbb{C}}$, and put
\begin{equation*}
\mathfrak{s}_P:=\sum_{m=1}^{4}|X_m|,\qquad
\mathfrak{d}_P:=|X_1+X_2-X_3-X_4|,\qquad
\mathfrak{b}_P:=\mathfrak{d}_P+\tfrac12\,\mathfrak{s}_P^{2}\,e^{\mathfrak{s}_P}.
\end{equation*}
If $\mathfrak{b}_P\le\frac12$, then
\begin{equation}\label{4.24:eq-plaquette-holonomy-a}
\operatorname{dis}\bigl(e^{X_1}e^{X_2}e^{-X_3}e^{-X_4}\bigr)\le 2\,\mathfrak{b}_P .
\end{equation}
The lemma is applied with the four bonds $b_1,b_2,b_3,b_4$ of $P$ taken in order around $P$,
in a cube gauge as in the structural hypotheses \eqref{4.24:eq-walk-hypotheses-c}:
there $X_m=i\eta B(b_m)\in i\mathfrak{g}$, where $\eta$ is the lattice spacing and $B$ the cube
potential, and, when the lattice curl $\operatorname{curl}B$ is formed with plain forward
differences, $\mathfrak{d}_P=|i\eta^{2}(\operatorname{curl}B)(P)|$.
\end{lemma}

\begin{proof}
Put $h:=e^{X_1}e^{X_2}e^{-X_3}e^{-X_4}$. Each factor has determinant $e^{\pm N\operatorname{tr}X_m}=1$,
so $h\in G^{\mathbb{C}}$ and $\operatorname{dis}h$ is defined as in
Lemma~\ref{4.24:lem-group-distance}.

Expanding each exponential in its power series, $h$ is the absolutely convergent sum, over all
$n_1,n_2,n_3,n_4\ge0$, of the monomials
\begin{equation*}
\frac{X_1^{n_1}X_2^{n_2}(-X_3)^{n_3}(-X_4)^{n_4}}{n_1!\,n_2!\,n_3!\,n_4!}.
\end{equation*}
The monomial of total degree $0$ is $\mathbf{1}$, and those of total degree $1$ add up to
$X_1+X_2-X_3-X_4$. Hence
\begin{equation*}
h-\mathbf{1}=(X_1+X_2-X_3-X_4)+R,
\end{equation*}
where $R$ is the sum of the monomials of total degree at least $2$. For a product of $k\ge1$ of
the $X_m$, the property \eqref{4.24:eq-colour-norms-a} of the norm $|\cdot|$ (applied with the
identity on the left, $X_{m_1}$ in the middle and $X_{m_2}\cdots X_{m_k}$ on the right), the
sub-multiplicativity of $\|\cdot\|$ and the inequality $\|\cdot\|\le|\cdot|$ give
\begin{equation*}
|X_{m_1}X_{m_2}\cdots X_{m_k}|
\le|X_{m_1}|\prod_{i=2}^{k}\|X_{m_i}\|
\le\prod_{i=1}^{k}|X_{m_i}|.
\end{equation*}
Summing these bounds over the monomials of total degree at least $2$, and noting that the
resulting sum consists of the terms of degree at least $2$ of the expansion of
$\prod_m e^{|X_m|}=e^{\mathfrak{s}_P}$, we obtain
\begin{equation*}
|R|\le\prod_{m=1}^{4}e^{|X_m|}-1-\mathfrak{s}_P
=e^{\mathfrak{s}_P}-1-\mathfrak{s}_P
\le\tfrac12\,\mathfrak{s}_P^{2}\,e^{\mathfrak{s}_P};
\end{equation*}
the last inequality holds because
\begin{equation*}
e^{\mathfrak{s}_P}-1-\mathfrak{s}_P=\sum_{k\ge2}\frac{\mathfrak{s}_P^{k}}{k!}
=\frac12\,\mathfrak{s}_P^{2}\sum_{j\ge0}\frac{2\,\mathfrak{s}_P^{j}}{(j+2)!}
\end{equation*}
and $2/(j+2)!\le1/j!$ for every $j\ge0$.

Put $\mathcal{E}_P:=h-\mathbf{1}$. By the two previous displays,
\begin{equation*}
|\mathcal{E}_P|\le\mathfrak{d}_P+|R|\le\mathfrak{b}_P\le\tfrac12<1 .
\end{equation*}
For $n\ge1$, \eqref{4.24:eq-colour-norms-a} (with $\mathcal{E}_P^{\,n-1}$ on the left and
$\mathcal{E}_P$ in the middle), sub-multiplicativity and $\|\cdot\|\le|\cdot|$ give
\begin{equation*}
|\mathcal{E}_P^{\,n}|\le\|\mathcal{E}_P\|^{n-1}|\mathcal{E}_P|\le|\mathcal{E}_P|^{n}.
\end{equation*}
Therefore the series
\begin{equation*}
\log h:=\sum_{n\ge1}\frac{(-1)^{n+1}}{n}\,\mathcal{E}_P^{\,n}
\end{equation*}
converges absolutely in $|\cdot|$, and, since $\|\mathcal{E}_P\|\le|\mathcal{E}_P|<1$, the power
series identity $\exp(\log(\mathbf{1}+Y))=\mathbf{1}+Y$, valid for every $Y\in M_N(\mathbb{C})$
with $\|Y\|<1$, gives
$e^{\log h}=h$. Moreover $\log h\in\mathfrak{sl}(N,\mathbb{C})=\mathfrak{g}^{\mathbb{C}}$. To see
this, let $t\in[0,1]$ and put $h_t:=e^{tX_1}e^{tX_2}e^{-tX_3}e^{-tX_4}$, so that $h_1=h$ and
$h_0=\mathbf{1}$. Formed with $tX_1,\dots,tX_4$ in place of $X_1,\dots,X_4$, the three quantities
of the statement become $t\,\mathfrak{s}_P$, $t\,\mathfrak{d}_P$ and
\begin{equation*}
t\,\mathfrak{d}_P+\tfrac12\,t^{2}\mathfrak{s}_P^{2}\,e^{t\mathfrak{s}_P}
\le\mathfrak{b}_P\le\tfrac12,
\end{equation*}
because $t^{2}\le t\le1$ and $e^{t\mathfrak{s}_P}\le e^{\mathfrak{s}_P}$. Hence everything proved
so far applies to $tX_1,\dots,tX_4$: for every $t\in[0,1]$ one has $\det h_t=1$ and
$|h_t-\mathbf{1}|\le\frac12$, the series $\log h_t$ (the series above with $h_t-\mathbf{1}$ in
place of $\mathcal{E}_P$) converges, $e^{\log h_t}=h_t$, and
\begin{equation*}
e^{N\operatorname{tr}\log h_t}=\det e^{\log h_t}=\det h_t=1,
\end{equation*}
that is, $N\operatorname{tr}\log h_t\in2\pi i\mathbb{Z}$. The map $t\mapsto h_t-\mathbf{1}$ is
continuous, and on $\{Y\in M_N(\mathbb{C}):|Y|\le\frac12\}$ the $n$-th term of the logarithmic
series is bounded in $|\cdot|$ by $2^{-n}/n$, so the series converges uniformly there; hence
$t\mapsto N\operatorname{tr}\log h_t$ is continuous on the connected interval $[0,1]$. It takes
values in $2\pi i\mathbb{Z}$ and vanishes at $t=0$, where $h_0-\mathbf{1}=0$; therefore it vanishes
at $t=1$, that is, $\operatorname{tr}\log h=0$.

Finally, using $|\mathcal{E}_P|\le\frac12$,
\begin{equation*}
|\log h|\le\sum_{n\ge1}\frac{|\mathcal{E}_P|^{n}}{n}
\le\frac{|\mathcal{E}_P|}{1-|\mathcal{E}_P|}
\le2|\mathcal{E}_P|\le2\,\mathfrak{b}_P ,
\end{equation*}
and $\operatorname{dis}h\le|\log h|$ by clause (d2) of \eqref{4.24:eq-group-distance-a}, since
$\log h\in\mathfrak{g}^{\mathbb{C}}$ is a logarithm of $h$. This is
\eqref{4.24:eq-plaquette-holonomy-a}.
\end{proof}

\begin{remark}
Let $P$ lie in a cube whose unit is $\ell_P$, and let the pair $(\mathbb{U},\mathbb{J})$ be in
cube form on that cube, as in
\eqref{4.24:eq-walk-hypotheses-c}. The bounds \cite[(3.35)]{B-CMP99b}, read in the cube's unit,
give
\begin{equation*}
\eta|B|\le O(1)\,M\tilde\alpha_0\,\frac{\eta}{\ell_P},\qquad
\eta^{2}|\nabla B|\le O(1)\,M\tilde\alpha_0\,\frac{\eta^{2}}{\ell_P^{2}},
\end{equation*}
so that $\mathfrak{b}_P\le c\,M\tilde\alpha_0\,\eta^{2}/\ell_P^{2}$ with an absolute constant
$c$. Thus, by Lemma~\ref{4.24:lem-plaquette-holonomy}, provided $\mathfrak{b}_P\le\frac12$, the
holonomy $h_P:=e^{X_1}e^{X_2}e^{-X_3}e^{-X_4}$ of one plaquette $P$ of a cube of unit $\ell_P$
satisfies
\begin{equation*}
\operatorname{dis}(h_P)\le c\,M\tilde\alpha_0\,(\eta/\ell_P)^{2}.
\end{equation*}
The covariant gradient $\nabla_U$ appearing in \cite[(3.35)]{B-CMP99b} differs from the plain
forward difference in the cube gauge by terms bounded by $(e^{2\eta|B|}-1)|B|$; these terms are
contained in the member $\frac12\mathfrak{s}_P^{2}e^{\mathfrak{s}_P}$ of $\mathfrak{b}_P$.
\end{remark}

\begin{corollary}[Reduction of the holonomy hypothesis to a flux bound]\label{4.24:cor-flux-reduction}
Let $\omega$ be a walk and $\Gamma_{\omega}$ its threading loop, as in the holonomy hypothesis on
threading loops (Definition~\ref{4.24:def-holonomy-hypothesis}). Suppose that $\Gamma_{\omega}$
admits a contracting sequence, that is, a finite sequence of moves of the kinds (m1) and (m2) of
Lemma~\ref{4.24:lem-loop-holonomy} carrying $\Gamma_{\omega}$ to a constant loop. Suppose further
that every plaquette $P$ of an (m2)-move of this sequence lies in a cube, of unit $\ell_P$, on which
there is a gauge transformation $u$ of the cube such that the transformed configuration $U^{u}$
satisfies $U^{u}(b)=e^{i\eta B(b)}$ for every bond $b$ of the cube, with $\eta$ and the cube
potential $B$ as in Lemma~\ref{4.24:lem-plaquette-holonomy}, and
\begin{equation}\label{4.24:eq-flux-reduction-2}
\eta|B|\le O(1)\,M\tilde\alpha_0\,\frac{\eta}{\ell_P},
\qquad
\eta^{2}|\nabla_U B|\le O(1)\,M\tilde\alpha_0\,\frac{\eta^{2}}{\ell_P^{2}},
\end{equation}
where $\nabla_U$ is the covariant difference of \cite[(3.35)]{B-CMP99b} and $O(1)$ stands for a
fixed absolute constant. The gauge may be chosen separately for each
plaquette, since $\operatorname{dis} U(\partial P)$ depends neither on the base point nor on the
gauge. Suppose finally that $\eta\le\ell_P$ for every such plaquette $P$ and that
$M\tilde\alpha_0\le\varepsilon_0$, where $\varepsilon_0\in(0,1]$ is the absolute constant fixed in
the proof below. Then
\begin{equation}\label{4.24:eq-flux-reduction-a}
\operatorname{dis}\operatorname{Hol}_U(\Gamma_{\omega})
\le c\,M\tilde\alpha_0\sum_{(\mathrm{m2})\text{-moves}}\Bigl(\frac{\eta}{\ell_P}\Bigr)^{2}
=:(\mathrm{FLX}_{\omega}),
\end{equation}
with an absolute constant $c$. We call $(\mathrm{FLX}_{\omega})$, the right-hand side of
\eqref{4.24:eq-flux-reduction-a}, the flux bound of $\Gamma_{\omega}$; it is the
\emph{flux} of the curvature budget \eqref{4.24:eq-flux-reduction-2} through a spanning disc of
$\Gamma_{\omega}$, counted with multiplicity. A plaquette used in several (m2)-moves contributes
once for each of them.
\end{corollary}

\begin{proof}
Fix a contracting sequence as in the hypothesis. By the bound of loop holonomy by plaquette
holonomies (Lemma~\ref{4.24:lem-loop-holonomy}), applied to the closed loop
$\Gamma=\Gamma_{\omega}$ and to this sequence of moves,
\begin{equation}\label{4.24:eq-flux-reduction-1}
\operatorname{dis}\operatorname{Hol}_U(\Gamma_{\omega})
\le\sum_{(\mathrm{m2})\text{-moves}}\operatorname{dis}U(\partial P),
\end{equation}
with one term for each (m2)-move. Fix one such move and its plaquette $P$. Let $\ell_P$ be the unit of
the cube containing $P$ on which the bounds \eqref{4.24:eq-flux-reduction-2} hold in some gauge.
Because
$\operatorname{dis}U(\partial P)$ is independent of the base point and of the gauge, we may compute it
in that gauge. There the variables $U^{u}(b_1),\dots,U^{u}(b_4)$ on the four bonds of $P$, taken in
order, are $e^{X_m}$ with $X_m=i\eta B(b_m)\in i\mathfrak{g}$. Hence $U^{u}(\partial P)$ equals
\begin{equation*}
e^{X_1}e^{X_2}e^{-X_3}e^{-X_4},
\end{equation*}
and $\operatorname{dis}U(\partial P)=\operatorname{dis}U^{u}(\partial P)$.
This is the situation of the holonomy of one plaquette in a cube gauge
(Lemma~\ref{4.24:lem-plaquette-holonomy}), with $\mathfrak{s}_P$, $\mathfrak{d}_P$ and
$\mathfrak{b}_P$ as defined there.

The bounds \eqref{4.24:eq-flux-reduction-2} are taken in the cube's unit $\ell_P$.
The first bound controls $\mathfrak{s}_P=\sum_m|X_m|$, and with it the member
$\tfrac12 \mathfrak{s}_P^2e^{\mathfrak{s}_P}$ of $\mathfrak{b}_P$.
The second bound controls $\mathfrak{d}_P$, which is $\eta^{2}$ times the plain forward-difference
curl, $\mathfrak{d}_P=|i\eta^{2}\operatorname{curl}B(P)|$. The covariant difference $\nabla_U$ of
\cite[(3.35)]{B-CMP99b} in the second bound differs from this plain difference in the cube gauge by
terms of size at most $(e^{2\eta|B|}-1)|B|$; these terms are bounded by an absolute multiple of
$\mathfrak{s}_P^{2}e^{\mathfrak{s}_P}$, and so are absorbed into the second member of
$\mathfrak{b}_P$ at the cost of the constant.
Thus, with an absolute constant $C$,
\begin{equation*}
\mathfrak{d}_P\le\eta^{2}|\nabla_U B|+C\,\mathfrak{s}_P^{2}e^{\mathfrak{s}_P}
\le O(1)\,M\tilde\alpha_0\Bigl(\frac{\eta}{\ell_P}\Bigr)^{2}+C\,\mathfrak{s}_P^{2}e^{\mathfrak{s}_P}.
\end{equation*}
Since $\mathfrak{s}_P$ is a sum
of four terms $\eta|B(b_m)|$, the first bound gives
$\mathfrak{s}_P\le 4\,O(1)\,M\tilde\alpha_0\,\eta/\ell_P$.
By the hypotheses $\eta\le\ell_P$ and $M\tilde\alpha_0\le\varepsilon_0\le 1$ this gives
$\mathfrak{s}_P\le 4\,O(1)$ and $(M\tilde\alpha_0)^{2}\le M\tilde\alpha_0$, and
hence
\begin{equation*}
\begin{split}
\tfrac12 \mathfrak{s}_P^{2}e^{\mathfrak{s}_P}
&\le 8\,O(1)^{2}\,(M\tilde\alpha_0)^{2}\,e^{\mathfrak{s}_P}\Bigl(\frac{\eta}{\ell_P}\Bigr)^{2}\\
&\le 8\,O(1)^{2}\,e^{4\,O(1)}\,M\tilde\alpha_0\Bigl(\frac{\eta}{\ell_P}\Bigr)^{2},
\end{split}
\end{equation*}
and in the same way
$C\,\mathfrak{s}_P^{2}e^{\mathfrak{s}_P}\le 16\,C\,O(1)^{2}e^{4\,O(1)}M\tilde\alpha_0(\eta/\ell_P)^{2}$.
Together these
estimates give
\begin{equation*}
\mathfrak{b}_P\le c'\,M\tilde\alpha_0\Bigl(\frac{\eta}{\ell_P}\Bigr)^{2},
\qquad
c':=O(1)+(16C+8)\,O(1)^{2}\,e^{4\,O(1)},
\end{equation*}
and $c'$ is an absolute constant. We fix $\varepsilon_0:=\min\{1,1/(2c')\}$. Since
$\eta\le\ell_P$ and $M\tilde\alpha_0\le\varepsilon_0$, we get
$\mathfrak{b}_P\le c'\varepsilon_0\le\tfrac12$, so
the hypothesis $\mathfrak{b}_P\le\tfrac12$ of Lemma~\ref{4.24:lem-plaquette-holonomy} holds, and
that lemma yields
\begin{equation*}
\operatorname{dis}U(\partial P)\le 2\mathfrak{b}_P
\le 2c'\,M\tilde\alpha_0\Bigl(\frac{\eta}{\ell_P}\Bigr)^{2}.
\end{equation*}
Inserting this bound into each term of \eqref{4.24:eq-flux-reduction-1} gives
\eqref{4.24:eq-flux-reduction-a} with $c:=2c'$.
\end{proof}

\begin{remark}[When the flux bound has the required form]\label{4.24:rem-flux-form}
Let $\omega$ be a walk of the expansion, with factor count $q=q(\omega)$ and threading loop
$\Gamma_{\omega}$, and let $\eta$ and $c$ be as in the reduction of the holonomy hypothesis to a
flux bound (Corollary~\ref{4.24:cor-flux-reduction}). In what follows, the \emph{per-cube
condition} is the per-cube smallness condition on the curvature that
Corollary~\ref{4.24:cor-flux-reduction} requires of the cubes containing the plaquettes of a
contracting sequence, each cube being taken in a gauge of its own. Suppose that all
the domains of $\omega$, together with a disc spanning $\Gamma_{\omega}$, lie at one level, with
unit $\ell$, inside one region in which the per-cube condition holds.

The threading loop $\Gamma_{\omega}$ then consists of at most $c'(q+1)\,M\ell/\eta$ bonds, where
$c'$ is an absolute constant. We contract $\Gamma_{\omega}$ inside the spanning disc by coning to
a point. The number of plaquette moves needed is at most the square of the length of the loop, up
to an absolute factor: there is an absolute constant $c''$ such that at most
\begin{equation}\label{4.24:eq-flux-form-2}
c''\,(q+1)^2\,\Bigl(\frac{M\ell}{\eta}\Bigr)^2
\end{equation}
plaquette moves are needed. Every plaquette $P$ of this contracting sequence lies in a cube of the
given region, so the per-cube condition holds on that cube, at the common level, and
$\ell_P=\ell$. Since $\operatorname{dis}U(\partial P)$ does not depend on the base point or on the
gauge, the gauge may be chosen plaquette by plaquette, and the hypothesis of
Corollary~\ref{4.24:cor-flux-reduction} is satisfied. In the flux bound
\eqref{4.24:eq-flux-reduction-a} that the corollary gives, every move contributes the same summand
$(\eta/\ell)^2$, and the number of moves is at most \eqref{4.24:eq-flux-form-2}. Hence
\begin{align}
\operatorname{dis}\operatorname{Hol}_U(\Gamma_{\omega})
&\le c\,M\tilde\alpha_0\cdot c''(q+1)^2\Bigl(\frac{M\ell}{\eta}\Bigr)^2
\Bigl(\frac{\eta}{\ell}\Bigr)^2 \notag\\
&= c'''\,(q+1)^2\,M^2\cdot M\tilde\alpha_0 ,\label{4.24:eq-flux-form-1}
\end{align}
where $c''':=c\,c''$.

The right member of \eqref{4.24:eq-flux-form-1} is polynomial in $|\omega|$. It is therefore of
the form required in \eqref{4.24:eq-holonomy-hypothesis-a} with an absolute constant
$C_{\mathfrak{w}}$. One may take $C_{\mathfrak{w}}=9$: writing $p=|\omega|$, the factor count
satisfies $q(\omega)\le 2p$ (as in Proposition~\ref{4.24:prop-field-dependence}), and
$(2p+1)^2\le 9^{p}$ for every integer $p\ge0$, since both sides
equal $1$ at $p=0$ and $9$ at $p=1$, while for $p\ge1$ the ratio of the right side at $p+1$ to
that at $p$ is $9$ and the ratio of the left side is
$\bigl((2p+3)/(2p+1)\bigr)^2\le(5/3)^2<9$. The right member of \eqref{4.24:eq-flux-form-1} is
thus at most $\mathfrak{w}_0\,9^{|\omega|}$, with $\mathfrak{w}_0$ taken equal to
$c'''M^2\cdot M\tilde\alpha_0$. Beside the factor $M\tilde\alpha_0$ that the per-cube condition
itself carries, however, it contains two further powers of $M$.

The radius $\tilde\alpha_0$ is subject to an upper bound depending on $\gamma$. This bound comes
from the requirement, taken by statement in Proposition~\ref{4.24:prop-field-dependence}, that the
bracket
\begin{equation*}
c_{G}'\,O(1)\bigl[5w_{*}(BCM\tilde\alpha_0+\tilde\alpha_1)+K_{\Delta}\gamma_0\bigr]
\end{equation*}
be at most $\vartheta$; here $B$, $C$, $w_{*}$, $K_{\Delta}$, $\gamma_0$ and $c_{G}'$ are the
constants of that bracket, with the meaning they carry there. The extra factor, a fixed power of
$M$,
multiplies a constant fixed before $\gamma$; the $\gamma$-dependent ceiling on $\tilde\alpha_0$ is
accordingly taken again, in sup form, with this factor included. This strengthening is one-sided;
it is not a cap in the sense of Definition~\ref{2.3:def-cap}, and in particular it places no upper
bound on $M$.

The member $BCM\tilde\alpha_0$ of that bracket, in terms of which the budget
$\mathfrak{w}_0\le O(1)\cdot 5w_{*}BCM\tilde\alpha_0$ of
Definition~\ref{4.24:def-holonomy-hypothesis} is written, is reproduced as written only if the
product $BC$ already dominates $c'''M^2$. Whether it does is not decided by the present argument.

For walks of the kind just described, and only for these, the flux bound would yield the holonomy
hypothesis of Definition~\ref{4.24:def-holonomy-hypothesis} at the price just described. This
partial discharge is not asserted here. Per-cube data on the cubes that a walk term reads cannot
supply the flux bound \eqref{4.24:eq-flux-reduction-a} in general, for the following reasons.
\begin{itemize}
\item[(a)] A spanning disc may leave the domain $D(\omega)$ read by the term. Its plaquettes then
lie in cubes that the term never reads, possibly at finer levels: nested large-field collars may
lie inside the loop. There the unit $\ell_P$ of the per-cube condition is smaller. The
contribution $(\eta/\ell_P)^2$ summed over such plaquettes equals the physical area of that part
of the disc divided by $\ell_P^2$, and it grows like $L^2$ raised to the number of levels by which
the finer cubes lie below the level of $\omega$. This growth is not polynomial in $|\omega|$.
\item[(b)] The plaquettes of every spanning disc may lie where the per-cube condition does not hold
at all. This happens inside an excluded large-field component, for instance the component arrested
in the arrest theorem, whose enlarged neighbourhood may be entered by the enlarged localisation
domains of the components that are not arrested. Alternatively, $\Gamma_{\omega}$ may be
non-contractible, winding around a cycle of the torus. In either case no flux bound
\eqref{4.24:eq-flux-reduction-a} exists, and the situation of
Proposition~\ref{4.24:prop-flat-holonomy} (a flat background with non-central holonomy) is
realised.
\end{itemize}

Hence one of two things is required. The first is the holonomy hypothesis
\eqref{4.24:eq-holonomy-hypothesis-a}, taken as a genuine additional input about the background
$U$: a bound on its Wilson loops around the threading loops of the walks, with multiplicity
$C_{\mathfrak{w}}^{|\omega|}$. The second is a reading of the field-dependence requirement taken
by statement in Proposition~\ref{4.24:prop-field-dependence}, namely that the bracket displayed
above be at most $\vartheta$, which tolerates, for the threading walks, the holonomy members
\begin{equation*}
s^{0}_{\omega}\operatorname{tr}_{\mathfrak{g}}\bigl(\operatorname{Ad}_{\mathfrak{h}_{\omega}}
-1\bigr)\cdot m(\omega)
\end{equation*}
(cf.\ Proposition~\ref{4.24:prop-flat-holonomy}). Which of the two is supplied is not decided
here.
\end{remark}

\begin{definition}[Conventions on bond sets, collar variables and dimensions]
\label{4.24:not-letter-conventions}
The following four conventions on letters are in force in this section.
\begin{enumerate}
\item The set of integration bonds at level $j$ is written $\mathfrak{Bo}$. It is the disjoint
union of the set $s$ of survivors' bonds and the set $\mathcal{W}$ of arrested bonds:
\begin{equation*}
\mathfrak{Bo}=s\sqcup\mathcal{W}.
\end{equation*}
The letter $\mathfrak{b}$ is never used for this set. The symbols $\mathfrak{b}_j$ and
$\mathfrak{b}'_j$ denote radii.
\item The new collar variable on the bonds of the arrested component $W$ is written $x$. A
point at which a sum of norms is rooted is written $y$. Thus a supremum over the root point
of such a sum is a supremum over $y$, and a quantity attached to the root point is written
as a function of $y$.
\item We put $n_{\mathfrak{g}}:=\dim\mathfrak{g}$, the dimension of the Lie algebra
$\mathfrak{g}$ of the gauge group. The letter $n$ without subscript is the
stage index.
\item The half-line variable of \cite[(3.36)]{B-CMP119}, written $\lambda$ there, is written
$\xi\ge 0$ here. The letter $\lambda$ is reserved for an exterior parameter of this chapter
and denotes nothing else. The
symbol $\lambda_0$ denotes the fixed positive number
so denoted in \cite{B-CMP119}.
\end{enumerate}
\end{definition}

\begin{definition}[Scales, filing levels, and the rim and collar of a stage]
\label{4.24:not-scales-rim-collar}
Fix the step $k$ and write $\eta=L^{-k}$; in the notation introduced below, $\eta=s_0$ is the side
of a level-$0$ cell.

\emph{Scales.} For a level $m$ (here $m$ is a level, not the index of the torus $\mathbb{T}_m$) put
\begin{equation*}
s_m:=L^{m-k},
\end{equation*}
the side of a level-$m$ cell; thus $s_k=1$. We write $\pi_m$ for the set of cells of side $Ms_m$ of the
standard nested grids with common origin, and $\mathbb{D}_m$ for the set of face-connected unions of
$\pi_m$-cells. For a set $X$, $\tau(X)$ (not to be confused with $\tau_h$) is the length of a shortest
tree contained in $X$ which meets every closed cell of $X$, and
\begin{equation*}
d_m(X):=\frac{\tau(X)}{Ms_m}.
\end{equation*}
For a set $Y$, $[Y]_m$ is the smallest $\mathbb{D}_m$-domain containing $Y$.

\emph{Filing levels.} Following \cite[(1.63)]{B-CMP122a}, the filing level of a region $X$ relative to
the level $j$ of the stage (fixed below as $j=h+n$) is
\begin{equation}\label{4.24:eq-scales-rim-collar-1}
m(X):=\min\{\,m\ge j+1:\ X\subset\Omega^{c}_{m+1}\,\}.
\end{equation}
In \eqref{4.24:eq-scales-rim-collar-1} the small-field regions $\Omega_m$ are those of step $k$,
formed with one large-field region $Z$ held fixed for all configurations of the class that the stage
integration treats together, as in the freezing convention of the stage integration.

\emph{Stages.} From here on, $N$ denotes a fixed depth, not the $N$ of $\mathrm{SU}(N)$, and
\begin{equation*}
h:=k-N
\end{equation*}
is the level $N$ steps below $k$; $h$ is a level, not a history.
For the stage index $n$ put $j:=h+n$. Stage $n+1$ (whose index is also written $l=n+1$) integrates the
pair
\begin{equation*}
\phi:=(A_j,B_j).
\end{equation*}
\begin{itemize}
\item The variable $A_j$ is integrated on the \emph{rim}
\begin{equation*}
\mathsf{R}_j:=Z_{j+1}\cap\Omega_{j+1}\cap Z
\end{equation*}
under the characteristic function $\chi^{(j)}$ of the box
\begin{equation*}
\bigl\{\sup|A_j|<\mathfrak{b}_j\bigr\},\qquad \mathfrak{b}_j:=g_j^{-1}\delta_j .
\end{equation*}
\item The variable $B_j$ is integrated on the starred \emph{collar} bonds
\begin{equation*}
\mathcal{C}_j:=\bigl(\Omega^{c}_{j+1}\setminus Z''_{j+1}\bigr)^{(j)*}
\end{equation*}
under the characteristic function $\chi'^{(j)}$, whose threshold is
\begin{equation*}
\mathfrak{b}'_j:=\frac{\delta'_j}{g_j};
\end{equation*}
this threshold is the one written $T_j$ in the description of the stages of the collar integration.
\end{itemize}
Here the regions $Z_{j+1}$, $Z''_{j+1}$ and the starring operation $(\cdot)^{(j)*}$ on bond sets are
those of \cite[p.~188, (1.53)--(1.63)]{B-CMP122a}; the box radius $\mathfrak{b}_j$, with its constant
$\delta_j$, is the one fixed in the statement on the rim box of the stage integration, and the
threshold $\mathfrak{b}'_j$, with its constant $\delta'_j$, the one fixed in the description of the
stages of the collar integration. In
accordance with the conventions on bond sets, collar variables and dimensions
(Definition~\ref{4.24:not-letter-conventions}), $\mathfrak{b}_j$ and $\mathfrak{b}'_j$ are radii and
never sets of bonds.

\emph{The band of a stage.} Put
\begin{equation*}
R_n:=Z_{j+1}\setminus Z''_{j+1}.
\end{equation*}
The band $R_n$ contains the rim and the collar, and is contained in the small-field complement of
level $j+2$; both inclusions are part of the description of the stages of the collar integration:
\begin{equation}\label{4.24:eq-scales-rim-collar-2}
\mathsf{R}_j\cup\mathcal{C}_j\subset R_n\subset\Omega^{c}_{j+2}.
\end{equation}
Finally
\begin{equation*}
R^{(n+1)}:=R_n\cap Z^{[n+1]},
\end{equation*}
where $Z^{[n+1]}\subset Z$ is the union of the components of $Z$ in the run at stage $n+1$, as fixed by
the freezing convention of the stage integration.
\end{definition}

\begin{definition}[Arrested components, survivors and integration bonds]\label{4.24:def-survivors-bonds}
Fix the step $k$ and the stage $n+1$. The level $h$, the index $j=h+n$, the rim $\mathsf{R}_j$, the collar $\mathcal{C}_j$, the characteristic functions $\chi^{(j)}$ and $\chi'^{(j)}$, the radii $\mathfrak{b}_j$ and $\mathfrak{b}'_j$, the band $R_n$, the field $\phi=(A_j,B_j)$ integrated at this stage and the cell partitions $\pi_m$ are as in Definition~\ref{4.24:not-scales-rim-collar}. For a function $f$ on bonds and a set $Y$ we write $f|_Y$ for its restriction to $Y$, and likewise $\mathbb{B}|_Y$ for the restriction to $Y$ of a background structure $\mathbb{B}$.

\begin{enumerate}
\item[(a)] \emph{Components in the run.} Let $Z$ denote the set of all components of the classes (i) and (ii) of \cite{B-CMP122a} that enter the first stage. A component is \emph{in the run} at stage $l$ if it has not been excluded at any of the stages $1,\dots,l-1$. In \cite{B-CMP122a} the symbol $Z$ is used again, after each exclusion of components, for the components that remain. Here we write instead $Z^{[l]}$ for the set of components still in the run at stage $l$.

\item[(b)] \emph{Freezing.} Let $W$ be a component that leaves the run after $n^+$ performed stages. For every $l>n^+$ it keeps the background structure
\begin{equation*}
\mathfrak{B}|_W:=\mathbb{B}^{(n^+)}_k|_W ,
\end{equation*}
and for $l>n^+$ the expression carried on $W$ is the one that the freezing statement for components out of the run assigns to such a component. For $l\le n^+$, in particular at the stage at which $W$ leaves the run, the expression carried on $W$ is that of \cite[(1.64)--(1.69)]{B-CMP122a} at stage $l$; we denote it by $\mathfrak{L}_l$.

\item[(c)] \emph{Components reaching the split.} $Z^{\sharp}$ is the set of components of $Z^{[n+1]}$ that reach the $\chi'^{(j)}$-split \cite[(1.55)]{B-CMP122a} at stage $n+1$, that is, the components on which the operations \cite[(1.53)--(1.55)]{B-CMP122a} are carried out. A component may leave the run at this stage before \cite[(1.53)]{B-CMP122a} is carried out, namely when the test of \cite[(1.27)]{B-CMP122a} applied to that component selects one of the branches of the case distinction of the stage in which no prepared stage is performed. Such a component does not belong to $Z^{\sharp}$.

\item[(d)] \emph{Arrested components.} $\mathfrak{W}\subset Z^{\sharp}$ is the set of components that take the factor $1-\chi'^{(j)}$ in \cite[(1.55)]{B-CMP122a}. To $W\in\mathfrak{W}$ corresponds the arrest datum
\begin{equation*}
\mathfrak{a}=(W,k,n^+,\beta),
\end{equation*}
which records the component, the step, the number $n^+=n+1$ of performed stages and the branch $\beta$ taken, here the branch that carries the factor $1-\chi'^{(j)}$. One term of the decomposition of unity in \cite[(1.55)]{B-CMP122a} is fixed throughout, so that $\mathfrak{W}$ is a fixed datum.

\item[(e)] \emph{Data attached to $W\in\mathfrak{W}$.} $Z_W\supset W$ is the component of $Z$ containing $W$. We put
\begin{equation*}
S^P_W:=\mathcal{C}_j\cap W,\qquad S^P:=\bigcup_{W\in\mathfrak{W}}S^P_W .
\end{equation*}
The letter $x$ denotes the new collar variable on $S^P$, that is, the variable introduced on the collar by the change of variables of the arrested stage. We write $\sigma_0(Z_W)$ for the family of sets attached to $Z_W$ in the description of the arrested branch that fixes the datum $\mathfrak{a}$, and we put
\begin{equation*}
Z_W^{+}:=\bigcup\sigma_0(Z_W),\qquad \mathfrak{B}^{\circ}:=\mathbb{B}^{(n+1)}_k|_{Z_W}.
\end{equation*}

\item[(f)] \emph{Survivors.} The survivors' region is
\begin{equation*}
Z^{s}:=Z^{\sharp}\setminus\bigcup\mathfrak{W},
\end{equation*}
the region whose level-$j$ variables are integrated at this stage by \cite[(1.60)]{B-CMP122a}. In the notation of item (a), $Z^{s}=Z^{[n+2]}$; this is shown below.

\item[(g)] \emph{Bond sets and variables.} $\mathfrak{Bo}$ is the set of level-$j$ integration bonds of $Z^{\sharp}$, namely the rim bonds of $\mathsf{R}_j\cap Z^{\sharp}$ and the starred collar bonds of $\mathcal{C}_j\cap Z^{\sharp}$. It is the disjoint union
\begin{equation*}
\mathfrak{Bo}=s\sqcup\mathcal{W},
\end{equation*}
where $s$ is the set of those bonds that belong to $Z^{s}$ and $\mathcal{W}$ the set of those that belong to $\bigcup\mathfrak{W}$. We put $\phi_s:=\phi|_s$ and, for $W\in\mathfrak{W}$,
\begin{equation*}
\phi_W:=\bigl(A_j|_W,\;x|_{S^P_W}\bigr),
\end{equation*}
held as exterior variables of the stage. The variables $\phi_W$ are real; $A_j|_W$ lies in the support of $\chi^{(j)}$, so $|A_j|<\mathfrak{b}_j$. On the support of \cite[(1.59)]{B-CMP122a} the new collar variable satisfies $|x|\le 2.2\,\mathfrak{b}'_j$, and the old collar variable satisfies $|B_j|\le 2.4\,\mathfrak{b}'_j$, by the change-of-variables lemma on the collar of the arrested stage. We put
\begin{equation*}
\varrho^{B}_j:=6.6\,\mathfrak{b}'_j ,
\end{equation*}
the radius of the disc $|\mathsf{B}|<\varrho^{B}_j$ in the collar variable $\mathsf{B}$ on which the analyticity statement for an arrested stage places every element variable $v_e$ when the arrest condition of that statement holds.

\item[(h)] \emph{Cubes.} $\mathsf{Q}$ is the set of $\pi_{j+1}$-cells meeting $\mathfrak{Bo}$; $\mathsf{Q}^{s}$ is the set of those meeting $s$ and $\mathsf{Q}^{\mathfrak{W}}$ the set of those meeting $\mathcal{W}$. For $b\in\mathfrak{Bo}$, $\operatorname{cube}(b)$ is the cell of $\mathsf{Q}$ holding $b$, chosen by a fixed rule, and $h(b)$ is the host cube (item (i) below) holding $b$. Both $\operatorname{cube}(b)$ and $h(b)$ contain an endpoint of $b$.

\item[(i)] \emph{Host cubes and counting constants.} The \emph{host cubes} are the cells of $\mathbb{B}^{(n)}_k$ in proper scales in the sense of \cite{B-CMP119}. Each host cube is a $\pi_t$-cell for some $t\le k$, and the decoupling parameters $\sigma'$ of the collar kernel are attached to all host cubes of $\mathbb{B}^{(n)}_k$. The constant $\mathsf{c}_{\mathrm{an}}=\mathsf{c}_{\mathrm{an}}(d,L)$ is such that, for every integer $\nu\ge1$ and every host cube, the number of wall-connected families of $\nu$ host cubes containing the given one is at most $e^{\mathsf{c}_{\mathrm{an}}\nu}$. The adjacency number is
\begin{equation*}
\mathsf{n}_{\mathrm{adj}}:=1+2dL^{d-1}.
\end{equation*}

\item[(j)] \emph{Scale constants.} For every level $i$ let $R_i$ be as in Lemma~\ref{4.4:lem-collar-avoidance}. We put
\begin{equation*}
\check R:=\max_{h<i\le k}R_i,\qquad R_k:=\min_{h<i\le k}R_i
\end{equation*}
(an abuse of notation: the minimum is again written $R_k$). This minimum satisfies $e^{-R_k}\le E(x_k)$, by the lower bound on the $R_i$ in terms of the exponential smallness $E(x_k)$ that accompanies their definition. Let
\begin{equation*}
V:=\sup_{n,\bar c}\;\sum_{\square}e^{-\frac12\delta\, d_{\mathrm{sc}}(\square,\,R_n\cap\bar c)},
\end{equation*}
the supremum being taken over the stages $n$ and the components $\bar c$ of $Z$; for each $n$ the sum runs over the $\pi_{j+1}$-cells $\square$, with $j=h+n$, $d_{\mathrm{sc}}(\square,Y)$ denotes the scaled distance between the cell $\square$ and a set $Y$, which vanishes when $\square$ meets $Y$, and $\delta>0$ is the decay constant of this cube sum. With a constant $C_V=C_V(d)$ depending on the dimension only, we put
\begin{equation*}
\bar V:=C_V\,(100L\check R)^d .
\end{equation*}
By the cube-sum bound for the components of $Z$, $V\le C_V(d)(100L\check R)^d=\bar V$. Every term of the sum defining $V$ is positive, and every cell of $\mathsf{Q}$ meeting $R_n\cap\bar c$ is at scaled distance zero from $R_n\cap\bar c$ and so contributes $e^{0}=1$. Hence, for every component $\bar c$ of $Z$,
\begin{equation}\label{4.24:eq-survivors-bonds-1}
\bar V\;\ge\;V\;\ge\;\#\{\square\in\mathsf{Q}:\ \square\ \text{meets}\ R_n\cap\bar c\}.
\end{equation}
The bound \eqref{4.24:eq-survivors-bonds-1} is used in this form for the arrested components. Finally,
\begin{equation*}
\varepsilon_Y:=2e^{-6R_k}.
\end{equation*}
\end{enumerate}
\end{definition}

\begin{proof}
Two assertions of the definition require an argument: the identity $Z^{s}=Z^{[n+2]}$ in item (f), and the placing of the collar values of item (g) in the disc of radius $\varrho^{B}_j$.

\emph{The identity $Z^{s}=Z^{[n+2]}$.} Every exclusion of components at stage $n+1$ precedes the integration \cite[(1.60)]{B-CMP122a}. In \cite[p.~188]{B-CMP122a} the displays \cite[(1.55)--(1.59)]{B-CMP122a}, at which components leave, come before \cite[(1.60)]{B-CMP122a}, and the same order holds at the first stage. No exclusion occurs in \cite[(1.61)--(1.63)]{B-CMP122a}.

A component of $Z^{[n+1]}$ leaves the run at stage $n+1$ in one of two ways. Either it leaves before \cite[(1.53)]{B-CMP122a}, and then it is not in $Z^{\sharp}$ by item (c). Or it reaches the split \cite[(1.55)]{B-CMP122a} and takes the factor $1-\chi'^{(j)}$ in the fixed term of the decomposition of unity, and then it lies in $\mathfrak{W}$ by item (d). Hence the components still in the run at stage $n+2$ are exactly the components of $Z^{\sharp}$ not in $\mathfrak{W}$. Their union is $Z^{\sharp}\setminus\bigcup\mathfrak{W}=Z^{s}$.

\emph{The collar values.} By item (g) the new collar variable satisfies $|x|\le 2.2\,\mathfrak{b}'_j$ on the support of \cite[(1.59)]{B-CMP122a}, and the old collar variable satisfies $|B_j|\le 2.4\,\mathfrak{b}'_j$. Since $\mathfrak{b}'_j=\delta'_j/g_j>0$, both sets lie in $\{|x|\le 2.4\,\mathfrak{b}'_j\}$. Because
\begin{equation*}
2.2\,\mathfrak{b}'_j\;\le\;2.4\,\mathfrak{b}'_j\;<\;6.6\,\mathfrak{b}'_j=\varrho^{B}_j ,
\end{equation*}
the set $\{|x|\le 2.4\,\mathfrak{b}'_j\}$ is contained in the disc $|\mathsf{B}|<\varrho^{B}_j$.
\end{proof}

\begin{lemma}[Separation of components and cells meeting one component]
\label{4.24:lem-component-separation}
Work at one stage of the large-field step at level $k$, at which level-$j$ variables are
integrated, $j+1\le k$, in the notation of
Definition~\ref{4.24:def-survivors-bonds}. In particular we use the components $Z^{\sharp}$, the
arrested components $\mathfrak{W}\subset Z^{\sharp}$, the survivors' region
$Z^{s}=Z^{\sharp}\setminus\bigcup\mathfrak{W}$ and the integration bonds
$\mathfrak{Bo}=s\sqcup\mathcal{W}$. We also use the cells $\mathsf{Q}$,
$\mathsf{Q}^{s}$, $\mathsf{Q}^{\mathfrak{W}}$, the host cube $h(b)$ of a bond $b$, and
the separation radius $R_k$ as fixed in Definition~\ref{4.24:def-survivors-bonds}.

An \emph{$M$-cube} is a $\pi_k$-cell of the level-$k$ lattice; its side is $M$ in level-$k$
units. Distances, written $\operatorname{dist}$, are taken in one of two metrics:
\begin{itemize}
\item[$(\infty)$] the $\ell^\infty$ distance;
\item[(E)] the Euclidean distance.
\end{itemize}
In either metric, the distance between two sets is the infimum of the distances between their
points, and the diameter of a set is the supremum of the distances between its points.
Hypotheses and conclusion are read in the same metric, and the lemma holds for either choice of
the metric.

For integration bonds $b,b'\in\mathfrak{Bo}$ write $\bar c$ and $\bar c'$ for the components of
$Z^{\sharp}$ containing them. A family $\mathcal{F}$ of host cubes is \emph{wall-connected} if any
two of its members are joined by a chain of members, consecutive members having
$(d-1)$-dimensional faces one of which is contained in the other, $d=4$ (such a face is a
\emph{wall}). Assume the following.
\begin{itemize}
\item[(A)] The separation property of the components: distinct components are at least $R_k$
$M$-cubes apart, the $M$-cubes being the $\pi_k$-cells. It is used in the following two forms.
First, for all integration bonds $b,b'$ with $\bar c\neq\bar c'$ and every wall-connected family
$\mathcal{F}$ of host cubes,
\begin{equation}\label{4.24:eq-component-separation-a}
h(b),h(b')\in\mathcal{F}
\quad\Longrightarrow\quad \#\mathcal{F}\ge R_k .
\end{equation}
Second, let $\square,\square'$ be cells of $\mathsf{Q}$ (or host cubes) such that $\square$
meets an integration bond of a component $\bar c$ and $\square'$ meets an
integration bond of a component $\bar c'\neq\bar c$. Then
\begin{equation}\label{4.24:eq-component-separation-b}
\operatorname{dist}(\square,\square')\ \ge\ R_k\,M .
\end{equation}
\item[(B)] Among the conditions on the parameters of the step, the separation radius satisfies
$R_k\ge 3$. This condition is imposed through a ceiling on the reference coupling $g$ only (it
bounds $g$ from above), and $R_k$ is non-decreasing as $g_k\downarrow 0$.
\end{itemize}
Let $P$ be either a single $\pi_{j+1}$-cell or the union of two wall-adjacent
$\pi_{j+1}$-cells. Then
\begin{equation}\label{4.24:eq-component-separation-c}
P\ \text{meets the integration bonds of at most one component.}
\end{equation}
In particular $\mathsf{Q}^{s}\cap\mathsf{Q}^{\mathfrak{W}}=\emptyset$.
\end{lemma}

\begin{proof}
Since $j+1\le k$, the $\pi_{j+1}$-cells are finer than or equal to the $\pi_k$-cells. Each
$\pi_{j+1}$-cell therefore has side at most $M$ in level-$k$ units. Two wall-adjacent
$\pi_{j+1}$-cells have union contained in a box whose sides are at most $2M,M,M,M$. We bound the
diameter of $P$ in each metric.

In the $\ell^\infty$ metric, the diameter of $P$ is at most $2M$. By hypothesis (B),
\begin{equation*}
2M<3M\le R_kM .
\end{equation*}

In the Euclidean metric, the diameter of $P$ is at most $\sqrt{4+1+1+1}\,M=\sqrt7\,M$; this bound
also covers a single cell, whose Euclidean diameter is at most $\sqrt4\,M$. Since $7<9$,
hypothesis (B) gives
\begin{equation*}
\sqrt7\,M<3M\le R_kM .
\end{equation*}
In both metrics, therefore, the diameter of $P$ is strictly less than $R_kM$.

Suppose that $P$ meets an integration bond $b$ of a component $\bar c$ and an integration bond
$b'$ of a component $\bar c'\neq\bar c$. Let $\square$ be a $\pi_{j+1}$-cell contained in $P$
that meets $b$, and let $\square'$ be one that meets $b'$; if $P$ is a single cell,
$\square=\square'=P$. Both $\square$ and $\square'$ meet $\mathfrak{Bo}$, so both are cells of
$\mathsf{Q}$. Since $\square$ meets $b$ and $\square'$ meets $b'$, integration bonds of the
distinct components $\bar c$ and $\bar c'$, the second form of (A) applies, and
\eqref{4.24:eq-component-separation-b} gives
\begin{equation*}
\operatorname{dist}(\square,\square')\ge R_kM .
\end{equation*}
On the other hand $\square,\square'\subset P$, so $\operatorname{dist}(\square,\square')$ is at
most the diameter of $P$, which is less than $R_kM$. This contradiction proves
\eqref{4.24:eq-component-separation-c}.

Finally, let $\square\in\mathsf{Q}^{s}\cap\mathsf{Q}^{\mathfrak{W}}$. Then $\square$ meets a bond
of $s$, which is an integration bond of a component of $Z^{s}=Z^{\sharp}\setminus\bigcup\mathfrak{W}$.
It also meets a bond of $\mathcal{W}$, which is an integration bond of a component in
$\mathfrak{W}$. These two components are distinct, since no component of $Z^{s}$ lies in
$\mathfrak{W}$. This contradicts \eqref{4.24:eq-component-separation-c} applied to the single
$\pi_{j+1}$-cell $P=\square$. Hence $\mathsf{Q}^{s}\cap\mathsf{Q}^{\mathfrak{W}}=\emptyset$.
\end{proof}

\begin{definition}[Fibres, pairings and gauge action on bond kernels]\label{4.24:not-fibres-pairings}
Throughout this section $\mathfrak{Bo}=s\sqcup\mathcal{W}$ is the set of level-$j$ integration bonds
of $Z^{\sharp}$ fixed in Definition~\ref{4.24:def-survivors-bonds}, and $G$ denotes the gauge group.
For a bond $b$ we write $b_-$ for its initial site.

\emph{The fibre.} Let $\mathfrak{g}:=\mathfrak{su}(2)$, the Lie algebra of $G$, and let $(\cdot,\cdot)$
be the $\operatorname{Ad}$-invariant positive inner product on $\mathfrak{g}$, extended
$\mathbb{C}$-bilinearly to the complexification $\mathfrak{g}^{\mathbb{C}}$. Let
$X\mapsto\bar X$ be the conjugation of $\mathfrak{g}^{\mathbb{C}}$ with respect to $\mathfrak{g}$,
and put
\begin{equation*}
\langle X,Y\rangle:=(\bar X,Y),\qquad X,Y\in\mathfrak{g}^{\mathbb{C}};
\end{equation*}
this is a positive Hermitian form. We write $n_{\mathfrak{g}}:=\dim\mathfrak{g}$; for
$A\in\operatorname{End}(\mathfrak{g}^{\mathbb{C}})$, $|A|$ is the operator norm of $A$ on
$(\mathfrak{g}^{\mathbb{C}},\langle\cdot,\cdot\rangle)$, $A^{t}$ its $(\cdot,\cdot)$-transpose,
defined by $(A^{t}X,Y)=(X,AY)$, and $\operatorname{tr}_{\mathfrak{g}}A$ its trace. Then
$|\operatorname{tr}_{\mathfrak{g}}A|\le n_{\mathfrak{g}}|A|$.

\emph{The bond space.} Let $\mathcal{H}:=\bigoplus_{b\in\mathfrak{Bo}}\mathfrak{g}^{\mathbb{C}}_b$,
each $\mathfrak{g}^{\mathbb{C}}_b$ a copy of $\mathfrak{g}^{\mathbb{C}}$, with
\begin{equation*}
(\phi,\psi):=\sum_{b}(\phi(b),\psi(b)),\qquad
\langle\phi,\psi\rangle:=(\bar\phi,\psi),\qquad
\|\phi\|^{2}:=\langle\phi,\phi\rangle,
\end{equation*}
where $\bar\phi(b):=\overline{\phi(b)}$. For $E\subset\mathfrak{Bo}$ let
$\mathcal{H}_E:=\{\phi\in\mathcal{H}:\phi(b)=0\ \text{for}\ b\notin E\}$, and let $P_E$ be the
coordinate projection, $(P_E\phi)(b)=\phi(b)$ for $b\in E$ and $(P_E\phi)(b)=0$ otherwise. Then $P_E$ is
$(\cdot,\cdot)$-symmetric and $\langle\cdot,\cdot\rangle$-orthogonal, $\|P_E\|\le1$, and
$P_EP_{E'}=P_{E\cap E'}$.

\emph{Kernels.} A kernel is a family $K=(K(b,b'))_{b,b'\in\mathfrak{Bo}}$ with
$K(b,b')\in\operatorname{End}(\mathfrak{g}^{\mathbb{C}})$; it acts on $\mathcal{H}$ by
$(K\phi)(b):=\sum_{b'}K(b,b')\phi(b')$. We put $K_{EE'}:=P_EKP_{E'}$, regarded as a map
$\mathcal{H}_{E'}\to\mathcal{H}_E$, and $K^{T}(b,b'):=K(b',b)^{t}$; $\|K\|$ is the operator norm of $K$ on
$(\mathcal{H},\langle\cdot,\cdot\rangle)$.

\emph{Gauge action.} For $u\in G$, $\operatorname{Ad}_u$ (extended $\mathbb{C}$-linearly to
$\mathfrak{g}^{\mathbb{C}}$) is $(\cdot,\cdot)$-orthogonal and commutes with conjugation, hence is
$\langle\cdot,\cdot\rangle$-unitary. For a $G$-valued function $u$ on sites put
$\mathcal{U}_u:=\bigoplus_{b}\operatorname{Ad}_{u(b_-)}$, that is
$(\mathcal{U}_u\phi)(b)=\operatorname{Ad}_{u(b_-)}\phi(b)$. Then $\mathcal{U}_u$ is unitary and
bond-diagonal, and
\begin{equation}\label{4.24:eq-fibres-pairings-1}
P_E\,\mathcal{U}_u=\mathcal{U}_u\,P_E\qquad\text{for every }E\subset\mathfrak{Bo}.
\end{equation}

\emph{The domain of pairs.} Let $\mathfrak{K}_P=\mathfrak{K}_P(\mathbb{B}^{(n)}_k;\gamma_0)$ denote the
space of complex background pairs of \cite{B-CMP109} built on the background structure
$\mathbb{B}^{(n)}_k$ with radius $\gamma_0$. We
write $\mathfrak{D}^{\flat}_{n+1}(X)$ for the analyticity domain of the stage-$(n+1)$ densities; it
is the intersection
\begin{equation*}
\mathfrak{D}^{\flat}_{n+1}(X)=\mathfrak{N}^{(n+1)}\cap\mathfrak{L}_{n+1}\cap\mathfrak{I}_{n+1}
\end{equation*}
of the three domains $\mathfrak{N}^{(n+1)}$, $\mathfrak{L}_{n+1}$ and $\mathfrak{I}_{n+1}$ of the
stage-$(n+1)$ construction, with $\mathfrak{L}_{n+1}$ the domain carried on arrested components as in
Definition~\ref{4.24:def-survivors-bonds}, and with membership of a pair in $\mathfrak{K}_P$ read in the
ambient sense: a pair on a set $Y$ belongs to $\mathfrak{K}_P$ if it is the
restriction to $Y$ of a pair of $\mathfrak{K}_P$ defined on the whole domain of $\mathfrak{B}$. We write
$\mathfrak{D}^{\flat,W}_{n+1}(X)\subset\mathfrak{D}^{\flat}_{n+1}(X)$ for the subset obtained by
imposing, in addition, a declared restriction on the admitted pairs relative to the arrested
component $W$; this restricts the admitted pairs and changes nothing else.

A \emph{real point} $z_{\mathbb{R}}$ of $\mathfrak{D}^{\flat}_{n+1}(X)$ is a pair
$(\mathbb{U},\mathbb{J})$ with $\mathbb{U}$ $G$-valued and $\mathbb{J}$ $\mathfrak{g}$-valued (here
$\mathbb{U}$ is the background configuration of the stage determined by the block field $V$, and
$\mathbb{J}$ the current determined by $\mathbb{U}$); at a real
point every kernel of this section is real, that is, each $K(b,b')$ maps $\mathfrak{g}$ into
$\mathfrak{g}$. The trivial pair is $(1,0)$: $\mathbb{U}\equiv1$, $\mathbb{J}\equiv0$.
\end{definition}

\begin{proof}
\emph{The fibre form.} Write $X=X_1+iX_2$ with $X_1,X_2\in\mathfrak{g}$, so that $\bar X=X_1-iX_2$. The
form $\langle X,Y\rangle=(\bar X,Y)$ is linear in $Y$ and conjugate-linear in $X$. Since $(\cdot,\cdot)$
is real on $\mathfrak{g}\times\mathfrak{g}$ and extended bilinearly,
$\overline{(X,Y)}=(\bar X,\bar Y)$ for $X,Y\in\mathfrak{g}^{\mathbb{C}}$;
hence, by symmetry of $(\cdot,\cdot)$,
\begin{equation*}
\overline{\langle X,Y\rangle}=(X,\bar Y)=(\bar Y,X)=\langle Y,X\rangle.
\end{equation*}
Moreover
\begin{equation*}
\langle X,X\rangle=(X_1-iX_2,X_1+iX_2)=(X_1,X_1)+(X_2,X_2),
\end{equation*}
the cross terms $i(X_1,X_2)-i(X_2,X_1)$ cancelling, and this is positive unless $X=0$. The same
computation bond by bond shows that $\langle\cdot,\cdot\rangle$ on $\mathcal{H}$ is Hermitian and
positive, since $\langle\phi,\psi\rangle=\sum_b\langle\phi(b),\psi(b)\rangle$.

\emph{The trace bound.} Let $(e_i)_{i=1}^{n_{\mathfrak{g}}}$ be a
$\langle\cdot,\cdot\rangle$-orthonormal basis of $\mathfrak{g}^{\mathbb{C}}$, whose complex dimension is
$n_{\mathfrak{g}}$. Then $Ae_i=\sum_j\langle e_j,Ae_i\rangle e_j$, so
$\operatorname{tr}_{\mathfrak{g}}A=\sum_i\langle e_i,Ae_i\rangle$, and by the Cauchy--Schwarz inequality
$|\langle e_i,Ae_i\rangle|\le|A|$ for each $i$; summing gives
$|\operatorname{tr}_{\mathfrak{g}}A|\le n_{\mathfrak{g}}|A|$.

\emph{The projections.} We have $(P_E\phi,\psi)=\sum_{b\in E}(\phi(b),\psi(b))=(\phi,P_E\psi)$, so
$P_E$ is $(\cdot,\cdot)$-symmetric. The projection $P_E$ commutes with bondwise conjugation, so
\begin{equation*}
\langle P_E\phi,\psi\rangle=(P_E\bar\phi,\psi)=(\bar\phi,P_E\psi)=\langle\phi,P_E\psi\rangle;
\end{equation*}
with
$P_E^2=P_E$ this makes $P_E$ a $\langle\cdot,\cdot\rangle$-orthogonal projection, and
$\|P_E\phi\|^2=\sum_{b\in E}\langle\phi(b),\phi(b)\rangle\le\|\phi\|^2$ gives $\|P_E\|\le1$. Finally,
writing $\mathbf{1}_E$ for the indicator function of $E\subset\mathfrak{Bo}$,
$P_EP_{E'}$ multiplies $\phi(b)$ by $\mathbf{1}_E(b)\mathbf{1}_{E'}(b)=\mathbf{1}_{E\cap E'}(b)$, so
$P_EP_{E'}=P_{E\cap E'}$.

\emph{The gauge action.} For $u\in G$ and $X,Y\in\mathfrak{g}$, $(\operatorname{Ad}_uX,\operatorname{Ad}_uY)=(X,Y)$
by $\operatorname{Ad}$-invariance; both sides are $\mathbb{C}$-bilinear in $(X,Y)$, so the identity holds on
$\mathfrak{g}^{\mathbb{C}}$. Since $\operatorname{Ad}_u$ maps $\mathfrak{g}$ into $\mathfrak{g}$,
$\operatorname{Ad}_u(X_1+iX_2)=\operatorname{Ad}_uX_1+i\operatorname{Ad}_uX_2$ with both terms in
$\mathfrak{g}$, and conjugating gives $\overline{\operatorname{Ad}_uX}=\operatorname{Ad}_u\bar X$. Therefore
\begin{equation*}
\langle\operatorname{Ad}_uX,\operatorname{Ad}_uY\rangle=(\operatorname{Ad}_u\bar X,\operatorname{Ad}_uY)
=(\bar X,Y)=\langle X,Y\rangle,
\end{equation*}
and $\operatorname{Ad}_u$ is invertible with inverse $\operatorname{Ad}_{u^{-1}}$; so it is unitary. Applying
this on each bond,
\begin{equation*}
\langle\mathcal{U}_u\phi,\mathcal{U}_u\psi\rangle
=\sum_b\langle\operatorname{Ad}_{u(b_-)}\phi(b),\operatorname{Ad}_{u(b_-)}\psi(b)\rangle
=\langle\phi,\psi\rangle,
\end{equation*}
and $\mathcal{U}_u$ has inverse $\mathcal{U}_{u^{-1}}$, where $u^{-1}$ is taken sitewise; hence
$\mathcal{U}_u$ is unitary. Its kernel is $\delta_{bb'}\operatorname{Ad}_{u(b_-)}$, so it is
bond-diagonal. Consequently, for every $b$,
\begin{equation*}
(P_E\mathcal{U}_u\phi)(b)=\mathbf{1}_E(b)\operatorname{Ad}_{u(b_-)}\phi(b)
=\operatorname{Ad}_{u(b_-)}\bigl(\mathbf{1}_E(b)\phi(b)\bigr)=(\mathcal{U}_uP_E\phi)(b),
\end{equation*}
which is \eqref{4.24:eq-fibres-pairings-1}.
\end{proof}

\begin{lemma}[Polarisation of a bond kernel]\label{4.24:lem-kernel-polarisation}
Use the fibres, pairings and gauge conjugation of Notation~\ref{4.24:not-fibres-pairings}. For
$A\in\operatorname{End}(\mathfrak{g}^{\mathbb{C}})$ let $A^{t}$ be its transpose with respect to the
bilinear form, defined by $(AX,Y)=(X,A^{t}Y)$ for all $X,Y\in\mathfrak{g}^{\mathbb{C}}$. Let $|A|$ be
its operator norm with respect to $\langle\cdot,\cdot\rangle$. For a kernel
$K=(K(b,b'))_{b,b'\in\mathfrak{Bo}}$ let $K^{T}$ be the kernel with entries
$K^{T}(b,b'):=K(b',b)^{t}$, let $\|K\|$ be the operator norm of $K$ on $\mathcal{H}$ with respect to
$\langle\cdot,\cdot\rangle$, and put
\begin{equation}\label{4.24:eq-kernel-polarisation-a}
  \mathcal{S}K:=\tfrac12\bigl(K+K^{T}\bigr).
\end{equation}
Then the following hold.
\begin{enumerate}
\item[(a)] \emph{Form.} $(\varphi,\mathcal{S}K\varphi)=(\varphi,K\varphi)$ for every $\varphi\in\mathcal{H}$.
  In particular, for $\psi$ in the real subspace $\bigoplus_{b}\mathfrak{g}_b$ the Gaussian weights
  $e^{-\frac12(\psi,K\psi)}$ and $e^{-\frac12(\psi,\mathcal{S}K\psi)}$ coincide, and so do their
  integrals over $\psi$.
\item[(b)] \emph{Holomorphy.} $\mathcal{S}$ is $\mathbb{C}$-linear in the entries. Hence if the entries
  of $K=K(z)$ depend holomorphically on a complex parameter $z$, so do those of $\mathcal{S}K(z)$.
\item[(c)] \emph{Gauge covariance.} For every gauge conjugation $\mathcal{U}_u$ of
  Notation~\ref{4.24:not-fibres-pairings}, with $u$ valued in $G$ or in $G^{\mathbb{C}}$, we have
  $\mathcal{U}_u^{T}=\mathcal{U}_u^{-1}$ and
  \begin{equation*}
    \bigl(\mathcal{U}_uK\mathcal{U}_u^{-1}\bigr)^{T}=\mathcal{U}_uK^{T}\mathcal{U}_u^{-1}.
  \end{equation*}
  Consequently
  $\mathcal{S}(\mathcal{U}_uK\mathcal{U}_u^{-1})=\mathcal{U}_u(\mathcal{S}K)\mathcal{U}_u^{-1}$.
\item[(d)] \emph{Kernel bound.} Let $d$ be the distance between bonds, and let $B_{\mathcal{A}}$ and
  $\delta$ be constants. Suppose $|K(b,b')|\le B_{\mathcal{A}}e^{-\delta d(b,b')}$ for all $b,b'$.
  Then $K^{T}$ and $\mathcal{S}K$ satisfy the same bound, with the same constants.
\item[(e)] \emph{Vanishing pattern.} Let $P$ be a set of pairs of bonds that is invariant under
  $(b,b')\mapsto(b',b)$. If $K(b,b')=0$ for all $(b,b')\in P$, then $(\mathcal{S}K)(b,b')=0$ for all
  $(b,b')\in P$. If a family $K(\sigma')$ satisfies $K(0)=0$, then $(\mathcal{S}K)(0)=0$.
\item[(f)] \emph{Coercivity of the real part.} Let $a\in\mathbb{R}$. If
  $\operatorname{Re}\langle\varphi,K\varphi\rangle\ge a\|\varphi\|^{2}$ for all $\varphi\in\mathcal{H}$,
  then $\operatorname{Re}\langle\varphi,K^{T}\varphi\rangle\ge a\|\varphi\|^{2}$ and
  $\operatorname{Re}\langle\varphi,\mathcal{S}K\varphi\rangle\ge a\|\varphi\|^{2}$ for all $\varphi$.
\item[(g)] \emph{Reality and norms.} If $K$ is real, that is, all entries map $\mathfrak{g}$ into
  $\mathfrak{g}$, then $K^{T}$ and $\mathcal{S}K$ are real. Moreover $\|K^{T}\|=\|K\|$ and
  $\|\mathcal{S}K\|\le\|K\|$. For a family $K(z)$,
  \begin{equation*}
    \|\mathcal{S}K(z)-\mathcal{S}K(z')\|\le\|K(z)-K(z')\|.
  \end{equation*}
\item[(h)] \emph{Blocks.} With $K_{ss}:=(K(b,b'))_{b,b'\in s}$ the block of $K$ on the survivors'
  bonds $s\subset\mathfrak{Bo}$, we have $(\mathcal{S}K)_{ss}=\mathcal{S}(K_{ss})$.
\end{enumerate}
Passing from $K$ to $\mathcal{S}K$ changes the determinant in general. Let $t$ be real with $t\neq0$,
and let
\begin{equation*}
  B_t:=\begin{pmatrix}1&t\\-t&1\end{pmatrix}
\end{equation*}
act on $\mathbb{C}^{2}$ with the standard bilinear and Hermitian pairings. Then
$\operatorname{Re}\langle\varphi,B_t\varphi\rangle=\|\varphi\|^{2}$ and $\mathcal{S}B_t=1$ (the identity
of $\mathbb{C}^{2}$), whereas $\det B_t=1+t^{2}\neq1$.
\end{lemma}

\begin{proof}
We first unwind the definition of $K^{T}$. For $\varphi,\psi\in\mathcal{H}$,
\begin{equation}\label{4.24:eq-kernel-polarisation-1}
  (\varphi,K^{T}\psi)=\sum_{b,b'}\bigl(\varphi(b),K(b',b)^{t}\psi(b')\bigr)
  =\sum_{b,b'}\bigl(K(b',b)\varphi(b),\psi(b')\bigr)=(K\varphi,\psi).
\end{equation}
The last equality holds because $\sum_{b}K(b',b)\varphi(b)=(K\varphi)(b')$. Since $(\cdot,\cdot)$ is
non-degenerate on the finite-dimensional space $\mathcal{H}$, \eqref{4.24:eq-kernel-polarisation-1}
characterises $K^{T}$. It follows that $(K^{T})^{T}=K$ and that $(K_1K_2)^{T}=K_2^{T}K_1^{T}$, since
$(K_1K_2\varphi,\psi)=(K_2\varphi,K_1^{T}\psi)=(\varphi,K_2^{T}K_1^{T}\psi)$.

(a) Set $\psi=\varphi$ in \eqref{4.24:eq-kernel-polarisation-1}. By the symmetry of $(\cdot,\cdot)$,
$(\varphi,K^{T}\varphi)=(K\varphi,\varphi)=(\varphi,K\varphi)$. Hence
\begin{equation*}
  (\varphi,\mathcal{S}K\varphi)=\tfrac12(\varphi,K\varphi)+\tfrac12(\varphi,K^{T}\varphi)
  =(\varphi,K\varphi)
\end{equation*}
for all $\varphi$, and in particular for real $\psi$. The two Gaussian weights are therefore the same
function of $\psi$, and they have the same integral.

(b) We have $(\mathcal{S}K)(b,b')=\frac12\bigl(K(b,b')+K(b',b)^{t}\bigr)$. The map $A\mapsto A^{t}$ is
$\mathbb{C}$-linear, because $(\cdot,\cdot)$ is $\mathbb{C}$-bilinear. So each entry of
$\mathcal{S}K$ is a $\mathbb{C}$-linear function of the entries of $K$, and holomorphic dependence on
$z$ is preserved.

(c) Let $u\in G$. The $\operatorname{Ad}$-invariance of $(\cdot,\cdot)$ gives
$(\operatorname{Ad}_uX,\operatorname{Ad}_uY)=(X,Y)$, that is,
$\operatorname{Ad}_u^{t}=\operatorname{Ad}_u^{-1}$. For complex elements we argue through the Lie
algebra. Differentiating the invariance shows that $\operatorname{ad}_X$ is $(\cdot,\cdot)$-skew for
$X\in\mathfrak{g}$. By $\mathbb{C}$-bilinearity, $\operatorname{ad}_Z$ is then $(\cdot,\cdot)$-skew for
every $Z\in\mathfrak{g}^{\mathbb{C}}$. Hence
$\exp(\operatorname{ad}_Z)^{t}=\exp(-\operatorname{ad}_Z)=\exp(\operatorname{ad}_Z)^{-1}$. The group
$G^{\mathbb{C}}$ is connected, so every $\operatorname{Ad}_u$ with $u\in G^{\mathbb{C}}$ is a finite
product of such exponentials. A product of operators whose transpose is their inverse again has this
property, so $\operatorname{Ad}_u^{t}=\operatorname{Ad}_u^{-1}$ for $u\in G^{\mathbb{C}}$ as well.

The conjugation $\mathcal{U}_u$ is diagonal in the bonds, with entries $\operatorname{Ad}$ of elements
of $G$, respectively of $G^{\mathbb{C}}$. Its transpose kernel is therefore diagonal with the
transposed entries, so $\mathcal{U}_u^{T}=\mathcal{U}_u^{-1}$. By the product rule for transposes and
$(\mathcal{U}_u^{-1})^{T}=(\mathcal{U}_u^{T})^{-1}=\mathcal{U}_u$, we get
\begin{equation*}
  (\mathcal{U}_uK\mathcal{U}_u^{-1})^{T}=(\mathcal{U}_u^{-1})^{T}K^{T}\mathcal{U}_u^{T}
  =\mathcal{U}_uK^{T}\mathcal{U}_u^{-1}.
\end{equation*}
Averaging this identity with $\mathcal{U}_uK\mathcal{U}_u^{-1}$ gives the identity for $\mathcal{S}$.

(d) Let $A^{*}$ be the $\langle\cdot,\cdot\rangle$-adjoint of $A$, and write $X\mapsto\bar X$ for the
conjugation. For $X,Y\in\mathfrak{g}^{\mathbb{C}}$,
\begin{equation*}
  \langle X,AY\rangle=(\bar X,AY)=(A^{t}\bar X,Y)=\bigl\langle\overline{A^{t}\bar X},Y\bigr\rangle .
\end{equation*}
Hence $A^{*}=\mathrm{conj}\circ A^{t}\circ\mathrm{conj}$. The conjugation is a
$\langle\cdot,\cdot\rangle$-isometry, since
\begin{equation*}
  \langle\bar X,\bar X\rangle=(X,\bar X)=(\bar X,X)=\langle X,X\rangle .
\end{equation*}
Summing over bonds, the fibrewise conjugation $\varphi\mapsto\bar\varphi$ is a
$\langle\cdot,\cdot\rangle$-isometry of $\mathcal{H}$, so $\|\bar\varphi\|=\|\varphi\|$.
Therefore $|A^{t}|=|A^{*}|=|A|$. Consequently
\begin{equation*}
  |K^{T}(b,b')|=|K(b',b)^{t}|=|K(b',b)|\le B_{\mathcal{A}}e^{-\delta d(b',b)}
  =B_{\mathcal{A}}e^{-\delta d(b,b')},
\end{equation*}
where the last step uses the symmetry of $d$. The bound for $\mathcal{S}K$ follows from the triangle
inequality applied to $\frac12K(b,b')+\frac12K^{T}(b,b')$.

(e) Let $(b,b')\in P$. Then $(b',b)\in P$, so $K^{T}(b,b')=K(b',b)^{t}=0$, and also $K(b,b')=0$.
Hence $(\mathcal{S}K)(b,b')=0$. If $K(0)=0$, then $(\mathcal{S}K)(0)=\mathcal{S}(0)=0$ by linearity.

(f) By \eqref{4.24:eq-kernel-polarisation-1} and the symmetry of $(\cdot,\cdot)$,
\begin{equation*}
  \langle\varphi,K^{T}\varphi\rangle=(\bar\varphi,K^{T}\varphi)=(K\bar\varphi,\varphi)
  =(\varphi,K\bar\varphi)=\langle\bar\varphi,K\bar\varphi\rangle .
\end{equation*}
The hypothesis holds for all vectors, so we may apply it at $\bar\varphi$. Together with
$\|\bar\varphi\|=\|\varphi\|$, shown in (d), this gives
\begin{equation*}
  \operatorname{Re}\langle\varphi,K^{T}\varphi\rangle
  =\operatorname{Re}\langle\bar\varphi,K\bar\varphi\rangle\ge a\|\bar\varphi\|^{2}=a\|\varphi\|^{2}.
\end{equation*}
Averaging with $\operatorname{Re}\langle\varphi,K\varphi\rangle\ge a\|\varphi\|^{2}$ gives
$\operatorname{Re}\langle\varphi,\mathcal{S}K\varphi\rangle\ge a\|\varphi\|^{2}$.

(g) Suppose $A$ maps $\mathfrak{g}$ into $\mathfrak{g}$. The form $(\cdot,\cdot)$ is real on
$\mathfrak{g}$, so $A|_{\mathfrak{g}}$ has a transpose $T$ on $\mathfrak{g}$. Its $\mathbb{C}$-linear
extension satisfies $(AX,Y)=(X,TY)$ on $\mathfrak{g}^{\mathbb{C}}$ by bilinearity, and hence equals
$A^{t}$. So $A^{t}$ maps $\mathfrak{g}$ into $\mathfrak{g}$. Applied to each entry, this shows that
$K^{T}$ is real when $K$ is, and then so is $\mathcal{S}K$.

For the norms, we repeat the computation of (d) on $\mathcal{H}$, using
\eqref{4.24:eq-kernel-polarisation-1} and the fibrewise conjugation. It gives
$K^{*}=\mathrm{conj}\circ K^{T}\circ\mathrm{conj}$, and the fibrewise conjugation is a
$\langle\cdot,\cdot\rangle$-isometry of $\mathcal{H}$ by (d). Hence $\|K^{T}\|=\|K^{*}\|=\|K\|$, and
\begin{equation*}
  \|\mathcal{S}K\|\le\tfrac12\|K\|+\tfrac12\|K^{T}\|=\|K\|.
\end{equation*}
By linearity (b), $\mathcal{S}K(z)-\mathcal{S}K(z')=\mathcal{S}\bigl(K(z)-K(z')\bigr)$. The previous
bound applied to $K(z)-K(z')$ gives the Lipschitz estimate.

(h) For $b,b'\in s$,
\begin{equation*}
  (\mathcal{S}K)_{ss}(b,b')=\tfrac12\bigl(K(b,b')+K(b',b)^{t}\bigr)=\mathcal{S}(K_{ss})(b,b'),
\end{equation*}
since the transpose of the block $K_{ss}$ has entries $K(b',b)^{t}$, $b,b'\in s$.

Finally, the matrix $B_t$. We have $B_t^{T}=B_{-t}$, so $\mathcal{S}B_t=1$, and
$\det B_t=1+t^{2}$. For $\varphi=(\varphi_1,\varphi_2)\in\mathbb{C}^{2}$,
\begin{equation*}
  \langle\varphi,B_t\varphi\rangle
  =\|\varphi\|^{2}+t\bigl(\bar\varphi_1\varphi_2-\bar\varphi_2\varphi_1\bigr).
\end{equation*}
The bracket equals $\bar\varphi_1\varphi_2-\overline{\bar\varphi_1\varphi_2}$, which is purely
imaginary. Hence $\operatorname{Re}\langle\varphi,B_t\varphi\rangle=\|\varphi\|^{2}$.
\end{proof}

The lemma bears on Gaussian integrals in the following way. Let $B\in\operatorname{End}(\mathcal{H})$
be non-symmetric with positive Hermitian part. By (a),
\begin{equation*}
  \int e^{-\frac12(\psi,B\psi)}\,d\psi=\int e^{-\frac12(\psi,\mathcal{S}B\psi)}\,d\psi,
\end{equation*}
and this integral is governed by $\det(\mathcal{S}B)$, which differs from $\det B$ in general, as the
matrix $B_t$ shows. For this reason the polarisation is taken before any determinant is formed.

At real arguments the stage form kernel $K_{\sharp}$ is real and self-adjoint, so the only question is
how it is continued to complex arguments. Setting $\mathcal{A}:=\mathcal{S}K_{\sharp}$ makes this
continuation the one whose determinant governs the Gaussian integral of (a), by the preceding
paragraph. Every property of the kinds treated in items (b)--(h) that is asserted of $K_{\sharp}$ in
this chapter, namely holomorphy, gauge covariance, an exponential kernel bound, a symmetric vanishing
pattern, coercivity of the real part, reality, and norm and Lipschitz bounds, holds of $\mathcal{A}$
with the same constants, by those items.

\begin{definition}[The stage quadratic form and its blocks]\label{4.24:def-stage-form}
We work at one stage of the construction, with the components, the survivors and the integration
bonds of Definition~\ref{4.24:def-survivors-bonds}, so that
$\mathfrak{Bo}=s\sqcup\mathcal{W}$, and with the fibres, pairings, kernels, coordinate projections
$P_E$, blocks $K_{EE'}=P_EKP_{E'}$, transposes $K^T$, the domain
$\mathfrak{D}^{\flat}_{n+1}(X)$ and its real points as in
Notation~\ref{4.24:not-fibres-pairings}. The complex background pair is
$(\mathbb{U},\mathbb{J})\in\mathfrak{K}_P$, where $\mathfrak{K}_P$ is the space of pairs over
$\mathbb{B}^{(n)}_k$ with radius $\gamma_0$ fixed in Notation~\ref{4.24:not-fibres-pairings}.

\emph{Block pairs and the elimination map.} A \emph{block pair} $c$ consists of two face-adjacent
level-$(j+1)$ blocks $B(c_-)$, $B(c_+)$, and $b_0(c)$ is the bond eliminated on it. We write $C$
for the linear map of
\cite{B-CMP119} that, in the notation there ($A'(b)$ the field variable on the bond $b$), determines
the variable $A'(b_0(c))$ as a linear function of the remaining
variables $A'(b)$, $b\subset B(c_-)\cup B(c_+)$. Thus $(C\mathsf{B})(b)=\mathsf{B}(b)$ on starred
bonds, and $(C\mathsf{B})(b_0(c))$ is linear in the restriction of $\mathsf{B}$ to
$B(c_-)\cup B(c_+)$. On complex pairs $C=C(\mathbb{U})$ is the holomorphic continuation of
Ba{\l}aban's real map, and we put $C^{*}:=C^{T}$, the $(\cdot,\cdot)$-transpose, which is the
holomorphic continuation of the real adjoint.

\emph{The collar kernel.} $\mathsf{B}$ denotes the collar variable on the starred collar bonds, and
$\varrho^{B}_j$ is its radius. $\Delta^{(j)}=\Delta^{(j)}(\sigma')$ is the decoupled operator of
\cite[(3.156)]{B-CMP99b} on the pair, where $\sigma'$ are the decoupling parameters on the host
cubes of $\mathbb{B}^{(n)}_k$; $a$ is the constant of that display. Ba{\l}aban defines the operator
through its quadratic form,
\begin{equation}\label{4.24:eq-stage-form-1}
\langle B,(QG_1Q^{*})^{-1}B\rangle-a\langle B,B\rangle-2\langle H_1D^{(2)}(B),J\rangle
=\langle B,\Delta_kB\rangle ,
\end{equation}
with $Q$, $G_1$, $H_1$, $D^{(2)}$ as in \cite{B-CMP99b}. Here $B$ is the field variable and $J$
the current of that display, and $\Delta^{(j)}(\sigma')$ is the operator $\Delta_k$ so defined,
taken at level $j$ and decoupled by the parameters $\sigma'$. We put
\begin{equation}\label{4.24:eq-stage-form-2}
K(\sigma'):=C^{*}\Delta^{(j)}(\sigma')\,C=\mathcal{A}^{w}(\sigma')-a\,C^{*}C ,
\end{equation}
which defines $\mathcal{A}^{w}(\sigma')$. Below, $\delta>0$ denotes the decay rate of the collar
kernel in hypothesis~(a) and $d(b,b')$ the
distance of bonds. We also use $n_b:=d\,(LM)^{d}$, with $d=4$ the dimension, the constant $n_C$
attached to $C$, the operator norm $\|C\|$ of $C$, and
\begin{equation*}
C_\rho:=\sup_b\sum_{b'}e^{-\delta d(b,b')} .
\end{equation*}
The constants $c_J$, $c_\flat$ are those of the element weights $w(e)$, and $\kappa_1$ is
Ba{\l}aban's auxiliary parameter.

\emph{Hypotheses.} The assertions below are made under the following two hypotheses.
\begin{enumerate}
\item[(a)] \emph{The collar form.} Hypothesis~(a) is the following statement on the decoupled
collar kernel $\mathcal{A}^{w}(\sigma')$, taken here by statement. There are
constants $B_{\mathcal{A}}$ and $\delta>0$ with the following
properties. $\mathcal{A}^{w}(\sigma')$ is holomorphic in $(\mathbb{U},\mathbb{J},\sigma')$ on
$\mathfrak{K}_P\times\{|\sigma'|\le e^{\kappa_1}\}$ and covariant under the gauge action of
Notation~\ref{4.24:not-fibres-pairings}. For every ordered pair $(b,b')$ and every non-empty
wall-connected family $Y'$ of host cubes holding $b$ and $b'$ we put
\begin{equation*}
v_{e_a^0(b,b')}(\mathsf{B}):=\tfrac12 a\,\langle\mathsf{B}(b),(C^{*}C)(b,b')\mathsf{B}(b')\rangle ,
\end{equation*}
\begin{equation*}
v_{e_a(b,b',Y')}(\mathsf{B}):=\prod_{\Delta\in Y'}\int_0^1\! d\sigma'(\Delta)\,
\partial_{\sigma'(\Delta)}
\Bigl[-\tfrac12\langle\mathsf{B}(b),\mathcal{A}^{w}(\sigma')(b,b')\mathsf{B}(b')\rangle
\Bigr]_{\sigma'=0\text{ off }Y'} ,
\end{equation*}
which is the fundamental theorem of calculus in $\sigma'$ applied to the entry. With these members,
\begin{equation}\label{4.24:eq-stage-form-a}
\begin{aligned}
&\text{(i)}\quad \mathcal{A}^{w}(0)=0,\qquad
|\mathcal{A}^{w}(\sigma')(b,b')|\le B_{\mathcal{A}}e^{-\delta d(b,b')};\\
&\text{(ii)}\quad \text{for every set }Y'\text{ of host cubes: }\sigma'=0\text{ off }Y'\
\Longrightarrow\\
&\qquad\quad\mathcal{A}^{w}(\sigma')(b,b')=0\text{ unless }b,b'\text{ lie in one component of }
Y';\\
&\text{(iii)}\quad -\tfrac12\langle\mathsf{B}(b),K(\mathbf{1})(b,b')\mathsf{B}(b')\rangle
=v_{e_a^0(b,b')}(\mathsf{B})+\sum_{Y'}v_{e_a(b,b',Y')}(\mathsf{B});\\
&\nu^{\natural}(e_a):=\tfrac12B_{\mathcal{A}}(\varrho^{B}_j)^2e^{-\delta d(b,b')}
e^{-(\kappa_1-1)|Y'|},\qquad
\nu^{\natural}(e_a^0):=\tfrac12a\|C\|^2(\varrho^{B}_j)^2;\\
&K_a:=K_a^{(1)}+K_a^{(2)},\qquad
K_a^{(1)}:=2e^{c_J+2c_\flat+2}\,n_b\,C_\rho\,B_{\mathcal{A}},\\
&\phantom{K_a:=K_a^{(1)}+K_a^{(2)},\qquad}
K_a^{(2)}:=2e^{3c_J+2c_\flat+2}\,n_b\,n_C\,a\|C\|^2 .
\end{aligned}
\end{equation}
In (ii) the entry is read on that component only; in (iii) $K(\mathbf{1})$ is $K(\sigma')$ at
$\sigma'\equiv1$, and $Y'$ runs over the non-empty wall-connected families of host cubes holding $b$
and $b'$. The members carry the cube sets
$\mathsf{b}(e_a):=\{\operatorname{cube}(b),\operatorname{cube}(b')\}$, with $\operatorname{cube}(b)$
the cell of $\mathsf{Q}$, the set of $\pi_{j+1}$-cells meeting $\mathfrak{Bo}$, that contains $b$,
so $|\mathsf{b}(e_a)|\le2$, and the
domain $\mathsf{d}(e_a):=Y'$. Their weights satisfy $w(e_a)\le c_J|Y'|$, the cubes of $Y'$ being
counted as objects of weight zero, and $w(e_a^0)\le2c_J$.
\item[(b)] \emph{Symmetry of the decoupled operator.} At every point of
$\mathfrak{D}^{\flat}_{n+1}(X)$, $\Delta^{(j)}$ is the $(\cdot,\cdot)$-symmetric operator of the
quadratic form by which \cite[(3.156)]{B-CMP99b} defines it, see \eqref{4.24:eq-stage-form-1}; that
is, $(\Delta^{(j)})^{T}=\Delta^{(j)}$ on the whole domain. Hypothesis~(b) is the identification on
which the normalisation of the stage Gaussian rests; it is assumed here and is not derived.
\end{enumerate}

\emph{Elements.} We put $\mathfrak{E}^{+}:=\mathfrak{E}\cup\{e_a\}\cup\{e_a^0\}$, so that
$\mathcal{A}ct^{P}=\sum_{e\in\mathfrak{E}^{+}}v_e$. Here $\mathfrak{E}$ consists of the elements
inherited from the large-field operation, together with the pieces of the prepared action, among
them the pieces $\mathbb{P}^{(j)}$. It is assumed throughout, as the premise of the cluster
expansion of this section under the conditions imposed on the arrested components of
Definition~\ref{4.24:def-survivors-bonds}, that each $v_e$ is
holomorphic in the collar variable on
$|\mathsf{B}|<\varrho^{B}_j$, jointly with $(\mathbb{U},\mathbb{J})$, the weak slots and $g$, and
that $|v_e|\le\nu^{\natural}(e)$ there. This assumption is not used in assertions (1)--(4) below.

\emph{The stage form.} $K_\sharp$ is the kernel on $\mathcal{H}$ of the quadratic form in the
exponent of the stage Gaussian, the Gaussian normalised at this stage, whose weight in the rim and
collar variables is displayed below. Here $A_j$ denotes the
variables on the rim $\mathsf{R}_j$, $B_j=\mathsf{B}$ the
collar variable, and $\Delta^{(j)}_{\Lambda_{j+1}}$ the Dirichlet restriction of $\Delta^{(j)}$ to
the bonds of $\Lambda_{j+1}$, the region of level $j+1$ carrying the rim variables. That Gaussian is
\begin{equation*}
\exp\Bigl(-\tfrac12\langle A_j,C^{*}\Delta^{(j)}_{\Lambda_{j+1}}CA_j\rangle
-\tfrac12\langle B_j,C^{*}\Delta^{(j)}CB_j\rangle-(\text{cross terms})\Bigr),
\end{equation*}
with complex covariance $\mathcal{C}:=(C^{*}\Delta^{(j)}C)^{-1}$ on the joint $(A_j,B_j)$-space
over the union of the collar and the $A_j$-region. This $\mathcal{C}$ is itself a conditional
covariance on a subdomain.
\begin{itemize}
\item The collar$\times$collar block of $K_\sharp$ is the kernel $K(\mathbf{1})$ of
hypothesis~(a). The members $e_a$, $e_a^0$ with both bonds off the arrested components make up the
survivors' Gaussian, which the arrest leaves unchanged.
\item The rim block of $K_\sharp$ is the form of \cite{B-CMP119}, which couples only the fields
$A_j$ in the same component. There the other terms of this quadratic form are included into the
effective action as $\mathbf{B}$-terms. This form is carried to the integral in
\cite[(1.60)]{B-CMP122a}, where it is stated that ``The integral in (1.60) is of the same type as
the integral studied in Sect.~3 [III]'', the reference [III] there being \cite{B-CMP119}.
\item The remnants of the rim block that couple distinct components are explicit quadratic
expressions bounded by $e^{-R_j}$. They are treated as $\mathbf{B}$-terms of the large-field
operation, with data in the weak slot of the $A$-variables and with no strong sets, by part (b) of
the proposition of this chapter on the $\mathbf{B}$-terms of the large-field operation in the
$A$-variables.
\item The cross terms in the display above are the $A_j\times B_j$ off-diagonal part of the joint
operator. The covariance $\mathcal{C}$ above is written as one operator on the joint space over the
union of the collar and the
$A_j$-region and refers to no component. \cite[(1.60)]{B-CMP122a} writes the two Gaussians as
separate factors and places the coupling between $A_j$ and $B_j$ in the pieces $\mathbb{P}^{(j)}$.
\end{itemize}

\emph{Two readings of the cross terms.} The notation above admits two readings of the cross terms
between distinct components. They are kept side by side, and neither is chosen here:
\begin{equation}\label{4.24:eq-stage-form-b}
\begin{aligned}
&\text{(I)}\quad K_a\ \text{unchanged};\\
&\text{(II)}\quad r_{\mathrm{rim}}:=\max\{\mathfrak{b}_j,\varrho_j,\varrho^{B}_j\}
=\max\{\varrho_j,\varrho^{B}_j\},\qquad
K_a^{\sharp}:=K_a\,(r_{\mathrm{rim}}/\varrho^{B}_j)^2 .
\end{aligned}
\end{equation}
\begin{enumerate}
\item[(I)] \emph{Per-component reading}, that of \cite[(1.60)]{B-CMP122a} for the two Gaussians
and of \cite{B-CMP119} for the rim. The $A_j\times B_j$ terms between distinct components are
pieces $\mathbb{P}^{(j)}$ or $\mathbf{B}$-terms. They are therefore elements of $\mathfrak{E}$, of
the two kinds of cross elements whose majorants are carried in the slot $\mathfrak{p}_j$ of
$(Q^{+})'$ and never in $K_a$, and nothing below changes.
\item[(II)] \emph{Joint reading}: the display above read literally as one operator on the joint
space.
\begin{itemize}
\item The rim$\times$collar and collar$\times$rim entries between distinct components are then
cross members with one end on the rim, of the type in \eqref{4.24:eq-stage-form-a}(iii). The
third kind of cross element, reserved for them among the kinds of elements of this section, then
occurs.
\item The conclusions \eqref{4.24:eq-stage-form-a} of hypothesis~(a), read on
$\mathfrak{Bo}=\text{rim}\cup\text{collar}$, become an additional input. It is not derived from
hypothesis~(a) as stated, which concerns the collar bonds.
\item The radius used on the rim must cover the exterior real box $|A_j|<\mathfrak{b}_j$. This
gives $r_{\mathrm{rim}}$ as in \eqref{4.24:eq-stage-form-b}, where $\varrho_j$ is the radius of the
variables $A_j$ in the large-field operation; the second equality holds because
$\varrho_j\ge2\mathfrak{b}_j$ by the conditions on the radii of that operation.
\item $K_a$ is replaced by $K_a^{\sharp}$, with $C_\rho$ read on $\mathfrak{Bo}$, and
$\vartheta_a$, $\vartheta_{\mathrm{cv}}$ by the corresponding quantities $\vartheta_a^{\sharp}$,
$\vartheta_{\mathrm{cv}}^{\sharp}$.
\item The smallness condition on the change of variables and the slot $\mathfrak{p}_j$ of
$(Q^{+})'$ are increased by the factor $(r_{\mathrm{rim}}/\varrho^{B}_j)^2$, a fixed power of
$\log g_j^{-2}$; in both, this factor multiplies the exponentially small quantity
$\varepsilon_Y\le2E(x_k)^6$.
\item The constants involved are fixed before $\gamma$. The increase imposes no cap in the sense of
Definition~\ref{2.3:def-cap}. The constant $C_0^{\sharp}$ and the coefficient $5$ in the bound for
members joining an arrested component to another component are unchanged.
\end{itemize}
\end{enumerate}
In either reading each such term enters exactly once. It is immaterial for clause (i) of the
determinant lemma for the survivors, which forms $\det\mathcal{A}_{ss}$ only. None of the
statements on which that clause depends refers to the type of an entry, and every region used there
is a union of the rims and collars of whole components. The assertions below hold in both readings.

\emph{The operator of the form.} We put
\begin{equation}\label{4.24:eq-stage-form-3}
\mathcal{A}:=\mathcal{S}K_\sharp=\tfrac12\bigl(K_\sharp+K_\sharp^{T}\bigr),
\end{equation}
the operator of the quadratic form, and its blocks
\begin{equation*}
\mathcal{A}_{ss}:=P_s\mathcal{A}P_s\ \text{on }\mathcal{H}_s,\qquad
\mathcal{A}_{s\mathcal{W}},\qquad\mathcal{A}_{\mathcal{W}s},\qquad
\mathcal{A}_{\mathcal{W}\mathcal{W}} ,
\end{equation*}
the last three in the notation $K_{EE'}$. We set $\mathcal{C}_s:=\mathcal{A}_{ss}^{-1}$, which
exists by the invertibility of $\mathcal{A}_{ss}$ established in this section. For an arrested
component $W$ we set
\begin{equation*}
\mathsf{q}_W(x):=-\tfrac12\langle x,\mathcal{A}_{WW}x\rangle ,
\end{equation*}
where $x$ is the new collar variable and $\mathcal{A}_{WW}$ the block of $\mathcal{A}$ on the
bonds of $W$ on which $x$ lives.

\emph{Assertions.} Under hypotheses (a) and (b):
\begin{enumerate}
\item[(1)] with $C_s:=P_sCP_s$ and $C_{\mathcal{W}}:=P_{\mathcal{W}}CP_{\mathcal{W}}$, one has
$C=C_s\oplus C_{\mathcal{W}}$ and
$C^{*}C=C_s^{*}C_s\oplus C_{\mathcal{W}}^{*}C_{\mathcal{W}}$;
\item[(2)] $(\phi,K_\sharp\phi)=(\phi,\mathcal{A}\phi)$ for all $\phi\in\mathcal{H}$;
\item[(3)] $\mathcal{A}=K_\sharp$ at the real points;
\item[(4)] $\mathcal{A}=K_\sharp$ on the whole of $\mathfrak{D}^{\flat}_{n+1}(X)$.
\end{enumerate}
\end{definition}

\begin{proof}
\emph{Assertion (1).} By the stencil of $C$, $C(b'',b)\neq0$ only if $b''$ and $b$ lie in one block
pair. Indeed, if $b''$ is starred then $C(b'',b)$ is the identity for $b=b''$ and vanishes
otherwise. If $b''=b_0(c)$, the row is supported on bonds of $B(c_-)\cup B(c_+)$. Now
\begin{equation*}
(C^{*}C)(b,b')=\sum_{b''}C(b'',b)^{t}\,C(b'',b') .
\end{equation*}
A summand is non-zero only if $b$ and $b'$ both lie in the support of the row $b''$. Hence
$(C^{*}C)(b,b')=0$ unless $b$ and $b'$ lie in one block pair.

A block pair consists of two face-adjacent level-$(j+1)$ blocks. These are cells of partitions
compatible with the $\pi_{j+1}$-cells, so a block pair lies in at most two wall-adjacent
$\pi_{j+1}$-cells. By \eqref{4.24:eq-component-separation-c} of
Lemma~\ref{4.24:lem-component-separation}, which is proved as an implication from the statements
named there, all integration bonds lying in such a pair of cells belong to one component. Hence all
bonds of a block pair belong to one component.

By Definition~\ref{4.24:def-survivors-bonds}, $s$ and $\mathcal{W}$ are unions of the bonds of
whole components (the survivors, respectively the arrested ones). Therefore $C$ maps
$\mathcal{H}_s$ into $\mathcal{H}_s$ and $\mathcal{H}_{\mathcal{W}}$ into
$\mathcal{H}_{\mathcal{W}}$, that is, $P_sC=CP_s$ and $P_{\mathcal{W}}C=CP_{\mathcal{W}}$. With
$C_s:=P_sCP_s$ and $C_{\mathcal{W}}:=P_{\mathcal{W}}CP_{\mathcal{W}}$ this gives
$C=C_s\oplus C_{\mathcal{W}}$. The projections are $(\cdot,\cdot)$-symmetric, so taking transposes
gives $C^{*}P_s=P_sC^{*}$ and $C^{*}P_{\mathcal{W}}=P_{\mathcal{W}}C^{*}$. Since
$P_s+P_{\mathcal{W}}=\mathbf{1}$, it follows that
$C^{*}C=C_s^{*}C_s\oplus C_{\mathcal{W}}^{*}C_{\mathcal{W}}$.

\emph{Assertion (2).} This is \eqref{4.24:eq-kernel-polarisation-a}(a) of
Lemma~\ref{4.24:lem-kernel-polarisation}, applied to $K=K_\sharp$: for all $\phi$,
$(\phi,\mathcal{S}K_\sharp\phi)=(\phi,K_\sharp\phi)$. In particular, for real integration variables
the Gaussian weight and its integral are the same whether written with $K_\sharp$ or with
$\mathcal{A}$. Thus $\mathcal{A}$ is Ba{\l}aban's own object: $\Delta_k$ is defined by its
quadratic form \eqref{4.24:eq-stage-form-1}, and the logarithm of the determinant entering the
normalisation is that of the Gaussian's quadratic form.

\emph{Assertion (3).} At a real point every kernel maps $\mathfrak{g}$ to $\mathfrak{g}$, by
Notation~\ref{4.24:not-fibres-pairings}. On real kernels the $(\cdot,\cdot)$-transpose is the
adjoint for the real inner product. $K_\sharp$ is real and self-adjoint there, so
$K_\sharp^{T}=K_\sharp$, and \eqref{4.24:eq-stage-form-3} gives $\mathcal{A}=K_\sharp$.

\emph{Assertion (4).} We show $K_\sharp^{T}=K_\sharp$ on the whole of
$\mathfrak{D}^{\flat}_{n+1}(X)$. By hypothesis~(b), $\Delta^{(j)}$ is $(\cdot,\cdot)$-symmetric
there. $K_\sharp$ is obtained from $\Delta^{(j)}$ by the operations below, each of which preserves
$(\cdot,\cdot)$-symmetry.
\begin{itemize}
\item The elimination $C^{*}(\cdot)C$, with $C^{*}=C^{T}$, is a congruence:
$(C^{T}\Delta C)^{T}=C^{T}\Delta^{T}C=C^{T}\Delta C$ whenever $\Delta^{T}=\Delta$.
\item The Dirichlet restriction to the bonds of a region $\Lambda'$, in particular to those of
$\Lambda_{j+1}$, is the congruence $\Delta\mapsto P\Delta P$ by the
$(\cdot,\cdot)$-symmetric coordinate projection $P$ onto the bonds of $\Lambda'$.
\item Gauge conjugation is a congruence. Let $\mathcal{U}_u=\bigoplus_b\operatorname{Ad}_{u(b_-)}$
be as in Notation~\ref{4.24:not-fibres-pairings}, now for $G^{\mathbb{C}}$-valued $u$. By
\eqref{4.24:eq-kernel-polarisation-a}(c) of
Lemma~\ref{4.24:lem-kernel-polarisation}, $\operatorname{Ad}_u^{t}=\operatorname{Ad}_u^{-1}$ also
for $u\in G^{\mathbb{C}}$, so $\mathcal{U}_u^{T}=\mathcal{U}_u^{-1}$. Hence the conjugation
$K\mapsto\mathcal{U}_uK\mathcal{U}_u^{-1}=\mathcal{U}_uK\,\mathcal{U}_u^{T}$ is a congruence.
\item Inverses and resolvents of $(\cdot,\cdot)$-symmetric operators are $(\cdot,\cdot)$-symmetric:
$(A^{-1})^{T}=(A^{T})^{-1}=A^{-1}$, and likewise $((A-z)^{-1})^{T}=(A-z)^{-1}$.
\end{itemize}
Hence $K_\sharp^{T}=K_\sharp$ on the whole domain, and $\mathcal{A}=K_\sharp$ there. The
symmetrisation $\mathcal{S}$ in \eqref{4.24:eq-stage-form-3} therefore acts as the identity on
$K_\sharp$; it is retained in the definition for safety.

The argument does not use the identity theorem. One could observe that
$K_\sharp-K_\sharp^{T}$ is holomorphic on the domain and, by assertion (3), vanishes at its real
points. Concluding from this would require that every connected component of the domain meet the
real locus in a relatively open subset of a maximal totally real submanifold, and that premise is
not available.

\emph{The two readings.} The reading in
\eqref{4.24:eq-stage-form-b} does not affect assertions (1)--(4). Assertion (1) concerns $C$ only.
Assertions (2)--(4) concern $K_\sharp$ as a whole, whichever family of elements its cross terms
between distinct components are assigned to.
\end{proof}

\begin{definition}[Constants of the Gaussian normalisation and parameter conditions]
\label{4.24:def-normalisation-constants}
The scales $s_m=L^{m-k}$, the cell partitions $\pi_m$, the domains $\mathbb{D}_m$, the tree length
$\tau$, the distances $d_m(X)=\tau(X)/(Ms_m)$, the hulls $[Y]_m$ and the filing level $m(X)$ are
those of Notation~\ref{4.24:not-scales-rim-collar}. The background structure
$\mathbb{B}^{(n)}_k$ and the space of pairs
$\mathfrak{K}_P=\mathfrak{K}_{109}(\mathbb{B}^{(n)}_k;\gamma_0)$ are those of
Notation~\ref{4.24:not-fibres-pairings}. We write $x_m=g_m^{-2}$ and
$\mathsf{t}_k=\log g_k^{-2}$.

\smallskip
\emph{Hypotheses.} The following statements are used as stated.
\begin{enumerate}
\item[(a)] The lemma of this chapter bounding the Gaussian normalisation density
$\mathbb{C}_G$. It provides a constant $c_G=c_G(d,L,M)$, of order $(LM)^4$, which does not
depend on $k$ and $n$ (in particular, the constant $c_G+1$ entering the same lemma is uniform
in $k$ and $n$). It is stated for rates in the admissible range $a_m\le 4\kappa L$.
\item[(b)] The lemma of this chapter on domains, junctions and animals. It has three
clauses.
\begin{enumerate}
\item[(i)] For levels $\ell\le m$ and every $X$,
\begin{equation}\label{4.24:eq-normalisation-constants-1}
d_m([X]_m)\le L^{-(m-\ell)}\,d_\ell(X).
\end{equation}
\item[(ii)] Each contact at which two domains are joined adds at most $\sqrt{d}\,Ms_m$ to
the tree length $\tau$. The polygon of a walk is measured in the sides of the walk's own
cubes.
\item[(iii)] For $\Delta\in\pi_m$ and $D\ge 0$,
\begin{equation}\label{4.24:eq-normalisation-constants-2}
\#\{X\in\mathbb{D}_m:\ X\ni\Delta,\ d_m(X)\le D\}\le e^{c_d(D+2)},
\end{equation}
where $c_d$ is the number defined in condition~(c1) below.
\end{enumerate}
\item[(c)] The assumption $\tfrac12(\kappa_1-1)\ge 2L\kappa$ made in \cite{B-CMP116}.
\item[(d)] The convergence floor on $\kappa_1$. It reads
\begin{equation}\label{4.24:eq-normalisation-constants-3}
\kappa_1-1\ge c_J+c_\flat+\mathsf{c}_{\mathrm{an}}+7,
\qquad c_J:=2\sqrt{d}\,(a_k+1).
\end{equation}
Here $a_k$ is the rate at the top level $k$ of the weighted norm in which that floor is
stated. The constant $c_\flat$ is the constant of the element weights of the cluster
expansion, and $\mathsf{c}_{\mathrm{an}}$ is the animal constant. The floor is a lower bound
on $\kappa_1$, imposed after $\kappa$ and before $M$; its right-hand side is fixed before
$\kappa_1$ (after $L$, $\kappa$ and the positive rate parameter $\beta$ of the construction,
which is subject to~(c1) and~(c4) below). In that norm the top rate is
$a_k=(1+3\beta)\kappa$, so that the number $c_J$ depends on $(\kappa,\beta)$ only; the
statements of hypothesis~(a) require only rates $a_m\le 4\kappa L$.
\item[(e)] The majorant of the radius of the space of pairs that depends on $\mathsf{t}_k$
alone and not on the couplings at earlier levels, stated in this
chapter and displayed in~\eqref{4.24:eq-normalisation-constants-5} below.
\item[(f)] The bound of this chapter on the normalisation density and its field dependence,
displayed in~\eqref{4.24:eq-normalisation-constants-6} below.
\item[(g)] The statement of this chapter on the radius of the space of pairs, which is
applied under the condition $K_\Delta\gamma_0\le\tfrac12$, with $K_\Delta$ the constant of the
bound of hypothesis~(f) described in~(n8) below.
\end{enumerate}

\smallskip
\emph{Normalisation objects.}
\begin{enumerate}
\item[(n1)] $z^{(j)}=\prod_c z^{(j)}(c)$. The product runs over the cells $c$ of the region
where the $\delta$-functions were eliminated, and $z^{(j)}(c)$ are the normalisation
factors of \cite{B-CMP119}, where the level index is written $k$, so that the same product
appears there as $\prod_c z^{(k)}(c)$.
On the survivors' region we put
\begin{equation*}
z^{(j)}_s:=\prod_{c\subset R^{(n+1)}\cap Z^{s}}z^{(j)}(c).
\end{equation*}
\item[(n2)] $E_0^{(j)}(\cdot)$ is the vacuum-energy counterterm of the prefactor in
\cite{B-CMP119}. In \cite{B-CMP119} it is a number which
contributes to the vacuum energy renormalisation. It is an additive set function of its
region.
\item[(n3)] $\lambda_0$ is the positive number of \cite{B-CMP119}, taken there with
$\lambda_0>r$ but small enough, $r$ denoting the parameter of \cite{B-CMP119} at that place.
The number $\lambda_0$ does not depend on the region.
\item[(n4)] $c_G':=(LM)^4\cdot O(1)$ is the number of densities per $\pi_{j+1}$-cell, and
$c_G=c_G(d,L,M)$ is the constant of hypothesis~(a).
\item[(n5)] The field-smallness parameter is
\begin{equation*}
\vartheta:=K_\vartheta\,\frac{A_1}{C_*}\Bigl(\tfrac12\log g_k^{-2}\Bigr)^{p_1-q_*}.
\end{equation*}
Here $A_1$ is the auxiliary parameter of that name in the tuple $\mathfrak{p}$
(Notation~\ref{2.1:not-prefix}); $K_\vartheta$ and $C_*$ are positive constants and $p_1$,
$q_*$ exponents of the construction, the two exponents being subject to the inequality
of~(c5) below.
\item[(n6)] Let $\mathsf{L}_0$ denote the auxiliary integer subject to
$\mathsf{L}_0^{\,2}\le L/3$ in~(c4) below, which enters $w_*$. Let $N_0(j)$ be
Ba{\l}aban's integer $N_0$ of \cite{B-CMP122a} at density $j$. For levels $m$ put
$\alpha_{i,m}:=C_i\,x_m^{-1/2}(\log x_m)^q$, where the index $i$ runs over the radii of
Ba{\l}aban's construction and $C_i$ is the constant of the $i$-th radius. Define
\begin{equation*}
w_*(j):=\mathsf{L}_0^{\,2N_0(j)},\qquad
\bar\alpha_j:=\max_i\max_{1\le m\le j}\alpha_{i,m},
\end{equation*}
so that $\bar\alpha_j$ is the maximum over the levels $m\le j$. The parameter $\gamma_0$ of
$\mathfrak{K}_P$ is taken at the current density $k$:
\begin{equation}\label{4.24:eq-normalisation-constants-4}
\gamma_0:=\gamma_0^{(k)}:=5\max_{j\le k}w_*(j)\,\bar\alpha_j .
\end{equation}
Thus $\gamma_0$ is five times the largest of the products of $w_*(j)$ with an amplitude built
from the couplings through $x_m^{-1/2}=g_m$. By hypothesis~(e), with $C_{\max}:=\max_iC_i$,
\begin{equation}\label{4.24:eq-normalisation-constants-5}
\gamma_0^{(k)}\le\bar\gamma_0(\mathsf{t}_k):=5(1+\beta_0)^2\Psi(\mathsf{t}_k),\qquad
\Psi(\mathsf{t}):=L^2(L+1)\,C_{\max}\,\mathsf{t}^{\,r+q}e^{-\mathsf{t}/2},
\end{equation}
with $\beta_0$ the constant of hypothesis~(e).
\item[(n7)] For a function $\mathbb{F}$ on domains, the filing norm at rates
$a_m\le 4\kappa L$ is
\begin{equation*}
\|\mathbb{F}\|:=\sup_m\sup_{\Delta\in\pi_m}\sum_{X\ni\Delta,\ m(X)=m}
\|\mathbb{F}(X)\|_\infty\,e^{a_m d_m(X)} ,
\end{equation*}
the sum running over the domains containing $\Delta$ that are filed at level $m$.
The stage norm at the rates $a^{(n+1)}_m$ is
\begin{equation*}
\|\mathbb{F}\|_{(n+1)}:=\sup_y\sum_{X\ni y}\|\mathbb{F}(X)\|_\infty\,
e^{a^{(n+1)}_{m(X)}\,d_{m(X)}(X)} ,
\end{equation*}
where $y$ runs over the points of the lattice and $a^{(n+1)}_m$ are the rates of
stage $n+1$.
\item[(n8)] Hypothesis~(f) states two bounds:
\begin{equation}\label{4.24:eq-normalisation-constants-6}
\begin{gathered}
\|\mathbb{C}_G\|_{(n+1)}\le 3c_G,\\
\text{(field dependence of }\mathbb{C}_G)\le c_G'\,O(1)\bigl[\,5w_*(BCM\tilde\alpha_0
+\tilde\alpha_1)+K_\Delta\gamma_0\,\bigr]\le\vartheta .
\end{gathered}
\end{equation}
In~\eqref{4.24:eq-normalisation-constants-6} the symbols $w_*$ (without level index),
$\tilde\alpha_0$, $\tilde\alpha_1$, $B$ and $C$ are those of the bound of hypothesis~(f):
$w_*$ is its weight $\mathsf{L}_0^{\,2N_0}$, with Ba{\l}aban's integer $N_0$ as in~(n6);
$\tilde\alpha_0$, $\tilde\alpha_1$ are smallness parameters of Ba{\l}aban's construction;
and $B$, $C$ are constants of Ba{\l}aban's construction entering the field-dependence bound.
The constant $K_\Delta$ is a number depending on $(d,L,\kappa_1)$ only. Besides
\eqref{4.24:eq-normalisation-constants-6}, the radius $\gamma_0$ enters the bounds of this
section only through the smallness parameter $\alpha_5$ of Ba{\l}aban's construction. We put
\begin{equation*}
C_0^{\sharp}:=3c_G+1 .
\end{equation*}
\end{enumerate}

\smallskip
\emph{Parameter conditions.} The parameters satisfy the following conditions.
\begin{enumerate}
\item[(c1)] $\beta\kappa\ge c_d+2\sqrt{d}\,(2+4\beta)$, where $c_d:=2^d\log(2de)+1$.
\item[(c2)] $(M/M_1)\cdot\min\{\delta_0,\delta\}\ge 8\kappa\sqrt{d}$. Here $\delta_0$ is the
random-walk rate of \cite{B-CMP99b}, and $\delta$ is the decay rate of the exponential
bound of this chapter under which these conditions are stated.
\item[(c3)] The restrictions of \cite{B-CMP116} hold, with $\varpi:=\beta/40$ in place of
the parameter written $\delta$ there. The parameter $\kappa_1$ is subject to the floor
\begin{equation}\label{4.24:eq-normalisation-constants-7}
\kappa_1\ge\max\{16\kappa+1,\;4L\kappa+1\}.
\end{equation}
The first entry, $16\kappa+1$, is the value substituted for $\kappa_1$ in these
restrictions, kept here only as a lower bound. The second entry is the assumption of
hypothesis~(c), since
$\tfrac12(\kappa_1-1)\ge 2L\kappa$ is equivalent to $\kappa_1\ge 4L\kappa+1$.
\item[(c4)] $L\ge 8$, $\mathsf{L}_0^{\,2}\le L/3$, $\beta\le\tfrac14$.
\item[(c5)] $q_*\ge p_1+1$ and $p_0\ge 5r$.
\item[(c6)] The coupling threshold $\gamma$ is so small that the smallness clauses on the
coupling that belong to these conditions hold for all $g_k\le\gamma$. This $\gamma$ is a
threshold on the coupling and is distinct from the radius $\gamma_0$
of~\eqref{4.24:eq-normalisation-constants-4}.
\end{enumerate}
Besides these, $\kappa_1$ satisfies the convergence floor
\eqref{4.24:eq-normalisation-constants-3}.

\smallskip
\emph{Conclusions.}
\begin{enumerate}
\item[(A)] The exponent of $\mathsf{t}_k$ in $\vartheta$ is at most $-1$.
\item[(B)] The floors \eqref{4.24:eq-normalisation-constants-7} and
\eqref{4.24:eq-normalisation-constants-3} are lower bounds on the integer $\kappa_1$, fixed
after $(L,\kappa)$ and before $M$. They are jointly solvable for every $L$ and $d$, and they
impose no upper bound on $L$.
\item[(C)] If the coupling threshold $\gamma$ is small enough, then the condition
$K_\Delta\gamma_0\le\tfrac12$, under which the statement of hypothesis~(g) is applied, holds
for all $g_k\le\gamma$. This condition is an upper bound on $\gamma_0$, and it is met through
an upper bound on the coupling.
\end{enumerate}
\end{definition}

\begin{proof}
\emph{Conclusion~\textup{(A)}.} By (c5), $q_*\ge p_1+1$, so $p_1-q_*\le -1$. Hence
$\vartheta$ is a constant multiple of $\bigl(\tfrac12\mathsf{t}_k\bigr)^{p_1-q_*}$, a
negative power of $\mathsf{t}_k=\log g_k^{-2}$.

\emph{Conclusion~\textup{(B)}.} Put
\begin{equation*}
F_1:=\max\{16\kappa+1,\,4L\kappa+1\},\qquad F_2:=c_J+c_\flat+\mathsf{c}_{\mathrm{an}}+8 .
\end{equation*}
Then \eqref{4.24:eq-normalisation-constants-7} reads $\kappa_1\ge F_1$, and
\eqref{4.24:eq-normalisation-constants-3} reads $\kappa_1\ge F_2$. The number $F_1$ depends
on $(L,\kappa)$ only. By hypothesis~(d), the number $F_2$ is fixed before $\kappa_1$ (after
$L$, $\kappa$, $\beta$). Neither $F_1$ nor $F_2$ involves $M$. Hence every integer
$\kappa_1\ge\max\{F_1,F_2\}$ satisfies both floors. Such an integer exists for every $L$ and
every $d$. Each of the two conditions bounds $\kappa_1$ from below once $L$ is given, so
neither restricts $L$ from above.

The size of $F_2$ depends on the value of the top rate $a_k$. We bound $F_2$ first on the
whole admissible range $a_k\le 4\kappa L$ and then at the value $a_k=(1+3\beta)\kappa$.

\emph{On the admissible range.} Hypothesis~(a) only uses rates
$a_m\le 4\kappa L$. Within this range $c_J\le 2\sqrt{d}\,(4\kappa L+1)$, so
\begin{equation*}
F_2\le 2\sqrt{d}\,(4\kappa L+1)+c_\flat+\mathsf{c}_{\mathrm{an}}+8 .
\end{equation*}
At $d=4$ this bound equals $16\kappa L+12+c_\flat+\mathsf{c}_{\mathrm{an}}$.

\emph{At the top rate $a_k=(1+3\beta)\kappa$.} This is a number depending on
$(\kappa,\beta)$. At $d=4$,
\begin{equation*}
c_J=4\bigl((1+3\beta)\kappa+1\bigr),\qquad
F_2=4(1+3\beta)\kappa+c_\flat+\mathsf{c}_{\mathrm{an}}+12 .
\end{equation*}

In both evaluations $F_2$ is finite and is fixed before $\kappa_1$. So
\eqref{4.24:eq-normalisation-constants-3} is a one-sided floor, jointly satisfiable with
\eqref{4.24:eq-normalisation-constants-7}.

Were \eqref{4.24:eq-normalisation-constants-7} imposed as the equality
$\kappa_1=16\kappa+1$, the floor \eqref{4.24:eq-normalisation-constants-3} could fail; hence
both are imposed as lower bounds.

The construction of this section and the statements of hypotheses~(a)--(g) use $\kappa_1$ in
two ways only: through lower bounds on $\kappa_1$, among them
\eqref{4.24:eq-normalisation-constants-3} and the assumption of hypothesis~(c); and through
constants that increase with $\kappa_1$, such as the factors $e^{16\kappa_1}$ and
$e^{4\kappa_1}$ and an exponential of a fixed multiple of $\kappa_1$, which are fixed after
$\kappa_1$ and before $M$ and $\gamma$. The only condition linking $\kappa_1$ and $M$ is a
lower bound on $M$. Therefore \eqref{4.24:eq-normalisation-constants-7}, with
\eqref{4.24:eq-normalisation-constants-3}, is the full set of conditions on $\kappa_1$ in
these statements, and the bounds of hypotheses~(a) and~(f) impose no further condition on
$\kappa_1$.

\emph{Conclusion~\textup{(C)}.} In hypothesis~(f), $\gamma_0$ enters the second line of
\eqref{4.24:eq-normalisation-constants-6} through $K_\Delta\gamma_0$. Since $K_\Delta$
depends on $(d,L,\kappa_1)$ only, it is fixed before $\gamma$, and the condition
$K_\Delta\gamma_0\le\frac12$ of hypothesis~(g) is an upper bound on $\gamma_0$. By
\eqref{4.24:eq-normalisation-constants-5},
\begin{equation*}
K_\Delta\gamma_0^{(k)}\le K_\Delta\,\bar\gamma_0(\mathsf{t}_k),\qquad
\bar\gamma_0(\mathsf{t})=5(1+\beta_0)^2L^2(L+1)\,C_{\max}\,\mathsf{t}^{\,r+q}e^{-\mathsf{t}/2}.
\end{equation*}
The factor in front of $\mathsf{t}^{\,r+q}e^{-\mathsf{t}/2}$ does not depend on
$\mathsf{t}$, and $\mathsf{t}^{\,r+q}e^{-\mathsf{t}/2}\to0$ as $\mathsf{t}\to\infty$, since
$e^{-\mathsf{t}/2}$ decays faster than every power of $\mathsf{t}$. Hence there is a number
$T$ with $K_\Delta\bar\gamma_0(\mathsf{t})\le\frac12$ for all $\mathsf{t}\ge T$. Now
$g_k\le\gamma$ is equivalent to $\mathsf{t}_k=\log g_k^{-2}\ge\log\gamma^{-2}$. If
$\gamma\le e^{-T/2}$, then $\log\gamma^{-2}\ge T$, so every $g_k\le\gamma$ has
$\mathsf{t}_k\ge T$ and therefore $K_\Delta\gamma_0^{(k)}\le\frac12$. The condition of
hypothesis~(g) is thus met through an upper bound on the coupling, of the one-sided form
of~(c6).
\end{proof}

\begin{definition}[Weights and totals of the fluctuation cluster expansion]
\label{4.24:def-cluster-weights}
Fix the stage, with its integration variables $\phi$, its complex background pair
$(\mathbb{U},\mathbb{J})$, the set $\mathsf{Q}$ of $\pi_{j+1}$-cells meeting the integration
bonds $\mathfrak{Bo}$, the nested cell partitions $\pi_m$ with the domains $\mathbb{D}_m$, and,
for a set $Y$, the smallest $\mathbb{D}_m$-domain $[Y]_m$ containing $Y$. Let
$\langle\cdot\rangle$ be the normalised Gaussian expectation of the stage, the one whose
normalisation constants are fixed in Definition~\ref{4.24:def-normalisation-constants}, and let
$\mathbb{C}_I(X)$ denote the localised interaction density on the region $X$.

\emph{Hypothesis.} Assume that the interaction in the exponent of the fluctuation integral
is a sum $\sum_e v_e$ over \emph{elements} $e$. Each element
carries a set $\mathsf{b}(e)\subset\mathsf{Q}$, a connected domain $\mathsf{d}(e)$ and a number
$\nu(e)\ge0$. The functional $v_e$ depends on $\phi$ only through its values on the cells of
$\mathsf{b}(e)$, and on $(\mathbb{U},\mathbb{J})$ only through their values on $\mathsf{d}(e)$.
It satisfies $|v_e|\le\nu(e)$, and $v_e=0$ if $\mathsf{b}(e)=\emptyset$.

\emph{Constants.} Let $c_F\ge0$ be the constant of the cluster expansion of \cite{B-CMP116}, and
let $a_k$ and $a^{(n+1)}_m$ be the coefficients of the exponential weights of the stage, fixed
with the stage data. Put
\begin{equation*}
c_\flat:=c_F+2\cdot 3^d\ \ge\ 2\cdot3^d,
\qquad
c_J:=2\sqrt d\,(a_k+1)\ \ge\ 0 .
\end{equation*}

\emph{Weights.} For an element $e$ put $w(e):=a^{(n+1)}_m\,d_m\bigl([\mathsf{d}(e)]_m\bigr)$,
where $m$ is the index of the cell partition $\pi_m$ used for this weight, fixed with the stage
(it is not the exponent of the torus), and
\begin{equation}\label{4.24:eq-cluster-weights-1}
\tilde\nu(e):=2\,\nu(e)\,e^{\,w(e)+c_J+c_\flat|\mathsf{b}(e)|},
\qquad
\tilde\nu'(e):=\tilde\nu(e)\,e^{|\mathsf{b}(e)|}.
\end{equation}

\emph{Totals.} Put
\begin{equation}\label{4.24:eq-cluster-weights-2}
\mu:=\sup_{q\in\mathsf{Q}}\ \sum_{e:\ \mathsf{b}(e)\ni q}\tilde\nu(e),
\qquad
\mathsf{M}_*:=\sup_{y}\ \sum_{e:\ [\mathsf{d}(e)]_k\ni y}\tilde\nu(e),
\end{equation}
and let $\mu'$ and $\mathsf{M}'_*$ be defined by the same two formulas with $\tilde\nu'$ in place
of $\tilde\nu$. Put further
\begin{equation}\label{4.24:eq-cluster-weights-3}
\tilde{\mathfrak{l}}:=2d\,(LM)^4\,e^{-\frac12\gamma_2T_j^2+c_J+c_\flat}.
\end{equation}
Here $\gamma_2$ and $T_j$ are numbers fixed with the stage data, as are the coefficients $a_k$
and $a^{(n+1)}_m$ above; they enter this definition only through \eqref{4.24:eq-cluster-weights-3}.
Finally, $V$ denotes the supremum over the components of the stage that enters the bound
\textup{(ii)} below.

\emph{Second hypothesis.} Assume the following statement on the fluctuation integral,
formulated in terms of the quantities above: if $\nu(e)\le\frac12$
for every element $e$, if $\mu+\tilde{\mathfrak{l}}\le\frac12$, and if the flat smallness
condition of the stage with parameter $\kappa$ and hypothesis~(H2) of that statement hold,
then
\begin{enumerate}
\item[(i)] one has
\begin{equation*}
\log\Bigl\langle \chi^{(j)}\chi'^{(j)}\exp\sum_e v_e\Bigr\rangle=\sum_X\mathbb{C}_I(X).
\end{equation*}
Here $X$ runs over the regions admitted by the localisation of the stage with
$X\cap R_n\neq\emptyset$, and each $\mathbb{C}_I(X)$ is holomorphic and local;
\item[(ii)] $\|\mathbb{C}_I\|_{(n+1)}\le 4\bigl(\mathsf{M}_*+2^dV\tilde{\mathfrak{l}}\bigr)$,
where $\|\cdot\|_{(n+1)}$ is the norm of localised densities of the stage at level $n+1$;
\item[(iii)] the corresponding statement in the dial variable $\mathfrak{z}$ holds, with the
constants $8$ and $32$.
\end{enumerate}
\end{definition}

\begin{proof}
The number $c_F$ is nonnegative, so $c_\flat=c_F+2\cdot3^d\ge2\cdot3^d$. The inequality
$c_J\ge0$ is taken with the stage data, together with the coefficient $a_k$; since $\sqrt d>0$,
it amounts to $a_k\ge-1$.

The weights \eqref{4.24:eq-cluster-weights-1} and the totals \eqref{4.24:eq-cluster-weights-2},
\eqref{4.24:eq-cluster-weights-3} are formed only from the data $(\mathsf{b}(e),\mathsf{d}(e),
\nu(e))$ supplied by the first hypothesis, together with the constants of the stage. They are
therefore well defined as soon as that hypothesis holds.

Statements \textup{(i)}--\textup{(iii)} are the conclusions of the second hypothesis,
read with the quantities defined above. They hold whenever the hypotheses displayed
before \textup{(i)} hold.
\end{proof}

\begin{definition}[Stored terms and parameter-dependent elements]\label{4.24:def-stored-terms}
Fix a level $j'$, and let $j$ be the step of the expansion, whose small-field characteristic
functions are $\chi^{(j)}$ and $\chi'^{(j)}$. Elements $e$, with their data $\mathsf{b}(e)$,
$\mathsf{d}(e)$, $\nu(e)$ and their functions $v_e$, are those of
Definition~\ref{4.24:def-cluster-weights}.
\begin{enumerate}
\item[(a)] A \emph{stored term} is a triple $(c,S,\mathcal{S})$, where $S$ is the set on which
the term is stored, the argument of both domains below, $\mathcal{S}$ is a
family of strong cubes and $c$ is a function on the product
\begin{equation*}
\mathfrak{U}_c(S)\times\mathfrak{D}_{\mathcal{S}}(S).
\end{equation*}
Among the arguments of $c$ there is a distinguished family, indexed by cubes, called the
\emph{A-slots} of $c$. On the cubes of $\mathcal{S}$ each A-slot ranges over the real box of
radius $\mathfrak{b}_{j'}$, and on the cubes not in $\mathcal{S}$ over the complex polydisc of
radius $\varrho_{j'}$; the latter is called the \emph{A-polydisc}. We put
\begin{equation}\label{4.24:eq-stored-terms-1}
\mathfrak{N}(c):=\sup_{\mathfrak{U}_c(S)\times\mathfrak{D}_{\mathcal{S}}(S)}|c| .
\end{equation}
\item[(b)] An element $e$ is \emph{parameter-dependent} if $v_e$ depends on a parameter
$\lambda\in D_e$, where $D_e$ is the product of a real box and a complex polydisc, continuously
in $\lambda$ and holomorphically in the polydisc coordinates, and if
\begin{equation*}
|v_e|\le\nu(e)\quad\text{for all }\lambda\in D_e .
\end{equation*}
\item[(b$'$)] Let $\gamma$ be a polymer at level $j'$, regarded as a set of elements. The
strong set of $\gamma$ at level $j'$ is
\begin{equation}\label{4.24:eq-stored-terms-2}
\mathcal{S}_{j'}(\gamma):=\bigcup_{e\in\gamma}\mathcal{S}_{j'}\bigl(c(e)\bigr),
\end{equation}
where each element $e\in\gamma$ carries a stored term $(c(e),S(e),\mathcal{S}(e))$ in the sense
of (a), and $\mathcal{S}_{j'}(c(e))$ denotes its family $\mathcal{S}(e)$ of strong cubes.
\end{enumerate}
The following three statements are used as stated: (c) is the lemma on parameter-dependent
elements, (d) is its corollary for the exterior slots, and (e) is the statement that the
A-polydisc never shrinks.
\begin{enumerate}
\item[(c)] Let the elements be parameter-dependent in the sense of (b), and suppose that none of
the following depends on $\lambda$: the pieces of the third and fourth kinds in the
classification of the pieces of the step, the large-field bonds, the decoupled cubes, the
Gaussian integration together with its
decoupling in the parameters $\mathsf{s}(\Delta)$, the characteristic functions
$\chi^{(j)}\chi'^{(j)}$, the partitions, and the contours. Then clauses (ii) and (iii) of the
fluctuation cluster-expansion bound hold, with the supremum over the parameters $\lambda$, each
ranging over its domain $D_e$, taken inside the norm $\|\cdot\|$ of that bound.
\item[(d)] Statement (c) applies, and its conclusion holds, with $\lambda$ taken to be the
exterior slots of the function $c(e)$ of the stored term carried by the element $e$, $D_e$ being
a product of a real box and a complex polydisc as in (b).
\item[(e)] The complex polydisc of radius $\varrho_{j'}$ that (a) assigns to the A-slots,
the A-polydisc, is never replaced by a polydisc of smaller radius.
\end{enumerate}
\end{definition}

\begin{definition}[Exterior variables, frozen and mixed members, and the stage bound]
\label{4.24:def-exterior-variables}
Fix the step to level $k$ and, within it, the stage $n+1$, which is performed at level $j$. We use
the following objects of the stage:
the arrested components $\mathfrak{W}$; the survivors' region $Z^{s}$; the rim $\mathsf{R}_j$; the
survivors' and arrested bonds $s$, $\mathcal{W}$; and their variables $\phi_s$ and $\phi_W$
($W\in\mathfrak{W}$). For $W\in\mathfrak{W}$, $\mathcal{A}_{WW}$ is the $W$-diagonal block of the
symmetrised stage operator $\mathcal{A}$, and
\begin{equation*}
\mathsf{q}_W(x)=-\tfrac12\langle x,\mathcal{A}_{WW}x\rangle
\end{equation*}
is the collar quadratic of $W$. The weights and totals $\tilde\nu$, $\mathsf{M}_*$, $V$,
$\tilde{\mathfrak{l}}$, $c_J$, $c_\flat$ are those of Definition~\ref{4.24:def-cluster-weights}, and
the constants of the Gaussian normalisation are those of
Definition~\ref{4.24:def-normalisation-constants}. A cube of $Z^{s}$ is called \emph{surviving}. We put
\begin{equation}\label{4.24:eq-exterior-variables-1}
\begin{gathered}
\vartheta_a:=\varepsilon_Y\bar V K_a(\varrho^B_j)^2,\qquad
\theta_{\mathrm{cv}}:=4H_{\max}/g_j,\\
\vartheta_{\mathrm{cv}}:=\theta_{\mathrm{cv}}\,\varepsilon_Y\bigl[1+2\bar V\bigl(1+K_a(\varrho^B_j)^2
+n_{\mathfrak{g}}n_b\bigr)\bigr],
\end{gathered}
\end{equation}
where $n_b$ is a constant of the stage; $\vartheta_{\mathrm{cv}}$ is the
quantity of clause~(e) of the stage theorem.
Assume the following.
\begin{enumerate}
\item[(1)] The parameters of the stage satisfy
\begin{equation}\label{4.24:eq-exterior-variables-2}
\begin{aligned}
&5(\vartheta_{\mathrm{cv}}+\vartheta_a)\le\min\{\tfrac12-\vartheta,\,g_k\},\\
&4\mathsf{n}_{\mathrm{adj}}\,e^{c_J+2c_\flat}\le e^{R_k},\qquad R_k\ge3,\\
&\kappa_1-1\ge c_J+c_\flat+\mathsf{c}_{\mathrm{an}}+7.
\end{aligned}
\end{equation}
They also satisfy the condition $(Q^{+})'$, defined as the condition $(Q^{+})$ of the
fluctuation cluster expansion whose weights are those of Definition~\ref{4.24:def-cluster-weights},
taken with $c_\flat$ replaced by $c_\flat+1$ and with $\mathfrak{p}_j$
replaced by $\mathfrak{p}_j+\vartheta_{\mathrm{cv}}+\vartheta_a$.
\item[(2)] The smallness number $\varepsilon_1$ of the summation step of the change-of-variables
estimate is the number for which that step gives
\begin{equation*}
\sum_Y\tilde\nu(e,Y)\le\theta_{\mathrm{cv}}\bigl(2e^{c_J}\varepsilon_1\bigr)\tilde\nu'(e).
\end{equation*}
This $\varepsilon_1$ is a number of that step, distinct from the radius-type parameters written
with the same letter elsewhere. It satisfies
\begin{equation}\label{4.24:eq-exterior-variables-3}
\varepsilon_1\le4\mathsf{n}_{\mathrm{adj}}\,e^{-7R_k}.
\end{equation}
\item[(3)] Under the conditions of item~(1) and with the value \eqref{4.24:eq-exterior-variables-3},
the proposition bounding the output of the stage gives the bound
\eqref{4.24:eq-exterior-variables-4} below for the total density defined in part~(b).
\end{enumerate}
We define the following.

\smallskip
(a) \emph{Exterior variables.} For $W\in\mathfrak{W}$ the variables $\phi_W$ are \emph{exterior}, and
$\phi_W$ is among the variables collected in the exterior argument $\lambda$. The family $u$, the
strong sets $\mathcal{S}_j(c)$ and the terms
$e_{\mathrm{str}}$ are indexed by the surviving cubes $c$ of the rim $\mathsf{R}_j$. For each such
cube $c$, the datum of $c$ holds the restriction to $W$ of the level-$j$ field $A_j$ in the
stored-term domain
$\mathfrak{D}_{\mathcal{S}}$. Elements that are not shifted depend on the argument $x$ of
$\mathsf{q}_W$ only weakly, with smallness parameter $e^{-2\varepsilon_\star}\varrho^B_j$, where
$\varepsilon_\star$ is a constant of the stage; this part of the
definition imposes no requirement on the change-of-variables objects $\mathfrak{O}$. Every variable
of $W$ that is not integrated by a stage already performed is an exterior real variable of the
large-field operation $\mathbf{T}'_k(W)$ in the sense of \cite{B-CMP122a}. It is confined by
background-free functions, as in \cite[(1.59)]{B-CMP122a}.

\smallskip
(b) \emph{Frozen and mixed members; the total density.} Let a member of the stage's expansion read
some $\phi_W$, and let $\mathsf{b}$ be the set of cubes in which it reads integration variables. The
member is \emph{frozen} if $\mathsf{b}$ contains no surviving cube. It is \emph{mixed} otherwise. The
elements $e_a$, $e_a^{0}$ both of whose bonds lie off $\bigcup\mathfrak{W}$ belong to the survivors'
Gaussian and are left untouched. The members are sorted as follows.
\begin{enumerate}
\item[(i)] Mixed members join the fluctuation cluster expansion of the survivors
(Definition~\ref{4.24:def-cluster-weights}) as elements.
\item[(ii)] Each frozen member is assigned to the first $W\in\mathfrak{W}$, in the order of
$\mathfrak{W}$ fixed by the stage. The frozen members assigned to $W$ make up the frozen
density $\mathbb{C}^{\mathrm{fr}}_W$.
\end{enumerate}
The quadratic $\mathsf{q}_W(x)$ stays in the exponent of the integrand of $\mathbf{T}'_k(W)$, by the
statement of this chapter that governs the integrand of $\mathbf{T}'_k(W)$. Write
$\mathbb{C}^{s}_{\mathrm{out}}$ for the output of the survivors' cluster expansion. The
total density of the stage is
\begin{equation*}
\mathbb{C}^{(n+1)}_k:=\mathbb{C}^{s}_{\mathrm{out}}
+\sum_{W\in\mathfrak{W}}\mathbb{C}^{\mathrm{fr}}_W,
\end{equation*}
and the output of the stage is written $\mathfrak{F}(\mathfrak{a})$.

\smallskip
(c) \emph{The stage bound.} Under hypotheses (1)--(3),
\begin{equation}\label{4.24:eq-exterior-variables-4}
\begin{aligned}
\bigl\|\mathbb{C}^{(n+1)}_k\bigr\|^{A,\flat}_{(n+1)}
&\le3c_G+\vartheta+4\bigl(\mathsf{M}_*+2^dV\tilde{\mathfrak{l}}\bigr)
+5(\vartheta_{\mathrm{cv}}+\vartheta_a)\\
&\le C_0^{\sharp}.
\end{aligned}
\end{equation}
\end{definition}

\begin{proof}
Part~(a) is a definition: it fixes which variables of $W$ are exterior and the properties asked of
them.

In part~(b), the two cases are exclusive and exhaust the members reading some $\phi_W$: either
$\mathsf{b}$ contains a surviving cube or it does not. A frozen member reads $\phi_W$ for at least one
$W\in\mathfrak{W}$, so $\mathfrak{W}$ is not empty. Hence the first $W$ in the order of
$\mathfrak{W}$ fixed by the stage exists and is unique. So each frozen member is assigned to exactly
one arrested component, and
$\mathbb{C}^{\mathrm{fr}}_W$ is defined for every $W$. The mixed members enter the survivors'
cluster expansion as elements, by sorting rule~(i), so the output of that expansion is
defined. Therefore $\mathbb{C}^{(n+1)}_k$ is defined.

For part~(c), the parameters of the stage satisfy the conditions
\eqref{4.24:eq-exterior-variables-2} and $(Q^{+})'$ by hypothesis~(1). The smallness number of the
summation step satisfies \eqref{4.24:eq-exterior-variables-3} by hypothesis~(2). These are the
conditions and the value under which hypothesis~(3) applies. Hypothesis~(3), applied to the density
$\mathbb{C}^{(n+1)}_k$ of part~(b), gives both inequalities of
\eqref{4.24:eq-exterior-variables-4}.
\end{proof}

\begin{definition}[Elements of the survivors' expansion and cross totals]\label{4.24:def-survivor-elements}
Let $\mathfrak{E}$ be the family of elements $e$ of the fluctuation cluster expansion, with data
$(\mathsf{b}(e),\mathsf{d}(e),\nu(e))$ and weights $\tilde\nu(e)$, $\tilde\nu'(e)$, as in
Definition~\ref{4.24:def-cluster-weights}. Elements are read with the $u$-family and the exterior
variables of Definition~\ref{4.24:def-exterior-variables}. Let $s$ and $\mathcal{W}$ be the
survivors' and the arrested integration bonds, and let $\mathsf{Q}^{s}$ be the set of
$\pi_{j+1}$-cells meeting $s$.
The \emph{element list of the survivors' stage} is the disjoint union
\begin{equation}\label{4.24:eq-survivor-elements-1}
  \mathfrak{E}^{\mathrm{StB}} := \mathfrak{E}^{s}\sqcup\mathfrak{E}^{\times}\sqcup\mathfrak{O}^{s}
\end{equation}
of the following three families; the superscript $\mathrm{StB}$ marks it as the list of the
survivors' stage.
\begin{itemize}
\item $\mathfrak{E}^{s}$ is the set of those $e\in\mathfrak{E}$ with
  $\mathsf{b}(e)\cap\mathsf{Q}^{s}\neq\emptyset$. It contains the survivors' own members and the
  mixed members of Definition~\ref{4.24:def-exterior-variables}.
\item $\mathfrak{E}^{\times}$, the \emph{collar pairs}, is the set of the collar elements
  $e_a(b,b',Y')$ of clause~(iii) of the lemma on the decoupled collar kernel for which exactly
  one of $b,b'$ lies in $s$ and the other lies in $\mathcal{W}$. For such an $e$ we write
  $b_s(e)$ for the bond in $s$, the \emph{survivor bond}, and $b_{\mathfrak{W}}(e)$ for the bond
  in $\mathcal{W}$. No element $e^{0}_{a}$ of that lemma enters $\mathfrak{E}^{\times}$, as is
  shown in the statement of this chapter that treats the elements $e^{0}_{a}$; thus
  $\mathfrak{E}^{\times}$ consists of elements $e_a$ only.
\item $\mathfrak{O}^{s}$ is the set of those change-of-variables objects $o\in\mathfrak{O}$ with
  $\mathsf{b}(o)\cap\mathsf{Q}^{s}\neq\emptyset$. Here $\mathfrak{O}$ is the family provided by
  clause~(c) of the theorem that introduces the change-of-variables objects.
\end{itemize}
Each member of $\mathfrak{E}^{\mathrm{StB}}$ carries the triple
$(\mathsf{b}\cap\mathsf{Q}^{s},\ \mathsf{d},\ \nu^{\natural})$, where $\nu^{\natural}$ is its
cross-member majorant.

Let $e=e_a(b,b',Y')\in\mathfrak{E}^{\times}$. Then $b_s(e)$ and $b_{\mathfrak{W}}(e)$ are
integration bonds of two distinct components. Hence the separation of components
\eqref{4.24:eq-component-separation-a} applies to them, and
\begin{equation}\label{4.24:eq-survivor-elements-2}
  |Y'|\ge R_k .
\end{equation}

The \emph{cross totals} are
\begin{align}
  \mu^{\times}
  &:= \sup_{q\in\mathsf{Q}^{s}}
     \sum_{\substack{e\in\mathfrak{E}^{\times}\\ \mathsf{b}(e)\cap\mathsf{Q}^{s}\ni q}}
     \tilde\nu'(e),
  \label{4.24:eq-survivor-elements-3}\\
  \mathsf{M}^{\times}
  &:= \sup_{y}
     \sum_{\substack{e\in\mathfrak{E}^{\times}\\ {[\mathsf{d}(e)]_k}\ni y}}
     \tilde\nu(e),
  \label{4.24:eq-survivor-elements-4}
\end{align}
and $\mathsf{M}'^{\times}$ is defined as in \eqref{4.24:eq-survivor-elements-4} with
$\tilde\nu'$ in place of $\tilde\nu$. Here $[\mathsf{d}(e)]_k$ is the smallest
$\mathbb{D}_k$-domain containing $\mathsf{d}(e)$.
\end{definition}

\begin{remark}[The normalisation as a resolvent expansion about a scalar]
\label{4.24:rem-scalar-resolvent}
Let $C^{*}\Delta^{(j)}C$ be the operator of the stage quadratic form of
Definition~\ref{4.24:def-stage-form}, and let $z^{(j)}(c)$, $E_0^{(j)}$ and $\lambda_0$ be the
constants of Definition~\ref{4.24:def-normalisation-constants}. In the Gaussian normalisation
lemma, the normalisation of the stage Gaussian and its localised densities $\mathbb{C}_G$ have the
following structure.
\begin{enumerate}
\item[(a)] The logarithm of the normalisation is the sum
\begin{equation}\label{4.24:eq-scalar-resolvent-1}
-\tfrac12\log\det\bigl(C^{*}\Delta^{(j)}C\bigr)+\sum_{c}\log z^{(j)}(c)+\mathrm{const}+E_0^{(j)} ,
\end{equation}
in which the last three summands are numbers and local holomorphic functions attached to bonds
and blocks of the collar, that is, single-cube terms of $\mathbb{D}_{j+1}$.
\item[(b)] The summand $-\tfrac12\log\det\bigl(C^{*}\Delta^{(j)}C\bigr)$ is represented by a
contour integral of $\log z$ times the trace of the resolvent of $C^{*}\Delta^{(j)}C$, and that
resolvent is expanded by the first-order resolvent identity about
the scalar reference $\lambda_0 I$, where $I$ is the identity. The first member gives a field-free
number, which contributes to $E_0^{(j)}$; the second member, collapsed onto the half-line
$z=-\xi$, $\xi\ge0$, is a sum of bond densities.
\item[(c)] The operator then expanded is the resolvent
$\bigl(C^{*}\Delta^{(j)}C+\xi I\bigr)^{-1}$, $\xi\ge0$, by the generalised random walk expansion of
\cite[Theorem~3.10]{B-CMP99b} and, on a sub-region, of \cite[Theorem~3.15]{B-CMP99b}; its building
blocks are local inverses on cubes of the proper scales. This yields, as in the Gaussian
normalisation lemma, $\|\mathbb{C}_G\|\le c_G$, and the bound of the field-dependent part by
$\vartheta$, \eqref{4.24:eq-scalar-resolvent-3}.
\item[(d)] Consequently the logarithm is never expanded; the only reference subtracted is the
scalar $\lambda_0 I$; no series of $\operatorname{Tr}\log(I+T)$, $T$ a perturbation, about an
operator reference occurs; and no decoupling or interpolation parameter ($\sigma'$, $\mathsf{s}$,
or the interpolation parameter $t$ of \cite[(3.30)--(3.32)]{B-CMP119}) enters the
kernel whose determinant is taken. Decoupling parameters enter only the interaction and the walks
of the covariance after whitening, and positivity is used once, for the undecoupled form, at the
large-field bonds.
\end{enumerate}
\end{remark}

\begin{proof}
\emph{Part} (a). By \cite[(3.33)]{B-CMP119}, the logarithm of the normalisation equals
\eqref{4.24:eq-scalar-resolvent-1}, where $c$ runs over the cells of the region on which the
$\delta$-functions were eliminated and $z^{(j)}=\prod_c z^{(j)}(c)$. The factors $z^{(j)}(c)$ and
the constant arise from the elimination of the $\delta$-functions and from the gauge-fixing
normalisations, and
$E_0^{(j)}$ is the number that \cite[(3.33)]{B-CMP119} assigns to the vacuum-energy
renormalisation. Each of these is a number or a local holomorphic function attached to a bond or a
block of the collar, hence a single-cube term of $\mathbb{D}_{j+1}$. What remains is the first
summand.

\emph{Part} (b). In \cite{B-CMP119} the level is called $k$; we write $j$ throughout. Following
\cite[(3.34)]{B-CMP119},
\[
-\tfrac12\log\det\bigl(C^{*}\Delta^{(j)}C\bigr)=-\tfrac12\operatorname{Tr}\log\bigl(C^{*}\Delta^{(j)}C\bigr),
\]
and this is written as a contour integral over a contour $\gamma$ surrounding the spectrum. The
integrand is $\log z\,\operatorname{Tr}\bigl(C^{*}\Delta^{(j)}C-zI\bigr)^{-1}$, and $\gamma$
contains an interval of the line $\operatorname{Re}z=r$ with $r>0$ small. With
$\lambda_0>r$ small enough, the resolvent in the integrand is expanded as in
\cite[(3.35)]{B-CMP119}, that is, by the first-order resolvent identity about the scalar
$\lambda_0 I$, which we write in the form
\begin{equation}\label{4.24:eq-scalar-resolvent-2}
\begin{split}
\bigl(C^{*}\Delta^{(j)}C-zI\bigr)^{-1}
&=(\lambda_0-z)^{-1}I\\
&\quad-(\lambda_0-z)^{-1}\bigl(C^{*}\Delta^{(j)}C-\lambda_0 I\bigr)
\bigl(C^{*}\Delta^{(j)}C-zI\bigr)^{-1}.
\end{split}
\end{equation}
The identity holds because $(\lambda_0-z)I-\bigl(C^{*}\Delta^{(j)}C-zI\bigr)$ equals
$-\bigl(C^{*}\Delta^{(j)}C-\lambda_0I\bigr)$. Multiplying this equation on the left by
$(\lambda_0-z)^{-1}$ and on the right by $\bigl(C^{*}\Delta^{(j)}C-zI\bigr)^{-1}$ gives
\eqref{4.24:eq-scalar-resolvent-2}.

By the sentence following \cite[(3.35)]{B-CMP119}, the integral of the first member of
\eqref{4.24:eq-scalar-resolvent-2} equals
$-\tfrac12\log\lambda_0\,\operatorname{Tr}I$. This number does not depend on the field, and it
contributes to the vacuum-energy counterterm $E_0^{(j)}$.

The second member is $O(|z|^{-2})$. For this reason the text following \cite[(3.35)]{B-CMP119}
replaces the contour and takes
the limit $r\to0$, which collapses the integral onto the half-line $z=-\xi$, $\xi\ge0$. The result,
\cite[(3.36), (3.37)]{B-CMP119}, is a sum of bond densities indexed by the level-$j$ bonds $b'$ of
the region over which the trace in \cite[(3.37)]{B-CMP119} is taken. The integrand of each density
is built from the operator
\[
\bigl(C^{*}\Delta^{(j)}C-\lambda_0I\bigr)\bigl(C^{*}\Delta^{(j)}C+\xi I\bigr)^{-1},
\qquad \xi\ge0,
\]
under the trace.

\emph{Part} (c). The operator expanded in the bond densities is the resolvent
$\bigl(C^{*}\Delta^{(j)}C+\xi I\bigr)^{-1}$ with $\xi\ge0$. It is expanded by the generalised random
walk expansion for the sequence of domains $\{\Omega'_m\}$ on which that expansion is built,
restricted to the integration region.

By \cite[Theorem~3.10]{B-CMP99b}, for $M$ sufficiently large and a
configuration satisfying \cite[(3.95)]{B-CMP99b}, the propagator is the sum
$\sum R(X_0)W(X_1)\cdots R(X_n)$ of \cite[(3.107)]{B-CMP99b}, in which $R(X_a)$ is the local
inverse on the cube $X_a$ and $W(X_a)$ is the operator attached to the cube $X_a$ in that display.
A term of this sum depends on the
configuration only through its restriction to $X_0\cup X_1\cup\dots\cup X_n$. The constant $O(1)$
in the bounds depends on $d$ and $L$ only.

For the propagator on a sub-region $\Lambda'$ the same holds by \cite[Theorem~3.15]{B-CMP99b}: the
constants $B_0,\delta_0$ there depend on $d$ and $L$ only, and this propagator has a convergent
random walk expansion of the same type.

The operators $R(X_a)$ are local inverses on cubes of the proper scales. These two theorems are used
in their full generality. A walk $\omega=(X_0,\dots,X_p)$ through $M_1$-cubes of the proper scales
contributes a term that depends on the background $\mathbb{U}$ only on $\bigcup_aX_a^{\sim}$. Its
bound is
\[
O(1)\bigl(O(1)M^{-1/2}\bigr)^{p}e^{-\delta_0 d(\omega)} ,
\]
where $d(\omega)$ is the length attached to the walk $\omega$ in the bound of
\cite[Theorem~3.10]{B-CMP99b}, not the dimension $d$.

The walks are then resummed by their hulls. The number of densities attached to one
$\pi_{j+1}$-cube is $(LM)^4\cdot O(1)=:c_G'(L,M)$, whence, as in the Gaussian normalisation lemma,
$\|\mathbb{C}_G\|\le c_G$.

For the dependence on the field, note two facts. The operators are covariant under gauge
transformations with values in the complexified gauge group $G^{\mathbb{C}}$, and the determinant
is invariant under similarity. The field dependence is therefore carried by a difference of
resolvents, whose contribution to each density is bounded, as in the Gaussian normalisation lemma,
by
\begin{equation}\label{4.24:eq-scalar-resolvent-3}
c_G'\cdot O(1)\bigl(B\,\hat C\,M\,\tilde\alpha_0+\tilde\alpha_1\bigr),
\end{equation}
where $B$ and $\hat C$ denote the two constants of the Gaussian normalisation lemma that multiply
$M\tilde\alpha_0$ in this bound; and this quantity is at most $\vartheta$.

\emph{Part} (d). Parts (a)--(c) describe the whole treatment of the normalisation. In it the
logarithm enters only through the contour representation of part~(b), and it is not expanded. The
only reference operator subtracted is the scalar $\lambda_0 I$ in
\eqref{4.24:eq-scalar-resolvent-2}. No series of $\operatorname{Tr}\log(I+T)$ and no operator
reference appears.

The integrals $\int dt$ that occur in \cite[(3.37)]{B-CMP119} are the $t$-interpolated expectations
of the characteristic functions and of the interaction from \cite[(3.30)--(3.32)]{B-CMP119}. In
the division of the step adopted in this chapter they belong to the logarithm treated in the
interaction-logarithm proposition, not to the
normalisation of the Gaussian normalisation lemma.
There, $t$ multiplies $g_j$ in the characteristic function $\chi^{(j)}$ and in the action. It never
multiplies the Gaussian kernel, whose logarithm separates free of $t$ in
\cite[(3.30)]{B-CMP119}.

Decoupling parameters occur in two places only. First, in the interaction: in the first step of the
proof of the interaction-logarithm proposition, each walk of the random walk expansions of the
operators entering the interaction is multiplied by $\prod_{\Delta}\mathsf{s}(\Delta)$, the product
running over the cubes $\Delta$ that the walk meets (here $\Delta$ denotes a cube, not to be
confused with the operator $\Delta^{(j)}$). Second, in the walks of the covariance after
whitening: in the second step of that proof, the Gaussian is conditioned on the complement
$Z_0^{c}$, and the covariance $\mathcal{C}=\bigl(C^{*}\Delta^{(j)}C\bigr)^{-1}$ of
Definition~\ref{4.24:def-stage-form} is itself a conditional covariance on a subdomain by
\cite[Theorem~3.15]{B-CMP99b}. The Gaussian is then whitened with $\mathcal{C}^{1/2}$ through the
resolvent, and the walks of $\mathcal{C}$ are decoupled with $\mathsf{s}$ on cubes of the filing
partitions of that proposition, that is, of the partitions into cubes on which those walks are
decoupled.

Positivity is used once, for the undecoupled form, at the large-field bonds, where the absorption
into the Gaussian uses
$\operatorname{Re}\bigl(C^{*}\Delta^{(j)}C\bigr)\ge\|\mathcal{C}\|^{-1}$ and
$\alpha_5\|\mathcal{C}\|<\tfrac12$ from \cite[(2.31)]{B-CMP116}. The complex background enters
$\alpha_5$ through $O(\tilde\alpha_0+\tilde\alpha_1)$. Hence no parameter $\sigma'$, $\mathsf{s}$
or $t$ enters the kernel $C^{*}\Delta^{(j)}C$ whose determinant is taken, which proves~(d).
\end{proof}

\begin{theorem}[Normalisation and cluster expansion at the survivors' block]
\label{4.24:thm-survivor-block-expansion}
Work at level $j$, at stage $n+1$ of the run of Definition~\ref{4.24:def-survivors-bonds}, and
at a fixed term of the expansion \cite[(1.55)]{B-CMP122a} whose set $\mathfrak{W}$ of
arrested components is non-empty. The objects used are those of the following statements:
\begin{itemize}
\item Definition~\ref{4.24:def-survivors-bonds}: the components, the survivors' region $Z^{s}$,
the survivors' bonds $s$ and the arrested bonds $\mathcal{W}$;
\item Notation~\ref{4.24:not-fibres-pairings}: the fibres, the pairings $(\cdot,\cdot)$ and
$\langle\cdot,\cdot\rangle$, and the gauge action;
\item Definition~\ref{4.24:def-stage-form}: the stage form $K_{\sharp}$, its symmetrised
operator $\mathcal{A}$ and the blocks $\mathcal{A}_{ss}$, $\mathcal{A}_{s\mathcal{W}}$,
$\mathcal{A}_{\mathcal{W}\mathcal{W}}$;
\item Definition~\ref{4.24:def-normalisation-constants}: the constants $z^{(j)}_s$,
$E_0^{(j)}(\cdot)$ and $c_G$;
\item Definition~\ref{4.24:def-cluster-weights}: the weights $\tilde\nu,\tilde\nu'$ and the
totals $\mu,\mu',\mathsf{M}_*,\mathsf{M}'_*$, together with $\tilde{\mathfrak{l}}$, $V$,
$\gamma_2$, $c_\flat$, $c_J$, $c_F$ and $K(d)$;
\item Definition~\ref{4.24:def-stored-terms}: the stored terms, the domains
$\mathfrak{D}_{\mathcal{S}}(\cdot)$, the parameter domains $D_e$ and the element parameter
$\lambda\in D_e$, with its box and polydisc coordinates;
\item Definition~\ref{4.24:def-exterior-variables}: the exterior variables $\phi_W$, the
smallnesses $\vartheta_a$ and $\vartheta_{\mathrm{cv}}$, and the stage bound with its
clauses (F1) and (F2);
\item Definition~\ref{4.24:def-survivor-elements}: the element list
$\mathfrak{E}^{\mathrm{StB}}$, the cubes $\mathsf{Q}^{s}$ and the majorants $\nu^{\natural}(e)$;
\item the statement on totals uniform in the number of arrested components, whose bounds are
\eqref{4.24:eq-uniform-totals-a}: the total $\mathsf{M}^{s}_*$ and the totals
$\mu^{\mathrm{StB}}$, $\mathsf{M}^{\mathrm{StB}}_*$.
\end{itemize}
For a domain $X$ let $\mathfrak{D}^{\flat}_{n+1}(X)$ be the analyticity domain of the stage
densities, and consider its subdomain on which the ambient condition attached to the arrested
components is imposed; this condition is void unless $X$ meets some $Z_W^{+}$. We call this
subdomain the ambient subdomain of $\mathfrak{D}^{\flat}_{n+1}(X)$. For rates $a=(a_m)_m$ let
$\|\cdot\|_{a}$ be the weighted norm of localised densities of
Definition~\ref{4.24:def-normalisation-constants}, taken at the rates $a$; it is the norm in
which the constant $c_G$ of that definition is measured. Let
$\|\cdot\|_{(n+1)}$ be the stage norm of localised densities; it is a norm of a different form
from $\|\cdot\|_{a}$, and the passage from a bound in $\|\cdot\|_{a}$ to a bound in
$\|\cdot\|_{(n+1)}$ is the norm-conversion proposition of clause~(i).

Assume the following.
\begin{itemize}
\item[(h1)] The standing conditions of the stage under which the stage bound is constructed,
namely two conditions on the constants of the stage, six conditions (i)--(vi) on the
change-of-variables parameters, and the condition
$(Q^{+})'$; the two standing hypotheses of the abstract cluster-expansion proposition; the
window condition; and the parameter condition
\begin{equation*}
\kappa_1\ \ge\ \max\{16\kappa+1,\ 4L\kappa+1\}
\end{equation*}
of Definition~\ref{4.24:def-normalisation-constants}.
\item[(h2)] The following three statements of this chapter: the premise statement of the stage,
which fixes the hypotheses under which the stage is performed; the statement that fixes the
analyticity domains $\mathfrak{D}^{\flat}_{n+1}(X)$, the displays belonging to them and their
ambient condition; and the arrest theorem.
\item[(h3)] The normalisation lemma at the full stage and the fluctuation cluster expansion at
the full stage, together with the three statements of the full stage on which these two rest:
the Gaussian lemma at the full stage, the fourth structural clause at the full stage, and the
supplementary section of the full-stage analysis.
\item[(h4)] \cite[Theorems~3.10 and~3.15]{B-CMP99b} and the analysis of the Gaussian integrals
in \cite[(3.33)--(3.37)]{B-CMP119}, each used as stated there. In addition, the statements of
\cite{B-CMP119} on the dependence on the parameters $\lambda$ and on symmetry, used as stated
there.
\item[(h5)] The coercivity of the Hermitian part of the stage form on its complex domain,
together with the tolerance statement used by the fluctuation cluster expansion at the full
stage, both under the standing hypotheses of the fluctuation cluster expansion at the full
stage.
\item[(h6)] The abstract cluster-expansion proposition, with its second clause in pointwise
form, and the separation of components (Lemma~\ref{4.24:lem-component-separation}).
\item[(h7)] The parameter lemma for cluster elements and its corollary, with the smallness
condition on stored terms attached to them.
\item[(h8)] The collar-kernel lemma, the construction of the exterior variables and of the stage
bound in Definition~\ref{4.24:def-exterior-variables}, and the smallness $\varepsilon_1$ fixed
in the construction of the change of variables.
\end{itemize}
Then the following hold.

\smallskip\noindent
(i) \emph{The normalisation at the block.} Let $d\phi_s$ be the Lebesgue measure on
$\bigoplus_{b\in s}\mathfrak{g}$ in $(\cdot,\cdot)$-orthonormal coordinates; Ba{\l}aban's
measure constants are contained in the constant of \cite[(3.33)]{B-CMP119}. Let $\phi_s$ range
over the survivors' rim and collar variables, with the joint form of
Definition~\ref{4.24:def-stage-form}. Then
\begin{equation}\label{4.24:eq-survivor-block-expansion-1}
\begin{split}
\log\Bigl[\,z^{(j)}_s\int d\phi_s\,
\exp\bigl(-\tfrac12\langle\phi_s,\mathcal{A}_{ss}\phi_s\rangle\bigr)\Bigr]
&+E_0^{(j)}\bigl((\Omega^{c}_{j+1}\setminus Z''_{j+1})\cap Z^{s}\bigr)\\
&=\sum_X\mathbb{C}^{s}_G(X).
\end{split}
\end{equation}
Here $Z''_{j+1}$ is the region so denoted in \cite{B-CMP122a}, $X$ ranges over domains of the
class \cite[(1.63)]{B-CMP122a}, and every such $X$ meets
$R_n\cap Z^{s}$. In particular no $X$ is anchored on $\bigcup_{W\in\mathfrak{W}}W$, although a
hull may enter some $Z_W^{+}$ and depend there on $(\mathbb{U},\mathbb{J})$.

Each $\mathbb{C}^{s}_G(X)$ has the following properties.
\begin{itemize}
\item It is holomorphic on $\mathfrak{D}^{\flat}_{n+1}(X)$, hence on the ambient subdomain of
$\mathfrak{D}^{\flat}_{n+1}(X)$.
\item It is gauge invariant.
\item It depends on $(\mathbb{U},\mathbb{J})$ inside $X$ only. It depends on no integration
variable $\phi$, on no exterior variable $\phi_W$, on no stored term, on no exterior field
$A_{j'}$ at any level $j'$ and on no shift.
\end{itemize}
At any rates with $a_m\le 4\kappa L$,
\begin{equation*}
\|\mathbb{C}^{s}_G\|_{a}\ \le\ c_G ,
\end{equation*}
with the same constant $c_G=c_G(d,L,M)$ as in the normalisation lemma at the full stage.

Two further bounds hold under additional hypotheses.
\begin{itemize}
\item Assume in addition the norm-conversion proposition. This is
the conversion, for the family $\mathbb{C}^{s}_G$, of a bound in $\|\cdot\|_{a}$ into a bound in
$\|\cdot\|_{(n+1)}$, as it is carried out for the stage densities. Then
$\|\mathbb{C}^{s}_G\|_{(n+1)}\le 3c_G$.
\item Assume in addition the holonomy bound. Then
\begin{equation*}
\bigl\|\mathbb{C}^{s}_G-\mathbb{C}^{s}_G\big|_{(\mathbb{U},\mathbb{J})=(1,0)}\bigr\|
\ \le\ \vartheta .
\end{equation*}
\end{itemize}

\smallskip\noindent
(ii) \emph{The cluster expansion at the block.} Let $\langle\cdot\rangle_s$ be the normalised
centred Gaussian expectation in $\phi_s$ with precision $\mathcal{A}_{ss}$, that is, with
covariance $\mathcal{C}_s$. Put
\begin{equation*}
\chi^{s}:=\prod_{c\subset Z^{s}}\chi^{(j)}_c\,\chi'^{(j)}_c ,
\end{equation*}
and take the cubes $\mathsf{Q}^{s}$ and the element list $\mathfrak{E}^{\mathrm{StB}}$ with data
$(\mathsf{b}(e)\cap\mathsf{Q}^{s},\mathsf{d}(e),\nu^{\natural}(e))$. The exterior variables
$\phi_W$, $W\in\mathfrak{W}$, are held in their boxes. Then the
following hold.
\begin{itemize}
\item[(a)] We have
\begin{equation*}
\log\Bigl\langle\chi^{s}\exp\sum_{e\in\mathfrak{E}^{\mathrm{StB}}}v_e\Bigr\rangle_s
=\sum_X\mathbb{C}^{s}_I(X),
\end{equation*}
where $v_e$ is the function carried by the element $e$ as in
Definition~\ref{4.24:def-cluster-weights}, and $X$ ranges over domains of the class
\cite[(1.63)]{B-CMP122a} such that
$X\cap R_n\cap Z^{s}\neq\emptyset$; in particular no $X$ is anchored on
$\bigcup_{W\in\mathfrak{W}}W$. Each $\mathbb{C}^{s}_I(X)$ is holomorphic on the ambient
subdomain of $\mathfrak{D}^{\flat}_{n+1}(X)$, and local in $X$. It is continuous in the
coordinates of $\lambda$ that range over boxes and holomorphic in those that range over
polydiscs.
\item[(b)] The densities satisfy
\begin{equation*}
\begin{split}
\|\mathbb{C}^{s}_I\|_{(n+1)}
&\le 4\bigl(\mathsf{M}^{\mathrm{StB}}_*+2^{d}V\tilde{\mathfrak{l}}\bigr)\\
&\le 4\bigl(\mathsf{M}^{s}_*+\vartheta_{\mathrm{cv}}+\vartheta_a+2^{d}V\tilde{\mathfrak{l}}\bigr).
\end{split}
\end{equation*}
The totals are computed with the weights $\tilde\nu'$, which cover the unprimed total $\mu$ of
the abstract cluster-expansion proposition and the primed total $\mu'$ of $(Q^{+})'$ at once.
They obey the bounds \eqref{4.24:eq-uniform-totals-a},
\begin{equation*}
\mu^{\mathrm{StB}}\le\mu'+\vartheta_a+\vartheta_{\mathrm{cv}},\qquad
\mathsf{M}^{\mathrm{StB}}_*\le\mathsf{M}'_*+\vartheta_a+\vartheta_{\mathrm{cv}},
\end{equation*}
and $\mu^{\mathrm{StB}}+\tilde{\mathfrak{l}}\le\frac12$ under $(Q^{+})'$. The constants
$\gamma_2$, $c_F$, $c_\flat$, $c_J$, $\tilde{\mathfrak{l}}$, $V$, $K(d)$ and the numerical
factors $4$, $8$ and $32$ are the same as in the fluctuation cluster expansion at the full stage.
\item[(c)] Clause (iii) of the abstract cluster-expansion proposition holds for this expansion
word for word. That clause concerns the dependence on the dial variable $\mathfrak{z}$.
\item[(d)] Consequently clause (F2) of the stage bound holds unchanged, with the same symbols:
\begin{equation*}
3c_G+\vartheta+\tfrac12+5(\vartheta_{\mathrm{cv}}+\vartheta_a)\ \le\ 3c_G+1\ =\ C_0^{\sharp}.
\end{equation*}
The inequality holds by condition (iv) on the change-of-variables parameters in
\textup{(h1)}.
\end{itemize}

\smallskip\noindent
(iii) \emph{The parameter dependence at the block.} For each $W\in\mathfrak{W}$ let
$\mathfrak{X}_W$ be the product of two factors:
\begin{itemize}
\item the set of fields $A_j$ on the rim cubes of $W$, each slot ranging over the box or
polydisc that Definition~\ref{4.24:def-stored-terms} assigns to it;
\item the set of real $x$ with $|x|\le 2.4\,\mathfrak{b}'_j$ on the set $S^{P}_W$ of slots of
$W$ that carry exterior real variables in the sense of
Definition~\ref{4.24:def-exterior-variables}.
\end{itemize}
Let the parameter $\lambda$, which contains $(\phi_W)_{W\in\mathfrak{W}}$, range over $D_e$ with
its exterior components in $\prod_{W\in\mathfrak{W}}\mathfrak{X}_W$. Then every member
$e$ of $\mathfrak{E}^{\mathrm{StB}}$ has on $D_e$ a majorant $\nu^{\natural}(e)$ that is free
of $\lambda$ and of $(\mathbb{U},\mathbb{J})$. No condition beyond the arrest condition of
Definition~\ref{4.24:def-survivors-bonds} is used.
Write $\phi_{\mathfrak{W}}:=(\phi_W)_{W\in\mathfrak{W}}$, and write
$\mathbb{C}^{s}_I(X,\mathcal{S};\phi_{\mathfrak{W}})$ for $\mathbb{C}^{s}_I(X)$ with its
dependence on a strong set $\mathcal{S}$ and on $\phi_{\mathfrak{W}}$ displayed. The corollary
of the parameter lemma places $\mathbb{C}^{s}_I(X,\mathcal{S};\phi_{\mathfrak{W}})$ on the
product of the ambient subdomain of $\mathfrak{D}^{\flat}_{n+1}(X)$ with
\begin{equation*}
\mathfrak{D}_{\mathcal{S}}(X)\times\prod_{W\in\mathfrak{W}}\mathfrak{X}_W .
\end{equation*}
This is the domain of clause (F1) of the stage bound, and the polydisc of the $A$-variables is
never shrunk.

\smallskip\noindent
(iv) \emph{Uniformity.} No constant and no condition in (i)--(iii) involves any of the
following:
\begin{itemize}
\item $|\mathfrak{W}|$, $n^{+}$, the number of components of $Z$ or of $Z^{s}$;
\item $N$, $N_0$, $k$, or $\check{R}$ other than through the component budget $\bar V$ already
present;
\item the age of any region, or the volume.
\end{itemize}
No condition bounding $g$ from below is used, no new lower bound on any parameter is introduced,
and the order in which the parameters are chosen is unchanged.
\end{theorem}

In the second of the two further bounds of clause~(i), each factor entering the field
dependence of $\mathbb{C}^{s}_G$, taken in the gauge of its own cube, is bounded as in
\eqref{4.24:eq-field-dependence-a-bis}, with $\theta$ and $\epsilon$ as
in \eqref{4.24:eq-field-dependence-b}; that bound rests on the short-path connectedness clause
of the index-zero rigidity lemma. Hence the field-dependent part of
$\mathbb{C}^{s}_G$ is at most
\begin{equation*}
c_G'\cdot O(1)\cdot\bigl[\,dCM\cdot 5w_*\bigl(BCM\tilde\alpha_0+\tilde\alpha_1\bigr)
+K_\Delta\gamma_0+\mathfrak{h}\,\bigr],
\end{equation*}
where $\mathfrak{h}$ is the holonomy term of the holonomy bound, $dCM$ is the factor
displayed in $\theta$ in \eqref{4.24:eq-field-dependence-b}, and $B$ and $C$ are the constants
of the corresponding bracket for the stage densities. The holonomy bound is the closing step
that turns this into the bound $\le\vartheta$; the normalisation lemma at the full stage takes
the same step.

The proof of Theorem~\ref{4.24:thm-survivor-block-expansion} is given in the lemmas that follow
in this section. Clause~(i) is proved first on the whole survivors' region and then a second
time component by component. Clauses~(ii), (iii) and~(iv) are proved in turn after it.

\begin{lemma}[Principal blocks: what restricts and what does not]\label{4.24:lem-principal-blocks}
We use the spaces $\mathcal{H}$, $\mathcal{H}_E$, the pairings $(\cdot,\cdot)$ and $\langle\cdot,\cdot\rangle$,
the norm $\|\cdot\|$, the coordinate projections $P_E$ and the gauge conjugation $\mathcal{U}_u$ of
Notation~\ref{4.24:not-fibres-pairings}. Norms of operators are operator norms for $\|\cdot\|$.

Let $\mathcal{A}$ be the symmetrised stage kernel of Definition~\ref{4.24:def-stage-form}. Let
$\mathfrak{Bo}=s\sqcup\mathcal{W}$ be the split into survivors' and arrested bonds. For a kernel $K$
and $E,E'\subset\mathfrak{Bo}$ write $K_{EE'}$ for $P_EKP_{E'}$, regarded as an operator from
$\mathcal{H}_{E'}$ to $\mathcal{H}_E$. In particular $\mathcal{A}_{ss}=P_s\mathcal{A}P_s$ on
$\mathcal{H}_s$.
\begin{enumerate}
\item[(a)] For $b,b'\in s$ one has $\mathcal{A}_{ss}(b,b')=\mathcal{A}(b,b')$.
 \begin{itemize}
 \item Hence holomorphy, the kernel bound and component-locality, which hold for $\mathcal{A}$
 (Definition~\ref{4.24:def-stage-form}), hold for $\mathcal{A}_{ss}$ verbatim.
 \item Gauge covariance also holds: for every gauge transformation $u$, with
 $(\mathbb{U}^u,\operatorname{Ad}_u\mathbb{J})$ the gauge-transformed background pair,
 \begin{equation*}
 \mathcal{A}_{ss}[\mathbb{U}^u,\operatorname{Ad}_u\mathbb{J}]=\mathcal{U}_u\mathcal{A}_{ss}\mathcal{U}_u^{-1}.
 \end{equation*}
 \item $(\mathcal{A}_{ss})^T=P_s\mathcal{A}^TP_s=\mathcal{A}_{ss}$.
 \item For $\xi\in\mathbb{C}$, $(\mathcal{A}+\xi)_{ss}=\mathcal{A}_{ss}+\xi 1_{\mathcal{H}_s}$.
 \item Compression is transitive: for $s\subset R\subset\mathfrak{Bo}$ and every kernel $K$,
 $P_s(P_RKP_R)P_s=P_sKP_s$.
 \item For every kernel $K$ and all $E,E'$, $\|K_{EE'}\|\le\|K\|$. In particular, for the values
 $\mathcal{A}(z),\mathcal{A}(z')$ of the kernel at two values $z,z'$ of the parameter on which it
 depends holomorphically (Definition~\ref{4.24:def-stage-form}),
 $\|(\mathcal{A}(z)-\mathcal{A}(z'))_{ss}\|\le\|\mathcal{A}(z)-\mathcal{A}(z')\|$.
 \end{itemize}
\item[(b)] Let $\phi\in\mathcal{H}_s$. Let $B$ be an operator on a finite-dimensional subspace of
 $\mathcal{H}$ (for instance some $\mathcal{H}_E$) with
 $\operatorname{Re}\langle\psi,B\psi\rangle\ge a\|\psi\|^2$ for all $\psi$ in that subspace, for some
 $a>0$. Let $\xi\ge0$. In addition, consider the instance with index set $\{1,2\}$ and scalar fibres,
 in which the letters $s=\{1\}$ and $\mathcal{W}=\{2\}$ are reused for this instance only: let
 $t\in\mathbb{C}$ with $0<|t|<1$, and let $\mathcal{A}^{(t)}$ be the matrix
 $\bigl(\begin{smallmatrix}1&t\\ t&1\end{smallmatrix}\bigr)$ on $\mathbb{C}^2$. Then
 \begin{equation}\label{4.24:eq-principal-blocks-a}
 \begin{gathered}
 \langle\phi,\mathcal{A}_{ss}\phi\rangle=\langle\phi,\mathcal{A}\phi\rangle,\\
 B\ \text{is bijective},\quad \|B^{-1}\|\le a^{-1},\quad
 \|(B+\xi)^{-1}\|\le(a+\xi)^{-1},\\
 \operatorname{Re}\zeta\ge a\ \text{for every eigenvalue } \zeta \text{ of } B,\\
 (\mathcal{A}^{(t)}_{ss})^{-1}=1\ne(1-t^2)^{-1}=((\mathcal{A}^{(t)})^{-1})_{ss},\\
 \log\det\mathcal{A}^{(t)}_{ss}=0\ne\log(1-t^2)
 =\log\det\mathcal{A}^{(t)}-\log\det\mathcal{A}^{(t)}_{\mathcal{W}\mathcal{W}}.
 \end{gathered}
 \end{equation}
 Here $\log$ is the principal branch.
\end{enumerate}
By the first line of \eqref{4.24:eq-principal-blocks-a}, the Rayleigh quotients (the numerical range)
of $\mathcal{A}_{ss}$ are among those of $\mathcal{A}$.

By the last two lines, the inverse and the logarithm of the determinant of a block are formed of the
block itself. In general they differ from the compression $(\mathcal{A}^{-1})_{ss}$ of the inverse and
from the difference $\log\det\mathcal{A}-\log\det\mathcal{A}_{\mathcal{W}\mathcal{W}}$, as the last two
lines of \eqref{4.24:eq-principal-blocks-a} show; they are always to be formed from $\mathcal{A}_{ss}$
itself.
\end{lemma}

\begin{proof}
(a) For $b\in E$ and $b'\in E'$, the entry of $P_EKP_{E'}$ at $(b,b')$ is $K(b,b')$, since $P_E$ and
$P_{E'}$ are the coordinate projections. Taking $E=E'=s$ and $K=\mathcal{A}$ gives the entry identity.
Holomorphy, the kernel bound and component-locality are properties of the entries $\mathcal{A}(b,b')$.
For $b,b'\in s$ these entries are the entries of $\mathcal{A}_{ss}$, so the three properties pass to
$\mathcal{A}_{ss}$ unchanged.

For gauge covariance we use $P_s\mathcal{U}_u=\mathcal{U}_uP_s$. This gives
$\mathcal{U}_u^{-1}P_s=P_s\mathcal{U}_u^{-1}$, and therefore, using the gauge covariance of
$\mathcal{A}$ at the first equality,
\begin{equation*}
\mathcal{A}_{ss}[\mathbb{U}^u,\operatorname{Ad}_u\mathbb{J}]
=P_s\,\mathcal{U}_u\mathcal{A}\,\mathcal{U}_u^{-1}P_s
=\mathcal{U}_u\,P_s\mathcal{A}P_s\,\mathcal{U}_u^{-1}
=\mathcal{U}_u\mathcal{A}_{ss}\mathcal{U}_u^{-1}.
\end{equation*}

The transpose with respect to $(\cdot,\cdot)$ reverses products. Moreover $P_s^T=P_s$, because $P_s$
is $(\cdot,\cdot)$-symmetric. Hence $(P_s\mathcal{A}P_s)^T=P_s\mathcal{A}^TP_s$. Since $\mathcal{A}$ is
a symmetrisation $\tfrac12(K+K^T)$, one has $\mathcal{A}^T=\mathcal{A}$, and so
$P_s\mathcal{A}^TP_s=\mathcal{A}_{ss}$.

Next, $P_s(\mathcal{A}+\xi)P_s=P_s\mathcal{A}P_s+\xi P_s^2$. On $\mathcal{H}_s$ the operator $P_s^2=P_s$
is the identity $1_{\mathcal{H}_s}$.

For $s\subset R$, the product rule $P_EP_{E'}=P_{E\cap E'}$ gives $P_sP_R=P_{s\cap R}=P_s$ and
$P_RP_s=P_s$. Therefore $P_sP_RKP_RP_s=P_sKP_s$.

Since $\|P_E\|\le1$, we get
\begin{equation*}
\|K_{EE'}\|\le\|P_E\|\,\|K\|\,\|P_{E'}\|\le\|K\|.
\end{equation*}
Compression is linear in $K$, so
$(\mathcal{A}(z)-\mathcal{A}(z'))_{ss}=\mathcal{A}(z)_{ss}-\mathcal{A}(z')_{ss}$. The last assertion of
(a) is the case $E=E'=s$, $K=\mathcal{A}(z)-\mathcal{A}(z')$.

(b) Rayleigh identity. $P_s$ is $\langle\cdot,\cdot\rangle$-orthogonal, hence self-adjoint, and
$P_s\phi=\phi$ for $\phi\in\mathcal{H}_s$. Therefore
\begin{equation*}
\langle\phi,\mathcal{A}_{ss}\phi\rangle=\langle\phi,P_s\mathcal{A}P_s\phi\rangle
=\langle P_s\phi,\mathcal{A}P_s\phi\rangle=\langle\phi,\mathcal{A}\phi\rangle .
\end{equation*}
Taking $\|\phi\|=1$ shows that every Rayleigh quotient of $\mathcal{A}_{ss}$ is one of $\mathcal{A}$.

Lax--Milgram bounds. By the Cauchy--Schwarz inequality and the hypothesis on $B$,
\begin{equation*}
\|B\psi\|\,\|\psi\|\ge|\langle\psi,B\psi\rangle|\ge\operatorname{Re}\langle\psi,B\psi\rangle
\ge a\|\psi\|^2 ,
\end{equation*}
so $\|B\psi\|\ge a\|\psi\|$. Thus $B$ is injective. Since the space on which $B$ acts is
finite-dimensional, $B$ is bijective. Putting $\psi=B^{-1}\chi$ gives
$\|B^{-1}\chi\|\le a^{-1}\|\chi\|$, that is, $\|B^{-1}\|\le a^{-1}$.

If $B\psi=\zeta\psi$ with $\|\psi\|=1$, then $\langle\psi,B\psi\rangle=\zeta$, since
$\langle\cdot,\cdot\rangle$ is linear in the second slot. Hence
$\operatorname{Re}\zeta=\operatorname{Re}\langle\psi,B\psi\rangle\ge a$.

For $\xi\ge0$, the reality of $\xi$ gives
\begin{equation*}
\operatorname{Re}\langle\psi,(B+\xi)\psi\rangle=\operatorname{Re}\langle\psi,B\psi\rangle+\xi\|\psi\|^2
\ge(a+\xi)\|\psi\|^2 .
\end{equation*}
Applying the bound just proved to $B+\xi$, with $a+\xi$ in place of $a$, gives
$\|(B+\xi)^{-1}\|\le(a+\xi)^{-1}$.

Schur witness. We have $\mathcal{A}^{(t)}_{ss}=1$ and $\mathcal{A}^{(t)}_{\mathcal{W}\mathcal{W}}=1$.
Also $\det\mathcal{A}^{(t)}=1-t^2\ne0$, and
\begin{equation*}
(\mathcal{A}^{(t)})^{-1}=(1-t^2)^{-1}\begin{pmatrix}1&-t\\-t&1\end{pmatrix}.
\end{equation*}
Hence $(\mathcal{A}^{(t)}_{ss})^{-1}=1$ and $((\mathcal{A}^{(t)})^{-1})_{ss}=(1-t^2)^{-1}$. These differ
because $t^2\ne0$.

Likewise $\log\det\mathcal{A}^{(t)}_{ss}=\log1=0$. On the other side,
$\log\det\mathcal{A}^{(t)}-\log\det\mathcal{A}^{(t)}_{\mathcal{W}\mathcal{W}}=\log(1-t^2)-0$.
Since $|t^2|<1$, the number $1-t^2$ has positive real part, so its principal logarithm is defined.
That logarithm is non-zero because $1-t^2\ne1$.

So neither the inverse nor the logarithm of the determinant of the block $s$ is obtained by
restricting the corresponding object of $\mathcal{A}^{(t)}$. Both are formed of
$\mathcal{A}^{(t)}_{ss}$.
\end{proof}

\begin{lemma}[Coercivity of the survivors' block]\label{4.24:lem-block-coercivity}
Let $\mathcal{H}=\bigoplus_{b\in\mathfrak{Bo}}\mathfrak{g}^{\mathbb{C}}$ be the bond Hilbert space of the stage, $s\subset\mathfrak{Bo}$ the survivors' bonds, and $P_s$ and $\mathcal{H}_s=P_s\mathcal{H}$ the coordinate projection and subspace. For a point $z=(\mathbb{U},\mathbb{J})$ of the domain of complex background pairs on which the stage form is defined, let $\mathcal{A}(z)$ be the symmetrised stage operator on $\mathcal{H}$ (Definition~\ref{4.24:def-stage-form}), and let $\mathcal{A}_{ss}(z)=P_s\mathcal{A}(z)P_s$ be its principal block on $\mathcal{H}_s$ (Lemma~\ref{4.24:lem-principal-blocks}). We write $(\cdot,\cdot)$ for the bilinear extension to $\mathcal{H}$ of the real inner product of $\bigoplus_{b}\mathfrak{g}$, $\langle\varphi,\psi\rangle=(\bar\varphi,\psi)$ for the Hermitian inner product, $\|\cdot\|$ for its norm and for the operator norm, and $\lambda_{\min}$ for the least eigenvalue of a real symmetric operator. A \emph{real point} is a point $z_{\mathbb{R}}$ of this domain whose first entry is $G$-valued and whose second entry is $\mathfrak{g}$-valued. Let $B_0=B_0(d,L)$ be the constant, fixed before $\gamma$, of the covariance bound of \cite[Theorem~3.15, (3.187)]{B-CMP99b}, and let $\alpha_5$ be the smallness parameter of \cite[(2.23)]{B-CMP116}; it is of order $\tilde\alpha_0+\tilde\alpha_1+e^{-\delta_0M/3}$, the complex background entering through the first two terms. Assume:
\begin{itemize}
\item[(a)] at every real point $z_{\mathbb{R}}$ the operator $\mathcal{A}(z_{\mathbb{R}})$ maps $\bigoplus_{b\in\mathfrak{Bo}}\mathfrak{g}$ into itself, is $(\cdot,\cdot)$-symmetric and positive definite there (it is the precision of the convergent fluctuation Gaussian of Ba{\l}aban's step at $z_{\mathbb{R}}$), and its inverse $\mathcal{C}(z_{\mathbb{R}})$ satisfies $\|\mathcal{C}(z_{\mathbb{R}})\|\le 2B_0$, the form in which \cite[Theorem~3.15, (3.187)]{B-CMP99b} is used here;
\item[(b)] the Hermitian-part coercivity of the stage form: for every $z$ in the domain and every $\varphi\in\mathcal{H}$,
\begin{equation}\label{4.24:eq-block-coercivity-1}
\operatorname{Re}\langle\varphi,\mathcal{A}(z)\varphi\rangle\ge\bigl((2B_0)^{-1}-\alpha_5\bigr)\|\varphi\|^2 ;
\end{equation}
\item[(c)] $2B_0\,\alpha_5<\tfrac12$, that is $\alpha_5<(4B_0)^{-1}$; this is the form, through the bound $\|\mathcal{C}(z_{\mathbb{R}})\|\le 2B_0$ of (a), of the condition $\alpha_5\|C^{(k)}(Z_0,0)\|<\tfrac12$ on Ba{\l}aban's covariance $C^{(k)}(Z_0,0)$ assumed in \cite[p.~17, the sentence following (2.24)]{B-CMP116}; since $B_0$ depends on $d$ and $L$ only, it is a smallness condition on $\alpha_5$ imposed after $B_0$ is fixed.
\end{itemize}
Put
\begin{equation}\label{4.24:eq-block-coercivity-a}
a_{\mathbb{C}}:=(4B_0)^{-1}.
\end{equation}
Then for every $z$ in the domain:
\begin{itemize}
\item[(i)] $\operatorname{Re}\langle\varphi,\mathcal{A}_{ss}(z)\varphi\rangle\ge a_{\mathbb{C}}\|\varphi\|^2$ for all $\varphi\in\mathcal{H}_s$;
\item[(ii)] the numerical range and the spectrum of $\mathcal{A}_{ss}(z)$ lie in $\{\zeta\in\mathbb{C}:\operatorname{Re}\zeta\ge a_{\mathbb{C}}\}$;
\item[(iii)] for every $\xi\ge0$ the operator $\mathcal{A}_{ss}(z)+\xi$ on $\mathcal{H}_s$ is invertible, with
$\|(\mathcal{A}_{ss}(z)+\xi)^{-1}\|\le(a_{\mathbb{C}}+\xi)^{-1}$;
\item[(iv)] the survivors' covariance $\mathcal{C}_s(z)=\mathcal{A}_{ss}(z)^{-1}$ (Definition~\ref{4.24:def-stage-form}) exists and $\|\mathcal{C}_s(z)\|\le a_{\mathbb{C}}^{-1}=4B_0$;
\item[(v)] at every real point,
\begin{equation}\label{4.24:eq-block-coercivity-2}
\lambda_{\min}\bigl(\mathcal{A}_{ss}(z_{\mathbb{R}})\bigr)\ge\lambda_{\min}\bigl(\mathcal{A}(z_{\mathbb{R}})\bigr)
=\|\mathcal{C}(z_{\mathbb{R}})\|^{-1}\ge(2B_0)^{-1}.
\end{equation}
\end{itemize}
\end{lemma}

\begin{proof}
Fix $z$ and write $\mathcal{A}=\mathcal{A}(z)$, $\mathcal{A}_{ss}=\mathcal{A}_{ss}(z)$.

\emph{The bound on $\mathcal{H}$.} By (c), $\alpha_5<(4B_0)^{-1}$, hence
\begin{equation*}
(2B_0)^{-1}-\alpha_5>(2B_0)^{-1}-(4B_0)^{-1}=(4B_0)^{-1}=a_{\mathbb{C}} .
\end{equation*}
Inserting this into \eqref{4.24:eq-block-coercivity-1} gives, pointwise in $\varphi\in\mathcal{H}$,
\begin{equation}\label{4.24:eq-block-coercivity-3}
\operatorname{Re}\langle\varphi,\mathcal{A}\varphi\rangle\ge a_{\mathbb{C}}\|\varphi\|^2 .
\end{equation}
If the lower bound (b) is given for the stage kernel $K_{\sharp}(z)$ of Definition~\ref{4.24:def-stage-form} rather than for $\mathcal{A}=\mathcal{S}K_{\sharp}(z)$, clause (f) of Lemma~\ref{4.24:lem-kernel-polarisation} transfers it to $\mathcal{A}$ with the same constant.

\emph{Proof of} (i). Let $\varphi\in\mathcal{H}_s$. By the Rayleigh identity of Lemma~\ref{4.24:lem-principal-blocks}, $\langle\varphi,\mathcal{A}_{ss}\varphi\rangle=\langle P_s\varphi,\mathcal{A}P_s\varphi\rangle=\langle\varphi,\mathcal{A}\varphi\rangle$. Since $\mathcal{H}_s\subset\mathcal{H}$, \eqref{4.24:eq-block-coercivity-3} applies to $\varphi$ and gives (i). Only the inclusion $\mathcal{H}_s\subset\mathcal{H}$ was used; hence the constant in (i) is the same $a_{\mathbb{C}}$ whatever the ambient bond set on which $\mathcal{A}$ and hypotheses (a)--(c) are given --- the bonds of the run or any larger set of bonds --- and whatever the subset $s$.

\emph{Proof of} (ii). The numerical range is $\{\langle\varphi,\mathcal{A}_{ss}\varphi\rangle:\varphi\in\mathcal{H}_s,\ \|\varphi\|=1\}$, and by (i) every such number has real part at least $a_{\mathbb{C}}$. The space $\mathcal{H}_s$ is finite-dimensional, so the spectrum of $\mathcal{A}_{ss}$ consists of eigenvalues. If $\mathcal{A}_{ss}\varphi=\lambda\varphi$ with $\|\varphi\|=1$, then, $\langle\cdot,\cdot\rangle$ being linear in its second argument, $\langle\varphi,\mathcal{A}_{ss}\varphi\rangle=\lambda\|\varphi\|^2=\lambda$, so $\lambda$ lies in the numerical range, and $\operatorname{Re}\lambda\ge a_{\mathbb{C}}$.

\emph{Proof of} (iii). Let $\xi\ge0$. By Lemma~\ref{4.24:lem-principal-blocks}, $(\mathcal{A}+\xi)_{ss}=\mathcal{A}_{ss}+\xi 1_{\mathcal{H}_s}$. For $\varphi\in\mathcal{H}_s$, the Cauchy--Schwarz inequality and (i) give
\begin{align*}
\|(\mathcal{A}_{ss}+\xi)\varphi\|\,\|\varphi\|
&\ge\bigl|\langle\varphi,(\mathcal{A}_{ss}+\xi)\varphi\rangle\bigr|
\ge\operatorname{Re}\langle\varphi,(\mathcal{A}_{ss}+\xi)\varphi\rangle\\
&=\operatorname{Re}\langle\varphi,\mathcal{A}_{ss}\varphi\rangle+\xi\|\varphi\|^2
\ge(a_{\mathbb{C}}+\xi)\|\varphi\|^2 ,
\end{align*}
hence $\|(\mathcal{A}_{ss}+\xi)\varphi\|\ge(a_{\mathbb{C}}+\xi)\|\varphi\|$. Since $a_{\mathbb{C}}+\xi>0$, the operator $\mathcal{A}_{ss}+\xi$ is injective on the finite-dimensional space $\mathcal{H}_s$, hence invertible. For $\psi\in\mathcal{H}_s$ put $\varphi=(\mathcal{A}_{ss}+\xi)^{-1}\psi$; the last inequality reads $\|\psi\|\ge(a_{\mathbb{C}}+\xi)\|(\mathcal{A}_{ss}+\xi)^{-1}\psi\|$, which is the asserted norm bound. This is the Lax--Milgram argument in finite dimension.

\emph{Proof of} (iv). This is (iii) at $\xi=0$, with $a_{\mathbb{C}}^{-1}=4B_0$ by \eqref{4.24:eq-block-coercivity-a}.

\emph{Proof of} (v). Let $z_{\mathbb{R}}$ be a real point. By Lemma~\ref{4.24:lem-principal-blocks}, $\mathcal{A}_{ss}(z_{\mathbb{R}})(b,b')=\mathcal{A}(z_{\mathbb{R}})(b,b')$ for $b,b'\in s$, so by (a) the operator $\mathcal{A}_{ss}(z_{\mathbb{R}})$ maps $\bigoplus_{b\in s}\mathfrak{g}$ into itself; by the same lemma $(\mathcal{A}_{ss})^T=\mathcal{A}_{ss}$, so it is $(\cdot,\cdot)$-symmetric there. The projection $P_s$ is real and $(\cdot,\cdot)$-symmetric, so $(\varphi,\mathcal{A}_{ss}(z_{\mathbb{R}})\varphi)=(\varphi,\mathcal{A}(z_{\mathbb{R}})\varphi)$ for real $\varphi\in\bigoplus_{b\in s}\mathfrak{g}$. By the variational characterisation of the least eigenvalue,
\begin{align*}
\lambda_{\min}\bigl(\mathcal{A}_{ss}(z_{\mathbb{R}})\bigr)
&=\min\Bigl\{(\varphi,\mathcal{A}(z_{\mathbb{R}})\varphi):\varphi\in\textstyle\bigoplus_{b\in s}\mathfrak{g},\ \|\varphi\|=1\Bigr\}\\
&\ge\min\Bigl\{(\varphi,\mathcal{A}(z_{\mathbb{R}})\varphi):\varphi\in\textstyle\bigoplus_{b\in\mathfrak{Bo}}\mathfrak{g},\ \|\varphi\|=1\Bigr\}
=\lambda_{\min}\bigl(\mathcal{A}(z_{\mathbb{R}})\bigr),
\end{align*}
the minimum over a subspace being at least the minimum over the whole space. By (a) and the spectral theorem, $\mathcal{A}(z_{\mathbb{R}})$ has positive eigenvalues, $\mathcal{C}(z_{\mathbb{R}})$ has their reciprocals as eigenvalues, and the norm of the real symmetric operator $\mathcal{C}(z_{\mathbb{R}})$, on the real space or on its complexification $\mathcal{H}$, is its largest eigenvalue $\lambda_{\min}(\mathcal{A}(z_{\mathbb{R}}))^{-1}$. Hence $\lambda_{\min}(\mathcal{A}(z_{\mathbb{R}}))=\|\mathcal{C}(z_{\mathbb{R}})\|^{-1}\ge(2B_0)^{-1}$, the inequality being (a). This proves \eqref{4.24:eq-block-coercivity-2}.

By \eqref{4.24:eq-block-coercivity-a}, the constant $a_{\mathbb{C}}$ depends on $d$ and $L$ only, through $B_0$; it does not enter the normalisation constant $c_G$ or the constant $C_0^{\sharp}=3c_G+1$.
\end{proof}

\begin{remark}[A sufficient condition for hypothesis (b)]
What Lemma~\ref{4.24:lem-block-coercivity} uses of hypothesis (b) is the lower bound on the Hermitian part; no inequality of the form $\operatorname{Re}A\ge\|A^{-1}\|^{-1}$ for a non-normal complex operator $A$ is used. The bound \eqref{4.24:eq-block-coercivity-1} is the one needed to absorb factors $e^{+\frac12\gamma_2|\phi(b)|^2}$ into a complex Gaussian. One way in which (b) arises is the following. Suppose every point $z$ of the domain has a real point $z_{\mathbb{R}}$ with
\begin{equation*}
\|\mathcal{A}(z)-\mathcal{A}(z_{\mathbb{R}})\|\le\alpha_5 ,
\end{equation*}
the complex deformations admitted in the domain being small perturbations of a real regular configuration. For $\varphi\in\mathcal{H}$ write $\varphi=\operatorname{Re}\varphi+i\operatorname{Im}\varphi$. Then
\begin{align*}
\langle\varphi,\mathcal{A}(z_{\mathbb{R}})\varphi\rangle
&=(\operatorname{Re}\varphi,\mathcal{A}(z_{\mathbb{R}})\operatorname{Re}\varphi)
+(\operatorname{Im}\varphi,\mathcal{A}(z_{\mathbb{R}})\operatorname{Im}\varphi)\\
&\quad+i\bigl[(\operatorname{Re}\varphi,\mathcal{A}(z_{\mathbb{R}})\operatorname{Im}\varphi)
-(\operatorname{Im}\varphi,\mathcal{A}(z_{\mathbb{R}})\operatorname{Re}\varphi)\bigr],
\end{align*}
and the bracket vanishes by the $(\cdot,\cdot)$-symmetry of $\mathcal{A}(z_{\mathbb{R}})$. The remaining two terms are at least $\lambda_{\min}(\mathcal{A}(z_{\mathbb{R}}))(\|\operatorname{Re}\varphi\|^2+\|\operatorname{Im}\varphi\|^2)$, which equals $\lambda_{\min}(\mathcal{A}(z_{\mathbb{R}}))\|\varphi\|^2$ and is at least $(2B_0)^{-1}\|\varphi\|^2$ by (a). Since
$|\langle\varphi,(\mathcal{A}(z)-\mathcal{A}(z_{\mathbb{R}}))\varphi\rangle|\le\alpha_5\|\varphi\|^2$,
\begin{equation*}
\operatorname{Re}\langle\varphi,\mathcal{A}(z)\varphi\rangle
\ge\langle\varphi,\mathcal{A}(z_{\mathbb{R}})\varphi\rangle-\alpha_5\|\varphi\|^2
\ge\bigl((2B_0)^{-1}-\alpha_5\bigr)\|\varphi\|^2 ,
\end{equation*}
which is \eqref{4.24:eq-block-coercivity-1}. Here $z_{\mathbb{R}}$ is the pair's own real part: the $G$-valued factor $U$ in Ba{\l}aban's factorisation of $\mathbb{U}$ in \cite{B-CMP122a}, and the $\mathfrak{g}$-part of $\mathbb{J}$. With this choice $z$ and $z_{\mathbb{R}}$ transform together under $G$-valued gauge transformations $u$, and since $\mathcal{U}_u$ is then unitary,
\begin{equation*}
\|\mathcal{A}(z)-\mathcal{A}(z_{\mathbb{R}})\|
=\|\mathcal{U}_u(\mathcal{A}(z)-\mathcal{A}(z_{\mathbb{R}}))\mathcal{U}_u^{-1}\| .
\end{equation*}
The quantity is therefore gauge-invariant; it measures the complex deformation only, not the distance of $U$ from $1$, and it is not the quantity $\|\mathcal{A}(\mathbb{U})-\mathcal{A}(1)\|$. This sufficient condition is not used in the proof of Lemma~\ref{4.24:lem-block-coercivity}.
\end{remark}

\begin{remark}[The block under a decomposition into components]
Lemma~\ref{4.24:lem-block-coercivity} does not use the following. Suppose $s$ is the disjoint union of components $c$ with coordinate projections $P_c$ and subspaces $\mathcal{H}_c=P_c\mathcal{H}$; any two components are at distance at least $R_kM$, and $\mathrm{gap}$ denotes the least such distance. When the stage operator is given component by component, so that the lower bound of (b) is available for each block $P_c\mathcal{A}P_c$ on $\mathcal{H}_c$ rather than for $\mathcal{A}$ on $\mathcal{H}$, coercivity of the survivors' block is recovered with a halved constant, as follows. Write $\operatorname{dist}(b,b')$ for the lattice distance between the bonds $b$ and $b'$. Suppose the entries of $\mathcal{A}$ obey $\|\mathcal{A}(b,b')\|\le B_{\mathcal{A}}e^{-\delta\operatorname{dist}(b,b')}$, with $B_{\mathcal{A}}$ and $C_\rho$ the constants of the decoupled collar kernel and $\delta>0$ its decay rate. Since the components are at mutual distance at least $\mathrm{gap}$, for every $b$ the sum of $e^{-\delta\operatorname{dist}(b,b')}$ over bonds $b'$ in components other than that of $b$ is at most $C_\rho e^{-\delta\,\mathrm{gap}/2}$. Then
\begin{equation*}
P_s\mathcal{A}P_s=\bigoplus_c P_c\mathcal{A}P_c+\mathsf{X},
\end{equation*}
where $\mathsf{X}$ has the entries of $\mathcal{A}$ between bonds of distinct components. By Schur's test, $\|\mathsf{X}\|^2$ is at most the product of the supremum of the row sums and the supremum of the column sums of the entry norms; both are bounded by the same expression, so
\begin{equation*}
\|\mathsf{X}\|\le B_{\mathcal{A}}\sup_b\sum_{b'}e^{-\delta\operatorname{dist}(b,b')}\le B_{\mathcal{A}}C_\rho e^{-\delta\,\mathrm{gap}/2}
\le B_{\mathcal{A}}C_\rho e^{-\delta R_kM/2},
\end{equation*}
the sum running over $b'$ in components other than that of $b$, and the last step using $\mathrm{gap}\ge R_kM$. Suppose now that each per-component block is coercive, $\operatorname{Re}\langle\psi,P_c\mathcal{A}P_c\psi\rangle\ge a_{\mathbb{C}}\|\psi\|^2$ for $\psi\in\mathcal{H}_c$. For $\varphi\in\mathcal{H}_s$, the $P_c$ being orthogonal projections with $\sum_c\|P_c\varphi\|^2=\|\varphi\|^2$,
\begin{align*}
\operatorname{Re}\langle\varphi,P_s\mathcal{A}P_s\varphi\rangle
&=\sum_c\operatorname{Re}\langle P_c\varphi,\mathcal{A}P_c\varphi\rangle
+\operatorname{Re}\langle\varphi,\mathsf{X}\varphi\rangle\\
&\ge\bigl(a_{\mathbb{C}}-\|\mathsf{X}\|\bigr)\|\varphi\|^2 .
\end{align*}
If moreover $B_{\mathcal{A}}C_\rho e^{-\delta R_kM/2}\le(8B_0)^{-1}=\tfrac12a_{\mathbb{C}}$, a smallness condition of the same sup-form as (c), then $a_{\mathbb{C}}-\|\mathsf{X}\|\ge\frac12a_{\mathbb{C}}$.
\end{remark}

\begin{remark}[Combes--Thomas decay of the block covariance. Proof outstanding]
\label{4.24:rem-covariance-decay}
Let $s$ be the set of survivors' bonds, $\Lambda'_s$ the survivors' integration region, and
$\mathcal{C}_s$ the covariance of the survivors' block. For bonds $b,b'$ let $d(b,b')$ be the
distance between them, measured as in Notation~\ref{4.4:not-units-distances}. There is a decay
rate $\mu_{\mathbb{C}}>0$ such that for all bonds $b,b'\in s$
\begin{equation}\label{4.24:eq-covariance-decay-1}
  \bigl|\mathcal{C}_s(b,b')\bigr|\;\le\;\frac{2}{a_{\mathbb{C}}}\,
  e^{-\mu_{\mathbb{C}}\,d(b,b')},
\end{equation}
where the constants $2/a_{\mathbb{C}}$ and $\mu_{\mathbb{C}}$ are independent of $s$. On
$\Lambda'_s$ this bound could replace Ba{\l}aban's covariance bounds in every application except
the walk expansion itself. It is not used in the proof of clause~(i) of
Theorem~\ref{3.1:thm-main}.

Proof outstanding.

\emph{Route.} The bound follows by the argument of \cite{CT73} from two inputs: the coercivity
constant $a_{\mathbb{C}}$ of the survivors' block (Lemma~\ref{4.24:lem-block-coercivity}) and the
exponential bound on the kernel of the form whose inverse is $\mathcal{C}_s$. For a fixed bond
$b_0\in s$ one conjugates that form by the weights $e^{\pm\mu_{\mathbb{C}}\,d(b_0,\cdot)}$ and uses
the kernel bound to keep the conjugated form coercive for $\mu_{\mathbb{C}}$ small enough, depending
only on $a_{\mathbb{C}}$ and on the constants of the kernel bound; the constants in
\eqref{4.24:eq-covariance-decay-1} then depend only on these quantities.
\end{remark}

\begin{lemma}[The logarithm of the block as a half-line resolvent integral]
\label{4.24:lem-log-resolvent}
Let $s$ be the set of survivors' bonds and $\mathcal{H}_s$ the corresponding coordinate subspace of
$\mathcal{H}$. An element of $\mathcal{H}_s$ assigns to every bond $b'\in s$ a vector of the fibre
$\mathfrak{g}^{\mathbb{C}}$, of complex dimension $n_{\mathfrak{g}}$, so that
$\dim_{\mathbb{C}}\mathcal{H}_s=\mathcal{N}:=n_{\mathfrak{g}}|s|$. Write
$\langle\cdot,\cdot\rangle$ for the Hermitian inner product of $\mathcal{H}_s$, $\|\cdot\|$ for
its norm and for the associated operator norm, $(\cdot,\cdot)$ for the complex-bilinear form of
Lemma~\ref{4.24:lem-kernel-polarisation}, $\operatorname{Tr}$ for the trace on
$\operatorname{End}(\mathcal{H}_s)$, $\operatorname{tr}_{\mathfrak{g}}$ for the trace on the
fibre, and $1$ for the identity of $\mathcal{H}_s$. An operator $K$ on $\mathcal{H}_s$ is
given by its kernel $K(b',b'')$, $b',b''\in s$. Put
\begin{equation*}
\mathfrak{P}:=\{B\in\operatorname{End}(\mathcal{H}_s):\ \operatorname{Re}\langle\varphi,B\varphi\rangle>0
\ \text{for all}\ \varphi\neq0\},
\end{equation*}
and, for $a>0$, let $\mathfrak{P}(a)$ be the set of $B$ with
$\operatorname{Re}\langle\varphi,B\varphi\rangle\ge a\|\varphi\|^2$ for all $\varphi\in\mathcal{H}_s$.
Let $\operatorname{Log}$ denote the principal branch of the logarithm on
$\mathbb{C}\setminus(-\infty,0]$. Fix $\lambda_0>0$.
\begin{enumerate}
\item[(a)] $\mathfrak{P}=\bigcup_{a>0}\mathfrak{P}(a)$; if $B\in\mathfrak{P}(a)$ and
$\|B'-B\|<a/2$, then $B'\in\mathfrak{P}(a/2)$, so that $\mathfrak{P}$ is open; and
$\mathfrak{P}$ is convex. For $B\in\mathfrak{P}$ the integral
\begin{equation}\label{4.24:eq-log-resolvent-a}
F(B):=\log\lambda_0\cdot 1+\int_0^\infty d\xi\,
\bigl[(\lambda_0+\xi)^{-1}1-(B+\xi)^{-1}\bigr]
\end{equation}
converges absolutely in $\operatorname{End}(\mathcal{H}_s)$.
\item[(b)] $F$ is holomorphic on $\mathfrak{P}$ as a function of the entries of $B$, and for every
$\dot B\in\operatorname{End}(\mathcal{H}_s)$
\begin{equation*}
DF(B)[\dot B]=\int_0^\infty(B+\xi)^{-1}\dot B\,(B+\xi)^{-1}\,d\xi,\qquad
d(\operatorname{Tr}F)(B)[\dot B]=\operatorname{Tr}(B^{-1}\dot B).
\end{equation*}
\item[(c)] $F(B)=\operatorname{Log}B$, the principal logarithm of $B$ in the holomorphic functional
calculus. Consequently $e^{F(B)}=B$ and $e^{\operatorname{Tr}F(B)}=\det B$; $F$ does not depend on
$\lambda_0$; the function $e^{-\frac12\operatorname{Tr}F(B)}$ is single-valued and holomorphic on
$\mathfrak{P}$ with square $1/\det B$; and
\begin{equation}\label{4.24:eq-log-resolvent-1}
-\tfrac12\operatorname{Tr}F(B)=-\tfrac12|s|\,n_{\mathfrak{g}}\log\lambda_0
-\tfrac12\sum_{b'\in s}\int_0^\infty d\xi\,\operatorname{tr}_{\mathfrak{g}}
\bigl[(\lambda_0+\xi)^{-1}1-(B+\xi)^{-1}\bigr](b',b').
\end{equation}
\item[(d)] Choose a basis of the real subspace $\bigoplus_{b'\in s}\mathfrak{g}\subset\mathcal{H}_s$
orthonormal for $(\cdot,\cdot)$, and identify $\mathcal{H}_s$ with $\mathbb{C}^{\mathcal{N}}$
through the resulting coordinates; in them $(\cdot,\cdot)$ is the standard bilinear form,
$\langle u,v\rangle=(\bar u,v)$ is the standard Hermitian form, and $\psi$ is called real when
$\bar\psi=\psi$. Let
$\operatorname{Sym}_{\mathcal{N}}(\mathbb{C})$ be the complex symmetric
$\mathcal{N}\times\mathcal{N}$ matrices, $\operatorname{Sym}_{\mathcal{N}}(\mathbb{R})$ the real
symmetric ones and $\operatorname{Sym}^{+}_{\mathcal{N}}(\mathbb{R})$ the real symmetric positive
definite ones. For $B\in\mathfrak{P}\cap\operatorname{Sym}_{\mathcal{N}}(\mathbb{C})$ put
\begin{equation*}
G(B):=\int_{\mathbb{R}^{\mathcal{N}}}e^{-\frac12(\psi,B\psi)}\,d\psi,\qquad
H(B):=(2\pi)^{\mathcal{N}/2}e^{-\frac12\operatorname{Tr}F(B)}.
\end{equation*}
Then for every such $B$ the integral defining $G(B)$ converges absolutely, locally uniformly in
$B$, and $G(B)=H(B)$ on $\mathfrak{P}\cap\operatorname{Sym}_{\mathcal{N}}(\mathbb{C})$.
\item[(e)] In the setting and under the hypotheses of Lemma~\ref{4.24:lem-block-coercivity}, at
every point $z$ of its domain of complex parameters, with $d\phi_s$ Lebesgue measure on the real
variables $\phi_s$ in the coordinates of \textup{(d)},
\begin{equation}\label{4.24:eq-log-resolvent-2}
\int d\phi_s\,e^{-\frac12\langle\phi_s,\mathcal{A}_{ss}(z)\phi_s\rangle}
=(2\pi)^{\mathcal{N}/2}e^{-\frac12\operatorname{Tr}F(\mathcal{A}_{ss}(z))},
\end{equation}
and both sides are holomorphic in $z$.
\end{enumerate}
\end{lemma}

\begin{proof}
\emph{Step 1: the set $\mathfrak{P}$.} If $B\in\mathfrak{P}(a)$ with $a>0$, then
$\operatorname{Re}\langle\varphi,B\varphi\rangle\ge a\|\varphi\|^2>0$ for $\varphi\neq0$, so
$\mathfrak{P}(a)\subset\mathfrak{P}$.
Conversely, let $B\in\mathfrak{P}$. The function $\varphi\mapsto\operatorname{Re}\langle\varphi,B\varphi\rangle$
is continuous and positive on the unit sphere of $\mathcal{H}_s$, which is compact since
$\mathcal{H}_s$ is finite-dimensional; so it has a minimum $a>0$ there, and by homogeneity of degree
two $\operatorname{Re}\langle\varphi,B\varphi\rangle\ge a\|\varphi\|^2$ for all $\varphi$, that is,
$B\in\mathfrak{P}(a)$. If $B\in\mathfrak{P}(a)$ and $\|B'-B\|<a/2$, then
\begin{equation*}
\operatorname{Re}\langle\varphi,B'\varphi\rangle\ge\operatorname{Re}\langle\varphi,B\varphi\rangle
-\|B'-B\|\,\|\varphi\|^2\ge\tfrac a2\|\varphi\|^2,
\end{equation*}
so $B'\in\mathfrak{P}(a/2)\subset\mathfrak{P}$; hence $\mathfrak{P}$ is open. For
$B,B'\in\mathfrak{P}$ and $t\in[0,1]$ the quantity
$\operatorname{Re}\langle\varphi,(tB+(1-t)B')\varphi\rangle$ equals
$t\operatorname{Re}\langle\varphi,B\varphi\rangle+(1-t)\operatorname{Re}\langle\varphi,B'\varphi\rangle$,
which is positive for $\varphi\ne0$; so $\mathfrak{P}$ is convex.

Two consequences of $B\in\mathfrak{P}(a)$ are used throughout. First, for $\xi\ge0$ we have
$B+\xi\in\mathfrak{P}(a+\xi)$, hence
\begin{equation*}
\|(B+\xi)\varphi\|\,\|\varphi\|\ge\operatorname{Re}\langle\varphi,(B+\xi)\varphi\rangle
\ge(a+\xi)\|\varphi\|^2;
\end{equation*}
thus $B+\xi$ is injective, hence invertible, and (Lax--Milgram)
\begin{equation}\label{4.24:eq-log-resolvent-3}
\|(B+\xi)^{-1}\|\le(a+\xi)^{-1}.
\end{equation}
Second, if $B\varphi=\mu\varphi$ with $\varphi\ne0$, then
$\operatorname{Re}\mu\,\|\varphi\|^2=\operatorname{Re}\langle\varphi,B\varphi\rangle\ge a\|\varphi\|^2$;
so the spectrum of $B$ lies in $\{\operatorname{Re}z\ge a\}$.

\emph{Step 2: convergence of \eqref{4.24:eq-log-resolvent-a}.} Let $B\in\mathfrak{P}(a)$ and $\xi\ge0$.
Then
\begin{align*}
(\lambda_0+\xi)^{-1}1-(B+\xi)^{-1}
&=(B+\xi)^{-1}\bigl[(B+\xi)-(\lambda_0+\xi)\bigr](\lambda_0+\xi)^{-1}\\
&=(B+\xi)^{-1}(B-\lambda_0)(\lambda_0+\xi)^{-1},
\end{align*}
and by \eqref{4.24:eq-log-resolvent-3} its norm is at most
$\|B-\lambda_0\|\,(a+\xi)^{-1}(\lambda_0+\xi)^{-1}$, which is $O(\xi^{-2})$ as $\xi\to\infty$. The
integrand is continuous in $\xi$ (inversion is continuous), so the integral converges absolutely.
This factorised integrand, of order $\xi^{-2}$, is the one appearing in \cite[(3.36)]{B-CMP119}, and
$\lambda_0$ plays the role of the free small positive parameter there.

\emph{Step 3: holomorphy and the derivatives.} Fix $B\in\mathfrak{P}(a)$ and let
$\|B'-B\|\le a/4$. By Step~1, $B'\in\mathfrak{P}(a/2)$, and
$\|B'-\lambda_0\|\le\|B\|+a+\lambda_0$. By Step~2 applied to $B'$ with $a/2$ in place of $a$, the
integrand at $B'$ is bounded in norm by
\begin{equation*}
\frac{\|B\|+a+\lambda_0}{(a/2+\xi)(\lambda_0+\xi)},
\end{equation*}
which is integrable on $[0,\infty)$ and independent of $B'$. For fixed $\xi$ each entry of
$(B'+\xi)^{-1}$ is, by Cramer's rule, a rational function of the entries of $B'$ whose denominator
$\det(B'+\xi)$ does not vanish on the ball; so the integrand is holomorphic in $B'$ and jointly
continuous in $(B',\xi)$. Hence, for every $R>0$, the truncated integral $\int_0^R$ is holomorphic
in the entries of $B'$: it is continuous, and in each entry separately its integrals over closed
triangles vanish by Fubini's theorem and Cauchy's theorem, so by Morera's theorem it is holomorphic
in each entry separately and, being continuous, jointly holomorphic in the entries (Osgood's
lemma). By the $B'$-free domination, the truncated integrals converge uniformly on the ball as
$R\to\infty$, and by Weierstrass' theorem the limit $F$ is holomorphic there, with derivatives the
limits of the derivatives of the truncated integrals. Since
$D[(B+\xi)^{-1}][\dot B]=-(B+\xi)^{-1}\dot B(B+\xi)^{-1}$, the derivative of the truncated
integral is $\int_0^R(B+\xi)^{-1}\dot B(B+\xi)^{-1}d\xi$, whose integrand has norm at most
$\|\dot B\|(a+\xi)^{-2}$ by \eqref{4.24:eq-log-resolvent-3}; letting $R\to\infty$ gives the formula for
$DF(B)[\dot B]$ in \textup{(b)}. Applying the linear functional $\operatorname{Tr}$, which commutes with
the convergent integral, and using cyclicity of the trace,
\begin{equation*}
d(\operatorname{Tr}F)(B)[\dot B]=\operatorname{Tr}\Bigl(\int_0^\infty(B+\xi)^{-2}d\xi\;\dot B\Bigr).
\end{equation*}
Since $\frac{d}{d\xi}(B+\xi)^{-1}=-(B+\xi)^{-2}$, we have
$\int_0^R(B+\xi)^{-2}d\xi=B^{-1}-(B+R)^{-1}$, and $\|(B+R)^{-1}\|\le(a+R)^{-1}\to0$; hence
$\int_0^\infty(B+\xi)^{-2}d\xi=B^{-1}$ and $d(\operatorname{Tr}F)(B)[\dot B]=\operatorname{Tr}(B^{-1}\dot B)$.

\emph{Step 4: the scalar identity.} Let $\operatorname{Re}\zeta>0$ and $\Lambda>0$. For
$\xi\in[0,\Lambda]$ the point $\zeta+\xi$ stays in the right half-plane, where $\operatorname{Log}$ is
holomorphic with derivative $w\mapsto1/w$. Therefore
\begin{equation*}
\int_0^\Lambda\bigl[(\lambda_0+\xi)^{-1}-(\zeta+\xi)^{-1}\bigr]d\xi
=\bigl[\log(\lambda_0+\Lambda)-\log\lambda_0\bigr]-\bigl[\operatorname{Log}(\zeta+\Lambda)
-\operatorname{Log}\zeta\bigr].
\end{equation*}
Since $\lambda_0+\Lambda>0$, division by it does not change the argument, so
$\log(\lambda_0+\Lambda)-\operatorname{Log}(\zeta+\Lambda)
=-\operatorname{Log}\bigl((\zeta+\Lambda)/(\lambda_0+\Lambda)\bigr)$. As $\Lambda\to\infty$,
$|(\zeta+\Lambda)/(\lambda_0+\Lambda)-1|=|\zeta-\lambda_0|/(\lambda_0+\Lambda)\to0$, uniformly for
$\zeta$ in a compact set, and $\operatorname{Log}$ is continuous at $1$ with $\operatorname{Log}1=0$.
Hence the integral tends to $\operatorname{Log}\zeta-\log\lambda_0$, uniformly for $\zeta$ in compact
subsets of the right half-plane.

\emph{Step 5: $F(B)=\operatorname{Log}B$.} Let $B\in\mathfrak{P}(a)$, and let $\Gamma$ be the
positively oriented boundary of the rectangle
$\{a/2\le\operatorname{Re}w\le\|B\|+1,\ |\operatorname{Im}w|\le\|B\|+1\}$. It lies, with its
interior, in $\{\operatorname{Re}w>0\}$, and it encloses the spectrum of $B$, which by Step~1 lies in
$\{\operatorname{Re}w\ge a\}\cap\{|w|\le\|B\|\}$. In the holomorphic functional calculus
\begin{equation*}
\operatorname{Log}B:=\frac1{2\pi i}\oint_\Gamma\operatorname{Log}w\,(w-B)^{-1}dw .
\end{equation*}
For $\xi\ge0$ the function $w\mapsto(w+\xi)^{-1}$ is holomorphic on $\{\operatorname{Re}w>0\}$;
since the holomorphic functional calculus is a unital algebra homomorphism sending $w\mapsto w+\xi$
to $B+\xi$, it sends $w\mapsto(w+\xi)^{-1}$ to $(B+\xi)^{-1}$ and the constant $1$ to $1$; hence
\begin{equation*}
(B+\xi)^{-1}=\frac1{2\pi i}\oint_\Gamma(w+\xi)^{-1}(w-B)^{-1}dw,\qquad
1=\frac1{2\pi i}\oint_\Gamma(w-B)^{-1}dw .
\end{equation*}
Let $R>0$. The function $(w,\xi)\mapsto[(\lambda_0+\xi)^{-1}-(w+\xi)^{-1}](w-B)^{-1}$ is continuous
on the compact set $\Gamma\times[0,R]$, so Fubini's theorem and Step~4 give
\begin{align*}
&\int_0^R\bigl[(\lambda_0+\xi)^{-1}1-(B+\xi)^{-1}\bigr]d\xi\\
&\qquad=\frac1{2\pi i}\oint_\Gamma\bigl[\log(\lambda_0+R)-\log\lambda_0-\operatorname{Log}(w+R)
+\operatorname{Log}w\bigr](w-B)^{-1}dw .
\end{align*}
As $R\to\infty$ the bracket converges to $\operatorname{Log}w-\log\lambda_0$ uniformly on $\Gamma$
(Step~4), while $\Gamma$ has finite length and $(w-B)^{-1}$ is bounded on $\Gamma$. Therefore
\begin{equation*}
\int_0^\infty\bigl[(\lambda_0+\xi)^{-1}1-(B+\xi)^{-1}\bigr]d\xi
=\frac1{2\pi i}\oint_\Gamma\bigl[\operatorname{Log}w-\log\lambda_0\bigr](w-B)^{-1}dw
=\operatorname{Log}B-\log\lambda_0\cdot1,
\end{equation*}
and adding $\log\lambda_0\cdot1$ gives $F(B)=\operatorname{Log}B$. By the composition rule of the
functional calculus, $e^{F(B)}=(\exp\circ\operatorname{Log})(B)=B$.

\emph{Step 6: the trace, the determinant, \eqref{4.24:eq-log-resolvent-1} and $\lambda_0$.} By Schur
triangularisation choose a unitary $U$ with $T'=U^*BU$ upper triangular, with diagonal entries
$t_{ii}$. Since $\langle\varphi,T'\varphi\rangle=\langle U\varphi,BU\varphi\rangle$ and
$\|U\varphi\|=\|\varphi\|$, we have $T'\in\mathfrak{P}(a)$; in particular
$\operatorname{Re}t_{ii}=\operatorname{Re}\langle e_i,T'e_i\rangle\ge a>0$ for the basis vectors
$e_i$ in which $T'$ is triangular. As $(T'+\xi)^{-1}=U^*(B+\xi)^{-1}U$ and the integral is linear,
$F(T')=U^*F(B)U$. The matrix $(T'+\xi)^{-1}$ is upper triangular with diagonal entries
$(t_{ii}+\xi)^{-1}$, so the integrand of \eqref{4.24:eq-log-resolvent-a} at $T'$ is upper triangular
with diagonal entries $(\lambda_0+\xi)^{-1}-(t_{ii}+\xi)^{-1}$, and by Step~4 $F(T')$ is upper
triangular with diagonal entries $\log\lambda_0+\operatorname{Log}t_{ii}-\log\lambda_0
=\operatorname{Log}t_{ii}$. Hence
\begin{equation*}
\operatorname{Tr}F(B)=\operatorname{Tr}F(T')=\sum_i\operatorname{Log}t_{ii},\qquad
e^{\operatorname{Tr}F(B)}=\prod_it_{ii}=\det T'=\det B .
\end{equation*}
For \eqref{4.24:eq-log-resolvent-1}: for an operator $K$ on $\mathcal{H}_s$,
$\operatorname{Tr}K=\sum_{b'\in s}\operatorname{tr}_{\mathfrak{g}}K(b',b')$, and
$\operatorname{Tr}(\log\lambda_0\cdot1)=\mathcal{N}\log\lambda_0=|s|\,n_{\mathfrak{g}}\log\lambda_0$.
Each map $K\mapsto\operatorname{tr}_{\mathfrak{g}}K(b',b')$ is linear and satisfies
$|\operatorname{tr}_{\mathfrak{g}}K(b',b')|\le n_{\mathfrak{g}}\|K(b',b')\|\le n_{\mathfrak{g}}\|K\|$,
the last inequality because a diagonal block is a compression of $K$ and has norm at most $\|K\|$
(Lemma~\ref{4.24:lem-principal-blocks}); so it commutes with the absolutely convergent integral of
Step~2, and the finite sum over $b'\in s$ commutes with it as well. Multiplying by $-\frac12$ gives
\eqref{4.24:eq-log-resolvent-1}.

$F$ does not depend on $\lambda_0$: the derivative of the integrand in $\lambda_0$ is
$-(\lambda_0+\xi)^{-2}1$, dominated by $(\lambda_1+\xi)^{-2}$ for $\lambda_0\ge\lambda_1>0$, so
differentiation under the integral is allowed and
\begin{equation*}
\frac{d}{d\lambda_0}F(B)=\lambda_0^{-1}1-\int_0^\infty(\lambda_0+\xi)^{-2}d\xi\cdot1
=\lambda_0^{-1}1-\lambda_0^{-1}1=0 .
\end{equation*}
Finally, $\operatorname{Tr}F$ is holomorphic on $\mathfrak{P}$ by Step~3, so
$e^{-\frac12\operatorname{Tr}F(B)}$ is a single-valued holomorphic function on $\mathfrak{P}$, and its
square is $e^{-\operatorname{Tr}F(B)}=1/\det B$. No branch of a square root of $\det B$ is chosen
anywhere. This proves \textup{(a)}, \textup{(b)} and \textup{(c)}.

\emph{Step 7: holomorphy of $G$ and $H$.} In the coordinates of \textup{(d)}, a real vector $\psi$
satisfies $\bar\psi=\psi$, so $(\psi,B\psi)=\langle\psi,B\psi\rangle$. Hence for
$B\in\mathfrak{P}(a)$ the integrand of $G$ has modulus
$e^{-\frac12\operatorname{Re}\langle\psi,B\psi\rangle}\le e^{-\frac a2|\psi|^2}$, and $G(B)$ converges
absolutely. The set $\mathfrak{P}\cap\operatorname{Sym}_{\mathcal{N}}(\mathbb{C})$ is open in
$\operatorname{Sym}_{\mathcal{N}}(\mathbb{C})\cong\mathbb{C}^{\mathcal{N}(\mathcal{N}+1)/2}$ (coordinates
the entries $B_{ij}$, $i\le j$), because $\mathfrak{P}$ is open (Step~1). If
$B\in\mathfrak{P}(a)\cap\operatorname{Sym}_{\mathcal{N}}(\mathbb{C})$ and $\|B'-B\|<a/2$, then
$B'\in\mathfrak{P}(a/2)$ and the integrand at $B'$ is dominated by $e^{-\frac a4|\psi|^2}$, which is
integrable and $B'$-free. For each $\psi$ the integrand $e^{-\frac12\sum_{i,j}B_{ij}\psi_i\psi_j}$ is
entire in the entries; as in Step~3, the integrals over balls $|\psi|\le R$ are holomorphic, they
converge locally uniformly as $R\to\infty$ by the domination, and $G$ is holomorphic on
$\mathfrak{P}\cap\operatorname{Sym}_{\mathcal{N}}(\mathbb{C})$. The function $H$ is holomorphic on
$\mathfrak{P}$ by Step~3, hence on $\mathfrak{P}\cap\operatorname{Sym}_{\mathcal{N}}(\mathbb{C})$, the
inclusion being linear.

\emph{Step 8: $G=H$ on real positive definite matrices.} Let
$B\in\operatorname{Sym}^{+}_{\mathcal{N}}(\mathbb{R})$. Being real symmetric, $B$ is Hermitian and
$\operatorname{Re}\langle\varphi,B\varphi\rangle=\langle\varphi,B\varphi\rangle>0$ for
$\varphi\ne0$, so $B\in\mathfrak{P}$. Choose an orthogonal $O$ with
$O^TBO=\operatorname{diag}(\beta_1,\dots,\beta_{\mathcal{N}})$, $\beta_i>0$. The substitution
$\psi=O\chi$ has Jacobian one, so by the one-dimensional Gaussian integral
\begin{equation*}
G(B)=\prod_i\int_{\mathbb{R}}e^{-\frac12\beta_i\chi_i^2}d\chi_i
=(2\pi)^{\mathcal{N}/2}\prod_i\beta_i^{-1/2}.
\end{equation*}
Since $O$ is unitary and $O^TBO$ is diagonal, hence upper triangular, Step~6 gives
$\operatorname{Tr}F(B)=\sum_i\operatorname{Log}\beta_i=\sum_i\log\beta_i$, so
\begin{equation*}
H(B)=(2\pi)^{\mathcal{N}/2}e^{-\frac12\sum_i\log\beta_i}=(2\pi)^{\mathcal{N}/2}\prod_i\beta_i^{-1/2}.
\end{equation*}

\emph{Step 9: the identity theorem.} Write
$B\in\operatorname{Sym}_{\mathcal{N}}(\mathbb{C})$ as $B=X+iY$ with $X,Y\in\operatorname{Sym}_{\mathcal{N}}(\mathbb{R})$
the real and imaginary parts of the entries. Then $B^*=X-iY$ and
$\operatorname{Re}\langle\varphi,B\varphi\rangle=\langle\varphi,\tfrac12(B+B^*)\varphi\rangle
=\langle\varphi,X\varphi\rangle$; a real symmetric matrix is positive definite on $\mathbb{C}^{\mathcal{N}}$
exactly when it is so on $\mathbb{R}^{\mathcal{N}}$. Hence
\begin{equation*}
\mathfrak{P}\cap\operatorname{Sym}_{\mathcal{N}}(\mathbb{C})
=\{X+iY:\ X\in\operatorname{Sym}^{+}_{\mathcal{N}}(\mathbb{R}),\ Y\in\operatorname{Sym}_{\mathcal{N}}(\mathbb{R})\},
\end{equation*}
which is convex, hence connected. In the coordinates $B_{ij}$, $i\le j$, the real points
$\operatorname{Sym}_{\mathcal{N}}(\mathbb{R})$ form a maximal totally real subspace, and
$\operatorname{Sym}^{+}_{\mathcal{N}}(\mathbb{R})$ is a nonempty (it contains $1$) open subset of it, on
which the holomorphic function $G-H$ vanishes by Step~8. Let
$X_0\in\operatorname{Sym}^{+}_{\mathcal{N}}(\mathbb{R})$. Since $G-H$ vanishes on a real neighbourhood of
$X_0$, all iterated derivatives of $G-H$ along real directions vanish at $X_0$; for a holomorphic
function the complex partial derivative $\partial/\partial w_j$ equals the real partial derivative
$\partial/\partial x_j$ along the real part of the coordinate $w_j=x_j+iy_j$, so all complex partial
derivatives of $G-H$ vanish at $X_0$. The Taylor series of $G-H$ at $X_0$ is therefore zero, and
$G-H\equiv0$ on a polydisc around $X_0$. By the identity theorem on the connected open set
$\mathfrak{P}\cap\operatorname{Sym}_{\mathcal{N}}(\mathbb{C})$, $G-H\equiv0$ there. This proves
\textup{(d)}.

Symmetry of $B$ is needed in \textup{(d)}: for a general $B\in\mathfrak{P}$ the Gaussian integral
depends on $B$ only through $\mathcal{S}B$ (Lemma~\ref{4.24:lem-kernel-polarisation}), while $\det B$
does not. This is why the precision $\mathcal{A}$ is taken to be the symmetrisation
$\mathcal{S}K_{\sharp}$ of the stage form kernel $K_{\sharp}$.

\emph{Step 10: application to $\mathcal{A}_{ss}$.} Assume the hypotheses of
Lemma~\ref{4.24:lem-block-coercivity}, let $z$ be a point of its domain of complex parameters, and
put $B:=\mathcal{A}_{ss}(z)$. By
Lemma~\ref{4.24:lem-principal-blocks}, $B^T=B$, that is,
$B\in\operatorname{Sym}_{\mathcal{N}}(\mathbb{C})$ in the coordinates of \textup{(d)}. By
Lemma~\ref{4.24:lem-block-coercivity} (coercivity of the survivors' block), whose hypotheses are in
force, $B\in\mathfrak{P}(a_{\mathbb{C}})$ with $a_{\mathbb{C}}$
as in \eqref{4.24:eq-block-coercivity-a}. For real $\phi_s$,
$\langle\phi_s,B\phi_s\rangle=(\phi_s,B\phi_s)$, so the left side of \eqref{4.24:eq-log-resolvent-2}
is $G(B)$, and \textup{(d)} gives \eqref{4.24:eq-log-resolvent-2}. The map
$z\mapsto\mathcal{A}_{ss}(z)$ is holomorphic (Lemma~\ref{4.24:lem-principal-blocks}) with values in
$\mathfrak{P}$, and $B\mapsto e^{-\frac12\operatorname{Tr}F(B)}$ is holomorphic on $\mathfrak{P}$ by
\textup{(c)}; their composition, the right side of \eqref{4.24:eq-log-resolvent-2}, is therefore
holomorphic in $z$, and so is the left side, which equals it. This proves \textup{(e)}.
\end{proof}

\begin{lemma}[Conditional and marginal precision of a Gaussian]\label{4.24:lem-conditioning-marginal}
Let $\mathfrak{Bo}=s\sqcup\mathcal{W}$, so that
$\mathcal{H}_{\mathfrak{Bo}}=\mathcal{H}_s\oplus\mathcal{H}_{\mathcal{W}}$ with coordinate projections
$P_s$, $P_{\mathcal{W}}$. Write $(\cdot,\cdot)$ for the bilinear pairing and
$\langle\cdot,\cdot\rangle$ for the Hermitian inner product. For $E,E'\in\{s,\mathcal{W}\}$ let
$\mathcal{A}_{EE'}=P_E\mathcal{A}P_{E'}$, regarded as a map $\mathcal{H}_{E'}\to\mathcal{H}_E$. Let
$\mathcal{A}$ be symmetric, $\mathcal{A}^T=\mathcal{A}$, and coercive in the sense that
$\operatorname{Re}\langle\phi,\mathcal{A}\phi\rangle\ge a\|\phi\|^2$ for all
$\phi\in\mathcal{H}_{\mathfrak{Bo}}$ and some $a>0$; that is, $\mathcal{A}$ lies in the class
$\mathfrak{P}(a)$ of Lemma~\ref{4.24:lem-log-resolvent}, taken on $\mathcal{H}_{\mathfrak{Bo}}$. For
$\phi=\phi_s\oplus\phi_{\mathcal{W}}$,
\begin{equation}\label{4.24:eq-conditioning-marginal-1}
(\phi,\mathcal{A}\phi)=(\phi_s,\mathcal{A}_{ss}\phi_s)+2(\phi_s,\mathcal{A}_{s\mathcal{W}}\phi_{\mathcal{W}})
+(\phi_{\mathcal{W}},\mathcal{A}_{\mathcal{W}\mathcal{W}}\phi_{\mathcal{W}}).
\end{equation}
\begin{itemize}
\item[(a)] Fix $\phi_{\mathcal{W}}$ (in particular, $\phi_{\mathcal{W}}=0$ is allowed). Keep the cross form
$\Psi(\phi_s):=(\phi_s,\mathcal{A}_{s\mathcal{W}}\phi_{\mathcal{W}})$ as a separate part of the integrand.
Then the $d\phi_s$-weight is $e^{-\frac12(\phi_s,\mathcal{A}_{ss}\phi_s)}$. Its precision is
$\mathcal{A}_{ss}$ and its covariance is $(\mathcal{A}_{ss})^{-1}$. It has one normalisation factor,
the right-hand side of the first identity in \eqref{4.24:eq-conditioning-marginal-a}, where
the integral is over real $\phi_s\in\mathcal{H}_s$ with Lebesgue measure and
$\mathcal{N}=n_{\mathfrak{g}}|s|$ is the real dimension, as in the Gaussian identity of
Lemma~\ref{4.24:lem-log-resolvent}.
This factor is a holomorphic functional of $\mathcal{A}_{ss}$ alone. It is the same for every
$\phi_{\mathcal{W}}$ and every cross kernel $\mathcal{A}_{s\mathcal{W}}$.
\item[(b)] Integrating $\phi_{\mathcal{W}}$ Gaussianly (the marginal) gives $\phi_s$ the precision
\begin{equation*}
S:=\mathcal{A}_{ss}-\mathcal{A}_{s\mathcal{W}}\mathcal{A}_{\mathcal{W}\mathcal{W}}^{-1}
\mathcal{A}_{\mathcal{W}s},
\end{equation*}
and the second and third identities in \eqref{4.24:eq-conditioning-marginal-a} hold. In particular
$(\mathcal{A}^{-1})_{ss}\ne(\mathcal{A}_{ss})^{-1}$ whenever
$\mathcal{A}_{s\mathcal{W}}\mathcal{A}_{\mathcal{W}\mathcal{W}}^{-1}\mathcal{A}_{\mathcal{W}s}\ne0$.
\item[(c)] Completing the square in (a), as in the last line of
\eqref{4.24:eq-conditioning-marginal-a}, would produce the quadratic form
$-(\phi_{\mathcal{W}},\mathcal{A}_{\mathcal{W}s}\mathcal{A}_{ss}^{-1}\mathcal{A}_{s\mathcal{W}}
\phi_{\mathcal{W}})$ in the variable $\phi_{\mathcal{W}}$, which is not integrated. Its kernel
$\mathcal{A}_{\mathcal{W}s}\mathcal{A}_{ss}^{-1}\mathcal{A}_{s\mathcal{W}}$ is non-local and anchored on
$\mathcal{W}$; it is the remnant of an integration over $\mathcal{W}$ that is not carried out.
The rearrangement is incompatible with the factors of the integrand in $\phi_s$ that are
not exponentials of quadratic forms, such as the characteristic functions $\chi^{s}$ and the
group-valued factor, and it is not used.
\end{itemize}
The identities referred to in (a), (b) and (c), with $S$ as defined in (b) and
$\phi_s'=\phi_s+\mathcal{A}_{ss}^{-1}\mathcal{A}_{s\mathcal{W}}\phi_{\mathcal{W}}$, are
\begin{equation}\label{4.24:eq-conditioning-marginal-a}
\begin{gathered}
\int e^{-\frac12(\phi_s,\mathcal{A}_{ss}\phi_s)}\,d\phi_s
=(2\pi)^{\mathcal{N}/2}e^{-\frac12\operatorname{Tr}F(\mathcal{A}_{ss})},\\
(\mathcal{A}^{-1})_{ss}=S^{-1},\qquad
\det\mathcal{A}=\det S\cdot\det\mathcal{A}_{\mathcal{W}\mathcal{W}},\\
(\phi_s,\mathcal{A}_{ss}\phi_s)+2\Psi(\phi_s)
=(\phi_s',\mathcal{A}_{ss}\phi_s')
-(\phi_{\mathcal{W}},\mathcal{A}_{\mathcal{W}s}\mathcal{A}_{ss}^{-1}\mathcal{A}_{s\mathcal{W}}\phi_{\mathcal{W}}).
\end{gathered}
\end{equation}
\end{lemma}

\begin{proof}
\emph{Symmetry of the blocks.} Coordinate projections satisfy $P_E^T=P_E$. Since
$\mathcal{A}^T=\mathcal{A}$, it follows, as in Lemma~\ref{4.24:lem-principal-blocks} for
$\mathcal{A}_{ss}$, that
\begin{equation*}
\mathcal{A}_{s\mathcal{W}}^T=P_{\mathcal{W}}\mathcal{A}^TP_s=\mathcal{A}_{\mathcal{W}s},\qquad
\mathcal{A}_{EE}^T=\mathcal{A}_{EE}\quad(E\in\{s,\mathcal{W}\}).
\end{equation*}
Expanding $(\phi,\mathcal{A}\phi)$ by bilinearity gives four terms. The two cross terms are equal,
because
$(\phi_{\mathcal{W}},\mathcal{A}_{\mathcal{W}s}\phi_s)=(\phi_s,\mathcal{A}_{s\mathcal{W}}\phi_{\mathcal{W}})$.
This proves \eqref{4.24:eq-conditioning-marginal-1}.

\emph{Coercivity of the diagonal blocks.} Since $P_E\varphi=\varphi$ for
$\varphi\in\mathcal{H}_E$, the Rayleigh identity
\eqref{4.24:eq-principal-blocks-a} of Lemma~\ref{4.24:lem-principal-blocks} holds for either
coordinate subspace: every $\varphi\in\mathcal{H}_E$ satisfies
$\langle\varphi,\mathcal{A}_{EE}\varphi\rangle=\langle\varphi,\mathcal{A}\varphi\rangle$. Hence
$\mathcal{A}_{ss}\in\mathfrak{P}(a)$ on $\mathcal{H}_s$ and
$\mathcal{A}_{\mathcal{W}\mathcal{W}}\in\mathfrak{P}(a)$ on $\mathcal{H}_{\mathcal{W}}$. An operator
$B$ in $\mathfrak{P}(a)$ on any of these spaces satisfies
\begin{equation*}
a\|\varphi\|^2\le\operatorname{Re}\langle\varphi,B\varphi\rangle\le\|\varphi\|\,\|B\varphi\|,
\end{equation*}
so $\|B\varphi\|\ge a\|\varphi\|$ and $B$ is
injective, and therefore invertible on the finite-dimensional space. In particular
$\mathcal{A}_{ss}^{-1}$ and $\mathcal{A}_{\mathcal{W}\mathcal{W}}^{-1}$ exist, and so does $\mathcal{A}^{-1}$.

\emph{(a).} By \eqref{4.24:eq-conditioning-marginal-1},
\begin{equation*}
e^{-\frac12(\phi,\mathcal{A}\phi)}=e^{-\frac12(\phi_s,\mathcal{A}_{ss}\phi_s)}\,e^{-\Psi(\phi_s)}\,
e^{-\frac12(\phi_{\mathcal{W}},\mathcal{A}_{\mathcal{W}\mathcal{W}}\phi_{\mathcal{W}})}.
\end{equation*}
At fixed $\phi_{\mathcal{W}}$ the last factor is a constant. The factor $e^{-\Psi(\phi_s)}$ is kept in
the integrand. What remains as the weight of $d\phi_s$ is
$e^{-\frac12(\phi_s,\mathcal{A}_{ss}\phi_s)}$, whose precision is $\mathcal{A}_{ss}$ and whose
covariance is $(\mathcal{A}_{ss})^{-1}$.

Since $\mathcal{A}_{ss}\in\mathfrak{P}$, the Gaussian identity of Lemma~\ref{4.24:lem-log-resolvent}
applies with $B=\mathcal{A}_{ss}$. It gives the first line of
\eqref{4.24:eq-conditioning-marginal-a}. The right-hand side is a holomorphic functional of
$\mathcal{A}_{ss}$, because $F$ is holomorphic on $\mathfrak{P}$ by the same lemma. It involves neither
$\phi_{\mathcal{W}}$ nor $\mathcal{A}_{s\mathcal{W}}$, so it is one factor, the same for every
$\phi_{\mathcal{W}}$ and every cross kernel.

\emph{(b).} Put $\eta=\phi_{\mathcal{W}}+\mathcal{A}_{\mathcal{W}\mathcal{W}}^{-1}\mathcal{A}_{\mathcal{W}s}\phi_s$.
Expanding $(\eta,\mathcal{A}_{\mathcal{W}\mathcal{W}}\eta)$ uses three facts: the symmetry of
$\mathcal{A}_{\mathcal{W}\mathcal{W}}$ and of its inverse, and
$\mathcal{A}_{\mathcal{W}s}^T=\mathcal{A}_{s\mathcal{W}}$. The expansion gives
\begin{equation*}
(\eta,\mathcal{A}_{\mathcal{W}\mathcal{W}}\eta)
=(\phi_{\mathcal{W}},\mathcal{A}_{\mathcal{W}\mathcal{W}}\phi_{\mathcal{W}})
+2(\phi_s,\mathcal{A}_{s\mathcal{W}}\phi_{\mathcal{W}})
+(\phi_s,\mathcal{A}_{s\mathcal{W}}\mathcal{A}_{\mathcal{W}\mathcal{W}}^{-1}\mathcal{A}_{\mathcal{W}s}\phi_s).
\end{equation*}
Combined with \eqref{4.24:eq-conditioning-marginal-1}, this yields
\begin{equation*}
(\phi,\mathcal{A}\phi)=(\eta,\mathcal{A}_{\mathcal{W}\mathcal{W}}\eta)+(\phi_s,S\phi_s).
\end{equation*}
The Gaussian integration over $\phi_{\mathcal{W}}$ absorbs the first summand. Indeed, write
$c=\mathcal{A}_{\mathcal{W}\mathcal{W}}^{-1}\mathcal{A}_{\mathcal{W}s}\phi_s\in\mathcal{H}_{\mathcal{W}}$
and let $I(c)$ be the integral of
$e^{-\frac12(\phi_{\mathcal{W}}+c,\,\mathcal{A}_{\mathcal{W}\mathcal{W}}(\phi_{\mathcal{W}}+c))}$
over real $\phi_{\mathcal{W}}$. For real $\phi_{\mathcal{W}}$ the bilinear and the Hermitian pairings
agree, so $\operatorname{Re}(\phi_{\mathcal{W}},\mathcal{A}_{\mathcal{W}\mathcal{W}}\phi_{\mathcal{W}})
\ge a\|\phi_{\mathcal{W}}\|^2$, and the real part of the exponent is bounded above by
\begin{equation*}
-\tfrac12a\|\phi_{\mathcal{W}}\|^2+\|\mathcal{A}_{\mathcal{W}\mathcal{W}}\|\,\|c\|\,
\|\phi_{\mathcal{W}}\|+\tfrac12\|\mathcal{A}_{\mathcal{W}\mathcal{W}}\|\,\|c\|^2 .
\end{equation*}
Hence the integral converges locally uniformly in $c\in\mathcal{H}_{\mathcal{W}}$, and $I$ is an entire
function of $c$. For real $c$, $I(c)=I(0)$ by translation invariance of $d\phi_{\mathcal{W}}$. An entire
function on $\mathcal{H}_{\mathcal{W}}$ that is constant on the real subspace is constant, so
$I(c)=I(0)$ for every $c$. The value $I(0)$ is given by the Gaussian identity of
Lemma~\ref{4.24:lem-log-resolvent}, taken on $\mathcal{H}_{\mathcal{W}}$ with
$B=\mathcal{A}_{\mathcal{W}\mathcal{W}}$, and does not depend on $\phi_s$. The integration thus leaves
the weight
$e^{-\frac12(\phi_s,S\phi_s)}$, so the precision of $\phi_s$ is $S$.

For the remaining identities, multiplying out the block
matrices gives the block LU factorisation
\begin{equation*}
\mathcal{A}=\begin{pmatrix}1&\mathcal{A}_{s\mathcal{W}}\mathcal{A}_{\mathcal{W}\mathcal{W}}^{-1}\\0&1
\end{pmatrix}
\begin{pmatrix}S&0\\0&\mathcal{A}_{\mathcal{W}\mathcal{W}}\end{pmatrix}
\begin{pmatrix}1&0\\\mathcal{A}_{\mathcal{W}\mathcal{W}}^{-1}\mathcal{A}_{\mathcal{W}s}&1\end{pmatrix}.
\end{equation*}
Indeed, the $ss$-block of the product is
$S+\mathcal{A}_{s\mathcal{W}}\mathcal{A}_{\mathcal{W}\mathcal{W}}^{-1}\mathcal{A}_{\mathcal{W}s}
=\mathcal{A}_{ss}$, and its other blocks are $\mathcal{A}_{s\mathcal{W}}$, $\mathcal{A}_{\mathcal{W}s}$
and $\mathcal{A}_{\mathcal{W}\mathcal{W}}$.

The triangular factors have determinant $1$, which gives
$\det\mathcal{A}=\det S\cdot\det\mathcal{A}_{\mathcal{W}\mathcal{W}}$. Because $\mathcal{A}$ and
$\mathcal{A}_{\mathcal{W}\mathcal{W}}$ are invertible, so is $S$. Inverting the factorisation,
\begin{equation*}
\mathcal{A}^{-1}=\begin{pmatrix}1&0\\-\mathcal{A}_{\mathcal{W}\mathcal{W}}^{-1}\mathcal{A}_{\mathcal{W}s}&1
\end{pmatrix}
\begin{pmatrix}S^{-1}&0\\0&\mathcal{A}_{\mathcal{W}\mathcal{W}}^{-1}\end{pmatrix}
\begin{pmatrix}1&-\mathcal{A}_{s\mathcal{W}}\mathcal{A}_{\mathcal{W}\mathcal{W}}^{-1}\\0&1\end{pmatrix}.
\end{equation*}
The $ss$-block of this product is $S^{-1}$. Moreover $S^{-1}=(\mathcal{A}_{ss})^{-1}$ holds if and only
if $S=\mathcal{A}_{ss}$, that is, if and only if
$\mathcal{A}_{s\mathcal{W}}\mathcal{A}_{\mathcal{W}\mathcal{W}}^{-1}\mathcal{A}_{\mathcal{W}s}=0$. This
is the Schur-complement relation of \eqref{4.24:eq-principal-blocks-a}.

\emph{(c).} The same expansion as in (b), with the roles of $s$ and $\mathcal{W}$ exchanged, gives
the last line of \eqref{4.24:eq-conditioning-marginal-a}. The second term on its right-hand side is
a quadratic form in the variable $\phi_{\mathcal{W}}$, which is not integrated. Its kernel is
$\mathcal{A}_{\mathcal{W}s}\mathcal{A}_{ss}^{-1}\mathcal{A}_{s\mathcal{W}}$, which is non-local and
acts on $\mathcal{H}_{\mathcal{W}}$. This term is what an integration over $\mathcal{W}$ would leave
behind, and no such integration is performed here.

Using $\phi_s'$ as the integration variable turns every factor of the integrand that depends on
$\phi_s$ and is not the exponential of a quadratic form into a function of $\phi_s'$ and
$\phi_{\mathcal{W}}$ jointly, through the
non-local operator $\mathcal{A}_{ss}^{-1}\mathcal{A}_{s\mathcal{W}}$. For this reason the rearrangement
is not used, and the cross form is kept as the integrand $\Psi$, as in (a).
\end{proof}

\begin{lemma}[The walk expansion on unions of whole components]\label{4.24:prf-region-walk-expansion}
Fix the step $k$ and a stage $n+1$ of the large-field operation at that step, at which the
level-$j$ fluctuation fields $A_j,B_j$ are integrated. Let $Z^{s}$ be the survivors' region of
the stage (Definition~\ref{4.24:def-survivors-bonds}).
\begin{enumerate}
\item[(a)] The region over which the integration $dA_j\,dB_j$ of \cite[(1.60)]{B-CMP122a}
runs at this stage is $(Z_{j+1}\setminus Z''_{j+1})\cap Z^{s}$, with $Z_{j+1}$ and $Z''_{j+1}$
the large-field regions of \cite[p.~188]{B-CMP122a}. We write $\Lambda'_s$ for the
level-$j$ integration region of $Z^{s}$, that is, its rim together with its collar.
\item[(b)] Let $\Lambda'$ be any union of whole components' stage-$(n+1)$ integration regions
inside the structure of step $k$ common to the whole class of components of
Definition~\ref{4.24:def-survivors-bonds}. Let $\mathcal{A}(\Lambda')$
be the compression to the level-$j$ integration bonds of $\Lambda'$ of the global level-$j$
fluctuation precision of the $(j+1)$-st application of the operation $\mathbf{T}$. Here the fields
off $\Lambda'$ are held at $0$, and the $\delta$-function constraint is eliminated by the operator $C$
of \cite[(3.156)--(3.158)]{B-CMP99b}. Assume that $M$ is sufficiently large and $M\alpha_0$
sufficiently small,
and that the background $\mathbb{U}$ satisfies \cite[(3.95)]{B-CMP99b} on the cubes visited by the
walks below. Then for every $\xi\ge0$
\begin{equation}\label{4.24:eq-region-walk-expansion-a}
\begin{gathered}
(\mathcal{A}(\Lambda')+\xi)^{-1}=\sum_{\omega}t^{\xi}_{\omega},\qquad
D(\omega):=X_0\cup\bigcup_{a\ge1}X_a^{\sim},\\
\|t^{\xi}_{\omega}\|\le m(\omega):=O(1)\bigl(O(1)M^{-1/2}\bigr)^{p}e^{-\delta_0 d(\omega)},
\end{gathered}
\end{equation}
with the following properties.
\begin{itemize}
\item The sum converges. It is the generalised random walk expansion of
\cite[(3.107)]{B-CMP99b}.
\item It runs over walks $\omega=(X_0,\dots,X_p)$ of localisation domains in proper scales,
consecutive domains intersecting.
\item Each $t^{\xi}_{\omega}$ depends on $\mathbb{U}$ only through its restriction to $D(\omega)$,
where $X_a^{\sim}$ and $d(\omega)$ are as in \cite[Theorem~3.10]{B-CMP99b}.
\item Each $t^{\xi}_{\omega}$ is holomorphic on the complex extension constructed in
\cite{B-CMP99b}.
\item The bound in \eqref{4.24:eq-region-walk-expansion-a} holds in each of the norms of
\cite[(3.42)--(3.47)]{B-CMP99b}.
\item The constants $O(1)$ and $\delta_0$ of $m(\omega)$, and the covariance constant $B_0$ of
\cite[Theorems~3.10, 3.15]{B-CMP99b}, depend on $(d,L)$ only. They do not depend on
the sequence $\{\Omega_j\}$ of small-field regions, and they do not depend on $\Lambda'$.
\end{itemize}
\item[(c)] Under the same hypotheses, consider the bond density of \cite[(3.36)--(3.37)]{B-CMP119}.
In it the class of walks with $|\omega|=0$, that is, of length $p=0$, is paired with
$(\lambda_0+\xi)^{-1}\cdot 1$, $\lambda_0$ being the number
that appears in that density. For a walk $\omega$, let $\hat t^{\xi}_{\omega}$ be the term of that
density indexed by $\omega$ before the $\xi$-integration, built from $t^{\xi}_{\omega}$; for the class
$|\omega|=0$ it is $t^{\xi}_{\omega}$ paired with $(\lambda_0+\xi)^{-1}\cdot1$. Then there is a majorant
$\bar m(\omega)$, independent of $\xi$ and having the same form as $m(\omega)$ in
\eqref{4.24:eq-region-walk-expansion-a}, with
\begin{equation}\label{4.24:eq-region-walk-expansion-b}
\Bigl\|\int \hat t^{\xi}_{\omega}\,d\xi\Bigr\|\le \bar m(\omega),
\end{equation}
the integral being the $\xi$-integration of \cite[(3.36)--(3.37)]{B-CMP119}.
\item[(d)] In particular, (b) and (c) hold at $\Lambda'=\Lambda'_s$, with the same constants as at
the integration region of the whole stage.
\end{enumerate}
\end{lemma}

\begin{proof}
\emph{Part (a): the region.} We first recall how \cite[p.~188]{B-CMP122a} proceeds at every stage.
\begin{itemize}
\item In the integral obtained so far, the decomposition of unity
$1=\chi'^{(j)}+(1-\chi'^{(j)})$ is introduced in each component of $Z$.
\item At \cite[(1.55)]{B-CMP122a}, the components carrying the large-field function
$1-\chi'^{(j)}$ are excluded from $Z$.
\item The excluded components are treated in \cite[(1.56)--(1.59)]{B-CMP122a}.
\item Only then is \cite[(1.60)]{B-CMP122a} written. It carries the characteristic function
$\chi'^{(j)}$ alone.
\end{itemize}
At the first stage \cite{B-CMP122a} does the same. It introduces the decomposition of unity
$1=\chi'_j+(1-\chi'_j)$ for each component of $Z$, excludes from $Z$ the components with the
large-field functions $1-\chi'_j$, and denotes the remaining region anew. Thus the letter $Z$ is
re-denoted after each exclusion. This is the convention by which
Definition~\ref{4.24:def-survivors-bonds} writes $Z^{[l]}$ for the components still in the run at
stage $l$.

Consequently, the integral \cite[(1.60)]{B-CMP122a} at the present stage is written after the
per-component split and after the exclusion. Its integration $dA_j\,dB_j$ therefore runs over
$(Z_{j+1}\setminus Z''_{j+1})\cap Z^{s}$. This proves (a). Three consequences follow.
\begin{itemize}
\item The band $R_n$ of the stage, read at this stage, is this region.
\item The condition of the normalisation lemma of the stage that every domain $X$ meet $R_n$
becomes, at this stage, the condition of Theorem~\ref{4.24:thm-survivor-block-expansion} that
$X$ meet $R_n\cap Z^{s}$.
\item The integration region at which the proof of the normalisation lemma of the stage uses the
walk expansions of \cite[Theorems~3.10, 3.15]{B-CMP99b} is, at this stage, $\Lambda'_s$.
\end{itemize}

Independently of the letter used for the region, in the fixed term of \cite[(1.55)]{B-CMP122a}
every factor other than the interaction splits by components, up to the cross pieces of the
collar form. We check the factors one at a time.
\begin{itemize}
\item The characteristic functions are local to components. By \cite{B-CMP122a}, all the
characteristic functions introduced depend only on the field variables localised in the
corresponding components of the large-field region $Z_k$.
\item The normalisation factor $z^{(j)}$ is a product over cells \cite{B-CMP119}.
\item The term $E_0^{(j)}$ is additive in the region on which it is evaluated.
\item The rim form is written component by component in \cite{B-CMP119}.
\item The entries of the collar form between different components enter only through elements
$e_a$ of that form (Definition~\ref{4.24:def-stage-form}) whose domain $Y'$ satisfies
$|Y'|\ge R_k$.
\item The pieces of the interaction that read both $\phi_s$ and $\phi_W$ are its mixed elements,
that is, the elements of the interaction reading variables on both bond sets $s$ and
$\mathcal{W}$; in them $\phi_W$ ranges over the set $\lambda$ admitted for the arrested
variables in these elements.
\end{itemize}
Finally, \cite{B-CMP122a} applies the analysis of \cite[\S3]{B-CMP119} to the integral
\cite[(1.60)]{B-CMP122a} over the region of the stage, without any condition on that region. In
particular it applies that analysis at the present region.

\emph{Part (b): the walk expansion on $\Lambda'$.} We first show that the Dirichlet hypotheses are
inherited by unions of whole components. The operators with Dirichlet boundary conditions are
introduced in \cite[\S E]{B-CMP99b} in the following words (the letters are those of
\cite{B-CMP99b}; in the application below Ba{\l}aban's level $k$ is the present level $j$, and
$\Lambda=\Lambda'$):
\begin{quote}
We need operators with Dirichlet boundary conditions outside some domain $\Lambda$ of the unit
lattice. We assume that $\Lambda\subset\Lambda_k=\Omega_k^{(k)}$, $\Lambda$ is a union of big blocks,
and a distance between $\Lambda$ and $\Lambda_k^c$ is bigger than $RM$.
\end{quote}
Connectedness of $\Lambda$
is not among these assumptions. Each assumption is a condition on the points, or on the big
blocks, of $\Lambda$ taken separately:
\begin{itemize}
\item containment in $\Lambda_k$ holds for a union as soon as it holds for each member;
\item the distance to $\Lambda_k^c$ is the minimum of the distances of the members;
\item the assumption that $\Lambda$ be a union of big blocks passes, like the other two, to any
sub-union of whole components' integration regions.
\end{itemize}
The integration region of the whole stage satisfies these assumptions, since it is the region at
which the normalisation lemma of the stage applies \cite[Theorems~3.10, 3.15]{B-CMP99b}, and it is
itself a union of whole components' integration regions. Since each assumption is a condition on
points or big blocks taken separately, every sub-union $\Lambda'$ of whole components' integration
regions satisfies them as well.

Next we identify the conditioned operators with the compression $\mathcal{A}(\Lambda')$. In
\cite[(3.156)]{B-CMP99b} the form is considered on the subspace of fields $B$ with $B=0$ on
$\Lambda^{c}$. Then \cite[(3.158)]{B-CMP99b} identifies $(C^{*}\Delta_kC)^{-1}$ on that subspace,
where $\Delta_k$ is the operator of the form considered in \cite[(3.156)]{B-CMP99b},
with the conditioned covariance $C^{(k)}(\Lambda)$. Thus conditioning to $\Lambda$ is the
compression of the precision to the bonds of $\Lambda$, with the fields off $\Lambda$ held at $0$.
This is the convention of \cite[(3.185)]{B-CMP99b} and of \cite[p.~1]{B-CMP116} for integrals
conditioned to subdomains of the lattice. It agrees with Lemma~\ref{4.24:lem-conditioning-marginal}(a):
at $\phi_W=0$ (the variables written $\phi_{\mathcal{W}}$ in that lemma) the precision of the
remaining variables is the principal block. The
stage covariance of Definition~\ref{4.24:def-stage-form} is read as a conditional covariance on a
subdomain in this sense. Applied with $\Lambda=\Lambda'$, this identifies the operator conditioned
to $\Lambda'$ with the compression $\mathcal{A}(\Lambda')$ of (b).

We now collect the walk expansion from the cited theorems. In \cite{B-CMP119},
\cite[Theorems~3.7--3.10]{B-CMP99b} are applied in their full generality, for families
$\sigma_0=\{\Delta\}$ of cubes
$\Delta$ in proper scales (here $\Delta$ denotes a cube, not the operator $\Delta_k$). With the
identification above, \cite[Theorem~3.10]{B-CMP99b} and
\cite[Theorem~3.15]{B-CMP99b} give the following.
\begin{itemize}
\item For every $\xi\ge0$, $(\mathcal{A}(\Lambda')+\xi)^{-1}$ has the convergent generalised random
walk expansion of \cite[(3.107)]{B-CMP99b}. The walks $\omega=(X_0,\dots,X_p)$ are walks of
localisation domains in proper scales, consecutive domains intersecting.
\item Each term $t^{\xi}_{\omega}$ depends on $\mathbb{U}$ only through its restriction to
$D(\omega)=X_0\cup\bigcup_{a\ge1}X_a^{\sim}$ \cite[Theorem~3.10]{B-CMP99b}.
\item Each term is bounded, in all the norms of \cite[(3.42)--(3.47)]{B-CMP99b}, by $m(\omega)$ of
\eqref{4.24:eq-region-walk-expansion-a}, where $O(1)$, $\delta_0$ and the covariance constant $B_0$
are numbers depending on
$(d,L)$ (\cite[Theorem~3.10]{B-CMP99b}, \cite[Theorem~3.15]{B-CMP99b}).
\item The terms are holomorphic on the complex extension constructed in \cite{B-CMP99b}.
\end{itemize}
The hypotheses of these theorems are met as follows.
\begin{itemize}
\item $M$ sufficiently large is the largeness condition imposed on $M$ when the auxiliary
parameters $\mathfrak{p}$ are chosen (Notation~\ref{2.1:not-prefix}), and $M\alpha_0$
sufficiently small is the smallness condition on $\gamma$.
\item The hypothesis that $\mathbb{U}$ satisfy \cite[(3.95)]{B-CMP99b} on the cubes visited is, by
\cite{B-CMP119}, the condition that the background pair of the stage belong to its space of pairs
$\mathfrak{K}_P$.
\end{itemize}
It remains to check that the constants are uniform. By \cite{B-CMP99b}, they do not depend on the
sequence $\{\Omega_j\}$. The statement of \cite[Theorem~3.15]{B-CMP99b} displays no hypothesis on
the domain of conditioning beyond those quoted above, and these hold for $\Lambda'$ by the
inheritance argument. Hence the same constants serve for every union $\Lambda'$ of whole
components. This proves (b).

\emph{Part (c): the $\xi$-integrated terms.} In \cite{B-CMP119} the operators entering the bond
density of \cite[(3.36)--(3.37)]{B-CMP119} are represented by the generalised random walk
expansions. After the integration in $\xi$, one obtains a representation of the form
\cite[(3.47)]{B-CMP119} whose terms satisfy \cite[(2.42)]{B-CMP119}. In this representation the
class $|\omega|=0$ is paired with $(\lambda_0+\xi)^{-1}\cdot1$ as in \cite{B-CMP119}. The bound
\cite[(2.42)]{B-CMP119} on the $\xi$-integrals of the terms $\hat t^{\xi}_{\omega}$ is a $\xi$-free
majorant $\bar m(\omega)$ of the form of $m(\omega)$. This is \eqref{4.24:eq-region-walk-expansion-b}.
The $\xi$-integration of \cite{B-CMP119} uses only the expansion and the bounds
\eqref{4.24:eq-region-walk-expansion-a}, which by (b) hold at every union $\Lambda'$ of whole
components with the same constants. Hence \eqref{4.24:eq-region-walk-expansion-b} holds at every
such $\Lambda'$. This proves (c).

\emph{Part (d).} By (a) and Definition~\ref{4.24:def-survivors-bonds}, $Z^{s}$ is a union of whole
components of the stage. Hence $\Lambda'_s$ is a union of whole components' stage-$(n+1)$
integration regions, and (b) and (c) apply at $\Lambda'=\Lambda'_s$. The integration region of the
whole stage is itself a union of whole components. So \eqref{4.24:eq-region-walk-expansion-a} and
\eqref{4.24:eq-region-walk-expansion-b} are exactly the inputs on which the normalisation lemma of
the stage depends, at every stage. At $\Lambda'_s$ they are the same quantified statements, with the
same constants. Their use at $\Lambda'_s$ therefore introduces no hypothesis beyond those under
which the normalisation lemma of the stage uses them.
\end{proof}

\begin{lemma}[The survivors' block as a conditioned precision]\label{4.24:prf-block-conditioned}
Let $\mathcal{A}$ be the symmetrised operator of the stage quadratic form of
Definition~\ref{4.24:def-stage-form}. Write $\mathcal{A}_{ss}$, $\mathcal{A}_{s\mathcal{W}}$ and
$\mathcal{A}_{\mathcal{W}\mathcal{W}}$ for its blocks with respect to the splitting
$\mathfrak{Bo}=s\sqcup\mathcal{W}$ of the level-$j$ integration bonds of $Z^{\sharp}$ into the
survivors' bonds $s$ and the arrested bonds $\mathcal{W}$. Let $\Lambda'_s$ be the survivors'
integration region, and let $\mathcal{A}(\Lambda'_s)$ be the compressed precision at the region
$\Lambda'_s$ in the walk expansion on unions of whole components,
\eqref{4.24:eq-region-walk-expansion-a}. Then
\begin{equation}\label{4.24:eq-block-conditioned-1}
\mathcal{A}_{ss}=\mathcal{A}(\Lambda'_s),
\end{equation}
and $\mathcal{C}_s=\mathcal{A}_{ss}^{-1}$ is the covariance $C^{(j)}(\Lambda'_s)$ of
\cite[Theorem~3.15]{B-CMP99b}. Moreover:
\begin{enumerate}
\item[(a)] for real $\phi_s$, the members $e_a^{0}(b,b')$ and $e_a(b,b',Y')$ of
\eqref{4.24:eq-stage-form-a} with $b,b'\in s$ sum, on the collar, to
$-\tfrac12\langle\phi_s,\mathcal{A}_{ss}\phi_s\rangle$;
\item[(b)] the rim part of $\mathcal{A}_{ss}$ is the per-component rim form of
\cite{B-CMP119} on $\mathsf{R}_j\cap Z^{s}$, and it does not depend on $\phi_W$.
\end{enumerate}
\end{lemma}

\begin{proof}
Write $\Delta_{\mathrm{glob}}$ for the operator of the global fluctuation form of the
$(j+1)$-st step of the operation $\mathbf{T}$. Write $C$ for the linear map through which the
integration variables enter the operator $\Delta^{(j)}$ of the stage form, so that the exponent in
Definition~\ref{4.24:def-stage-form} has the form $C^{*}\Delta^{(j)}C$.
In these terms the kernel $K_{\sharp}$ of Definition~\ref{4.24:def-stage-form} is
$K_{\sharp}=C^{*}\Delta^{(j)}C$; by the first paragraph below it is the compression
$P_{\mathfrak{Bo}}\Delta_{\mathrm{glob}}P_{\mathfrak{Bo}}$ of the global form to the bonds
$\mathfrak{Bo}$ of the stage; and $\mathcal{A}=\mathcal{S}K_{\sharp}$ with
$\mathcal{S}K=\tfrac12(K+K^{T})$. Since $P_E^{T}=P_E$ for every set of bonds $E$, we have
$(P_EKP_E)^{T}=P_EK^{T}P_E$, hence $\mathcal{S}(P_EKP_E)=P_E(\mathcal{S}K)P_E$: symmetrisation
commutes with compression.

\emph{The stage Gaussian is a compression.} The Gaussian \cite[(1.60)]{B-CMP122a} of the stage is
the global fluctuation form of the $(j+1)$-st step of $\mathbf{T}$, compressed to the variables
that remain for this stage. The reason is that the variables not integrated at this stage are
parameters of the integrand: in the notation of Definition~\ref{4.24:def-stage-form}, the fields
of the levels exterior to this stage have not yet been integrated, and the field of level $k$
enters only as a parameter of the integrand. The dependence on these parameters that the
integration induces is carried by the boundary terms $\mathbf{B}$, in the form stated in
\cite{B-CMP119}. By
Definition~\ref{4.24:def-stage-form}, the covariance of the stage is a conditional covariance on a
subdomain in the sense of \cite[Theorem~3.15]{B-CMP99b}.

\emph{Conditioning compresses the precision.} Apply the statement on conditional and marginal
precision of a Gaussian, clause (a) of
\eqref{4.24:eq-conditioning-marginal-a} (Lemma~\ref{4.24:lem-conditioning-marginal}), to
$\phi=\phi_s\oplus\phi_W$ at $\phi_W=0$. The density in $\phi_s$ then has precision
$P_s\mathcal{A}P_s=\mathcal{A}_{ss}$. It is not the Schur complement, which is the precision of
the marginal.

\emph{Transitivity of compression.} Since $s\subset\mathfrak{Bo}$, the transitivity of
compression in Lemma~\ref{4.24:lem-principal-blocks}, applied with
$R=\mathfrak{Bo}$ and $K=\Delta_{\mathrm{glob}}$, gives
\begin{equation}\label{4.24:eq-block-conditioned-2}
P_s\bigl(P_{\mathfrak{Bo}}\,\Delta_{\mathrm{glob}}\,P_{\mathfrak{Bo}}\bigr)P_s
=(\Delta_{\mathrm{glob}})_{ss}.
\end{equation}
Compressing first to the bonds $\mathfrak{Bo}$ of the stage and then to $s$ is therefore the same
as compressing the global form directly to $s$.

\emph{The block in the region's variables.} Let $\bar s$ be the set of all bonds of $\Lambda'_s$,
none excluded. According to the splitting $\phi=\phi_s\oplus\phi_W$, the map $C$ is the direct
sum $C=C_s\oplus C_{\mathcal{W}}$ of a map $C_s$, which acts on $\phi_s$ alone and takes values
on the bonds $\bar s$, and a map $C_{\mathcal{W}}$, which acts on $\phi_W$ alone. Hence $CP_s=C_s$
has range in $\mathcal{H}_{\bar s}$, and
\begin{equation}\label{4.24:eq-block-conditioned-3}
\bigl(C^{*}\Delta^{(j)}C\bigr)_{ss}=C_s^{*}\,\Delta^{(j)}_{\bar s\bar s}\,C_s .
\end{equation}

\emph{Identification.} By \eqref{4.24:eq-block-conditioned-2} and
\eqref{4.24:eq-block-conditioned-3}, $\mathcal{A}_{ss}$ is the precision compressed to the region
$\Lambda'_s$. This is \eqref{4.24:eq-block-conditioned-1}: $\mathcal{A}_{ss}$ is the operator
$\mathcal{A}(\Lambda'_s)$ of \eqref{4.24:eq-region-walk-expansion-a}. It is the same kind of
object as the stage operator $\mathcal{A}=\mathcal{A}(\Lambda'_{Z^{\sharp}})$, taken at the region
$\Lambda'_s$ instead of $\Lambda'_{Z^{\sharp}}$. Consequently
$\mathcal{C}_s=\mathcal{A}_{ss}^{-1}$ is the covariance $C^{(j)}(\Lambda'_s)$ of
\cite[Theorem~3.15]{B-CMP99b}, and so is itself a conditional covariance on the subdomain
$\Lambda'_s$.

\emph{Clause} (a). The expression \eqref{4.24:eq-stage-form-a} of the stage form, clause (iii),
is termwise in the pair $(b,b')$. Take real $\phi_s$, and keep the members $e_a^{0}(b,b')$ and
$e_a(b,b',Y')$ with $b,b'\in s$. By clauses (a) and (h) of the polarisation of a bond kernel,
\eqref{4.24:eq-kernel-polarisation-a}, these members sum on the collar to
$-\tfrac12\langle\phi_s,\mathcal{A}_{ss}\phi_s\rangle$. Thus the survivors' collar Gaussian is
left untouched at this stage.

\emph{Clause} (b). The rim part of $\mathcal{A}_{ss}$ is the per-component rim form of
\cite{B-CMP119}, taken on $\mathsf{R}_j\cap Z^{s}$. As a block of $\mathcal{A}$ on survivors'
bonds, it reads no arrested variable $\phi_W$.
\end{proof}

\begin{lemma}[The only determinant of the stage]\label{4.24:prf-single-determinant}
Work with the stage quadratic form of Definition~\ref{4.24:def-stage-form} and its blocks
\eqref{4.24:eq-stage-form-a}. Recall that $\mathfrak{Bo}=s\sqcup\mathcal{W}$ is the set of
level-$j$ integration bonds of $Z^{\sharp}$, split into the survivors' bonds $s$ and the arrested
bonds $\mathcal{W}$, with variables $\phi_s$ and $\phi_W$. Let $\mathcal{N}=n_{\mathfrak{g}}|s|$
and $\mathcal{A}_{ss}=P_s\mathcal{A}P_s$. Then:
\begin{enumerate}
\item[(a)] no element $e^{0}_a(b,b')$ of the collar part of \eqref{4.24:eq-stage-form-a} has one
bond in $s$ and the other in $\mathcal{W}$, and none joins two distinct arrested components;
\item[(b)] every element $e_a(b,b',Y')$ of that collar part with one bond in $s$ and the other in
$\mathcal{W}$ satisfies $|Y'|\ge R_k$, and it is a member of the family $\mathfrak{E}^{\times}$ of
Definition~\ref{4.24:def-survivor-elements};
\item[(c)] at fixed real $\phi_W$ the $d\phi_s$-Gaussian of the stage has precision
$\mathcal{A}_{ss}$; no Schur complement in $\phi_W$ is formed, and $\mathcal{A}_{ss}$ is the
conditional, not the marginal, precision of Lemma~\ref{4.24:lem-conditioning-marginal}; the only
determinant of the stage is $\det(\mathcal{A}_{ss})^{-1/2}$, and the normalisation factor
$(2\pi)^{\mathcal{N}/2}\det(\mathcal{A}_{ss})^{-1/2}$ of that Gaussian is written without a choice
of branch as
\begin{equation}\label{4.24:eq-single-determinant-1}
(2\pi)^{\mathcal{N}/2}\,e^{-\frac12\operatorname{Tr}F(\mathcal{A}_{ss})},
\end{equation}
with $F$ as in \eqref{4.24:eq-log-resolvent-a}.
\end{enumerate}
\end{lemma}

\begin{proof}
No Schur complement is formed at this stage. The restriction of $\phi_W$ to the bonds of each
arrested component in $\mathfrak{W}$ is an exterior real box coordinate, in the sense of the
statement fixing the exterior coordinate of an arrested component;
\cite[(1.59)]{B-CMP122a} states that this coordinate does not depend at all on the background
fields. The level-$j$ variables of $\mathcal{W}$ are not integrated at this stage.

The collar part of \eqref{4.24:eq-stage-form-a} resolves the collar form of $Z^{\sharp}$ into
ordered pairs of bonds of three kinds. In the first kind both bonds lie in $s$, in the second both
lie in $\mathcal{W}$, and in the third, the mixed pairs, one bond lies in $s$ and the other in
$\mathcal{W}$. We treat the elements of the collar part in turn.

\emph{Part (a).} Let $C$ be the averaging operator appearing in \eqref{4.24:eq-stage-form-a}, so
that the kernel $(C^{*}C)(b,b')$ is the one appearing there. An element $e^{0}_a(b,b')$ is present
only if $(C^{*}C)(b,b')\ne0$. By \eqref{4.24:eq-stage-form-a} this means that $b$ and $b'$ lie in
one block pair, and such a pair is contained in at most two wall-adjacent $\pi_{j+1}$-cells. By
\eqref{4.24:eq-component-separation-c}, in both of the forms given there, two wall-adjacent
$\pi_{j+1}$-cells meeting the integration bonds meet the bonds of one component only. Hence $b$ and
$b'$ belong to one component. Consequently no $e^{0}_a$ is mixed, and no $e^{0}_a$ joins two
arrested components.

\emph{Part (b).} Let $e_a(b,b',Y')$ be mixed, with $b\in s$ and $b'\in\mathcal{W}$; for a bond $c$
write $h(c)$ for its host cube, as in Lemma~\ref{4.24:lem-component-separation}. Then $Y'$ is a
wall-connected family of host cubes containing
$h(b)$ and $h(b')$. Thus $Y'$ meets the integration bonds of two distinct components, and
\eqref{4.24:eq-component-separation-a} gives $|Y'|\ge R_k$.

By clause~(1) of the classification of collar elements, $e_a$ is a mixed element. It is bilinear in
the pair $(\phi_s(b),x(b'))$, where $x(b')$ is the coordinate of $b'$ in the real box of its
arrested component. That box lies in the disc of radius $\varrho^{B}_j$. The majorant
$\nu^{\natural}(e_a)$ is the one given in \eqref{4.24:eq-stage-form-a}. These are exactly the
members of the family $\mathfrak{E}^{\times}$ of Definition~\ref{4.24:def-survivor-elements}.

\emph{The pairs inside $\mathcal{W}$.} The terms with both bonds in $\mathcal{W}$ are of two kinds:
those inside one arrested component $W$, which together make up its collar quadratic
$\mathsf{q}_W$, so that these terms sum to $\sum_{W}\mathsf{q}_W$; and the elements $e_a$ joining
two distinct arrested components, which have $|Y'|\ge R_k$ by
\eqref{4.24:eq-component-separation-a}. All of them are frozen, by clause~(2) of
the classification of collar elements and by the freezing step for the collar forms of the
arrested components.

\emph{The rim.} The entries of the rim $\mathsf{R}_j$ between distinct components are Ba{\l}aban's
$\mathbf{B}$-terms. They are the parent terms of the bound displayed in clause~(b) of the
proposition on the boundary terms $\mathbf{B}$, and they are never entries of $\mathcal{A}_{ss}$.

\emph{Part (c).} Combining parts (a), (b) and the two preceding paragraphs, the exponent of the
term of the stage integrand carrying the form of Definition~\ref{4.24:def-stage-form}, as a
function of $\phi_s$ at fixed real $\phi_W$, is
\begin{equation}\label{4.24:eq-single-determinant-2}
-\tfrac12\langle\phi_s,\mathcal{A}_{ss}\phi_s\rangle
+\sum(\text{elements, some of which read }\phi_W)
+(\text{terms independent of }\phi_s).
\end{equation}
Nothing in \eqref{4.24:eq-single-determinant-2} is completed to a square, so clause~(c) of
Lemma~\ref{4.24:lem-conditioning-marginal} is not used. We are in the situation of clause~(a) of
Lemma~\ref{4.24:lem-conditioning-marginal}: the $d\phi_s$-Gaussian has precision
$\mathcal{A}_{ss}=P_s\mathcal{A}P_s$. It is not the marginal precision of clause~(b) of
Lemma~\ref{4.24:lem-conditioning-marginal}. That precision would arise only if $\phi_W$
were integrated as a Gaussian variable at this stage, and it is not.

The normalisation of this Gaussian therefore contains the only determinant of the stage,
$\det(\mathcal{A}_{ss})^{-1/2}$, and equals $(2\pi)^{\mathcal{N}/2}\det(\mathcal{A}_{ss})^{-1/2}$,
with $\mathcal{N}=n_{\mathfrak{g}}|s|$. By Lemma~\ref{4.24:lem-log-resolvent},
$(2\pi)^{\mathcal{N}/2}\det(\mathcal{A}_{ss})^{-1/2}$ equals
\eqref{4.24:eq-single-determinant-1}, without a choice of branch.

\emph{No Jacobi or Schur factorisation.} The factorisation
\begin{equation}\label{4.24:eq-single-determinant-3}
\det\mathcal{A}_{ss}=\det\mathcal{A}\cdot\det\bigl((\mathcal{A}^{-1})_{\mathcal{W}\mathcal{W}}\bigr)
\end{equation}
is not used. Using it would introduce normalisation terms anchored on the arrested components,
belonging to an integration over $\phi_W$ that is not performed at this stage.
\end{proof}

\begin{lemma}[The normalisation as a sum of bond densities]\label{4.24:prf-bond-densities}
Assume the hypotheses of Lemma~\ref{4.24:lem-block-coercivity}.
Let $\lambda_0>0$, the factors $z^{(j)}(c)$ and their product
$z^{(j)}_s=\prod_{c\subset R^{(n+1)}\cap Z^s}z^{(j)}(c)$ be as in
Definition~\ref{4.24:def-normalisation-constants}. Let $\mathcal{A}_{ss}$ be the block of $\mathcal{A}$
on the survivors' bonds $s$, as in Lemma~\ref{4.24:lem-principal-blocks}, and let $F$ be the function
\eqref{4.24:eq-log-resolvent-a} of Lemma~\ref{4.24:lem-log-resolvent}.
Then
\begin{multline}\label{4.24:eq-bond-densities-1}
z^{(j)}_s\int d\phi_s\,e^{-\frac12\langle\phi_s,\mathcal{A}_{ss}\phi_s\rangle}\\
=\exp\Bigl[-\tfrac12\operatorname{Tr}F(\mathcal{A}_{ss})
+\sum_{c\subset R^{(n+1)}\cap Z^s}\log z^{(j)}(c)+\mathrm{const}_s\Bigr],
\end{multline}
where $\mathrm{const}_s$ is a number, independent of the fields, and
\begin{equation}\label{4.24:eq-bond-densities-2}
-\tfrac12\operatorname{Tr}F(\mathcal{A}_{ss})=\sum_{b'\in s}\mathfrak{e}(b'),
\end{equation}
with, for $b'\in s$,
\begin{equation}\label{4.24:eq-bond-densities-a}
\mathfrak{e}(b'):=-\tfrac12 n_{\mathfrak{g}}\log\lambda_0
-\tfrac12\int_0^\infty d\xi\,\operatorname{tr}_{\mathfrak{g}}
\bigl[(\lambda_0+\xi)^{-1}1-(\mathcal{A}_{ss}+\xi)^{-1}\bigr](b',b').
\end{equation}
\end{lemma}

\begin{proof}
\emph{The precision of the integral.} The variable $\phi_s$ ranges over the coordinate subspace
$\mathcal{H}_s=P_s\mathcal{H}$. For $\phi\in\mathcal{H}_s$ the Rayleigh identity of
Lemma~\ref{4.24:lem-principal-blocks} gives $\langle\phi,\mathcal{A}_{ss}\phi\rangle=\langle\phi,\mathcal{A}\phi\rangle$,
so the quadratic form in the exponent on the left of \eqref{4.24:eq-bond-densities-1} is the form of
$\mathcal{A}$ restricted to the survivors' variables, and its operator on $\mathcal{H}_s$ is $\mathcal{A}_{ss}$.
The same lemma gives $(\mathcal{A}_{ss})^T=\mathcal{A}_{ss}$ and $(\mathcal{A}+\xi)_{ss}=\mathcal{A}_{ss}+\xi 1_{\mathcal{H}_s}$
for $\xi\ge0$. By Lemma~\ref{4.24:lem-block-coercivity}, whose hypotheses are assumed,
$\mathcal{A}_{ss}$ is coercive: $\mathcal{A}_{ss}\in\mathfrak{P}(a_{\mathbb{C}})$ with
$a_{\mathbb{C}}=(4B_0)^{-1}$, so in particular $\mathcal{A}_{ss}\in\mathfrak{P}$, the domain on which $F$ is defined.

\emph{The normalisation formula.} In \cite[(3.33)]{B-CMP119} the normalising integral is written, in the
notation used there, in the form
\begin{equation*}
\exp\Bigl[-\tfrac12\log\det\bigl(C^*\Delta^{(k)}C\bigr)+\sum_c\log z^{(k)}(c)+\mathrm{const}\Bigr],
\end{equation*}
where $C$ is the map by which the $\delta$-functions are eliminated; the same place states that the first
term already has the localised form.
We read this at level $j$ for the survivors'
integration region $\Lambda'_s$. Read for $\Lambda'_s$, the quadratic form $C^*\Delta^{(k)}C$ of
\cite[(3.33)]{B-CMP119} is the form $\phi_s\mapsto\langle\phi_s,\mathcal{A}_{ss}\phi_s\rangle$ of the
integral on the left of \eqref{4.24:eq-bond-densities-1}, by the description of the elimination of the
$\delta$-functions on $\Lambda'_s$ and of the stage form kernel $K_{\sharp}$. The cells over which
the sum runs are the cells $c\subset R^{(n+1)}\cap Z^s$, and the corresponding product of the factors
$z^{(j)}(c)$ is $z^{(j)}_s$, by Definition~\ref{4.24:def-normalisation-constants}; this is the prefactor
on the left of \eqref{4.24:eq-bond-densities-1}, and on the right it contributes
$\sum_{c\subset R^{(n+1)}\cap Z^s}\log z^{(j)}(c)$. By the first paragraph the precision of the
Gaussian integral is $\mathcal{A}_{ss}$, which lies in $\mathfrak{P}$. By the Gaussian identity of
Lemma~\ref{4.24:lem-log-resolvent}, the Gaussian integral with a precision $B\in\mathfrak{P}$ is
$(2\pi)^{\mathcal{N}/2}\exp[-\frac12\operatorname{Tr}F(B)]$ with
$\mathcal{N}=\dim\mathcal{H}_s=n_{\mathfrak{g}}|s|$; we apply
it with $B=\mathcal{A}_{ss}$. Thus the term $-\frac12\log\det$ is realised by
$-\frac12\operatorname{Tr}F(\mathcal{A}_{ss})$, and \eqref{4.24:eq-bond-densities-1} holds with
$\mathrm{const}_s$ equal to $(\mathcal{N}/2)\log 2\pi$ plus the constants of the measure of
\cite[(3.33)]{B-CMP119} for the cells of $Z^s$. These are numbers, and they are distributed over the
cells and bonds of $Z^s$.

The contour argument of \cite{B-CMP119} that accompanies the logarithm of the determinant uses of the
integrand only its decay like the inverse square of the spectral variable. For $B=\mathcal{A}_{ss}$ this
decay is supplied by the well-definedness part of Lemma~\ref{4.24:lem-log-resolvent}: for
$B\in\mathfrak{P}(a)$,
\begin{equation*}
(\lambda_0+\xi)^{-1}1-(B+\xi)^{-1}=(B+\xi)^{-1}(B-\lambda_0)(\lambda_0+\xi)^{-1},
\end{equation*}
whose norm is at most $\|B-\lambda_0\|/((a+\xi)(\lambda_0+\xi))$, which is $O(\xi^{-2})$ as
$\xi\to\infty$.

\emph{The trace as a sum over bonds.} By \eqref{4.24:eq-log-resolvent-a},
\begin{equation*}
F(\mathcal{A}_{ss})=\log\lambda_0\cdot 1
+\int_0^\infty d\xi\,\bigl[(\lambda_0+\xi)^{-1}1-(\mathcal{A}_{ss}+\xi)^{-1}\bigr],
\end{equation*}
the integral converging in norm by the bound just recalled. The space $\mathcal{H}_s$ is the direct sum,
over $b'\in s$, of one copy of $\mathfrak{g}$ (complexified) attached to $b'$. For an operator $K$ on
$\mathcal{H}_s$ write $K(b',b')$ for its diagonal block on the copy attached to $b'$; then
$\operatorname{Tr}K=\sum_{b'\in s}\operatorname{tr}_{\mathfrak{g}}K(b',b')$, because the trace is the sum
of the diagonal matrix elements in a basis adapted to the direct sum. The map
$K\mapsto\operatorname{tr}_{\mathfrak{g}}K(b',b')$ is a finite linear combination of matrix elements of
$K$, hence commutes with the norm-convergent integral over $\xi$. Applied to the two terms of
$F(\mathcal{A}_{ss})$ this gives, first,
$\operatorname{Tr}(\log\lambda_0\cdot1)=\sum_{b'\in s}n_{\mathfrak{g}}\log\lambda_0$, since
$\operatorname{tr}_{\mathfrak{g}}1=n_{\mathfrak{g}}=\dim\mathfrak{g}$; and, second,
\begin{align*}
&\operatorname{Tr}\int_0^\infty d\xi\,\bigl[(\lambda_0+\xi)^{-1}1-(\mathcal{A}_{ss}+\xi)^{-1}\bigr]\\
&\qquad=\sum_{b'\in s}\int_0^\infty d\xi\,\operatorname{tr}_{\mathfrak{g}}
\bigl[(\lambda_0+\xi)^{-1}1-(\mathcal{A}_{ss}+\xi)^{-1}\bigr](b',b').
\end{align*}
Multiplying by $-\frac12$ and comparing with \eqref{4.24:eq-bond-densities-a} gives
\eqref{4.24:eq-bond-densities-2}.

Thus the right side of \eqref{4.24:eq-bond-densities-2} carries exactly one density for each bond
$b'\in s$ and none for the arrested bonds of $\mathcal{W}$. These are the terms of
\cite[(3.36)]{B-CMP119} and the resolvent term of \cite[(3.37)]{B-CMP119}, read with the set of bonds
$\Lambda^{(k)*}$ indexing $b'$ there replaced by $s$. The terms of \cite[(3.37)]{B-CMP119} carrying an
integral over $t$ are not members of \eqref{4.24:eq-bond-densities-1}.

\emph{Assignment of the members to domains.} The first term $-\frac12n_{\mathfrak{g}}\log\lambda_0$ of
$\mathfrak{e}(b')$ is a field-free number; it is the fibre trace at $b'$ of the term of
$-\frac12F(\mathcal{A}_{ss})$ that \cite{B-CMP119} writes as
$-\frac12\log\lambda_0 I$, and $\lambda_0$ does not depend on the region
(Definition~\ref{4.24:def-normalisation-constants}). For each $b'\in s$ we assign it, together with the
part of $\mathrm{const}_s$ attached to $b'$, to the domain $X=[\square(b')]_{j+1}$, the smallest
$\mathbb{D}_{j+1}$-domain containing the cube $\square(b')$ of the bond $b'$, as a single-cube term. This
is the same assignment as in the lemma on the Gaussian normalisation over the whole stage, whose constants
are those of Definition~\ref{4.24:def-normalisation-constants}: there the additive constant
$\mathrm{const}$ of \cite[(3.33)]{B-CMP119} and
the vacuum-energy counterterm $E^{(j)}_0$ are single-cube terms of $\mathbb{D}_{j+1}$. Together with the
value of the counterterm $E^{(j)}_0$ of Definition~\ref{4.24:def-normalisation-constants} on the set
$(\Omega^c_{j+1}\setminus Z''_{j+1})\cap Z^s$, the part of the survivors' region $Z^s$ that lies in the
small-field complement $\Omega^c_{j+1}$ and outside the region $Z''_{j+1}$, these numbers contribute to the
vacuum-energy renormalisation, as stated in \cite{B-CMP119}, that is, to the vacuum-energy term written
$E_k(Z)$ in \cite{B-CMP122b}, where $k$ denotes the level in Ba{\l}aban's notation.

The remaining members of the objects of Definition~\ref{4.24:def-normalisation-constants} are those of
\cite[(3.33)]{B-CMP119}: numbers, and local holomorphic functions attached to bonds and blocks of the
collar (the Jacobians of the elimination of the $\delta$-functions and the normalisations of the gauge
fixing), all single-cube terms of $\mathbb{D}_{j+1}$. Here they are indexed by $b'\in s$ and by the cells
$c\subset R^{(n+1)}\cap Z^s$. Each of them is the same local function as in the full stage; in
particular $z^{(j)}(c)$ depends on $\mathbb{U}$ only through its values on the block pair of $c$. Each is
assigned to the domain $X$ equal to the smallest $\mathbb{D}_{j+1}$-domain containing its block pair.
Since $R_n\subset\Omega^c_{j+2}$, such a domain is of level $j+1$ and $d_{j+1}(X)\le1$.
\end{proof}

\begin{lemma}[Termwise walk expansion of the block resolvent]\label{4.24:prf-block-walk-expansion}
Work at the stage $n+1$ of the construction.

\begin{itemize}
\item Let $s$ be the set of survivors' bonds and $\Lambda'_s$ the survivors' integration region.
\item Let $\mathcal{A}$ be the symmetrised operator of the stage form kernel $K_{\sharp}$, and $\mathcal{A}_{ss}$ its block on $s$.
\item Let $(\mathbb{U},\mathbb{J})$ be the complex background pair of the stage.
\item Let $\{\Omega_m\}$ be the sequence of small-field domains fixed for the whole class of configurations under consideration.
\end{itemize}

Then for every $\xi\ge0$ and every $b'\in s$
\begin{equation}\label{4.24:eq-block-walk-expansion-1}
(\mathcal{A}_{ss}+\xi)^{-1}(b',b')=\sum_{\omega}t^{\xi}_{\omega}(b',b'),
\end{equation}
where the sum runs over walks $\omega=(X_0,\dots,X_p)$ of domains with the following properties:
\begin{itemize}
\item $b'\in X_0$;
\item $X_{i-1}\cap X_i\neq\emptyset$ for $1\le i\le p$;
\item every $X_i$ is a domain in a proper scale of the sequence $\{\Omega_m\}$.
\end{itemize}
Each term $t^{\xi}_{\omega}$ has the following properties:
\begin{itemize}
\item[(a)] it depends on $\mathbb{U}$ only through its restriction to the domain $D(\omega)$ of the walk, and it is holomorphic there;
\item[(b)] it is covariant under the adjoint action of gauge transformations, being a product of factors each of which is so covariant;
\item[(c)] it obeys $|t^{\xi}_{\omega}|\le m(\omega)$, where $m(\omega)$ is the majorant attached to $\omega$ in Lemma~\ref{4.24:prf-region-walk-expansion}. The constants of this majorant depend only on $(d,L)$; they depend neither on the sequence $\{\Omega_m\}$ nor on the region $\Lambda'_s$.
\end{itemize}
\end{lemma}

\begin{proof}
\emph{The operator.} By Lemma~\ref{4.24:prf-block-conditioned}, which presents the survivors' block as a conditioned precision, we have
\begin{equation*}
(\mathcal{A}_{ss}+\xi)^{-1}=(\mathcal{A}(\Lambda'_s)+\xi)^{-1}.
\end{equation*}
This is the $\xi$-shifted conditional covariance of the region $\Lambda'_s$ over which the survivors' variables are integrated at this stage in Ba{\l}aban's representation of the stage's integral. It is therefore the object of \cite[Theorem~3.15]{B-CMP99b} with that theorem's parameter $\lambda$ equal to $0$.

Consequently the walk expansion on unions of whole components, Lemma~\ref{4.24:prf-region-walk-expansion} with the expansion \eqref{4.24:eq-region-walk-expansion-a}, applies to this operator once its three hypotheses are verified. These are:
\begin{itemize}
\item the conditions on $M$ and on $M\alpha_0$;
\item the regularity condition \cite[(3.95)]{B-CMP99b};
\item the admissibility of a complex background.
\end{itemize}
We verify them one at a time. The mapping of each hypothesis is the same as in the localisation of $-\tfrac12\operatorname{Tr}F$ through the resolvent's walk expansion, where the region is $\Lambda'_{Z^{\sharp}}$.

\emph{The conditions on $M$.} Lemma~\ref{4.24:prf-region-walk-expansion} requires two things:
\begin{itemize}
\item $M$ sufficiently large. This holds by the lower bound $M\ge M_0(\kappa,L,d)$ of the largeness condition on $M$ under which the localisation of $-\tfrac12\operatorname{Tr}F$ is carried out; here $M_0(\kappa,L,d)$ is the threshold of that condition, which depends only on $\kappa$, $L$ and $d$.
\item $M\alpha_0$ sufficiently small. This follows from the choice of $\gamma$.
\end{itemize}

\emph{The regularity condition.} The condition \cite[(3.95)]{B-CMP99b} is regularity of the background at the scale of the shells. Equivalently, on the cubes that a term of the expansion visits, the background belongs to the space $\tilde U^{(n+1)c,\Sigma}_k$ of regular configurations of the stage, in the sense of \cite{B-CMP119}. We verify this cube by cube.

\begin{itemize}
\item \emph{Cubes meeting $Z^{s}$.} Here the required membership is membership in $\mathfrak{N}^{(n+1)}$. That the walks of the stage live on $\mathfrak{N}^{(n+1)}$, as in \cite[Theorem~3.10]{B-CMP99b}, and that $\mathfrak{N}^{(n+1)}$ is contained, through the corresponding space of the inductive hypothesis, in the space of pairs $\mathfrak{K}_P$ of \cite{B-CMP109},
\begin{equation*}
\mathfrak{K}_P=\mathfrak{K}(\mathbb{B}^{(n)}_k;\gamma_0),
\end{equation*}
on which every estimate used at this stage holds, are both part of the statement that fixes the analyticity domains of the stage.
\item \emph{Cubes of $Z_W^{+}$.} Here the required membership is the condition in force at this stage on the arrested component. At the arrest stage $n+1=n^{+}$ these are still Ba{\l}aban's conditions $\mathfrak{L}_{n+1}$; the arrest condition replaces them only at stages $l>n^{+}$. They hold together with the global complex pair $(\mathbb{U},\mathbb{J})$ of the stage.
\end{itemize}

Thus the regularity clauses required at each cube are the same clauses, at the same cubes, as for the full stage, whichever component the cube meets.

\emph{The complex background.} That the walk expansion holds for complex backgrounds of the class considered is the complex extension in \cite{B-CMP99b} together with \cite[p.~15]{B-CMP116}.

\emph{Conclusion.} With its three hypotheses verified, Lemma~\ref{4.24:prf-region-walk-expansion} applied to $(\mathcal{A}(\Lambda'_s)+\xi)^{-1}$ gives the expansion \eqref{4.24:eq-block-walk-expansion-1} for every $\xi\ge0$ and $b'\in s$. The walks have the stated properties:
\begin{itemize}
\item their first domain contains $b'$;
\item consecutive domains intersect;
\item every domain lies in a proper scale of $\{\Omega_m\}$.
\end{itemize}

The properties of the terms follow from that lemma:
\begin{itemize}
\item Property (a), locality on $D(\omega)$ and holomorphy there, is the locality and analyticity that Lemma~\ref{4.24:prf-region-walk-expansion} gives for each term.
\item Property (b) holds because each term is a product of factors, each of which is covariant under the adjoint action by clause (c) of the weighted localisation statement and by the covariance assertion of clause (i) of the operator statement of the stage; a product of covariant factors is covariant.
\item Property (c) is the majorant of that lemma. Its constants depend on $(d,L)$ only. In particular they do not depend on the sequence $\{\Omega_m\}$, and they do not depend on the region: the same constants serve for $\Lambda'_s$ as for any other union of whole components.
\end{itemize}

\emph{Terms and majorants.} The terms of the expansion \eqref{4.24:eq-block-walk-expansion-1} of $(\mathcal{A}_{ss}+\xi)^{-1}$ differ from the terms of the walk expansion of $(\mathcal{A}+\xi)^{-1}$ near the union of the arrested components $\mathfrak{W}$. The reason is that the local inverses entering them are those of the region $\Lambda'_s$, taken on its own cubes.

The two expansions have majorants of the same form $m(\omega)$ with the same constants; the expansion of $(\mathcal{A}_{ss}+\xi)^{-1}$ has fewer anchoring cubes, namely those of $\Lambda'_s$. Only the majorants enter what follows; it is not asserted that the terms coincide.

This suffices for the localisation of the bond density $\mathfrak{e}(b')$ of $-\tfrac12\operatorname{Tr}F(\mathcal{A}_{ss})$ by the resolvent's walk expansion, because that localisation uses only the following features of the expansion:
\begin{itemize}
\item the locality of each term on $D(\omega)$;
\item a majorant of the fixed form $m(\omega)$;
\item the covariance of each term;
\item the count of walks anchored in a given cube.
\end{itemize}

Finally, no interpolation parameter acts on the kernel $\mathcal{A}_{ss}$ inside $-\tfrac12\operatorname{Tr}F(\mathcal{A}_{ss})$: neither the decoupling parameters $\sigma'(\Delta)$ of the collar kernel nor the decoupling parameters $\mathsf{s}(\Delta)$ of the covariance's walks. Two facts give this:
\begin{itemize}
\item The localisation of $-\tfrac12\operatorname{Tr}F$ uses the walk expansion of the resolvent term by term, and each term is already local.
\item The parameters $\mathsf{s}(\Delta)$ decouple only the walks of the covariance, in the estimate where the covariance is localised.
\end{itemize}
Neither is a decoupling of a form appearing under a logarithm.
\end{proof}

\begin{lemma}[Hulls, filing, anchors and holomorphy of the densities]\label{4.24:prf-hulls-anchors}
Work at the stage fixed in Notation~\ref{4.24:not-scales-rim-collar}, with the survivors' bonds $s$, the
arrested bonds $\mathcal{W}$, the arrested components $\mathfrak{W}$, the complex background pair
$(\mathbb{U},\mathbb{J})$ and the class-wide regions $\Omega_m$ of the step; the filing level $m(\cdot)$
and the hull $[\,\cdot\,]_m$ are those of Notation~\ref{4.24:not-scales-rim-collar}. For $b'\in s$ let
$\operatorname{cube}(b')\in\mathsf{Q}^{s}$ be the $\pi_{j+1}$-cell containing $b'$, and let $\omega$ run over the walks of the
termwise walk expansion of the block resolvent (Lemma~\ref{4.24:prf-block-walk-expansion}), with domain
$D(\omega)$, length $d(\omega)$ (not to be confused with the dimension $d$) and terms $t^{\xi}_{\omega}$;
write $\omega_0$ for a walk of the class
$|\omega|=0$, and $\mathbf{1}_{\omega_0}$ for the share of the identity operator carried by $\omega_0$ in that
expansion. Put
\begin{equation}\label{4.24:eq-hulls-anchors-a}
\begin{gathered}
X(\omega):=[\operatorname{cube}(b')\cup D(\omega)]_m,\qquad m:=m\bigl(\operatorname{cube}(b')\cup D(\omega)\bigr),\\
\hat t^{\xi}_{\omega}:=t^{\xi}_{\omega}\quad(|\omega|\ge1),\qquad
\hat t^{\xi}_{\omega_0}:=t^{\xi}_{\omega_0}-(\lambda_0+\xi)^{-1}\mathbf{1}_{\omega_0}\quad(|\omega_0|=0),\\
\mathbb{C}^{s}_G(X):=\sum_{b'\in s}\ \sum_{\omega:\,X(\omega)=X}\frac12\int_0^\infty d\xi\,
\operatorname{tr}_{\mathfrak{g}}\hat t^{\xi}_{\omega}(b',b')\\
\qquad+\bigl(\text{the single-cube members filed at }X\bigr),
\end{gathered}
\end{equation}
where the single-cube members are those of the normalisation written as a sum of bond densities
(Lemma~\ref{4.24:prf-bond-densities}), each filed at its hull as there; among them, for each $b'\in s$ with
$[\operatorname{cube}(b')]_{j+1}=X$, is the number $-\tfrac12 n_{\mathfrak{g}}\log\lambda_0$ of
$\mathfrak{e}(b')$ in \eqref{4.24:eq-bond-densities-a}. The \emph{anchor} of a term of
\eqref{4.24:eq-hulls-anchors-a} is its index: the bond $b'\in s$ for a resolvent density, and the bond of $s$
or the cell $c\subset R^{(n+1)}\cap Z^{s}$ indexing a single-cube member; $X$ is \emph{anchored on} a set if
some term filed at $X$ has its anchor in that set. Then:
\begin{enumerate}
\item[(a)] every hull $X(\omega)$ is face-connected; every $X$ carrying a term of
\eqref{4.24:eq-hulls-anchors-a} meets $R_n\cap Z^{s}$, and no such $X$ is anchored on $\bigcup\mathfrak{W}$;
\item[(b)] the sum over $X$ of the terms in \eqref{4.24:eq-hulls-anchors-a} contains
$\sum_{b'\in s}\mathfrak{e}(b')$ in full; the constants $-\tfrac12 n_{\mathfrak{g}}\log\lambda_0$ do not occur
in $\mathbb{C}^{s}_G-\mathbb{C}^{s}_G|_{(1,0)}$, where $\mathbb{C}^{s}_G|_{(1,0)}$ denotes $\mathbb{C}^{s}_G$
evaluated at the pair $(\mathbb{U},\mathbb{J})=(1,0)$, and their total per cube lies within the single-cube
budget $c_{\mathrm{sc}}$, of size $O(1)\cdot(LM)^4$;
\item[(c)] $\mathbb{C}^{s}_G(X)$ depends on $(\mathbb{U},\mathbb{J})$ only through its values inside $X$;
\item[(d)] $\mathbb{C}^{s}_G(X)$ is holomorphic on $\mathfrak{D}^{\flat}_{n+1}(X)$, and hence on the subdomain
$\mathfrak{D}^{\flat}_{n+1}[(\mathrm{amb})_W](X)\subset\mathfrak{D}^{\flat}_{n+1}(X)$ attached to an arrested
component $W\in\mathfrak{W}$;
\item[(e)] each term $\operatorname{tr}_{\mathfrak{g}}\hat t^{\xi}_{\omega}(b',b')$ is gauge invariant;
\item[(f)] for every walk $\omega$ with hull $X=X(\omega)$ filed at level $m\ge j+1$,
\begin{equation}\label{4.24:eq-hulls-anchors-1}
e^{-\delta_0 d(\omega)}\le e^{-(\delta_0 M_1^{-1}M/\sqrt d)\,d_j([X]_j)+c_{\mathrm{cv}}},
\end{equation}
where $c_{\mathrm{cv}}$ is the additive constant of the conversion of walk length into distance in filing
units; the decay
rate in filing units is at least $L\delta_0 M/(\sqrt d M_1)$, and
\begin{equation}\label{4.24:eq-hulls-anchors-2}
\frac{L\delta_0 M}{\sqrt d\,M_1}\ \ge\ 8\kappa L\ \ge\ a_m+c_d+1\qquad\text{for every }a_m\le 4\kappa L ,
\end{equation}
where $c_d$ is the constant of the parameter conditions of
Definition~\ref{4.24:def-normalisation-constants}.
\end{enumerate}
\end{lemma}

\begin{proof}
\emph{Face-connectedness.} The hull $X(\omega)$ is the smallest $\mathbb{D}_m$-domain containing
$\operatorname{cube}(b')\cup D(\omega)$. Consecutive domains of the walk $\omega$ intersect, and the first of
them contains $b'$, hence meets $\operatorname{cube}(b')$; so $\operatorname{cube}(b')\cup D(\omega)$ is
connected, and the smallest $\mathbb{D}_m$-domain containing it is a face-connected union of $\pi_m$-cells.

\emph{The definition collects the bond densities.} By Lemma~\ref{4.24:prf-bond-densities},
$-\tfrac12\operatorname{Tr}F(\mathcal{A}_{ss})=\sum_{b'\in s}\mathfrak{e}(b')$ with $\mathfrak{e}(b')$ as in
\eqref{4.24:eq-bond-densities-a}. For each $b'\in s$ the resolvent part of $\mathfrak{e}(b')$ in
\eqref{4.24:eq-bond-densities-a}, namely its $\xi$-integral, is expanded termwise by
Lemma~\ref{4.24:prf-block-walk-expansion}; the subtraction of $(\lambda_0+\xi)^{-1}\mathbf{1}$ is carried by
the class $|\omega|=0$, which is why $\hat t^{\xi}_{\omega_0}$ is defined as in
\eqref{4.24:eq-hulls-anchors-a}, and the $|\omega|=0$ term is paired with this subtraction so that its
$\xi$-integral converges, as in the statement on the convergence of the $\xi$-integrals. Every walk
$\omega$ has exactly one hull $X(\omega)$, so summing \eqref{4.24:eq-hulls-anchors-a} over $X$ recovers the
whole resolvent part; the constant $-\tfrac12 n_{\mathfrak{g}}\log\lambda_0$ of $\mathfrak{e}(b')$ is a
single-cube member with hull $[\operatorname{cube}(b')]_{j+1}$. Hence $\sum_X$ collects
$\sum_{b'\in s}\mathfrak{e}(b')$ whole. The constant $-\tfrac12 n_{\mathfrak{g}}\log\lambda_0$ does not depend
on $(\mathbb{U},\mathbb{J})$, so it cancels in $\mathbb{C}^{s}_G-\mathbb{C}^{s}_G|_{(1,0)}$; its size is
contained in the single-cube budget $c_{\mathrm{sc}}$, which allows $O(1)\cdot(LM)^4$ per cube in the
field-dependence bound of the densities. This proves (b).

\emph{Anchors.} Every resolvent density in \eqref{4.24:eq-hulls-anchors-a} is indexed by a bond $b'\in s$,
and every single-cube member by a bond of $s$ or by a cell $c\subset R^{(n+1)}\cap Z^{s}$. The cell
$\operatorname{cube}(b')$ belongs to $\mathsf{Q}^{s}$, and the block pair of such a cell $c$
(Lemma~\ref{4.24:prf-bond-densities}) meets
$R_n\cap Z^{s}$; in the first case $\operatorname{cube}(b')$ meets $R_n\cap Z^{s}$ as well. Since
$X\supseteq\operatorname{cube}(b')$ (respectively contains the block pair of $c$), every $X$ meets
$R_n\cap Z^{s}$. The index set of the resolvent densities is $s$, which contains no bond of $\mathcal{W}$:
the variables $\phi_W$ on $\mathcal{W}$ are not integrated at this stage. Hence no $X$ is anchored on
$\bigcup\mathfrak{W}$, and (a) is proved.

\emph{Locality in the pair.} A hull may enter $Z_W^{+}$, since the walks of the class-wide structure are
not confined away from it; there the terms depend on $(\mathbb{U},\mathbb{J})$ as defined at this stage
(Lemma~\ref{4.24:prf-block-walk-expansion}). Each $t^{\xi}_{\omega}$ depends on $\mathbb{U}$ only through its
values on $D(\omega)\subset X$, the subtracted term $(\lambda_0+\xi)^{-1}\mathbf{1}_{\omega_0}$ is constant, and
each single-cube member depends on the pair only through its values on its block pair, which lies in $X$.
Hence $\mathbb{C}^{s}_G(X)$ depends on $(\mathbb{U},\mathbb{J})$ inside $X$ only, and on nothing else; this
is (c).

\emph{Holomorphy.} For each walk $\omega$, the integrand $\operatorname{tr}_{\mathfrak{g}}\hat
t^{\xi}_{\omega}(b',b')$ is holomorphic in the complex coordinates $z$ of the pair, jointly continuous in
$(\xi,z)$, and dominated by a $\xi$-integrable function locally uniformly in $z$, by the statement on the
convergence of the $\xi$-integrals; hence $\int_0^\infty d\xi\,\operatorname{tr}_{\mathfrak{g}}\hat
t^{\xi}_{\omega}(b',b')$ is holomorphic in $z$. The sum over $\{\omega:X(\omega)=X\}$ converges absolutely and
locally uniformly, being dominated by the majorants of the field-dependence bound of the densities, which are
free of $(\mathbb{U},\mathbb{J})$. By Weierstrass's theorem a locally uniformly convergent series of holomorphic
functions is holomorphic, so $\mathbb{C}^{s}_G(X)$ is holomorphic on $\mathfrak{D}^{\flat}_{n+1}(X)$, and a
fortiori on the subset $\mathfrak{D}^{\flat}_{n+1}[(\mathrm{amb})_W](X)$. No further smallness is needed for
the smaller domain; this includes the hulls $X$ that meet $Z_W^{+}$. This is (d).

\emph{Gauge invariance.} The diagonal block $\hat t^{\xi}_{\omega}(b',b')$ transforms covariantly under gauge
transformations, by conjugation with $\operatorname{Ad}$, and $\operatorname{tr}_{\mathfrak{g}}$ is invariant
under conjugation. This is (e).

\emph{Junctions.} We use clause (ii) of the lemma on junctions and filing units. The tree length
$\tau(X(\omega))$ is bounded by the length of the polygon of the walk, measured in the walk's own cube sides,
and each contact of consecutive domains contributes at most $\sqrt d\,Ms_m$. The cubes of a hull filed at
level $m$ lie in $\Omega^{c}_{m+1}$; therefore these cubes and the domains of the walk have levels at most $m$
and sides at most $Ms_m$, whatever component they cross. Here levels are those of the class-wide structure:
on an excluded $W$ the structure only freezes or coarsens (Notation~\ref{4.24:not-scales-rim-collar}), and at
this stage it is the current one; cubes inside $W$ that are finer only lengthen $d(\omega)$ per unit of
Euclidean length. With these facts the conversion of walk length into distance in filing units applies word
for word to the class-wide structure and yields \eqref{4.24:eq-hulls-anchors-1}. By clause (i) of the same
lemma, the rate in filing units, $m\ge j+1$, is at least $L\delta_0 M/(\sqrt d M_1)$.

It remains to check \eqref{4.24:eq-hulls-anchors-2}. The parameter conditions of
Definition~\ref{4.24:def-normalisation-constants}, with the constants $\beta$ and $\delta$ fixed there,
contain $(M/M_1)\min\{\delta_0,\delta\}\ge 8\kappa\sqrt d$.
Since $\delta_0\ge\min\{\delta_0,\delta\}$, this gives $\delta_0 M/(\sqrt d M_1)\ge 8\kappa$, and multiplying
by $L$ gives the first inequality. The same conditions contain
$\beta\kappa\ge c_d+2\sqrt d(2+4\beta)$ and $L\ge 8$; these give
$4\kappa L\ge c_d+1$. Then, for every $a_m\le 4\kappa L$,
\begin{equation*}
8\kappa L=4\kappa L+4\kappa L\ \ge\ a_m+(c_d+1).
\end{equation*}
This proves (f).

None of the clauses (i)--(iii) of the lemma on junctions and filing units, nor the filing level, nor the
arithmetic above, depends on the region, on the boundary of $\Lambda'_s$ near $\bigcup\mathfrak{W}$, or on
which anchors occur. In particular the arrested components produce no new kind of junction.
\end{proof}

\begin{lemma}[The normalisation constant is unchanged]\label{4.24:prf-constant-unchanged}
Fix the stage $n$ with its level $j$. Let $\pi_m$ be the cell partitions and let $\kappa$ be Ba{\l}aban's auxiliary parameter. For a region $X$ let $m(X)$ be the scale in which $X$ is filed and $d_m(X)$ its length in scale $m$, as in \eqref{4.24:eq-hulls-anchors-a}. For a rate family $a=(a_m)_m$ and a family $\mathbb{F}=(\mathbb{F}(X))_X$ of densities put
\begin{equation}\label{4.24:eq-constant-unchanged-1}
\|\mathbb{F}\|_{a}:=\sup_m\ \sup_{\Delta\in\pi_m}\ \sum_{X\ni\Delta,\ m(X)=m}
\|\mathbb{F}(X)\|_\infty\, e^{a_m d_m(X)} ,
\end{equation}
the supremum norm $\|\mathbb{F}(X)\|_\infty$ being taken over the analyticity domain $\mathfrak{D}^{\flat}_{n+1}(X)$.
This is the norm in which the normalisation lemma, whose objects and constants are fixed in Definition~\ref{4.24:def-normalisation-constants}, is stated.

For a $\pi_{j+1}$-cell $q$ let $\operatorname{Walk}(q)$ be the class of all walks starting in $q$. Here a walk is a sequence of admissible localisation domains in proper scales of the class-wide structure, any two consecutive ones intersecting. Let $\omega\mapsto X(\omega)$ be the hull map and $\bar m(\omega)$ the $\xi$-integrated majorant of the walk term of $\omega$, see \eqref{4.24:eq-hulls-anchors-a} and \eqref{4.24:eq-region-walk-expansion-b}. Put
\begin{equation}\label{4.24:eq-constant-unchanged-2}
\mathfrak{S}:=\sup_m\ \sup_{\Delta\in\pi_m}\ \sum_{q}\
\sum_{\substack{\omega\in\operatorname{Walk}(q):\\ X(\omega)\ni\Delta,\ m(X(\omega))=m}}
\bar m(\omega)\, e^{a_m d_m(X(\omega))} ,
\end{equation}
the sum over $q$ running over the $\pi_{j+1}$-cells. Let $c_{\mathrm{sc}}$ be the per-cube bound $O(1)\cdot(LM)^4e^{a_{j+1}}$ for the members whose region is a single cube. Both $\mathfrak{S}$ and $c_{\mathrm{sc}}$ depend on the rate family $a$; the majorant $\bar m(\omega)$ does not, so, since $d_m(X(\omega))\ge0$, $\mathfrak{S}$ is nondecreasing in every $a_m$, and $c_{\mathrm{sc}}$ is nondecreasing in $a_{j+1}$. By the normalisation lemma, its constant $c_G=c_G(d,L,M)$ is given by the universal majorant series $\mathfrak{S}$ as
\begin{equation}\label{4.24:eq-constant-unchanged-a}
c_G = c_G'\,\mathfrak{S} + c_{\mathrm{sc}},
\end{equation}
with $\mathfrak{S}$ and $c_{\mathrm{sc}}$ taken at the rates $a_m=4\kappa L$, where $c_G'$ is the per-cell count of Definition~\ref{4.24:def-normalisation-constants}. Moreover:
\begin{enumerate}
\item[(a)] For every rate family with $a_m\le 4\kappa L$ for all $m$ one has $\mathfrak{S}<\infty$ and $c_G'\,\mathfrak{S}+c_{\mathrm{sc}}\le c_G$, and the survivors' normalisation densities satisfy $\|\mathbb{C}^{s}_G\|_{a}\le c_G$, with the same constant $c_G(d,L,M)$.
\item[(b)] Assume the conversion to the stage norm, a statement whose proof is outstanding, which asserts $\|\mathbb{F}\|_{(n+1)}\le 3c_G$ for every family $\mathbb{F}$ of densities that has the hulls, the walk majorants and the anchoring of the family of the normalisation lemma and satisfies $\|\mathbb{F}\|_a\le c_G$ at every rate family with $a_m\le 4\kappa L$. Here $\|\cdot\|_{(n+1)}$ is the norm of stage $n+1$, in which the densities are weighted at the rates $a^{(n+1)}_m$ and summed over their filing levels. Under this assumption, $\|\mathbb{C}^{s}_G\|_{(n+1)}\le 3c_G$. In particular the constant $C_0^{\sharp}=3c_G+1$ serves for $\mathbb{C}^{s}_G$ unchanged.
\end{enumerate}
\end{lemma}

\begin{proof}
\emph{Finiteness of $\mathfrak{S}$.} Let $a$ be a rate family with $a_m\le 4\kappa L$ for all $m$. The proof of the normalisation lemma establishes $\mathfrak{S}<\infty$ at such rates by three controls.
\begin{itemize}
\item The length of a hull $X(\omega)$ is bounded by the lengths of the polygons of the walk $\omega$ plus the junctions between consecutive domains.
\item The walk majorants decay at a rate at least $8\kappa L$ per filing unit. Here $8\kappa L\ge a_m+c_d+1$, where $c_d$ is the constant entering the rate condition of the normalisation lemma. This rate pays for the decay in the position of the anchor cell $q$ relative to $\Delta$ and for the factor $e^{a_md_m(X(\omega))}$.
\item A lattice-animal count controls the combinatorics of the hulls.
\end{itemize}
Since $\mathfrak{S}$ is nondecreasing in every $a_m$, $c_{\mathrm{sc}}$ is nondecreasing in $a_{j+1}$, and $a_m\le 4\kappa L$, their values at $a$ are at most their values at the rates $a_m=4\kappa L$; by \eqref{4.24:eq-constant-unchanged-a}, $c_G'\,\mathfrak{S}+c_{\mathrm{sc}}\le c_G$ at the rate family $a$.

That proof bounds its family of walks only through the quantity $\mathfrak{S}$, taken over all of $\operatorname{Walk}(q)$. It never uses which walks actually occur. Since every term of \eqref{4.24:eq-constant-unchanged-2} is non-negative, the following holds. For every choice of sub-collections $\mathcal{V}(q)\subset\operatorname{Walk}(q)$, one for each $\pi_{j+1}$-cell $q$, and for every $m$ and every $\Delta\in\pi_m$,
\begin{equation}\label{4.24:eq-constant-unchanged-3}
\sum_q\ \sum_{\substack{\omega\in\mathcal{V}(q):\\ X(\omega)\ni\Delta,\ m(X(\omega))=m}}
\bar m(\omega)\,e^{a_md_m(X(\omega))}\ \le\ \mathfrak{S}.
\end{equation}

\emph{The survivors' family.} Fix $m$ and $\Delta\in\pi_m$. The members of $\mathbb{C}^{s}_G$ whose region is not a single cube are indexed as in \eqref{4.24:eq-hulls-anchors-a}. Each is the density of a walk $\omega$ of the survivors' walk expansion on $\Lambda'_s$. This walk starts from a bond $b'\in q\cap s$ of some $\pi_{j+1}$-cell $q$, the member's region is $X=X(\omega)$, and its sup norm is bounded by $\bar m(\omega)$, by \eqref{4.24:eq-region-walk-expansion-a} and \eqref{4.24:eq-region-walk-expansion-b}. Group the members according to the cell $q$ in which they are anchored. Bound the sum over the members anchored in $q$ by the number of densities of $\mathbb{C}^{s}_G$ anchored in $q$, times the supremum over the anchor bond. The remaining members are those whose region is a single cube. We obtain
\begin{equation}\label{4.24:eq-constant-unchanged-4}
\begin{split}
\sum_{X\ni\Delta,\ m(X)=m}\|\mathbb{C}^{s}_G(X)\|_\infty e^{a_md_m(X)}
&\le \sum_q \bigl[\#\text{densities of }\mathbb{C}^{s}_G\text{ anchored in }q\bigr]\\
&\qquad\cdot \sup_{b'\in q\cap s}
\sum_{\substack{\omega\text{ from }b':\\ X(\omega)\ni\Delta,\ m(X(\omega))=m}}
\bar m(\omega)\,e^{a_md_m(X(\omega))}\\
&\qquad + (\text{single-cube members}).
\end{split}
\end{equation}
Two facts bound the right-hand side.

$(\alpha)$ The walks from $b'$ of the survivors' expansion form a sub-collection of $\operatorname{Walk}(q)$, where $q$ is the $\pi_{j+1}$-cell containing $b'$. Their terms are dominated termwise by the same $\bar m(\omega)$. This holds because the constants in \eqref{4.24:eq-region-walk-expansion-a} and \eqref{4.24:eq-region-walk-expansion-b} do not depend on $\Lambda'$. Hence the supremum over $b'\in q\cap s$ is at most the sum over $\operatorname{Walk}(q)$. By \eqref{4.24:eq-constant-unchanged-3}, the resulting sum over $q$ is at most $\mathfrak{S}$, provided the bracket is bounded uniformly in $q$.

$(\beta)$ The bracket is bounded by $c_G'$, uniformly in $q$ and independently of the region. A cell $q$ contains $(LM)^d$ sites of level $j$, and each site is the initial point of $d$ bonds. So $q$ carries at most $d(LM)^d=4(LM)^4$ bond densities $\mathfrak{e}(b')$ with $b'\in q\cap s$ (here $d=4$). Next, by \eqref{4.24:eq-component-separation-c} of Lemma~\ref{4.24:lem-component-separation}, the cell $q$ meets at most one component. It therefore carries in addition at most $O(1)\cdot(LM)^4$ single-cube members of that one component. Together this is the count $(LM)^4\cdot O(1)$ entering the constant of the normalisation lemma. A cell of $\mathsf{Q}\setminus\mathsf{Q}^{s}$ carries no density of $\mathbb{C}^{s}_G$ at all. Thus the survivors' family has fewer anchors than the family of the normalisation lemma, never more, and the bracket is at most $c_G'$.

The single-cube members contribute at most $c_{\mathrm{sc}}$ per cube, by the definition of $c_{\mathrm{sc}}$. Inserting $(\alpha)$ and $(\beta)$ into \eqref{4.24:eq-constant-unchanged-4} and using the inequality $c_G'\,\mathfrak{S}+c_{\mathrm{sc}}\le c_G$ obtained above from \eqref{4.24:eq-constant-unchanged-a}, we get
\begin{equation*}
\sum_{X\ni\Delta,\ m(X)=m}\|\mathbb{C}^{s}_G(X)\|_\infty e^{a_md_m(X)}
\ \le\ c_G'\,\mathfrak{S}+c_{\mathrm{sc}}\ \le\ c_G .
\end{equation*}
The right-hand side does not depend on $m$ or $\Delta$. Taking the suprema in \eqref{4.24:eq-constant-unchanged-1} gives $\|\mathbb{C}^{s}_G\|_a\le c_G$, with the constant $c_G(d,L,M)$ of the normalisation lemma. Nothing above used the rates beyond $a_m\le 4\kappa L$, so the bound holds for every such rate family. This proves (a).

\emph{The stage norm.} Assume the conversion to the stage norm of (b). We check that $\mathbb{F}=\mathbb{C}^{s}_G$ has the hulls, the walk majorants and the anchoring of the family of the normalisation lemma, and the bounds at every admissible rate.
\begin{itemize}
\item Hulls: its regions are single cubes or hulls $X(\omega)$ of walks $\omega$ of the survivors' expansion, indexed as in \eqref{4.24:eq-hulls-anchors-a}, and these walks belong to $\operatorname{Walk}(q)$ by $(\alpha)$.
\item Walk majorants: termwise domination by $\bar m(\omega)$ is $(\alpha)$, so the majorant series is the same $\mathfrak{S}$; by $(\beta)$ the family has fewer anchors.
\item Anchoring: every region of $\mathbb{C}^{s}_G$ meets $R_n\cap Z^{s}\subset R_n$.
\item Rates: the bound $\|\mathbb{C}^{s}_G\|_a\le c_G$ at every rate family with $a_m\le 4\kappa L$ is (a).
\end{itemize}
Hence the conversion applies to $\mathbb{C}^{s}_G$ and gives $\|\mathbb{C}^{s}_G\|_{(n+1)}\le 3c_G$, so $C_0^{\sharp}=3c_G+1$ serves for $\mathbb{C}^{s}_G$ unchanged. This proves (b).
\end{proof}

The conversion to the stage norm reads only the hulls, the walk majorants, the anchoring and the bounds at every rate $a_m\le 4\kappa L$; it therefore applies in the same way to the joint family formed by $\mathbb{C}^{s}_G$ together with the densities $\mathbb{C}^{G}_W$ of the arrested components $W\in\mathfrak{W}$, which is introduced in the statement on the factors of the arrested components.

\begin{proposition}[Conversion of per-rate bounds to the stage norm. Proof outstanding]
\label{4.24:prop-norm-conversion}
Work at stage $n+1$ of the step, with the objects of Definition~\ref{4.24:def-normalisation-constants}.
For a family $\mathbb{F}=\{\mathbb{F}(X)\}$ of localised densities, each region $X$ is filed at a level
$m$, the index of the partition $\pi_m$ at which its weight in \eqref{4.24:eq-norm-conversion-1} is
read, as for the localised densities of Definition~\ref{4.24:def-normalisation-constants}; let
$d_m(X)$ be as in the glossary
(Notation~\ref{0.2:not-glossary}). For a sequence
of rates $\mathbf{a}=(a_m)_m$ put
\begin{equation}\label{4.24:eq-norm-conversion-1}
\|\mathbb{F}\|_{\mathbf{a}}
:=\sup_{m}\ \sup_{\Delta\in\pi_m}\
\sum_{X\ni\Delta,\ X\ \text{filed at level}\ m}\|\mathbb{F}(X)\|_{\infty}\,e^{a_m d_m(X)} .
\end{equation}
Let $\|\cdot\|_{(n+1)}$ be the stage norm of the localised densities at stage $n+1$, with rates
$a^{(n+1)}_m$. Let $c_G$ be the constant of \eqref{4.24:eq-constant-unchanged-a}, and $\mathfrak{S}$ the
universal majorant series defined in that display. Suppose that $\mathbb{F}$ has the following three
properties:
\begin{itemize}
\item[(a)] $\|\mathbb{F}\|_{\mathbf{a}}\le c_G$ for every sequence $\mathbf{a}$ with
$a_m\le 4\kappa L$ for all $m$;
\item[(b)] apart from single-cube members, whose contribution is bounded per cube by the budget
$c_{\mathrm{sc}}$ entering \eqref{4.24:eq-constant-unchanged-a}, each member of $\mathbb{F}$ anchored
at a bond $b'$ is indexed by a walk $\omega$ of the class $\operatorname{Walk}(q)$ of the cell
$q\in\mathsf{Q}$ containing $b'$ (the $\pi_{j+1}$-cells over which $\mathfrak{S}$ is summed in
\eqref{4.24:eq-constant-unchanged-a}), has region $X(\omega)$, and is dominated termwise by the walk
majorant $\bar m(\omega)$, so that the series of these majorants is dominated by the universal
series $\mathfrak{S}$;
\item[(c)] every region $X$ of the family meets $R_n$.
\end{itemize}
Then
\begin{equation}\label{4.24:eq-norm-conversion-2}
\|\mathbb{F}\|_{(n+1)}\le 3c_G .
\end{equation}
In particular \eqref{4.24:eq-norm-conversion-2} holds for the survivors' normalisation family
$\mathbb{C}^{s}_G$: every region of it meets $R_n\cap Z^{s}$, its majorants are those of
$\mathbb{C}_G$, with fewer anchors, and hypothesis (a) for it is the per-rate bound of the statement
that the normalisation constant is unchanged, whose constant $c_G$ is that of
\eqref{4.24:eq-constant-unchanged-a}. It also holds for the joint family formed by $\mathbb{C}^{s}_G$ and
the Jacobian densities $\mathbb{C}^{G}_W$, $W\in\mathfrak{W}$. For these families, therefore, the
constant $C_0^{\sharp}=3c_G+1$ is the same as for $\mathbb{C}_G$.
\end{proposition}

Proof outstanding.

\emph{Route.} A proof would have to carry out the passage from the bounds (a) to the stage norm
$\|\cdot\|_{(n+1)}$ using only the data (a)--(c). For $\mathbb{F}=\mathbb{C}_G$ the bound
\eqref{4.24:eq-norm-conversion-2} is the bound $\|\mathbb{C}_G\|_{(n+1)}\le 3c_G$ stated for the
normalisation family itself; the mechanism below is the clause treating $\mathbb{C}_G$ in the proof of
the theorem bounding the localised densities of the stage in the stage norm. The mechanism is the
following. The bounds (a) hold
at every rate $a_m\le 4\kappa L$, not only at the rates $a^{(n+1)}_m$ of the stage norm. The rate surplus
$4\kappa L-a^{(n+1)}_m$ of the walk majorants $\bar m(\omega)$ then pays the sum over the filing levels
$m$ that the stage norm collects. What is outstanding is this summation over filing levels, carried
to the constant $3$ in \eqref{4.24:eq-norm-conversion-2}.

\begin{lemma}[Field dependence of the densities in cube gauges. Stated; proof incomplete at the step marked $(\ast)$]\label{4.24:prf-field-dependence}
Use the fibres, pairings and gauge action of Notation~\ref{4.24:not-fibres-pairings}, the termwise walk
expansion of the block resolvent (Lemma~\ref{4.24:prf-block-walk-expansion}) and the densities
$\hat t^{\xi}_{\omega}$, $\mathbb{C}^{s}_{G}$ of \eqref{4.24:eq-hulls-anchors-a}. Write
$\mathbb{C}^{s}_{G}{\restriction}_{(1,0)}$ for the same family at the trivial pair $(\mathbb{U},\mathbb{J})=(\mathbf{1},0)$,
and $\mathcal{A}_{ss}{\restriction}_{(1,0)}$ for the block at that pair.

Let $X$ be the region on which the clauses of \cite[(1.64)--(1.69)]{B-CMP122a} are imposed, $\eta$ the
lattice unit of \cite{B-CMP122a}, and, for a bond $b$, $s(b)$ the side unit of the shell containing $b$,
that is, $L^{i}\eta$ when that shell has index $i$. Let $B$ and $C$ be the constants of those clauses, the
cube of a clause in the shell of index $i$ having side
$CML^{i}\eta$, and $C\le 2$. Write $U$ for the $G$-valued factor of $\mathbb{U}$, so that
$\mathbb{U}=e^{i\eta A'}U$, where $A'$ is the bond field of the complex factor of the pair in
\cite{B-CMP122a}, and $Z(G)$ for the centre of $G$ ($Z(G)=\{\pm1\}$ for $G=SU(2)$). Put
\begin{equation}\label{4.24:eq-field-dependence-b}
\begin{gathered}
\theta:=dCM\,\varepsilon_{\square},\qquad
\varepsilon_{\square}:=c_1\,\mathfrak{br}\,(BCM\tilde\alpha_0+\tilde\alpha_1),\\
\|R_a(\mathbb{U}^{u_a})-R_a(\mathbf{1})\|\le\epsilon\,r_a,\\
\epsilon=O(1)(BCM\tilde\alpha_0+\tilde\alpha_1)\quad\text{or}\\
\epsilon=O(1)\bigl[5w_*(j)(BCM\tilde\alpha_0+\tilde\alpha_1)+K_\Delta\gamma_0^{(k)}\bigr],
\end{gathered}
\end{equation}
where $c_1$ is absolute and $\mathfrak{br}$ is the largest of the brackets and powers of the auxiliary length
parameter of \cite[(1.65)--(1.69)]{B-CMP122a} multiplying $BCM\tilde\alpha_0$ and $\tilde\alpha_1$ there. The
second line of \eqref{4.24:eq-field-dependence-b} is the smallness of the factor $R_a$ of a walk term,
in the gauge $u_a$ of its clause cube, against its per-factor majorant $r_a$ (these objects are
introduced in the proof); $\epsilon$ is its constant. In the field-dependence bound for the normalisation
densities at the full stage, $\epsilon=O(1)(BCM\tilde\alpha_0+\tilde\alpha_1)$; in the form in which the
parameter conditions of Definition~\ref{4.24:def-normalisation-constants} carry it, $\epsilon$ is the last
line of \eqref{4.24:eq-field-dependence-b}.
Part (b) holds with either form; part (c) is written with the second.
\begin{enumerate}
\item[(a)] Every density $\operatorname{tr}_{\mathfrak{g}}\hat t^{\xi}_{\omega}(b',b')$ is a gauge-invariant
function of $(\mathbb{U},\mathbb{J})$.
\item[(b)] For every $b'\in s$, every $\xi\ge0$ and every walk $\omega=(X_0,\dots,X_p)$ of the expansion
with $X_0\ni b'$,
\begin{equation}\label{4.24:eq-field-dependence-a-bis}
\begin{split}
&\bigl|\operatorname{tr}_{\mathfrak{g}}t^{\xi}_{\omega}(\mathbb{U})(b',b')
-\operatorname{tr}_{\mathfrak{g}}t^{\xi}_{\omega}(\mathbf{1})(b',b')\bigr|\\
&\qquad\le n_{\mathfrak{g}}\Bigl[(p+1)\epsilon+(6p+2)\theta
+2\min_{z\in Z(G)}\bigl|U(\ell_\omega)-z\bigr|\Bigr]\,m(\omega),
\end{split}
\end{equation}
where $\ell_\omega\subset D(\omega)$ is the closed lattice loop through the chain
$X_0,\dots,X_p,X_0$ constructed in Step~7 of the proof and $U(\ell_\omega)$ is the transport of $U$
around it.
\item[(c)] Assume the holonomy bound for walks across several cubes,
Proposition~\ref{4.24:prop-holonomy-bound}, whose proof is outstanding, with the constant $\mathfrak{h}$ of
that proposition. Then
\begin{equation}\label{4.24:eq-field-dependence-1-bis}
\bigl\|\mathbb{C}^{s}_{G}-\mathbb{C}^{s}_{G}{\restriction}_{(1,0)}\bigr\|
\le c_G'\,O(1)\Bigl[dCM\bigl(5w_*(j)(BCM\tilde\alpha_0+\tilde\alpha_1)\bigr)
+K_\Delta\gamma_0^{(k)}+\mathfrak{h}\Bigr]\le\vartheta,
\end{equation}
the $O(1)$ carrying $c_1\mathfrak{br}$; the last inequality holds for all $g_k\le\gamma$ whenever $\gamma$ is
small enough that
\begin{equation}\label{4.24:eq-field-dependence-9}
c_G'\,O(1)\Bigl[dCM\bigl(5w_*(j)(BCM\tilde\alpha_0+\tilde\alpha_1)\bigr)+K_\Delta\gamma_0^{(k)}\Bigr]
\le\tfrac34\,\vartheta\qquad\text{for all }g_k\le\gamma,
\end{equation}
a condition that bounds the coupling from above only.
\end{enumerate}
\end{lemma}

\begin{proof}
\emph{Part (a).} By the gauge covariance of the stage form, the block $\mathcal{A}_{ss}$ is gauge covariant,
$\mathcal{A}_{ss}[\mathbb{U}^{u},\mathbb{J}^{u}]=\mathcal{U}_u\mathcal{A}_{ss}[\mathbb{U},\mathbb{J}]\mathcal{U}_u^{-1}$.
For a unitary $\mathcal{U}$, any $T\in\mathfrak{P}$ and $\xi\ge0$ one has
$(\mathcal{U}T\mathcal{U}^{-1}+\xi)^{-1}=\mathcal{U}(T+\xi)^{-1}\mathcal{U}^{-1}$, and conjugation by a unitary
preserves the class $\mathfrak{P}$ of operators with positive Hermitian part; hence
$F(\mathcal{U}\mathcal{A}_{ss}\mathcal{U}^{-1})=\mathcal{U}F(\mathcal{A}_{ss})\mathcal{U}^{-1}$. Each walk term
$t^{\xi}_{\omega}$ is a product of covariant factors (see below), so
$t^{\xi}_{\omega}(\mathbb{U}^{u})=\mathcal{U}_u t^{\xi}_{\omega}(\mathbb{U})\mathcal{U}_u^{-1}$. Since $\mathcal{U}_u$ is
bond-diagonal with entry $\operatorname{Ad}_{u(b'_-)}$ at $b'$, the diagonal entry transforms as
\begin{equation*}
t^{\xi}_{\omega}(\mathbb{U}^{u})(b',b')
=\operatorname{Ad}_{u(b'_-)}\,t^{\xi}_{\omega}(\mathbb{U})(b',b')\,\operatorname{Ad}_{u(b'_-)}^{-1},
\end{equation*}
and its trace $\operatorname{tr}_{\mathfrak{g}}$ is unchanged. The subtraction that turns $t^{\xi}_{\omega}$ into
$\hat t^{\xi}_{\omega}$ in \eqref{4.24:eq-hulls-anchors-a} is a field-free multiple of the identity, so
$\operatorname{tr}_{\mathfrak{g}}\hat t^{\xi}_{\omega}(b',b')$ is gauge invariant.

\emph{Elementary inequalities.} Norms are those of Notation~\ref{4.24:not-fibres-pairings}: $|\cdot|$ on
the fibre and on $2\times2$ matrices, $\|\cdot\|$ on $\mathcal{H}$. For a $G$-valued gauge transformation
$u$ on a set of sites, $\mathcal{U}_u:=\bigoplus_b\operatorname{Ad}_{u(b_-)}$ is unitary and bond-diagonal,
and $P_E\mathcal{U}_u=\mathcal{U}_uP_E$; for a constant $g\in G$ we write
$\mathcal{U}_g:=\bigoplus_b\operatorname{Ad}_g$. We use the following elementary inequalities. For a unitary $V$ on a
finite-dimensional Hilbert space and any operator $T$,
\begin{equation}\label{4.24:eq-field-dependence-2-bis}
\|VTV^{-1}-T\|\le2\|V-1\|\,\|T\|,
\end{equation}
since $VTV^{-1}-T=(V-1)TV^{-1}+T(V^{-1}-1)$ and $\|V^{-1}-1\|=\|V-1\|$. For $h\in G=SU(2)\subset U(2)$,
$|\operatorname{Ad}_h-1|\le2|h-1|$: indeed $\operatorname{Ad}_hY-Y=(h-1)Yh^{-1}+Y(h^{-1}-1)$, and $h-1$,
$h^{-1}-1$ are scalar multiples of unitaries with the scalar of modulus $|h-1|$, because
$(h-1)^{*}(h-1)=2-h-h^{*}$ is a multiple of the identity for $h\in SU(2)$. For unitaries $h_1,\dots,h_n$,
\begin{equation}\label{4.24:eq-field-dependence-3-bis}
|h_1\cdots h_n-1|\le\sum_{i=1}^{n}|h_i-1|,
\end{equation}
by telescoping, each partial product being unitary. In the same way, for unitaries $V_i,W_i$,
\begin{equation*}
\prod_{i}W_iV_i-\prod_iV_i
=\sum_i\Bigl(\prod_{l<i}W_lV_l\Bigr)(W_i-1)V_i\Bigl(\prod_{l>i}V_l\Bigr),
\end{equation*}
so that $|\prod_iW_iV_i-\prod_iV_i|\le\sum_i|W_i-1|$.

\emph{The walk terms.} By Lemma~\ref{4.24:prf-block-walk-expansion} and
\cite[(3.107)]{B-CMP99b}, for $\xi\ge0$ and $b'\in s$,
$(\mathcal{A}_{ss}+\xi)^{-1}(b',b')=\sum_\omega t^{\xi}_{\omega}(b',b')$, the sum over walks
$\omega=(X_0,\dots,X_p)$ with $X_0\ni b'$ and $X_{a-1}\cap X_a\ne\emptyset$, and
$t^{\xi}_{\omega}=R_0R_1\cdots R_p$, where $R_a=R_a(X_a;\mathbb{U})$ (its dependence on $\xi$ suppressed) is an
operator on $\mathcal{H}$ whose kernel is supported on bonds of $X_a^{\sim}\times X_a^{\sim}$ and which
depends on $\mathbb{U}$ restricted to $X_a^{\sim}$ only, as in \cite[(3.107)]{B-CMP99b}. Each $R_a$ is
Ad-covariant, $R_a(\mathbb{U}^{u})=\mathcal{U}_uR_a(\mathbb{U})\mathcal{U}_u^{-1}$, as are the operators of
\cite[(3.34)]{B-CMP99b}. By \cite[Theorem~3.10, (3.108)]{B-CMP99b},
\begin{equation*}
|t^{\xi}_{\omega}|\le m(\omega)=O(1)\bigl(O(1)M^{-1/2}\bigr)^{p}e^{-\delta_0d(\omega)}.
\end{equation*}

\emph{Cube gauges.} By \cite[(1.64)--(1.69)]{B-CMP122a}, for every cube $\square$ of the family of the
clauses there is a gauge transformation $u$ on $\square\cap X$ with $U^{u}=\exp i\eta A$ for a
$\mathfrak{g}$-valued bond field $A$, and by
\cite{B-CMP99b} each domain $X_a^{\sim}$ is contained in one of the cubes on which these conditions
hold. Thus for each $a$ there is a clause cube $\square_a$ with $X_a^{\sim}\subset\square_a$ and a
$G$-valued $u_a$ on $\square_a\cap X$. We extend $u_a$ to all sites by $u_a:=1$ off $\square_a\cap X$;
since the kernel of $R_a$ lives on $X_a^{\sim}\times X_a^{\sim}\subset(\square_a\cap X)^2$, the
covariance identity and the bound $(\ast)$(ii) below read $u_a$ on $\square_a\cap X$ only.

\emph{Per-bond smallness in the cube gauge.} In the gauge $u_a$, on $\square_a\cap X$,
$\mathbb{U}^{u_a}(b)=(U')^{u_a}(b)\,U^{u_a}(b)$ with $U^{u_a}(b)=\exp i\eta A(b)$, where
$L^{i}\eta|A|$ is less than a bracket times $BCM\tilde\alpha_0$ by \cite[(1.65), (1.67), (1.69)]{B-CMP122a},
and $U'=\exp i\eta A'$, with $|A'|$ at most a bracket times $\tilde\alpha_1/(L^{i}\eta)$ by the
$|A'|$-members of \cite[(1.64), (1.66)]{B-CMP122a}. With $|e^{iY}-1|\le|Y|e^{|Y|}$ this gives, for every
bond $b\subset\square_a\cap X$,
\begin{equation}\label{4.24:eq-field-dependence-4}
|\mathbb{U}^{u_a}(b)-1|\le\varepsilon_{\square}\,\eta/s(b),\qquad
|U^{u_a}(b)-1|\le\varepsilon_{\square}\,\eta/s(b).
\end{equation}
The bracket of \cite[(1.66)]{B-CMP122a} is at most $1$; the powers of the auxiliary length parameter of
\cite[(1.65)--(1.69)]{B-CMP122a} entering $\mathfrak{br}$ are at most its $2N_0$-th power, $N_0$ the integer
of \cite[p.~192]{B-CMP122a}, and by \cite[p.~192]{B-CMP122a} this is less than $L(L+1)N_0^{-1}R_k$, which
is bounded by a constant times a power of $\log g_k^{-2}$; so $\varepsilon_{\square}$ is at most $g_k$ times
a constant times a power of $\log g_k^{-2}$, and hence
polynomially small in $g_k$. The set $\square_a\cap X$
has $\ell^1$-diameter at most $dCML^{i}$ bonds, $L^{i}\eta$ being the unit with which the same clause
measures $A$ (cube side $CML^{i}\eta$
paired with $L^{i}\eta|A|$ in \cite[(1.65)]{B-CMP122a}, $CML^{i}\eta$ with $L^{i}\eta|A|$ in
\cite[(1.67)]{B-CMP122a}, and likewise in \cite[(1.68)--(1.69)]{B-CMP122a}). For a lattice path
$\gamma\subset\square_a\cap X$ of $n$ bonds, \eqref{4.24:eq-field-dependence-3-bis} and
\eqref{4.24:eq-field-dependence-4} give $|U^{u_a}(\gamma)-1|\le n\varepsilon_{\square}\eta/s_{\min}(\gamma)$,
where $s_{\min}(\gamma):=\min_{b\in\gamma}s(b)$; for a path inside one cube crossing it once,
$n\eta/s_{\min}(\gamma)\le dCM$, so
\begin{equation}\label{4.24:eq-field-dependence-5}
|U^{u_a}(\gamma)-1|\le dCM\,\varepsilon_{\square}=\theta
\end{equation}
for every path $\gamma$ inside $\square_a\cap X$ of $\ell^1$-length at most the diameter of the cube.
The quantity $\theta$ is the same on every shell and carries one factor $M$ more than the term
$BCM\tilde\alpha_0$ of $\epsilon$: it is the integral of a field of size $BCM\tilde\alpha_0/s(b)$ over a
length $Ms(b)$. The factor is present: a constant $A$ with $|A|=BCM\tilde\alpha_0/s(b)$ is admitted by the
clause and rotates by $M\cdot BCM\tilde\alpha_0$ across the cube.

\emph{The step marked $(\ast)$.} The proof is incomplete at the following three statements. The first two
are used as stated, with the references given beside them. The third, the short-path connectedness of the
overlaps and of the walk domains, is a geometric statement about the region $X$, the clause cubes and the
domains of the walk expansion; it is proved below only in the cases noted after it, and in general it is
the step at which the proof is incomplete.
\begin{enumerate}
\item[(i)] \emph{Per-factor majorants.} There are majorants $r_a$ of the factors, $\|R_a\|\le r_a$ at every
pair entering below (the trivial pair included), with $\prod_{a=0}^{p}r_a=m(\omega)$. This is the form in
which the proof of \cite[Theorem~3.10]{B-CMP99b} builds the bound \cite[(3.108)]{B-CMP99b}: the
operators constructed for a sequence of cubes satisfy the inequalities of
\cite[Theorems~3.1--3.3]{B-CMP99b}, and the estimates of the expansion are products of these; it is used
here as stated in the proof of \cite[Theorem~3.10]{B-CMP99b}.
Since $\mathcal{U}_{u_a}$ is unitary, $\|R_a(\mathbb{U}^{u_a})\|=\|R_a(\mathbb{U})\|\le r_a$.
\item[(ii)] \emph{Per-factor smallness in the cube gauge.} For each $a$ the second line of
\eqref{4.24:eq-field-dependence-b}, $\|R_a(\mathbb{U}^{u_a})-R_a(\mathbf{1})\|\le\epsilon\,r_a$, holds.
This is the cube-wise bound
\begin{equation*}
\|\Delta^{(j)}(\mathbb{U})-\Delta^{(j)}(\mathbf{1})\|\le O(1)(BCM\tilde\alpha_0+\tilde\alpha_1)
\end{equation*}
of the field-dependence bound for the normalisation densities at the full stage, $\Delta^{(j)}$ being the
operator of the stage from which, with the local inverses of the cube, the walk factor is built; it is
carried to that walk factor in the way in which the remainders of
\cite{B-CMP99b} satisfy the bounds of \cite[Theorem~3.1]{B-CMP99b} with an additional small factor
$O(1)\alpha_1$, and it is used here as stated in the full-stage bound.
\item[(iii)] \emph{Short-path connectedness of the overlaps and of the walk domains.} For each
$a=1,\dots,p$ and each site $x$ of $O_a$ (defined in Step~1), the site $x_a$ chosen in Step~3 is joined
to $x$ by a lattice path of bonds of $(\square_{a-1}\cap X)\cap(\square_a\cap X)$ of $\ell^1$-length at
most the smaller of the $\ell^1$-diameters of the two clause cubes $\square_{a-1}$, $\square_a$ (note
$O_a\subset X_{a-1}^{\sim}\cap X_a^{\sim}\subset(\square_{a-1}\cap X)\cap(\square_a\cap X)$, so paths inside
$O_a$ qualify); and for each $a=1,\dots,p$ the sites $x_a,x_{a+1}$ chosen in Step~3 are joined inside
$X_a^{\sim}$ by a lattice path of $\ell^1$-length at most the $\ell^1$-diameter of $\square_a$.
\end{enumerate}
The corresponding statement for the first domain holds. A lattice box is a product of $d$ lattice
intervals; two of its sites $x,y$ are joined inside it, changing the coordinates of $x$ into those of $y$
one axis at a time and one bond at a time, by a lattice path of $\ell^1$-length $|x-y|_1$, which is at most
the $\ell^1$-diameter of the box; a box contained in another has no larger $\ell^1$-diameter, and the
intersection of two lattice boxes is a lattice box (or empty) contained in both. The first domain $X_0$ of
a walk is a single cube of the cover of \cite{B-CMP99b}, hence a lattice box, and
$X_0\subset X_0^{\sim}\subset\square_0$; so the sites $x_{p+1},x_1$ of $X_0$ chosen in Step~3 are joined
inside $X_0$ by a lattice path of $\ell^1$-length at most the $\ell^1$-diameter of $\square_0$. By the same
argument the second clause of (iii) holds whenever $X_a^{\sim}$ is a lattice box, and the first whenever
$X_{a-1}^{\sim}$ and $X_a^{\sim}$ are lattice boxes, $O_a$ being then a lattice box contained in both.

\emph{Step 1 (insertion of the cube gauges).} Fix $b'\in s$, $\xi\ge0$ and a walk
$\omega=(X_0,\dots,X_p)$ with $X_0\ni b'$; because the entry is diagonal, the right index of $R_p$ is
$b'$, so $X_p^{\sim}\ni b'$ as well. Put $O_a:=X_{a-1}^{\sim}\cap X_a^{\sim}\ne\emptyset$ for $a=1,\dots,p$
and $O_{p+1}:=X_p^{\sim}\cap X_0^{\sim}\ni b'$. Abbreviate $R_a:=R_a(\mathbb{U})$,
$\tilde R_a:=R_a(\mathbb{U}^{u_a})$, $R_a^{0}:=R_a(\mathbf{1})$, $\mathcal{U}_a:=\mathcal{U}_{u_a}$. By covariance
$R_a=\mathcal{U}_a^{-1}\tilde R_a\mathcal{U}_a$, hence
\begin{equation*}
t^{\xi}_{\omega}=R_0R_1\cdots R_p
=\mathcal{U}_0^{-1}\,\tilde R_0\,T_1\,\tilde R_1\,T_2\cdots T_p\,\tilde R_p\,\mathcal{U}_p,
\end{equation*}
with $T_a:=\mathcal{U}_{a-1}\mathcal{U}_a^{-1}=\mathcal{U}_{h_a}$ and $h_a:=u_{a-1}u_a^{-1}$.

\emph{Step 2 (only the overlaps are read).} In $\tilde R_{a-1}T_a\tilde R_a$ the middle index runs over
bonds in the support of the right index of $\tilde R_{a-1}$ intersected with the support of the left
index of $\tilde R_a$, a subset of $O_a$. Since $T_a$ is bond-diagonal, it may be replaced by
$P_{O_a}\mathcal{U}_{h_a}P_{O_a}$: only the values of $h_a$ at the sites of $O_a$ enter. Likewise
$\mathcal{U}_0^{-1}$ and $\mathcal{U}_p$ enter the diagonal entry only at $b'$:
\begin{equation*}
t^{\xi}_{\omega}(b',b')=\operatorname{Ad}_{u_0(b'_-)}^{-1}
\bigl[\tilde R_0T_1\tilde R_1\cdots T_p\tilde R_p\bigr](b',b')\,\operatorname{Ad}_{u_p(b'_-)}.
\end{equation*}
By cyclicity of $\operatorname{tr}_{\mathfrak{g}}$ and
$\operatorname{Ad}_{u_p(b'_-)}\operatorname{Ad}_{u_0(b'_-)}^{-1}=\operatorname{Ad}_{h_{p+1}(b'_-)}$,
\begin{equation}\label{4.24:eq-field-dependence-7}
\begin{split}
\operatorname{tr}_{\mathfrak{g}}t^{\xi}_{\omega}(b',b')
&=\operatorname{tr}_{\mathfrak{g}}\Bigl(\bigl[\tilde R_0T_1\tilde R_1\cdots T_p\tilde R_p\bigr](b',b')\,
\operatorname{Ad}_{h_{p+1}(b'_-)}\Bigr),\\
h_{p+1}&:=u_pu_0^{-1}\quad\text{on }O_{p+1}.
\end{split}
\end{equation}

\emph{Step 3 (the transition functions are nearly constant on overlaps).} Here and in Step~7 the
$G$-valued factor $U$ is used in place of $\mathbb{U}$; the complex factor $e^{i\eta A'}$ is carried inside
each factor $\tilde R_a$. Gauge transformations act on $U$ by
$U^{u}(b)=u(x)U(b)u(y)^{-1}$ for $b=\langle x,y\rangle$, and $u_{a-1}=h_au_a$ at every site. Hence for
sites $x,y$ joined by a bond $b=\langle x,y\rangle$ of $(\square_{a-1}\cap X)\cap(\square_a\cap X)$,
$U^{u_{a-1}}(b)=h_a(x)U^{u_a}(b)h_a(y)^{-1}$, and therefore
\begin{equation*}
h_a(x)h_a(y)^{-1}=U^{u_{a-1}}(b)\,\bigl(h_a(y)U^{u_a}(b)h_a(y)^{-1}\bigr)^{-1}.
\end{equation*}
By \eqref{4.24:eq-field-dependence-3-bis}, conjugation invariance of $|\cdot|$ and
\eqref{4.24:eq-field-dependence-4}, applicable because $b$ lies in $(\square_{a-1}\cap X)\cap(\square_a\cap X)$,
\begin{equation*}
|h_a(x)h_a(y)^{-1}-1|\le|U^{u_{a-1}}(b)-1|+|U^{u_a}(b)-1|\le2\varepsilon_{\square}\eta/s(b).
\end{equation*}
Pick a site $x_1$ of $X_0\cap X_1\subset O_1$, which is non-empty since consecutive domains of a walk
meet, a site $x_a$ of $O_a$ for $2\le a\le p$, and $x_{p+1}:=b'_-$, a site of $X_0$; put
$g_a:=h_a(x_a)\in G$ and $k_a(x):=g_a^{-1}h_a(x)$. For $a=p+1$ only the identity
$g_{p+1}=h_{p+1}(x_{p+1})$ is used below.
For $a=1,\dots,p$ and a site $x$ of $O_a$ take, by the short-path connectedness statement $(\ast)$(iii)
of the overlaps, a path
$x_a=y_0,y_1,\dots,y_n=x$ of bonds of $(\square_{a-1}\cap X)\cap(\square_a\cap X)$ of $\ell^1$-length at most
the smaller of the diameters of $\square_{a-1}$ and $\square_a$. With $i_{\min}$ the least shell index of its
bonds, such a path has at most $dCML^{i_{\min}}$ bonds, and by \eqref{4.24:eq-field-dependence-4} each of
its bonds $b$ costs at most $\varepsilon_{\square}\eta/s(b)\le\varepsilon_{\square}L^{-i_{\min}}$ in either
of the two gauges $u_{a-1}$, $u_a$: its bonds lie in both clause domains, where both bounds hold, whereas
$u_{a-1}$ is extended by $1$ off $\square_{a-1}\cap X$, where no smallness is available. Hence in each of the
two gauges the sum of the costs along the path is at most $dCM\varepsilon_{\square}=\theta$. Since
$k_a(x)$ is conjugate to $h_a(x)h_a(x_a)^{-1}=\prod_{i}h_a(y_{i+1})h_a(y_i)^{-1}$,
\eqref{4.24:eq-field-dependence-3-bis}, the last display and these two sums give
\begin{equation*}
\sup_{x\in O_a}|k_a(x)-1|\le2\theta,\qquad
\bigl\|P_{O_a}(\mathcal{U}_{h_a}-\mathcal{U}_{g_a})P_{O_a}\bigr\|
=\sup_{x}\bigl|\operatorname{Ad}_{g_a}(\operatorname{Ad}_{k_a(x)}-1)\bigr|\le4\theta,
\end{equation*}
the supremum over the initial sites $x$ of bonds of $O_a$, using $|\operatorname{Ad}_h-1|\le2|h-1|$.

\emph{Step 4 (transitions replaced by constants).} Let $\Pi$ be the product
$\tilde R_0T_1\tilde R_1\cdots T_p\tilde R_p$ followed at the entry by
$T'_{p+1}:=\operatorname{Ad}_{h_{p+1}(b'_-)}$, so that $\operatorname{tr}_{\mathfrak{g}}\Pi(b',b')$ is the
right side of \eqref{4.24:eq-field-dependence-7}, and let
\begin{equation*}
\Pi^{c}:=\tilde R_0\,\mathcal{U}_{g_1}\tilde R_1\,\mathcal{U}_{g_2}\cdots\mathcal{U}_{g_p}\tilde R_p\cdot
\operatorname{Ad}_{g_{p+1}}.
\end{equation*}
Since $h_{p+1}(x_{p+1})=g_{p+1}$, the last replacement is exact. Replace
$T_a$ by $\mathcal{U}_{g_a}$ one at a time, $a=1,\dots,p$. By Step~2, applied equally to the bond-diagonal
$\mathcal{U}_{g_a}$, each difference is $P_{O_a}(\mathcal{U}_{h_a}-\mathcal{U}_{g_a})P_{O_a}$, of norm at most
$4\theta$ by Step~3, sandwiched between the untouched factors $\tilde R_{a'}$, $a'\ne a$, of norm at most
$r_{a'}$ by $(\ast)$(i), and unitaries of norm $1$. With $|\operatorname{tr}_{\mathfrak{g}}A|\le n_{\mathfrak{g}}|A|$
and $|K(b',b')|\le\|K\|$,
\begin{equation*}
\bigl|\operatorname{tr}_{\mathfrak{g}}\Pi(b',b')-\operatorname{tr}_{\mathfrak{g}}\Pi^{c}(b',b')\bigr|
\le n_{\mathfrak{g}}\cdot p\cdot4\theta\cdot\prod_{a=0}^{p}r_a .
\end{equation*}

\emph{Step 5 (each factor to the trivial pair in its own gauge).} In $\Pi^{c}$ replace $\tilde R_a$ by
$R_a^{0}$ one at a time, $a=0,\dots,p$. By the second line of \eqref{4.24:eq-field-dependence-b} each
replacement costs $\epsilon r_a$ against factors of norm at most $r_{a'}$ or $1$; hence
\begin{gather*}
\bigl|\operatorname{tr}_{\mathfrak{g}}\Pi^{c}(b',b')-\operatorname{tr}_{\mathfrak{g}}\Pi^{0}(b',b')\bigr|
\le n_{\mathfrak{g}}\cdot(p+1)\cdot\epsilon\cdot\prod_{a=0}^{p}r_a,\\
\Pi^{0}:=R_0^{0}\,\mathcal{U}_{g_1}R_1^{0}\,\mathcal{U}_{g_2}\cdots\mathcal{U}_{g_p}R_p^{0}\cdot
\operatorname{Ad}_{g_{p+1}}.
\end{gather*}

\emph{Step 6 (constants commute with the trivial-pair factors).} For a constant $g\in G$ one has
$\mathbf{1}^{g}=\mathbf{1}$, since $g\,1\,g^{-1}=1$, so covariance gives
$R_a^{0}=R_a(\mathbf{1}^{g})=\mathcal{U}_gR_a^{0}\mathcal{U}_g^{-1}$: every $R_a^{0}$ commutes with every
$\mathcal{U}_g$. Moving the $\mathcal{U}_{g_a}$ to the right and using $\mathcal{U}_g\mathcal{U}_{g'}=\mathcal{U}_{gg'}$,
$\Pi^{0}=R_0^{0}R_1^{0}\cdots R_p^{0}\,\mathcal{U}_{H_\omega}$ at the entry, that is,
\begin{gather*}
\operatorname{tr}_{\mathfrak{g}}\Pi^{0}(b',b')
=\operatorname{tr}_{\mathfrak{g}}\bigl(t^{\xi}_{\omega}(\mathbf{1})(b',b')\,\operatorname{Ad}_{H_\omega}\bigr),\\
H_\omega:=g_1g_2\cdots g_{p+1}=\prod_{a=1}^{p+1}u_{a-1}(x_a)u_a(x_a)^{-1},
\end{gather*}with $u_{p+1}:=u_0$.

\emph{Step 7 (the holonomy).} Regrouping,
\begin{equation*}
H_\omega=u_0(x_1)\cdot\prod_{a=1}^{p}\bigl[u_a(x_a)^{-1}u_a(x_{a+1})\bigr]\cdot u_0(x_{p+1})^{-1}.
\end{equation*}
Both $x_a$ and $x_{a+1}$ lie in $X_a^{\sim}$, since $O_a,O_{a+1}\subset X_a^{\sim}$;
by the short-path connectedness statement $(\ast)$(iii) of the walk domains choose a lattice path
$\gamma_a\subset X_a^{\sim}$ from $x_a$ to $x_{a+1}$ of
$\ell^1$-length at most the diameter of $\square_a$, and, by the box argument following $(\ast)$(iii), a
path $\gamma_0\subset X_0$ from
$x_{p+1}$ to $x_1$ of $\ell^1$-length at most the diameter of $\square_0$ (both lie in $X_0$). Since
$X_a^{\sim}\subset\square_a\cap X$ for every $a$ and $X_0\subset X_0^{\sim}$, these paths lie in
$\square_a\cap X$, respectively $\square_0\cap X$, and \eqref{4.24:eq-field-dependence-5} applies to them. From
$U^{u_a}(\gamma_a)=u_a(x_a)U(\gamma_a)u_a(x_{a+1})^{-1}$,
\begin{equation*}
u_a(x_a)^{-1}u_a(x_{a+1})=W_a\,U(\gamma_a),\qquad
W_a:=u_a(x_a)^{-1}\bigl[U^{u_a}(\gamma_a)\bigr]^{-1}u_a(x_a),
\end{equation*}
and $|W_a-1|=|U^{u_a}(\gamma_a)-1|\le\theta$ by \eqref{4.24:eq-field-dependence-5}; in the same way
$u_0(x_{p+1})^{-1}u_0(x_1)=W_0U(\gamma_0)$ with $|W_0-1|\le\theta$. Let $\ell_\omega:=\gamma_1\gamma_2\cdots
\gamma_p\gamma_0$; it lies in $X_0\cup\bigcup_{a\ge1}X_a^{\sim}=D(\omega)$ and is the closed lattice loop from
$x_1$ around the chain of cubes back to
$x_1$, so that $U(\ell_\omega)=U(\gamma_1)\cdots U(\gamma_p)U(\gamma_0)$. Then
\begin{equation*}
H'_\omega:=u_0(x_1)^{-1}H_\omega u_0(x_1)=\Bigl[\prod_{a=1}^{p}W_aU(\gamma_a)\Bigr]W_0U(\gamma_0),
\end{equation*}
and the telescoping identity for products of unitaries gives
$|H'_\omega-U(\ell_\omega)|\le(p+1)\theta$. Since $|\cdot|$ is conjugation invariant,
$|\operatorname{Ad}_{H_\omega}-1|=|\operatorname{Ad}_{H'_\omega}-1|$. For $z\in Z(G)$,
$\operatorname{Ad}_z=1$, so $\operatorname{Ad}_{H'_\omega}=\operatorname{Ad}_{H'_\omega z^{-1}}$ and
$|H'_\omega z^{-1}-1|=|H'_\omega-z|$; hence
\begin{equation*}
|\operatorname{Ad}_{H_\omega}-1|\le2|H'_\omega-z|
\le2(p+1)\theta+2|U(\ell_\omega)-z|,
\end{equation*}
and we minimise over $z\in Z(G)$. The choice $z=1$ gives the weaker member $2|U(\ell_\omega)-1|$.

\emph{Step 8 (assembly).} By Steps~1--2,
$\operatorname{tr}_{\mathfrak{g}}t^{\xi}_{\omega}(\mathbb{U})(b',b')=\operatorname{tr}_{\mathfrak{g}}\Pi(b',b')$.
For $T:=t^{\xi}_{\omega}(\mathbf{1})(b',b')$ one has
$|\operatorname{tr}_{\mathfrak{g}}(T\operatorname{Ad}_{H})-\operatorname{tr}_{\mathfrak{g}}T|
\le n_{\mathfrak{g}}|T|\,|\operatorname{Ad}_H-1|$ and $|T|\le\prod_ar_a$ by $(\ast)$(i). Adding the bounds
of Steps~4, 5 and 7 and using $\prod_ar_a=m(\omega)$,
\begin{equation*}
\bigl|\operatorname{tr}_{\mathfrak{g}}t^{\xi}_{\omega}(\mathbb{U})(b',b')
-\operatorname{tr}_{\mathfrak{g}}t^{\xi}_{\omega}(\mathbf{1})(b',b')\bigr|
\le n_{\mathfrak{g}}\Bigl[(p+1)\epsilon+4p\theta+2(p+1)\theta
+2\min_{z\in Z(G)}|U(\ell_\omega)-z|\Bigr]m(\omega),
\end{equation*}
which is \eqref{4.24:eq-field-dependence-a-bis}, since $4p+2(p+1)=6p+2$. For $p=0$ there is one cube, there
is no transition, $h_1=u_0u_0^{-1}=1$ and $H_\omega=1$, so Steps~4, 6 and 7 contribute nothing and Step~5
alone gives the bound $n_{\mathfrak{g}}\epsilon\,m(\omega)$,
the single-cube form of the field-dependence bound at the full stage. Replacing $t$ by $\hat t$ changes
nothing, the subtracted term being field-free. Every quantity in \eqref{4.24:eq-field-dependence-a-bis} other
than the loop member is controlled by the per-cube clauses and by $(\ast)$(ii). The loop member is not:
a flat pair with non-central holonomy around an annular hull meets every per-cube clause with
smallness zero and moves $\operatorname{tr}_{\mathfrak{g}}\hat t^{\xi}_{\omega}$ at zeroth order. It is bounded
by Proposition~\ref{4.24:prop-holonomy-bound}, whose proof is outstanding, for walks with $p\ge1$ whose
domain is not contained in one clause cube. The single-cube members of the family, those counted by
$c_{\mathrm{sc}}$ in \eqref{4.24:eq-constant-unchanged-a}, are local holomorphic gauge-covariant functions of
$\mathbb{U}$ on one block pair; their field dependence is inside the full-stage field-dependence bound, by the
same passage to the gauge of the cube.

\emph{The block.} With $\delta\mathcal{A}:=\mathcal{A}-\mathcal{A}{\restriction}_{(1,0)}$ and
$(\delta\mathcal{A})_{ss}=P_s\,\delta\mathcal{A}\,P_s$,
\begin{equation*}
(\mathcal{A}_{ss}+\xi)^{-1}-(\mathcal{A}_{ss}{\restriction}_{(1,0)}+\xi)^{-1}
=-(\mathcal{A}_{ss}+\xi)^{-1}(\delta\mathcal{A})_{ss}(\mathcal{A}_{ss}{\restriction}_{(1,0)}+\xi)^{-1}.
\end{equation*}
By Notation~\ref{4.24:not-fibres-pairings}, $P_s$ commutes with every $\mathcal{U}_u$ and every cube
projection $P_\square$, so each cube-wise bound in the gauge of $\square$,
\begin{equation*}
\bigl\|P_\square\bigl(\mathcal{A}[\mathbb{U}^{u_\square},\mathbb{J}^{u_\square}]
-\mathcal{A}[\mathbf{1},0]\bigr)P_\square\bigr\|
\le O(1)(BCM\tilde\alpha_0+\tilde\alpha_1),
\end{equation*}
holds for the block with the same constant;
the global $\|\delta\mathcal{A}\|$ depends on the gauge and is not used. Both resolvents are expanded as in
Lemma~\ref{4.24:prf-block-walk-expansion}, and \eqref{4.24:eq-field-dependence-a-bis} holds term by term
for the walks of the block as for those of the full stage.

\emph{Part (c): absorption of the polynomial in $p$.} The count of
Lemma~\ref{4.24:prf-constant-unchanged} sums the walk majorants
$\bar m(\omega)e^{a_md_m(X(\omega))}$ of \eqref{4.24:eq-constant-unchanged-a} over walks, with
$c_G'$ densities per $\pi_{j+1}$-cell, by the geometric factor $(O(1)M^{-1/2})^{p}$ against the branching
number per step. Write the resulting convergent majorant series as $\sum_pN_px^{p}$, with $x$ the product of
the branching number per step and the geometric factor $O(1)M^{-1/2}$, and suppose that the lower bound on
$M$ under which
$\mathfrak{S}<\infty$ gives $x\le\frac12$. Since $(p+2)^2\le4(p+1)^2$,
\begin{equation*}
\sum_{p\ge0}(p+2)^2x^{p}\le4\sum_{p\ge0}(p+1)^2x^{p}=4(1+x)(1-x)^{-3}\le48
\quad\text{at }x\le\tfrac12,
\end{equation*}
so replacing $\bar m(\omega)$ by $(p+2)\bar m(\omega)$ or $(p+2)^2\bar m(\omega)$ multiplies the majorant of
$\mathfrak{S}$ by at most $48$; if the lower bound on $M$ gives only $x<1$, the constant is
$4(1+x)(1-x)^{-3}$, a number fixed together with $M$ before $\gamma$. Since $(p+1)\epsilon\le(p+2)\epsilon$
and $(6p+2)\theta\le6(p+2)\theta$, the first two members of \eqref{4.24:eq-field-dependence-a-bis} give,
through the count unchanged, a contribution to
$\|\mathbb{C}^{s}_{G}-\mathbb{C}^{s}_{G}{\restriction}_{(1,0)}\|$ which, with $\epsilon$ in its first form
$\epsilon=O(1)(BCM\tilde\alpha_0+\tilde\alpha_1)$, is at most
\begin{equation}\label{4.24:eq-field-dependence-6}
c_G'\cdot48\cdot O(1)\cdot n_{\mathfrak{g}}(\epsilon+6\theta)
\le c_G'\,O(1)\,dCM\,(BCM\tilde\alpha_0+\tilde\alpha_1),
\end{equation}
to which the contribution of the loop member is added.
This exceeds the bracket of the full-stage field-dependence bound by the absolute factor
$48n_{\mathfrak{g}}\cdot6dC$ and one factor $M$, coming from $\theta$. For a walk with $p\ge1$ whose
domain $D(\omega)$ lies in one clause cube $\square$, all $\square_a$ may be taken equal to $\square$ and all
$u_a$ equal to its gauge $u$; then $h_a=1$ for every $a$, so $T_a=1$, $g_a=1$ and $H_\omega=1$, and
Steps~4--7 show that \eqref{4.24:eq-field-dependence-a-bis} holds for such a walk with the $\theta$-members and
the loop member absent; its contribution is therefore included in the right side of
\eqref{4.24:eq-field-dependence-6}. For the remaining walks, under
Proposition~\ref{4.24:prop-holonomy-bound}, whose proof is outstanding, the loop member is at most
$2(p+1)^2\mathfrak{h}\,n_{\mathfrak{g}}\bar m(\omega)\le2(p+2)^2\mathfrak{h}\,n_{\mathfrak{g}}\bar m(\omega)$ per
term, and the same absorption bounds its contribution by $c_G'O(1)\mathfrak{h}$, with
$c_G'O(1)\mathfrak{h}\le\frac14\vartheta$ for all $g_k\le\gamma$. With $\epsilon$ in its second form, the
form in which the parameter conditions of Definition~\ref{4.24:def-normalisation-constants} carry it, the
member $K_\Delta\gamma_0^{(k)}$ is added inside the bracket, which gives the
first inequality of \eqref{4.24:eq-field-dependence-1-bis}, the $O(1)$ carrying $c_1\mathfrak{br}$.

\emph{The last inequality of \eqref{4.24:eq-field-dependence-1-bis} is a ceiling on the coupling.} By
\cite[(2.28)]{B-CMP119},
$\alpha_{0,l}=g_lC_0(\log g_l^{-2})^{q_0}$ and $\alpha_{1,l}=g_lC_1(\log g_l^{-2})^{q_1}$, and
$\tilde\alpha_0,\tilde\alpha_1$ are these with the brackets of \cite[(1.66)]{B-CMP122a}, at most $1$; so they
are polynomially small in $g$, while $\vartheta$ is $K_\vartheta$ times a constant times a negative power
of $\tfrac12\log g_k^{-2}$, that is, a constant multiple of a negative power of
$\mathsf{t}_k=\log g_k^{-2}$, and hence only logarithmically
small. The condition
\begin{equation*}
c_G'\,O(1)\,dCM\,\mathfrak{br}\,(BCM\tilde\alpha_0+\tilde\alpha_1)\le\tfrac12\vartheta
\end{equation*}
(here $O(1)$ denotes the constant of \eqref{4.24:eq-field-dependence-1-bis} with the factor $\mathfrak{br}$
taken out, so that it is free of the coupling)
therefore holds for all sufficiently small $g_k$, since $g$ times a power of $\log g^{-2}$ is at most any
prescribed negative power of $\log g^{-2}$ once $g$ is small enough; its constants $d,L,M,C,B$ are fixed before
$\gamma$, it bounds the coupling from above only, imposes no cap in the sense of
Definition~\ref{2.3:def-cap}, and leaves the choice of $L_0$ unchanged, $L$ entering only through
$c_G'=(LM)^4O(1)$. For the member $K_\Delta\gamma_0^{(k)}$, with
$\gamma_0^{(k)}=5\max_{l\le k}w_*(l)\bar\alpha_l$ as in Definition~\ref{4.24:def-normalisation-constants}
and $\bar\alpha_l$ a coupling amplitude, of the form $C_\iota\,x_l^{-1/2}(\log x_l)^{q_\iota}$ with
$\iota\in\{0,1\}$ and $x_l=g_l^{-2}$, as in \cite[(2.28)]{B-CMP119}, the full-stage
condition reads
\begin{equation*}
c_G'O(1)(BCM+1+K_\Delta)\gamma_0^{(k)}\le c_G'O(1)(BCM+1+K_\Delta)\bar\gamma_0(\mathsf{t}_k)
\le\vartheta(\mathsf{t}_k),
\end{equation*}
which follows from
\begin{equation}\label{4.24:eq-field-dependence-8}
\sup_{\mathsf{t}\ge\mathsf{t}_\gamma}
\frac{c_G'O(1)(BCM+1+K_\Delta)\,\bar\gamma_0(\mathsf{t})}{\vartheta(\mathsf{t})}\le1,
\qquad\mathsf{t}_\gamma:=\log\gamma^{-2},
\end{equation}
a condition on one variable $\mathsf{t}$: the ratio $\bar\gamma_0(\mathsf{t})/\vartheta(\mathsf{t})$ is a power
of $\mathsf{t}$ times $e^{-\mathsf{t}/2}$ and tends to zero faster than every power of $\mathsf{t}^{-1}$, so
\eqref{4.24:eq-field-dependence-8} holds for $\mathsf{t}_\gamma$ large, that is, for $\gamma$ small. Because
$\gamma_0^{(k)}$ is the coupling-dependent radius and not a number free of $g$, the condition is a ceiling
on the coupling and never a lower bound on it; this uses that $\vartheta$ is a power of
$\mathsf{t}^{-1}$. With the extra factor $dCM$ of
\eqref{4.24:eq-field-dependence-1-bis} (the member $\mathfrak{h}$ being bounded separately, by
$\frac14\vartheta$, below), the same ratio condition is read with its constant enlarged so that the
numerator is the whole bracket of \eqref{4.24:eq-field-dependence-9}, at the full stage as well as for the
block: this is condition
\eqref{4.24:eq-field-dependence-9} of part~(c). Each member of its left side is of one of the two kinds just
described, polynomially small in the coupling against the logarithmically small $\vartheta$, or a multiple of
$\gamma_0^{(k)}$, whose ratio to $\vartheta$ tends to zero faster than every power of $\mathsf{t}^{-1}$; so
\eqref{4.24:eq-field-dependence-9} is a ceiling on the coupling of the same kind as
\eqref{4.24:eq-field-dependence-8}, holding for $\gamma$ small and bounding the coupling from above only.
Under \eqref{4.24:eq-field-dependence-9} and the bound $c_G'O(1)\mathfrak{h}\le\frac14\vartheta$ of
Proposition~\ref{4.24:prop-holonomy-bound}, whose proof is outstanding, the right side of the first
inequality of \eqref{4.24:eq-field-dependence-1-bis} is at most $\frac34\vartheta+\frac14\vartheta=\vartheta$.
\end{proof}

\begin{proposition}[Holonomy bound for walks across several cubes; proof outstanding]
\label{4.24:prop-holonomy-bound}
Let $G=SU(2)$ be the gauge group and $Z(G)=\{\pm 1\}$ its centre. For a pair of the
stage write $\mathbb{U}=e^{i\eta A'}U$ as in \cite{B-CMP122a}, where $\eta$ is the lattice spacing
in which the conditions \cite[(1.64)--(1.69)]{B-CMP122a} are written, $U$ is the $G$-valued factor,
and $e^{i\eta A'}$ is the complex factor, $A'$ being its potential as in \cite{B-CMP122a}. The
complex factor is carried inside the form of each cube.
For a closed lattice loop $\ell$ let $U(\ell)$ be the transport of $U$ around $\ell$, that is,
the ordered product of $U$ along $\ell$. A \emph{clause cube} is a cube of the family on which the
conditions \cite[(1.64)--(1.69)]{B-CMP122a} are imposed; its side is $CML^m\eta$ with $C\le 2$,
where $m$ is the index of the shell in which the cube lies (a local index, distinct from the torus
exponent). We write $B$ for the constant multiplying $CM\tilde\alpha_0$ in those conditions, and
$N_0$ for the integer of \cite[p.~192]{B-CMP122a}: the numerical brackets multiplying
$BCM\tilde\alpha_0$ and $\tilde\alpha_1$ in \cite[(1.65)--(1.69)]{B-CMP122a} are either at most $1$
or powers of one fixed constant of that paper, of exponent at most $2N_0$.

There is a quantity $\mathfrak{h}=\mathfrak{h}(d,L,M,M_1,C,B;g_k)$, depending on the level $k$
only through $g_k$ and not on $N$, $N_0$, $|\mathfrak{W}|$ or the position of the walk, with
\begin{equation*}
c_G'\,O(1)\,\mathfrak{h}\le\tfrac14\vartheta\qquad\text{for all } g_k\le\gamma ,
\end{equation*}
where $O(1)$ denotes the constant, fixed before $\gamma$, with which the members of
\eqref{4.24:eq-field-dependence-a-bis}, summed over the walks and over the $c_G'$ densities of a cell,
enter the bound on the difference between $\mathbb{C}^{s}_G$ and the same family of densities at
the trivial pair, and
the constants $d,L,M,M_1,C,B$ being fixed before $\gamma$ in the fixed order of choice, such that
the following holds.

Let $b'$ be an anchor in the sense of Lemma~\ref{4.24:prf-hulls-anchors}. Let
$\omega=(X_0,\dots,X_p)$ with $p\ge1$ be a walk of the termwise walk
expansion of the diagonal entry at $b'$ of the block resolvent
(Lemma~\ref{4.24:prf-block-walk-expansion}), so that $X_0\ni b'$; this is an
expansion of the form \cite[(3.107)]{B-CMP99b}. Suppose that the domain $D(\omega)$ is not
contained in one clause cube. Let $\ell_\omega\subset D(\omega)$ be a closed lattice loop that
visits $X_0^{\sim},X_1^{\sim},\dots,X_p^{\sim},X_0^{\sim}$ in this order, with, for each $a$, one
segment inside $X_a^{\sim}$ whose $\ell^1$-length is at most the diameter of a clause cube
$\square_a$ containing $X_a^{\sim}$.

Consider every pair of $\mathfrak{N}^{(n+1)}(X(\omega))$, taken in the sense of the convention for
pairs on a hull, that is, as the restriction to $X(\omega)$ of a pair of
$\mathfrak{K}_P=\mathfrak{K}_{109}(\mathbb{B}^{(n)}_k;\gamma_0)$
on the whole domain of $\mathfrak{B}$. The $G$-valued factor $U$ of every such pair satisfies
\begin{equation}\label{4.24:eq-holonomy-bound-1}
\min_{z\in Z(G)}\bigl|U(\ell_\omega)-z\bigr|\le (p+1)^2\,\mathfrak{h}.
\end{equation}
\end{proposition}

Proof outstanding.

\smallskip
\noindent\textit{Route.}
Since only $\operatorname{Ad}_{U(\ell_\omega)}$ enters the field-dependence bound
\eqref{4.24:eq-field-dependence-a-bis}, and $\operatorname{Ad}_z=1$ for $z\in Z(G)$, it is the
distance of $U(\ell_\omega)$ to the centre, the left side of \eqref{4.24:eq-holonomy-bound-1},
that must be small.

Inequality \eqref{4.24:eq-holonomy-bound-1} bounds the loop member
$2\min_{z\in Z(G)}|U(\ell_\omega)-z|$ of \eqref{4.24:eq-field-dependence-a-bis}. Every other member
of that bound is controlled by estimates on single clause cubes, each factor of the walk term
being taken in the gauge of its own cube.

The same step, for walk terms that span several clause cubes, is taken in the field-dependence
bound of the normalisation lemma for the densities $\mathbb{C}_G$ of the full stage, whose
objects are fixed in Definition~\ref{4.24:def-normalisation-constants}, where that lemma rests on
the present proposition. That lemma is named here as the place where the step occurs, not as a
justification of it.

When $\ell_\omega$ bounds on the torus, let $S$ be a lattice $2$-chain spanning $\ell_\omega$.
The lattice non-abelian Stokes formula reduces \eqref{4.24:eq-holonomy-bound-1} to two inputs.

The first input is a bound on the curvature $|U(\partial P)-1|$ for every plaquette $P$ of $S$
(the letter $P$ is used for plaquettes, $p$ being the length of the walk).
The second input is the count
\begin{equation*}
\sum_{P\in S}L^{-2m_P}\le C_{\mathrm{S}}\,(p+1)^2 .
\end{equation*}
Here $m_P$ is the index of the shell containing $P$, so that $L^{m_P}\eta$ is the length unit of
that shell, and $C_{\mathrm{S}}$ is a constant fixed before $\gamma$. For the
pairs of $\mathfrak{N}^{(n+1)}(X(\omega))$ the conditions \cite[(1.64)--(1.69)]{B-CMP122a} are
imposed with the domain written $X$ there taken to be the hull $X(\omega)$. The
shells are the sets $(\Omega_m\setminus\Omega_{m+1})\cap X(\omega)$, where
$\Omega_m\supset\Omega_{m+1}$ are the nested domains entering
\cite[(1.64), (1.66), (1.68)]{B-CMP122a}. Two cases arise.

\emph{Case $(\alpha)$: $\ell_\omega$ bounds a $2$-chain $S$ inside $X(\omega)$.} The plaquette
bound on $S$ is then available. It is the member of
\cite[(1.64), (1.66), (1.68)]{B-CMP122a} bounding $|U(\partial P)-1|$ on the shells
$(\Omega_m\setminus\Omega_{m+1})\cap X(\omega)$.
It enters through the curvature entry of the space $\mathfrak{N}^{(n+1)}$.

What remains in this case is the count above, adapted to the shells. The cone over $\ell_\omega$
is not a sufficient choice of $S$, because its count fails when $S$ crosses several shells.

\emph{Case $(\beta)$: $\ell_\omega$ does not bound inside $X(\omega)$.} This happens for hulls of
$\mathbb{D}_m$-domains with holes, which do occur, and for loops winding a period of the torus, if
such loops are reached; a loop winding a period of the torus bounds no $2$-chain on the torus, and
the reduction above does not apply to it as stated. In this case, when $\ell_\omega$ still bounds
on the torus and $S$ is a spanning $2$-chain, no plaquette bound on the cells
of $S$ outside $X(\omega)$ is supplied by the following:
\begin{itemize}
\item the conditions \cite[(1.64)--(1.69)]{B-CMP122a} and \cite[(3.35)]{B-CMP99b}, which are
conditions on $X(\omega)$ only;
\item the bounds in the gauge of a single cube, which are per cube.
\end{itemize}
A proof must therefore draw the curvature bound on the cells of $S$ outside $X(\omega)$ from the
pair of $\mathfrak{K}_P$ on the whole domain of $\mathfrak{B}$, in the unit of the walk's scale and
across shells. It must also establish the count. For loops winding a period of the torus a route
other than a spanning $2$-chain is needed.

Per-cube bounds alone cannot give \eqref{4.24:eq-holonomy-bound-1}. A flat pair whose holonomy
around an annular hull is not central meets every condition on a single clause cube with
smallness $0$. Yet for such a pair the density
$\operatorname{tr}_{\mathfrak{g}}\hat t^{\xi}_{\omega}(b',b')$, $\xi\ge0$ being the resolvent
parameter of Lemma~\ref{4.24:prf-block-walk-expansion}, differs from its value at the trivial pair
by a term of zeroth order in the smallness parameters.

Suppose that walk terms winding a hole occur among the walks of
Lemma~\ref{4.24:prf-block-walk-expansion}, and that the interior of the hole
carries no small-curvature bound commensurate with the walk's scale. Then no bound of the form
\eqref{4.24:eq-holonomy-bound-1} with $\mathfrak{h}$ as above can be obtained for those terms. For
them the field-dependence bound
holds only up to the trivial bound $2n_{\mathfrak{g}}\bar m(\omega)$. Summed over such terms,
this gives a fixed tail that is small in $M/M_1$ but not small in $\gamma$.

\begin{lemma}[Termwise integrability on the half-line]\label{4.24:prf-halfline-integrability}
Assume the hypotheses of Lemma~\ref{4.24:lem-block-coercivity}. Let $\lambda_0>0$ be the constant
fixed in Lemma~\ref{4.24:lem-log-resolvent}, and for $b'\in s$ let $\mathfrak{e}(b')$ be the bond
density \eqref{4.24:eq-bond-densities-a}.
\begin{enumerate}
\item[(a)] For every $\xi\ge 0$ the integrand of $\mathfrak{e}(b')$ is
$-\tfrac12\operatorname{tr}_{\mathfrak{g}}$ of the $(b',b')$ block of
$(\mathcal{A}_{ss}+\xi)^{-1}(\mathcal{A}_{ss}-\lambda_0)(\lambda_0+\xi)^{-1}$, whose norm satisfies
\begin{equation}\label{4.24:eq-halfline-integrability-1}
\bigl\|(\mathcal{A}_{ss}+\xi)^{-1}(\mathcal{A}_{ss}-\lambda_0)(\lambda_0+\xi)^{-1}\bigr\|
\le\frac{\|\mathcal{A}_{ss}\|+\lambda_0}{(a_{\mathbb{C}}+\xi)(\lambda_0+\xi)}
=O(\xi^{-2}).
\end{equation}
The $\xi$-integral in $\mathfrak{e}(b')$ therefore converges absolutely, locally uniformly on the
domain of the variables on which Lemma~\ref{4.24:lem-block-coercivity} holds.
\item[(b)] The same holds termwise in the walk expansion \eqref{4.24:eq-region-walk-expansion-b}
of the resolvent $(\mathcal{A}_{ss}+\xi)^{-1}$ entering $\mathfrak{e}(b')$.
\begin{itemize}
\item The class of walks $\omega$ with $|\omega|=0$ is kept paired with
$(\lambda_0+\xi)^{-1}\cdot 1$, as in \cite[(3.36)--(3.37)]{B-CMP119}. Its $\xi$-integrand is then
$O(\xi^{-2})$.
\item Every term with $|\omega|\ge 1$ is bounded in norm by $(a_{\mathbb{C}}+\xi)^{-2}$ times a
factor independent of $\xi$.
\item The summable majorants $\bar m(\omega)$, uniform in $\xi$, that control the terms after the
$\xi$-integration are those of \eqref{4.24:eq-region-walk-expansion-b}.
\end{itemize}
\end{enumerate}
\end{lemma}

\begin{proof}
\emph{Coercivity and the resolvent bound.} By Lemma~\ref{4.24:lem-block-coercivity}, the
coercivity bound \eqref{4.24:eq-block-coercivity-a} holds with $a_{\mathbb{C}}=(4B_0)^{-1}$:
$\operatorname{Re}\langle\phi,\mathcal{A}_{ss}\phi\rangle\ge a_{\mathbb{C}}\|\phi\|^2$ for all
$\phi\in\mathcal{H}_s$. In the notation of Lemma~\ref{4.24:lem-log-resolvent} this says
$\mathcal{A}_{ss}\in\mathfrak{P}(a_{\mathbb{C}})$.

Let $B\in\mathfrak{P}(a)$ and $\xi\ge0$. Then
$\operatorname{Re}\langle\phi,(B+\xi)\phi\rangle\ge(a+\xi)\|\phi\|^2$, and the Cauchy--Schwarz
inequality gives $\|(B+\xi)\phi\|\ge(a+\xi)\|\phi\|$. This is the Rayleigh-quotient and
Lax--Milgram argument of \eqref{4.24:eq-principal-blocks-a}, here for a general
$B\in\mathfrak{P}(a)$. It yields invertibility of $B+\xi$ and the bound
$\|(B+\xi)^{-1}\|\le(a+\xi)^{-1}$.

The bound is inherited by every compression $P_E B P_E$ of $B$ to a coordinate subspace
$\mathcal{H}_E$, because $\langle\phi,P_EBP_E\phi\rangle=\langle\phi,B\phi\rangle$ for
$\phi\in\mathcal{H}_E$.

\emph{Part} (a). By the resolvent identity of Lemma~\ref{4.24:lem-log-resolvent}, for
$B\in\mathfrak{P}$ and $\xi\ge0$,
\begin{equation}\label{4.24:eq-halfline-integrability-2}
(\lambda_0+\xi)^{-1}\cdot 1-(B+\xi)^{-1}
=(B+\xi)^{-1}\bigl[(B+\xi)-(\lambda_0+\xi)\bigr](\lambda_0+\xi)^{-1}
=(B+\xi)^{-1}(B-\lambda_0)(\lambda_0+\xi)^{-1}.
\end{equation}
The integrand of $\mathfrak{e}(b')$ in \eqref{4.24:eq-bond-densities-a} is
$-\tfrac12\operatorname{tr}_{\mathfrak{g}}$ of the $(b',b')$ block of the left side of
\eqref{4.24:eq-halfline-integrability-2} with $B=\mathcal{A}_{ss}$. The $(b',b')$ block of an
operator on $\mathcal{H}_s$ has norm at most the norm of that operator. Combining
$\|\mathcal{A}_{ss}-\lambda_0\|\le\|\mathcal{A}_{ss}\|+\lambda_0$ with the resolvent bound at
$a=a_{\mathbb{C}}$ gives \eqref{4.24:eq-halfline-integrability-1}.

Since $(a_{\mathbb{C}}+\xi)(\lambda_0+\xi)\ge(\min\{a_{\mathbb{C}},\lambda_0\}+\xi)^2$ and
\begin{equation*}
\int_0^\infty\frac{d\xi}{(\min\{a_{\mathbb{C}},\lambda_0\}+\xi)^2}
=\frac{1}{\min\{a_{\mathbb{C}},\lambda_0\}}<\infty,
\end{equation*}
the integral of the block converges absolutely. The same holds for
$-\tfrac12\operatorname{tr}_{\mathfrak{g}}$ of it, a linear functional on a finite-dimensional
space.

For the local uniformity, note that $\mathcal{A}_{ss}=P_s\mathcal{A}P_s$ on $\mathcal{H}_s$, so
$\|\mathcal{A}_{ss}\|\le\|\mathcal{A}\|$. The norm $\|\mathcal{A}\|$ is bounded, locally uniformly
on the domain of the variables on which Lemma~\ref{4.24:lem-block-coercivity} holds, by the
kernel bound for the stage form kernel $K_\sharp$ together with Schur's test, as in
\eqref{4.24:eq-principal-blocks-a}. With $a_{\mathbb{C}}$ and $\lambda_0$ fixed, the right side of
\eqref{4.24:eq-halfline-integrability-1} therefore has an integrable majorant that is locally
uniform on that domain, and the convergence is locally uniform there.

\emph{Part} (b), the class $|\omega|=0$. This class stays paired with $(\lambda_0+\xi)^{-1}\cdot 1$
exactly as in \cite[(3.36)--(3.37)]{B-CMP119}. Its $\xi$-integrand is therefore a difference
of the form \eqref{4.24:eq-halfline-integrability-2}, with $B$ a compression
$P_E\mathcal{A}_{ss}P_E$ of $\mathcal{A}_{ss}$ to a local block. Such a compression satisfies
$\operatorname{Re}\langle\phi,P_E\mathcal{A}_{ss}P_E\phi\rangle\ge a_{\mathbb{C}}\|\phi\|^2$ for
$\phi\in\mathcal{H}_E$ by the inheritance noted above, and has norm at most
$\|\mathcal{A}_{ss}\|$. The argument of part (a) then bounds this integrand by
\begin{equation*}
\frac{\|\mathcal{A}_{ss}\|+\lambda_0}{(a_{\mathbb{C}}+\xi)(\lambda_0+\xi)}=O(\xi^{-2}).
\end{equation*}

\emph{Part} (b), the terms with $|\omega|\ge1$. By \cite[(3.107)]{B-CMP99b}, each such term
contains at least two local resolvent factors. Each of these is the inverse of a $\xi$-shifted
compression of the coercive form, so its norm is at most $(a_{\mathbb{C}}+\xi)^{-1}$ by the
Lax--Milgram estimate above, that is, by the argument of \eqref{4.24:eq-principal-blocks-a}. The
term is therefore $O\bigl((a_{\mathbb{C}}+\xi)^{-2}\bigr)$ times factors independent of $\xi$, and
$(a_{\mathbb{C}}+\xi)^{-2}$ is integrable on $[0,\infty)$ with integral $a_{\mathbb{C}}^{-1}$.

\emph{Part} (b), the majorants. The majorants $\bar m(\omega)$ of the terms after the
$\xi$-integration, uniform in $\xi$ and summable over $\omega$, are those supplied by the walk
expansion \eqref{4.24:eq-region-walk-expansion-b}.

No positivity of individual walk terms is used in this proof; within it, positivity enters only
through the coercivity of Lemma~\ref{4.24:lem-block-coercivity}.
\end{proof}

\begin{lemma}[No new condition on the parameters]\label{4.24:prf-no-new-condition}
The bounds of clause~(i) of the main statement of this section hold with constants that do
not depend on the region of integration, on the number of components, on the family
$\mathfrak{W}$ of arrested components, on the stage, or on the levels. Their proof imposes
no condition on the parameters beyond those already in force. In particular, it introduces
no cap in the sense of Definition~\ref{2.3:def-cap}, no lower bound on the running couplings
$g_k$ and no new floor, and it leaves the order of choice unchanged.
\end{lemma}

\begin{proof}
We go through every quantity that could produce a volume factor, a dependence on the
components or on the stage, or a condition on the parameters, and check each in turn.

\emph{The index set of the densities.} The densities are indexed by the survivors' bonds
$s$. They are counted per cell, through the per-cell count $c_G'$. They are never summed
over the region, over the components, or over $\mathfrak{W}$.

\emph{Walk and animal sums.} Every walk sum is taken per anchor, and every animal sum per
cell, by the per-cell form of the animal bound. The norm in which
clause~(i) is stated is a supremum over cells. Hence no volume factor arises.

\emph{Components.} The components enter only through the separation of components and of
cells meeting one component (Lemma~\ref{4.24:lem-component-separation}), and only through
two of its assertions. One of these assertions is used in the bound on the Gaussian
normalisation and in two of the estimates on the localised densities. The other assertion,
namely that every
wall-connected family of host cubes containing the host cubes of integration bonds of two
distinct components has at least $R_k$ members, is used in one of these estimates, and
there it enters as the exponent of a smallness factor. In neither use does the number of
components appear.

\emph{The family $\mathfrak{W}$ and the stage.} The family $\mathfrak{W}$ is a fixed datum
of the term of the expansion under consideration. No sum over $\mathfrak{W}$, over arrests
or over stages occurs, and $|\mathfrak{W}|$ appears nowhere. The index $n^{+}=n+1$ serves
only as a label of the stage.

\emph{Levels and the radii $R_i$.} No dependence on the level $k$ or on the number $K$ of
renormalisation steps is introduced. The
constant $c_G$ of Definition~\ref{4.24:def-normalisation-constants} and
Lemma~\ref{4.24:prf-constant-unchanged} is uniform in $k$ and $n$, by the uniformity
statement for the normalisation constants on which
Definition~\ref{4.24:def-normalisation-constants} rests. The filing levels
$m\in\{j+1,\dots,k\}$ enter only under $\sup_m$, as in the definition of $\mathfrak{S}$ in
Lemma~\ref{4.24:prf-constant-unchanged}. The quantity $\check{R}$ appears in exactly two
places: in the total $V$ of the cluster expansion, through its $\tilde{\mathfrak{l}}$-term,
and in the component budget $\bar V$. It does not appear in clause~(i). No quantity
depending on the number of steps since a large-field region was created enters.

\emph{Conditions involving $g$.} Exactly three conditions involving $g$ are used, and each
is a ceiling of the kind already imposed through the parameter $\gamma$, already in force.
\begin{itemize}
\item[(a)] The condition $2\alpha_5B_0<\tfrac12$ of clause~(c) of the coercivity of the
survivors' block (Lemma~\ref{4.24:lem-block-coercivity}).
\item[(b)] The requirement in Lemma~\ref{4.24:prf-field-dependence} (stated; proof
incomplete at the step marked $(\ast)$) that the bracket of field-dependence quantities be
at most $\vartheta$. The member $K_\Delta\gamma_0$ of that bracket is bounded by means of a
ratio ceiling that is among the conditions already in
force. Here $\gamma_0=\gamma_0^{(k)}$ is the radius fixed in
Definition~\ref{4.24:def-normalisation-constants}, and this ceiling also bounds the running
couplings from above.
\item[(c)] The condition $R_k\ge 3$ used in Lemma~\ref{4.24:lem-component-separation}.
\end{itemize}
Each of these is an upper bound on the running couplings $g_k$. No lower bound on the $g_k$
is used. No upper bound on $L$, $M$, $\kappa$, $\kappa_1$, or on the parameter
$\mathfrak{M}$ that the order of choice fixes after $L$ as a function $\mathfrak{M}(R_1)$ of
$R_1$, is used, and no cap in the sense of Definition~\ref{2.3:def-cap} is introduced.

\emph{Floors and the order of choice.} The only lower bounds on parameters that are used are
the floors already imposed among the parameter conditions of
Definition~\ref{4.24:def-normalisation-constants}. No new floor is introduced. The order of
choice ($L$ first, then $\mathfrak{M}(R_1)$, then $\kappa$, $\kappa_1$, $M_1$, $M$ and
finally $\gamma$) is therefore untouched.

\emph{The coercivity constant.} The constant $a_{\mathbb{C}}=(4B_0)^{-1}$ of
\eqref{4.24:eq-block-coercivity-a} involves $B_0=B_0(d,L)$, which depends only on $d$ and
$L$. It is independent of the region, and it enters no constant of the statement.

Together, these observations give the three assertions of the lemma: the constants are
uniform in the listed data, no condition on the parameters is added, and the order of
choice is unchanged.
\end{proof}

\begin{lemma}[Elimination factors carried by the arrested components]\label{4.24:prf-elimination-factors}
Consider stage $n+1$, performed at level $j$, with $\mathfrak{W}$ the set of arrested components and
$Z^{s}$ the survivors'
region, and let $W\in\mathfrak{W}$. For a function $f$ of the complex background pair
$(\mathbb{U},\mathbb{J})$ write $f|_{(1,0)}$ for its value at the trivial pair
$(\mathbb{U},\mathbb{J})=(1,0)$. Let $c$ range over the cells of
Definition~\ref{4.24:def-normalisation-constants} with $c\subset R^{(n+1)}\cap W$.
\begin{enumerate}
\item[(a)] The product $\prod_{c}z^{(j)}(c)$ multiplies the integral over $W$, which is not performed
at this stage. The factors $z^{(j)}(c)$ are single-cube local functions,
independent of the further argument $x$ on which densities of the stage may depend. They are not
factors of the product $z^{(j)}$ of \cite[(1.60)]{B-CMP122a}.
Hence they do not enter the product $z^{(j)}_s$ of
Definition~\ref{4.24:def-normalisation-constants}, nor the left-hand side of clause~(i) of
Theorem~\ref{4.24:thm-survivor-block-expansion}.
\item[(b)] Assigning the field parts
\begin{equation*}
\log z^{(j)}(c)-\log z^{(j)}(c)|_{(1,0)},\qquad c\subset R^{(n+1)}\cap W,
\end{equation*}
as members of the family $\mathbb{C}^{\mathrm{fr}}_W$ of frozen densities of $W$ is an assignment of
the same kind as those made where that family is defined, and changes no density.
The family of these members is denoted $\mathbb{C}^{G}_W$. Each member is a single-block term indexed
by a domain $X\in\mathbb{D}_{j+1}$. It is independent of $x$, of the integration field $A_j$ of the
unperformed integral over $W$ and of the further argument $\lambda$ on which densities of the stage
may depend, has no strong set
($\mathcal{S}=\emptyset$), carries no cross-member majorant $\nu^{\natural}$, and has analyticity
domain $\mathfrak{D}^{\flat}_{n+1}[(\mathrm{amb})_W](X)$, the analyticity domain of the stage
densities on $X$ attached to $W$. The numbers
$\sum_c\log z^{(j)}(c)|_{(1,0)}$ and $E_0^{(j)}|_W$ are field-independent constants of
$\mathbf{T}'_k(W)$; here $\mathbf{T}'_k(W)$ is Ba{\l}aban's large-field operation $\mathbf{T}'_k$
(Notation~\ref{0.2:not-glossary}) on $W$, and $E_0^{(j)}|_W$ is the value of the additive set
function $E_0^{(j)}$ on the part of its region lying in $W$. These constants are among the constants
$E_k(Z)$ of \cite{B-CMP122b}, and they are estimated in no norm.
\item[(c)] With this assignment, each $\pi_{j+1}$-cell carries at most $c_G'$ members of
$\mathbb{C}^{s}_G\sqcup\bigsqcup_{W\in\mathfrak{W}}\mathbb{C}^{G}_W$, with the majorants of
$\mathbb{C}^{s}_G$. This family satisfies, with the same constant $c_G$, the weighted per-cell bound
satisfied by $\mathbb{C}^{s}_G$. Granted the conversion of per-rate bounds to the stage norm
(Proposition~\ref{4.24:prop-norm-conversion}, whose proof is outstanding), it is bounded by $3c_G$ in
the stage norm $\|\cdot\|_{(n+1)}$. Granted the holonomy bound
(Proposition~\ref{4.24:prop-holonomy-bound}, whose proof is outstanding) and
Lemma~\ref{4.24:prf-field-dependence}, whose proof is incomplete at one step, its field part is
bounded by $\vartheta$. In particular, under these antecedents the joint bound $3c_G+\vartheta$ for
these densities is the same as for $\mathbb{C}^{s}_G$ alone.
\item[(d)] In Ba{\l}aban's own placement, $\mathbb{C}^{G}_W:=0$ and the product
$\prod_c z^{(j)}(c)$ stays whole as a factor of $\mathbf{T}'_k(W)$ beside $\mathsf{q}_W$, joining the
list of such factors of $\mathbf{T}'_k(W)$. Under it every count in~(c) decreases or stays the same.
\end{enumerate}
\end{lemma}

\begin{proof}
(a) We use the order of the operations of the stage, as fixed where the operations of the stage are
described. In
\cite[p.~188, preceding (1.55)]{B-CMP122a} the $\delta$-functions are removed using the operator $C$
introduced there. Then the scaling transformation $B'_j=g_jB_j$ of the field variables
$B_j$ is performed. Both operations are carried out on all of $Z^{\sharp}$, and they precede
the split into survivors and arrested components.

By Definition~\ref{4.24:def-normalisation-constants}, the elimination produces one factor
$z^{(j)}(c)$ for each cell $c$ of the region in which the $\delta$-functions are eliminated. Each
factor is a local function of the single cube $c$ and does not depend on $x$. Since $W$ is a
component of $Z^{\sharp}$, every cell $c\subset R^{(n+1)}\cap W$ contributes such a factor. This
factor multiplies the integral over $W$, which is not performed at this stage.

The product $z^{(j)}$ of \cite[(1.60)]{B-CMP122a} runs over the region written $Z$ there, which after
the split is the survivors' region $Z^{s}$, by the re-denotation $Z=Z^{s}$ fixed with the notation of
the stage. The arrested component $W$ is disjoint from $Z^{s}$.
Hence no cell of $R^{(n+1)}\cap W$ occurs in $z^{(j)}$. Consequently no such cell occurs in
$z^{(j)}_s=\prod_{c\subset R^{(n+1)}\cap Z^{s}}z^{(j)}(c)$. Therefore no such cell occurs in the
object on the left-hand side of clause~(i) of Theorem~\ref{4.24:thm-survivor-block-expansion}.

In Ba{\l}aban's construction these factors belong to the operation $\mathbf{T}'_k(W)$, which is not
performed on $W$ at this stage \cite{B-CMP122a}. They stand beside $\mathsf{q}_W$ and the factor of
\cite[(1.59)]{B-CMP122a}.

(b) For each $c$ we have the identity
\begin{equation*}
\log z^{(j)}(c)=\bigl(\log z^{(j)}(c)-\log z^{(j)}(c)|_{(1,0)}\bigr)
+\log z^{(j)}(c)|_{(1,0)} .
\end{equation*}
Summing over $c\subset R^{(n+1)}\cap W$ and exponentiating gives the factor $\prod_c z^{(j)}(c)$
as a product of two parts. The first is the exponential of the sum of the members of
$\mathbb{C}^{G}_W$. The second is the exponential of the constant $\sum_c\log z^{(j)}(c)|_{(1,0)}$.
The constant $E_0^{(j)}|_W$ is left unchanged. The assignment therefore only fixes the family in
which each part is counted, and it changes no density.

Each factor $z^{(j)}(c)$ is a local function on the single cube $c$
(Definition~\ref{4.24:def-normalisation-constants}); it is independent of $x$, involves no
integration variable $A_j$ and no $\lambda$, and for the constant part it is evaluated at the fixed
trivial pair. The members have empty strong set, $\mathcal{S}=\emptyset$, and, reading no element,
carry no cross-member majorant
$\nu^{\natural}$. The constant parts do not depend on the field,
so they are counted among the constants $E_k(Z)$ of \cite{B-CMP122b}, as stated where the outputs of
the stage are described, and no norm is applied to them.

(c) We count block by block, that is, per $\pi_{j+1}$-cell. Each such cell meets at most one
component, so there are three cases.
\begin{itemize}
\item A survivors' cell carries at most $d(LM)^d$ bond densities and at most $dM^d$ single-block
terms of $\mathbb{C}^{s}_G$.
\item A cell holding blocks of $R^{(n+1)}\cap W$ carries at most $dM^d$ members of
$\mathbb{C}^{G}_W$ and no bond density.
\item Every other cell carries no member.
\end{itemize}
The count in a cell of the second kind is therefore at most the bound for a survivors' cell. By
Definition~\ref{4.24:def-normalisation-constants}, the bound for a survivors' cell lies within the
per-cell count $c_G'$. Hence the per-cell count of the enlarged family is at most $c_G'$.

Each $z^{(j)}(c)$, $c\subset R^{(n+1)}\cap W$, is the same single-cube function as the factors
entering $z^{(j)}_s$ (Definition~\ref{4.24:def-normalisation-constants}), so its field part, read on
the domain $\mathfrak{D}^{\flat}_{n+1}[(\mathrm{amb})_W](X)$, has the majorant of the single-block
terms of $\mathbb{C}^{s}_G$. The constant $c_G=c_G'\,\mathfrak{S}+c_{\mathrm{sc}}$ of
Lemma~\ref{4.24:prf-constant-unchanged} depends only on the per-cell count $c_G'$, the universal
series $\mathfrak{S}$ and the single-cube budget $c_{\mathrm{sc}}$, none of which the enlarged family
changes. Hence the family
$\mathbb{C}^{s}_G\sqcup\bigsqcup_{W}\mathbb{C}^{G}_W$ satisfies, with the same constant $c_G$, the
weighted per-cell bound satisfied by $\mathbb{C}^{s}_G$.

The bound $3c_G$ in $\|\cdot\|_{(n+1)}$ follows from the conversion of per-rate bounds to the stage
norm (Proposition~\ref{4.24:prop-norm-conversion}, whose proof is outstanding). It is applied to the
enlarged family exactly as to $\mathbb{C}^{s}_G$.

The bound $\vartheta$ on the field part follows as for $\mathbb{C}^{s}_G$ in
Lemma~\ref{4.24:prf-field-dependence}, whose proof is incomplete at one step. That argument uses the
holonomy bound (Proposition~\ref{4.24:prop-holonomy-bound}, whose proof is outstanding). A
single-block term depends only on the variables of one block pair.

Since $c_G$ and $\vartheta$ are unchanged, so is the joint bound $3c_G+\vartheta$.

(d) Under Ba{\l}aban's own placement $\mathbb{C}^{G}_W:=0$, a cell holding blocks of
$R^{(n+1)}\cap W$ carries no member, and every other cell carries what it carries in~(c). Hence every
per-cell count, and with it every count entering the bounds of~(c), decreases or stays the same.
\end{proof}

\begin{lemma}[Conclusion of the block normalisation]\label{4.24:prf-normalisation-conclusion}
Work in the setting of the stage, with $E_0^{(j)}$ and $Z''_{j+1}$ as fixed in its construction
and with $z^{(j)}_s$ and $z^{(j)}(c)$ as in Lemma~\ref{4.24:prf-bond-densities}. Then
\begin{equation}\label{4.24:eq-normalisation-conclusion-1}
\begin{split}
&\log\Big[z^{(j)}_s\int d\phi_s\,
 e^{-\frac12\langle\phi_s,\mathcal{A}_{ss}\phi_s\rangle}\Big]
 +E_0^{(j)}\big((\Omega^{c}_{j+1}\setminus Z''_{j+1})\cap Z^{s}\big)\\
&\qquad=\sum_{b'\in s}\mathfrak{e}(b')
 +\sum_{c\subset R^{(n+1)}\cap Z^{s}}\log z^{(j)}(c)
 +\big(\text{numbers per cell of }Z^{s}\big)\\
&\qquad=\sum_{X}\mathbb{C}^{s}_G(X),
\end{split}
\end{equation}
where the third term of the middle member is a sum of numbers, one attached to each cell of
$Z^{s}$, and the last sum runs over the hulls $X$ of \eqref{4.24:eq-hulls-anchors-a};
the family $\{\mathbb{C}^{s}_G(X)\}_X$ satisfies every clause of part~(i) of the block
normalisation statement, the clause on the field dependence of the densities in cube gauges
being the conclusion of Lemma~\ref{4.24:prf-field-dependence}, stated there with proof
incomplete at the step marked $(\ast)$. Clause~(iv) holds for this family.
\end{lemma}

\begin{proof}
The first equality in \eqref{4.24:eq-normalisation-conclusion-1} is obtained by collecting the
statements of this section. The normalisation as a sum of bond densities
(Lemma~\ref{4.24:prf-bond-densities}) writes the first term of the left member of
\eqref{4.24:eq-normalisation-conclusion-1} as
$-\frac12\operatorname{Tr}F(\mathcal{A}_{ss})$ plus
$\sum_{c\subset R^{(n+1)}\cap Z^{s}}\log z^{(j)}(c)$ plus a constant. It also writes
$-\frac12\operatorname{Tr}F(\mathcal{A}_{ss})$ as the sum over $b'\in s$ of the bond densities
$\mathfrak{e}(b')$ of \eqref{4.24:eq-bond-densities-a}. The constant of that lemma and the term
$E_0^{(j)}$ evaluated on $(\Omega^{c}_{j+1}\setminus Z''_{j+1})\cap Z^{s}$ are collected into the
numbers attached to the cells of $Z^{s}$ in the middle member.

The second equality and the clauses of part~(i) are collected from four further statements.
\begin{itemize}
\item Lemma~\ref{4.24:prf-hulls-anchors}, with the hulls, filing and anchors of
\eqref{4.24:eq-hulls-anchors-a}, files the three terms of the middle member into
$\sum_X\mathbb{C}^{s}_G(X)$ and gives the holomorphy of the densities.
\item Lemma~\ref{4.24:prf-constant-unchanged} shows that the normalisation constant is
unchanged.
\item Lemma~\ref{4.24:prf-elimination-factors} gives the elimination factors carried by the
arrested components.
\item The field dependence of the densities in cube gauges is the conclusion of
Lemma~\ref{4.24:prf-field-dependence}, stated there with proof incomplete at the step
marked $(\ast)$.
\end{itemize}
Clause~(iv) for this family is Lemma~\ref{4.24:prf-no-new-condition}: no new condition on the
parameters arises.
\end{proof}

\begin{remark}[Filing component by component gives the same bound]\label{4.24:rem-per-component-filing}
Let $c$ range over the components of $Z^{s}$. For each $c$ put
\[
\rho(c):=(\mathsf{R}_j\cup\mathcal{C}_j)\cap c ,
\]
let $s_c$ be the set of integration bonds of $\rho(c)$, let $\phi_c:=\phi|_{s_c}$ be the restriction
of the integration variable to $s_c$, put $\mathcal{A}_{cc}:=P_{s_c}\mathcal{A}P_{s_c}$, and let
$z^{(j)}_c$ be the product of the factors $z^{(j)}$ over the cells contained in $\rho(c)$. Let
\[
\mathfrak{E}^{\times\times}:=\bigl\{e_a(b,b',Y'):\ b\in s_{c_1},\ b'\in s_{c_2},\
c_1\neq c_2\subset Z^{s}\bigr\}
\]
be the family of collar members joining two distinct surviving components. In the proof of
clause (i) of the statement of this section on the survivors' integration, the survivors'
Gaussian has the precision $\mathcal{A}_{ss}$ on all of $s$, and only the entries between $s$ and
$\mathcal{W}$ are assigned to elements. Assign instead the entries of $\mathcal{A}_{ss}$ between
distinct surviving components, that is, the off-block remainder $\mathsf{X}$, to the element list
$\mathfrak{E}^{\times\times}$, so that the survivors' Gaussian becomes the product over $c$ of the
Gaussians with precisions $\mathcal{A}_{cc}$. This assignment leads to the same bounds. The
logarithm of the Gaussian normalisation is literally a sum over $c$. The union
$\mathbb{C}^{s}_G:=\bigsqcup_c\mathbb{C}^{c}_G$ of the per-component families $\mathbb{C}^{c}_G$
of normalisation densities obeys the bound $c_G'\mathfrak{S}+c_{\mathrm{sc}}\le c_G$
of \eqref{4.24:eq-constant-unchanged-a}. Under the same hypotheses as for the survivors' block
(the holonomy-certificate hypothesis for the field-dependent part, the convergence condition on
the stage norms for the norm of stage $n+1$), its field-dependent part is at most $\vartheta$
(by the field-dependence bound \eqref{4.24:eq-field-dependence-a-bis}, which is stated with proof
incomplete at the step marked $(\ast)$) and its norm of stage $n+1$ is at most $3c_G$. The
enlarged cross family
$\mathfrak{E}^{\times}\sqcup\mathfrak{E}^{\times\times}$
has $\mu$-total at most $\varepsilon_YK_a(\varrho^{B}_j)^2\le\vartheta_a$, where $K_a$ is the
constant of \eqref{4.24:eq-stage-form-a} and $\varrho^{B}_j$ is the constant of the same name in the
bound for $\mathfrak{E}^{\times}$ in the proof of clause (ii), and $\mathsf{M}_*$-total at
most $\vartheta_a$. The coefficient $5$ and the constant $C_0^{\sharp}$ of the element bound in which
this family enters are unchanged.

The lemma that localises the logarithm of the normalisation of the stage Gaussian into the
densities $\mathbb{C}_G$ (in what follows, the normalisation lemma) is itself also a
per-component statement, for three reasons. Its
integration region is a union of whole components. Its requirement that every hull $X$ meet $R_n$
is a condition on each anchor separately. Its constant $c_G$ is counted per $\pi_{j+1}$-cell.

The assignment on the survivors' block, used in the proof of clause (i), needs no structural
statement about the entries of $\Delta^{(j)}_{\Lambda_{j+1}}$ between distinct components. Whatever
such entries the operator $\mathcal{A}$ of Definition~\ref{4.24:def-stage-form} keeps, that proof
runs on the principal block and assigns the entries between $s$ and $\mathcal{W}$ to elements. It
also leaves unchanged the assignment, recorded in Definition~\ref{4.24:def-stage-form}, of the members
$e_a$ and $e_a^{0}$ with both bonds off $\bigcup\mathfrak{W}$ to the survivors' Gaussian.

The per-component assignment is the one of Ba{\l}aban's construction, and it makes the logarithm
of the Gaussian normalisation literally a sum over the components (Step~1 of the proof below).
It changes that assignment for the members $e_a$ joining two distinct
surviving components, which become elements of $\mathfrak{E}^{\times\times}$. The bounds of this
section are proved with the assignment on the survivors' block; the per-component assignment is
the equivalent one recorded in this remark.
\end{remark}

\begin{proof}
By \eqref{4.24:eq-component-separation-a} (Lemma~\ref{4.24:lem-component-separation}), every member
$e_a(b,b',Y')\in\mathfrak{E}^{\times\times}$ has $|Y'|\ge R_k$.

\emph{Step 1: an identity, with no Schur complement.} Consider clause (iii) of
\eqref{4.24:eq-stage-form-a} (Definition~\ref{4.24:def-stage-form}). It writes the stage quadratic
form as a sum of members
indexed by pairs of bonds $(b,b')$, and this identity holds pair by pair. We group the pairs into
four classes:
\begin{itemize}
\item both bonds in one surviving component;
\item the two bonds in two distinct surviving components (\emph{cross});
\item one bond in a surviving component and one in $\mathcal{W}$ (\emph{mixed});
\item both bonds in $\mathcal{W}$.
\end{itemize}
For an element $e$ write $v_e$ for its contribution to the exponent. The grouping gives
\begin{equation}\label{4.24:eq-per-component-filing-1}
\begin{aligned}
-\tfrac12\langle\mathsf{B},\mathcal{A}\mathsf{B}\rangle
&=\sum_{c\subset Z^{s}}\Bigl[-\tfrac12\bigl\langle\mathsf{B}|_{s_c},
\mathcal{A}_{cc}\,\mathsf{B}|_{s_c}\bigr\rangle\Bigr]
+\sum_{(b,b')\ \text{cross or mixed}}\ \sum_{Y'}v_{e_a(b,b',Y')}\\
&\quad+\sum_{(b,b')\in\mathcal{W}^2}\Bigl[v_{e_a^{0}(b,b')}+\sum_{Y'}v_{e_a(b,b',Y')}\Bigr].
\end{aligned}
\end{equation}
No member $e_a^{0}(b,b')$ has its two bonds in distinct components; this is shown for the members
$e_a^{0}$ in the proof of the normalisation lemma. Every cross or
mixed member $e_a(b,b',Y')$ has $|Y'|\ge R_k$, by
\eqref{4.24:eq-component-separation-a}.

Three further pieces of the integrand also split by components. The characteristic functions
$\chi^{(j)}$ and $\chi'^{(j)}$, the per-cell factors $z^{(j)}$ and the factor $E_0$ of the integrand
factorise over the
components of $Z^{s}$, by clause (i) of the lemma of this section on the normalisation ratios. The
rim members are attached to single components:
in Ba{\l}aban's construction the off-diagonal quadratic terms are placed in the action (among the
terms that Ba{\l}aban calls $B$-terms, \cite{B-CMP119}) and not in the Gaussian measure.

Read \eqref{4.24:eq-per-component-filing-1} in the integration variables, where
$\mathsf{B}|_{s_c}=\phi_c$. The integrand of the $d\phi_s$-integral is then
\begin{equation}\label{4.24:eq-per-component-filing-2}
\prod_{c\subset Z^{s}}\Bigl[z^{(j)}_c\,\chi^{(j)}_c\,\chi'^{(j)}_c\,
e^{-\frac12\langle\phi_c,\mathcal{A}_{cc}\phi_c\rangle}\Bigr]
\cdot\exp\Bigl\{\sum_{e\in\mathfrak{E}^{\mathrm{StB}}\sqcup\mathfrak{E}^{\times\times}}
v_e(\phi_s;\lambda)\Bigr\}\cdot F_{\mathrm{ext}}(\lambda).
\end{equation}
Here $\chi^{(j)}_c$ and $\chi'^{(j)}_c$ are the factors of $\chi^{(j)}$ and $\chi'^{(j)}$ belonging
to $c$, $\lambda$ stands for the variables of the elements other than $\phi_s$, and
$F_{\mathrm{ext}}(\lambda)$ collects the factors that do not depend on $\phi_s$.

The only Gaussian measures in \eqref{4.24:eq-per-component-filing-2} are therefore the
per-component ones. The normalisation of the one belonging to $c$ is a power of $2\pi$ depending
only on $|s_c|$, times $e^{-\frac12\operatorname{Tr}F(\mathcal{A}_{cc})}$, with $F$ the function
of Lemma~\ref{4.24:lem-log-resolvent}. The logarithm of the product of these normalisations is
therefore, literally, the sum over $c$ of the logarithms of the factors.

\emph{Step 2: the normalisation lemma read with integration region $\rho(c)$.} We read the
normalisation lemma with its integration region taken to be $\rho(c)$. For each $c$, three inputs
of that lemma hold with $s$ replaced by $s_c$, with no change in the argument:
\begin{itemize}
\item the gauge covariance of the block;
\item the coercivity of Lemma~\ref{4.24:lem-block-coercivity};
\item the properties of $F$ in Lemma~\ref{4.24:lem-log-resolvent}.
\end{itemize}
For the coercivity this is the restriction of a Rayleigh quotient from $\mathcal{H}$ to
$\mathcal{H}_{s_c}\subset\mathcal{H}$, with the same constant $a_{\mathbb{C}}$ of
\eqref{4.24:eq-block-coercivity-a}.

Two remarks relate coercivity of the block $P_s\mathcal{A}P_s$ to coercivity of
$\bigoplus_c\mathcal{A}_{cc}$. Their difference is the off-block remainder $\mathsf{X}$.
The remark of this section following Lemma~\ref{4.24:lem-block-coercivity} bounds it by
\[
\|\mathsf{X}\|\le B_{\mathcal{A}}C_\rho e^{-\delta D/2},
\]
where $D$ is the gap between the bond sets of distinct components and $\delta$ is the decay rate,
both as in that remark; so $\mathsf{X}$ is an operator of size $e^{-c'R_k}$ for a constant $c'>0$.
\begin{itemize}
\item If one of the two operators, $B$, satisfies $\operatorname{Re}\langle\varphi,B\varphi\rangle\ge
a_{\mathbb{C}}\|\varphi\|^2$ and $\|\mathsf{X}\|\le a_{\mathbb{C}}/2$, then the other, $B\pm\mathsf{X}$,
satisfies
\[
\operatorname{Re}\langle\varphi,(B\pm\mathsf{X})\varphi\rangle\ge(a_{\mathbb{C}}-\|\mathsf{X}\|)
\|\varphi\|^2\ge\tfrac12a_{\mathbb{C}}\|\varphi\|^2 .
\]
So coercivity of either operator implies coercivity of the other with constant $a_{\mathbb{C}}/2$.
\item In the direction from $\mathcal{A}$ to $\bigoplus_c\mathcal{A}_{cc}$ there is no loss at all.
Write $P_c:=P_{s_c}$. The sets $s_c$ partition $s$, so the coordinate subspaces $\mathcal{H}_{s_c}$
are mutually orthogonal with orthogonal sum $\mathcal{H}_s$. Since
$\mathcal{A}_{cc}=P_c\mathcal{A}P_c$, for $\varphi\in\mathcal{H}_s$ we get
\begin{equation}\label{4.24:eq-per-component-filing-3}
\operatorname{Re}\Bigl\langle\varphi,\bigoplus_c\mathcal{A}_{cc}\varphi\Bigr\rangle
=\sum_c\operatorname{Re}\langle P_c\varphi,\mathcal{A}P_c\varphi\rangle
\ge a_{\mathbb{C}}\sum_c\|P_c\varphi\|^2=a_{\mathbb{C}}\|\varphi\|^2 .
\end{equation}
The inequality is the coercivity of Lemma~\ref{4.24:lem-block-coercivity}, applied to each
$P_c\varphi\in\mathcal{H}_{s_c}\subset\mathcal{H}_s$.
\end{itemize}

The remaining steps of the normalisation lemma also run per component: the walk expansion of
$-\frac12\operatorname{Tr}F$ and the bounds on its densities, up to and including the
field-dependence bound \eqref{4.24:eq-field-dependence-a-bis}, which is stated with proof incomplete
at the step marked $(\ast)$.
For them we take the region $\Lambda'$ of the walk expansion on unions of
whole components, \eqref{4.24:eq-region-walk-expansion-a}, to be $\rho(c)$. This is the smallest
region of the kind to which that expansion applies. The covariance is
\[
\mathcal{C}_c:=\mathcal{A}_{cc}^{-1}=C^{(j)}(\rho(c)),
\]
that is, the covariance of Definition~\ref{4.24:def-stage-form} conditioned to the subdomain
$\rho(c)$ (the letter $\mathcal{C}_c$ denotes a covariance, not the collar bond set $\mathcal{C}_j$),
which is Ba{\l}aban's object \cite[Theorem~3.15]{B-CMP99b}.
Each term of the expansion carries one walk, so no junction of hulls
inside the normalisation lemma joins two components. The anchor bond $b'$ of every term lies in
$s_c$; hence every hull $X$ meets $R_n\cap c$, and no term is anchored on $\bigcup\mathfrak{W}$.

\emph{Step 3: the union family.} Let $\mathbb{C}^{c}_G$ be the family of densities produced in
Step~2 for the component $c$, and put $\mathbb{C}^{s}_G:=\bigsqcup_c\mathbb{C}^{c}_G$. Let
$\mathsf{Q}(c)$ be the set of $\pi_{j+1}$-cells meeting $s_c$. Fix $(m,\Delta)$. The sum over
anchor cells in the weighted norm in which $c_G$ of \eqref{4.24:eq-constant-unchanged-a} is
measured runs over
\[
q\in\mathsf{Q}^{s}=\bigsqcup_c\mathsf{Q}(c).
\]
By \eqref{4.24:eq-component-separation-c}, each such cell meets $s_c$ for exactly one component
$c$, so the union is disjoint and no
term is counted twice. The sum is majorised by the universal majorant series $\mathfrak{S}$ of
\eqref{4.24:eq-constant-unchanged-a}, and in that norm
\begin{equation}\label{4.24:eq-per-component-filing-4}
\|\mathbb{C}^{s}_G\|\le c_G'\,\mathfrak{S}+c_{\mathrm{sc}}\le c_G .
\end{equation}
Two further bounds hold on the same footing.
\begin{itemize}
\item The field-dependent part of $\mathbb{C}^{s}_G$ is at most $\vartheta$. This follows from the
field-dependence bound \eqref{4.24:eq-field-dependence-a-bis} applied per component, under the
holonomy-certificate hypothesis on which that bound rests; that bound is stated with proof
incomplete at the step marked $(\ast)$.
\item $\mathbb{C}^{s}_G$ is at most $3c_G$ in the norm of stage $n+1$, under the convergence
condition on the stage norms on which the same bound for the survivors' block rests.
\end{itemize}

\emph{Step 4: the cross family re-assigned to elements costs nothing new.} The two assignments
differ in one respect only: whether the cross-survivor entries $\mathsf{X}$ stand in the Gaussian
or in the element list $\mathfrak{E}^{\times\times}$. The change from one to the other is an exact
re-assignment of terms.

Consider the bounds for $\mathfrak{E}^{\times}$ proved in the estimate of the cross family in the
proof of clause (ii). They are anchored at a survivor bond. The partner bond is summed over all
bonds through $C_\rho$, and $Y'$ over all families with $|Y'|\ge R_k$. Whether the partner lies in
$\mathcal{W}$ or in another surviving component makes no difference to the majorant. By the same
computation, therefore, the enlarged family
$\mathfrak{E}^{\times}\sqcup\mathfrak{E}^{\times\times}$ has
\[
\mu\text{-total}\le\varepsilon_YK_a(\varrho^{B}_j)^2\le\vartheta_a,
\qquad
\mathsf{M}_*\text{-total}\le\vartheta_a ,
\]
with $K_a$ and $\varrho^{B}_j$ as in the bound for $\mathfrak{E}^{\times}$.
That computation packs over $R_k$-separated components, using \eqref{4.24:eq-packing-count-a}
(Lemma~\ref{4.24:lem-packing-count}). The coefficient $5$ and the constant $C_0^{\sharp}$ of the
element bound in which this family enters are unmoved.

Clause (ii) is then read at the product $\bigotimes_c\langle\cdot\rangle_c$ of the per-component
expectations. By clause (i) of the lemma of this section on the normalisation ratios, at
decoupling parameter $\mathsf{s}=0$ off $Z$ the Gaussian structure is
a product over the components of $Z$; hence every estimate is made component by component. This
proves clause (i) with the per-component assignment of the cross-survivor entries, and the
assertions of the remark.
\end{proof}

\begin{lemma}[Square root of the block covariance]\label{4.24:lem-covariance-root}
Let $\mathfrak{P}$ and $\mathfrak{P}(a)$, $a>0$, be the sets of operators on $\mathcal{H}_s$ introduced in
Lemma~\ref{4.24:lem-log-resolvent}. Let $\langle\cdot,\cdot\rangle$ and $\|\cdot\|$ be the Hermitian inner
product and norm on $\mathcal{H}_s$ in terms of which $\mathfrak{P}$ and $\mathfrak{P}(a)$ are defined, and let
$(\cdot,\cdot)$ and ${}^T$ be the bilinear pairing and the transpose of
Lemma~\ref{4.24:lem-principal-blocks}. For $a>0$ and $B\in\mathfrak{P}(a)$ put
\begin{equation}\label{4.24:eq-covariance-root-1}
S(B):=\frac{1}{\pi}\int_0^\infty \xi^{-1/2}\,(B+\xi)^{-1}\,d\xi .
\end{equation}
\begin{enumerate}
\item[(a)] The integral converges absolutely in norm, and $\|S(B)\|\le a^{-1/2}$.
\item[(b)] $S(B)=B^{-1/2}$, the inverse square root in the holomorphic functional calculus with the
principal branch, and $S(B)^2=B^{-1}$.
\item[(c)] If $B^T=B$, then $S(B)^T=S(B)$.
\item[(d)] $S$ is holomorphic on $\mathfrak{P}$.
\item[(e)] In the setting and under the hypotheses of Lemma~\ref{4.24:lem-block-coercivity}, the block
$\mathcal{A}_{ss}$ belongs to $\mathfrak{P}(a_{\mathbb{C}})$, with $a_{\mathbb{C}}$ as in
\eqref{4.24:eq-block-coercivity-a}, and is symmetric. With
$\mathcal{C}_s^{1/2}:=S(\mathcal{A}_{ss})$ one has $(\mathcal{C}_s^{1/2})^2=\mathcal{A}_{ss}^{-1}$,
$\|\mathcal{C}_s^{1/2}\|\le a_{\mathbb{C}}^{-1/2}$, and for every $\phi'\in\mathcal{H}_s$ the
substitution $\phi_s=\mathcal{C}_s^{1/2}\phi'$ gives
\begin{equation}\label{4.24:eq-covariance-root-2}
(\phi_s,\mathcal{A}_{ss}\phi_s)=(\phi',\phi').
\end{equation}
\end{enumerate}
\end{lemma}

\begin{proof}
(a) Let $\xi\ge0$ and $\varphi\in\mathcal{H}_s$. Since $B\in\mathfrak{P}(a)$,
\begin{equation*}
(a+\xi)\|\varphi\|^2\le \operatorname{Re}\langle\varphi,(B+\xi)\varphi\rangle
\le\|\varphi\|\,\|(B+\xi)\varphi\| .
\end{equation*}
Hence $B+\xi$ is injective, so invertible on the finite-dimensional space $\mathcal{H}_s$, and
$\|(B+\xi)^{-1}\|\le(a+\xi)^{-1}$. The norm of the integrand in \eqref{4.24:eq-covariance-root-1} is
therefore at most $\xi^{-1/2}(a+\xi)^{-1}$, which is integrable on $(0,\infty)$. Consequently
\begin{equation*}
\|S(B)\|\le\frac{1}{\pi}\int_0^\infty\frac{\xi^{-1/2}\,d\xi}{a+\xi}.
\end{equation*}
The substitution $\xi=au^2$, $d\xi=2au\,du$, turns the last integral into
\begin{equation*}
2a^{-1/2}\int_0^\infty\frac{du}{1+u^2}=\pi a^{-1/2},
\end{equation*}
so $\|S(B)\|\le a^{-1/2}$.

(b) The same substitution with $a$ replaced by $\zeta>0$ gives
\begin{equation}\label{4.24:eq-covariance-root-3}
\frac{1}{\pi}\int_0^\infty\xi^{-1/2}(\zeta+\xi)^{-1}\,d\xi=\zeta^{-1/2}
\end{equation}
for $\zeta\in(0,\infty)$. The identity extends to $\operatorname{Re}\zeta>0$, with the principal
branch on the right. Both sides are holomorphic there. For the right side this is the choice of branch.
For the left side, let $a'>0$. On $\{\operatorname{Re}\zeta\ge a'\}$ one has
$|\zeta+\xi|\ge a'+\xi$. Hence the integrals over $[1/n,n]$, which are holomorphic in $\zeta$,
converge uniformly there. By Weierstrass' theorem the limit is holomorphic. The two sides agree on
$(0,\infty)$, so they agree on the right half-plane by the identity theorem.

We now lift \eqref{4.24:eq-covariance-root-3} to $B$ by the holomorphic functional calculus used in the
proof of Lemma~\ref{4.24:lem-log-resolvent}. Let $\lambda$ be an eigenvalue of $B$ with eigenvector
$\varphi$. Then
$\operatorname{Re}\lambda\,\|\varphi\|^2=\operatorname{Re}\langle\varphi,B\varphi\rangle\ge a\|\varphi\|^2$,
so the spectrum of $B$ lies in $\{\operatorname{Re}\zeta\ge a\}$, inside the half-plane on which both sides
of \eqref{4.24:eq-covariance-root-3} are holomorphic. For each $\xi\ge0$ the function
$\zeta\mapsto(\zeta+\xi)^{-1}$ is holomorphic on $\operatorname{Re}\zeta>-\xi$, and its lift is
$(B+\xi)^{-1}$. Applying the calculus to \eqref{4.24:eq-covariance-root-3} therefore gives
$S(B)=B^{-1/2}$, and, the calculus being multiplicative, $S(B)^2=B^{-1}$.

(c) Transposition is linear and continuous on the finite-dimensional space of operators on
$\mathcal{H}_s$. It therefore passes inside the absolutely convergent integral
\eqref{4.24:eq-covariance-root-1}. If $B^T=B$, then
$((B+\xi)^{-1})^T=((B+\xi)^T)^{-1}=(B+\xi)^{-1}$. Hence $S(B)^T=S(B)$.

(d) Let $S_n(B)$ be the integral \eqref{4.24:eq-covariance-root-1} restricted to $[1/n,n]$. The map
$B\mapsto(B+\xi)^{-1}$ is holomorphic on $\mathfrak{P}$. Together with its derivative
$H\mapsto-(B+\xi)^{-1}H(B+\xi)^{-1}$ it is continuous on $[1/n,n]\times\mathfrak{P}$. Differentiating
under the integral sign shows that $S_n$ is holomorphic on $\mathfrak{P}$. For $B\in\mathfrak{P}(a)$,
part (a) gives
\begin{equation*}
\|S(B)-S_n(B)\|\le\frac1\pi\Bigl(\int_0^{1/n}\frac{\xi^{-1/2}}{a}\,d\xi+\int_n^\infty\xi^{-3/2}\,d\xi\Bigr)
=\frac{2}{\pi}\bigl(a^{-1}+1\bigr)n^{-1/2},
\end{equation*}
uniformly in $B\in\mathfrak{P}(a)$. By Lemma~\ref{4.24:lem-log-resolvent}, every point of $\mathfrak{P}$
lies in some $\mathfrak{P}(a)$ and has a neighbourhood contained in $\mathfrak{P}(a/2)$. So
$S_n\to S$ locally uniformly on $\mathfrak{P}$, and $S$ is holomorphic there by Weierstrass' theorem.

(e) By Lemma~\ref{4.24:lem-principal-blocks}, $\mathcal{A}_{ss}^T=\mathcal{A}_{ss}$. By
Lemma~\ref{4.24:lem-block-coercivity},
$\mathcal{A}_{ss}\in\mathfrak{P}(a_{\mathbb{C}})$ with $a_{\mathbb{C}}$ as in
\eqref{4.24:eq-block-coercivity-a}. Parts (a)--(c) with $B=\mathcal{A}_{ss}$ and $a=a_{\mathbb{C}}$ give
$(\mathcal{C}_s^{1/2})^2=\mathcal{A}_{ss}^{-1}$,
$\|\mathcal{C}_s^{1/2}\|\le a_{\mathbb{C}}^{-1/2}$ and
$(\mathcal{C}_s^{1/2})^T=\mathcal{C}_s^{1/2}$. Each $(\mathcal{A}_{ss}+\xi)^{-1}$ commutes with
$\mathcal{A}_{ss}$, so $\mathcal{C}_s^{1/2}$ does too. Therefore, for $\phi'\in\mathcal{H}_s$ and
$\phi_s=\mathcal{C}_s^{1/2}\phi'$,
\begin{align*}
(\phi_s,\mathcal{A}_{ss}\phi_s)
&=(\phi',(\mathcal{C}_s^{1/2})^T\mathcal{A}_{ss}\mathcal{C}_s^{1/2}\phi')
=(\phi',\mathcal{A}_{ss}(\mathcal{C}_s^{1/2})^2\phi')\\
&=(\phi',\mathcal{A}_{ss}\mathcal{A}_{ss}^{-1}\phi')=(\phi',\phi'),
\end{align*}
which is \eqref{4.24:eq-covariance-root-2}.
\end{proof}

\begin{lemma}[The Cauchy bound per extracted cube]\label{4.24:lem-cube-extraction-bound}
Let $\kappa_1>0$ and put $R:=e^{\kappa_1}$. Let $g$ be a function of finitely many complex
variables $\mathsf{s}(\Delta)$, one for each cube $\Delta$ of a finite set $\mathcal{D}$ of cubes,
holomorphic on an open neighbourhood of the closed polydisc
$\{\mathsf{s}: |\mathsf{s}(\Delta)|\le R \text{ for all } \Delta\in\mathcal{D}\}$, and let $G\ge 0$
be such that $|g|\le G$ on this polydisc. Let $Y\subset\mathcal{D}$. Then
\begin{equation}\label{4.24:eq-cube-extraction-bound-a}
\begin{split}
\Bigl|\Bigl(\prod_{\Delta\in Y}\int_0^1 d\mathsf{s}(\Delta)\,\partial_{\mathsf{s}(\Delta)}\Bigr)
 g\,\Big|_{\mathsf{s}(\Delta)=0,\ \Delta\notin Y}\Bigr|
 &\le G\,\bigl(R(R-1)^{-2}\bigr)^{|Y|}\\
 &\le G\,e^{-(\kappa_1-1)|Y|},
\end{split}
\end{equation}
where the second inequality holds for $\kappa_1\ge 1$.
\end{lemma}

\begin{proof}
Put $\mathsf{s}_Y=(\mathsf{s}(\Delta))_{\Delta\in Y}$ and let $h(\mathsf{s}_Y)$ be the value of $g$
at the point with coordinates $\mathsf{s}_Y$ on $Y$ and $\mathsf{s}(\Delta)=0$ for
$\Delta\notin Y$. This point lies in the polydisc of the hypothesis whenever
$|\mathsf{s}(\Delta)|\le R$ for all $\Delta\in Y$, so $h$ is holomorphic on an open neighbourhood
of the closed polydisc of radius $R$ in the variables $\mathsf{s}_Y$, and $|h|\le G$ there.
Since $h$ is holomorphic, its partial derivatives are continuous and commute, and
differentiation in one variable commutes with integration over $[0,1]$ in the others. Hence the
quantity inside the absolute value in \eqref{4.24:eq-cube-extraction-bound-a} equals
\begin{equation*}
\int_{[0,1]^{Y}}\Bigl(\prod_{\Delta\in Y}\partial_{\mathsf{s}(\Delta)}\Bigr)h(\mathsf{s}_Y)\,
d\mathsf{s}_Y .
\end{equation*}

Fix $\mathsf{s}_Y\in[0,1]^{Y}$. Applying the one-variable Cauchy formula for the first
derivative successively in each variable $\mathsf{s}(\Delta)$, $\Delta\in Y$, on the circle
$|\zeta_\Delta|=R$, with $\zeta_Y=(\zeta_\Delta)_{\Delta\in Y}$, gives
\begin{equation*}
\Bigl(\prod_{\Delta\in Y}\partial_{\mathsf{s}(\Delta)}\Bigr)h(\mathsf{s}_Y)
=\prod_{\Delta\in Y}\Bigl(\frac{1}{2\pi i}\oint_{|\zeta_\Delta|=R}d\zeta_\Delta\Bigr)\,
\frac{h(\zeta_Y)}{\prod_{\Delta\in Y}(\zeta_\Delta-\mathsf{s}(\Delta))^{2}} .
\end{equation*}
On the contours
\begin{equation*}
|\zeta_\Delta-\mathsf{s}(\Delta)|\ \ge\ R-|\mathsf{s}(\Delta)|\ \ge\ R-1\ >\ 0
\qquad(\Delta\in Y),
\end{equation*}
because $0\le\mathsf{s}(\Delta)\le 1<R$; moreover $|h(\zeta_Y)|\le G$, and each circle has length
$2\pi R$.
Therefore
\begin{equation*}
\Bigl|\Bigl(\prod_{\Delta\in Y}\partial_{\mathsf{s}(\Delta)}\Bigr)h(\mathsf{s}_Y)\Bigr|
\le G\prod_{\Delta\in Y}\frac{2\pi R}{2\pi (R-1)^{2}}
= G\,\bigl(R(R-1)^{-2}\bigr)^{|Y|}.
\end{equation*}
Integrating this bound over $[0,1]^{Y}$, which has volume $1$, gives the first inequality of
\eqref{4.24:eq-cube-extraction-bound-a}. This first inequality uses only $R>1$, that is,
$\kappa_1>0$.

For the second inequality it suffices that $R(R-1)^{-2}\le e^{-(\kappa_1-1)}=e\,R^{-1}$. This is
equivalent to $R^{2}\le e\,(R-1)^{2}$, that is, since $R>1$, to $R\le e^{1/2}(R-1)$, that is, to
$R(1-e^{-1/2})\ge 1$, and hence to
\begin{equation*}
\kappa_1\ \ge\ \ln\frac{1}{1-e^{-1/2}}=0.9327\ldots .
\end{equation*}
If $\kappa_1\ge 1$ this holds: indeed $R\ge e$, and $e^{1/2}<17/10$ because
$(17/10)^{2}=289/100>e$; since moreover $e>27/10$, we get
\begin{equation*}
R(1-e^{-1/2})\ \ge\ e(1-e^{-1/2})\ =\ e-e^{1/2}\ >\ \tfrac{27}{10}-\tfrac{17}{10}\ =\ 1 .
\end{equation*}
Raising $R(R-1)^{-2}\le e^{-(\kappa_1-1)}$ to the power
$|Y|$ and multiplying by $G$ gives the second inequality. In the application of this lemma the
convergence condition imposed on the auxiliary parameters gives $\kappa_1\ge 8$, so the
hypothesis $\kappa_1\ge 1$ is met.
\end{proof}

The bound \eqref{4.24:eq-cube-extraction-bound-a} is the Cauchy estimate on the circles
$|\mathsf{s}(\Delta)|=e^{\kappa_1}$ by which the decoupling parameters $\mathsf{s}(\Delta)$ of the
covariance's walks are extracted, cube by cube, in the estimates of this chapter that use it.

\begin{lemma}[Gaussian inputs at the block: conditioning, absorption and covariance bounds]
\label{4.24:prf-gaussian-inputs}
Work at a stage of the step with the following data:
\begin{itemize}
\item the symmetrised operator $\mathcal{A}$ on $\mathcal{H}$ of the stage form (Definition~\ref{4.24:def-stage-form});
\item the survivors' bonds $s$ and the arrested bonds $\mathcal{W}$, with $\mathfrak{Bo}=s\sqcup\mathcal{W}$;
\item the survivors' block $\mathcal{A}_{ss}=P_s\mathcal{A}P_s$ on $\mathcal{H}_s$ and the survivors' covariance
$\mathcal{C}_s$, whose precision is $\mathcal{A}_{ss}$ by Lemma~\ref{4.24:prf-block-conditioned}, so that
$\mathcal{C}_s=\mathcal{A}_{ss}^{-1}$;
\item the survivors' integration region $\Lambda'_s$, the cells $\mathsf{Q}^{s}$, the survivors' characteristic functions $\chi^{s}$ and the expectation $\langle\cdot\rangle_s$;
\item the coercivity constant $a_{\mathbb{C}}=(4B_0)^{-1}$ of \eqref{4.24:eq-block-coercivity-a}.
\end{itemize}
Then the following Gaussian-dependent hypotheses of the Gaussian integration proposition at the full stage, of the cluster-expansion proposition at the stage and of the tolerance bound on $\mathfrak{N}(c)$ hold at the block, with the same constants.
\begin{enumerate}
\item[(i)] \emph{Conditioning.} The survivors' Gaussian is centred with precision $\mathcal{A}_{ss}$. For every union of cubes $Z_1\subset\mathsf{Q}^{s}$, the Gaussian object attached to $Z_1$ by the conditioning step is one of the following:
\begin{itemize}
\item the principal block $(\mathcal{A}_{ss})_{Z_1Z_1}$, which equals $\mathcal{A}_{Z_1Z_1}$;
\item a restriction $(\mathcal{C}_s)_{Z_1Z_1}$ of the covariance or of its decoupled family, which enters only through bounds.
\end{itemize}
\item[(ii)] \emph{Absorption.} The Hermitian-part bound used to absorb the quadratic factors of the bonds of a bond set $P\subset s$ holds on $\mathcal{H}_s$ with the same $a_{\mathbb{C}}$. The absorption therefore holds with the same $\gamma_2$, the same $T_j$ and the same margin $\alpha_5\|\mathcal{C}\|<\tfrac12$.
\item[(iii)] \emph{Covariance bounds.} $\|\mathcal{C}_s\|\le 2B_0$, and independently $\|\mathcal{C}_s\|\le 4B_0$.
\item[(iv)] \emph{Whitening.} The square root $\mathcal{C}_s^{1/2}$ exists, satisfies $\|\mathcal{C}_s^{1/2}\|\le a_{\mathbb{C}}^{-1/2}$ and is holomorphic in the pair.
\item[(v)] \emph{Decoupling of walks.} Each resolvent $(\mathcal{A}_{ss}+\xi)^{-1}$, $\xi\ge0$, carries a generalised random walk expansion with majorants of the same form and the same $(d,L)$-constants as at the full stage. The decoupled family is obtained by multiplying each walk by the decoupling parameters of the cells it meets.
\end{enumerate}
In particular, for every union of cubes $Z_1\subset\mathsf{Q}^{s}$ and every $\varphi\in\mathcal{H}_s$,
\begin{equation}\label{4.24:eq-gaussian-inputs-a}
\begin{gathered}
(\mathcal{A}_{ss})_{Z_1Z_1}=\mathcal{A}_{Z_1Z_1},\qquad
\operatorname{Re}\langle\varphi,\mathcal{A}_{ss}\varphi\rangle\ge a_{\mathbb{C}}\|\varphi\|^2,\\
\|\mathcal{C}_s\|\le 2B_0,\qquad \|\mathcal{C}_s^{1/2}\|\le a_{\mathbb{C}}^{-1/2}.
\end{gathered}
\end{equation}
\end{lemma}

\begin{proof}
We treat the five inputs in turn.

\emph{Conditioning.} The stage form couples $s$ and $\mathcal{W}$ through its $s\times\mathcal{W}$ part. That part is the element family $\mathfrak{E}^{\times}$ of the cluster expansion and does not stand in the exponent of the survivors' Gaussian. Hence the exponent of the $\phi_s$-integral is $-\tfrac12\langle\phi_s,\mathcal{A}_{ss}\phi_s\rangle$. It contains no linear term in $\phi_s$, so the survivors' Gaussian has no mean and no Schur complement is formed. It is therefore centred with precision $\mathcal{A}_{ss}$.

The conditioning step of the cluster expansion, \cite[(2.5)]{B-CMP116} together with \cite[(2.6)--(2.8)]{B-CMP116}, is used with $Z_0$ the set of cubes of $\mathsf{Q}$ containing $Y_0\cup P$. Here $Y_0$ and $P$ are the bond sets of a term $\mathcal{T}(D,P,Z)$ of the Mayer expansion. This step consists of algebra and Fubini's theorem only. It is therefore an identity for any finite index set and any union of cubes $Z_0\subset\mathsf{Q}^{s}$; in general $Z_0$ is a union of connected components.

Let $Z_1\subset\mathsf{Q}^{s}$ be a cube set. By \eqref{4.24:eq-conditioning-marginal-a}, the Gaussian object that the step attaches to $Z_1$ is one of two things.
\begin{itemize}
\item[(c1)] In the case that this object is the conditional precision, obtained by conditioning $\phi$ on the bonds of $s$ off $Z_1$ to $0$ (the case used by the Gaussian integration proposition at the full stage), it is the principal block of the precision on the bonds of $Z_1$. Since the bonds of $Z_1$ in question are bonds of $s$, transitivity of compression (Lemma~\ref{4.24:lem-principal-blocks}) gives
\begin{equation*}
(\mathcal{A}_{ss})_{Z_1Z_1}=P_{Z_1}P_s\mathcal{A}P_sP_{Z_1}=\mathcal{A}_{Z_1Z_1}.
\end{equation*}
This is the full stage's own matrix on those bonds.
\item[(c2)] In the case that this object is the marginal covariance, it is the restriction $(\mathcal{C}_s)_{Z_1Z_1}$ of the covariance or of its decoupled family. It is never used through its value. It enters only through the covariance bounds and the walk decoupling proved below, and through the further bound on $\mathcal{C}_s$ established after the present lemma in this section. All of these hold for $\mathcal{C}_s$ with the same numbers.
\end{itemize}
This proves (i) and the first identity in \eqref{4.24:eq-gaussian-inputs-a}.

\emph{Absorption.} In the Gaussian integration proposition at the full stage, each bond $b$ of $P$ costs the factor
\begin{equation*}
e^{-\frac12\gamma_2T_j^2+\frac12\gamma_2|\mathsf{B}(b)|^2},
\end{equation*}
where $\mathsf{B}$ is the collar variable.
The second factor is absorbed into the Gaussian by means of the lower bound $\|\mathcal{C}\|^{-1}$ on the Hermitian part of the precision.

We first show that $P\subset s$. The characteristic functions of the arrested components $W\in\mathfrak{W}$ depend on $\phi_W$ alone \cite{B-CMP122a}. By \cite[(1.59)]{B-CMP122a}, they do not depend at all on the background fields. They are factors of $\mathbf{T}'_k(W)$ and stand outside the $\phi_s$-integral. Hence only bonds of $s$ carry such a factor inside the $\phi_s$-integral, and $P\subset s$.

Consequently the absorbed quadratic form acts on $\mathcal{H}_s$: on the bonds of $P$ the variable $\mathsf{B}(b)$ is the integration variable $\phi(b)$, and for $\varphi=\phi_s\in\mathcal{H}_s$,
\begin{equation*}
\tfrac12\gamma_2\sum_{b\in P}|\varphi(b)|^2=\tfrac12\langle\varphi,\gamma_2P_P\varphi\rangle.
\end{equation*}
Lemma~\ref{4.24:lem-block-coercivity} gives
\begin{equation*}
\operatorname{Re}\langle\varphi,\mathcal{A}\varphi\rangle\ge a_{\mathbb{C}}\|\varphi\|^2\quad\text{for all }\varphi\in\mathcal{H}.
\end{equation*}
This inequality holds pointwise in $\varphi$ and $\mathcal{H}_s\subset\mathcal{H}$. By the Rayleigh identity \eqref{4.24:eq-principal-blocks-a}, $\langle\varphi,\mathcal{A}_{ss}\varphi\rangle=\langle\varphi,\mathcal{A}\varphi\rangle$ for $\varphi\in\mathcal{H}_s$. Hence the Hermitian-part bound that the absorption needs holds on $\mathcal{H}_s$ with the same $a_{\mathbb{C}}$. This is the second relation in \eqref{4.24:eq-gaussian-inputs-a}. The absorption therefore goes through with the same $\gamma_2$, the same $T_j$ and the same margin $\alpha_5\|\mathcal{C}\|<\tfrac12$ of \cite[(2.31)]{B-CMP116}. This proves (ii).

\emph{Covariance bounds.} By \cite[Theorem~3.15]{B-CMP99b}, used at the region $\Lambda'$ as stated in the Gaussian integration proposition at the full stage, $\|\mathcal{C}\|\le 2B_0$, with $B_0$ and $\delta_0$ depending on $(d,L)$ only. The region $\Lambda'_s$ is a union of components of that region. By Lemma~\ref{4.24:prf-block-conditioned}, $\mathcal{C}_s$ is the conditional covariance that \cite[Theorem~3.15]{B-CMP99b} treats at $\Lambda'_s$. The same use at $\Lambda'_s$ therefore gives $\|\mathcal{C}_s\|\le 2B_0$.

Independently, the coercivity on $\mathcal{H}_s$ just proved and the Lax--Milgram bound in \eqref{4.24:eq-principal-blocks-a} give
\begin{equation*}
\|\mathcal{C}_s\|\le a_{\mathbb{C}}^{-1}=4B_0.
\end{equation*}
Using this second bound instead of the first changes one of the ceilings on $\gamma$ by a factor $2$ and adds no cap in the sense of Definition~\ref{2.3:def-cap}.

The constant $\alpha_5$ enters the present argument in only two ways. The first is the Hermitian-part lower bound of Lemma~\ref{4.24:lem-block-coercivity}, together with the ceiling on $\gamma$ stated there. The second is by statement, as a displayed hypothesis of the tolerance bound on the element parameter domain $\mathfrak{N}(c)$.

The lower bound is a gauge-invariant property of the pair. Let $u$ be $G$-valued. Then
\begin{equation*}
\mathcal{A}(\mathbb{U}^u,\mathbb{J}^u)=\mathcal{U}_u\mathcal{A}\,\mathcal{U}_u^{-1}
\end{equation*}
by the gauge covariance of the stage operator $\mathcal{A}$, and $\mathcal{U}_u$ is unitary. The block $\mathcal{A}_{ss}$ has the same covariance by Lemma~\ref{4.24:lem-principal-blocks}. For $\varphi\in\mathcal{H}$ put $\psi=\mathcal{U}_u^{-1}\varphi$. Then
\begin{equation*}
\langle\varphi,\mathcal{U}_u\mathcal{A}\,\mathcal{U}_u^{-1}\varphi\rangle=\langle\psi,\mathcal{A}\psi\rangle
\quad\text{and}\quad \|\psi\|=\|\varphi\|.
\end{equation*}
Hence the condition
\begin{equation*}
\operatorname{Re}\langle\varphi,\mathcal{A}\varphi\rangle\ge a\|\varphi\|^2\quad\text{for all }\varphi\in\mathcal{H}
\end{equation*}
holds in one gauge if and only if it holds in every gauge. By the Rayleigh identity \eqref{4.24:eq-principal-blocks-a}, it restricts to $\mathcal{H}_s$ with the same constant, whatever the derivation of $\alpha_5$. Thus $\alpha_5$, which enters the constant $c_F$ of the full-stage bounds as one of its ingredients, is unchanged at the block.

How $\alpha_5$ depends on the background is taken by statement from \cite[(2.31)]{B-CMP116}, from the Gaussian integration proposition at the full stage and from the tolerance bound, as Lemma~\ref{4.24:lem-block-coercivity}(b) records it. It is not re-derived here. In particular it is not obtained here from the global norm $\|\mathcal{A}(\mathbb{U})-\mathcal{A}(\mathbf{1})\|$. That norm depends on the gauge and is not small on the space of pairs: a pure-gauge pair changes it at zeroth order while leaving every gauge-invariant functional unchanged.

As an alternative to the appeal to \cite[Theorem~3.15]{B-CMP99b}, one may use Remark~\ref{4.24:rem-covariance-decay}, whose proof is outstanding. With $\mu_{\mathbb{C}}>0$ the decay rate fixed there, it gives
\begin{equation*}
|\mathcal{C}_s(b,b')|\le 2a_{\mathbb{C}}^{-1}e^{-\mu_{\mathbb{C}}d(b,b')},
\end{equation*}
with constants independent of $s$. This suffices for every use except the walk expansion itself, which is taken from \cite[Theorem~3.15]{B-CMP99b} at either region. This proves (iii) and the third relation in \eqref{4.24:eq-gaussian-inputs-a}.

\emph{Whitening.} By coercivity on $\mathcal{H}_s$, $\mathcal{A}_{ss}\in\mathfrak{P}(a_{\mathbb{C}})$. Lemma~\ref{4.24:lem-covariance-root} at $B=\mathcal{A}_{ss}$ then gives three facts:
\begin{itemize}
\item $S(\mathcal{A}_{ss})=\mathcal{A}_{ss}^{-1/2}=\mathcal{C}_s^{1/2}$ exists;
\item $\|\mathcal{C}_s^{1/2}\|\le a_{\mathbb{C}}^{-1/2}$;
\item $\mathcal{C}_s^{1/2}$ is holomorphic in the pair, since $a_{\mathbb{C}}$ is uniform over the pairs considered.
\end{itemize}
Write $\operatorname{Num}(B)=\{\langle\varphi,B\varphi\rangle:\|\varphi\|=1\}$ for the numerical range. For a unit vector $\varphi\in\mathcal{H}_s$, the Rayleigh identity gives $\langle\varphi,\mathcal{A}_{ss}\varphi\rangle=\langle\varphi,\mathcal{A}\varphi\rangle$. Coercivity and $|\langle\varphi,\mathcal{A}\varphi\rangle|\le\|\mathcal{A}\|$ then give
\begin{equation*}
\operatorname{Num}(\mathcal{A}_{ss})\subset\operatorname{Num}(\mathcal{A})
\subset\{\zeta:\operatorname{Re}\zeta\ge a_{\mathbb{C}}\}\cap\{\zeta:|\zeta|\le\|\mathcal{A}\|\}.
\end{equation*}
So $\mathcal{A}_{ss}$ lies in the same sector as $\mathcal{A}$. The substitution through $\mathcal{C}_s^{1/2}$ in \cite[(2.7)]{B-CMP116}, as used by the Gaussian integration proposition at the full stage, consists of algebra and Fubini's theorem. This proves (iv) and the last relation in \eqref{4.24:eq-gaussian-inputs-a}.

\emph{Decoupling of walks.} We follow \cite[(2.8)]{B-CMP116}, where the term of the walk expansion belonging to a walk $\omega$ is multiplied by the product of the decoupling parameters of the cubes the walk meets. The same decoupling is used by the Gaussian integration proposition at the full stage, and by the tolerance bound on $\mathfrak{N}(c)$ through its conditions on the decoupling parameters.

The objects decoupled are the resolvents $(\mathcal{A}_{ss}+\xi)^{-1}$ inside the integral defining $S(\mathcal{A}_{ss})$ in Lemma~\ref{4.24:lem-covariance-root}. Each of them carries a generalised random walk expansion of the type of \cite[Theorem~3.15]{B-CMP99b} at $\Lambda'_s$. The shift $\xi\ge0$ only enlarges the coercivity. This is the same statement of \cite{B-CMP99b} that the Gaussian integration proposition at the full stage uses at $\Lambda'$, in the form \eqref{4.24:eq-region-walk-expansion-a} of Lemma~\ref{4.24:prf-region-walk-expansion}. It gives majorants $m(\omega)$ of the same form, with the same $(d,L)$-constants.

Each walk is multiplied by $\prod_\Delta\mathsf{s}(\Delta)$, the product over the cells $\Delta$ of the nested partitions that the walk meets, wherever it goes. These cells are:
\begin{itemize}
\item $\pi_{j+1}$-cells on the collar region;
\item $\pi_m$-cells on the region $\Sigma_m$;
\item the collar cells of the arrested components, included as background cubes.
\end{itemize}
The object cubes entering the weight $\mu$ are the cells of $\mathsf{Q}^{s}$. By the stencil of $C$ (Definition~\ref{4.24:def-stage-form}), the block of $C^*\Delta^{(j)}C$ on the survivors' bonds, formed from the restriction of $C$ to $\mathcal{H}_s$ and the block of $\Delta^{(j)}$ on the bonds this restriction reaches, equals $\mathcal{A}_{ss}$.

This decoupling, of the walk expansion of the covariance term by term, is the only one formed in the survivors' integration of the present stage. No kernel of a determinant is decoupled: the normalisation is the density $\mathbb{C}_G$ of the normalisation lemma of the stage, which stands outside $\langle\cdot\rangle_s$. Hence the hypothesis of the cluster-expansion proposition concerning the decoupling of determinant kernels has no object here.

The conclusion holds whichever connectedness the region $\Lambda'$ of \cite[(3.185)]{B-CMP99b} is required to have.
\begin{itemize}
\item If $\Lambda'$ may be disconnected, the appeal to \cite[Theorem~3.15]{B-CMP99b} at $\Lambda'_s$ is the appeal at $\Lambda'$ verbatim.
\item If $\Lambda'$ must be connected, the covariance $\mathcal{C}$ of the full stage is a direct sum over components \cite{B-CMP119}. Then $\mathcal{A}$ is block-diagonal with respect to components, and
\begin{equation*}
\mathcal{A}_{ss}=\bigoplus_{c\subset Z^{s}}\mathcal{A}_{cc},\qquad
\mathcal{C}_s=\bigoplus_{c\subset Z^{s}}C^{(j)}(\rho(c)),
\end{equation*}
where, for each component $c$ of $Z^{s}$, $C^{(j)}(\rho(c))$ is the covariance of \cite[Theorem~3.15]{B-CMP99b} on the region $\rho(c)$ of $c$.
\end{itemize}
Either case yields the same expansion with the same constants. This proves (v).
\end{proof}

\begin{lemma}[Decoupling, reference numbers and normalisation ratios]
\label{4.24:prf-normalisation-ratios}
Work in the setting of the Gaussian inputs at the block,
\eqref{4.24:eq-gaussian-inputs-a}:
$\mathfrak{Bo}=s\sqcup\mathcal{W}$ is the set of integration bonds of the stage, $s$ the
survivors' bonds, $\mathcal{A}$ the symmetrised operator of the stage form, $\mathcal{A}_{ss}$
its principal block on $s$, $\mathcal{C}_s=\mathcal{A}_{ss}^{-1}$ the survivors' covariance,
$\mathsf{s}(\Delta)$ the decoupling parameters of the walks of $\mathcal{C}_s$, and
$a_{\mathbb{C}}=(4B_0)^{-1}$. Let $\delta_1$ be the constant of the walk count of
\cite{B-CMP116}, let $\kappa_1\ge1$, and let the auxiliary parameters satisfy
\begin{equation}\label{4.24:eq-normalisation-ratios-1}
\delta_1M\ge16\kappa_1 .
\end{equation}
Then the following hold.
\begin{enumerate}
\item[(i)] \emph{Decoupling.} The decoupled family $\mathcal{C}_s(\mathsf{s})$,
$\mathcal{C}_s^{1/2}(\mathsf{s})$ is holomorphic on the closed polydisc
$\{|\mathsf{s}(\Delta)|\le e^{\kappa_1}\}$, with termwise majorants independent of the parameter
$\xi$ carried by the walk terms $t^{\xi}_{\omega}$ of \eqref{4.24:eq-region-walk-expansion-a},
and each cube extracted from it
contributes the factor $e^{-(\kappa_1-1)}$ to the bound. If $\mathsf{s}=0$ off a cube
set $Z$, the decoupled covariance is block-diagonal over the components of $Z$, and the whitened
integrand factorises against the product measure of the whitened variable $\phi'$.
\item[(ii)] \emph{Reference numbers.} The per-cube and per-bond reference numbers of the
survivors' integration, namely
\begin{itemize}
\item the cost, at most $\mathfrak{l}_j^2$, of a large cube in meeting the conditions
\cite[(2.31)--(2.34)]{B-CMP116}, $\mathfrak{l}_j$ the large-cube weight;
\item the factor $e^{-\frac12\gamma_2T_j^2}$ per bond of the bond set $P$ of a term
$\mathcal{T}(D,P,Z)$ of the Mayer expansion;
\item the count of $2d(LM)^4$ bonds per cell, $d=4$ the dimension;
\item the weight $\tilde{\mathfrak{l}}$;
\item the constant $c_F$, equal to the constant $\alpha_*$ of the Gaussian normalisation-ratio
bound times twice the number of bonds per cube;
\item the per-bond reference bounds;
\item the condition $\gamma_2T_j^2\ge320\,\kappa(LM)^4$,
\end{itemize}
are each attached to one cell or one bond of $\mathsf{Q}^{s}$, and do not depend on the region.
\item[(iii)] \emph{Normalisation ratios.} For every term $\mathcal{T}(D,P,Z)$ of the Mayer
expansion of the survivors' integration, the constants $\alpha_*$ and $c_F$ of the Gaussian
normalisation-ratio bound, applied to the ratios formed for that term, are the same as for the
full stage.
\item[(iv)] \emph{The complex pair.} $\mathcal{A}_{ss}$, $\mathcal{C}_s$, their walks and the
elements of the expansion are all evaluated at the same pair
$(\mathbb{U},\mathbb{J})\in\mathfrak{D}^{\flat}_{n+1}(X)$; the additional condition attached to
the ambient domain on $Z_W^{+}$ enters no estimate.
\item[(v)] \emph{Classification.} Every Gaussian-dependent hypothesis displayed in the inputs of
Lemma~\ref{4.24:lem-block-coercivity}, in the Gaussian normalisation-ratio bound, in the
per-bond reference bounds and in the stage estimates belongs to one of the following five
classes:
\begin{enumerate}
\item[(a)] a coercivity constant for the Hermitian part, pointwise in $\phi\in\mathcal{H}$;
\item[(b)] upper bounds on covariance objects, read off majorants;
\item[(c)] holomorphy and termwise bounds of the decoupled walk family on the closed polydisc
$\{|\mathsf{s}(\Delta)|\le e^{\kappa_1}\}$, with the factor $e^{-(\kappa_1-1)}$ per extracted
cube;
\item[(d)] per-cube and per-bond reference numbers;
\item[(e)] normalisation ratios formed through the sector.
\end{enumerate}
Each such hypothesis holds for the principal block, or on the sub-region, with the
same number as for the full stage.
\end{enumerate}
\end{lemma}

\begin{proof}
Here $\delta_0$ and $B_0$ are Ba{\l}aban's covariance constants; for a walk $\omega$,
$|\omega|$ is its length and
$n(\omega)$ the number of cubes it meets; $\alpha_*$ and $c_F$ are the constants of the Gaussian
normalisation-ratio bound; $\phi'$ is the whitened variable of the survivors' integration, the
variable with respect to which $\mathcal{C}_s^{1/2}(\mathsf{s})$ whitens the integrand.

\emph{Proof of (i).} The estimate of the decoupled walk family on the polydisc is the counting
argument of \cite{B-CMP116}. On the polydisc $\{|\mathsf{s}(\Delta)|\le e^{\kappa_1}\}$ the term
of a walk $\omega$ is multiplied by at most $e^{\kappa_1n(\omega)}$. If $n(\omega)>2^4$, then
\begin{equation}\label{4.24:eq-normalisation-ratios-2}
\delta_0\,|\omega|\ \ge\ \delta_1\,n(\omega)\,M ,
\end{equation}
which \cite{B-CMP116} uses under $\delta_1M\ge\kappa_1$. Under
\eqref{4.24:eq-normalisation-ratios-1} this factor is compensated by one sixteenth of the
walk's decay rate. Indeed,
\eqref{4.24:eq-normalisation-ratios-2} and \eqref{4.24:eq-normalisation-ratios-1} give
\begin{equation*}
\kappa_1\,n(\omega)\ \le\ \tfrac{1}{16}\,\delta_1M\,n(\omega)\ \le\
\tfrac{1}{16}\,\delta_0\,|\omega|.
\end{equation*}
If $n(\omega)\le2^4$, the factor is at most the constant $e^{16\kappa_1}$.

The count is made walk by walk. It depends only on the walk and on the number of cubes the walk
meets, never on the region on which the family lives. It therefore applies verbatim to an
arbitrary
subfamily of walks. In particular it applies to the walk family of $\mathcal{C}_s(\mathsf{s})$
and of $\mathcal{C}_s^{1/2}(\mathsf{s})$, whose walks are among those of the walk expansion on
unions of whole components, \eqref{4.24:eq-region-walk-expansion-a}. The decoupled family is
thus a sum of holomorphic terms with majorants independent of $\xi$ that are summable uniformly
on the closed polydisc $\{|\mathsf{s}(\Delta)|\le e^{\kappa_1}\}$. Hence the family is
holomorphic there, with those majorants termwise.

Apply the Cauchy bound per extracted cube (Lemma~\ref{4.24:lem-cube-extraction-bound},
\eqref{4.24:eq-cube-extraction-bound-a}) with $R=e^{\kappa_1}$. That lemma is stated for a
function holomorphic near the closed polydisc $\{|\mathsf{s}(\Delta)|\le R\}$ and bounded there
by $G$. Its proof uses only the Cauchy circles $|\zeta-\mathsf{s}(\Delta)|=R-1$,
$\mathsf{s}(\Delta)\in[0,1]$, which lie in $\{|\zeta|\le R\}$; so its conclusion holds for a
function holomorphic in the open polydisc and continuous on the closed one, bounded there by
$G$, and both properties were just obtained. Since $\kappa_1\ge1$, that lemma bounds the
contribution of each extracted cube by
$e^{-(\kappa_1-1)}$.

For the product structure, consider the terms at $\mathsf{s}=0$ off $Z$. Termwise, a walk
joining two distinct components of $Z$ meets a cube off $Z$, where $\mathsf{s}=0$, so its term
vanishes; this is the statement of \cite{B-CMP116} that the decoupled kernels vanish unless both
of their arguments lie in the interior of one component of the region on which the decoupling
parameters are not set to zero. Hence at
$\mathsf{s}=0$ off $Z$ the decoupled covariance has no matrix elements between different
components of $Z$: it is block-diagonal over the components of $Z$. The same holds for
$\mathcal{C}_s^{1/2}(\mathsf{s})$. The whitened integrand, which depends on $\phi'$ through these
block-diagonal operators, is therefore a product over the components of $Z$. This product is
taken against the product measure of $\phi'$.

\emph{Proof of (ii).} The numbers in question are those listed in part (ii) of the statement.
Each of these is attached to one cell or to one bond of $\mathsf{Q}^{s}$ and, by inspection,
involves no datum of the region. Hence each takes the same value on every cell or bond of
$\mathsf{Q}^{s}$ as for the full stage.

\emph{Proof of (iii).} A term $\mathcal{T}(D,P,Z)$ lives on cubes $Z_0\subset\mathsf{Q}^{s}$ and
on $Z\setminus Z_0$. The latter is contained in the union of two families of cubes:
\begin{itemize}
\item the cubes of the decoupling partition met by the walks of $\mathcal{C}_s$, including the
walks that enter the cubes of the arrested component $W$, as in the walk-family item of
\eqref{4.24:eq-gaussian-inputs-a};
\item the cubes of the proper scale in the survivors' region $Z^{s}$.
\end{itemize}
Any Gaussian normalisation ratio formed for this term is a function of objects of two kinds.

The first kind consists of the principal blocks of $\mathcal{A}$ on $(\text{bonds of }Z)\cap s$
and on the bonds of $Z_0$. Since $P_E$ is the coordinate projection onto the bonds of $E$,
\begin{equation*}
P_ZP_s=P_{Z\cap s}.
\end{equation*}
By the compression identities of Lemma~\ref{4.24:lem-principal-blocks}
($\mathcal{A}_{ss}(b,b')=\mathcal{A}(b,b')$ for $b,b'\in s$, and compression is transitive),
these blocks are identical to the full stage's matrices on the same bonds.

The second kind consists of the restrictions $\mathcal{C}_s(\mathsf{s})|_{Z}$.
\begin{itemize}
\item Their numerical range and norm lie in the same sector as for the full stage. This follows
from the Rayleigh identity \eqref{4.24:eq-principal-blocks-a} of
Lemma~\ref{4.24:lem-principal-blocks}, together with the coercivity of the survivors' block,
\eqref{4.24:eq-block-coercivity-a} of Lemma~\ref{4.24:lem-block-coercivity}, with the same
$a_{\mathbb{C}}$.
\item Their walk family has the same majorants. This follows from the walk-family bounds of
\eqref{4.24:eq-gaussian-inputs-a} and from part (i).
\end{itemize}

The constants $\alpha_*$ and $c_F$ are computed from these sector data and from the numbers of
part (ii), and these are the same as for the full stage. Hence $\alpha_*$ and $c_F$ are
unchanged.

The hypotheses displayed for the Gaussian normalisation-ratio bound are the following:
\begin{itemize}
\item the parameter $a$ there equal to $\gamma_2$;
\item the smallness $\varepsilon$ obtained from the per-bond reference bounds;
\item Cauchy estimates in the decoupling parameter $\sigma$ of that bound;
\item membership of the pair in the space of pairs $\mathfrak{N}^{(n+1)}$ of the stage, which
satisfies $\mathfrak{N}^{(n+1)}\subset\mathfrak{K}_P$, where $\mathfrak{K}_P$ denotes the
space of complex background pairs of the stage, of radius $\gamma_0^{(k)}$ about the
background structure $\mathbb{B}^{(n)}_k$;
\item the smallness conditions on the pair imposed there, and the smallness of $\alpha_5$.
\end{itemize}
Each of these is of one of the classes (a)--(e) of part (v) of the statement.

\emph{Proof of (iv).} The operators $\mathcal{A}_{ss}$ and $\mathcal{C}_s$, their walks and the
elements of the expansion are all evaluated at the same pair $(\mathbb{U},\mathbb{J})\in
\mathfrak{D}^{\flat}_{n+1}(X)$ on the cubes of $Z^{s}$. Through hulls and domains reaching
$Z_W^{+}$, they are evaluated at it also on $Z_W^{+}$. There the pair is taken in the ambient
domain $\mathfrak{D}^{\flat}_{n+1}[(\mathrm{amb})_W](X)$, under the additional condition on
$Z_W^{+}$ imposed where that ambient domain is introduced. That condition restricts only the
verification that the pair belongs to the domain: the terms remain local in a region, and only
the verification of membership involves that region's enlargement. No estimate depends on it.

\emph{Proof of (v).} By inspection of the statements on which the survivors' integration depends,
the Gaussian-dependent hypotheses displayed in them are exhausted by the items of
\eqref{4.24:eq-gaussian-inputs-a} and by parts (i)--(iv). These statements are the inputs of
Lemma~\ref{4.24:lem-block-coercivity}, the Gaussian normalisation-ratio bound, the per-bond
reference bounds and the stage estimates. The hypotheses fall into the five classes of part (v)
of the statement, as follows.
\begin{enumerate}
\item[(a)] A coercivity constant for the Hermitian part, pointwise in $\phi\in\mathcal{H}$.
It holds for $\phi\in\mathcal{H}_s$, that is for the
principal block, by the Rayleigh identity \eqref{4.24:eq-principal-blocks-a}, with the same
$a_{\mathbb{C}}$ and $\gamma_2$.
\item[(b)] Upper bounds on covariance objects, read off majorants:
$\|\mathcal{C}_s\|\le a_{\mathbb{C}}^{-1}$; the kernel bound $B_0e^{-\delta_0|y-y'|}$ for the
covariance
\begin{equation*}
\mathcal{C}_s=\mathcal{A}_{ss}^{-1}=\mathcal{A}(\Lambda'_s)^{-1},
\end{equation*}
where $\mathcal{A}(\Lambda'_s)$ is the compressed precision of
\eqref{4.24:eq-region-walk-expansion-a} on the survivors' integration region $\Lambda'_s$; or
the block bound for the survivors' covariance, an upper bound on $\mathcal{C}_s$ stated in this
section. The first follows from the coercivity
\eqref{4.24:eq-block-coercivity-a} of $\mathcal{A}_{ss}$ with constant $a_{\mathbb{C}}$, by the
Lax--Milgram bound of \eqref{4.24:eq-principal-blocks-a}; the second is Ba{\l}aban's
exponential-decay bound for the covariance of the fluctuation integration, applied with the
region $\Lambda'_s$ in place of the full integration region; the third is used as that block
bound states it. Restrictions $\mathcal{C}_s(\mathsf{s})|_{Z}$ then inherit the norm and
kernel bounds by the norm inequality $\|K_{EE'}\|\le\|K\|$ of
Lemma~\ref{4.24:lem-principal-blocks} and by the termwise majorants of part (i).
\item[(c)] Holomorphy and termwise bounds of the decoupled walk family on
$|\mathsf{s}|\le e^{\kappa_1}$, with the factor $e^{-(\kappa_1-1)}$ per cube. These are given by
the count of \cite{B-CMP116} and Lemma~\ref{4.24:lem-cube-extraction-bound}, which are
independent of the region by part (i).
\item[(d)] Per-cube and per-bond reference numbers, as in part (ii).
\item[(e)] Normalisation ratios formed through the sector. Principal blocks on subsets of
$\mathsf{Q}^{s}$ are the full stage's matrices, and the restrictions of the covariance obey the
same sector and majorants, as in part (iii).
\end{enumerate}
Each hypothesis of these classes holds for the principal block, or on the sub-region, with its
number, by the part of the proof cited beside it.

The hypotheses that would not pass to a principal block are of three kinds. The first is an
identity over the whole bond set, such as a row-sum or Ward-type identity, which is false on a
sub-block. The second is a lower bound on off-diagonal structure across a cut. The third is the
value of $\log\det$, or of $(\mathcal{A}^{-1})_{ZZ}$; for the latter see the Schur witness in
\eqref{4.24:eq-principal-blocks-a}. None of these is displayed among the Gaussian-dependent
hypotheses of the statements listed above. Every displayed Gaussian input is either an
inequality of one of the classes (a)--(e), or an identity valid for an arbitrary finite index
set; such an identity holds on the bonds of $s$ as on any other finite set.

The restriction claim rests on
(i)--(iv) and the items of \eqref{4.24:eq-gaussian-inputs-a}, under the hypotheses of
Lemma~\ref{4.24:lem-block-coercivity} and of the stage estimates.
\end{proof}

The stage estimates are formulated at every stage $l$ on the intersection of the stage's region
with $Z^{[l]}$, where $Z^{[l]}$ is an arbitrary sub-collection of the components in the run at
stage $l$, and the survivors' region $Z^{s}$ is exactly the sub-collection $Z^{[n+2]}$; since
the stage estimates at stage $l$ refer to the level $\ell_l$ attached to stage $l$, this
observation is not used in the proof above.

\begin{lemma}[Each term of the exponent counted once]\label{4.24:prf-counted-once}
Work at the stage of Definition~\ref{4.24:def-survivor-elements}. Let $s$ and $\mathcal{W}$ be the survivors' and the arrested
integration bonds, with variables $\phi_s$ and $\phi_W$, and let $\mathcal{A}_{ss}$ and $\mathsf{q}_W$ be as in
Definition~\ref{4.24:def-stage-form}. Write $x$ for the real exterior variable on arrested collar bonds and $\lambda$ for the
exterior variables of the stage. Then the following hold.
\begin{enumerate}
\item[(a)] Every summand of the exponent of the stage appears exactly once: in
$-\tfrac12\langle\phi_s,\mathcal{A}_{ss}\phi_s\rangle$, in a member of $\mathfrak{E}^{\mathrm{StB}}$, in
$\sum_W\mathbb{C}^{\mathrm{fr}}_W$, or in $\mathsf{q}_W$ for one arrested component $W$.
\item[(b)] The members of $\mathfrak{E}^{\mathrm{StB}}$ are those of the survivors' list of the assignment rule, member by
member, with $\mathsf{b}(e)$ replaced by $\mathsf{b}(e)\cap\mathsf{Q}^s$ and with $\mathsf{d}(e)$, $\nu^\natural(e)$, $w(e)$
unchanged. For every member, $\tilde\nu(e)$ at the block is at most $\tilde\nu(e)$ in the whole run. A member whose mixed
difference reads no bond of $\phi_s$ vanishes.
\item[(c)] The first step of the cluster expansion for the survivors, applied to the centred Gaussian with precision
$\mathcal{A}_{ss}$, whose expectation is $\langle\cdot\rangle_s$, yields
$\sum_{D,P,Z}\mathcal{T}(D,P,Z)$. Each connected component of $Z$ contains a component of $Z_0$, and at $\mathsf{s}=0$ off $Z$
each term is a product over the connected components of $Z$. These identities hold at each fixed value of $\lambda$ and of the
complex background pair.
\end{enumerate}
\end{lemma}

\begin{proof}
\emph{Step 1: no pair is counted twice.}
Consider the collar members $e_a(b,b',Y')$ of item A(iii) of the stage form \eqref{4.24:eq-stage-form-a}. Their index set
consists of all
triples $(b,b';Y')$, together with the triples $(b,b';\emptyset)$ in which $b,b'$ lie in one block pair; these are the
collar pairs. We partition this set according to where $b$ and $b'$ lie.

\begin{itemize}
\item \emph{Case $(s,s)$.} Here $Y'$ is arbitrary. The domain $Y'$ may pass through host cubes of $\bigcup\mathfrak{W}$; there it
reads only the background and no exterior variable $\lambda$. Every $e^0_a$ with both bonds in $s$ also belongs to this class.
Summed, these terms give $-\tfrac12\langle\phi_s,\mathcal{A}_{ss}\phi_s\rangle$ on the collar. This is the survivors' Gaussian
exponent and not an element.

For these terms the change-of-variables objects they would generate have $\mathsf{b}^{P}=\emptyset$, where, in the notation of
the statement that introduces the change-of-variables objects, $S^{P}$ denotes the bond set attached to these objects and
$\mathsf{b}^{P}$ the set of cubes in which such an object reads bonds of $S^{P}$. Indeed, these terms read no bond of $S^{P}$,
and $S^{P}\subset\bigcup\mathfrak{W}$.

\item \emph{Case $(s,\mathcal{W})\cup(\mathcal{W},s)$.} These are the members of $\mathfrak{E}^{\times}$. Each is bilinear in
$(\phi_s(b),x(b'))$; the real exterior variable $x$ enters linearly and is never squared in these members. Since $b$ and $b'$
lie in distinct
components, $|Y'|\ge R_k$ by \eqref{4.24:eq-component-separation-a} of Lemma~\ref{4.24:lem-component-separation}. By
Definition~\ref{4.24:def-survivor-elements}, no $e^0_a$ is among these members.

\item \emph{Case $(\mathcal{W},\mathcal{W})$, one component.} If $b,b'$ lie in one arrested component $W$, the term belongs to
$\mathsf{q}_W$ and is kept in the exponent of $\mathbf{T}'_k(W)$.

\item \emph{Case $(\mathcal{W},\mathcal{W})$, two components.} If $b\subset W_1$ and $b'\subset W_2\neq W_1$, then
$|Y'|\ge R_k$ by \eqref{4.24:eq-component-separation-a}. The term is frozen and is assigned to the first of the two components,
by the assignment rule for members joining two arrested components.
\end{itemize}

\emph{Step 2: the remaining contributions.}
The other contributions are the elements of the localised interaction, the stored terms of the third and fourth kinds
and the $A$--$B$ coupling pieces $\mathbb{P}^{(j)}$, and the change-of-variables objects
$o\in\mathfrak{O}$. Each is assigned by one alternative:
\begin{itemize}
\item $\mathsf{b}$ contains a surviving cube, and then the contribution belongs to $\mathfrak{E}^{s}$, or, if it is a
change-of-variables object, to $\mathfrak{O}^{s}$;
\item $\mathsf{b}$ contains no surviving cube, and then the contribution is frozen and enters $\mathbb{C}^{\mathrm{fr}}_W$.
\end{itemize}
These alternatives are exclusive and exhaustive, so this is a partition.

A bilinear piece $\tfrac12\langle Q(Y)\phi,\phi\rangle$ of the fourth kind, where $Q(Y)$ is the kernel of the piece and not the
linearised block-averaging map, is
one piece with one $\mathsf{b}$. It is never split into $s\times s$, $s\times\mathcal{W}$ and $\mathcal{W}\times\mathcal{W}$
parts.

Terms reading no field variable at all, neither $\phi_s$ nor $\phi_W$, were removed beforehand, at the subtraction point of the
extension of the element list to
the stage. They are the value $A^{(n)}_k(U^{(n+1)}_k,A_j|_{\mathrm{surv}}=0)$ of the action at vanishing surviving variables.

\emph{Step 3: the rim.}
The quadratic terms on the rim that join distinct components are boundary terms $\mathbf{B}$ of Ba{\l}aban's construction
\cite{B-CMP119}. In the present notation they are parents of elements of the localised interaction whose $A$-slot carries weak
level-$j$ data. They therefore lie in $\mathfrak{E}$ already, and none of them is a second copy of a term of $\mathcal{A}$.

Steps~1--3 together prove~(a).

\emph{Step 4: the element side.}
The localisation lemma for the interaction acts on the interaction only. It is formulated with an index $c$, on which its
operator $\mathcal{A}_c$, its $u$-family $\mathcal{Q}^{w}_{c}$ and its weight $\mathfrak{w}(c)$ depend, in the notation of
that lemma. It consists of three operations:
\begin{itemize}
\item M\"obius inversion in the interpolation parameters $t_{\square}$ over the blocks of the operator $\mathcal{A}_c$ of that
lemma;
\item Cauchy's formula in $(t,u)$;
\item localisation in the decoupling parameters $\mathsf{s}$ of the walks of the response operators.
\end{itemize}
The parameters $\mathsf{s}$ multiply walks of the response operators only, and $c$ is never among the objects they multiply.
None of these operations reads the Gaussian.

The extension of the element list to the stage changes only which rim cubes carry the $u$-family
$\mathcal{Q}^{w}_{c}$ and the terms $e_{\mathrm{str}}$: they are carried by the surviving rim cubes. After this extension,
items (i)--(iii) and property (D2) of the localisation lemma for the interaction hold verbatim when they are used together
with the lemma that accompanies them in the extension of the element list to the stage, the majorants being those of a
sub-family under the same weight $\mathfrak{w}(c)$.

The collar datum map $b(\cdot)$ (the letter $b$ here is not a bond) is
\begin{equation}\label{4.24:eq-counted-once-1}
b(\mathsf{B})=g_jC\mathsf{B}-h\tilde D\bigl(g_jC\mathsf{B}\bigr),
\end{equation}
with $C$, $h$ and $\tilde D$ the operators of \cite[(1.53), (1.60)]{B-CMP122a}.

The map \eqref{4.24:eq-counted-once-1} is diagonal in the components, by the stencil of $C$ and by the location of
$h\tilde D$ at a point of the block pair determined by $c$ \cite{B-CMP109}. Hence the interpolated functional
$\Phi_{t,\sigma}(b(\mathsf{B}))$ reads $x$ only inside members whose $\mathsf{b}$ meets the collar cubes of $W$.

Consequently $\mathfrak{E}^{\mathrm{StB}}$ coincides, member by member, with the survivors' list of the assignment rule, with
the following data.
\begin{itemize}
\item $\mathsf{b}\mapsto\mathsf{b}\cap\mathsf{Q}^s$. These are the cubes where the member reads integrated variables. They are
the only datum for which the first three steps of the survivors' cluster expansion use $\mathsf{b}$, namely the sets $Y_0$,
$Z_0$ and the touching relation.
\item $\mathsf{d}$ is unchanged, and it is connected by item~(ii) of the localisation lemma for the interaction.
\item $\nu^\natural$ is unchanged, since it is a supremum over the collar disc and the slot ranges.
\item $w$ is unchanged, since it reads only $\mathsf{d}$.
\end{itemize}

Of the data entering $\tilde\nu$, therefore, only the weight $e^{c_\flat|\mathsf{b}|}$ changes. Since
$|\mathsf{b}\cap\mathsf{Q}^s|\le|\mathsf{b}|$, this weight can only decrease, so $\tilde\nu$ at the block is at most
$\tilde\nu$ in the whole run, for every member.

A member of $\mathfrak{E}^{\mathrm{StB}}$ has $\mathsf{b}\cap\mathsf{Q}^s\neq\emptyset$ by definition. Suppose its mixed
difference reads no bond of $\phi_s$. Apply item~(ii) of the localisation lemma for the interaction, which says that a mixed
difference vanishes if its $\mathsf{b}$ is empty, to the $\phi_s$-slot with $\phi_W$ held fixed. The member then vanishes. This
proves~(b).

\emph{Step 5: the identities.}
We apply the first step of the cluster expansion for the survivors to the centred Gaussian with precision $\mathcal{A}_{ss}$,
with no change. It consists of the following identities.
\begin{itemize}
\item The Mayer expansion \cite[(2.1)]{B-CMP116} over finite sets $D\subset\mathfrak{E}^{\mathrm{StB}}$.
\item The decomposition \cite[(2.3)]{B-CMP116}:
\begin{equation}\label{4.24:eq-counted-once-2}
\chi^{s}=\chi_{Y_0}\sum_{P}(-1)^{|P|}\chi^{c}_P ,
\end{equation}
where $P\subset s$ and $Y_0$ is the set of bonds of $\bigcup_{e\in D}\bigl(\mathsf{b}(e)\cap\mathsf{Q}^s\bigr)$.
\item The identities \cite[(2.5)--(2.8)]{B-CMP116}. These are conditioning on $Z_0^{c}$ inside $\mathsf{Q}^s$, whitening by
$\mathcal{C}_s^{1/2}$, and decoupling in $\mathsf{s}$ of the walks of $\mathcal{C}_s$. The conditioning, whitening and
decoupling statements used are those of the Gaussian inputs \eqref{4.24:eq-gaussian-inputs-a}.
\end{itemize}
These identities are algebra and Fubini's theorem, valid for any union of cubes $Z_0$.

The result is $\sum_{D,P,Z}\mathcal{T}(D,P,Z)$, and each connected component of $Z$ contains a component of $Z_0$. At
$\mathsf{s}=0$ off $Z$, the product statement of \eqref{4.24:eq-gaussian-inputs-a} shows that each term is a product over the
connected components of $Z$.

None of these steps uses $\lambda$ or the complex background pair except as fixed parameters. The identities therefore hold at
each fixed value of $\lambda$ and of the pair. This proves~(c).
\end{proof}

\begin{lemma}[Bound on a single term of the expansion]\label{4.24:prf-single-term}
Let $\mathcal{T}=\mathcal{T}(D,P,Z)$ be one term of the Mayer expansion of the survivors' expectation
$\langle\cdot\rangle_s$. Here $D$ is its set of elements, $P\subset s$ its set of bonds, $Z$ its
set of cubes, $Z_0\subset Z$ the cubes whose bonds form the integration bond set, and $Y_0$ the
bond set carrying the small-field characteristic function $\chi_{Y_0}$. Let the exterior variables
be real and in their boxes: the collar variables $x$ of the arrested components satisfy
$|x|\le2.4\,\mathfrak{b}'_j$, and the rim variables $A_j$ satisfy $|A_j|<\mathfrak{b}_j$. Then
\begin{equation}\label{4.24:eq-single-term-1}
|\mathcal{T}|\le\prod_{e\in D}2\nu^{\natural}(e)\cdot e^{-\frac12\gamma_2T_j^2|P|}
\cdot e^{-(\kappa_1-1)|Z\setminus Z_0|}\cdot e^{c_F|Z|}.
\end{equation}
The constant $c_F$ is the constant of the same name in the one-term bound of the polymer
expansion of the stage, and it is not changed here. Call the elements $e\in D$, the bonds of $P$
and the decoupled cubes of $Z$ the \emph{objects} $o$ of $\mathcal{T}$. Write $\operatorname{supp}o$
for the set of cubes of $Z$ attached to $o$: for an element $e$ this is $\mathsf{b}(e)\cap\mathsf{Q}^{s}$,
for a bond of $P$ it is the cube holding it, and for a decoupled cube it is the cube itself. Then
\begin{equation}\label{4.24:eq-single-term-2}
c_F|Z|\le\sum_{o}c_F|\operatorname{supp}o|,
\end{equation}
where the sum runs over the objects $o$ of $\mathcal{T}$.
\end{lemma}

\begin{proof}
\emph{The form of the term.} For every element $e\in D$, let $v_e$ be the function whose Mayer
factor $e^{v_e}-1$ appears in $\mathcal{T}$. Let $\chi^{c}_P$ be the product over $b\in P$ of the
characteristic functions complementary to the survivors' condition $|\phi(b)|<T_j$. Put
\begin{equation}\label{4.24:eq-single-term-3}
G_{\mathcal{T}}:=(-1)^{|P|}\,\chi_{Y_0}\,\chi^{c}_P\prod_{e\in D}\bigl(e^{v_e}-1\bigr).
\end{equation}
Then $\mathcal{T}$ has the form treated by the general one-term bound, that is, the bound for a
term specified by an integration bond set $I$, a bond set $E$ and an integrand $G$, which is used
here as stated in this chapter. Its integration bond set $I$ is the set of bonds of $Z_0$, its bond
set $E$ is the set of bonds of $Z$, and its integrand is $G_{\mathcal{T}}$. We now bound the
factors of $G_{\mathcal{T}}$ on the real support of $\chi^{s}$, with the exterior variables fixed,
real and in their boxes.

\emph{The Mayer factors.} Fix $e\in D$. Two properties are used.

First, $\nu^{\natural}(e)$ bounds $|v_e|$ on the whole collar disc
$|\mathsf{B}|<\varrho^{B}_j=6.6\,\mathfrak{b}'_j$ and on the slot ranges. This is the premise under
which the cross-member majorant $\nu^{\natural}$ is constructed. It is supplied by the lemma and
the theorem of this chapter that bound $|v_e|$ by $\nu^{\natural}(e)$ on the collar disc and on the
slot ranges.

Second, the real supports lie inside these ranges. The survivors' variables satisfy
$|B_j|<T_j=\mathfrak{b}'_j$. The collar variables $x$ of an arrested component $W$ satisfy
$|x|\le2.4\,\mathfrak{b}'_j$; this is the bound on the collar variables in the lemma on the
arrested components. The rim variables $A_j$, which are exterior, range in the rim box
$|A_j|<\mathfrak{b}_j$. The first two ranges lie inside the collar disc, since
$\mathfrak{b}'_j<6.6\,\mathfrak{b}'_j$ and $2.4\,\mathfrak{b}'_j<6.6\,\mathfrak{b}'_j$; the rim range
$|A_j|<\mathfrak{b}_j$ lies inside the slot ranges on which $\nu^{\natural}$ bounds $v_e$. Hence
$|v_e|\le\nu^{\natural}(e)$ on the real support of $\chi^{s}$.

Moreover $\nu^{\natural}(e)\le\frac14$. This follows from the bound
$\tilde\nu\le\mu^{\mathrm{StB}}\le\frac12$ on the element weights, proved for this expansion in the
statement on the summation of the element weights.

Since $|v_e|\le\frac14$ on this support, the exponential series gives
\begin{equation}\label{4.24:eq-single-term-4}
|e^{v_e}-1|\le\sum_{n\ge1}\frac{|v_e|^n}{n!}\le|v_e|\,e^{|v_e|}\le e^{1/4}|v_e|\le2\,\nu^{\natural}(e),
\end{equation}
because $e^{1/4}<2$. Taking the product over $e\in D$, we get
$\prod_{e\in D}|e^{v_e}-1|\le\prod_{e\in D}2\nu^{\natural}(e)$ on the real support of $\chi^{s}$.

\emph{The large-field indicators of $P$.} On the support of $\chi^{c}_P$ we have $|\phi(b)|\ge T_j$
for every $b\in P$. Hence
$-\frac12\gamma_2T_j^2+\frac12\gamma_2|\phi(b)|^2\ge0$ for every $b\in P$, and therefore
\begin{equation}\label{4.24:eq-single-term-5}
\chi^{c}_P\le\prod_{b\in P}e^{-\frac12\gamma_2T_j^2+\frac12\gamma_2|\phi(b)|^2}
=e^{-\frac12\gamma_2T_j^2|P|}\prod_{b\in P}e^{\frac12\gamma_2|\phi(b)|^2}.
\end{equation}
The second factor, $\prod_{b\in P}e^{\frac12\gamma_2|\phi(b)|^2}$, is a Gaussian factor in the
integration variables of the bonds of $P$. It is absorbed into the survivors' Gaussian measure by the
absorption bound among the Gaussian inputs of
Lemma~\ref{4.24:prf-gaussian-inputs}, see~\eqref{4.24:eq-gaussian-inputs-a}. The factors $(-1)^{|P|}$
and $\chi_{Y_0}$ have modulus at most one.

\emph{Application of the term bound.} We apply the general one-term bound to
$\mathcal{T}$ with $I$, $E$, $G_{\mathcal{T}}$ as above and with Gaussian weight parameter $a=\gamma_2$.
Its inputs are supplied as follows.
\begin{itemize}
\item The reference numbers are those of Lemma~\ref{4.24:prf-normalisation-ratios}, among the
inputs collected in~\eqref{4.24:eq-gaussian-inputs-a}.
\item The covariance inputs are the conditioning and covariance bounds of
Lemma~\ref{4.24:prf-gaussian-inputs}, see~\eqref{4.24:eq-gaussian-inputs-a}.
\item The normalisation ratios are those of Lemma~\ref{4.24:prf-normalisation-ratios}, among the
inputs collected in~\eqref{4.24:eq-gaussian-inputs-a}.
\item The extraction of the cubes of $Z\setminus Z_0$ in the decoupling parameters $\mathsf{s}$ is
priced by the Cauchy bound per extracted cube, Lemma~\ref{4.24:lem-cube-extraction-bound}, in the
form~\eqref{4.24:eq-cube-extraction-bound-a}.
\end{itemize}
The integrand $G_{\mathcal{T}}$ is bounded by the product of the bounds obtained above for its
factors. With these inputs the general one-term bound yields exactly the one-term bound of the
polymer expansion of the stage, that is~\eqref{4.24:eq-single-term-1}, the factor
$e^{-(\kappa_1-1)}$ for each cube of $Z\setminus Z_0$ coming from
Lemma~\ref{4.24:lem-cube-extraction-bound}; the constant $c_F$ is the same as in that bound.

\emph{Distribution of $c_F|Z|$ over the objects.} Every cube of $Z$ falls into one of three
classes:
\begin{itemize}
\item it lies in $\mathsf{b}(e)\cap\mathsf{Q}^{s}$ for some $e\in D$;
\item it holds a bond of $P$;
\item it is decoupled.
\end{itemize}
In each case the cube belongs to $\operatorname{supp}o$ for at least one object $o$ of $\mathcal{T}$.
Hence $Z\subset\bigcup_o\operatorname{supp}o$, so $|Z|\le\sum_o|\operatorname{supp}o|$. Multiplying
by $c_F$ gives~\eqref{4.24:eq-single-term-2}.
\end{proof}

\begin{lemma}[Anchored sum of cross weights at a survivor cell]\label{4.24:lem-cross-cell}
Work at the stage and with the survivors' cross family $\mathfrak{E}^{\times}$ of
Definition~\ref{4.24:def-survivor-elements}, with the scales, the partitions $\pi_m$, the
domains $\mathbb{D}_m$ and the hulls $[Y]_m$ of
Notation~\ref{4.24:not-scales-rim-collar}. Two host cubes are \emph{wall-adjacent} if they
share a $(d-1)$-dimensional face, and a family of host cubes is \emph{wall-connected} if it is
connected for this relation. For a site $y$ write $\Delta_k(y)$ for the $\pi_k$-cell
containing $y$. On $\pi_k$ put $\operatorname{dist}_{\pi_k}$, the graph distance of the
face-adjacency graph of $\pi_k$-cells. Let $\mathsf{c}_{\mathrm{an}}$ be the animal constant:
for every host cube $q_0$ and every $m\ge1$, at most $e^{\mathsf{c}_{\mathrm{an}}m}$
wall-connected families of $m$ host cubes contain $q_0$. Assume the convergence condition
\begin{equation*}
\kappa_1-1-c_J-\mathsf{c}_{\mathrm{an}}\ \ge\ c_\flat+7 .
\end{equation*}
\begin{enumerate}
\item[(a)] Every $\pi_t$-cell with $t\le k$, in particular every host cube and every cell of
$\mathsf{Q}$, lies in exactly one $\pi_k$-cell, written $\Delta_k(q)$ for the cell $q$; the
closed hulls nest in the same way.
\item[(b)] Wall-adjacent host cubes lie in equal or face-adjacent $\pi_k$-cells.
\item[(c)] For every wall-connected family $Y'$ of host cubes,
$[Y']_k=\bigcup\{\Delta_k(q):q\in Y'\}$, and for every site $y$ one has $y\in[Y']_k$ if and
only if $\Delta_k(y)$ contains a cube of $Y'$.
\item[(d)] A wall-chain of host cubes from a cube inside $Q\in\pi_k$ to a cube inside
$Q'\in\pi_k$ with $\operatorname{dist}_{\pi_k}(Q,Q')=D$ has at least $D+1$ members; and
for every integration bond $b$ of the stage,
$\operatorname{dist}_{\pi_k}\bigl(\Delta_k(h(b)),\Delta_k(\operatorname{cube}(b))\bigr)\le1$.
\item[(e)] For every host cube $q_0$ and every $\ell\ge1$,
\begin{equation}\label{4.24:eq-cross-cell-a}
\sum_{\substack{Y'\ \text{wall-connected},\ Y'\ni q_0\\ |Y'|\ge\ell}}
e^{-(\kappa_1-1-c_J)|Y'|}\;\le\;1.001\,e^{-7\ell}\,e^{-c_\flat\ell}\;\le\;1.001\,e^{-7\ell}.
\end{equation}
\item[(f)] For every $q\in\mathsf{Q}^{s}$,
\begin{equation}\label{4.24:eq-cross-cell-1}
\begin{split}
\sum_{\substack{e\in\mathfrak{E}^{\times}\\ \mathsf{b}(e)\cap\mathsf{Q}^{s}\ni q}}\tilde\nu'(e)
&\le 2n_bC_\rho B_{\mathcal{A}}(\varrho^{B}_j)^2e^{c_J+2c_\flat+2}\cdot1.001\,e^{-7R_k}
= 1.001\,e^{-7R_k}K_a^{(1)}(\varrho^{B}_j)^2\\
&\le \varepsilon_YK_a(\varrho^{B}_j)^2\le\vartheta_a ,
\end{split}
\end{equation}
where $K_a^{(1)}=2n_bC_\rho B_{\mathcal{A}}e^{c_J+2c_\flat+2}$. In particular
$\mu^{\times}\le\vartheta_a$.
\end{enumerate}
Here, for an integration bond $b$ of the stage, $h(b)$ is its host cube and
$\operatorname{cube}(b)$ its cell in $\mathsf{Q}$, each containing an endpoint of $b$; for
every cell $q\in\mathsf{Q}$, at most $n_b$ integration bonds $b$ of the stage have
$\operatorname{cube}(b)=q$. The constants $B_{\mathcal{A}}$, $C_\rho$, $K_a$ and the majorant
$\nu^{\natural}$ are those of \eqref{4.24:eq-stage-form-a}, $\delta$ being the decay rate and
$d(b,b')$ the bond distance in the bound on $\nu^{\natural}(e_a)$ there; $\vartheta_a$ is the
cross smallness of Definition~\ref{4.24:def-exterior-variables}; and
$\varepsilon_Y=2e^{-6R_k}$.
\end{lemma}

\begin{proof}
(a) Let $t\le k$. The $\pi_t$-cells have side $a_t=Ms_t$, and $N'=L^{k-t}$ is a positive
integer with $M=N'a_t$. The grids have a common origin, so the $\pi_t$-cell of index
$z\in\mathbb{Z}^d$ consists of the sites $y$ with $\lfloor y_\nu/a_t\rfloor=z_\nu$ for every
$\nu$. For such $y$, the identity $\lfloor x/n\rfloor=\lfloor\lfloor x\rfloor/n\rfloor$ for
real $x$ and positive integers $n$, applied with $x=y_\nu/a_t$ and $n=N'$, gives
\begin{equation*}
\Bigl\lfloor \frac{y_\nu}{M}\Bigr\rfloor
=\Bigl\lfloor\frac{\lfloor y_\nu/a_t\rfloor}{N'}\Bigr\rfloor
=\Bigl\lfloor \frac{z_\nu}{N'}\Bigr\rfloor .
\end{equation*}
The right side does not depend on $y$, so all sites of the cell have the same $\pi_k$-index,
and the cell lies in exactly one $\pi_k$-cell. Taking closures, the closed hull of the
$\pi_t$-cell lies in the closed hull of that $\pi_k$-cell. Host cubes and cells of
$\mathsf{Q}$ are cells of partitions $\pi_t$ with $t\le k$, so they are covered.

(b) Let two host cubes be wall-adjacent. By (a) each lies, with its closed hull, in the
closed hull of its $\pi_k$-cell; hence their shared $(d-1)$-dimensional face lies in the
intersection of the two closed $\pi_k$-hulls. If the $\pi_k$-indices are $w\neq w'$, this
intersection is the product over $\nu$ of the intersections of the closed intervals
$[w_\nu M,(w_\nu+1)M]$ and $[w'_\nu M,(w'_\nu+1)M]$; each factor is empty if
$|w_\nu-w'_\nu|\ge2$, a point if $|w_\nu-w'_\nu|=1$, and an interval of length $M$ if
$w_\nu=w'_\nu$. The intersection is therefore $(d-1)$-dimensional only if $w$ and $w'$
differ in exactly one coordinate, by one, that is, $|w-w'|_1=1$: the two $\pi_k$-cells are
face-adjacent.

(c) Let $Y'$ be wall-connected and put $U=\bigcup\{\Delta_k(q):q\in Y'\}$. By (b), cubes of
$Y'$ that are wall-adjacent lie in equal or face-adjacent $\pi_k$-cells, so the
wall-connectedness
of $Y'$ makes $U$ face-connected; thus $U$ is a $\mathbb{D}_k$-domain. By (a), $q\subset
\Delta_k(q)$ for every $q\in Y'$, so $U\supset\bigcup Y'$. Conversely, a
$\mathbb{D}_k$-domain containing $\bigcup Y'$ is a union of $\pi_k$-cells containing each
$q\in Y'$; since $\pi_k$ is a partition and $q\subset\Delta_k(q)$, the only $\pi_k$-cell
meeting $q$ is $\Delta_k(q)$, so the domain contains every $\Delta_k(q)$, hence $U$. Thus $U$
is the smallest $\mathbb{D}_k$-domain containing $\bigcup Y'$, that is $U=[Y']_k$. For a
site $y$, $y\in[Y']_k$ means $\Delta_k(y)=\Delta_k(q)$ for some $q\in Y'$, which by (a) is the
same as saying that $\Delta_k(y)$ contains a cube of $Y'$.

(d) In the face-adjacency graph of $\pi_k$ the graph distance is the $\ell^1$-distance of the
indices. Along a wall-chain of host cubes, consecutive members are wall-adjacent, so by (b)
the $\pi_k$-index of the member moves by an $\ell^1$-step of length at most one. To pass from
$Q$ to $Q'$ at distance $D$ it therefore needs at least $D$ steps, and the chain has at least
$D+1$ members. Next, the two endpoints of a level-$j$ bond differ in a single coordinate, by
the bond's length, which does not exceed $M$; so their $\pi_k$-indices agree except in that
coordinate, where they differ by at most one, and the $\pi_k$-cells of the two endpoints are
equal or face-adjacent. A cell containing an endpoint lies, by (a), in the $\pi_k$-cell of
that endpoint. Applied to $h(b)$ and $\operatorname{cube}(b)$, this gives
$\operatorname{dist}_{\pi_k}(\Delta_k(h(b)),\Delta_k(\operatorname{cube}(b)))\le1$.

(e) By the definition of the animal constant, the number of wall-connected families $Y'$ of
host cubes with $|Y'|=m$ containing $q_0$ is at most $e^{\mathsf{c}_{\mathrm{an}}m}$.
Put $x=e^{-(c_\flat+7)}$; since $c_\flat\ge0$ (Definition~\ref{4.24:def-cluster-weights}
gives $c_\flat\ge2\cdot3^d$), $x\le e^{-7}$. Grouping the families by size $m$, summing over
the integers $m\ge\ell$ (the term $m=\ell$ included) and using the convergence condition
assumed in the lemma,
\begin{equation*}
\begin{split}
\sum_{\substack{Y'\ni q_0\\ |Y'|\ge\ell}}e^{-(\kappa_1-1-c_J)|Y'|}
&\le\sum_{m\ge\ell}e^{\mathsf{c}_{\mathrm{an}}m}e^{-(\kappa_1-1-c_J)m}
\le\sum_{m\ge\ell}e^{-(c_\flat+7)m}\\
&\le\frac{e^{-(c_\flat+7)\ell}}{1-e^{-7}}
\le 1.001\,e^{-7\ell}e^{-c_\flat\ell}\le1.001\,e^{-7\ell}.
\end{split}
\end{equation*}
The third inequality is the geometric sum $\sum_{m\ge\ell}x^m\le x^\ell/(1-x)$ together with
$1/(1-x)\le1/(1-e^{-7})$; the fourth uses $1/(1-e^{-7})\le1.001$, which holds because
$e^{-7}<0.000912$; the last uses $c_\flat\ell\ge0$. This is \eqref{4.24:eq-cross-cell-a}.

(f) Fix $q\in\mathsf{Q}^{s}$. By Definition~\ref{4.24:def-survivor-elements}, a member
$e=e_a(b,b',Y')\in\mathfrak{E}^{\times}$ has exactly one of its bonds, the survivor bond
$b_s(e)$, in $s$ and the other, $b_{\mathfrak{W}}(e)$, in $\mathcal{W}$, and
$\mathsf{b}(e)=\{\operatorname{cube}(b_s),\operatorname{cube}(b_{\mathfrak{W}})\}$ with
$\operatorname{cube}(b_{\mathfrak{W}})\in\mathsf{Q}^{\mathfrak{W}}$. Since
$\mathsf{Q}^{s}\cap\mathsf{Q}^{\mathfrak{W}}=\emptyset$, we get
$\mathsf{b}(e)\cap\mathsf{Q}^{s}=\{\operatorname{cube}(b_s)\}$, and the members counted at $q$
are exactly those with $\operatorname{cube}(b_s(e))=q$. They are parametrised injectively by
the following four parameters: the slot (first or second) of the survivor bond in the ordered
pair $(b,b')$; the
bond $b_s$ with $\operatorname{cube}(b_s)=q$, of which there are at most $n_b$; the partner
$b_{\mathfrak{W}}$; and the wall-connected family $Y'$ containing $h(b_s)$ and
$h(b_{\mathfrak{W}})$, which by \eqref{4.24:eq-component-separation-a} has $|Y'|\ge R_k$.

For such $e$, the majorant of \eqref{4.24:eq-stage-form-a} satisfies
\begin{equation*}
\nu^{\natural}(e_a)\le\tfrac12B_{\mathcal{A}}(\varrho^{B}_j)^2
e^{-\delta d(b_s,b_{\mathfrak{W}})}e^{-(\kappa_1-1)|Y'|}.
\end{equation*}
Inserting this in the weight $\tilde\nu'(e)$ of
Definition~\ref{4.24:def-survivor-elements}, with $|\mathsf{b}(e)|\le2$ and the bound
$w(e)\le c_J|Y'|$ on the element's weight, gives
\begin{equation*}
\begin{split}
\tilde\nu'(e)
&\le 2\cdot\tfrac12B_{\mathcal{A}}(\varrho^{B}_j)^2e^{-\delta d(b_s,b_{\mathfrak{W}})}
e^{-(\kappa_1-1)|Y'|}\cdot e^{c_J|Y'|+c_J+2c_\flat+2}\\
&=B_{\mathcal{A}}(\varrho^{B}_j)^2e^{c_J+2c_\flat+2}\,e^{-\delta d(b_s,b_{\mathfrak{W}})}\,
e^{-(\kappa_1-1-c_J)|Y'|}.
\end{split}
\end{equation*}
We now sum over the four parameters, dropping the requirement $Y'\ni h(b_{\mathfrak{W}})$,
which can only enlarge the sum. The slot contributes a factor $2$; the bond $b_s$ a factor
$n_b$; the partner, summed over all bonds $b'$ of the stage, a factor
$\sum_{b'}e^{-\delta d(b_s,b')}\le C_\rho$; and the family $Y'$, wall-connected, containing
$h(b_s)$ and with $|Y'|\ge R_k$, a factor at most $1.001\,e^{-7R_k}$ by
\eqref{4.24:eq-cross-cell-a} with $q_0=h(b_s)$ and $\ell=R_k$. This gives the first
inequality in \eqref{4.24:eq-cross-cell-1}, and the equality following it is the definition
of $K_a^{(1)}$. The remaining two inequalities rest on two facts: first,
$1.001\,e^{-7R_k}\le2e^{-6R_k}=\varepsilon_Y$, which is equivalent to $e^{R_k}\ge0.5005$;
second, $\bar V\ge1$, $\bar V$ being the component budget. The second fact enters through the
expressions for $K_a$ in \eqref{4.24:eq-stage-form-a} and for $\vartheta_a$ in
Definition~\ref{4.24:def-exterior-variables}. Taking the supremum over
$q\in\mathsf{Q}^{s}$ in the definition of $\mu^{\times}$ gives $\mu^{\times}\le\vartheta_a$.

In this bound the anchor is the survivor cell $q$, all partners are summed through
$C_\rho$, and no summation over components occurs, so that the bound does not depend on the
number of arrested components; the families with $|Y'|=R_k$ exactly are included. Finally,
$\tilde\nu'\ge\tilde\nu$, and $|\mathsf{b}(e)|\le2$ holds throughout the run, so the bound
\eqref{4.24:eq-cross-cell-1} holds also with $\tilde\nu$ in place of $\tilde\nu'$.
\end{proof}

\begin{lemma}[Packing count of separated components at a point]\label{4.24:lem-packing-count}
Fix the level $k$, the stage $j$ and a lattice point $y$, with $\Delta_k(y)$ the $\pi_k$-cell
containing it, and use the objects of the anchored sum of
Lemma~\ref{4.24:lem-cross-cell} and of the stage quadratic form \eqref{4.24:eq-stage-form-a}:
the cross elements $e\in\mathfrak{E}^{\times}$ with their weights $\tilde\nu(e)$ and
$\tilde\nu'(e)$, the constants $n_b$, $C_\rho$, $B_{\mathcal{A}}$, $K_a$, $c_J$, $c_\flat$, the
field-size parameter $\varrho^B_j$ of the stage, the component budget $\bar V$, and
\begin{equation*}
K_a^{(1)}=2n_bC_\rho B_{\mathcal{A}}e^{c_J+2c_\flat+2},\qquad
\varepsilon_Y=2e^{-6R_k},\qquad
\vartheta_a=2e^{-6R_k}\bar V K_a(\varrho^B_j)^2 .
\end{equation*}
Write $\mathsf{M}^{\times}(y)$ and $\mathsf{M}'^{\times}(y)$ for the totals, over the elements
$e\in\mathfrak{E}^{\times}$ with $y\in[\mathsf{d}(e)]_k$, of $\tilde\nu(e)$ and of
$\tilde\nu'(e)$ respectively. For a survivor component $c$ put
\begin{equation*}
\mathcal{P}_c:=\{\Delta_k(q): q \text{ a collar cube of } c\},\qquad
D_c:=\operatorname{dist}_{\pi_k}\bigl(\mathcal{P}_c,\Delta_k(y)\bigr),
\end{equation*}
and group the survivor components into the shells $\mathfrak{c}_0:=\{c: D_c\le R_k\}$ and
$\mathfrak{c}_m:=\{c: mR_k<D_c\le(m+1)R_k\}$, $m\ge1$. Assume:
\begin{itemize}
\item[(R)] the condition on $R_k$ imposed in the choice of constants of the stage,
$4\mathsf{n}_{\mathrm{adj}}e^{c_J+2c_\flat}\le e^{R_k}$, where $c_\flat\ge2\cdot3^d$,
$\mathsf{n}_{\mathrm{adj}}\ge1$ and $c_J\ge0$;
\item[(S)] the separation of components, Lemma~\ref{4.24:lem-component-separation}, in the form
\eqref{4.24:eq-component-separation-a} and in the form \eqref{4.24:eq-component-separation-b};
the latter gives, for collar cubes of distinct components, an
$\ell^\infty$-separation of at least $\rho$ in $\pi_k$-units, where $\rho=R_k$ if that separation is
measured in the $\ell^\infty$-metric (\emph{reading} $\infty$) and $\rho=R_k/\sqrt d$ if it is measured
in the Euclidean metric (\emph{reading} E).
\end{itemize}
Define, for $m\ge0$,
\begin{equation}\label{4.24:eq-packing-count-a}
\begin{gathered}
N^{\mathrm{pk}}_m:=(2m+4)^d\ \ \text{(reading $\infty$)},\qquad
N^{\mathrm{pk}}_m:=\bigl((2m+3)\sqrt d+1\bigr)^d\ \ \text{(reading E)},\\
C_{\mathrm{pack}}:=N^{\mathrm{pk}}_0+N^{\mathrm{pk}}_1+\sum_{m\ge2}N^{\mathrm{pk}}_m e^{-7(m-1)} .
\end{gathered}
\end{equation}
At $d=4$ reading E gives $N^{\mathrm{pk}}_m=(4m+7)^4$. Then the following hold.
\begin{itemize}
\item[(a)] For every $m\ge0$, $\#\mathfrak{c}_m\le N^{\mathrm{pk}}_m$. The bound depends on $d$ and
$m$ only: it involves neither $|\mathfrak{W}|$ nor the number nor the sizes of the components.
\item[(b)] $\sum_c e^{-7\max\{R_k,D_c\}}\le C_{\mathrm{pack}}e^{-7R_k}$, and at $d=4$
\begin{equation*}
C_{\mathrm{pack}}\le 256+1296+4\le 2\cdot6^4\ \ \text{(reading $\infty$)},\qquad
C_{\mathrm{pack}}\le 2401+14641+47\le 1.8\cdot10^4\ \ \text{(reading E)}.
\end{equation*}
\item[(c)] Under either reading, $\mathsf{M}^{\times}(y)\le\vartheta_a$ and
$\mathsf{M}'^{\times}(y)\le\vartheta_a$.
\item[(d)] Assume in addition
\begin{itemize}
\item[($\varepsilon$)] the bound $\varepsilon_1\le4\mathsf{n}_{\mathrm{adj}}e^{-7R_k}$ on the cross
smallness $\varepsilon_1$, established in the fifth step of the change-of-variables estimate of
this section.
\end{itemize}
Then, under either reading, the total of the clause of $\mathsf{M}^{\mathrm{cv}}_*$ carrying the
condition $y\in[Y']_k$, over all arrested components $W$ and all shells, is at most
$\theta_{\mathrm{cv}}\varepsilon_Y\bar VK_a(\varrho^B_j)^2$. In particular, clause~(e) of that
estimate holds, under either reading, with the constant
\begin{equation*}
\vartheta_{\mathrm{cv}}:=\theta_{\mathrm{cv}}\varepsilon_Y
\bigl[1+2\bar V\bigl(1+K_a(\varrho^B_j)^2+n_{\mathfrak{g}}n_b\bigr)\bigr]
\end{equation*}
defined there.
\end{itemize}
No condition on the constants other than \textup{(R)} is used.
\end{lemma}

The packing numbers $N^{\mathrm{pk}}_m$ depend on $d$ and $m$ only; they are not to be confused with
the integer $N$ of $SU(N)$, nor with the number $N_0(k)$ of Ba{\l}aban's construction.

\begin{proof}
Throughout, (R) gives $e^{R_k}\ge e^{2c_\flat}$, and $c_\flat\ge2\cdot3^4=162$ at $d=4$, so
$R_k\ge2c_\flat\ge324$; in particular $R_k\ge3$ and $R_k\ge1$.

\emph{Rooting.} Root each $e\in\mathfrak{E}^{\times}$ at $q_s(e)$, the cube of the survivor bond
$b_s(e)$ of $e$; this is a collar cube of $\mathsf{Q}^{s}$ belonging to the survivor component
$c(e)$, since the stage form \eqref{4.24:eq-stage-form-a} is indexed by collar cubes. Let $Y'$ be the
domain of $e$ as in \eqref{4.24:eq-cross-cell-a}. If $y\in[Y']_k$, then by criterion~(c) of
\eqref{4.24:eq-cross-cell-a} the cell $\Delta_k(y)$ contains a cube of $Y'$. A wall-chain inside $Y'$
from the host cube $h(b_s)$ of $b_s=b_s(e)$ to that cube has, by criterion~(d) of
\eqref{4.24:eq-cross-cell-a}, at least
\begin{equation*}
\operatorname{dist}_{\pi_k}\bigl(\Delta_k(h(b_s)),\Delta_k(y)\bigr)+1\ge D_c
\end{equation*}
members, where $c=c(e)$; the inequality holds because $h(b_s)$ is the cube of $b_s$, that is, the
root cube $q_s(e)$, a collar cube of $c$, so that $\Delta_k(h(b_s))\in\mathcal{P}_c$. Together with
\eqref{4.24:eq-component-separation-a}, which gives $|Y'|\ge R_k$, this yields
\begin{equation}\label{4.24:eq-packing-count-1}
|Y'|\ge\ell_c:=\max\{R_k,D_c\}.
\end{equation}

\emph{The sum at one root cube.} Fix a collar cube $q_s$ of $c$. With the weight $\tilde\nu$, which
lacks the factor $e^{|\mathsf{b}|}$ present in $\tilde\nu'$, the elements rooted at $q_s$ and having
$y\in[\mathsf{d}(e)]_k$ satisfy, by the anchored sum of Lemma~\ref{4.24:lem-cross-cell} with the
lower bound \eqref{4.24:eq-packing-count-1} in place of $|Y'|\ge R_k$,
\begin{align*}
\sum_{\substack{e \text{ rooted at } q_s\\ y\in[\mathsf{d}(e)]_k}}\tilde\nu(e)
&\le 2n_bC_\rho B_{\mathcal{A}}(\varrho^B_j)^2e^{c_J+2c_\flat}\cdot1.001\,e^{-7\ell_c}\\
&=1.001\,e^{-2}K_a^{(1)}(\varrho^B_j)^2e^{-7\ell_c}
\le 0.14\,K_a(\varrho^B_j)^2e^{-7\max\{R_k,D_c\}} .
\end{align*}
Here $1.001$ bounds $(1-e^{-7})^{-1}$: since $e^7\ge1096$, one has
$(1-e^{-7})^{-1}\le1096/1095\le1.001$. The equality is the definition of $K_a^{(1)}$, and the last
inequality uses $K_a^{(1)}\le K_a$ together with $e^{-2}\le0.1354$, whence
$1.001\,e^{-2}\le0.1356\le0.14$. With the weight $\tilde\nu'$ in place of $\tilde\nu$ the factor
$e^{-2}$ is absent, and the same sum is at most
$1.001\,K_a(\varrho^B_j)^2e^{-7\max\{R_k,D_c\}}$. By the definition of the component budget, $c$ has
at most $\bar V$ collar cubes, so
\begin{equation}\label{4.24:eq-packing-count-2}
\mathsf{M}^{\times}(y)\le0.14\,\bar VK_a(\varrho^B_j)^2\sum_c e^{-7\max\{R_k,D_c\}},
\end{equation}
and the same holds for $\mathsf{M}'^{\times}(y)$ with $1.001$ in place of $0.14$.

\emph{Packing; proof of \textup{(a)}.} For each survivor component $c$ pick a witness collar cube
$q_c$ with $\Delta_k(q_c)\in\mathcal{P}_c$ at $\pi_k$-distance $D_c$ from $\Delta_k(y)$. Let
$c\in\mathfrak{c}_m$, $m\ge0$. Since $D_c\le(m+1)R_k$ and $R_k\ge1$, in $\pi_k$-units the cube $q_c$
lies in the cube of side at most $(2m+3)R_k$ centred at $\Delta_k(y)$. By
\eqref{4.24:eq-component-separation-b}, the witness cubes of distinct components are at
$\ell^\infty$-distance at least $\rho$, where $\rho:=R_k$ under reading $\infty$ and
$\rho:=R_k/\sqrt d$ under reading E. Hence the open $\ell^\infty$-balls of radius $\rho/2$ about the
witness cubes of the members of $\mathfrak{c}_m$ are pairwise disjoint, and all of them lie in the
cube of side $(2m+3)R_k+\rho$ centred at $\Delta_k(y)$. Each such ball has volume $\rho^d$, so
comparing volumes gives
\begin{equation*}
\#\mathfrak{c}_m\le\Bigl(\frac{(2m+3)R_k+\rho}{\rho}\Bigr)^d .
\end{equation*}
Under reading $\infty$ the right side is $(2m+4)^d$. Under reading E it is
$((2m+3)\sqrt d+1)^d$, which equals $(4m+7)^4$ at $d=4$ because $\sqrt4=2$. This is
$N^{\mathrm{pk}}_m$ of \eqref{4.24:eq-packing-count-a}. The argument uses only
\eqref{4.24:eq-component-separation-b} in the level-$k$ lattice and the volume comparison. It involves
neither $|\mathfrak{W}|$ nor the number nor the sizes of the components.

\emph{Proof of \textup{(b)}.} For $c\in\mathfrak{c}_0$ one has $\max\{R_k,D_c\}\ge R_k$. For
$c\in\mathfrak{c}_m$ with $m\ge1$ one has $D_c>mR_k$, so
$e^{-7\max\{R_k,D_c\}}\le e^{-7R_k}e^{-7(m-1)R_k}$. By (a),
\begin{equation*}
\sum_c e^{-7\max\{R_k,D_c\}}
\le e^{-7R_k}\Bigl[N^{\mathrm{pk}}_0+\sum_{m\ge1}N^{\mathrm{pk}}_m e^{-7(m-1)R_k}\Bigr]
\le C_{\mathrm{pack}}e^{-7R_k},
\end{equation*}
where the second inequality uses $R_k\ge1$. At $d=4$ the numerical bounds are as follows.

Under reading $\infty$: $N^{\mathrm{pk}}_0=4^4=256$ and $N^{\mathrm{pk}}_1=6^4=1296$. The first term
of the tail is $8^4e^{-7}\le3.74$. For $m\ge2$ the ratio of consecutive tail terms is
\begin{equation*}
\frac{N^{\mathrm{pk}}_{m+1}}{N^{\mathrm{pk}}_m}\,e^{-7}\le\frac{(5/4)^4}{1096}\le0.0023,
\end{equation*}
so $\sum_{m\ge2}\le3.75\le4$. This gives $C_{\mathrm{pack}}\le256+1296+4\le2\cdot6^4$.

Under reading E: $N^{\mathrm{pk}}_0=7^4=2401$ and $N^{\mathrm{pk}}_1=11^4=14641$. The first tail term
is $15^4e^{-7}\le46.2$, and for $m\ge2$ the ratio of consecutive tail terms is at most
$(19/15)^4/1096\le0.0024$. Hence $\sum_{m\ge2}\le46.4\le47$, and
\begin{equation*}
C_{\mathrm{pack}}\le2401+14641+47=17089\le1.8\cdot10^4 .
\end{equation*}

\emph{Absorption; proof of \textup{(c)}.} By \eqref{4.24:eq-packing-count-2} and (b),
\begin{equation*}
\mathsf{M}^{\times}(y)\le0.14\,C_{\mathrm{pack}}e^{-7R_k}\bar VK_a(\varrho^B_j)^2 .
\end{equation*}
The right side is at most $\vartheta_a=2e^{-6R_k}\bar VK_a(\varrho^B_j)^2$ if and only if
$0.07\,C_{\mathrm{pack}}\le e^{R_k}$. This holds under either reading, because
\begin{equation*}
0.07\,C_{\mathrm{pack}}\le1.3\cdot10^3\le e^{8}\le e^{2c_\flat}
\le4\mathsf{n}_{\mathrm{adj}}e^{c_J+2c_\flat}\le e^{R_k}.
\end{equation*}
The individual steps are justified as follows:
\begin{itemize}
\item $0.07\cdot17089\le1197\le1.3\cdot10^3$, and reading $\infty$ gives a smaller value;
\item $e^8\ge2980$;
\item $2c_\flat\ge324\ge8$;
\item $\mathsf{n}_{\mathrm{adj}}\ge1$ and $c_J\ge0$;
\item the last inequality is (R).
\end{itemize}
For $\mathsf{M}'^{\times}(y)$ the same argument, with $1.001$ in place of $0.14$, requires
$1.001\,C_{\mathrm{pack}}e^{-7R_k}\le2e^{-6R_k}$. This follows from
\begin{equation*}
0.5005\,C_{\mathrm{pack}}\le10^4\le e^{10}\le e^{2c_\flat}\le e^{R_k},
\end{equation*}
where $0.5005\cdot17089\le8554$, $e^{10}\ge2.2\cdot10^4$, $2c_\flat\ge324$, and the last step is
(R) with $\mathsf{n}_{\mathrm{adj}}\ge1$ and $c_J\ge0$. No condition beyond (R) enters.

Two variants are covered by the same margin. First, suppose $y$ ranges over continuum points on cell
boundaries. Then criterion~(c) of \eqref{4.24:eq-cross-cell-a} becomes: one of the at most $2^d$
closed cells at $y$ contains a cube of $Y'$. This costs a factor $16$ at $d=4$, and the factor is
absorbed by the same $e^{2c_\flat}$, since $16\cdot1.3\cdot10^3\le e^{2c_\flat}$. Second, suppose the
separation of components separates collar cubes by $R_kM/\alpha$ for a bounded
$\alpha\ge1$. Then $N^{\mathrm{pk}}_m$ changes to numbers depending on $d$, $m$ and $\alpha$ only,
and the resulting $C_{\mathrm{pack}}$ is still at most $e^{2c_\flat}/0.07$.

\emph{Proof of \textup{(d)}.} In the clause of $\mathsf{M}^{\mathrm{cv}}_*$ with $y\in[Y']_k$, the
objects are pairs $o=(e,Y)$ whose parents satisfy $e\in\{e_a,e_a^0\}$ and
$\mathsf{b}^P(e)\ne\emptyset$. Each is rooted at the collar cube of its first $S^P$-bond, which is
one of the at most $\bar V$ collar cubes of $W$.

For an arrested component $W$ put
\begin{equation*}
\mathcal{P}_W:=\{\Delta_k(q): q \text{ a collar cube of } W\},\qquad
D_W:=\operatorname{dist}_{\pi_k}\bigl(\mathcal{P}_W,\Delta_k(y)\bigr).
\end{equation*}
By \eqref{4.24:eq-cross-cell-a}, $|Y'|\ge D_W$ for $e_a$. For $e_a^0$ one has $D_W\le2<R_k$. The
shells of arrested components are defined by the same inequalities as $\mathfrak{c}_m$, with $D_W$
in place of $D_c$. The same
packing as in (a), which holds whatever $|\mathfrak{W}|$ is, shows:
\begin{itemize}
\item at most $N^{\mathrm{pk}}_0$ arrested components have $D_W\le R_k$;
\item at most $N^{\mathrm{pk}}_m$ lie in shell $m\ge1$, and these carry the additional factor
$1.001\,e^{-7mR_k}$.
\end{itemize}

The fifth step of the change-of-variables estimate bounds
\begin{equation*}
\sum_Y\tilde\nu(e,Y)\le\theta_{\mathrm{cv}}\bigl(2e^{c_J}\varepsilon_1\bigr)\tilde\nu'(e).
\end{equation*}
By ($\varepsilon$) and then (R),
\begin{align*}
2e^{c_J}\varepsilon_1\le8\mathsf{n}_{\mathrm{adj}}e^{c_J}e^{-7R_k}
&=2\bigl(4\mathsf{n}_{\mathrm{adj}}e^{c_J+2c_\flat}\bigr)e^{-2c_\flat}e^{-7R_k}\\
&\le2e^{R_k}e^{-2c_\flat}e^{-7R_k}=\varepsilon_Y\,e^{-2c_\flat}.
\end{align*}
Thus the derivation of $\varepsilon_Y$ leaves a spare factor $e^{-2c_\flat}$. Summing over the
arrested components $W$ and the shells, with the counts above and the at most $\bar V$ root cubes of
each $W$, the total of the clause is at most
\begin{equation*}
\theta_{\mathrm{cv}}\varepsilon_Ye^{-2c_\flat}\Bigl(N^{\mathrm{pk}}_0
+\sum_{m\ge1}N^{\mathrm{pk}}_m\,1.001\,e^{-7mR_k}\Bigr)\bar VK_a(\varrho^B_j)^2 .
\end{equation*}

The sum over $m\ge1$ is at most $1$ under either reading, since $R_k\ge3$. Under reading E:
\begin{itemize}
\item $11^4\cdot1.001\le14656$ and $e^{-21}\le7.6\cdot10^{-10}$, so the term $m=1$ is at most
$1.2\cdot10^{-5}$;
\item for $m\ge1$ the ratio of consecutive terms is at most
$(N^{\mathrm{pk}}_{m+1}/N^{\mathrm{pk}}_m)e^{-7R_k}\le(15/11)^4e^{-21}\le10^{-8}$;
\item hence the sum is at most $1.3\cdot10^{-5}$.
\end{itemize}
Under reading $\infty$ the sum is smaller still, because $6^4\le11^4$ and $8/6\le15/11$. The total is
therefore at most
\begin{equation*}
\theta_{\mathrm{cv}}\varepsilon_Ye^{-2c_\flat}\bigl(N^{\mathrm{pk}}_0+1\bigr)\bar VK_a(\varrho^B_j)^2
\le\theta_{\mathrm{cv}}\varepsilon_Y\bar VK_a(\varrho^B_j)^2 .
\end{equation*}
The inequality uses, under either reading,
\begin{equation*}
N^{\mathrm{pk}}_0+1\le7^4+1=2402\le e^{8}\le e^{2c_\flat},
\end{equation*}
where $e^8\ge2980$ and $2c_\flat\ge324$. This total is below what the fifth step assigns, inside
$\vartheta_{\mathrm{cv}}$, to the nearest arrested component alone. Hence clause~(e) of the
change-of-variables estimate holds, under either reading, with the constant
$\vartheta_{\mathrm{cv}}$ defined there, and so does every bound in which $\vartheta_{\mathrm{cv}}$
enters. In each absorption in (c) and in the last bound of the total in (d), the
margin is at least $e^{2c_\flat-10}\ge e^{314}$, so none of the decimal bounds decides anything.

The neighbouring clause, in which $y\in[Y]_k$ determines an arrested component $W$, is treated as
follows. Under reading $\infty$, since $R_k\ge3$, the cell $\Delta_k(y)$ meets at most one
enlargement $Z_W^{+}$, as the fifth step uses. Under reading E, the same disjoint-balls count shows
that $\Delta_k(y)$ meets at most $3^d$ of them. This number depends on $d$ only and is absorbed by the
spare factor $e^{-c_\flat}$ in the line of $\mu^{\mathrm{cv}}$ indexed by $(b,Y)$, since
$c_\flat\ge2\cdot3^d$. So the fifth step holds under either reading.
\end{proof}

\begin{lemma}[Totals uniform in the number of arrested components]\label{4.24:prf-uniform-totals}
Let
\begin{equation*}
\mathfrak{E}^{\mathrm{StB}}=\mathfrak{E}^{s}\sqcup\mathfrak{E}^{\times}\sqcup\mathfrak{O}^{s}
\end{equation*}
be the element list of Definition~\ref{4.24:def-survivor-elements}. Each member $e$ carries
$\mathsf{b}(e)\cap\mathsf{Q}^{s}$, $\mathsf{d}(e)$ and $\nu^{\natural}(e)$. Give each member its primed weight
$\tilde\nu'(e)$, in which $|\mathsf{b}|$ takes its whole-run value; that value is at most $2$. Put
\begin{equation}\label{4.24:eq-uniform-totals-1}
\mu^{\mathrm{StB}}:=\sup_{q\in\mathsf{Q}^{s}}
\sum_{\substack{e\in\mathfrak{E}^{\mathrm{StB}}\\ q\in\mathsf{b}(e)\cap\mathsf{Q}^{s}}}\tilde\nu'(e),
\qquad
\mathsf{M}^{\mathrm{StB}}_*:=\sup_{y}
\sum_{\substack{e\in\mathfrak{E}^{\mathrm{StB}}\\ y\in[\mathsf{d}(e)]_k}}\tilde\nu'(e).
\end{equation}
These totals dominate two other sets of totals:
\begin{itemize}
\item the totals formed with the unprimed weights $\tilde\nu$ in the cluster-expansion theorem applied
to $\mathfrak{E}^{\mathrm{StB}}$;
\item the primed totals entering the condition $(Q^{+})'$.
\end{itemize}
Let $\mu'$ and $\mathsf{M}'_*$ be the corresponding primed totals of the whole-run family, that is, of the
family formed with every arrested component $W$ of the run. Then, uniformly in the number
$|\mathfrak{W}|$ of arrested components,
\begin{equation}\label{4.24:eq-uniform-totals-a}
\mu^{\mathrm{StB}}\le\mu'+\vartheta_a+\vartheta_{\mathrm{cv}},
\qquad
\mathsf{M}^{\mathrm{StB}}_*\le\mathsf{M}'_*+\vartheta_a+\vartheta_{\mathrm{cv}}.
\end{equation}
Under the condition $(Q^{+})'$,
\begin{equation}\label{4.24:eq-uniform-totals-2}
\mu^{\mathrm{StB}}+\tilde{\mathfrak{l}}\le\tfrac12,
\qquad
4\bigl(\mathsf{M}^{\mathrm{StB}}_*+2^{d}V\tilde{\mathfrak{l}}\bigr)\le\tfrac12 .
\end{equation}
Under the same condition, the remaining hypotheses of the cluster-expansion theorem hold for
$\mathfrak{E}^{\mathrm{StB}}$:
\begin{itemize}
\item the hypothesis $\nu^{\natural}(e)\le\frac12$ holds for every member $e$, indeed
$\nu^{\natural}(e)\le\frac14$, with no separate ceiling imposed;
\item the total $\mu_2$ of the cluster-expansion theorem satisfies $\mu_2\le\frac12+e^{-4}$;
\item its floor hypotheses are unaffected.
\end{itemize}
\end{lemma}

\begin{proof}
Write $\mu^{s}$ and $\mathsf{M}^{s}_*$ for the suprema in \eqref{4.24:eq-uniform-totals-1} with the sums
restricted to $e\in\mathfrak{E}^{s}$. Write $\mu^{\mathrm{cv}}$ and $\mathsf{M}^{\mathrm{cv}}_*$ for the
same suprema with the sums restricted to $e\in\mathfrak{O}^{s}$. Let $\mu^{\times}$ and
$\mathsf{M}^{\times}$ be the cross totals of Definition~\ref{4.24:def-survivor-elements}, which are formed
with the primed weights. The family $\mathfrak{E}^{\mathrm{StB}}$ is the disjoint union of the three
sub-families, and all weights are non-negative. So for each fixed $q$, and for each fixed $y$, the sum
in \eqref{4.24:eq-uniform-totals-1} is the sum of the three restricted sums. The supremum of a sum is at
most the sum of the suprema, hence
\begin{equation}\label{4.24:eq-uniform-totals-3}
\mu^{\mathrm{StB}}\le\mu^{s}+\mu^{\times}+\mu^{\mathrm{cv}},
\qquad
\mathsf{M}^{\mathrm{StB}}_*\le\mathsf{M}^{s}_*+\mathsf{M}^{\times}+\mathsf{M}^{\mathrm{cv}}_* .
\end{equation}
We bound the three pieces in turn.

\emph{(a) The members of $\mathfrak{E}^{s}$.} The family $\mathfrak{E}^{s}$ is a sub-family of the
whole-run family. Member by member, its primed weights $\tilde\nu'$ are smaller than or equal to those of
the whole-run family, by the comparison of the survivors' weights with the whole-run weights. The totals
$\mu^{s}$ and $\mathsf{M}^{s}_*$ are suprema of sums of non-negative terms over this sub-family, with
terms no larger than the corresponding terms of $\mu'$ and $\mathsf{M}'_*$. Therefore
\begin{equation*}
\mu^{s}\le\mu' \qquad\text{and}\qquad \mathsf{M}^{s}_*\le\mathsf{M}'_* .
\end{equation*}
The right-hand sides are bounded by the two displays of the whole-run totals estimate, the bound on
$\mu'$ and the bound on $\mathsf{M}'_*$. These displays are used with the constant $c_\flat$ replaced by
$c_\flat+1$, as the condition $(Q^{+})'$ prescribes.

\emph{(b) The cross members $\mathfrak{E}^{\times}$.} Fix $q\in\mathsf{Q}^{s}$. By
Lemma~\ref{4.24:lem-cross-cell}, the anchored sum of the cross weights at $q$ satisfies
\begin{equation*}
\sum_{\substack{e\in\mathfrak{E}^{\times}\\ q\in\mathsf{b}(e)\cap\mathsf{Q}^{s}}}\tilde\nu'(e)
\le\varepsilon_Y K_a(\varrho^{B}_j)^2\le\vartheta_a .
\end{equation*}
That lemma rests on the packing count of Lemma~\ref{4.24:lem-packing-count}, which is proved as an
implication from the statements named there. Taking the supremum over $q$ gives
\begin{equation*}
\mu^{\times}\le\varepsilon_Y K_a(\varrho^{B}_j)^2\le\vartheta_a .
\end{equation*}
The companion bound $\mathsf{M}^{\times}\le\vartheta_a$ is obtained from the same packing count of
Lemma~\ref{4.24:lem-packing-count}. The packing count is a statement about
geometry only: given the separation of the components in the step-$k$ lattice, it reads neither
$|\mathfrak{W}|$, nor the number of the components, nor their sizes. Hence both bounds hold uniformly in
$|\mathfrak{W}|$.

\emph{(c) The change-of-variables members $\mathfrak{O}^{s}$.} Clause (e) of the change-of-variables
theorem gives
\begin{equation*}
\mu^{\mathrm{cv}}\le\vartheta_{\mathrm{cv}}
\qquad\text{and}\qquad
\mathsf{M}^{\mathrm{cv}}_*\le\vartheta_{\mathrm{cv}},
\end{equation*}
uniformly in $|\mathfrak{W}|$ and in the number of branches. The sums of that clause are anchored at
$\mathsf{b}(e)$ and count no arrested component $W$. The one point of its proof that depends on
$|\mathfrak{W}|$ is dealt with in the packing count of Lemma~\ref{4.24:lem-packing-count}. Clause (e)
therefore applies to $\mathfrak{O}^{s}$ without modification of its hypotheses.

\emph{Conclusion of the totals.} Insert (a), (b) and (c) into \eqref{4.24:eq-uniform-totals-3}. This
gives \eqref{4.24:eq-uniform-totals-a}, and every bound used is uniform in $|\mathfrak{W}|$.

\emph{The two smallness lines.} Recall that the condition $(Q^{+})'$ is the smallness condition of the
whole-run totals estimate with the constant $c_\flat$ replaced by $c_\flat+1$ and $\mathfrak{p}_j$
replaced by $\mathfrak{p}_j+\vartheta_{\mathrm{cv}}+\vartheta_a$. By \eqref{4.24:eq-uniform-totals-a}, the
totals of $\mathfrak{E}^{\mathrm{StB}}$ exceed $\mu'$ and $\mathsf{M}'_*$ by at most
$\vartheta_a+\vartheta_{\mathrm{cv}}$. The condition $(Q^{+})'$ is exactly the re-maximised condition
under which the two lines of the whole-run totals estimate give \eqref{4.24:eq-uniform-totals-2}.

\emph{The remaining hypotheses of the cluster-expansion theorem.}

First, the hypothesis $\nu^{\natural}(e)\le\frac12$ needs no separate ceiling. Every member $e$ of
$\mathfrak{E}^{\mathrm{StB}}$ has $\mathsf{b}(e)\cap\mathsf{Q}^{s}\ne\emptyset$. So there is a cell
$q\in\mathsf{b}(e)\cap\mathsf{Q}^{s}$, and $\tilde\nu'(e)$ is one of the non-negative terms of the sum
defining $\mu^{\mathrm{StB}}$ at that $q$. Since the primed weights dominate the unprimed ones,
$\tilde\nu(e)\le\tilde\nu'(e)\le\mu^{\mathrm{StB}}$. The whole-run totals estimate obtains
$\nu^{\natural}(e)\le\frac12$ from its bound on $\mu'$ through the inequality
$\nu^{\natural}(e)\le\frac12\tilde\nu(e)$. The same inequality, together with the first bound of
\eqref{4.24:eq-uniform-totals-2} and $\tilde{\mathfrak{l}}\ge0$, gives
\begin{equation*}
\nu^{\natural}(e)\le\tfrac12\tilde\nu(e)\le\tfrac12\mu^{\mathrm{StB}}\le\tfrac14 .
\end{equation*}

Second, with the constants as the condition $(Q^{+})'$ prescribes, the total $\mu_2$ of the
cluster-expansion theorem satisfies
\begin{equation*}
\mu_2\le\mu^{\mathrm{StB}}+\tilde{\mathfrak{l}}+e^{-(\kappa_1-1)+c_J+c_\flat+1}
\le\tfrac12+e^{-4}.
\end{equation*}
In the second inequality, the first bound of \eqref{4.24:eq-uniform-totals-2} is used. The lower bound
on $\kappa_1$ required by the change-of-variables theorem gives $-(\kappa_1-1)+c_J+c_\flat+1\le-4$.

Third, the floor hypotheses of the cluster-expansion theorem do not involve the totals. Passing from the
whole-run family to $\mathfrak{E}^{\mathrm{StB}}$ leaves them untouched.

Finally, for later use, the majorants of the cross members also tend to zero along the flow. Let
$e=e_a\in\mathfrak{E}^{\times}$. Then
\begin{equation}\label{4.24:eq-uniform-totals-4}
\begin{aligned}
\nu^{\natural}(e_a)
&\le\tfrac12 B_{\mathcal{A}}(\varrho^{B}_j)^2e^{-(\kappa_1-1)R_k}\\
&\le\tfrac12 B_{\mathcal{A}}(6.6A_1)^2\bigl(2\log g_k^{-2}\bigr)^{2p_1}e^{-7(\log x_k)^{r}},
\end{aligned}
\end{equation}
and the right-hand side tends to $0$ as $x_k\to\infty$. The second line uses three facts:
\begin{itemize}
\item $\kappa_1-1\ge7$, by the lower bound on $\kappa_1$ required by the change-of-variables theorem,
which is fixed before $\gamma$;
\item $R_k\ge(\log g_k^{-2})^{r}$;
\item $\log x_j\le2\log x_k$, by the coupling window of the stage.
\end{itemize}
In the last fact, the depth, written $N$ in \cite{B-CMP122a} and equal to $k-h$ here, depends on $g$.
This is used here as stated in \cite[p.~198, after (1.94)]{B-CMP122a}:
\begin{quotation}
The last inequality holds under two restrictions on $N$. At first, we assume that
$N \le O(1)(\log g_k^{-2})^{\nu}$ \dots{} The second is that $N$ has to be sufficiently large \dots{}
These two conditions can be satisfied by $N$ to the positive power of $\log g_k^{-2}$.
\end{quotation}
It is also used as stated in \cite[p.~179]{B-CMP122a}:
\begin{quotation}
Take the smallest positive integer $N_0$ such, that $L^{-N_0+1}MR_{k-N_0+1} = M$. \dots{} We assume that
$N > N_0$ \dots
\end{quotation}
With $N$ read in this way, the coupling window of the stage is the coupling window of
\cite{B-CMP122a}. It is a one-sided condition, a ceiling involving $\gamma$ that bounds the coupling $g$
from above only, and not a cap on the depth $N$.
\end{proof}

\begin{lemma}[Holomorphy, locality and the pointwise stage constant]\label{4.24:prf-holomorphy-locality}
Work in the setting of Theorem~\ref{4.24:thm-survivor-block-expansion}, with the complex background pair
$(\mathbb{U},\mathbb{J})$ of the stage, the survivors' region $Z^{s}$, the band $R_n$, and the element family
$\mathfrak{E}$ of the stage together with its weights $\tilde\nu\le\tilde\nu'$. Assume the following four
statements.
\begin{enumerate}
\item[(a)] The convergence statement for the Gaussian normalisation densities of the stage in the stage norm
$\|\cdot\|_{(n+1)}$, including the bound on its remainder term $\Delta W$ on $Z^{s}$.
\item[(b)] The holonomy bound for walks across several cubes, Proposition~\ref{4.24:prop-holonomy-bound},
whose proof is outstanding.
\item[(c)] The pointwise bound reached at the sixth step of the proof of clause (ii) of the cluster-expansion
proposition for the element families of the stage. For a point $y$ write $\mathsf{M}^{s}(y)$ for the
pointwise total at $y$ of the survivors' family $\mathfrak{E}^{s}$, that is, for the sum of its weights over
the elements $e$ with $y\in[\mathsf{d}(e)]_k$. The bound states: for every point $y$, the contribution of the
survivors' interaction densities $\mathbb{C}^{s}_I$ to $\mathbb{C}^{(n+1)}_k(y)$ is at most
\begin{equation*}
4\bigl(\mathsf{M}^{s}(y)+\vartheta_{\mathrm{cv}}+\vartheta_a+2^{d}V\tilde{\mathfrak{l}}\bigr).
\end{equation*}
\item[(d)] The conversion of per-rate bounds to the stage norm, Proposition~\ref{4.24:prop-norm-conversion},
whose proof is outstanding.
\end{enumerate}
Then the following hold.
\begin{enumerate}
\item[(1)] For every domain $X$ of the expansion of clause (ii) of
Theorem~\ref{4.24:thm-survivor-block-expansion}, each term of the Mayer--polymer--Ursell series defining the
density at $X$ is holomorphic in $(\mathbb{U},\mathbb{J})$ on
$\mathfrak{D}^{\flat}_{n+1}[(\mathrm{amb})_W](X)$, the analyticity domain of the stage densities on $X$ taken
relative to the ambient components $Z_W$ of the arrested components $W\in\mathfrak{W}$; the series converges
uniformly there and its sum is holomorphic; the sum depends on the pair only through its values in the cubes
of $X$; and $X\cap R_n\cap Z^{s}\neq\emptyset$.
\item[(2)] Clause (iii) of the cluster-expansion proposition holds for the survivors' family, with the
interpolated elements
\begin{equation*}
v_e^{\mathfrak{z}} = v_e^{0}+\mathfrak{z}\,(v_e^{1}-v_e^{0})
\end{equation*}
in the dial variable $\mathfrak{z}$ and with $\Lambda := 1/\bigl(4\mu^{\mathrm{StB}}[\nu^{m}]\bigr)$, where
$\mu^{\mathrm{StB}}$ is the total of \eqref{4.24:eq-uniform-totals-a}, here taken with the weights $\nu^{m}$.
\item[(3)] For every point $y$,
\begin{equation}\label{4.24:eq-holomorphy-locality-1}
\mathbb{C}^{(n+1)}_k(y)\le C_0^{\sharp}=3c_G+1 .
\end{equation}
Moreover $\vartheta_a/g_k\to0$ and $\vartheta_{\mathrm{cv}}/g_k\to0$ as $g_k\to0$.
\end{enumerate}
\end{lemma}

\begin{proof}
\emph{Holomorphy.} By construction, each term of the Mayer--polymer--Ursell series is a finite sum of products
of contour integrals, in the decoupling parameters $\mathsf{s}$ and the interpolation parameters $\tau$, $t$,
$u$ of the expansion, of three
kinds of factors: finitely many elements $v_e$; decoupled walk terms of the covariance $\mathcal{C}_s$; and the
Gaussian factors of the one-term representation of Lemma~\ref{4.24:prf-single-term}, which do not depend on the
exterior parameter $\lambda$. Each factor is holomorphic in the pair on
$\mathfrak{D}^{\flat}_{n+1}[(\mathrm{amb})_W](X)$. For the elements, $v_e$ is holomorphic on
$\mathfrak{D}^{\flat}_{n+1}(\mathsf{d}(e))$ by the analyticity statements for the elements of the stage, and
the definition of the domains of the stage densities restricts it to
$\mathfrak{D}^{\flat}_{n+1}[(\mathrm{amb})_W](X)$. For the Gaussian
data, $\mathcal{A}_{ss}$ and $\mathcal{C}_s(\mathsf{s})$ are holomorphic by the conditioning, absorption and
covariance bounds among the Gaussian inputs at the block
collected in \eqref{4.24:eq-gaussian-inputs-a}. A contour integral of a function holomorphic in the pair is
holomorphic in the pair, and so are finite sums and products of such integrals; hence every term is
holomorphic.

\emph{Convergence.} The series is dominated term by term by the majorants of
Lemma~\ref{4.24:prf-single-term}, of the anchoring step of the proof of clause (ii) of
Theorem~\ref{4.24:thm-survivor-block-expansion}, and of Lemma~\ref{4.24:prf-uniform-totals}. These majorants do
not depend on $(\mathbb{U},\mathbb{J})$ or on $\lambda$, and their sum is finite. The series therefore converges
uniformly on $\mathfrak{D}^{\flat}_{n+1}[(\mathrm{amb})_W](X)$, and by the Weierstrass theorem on uniform limits
of holomorphic functions its sum is holomorphic there.

\emph{Locality in $X$.} Every ingredient depends on the pair only through its values inside its own cubes: for
the walk terms of the covariance this is \cite[p.~5]{B-CMP116} and \cite[(3.107)]{B-CMP99b}, the elements $v_e$
depend on the pair on $\mathsf{d}(e)$ only, and the Gaussian factors of Lemma~\ref{4.24:prf-single-term} involve
$\mathcal{A}_{ss}$, whose bond densities depend on the pair inside their cubes by the same two references. Every
domain $X$ of the series meets $R_n\cap Z^{s}$ by the anchoring step of the proof of
clause (ii) of Theorem~\ref{4.24:thm-survivor-block-expansion}. This is clause (i) of the cluster-expansion
proposition for the survivors' family, including its locality statement, and proves (1).

\emph{The dial.} Write $v_e^{\mathfrak{z}}$ as in (2). Here $\mathfrak{z}$ is the dial variable, not the exterior
parameter $\lambda$. Put $\Lambda := 1/(4\mu^{\mathrm{StB}}[\nu^{m}])$ and apply the Schwarz lemma in
$\mathfrak{z}$ exactly as in the proof of clause (iii) of the cluster-expansion proposition. That proof uses
only totals of the element weights, and these are bounded by \eqref{4.24:eq-uniform-totals-a}. This proves (2).

\emph{The pointwise stage constant.} Fix $y$. Write $\mathsf{M}^{\mathrm{fr}}(y)$ for the pointwise total at $y$
of the frozen family $\mathfrak{E}^{\mathrm{fr}}$, defined as $\mathsf{M}^{s}(y)$ in hypothesis (c) with
$\mathfrak{E}^{\mathrm{fr}}$ in place of $\mathfrak{E}^{s}$.
The density $\mathbb{C}^{(n+1)}_k(y)$ is bounded by the sum of four contributions.
\begin{itemize}
\item[(G)] The normalisation densities $\mathbb{C}^{s}_G$ and the Jacobian densities $\mathbb{C}^{G}_W$, over all
arrested components $W$,
are counted jointly per cell. Their stage norm $\|\cdot\|_{(n+1)}$ is controlled by hypotheses (a) and (d), the
latter being the conversion of per-rate bounds to the stage norm, Proposition~\ref{4.24:prop-norm-conversion},
whose proof is
outstanding. Their normalisation constant $c_G$ is the one of \eqref{4.24:eq-constant-unchanged-a}, and they
contribute at most $3c_G$.
\item[($\vartheta$)] The remainder $\Delta W$ of hypothesis (a) on $Z^{s}$ contributes at most $\vartheta$. This
is the $\vartheta$-clause of Theorem~\ref{4.24:thm-survivor-block-expansion}(i), which rests on hypothesis (b),
Proposition~\ref{4.24:prop-holonomy-bound}, whose proof is outstanding.
\item[(s)] The survivors' interaction densities contribute at most
\begin{equation*}
4\bigl(\mathsf{M}^{s}(y)+\vartheta_{\mathrm{cv}}+\vartheta_a+2^{d}V\tilde{\mathfrak{l}}\bigr),
\end{equation*}
by hypothesis (c).
\item[(fr)] The frozen family $\mathbb{C}^{\mathrm{fr}}_W$ contributes at most
$\tfrac12\mathsf{M}^{\mathrm{fr}}(y)+\vartheta_{\mathrm{cv}}+\vartheta_a$. The Jacobian densities
$\mathbb{C}^{G}_W$ are excluded from this family: they carry no cross-member majorant $\nu^{\natural}$, and they
are already counted in (G).
\end{itemize}

The decomposition $\mathfrak{E}^{s}\sqcup\mathfrak{E}^{\mathrm{fr}}=\mathfrak{E}$ is a partition, so each element
is counted once, at the same $y$. Moreover the survivors' weights satisfy
$\tilde\nu^{s}\le\tilde\nu\le\tilde\nu'$. Hence
\begin{equation}\label{4.24:eq-holomorphy-locality-2}
4\mathsf{M}^{s}(y)+\mathsf{M}^{\mathrm{fr}}(y)
\le 4\sum_{e\in\mathfrak{E}:\,[\mathsf{d}(e)]_k\ni y}\tilde\nu'(e)
\le 4\mathsf{M}'_* .
\end{equation}
The right-hand side is bounded by the right-hand member of the whole-run bound on the primed totals, taken by
statement at the primed prescription $(Q^{+})'$ with $W$ counted among the components in the run, and quoted in
\eqref{4.24:eq-uniform-totals-a}; it gives $4(\mathsf{M}'_*+2^{d}V\tilde{\mathfrak{l}})\le\tfrac12$.

Now add the four contributions. Use $\tfrac12\mathsf{M}^{\mathrm{fr}}(y)\le\mathsf{M}^{\mathrm{fr}}(y)$ and
\eqref{4.24:eq-holomorphy-locality-2}, and collect the four copies of $\vartheta_{\mathrm{cv}}+\vartheta_a$
from (s) with the one from (fr). This gives
\begin{align*}
\mathbb{C}^{(n+1)}_k(y)
&\le 3c_G+\vartheta+4\bigl(\mathsf{M}'_*+2^{d}V\tilde{\mathfrak{l}}\bigr)
  +5(\vartheta_{\mathrm{cv}}+\vartheta_a)\\
&\le 3c_G+\vartheta+\tfrac12+5(\vartheta_{\mathrm{cv}}+\vartheta_a)\\
&\le 3c_G+\vartheta+\tfrac12+\bigl(\tfrac12-\vartheta\bigr)=3c_G+1=C_0^{\sharp}.
\end{align*}
The last inequality is clause (iv) of the change-of-variables smallness conditions of the stage,
$5(\vartheta_{\mathrm{cv}}+\vartheta_a)\le\tfrac12-\vartheta$. This proves
\eqref{4.24:eq-holomorphy-locality-1}.

\emph{The quotients $\vartheta_a/g_k$ and $\vartheta_{\mathrm{cv}}/g_k$.} Each of $\vartheta_a$ and
$\vartheta_{\mathrm{cv}}$ is bounded by a polynomial in $\log g_j^{-2}$ times $e^{-6R_k}$. The reference-window
condition gives $\log x_j\le 2\log x_k$, so each is bounded by a polynomial in $\log g_k^{-2}$ times
$e^{-6R_k}$. The
number of steps
between the creation level of the region and the level $k$ is Ba{\l}aban's $g$-dependent depth, used here
as stated in \cite[p.~198]{B-CMP122a} and \cite[p.~179]{B-CMP122a}. Hence both quotients tend to $0$ as
$g_k\to0$, and the ceiling $C_0^{\sharp}$ in \eqref{4.24:eq-holomorphy-locality-1} holds in supremum form
over $y$.
\end{proof}

This completes the proof of clause (ii) of Theorem~\ref{4.24:thm-survivor-block-expansion}.

\begin{lemma}[Independence of the Gaussian objects from exterior data]
\label{4.24:prf-exterior-independence}
Let $\lambda$ be the tuple formed by the following variables:
\begin{itemize}
\item the exterior slots $(A_{j'}\lceil S(e))_{j'>j}$ of the parents of the members $e$ of
$\mathfrak{E}^{\mathrm{StB}}$, where $S(e)$ is the set $S$ of the stored term that is the
parent of $e$ (Definition~\ref{4.24:def-stored-terms});
\item the un-integrated level-$j$ slots of the components arrested at stages $\le n$, or leaving
the run at the present stage before \cite[(1.53)]{B-CMP122a}; by
Definition~\ref{4.24:def-survivors-bonds} these components lie outside $Z^{\sharp}$;
\item for each $W\in\mathfrak{W}$, the variable
$\phi_W=(A_j\lceil(\mathsf{R}_j\cap W),\,x\lceil S^{P}_{W})$.
\end{itemize}
Each coordinate of the first two kinds ranges over the datum of its parent stored term
(a real box or a complex polydisc, as in Definition~\ref{4.24:def-stored-terms}). Each
$\phi_W$ ranges over the exterior box $\mathfrak{X}_W$. Let $D$ be the product of these
ranges. Consider the two conditions:
\begin{enumerate}
\item[(1)] the data of the survivors' expansion that are not members of
$\mathfrak{E}^{\mathrm{StB}}$ are independent of $\lambda$. These are the Gaussian objects,
the characteristic function $\chi^{s}$, the sets $P$ and $Z_0$, the decoupled cubes, the
partitions and the contours;
\item[(2)] the value $v_e$ of each member $e\in\mathfrak{E}^{\mathrm{StB}}$ is continuous on
$D_e$ and holomorphic in the polydisc coordinates of $D_e$, and $|v_e|\le\nu^{\natural}(e)$
on $D_e$ for a majorant $\nu^{\natural}(e)$ that is independent of $\lambda$ and of the pair
$(\mathbb{U},\mathbb{J})$.
\end{enumerate}
Then:
\begin{enumerate}
\item[(a)] if \textup{(1)} and \textup{(2)} hold, then the conclusion of the
parameter-dependence lemma stated with the stored terms in
Definition~\ref{4.24:def-stored-terms} holds for the survivors' expansion, with parameter
$\lambda\in D$. Precisely: the expansion identity holds at every $\lambda\in D$; the
bounds on the terms and their totals hold with majorants independent of $\lambda$; and
every term, as well as the sum, is continuous on $D$ and holomorphic in the polydisc
coordinates;
\item[(b)] condition \textup{(1)} holds.
\end{enumerate}
\end{lemma}

\begin{proof}
\emph{Part (a).} We run the proof of the parameter-dependence lemma of
Definition~\ref{4.24:def-stored-terms} again, in the form used
here. That proof has three parts.

The first part concerns the first step of the cluster-expansion bound for the survivors.
This step is an identity at each value of the parameter. Here the expansion identity of
the survivors' expansion holds at each fixed $\lambda\in D$, for two reasons.

First, the Gaussian objects, $\chi^{s}$, $P$, $Z_0$, the partitions and the contours do
not move with $\lambda$; this is condition~(1).

Second, the members enter only through their values $v_e(\lambda)$. At fixed $\lambda$
these are bounded functions inserted into an identity, and an identity that holds for
arbitrary bounded inserted functions holds in particular for these.

The second part of the proof concerns the remaining steps. They bound every term of the
expansion by a majorant built from the weights $\nu(e)$ and from $\lambda$-free data, and
then sum the majorants. Here these steps are the bound on a single term of the expansion
(Lemma~\ref{4.24:prf-single-term}) and the summation of the resulting bounds, with totals
uniform in the number of arrested components (Lemma~\ref{4.24:prf-uniform-totals}). These
bounds are built from the majorants $\nu^{\natural}(e)$ and from the data listed in~(1).
They are therefore independent of $\lambda$ as soon as each member's majorant
$\nu^{\natural}(e)$ is independent of $\lambda$, which is condition~(2).

The third part of the proof concerns regularity. Each term of the expansion is a finite sum
of products. Each factor is either a contour integral, over a $\lambda$-free contour, of
finitely many of the $v_e$, or a $\lambda$-free factor. By~(2) each $v_e$, read as a
function on $D$ through the coordinates of $\lambda$ that $e$ reads, is continuous on $D$
and holomorphic in the common polydisc coordinates. A contour integral over a fixed
contour of such functions is again continuous on $D$ and holomorphic in the common polydisc
coordinates, and finite sums and products preserve both properties.

The majorants of the second part are independent of $\lambda$ and summable. Hence the
expansion converges uniformly on $D$. A uniform limit of continuous functions is
continuous, and by Weierstrass's theorem a uniform limit of holomorphic functions is
holomorphic. Therefore the sum is continuous on $D$ and holomorphic in the polydisc
coordinates.

Nothing in this argument distinguishes among the kinds of members of
$\mathfrak{E}^{\mathrm{StB}}$. The members of $\mathfrak{E}^{s}$ drawn from the element
family $\mathfrak{E}$, the (S3)- and (S4)-pieces of Definition~\ref{4.24:def-stored-terms}
and the cross pieces of $\mathfrak{E}^{\times}$ are all treated alike.

The proviso of the parameter-dependence lemma requires the (S3)- and (S4)-pieces to be
independent of the parameter. It serves the particular parameter of the corollary of that
lemma stated with Definition~\ref{4.24:def-stored-terms},
namely the exterior slots with $j'>j$, which are read neither by the A--B coupling pieces
$\mathbb{P}^{(j)}$ of Ba{\l}aban's construction nor by the other pieces of the stage that
enter that corollary. The proof does not use the proviso once the (S3)- and (S4)-pieces are
entered as members of $\mathfrak{E}^{\mathrm{StB}}$ with $\lambda$-free majorants.
What the argument needs is therefore exactly~(1) and~(2), and part~(a) is proved.

\emph{Part (b).} We first treat the covariance.

The operator $\mathcal{A}_{ss}$ is made of two parts: the collar block of
$\mathcal{A}(\mathbf{1})$, the operator associated with the stage form of
Definition~\ref{4.24:def-stage-form} at decoupling parameters $\sigma'\equiv 1$; and, for
each component, the rim form of
Ba{\l}aban's construction on $\mathsf{R}_j\cap Z^{s}$. It is a kernel in the pair
$(\mathbb{U},\mathbb{J})$. Indeed, $\Delta^{(j)}(\mathbb{U},\mathbb{J})$, the operator
entering the form $C^{*}\Delta^{(j)}C$ of Definition~\ref{4.24:def-stage-form}, is the
operator of the statement fixing the stage operators, built on the pair over
$\mathbb{B}^{(n)}_k$, and $C=C(\mathbb{U})$; both are functions of the pair alone.

The stage-$(n+1)$ pair is a function of the data $V$ through $U^{(n+1)}_k(V)$. This is the
corollary of the parameter-dependence lemma stated with
Definition~\ref{4.24:def-stored-terms}: the backgrounds are functions of the data $V$
and not of the fluctuation variables $A$. Since $U^{(n+1)}_k$ carries only levels
$\ge j+1$, the pair carries levels $\ge j+1$ and is independent of the level-$j$ collar
variable $x$; this is the $x$-freeness clause of the lemma on the collar variable of
an arrested component, which states that the objects of the arrested component that are
functions of $U^{(n+1)}_k$ are independent of $x$.

Every coordinate of $\lambda$ is either a slot of a fluctuation variable $A_{j'}$ with
$j'\ge j$ or a restriction of the collar variable $x$. Since the pair depends on neither,
it is independent of $\lambda$. Hence $\mathcal{A}_{ss}$, a kernel in the pair, is
independent of $\lambda$, and so is $\mathcal{C}_s$. So are
its resolvents, their walk expansions and the family decoupled in the parameter
$\mathsf{s}$, since all of these are built from $\mathcal{A}_{ss}$.

We now treat the remaining data of~(1). The characteristic function $\chi^{s}$ reads only
the survivors' variables $\phi_s$ \cite{B-CMP122a}, and none of these is a coordinate of
$\lambda$. The set $P$ is contained in $s$. Finally, none of the following depends on
$\lambda$; they are fixed:
\begin{itemize}
\item the decoupled cubes;
\item the partition of unity $\{\zeta_{\square}\}$;
\item the filing partitions of the expansion;
\item the contours of integration in $\mathsf{s}$, $\tau$, $t$ and $u$.
\end{itemize}
This proves~(1).
\end{proof}

\begin{lemma}[Exterior-free majorants of the members, by kind]\label{4.24:prf-member-majorants}
Work at stage $n+1$ of the step at level $j$, in the setting of
Definitions~\ref{4.24:def-stage-form}, \ref{4.24:def-stored-terms},
\ref{4.24:def-exterior-variables} and~\ref{4.24:def-survivor-elements}.

Here $W$ runs over the arrested components $\mathfrak{W}$, and $\phi_W$ denotes the exterior
variables of $W$, which are among the parameters $\lambda$ of the elements. We write $x$ for
the collar coordinate of $\phi_W$ and $\mathsf{B}$ for the collar variable, both as in
Definition~\ref{4.24:def-exterior-variables}. We write $\varrho^{B}_j=6.6\,\mathfrak{b}'_j$
for the radius of the collar disc, and $\varrho_{j'}$ for the radius of the complex polydisc
of the stored-term domains of Definition~\ref{4.24:def-stored-terms}.

Assume the following.
\begin{enumerate}
\item[(a)] Put $r_{A,j'}:=\varrho_{j'}/\mathfrak{b}_{j'}$ for each level $j'$ in the range
$h\le j'<k$ of the stored-term domains of Definition~\ref{4.24:def-stored-terms}, where $h$
denotes the lowest level of that range; the level $j$ of the stage lies in this range. There is
a number $r_A$ with
\begin{equation}\label{4.24:eq-member-majorants-1}
r_A\le\min_{h\le j'<k}r_{A,j'},\qquad r_A\ge 2 .
\end{equation}
This holds under the standing smallness condition on $\gamma$ under which the stored-term
domains of Definition~\ref{4.24:def-stored-terms} are set up, and is assumed here.
\item[(b)] The kernel entries $Q(Y;b,b')$ of the quadratic mixed pieces of the stage satisfy
\eqref{4.24:eq-member-majorants-2}; the value of such a piece is then bounded by the right side
of \eqref{4.24:eq-member-majorants-2} times the radii of the ranges of its two arguments. The
pieces $\mathbb{V}''(Y)$ satisfy \eqref{4.24:eq-member-majorants-3}. Both bounds hold on the
range $|A_j|\le g_j^{-1}\delta_j$ of the rim arguments:
\begin{align}
|Q(Y;b,b')|&\le C_3C_1\delta'_jM^4e^{C_2\kappa_1}
  e^{-\frac18(\kappa_1-1)d(Y)-\frac12(\kappa_1-1)M^{-1}|b-b'|},
  \label{4.24:eq-member-majorants-2}\\
|\mathbb{V}''(Y)|&\le\mathfrak{p}_j\,e^{-(1-2\varpi)\kappa d(Y)} .
  \label{4.24:eq-member-majorants-3}
\end{align}
Here $C_1$, $C_2$, $C_3$ and $\varpi$ are the constants of the statement bounding the mixed
pieces $Q(Y)$ and $\mathbb{V}''(Y)$ of the stage, which is taken here by statement, and $d(Y)$
is the size of the domain $Y$ as measured in those bounds.
\item[(c)] The following holds under the arrest condition (the condition of the stage under
which the components in $\mathfrak{W}$ are arrested), for the mixed pieces of the stage (the
pieces of $\mathbb{P}^{(j)}$ among them, \cite[Lemma~2]{B-CMP116}) and for every member of kind
\textup{(}$\alpha$\textup{)} in \textup{(3)} below whose cubes $\mathsf{b}(e)$ meet the collar
of $W$. Each value $v_e$
is holomorphic in the collar variable on $|\mathsf{B}|<\varrho^{B}_j$, jointly with the further
variables named in the holomorphy premise of the stage bound of
Definition~\ref{4.24:def-exterior-variables}, and $|v_e|\le\nu^{\natural}(e)$
there.
\item[(d)] This hypothesis is used only under lettering (II) of \eqref{4.24:eq-stage-form-b}.
The statements of Definition~\ref{4.24:def-stage-form} hold with the kernel $\mathcal{A}$ read
on $\mathfrak{Bo}$ (rim and collar together), and with $C_\rho$ read on $\mathfrak{Bo}$.
\end{enumerate}

Then the following hold.
\begin{enumerate}
\item[(1)] Work under the collar lettering (I) of \eqref{4.24:eq-stage-form-b}.
  \begin{itemize}
  \item Every cross member $e_a\in\mathfrak{E}^{\times}$ reads $\phi_W$ only through $x(b')$,
  in the box $|x|\le 2.4\,\mathfrak{b}'_j$.
  \item Its majorant $\nu^{\natural}(e_a)$ in \eqref{4.24:eq-stage-form-a} is independent of
  $\lambda$ as it stands.
  \item $K_a$ is that of \eqref{4.24:eq-stage-form-a}, and $K_a$, $\vartheta_a$ and the
  change-of-variables ceiling reading $K_a$ are unchanged.
  \item No element of the stage form reads $A_j$ restricted to the rim of $W$.
  \item No shift of the change of variables acts on any $A_j$.
  \end{itemize}
\item[(2)] Let $c$ be a stored term and consider the rim cubes of $W$ that $c$ reads. On each
of them the exterior rim box $\{|A_j|<\mathfrak{b}_j\}$ lies in $\mathfrak{D}_{\mathcal{S}}(S)$,
both on cubes strong for $c$ and on cubes weak for $c$. Thus the containment clause of
Definition~\ref{4.24:def-exterior-variables} holds on strong and weak cubes alike.
\item[(3)] Besides the cross members of \textup{(1)}, the elements of
$\mathfrak{E}^{\mathrm{StB}}$ reading $\phi_W$ are of the following four kinds:
  \begin{itemize}
  \item[\textup{(}$\alpha$\textup{)}] members whose parent stored term reads un-integrated
  slots of $W$, or of a component excluded at an earlier stage;
  \item[\textup{(}$\beta$\textup{)}] mixed pieces $Q(Y)$ and $\mathbb{V}''(Y)$ with an argument
  on $W$;
  \item[\textup{(}$\gamma$\textup{)}] cross members $e_a$ ending on the rim;
  \item[\textup{(}$\delta$\textup{)}] change-of-variables objects $o\in\mathfrak{O}^{s}$.
  \end{itemize}
Kind \textup{(}$\gamma$\textup{)} is empty under lettering (I). For all these elements:
  \begin{itemize}
  \item the majorant $\nu^{\natural}(e)$ is independent of $\lambda$ and of
  $(\mathbb{U},\mathbb{J})$;
  \item $D_e$ is the product of a box and a polydisc.
  \end{itemize}
Consequently the parameter clause of Definition~\ref{4.24:def-stored-terms} applies to
$\mathfrak{E}^{\mathrm{StB}}$.
\item[(4)] Under lettering (II) of \eqref{4.24:eq-stage-form-b}, together with hypothesis
\textup{(d)}, put $r_{\mathrm{rim}}:=\max\{\varrho_j,\varrho^{B}_j\}$.
  \begin{itemize}
  \item The majorants of the cross members of kind \textup{(}$\gamma$\textup{)} are
  \begin{equation}\label{4.24:eq-member-majorants-4}
  \nu^{\natural}(e_a)=\tfrac12B_{\mathcal{A}}\,r_{\mathrm{rim}}^2\,e^{-\delta d}
  e^{-(\kappa_1-1)|Y'|},
  \end{equation}
  with $B_{\mathcal{A}}$, $\delta$ and $d$ as in \eqref{4.24:eq-stage-form-a}.
  \item $K_a$ is replaced by $K_a^{\sharp}:=K_a(r_{\mathrm{rim}}/\varrho^{B}_j)^2$.
  \item $\vartheta_a$ and $\vartheta_{\mathrm{cv}}$ are replaced by the quantities
  $\vartheta_a^{\sharp}$ and $\vartheta_{\mathrm{cv}}^{\sharp}$, obtained by this replacement;
  the replacement brings in a fixed power of $\log g_j^{-2}$, which is set against
  $\varepsilon_Y\le 2E(x_k)^6$.
  \item The change-of-variables ceiling reading $K_a$ and the $\mathfrak{p}_j$-term of
  $(Q^{+})'$ are re-maximised accordingly.
  \item This is a condition on $\gamma$, and it imposes no cap in the sense of
  Definition~\ref{2.3:def-cap}.
  \item $C_0^{\sharp}$ and the coefficient $5$ of clause (F2) of the stage bound of
  Definition~\ref{4.24:def-exterior-variables} are unchanged.
  \end{itemize}
\item[(5)] The output's domain is
\begin{equation}\label{4.24:eq-member-majorants-5}
\mathfrak{D}^{\flat}_{n+1}[(\mathrm{amb})_W](X)\times\mathfrak{D}_{\mathcal{S}}(X)
\times\prod_W\mathfrak{X}_W .
\end{equation}
This is the domain of clause (F1) of the stage bound of
Definition~\ref{4.24:def-exterior-variables}. The additional clause of (F1) for domains $X$
meeting an enlarged ambient component $Z_W^{+}$ applies to exactly those $X$ that meet some
$Z_W^{+}$.
\end{enumerate}
\end{lemma}

\begin{proof}
\emph{Scope on the rim; the cross members.}
Under lettering (I) of \eqref{4.24:eq-stage-form-b} the kernel $\mathcal{A}$ is lettered on
the starred collar bonds $\mathcal{C}_j$. Hence the indices $(b,b',Y')$ of the cross members in
the third clause of \eqref{4.24:eq-stage-form-a} are pairs of collar bonds. A cross member
$e_a\in\mathfrak{E}^{\times}$ therefore reads $\phi_W$ only through the collar coordinate
$x(b')$.

This coordinate ranges over the box $|x|\le 2.4\,\mathfrak{b}'_j$, since the supports
$2.2\,\mathfrak{b}'_j$ and $2.4\,\mathfrak{b}'_j$ in clause (iv) of the change-of-variables
lemma, and the range
$2.3\,\mathfrak{b}'_j$ in clause (d) of the theorem on the frozen members and
change-of-variables objects, all lie inside $2.4\,\mathfrak{b}'_j$. The box lies in the disc
$\{|\mathsf{B}|<\varrho^{B}_j=6.6\,\mathfrak{b}'_j\}$, on which, by
Definition~\ref{4.24:def-exterior-variables}, $x$ is weak at
$e^{-2\varepsilon_\star}\varrho^{B}_j$ for un-shifted elements.

By the first clause of \eqref{4.24:eq-stage-form-a}, $e_a$ is bilinear and holomorphic on
$\mathfrak{D}^{\flat}_{n+1}(Y')\subset\mathfrak{K}_P$. By \eqref{4.24:eq-stage-form-a},
\begin{equation*}
|v_{e_a}|\le\nu^{\natural}(e_a)
=\tfrac12B_{\mathcal{A}}(\varrho^{B}_j)^2e^{-\delta d}e^{-(\kappa_1-1)|Y'|}
\end{equation*}
on the product of the two collar discs $\{|\mathsf{B}|<\varrho^{B}_j\}$, one for each argument,
which contains the product of the disc and the box $|x|\le 2.4\,\mathfrak{b}'_j$.

The factor $e^{-(\kappa_1-1)|Y'|}$ is the Cauchy price of extracting the cubes of $Y'$, bounded
in Lemma~\ref{4.24:lem-cube-extraction-bound} by \eqref{4.24:eq-cube-extraction-bound-a}. It
depends on $|Y'|$ only and not on $\lambda$. The same formula holds for every $Y'$, including
$|Y'|=R_k$ exactly. So $\nu^{\natural}(e_a)$ is independent of $\lambda$ as it stands, and
$K_a$, $\vartheta_a$ and the change-of-variables ceiling reading $K_a$ are unchanged under
lettering (I). This is the first half of \textup{(1)}.

The survivors' rim variables are integrated against Ba{\l}aban's rim form, which is
component-diagonal by its definition in \cite{B-CMP119}. Its inter-component remnants are
boundary terms $\mathbf{B}$ created at step $j+1$. By Definition~\ref{4.24:def-stored-terms}
these are:
\begin{itemize}
\item explicit quadratics, entire in $A$;
\item bounded by $e^{-R_j}$ on the polydisc of radius $\varrho_j$;
\item weak, with empty strong set $\mathcal{S}=\emptyset$.
\end{itemize}
At stage $n+1$ they are read as parents $c$ of members of kind \textup{(}$\alpha$\textup{)}
below. Hence no element of the stage form reads $A_j$ restricted
to the rim of $W$. No shift of the change of variables acts on any $A_j$ either, since by
clause (iv) of the change-of-variables lemma the shifts $\delta_b$ of the change of variables
are supported on $S^{P}$, the region of the prepared data on which that lemma places them.
This completes \textup{(1)}.

We now take the kinds of elements that read $\phi_W$ in turn.

\emph{Kind \textup{(}$\alpha$\textup{)}: members whose parent $c$ is a stored term reading
un-integrated slots of $W$, or of a component excluded at an earlier stage.}
By Definition~\ref{4.24:def-stored-terms}:
\begin{itemize}
\item $c$ is defined on $\mathfrak{U}_c(S)\times\mathfrak{D}_{\mathcal{S}}(S)$ and is
continuous there;
\item $c$ is holomorphic in $(\mathbb{U},\mathbb{J})$ and in its weak slots for fixed real
strong part;
\item $\mathfrak{N}(c)$ is the supremum of $|c|$ over that product.
\end{itemize}
By construction, then, $\mathfrak{N}(c)$ does not depend on where in
$\mathfrak{D}_{\mathcal{S}}(S)$ the exterior slots lie, provided they lie in it.

The exterior rim box of $W$ is $\{|A_j|<\mathfrak{b}_j\}$, the support of $\chi^{(j)}$. That it
lies inside $\mathfrak{D}_{\mathcal{S}}(S)$ on the rim cubes of $W$ that $c$ reads is the
containment clause of Definition~\ref{4.24:def-exterior-variables}: the datum of $c$ holds
$A_j$ restricted to $W$ in $\mathfrak{D}_{\mathcal{S}}$.
\begin{itemize}
\item On a rim cube strong for $c$ the containment is literal. There
$\mathfrak{D}_{\mathcal{S}}(S)$ is, by Definition~\ref{4.24:def-stored-terms}, the real box of
radius $\mathfrak{b}_j$.
\item On a rim cube of $W$ weak for $c$, the domain $\mathfrak{D}_{\mathcal{S}}(S)$ is the
complex polydisc of radius $\varrho_j$. It contains the real box of radius $\mathfrak{b}_j$,
because hypothesis (a), applied at $j'=j$, gives
\begin{equation}\label{4.24:eq-member-majorants-6}
\varrho_j=r_{A,j}\,\mathfrak{b}_j\ge r_A\,\mathfrak{b}_j\ge 2\,\mathfrak{b}_j .
\end{equation}
\end{itemize}
This holds under the standing smallness condition on $\gamma$ of hypothesis (a), which this
step inherits, with its constants, together with Definition~\ref{4.24:def-stored-terms}, its
parameter clause and the filing rule for outputs of stored terms, as the inheritance statement
of the stage provides. Hence the containment clause
holds on strong and weak cubes alike, which is \textup{(2)}.

The enlargement of the radius to $r_{\mathrm{rim}}$, used under lettering (II) below, plays no
part in this step. It enlarges the radius of a majorant of the bilinear cross members of the
stage form. Here the question is the domain of definition of a stored term itself, which no
enlargement of a majorant's radius changes.

The majorant $\nu^{\natural}(e)$ of the member is built, by the majorant lemma for
stored-term members, from $\mathfrak{N}(c)$ and from numbers independent of $\lambda$. These
numbers are the factors $K_Q/r_A$, $\varepsilon_S$ and $e^{\kappa_A|\mathcal{S}_{>j}|}$, as
the filing rule for outputs of stored terms enters them; here $\mathcal{S}_{>j}$ is the set of
strong cubes of $c$ at levels above $j$, and $K_Q$, $\varepsilon_S$ and $\kappa_A$ are the
constants of the stored-term bounds of Definition~\ref{4.24:def-stored-terms}.

The parameter domain is
$D_e:=\mathfrak{D}_{\mathcal{S}(c)}(S)$ at the exterior levels, multiplied by the $x$-box when
$\mathsf{b}(e)$ meets the collar of $W$. This is exactly the lettering of the parameter clause
of Definition~\ref{4.24:def-stored-terms} and of the filing rule.

The functional $\Phi_{t,\sigma}(b(\mathsf{B}))$ of the change of variables on the collar, with
its parameters $t,\sigma$ and its argument $b(\mathsf{B})$ as in the construction of that change
of variables, reads $x$ only inside such members, by the support statement of that
construction. Each of them stands under its own majorant
$\nu^{\natural}$, the supremum over $|\mathsf{B}|<\varrho^{B}_j$ being supplied by
hypothesis (c).

\emph{Kind \textup{(}$\beta$\textup{)}: mixed pieces $Q(Y)$ and $\mathbb{V}''(Y)$ with an
argument on $W$} (the pieces of $\mathbb{P}^{(j)}$ among them). Hypothesis (c) covers them;
for $\mathbb{P}^{(j)}$ it rests on \cite[Lemma~2]{B-CMP116}, under the arrest condition.

Their majorant $\nu^{\natural}$ is defined as the supremum, over the disc $\times$ the slot
ranges of $W$, of the bounds \eqref{4.24:eq-member-majorants-2}, multiplied by the two radii,
and \eqref{4.24:eq-member-majorants-3}. It is independent of $\lambda$. It enters the
$\mathfrak{p}_j$-term of $(Q^{+})'$, and not $K_a$.

A collar argument on $W$ ranges over the $x$-box, which lies in the disc. A rim argument ranges
over $|A_j|<\mathfrak{b}_j=g_j^{-1}\delta_j$, which lies in the stated range
$|A_j|\le g_j^{-1}\delta_j$ of hypothesis (b).

\emph{Kind \textup{(}$\gamma$\textup{)}: cross members $e_a$ ending on the rim.} Under
lettering (I) there are none, since by the first paragraph all indices of $\mathcal{A}$ are
collar bonds.

Now suppose $\mathcal{A}$ is lettered by lettering (II), that is, as one joint kernel on rim
and collar bonds with inter-component rim entries inside it. This is not the lettering of
\cite{B-CMP119}, whose rim form is component-diagonal. Then hypothesis (d), the statements of
Definition~\ref{4.24:def-stage-form} read on $\mathfrak{Bo}$, is an input.

Cross members ending on the rim read $A_j$ restricted to the rim of $W$ on the exterior real
box $|A_j|<\mathfrak{b}_j$. The weak polydisc covers this box by
\eqref{4.24:eq-member-majorants-6}, exactly as in kind \textup{(}$\alpha$\textup{)}. The rim
radius is
\begin{equation*}
r_{\mathrm{rim}}=\max\{\varrho_j,\varrho^{B}_j\}=\max\{\mathfrak{b}_j,\varrho_j,\varrho^{B}_j\},
\end{equation*}
where the second equality holds because $\varrho_j\ge2\mathfrak{b}_j\ge\mathfrak{b}_j$ by
\eqref{4.24:eq-member-majorants-6}.

Since $r_{\mathrm{rim}}\ge\varrho^{B}_j$, the majorant becomes
\eqref{4.24:eq-member-majorants-4}, that is, $K_a$ is replaced by
$K_a^{\sharp}=K_a(r_{\mathrm{rim}}/\varrho^{B}_j)^2$, with
$C_\rho$ read on $\mathfrak{Bo}$. This replacement has the following effects.
\begin{itemize}
\item $\vartheta_a$ becomes $\vartheta_a^{\sharp}$.
\item Clause (e) of the theorem on the frozen members and change-of-variables objects sums the
parents' majorants in the same shape, so $\vartheta_{\mathrm{cv}}$ becomes
$\vartheta_{\mathrm{cv}}^{\sharp}$.
\item The replacement brings in a fixed power of $\log g_j^{-2}$, set against
$\varepsilon_Y\le2E(x_k)^6$. This is a condition on $\gamma$, all other constants being chosen
before $\gamma$.
\item The change-of-variables ceiling reading $K_a$ and the $\mathfrak{p}_j$-term of $(Q^{+})'$
are re-maximised.
\item $C_0^{\sharp}$ and the coefficient $5$ of clause (F2) are unchanged.
\end{itemize}
Only $\vartheta_a$, $\vartheta_{\mathrm{cv}}$ and the ceilings reading them change, and no cap
in the sense of Definition~\ref{2.3:def-cap} is imposed under either lettering. This is
\textup{(4)}.

\emph{Kind \textup{(}$\delta$\textup{)}: change-of-variables objects $o\in\mathfrak{O}^{s}$.}
\begin{itemize}
\item By clause (d) of the theorem on the frozen members and change-of-variables objects, $x$
is a silent slot of $o$.
\item The data $H_\sigma$, $\mathsf{g}'_j$ and $\sigma$ of the change of variables, as fixed in
that theorem, are independent of $\lambda$.
\item The bounds of clause (e) of that theorem are uniform on the real box
$|x|\le2.3\,\mathfrak{b}'_j$ and independent of the background pair $(\mathbb{U},\mathbb{J})$.
\item The other slots of the underlying element $e$ keep the datum of $e$, as treated under
kinds \textup{(}$\alpha$\textup{)} and \textup{(}$\beta$\textup{)}.
\end{itemize}
A silent slot is a box coordinate in the sense of the parameter clause of
Definition~\ref{4.24:def-stored-terms}, on which only continuity is required.

\emph{Conclusion.} The cross members of \textup{(1)} and the four kinds have been covered. In
each case the majorant is independent of $\lambda$ and of $(\mathbb{U},\mathbb{J})$, and $D_e$
is the product of a box and a polydisc. Hence the parameter clause of
Definition~\ref{4.24:def-stored-terms} applies to $\mathfrak{E}^{\mathrm{StB}}$, which is
\textup{(3)}.

The filing rule for outputs of stored terms then applies without change.
\begin{itemize}
\item For a family of elements entering one output and $j'>j$, the strong set of the output at
level $j'$ is the union $\bigcup_{e}\mathcal{S}_{j'}(c(e))$ of the parents' strong sets, over
the elements $e$ of the family.
\item The level-$j$ structure on $W$ is carried as exterior data. By
Definition~\ref{4.24:def-exterior-variables}, the level-$j$ strong weight of $W$ stays in
the weight $\mathfrak{N}^{\natural}$ of that definition.
\item A level-$j'$ cube is weak for the output if and only if it is weak for every parent
reading it, and its radius $\varrho_{j'}$ is unchanged. So the polydisc in $A$ never shrinks,
as the remark on the radii of the outputs' polydiscs records.
\item The $x$-coordinates are box coordinates of $\prod_W\mathfrak{X}_W$.
\end{itemize}
The output's domain is therefore \eqref{4.24:eq-member-majorants-5}. This is the domain of
clause (F1) of the stage bound of Definition~\ref{4.24:def-exterior-variables}, in
which the factor $\prod_W\mathfrak{X}_W$ consists of the boxes of the $\phi_W$. The
additional clause of (F1) for domains $X$ meeting an enlarged ambient component $Z_W^{+}$
applies to exactly those $X$ that meet some $Z_W^{+}$. This is \textup{(5)}.

Finally, the rim Gaussian of $W$ itself is the weight of the un-performed $A_j$-integral of $W$
inside $\mathbf{T}'_k(W)$, beside $\mathsf{q}_W$. It is not an object of the survivors' side
at stage $n+1$.
\end{proof}

\begin{lemma}[Constants free of volume, age and component count]\label{4.24:prf-constant-uniformity}
This is part~(iv) of the main statement of this section.
Among the constants entering the bounds of this section are the following:
\begin{itemize}
\item the per-cell count $c_G'$ and the series $\mathfrak{S}$ of Lemma~\ref{4.24:prf-constant-unchanged};
\item the coercivity constant $a_{\mathbb{C}}$;
\item the cross totals $\mu^{\times}$, $\mathsf{M}^{\times}$ and the change-of-variables total
$\mathsf{M}^{\mathrm{cv}}$ of the cluster expansion;
\item the bound on the sum $\sum_{W\in\mathfrak{W}}\mathbb{C}^{\mathrm{fr}}_W$ of the frozen densities.
\end{itemize}
These constants have the following properties.
\begin{enumerate}
\item[(a)] None of them depends on the number $|\mathfrak{W}|$ of arrested components, or on the number
of components.
\item[(b)] None of them depends on the volume or on the ages of the large-field regions.
\item[(c)] They depend on the level $k$ only through the couplings $g_j$ and through $R_k$, and on the
number $N_0$ of Ba{\l}aban's construction, which is distinct from the packing numbers
$N^{\mathrm{pk}}_m$ of Lemma~\ref{4.24:lem-packing-count}, and the number $N$ of
\cite[pp.~176--177]{B-CMP122a} only through conditions imposed before this section. They do not
depend on $n^{+}$.
\item[(d)] No lower bound on $g$ and no upper bound on $L$ is used, and no new lower bound on the
auxiliary parameters is imposed. The order of choice
\begin{equation}\label{4.24:eq-constant-uniformity-2}
L\to\mathfrak{M}(R_1)\to\kappa\to\kappa_1\to M_1\to M\to\gamma,
\end{equation}
in which $\kappa,\kappa_1,M_1,M$ are among the auxiliary parameters $\mathfrak{p}$ of
Notation~\ref{2.1:not-prefix} and the quantity $\mathfrak{M}(R_1)$, depending on $R_1$, is chosen
after $L$ and before $\kappa$, is unchanged.
\end{enumerate}
\end{lemma}

\begin{proof}
We go through the constants one at a time, and then through each possible source of dependence.

\emph{The per-cell count $c_G'$.} The count $c_G'$ is taken per $\pi_{j+1}$-cell. By
\eqref{4.24:eq-component-separation-c} of Lemma~\ref{4.24:lem-component-separation}, a cell meets at
most one component. Hence no count of components enters $c_G'$. This is the count used in the
formula $c_G=c_G'\,\mathfrak{S}+c_{\mathrm{sc}}$ of Lemma~\ref{4.24:prf-constant-unchanged}.

\emph{The hull sum.} The series $\mathfrak{S}$ of Lemma~\ref{4.24:prf-constant-unchanged} is a sum
taken per anchor cell. Its terms are the universal walk majorants of
Lemma~\ref{4.24:prf-constant-unchanged}, which do not depend on the region. Hence the product
$c_G'\,\mathfrak{S}$ in $c_G$ inherits the
properties (a)--(c) from $c_G'$ and from $\mathfrak{S}$.

\emph{The joint count.} The densities $\mathbb{C}^{s}_G$ of the survivors and $\mathbb{C}^{G}_W$ of the
arrested components are counted jointly, per cell. This is the count of
Lemma~\ref{4.24:prf-elimination-factors}. It therefore involves no number of components either.

\emph{The coercivity constant.} The constant $a_{\mathbb{C}}=(4B_0)^{-1}$ of
\eqref{4.24:eq-block-coercivity-a} is supplied by Lemma~\ref{4.24:lem-block-coercivity}. It follows
from an upper bound on $\gamma$ imposed before this section, through the condition
$\alpha_5\|\mathcal{C}\|<\tfrac12$ of \cite[(2.31)]{B-CMP116}. Here $B_0=B_0(d,L)$. In part~(i) of
the main statement of this section $a_{\mathbb{C}}$ is used only qualitatively. In part~(ii) of
that statement it is used only through a margin.

\emph{The conditions on $R_k$.} The conditions used on $R_k$ are
\begin{equation*}
R_k\ge 3
\qquad\text{and}\qquad
0.07\,C_{\mathrm{pack}}(d)\le e^{2c_\flat}\le e^{R_k}.
\end{equation*}
Both are already among the conditions, of supremum form, of part~(v) of the smallness condition
for the change of variables, which was imposed before this section; no condition is added. That
condition bounds $g$ from above only.

\emph{The smallness of the weight $\nu$.} The smallness bound $\nu\le\tfrac12$ of the cluster
expansion of this section follows from the bound $\mu^{\mathrm{StB}}\le\tfrac12$ of
Lemma~\ref{4.24:prf-uniform-totals}. No separate upper bound on the coupling is needed for it.

\emph{The total $\mu^{\times}$.} This total is anchored at a survivor cell, and all partners are summed.
Hence $|\mathfrak{W}|$ does not enter it.

\emph{The totals $\mathsf{M}^{\times}$ and $\mathsf{M}^{\mathrm{cv}}$.} These totals are bounded by the
packing count $N^{\mathrm{pk}}_m$ and the constant $C_{\mathrm{pack}}$ of
\eqref{4.24:eq-packing-count-a} in Lemma~\ref{4.24:lem-packing-count}. That count is purely
geometric: it holds for either of the two distances in which it is stated, and it is taken at
$d=4$. The component budget
\begin{equation*}
\bar V=C_V(100L\check{R})^d,
\end{equation*}
where $C_V$ is its constant prefactor, is a supremum over components. It is already present in the
smallnesses $\vartheta_a$ and $\vartheta_{\mathrm{cv}}$, through a condition imposed before this
section, and its size is
compensated by the smallness factor $\varepsilon_Y=2e^{-6R_k}$.

\emph{The sum over arrested components.} Fix a point $y$. The sum
$\sum_{W\in\mathfrak{W}}\mathbb{C}^{\mathrm{fr}}_W$ at $y$ is a repartition of a sum that does not
depend on $|\mathfrak{W}|$. Indeed, each frozen member is assigned to exactly one $W$. Moreover,
$\mathbb{C}^{G}_W(y)\ne0$ for at most one $W$, namely the one whose block contains $y$. No constant
depends on the order in which the members are assigned.

\emph{Dependence on $N_0$ and $N$.} These numbers enter only through two conditions imposed before
this section. The first is the condition on the totals of the cluster expansion. The second is
the window condition of the large-field step. The dependence on $N$ is the one already present in
\cite[pp.~176--177]{B-CMP122a}.

\emph{Dependence on $k$, $n^{+}$, ages and volume.} The level $k$ enters only through the couplings
$g_j$ and through $R_k$. The index $n^{+}$ occurs in one statement only, and that statement enters
no constant of this section. No constant depends on the ages of the large-field regions or on the
volume.

\emph{No lower bound on $g$.} The following conditions are suprema over the inverse couplings
$x_k=g_k^{-2}\ge\gamma^{-2}$:
\begin{itemize}
\item parts~(iv) and~(v) of the smallness condition for the change of variables;
\item the arrest condition;
\item the primed totals condition of the cluster expansion;
\item the condition $\alpha_5\|\mathcal{C}\|<\tfrac12$.
\end{itemize}
The smallness condition of the cluster expansion, imposed before this section, has among its
members the product $K_\Delta\gamma_0$. This member is controlled by a bound on it established
earlier in this chapter, whose constant is a supremum over
$\mathsf{t}\ge\mathsf{t}_\gamma$, where $\mathsf{t}_\gamma:=\log\gamma^{-2}$.
Here
\begin{equation}\label{4.24:eq-constant-uniformity-1}
\gamma_0=\gamma_0^{(k)}=5\max_{j\le k}w_*(j)\,\bar\alpha_j ,
\end{equation}
with $\bar\alpha_j$ a coupling amplitude. Hence no lower bound on $g$ is used anywhere.

\emph{No new floor.} The lower bound on $M_1$ imposed earlier is not modified. The only statement
that would modify it is not used in this section. The condition $\kappa_1\ge8$ is implied by the floor
already imposed for the change of variables.

\emph{No upper bound on $L$, and the order of choice.} No upper bound on $L$ is used at any point. The
order of choice \eqref{4.24:eq-constant-uniformity-2} is
therefore unchanged.
\end{proof}

\begin{remark}[General compact gauge group]\label{4.24:rem-gauge-group}
Let $N\ge3$ and $G=\mathrm{SU}(N)$, with Lie algebra $\mathfrak{g}$, $n_{\mathfrak{g}}=N^2-1$, and the
$\operatorname{Ad}$-invariant inner product $(\cdot,\cdot)$ on $\mathfrak{g}$. Nothing in the
normalisation and cluster expansion at the survivors' block
(Theorem~\ref{4.24:thm-survivor-block-expansion}) is specific to the value of $N$: apart from the
constants listed below, $n_{\mathfrak{g}}$ enters only through three quantities: the term
$-\tfrac12 n_{\mathfrak{g}}|s|\log\lambda_0$, where $|s|$ is the number of survivors' bonds and
$\lambda_0>0$ is the constant of
Lemma~\ref{4.24:lem-log-resolvent}, the fibre traces $\operatorname{tr}_{\mathfrak{g}}$, and the
number written $O(1)$ in the per-cell count $c_G'$. Suppose that the constants taken from
Ba{\l}aban's papers on which Theorem~\ref{4.24:thm-survivor-block-expansion} and the statement that
opens the present section depend, which those papers give at $\mathrm{SU}(2)$, are re-established at
$\mathrm{SU}(N)$; among them are his covariance constants $B_0$, $\delta_0$ and his constants $a$,
$c_4$, $K_u$, $C_1$, $C_2$, $C_3$ (here $a$ denotes Ba{\l}aban's constant, not the parameter of
$\mathfrak{P}(a)$). Then
Theorem~\ref{4.24:thm-survivor-block-expansion} and the statement that opens the present section
hold verbatim,
as the same implication from the same named statements. The numbers that depend on $(d,L,N)$ are
absorbed by the same upper bounds imposed on the parameters $\gamma$.
\end{remark}

\begin{proof}
Since the space $\bigoplus_{b\in s}\mathfrak{g}$ of the survivors' variables has dimension
$n_{\mathfrak{g}}|s|$, the term $\log\lambda_0\cdot1$ of $F(B)$ in
Lemma~\ref{4.24:lem-log-resolvent} contributes
\begin{equation*}
-\tfrac12\operatorname{Tr}(\log\lambda_0\cdot 1)=-\tfrac12 n_{\mathfrak{g}}|s|\log\lambda_0
\end{equation*}
to $-\tfrac12\operatorname{Tr}F(\mathcal{A}_{ss})$.

The holonomy hypothesis in the statement of the field dependence of the densities in cube gauges
(the statement containing \eqref{4.24:eq-field-dependence-a-bis}; stated; proof incomplete at the step
marked $(\ast)$) is stated for a general $G$ through its centre $Z(G)$; for example,
$Z(\mathrm{SU}(2))=\{\pm1\}$.

Beside that hypothesis the statement exhibits a witness for it; for a walk $\omega$ with holonomy
$H_\omega\in G$ the defect of this witness is
$\tau_\omega(\chi_{\mathrm{adj}}(H_\omega)-n_{\mathfrak{g}})$, with $\tau_\omega$ as in that
statement and $\chi_{\mathrm{adj}}$ the character of the adjoint
representation; the defect vanishes at trivial holonomy, since
$\chi_{\mathrm{adj}}(1)=n_{\mathfrak{g}}$. At $N=2$, let $\varphi$ be the angle of the rotation
$\operatorname{Ad}_{H_\omega}$ of
$\mathfrak{g}\cong\mathbb{R}^3$. Then $\chi_{\mathrm{adj}}(H_\omega)=1+2\cos\varphi$ and
$n_{\mathfrak{g}}=3$,
so the defect reads
\begin{equation*}
\tau_\omega(1+2\cos\varphi-3)=-(2-2\cos\varphi)\tau_\omega .
\end{equation*}

The following ingredients use only that $G$ is compact and that $(\cdot,\cdot)$ is
$\operatorname{Ad}$-invariant:
\begin{itemize}
\item the $\operatorname{Ad}$-invariant form $(\cdot,\cdot)$ on $\mathfrak{g}$;
\item orthogonality of $\operatorname{Ad}_u$ with respect to $(\cdot,\cdot)$, which is the statement of
Lemma~\ref{4.24:lem-kernel-polarisation}(c) and holds also for $u$ in the complexification
$G^{\mathbb{C}}$ of $G$;
\item coercivity of the survivors' block (Lemma~\ref{4.24:lem-block-coercivity});
\item the resolvent representation of the logarithm (Lemma~\ref{4.24:lem-log-resolvent});
\item the square root of the block covariance (Lemma~\ref{4.24:lem-covariance-root});
\item the packing count (Lemma~\ref{4.24:lem-packing-count});
\item the bounds of the Mayer expansion.
\end{itemize}
Each of these therefore holds for every compact $G$ carrying an $\operatorname{Ad}$-invariant inner
product, and in particular for $\mathrm{SU}(N)$.

With the listed constants at $\mathrm{SU}(N)$ every input of
Theorem~\ref{4.24:thm-survivor-block-expansion} is available in the same form.
\end{proof}

\begin{definition}[The arrested region and its collar weight]\label{4.25:not-collar-weight}
The following notation is in force in this section.

\emph{The arrest.} Fix a region $W$ and a step $k_W$ of the renormalisation group. In the
notation of \cite{B-CMP122a}, stage $n_{\mathfrak{a}}+1$ of the large-field operation at step
$k_W$ integrates the level
\begin{equation*}
j_W=h+n_{\mathfrak{a}} ,
\end{equation*}
where $h$ is the integer fixed in \cite{B-CMP122a} for the operation at step $k_W$; and the
region $W$ is excluded at that stage by the cut-off factor $1-\chi'^{(j_W)}$.
The \emph{arrest} is the datum
\begin{equation*}
\mathfrak{a}=(W,k_W,n_{\mathfrak{a}}+1,\iota_{\mathfrak{a}}),
\end{equation*}
where $\iota_{\mathfrak{a}}$ denotes the type of the arrest. The arrest is \emph{live}, that is,
its cut-off factor is still carried, until it is dissolved according to the dissolution rule
for arrests. We write
$\mathfrak{B}_{\mathfrak{a}}:=\mathbb{B}^{(n_{\mathfrak{a}})}_{k_W}$ for the determining
family frozen at the arrest.

\emph{Cells and blocks.} For a level $m$ put $\ell_m:=L^m\eta$, so that $\eta=\ell_0$ is the
length unit of level $0$. The \emph{cells} of
$\mathfrak{a}$ are its host cubes: closed $M$-cubes of the local scale, a cell of level $m$
having side $M\ell_m$, the grids of the different levels being nested $L$-adic grids. For a
cell $c$ we write $\operatorname{lev}c$ for its level. The \emph{blocks} are the cubes
$\Delta(y)$ of side $\ell_y$, each with its enlargement $\tilde\Delta(y)$ as in
\cite{B-CMP99b}; for a set $S$ we write $y\in S$ when $\Delta(y)\subset S$. Two distinct cells
\emph{touch} if their closures meet. The \emph{cell-distance} between two cells is their
distance in the graph whose vertices are the cells and whose edges join touching cells;
$\operatorname{dist}_\infty(A,B)$ denotes the distance between point sets $A,B$ in the
maximum norm.

\emph{The collar variable.} Let $x_m$ denote the coupling parameter of level $m$, so that
$x_m\asymp g_m^{-2}\ge\gamma^{-2}$, and let $\mathfrak{b}'_j=A_1(\log x_{j_W})^{p_1}$ as in
\cite{B-CMP122a}. Let $S^P_W$ be the collar set of $W$, a set of bonds of $W$, and let
$x$ be the collar variable of $W$ on $S^P_W$, a function $b\mapsto x(b)$ on the bonds
$b\in S^P_W$. It is real, and
\begin{equation*}
|x(b)|\le\mathfrak{b}_x:=2.3\,\mathfrak{b}'_j ,
\end{equation*}
the law of $x(b)$ being attached to the site $b_-$ of the bond $b$, in the convention of
\cite{B-CMP96}. Put $\varrho:=6.6\,\mathfrak{b}'_j$ and
\begin{equation*}
P_{\mathfrak{b}}(t):=(6.6A_1)^2(\log t)^{2p_1},
\qquad\text{so that}\qquad
\varrho^2=P_{\mathfrak{b}}(x_{j_W}).
\end{equation*}

\emph{The collar weight.} Let $\mathcal{A}:=C^*\Delta^{(j)}C$ be the kernel of the quadratic
collar form; on the coupling-free regularity class $\mathfrak{K}_G$ of pairs the decomposition
lemma for the collar form, whose proof on that class is outstanding at one step, gives
\begin{equation*}
\mathcal{A}=\mathcal{A}^{w}(\mathbf{1})-a\,C^*C .
\end{equation*}
For a presented gauge field $\mathbb{U}$ and current $\mathbb{J}$ the \emph{quadratic collar
weight} is
\begin{equation*}
\mathsf{q}_W(x;\mathbb{U},\mathbb{J}):=
-\tfrac12\big\langle x,\mathcal{A}_{WW}(\mathbb{U},\mathbb{J})\,x\big\rangle ,
\end{equation*}
where $\mathcal{A}_{WW}(\mathbb{U},\mathbb{J})$ is the restriction of $\mathcal{A}$, taken at
$(\mathbb{U},\mathbb{J})$, to pairs of bonds of $W$, $x$ being extended by zero to the bonds
of $W$ outside $S^P_W$. The weight $\mathsf{q}_W$ is kept in the
exponential under the operation $\mathbf{T}'_{k_W}(W)$ of
\cite{B-CMP122a}.

\emph{The pieces.} For a bond $b\in S^P_W$, the \emph{cubes of $b$} are the cells
containing $b$. The index set $\mathfrak{E}^G_W$ is the union of the set of all
$e^0_a(b,b')$ with $b,b'\in S^P_W$ lying in one pair of blocks, and of the set of all
$e_a(b,b',Y')$ with $b,b'\in S^P_W$ and $Y'\neq\emptyset$ a family of cells, connected in the
graph of touching cells, holding the cubes of $b$ and of $b'$. For $e\in\mathfrak{E}^G_W$:
\begin{itemize}
\item if $e=e_a(b,b',Y')$, let $\tilde S(e)$ be $\bigcup Y'$; if $e=e^0_a(b,b')$, let
$\tilde S(e)$ be the union of the at most two cells meeting the pair of blocks; the bonds of
$\tilde S(e)$ are the bonds with an end in it;
\item let $\mathsf{b}(e)$ be the set of cubes of $b$ and of $b'$ (the \emph{root cubes} of $e$);
\item put $|Y'|:=0$ if $e=e^0_a(b,b')$, and $n(e):=|Y'|+2$;
\item let $S^+(e)$ be the union of $\tilde S(e)$ with the cells at cell-distance at most $2$
from it.
\end{itemize}
Further,
\begin{equation*}
\mathsf{n}_{nb}:=3^dL^{d-1},\qquad
\mathsf{c}_{nb}:=\ln(1+\mathsf{n}_{nb}),\qquad
\mathsf{n}_2:=(1+\mathsf{n}_{nb})^2 ,
\end{equation*}
and $\mathfrak{Q}(\mathfrak{a}):=\{v_e:\ e\in\mathfrak{E}^G_W\}$, where $v_e$ is the piece of
$\mathsf{q}_W$ indexed by $e$ in the decomposition lemma for the collar form.

\emph{Standing conditions.} The blocking factor $L$ of Definition~\ref{1.1:def-lattice}, an
odd integer, satisfies $L\ge3$. The auxiliary
parameter $M$, chosen after $L$, satisfies the one-sided condition
\begin{equation}\label{4.25:eq-collar-weight-b}
M\ge 2L^2 .
\end{equation}

\emph{Properties.} Let the sets $\Omega_m$ be those of \cite{B-CMP122a} in force at the stage
considered; by the freezing lemma for arrests they are frozen on the live pieces of the arrest.
The following hold.
\begin{enumerate}
\item[(a)] For cells $c,c''$ of these sets,
\begin{equation}\label{4.25:eq-collar-weight-a}
\operatorname{lev}c''\ge\operatorname{lev}c+2
\ \Longrightarrow\
\operatorname{dist}_\infty(c,c'')\ge M\ell_{\operatorname{lev}c''-1} ;
\end{equation}
in particular touching cells differ by at most one level, contact at corners included.
\item[(b)] Every cell of these sets is touched by at most $\mathsf{n}_{nb}$ cells.
\item[(c)] For every $e\in\mathfrak{E}^G_W$ the number of cells in $S^+(e)$ satisfies
$|S^+(e)|\le\mathsf{n}_2\,n(e)$.
\item[(d)] For $(\mathbb{U},\mathbb{J})$ in $\mathfrak{K}_G$,
$\mathsf{q}_W=\sum_{e\in\mathfrak{E}^G_W}v_e$ exactly.
\end{enumerate}
\end{definition}

\begin{proof}
(a) The bound \eqref{4.25:eq-collar-weight-a} is the separation property of the sets
$\Omega_m$ of \cite{B-CMP122a}; it is applied only to the sets in force at the stage
considered, which, as stated above, are frozen on the live pieces of the arrest.
Let $m''=\operatorname{lev}c''\ge\operatorname{lev}c+2$. Then $c''\subset\Omega_{m''}$
and $c\subset\Omega^c_{m''-1}$, and between $\Omega_{m''}$ and $\Omega^c_{m''-1}$ there
lies one full layer of $M$-cubes of scale $m''-1$. Each such cube has side $M\ell_{m''-1}$,
so $\operatorname{dist}_\infty(c,c'')\ge M\ell_{m''-1}$, which is
\eqref{4.25:eq-collar-weight-a}. Since $M\ell_{m''-1}>0$, the closures of $c$ and $c''$
are disjoint, so the two cells do not touch, even at a corner. Hence two touching cells
$c,c''$ satisfy $|\operatorname{lev}c-\operatorname{lev}c''|\le1$.

(b) Let $c$ be a cell of level $m$, and suppose first that $m$ is not the lowest level.
Then $c$ is the union of $L^d$ cubes of side $M\ell_{m-1}$ of
the grid of level $m-1$, and the cube concentric with $c$ of side $(L+2)M\ell_{m-1}$ is the
union of $(L+2)^d$ such cubes. By (a) every cell touching $c$ has level at least $m-1$, and
is therefore, the grids being nested, a union of cubes of the grid of level $m-1$. Since its
closure meets $c$, one of these cubes touches $c$; this cube is not contained in $c$, because
distinct cells have disjoint interiors, and for the same reason distinct cells touching $c$
yield distinct such cubes. Every cube of the grid of level $m-1$ touching $c$ lies in the
concentric cube of side $(L+2)M\ell_{m-1}$. Hence the number of cells touching $c$ is at most
$(L+2)^d-L^d$,
the number of cubes of that grid which touch $c$ and are not contained in it. By the
binomial formula,
\begin{equation*}
(L+2)^d-L^d=\sum_{i=1}^{d}\binom{d}{i}2^iL^{d-i} .
\end{equation*}
Since $L\ge1$ we have $L^{d-i}\le L^{d-1}$ for $1\le i\le d$, and
$\sum_{i=1}^{d}\binom{d}{i}2^i=3^d-1$. Therefore
\begin{equation*}
(L+2)^d-L^d\le(3^d-1)L^{d-1}\le 3^dL^{d-1}=\mathsf{n}_{nb} .
\end{equation*}
If $m$ is the lowest level, then by (a) every cell touching $c$ has level $m$ or $m+1$, and
is therefore, the grids being nested, a union of cubes of the grid of level $m$. Since its
closure meets $c$, one of these cubes touches $c$; this cube is not $c$, because distinct
cells have disjoint interiors, and for the same reason distinct cells touching $c$ yield
distinct such cubes. At most $3^d-1$ cubes of the grid of level $m$ touch $c$ and are
different from $c$, and $3^d-1\le3^dL^{d-1}=\mathsf{n}_{nb}$ since $L\ge1$.

(c) By (b) each cell has at most $1+\mathsf{n}_{nb}$ cells at cell-distance at most $1$
(itself and the cells touching it). Each of these again has at most $1+\mathsf{n}_{nb}$
cells at cell-distance at most $1$, so each cell has at most $(1+\mathsf{n}_{nb})^2=\mathsf{n}_2$
cells at cell-distance at most $2$. Summing over the cells of $\tilde S(e)$ gives
$|S^+(e)|\le\mathsf{n}_2|\tilde S(e)|$. If $e=e_a(b,b',Y')$, then $\tilde S(e)=\bigcup Y'$
consists of $|Y'|\le|Y'|+2=n(e)$ cells. If $e=e^0_a(b,b')$, then $\tilde S(e)$ consists of at
most $2=0+2=n(e)$ cells. In both cases $|S^+(e)|\le\mathsf{n}_2\,n(e)$.

(d) This is clause (iii) of the decomposition lemma for the collar form, applied on the
class $\mathfrak{K}_G$; on that class the proof of the lemma is outstanding at one step, the
existence of a finite threshold among the constants of its proof.
\end{proof}

\begin{definition}[Constants and the coupling-free regularity class]\label{4.25:not-regularity-class}
The setting, the arrest $\mathfrak{a}$, the determining family $\mathfrak{B}_{\mathfrak{a}}$, the cells,
the lattice spacing $\eta$ (so that $\ell_m=L^m\eta$),
the blocks $\Delta(y)$ with their stencils $\tilde\Delta(y)$, the collar variable on the set $S^P_W$
with its radius
$\varrho$, the index set $\mathfrak{E}^G_W$ of the pieces $e^0_a(b,b')$ and $e_a(b,b',Y')$, the root
cubes $\mathsf{b}(e)$ (which are cells), the size $|Y'|$ (with $|Y'|:=0$ for $e^0_a$) and the neighbour
counts $\mathsf{n}_{nb}$, $\mathsf{c}_{nb}$, $\mathsf{n}_2$ are those of
Definition~\ref{4.25:not-collar-weight}.

\emph{Constants.} The constants $c_{wt}$, $\mathsf{u}$, $c_*$ and $K_a^{\circ}$ defined below, and the
constants from which they are built, depend only on $d$, $L$, $\kappa$, $\kappa_1$, and they are fixed
before $\gamma$; the sizes $\nu^{\circ}(e)$, $\nu^{\natural}(e)$, $\mathfrak{n}^{\circ\circ}(e)$ and
$\mathfrak{n}^{\circ}(e)$ are functions of the piece $e$, and $\nu^{\natural}(e)$ carries in addition
the factor $\varrho^2$. The constant $\mathsf{c}_{an}$ of the count of connected families of cells is
such that, for every cube $q$ (here and below a cube is a cell of the setting) and
every $m\ge1$, the number of connected families $Y'$ of cells with $q\in Y'$ and $|Y'|=m$ is at most
$e^{\mathsf{c}_{an}m}$. The constants $c_J$ and $c_\flat\ge0$ likewise depend only on $d$, $L$,
$\kappa$, $\kappa_1$ and are fixed before $\gamma$; $c_J$ enters the constant $c_{wt}$ and, together
with $c_\flat$, the exponent in \eqref{4.25:eq-regularity-class-2} below, and $c_\flat$ enters the
floor \eqref{4.25:eq-regularity-class-b}.
Let $\delta>0$ be the decay rate, and $\operatorname{dist}(b,b')$ the distance
between bonds, in the decay bound of the collar kernel $\mathcal{A}$. The constants $B_{\mathcal{A}}$,
$\|C\|$, $C_\rho$, $a$, $n_b$, $n_C$ are those of the bounds on the collar kernel, each read as its
supremum over the class $\mathfrak{K}_G$ defined below, and not over the class $\mathfrak{K}_P^{(k')}$;
that these bounds hold with the suprema taken over $\mathfrak{K}_G$ is part of the input taken from
the bounds on the collar kernel, and the proof of that input is outstanding at one step, the
finiteness of the constant $K'_\Delta$ introduced below. Of these, the proof of
\eqref{4.25:eq-regularity-class-a} uses only the following properties: every cube contains at most
$n_b$ bonds of $S^P_W$; for every bond $b$,
\[
\sum_{b'\in S^P_W}e^{-\delta\operatorname{dist}(b,b')}\le C_\rho ;
\]
for every bond $b$, at most $n_C$ bonds $b'\in S^P_W$ lie in one block pair with $b$. Put
\[
c_{wt}:=\mathsf{n}_2\bigl[c_J+(2\kappa+5)\sqrt d\,\bigr],\qquad
\mathsf{u}:=2\bigl(\mathsf{c}_{an}+\mathsf{c}_{nb}+d\ln L+2\bigr),\qquad
c_*:=c_{wt}+\mathsf{u}+2 .
\]
For the pieces of $\mathfrak{E}^G_W$ put
\begin{equation}\label{4.25:eq-regularity-class-1}
\begin{aligned}
\nu^{\circ}(e_a(b,b',Y'))&:=\tfrac12 B_{\mathcal{A}}\,e^{-\delta\operatorname{dist}(b,b')}\,
e^{-(\kappa_1-1-c_*)|Y'|},\\
\nu^{\circ}(e^0_a(b,b'))&:=\tfrac12 a\|C\|^2 ,
\end{aligned}
\end{equation}
and, for every $e\in\mathfrak{E}^G_W$, $\nu^{\natural}(e):=\varrho^2\,\nu^{\circ}(e)\,e^{-c_*|Y'|}$.
On $\kappa_1$ we impose the floor
\begin{equation}\label{4.25:eq-regularity-class-b}
\kappa_1-1-c_*\ \ge\ \mathsf{c}_{an}+2c_\flat+3 ,
\end{equation}
a one-sided condition imposed in the order of choice after $L$ and $\kappa$, at the place in the
order of choice already held by the earlier floor on $\kappa_1$. Further, put
\begin{equation}\label{4.25:eq-regularity-class-2}
\mathfrak{n}^{\circ\circ}(e):=2\,\nu^{\circ}(e)\,e^{c_J+(c_\flat+1)|\mathsf{b}(e)|},\qquad
\mathfrak{n}^{\circ}(e):=\mathfrak{n}^{\circ\circ}(e)\,e^{-\mathsf{u}|Y'|},
\end{equation}
and
\[
K_a^{\circ}:=2e^{c_J+2c_\flat+2}\,n_b\bigl[C_\rho B_{\mathcal{A}}+n_C\,a\|C\|^2\bigr].
\]
With these constants the pile bound
\begin{equation}\label{4.25:eq-regularity-class-a}
\sup_{q}\ \sum_{e\in\mathfrak{E}^G_W:\ q\in\mathsf{b}(e)}\mathfrak{n}^{\circ\circ}(e)\ \le\ K_a^{\circ}
\end{equation}
holds, the supremum being over all cubes $q$.

\emph{The class.} For a determining family $\mathfrak{B}$ and $\gamma_0>0$, we recall the class
$\mathfrak{K}_{109}(\mathfrak{B};\gamma_0)$ of pairs formed over the determining family
$\mathfrak{B}$: it consists of the pairs $(\mathcal{W},J)$
with $\mathcal{W}=e^{i\eta A'}U$, where $U$ satisfies the conditions
\cite[(3.35)]{B-CMP99b}, $A'$ satisfies \cite[(3.37)]{B-CMP99b}, and $J$ is any
$\mathfrak{g}^{\mathbb{C}}$-valued $1$-form with $\sup\ell^3|J|\le\gamma_0$, $\ell$ being the local
scale ($\ell=\ell_m$ at level $m$). The conditions on $U$ are taken
with radii $\alpha_0$, $\alpha_1$ such that
\[
M\alpha_0\le a_0,\qquad \alpha_1\le\alpha_1^{SL}:=\min\{a_1,c_4\},
\]
where $a_0$, $a_1$ and $c_4$ are the ceilings on the radii under which the conditions on $U$ are
taken in the cited class; they depend only on $(d,L,\kappa_1)$.
Set
\[
\mathfrak{K}_G:=\mathfrak{K}_{109}(\mathfrak{B}_{\mathfrak{a}};\gamma_{reg}),\qquad
\gamma_{reg}:=\min\Bigl\{\frac{1}{2K'_\Delta},\ \frac{1}{6dK_\sharp^2}\Bigr\}.
\]
Here $K_\Delta$ is the constant, depending only on $(d,L,\kappa_1)$, of the convergence statement for
the series defining $\Delta^{(j)}$. That statement says: if $K_\Delta\gamma_0\le\frac12$, the series
converge on the product of the class of pairs at the radius $\gamma_0$ of the hypothesis with the
disc $\{|\sigma|\le e^{\kappa_1}\}$ of the parameter $\sigma$ of
the series, and $\Delta^{(j)}$ is holomorphic,
bilinear and symmetric; moreover, at the pair formed by $U$ and the current $\mathcal{J}(U)$
attached to $U$, and at $\sigma\equiv1$, $\Delta^{(j)}$ is by definition the operator of
Ba{\l}aban's construction. The constant $K'_\Delta$ is the threshold below which the Neumann series
for the individual factors of the expansion of $\Delta^{(j)}$ used in the bounds on the collar kernel
converge in the kernel norms weighted by $e^{-\delta\operatorname{dist}(\cdot,b')}$, factor by factor
and uniformly in the bond $b'$; replacing $K'_\Delta$ by $\max\{K'_\Delta,K_\Delta\}$ if necessary,
we may and do assume $K'_\Delta\ge K_\Delta$, and $K'_\Delta$ is assumed finite and
dependent only on $(d,L,\kappa_1)$. We use the convergence statement with the determining family
$\mathfrak{B}_{\mathfrak{a}}$ and $\gamma_0=\gamma_{reg}$, for which
\[
K_\Delta\gamma_{reg}\le K'_\Delta\gamma_{reg}\le\tfrac12 ,
\]
the hypothesis is met and the conclusion holds on $\mathfrak{K}_G$; and $K'_\Delta\gamma_{reg}\le\frac12$
is the condition under which, once the finiteness of $K'_\Delta$ (the outstanding step named above)
is established, the series converge in the weighted kernel norms just described.
Further, $K_\sharp$ is the constant, depending only on
$(d,L,\kappa_1)$, of the lemma stated together with the class $\mathfrak{K}_{109}$ whose hypothesis
on the radius $\gamma_0$ is met when $6dK_\sharp^2\gamma_0\le1$; the second member of the minimum
gives $6dK_\sharp^2\gamma_{reg}\le1$, so that this hypothesis is met should the first clause of that
lemma be read on $\mathfrak{K}_G$. All radii entering $\mathfrak{K}_G$ depend only on
$(d,L,\kappa_1)$. In particular the radii defining $\mathfrak{K}_G$ do not depend on the coupling;
this is the sense in which $\mathfrak{K}_G$ is called coupling-free. Put further
\[
\varepsilon_{reg}:=\tfrac13\min\{\alpha_1^{SL},10^{-1}\}.
\]
For a region $S$, let $\mathfrak{K}_G(S)$ be the set of restrictions to the bonds of $S$ of the
$P\in\mathfrak{K}_G$ defined on the whole domain. For a level $k'$, let
\[
\mathfrak{K}_P^{(k')}:=\mathfrak{K}_{109}(\mathfrak{B}_{cur};\gamma_0^{(k')}),\qquad
\gamma_0^{(k')}\le\bar\gamma_0(x_{k'}),
\]
where $\mathfrak{B}_{cur}$ is the present determining family and $\bar\gamma_0(\cdot)$ is the
ceiling, depending on the coupling, on the radius $\gamma_0$ that bounds the $1$-form $J$.

For a block $\Delta(y)$ of side $\ell_y$ and a $1$-form $A$, the block size of $A$ is
\begin{equation}\label{4.25:eq-regularity-class-3}
N_y(A):=\sup_{\tilde\Delta(y)}\max\bigl\{\ell_y|A|,\ \ell_y^2|\nabla_UA|\bigr\}.
\end{equation}
Here $U$ is the $G$-valued factor of the base in question, and $\nabla_U$ is the covariant derivative
of \cite[(3.37)]{B-CMP99b}. In the statements from which bounds are taken in what follows, $N_y$ is
read with $\nabla_{\mathcal{W}}$ in place of $\nabla_U$. Once $\ell|A'|\le 1/40$, by the comparison
statement for the covariant derivatives $\nabla_U$ and $\nabla_{\mathcal{W}}$, each of the two readings
of $N_y(A)$ is at most $9/8$ times the other.
\end{definition}

\begin{proof}[Proof of \eqref{4.25:eq-regularity-class-a}]
Fix a cube $q$, and let $e\in\mathfrak{E}^G_W$ with $q\in\mathsf{b}(e)$. The set $\mathsf{b}(e)$
consists of the cubes of $b$ and $b'$, so $|\mathsf{b}(e)|\le 2$. Since $c_\flat\ge0$, we have
$(c_\flat+1)|\mathsf{b}(e)|\le 2c_\flat+2$. By \eqref{4.25:eq-regularity-class-2},
\begin{equation}\label{4.25:eq-regularity-class-4}
\mathfrak{n}^{\circ\circ}(e)\le 2e^{c_J+2c_\flat+2}\,\nu^{\circ}(e).
\end{equation}

The condition $q\in\mathsf{b}(e)$ means that $q$ is the cube of $b$ or the cube of $b'$. We bound
the sum over the pieces with $b\in q$; the sum over those with $b'\in q$ is bounded in the same way,
with the roles of $b$ and $b'$ exchanged. This exchange is legitimate for three reasons:
$\operatorname{dist}(b,b')=\operatorname{dist}(b',b)$; the relation ``$b$ and $b'$ lie in one block
pair'' is symmetric; and $Y'$ holds
the cubes of both bonds.

\emph{Pieces $e^0_a(b,b')$ with $b\in q$.} By \eqref{4.25:eq-regularity-class-1} and
\eqref{4.25:eq-regularity-class-4}, each such piece satisfies
\[
\mathfrak{n}^{\circ\circ}(e)\le e^{c_J+2c_\flat+2}\,a\|C\|^2 .
\]
There are at most $n_b$ bonds $b\in S^P_W$ in $q$, and for each of them at most $n_C$ bonds
$b'\in S^P_W$ lie in one
block pair with $b$. The sum over these pieces is therefore at most
$e^{c_J+2c_\flat+2}\,n_b\,n_C\,a\|C\|^2$.

\emph{Pieces $e_a(b,b',Y')$ with $b\in q$.} By \eqref{4.25:eq-regularity-class-1} and
\eqref{4.25:eq-regularity-class-4}, each such piece satisfies
\[
\mathfrak{n}^{\circ\circ}(e)\le e^{c_J+2c_\flat+2}\,B_{\mathcal{A}}\,e^{-\delta\operatorname{dist}(b,b')}\,
e^{-(\kappa_1-1-c_*)|Y'|}.
\]
By \eqref{4.25:eq-regularity-class-b} and $c_\flat\ge0$,
\[
e^{-(\kappa_1-1-c_*)|Y'|}\le e^{-(\mathsf{c}_{an}+3)|Y'|}.
\]
The family $Y'$ is non-empty and connected, and it holds the cube of $b$. Grouping the families by
$m=|Y'|\ge1$, and using that at most $e^{\mathsf{c}_{an}m}$ of them have size $m$, we obtain
\[
\sum_{Y'}e^{-(\kappa_1-1-c_*)|Y'|}
\le\sum_{m\ge1}e^{\mathsf{c}_{an}m}\,e^{-(\mathsf{c}_{an}+3)m}
=\frac{1}{e^3-1}\le 1 .
\]
Summing next over $b'$ with $\sum_{b'\in S^P_W}e^{-\delta\operatorname{dist}(b,b')}\le C_\rho$, and then
over the at most $n_b$
bonds $b\in S^P_W$ in $q$, shows that the sum over these pieces is at most
$e^{c_J+2c_\flat+2}\,n_b\,C_\rho B_{\mathcal{A}}$.

Adding the two kinds of pieces, the pieces with $b\in q$ contribute at most
\[
e^{c_J+2c_\flat+2}\,n_b\bigl[C_\rho B_{\mathcal{A}}+n_C\,a\|C\|^2\bigr].
\]
The pieces with $b'\in q$ contribute at most the same amount. The total is therefore at most
$K_a^{\circ}$, uniformly in $q$, which is \eqref{4.25:eq-regularity-class-a}.
\end{proof}

\begin{definition}[Coupling-flow inequalities and block thresholds]\label{4.25:not-flow-inequalities}
We work in the setting of Definition~\ref{4.25:not-collar-weight}, with its cells, its blocks
$\Delta(y)$ of side $\ell_y$, their doubled stencils $\tilde\Delta(y)$, and its pieces $e$
with supports $\tilde S(e)$ and their enlargements $S^+(e)$. We write $x_i=g_i^{-2}$ for the
inverse running coupling at level $i$.
In this section the following facts about the sequence of renormalisation steps are used as
inputs. Each is one-sided as written.

\emph{The coupling flow.} For all levels $i\le i'$,
\begin{align}
x_i &\ge x_{i'}-2c_2, \label{4.25:eq-flow-inequalities-a}\\
x_i &\le x_{i'}+B_\sharp\,(i'-i). \label{4.25:eq-flow-inequalities-b}
\end{align}
These two inequalities are inputs from the analysis of the coupling flow of Ba{\l}aban's step,
for the running inverse couplings of Definition~\ref{1.2:def-couplings}.

\emph{Ratio of radii.} Write $\alpha_{*,k}$ for either of the radii $\alpha_{0,k},\alpha_{1,k}$ of
Lemma~\ref{4.4:lem-radii-ratio}; the thresholds
below are built from its envelope $\bar\alpha$, introduced after
\eqref{4.25:eq-flow-inequalities-c}. For all levels $m\le k'$,
\begin{equation}\label{4.25:eq-flow-inequalities-c}
\frac{\alpha_{*,k'}}{\alpha_{*,m}}\le 2\Bigl(1+\frac{B_\sharp\,(k'-m)}{x_{k'}}\Bigr)^{1/2}.
\end{equation}
This is the ratio of radii along the coupling flow, Lemma~\ref{4.4:lem-radii-ratio}.
Moreover there is a function $\bar\alpha$ of the inverse coupling such that $\bar\alpha(x)$ is
$x^{-1/2}$ times a polylogarithmic factor in $x$, and $\alpha_{*,k}\le\bar\alpha(x_k)$ at every
level $k$. Under the condition $q_0=q_1$ on the two exponents $q_0,q_1$ that enter the radii,
imposed where the radii are chosen, $\bar\alpha(x_k)$ is a fixed multiple of $\alpha_{*,k}$, and
the bound \eqref{4.25:eq-flow-inequalities-c} holds also with $\bar\alpha(x_{k'})/\bar\alpha(x_m)$
on its left-hand side.

\emph{The separation integer.} For every level $k$ let $\check{R}_k$ denote the separation
integer of step $k$: it is the integer written $R_k$ in \cite[(3.20)]{B-CMP119} and
\cite[(1.10)]{B-CMP122a} at the same step, so that a separation that Ba{\l}aban states in units
of $R_k$ is here in units of $\check{R}_k$. In terms of the inverse couplings it is given by
$\check{R}_k=R(\min_{i\le k}x_i)$, where $R(\cdot)$ is the separation function, written $R(g')$
in the glossary and read here as a function of $x=g'^{-2}$. It is non-increasing in $k$, and
\begin{equation*}
\bigl(\log(x_k-c_5)\bigr)^r\le\check{R}_k ,
\end{equation*}
where $c_5=2c_2$ and $r$ is the fixed exponent of the separation integer. These facts are those
fixed with the arrest theorem.
Set $E(x):=\exp\bigl[-\bigl(\log(x-c_5)\bigr)^r\bigr]$ for $x>c_5+1$; then
\begin{equation}\label{4.25:eq-flow-inequalities-d}
e^{-\check{R}_k}\le E(x_k),\qquad E(x):=\exp\bigl[-\bigl(\log(x-c_5)\bigr)^r\bigr].
\end{equation}

\emph{The window condition.} Let $N_k$ be the depth of the window of the large-field operation
$\mathbf{R}$ at step $k$: the integer $N=k-h$ of \cite[p.~178]{B-CMP122a}, where $h=k-N$ and the
last $N$ one-step operations are written out there. It is determined by the coupling $g_k$;
\cite[p.~198]{B-CMP122a} states an upper restriction on $N$ in terms of $\log g_k^{-2}$, and
\cite[p.~179]{B-CMP122a} a lower restriction on $N$. The inequality
\begin{equation}\label{4.25:eq-flow-inequalities-e}
B_\sharp\,N_k\le x_k
\end{equation}
is an upper bound on $N_k$; it is the form in which the restriction on $N$ stated in
\cite[p.~198]{B-CMP122a} enters this section. Since $N_k$ is determined by $g_k$, it is read as a
one-sided condition on $\gamma$: the inequality is required for every value $x_k\ge\gamma^{-2}$ of
the inverse coupling, so that it holds at every level; it bounds the coupling from above only, and
it is not a cap in the sense of Definition~\ref{2.3:def-cap}.

\emph{Block thresholds.} Let $y$ be a block whose current level is $\ell'(y)$ and whose weight
is $\mathsf{w}(y)\ge1$. Each space of configurations in use imposes at $y$ smallness conditions
of the form
\begin{equation}\label{4.25:eq-flow-inequalities-1}
Q_a(y)<5\,\mathsf{w}(y)\,r_a(\ell'(y)),\qquad r_a(\ell)\le\bar\alpha(x_\ell).
\end{equation}
Here $Q_a(y)$ is the size of the configuration at the block $y$ that enters the $a$-th
condition, and $r_a(\ell)$ is the corresponding radius at level $\ell$, measured in the units of
level $\ell$; these are the block smallness conditions of the spaces of configurations in use,
compare \cite[(2.34)--(2.39)]{B-CMP119}. The weight is bounded by
\begin{equation}\label{4.25:eq-flow-inequalities-2}
\mathsf{w}(y)\le P_w\bigl(x_{k''(y)}\bigr),\qquad k''(y)\ge\ell'(y),
\end{equation}
where $P_w$ is an increasing polynomial in the logarithm of its argument and $k''(y)$ is the
step at which the condition at $y$ is written. We define
\begin{equation}\label{4.25:eq-flow-inequalities-f}
\begin{gathered}
\theta(y):=5\,\mathsf{w}(y)\,\bar\alpha\bigl(x_{\ell'(y)}\bigr),\qquad
\bar\theta(t):=5\,P_w(t+2c_2)\,\bar\alpha(t),\\
\bar\theta^{\downarrow}(t):=\sup_{t'\ge t}\bar\theta(t'),\qquad
\Theta(e):=\sup_{y\in\mathcal{Y}(e)}\theta(y).
\end{gathered}
\end{equation}
Here $t$ ranges over $[\gamma^{-2},\infty)$, which contains every inverse coupling $x_m$, and
$\mathcal{Y}(e)$ is the set of blocks of the current step that meet the enlarged support
$S^+(e)$ of the piece $e$.

\emph{Nesting.} Along the sequence of renormalisation steps the level of a block never
decreases in time, the block partitions of successive steps are nested, and the doubled stencil
of a block $y_0$ of an earlier step lies in the doubled stencil of the current block $y$
containing it: $\tilde\Delta(y_0)\subset\tilde\Delta(y)$. The block size of a field $A$ at a
block $y$ is
\begin{equation*}
N_y(A)=\sup_{\tilde\Delta(y)}\max\bigl\{\ell_y|A|,\ \ell_y^2|\nabla_UA|\bigr\},
\end{equation*}
where $U$ is the $G$-valued factor, $G$ the gauge group, of the base pair at which the size is
measured. The following hold:
\begin{enumerate}
\item[(i)] $\bar\theta(x_{\ell'(y)})\ge\theta(y)$ for every block $y$;
\item[(ii)] $\bar\theta^{\downarrow}$ is non-increasing and $\bar\theta^{\downarrow}(t)\to0$ as
$t\to\infty$;
\item[(iii)] the block size of a field measured in the units of an old block does not exceed its
block size measured in the units of the current block:
\begin{equation}\label{4.25:eq-flow-inequalities-3}
N_{y_0}(A)\le N_y(A).
\end{equation}
\end{enumerate}
\end{definition}

\begin{proof}
\emph{The bound \eqref{4.25:eq-flow-inequalities-d}.} The map $s\mapsto e^{-s}$ is decreasing.
Applying it to the inequality $(\log(x_k-c_5))^r\le\check{R}_k$ gives
$e^{-\check{R}_k}\le\exp[-(\log(x_k-c_5))^r]$, and the right-hand side is $E(x_k)$ by
definition.

\emph{Claim \textup{(i)}.} Fix a block $y$ and set $i=\ell'(y)$ and $i'=k''(y)$. By
\eqref{4.25:eq-flow-inequalities-2} we have $i\le i'$. Then \eqref{4.25:eq-flow-inequalities-a}
gives $x_{\ell'(y)}\ge x_{k''(y)}-2c_2$, that is,
\begin{equation*}
x_{k''(y)}\le x_{\ell'(y)}+2c_2 .
\end{equation*}
Since $P_w$ is increasing, \eqref{4.25:eq-flow-inequalities-2} yields
\begin{equation*}
\mathsf{w}(y)\le P_w\bigl(x_{k''(y)}\bigr)\le P_w\bigl(x_{\ell'(y)}+2c_2\bigr).
\end{equation*}
Multiply by $5\,\bar\alpha(x_{\ell'(y)})$, which is non-negative. By the definitions
\eqref{4.25:eq-flow-inequalities-f} this gives
\begin{equation*}
\theta(y)=5\,\mathsf{w}(y)\,\bar\alpha\bigl(x_{\ell'(y)}\bigr)
\le5\,P_w\bigl(x_{\ell'(y)}+2c_2\bigr)\,\bar\alpha\bigl(x_{\ell'(y)}\bigr)
=\bar\theta\bigl(x_{\ell'(y)}\bigr).
\end{equation*}

\emph{Claim \textup{(ii)}.} If $t_1\le t_2$, then $\{t'\ge t_2\}\subset\{t'\ge t_1\}$. The
supremum over the smaller set is not larger, so
$\bar\theta^{\downarrow}(t_2)\le\bar\theta^{\downarrow}(t_1)$. Next, $P_w(t+2c_2)$ is a
polynomial in $\log(t+2c_2)$, and $\bar\alpha(t)$ is $t^{-1/2}$ times a polylogarithmic factor
in $t$. Both factors are continuous on $[\gamma^{-2},\infty)$, so $\bar\theta$ is continuous
there; and since every power of $\log(t+2c_2)$ and of $\log t$ grows more slowly than every
positive power of $t$, the factor $t^{-1/2}$ forces $\bar\theta(t)\to0$ as $t\to\infty$. Given
$\varepsilon>0$, choose $T\ge\gamma^{-2}$ such that $\bar\theta(t')\le\varepsilon$ for all
$t'\ge T$. For $\gamma^{-2}\le t<T$,
\begin{equation*}
\bar\theta^{\downarrow}(t)
=\max\Bigl\{\max_{t'\in[t,T]}\bar\theta(t'),\ \sup_{t'\ge T}\bar\theta(t')\Bigr\},
\end{equation*}
where the first maximum is finite by continuity of $\bar\theta$ on the closed interval $[t,T]$
and the second supremum is at most $\varepsilon$; hence $\bar\theta^{\downarrow}(t)$ is finite.
For $t\ge T$ we have $\bar\theta^{\downarrow}(t)\le\varepsilon$. Since $\varepsilon>0$ was
arbitrary, $\bar\theta^{\downarrow}(t)\to0$.

\emph{Claim \textup{(iii)}.} Levels do not decrease in time, so the current block $y$ has level
at least that of $y_0$. Since $\ell_y=\ell_m$ for a block $y$ of level $m$
(Definition~\ref{4.25:not-collar-weight}), and $\ell_m=L^m\eta$, with $\eta$ as there, increases
with $m$, we have $\ell_{y_0}\le\ell_y$. It follows that at every point
\begin{equation*}
\max\bigl\{\ell_{y_0}|A|,\ell_{y_0}^2|\nabla_UA|\bigr\}
\le\max\bigl\{\ell_y|A|,\ell_y^2|\nabla_UA|\bigr\}.
\end{equation*}
The supremum over $\tilde\Delta(y_0)$ of the left-hand side is $N_{y_0}(A)$. Because
$\tilde\Delta(y_0)\subset\tilde\Delta(y)$, this supremum is at most the supremum of the
right-hand side over $\tilde\Delta(y)$, which is $N_y(A)$. This proves
\eqref{4.25:eq-flow-inequalities-3}.
\end{proof}

\begin{lemma}[A word-length distance on the complexified group]\label{4.25:lem-word-length-distance}
Let $G=SU(N)$ be the gauge group, so that every element of $G$ is unitary, let
$G^{\mathbb{C}}\subset M_N(\mathbb{C})$ be its complexification and
$\mathfrak{g}^{\mathbb{C}}$ the Lie algebra of $G^{\mathbb{C}}$. Write $\|\cdot\|$ for the operator norm on
$M_N(\mathbb{C})$. Let $|\cdot|$ be a norm on $M_N(\mathbb{C})$ which satisfies, for all
$X,A,Y\in M_N(\mathbb{C})$,
\begin{equation}\label{4.25:eq-word-length-distance-2}
|XAY|\le\|X\|\,|A|\,\|Y\|,\qquad \|A\|\le|A| .
\end{equation}
For $g\in G^{\mathbb{C}}$ put $\operatorname{pol}(g):=\log\max\{\|g\|,\|g^{-1}\|\}$, and for
$h\in G^{\mathbb{C}}$ put
\begin{equation}\label{4.25:eq-word-length-distance-1}
\operatorname{dis}(h):=\inf\Bigl\{\sum_{i=1}^{r}|Z_i|\;:\;h=e^{Z_1}\cdots e^{Z_r},\
Z_i\in\mathfrak{g}^{\mathbb{C}},\ r\ge1\Bigr\}.
\end{equation}
The set on the right is non-empty, since $G^{\mathbb{C}}$ is connected. For all
$h,h_1,h_2,g\in G^{\mathbb{C}}$:
\begin{equation}\label{4.25:eq-word-length-distance-a}
\begin{aligned}
&\text{(d1)}\quad \|h-1\|\le e^{\operatorname{dis}(h)}-1,
\qquad \operatorname{dis}(h)=0\iff h=1;\\
&\text{(d2)}\quad \operatorname{dis}(h_1h_2)\le\operatorname{dis}(h_1)+\operatorname{dis}(h_2),
\qquad \operatorname{dis}(h^{-1})=\operatorname{dis}(h),\\
&\qquad\quad\ \operatorname{dis}(ghg^{-1})\le e^{2\operatorname{pol}(g)}\operatorname{dis}(h),
\qquad \operatorname{dis}(h)\le|Z|\ \text{whenever } Z\in\mathfrak{g}^{\mathbb{C}},\ e^{Z}=h;\\
&\text{(d3)}\quad \operatorname{pol}(h)\le\operatorname{dis}(h);\\
&\text{(d4)}\quad \|\operatorname{Ad}_h-1\|\le e^{2\operatorname{dis}(h)}-1,
\end{aligned}
\end{equation}
where $\operatorname{Ad}_hB=hBh^{-1}$ and, here and below, for a linear map $T$ of
$(\mathfrak{g}^{\mathbb{C}},|\cdot|)$ into itself, $\|T\|$ denotes its operator norm.
\end{lemma}

\begin{proof}
The set in \eqref{4.25:eq-word-length-distance-1} is non-empty because a connected Lie group is
generated by the image of its exponential map. We use three elementary facts. The first is
\eqref{4.25:eq-word-length-distance-2}. The second: for $Z\in\mathfrak{g}^{\mathbb{C}}$ one has
$\|e^{\pm Z}\|\le e^{\|Z\|}\le e^{|Z|}$, hence $\operatorname{pol}(e^{Z})\le|Z|$; and $\operatorname{pol}$ is
sub-additive, because $\|g_1g_2\|\le\|g_1\|\|g_2\|$ and
$\|(g_1g_2)^{-1}\|\le\|g_2^{-1}\|\|g_1^{-1}\|$, so that
$\max\{\|g_1g_2\|,\|(g_1g_2)^{-1}\|\}$ is at most the product of the two maxima. The third: for
$g\in G^{\mathbb{C}}$ and $Z\in\mathfrak{g}^{\mathbb{C}}$, by \eqref{4.25:eq-word-length-distance-2},
\begin{equation}\label{4.25:eq-word-length-distance-3}
|\operatorname{Ad}_gZ|\le\|g\|\,|Z|\,\|g^{-1}\|\le e^{2\operatorname{pol}(g)}|Z| ,
\end{equation}
that is, $\|\operatorname{Ad}_g\|\le e^{2\operatorname{pol}(g)}$.

\emph{(d1).} For $Z\in M_N(\mathbb{C})$ the power series gives
$\|e^{Z}-1\|\le\sum_{n\ge1}\|Z\|^n/n!=e^{\|Z\|}-1$. If $\|a-1\|\le e^{\alpha}-1$ and
$\|b-1\|\le e^{\beta}-1$, the identity $ab-1=(a-1)(b-1)+(a-1)+(b-1)$ gives
\begin{equation*}
\|ab-1\|\le(e^{\alpha}-1)(e^{\beta}-1)+(e^{\alpha}-1)+(e^{\beta}-1)=e^{\alpha+\beta}-1 .
\end{equation*}
By induction on $r$,
\begin{equation*}
\|e^{Z_1}\cdots e^{Z_r}-1\|\le e^{\sum_i\|Z_i\|}-1\le e^{\sum_i|Z_i|}-1,
\end{equation*}
the last step by $\|Z_i\|\le|Z_i|$. This holds for every presentation of $h$ in
\eqref{4.25:eq-word-length-distance-1}; since $t\mapsto e^{t}-1$ is increasing and continuous, taking
the infimum gives $\|h-1\|\le e^{\operatorname{dis}(h)}-1$. If $\operatorname{dis}(h)=0$, then $\|h-1\|\le0$,
so $h=1$. Conversely $1=e^{0}$ gives $0\le\operatorname{dis}(1)\le0$.

\emph{(d2).} If $h_1=e^{Z_1}\cdots e^{Z_r}$ and $h_2=e^{Z'_1}\cdots e^{Z'_s}$, the concatenated word
presents $h_1h_2$ with sum $\sum_i|Z_i|+\sum_j|Z'_j|$; taking infima over both presentations gives
sub-additivity. If $h=e^{Z_1}\cdots e^{Z_r}$, then $h^{-1}=e^{-Z_r}\cdots e^{-Z_1}$ with the same sum,
so $\operatorname{dis}(h^{-1})\le\operatorname{dis}(h)$; applied to $h^{-1}$ this gives equality. For
conjugation, $ge^{Z}g^{-1}=e^{\operatorname{Ad}_gZ}$ by the power series, and
$\operatorname{Ad}_gZ\in\mathfrak{g}^{\mathbb{C}}$. Hence
$ghg^{-1}=e^{\operatorname{Ad}_gZ_1}\cdots e^{\operatorname{Ad}_gZ_r}$, and by
\eqref{4.25:eq-word-length-distance-3} the sum of this presentation is at most
$e^{2\operatorname{pol}(g)}\sum_i|Z_i|$; the infimum gives the bound. Finally, if $e^{Z}=h$ with
$Z\in\mathfrak{g}^{\mathbb{C}}$, this is a presentation with $r=1$, so $\operatorname{dis}(h)\le|Z|$.

\emph{(d3).} For $h=e^{Z_1}\cdots e^{Z_r}$, the sub-additivity of $\operatorname{pol}$ and
$\operatorname{pol}(e^{Z_i})\le|Z_i|$ give $\operatorname{pol}(h)\le\sum_i|Z_i|$; take the infimum.

\emph{(d4).} Let $Z,B\in\mathfrak{g}^{\mathbb{C}}$. Since
$\frac{d}{dt}\bigl(e^{tZ}Be^{-tZ}\bigr)=e^{tZ}[Z,B]e^{-tZ}$,
\begin{equation*}
(\operatorname{Ad}_{e^{Z}}-1)B=\int_0^1e^{tZ}[Z,B]e^{-tZ}\,dt .
\end{equation*}
By \eqref{4.25:eq-word-length-distance-2}, $|ZB|\le\|Z\||B|$ and $|BZ|\le|B|\|Z\|$, so
$|[Z,B]|\le2\|Z\||B|$. Also $\|e^{\pm tZ}\|\le e^{t\|Z\|}$. Hence
\begin{equation*}
|(\operatorname{Ad}_{e^{Z}}-1)B|\le\int_0^1e^{2t\|Z\|}\,2\|Z\|\,|B|\,dt
=\bigl(e^{2\|Z\|}-1\bigr)|B| .
\end{equation*}
By \eqref{4.25:eq-word-length-distance-2}, moreover,
$\|\operatorname{Ad}_{e^{Z}}\|\le\|e^{Z}\|\|e^{-Z}\|\le e^{2\|Z\|}$.

We prove by induction on $r$ that, for $h=e^{Z_1}\cdots e^{Z_r}$,
$\|\operatorname{Ad}_h-1\|\le e^{2\sum_i\|Z_i\|}-1$. The case $r=1$ is the bound just shown. For $r\ge2$,
put $g_1=e^{Z_1}\cdots e^{Z_{r-1}}$, $g_2=e^{Z_r}$, $a=\sum_{i<r}\|Z_i\|$ and $b=\|Z_r\|$. Since
$\operatorname{Ad}$ is multiplicative, $\|\operatorname{Ad}_{g_1}\|\le\prod_{i<r}e^{2\|Z_i\|}=e^{2a}$. The
identity
$\operatorname{Ad}_{g_1g_2}-1=\operatorname{Ad}_{g_1}(\operatorname{Ad}_{g_2}-1)+(\operatorname{Ad}_{g_1}-1)$,
together with the induction hypothesis for $g_1$ and the case $r=1$ for $g_2$, gives
\begin{equation*}
\|\operatorname{Ad}_h-1\|\le e^{2a}\bigl(e^{2b}-1\bigr)+\bigl(e^{2a}-1\bigr)=e^{2(a+b)}-1 .
\end{equation*}
Using $\|Z_i\|\le|Z_i|$, we get $\|\operatorname{Ad}_h-1\|\le e^{2\sum_i|Z_i|}-1$ for every presentation
of $h$. Taking the infimum, as in (d1), gives (d4).
\end{proof}

Since every element of $G$ is unitary, $\|g\|=\|g^{-1}\|=1$ for $g\in G$, so $\operatorname{pol}$
vanishes identically on $G$ and measures the departure of an element of $G^{\mathbb{C}}$ from $G$;
in what follows $\operatorname{pol}$ enters only through the bound
$\|\operatorname{Ad}_g\|\le e^{2\operatorname{pol}(g)}$ of \eqref{4.25:eq-word-length-distance-3}.

\begin{lemma}[The cut-off of a piece: plateau, support, slope]\label{4.25:lem-piece-cut-off}
We use the cells, their levels, the blocks $\Delta(y)$ of side $\ell_y$ with their doubled stencils
$\tilde\Delta(y)$, the sets $\tilde S(e)$ and $S^+(e)$ attached to a piece $e\in\mathfrak{E}^G_W$,
the bonds of $\tilde S(e)$ (those with an end in $\tilde S(e)$), and the touching relation and the
cell-distance, all as in Definition~\ref{4.25:not-collar-weight}; of the doubled stencil we use only
that each of its bonds has its initial site within $\ell^\infty$-distance $2\ell_y$ of $\Delta(y)$.
For a cell $C$ we write
$\operatorname{lev}C$ for its level and $\ell_C:=\ell_{\operatorname{lev}C}$, so that $C$ is a closed
cube of side $M\ell_C$, and $\operatorname{dist}_\infty$ denotes $\ell^\infty$-distance in
$\mathbb{R}^d$. For a cell $C$ and $z\in\mathbb{R}^d$ put
\begin{equation}\label{4.25:eq-piece-cut-off-1}
\psi_C(z):=\min\Bigl\{1,\Bigl(2-\frac{2L\operatorname{dist}_\infty(z,C)}{M\ell_C}\Bigr)^{+}\Bigr\},
\qquad
\chi_e:=\max_{C\subset\tilde S(e)}\psi_C ,
\end{equation}
the maximum running over the cells contained in $\tilde S(e)$; on a bond $b$ the function $\chi_e$ is
read at the initial site $b_-$ of $b$, that is, $\chi_e(b):=\chi_e(b_-)$. Assume the level condition
\eqref{4.25:eq-collar-weight-a} and the floor \eqref{4.25:eq-collar-weight-b}. Then, for every cell $C$,
every point $z$ and every closed cell $c$ containing $z$ in $(\chi0)$; for every block $y$ and every
pair of lattice neighbours $z,z'$ with $z$ within $\ell^\infty$-distance $2\ell_y$ of $\Delta(y)$ in
$(\chi3)$; and for every $1$-form $A$, every connection $V$ and every block $y$ in $(\chi4)$, with
the suprema taken over the bonds of $\tilde\Delta(y)$ and $\nabla_V$ the covariant difference along
$V$:
\begin{equation}\label{4.25:eq-piece-cut-off-a}
\begin{aligned}
&(\chi0)\quad \operatorname{dist}_\infty(z,C)<M\ell_C/L\ \Longrightarrow\
 c=C\ \text{or}\ c\ \text{touches}\ C;\\
&(\chi1)\quad \chi_e(b)=1\ \text{for every bond $b$ of}\ \tilde S(e);\\
&(\chi2)\quad \chi_e=0\ \text{on every cell at cell-distance}\ \ge2\ \text{from}\ \tilde S(e);\\
&(\chi3)\quad |\chi_e(z)-\chi_e(z')|\le \eta/\ell_y;\\
&(\chi4)\quad \sup_{\tilde\Delta(y)}\max\bigl\{\ell_y|\chi_eA|,\ \ell_y^2|\nabla_V(\chi_eA)|\bigr\}
 \le 2\sup_{\tilde\Delta(y)}\max\bigl\{\ell_y|A|,\ \ell_y^2|\nabla_VA|\bigr\},\\
&\phantom{(\chi4)}\quad \text{and the left-hand side of }(\chi4)\text{ is }0\text{ unless }y\in S^+(e).
\end{aligned}
\end{equation}
\end{lemma}

\begin{proof}
Throughout, \eqref{4.25:eq-collar-weight-a} is used in the form: if two cells $c,c''$ satisfy
$\operatorname{lev}c''\ge\operatorname{lev}c+2$, then
$\operatorname{dist}_\infty(c,c'')\ge M\ell_{\operatorname{lev}c''-1}$; and
\eqref{4.25:eq-collar-weight-b} is the floor $M\ge2L^2$, with $L\ge3$. We also note from
\eqref{4.25:eq-piece-cut-off-1} that $0\le\psi_C\le1$, that $\psi_C(z)>0$ if and only if
$\operatorname{dist}_\infty(z,C)<M\ell_C/L$, and that $\psi_C(z)=1$ if and only if
$\operatorname{dist}_\infty(z,C)\le M\ell_C/(2L)$.

\emph{Proof of $(\chi0)$.} Write $C=\prod_i[a_i,a_i+s]$ with $s:=M\ell_C$. The grids carrying the cells
are nested $L$-adic grids, so the grid of mesh $s/L$ refining the grid of $C$ has, in the $i$-th
coordinate, the points $a_i+(s/L)\mathbb{Z}$. Let $\gamma=\prod_i[g_i,g_i+s/L]$ be a closed cube of
this grid with $z\in\gamma$. The hypothesis $\operatorname{dist}_\infty(z,C)<s/L$ gives, for every $i$,
\begin{equation*}
a_i-s/L<z_i<a_i+s+s/L ,
\end{equation*}
and $z\in\gamma$ gives $g_i\le z_i\le g_i+s/L$; hence $a_i-2s/L<g_i<a_i+s+s/L$, and since $g_i$ is a
grid point, $a_i-s/L\le g_i\le a_i+s$. Thus in every coordinate the interval $[g_i,g_i+s/L]$ either lies
in $[a_i,a_i+s]$ or meets it in an endpoint, so $\gamma\cap C\neq\emptyset$: either $\gamma\subset C$
or $\gamma$ touches $C$.

Now let $c$ be a closed cell containing $z$. The cubes of the mesh-$(s/L)$ grid that contain $z$ cover
a neighbourhood of $z$, and $z$ lies in the closure of the interior of $c$; so we may take such a cube
$\gamma\ni z$ whose interior meets the interior of $c$. Since the grids are nested, either
$c\supseteq\gamma$ or $c\subsetneq\gamma$. If $c\supseteq\gamma$, then $c$ meets $C$ because $\gamma$
does, so $c=C$ or $c$ touches $C$. If $\gamma\subset C$ then, $C$ being a cell and the grids nested,
no cell lies strictly inside $\gamma$, so $c\supseteq\gamma$. There remains the case
$c\subsetneq\gamma\not\subset C$. Then, the grids being
nested, $\gamma$ is a union of cells of level at most $\operatorname{lev}C-2$, and since the closed
cube $\gamma$ meets $C$, one of these cells, say $c'$, meets $C$, i.e.\ touches it. Then
$\operatorname{lev}C\ge\operatorname{lev}c'+2$ while
$\operatorname{dist}_\infty(c',C)=0<M\ell_{\operatorname{lev}C-1}$, against
\eqref{4.25:eq-collar-weight-a} applied to $(c,c'')=(c',C)$. So this case does not occur, and
$(\chi0)$ is proved.

\emph{Proof of $(\chi1)$.} By the remark above, the plateau $\{\psi_C=1\}$ is the set of points at
$\ell^\infty$-distance at most $M\ell_C/(2L)$ from $C$, and by \eqref{4.25:eq-collar-weight-b}
\begin{equation*}
\frac{M\ell_C}{2L}\ge L\ell_C\ge\eta ,
\end{equation*}
since $\ell_C=L^{\operatorname{lev}C}\eta\ge\eta$. Let $b$ be a bond of $\tilde S(e)$. One end of $b$
lies in $\tilde S(e)$, which is a union of cells, hence in some cell $C\subset\tilde S(e)$; the initial
site $b_-$ is that end or lies at distance $\eta$ from it, so $\operatorname{dist}_\infty(b_-,C)\le\eta$
and $\psi_C(b_-)=1$. As every $\psi_{C'}\le1$, we get $\chi_e(b)=\chi_e(b_-)=1$.

\emph{Proof of $(\chi2)$.} Let $z$ lie in a cell $c$ and suppose $\chi_e(z)>0$. Then $\psi_C(z)>0$ for
some cell $C\subset\tilde S(e)$, that is, $\operatorname{dist}_\infty(z,C)<M\ell_C/L$, and by $(\chi0)$
the cell $c$ equals $C$ or touches $C$. Hence $c$ is at cell-distance at most $1$ from $\tilde S(e)$.
Consequently $\chi_e$ vanishes on every cell at cell-distance $\ge2$ from $\tilde S(e)$.

\emph{Proof of $(\chi3)$.} Since the maxima over $C$ of two families of real numbers differ by at most
the maximum over $C$ of the differences, and cells with
$\psi_C(z)=\psi_C(z')=0$ contribute $0$, we have $|\chi_e(z)-\chi_e(z')|\le\max|\psi_C(z)-\psi_C(z')|$,
the maximum over the cells $C\subset\tilde S(e)$ with $\psi_C(z)+\psi_C(z')>0$ (the difference is $0$
if there is none). The function $\min\{1,(2-\cdot)^+\}$ on $\mathbb{R}$ is $1$-Lipschitz and
$\operatorname{dist}_\infty(\cdot,C)$ is $1$-Lipschitz for the $\ell^\infty$-norm, while lattice
neighbours satisfy $|z-z'|_\infty=\eta$; therefore
\begin{equation*}
|\psi_C(z)-\psi_C(z')|\le\frac{2L}{M\ell_C}\,|z-z'|_\infty=\frac{2L\eta}{M\ell_C}.
\end{equation*}
Let $c_y$ be the cell of $y$, so that $\Delta(y)\subset c_y$ and $\ell_{c_y}=\ell_y$. We claim that
every $C$ with $\psi_C(z)+\psi_C(z')>0$ has $\operatorname{lev}C\ge\operatorname{lev}c_y-1$. Otherwise
$\operatorname{lev}c_y\ge\operatorname{lev}C+2$, so $\ell_C\le\ell_y/L^2$, and
\eqref{4.25:eq-collar-weight-a} gives
$\operatorname{dist}_\infty(c_y,C)\ge M\ell_{\operatorname{lev}c_y-1}=M\ell_y/L$. On the other hand one
of the two points $z,z'$, say $z''$, satisfies $\operatorname{dist}_\infty(z'',C)<M\ell_C/L\le M\ell_y/L^3$,
and $z''$ lies within $2\ell_y+\eta\le3\ell_y$ of $\Delta(y)\subset c_y$. Hence
\begin{equation*}
\frac{M\ell_y}{L}\le\operatorname{dist}_\infty(c_y,C)\le3\ell_y+\frac{M\ell_y}{L^3},
\qquad\text{i.e.}\qquad M\le\frac{3L^3}{L^2-1}<2L^2 ,
\end{equation*}
where the last inequality is $3L<2L^2-2$, i.e.\ $(2L+1)(L-2)>0$, true for $L\ge3$. This contradicts
\eqref{4.25:eq-collar-weight-b}, and the claim follows. So $\ell_C\ge\ell_y/L$ for all such $C$, and
by \eqref{4.25:eq-collar-weight-b}
\begin{equation*}
|\chi_e(z)-\chi_e(z')|\le\frac{2L^2\eta}{M\ell_y}\le\frac{\eta}{\ell_y}.
\end{equation*}

\emph{Proof of $(\chi4)$.} Write $\chi=\chi_e$, let $b$ be a bond of $\tilde\Delta(y)$, $\mu$ a
direction, and $b':=b+\eta e_\mu$, with $e_\mu$ the unit vector in the direction $\mu$, the shifted
bond. Since $\chi$ is real-valued and $\nabla_{V,\mu}$ is
the difference quotient of the parallel-transported value of $A$ at $b'$ and $A(b)$, inserting
$\pm\chi(b')A(b)$ gives
\begin{equation}\label{4.25:eq-piece-cut-off-2}
\nabla_{V,\mu}(\chi A)(b)=\chi(b')\,\nabla_{V,\mu}A(b)+\eta^{-1}\bigl(\chi(b')-\chi(b)\bigr)A(b),
\end{equation}
the second member carrying no adjoint action. The initial sites of $b$ and $b'$ are lattice neighbours,
that of $b$ within $2\ell_y$ of $\Delta(y)$, so $(\chi3)$ gives $|\chi(b')-\chi(b)|\le\eta/\ell_y$;
with $0\le\chi\le1$, \eqref{4.25:eq-piece-cut-off-2} yields
\begin{equation*}
\ell_y^2|\nabla_V(\chi A)(b)|\le\ell_y^2|\nabla_VA(b)|+\ell_y|A(b)|
\le2\max\bigl\{\ell_y|A(b)|,\ \ell_y^2|\nabla_VA(b)|\bigr\},
\end{equation*}
and trivially $\ell_y|\chi A(b)|\le\ell_y|A(b)|$. Taking suprema over $\tilde\Delta(y)$ gives the
inequality in $(\chi4)$.

For the vanishing, suppose the left-hand side of $(\chi4)$ is non-zero. By
\eqref{4.25:eq-piece-cut-off-2} there is a bond $b$ of $\tilde\Delta(y)$ such that $\chi_e>0$ at $b$ or
at its shift $b'$; let $z''$ be the corresponding initial site. It lies within $2\ell_y+\eta\le3\ell_y$
of $\Delta(y)$, hence of $c_y\supset\Delta(y)$, and $3\ell_y<2L\ell_y\le M\ell_y/L$ by
\eqref{4.25:eq-collar-weight-b}; so
$\operatorname{dist}_\infty(z'',c_y)<M\ell_{c_y}/L$, and by $(\chi0)$ the cell containing $z''$ equals
$c_y$ or touches $c_y$. Since $\chi_e(z'')>0$, by $(\chi2)$ that cell is at cell-distance at most $1$
from $\tilde S(e)$. Hence $c_y$ is at cell-distance at most $2$ from $\tilde S(e)$, so
$\Delta(y)\subset c_y\subset S^+(e)$, that is, $y\in S^+(e)$.
\end{proof}

\begin{lemma}[Gaussian pieces: domain, invariance and continuation. Stated; proof incomplete at the
step marked $(\ast)$]\label{4.25:lem-gaussian-pieces}
The following notation is used.
\begin{itemize}
\item The arrest $\mathfrak{a}$, its frozen family $\mathfrak{B}_{\mathfrak{a}}$, the cells, and the blocks
$\Delta(y)$ of side $\ell_y$ with their stencils $\tilde\Delta(y)$ are as in
Definition~\ref{4.25:not-collar-weight}.
\item The constants, the coupling-free class $\mathfrak{K}_G=\mathfrak{K}_{109}(\mathfrak{B}_{\mathfrak{a}};\gamma_{reg})$
with its radii $\gamma_{reg}$, $\varepsilon_{reg}$, $\alpha_1^{SL}$, the restricted classes
$\mathfrak{K}_G(S)$, the coupling-small classes $\mathfrak{K}_P^{(k')}$ with their current radii
$\gamma_0^{(k')}$, and the block sizes $N_y$ are as in Definition~\ref{4.25:not-regularity-class}.
\item In a current condition $\ell^3|J|$ the letter $\ell$ is the local unit at the bond, as in the
current clause of $\mathfrak{K}_{109}$ in Definition~\ref{4.25:not-regularity-class}; on the stencil
of a block $y$ it is written $\ell_y$.
\item The functions $v_e$, $e\in\mathfrak{E}^G_W$, are the pieces of the quadratic collar weight
$\mathsf{q}_W$, with supports $\tilde S(e)$ and neighbourhoods $S^+(e)$ as in
Definition~\ref{4.25:not-collar-weight}, and with cut-offs $\chi_e$ as in
Lemma~\ref{4.25:lem-piece-cut-off}.
\item The collar variable $x$ ranges over the box $\mathfrak{b}_x$ of
Definition~\ref{4.25:not-collar-weight}, $\varrho$ is the radius of that box, and
$\sigma'=(\sigma'(\Delta))_{\Delta}$ denotes the decoupling parameters, with components $\sigma'(\Delta)$.
\item A pair is \emph{global} if it is defined on all bonds of the domain. For a global pair $P$ and a
region $S$, $P|_S$ denotes its restriction to the bonds of $S$.
\item For a pair $p=(\mathbb{U},\mathbb{J})$ and a gauge transformation $u$ we put
$p^u:=(\mathbb{U}^u,\operatorname{Ad}_u\mathbb{J})$.
\item For a bond $b$, $b_-$ is its initial site.
\item For an assignment $g$ of an element of $G^{\mathbb{C}}$ to every site we put
$x^{[g]}(b):=\operatorname{Ad}_{g(b_-)}x(b)$.
\end{itemize}
Then the following hold.
\begin{enumerate}
\item[(i)] Each $v_e$ is a well-defined function on
$\mathfrak{K}_G(\tilde S(e))\times\{\sigma':|\sigma'(\Delta)|\le e^{\kappa_1}\text{ for every }\Delta\}$.
It is holomorphic along holomorphic
families of global pairs in $\mathfrak{K}_G$. It is jointly $\mathcal{G}$-invariant:
\begin{equation}\label{4.25:eq-gaussian-pieces-1}
v_e\bigl(x^{[u]};\mathbb{U}^u,\operatorname{Ad}_u\mathbb{J}\bigr)=v_e(x;\mathbb{U},\mathbb{J})
\qquad\text{for every $G$-valued $u$.}
\end{equation}
It is a quadratic form, $v_e(x)=Q_e(x,x)$ with $Q_e$ bilinear, and on $\mathfrak{K}_G(\tilde S(e))$
it satisfies, for all collar vectors $x_1,x_2$,
\begin{equation}\label{4.25:eq-gaussian-pieces-2}
|Q_e(x_1,x_2)|\le \nu^{\natural}(e)\,|x_1|_\infty|x_2|_\infty/\varrho^2 .
\end{equation}
\item[(ii)] (The orbit datum is entire.) Put
\begin{equation}\label{4.25:eq-gaussian-pieces-3}
v_e^{\sharp}(x;p;g):=v_e\bigl(x^{[g]};p\bigr).
\end{equation}
Then $v_e^{\sharp}$ is a polynomial in the entries of $\operatorname{Ad}_g$ and is defined for every
$g\in(G^{\mathbb{C}})^{\text{sites}}$. At $G$-valued $g$ it equals $v_e(x;p^{g^{-1}})$.
\item[(iii)] (Containment.) Let $\mathfrak{O}$ be a later operation, taken at a step $k'$, and let
$\mathfrak{T}$ be a space that $\mathfrak{O}$ presents over a region $Y\supseteq S^+(e)$; every member
of $\mathfrak{T}$ is the restriction of a global pair in $\mathfrak{K}_P^{(k')}$ that satisfies the
smallness conditions which the presentation of $\mathfrak{T}$ imposes on the enlargement $Y^+$ of $Y$
(in what follows, the presentation conditions on $Y^+$). Let $\tilde P\in\mathfrak{K}_P^{(k')}$ be a
global pair that satisfies the presentation conditions on $Y^+$ and whose restriction is a member of
$\mathfrak{T}$. Write $A'_{\tilde P}$ and
$\mathbb{J}_{\tilde P}$ for the one-form and the current of $\tilde P$, let $\mathcal{Y}(e)$ be the set
of current blocks meeting $S^+(e)$, and let $\theta(y)$ be as in
\eqref{4.25:eq-flow-inequalities-f}. Then $\tilde P\in\mathfrak{K}_G$,
\begin{equation}\label{4.25:eq-gaussian-pieces-a}
N_y(A'_{\tilde P})<\theta(y)\qquad\text{for } y\in\mathcal{Y}(e),
\end{equation}
and $\ell^3|\mathbb{J}_{\tilde P}|\le\gamma_0^{(k')}$ everywhere.
\item[(iv)] (Cut-off continuation.) Let $P=(e^{i\eta A'}U,\mathbb{J})\in\mathfrak{K}_G$ be global. Assume
that $N_y(A')\le\varepsilon_{reg}/5$ and $\ell_y^3|\mathbb{J}|\le\frac12\gamma_{reg}$ on $S^+(e)$, and
write $P^U:=e^{i\eta A'}U$ for its gauge-field entry. Let $A''$, $J''$ be
$\mathfrak{g}^{\mathbb{C}}$-valued $1$-forms on the bonds of the stencils $\tilde\Delta(y)$,
$y\in S^+(e)$, with
\begin{equation}\label{4.25:eq-gaussian-pieces-4}
\mathfrak{a}'':=\sup_{S^+(e)}N_y(A'')\le\tfrac15\varepsilon_{reg},\qquad
\mathfrak{j}'':=\sup_{S^+(e)}\ell^3|J''|\le\tfrac14\gamma_{reg}.
\end{equation}
Put
\begin{equation}\label{4.25:eq-gaussian-pieces-b}
\begin{gathered}
Q(\tau,t):=\bigl(e^{i\eta\tau\chi_eA''}e^{i\eta A'}U,\ \mathbb{J}+t\chi_eJ''\bigr),\\
T_U:=\varepsilon_{reg}/\mathfrak{a}''\ (\ge 5),\qquad
T_J:=\tfrac12\gamma_{reg}/\mathfrak{j}''\ (\ge 2).
\end{gathered}
\end{equation}
Then $Q(\tau,t)$ is global and entire in $(\tau,t)$. It lies in $\mathfrak{K}_G$ for $|\tau|<T_U$,
$|t|<T_J$. On the bonds of $\tilde S(e)$ it equals $(e^{i\eta\tau A''}P^U,\mathbb{J}+tJ'')$. The
function $(\tau,t)\mapsto v_e(x;Q(\tau,t)|_{\tilde S(e)})$ is holomorphic on the bidisc
$\{|\tau|<T_U\}\times\{|t|<T_J\}$, and
\begin{equation}\label{4.25:eq-gaussian-pieces-5}
p_2:=\bigl(e^{i\eta A''}P^U,\ \mathbb{J}+J''\bigr)\big|_{\tilde S(e)}\in\mathfrak{K}_G(\tilde S(e)).
\end{equation}
\item[(iv)$'$] If $\mathfrak{a}''\le\varepsilon_{reg}/20$ and $N_y(A')\le\varepsilon_{reg}/20$ on
$S^+(e)$, and if
$\ell^3|\mathbb{J}|\le\frac14\gamma_{reg}$ on $S^+(e)$ and $\mathfrak{j}''\le\frac14\gamma_{reg}$, then
$Q(1,1)$ is again a pair $P$ satisfying the hypotheses of \textup{(iv)}.
\end{enumerate}
\end{lemma}

\begin{proof}
\emph{Proof of (i).} The decomposition lemma for the collar weight, which defines the pieces
$v_e$ and proves their properties, is stated on $\mathfrak{K}_{109}(\mathfrak{B}_{\mathfrak{a}};\gamma_0)$.
It is stated under the floor condition relating $M$ and $\kappa_1$, under its condition on the levels,
and under $K_\Delta\gamma_0\le\frac12$, where $K_\Delta$ is the $(d,L,\kappa_1)$-number in the
hypothesis of the convergence statement for the walk series of the operator $\Delta^{(j)}$ described
below (Definition~\ref{4.25:not-regularity-class}).
Its radius $\gamma_0$ is the coupling-small radius of the classes $\mathfrak{K}_P^{(k')}$, and
$\gamma_0\le\bar\gamma_0(x_{k'})\le\frac14\gamma_{reg}$; here $\bar\gamma_0(x_{k'})$ is the
coupling-dependent ceiling on that radius fixed in Definition~\ref{4.25:not-regularity-class}, and the
second inequality is a one-sided condition on the coupling $x_{k'}$ in force at step $k'$. At equal
partition and equal conditions on the gauge field, the current
clause of $\mathfrak{K}_{109}(\mathfrak{B};\gamma)$ admits every $\mathfrak{g}^{\mathbb{C}}$-valued
one-form $J$ with $\sup\ell^3|J|\le\gamma$, so these classes grow with $\gamma$. Hence
$\mathfrak{K}_{109}(\mathfrak{B}_{\mathfrak{a}};\gamma_0)$ is the smaller of the two classes, no
containment of $\mathfrak{K}_G$ in it is available, and the statement of the decomposition lemma does
not cover $\mathfrak{K}_G$.

Its proof uses four inputs. The first is the convergence statement for the walk series of the operator
$\Delta^{(j)}$; its only hypothesis on the radius $\gamma$ of the current is $K_\Delta\gamma\le\frac12$,
and its series are expansions around an operator supplied by a further lemma, which asks nothing of
$\gamma$.
The second is the localisation statement for the walk-expanded kernel, which carries no hypothesis on
$\gamma$ at all; it holds under the same floor condition relating $M$ and $\kappa_1$ (a lower bound on
$M$ proportional to $\kappa_1$) and for $|\sigma'(\Delta)|\le e^{\kappa_1}$. The
third consists of the two bounds on the walk expansion of the kernel. The fourth is the bound on the
operator $C=C(\mathbb{U})$.
If the third input, which is applied factor by factor as in clause (a) of the lemma bounding walk
expansions with the $(d,L,\kappa_1)$-number $K_\sharp$, is read as an appeal to that lemma's statement
rather than to its method, that lemma's hypothesis $6dK_\sharp^2\gamma\le1$ is also a pure number
condition, met at $\gamma_{reg}$ by the second member of the minimum defining $\gamma_{reg}$ in
Definition~\ref{4.25:not-regularity-class}; in either reading no input asks the radius to be
coupling-small.

The proof of the decomposition lemma uses the first input also in the weighted kernel norms of its
kernel clause, factor by factor. Let $K_\Delta'\ge K_\Delta$ be the
threshold, fixed in Definition~\ref{4.25:not-regularity-class}, below which the walk series of the
first input converge in those weighted kernel norms; the existence of such a finite threshold is part
of the step marked $(\ast)$ below. By the first member of the minimum defining $\gamma_{reg}$ in
Definition~\ref{4.25:not-regularity-class}, $K_\Delta'\gamma_{reg}\le\frac12$, and a fortiori
$K_\Delta\gamma_{reg}\le\frac12$.
$(\ast)$ Under the floor condition relating $M$ and $\kappa_1$, the condition on the levels,
$K_\Delta\gamma_{reg}\le\frac12$ and, for the weighted kernel norms, $K_\Delta'\gamma_{reg}\le\frac12$,
each of the four inputs, and the lemma on which the series of the first input are built, hold on
$\mathfrak{K}_G$; the walk series of the first input converge in the weighted kernel norms of the
kernel clause whenever $K_\Delta'\gamma\le\frac12$, for a finite $(d,L,\kappa_1)$-number
$K_\Delta'\ge K_\Delta$, an estimate on the proof of that input which its statement does not display
and which is not supplied here; and conclusions (i)--(iii) of the decomposition
lemma hold for pairs
$(\mathbb{U},\mathbb{J})\in\mathfrak{K}_G$ and $|\sigma'(\Delta)|\le e^{\kappa_1}$, with the constants
$B_{\mathcal{A}}$, $\|C\|$, $a$, $C_\rho$, $n_b$, $n_C$ taken as their suprema over $\mathfrak{K}_G$, as
in Definition~\ref{4.25:not-regularity-class}. What is not carried out here is the verification that
the proof of the decomposition lemma invokes only these inputs, none of which requires the current
radius to be coupling-small, so that this proof transfers without change from
$\mathfrak{K}_{109}(\mathfrak{B}_{\mathfrak{a}};\gamma_0)$ to $\mathfrak{K}_G$; that proof is not
reproduced here.

\emph{Well-definedness on restrictions.} We use clause (ii) of the decomposition lemma together with the
margin in the locality property of the block pairs \cite{B-CMP116,B-CMP98,B-CMP109}.
At $\sigma'=0$ off the polymer $Y'$
of the piece, the walks that survive are those of $\bigcup Y'$, and they depend on the pair only through
its values there. The piece $v_{e_a^0}$ depends on $\mathbb{U}$ only through the block pair, again by
that locality property. Hence
$v_e$ depends on a global pair only through its restriction to $\tilde S(e)$.

\emph{The bound on $Q_e$.} Let $K_e$ be the piece, attached to $e$, of the walk expansion of the kernel
$\mathcal{A}=C^*\Delta^{(j)}C$, and let $b,b'$ be the two bonds and $Y'$ the polymer of $e$. Let
$\delta$ be the decay rate and $d(b,b')$ the distance of bonds in the kernel bound of the decomposition
lemma, which enter the constants of Definition~\ref{4.25:not-regularity-class}. Cauchy's
inequality on the circles $|\sigma'(\Delta)|=e^{\kappa_1}$ gives
\begin{equation}\label{4.25:eq-gaussian-pieces-6}
|K_e(b,b')|\le B_{\mathcal{A}}\,e^{-\delta d(b,b')}\,e^{-(\kappa_1-1)|Y'|}.
\end{equation}
We put $Q_e(x_1,x_2):=-\frac12\langle x_1(b),K_e(b,b')\,x_2(b')\rangle$, which is bilinear, and
$v_e(x)=Q_e(x,x)$. By \eqref{4.25:eq-gaussian-pieces-6},
\begin{equation*}
|Q_e(x_1,x_2)|\le\tfrac12B_{\mathcal{A}}e^{-\delta d(b,b')}e^{-(\kappa_1-1)|Y'|}\,|x_1|_\infty|x_2|_\infty .
\end{equation*}
For $e=e_a$ the definitions of Definition~\ref{4.25:not-regularity-class} give
\begin{align*}
\frac{\nu^{\natural}(e)}{\varrho^2}&=\nu^{\circ}(e)e^{-c_*|Y'|}\\
&=\tfrac12B_{\mathcal{A}}e^{-\delta d(b,b')}e^{-(\kappa_1-1-c_*)|Y'|}e^{-c_*|Y'|}\\
&=\tfrac12B_{\mathcal{A}}e^{-\delta d(b,b')}e^{-(\kappa_1-1)|Y'|},
\end{align*}
which is \eqref{4.25:eq-gaussian-pieces-2} for the pieces $e=e_a$.
For $e=e_a^0$ the decomposition lemma gives the block-pair term
$v_e(x)=\frac12a\langle x(b),(C^*C)(b,b')x(b')\rangle$, and we put
\begin{equation*}
Q_e(x_1,x_2):=\tfrac12a\langle x_1(b),(C^*C)(b,b')x_2(b')\rangle .
\end{equation*}
The bound on $C(\mathbb{U})$, read on $\mathfrak{K}_G$, gives $|(C^*C)(b,b')|\le\|C\|^2$, so
\begin{equation*}
|Q_e(x_1,x_2)|\le\tfrac12a\|C\|^2|x_1|_\infty|x_2|_\infty=\nu^{\circ}(e_a^0)\,|x_1|_\infty|x_2|_\infty .
\end{equation*}
Since this piece carries no polymer, $\nu^{\natural}(e)/\varrho^2=\nu^{\circ}(e)$, and this is
\eqref{4.25:eq-gaussian-pieces-2} for $e=e_a^0$.

\emph{Invariance.} Let $v$ be $G$-valued. Clause (c) of the localisation statement gives
\begin{equation*}
\mathcal{A}[\mathbb{U}^v,\operatorname{Ad}_v\mathbb{J}](b,b')
=\operatorname{Ad}_{v(b_-)}\,\mathcal{A}[\mathbb{U},\mathbb{J}](b,b')\,\operatorname{Ad}_{v(b'_-)}^{-1}.
\end{equation*}
By the $\operatorname{Ad}$-covariance in clause (i) of the decomposition lemma, the same relation holds
for the walk-expanded kernel at every value of $\sigma'$; the piece $K_e$ is obtained from that kernel
by differentiations and integrations in $\sigma'$, which commute with the two fixed conjugations, so the
relation holds for $K_e$.
We insert this in $Q_e(x^{[v]},x^{[v]})$. The factor $\operatorname{Ad}_{v(b'_-)}^{-1}$ cancels the
rotation of $x(b')$. The factor $\operatorname{Ad}_{v(b_-)}$ preserves the pairing, since $v$ is
$G$-valued, and so cancels against the rotation of $x(b)$. This gives
\eqref{4.25:eq-gaussian-pieces-1}.

Holomorphy along holomorphic families of global pairs in $\mathfrak{K}_G$ is part of the conclusion of
the decomposition lemma, read on $\mathfrak{K}_G$ as above.

\emph{Proof of (ii).} Let $g$ be $G$-valued. Apply \eqref{4.25:eq-gaussian-pieces-1} with $u=g$ to the
pair $p^{g^{-1}}$, and use $(p^{g^{-1}})^g=p$. This gives
\begin{equation*}
v_e\bigl(x^{[g]};p\bigr)=v_e\bigl(x^{[g]};(p^{g^{-1}})^g\bigr)=v_e\bigl(x;p^{g^{-1}}\bigr).
\end{equation*}
By (i), $v_e(x;p)=Q_e(x,x)$ with $Q_e$ bilinear. Moreover $x^{[g]}$ is linear in
$x$, with coefficients the entries of $\operatorname{Ad}_{g(b_-)}$. Hence $g\mapsto Q_e(x^{[g]},x^{[g]})$
is a polynomial of degree at most two in these entries, and it is defined for every
$g\in(G^{\mathbb{C}})^{\text{sites}}$.

\emph{Proof of (iii).} Let $\tilde P\in\mathfrak{K}_P^{(k')}=\mathfrak{K}_{109}(\mathfrak{B}_{cur};\gamma_0^{(k')})$
be a global pair as in the statement, where $\mathfrak{B}_{cur}$ is the family that determines the
levels at step $k'$. This family is, cell by cell, coarser than or equal to $\mathfrak{B}_{\mathfrak{a}}$,
and the two are nested.

\emph{The conditions on $A'$ and $\mathbb{J}$.} At a cell let $m_1$ be the level assigned by
$\mathfrak{B}_{\mathfrak{a}}$ and $m_2\ge m_1$ that assigned by $\mathfrak{B}_{cur}$. Since
$\ell_{m_1}\le\ell_{m_2}$, the condition
of \cite[(3.37)]{B-CMP99b} at level $m_2$, $\ell_{m_2}|A'|<\alpha_1$, implies
$\ell_{m_1}|A'|<\alpha_1$. The
same holds for the gradient conditions and for $\ell^3|\mathbb{J}|$; since moreover
$\gamma_0^{(k')}\le\frac14\gamma_{reg}<\gamma_{reg}$, the current clause of $\mathfrak{K}_G$ holds.

\emph{The condition on $U$.} We use the extension of the condition \cite[(3.35)]{B-CMP99b} to cubes
that cross a seam between levels and have side at most $20M$ times the coarser unit, with level-wise
bounds. Take a cube of the old class, of index $\iota$, and complete
it outward to the mesh of its least current level $\iota''\ge\iota$. Since $\ell_\iota\le\ell_{\iota''}$,
the completed cube has side at most
\begin{equation*}
10M\ell_\iota+2M\ell_{\iota''}\le12M\ell_{\iota''}\le20M\ell_{\iota''},
\end{equation*}
so it is a current cube of that class. Its gauge restricts to the old cube; the constants leave room
for this restriction by a part of the same one-sided condition on the coupling in force at step $k'$
that bounds $\bar\gamma_0(x_{k'})$. Together with the conditions on $A'$
and $\mathbb{J}$, this gives $\tilde P\in\mathfrak{K}_G$.

\emph{The bound \eqref{4.25:eq-gaussian-pieces-a}.} These are the sixth and seventh of the
presentation conditions on $Y^+$. They apply because $Y^+$ contains the stencils
$\tilde\Delta(y)$, $y\in\mathcal{Y}(e)$, which holds since the cells of the arrested region coincide
with the cubes on which the presentation conditions are imposed, by clause (i) of the lemma comparing
these two families of cubes.

\emph{The current.} The bound $\ell^3|\mathbb{J}_{\tilde P}|\le\gamma_0^{(k')}$ everywhere is the
current clause in the definition of $\mathfrak{K}_P^{(k')}$.

\emph{Proof of (iv).} We first record an estimate for products of exponentials. Put
$r_{\Psi}:=0.04$. For matrices $\Phi_1,\Phi_2$ with $|\Phi_1|+|\Phi_2|\le r_{\Psi}$, put
$\Gamma:=e^{\Phi_1}e^{\Phi_2}-1$ and $\Psi(\Phi_1,\Phi_2):=\log(1+\Gamma)$, the
logarithm given by its power series.

The norm $|\cdot|$ is sub-multiplicative. Expanding both exponentials therefore gives
\begin{equation*}
|\Gamma|\le e^{|\Phi_1|+|\Phi_2|}-1\le e^{r_{\Psi}}-1<1,
\end{equation*}
so $\Psi$ is defined. Term by term in the power series, the derivative of $\log(1+\Gamma)$ in
a direction $H_1$ satisfies
\begin{equation*}
\bigl|D\log(1+\Gamma)[H_1]\bigr|\le\frac{|H_1|}{1-|\Gamma|}\le\frac{|H_1|}{2-e^{r_{\Psi}}}.
\end{equation*}
The derivative of $(\Phi_1,\Phi_2)\mapsto e^{\Phi_1}e^{\Phi_2}$ in a direction $(H_1,H_2)$ satisfies
\begin{equation*}
\bigl|D(e^{\Phi_1}e^{\Phi_2})[H_1,H_2]\bigr|\le e^{r_{\Psi}}(|H_1|+|H_2|).
\end{equation*}
By the chain rule and $e^{r_{\Psi}}\le1.0409$, $2-e^{r_{\Psi}}\ge0.9591$,
\begin{equation*}
\bigl|D\Psi(\Phi_1,\Phi_2)[H_1,H_2]\bigr|
\le\frac{e^{r_{\Psi}}}{2-e^{r_{\Psi}}}(|H_1|+|H_2|)
\le\frac{1.0409}{0.9591}(|H_1|+|H_2|)\le1.1(|H_1|+|H_2|).
\end{equation*}
The set $\{|\Phi_1|+|\Phi_2|\le r_{\Psi}\}$ is convex. Integrating along segments therefore shows that
$\Psi$ is $1.1$-Lipschitz in $|\Phi_1|+|\Phi_2|$ on it:
\begin{equation}\label{4.25:eq-gaussian-pieces-7}
|\Psi(\Phi_1,\Phi_2)-\Psi(\Phi_1',\Phi_2')|\le1.1\bigl(|\Phi_1-\Phi_1'|+|\Phi_2-\Phi_2'|\bigr).
\end{equation}
From the power series, $\Psi(0,0)=0$, $\Psi(0,\Phi_2)=\Phi_2$, and
$\Psi(\operatorname{Ad}_U\Phi_1,\operatorname{Ad}_U\Phi_2)=\operatorname{Ad}_U\Psi(\Phi_1,\Phi_2)$.

\emph{Application bondwise.} On each bond of the stencils $\tilde\Delta(y)$, $y\in S^+(e)$, put
$\Phi_1:=i\eta\tau\chi_eA''$ and $\Phi_2:=i\eta A'$, and
define $A'''$ by $i\eta A''':=\Psi(\Phi_1,\Phi_2)$, so that
$e^{i\eta\tau\chi_eA''}e^{i\eta A'}=e^{i\eta A'''}$.

Let $y$ be a block in $S^+(e)$. There $|\chi_e|\le1$, $\ell_y|A''|\le\mathfrak{a}''$ and
$\ell_y|A'|\le\varepsilon_{reg}/5$. For $|\tau|<T_U$ we have $|\tau|\mathfrak{a}''<\varepsilon_{reg}$.
Moreover $\eta\le\ell_y$, and $\varepsilon_{reg}\le10^{-1}/3$ by its definition. Hence
\begin{equation*}
|\Phi_1|+|\Phi_2|\le\frac{\eta}{\ell_y}\Bigl(|\tau|\mathfrak{a}''+\frac{\varepsilon_{reg}}5\Bigr)
\le1.2\,\varepsilon_{reg}\le r_{\Psi} .
\end{equation*}

\emph{The value.} By \eqref{4.25:eq-gaussian-pieces-7} and $\Psi(0,0)=0$,
\begin{equation*}
\ell_y|A'''|\le1.1\bigl(|\tau|\,\ell_y|\chi_eA''|+\ell_y|A'|\bigr).
\end{equation*}

\emph{The gradient.} Let $b'$ be the bond obtained from $b$ by the lattice shift in the direction $\mu$,
and let $\operatorname{Ad}_U$ be the parallel transport by $U$ entering $\nabla_{U,\mu}$. Equivariance
of $\Psi$ gives
\begin{equation*}
i\eta^2\nabla_{U,\mu}A'''(b)
=\Psi\bigl(\operatorname{Ad}_U(\Phi_1,\Phi_2)(b')\bigr)-\Psi\bigl((\Phi_1,\Phi_2)(b)\bigr).
\end{equation*}
Since $|\operatorname{Ad}_U\cdot|=|\cdot|$ for $G$-valued $U$, the arguments again lie in the domain of
\eqref{4.25:eq-gaussian-pieces-7}. That inequality yields
\begin{equation*}
\ell_y^2|\nabla_{U,\mu}A'''|\le1.1\bigl(|\tau|\,\ell_y^2|\nabla_{U,\mu}(\chi_eA'')|
+\ell_y^2|\nabla_{U,\mu}A'|\bigr).
\end{equation*}

\emph{The block size.} We take the supremum over $\tilde\Delta(y)$ and apply clause $(\chi4)$ of
Lemma~\ref{4.25:lem-piece-cut-off}, see \eqref{4.25:eq-piece-cut-off-a}, with the connection $U$. It
bounds the contribution of $\chi_eA''$ by twice that of $A''$, that is, by $2\mathfrak{a}''$. Therefore,
on $S^+(e)$,
\begin{equation*}
N_y(A''')\le1.1\Bigl(2|\tau|\mathfrak{a}''+\frac{\varepsilon_{reg}}5\Bigr)<2.42\,\varepsilon_{reg}
<\alpha_1^{SL},
\end{equation*}
where the last step uses $\varepsilon_{reg}\le\alpha_1^{SL}/3$.

\emph{Away from $S^+(e)$.} On the bonds where $\chi_e$ vanishes together with its shifts, $\Phi_1=0$
and $e^{\Phi_1}e^{i\eta A'}=e^{i\eta A'}$; there we put $A''':=A'$, which agrees with the definition by
$\Psi$ on those of these bonds that lie in the stencils, since $\Psi(0,\Phi_2)=\Phi_2$. By $(\chi4)$,
$\chi_eA''$ contributes nothing to $N_y$ unless $y\in S^+(e)$. The
$G$-valued factor $U$ is untouched.

\emph{The current.} On $S^+(e)$, for $|t|<T_J$, we have $|t|\,\ell^3|J''|<T_J\mathfrak{j}''=\frac12\gamma_{reg}$,
and therefore
\begin{equation*}
\ell^3|\mathbb{J}+t\chi_eJ''|<\tfrac12\gamma_{reg}+\tfrac12\gamma_{reg}=\gamma_{reg}.
\end{equation*}
Off $S^+(e)$ the current is that of $P$.

\emph{Conclusion for the bidisc.} Hence $Q(\tau,t)\in\mathfrak{K}_G$ on the bidisc. It is global, and it
is entire in $(\tau,t)$, since it is built from the exponential of a bondwise linear function of $\tau$
and from an affine function of $t$.

By clause $(\chi1)$ of Lemma~\ref{4.25:lem-piece-cut-off}, $\chi_e\equiv1$ on the bonds of
$\tilde S(e)$. There $Q(\tau,t)=(e^{i\eta\tau A''}P^U,\mathbb{J}+tJ'')$.

The map $(\tau,t)\mapsto Q(\tau,t)$ is a holomorphic family of global pairs in $\mathfrak{K}_G$ on the
bidisc. By (i), $(\tau,t)\mapsto v_e(x;Q(\tau,t)|_{\tilde S(e)})$ is holomorphic there.

Finally, $T_U\ge5$ and $T_J\ge2$ by \eqref{4.25:eq-gaussian-pieces-4}, so $(1,1)$ lies in the bidisc.
By the identity on the bonds of $\tilde S(e)$ just obtained, $p_2=Q(1,1)|_{\tilde S(e)}$, which
belongs to $\mathfrak{K}_G(\tilde S(e))$ by the definition of that class.

\emph{Proof of (iv)$'$.} At $\tau=1$, taking the supremum over $\tilde\Delta(y)$ in the bounds for the
value and the gradient of $A'''$ obtained in the proof of (iv), and applying $(\chi4)$ as there, gives,
on $S^+(e)$,
\begin{equation*}
N_y(A''')\le1.1\bigl(2\mathfrak{a}''+N_y(A')\bigr).
\end{equation*}
Under the stronger hypotheses this is at most
\begin{equation*}
1.1\Bigl(\frac{2}{20}+\frac1{20}\Bigr)\varepsilon_{reg}
=\frac{1.1\cdot3}{20}\,\varepsilon_{reg}\le\frac15\,\varepsilon_{reg}.
\end{equation*}
For the current at $t=1$,
\begin{equation*}
\ell^3|\mathbb{J}+\chi_eJ''|\le\tfrac14\gamma_{reg}+\tfrac14\gamma_{reg}=\tfrac12\gamma_{reg}.
\end{equation*}
By (iv), $Q(1,1)=(e^{i\eta A'''}U,\mathbb{J}+\chi_eJ'')$ is a global pair in $\mathfrak{K}_G$. It
therefore satisfies the hypotheses imposed on $P$ in (iv).
\end{proof}

\begin{remark}
The lemma is formulated on $\mathfrak{K}_G$ and not on $\mathfrak{K}_P^{(k')}$ for the following reason.
A Schwarz estimate carried out at the coupling-small radius $\gamma_0^{(k')}$ would
return a coupling-dependent polylogarithmic factor, whereas all radii of $\mathfrak{K}_G$ are
coupling-free. Clause (iv) asks no global smallness of $(A'',J'')$: only their sizes
$\mathfrak{a}''$ and $\mathfrak{j}''$ on $S^+(e)$ enter.
\end{remark}

\begin{lemma}[Oscillation of a Gaussian piece under small moves. Stated; proof incomplete at the
step marked $(\ast)$]\label{4.25:lem-piece-oscillation}
Let $e\in\mathfrak{E}^G_W$, and let $P$, $A''$, $J''$, $\mathfrak{a}''$, $\mathfrak{j}''$ be as in
clause (iv) of Lemma~\ref{4.25:lem-gaussian-pieces} (stated; proof incomplete), displayed in
\eqref{4.25:eq-gaussian-pieces-b}. Put $p_1:=P\restriction\tilde S(e)$, the restriction of $P$ to
the bonds of $\tilde S(e)$, and let $p_2$ be the pair defined in that clause. Let $g_1,g_2$ be
maps from the lattice sites to $G^{\mathbb{C}}$ such that
$\sup_s\operatorname{pol}(g_i(s))\le \tfrac13$ for $i=1,2$ and
\begin{equation*}
\mathfrak{g}:=\sup_s\operatorname{dis}\bigl(g_2(s)g_1(s)^{-1}\bigr)\le\tfrac13,
\end{equation*}
and let $x$ be a collar variable with $|x|_\infty\le\mathfrak{b}_x$. Then
\begin{equation}\label{4.25:eq-piece-oscillation-a}
\bigl|v_e^{\sharp}(x;p_2;g_2)-v_e^{\sharp}(x;p_1;g_1)\bigr|
\le\nu^{\natural}(e)\bigl[\mathsf{c}_U\mathfrak{a}''+\mathsf{c}_J\mathfrak{j}''+3\mathfrak{g}\bigr],
\end{equation}
where $\mathsf{c}_U:=1/\varepsilon_{reg}$, and $\mathsf{c}_J:=2/\gamma_{reg}$ for the pieces $e=e_a$,
while $\mathsf{c}_J:=0$ for
the pieces $e=e_a^0$; here $\varepsilon_{reg}$ and $\gamma_{reg}$ are the coupling-free radii of
the class $\mathfrak{K}_G$ that enter clause (iv) of that lemma.
The constants $\mathsf{c}_U$ and $\mathsf{c}_J$ are numbers depending only on
$(d,L,\kappa_1)$; they do not depend on the coupling. The right member of
\eqref{4.25:eq-piece-oscillation-a} is linear in $(\mathfrak{a}'',\mathfrak{j}'',\mathfrak{g})$.
\end{lemma}

\begin{proof}
For $i=1,2$ let $\Lambda_i$ be the operator acting on $\mathfrak{g}^{\mathbb{C}}$-valued functions
$y$ on the bonds by $(\Lambda_iy)(b):=\operatorname{Ad}_{g_i(b_-)}y(b)$, where $b_-$ denotes the
initial site of the bond $b$, so that
$\Lambda_ix=x^{[g_i]}$ is the rotated variable of clause (ii) of
Lemma~\ref{4.25:lem-gaussian-pieces} (stated; proof incomplete), and, by the
definition of the orbit datum in that clause,
\begin{equation}\label{4.25:eq-piece-oscillation-1}
v_e^{\sharp}(x;p;g_i)=v_e(\Lambda_ix;p)\qquad\text{for every pair }p\in\mathfrak{K}_G(\tilde S(e)).
\end{equation}
Norms of the operators $\Lambda_i$ are taken with respect to $|\cdot|_\infty$.

\emph{The rotations.} Since $\sup_s\operatorname{pol}(g_i(s))\le\frac13$, we have
$\|\Lambda_i\|\le e^{2/3}\le2$. Indeed, for $g\in G^{\mathbb{C}}$ we have
$\operatorname{Ad}_gy=gyg^{-1}$; with $|\cdot|$ the norm on $\mathfrak{g}^{\mathbb{C}}$ of
Lemma~\ref{4.25:lem-word-length-distance}, for which $|gyg^{-1}|\le\|g\|\,|y|\,\|g^{-1}\|$, and
$\|g\|\,\|g^{-1}\|\le e^{2\operatorname{pol}(g)}$ by the definition
$\operatorname{pol}(g)=\log\max\{\|g\|,\|g^{-1}\|\}$; moreover $e^{2/3}<2$. At every site $s$,
$\operatorname{Ad}_{g_2(s)}=\operatorname{Ad}_{g_2(s)g_1(s)^{-1}}\operatorname{Ad}_{g_1(s)}$, hence
\begin{equation*}
\Lambda_2-\Lambda_1=\bigl(\operatorname{Ad}_{g_2g_1^{-1}}-1\bigr)\Lambda_1,
\end{equation*}
where $\operatorname{Ad}_{g_2g_1^{-1}}$ acts at each bond $b$ through the site $b_-$. By property
(d4) of Lemma~\ref{4.25:lem-word-length-distance}, displayed in
\eqref{4.25:eq-word-length-distance-a}, applied at each site $s$ to the element
$g_2(s)g_1(s)^{-1}$ of $G^{\mathbb{C}}$, whose word-length distance is at most $\mathfrak{g}$,
we have
\begin{equation*}
\bigl\|\operatorname{Ad}_{g_2(s)g_1(s)^{-1}}-1\bigr\|
\le e^{2\operatorname{dis}(g_2(s)g_1(s)^{-1})}-1\le e^{2\mathfrak{g}}-1,
\end{equation*}
because $\lambda\mapsto e^{2\lambda}-1$ is increasing. The
function $\lambda\mapsto e^{2\lambda}-1$ is convex and vanishes at $\lambda=0$; hence on
$[0,\frac13]$ it lies below its chord, whose slope is $3(e^{2/3}-1)<3$, and so
$e^{2\lambda}-1\le3\lambda$ there. As
$\mathfrak{g}\le\frac13$, we obtain
\begin{equation}\label{4.25:eq-piece-oscillation-2}
\|\Lambda_2-\Lambda_1\|\le2\bigl(e^{2\mathfrak{g}}-1\bigr)\le6\mathfrak{g}.
\end{equation}

\emph{Decomposition.} Write $Q_2:=Q_e(\cdot,\cdot;p_2)$ for the bilinear form of $v_e$ at the pair
$p_2$, so that $v_e(y;p_2)=Q_2(y,y)$ by clause (i) of Lemma~\ref{4.25:lem-gaussian-pieces}
(stated; proof incomplete); the pair $p_2$ lies in $\mathfrak{K}_G(\tilde S(e))$ by clause
(iv) of the same lemma. Put $x_1:=\Lambda_1x$. By \eqref{4.25:eq-piece-oscillation-1},
\begin{equation*}
v_e^{\sharp}(x;p_2;g_2)-v_e^{\sharp}(x;p_1;g_1)=D_1+D_2,
\end{equation*}
where
\begin{align*}
D_1&:=Q_2(\Lambda_2x,\Lambda_2x)-Q_2(\Lambda_1x,\Lambda_1x),\\
D_2&:=v_e(x_1;p_2)-v_e(x_1;p_1).
\end{align*}

\emph{The term $D_1$.} By bilinearity,
\begin{equation*}
D_1=Q_2\bigl((\Lambda_2-\Lambda_1)x,\Lambda_2x\bigr)+Q_2\bigl(\Lambda_1x,(\Lambda_2-\Lambda_1)x\bigr),
\end{equation*}
since the right member equals
\begin{equation*}
Q_2(\Lambda_2x,\Lambda_2x)-Q_2(\Lambda_1x,\Lambda_2x)+Q_2(\Lambda_1x,\Lambda_2x)
-Q_2(\Lambda_1x,\Lambda_1x).
\end{equation*}

$(\ast)$ We now use the kernel bound of clause (i) of Lemma~\ref{4.25:lem-gaussian-pieces},
namely $|Q_e(y,z)|\le\nu^{\natural}(e)|y|_\infty|z|_\infty/\varrho^2$, on the coupling-free class
$\mathfrak{K}_G(\tilde S(e))$ to which $p_2$ belongs. Its proof on the class
$\mathfrak{K}_G(\tilde S(e))$ is the step not completed here.

By $\|\Lambda_i\|\le2$ and \eqref{4.25:eq-piece-oscillation-2} we have
$|\Lambda_ix|_\infty\le2\mathfrak{b}_x$ and
$|(\Lambda_2-\Lambda_1)x|_\infty\le6\mathfrak{g}\mathfrak{b}_x$. Each of the two terms of $D_1$ is
therefore at most $\nu^{\natural}(e)\cdot6\mathfrak{g}\mathfrak{b}_x\cdot2\mathfrak{b}_x/\varrho^2$,
and
\begin{equation*}
|D_1|\le24\,\mathfrak{g}\,\nu^{\natural}(e)\Bigl(\frac{\mathfrak{b}_x}{\varrho}\Bigr)^2
\le3\,\mathfrak{g}\,\nu^{\natural}(e),
\end{equation*}
because $\mathfrak{b}_x/\varrho\le2.3/6.6$, where $\mathfrak{b}_x$ is the bound defining the box of
the collar variable and $\varrho$ its radius, both fixed with the collar variable, and
$(2.3/6.6)^2\le\frac18$.

\emph{A two-point Schwarz bound.} Let $h$ be holomorphic on the disc $|\tau|<T$ with
$|h|\le B$ there, and let $T\ge2$. The closed disc of radius $r<T-1$ about a point
$\tau\in[0,1]$ lies in the disc of radius $T$ about $0$, so Cauchy's estimate gives
$|h'(\tau)|\le B/r$; letting $r\uparrow T-1$,
\begin{equation}\label{4.25:eq-piece-oscillation-3}
|h(1)-h(0)|\le\int_0^1|h'(\tau)|\,d\tau\le\frac{B}{T-1}\le\frac{2B}{T},
\end{equation}
the last inequality because $T\ge2$ implies $T-1\ge T/2$.

\emph{The term $D_2$.} Let $Q(\tau,t)$ be the continuation of clause (iv) of
Lemma~\ref{4.25:lem-gaussian-pieces} (stated; proof incomplete), with the radii
$T_U:=\varepsilon_{reg}/\mathfrak{a}''\ge5$ and $T_J:=\frac12\gamma_{reg}/\mathfrak{j}''\ge2$. By that
clause, $(\tau,t)\mapsto v_e(x_1;Q(\tau,t)\restriction\tilde S(e))$ is holomorphic on the bidisc
$|\tau|<T_U$, $|t|<T_J$, and $Q(\tau,t)\restriction\tilde S(e)\in\mathfrak{K}_G(\tilde S(e))$
there. We have $|x_1|_\infty\le2\mathfrak{b}_x$, so, by the kernel bound used at $(\ast)$, every
value of $v_e(x_1;\cdot)$ occurring below satisfies
\begin{equation*}
|v_e(x_1;Q(\tau,t)\restriction\tilde S(e))|
\le\frac{\nu^{\natural}(e)}{\varrho^2}(2\mathfrak{b}_x)^2\le\tfrac12\nu^{\natural}(e),
\end{equation*}
since $4(\mathfrak{b}_x/\varrho)^2\le\frac48$.

Let $f(\tau):=v_e(x_1;Q(\tau,0)\restriction\tilde S(e))$ for $|\tau|<T_U$. Applying
\eqref{4.25:eq-piece-oscillation-3} with $B=\frac12\nu^{\natural}(e)$ and $T=T_U\ge2$,
\begin{equation*}
|f(1)-f(0)|\le\frac{\nu^{\natural}(e)}{T_U}=\frac{\nu^{\natural}(e)\,\mathfrak{a}''}{\varepsilon_{reg}}.
\end{equation*}
Let $h(t):=v_e(x_1;Q(1,t)\restriction\tilde S(e))$ for $|t|<T_J$. For the pieces $e=e_a$,
applying \eqref{4.25:eq-piece-oscillation-3} with $B=\frac12\nu^{\natural}(e)$ and $T=T_J\ge2$
gives
\begin{equation*}
|h(1)-h(0)|\le\frac{\nu^{\natural}(e)}{T_J}
=\frac{2\,\nu^{\natural}(e)\,\mathfrak{j}''}{\gamma_{reg}}
=\mathsf{c}_J\,\nu^{\natural}(e)\,\mathfrak{j}'',
\end{equation*}
by the definition of $\mathsf{c}_J$ for these pieces. For the pieces $e=e_a^0$ the function
$v_e(x_1;\cdot)$ depends on the gauge-field component $\mathbb{U}$ of the pair alone, and $t$
changes only the current component; hence $h$ is constant and $h(1)=h(0)$. Thus
$|h(1)-h(0)|\le\mathsf{c}_J\,\nu^{\natural}(e)\,\mathfrak{j}''$ in both cases.

Finally, $Q(0,0)=P$ by the definition of the continuation, so $f(0)=v_e(x_1;p_1)$; both $f(1)$
and $h(0)$
are the value at $Q(1,0)$, so $f(1)=h(0)$; and on the bonds of $\tilde S(e)$ the pair $Q(1,1)$
equals $p_2$ by clause (iv) of Lemma~\ref{4.25:lem-gaussian-pieces} (stated; proof
incomplete), so $h(1)=v_e(x_1;p_2)$.
Therefore
\begin{equation*}
D_2=\bigl(h(1)-h(0)\bigr)+\bigl(f(1)-f(0)\bigr),\qquad
|D_2|\le\nu^{\natural}(e)\bigl[\mathsf{c}_U\mathfrak{a}''+\mathsf{c}_J\mathfrak{j}''\bigr].
\end{equation*}
Adding the bounds on $|D_1|$ and $|D_2|$ gives \eqref{4.25:eq-piece-oscillation-a}.
\end{proof}

\begin{remark}
Since $\gamma_{reg}$ depends only on $(d,L,\kappa_1)$ and not on the coupling, the loss
$2\mathfrak{j}''/\gamma_{reg}$ in the current direction is the absolute increment $\mathfrak{j}''$
times a coupling-free constant.
\end{remark}

\begin{remark}
The requirement that $p_2$ be the pair of clause (iv) of
Lemma~\ref{4.25:lem-gaussian-pieces} (stated; proof incomplete) cannot be dropped,
and $\operatorname{dis}$ cannot be replaced by $\operatorname{pol}$ in the definition of $\mathfrak{g}$.
Let $u$ be $G$-valued, let $e=e_a^0(b,b')$ be the piece $e_a^0$ attached to the pair of bonds
$(b,b')$, and suppose $u(b_-)=1$,
$\operatorname{Ad}_{u(b'_-)}\xi=-\xi$, $x(b)=x(b')=\xi$ and $|\xi|=\mathfrak{b}_x$.
\begin{itemize}
\item[($\alpha$)] Take $g_1=g_2=1$ and $p_2:=p_1^u$, the gauge transform of $p_1$ by $u$. The
difference of the two values is
$a|\langle\xi,(C^*C)(b,b')\xi\rangle|$, where $a$ is the scalar coefficient carried by the piece
$e_a^0(b,b')$,
$C$ is the operator through which the kernel $\mathcal{A}=C^*\Delta^{(j)}C$ of the quadratic
collar form is written, $(C^*C)(b,b')$ is the $(b,b')$ entry of $C^*C$, and
$\langle\cdot,\cdot\rangle$ is the pairing of the quadratic collar weight $\mathsf{q}_W$. Since the
presented spaces are $G$-invariant, $p_1^u$ lies in the same space as $p_1$, while the
difference above is of the order of $\nu^{\natural}(e)$; hence no bound of the form
\eqref{4.25:eq-piece-oscillation-a} holds between the values at two arbitrary points of such a
space, and $p_2$ must be the pair of clause (iv).
\item[($\beta$)] Take $A''=J''=0$, $g_1=1$, $g_2=u$. The difference is the same number, by
clause (ii) of Lemma~\ref{4.25:lem-gaussian-pieces} (stated; proof incomplete). The
quantity
$\sup_s\operatorname{pol}(g_2(s)g_1(s)^{-1})$ vanishes here, since $u$ is $G$-valued, so a bound
with $\operatorname{pol}$ in place of $\operatorname{dis}$ would have right member zero. With
$\operatorname{dis}$, property (d4) of Lemma~\ref{4.25:lem-word-length-distance} gives
\begin{equation*}
2|\xi|=\bigl|(\operatorname{Ad}_{u(b'_-)}-1)\xi\bigr|
\le\bigl(e^{2\operatorname{dis}(u(b'_-))}-1\bigr)|\xi|,
\end{equation*}
which forces $\operatorname{dis}(u(b'_-))\ge\frac12\log3>\frac13$. This configuration is therefore
excluded by the hypothesis $\mathfrak{g}\le\frac13$. More generally, every element of
$G^{\mathbb{C}}$ whose adjoint action maps some non-zero $\xi\in\mathfrak{g}^{\mathbb{C}}$ to
$-\xi$ has word-length distance at least $\frac12\log3>\frac13$, by the same display with that
element in place of $u(b'_-)$, so every such rotation is
excluded by the hypothesis $\mathfrak{g}\le\frac13$.
\end{itemize}
\end{remark}

\begin{lemma}[The current partner of a difference of layers]\label{4.25:lem-current-partner}
Fix one kernel $(\mathcal{W},\sigma)$ with $\mathsf{p}(\mathcal{W})\le 1$. For $i=1,2$ let
$(A'_i,X_i,\mathbb{H}_i)$ be data
as in the representation statement for the partner current of a layer, both taken at this one
kernel $(\mathcal{W},\sigma)$, so that the partner current of the layer is given by the
representation
\begin{equation}\label{4.25:eq-current-partner-1}
\mathfrak{s}^{\sharp}_{\sigma}[\mathbb{H}]
=\bigl(\Delta_{\mathcal{W}}-P_1(\sigma)+\mathsf{K}\bigr)A'-\mathfrak{m}_H(\sigma)X+F(\mathcal{W},\mathbb{H}),
\end{equation}
that is, by the equations \cite[(3.11)--(3.13)]{B-CMP109} written for the pair, and assume
$N_y(\mathbb{H}_i)\le 1/40$ for $i=1,2$. For any of these quantities write
$\delta\cdot:=\cdot_1-\cdot_2$, and set
\begin{equation}\label{4.25:eq-current-partner-2}
\mathsf{N}_y(\delta):=\max\bigl\{N^{\Delta}_y(\delta A'),\,N_y(\delta\mathbb{H}),\,
\ell^3|P_1(\sigma)\delta A'|,\,\ell^3|\mathfrak{m}_H(\sigma)\delta X|\bigr\},
\end{equation}
the four entries of \cite[(3.14)]{B-CMP109}; here $N^{\Delta}_y(\delta A')$ denotes the entry of
\cite[(3.14)]{B-CMP109} attached to the term $\Delta_{\mathcal{W}} A'$, and $\ell=\ell_y$ is the
length in the
definition of $N_y$. Then
\begin{equation}\label{4.25:eq-current-partner-a}
\ell_y^3\bigl|\mathfrak{s}^{\sharp}_{\sigma}[\mathbb{H}_1]-\mathfrak{s}^{\sharp}_{\sigma}[\mathbb{H}_2]\bigr|
\le K_{\mathsf{j}}\,\mathsf{N}_y(\delta),
\qquad K_{\mathsf{j}}:=3+18d+9K_J ,
\end{equation}
where $K_J$ denotes the constant of the layer bound for the functional $F(\mathcal{W},\cdot)$.
\end{lemma}

\begin{proof}
Both data are taken at the same kernel $(\mathcal{W},\sigma)$, so the operators
$\Delta_{\mathcal{W}}$, $P_1(\sigma)$,
$\mathsf{K}$ and $\mathfrak{m}_H(\sigma)$ in \eqref{4.25:eq-current-partner-1} are the same for
$i=1$ and $i=2$. Subtracting the two instances of \eqref{4.25:eq-current-partner-1} and using the
linearity of these operators gives
\begin{align*}
\mathfrak{s}^{\sharp}_{\sigma}[\mathbb{H}_1]-\mathfrak{s}^{\sharp}_{\sigma}[\mathbb{H}_2]
&=\Delta_{\mathcal{W}}\delta A'-P_1(\sigma)\delta A'+\mathsf{K}\delta A'-\mathfrak{m}_H(\sigma)\delta X\\
&\quad+\bigl(F(\mathcal{W},\mathbb{H}_1)-F(\mathcal{W},\mathbb{H}_2)\bigr).
\end{align*}
We bound the five terms separately.

\emph{The linear members.} The entry $N^{\Delta}_y(\delta A')$ of \eqref{4.25:eq-current-partner-2}
bounds the first term:
$\ell^3|\Delta_{\mathcal{W}}\delta A'|\le N^{\Delta}_y(\delta A')\le\mathsf{N}_y(\delta)$.
By the lemma bounding the operator $\mathsf{K}$, the third term satisfies
\begin{equation*}
\ell^3|\mathsf{K}\delta A'|\le 18d\,\mathsf{p}(\mathcal{W})\,N_y(\delta\mathbb{H})\le 18d\,\mathsf{N}_y(\delta),
\end{equation*}
where $\mathsf{p}(\mathcal{W})\le 1$ was used. The second and fourth terms are bounded by
$\mathsf{N}_y(\delta)$,
since $\ell^3|P_1(\sigma)\delta A'|$ and $\ell^3|\mathfrak{m}_H(\sigma)\delta X|$ are themselves entries of
\eqref{4.25:eq-current-partner-2}. The linear members together contribute at most
$(3+18d)\,\mathsf{N}_y(\delta)$.

\emph{The nonlinear member.} The functional $F(\mathcal{W},\cdot)$ is local \cite[\S3]{B-CMP109}
and holomorphic,
and for a layer $h$ the layer bound gives
\begin{equation*}
\ell^3|F(\mathcal{W},h)|\le K_J\,(\mathsf{p}(\mathcal{W})+\Theta_0+\Theta_1)(\Theta_0+\Theta_1),
\end{equation*}
where $\Theta_0$ and $\Theta_1$ are the two size quantities of the layer $h$ on which that bound depends.
Put $D:=N_y(\delta\mathbb{H})$ and $\varphi(z):=F(\mathcal{W},\mathbb{H}_2+z\,\delta\mathbb{H})$
for complex $z$,
so that $\varphi(0)=F(\mathcal{W},\mathbb{H}_2)$ and $\varphi(1)=F(\mathcal{W},\mathbb{H}_1)$.
On the disc $|z|\le R$,
$R:=1/(40D)$, the layer bound gives $\ell^3|\varphi(z)|\le 0.11\,K_J$. In particular, at $z=0$ this is the
bound $\ell^3|F(\mathcal{W},\mathbb{H}_2)|\le 0.11\,K_J$. The hypotheses are symmetric in the
indices $1,2$, so the
same argument with the roles of $\mathbb{H}_1$ and $\mathbb{H}_2$ exchanged gives
$\ell^3|F(\mathcal{W},\mathbb{H}_1)|\le 0.11\,K_J$.

Suppose first that $R\ge 1$. The function $\varphi-\varphi(0)$ is holomorphic on $|z|<R$, vanishes at
$z=0$, and is bounded there by $0.22\,K_J$ in the norm $\ell^3|\cdot|$. The Schwarz lemma (for a function
with values in a normed space one applies it to $\lambda\circ(\varphi-\varphi(0))$ for every linear
functional $\lambda$ of norm one) gives $\ell^3|\varphi(z)-\varphi(0)|\le 0.22\,K_J|z|/R$ for $|z|<R$, and
by continuity for $|z|\le R$. At $z=1$, which lies in the disc since $R\ge 1$, and with $1/R=40D$:
\begin{equation*}
\ell^3|F(\mathcal{W},\mathbb{H}_1)-F(\mathcal{W},\mathbb{H}_2)|\le\frac{0.22\,K_J}{R}=8.8\,K_J D\le 9K_J D .
\end{equation*}
Suppose next that $R<1$, that is, $40D>1$. The two sup bounds obtained above give
\begin{equation*}
\ell^3|F(\mathcal{W},\mathbb{H}_1)-F(\mathcal{W},\mathbb{H}_2)|\le 0.22\,K_J<0.22\cdot 40\,D\,K_J\le 9K_J D .
\end{equation*}
In both cases the nonlinear member is at most $9K_J D=9K_J N_y(\delta\mathbb{H})\le 9K_J\mathsf{N}_y(\delta)$.

Adding the contribution $(3+18d)\,\mathsf{N}_y(\delta)$ of the linear members gives
\eqref{4.25:eq-current-partner-a} with $K_{\mathsf{j}}=3+18d+9K_J$.
\end{proof}

\begin{remark}
The size $\mathsf{N}_y(\delta)$ of \eqref{4.25:eq-current-partner-2} is the quantity called the unit
shift in \cite[\S3]{B-CMP109}. Consequently, at a site whose partner
current is given by the representation \eqref{4.25:eq-current-partner-1}, the Lipschitz bound
\eqref{4.25:eq-current-partner-a} is available as soon as the layer at that site is bounded in
$\mathsf{N}_y$.
\end{remark}

\begin{lemma}[The current partner of a cut layer. Stated; the proof below is incomplete at the
step marked $(\ast)$]\label{4.25:lem-cut-layer-partner}
The lattice operators are written as follows, consistently with \cite[(3.10)--(3.13)]{B-CMP109}.
The lattice spacing of the layer is $\eta$, $e_\nu$ is the unit vector in the direction $\nu$, a
bond $b$ has initial point $b_-$, and for the bond from $x$ to $x+\eta e_\nu$ we write
$R_{x,\nu}$ for $R_b$. Scalar functions act at the base point, $(fB)_\mu(x):=f(x)B_\mu(x)$, and
\begin{gather*}
R_b:=\operatorname{Ad}_{W(b)},\qquad
(\nabla_\nu B)_\mu(x):=\eta^{-1}\bigl[R_{x,\nu}B_\mu(x+\eta e_\nu)-B_\mu(x)\bigr],\\
(D_WB)_{\mu\nu}:=(\nabla_\mu B)_\nu-(\nabla_\nu B)_\mu,\qquad
\operatorname{st}(b):=\{z:\ |z-b_-|_\infty\le\eta\}.
\end{gather*}
We write $D^*_W$ for the transpose of $D_W$ with respect to the bilinear pairings, and
$\nabla_W B$ for the family of the $(\nabla_\nu B)_\mu$. Put
\[
\tau:=\sup_b\,\max\bigl\{\|R_b\|,\|R_b^{-1}\|\bigr\},
\]
the supremum being taken over the bonds of the stencil. Then $\tau=1$ when $W$ is $G$-valued,
$G$ the gauge group, and $\tau\le e^{2\eta\sup|A^W|}\le 5/4$ whenever
$\eta\sup|A^W|\le 1/10$, where $A^W$ is the potential of $W$; in particular, in the regime
$\ell|A^W|\le 1/40$, where $\ell\ge\eta$ is the unit length of that regime, one has
$\tau\le e^{1/20}<9/8$.

Let $(W,\sigma)$ be a kernel whose irregularity $\mathsf{p}(W)$, as in
Lemma~\ref{4.25:lem-current-partner}, satisfies $\mathsf{p}(W)\le 1$, and let $(A',X,\mathbb{H})$
be the layer datum at $(W,\sigma)$ furnished by the layer theorem of this chapter. The operators
$\Delta_W$, $P_1(\sigma)$, $\mathsf{K}$, $\mathfrak{m}_H(\sigma)$, $H(\sigma)$, the remainder
$F(W,\cdot)$, the block sizes $N_y$, $N^{\Delta}_y$ and the constant $K_J$ are those of
Lemma~\ref{4.25:lem-current-partner}, and $K_{\mathsf{j}}$ is the constant
\eqref{4.25:eq-current-partner-a}. Let $\square$ be the response cube, a cube of the cover with
centre $y_\square$ and level $\operatorname{lev}\square$; put
$\ell_\square:=\ell_{\operatorname{lev}\square}$. The cube $\square$ is the union of $2^d$ cubes
of side $M\ell_\square$ of the partition $\pi_{\operatorname{lev}\square}$, namely
$\square=\{z:\ |z-y_\square|_\infty\le M\ell_\square\}$. Its cut-off is
\[
\hat\chi_{\square}(x):=\prod_{\nu}\zeta\Bigl(\frac{x_\nu-y_{\square,\nu}}{M\ell_\square}\Bigr),
\]
where $\zeta$ is the $C^2$ function of the cover of \cite[\S3]{B-CMP109}: $\zeta=1$ on
$[-\tfrac13,\tfrac13]$, $\zeta=0$ outside $(-\tfrac23,\tfrac23)$, and $|\zeta'|,|\zeta''|\le 5$;
we take moreover $0\le\zeta\le 1$, as a partition of unity allows (any finite bound on $\zeta$
would serve). For $t\in[0,1]$ and the cut parameter $\mathsf{t}$ put
$\chi:=t+\mathsf{t}\,\hat\chi_{\square}$. Fix a block $\Delta(y)$, of side $\ell_y$, with doubled
stencil $\tilde\Delta(y)$, assume $N_y(\chi\mathbb{H})\le 1/40$, and define the five-entry size
\begin{equation}\label{4.25:eq-cut-layer-partner-2}
\begin{split}
\mathsf{N}^{\chi}_y:=\max\bigl\{&N^{\Delta}_y(\chi A'),\ N_y(\chi\mathbb{H}),\
\ell_y^3|P_1(\sigma)(\chi A')|,\ \ell_y^3|\chi\,\mathfrak{m}_H(\sigma)X|,\\
&|\mathsf{t}|\,N_y(H(\sigma)X)\,\mathbf{1}[\tilde\Delta(y)\cap\square\neq\emptyset]\bigr\},
\end{split}
\end{equation}
whose entries are referred to as the first, \dots, fifth entry. Then the partner current of the
cut layer,
\begin{equation}\label{4.25:eq-cut-layer-partner-1}
\mathcal{J}^{\sharp}[\chi]:=(\Delta_W-P_1(\sigma)+\mathsf{K})(\chi A')
-\chi\,\mathfrak{m}_H(\sigma)X-[D^*_WD_W,\chi]\,H(\sigma)X+F(W,\chi\mathbb{H}),
\end{equation}
satisfies
\begin{equation}\label{4.25:eq-cut-layer-partner-3}
\begin{split}
\ell_y^3\,|\mathcal{J}^{\sharp}[\chi]|
&\le K_{\mathsf{j}}\max\bigl\{N^{\Delta}_y(\chi A'),\,N_y(\chi\mathbb{H}),\,
\ell_y^3|P_1(\sigma)(\chi A')|,\,\ell_y^3|\chi\,\mathfrak{m}_H(\sigma)X|\bigr\}\\
&\quad+c_\zeta\,|\mathsf{t}|\,N_y(H(\sigma)X)\,\mathbf{1}[\tilde\Delta(y)\cap\square\neq\emptyset]
\ \le\ (K_{\mathsf{j}}+c_\zeta)\,\mathsf{N}^{\chi}_y,
\end{split}
\end{equation}
where
\begin{equation}\label{4.25:eq-cut-layer-partner-a}
c_\zeta:=2c_d,\qquad c_d:=(2+4\tau)(d-1).
\end{equation}
Thus $c_d=6(d-1)$ at $G$-valued kernels, and $c_d\le 7(d-1)$ on the regularity class of pairs
cited from \cite{B-CMP109}, on which $\tau\le 5/4$; for $d=4$ these values are $18$ and $21$.
\end{lemma}

\begin{proof}
We estimate the four kinds of terms of \eqref{4.25:eq-cut-layer-partner-1} in turn.

\emph{The terms with $\Delta_W$, $\mathsf{K}$, $P_1(\sigma)$, $\mathfrak{m}_H(\sigma)$ and the
remainder.} These are estimated by the lines of the proof of
Lemma~\ref{4.25:lem-current-partner}, applied with the first datum equal to the cut datum and
the second datum equal to zero, so that the differences there are read at
$(\chi A',\chi\mathbb{H})$; the hypothesis $N_y(\chi\mathbb{H})\le 1/40$ is the one that lemma
requires of the first datum, and the zero datum satisfies it. The first of those lines gives
$\ell_y^3|\Delta_W(\chi A')|\le N^{\Delta}_y(\chi A')$, which is the first entry. The line for
$\mathsf{K}$ gives
\[
\ell_y^3|\mathsf{K}(\chi A')|\le 18d\,\mathsf{p}(W)\,N_y(\chi A'),
\]
and since $\mathsf{p}(W)\le 1$ and $N_y(\chi A')\le N^{\Delta}_y(\chi A')$, this is at most
$18d$ times the first entry. The terms
$P_1(\sigma)(\chi A')$ and $\chi\,\mathfrak{m}_H(\sigma)X$ contribute exactly the third and the
fourth entry. For the remainder, $F(W,0)=0$ by \cite[(3.11)]{B-CMP109} (at $A=0$ that formula
gives $J'=D^*\xi^{-2}\operatorname{Im}\partial U$, in the notation of the cited paper), and the
line of the proof of Lemma~\ref{4.25:lem-current-partner} for $F$ gives
\[
\ell_y^3|F(W,\chi\mathbb{H})|\le K_J\bigl(1+\tfrac1{20}\bigr)\cdot 2N_y(\chi\mathbb{H})
\le 9K_J\,N_y(\chi\mathbb{H}),
\]
that is, at most $9K_J$ times the second entry.

\emph{The commutator: product rules.} For a scalar function $f$ let
\[
(\partial_\nu f)(x):=\eta^{-1}[f(x+\eta e_\nu)-f(x)],\qquad
(\partial^-_\nu f)(x):=\eta^{-1}[f(x)-f(x-\eta e_\nu)]
\]
be its forward and backward lattice differences. Inserting
$f(x+\eta e_\nu)=f(x)+\eta(\partial_\nu f)(x)$ into the definition of $\nabla_\nu$ gives the
lattice product rule
\begin{equation}\label{4.25:eq-cut-layer-partner-4}
\nabla_\nu(fB)_\mu(x)=f(x)(\nabla_\nu B)_\mu(x)+(\partial_\nu f)(x)\,R_{x,\nu}B_\mu(x+\eta e_\nu).
\end{equation}
This rule and its transposed form give $[D_W,\chi]=E_\chi$ and $[D^*_W,\chi]=E'_\chi$, where
\begin{equation}\label{4.25:eq-cut-layer-partner-5}
\begin{split}
(E_\chi B)_{\mu\nu}(x)&:=(\partial_\mu\chi)(x)\,R_{x,\mu}B_\nu(x+\eta e_\mu)
-(\partial_\nu\chi)(x)\,R_{x,\nu}B_\mu(x+\eta e_\nu),\\
(E'_\chi \Phi)_\nu(x)&:=-\sum_{\mu\neq\nu}(\partial^-_\mu\chi)(x)\,
R^{-1}_{x-\eta e_\mu,\mu}\,\Phi_{\mu\nu}(x-\eta e_\mu),
\end{split}
\end{equation}
with $\Phi$ a field on plaquettes. Since the commutator with $\chi$ is a derivation,
$[D^*_WD_W,\chi]=D^*_W[D_W,\chi]+[D^*_W,\chi]D_W$, and therefore
\begin{equation}\label{4.25:eq-cut-layer-partner-6}
[D^*_WD_W,\chi]B=D^*_W(E_\chi B)+E'_\chi(D_WB).
\end{equation}

\emph{The commutator: the coefficient count.} $(\ast)$ We expand $D^*_W(E_\chi B)$ once more by
the product rule \eqref{4.25:eq-cut-layer-partner-4}. At a bond $b$ of direction $\nu$, each
pair $(\mu,\nu)$ then contributes terms of first order in the differences of $\chi$ with
coefficients $1$, $1+2\tau$ and $2\tau$, and terms containing second differences of $\chi$ with
coefficients $1$ and $\tau^2$; all differences of $\chi$ that occur are forward differences
based at points of $\operatorname{st}(b)$, and all values of $B$ and of $\nabla_WB$ that occur
are read on $\operatorname{st}(b)$. Summing over the $d-1$ values of $\mu\neq\nu$ and using
$1+\tau^2\le 2+4\tau$ we obtain
\begin{equation}\label{4.25:eq-cut-layer-partner-7}
\bigl|([D^*_WD_W,\chi]B)(b)\bigr|\le c_d\Bigl[s_1(b)\sup_{\operatorname{st}(b)}|\nabla_WB|
+s_2(b)\sup_{\operatorname{st}(b)}|B|\Bigr],
\end{equation}
where $s_1(b)$ and $s_2(b)$ are the maxima over $\operatorname{st}(b)$ of the first and of the
second lattice differences of $\chi$, and $c_d$ is given by \eqref{4.25:eq-cut-layer-partner-a}.

\emph{The differences of the cut-off.} With $\chi=t+\mathsf{t}\,\hat\chi_{\square}$ every
difference of $\chi$ equals $\mathsf{t}$ times the same difference of $\hat\chi_\square$, the
constant $t$ being annihilated. A forward difference based at $z$ reads $\hat\chi_\square$
within $\{|\cdot-z|_\infty\le\eta\}$. Since $\hat\chi_\square$ vanishes wherever some coordinate
satisfies $|x_\nu-y_{\square,\nu}|\ge\frac23M\ell_\square$, such a difference vanishes unless
\[
|z-y_\square|_\infty\le\tfrac23M\ell_\square+2\eta\le M\ell_\square,
\]
the last inequality because $M\ell_\square\ge 18\eta$ (by \eqref{4.25:eq-collar-weight-b} and
$L\ge 3$ one has $M\ge 2L^2\ge 18$, and $\ell_\square\ge\eta$); that is, it vanishes unless
$z\in\square$. This is why the fifth entry carries the indicator of
$\tilde\Delta(y)\cap\square\neq\emptyset$ for the cube $\square$ itself: if
$\tilde\Delta(y)\cap\square=\emptyset$, then $s_1(b)=s_2(b)=0$ for every bond $b$ with
$b_-\in\Delta(y)$, since $\operatorname{st}(b)\subset\tilde\Delta(y)$ for such $b$, and the
commutator term vanishes there. By the mean-value property of difference quotients of the $C^2$
function $\hat\chi_\square$, every first difference is a first partial derivative of
$\hat\chi_\square$ at an intermediate point and every second difference a second partial
derivative at an intermediate point. With $0\le\zeta\le 1$ and $|\zeta'|,|\zeta''|\le 5$ the
first partial derivatives of $\hat\chi_\square$ are bounded by $5/(M\ell_\square)$, and the
second ones by $5/(M\ell_\square)^2$ (repeated index) or $25/(M\ell_\square)^2$ (distinct
indices). Hence
\begin{equation}\label{4.25:eq-cut-layer-partner-8}
s_1(b)\le\frac{5|\mathsf{t}|}{M\ell_\square},\qquad s_2(b)\le\frac{25|\mathsf{t}|}{(M\ell_\square)^2}.
\end{equation}

\emph{Comparison of the scales: $\ell_y\le L\ell_\square$.} We may assume
$\tilde\Delta(y)\cap\square\neq\emptyset$. The cube $\square$ contains a current cell
$c_\square$ of level $\operatorname{lev}\square$. Indeed, in the construction of the cover the
member $\zeta_\square$ of the partition of unity attached to $\square$ satisfies
$\operatorname{supp}\zeta_\square\subset\{|\cdot-y_\square|_\infty\le\frac23M\ell_\square\}$, so
a point of the layer at which $\zeta_\square\neq0$ lies in one of the $2^d$ cubes of side
$M\ell_\square$ of the partition $\pi_{\operatorname{lev}\square}$ that make up $\square$, and
that cube is then a cell of the layer; the cover at the cut cube is of the same kind as the cover
used at a stage of the step, \cite[\S3]{B-CMP109}, \cite[p.~276]{B-CMP119}. Let $c_y$ be the
current cell containing $y$; its level is that of the block, so that
$\ell_y=\ell_{\operatorname{lev}c_y}$.
Suppose that $\operatorname{lev}c_y\ge\operatorname{lev}\square+2$. Then
\eqref{4.25:eq-collar-weight-a} gives
\[
\operatorname{dist}_\infty(c_y,c_\square)\ge M\ell_{\operatorname{lev}c_y-1}=M\ell_y/L.
\]
On the other hand let $z\in\tilde\Delta(y)\cap\square$. Since $y\in c_y$, the point $z$ lies
within $\ell_y$ of $c_y$. Since $z\in\square$ and $\square$ has side $2M\ell_\square$, the point
$z$ lies within $M\ell_\square$ of every sub-cube of $\square$ of side $M\ell_\square$, in
particular of $c_\square$, whichever cube of $\pi_{\operatorname{lev}\square}$ in $\square$
contains $z$; and $M\ell_\square\le M\ell_y/L^2$ because
$\operatorname{lev}c_y\ge\operatorname{lev}\square+2$. Hence
\[
\operatorname{dist}_\infty(c_y,c_\square)\le\ell_y+\frac{M\ell_y}{L^2}<\frac{M\ell_y}{L},
\]
where the strict inequality is equivalent to $L^2+M<ML$, which holds under
\eqref{4.25:eq-collar-weight-b}, $M\ge 2L^2$ (indeed $M(L-1)\ge 2L^2(L-1)>L^2$ for $L\ge3$).
The two bounds on $\operatorname{dist}_\infty(c_y,c_\square)$ contradict each other; therefore
$\operatorname{lev}c_y\le\operatorname{lev}\square+1$ and $\ell_y\le L\ell_\square$.

\emph{The weighted differences.} From \eqref{4.25:eq-cut-layer-partner-8}, $\ell_y\le
L\ell_\square$ and $M\ge 2L^2$,
\[
\ell_y s_1(b)\le|\mathsf{t}|\,\frac{5L}{M}\le|\mathsf{t}|\,\frac{5}{2L},\qquad
\ell_y^2 s_2(b)\le|\mathsf{t}|\,\frac{25L^2}{M^2}\le|\mathsf{t}|\,\frac{25}{4L^2},
\]
and, since $L\ge3$,
\begin{equation}\label{4.25:eq-cut-layer-partner-9}
\ell_y s_1(b)+\ell_y^2 s_2(b)\le|\mathsf{t}|\Bigl(\frac56+\frac{25}{36}\Bigr)
=\frac{55}{36}\,|\mathsf{t}|<2|\mathsf{t}|.
\end{equation}
Each of the two terms is in fact less than $|\mathsf{t}|$, but this is not used; only the
finiteness of the bound enters $c_\zeta$.

\emph{The commutator term.} Put $B:=H(\sigma)X$. For a bond $b$ with $b_-\in\Delta(y)$ one has
$\operatorname{st}(b)\subset\tilde\Delta(y)$, and
$\ell_y^2\sup_{\operatorname{st}(b)}|\nabla_WB|\le N_y(B)$,
$\ell_y\sup_{\operatorname{st}(b)}|B|\le N_y(B)$. Multiplying
\eqref{4.25:eq-cut-layer-partner-7} by $\ell_y^3$ and using
\eqref{4.25:eq-cut-layer-partner-9}, together with the vanishing established above when
$\tilde\Delta(y)\cap\square=\emptyset$,
\begin{align*}
\ell_y^3\,\bigl|([D^*_WD_W,\chi]H(\sigma)X)(b)\bigr|
&\le c_d\bigl[\ell_ys_1(b)\cdot\ell_y^2\sup_{\operatorname{st}(b)}|\nabla_WB|
+\ell_y^2s_2(b)\cdot\ell_y\sup_{\operatorname{st}(b)}|B|\bigr]\\
&\le c_d\cdot 2|\mathsf{t}|\,N_y(H(\sigma)X)\,\mathbf{1}[\tilde\Delta(y)\cap\square\neq\emptyset],
\end{align*}
which is $c_\zeta$ times the fifth entry, by \eqref{4.25:eq-cut-layer-partner-a}.

\emph{Summation.} Adding the bounds obtained for the terms of
\eqref{4.25:eq-cut-layer-partner-1},
\begin{align*}
\ell_y^3|\mathcal{J}^{\sharp}[\chi]|
&\le(1+1+18d)\max\{\text{first entry},\ \text{third entry}\}+\text{fourth entry}\\
&\quad+9K_J\cdot\text{second entry}+c_\zeta\cdot\text{fifth entry}\\
&\le K_{\mathsf{j}}\max\{\text{first to fourth entries}\}+c_\zeta\cdot\text{fifth entry},
\end{align*}
because $(1+1+18d)+1+9K_J=3+18d+9K_J=K_{\mathsf{j}}$ by \eqref{4.25:eq-current-partner-a}.
This is the first inequality of \eqref{4.25:eq-cut-layer-partner-3}; the second holds because
each entry is at most $\mathsf{N}^{\chi}_y$.
\end{proof}

\begin{remark}
The fifth entry of \eqref{4.25:eq-cut-layer-partner-2} is controlled by the uncut layer. The
layer datum satisfies $\mathbb{H}=A'-H(\sigma)X$, so
$N_y(H(\sigma)X)\le N_y(A')+N_y(\mathbb{H})$, two entries of the uncut layer. Writing
$\mathbf{1}_\square:=\mathbf{1}[\tilde\Delta(y)\cap\square\neq\emptyset]$,
\eqref{4.25:eq-cut-layer-partner-3} then gives the bound with one explicit constant,
\[
\ell_y^3|\mathcal{J}^{\sharp}[\chi]|\le(K_{\mathsf{j}}+2c_\zeta)\max\bigl\{
\text{first to fourth entries},\ |\mathsf{t}|N_y(A')\mathbf{1}_\square,\
|\mathsf{t}|N_y(\mathbb{H})\mathbf{1}_\square\bigr\}.
\]
At the seams of the cover, \cite[p.~276]{B-CMP119} specifies only the corresponding partition of
unity.
If $\hat\chi_\square$ is replaced there by the normalised cut-off
$\zeta_\square/\sum_{\square'}\zeta_{\square'}$, the support is unchanged and, by the quotient
rule, the bounds $5$ and $25$ in \eqref{4.25:eq-cut-layer-partner-8} are replaced by numbers
depending only on $d$, measured in the finer unit; only $c_\zeta$ changes. Likewise, if the
transport is placed differently, $D'_W=D_W+\mathsf{V}$ with $\mathsf{V}$ an operator of order zero
on the plaquette stencil, then $c_d$ changes by a number depending only on $d$ times
$\mathsf{p}(W)$, since every commutator with $\chi$ carries a difference of $\chi$ across a
stencil of size $\eta$.
\end{remark}

\begin{definition}[The increment hypothesis at later presentations]\label{4.25:def-increment-hypothesis}
Let $\mathfrak{a}$ be the arrest, with region $W$ and step $k_W$. Let $\mathfrak{Q}(\mathfrak{a})$ be the
family of pieces $e$ of the quadratic collar weight $\mathsf{q}_W$, with functions $v_e$. For a piece $e$
write $S^+(e)$ for the $2$-neighbourhood of its support $\tilde S(e)$, and $\mathcal{Y}(e)$ for the
set of current blocks $y$ meeting $S^+(e)$. The later operations include the six entries $O1$, $O2$,
$O3$, $O4$, $O4'$ and $O5$ of the list of later presentations fixed together with the arrest
theorem. The sites listed at the end of this definition say at which of these entries each part
of the hypothesis is used. Let $\mathfrak{O}$ be a later operation, taken at a
step $k'\ge k_W$; any of the six entries is allowed. Let $\mathfrak{T}$ be the space presented
by $\mathfrak{O}$, and let
$\mathcal{Z}_{\mathfrak{O}}$ be the parameter set of its family. Throughout, $\eta$ is the
lattice spacing of the presentation.

The \emph{increment hypothesis} for $\mathfrak{Q}(\mathfrak{a})$ is the following requirement, for
every piece $e$ of $\mathfrak{Q}(\mathfrak{a})$.
At every presentation of $\mathfrak{Q}(\mathfrak{a})$ by such an $\mathfrak{O}$, and for every
member $z\in\mathcal{Z}_{\mathfrak{O}}$ of its family, the pair over $S^+(e)$ has the form
\begin{equation}\label{4.25:eq-increment-hypothesis-1}
p(z)=\bigl(e^{i\eta A''(z)}\,\mathbb{U}_{\mathfrak{O}},\ \mathbb{J}_{\mathfrak{O}}+J''(z)\bigr).
\end{equation}
Here the following holds.
\begin{itemize}
\item The pair is built over one base $P_{\mathfrak{O}}\in\mathfrak{T}$, with gauge field
$\mathbb{U}_{\mathfrak{O}}$ and current $\mathbb{J}_{\mathfrak{O}}$, and this base does not
depend on $z$.
\item The base $P_{\mathfrak{O}}$ is the subtracted reference pair of the presentation.
\item The increments $(A'',J'')$ are holomorphic in $z$.
\end{itemize}
Moreover
\begin{equation}\label{4.25:eq-increment-hypothesis-a}
\begin{gathered}
N_y\bigl(A''(z)\bigr)\le C_{inc}\,\theta(y),\qquad
\ell_y^3\,\bigl\lvert J''(z)\bigr\rvert\le C_{inc}\,\theta(y)
\qquad\bigl(y\in\mathcal{Y}(e),\ z\in\mathcal{Z}_{\mathfrak{O}}\bigr),\\
C_{inc}:=2\,\bigl(K_{\mathsf{j}}+2c_\zeta\bigr)\,\hat B_H\,C'_{age}.
\end{gathered}
\end{equation}
The symbols in \eqref{4.25:eq-increment-hypothesis-a} are as follows.
\begin{itemize}
\item $\theta(y)$ is the block threshold defined in \eqref{4.25:eq-flow-inequalities-f}, among
the coupling-flow inequalities and block thresholds (Notation~\ref{4.25:not-flow-inequalities}).
\item $K_{\mathsf{j}}$ is the constant defined in \eqref{4.25:eq-current-partner-a}, in the
current partner of a difference of layers (Lemma~\ref{4.25:lem-current-partner}).
\item $c_\zeta=2c_d$ is the constant defined in \eqref{4.25:eq-cut-layer-partner-a}, in the
current partner of a cut layer (Lemma~\ref{4.25:lem-cut-layer-partner}), whose proof is
incomplete at one step.
\item $\hat B_H:=\max\{1,B_H\}$, where $B_H$ is the tube constant.
\item $C'_{age}$ is the age constant defined in the increment lemma below.
\end{itemize}
At every statement that uses the hypothesis, each of the three region parameters on which that
statement is read (written $S$, $S_v$ and $X$ there) is taken equal to $S^+(e)$.

The sites at which the hypothesis is used, and the statements that supply it there, are the
following.
\begin{itemize}
\item[(P)] The operations $O1$, $O4$, $O4'$. The parameter $t=(t_\square)$ ranges over the
product of the closed discs $\prod_\square\{\lvert t_\square\rvert\le r_\square\}$ of radii
$r_\square$, and the reference is $t=0$. The presented pair is
\begin{equation}\label{4.25:eq-increment-hypothesis-2}
\Phi_t=\bigl(e^{i\eta\chi_t\mathbb{H}}\,\mathbb{U},\ \mathbb{J}+\chi_t\,\mathcal{J}^{\sharp}[1]\bigr).
\end{equation}
Here $\chi_t$ is the cut-off at parameter $t$. The current $\mathcal{J}^{\sharp}[1]$ is the
current $\mathcal{J}^{\sharp}[\chi]$ defined in the statement of
Lemma~\ref{4.25:lem-cut-layer-partner} (whose proof is incomplete at one step), taken at
$\chi\equiv1$. The second component is the local one, $\chi_t\mathcal{J}^{\sharp}[1]$, and not
$\mathcal{J}^{\sharp}[\chi_t]$.

The lemma bounding the cut layer and its local current at the presentations $O1$, $O4$, $O4'$
supplies the following bounds, for all blocks at once and at every location,
the excluded region $W$ included:
\begin{equation}\label{4.25:eq-increment-hypothesis-3}
\begin{gathered}
N_y(\chi_t\mathbb{H})\le\mathsf{k}_{\mathsf{w}}(y)
\le\mathsf{w}_0\,\varsigma_{m(y)}\,\alpha_{\cdot,m(y)},\\
\ell^3\,\bigl\lvert\chi_t\,\mathcal{J}^{\sharp}[1]\bigr\rvert
\le\tfrac{11}{2^9}\,\mathsf{w}_0\,\varsigma_{m(y)}\,\alpha_{0,m(y)},\qquad m(y)=\ell'(y).
\end{gathered}
\end{equation}
In this display:
\begin{itemize}
\item $\alpha_{\cdot,m(y)}$ denotes either of the radii $\alpha_{0,m(y)},\alpha_{1,m(y)}$;
\item $\mathsf{k}_{\mathsf{w}}(y)$ is the bound that lemma gives for the block size of the cut
layer on the block $y$; $\mathsf{w}_0$ is the amplitude constant of that lemma's bound on the
layer, and $\varsigma_{m(y)}$ its level factor at the level $m(y)=\ell'(y)$ of the block $y$;
all three are fixed in that lemma.
\end{itemize}
On its own radii, the uniform bound on the increment size $\mathfrak{q}$ gives, uniformly over
shells and base points,
\begin{equation*}
\mathfrak{q}\le\mathfrak{q}_*:=16\,N_{cov}\,c_R^0 ,
\end{equation*}
where $N_{cov}$ is the covering number and $c_R^0$ the constant in the radii of that bound.
This is at most one quarter of the unit of room of that statement, for the gauge-field component
and for the current component of the pair alike.

At this site the bounds \eqref{4.25:eq-increment-hypothesis-a}, with the constant $C_{inc}$, are
obtained from \eqref{4.25:eq-increment-hypothesis-3}; this rests on the inequality
$\mathsf{w}_0\varsigma_{m(y)}\le 5C_{inc}$, which is established for the presentations
$O1$, $O4$, $O4'$ by the statement that delivers the bounds of this site with the constant
$C_{inc}$ (the delivery lemma of the (P) site), and the choice of $C_{inc}$ as the enlarged
value in \eqref{4.25:eq-increment-hypothesis-a} only weakens this requirement.
\item[(T)] The operation $O3$. The hypothesis is supplied by the increment lemma for the (T) and
far (L) sites, which treats this site and the far (L) site together.
\item[\textup{(L), far}] The operation $O2$, with $S^+(e)\cap Z=\emptyset$. The hypothesis is
supplied by the same increment lemma.
\item[\textup{(L), deep}] The operation $O5$, with $S^+(e)\subset\mathsf{C}_C$, the core of the
healer. The hypothesis is supplied by the inputs of the three-field bound.
\item[\textup{(L), other}] For the other (L)-pieces, the norm bound $2\nu^{\natural}(e)$ of the
shell estimate, on the domain fixed by the domain condition of this chapter, supplies the
hypothesis.
\end{itemize}
On $W$ itself, and with no dilation applied, the hypothesis agrees with the following bounds:
\begin{itemize}
\item the third of the bounds that the layer lemma of the arrest theorem lists among the
remaining ones, a bound linear in the layer;
\item the two displayed bounds of the statement of which the current partner of a cut layer
(Lemma~\ref{4.25:lem-cut-layer-partner}, whose proof is incomplete at one step) is one
clause.
\end{itemize}
\end{definition}

We work in the setting of the arrested region and its collar weight
(Notation~\ref{4.25:not-collar-weight}). The general reading $y\in S$, meaning $\Delta(y)\subset S$,
is not the right one in what follows. Suppose a
cell of $S^+(e)$ has level $m$ while a current block $y$ meeting it has level $\ell'(y)$
with $L^{\ell'(y)-m}>M$. The side $\ell_y=\ell_{\ell'(y)}=L^{\ell'(y)-m}\ell_m$ of $\Delta(y)$
(Notation~\ref{4.25:not-collar-weight}) then exceeds the side $M\ell_m$ of the cell, so the block
is coarser than the cell itself. Where this happens for the cells of $S^+(e)$, no current block
lies inside $S^+(e)$. Under the general reading, the supremum defining $\Theta(e)$, taken
over the blocks $y\in S^+(e)$, would then range over an empty set. The increment size
$\mathfrak{a}''$, a supremum over the blocks written at earlier steps, is consumed through
$\Theta(e)$ and would then likewise carry an empty hypothesis.
Every hypothesis quantified over such blocks would then be vacuous. Under the reading
$y\in\mathcal{Y}(e)$ of Definition~\ref{4.25:def-enlarged-support-blocks} below, the blocks that
meet $S^+(e)$, and in particular the coarser current blocks over such cells, are the ones at which
the hypotheses are imposed.

\begin{definition}[Blocks meeting the enlarged support of a piece]\label{4.25:def-enlarged-support-blocks}
We work in the setting of the arrested region and its collar weight
(Notation~\ref{4.25:not-collar-weight}). The regions $\Delta(y)$ and $S^+(e)$ below are read as
sets of sites; bonds are attached to sites as in \cite{B-CMP96}. For a current block $y$ let
$\ell'(y)$ denote its current level. The
block $\Delta(y)$ is the set of those sites whose block of level $\ell'(y)$ is the block of $y$;
in Ba{\l}aban's notation this is the block $B^{j}(y)$ with $j=\ell'(y)$.
For a piece $e\in\mathfrak{E}^G_W$, the family of cells $S^+(e)$ of
Notation~\ref{4.25:not-collar-weight}, that is, the cells of $\tilde S(e)$ together with the cells
at cell-distance at most $2$ from them, is identified with the set of the sites of its cells. The
blocks $\Delta(y)$ of the current blocks $y$ partition the set of sites, and the partitions into
current blocks at successive steps are nested
(Notations~\ref{4.25:not-collar-weight} and~\ref{4.25:not-flow-inequalities}). For
a site $p$ we write $B^{top}(p)$ for the unique current block $y$ with $p\in\Delta(y)$.

The class of current blocks attached to $e$ is
\begin{equation}\label{4.25:eq-enlarged-support-blocks-1}
  \begin{aligned}
    \mathcal{Y}(e)
      &:= \{\,y \text{ a current block} : \Delta(y)\cap S^+(e)\neq\emptyset\,\}\\
      &\phantom{:}= \{\,B^{top}(p) : p \text{ a site of } S^+(e)\,\}.
  \end{aligned}
\end{equation}
It consists of the current blocks that meet $S^+(e)$. Its definition uses the block $\Delta(y)$ itself,
not the doubled stencil $\tilde\Delta(y)$.

The envelope $\Theta(e)$ of Notation~\ref{4.25:not-flow-inequalities}, built from the block
thresholds $\theta(y)$ of that notation, is, with the class
\eqref{4.25:eq-enlarged-support-blocks-1} now defined, the supremum over it:
\begin{equation}\label{4.25:eq-enlarged-support-blocks-2}
  \Theta(e) = \sup_{y\in\mathcal{Y}(e)} \theta(y).
\end{equation}
The following reading convention is in force in the definition
\eqref{4.25:eq-enlarged-support-blocks-2} of $\Theta(e)$ and, further, in the base comparison of
this section, in the increment bound and in its cut form, in the rotated-increment bound, and in
items (b) and (e) of the lemma on the current blocks over the cells of a piece.
There the condition $y\in S^+(e)$
means $y\in\mathcal{Y}(e)$. It does not mean $\Delta(y)\subset S^+(e)$, which is the general reading
of $y\in S$ fixed in Notation~\ref{4.25:not-collar-weight}.
\end{definition}

\begin{proof}[Proof of the equality in~\eqref{4.25:eq-enlarged-support-blocks-1}]
Let $p$ be a site of $S^+(e)$ and put $y=B^{top}(p)$. By the definition of $B^{top}(p)$, the block
$y$ is a current block with $p\in\Delta(y)$. Hence $p\in\Delta(y)\cap S^+(e)$, this intersection is
non-empty, and $y\in\mathcal{Y}(e)$.

Conversely, let $y\in\mathcal{Y}(e)$ and choose a site $p\in\Delta(y)\cap S^+(e)$. Then $y$ is a
current block with $p\in\Delta(y)$. By the uniqueness of the current block whose block contains $p$,
we get $y=B^{top}(p)$, with $p$ a site of $S^+(e)$.

The two sets in~\eqref{4.25:eq-enlarged-support-blocks-1} therefore coincide.
\end{proof}

\begin{lemma}[Domination of old blocks by current blocks]\label{4.25:lem-domination-old-blocks}
Let $e$ be a piece, with enlarged support $S^+(e)$ and with the set $\mathcal{Y}(e)$ of current blocks
meeting it, as in Definition~\ref{4.25:def-enlarged-support-blocks}. Cells, blocks, their levels, the
units $\ell_m$ and the top block $B^{top}(p)$ of a site $p$ are those of the setting of
Notation~\ref{4.25:not-collar-weight}. For a cell $c$ we write $\hat c$ for its \emph{hat}, the current
cell containing $c$. For a block $y$ of side $\ell_y$, a $1$-form $A$ and a connection $V$, recall that
\begin{equation*}
N_y(A)=\sup_{\tilde\Delta(y)}\max\bigl\{\ell_y|A|,\ \ell_y^2|\nabla_V A|\bigr\},
\end{equation*}
the covariant derivative being taken in the connection $V$.
\begin{enumerate}
\item[(i)] The current cells of the blocks of $\mathcal{Y}(e)$ are exactly the hats of the cells of
$S^+(e)$.
\item[(ii)] Let $y_0$ be an old block of level $m$ with $\Delta(y_0)\subset S^+(e)$, and put
$y_1:=B^{top}(y_0)$. Then $y_1\in\mathcal{Y}(e)$, the level $m'$ of $y_1$ satisfies $m'\ge m$, and
$\tilde\Delta(y_0)\subset\tilde\Delta(y_1)$. Hence, for every $1$-form $A$ and every connection $V$,
\begin{equation}\label{4.25:eq-domination-old-blocks-1}
N_{y_0}(A)\le N_{y_1}(A),
\end{equation}
where the left side is computed in the old units $\ell_m$ and the right side in the current units
$\ell_{m'}$.
\item[(iii)] Every block of $\mathcal{Y}(e)$ has the level of its current cell, and that cell is the
hat of a cell of $S^+(e)$. Consequently the level bound
\eqref{4.25:eq-pile-bound-a} on the hats, and every item of the conversion of old polylogarithmic
factors, \eqref{4.25:eq-conversion-old-factors-a}, hold verbatim with the blocks $y\in\mathcal{Y}(e)$ in
place of the blocks $y$ with $\Delta(y)\subset S^+(e)$. None of their constants changes.
\end{enumerate}
\end{lemma}

\begin{proof}
\emph{(i).} The current cells partition the sites, and $S^+(e)$ is the set of sites of its cells.

Let $y\in\mathcal{Y}(e)$. By Definition~\ref{4.25:def-enlarged-support-blocks}, $y=B^{top}(p)$ for some
site $p$ of $S^+(e)$. This site lies in an old cell $c\subset S^+(e)$. The partitions into cells are
nested (Notation~\ref{4.25:not-collar-weight}), so $p$ lies in the hat $\hat c$, and
$B^{top}(p)\subset\hat c$. Thus the current cell of $y$ is the hat of a cell of $S^+(e)$.

Conversely, let $c$ be a cell of $S^+(e)$. Its hat $\hat c$ contains $c$, so it holds a site $p$ of
$S^+(e)$. The block $B^{top}(p)$ belongs to $\mathcal{Y}(e)$, again by
Definition~\ref{4.25:def-enlarged-support-blocks}, and its current cell is $\hat c$. This proves (i).

\emph{(ii).} The site $y_0$ lies in $\Delta(y_0)$, hence in $S^+(e)$. Therefore
$y_1=B^{top}(y_0)\in\mathcal{Y}(e)$, by the second description of $\mathcal{Y}(e)$ in
Definition~\ref{4.25:def-enlarged-support-blocks}. Levels never decrease
(Notation~\ref{4.25:not-collar-weight}), so $m'\ge m$.

The grids are nested. Hence $y_0$ is a point of the lattice of spacing $\ell_m$ that lies in the block
of side $\ell_{m'}$ at $y_1$. Whether blocks are indexed by their corner or by their centre, such a
lattice point is at offset at most $\ell_{m'}-\ell_m$ from the index point in each coordinate, so
$|y_0-y_1|_\infty\le\ell_{m'}-\ell_m$.
Now let $z\in y_0+[-\ell_m,\ell_m]^d$. By the triangle inequality,
\begin{equation*}
|z-y_1|_\infty\le|z-y_0|_\infty+|y_0-y_1|_\infty\le\ell_m+(\ell_{m'}-\ell_m)=\ell_{m'} .
\end{equation*}
Hence
\begin{equation*}
y_0+[-\ell_m,\ell_m]^d\subset y_1+[-\ell_{m'},\ell_{m'}]^d .
\end{equation*}
In Ba{\l}aban's construction $\tilde\Delta(y)$ is the cube of size $2\ell_y$ centred at $y$
\cite{B-CMP99b}. The last inclusion is therefore $\tilde\Delta(y_0)\subset\tilde\Delta(y_1)$.

Since $\ell_{m'}=L^{m'-m}\ell_m$ and $m'\ge m$, we have $\ell_m\le\ell_{m'}$. A supremum over a larger
set is not smaller, so
\begin{align*}
\sup_{\tilde\Delta(y_0)}\ell_m|A|&\le\sup_{\tilde\Delta(y_1)}\ell_{m'}|A|,\\
\sup_{\tilde\Delta(y_0)}\ell_m^2|\nabla_V A|&\le\sup_{\tilde\Delta(y_1)}\ell_{m'}^2|\nabla_V A|.
\end{align*}
By definition, $N_y(A)$ is the larger of the two suprema
$\sup_{\tilde\Delta(y)}\ell_y|A|$ and $\sup_{\tilde\Delta(y)}\ell_y^2|\nabla_V A|$. Taking the maximum of
the two inequalities therefore gives \eqref{4.25:eq-domination-old-blocks-1}.

\emph{(iii).} The level bound \eqref{4.25:eq-pile-bound-a} bounds the levels of hats. These levels are
at most $j_W+n(e)+2$, and their spread is at most $n(e)+3$. By (i), the hats are exactly the current
cells of the blocks of $\mathcal{Y}(e)$, so this bound applies unchanged to the blocks of
$\mathcal{Y}(e)$.

The first three items of the conversion \eqref{4.25:eq-conversion-old-factors-a} evaluate the inverse
coupling $x_k$ at the level $k$ of a block $y$. For $y\in\mathcal{Y}(e)$ this is the level of a hat,
by (i). These items are therefore unchanged.

Here it matters that $\mathcal{Y}(e)$ is selected by $\Delta(y)$ meeting $S^+(e)$ on sites, and not by
$\tilde\Delta(y)$ meeting it. A block whose stencil $\tilde\Delta(y)$ merely touches a hat from outside
does not belong to $\mathcal{Y}(e)$. Hence no level other than a hat's level enters, and no constant of
\eqref{4.25:eq-pile-bound-a} or \eqref{4.25:eq-conversion-old-factors-a} moves.
\end{proof}

\begin{remark}
Recall that the increment size $\mathfrak{a}''$ of the statement on Gaussian pieces is the
supremum, in
old units, of $N_y(A'')$ over the old blocks $y$ with $\Delta(y)\subset S^+(e)$, $A''$ being the
increment $1$-form of the increment hypothesis \eqref{4.25:eq-increment-hypothesis-a}. Part (ii) of
Lemma~\ref{4.25:lem-domination-old-blocks} and the increment hypothesis on $\mathcal{Y}(e)$ give
\begin{equation*}
\mathfrak{a}''\le\sup_{y\in\mathcal{Y}(e)}N_y(A'')\le C_{inc}\,\Theta(e).
\end{equation*}
This bound requires no condition on the size of $C_{inc}$.

A condition on the size of $C_{inc}$ enters only through the hypothesis on the base $1$-form $A'$, not
through the increment hypothesis: namely through the two base clauses of the statement on Gaussian
pieces, those resting on \eqref{4.25:eq-gaussian-pieces-a}. There
\begin{equation*}
N_y(A')<\theta(y)\le\Theta(e)\le\tfrac{8}{9}\,\varpi_G/C_{inc}.
\end{equation*}
This asks for $C_{inc}\ge 2/9$ in the first of these clauses and $C_{inc}\ge 8/9$ in the second.
\end{remark}

\begin{lemma}[Increment bound at tube and far sites]\label{4.25:lem-increment-bound}
Let $e$ be a piece of the family $\mathfrak{Q}(\mathfrak{a})$, let $\mathcal{Y}(e)$ be the set of current
blocks of Definition~\ref{4.25:def-enlarged-support-blocks}, let $S_v=S^+(e)$ as in
Definition~\ref{4.25:def-increment-hypothesis}, and let $k'$ be a step. Let $\mathfrak{s}$ be a site of
one of the following two kinds: the tube site of step $k'$, labelled (T), or a far site, labelled
(L)-far. Each kind fixes a region $\mathfrak{F}_{\mathfrak{s}}$, from which the distance $d_{sc}$ below is
measured, a decay rate $\delta_{\mathfrak{s}}>0$, a unit $\mathsf{U}_{\mathfrak{s}}$ and a radius
$\mathsf{w}_{\mathfrak{s}}\ge0$.
\begin{itemize}
\item[(T)] $\mathfrak{s}$ is the tube site of step $k'$. Then $\mathfrak{F}_{\mathfrak{s}}$ is the reference
region of the tube site, $\delta_{\mathfrak{s}}:=\delta_1^{sc}$ is the scaled decay rate of the tube,
$\mathsf{U}_{\mathfrak{s}}:=\mathsf{U}_T$, and there is one layer.
\item[(L)-far] $\mathfrak{s}$ is a far site, that is, a site of the large-field kind whose treated
large-field region $Z$ satisfies $S^+(e)\cap Z=\emptyset$. Then $\mathfrak{F}_{\mathfrak{s}}$ is the enlarged
reference region of the far site, $\delta_{\mathfrak{s}}:=\delta_\natural$ is the decay rate of the far site,
$\mathsf{U}_{\mathfrak{s}}:=\mathsf{U}_L$, where
\begin{equation*}
\mathsf{U}_L(x):=C_LK_{\mathfrak{Q}}\,\alpha^+(x)/K_u ,
\end{equation*}
with $C_L$ and $K_{\mathfrak{Q}}$ constants fixed with the far site and $\alpha^+$ a function of the inverse
coupling $x$; the local constant $c_\chi$ of the cut (a number depending on $d$) is placed inside
$\mathsf{U}_L$, so that condition~(i) of \eqref{4.25:eq-frozen-collar-weight-a} is read with the same
$\mathsf{U}_L$; that display, which fixes the frozen Gaussian collar weight $\varpi_G$ of an arrested
region and the conditions imposed on it, is printed later in this section. There are $r\le4$ layers,
the links $a_1,\dots,a_r$ of the far site.
\end{itemize}
The base pair is $(\mathbf{U},\mathbf{J})$. The layers are $\mathcal{K}_i(z)$, $z=(\mathsf{t},\sigma)$; exactly
one of them depends on $\mathsf{t}$, and it does so linearly \cite{B-CMP109}. The letters of the layers are
as follows.

At (T), every $\mathcal{K}_i$ is a layer datum of the kind to which
Lemma~\ref{4.25:lem-current-partner} applies, with partner current $\mathfrak{s}^{\sharp}[\mathcal{K}_i]$. The cut
acts on the datum, $A'\mapsto A'+t_\square\delta A'$, $X\mapsto X+t_\square\delta X$, and
Lemma~\ref{4.25:lem-current-partner} applies to the layers so obtained. The parameter $\mathsf{t}$ of the
layers is the cut parameter $t_\square$. We write $\delta\mathbb{H}_p$ for the derivative in
$t_\square$, and $\mathsf{N}_y$ for the size with the four entries of
Lemma~\ref{4.25:lem-current-partner}.

At (L)-far, the lower layers are layer data with partner currents $\mathfrak{s}^{\sharp}$ as at (T). The
$\mathsf{t}$-holding layer is the field cut
\begin{equation}\label{4.25:eq-increment-bound-1}
\mathcal{K}(z)=\chi\,\mathbb{H}_s,\qquad \chi:=t+\mathsf{t}\,\hat\chi_{\square},\qquad t\in[0,1],
\end{equation}
where $\hat\chi_{\square}$ is the cut-off of the response cube $\square$ \cite{B-CMP109}, as in
Lemma~\ref{4.25:lem-cut-layer-partner}, whose proof is incomplete at one step. This layer is a layer
function attached to no datum. For it we set $\delta\mathbb{H}_p:=\hat\chi_{\square}\mathbb{H}_s$. Its partner
current is $\mathcal{J}^{\sharp}[\chi]$ of Lemma~\ref{4.25:lem-cut-layer-partner}, and its size is the
five-entry size $\mathsf{N}^{\chi}_y$ of that lemma. We further set
\begin{equation}\label{4.25:eq-increment-bound-2}
\mathsf{N}_y(\delta\mathbb{H}_p):=\max\bigl\{N^{\Delta}_y(\hat\chi_{\square}A'_s),\ N_y(\hat\chi_{\square}\mathbb{H}_s),\
\ell^3|P_1(\sigma)(\hat\chi_{\square}A'_s)|,\ \ell^3|\hat\chi_{\square}\mathfrak{m}_H(\sigma)X_s|\bigr\}.
\end{equation}
These are four cut entries: $\hat\chi_{\square}$ stands inside $\Delta_W$, $\nabla_W$ and $P_1(\sigma)$. The third
entry is the non-local tail of $P_1(\sigma)$, and the other three are local multiples of the uncut
entries. Write $\mathsf{N}_y(\mathbb{H}_s)$ for the four-entry size of the uncut datum
$(A'_s,X_s,\mathbb{H}_s)$. Then each of the first four entries of $\mathsf{N}^{\chi}_y(\mathcal{K}(z))$ obeys
\begin{equation}\label{4.25:eq-increment-bound-3}
\text{(entry $j$ of $\mathsf{N}^{\chi}_y(\mathcal{K}(z))$)}\le\mathsf{N}_y(\mathbb{H}_s)+|\mathsf{t}|\,
\mathsf{N}_y(\delta\mathbb{H}_p),\qquad j=1,2,3,4 .
\end{equation}

In either kind, the $\mathsf{t}$-holding layer, being linear in $\mathsf{t}$, has the form
$\mathcal{K}=\mathcal{K}^{0}+\mathsf{t}\,\delta\mathbb{H}_p$ with $\mathcal{K}^{0}$ independent of $\mathsf{t}$; at
(L)-far, $\mathcal{K}^{0}=t\,\mathbb{H}_s$ by \eqref{4.25:eq-increment-bound-1}. We call $\mathcal{K}^{0}$ the
undilated part and $\mathsf{t}\,\delta\mathbb{H}_p$ the dilated part of this layer. The
undilated layers are the layers not holding $\mathsf{t}$ and the undilated part of the $\mathsf{t}$-holding
layer; at (L)-far the latter is represented by $\mathbb{H}_s$.

For $y\in\mathcal{Y}(e)$ and $|\sigma|\le e^{\kappa_1}$ write $\ell'(y)\le k'$ for the level of the current
block $y$, put
$\mathfrak{d}:=d_{sc}(y,\mathfrak{F}_{\mathfrak{s}})\ge0$, and consider the inequalities
\begin{equation}\label{4.25:eq-increment-bound-a}
\begin{aligned}
&\text{(e1)}\quad \mathsf{N}_y(\delta\mathbb{H}_p)\le\eta_p\,\mathsf{U}_{\mathfrak{s}},\\
&\text{(e2)}\quad \mathsf{N}_y(\delta\mathbb{H}_p)\ \text{and}\ \mathsf{N}_y(\mathcal{K})\
\text{for every undilated layer } \mathcal{K}\ \text{are}\
\le\mathsf{U}_{\mathfrak{s}}e^{-\delta_{\mathfrak{s}}(\mathfrak{d}-1)},\\
&\text{(e3)}\quad (1+\mathsf{w}_{\mathfrak{s}})\,\mathsf{U}_{\mathfrak{s}}\,
e^{-\frac14\delta_{\mathfrak{s}}(\mathfrak{d}-1)^+}\le 2\hat B_H\,\bar\alpha(x_{k'}),\\
&\text{(e4)}\quad |\mathsf{t}|\,N_y(H(\sigma)X_s)\,\mathbf{1}[\tilde\Delta(y)\cap\square\neq\emptyset]
\le 2\mathsf{w}_{\mathfrak{s}}\mathsf{U}_{\mathfrak{s}}e^{-\delta_{\mathfrak{s}}(\mathfrak{d}-1)} .
\end{aligned}
\end{equation}
Assume:
\begin{itemize}
\item (e1), (e2) and (e3) of \eqref{4.25:eq-increment-bound-a} hold at every $y\in\mathcal{Y}(e)$ and every
$|\sigma|\le e^{\kappa_1}$;
\item $\eta_p:=e^{-\delta_1M\operatorname{dist}(\square,S_v)}$, where $\delta_1$ is a fixed decay constant and
the distance is counted in cubes, so that
$\operatorname{dist}(\square,S_v)=0$ if and only if $\tilde\square\cap S_v\neq\emptyset$;
\item the coupling-flow inequality \eqref{4.25:eq-flow-inequalities-c} of
Definition~\ref{4.25:not-flow-inequalities}, for $\bar\alpha$;
\item for every $y\in\mathcal{Y}(e)$, every contour from $y$ to $\mathfrak{F}_{\mathfrak{s}}$ (a chain of cells
along which $d_{sc}(y,\mathfrak{F}_{\mathfrak{s}})$ is measured) ends in cells of level $k'$ and crosses the
levels $\ell'(y)+1,\dots,k'-1$, which are consecutive, each over at least one $M$-cube of its own
scale;
\item $\mathsf{p}(W)\le1$ at the kernel $W$ of each layer (at the far site this is supplied by the
three-field bound, at the tube site by the construction of the tube).
\end{itemize}
Define the disc
\begin{equation}\label{4.25:eq-increment-bound-c}
|\mathsf{t}|\le\rho_A:=\mathsf{w}_{\mathfrak{s}}/\eta_p^{1/2} .
\end{equation}
Then the following hold.
\begin{enumerate}
\item At (L)-far, (e2) yields (e4) of \eqref{4.25:eq-increment-bound-a} on the disc
\eqref{4.25:eq-increment-bound-c}. The left side of (e4) is the fifth entry of
$\mathsf{N}^{\chi}_y(\mathcal{K}(z))$.
\item On the disc \eqref{4.25:eq-increment-bound-c}, at every $y\in\mathcal{Y}(e)$ and
$|\sigma|\le e^{\kappa_1}$, the \emph{layer-count bound}
\begin{equation}\label{4.25:eq-increment-bound-b}
\begin{gathered}
\mathsf{N}_y(\mathcal{K}_i(z))\le 2\hat B_H\,C'_{age}\,\bar\alpha(x_{\ell'(y)})
\quad\text{for every layer } \mathcal{K}_i,\\
C'_{age}:=2e^{\delta_{\mathfrak{s}}}\sup_{a\ge0}\,(1+B_\sharp a)^{1/2}
e^{-\frac14\delta_{\mathfrak{s}}(a-1)^+}.
\end{gathered}
\end{equation}
For the cut layer, $\mathsf{N}_y$ in the layer-count bound denotes the maximum of the first four entries of
$\mathsf{N}^{\chi}_y$, and its fifth entry is at most twice the right side of the layer-count bound.
\item The bound \eqref{4.25:eq-increment-hypothesis-a} of the increment hypothesis of
Definition~\ref{4.25:def-increment-hypothesis} holds on $\mathcal{Y}(e)$ at the presentation by the site
$\mathfrak{s}$, with
\begin{equation}\label{4.25:eq-increment-bound-4}
C_{inc}:=2(K_{\mathsf{j}}+2c_\zeta)\hat B_H C'_{age} .
\end{equation}
The constant $C_{inc}$ is a single number fixed before $\gamma$. Condition~(i) of
\eqref{4.25:eq-frozen-collar-weight-a} is imposed with this value of $C_{inc}$.
\end{enumerate}
\end{lemma}

\begin{proof}
Throughout, $y\in\mathcal{Y}(e)$ and $|\sigma|\le e^{\kappa_1}$. We abbreviate $\delta:=\delta_{\mathfrak{s}}$,
$\mathsf{U}:=\mathsf{U}_{\mathfrak{s}}$ and $\mathsf{w}:=\mathsf{w}_{\mathfrak{s}}$.

We begin with \eqref{4.25:eq-increment-bound-3} at (L)-far. Since $\chi=t+\mathsf{t}\hat\chi_{\square}$, the
four arguments of the entries decompose as follows:
\begin{align*}
\chi A'_s&=tA'_s+\mathsf{t}(\hat\chi_{\square}A'_s),\qquad
\chi\mathbb{H}_s=t\mathbb{H}_s+\mathsf{t}(\hat\chi_{\square}\mathbb{H}_s),\\
P_1(\sigma)(\chi A'_s)&=tP_1(\sigma)A'_s+\mathsf{t}P_1(\sigma)(\hat\chi_{\square}A'_s),\\
\chi\mathfrak{m}_H(\sigma)X_s&=t\,\mathfrak{m}_H(\sigma)X_s+\mathsf{t}\,\hat\chi_{\square}\mathfrak{m}_H(\sigma)X_s .
\end{align*}
The third identity holds because $P_1(\sigma)$ is linear. The fourth holds because $\chi$ multiplies
from outside $\mathfrak{m}_H(\sigma)$, as in the definition of $\mathcal{J}^{\sharp}[\chi]$ in
Lemma~\ref{4.25:lem-cut-layer-partner}, whose proof is incomplete at one step. Each entry is a seminorm of
its argument, and $0\le t\le1$. Hence the $j$-th entry of $\mathsf{N}^{\chi}_y(\mathcal{K}(z))$ is at most
$t$ times the $j$-th entry of $\mathsf{N}_y(\mathbb{H}_s)$ plus $|\mathsf{t}|$ times the $j$-th entry of
\eqref{4.25:eq-increment-bound-2}. This is at most the right side of \eqref{4.25:eq-increment-bound-3}.

We next derive (e4) at (L)-far, on the disc \eqref{4.25:eq-increment-bound-c}.

First, the definition of the layer datum, $\mathbb{H}_s:=A'_s-H(\sigma)X_s$, gives
$H(\sigma)X_s=A'_s-\mathbb{H}_s$. Since $N_y$ is a seminorm and $N_y\le N^{\Delta}_y$, we obtain
\begin{equation*}
\begin{split}
N_y(H(\sigma)X_s)&\le N_y(A'_s)+N_y(\mathbb{H}_s)\le N^{\Delta}_y(A'_s)+N_y(\mathbb{H}_s)\\
&\le 2\mathsf{N}_y(\mathbb{H}_s)\le 2\mathsf{U}e^{-\delta(\mathfrak{d}-1)} .
\end{split}
\end{equation*}
The inequality $N_y\le N^{\Delta}_y$, the domination of the block size by the norm $N^{\Delta}_y$, follows
from the definition of the norm $N^{\Delta}_y$ and part~(i) of the corollary to the layer theorem; it is
used in the same way in
Lemma~\ref{4.25:lem-current-partner}. All entries are taken in the convention of $\nabla_W$. The
third inequality holds because $N^{\Delta}_y(A'_s)$ and $N_y(\mathbb{H}_s)$ are the first two entries of
$\mathsf{N}_y(\mathbb{H}_s)$. The last is (e2) for the undilated layer $\mathbb{H}_s$.

Second, write $y_\square$ for the base point and $\ell_\square$ for the scale of the response cube
$\square$, the quantities in terms of which its cut-off $\hat\chi_{\square}$ is defined in
Lemma~\ref{4.25:lem-cut-layer-partner}. Suppose that $\tilde\Delta(y)$ meets $\square$ at a point $q$.
Since $y\in\mathcal{Y}(e)$, the
block $\Delta(y)$ contains a site $p$ of $S_v=S^+(e)$. As $p\in\Delta(y)$ and $q\in\tilde\Delta(y)$,
\begin{equation*}
\operatorname{dist}_\infty(p,\square)\le|p-y|_\infty+|y-q|_\infty\le2\ell_y\le2L\ell_\square<M\ell_\square .
\end{equation*}
Here $\ell_y\le L\ell_\square$ by the argument based on \eqref{4.25:eq-collar-weight-b} in
Lemma~\ref{4.25:lem-cut-layer-partner}, whose proof is incomplete at one step. The last inequality
holds because $M\ge2L^2>2L$, by the choice of the auxiliary parameter $M$.

Now $\tilde\square$ is $\square$ together with one layer of $\pi$-cubes touching it \cite{B-CMP119}; here
the $\pi$-cubes are cubes of side $M\ell_\square$. As a
point set $\tilde\square$ equals $\{\operatorname{dist}_\infty(\cdot,\square)\le M\ell_\square\}$. We verify
the inclusion $\supseteq$ for points directly. If
$|p-y_\square|_\infty\le2M\ell_\square$, then $p$ lies in a $\pi$-cube
\begin{equation*}
\prod_{\nu}\bigl[y_{\square,\nu}+n_\nu M\ell_\square,\ y_{\square,\nu}+(n_\nu+1)M\ell_\square\bigr],
\qquad n_\nu\in\{-2,-1,0,1\},
\end{equation*}
and every such cube touches $\square$, corners included. The point $p$ found above satisfies
$|p-y_\square|_\infty\le2M\ell_\square$. Hence $p\in\tilde\square$, and so
$\tilde\square\cap S_v\neq\emptyset$.

Third, this is the one step taken from the convention for $\eta_p$ in the hypotheses. Since
$\tilde\square\cap S_v\neq\emptyset$, the cube-counted distance vanishes, so $\eta_p=1$. The disc
\eqref{4.25:eq-increment-bound-c} then reads $|\mathsf{t}|\le\rho_A=\mathsf{w}$.

Finally, if the indicator in (e4) vanishes there is nothing to prove. Otherwise we multiply
$|\mathsf{t}|\le\mathsf{w}$ by the bound of the first step, which gives (e4).

Any other distance convention for $\eta_p$ changes only a number. Wherever the fifth entry is non-zero,
the second step gives $\operatorname{dist}_\infty(\square,S_v)\le2L\ell_\square$ uniformly. Hence $\eta_p$
is bounded below by a number $\eta_*$ satisfying
\begin{equation*}
\eta_p\ge\eta_*\ge e^{-\delta_1M(2L+1)} .
\end{equation*}
It then suffices to replace $c_\zeta$ by $c_\zeta\eta_*^{-1/2}$ inside $C_{inc}$; this is a number
depending on $(d,L,\kappa_1,M)$ and fixed before $\gamma$. The first four entries are insensitive to the
convention, since they reach the edge of the disc through the geometric mean of (e1) and (e2), as
below.

We now compare depth with age. Put $a:=k'-\ell'(y)\ge0$. By the fourth hypothesis, a contour from $y$ to
$\mathfrak{F}_{\mathfrak{s}}$ ends in cells of level $k'$. It crosses the levels
$\ell'(y)+1,\dots,k'-1$, each over at least one $M$-cube of its own scale. Hence
\begin{equation}\label{4.25:eq-increment-bound-5}
\mathfrak{d}\ge(a-1)^+ .
\end{equation}

Next we bound the dilated part. Using $\min\{u,v\}\le(uv)^{1/2}$, the disc
\eqref{4.25:eq-increment-bound-c}, (e1) and (e2) give
\begin{equation*}
|\mathsf{t}|\,\mathsf{N}_y(\delta\mathbb{H}_p)\le\frac{\mathsf{w}}{\eta_p^{1/2}}
\bigl(\eta_p\mathsf{U}\cdot\mathsf{U}e^{-\delta(\mathfrak{d}-1)}\bigr)^{1/2}
=\mathsf{w}\,\mathsf{U}e^{-\frac12\delta(\mathfrak{d}-1)} .
\end{equation*}
By (e4), the fifth entry is at most $2\mathsf{w}\mathsf{U}e^{-\delta(\mathfrak{d}-1)}$.

The $\mathsf{t}$-holding layer is the sum of its undilated and dilated parts. It is linear in the
parameter \cite{B-CMP109}, and at (L)-far it has the form \eqref{4.25:eq-increment-bound-1}. Since each
entry is a seminorm, its size is bounded jointly, by the sum of two contributions:
\begin{itemize}
\item its undilated part contributes at most $\mathsf{U}e^{-\delta(\mathfrak{d}-1)}$, by (e2);
\item its dilated part contributes at most $\mathsf{w}\mathsf{U}e^{-\frac12\delta(\mathfrak{d}-1)}$.
\end{itemize}
At (L)-far this is the content of \eqref{4.25:eq-increment-bound-3}.
The factor $(1+\mathsf{w})$ of (e3) controls the sum. Every other layer $\mathcal{K}_i$ is undilated,
and (e2) alone gives $\mathsf{N}_y(\mathcal{K}_i)\le\mathsf{U}e^{-\delta(\mathfrak{d}-1)}$. This is the first
summand, so it obeys each of the bounds below.

Case $\mathfrak{d}\ge1$. Here $(\mathfrak{d}-1)^+=\mathfrak{d}-1$. By (e3),
\begin{equation*}
\begin{split}
\mathsf{U}e^{-\delta(\mathfrak{d}-1)}+\mathsf{w}\mathsf{U}e^{-\frac12\delta(\mathfrak{d}-1)}
&\le(1+\mathsf{w})\mathsf{U}e^{-\frac12\delta(\mathfrak{d}-1)}
=\bigl[(1+\mathsf{w})\mathsf{U}e^{-\frac14\delta(\mathfrak{d}-1)}\bigr]e^{-\frac14\delta(\mathfrak{d}-1)}\\
&\le2\hat B_H\bar\alpha(x_{k'})\,e^{-\frac14\delta(\mathfrak{d}-1)} .
\end{split}
\end{equation*}
The fifth entry is at most $2\mathsf{w}\mathsf{U}e^{-\frac12\delta(\mathfrak{d}-1)}$, which is twice the
dilated bound. By \eqref{4.25:eq-increment-bound-5}, $\mathfrak{d}-1\ge(a-1)^+-1$. Hence
\begin{equation*}
e^{-\frac14\delta(\mathfrak{d}-1)}\le e^{\frac14\delta}e^{-\frac14\delta(a-1)^+} .
\end{equation*}

Case $0\le\mathfrak{d}<1$. Here $(\mathfrak{d}-1)^+=0$. By \eqref{4.25:eq-increment-bound-5},
$(a-1)^+\le\mathfrak{d}<1$. Since $a$ is an integer, this forces $a\le1$, and then
$e^{-\frac14\delta(a-1)^+}=1$. In this case (e3) reads $(1+\mathsf{w})\mathsf{U}\le2\hat B_H\bar\alpha(x_{k'})$.
Since $1-\mathfrak{d}\le1$,
\begin{equation*}
\mathsf{U}e^{-\delta(\mathfrak{d}-1)}+\mathsf{w}\mathsf{U}e^{-\frac12\delta(\mathfrak{d}-1)}
\le\mathsf{U}e^{\delta}+\mathsf{w}\mathsf{U}e^{\frac12\delta}\le(1+\mathsf{w})\mathsf{U}e^{\delta}
\le2\hat B_H\bar\alpha(x_{k'})\,e^{\delta} .
\end{equation*}
The fifth entry is at most $2\mathsf{w}\mathsf{U}e^{\delta}$.

In both cases every layer satisfies
\begin{equation}\label{4.25:eq-increment-bound-6}
\mathsf{N}_y(\text{layer})\le2\hat B_H\bar\alpha(x_{k'})\,e^{\delta}e^{-\frac14\delta(a-1)^+} .
\end{equation}
In the first case this uses $e^{\frac14\delta}\le e^{\delta}$. The fifth entry is at most twice the right
side of \eqref{4.25:eq-increment-bound-6}.

We now pass from level $k'$ to level $\ell'(y)$. Apply \eqref{4.25:eq-flow-inequalities-c} with
$m=\ell'(y)\le k'$, so that $k'-m=a$, and use $x_{k'}\ge1$. This gives
\begin{equation*}
\bar\alpha(x_{k'})\le2(1+B_\sharp a)^{1/2}\,\bar\alpha(x_{\ell'(y)}) .
\end{equation*}
Hence
\begin{equation*}
\mathsf{N}_y(\mathcal{K}_i(z))\le2\hat B_H\bigl[2e^{\delta}(1+B_\sharp a)^{1/2}
e^{-\frac14\delta(a-1)^+}\bigr]\bar\alpha(x_{\ell'(y)})\le2\hat B_HC'_{age}\,\bar\alpha(x_{\ell'(y)}) .
\end{equation*}
The supremum in $C'_{age}$ is finite because $\delta>0$. This proves the layer-count bound in
\eqref{4.25:eq-increment-bound-b}. The fifth entry of the cut layer is at most twice its right side. At
$\mathfrak{d}=0$, $a=1$, for an undilated layer, the factor $e^{\delta_{\mathfrak{s}}}$ of $C'_{age}$ is used
in full.

Before bounding the increment, we verify the smallness conditions needed below. The derivation of the
layer-count bound above invokes neither Lemma~\ref{4.25:lem-current-partner} nor
Lemma~\ref{4.25:lem-cut-layer-partner}, so using the layer-count bound to verify their hypotheses is not
circular. Put $b:=\hat B_HC'_{age}$ and abbreviate
$\bar\alpha:=\bar\alpha(x_{\ell'(y)})$. We use:
\begin{itemize}
\item the layer-count bound \eqref{4.25:eq-increment-bound-b};
\item the relation $\bar\alpha(x_{\ell'(y)})\le\theta(y)/5$ between the envelope and the block threshold
$\theta(y)$ of \eqref{4.25:eq-flow-inequalities-f};
\item $C_{inc}\ge6b$, which holds because $K_{\mathsf{j}}\ge3$ and $c_\zeta\ge0$;
\item $\theta(y)\le\Theta(e)$;
\item $\Theta(e)\le\bar\theta^{\downarrow}$, by (1) and (2) of \eqref{4.25:eq-conversion-old-factors-a};
\item condition~(i) of \eqref{4.25:eq-frozen-collar-weight-a}, which gives
$\tfrac98C_{inc}\bar\theta^{\downarrow}\le\varpi_G$;
\item $\varpi_G\le\varepsilon_{reg}/20\le1/600$, from the definition of $\varpi_G$ in
\eqref{4.25:eq-frozen-collar-weight-a}.
\end{itemize}
These give
\begin{equation*}
\begin{split}
N_y(\mathcal{K}_i)&\le2b\bar\alpha\le\tfrac25\,b\,\theta(y)\le\tfrac1{15}C_{inc}\theta(y)
\le\tfrac1{15}C_{inc}\Theta(e)\le\tfrac{8}{135}\varpi_G\\
&\le\tfrac1{600}\cdot\tfrac{8}{135}<10^{-4}<\tfrac1{40} .
\end{split}
\end{equation*}
Here $N_y(\mathcal{K}_i)$ is the second entry of $\mathsf{N}_y$, resp.\ of $\mathsf{N}^{\chi}_y$. This is the
hypothesis $N_y\le1/40$ of Lemma~\ref{4.25:lem-current-partner}, and of
Lemma~\ref{4.25:lem-cut-layer-partner} (whose proof is incomplete at one step), for every layer. Where
the local constant $c_\chi$ of the cut enters the second entry of the cut layer as a factor, the chain
gives
\begin{equation*}
N_y\le c_\chi\cdot\tfrac{8}{135}\varpi_G\le c_\chi\cdot\tfrac{8}{81000},
\end{equation*}
which is less than $1/40$ for $c_\chi\le253$. For larger $c_\chi$ the condition
$\varpi_G\le135/(320c_\chi)$ is added to the conditions whose minimum defines $\varpi_G$; it gives
$c_\chi\cdot\tfrac{8}{135}\varpi_G\le\tfrac1{40}$.

The same bound supplies the room $0.04$ of the composition map $\Psi$ in \eqref{4.25:eq-gaussian-pieces-b}
of Lemma~\ref{4.25:lem-gaussian-pieces}, whose proof is incomplete at one step. This room is needed for
the at most four compositions $W_{i+1}=e^{i\eta\mathcal{K}_i}W_i$, where $W_1$ is the kernel of the base
pair; write $A'_{W_i}$ for the potential of the kernel $W_i$. By induction on $i$,
\begin{equation*}
N_y(A'_{W_i})\le(5/4)^{i-1}\bigl[\Theta(e)+(i-1)\cdot10^{-4}\bigr]<6.1\cdot10^{-4} .
\end{equation*}
The sizes of the two arguments of $\Psi$ therefore sum to less than $10^{-3}\le0.04$.

For the gauge-field slot, set $e^{i\eta A''}:=\prod_ie^{i\eta\mathcal{K}_i}$. Apply the composition estimate
of \eqref{4.25:eq-gaussian-pieces-b} three times, in its regime, as just verified; Lemma~\ref{4.25:lem-gaussian-pieces}
has a proof incomplete at one step. Using the layer-count bound for the second entries, this gives
\begin{equation*}
N_y(A'')\le1.1^3(9/8)^2\cdot4\max_iN_y(\mathcal{K}_i)\le6.8\cdot2b\bar\alpha .
\end{equation*}
The factor $9/8$ comparing $\nabla_U$ with $\nabla_W$ is accounted for at this point only.

For the current slot, $J''=\sum_i(\text{partner of }\mathcal{K}_i)$. A layer that is a layer datum is
treated by Lemma~\ref{4.25:lem-current-partner} at its own kernel, applied against the zero datum with
$\mathsf{p}\le1$; together with the layer-count bound this gives
\begin{equation*}
\ell_y^3|\mathfrak{s}^{\sharp}[\mathcal{K}_i]|\le K_{\mathsf{j}}\mathsf{N}_y(\mathcal{K}_i)\le2K_{\mathsf{j}}\,b\,\bar\alpha .
\end{equation*}
The cut layer at (L)-far is treated by
Lemma~\ref{4.25:lem-cut-layer-partner}, whose proof is incomplete at one step. By that lemma, the
layer-count bound and the bound on the fifth entry,
\begin{equation*}
\ell_y^3|\mathcal{J}^{\sharp}[\chi]|\le K_{\mathsf{j}}\cdot(\text{first four entries})
+c_\zeta\cdot(\text{fifth entry})\le(K_{\mathsf{j}}+2c_\zeta)\cdot2b\bar\alpha .
\end{equation*}
Since there are at most four layers, $\ell_y^3|J''|\le4(K_{\mathsf{j}}+2c_\zeta)\cdot2b\bar\alpha$.

Finally, $\bar\alpha\le\theta(y)/5$ gives
\begin{equation*}
N_y(A'')\le2.72\,b\,\theta(y),\qquad
\ell_y^3|J''|\le1.6\,(K_{\mathsf{j}}+2c_\zeta)\,b\,\theta(y) .
\end{equation*}
The constant $C_{inc}=2(K_{\mathsf{j}}+2c_\zeta)b$ of \eqref{4.25:eq-increment-bound-4} dominates both
factors. Indeed, $K_{\mathsf{j}}\ge3$ gives $2.72\le6\le2(K_{\mathsf{j}}+2c_\zeta)$, and $1.6\le2$. This is the
bound \eqref{4.25:eq-increment-hypothesis-a} of Definition~\ref{4.25:def-increment-hypothesis} on
$\mathcal{Y}(e)$, at the presentation by the site $\mathfrak{s}$.
\end{proof}

\begin{remark}[the tube constant read as at least one]\label{4.25:rem-tube-constant}
The tube constant $B_H$ enters this section through the supply bound of the tube theorem
\begin{equation*}
(1+\mathsf{w}_T)\,\mathsf{U}_T\le B_H\,\alpha_*,
\end{equation*}
and nothing fixes a lower bound for it. Three steps of the section nevertheless use one.
\begin{itemize}
\item The base bound $\theta(y)\le\Theta(e)$ of \eqref{4.25:eq-gaussian-pieces-a}, combined
with the first inequality of member (i) of \eqref{4.25:eq-frozen-collar-weight-a}, gives
$\Theta(e)\le\tfrac{8}{9}\varpi_G/C_{inc}$; this must be at most $\varepsilon_{reg}/5$ where
the second Gaussian-piece bound is applied through clause (iv) of the lemma on Gaussian pieces
(domain, invariance and continuation).
Given the ceiling $\varpi_G\le\varepsilon_{reg}/20$, this holds as soon as $C_{inc}\ge 2/9$.
\item The same bound must be at most $\varepsilon_{reg}/20$ where the primed form of clause
(iv) of that lemma is applied at the base level of the three-field construction. Given
the ceiling $\varpi_G\le\varepsilon_{reg}/20$, this holds as soon as $C_{inc}\ge 8/9$.
\item The base class of the links in the rotated increment bound requires
$B_H C'_{age}\ge 0.2583$.
\end{itemize}
Accordingly, throughout this section every occurrence of $B_H$ is read as the number
$\hat B_H$ of the first line of the following display, in which $b$ is also defined. Under
this reading the bounds in its second line hold; and, with the base bound
\eqref{4.25:eq-gaussian-pieces-a}, the first inequality of member (i) of
\eqref{4.25:eq-frozen-collar-weight-a} and the ceiling $\varpi_G\le\varepsilon_{reg}/20$, so
does the chain in its third line, for every piece $e$ of $\mathfrak{Q}(\mathfrak{a})$:
\begin{equation}\label{4.25:eq-tube-constant-a}
\begin{gathered}
\hat B_H:=\max\{1,B_H\},\qquad b:=\hat B_H C'_{age},\\
C_{inc}\ge 2K_{\mathsf{j}}C'_{age}\ge 300,\qquad b\ge 2,\\
\Theta(e)\le\frac{8}{9}\,\frac{\varpi_G}{C_{inc}}\le\frac{\varpi_G}{337.5}
\le\frac{\varepsilon_{reg}}{6750}.
\end{gathered}
\end{equation}
Under this reading every hypothesis in which the tube constant occurs is weakened or
unchanged, every conclusion in which it stands is read with the larger majorant, and the
ceilings on $\gamma$ keep their form. The reading introduces neither a cap
(Definition~\ref{2.3:def-cap}) nor a lower bound on any auxiliary parameter. If $B_H\ge 1$,
then $\hat B_H=B_H$ and the reading changes nothing.
In particular the three requirements listed above are met.
\end{remark}

\begin{proof}
\emph{The reading is admissible.} In this section $B_H$ stands only on majorant sides, in
the following places:
\begin{itemize}
\item the right member of hypothesis (e3)$^+$ of the increment bound at tube and far sites
(Lemma~\ref{4.25:lem-increment-bound});
\item the right member $2B_H\bar\alpha$ of the site-unit inequality of member (i)
of \eqref{4.25:eq-frozen-collar-weight-a};
\item the tube supply bound displayed above;
\item the conclusions \eqref{4.25:eq-increment-bound-b} of
Lemma~\ref{4.25:lem-increment-bound} and of the rotated increment bound;
\item the increment constant $C_{inc}$ of \eqref{4.25:eq-increment-hypothesis-a}, and hence
the constants $C_G$, $\mathfrak{c}_G$, $b_\delta^G$, $\mathsf{n}_G^J$, $\mathsf{n}_G^{blk}$ and
$\varepsilon_G$, which are built from $C_{inc}$.
\end{itemize}
Since $\hat B_H\ge B_H$, replacing $B_H$ by $\hat B_H$ enlarges the right members of
(e3)$^+$ and of that inequality of member (i). Both hypotheses are therefore weakened.

The input that supplies (e3)$^+$ at the tube site is the tube supply bound. That bound
implies (e3)$^+$ with $B_H$, and it still implies (e3)$^+$ with $\hat B_H$, because
$B_H\alpha_*\le\hat B_H\alpha_*$.

The ceilings on $\gamma$ that members (i)--(iii) of \eqref{4.25:eq-frozen-collar-weight-a}
impose keep their
form. Their left members are numbers fixed before $\gamma$, re-maximised with $\hat B_H$ in
place of $B_H$, and they multiply explicit functions tending monotonically to zero. Each
ceiling therefore remains a condition that holds for $\gamma$ small enough. No condition on
the auxiliary parameters changes its kind, so no cap and no lower bound arise.

Finally, if $B_H\ge 1$ then $\max\{1,B_H\}=B_H$, and every statement is literally unchanged.

\emph{The bound on $C_{inc}$.} By \eqref{4.25:eq-increment-hypothesis-a},
$C_{inc}=2(K_{\mathsf{j}}+2c_\zeta)\hat B_H C'_{age}$. Since $c_\zeta\ge 0$ and
$\hat B_H\ge 1$, this gives $C_{inc}\ge 2K_{\mathsf{j}}C'_{age}$.

By \eqref{4.25:eq-current-partner-a}, $K_{\mathsf{j}}=3+18d+9K_J$, with $K_J\ge 0$ the
constant of the current partner in \eqref{4.25:eq-current-partner-a}. At $d=4$ this gives
$K_{\mathsf{j}}\ge 3+72=75$.

By \eqref{4.25:eq-increment-bound-b},
\begin{equation*}
C'_{age}=2e^{\delta_{\mathfrak{s}}}\sup_{a\ge 0}\,(1+B_\sharp a)^{1/2}\,
e^{-\frac14\delta_{\mathfrak{s}}(a-1)^+},
\end{equation*}
where $B_\sharp$ is the flow constant of the coupling flow
(Definition~\ref{1.2:def-flow-constants}).
Evaluate the supremum at $a=1$, where $(a-1)^+=0$. This gives
\begin{equation*}
C'_{age}\ge 2e^{\delta_{\mathfrak{s}}}(1+B_\sharp)^{1/2}\ge 2,
\end{equation*}
since $\delta_{\mathfrak{s}}\ge 0$ and $B_\sharp\ge 0$. The inequality is strict as soon as
$\delta_{\mathfrak{s}}>0$. Hence
\begin{equation*}
C_{inc}\ge 2K_{\mathsf{j}}C'_{age}\ge 2\cdot 75\cdot 2=300.
\end{equation*}

\emph{The bound on $b$.} From $\hat B_H\ge 1$ and $C'_{age}\ge 2$ we get
$b=\hat B_H C'_{age}\ge 2$.

\emph{The bound on $\Theta(e)$.} The base bound \eqref{4.25:eq-gaussian-pieces-a}, taken
together with the first inequality of member (i) of \eqref{4.25:eq-frozen-collar-weight-a},
gives
$\Theta(e)\le\tfrac{8}{9}\varpi_G/C_{inc}$. Inserting $C_{inc}\ge 300$ gives
\begin{equation*}
\frac{8}{9}\,\frac{\varpi_G}{C_{inc}}\le\frac{8\,\varpi_G}{2700}=\frac{\varpi_G}{337.5}.
\end{equation*}
The last inequality of \eqref{4.25:eq-tube-constant-a} is equivalent to the ceiling
$\varpi_G\le\varepsilon_{reg}/20$.

\emph{The three requirements.} We have
\begin{equation*}
C_{inc}\ge 300>\tfrac{8}{9}>\tfrac{2}{9}\qquad\text{and}\qquad b\ge 2>0.2583.
\end{equation*}
So the three requirements listed in the remark are met.

\emph{The max-form of member (i).} The first inequality of member (i) of
\eqref{4.25:eq-frozen-collar-weight-a} is written in the max-form
\begin{equation*}
\max\Bigl\{\tfrac{9}{8}C_{inc},\,2\Bigr\}\,\bar\theta^{\downarrow}(x-2c_2)\le\varpi_G .
\end{equation*}
Under the reading, $\tfrac{9}{8}C_{inc}\ge\tfrac{9}{8}\cdot 300>2$. The condition therefore
coincides with $\tfrac{9}{8}C_{inc}\,\bar\theta^{\downarrow}(x-2c_2)\le\varpi_G$, which is the
form used above.

Without the reading, the max-form alone still gives $\Theta(e)\le\tfrac12\varpi_G$ for every
$B_H>0$. Given the ceiling $\varpi_G\le\varepsilon_{reg}/20$, this yields
$\Theta(e)\le\varepsilon_{reg}/40$. That meets both requirements $\varepsilon_{reg}/5$ and
$\varepsilon_{reg}/20$ of the base bound. For the base class of the links in the rotated
increment bound, with $b=B_HC'_{age}$, the line of its display then reads
\begin{equation*}
1.331\,\bigl(\Theta(e)+1.2\,b\,\theta(y)\bigr)\le 1.331\,\bigl(\tfrac12+\tfrac{8}{45}\bigr)\,
\varpi_G=0.9021\,\varpi_G .
\end{equation*}
\end{proof}

\begin{remark}[Supply of the depth hypothesis]\label{4.25:rem-depth-hypothesis}
Let $e$ be a piece and let $\mathfrak{s}$ be a site of Lemma~\ref{4.25:lem-increment-bound}, whose
notation is used throughout. Either $\mathfrak{s}$ is the tube site (T) of step $k'$, with
$\mathfrak{F}_{\mathfrak{s}}=\Omega$, $\delta_{\mathfrak{s}}=\delta_1^{sc}$ and
$\mathsf{U}_{\mathfrak{s}}=\mathsf{U}_T$; or $\mathfrak{s}$ is a far (L)-site, that is, the treated
class $Z$ of the (L)-site satisfies $S^+(e)\cap Z=\emptyset$, and then
$\mathfrak{F}_{\mathfrak{s}}=\Lambda^{\sim}$, $\delta_{\mathfrak{s}}=\delta_\natural$ and
$\mathsf{U}_{\mathfrak{s}}=\mathsf{U}_L$, with $S_v=S^+(e)$.

For $y\in\mathcal{Y}(e)$ (Definition~\ref{4.25:def-enlarged-support-blocks}) put
$\mathfrak{d}:=d_{sc}(y,\mathfrak{F}_{\mathfrak{s}})\ge0$, and write $\hat B_H:=\max\{1,B_H\}$.
Assume the following.
\begin{enumerate}
\item[(A)] \emph{The tube bound.} At (T), the tube bound of step $k'$,
\begin{equation}\label{4.25:eq-depth-hypothesis-1}
(1+\eta_p|\mathsf{t}|)\,K_R\,g_{k'}\sup|\mathsf{B}|\le\alpha_{*,k'},
\end{equation}
holds on its own disc in $\mathsf{t}$. Here $K_R$ and $\mathsf{B}$ are the constant and the field that
enter that bound. The number $\mathsf{w}_T$ is the radius of this disc, and on it the bound gives
$(1+\mathsf{w}_T)\mathsf{U}_T\le B_H\alpha_{*,k'}$.
\item[(B)] \emph{Far-bond separation.} Let $\mathcal{Z}_Z$ denote $Z$ with the outer five layers
of cubes of side $M\check R_{k'}$ of its last shell removed. Every bond $b$ of the level-$k'$ lattice
with $b\notin Z$ satisfies $d_{sc}(b,\mathcal{Z}_Z)\ge5\check R_{k'}$. Moreover
$\Lambda\subset\mathcal{Z}_Z$ and $\operatorname{dist}(\partial\mathcal{Z}_Z,\Lambda)\ge5M\check R_{k'}$.
The separation integer of these statements is the integer $\check R_{k'}$ of the run.
\item[(C)] \emph{The far-term support statement.} This is the positional fact from which, together
with (B), the bound on the far terms is proved: the arguments of the far terms have their supports
in $\Lambda^{\sim}$.
Here $\Lambda^{\sim}$ enlarges $\Lambda$ by at most $4\check R_{k'}+1$ layers of level-$k'$ cubes of
side $M$. This count includes the cost of passing from a site to a block. That cost is $0$ for a block
off $Z$, and at most $1$ under the convention that a site is represented by its top block
$B^{top}(\cdot)$.
\item[(D)] The coupling-flow inequality \eqref{4.25:eq-flow-inequalities-d},
$e^{-\check R_k}\le E(x_k)$.
\item[(E)] \emph{The $E$-condition on the site units.} For every value $x$ of the running inverse
couplings of the run, in particular for $x=x_{k'}$,
\begin{equation}\label{4.25:eq-depth-hypothesis-2}
(1+\mathsf{w}_L)\,\mathsf{U}_L(x)\,e^{\delta_\natural}E(x)^{\delta_\natural/4}\le2\hat B_H\bar\alpha(x).
\end{equation}
\item[(F)] \emph{The convention on the block sizes $\mathsf{N}_y$.}
\begin{itemize}
\item At (T), the Lipschitz bound of the layer map is used. For every weight $p$ of the class
$\mathfrak{M}_\delta$ of weights admitted in that bound, $\delta>0$ being the decay rate indexing
this class, and any two data of the layer map, numbered $1$ and $2$, it
bounds each of the five weighted norms
\begin{equation*}
\begin{gathered}
\|A'_1-A'_2\|^{\Delta}_p,\quad \|\mathbb{H}_1-\mathbb{H}_2\|^{\Delta}_p,\quad \|X_1-X_2\|_p,\\
\|P_1(\sigma)(A'_1-A'_2)\|_p,\quad \|\mathfrak{m}_H(\sigma)(X_1-X_2)\|_p
\end{gathered}
\end{equation*}
by the Lipschitz constant $B_{\mathbb{H}}$ of the layer map times the $p$-weighted norm $\|\cdot\|_p$
of the difference of the two data. Here
$A'_i$, $\mathbb{H}_i$ and $X_i$ are the potential, the layer datum and the auxiliary field obtained
from the $i$-th datum, and
$\|\cdot\|^{\Delta}_p$, $\|\cdot\|_p$ are the weighted block norms of that bound. These norms contain
the four entries of $\mathsf{N}_y$;
the norm $\|X_1-X_2\|_p$ is idle. The bound is applied at the weight
$p=e^{-\delta d^{\wedge}(\cdot,\Omega)}$, built on the distance $d^{\wedge}(\cdot,\Omega)$ to $\Omega$
used at the tube site, and its companion bound at the weight
$p=e^{-\delta\operatorname{dist}(\cdot,\square)}$. In addition: $\eta_p$ has rate at most $\delta$;
these weights belong to $\mathfrak{M}_\delta$; and the relation of $B_H$ to $B_{\mathbb{H}}$ is part of this
convention, only $\hat B_H$ entering the bounds below.
\item At (L), the cut layer $(t+\mathsf{t}\hat\chi_\square)\mathbb{H}_s$, with partner
$\mathcal{J}^{\sharp}[\chi]$, has the five-entry size $\mathsf{N}^{\chi}_y$. Its entries are the four
uncut entries together with the $P_1$-tail, and their bounds are invoked only after
$(\mathrm{e}3)^+$ has been established.
Entries $1$, $2$ and $4$ of the cut size $\mathsf{N}_y(\delta\mathbb{H}_p)$, where
$\delta\mathbb{H}_p=\hat\chi_\square\mathbb{H}_s$, are at most the corresponding entries of
$\mathsf{N}_y(\mathbb{H}_s)$ multiplied by
$c_\chi\,\mathbf{1}[\tilde\Delta(y)\cap\square\neq\emptyset]$, where
$c_\chi:=1+2c_d$ and $\mathbf{1}[\cdot]$ is the indicator, the number $c_\chi$ being placed inside
$\mathsf{U}_L$.
The tail is taken in the pointwise unit-free product form: at every current block $y$,
\begin{equation}\label{4.25:eq-depth-hypothesis-3}
\ell_y^3\,\bigl|P_1(\sigma)(\hat\chi_\square A'_s)\bigr|\le\eta_p\,\mathsf{f}(y)\,\mathsf{U}_L,
\qquad \mathsf{f}(y):=e^{-\delta_\natural(d_{sc}(y,\Lambda^{\sim})-1)}.
\end{equation}
\item At (L): $\eta_p=1$ whenever the enlarged cube $\tilde\square$, namely the cube $\square$ of the
cut-off $\hat\chi_\square$ together with one layer of partition cubes touching it, meets $S_v$.
\item At both sites: $\mathsf{w}_T$ is the radius in (A), and $\mathsf{p}\le1$ at the kernel of each
layer.
\end{itemize}
\end{enumerate}
Then the following hold.
\begin{enumerate}
\item[(i)] At both sites, hypothesis $(\mathrm{e}3)^+$ of Lemma~\ref{4.25:lem-increment-bound}
(displayed in \eqref{4.25:eq-increment-bound-a}) holds at every $y\in\mathcal{Y}(e)$.
\item[(ii)] At a far (L)-site, every block of $\mathcal{Y}(e)$ is off $Z$, and
$\mathfrak{d}-1\ge\check R_{k'}-2$ for every $y\in\mathcal{Y}(e)$.
\item[(iii)] No bound on the block sizes beyond $(\mathrm{e}1)$, $(\mathrm{e}2)$, $(\mathrm{e}3)^+$
and the convention (F) enters the proof of Lemma~\ref{4.25:lem-increment-bound}; in particular the
bound $(\mathrm{e}5)$ is a consequence of $(\mathrm{e}2)$.
\end{enumerate}
If at the blocks of a far (L)-site the disc $|\mathsf{t}|\le\mathsf{w}_L/\eta_p^{1/2}$ were replaced by
the larger disc
$|\mathsf{t}|\le\mathsf{w}_L/(\eta_p\mathsf{f}_v)^{1/2}$, with $\mathsf{f}_v$ the value of the profile
$\mathsf{f}$ entering that disc, then $(\mathrm{e}1)$ and $(\mathrm{e}2)$ must be read together as the
product bound $\mathsf{N}_y(\delta\mathbb{H}_p)\le\eta_p\,\mathsf{f}(y)\,\mathsf{U}_L$, and the proof of
Lemma~\ref{4.25:lem-increment-bound} runs unchanged.
\end{remark}

\begin{proof}
\emph{The tube site.} By (A), on its own disc, with $\mathsf{w}_T$ read as its radius, the bound
\eqref{4.25:eq-depth-hypothesis-1} gives $(1+\mathsf{w}_T)\mathsf{U}_T\le B_H\alpha_{*,k'}$.
Three facts finish the estimate: $B_H\le\hat B_H$; $\alpha_{*,k'}\le\bar\alpha(x_{k'})$, one of the
coupling-flow inequalities of Notation~\ref{4.25:not-flow-inequalities}; and $\hat B_H\bar\alpha\ge0$.
They give
\begin{equation*}
(1+\mathsf{w}_T)\mathsf{U}_T\le B_H\alpha_{*,k'}\le2\hat B_H\bar\alpha(x_{k'}).
\end{equation*}
Since $(\mathfrak{d}-1)^+\ge0$, we have $e^{-\frac14\delta_1^{sc}(\mathfrak{d}-1)^+}\le1$. Multiplying
the display by this factor yields $(\mathrm{e}3)^+$ at (T).

\emph{A far (L)-site: position of the blocks.} Let $y\in\mathcal{Y}(e)$. By
Definition~\ref{4.25:def-enlarged-support-blocks}, $\Delta(y)$ contains a site of $S^+(e)$. The region
$Z$ is a union of current cells \cite{B-CMP122a}. If $\Delta(y)$ lay in $Z$, it would contain a site
of $S^+(e)$ inside $Z$, against $S^+(e)\cap Z=\emptyset$. Hence, $Z$ being a union of current cells,
$\Delta(y)$ is off $Z$: every block of $\mathcal{Y}(e)$ is off $Z$, and every bond of it lies
outside $Z$. This proves the first half of (ii).

\emph{A far (L)-site: the depth.} The lower bound $\check R_{k'}-2$ in
$\mathfrak{d}-1\ge\check R_{k'}-2$ is the one prescribed at law sites by the site geometry of the
step; at the blocks $y\in\mathcal{Y}(e)$ it is
established directly from (B) and the support property (C), as
follows.

Let $b$ be a bond of $\Delta(y)$. Then $b\notin Z$, so (B) gives $d_{sc}(b,\mathcal{Z}_Z)\ge5\check R_{k'}$.
By (B), $\Lambda$ lies inside $\mathcal{Z}_Z$, at distance at least $5M\check R_{k'}$ from its boundary.
Hence every point of $\Lambda$ is at scaled distance at least $5\check R_{k'}$ from $b$. Passing from
$\Lambda$ to $\mathfrak{F}_{\mathfrak{s}}=\Lambda^{\sim}$ lowers this distance by at most the number
of layers added. By (C) that number is at most $4\check R_{k'}+1$, the cost of passing from a site to a
block included; for the block $\Delta(y)$, which is off $Z$, this cost is $0$. Therefore
\begin{equation*}
\mathfrak{d}\ge5\check R_{k'}-(4\check R_{k'}+1)=\check R_{k'}-1,
\end{equation*}
that is, $\mathfrak{d}-1\ge\check R_{k'}-2$, which is the second half of (ii). In particular
$(\mathfrak{d}-1)^+\ge\mathfrak{d}-1\ge\check R_{k'}-2$.

\emph{A far (L)-site: $(\mathrm{e}3)^+$.} By (D) at $k=k'$,
$e^{-\check R_{k'}}\le E(x_{k'})$. Combined with the lower bound on $(\mathfrak{d}-1)^+$, this gives
\begin{equation*}
\begin{aligned}
e^{-\frac14\delta_\natural(\mathfrak{d}-1)^+}
&\le e^{-\frac14\delta_\natural(\check R_{k'}-2)}
 =e^{\frac12\delta_\natural}\bigl(e^{-\check R_{k'}}\bigr)^{\delta_\natural/4}\\
&\le e^{\frac12\delta_\natural}E(x_{k'})^{\delta_\natural/4}
 \le e^{\delta_\natural}E(x_{k'})^{\delta_\natural/4}.
\end{aligned}
\end{equation*}
Multiply by $(1+\mathsf{w}_L)\mathsf{U}_L(x_{k'})$ and apply \eqref{4.25:eq-depth-hypothesis-2} at
$x=x_{k'}$. This gives
\begin{equation*}
(1+\mathsf{w}_L)\mathsf{U}_L(x_{k'})\,e^{-\frac14\delta_\natural(\mathfrak{d}-1)^+}
\le2\hat B_H\bar\alpha(x_{k'}),
\end{equation*}
which is $(\mathrm{e}3)^+$ at the far (L)-site. Assumption (E) used here is a bound uniform in $x$,
and it carries the factor $E(x)^{\delta_\natural/4}$, where $E(x)=\exp[-(\log(x-c_5))^r]$.
This completes (i).

\emph{The $P_1$-tail.} At (L), work on the disc $|\mathsf{t}|\le\rho_A=\mathsf{w}_L/\eta_p^{1/2}$ of
\eqref{4.25:eq-increment-bound-c}, and write $\mathfrak{d}=d_{sc}(y,\Lambda^{\sim})$. By
\eqref{4.25:eq-depth-hypothesis-3} and $\eta_p\le1$ (in the notation of
Lemma~\ref{4.25:lem-increment-bound}, $\eta_p$ is the exponential of minus a non-negative multiple of
a distance),
\begin{equation*}
\begin{aligned}
|\mathsf{t}|\,\ell_y^3\bigl|P_1(\sigma)(\hat\chi_\square A'_s)\bigr|
&\le\frac{\mathsf{w}_L}{\eta_p^{1/2}}\,\eta_p\,\mathsf{f}(y)\,\mathsf{U}_L\\
&=\mathsf{w}_L\,\eta_p^{1/2}\,\mathsf{U}_L\,e^{-\delta_\natural(\mathfrak{d}-1)}
\le\mathsf{w}_L\,\mathsf{U}_L\,e^{-\delta_\natural(\mathfrak{d}-1)}.
\end{aligned}
\end{equation*}
For $\mathfrak{d}\ge1$ the right side is at most
$\mathsf{w}_L\mathsf{U}_Le^{-\frac12\delta_\natural(\mathfrak{d}-1)}$. For
$0\le\mathfrak{d}<1$ it is at most $\mathsf{w}_L\mathsf{U}_Le^{\delta_\natural}$. Both case bounds used
in the proof of Lemma~\ref{4.25:lem-increment-bound} therefore hold.

\emph{The local cut entries.} As stated in (F), entries $1$, $2$ and $4$ of the cut size
$\mathsf{N}_y(\delta\mathbb{H}_p)$ are at most the corresponding entries of $\mathsf{N}_y(\mathbb{H}_s)$
multiplied by
$c_\chi\,\mathbf{1}[\tilde\Delta(y)\cap\square\neq\emptyset]$. This form follows from the same lattice
product rule that controls the commutator in the construction of the cut partner
$\mathcal{J}^{\sharp}[\chi]$. These entries could be derived from $(\mathrm{e}2)$ exactly as
$(\mathrm{e}5)$ is, the number $c_\chi$ being placed inside $\mathsf{U}_L$. In
Lemma~\ref{4.25:lem-increment-bound} they are taken from assumption (F), not derived from
$(\mathrm{e}2)$.

\emph{The fifth entry and (iii).} The fifth entry is not assumed. It is the bound $(\mathrm{e}5)$ of
\eqref{4.25:eq-increment-bound-a}, derived from the two uncut entries in $(\mathrm{e}2)$. With (A)--(F)
supplying $(\mathrm{e}3)^+$ by (i), and the $P_1$-tail and the local cut entries entering through (F),
no bound on the block sizes other than $(\mathrm{e}1)$, $(\mathrm{e}2)$, $(\mathrm{e}3)^+$ and the
convention (F) enters the proof of Lemma~\ref{4.25:lem-increment-bound}. This is (iii).
\end{proof}

\begin{lemma}[The reading rotation, per link and layer]\label{4.25:lem-reading-rotation}
Let $\mathfrak{O}$ be a later operation, of step $k'$, presenting $\mathfrak{Q}(\mathfrak{a})$ as in
Definition~\ref{4.25:def-increment-hypothesis}, let $e$ be a piece of $\mathfrak{Q}(\mathfrak{a})$ and
$z\in\mathcal{Z}_{\mathfrak{O}}$.
Let $\mathcal{Y}(e)$ be as in Definition~\ref{4.25:def-enlarged-support-blocks}, and
$\theta(y)$, $\bar\theta$, $\bar\theta^{\downarrow}$ and $\Theta(e)=\sup_{y\in\mathcal{Y}(e)}\theta(y)$ as in
\eqref{4.25:eq-flow-inequalities-f}. Put $b:=\hat B_H C'_{age}$, and let
$C_{inc}=2(K_{\mathsf{j}}+2c_\zeta)\hat B_H C'_{age}$ be the constant of
\eqref{4.25:eq-increment-hypothesis-a}. Consider a term of the summation of boundary terms at
finer zones, and the two sites at which its law is read (its law sites), written $b_-$ and $b'_-$
(the letter unrelated to the number $b$); these are sites of the set $S^P_W$ of sites attached to
the region $W$ of the arrest $\mathfrak{a}$, a subset of $S^+(e)$. A site of $\mathfrak{O}$ is of
kind $\mathrm{P}$ (a site at which, as at a
tube site, the presentation carries one layer), of kind $\mathrm{T}$ (a tube
site), or of kind $\mathrm{L}$; a site of kind $\mathrm{L}$ is either far, when the class it treats
does not meet $S^+(e)$, as in Lemma~\ref{4.25:lem-increment-bound}, or deep, the case treated by the
three-field lemma in \textup{(g)}. With the constant $c_4$ entering the ceiling statement of
\textup{(d)}, and with $\alpha_1^{\mathfrak{u}}$ the base
radius of the unitary-factor lemma, put
\begin{equation*}
\varrho_0:=\min\{10^{-2}/K_u,\ c_4,\ \alpha_1^{\mathfrak{u}}\}.
\end{equation*}
Assume the following.
\begin{enumerate}
\item[(a)] The small-coupling hypothesis on the presented pair, entering this lemma only
through the logarithmic bound it gives together with clause (i) of the unitary-factor lemma: at a
site of kind $\mathrm{P}$ or $\mathrm{T}$, for the factor $\mathfrak{u}$ of \textup{(b)} and its one
layer $\mathcal{K}$,
\begin{equation*}
|\log\mathfrak{u}(s)|\le K_u\sup_{B^{top}(s)}N_y(\mathcal{K});
\end{equation*}
and, for the unitary factor $\mathfrak{u}[e^{i\eta\mathcal{K}'}W;W]$ of a layer $\mathcal{K}'$ at a
base $W$ in the base class of \textup{(c)} below, the same bound with $\mathcal{K}'$ in place of
$\mathcal{K}$.
\item[(b)] The formula of clause (i) of the unitary-factor lemma for the factor: at a site of
kind $\mathrm{P}$ or $\mathrm{T}$, where the presentation carries one layer $\mathcal{K}$ over the one
base $P_{\mathfrak{O}}$, independent of $z$, of Definition~\ref{4.25:def-increment-hypothesis}, the
rotation at which the term is read is $g(z)=\mathfrak{r}(z)g_{\mathfrak{O}}$ with $g_{\mathfrak{O}}$
independent of $z$ (the letters $g$, $g_{\mathfrak{O}}$, and $g'$ below, denote group-valued
rotations, not couplings), and
$\mathfrak{r}=\mathfrak{u}[e^{i\eta\mathcal{K}}\mathbf{K};\mathbf{K};\mathbf{B}]$, with $\eta$ the scalar
multiplying the increment in the exponent of the presented pair $p(z)$ of
Definition~\ref{4.25:def-increment-hypothesis}, and $\mathbf{K},\mathbf{B}$ the two configurations
entering the unitary-factor map $\mathfrak{u}$ in clause (i) (here $\mathbf{B}$ is a configuration,
not the large-field operation).
\item[(c)] The base class of the unitary-factor lemma, of base radius $\alpha_1^{\mathfrak{u}}$ (a number
depending on $d$ and $L$ only); it contains every base $W$ whose potential $A'_W$ satisfies
$N_y(A'_W)\le\varrho_0$ for every $y\in\mathcal{Y}(e)$.
\item[(d)] Condition (i) of the frozen Gaussian collar weight of an arrested region,
\eqref{4.25:eq-frozen-collar-weight-a}, with its number $\varpi_G$; the number $\alpha_1^{\mathfrak{u}}$
is among those whose minimum is $\varpi_G$; and the ceiling statement for the links' bases: for
every base $W_i(z)$ of \textup{(f)} and every $y\in\mathcal{Y}(e)$,
\begin{equation*}
N_y(A'_{W_i})\le1.1^3\bigl(\theta(y)+3\max_iN_y(\mathcal{K}_i)\bigr),\qquad
\varpi_G\le\varrho_0 .
\end{equation*}
\item[(e)] The per-layer bound \eqref{4.25:eq-increment-bound-b} of
Lemma~\ref{4.25:lem-increment-bound}, for each of the layers $\mathcal{K}_i$, $i\le r$, of
\textup{(f)} at far sites, and for the one layer at tube sites.
\item[(e$'$)] At sites of kind $\mathrm{P}$ and $\mathrm{T}$, the increment hypothesis
\eqref{4.25:eq-increment-hypothesis-a} for the one layer $\mathcal{K}=A''(z)$ of the presented pair,
with the constant $C_{inc}$, so that $N_y(\mathcal{K})\le C_{inc}\theta(y)$ for $y\in\mathcal{Y}(e)$;
at tube sites this is the conclusion of Lemma~\ref{4.25:lem-increment-bound} with $r=1$.
\item[(f)] The links at far sites: there are $r\le4$ bases $W_i=W_i(z)$ and layers $\mathcal{K}_i$,
the rotation at which the term is read is $g(z)=\mathfrak{r}(z)g_{\mathfrak{O}}$ with
$g_{\mathfrak{O}}$ independent of $z$, and $\mathfrak{r}=a_1\cdots a_r$ with
$a_i:=\mathfrak{u}[e^{i\eta\mathcal{K}_i}W_i;W_i]$.
\item[(g)] At deep sites: the data $\mathsf{u}'$ (a group element), $g'$ and $\lambda$ of the deep
site, with $\mathsf{u}'g'$ independent of $z$, and the field $Z$ of the three-field lemma, with
values in the complexified Lie algebra (not a large-field region); the
rotation at which the term is read there is $g(z)=\mathsf{u}'e^{\lambda Z}g'$.
\end{enumerate}
Then the rotation at which the term is read is $g(z)=\mathfrak{r}(z)g_{\mathfrak{O}}$ with
$g_{\mathfrak{O}}$ independent of $z$ (at deep sites $g_{\mathfrak{O}}=\mathsf{u}'g'$ and
$\mathfrak{r}=\mathsf{u}'e^{\lambda Z}\mathsf{u}'^{-1}$), and, for $s\in\{b_-,b'_-\}$:
\begin{enumerate}
\item[(1)] at sites of kind $\mathrm{P}$ and $\mathrm{T}$ and at far sites,
\begin{equation}\label{4.25:eq-reading-rotation-a}
\operatorname{dis}(\mathfrak{r}(z)(s))\le C_uC_{inc}\Theta(e),\qquad C_u:=1.1K_u;
\end{equation}
\item[(2)] at deep sites,
\begin{equation*}
\operatorname{dis}(\mathfrak{r}(z)(s))\le e^{2\operatorname{pol}\mathsf{u}'}|\lambda Z(s)|.
\end{equation*}
\end{enumerate}
\end{lemma}

\begin{proof}
Throughout, $\operatorname{dis}$ and $\operatorname{pol}$ are as in
Lemma~\ref{4.25:lem-word-length-distance}, and we use its clause (d2) in the display
\eqref{4.25:eq-word-length-distance-a}: subadditivity under products, the bound of
$\operatorname{dis}$ of a conjugate by $e^{2\operatorname{pol}}$ of the conjugating element times
$\operatorname{dis}$, and $\operatorname{dis}(h)\le|\log h|$ for every logarithm of $h$.

\emph{The law sites.} Each $s\in\{b_-,b'_-\}$ is a site of $S^P_W\subset S^+(e)$. By
Definition~\ref{4.25:def-enlarged-support-blocks}, $\mathcal{Y}(e)$ is the set of top blocks
$B^{top}(p)$ of the sites $p$ of $S^+(e)$; hence $B^{top}(s)\in\mathcal{Y}(e)$, and every supremum
over $B^{top}(s)$ below is bounded by the corresponding supremum over $y\in\mathcal{Y}(e)$.

\emph{Sites of kind $\mathrm{P}$ and $\mathrm{T}$.} Here the presentation carries the one layer
$\mathcal{K}$ of the increment hypothesis, and
$\mathfrak{r}=\mathfrak{u}[e^{i\eta\mathcal{K}}\mathbf{K};\mathbf{K};\mathbf{B}]$ by (b). In the
display below the first inequality is $\operatorname{dis}(h)\le|\log h|$ of (d2), the second is the
bound of (a), the third is hypothesis (e$'$) together with $B^{top}(s)\in\mathcal{Y}(e)$, the
equality is the definition of $\Theta(e)$, and the last inequality is $K_u\le1.1K_u=C_u$:
\begin{align*}
\operatorname{dis}(\mathfrak{r}(z)(s))&\le|\log\mathfrak{u}(s)|
\le K_u\sup_{B^{top}(s)}N_y(\mathcal{K})\\
&\le K_u\,C_{inc}\sup_{y\in\mathcal{Y}(e)}\theta(y)=K_uC_{inc}\Theta(e)\le C_uC_{inc}\Theta(e),
\end{align*}
which is \eqref{4.25:eq-reading-rotation-a}.

\emph{The link bases lie in the base class.} At far sites the bound of (a) is applied to the link
factors $a_i$, and there it requires their bases to lie in the base class of (c). For the bases
$W_i(z)$ of the links used at far sites below, this follows from the
chain
\begin{align*}
N_y(A'_{W_i})&\le 1.1^3\bigl(\theta(y)+3\max_iN_y(\mathcal{K}_i)\bigr)
\le 1.331(1+1.2b)\,\theta(y)\\
&\le 6.75\,b\,\bar\theta(t_y)=\tfrac98\cdot6b\cdot\bar\theta(t_y)
\le\tfrac98\,C_{inc}\,\bar\theta^{\downarrow}(x_{k'}-2c_2)\le\varpi_G\le\varrho_0 ,
\end{align*}
in which $t_y$ is the argument at which the envelope $\bar\theta$ is evaluated for the block $y$ in
items (1) and (2) of \eqref{4.25:eq-conversion-old-factors-a}.
The first inequality and the last one, $\varpi_G\le\varrho_0$, are the two inequalities of the
ceiling statement in (d), the member $\alpha_1^{\mathfrak{u}}$ of $\varrho_0$ being among the
numbers whose minimum is $\varpi_G$; this radius is the base radius of the base class in (c), not the
radius $\alpha_1^{SL}$ of the class $\mathfrak{K}_G$. The second inequality uses $1.1^3=1.331$ and
$\max_iN_y(\mathcal{K}_i)\le\frac25b\,\theta(y)$, which follows from \eqref{4.25:eq-increment-bound-b},
that is $\mathsf{N}_y(\mathcal{K}_i)\le2b\,\bar\alpha(x_{\ell'(y)})$, from
$N_y(\mathcal{K}_i)\le\mathsf{N}_y(\mathcal{K}_i)$ ($N_y$ being one of the entries of the size
$\mathsf{N}_y$), and from $\bar\alpha(x_{\ell'(y)})\le\theta(y)/5$, a property of the block threshold
$\theta(y)$ of \eqref{4.25:eq-flow-inequalities-f}. The
third inequality uses $\theta(y)\le\bar\theta(t_y)$ together with
$1.331(1+1.2b)\le6.75\,b$; the latter is equivalent to $1.331\le(6.75-1.5972)\,b=5.1528\,b$, that is,
to $b\ge0.2583$, and here $b=\hat B_HC'_{age}\ge2$ by \eqref{4.25:eq-tube-constant-a}. The equality
is $\frac98\cdot6=6.75$. The fifth inequality uses $C_{inc}=2(K_{\mathsf{j}}+2c_\zeta)\,b\ge6b$, valid
because $K_{\mathsf{j}}+2c_\zeta\ge3$ ($K_{\mathsf{j}}$ and $c_\zeta$ being the constants entering
$C_{inc}$ in \eqref{4.25:eq-increment-hypothesis-a}), and
\begin{equation*}
\theta(y)\le\bar\theta(t_y)\le\bar\theta^{\downarrow}(x_{k'}-2c_2),
\end{equation*}
which are items (1) and (2) of the conversion of old polylogarithmic factors,
\eqref{4.25:eq-conversion-old-factors-a}. The sixth inequality is condition (i) of the frozen
Gaussian collar weight, \eqref{4.25:eq-frozen-collar-weight-a}, of hypothesis (d).

Directly also: for $y\in\mathcal{Y}(e)$ the sixth inequality of the chain, together
with $\theta(y)\le\bar\theta^{\downarrow}(x_{k'}-2c_2)$, gives $\frac98C_{inc}\theta(y)\le\varpi_G$, hence
$\Theta(e)\le\frac89\varpi_G/C_{inc}$. With $\max_iN_y(\mathcal{K}_i)\le\frac25b\,\Theta(e)$ and
$C_{inc}\ge6b$ this gives
\begin{equation*}
\max_iN_y(\mathcal{K}_i)\le\tfrac25\cdot\tfrac16\cdot\tfrac89\,\varpi_G=\tfrac{8}{135}\,\varpi_G ,
\end{equation*}
and with $C_{inc}\ge300$ from \eqref{4.25:eq-tube-constant-a} it gives
$\Theta(e)\le\frac{8}{2700}\varpi_G$. The first inequality of the chain then yields
\begin{align*}
N_y(A'_{W_i})&\le1.331\bigl(\Theta(e)+3\cdot\tfrac{8}{135}\varpi_G\bigr)
\le1.331\bigl(\tfrac{8}{2700}+\tfrac{24}{135}\bigr)\varpi_G\le0.2406\,\varpi_G ,
\end{align*}
since $\frac{8}{2700}+\frac{24}{135}\le0.18075$ and $1.331\cdot0.18075\le0.2406$. Thus the link bases
satisfy $N_y(A'_{W_i})\le\varpi_G$, and it is to place them in the base class of (c), of radius
$\alpha_1^{\mathfrak{u}}$, that $\alpha_1^{\mathfrak{u}}$, a number depending on $d$ and $L$ only, is among the
numbers whose minimum is $\varpi_G$.

\emph{Far sites.} By (f), $\mathfrak{r}=a_1\cdots a_r$ with
$a_i=\mathfrak{u}[e^{i\eta\mathcal{K}_i}W_i;W_i]$ and $r\le4$. Subadditivity in (d2) gives the first
inequality below and $\operatorname{dis}(h)\le|\log h|$ the second; the third is the bound of (a) at
the base $W_i$, which lies in the base class by the chain above; the fourth is
\eqref{4.25:eq-increment-bound-b} applied to each of the $r$ layers separately, with
$B^{top}(s)\in\mathcal{Y}(e)$:
\begin{align*}
\operatorname{dis}(\mathfrak{r}(z)(s))&\le\sum_{i\le r}\operatorname{dis}(a_i(s))
\le\sum_{i\le r}|\log a_i(s)|
\le\sum_{i\le r}K_u\sup_{B^{top}(s)}N_y(\mathcal{K}_i(z))\\
&\le r\cdot K_u\cdot2b\sup_{y\in\mathcal{Y}(e)}\bar\alpha(x_{\ell'(y)})
\le4K_u\cdot\tfrac25\,b\,\Theta(e)\le C_uC_{inc}\Theta(e).
\end{align*}
The fifth inequality uses $r\le4$ and $\bar\alpha(x_{\ell'(y)})\le\theta(y)/5$, the property of the
block threshold of \eqref{4.25:eq-flow-inequalities-f} used above. The last one is
\begin{equation*}
1.6K_u\,b\le1.1K_u\cdot2(K_{\mathsf{j}}+2c_\zeta)\,b=C_uC_{inc},
\end{equation*}
which holds because $2(K_{\mathsf{j}}+2c_\zeta)\ge6$, as $K_{\mathsf{j}}+2c_\zeta\ge3$, so the
right side is at least $6.6K_ub$. This is
\eqref{4.25:eq-reading-rotation-a} at far sites. The bound is read link by link and layer by layer:
each link is estimated through its own layer by \eqref{4.25:eq-increment-bound-b}, not through the
net increment of the presented pair, which may cancel between links; no cancellation between the
links is used, and the product $a_1\cdots a_r$ is not recombined into a single exponential.

\emph{Deep sites.} By (g),
$g(z)=\mathsf{u}'e^{\lambda Z}g'=(\mathsf{u}'e^{\lambda Z}\mathsf{u}'^{-1})(\mathsf{u}'g')$; thus
$g_{\mathfrak{O}}:=\mathsf{u}'g'$, which is independent of $z$, and
$\mathfrak{r}=\mathsf{u}'e^{\lambda Z}\mathsf{u}'^{-1}$, with $Z$ the field supplied by the three-field
lemma. The conjugation bound of (d2) gives
\begin{equation*}
\operatorname{dis}(\mathfrak{r}(z)(s))\le e^{2\operatorname{pol}\mathsf{u}'}
\operatorname{dis}(e^{\lambda Z(s)}),
\end{equation*}
and since $\lambda Z(s)$ is a logarithm of $e^{\lambda Z(s)}$, (d2) gives
$\operatorname{dis}(e^{\lambda Z(s)})\le|\lambda Z(s)|$. This is clause (2).

In each case $g(z)=\mathfrak{r}(z)g_{\mathfrak{O}}$ with $g_{\mathfrak{O}}$ independent of $z$, by
(b) at sites of kind $\mathrm{P}$ and $\mathrm{T}$, by (f) at far sites and by (g) at deep sites, and
the lemma is proved.
\end{proof}

\begin{corollary}[The oscillation datum at later presentations]\label{4.25:cor-oscillation-datum}
Let $\mathfrak{a}$ be an arrest, $\mathfrak{Q}(\mathfrak{a})$ its family of pieces and
$e\in\mathfrak{Q}(\mathfrak{a})$ a piece, and assume, for every such $e$, the following.
\begin{itemize}
\item The increment hypothesis (Definition~\ref{4.25:def-increment-hypothesis}) holds on
$\mathcal{Y}(e)$. At a presentation of $\mathfrak{Q}(\mathfrak{a})$ by a later operation
$\mathfrak{O}$, $p(z)$ denotes the pair over $S^+(e)$ of that definition. Its base
$P_{\mathfrak{O}}=(\mathbb{U}_{\mathfrak{O}},\mathbb{J}_{\mathfrak{O}})$ does not depend on $z$,
and its increments $(A''(z),J''(z))$ are holomorphic in $z$. The constant $C_{inc}$ is that of
\eqref{4.25:eq-increment-hypothesis-a}.
\item The hypotheses of Lemma~\ref{4.25:lem-reading-rotation} hold, among them condition~(i) of
\eqref{4.25:eq-frozen-collar-weight-a}. So do the inputs from which the bound of that lemma is
derived; at the far sites of Lemma~\ref{4.25:lem-increment-bound} this means the per-layer bound
\eqref{4.25:eq-increment-bound-b}. The rotation at which a piece is read is
$g(z)=\mathfrak{r}(z)g_{\mathfrak{O}}$, with $g_{\mathfrak{O}}$ independent of $z$, as in that
lemma, and $C_u=1.1K_u$.
\item $\operatorname{pol}(g_{\mathfrak{O}})\le\frac16$. At chart points the condition in force on
the polar size is $\operatorname{pol}<\varepsilon_\star$, where $\varepsilon_\star$ denotes the
chart's threshold for the polar size.
\end{itemize}
Let $\mathsf{c}_U,\mathsf{c}_J$ be the constants of \eqref{4.25:eq-piece-oscillation-a}, and let $\Theta(e)$ be
as in \eqref{4.25:eq-flow-inequalities-f}. Write $P{\restriction}_{\tilde S(e)}$ for the restriction
of a pair $P$ to the bonds of $\tilde S(e)$. Then the following holds for every
$e\in\mathfrak{Q}(\mathfrak{a})$, every later operation $\mathfrak{O}$, every
$z\in\mathcal{Z}_{\mathfrak{O}}$ and every collar variable $x$ with $|x|\le\mathfrak{b}_x$:
\begin{equation}\label{4.25:eq-oscillation-datum-a}
\begin{gathered}
\bigl|v_e^{\sharp}(x;p(z);g(z))
-v_e^{\sharp}(x;P_{\mathfrak{O}}{\restriction}_{\tilde S(e)};g_{\mathfrak{O}})\bigr|
\le\mathfrak{N}^G(e;\mathfrak{O}):=C_G\,\nu^{\natural}(e)\,\Theta(e),\\
C_G:=(\mathsf{c}_U+\mathsf{c}_J+3C_u)\,C_{inc}.
\end{gathered}
\end{equation}
Moreover, $z\mapsto v_e^{\sharp}(x;p(z);g(z))$ is holomorphic on $\mathcal{Z}_{\mathfrak{O}}$.
\end{corollary}

\begin{proof}
Fix $e$, $\mathfrak{O}$ (at step $k'$), $z\in\mathcal{Z}_{\mathfrak{O}}$ and $x$ as in the statement.
The proof has six steps.

\emph{The base.} We apply Lemma~\ref{4.25:lem-gaussian-pieces}(iii), whose proof is incomplete at
one step, to $\mathfrak{O}$ and to the space $\mathfrak{T}$ containing $P_{\mathfrak{O}}$. It gives
three things:
\begin{itemize}
\item $P_{\mathfrak{O}}\in\mathfrak{K}_G$ as a global pair;
\item the base bound on $\mathcal{Y}(e)$;
\item the bound $\ell^3|\mathbb{J}_{\mathfrak{O}}|\le\gamma_0^{(k')}$ everywhere, with the constant
$\gamma_0^{(k')}$ of that part.
\end{itemize}
Write $\mathbb{U}_{\mathfrak{O}}=e^{i\eta A'_{P_{\mathfrak{O}}}}U$, in the form of the pairs of part
(iv) of the same lemma. Remark~\ref{4.25:rem-tube-constant} gives
$\Theta(e)\le\varepsilon_{reg}/6750$ (see \eqref{4.25:eq-tube-constant-a}). Hence, for
$y\in\mathcal{Y}(e)$,
\begin{equation}\label{4.25:eq-oscillation-datum-1}
N_y(A'_{P_{\mathfrak{O}}})<\theta(y)\le\Theta(e)\le\frac{\varepsilon_{reg}}{6750}
<\frac{\varepsilon_{reg}}{5}.
\end{equation}
For the current, with $\bar\gamma_0$ the ceiling of the constants $\gamma_0^{(k')}$,
\begin{equation}\label{4.25:eq-oscillation-datum-2}
\ell^3|\mathbb{J}_{\mathfrak{O}}|\le\gamma_0^{(k')}\le\bar\gamma_0\le\tfrac14\gamma_{reg}
\le\tfrac12\gamma_{reg}.
\end{equation}
These are the bounds on the base pair required in Lemma~\ref{4.25:lem-gaussian-pieces}(iv), with
$P:=P_{\mathfrak{O}}$.

\emph{The increments.} Let $\mathfrak{a}''$ and $\mathfrak{j}''$ be the sizes of
Lemma~\ref{4.25:lem-gaussian-pieces}(iv), whose proof is incomplete at one step, formed with
$A''=A''(z)$ and $J''=J''(z)$. They are suprema over $S^+(e)$ taken in the units of the old blocks.
By Lemma~\ref{4.25:lem-domination-old-blocks}(ii), such old-unit suprema are bounded by the
corresponding current-unit suprema over the blocks of $\mathcal{Y}(e)$. By
\eqref{4.25:eq-increment-hypothesis-a}, these are at most $C_{inc}\theta(y)\le C_{inc}\Theta(e)$.
Condition~(i) of \eqref{4.25:eq-frozen-collar-weight-a} then yields
\begin{equation}\label{4.25:eq-oscillation-datum-3}
\mathfrak{a}'',\ \mathfrak{j}''\le C_{inc}\Theta(e)\le\tfrac89\varpi_G\le\varpi_G
\le\min\bigl\{\tfrac15\varepsilon_{reg},\ \tfrac14\gamma_{reg}\bigr\}.
\end{equation}
So the hypotheses of Lemma~\ref{4.25:lem-gaussian-pieces}(iv) hold with $P=P_{\mathfrak{O}}$.
Moreover, $T_U=\varepsilon_{reg}/\mathfrak{a}''\ge5$ and $T_J=\frac12\gamma_{reg}/\mathfrak{j}''\ge2$.

\emph{The rotation.} Put $g_1:=g_{\mathfrak{O}}$ and $g_2:=g(z)$. Then $g_2g_1^{-1}=\mathfrak{r}(z)$.
Lemma~\ref{4.25:lem-reading-rotation}, whose hypotheses and inputs are assumed, bounds
$\operatorname{dis}(\mathfrak{r}(z)(s))$ by $C_uC_{inc}\Theta(e)$. Hence the quantity
$\mathfrak{g}$ of Lemma~\ref{4.25:lem-piece-oscillation} satisfies
\begin{equation}\label{4.25:eq-oscillation-datum-4}
\begin{split}
\mathfrak{g}&\le C_uC_{inc}\Theta(e)\le 1.1K_u\cdot\tfrac89\cdot\frac{10^{-2}}{K_u}\\
&=\tfrac{8.8}{9}\cdot10^{-2}\le0.978\cdot10^{-2}\le\tfrac16 .
\end{split}
\end{equation}
Here we used $C_{inc}\Theta(e)\le\frac89\varpi_G$ from \eqref{4.25:eq-oscillation-datum-3},
together with the ceiling $\varpi_G\le\varrho_0\le10^{-2}/K_u$ that accompanies
\eqref{4.25:eq-frozen-collar-weight-a}, $\varrho_0$ being a minimum of numbers among which
is $10^{-2}/K_u$.

\emph{The group elements.} By hypothesis, $\operatorname{pol}(g_1)=\operatorname{pol}(g_{\mathfrak{O}})\le\frac16$,
which is at most $\frac13$.

For $g_2$, apply $\|gh\|\le\|g\|\,\|h\|$ to $\mathfrak{r}g_{\mathfrak{O}}$ and to
$(\mathfrak{r}g_{\mathfrak{O}})^{-1}=g_{\mathfrak{O}}^{-1}\mathfrak{r}^{-1}$. Since
$\operatorname{pol}$ is the logarithm of $\max\{\|h\|,\|h^{-1}\|\}$, this gives
$\operatorname{pol}(\mathfrak{r}g_{\mathfrak{O}})\le\operatorname{pol}(\mathfrak{r})
+\operatorname{pol}(g_{\mathfrak{O}})$.

By (d3) of \eqref{4.25:eq-word-length-distance-a} (Lemma~\ref{4.25:lem-word-length-distance}),
$\operatorname{pol}(\mathfrak{r})\le\operatorname{dis}(\mathfrak{r})$. Therefore
\begin{equation}\label{4.25:eq-oscillation-datum-5}
\operatorname{pol}(g(z))\le\operatorname{dis}(\mathfrak{r}(z))+\operatorname{pol}(g_{\mathfrak{O}})
\le0.00978+\tfrac16\le\tfrac13 .
\end{equation}

\emph{The oscillation.} By Lemma~\ref{4.25:lem-gaussian-pieces}(iv), whose proof is incomplete at one
step, and with $P=P_{\mathfrak{O}}$, $A''=A''(z)$, $J''=J''(z)$, the pairs of
Lemma~\ref{4.25:lem-piece-oscillation} are
\begin{equation*}
p_1=P_{\mathfrak{O}}{\restriction}_{\tilde S(e)},\qquad
p_2=p(z){\restriction}_{\tilde S(e)},
\end{equation*}
because $P^U=\mathbb{U}_{\mathfrak{O}}$ and $\mathbb{J}=\mathbb{J}_{\mathfrak{O}}$.

The remaining hypotheses of Lemma~\ref{4.25:lem-piece-oscillation}, whose proof is incomplete at one
step, also hold:
\begin{itemize}
\item $\operatorname{pol}(g_i)\le\frac13$ for $i=1,2$, by \eqref{4.25:eq-oscillation-datum-5} and
the bound on $g_1$;
\item $\mathfrak{g}\le\frac13$, by \eqref{4.25:eq-oscillation-datum-4};
\item the collar variable obeys $|x|\le\mathfrak{b}_x$.
\end{itemize}
The right member of that lemma is linear in $(\mathfrak{a}'',\mathfrak{j}'',\mathfrak{g})$ with
non-negative coefficients. Inserting \eqref{4.25:eq-oscillation-datum-3} and
\eqref{4.25:eq-oscillation-datum-4} therefore gives
\begin{equation*}
\begin{split}
\bigl|v_e^{\sharp}(x;p_2;g_2)-v_e^{\sharp}(x;p_1;g_1)\bigr|
&\le\nu^{\natural}(e)\bigl[\mathsf{c}_U\mathfrak{a}''+\mathsf{c}_J\mathfrak{j}''+3\mathfrak{g}\bigr]\\
&\le\nu^{\natural}(e)\,(\mathsf{c}_U+\mathsf{c}_J+3C_u)\,C_{inc}\,\Theta(e).
\end{split}
\end{equation*}
This is \eqref{4.25:eq-oscillation-datum-a}.

\emph{Holomorphy.} Put
\begin{equation*}
Q_z:=\bigl(e^{i\eta\chi_eA''(z)}\mathbb{U}_{\mathfrak{O}},\ \mathbb{J}_{\mathfrak{O}}+\chi_eJ''(z)\bigr).
\end{equation*}
This is the pair $Q(1,1)$ of Lemma~\ref{4.25:lem-gaussian-pieces}(iv), whose proof is incomplete at
one step, formed with $A''(z)$ and $J''(z)$. Since $T_U\ge5$ and $T_J\ge2$, the point $(1,1)$ lies
in the bidisc of that part. Hence $Q_z$ is global, lies in $\mathfrak{K}_G$, and coincides with
$p(z)$ on the bonds of $\tilde S(e)$:
\begin{equation*}
p(z){\restriction}_{\tilde S(e)}=Q_z{\restriction}_{\tilde S(e)}.
\end{equation*}
The increments $(A''(z),J''(z))$ are holomorphic in $z$ by the increment hypothesis. So
$z\mapsto Q_z$ is a holomorphic family of global $\mathfrak{K}_G$-pairs.

Part~(i) of the same lemma therefore makes $z\mapsto v_e(\,\cdot\,;Q_z{\restriction}_{\tilde S(e)})$
holomorphic. Part~(ii) expresses $v_e^{\sharp}(x;p;g)=v_e(x^{[g]};p)$ as a polynomial in the entries
of $\operatorname{Ad}_g$. Together these give the holomorphy of
$z\mapsto v_e^{\sharp}(x;p(z);g(z))$ on $\mathcal{Z}_{\mathfrak{O}}$.
\end{proof}

\begin{lemma}[Piece geometry and the pile bound]\label{4.25:lem-pile-bound}
Let $\mathfrak{a}$ be an arrest as in Definition~\ref{4.25:not-collar-weight}, with region $W$
and integrated level $j_W$.

Fix a step after the arrest. We use the following terms.
\begin{itemize}
\item The \emph{current cells} are the cells of the family of that step, and $\mathfrak{F}$ denotes
the set of current cells.
\item For a cell or a cube $c$ of an earlier family, $\hat c\in\mathfrak{F}$ denotes the current
cell containing it (its \emph{hat}).
\item $\operatorname{lev}c$ denotes the level of a cell $c$.
\item For a piece $e$ of an arrest $\mathfrak{a}'$ (in particular for $e\in\mathfrak{E}^G_W$, with
$\mathfrak{a}'=\mathfrak{a}$) the cells of $\tilde S(e)$ and of $S^+(e)$ are cells of
$\mathfrak{B}_{\mathfrak{a}'}$. The \emph{current levels on
$S^+(e)$} are the levels $\operatorname{lev}\hat c$, where $c$ runs over the cells of $S^+(e)$. By
Lemma~\ref{4.25:lem-domination-old-blocks}\,(i) these are the levels of the current cells of the
blocks of $\mathcal{Y}(e)$.
\item For an arrest $\mathfrak{a}'$ we write $e\in\mathfrak{Q}(\mathfrak{a}')$ when
$v_e\in\mathfrak{Q}(\mathfrak{a}')$.
\item $\pi_\ell$ is the partition into the cubes of the level-$\ell$ grid of cells.
\item For $e\in\mathfrak{Q}(\mathfrak{a}')$, $m_0(e)$ is the top level of the cells of
$\mathfrak{B}_{\mathfrak{a}'}$ in $S^+(e)$.
\end{itemize}

Then the following hold.
\begin{enumerate}
\item[(a)] For every $e\in\mathfrak{E}^G_W$ and all cells $c,c'$ of $S^+(e)$, the levels
$\operatorname{lev}c$ and $\operatorname{lev}\hat c$ are at most $j_W+n(e)+2$. The current levels on
$S^+(e)$ differ by at most $n(e)+3$. This is the first line of \eqref{4.25:eq-pile-bound-a}.
\item[(b)] Let $\mathfrak{A}\subset\mathfrak{F}$ be a finite set of current cells. The sum over live
arrests $\mathfrak{a}'$ of the kind of Definition~\ref{4.25:not-collar-weight}, and over the pieces
$e$ whose set $S^+(e)$ meets $\bigcup\mathfrak{A}$, obeys the second line of
\eqref{4.25:eq-pile-bound-a}.
\item[(c)] For every level $\ell$ and every cube $\square'$ of $\pi_\ell$, the third line of
\eqref{4.25:eq-pile-bound-a} holds.
\end{enumerate}
\begin{equation}\label{4.25:eq-pile-bound-a}
\begin{gathered}
\max\{\operatorname{lev}c,\ \operatorname{lev}\hat c\}\le j_W+n(e)+2,\qquad
|\operatorname{lev}\hat c-\operatorname{lev}\hat c'|\le n(e)+3,\\
\sum_{\mathfrak{a}'\ \mathrm{live}}\
\sum_{\substack{e\in\mathfrak{Q}(\mathfrak{a}')\\ S^+(e)\cap\bigcup\mathfrak{A}\neq\emptyset}}
\mathfrak{n}^{\circ}(e)\le C_\alpha|\mathfrak{A}|K_a^{\circ},
\qquad C_\alpha:=2L^de^{4\mathsf{u}},\\
\sum_{\mathfrak{a}'}\
\sum_{\substack{e\in\mathfrak{Q}(\mathfrak{a}')\\ m_0(e)=\ell,\ S^+(e)\cap\square'\neq\emptyset}}
\mathfrak{n}^{\circ}(e)\le C_\beta K_a^{\circ},
\qquad C_\beta:=2\cdot3^dL^{9d}e^{4\mathsf{u}}.
\end{gathered}
\end{equation}
None of these bounds carries a factor counting arrests, nesting depths or scales.
\end{lemma}

\begin{proof}
\emph{Roots.} Let $\mathfrak{a}'$ be a live arrest with region $W'$ and integrated level $j_{W'}$,
and let $e\in\mathfrak{Q}(\mathfrak{a}')$. Choose a cube $q_e\in\mathsf{b}(e)$. It is a cube of one
of the bonds $b,b'$ of the set $S^P_{W'}\subset W'$ carrying the collar variable of $W'$
(Definition~\ref{4.25:not-collar-weight}), hence a collar cube, contained in $W'$.

Two earlier statements now apply. The first is the lemma according to which the bands of all later
steps (the bands being the sets on which those steps integrate) lie off $X_{\mathfrak{a}'}$, the
large-field region attached to the live arrest $\mathfrak{a}'$. The second is the rule according to
which the band of an arrest of this kind is frozen at level one above its integrated level.
Together they give three facts about $\hat q_e$:
\begin{itemize}
\item $\hat q_e$ is a cell of the family $\mathfrak{B}_{\mathfrak{a}'}$ itself, namely the cell of
that family holding $q_e$, which is still a current cell;
\item its level lies between $j_{W'}$ and $j_{W'}+1$;
\item it contains collar cubes of $W'$ alone, and at most $L^d$ of them.
\end{itemize}

Fix $c\in\mathfrak{F}$. All pieces $e$ with $\hat q_e=c$ belong to the single arrest whose collar
cubes $c$ contains, and their cubes $q_e$ lie among at most $L^d$ cubes. The estimate
\eqref{4.25:eq-regularity-class-a} bounds the sum of $\mathfrak{n}^{\circ\circ}$ over the pieces
rooted at one cube by $K_a^{\circ}$, uniformly in the arrest. Hence
\begin{equation}\label{4.25:eq-pile-bound-1}
\sum_{\mathfrak{a}'\ \mathrm{live}}\
\sum_{\substack{e\in\mathfrak{Q}(\mathfrak{a}')\\ \hat q_e=c}}\mathfrak{n}^{\circ\circ}(e)
\le L^dK_a^{\circ}\qquad (c\in\mathfrak{F}).
\end{equation}

Let $D$ denote the distance in the touching graph of current cells
(Definition~\ref{4.25:not-collar-weight}). Touching cells have touching or equal hats. Indeed, if
the closures of $c$ and $c'$ meet, then so do the closures of $\hat c\supset c$ and
$\hat c'\supset c'$.

Let $\hat Y(e)$ be the set of current cells containing cells of $\tilde S(e)$. It contains
$\hat q_e$, because $q_e$ lies in a cell of $\tilde S(e)$. Since $\tilde S(e)$ is a connected family
of cells, $\hat Y(e)$ is connected in the touching graph. Moreover $|\hat Y(e)|\le n(e)$, since
$\tilde S(e)$ consists of at most $\max\{|Y'|,2\}\le n(e)$ cells. Every cell of $S^+(e)$ lies at
cell-distance at most $2$ from a cell of $\tilde S(e)$, so every hat of a cell of $S^+(e)$ lies at
distance $D\le2$ from $\hat Y(e)$.

Every cell of $\hat Y(e)$ is joined to $\hat q_e$ by a path in $\hat Y(e)$ of at most
$|\hat Y(e)|-1$ steps, and by the previous paragraph every hat $\hat c$ of a cell of $S^+(e)$ is at
most two further steps away. This gives, for every cell $c$ of $S^+(e)$,
\begin{equation}\label{4.25:eq-pile-bound-2}
D(\hat q_e,\hat c)\le(|\hat Y(e)|-1)+2\le n(e)+1 .
\end{equation}

\emph{Part (a).} Apply the above with $\mathfrak{a}'=\mathfrak{a}$. By
\eqref{4.25:eq-collar-weight-a}, applied to the current cells, each step of the chain in
\eqref{4.25:eq-pile-bound-2} changes the level by at most one. Hence
\begin{equation*}
\operatorname{lev}\hat c\le\operatorname{lev}\hat q_e+n(e)+1\le j_W+n(e)+2 .
\end{equation*}

Now take two cells $c,c'$ of $S^+(e)$. Their hats are each at distance at most $2$ from
$\hat Y(e)$, and $\hat Y(e)$ has diameter at most $|\hat Y(e)|-1\le n(e)-1$. So
$D(\hat c,\hat c')\le n(e)+3$, and the same rule bounds $|\operatorname{lev}\hat c-\operatorname{lev}\hat c'|$
by $n(e)+3$.

The levels $\operatorname{lev}c$ of the cells of $\mathfrak{B}_{\mathfrak{a}}$ in $S^+(e)$ are bounded
in the same way. The cell of $\mathfrak{B}_{\mathfrak{a}}$ holding $q_e$ is $\hat q_e$ by the roots
paragraph, so its level is at most $j_W+1$. The chain is now run through the cells of
$\mathfrak{B}_{\mathfrak{a}}$, from this cell through $\tilde S(e)$ and then at most two further
steps, at most $n(e)+1$ steps in all; \eqref{4.25:eq-collar-weight-a}, applied to the family
$\mathfrak{B}_{\mathfrak{a}}$, bounds the level change along each step by one. Hence
$\operatorname{lev}c\le j_W+n(e)+2$.

\emph{Part (b).} Let $e\in\mathfrak{Q}(\mathfrak{a}')$ with $S^+(e)\cap\bigcup\mathfrak{A}\neq\emptyset$.
By \eqref{4.25:eq-pile-bound-2} this forces
\begin{equation*}
D(\hat q_e,\mathfrak{A})\le n(e)+1=|Y'|+3 .
\end{equation*}
Since $\mathfrak{n}^{\circ}(e)=\mathfrak{n}^{\circ\circ}(e)e^{-\mathsf{u}|Y'|}$, we get
\begin{equation*}
\mathfrak{n}^{\circ}(e)\le e^{3\mathsf{u}}\mathfrak{n}^{\circ\circ}(e)e^{-\mathsf{u}D(\hat q_e,\mathfrak{A})}
\le e^{3\mathsf{u}}\mathfrak{n}^{\circ\circ}(e)\sum_{c'\in\mathfrak{A}}e^{-\mathsf{u}D(\hat q_e,c')} .
\end{equation*}

A cell touches at most $\mathsf{n}_{nb}$ cells (Definition~\ref{4.25:not-collar-weight}). Hence the
number of cells at distance $D=m$ from a given cell is at most
$\mathsf{n}_{nb}^m\le e^{\mathsf{c}_{nb}m}$. Group the pieces according to their root cell
$\hat q_e$, and bound the sum over pieces with a given root cell by \eqref{4.25:eq-pile-bound-1}.
This yields
\begin{equation*}
\sum_{\mathfrak{a}'\ \mathrm{live}}\sum_{e}\mathfrak{n}^{\circ}(e)
\le|\mathfrak{A}|L^de^{3\mathsf{u}}K_a^{\circ}\sum_{m\ge0}e^{-(\mathsf{u}-\mathsf{c}_{nb})m} .
\end{equation*}

By the definition of $\mathsf{u}$ in Definition~\ref{4.25:not-regularity-class}, and because
$\mathsf{c}_{an}$, $\mathsf{c}_{nb}$ and $\ln L$ are non-negative, we have
$\mathsf{u}-\mathsf{c}_{nb}\ge1$. The last sum is therefore at most $e/(e-1)<2$. The total is
bounded by $2L^de^{3\mathsf{u}}|\mathfrak{A}|K_a^{\circ}\le C_\alpha|\mathfrak{A}|K_a^{\circ}$.

\emph{Part (c).} Let $e$ be a piece counted in the third line of \eqref{4.25:eq-pile-bound-a}, and
write $n=n(e)\ge2$.

\emph{Level of the root.} Part (a), applied to $\mathfrak{a}'$, gives $\ell=m_0(e)\le j_{W'}+n+2$.
By the roots paragraph,
\begin{equation*}
\operatorname{lev}\hat q_e\ge j_{W'}\ge\ell-n-2 .
\end{equation*}

\emph{A cell meeting $\square'$.} Since $S^+(e)$ meets $\square'$, some cell $c$ of $S^+(e)$ meets
$\square'$, and then $c_0:=\hat c$ is a current cell meeting $\square'$. By
\eqref{4.25:eq-pile-bound-2}, $D(\hat q_e,c_0)\le n+1$. By \eqref{4.25:eq-collar-weight-a},
\begin{equation*}
\operatorname{lev}c_0\ge\operatorname{lev}\hat q_e-(n+1)\ge\ell-2n-3 .
\end{equation*}

\emph{Counting.} There are at most $3^dL^{d(2n+3)}$ current cells $c_0$ of level at least
$\ell-2n-3$ meeting $\square'$. For each such $c_0$, the current cells at distance at most $n+1$
from $c_0$ number at most
\begin{equation*}
\sum_{m=0}^{n+1}e^{\mathsf{c}_{nb}m}\le2e^{\mathsf{c}_{nb}(n+1)},
\end{equation*}
since $e^{\mathsf{c}_{nb}}=1+\mathsf{n}_{nb}\ge2$. Each of these cells carries at most
$L^dK_a^{\circ}$ by \eqref{4.25:eq-pile-bound-1}.

\emph{Summation.} Finally $e^{-\mathsf{u}|Y'|}=e^{2\mathsf{u}}e^{-\mathsf{u}n}$, so
\begin{equation*}
\sum_{\mathfrak{a}'}\sum_{e}\mathfrak{n}^{\circ}(e)\le
\sum_{n\ge2}3^dL^{d(2n+3)}\cdot2e^{\mathsf{c}_{nb}(n+1)}\cdot L^dK_a^{\circ}e^{2\mathsf{u}-\mathsf{u}n} .
\end{equation*}
Put $\varsigma:=L^{2d}e^{\mathsf{c}_{nb}-\mathsf{u}}$. The right-hand side equals
\begin{equation*}
2\cdot3^dL^{4d}e^{\mathsf{c}_{nb}+2\mathsf{u}}K_a^{\circ}\sum_{n\ge2}\varsigma^n .
\end{equation*}

By the definition of $\mathsf{u}$,
\begin{equation*}
\mathsf{u}-\mathsf{c}_{nb}-2d\ln L=2\mathsf{c}_{an}+\mathsf{c}_{nb}+4\ge4,
\end{equation*}
that is, $\varsigma\le e^{-4}$. Hence $\sum_{n\ge2}\varsigma^n=\varsigma^2/(1-\varsigma)\le2\varsigma^2$,
and the sum is at most
$4\cdot3^dL^{8d}e^{3\mathsf{c}_{nb}}K_a^{\circ}$. Since $\mathsf{u}\ge\mathsf{c}_{nb}+4$, we have
$2e^{3\mathsf{c}_{nb}}\le e^{4\mathsf{u}}$; and $L^{8d}\le L^{9d}$. The bound is therefore at most
$C_\beta K_a^{\circ}$.

No step of the argument counts arrests, nesting depths or scales.
\end{proof}

\begin{proposition}[Conversion of old polylogarithmic factors]\label{4.25:prop-conversion-old-factors}
Let $e\in\mathfrak{Q}(\mathfrak{a})$ be a piece presented at step $k'$, and write $n:=n(e)$. Let
$y\in\mathcal{Y}(e)$ be a current block meeting the enlarged support of $e$
(Definition~\ref{4.25:def-enlarged-support-blocks}). Let $\ell'(y)$ be the current level of $y$,
and put $t_y:=x_{\ell'(y)}$. Let $j_W$ be the level entering part~(a) of
Lemma~\ref{4.25:lem-pile-bound} (piece geometry and the pile bound). Let $\theta(y)$,
$\bar\theta$ and $\Theta(e)$ be the block thresholds and envelopes
of~\eqref{4.25:eq-flow-inequalities-f}; let $C_G$ be the constant
of~\eqref{4.25:eq-oscillation-datum-a} and $c_{wt}$ that of
Definition~\ref{4.25:not-regularity-class}. Then, with $\mathfrak{c}_G$ as defined in the last row,
\begin{equation}\label{4.25:eq-conversion-old-factors-a}
\begin{aligned}
&\text{(1)}\quad t_y\ge x_{k'}-2c_2,\\
&\text{(2)}\quad \theta(y)\le\bar\theta(t_y),\\
&\text{(3)}\quad x_{j_W}\le t_y+2c_2+B_\sharp(n+2),\\
&\text{(4)}\quad P_{\mathfrak{b}}^{\uparrow}(t)
  :=\sup_{n'\ge0}e^{-n'}P_{\mathfrak{b}}\bigl(t+2c_2+B_\sharp(n'+2)\bigr)<\infty,\\
&\text{(5)}\quad \Pi_G(x):=\sup_{t\ge x-2c_2}P_{\mathfrak{b}}^{\uparrow}(t)\,\bar\theta(t)\\
&\phantom{\text{(5)}}\quad\text{is non-increasing in $x$ and tends to $0$ as $x\to\infty$},\\
&\phantom{\text{(5)}}\quad \mathfrak{c}_G:=C_G\,e^{3+2c_{wt}} .
\end{aligned}
\end{equation}
In~(4) the function $P_{\mathfrak{b}}^{\uparrow}$ is finite, and it is bounded by a
polylogarithmic function of $\max\{t,1\}$. Consequently
\begin{equation}\label{4.25:eq-conversion-old-factors-1}
\varrho^2\,\Theta(e)\,e^{-n(e)}\le\Pi_G(x_{k'}).
\end{equation}
The right-hand side is a function of $x_{k'}=g_{k'}^{-2}$, the inverse squared coupling of the step
at which $e$ is presented, and of nothing else.

Moreover, let $\mathsf{v}_1,\mathsf{v}_2,\mathsf{v}_3$ be real numbers satisfying
\begin{itemize}
\item[(a)] $\mathsf{v}_1\le c_J|S^+(e)|$, where $c_J$ is the constant entering $c_{wt}$ in
Definition~\ref{4.25:not-regularity-class};
\item[(b)] $\mathsf{v}_2\le(2\kappa+3)\sqrt d\,|S^+(e)|$;
\item[(c)] $\mathsf{v}_3\le\tfrac13+\tfrac32\sqrt d\,|S^+(e)|$.
\end{itemize}
Then for every later operation $\mathfrak{O}$ the oscillation
datum $\mathfrak{N}^G(e;\mathfrak{O})$ of~\eqref{4.25:eq-oscillation-datum-a} satisfies
\begin{equation}\label{4.25:eq-conversion-old-factors-2}
\begin{split}
\mathfrak{N}^G(e;\mathfrak{O})\,e^{\mathsf{v}_1+\mathsf{v}_2+\mathsf{v}_3}
&\le\mathfrak{c}_G\,\mathfrak{n}^{\circ}(e)\,\varrho^2\,\Theta(e)\,e^{-n(e)}\\
&\le\mathfrak{c}_G\,\mathfrak{n}^{\circ}(e)\,\Pi_G(x_{k'}) .
\end{split}
\end{equation}
\end{proposition}

The weights that the estimates using the oscillation datum attach to $e$ are of the form
considered in Proposition~\ref{4.25:prop-conversion-old-factors}. The first is $e^{\mathsf{v}_1}$
with $\mathsf{v}_1$ as in~(a). The second is the weight $e^{\mathsf{v}_2}$ attached to $e$ through
the blocking of $S^+(e)$ in the lemma on blockings of enlarged supports; its exponent is the
product of a factor at most $2\kappa+3$ with a non-negative length of that blocking, and by
part~(ii) of that lemma this length is at most $\sqrt d\,|S^+(e)|$; hence~(b). The third is the
exponential weight that a further estimate of this chapter attaches to $e$; its exponent is at
most $\tfrac13+\tfrac32\sqrt d\,|S^+(e)|$, which is~(c).

\begin{proof}
By Lemma~\ref{4.25:lem-domination-old-blocks}\,(iii) (domination of old blocks by current
blocks), the statements below may be read with the blocks of
$\mathcal{Y}(e)$ wherever blocks lying in $S^+(e)$ are meant. Under this reading no constant
changes. By part~(i) of the same lemma, the current cells of the blocks of $\mathcal{Y}(e)$ are
exactly the hats of the cells of $S^+(e)$. Hence $\ell'(y)$ is a current level on $S^+(e)$.

Two properties of the polylogarithmic function $P_{\mathfrak{b}}$, fixed with the box
$\mathfrak{b}_x$ of the collar variable and its radius $\varrho$, are used below:
$P_{\mathfrak{b}}$ is non-decreasing, and $\varrho^2\le P_{\mathfrak{b}}(x_{j_W})$.

\emph{Item~(1).} The current level of a block at step $k'$ satisfies $\ell'(y)\le k'$. Inequality
\eqref{4.25:eq-flow-inequalities-a} of the coupling-flow inequalities
(Definition~\ref{4.25:not-flow-inequalities}) states $x_i\ge x_{i'}-2c_2$ for $i\le i'$. Applied
with $i=\ell'(y)$ and $i'=k'$, it gives
\begin{equation*}
t_y=x_{\ell'(y)}\ge x_{k'}-2c_2 .
\end{equation*}

\emph{Item~(2).} By the definitions~\eqref{4.25:eq-flow-inequalities-f}, the envelope
$\bar\theta(t)$ dominates $\theta(y)$ at $t=x_{\ell'(y)}=t_y$, as recorded in
Definition~\ref{4.25:not-flow-inequalities}. This is~(2).

\emph{Item~(3).} There are two cases.

Suppose first $\ell'(y)\le j_W$. Inequality~\eqref{4.25:eq-flow-inequalities-a} with $i=\ell'(y)$
and $i'=j_W$ gives $x_{\ell'(y)}\ge x_{j_W}-2c_2$, that is
\begin{equation*}
x_{j_W}\le t_y+2c_2\le t_y+2c_2+B_\sharp(n+2).
\end{equation*}

Suppose next $\ell'(y)>j_W$. Inequality~\eqref{4.25:eq-flow-inequalities-b} states
$x_i\le x_{i'}+B_\sharp(i'-i)$ for $i\le i'$. Applied with $i=j_W$ and $i'=\ell'(y)$, it gives
\begin{equation*}
x_{j_W}\le t_y+B_\sharp\bigl(\ell'(y)-j_W\bigr).
\end{equation*}
Since $\ell'(y)$ is a current level on $S^+(e)$, part~(a) of
Lemma~\ref{4.25:lem-pile-bound}, displayed in~\eqref{4.25:eq-pile-bound-a}, gives
$\ell'(y)\le j_W+n(e)+2$. Hence $\ell'(y)-j_W\le n+2$. Therefore, since $c_2\ge0$,
\begin{equation*}
x_{j_W}\le t_y+B_\sharp(n+2)\le t_y+2c_2+B_\sharp(n+2).
\end{equation*}

\emph{Item~(4).} Since $P_{\mathfrak{b}}$ is non-decreasing, for $t<1$ each member of the
supremum in~(4) is at most the corresponding member at $t=1$; so it suffices to treat $t\ge1$.
For $a\ge1$ and $b\ge0$ we have $a+b\le a(1+b)$, and therefore
\begin{equation*}
\log(a+b)\le\log a+\log(1+b).
\end{equation*}
Apply this with $a=t\ge1$ and $b=2c_2+B_\sharp(n'+2)$. Every logarithm of the argument entering
the polylogarithmic function $P_{\mathfrak{b}}$ is then bounded by
$\log t+\log(1+2c_2+B_\sharp(n'+2))$. Hence
$P_{\mathfrak{b}}(t+2c_2+B_\sharp(n'+2))$ is bounded by a polynomial, with coefficients independent
of $n'$, in the two quantities $\log t$ and $\log(1+2c_2+B_\sharp(n'+2))$. Each power of the
second quantity grows at most like a power of $\log n'$, so its product with $e^{-n'}$ is bounded
uniformly in $n'\ge0$. Consequently the supremum in~(4) is finite, and for $t\ge1$ it is bounded by
a polynomial in $\log t$; by the reduction above, for every $t$ it is bounded by the same
polynomial in $\log\max\{t,1\}$.

\emph{Item~(5).} As $x$ increases, the set $\{t\ge x-2c_2\}$ shrinks. Hence $\Pi_G$ is
non-increasing. By Definition~\ref{4.25:not-flow-inequalities},
$\bar\theta(t)=5P_w(t+2c_2)\bar\alpha(t)$. There $P_w$ is an increasing polylogarithmic function,
and $\bar\alpha(t)$ is $t^{-1/2}$ times a polylogarithmic factor. By the logarithmic inequality
used in~(4), $P_w(t+2c_2)$ is bounded by a polylogarithmic function of $t$. By~(4),
$P_{\mathfrak{b}}^{\uparrow}(t)$ is also bounded by a polylogarithmic function of $t$. Thus
$P_{\mathfrak{b}}^{\uparrow}(t)\bar\theta(t)$ is at most $t^{-1/2}$ times a polylogarithmic
function of $t$, and this tends to $0$ as $t\to\infty$. Therefore $\Pi_G$, which is
non-increasing, is finite for all sufficiently large $x$ and tends to $0$ as $x\to\infty$.

\emph{The bound~\eqref{4.25:eq-conversion-old-factors-1}.} The polylogarithmic function
$P_{\mathfrak{b}}$ enters here through the two properties recalled at the beginning of the proof:
it is non-decreasing, and $\varrho^2\le P_{\mathfrak{b}}(x_{j_W})$. By~(3) and these two properties,
\begin{equation*}
\varrho^2e^{-n}\le e^{-n}P_{\mathfrak{b}}\bigl(t_y+2c_2+B_\sharp(n+2)\bigr)
\le P_{\mathfrak{b}}^{\uparrow}(t_y).
\end{equation*}
The last inequality holds because the middle term is the member $n'=n$ of the supremum in~(4).
Multiply by $\theta(y)\ge0$ and use~(2). Then use $t_y\ge x_{k'}-2c_2$, which is~(1), and the
definition~(5). This gives
\begin{equation*}
\varrho^2e^{-n}\theta(y)\le P_{\mathfrak{b}}^{\uparrow}(t_y)\,\bar\theta(t_y)\le\Pi_G(x_{k'}).
\end{equation*}
The right-hand side does not depend on $y$. Taking the supremum over $y\in\mathcal{Y}(e)$ and
recalling $\Theta(e)=\sup_{y\in\mathcal{Y}(e)}\theta(y)$ from~\eqref{4.25:eq-flow-inequalities-f}
proves~\eqref{4.25:eq-conversion-old-factors-1}.

\emph{The bound~\eqref{4.25:eq-conversion-old-factors-2}.} Adding the three bounds (a)--(c) gives
\begin{equation*}
\mathsf{v}_1+\mathsf{v}_2+\mathsf{v}_3
\le\tfrac13+\bigl[c_J+(2\kappa+\tfrac92)\sqrt d\,\bigr]|S^+(e)|.
\end{equation*}
Write $Y'$ for the region of the piece $e$ that enters the definitions of $\nu^{\circ}(e)$ and
$\mathfrak{n}^{\circ}(e)$ in Definition~\ref{4.25:not-regularity-class}. By the construction of
the pieces, the size of the piece is $n(e)=|Y'|+2$, and its enlarged support satisfies
$|S^+(e)|\le\mathsf{n}_2\,n(e)=\mathsf{n}_2(|Y'|+2)$. With the constant
$c_{wt}=\mathsf{n}_2[c_J+(2\kappa+5)\sqrt d\,]$ of Definition~\ref{4.25:not-regularity-class},
these give
\begin{equation*}
e^{\mathsf{v}_1+\mathsf{v}_2+\mathsf{v}_3}\le e^{1/3+2c_{wt}}\,e^{c_{wt}|Y'|}.
\end{equation*}
By Corollary~\ref{4.25:cor-oscillation-datum}, the datum is
$\mathfrak{N}^G(e;\mathfrak{O})=C_G\nu^{\natural}(e)\Theta(e)$. By
Definition~\ref{4.25:not-regularity-class},
\begin{equation*}
\nu^{\natural}(e)=\varrho^2\nu^{\circ}(e)e^{-c_*|Y'|},
\qquad c_*=c_{wt}+\mathsf{u}+2 .
\end{equation*}
Hence
\begin{align*}
\mathfrak{N}^G(e;\mathfrak{O})\,e^{\mathsf{v}_1+\mathsf{v}_2+\mathsf{v}_3}
&\le C_G\,e^{1/3+2c_{wt}}\,\nu^{\circ}(e)e^{-\mathsf{u}|Y'|}\,
  e^{-2|Y'|}\varrho^2\Theta(e)\\
&\le C_G\,e^{1/3+2c_{wt}}\,\nu^{\circ}(e)e^{-\mathsf{u}|Y'|}\,
  e^{-|Y'|}\varrho^2\Theta(e),
\end{align*}
where one factor $e^{-|Y'|}\le1$ has been discarded. Since $n(e)=|Y'|+2$ by the first counting
fact,
\begin{equation*}
e^{-|Y'|}\varrho^2\Theta(e)=e^{2}e^{-n(e)}\varrho^2\Theta(e).
\end{equation*}
By the
definitions of $\mathfrak{n}^{\circ\circ}$ and $\mathfrak{n}^{\circ}$ in
Definition~\ref{4.25:not-regularity-class},
\begin{equation*}
2\nu^{\circ}(e)e^{-\mathsf{u}|Y'|}\le\mathfrak{n}^{\circ}(e).
\end{equation*}
Combining these,
\begin{align*}
\mathfrak{N}^G(e;\mathfrak{O})\,e^{\mathsf{v}_1+\mathsf{v}_2+\mathsf{v}_3}
&\le\tfrac12\,C_G\,e^{7/3+2c_{wt}}\,\mathfrak{n}^{\circ}(e)\,\varrho^2\,\Theta(e)\,e^{-n(e)}\\
&\le C_G\,e^{3+2c_{wt}}\,\mathfrak{n}^{\circ}(e)\,\varrho^2\,\Theta(e)\,e^{-n(e)}.
\end{align*}
The right-hand side is $\mathfrak{c}_G\mathfrak{n}^{\circ}(e)\varrho^2\Theta(e)e^{-n(e)}$ by the
last line of~\eqref{4.25:eq-conversion-old-factors-a}; this is the first inequality
of~\eqref{4.25:eq-conversion-old-factors-2}. Since $\mathfrak{c}_G$ and $\mathfrak{n}^{\circ}(e)$
are non-negative, \eqref{4.25:eq-conversion-old-factors-1} bounds this right-hand side by
$\mathfrak{c}_G\mathfrak{n}^{\circ}(e)\Pi_G(x_{k'})$, which is the second inequality. This
proves~\eqref{4.25:eq-conversion-old-factors-2}.
\end{proof}

\begin{definition}[Cells meeting an enlarged response cube]\label{4.25:def-response-cube-cells}
Let $\mathfrak{a}$ be an arrest, and consider a presentation of it whose current cells form the
family $\mathfrak{F}$. Let $\square_0$ be a response cube of the M\"obius family of the stage, and
let $\tilde\square_0$ be its enlargement. The \emph{set of current cells meeting
$\tilde\square_0$} is
\begin{equation}\label{4.25:eq-response-cube-cells-a}
\mathfrak{A}(\tilde\square_0)
:=\{\,c\in\mathfrak{F}:\ c\cap\tilde\square_0\neq\emptyset\,\}.
\end{equation}
\end{definition}

The set \eqref{4.25:eq-response-cube-cells-a} serves to bound the anchored sum over the pieces
\begin{equation*}
c_e=(v_e,S^+(e),\emptyset)\qquad\text{with}\qquad S^+(e)\cap\tilde\square_0\neq\emptyset .
\end{equation*}
Here $e$ ranges over the pieces of the families $\mathfrak{Q}(\mathfrak{a}')$ of the live arrests
$\mathfrak{a}'$ of the same type as $\mathfrak{a}$, and $S^+(e)$ is the $2$-neighbourhood of the
support of $e$.

This sum is bounded by clause~$(\alpha)$ of the pile bound in
Lemma~\ref{4.25:lem-pile-bound}, display~\eqref{4.25:eq-pile-bound-a}, taken with
$\mathfrak{A}=\mathfrak{A}(\tilde\square_0)$:
\begin{equation}\label{4.25:eq-response-cube-cells-1}
\sum_{\substack{\mathfrak{a}'\ \mathrm{live}\\ \text{of the type of }\mathfrak{a}}}\
\sum_{\substack{e\in\mathfrak{Q}(\mathfrak{a}')\\
S^+(e)\cap\bigcup\mathfrak{A}(\tilde\square_0)\neq\emptyset}}
\mathfrak{n}^{\circ}(e)
\;\le\; C_\alpha\,\bigl|\mathfrak{A}(\tilde\square_0)\bigr|\,K_a^{\circ},
\end{equation}
where
\begin{equation*}
C_\alpha=2L^{d}e^{4\mathsf{u}}.
\end{equation*}
The cardinality $|\mathfrak{A}(\tilde\square_0)|$ enters this estimate only once, and linearly.

The set is used where the plain line~$(\gamma)$ of the response bound is read at a presentation of
$\mathfrak{a}$, namely at the sibling stages at step $k_W$ and at the stages of a later attempt
containing $\mathfrak{a}$.

\begin{lemma}[Geometry of a response cube]\label{4.25:lem-response-cube-geometry}
Let $\mathfrak{F}$ be the current cells of a presentation, as in
Definition~\ref{4.25:def-response-cube-cells}, with cells and the lengths
$\ell_m$ as in Definition~\ref{4.25:not-collar-weight}. Let $i_0$ be a level and
let $\pi_{i_0}$ be the partition into closed cubes of side $M\ell_{i_0}$ on the grid
of level $i_0$. We use the layer of level $i_0$ of the constructions in
\cite{B-CMP109} and \cite{B-CMP119}; every point of this layer lies in a current cell
of level $i_0$, and that cell is a cube of $\pi_{i_0}$. For $x\in\mathbb{R}^d$ and a
non-empty set $S$ write
\begin{equation*}
\operatorname{dist}_\infty(x,S)=\inf_{z\in S}|x-z|_\infty .
\end{equation*}
Let $y$ be a vertex of the grid of $\pi_{i_0}$ lying on the layer of level $i_0$. Let
$\square_0$ be the union of the $2^d$ closed cubes of $\pi_{i_0}$ having $y$ as a
vertex \cite{B-CMP109}, \cite{B-CMP119}, and let $\tilde\square_0$ be the union of
$\square_0$ with all cubes of $\pi_{i_0}$ that touch $\square_0$ \cite{B-CMP119}. Let
$\zeta_{\square_0}$ be the member attached to $\square_0$ of the partition of unity
$\{\zeta_{\square}\}$ of \cite{B-CMP119}; by its construction there, its support lies in
$\{x:|x-y|_\infty\le\tfrac23 M\ell_{i_0}\}$. Assume that $\zeta_{\square_0}$ is not
identically zero on the layer of level $i_0$. Then:
\begin{enumerate}
\item[(a)] $\square_0$ is a closed cube of side $2M\ell_{i_0}$ with centre $y$.
\item[(b)] One of the $2^d$ cubes of $\pi_{i_0}$ at $y$ is a current cell: there is
$c\in\mathfrak{F}$ of level $i_0$ with $c\subset\square_0$.
\item[(c)] With
\begin{equation*}
\begin{gathered}
Q(\square_0):=\{x:\operatorname{dist}_\infty(x,\square_0)\le M\ell_{i_0}\},\\
Q(c):=\{x:\operatorname{dist}_\infty(x,c)\le 2M\ell_{i_0}\},
\end{gathered}
\end{equation*}
we have
\begin{equation}\label{4.25:eq-response-cube-geometry-a}
\tilde\square_0\subset Q(\square_0)\subset Q(c).
\end{equation}
\item[(d)] $Q(\square_0)$ is a closed cube of side exactly $4M\ell_{i_0}$. In
particular, two closed sets which both meet $Q(\square_0)$ are at
$|\cdot|_\infty$-distance at most $4M\ell_{i_0}$ from each other.
\item[(e)] For the set $\mathcal{A}(\tilde\square_0)$ of
\eqref{4.25:eq-response-cube-cells-a},
\begin{equation}\label{4.25:eq-response-cube-geometry-1}
\begin{aligned}
\mathcal{A}(\tilde\square_0)
&\subset\{c'\in\mathfrak{F}: c'\cap Q(\square_0)\neq\emptyset\}\\
&\subset\{c'\in\mathfrak{F}: c'\cap Q(c)\neq\emptyset\},
\end{aligned}
\end{equation}
and $c\in\mathcal{A}(\tilde\square_0)$.
\end{enumerate}
\end{lemma}

\begin{proof}
Write $\ell=\ell_{i_0}$ and $y=(y_1,\dots,y_d)$. The vertices of the cubes of
$\pi_{i_0}$ are the points with every coordinate $x_\mu$ in $y_\mu+M\ell\mathbb{Z}$,
since $y$ is such a vertex. A cube of $\pi_{i_0}$
has $y$ as a vertex exactly when, for each $\mu$, its $\mu$-th side is
$[y_\mu-M\ell,y_\mu]$ or $[y_\mu,y_\mu+M\ell]$.

(a) The union of the $2^d$ products of such intervals is
$\prod_\mu[y_\mu-M\ell,y_\mu+M\ell]$, a closed cube of side $2M\ell$ and centre $y$.

(b) Since $\zeta_{\square_0}$ is not identically zero on the layer, there is a point
$x$ of the layer with $\zeta_{\square_0}(x)\neq0$. By the support property,
\begin{equation*}
|x-y|_\infty\le\tfrac23 M\ell<M\ell .
\end{equation*}
As $x$ lies on the layer, some current cell $c$ of level $i_0$ contains $x$, and
$c$ is a cube of $\pi_{i_0}$. Fix $\mu$ and let $[a,a+M\ell]$, with
$a\in y_\mu+M\ell\mathbb{Z}$, be the $\mu$-th side of $c$. From
$a\le x_\mu<y_\mu+M\ell$ we get $a\le y_\mu$. From
$a+M\ell\ge x_\mu>y_\mu-M\ell$ we get $a>y_\mu-2M\ell$. Hence $a=y_\mu-M\ell$ or
$a=y_\mu$, and in both cases $y_\mu$ is an endpoint of the side. Thus $y$ is a vertex
of $c$, so $c$ is one of the $2^d$ cubes of $\pi_{i_0}$ at $y$. By (a), $c\subset\square_0$.

(c) First inclusion. A point of $\tilde\square_0$ lies either in $\square_0$, where
its distance to $\square_0$ is $0$, or in a cube of $\pi_{i_0}$ touching
$\square_0$. In the second case pick a point $x''$ common to this cube and
$\square_0$. Every point $x'$ of this cube satisfies $|x'-x''|_\infty\le M\ell$,
because the cube has side $M\ell$. Hence
$\operatorname{dist}_\infty(x',\square_0)\le M\ell$. Therefore $\tilde\square_0\subset Q(\square_0)$.

By \cite{B-CMP119}, a point outside $\tilde\square_0$ lies at
$|\cdot|_\infty$-distance at least $M\ell$ from $\square_0$; this is the inclusion
$\{x:\operatorname{dist}_\infty(x,\square_0)<M\ell\}\subset\tilde\square_0$, the reverse of
the first inclusion of (c), and the first inclusion is the one used here.

Second inclusion. Let $u\in\square_0$ and fix $\mu$. The $\mu$-th side of $c$ is one
of the two halves of $[y_\mu-M\ell,y_\mu+M\ell]$, and $u_\mu$ lies in this interval of
length $2M\ell$. Hence the distance from $u_\mu$ to the $\mu$-th side of $c$ is at
most $M\ell$. Since $c$ is a product of intervals, $\operatorname{dist}_\infty(u,c)$ is
the maximum over $\mu$ of these coordinate distances, so
\begin{equation*}
\operatorname{dist}_\infty(u,c)\le M\ell\qquad (u\in\square_0).
\end{equation*}
Now let $x\in Q(\square_0)$. Since $\square_0$ is compact, there is $u\in\square_0$
with $|x-u|_\infty=\operatorname{dist}_\infty(x,\square_0)\le M\ell$. Then
\begin{equation*}
\operatorname{dist}_\infty(x,c)\le|x-u|_\infty+\operatorname{dist}_\infty(u,c)
\le 2M\ell .
\end{equation*}
So $x\in Q(c)$, which proves \eqref{4.25:eq-response-cube-geometry-a}.

(d) For a product of intervals, $\operatorname{dist}_\infty(x,\cdot)$ is the maximum
over $\mu$ of the distances from $x_\mu$ to the $\mu$-th interval. By (a), therefore,
\begin{equation*}
Q(\square_0)=\prod_\mu[y_\mu-2M\ell,\,y_\mu+2M\ell],
\end{equation*}
a closed cube of side $4M\ell$. If two closed sets meet it, at points $u$ and $v$
say, then $|u-v|_\infty\le4M\ell$, which bounds their distance.

(e) Let $c'\in\mathcal{A}(\tilde\square_0)$, so that $c'\cap\tilde\square_0\neq\emptyset$
by \eqref{4.25:eq-response-cube-cells-a}. By \eqref{4.25:eq-response-cube-geometry-a},
$c'$ meets $Q(\square_0)$, and then it also meets $Q(c)$. This gives
\eqref{4.25:eq-response-cube-geometry-1}. Finally, $c\in\mathfrak{F}$, and by (b) $c$ is
non-empty and contained in $\square_0$, which is contained in $\tilde\square_0$; so
$c\cap\tilde\square_0\neq\emptyset$, that is, $c\in\mathcal{A}(\tilde\square_0)$.
\end{proof}

\begin{definition}[Raised cells, attempts and window cells]\label{4.25:def-window-cells}
Throughout, $\ell_m$ is the length scale of level $m$ and cells are the closed $M$-cubes of
the conventions of~\ref{4.25:not-collar-weight}, and $\operatorname{dist}_\infty$ is the distance
induced by the maximum norm. The regions $\Omega_m$ and $\Lambda_k$ are those of
Ba{\l}aban's large-field construction: $\Lambda_k$ is the small-field region of step $k$, so that the
components of its complement $\Lambda^c_k$ are the large-field components of step $k$; and
$\Omega_m$ is the domain of level $m$ of that construction.

\emph{Attempts.} An \emph{$\mathbf{R}$-attempt} $w$ carries the following data:
\begin{itemize}
\item its \emph{class} $k_w$, which is its step;
\item its \emph{component} $Z_w$, a rectangular parallelepiped that is a connected component of
$\Lambda^c_{k_w}$;
\item the two levels
\begin{equation*}
h_w := k_w - N, \qquad k_0^w := k_w - N_0^{(k_w)},
\end{equation*}
where $N$ is the integer fixed in Ba{\l}aban's large-field step \cite{B-CMP122a} (in this
definition $N$ does not denote the rank of the gauge group), and $N_0^{(k_w)}$ is the number $N_0$ of
\cite[(1.10)]{B-CMP122a} taken at step $k_w$. Since $N > N_0^{(k_w)}$, we have $h_w \le k_0^w - 1$.
\end{itemize}
An attempt is \emph{present} at the current step if it is \emph{running}, that is, it is the
presenting attempt or a sibling of it at the current step, or if it is \emph{live-arrested},
that is, frozen. The enlarged component of $w$ is
\begin{equation*}
V_w := \{x : \operatorname{dist}_\infty(x, Z_w) \le 7M\ell_{k_w}\}.
\end{equation*}

\emph{Current cells and their levels.} The set of cells of the current determining family is
denoted $\mathfrak{F}$. These are closed $M$-cubes, of side $M\ell_m$ at level $m$, lying on nested
$L$-adic grids. The current determining family is the following, where $k_{\square}$ denotes the
current step and $n$ the current stage:
\begin{itemize}
\item on the component $Z$ of a running attempt, the family $\mathbb{B}^{(n)}_{k_{\square}}$
restricted to $Z$, as in \cite[(1.16)]{B-CMP122a};
\item on a live arrested piece (one still carrying its large-field factor at the current step),
the family $\mathbb{B}^{(n^+)}_{k_a}$ frozen at the arrest, where
$k_a$ is the step of the arrest and $n^+$ the stage index of the frozen family (the frozen
determining family of the arrest, cf.\ the conventions of~\ref{4.25:not-collar-weight});
\item on a component $Z_c$ healed at step $k_c$ (its large-field factor absorbed back into the
small-field format at step $k_c$), the one-level family $\{Z_c^{(k_c)}\}$, as in
\cite[(1.33)]{B-CMP122b};
\item elsewhere, the family $\mathbb{B}_{k_{\square}}$.
\end{itemize}
For $c \in \mathfrak{F}$, $\operatorname{lev} c$ denotes its level. The current cells satisfy the
level-gap property \eqref{4.25:eq-collar-weight-a} and the property that levels never decrease and
the partitions are nested, both stated in the conventions of~\ref{4.25:not-collar-weight}.

\emph{Raised and un-raised cells.} For a point $x$ define
\begin{equation*}
\begin{split}
m_{\mathbf{T}}(x) := \max\bigl(
 &\{m : x \in \Omega_m,\ \text{step } m \text{ not later than the current step}\}\\
 &\cup \{k_c : x \in Z_c,\ Z_c \text{ healed at a step } k_c
 \text{ not later than the current step}\}\bigr),
\end{split}
\end{equation*}
and, for each level $m$, set
\begin{equation*}
\mathfrak{O}_m := \{x : m_{\mathbf{T}}(x) \ge m\}.
\end{equation*}
A cell $c \in \mathfrak{F}$ is \emph{raised} if $\operatorname{lev} c > m_{\mathbf{T}}(x)$ for every
$x$ in the interior of $c$. Such a cell carries a level assigned at a raising stage of some step and
not overwritten since. A cell is \emph{un-raised} otherwise.
A raised cell is \emph{raised by} an attempt $w$ if its current level was assigned at a raising
stage of $w$.

\emph{Cube sides.} For each level $m$ let $R_m$ be the large-cube factor of Ba{\l}aban's
large-field step at level $m$ (not to be confused with the ball radius $R$), so that the
large cubes of level $m$ are the cubes of side $s_m := M R_m \ell_m$. These sides satisfy
\begin{equation*}
s_m \ge M\ell_m, \qquad s_{m+1} \ge s_m,
\end{equation*}
and, for every attempt $w$,
\begin{equation*}
s_{k_0^w+1} = M\ell_{k_w}.
\end{equation*}

\emph{Window cells.} A \emph{window cell} of an attempt $w$ is a cell raised by $w$ whose level
exceeds $k_0^w$. By clause (W4) of the window lemma (item (c) below), these are
the cells of the raised top zone $T_j$ of $w$, of level $j > k_0^w$, and of the thin strata
$\Gamma''_m$, $k_0^w < m < j$, defined there. The set of
window cells of $w$ is denoted $\mathfrak{W}_w$.

\emph{Collars of cells.} For a cell $c$, its collar is
\begin{equation*}
Q(c) := \{x : \operatorname{dist}_\infty(x, c) \le 2M\ell_{\operatorname{lev} c}\}.
\end{equation*}
\end{definition}

The following four statements are used below together with the objects just
defined.
\begin{enumerate}
\item[(a)] The lemma on enlargements of unions of closed grid cubes: let $X$ be a union of closed
grid cubes of side $s$, and write $X^{\sim n}$ for its $n$-fold enlargement by grid cubes of
side $s$ in the touching relation of the conventions of~\ref{4.25:not-collar-weight} (two closed
cubes touch if they meet, contact at corners included), that is, $X^{\sim 0} := X$ and
$X^{\sim (j+1)}$ is the union of $X^{\sim j}$ with all grid cubes of side $s$ touching it. Then
\begin{equation*}
X^{\sim n} = \{x : \operatorname{dist}_\infty(x, X) \le ns\}.
\end{equation*}
\item[(b)] The separation lemma for cells of low level: a cell $c'$ with
$\operatorname{lev} c' \le m-2$ satisfies
\begin{align*}
&\text{(i)}\quad \operatorname{dist}_\infty(c', \Omega_m) > 8s_m;\\
&\text{(ii)}\quad \operatorname{dist}_\infty(c', Z_c) > 10s_m
 \quad \text{for a component } Z_c \text{ healed at step } m;\\
&\text{(iii)}\quad \operatorname{dist}_\infty(c', \mathfrak{O}_m) > 8s_m.
\end{align*}
\item[(c)] The window lemma: for an attempt $w$, with the sets $\Omega^{\mathrm{term}}_m$ defined
there,
\begin{align*}
&\text{(W0)}\quad \Omega^{\mathrm{term}}_m \cap V_w \subset \mathfrak{O}_m
 \quad \text{for } h_w < m \le k_w;\\
&\text{(W1)}\quad \text{a cell of } Z_w \text{ of level} > k_0^w
 \text{ lies in } \bigl(\Omega^{\mathrm{term}}_{k_0^w+1}\bigr)^{\sim 7};\\
&\text{(W2)}\quad \text{a cell of } Z_w \text{ of level } k_w
 \text{ lies in } \bigl(\Omega^{\mathrm{term}}_{k_0^w+1}\bigr)^{\sim 5}
 \cup \Omega^{\mathrm{term}}_{k_w};\\
&\text{(W4)}\quad \text{the window cells of } w \text{ are the cells of the raised top zone }
 T_j \text{ of } w,\\
&\qquad\quad \text{of level } j > k_0^w, \text{ and of the thin strata } \Gamma''_m,
 \ k_0^w < m < j.
\end{align*}
The window lemma has one further clause, (W3).
\item[(d)] The attribution lemma for raised cells: every cell $c \in \mathfrak{F}$ is un-raised,
or is raised by exactly one present attempt $w$, and then $c \subset Z_w$ and
$\operatorname{lev} c \le k_w$.
\end{enumerate}

\begin{lemma}[The class entered at a response cube]\label{4.25:lem-entered-class}
Work at a fixed moment of the induction, with the current cells $\mathfrak{F}$, their levels
$\operatorname{lev}$, the current step $k_\square$, the present attempts $w$ with their classes
$k_w$, components $Z_w$, numbers $h_w=k_w-N$ and $k_0^w=k_w-N_0^{(k_w)}$, collars $V_w$, window
cells $\mathfrak{W}_w$, the function $m_{\mathbf{T}}$, the sets
$\mathfrak{O}_m=\{m_{\mathbf{T}}\ge m\}$, the regions $\Omega_m$, $\Omega^{\mathrm{term}}_m$ and
$\Lambda_m$, and the $R$-cubes of level $m$, of side $s_m$, all as in
Definition~\ref{4.25:def-window-cells}. For a union $X$ of closed grid cubes of one level and an
integer $n\ge 1$, $X^{\sim n}$ denotes $X$ enlarged by $n$ layers of cubes of that grid, cubes
touching at corners included, and a superscript $\mathrm{c}$ denotes the complement. By
\cite[(3.20)]{B-CMP119}, the small-field region $\Lambda_m$ is the complement of the union of
$(\Omega^{\mathrm{c}}_m)^{\sim2}$ with a further region; in particular
\begin{equation}\label{4.25:eq-entered-class-2}
\Omega^{\mathrm{c}}_m\subset\bigl(\Omega^{\mathrm{c}}_m\bigr)^{\sim2}\subset\Lambda^{\mathrm{c}}_m,
\qquad\text{so that}\qquad \Lambda_m\subset\Omega_m .
\end{equation}
Assume the
following.
\begin{itemize}
\item[(i)] The theorem bounding the levels of the current cells that meet the collar $Q(c)$ of a
current cell $c$, with its clauses (U), (D1), (D2) and (K). In particular, by (K), every
$c'\in\mathfrak{F}$ satisfies $\operatorname{lev}c'\le\kappa(c')$, where $\kappa(c'):=k_w$ if
$c'$ is a window cell of some present attempt $w$, and $\kappa(c'):=k_\square$ otherwise.
\item[(ii)] The raising lemma: every cell $c$ of $\mathfrak{F}$ is un-raised, or raised by exactly
one present attempt $w$, and then $c\subset Z_w$ and $\operatorname{lev}c\le k_w$.
\item[(iii)] The collar lemma for grid cubes: if $X$ is a union of closed grid cubes of side $s$,
cubes touching at corners included, then
$X^{\sim n}=\{\operatorname{dist}_\infty(\cdot,X)\le ns\}$.
\item[(iv)] Clause (W0) of the window lemma: for every present $w$ and every $m$ with
$h_w<m\le k_w$,
\begin{equation*}
\Omega^{\mathrm{term}}_m\cap V_w\subset\mathfrak{O}_m .
\end{equation*}
\item[(v)] Clauses (T1) and (T2) of the arrest-nesting lemma, in the form used below. For a step
$k$, let $Z^{\mathrm{in}}_{k-1}$ denote the region of that lemma which contains the component of
every arrest live at step $k$, and let $Z'_{k-1}$ denote the region of that lemma whose tenfold
collar lies in the complement of the small-field region of step $k$:
\begin{equation*}
\bigl(Z'_{k-1}\bigr)^{\sim10}\subset\Lambda^{\mathrm{c}}_k .
\end{equation*}
Then: (T1) at every step $k$,
$Z^{\mathrm{in}}_{k-1}\subset Z'_{k-1}$; (T2) if a live arrest $w_1$ lies inside the
component of a present attempt $w_2$ with $k_{w_1}<k_{w_2}$, then $k_{w_1}\le h_{w_2}$.
\item[(vi)] $N>N_0^{(k_w)}$ for every present attempt $w$.
\end{itemize}
Let $\square_0$ be a response cube at level $i_0$, with $Q(\square_0)$ as in
\eqref{4.25:eq-response-cube-geometry-a}, and let $c\in\mathfrak{F}$ be any current cell with
$\operatorname{lev}c=i_0$ and $c\subset\square_0$ (one exists wherever
$\zeta_{\square_0}\not\equiv0$ on the layer, by Lemma~\ref{4.25:lem-response-cube-geometry}).
Define
\begin{equation}\label{4.25:eq-entered-class-a}
\begin{gathered}
M(\square_0):=\{\,w\ \text{present}:\ \text{some cell of }\mathfrak{W}_w
\text{ meets }Q(\square_0)\,\},\\
\kappa(\square_0):=
\begin{cases}
k_w\ \text{for } w\in M(\square_0), & M(\square_0)\ne\emptyset,\\
k_\square, & M(\square_0)=\emptyset ,
\end{cases}
\end{gathered}
\end{equation}
the class \emph{entered} at $\square_0$. In the definition of $M(\square_0)$ a window cell of $w$,
that is, a cell raised by $w$ of level above $k_0^w$, may be a thin stratum or a cell of the
raised top zone of $w$ (Definition~\ref{4.25:def-window-cells}); clauses (D2) and (K) of
hypothesis (i) treat both alike. Then the following hold.
\begin{itemize}
\item[(P1)] Every $c'\in\mathfrak{F}$ has $\operatorname{lev}c'\le k_\square$.
\item[(P2)] Exactly one of the following holds: $c$ is un-raised; or $c$ is raised by exactly one
present attempt $w_c$, with $c\subset Z_{w_c}$ and $\operatorname{lev}c\le k_{w_c}$, and either
$i_0\le k_0^{w_c}$ or $i_0>k_0^{w_c}$. Moreover, for every present $w$ and every
$c'\in\mathfrak{W}_w$ one has $k_0^w<\operatorname{lev}c'\le k_w$.
\item[(P3)] Two present attempts are nested or disjoint: if $w_1\ne w_2$ are present with
$k_{w_1}\le k_{w_2}$, then either $Z_{w_1}\subset Z_{w_2}$, in which case $k_{w_1}<k_{w_2}$,
$w_1$ is a live arrest inside the component of $w_2$, and
\begin{equation*}
k_{w_1}\le h_{w_2}\le k_0^{w_2}-1 ;
\end{equation*}
or $Z_{w_1}\cap Z_{w_2}=\emptyset$.
\item[(P4)] The collar of a raised cell: if $w$ is present and $c'\subset Z_w$ is a cell raised by
$w$, then
\begin{equation}\label{4.25:eq-entered-class-1}
\mathcal{N}(c'):=\{\operatorname{dist}_\infty(\cdot,c')\le 2s_{k_w}\}\subset Z_w .
\end{equation}
\item[(S)] If $A$ and $B$ are compact, $a,b\ge0$, and the sets
$\{\operatorname{dist}_\infty(\cdot,A)\le a\}$ and $\{\operatorname{dist}_\infty(\cdot,B)\le b\}$
are disjoint, then $\operatorname{dist}_\infty(A,B)>a+b$.
\end{itemize}
\end{lemma}

\begin{proof}
\emph{(P1).} Let $c'\in\mathfrak{F}$. Clause (K) of hypothesis (i), applied with $c'$ in the role
of the cell $c$ there, gives $\operatorname{lev}c'\le\kappa(c')$. Either $\kappa(c')=k_\square$,
or $\kappa(c')=k_w$ for a present attempt $w$. A present attempt is either running, and then its
class is the current step, or live-arrested, and then its class is the step at which it was
taken, which does not exceed the current one; in both cases $k_w\le k_\square$. Hence
$\operatorname{lev}c'\le\kappa(c')\le k_\square$.

\emph{(P2).} The raising lemma, hypothesis (ii), applied to $c$, states that $c$ is un-raised or
is raised by exactly one present attempt, which we call $w_c$, and that in the second case
$c\subset Z_{w_c}$ and $\operatorname{lev}c\le k_{w_c}$. In that case the alternative between
$i_0\le k_0^{w_c}$ and $i_0>k_0^{w_c}$ separates case (D1) from case (D2) of hypothesis (i); in
the second alternative $c$ is, by definition, a window cell of $w_c$. For the last assertion let
$w$ be present and $c'\in\mathfrak{W}_w$. By definition of a window cell, $c'$ is raised by $w$
and $\operatorname{lev}c'>k_0^w$. By hypothesis (ii) applied to $c'$, the attempt raising $c'$
is unique, hence equal to $w$, and $\operatorname{lev}c'\le k_w$.

\emph{(P3).} Let $w_1\ne w_2$ be present with $k_{w_1}\le k_{w_2}$.

Suppose first $k_{w_1}=k_{w_2}$. Then $w_1$ and $w_2$ are distinct attempts of one step. Their
components are distinct components of $\Lambda^{\mathrm{c}}_{k_{w_1}}$, and distinct components
of one set are disjoint; so $Z_{w_1}\cap Z_{w_2}=\emptyset$.

Suppose next $k_{w_1}<k_{w_2}$. Then $w_1$ is live at step $k_{w_2}$, so that
$Z_{w_1}\subset Z^{\mathrm{in}}_{k_{w_2}-1}$. Using clause (T1) of hypothesis (v), the trivial
inclusion of a set in its collar, and the inclusion
$(Z'_{k_{w_2}-1})^{\sim10}\subset\Lambda^{\mathrm{c}}_{k_{w_2}}$ of hypothesis (v), we obtain
\begin{equation*}
Z_{w_1}\subset Z^{\mathrm{in}}_{k_{w_2}-1}\subset Z'_{k_{w_2}-1}
\subset\bigl(Z'_{k_{w_2}-1}\bigr)^{\sim10}\subset\Lambda^{\mathrm{c}}_{k_{w_2}} .
\end{equation*}
Since $Z_{w_1}$ is connected (it is a rectangular parallelepiped), it lies inside one component
$C$ of $\Lambda^{\mathrm{c}}_{k_{w_2}}$. Since $Z_{w_2}$ is itself a component of
$\Lambda^{\mathrm{c}}_{k_{w_2}}$, either $C=Z_{w_2}$ or $C\cap Z_{w_2}=\emptyset$. In the second
case $Z_{w_1}\cap Z_{w_2}=\emptyset$. In the first case $Z_{w_1}\subset Z_{w_2}$ with
$k_{w_1}<k_{w_2}$; as $k_{w_1}$ is not the current step, $w_1$ is not running, so it is a live
arrest, and it lies inside the component of $w_2$. Clause (T2) of hypothesis (v) then gives
$k_{w_1}\le h_{w_2}$. Finally, by hypothesis (vi),
\begin{equation*}
h_{w_2}=k_{w_2}-N\le k_{w_2}-N_0^{(k_{w_2})}-1=k_0^{w_2}-1 .
\end{equation*}

\emph{(P4).} Let $w$ be present and $c'\subset Z_w$ a cell raised by $w$; by hypothesis (ii),
$\operatorname{lev}c'\le k_w$. We claim
$\operatorname{int}c'\not\subset\Omega^{\mathrm{term}}_{k_w}$. Otherwise, since
$c'\subset Z_w\subset V_w$ and $m=k_w$ satisfies $h_w<m\le k_w$, clause (W0) of hypothesis (iv)
gives $\operatorname{int}c'\subset\mathfrak{O}_{k_w}$, that is,
$m_{\mathbf{T}}\ge k_w\ge\operatorname{lev}c'$ on $\operatorname{int}c'$. Then
$\operatorname{lev}c'>m_{\mathbf{T}}$ fails on $\operatorname{int}c'$, so $c'$ is un-raised, against
the assumption. Hence $c'$ lies in an $R$-cube of level $k_w$ which is not contained in
$\Omega^{\mathrm{term}}_{k_w}$, that is, in $\Omega^{\mathrm{c}}_{k_w}$. The $R$-cubes of level
$k_w$ have side $s_{k_w}$, so by the collar lemma for grid cubes, hypothesis (iii) with $n=2$ and
$s=s_{k_w}$, and then by the inclusion \eqref{4.25:eq-entered-class-2} with $m=k_w$,
\begin{equation*}
\mathcal{N}(c')\subset\{\operatorname{dist}_\infty(\cdot,\Omega^{\mathrm{c}}_{k_w})\le 2s_{k_w}\}
=\bigl(\Omega^{\mathrm{c}}_{k_w}\bigr)^{\sim2}\subset\Lambda^{\mathrm{c}}_{k_w} .
\end{equation*}
The set $\mathcal{N}(c')$ is connected, and it meets $Z_w$, because it contains $c'\subset Z_w$.
Since $Z_w$ is a component of $\Lambda^{\mathrm{c}}_{k_w}$, it follows that
$\mathcal{N}(c')\subset Z_w$, which is \eqref{4.25:eq-entered-class-1}.

\emph{(S).} Suppose $\operatorname{dist}_\infty(A,B)\le a+b$. By compactness there are $x\in A$
and $y\in B$ with $|x-y|_\infty\le a+b$. If $a+b=0$, then $x=y$ lies in both sets. Otherwise put
$t=a/(a+b)\in[0,1]$ and $z=x+t(y-x)$. Then
\begin{equation*}
|z-x|_\infty=t\,|y-x|_\infty\le a,\qquad |z-y|_\infty=(1-t)\,|y-x|_\infty\le b,
\end{equation*}
so $z$ lies in both sets. In either case the two sets are not disjoint. Hence disjointness
forces $\operatorname{dist}_\infty(A,B)>a+b$.
\end{proof}

\begin{corollary}[Three cases at a response cube]\label{4.25:cor-three-cases}
Let $\square_0$ be a response cube at level $i_0$, and let $c\in\mathfrak{F}$ be the current cell with
$\operatorname{lev}c=i_0$ and $c\subset\square_0$ provided by the geometry of a response cube
(Lemma~\ref{4.25:lem-response-cube-geometry}). Thus
$\tilde\square_0\subset Q(\square_0)\subset Q(c)$, the second inclusion being
\eqref{4.25:eq-response-cube-geometry-a}. Let $M(\square_0)$ and $\kappa(\square_0)$ be as in
\eqref{4.25:eq-entered-class-a}. Assume the following.
\begin{itemize}
\item The level-spread statement for current cells, in the vocabulary of
Definition~\ref{4.25:def-window-cells}, whose clauses (U), (D1), (D2), (K) we recall. Let
$e\in\mathfrak{F}$, put
$i:=\operatorname{lev}e$ and $Q(e):=\{x:\operatorname{dist}_\infty(x,e)\le 2M\ell_{i}\}$, and let
$e'\in\mathfrak{F}$ meet $Q(e)$.
\begin{itemize}
\item[(U)] $\operatorname{lev}e'\le i+1$.
\item[(D1)] If $e$ is un-raised, or raised by a present attempt $w$ with $i\le k_0^w$, then
$\operatorname{lev}e'\ge i-1$.
\item[(D2)] If $e$ is raised by a present attempt $w$ with $i>k_0^w$, that is $e\in\mathfrak{W}_w$, then
$i\le k_w$ and
\begin{equation*}
\operatorname{lev}e'\ \ge\ k_0^w\ =\ k_w-N_0^{(k_w)}\ \ge\ i-N_0^{(k_w)}.
\end{equation*}
Here $N_0^{(k)}$ is the integer of Definition~\ref{4.25:def-window-cells} with
$k_0^w=k_w-N_0^{(k_w)}$; the window cells of $w$ are its raised cells of level strictly above $k_0^w$.
\item[(K)] Put $\kappa(e):=k_w$ in the situation of (D2) and $\kappa(e):=k_{\square}$, the current step,
otherwise. Then $i\le\kappa(e)$, and every level $\lambda$ of a cell meeting $Q(e)$ satisfies
\begin{equation*}
i-N_0^{(\kappa(e))}\ \le\ \lambda\ \le\ i+1\ \le\ \kappa(e)+1 .
\end{equation*}
\end{itemize}
\item The four preliminaries (P1)--(P4) of the class entered at a response cube
(Lemma~\ref{4.25:lem-entered-class}).
\end{itemize}
Say that $c$ is \emph{of type} (D1) if it is un-raised, or raised by a present attempt $w$ with
$i_0\le k_0^w$. Below, $\lambda$ denotes the level of an arbitrary cell of $\mathfrak{F}$ meeting $Q(c)$.
Then exactly one of the three cases (A), (B), (none) below holds, and in addition statement (C) holds; by
(C), every window cell $c$ falls under (A), with $M(\square_0)$ a singleton.
\begin{itemize}
\item[(A)] $c\in\mathfrak{W}_w$ for a present attempt $w$. Then $i_0\le k_w$; every cell meeting $Q(c)$, in
particular every cell of $\mathcal{A}(\tilde\square_0)$ (see \eqref{4.25:eq-response-cube-cells-a}), has
level in $[k_0^w,i_0+1]$; $M(\square_0)=\{w\}$; and $\kappa(\square_0)=k_w\ge i_0$.
\item[(B)] $c$ is of type (D1) and $M(\square_0)\neq\emptyset$. Then every cell meeting $Q(c)$ has level in
$[i_0-1,i_0+1]$. For every $w\in M(\square_0)$ and every $c_1\in\mathfrak{W}_w$ meeting $Q(\square_0)$, the
level $i_1:=\operatorname{lev}c_1$ satisfies $i_0-1\le i_1$ and $i_1\le k_w$. Hence the bounds of case (B)
in \eqref{4.25:eq-three-cases-a} hold, and $\kappa(\square_0)=k_w$ whichever $w\in M(\square_0)$ is used in
its definition.
\item[(none)] $c$ is of type (D1) and $M(\square_0)=\emptyset$. Then $\kappa(\square_0)=k_{\square}\ge i_0$,
and every $\lambda$ lies in $[i_0-1,i_0+1]\cap[0,k_{\square}]$.
\item[(C)] If $c\in\mathfrak{W}_{w'}$ for a present attempt $w'$ and $w\in M(\square_0)$, then $w=w'$.
\end{itemize}
In short,
\begin{equation}\label{4.25:eq-three-cases-a}
\begin{aligned}
&\text{(A)}\quad k_0^{w}\le\lambda\le i_0+1,\qquad i_0\le k_w=\kappa(\square_0),
\qquad M(\square_0)=\{w\};\\
&\text{(B)}\quad i_0-1\le\lambda\le i_0+1\le k_w+2,\qquad i_0\le k_w+1,\qquad\kappa(\square_0)=k_w\\
&\qquad\qquad(\text{for every } w\in M(\square_0));\\
&\text{(none)}\quad \lambda\in[i_0-1,i_0+1]\cap[0,k_{\square}],\qquad \kappa(\square_0)=k_{\square}\ge i_0;\\
&\text{(C)}\quad c\in\mathfrak{W}_{w'},\ w\in M(\square_0)\ \Longrightarrow\ w=w'.
\end{aligned}
\end{equation}
\end{corollary}

\begin{proof}
\emph{The cases are exclusive and exhaustive.} By (P2) applied to $c$, exactly one of the following holds:
$c$ is un-raised; $c$ is raised by exactly one present attempt $w_c$, with $c\subset Z_{w_c}$ and
$i_0\le k_0^{w_c}$; $c$ is raised by exactly one present attempt $w_c$, with $c\subset Z_{w_c}$ and
$i_0>k_0^{w_c}$. The first two alternatives say that $c$ is of type (D1). The third says that $c$ is a window
cell of $w_c$, and since $c$ is raised by only one attempt, $c$ lies in $\mathfrak{W}_w$ for at most one
present $w$. Cells of type (D1) split according as $M(\square_0)$ is empty or not. So exactly one of (A),
(B), (none) occurs.

Throughout we use that a cell meeting $Q(\square_0)$ meets $Q(c)$, because
$Q(\square_0)\subset Q(c)$ by \eqref{4.25:eq-response-cube-geometry-a}, and that a cell of
$\mathcal{A}(\tilde\square_0)$ meets $\tilde\square_0\subset Q(\square_0)$, hence meets $Q(c)$.

\emph{Case} (A). Here $c\in\mathfrak{W}_w$, so (D2) at $c$ gives $i_0\le k_w$ and $\lambda\ge k_0^w$ for every
cell meeting $Q(c)$, while (U) at $c$ gives $\lambda\le i_0+1$. This applies in particular to the cells of
$\mathcal{A}(\tilde\square_0)$. The cell $c$ itself belongs to $\mathfrak{W}_w$ and meets $Q(\square_0)$,
since $c\subset\square_0$; hence $w\in M(\square_0)$. By (C), proved below, every element of $M(\square_0)$
equals $w$, so $M(\square_0)=\{w\}$. By the definition \eqref{4.25:eq-entered-class-a},
$\kappa(\square_0)=k_w$, and $k_w\ge i_0$ was shown above.

\emph{Case} (B). Since $c$ is of type (D1), (D1) and (U) at $c$ give $i_0-1\le\lambda\le i_0+1$ for every cell
meeting $Q(c)$. Let $w\in M(\square_0)$ and let $c_1\in\mathfrak{W}_w$ meet $Q(\square_0)$. Then $c_1$ meets
$Q(c)$, so $i_1:=\operatorname{lev}c_1\ge i_0-1$ by (D1) at $c$. By (P2), every cell of $\mathfrak{W}_w$ has
level at most $k_w$, so $i_1\le k_w$. Together these give $i_0\le i_1+1\le k_w+1$, and then
\begin{equation*}
\lambda\ \le\ i_0+1\ \le\ k_w+2
\end{equation*}
for every cell meeting $Q(c)$. None of these bounds depends on which $w\in M(\square_0)$ was chosen. Hence
$\kappa(\square_0)=k_w$, as defined in \eqref{4.25:eq-entered-class-a}, satisfies them for each
$w\in M(\square_0)$.

\emph{Case} (none). Here $M(\square_0)=\emptyset$, so $\kappa(\square_0)=k_{\square}$ by
\eqref{4.25:eq-entered-class-a}. Since $c$ is of type (D1), it is not in the situation of (D2), so
$\kappa(c)=k_{\square}$ in (K), and (K) gives $i_0\le k_{\square}$. For a cell meeting $Q(c)$, (D1) and (U)
give $\lambda\in[i_0-1,i_0+1]$. By (P1), every cell of $\mathfrak{F}$ has level at most $k_{\square}$, and
levels are non-negative; so $\lambda\in[0,k_{\square}]$ as well.

\emph{Proof of} (C). Let $c\in\mathfrak{W}_{w'}$, let $w\in M(\square_0)$, and suppose $w\neq w'$. Choose
$c_1\in\mathfrak{W}_w$ meeting $Q(\square_0)$ and put $i_1:=\operatorname{lev}c_1$. Then $c_1$ meets $Q(c)$.
So (D2) and (U) at $c$, which is a window cell of $w'$, give
\begin{equation*}
k_0^{w'}\ \le\ i_1\ \le\ i_0+1,\qquad i_0\le k_{w'} .
\end{equation*}
By (P2), applied to $c\in\mathfrak{W}_{w'}$ and to $c_1\in\mathfrak{W}_w$,
\begin{equation*}
k_0^{w'}<i_0,\qquad k_0^{w}\ <\ i_1\ \le\ k_w .
\end{equation*}
By (P3), the two distinct present attempts $w$ and $w'$ are nested or disjoint. If they are nested, (P3)
says that the attempt of smaller class, which we call the older, lies inside the other: for nested present
attempts with $Z_{w_1}\subset Z_{w_2}$ and $w_1$ the older, one has $k_{w_1}\le h_{w_2}\le k_0^{w_2}-1$.
Three cases remain.

(C1) $Z_w\subset Z_{w'}$, $w$ the older. Then $k_w\le h_{w'}\le k_0^{w'}-1$, so
\begin{equation*}
i_1\ \le\ k_w\ <\ k_0^{w'}\ \le\ i_1 ,
\end{equation*}
a contradiction.

(C2) $Z_{w'}\subset Z_w$, $w'$ the older. Then $k_{w'}\le h_w\le k_0^w-1$, so
\begin{equation*}
i_1\ \le\ i_0+1\ \le\ k_{w'}+1\ \le\ k_0^w\ <\ i_1 ,
\end{equation*}
a contradiction.

(C3) $Z_w\cap Z_{w'}=\emptyset$. By (P2), $c$ is raised by $w'$ and $c\subset Z_{w'}$, and $c_1$ is raised by
$w$ and $c_1\subset Z_w$. Now (P4), applied to $(w',c)$ and to $(w,c_1)$, gives
\begin{equation*}
N(c):=\{\operatorname{dist}_\infty(\cdot,c)\le 2s_{k_{w'}}\}\subset Z_{w'},\qquad
N(c_1):=\{\operatorname{dist}_\infty(\cdot,c_1)\le 2s_{k_w}\}\subset Z_w .
\end{equation*}
These two sets are therefore disjoint. We use the following elementary fact. If $A$, $B$ are compact and
$\{\operatorname{dist}_\infty(\cdot,A)\le a\}$ and $\{\operatorname{dist}_\infty(\cdot,B)\le b\}$ are disjoint,
then $\operatorname{dist}_\infty(A,B)>a+b$. Indeed, otherwise choose $x\in A$ and $y\in B$ with
$|x-y|_\infty\le a+b$. If $a+b>0$, the point $z=x+\tfrac{a}{a+b}(y-x)$ satisfies $|z-x|_\infty\le a$ and
$|z-y|_\infty\le b$; if $a+b=0$, the point $z=x=y$ does. In either case $z$ lies in both sets.

Hence
\begin{equation*}
\operatorname{dist}_\infty(c,c_1)\ >\ 2s_{k_{w'}}+2s_{k_w}\ \ge\ 2s_{k_{w'}}\ \ge\ 2M\ell_{k_{w'}}
\ \ge\ 2M\ell_{i_0}.
\end{equation*}
The third inequality holds because the separation integer $\check R_{k_{w'}}$ entering the collar width
$s_{k_{w'}}$ satisfies $\check R_{k_{w'}}\ge1$. The last holds because $i_0\le k_{w'}$ and the cell sides
$M\ell_m$ increase with the level $m$.
Thus $c_1$ misses $Q(c)=\{\operatorname{dist}_\infty(\cdot,c)\le 2M\ell_{i_0}\}$, contradicting the fact
that $c_1$ meets $Q(c)$. Hence $w=w'$.
\end{proof}

\begin{remark}
This remark is not used in what follows. Also in case (B), $|M(\square_0)|\le1$. Suppose
$w_1\neq w_2$ both lie in $M(\square_0)$, and let $c_1\in\mathfrak{W}_{w_1}$ and $c_2\in\mathfrak{W}_{w_2}$
meet $Q(\square_0)$. The cube $\square_0$ has side $2M\ell_{i_0}$ and
$Q(\square_0)=\{\operatorname{dist}_\infty(\cdot,\square_0)\le M\ell_{i_0}\}$, so any two points of
$Q(\square_0)$ are at $\infty$-distance at most $4M\ell_{i_0}$. Hence
$\operatorname{dist}_\infty(c_1,c_2)\le 4M\ell_{i_0}$. Since $c_1$ and $c_2$ meet $Q(c)$, (D1) at $c$ and
(P2) give, for $j=1,2$,
\begin{equation*}
i_0\ \le\ \operatorname{lev}c_j+1\ \le\ k_{w_j}+1 .
\end{equation*}
Relabel so that $k_{w_1}\le k_{w_2}$. By (P3), the attempts are nested or disjoint.

\emph{Nested}: $Z_{w_1}\subset Z_{w_2}$, with $w_1$ the older, that is, of smaller class. Put
$m:=k_0^{w_2}+1$. By (P2) and (P3),
\begin{equation*}
\operatorname{lev}c_1\ \le\ k_{w_1}\ \le\ h_{w_2}\ \le\ m-2 .
\end{equation*}
Clause (iii) of the lemma that separates cells of low level from the region $\mathfrak{O}_m$, applied at
level $m$, then places $c_1$ at $\infty$-distance greater than
$8s_m=8M\ell_{k_{w_2}}$ from $\mathfrak{O}_m$. By clauses (W1)/(W2) and (W0) of the lemma that locates
window cells relative to the region $\Omega^{\mathrm{term}}_m$ of Definition~\ref{4.25:def-window-cells},
the window cell $c_2$ of $w_2$ lies within $7s_m$ of $\Omega^{\mathrm{term}}_m\cap V_{w_2}\subset\mathfrak{O}_m$.
Hence
$\operatorname{dist}_\infty(c_1,c_2)>M\ell_{k_{w_2}}$. Comparing with the bound $4M\ell_{i_0}$ gives
$L^{k_{w_2}-i_0}<4$, so $k_{w_2}-i_0\le1$, as $L^2\ge4$. On the other hand, the bound $i_0\le k_{w_1}+1$
above, the nesting inequalities of (P3), and the identity $k_0^{w_2}=k_{w_2}-N_0^{(k_{w_2})}$ of (D2) give
\begin{equation*}
i_0\ \le\ k_{w_1}+1\ \le\ h_{w_2}+1\ \le\ k_0^{w_2}\ =\ k_{w_2}-N_0^{(k_{w_2})}.
\end{equation*}
Therefore $N_0^{(k_{w_2})}\le k_{w_2}-i_0\le1$. This is excluded by $N_0^{(k)}\ge2$ for every $k$, which is
equivalent to $\check R_k\ge L$, that is, to clause (iv) of the conditions \eqref{4.25:eq-frozen-collar-weight-a}.

\emph{Disjoint}: $Z_{w_1}\cap Z_{w_2}=\emptyset$. By (P2), $c_j$ is raised by $w_j$ and $c_j\subset Z_{w_j}$.
Then (P4), applied twice, and the elementary fact in the proof of (C) give
\begin{equation*}
\operatorname{dist}_\infty(c_1,c_2)\ >\ 2s_{k_{w_1}}+2s_{k_{w_2}}\ \ge\ 4M\ell_{i_0}.
\end{equation*}
The last inequality uses $\check R_k\ge L$: then $s_k\ge M\ell_{k+1}$, and $M\ell_{k_{w_j}+1}\ge M\ell_{i_0}$
because $i_0\le k_{w_j}+1$. This contradicts $\operatorname{dist}_\infty(c_1,c_2)\le 4M\ell_{i_0}$.

The argument uses $\check R_k\ge L$, clause (iv) of \eqref{4.25:eq-frozen-collar-weight-a}, in both cases.
\end{remark}

\begin{definition}[The window-number bound on the depth]\label{4.25:def-window-number-bound}
For every class $k$ of the run write $x_k=g_k^{-2}$, and let $N_0^{(k)}$ be the depth fixed at step $k$:
an attempt $w$ of class $k_w$ has $k_0^w=k_w-N_0^{(k_w)}$, with $k_0^w$ as in
Lemma~\ref{4.25:lem-entered-class}. In Ba{\l}aban's construction \cite{B-CMP122a}
the depth is a function of the separation integer of the run at step $k$, written $\check{R}_k$ here.
Let $\mathsf{l}\ge 1$ be the window number, a function of the inverse coupling defined on a set of
inverse couplings $x>1$, such that for every $x$ in its domain
\begin{equation*}
c\,(\log x)^{4r}\le\mathsf{l}(x)^2\le C\,(\log x)^{4r},
\end{equation*}
where $0<c\le C$ are constants and $r>0$ is the exponent in $E(x)$ of
Notation~\ref{4.25:not-flow-inequalities}.
For real $t$ such that the domain of $\mathsf{l}$ meets $(-\infty,t]$ put
\begin{equation*}
\mathsf{l}^{\uparrow}(t):=\sup\bigl\{\mathsf{l}(t'):\ t'\le t,\ t'\ \text{in the domain of}\ \mathsf{l}\bigr\}.
\end{equation*}
Assume the window-number bound on the depth, stated at every class $k$ of the run --- the first line
of \eqref{4.25:eq-window-number-bound-a} --- and define $\mathsf{l}^{\sharp}$ by its second line:
\begin{equation}\label{4.25:eq-window-number-bound-a}
\begin{gathered}
L^{4N_0^{(k)}}\le \mathsf{l}(x_k)^2\quad\text{for every class $k$ of the run},\\
\mathsf{l}^{\sharp}(t):=4\,\mathsf{l}^{\uparrow}\bigl(t+2c_2+2B_\sharp\bigr).
\end{gathered}
\end{equation}
Let $\square_0$ be a response cube and let $\kappa(\square_0)$ be the class entered at $\square_0$
(Lemma~\ref{4.25:lem-entered-class}, display~\eqref{4.25:eq-entered-class-a}).
Then $\mathsf{l}^{\uparrow}$ is finite and non-decreasing, $\mathsf{l}\le\mathsf{l}^{\uparrow}$ on the domain of
$\mathsf{l}$, $\mathsf{l}^{\sharp}$ is non-decreasing, and
\begin{equation*}
L^{4N_0^{(\kappa(\square_0))}}\le \mathsf{l}\bigl(x_{\kappa(\square_0)}\bigr)^2 .
\end{equation*}
When $\mathsf{l}$ is non-decreasing, $\mathsf{l}^{\uparrow}=\mathsf{l}$ on the domain of $\mathsf{l}$, and
$\mathsf{l}^{\sharp}(t)=4\,\mathsf{l}(t+2c_2+2B_\sharp)$ whenever $t+2c_2+2B_\sharp$ lies in that domain.
\end{definition}

\begin{proof}
The class $\kappa(\square_0)$ is, by its definition in Lemma~\ref{4.25:lem-entered-class}, the class of a step
of the run, either a past step or the current one. The hypothesis is stated at every class, so it may be applied
at $k=\kappa(\square_0)$. This is legitimate also when $\kappa(\square_0)$ is a past step: both sides of
the inequality, the depth $N_0^{(\kappa(\square_0))}$ and the inverse coupling $x_{\kappa(\square_0)}$, were
fixed at that step and are not altered by any later step of the run, by the arrest theorem (the data fixed
at a step are left untouched by the later steps of the run), and by \cite{B-CMP122a}. The bound is therefore
taken at the inverse coupling $x_{\kappa(\square_0)}$ of that step, and no dependence on the number of steps
elapsed since then enters. This gives the bound at $\kappa(\square_0)$ stated above.

For $t_1\le t_2$ the set of $t'$ in the domain of $\mathsf{l}$ with $t'\le t_1$ is contained in the set of those
with $t'\le t_2$, so the supremum over the first is at
most the supremum over the second; thus $\mathsf{l}^{\uparrow}(t_1)\le\mathsf{l}^{\uparrow}(t_2)$. Taking $t'=t$
in the supremum, for $t$ in the domain of $\mathsf{l}$, gives $\mathsf{l}(t)\le\mathsf{l}^{\uparrow}(t)$.
For finiteness, use the upper bound
$\mathsf{l}(x)^2\le C(\log x)^{4r}$ on the domain of $\mathsf{l}$. Every $t'$ in the domain with $t'\le t$
satisfies $1<t'\le t$, hence $0<\log t'\le\log t$, and since $r>0$ this gives
$\mathsf{l}(t')^2\le C(\log t)^{4r}$; therefore
\begin{equation*}
\mathsf{l}^{\uparrow}(t)\le C^{1/2}(\log t)^{2r}<\infty .
\end{equation*}
The function $\mathsf{l}^{\sharp}$ is four times
the composition of the non-decreasing function $\mathsf{l}^{\uparrow}$ with the translation
$t\mapsto t+2c_2+2B_\sharp$, hence non-decreasing. If $\mathsf{l}$ is non-decreasing and $t$ lies in the
domain of $\mathsf{l}$, the supremum defining $\mathsf{l}^{\uparrow}(t)$ is attained at $t'=t$, so
$\mathsf{l}^{\uparrow}(t)=\mathsf{l}(t)$; substituting this at the point $t+2c_2+2B_\sharp$, when it lies in
that domain, in the definition of $\mathsf{l}^{\sharp}$ gives $\mathsf{l}^{\sharp}(t)=4\,\mathsf{l}(t+2c_2+2B_\sharp)$.
\end{proof}

\begin{lemma}[The cell count at a response cube]\label{4.25:lem-cell-count}
Throughout, $\ell_m$ denotes the length scale of level $m$, with $\ell_{m+1}=L\ell_m$: a cell of level $m$ is a closed cube of side $M\ell_m$ of the level-$m$ grid, and the grids of the successive levels are nested and $L$-adic.
\begin{enumerate}
\item[(a)] Let $Q$ be a closed cube of side $s$, and let $m$ be a level. The closed cells of level $m$ meeting $Q$ number at most $(s/(M\ell_m)+2)^d$.
\item[(b)] Let $d=4$. Let $\square_0$ be a response cube, in the sense of Lemma~\ref{4.25:lem-response-cube-geometry}, at a presentation of one of the types $O1$, $O4$, $O4'$. Let $i_0$ be its level, and let $\kappa:=\kappa(\square_0)$ be its entered class, as defined in \eqref{4.25:eq-entered-class-a}; in this lemma $\kappa$ abbreviates $\kappa(\square_0)$ and does not denote the auxiliary parameter. Assume the following:
\begin{itemize}
\item the level-spread theorem for current cells, in the form of the three cases at a response cube (Corollary~\ref{4.25:cor-three-cases}, with its cases (A), (B), (none) as in \eqref{4.25:eq-three-cases-a}), together with the preliminary (P1) of \eqref{4.25:eq-entered-class-a};
\item the window-number bound \eqref{4.25:eq-window-number-bound-a} at the class $k=\kappa$.
\end{itemize}
Then the set $\mathcal{A}(\tilde\square_0)$ of current cells meeting $\tilde\square_0$ (Definition~\ref{4.25:def-response-cube-cells}, \eqref{4.25:eq-response-cube-cells-a}) satisfies the bounds of line (b) of \eqref{4.25:eq-cell-count-a} below, with the constant $C'_{cnt}$ defined there.
Here $m_-$ is defined case by case:
\begin{itemize}
\item in case (A), $m_-:=k_0^w$, where $w$ is the attempt of that case, and then $m_-\ge i_0-N_0^{(\kappa)}$;
\item in cases (B) and (none), $m_-:=i_0-1$.
\end{itemize}
In case (B) the bounds of line (b) of \eqref{4.25:eq-cell-count-a} hold alike with $\kappa=k_w$ for every attempt $w\in M(\square_0)$, that is, for every present attempt a window cell of which meets $Q(\square_0)$.
\item[(c)] Assume in addition the coupling-flow inequalities \eqref{4.25:eq-flow-inequalities-a} and \eqref{4.25:eq-flow-inequalities-b}. Under the hypotheses of \textup{(b)}, $i_0\le\kappa+1$. Every level $\lambda$ of a cell of $\mathcal{A}(\tilde\square_0)$ satisfies $\lambda\le\kappa+2$ and $x_\kappa\le x_\lambda+2c_2+2B_\sharp$. Hence $|\mathcal{A}(\tilde\square_0)|\le C_{cnt}\,\mathsf{l}^{\sharp}(x_\lambda)^2$ for every such level $\lambda$, with $C_{cnt}$ as in \eqref{4.25:eq-cell-count-a}, a number depending on $d$ and $L$ only.
\end{enumerate}
The bounds of (a), (b) and (c) are collected in
\begin{equation}\label{4.25:eq-cell-count-a}
\begin{gathered}
\text{(a)}\quad \#\{\text{closed cells of level } m \text{ meeting } Q\}\le\bigl(s/(M\ell_m)+2\bigr)^d,\\
\text{(b)}\quad |\mathcal{A}(\tilde\square_0)| \le 3\cdot 6^d L^{d(i_0-m_-)} \le C'_{cnt}\,\mathsf{l}(x_\kappa)^2,\\
\text{(c)}\quad x_\kappa\le x_\lambda+2c_2+2B_\sharp,\qquad
|\mathcal{A}(\tilde\square_0)|\le C_{cnt}\,\mathsf{l}^{\sharp}(x_\lambda)^2,\\
C'_{cnt}:=3\cdot 6^dL^d,\qquad C_{cnt}:=C'_{cnt}/16 .
\end{gathered}
\end{equation}
\end{lemma}

\begin{proof}
\emph{(a)} Put $h:=M\ell_m$. The cells of level $m$ are closed cubes of the grid of mesh $h$, with pairwise disjoint interiors. Hence distinct cells of level $m$ are distinct grid cubes, and it suffices to count the grid cubes meeting $Q$.

Choose coordinates in which the grid cubes are the cubes $\prod_{i=1}^d[n_ih,(n_i+1)h]$ with $n\in\mathbb{Z}^d$, and write $Q=\prod_{i=1}^d[a_i,a_i+s]$. Two closed products of intervals meet if and only if they meet on every axis. On axis $i$, the interval $[n_ih,(n_i+1)h]$ meets $[a_i,a_i+s]$ if and only if $n_ih\le a_i+s$ and $(n_i+1)h\ge a_i$, that is,
\begin{equation*}
a_i/h-1\le n_i\le (a_i+s)/h .
\end{equation*}
This is an interval of length $s/h+1$, so it contains at most $s/h+2$ integers. Taking the product over the $d$ axes, the number of grid cubes meeting $Q$ is at most $(s/h+2)^d=(s/(M\ell_m)+2)^d$.

\emph{(b)} By Lemma~\ref{4.25:lem-response-cube-geometry}, the response cube $\square_0$ contains a current cell $c$ of level $i_0$. The same lemma gives the inclusions \eqref{4.25:eq-response-cube-geometry-a}, namely $\tilde\square_0\subset Q(\square_0)\subset Q(c)$. Hence every cell of $\mathcal{A}(\tilde\square_0)$ meets $Q(\square_0)$, and in particular meets $Q(c)$.

By Corollary~\ref{4.25:cor-three-cases}, every cell meeting $Q(c)$ has level in $[m_-,i_0+1]$. The upper end $i_0+1$ is the upper bound of Corollary~\ref{4.25:cor-three-cases}, valid in each of its cases. The lower end is $k_0^w$ in case (A) and $i_0-1$ in cases (B) and (none). Put $T:=i_0-m_-\ge 0$.
\begin{itemize}
\item In case (A), $c$ is a window cell of $w$, so its level exceeds $k_0^w$. Thus $T=i_0-k_0^w\ge 1$.
\item In cases (B) and (none), $T=1$.
\end{itemize}

Now apply (a) to the closed cube $Q(\square_0)$, whose side is $s=4M\ell_{i_0}$. Since the grids are $L$-adic, $\ell_{i_0}/\ell_{i_0-t}=L^t$ and $\ell_{i_0}/\ell_{i_0+1}=L^{-1}$. For $t\in[0,T]$, the cells of $\mathcal{A}(\tilde\square_0)$ of level $m=i_0-t$ therefore number at most
\begin{equation*}
(4L^t+2)^d\le(6L^t)^d,
\end{equation*}
using $L^t\ge 1$. Those of level $i_0+1$ number at most
\begin{equation*}
(4/L+2)^d\le(10/3)^d\le 6^d,
\end{equation*}
using $L\ge3$, which holds since $L$ is an odd integer greater than $1$. Summing over the levels in $[m_-,i_0+1]$ gives
\begin{equation*}
|\mathcal{A}(\tilde\square_0)|\le 6^d\Bigl[\sum_{t=0}^{T}L^{dt}+1\Bigr].
\end{equation*}

We bound the bracket. First,
\begin{equation*}
\sum_{t=0}^{T}L^{dt}=L^{dT}\sum_{u=0}^{T}L^{-du}\le L^{dT}(1-L^{-d})^{-1}\le \tfrac{81}{80}\,L^{dT},
\end{equation*}
because $L^{-d}\le 3^{-4}=1/81$ for $L\ge 3$ and $d=4$. Second, $1\le L^{dT}$. Together (the estimate also holds for $T=0$, where the bracket equals $2$),
\begin{equation}\label{4.25:eq-cell-count-1}
|\mathcal{A}(\tilde\square_0)|\le 6^d\bigl[\tfrac{81}{80}+1\bigr]L^{dT}\le 3\cdot 6^dL^{dT}
=3\cdot 6^dL^{d(i_0-m_-)}.
\end{equation}
This is the first inequality of line (b) of \eqref{4.25:eq-cell-count-a}.

\emph{Case (A).} By Corollary~\ref{4.25:cor-three-cases}, $i_0\le k_w$ and $\kappa=\kappa(\square_0)=k_w$. Recall $k_0^w=k_w-N_0^{(k_w)}$. Hence
\begin{equation*}
m_-=k_0^w\ge i_0-N_0^{(k_w)}=i_0-N_0^{(\kappa)},
\end{equation*}
which is the lower bound claimed for $m_-$, and
\begin{equation*}
T=i_0-k_0^w\le k_w-k_0^w=N_0^{(k_w)}=N_0^{(\kappa)}.
\end{equation*}
Since $d=4$, \eqref{4.25:eq-cell-count-1} and the window-number bound \eqref{4.25:eq-window-number-bound-a} at $k=\kappa$ give
\begin{equation*}
|\mathcal{A}(\tilde\square_0)|\le 3\cdot 6^dL^{4N_0^{(\kappa)}}\le 3\cdot6^d\,\mathsf{l}(x_\kappa)^2
=L^{-d}C'_{cnt}\,\mathsf{l}(x_\kappa)^2\le C'_{cnt}\,\mathsf{l}(x_\kappa)^2 .
\end{equation*}

\emph{Cases (B) and (none).} Here $T=1$, and the summation above gives
\begin{equation*}
|\mathcal{A}(\tilde\square_0)|\le 6^d(L^d+1+1)\le 3\cdot 6^dL^d=C'_{cnt}\le C'_{cnt}\,\mathsf{l}(x_\kappa)^2,
\end{equation*}
since $\mathsf{l}\ge1$. In case (B), Corollary~\ref{4.25:cor-three-cases} gives $\kappa=k_w$ for any $w\in M(\square_0)$. The bound just obtained does not depend on which $w$ is chosen, so it holds alike for every $w\in M(\square_0)$. In case (none), $\kappa=k_{\square}$, the index of the current step.

This proves line (b) of \eqref{4.25:eq-cell-count-a} in all three cases.

\emph{(c)} Let $\lambda$ be the level of a cell of $\mathcal{A}(\tilde\square_0)$. By (b), this cell meets $Q(c)$, so the upper bound of Corollary~\ref{4.25:cor-three-cases} gives $\lambda\le i_0+1$. We show $i_0\le\kappa+1$ case by case.
\begin{itemize}
\item In case (A), $i_0\le k_w=\kappa$ by Corollary~\ref{4.25:cor-three-cases}.
\item In case (none), $\kappa=k_{\square}\ge i_0$ by the preliminary (P1) of \eqref{4.25:eq-entered-class-a}, which bounds the level of every current cell by $k_{\square}$.
\item In case (B), let $w\in M(\square_0)$ with $\kappa=k_w$, and let $c_1$ be a window cell of $w$ meeting $Q(\square_0)$, of level $i_1$. Corollary~\ref{4.25:cor-three-cases} gives $i_0\le i_1+1$ and $i_1\le k_w$. Hence $i_0\le i_1+1\le\kappa+1$.
\end{itemize}
Thus $i_0\le\kappa+1$ and $\lambda\le i_0+1\le\kappa+2$ in every case. The bound is sharper in two situations: $\lambda\le\kappa+1$ in case (A), and $\lambda\le\kappa$ in case (none) and whenever $\kappa=k_{\square}$, by (P1).

Next we compare $x_\kappa$ with $x_\lambda$.
\begin{itemize}
\item If $\lambda\le\kappa$, apply the coupling-flow inequality \eqref{4.25:eq-flow-inequalities-a}, $x_i\ge x_{i'}-2c_2$ for $i\le i'$, with $i=\lambda$ and $i'=\kappa$. It gives $x_\kappa\le x_\lambda+2c_2$.
\item If $\lambda\in\{\kappa+1,\kappa+2\}$, apply the coupling-flow inequality \eqref{4.25:eq-flow-inequalities-b}, $x_i\le x_{i'}+B_\sharp(i'-i)$ for $i\le i'$, with $i=\kappa$ and $i'=\lambda$. Since $\lambda-\kappa\le 2$, it gives $x_\kappa\le x_\lambda+2B_\sharp$.
\end{itemize}
In case (A), where $\lambda\le\kappa+1$, only one $B_\sharp$ is used; the full $2B_\sharp$ is needed in case (B). Since $c_2\ge0$ and $B_\sharp\ge0$, both alternatives give
\begin{equation*}
x_\kappa\le x_\lambda+2c_2+2B_\sharp .
\end{equation*}

Recall from \eqref{4.25:eq-window-number-bound-a} that $\mathsf{l}^{\uparrow}$ is non-decreasing, that $\mathsf{l}\le\mathsf{l}^{\uparrow}$, and that $\mathsf{l}^{\sharp}(t)=4\,\mathsf{l}^{\uparrow}(t+2c_2+2B_\sharp)$. Therefore
\begin{equation*}
\mathsf{l}(x_\kappa)\le\mathsf{l}^{\uparrow}(x_\kappa)\le\mathsf{l}^{\uparrow}(x_\lambda+2c_2+2B_\sharp)
=\tfrac14\,\mathsf{l}^{\sharp}(x_\lambda).
\end{equation*}
Since $\mathsf{l}\ge1$, both sides are positive, and squaring gives $\mathsf{l}(x_\kappa)^2\le\mathsf{l}^{\sharp}(x_\lambda)^2/16$. Inserting this into line (b) of \eqref{4.25:eq-cell-count-a} gives
\begin{equation*}
|\mathcal{A}(\tilde\square_0)|\le C'_{cnt}\,\mathsf{l}^{\sharp}(x_\lambda)^2/16=C_{cnt}\,\mathsf{l}^{\sharp}(x_\lambda)^2 .
\end{equation*}
This holds for every level $\lambda$ of a cell of $\mathcal{A}(\tilde\square_0)$, and $C_{cnt}=3\cdot6^dL^d/16$ depends on $d$ and $L$ only.
\end{proof}

\begin{remark}[Remarks on the cell count]\label{4.25:rem-cell-count-remarks}
The following four observations concern the cell count at a response cube
(Lemma~\ref{4.25:lem-cell-count}) and the three cases at a response cube
(Corollary~\ref{4.25:cor-three-cases}). Throughout, $\square_0$ is a response cube of level $i_0$,
$\kappa=\kappa(\square_0)$ is the class entered at $\square_0$, and $\lambda$ denotes a
\emph{met level}, that is, the level of a cell of $\mathcal{A}(\tilde\square_0)$. For
$k\ge0$ we write $N_0^{(k)}$ for the depth of a window opened at level $k$, whether by an attempt
or at an arrest, so that $k_w-k_0^w=N_0^{(k_w)}$ for an attempt $w$ at level $k_w$.
\begin{enumerate}
\item[(1)] The bounds of Lemma~\ref{4.25:lem-cell-count} depend on none of the following:
the arrests present, their nesting depth, the ages of large-field regions, the volume, $K$,
$k$. The number $N_0^{(\kappa)}$ enters only through the polylogarithmic bound on the window
depth assumed in Lemma~\ref{4.25:lem-cell-count}, that is, as a polylogarithm of the inverse
coupling $x_\kappa$ of the window entered at $\square_0$ itself. It leaves through
clause~(c) of Lemma~\ref{4.25:lem-cell-count} at the met level $\lambda$, which is the level at
which the convergence bound is applied to the piece of the collar weight rooted in a cell of
$\mathcal{A}(\tilde\square_0)$ of that level.
If the window entered at $\square_0$ is
that of an older arrest $\mathfrak{a}=(W,k_W,\dots)$, made at step $k_W$, then the possibly
larger number $N_0^{(k_W)}$ is paid through the inequality
\begin{equation*}
x_{k_W}\le x_\lambda+2c_2+2B_\sharp ,
\end{equation*}
in which no age occurs. The difference $k_\square-k_w$ between the class $k_\square$ entered in
case~(none) of Corollary~\ref{4.25:cor-three-cases}, where $\kappa=k_\square\ge i_0$, and the
level $k_w$ of an attempt occurs nowhere.
\item[(2)] The count is not vacuous, because the response cube, a cube of the cover, protrudes
beyond the thin zone it meets. A thin zone of level $i$
is a single layer of cubes of level $i+1$, that is, $L$ cells of level $i$ wide
\cite[p.~179, (1.10)]{B-CMP122a}. A cube on the rim overhangs the zone by one cube, of
width $2M\ell_{i_0}$; this overhang satisfies
\begin{equation}\label{4.25:eq-cell-count-remarks-1}
2M\ell_{i_0}>\frac{M\ell_{i_0}L}{L-1},
\end{equation}
the right-hand side being $\sum_{t\ge0}M\ell_{i_0}L^{-t}$, the combined width of one cell of
each level $i_0,i_0-1,\dots$, so that the overhang crosses a cell of every finer level. It
thus reaches the finer zones lying below a window cell of an attempt $w$: at the large-field
operation $\mathbf{T}$ of step $k$ it reaches every level down to $k_0^w$, the lower end of the
range of met levels in case~(A) of Corollary~\ref{4.25:cor-three-cases}, and each of these finer
cells may carry the roots of
pieces of a distinct small arrest. In \cite{B-CMP122a} this overlap is accounted for by the
bound ``at most $N_0$ such terms''. The upper bound of clause~(b) of
Lemma~\ref{4.25:lem-cell-count} is not affected by this description.
\item[(3)] The factor $\mathsf{l}^{\sharp}(x_\lambda)^2$ produced by the count is paid for by
the envelope $\bar\alpha$, proportional to $x^{-1/2}$, which enters $\Theta(e)$ inside the
conversion envelope $\Pi_G^{blk}$; the latter is given by its defining display as a supremum and
decreases to zero.
\item[(4)] With $\kappa=\kappa(\square_0)$, as in
Lemma~\ref{4.25:lem-cell-count}, the upper level bound, the two lower level bounds and the class
bound that enter the proof of Corollary~\ref{4.25:cor-three-cases} give in case~(B) only
$\lambda\le\kappa+2$. The possible extra level is absorbed by the term $2B_\sharp$ of
clause~(c), and the constants $C_{cnt}$, $C'_{cnt}$ are unchanged. The value
$\lambda=\kappa+2$ would be excluded by the inequality $i_0\le k_w$ for every
$w\in M(\square_0)$ in case~(B), which would give $\lambda\le i_0+1\le\kappa+1$; this inequality
is neither proved nor used here.
\end{enumerate}
\end{remark}

\begin{proof}
(1) By their definitions in Lemma~\ref{4.25:lem-cell-count}, $C'_{cnt}=3\cdot6^dL^d$ and
$C_{cnt}=C'_{cnt}/16$ are numbers depending on $d$ and $L$ alone. Apart from these constants,
the bound of clause~(b) contains only $\mathsf{l}(x_\kappa)$, which comes from the
polylogarithmic bound on the window depth assumed in Lemma~\ref{4.25:lem-cell-count}, applied at
$\kappa$. Clause~(c) replaces $x_\kappa$ by $x_\lambda+2c_2+2B_\sharp$, a quantity determined
by the met level $\lambda$ and the constants $c_2$, $B_\sharp$. When the window entered at
$\square_0$ is that of an arrest made at step $k_W$, its class is $\kappa=k_W$, and the
inequality of clause~(c) reads as displayed. None of the
quantities listed in~(1) occurs in these steps, nor does $k_\square-k_w$.

(2) Since $L$ is an odd blocking factor with $L\ge3$, we have $L-1\ge2$ and hence
\begin{equation*}
\frac{L}{L-1}=1+\frac{1}{L-1}\le\frac32<2,
\end{equation*}
which gives \eqref{4.25:eq-cell-count-remarks-1} after multiplication by $M\ell_{i_0}>0$. A cell
of level $i_0-t$ has side $M\ell_{i_0-t}=M\ell_{i_0}L^{-t}$, and summing the geometric series
gives
\begin{equation*}
\sum_{t\ge0}M\ell_{i_0}L^{-t}=\frac{M\ell_{i_0}}{1-L^{-1}}=\frac{M\ell_{i_0}L}{L-1},
\end{equation*}
which identifies the right-hand side of \eqref{4.25:eq-cell-count-remarks-1}. The
layer structure of a thin zone is that of \cite[p.~179, (1.10)]{B-CMP122a}. The proof of
clause~(b) of Lemma~\ref{4.25:lem-cell-count} uses only the grid count of clause~(a) on
$Q(\square_0)$ and the range of met levels given by Corollary~\ref{4.25:cor-three-cases}; it
does not use the description of the overlap in~(2), so the upper bound stands under any
description of it.

(3) The envelope $\Pi_G^{blk}$ is given by its defining display as a supremum in which
$\Theta(e)$, and through it the
factor $\bar\alpha$ proportional to $x^{-1/2}$, occurs; by its definition it decreases to zero,
and the factor $\mathsf{l}^{\sharp}(x_\lambda)^2$ of the count is placed against this factor.

(4) In case~(B), the upper level bound, the two lower level bounds and the class bound give for
a met level only $\lambda\le\kappa+2$. For
$\lambda\in\{\kappa+1,\kappa+2\}$ the proof of clause~(c) of Lemma~\ref{4.25:lem-cell-count}
bounds $x_\kappa$ by $x_\lambda+2B_\sharp$, the running bound on the inverse couplings applied
over at most two levels, and this term is already contained in $x_\lambda+2c_2+2B_\sharp$
because $c_2\ge0$; hence the estimate of clause~(c), and with it $C_{cnt}$ and $C'_{cnt}$, is the
same whether $\lambda\le\kappa+1$ or $\lambda\le\kappa+2$. Concerning the value
$\lambda=\kappa+2$: since in case~(B) one may take $\kappa=k_w$ for any $w\in M(\square_0)$, and
the upper level bound gives $\lambda\le i_0+1$, the inequality $i_0\le k_w$ for
$w\in M(\square_0)$ would give $\lambda\le\kappa+1$. That inequality is neither proved nor used
here, and no step above depends on it.
\end{proof}

\begin{proposition}[Conversion at a response cube]\label{4.25:prop-response-cube-conversion}
Let $e\in\mathfrak{Q}(\mathfrak{a})$ be a piece of the quadratic collar weight of the arrest
$\mathfrak{a}$, presented at step $k'$ at a presentation to which the cell count at a response
cube (Lemma~\ref{4.25:lem-cell-count}) applies. Let $\square_0$ be a response cube, with
enlargement $\tilde\square_0$, such that $S^+(e)\cap\tilde\square_0\neq\emptyset$ as sets of
sites. Put $\mathcal{A}:=\mathcal{A}(\tilde\square_0)$, the set of cells meeting
$\tilde\square_0$ (see \eqref{4.25:eq-response-cube-cells-a}), and $n:=n(e)$. Then
\begin{equation}\label{4.25:eq-response-cube-conversion-a}
\begin{gathered}
|\mathcal{A}|\,\varrho^{2}\,\Theta(e)\,e^{-n}\le C_{cnt}\,\Pi_G^{blk}(x_{k'}),\\
\Pi_G^{blk}(x):=\sup_{t\ge x-2c_2}\,\sup_{n'\ge 0}\,e^{-n'}\,
P_{\mathfrak{b}}\bigl(t+2c_2+B_\sharp(n'+2)\bigr)\\
\qquad\qquad\times\mathsf{l}^{\sharp}\bigl(t+2c_2+B_\sharp(n'+3)\bigr)^{2}\,\bar\theta(t).
\end{gathered}
\end{equation}
The function $\Pi_G^{blk}$ decreases to zero: it is a polylogarithmic factor times
$\bar\theta$, and $\bar\theta$ is of order $x^{-1/2}$ times a polylogarithm.
\end{proposition}

\begin{proof}
By hypothesis there is a site $p\in S^+(e)\cap\tilde\square_0$. Let $\hat c$ be the current cell
containing $p$; it is the hat of the cell of $S^+(e)$ containing $p$, and by part~(i) of the
lemma on domination of old blocks by current blocks
(Lemma~\ref{4.25:lem-domination-old-blocks}) the hats of the cells of $S^+(e)$ are exactly the
current cells of the blocks of $\mathcal{Y}(e)$. Since $p\in\hat c\cap\tilde\square_0$, the cell
$\hat c$ meets $\tilde\square_0$, so $\hat c\in\mathcal{A}$, and its level $\lambda$ is therefore
a level of a cell of $\mathcal{A}(\tilde\square_0)$. Part~(c) of the cell count
(Lemma~\ref{4.25:lem-cell-count}, see \eqref{4.25:eq-cell-count-a}), applied at the class
$\kappa(\square_0)$ entered at the response cube, gives a bound valid for every such level, and
its conclusion does not involve $\kappa(\square_0)$; at the level $\lambda$ it reads
\begin{equation}\label{4.25:eq-response-cube-conversion-1}
|\mathcal{A}|\le C_{cnt}\,\mathsf{l}^{\sharp}(x_\lambda)^{2}.
\end{equation}

Let $y^*\in\mathcal{Y}(e)$ be a block at which $\theta$ attains its maximum over
$\mathcal{Y}(e)$, so that $\Theta(e)=\theta(y^*)$, and let $\ell'(y^*)$ be its level. Put
$t^*:=x_{\ell'(y^*)}$. Items~(1) and~(2) of the conversion of old polylogarithmic factors
(Proposition~\ref{4.25:prop-conversion-old-factors}), applied at $y=y^*$, give
\begin{equation}\label{4.25:eq-response-cube-conversion-2}
\Theta(e)=\theta(y^*)\le\bar\theta(t^*),\qquad t^*\ge x_{k'}-2c_2 .
\end{equation}
Item~(3) of the same proposition at $y=y^*$ gives
$x_{j_W}\le t^*+2c_2+B_\sharp(n+2)$, and therefore
\begin{equation}\label{4.25:eq-response-cube-conversion-3}
\varrho^{2}=P_{\mathfrak{b}}(x_{j_W})\le P_{\mathfrak{b}}\bigl(t^*+2c_2+B_\sharp(n+2)\bigr).
\end{equation}

Next we compare $\lambda$ with $\ell'(y^*)$. By part~(b) of the lemma on piece geometry and the
pile bound (Lemma~\ref{4.25:lem-pile-bound}), current levels on $S^+(e)$ differ by at most
$n(e)+3$; by parts~(i) and~(iii) of Lemma~\ref{4.25:lem-domination-old-blocks} this holds for
the hats, which are the cells of the blocks of $\mathcal{Y}(e)$. Both $\hat c$ and the cell of
$y^*$ are such hats, hence $|\lambda-\ell'(y^*)|\le n+3$. If $\lambda\le\ell'(y^*)$, the second
coupling-flow inequality \eqref{4.25:eq-flow-inequalities-b}, with $i=\lambda$ and
$i'=\ell'(y^*)$, gives
\begin{equation*}
x_\lambda\le x_{\ell'(y^*)}+B_\sharp\bigl(\ell'(y^*)-\lambda\bigr)\le t^*+B_\sharp(n+3).
\end{equation*}
If $\lambda>\ell'(y^*)$, the first coupling-flow inequality \eqref{4.25:eq-flow-inequalities-a},
with $i=\ell'(y^*)$ and $i'=\lambda$, gives $t^*=x_{\ell'(y^*)}\ge x_\lambda-2c_2$, that is,
$x_\lambda\le t^*+2c_2$. Since $2c_2=c_5\ge0$ and $B_\sharp>0$, in both cases
\begin{equation*}
x_\lambda\le t^*+2c_2+B_\sharp(n+3).
\end{equation*}
The function $\mathsf{l}^{\sharp}$ is non-decreasing (Definition~\ref{4.25:def-window-number-bound},
see \eqref{4.25:eq-window-number-bound-a}), so
\begin{equation}\label{4.25:eq-response-cube-conversion-4}
\mathsf{l}^{\sharp}(x_\lambda)\le\mathsf{l}^{\sharp}\bigl(t^*+2c_2+B_\sharp(n+3)\bigr).
\end{equation}

Multiplying \eqref{4.25:eq-response-cube-conversion-1},
\eqref{4.25:eq-response-cube-conversion-2} and \eqref{4.25:eq-response-cube-conversion-3},
inserting \eqref{4.25:eq-response-cube-conversion-4} and the factor $e^{-n}$, we obtain
\begin{align*}
|\mathcal{A}|\,\varrho^{2}\,\Theta(e)\,e^{-n}
&\le C_{cnt}\,e^{-n}\,P_{\mathfrak{b}}\bigl(t^*+2c_2+B_\sharp(n+2)\bigr)\\
&\qquad\times\mathsf{l}^{\sharp}\bigl(t^*+2c_2+B_\sharp(n+3)\bigr)^{2}\,\bar\theta(t^*).
\end{align*}
The right-hand side is the expression under the suprema in the definition of
$\Pi_G^{blk}(x_{k'})$, evaluated at $n':=n\ge0$ and $t:=t^*$, and $t^*\ge x_{k'}-2c_2$ by
\eqref{4.25:eq-response-cube-conversion-2}; it is therefore at most
$C_{cnt}\,\Pi_G^{blk}(x_{k'})$. This is the inequality in
\eqref{4.25:eq-response-cube-conversion-a}.

Finally, $\Pi_G^{blk}$ is non-increasing, because the range $t\ge x-2c_2$ of the outer supremum
shrinks as $x$ increases; it tends to zero because it is a polylogarithmic factor times
$\bar\theta$, and $\bar\theta$ is of order $x^{-1/2}$ times a polylogarithm.
\end{proof}

\begin{proposition}[Anchored block sum of the pieces]\label{4.25:prop-anchored-block-sum}
Assume the following statements.
\begin{enumerate}
\item[(a)] Clause $(\alpha)$ of Lemma~\ref{4.25:lem-pile-bound} (piece geometry and the pile bound), with
its constant $C_\alpha$, as displayed in \eqref{4.25:eq-pile-bound-a}.
\item[(b)] The cell count at a response cube, Lemma~\ref{4.25:lem-cell-count}, read at the class
$\kappa(\square_0)$ entered at that cube, with its constant $C_{cnt}$ as in
\eqref{4.25:eq-cell-count-a}.
\item[(c)] Conversion at a response cube, Proposition~\ref{4.25:prop-response-cube-conversion}, in the
form \eqref{4.25:eq-response-cube-conversion-a}.
\item[(d)] Three statements taken as given:
\begin{itemize}
\item the weight bound of the anchored estimate, which attaches to a piece $c$ with support
$S=\tilde S(c)$ the
weight $\mathfrak{w}(c)\,e^{\Sigma_c+d_k(S)/N_{wt}}$, whose exponential factor is one of the weights
listed in Proposition~\ref{4.25:prop-conversion-old-factors}, and which bounds the factor
$\mathfrak{w}(c)$
by the oscillation datum, $\mathfrak{w}(c)\le\mathfrak{N}^G(c;\mathfrak{O})$, with
$\mathfrak{N}^G$ as in \eqref{4.25:eq-oscillation-datum-a} and $\mathfrak{O}$ the later operation
presenting $c$. Here $\Sigma_c$ is the number which the weight bound attaches
to $c$ in the exponent (it is not a summation); $d_k(S)$ is the weight bound's distance number of
$S$ at the level $k$ at which the weight bound is read; and $N_{wt}$ is the constant by which the
weight bound divides $d_k(S)$;
\item the totals of the anchored estimate, which fix the quantities $K_\flat$, $\mathsf{r}$ and $V$
entering $\varepsilon_G$ below;
\item the lower bound $R(x)\ge L$ on the separation function $R$ (not the ball radius of the
observables), so that $R_k=R(x_k)\ge L$ at every level $k$.
\end{itemize}
\end{enumerate}

Let $k'$ be a step, and fix a presentation site of step $k'$ of the kind to which (b) and (c)
apply, at which the pieces of the live arrests below are presented. Let $\square_0$ be a response
cube there, and let $\tilde\square_0$ be its enlargement. Let
$\mathcal{A}(\tilde\square_0)$ be the family of cells meeting $\tilde\square_0$, as in
\eqref{4.25:eq-response-cube-cells-a}. The union below runs over the live arrests $\mathfrak{a}'$ to
which (a) applies. Then
\begin{equation}\label{4.25:eq-anchored-block-sum-a}
\begin{gathered}
\sum_{\substack{c\in\bigcup_{\mathfrak{a}'}\mathfrak{Q}(\mathfrak{a}')\\
S\cap\tilde\square_0\neq\emptyset}}\mathfrak{w}(c)\,e^{\Sigma_c+d_k(S)/N_{wt}}
\le\mathsf{n}_G^{blk}(x_{k'}),\\
\mathsf{n}_G^{blk}(x_{k'}):=\mathfrak{c}_G\,C_\alpha\,C_{cnt}\,K_a^{\circ}\,\Pi_G^{blk}(x_{k'}),
\qquad
\varepsilon_G:=K_\flat\,\mathsf{r}\,V\,\mathsf{n}_G^{blk}.
\end{gathered}
\end{equation}
Here $\mathfrak{c}_G$ is the weighted constant of Proposition~\ref{4.25:prop-conversion-old-factors},
and $K_\flat$, $\mathsf{r}$ and $V$ are the quantities fixed in the totals of the anchored estimate
($V$ is not the block field $V_k$).
\end{proposition}

The allowance $\mathsf{n}_G^{blk}(x_{k'})$ is entered in the totals of the anchored estimate. The
allowance $\varepsilon_G$ is entered in the collected remainder bound, and with the factor $5$ in the
fourth of the convergence conditions.

Among the conditions \eqref{4.25:eq-frozen-collar-weight-a}:
\begin{itemize}
\item the third is the requirement $\Pi_G^{blk}\le1$;
\item the fourth is $R(x)\ge L$, the third statement of (d); it is not a new condition, being
already contained in the ceilings imposed by the correlation theorem and by the bound for
typical regions.
\end{itemize}
The condition $R_k=R(x_k)\ge L$ is equivalent to $N_0\ge2$ for the number $N_0$ of the window-number
bound \eqref{4.25:eq-window-number-bound-a}. This is the floor of the window in the hypotheses on
which the cell count (b) depends.

\begin{proof}
Fix $k'$ and $\square_0$, and write $\mathcal{A}:=\mathcal{A}(\tilde\square_0)$.
For a piece $c$ in the index set of \eqref{4.25:eq-anchored-block-sum-a}, write $e:=c$ and $n:=n(e)$.
The support $S=\tilde S(e)$ of $e$ is contained in its $2$-neighbourhood $S^+(e)$ by definition of
the latter.

\emph{Case of an empty index set.} The left side of \eqref{4.25:eq-anchored-block-sum-a} is then zero.
The right side is non-negative, since the constants are positive and $\Pi_G^{blk}$ is a supremum of
products of non-negative factors. Assume from now on that the index set is non-empty.

\emph{The weights.} Each summand is a piece $e\in\mathfrak{Q}(\mathfrak{a}')$ presented at the site of
step $k'$ fixed above. By hypothesis (d), and since the exponential factor is positive, its weight is
at most $\mathfrak{N}^G(e;\mathfrak{O})\,e^{\Sigma_e+d_k(S)/N_{wt}}$. The exponential factor is one of the
weights listed in Proposition~\ref{4.25:prop-conversion-old-factors}, where it is bounded by
$e^{1/3+(3/2)\sqrt{d}|S^+(e)|}$. For the weights listed there, that proposition bounds
$\mathfrak{N}^G(e;\mathfrak{O})$ times the weight by
$\mathfrak{c}_G\,\mathfrak{n}^{\circ}(e)\,\varrho^2\Theta(e)e^{-n}$, before it applies the bound
$\varrho^2\Theta(e)e^{-n}\le\Pi_G(x_{k'})$ which yields its final form. Hence each summand satisfies
\begin{equation}\label{4.25:eq-anchored-block-sum-1}
\mathfrak{w}(e)\,e^{\Sigma_e+d_k(S)/N_{wt}}\le\mathfrak{c}_G\,\mathfrak{n}^{\circ}(e)\cdot
\bigl(\varrho^2\Theta(e)e^{-n}\bigr).
\end{equation}

\emph{The hypotheses of (a) and (c) hold for every summand.} Let $e$ be a summand and take a site
$p\in S\cap\tilde\square_0$. Then $p\in S^+(e)\cap\tilde\square_0$, so $S^+(e)$ and $\tilde\square_0$
have a site in common; this is the hypothesis of Proposition~\ref{4.25:prop-response-cube-conversion}.

Since $p\in S^+(e)$, the current cell containing $p$ is one of the cells of the current blocks
$\mathcal{Y}(e)$ meeting $S^+(e)$, as in the proof of
Proposition~\ref{4.25:prop-response-cube-conversion}; since it contains $p$, it meets
$\tilde\square_0$. By \eqref{4.25:eq-response-cube-cells-a} it therefore belongs to $\mathcal{A}$.
Consequently:
\begin{itemize}
\item $|\mathcal{A}|\ge1$;
\item $S^+(e)\cap\bigcup\mathcal{A}\neq\emptyset$, which is the condition on the pieces counted in
clause $(\alpha)$ of Lemma~\ref{4.25:lem-pile-bound}.
\end{itemize}
The family $\mathcal{A}$ consists of current cells and is finite by the cell count (b), so clause
$(\alpha)$ applies to it.

\emph{The sum.} Sum \eqref{4.25:eq-anchored-block-sum-1} over the index set and multiply and divide
each term by $|\mathcal{A}|\ge1$:
\begin{equation}\label{4.25:eq-anchored-block-sum-2}
\sum_e\mathfrak{c}_G\,\mathfrak{n}^{\circ}(e)\cdot\bigl(\varrho^2\Theta(e)e^{-n}\bigr)
=|\mathcal{A}|^{-1}\sum_e\mathfrak{c}_G\,\mathfrak{n}^{\circ}(e)\cdot
\bigl(|\mathcal{A}|\,\varrho^2\Theta(e)e^{-n}\bigr).
\end{equation}

For each $e$, Proposition~\ref{4.25:prop-response-cube-conversion}, in the form
\eqref{4.25:eq-response-cube-conversion-a}, gives
\begin{equation*}
|\mathcal{A}|\,\varrho^2\Theta(e)e^{-n}\le C_{cnt}\,\Pi_G^{blk}(x_{k'}).
\end{equation*}
This bound does not depend on $e$. Since $\mathfrak{n}^{\circ}(e)\ge0$, it may be inserted termwise.
Clause $(\alpha)$ of Lemma~\ref{4.25:lem-pile-bound}, applied with this $\mathcal{A}$, bounds the
remaining sum: the sum of $\mathfrak{n}^{\circ}(e)$ over the summands is at most the sum over all pieces
of all live arrests whose $S^+(e)$ meets $\bigcup\mathcal{A}$, which is at most
$C_\alpha|\mathcal{A}|K_a^{\circ}$. Hence the right side of \eqref{4.25:eq-anchored-block-sum-2} is
bounded by
\begin{equation*}
|\mathcal{A}|^{-1}\bigl(C_\alpha|\mathcal{A}|K_a^{\circ}\bigr)\,\mathfrak{c}_G\,C_{cnt}\,
\Pi_G^{blk}(x_{k'})
=\mathfrak{c}_G\,C_\alpha\,C_{cnt}\,K_a^{\circ}\,\Pi_G^{blk}(x_{k'})
=\mathsf{n}_G^{blk}(x_{k'}).
\end{equation*}

The constant $C_{cnt}$ entering here is the one of the cell count, hypothesis (b). The cell count
enters through its use at the class $\kappa(\square_0)$ inside
Proposition~\ref{4.25:prop-response-cube-conversion}, where its conclusion does not depend on the
class. Combining with \eqref{4.25:eq-anchored-block-sum-1} gives \eqref{4.25:eq-anchored-block-sum-a}.
\end{proof}

\begin{lemma}[Non-deep pieces are small. Stated; proof incomplete at the step
marked $(\ast)$]\label{4.25:lem-non-deep-pieces}
Let $\mathfrak{a}$ be an arrest of the region $W$, with level $j_W$, as in
Notation~\ref{4.25:not-collar-weight}. Let $e\in\mathfrak{Q}(\mathfrak{a})$ be a piece of the
collar weight $\mathsf{q}_W$, with support $\tilde S(e)$, its $2$-neighbourhood $S^+(e)$ and its
size $n(e)$. Let $k$ be a step.
\begin{itemize}
\item[(i)] Suppose that at step $k$ there is a healer of a large-field region $Z_h$ of level $h$
(in clause (i) and in its proof, $h$ denotes this level)
with $W\subset Z_h$, the inclusion holding by the dissolution rule of arrests recalled in
Notation~\ref{4.25:not-collar-weight}. Let $\mathsf{C}_C$ be the core of this healer. If
$S^+(e)\not\subset\mathsf{C}_C$, then
\begin{equation*}
n(e)\ \ge\ \check R_k+(h-j_W-2)^+-2 .
\end{equation*}
\item[(ii)] Let $Z'$ be a treated class that does not contain $W$. If
$S^+(e)\cap Z'\neq\emptyset$, then
\begin{equation*}
n(e)\ \ge\ \check R_k+(k-j_W-2)^+-2 .
\end{equation*}
\end{itemize}
In both cases
\begin{equation}\label{4.25:eq-non-deep-pieces-a}
\varrho^2e^{-n(e)}\ \le\ \Pi^{sh}(x_k)
:=e^2\sup_{t\ge x_k}E(t)\,\sup_{a\ge 0}e^{-a}P_{\mathfrak{b}}\bigl(2t+B_\sharp(a+2)\bigr),
\end{equation}
where $\varrho$ is the radius of the box $\mathfrak{b}_x$ of the collar variable, $P_{\mathfrak{b}}$
is the polynomial in the logarithm attached to that box, and $E$ is the function of
\eqref{4.25:eq-flow-inequalities-d}. The function $\Pi^{sh}$ has degree $-\infty$.
Consequently $2\nu^{\natural}(e)$, multiplied by the weights carried by the piece $e$, is at most
\begin{equation*}
2e^{3+2c_{wt}}\,\mathfrak{n}^{\circ}(e)\,\Pi^{sh}(x_k),
\end{equation*}
with $\nu^{\natural}$, $\mathfrak{n}^{\circ}$ and $c_{wt}$ as in
Notation~\ref{4.25:not-regularity-class}.
\end{lemma}

\begin{proof}
Cells, their levels, and the relation of touching between cells are those of
Notation~\ref{4.25:not-collar-weight}; a cell of level $m$ is a closed cube of side
$M\ell_m$. In both cases the host cells of the cubes of $S^+(e)$ contain a path of at most
$n(e)+2$ cells, consecutive cells of the path touching. The path starts at the host cell $\hat q$
of a root cube $q\in\mathsf{b}(e)$, whose level is at
most $j_W+1$. It ends in the complement $\mathsf{C}_C^{c}$ of the core in case (i), and in $Z'$
in case (ii).

\emph{Case (i).} Consecutive touching cells of the path differ in level by at most one.
Since the path starts at level at most $j_W+1$, it therefore contains a cell of each of the
levels $j_W+2,\dots,h-1$, that is, at least $(h-j_W-2)^+$ cells. The path also contains at least
$\check R_k$ cells of level $h$ lying in the moat, that is, among the cells of level $h$ between
$Z_h$ and $\mathsf{C}_C^{c}$. Indeed,
by the moat-width bound for the core of a healer, which follows from the nesting condition on
the cores, the distance from $Z_h$ to $\mathsf{C}_C^{c}$, measured in cells of level $h$,
satisfies
\begin{equation*}
\operatorname{dist}_h\bigl(Z_h,\mathsf{C}_C^{c}\bigr)\ \ge\ L\check R_{h+1}\ \ge\ \check R_k .
\end{equation*}
The cells of level $h$ counted here are distinct from those of the levels $j_W+2,\dots,h-1$.
Hence $n(e)+2\ge\check R_k+(h-j_W-2)^+$, which is the bound of (i).

\emph{Case (ii).} The same level-by-level argument shows that the path contains a cell of each
of the levels $j_W+2,\dots,k-1$, that is, at least $(k-j_W-2)^+$ cells. It then contains at
least $\check R_k$ cells of level $k$. For this we read at level $k$ the separation statement of
\cite[\S3]{B-CMP119}, which says of the regions to which it applies that ``the distance between
their boundaries is at least equal to $2MR_j$''. A
distance of at least $2MR_k$ spans at least $2R_k$ cells of level $k$. Here $R_k$ is the
separation integer fixed in the setting of the renormalisation steps $0,\dots,K$, and it equals
$\check R_k$, the one separation integer of these steps, by the convention fixed with
\eqref{4.25:eq-flow-inequalities-d} in Notation~\ref{4.25:not-flow-inequalities}; so
$2R_k\ge\check R_k$.
Hence $n(e)+2\ge\check R_k+(k-j_W-2)^+$, which is the bound of (ii).

If $\check R_k$ is read otherwise than as this integer, the conclusion
$\varrho^2e^{-n(e)}\le\Pi^{sh}(x_k)$ still holds, with the lower bound
$R_k\ge(\log g_k^{-2})^r$ of the setting of the renormalisation steps, used by statement, in
place of \eqref{4.25:eq-flow-inequalities-d}.

\emph{The bound \eqref{4.25:eq-non-deep-pieces-a}.} In case (ii) we put $h:=k$. Let $a\ge0$ be
the positive part that occurs in the case at hand:
\begin{equation*}
a=(h-j_W-2)^+\ \text{in case (i)},\qquad a=(k-j_W-2)^+\ \text{in case (ii)}.
\end{equation*}
The lower bounds on $n(e)$ give $e^{-n(e)}\le e^2e^{-\check R_k}e^{-a}$. In both cases the
definition of $a$ gives $h-j_W\le a+2$.

$(\ast)$ Let $N_k$ denote the window depth of \eqref{4.25:eq-flow-inequalities-e}, that is, the
depth of the window of the large-field operation at step $k$. In case (i) the level $h$ is the
bottom level of this window, $h=k-N_k$, as in \cite[p.~178, after (1.2)]{B-CMP122a}, where the
window is written with ``$h=k-N$''; in case (ii), where $h=k$, we have $k-h=0\le N_k$. In both
cases, therefore, $k-j_W\le N_k+h-j_W$. By the upper flow inequality
\eqref{4.25:eq-flow-inequalities-b}, applied with $i=j_W$ and $i'=k$, and the inequality
\eqref{4.25:eq-flow-inequalities-e},
which bounds $B_\sharp$ times the window depth by $x_k$, we obtain, in both cases,
\begin{equation}\label{4.25:eq-non-deep-pieces-1}
x_{j_W}\ \le\ x_k+B_\sharp(N_k+h-j_W)\ \le\ 2x_k+B_\sharp(a+2).
\end{equation}
Here \eqref{4.25:eq-flow-inequalities-e} is used at step $k$ as stated in
Notation~\ref{4.25:not-flow-inequalities}. It is an upper bound on the window depth $N_k$,
whereas the conditions on $N_k$ in force there are one-sided, each bounding the coupling from
above only; it is not derived from them in this proof.
Finally, \eqref{4.25:eq-flow-inequalities-d} bounds $e^{-\check R_k}$ by $E(x_k)$. The bound
$e^{-n(e)}\le e^2E(x_k)e^{-a}$ and the inequality \eqref{4.25:eq-non-deep-pieces-1} are the
inputs, at $t=x_k$ and the above value of $a$ in the suprema of
\eqref{4.25:eq-non-deep-pieces-a}, from which $\varrho^2e^{-n(e)}\le\Pi^{sh}(x_k)$ is drawn.
\end{proof}

\begin{lemma}[Oscillation of deep pieces]\label{4.25:lem-deep-oscillation}
Consider the healing operation at step $k$, with core $\mathsf{C}_C$. Let
$e\in\mathfrak{Q}(\mathfrak{a})$ be a piece of the quadratic collar weight of the arrest
$\mathfrak{a}$, and put $X:=S^+(e)$; assume $X\subset\mathsf{C}_C$.

We use the following letters of the three-field analysis of this operation:
\begin{itemize}
\item the presented pair $(\mathbb{W},\mathbb{J})$ and the rotation $\mathsf{u}'$;
\item the increment $(\mathcal{K},\jmath,\mathcal{Z})$, consisting of a potential, a current and a
rotation generator;
\item the family $\mathfrak{P}$ and its parameter $\varpi$;
\item the boxes $\mathfrak{b}_3,\mathfrak{b}_4$ of the collar variable;
\item the functional $F^{\mathfrak{b}}(\,\cdot\,;\mathsf{g}')$ of the three-field oscillation
conditions, built from the functional slot $\mathfrak{E}$ of those conditions, here
$\mathfrak{E}=v_e$, and evaluated at the base
rotation $\mathsf{g}'$;
\item the exponent data $\delta_P$ and $D_X$ of the factor $\rho_X$ below; the constants
$r^0$, $\varphi_\flat$ and $K_D$; and the unit $\mathsf{U}$ in which the increment bounds are
expressed;
\item the level $\ell'(y)$ attached to a current block $y$;
\item the enlargement $X^+$ of $X$.
\end{itemize}
Let the members $\mathsf{m},\mathsf{m}'\in\mathfrak{P}$ have the same value of $\varpi$, let
$\mathfrak{b}\in\{\mathfrak{b}_3,\mathfrak{b}_4\}$, and put
$\rho_X:=r^0e^{\frac12\delta_PD_X}$.

Assume the following statements of the three-field analysis.
\begin{enumerate}
\item[(1)] The presentation identities of the healing operation. In particular, the base pair
$(W_3,J_3)(\mathsf{m}')$ of the presentation at $\mathsf{m}'$ lies, on the whole domain, in
the cited regularity class $\mathfrak{K}_{109}(\mathbf{B}''_k;\gamma_0^{(k)})$, whose
parameters are $\mathbf{B}''_k$ and the current radius $\gamma_0^{(k)}$. It coincides on
$(D_2,\mathsf{C})$ with the pair $(U_2,J_2)$ of the presentation, $D_2$ and $\mathsf{C}$ being
the domain and the region of that pair, and $U_2$ is real. The presented pair at
$\mathsf{m}'$ is obtained from $(W_3,J_3)(\mathsf{m}')$ by the move with potential
$A''_0:=\Psi(\mathbb{H}_4,\mathbb{H}_3)(\mathsf{m}')$, the image of the layer data
$\mathbb{H}_3,\mathbb{H}_4$ under the potential map $\Psi$ of the presentation, and with
current $J''_0:=\sum\mathfrak{s}^{\sharp}$, the sum of the partner currents of the layers.
\item[(2)] The amplitude bound, with amplitude $\mathsf{a}_\star$.
\item[(3)] Bound (b) of the layer-map theorem, with its constant $B_{\mathbb{H}}$.
\item[(4)] The profile lemma. The profile $\omega$ of the three-field analysis satisfies
\begin{equation*}
N_y(\mathcal{K}),\ \ell^3|\jmath|\le\mathsf{U}\,\omega(y),\qquad
|\mathcal{Z}(s)|\le K_u\,\mathsf{U}\sup_{B^{top}(s)}\omega .
\end{equation*}
\item[(5)] The bound established in the proof of the halving bound: the decay function
$\varsigma$ of the three-field analysis, a function on current blocks, satisfies
$\mathsf{U}\,\omega(y)\le\varsigma(y)/\rho_X$ on $X^+$.
\item[(6)] The age-parameter bound of the three-field analysis, and the definition of the
constant $r^0$.
\item[(7)] The holomorphy of the presented family in the parameters $\vec{\mathsf{s}}$, and the
Cauchy estimate of the three-field analysis.
\item[(8)] The second and third members of item (i) of
\eqref{4.25:eq-frozen-collar-weight-a}, and the bound $\Theta(e)\le\varepsilon_{reg}/6750$
displayed in \eqref{4.25:eq-tube-constant-a} of Remark~\ref{4.25:rem-tube-constant}.
\end{enumerate}

Then the three-field oscillation conditions (a), (b) and (e) hold for $\mathfrak{E}=v_e$, with
oscillation datum $\tfrac12\mathfrak{N}^{G,L}(e)$. In particular,
\begin{equation}\label{4.25:eq-deep-oscillation-1}
\bigl|F^{\mathfrak{b}}(\mathsf{m};\mathsf{g}')-F^{\mathfrak{b}}(\mathsf{m}';\mathsf{g}')\bigr|
\le\frac{\mathfrak{N}^{G,L}(e)}{\rho_X},
\end{equation}
where
\begin{equation}\label{4.25:eq-deep-oscillation-a}
\begin{aligned}
\mathfrak{N}^{G,L}(e)&:=(c_1+c_2+3C_u)\,\tfrac98\,\nu^{\natural}(e)\,\Theta^L(e),\\
\Theta^L(e)&:=\frac{\varphi_\flat}{2K_D}\sup_{y\in\mathcal{Y}(e)}\alpha_{*,\ell'(y)}
\le\Theta(e).
\end{aligned}
\end{equation}
Here $c_1=1/\varepsilon_{reg}$ and $c_2=2/\gamma_{reg}=\max\{4K_{\Delta}',12dK_\sharp^2\}$,
with $c_2=0$ for the pieces $e_a^0$, are the
constants of Lemma~\ref{4.25:lem-piece-oscillation}, whose proof is incomplete at one step,
$C_u=1.1K_u$ is the constant of
Lemma~\ref{4.25:lem-reading-rotation}, and $\mathcal{Y}(e)$ is the set of
Definition~\ref{4.25:def-enlarged-support-blocks}. The tube constant $B_H$ enters this
lemma only through hypothesis (8). The constant $B_{\mathbb{H}}$ is that of the layer-map theorem.
\end{lemma}

\begin{proof}
\emph{The base pair.}
By hypothesis (1), the pair $(W_3,J_3)(\mathsf{m}')$ lies in
$\mathfrak{K}_{109}(\mathbf{B}''_k;\gamma_0^{(k)})$ on the whole domain. By clause (iii) of
Lemma~\ref{4.25:lem-gaussian-pieces}, whose proof is incomplete at one step, it therefore lies in
$\mathfrak{K}_G$.

It carries the room \eqref{4.25:eq-gaussian-pieces-a} of that clause: its potential $A'$
satisfies $N_y(A')<\theta(y)$ for $y\in\mathcal{Y}(e)$. Moreover,
\begin{equation*}
\theta(y)\le\Theta(e)\le\frac{\varepsilon_{reg}}{6750}<\frac{\varepsilon_{reg}}{20}.
\end{equation*}
The first inequality holds because $\Theta(e)$ is the supremum of $\theta$ over $\mathcal{Y}(e)$
(Definition~\ref{4.25:def-enlarged-support-blocks}). The second is hypothesis (8), the bound on
$\Theta$ in \eqref{4.25:eq-tube-constant-a} of Remark~\ref{4.25:rem-tube-constant}. The current
satisfies
\begin{equation*}
\ell^3|J_3|\le\gamma_0^{(k)}\le\tfrac14\gamma_{reg}.
\end{equation*}
The first inequality is the current entry of the room \eqref{4.25:eq-gaussian-pieces-a}; the
second, $\gamma_0^{(k)}\le\frac14\gamma_{reg}$, orders the current radius of the cited class
below a quarter of the coupling-free radius $\gamma_{reg}$, and is used together with that room.
By hypothesis (1), the pair coincides with $(U_2,J_2)$ on $(D_2,\mathsf{C})$, and $U_2$ is real.

\emph{The move to the presented pair.}
By hypothesis (1), the presented pair at $\mathsf{m}'$ is $(W_3,J_3)(\mathsf{m}')$ moved by the
potential $A''_0=\Psi(\mathbb{H}_4,\mathbb{H}_3)(\mathsf{m}')$ and the current
$J''_0=\sum\mathfrak{s}^{\sharp}$. Three inputs bound these:
\begin{itemize}
\item bound (b) of the layer-map theorem (hypothesis (3)), taken with profile identically one;
\item the amplitude bound (hypothesis (2));
\item the bound on the partner current of a difference of layers,
Lemma~\ref{4.25:lem-current-partner}.
\end{itemize}
Together they give
\begin{equation*}
N_y(A''_0)\le 2.5\,B_{\mathbb{H}}\mathsf{a}_\star,\qquad
\ell^3|J''_0|\le 2K_{\mathsf{j}}B_{\mathbb{H}}\mathsf{a}_\star .
\end{equation*}
These are at most $\varpi_G$ by the third member of item (i) of
\eqref{4.25:eq-frozen-collar-weight-a} (hypothesis (8)).

Write $\mathfrak{a}''_0:=\sup_{S^+(e)}N_y(A''_0)$ and
$\mathfrak{j}''_0:=\sup_{S^+(e)}\ell^3|J''_0|$. Clause (iv)$'$ of
Lemma~\ref{4.25:lem-gaussian-pieces}, whose proof is incomplete at one step, requires
\begin{equation*}
\mathfrak{a}''_0,\ N_y(A')\le\frac{\varepsilon_{reg}}{20},\qquad
\ell^3|J_3|,\ \mathfrak{j}''_0\le\tfrac14\gamma_{reg}.
\end{equation*}
The bounds on $N_y(A')$ and $\ell^3|J_3|$ were obtained in the base step; granted
$\varpi_G\le\varepsilon_{reg}/20$ and $\varpi_G\le\frac14\gamma_{reg}$, the bounds on
$\mathfrak{a}''_0$ and $\mathfrak{j}''_0$ follow from the bounds by $\varpi_G$ just obtained.
We apply clause (iv)$'$ to the base pair with $A'':=A''_0$ and $J'':=J''_0$. Since
$2\mathfrak{a}''_0+N_y(A')\le3\varepsilon_{reg}/20$, the moved potential then has size at most
\begin{equation*}
1.1\bigl(2\mathfrak{a}''_0+N_y(A')\bigr)\le0.165\,\varepsilon_{reg}\le\frac{\varepsilon_{reg}}{5}.
\end{equation*}
Hence $P:=Q(1,1)$ is a global $\mathfrak{K}_G$-continuation of
$(\mathbb{W},\mathbb{J})\lceil\tilde S(e)$. It carries the rooms required of the pair $P$ in
clause (iv) of that lemma, displayed in \eqref{4.25:eq-gaussian-pieces-b}.

\emph{The increment.}
By the profile lemma, hypothesis (4), the increment satisfies
\begin{equation*}
N_y(\mathcal{K}),\ \ell^3|\jmath|\le\mathsf{U}\,\omega(y),\qquad
|\mathcal{Z}(s)|\le K_u\,\mathsf{U}\sup_{B^{top}(s)}\omega .
\end{equation*}
By hypothesis (5), $\mathsf{U}\,\omega(y)\le\varsigma(y)/\rho_X$ on $X^+$. The age-parameter
bound, taken with its age parameter equal to zero, together with the definition of $r^0$
(hypothesis (6)), gives
\begin{equation*}
\varsigma(y)\le\frac{\varphi_\flat}{2K_D}\,\alpha_{*,\ell'(y)} .
\end{equation*}

Take $A'':=\mathcal{K}$ and $J'':=\jmath$ in Lemmas~\ref{4.25:lem-gaussian-pieces}
and~\ref{4.25:lem-piece-oscillation}, whose proofs are incomplete at one step each, and let
$\mathfrak{g}$ be the rotation increment generated
by $\mathcal{Z}$. For $\mathfrak{g}$ we use, besides the bound on $|\mathcal{Z}(s)|$, the deep-case
bound of Lemma~\ref{4.25:lem-reading-rotation}, at $\lambda=1$:
\begin{equation*}
\operatorname{dis}\bigl(\mathsf{u}'e^{\mathcal{Z}}\mathsf{u}'^{-1}\bigr)(s)
\le e^{2\operatorname{pol}\mathsf{u}'}\,|\mathcal{Z}(s)|;
\end{equation*}
this is where the constant $C_u=1.1K_u$ of that lemma enters. The three bounds above, together
with this one, then yield
\begin{equation*}
\mathfrak{a}'',\ \mathfrak{j}''\le\frac98\,\frac{\Theta^L(e)}{\rho_X},\qquad
\mathfrak{g}\le C_u\,\frac{\Theta^L(e)}{\rho_X}.
\end{equation*}
Since $\rho_X\ge r^0$, the second member of item (i) of \eqref{4.25:eq-frozen-collar-weight-a}
(hypothesis (8)) places $\mathfrak{a}''$, $\mathfrak{j}''$ and $\mathfrak{g}$ in the ranges
required by Lemma~\ref{4.25:lem-piece-oscillation}, whose proof is incomplete at one step.

\emph{The oscillation.}
In the deep case the rotation increment at which the piece is read is
$\mathsf{u}'e^{\lambda\mathcal{Z}}\mathsf{u}'^{-1}$, applied to the rotation $\mathsf{u}'\mathsf{g}'$,
and depends on a parameter $\lambda$.
We apply Lemma~\ref{4.25:lem-piece-oscillation}, whose proof is incomplete at one step, directly at
$\lambda=1$. This is possible because the right member of \eqref{4.25:eq-piece-oscillation-a} is
linear in $(\mathfrak{a}'',\mathfrak{j}'',\mathfrak{g})$. Consequently the bound uses the
parameter only at $\lambda=1$, holds for every value of $\rho_X$, and involves only the three
sizes $\mathfrak{a}''$, $\mathfrak{j}''$ and $\mathfrak{g}$.

The pair $P$ of the base step plays the role of $P$ in that lemma, and the increment is
$(\mathcal{K},\jmath)$ with rotation generated by $\mathcal{Z}$. By the construction of
$F^{\mathfrak{b}}$ from $\mathfrak{E}=v_e$ in the three-field analysis, the left member of
\eqref{4.25:eq-deep-oscillation-1} is controlled, uniformly in the collar variable $x$ in the
box $\mathfrak{b}$, for which $|x|_\infty\le\mathfrak{b}_x$ as that lemma requires, by the
oscillation
$|v_e^{\sharp}(x;p_2;g_2)-v_e^{\sharp}(x;p_1;g_1)|$ of that lemma, with $p_1=P\lceil\tilde S(e)$,
$p_2$ the pair moved by $(\mathcal{K},\jmath)$, $g_1=\mathsf{u}'\mathsf{g}'$ and
$g_2=\mathsf{u}'e^{\mathcal{Z}}\mathsf{u}'^{-1}g_1$. The bounds of the previous step give
\begin{align*}
\bigl|F^{\mathfrak{b}}(\mathsf{m};\mathsf{g}')-F^{\mathfrak{b}}(\mathsf{m}';\mathsf{g}')\bigr|
&\le\nu^{\natural}(e)\bigl[c_1\mathfrak{a}''+c_2\mathfrak{j}''+3\mathfrak{g}\bigr]\\
&\le\nu^{\natural}(e)\Bigl[\tfrac98(c_1+c_2)+3C_u\Bigr]\frac{\Theta^L(e)}{\rho_X}\\
&\le(c_1+c_2+3C_u)\,\tfrac98\,\nu^{\natural}(e)\,\frac{\Theta^L(e)}{\rho_X}
=\frac{\mathfrak{N}^{G,L}(e)}{\rho_X}.
\end{align*}
The last inequality holds because $3C_u\le\frac98\cdot3C_u$ and all terms are non-negative. This
is \eqref{4.25:eq-deep-oscillation-1}. It is condition (a), read with the oscillation datum
$\tfrac12\mathfrak{N}^{G,L}(e)$.

\emph{Condition (b).}
This is the Cauchy estimate of the three-field analysis (hypothesis (7)), applied as it stands to
$F^{\mathfrak{b}}-c$, with $c$ an arbitrary constant. The holomorphy in $\vec{\mathsf{s}}$ that
the estimate requires comes from two sources:
\begin{itemize}
\item the holomorphy of the presented family in $\vec{\mathsf{s}}$ (hypothesis (7));
\item clauses (i) and (iv) of Lemma~\ref{4.25:lem-gaussian-pieces}, whose proof is incomplete at
one step.
\end{itemize}
\end{proof}

\begin{lemma}[The norm bound at non-deep large-field sites]\label{4.25:lem-norm-bound}
Let $\mathfrak{a}$ be an arrest as in Notation~\ref{4.25:not-collar-weight}. Let $e$ be a piece of the
family $\mathfrak{Q}(\mathfrak{a})$ of pieces of the collar weight $\mathsf{q}_W$, with support
$\tilde S(e)$, its $2$-neighbourhood $S^+(e)$, piece $v_e$, orbit datum $v_e^{\sharp}$ and constant
$\nu^{\natural}(e)$ as in Lemma~\ref{4.25:lem-gaussian-pieces}. Suppose that $e$ is of the third kind
in the case distinction of the presentation statements for the later large-field operations, and
that $e$ is presented, through a presentation of the class considered in that case, by a later
large-field operation at a site that is not deep. Such a presentation comes with a parameter $z$
of the presenting family, a presented point $p(z)$ and a presenting gauge transformation $g(z)$, a
function on the sites with values in the complexification $G^{\mathbb{C}}$ of the gauge group $G$ of
Definition~\ref{4.4:def-regular-backgrounds}. Let $\varepsilon_\star$ denote the ceiling on the
polar size $\operatorname{pol}$ of presenting gauge transformations that is fixed in the presentation
theorem for the later large-field operations; it satisfies $e^{4\varepsilon_\star}\le 2$. Let $x$ lie
in the box $\mathfrak{b}_x$ of the collar variable, so that $|x|_\infty\le\varrho$. Assume the
following \emph{domain hypothesis}:
\begin{itemize}
\item at every presentation of $e$ at a non-deep large-field site, the presented point $p(z)$
is the restriction to the bonds of $\tilde S(e)$ of a global pair $\tilde P$ with the following
property: a later operation $\mathfrak{O}$, taken at a step $k'$, presents a space $\mathfrak{T}$
over a region $Y\supseteq S^+(e)$ in the ambient sense of clause~(iii) of
Lemma~\ref{4.25:lem-gaussian-pieces}, the pair $\tilde P$ belongs to the coupling-small class
$\mathfrak{K}_P^{(k')}$ on the whole domain, and its restriction is a member of $\mathfrak{T}$. The
presenting gauge transformation $g(z)$ satisfies $\operatorname{pol} g(z)<\varepsilon_\star$ at every
site.
\end{itemize}
Then, at every such presentation, $p(z)\in\mathfrak{K}_G(\tilde S(e))$ and
\begin{equation}\label{4.25:eq-norm-bound-1}
|v_e^{\sharp}(x;p(z);g(z))|\le 2\nu^{\natural}(e).
\end{equation}
\end{lemma}

The domain hypothesis is the conclusion, at non-deep sites, of the presentation statements for the
later large-field operations: of clause~$(\beta)$ of the off-diagonal complement to the presentation
theorem for the later large-field operations, in its near and arrested cases and for the fifth later
operation; of the window statement; and of the statement fixing the ambient sense of a presented
space; the bound on $\operatorname{pol} g(z)$ is the polar-size bound of the presentation theorem for
the later large-field operations. It is an input of this lemma.

\begin{proof}
Fix a presentation at a non-deep large-field site, and let $\tilde P$, $\mathfrak{T}$,
$\mathfrak{O}$, $k'$ and $Y$ be as in the domain hypothesis.

We first show that $p(z)\in\mathfrak{K}_G(\tilde S(e))$. The restriction of the global pair
$\tilde P\in\mathfrak{K}_P^{(k')}$ is a member of the space $\mathfrak{T}$ that $\mathfrak{O}$
presents over $Y\supseteq S^+(e)$ in the ambient sense, so clause~(iii) of
Lemma~\ref{4.25:lem-gaussian-pieces} (a lemma stated with its proof incomplete at one marked step)
applies to $\mathfrak{T}$ and gives $\tilde P\in\mathfrak{K}_G$ for this global pair. We have
$\tilde S(e)\subset S^+(e)\subset Y$, and $p(z)$ is the restriction of $\tilde P$ to the bonds of
$\tilde S(e)$. Hence $p(z)\in\mathfrak{K}_G(\tilde S(e))$.

We next bound the rotated variable. By clause~(ii) of Lemma~\ref{4.25:lem-gaussian-pieces} (a lemma
stated with its proof incomplete at one marked step),
\begin{equation*}
v_e^{\sharp}(x;p(z);g(z))=v_e\bigl(x^{[g(z)]};p(z)\bigr),\qquad
x^{[g]}(b)=\operatorname{Ad}_{g(b_-)}x(b).
\end{equation*}
Fix a bond $b$ and write $h=g(z)(b_-)$. Since $\operatorname{Ad}_hx_1=hx_1h^{-1}$ and
$|hx_1h^{-1}|\le\|h\|\,|x_1|\,\|h^{-1}\|$ for $x_1\in\mathfrak{g}^{\mathbb{C}}$ and the operator norm
$\|\cdot\|$ entering $\operatorname{pol}h=\log\max\{\|h\|,\|h^{-1}\|\}$, the bounds
$\operatorname{pol}h<\varepsilon_\star$ and $|x(b)|\le\varrho$ give
\begin{equation*}
|x^{[g(z)]}(b)|\le\|h\|\,\|h^{-1}\|\,|x(b)|\le e^{2\operatorname{pol}h}\,|x(b)|
\le e^{2\varepsilon_\star}\varrho.
\end{equation*}
As $b$ was arbitrary,
\begin{equation}\label{4.25:eq-norm-bound-2}
|x^{[g(z)]}|_\infty\le e^{2\varepsilon_\star}\varrho.
\end{equation}

By the first part of the proof $p(z)\in\mathfrak{K}_G(\tilde S(e))$, so clause~(i) of
Lemma~\ref{4.25:lem-gaussian-pieces} (a lemma stated with its proof incomplete at one marked step)
applies at $p(z)$, that is, on its domain: there $v_e(x_1;p(z))=Q_e(x_1,x_1)$ with $Q_e$ bilinear and
$|Q_e(x_1,x_2)|\le\nu^{\natural}(e)|x_1|_\infty|x_2|_\infty/\varrho^2$. Taking $x_1=x_2=x^{[g(z)]}$
and using \eqref{4.25:eq-norm-bound-2} gives
\begin{equation*}
|v_e^{\sharp}(x;p(z);g(z))|
\le\frac{\nu^{\natural}(e)}{\varrho^2}\bigl(e^{2\varepsilon_\star}\varrho\bigr)^2
=e^{4\varepsilon_\star}\nu^{\natural}(e)
\le 2\nu^{\natural}(e),
\end{equation*}
the last inequality by the property $e^{4\varepsilon_\star}\le 2$ of $\varepsilon_\star$. This proves
\eqref{4.25:eq-norm-bound-1}.
\end{proof}

\begin{proposition}[Dissolution of the arrested region. Stated; proof incomplete at the step marked
$(\ast)$]\label{4.25:prop-dissolution}
Let $\mathfrak{a}$ be a live arrest of the region $W$, made at step $k_W$ in stage $n_{\mathfrak{a}}+1$,
as in Definition~\ref{4.25:not-collar-weight}, and let $X_{\mathfrak{a}}\supseteq W$ denote the region
attached to $\mathfrak{a}$ whose points, by clause (T3) of the nesting lemma for arrests, no operation
before the dissolution of $\mathfrak{a}$ integrates or re-levels. Consider the completed operation
$\mathbf{R}$, of the first or of the second class, at step $k_b$ on the component $C$ with
$C\supseteq X_{\mathfrak{a}}\supseteq W$; it is the later presentation of kind $O5$ for $\mathfrak{a}$,
and we write $O5$ for it. A piece $e\in\mathfrak{Q}(\mathfrak{a})$ is called
\emph{deep} if $S^+(e)\subset\mathsf{C}_C$, and \emph{non-deep} otherwise. Then the following hold.
\begin{enumerate}
\item[(g$'$1)] \emph{(The last presentation.)} At $O5$ every piece $e\in\mathfrak{Q}(\mathfrak{a})$
that is still stored is presented once more. It is presented as a large-field unit of type (L) of the
first part of item (i$'$) of the fourth row of the summation display of the correlation theorem, with
support $S_v:=S^+(e)$:
\begin{itemize}
\item under the deep case of clause (e) of the presentation statement at an $\mathbf{R}$-site if $e$ is
deep;
\item under the non-deep case of that clause, with norm $2\nu^{\natural}(e)$, otherwise.
\end{itemize}
The containment required by that statement holds.

A piece for which $S^+(e)$ met the component of a completed chain other than the one on $C$, that is,
the component on which another chain of large-field operations was completed, has
already been presented on that component with the norm of the non-deep case, and it ceased there. Its
slot for $x$ is carried by the output of that chain until $O5$.

The collar variable $x\in S^P_W$, where $S^P_W$ denotes the space of the collar variable $x$ of
$\mathsf{q}_W$, becomes an integration variable at $O5$, and not before. Indeed, the
Gaussian integration $\mathbf{T}'_{k_W}(W)$, schematically of the form
$\int dx\,e^{\mathsf{q}_W(x;\,\cdot\,)}(\cdots)$, which is kept under the arrest, is one of
Ba{\l}aban's intermediate operations $\mathbf{T}'_k$, and it is
carried by the cluster activities of $O5$. Before $O5$ the variable $x$ is not integrated, and it is real
at $O1$--$O4$ and at $O4'$. The box $|x|\le\mathfrak{b}_x$ is the one of \cite[(1.59)]{B-CMP122a}.
The later presentations of a piece of $\mathfrak{Q}(\mathfrak{a})$ are those at $O1$--$O5$, $O4'$
included.

\item[(g$'$2)] \emph{(The output of $O5$.)} The output of $O5$ on $C$ is the one described by the
emission clause of the correlation theorem's proposition on completed operations $\mathbf{R}$. It
consists of the following.
\begin{itemize}
\item The stored terms $\mathbf{B}'^{(k_b)}$ and $R'^{(k_b)}$ on the new $\mathbb{B}_{k_b}$, in the
notation of Definition~\ref{4.25:not-collar-weight}. These are holomorphic on the current analyticity
spaces.
\item For every region $Y_i$ of the second class, the functions $\mathbf{T}'_k(Y_i)1$ of
\cite{B-CMP122b} multiplying the action density. These depend only on the new field variables $V_k$
restricted to $Y_i$.
\end{itemize}
None of these objects is a member of $\mathfrak{Q}(\mathfrak{a})$. None has $\mathfrak{K}_G(\tilde S(e))$
as its domain, and none has $x$ as a variable. The norms, domains and variables of everything that
$O5$ emits are those of that clause at the current value $x_{k_b}=g_{k_b}^{-2}$.

\item[(g$'$3)] \emph{(Nothing stored.)} After $O5$ the arrest $\mathfrak{a}$ and every piece nested
with it are dropped. The family $\mathfrak{Q}(\mathfrak{a})$, which is part of the family of terms stored
under $\mathfrak{a}$, ceases. The collection of arrested left-overs does not contain
$\mathfrak{Q}(\mathfrak{a})$.

Consequently the phrase ``every later presentation'', wherever it occurs, ranges over $O1$--$O5$, $O4'$
included. This applies in clauses (iii)--(iv) of Lemma~\ref{4.25:lem-gaussian-pieces}, in the
increment lemma and the domain statement of this section, in the continuation corollary, and in
clauses (b)--(f) of the presentation statement. For each arrest this is a finite list.

No statement about operations after $O5$ involves $v_e$, $\nu^{\natural}(e)$,
$\varrho^2=P_{\mathfrak{b}}(x_{j_W})$ or $\mathfrak{K}_G$. The age of $W$ enters at $O5$ only through
the data of clause (e) of the presentation statement, converted there
by the conversion statement at $x_{k_b}$ and, for non-deep pieces, through $\Pi^{sh}(x_{k_b})$ of
\eqref{4.25:eq-non-deep-pieces-a}.

\item[(g$'$4)] \emph{(No piece survives $O5$.)} The stored terms of the fourth row of the correlation
theorem's summation display are the following:
\begin{itemize}
\item the terms $\mathbf{B}$;
\item the stage terms;
\item the terms $\mathbf{B}'^{(j)}$;
\item the frozen stage terms of the collection of arrested left-overs;
\item the terms frozen under arrests.
\end{itemize}
None of them is a piece $v_e$. If a member of $\mathfrak{Q}(\mathfrak{a})$ were carried past $O5$, its
factor $P_{\mathfrak{b}}$ would be an age factor, to be compensated by the own-level scaling of the
large-field bounds. That compensation is not invoked, because by (g$'$3) this case does not occur.

\item[(g$'$5)] \emph{(Value.)} The datum that the fourth row of the correlation theorem's summation
display records at the pieces of
$\mathfrak{Q}(\mathfrak{a})$ is an oscillation datum. It is consumed by the three quantities of the
correlation theorem that read the fourth row's data at $\mathfrak{Q}(\mathfrak{a})$, namely the cell
quantity $\mathsf{K}_{cell}$, the quantity $\mathsf{P}(\mathfrak{a})$ attached to the arrest, and
the quantity $\mathfrak{v}^{val}$; none of them reads it as a value.

The value $\mathsf{q}_W(x;U_2,J_2)$ at the frozen real own-current pair $(U_2,J_2)$ stays under the
exponential of $\mathbf{T}'_{k_W}(W)$. Inside that exponential it is used only through its sign, in the
inner operation of the modulus statement of the correlation theorem, written $\mathbf{T}_h$ in that
statement.

In formulas, put $P^{\mathbb{R}}:=(U_2,J_2)$ on $D_2$, the domain on which the frozen pair is given,
and let the sums run over
$e\in\mathfrak{Q}(\mathfrak{a})$. Then
\begin{equation}\label{4.25:eq-dissolution-a}
\begin{gathered}
\sum_{e\ \text{deep}}v_e(x;P^{\mathbb{R}})
=\mathsf{q}_W(x;P^{\mathbb{R}})-\sum_{e\ \text{non-deep}}v_e(x;P^{\mathbb{R}}),\\
\Bigl|\sum_{e\ \text{non-deep}}v_e(x;P^{\mathbb{R}})\Bigr|
\le\frac18\sum_{e\ \text{non-deep}}\nu^{\natural}(e).
\end{gathered}
\end{equation}
The right-hand side of the inequality is of degree $-\infty$, by clause (i) of
Lemma~\ref{4.25:lem-non-deep-pieces}, whose proof is incomplete at one step.
\end{enumerate}
\end{proposition}

\begin{proof}
\emph{Proof of \textup{(g$'$1)}.}

\emph{Step 1: the arrest is dissolved at $O5$.} By the convention on arrests, an arrest is live from the
step at which it is made until the first later completed operation $\mathbf{R}$ on a component
containing its region, and it is dissolved there. Since $C\supseteq X_{\mathfrak{a}}\supseteq W$ and
$O5$ is that operation for $\mathfrak{a}$, the arrest $\mathfrak{a}$ is live up to $O5$ and is dissolved
at $O5$.

By clause (T4) of the nesting lemma for arrests, arrests nested with $\mathfrak{a}$ are dissolved
together with it.

\emph{Step 2: the range of later presentations.} By the definition of the stored family of an arrest,
the later presentations of a stored piece range over $O1$--$O5$, $O4'$ included. A piece that has not
ceased before $O5$ is therefore still stored at $O5$, and it is presented there.

The same definition treats a piece for which $S^+(e)$ meets the component of a completed chain other
than the one on $C$. Such a piece is absorbed there into the term $V(Y)$ that the definition of the
stored family attaches to the region $Y$ of that chain, and ceases. It was
presented there under the non-deep case, with the norm of that case. This is clause (ii) of
Lemma~\ref{4.25:lem-non-deep-pieces}, whose proof is incomplete at one step, applied with $Z'$ the
treated class of that component. Its slot for $x$ is carried by the output of that chain until $O5$.

\emph{Step 3: the case of the presentation.} At an $\mathbf{R}$-site the list of cases of clause (e) of
the presentation statement is exhaustive. The far case of clause (e) does not occur, because
$W\subset Z$, where $Z$ is the class treated at $O5$. Hence every
piece still stored at $O5$ falls under one of the two remaining cases, as follows.
\begin{itemize}
\item If $e$ is deep, it falls under the deep case of clause (e). The oscillation bound is then that of
the oscillation of deep pieces, Lemma~\ref{4.25:lem-deep-oscillation}.
\item If $e$ is non-deep, it falls under the non-deep case of clause (e). The norm there is
$2\nu^{\natural}(e)$ by
the norm bound at non-deep large-field sites, Lemma~\ref{4.25:lem-norm-bound}. The smallness of such a
piece is clause (i) of Lemma~\ref{4.25:lem-non-deep-pieces}, whose proof is incomplete at one step.
That clause applies because the arrest is being healed at $O5$, so that $W$ lies in the treated class
$Z_h$ of the healing step, where $h$ denotes the level of that class, by the dropping clause of the
convention on arrests.
\end{itemize}
In both cases the support of the unit is $S_v:=S^+(e)$. Its containment is the domain statement of
this section.

\emph{Step 4: the collar variable.} The Gaussian integral $\mathbf{T}'_{k_W}(W)$ is kept under the
arrest in the form of \cite{B-CMP122a}, with the box $|x|\le\mathfrak{b}_x$ of
\cite[(1.59)]{B-CMP122a}. It is one of the operations that Ba{\l}aban introduces as
\begin{quote}
the intermediate operations \dots{} denoted by $\mathbf{T}'_k$
\end{quote}
in \cite{B-CMP122a}. These operations are carried by the cluster activities of the healing step
\cite[(1.98)--(1.99)]{B-CMP122b}, and they are named in \cite[(1.100)]{B-CMP122b}. Clause (d) of the
statement on the terms kept at step $k_b$ identifies them at $O5$.

By clause (T3) of the nesting lemma for arrests, no operation before $O5$ integrates or re-levels a point
of $X_{\mathfrak{a}}$. Hence at $O1$--$O4$ and at $O4'$ the variable $x$ is not integrated. It is real
there, by the first clause of the presentation statement. It becomes an integration variable at $O5$,
inside $\mathbf{T}'_{k_W}(W)$.

\emph{Proof of \textup{(g$'$2)}.}

\emph{Step 1: what $O5$ emits.} We take the correlation theorem's proposition on completed operations
$\mathbf{R}$ by its statement. Its emission clause (e) returns the second group of terms into the
family $\mathbf{R}^{w}$, and the first group $\mathbf{B}'^{(k+1)}$ into the family $\mathbf{B}^{w}$, of
the correlation theorem's sourced terms.
Its carrying clause (a), together with the corresponding clause of the correlation theorem's treatment
of arrested terms, carries the arrested left-overs of every live arrest unchanged at $O1$--$O4$. Thus
clause (a) carries and clause (e) emits.

The terms produced by the stages of the operation, which appear at the head of clause (a), are consumed
there. They are not emitted.

The correlation theorem indexes the first group as $\mathbf{B}'^{(k+1)}$. In \cite[(1.98)]{B-CMP122b}
and in clause (d) of the statement on the terms kept at step $k_b$, the same terms are written
$\mathbf{B}'^{(k_b)}$, $R'^{(k_b)}$ at step $k_b$. These are one object.

\emph{Step 2: holomorphy and the bound.} These terms are holomorphic on the current analyticity spaces
by \cite[(1.98)--(1.99)]{B-CMP122b}. In \cite{B-CMP122b} Ba{\l}aban writes
\begin{quote}
The exponentiation (1.98) completes the R-operation
\end{quote}
and we use \cite[(1.99)]{B-CMP122b} in the form
\begin{equation*}
\bigl|\mathbf{R}'^{(k)}(X,(\mathbf{U},\mathbf{J}))\bigr|
\le O(1)\,c_1\exp\Bigl(-\bigl(1+\tfrac12\beta\bigr)\kappa\,d_{k,\cup Y_i}(X)\Bigr),
\end{equation*}
where the constants $c_1$, $\beta$, $\kappa$ and the distance $d_{k,\cup Y_i}(X)$ are those of
\cite[(1.99)]{B-CMP122b}.
Here $\mathbf{R}'^{(k)}$, at $k=k_b$, is the term written $R'^{(k_b)}$ in Step~1.

\emph{Step 3: the second-class regions.} For regions $Y_i$ of the second class, Ba{\l}aban writes in
\cite{B-CMP122b} that
\begin{quote}
there are only the characteristic functions, the $\delta$-functions and the functions
$\mathbf{T}'_k(Y_i)1$ multiplying the action density
\end{quote}
and that these functions
\begin{quote}
depend only on the new field variables $V_k$ restricted to $Y_i$, therefore they do not participate in
operations connected with next renormalization transformations.
\end{quote}

\emph{Step 4: a different object.} The functional $V(Y,U_k)$ of \cite[(1.68)]{B-CMP122b} involves, inside
the integral, Ba{\l}aban's quantities $A$, $B$ and $V''$. It is not one of the functions
$\mathbf{T}'_k(Y_i)1$ of Step~3, nor one of the terms of Steps~1 and~2; this is the content of the
remark on $V(Y,U_k)$ accompanying the convention on arrests.

\emph{Step 5: conclusion.} None of the objects listed in Steps~1--3 is a piece $v_e$. None of them is
defined on $\mathfrak{K}_G(\tilde S(e))$, and none has $x$ as a variable. By the clause of the convention
on arrests that fixes the coupling of the output of a completed operation, their norms, domains and
variables are those of the emission clause at $x_{k_b}$.

\emph{Proof of \textup{(g$'$3)}.} By the dropping clause of the convention on arrests, after $O5$ the
arrest $\mathfrak{a}$ and every piece nested with it are dropped. By the definition of the stored family
of an arrest, $\mathfrak{Q}(\mathfrak{a})$ ceases with it, and the collection of arrested left-overs no
longer contains it.

Hence every later presentation of a piece of $\mathfrak{Q}(\mathfrak{a})$ is one of $O1$--$O5$ or $O4'$.
This is a finite list, and it is the range of the phrase in each of the statements listed in (g$'$3).
After $O5$ no piece $v_e$ is present, so nothing downstream involves $v_e$, $\nu^{\natural}(e)$,
$\varrho^2=P_{\mathfrak{b}}(x_{j_W})$ or $\mathfrak{K}_G$.

The dependence on the age of $W$ at $O5$ sits in the data of clause (e) of the presentation statement,
presented there. By the conversion statement at $x_{k_b}$ it is converted there, and for non-deep
pieces it is bounded through $\Pi^{sh}(x_{k_b})$ of
\eqref{4.25:eq-non-deep-pieces-a}.

\emph{Proof of \textup{(g$'$4)}.} The list of the stored terms of the fourth row of the summation display
is the one given in the
correlation theorem's statement on stored terms, and it contains no piece $v_e$. The age compensation
that would be needed for a member of $\mathfrak{Q}(\mathfrak{a})$ carried past $O5$ is the one named in
clause (d) of the statement on the terms kept at step $k_b$. By (g$'$3) no such member exists,
so that compensation is not used.

\emph{Proof of \textup{(g$'$5)}.}

\emph{Step 1: the value enters only through its sign.} The datum that the fourth row of the summation
display records at the
pieces of $\mathfrak{Q}(\mathfrak{a})$ is an oscillation datum: it enters the three quantities of the
correlation theorem that consume it, $\mathsf{K}_{cell}$,
$\mathsf{P}(\mathfrak{a})$ and $\mathfrak{v}^{val}$, only as such, never as a value. The value
$\mathsf{q}_W(x;U_2,J_2)$ remains in
the exponent of $\mathbf{T}'_{k_W}(W)$. There it enters only through its sign, in the inner operation
of the modulus statement of the correlation theorem, written $\mathbf{T}_h$ there, which we take by its
statement.

The pair $P^{\mathbb{R}}=(U_2,J_2)$ on $D_2$ is the frozen real own-current pair of the three-field
construction, and $U_2$ is real.

\emph{Step 2: the identity.} $(\ast)$ Property (d) of Definition~\ref{4.25:not-collar-weight} gives the
decomposition $\mathsf{q}_W=\sum_{e\in\mathfrak{Q}(\mathfrak{a})}v_e$ of the collar weight into its
pieces on the domain on which the pieces are defined, namely that of clause (i) of
Lemma~\ref{4.25:lem-gaussian-pieces}, whose proof is incomplete at one step. We use it at the point
$(x;P^{\mathbb{R}})$. Its evaluation at the frozen pair $P^{\mathbb{R}}$, which is not shown here to lie
in $\mathfrak{K}_G(\tilde S(e))$ for every $e\in\mathfrak{Q}(\mathfrak{a})$, is the step that is not
completed.
Splitting the sum into deep and non-deep pieces and moving the non-deep part to the right gives the
first line of \eqref{4.25:eq-dissolution-a}.

\emph{Step 3: the bound.} Let $e$ be non-deep. By clause (i) of Lemma~\ref{4.25:lem-gaussian-pieces},
whose proof is incomplete at one step, $v_e(x)=Q_e(x,x)$, and
$|Q_e(x,x)|\le\nu^{\natural}(e)|x|_\infty^2/\varrho^2$ on $\mathfrak{K}_G(\tilde S(e))$; we apply this
bound at $P^{\mathbb{R}}$ under the membership left open at $(\ast)$. Since $|x|\le\mathfrak{b}_x$ by
(g$'$1),
\begin{equation}\label{4.25:eq-dissolution-1}
|v_e(x;P^{\mathbb{R}})|
\le\nu^{\natural}(e)\Bigl(\frac{\mathfrak{b}_x}{\varrho}\Bigr)^2
\le\frac18\,\nu^{\natural}(e).
\end{equation}
Summing \eqref{4.25:eq-dissolution-1} over the non-deep pieces and using the triangle inequality gives
the second line of \eqref{4.25:eq-dissolution-a}.

\emph{Step 4: degree.} That the right-hand side is of degree $-\infty$ is clause (i) of
Lemma~\ref{4.25:lem-non-deep-pieces}, whose proof is incomplete at one step, through the envelope
\eqref{4.25:eq-non-deep-pieces-a}.

Thus (g$'$1)--(g$'$4) follow from the convention on arrests, the nesting lemma, the definition of the
stored family and the emission clause of the correlation theorem, each taken by statement. The value
part of (g$'$5) is the modulus statement of the correlation theorem, and both lines of
\eqref{4.25:eq-dissolution-a} rest on the step $(\ast)$.
\end{proof}

\begin{theorem}[Frozen Gaussian collar weight of an arrested region]
\label{4.25:thm-frozen-collar-weight}
Use the setting and notation of Notation~\ref{4.25:not-collar-weight}, the constants of
Notation~\ref{4.25:not-regularity-class}, the flow inequalities and block thresholds of
Notation~\ref{4.25:not-flow-inequalities}, and the sets $\mathcal{Y}(e)$ of current blocks
of Definition~\ref{4.25:def-enlarged-support-blocks}. In every bound below, a statement made
for ``$y\in S^+(e)$'' means $y\in\mathcal{Y}(e)$.

A \emph{later presentation} of the family $\mathfrak{Q}(\mathfrak{a})$ is a presentation at a
step $k'\ge k_W$, either at a later stage $l'>n_{\mathfrak{a}}+1$ of the arresting site or by
one of the operations $O1\text{--}O5$, $O4'$. It is of kind (P) at the
later stages $l'>n_{\mathfrak{a}}+1$ of the arresting site and at $O1$, $O4$, $O4'$, of kind
(T) at $O3$, and of kind (L) at a large-field operation
$\mathbf{R}$ ($O2$, $O5$). A presentation of kind (L) is \emph{far} for a piece $e$ if
$S^+(e)\cap Z=\emptyset$ for the treated class $Z$ of the site. It is \emph{deep} for $e$ if
$S^+(e)\subset\mathsf{C}_C$, the core of the healing operation.

Assume the following.
\begin{itemize}
\item The hypotheses of the statement that freezes the Gaussian collar weight at an arrest.
\item The floor \eqref{4.25:eq-regularity-class-b} on $\kappa_1$, and the floor $M\ge 2L^2$
of \eqref{4.25:eq-collar-weight-b}.
\item The ceiling \eqref{4.25:eq-frozen-collar-weight-a} of clause~(h) below, in which the tube
constant enters through $\hat B_H:=\max\{1,B_H\}$ (Remark~\ref{4.25:rem-tube-constant}).
\item The increment hypothesis at later presentations
(Definition~\ref{4.25:def-increment-hypothesis}), on $\mathcal{Y}(e)$, at every later
presentation of kind (P), of kind (T) and of far kind (L). At far presentations of kind (L)
we assume in addition the per-layer bound \eqref{4.25:eq-increment-bound-b}. These are
supplied as follows.
\begin{itemize}
\item At kind (P), the increment hypothesis is supplied by the bounds of the response stages
on the layer and its partner current, together with the uniform bound that the anchored moment
bound gives on its own radii.
\item At kind (T) and at far kind (L), it is the conclusion of Lemma~\ref{4.25:lem-increment-bound}.
That conclusion is drawn in both cases $d_{sc}(y,\mathfrak{F}_{\mathfrak{s}})\ge1$ and
$d_{sc}(y,\mathfrak{F}_{\mathfrak{s}})<1$ of that lemma, the two parts of the layer holding
$\mathsf{t}$ being taken jointly through the factor $(1+\mathsf{w}_{\mathfrak{s}})$, on the
disc \eqref{4.25:eq-increment-bound-c}. It is drawn from the assumptions
(e1)--(e3)$^+$ of \eqref{4.25:eq-increment-bound-a} of that lemma.
\item Those assumptions in turn follow from the following bounds on the layer sizes at the
sites of the correlation theorem. At kind (T), part~(b) of the layer-size theorem, which bounds
the four-entry size $\mathsf{N}_y$ of a layer datum $\mathbb{H}$, holds
with the exponential weight of the scaled distance to the small-field region of the step, and
the weighted bound stated together with that theorem holds with the exponential weight of the
distance to the response cube $\square$; both weights lie in the admissible weight class of that
theorem. At kind (L), the cut layer $\chi\mathbb{H}_s$, $\chi:=t+\mathsf{t}\hat\chi_{\square}$
with $t\in[0,1]$ fixed, of the layer datum $\mathbb{H}_s$ and its potential $A'_s$, with
partner current $\mathcal{J}^{\sharp}[\chi]$, is measured by five entries, namely the four
uncut ones and the tail of $P_1$; the tail is given in the pointwise product profile, which
carries no units,
\begin{equation*}
\ell_y^3\,\bigl|P_1(\sigma)(\hat\chi_{\square}A'_s)\bigr|\le\eta_p\,\mathsf{f}(y)\,\mathsf{U}_L ,
\end{equation*}
and the entries are read after the bound (e3)$^+$ of \eqref{4.25:eq-increment-bound-a}.
Further, $\eta_p=1$ when $\tilde\square$ meets $S^+(e)$; these bounds involve both the tube
constant $B_H$ and the layer-map constant $B_\sharp$ (see Remark~\ref{4.25:rem-tube-constant});
$\mathsf{w}_T$ is the radius of the homogeneity bound of the tube; and
the kernel irregularity $\mathsf{p}$ of the kernel of each layer is at most $1$.
\end{itemize}
\item At every deep presentation of kind (L): the inputs of
Lemma~\ref{4.25:lem-deep-oscillation}, that is, the three-field bounds which that lemma
takes by statement, among them the amplitude bound, those on the moat being consequences of
the nesting statement.
\item At every other presentation of kind (L): the domain statement for large-field
presentations of the correlation theorem.
\item At every presentation of kind (P): the statement on the levels met at a response cube,
whose window floor is member~(iv) of the ceiling below, and the condition at the
class $\kappa(\square_0)$, the two hypotheses under which Lemma~\ref{4.25:lem-cell-count} is
stated. Here
$\kappa(\square_0)$ is the class of the one window into which $Q(\square_0)$ enters.
\item The base-class hypothesis on the unitary factor map $\mathfrak{u}$, one of the two
hypotheses under which Lemma~\ref{4.25:lem-reading-rotation} is stated.
\item The statements taken from other chapters in the proof below:
\begin{itemize}
\item the arrest theorem, namely its storage definitions and its clauses on the creation, life
and dissolution of an arrest;
\item the anchored moment bound and its totals;
\item the block-format proposition;
\item the complex interpolation corollary of the relative small-field machine;
\item from the correlation theorem: the oscillation slot of its lemma that bounds stored terms
of the fourth class, its sign lemma, and the carrying part and the emitting part of its output
proposition.
\end{itemize}
\end{itemize}

Let $\mathfrak{a}$ be a live arrest of Gaussian-collar type, of the region $W$ at step $k_W$
and stage $n_{\mathfrak{a}}+1$. Then, with ``later presentation'' ranging over
$O1\text{--}O5$, $O4'$ included, together with the later stages $l'>n_{\mathfrak{a}}+1$ of the
arresting site (see (g$'$1) and (g$'$3) of
Proposition~\ref{4.25:prop-dissolution}, whose proof is incomplete at one step), the following
hold.
\begin{itemize}
\item[(a)] \emph{Storage.} The family $\mathfrak{Q}(\mathfrak{a})$ belongs to the family of
terms stored under $\mathfrak{a}$. Its members are stored terms with the following datum.
\begin{itemize}
\item Domain: $\mathfrak{K}_G(\tilde S(e))$.
\item Declared support, and region on which its bounds are asserted: $S^+(e)$, the
hypotheses being read on $\mathcal{Y}(e)$.
\item Label: $A_e$, the union of the $\pi_{m_0(e)}$-cubes meeting $S^+(e)$, assigned at the
creation of the term.
\item A real variable $x$ in the box $|x|\le\mathfrak{b}_x$ of \cite[(1.59)]{B-CMP122a},
not integrated before $O5$.
\item The orbit datum $v_e^{\sharp}$ of Lemma~\ref{4.25:lem-gaussian-pieces}(ii), whose proof
is incomplete at one step; it is entire in the rotation.
\item No strong cube, no run slot and no anchor slot.
\end{itemize}
At stage $n_{\mathfrak{a}}+1$ the weight $\mathsf{q}_W$ depends on the variable $\lambda$ of
the statement that freezes the Gaussian collar weight at an arrest and on nothing else, as that
statement provides; accordingly clause~(F2) of the proposition describing the frozen weight
holds unchanged.

\item[(b)] \emph{Oscillation datum.} Assume the increment hypothesis on $\mathcal{Y}(e)$,
the sources of the reading rotation of Lemma~\ref{4.25:lem-reading-rotation} (at far kind
(L): the per-layer bound \eqref{4.25:eq-increment-bound-b}), the base-class hypothesis on
$\mathfrak{u}$, and member~(i) of \eqref{4.25:eq-frozen-collar-weight-a}. These, together with
the bound $\operatorname{pol}(g_{\mathfrak{O}})\le\tfrac16$ on the base rotation
$g_{\mathfrak{O}}$ of the presentation, are the hypotheses under which
Corollary~\ref{4.25:cor-oscillation-datum} is stated; the last of them holds at every later
presentation, the base rotation $g_{\mathfrak{O}}$ being a chart point, so that
$\operatorname{pol}(g_{\mathfrak{O}})<\varepsilon_\star$, where $\varepsilon_\star$ denotes
the bound on $\operatorname{pol}$ at chart points. Under these, the
conclusion of Corollary~\ref{4.25:cor-oscillation-datum},
\eqref{4.25:eq-oscillation-datum-a}, holds at every later presentation of kind (P), (T) and
far kind (L); through Lemma~\ref{4.25:lem-deep-oscillation} it also holds at the deep ones.
Here $\Theta(e):=\sup_{y\in\mathcal{Y}(e)}\theta(y)$. For every weight by which a later bound
multiplies the datum, the product of $\mathfrak{N}^G(e;\mathfrak{O})$ with that weight is at
most $\mathfrak{c}_G\mathfrak{n}^{\circ}(e)\Pi_G(x_{k'})$, by the weighted bound of
Proposition~\ref{4.25:prop-conversion-old-factors}. No oscillation bound is asserted at the pieces
of~(e.iii) below; there the norm is used, on the domain given by the domain statement.

\item[(c)--(d)] \emph{Presentations of kind (P).} This clause concerns the stages
$l'>n_{\mathfrak{a}}+1$ of the site, the operation $O1$, and the operations $O4$, $O4'$
carried out by the healing $\mathbf{R}$ (neither integrates $x$; no operation before $O5$
does), their adjacent first stage included. No
distance is used. At these presentations each $v_e$ enters the anchored moment bound as the
element $c=(v_e,S^+(e),\mathcal{S})$, where
\begin{equation*}
\mathcal{S}:=\emptyset,\qquad
\mathcal{A}_c=\{\square:\tilde\square\cap S^+(e)\ne\emptyset\},\qquad
\mathfrak{N}(c):=\mathfrak{N}^G(e;l').
\end{equation*}
The entry is legitimate, for the following reasons.
\begin{itemize}
\item $v_e$ depends on no integration variable, so the sets of cubes that the anchored moment
bound attaches to the integration variables on which $c$ depends are all empty.
There is no rim anchor, and every element has a non-empty cube set $A$.
\item Let $F$ be the orbit datum $v_e^{\sharp}$, regarded as a function of the interpolation
parameters $(t_A)_A$ of the anchored moment bound, one for each cube set $A$: it is evaluated
at $x$, at the pair interpolated with these parameters, restricted to the domain of $v_e$, and
at the rotation obtained by composing the unitary factor at these parameters with the base
rotation of the presentation; the interpolated pair and the unitary factor are those of the
complex interpolation corollary of the relative small-field machine.
Write $F(0)$ for its value at $t_A=0$ for all $A$. The coefficient
$c_{A,\emptyset}=\int\partial_{t_A}F$ annihilates $F(0)$.
\item The Cauchy estimate of clause~(ii) of the anchored moment bound is therefore applied to
$F-F(0)$, which is at most $\mathfrak{N}^G$ on the polydisc. Hence every element of the
expansion of $c$ has size $\nu$ bounded by
\begin{equation}\label{4.25:eq-frozen-collar-weight-1}
\nu\le\mathfrak{N}^G(e)\prod_{A}K_A\,\varsigma(\square)\,e^{-(\kappa_1-1)|\mathsf{y}|},
\end{equation}
where, in the notation of the anchored moment bound, $K_A$ are the cube constants,
$\varsigma(\square)$ is the polymer weight of the cube $\square$, and $|\mathsf{y}|$ is the
length of the anchoring tree.
\end{itemize}
The sums are taken as follows.
\begin{itemize}
\item[(c1)] In the complexified setting the statement that uses these terms is the block-format
proposition. Its format class contains $\mathsf{q}_W$ and uses it only through clause~(f);
the block and window sums are additive lemmas of that proposition.
\item[(c2)] The anchored block sum (Proposition~\ref{4.25:prop-anchored-block-sum}) holds;
here $S=S^+(e)$ is the declared support of $c$, $\Sigma_c$ is the exponent attached to the
element $c$, and $N$ is the divisor in the size weight $e^{d_k(S)/N}$ of the anchored moment
bound; this $N$ is not the rank of the gauge group:
\begin{equation*}
\sum_{\substack{c\in\bigcup_{\mathfrak{a}'}\mathfrak{Q}(\mathfrak{a}')\\
S\cap\tilde\square_0\ne\emptyset}}
\mathfrak{w}(c)\,e^{\Sigma_c+d_k(S)/N}\le\mathsf{n}_G^{blk}(x_{k'})
:=\mathfrak{c}_GC_\alpha C_{cnt}K_a^{\circ}\,\Pi_G^{blk}(x_{k'}).
\end{equation*}
The quantity $\mathsf{n}_G^{blk}$ is added to the totals of the anchored moment bound.
With $\varepsilon_G:=K_\flat\mathsf{r}V\mathsf{n}_G^{blk}$, the quantity $\varepsilon_G$ is
added to the augmented smallness condition of the anchored moment bound, and $5\varepsilon_G$
to the fourth member of the convergence ceiling.
\end{itemize}

\item[(e)] \emph{The operations $O2$, $O3$, $O5$ of the correlation theorem.} The first part
of the fourth class of stored terms of the correlation theorem receives $\mathfrak{Q}(\mathfrak{a})$
of every live arrest of Gaussian-collar type, with the following data: $S_v:=S^+(e)$;
hypotheses on $\mathcal{Y}(e)$; domain $\mathfrak{K}_G(\tilde S(e))$; and $x$ real and not
integrated before $O5$, where it is integrated.
\begin{itemize}
\item[(e.i)] At kind (T), and at kind (L) with $S^+(e)\cap Z=\emptyset$: the oscillation slot
of the lemma on stored terms of the fourth class receives
$\mathfrak{N}(v):=\mathfrak{N}^G(e;\mathfrak{O})$, $\mathfrak{O}$ the operation of the site.
That slot bounds, on the disc $|\mathsf{t}|<\rho_A$ of
\eqref{4.25:eq-increment-bound-c}, the deviation of the term read at the presented pair from
its value at the base pair; it is filled by Corollary~\ref{4.25:cor-oscillation-datum}, whose
reference is $(P_{\mathfrak{O}},g_{\mathfrak{O}})$. The pile
satisfies
\begin{equation}\label{4.25:eq-frozen-collar-weight-2}
\begin{aligned}
\mathsf{W}_G&:=\sup_\Delta\sum_{e:\,S^+(e)\cap\Delta\ne\emptyset}
4\mathsf{w}^{-1}\mathfrak{N}^G(e)\,e^{(1-\frac12\delta)\kappa d(\bar S_v)}\\
&\le\mathsf{n}_G^J(x_{k'}):=4\mathsf{w}^{-1}\mathfrak{c}_GC_\alpha K_a^{\circ}\,\Pi_G(x_{k'}),
\end{aligned}
\end{equation}
where $\Delta$ ranges over the cells of the multiscale cover, $\mathsf{w}=\mathsf{w}_{\mathfrak{s}}$
is the disc radius of the site, and $\delta$ and $d(\bar S_v)$ are the decay exponent and the
tree length of the lemma on stored terms of the fourth class.
\item[(e.ii)] At kind (L) with $S^+(e)\subset\mathsf{C}_C$: Lemma~\ref{4.25:lem-deep-oscillation}
holds, with $\mathfrak{N}$ read as $\tfrac12\mathfrak{N}^{G,L}$.
\item[(e.iii)] This item covers all other pieces at a site of $\mathbf{R}$: the interior shallow
pieces, which carry no cut parameter $\mathsf{t}$ and are bounded in norm; the pieces whose
unit has a region not contained in the treated class $Z$; the pieces taken at
links two to four; and the pieces meeting the component of a foreign completed chain, which
are bounded there and cease to be stored there. For these,
$\mathfrak{N}(v):=2\nu^{\natural}(e)$, which bounds the term on the domain that the domain
statement supplies (Lemma~\ref{4.25:lem-norm-bound}) and is small by
Lemma~\ref{4.25:lem-non-deep-pieces}, whose proof is incomplete at one step. Their piles are
bounded by ($\alpha$) and ($\beta$) of \eqref{4.25:eq-pile-bound-a}, with $\Pi^{sh}$.
\end{itemize}

\item[(f)] \emph{The block format.} Put $F_A:=\sum_{A_e=A}v_e$, with
\begin{equation*}
A_F:=A,\qquad j_F:=m_0(e),\qquad X_F:=\bigcup_{A_e=A}S^+(e),\qquad \mathcal{P}_F:=\emptyset .
\end{equation*}
The two requirements of the block format on the label, namely $A\in\mathbf{D}_j$ and that $A$
contains a level-$j$ cell contained in the small-field region of level $j$ and a collar cube
of $W$, are carried with $j=j_F$; they concern the form of the label only. The datum
$\mathfrak{N}^A$ is the oscillation datum $\mathfrak{N}^G(e;\mathfrak{O})$.
On non-deep pieces of kind (L) it is the norm $2\nu^{\natural}(e)$, by
Lemma~\ref{4.25:lem-non-deep-pieces}, whose proof is incomplete at one step, and the domain
statement. Elsewhere the norm form is
false at the level of numbers, since $\nu^{\natural}\propto(\log x_{j_W})^{2p_1}$ for a fixed
exponent $p_1$; the
statements that legitimately use $F$ use it only through differences
$F(\text{moved})-F(\text{base})$.
Consider every statement that uses the term, at a site of kind (T), (P) or (L) of step $k'$,
every $\ell$, every $\square'\in\pi_\ell$, and every $r'\le2\kappa+3$. Here $e^{\mathfrak{h}(e)}$,
with $\mathfrak{h}(e)\le c_J|S^+(e)|$, is the volume weight of the block format. For all of
these,
\begin{equation}\label{4.25:eq-frozen-collar-weight-3}
\begin{aligned}
\sum_{\mathfrak{a}'}\sum_{\substack{e:\,m_0(e)=\ell\\ A_e\ni\square'}}
\mathfrak{N}^A(e)\,e^{\mathfrak{h}(e)}e^{r'd_\ell(A_e)}
&\le b_\delta^G\cdot\max\{\Pi_G,\Pi^{sh}\}(x_{k'})\le b_\delta^G,\\
b_\delta^G&:=\max\{\mathfrak{c}_G,2e^{3+2c_{wt}}\}\,C_\beta K_a^{\circ}.
\end{aligned}
\end{equation}
Here $C_\beta$ is the pile constant of clause~$(\beta)$ of Lemma~\ref{4.25:lem-pile-bound},
not the flow constant of Definition~\ref{1.2:def-flow-constants}.
The range $r'\le2\kappa+3$ covers the exponents with which the block format weights the label,
which are at most $2\kappa$ and $\tfrac85\kappa+3$. The constant $b_\delta$ of the block
format is taken to be $b_\delta^G$, a number fixed before $\gamma$.

\item[(g)] \emph{Dissolution.} For every member $v_e$ of $\mathfrak{Q}(\mathfrak{a})$, the
conclusion of the
arrest theorem is clauses (iii)--(iv) of Lemma~\ref{4.25:lem-gaussian-pieces}, whose proof is
incomplete at one step, with the
following inputs: the increment hypothesis on $\mathcal{Y}(e)$ at kinds (P), (T) and far (L);
the inputs of Lemma~\ref{4.25:lem-deep-oscillation} at deep kind (L); and the domain statement
elsewhere. Clauses (g$'$1)--(g$'$5) of
Proposition~\ref{4.25:prop-dissolution}, whose proof is incomplete at one step, hold:
\begin{itemize}
\item at $O5$ the still-stored family is read a last time as in~(e);
\item $x$ is integrated inside the activity of the healing operation;
\item the output of $O5$ is that of the emitting part of the output proposition of the
correlation theorem;
\item the value of the form is read by sign inside $\mathbf{T}'_{k_W}(W)$, in the correlation
theorem;
\item thereafter no member of $\mathfrak{Q}(\mathfrak{a})$ is stored.
\end{itemize}

\item[(h)] \emph{The ceiling.} The ceiling assumed above is a single ceiling in $\gamma$.
In it, $x\ge\gamma^{-2}$ denotes a value of the inverse coupling $x_k=g_k^{-2}$, not the
collar variable of clause~(a).
Each of its members is a supremum over $x\ge\gamma^{-2}$:
\begin{equation}\label{4.25:eq-frozen-collar-weight-a}
\begin{aligned}
&\varpi_G:=\min\Bigl\{\tfrac{1}{20}\varepsilon_{reg},\ \tfrac14\gamma_{reg},\ c_4,\
10^{-2}/K_u,\ \alpha_1^{\mathfrak{u}}\Bigr\},\\
&\text{(i)}\ \ \max\{\tfrac98C_{inc},2\}\,\bar\theta^{\downarrow}(x-2c_2)\le\varpi_G,
\qquad C_{inc}:=2(K_{\mathsf{j}}+2c_\zeta)\hat B_HC'_{age},\\
&\phantom{\text{(i)}\ \ }\tfrac98\,\frac{\varphi_\flat}{2K_Dr^0}\sup_{t\ge x-2c_2}\alpha_*(t)
\le\varpi_G,\qquad 2.5\,K_{\mathsf{j}}B_\sharp C_{\mathsf{a}}\,\alpha_*(x)\le\varpi_G,\\
&\phantom{\text{(i)}\ \ }\bar\gamma_0(x)\le\tfrac14\gamma_{reg},\qquad
BCM\tilde\alpha_0(x)\le a_0,\\
&\phantom{\text{(i)}\ \ }(1+\mathsf{w}_L)\,\mathsf{U}_L(x)\,e^{\delta_\natural}E(x)^{\delta_\natural/4}
\le2\hat B_H\bar\alpha(x),\qquad \mathsf{U}_L(x):=C_LK_{\mathfrak{Q}}\alpha^+(x)/K_u,\\
&\text{(ii)}\ \ \mathsf{n}_G^J(x)\le\mathfrak{s}_G^J,\\
&\text{(iii)}\ \ \Pi_G,\ \Pi_G^{blk},\ \Pi^{sh}\le1,\\
&\text{(iv)}\ \ R(x)\ge L.
\end{aligned}
\end{equation}
Notes on the members:
\begin{itemize}
\item The member $BCM\tilde\alpha_0(x)\le a_0$ is the condition \cite[(2.38)]{B-CMP119}, with
the constants $B$, $C$, $a_0$ and the radius $\tilde\alpha_0$ of that display, and $M$ as
above.
\item The local number $c_\chi$ of the cut is contained in $\mathsf{U}_L$.
\item $\delta_\natural$ is the decay exponent $\delta_{\mathfrak{s}}$ of
Lemma~\ref{4.25:lem-increment-bound} at far sites of kind (L).
\item $\alpha_1^{\mathfrak{u}}$ is the base radius of the lemma on the unitary factor map
$\mathfrak{u}$, and $c_4$ is the further constant entering, with $10^{-2}/K_u$ and
$\alpha_1^{\mathfrak{u}}$, the radius $\varrho_0:=\min\{10^{-2}/K_u,c_4,\alpha_1^{\mathfrak{u}}\}$
of Lemma~\ref{4.25:lem-reading-rotation}; $C_{\mathsf{a}}$ is the constant of
the amplitude bound among the three-field
bounds; $\bar\gamma_0$ is the envelope of the current radii $\gamma_0^{(k)}$; $\mathsf{w}_L$
is the disc radius at large-field sites; $C_L$, $K_{\mathfrak{Q}}$ and $\alpha^+$ are the
constants and the radius from which the site unit $\mathsf{U}_L$ is built.
\item $\mathfrak{s}_G^J$ is the allowance of the correlation theorem for stored terms of the
fourth class.
\item Under member~(iii) the quantity $\mathsf{n}_G^{blk}$ is also added to the bracket of the
augmented smallness condition of the anchored moment bound.
\item $\Pi_G^{blk}$ contains the square of $\mathsf{l}^{\sharp}$.
\item In member~(iv), $R(x)$ denotes the separation integer $R_k$ of the run at inverse
coupling $x_k=x$; it is not the ball radius $R$ of the observables.
\item Member~(iv) is the window floor of the statement on the levels met at a response cube,
already imposed by the ceilings of the correlation theorem; it is not new in kind.
\item By Remark~\ref{4.25:rem-tube-constant} and \eqref{4.25:eq-tube-constant-a},
$C_{inc}\ge300$. The maximum in the first member of~(i) therefore equals
$\tfrac98C_{inc}$, and that member reads
$\tfrac98C_{inc}\bar\theta^{\downarrow}(x-2c_2)\le\varpi_G$.
\end{itemize}
General properties of the ceiling:
\begin{itemize}
\item The left members are proportional to $x^{-1/2}$ times a polylogarithm, or to $E(x)^c$,
and tend to zero.
\item No constant depends on $K$, $k$, $k'-k_W$, the ages of large-field regions, the number
of window cells, the number or nesting depth of arrests, or the volume.
\item The floors \eqref{4.25:eq-regularity-class-b} and \eqref{4.25:eq-collar-weight-b} are
one-sided and are imposed after $L$.
\item The order in which the constants are chosen (Notation~\ref{2.1:not-prefix}) is
unchanged: after $L$ come, in this order, the class $\mathfrak{M}$ of admissible weights, then
$\kappa$, $\kappa_1$, $M_1$, $M$, then the pair $(C_\Sigma,p_\Sigma)$, then the number
$\mathfrak{n}^{kept}$, and $\gamma$ is chosen last.
\item $C_{cnt}$, $C'_{cnt}$, $C'_{age}$, $c_d$, $c_\zeta$, $c_\chi$, $C_{inc}$,
$K_{\mathsf{j}}$, $\alpha_1^{\mathfrak{u}}$, $C_G$, $\mathfrak{c}_G$, $b_\delta^G$ are numbers fixed
before $\gamma$.
\end{itemize}
\end{itemize}
\end{theorem}

\begin{proof}
We prove the clauses in order.

\emph{Clause (a).} Membership of $\mathfrak{Q}(\mathfrak{a})$ in the family stored under
$\mathfrak{a}$ is the storage definition of the arrest theorem for arrests of Gaussian-collar
type. The domain $\mathfrak{K}_G(\tilde S(e))$ of the storage datum is the class on which
clause~(iii) of the lemma named as the input of Lemma~\ref{4.25:lem-gaussian-pieces}(i) holds;
the proof of Lemma~\ref{4.25:lem-gaussian-pieces} itself is incomplete at one step. The
remaining entries of the datum are supplied as follows.
\begin{itemize}
\item By Lemma~\ref{4.25:lem-gaussian-pieces}, whose proof is incomplete at one step, each
$v_e$ is a well-defined quadratic form on $\mathfrak{K}_G(\tilde S(e))$, jointly invariant,
with orbit datum $v_e^{\sharp}$. This gives the domain and the orbit datum.
\item By clause (g$'$1) of Proposition~\ref{4.25:prop-dissolution}, whose proof is incomplete
at one step, the variable $x$ is real, in the box $|x|\le\mathfrak{b}_x$, and not integrated
before $O5$.
\end{itemize}

\emph{Clause (b).} Fix a later presentation $\mathfrak{O}$ of kind (P), (T) or far kind (L),
and a piece $e\in\mathfrak{Q}(\mathfrak{a})$. We check the hypotheses of
Corollary~\ref{4.25:cor-oscillation-datum}.
\begin{itemize}
\item The increment hypothesis on $\mathcal{Y}(e)$ is assumed.
\item The sources of the rotation are given by Lemma~\ref{4.25:lem-reading-rotation}. That
lemma applies because the base-class hypothesis on $\mathfrak{u}$ and member~(i) of
\eqref{4.25:eq-frozen-collar-weight-a} are assumed. At far kind (L), the lemma is applied per
link and per layer, from the per-layer bound \eqref{4.25:eq-increment-bound-b}, which is
assumed there. Its conclusion is \eqref{4.25:eq-reading-rotation-a}.
\item The base-class hypothesis and member~(i) of the ceiling are assumed directly.
\item The bound $\operatorname{pol}(g_{\mathfrak{O}})\le\tfrac16$ on the base rotation holds
because $g_{\mathfrak{O}}$ is a chart point, as stated in clause~(b).
\end{itemize}
Lemma~\ref{4.25:lem-domination-old-blocks}(ii) bounds every block size in old units on a block
$y_0$ with $\Delta(y_0)\subset S^+(e)$ by the size in current units on the current block
$B^{top}(y_0)\in\mathcal{Y}(e)$. This is the form in which
Lemma~\ref{4.25:lem-gaussian-pieces}(iv), whose proof is incomplete at one step, receives the
increment.

Corollary~\ref{4.25:cor-oscillation-datum} therefore gives \eqref{4.25:eq-oscillation-datum-a},
that is, the bound by $\mathfrak{N}^G(e;\mathfrak{O})=C_G\nu^{\natural}(e)\Theta(e)$ together
with holomorphy in the family parameter.

At a deep presentation of kind (L), Lemma~\ref{4.25:lem-deep-oscillation} applies, its inputs
being assumed. It gives the bound \eqref{4.25:eq-deep-oscillation-a}, with $\mathfrak{N}$ read
as $\tfrac12\mathfrak{N}^{G,L}(e)$ and with $\Theta^L(e)\le\Theta(e)$.

Finally, Proposition~\ref{4.25:prop-conversion-old-factors}
gives $\varrho^2\Theta(e)e^{-n(e)}\le\Pi_G(x_{k'})$.
Its weighted bound gives, for every weight by which a later bound multiplies it, the bound of
$\mathfrak{N}^G(e;\mathfrak{O})$
times that weight by $\mathfrak{c}_G\mathfrak{n}^{\circ}(e)\Pi_G(x_{k'})$.

\emph{Clauses (c)--(d).} Clauses (i) and (iv) of the anchored moment bound are applied to the
element $c=(v_e,S^+(e),\emptyset)$. For this element the rim-anchor sets are empty, as stated,
because $v_e$ depends on no integration variable. Clause~(ii) of the anchored moment bound, its
Cauchy estimate, is applied to $F-F(0)$. This is legitimate because $c_{A,\emptyset}=\int\partial_{t_A}F$
does not see $F(0)$. The estimate is linear in the function to which it is applied, and by
clause~(b) the modulus of $F-F(0)$ on the polydisc is at most $\mathfrak{N}^G(e;l')$. This
gives \eqref{4.25:eq-frozen-collar-weight-1}.

The sum (c2) is Proposition~\ref{4.25:prop-anchored-block-sum}. That proposition is proved from
three statements: clause~($\alpha$) of Lemma~\ref{4.25:lem-pile-bound} at
$\mathcal{A}=\mathcal{A}(\tilde\square_0)$; Lemma~\ref{4.25:lem-cell-count}; and
Proposition~\ref{4.25:prop-response-cube-conversion}. The hypotheses of
Lemma~\ref{4.25:lem-cell-count} are the statement on the levels met at a response cube and the
condition at $\kappa(\square_0)$ under which that lemma is stated, and both are assumed at every
presentation of kind (P).

\emph{Clause (e).} In case~(e.i) the lemma of the correlation theorem on stored terms of the
fourth class requires five pieces of information about each term: its domain, holomorphy,
invariance, the datum $\mathfrak{N}$, and a pile bound. We supply them, and then treat
cases~(e.ii) and~(e.iii).
\begin{itemize}
\item In case~(e.i):
\begin{itemize}
\item Domain, holomorphy and invariance are those of Lemma~\ref{4.25:lem-gaussian-pieces},
whose proof is incomplete at one step.
\item The datum is clause~(b).
\item The pile \eqref{4.25:eq-frozen-collar-weight-2} follows from the weight bound of
clause~(b), which turns each summand into at most
$4\mathsf{w}^{-1}\mathfrak{c}_G\mathfrak{n}^{\circ}(e)\Pi_G(x_{k'})$. It also uses
clause~($\alpha$) of \eqref{4.25:eq-pile-bound-a} with $|\mathcal{A}|=1$, $\Delta$ being a cell
of the multiscale cover, which bounds the remaining sum of $\mathfrak{n}^{\circ}(e)$ by
$C_\alpha K_a^{\circ}$.
\end{itemize}
\item In case~(e.ii), Lemma~\ref{4.25:lem-deep-oscillation} supplies clauses (a), (b) and (e)
of the three-field bounds for $v_e$, with $\mathfrak{N}$ read as
$\tfrac12\mathfrak{N}^{G,L}(e)$.
\item In case~(e.iii), the norm $2\nu^{\natural}(e)$ and its weighted form are supplied by
Lemma~\ref{4.25:lem-non-deep-pieces}, whose proof is incomplete at one step. The domain is the
one that the domain statement supplies, on which Lemma~\ref{4.25:lem-norm-bound} gives the bound
$|v_e^{\sharp}|\le2\nu^{\natural}(e)$. The piles follow from ($\alpha$) and ($\beta$) of
\eqref{4.25:eq-pile-bound-a}.
\end{itemize}

\emph{Clause (f).} We treat the two kinds of piece separately.
\begin{itemize}
\item For a piece carrying the oscillation datum, clause~(b) bounds
$\mathfrak{N}^A(e)e^{\mathfrak{h}(e)}e^{r'd_\ell(A_e)}$ by
$\mathfrak{c}_G\mathfrak{n}^{\circ}(e)\Pi_G(x_{k'})$. This uses the weighted bound of
Proposition~\ref{4.25:prop-conversion-old-factors}, whose weights include $e^{\mathfrak{h}(e)}$ and
$e^{r'd_\ell(A_e)}$ with $r'\le2\kappa+3$.
\item For a non-deep piece of kind (L), Lemma~\ref{4.25:lem-non-deep-pieces}, whose proof is
incomplete at one step, bounds the same product by
$2e^{3+2c_{wt}}\mathfrak{n}^{\circ}(e)\Pi^{sh}(x_{k'})$.
\end{itemize}
In both cases the summand is at most
\begin{equation*}
\max\{\mathfrak{c}_G,2e^{3+2c_{wt}}\}\,\mathfrak{n}^{\circ}(e)\,\max\{\Pi_G,\Pi^{sh}\}(x_{k'}).
\end{equation*}
Clause~($\beta$) of \eqref{4.25:eq-pile-bound-a} bounds the sum of $\mathfrak{n}^{\circ}(e)$
over $\mathfrak{a}'$ and over $e$ with $m_0(e)=\ell$ and $A_e\ni\square'$ by $C_\beta K_a^{\circ}$;
here $A_e\ni\square'$ means that $S^+(e)$ meets $\square'$. This gives the first inequality of
\eqref{4.25:eq-frozen-collar-weight-3}. The second follows from
$\Pi_G,\Pi^{sh}\le1$, which is member~(iii) of \eqref{4.25:eq-frozen-collar-weight-a}.

\emph{Clause (g).} The identification of the conclusion of the arrest theorem for
every member $v_e$ of $\mathfrak{Q}(\mathfrak{a})$ with clauses (iii)--(iv) of
Lemma~\ref{4.25:lem-gaussian-pieces} follows from that lemma, whose proof is incomplete at one
step. It is read under the inputs listed, which are assumed at the respective presentations.
The life of $\mathfrak{a}$ is given by the storage definitions of the arrest theorem and by its
clauses on the creation, life and dissolution of an arrest: $\mathfrak{a}$ lives from its
creation until the first later completed $\mathbf{R}$ on the component containing $W$. The
statements on the last reading, the integration of $x$, the output, the reading by sign and the
absence of stored members afterwards are clauses (g$'$1)--(g$'$5) of
Proposition~\ref{4.25:prop-dissolution}, whose proof is incomplete at one step.

\emph{Clause (h).} This clause collects the members of the ceiling under which the preceding
clauses were obtained.
\begin{itemize}
\item Member~(i) is the hypothesis of Lemma~\ref{4.25:lem-reading-rotation} and
Corollary~\ref{4.25:cor-oscillation-datum} used in~(b), and of
Lemma~\ref{4.25:lem-deep-oscillation} used in (b) and~(e.ii).
\item Member~(ii) bounds the pile \eqref{4.25:eq-frozen-collar-weight-2} of~(e.i) by the
allowance of the correlation theorem.
\item Member~(iii) was used in~(c2) and in~(f).
\item Member~(iv) is the window floor of the statement on the levels met at a response cube
entering Lemma~\ref{4.25:lem-cell-count}.
\end{itemize}
The statements on the size of the left members, on the independence of the constants, and on
the order of choice are read off from the definitions of the constants in
Notations~\ref{4.25:not-regularity-class} and~\ref{4.25:not-flow-inequalities}, and in the
statements cited above.
\end{proof}

\begin{definition}[Stages, regions, and the frozen reading on excluded components]
\label{4.26:def-stages-regions}
Throughout this definition the level $k$ is fixed, and we put
\begin{equation*}
\eta=L^{-k},\qquad h=k-N,\qquad k_0=k-N_0,
\end{equation*}
where $N$ is the number of stages of one application of Ba{\l}aban's large-field operation
$\mathbf{R}$ of \cite{B-CMP122a}. In this definition and in the statements that use it, $N$
denotes this number of stages and not the rank of the gauge group. The integer $N_0$ is
Ba{\l}aban's integer of \cite{B-CMP122a} which fixes the weights of \cite[(1.64)]{B-CMP122a}:
the weight $L_0^{2\max\{0,m-k_0\}}$ at level $m$ equals $1$ exactly when $m\le k_0=k-N_0$, so
that the weights exceed $1$ precisely at the $N_0$ levels $k_0<m\le k$. We write $L_0$
for the base of the weights of \cite{B-CMP122a}, $w_\ast=L_0^{2N_0}$, and
$\beta:=\tfrac15$, the value of Ba{\l}aban's constant $\beta$ in \cite{B-CMP122a}.
In this section the letters $h$, $L_0$, $\mathfrak{T}$, $\beta$, $g(\cdot)$, $\rho^{(l)}_m$ and
$R^{(l)}$ have the meanings fixed in this
definition and not those of the same letters in the glossary (Notation~\ref{0.2:not-glossary}).

\emph{Stages.} For $l\in\{1,\dots,N\}$ write also $n=l-1$, so that stage $l$ is also called
stage $n+1$; the letter $n$ is used in this sense only in the superscript of
$\mathfrak{N}^{(n)}$ below. Stage $l$ integrates the field $\phi=(A_j,B_j)$ of level
\begin{equation}\label{4.26:eq-stages-regions-1}
j=\ell_l:=h+l-1,
\end{equation}
the variable $B_j$ on the collar bonds and the variable $A_j$ on the rim.

\emph{The class-wide region and the run.} Let $Z$ be the union of all components of the
classes \textup{(i)} and \textup{(ii)} of \cite{B-CMP122a} which enter stage $1$ at the present
application of $\mathbf{R}$. We assume moreover the following nesting condition: during the $N$
steps of this application no components are joined inside $Z$. In \cite{B-CMP122a} a
decomposition of unity
$1=\chi'_j+(1-\chi'_j)$ is introduced for each component of $Z$, the components carrying the
large-field functions $1-\chi'_j$ are excluded from $Z$, and the remaining region is denoted
again by $Z$. From here on, in the stages of the present application of $\mathbf{R}$, the symbol
$Z$ denotes the class-wide union just defined and is not re-denoted after exclusions;
we write $Z^{[l]}$ for the union of the components of $Z$ not excluded at any stage
$\le l-1$, the components still in the run at stage $l$. Here the \emph{run} of a component at
an application of $\mathbf{R}$ is the succession of stages of that application performed on it,
and a component is \emph{in the run} at stage $l$ if stage $l$ is performed on it.

\emph{The regions of a stage.} With $j=\ell_l$ as in \eqref{4.26:eq-stages-regions-1} and
with Ba{\l}aban's regions $Z_{j+1}$, $Z''_{j+1}$ and $\Omega_{j+1}$ of \cite{B-CMP122a}, put
\begin{align*}
R^{(l)}&:=R_{l-1}:=(Z_{j+1}\setminus Z''_{j+1})\cap Z^{[l]},\\
\mathsf{R}_j&:=Z_{j+1}\cap\Omega_{j+1}\cap Z^{[l]} .
\end{align*}
The region $R^{(l)}$ contains the collar and the rim of stage $l$. We write
$\mathcal{C}^{(l)}:=\mathcal{C}_j\subset R^{(l)}$ for the set of level-$j$ bonds of the collar;
the variable $b=B'_j$ of stage $l$ lives on $\mathcal{C}^{(l)}$.

\emph{The untouched region and the region on the run.} In the setting of
\cite[(1.14)--(1.17)]{B-CMP122a}, put
\begin{equation}\label{4.26:eq-stages-regions-2}
\mathfrak{T}:=\Omega_k\cup Z^c,\qquad Z^{\mathrm{on}}:=Z\cap\Omega_k^c .
\end{equation}
Since $Z$ is the class-wide union, no stage of the present application of $\mathbf{R}$ changes
the determining set on $\mathfrak{T}$: the operation $\mathbf{T}_k(Z_k\cap Z^c)$ is left
unchanged by the stages, as stated in \cite{B-CMP122a}.

\emph{Excluded components.} In the terminology of the definition of arrested components, the
run of a component $W$ of $Z\cap\Omega_k^c$ may end, after $n_{\mathfrak{a}}$ performed stages,
on one of the
four large-field branches B1--B4; the component is then called \emph{arrested}, and the
\emph{arrest} is the quadruple $\mathfrak{a}=(W_{\mathfrak{a}},k_a,n^+_a,\text{branch})$ formed
by the component $W_{\mathfrak{a}}=W$, the level $k_a$ of the application of $\mathbf{R}$ at
which its run ended, the number $n^+_a$ of stages whose clause sets are in use on $W$ (the
number $n_{\mathfrak{a}}$ of performed stages on the branches B1--B3, and $n_{\mathfrak{a}}+1$
on the branch B4), and the
branch. The subscripts $a$ in $k_a$, $n^+_a$ (and in $h_a$, $k_0^a$, $N_0^a$ below) refer to
the arrest $\mathfrak{a}$; the index $a$ of the shell radii $\mathbf{r}_{a,m}$, of the radii
$r_a(y)$ and of the threshold $\theta_{\mathrm{src},a}$ below is the index of the shell radii
and does not refer to an arrest.
For an arrest at the present application, $k_a=k$; we write $n^+:=n^+_a$
and put $j^+:=h+n^+$. Such a $W$ is called an \emph{excluded component}; it is \emph{live} (it
still carries its large-field factor at level $k$) but out of the run for $l>n^+$, and
$W\subset Z^{\mathrm{on}}$. For all $l>n^+$ the determining set on $W$ is frozen:
\begin{equation}\label{4.26:eq-stages-regions-3}
\mathfrak{B}\restriction W:=\mathbb{B}^{(n^+)}_k\restriction W \qquad (l>n^+),
\end{equation}
that is, the level, the weight, the multiplicity and the cube family at every $x\in W$ are
those which the attempt space $\mathfrak{S}^{(k)}_{n^+}$ of stage $n^+$ at level $k$ (the
attempt space whose clauses form the clause set of that stage) assigns. In \cite{B-CMP122a} it
is stated of the components excluded at the first decomposition that they are not changed by
the subsequent operations; for every branch, \eqref{4.26:eq-stages-regions-3} is the freeze of
the structure of an arrested component at the stage of its arrest, given by the statements on
arrested components for the later decompositions of \cite{B-CMP122a}.

On an excluded component the clause set of stage $l>n^+$ is the arrested one: for every set
$Y$,
\begin{equation*}
\mathfrak{L}_l\restriction(Y\cap W):=\text{the clauses on $W$ of the arrest-stage set
$\mathfrak{S}^{(k),\mathrm{arrest}}_l$},
\end{equation*}
the counterpart at an arrest of the attempt space $\mathfrak{S}^{(k)}_l$. Its
clauses are taken on the frozen structure \eqref{4.26:eq-stages-regions-3}, with the factor
$\varphi^{(l),(k)}_{\ell_{\mathfrak{a}}(x)}$, $\ell_{\mathfrak{a}}(x)$ the frozen level at $x$;
this factor carries the stage-dependent subtraction of \cite{B-CMP122a},
\begin{equation*}
\beta\sum_{i=h+1}^{h+l}2^{-|\ell_{\mathfrak{a}}(x)-i|},
\end{equation*}
and pieces of the same class need no additional correction term $\sigma_{\mathfrak{a}}$ inside
the run. The range of the arrest is $I(\mathfrak{a}):=\{h+1,\dots,k+1\}$.

\emph{Slots of an excluded component.} Every slot of $W$ which no performed stage integrates
--- the slots of levels $\ge j^+$, and on the branch B4 also the field $\phi_W$ of level
$j^+-1$, which is parametrised by \cite[(1.53)]{B-CMP122a} and never integrated --- remains
exterior at every later stage: it is a coordinate of the parameter $\lambda$ of the slot domain
$\mathfrak{D}_{\mathcal{S}}$, and its strong weight is carried in the bookkeeping quantity
$\mathfrak{N}^{\natural}$ of the elements.
This is consistent with the stage integration: the bound for the integral of a stage (the
stage-integration lemma) holds for every value of $\lambda$; its corollary, which treats the
exterior coordinates of levels $j'>j$, $j=\ell_l$, holds with this set of coordinates enlarged
by the slots just described; and the Gaussian integral of a stage $l>n^+$ lives on
$Z^{[l]}$ and does not depend on $\lambda$. For the frozen terms $\mathbb{C}^{(i)}$,
$i\le n_{\mathfrak{a}}$, of the performed stages on $W$,
every element quantity $\nu^{\natural}(e)$ of that bookkeeping is a supremum over
$\mathfrak{D}_{\mathcal{S}}$, so that no estimate
changes; every other use of these slots is governed by the bound for the frozen terms on
the branch B4.

\emph{Structure and clause set off the excluded components.} Off the excluded components we put
$\mathfrak{B}:=\mathbb{B}^{(l-1)}_k$, and the clause set of stage $l$ is that of
\cite[(1.64)--(1.69)]{B-CMP122a} at stage $l$, in the convention $(\Sigma^{\ast})$ for the
derived entries. On the older live pieces, that is, on the components $W_{\mathfrak{a}}$ of live
arrests $\mathfrak{a}$ with $k_a<k$, the frozen structure and the arrested clauses are
used; they are fixed below under the heading \emph{older live pieces}.

\emph{Shell index and radii on $\mathfrak{T}$.} On $\mathfrak{T}$ the shell index $m(y)$
satisfies $m(y)\le k$, and the radii $r_a(y):=\mathbf{r}_{a,m(y)}$, built from the shell radii
$\mathbf{r}_{a,m}$, do not depend on the stage. On an older live piece lying in
$\mathfrak{T}$ one has $m=\ell_{\mathfrak{a}}$ and $r_a=\omega\,\mathbf{r}_{a,\ell_{\mathfrak{a}}}$
(see below), which do not depend on the stage either.

\emph{One-step rooms.} For every integer $m$ and every stage $l$ put
\begin{equation}\label{4.26:eq-stages-regions-4}
\rho^{(l)}_m:=\beta\,2^{-|m-\ell_l-1|}.
\end{equation}
Then
\begin{equation}\label{4.26:eq-stages-regions-5}
\sum_{l=1}^{N}\rho^{(l)}_m\le 3\beta\quad\text{for every $m$},\qquad
\sum_{l=1}^{N}\rho^{(l)}_k<2\beta .
\end{equation}

\emph{Strata count.} On a large-field region $W'\subset Z_k\setminus Z$ which the stages of the
present application of $\mathbf{R}$ leave untouched we put $n_W(y):=k-1-m(y)$; on $\Omega_k$
we put $n_W(y):=0$.

\emph{The class-wide choice of $Z$.} On an excluded component $W$ and for $l>n^+$, the region
$R^{(l)}$, the structure $\mathfrak{B}$ and the clause set $\mathfrak{L}_l$ are the frozen ones
defined above. Let
$\mathsf{m}$ denote the level of a frame structure and $\mathsf{m}^{\circ}$ the level of its
parent structure in the construction of the position-dependent radii $r_e(m,l')$ attached to
the cover of the stage, which is used later in this section. The frame
condition $\mathsf{m}^{\circ}\le\mathsf{m}$ holds on $W$, because the structure of $W$ is
frozen at stage $n^+\le l-1$ and stages only coarsen. If $Z$ were replaced by the region which
remains after the exclusions, an excluded $W$ would lie in $\mathfrak{T}$, although the stages
$\le n^+$ have changed the levels on $W$ and $\mathcal{C}^{(l)}\cap W\neq\emptyset$.
With the class-wide $Z$ the region $\mathfrak{T}$ is untouched by every stage, and none of the
statements of the following sections at points of $\mathfrak{T}$ changes.

\emph{Older live pieces.} The statements on arrested components (the freeze of the structure
and the arrest-stage sets) are in force. For an arrest $\mathfrak{a}$ as defined above we
write $h_a$, $k_0^a$, $N_0^a$ for the integers $h$, $k_0$, $N_0$ of the application of
$\mathbf{R}$ at level $k_a$. The region $\mathfrak{T}$ is untouched by the
stages of the present application of $\mathbf{R}$, but it is not ordinary: it contains the
component $W_{\mathfrak{a}}$ of every arrest $\mathfrak{a}$ which is live at level $k$, with
$k_a<k$, and whose component does not lie in $Z$; the live older pieces whose components lie
in $Z$ are treated in case (a) below. In \cite{B-CMP122a} it is stated, at
the first decomposition, that regions in which
new large fields have been introduced do not satisfy condition \textup{(ii)}, are excluded from
$Z$, are not changed by the subsequent operations, and are treated in the same way as the
remaining large-field regions; for $n^+\ge1$ the corresponding statement is also in
\cite{B-CMP122a}. In \cite{B-CMP122b} it is stated that, if a new large field was introduced in
the preparatory operations, then, in the notation there, $K=R_{j+1}$ for $Z$.
A component joined to others remains outside $Z$ however many
levels below $k$ it was arrested: nothing
below uses $k-k_a$, $N$, or a width of the shells of $W_{\mathfrak{a}}$.

At a point $y\in W_{\mathfrak{a}}$ every object of this section and of the following sections
is evaluated on the frozen structure of the earliest live arrest $\mathfrak{a}$ whose component
contains $y$, namely the structure which the freeze \eqref{4.26:eq-stages-regions-3} of the
application of $\mathbf{R}$ at level $k_a$ fixes on $W_{\mathfrak{a}}$. Writing
$\ell_{\mathfrak{a}}(y)$ and
$w_{\mathfrak{a}}(y)$ for the level and the weight of that frozen structure at $y$, we put
\begin{equation}\label{4.26:eq-stages-regions-6}
m(y):=\ell_{\mathfrak{a}}(y),\qquad \omega(y):=w_{\mathfrak{a}}(y).
\end{equation}
The constant $c$ of \cite[(1.68)]{B-CMP122a} takes on $W_{\mathfrak{a}}$ the value
$c_{\mathfrak{a}}=1$: the value $c=3$ appears there only for $m=k$ and on $\Omega_k\cap X$,
with $X$ the region of that display; read at the level $k_a$ of the arrest, this set is
contained in $\Omega_{k_a}$, whereas $W_{\mathfrak{a}}\subset\Omega^c_{k_a}$.
The radii at $y$ are
\begin{align*}
r_a(y)&:=\omega(y)\,\mathbf{r}_{a,\ell_{\mathfrak{a}}(y)},\\
r_e^{(p)}(y)&:=\omega(y)\,r_e\bigl(\ell_{\mathfrak{a}}(y),\min\{p,\lambda_{\mathfrak{s}}(y)\}\bigr),
\end{align*}
the second at the triple $(p,X',\mathfrak{s})$ of the convention $(\Sigma^{\ast})$ in use, where
$r_e(m,l')$ is the position-dependent radius of the construction recalled above and
$\lambda_{\mathfrak{s}}(y)$ is the level
at $y$ of the structure $\mathfrak{s}$ of the triple. One has
$\lambda_{\mathfrak{s}}(y)=\ell_{\mathfrak{a}}(y)$ at the frozen structure and at every
structure of the present application of $\mathbf{R}$; $\lambda_{\mathfrak{s}}(y)<
\ell_{\mathfrak{a}}(y)$ occurs only at a retained restriction of $\mathbb{B}^{(n')}_{k_a}$ in
the version $(\Sigma^{\ast})^{\mathfrak{a}}$ of that convention on the live piece, whose level
at $y$ a later stage of the run of the piece has
raised. The source threshold is
\begin{equation*}
\theta_{\mathrm{src},a}(y):=F^{(k)}(y)\,r_a(y),
\end{equation*}
with $F^{(k)}$ the source factor of the arrested clause set and with the range
$I(\mathfrak{a})=\{h_a+1,\dots,k_a+1\}$. Finally we put
\begin{equation}\label{4.26:eq-stages-regions-7}
n_W(y):=k-1-\ell_{\mathfrak{a}}(y),\qquad
\mathfrak{c}(y):=\ell_{\mathfrak{a}}(y)-\tfrac12\log_{L_0}\omega(y),
\end{equation}
so that
\begin{equation}\label{4.26:eq-stages-regions-8}
\omega(y)=L_0^{2(\ell_{\mathfrak{a}}(y)-\mathfrak{c}(y))}.
\end{equation}
By the weight table \eqref{4.26:eq-frozen-structure-a} (level $m$ with weight
$L_0^{2\max\{0,m-k_0^a\}}$ on the shells of $\Omega''$, as in \cite[(1.64)]{B-CMP122a}; level
$j^+_a:=h_a+n^+_a$ with weight $L_0^{2(j^+_a-m)}$ on the old shell $m$, as in
\cite[(1.66)]{B-CMP122a}), one
has $\mathfrak{c}=k_0^a$, respectively $\mathfrak{c}=m$, where $\omega>1$, and
$\mathfrak{c}=\ell_{\mathfrak{a}}$ where $\omega=1$. The point $y$ is called \emph{lifted} if
$g(y):=\ell_{\mathfrak{a}}(y)-\mathfrak{c}(y)\ge1$.

For a term $T$ with level $\ell_T$, the radii $\mathbf{r}^T_a(y)$ and $\mathbf{r}^{T,(p)}_a(y)$
are the clause radii of $T$ itself, in the sense of the definition of the clause radii of a
term. If the clauses of $T$
are those of an arrest-stage set --- that is, $T$ is a frozen term of the arrest or a
term born at a later level $k'>k_a$ --- the radius carries the weight $\omega(y)$ of $T$
itself. If the clauses of $T$
were fixed before the arrest (the terms of \cite{B-CMP119} of level $\le k_a$), the radius is
unweighted and of a level $\le\min\{\ell_T,\mathfrak{c}(y)\}$.

Two locations of an older live piece are distinguished.
\begin{itemize}
\item[(a)] $W_{\mathfrak{a}}\subset Z$. Then the piece lies in the core, $h\ge k_a$, by the
lemma placing an arrested component that lies in $Z$ in the core, and $W_{\mathfrak{a}}$ is a
location of $Z^{\mathrm{on}}$ with the clauses of the arrest-stage set
$\mathfrak{S}^{(k),\mathrm{arrest}}_l$, as on an excluded component: $\sigma_{\mathfrak{a}}$
does not depend on $l$. There the fat factor of $\mathfrak{N}^{(n)}$ is
$5\,\mathsf{w}^{\mathrm{on}}(x)=5\max\{w_\ast,\omega(x)\}$, taken on the unweighted radius
$\mathbf{r}_{a,\ell_{\mathfrak{a}}(x)}$, so that only one weight enters. The composite clauses of
$\mathfrak{N}^{(n)}$ at $Z^{\mathrm{on}}$ are those of the clause sets, on the triples of
$(\Sigma^{\ast})^{\mathfrak{a}}$: the arrest-stage set
$\mathfrak{S}^{(k_a),\mathrm{arrest}}_{n'}$ of a stage $n'$ of the application of $\mathbf{R}$ at
level $k_a$ is
taken at the birth radius of $\mathfrak{s}$ and with its threshold at $\lambda_{\mathfrak{s}}$;
the clause set of every later density --- the present $\mathfrak{L}_l$ included --- is taken at
the frozen pair $(m,\omega)$, with the block side of the construction of the radii
$r_e(m,l')$ and with its own
threshold at $m$. At boxes which meet the lifted set of such a piece, the near-part estimate is
not the one of this section but its version for pieces in the core, with the accompanying
lemmas on $\mathfrak{c}$, on the weight, on the distance function and on the cover sums at
points of $Z^{\mathrm{on}}$.
\item[(b)] $W_{\mathfrak{a}}\subset Z^c$. Then $W_{\mathfrak{a}}$ is a location of
$\mathfrak{T}$.
\end{itemize}
\end{definition}

\begin{proof}[Verification of the elementary assertions in
Definition~\ref{4.26:def-stages-regions}]
\emph{The partition into $\mathfrak{T}$ and $Z^{\mathrm{on}}$; excluded components.} By
\eqref{4.26:eq-stages-regions-2}, the complement of $Z^{\mathrm{on}}=Z\cap\Omega_k^c$ is
$Z^c\cup\Omega_k=\mathfrak{T}$; hence $\mathfrak{T}$ and $Z^{\mathrm{on}}$ are disjoint and
their union is the whole lattice. An excluded component $W$ is by definition a component of
$Z\cap\Omega_k^c=Z^{\mathrm{on}}$, so $W\subset Z^{\mathrm{on}}$, and therefore
$W\cap\Omega_k=\emptyset$.

\emph{The two sums \eqref{4.26:eq-stages-regions-5}.} By \eqref{4.26:eq-stages-regions-1},
$\ell_l=h+l-1$, so for fixed $m$ the exponent in \eqref{4.26:eq-stages-regions-4} is
$|m-\ell_l-1|=|m-h-l|$. The map $l\mapsto m-h-l$ is injective on $\{1,\dots,N\}$, because
the integers $\ell_l$ are distinct. Hence the numbers $m-h-l$, $1\le l\le N$, are distinct
integers, and
\begin{equation*}
\sum_{l=1}^N 2^{-|m-\ell_l-1|}\le\sum_{q\in\mathbb{Z}}2^{-|q|}
=1+2\sum_{q\ge1}2^{-q}=3 .
\end{equation*}
Multiplying by $\beta$ gives the first bound. At $m=k$ we have $k-h-l=N-l$, which runs through
$\{0,1,\dots,N-1\}$ as $l$ runs through $\{1,\dots,N\}$; all these exponents are
non-negative, so
\begin{equation*}
\sum_{l=1}^N 2^{-|k-\ell_l-1|}=\sum_{q=0}^{N-1}2^{-q}=2-2^{1-N}<2,
\end{equation*}
and multiplying by $\beta$ gives the second bound.

\emph{The identity \eqref{4.26:eq-stages-regions-8} and the values of $\mathfrak{c}$.} By
\eqref{4.26:eq-stages-regions-7},
$2\bigl(\ell_{\mathfrak{a}}(y)-\mathfrak{c}(y)\bigr)=\log_{L_0}\omega(y)$, and exponentiating to
the base $L_0$ gives \eqref{4.26:eq-stages-regions-8}. On a shell of $\Omega''$ of the weight
table \eqref{4.26:eq-frozen-structure-a} the level is $\ell_{\mathfrak{a}}=m$ and the weight is
$L_0^{2\max\{0,m-k_0^a\}}$. Where $\omega>1$ we have $m>k_0^a$, so
$\tfrac12\log_{L_0}\omega=m-k_0^a$ and $\mathfrak{c}=m-(m-k_0^a)=k_0^a$. On the old shell $m$
the level is $\ell_{\mathfrak{a}}=j^+_a$ and the weight is $L_0^{2(j^+_a-m)}$, so
$\tfrac12\log_{L_0}\omega=j^+_a-m$ and $\mathfrak{c}=j^+_a-(j^+_a-m)=m$. Where $\omega=1$ we
have $\log_{L_0}\omega=0$ and $\mathfrak{c}=\ell_{\mathfrak{a}}$. In particular
$g(y)=\tfrac12\log_{L_0}\omega(y)$, so the lifted points are those at which
$\omega(y)\ge L_0^2$.
\end{proof}

\begin{lemma}[The domain statements at points of older live pieces]\label{4.26:lem-older-pieces-domain}
Let the level $k$, the stages $l=1,\dots,N$, the region $Z$, the set $Z^{\mathrm{on}}=Z\cap\Omega_k^c$ and the
region $\mathfrak{T}=\Omega_k\cup Z^c$ be as in Definition~\ref{4.26:def-stages-regions}, and let
$\beta=1/5$ and $\mathsf{w}_0=1/24$. Let $\mathfrak{a}=(W_{\mathfrak{a}},k_a,n^+_a,\text{branch})$ be an arrest
that is live at level $k$, with $k_a<k$, and whose component is not in the class of $Z$. Let
$y\in W_{\mathfrak{a}}\cap\mathfrak{T}$. Every object at $y$ is read on the frozen structure of the earliest
live arrest holding $y$, as in Definition~\ref{4.26:def-stages-regions}. This fixes the shell index $m(y)$, the
weight $\omega(y)=w_{\mathfrak{a}}(y)\ge1$, the radii $r_a(y)$ and $r_e^{(p)}(y)$, the source threshold
$\theta_{\mathrm{src},a}(y)$, and the numbers $n_W(y)$, $\mathfrak{c}(y)$ and $g(y)$. In particular
$m(y)=\ell_{\mathfrak{a}}(y)$, and we put $\ell:=\ell_{\mathfrak{a}}(y)$. The point $y$ is called \emph{lifted} if
$g(y)\ge1$. We write $\alpha_{\cdot,j}:=\min\{\alpha_{0,j},\alpha_{1,j}\}$.

Let $l$ be a stage of the site and put $n=l-1$; $\mathfrak{N}^{(n)}$ denotes the regularity space
of~\eqref{4.26:eq-regularity-space-a} at that stage.
In the display below, $w$ is the width constant of the second separation property of an admissible
determining set. The letter $P$ denotes the constant of the tube proposition, in which the numerical factor
$C_{1/2}=\sqrt2$, defined in Step~(6) of the proof below, appears; it is the same constant as in the
condition $16\max\{P,P',P''\}\le\beta\theta_S$ of~\eqref{4.26:eq-source-image-a}. Finally,
$\Delta_{\mathrm{tube}}(y)$ denotes the increment of the tube around the source at $y$. Rooms and
$\Delta_{\mathrm{tube}}(y)$ are measured in units of $r_a(y)$. Let $\delta_1$ be the decay rate entering the
tube proposition. Assume the conditions~\eqref{4.26:eq-inherited-conditions-a}, the condition
$16\max\{P,P',P''\}\le\beta\theta_S$ of~\eqref{4.26:eq-source-image-a}, and the condition
$\delta_1M/M_1\ge2$ on the auxiliary parameters for arrested pieces. Assume further that part~(a) of the
inherited-domain theorem and the corollary on the positive room of the inherited domain over the source,
\eqref{4.26:eq-inherited-room-a}, hold at every site of level smaller than $k$. Then the following hold.
\begin{equation}\label{4.26:eq-older-pieces-domain-a}
\begin{gathered}
F^{(k)}(y)\le 1<1+2\beta-2\mathsf{w}_0\beta,\qquad
(1+2\beta)\,\omega(y)\,\alpha_{\cdot,\ell}<\gamma_0^{(k)},\\
d^{\wedge}\bigl(y,\mathcal{C}^{(l)}\bigr)\ge w\,n_W(y),\qquad
r_{\min}(y)\ge\beta\theta_S\,2^{-(k-\ell)},\\
\Delta_{\mathrm{tube}}(y)\le P\,2^{-n_W(y)}\le\tfrac14\,\mathrm{room}^{\flat}(y).
\end{gathered}
\end{equation}
More precisely:
\begin{enumerate}
\item[(1)] The arrested source $\tilde{U}^{\mathrm{arrest}}_k$ lies in $\mathfrak{N}^{(n)}$ at $y$. Moreover
$\mathfrak{N}^{(n)}\subset\mathfrak{K}_{109}(\mathfrak{B};\gamma_0^{(k)})$ at $y$, where
$\mathfrak{K}_{109}(\mathfrak{B};\gamma_0^{(k)})$ denotes the kernel-regularity space of \cite{B-CMP109}
(to which the index $109$ refers), taken over the determining set $\mathfrak{B}$ with size $\gamma_0^{(k)}$;
the bounds in the first line of~\eqref{4.26:eq-older-pieces-domain-a} hold.
\item[(2)] At $y$ the stage increments of direct entries are at most $\iota_l\,r_a(y)$. Those of
composite entries are at most $\iota_l\,r_e^{(p)}(y)$.
\item[(3)] The distance bound in the second line of~\eqref{4.26:eq-older-pieces-domain-a} holds, with $w$
as above. Hence the bounds on
$\sum\iota$ and $\sum\iota^1$ of~\eqref{4.26:eq-inherited-expansion-a} hold at $y$ as stated there.
\item[(4)] $r_{\min}(y)\ge\beta\theta_S2^{-(k-\ell)}$. This is the quantity $\beta e_\ast(\ell)$ of the room
table~\eqref{4.26:eq-room-table-a}. Thus the bound (k4) of~\eqref{4.26:eq-inherited-room-a} holds at
$n=n_W(y)$, as stated there.
\item[(5)] Consider the terms of kind $\mathbb{E}$ and $\mathbb{R}$, and every term whose clauses were fixed
before the arrest of $\mathfrak{a}$. For such a term $T$ of level $\ell_T$: if $g(y)=0$ and $\ell_T=\ell$, the
bound (k5) of~\eqref{4.26:eq-inherited-room-a} holds at $y$ as stated; if $g(y)=0$ and $\ell_T<\ell$,
entry~(b) of the room table~\eqref{4.26:eq-room-table-a} holds at $y$ as stated; if $g(y)\ge1$, entry~(b)
holds at $y$ with (k6) replaced by (k6)$''$.
\item[(6)] The tube increment $\Delta_{\mathrm{tube}}(y)$ satisfies the third line
of~\eqref{4.26:eq-older-pieces-domain-a}. The exponent on both sides is $n_W(y)$; the bound involves neither
$k-k_a$, nor the number $N$ of stages of the site, nor the width of the shells of $W_{\mathfrak{a}}$.
\item[(7)] Consequently, the following statements hold at $y$ with the arrested source in the role of the
source:
\begin{itemize}
\item the three lines of the corollary on the positive room of the inherited domain over the source,
\eqref{4.26:eq-inherited-room-a}, together with its concluding inequality;
\item part~$(\beta)$ of the claim that the image of the source lies in the inherited domain,
\eqref{4.26:eq-source-image-a};
\item parts (1)--(4)$(\gamma)$ of the inherited-domain theorem, whose conditions
are~\eqref{4.26:eq-inherited-conditions-a}.
\end{itemize}
None of these statements uses any further property of $W_{\mathfrak{a}}$.
\end{enumerate}
\end{lemma}

\begin{proof}
Throughout, $y\in W_{\mathfrak{a}}\cap\mathfrak{T}$ and $\ell=\ell_{\mathfrak{a}}(y)$ are fixed. All objects at
$y$ are read on the frozen structure described in the statement.

\emph{Step (1): the regularity space and its class.} We first check membership in the regularity space
$\mathfrak{N}^{(n)}$, one clause at a time. The source factor satisfies $F^{(k)}(y)\le1$. With $\beta=1/5$
and $\mathsf{w}_0=1/24$ we have
\[
1+2\beta-2\mathsf{w}_0\beta=1.4-\tfrac{1}{60}=1.383\ldots .
\]
The factor that $\mathfrak{N}^{(n)}$ allows in each clause, as defined
in~\eqref{4.26:eq-regularity-space-a}, is $1+2\beta-\sum\iota^1$. By the consumption of room by the proof
family, \eqref{4.26:eq-family-consumption-a}, it is at least $1.383$. The direct entries of both the source
and $\mathfrak{N}^{(n)}$ are measured in the unit $r_a(y)$.

The composite entries are compared at every $(\Sigma^{\ast})^{\mathfrak{a}}$-triple $(p,X',\mathfrak{s})$. There
both the source and $\mathfrak{N}^{(n)}$ are measured in the same unit $r_e^{(p)}(y)$ of
Definition~\ref{4.26:def-stages-regions}, whose block side is $\min\{p,\lambda_{\mathfrak{s}}(y)\}$: for
$\mathfrak{N}^{(n)}$ by the block-side clause of the construction of the regularity space, and for the arrested
source by the corresponding block-side clause of the comparison of the arrested source with the attempt spaces.
Since $F^{(k)}(y)\le1<1.383$ in every clause, the arrested source lies in $\mathfrak{N}^{(n)}$ at $y$.

Next we show $\mathfrak{N}^{(n)}\subset\mathfrak{K}_{109}(\mathfrak{B};\gamma_0^{(k)})$ at $y$; this space is
formed with the density-indexed size
$\gamma_0^{(k)}$ of~\eqref{4.26:eq-regularity-radius-a}. The weight at $y$ satisfies
\[
\omega(y)\le L_0^{2N_0^a}=w_\ast(k_a),
\]
since $w_\ast(k_a)$ is the highest power of the weight that occurs at density $k_a$ in the construction of
\cite{B-CMP122a}. The nesting lemma for arrests gives $\ell\le k_a\le k$. Since $\bar{\alpha}_{k_a}$ is the
maximum of the radii $\alpha_{i,m}$ over the levels $m\le k_a$, it follows from its definition that
$\alpha_{\cdot,\ell}\le\bar{\alpha}_{k_a}$. The definition of $\gamma_0^{(k)}$
in~\eqref{4.26:eq-regularity-radius-a} bounds $5\,w_\ast(k_a)\bar{\alpha}_{k_a}$ by $\gamma_0^{(k)}$. Together
with $1+2\beta=1.4$ this gives
\[
(1+2\beta)\,\omega(y)\,\alpha_{\cdot,\ell}
\le 1.4\,w_\ast(k_a)\,\bar{\alpha}_{k_a}
\le \tfrac{1.4}{5}\,\gamma_0^{(k)}
<\gamma_0^{(k)} .
\]
This holds at every age $k-k_a$ of the piece. The remaining condition in the definition of
$\mathfrak{K}_{109}(\mathfrak{B};\gamma_0^{(k)})$, the
one carrying the
constant $600$, holds because $120\,\gamma_0^{(k)}\le1$. This is the condition on $\gamma_0^{(k)}$ among the
conditions~\eqref{4.26:eq-inherited-conditions-a}, assumed in the statement. This proves the first line
of~\eqref{4.26:eq-older-pieces-domain-a}, and hence~(1).

\emph{Step (2): increments.} Consider part~(i) of the increment lemma for the stage layers. It holds for every
determining set $\mathfrak{B}$, because it uses $\mathfrak{B}$ only through its admissibility properties,
including its two separation properties. Its increments are stated in unweighted shell units.

For a direct entry the increment is therefore at most $\iota_l\mathbf{r}_{a,\ell}$. Since $\omega(y)\ge1$ and
$r_a(y)=\omega(y)\mathbf{r}_{a,\ell}$, this is at most $\iota_l\,r_a(y)$.

For a composite entry at a $(\Sigma^{\ast})^{\mathfrak{a}}$-triple $(p,X',\mathfrak{s})$ write
$r_e:=r_e\bigl(\ell,\min\{p,\lambda_{\mathfrak{s}}(y)\}\bigr)$, so that $\omega(y)\,r_e=r_e^{(p)}(y)$. Let $Q_e$
denote the quantity through which part~(i) of the increment lemma bounds the increment of a composite entry.
It satisfies $Q_e\le1.4\,\omega r_e$. The increment of the composite entry is then at most the first member of
the chain
\[
(1.4\,\pi_\times\omega+\pi)\,r_e\le 2.4\,\pi\,\omega\,r_e\le\iota_l\,\omega\,r_e=\iota_l\,r_e^{(p)}(y) .
\]
Here $\pi$ and $\pi_\times$ are the increment constants entering~\eqref{4.26:eq-family-consumption-a}. The
middle inequality uses $\omega(y)\ge1$ together with $\pi_\times\le\pi$, and the last uses $2.4\,\pi\le\iota_l$,
relations between the constants of~\eqref{4.26:eq-family-consumption-a}. This proves~(2).

\emph{Step (3): the expansion on the inherited domain.} The sequence $\{\Omega''_j\}$ is stated in
\cite{B-CMP122a} to be an admissible sequence. Hence the strata $\ell+1,\dots,k-1$ of $\mathfrak{B}$ separate
$y$ from $\Lambda_k$. Each of these strata is at least one of its own $M$-cubes thick, whichever step created
it. By the second separation property of the admissible determining set $\mathfrak{B}$,
\[
d^{\wedge}\bigl(y,\mathcal{C}^{(l)}\bigr)\ge w\,n_W(y).
\]
This is the distance that the bounds on $\sum\iota$ and $\sum\iota^1$
in~\eqref{4.26:eq-inherited-expansion-a} require. These bounds therefore hold at $y$ as stated there, and (3)
is proved.

\emph{Step (4): rooms over the source threshold.} For every term $T$ whose clauses are read at $y$, we bound
the room of $T$ over the source threshold $\theta_{\mathrm{src},a}(y)$, measured in units of $r_a(y)$. The
cases are the following.
\begin{itemize}
\item[(a)] \emph{Clause sets of class $k_a$.} These are the frozen terms $\mathfrak{F}(\mathfrak{a})$ of
$\mathfrak{a}$, the frozen terms $\mathfrak{F}(\mathfrak{b})$ of its siblings $\mathfrak{b}$, and the clause
sets of the stage structure $\mathbb{B}''_{k_a}$ on an untouched region $W'\subset Z_k\setminus Z$. The
extended room lemma for frozen clauses gives a room of at least
\[
\beta2^{-(k_a+1-\ell)}\ge\beta2^{-(k-\ell)} .
\]
The inequality holds because $k_a+1\le k$.
\item[(b)] \emph{The terms $\mathbf{B}^{(k')}$ on the clauses of the arrested attempt spaces.} These are
the terms born by $\mathbf{T}$ at a level $k'$ with $k_a<k'<k$, or born by $\mathbf{R}$ at a level $k'$ with
$k_a\le k'<k$. Their room is
\[
F^{(k')}-F^{(k)}\ge s^{(k')}_\ell-s^{(k)}_\ell
=\beta\bigl(2^{-(k'-\ell)}-2^{-(k-\ell)}\bigr)\ge\beta2^{-(k-\ell)} .
\]
The last inequality holds because $k'\le k-1$ gives $2^{-(k'-\ell)}\ge2\cdot2^{-(k-\ell)}$. Arrests made live
at levels between $k'$ and $k$ can only lower $F^{(k)}$, so they do not decrease this room.
\item[(c)] \emph{Terms born by $\mathbf{T}$ at level $k$.} Their room is $\theta_S\beta2^{-(k-\ell)}$, by the
schedule for newborn terms in the room table~\eqref{4.26:eq-room-table-a}.
\item[(d)] \emph{Frozen terms of a later foreign arrest.} Let $\mathfrak{b}$ be an arrest with
$k_a<k_b<k$ whose frozen terms reach $y$, with $y\in\mathfrak{T}$ at the site $k_b$. Here the site $k_b$ is the
application of the large-field operation $\mathbf{R}$ at level $k_b$. The domain of these terms is the
inherited domain of the site $k_b$. The sites are treated in increasing order of level. Since $k_b<k$,
part~(a) of the inherited-domain theorem and the corollary on the positive room of the inherited domain over
the source, \eqref{4.26:eq-inherited-room-a}, hold at the site $k_b$ by hypothesis. They give a room of at
least
\[
\tfrac12\beta\theta_S2^{-(k_b-\ell)}\ge\beta\theta_S2^{-(k-\ell)} .
\]
The inequality holds because $k_b\le k-1$.
\item[(d$'$)] \emph{Frozen terms of a nested later arrest.} Let $\mathfrak{b}\supset\mathfrak{a}$ be a later
arrest, with $k_a<k_b<k$. By the nesting lemma for arrests, $W_{\mathfrak{a}}\subset W_{\mathfrak{b}}$ and
$h_b\ge k_a$. Then $y$
lies on the piece of $\mathfrak{b}$, which is a $Z^{\mathrm{on}}$-location of the site $k_b$. There the frozen
terms of $\mathfrak{b}$ are clause sets of class $k_b$. The extended room lemma for frozen clauses, applied at
class $k_b$, gives a room of at least
\[
\beta2^{-(k_b+1-\ell)}\ge\beta2^{-(k-\ell)} .
\]
The inequality holds because $k_b+1\le k$. This is the bound of~(a), one class up.
\item[(d$''$)] \emph{Frozen terms of a sibling of a nested later arrest.} Let $T\in\mathfrak{F}(\mathfrak{b}'')$,
where $\mathfrak{b}''$ is a sibling of a nested arrest $\mathfrak{b}'\supset\mathfrak{a}$ with $k_a<k_{b'}<k$.
Then $y$ is a
$Z^{\mathrm{on}}$-location at the level $k_{b'}$. The terms covered by the extended room lemma for frozen
clauses include the frozen terms of siblings. That lemma therefore gives the same bound,
\[
\beta2^{-(k_{b'}+1-\ell)}\ge\beta2^{-(k-\ell)} .
\]
Case~(d) concerns only foreign arrests $\mathfrak{b}$, that is, those with $y\in\mathfrak{T}$ at the site $k_b$.
\item[(e)] \emph{Terms whose clauses were fixed before the arrest of $\mathfrak{a}$.} These are the terms of
\cite{B-CMP119} and the frozen terms $\mathfrak{F}(\mathfrak{b})$ with $k_b<k_a$. Suppose first that $y$ is
not lifted. Then there is the chain of inclusions
\[
\mathfrak{U}_T\supseteq\mathfrak{S}^{(k_a)}_0\supseteq\mathfrak{S}^{(k_a)}_{n^+_a}
\supseteq\tilde{U}^{\mathrm{arrest}}_k .
\]
The chain rests on the descent of the attempt spaces in \cite{B-CMP122a}, on the extended room lemma for
frozen clauses at attempt number $0$ for the terms of \cite{B-CMP119}, and on the arrested-source theorem for
the terms $\mathfrak{F}(\mathfrak{b})$. This chain yields at
least the bound of~(a). If $y$ is lifted, the room follows from the ratio bound obtained in Step~(5) below.
\end{itemize}
Taking the minimum over these cases gives
\[
r_{\min}(y)\ge\beta\theta_S2^{-(k-\ell)} .
\]
The right-hand side is the quantity $\beta e_\ast(\ell)$ of the room table~\eqref{4.26:eq-room-table-a}.
Hence the bound (k4) of~\eqref{4.26:eq-inherited-room-a} holds at $n=n_W(y)$, in the form stated there. This
proves~(4).

\emph{Step (5): terms of kind $\mathbb{E}$ and $\mathbb{R}$, and every term of case~(e), at $y$.} We
distinguish three cases. For such a term $T$ we write $\ell_T$ for its level and $r_T$ for its room at $y$ over
the source threshold, in units of $r_a(y)$.

\emph{Case $g(y)=0$ and $\ell_T=\ell$.} Here the room satisfies
\[
r_T\ge1-F^{(k)}(y)\ge1-s^{(k)}_\ell\ge\beta/2 .
\]
This is the bound (k5) of~\eqref{4.26:eq-inherited-room-a}, as stated there.

\emph{Case $g(y)=0$ and $\ell_T<\ell$.} Since $\omega(y)=L_0^{2g(y)}$, here $\omega(y)=1$, and $y$ lies among
the bands $\Gamma''_m$ with $m\le k_0^a$. On these bands the weight is $1$, and the frozen structure changes the
level of each of these bands by one. Entry~(b) of the room table~\eqref{4.26:eq-room-table-a} therefore applies
as stated.

\emph{Case $g:=g(y)\ge1$.} The level of the clause of $T$ at $y$ is at most $\mathfrak{c}(y)$. For a term of
kind $\mathbb{E}$ or $\mathbb{R}$ this holds because $X_T\subset\Lambda_{\ell_T}$ and
$y\notin\Omega_{\mathfrak{c}(y)+1}$. For a boundary term of case~(e), the level of its clause at $y$ is the
level that $y$ had at the birth of $T$.

Among the radii, the least favourable entry is $a=6$. For it,
\[
\frac{\mathbf{r}^T_6}{r_6(y)}\ge\Bigl(\frac{L}{(1+\beta_0)L_0^2}\Bigr)^{g}\ge 2.625^{\,g}.
\]
The second inequality uses $L_0^2\le L/3$ and $\beta_0=1/7$, as in \cite{B-CMP122a}. With these values,
\[
\frac{L}{(1+\beta_0)L_0^2}\ge\frac{3}{8/7}=\frac{21}{8}=2.625 .
\]
Hence, writing $\theta_T$ for the threshold of the clause of $T$ at $y$, and since $g\ge1$ and
$0.8\cdot2.625=2.1$,
\[
\theta_T\ge0.8\cdot2.625^{\,g}\,r_a(y)\ge2.1\,r_a(y) .
\]
This is entry~(b) of the room table~\eqref{4.26:eq-room-table-a}, with (k6) replaced by (k6)$''$
of~\eqref{4.26:eq-inherited-room-a}. This proves~(5), and it also supplies the bound used in case~(e) of
Step~(4) at lifted points.

\emph{Step (6): the tube.} We apply the concluding estimate of the tube proposition, which bounds the increment
of the tube around the source. The far-separation hypothesis of that proposition is supplied by Step~(3).

By the scale-free separation property of admissible sequences, applied to the admissible sequences of $Z$ and
of $W_{\mathfrak{a}}$, the scaled distance $d_{\mathrm{sc}}(y,\mathcal{Z})$ is at least $5R_k$ plus, for each
stratum crossed, the side of one of that stratum's own $M$-cubes.
Here $\mathcal{Z}$ is the set from which the far-separation hypothesis of the tube proposition measures
distances, and $R_k$ is the distance parameter of \cite{B-CMP122a} at level $k$. The number of strata crossed
is $n_W(y)$.

By hypothesis, $\delta_1M/M_1\ge2$. Each stratum crossed is at least one own $M$-cube
thick, and hence contributes a factor at most
\[
e^{-\delta_1M/M_1}\le e^{-2}\le\tfrac12\,e^{-0.6}.
\]
Together the $n_W(y)$ strata contribute a factor at most $2^{-n_W(y)}e^{-0.6\,n_W(y)}$.

The factor $e^{-0.6\,n_W(y)}$ absorbs the ratio of couplings. The increment of the tube proposition carries the
ratio $\alpha_{0,k}/\alpha_{\cdot,\ell}$. This is a coarse radius over a fine one, since $\ell<k$. For $i=0,1$ the
radii have the form
\[
\alpha_{i,j}=C_i\,g_j(\log g_j^{-2})^q ,
\]
and $\alpha_{\cdot,j}=\min\{\alpha_{0,j},\alpha_{1,j}\}$. Therefore
\[
\frac{\alpha_{0,k}}{\alpha_{\cdot,\ell}}
=\frac{C_0}{\min\{C_0,C_1\}}\cdot\frac{g_k}{g_\ell}
\cdot\Bigl(\frac{\log g_k^{-2}}{\log g_\ell^{-2}}\Bigr)^{q}.
\]
To bound the two factors on the right we use the first chain of inequalities of
\cite[(2.6), p.~255]{B-CMP119}. It states, for $n>m$,
\[
g_n\le\bigl(1+g_n^2\beta'(n-m)\bigr)^{1/2}g_m
\le(1+g_n^2\beta')^{1/2}(n-m)^{1/2}g_m
\le(1+\beta_0)(n-m)^{1/2}g_m .
\]
Here $\beta'$ is the constant of that display. The first inequality of \cite[(2.7), p.~255]{B-CMP119} states,
for $n>m$,
\[
(\log g_n^{-2})^p\le(1+\beta_0)(\log g_m^{-2})^p .
\]
Apply the first with $(n,m)=(k,\ell)$, and the second with $(n,m)=(k,\ell)$ and $p=q$. This gives
\begin{equation}\label{4.26:eq-older-pieces-domain-1}
\alpha_{\cdot,k}\le(1+\beta_0)^2(k-\ell)^{1/2}\alpha_{\cdot,\ell},
\end{equation}
and, since $1+\beta_0>1$,
\[
\frac{\alpha_{0,k}}{\alpha_{\cdot,\ell}}
\le\frac{C_0}{\min\{C_0,C_1\}}\,(1+\beta_0)^{\max\{2,1+q\}}\,(k-\ell)^{1/2} ;
\]
the constant in front of $(k-\ell)^{1/2}$ is absorbed in $P$.

The radii $\alpha_{\cdot,j}$ are of the same form as the quantities
\[
\varepsilon_j=g_jA_0(\log g_j^{-2})^{p_0}
\]
of \cite[(2.4)]{B-CMP119}. For these, \cite[(2.8), p.~256]{B-CMP119} states as its first inequality
\[
\varepsilon_n\le(1+\beta_0)(n-m)^{1/2}\varepsilon_m .
\]
The same use of the square-root bound against the decay of the shells is made in \cite{B-CMP122a}. The second
inequality of \cite[(2.8)]{B-CMP119} bounds the finer quantity by the coarser one. It is not the bound used
here.

Since $n_W(y)=k-1-\ell$, we have $k-\ell=1+n_W(y)$. Put $f(x)=(2+x)^{1/2}e^{-0.6x}$. On $x\ge0$,
\[
\frac{f'(x)}{f(x)}=\frac{1}{2(2+x)}-0.6\le\frac14-0.6<0,
\]
so $f$ is decreasing on $[0,\infty)$ and attains its supremum at $x=0$. Define
\[
C_{1/2}:=\sup_{x\ge0}(2+x)^{1/2}e^{-0.6x}=\sqrt2 .
\]
Then
\begin{align*}
(k-\ell)^{1/2}e^{-0.6\,n_W(y)}
&=\bigl(1+n_W(y)\bigr)^{1/2}e^{-0.6\,n_W(y)}\\
&\le\bigl(2+n_W(y)\bigr)^{1/2}e^{-0.6\,n_W(y)}\le C_{1/2}.
\end{align*}

The constants enter $P$ in the form in which the tube proposition carries them. These are
$(1+\varrho_k)(1+\beta_0)^2$, with $\varrho_k=\alpha_{1,k}/\alpha_{0,k}$, the constant
$C_0/\min\{C_0,C_1\}$ of the ratio, and the factor $C_{1/2}$.

There is a second route to the same bound, which uses only the conditions of the inherited-domain theorem.
Apply the adjacent-ratio condition of~\eqref{4.26:eq-inherited-conditions-a} at each of the $k-\ell$ adjacent
pairs of levels between $\ell$ and $k$. This gives
\[
\frac{\alpha_{\cdot,k}}{\alpha_{\cdot,\ell}}\le(1+\beta_0)^{k-\ell}=(1+\beta_0)(1+\beta_0)^{n_W(y)} .
\]
Since $e^{-0.6}=0.5488\ldots$ and $\beta_0=1/7$,
\[
(1+\beta_0)e^{-0.6}\le\tfrac87\cdot0.55<0.63 .
\]
Therefore
\begin{align*}
(1+\beta_0)^{k-\ell}e^{-0.6\,n_W(y)}
&=(1+\beta_0)\bigl((1+\beta_0)e^{-0.6}\bigr)^{n_W(y)}\\
&\le1+\beta_0<(1+\beta_0)^2C_{1/2},
\end{align*}
and the same constant $P$ serves.

By either route the estimate involves $n_W(y)$ on both sides, and it contains none of $k-k_a$, $N$, or a width
of the shells of $W_{\mathfrak{a}}$. The increment of the tube at $y$ is therefore at most $P\,2^{-n_W(y)}$ in
the unweighted units $\mathbf{r}_{a,\ell}$. Since $\omega(y)\ge1$, we have $\mathbf{r}_{a,\ell}\le r_a(y)$, so
that in units of $r_a(y)$
\[
\Delta_{\mathrm{tube}}(y)\le P\,2^{-n_W(y)} .
\]
It is to be compared with a quarter of the room, also measured in units of $r_a(y)$. By Step~(4) and the
definition of $\mathrm{room}^{\flat}$ in~\eqref{4.26:eq-inherited-room-a},
\[
\tfrac14\,\mathrm{room}^{\flat}(y)\ge\tfrac1{16}\,\beta\theta_S\,2^{-n_W(y)} .
\]
Hence the third line of~\eqref{4.26:eq-older-pieces-domain-a} follows from $16P\le\beta\theta_S$. This is part of
the condition $16\max\{P,P',P''\}\le\beta\theta_S$ in~\eqref{4.26:eq-source-image-a}, which is assumed in the
statement. This proves~(6).

\emph{Step (7): conclusion.} Steps~(1)--(6) supply at $y$ every property on which the following statements
depend, with the arrested source in the role of the source:
\begin{itemize}
\item the three lines of the corollary on the positive room of the inherited domain over the source,
\eqref{4.26:eq-inherited-room-a}, and its concluding inequality;
\item part~$(\beta)$ of the claim that the image of the source lies in the inherited domain,
\eqref{4.26:eq-source-image-a};
\item parts (1)--(4)$(\gamma)$ of the inherited-domain theorem.
\end{itemize}
These statements use no further property of $W_{\mathfrak{a}}$. Hence they hold at $y$, which proves~(7).

Part (4)$(\delta)$ of the inherited-domain theorem, the near-pair case at boxes that meet lifted points, is not
among the conclusions of this lemma. It rests on the weighted scale comparison for boxes meeting the lifted
set.
\end{proof}

\begin{definition}[The cover, the regularity radius and its run-free majorant]
\label{4.26:def-regularity-radius}
Throughout, the current density $k$ is fixed. The stages, the regions, and the notation $\eta=L^{-k}$,
$h=k-N$, $k_0=k-N_0$ are those of Definition~\ref{4.26:def-stages-regions}. We write $\mathfrak{B}$ for
the current determining-set structure. For every level $m$ we write $x_m=g_m^{-2}$ and
$\mathsf{t}_m:=\log g_m^{-2}=\log x_m$. We write $\mathsf{t}_\gamma:=\log\gamma^{-2}$ for the threshold on
$\mathsf{t}$; hypothesis (g) of the stage construction gives $\mathsf{t}_j\ge\mathsf{t}_\gamma$ at every
level $j\le k$.

The rate of the scaled distance is $\delta:=\tfrac18\delta_0$. It is distinct from two other quantities
written $\delta$: the parameter written $\delta$ in \cite{B-CMP116}, which enters here at one place only,
in the order in which the constants are chosen; and the window parameter $\delta\in(0,\delta_J]$ of the
correlation theorem. We put $w:=M/M_1$.

Let $d^{\wedge}:=d^{\wedge}_{\mathfrak{B}}$ be the scaled distance attached to $\mathfrak{B}$. It
satisfies $d^{\wedge}\le d$, and two further properties fixed in its construction; by these
properties, a path that crosses a stratum of $\mathfrak{B}$ is charged at least $w$. The notation
$d_{\mathrm{sc}}$ means $d^{\wedge}$. The rate and $w$ are subject to the level condition
\begin{equation}\label{4.26:eq-regularity-radius-1}
\delta w\ge 16\ln L+4 .
\end{equation}

The cover of the stage is the family $\mathcal{C}=\{\square\}$ of cubes built on $\mathfrak{B}$. It comes
with the partition of unity $\{\zeta_\square\}$, the enlargements $\tilde{\square}$, and the sets
$\square^{\zeta}\subset\square$ of its construction. The estimates on the terms of the class
$\mathcal{T}_1$ use instead the graded cover $\mathcal{C}^1$, built on the graded refinement
$\mathfrak{B}^{\mathfrak{c}}$ of $\mathfrak{B}$ as in \eqref{4.26:eq-graded-cover-a}. We write $N_y$ and
$N^{\Delta}_y$ for the local norms of the stage construction.

Let $\sigma$ be the complex parameters of the stage, with $|\sigma(\Delta)|\le e^{\kappa_1}$ for every
cube $\Delta$, and let $Y(\sigma)$ be the region attached to them. Let $b$ be a complex function on the
level-$j$ collar bonds $\mathcal{C}_j$ (not to be confused with the cover $\mathcal{C}$), where
$j=\ell_l$ is the level integrated at the stage (Definition~\ref{4.26:def-stages-regions}), with
\begin{equation*}
\|b\|_\infty\le 1.1\,C_1\delta'_j .
\end{equation*}
We write $(A',X,\mathbb{H})$ for the three objects furnished by the theorem that constructs the stage
layer $\mathbb{H}(\sigma;b)$. That theorem is applied at the data
$(\mathfrak{B};\mathbb{U},\mathbb{J};\sigma;(0,b))$, where $(\mathbb{U},\mathbb{J})$ is the background
pair. The space on which it is applied is the kernel-regularity space
$\mathfrak{K}_P:=\mathfrak{K}(\mathfrak{B};\gamma_0)$ of \cite{B-CMP109}, built on $\mathfrak{B}$, of size
$\gamma_0:=\gamma_0^{(k)}$. Here $\gamma_0^{(k)}$ is the \emph{regularity radius} at density $k$:
\begin{equation}\label{4.26:eq-regularity-radius-a}
\gamma_0^{(k)}:=5\max_{j\le k}\,w_\ast(j)\,\bar{\alpha}_j .
\end{equation}
In \eqref{4.26:eq-regularity-radius-a} and in item (b) below the letter $j$ runs over densities
$j\le k$; it is not the level $\ell_l$ of the stage.
The ingredients of \eqref{4.26:eq-regularity-radius-a} are as follows.
\begin{itemize}
\item The maximum runs over those densities $j\le k$ at which Ba{\l}aban's integer $N_0(j)$ of
\cite{B-CMP122a} is defined. At every site of the operation $\mathbf{R}$ these include the current
density $k$ and the density $k_a$ of every live piece.
\item $w_\ast(j):=L_0^{2N_0(j)}$, where $N_0(j)$ is Ba{\l}aban's integer at density $j$ and $L_0$ is the
base of Ba{\l}aban's weight $\omega$ in \cite{B-CMP122a}, subject to $L_0^2\le L/3$; the integer $N_0(j)$
satisfies $L^{N_0(j)-1}=R_{j-N_0(j)+1}$ \cite{B-CMP122a}.
\item $\bar{\alpha}_j$ is the maximum over the levels $m\le j$:
\begin{equation*}
\bar{\alpha}_j:=\max_i\max_{1\le m\le j}\alpha_{i,m},\qquad
\alpha_{i,m}=C_i\,x_m^{-1/2}(\log x_m)^q .
\end{equation*}
\item $\alpha_{\cdot,m}$ stands for any of the $\alpha_{i,m}$, $i\in\{0,1\}$; these are the two radii of
Lemma~\ref{4.4:lem-radii-ratio}.
\item We write $w_\ast:=w_\ast(k)$ when the weight at the current density is meant.
\end{itemize}
Wherever $\mathfrak{K}_P$ is formed, $\gamma_0$ means $\gamma_0^{(k)}$.

\smallskip
(a) Let $\mathfrak{a}$ be an arrest live at the current density whose piece $W_{\mathfrak{a}}$ is older,
$k_a\le k$, and put $N_0^a:=N_0(k_a)$. Then for every $y\in W_{\mathfrak{a}}$,
\begin{equation}\label{4.26:eq-regularity-radius-2}
\omega(y)\le L_0^{2N_0^a}=w_\ast(k_a),\qquad \ell_{\mathfrak{a}}(y)\le k_a,
\end{equation}
and hence
\begin{equation}\label{4.26:eq-regularity-radius-3}
(1+2\beta)\,\omega(y)\,\alpha_{\cdot,\ell_{\mathfrak{a}}(y)}
\le 1.4\,w_\ast(k_a)\,\bar{\alpha}_{k_a}
\le \tfrac{1.4}{5}\,\gamma_0^{(k)} .
\end{equation}
At points $x\in Z^{\mathrm{on}}$ we have $5\,\mathsf{w}^{\mathrm{on}}(x)\,\alpha_{\cdot,m(x)}\le
\gamma_0^{(k)}$, for either branch of the weight $\mathsf{w}^{\mathrm{on}}(x)$ of the fat factor. No
comparison of couplings at two different densities enters these bounds.

\smallskip
(b) Put $C:=\max_iC_i$, $c_\beta:=2\ln(1+\beta_0)$, where $\beta_0$ is the constant entering the
inequalities \cite[(2.6)]{B-CMP119} used in the proof below, and
\begin{equation}\label{4.26:eq-regularity-radius-4}
\Psi(\mathsf{t}):=L^2(L+1)\,C\,\mathsf{t}^{r+q}e^{-\mathsf{t}/2},\qquad
\bar{\gamma}_0(\mathsf{t}):=5(1+\beta_0)^2\,\Psi(\mathsf{t}).
\end{equation}
Assume hypothesis (g) and the range clause
\begin{equation}\label{4.26:eq-regularity-radius-5}
\mathsf{t}_\gamma\ge 2(r+q)+c_\beta .
\end{equation}
Then the regularity radius has the \emph{run-free majorant}
\begin{equation}\label{4.26:eq-regularity-radius-6}
\gamma_0^{(k)}\le\bar{\gamma}_0(\mathsf{t}_k).
\end{equation}
The function $\bar{\gamma}_0$ is decreasing in the single variable $\mathsf{t}$ on
$[2(r+q),\infty)$. It is of degree $-\infty$: $\mathsf{t}^n\bar{\gamma}_0(\mathsf{t})\to0$ as
$\mathsf{t}\to\infty$ for every $n$. The functions $\Psi$ and $\alpha_i$ may be read as their
non-increasing envelopes $\Psi^{\downarrow}(\mathsf{t}):=\sup_{s\ge\mathsf{t}}\Psi(s)$ and
$\alpha_i^{\downarrow}(\mathsf{t}):=\sup_{s\ge\mathsf{t}}\alpha_i(s)$, which coincide with them on the
range used. With this reading, the range clause \eqref{4.26:eq-regularity-radius-5} is weakened to
$\mathsf{t}_\gamma\ge\max\{2r,c_\beta\}$. The proof below uses the range clause in the form
\eqref{4.26:eq-regularity-radius-5}.

\smallskip
(c) The theorem constructing the stage layer has the following hypotheses on $\gamma_0$:
\begin{itemize}
\item the condition $K_\Delta\gamma_0\le\tfrac12$, with the constant $K_\Delta$ of that theorem;
\item a second smallness condition on $\gamma_0$ of that theorem;
\item the smallness hypothesis $(S^{\mathrm{b}})$ of the bump lemma below;
\item the member of its list of conditions that carries the numerical constant $600$, which in the
present notation reads $120\,\gamma_0^{(k)}\le1$.
\end{itemize}
Each of these hypotheses is imposed with $\gamma_0$ replaced by $\bar{\gamma}_0(\mathsf{t}_\gamma)$. Each
is therefore a condition on the supremum, over the single variable $\mathsf{t}\ge\mathsf{t}_\gamma$, of a
function of degree $-\infty$.

\smallskip
The decaying majorants of the stage are
\begin{equation}\label{4.26:eq-regularity-radius-7}
p(y):=\Bigl(\sup_{\mathcal{C}_j}|b|\Bigr)e^{-\delta d^{\wedge}(y,\mathcal{C}_j)},\qquad
q_\square(y):=\sup_{y'\in\square^{\zeta}}p(y')\,e^{-2\delta d(y,y')} .
\end{equation}
The following bounds of the stage construction are used in the form stated.
\begin{itemize}
\item The layer bound: for $|\sigma|\le e^{\kappa_1}$, each of $N^{\Delta}_y(\mathbb{H})$,
$N^{\Delta}_y(A')$, and the corresponding norms of the other objects furnished by the stage-layer
theorem, is at most $B_{\mathbb{H}}\,p(y)$, where $B_{\mathbb{H}}$ denotes the constant of this bound.
\item The absorption inequality:
\begin{equation}\label{4.26:eq-regularity-radius-8}
C_1\delta'_j\,e^{-\frac14\delta d^{\wedge}(y,\mathcal{C}_j)}
\le\vartheta\,\rho_{m(y)}\,\alpha_{\cdot,m(y)} .
\end{equation}
\item The local current slot: $\mathcal{J}^{\sharp}[1](\sigma;b)$ is the local current slot of the
stage layer. Clause (iv) of the recursive layer bounds, applied with $\chi\equiv1$ and with the
majorant $p$, gives $\ell^3|\mathcal{J}^{\sharp}[1](\sigma;b)|\le2K_\chi\,p$.
\item Bump functions: for $\chi=\sum_\square c_\square\zeta_\square$ put
\begin{equation*}
a(y):=\sup_{\tilde{\Delta}(y)}\max\{|\chi|,\ell|\nabla\chi|,\ell^2|\nabla\nabla\chi|\},\qquad
\mathsf{q}:=10\sum_\square|c_\square|\,q_\square .
\end{equation*}
Then $\mathsf{q}\in\mathfrak{M}_{2\delta}$, and clause (v) of the same bounds gives
$a\cdot p\le\mathsf{q}$.
\end{itemize}

The constants of the radii are
\begin{align*}
\mathsf{c}_\square&:=\sup_y\sum_\square e^{-\delta d(y,\tilde{\square})},&
\mathsf{K}_r&:=2^9\,\mathsf{c}_\square\,\mathsf{c}_\vartheta
\max\{K_\chi,\;K_{S^{\mathrm{b}}}B_{\mathbb{H}}\},\\
R_0&:=\frac{e^{-C_2\kappa_1}}{2\,\mathsf{K}_r\,\vartheta},&
K_A&:=4\,\mathsf{K}_r\,e^{C_2\kappa_1}.
\end{align*}
Here $K_\chi$ is the constant of the recursive layer bounds used above and $K_{S^{\mathrm{b}}}$ the
constant of the bump lemma below.
For a cube $\square$ of the cover we put
\begin{equation*}
d^-_\square:=d_{\mathrm{sc}}(\square,R^{(l)}),
\end{equation*}
which satisfies $d^-_\square\le\inf_{\square^{\zeta}}d^{\wedge}(\cdot,\mathcal{C}_j)$. We also put
\begin{equation}\label{4.26:eq-regularity-radius-9}
r^\sharp_\square:=R_0\,e^{\frac12\delta d^-_\square},\qquad
\vartheta(\square):=\vartheta\,e^{-\frac12\delta d^-_\square},\qquad\text{so that}\qquad
\frac{2}{r^\sharp_\square}=K_A\,\vartheta(\square).
\end{equation}

\smallskip
The bump lemma of the stage construction is used in the following form, in its notation. Let
$(X',\mathfrak{A},
\mathfrak{B}_\ast)$ be any frame, and let $\mathcal{K}$ be any $1$-form with $N_y(\mathcal{K})\le
\varkappa(y)$ for some $\varkappa\in\mathfrak{M}^{\wedge}_{\delta/2}(\mathfrak{A})$. Assume the
hypotheses (D0)--(D2) of the bump lemma on the frame and its smallness hypothesis
$(S^{\mathrm{b}})$. Then at every $x\in X'^{\flat}$, for $e\in\{2,4\}$,
\begin{equation}\label{4.26:eq-regularity-radius-10}
E_e\bigl(e^{i\eta\mathcal{K}}\mathbb{W}\bigr)\le(1+\pi_\times)\,E_e(\mathbb{W})+\pi\,r_e,
\qquad \pi=\frac{K_{S^{\mathrm{b}}}\,\varkappa(x)}{\tilde{\alpha}_{0,m}},\quad \pi_\times\le\pi,
\end{equation}
pointwise in $x$. Here $m=m(x)$ is the level at $x$, and $r_e=r_e(m,l_x)$ is the shell
radius of entry $e$ at $x$: $r_2(m)$ for $e=2$ and $r_4(m,q)$ for $e=4$. These coincide with the
unweighted shell radii $\mathbf{r}_{e,m}$ at the level of the triple.

The same bound holds in shell form. Let $\mathfrak{h}$ be any $1$-form with
$N_\cdot(\mathfrak{h})\le\varkappa$ for some $\varkappa\in\mathfrak{M}^{\wedge}_{\delta/2}(\mathfrak{B})$,
and assume $(S^{\mathrm{b}})$. Then for any $(\Sigma^{\ast})$-triple (see
\eqref{4.26:eq-graded-cover-a}), any $x\in X'^{\flat}$ and any $e\in\{2,4\}$,
\begin{equation}\label{4.26:eq-regularity-radius-11}
Q_e^+(x)\le Q_e(x)\,(1+\pi_\times)+\pi\,\mathbf{r}_{e,m(x)} .
\end{equation}
\end{definition}

\begin{proof}
We prove items (a) and (b), the consequence of (b) stated in (c), and the identity in
\eqref{4.26:eq-regularity-radius-9}.

\emph{Item (a).} Let $y\in W_{\mathfrak{a}}$. The weight $\omega$ on a live piece is defined through
Ba{\l}aban's construction \cite{B-CMP122a} at the density $k_a$ of the piece. By the bound on the weight
of a live piece, read at the density $k_a$, we have $\omega(y)\le L_0^{2N_0^a}$. By the definition of
$w_\ast$ at the density $k_a$, this right-hand side
equals $w_\ast(k_a)$; this is the first half of \eqref{4.26:eq-regularity-radius-2}. The second half,
$\ell_{\mathfrak{a}}(y)\le k_a$, is a property of the frozen level provided by the relative small-field
machine.

Since $\ell_{\mathfrak{a}}(y)\le k_a$, the radius $\alpha_{\cdot,\ell_{\mathfrak{a}}(y)}$ is one of
the $\alpha_{i,m}$ with $m\le k_a$. It is therefore at most $\bar{\alpha}_{k_a}$. Multiplying by
$\omega(y)\le w_\ast(k_a)$ and using $1+2\beta\le1.4$ for the schedule constant $\beta$ of
Definition~\ref{4.4:def-post-step-space} gives the first inequality of
\eqref{4.26:eq-regularity-radius-3}.

The density $k_a\le k$ is one of the densities over which the maximum in
\eqref{4.26:eq-regularity-radius-a} is taken. Hence $w_\ast(k_a)\bar{\alpha}_{k_a}\le\gamma_0^{(k)}/5$,
which is the second inequality of \eqref{4.26:eq-regularity-radius-3}.

At points $x\in Z^{\mathrm{on}}$, the inequality $5\mathsf{w}^{\mathrm{on}}(x)\alpha_{\cdot,m(x)}\le
\gamma_0^{(k)}$ is the bound of the stage construction for the weight $\mathsf{w}^{\mathrm{on}}$, taken
for each of its two branches. In none of these steps are couplings at two different densities compared:
each product $w_\ast(j)\bar{\alpha}_j$ is bounded by the maximum in
\eqref{4.26:eq-regularity-radius-a} at its own density $j$.

\emph{Item (b): the functions $\alpha_i$ and $\Psi$.} For $\mathsf{t}>0$ put
$\alpha_i(\mathsf{t}):=C_i\,\mathsf{t}^q e^{-\mathsf{t}/2}$ and
$\alpha(\mathsf{t}):=\max_i\alpha_i(\mathsf{t})=C\,\mathsf{t}^qe^{-\mathsf{t}/2}$. Since
$x_m^{-1/2}=e^{-\mathsf{t}_m/2}$ and $\log x_m=\mathsf{t}_m$, we have
$\alpha_{i,m}=\alpha_i(\mathsf{t}_m)$.

Two elementary facts about $\alpha_i$ are used.
\begin{itemize}
\item Monotonicity: $\frac{d}{d\mathsf{t}}\log\alpha_i(\mathsf{t})=\frac{q}{\mathsf{t}}-\frac12$, which
is $\le0$ for $\mathsf{t}\ge2q$. So $\alpha_i$ is decreasing on $[2q,\infty)$.
\item Shift: for $\mathsf{t}>c_\beta$,
\begin{equation*}
\frac{\alpha_i(\mathsf{t}-c_\beta)}{\alpha_i(\mathsf{t})}
=e^{c_\beta/2}\Bigl(1-\frac{c_\beta}{\mathsf{t}}\Bigr)^{q}\le e^{c_\beta/2}=1+\beta_0 ,
\end{equation*}
by the definition $c_\beta=2\ln(1+\beta_0)$.
\end{itemize}
In the same way, $\frac{d}{d\mathsf{t}}\log\Psi(\mathsf{t})=\frac{r+q}{\mathsf{t}}-\frac12$. Hence $\Psi$
is decreasing on $[2(r+q),\infty)$, and for $\mathsf{t}>c_\beta$,
\begin{equation}\label{4.26:eq-regularity-radius-12}
\Psi(\mathsf{t}-c_\beta)\le(1+\beta_0)\,\Psi(\mathsf{t}).
\end{equation}

\emph{Item (b): the range of $\mathsf{t}$.} By hypothesis (g), $\mathsf{t}_j\ge\mathsf{t}_\gamma$ at every
level $j\le k$, in particular at the densities $j$ entering \eqref{4.26:eq-regularity-radius-a} and at
$j=k$. Since $r\ge0$, the range
clause \eqref{4.26:eq-regularity-radius-5} then gives
\begin{equation*}
\mathsf{t}_j-c_\beta\ge2q\quad\text{and}\quad \mathsf{t}_k-c_\beta\ge2(r+q).
\end{equation*}
For the running couplings of the present construction, both inequalities of \cite[(2.6)]{B-CMP119} are
consequences of the two-sided flow inequalities for the running couplings, under the conditions
$\gamma^2\le15/(64c_5)$ and $B_\sharp\gamma^2\le15/49$. The first inequality of \cite[(2.6)]{B-CMP119}
enters once, through \eqref{4.26:eq-regularity-radius-13} below. The second is applied once at the pair
$(m,j)$ and once at the pair $(j,k)$, for levels $m\le j\le k$, and is never iterated along the levels
between them; at a pair $(m,j)$ with $m\le j$ it gives $\mathsf{t}_m\ge\mathsf{t}_j-c_\beta$.

\emph{Item (b), first step: the weight.} For each $j\le k$,
\begin{equation}\label{4.26:eq-regularity-radius-13}
w_\ast(j)<0.37\,L^2(L+1)\,\mathsf{t}_j^{\,r}.
\end{equation}
This is the lemma bounding the weight $w_\ast(j)$ by a power of $\mathsf{t}_j$; its proof uses only the
following inputs:
\begin{itemize}
\item the equality $L^{N_0(j)-1}=R_{j-N_0(j)+1}$ defining $N_0(j)$;
\item the minimality in \cite[(2.5)]{B-CMP119} at the densities $j$ and $j-N_0(j)+1$;
\item the inequality $L_0^2\le L/3$;
\item the first inequality of \cite[(2.6)]{B-CMP119}, applied once, at the pair of levels
$(j-N_0(j)+1,\,j)$.
\end{itemize}
No use is made of \cite[(2.9)]{B-CMP119}, nor of any rule governing how the sizes $R_i$ change with the
level, and a factor logarithmic in the difference of the two levels
is absorbed into the constant $0.37$. Ba{\l}aban's own route in \cite{B-CMP122a}, through the second
inequality of \cite[(2.9)]{B-CMP119} applied once, gives the same bound.

\emph{Item (b), second step: the radii.} For each $j\le k$,
\begin{equation}\label{4.26:eq-regularity-radius-14}
\bar{\alpha}_j\le(1+\beta_0)\,\alpha(\mathsf{t}_j).
\end{equation}
Indeed, let $1\le m\le j$ and let $i$ be any index. The second inequality of \cite[(2.6)]{B-CMP119},
applied once at the pair $(m,j)$, gives $\mathsf{t}_m\ge\mathsf{t}_j-c_\beta$. Both $\mathsf{t}_m$ and
$\mathsf{t}_j-c_\beta$ lie in $[2q,\infty)$, where $\alpha_i$ is decreasing. Hence
\begin{equation*}
\alpha_{i,m}=\alpha_i(\mathsf{t}_m)\le\alpha_i(\mathsf{t}_j-c_\beta)
\le(1+\beta_0)\,\alpha_i(\mathsf{t}_j)\le(1+\beta_0)\,\alpha(\mathsf{t}_j),
\end{equation*}
the second inequality being the shift bound for $\alpha_i$. Taking the maximum over $i$ and over
$1\le m\le j$ gives \eqref{4.26:eq-regularity-radius-14}.

\emph{Item (b), third step: the product.} Multiplying \eqref{4.26:eq-regularity-radius-13} by
\eqref{4.26:eq-regularity-radius-14} and inserting $\alpha(\mathsf{t}_j)=C\mathsf{t}_j^qe^{-\mathsf{t}_j/2}$
gives
\begin{equation*}
w_\ast(j)\,\bar{\alpha}_j<0.37\,(1+\beta_0)\,L^2(L+1)\,C\,\mathsf{t}_j^{\,r+q}e^{-\mathsf{t}_j/2}
\le(1+\beta_0)\,\Psi(\mathsf{t}_j).
\end{equation*}
The second inequality of \cite[(2.6)]{B-CMP119}, applied once at the pair $(j,k)$, gives
$\mathsf{t}_j\ge\mathsf{t}_k-c_\beta$. Since $\mathsf{t}_k-c_\beta\ge2(r+q)$ and $\Psi$ is decreasing on
$[2(r+q),\infty)$, we get $\Psi(\mathsf{t}_j)\le\Psi(\mathsf{t}_k-c_\beta)$. By
\eqref{4.26:eq-regularity-radius-12}, the last quantity is at most $(1+\beta_0)\Psi(\mathsf{t}_k)$.
Altogether,
\begin{equation*}
w_\ast(j)\,\bar{\alpha}_j<(1+\beta_0)\,\Psi(\mathsf{t}_j)
\le(1+\beta_0)\,\Psi(\mathsf{t}_k-c_\beta)\le(1+\beta_0)^2\,\Psi(\mathsf{t}_k).
\end{equation*}

\emph{Item (b), conclusion.} The bound just obtained holds for each product $w_\ast(j)\bar{\alpha}_j$
separately, so the maximum in \eqref{4.26:eq-regularity-radius-a} is taken over products each bounded by
$(1+\beta_0)^2\Psi(\mathsf{t}_k)$. Multiplying by $5$ gives \eqref{4.26:eq-regularity-radius-6}. The
argument does not use the differences $k-j$ or $k-k_a$, and uses the integers $N_0(j)$ only through
\eqref{4.26:eq-regularity-radius-13}.

The function $\bar{\gamma}_0$ is a positive multiple of $\Psi$, so it is decreasing on $[2(r+q),\infty)$.
For every $n$, $\mathsf{t}^{\,n+r+q}e^{-\mathsf{t}/2}\to0$ as $\mathsf{t}\to\infty$, so $\bar{\gamma}_0$
is of degree $-\infty$.

\emph{Item (c).} By \eqref{4.26:eq-regularity-radius-5}, $\mathsf{t}_\gamma\ge2(r+q)$, and by hypothesis
(g), $\mathsf{t}_k\ge\mathsf{t}_\gamma$. Since $\bar{\gamma}_0$ is decreasing on $[2(r+q),\infty)$,
\eqref{4.26:eq-regularity-radius-6} gives
\begin{equation*}
\gamma_0^{(k)}\le\bar{\gamma}_0(\mathsf{t}_k)\le\bar{\gamma}_0(\mathsf{t}_\gamma).
\end{equation*}
Imposing each hypothesis in (c) at $\bar{\gamma}_0(\mathsf{t}_\gamma)$ therefore imposes it at a value not
smaller than $\gamma_0^{(k)}$, at every density $k$.

\emph{The identity in \eqref{4.26:eq-regularity-radius-9}.} By the definitions of $r^\sharp_\square$ and
$R_0$,
\begin{equation*}
\frac{2}{r^\sharp_\square}=\frac{2\cdot2\,\mathsf{K}_r\,\vartheta}{e^{-C_2\kappa_1}}\,
e^{-\frac12\delta d^-_\square}
=4\,\mathsf{K}_r\,e^{C_2\kappa_1}\cdot\vartheta\,e^{-\frac12\delta d^-_\square}
=K_A\,\vartheta(\square),
\end{equation*}
by the definitions of $K_A$ and $\vartheta(\square)$.
\end{proof}

\begin{definition}[Notation for the slot expansion and its collected constant]
\label{4.26:not-slot-notation}
In this section the following letters are used with exactly the meanings fixed in the definition
of the slot expansion:
\begin{itemize}
\item the letter $\mathfrak{N}(c)$ and the norms
  $\|\cdot\|^{A}_{(n)}$;
\item the radius $r_A$ and the letters $K_Q$ and $\varepsilon_S$;
\item $\kappa_A$, with $\kappa_A := \kappa_A(k)$ at level $k$;
\item the letters $A_j^u$, $\mathcal{A}_c$ and $\mathcal{Q}^w_c$;
\item the elements $e = (c, A, Q, \mathsf{y})$, the element $e_{\mathrm{str}}$ and the data
  $(\mathsf{b}, \mathsf{d}, \nu)$ attached to an element, where $\mathsf{b}(e)$ is the set of
  blocks read by $e$, $\mathsf{d}(e)$ its domain and $\mathsf{y}$ its animal;
\item the counts $V$, $\mathsf{p}(S)$ and $\mathsf{r}$;
\item the constants $K_2$, $K_{\flat}$, $c_{\flat}$, $c_J$, and
  $C_0^{\sharp} := 3c_G + 1$.
\end{itemize}
Wherever the separately stated convention governing the letters $\chi^{(j)}$ is in force,
$\chi^{(j)}$ stands for $\chi^{(j)\sharp}$.

We put
\begin{equation}\label{4.26:eq-slot-notation-1}
K_{\flat}' := 2\cdot 52^{d}\cdot (2L)^{d}\cdot e^{2c_{\flat}(16L)^{d}}
\cdot e^{(a_k + 1 + c_J)16^{d}}\cdot K_{\flat}.
\end{equation}
The number $K_{\flat}'$ depends on $L$ and $\kappa$ only. It collects into a single constant the
hull constant, the factor $2\cdot 52^{d}$, the factor $(2L)^{d}$, the factor produced by the bound
$|\mathsf{b}(e)| \le 2(16L)^{d} + |\mathsf{y}|$ on the number of blocks read by an
element $e$, and the factor $e^{c_{\flat}(12L)^{d}}$. It is fixed before $\gamma$ is chosen in
the order of choice of the constants, and no power of $\mathsf{t}$ enters
\eqref{4.26:eq-slot-notation-1}.
\end{definition}

\begin{definition}[Terms, clause entries, radii and the unit rule]\label{4.26:def-clauses-units}
Fix the step $k$ and use the stages and regions of Definition~\ref{4.26:def-stages-regions}; in
particular $\eta=L^{-k}$. For a level $\nu$ we write $\ell_\nu$ for the length, in absolute units,
of the blocks of level $\nu$ (a length, not to be confused with the levels $\ell_l$ of
Definition~\ref{4.26:def-stages-regions} and $\ell_T$ below), so that
$\ell_{\nu}=L^{\nu-\lambda}\ell_{\lambda}$ for levels
$\lambda\le\nu$. We write $\alpha_{0,\nu}$ for the radius of level $\nu$ of \cite[(2.28)]{B-CMP119}
(as in Lemma~\ref{4.4:lem-twist-stability}), $m(y)$ for the level of the shell at the point $y$, and
$\beta_0$ for the constant of the adjacency condition on the radii $\alpha_{0,\cdot}$ among the
conditions~\eqref{4.26:eq-inherited-conditions-a}.
\begin{enumerate}
\item[(i)] \emph{Terms.} $\mathcal{T}$ is the set of terms of the representation of the effective
density after the large-field operation $\mathbf{T}$ at step $k$, that is, the terms of
$\mathbb{E}_k$, $\mathbb{R}_k$ and $\mathbb{B}_k$. It is the union of two classes:
$\mathcal{T}^{\mathbb{E}\mathbb{R}}$, the terms of $\mathbb{E}_k$ and $\mathbb{R}_k$, which carry no
argument $\mathbf{A}$ \cite{B-CMP119}; and $\mathcal{T}^{\mathbb{B}}$, the terms of $\mathbb{B}_k$. Each
$T\in\mathcal{T}^{\mathbb{B}}$ is a sum $T=\sum_{\mathcal{S}}T(X_T,\mathcal{S})$ over structures
$\mathcal{S}$ on its localisation domain $X_T$, all summands being defined on one space
$\mathfrak{U}_T(X_T)$ of pairs $(\mathbb{U},\mathbb{J})$.
\item[(ii)] \emph{Clause entries.} Each $T\in\mathcal{T}$ has a clause set
$\operatorname{Cl}(T)\subset X_T\times\{1,\dots,8\}$, and $\mathfrak{U}_T$ is the set of arguments
at which $Q_a(y)<\theta_{T,a}(y)$ for every $(y,a)\in\operatorname{Cl}(T)$, where the entries
$Q_1,\dots,Q_7$ at $y$ are the components of the clause vector of the small-field bounds at a
triple, whose notation ($U_{p,X}$, $J_{p,X}$, $M\bullet$, $U$, $A'$, $\nabla$) is used here as
defined with that vector:
\begin{equation*}
\begin{gathered}
Q_1=|\partial\mathbb{U}-1|,\quad Q_2=|\partial U_{p,X}(M\bullet\mathbb{U})-1|,\quad
Q_3=|\mathbb{J}|,\quad Q_4=|J_{p,X}(M\bullet\mathbb{U})|,\\
Q_5=|\partial U-1|,\quad Q_6=|A'|,\quad Q_7=|\nabla A'|,
\end{gathered}
\end{equation*}
and the index $8$ labels the cube clause. The entries with $a\in\{2,4\}$ are the \emph{composite}
entries; their clauses are imposed at the structures of $T$, namely at the triples
$(p,X',\mathfrak{s})$ of the family $(\Sigma^{\ast})$ with $X'\subseteq X_T$, where $p$ is a level
(not a plaquette) and $\mathfrak{s}$ a structure on $X'$. Each $\mathfrak{U}_T$ is open. For
$T\in\mathcal{T}^{\mathbb{E}\mathbb{R}}$ it is a union of orbits of the complexified group $G^c$
acting by the gauge transformations of \cite[(1.10)]{B-CMP109}; \cite{B-CMP119} states that these
domains are invariant with respect to those gauge transformations. For $T\in\mathcal{T}^{\mathbb{B}}$,
$\mathfrak{U}_T$ is invariant under the real gauge group $G$ \cite{B-CMP119}.
\item[(iii)] \emph{Radii of a term.} Let $\ell_T$ be the level of $T$. For $y\in X_T$ and
$a\in\{1,\dots,8\}$, $\mathbf{r}^T_a(y)$ is the radius of the clause of $T$ at $y$; it is of level
$\min\{\ell_T,m(y)\}$ (on an older live piece the level is read according to the two cases of the
frozen reading of Definition~\ref{4.26:def-stages-regions}). With $r_a(y)$ the radius of the shell
at $y$, one has
$\mathbf{r}^T_a(y)=r_a(y)$ if $\ell_T\ge m(y)$ and $\mathbf{r}^T_a(y)\ge r_a(y)$ if
$\ell_T<m(y)$, the second by the comparison of the radii of terms with the radii of the shell that
accompanies the shell construction. Thresholds satisfy $\theta_{T,a}\le 2\mathbf{r}^T_a$.
\item[(iv)] \emph{Units at a triple.} Entries and radii are read in absolute lengths, in the units
of the shell construction, whose radii are
\begin{equation*}
r_2(m)=\alpha_{0,m}\,\eta^2\ell_m^{-2},\qquad r_4(m,q)=\alpha_{0,m}\,\ell_m^{-2}\ell_q^{-1};
\end{equation*}
the clause vector of the small-field bounds at a triple writes the same clause in units of the
lattice of level $p$; here it is read in absolute lengths. At a triple $(p,X',\mathfrak{s})$ and for
$a\in\{2,4\}$ put $\lambda:=\min\{\ell_T,m(y)\}$ (on
an older live piece, $\lambda$ is the level of the clause of $T$ at $y$, and the radii
$\mathbf{r}^{T,(p)}_a$, $r^{(p)}_a$ below carry the weight $\omega$ of $T$, respectively of the shell,
according to the frozen reading of Definition~\ref{4.26:def-stages-regions}). The radius of $T$ at
the triple is the radius of level $\lambda$ read at the level $p$ of the triple:
\begin{equation*}
\mathbf{r}^{T,(p)}_2(y):=\alpha_{0,\lambda}\,\eta^2\ell_\lambda^{-2},\qquad
\mathbf{r}^{T,(p)}_4(y):=\alpha_{0,\lambda}\,\ell_\lambda^{-2}\ell_{\min\{p,\lambda\}}^{-1}.
\end{equation*}
Thus $\mathbf{r}^{T,(p)}_a(y)=\mathbf{r}^T_a(y)$ for $p\ge\lambda$, and
$\mathbf{r}^{T,(p)}_4(y)=L^{\lambda-p}\mathbf{r}^T_4(y)$ for $p<\lambda$; this is the scaling of the
current condition $|J_n(M^{\cdot}(\mathbf{U}))|<\alpha_0(L^n\xi)^2$ of \cite{B-CMP109}, see also
\cite[(2.37)]{B-CMP119}. The shell's radius at the triple is
$r^{(p)}_a(y):=r_a(m(y),\min\{p,m(y)\})$ (for $a=2$, read $r_2(m(y))$); here $\min\{p,m(y)\}$ is the
length index which the shell construction uses at $y$ on $\mathfrak{T}$.
\item[(v)] \emph{The unit rule.} For each $(y,a)$ there is one factor: $s_T(y)$, the factor of $T$ at
$y$ (on a drop of scale, the factor of the top shell of $T$), respectively
$s^{(n)}_{T,a}(y):=\theta^{(n)}_{T,a}(y)/\mathbf{r}^T_a(y)$, where $\theta^{(n)}_{T,a}$ is the
threshold of the clause of $T$ at stage $n$. A threshold is a factor times a unit:
\begin{equation}\label{4.26:eq-clauses-units-b}
\begin{gathered}
\theta_{T,a}(y):=s_T(y)\,\mathbf{r}^T_a(y)\qquad
\bigl(a\in\{2,4\},\ y\in X_T\setminus X_T^{\sim2}\bigr),\\
\theta_{T,a}(y;p):=s_T(y)\,\mathbf{r}^{T,(p)}_a(y),\qquad
\theta^{(n)}_{T,a}(y;p):=s^{(n)}_{T,a}(y)\,\mathbf{r}^{T,(p)}_a(y).
\end{gathered}
\end{equation}
The first line concerns the points $y\in X_T$ at which $\operatorname{Cl}(T)$ contains no composite
clause; in Ba{\l}aban's construction these are the points of the two outer layers
$X_T\setminus X_T^{\sim2}$, where $X_T^{\sim2}$ is $X_T$ with its two outer layers removed. This
fixes the threshold only for the composite clause that the present construction imposes on the
boxes of its cover; it says nothing about the thresholds of the clauses in Ba{\l}aban's construction.
\item[(vi)] \emph{Comparison of units.} With $m=m(y)$, at $a=4$,
\begin{equation}\label{4.26:eq-clauses-units-a}
\frac{r^{(p)}_4(y)}{\mathbf{r}^{T,(p)}_4(y)}
=\frac{\alpha_{0,m}}{\alpha_{0,\lambda}}\,L^{-2(m-\lambda)}\,
\frac{\ell_{\min\{p,\lambda\}}}{\ell_{\min\{p,m\}}}
\le\Bigl(\frac{1+\beta_0}{L^2}\Bigr)^{m-\lambda}\le1,
\end{equation}
and at $a=2$ the same holds without the last factor of the middle expression. Let $\mathfrak{a}$ be
an arrest live at the current density, with region $W_{\mathfrak{a}}$, step $k_a$, final stage
$n^+_a$, frozen terms $\mathfrak{F}(\mathfrak{a})$ and level $\mathfrak{c}(y)$ at
$y$, and let $\omega(y)$ be the weight at $y$; the base $L_0$ of the weights is subject to the
condition $L_0^2\le L/3$. On $W_{\mathfrak{a}}$: for a term $T$ created at a step $k'>k_a$ whose
clause set is an arrest-stage set $\mathfrak{S}^{(k'),\mathrm{arrest}}_l$ (the set held at stage $l$
of the application of $\mathbf{R}$ at step $k'$ when that application is arrested), both radii carry
$\omega(y)$ and the quotient is the one
in~\eqref{4.26:eq-clauses-units-a}. For a frozen term
$T=\mathbb{C}^{(i)}_{k_a}\in\mathfrak{F}(\mathfrak{a})$, $i<n^+_a$, whose clauses were fixed at stage
$i$ on the level $\lambda\le m$ that $y$ had at that stage, with the weight
$L_0^{2(\lambda-\mathfrak{c}(y))^+}$, and with $p_a\ge1$ the exponent of $L$ in the unweighted
quotient of the radii at that stage,
\begin{equation}\label{4.26:eq-clauses-units-1}
\begin{aligned}
\frac{r^{(p)}_a(y)}{\mathbf{r}^{T,(p)}_a(y)}
&\le L_0^{2(m-\max\{\lambda,\mathfrak{c}(y)\})}\Bigl(\frac{1+\beta_0}{L^{p_a}}\Bigr)^{m-\lambda}\\
&\le\Bigl(\frac{(1+\beta_0)L_0^2}{L}\Bigr)^{m-\lambda}\le 0.381^{m-\lambda}\le1 .
\end{aligned}
\end{equation}
For a term $T$ whose clauses were fixed before the arrest, with $\lambda\le\mathfrak{c}(y)$,
\begin{equation}\label{4.26:eq-clauses-units-2}
\frac{r^{(p)}_a(y)}{\mathbf{r}^{T,(p)}_a(y)}
\le\Bigl(\frac{(1+\beta_0)L_0^2}{L^2}\Bigr)^{m-\mathfrak{c}(y)}
\Bigl(\frac{1+\beta_0}{L^2}\Bigr)^{\mathfrak{c}(y)-\lambda}\le1 .
\end{equation}
\item[(vii)] \emph{Reading at triples.} Wherever, in this chapter, a composite entry at a triple is
compared with $\mathbf{r}^T_a$, $\theta_{T,a}$, $\theta^{(n)}_{T,a}$ or $r_a$, these are read as
$\mathbf{r}^{T,(p)}_a$, $\theta_{T,a}(y;p)$, $\theta^{(n)}_{T,a}(y;p)$ and $r^{(p)}_a$ respectively.
\item[(viii)] \emph{The source.} The source space is
$\mathfrak{S}_{\mathrm{src}}=\tilde{U}^c_k(Y_4,\tilde{\alpha}_0,\tilde{\alpha}_1)^{(\Sigma^{\ast})}$,
the space $\tilde{U}^c_k(Y_4,\tilde{\alpha}_0,\tilde{\alpha}_1)$ on which the derived-entry family
$(\Sigma^{\ast})$ is defined (with $Y_4$, $\tilde{\alpha}_0$, $\tilde{\alpha}_1$ as fixed together
with that family), read through the derived-entry family $(\Sigma^{\ast})$,
with thresholds $\theta_{\mathrm{src},a}:=s^{(k)}_m r_a$ for $a\le7$ and
$\theta_{\mathrm{src},8}:=r_8$. On live pieces the source is its arrested version
$\tilde{U}^{\mathrm{arrest}}_k$, with $\theta_{\mathrm{src},a}=F^{(k)}r_a$ there, according to the
frozen reading of Definition~\ref{4.26:def-stages-regions}.
\item[(ix)] \emph{The newborn number.} The number $\theta_S$ of the schedule for newborn terms
in~\eqref{4.26:eq-room-table-a} is written $\mathsf{c}_e:=\theta_S$; it is distinct from the size
constants $c^{\mathrm{sz}}_T$ of the terms.
\item[(x)] \emph{Stage classes.} At stage $l$, the class $\mathcal{T}_2$ consists of
the terms of $\mathcal{T}^{\mathbb{B}}$, the chain terms $\mathbb{C}^{(i)}(S,\mathcal{S}')$ with
$i<l$, and the tails of $\mathcal{T}^{\mathbb{E}\mathbb{R}}$, that is, the terms
$T\in\mathcal{T}^{\mathbb{E}\mathbb{R}}$ with $X_T$ contained in no $\tilde{\square}^2$,
$\square\in\mathcal{C}^1$; $\mathcal{T}_1$ consists of the other terms of
$\mathcal{T}^{\mathbb{E}\mathbb{R}}$. The union of $\mathcal{T}^{\mathbb{B}}$ and
the chain terms is contained in $\mathcal{T}_2$.
\end{enumerate}
\end{definition}

\begin{proof}[Proof of \eqref{4.26:eq-clauses-units-a}, \eqref{4.26:eq-clauses-units-1}
and \eqref{4.26:eq-clauses-units-2}]
We first record two numerical facts. The condition
$L_0^2\le L/3$ and the numerical bound of the
unit-conversion condition in~\eqref{4.26:eq-graded-box-geometry-a} give
\begin{equation}\label{4.26:eq-clauses-units-3}
\frac{(1+\beta_0)L_0^2}{L}\le 0.381 .
\end{equation}
Since $L\ge1$, it follows that
\begin{equation*}
\frac{(1+\beta_0)L_0^2}{L^2}\le\frac{(1+\beta_0)L_0^2}{L}\le0.381<1 .
\end{equation*}

In all cases $\lambda\le m$: off the live pieces because $\lambda=\min\{\ell_T,m(y)\}$; for the
frozen terms by the choice of $\lambda$ as a level $\lambda\le m$ that $y$ had at stage $i$; and on a
live piece because the level of the clause of $T$ at $y$ does not exceed the shell's level $m(y)$ in
the frozen reading of Definition~\ref{4.26:def-stages-regions}.

\emph{The comparison~\eqref{4.26:eq-clauses-units-a}.} Let $a=4$. By (iv),
\begin{gather*}
r^{(p)}_4(y)=r_4(m,\min\{p,m\})=\alpha_{0,m}\ell_m^{-2}\ell_{\min\{p,m\}}^{-1},\\
\mathbf{r}^{T,(p)}_4(y)=\alpha_{0,\lambda}\ell_\lambda^{-2}\ell_{\min\{p,\lambda\}}^{-1},
\end{gather*}
so that
\begin{equation*}
\frac{r^{(p)}_4(y)}{\mathbf{r}^{T,(p)}_4(y)}
=\frac{\alpha_{0,m}}{\alpha_{0,\lambda}}\Bigl(\frac{\ell_\lambda}{\ell_m}\Bigr)^2
\frac{\ell_{\min\{p,\lambda\}}}{\ell_{\min\{p,m\}}} .
\end{equation*}
Since $\lambda\le m$ and $\ell_m=L^{m-\lambda}\ell_\lambda$, we have
$(\ell_\lambda/\ell_m)^2=L^{-2(m-\lambda)}$, which is the equality
in~\eqref{4.26:eq-clauses-units-a}. At $a=2$ the factor $\eta^2$ appears in both radii and cancels,
and the quotient is $(\alpha_{0,m}/\alpha_{0,\lambda})L^{-2(m-\lambda)}$, the same expression without
the last factor. For the first inequality: from $\lambda\le m$ we get
$\min\{p,\lambda\}\le\min\{p,m\}$, and since block lengths increase with the level, the last factor
is at most $1$; the adjacency condition on the radii in~\eqref{4.26:eq-inherited-conditions-a} gives
$\alpha_{0,m}/\alpha_{0,\lambda}\le(1+\beta_0)^{m-\lambda}$. Multiplying,
the quotient is at most $\bigl((1+\beta_0)/L^2\bigr)^{m-\lambda}$. For the last inequality, since
$m-\lambda\ge0$ it suffices that $(1+\beta_0)/L^2\le1$; since $L_0\ge1$, the numerical facts above
give $(1+\beta_0)/L^2\le(1+\beta_0)L_0^2/L^2\le0.381$.

\emph{Terms with an arrest-stage clause set on $W_{\mathfrak{a}}$.} For a term created at a step
$k'>k_a$ whose clause set is an arrest-stage set $\mathfrak{S}^{(k'),\mathrm{arrest}}_l$, both
$r^{(p)}_a(y)$ and $\mathbf{r}^{T,(p)}_a(y)$
carry the factor $\omega(y)$, which cancels in the quotient; the quotient is therefore the unweighted
one, and~\eqref{4.26:eq-clauses-units-a} applies to it.

\emph{Frozen terms, \eqref{4.26:eq-clauses-units-1}.} Let
$T=\mathbb{C}^{(i)}_{k_a}\in\mathfrak{F}(\mathfrak{a})$, $i<n^+_a$. At stage $i$ its clauses were fixed
on the level $\lambda\le m$ that $y$ had at that stage, and the weight at that stage is
$L_0^{2(\nu-\mathfrak{c})}$ at level $\nu$ whenever this number exceeds $1$; hence the radius of the
term carries the weight $L_0^{2(\lambda-\mathfrak{c}(y))^+}$. The shell's radius at $y$ carries
$\omega(y)$, and by the definition of $\mathfrak{c}(y)$ in Definition~\ref{4.26:def-stages-regions},
$\omega(y)=L_0^{2(m-\mathfrak{c}(y))}$. Since
$2(m-\mathfrak{c}(y))-2(\lambda-\mathfrak{c}(y))^+=2(m-\max\{\lambda,\mathfrak{c}(y)\})$ (for
$\lambda\ge\mathfrak{c}(y)$ both sides equal $2(m-\lambda)$, for $\lambda<\mathfrak{c}(y)$ both equal
$2(m-\mathfrak{c}(y))$), the quotient of the weights is
\begin{equation*}
\omega(y)\,L_0^{-2(\lambda-\mathfrak{c}(y))^+}=L_0^{2(m-\max\{\lambda,\mathfrak{c}(y)\})}.
\end{equation*}
This, combined with the quotient of the unweighted radii at the
frozen stage, in which $L^{p_a}$ stands in place of $L^2$, gives the first inequality
in~\eqref{4.26:eq-clauses-units-1}. Next, $\max\{\lambda,\mathfrak{c}(y)\}\ge\lambda$ and $L_0\ge1$
give $L_0^{2(m-\max\{\lambda,\mathfrak{c}(y)\})}\le L_0^{2(m-\lambda)}$, and $p_a\ge1$, $L\ge1$ give
$L^{p_a}\ge L$; hence the first expression is at most
$\bigl((1+\beta_0)L_0^2/L\bigr)^{m-\lambda}$. By~\eqref{4.26:eq-clauses-units-3} this is at most
$0.381^{m-\lambda}$, which is at most $1$ because $m-\lambda\ge0$.

\emph{Terms fixed before the arrest, \eqref{4.26:eq-clauses-units-2}.} Let $\lambda\le
\mathfrak{c}(y)$. The radius of $T$ carries no weight exceeding $1$, while the shell's radius carries
$\omega(y)=L_0^{2(m-\mathfrak{c}(y))}$, which is the value of the weight quotient
$L_0^{2(m-\max\{\lambda,\mathfrak{c}(y)\})}$ displayed in the preceding case at
$\lambda\le\mathfrak{c}(y)$. The unweighted
quotient is at most $\bigl((1+\beta_0)/L^2\bigr)^{m-\lambda}$
by~\eqref{4.26:eq-clauses-units-a}. Hence the quotient is at most
\begin{equation*}
L_0^{2(m-\mathfrak{c}(y))}\Bigl(\frac{1+\beta_0}{L^2}\Bigr)^{m-\lambda}
=\Bigl(\frac{(1+\beta_0)L_0^2}{L^2}\Bigr)^{m-\mathfrak{c}(y)}
\Bigl(\frac{1+\beta_0}{L^2}\Bigr)^{\mathfrak{c}(y)-\lambda},
\end{equation*}
where the equality comes from writing $m-\lambda=(m-\mathfrak{c}(y))+(\mathfrak{c}(y)-\lambda)$. This
is the first inequality in~\eqref{4.26:eq-clauses-units-2}. For the second, if
$m\ge\mathfrak{c}(y)$, both exponents are non-negative; the first base is at most $0.381$ by the
numerical facts above and the second is at most $1$ as in the proof
of~\eqref{4.26:eq-clauses-units-a}, so the product is at most $1$. If $m<\mathfrak{c}(y)$, the
left-hand form
$L_0^{2(m-\mathfrak{c}(y))}\bigl((1+\beta_0)/L^2\bigr)^{m-\lambda}$ is a product of a number at most
$1$ (as $L_0\ge1$) and a non-negative power of a number at most $1$ (as $m\ge\lambda$), hence again at
most $1$.
\end{proof}

\begin{definition}[The spend, the four radii and the window condition]\label{4.26:def-radii-window}
We use the region $Y$, the cover $\{\square\}$, the partition of unity $\zeta_\square$, the enlargements
$\tilde{\square}$, $\tilde{\square}^2$, $\square_0=\tilde{\square}^5$ and $\square^{\zeta}\subset\square$,
the rate $\delta$, the ratio $w=M/M_1$, the distances $d^-_\square$ and $d^{\wedge}$, the seminorms
$N_y$, and the run-free majorant $r^\sharp_\square$ of the regularity radius, all as in
Definition~\ref{4.26:def-regularity-radius}. We also use the terms $T$, their localization domains
$X_T$, their clause sets $\operatorname{Cl}(T)$, thresholds $\theta_{T,a}$ and radii
$\mathbf{r}^T_a$, all as in Definition~\ref{4.26:def-clauses-units}. Here $\vartheta$ is the
smallness parameter of the expansion,
$\vartheta(\square):=\vartheta e^{-\frac12\delta d^-_\square}$ is its box-dependent form,
$\ell=\ell_T$ is the level of the term $T$, and $d_\ell(X,X')$ denotes the distance of two sets
$X$, $X'$ measured in the unit of level $\ell$.
\begin{itemize}
\item[(a)] \emph{The spend.} Let $\mathsf{c}_e=\theta_S$ be the newborn number; in the order of
choice of the constants of this section it satisfies $\theta_S\ge\frac12$. Put
\begin{equation*}
\mathsf{w}_0:=\tfrac{1}{12}\min\{\tfrac12,\mathsf{c}_e\},\qquad
\hat{c}:=\min\{\tfrac12,\mathsf{c}_e\},\qquad \theta:=\tfrac12,\qquad
\vartheta^{\natural}:=\vartheta/\mathsf{w}_0 .
\end{equation*}
The number $\mathsf{w}_0$ is \emph{the spend}.
\item[(b)] \emph{The four radii.} Let $t=(t_\square)_\square$ and $\tau$ be the complex parameters
of the family $\mathcal{W}_T(t,\tau)$; the remainder pieces of a near pair carry in addition a
parameter $\tau_1$. These parameters range over the following discs:
\begin{equation}\label{4.26:eq-radii-window-a}
\begin{aligned}
&\text{(i)} && |t_\square|\le R^{\natural}_\square
  :=\mathsf{w}_0R_0e^{\frac12\delta d^-_\square}=\mathsf{w}_0r^\sharp_\square,\\
&\text{(ii)} && |\tau|\le\mathsf{w}_0r_T(\square),\qquad
  r_T(\square):=\frac{e^{\frac14\delta d_\ell(X_T,\square^{\zeta})}}{14K^+\vartheta(\square)},\\
&\text{(iii)} && |\tau|\le\mathsf{w}_0r^{\mathrm{nr}}_\square,\qquad
  r^{\mathrm{nr}}_\square:=\frac{1}{K^{\mathrm{nr}}\vartheta(\square)},\\
&\text{(iv)} && |\tau|\le\mathsf{w}_0r^{\mathrm{rem}}_T(\square),\qquad
  r^{\mathrm{rem}}_T(\square):=\frac{e^{\frac14\delta D}}{14K^{\mathrm{rem}}\vartheta(\square)},\\
& && |\tau_1|\le\frac{\mathsf{w}_0e^{\frac14\delta D}}{28K^{\mathrm{rem}}\vartheta}.
\end{aligned}
\end{equation}
In (i), which is the radius for the terms of the M-class, the parameters $t_\square$ of all boxes
with $\tilde{\square}\subset Y$ range over their discs
simultaneously, and $t_\square:=0$ for every other box. With
$\vartheta^{\natural}(\square):=\vartheta^{\natural}e^{-\frac12\delta d^-_\square}$ one has
$2/R^{\natural}_\square=K_A\vartheta^{\natural}(\square)$.

In (ii) and (iii), $T\in\mathcal{T}_1$ with $X_T\not\subset\tilde{\square}^2$ and with
$X_T\subset\tilde{\square}^2$ respectively. Here $K^+$ is the constant of the off-box estimate for
the terms of $\mathcal{T}_1$, that is, the estimate for the terms with $X_T\not\subset\tilde{\square}^2$
in which the radius $r_T(\square)$ of (ii) is read; that estimate is used under the condition
$17K^+\vartheta\le1$.

In (iv), $T$ runs over the remainder pieces of a near pair. We put
\begin{equation*}
D:=d_\ell\bigl(X_T,\operatorname{supp}(1-\zeta'_\square)\bigr),
\end{equation*}
and for these pieces $D\ge wL^{(i-\ell)^+}$.

All radii in \eqref{4.26:eq-radii-window-a} are numbers.
\item[(c)] \emph{The constants and the window.} Let
$K_{\mathrm{cv}}=K_{\mathrm{cv}}(d,L,B_3)\ge1$ be the constant of Ba{\l}aban's entry conversion
\cite[(3.48)--(3.52)]{B-CMP109}: a layer on $\square_0$ of size at most $s\alpha_{\cdot,m}$, $s\le1$,
moves each entry $a$ of the box pair by at most $K_{\mathrm{cv}}s\,\mathbf{r}^T_a$. Put
$K^{\mathrm{nr}}:=66B_\sharp K_{\mathrm{cv}}$. Let $K^{\mathrm{rem}}$ be any constant with which the
near-pair estimate holds on the remainder data of a near pair. Let
$K_{\min}:=\max\{K_A,\,28K^+,\,6K^{\mathrm{nr}},\,28K^{\mathrm{rem}}\}$. The \emph{window condition} is
\begin{equation}\label{4.26:eq-radii-window-b}
K_{\min}\vartheta\le\mathsf{w}_0 ;
\end{equation}
it is a ceiling on $\gamma$.
\item[(d)] \emph{The third radius of a term.} For a term $T$ with at least one clause
$(y,a)\in\operatorname{Cl}(T)$ with $y\in\mathfrak{T}$, put
\begin{equation}\label{4.26:eq-radii-window-c}
\begin{aligned}
\alpha_3(T)&:=\min\Bigl\{\frac{\operatorname{room}_T\,\alpha_{\cdot,\ell_T}}{12B_3},\
  \frac{\beta L^{-2}\alpha_{0,\cdot}}{2B_3},\ \frac{\alpha_{1,\cdot}}{3B_3}\Bigr\},\\
\operatorname{room}_T&:=\inf_{\substack{(y,a)\in\operatorname{Cl}(T),\ y\in\mathfrak{T}\\ a\notin\{5,8\}}}
  \frac{\theta_{T,a}-\theta_{\mathrm{src},a}}{\mathbf{r}^T_a}\ \ge\ \beta\hat{c},
\end{aligned}
\end{equation}
where an infimum over the empty set is $+\infty$.
The radii $\alpha$ in this display are those of the local shell in which the box sits. For a term
with no clause $(y,a)\in\operatorname{Cl}(T)$ with $y\in\mathfrak{T}$, $\alpha_3(T)$ is the third
radius of the term-wise form of Ba{\l}aban's bound
\cite[(3.54)]{B-CMP109}.
\end{itemize}
The following properties hold.
\begin{itemize}
\item[(e1)] $\mathsf{w}_0=\frac1{24}$ and $\hat{c}=\frac12$.
\item[(e2)] Under \eqref{4.26:eq-radii-window-b} the following hold:
$R^{\natural}_\square\ge2$, $\mathsf{w}_0r_T(\square)\ge2$, $\mathsf{w}_0r^{\mathrm{nr}}_\square\ge6$,
$\mathsf{w}_0r^{\mathrm{rem}}_T(\square)\ge2$, and $17K^+\vartheta\le1$.
\item[(e3)] Assume \eqref{4.26:eq-radii-window-b}. Let $\square$ be a box of $\mathfrak{T}$ and
$T\in\mathcal{T}_1$ with $X_T\subset\tilde{\square}^2$. For $|\tau|\le\mathsf{w}_0r^{\mathrm{nr}}_\square$,
the layer moves each entry $a$ of the box pair by at most $\frac13\mathsf{w}_0\rho_m\mathbf{r}^T_a$.
For $|\tau|\le1$ it moves each entry by at most $\frac{\mathsf{w}_0}{18}\rho_m\mathbf{r}^T_a$.
\item[(e4)] $\alpha_3(T)$ meets the restrictions $B_3\alpha_3<\frac12\alpha_1$ and $2B_3\alpha_3\le\beta L^{-2}\alpha_0$ imposed on the ball $B'$ in \cite{B-CMP109}. Moreover $2B_3\alpha_3(T)\le\frac16\operatorname{room}_T\alpha_{\cdot,\ell_T}$.
\item[(e5)] With $\alpha_3(T)$ as in (d), the constant of \cite[(3.35)]{B-CMP109} that enters
\cite[(3.54)]{B-CMP109} becomes
\begin{equation*}
\Bigl(O(1)MB_3^2\mathsf{a}\cdot\max\bigl\{\tfrac{24}{\beta},\tfrac{2L^2}{\beta},3\bigr\}\Bigr)^5 ,
\end{equation*}
a number depending on $(L,M,B_3)$; it is included in $C^{\natural}_R$ and fixed before $\gamma$.
\end{itemize}
\end{definition}

\begin{proof}
\emph{The spend.} Since $\mathsf{c}_e=\theta_S\ge\frac12$, we have $\min\{\frac12,\mathsf{c}_e\}=\frac12$. Hence $\hat{c}=\frac12$ and $\mathsf{w}_0=\frac12\cdot\frac1{12}=\frac1{24}$. This is (e1).

\emph{The first radius.} By the definitions of $R^{\natural}_\square$ and $\vartheta^{\natural}(\square)$,
\begin{equation*}
\frac{2}{R^{\natural}_\square}=\frac{2e^{-\frac12\delta d^-_\square}}{\mathsf{w}_0R_0}
\quad\text{and}\quad
K_A\vartheta^{\natural}(\square)=\frac{K_A\vartheta e^{-\frac12\delta d^-_\square}}{\mathsf{w}_0}.
\end{equation*}
Thus the relation $2/R^{\natural}_\square=K_A\vartheta^{\natural}(\square)$ stated in (b) is the same as
$K_A\vartheta R_0=2$, where $R_0$ is the constant of the run-free majorant
$r^\sharp_\square=R_0e^{\frac12\delta d^-_\square}$ of Definition~\ref{4.26:def-regularity-radius}.
From the relation of (b) we obtain
\begin{equation*}
R^{\natural}_\square=\frac{2}{K_A\vartheta^{\natural}(\square)}
=\frac{2\mathsf{w}_0e^{\frac12\delta d^-_\square}}{K_A\vartheta}\ge\frac{2\mathsf{w}_0}{K_A\vartheta}\ge2 .
\end{equation*}
The first inequality holds because $d^-_\square\ge0$. The second holds because $K_A\vartheta\le K_{\min}\vartheta\le\mathsf{w}_0$ by \eqref{4.26:eq-radii-window-b}.

\emph{The other radii.} The exponential factors in (ii) and (iv) of \eqref{4.26:eq-radii-window-a} are at least $1$, since the distances are non-negative. We also use $\vartheta(\square)\le\vartheta$. Then \eqref{4.26:eq-radii-window-b} gives $28K^+\vartheta\le\mathsf{w}_0$, $6K^{\mathrm{nr}}\vartheta\le\mathsf{w}_0$ and $28K^{\mathrm{rem}}\vartheta\le\mathsf{w}_0$. Therefore
\begin{equation*}
\mathsf{w}_0r_T(\square)\ge\frac{\mathsf{w}_0}{14K^+\vartheta}\ge2,\qquad
\mathsf{w}_0r^{\mathrm{nr}}_\square\ge\frac{\mathsf{w}_0}{K^{\mathrm{nr}}\vartheta}\ge6,\qquad
\mathsf{w}_0r^{\mathrm{rem}}_T(\square)\ge\frac{\mathsf{w}_0}{14K^{\mathrm{rem}}\vartheta}\ge2 .
\end{equation*}
Finally
\begin{equation*}
17K^+\vartheta\le28K^+\vartheta\le\mathsf{w}_0=\tfrac1{24}\le1 .
\end{equation*}
This proves (e2).

\emph{The constant $K^{\mathrm{nr}}$: the first factor $11B_\sharp$.} The constant $11B_\sharp$ is the first-order constant of the first-order layer bound at pairs of the kernel-regularity space. Let $y\in\square^{\zeta}$, $m=m(y)$ and $\|b\|_\infty\le1.1C_1\delta'_j$. Then
\begin{equation}\label{4.26:eq-radii-window-1}
N_y(\zeta_\square\mathbb{H}^w)\le2\cdot5\cdot B_\sharp\,p(y)
\le11B_\sharp\,\vartheta(\square)\rho_m\alpha_{\cdot,m},
\end{equation}
where $\mathbb{H}^w$ is the stage layer and $p(y)$ is the pointwise majorant of that bound.
This bound comes from the first-order layer bound, applied at
$(\mathbb{U}^w_\square,\mathbb{J})\in\mathfrak{K}_P$. It also uses the bound $a_\square\le5$, the absolute bound on the layer, and the inequality $d^-_\square\le d^{\wedge}(\square^{\zeta},\mathcal{C}_j)$.

The membership $(\mathbb{U}^w_\square,\mathbb{J})\in\mathfrak{K}_P$ holds at boxes in $\mathfrak{T}$, for two reasons. First, the space $\mathfrak{N}^{(n)}$ bounds the composite entries of the box pair at $(p_\square,\square_0,\mathfrak{s})$; this is the composite shell clause in \eqref{4.26:eq-regularity-space-a}. Second, $\mathfrak{K}_P=\mathfrak{K}_{109}(\mathfrak{B};\gamma_0)$ imposes smallness conditions only.

\emph{The entry conversion.} We use Ba{\l}aban's entry conversion \cite[(3.48)--(3.52)]{B-CMP109} in the form stated there, as it is used in the term-wise form of Ba{\l}aban's bound \cite[(3.54)]{B-CMP109}. Its content is the following. Let $\mathbf{A}$ be a layer on $\square_0$ with $N_y(\mathbf{A})\le s\alpha_{\cdot,m}$, where $s\le1$. Then $\mathbf{A}$ obeys \cite[(3.31)]{B-CMP109} with $s\alpha_2$ in place of $\alpha_2$, and it moves each entry $a$ of the box pair by at most $K_{\mathrm{cv}}\,s\,\mathbf{r}^T_a$. Here $K_{\mathrm{cv}}=K_{\mathrm{cv}}(d,L,B_3)\ge1$ is a number fixed before $\gamma$.

\emph{The factor $2$: variation of the one-step rooms over the box.} A box $\square_0$ may meet two
adjacent shells (clause (a) of \eqref{4.26:eq-near-box-geometry-b}) or two adjacent frozen levels (the
frozen-level analogue of that clause). In either case the levels $m(y')$, $y'\in\square_0$, lie in $\{m-1,m,m+1\}$, and the size of the layer over $\square_0$ carries
\begin{equation*}
\bar{\rho}_\square:=\max_{y'\in\square_0}\rho_{m(y')}\le2\rho_m ,
\end{equation*}
because the one-step rooms satisfy $\rho^{(l)}_{m\pm1}\le2\rho^{(l)}_m$.

\emph{The uniform parameter.} We apply $K_{\mathrm{cv}}$ at the uniform value
$s:=\sup_{\square_0}N/\alpha$, measured in the unit of the level $\ell$ of $T$. Passing to this unit only improves the bound. Indeed, $\ell\le\mathfrak{c}\le$ the levels on $\square_0$, and by the change-of-level property of $N_y$ the passage from the level $\mathfrak{c}(y')$ to the level $\ell$ carries the factor
\begin{equation*}
\Bigl(\frac{1+\beta_0}{L}\Bigr)^{\mathfrak{c}(y')-\ell}\le1 .
\end{equation*}

\emph{The factor $3$ and the conclusion (e3).} The layer attached to $T$ at the parameter $\tau$ has, over $\square_0$, size at most $|\tau|$ times the right-hand side of \eqref{4.26:eq-radii-window-1} with $\rho_m$ replaced by $\bar{\rho}_\square$. Hence, for $|\tau|\le\mathsf{w}_0r^{\mathrm{nr}}_\square$,
\begin{equation*}
s\le\frac{\mathsf{w}_0}{K^{\mathrm{nr}}\vartheta(\square)}\cdot11B_\sharp\vartheta(\square)\cdot2\rho_m
=\frac{22B_\sharp\,\mathsf{w}_0\rho_m}{66B_\sharp K_{\mathrm{cv}}}
=\frac{\mathsf{w}_0\rho_m}{3K_{\mathrm{cv}}} .
\end{equation*}
By the entry conversion this member moves each entry of the box pair by at most
$K_{\mathrm{cv}}s\,\mathbf{r}^T_a\le\frac13\mathsf{w}_0\rho_m\mathbf{r}^T_a$. This explains the decomposition $66=2\cdot3\cdot11$.

Now consider the undilated member of the family $\mathcal{W}_T$, whose parameter ranges over
$|\tau|\le1$. By (e2), $|\tau|\le1\le\frac16\mathsf{w}_0r^{\mathrm{nr}}_\square$, so $s$ is at most one sixth of the bound just obtained. The move is therefore at most
\begin{equation*}
\tfrac16\cdot\tfrac13\,\mathsf{w}_0\rho_m\mathbf{r}^T_a=\tfrac{\mathsf{w}_0}{18}\,\rho_m\mathbf{r}^T_a .
\end{equation*}
This proves (e3).

\emph{The constant $K^{\mathrm{rem}}$.} The remainder data of a near pair are of two kinds. They are either data bounded by $2\varrho$ times the first-order layer bound, or the quantity $Q$ of the layer of the near pair; in both cases $\hat{p}$ is re-regularised on the set of $\square_0$. The near-pair estimate applies to these data as it stands. We take for $K^{\mathrm{rem}}$ any constant with which that estimate holds on these data. It is a number fixed before $\gamma$, and it enters $K_{\min}$.

\emph{The third radius.} The restrictions placed on $\alpha_3$ in \cite{B-CMP109} for the ball $B'$ are $B_3\alpha_3<\frac12\alpha_1$ and $2B_3\alpha_3\le\beta L^{-2}\alpha_0$; the estimate there is stated for $\alpha_3$ satisfying all the restrictions. By \eqref{4.26:eq-radii-window-c},
\begin{equation*}
B_3\alpha_3(T)\le\tfrac13\alpha_{1,\cdot}<\tfrac12\alpha_{1,\cdot},\qquad
2B_3\alpha_3(T)\le\beta L^{-2}\alpha_{0,\cdot}.
\end{equation*}
Both restrictions therefore hold. The minimum in \eqref{4.26:eq-radii-window-c} is needed. If one took $\alpha_3(T)=\operatorname{room}_T\alpha_{\cdot,\ell_T}/(12B_3)$ alone, the second restriction would read
\begin{equation*}
L^2\le\frac{6\beta\alpha_0}{\operatorname{room}_T\alpha_{\cdot,\ell_T}} ,
\end{equation*}
which is an upper bound on $L$, that is, a cap in the sense of Definition~\ref{2.3:def-cap}. The minimum avoids this cap.

The ball $B'$ still takes at most
\begin{equation*}
2B_3\alpha_3(T)\le\frac{2B_3\operatorname{room}_T\alpha_{\cdot,\ell_T}}{12B_3}
=\tfrac16\operatorname{room}_T\alpha_{\cdot,\ell_T},
\end{equation*}
which completes (e4).

With this $\alpha_3(T)$, the constant of \cite[(3.35)]{B-CMP109} that enters \cite[(3.54)]{B-CMP109} becomes
\begin{equation*}
\Bigl(O(1)MB_3^2\mathsf{a}\cdot\max\bigl\{\tfrac{24}{\beta},\tfrac{2L^2}{\beta},3\bigr\}\Bigr)^5 .
\end{equation*}
The factor $(s_\ell/s_i)^{4+\alpha'}$ of the term-wise form of Ba{\l}aban's bound \cite[(3.54)]{B-CMP109} carries the ratios of the radii $\alpha$ as in \cite{B-CMP109}, unchanged. The constant displayed above is a number depending on $(L,M,B_3)$. It is included in $C^{\natural}_R$ and fixed before $\gamma$. This is (e5).

\emph{The room.} If $T$ has no clause $(y,a)\in\operatorname{Cl}(T)$ with $y\in\mathfrak{T}$ and
$a\notin\{5,8\}$, then $\operatorname{room}_T=+\infty$ and there is nothing to prove. Let
$(y,a)\in\operatorname{Cl}(T)$ with $y\in\mathfrak{T}$ and $a\notin\{5,8\}$. We use the table of
rooms over the source, \eqref{4.26:eq-room-table-a}, and distinguish two cases.

If $\ell_T=m(y)$, the table gives room at least $\beta\min\{\frac12,\mathsf{c}_e\}=\beta\hat{c}$.

If $\ell_T<m(y)$, the table gives $\theta_{T,a}\ge2r_a\ge2\theta_{\mathrm{src},a}$. Hence
$\theta_{T,a}-\theta_{\mathrm{src},a}\ge\frac12\theta_{T,a}$, and the table bounds
$\theta_{T,a}/\mathbf{r}^T_a$ below by $1-\beta$. Here $\beta$ is the schedule constant of the layered
analyticity spaces (Definition~\ref{4.4:def-post-step-space}), and $\beta=\frac15$, so that
$\frac12(1-\beta)=\frac25$ and $\beta\hat{c}=\frac1{10}$. Therefore
\begin{equation*}
\frac{\theta_{T,a}-\theta_{\mathrm{src},a}}{\mathbf{r}^T_a}\ge\frac12\,\frac{\theta_{T,a}}{\mathbf{r}^T_a}
\ge\tfrac12(1-\beta)=\tfrac25\ge\tfrac1{10}=\beta\hat{c} .
\end{equation*}

Taking the infimum over these clauses gives $\operatorname{room}_T\ge\beta\hat{c}$, which is the inequality in \eqref{4.26:eq-radii-window-c}.
\end{proof}

\begin{definition}[The proof family: block shifts, dilations, near pairs and remainders]
\label{4.26:def-proof-family}
Fix the stage $l$, a region $Y$ and the input pair $(\mathbb{U},\mathbb{J})$ (gauge field and
current) on $Y$. The radii $R^{\natural}_{\square}$, $r_T(\square)$, $r^{\mathrm{nr}}_{\square}$, the
remainder radii and the spend $\mathsf{w}_0$ are those of
Definition~\ref{4.26:def-radii-window}, displayed in \eqref{4.26:eq-radii-window-a}. The
\emph{proof family} $\mathfrak{W}^{(l),\natural}(Y)$ at $(\mathbb{U},\mathbb{J})$ on $Y$ consists of
the following four families.

\begin{enumerate}
\item[(a)] \emph{Block shifts.} For a family $t=(t_{\square})$ indexed by the cubes $\square$ of the
cover, put
\begin{equation}\label{4.26:eq-proof-family-1}
\chi_t:=\sum_{\square}t_{\square}\,\zeta_{\square},
\end{equation}
where the parameters satisfy $|t_{\square}|\le R^{\natural}_{\square}$ for every $\square$ with
$\tilde{\square}\subset Y$, and $t_{\square}=0$ otherwise. The block shift of the pair is
\begin{equation}\label{4.26:eq-proof-family-2}
\Phi_{t,\sigma}(b):=\Bigl(e^{i\eta\chi_t\mathbb{H}(\sigma;b)}\,\mathbb{U},\;
\mathbb{J}+\chi_t\,\mathcal{J}^{\sharp}[1](\sigma;b)\Bigr),
\end{equation}
where $\eta$ is the scale factor of Ba{\l}aban's construction that multiplies the layer field in
the exponent, here and in \cite[(3.30), (3.48)]{B-CMP109}. The second entry of
$\Phi_{t,\sigma}(b)$ is local: $\mathcal{J}^{\sharp}[1](\sigma;b)$ is the local current slot of
$\mathbb{H}(\sigma;b)$, supplied by the corollary on the local current slot of the stage layer. The
family is taken together with the slot data $(A^u_j,\lambda)$, where $u=(u_q)_{q\in\mathcal{Q}^w_c}$
are the parameters of the slot datum $A^u_j$, subject to the following conditions:
$|u_q|\le r_A$ for $q\in\mathcal{Q}^w_c$; the strong part of $A_j$ is real, in
$\mathfrak{X}^{\sharp}_j$; and $\lambda$ denotes the exterior slots in the slot domain
$\mathfrak{D}_{\mathcal{S}}$.

\item[(b)] \emph{Dilations.} For a cube $\square$ of the cover and a term $T\in\mathcal{T}_1$ with
$X_T\not\subset\tilde{\square}^2$,
the family is $\mathcal{W}_T(t,\tau)$ of the dilation statement for single terms. The parameters are
$t\in[0,1]$ and $|\tau|\le\mathsf{w}_0\,r_T(\square)$, the radius condition of
\eqref{4.26:eq-radii-window-a} for such terms.

\item[(c)] \emph{Near pairs.} Let $\square$ be a cube with $\square_0:=\tilde{\square}^5\subset Y$,
and let $T\in\mathcal{T}_1$ be a term with $X_T\subset\tilde{\square}^2$, level $\ell:=\ell_T$ and
localisation domain $X_T$. The stage layer at $(t,\tau)$ is
\begin{equation}\label{4.26:eq-proof-family-3}
\mathbf{A}(t,\tau):=\bigl(t\,\zeta'_{\square}+\tau\,\zeta_{\square}\bigr)\,\mathbb{H}^w ,
\end{equation}
where $\mathbb{H}^w$ is the layer field in the form in which it enters the statement on near pairs
at the stage layer, and $\mathbf{A}(t,\tau)$ weights it by the functions $\zeta'_{\square}$ and
$\zeta_{\square}$ of the partitions of unity, with the parameters $t$ and $\tau$.
The near pairs are the pairs of \cite[(3.53)]{B-CMP109} taken at this layer on $\square_0$:
\begin{equation}\label{4.26:eq-proof-family-4}
\begin{split}
P(\varsigma;t,\tau;B'):=\Bigl(&U_{\ell}\bigl(\square_0,\exp i(\varsigma B(t,\tau)+B')\bigr),\\
&J_{\ell}\bigl(\square_0,\exp i(\varsigma B(t,\tau)+B')\bigr)\Bigr)\restriction X_T ,
\end{split}
\end{equation}
where
\begin{equation}\label{4.26:eq-proof-family-5}
B(t,\tau):=Q_{\ell}\bigl(\eta\cdot A(\mathbf{A}(t,\tau),\mathbf{H}^{\square})\bigr).
\end{equation}
Here $U_{\ell}(\square_0,\cdot)$ and $J_{\ell}(\square_0,\cdot)$ are the box minimiser on
$\square_0$ at level $\ell$ and its current, $Q_{\ell}$ is the averaging operator at level $\ell$,
and $A(\mathbf{A},\mathbf{H})$ is the potential composed from a layer $\mathbf{A}$ and a box
function $\mathbf{H}$, all as in \cite[(3.30), (3.48), (3.53)]{B-CMP109}.
As in \cite[(3.30), (3.48)]{B-CMP109}, $B(t,\tau)$ splits as
\begin{equation}\label{4.26:eq-proof-family-11}
B(t,\tau)=Q_{\ell}\bigl(\eta\cdot\mathbf{H}^{\square}\bigr)+Q'(t,\tau),
\end{equation}
which defines the remainder $Q'(t,\tau)$. The function
$\mathbf{H}^{\square}$ is the box function of the box configuration
\begin{equation}\label{4.26:eq-proof-family-6}
\mathbb{U}^{\mathfrak{s}}_{\square}:=U_{p_{\square}}\bigl(\square_0,M^{\mathfrak{s}}(\mathbb{U})\bigr),
\end{equation}
where $p_{\square}$ is the top level of the structure $\mathfrak{s}$ on $\square_0$. The structure
$\mathfrak{s}$ is chosen as follows.
\begin{itemize}
\item In general, $\mathfrak{s}:=\mathbb{B}^{(l)}_k\restriction\square_0$ is the structure of the
input. The unit set $\mathfrak{B}=\mathbb{B}^{(l-1)}_k$ is one stage finer.
\item On a live older piece of $\mathfrak{T}$,
$\mathfrak{s}:=\mathfrak{B}^{\mathfrak{c}}\restriction\square_0$, the graded structure of
\eqref{4.26:eq-graded-cover-a}. It is a predecessor structure of the piece's own run, in the
sense of the structures on a live piece displayed there.
\end{itemize}
The configuration \eqref{4.26:eq-proof-family-6} corresponds to the configuration
$U_{k+1}(\square_0,M'(\mathbf{U}))$ of \cite{B-CMP109}. It is the inner box minimiser
$\mathbb{U}^w_{\square}$ of the statement on near pairs at the stage layer, which is retained.

The parameters are as follows: $\varsigma\in[0,1]$ is the Taylor parameter of
\cite[(3.34), (3.53)]{B-CMP109}; $\tau$ satisfies
$|\tau|\le\mathsf{w}_0\,r^{\mathrm{nr}}_{\square}$, the radius condition of
\eqref{4.26:eq-radii-window-a} for such terms;
and $|B'|<\alpha_3(T)$, where $\alpha_3(T)$ is the radius $\alpha_3$ among Ba{\l}aban's radii
$\alpha$ of the auxiliary parameters, taken for the term $T$.

\item[(c$'$)] \emph{Remainders.} These are the two remainder configurations of the near pair, which
accompany the near pairs as the remainders of the expansion in \cite{B-CMP109}. For a parameter
$t_{\square}$ let $\mathsf{D}_{t_{\square}}$ be the far datum of \cite{B-CMP109},
\begin{equation}\label{4.26:eq-proof-family-7}
\mathsf{D}_{t_{\square}}:=(1-\zeta'_{\square})\cdot
Q\bigl(\xi\,\mathbf{H}_{\ell}((t+t_{\square}\zeta_{\square})\mathbb{H}^w)\bigr),
\end{equation}
where $\xi$ is the scale factor of Ba{\l}aban's construction that multiplies the box field in the
exponents of the remainder configurations below, $Q$ is the averaging operator, and
$\mathbf{H}_{\ell}$ is the box field of \cite{B-CMP109} at level $\ell$, taken globally in
\eqref{4.26:eq-proof-family-7} and on $\square_0$ with a datum in \eqref{4.26:eq-proof-family-8}.
It is supported at $\ell$-distance $d_{\ell}\ge D$ from $X_T$, where $D$ is the distance
threshold of the far datum. Put
$\mathsf{D}_0:=\mathsf{D}_{t_{\square}}$ at $t_{\square}=0$, and
\begin{equation}\label{4.26:eq-proof-family-8}
\delta\mathbf{H}_{\ell}:=\partial_{t_{\square}}\big|_{t_{\square}=0}
\mathbf{H}_{\ell}\bigl(\square_0,\mathsf{D}_{t_{\square}}\bigr).
\end{equation}
The two remainder configurations are
\begin{equation}\label{4.26:eq-proof-family-9}
e^{i\xi[\mathbf{H}_{\ell}(\square_0,\mathsf{D}_0)+\tau\,\delta\mathbf{H}_{\ell}]}
\cdot P(\varsigma;t,0;B')
\end{equation}
and
\begin{equation}\label{4.26:eq-proof-family-10}
e^{i\xi\tau_1\mathbf{H}_{\ell}(\square_0,\mathsf{D}_0)}\cdot P(\varsigma;t,\tau_2;B');
\end{equation}
the term $T$ is evaluated at each of them.
The parameters are subject to the remainder radii of \eqref{4.26:eq-radii-window-a} and to
$|\tau_2|\le\mathsf{w}_0\,r^{\mathrm{nr}}_{\square}$; the parameters $\varsigma$ and $B'$ are as
in (c).
\end{enumerate}
\end{definition}

In what follows $l$ is the stage, $m(y)$ is the level at a location $y$, $\rho^{(l)}_m$ is the one-step room at level $m$, and $\mathcal{C}^{(l)}=\mathcal{C}_j$ is the set of level-$j$ bonds of the collar of the stage, the level $j$ being the one fixed in Definition~\ref{4.26:def-regularity-radius}. For $y\in Y$ we put
\begin{equation}\label{4.26:eq-family-consumption-a}
\iota_l(y):=\tfrac16\,\mathsf{w}_0\,\rho^{(l)}_{m(y)}\,e^{-\frac14\delta d^{\wedge}(y,\mathcal{C}^{(l)})},
\qquad
\iota^1_l(y):=\tfrac23\,\mathsf{w}_0\,\rho^{(l)}_{m(y)} .
\end{equation}
Since $\frac16=\frac14\cdot\frac23$ and the exponential factor is at most $1$, we have $\iota_l\le\frac14\iota^1_l$. Since $\frac23\mathsf{w}_0\le1$, we have $\iota^1_l(y)\le\rho^{(l)}_{m(y)}$.

\begin{lemma}[Room consumed by the proof family, point by point]\label{4.26:lem-family-consumption}
Let $1\le l\le N$ and let $Y$ be a union of cubes. Let $(\mathbb{U},\mathbb{J})\in\mathfrak{N}(Y)\subset\mathfrak{K}_P$, where $\mathfrak{N}(Y)$ is the fat regularity space of \eqref{4.26:eq-regularity-space-a}. Assume the floors and the $\gamma$-ceilings of the parameter choice, and the window condition \eqref{4.26:eq-radii-window-b}. For $y\in Y$ write $m:=m(y)$ and $\rho_m:=\rho^{(l)}_m$. For the proof family $\mathfrak{W}^{(l),\natural}(Y)$ of Definition~\ref{4.26:def-proof-family} the following hold.
\begin{enumerate}
\item[(i)] \emph{All blocks at once, at every location $y\in Y$.} We have $N_y(\chi_t\mathbb{H})\le\varkappa_{\mathsf{w}}(y)$, where
\begin{equation*}
\varkappa_{\mathsf{w}}(y):=22B_\sharp\,\mathsf{c}_{\square}\,\mathsf{w}_0R_0C_1\delta'_j\,
e^{-\frac12\delta d^{\wedge}(y,\mathcal{C}_j)}\in\mathfrak{M}^{\wedge}_{\delta/2}(\mathfrak{B}).
\end{equation*}
Put
\begin{equation*}
\Delta_a:=Q_a(\Phi_{t,\sigma}(b))(y)-Q_a(\mathbb{U},\mathbb{J})(y).
\end{equation*}
Then $\Delta_5=\Delta_8=0$, and $\Delta_a\le\iota_l(y)r_a(y)$ for $a\in\{1,3,6,7\}$, in the shell units.

Now let $a\in\{2,4\}$. Consider every frame whose unit set is $\mathfrak{B}$: the $(\Sigma^{\flat})$-structures, a term $T$'s own structures, and every $(\Sigma^{\ast})$-triple $(p,X',\mathfrak{s})$ with $X'\subset Y$. In each such frame, at the points $x\in X'^{\flat}$ where composite entries are sited,
\begin{equation*}
Q_a(\Phi)\le(1+\pi_\times)Q_a+\pi\,r_a(m,l_x),\qquad \pi_\times\le\pi\le\tfrac13\iota_l(y).
\end{equation*}
Hence $\Delta_a\le\iota_l(y)\mathbf{r}^T_a(y)$ wherever $Q_a\le2\mathbf{r}^T_a$ (the unit of $T$). Likewise $\Delta_a\le\iota_l(y)r_a(m,l_x)$ wherever $Q_a\le1.4\,r_a(m,l_x)$ (the shell unit; here $2.4\pi\le\iota_l$).

The map $(\mathbb{U},\mathbb{J},b,\sigma,t)\mapsto\Phi$ is holomorphic. The bounds at $y$ read $(\mathbb{U},\mathbb{J})$ only on the component of $Y(\sigma)$ containing $y$.

\item[(ii)] \emph{On $Z^{\mathrm{on}}$, location by location, excluded components included.} Suppose moreover that $(\mathbb{U},\mathbb{J})$ obeys the lettered clauses $\mathfrak{L}_l$ at the locations of a set $A\subset Y\cap Z^{\mathrm{on}}$. Here $\mathfrak{L}_l$ means Ba{\l}aban's clauses at stage $l$; on an excluded component $W$ with $l>n^+$ it means the clauses of the arrest-stage set $\mathfrak{S}^{(k),\mathrm{arrest}}_l$, that is, of the set held at stage $l$ of the application of $\mathbf{R}$ at scale $k$ when that application is arrested. Then $\Phi_{t,\sigma}(b)$ obeys $\mathfrak{L}_{l-1}$ at the locations of $A$, using at most $\frac16\mathsf{w}_0$ of each one-step room.

\item[(iii)] \emph{The second slot is free.} $\Phi_{t,\sigma}(b)$ depends neither on $u$, nor on $A_j$, nor on $\lambda$. Thus $(t,\sigma,b)$ and $(u,\lambda)$ range over a product. For real $A_j\in\operatorname{supp}\chi^{(j)\sharp}$ and $|u_q|\le r_A$, the level-$j$ slot datum $A_j^u$ lies in $\mathfrak{D}_{\mathcal{S}}$. The $u$-family moves no $Q_a$, so it spends none of the room of $(\mathbb{U},\mathbb{J})$.

\item[(iv)] \emph{Pieces.} Let $(c,S,\mathcal{S})$ be a stored term, and let $\mathfrak{U}$ be a set on which $c$ is holomorphic. Put
\begin{equation*}
\mathfrak{N}_{\mathfrak{U}}(c):=\sup_{\mathfrak{U}\times\mathfrak{D}_{\mathcal{S}}}|c|,
\end{equation*}
and assume $\Phi_{t,\sigma}(b)\restriction S\in\mathfrak{U}$ for the whole family of clause (a) of Definition~\ref{4.26:def-proof-family}. Then for $A\subset\mathcal{A}_c$ and $Q\subset\mathcal{Q}^w_c$,
\begin{equation*}
|c_{A,Q}|\le\mathfrak{N}_{\mathfrak{U}}(c)\prod_{\square\in A}K_A\vartheta^{\natural}(\square)
\prod_{q\in Q}\frac{2}{r_A},\qquad |c_{\mathrm{str}}|\le2\,\mathfrak{N}_{\mathfrak{U}}(c).
\end{equation*}

\item[(v)] \emph{Terms of $\mathcal{T}_1$.} Let $T\in\mathcal{T}_1$.

First let $X_T\not\subset\tilde{\square}^2$ and $|\tau|\le\mathsf{w}_0r_T(\square)$, with $r_T(\square)$ as in \eqref{4.26:eq-radii-window-a}. At $y\in X_T$, each entry of $\mathcal{W}_T(t,\tau)$ differs from the input's by at most $\iota^1_l(y)\mathbf{r}^T_a(y)$, wherever the input's entries $2$ and $4$ are at most $(1+\beta)\mathbf{r}^T_a$; entries $5$ and $8$ differ by $0$. If moreover $\mathcal{W}_T(t,\tau)\restriction X_T\in\mathfrak{U}_T$ on the disc, then
\begin{equation*}
|\operatorname{piece}(T,\square)|\le\|T\|_\infty K''\vartheta^{\natural}
e^{-\frac12\delta d^-_{\square}}e^{-\frac14\delta d_\ell(X_T,\square^{\zeta})},
\qquad K'':=28K^+ .
\end{equation*}

Next let $X_T\subset\tilde{\square}^2$ and $|\tau|\le\mathsf{w}_0r^{\mathrm{nr}}_\square$, with $r^{\mathrm{nr}}_\square$ as in \eqref{4.26:eq-radii-window-a}. The stage layer moves each entry of the box pair $P$ by at most $\iota^1_l(y)\mathbf{r}^T_a(y)$. In units of $\mathbf{r}^T_a(y)$, the dilated member moves by at most $\frac13\mathsf{w}_0\rho_m$ and the undilated member by at most $\mathsf{w}_0\rho_m/18$. The remainder configurations of clause (c$'$) of Definition~\ref{4.26:def-proof-family}, at the radii of \eqref{4.26:eq-radii-window-a}, are moved by at most $0.617\,\mathsf{w}_0\rho_m\mathbf{r}^T_a$ and $0.605\,\mathsf{w}_0\rho_m\mathbf{r}^T_a$ respectively. Both are smaller than $\iota^1_l(y)\mathbf{r}^T_a$.
\end{enumerate}
\end{lemma}

\begin{proof}
Part (i) is proved at amplitude $\mathsf{w}_0R_0$, with $\theta=\frac12$, on $\mathfrak{K}_P$, in five steps.

\emph{Step 1.} Apply clause (v) of the recursive stage-layer estimates, recorded in Definition~\ref{4.26:def-regularity-radius}, to $\chi=\chi_t=\sum_\square t_\square\zeta_\square$, that is, with $c_\square=t_\square$. With $a(y)$ as defined there, this gives
\begin{equation*}
q:=10\sum_\square|t_\square|\,q_\square\in\mathfrak{M}_{2\delta},\qquad a\cdot p\le q .
\end{equation*}
We bound $q$ at every point, taking all blocks at once. Clause (vi) of the recursive stage-layer estimates is homogeneous in the amplitude; we apply it to one term $t_\square\zeta_\square$ at a time and sum over the blocks.

Fix $\square$ and $y'\in\square^{\zeta}$. By Definition~\ref{4.26:def-regularity-radius}, $d^-_\square\le d^{\wedge}(y',\mathcal{C}_j)$. By \eqref{4.26:eq-radii-window-a}, $|t_\square|\le R^\natural_\square=\mathsf{w}_0R_0e^{\frac12\delta d^-_\square}$. By definition, $p(y')=\sup|b|\,e^{-\delta d^{\wedge}(y',\mathcal{C}_j)}$. Together these give
\begin{equation*}
|t_\square|\,p(y')\le\mathsf{w}_0R_0\sup|b|\,e^{-\frac12\delta d^{\wedge}(y',\mathcal{C}_j)} .
\end{equation*}
By property (g1) of $d^{\wedge}$ and by $d^{\wedge}\le d$ (both in Definition~\ref{4.26:def-regularity-radius}),
\begin{equation*}
e^{-\frac12\delta d^{\wedge}(y',\mathcal{C}_j)}\le e^{-\frac12\delta d^{\wedge}(y,\mathcal{C}_j)}\,e^{\frac12\delta d(y,y')} .
\end{equation*}
Multiply by $e^{-2\delta d(y,y')}$ and take the supremum over $y'\in\square^{\zeta}$, which is the definition of $q_\square(y)$. Since $\square^{\zeta}\subset\square\subset\tilde{\square}$, we have $d(y,y')\ge d(y,\tilde{\square})$. Therefore
\begin{equation*}
|t_\square|\,q_\square(y)\le\mathsf{w}_0R_0\sup|b|\,e^{-\frac12\delta d^{\wedge}(y,\mathcal{C}_j)}\,
e^{-\frac32\delta d(y,\tilde{\square})} .
\end{equation*}
Now sum over $\square$. We use
\begin{equation*}
\sum_\square e^{-\frac32\delta d(y,\tilde{\square})}\le\sum_\square e^{-\delta d(y,\tilde{\square})}\le\mathsf{c}_\square
\end{equation*}
and $\sup|b|\le1.1\,C_1\delta'_j$. The result holds at every $y$:
\begin{equation*}
q(y)\le\bar{q}(y):=11\,\mathsf{c}_\square\mathsf{w}_0R_0C_1\delta'_j\,e^{-\frac12\delta d^{\wedge}(y,\mathcal{C}_j)} .
\end{equation*}
No cancellation among the $\zeta_\square$ is used, and the phases of the $t_\square$ are arbitrary. The bound holds because the amplitudes grow at rate $\frac12\delta$ while the layer decays at rate $\delta$. The tails of far blocks are summed by $\mathsf{c}_\square$, which is a factor of $\mathsf{K}_r$.

\emph{Step 2.} The layer bound of Definition~\ref{4.26:def-regularity-radius} states $N^\Delta_y(\mathbb{H})\le B_\sharp p(y)$ on $|\sigma|\le e^{\kappa_1}$. With it and Step~1,
\begin{equation*}
N_y(\chi_t\mathbb{H})\le2a(y)N_y(\mathbb{H})\le2B_\sharp\,a(y)p(y)\le2B_\sharp\,\bar{q}(y)=\varkappa_{\mathsf{w}}(y).
\end{equation*}
By (g1), $\varkappa_{\mathsf{w}}\in\mathfrak{M}^{\wedge}_{\delta/2}(\mathfrak{B})$.

The absorption inequality of Definition~\ref{4.26:def-regularity-radius} reads
\begin{equation*}
C_1\delta'_j e^{-\frac14\delta d^{\wedge}(y,\mathcal{C}_j)}\le\vartheta\rho_{m(y)}\alpha_{\cdot,m(y)} .
\end{equation*}
The definition of $R_0$ gives $R_0\vartheta\le1/(2\mathsf{K}_r)$. From these two facts,
\begin{equation}\label{4.26:eq-family-consumption-b}
\varkappa_{\mathsf{w}}(y)\le\frac{11B_\sharp\mathsf{c}_\square}{\mathsf{K}_r}\,\mathsf{w}_0\,\rho_m\,\alpha_{\cdot,m}\,
e^{-\frac14\delta d^{\wedge}(y,\mathcal{C}_j)} .
\end{equation}
The decay rate $\delta$ of the layer is split as follows: $\frac12\delta$ pays for the amplitude, $\frac14\delta$ is used in the absorption inequality, and $\frac14\delta$ is kept.

\emph{Step 3.} We treat the entries one at a time.

Entries $5$ and $8$: the decomposition keeps $U$ and the cube gauges for every shift, so $\Delta_5=\Delta_8=0$.

Entries $1$, $6$, $7$ read first-order sizes only; the multiplicative pieces are absorbed once $600\,w_\ast\sup\alpha\le1$, which holds under the assumed $\gamma$-ceilings, since
\begin{equation*}
600\,w_\ast\sup\alpha=120\,\gamma_0^{(k)}\le1 .
\end{equation*}
By \eqref{4.26:eq-family-consumption-b} they move by at most
\begin{equation*}
3\varkappa_{\mathsf{w}}/\alpha_{\cdot,m}\le\tfrac{33}{2^9}\,\mathsf{w}_0\rho_m\,e^{-\frac14\delta d^{\wedge}},
\end{equation*}
in the units $r_a(y)$. Here we used $\mathsf{K}_r\ge2^9\mathsf{c}_\square B_\sharp$. This holds by the definition of $\mathsf{K}_r$ in Definition~\ref{4.26:def-regularity-radius}, because $\mathsf{c}_\vartheta\ge1$ and $K_S^b\ge1$.

Entries $2$ and $4$: apply the lemma on exponentiated one-forms in a frame, whose statement is recorded in Definition~\ref{4.26:def-regularity-radius}, with $\mathcal{K}:=\chi_t\mathbb{H}$ and $\varkappa:=\varkappa_{\mathsf{w}}$. Its hypotheses hold as follows.
\begin{itemize}
\item (D0) holds because $\mathbb{U}\in\mathfrak{K}_P$, with the ambient clause imposed off $Y$.
\item (D1) holds on the unit set $\mathfrak{B}$ itself.
\item (D2) is $\sup\varkappa_{\mathsf{w}}\le\alpha_{\cdot,j}$, which holds by \eqref{4.26:eq-family-consumption-b} on $\mathcal{C}_j$.
\item Hypothesis $(\mathrm{S}^b)$ is a smallness condition on $\gamma_0$ times a polylogarithmic factor; it holds under the $\gamma$-ceilings assumed in the lemma.
\end{itemize}
The lemma then gives
\begin{equation*}
\pi=\frac{K_S^b\varkappa_{\mathsf{w}}(x)}{\tilde{\alpha}_{0,m}}\le\tfrac{22}{2^9}\,\mathsf{w}_0\rho_m\,e^{-\frac14\delta d^{\wedge}} .
\end{equation*}
A factor $2$ is spent in comparing $\alpha_{\cdot,m}$ with $\tilde{\alpha}_{0,m}$, and we used $\mathsf{K}_r\ge2^9\mathsf{c}_\square K_S^bB_\sharp$.

Entry $3$: the second slot is local. Apply clause (iv) of the recursive stage-layer estimates at $\chi\equiv1$. Its hypothesis $2B_\sharp p\le1/10$ holds under the assumed $\gamma$-ceilings, since $p\le\sup|b|\le1.1\,C_1\delta'_j$. Then use $|\chi_t|\le a$, the bound $a\cdot p\le q\le\bar{q}$ of Step~1, and the absorption inequality, to get
\begin{equation*}
\ell^3|\chi_t\mathcal{J}^{\sharp}[1]|\le a(y)\cdot2K_\chi p(y)\le2K_\chi\bar{q}(y)
\le\tfrac{11}{2^9}\,\mathsf{w}_0\rho_m\alpha_{0,m}\,e^{-\frac14\delta d^{\wedge}} .
\end{equation*}
The last inequality uses the absorption inequality, $R_0\vartheta\le1/(2\mathsf{K}_r)$ and $\mathsf{K}_r\ge2^9\mathsf{c}_\square K_\chi$.

The numerical comparisons needed are
\begin{equation}\label{4.26:eq-family-consumption-c}
\begin{gathered}
\tfrac{33}{2^9}<\tfrac16,\qquad \tfrac{22}{2^9}<\tfrac1{18},\qquad \tfrac{11}{2^9}<\tfrac16,\\
3\cdot\tfrac{22}{2^9}=0.129<\tfrac16,\qquad 2.4\cdot\tfrac{22}{2^9}=0.103<\tfrac16 .
\end{gathered}
\end{equation}
Since $\mathcal{C}^{(l)}=\mathcal{C}_j$, the first and third comparisons give $\Delta_a\le\iota_l(y)r_a(y)$ for $a\in\{1,6,7\}$ and for $a=3$. The second gives $\pi\le\frac13\iota_l(y)$.

\emph{Step 4.} We derive the last clauses of (i). In a frame with top level $p$, apply the unit rule \eqref{4.26:eq-clauses-units-a} at $p$, which bounds $r_a(m,l_x)$ by $\mathbf{r}^T_a$, together with $\pi_\times\le\pi$. Where $Q_a\le2\mathbf{r}^T_a$,
\begin{equation*}
\Delta_a\le\pi_\times Q_a+\pi\,r_a(m,l_x)\le\pi(2\mathbf{r}^T_a+\mathbf{r}^T_a)=3\pi\,\mathbf{r}^T_a .
\end{equation*}
Where $Q_a\le1.4\,r_a(m,l_x)$, the same argument gives $\Delta_a\le\pi(1.4+1)r_a(m,l_x)$. By Step~3 and the fourth and fifth comparisons in \eqref{4.26:eq-family-consumption-c}, both bounds are at most $\iota_l(y)$ times the respective unit.

\emph{Step 5.} Holomorphy of $\Phi$ follows from clause (ii) of the recursive stage-layer estimates, since $\chi_t$ is linear in $t$. Locality follows from the locality clause of the shell construction, together with the pointwise form of the lemma on exponentiated one-forms. Both slots of $\Phi$ read $\chi_t$ only at the point. Hence, for every stored term $(c,S,\mathcal{S})$ as in (iv), $\Phi\restriction S$ reads $t_\square$ only for $\square\in\mathcal{A}_c$, and the bounds at $y$ read $(\mathbb{U},\mathbb{J})$ only on the component of $Y(\sigma)$ containing $y$. This proves (i).

\emph{Part (ii).} Membership in the lettered clauses is checked separately at each structure and each location. At each such pair, the clause of $\mathfrak{L}_l$ (the source) and the clause of $\mathfrak{L}_{l-1}$ (the target) read the same functional.

Consider a $(\Sigma^{\flat})$-structure with birth level $m_{\mathrm{ref}}$, and put $\Delta m:=m-m_{\mathrm{ref}}\ge0$. The room there is $\rho_{m_{\mathrm{ref}}}\mathbf{r}$, where $\mathbf{r}$ denotes the radius of the entry in question. We use the comparison of shell radii and rooms across levels:
\begin{equation*}
r_e(m,\cdot)\le\bigl(2.2/L^2\bigr)^{\Delta m}r_e(m_{\mathrm{ref}},\cdot),\qquad
\rho_m\le2^{\Delta m}\rho_{m_{\mathrm{ref}}},\qquad \pi_\times\le\pi L^{-2\Delta m}.
\end{equation*}
With these and part (i), the entry moves by at most $3\cdot\frac{22}{2^9}\mathsf{w}_0\rho_{m_{\mathrm{ref}}}\mathbf{r}$ when $\Delta m=0$. When $\Delta m\ge1$ it moves by at most
\begin{equation*}
2\cdot\tfrac{22}{2^9}\,\mathsf{w}_0\,(4.4/L^2)^{\Delta m}\rho_{m_{\mathrm{ref}}}\mathbf{r}.
\end{equation*}
The first bound is at most $\frac16\mathsf{w}_0\rho_{m_{\mathrm{ref}}}\mathbf{r}$ by \eqref{4.26:eq-family-consumption-c}. The second is at most $\frac16\mathsf{w}_0\rho_{m_{\mathrm{ref}}}\mathbf{r}$ for every $\Delta m\ge1$, because $2\cdot22/2^9<\frac16$ and $(4.4/L^2)^{\Delta m}\le1$.

For the direct entries, the gap is $\rho_m$ at locations of class S and the ratio room at raised locations. Each gap is at least the increment of part (i).

In this way the argument of part (i) is used on a sub-polydisc, location by location, with input in the lettered space. This is legitimate for two reasons. First, each step of the proof of (i) is a pointwise increment on $\mathfrak{K}_P$, together with a comparison with the input's own entry at the same structure and location. Second, nothing is required of the input off $A$ beyond hypothesis (D0) of the lemma on exponentiated one-forms on the stencils, that is, fat regularity.

Now let $W$ be an excluded component with $l>n^+$. There the source clauses $\mathfrak{L}_l$ and the target clauses $\mathfrak{L}_{l-1}$ are those of the sibling arrest-stage sets $\mathfrak{S}^{(k),\mathrm{arrest}}_l$ and $\mathfrak{S}^{(k),\mathrm{arrest}}_{l-1}$ on $W$. These read one frozen structure at every $x\in W$, by the definition of the arrest-stage sets, as in the construction of \cite{B-CMP122a}. Consequently:
\begin{itemize}
\item every location of $W$ is of class S;
\item every entry is compared with the entry of the same kind;
\item for entries $2$ and $4$, the lemma on exponentiated one-forms is applied in the frame given by the frozen cube family of $W$.
\end{itemize}
The one-step room at a location of level $\ell$ is the difference of Ba{\l}aban's subtractions at stages $l$ and $l-1$. Here $h$ is the offset of these subtraction sums and $\ell_l=h+l-1$ is the level entering $\rho^{(l)}$, so the difference is the single term $i=h+l$ of the first sum:
\begin{equation*}
\beta\sum_{i=h+1}^{h+l}2^{-|\ell-i|}-\beta\sum_{i=h+1}^{h+l-1}2^{-|\ell-i|}
=\beta\,2^{-|\ell-\ell_l-1|}=\rho^{(l)}_\ell .
\end{equation*}
The $\sigma$-terms of older enclosed pieces do not depend on $l$ and cancel in this difference.

Part (i) holds at any determining set. It reads $\mathfrak{B}$ only through the layer bound, the absorption inequality and properties (g1)--(g2), and all of these are stated in Definition~\ref{4.26:def-regularity-radius} for an arbitrary admissible $\mathfrak{B}$. Hence the increment is at most $\iota_l\le\frac16\mathsf{w}_0\rho^{(l)}_\ell$. This proves (ii).

\emph{Part (iii).} Consider \cite[(1.60), pp.~188--189]{B-CMP122a}. The first bracket of the integrand there is
\begin{equation*}
(\mathbb{E}_k+\mathbb{R}_k+\mathbb{B}_k+\mathbb{B}^{(n)}_k)
\bigl(U^{(n)}_k(\exp i[g_jCB_j-h\tilde{D}(g_jCB_j)]V^{(j)})\bigr),
\end{equation*}
and it is subtracted at $(U^{(n+1)}_k,\,A_j=0)$. Thus the background is a function of $B_j$ alone.

The variable $A_j$ enters in four places: through the slot of the terms (a term depends on $U_k$ and on $A$ restricted to $X$, \cite{B-CMP119}), through $\mathbb{P}^{(j)}(g_j,A_j,B_j)$, through $\chi^{(j)}$, and through its own Gaussian measure. So $b=B'_j$, $\mathbb{H}(\sigma;b)$, $\mathcal{J}^{\sharp}[1](\sigma;b)$ and $\Phi$ are free of $A_j$. The backgrounds are free of $\lambda$: the configuration enters the terms through the background field $U_k(V)$ \cite{B-CMP119}, which does not contain the exterior slots $\lambda$; these are separate arguments of the terms, ranging over $\mathfrak{D}_{\mathcal{S}}$.

The domain of $c$ is a product. For leaves this holds by the definition of the leaves' domain; for chain terms it holds by the corollary on chain terms.

On $\operatorname{supp}\chi^{(j)\sharp}$ one has $T_e\le\mathfrak{b}_j$ at every bond $e$, by the first of the radius clauses that fix $\mathfrak{b}_j$; here $T_e$ is the quantity of that clause which bounds $|A_j(e)|$ at the bond $e$. Hence, at every bond $e$ of a weak cube, $|u_qA_j(e)|<r_A\mathfrak{b}_j\le\varrho_j$, which is the weak-cube condition of $\mathfrak{D}_{\mathcal{S}}$. The strong part of $A_j$ is real and lies in $\mathfrak{X}^{\sharp}_j\supset\operatorname{supp}\chi^{(j)\sharp}$. Therefore $A_j^u\in\mathfrak{D}_{\mathcal{S}}$.

Since $\Phi$ does not depend on $u$, the $u$-family moves no $Q_a$. This proves (iii).

\emph{Part (iv).} Put
\begin{equation*}
F(t,u):=c\bigl(S;\Phi_{t,\sigma}(b)\restriction S,A_j^u\bigr).
\end{equation*}
This function is holomorphic on
\begin{equation*}
\prod_\square\{|t_\square|\le R^\natural_\square\}\times\prod_{q\in\mathcal{Q}^w_c}\{|u_q|\le r_A\},
\end{equation*}
and there $|F|\le\mathfrak{N}_{\mathfrak{U}}(c)$. Four facts give this: the hypothesis $\Phi\restriction S\in\mathfrak{U}$; part (iii), which puts $A_j^u$ in $\mathfrak{D}_{\mathcal{S}}$; the definition of the domain of $c$; and the holomorphy of $\Phi$ in $t$ (part (i)) together with the linearity of $A_j^u$ in $u$.

We have
\begin{equation*}
c_{A,Q}=\int_{[0,1]^{A\sqcup Q}}\partial_{t_A}\partial_{u_Q}F(t\mathbf{1}_A,u\mathbf{1}_Q)\,dt\,du .
\end{equation*}
For $t_\square,u_q\in[0,1]$, apply Cauchy's formula on the circles
\begin{equation*}
|\tau_\square-t_\square|=R^\natural_\square-1,\qquad |\upsilon_q-u_q|=r_A-1 .
\end{equation*}
These circles lie in the polydisc. Moreover $R^\natural_\square-1\ge\frac12R^\natural_\square$ by \eqref{4.26:eq-radii-window-b}, and $r_A-1\ge\frac12r_A$ because $r_A\ge2$. Cauchy's formula therefore bounds the integrand by
\begin{equation*}
\mathfrak{N}_{\mathfrak{U}}(c)\prod_{\square\in A}\bigl(2/R^\natural_\square\bigr)\prod_{q\in Q}\bigl(2/r_A\bigr),
\end{equation*}
and $2/R^\natural_\square=K_A\vartheta^{\natural}(\square)$ by \eqref{4.26:eq-radii-window-a}. Finally, $c_{\mathrm{str}}$ is the difference of two values of $F$ at $t=0$, so $|c_{\mathrm{str}}|\le2\mathfrak{N}_{\mathfrak{U}}(c)$.

\emph{Part (v).} Let $X_T\not\subset\tilde{\square}^2$. The entry-move bound for the interpolation $\mathcal{W}_T(t,\tau)$ of Definition~\ref{4.26:def-proof-family} holds on $\mathfrak{K}_P$. By it, every entry moves by at most
\begin{equation*}
K^+\bigl[5.5\,\vartheta+4.4\,|\tau|\,\vartheta(\square)e^{-\frac14\delta d_\ell(y,\square^{\zeta})}\bigr]\rho_m\mathbf{r}^T_a .
\end{equation*}
Take $|\tau|\le\mathsf{w}_0r_T(\square)$ and $y\in X_T$. Then $d_\ell(y,\square^{\zeta})\ge d_\ell(X_T,\square^{\zeta})$. With the definition of $r_T(\square)$ in \eqref{4.26:eq-radii-window-a}, this bounds the second member by $4.4\,\mathsf{w}_0/14$. The first member is at most $5.5\,\mathsf{w}_0/28$ by \eqref{4.26:eq-radii-window-b}.

The entry-move bound uses $\pi(E+r)\le2\pi\mathbf{r}^T_a$, where $E$ is the input's entry and $r$ the shell radius, bounded by $\mathbf{r}^T_a$; it therefore presumes $E\le\mathbf{r}^T_a$. The entry of a newborn term may reach $(1+\beta)\mathbf{r}^T_a$; then the factor is $2.2\,\pi$ in place of $2\pi$, and the bound is multiplied by $1.1$. Since
\begin{equation}\label{4.26:eq-family-consumption-d}
(0.315+0.197)\cdot1.1=0.563<\tfrac23,
\end{equation}
each entry moves by at most $\iota^1_l(y)\mathbf{r}^T_a(y)$.

The piece is bounded by Cauchy's estimate at radius $\mathsf{w}_0r_T(\square)$, under the Cauchy-estimate convention for the pieces of small terms. By the definition of $r_T(\square)$ this gives the factor
\begin{equation*}
\frac{2}{\mathsf{w}_0r_T(\square)}=\frac{K''\vartheta(\square)\,e^{-\frac14\delta d_\ell(X_T,\square^{\zeta})}}{\mathsf{w}_0},
\qquad K''=28K^+ .
\end{equation*}
Since $\vartheta(\square)=\vartheta e^{-\frac12\delta d^-_\square}$ and $\vartheta^{\natural}=\vartheta/\mathsf{w}_0$, we have $\vartheta(\square)/\mathsf{w}_0=\vartheta^{\natural}e^{-\frac12\delta d^-_\square}$. This gives the stated bound on $|\operatorname{piece}(T,\square)|$.

Now let $X_T\subset\tilde{\square}^2$ (near pairs). By the formula for $K^{\mathrm{nr}}$ in \eqref{4.26:eq-radii-window-a}, for $|\tau|\le\mathsf{w}_0r^{\mathrm{nr}}_\square$,
\begin{equation*}
N_y(\tau\zeta_\square\mathbb{H}^w)\le\mathsf{w}_0\frac{11B_\sharp\bar{\rho}_\square\alpha_{\cdot,m}}{K^{\mathrm{nr}}}
=\frac{\mathsf{w}_0\bar{\rho}_\square\alpha_{\cdot,m}}{6K_{\mathrm{cv}}}
\le\frac{\mathsf{w}_0\rho_m\alpha_{\cdot,m}}{3K_{\mathrm{cv}}} .
\end{equation*}
By the definition of $K_{\mathrm{cv}}$, the dilated member therefore moves by at most $\frac13\mathsf{w}_0\rho_m\mathbf{r}^T_a$.

For the undilated member $t\zeta'_\square\mathbb{H}^w$ we use $t\le1$, the size bound $a'\le5$ for its cut-off function $\zeta'_\square$, and the absorption inequality without the block's decay. These give
\begin{equation*}
N_y\le11B_\sharp\vartheta\bar{\rho}_\square\alpha_{\cdot,m}\le22B_\sharp\vartheta\rho_m\alpha_{\cdot,m}
=\frac{K^{\mathrm{nr}}}{3K_{\mathrm{cv}}}\,\vartheta\rho_m\alpha_{\cdot,m}.
\end{equation*}
So the undilated member moves by at most $\frac13K^{\mathrm{nr}}\vartheta\rho_m\mathbf{r}^T_a$, which is at most $(\mathsf{w}_0/18)\rho_m\mathbf{r}^T_a$ by \eqref{4.26:eq-radii-window-b}. Since $\frac13+\frac1{18}<\frac23$, each entry of the box pair $P$ moves by at most $\iota^1_l(y)\mathbf{r}^T_a(y)$.

The bounds $0.617\,\mathsf{w}_0\rho_m\mathbf{r}^T_a$ and $0.605\,\mathsf{w}_0\rho_m\mathbf{r}^T_a$ for the remainder configurations are the bounds \eqref{4.26:eq-near-remainders-a} for the remainder pieces of a near pair. Both are smaller than $\frac23\mathsf{w}_0\rho_m\mathbf{r}^T_a=\iota^1_l(y)\mathbf{r}^T_a$.
\end{proof}

\begin{remark}[Remarks on the room consumed by the proof family]\label{4.26:rem-consumption-remarks}
The following three remarks concern Lemma~\ref{4.26:lem-family-consumption} (room consumed by the
proof family, point by point).
\begin{enumerate}
\item Lemma~\ref{4.26:lem-family-consumption} treats all cubes at once, not one active cube at a
time. The cover-multiplicity constant that enters its bounds is $\mathsf{c}_{\square}$, which is among
the constants collected in $\mathsf{K}_r$. For $a\in\{2,4\}$ the increments
$\Delta_a$ are controlled through the bump structure $S^{\mathrm{bump}}$, in the frames of unit set
$\mathfrak{B}$, as in the statement of Lemma~\ref{4.26:lem-family-consumption}, and never through a
smearing, in shell units, of the entries $Q_2$ and $Q_4$. The distance used throughout is
$d^{\wedge}$. The condition $\delta_1 M/M_1\ge 4$, in which $M$ and $M_1$ are among the auxiliary
parameters $\mathfrak{p}$, is not used.
\item At unit weight, $\mathsf{w}=1$ in the notation of Lemma~\ref{4.26:lem-family-consumption}, the
clause of that lemma on the sealed radius is a localised bound in which the factor
$e^{-\frac14\delta d^{\wedge}}$ is retained. With $\rho_m$, $m=m(y)$, the one-step room in the
notation of Lemma~\ref{4.26:lem-family-consumption}, the sealed radius consumes at most
\begin{equation*}
  \tfrac16\,\rho_m\, e^{-\frac14\delta d^{\wedge}},
\end{equation*}
and not $\tfrac12\rho_m$ at every cube, however far. The quantity $\tfrac12\rho_m$ at every cube
bounds the budget of the first step of the argument that spends this room. It is a
majorant of the consumption, with the factor $e^{-\frac14\delta d^{\wedge}}$ to spare, and it is
false as an equality, since the room actually consumed at a cube carries the decaying factor
$e^{-\frac14\delta d^{\wedge}}$ displayed above. What the sealed radius must give up is a constant,
independent of $d^{\wedge}$.
\item Direct entries are measured in shell units. Composite entries are measured in the unit of
the term $T$ at the structures of $T$, and in the shell unit at the fat triples. The space
$\mathfrak{N}(Y)$ contains both the direct entries and the box clause in shell
units of the definition of the fat regularity space. For each
clause $a$ there is an exponent
$e_a\in\{1,2,3\}$ such that the ratio of the clause radius $\mathbf{r}^T_a$ of the term $T$ to
the weighted shell radius $r_a$ satisfies, as part of the construction of the shell radii,
\begin{equation*}
  \mathbf{r}^T_a/r_a\asymp L^{e_a(m-\ell_T)},
\end{equation*}
where $\ell_T$ is the level of $T$, so that $m-\ell_T$ is its age. This ratio is unbounded in
the age. Nothing in the proof of Lemma~\ref{4.26:lem-family-consumption} requires it to be
bounded, and no clause of Ba{\l}aban's displayed bounds on the terms is transcribed into the
unit of $T$.
\end{enumerate}
\end{remark}

\begin{lemma}[The cover-cube count and its bound]\label{4.26:lem-cube-count}
Let $k$, $h$, the stages, and the set $Z$ be as in Definition~\ref{4.26:def-stages-regions}, and let
$\bar c$ range over the components of $Z$. Put
\begin{equation*}
\check R:=\max_{h<i\le k}R_i ,
\end{equation*}
and recall that $R_k=\min_iR_i$. Define
\begin{equation}\label{4.26:eq-cube-count-1}
\begin{gathered}
V^0:=\sup_{n,\bar c}\ \sum_{\square\in\mathcal{C}}
e^{-\frac12\delta\, d_{\mathrm{sc}}(\square,\,R_n\cap\bar c)},\qquad
V^A:=\sup_{n,\bar c}\ \#\{\text{rim }\pi_{j+1}\text{-cubes of }\bar c\},\\
V:=\max\{V^0,V^A,V^{\mathrm{fr}}\}.
\end{gathered}
\end{equation}
The suprema run over $n=0,\dots,N-1$ and over the components $\bar c$ of $Z$, the stage $l=n+1$
integrating the level $j=h+n$. In the definition of $V^0$, $\delta$ is the decay constant of the
cover sums of this chapter, and $R_n$ denotes the set attached to stage $n$ in
Definition~\ref{4.26:def-stages-regions}; this set is distinct from the integers $R_i$. The count
$V^{\mathrm{fr}}$ is the count on lifted sets in \eqref{4.26:eq-lifted-counts-a}. Put
\begin{equation*}
C_F:=12(2L^2)^d,\qquad C_V:=C_F\,C_V(d),
\end{equation*}
where $C_V(d)$, which depends on $d$ only and is distinct from $C_V$, is the constant of the
counting statement for the stage cover, read on the thresholded form $V^{0,\mathrm{th}}$ of the sum
$V^0$. Then $V\le\bar V$, and for every set $S$ and every $c$ and $w$ as in the counting statement
for the stage cover, the two bounds of that statement hold with this $V$:
\begin{equation}\label{4.26:eq-cube-count-a}
\begin{gathered}
V\le \bar V:=C_V\,(100L\check R)^d,\\
\sum_{\square:\,\tilde\square\cap S\neq\emptyset}e^{-\frac12\delta\, d^-_{\square}}
\le V\,\mathsf{p}(S),\qquad
|\mathcal{Q}^w_c|\le V\,\mathsf{p}(S).
\end{gathered}
\end{equation}
\end{lemma}

\begin{proof}
\emph{The level of a stage.} Fix a stage, with its level $j$, and fix a component $\bar c$ of $Z$.
By Definition~\ref{4.26:def-stages-regions}, the stage $l=n+1$, with $n\in\{0,\dots,N-1\}$,
integrates the level $j=h+n$, where $h=k-N$. Hence $h\le j\le k-1$, so that $h<j+1\le k$. The
index $j+1$ therefore lies in the range over which $\check R$ is the maximum, and
\begin{equation*}
R_{j+1}\le\check R .
\end{equation*}

\emph{The nesting of the regions.} Write $Z_j$ and $R_j$ for the large-field regions and the integers
of the construction in \cite{B-CMP122a}. The definition of the class of large-field regions in
\cite{B-CMP122a} contains a condition (i) on components. Its condition (ii) has a second clause,
which reads
\begin{quote}
``and the previous regions contained in it satisfy the condition (i) on the corresponding
scales''.
\end{quote}
We combine this clause with the single-nest condition imposed on $Z$ in
Definition~\ref{4.26:def-stages-regions}. The nesting lemma then gives, at every level $i$ of the
run, the following two facts.
\begin{itemize}
\item The set $Z\cap Z_i$ is the only component of $Z_i$ meeting $Z$.
\item This set is a rectangular parallelepiped contained in a cube of side $100MR_i$.
\end{itemize}
We apply this at the level $j+1$ and to the component $\bar c$. It follows that $Z_{j+1}\cap\bar c$
is one rectangular parallelepiped. At its own scale, that of level $j+1$, this parallelepiped lies
inside a cube of side $100MR_{j+1}$.

\emph{Counting the cubes inside the parallelepiped.} A cube of side $100MR_{j+1}$ at the scale of
level $j+1$ contains at most $(100R_{j+1})^d$ cubes of the partition $\pi_{j+1}$. The cubes of
$\pi_j$ are finer by the factor $L$, so the same cube contains at most $(100LR_{j+1})^d$ cubes of
$\pi_j$. Since $L\ge1$ and $R_{j+1}\le\check R$, both numbers are at most $(100L\check R)^d$.

The collar and the rim of $\bar c$ at this stage lie in the parallelepiped $Z_{j+1}\cap\bar c$. In
particular, the number of rim $\pi_{j+1}$-cubes of $\bar c$ is at most
$(100R_{j+1})^d\le(100L\check R)^d$. The stage and the component were arbitrary, so this bound
holds for $V^A$.

\emph{Cubes off the parallelepiped.} The cubes of the cover lying off the parallelepiped contribute
tails. The tail estimate for cover sums bounds the contribution of the cubes at distance $D$ from the
parallelepiped by $P(D)\,e^{-\frac12\delta D}$, where $P$ is a polynomial. We feed the count of at
most $(100L\check R)^d$ cubes inside the parallelepiped, obtained in the previous paragraph, and
these tails off it into the bound of the counting statement for the stage cover, which produces the
constants $C_F$ and $C_V(d)$. This gives
\begin{equation*}
V^0,\ V^A\le C_F\,C_V(d)\,(100L\check R)^d=C_V(100L\check R)^d .
\end{equation*}
The same statement, applied to the count on lifted sets \eqref{4.26:eq-lifted-counts-a}, bounds
$V^{\mathrm{fr}}$ by the same number. This is the first line of \eqref{4.26:eq-cube-count-a}.

\emph{The second line of \eqref{4.26:eq-cube-count-a}.} The two inequalities follow from the argument
of the counting statement for the stage cover. That argument applies without change, with the count
$V$ of \eqref{4.26:eq-cube-count-1}. Its last step uses the inequality
\begin{equation*}
e^{-\delta wR_k/8}\le (4^dC_dR_k)^{-1},
\end{equation*}
where $w$ is the weight and $C_d$ the dimensional constant of that statement, and this inequality is
secured by the choice of $\gamma$. This completes the proof of \eqref{4.26:eq-cube-count-a}.

\emph{Both inputs are used.} Neither the second clause of condition (ii) nor the single-nest
condition can be omitted from the argument above.

$(\alpha)$ Suppose only condition (i) of \cite{B-CMP122a} is used. This condition admits a
rectangular parallelepiped $Z$ of side of the order of $100MR_k s_k$, where $s_k$ is the factor with
which this side occurs in condition (i) of \cite{B-CMP122a}, whose large-field region at level $h+1$
is $Z$ less a thin layer. For such a $Z$,
\begin{equation*}
V\ge 2^{-d}L^{(d-1)(N-1)} .
\end{equation*}

$(\beta)$ Suppose conditions (i) and (ii) are used, but the single-nest condition is not. Three
features of the definitions then leave the count unbounded.
\begin{itemize}
\item Condition (i) concerns one component at a time. In the wording of \cite{B-CMP119}, it reads:
``if a component of $Z_j$ is contained in a cube of the size $100\,MR_j$ \dots, then it is a
rectangular parallelepiped''. The same condition appears in \cite{B-CMP122a}.
\item The second clause of condition (ii) bounds each previous component, not the number of
previous components.
\item The first clause of condition (ii) is a prohibition of creation; it does not forbid the
joining of components. By \cite{B-CMP122b},
\begin{quote}
``$Z$ is obtained from some number of components of $Z_j$, and some number of new large field
regions, joined together into the one component of $Z_{j+1}$ \dots\ by adding layers of
$MR_{j+1}$-cubes''.
\end{quote}
\end{itemize}
The cluster geometry of the statement on joined large-field clusters, formed by such joinings, gives
$V$ at least of the order of $(1.6L)^{d(N-1)}$. The number of stages satisfies
$N=\check R_k\ge c\,\mathsf{t}^r$, where $\check R_k$ is the envelope of $R_k$,
$\mathsf{t}=\log g_k^{-2}$, and $c$, $r$ are the constants of that envelope. For these geometries,
therefore, $N\,V\,\vartheta\to\infty$ as $\mathsf{t}\to\infty$.
\end{proof}

For $j>k_0$ the original levels $m_0$ of the cubes of the collar region lie in $[k_0,j]$, by the
statement on the raised zone and \cite{B-CMP122a}. The original level $m_0$ of a cube and the
original cover are those fixed in the statement on the raised zone; the raised zone is the part of
the collar region on which the regularity conditions carry the weights described in the lemma below.
The boundary terms $\mathbb{B}^{(m)}(X)$ with $k_0<m<j$ number $(2L^{j-m})^d$ under one cover
cube. The anchor count of the cover-sum estimate, which rests on the near-pair bounds on the
original shells, therefore does not apply on the raised stratum. These contributions cost nothing in
the degree, provided they are added to the volume factor $V$ and not multiplied with it. The
following lemma carries this out on the elements of the slot expansion.

\begin{lemma}[The summation lemma, near groups included]\label{4.26:lem-summation}
The notation is that of Notation~\ref{4.26:not-slot-notation},
Definition~\ref{4.26:def-radii-window} and Lemma~\ref{4.26:lem-cube-count}.

Let $l$ be a stage and let $j=\ell_l$ be its level. Let $k_0$, $N_0$ and the base $L_{\circ}$ of the
raised-zone weights be as in the statement on the raised zone; $L_{\circ}$ satisfies
$L_{\circ}^2\le L/3$. For $j>k_0$, let $\mathsf{l}_j^2$ be the tile number of the statement on the
raised zone, which satisfies $(2L^{N_0})^d\le\mathsf{l}_j^2$; for $j\le k_0$ put
$\mathsf{l}_j^2:=1$. Put
\[
s^{\natural}:=\vartheta^{\natural}+K_Q/r_A+\varepsilon_S .
\]
In this lemma and its proof $N$ denotes $\check{R}_k$, Ba{\l}aban's letter in the weight
$e^{d_k(S)/N}$, and not the rank of the gauge group. Let $K_S$, $K(d)$, $A_{\flat}$, $E_0$ and
$B_0$ be the constants of the bound of the cover-sum estimate at one cover cube.

On the raised zone, the regularity conditions \cite[(1.64)--(1.69)]{B-CMP122a}, read at stage
$l-1$ as in the statement on the raised zone, carry on each cube a weight which is a power of
$L_{\circ}$. Half its exponent is:
\begin{itemize}
\item $j-m_0$ on the level-$j$ stratum $(\Omega^c_{j+1}\setminus Z''_j)^{(j)}$, where $m_0$ is
the original level (by \cite[(1.66)]{B-CMP122a}, and \cite[(1.64)]{B-CMP122a} at $m=j$);
\item $(m-k_0)^+$ on $\Gamma''_m\subset Z''_j$ for $m<j$ (by \cite[(1.64)]{B-CMP122a});
\item $0$ on the shells of level greater than $j$ (by \cite[(1.68)]{B-CMP122a}).
\end{itemize}
These numbers do not exceed $N_0-1$. For a cover cube $\square_0$, let $s(\square_0)$ be the
maximum of this number over the cubes of $\tilde{\square}_0^5$. Thus $L_{\circ}^{2s(\square_0)}$ is
the largest of these weights on $\tilde{\square}_0^5$.

Then at stage $l$ the following hold for every cover cube $\square_0$, rim cubes included.
\begin{itemize}
\item[(a)] The sum of the boundary terms of the cover-sum estimate at $\square_0$ is at
most $\mathsf{l}_j^2K_SK(d)A_{\flat}(E_0+B_0+1)$. The near group at $\square_0$ is the total
bounded in the near-pair bounds at one cover cube of the stage, together with its remainder terms.
It carries the size factor $L_{\circ}^{6s(\square_0)}\le\mathsf{l}_j^2$.
\item[(b)] For $\gamma$ sufficiently small,
\[
\sum_{c:\,S\cap\tilde{\square}_0\neq\emptyset}\mathfrak{w}(c)\,e^{d_k(S)/N}
\le\bigl[(5L)^d(l-1)+(2L)^d\mathsf{l}_j^2(N_0+3)+1\bigr]C_0^{\sharp}.
\]
Here $S$ is the set attached to $c$ in Notation~\ref{4.26:not-slot-notation}, and $d_k(S)$ is its
size in units of $\pi_k$.
\item[(c)] When summed over the cover cubes $\square_0$ of one component with weights
$\vartheta^{\natural}(\square_0)$, the parts of the sums in (a) and (b) that come from objects
assigned to levels $m\in[k_0,j)$ (the fine parts) are at most
\[
4^d(N_0+1)(100L\check{R})^d\bigl[(N_0+3)C_0^{\sharp}+E_0+B_0+1\bigr]\vartheta^{\natural},
\]
multiplied by the constants already present in (a) and (b). In particular this bound contains no
product of $V$ with $\mathsf{l}_j^2$.
\item[(c)$^{\mathrm{nr}}$] This clause, whose superscript refers to the near groups, bounds their
anchor sum. For $j>k_0$ and every component $\bar c$ of $Z$,
\begin{equation}\label{4.26:eq-summation-c}
\sum_{\square:\,\tilde{\square}^5\cap\bar c\neq\emptyset}\vartheta^{\natural}(\square)\,
L_{\circ}^{6s(\square)}
\le\bigl[V+(2^d+L_{\circ}^6)(100L\check{R})^d+(2\cdot 52^d+c_L)\mathsf{l}_j^2\bigr]
\vartheta^{\natural},
\end{equation}
where $c_L:=8\cdot 102^3L^4$. For $j\le k_0$ the left side is at most $V\vartheta^{\natural}$.
\item[(fr)] Clause (fr) completes this list; it is stated, with its proof, as a separate lemma of
this section.
\end{itemize}
Let $\mu$ denote the per-anchor total of the slot expansion
(Notation~\ref{4.26:not-slot-notation}) into which the bounds (a) and (b) enter at its anchors.
Consequently, with the constant $K_{\flat}'$ of Notation~\ref{4.26:not-slot-notation} in place
of $K_{\flat}$,
\begin{equation}\label{4.26:eq-summation-a}
\mu\le\tfrac14K_{\flat}'s^{\natural}\bigl[\mathsf{l}_j^2\bigl(E_0+B_0+1+(N_0+3)C_0^{\sharp}\bigr)
+NC_0^{\sharp}\bigr]+\tfrac14K_{\flat}'\mathfrak{p}_j
\end{equation}
and
\begin{equation}\label{4.26:eq-summation-b}
\begin{split}
4\bigl(\mathsf{M}_{\ast}+2^dV\tilde{\mathfrak{l}}\bigr)
\le K_{\flat}'\mathsf{r}\Bigl\{&V\bigl(NC_0^{\sharp}+E_0+B_0+1\bigr)\\
&+4^d(100L\check{R})^d(N_0+1)\bigl[(N_0+3)C_0^{\sharp}+E_0+B_0+1\bigr]\\
&+\bigl[(2^d+L_{\circ}^6)(100L\check{R})^d+(2\cdot 52^d+c_L)\mathsf{l}_j^2\bigr](E_0+1)
\Bigr\}s^{\natural}\\
&+K_{\flat}'\mathsf{r}V\bigl(\mathfrak{p}_j+\tilde{\mathfrak{l}}\bigr).
\end{split}
\end{equation}
\end{lemma}

\begin{proof}
The proof treats (a), (b), (c) and (c)$^{\mathrm{nr}}$ in turn, and then the two displays.

\emph{Proof of (a): the boundary terms and the chain terms of the cover-sum estimate.} By the
statement on the raised zone, the enlarged cube $\tilde{\square}_0$ is tiled by at most
$\mathsf{l}_j^2$ cubes $\square''$ of the original cover. Every object entering the sum that meets
$\tilde{\square}_0$ meets $\tilde{\square}''$ for some tile $\square''$. At each tile the two parts
of the near-pair bounds that concern the original shells hold as stated, because they are statements
on those shells. At each tile the sum of the boundary terms of the cover-sum estimate is therefore
bounded by $K_SK(d)A_{\flat}(E_0+B_0+1)$. Summing over the at most $\mathsf{l}_j^2$ tiles gives the
bound of (a) for the boundary terms. The chain terms of the cover-sum estimate are treated by the
same tiling.

\emph{Proof of (a): the near group.} By the cover-sum estimate, a near group is one resummed
element for each pair $(\square_0,\mathsf{y})$. The near-pair bound for this total is a statement
on a cover cube of the stage. It holds after the cancellations of \cite[\S\S4--5]{B-CMP109}. For
this reason it can be split neither over the original tiles nor over the individual terms $T$. The
effect of the raised zone on it is the size factor at $\square_0$, as follows.
\begin{itemize}
\item By the statement on the raised zone, the data constant on the raised zone is
$L_{\circ}^{2(j-m_0)}\mathfrak{a}_j$.
\item By the same statement, the near bounds are at most cubic in the data constant. This
produces the factor $L_{\circ}^{6s(\square_0)}$, with $s(\square_0)$ read on the regularity
conditions in the units of that statement.
\end{itemize}
This matches \cite{B-CMP122a}, which states that the regularity conditions are changed by a
factor that is a power of some number, the power being proportional to a number of overlapping
regions. The profile of the regularity parameters $\varepsilon$ entering here is the one of
\cite[(1.24), p.~182]{B-CMP122a}. Finally, since $s(\square_0)\le N_0$, the statement on the
raised zone gives
\[
L_{\circ}^{6s(\square_0)}\le L_{\circ}^{6N_0}\le L^{4N_0}\le\mathsf{l}_j^2 .
\]

\emph{Proof of (b).} Let $m(S)$ be the level of the cubes of which $S$ is a union. We cut at the
level $j$.

\emph{Case $m(S)\ge j$.} The cubes of $S$ are then not finer than those of the cover. The anchor
count of the cover-sum estimate bounds the contribution of these elements by
$(5L)^d(l-1)C_0^{\sharp}$.

\emph{Case $m(S)\le j-1$.} Write $m_S:=m(S)$, and let $d_{m_S}(S)$ be the size of $S$ in units of
$\pi_{m_S}$. Every link of the hull of $S$ climbs to a coarser level. Hence, by the climbing-link
clause of the near-pair bounds, the hull uses at most $(a_m/L)\,d_{m_S}(S)$ of the decay, where
$a_m$ is the rate at the coarser level $m$ to which the links climb. The surplus
$(a_{m_S}-a_m/L)\,d_{m_S}(S)$ is therefore free, and
\[
(a_{m_S}-a_m/L)\,d_{m_S}(S)\ge\tfrac12a_{m_S}d_{m_S}(S).
\]
This last inequality is equivalent to $a_m/L\le\tfrac12a_{m_S}$. For the rates of
Ba{\l}aban's construction, with $\beta=1/5$ and $\kappa\ge 1$, and using of $L$ only the bound
$1/L\le 1/8$, it is the inequality
\begin{equation}\label{4.26:eq-summation-1}
\frac{1.6\kappa+1}{8}\le 0.4\kappa=\tfrac12a_{m_S}.
\end{equation}
For the rates of the hull construction, the same inequality is one of the hypotheses of that
construction.

\emph{The cut at $8R_k/L$.} Let $S$ be created at stage $i'$, at level $j'$.
\begin{itemize}
\item By the inheritance statement for sets created at a stage, $S$ meets the region $R^{(i')}$
of stage $i'$, and $R^{(i')}\subset\Omega^c_{j'+2}$.
\item Since $S$ is assigned to level $m_S$, it also meets $\Omega_{m_S}$.
\item The regions are nested, $\Omega_i\supseteq\Lambda_i\supseteq\Omega_{i+1}$ for every $i$,
as stated in \cite[(2.1), p.~254]{B-CMP119}.
\item The large-field region of a term is $\Lambda^c$: \cite{B-CMP122a} states that each term of
the expansion has a large-field region, the complement $\Lambda^c$ of the small-field region of its
own step, which is a union of connected components.
\end{itemize}
Suppose $j'+2\le m_S-2$. Then $S$ meets $\Omega^c_{m_S-2}\subset\Lambda^c_{m_S-2}$, and
$\Omega_{m_S}\subset\Lambda_{m_S-1}$. So the connected hull of $S$ crosses the layers by which
$\Lambda^c_{m_S-1}$ exceeds $\Lambda^c_{m_S-2}$. Their thickness is controlled as follows.
\begin{itemize}
\item \cite{B-CMP122a} states that each renormalization step adds at least ten layers of
$MR$-cubes of the scale of that step, so that the size must exceed twenty times $MR$ at that
scale; we read this at step $m_S-1$.
\item The cubes are those of the $(m_S-1)$-lattice: \cite{B-CMP122b} states that the components
are joined by adding layers of $MR$-cubes of the next scale.
\item By \cite{B-CMP109}, after division by $M$ the cubes of the block partition become unit cubes.
\end{itemize}
Ten layers of $MR_{m_S-1}$-cubes are therefore $10R_{m_S-1}$ units of $\pi_{m_S-1}$. Since
$R_{m_S-1}\ge R_k$, this gives
\[
10R_{m_S-1}\ \text{units of }\pi_{m_S-1}
\ \ge\ 10R_{m_S-1}/L\ \ge\ 10R_k/L\ \text{units of }\pi_{m_S}.
\]
For $\gamma$ small we have $R_k\ge L$, and hence
\[
d_{m_S}(S)\ge 10R_k/L-2\ge 8R_k/L .
\]

\emph{Short elements.} Call $S$ short if $d_{m_S}(S)<8R_k/L$. By the preceding paragraph, a short
$S$ has $j'\ge m_S-3$. By the statement on the raised zone, a term is never finer than the
original level $m_0\ge k_0$ of the place where it lies, so $m_S\ge k_0$ and $j'\ge k_0-3$.
Moreover $j'\le j-2$. Hence at most $N_0+3$ stages contribute.

The anchors of $S$ under $\tilde{\square}_0$ are $\pi_{j'+1}$-cubes, by \cite[(1.63)]{B-CMP122a}.
Since $j'\ge k_0-3$ and $j-k_0\le N_0-1$ (the half-exponent $j-k_0$ at the foot of the stratum
does not exceed $N_0-1$), we have $j-j'-1\le N_0+1$. Since $\tilde{\square}_0$ has four
$\pi_j$-sides, and since $(2L^{N_0})^d\le\mathsf{l}_j^2$, their number satisfies
\[
\bigl(4L^{j-j'-1}\bigr)^d\le\bigl(4L^{N_0+1}\bigr)^d=(2L)^d(2L^{N_0})^d\le(2L)^d\mathsf{l}_j^2 .
\]
Each pair (stage, anchor) contributes at most
$\|\mathbb{C}^{(i')}\|^{A}_{(i')}\le C_0^{\sharp}$. The short elements thus contribute at most
$(N_0+3)(2L)^d\mathsf{l}_j^2C_0^{\sharp}$.

\emph{Long elements.} Call $S$ long if $d_{m_S}(S)\ge 8R_k/L$. The free surplus is then at least
$\tfrac12a_{m_S}\cdot 8R_k/L=4a_{m_S}R_k/L$. The long elements therefore contribute at most
\[
N(2L)^d\mathsf{l}_j^2e^{-4a_{m_S}R_k/L}C_0^{\sharp}\le C_0^{\sharp}.
\]
Here the inequality uses the smallness condition on $\gamma$ of the cover-sum estimate,
$Ne^{-\sigma_{\ast}R_{\ast}/L}\le 1$, with $N=\check{R}_k$ and with the rate $\sigma_{\ast}$ and
the radius $R_{\ast}$ of that estimate. This contribution is paid by the surplus and not by the rate.

\emph{Conclusion of (b).} Adding the three contributions, $(5L)^d(l-1)C_0^{\sharp}$ from
$m(S)\ge j$, $(N_0+3)(2L)^d\mathsf{l}_j^2C_0^{\sharp}$ from the short elements and $C_0^{\sharp}$
from the long ones, gives (b).

\emph{Proof of (c).} Every $S$, or $X_T$, assigned to level $m_S$ lies in $Z_{m_S+1}\cap\bar c$.
By Lemma~\ref{4.26:lem-cube-count}, this set contains at most $(100L\check{R})^d$ cubes of
$\pi_{m_S}$. The fine levels $m_S\in[k_0,j)$ are at most $N_0+1$ in number. Anchor each fine
object at one of its own cubes and sum it once. Per anchor this costs $(N_0+3)C_0^{\sharp}$ for
the objects of (b), and the constants of (a) times $(E_0+B_0+1)$ for those of (a). A fine object
meets at most $4^d$ cover cubes per own cube. The cover cubes met by a long object are summed by
the surplus of the proof of (b). Finally $\vartheta^{\natural}(\square_0)\le\vartheta^{\natural}$.
Multiplying these factors gives the bound of (c). This count covers the boundary terms and the
chain terms. A near group is not such an object; near groups are counted in
(c)$^{\mathrm{nr}}$.

\emph{Proof of \textup{(c)$^{\mathrm{nr}}$}: the cubes with $s(\square)=0$.} By
Definition~\ref{4.26:def-radii-window},
$\vartheta^{\natural}(\square)=\vartheta^{\natural}e^{-\frac12\delta d^-_{\square}}$. By the
definition of $V$ in Lemma~\ref{4.26:lem-cube-count}, the cover cubes with $s(\square)=0$
therefore contribute at most $V\vartheta^{\natural}$. For $j\le k_0$ this is the whole sum. For
$j>k_0$, every other cover cube lies in one stratum of $\mathfrak{B}$, and we treat the strata in
turn.

\emph{The stratum of level $j$: counting.} Let $s\ge 1$ be an integer. We have
\[
\{m_0\le j-s\}\cap\bar c=\Omega^c_{j-s+1}\cap\bar c\subset Z_{j-s+1}\cap\bar c,
\]
by the convention of \cite{B-CMP122a}, which is the one used in
Lemma~\ref{4.26:lem-cube-count}. By that lemma, the set $Z_{j-s+1}\cap\bar c$ is one
parallelepiped contained in a cube of side $100MR_{j-s+1}$ at scale $j-s+1$.
\begin{itemize}
\item The cubes of $\pi_j$ are $L^{s-1}$ times coarser than those of $\pi_{j-s+1}$. So the sides of
the parallelepiped are at most $100\check{R}L^{1-s}$ $\pi_j$-cubes long. At $s=0$ this agrees
with the count of $(100LR_{j+1})^d$ cubes of $\pi_j$ in Lemma~\ref{4.26:lem-cube-count}.
\item A cover cube $\square$ whose enlargement $\tilde{\square}^5$ meets the parallelepiped lies
within $13$ side lengths of $\pi_j$-cubes of it.
\end{itemize}
Therefore
\begin{equation}\label{4.26:eq-summation-2}
\#\{\square:\ s(\square)\ge s\}\le\bigl(100\check{R}L^{1-s}+26\bigr)^d
\le 2^d\bigl[(100L\check{R})^dL^{-ds}+26^d\bigr].
\end{equation}

\emph{The stratum of level $j$: summing.} Since
$L_{\circ}^{6s(\square)}\le\sum_{1\le s\le s(\square)}L_{\circ}^{6s}$ and
$\vartheta^{\natural}(\square)\le\vartheta^{\natural}$, we may interchange the sums over $\square$
and over $s$ and insert \eqref{4.26:eq-summation-2}. The sum over $\{s(\square)\ge 1\}$ is then at
most
\[
\vartheta^{\natural}\sum_{s=1}^{N_0}L_{\circ}^{6s}\,2^d\bigl[(100L\check{R})^dL^{-ds}+26^d\bigr].
\]
From $L_{\circ}^2\le L/3$ and $d=4$ we get $L_{\circ}^6L^{-d}\le 1/(27L)$. Since
$\sum_{s\ge1}(27L)^{-s}<1$, the first part is at most $2^d(100L\check{R})^d\vartheta^{\natural}$.
The second part is at most
\[
2^d26^d\cdot 2L_{\circ}^{6N_0}\vartheta^{\natural}\le 2\cdot 52^d\mathsf{l}_j^2\vartheta^{\natural},
\]
by the inequality $L_{\circ}^{6N_0}\le\mathsf{l}_j^2$ of the statement on the raised zone.

\emph{Inside $Z''_j$: the size of $Z''_k$.} Here the cover follows the scale of the stratum, as in
\cite{B-CMP119}: its cubes are finer than those of $\pi_j$, and the count by $\pi_j$-sides above
does not apply to them. By the definition of $N_0$, $L^{-N_0+1}R_{k_0+1}=1$, as in
\cite{B-CMP122a}. With Lemma~\ref{4.26:lem-cube-count}, $Z''_k\subset Z_{k_0+1}\cap\bar c$ then
has sides of at most $100$ $M$-cubes of scale $k$.

This rests on \cite{B-CMP122a}, which states
\[
Z''_k=(\Omega^{\sim5}_{k_0+1})^c\cap Z,\qquad \Lambda=(\Omega^{\sim4}_{k_0+1})^c\cap Z ,
\]
the second identity being \cite[(1.73)]{B-CMP122a}. The accompanying text of \cite{B-CMP122a}
states that the latter domain, obtained by adding one layer of $M$-cubes to $Z''_k$, is a union of
$M$-cubes and, by condition (i) there, a rectangular parallelepiped contained in a cube of size
$100M$. The operation ${}^{\sim}$ enlarges $\Omega$, so
$Z''_k\subset\Omega^c_{k_0+1}\cap Z\subset Z_{k_0+1}$.

\emph{Inside $Z''_j$: the thin layers, $k_0<m<j$.} By \cite{B-CMP122a}, $Z''_{m+1}\setminus Z''_m$
is one layer of $M$-cubes of scale $m+1$. It surrounds a parallelepiped with sides of at most
$100L^{k-m-1}$ such cubes. So the layer consists of at most $2d(102L^{k-m-1})^{d-1}$ cubes of scale
$m+1$, that is, with $d=4$, at most $8\cdot102^3L\cdot L^{3(k-m)}$ cubes of scale $m$.

The layer is $L>5$ own cubes thick. So $\tilde{\square}^5$ reaches layer $m+1$ and no further; for
$m=j-1$ this is the foot of the stratum, with exponent $j-k_0$. Hence
$s(\square)\le t:=m+1-k_0$. By $L_{\circ}^6\le L^3/27$,
\[
L^{3(k-m)}L_{\circ}^{6t}\le L^{3(N_0+1)}27^{-t}.
\]
Since $\sum_{t\ge2}27^{-t}<1$, and using $L^{4N_0}\le\mathsf{l}_j^2$, the thin layers contribute in
total at most
\[
8\cdot102^3L^4\cdot L^{3N_0}\,\vartheta^{\natural}\le c_L\mathsf{l}_j^2\vartheta^{\natural}.
\]

\emph{The stratum $\Gamma''_{k_0}$.} Here $s\le 1$. By Lemma~\ref{4.26:lem-cube-count}, there are
at most $(100L\check{R})^d$ cubes of $\pi_{k_0}$ in $Z_{k_0+1}\cap\bar c$. The contribution is at
most $L_{\circ}^6(100L\check{R})^d\vartheta^{\natural}$.

\emph{The strata $\Gamma''_m$, $m<k_0$.} By the statement on the raised zone,
$Z''_{k_0+1}\setminus Z''_{k_0}\supseteq\Omega_{k_0}\setminus\Omega^{\sim7}_{k_0+1}$. This set is
at least one layer of $MR_{k_0+1}$-cubes thick (eight layers less seven), that is, at least
$LR_{k_0+1}>10$ sides of $\pi_{k_0}$-cubes. So $\tilde{\square}^5$ cannot reach the raised zone
from these strata, and $s=0$ there.

\emph{The shells of level greater than $j$.} The old shell of level $j$ consists of at least eight
layers of $MR_{j+1}$-cubes. It separates these shells from the raised zone, so $s=0$ there.

\emph{Conclusion of \textup{(c)$^{\mathrm{nr}}$}.} The contributions are:
\begin{itemize}
\item $V\vartheta^{\natural}$ from $s=0$;
\item $2^d(100L\check{R})^d\vartheta^{\natural}$ and $2\cdot52^d\mathsf{l}_j^2\vartheta^{\natural}$
from the stratum of level $j$;
\item $c_L\mathsf{l}_j^2\vartheta^{\natural}$ from the thin layers;
\item $L_{\circ}^6(100L\check{R})^d\vartheta^{\natural}$ from $\Gamma''_{k_0}$.
\end{itemize}
Their sum is the right side of \eqref{4.26:eq-summation-c}.

\emph{The displays \eqref{4.26:eq-summation-a} and \eqref{4.26:eq-summation-b}.} They follow from
the totals of the slot expansion, with $K_{\flat}'$ of Notation~\ref{4.26:not-slot-notation} in
place of $K_{\flat}$, by inserting:
\begin{itemize}
\item (a) and (b) at the anchors of $\mu$, which gives \eqref{4.26:eq-summation-a};
\item (c) in the anchor sum of $\mathsf{M}_{\ast}$;
\item the anchor sum of the near groups,
$\sum_{\square_0}\vartheta^{\natural}(\square_0)L_{\circ}^{6s(\square_0)}$, bounded by
\eqref{4.26:eq-summation-c}. It enters once with the factor
$c_{\alpha}C_R^{\natural}(E_0+1)e^{C_2\kappa_1}$ for the near group proper, and once with the
constant $K''$ of the remainders; here $c_{\alpha}$, $C_R^{\natural}$, $C_2$ and $K''$ are the
constants of the near group and of its remainders in the near-pair bounds.
\end{itemize}
Accordingly, the bracket in \eqref{4.26:eq-summation-b} has three members:
\begin{itemize}
\item the volume total of the slot expansion, $V(NC_0^{\sharp}+E_0+B_0+1)$, which enters
$\mathsf{M}_{\ast}$;
\item the fine parts bounded in (c);
\item the members of \eqref{4.26:eq-summation-c} other than $V$, multiplied by $E_0+1$.
\end{itemize}
If the near groups were anchored through the bound
$\sum_{\square_0}\vartheta(\square_0)\le 2^dV\vartheta$ of the cover-sum estimate, with the factor
of (a) at each cube, the product $V\mathsf{l}_j^2$ would appear, a factor of degree
$10r-\mathsf{m}$ in $\mathsf{t}:=\log g_k^{-2}$, which is positive under the standing conditions on
the degrees.
By \eqref{4.26:eq-summation-c} it is replaced by the sum
$V+(2^d+L_{\circ}^6)(100L\check{R})^d+(2\cdot52^d+c_L)\mathsf{l}_j^2$.
\end{proof}

\emph{Degrees in $\mathsf{t}$.} Below, $+0^{+}$ in a degree records the factors
polynomial in $N_0$. The members have the following degrees:
\begin{itemize}
\item $\mathsf{l}^2(N_0+3)$ has degree $4r+0^{+}$, where $\mathsf{l}^2$ is the fixed size
parameter of the statement on the raised zone, with $\mathsf{l}_j^2\le16\mathsf{l}^2$;
\item $(100L\check{R})^d(N_0+1)(N_0+3)$ has degree at most $6r+0^{+}$;
\item $N\mathsf{r}\bar{V}$ has degree at most $7r$;
\item the members of \eqref{4.26:eq-summation-c} have degrees at most $6r$ and $4r$, and both
kinds already occur in the bracket.
\end{itemize}
For $\gamma$ sufficiently small,
\[
\bigl[4^d(N_0+1)(N_0+3)+2^d+L_{\circ}^6\bigr](100L\check{R})^d
\le\frac{5^d+L_{\circ}^6}{C_V}\,N\bar{V},
\]
so these members are contained in the term $K'N\mathsf{r}\bar{V}$, where $K'$ is the generic
prefactor of the bracket in which these totals are consumed. Further,
$(2\cdot52^d+c_L)\mathsf{l}_j^2$ is contained in $K'(N_0+3)\mathsf{l}^2$. This uses the bound
$\mathsf{l}_j^2\le16\mathsf{l}^2$ of the statement on the raised zone, and requires replacing the
generic prefactor $K'$ by $16(1+2\cdot52^d+c_L)K'$.

In what follows $\mathfrak{a}$ is a live piece at the site $k$ as in Definition~\ref{4.26:def-frozen-structure}: $k_a<k$ and $W_{\mathfrak{a}}\subset Z^c$. Its own letters are $N_a$, $h_a$, $N_0^a$, $k_0=k_a-N_0^a$ and $j^+\le k_a$. The sets $\Omega_i$, $\Lambda_i$, $Z''_i$, $\Gamma''_m$ are those fixed there, cut to $W_{\mathfrak{a}}$, and we assume $j^+>k_0$.

The remaining letters are as follows.
\begin{itemize}
\item $\mathfrak{c}$ is the original-level function on $W_{\mathfrak{a}}$. Its strata are described in Lemma~\ref{4.26:lem-original-level-strata}, and its values on the strata of the frozen structure, together with the frozen level $\ell_{\mathfrak{a}}$, are listed in the weight table \eqref{4.26:eq-frozen-structure-a}.
\item $g$ is the lift on $W_{\mathfrak{a}}$.
\item $\ell_i$ is the spacing of the lattice of level $i$, so that $\ell_{i-1}=\ell_i/L$ and a $\pi_i$-cube has side $M\ell_i$. A length of $t$ $\pi_i$-sides means the length $tM\ell_i$. The spacing $\ell_i$ is a length and is not to be confused with the frozen level $\ell_{\mathfrak{a}}$, which is a level.
\item $\mathcal{C}$ is the cover of the stage and $\mathcal{C}^1$ the graded cover, both as in Definition~\ref{4.26:def-graded-cover}.
\item $d_{\mathrm{sc}}$ is the distance $d^{\wedge}_{\mathfrak{B}}$, and $\operatorname{len}_{\mathfrak{B}}\Gamma$ is the length of a contour $\Gamma$ measured with it.
\item $w$ is the lower bound, fixed with $d^{\wedge}_{\mathfrak{B}}$, for the length in this distance of a crossing of one cube of side $M$ at its own level.
\item $\delta$, the stages $n$, the regions $R_n$, the components $\bar c$ of $Z$ and the count $V^0$ are those of the cover-cube count (Lemma~\ref{4.26:lem-cube-count}, display \eqref{4.26:eq-cube-count-a}). The regions $R_n$ are indexed by the stage $n$ and are not to be confused with the numbers $R_i$ of the levels $i$.
\end{itemize}

We put
\begin{equation*}
\mathsf{L}_{\mathfrak{a}}:=\{g\ge 1\}\cap W_{\mathfrak{a}}
=(\Omega^c_{j^+}\setminus Z''_{k_0+1})\cap W_{\mathfrak{a}},
\qquad
\varrho^F_m:=\frac{R_{m+1}\ell_{m+1}}{\ell_{j^+}}\quad(k_0\le m<j^+),
\end{equation*}
and call $\mathsf{L}_{\mathfrak{a}}$ the \emph{lifted set} of $\mathfrak{a}$.

For a cube $\square\in\mathcal{C}\cup\mathcal{C}^1$ meeting $\mathsf{L}_{\mathfrak{a}}$ we put $u(\square):=j^+-\min_{\square}\mathfrak{c}$. Then $u(\square)\ge 1$, and $g\le u(\square)$ on $\square$.

The multiplicity weights are
\begin{equation*}
\mathsf{m}_F(\square):=\bigl(2L^{u(\square)+1}\bigr)^d\bigl(u(\square)+4\bigr)
\end{equation*}
for $\square\in\mathcal{C}$ meeting a lifted set, and $\mathsf{m}_F(\square):=1$ for every other cube of $\mathcal{C}$ and for every cube of $\mathcal{C}^1$.

Let $\partial_{\mathfrak{a}}$ be the set of $\pi_{j^+}$-cubes of $\Omega_{j^+}$ touching $\Omega^c_{j^+}\cap W_{\mathfrak{a}}$. For a stage $n$ and a component $\bar c$ of $Z$ we put
\begin{equation*}
D_F(\square):=d_{\mathrm{sc}}(\square,R_n\cap\bar c),\qquad
D(\Delta):=d_{\mathrm{sc}}(\Delta,R_n\cap\bar c),\qquad
d^-_{\mathfrak{a}}:=\min_{\Delta\in\partial_{\mathfrak{a}}}D(\Delta).
\end{equation*}

\begin{lemma}[Counts on lifted sets paid by the lifted shells]\label{4.26:lem-lifted-counts}
Let $\mathfrak{a}$ be a live piece at the site $k$ as above, with $k_a<k$, $W_{\mathfrak{a}}\subset Z^c$ and $j^+>k_0$, and let $d=4$, $L\ge 12$ and $R_{\min}\ge 5L$. Assume moreover
\begin{equation*}
\delta w\ge 16\ln L+4 .
\end{equation*}

Statement (F0) consists of the first line of \eqref{4.26:eq-lifted-counts-a} together with two geometric facts. First, for $k_0<m<j^+$ the old shell $m$ is at least $8\varrho^F_m$ $\pi_{j^+}$-sides wide, and it carries frozen level $j^+$. Second, for $k_0\le m<j^+$, $\Omega^c_{m+1}\cap W_{\mathfrak{a}}$ lies in a cube of $100\varrho^F_m$ $\pi_{j^+}$-sides; this second fact rests on the second clause of condition~(ii) imposed in \cite{B-CMP122a} on the components to which the operation $\mathbf{R}$ is applied, read at the site $k_a$, and on the nesting condition at the site $k_a$.

In (F4) we put
\begin{equation*}
V^{\mathrm{fr}}:=\sup_{n,\bar c}\Bigl[\sum_{\square\in\mathcal{C}^1}e^{-\frac12\delta D_F(\square)}
+\sum_{\square\in\mathcal{C}}\mathsf{m}_F(\square)\,e^{-\frac12\delta D_F(\square)}\Bigr],
\end{equation*}
and $V^{0,\mathrm{th}}$ denotes the supremum over $n,\bar c$ of the sum defining $V^0$, restricted to the cubes $\square\in\mathcal{C}$ that meet no lifted set or have $u(\square)\le 1$.

Then the following hold, where (F1) and (F2) are asserted for every stage $n$, every component $\bar c$ of $Z$ and every $\square\in\mathcal{C}\cup\mathcal{C}^1$ meeting $\mathsf{L}_{\mathfrak{a}}$, with $m:=\min_{\square}\mathfrak{c}$ and $u:=u(\square)$, and (F3) is asserted for every stage $n$ and every component $\bar c$ of $Z$:
\begin{equation}\label{4.26:eq-lifted-counts-a}
\begin{aligned}
&\text{(F0)}\quad 1\le\varrho^F_m\le\varrho^F_{m'}\qquad(k_0\le m\le m'<j^+),\\
&\text{(F1)}\quad D_F(\square)\ge d^-_{\mathfrak{a}}+6w(u-1)\varrho^F_m,
\qquad d^-_{\mathfrak{a}}\ge 2R_kw,\\
&\text{(F2)}\quad \bigl(3L^{u+1}\bigr)^{2d}(u+4)\,e^{-\frac12\delta D_F(\square)}\le 1,\\
&\text{(F3)}\quad \sum_{\substack{\square\in\mathcal{C}\cup\mathcal{C}^1:\\
\square\cap\mathsf{L}_{\mathfrak{a}}\neq\emptyset,\ u(\square)\ge 2}}
\mathsf{m}_F(\square)\,e^{-\frac12\delta D_F(\square)}
\le L^{-20}e^{-\frac12\delta d^-_{\mathfrak{a}}},\\
&\text{(F4)}\quad V^{\mathrm{fr}}\le 6(2L^2)^dV^0,\qquad
V^0\le\bigl(1+2^dL^{-20}\bigr)V^{0,\mathrm{th}}\le 2V^{0,\mathrm{th}},\\
&\phantom{\text{(F4)}}\quad \max\{V^0,V^{\mathrm{fr}}\}\le C_FV^{0,\mathrm{th}},
\qquad C_F:=12(2L^2)^d .
\end{aligned}
\end{equation}

The five statements (F0)--(F4) are absolute. None of them depends on $k-k_a$, $N_a$, $N_0^a$, $j^+-k_0$, or on the coupling.
\end{lemma}

\begin{proof}
\emph{Proof of (F0): monotonicity of $\varrho^F$.} We use the inequality
\begin{equation*}
LR_{i+1}\ge R_i\qquad\text{for every level } i ,
\end{equation*}
which is among the properties of the envelope $\check{R}$ of the numbers $R_i$ used in the cover-cube count (Lemma~\ref{4.26:lem-cube-count}). It rests on two facts. The number $R_j$ is defined in \cite[(2.5)]{B-CMP119} as the smallest number of the form $L^i$ with $R_j\ge(\log g_j^{-2})^r$. In \cite{B-CMP122a} the passage to the next step is described in the words ``Passing to the next step it is usually rescaled by $L^{-1}$, but in some steps the number $R_k$ decreases by the factor $L^{-1}$''.
Since $\ell_{i+1}=L\ell_i$, it follows that
\begin{equation*}
R_{i+1}\ell_{i+1}=LR_{i+1}\ell_i\ge R_i\ell_i ,
\end{equation*}
so $i\mapsto R_i\ell_i$ is non-decreasing. Dividing by $\ell_{j^+}$ gives $\varrho^F_m\le\varrho^F_{m'}$ for $k_0\le m\le m'<j^+$.

The number $N_0$ is the smallest positive integer with $L^{-N_0+1}MR_{k-N_0+1}=M$, that is, with $L^{-N_0+1}R_{k-N_0+1}=1$ \cite{B-CMP122a}. Applied at $k=k_a$, with $k_0+1=k_a-N_0^a+1$ and $\ell_{k_0+1}=L^{-N_0^a+1}\ell_{k_a}$, this gives
\begin{equation*}
R_{k_0+1}\ell_{k_0+1}=L^{-N_0^a+1}R_{k_a-N_0^a+1}\,\ell_{k_a}=\ell_{k_a}.
\end{equation*}
Hence $\varrho^F_m\ge\varrho^F_{k_0}=\ell_{k_a}/\ell_{j^+}$, and this ratio is at least $1$ because $j^+\le k_a$. This proves the first line of \eqref{4.26:eq-lifted-counts-a}.

\emph{Proof of (F0): width and level of the old shells.} Let $k_0<m<j^+$. By the weight table \eqref{4.26:eq-frozen-structure-a}, on the old shell $(\Omega_m\setminus\Omega_{m+1})\cap W_{\mathfrak{a}}$ one has $\mathfrak{c}=m$ and frozen level $\ell_{\mathfrak{a}}=j^+$. So this shell is the $\mathfrak{c}$-shell $m$, and it carries frozen level $j^+$. By Lemma~\ref{4.26:lem-original-level-strata}(iii) its width is at least eight layers of $MR_{m+1}$-cubes of the $(m+1)$-lattice. That is, it is at least
\begin{equation*}
8MR_{m+1}\ell_{m+1}=8\varrho^F_m\,M\ell_{j^+},
\end{equation*}
which is $8\varrho^F_m$ $\pi_{j^+}$-sides.

\emph{Proof of (F0): the container.} Let $k_0\le m<j^+$. We have $\Omega^c_{m+1}\subset\Lambda^c_{m+1}=Z_{m+1}$, where $Z_{m+1}$ is the large-field region of level $m+1$. Write $X_{\mathfrak{a}}$ for the component to which the operation $\mathbf{R}$ was applied at the site $k_a$ and which carried $\mathfrak{a}$; it satisfies conditions~(i) and~(ii) of \cite{B-CMP122a}.

The piece $W_{\mathfrak{a}}$ entered the first stage at the site $k_a$. The nesting condition imposed there states that during the $N_a$ steps preceding $k_a$ no components were joined inside the union of the components entering the first stage at the site $k_a$; this union contains $X_{\mathfrak{a}}$. Consequently, by the lemma on components under the nesting condition, $Z_{m+1}\cap X_{\mathfrak{a}}$ is a single component. It is the only component of $Z_{m+1}$ meeting $X_{\mathfrak{a}}$.

By the second clause of condition~(ii) of \cite{B-CMP122a}, each previous component is one parallelepiped contained in a cube of side $100MR$ at its own scale \cite{B-CMP119}. For $Z_{m+1}$ this cube has side
\begin{equation*}
100MR_{m+1}\ell_{m+1}=100\varrho^F_m\,M\ell_{j^+},
\end{equation*}
and so $\Omega^c_{m+1}\cap W_{\mathfrak{a}}$ lies in a cube of $100\varrho^F_m$ $\pi_{j^+}$-sides.

Neither of the two conditions can be dropped here. The second clause of condition~(ii) bounds each previous component but not their number, and the first clause forbids the creation of new large-field regions but not the joining of components (the two clauses of condition~(ii) being those stated in \cite{B-CMP122b}). The argument for (F1) below does not use the nesting condition. It uses only the fact, stated in \cite{B-CMP122a}, that the previous regions contained in the component satisfy condition~(i) on the corresponding scales.

\emph{Proof of (F1): the first inequality.} Fix a stage $n$, a component $\bar c$ of $Z$, and $\square\in\mathcal{C}\cup\mathcal{C}^1$ meeting $\mathsf{L}_{\mathfrak{a}}$, with $m=\min_{\square}\mathfrak{c}$ and $u=u(\square)$. We have $R_n\subset Z$ and $W_{\mathfrak{a}}\subset Z^c$, and $\square$ meets $\mathsf{L}_{\mathfrak{a}}\subset\Omega^c_{j^+}\cap W_{\mathfrak{a}}$. Therefore every contour $\Gamma$ from $\square$ to $R_n\cap\bar c$ leaves $\Omega^c_{j^+}\cap W_{\mathfrak{a}}$ through some $\Delta\in\partial_{\mathfrak{a}}$, and
\begin{equation}\label{4.26:eq-lifted-counts-1}
\operatorname{len}_{\mathfrak{B}}\Gamma\ge
\bigl(\text{length of the part of }\Gamma\text{ inside }\Omega^c_{j^+}\bigr)+D(\Delta).
\end{equation}

The cube $\square$ meets at most the $\mathfrak{c}$-shells $m$ and $m+1$.
\begin{itemize}
\item For $\square\in\mathcal{C}^1$ this is Lemma~\ref{4.26:lem-graded-box-geometry}(a), applied to $\square_0=\tilde{\square}^5\supset\square$.
\item For $\square\in\mathcal{C}$ it holds because the side of $\square$ is at most $2M\ell_{j^+}$, which is less than the width $8\varrho^F_{m+1}M\ell_{j^+}\ge 8M\ell_{j^+}$ of the shell $m+1$ given by (F0), when $m+1<j^+$; for $m+1=j^+$, that is $u=1$, the claim is not needed below.
\end{itemize}

In either case $\square$ reaches at most two $\pi_{j^+}$-sides into the shell $m+1$.

Let $u\ge2$. The part of $\Gamma$ inside $\Omega^c_{j^+}$ crosses the rest of the shell $m+1$ and then the shells $m+2,\dots,j^+-1$. All of these are old shells above $k_0$ at frozen level $j^+$. There are $u-1$ indices $m'$ with $m<m'<j^+$, and $\varrho^F_{m'}\ge\varrho^F_m$ by (F0). Using the widths of (F0) we obtain at least
\begin{equation*}
8\sum_{m<m'<j^+}\varrho^F_{m'}-2\;\ge\;8(u-1)\varrho^F_m-2\;\ge\;6(u-1)\varrho^F_m
\end{equation*}
cubes of side $M\ell_{j^+}$ at their own level $j^+$. The last inequality holds because $2(u-1)\varrho^F_m\ge2$ for $u\ge2$ and $\varrho^F_m\ge1$.

By the conversion of lengths into the units of $d^{\wedge}_{\mathfrak{B}}$, each such crossing contributes at least $w$. So the inside part has length at least $6w(u-1)\varrho^F_m$. For $u=1$ this lower bound is $0$, which is trivially true. Taking the infimum over $\Gamma$ in \eqref{4.26:eq-lifted-counts-1} and using $D(\Delta)\ge d^-_{\mathfrak{a}}$ gives the first inequality of (F1).

\emph{Proof of (F1): the second inequality.} Let $\Delta\in\partial_{\mathfrak{a}}$. We have $\Delta\subset\Omega^c_{j^++1}$; and $\Omega^c_{j^++1}\subset\Omega^c_k$, since $j^+\le k_a<k$ gives $j^++1\le k$. Hence $\Delta\subset\Omega^c_k$.

Let $W'$ be the component of $\Lambda^c_k$ containing $W_{\mathfrak{a}}$. Since $W_{\mathfrak{a}}\subset Z^c$, the set $W'$ is not $\bar c$. A contour from $\Delta$ to $R_n\cap\bar c$ therefore passes through $\Lambda_k$, and before that through the band $W'\cap\Omega_k$.

This band is at least two layers of $MR_k$-cubes of level $k$ wide. That is \cite[(3.20)]{B-CMP119}, together with the separation statement of \cite{B-CMP119}, whose proviso is met at $j=k$. The band is not touched by the stages of the present site. By the same conversion of lengths, $D(\Delta)\ge 2R_kw$, and hence $d^-_{\mathfrak{a}}\ge2R_kw$.

\emph{Proof of (F2).} By (F1) and $\varrho^F_m\ge1$,
\begin{equation*}
\tfrac12\delta D_F(\square)\ge\delta wR_k+3\delta w(u-1).
\end{equation*}
Together with $\delta w\ge16\ln L+4$ this gives
\begin{equation*}
e^{-\frac12\delta D_F(\square)}\le e^{-\delta wR_k}e^{-3\delta w(u-1)}
\le L^{-16R_k}e^{-4R_k}\cdot L^{-48(u-1)}e^{-12(u-1)}.
\end{equation*}
For $d=4$ we have $(3L^{u+1})^{2d}(u+4)=3^8L^{8u+8}(u+4)$, so the left side of (F2) is at most
\begin{equation*}
L^{8u+8-16R_k-48(u-1)}\cdot 3^8(u+4)e^{-12(u-1)}\cdot e^{-4R_k}.
\end{equation*}

We bound the power of $L$ and the product of the remaining two factors separately.
\begin{itemize}
\item The power of $L$ is at most $1$. Indeed $8u+8\le48(u-1)+16$, which is equivalent to $40\le40u$. Moreover $48(u-1)+16\le48(u-1)+16R_k$, because $R_k\ge1$.
\item The product of the remaining two factors is at most $1$. The function $u\mapsto(u+4)e^{-12(u-1)}$ is decreasing for $u\ge1$. Hence
\begin{equation*}
3^8(u+4)e^{-12(u-1)}\le3^8\cdot5<e^{11}.
\end{equation*}
Finally $e^{11}\le e^{4R_k}$, since $R_k\ge R_{\min}\ge5L\ge60$.
\end{itemize}

The product is therefore at most $1$.

\emph{Proof of (F3): counting the cubes.} Fix $m$ with $u=j^+-m\ge2$, and let $S_m$ be the part of the sum in (F3) over the cubes with $\min_{\square}\mathfrak{c}=m$.

The set $\{\mathfrak{c}=m\}$ is a union of $\pi_m$-cubes, and all the partitions involved are compatible. So a cube $\square$ with $\min_{\square}\mathfrak{c}=m$ contains a $\pi_m$-cube of $\Omega^c_{m+1}\cap W_{\mathfrak{a}}$.

A $\pi_m$-cube lies in at most $2^d$ cubes of a cover at any one scale. It lies in cubes of at most two scales of $\mathcal{C}^1$ and of at most four scales of $\mathcal{C}$, for the following reason.
\begin{itemize}
\item A cube of the family attached to a stratum sticks out of that stratum by at most one of its own cubes.
\item A thin stratum $\Gamma''_{n'}$ is one $\pi_{n'+1}$-cube thick.
\item Hence the only families reaching a point of $\Gamma''_{n'}$ are those of indices $n'-1$, $n'$, $n'+1$ and $n'+2$. The family of index $n'+2$ does so because its cubes stick out by one $\pi_{n'+2}$-cube, which is the width of $\Gamma''_{n'+1}$, down to $\partial\Gamma''_{n'}$.
\end{itemize}

Thus a $\pi_m$-cube lies in at most $6\cdot2^d\le2^{d+3}$ cubes of $\mathcal{C}\cup\mathcal{C}^1$ (for $d=4$, $96\le128$).

By (F0), $\Omega^c_{m+1}\cap W_{\mathfrak{a}}$ lies in a cube of $100\varrho^F_m$ $\pi_{j^+}$-sides. Since $\ell_{j^+}=L^u\ell_m$, this is $100\varrho^F_mL^u$ $\pi_m$-sides. It therefore contains at most $(100\varrho^F_mL^u)^d$ cubes of $\pi_m$, and the number of cubes in $S_m$ is at most
\begin{equation*}
2^{d+3}(100\varrho^F_mL^u)^d .
\end{equation*}

\emph{Proof of (F3): bounding each term.} Each cube in $S_m$ satisfies
\begin{equation*}
\mathsf{m}_F(\square)\le(2L^{u+1})^d(u+4),
\end{equation*}
with equality on $\mathcal{C}$ and the value $1$ on $\mathcal{C}^1$. By (F1) and $\delta w\ge16\ln L+4$,
\begin{equation*}
e^{-\frac12\delta D_F(\square)}\le e^{-\frac12\delta d^-_{\mathfrak{a}}}\,
L^{-48(u-1)\varrho^F_m}e^{-12(u-1)\varrho^F_m}.
\end{equation*}
Therefore
\begin{equation*}
\begin{split}
S_m\le{}&2^{d+3}(100\varrho^F_m)^dL^{du}\cdot(2L^{u+1})^d(u+4)\\
&\cdot L^{-48(u-1)\varrho^F_m}e^{-12(u-1)\varrho^F_m}\cdot e^{-\frac12\delta d^-_{\mathfrak{a}}}.
\end{split}
\end{equation*}

Let $d=4$. Using $\varrho^F_m\ge1$ in the exponent of $L$, the powers of $L$ combine to at most
\begin{equation*}
L^{8u+4-48(u-1)}=L^{52-40u}=L^{-28-40(u-2)} .
\end{equation*}
The numerical factor is $2^{11}\cdot10^8\,(\varrho^F_m)^4(u+4)e^{-12(u-1)\varrho^F_m}$. The function $\varrho\mapsto\varrho^4e^{-12(u-1)\varrho}$ is decreasing for $\varrho\ge1$, and $u\mapsto(u+4)e^{-12(u-1)}$ is decreasing for $u\ge2$. So the numerical factor is at most
\begin{equation*}
2^{11}\cdot10^8\cdot6e^{-12}<8\cdot10^6<12^7\le L^7 .
\end{equation*}
Hence
\begin{equation*}
S_m\le L^{-21-40(u-2)}e^{-\frac12\delta d^-_{\mathfrak{a}}}.
\end{equation*}

Distinct $m$ give distinct $u\ge2$. Summing,
\begin{equation*}
\sum_{u\ge2}L^{-21-40(u-2)}=\frac{L^{-21}}{1-L^{-40}}\le L^{-20},
\end{equation*}
which is (F3).

\emph{Proof of (F4): the splitting.} Fix $n$ and $\bar c$, and split the cubes of the two sums defining $V^{\mathrm{fr}}$ into three classes.

$(\alpha)$ \emph{Cubes meeting no lifted set.} There $\mathcal{C}^1$ coincides with $\mathcal{C}$ (Definition~\ref{4.26:def-graded-cover}) and $\mathsf{m}_F=1$. Each of the two sums is then at most $V^0$, so their contribution is at most $2V^0$.

$(\beta)$ \emph{Cubes with $u(\square)=1$, of $\mathcal{C}^1$.} Such a cube has scale $j^+-1$, or scale $j^+$, and in the latter case it also belongs to $\mathcal{C}$. We show that it lies in a cube $\square'\in\mathcal{C}$ of scale $j^+$.

The cover $\mathcal{C}$ at scale $j^+$ consists of cubes that are unions of $2^d$ neighbouring cubes from $\pi_{j^+}$ \cite{B-CMP119}. Every such union occurs, because the construction of the partition of unity in \cite{B-CMP109} has $\zeta=0$ for $|t|\ge\frac23$ and $\sum\zeta_{\square}=1$, which places a centre at every vertex. A cube of side at most $M\ell_{j^+}$ lies in the $2^d$-union based at the $\pi_{j^+}$-cube of its lowest corner. Our cube qualifies, since $2M\ell_{j^+-1}=2M\ell_{j^+}/L\le M\ell_{j^+}$.

Assign each cube of $\mathcal{C}^1$ with $u=1$ to one such $\square'$. Then $D_F(\square)\ge D_F(\square')$, because $\square\subset\square'$. At most $(2L-1)^d\le(2L)^d$ cubes are assigned to each $\square'$. The contribution of these cubes is therefore at most $(2L)^dV^0$.

$(\beta')$ \emph{Cubes with $u(\square)=1$, of $\mathcal{C}$.} For these $\mathsf{m}_F=5(2L^2)^d$, so their contribution is at most $5(2L^2)^dV^0$.

$(\gamma)$ \emph{Cubes with $u(\square)\ge2$.} By (F3), summed over the pieces, their contribution is at most $L^{-20}\sum_{\mathfrak{a}}e^{-\frac12\delta d^-_{\mathfrak{a}}}$. We bound this sum by $2^dV^{0,\mathrm{th}}$.

The minimising $\Delta\in\partial_{\mathfrak{a}}$ lies in a cube $\square_{\mathfrak{a}}\in\mathcal{C}$ of scale $j^+$ off the lifted set. Then $D_F(\square_{\mathfrak{a}})\le D(\Delta)=d^-_{\mathfrak{a}}$, because $\Delta\subset\square_{\mathfrak{a}}$. A cube of $\mathcal{C}$ of scale $J$ serves at most $2^d$ pieces, for two reasons.
\begin{itemize}
\item Nested live pieces $\mathfrak{a}$ inside $\mathfrak{b}$ have different values of $j^+$. By the lemma on nested arrests, $j^+_a\le k_a\le h_b<j^+_b$.
\item Disjoint pieces have disjoint sets $\Omega^c_J\cap W$. These are unions of aligned cubes of side $MR_J\ell_J>4M\ell_J$, and at most $2^d$ of them touch a set of diameter $4M\ell_J$.
\end{itemize}

Hence the contribution of $(\gamma)$ is at most $2^dL^{-20}V^{0,\mathrm{th}}$. This is at most $V^0$, since $2^dL^{-20}\le1$ and $V^{0,\mathrm{th}}\le V^0$.

\emph{Proof of (F4): the totals.} Adding the four contributions and taking the supremum over $n,\bar c$,
\begin{equation*}
V^{\mathrm{fr}}\le\bigl(2+(2L)^d+5(2L^2)^d+1\bigr)V^0\le6(2L^2)^dV^0,
\end{equation*}
because $3+(2L)^d\le(2L^2)^d$.

For the bound on $V^0$, split its sum into the cubes counted in $V^{0,\mathrm{th}}$ and the cubes of $\mathcal{C}$ meeting a lifted set with $u\ge2$. On the latter class $e^{-\frac12\delta D_F}\le\mathsf{m}_F\,e^{-\frac12\delta D_F}$, since $\mathsf{m}_F\ge1$. So by $(\gamma)$ that class contributes at most $2^dL^{-20}V^{0,\mathrm{th}}$. Hence
\begin{equation*}
V^0\le(1+2^dL^{-20})V^{0,\mathrm{th}}\le2V^{0,\mathrm{th}},
\end{equation*}
and then
\begin{equation*}
\max\{V^0,V^{\mathrm{fr}}\}\le6(2L^2)^d\cdot2V^{0,\mathrm{th}}=C_FV^{0,\mathrm{th}} .
\end{equation*}

\emph{Independence of the constants.} The constants appearing in (F0)--(F4), namely $8$, $100$, $6$, $2$, $L^{-20}$, $6(2L^2)^d$ and $C_F$, depend on $L$ and $d$ only. None of them involves $k-k_a$, $N_a$, $N_0^a$, $j^+-k_0$, or the coupling.
\end{proof}

\begin{corollary}[The summation lemma on lifted sets of older pieces]\label{4.26:cor-summation-lifted}
Let $\mathfrak{a}$ be a live older piece of $\mathfrak{T}$, with region $W_{\mathfrak{a}}$, levels
$k_a$ and $j^+$, and lifted set $\mathsf{L}_{\mathfrak{a}}$. Let $\square_0\in\mathcal{C}$ meet
$\mathsf{L}_{\mathfrak{a}}$. Let $u=u(\square_0)$ and the multiplicity weight
$\mathsf{m}_F(\square_0)$ be those of Lemma~\ref{4.26:lem-lifted-counts}. Let the constants
$K_{\mathrm{SK}}(d)$, $A_\flat$, $E_0$, $B_0$, $C_0^{\sharp}$, the tile number $\mathsf{l}_j^2$, the
stratum level $m_0$ and the function $s(\cdot)$ be those of Lemma~\ref{4.26:lem-summation}.
\begin{enumerate}
\item[(i)] The leaf sum of clause (a) of Lemma~\ref{4.26:lem-summation} at $\square_0$ comprises the
leaves from $\mathbb{B}$, the frozen terms of the piece and the tails of item (Rd4) of
\eqref{4.26:eq-cover-lifted-uses-a}. It satisfies
\begin{equation*}
\text{leaf sum at }\square_0\;\le\;\mathsf{m}_F(\square_0)\,K_{\mathrm{SK}}(d)\,A_\flat\,
\bigl(E_0+B_0+1+C_0^{\sharp}\bigr);
\end{equation*}
thus $\mathsf{m}_F(\square_0)$ replaces $\mathsf{l}_j^2$ in the bound of that clause, and the frozen
terms of the piece add the summand $C_0^{\sharp}$.
Clause (b) of Lemma~\ref{4.26:lem-summation} holds at $\square_0$ as stated there. Let
$\square\in\mathcal{C}^1$ lie on a lifted set. Then two quantities at $\square$ are bounded as
stated, with $s(\square)=0$: the sum over the terms $T$ localised near $\square$ in clause (iv) of
the strengthened kernel bound at cover cubes, and the total of the near group.
\item[(ii)] One has
\begin{equation}\label{4.26:eq-summation-lifted-1}
(3L^{u})^d\,\mathsf{m}_F(\square_0)\,\vartheta^{\natural}(\square_0)\le\vartheta^{\natural}.
\end{equation}
Consequently \eqref{4.26:eq-summation-a} holds unchanged.
\item[(iii)] The anchor sums are bounded by $2^dV^{\mathrm{fr}}\vartheta^{\natural}$ times the
per-cube constants of Lemma~\ref{4.26:lem-summation}. Consequently \eqref{4.26:eq-summation-b}
holds as stated, with
\begin{equation}\label{4.26:eq-summation-lifted-2}
V:=\max\{V^{0},V^{A},V^{\mathrm{fr}}\}\le\bar{V}:=C_F\,C_V(d)\,(100L\check{R})^d .
\end{equation}
\end{enumerate}
\end{corollary}

\begin{proof}
\emph{Proof of \textup{(i)}.} We use the tiling statement for enlarged cover cubes by cubes of the
original cover, which says the following. An enlarged cube $\tilde{\square}'$ of level $i'$, on a
stratum of level $m_0$, is tiled by at most $(2L^{i'-m_0})^d$ cubes $\square''$ of the original
cover. Every term $T$ meeting $\tilde{\square}'$ meets some $\tilde{\square}''$. At each
$\square''$ clauses $(\beta)$ and $(\gamma)$ of the statement on per-cube bounds at a cube of the
original cover hold as stated.

We apply it with the graded cover $\mathcal{C}^1$ of \eqref{4.26:eq-graded-cover-a} in the role of
the original cover. The enlarged cube $\tilde{\square}_0$ is four sides wide, each side being at
most $M\ell_{j^+}$, the side of a cover cube at level $j^+$; we write $M\ell_i$ for the side of a
cover cube at level $i$. On $\tilde{\square}_0$ one has $\min_{\tilde{\square}_0}\mathfrak{c}\ge j^+-u-1$,
since $\tilde{\square}_0$ reaches at most one $\mathfrak{c}$-shell below that of $\square_0$. Hence
$\tilde{\square}_0$ is covered by at most
\begin{equation*}
\bigl(4L^{u+1}/2\bigr)^d=(2L^{u+1})^d
\end{equation*}
cubes $\square''\in\mathcal{C}^1$.

At each such $\square''$, clauses \textup{(S5)}$(\alpha)$ and \textup{(S5)}$(\beta)$ of the same
statement on per-cube bounds at a cube of the original cover hold as stated. Clause
\textup{(S5)}$(\alpha)$ concerns the tails: by item (Rd4) of \eqref{4.26:eq-cover-lifted-uses-a}, a
tail has $\operatorname{diam}X_T\ge 2M\ell_{\operatorname{scale}(\square'')}$. Clause
\textup{(S5)}$(\beta)$ concerns the shell widths: by Lemma~\ref{4.26:lem-original-level-strata}(iii),
every $\mathfrak{c}$-shell $i\ge h_a$ of $W_{\mathfrak{a}}$ is at least three layers of
$MR_{i+1}$-cubes of the $(i+1)$-lattice wide, hence at least $3R_{\min}$ sides of
$\pi_{i+1}$-cubes. Both clauses are statements
on thick shells of an admissible sequence, so they apply at $\square''$.

We turn to the frozen terms $\mathbb{C}^{(n')}_{k_a}(X)$ of the piece, born at stage $n'$. By
\cite[(1.63)]{B-CMP122a}, such a term is never finer than $m_0$ where it lies. Consider the short
frozen terms at $\square''$. By the cut in the proof of clause (b) of
Lemma~\ref{4.26:lem-summation}, they come from stages $j'\in[\mathfrak{c}-3,\,j^+-2]$. These are at
most $u+4$ stages, and each contributes at most $C_0^{\sharp}$.

Now consider the leaves born at levels $k'>k_a$. Their domains are made of cubes of the frozen
scale on $W_{\mathfrak{a}}$. They are therefore not finer than $\square_0$, and they are bounded as
stated in Lemma~\ref{4.26:lem-summation}.

Together, these give the bound of (i).

For clause (b), consider a chain term of the current run, born at stage $n'$. It meets the stage's
region $R^{(n')}\subset Z$. By
\cite[(1.63)]{B-CMP122a}, it is made of cubes of the frozen scale on $W_{\mathfrak{a}}$. Hence the
count of objects not finer than the cover's cubes applies to it as stated.

\emph{Proof of \textup{(ii)}.} Inequality \eqref{4.26:eq-summation-lifted-1} follows from item (F2)
of \eqref{4.26:eq-lifted-counts-a}, which reads
\begin{equation*}
(3L^{u+1})^{2d}(u+4)\,e^{-\frac12\delta D_F(\square)}\le 1 .
\end{equation*}
We use it together with the elementary inequality
\begin{equation*}
(3L^{u})^d\,(2L^{u+1})^d\le(3L^{u+1})^{2d}.
\end{equation*}

We now show that \eqref{4.26:eq-summation-a} is not moved. Consider the anchored summation step
over elements. There, every element $e$ with a block or a cube on a
lifted set carries $K_A\vartheta^{\natural}(\square_0)$ when read at $\square_0\in\mathcal{C}$,
respectively a factor proportional to $\vartheta^{\natural}(\square)$ when read at
$\square\in\mathcal{C}^1$. Two facts give this:
\begin{itemize}
\item the factor $v_e$ of the element vanishes if $\mathsf{b}(e)=\emptyset$;
\item if $\mathsf{Q}\subset Z$, the fluctuation field $\phi$ is read through the animal
$\mathsf{y}$ hung on the block of $e$ that lies on the lifted set.
\end{itemize}
A $\sigma_0$-cube of an animal, which has the scale of $\mathfrak{B}$, touches at most $4^d$ cubes
of $\mathcal{C}$ (the factor $4^d$ being contained in the constant $K_2$ of the anchored summation
step over elements) and at most $(3L^{u})^d$ cubes of $\mathcal{C}^1$. By
\eqref{4.26:eq-summation-lifted-1} we therefore have, with the per-cube sum on the right taken from
Lemma~\ref{4.26:lem-summation} with $\mathsf{l}_j\mapsto 1$,
\begin{equation*}
\begin{split}
[\text{multiplicity}]&\times[\text{per-cube sum}]\times\vartheta^{\natural}(\square)\\
&\le[\text{per-cube sum with }\mathsf{l}_j\mapsto 1]\times\vartheta^{\natural}.
\end{split}
\end{equation*}
Hence \eqref{4.26:eq-summation-a} holds as stated.

\emph{Proof of \textup{(iii)}.} In the anchored summation step over elements, the anchor at
$\square_0$ is estimated by
\begin{equation*}
\sum_{\square_0}\vartheta(\square_0)\le 2^dV\vartheta .
\end{equation*}
Taken with the weights $\mathsf{m}_F$, this estimate is item (F4) of
\eqref{4.26:eq-lifted-counts-a}. This gives the anchor sums of (iii) with the factor
$2^dV^{\mathrm{fr}}\vartheta^{\natural}$.

It remains to bound the members of $V$. We have
\begin{equation*}
V^{0}\le C_V(d)(100L\check{R})^d .
\end{equation*}
This is the bound of the cover-cube count \eqref{4.26:eq-cube-count-a}. It is used as stated, for
the count $V^0$ of cover cubes in the thick strata. The last clause of item (F4) supplies the
remaining members. This gives \eqref{4.26:eq-summation-lifted-2}, and with it
\eqref{4.26:eq-summation-b} as stated.
\end{proof}

\begin{remark}
Clause (i) is needed independently of the bounds on the near groups. Clause (a) of
Lemma~\ref{4.26:lem-summation} is stated for every cover cube $\square_0$, with
$\mathsf{l}_j^2:=1$ for $j\le k_0$, where $k_0$ is the base level of the present level $j$. Write
$k_0^a$ for the base level of the piece $\mathfrak{a}$, and consider the sets
$\mathbb{B}^{(m)}(X)$ with $k_0^a<m<j^+$ and
$d_m(X)$ bounded by a constant; a cube can carry as many as $(2L^{j-m})^d$ of them. At every stage
of the present site they lie on the lifted stratum of the older piece $\mathfrak{a}$, foreign to
the strata of that stage, and they are counted by the piece's own $N_0^a$.
At such a cube the tile number $\mathsf{l}_j^2$ of the present level, which is $1$ for $j\le k_0$,
does not count these sets: their count is governed by the strata of $\mathfrak{a}$, and it is this
count that clause (i) supplies through $\mathsf{m}_F(\square_0)$, the excess being absorbed by
item (F2) of \eqref{4.26:eq-lifted-counts-a}.

The tile number $\mathsf{l}_j^2$ plays the same role where a sum over neighbouring cover cubes
$\square'$ is bounded by the
supremum over $\square'$ of the per-cube bounds. There, a term $T$ at a lifted cube $\square'$ is
read, by the comparison of the distance $d(\tilde{\square},X_T)$ with $D_F(\square')$ on lifted
cubes, with
\begin{equation*}
\vartheta(\square)\,e^{-\delta d(\tilde{\square},X_T)}
\le\vartheta\,e^{-\frac12\delta D_F(\square')} ,
\end{equation*}
and this pays $\mathsf{m}_F(\square')$ by item (F2).
\end{remark}

\begin{definition}[Stage systems and their stability hypothesis]\label{4.26:def-stage-systems}
Let $l$ be a stage of the run. Throughout, $\mathfrak{T}=\Omega_k\cup Z^c$ is the region untouched by the stages, $Z^{\mathrm{on}}=Z\cap\Omega_k^c$, and $\mathfrak{N}^{(n)}(Y)$ is the regularity space of stage $n$ on $Y$, with its composite shell clause and the bound clause accompanying that clause, as in \eqref{4.26:eq-regularity-space-a}.

\emph{Systems.} A \emph{stage-$l$ system} $\mathfrak{D}_l$ assigns to every union $Y$ of cubes an open, $\mathcal{G}$-invariant set
\[
\mathfrak{D}_l(Y)\subset\mathfrak{N}^{(l)}(Y)
\]
of pairs $(\mathbb{U},\mathbb{J})$ on $Y$, subject to the restriction property
\begin{equation}\label{4.26:eq-stage-systems-a}
Y'\subset Y\quad\Longrightarrow\quad \mathfrak{D}_l(Y)\restriction Y'\subset\mathfrak{D}_l(Y').
\end{equation}

\emph{The norm relative to a system.} For a family $C$ of terms $C(X,\mathcal{S})$, the norm $\|C\|^{A,\mathfrak{D}}_{(l)}$ is obtained from the norm $\|C\|^{A}_{(l)}$ of the terms at stage $l$ by one change. In the formula for $\|C\|^{A}_{(l)}$ the size of each term enters as $\mathfrak{N}(C(X,\mathcal{S}))$; this size functional is a different object from the regularity space $\mathfrak{N}^{(n)}(Y)$. Every such size is replaced by
\[
\sup_{\mathfrak{D}_l(X)\times\mathfrak{D}_{\mathcal{S}}(X)}|C(X,\mathcal{S})|,
\]
where $\mathfrak{D}_{\mathcal{S}}$ is the slot domain.

\emph{Domains.} For a term $T\in\mathcal{T}$ we put $\operatorname{dom}(T):=\mathfrak{U}_T$. Here $\mathfrak{U}_T$ is the space of $T$ fixed by its clause set $\operatorname{Cl}(T)$, as in Definition~\ref{4.26:def-clauses-units}. For a term $\mathbb{C}^{(i)}(S,\mathcal{S}')$ of the representation we put $\operatorname{dom}(\mathbb{C}^{(i)}(S,\mathcal{S}')):=\mathfrak{D}_i(S)$.

\emph{The stability hypothesis.} A stage-$l$ system $\mathfrak{D}_l$ satisfies hypothesis $(\mathrm{H}^{\natural})_l$ if the following holds. Take every union $Y$ of cubes, every $(\mathbb{U},\mathbb{J})\in\mathfrak{D}_l(Y)$, and all pairs $(b,\sigma)$ in the range fixed for the stage, that is, the range over which the proof family $\mathfrak{W}^{(l),\natural}(Y)$ at $(\mathbb{U},\mathbb{J})$ on $Y$ of Definition~\ref{4.26:def-proof-family} is formed. Then:
\begin{enumerate}
\item[(a)] every $\Phi$ of family (a) of Definition~\ref{4.26:def-proof-family} lies in $\mathfrak{N}^{(l-1)}(Y)$, and $\Phi\restriction S\in\operatorname{dom}(c)$ for every $M$-class $c$ with $S\subset Y$;
\item[(b)] the following holds for every $T\in\mathcal{T}_1$ with $X_T\subset Y$ and $X_T\not\subset\tilde{\square}^2$. For all parameters $(t,\tau)$ in the range fixed for family (b) in Definition~\ref{4.26:def-proof-family}, one has $\mathcal{W}_T(t,\tau)\restriction X_T\in\mathfrak{U}_T$;
\item[(c)] the following holds for every near pair $(T,\square)$, that is, with $X_T\subset\tilde{\square}^2$, for which $\square_0=\tilde{\square}^5\subset Y$. The input hypotheses of \cite[Lemma~4]{B-CMP109} hold on $\square_0$ in the form described below, and the families (c) and (c$'$) of Definition~\ref{4.26:def-proof-family} are contained in $\mathfrak{U}_T$. In \cite{B-CMP109} these input hypotheses read $U\in U^c_{k+1}(\square_0,(1+2\beta)\alpha_0,(1+2\beta)\alpha_1)$. That space is defined by the same conditions (i)--(iv), the only difference being that the configurations $U,J$ are defined and satisfy (i)--(iii) on the whole lattice. At the scale of the shell the corresponding conditions are those of \cite{B-CMP119}. The form required here is:
\begin{itemize}
\item conditions (i)--(iii), the direct entries, hold with factor at most $1+2\beta$ on $\square_0\cap\mathfrak{T}$, and hold by the lettered clauses $\mathfrak{L}_l$ on $\square_0\cap Z^{\mathrm{on}}$;
\item in place of condition (iv), on $\tilde{\square}^3\cap Z^{\mathrm{on}}$ the $(\Sigma^{\ast})$-triples of the lettered clauses hold, in the form fixed by the standing assumption on the derived-entry family on $Z^{\mathrm{on}}$;
\item in place of condition (iv), on $\tilde{\square}^3\cap\mathfrak{T}$ the composite shell clause of \eqref{4.26:eq-regularity-space-a} holds with $X'=\square_0$. This means that \cite[(1.16)]{B-CMP109} holds at every $p\le p_{\square}$ with factor $<1+2\beta$, measured in the unit of the shell clause, and that the bound clause accompanying the composite shell clause in \eqref{4.26:eq-regularity-space-a} holds literally;
\item at every location $y\in X_T\cap\mathfrak{T}$ of a clause of $T$, moreover, the box condition \eqref{4.26:eq-inherited-domain-a} holds at the box-level triple $(p_{\square},\square_0,\mathbb{B}^{(l)}_k)$; this is the diminished threshold of $T$, the threshold with which the near-pair analysis of $T$ begins.
\end{itemize}
\end{enumerate}
In formulas, the membership relations asserted by $(\mathrm{H}^{\natural})_l$ are
\begin{equation}\label{4.26:eq-stage-systems-b}
\begin{gathered}
\text{(a)}\quad \Phi\in\mathfrak{N}^{(l-1)}(Y),\qquad \Phi\restriction S\in\operatorname{dom}(c)\quad(S\subset Y);\\
\text{(b)}\quad \mathcal{W}_T(t,\tau)\restriction X_T\in\mathfrak{U}_T
\quad(T\in\mathcal{T}_1,\ X_T\subset Y,\ X_T\not\subset\tilde{\square}^2);\\
\text{(c)}\quad \text{family }(\mathrm{c})\cup\text{family }(\mathrm{c}')\subset\mathfrak{U}_T
\quad(X_T\subset\tilde{\square}^2,\ \square_0=\tilde{\square}^5\subset Y),
\end{gathered}
\end{equation}
together with the conditions on $\square_0$ and $\tilde{\square}^3$ listed in (c).

On a live older piece of $\mathfrak{T}$, two readings apply in (c). First, the unit of the shell clause is the frozen weighted unit of the shell. This unit implies the unweighted unit of the shell built on $\mathfrak{c}$, by the unit conversion \eqref{4.26:eq-graded-box-geometry-a}. Second, the box-level triple is $(p_{\square},\square_0,\mathfrak{B}^{\mathfrak{c}}\restriction\square_0)$ in place of $(p_{\square},\square_0,\mathbb{B}^{(l)}_k)$.
\end{definition}

\begin{remark}\label{4.26:lem-stage-expansion}
The statement of this lemma is not written out here; parts of its proof are printed below.
\end{remark}

\begin{lemma}[The stage expansion: small terms away from their cube]
\label{4.26:prf-expansion-small-terms}
Assume the hypotheses of Lemma~\ref{4.26:lem-stage-expansion} at a stage $l$, $1\le l\le N$, and
use the radii, the spend $\mathsf{w}_0$ and the window condition of
Definition~\ref{4.26:def-radii-window}. In the expansion of that lemma, a small term
$T\in\mathcal{T}_1$ enters through elements $e$, each with a domain $\mathsf{d}(e)$, a set
$\mathsf{b}(e)$ of read fluctuation cubes and a weight $\nu(e)$. These are given by
\eqref{4.26:eq-expansion-small-terms-1} and \eqref{4.26:eq-expansion-small-terms-2} for the
pairs $(T,\square)$ with $X_T\not\subset\tilde{\square}^2$; for the near groups the domains and
the count of read blocks are given by \eqref{4.26:eq-expansion-small-terms-7}. For the pairs with
$X_T\not\subset\tilde{\square}^2$ the sums over elements at every anchor are bounded by a
constant depending on $L$ only times $K''\vartheta^{\natural}(E_0+1)$; this bound lies inside the
total $(\Sigma\text{-}\mu)$ of the dictionary of Lemma~\ref{4.26:lem-stage-expansion}. For the
near groups the total per cover cube is bounded by \eqref{4.26:eq-expansion-small-terms-8}. The
pieces of clause (S4) of the expansion theorem carry no radius.
\end{lemma}

\begin{proof}
Throughout, $\mathsf{Q}$ is the set of fluctuation cubes and $\square$ runs over the cubes of
the cover. The set $\tilde{X}_T^4$ is $X_T$ enlarged by four layers of cubes, which is the same
margin that replaces $\tilde{\square}$ by $\tilde{\square}^4$ in the cluster expansion of the
small terms. The distances $d_\ell(X_T,\square^{\zeta})$ and $d^-_{\square}$ are those of the
radii in \eqref{4.26:eq-radii-window-a}. The rates $a_k$, $a_m$ are those of clause~(c) of
Lemma~\ref{4.26:lem-stage-expansion}, $d_m$ denotes the hull length at the scale $m$, so that
$e^{a_md_m}$ is the hull weight; here and below $m=m(y)$ is the level that
Lemma~\ref{4.26:lem-family-consumption} attaches to a location $y\in Y$. The constant
$K_\flat'$ is the constant of the dictionary of Lemma~\ref{4.26:lem-stage-expansion}. The
remaining letters are
constants of the cluster expansion of the small terms and of the bookkeeping, with the
following roles: $K''$ is the constant of clause~(v) of Lemma~\ref{4.26:lem-family-consumption}
entering the weights below; $K_S$ and $E_0$ are the constants of the bound on the sum over small
terms at a fixed cube; $c_J$ and $c_\flat$ are the per-cube charges of the bookkeeping, and $c_F$
is the charge per read cube of the activity bound; $C_d$ depends on $d$ only; $w$ is the lower
bound on relative
distances of cubes at different levels (it is not the twist parameter of
Definition~\ref{4.4:def-twisted-argument}); $v_e$ is the activity of an element $e$;
$L_\circ$, $s(\square)$, $c_\alpha$, $C_2$, $C_R^{\natural}$, $B_3$ and the power $M^q$ are the
constants and the scale index of the near-range bound of the cluster expansion of the small
terms, $C_R^{\natural}$ and $M^q$ being fixed before $\gamma$; $\mathbb{P}$ and $W$ name the two
types of pieces of clause~(S4) of the expansion theorem. The bookkeeping proposition of this
section reads, for each element $e$, the triple $(\mathsf{b}(e),\mathsf{d}(e),\nu(e))$. We treat
first the pairs that are not near and then the near groups.

\emph{The elements of a pair that is not near.} Let $T\in\mathcal{T}_1$ and let $\square$ be a
cover cube with $X_T\not\subset\tilde{\square}^2$. The elements attached to $(T,\square)$ are the
triples $e=(T,\square,\mathsf{y})$, where $\mathsf{y}$ is a set of cubes (an animal). Their
domains and read blocks are
\begin{equation}\label{4.26:eq-expansion-small-terms-1}
\begin{gathered}
\mathsf{d}(e):=\tilde{X}_T^4\cup\tilde{\square}^4\cup\mathsf{y}\quad\text{(required to be
connected)},\\
\mathsf{b}(e):=\bigl\{q\in\mathsf{Q}:\ q\cap\bigl(\tilde{X}_T^4\cup\tilde{\square}^4\cup
\mathsf{y}\bigr)\neq\emptyset\bigr\},\qquad |\mathsf{b}(e)|\le 2(16L)^d+|\mathsf{y}|,
\end{gathered}
\end{equation}
and their weights are
\begin{equation}\label{4.26:eq-expansion-small-terms-2}
\nu(e):=\|T\|_\infty K''\vartheta^{\natural}e^{-\frac12\delta d^-_{\square}}\,
e^{-\frac14\delta d_\ell(X_T,\square^{\zeta})}\,e^{-(\kappa_1-1)|\mathsf{y}|}.
\end{equation}
By \eqref{4.26:eq-radii-window-a}, $\vartheta^{\natural}e^{-\frac12\delta d^-_{\square}}
=\vartheta^{\natural}(\square)$.

\emph{The domain.} The domain $\mathsf{d}(e)$ has to contain the cubes of $\mathsf{b}(e)$, for
the following reason. The bookkeeping proposition forms the localisation domain of a polymer
$\gamma$ as $X:=[\bigcup_o\mathsf{d}(o)]_m$, the union of the domains of its elements $o$
rounded at the scale $m$. The activity $z(\gamma)$, in the form given by the cluster expansion
of the small terms, reads the pair $(\mathbb{U},\mathbb{J})$ on a set $Z_\gamma\supset
\mathsf{b}(e)$. If the domain did not contain these cubes, the locality in $X$ required in
clause~(b) of Lemma~\ref{4.26:lem-stage-expansion} would fail.

Let $\square'$ be the cover cube nearest to $X_T$. Then $X_T\subset\tilde{\square}'^2$ and
$\tilde{X}_T^4\subset\tilde{\square}'^6$. Including $\tilde{X}_T^4$ therefore enlarges the hull
by at most $16^d$ cubes. This costs at most $e^{(a_k+1+c_J)16^d}$, which is absorbed into
$K_\flat'$.

Hypothesis~(b) of $(\mathrm{H}^{\natural})_l$ in \eqref{4.26:eq-stage-systems-b} is used with
its scope $X_T\subset Y$ exactly as stated there. Here it is applied on sets
$Y\supseteq\mathsf{d}(e)$. Since $X_T\subset\tilde{X}_T^4\subset\mathsf{d}(e)\subseteq Y$, each
such application is a case of that scope.

The mechanism that makes such elements affordable is the one used for far pieces in the cluster
expansions of \cite{B-CMP116}. For a piece of this kind the animal $\mathsf{y}$ joins
$\tilde{\square}$ to $X_T$, and each cube $\Delta\subset\mathsf{y}$ pays a factor
$e^{-(\kappa_1-1)}$.

\emph{The read blocks.} The bound on $|\mathsf{b}(e)|$ in
\eqref{4.26:eq-expansion-small-terms-1} is obtained as follows. Because
$X_T\subset\tilde{\square}'^2$, the fluctuation cubes meeting $\tilde{X}_T^4$ number at most
$(16L)^d$. The same count $(16L)^d$ bounds those meeting $\tilde{\square}^4$. The cubes met by
$\mathsf{y}$ contribute $|\mathsf{y}|$.

The cubes meeting $\tilde{X}_T^4$ must be read, because of the form of the derivative piece,
written in the notation of \cite[(3.4)]{B-CMP109} with $\mathbb{H}$ the stage layer:
\begin{equation}\label{4.26:eq-expansion-small-terms-3}
\int_0^1dt\,\partial_{t_\square}\big|_{t_\square=0}\,
T\Bigl(X_T;\,U_\ell\bigl(\exp iQ_\ell(\eta(t+t_\square\zeta_\square)\mathbb{H})\,
\bar{U}^{\ell}\bigr)\Bigr).
\end{equation}
In this piece the parameter $t$ is one global parameter, as in \cite[(3.4)]{B-CMP109}. In the
cluster expansion of the small terms, the piece attached to $\square$ of a term on a set $X$
reads $t\mathbf{H}_k(B')$, in the notation of \cite{B-CMP109}, on all of $X$. Consequently the
activity $v_e$ reads the fluctuation
field $\phi$ on every fluctuation cube meeting the neighbourhood of $X_T$. The animal
$\mathsf{y}$ touches that neighbourhood but need not cover it.

A read-block set that omitted the cubes at $X_T$ would fail on two counts. The conditioning
identity of the bookkeeping proposition would be false. The charge $c_F$ levied per read cube
would also be left unpaid on $X_T$. The cubes of $\mathsf{b}(e)$ not met by $\mathsf{y}$ cost at
most $e^{2c_\flat(16L)^d}$, and this factor is absorbed into $K_\flat'$.

\emph{The weight.} The bound \eqref{4.26:eq-expansion-small-terms-2} follows from three
ingredients:
\begin{itemize}
\item clause~(v) of Lemma~\ref{4.26:lem-family-consumption};
\item hypothesis~(b) of $(\mathrm{H}^{\natural})_l$ in \eqref{4.26:eq-stage-systems-b}, applied
on $Y\supseteq\mathsf{d}(e)$ as explained above;
\item the Cauchy estimate in the decoupling parameters $\sigma(\Delta)$ on the circles
$|\sigma(\Delta)|=e^{\kappa_1}$.
\end{itemize}

\emph{Three mechanisms, three costs.} Three mechanisms pay three distinct costs.

First, the hull weight $e^{a_md_m}$ over the gap between $\tilde{\square}$ and $X_T$ is paid
cube by cube of $\mathsf{y}$, by the factor $e^{-(\kappa_1-1)}$. It is not paid by
$e^{-\frac14\delta d_\ell(X_T,\square^{\zeta})}$. Indeed, the floor on the separation $w$ among
the floors assumed in Lemma~\ref{4.26:lem-family-consumption} gives only
$\frac14\delta w\ge 4\ln L+1$, and $4\ln L+1<a_k$.

Second, fix an anchor $q\in\mathsf{Q}$ meeting $\tilde{\square}^4\cup\mathsf{y}$. The sum over
$T$ at fixed $\square$ is the sum already bounded for the small terms at a fixed cube:
\begin{equation}\label{4.26:eq-expansion-small-terms-4}
\sum_T\|T\|_\infty\,e^{-\frac14\delta d_\ell(X_T,\square^{\zeta})}\le K_S(E_0+1).
\end{equation}
This bound applies without change because the radius (r-1) of
\eqref{4.26:eq-radii-window-a}, namely
$r_T(\square)=e^{\frac14\delta d_\ell(X_T,\square^{\zeta})}/(14K^+\vartheta(\square))$, carries
the factor $e^{\frac14\delta d_\ell(X_T,\square^{\zeta})}$. Hence the weight
\eqref{4.26:eq-expansion-small-terms-2} keeps the decaying factor
$e^{-\frac14\delta d_\ell(X_T,\square^{\zeta})}$.

Third, fix an anchor $q$ meeting $\tilde{X}_T^4$, that is, one of the cubes added to
$\mathsf{b}(e)$ above. Write $\tilde{\nu}(e)$ for $\nu(e)$ multiplied by the per-cube charges
$e^{c_J}$ and $e^{c_\flat}$ of the bookkeeping on the cubes of $\mathsf{y}$.
The sum to be bounded is $\sum_T\sum_{(\square,\mathsf{y})}\tilde{\nu}(e)$.

At fixed $T$, the pairs $(\square,\mathsf{y})$ are summed through the animals $\mathsf{y}$ that
join $\tilde{\square}$ to $X_T$. This gives
\begin{equation}\label{4.26:eq-expansion-small-terms-5}
\sum_{(\square,\mathsf{y})}e^{-(\kappa_1-1-c_J-c_\flat)|\mathsf{y}|}
\le C_{\mathrm{an}},
\end{equation}
where $C_{\mathrm{an}}$ is an animal constant for animals rooted at $\tilde{X}_T^4$. It is
obtained with the factor $e^{-4}$ per cube of the cluster expansion of the small terms. The sum
includes the at most $16^d$ blocks for which $\mathsf{y}=\emptyset$.

The sum over $T$ at fixed $q$ is paid by the series of the second case, uniformly in $\square$.
Suppose $X_T\cap\tilde{\square}=\emptyset$. Then $d_\ell(X_T,\square^{\zeta})\ge
wL^{(i-\ell)^+}$, where $\ell$ is the level of $T$ and $i$ is the level at which $T$ is counted,
namely that of $q$. Therefore
\begin{equation}\label{4.26:eq-expansion-small-terms-6}
\sum_{\ell\le i}C_d\bigl(4L^{i-\ell}\bigr)^dE_0\,e^{-\frac14\delta wL^{i-\ell}}
\le K_S(E_0+1).
\end{equation}
The series \eqref{4.26:eq-expansion-small-terms-6} is run at the stratum of the contour joining
$\tilde{\square}$ to $X_T$. When $i$ is smaller by at least two than the level of that stratum,
the contour crosses the stratum one level below it entirely. The series is then run at that
lower stratum, at the cost of a factor
$L^{2d}$. This is a number depending on $L$ only, and it is absorbed into the constant depending
on $L$ only that multiplies $K''$ below.

Suppose instead that $X_T\cap\tilde{\square}\neq\emptyset$ and $X_T\not\subset\tilde{\square}^2$.
Then the pair belongs to the class treated in clause~(S5)($\alpha$) of the expansion theorem
whose clauses Lemma~\ref{4.26:lem-stage-expansion} establishes at stage $l$.

Collecting the three contributions, the total is at most a constant depending on $L$ only times
$K''\vartheta^{\natural}(E_0+1)$. This lies
inside the corresponding member of the total $(\Sigma\text{-}\mu)$.

\emph{Near groups.} For the near groups, that is, the pairs with $X_T\subset\tilde{\square}^2$,
which clause (S3) of the expansion theorem treats under the dictionary of
Lemma~\ref{4.26:lem-stage-expansion}, we put
\begin{equation}\label{4.26:eq-expansion-small-terms-7}
\begin{gathered}
\mathsf{d}:=\tilde{\square}^5\cup\mathsf{y},\qquad |\mathsf{b}|\le(12L)^d+|\mathsf{y}|,\\
\text{on lifted sets:}\quad\mathsf{d}:=[\tilde{\square}^5]_{\mathfrak{B}}\cup\mathsf{y}.
\end{gathered}
\end{equation}
On a lifted set, the hull of $\tilde{\square}^5$ at the frozen scale has at most $16^d$ cubes,
by (Rd5) of \eqref{4.26:eq-cover-lifted-uses-a}. The cluster expansion of the small terms
places the near pieces on $\tilde{\square}^4\cup\mathsf{y}$, and the box is
$\square_0=\tilde{\square}^5$. This gives the count $(12L)^d$ and a cost
$e^{c_\flat(12L)^d}$, which is absorbed into $K_\flat'$.

Ba{\l}aban's radius $\alpha_2$ is kept without change. So is the statement of
\cite{B-CMP109} that $H_j(B(t))$, with Ba{\l}aban's index $j$ of \cite{B-CMP109}, satisfies
\cite[(3.14)]{B-CMP109} with $\frac13\alpha_2$.

The circle in $t_\square$ has radius $\mathsf{w}_0r^{\mathrm{nr}}_\square$, with $r^{\mathrm{nr}}_\square$ from
(r-n) of \eqref{4.26:eq-radii-window-a}. For the remainder pieces its radius is that of (r-rem)
of \eqref{4.26:eq-radii-window-a}. This circle is contained in, and has $\mathsf{w}_0$ times the
radius of, the $t_\square$-circle on which the near-range bound of the cluster expansion of the
small terms is stated.

Analyticity in $t_\square$ is used once, by one Cauchy estimate, as in \cite{B-CMP116}. The
bound therefore degrades by the single factor $1/\mathsf{w}_0$. The outer prefactor
$O(1)\alpha_2^{-1}\varepsilon_1$, which
\cite{B-CMP109} obtains from this differentiation, becomes proportional to
$\vartheta^{\natural}(\square)$.

\cite[Lemma~4]{B-CMP109} is applied at $\alpha_3(T)$, the minimum in
\eqref{4.26:eq-radii-window-c}. The bound \cite[(3.54)]{B-CMP109} is applied at
$r=\alpha_3(T)/\sup|B|$. The bound \cite[(3.35)]{B-CMP109} keeps its factor
$(s_\ell/s_i)^5$, in Ba{\l}aban's notation there, now with the constant of
Definition~\ref{4.26:def-radii-window}. The resulting
factors $M^5$, $(L^2)^5$, $\mathsf{a}^5$ and $B_3^{10}$ join the constants $M^q$ and
$C_R^{\natural}$, which are fixed before $\gamma$.

\emph{The hypotheses and the conclusion of \cite[Lemma~4]{B-CMP109}.} The hypotheses of
\cite[Lemma~4]{B-CMP109} are supplied in the way stated in the third hypothesis of
$(\mathrm{H}^{\natural})_l$ in \eqref{4.26:eq-stage-systems-b}. Its conclusion, that the union of
the families labelled (c) and (c$'$) in \cite[Lemma~4]{B-CMP109} is contained in
$\mathfrak{U}_T$, is obtained in one of two ways.
\begin{itemize}
\item[($\alpha$)] At the clauses of $T$ on $Z^{\mathrm{on}}$, the conclusion is taken from the
hypothesis standing in clause~(S3) of the expansion theorem.
\item[($\beta$)] At every clause of $T$ on $\mathfrak{T}$, the conclusion is taken from the two
propositions on the region $\mathfrak{T}$, the second of them being the frozen variant, on the
cover $\mathcal{C}^1$. These propositions re-run \cite[(3.37)--(3.52)]{B-CMP109}, starting from
the box condition at the entry of the box level. The re-run is valid at every level
$\ell_T\le m(y)$, where $\ell_T$ is the level of the term $T$ (the index $\ell$ of
\eqref{4.26:eq-expansion-small-terms-6}), and involves no factor $L^{-2}$.
\end{itemize}
There is a second derivation in part of the range of ($\beta$). Where $\ell_T<m(y)$ for every
$y\in X_T$ and the box lies off the raised zone, \cite[Lemma~4]{B-CMP109} itself yields the same
conclusion, because its four hypotheses hold literally there. Among them are the shell form of
the box condition, and the condition that Ba{\l}aban's index $j$ of \cite{B-CMP109} satisfies
$j\le k$, hence $L^j\eta\le1$
and $(1+7\beta)L^{-2}\le1$. The last inequality holds because $(1+8\beta)L^{-2}<1$ for the
values of $\beta$ and of the blocking factor $L$ admitted in this chapter.
This second derivation agrees with ($\beta$) and is not needed for it.

The identities and operator bounds of Ba{\l}aban that are used inside the propositions on
$\mathfrak{T}$ are applied by statement, at the scale of the shell, in the form stated in
\cite{B-CMP119}. The same is done throughout this part of the proof.

\emph{The total per cube.} For the near groups, the total per cover cube is at most
\begin{equation}\label{4.26:eq-expansion-small-terms-8}
L_\circ^{6s(\square)}\,c_\alpha C_R^{\natural}(E_0+1)\,e^{C_2\kappa_1}\,
\vartheta^{\natural}(\square)\,\bigl(1+\text{remainders}\bigr),
\end{equation}
where the remainders are the contributions of the remainder pieces, terms bounded
by $K''e^{-\frac14\delta w}$. This is the per-cube member of the near-range total, and its sum
over the cover cubes is the near-range sum entering the totals of the dictionary of
Lemma~\ref{4.26:lem-stage-expansion}.

\emph{The pieces of clause (S4).} The pieces of types $\mathbb{P}$ and $W$ in clause~(S4) of the
expansion theorem need no radius. Such pieces read only the following objects:
\begin{itemize}
\item the fluctuation field $\phi$ on bounded supports;
\item the input pair $(\mathbb{U},\mathbb{J})$ through kernels holomorphic on $\mathfrak{K}_P$,
as in \cite{B-CMP116};
\item Wilson plaquettes of $\Phi\in\mathfrak{N}^{(l-1)}$.
\end{itemize}
This completes the treatment of the small terms away from their cube and of the near groups in
the proof of Lemma~\ref{4.26:lem-stage-expansion}.
\end{proof}

\begin{proof}[Proof of Lemma~\ref{4.26:lem-stage-expansion}, continued]
\label{4.26:prf-expansion-leaves}
This part of the proof completes the transfer to stage $l$ of the expansion theorem whose clauses
(a)--(d) Lemma~\ref{4.26:lem-stage-expansion} asserts, under the dictionary fixed in that lemma.
It treats, in order: the size bound on the boundary terms and the chain estimate; the polymer
proposition; the majorant lemma; the Gaussian term $\mathbb{C}_G$ and the correction $\Delta W$;
the totals; and gauge invariance together with clause (d).

\emph{Boundary terms.} The size bound
\begin{equation}\label{4.26:eq-expansion-leaves-1}
\sum_{\mathcal{S}} e^{\kappa_A|\mathcal{S}|}\,\mathfrak{N}\bigl(T(X_T,\mathcal{S})\bigr)
\le c_T^{\mathrm{sz}}\,e^{-\kappa_T d_\ell(X_T)}
\end{equation}
is a supremum over $\mathfrak{U}_T\times\mathfrak{D}_{\mathcal{S}}$, where
$\mathfrak{N}\bigl(T(X_T,\mathcal{S})\bigr)$ denotes the size of the term $T(X_T,\mathcal{S})$ as in
the size hypothesis of Lemma~\ref{4.26:lem-stage-expansion} on the terms $T(X_T,\mathcal{S})$ (a
functional of the
term, not the fat regularity space $\mathfrak{N}$ below); it is unchanged by the
transfer. The bound \eqref{4.26:eq-expansion-leaves-1} is the content of that size hypothesis
of Lemma~\ref{4.26:lem-stage-expansion}, in the plain length $d_\ell$,
for every term of Ba{\l}aban's boundary-term sets $\mathbb{B}_k$. For the terms of the first group
created by the operation $\mathbf{R}$, namely the new boundary terms, whose localization domains
$X$ intersect the large-field region, \cite[(1.99)]{B-CMP122b} states the bound in a relative
length. The replacement of the linear size of the domain $X$ in \cite[(1.99)]{B-CMP122b} by
$d_k(X)$ is stated in \cite{B-CMP122b} for the terms of the second group.

Let $\mathcal{K}_{\mathsf{d}}$ be the least class of terms that contains every term defined over a
strip and that contains every term built from terms of the class. For $T\in\mathcal{K}_{\mathsf{d}}$
the length $d_\ell(X_T)$ in \eqref{4.26:eq-expansion-leaves-1} is replaced by the relative length
$\mathsf{d}_\ell(X_T)$ of the relative bound for boundary terms, which satisfies
$\mathsf{d}_\ell\le d_\ell$. The bound \eqref{4.26:eq-expansion-leaves-1} with
$\mathsf{d}_\ell(X_T)$ in place of $d_\ell(X_T)$ is not derived here. It is taken by statement from
the relative bound for boundary terms. For terms outside $\mathcal{K}_{\mathsf{d}}$ nothing
changes. The surplus cases, and the constants $A_\flat$ and $\sigma_\ast$ of the estimate on these
terms in the transferred proof ($\sigma_\ast$ being the surplus rate defined in the Gaussian part
below), are
unchanged. The sum over cover cubes is part (a) of Lemma~\ref{4.26:lem-summation}.

\emph{The chain estimate.} The chain estimate of the polymer proposition, read in the norm
$\|\cdot\|^{A,\mathfrak{D}}$, is supplied by part (b) of Lemma~\ref{4.26:lem-summation}.

\emph{Polymers.} We follow the proof of the polymer proposition.

In the first step, the identities hold at each background pair $(\mathbb{U},\mathbb{J})$, with the
enlarged set of read blocks $\mathsf{b}(e)$.

In the second step, the exponent $v_e$ of the element $e$ in the polymer proposition satisfies
$|e^{v_e}-1|\le 2\nu^{\natural}(e)$, with $\nu^{\natural}(e)$ as in
the corresponding display, on the real support of the cut-off functions $\chi$ of the
expansion. The tolerance of the second
step on the fat regularity space $\mathfrak{N}$, which is contained in the polydisc
$\mathfrak{K}_P$, is obtained as follows. Two ingredients enter it.
\begin{itemize}
\item The first ingredient is the set of bounds on the operator
$\Delta^{(j)}(\sigma;\mathbb{U},\mathbb{J})$ of the background pair $(\mathbb{U},\mathbb{J})$;
$K_\Delta$ denotes the constant in these bounds. These bounds are:
\begin{itemize}
\item[($\Delta$-1)] convergence of the expansion defining $\Delta^{(j)}(\sigma;\mathbb{U},\mathbb{J})$
on the product of
$\mathfrak{K}_P$ with $\{|\sigma|\le e^{\kappa_1}\}$, under $K_\Delta\gamma_0\le\frac12$, where here
and below $\gamma_0=\gamma_0^{(k)}$ is the size of the polydisc $\mathfrak{K}_P$;
\item[($\Delta$-2)] the bound displayed below;
\item[($\Delta$-3)] two increments: the bound
\begin{equation*}
\|\Delta^{(j)}(\mathbb{U})-\Delta^{(j)}(1)\|\le O(1)(BCM\tilde{\alpha}_0+\tilde{\alpha}_1)
\end{equation*}
of the walk lemma, with Ba{\l}aban's constants $B$, $C$, $M$ of the clause member
$BCM\alpha_{0,n}$ of \cite[(2.34)--(2.39)]{B-CMP119}, is increased by $K_\Delta\gamma_0$, and the
constant $\alpha_5$ of the second step is increased by $K_\Delta\gamma_0$.
\end{itemize}
The bound ($\Delta$-2) reads
\begin{equation}\label{4.26:eq-expansion-leaves-2}
\bigl|[\Delta^{(j)}(\sigma;\mathbb{U},\mathbb{J})-\Delta^{(j)}(\sigma;\mathbb{U},0)](b,b')\bigr|
\le K_\Delta\gamma_0\,e^{-\delta d(b,b')}.
\end{equation}
\item The second ingredient is the bound on $\alpha_5$. The inherited-domain theorem bounds
$\alpha_5$ by a member $O(5w_\ast(\tilde{\alpha}_0+\tilde{\alpha}_1))$, to which ($\Delta$-3) adds
$K_\Delta\gamma_0$; since each product $w_\ast\tilde{\alpha}_i$ is at most $\gamma_0^{(k)}/5$ (see
below), this yields
\begin{equation}\label{4.26:eq-expansion-leaves-3}
\alpha_5\le\bigl(2\,O(1)+K_\Delta\bigr)\gamma_0^{(k)}
\le\bigl(2\,O(1)+K_\Delta\bigr)\bar{\gamma}_0(\mathsf{t}_k).
\end{equation}
Indeed, every product of a weight with one of the radii $\tilde{\alpha}_i$ is at most
$\gamma_0^{(k)}/5$ by the definition \eqref{4.26:eq-regularity-radius-a} of $\gamma_0^{(k)}$. This
includes the products with $\omega\le w_\ast(k_a)$ on live pieces, and it is a smallness condition.
The second inequality is the run-free majorant of Definition~\ref{4.26:def-regularity-radius}.
\end{itemize}

The bounds ($\Delta$-1)--($\Delta$-3) involve no gauge fixing on individual cubes and do not use the
holonomy bound for closed chains of cubes introduced in the Gaussian part below.

The member $O(5w_\ast(\tilde{\alpha}_0+\tilde{\alpha}_1))$ of $\alpha_5$ carries no
factor $BCM$, the factor of the clause of the definition of Ba{\l}aban's space that carries the
member $BCM\alpha_{0,n}$ \cite[(2.34)--(2.39)]{B-CMP119}, in contrast with the member
$BCM\tilde{\alpha}_0+\tilde{\alpha}_1$ of
the walk lemma. It is therefore estimated from the gauge-invariant clause bounds on each cube
separately. Ba{\l}aban's share
of $\alpha_5$ in \cite[(2.16)--(2.17)]{B-CMP116} keeps the $G$-valued $U$ as background and perturbs
only $(A',\mathbb{J})$, where $A'$ is the field entering the clauses
\cite[(1.64), (1.66), (1.68)]{B-CMP122a}; no cube-wise gauge and no holonomy enter it. Hence the
second step does not
use the holonomy bound, and no estimate of $\alpha_5$ through cube-wise gauges on a support of more
than one clause cube is used.

The constant $c_F$ of the polymer proposition is unchanged. The remaining four steps of the polymer
proposition use only the
data $(\mathsf{b},\mathsf{d},\tilde{\nu})$ of the elements: the read blocks $\mathsf{b}(e)$, the
domains $\mathsf{d}(e)$ (unrelated to the relative length $\mathsf{d}_\ell$) and the numbers
$\tilde{\nu}$ attached to them.

For the analyticity clause, the series converge uniformly on $\mathfrak{D}_l(X)$ under majorants
that do not depend on $(\mathbb{U},\mathbb{J})$. Holomorphy follows from the Weierstrass theorem.
Locality in $X$ and the condition $X\cap R^{(l)}\neq\emptyset$, where $R^{(l)}$ is the region of
stage $l$, hold term by term in the representation of the polymer proposition and pass to the
uniform limit.

\emph{Majorants.} In the majorant lemma and its corollary the majorants are independent of
the parameter $\lambda$ of that lemma and of $(\mathbb{U},\mathbb{J})$ on $\mathfrak{D}_l$. For a
collection $\mathcal{E}$ of elements $e$, with $c(e)$ the label attached to the element $e$, we put
\begin{equation*}
\mathcal{S}_{j'}(\mathcal{E}):=\bigcup_{e\in\mathcal{E}}\mathcal{S}_{j'}(c(e)).
\end{equation*}
The weights are submultiplicative, and the $A$-polydisc never shrinks. The clauses (r1)--(r3) of the
flat form of the majorant lemma do not depend on the radii or on the domains.

\emph{The Gaussian part.} The walk lemma uses no stored term, no variable $A$ and no shift. It
expands in walks on $\mathfrak{N}$. Its inputs are:
\begin{itemize}
\item the random-walk expansion of \cite[Theorem~3.10]{B-CMP99b};
\item the bounds ($\Delta$-1)--($\Delta$-3);
\item the holonomy bound for closed chains of cubes, applied to the background pair in
$\mathfrak{K}_P$ taken on the whole domain of $\mathfrak{B}$, not cube by cube.
\end{itemize}
The holonomy bound is proved as an implication, in three cases according to the closed chain of
cubes: chains that do not wind; chains that wind around the torus, for which it rests on the
hypotheses (H3) and (H5) of Notation~\ref{2.2:not-hypotheses} and on the setting; and chains that
encircle a hole of the domain, for which it rests on statements of Ba{\l}aban's papers used as
stated there (the size hypothesis of one of them being met by \cite[p.~177, (ii)]{B-CMP122a}) and
on the condition $d\ge3$, together with the further antecedents listed in the statement of the
holonomy bound. It bounds
the contribution of a walk whose closed chain of cubes has length $p$ by $(p+1)^2\mathfrak{h}$,
where $\mathfrak{h}$, its holonomy member, does not depend on $p$; no ceiling on the length of the
chain is assumed.

\emph{The bound on $\mathbb{C}_G$.} The walk lemma bounds $\mathbb{C}_G$ in its own norm, which
treats one filing level $m$ at a time: it is a supremum over $m$ and over the anchor cube
$\Delta\in\pi_m$, where $\pi_m$ is Ba{\l}aban's partition into cubes at level $m$, and the walk
lemma gives
\begin{equation*}
\|\mathbb{C}_G\|\le c_G(d,L,M)\qquad\text{at any rates }a_m\le4\kappa L.
\end{equation*}
The norm $\|\cdot\|^{A,\mathfrak{D}}_{(l)}$ of Lemma~\ref{4.26:lem-stage-expansion}, on the other
hand, adds at one site $x$ the terms of every filing level whose domains contain $x$, all filing
levels in one sum; for $\mathbb{C}_G$ the set of slots is empty, $\mathcal{S}=\emptyset$. We pass
from the first bound to the second. Let $i(x)$ be the least level at which a domain containing $x$
can be filed by the filing rule \cite[(1.63)]{B-CMP122a}; a domain $X\ni x$ filed at level $m$ has
$m\ge i(x)$. The levels $m\ge i(x)$ fall into
three groups.
\begin{itemize}
\item The level $i(x)$ itself.
\item The level $i(x)+1$, at which no distance is forced: a domain consisting of the cube
$\Delta_{i(x)+1}(x)$ of $\pi_{i(x)+1}$ containing $x$ and a cube $\Delta''$ meeting
$\Omega_{i(x)+1}$ is filed there with a length of order $\sqrt{d}$, $d$ being the dimension.
\item All levels $m\ge i(x)+2$ together.
\end{itemize}
Each of the first two groups contributes at most $c_G$, by the walk lemma's bound read at the rate
vector $a^{(l)}+\sigma_\ast$. Here $a^{(l)}$ are the running rates of clause (c) of
Lemma~\ref{4.26:lem-stage-expansion}, and
\begin{equation*}
\sigma_\ast:=\min\{\beta\kappa,\ \kappa/(4c^\sharp)\},
\end{equation*}
$c^\sharp$ being the constant entering those rates; each entry of the rate vector is at most
$(1+4\beta)\kappa+1<4\kappa L$, so that the walk lemma applies at these rates.

For the third group, a domain $X$ filed at a level $m\ge i(x)+2$ crosses the whole shell
$\Sigma_{m-1}:=\Omega_{m-1}\setminus\Omega_m$, where the $\Omega_m$ are the regions on which the
filing by levels is based. Write $d_m(X)$ for the length of $X$ at level $m$ in the walk lemma's
norm. The shell $\Sigma_{m-1}$ is of one of the following two kinds, and we treat each.
\begin{itemize}
\item[(I)] $\Sigma_{m-1}$ is a shell made by the operation $\mathbf{T}$ \cite[\S3]{B-CMP119}. Its
width is at least $2R_{m-1}/L\ge2R_{\min}/L$ cubes of $\pi_m$, where $R_{m-1}$ is the number of
\cite[(2.5)]{B-CMP119} at level $m-1$; this follows from the width clause for the shells
$\Sigma_{m'}$ of the cover lemma underlying the walk lemma's norm, together with the nesting of the
regions $\Omega_m$ in \cite{B-CMP119}. Hence
\begin{equation*}
d_m(X)\ge\frac{2R_{\min}}{L}-2\sqrt{d}\ge\frac{R_{\min}}{L},
\end{equation*}
the last inequality because $R_{\min}/L\ge2\sqrt{d}=4$ at $d=4$, which follows from
$R_{\min}\ge5L$ in \eqref{4.26:eq-inherited-conditions-a}.
\item[(II)] $\Sigma_{m-1}$ is a shell of width one cube of a live older piece. Such a piece carries
filing levels at least $\ell_l+2$, where $\ell_l$ is the level of stage $l$ in
Lemma~\ref{4.26:lem-summation}; it is therefore younger than the window, $k_a>h$, hence lies outside
the core, in another connected component of $\Lambda_k^c$. It lies at a distance of at least $R_k$
cubes of $\pi_k$ from the region
$R^{(l)}\subset Z^{\mathrm{on}}=Z\cap\Omega_k^c$ of stage $l$, which every domain of the walk lemma
meets. Hence
\begin{equation*}
d_m(X)\ge R_kL^{k-m}-2\sqrt{d}\ge\frac{R_{\min}}{L}.
\end{equation*}
\end{itemize}
In either case no domain filed at a level $m\ge i(x)+2$ has $d_m(X)<R_{\min}/L$. The surplus rate
$\sigma_\ast$ therefore damps the contribution of each such level by the factor
\begin{equation*}
e^{-\sigma_\ast R_{\min}/L}\le e^{-\sigma_\ast R_\ast/L},
\end{equation*}
using $R_{\min}\ge R_\ast$, where $R_\ast$ is the radius of the smallness condition in summed form
\eqref{4.26:eq-summed-smallness-a}, which requires $R_{\min}\ge R_\ast$. There are at most
$\check{R}_k$ such levels, $\check{R}_k$ being the
envelope of Ba{\l}aban's number $R_k$; together they contribute at most
$\check{R}_ke^{-\sigma_\ast R_\ast/L}c_G\le c_G$, by the condition
$\check{R}_ke^{-\sigma_\ast R_\ast/L}\le1$ among the conditions on the constants in the hypotheses
of Lemma~\ref{4.26:lem-stage-expansion}. Hence
\begin{equation*}
\|\mathbb{C}_G\|^{A,\mathfrak{D}}_{(l)}
\le\bigl(1+1+\check{R}_ke^{-\sigma_\ast R_\ast/L}\bigr)c_G\le3c_G.
\end{equation*}
The constant obtained is $(2+\check{R}_ke^{-\sigma_\ast R_\ast/L})c_G$; the bound $3c_G$ depends on
none of $n$, $N_0$, $k$, $\check{R}$, the volume or $g$, $\check{R}_k$ entering only through the
condition just used. The correction $\Delta W$, which consists of single cubes, whose length is $0$,
takes no
factor for levels, and the constant $C_0^{\sharp}=3c_G+1$ is not affected.

Write $dCM$ for the cost of the transitions between
the gauges of adjacent clause cubes; it is kept explicit rather than absorbed into $O(1)$, and $d$,
$C$, $M$ are fixed before $\gamma$.
The per-walk weight $(p+1)^2$ is summed against the geometric walk series of the walk lemma at its
unchanged ratio $x<1$:
\begin{equation*}
\sum_{p\ge0}(p+1)^2x^p=(1+x)(1-x)^{-3},
\end{equation*}
a number absorbed in $O(1)$. No floor on $M$ is changed, and no ceiling on $p$ is imposed. The field
dependence of $\mathbb{C}_G$, whose size we denote by $\mathsf{F}_G$, is bounded as follows, $c_G'$
being the constant of the walk lemma's bound on field dependence:
\begin{equation}\label{4.26:eq-expansion-leaves-4}
\begin{aligned}
\mathsf{F}_G&\le c_G'\,O(1)\bigl[dCM\cdot 5w_\ast(BCM\tilde{\alpha}_0+\tilde{\alpha}_1)
+K_\Delta\gamma_0^{(k)}+\mathfrak{h}\bigr]\\
&\le c_G'\,O(1)\bigl[dCM\,(BCM+1)+K_\Delta\bigr]\gamma_0^{(k)}+c_G'\,O(1)\,\mathfrak{h}\\
&\le c_G'\,O(1)\bigl[dCM\,(BCM+1)+K_\Delta\bigr]\bar{\gamma}_0(\mathsf{t}_k)
+\tfrac14\vartheta(\mathsf{t}_k)\\
&\le \vartheta(\mathsf{t}_k).
\end{aligned}
\end{equation}
The steps are justified as follows.
\begin{itemize}
\item The first inequality is the walk lemma's field-dependence bound
$c_G'\,O(1)\cdot5w_\ast(BCM\tilde{\alpha}_0+\tilde{\alpha}_1)$ with the increment
$K_\Delta\gamma_0$ of ($\Delta$-3); the factor $dCM$ is the cost, introduced above, of passing
between the gauges of adjacent clause cubes, and the holonomy bound enters only through the member
$BCM$ and through the separate member $\mathfrak{h}$. It is used for every walk through more than
one cube, including walks whose chain of cubes winds around the torus.
\item The second inequality uses $5w_\ast\tilde{\alpha}_i\le\gamma_0^{(k)}$ for $i=0,1$.
\item The third inequality uses the run-free majorant $\gamma_0^{(k)}\le\bar{\gamma}_0(\mathsf{t}_k)$
of Definition~\ref{4.26:def-regularity-radius} and the clause
$c_G'\,O(1)\,\mathfrak{h}\le\frac14\vartheta$ of the holonomy bound.
\item The last inequality is the condition (a) of \eqref{4.26:eq-expansion-leaves-a} below,
evaluated at $\mathsf{t}=\mathsf{t}_k\ge\mathsf{t}_\gamma$.
\end{itemize}
Of the members $BCM$, $1$ and $K_\Delta$ of the bracket $dCM\,(BCM+1)+K_\Delta$ in
\eqref{4.26:eq-expansion-leaves-4}, only $BCM$ uses the holonomy bound; besides
it, the holonomy bound contributes the separate member $\mathfrak{h}$. The walk lemma uses the
holonomy bound in the same way as this step. The member $K_\Delta$ comes from the bounds
($\Delta$-1)--($\Delta$-3). The member $1$ comes from $A'$: the entries $|A'|$ and
$|\nabla^\eta_U A'|$, against $\alpha_{1,m}(L^m\eta)^{-1}$ and $\alpha_{1,m}(L^m\eta)^{-2}$, of
\cite[(1.64), (1.66), (1.68)]{B-CMP122a}, $\eta$ being the lattice spacing in
\cite[(1.64)]{B-CMP122a}, estimated on each cube separately with no holonomy.

For the cubes of $\Delta W$, let $R_k$ be Ba{\l}aban's number of \cite[(2.5)]{B-CMP119}, and let
$\beta_\star$ be the constant of the
walk lemma's bound for $\Delta W$. By the walk lemma, the contribution $\mathsf{F}_W$ of the cubes
of $\Delta W$ satisfies
\begin{equation}\label{4.26:eq-expansion-leaves-5}
\begin{aligned}
\mathsf{F}_W&\le\beta_\star R_k(LM)^4\cdot 4(5w_\ast)^2(2\tilde{\alpha}_0+\tilde{\alpha}_1)^2\\
&\le\beta_\star R_k(LM)^4\cdot 36\,\bar{\gamma}_0(\mathsf{t}_k)^2\le\vartheta(\mathsf{t}_k).
\end{aligned}
\end{equation}
Here
\begin{equation*}
R_k=R(\mathsf{t}_k)<\bar{N}(\mathsf{t}_k):=L\,\mathsf{t}_k^{\,r},
\end{equation*}
by the minimality in \cite[(2.5)]{B-CMP119}, evaluated at the same $\mathsf{t}$. Accordingly the
factor $\beta_\star R_k(LM)^4$ is majorised by $\beta_\star\bar{N}(\mathsf{t}_k)(LM)^4$. The last
inequality in \eqref{4.26:eq-expansion-leaves-5} follows from $R_k<\bar{N}(\mathsf{t}_k)$ and the
condition (b) of \eqref{4.26:eq-expansion-leaves-a}, evaluated at
$\mathsf{t}=\mathsf{t}_k\ge\mathsf{t}_\gamma$.

The constant $36=4\cdot3^2$ comes from
$5w_\ast(2\tilde{\alpha}_0+\tilde{\alpha}_1)\le 3\gamma_0^{(k)}$, which we now justify.
\begin{itemize}
\item The numbers $\tilde{\alpha}_0,\tilde{\alpha}_1$ are the arguments of Ba{\l}aban's space
$\tilde{U}^{(n)c}_k(X,\tilde{\alpha}_0,\tilde{\alpha}_1)$ of \cite[p.~190]{B-CMP122a}. For $n=0$
this space coincides with $\tilde{U}^{c}_k(X,\tilde{\alpha}_0,\tilde{\alpha}_1)$
\cite[p.~191]{B-CMP122a}.
\item They name the level families
\begin{equation*}
\alpha_{0,j}=g_jC_0(\log g_j^{-2})^{q_0},\qquad\alpha_{1,j}=g_jC_1(\log g_j^{-2})^{q_1}
\end{equation*}
of \cite[(2.28)]{B-CMP119}.
\item The clause members read, level by level and with no enlargement factor,
\begin{equation*}
(1-\beta(1-2^{-(j-n)}))\alpha_{0,n},\qquad BCM\alpha_{0,n},\qquad
(1-\beta(1-2^{-(j-n)}))\alpha_{1,n}
\end{equation*}
\cite[(2.34)--(2.39)]{B-CMP119}, and, in \cite[(1.65), (1.67), (1.69)]{B-CMP122a},
$BCM\alpha_{0,m}$ multiplied by the power of $L_0$ that appears in $w_\ast(j)=L_0^{2N_0(j)}$.
\item The walk lemma uses $\tilde{\alpha}_0,\tilde{\alpha}_1$ as they stand, with the factor
$5w_\ast$ outside.
\end{itemize}
Hence no enlargement factor $1+2\beta$, such as appears in the arguments of the spaces of
\cite[p.~277]{B-CMP109}, enters $\tilde{\alpha}_i$, so that
$\tilde{\alpha}_i\le\bar{\alpha}_{i,k}$ and $2\tilde{\alpha}_0+\tilde{\alpha}_1\le3\bar{\alpha}_k$.
Since $5w_\ast(k)\,\bar{\alpha}_k\le\gamma_0^{(k)}$ by the definition
\eqref{4.26:eq-regularity-radius-a}, it follows that
\begin{equation*}
5w_\ast(2\tilde{\alpha}_0+\tilde{\alpha}_1)\le5w_\ast(k)\cdot3\bar{\alpha}_k\le3\gamma_0^{(k)}.
\end{equation*}

The two conditions invoked above are
\begin{equation}\label{4.26:eq-expansion-leaves-a}
\begin{gathered}
\text{(a)}\quad\sup_{\mathsf{t}\ge\mathsf{t}_\gamma}
\frac{c_G'\,O(1)\bigl[dCM\,(BCM+1)+K_\Delta\bigr]\bar{\gamma}_0(\mathsf{t})}{\vartheta(\mathsf{t})}
\le\frac34,\\[2pt]
\text{(b)}\quad\sup_{\mathsf{t}\ge\mathsf{t}_\gamma}
\frac{\beta_\star\,\bar{N}(\mathsf{t})\,(LM)^4\cdot36\,\bar{\gamma}_0(\mathsf{t})^2}{\vartheta(\mathsf{t})}
\le1.
\end{gathered}
\end{equation}
Each is a supremum in the single variable $\mathsf{t}$. In (a) the constant is taken at the share
$\frac34$, the remaining
$\frac14\vartheta$ being the share of the holonomy member. The quotient in (a) has the profile
$\mathsf{t}^{\,r+2q-p_1}e^{-\mathsf{t}/2}$, and that in (b) the profile
$\mathsf{t}^{\,3r+3q-p_1}e^{-\mathsf{t}}$; both quotients therefore have degree $-\infty$. Hence
(a) and (b) are $\gamma$-ceilings of degree $-\infty$, of the same kind as the $\gamma$-ceilings
among the conditions \eqref{4.26:eq-inherited-conditions-a}; they hold once
$\mathsf{t}_\gamma=\log\gamma^{-2}$ exceeds a threshold depending only on the constants displayed in
them, so they are smallness conditions on $\gamma$, and they are imposed in this proof.

\emph{Totals.} The totals of the transferred proof are \eqref{4.26:eq-summation-a} and
\eqref{4.26:eq-summation-b}, taken at the constant $K_\flat'$. By
\eqref{4.26:eq-summed-smallness-a} and the condition on the counting constants $\mathsf{M}_\ast$,
$V$ and $\tilde{\mathfrak{l}}$ imposed with it,
\begin{equation*}
\mu+\tilde{\mathfrak{l}}\le\tfrac12,\qquad
4(\mathsf{M}_\ast+2^dV\tilde{\mathfrak{l}})\le\tfrac12.
\end{equation*}
With $\|\mathbb{C}_G\|\le3c_G$ and the Gaussian part bounded as above, the totals give
\begin{equation*}
\|\mathbb{C}^{(l)}\|^{A,\mathfrak{D}}_{(l)}\le 3c_G+\vartheta+\tfrac12\le C_0^{\sharp}.
\end{equation*}
The constant $C_0^{\sharp}=3c_G+1$ is fixed before these conditions are imposed, so the argument is
not circular. The bound $\nu^{\natural}(e)\le\frac12$ follows from the bound on $\mu$. For the
elements $e_{\mathrm{str}}$ we require
\begin{equation*}
\kappa_A\ge\log\bigl(8\max\{c_T^{\mathrm{sz}},C_0^{\sharp}\}\bigr),
\end{equation*}
a lower bound on $\kappa_A$ which we impose together with the conditions on the constants in the
hypotheses of Lemma~\ref{4.26:lem-stage-expansion}.

\emph{$\mathcal{G}$-invariance.} Each ingredient of the terms behaves under the group $\mathcal{G}$
of joint $G$-valued gauge transformations as follows.
\begin{itemize}
\item The stage layer $\mathbb{H}(\sigma;b)$ and its local current slot
$\mathcal{J}^{\sharp}[1](\sigma;b)$ are $\operatorname{Ad}$-covariant, by the covariance clause of
the stage layer.
\item The functions $\zeta_\square$ of the partition of unity, and the parameters $t$, $u$ and
$\sigma$ of the expansion, are scalars.
\item The functional $c$ is jointly invariant \cite{B-CMP119}.
\item The measure, the functions $\chi$, the boxes and the polydiscs are
$\operatorname{Ad}$-invariant.
\item The domain $\mathfrak{D}_l(X)$ is $\mathcal{G}$-invariant.
\end{itemize}
This gives the invariance asserted in clause (b).

\emph{Clause (d).} The constants entering the steps above, among them $c_T^{\mathrm{sz}}$,
$\kappa_T$, $\kappa_A$, $A_\flat$, $\sigma_\ast$, $c_F$, $K_\Delta$, $c_G$, $c_G'$, $\beta_\star$,
the factors $BCM$ and $dCM$, the numbers $O(1)$, $\mathsf{w}_0$, $K_\flat'$ and
$C_0^{\sharp}=3c_G+1$, together with $L$, $M$, $\beta$, the exponents $r$, $q$, $p_1$ and the
counting constants $\mathsf{M}_\ast$, $V$, $\tilde{\mathfrak{l}}$ of the totals, depend on none of
$n$, the number of stages, $N_0$, $\check{R}$, $k$, $K$, the volume or $g$,
which is clause (d); among them $\mathsf{w}_0$ and $K_\flat'$ are fixed before $\gamma$.
\end{proof}

\begin{definition}[The enlarged regularity space with its composite shell clause]
\label{4.26:def-regularity-space}
Fix the density $k$ and the $\mathbf{R}$-site of Definition~\ref{4.26:def-stages-regions}, with its
stages $l=1,\dots,N$, its regions $Z^{\mathrm{on}}=Z\cap\Omega_k^c$ and
$\mathfrak{T}=\Omega_k\cup Z^c$, its live arrests and their frozen levels and weights. Let
$\mathfrak{K}_P=\mathfrak{K}_{109}(\mathfrak{B};\gamma_0)$ be the analyticity polydisc of size
$\gamma_0^{(k)}$, the radius of \eqref{4.26:eq-regularity-radius-a} (Definition~\ref{4.26:def-regularity-radius}).
Let $Q_1,\dots,Q_8$ be the clause entries, $\mathbf{r}_{a,m}$ the unweighted shell radii and
$(p,X',\mathfrak{s})$ the $(\Sigma^{\ast})$-triples of Definition~\ref{4.26:def-clauses-units}; entries $2$
and $4$ are the composite entries, entry $8$ is the cube clause. Let $r_a(y)$ be the weighted radii of
Definition~\ref{4.26:def-stages-regions}, $w_\ast(j)$ the weight of
Definition~\ref{4.26:def-regularity-radius}, and $\beta$ the schedule constant of
Definition~\ref{4.4:def-post-step-space};
and let $\iota^1_l$ be the hereditary part of the stage increment $\iota_l$ that enters
Lemma~\ref{4.26:lem-family-consumption}.

Let $Y$ be a union of cubes and $0\le n\le N$. Every location at which a clause below is imposed is
assigned to $Z^{\mathrm{on}}$ or to $\mathfrak{T}$ according to which of the two regions contains it;
a cube of the cube clause that meets both regions is assigned to $Z^{\mathrm{on}}$.

\emph{Weight and level.} For a point $x$ that lies in a live piece, let $\mathfrak{a}^{\circ}$ be the
innermost live arrest containing $x$, and let $k^{\circ}\le k$ be its density, with $k^{\circ}=k$ for a
component excluded at this site. We put $\omega(x):=w_{\mathfrak{a}^{\circ}}(x)$, the frozen weight of
$\mathfrak{a}^{\circ}$ at $x$ (Definition~\ref{4.26:def-stages-regions}), and
$m(x):=\ell_{\mathfrak{a}^{\circ}}(x)$, its frozen level at $x$. At a
point of no live piece we put $\omega(x):=1$, and $m(x)$ is the level at $x$. The radii
$\mathbf{r}_{a,m(x)}$ are the unweighted ones; the weight enters only through
$\mathsf{w}^{\mathrm{on}}(x)$ below.

\emph{Shell radii of the composite entries.} For a $(\Sigma^{\ast})$-triple $(p,X',\mathfrak{s})$, a
point $y'$ and $m=m(y')$, the shell radius of entry $e\in\{2,4\}$ at $y'$ for the level $p$ is
\begin{equation}\label{4.26:eq-regularity-space-1}
r_2^{(p)}(y'):=\alpha_{0,m}\,L^{-2p}\bigl(L^{m}L^{-p}\bigr)^{-2},\qquad
r_4^{(p)}(y'):=\alpha_{0,m}\,L^{\,m-\min\{p,m\}}\bigl(L^{m}L^{-p}\bigr)^{-3}.
\end{equation}
If the structure $\mathfrak{s}$ of the triple has, at $y'$, the level
$\lambda_{\mathfrak{s}}(y'):=\mathsf{m}^{\circ}(y')<m(y')$ (this occurs on a live older piece, at a
retained structure of the piece's own run, that is, at a triple of the family
$(\Sigma^{\ast})^{\mathfrak{a}}$ of \eqref{4.26:eq-graded-cover-a}), then in
\eqref{4.26:eq-regularity-space-1} the exponent $\min\{p,m\}$ is read as
$\min\{p,\lambda_{\mathfrak{s}}(y')\}$, while the radius $\alpha_{0,m}$ and the weight remain those
frozen at the point $y'$.

\emph{The ambient domain.} For $x\in Y$ let $B(x)$ be the ball of radius
$\varrho_0\ell_{m(x)}$ on which the lemma propagating the shell clauses through a stage uses its
inputs, and put
\begin{equation}\label{4.26:eq-regularity-space-b}
Y^{+}:=Y\cup\bigcup_{x\in Y}B(x).
\end{equation}

\emph{The space.} We say that $(\mathbb{U},\mathbb{J})\in\mathfrak{N}^{(n)}(Y)$ if and only if
$(\mathbb{U},\mathbb{J})$ is the restriction to $Y$ of a pair which belongs to $\mathfrak{K}_P$ on the
whole domain of $\mathfrak{B}$, which satisfies the clauses \textup{(A)}--\textup{(C)} below at every
location of $Y^{+}$, and which satisfies the clauses \textup{(D)} and \textup{(E)} below for every
$(\Sigma^{\ast})$-triple $(p,X',\mathfrak{s})$ with $X'\subset Y$, every
$y'\in X'^{\sim 2}\cap\mathfrak{T}$ and every $e\in\{2,4\}$:
\begin{equation}\label{4.26:eq-regularity-space-a}
\begin{aligned}
&\mathsf{w}^{\mathrm{on}}(x):=\max\{w_\ast(k),\omega(x)\},\\
&\textup{(A)}\quad Q_a(x)<5\,\mathsf{w}^{\mathrm{on}}(x)\,\mathbf{r}_{a,m(x)}
&&(x\in Z^{\mathrm{on}},\ a\in\{1,3,5,6,7,8\}),\\
&\textup{(B)}\quad Q_a(y)<\Bigl(1+2\beta-\sum_{l\le n}\iota^1_l(y)\Bigr)\,r_a(y)
&&(y\in\mathfrak{T},\ a\in\{1,3,6,7\}),\\
&\textup{(C)}\quad Q_a(y)<(1+2\beta)\,r_a(y)
&&(y\in\mathfrak{T},\ a\in\{5,8\}),\\
&\textup{(D)}\quad Q_e^{(p,X',\mathfrak{s})}(y')
<\Bigl(1+2\beta-\sum_{l\le n}\iota^1_l(y')\Bigr)\,r_e^{(p)}(y')
&&(y'\in X'^{\sim 2}\cap\mathfrak{T},\ e\in\{2,4\}),\\
&\textup{(E)}\quad Q_e^{(p,X',\mathfrak{s})}[U](y')<(1+2\beta)\,r_e^{(p)}(y')
&&(y'\in X'^{\sim 2}\cap\mathfrak{T},\ e\in\{2,4\}).
\end{aligned}
\end{equation}
Here $Q_e^{(p,X',\mathfrak{s})}[U]$ denotes the same composite entry formed with the unitary part $U$
of $\mathbb{U}$ in place of $\mathbb{U}$. The clauses \textup{(A)}--\textup{(C)} are the
\emph{direct clauses}; \textup{(D)} is the \emph{composite shell clause}, and \textup{(E)} is its
\emph{companion clause for the unitary part}. This way of defining membership is the
\emph{ambient-domain convention}. We put $\mathfrak{N}(Y):=\mathfrak{N}^{(0)}(Y)$.

Then:
\begin{enumerate}
\item on $Z^{\mathrm{on}}$ the direct clauses are equivalent to
$Q_a(x)<5\max\{w_\ast(k)/\omega(x),1\}\,r_a(x)$, with the weighted radius $r_a(x)$;
\item at $x\in Z^{\mathrm{on}}$ the factor $5\,\mathsf{w}^{\mathrm{on}}(x)\,\alpha_{\cdot,m(x)}$ of the
threshold of \textup{(A)}, where $\alpha_{\cdot,m(x)}$ is the radius of level $m(x)$ from which
$\mathbf{r}_{a,m(x)}$ is formed, is bounded by the term $j=k$ or the term $j=k^{\circ}$ of
\eqref{4.26:eq-regularity-radius-a}, and $\mathfrak{N}^{(n)}(Y)\subset\mathfrak{K}_P$;
\item $\mathfrak{N}^{(n)}(Y)$ is open and $\mathcal{G}$-invariant, and it is stable under the
restriction operation.
\end{enumerate}
\end{definition}

\begin{proof}
\emph{The equivalent form on $Z^{\mathrm{on}}$.} On a live piece the weighted radius is the frozen weight
times the unweighted radius, $r_a(x)=\omega(x)\mathbf{r}_{a,m(x)}$. Since $\omega(x)>0$,
\begin{equation*}
5\max\Bigl\{\frac{w_\ast(k)}{\omega(x)},1\Bigr\}\,r_a(x)
=5\max\{w_\ast(k),\omega(x)\}\,\mathbf{r}_{a,m(x)}
=5\,\mathsf{w}^{\mathrm{on}}(x)\,\mathbf{r}_{a,m(x)},
\end{equation*}
so the two forms impose the same inequality. The weight is written once: either inside
$\mathsf{w}^{\mathrm{on}}$ with the unweighted radius, as in \eqref{4.26:eq-regularity-space-a}, or
inside $r_a$, with the ratio $w_\ast(k)/\omega(x)$ in the factor; the two conventions are not combined.

\emph{Containment in $\mathfrak{K}_P$.} The threshold of the direct clause \textup{(A)} at
$x\in Z^{\mathrm{on}}$ is $5\mathsf{w}^{\mathrm{on}}(x)\mathbf{r}_{a,m(x)}$, and the radius
$\mathbf{r}_{a,m(x)}$ is formed from the radius $\alpha_{0,m(x)}$ or $\alpha_{1,m(x)}$ of level $m(x)$
that measures the entry $a$ (Definition~\ref{4.26:def-clauses-units}); write $\alpha_{\cdot,m(x)}$ for
this radius. The comparison with the size $\gamma_0^{(k)}$ of the polydisc is made on these radii, that
is, between $5\mathsf{w}^{\mathrm{on}}(x)\alpha_{\cdot,m(x)}$ and the terms of
\eqref{4.26:eq-regularity-radius-a}. Recall that $\bar{\alpha}_j$ is the maximum of the radii
$\alpha_{i,m}$ over the levels $m\le j$. There are two branches of the maximum defining
$\mathsf{w}^{\mathrm{on}}(x)$.

If $\mathsf{w}^{\mathrm{on}}(x)=w_\ast(k)$, then, because $m(x)\le k$,
\begin{equation*}
5\,w_\ast(k)\,\alpha_{\cdot,m(x)}\le 5\,w_\ast(k)\,\bar{\alpha}_k ,
\end{equation*}
and the right-hand side is the term $j=k$ in the definition
\eqref{4.26:eq-regularity-radius-a} of $\gamma_0^{(k)}$.

If $\mathsf{w}^{\mathrm{on}}(x)=\omega(x)$ on a live piece, let $k^{\circ}$ be the density of the
innermost live arrest $\mathfrak{a}^{\circ}$ containing $x$. The frozen weight of an arrest of density
$k^{\circ}$ does not exceed $w_\ast(k^{\circ})$, and its frozen level does not exceed $k^{\circ}$;
hence $\omega(x)\le w_\ast(k^{\circ})$ and $m(x)\le k^{\circ}$, and
\begin{equation*}
5\,\omega(x)\,\alpha_{\cdot,m(x)}\le 5\,w_\ast(k^{\circ})\,\bar{\alpha}_{k^{\circ}},
\end{equation*}
which is the term $j=k^{\circ}$ of \eqref{4.26:eq-regularity-radius-a}. Since $k^{\circ}\le k$, this
term is one of those over which $\gamma_0^{(k)}$ is formed. In both branches the threshold is therefore
bounded by a term of $\gamma_0^{(k)}$; this is the content of the definition of $\gamma_0^{(k)}$, which
was made so that the direct clauses on $Z^{\mathrm{on}}$ lie within the polydisc of size
$\gamma_0^{(k)}$. Since, moreover, by the ambient-domain convention every element of
$\mathfrak{N}^{(n)}(Y)$ is the restriction of a pair of $\mathfrak{K}_P$, we have
$\mathfrak{N}^{(n)}(Y)\subset\mathfrak{K}_P$. Consequently every estimate of
Definition~\ref{4.26:def-regularity-radius} stated on $\mathfrak{K}_P$ holds on
$\mathfrak{N}^{(n)}(Y)$.

\emph{The clauses on $\mathfrak{T}$ and the conditions of \cite{B-CMP109}.} The direct clauses
\textup{(B)} and \textup{(C)} of \eqref{4.26:eq-regularity-space-a} are the conditions (i)--(iii) of
\cite{B-CMP109} on
$\mathfrak{T}$, at the factor $1+2\beta$ used there (see also \cite{B-CMP116}), diminished at the
entries $1,3,6,7$ by the hereditary increments $\iota^1_l$ of the stages $l\le n$. Condition (iv) of
\cite{B-CMP109} enters at the same factor on $\mathfrak{T}$ through the composite shell clause
\textup{(D)}. The constants in these conditions are those of \cite{B-CMP109}. Neither
$(\mathsf{w}^{\mathrm{on}})^5$ nor $5^5$ enters the constant $C_R^{\natural}$ of the estimates stated
on $\mathfrak{K}_P$; the factor $5\,\mathsf{w}^{\mathrm{on}}$ of the clause \textup{(A)} does not
propagate into it. The box
clause relative to a term, which enters the diminished term clauses $\mathfrak{I}_n$, is a different
condition and is not part of $\mathfrak{N}^{(n)}(Y)$.

\emph{The composite clause and its quantifier.} The quantifier of the composite shell clause
\textup{(D)} is that of the composite clause with the factor $f_n(b)$ imposed at stage $n$: for every
$(\Sigma^{\ast})$-triple and every derived bond $b$ of its structure, the entries $Q_2(b),Q_4(b)$ are
bounded by $f_n(b)$ times $\mathbf{r}_{2,m(b)}$, respectively
$\mathbf{r}_{4,m(b)}$. The radius $r_e^{(p)}(y')$ of \eqref{4.26:eq-regularity-space-1} coincides with
the radius $r_e(m,l_{y'})$ of the lemma propagating the shell clauses, where $m=m(y')$ and
$l_{y'}=\min\{p,\mathsf{m}^{\circ}(y')\}$, and with the radius $\mathbf{r}_{e,m(b)}$ of that composite
clause, where $b$ is the derived bond of the triple's structure at $y'$.

\emph{The unit of the composite clause.} The unit in the composite shell clause is the box's own
unit. In the notation of \cite{B-CMP109}, with its spacing $\xi$, its radius $\alpha_0$ and its
minimiser $M'(\mathbf{U})$ on the box, \cite[(1.16)]{B-CMP109} states that for every $n$ (here $n$ is
the level index of \cite{B-CMP109}, not the stage index of this definition),
\begin{equation}\label{4.26:eq-regularity-space-2}
\bigl|\partial U_n(M'(\mathbf{U}))-1\bigr|<\alpha_0\,\xi^2,\qquad
\bigl|J_n(M'(\mathbf{U}))\bigr|<\alpha_0\,(L^n\xi)^2\qquad\text{on }X^{\sim 2}.
\end{equation}
The composite shell clause is this condition, read in the shell radii
\eqref{4.26:eq-regularity-space-1}, at a factor less than $1+2\beta$. It is not the unit
$\mathbf{r}^{T,(p)}_e$ of a term's own clause (Definition~\ref{4.26:def-clauses-units}): at $p<m$
that unit is weaker than \eqref{4.26:eq-regularity-space-2} by the factor $L^{2(m-p)}$.

At a triple whose structure $\mathfrak{s}$ has level $\lambda_{\mathfrak{s}}(y')<m(y')$ at $y'$, the
block side of the lemma propagating the shell clauses governs; this is why $\min\{p,m\}$ is read as
$\min\{p,\lambda_{\mathfrak{s}}(y')\}$. The block side is taken at the structure's level
$\lambda_{\mathfrak{s}}(y')$, not at the level $p$ of the triple. With the literal exponent
$\min\{p,m\}$ the clause is not propagated by that lemma, and it breaks, at
such triples, the first of the conditions on the weighted radii of a live piece in
Definition~\ref{4.26:def-stages-regions}.

Nor is the unit the birth radius of a retained structure. That radius is the unit of the clauses of
the structure's own run, that is, of the family $(\Sigma^{\flat})$ and of the stage clauses
$\mathfrak{L}_n$ of the site $k_a$ on its excluded component $W$. Every clause written at a later
density, including the stage clauses $\mathfrak{L}_l$ of an enclosing run, is measured in the frozen
shell unit with the block side of the lemma propagating the shell clauses. The birth radius is larger
than the shell unit by a factor at least $(2.625L)^{m-\lambda_{\mathfrak{s}}}$.

\emph{The companion clause for the unitary part.} In \cite{B-CMP109} the bounds
\eqref{4.26:eq-regularity-space-2} are stated together with the same bounds for $U$ in place of
$\mathbf{U}$ (see \cite[(1.16)]{B-CMP109}). They are part of the hypotheses of the space on which the
estimates of \cite{B-CMP109} that use them are stated. No stage changes the unitary part: in
Lemma~\ref{4.26:lem-family-consumption} the increment of the entry $5$ is $\Delta_5=0$. Hence the
companion clause \textup{(E)} propagates from stage
to stage at its fixed factor $1+2\beta$, in the same way as the entries $5$ and $8$, whose thresholds in
the clause \textup{(C)} of \eqref{4.26:eq-regularity-space-a} carry no increment either. The source space
$\mathfrak{S}_{\mathrm{src}}$ and the tube $\mathfrak{D}_{\mathfrak{E}}$ carry it at a factor at most
$1.1$.

\emph{The ambient-domain convention: why $Y^{+}$.} The lemma propagating the shell clauses uses its
inputs on the ball $B(x)$ of radius $\varrho_0\ell_{m(x)}$, which, by the condition fixing the radius
$\varrho_0$, is smaller than one $M$-cube of the local scale. At points
$x\in Y\setminus Y^{\sim 2}$, and at locations where the scale drops, this ball may leave $Y$. The
direct clauses are therefore required on all of $Y^{+}$ of \eqref{4.26:eq-regularity-space-b}, for a
pair defined on the whole domain of $\mathfrak{B}$.

The existential quantifier stands where the ambient-domain convention places it. A member of
$\mathfrak{N}^{(n)}(Y)$ is the restriction of one pair that belongs to $\mathfrak{K}_P$ globally and
satisfies the direct clauses on $Y^{+}$. It is not a pair on $Y^{+}$ of which it is merely asserted that
some extension in $\mathfrak{K}_P$ exists. A pair satisfying the direct clauses on a region $Y^{+}$
that is not simply connected need not extend with small curvature.

The composite shell clause is imposed only at triples with $X'\subset Y$. Otherwise the enlargements
$X'^{\sim 2}$ of triples not contained in $Y$ would force the domain of the direct clauses to be
enlarged again, and so on.

\emph{Openness, invariance and restriction.} Every clause of \eqref{4.26:eq-regularity-space-a} is a
strict inequality and $\mathfrak{K}_P$ is open; in this sense the class $\mathfrak{N}^{(n)}(Y)$ is open.
It is $\mathcal{G}$-invariant and stable under the restriction operation.

\emph{Locality.} The terms remain local in $Y$; only the membership condition involves $Y^{+}$. The
proof family $\mathfrak{W}^{(l),\natural}(Y)$ of Lemma~\ref{4.26:lem-family-consumption} respects this.
Its parameters satisfy $t_{\square}=0$ unless $\tilde{\square}\subset Y$, by the definition of the
parameters of the family. Hence $\chi_t$ is supported in
the union of the enlarged cubes $\tilde{\square}$ of an element $c$ of the family, which lies in the
element's domain $\mathsf{d}(c)\subset Y$. Consequently the configuration of the family satisfies
$\Phi_{t,\sigma}(b)=(\mathbb{U},\mathbb{J})$ on $Y^{+}\setminus Y$, where the direct clauses are those
of the input pair.

\emph{Beyond $Y^{+}$.} At an unlocalised $\sigma$ the stage layer $\mathbb{H}(\sigma;b)$ depends on
$\mathbb{U}$ on the support of $\sigma$. Its holomorphy is obtained along the pair's own continuation in
$\mathfrak{K}_P$. This is the situation described in \cite{B-CMP116}: the quantity considered there
depends on the pair restricted to $Y$ only, and the conditions outside $Y$ are unessential.
The value of each localised term depends on the pair on its localisation domain $X$ only, by the
localisation property of the terms. In every use of the class (the proof family, the tube and the units
of the correlation theorem) the pair is a field on the whole domain of $\mathfrak{B}$.
\end{proof}

\begin{definition}[Lettered clauses, diminished term clauses and the inherited domain]
\label{4.26:def-inherited-domain}
Let $Y$ be a union of cubes and let $0\le n\le N$, where $N$ is the last stage, as in
Definition~\ref{4.26:def-regularity-space}. The terms $\mathcal{T}$, their two classes
$\mathcal{T}^{\mathbb{E}\mathbb{R}}$ and $\mathcal{T}^{\mathbb{B}}$, the localisation domain $X_T$
and clause set $\operatorname{Cl}(T)\subset X_T\times\{1,\dots,8\}$ of a term $T$, the entries $Q_a$,
the thresholds $\theta_{T,a}$ and $\theta_{\mathrm{src},a}$, the radii $\mathbf{r}^T_a$ and
$\mathbf{r}^{T,(p)}_a$, and the unit rule \eqref{4.26:eq-clauses-units-b} are those of
Definition~\ref{4.26:def-clauses-units}. The space $\mathfrak{N}^{(n)}(Y)$, the regions
$Z^{\mathrm{on}}$ and $\mathfrak{T}$, and the assignment of each clause location to $Z^{\mathrm{on}}$ or
to $\mathfrak{T}$ are those of Definition~\ref{4.26:def-regularity-space}. The composite entries
$a\in\{2,4\}$ are read at the triples $(p,X',\mathfrak{s})$ of the derived-entry family
$(\Sigma^{\ast})$, and $Q_a^{(p,X',\mathfrak{s})}(y)$ denotes entry $a$ at $y$ read at the triple
$(p,X',\mathfrak{s})$. Every set defined below is a set of background pairs
$(\mathbb{U},\mathbb{J})$.
\begin{enumerate}
\item[(a)] \emph{Lettered clauses.} For $A\subset Y\cap Z^{\mathrm{on}}$, $\mathfrak{L}_n(A)$ is the set
of pairs at which the lettered clauses of the space $\tilde{U}^{(n)c,\Sigma\flat}_k$, that is,
\cite[(1.64)--(1.69)]{B-CMP122a} at index $n$, with the composite entries read through
$(\Sigma^{\ast})$, hold at every location of $A$. On a component $W=W_{\mathfrak{a}}$ excluded at an
arrest $\mathfrak{a}$, and for $n>n^+_{\mathfrak{a}}$, where $n^+_{\mathfrak{a}}$ is the stage entry of
the arrest $\mathfrak{a}$, these clauses are replaced there by the clauses of the arrest-stage set
$\mathfrak{S}^{(k),\mathrm{arrest}}_n$.
\item[(b)] \emph{Diminished thresholds.} Put $\iota^a_l:=\iota_l$ and $\iota^{1,a}_l:=\iota^1_l$ for
$a\notin\{5,8\}$, and $\iota^a_l:=\iota^{1,a}_l:=0$ for $a\in\{5,8\}$. For $T\in\mathcal{T}$ and
$(y,a)\in\operatorname{Cl}(T)$,
\begin{equation}\label{4.26:eq-inherited-domain-1}
\theta^{(n)}_{T,a}:=
\begin{cases}
\theta_{T,a}-\sum_{l\le n}\iota^a_l\,\mathbf{r}^T_a, & T\in\mathcal{T}^{\mathbb{B}},\\[3pt]
\theta_{T,a}-\tfrac13\bigl(\theta_{T,a}-\theta_{\mathrm{src},a}\bigr)
-\sum_{l\le n}\iota^{1,a}_l\,\mathbf{r}^T_a, & T\in\mathcal{T}^{\mathbb{E}\mathbb{R}}.
\end{cases}
\end{equation}
We write $s^{(n)}_{T,a}(y)$ for the diminished factor of $T$ at $(y,a)$, that is, for
$\theta^{(n)}_{T,a}(y)$ expressed in the unit $\mathbf{r}^T_a$ of the unit rule
\eqref{4.26:eq-clauses-units-b}:
$\theta^{(n)}_{T,a}(y)=s^{(n)}_{T,a}(y)\,\mathbf{r}^T_a(y)$. For a localisation domain
$X'$, $X'^{\sim2}$ denotes $X'$ with its two outer layers of cubes removed, as in
\cite[p.~262]{B-CMP109}.
\item[(c)] \emph{Diminished term clauses.} For $A\subset Y\cap\mathfrak{T}$, $\mathfrak{I}_n(A)$ is the
set of pairs with the following two properties.
\begin{enumerate}
\item[(c1)] For every $T\in\mathcal{T}$ with $X_T\subset Y$ and every $(y,a)\in\operatorname{Cl}(T)$
with $y\in A$: $Q_a(y)<\theta^{(n)}_{T,a}(y)$.
\item[(c2)] For every $T\in\mathcal{T}^{\mathbb{E}\mathbb{R}}$, every $a\in\{2,4\}$ and every
$y\in X_T\cap A$: for every localisation domain $X'\subset Y$ with $X_T\subset X'^{\sim2}$, every
structure $\mathfrak{s}$ of $(\Sigma^{\ast})$ on $X'$, and every level $p$ up to the top level of
$\mathfrak{s}$ on $X'$,
\begin{equation}\label{4.26:eq-inherited-domain-a}
Q_a^{(p,X',\mathfrak{s})}(y)<s^{(n)}_{T,a}(y)\,\mathbf{r}^{T,(p)}_a(y).
\end{equation}
\end{enumerate}
\item[(d)] \emph{The inherited domain.} We put
\begin{equation}\label{4.26:eq-inherited-domain-b}
\mathfrak{D}^{\flat}_n(Y):=\mathfrak{N}^{(n)}(Y)\cap\mathfrak{L}_n\bigl(Y\cap Z^{\mathrm{on}}\bigr)
\cap\mathfrak{I}_n\bigl(Y\cap\mathfrak{T}\bigr),
\end{equation}
and write $\|\cdot\|^{A,\flat}_{(n)}$ for the stage-$n$ norm $\|\cdot\|^{A,\mathfrak{D}}_{(n)}$ of the
stage systems, whose stability hypothesis is \eqref{4.26:eq-stage-systems-a}, taken with the
domain $\mathfrak{D}=\mathfrak{D}^{\flat}_n$; the superscript $A$ is part of the name of this norm and
does not denote a subset of $Y$.
\end{enumerate}
The right side of \eqref{4.26:eq-inherited-domain-a} is the diminished factor of $T$ at $(y,a)$
times the clause radius of $T$ at the level $p$ of the triple, as the unit rule
\eqref{4.26:eq-clauses-units-b} prescribes; the clause is imposed at every pair $(p,\mathfrak{s})$.
At the box-level triple $(p_{\square},\square_0,\mathbb{B}^{(n)}_k)$ of the box analysis, in the range
of top levels $p_{\square}$ considered there, it is the
bound $\theta^{(n)}_{T,a}(y)$ written in the unit $\mathbf{r}^T_a$.

The reserve $\tfrac13(\theta_{T,a}-\theta_{\mathrm{src},a})$ in \eqref{4.26:eq-inherited-domain-1}
is imposed according to the class of $T$: for $T\in\mathcal{T}^{\mathbb{E}\mathbb{R}}$ and not for
$T\in\mathcal{T}^{\mathbb{B}}$, and does not depend on the stage. It is subtracted once, not once per stage.
The stage subtractions are measured in the unit $\mathbf{r}^T_a$ of
Definition~\ref{4.26:def-clauses-units}.

Clause (c2) is needed for the following reason. The own space $\mathfrak{U}_T$ of a term imposes the
composite entries only at triples with $X'\subseteq X_T$, by Definition~\ref{4.26:def-clauses-units}.
The function bounded in the box analysis,
\cite[(3.39)]{B-CMP109}, by contrast reads the triple at the box $\square_0\supsetneq X_T$, at its
top level. Where
the lettered clauses are imposed, the bound at the box is obtained from them, as in the derivation of
the box bound from the lettered clauses in the box analysis. The domain
$\mathfrak{D}^{\flat}_n(Y)$ imposes the lettered clauses only on $Y\cap Z^{\mathrm{on}}$, and on
$Y\cap\mathfrak{T}$ the bound is provided by \eqref{4.26:eq-inherited-domain-a}. This clause belongs
to the unit of $T$ and serves the start of the box analysis. The composite shell clause of
$\mathfrak{N}^{(n)}(Y)$ in \eqref{4.26:eq-regularity-space-a} belongs to the shell unit; it supplies
hypothesis (iv) of \cite[Lemma~4]{B-CMP109}, and the class of pairs that the box analysis requires at
boxes in $\mathfrak{T}$. These are two distinct
clauses with two distinct uses.

The set $\mathfrak{D}^{\flat}_n(Y)$ is open and $\mathcal{G}$-invariant; it decreases as $n$
increases; and it is stable under restriction in the sense of \eqref{4.26:eq-stage-systems-a}.
\end{definition}

\begin{proof}
\emph{Openness.} By Definition~\ref{4.26:def-regularity-space}, $\mathfrak{N}^{(n)}(Y)$ is defined by
finitely many strict inequalities. The sets $\mathfrak{L}_n(Y\cap Z^{\mathrm{on}})$ and
$\mathfrak{I}_n(Y\cap\mathfrak{T})$ are likewise defined by strict inequalities. Each of these
inequalities is a strict inequality between two continuous functions of the pair, the entry at one
location, read directly or at one triple, and its bound, so each cuts out an open set. The lattice is
finite, so there are only finitely many
locations, terms $T$, entries $a$, localisation domains $X'$ and pairs $(p,\mathfrak{s})$. Hence
only finitely many strict inequalities occur. The set on which finitely many strict inequalities
hold is a finite intersection of open sets, and is therefore open. By
\eqref{4.26:eq-inherited-domain-b}, $\mathfrak{D}^{\flat}_n(Y)$ is the intersection of the three sets,
and so it is open.

\emph{$\mathcal{G}$-invariance.} $\mathfrak{N}^{(n)}(Y)$ is $\mathcal{G}$-invariant by
Definition~\ref{4.26:def-regularity-space}. Each of $\mathfrak{L}_n(Y\cap Z^{\mathrm{on}})$ and
$\mathfrak{I}_n(Y\cap\mathfrak{T})$ is an intersection of sets each cut out by one clause, so it
suffices that every clause is $\mathcal{G}$-invariant. The clauses entering these two sets, among them
the lettered clauses \cite[(1.64)--(1.69)]{B-CMP122a}, are $\mathcal{G}$-invariant, in
the form stated in \cite{B-CMP122a} and \cite{B-CMP119}. The intersection
\eqref{4.26:eq-inherited-domain-b} of three $\mathcal{G}$-invariant sets is $\mathcal{G}$-invariant.

\emph{Monotonicity in $n$.} We show that each of the three sets in
\eqref{4.26:eq-inherited-domain-b} at stage $n+1$ is contained in the same set at stage $n$.
\begin{itemize}
\item \emph{The set $\mathfrak{L}_n$.} The lettered clauses \cite[(1.64)--(1.69)]{B-CMP122a} at index
$n+1$ imply those at index $n$, in the form stated in \cite{B-CMP122a} and by the monotonicity of the
stage spaces of a run. On
an excluded component $W$, the amount subtracted in the source factor $F^{(k)}$ of the arrest-stage
set grows with $n$. Hence the clauses of the arrest-stage set at $n+1$ imply those at $n$.
\item \emph{The set $\mathfrak{I}_n$.} Passing from $n$ to $n+1$ adds the term $l=n+1$ to each sum
$\sum_{l\le n}$ in \eqref{4.26:eq-inherited-domain-1}. Each added term is the increment
$\iota^a_{n+1}$ or $\iota^{1,a}_{n+1}$, which is $\ge0$, multiplied by $\mathbf{r}^T_a$. So
$\theta^{(n+1)}_{T,a}\le\theta^{(n)}_{T,a}$. The same holds for the diminished factor
$s^{(n)}_{T,a}$ attached to $\theta^{(n)}_{T,a}$ by the unit rule. Therefore clauses (c1) and (c2)
at stage $n+1$ imply those at stage $n$.
\item \emph{The set $\mathfrak{N}^{(n)}$.} The clauses of $\mathfrak{N}^{(n)}(Y)$ at
$Z^{\mathrm{on}}$-locations do not depend on $n$. At $\mathfrak{T}$-locations, the factors
$1+2\beta-\sum_{l\le n}\iota^1_l$ of Definition~\ref{4.26:def-regularity-space} decrease in $n$, for
the same reason as in the previous item.
\end{itemize}
Taking the intersection gives $\mathfrak{D}^{\flat}_{n+1}(Y)\subset\mathfrak{D}^{\flat}_n(Y)$.

\emph{Stability under restriction.} Let $Y'\subset Y$ be a union of cubes, and let the pair lie in
$\mathfrak{D}^{\flat}_n(Y)$. The space $\mathfrak{N}^{(n)}$ is stable under restriction by
Definition~\ref{4.26:def-regularity-space}, so the restriction of the pair to $Y'$ lies in
$\mathfrak{N}^{(n)}(Y')$. It remains to treat $\mathfrak{L}_n$ and $\mathfrak{I}_n$: we show that
every clause defining $\mathfrak{L}_n(Y'\cap Z^{\mathrm{on}})$ and $\mathfrak{I}_n(Y'\cap\mathfrak{T})$
is among the clauses defining $\mathfrak{L}_n(Y\cap Z^{\mathrm{on}})$ and
$\mathfrak{I}_n(Y\cap\mathfrak{T})$, read on the restricted pair.
\begin{itemize}
\item \emph{Fewer locations.} The assignment of a location to $Z^{\mathrm{on}}$ or to $\mathfrak{T}$
is by containment, so it is the same whether the location is regarded in $Y'$ or in $Y$. Thus the
locations of $Y'\cap Z^{\mathrm{on}}$ and of $Y'\cap\mathfrak{T}$ are among those of
$Y\cap Z^{\mathrm{on}}$ and of $Y\cap\mathfrak{T}$.
\item \emph{Fewer terms.} A term $T$ with $X_T\subset Y'$ satisfies $X_T\subset Y$.
\item \emph{Fewer localisation domains.} A localisation domain $X'\subset Y'$ satisfies
$X'\subset Y$, and the condition $X_T\subset X'^{\sim2}$ does not refer to $Y$ or $Y'$.
\item \emph{Derived entries.} The derived entries read at a triple on $Y'$ are the restrictions of
the derived entries on $Y$. This is the restriction of $(\Sigma^{\ast})$ to sub-domain triples, a
property of the derived-entry family $(\Sigma^{\ast})$ stated with that family.
\end{itemize}
Hence the restriction of the pair to $Y'$ lies in $\mathfrak{L}_n(Y'\cap Z^{\mathrm{on}})$ and in
$\mathfrak{I}_n(Y'\cap\mathfrak{T})$, and therefore in $\mathfrak{D}^{\flat}_n(Y')$. This is
\eqref{4.26:eq-stage-systems-a}.
\end{proof}

\begin{theorem}[The expansion on the inherited domain]\label{4.26:thm-inherited-expansion}
Assume the following.
\begin{itemize}
\item The hypothesis on the constants under which Lemma~\ref{4.26:lem-stage-expansion} is
  stated.
\item The input hypothesis of the expansion theorem for the stage terms, which is the flat form
  of its zeroth analyticity statement, for the terms
  $T\in\mathcal{T}$. Each space $\mathfrak{U}_T$ satisfies the newborn clause of
  \eqref{4.26:eq-room-table-a} at $\mathsf{c}_e=\theta_S$.
\item The implication hypothesis, asked at $Z^{\mathrm{on}}$-locations only. For every
  $T\in\mathcal{T}$ and every $n'\ge 0$, the lettered clauses $\mathfrak{L}_{n'}$ at the
  $Z^{\mathrm{on}}$-locations imply the clauses of $T$ at those locations, with the unused room
  $\rho_m$. On a component $W$ excluded at this site, $\mathfrak{L}_{n'}$ is read as in
  Definition~\ref{4.26:def-inherited-domain}; this reading rests on the extended form of the
  lemma on arrested clauses.
\end{itemize}
Then, for every $n\in\{1,\dots,N\}$, with $N$ the number of stages at this site, and every term
$\mathbb{C}^{(n)}_k(X,\mathcal{S})$ of the representation
\begin{equation*}
\mathbb{C}^{(n)}_k=\sum_{(X,\mathcal{S})}\mathbb{C}^{(n)}_k(X,\mathcal{S}),
\end{equation*}
the following hold.
\begin{enumerate}
\item[(a)] The function $\mathbb{C}^{(n)}_k(X,\mathcal{S})$ is a stored term whose $A$-slot
  datum lies on $\mathfrak{D}^{\flat}_n(X)\times\mathfrak{D}_{\mathcal{S}}(X)$. It has the
  following four properties.
  \begin{itemize}
  \item It is holomorphic in the pair $(\mathbb{U},\mathbb{J})$ and in the exterior weak
    coordinates.
  \item It is bounded on the real strong boxes.
  \item It is jointly $\mathcal{G}$-invariant.
  \item It is local in $X$.
  \end{itemize}
  The family $\{\mathbb{C}^{(n)}_k(X,\mathcal{S})\}_{(X,\mathcal{S})}$ satisfies
  \begin{equation*}
  \|\mathbb{C}^{(n)}_k\|^{A,\flat}_{(n)}\le C_0^{\sharp},
  \end{equation*}
  and consequently
  \begin{equation*}
  \sum_{\mathcal{S}}e^{\kappa_A|\mathcal{S}|}
  \sup_{\mathfrak{D}^{\flat}_n(X)\times\mathfrak{D}_{\mathcal{S}}(X)}
  \bigl|\mathbb{C}^{(n)}_k(X,\mathcal{S})\bigr|
  \le C_0^{\sharp}e^{-a_m d_m(X)} .
  \end{equation*}
  For terms of the class $\mathcal{K}_{\mathsf{d}}$ the weight $e^{-a_md_m(X)}$ is replaced by
  $e^{-a_m\mathsf{d}_m(X)}$, here and in the weight of the norm $\|\cdot\|^{A}$. This follows
  from the relative form of the large-field bound on which the expansion theorem for the stage
  terms rests.

  The conclusions (a)--(d) of the expansion theorem for the stage terms, and its corollary,
  hold on $\mathfrak{D}^{\flat}_n$. The constant $C_0^{\sharp}$ and the rates are the same,
  and $\vartheta$ is replaced by $\vartheta^{\natural}$ in clause (f).
\item[(b)] Consider a location of $\mathfrak{T}$ at which the lettered clause of Ba{\l}aban's
  space $\tilde{U}^{(n)c,\Sigma\flat}_k$ is unweighted and has $c=1$. At every such location,
  \begin{equation*}
  \mathfrak{D}^{\flat}_n(X)\supseteq\tilde{U}^{(n)c,\Sigma\flat}_k(X)\cap\mathfrak{N}^{(n)}(X).
  \end{equation*}
  Such locations fill all of $\mathfrak{T}$ under the reading of the derived-entry family
  adopted in this chapter, together with the reading of that clause at $m=k$ also adopted in
  this chapter. At a location of $\mathfrak{T}$ where
  the clause carries a weight $w>1$ or the constant $c=3$, the domain $\mathfrak{D}^{\flat}_n$
  imposes instead $\mathfrak{I}_n$ and the composite shell clause of
  \eqref{4.26:eq-regularity-space-a}.
\item[(c)] On $X\cap\mathfrak{T}$ the only constraints are of two kinds:
  \begin{itemize}
  \item those of $\mathfrak{N}^{(n)}$, the composite shell clause of
    \eqref{4.26:eq-regularity-space-a} included;
  \item the diminished term clauses, the box clause \eqref{4.26:eq-inherited-domain-a} among
    them.
  \end{itemize}
  For $y\in X\cap\mathfrak{T}$ the stage increments of \eqref{4.26:eq-family-consumption-a}
  satisfy
  \begin{equation}\label{4.26:eq-inherited-expansion-a}
  \begin{aligned}
  \sum_{l\le n}\iota_l(y)&\le\tfrac12\mathsf{w}_0\beta e^{-n_W(y)}
  &&\bigl(\le\tfrac13\mathsf{w}_0\beta\ \text{on }\Omega_k\bigr),\\
  \sum_{l\le n}\iota^1_l(y)&\le 2\mathsf{w}_0\beta
  &&\bigl(<\tfrac43\mathsf{w}_0\beta\ \text{on }\Omega_k\bigr).
  \end{aligned}
  \end{equation}
  These bounds do not depend on the number $N$ of stages.
\end{enumerate}
\end{theorem}

The qualifier ``with the unused room $\rho_m$'' in the implication hypothesis belongs to that
hypothesis as the expansion theorem for the stage terms states it. At stage $n$ it follows from
the hypothesis at $n'=n-1$ together with the room between the stages $n-1$ and $n$, namely
their gap and their ratio.

\begin{proof}[Proof of Theorem~\ref{4.26:thm-inherited-expansion}: the induction and the
first three steps]
\emph{The induction.} The proof proceeds by induction on the density index $k$ of the site at
which the operation $\mathbf{R}$ acts. At fixed $k$ it proceeds by induction on $n$.

The induction on $k$ is needed because of statement (4)(d) of
\eqref{4.26:eq-older-pieces-domain-a}. That statement uses part~(a) of the present theorem,
together with the corollary on the fat clauses. Both are used only at sites of
density index $k_b<k$, the near part at still older pieces included. At the first site it is
void.

At fixed $k$, the inductive hypothesis at $n-1$ reads: part~(a) holds for every $i\le n-1$. At
$n=1$ it is void. Under this hypothesis we verify the stability hypothesis
\eqref{4.26:eq-stage-systems-b} at stage $n$ for the systems
$\mathfrak{D}_i:=\mathfrak{D}^{\flat}_i$, $i\le n$.

Each $\mathfrak{D}^{\flat}_i$ is a stage-$i$ system in the sense of
Definition~\ref{4.26:def-stage-systems}. Indeed, by Definition~\ref{4.26:def-inherited-domain}
the set $\mathfrak{D}^{\flat}_i(Y)$ is open, $\mathcal{G}$-invariant, contained in
$\mathfrak{N}^{(i)}(Y)$, and stable under restriction.

Once the stability hypothesis at stage $n$ is established, Lemma~\ref{4.26:lem-stage-expansion}
applies with $l=n$. Its remaining hypotheses hold for the following reasons:
\begin{itemize}
\item the input hypothesis and the hypothesis on the constants are assumed;
\item for $i<n$, the terms $\mathbb{C}^{(i)}(S,\mathcal{S}')$ are stored terms with datum on
  $\mathfrak{D}^{\flat}_i(S)\times\mathfrak{D}_{\mathcal{S}'}(S)$ and norm at most $C_0^{\sharp}$,
  by the inductive hypothesis.
\end{itemize}
The norm $\|\cdot\|^{A,\flat}_{(n)}$ is by definition the norm $\|\cdot\|^{A,\mathfrak{D}}_{(n)}$
of this system. Lemma~\ref{4.26:lem-stage-expansion} therefore gives part~(a) at stage $n$.

Let $Y$ be a union of cubes, and let $(\mathbb{U},\mathbb{J})\in\mathfrak{D}^{\flat}_n(Y)$. Let
$\Phi$ be a member of the family~(a) of the stability hypothesis
\eqref{4.26:eq-stage-systems-b}, with $(b,\sigma)$ as there. By
\eqref{4.26:eq-inherited-domain-b},
\begin{equation*}
\mathfrak{D}^{\flat}_{n-1}(Y)=\mathfrak{N}^{(n-1)}(Y)\cap\mathfrak{L}_{n-1}(Y\cap Z^{\mathrm{on}})
\cap\mathfrak{I}_{n-1}(Y\cap\mathfrak{T}).
\end{equation*}
Steps~1, 2 and~3 below show that $\Phi$ lies in $\mathfrak{L}_{n-1}(Y\cap Z^{\mathrm{on}})$, in
$\mathfrak{I}_{n-1}(Y\cap\mathfrak{T})$ and in $\mathfrak{N}^{(n-1)}(Y)$, respectively. We write
\begin{equation*}
\Delta_a:=Q_a(\Phi)(y)-Q_a(\mathbb{U},\mathbb{J})(y),
\end{equation*}
as in Lemma~\ref{4.26:lem-family-consumption}.

\emph{Step 1: the locations of $Z^{\mathrm{on}}$.} As a point of $\mathfrak{D}^{\flat}_n(Y)$,
the pair $(\mathbb{U},\mathbb{J})$ lies in $\mathfrak{L}_n(Y\cap Z^{\mathrm{on}})$ and in
$\mathfrak{N}(Y)$. Apply Lemma~\ref{4.26:lem-family-consumption}(ii) at $l=n$. It applies alike
on the components still in the run at stage $n$ and on the excluded components, and it gives
$\Phi\in\mathfrak{L}_{n-1}(Y\cap Z^{\mathrm{on}})$.

\emph{Step 2: the diminished term clauses on $\mathfrak{T}$.} Let $T\in\mathcal{T}$ with
$X_T\subset Y$, and let $(y,a)\in\operatorname{Cl}(T)$ with $y\in Y\cap\mathfrak{T}$.

We first bound the change of $Q_a$ at $y$. Lemma~\ref{4.26:lem-family-consumption}(i) gives
$\Delta_5=\Delta_8=0$, and $\Delta_a\le\iota_n(y)r_a(y)$ for $a\in\{1,3,6,7\}$. Since
$r_a\le\mathbf{r}^T_a$, this yields $\Delta_a\le\iota^a_n(y)\mathbf{r}^T_a(y)$ for every
$a\notin\{2,4\}$, with $\iota^a_n=\iota_n$ for $a\in\{1,3,6,7\}$ and $\iota^a_n=0$ for
$a\in\{5,8\}$. For $a\in\{2,4\}$, the same part of that lemma gives the
same bound in its frames, since
\begin{equation*}
Q_a(y)<\theta_{T,a}(y)\le 2\mathbf{r}^T_a(y).
\end{equation*}
Hence in all cases
\begin{equation*}
Q_a(\Phi)(y)\le Q_a(\mathbb{U},\mathbb{J})(y)+\iota^a_n(y)\mathbf{r}^T_a(y).
\end{equation*}

Next we compare thresholds. Since $(\mathbb{U},\mathbb{J})\in\mathfrak{I}_n(Y\cap\mathfrak{T})$,
we have $Q_a(\mathbb{U},\mathbb{J})(y)<\theta^{(n)}_{T,a}(y)$. The definitions of the diminished
thresholds in Definition~\ref{4.26:def-inherited-domain} give the following. For
$T\in\mathcal{T}^{\mathbb{B}}$,
\begin{equation*}
\theta^{(n)}_{T,a}+\iota^a_n\mathbf{r}^T_a=\theta^{(n-1)}_{T,a}.
\end{equation*}
For $T\in\mathcal{T}^{\mathbb{E}\mathbb{R}}$ the reserve
$\tfrac13(\theta_{T,a}-\theta_{\mathrm{src},a})$ does not depend on the stage, so
\begin{equation*}
\theta^{(n)}_{T,a}+\iota^{1,a}_n\mathbf{r}^T_a=\theta^{(n-1)}_{T,a},
\end{equation*}
and $\iota^a_n\le\iota^{1,a}_n$ because $\iota_n\le\iota^1_n$; this inequality is used again in
Step~3. In both cases $Q_a(\Phi)(y)<\theta^{(n-1)}_{T,a}(y)$.

\emph{The box clause.} Let $T\in\mathcal{T}^{\mathbb{E}\mathbb{R}}$, $a\in\{2,4\}$ and
$y\in X_T\cap Y\cap\mathfrak{T}$. Let $(p,X',\mathfrak{s})$ be a $(\Sigma^{\ast})$-triple with
$X'\subset Y$ and $X_T\subset X'^{\sim2}$.

We check that the frame proposition for derived entries applies to this triple. That
proposition makes $(p,X',\mathfrak{s})$ a frame of unit set $\mathfrak{B}$. Its relative shell
bound holds for any $(\Sigma^{\ast})$-triple, any derived bond $b\in X'^{\flat}$ and
$e\in\{2,4\}$. The triples to be checked at stage $n-1$ have structures $\mathfrak{s}$ of stages
at most $n-1$ under $\mathfrak{B}=\mathbb{B}^{(n-1)}_k$, and $\mathsf{m}^{\circ}\le\mathsf{m}$
because successive stages only coarsen the structure. Moreover
\begin{equation*}
y\in X_T\subset X'^{\sim2}\subset X'^{\flat}.
\end{equation*}
The first of the frame's conditions on the pair holds, by a clause of the frame proposition
itself or by the convention \eqref{4.26:eq-regularity-space-b}.

Since $(\mathbb{U},\mathbb{J})$ satisfies the box clause \eqref{4.26:eq-inherited-domain-a} at
stage $n$,
\begin{equation*}
Q_a^{(p,X',\mathfrak{s})}(y)<\theta^{(n)}_{T,a}(y)\le 2\mathbf{r}^T_a(y).
\end{equation*}
By the relative shell bound, passing from $(\mathbb{U},\mathbb{J})$ to $\Phi$ changes this entry
by at most
\begin{equation*}
3\pi\,\mathbf{r}^T_a(y)\le\iota_n(y)\mathbf{r}^T_a(y)\le\iota^1_n(y)\mathbf{r}^T_a(y).
\end{equation*}
By the identity for $\mathcal{T}^{\mathbb{E}\mathbb{R}}$ above, the entry at $\Phi$ is therefore
smaller than $\theta^{(n-1)}_{T,a}(y)$. Together with the preceding paragraph this proves
$\Phi\in\mathfrak{I}_{n-1}(Y\cap\mathfrak{T})$.

\emph{Step 3: the fat regularity space.}

\emph{Direct entries at the locations of $\mathfrak{T}$.} Let $y\in Y\cap\mathfrak{T}$ and
$a\in\{1,3,6,7\}$. By \eqref{4.26:eq-regularity-space-a},
\begin{equation*}
Q_a(\mathbb{U},\mathbb{J})(y)<\Bigl(1+2\beta-\sum_{l\le n}\iota^1_l(y)\Bigr)r_a(y).
\end{equation*}
Lemma~\ref{4.26:lem-family-consumption}(i) in the shell units gives $\Delta_a\le\iota_n(y)r_a(y)$.
Since $\iota_n\le\iota^1_n$,
\begin{equation*}
\begin{aligned}
Q_a(\Phi)(y)
&<\Bigl(1+2\beta-\sum_{l\le n}\iota^1_l(y)+\iota_n(y)\Bigr)r_a(y)\\
&\le\Bigl(1+2\beta-\sum_{l\le n-1}\iota^1_l(y)\Bigr)r_a(y).
\end{aligned}
\end{equation*}
The entries $a=5,8$ are not moved, because $\Delta_5=\Delta_8=0$. Hence
$Q_a(\Phi)(y)<(1+2\beta)r_a(y)$ for them as well.

\emph{Direct entries at the locations of $Z^{\mathrm{on}}$.} By Step~1, $\Phi$ satisfies the
lettered clauses $\mathfrak{L}_{n-1}$ at $Y\cap Z^{\mathrm{on}}$. These clauses bound the direct
entries at $x$ by the factor $3\mathsf{w}^{\mathrm{on}}(x)$ times $\mathbf{r}_{a,m(x)}$, and
$3\mathsf{w}^{\mathrm{on}}(x)<5\mathsf{w}^{\mathrm{on}}(x)$. The bound by
$3\mathsf{w}^{\mathrm{on}}(x)$ holds for the following reasons.
\begin{itemize}
\item On the current structure, the product of the source factor $F$, the weight $w$ and
  Ba{\l}aban's constant $c$ satisfies
  \begin{equation*}
  F\cdot w\cdot c\le 1\cdot w_\ast(k)\cdot 3 ,
  \end{equation*}
  where $c=3$ occurs only at $m=k$ \cite{B-CMP122a}. This includes the lifted strata of the run
  itself; at a point of no live piece the weight is at most $w_\ast(k)$.
\item On a live older piece in the core of $Z$ one has $c\equiv1$. The source factor of the
  arrest satisfies $F^{(k)}\le1$, and the floor lemma for the source factor gives
  $F\ge\varphi_{\mathrm{ARR}}>0$. The weight is $\omega(x)$, so
  \begin{equation*}
  F\cdot\omega\cdot c\le\omega(x)\le\mathsf{w}^{\mathrm{on}}(x)<3\mathsf{w}^{\mathrm{on}}(x).
  \end{equation*}
\end{itemize}
Consequently the source space $\tilde{U}^{\mathrm{ARR}}_k$ is contained in $\mathfrak{N}^{(n)}$
at the locations of $Z^{\mathrm{on}}$, clause by clause, for the direct entries, at every age.

\emph{Composite clauses at the locations of $Z^{\mathrm{on}}$.} There the composite clauses are
the family $(\Sigma^{\ast})^{\mathfrak{a}}$ with the same factor. They are those of the lettered
clauses, in two cases.
\begin{itemize}
\item For the lettered clauses of the structure's own run, they are taken at the structure's
  birth radius with the threshold at level $\lambda_{\mathfrak{s}}$, under the reading
  $(\Sigma^{\flat})$.
\item For the lettered clauses of every run of density index larger than that of the
  structure's own run, the present run's clauses $\mathfrak{L}_l$
  included, they are taken at the frozen pair $(m,\omega)$ with the block side
  $\min\{p,\lambda_{\mathfrak{s}}\}$ of the shell radii of \eqref{4.26:eq-regularity-space-a},
  and with the clause's own threshold at $m$.
\end{itemize}
The composite shell clause of \eqref{4.26:eq-regularity-space-a} is quantified at the locations
of $\mathfrak{T}$ only. So no composite clause of $\mathfrak{N}^{(n)}$ is imported at the
locations of $Z^{\mathrm{on}}$.

\emph{The composite shell clause.} Let $(p,X',\mathfrak{s})$ be a $(\Sigma^{\ast})$-triple with
$X'\subset Y$. Let $y'\in X'^{\sim2}\cap\mathfrak{T}\subset X'^{\flat}$ and $e\in\{2,4\}$. Apply
the relative shell bound of the frame proposition for derived entries in the shell unit
$r_e^{(p)}(y')$; the conditions of the frame on the pair hold as in Step~2.

The composite shell clause for $(\mathbb{U},\mathbb{J})$ gives
$Q_e^{(p,X',\mathfrak{s})}(y')<1.4\,r_e^{(p)}(y')$. The increment is therefore at most
\begin{equation*}
(1.4\pi_\times+\pi)r_e^{(p)}(y')\le 2.4\pi\,r_e^{(p)}(y')\le\iota_n(y')r_e^{(p)}(y')
\le\iota^1_n(y')r_e^{(p)}(y').
\end{equation*}

The inequality $2.4\pi\le\iota_n(y')$ rests on \eqref{4.26:eq-family-consumption-c}, since
$2.4\cdot22/2^9=0.103<\tfrac16$; the last inequality is $\iota_n\le\iota^1_n$. Here $\pi$ and
$\iota_n$ carry the same factor
$\mathsf{w}_0\rho_me^{-\frac14\delta d^{\wedge}}$ at every $y'$. This holds because the majorant
$\varkappa_{\mathsf{w}}$ of Lemma~\ref{4.26:lem-family-consumption}(i) and the ratio $\pi(x)$ are
defined globally in $x$.

Adding this increment to the composite shell clause for $(\mathbb{U},\mathbb{J})$ at stage $n$,
\begin{equation*}
Q_e^{(p,X',\mathfrak{s})}(\mathbb{U},\mathbb{J})(y')
<\Bigl(1+2\beta-\sum_{l\le n}\iota^1_l(y')\Bigr)r_e^{(p)}(y'),
\end{equation*}
gives
\begin{equation*}
Q_e^{(p,X',\mathfrak{s})}(\Phi)(y')<\Bigl(1+2\beta-\sum_{l\le n-1}\iota^1_l(y')\Bigr)
r_e^{(p)}(y'),
\end{equation*}
which is the composite shell clause at stage $n-1$. Therefore $\Phi\in\mathfrak{N}^{(n-1)}(Y)$.

The proof continues with Step~4, which completes the verification of the stability hypothesis
\eqref{4.26:eq-stage-systems-b} at stage $n$, and with parts~(b) and~(c) below.
\end{proof}

\begin{lemma}[The inherited-domain expansion: the constituents]\label{4.26:prf-inherited-constituents}
We work in the setting of the proof of Theorem~\ref{4.26:thm-inherited-expansion}. The stage is
$n\in\{1,\dots,N\}$, with $N$ the number of stages in that theorem, and $Y$ is a union of cubes.
The pair $(\mathbb{U},\mathbb{J})$ lies in the inherited domain $\mathfrak{D}^{\flat}_n(Y)$ of
\eqref{4.26:eq-inherited-domain-b}, and the conclusions of the first three steps of that proof hold
for it. For every pair $(b,\sigma)$ of Definition~\ref{4.26:def-stage-systems}, $\Phi$
denotes a member of family (a) of the stability hypothesis \eqref{4.26:eq-stage-systems-b}
of that definition at $l=n$, and $\mathcal{W}_T(t,\tau)$ a member of its
family (b) attached to a term $T$. Then the following hold.
\begin{itemize}
\item[($\alpha$)] For every leaf $T$ of the representation in the $M$-class with $X_T\subset Y$,
namely every piece
$T(X_T,\mathcal{S})$ of type $\mathbb{B}$ and every tail of type $\mathbb{E}$ or $\mathbb{R}$, one has
$\Phi\restriction X_T\in\mathfrak{U}_T$.
\item[($\beta$)] Let $\mathbb{C}^{(i)}(S,\mathcal{S}')$ be a chain term with $i\le n-1$ and
$S\subset Y$. Then $\Phi\restriction S\in\mathfrak{D}^{\flat}_i(S)$. On this set
$\mathbb{C}^{(i)}(S,\mathcal{S}')$ is a stored term, bounded by its share of the norm
$\|\cdot\|^{A,\flat}_{(i)}$, and its slot domain $\mathfrak{D}_{\mathcal{S}'}(S)$ is unchanged. Together
with ($\alpha$), this gives clause (a) of the stability hypothesis.
\item[($\gamma$)] Let $T\in\mathcal{T}_1$ be in the non-near case. At every location
$y\in X_T\cap\mathfrak{T}$ of a clause $(y,a)\in\operatorname{Cl}(T)$ one has
$Q_a(\mathcal{W}_T(t,\tau))(y)<\theta_{T,a}$. At the $Z^{\mathrm{on}}$-locations of
$\operatorname{Cl}(T)$ the clauses of $T$ hold for $\mathcal{W}_T(t,\tau)$, at a smaller radius.
Clause (b) of the stability hypothesis holds.
\item[($\delta$)] Let $T\in\mathcal{T}_1$ be in the near case. Then the union of the families (c) and
(c$'$) of the stability hypothesis is contained in $\mathfrak{U}_T$, and clause (c) holds.
\item[($\varepsilon$)] On the stage-$n$ fluctuation region $R^{(n)}$, including the last $A$-region
$Z\cap\Omega_k\subset\mathfrak{T}$, the $W$-pieces are polynomials in the plaquette variables of
$\Phi\in\mathfrak{N}^{(n-1)}$; the polynomials $\mathbb{P}^{(j)}(g_j,A_j,B_j)$, the functional
$\mathbb{V}''$ and all the kernels are analytic functions of the input pair $(\mathbb{U},\mathbb{J})$
on the domain given by conditions (i)--(iii) of \cite{B-CMP109} with some radii $\alpha_0'$,
$\alpha_1'$, with bounds uniform on that domain, as stated in \cite{B-CMP116}.
\end{itemize}
\end{lemma}

\begin{proof}
We take the five constituents in turn.

\emph{($\alpha$) The leaves.} Let $T$ be a leaf in the $M$-class with $X_T\subset Y$. Each clause of
$T$ sits at a
location either in $\mathfrak{T}$ or in $Z^{\mathrm{on}}$, and we treat the two cases separately.
\begin{itemize}
\item The clauses at locations in $\mathfrak{T}$ hold by the second step of the proof of
Theorem~\ref{4.26:thm-inherited-expansion}.
\item The clauses at locations in $Z^{\mathrm{on}}$ hold by the first step of that proof, combined with
hypothesis (I3) of Theorem~\ref{4.26:thm-inherited-expansion}. That hypothesis is assumed only at
$Z^{\mathrm{on}}$-locations, and it is used only there. It states that, for every $T\in\mathcal{T}$
and every $n'\ge0$, the lettered clauses $\mathfrak{L}_{n'}$ at $Z^{\mathrm{on}}$-locations imply the
clauses of $T$ at those locations, with the unused room $\rho_m$.
\end{itemize}
Since every clause of $T$ holds, $\Phi\restriction X_T\in\mathfrak{U}_T$.

\emph{($\beta$) The chain terms.} Let $\mathbb{C}^{(i)}(S,\mathcal{S}')$ be a chain term with
$i\le n-1$ and $S\subset Y$. By Definition~\ref{4.26:def-stage-systems}, its domain is
$\operatorname{dom}(\mathbb{C}^{(i)}(S,\mathcal{S}'))=\mathfrak{D}_i(S)$, and in the proof of
Theorem~\ref{4.26:thm-inherited-expansion} the stage system is the inherited domain,
$\mathfrak{D}_i=\mathfrak{D}^{\flat}_i$ of \eqref{4.26:eq-inherited-domain-b}. We use four facts:
\begin{itemize}
\item the first three steps of the proof of Theorem~\ref{4.26:thm-inherited-expansion};
\item the monotonicity of the regularity spaces, the lettered clauses and the diminished term clauses
in the stage;
\item the restriction property (res) of Definition~\ref{4.26:def-stage-systems}: for $Y'\subset Y$ one
has $\mathfrak{D}_l(Y)\restriction Y'\subset\mathfrak{D}_l(Y')$;
\item on an excluded region $W$ with arrest stage $n^+$, monotonicity also across $n^+$.
\end{itemize}
The last fact rests on two inputs. First, the sibling arrest-stage sets
$\mathfrak{S}^{(k),\mathrm{arrest}}_i$ descend, by the display of the
arrest construction that defines the source factor of the arrest. Second, for
$i\le n^+<n-1$,
\begin{equation*}
\mathfrak{S}^{(k),\mathrm{arrest}}_i\supseteq\mathfrak{S}^{(k),\mathrm{arrest}}_{n-1},
\end{equation*}
which is the extended form of the statement on the arrest clause of an excluded region. Together
these give
\begin{equation}\label{4.26:eq-inherited-constituents-1}
\begin{split}
\Phi\restriction S&\in\mathfrak{N}^{(n-1)}(S)\cap\mathfrak{L}_{n-1}(S\cap Z^{\mathrm{on}})
\cap\mathfrak{I}_{n-1}(S\cap\mathfrak{T})\\
&\subset\mathfrak{D}^{\flat}_i(S)=\operatorname{dom}(\mathbb{C}^{(i)}(S,\mathcal{S}')),
\end{split}
\end{equation}
where $\mathfrak{D}^{\flat}_i(S)$ is defined in \eqref{4.26:eq-inherited-domain-b}. This is the domain
on which the term is to be read. The proof of Theorem~\ref{4.26:thm-inherited-expansion} proceeds by
induction on the stage, and the induction hypothesis at stage $i<n$ is clause (a) of that theorem at
stage $i$. By it, $\mathbb{C}^{(i)}(S,\mathcal{S}')$ is a stored term on the domain
$\mathfrak{D}^{\flat}_i(S)$ of \eqref{4.26:eq-inherited-constituents-1}, bounded by its share of
$\|\cdot\|^{A,\flat}_{(i)}$. Its slot
domain $\mathfrak{D}_{\mathcal{S}'}(S)$ is left unchanged by the proof family, by clause (iii) of
Lemma~\ref{4.26:lem-family-consumption}.

The scope of this treatment is as follows. The diminished term clauses $\mathfrak{I}_n$ of
Definition~\ref{4.26:def-inherited-domain} quantify only over the terms of the class $\mathcal{T}$,
that is, over the terms of the convergent renormalization expansion of \cite{B-CMP119}. The chain
terms are therefore treated here, in the way just
described. This includes the frozen terms $\mathbb{C}^{(i)}$ of an excluded region $W$, for every
$i$ up to the number of stages performed at its arrest.

On one branch of an arrest $\mathfrak{a}$, the set $\mathfrak{F}(\mathfrak{a})$ of frozen terms
contains further elements, for which we refer to \cite{B-CMP122a}:
\begin{itemize}
\item the change-of-variables terms;
\item the stage-$n^+$ elements of $W$, read at a configuration that the arrest construction fixes as
a function of $\mathbb{U}$, and never integrated.
\end{itemize}
These are not of the form $\mathbb{C}^{(i)}$, and the induction hypothesis does not cover them. They
are the subject of the separate statement on the frozen terms of that branch, and not of the present
clause.

Together with ($\alpha$), \eqref{4.26:eq-inherited-constituents-1} shows that every member $\Phi$ of
the first family lies where the leaves and the chain terms are defined. Hence clause (a) of the
stability hypothesis \eqref{4.26:eq-stage-systems-b} holds at $l=n$.

\emph{($\gamma$) Terms of $\mathcal{T}_1$, the non-near case.} Let $T\in\mathcal{T}_1$ be in the non-near
case, and let $(y,a)\in\operatorname{Cl}(T)$.

First let $y\in\mathfrak{T}$. We apply clause (v) of Lemma~\ref{4.26:lem-family-consumption}. For the
entries $a\in\{2,4\}$ that clause carries the proviso $Q_a<\theta_{T,a}\le(1+\beta)\mathbf{r}^T_a$,
and the lemma is applied under it; it gives the first inequality in the display below. The
diminished term clause $\mathfrak{I}_n$ of
Definition~\ref{4.26:def-inherited-domain} holds at $y$, because $y\in X_T\cap\mathfrak{T}$, and it
gives the second:
\begin{equation}\label{4.26:eq-inherited-constituents-2}
Q_a(\mathcal{W}_T(t,\tau))(y)\le Q_a(y)+\iota^{1,a}_n\mathbf{r}^T_a<\theta_{T,a}.
\end{equation}

Now let $y\in Z^{\mathrm{on}}$. There the clauses of $T$ hold for $\mathcal{W}_T(t,\tau)$; they
follow from the lettered clauses $\mathfrak{L}_n$ satisfied by the input. The link is clause (iv) of
the statement bounding the proof family at $Z^{\mathrm{on}}$-locations, combined with hypothesis (I3)
of Theorem~\ref{4.26:thm-inherited-expansion} and with two statements of the treatment of the floors
and $\gamma$-ceilings of the proof family, the second of them through its clause (c). The clauses so
obtained hold at a smaller radius.

This establishes clause (b) of the stability hypothesis \eqref{4.26:eq-stage-systems-b} at $l=n$.

\emph{($\delta$) Terms of $\mathcal{T}_1$, the near case: hypotheses.} Let $T\in\mathcal{T}_1$ be in the
near case, with cover cube $\square$, its enlargement $\square_0=\tilde{\square}^5$, and top level
$p_{\square}$ on the box. We first check the hypotheses of the near-pair bounds on the two parts of
$\square_0$.

On $\square_0\cap Z^{\mathrm{on}}$, the hypotheses are conditions (i)--(iv) contained in the lettered
clauses $\mathfrak{L}_n$. These hold there, and they are the standing assumption of the near-pair
statement at $Z^{\mathrm{on}}$-locations.

On $\square_0\cap\mathfrak{T}$ we argue condition by condition.
\begin{itemize}
\item Conditions (i)--(iii) are supplied by the regularity space $\mathfrak{N}^{(n)}$ of
\eqref{4.26:eq-regularity-space-a}: it bounds each direct entry by $1+2\beta$ times its radius.
\item Condition (iv) is required in \cite[(1.16)]{B-CMP109} on the cube $\tilde{\square}^3$, two
enlargements short of $\square_0=\tilde{\square}^5$; here $\square_0$ is the cube that
\cite{B-CMP109} describes as ``the cube $\square_0=\tilde{\square}^5$ as in the condition (iv)''.
Condition (iv), imposed literally, is the composite
shell clause $(\mathrm{box})^{\mathrm{sh}}$ of \eqref{4.26:eq-regularity-space-a} at $X'=\square_0$,
where $\square_0\subset Y$. That clause says that, for every $p\le p_{\square}$ and every
$y'\in\tilde{\square}^3\cap\mathfrak{T}$, each composite entry at $(p,y')$ is less than $1+2\beta$
times the unit radius of the box.
\end{itemize}

Condition (iv) does not follow from the direct entries. From conditions (i)--(iii) alone,
\cite{B-CMP109} states, in its own notation, the bound
\begin{equation*}
|\partial U_n(M'(\mathbf{U}))-1|<B_3\,2\alpha_0'L^{-2n}(L^n\xi)^2=2B_3\alpha_0'\xi^2,
\end{equation*}
which is condition (iv) with the constant $2B_3$. That bound places nothing below the threshold of a
term. For this reason we use the composite shell clause itself.

In addition, let $(y,a)\in\operatorname{Cl}(T)$ with $y\in X_T\cap\mathfrak{T}$. There the clause (box)
of \eqref{4.26:eq-inherited-domain-a}, taken at $(p_{\square},\square_0,\mathbb{B}^{(n)}_k)$, bounds
the box-level composite entry of index $a$ of the input by the diminished threshold
$\theta^{(n)}_{T,a}(y)$ of $T$.

\emph{($\delta$) Terms of $\mathcal{T}_1$, the near case: conclusion.} We show that the union of the
families (c) and (c$'$) lies in $\mathfrak{U}_T$.

At the $Z^{\mathrm{on}}$-clauses of $T$, this follows from the lettered clauses $\mathfrak{L}_n$,
through the near-pair statement at $Z^{\mathrm{on}}$-locations recalled above.

At the $\mathfrak{T}$-clauses of $T$, it follows from the near-pair proposition at
$\mathfrak{T}$-locations. Write
$\operatorname{room}_{T,a}:=(\theta_{T,a}-\theta_{\mathrm{src},a})/\mathbf{r}^T_a$ for the room of $T$
at the entry $a$ (a notation used only in this proof): the gap between the threshold of $T$ and the
source threshold, in units of
$\mathbf{r}^T_a$, the unit in which every quantity in the following account is measured (compare the
diminished threshold of a term of $\mathcal{T}^{\mathbb{E}\mathbb{R}}$ in
Definition~\ref{4.26:def-inherited-domain}). The two costs below that are proportional to the room of
$T$ are taken, in the account at the entry $a$, with the room $\operatorname{room}_{T,a}$. The account
has one margin and three costs. In it, as in \eqref{4.26:eq-inherited-constituents-2}, $\iota^1_l$
stands for the increment $\iota^{1,a}_l$ at the entry $a$, which is the increment carried by the
diminished threshold of a term in Definition~\ref{4.26:def-inherited-domain}.
\begin{itemize}
\item The starting point lies below $\theta_{T,a}$ by at least
$\tfrac13\operatorname{room}_{T,a}+\sum_{l\le n}\iota^1_l$. This is the clause (box) just recalled.
\item The passage from \cite[(3.39)]{B-CMP109} to \cite[(3.52)]{B-CMP109} costs at most
$\tfrac16\operatorname{room}_{T,a}$, under the $\gamma$-ceiling relating $\gamma$ to the constants of
\cite{B-CMP109}.
\item The stage layer $\mathbb{H}(\sigma;b)$ costs at most $\iota^1_n$.
\item The ball $B'$ entering the members $P(\varsigma;t,\tau;B')$ of the proof family costs at most
$\tfrac16\operatorname{room}_{T,a}$.
\end{itemize}
The three costs add up to
\begin{equation}\label{4.26:eq-inherited-constituents-3}
\iota^1_n+\tfrac16\operatorname{room}_{T,a}+\tfrac16\operatorname{room}_{T,a}
\le\sum_{l\le n}\iota^1_l+\tfrac13\operatorname{room}_{T,a}.
\end{equation}
The inequality holds because $\iota^1_n$ is one of the summands $\iota^1_l$, $l\le n$. The right-hand
side of \eqref{4.26:eq-inherited-constituents-3} does not exceed the starting margin. Hence the members
of the families (c) and (c$'$) satisfy the $\mathfrak{T}$-clauses of $T$.

Both kinds of clause of $T$ therefore hold, and the union of (c) and (c$'$) lies in $\mathfrak{U}_T$.
This establishes clause (c) of the stability hypothesis \eqref{4.26:eq-stage-systems-b} at $l=n$.

\emph{($\varepsilon$) The fluctuation-scale ingredients.} Consider the stage-$n$ fluctuation region
$R^{(n)}$, including the last $A$-region $Z\cap\Omega_k$, which lies in $\mathfrak{T}$. On it, the
argument for the fluctuation-scale ingredients is that of clause ($\gamma$) of the first
inherited-domain lemma, which applies without change:
\begin{itemize}
\item The $W$-pieces (the letter $W$ here names these pieces of the effective action, not an excluded
region) are polynomials in the plaquette variables of $\Phi\in\mathfrak{N}^{(n-1)}$.
\item The polynomials $\mathbb{P}^{(j)}(g_j,A_j,B_j)$, the
functional $\mathbb{V}''$ and all the kernels are functions of the input pair $(\mathbb{U},\mathbb{J})$.
\end{itemize}
In \cite{B-CMP116} they are ``analytic functions of the configurations $U$, $J$ in a domain given by the
conditions I.(i)-(iii) with some $\alpha_0'$, $\alpha_1'$, and all the bounds above are uniform on the
domain''; there I denotes \cite{B-CMP109}. We use this as stated there. This completes the treatment
of the five constituents.
\end{proof}

\begin{lemma}[Proof of the inherited-domain expansion: comparison with Ba{\l}aban's clauses]
\label{4.26:prf-inherited-comparison}
Under the hypotheses of Theorem~\ref{4.26:thm-inherited-expansion}, clauses (b) and (c) of that
theorem hold.
\end{lemma}

\begin{proof}
Throughout, $y\in X\cap\mathfrak{T}$ and $m=m(y)$ is the shell level of $y$. For an entry $a$ we
write $\theta_{\mathrm{let},a}(y)$ for the threshold at which the lettered clauses of Ba{\l}aban's
space $\tilde{U}^{(n)c,\Sigma\flat}_k(X)$ constrain $Q_a(y)$. The stage rooms are
$\rho^{(l)}_m=\beta\,2^{-|m-\ell_l-1|}$, and the stage increments $\iota_l$, $\iota^1_l$ are those
of \eqref{4.26:eq-family-consumption-a}.

\emph{Clause (b).} Let
$(\mathbb{U},\mathbb{J})\in\tilde{U}^{(n)c,\Sigma\flat}_k(X)\cap\mathfrak{N}^{(n)}(X)$. Let $y$ be
a location at which the lettered clause is unweighted and its constant $c$ equals $1$. At such a
point $\mathfrak{D}^{\flat}_n(X)$ imposes, besides $\mathfrak{N}^{(n)}(X)$, the diminished term
clauses $\mathfrak{I}_n$, the box clause among them (Definition~\ref{4.26:def-inherited-domain},
display \eqref{4.26:eq-inherited-domain-b}). We verify each of these from the lettered clause.

\emph{The lettered threshold.} At $y$ every stage $l\le n$ is of class S, in the sense in which the
stages are classified where the thresholds of the stage clauses are fixed. By the statement of this
chapter that fixes those thresholds, for stages of this class the threshold of the lettered clause
is given exactly by
\begin{equation}\label{4.26:eq-inherited-comparison-1}
  \theta_{\mathrm{let},a}(y)
  =\Bigl(s^{(k)}_m-\sum_{l\le n}\rho^{(l)}_m\Bigr)r_a(y)
  =\theta_{\mathrm{src},a}(y)-\sum_{l\le n}\rho^{(l)}_m\,r_a(y).
\end{equation}

\emph{Diminished clauses, terms with $\ell_T\ge m$.} Let $T\in\mathcal{T}$ with $\ell_T\ge m$. For
such a term $\theta_{T,a}\ge\theta_{\mathrm{src},a}$ (Lemma~\ref{4.26:lem-room-table}) and
$\mathbf{r}^T_a=r_a$. Moreover $\iota_l\le\iota^1_l\le\rho^{(l)}_m$ for every $l\le n$, by
\eqref{4.26:eq-family-consumption-a}. Replacing $\rho^{(l)}_m$ by the smaller
$\iota^1_l$ in \eqref{4.26:eq-inherited-comparison-1}, and then adding the non-negative quantity
$\frac23(\theta_{T,a}-\theta_{\mathrm{src},a})$, we obtain
\begin{equation}\label{4.26:eq-inherited-comparison-2}
  \theta_{\mathrm{let},a}
  \le\theta_{\mathrm{src},a}-\sum_{l\le n}\iota^1_l\,r_a
  \le\theta_{\mathrm{src},a}+\tfrac23(\theta_{T,a}-\theta_{\mathrm{src},a})
     -\sum_{l\le n}\iota^1_l\,r_a .
\end{equation}
We compare the right-hand side with the diminished threshold of
\eqref{4.26:eq-inherited-domain-b}, using $\mathbf{r}^T_a=r_a$; that threshold reads
\begin{align*}
  \theta^{(n)}_{T,a}&=\theta_{T,a}-\tfrac13(\theta_{T,a}-\theta_{\mathrm{src},a})
     -\sum_{l\le n}\iota^1_l\,r_a
  &&\text{if }T\in\mathcal{T}^{\mathbb{E}\mathbb{R}},\\
  \theta^{(n)}_{T,a}&=\theta_{T,a}-\sum_{l\le n}\iota_l\,r_a
  &&\text{if }T\in\mathcal{T}^{\mathbb{B}}.
\end{align*}

If $T\in\mathcal{T}^{\mathbb{E}\mathbb{R}}$, the diminished threshold equals
the right-hand side of \eqref{4.26:eq-inherited-comparison-2}.

If $T\in\mathcal{T}^{\mathbb{B}}$, two facts show that the right-hand side
of \eqref{4.26:eq-inherited-comparison-2} is at most the diminished threshold. First,
$\theta_{\mathrm{src},a}+\frac23(\theta_{T,a}-\theta_{\mathrm{src},a})\le\theta_{T,a}$, because
$\theta_{T,a}\ge\theta_{\mathrm{src},a}$. Second, $\sum_l\iota_l\le\sum_l\iota^1_l$.

In both cases $\theta_{\mathrm{let},a}\le\theta^{(n)}_{T,a}$. Hence
$Q_a(y)<\theta_{\mathrm{let},a}(y)\le\theta^{(n)}_{T,a}(y)$.

\emph{Diminished clauses, terms with $\ell_T<m$.} By the scale-drop line of the positive room of
the inherited domain over the source, \eqref{4.26:eq-inherited-room-a}, we have
$\theta^{(n)}_{T,a}\ge 1.5\,r_a$. Next, $s^{(k)}_m=1-\beta(1-2^{-(k-m)})\le 1$
(Lemma~\ref{4.26:lem-room-table}) and the rooms are non-negative. Hence, by
\eqref{4.26:eq-inherited-comparison-1}, $\theta_{\mathrm{let},a}\le r_a$. Therefore
$\theta^{(n)}_{T,a}\ge1.5\,r_a>\theta_{\mathrm{let},a}$, and the diminished clause follows from the
lettered one.

\emph{The box clause.} Under the reading $(\Sigma^{\ast})$ of the lettered clauses fixed in
Definition~\ref{4.26:def-inherited-domain}, the lettered clause imposes its condition
on every triple with $X'\subseteq X$, at the threshold $\theta_{\mathrm{let},a}$. Moreover
$\theta_{\mathrm{let},a}\le r_a\le\mathbf{r}^T_a$. The comparison of the two preceding paragraphs,
applied triple by triple, therefore yields the box clause of \eqref{4.26:eq-inherited-domain-b}.

\emph{The composite shell clause.} This clause of \eqref{4.26:eq-regularity-space-a} is a clause of
$\mathfrak{N}^{(n)}(X)$, so it stands on both sides of the inclusion. It holds by the assumption
$(\mathbb{U},\mathbb{J})\in\mathfrak{N}^{(n)}(X)$.

The lettered clause alone also gives it. Clause (c), proved below without use of clause (b), gives
$\sum_l\iota^1_l\le2\mathsf{w}_0\beta$. With this bound,
\begin{equation}\label{4.26:eq-inherited-comparison-3}
  \theta_{\mathrm{let},a}\le 1\cdot r^{(p)}_e
  <(1+2\beta-2\mathsf{w}_0\beta)\,r^{(p)}_e=1.383\,r^{(p)}_e .
\end{equation}
Since $\sum_l\iota^1_l\le2\mathsf{w}_0\beta$, the quantity $(1+2\beta-2\mathsf{w}_0\beta)\,r^{(p)}_e$
is at most the threshold of the composite shell clause of \eqref{4.26:eq-regularity-space-a}; hence
$Q_e<\theta_{\mathrm{let},e}$ implies that clause.

All constraints that $\mathfrak{D}^{\flat}_n(X)$ imposes at $y$ therefore hold. At the remaining
locations of $\mathfrak{T}$, where the lettered clause carries a weight larger than one or the
constant $c=3$, clause (b) asserts that $\mathfrak{D}^{\flat}_n(X)$ imposes $\mathfrak{I}_n$ and the
composite shell clause instead, and this is the content of
Definition~\ref{4.26:def-inherited-domain}. This proves clause (b).

\emph{Clause (c).} We prove the bounds \eqref{4.26:eq-inherited-expansion-a} of clause (c). Let $y$
lie in an untouched region $W'$, in shell $m$, and let $l\le n$. In this paragraph $w$ denotes the
lower bound on the $d^{\wedge}$-length of a crossing of one stratum, and $\delta$ the decay rate in
the exponential factor of $\iota_l$ in \eqref{4.26:eq-family-consumption-a}; the letter $w$ here is
not the weight of the lettered clause met in clause (b).
The surviving part of the collar bonds $\mathcal{C}^{(l)}$ lies in $Z$. By the statement of this
chapter that separates the surviving collar bonds from the untouched regions, read in the distance
$d^{\wedge}$, $\mathcal{C}^{(l)}$ and $W'$ lie in distinct components of
$\Lambda^c_k$. Hence every contour from $y$ to $\mathcal{C}^{(l)}$ passes through
$\Lambda_k\subset\Omega_k$ and crosses the strata of $W'$ at the levels $k-1,\dots,m+1$. Read in
$d^{\wedge}$ through the clause of that statement which counts the crossed strata, this gives
\begin{equation}\label{4.26:eq-inherited-comparison-4}
  d^{\wedge}\bigl(y,\mathcal{C}^{(l)}\bigr)\ge w\,n_W(y).
\end{equation}
By the condition on the levels under which the distance $d^{\wedge}$ is set up,
$\frac14\delta w\ge4\ln L+1\ge1$. Hence
$\frac14\delta\,d^{\wedge}(y,\mathcal{C}^{(l)})\ge n_W(y)$.

The rooms satisfy $\sum_l\rho^{(l)}_m\le3\beta$, and $\sum_l\rho^{(l)}_m<2\beta$ at $m=k$, that is,
for $y\in\Omega_k$.

By \eqref{4.26:eq-family-consumption-a} the increments satisfy
\begin{equation}\label{4.26:eq-inherited-comparison-5}
  \iota_l(y)\le\tfrac16\mathsf{w}_0\,\rho^{(l)}_m\,
     e^{-\frac14\delta\,d^{\wedge}(y,\mathcal{C}^{(l)})},
  \qquad
  \iota^1_l(y)\le\tfrac23\mathsf{w}_0\,\rho^{(l)}_m .
\end{equation}
By \eqref{4.26:eq-inherited-comparison-4} and $\frac14\delta w\ge1$, the exponential factor in
\eqref{4.26:eq-inherited-comparison-5} is at most $e^{-n_W(y)}$. Summing
\eqref{4.26:eq-inherited-comparison-5} over $l\le n$, inserting the bounds on the rooms, and using
the products $\frac16\cdot3=\frac12$, $\frac16\cdot2=\frac13$, $\frac23\cdot3=2$ and
$\frac23\cdot2=\frac43$, we obtain
\begin{align*}
  \sum_{l\le n}\iota_l(y)&\le\tfrac12\mathsf{w}_0\beta\,e^{-n_W(y)}
     \quad\bigl(\le\tfrac13\mathsf{w}_0\beta\ \text{on }\Omega_k\bigr),\\
  \sum_{l\le n}\iota^1_l(y)&\le2\mathsf{w}_0\beta
     \quad\bigl(<\tfrac43\mathsf{w}_0\beta\ \text{on }\Omega_k\bigr).
\end{align*}
None of these bounds depends on $N$.
\end{proof}

\begin{remark}
(1) If $m(X)<k$, then $X\subset Z^{\mathrm{on}}$ and
$\mathfrak{D}^{\flat}_n(X)=\mathfrak{N}^{(n)}\cap\mathfrak{L}_n$. The inherited domain then imposes
nothing beyond the lettered clauses, and its enlargement over Ba{\l}aban's space is genuine only for
terms of index $k$.

(2) The hypothesis of Theorem~\ref{4.26:thm-inherited-expansion} by which the lettered clauses
imply the term clauses is used at $Z^{\mathrm{on}}$-locations only. The constant $c$ of the
lettered clause at $m=k$ is never used.

(3) The function $\mathbb{C}^{(n)}_k(X,\mathcal{S})$ carries two bounds. The first is that of its
representation, with smallness parameter $\vartheta$, on the domain of the lettered clauses. The
second is that of Theorem~\ref{4.26:thm-inherited-expansion}, with
$\vartheta^{\natural}=\vartheta/\mathsf{w}_0$, on $\mathfrak{D}^{\flat}_n(X)$.

(4) The following are read exactly as in the representation of the terms
$\mathbb{C}^{(n)}_k(X,\mathcal{S})$ to which Theorem~\ref{4.26:thm-inherited-expansion} applies:
the strong sets,
the constants $\kappa_A$ and $r_A$, the radii $\varrho_j$, the slot sets
$\mathfrak{X}^{\sharp}_j$, the constant $C_{\mathrm{far}}$ and the scaffold rule
$\mathcal{S}(P)$.

(5) The element $e_{\mathrm{str}}$, and every member of the interpolation at $t=0$, are evaluated on
the input itself. They are covered, since
$\mathfrak{D}^{\flat}_n(X)\subset\mathfrak{D}^{\flat}_{n-1}(X)$ and the inclusions of the
interpolation established in the proof of Theorem~\ref{4.26:thm-inherited-expansion} hold at $t=0$.
\end{remark}

\begin{definition}[Ba{\l}aban's near-pair lemma read at the shell scale]\label{4.26:def-near-pair-dictionary}
Let $n$ be a stage and let $(T,\square)$ be a near pair. That is, $T$ is a term with localisation domain $X_T\subset\tilde{\square}^2$, the cube $\square$ is a cover cube of the shell $i$, and $\square_0=\tilde{\square}^5\subset Y$, with $Y$ the domain of Definition~\ref{4.26:def-radii-window}. We assume the hypotheses under which the geometry of a near pair's box, \eqref{4.26:eq-near-box-geometry-b}, is stated. Let $y\in X_T\cap\mathfrak{T}$. We write $m:=m(y)$ for the index of the shell containing $y$, and more generally $m(y')$ for the index of the shell containing $y'$. We write $\ell:=\ell_T$ for the level of $T$. A clause $(y,a)\in\operatorname{Cl}(T)$ with $y\in\mathfrak{T}$ is called a $\mathfrak{T}$-clause. Entries are numbered as in Definition~\ref{4.26:def-clauses-units}.

\begin{enumerate}
\item[(a)] \emph{Dictionary.} The analysis of \cite[(3.37)--(3.52)]{B-CMP109} is read with the following replacements.
\begin{itemize}
\item Ba{\l}aban's level $j$ becomes $\ell:=\ell_T$.
\item His level $k+1$ becomes $p_{\square}$, the top level of the structure
$\mathfrak{s}:=\mathbb{B}^{(n)}_k\restriction\square_0$; on a live older piece of $\mathfrak{T}$ one takes instead
$\mathfrak{s}:=\mathfrak{B}^{\mathfrak{c}}\restriction\square_0$.
\item His sets $\square_0$, $\tilde{\square}^3$ and $X$ become, respectively,
$\square_0=\tilde{\square}^5$, the second enlargement $\square_0^{\sim 2}$ of $\square_0$, and $X_T$, with $X_T\subset\square_0^{\sim 2}$.
\item His unit $\xi$ becomes the level-$\ell$ unit of $T$.
\item His unit $L^{j-1}\eta$ becomes the unit of the local shell, taken cube by cube. This follows the convention of \cite{B-CMP119} that, in all bounds connected with a cover cube, the unit of the current lattice is replaced by the unit of that cube's own level.
\end{itemize}
On $\square_0$ one has
\begin{equation}\label{4.26:eq-near-pair-dictionary-1}
\ell\le m(y')\le p_{\square}\qquad\text{for every } y'\in\square_0 .
\end{equation}

\item[(b)] \emph{The input's size constants.} For each cube of $\tilde{\square}^4$ and each of the entries $5$ and $1$, respectively $6$ and $7$, form the ratio of the bound that the input pair $(\mathbb{U},\mathbb{J})$ satisfies on the entry at that cube to the entry's level-$i$ radius, where $i$ is the shell of $\square$. The constant $\mathfrak{a}_0$ is $\alpha_{0,i}$ times the supremum of these ratios over the cubes of $\tilde{\square}^4$ and the entries $5,1$; the constant $\mathfrak{a}_1$ is $\alpha_{1,i}$ times the supremum over the cubes of $\tilde{\square}^4$ and the entries $6,7$. When the level is to be shown we write $\mathfrak{a}^{(i)}_0,\mathfrak{a}^{(i)}_1$. These are the size constants of the input pair's direct entries on $\tilde{\square}^4$ that enter \cite[(3.26)--(3.27)]{B-CMP109}. Entry $3$ is not included. We put
\begin{equation}\label{4.26:eq-near-pair-dictionary-2}
\mathsf{a}:=L^3(1+2\beta)(1+\beta_0)^4 .
\end{equation}
At every $\mathfrak{T}$-clause, on the graded cover $\mathcal{C}^1$ of \eqref{4.26:eq-graded-cover-a}, one has
\begin{equation}\label{4.26:eq-near-pair-dictionary-3}
\mathfrak{a}_0\le \mathsf{a}\,\alpha_{0,i},\qquad \mathfrak{a}_1\le\mathsf{a}\,\alpha_{1,i}.
\end{equation}
Here the constants and the radii are taken in the unit of the box, at the level $i$ of the shell of $\square$.

\item[(c)] \emph{Units at a $\mathfrak{T}$-clause.} Let $(y,a)\in\operatorname{Cl}(T)$ with $y\in\mathfrak{T}$. The unit at this clause is the clause radius $\mathbf{r}^T_a(y)$. We put
\begin{equation}\label{4.26:eq-near-pair-dictionary-4}
s_T:=\frac{\theta_{T,a}(y)}{\mathbf{r}^T_a(y)},\qquad
\mathrm{room}_a:=\frac{(\theta_{T,a}-\theta_{\mathrm{src},a})(y)}{\mathbf{r}^T_a(y)} ,
\end{equation}
and these satisfy
\begin{equation}\label{4.26:eq-near-pair-dictionary-5}
s_T\in[1-\beta,1+\beta],\qquad \mathrm{room}_a\ge\mathrm{room}_T\ge\beta\hat{c},
\end{equation}
where $\mathrm{room}_T$ denotes the room of the term $T$ in Definition~\ref{4.26:def-radii-window}.
Now let $T$ be a one-scale term. We write
$\mathbf{r}^T_{\partial}:=\mathbf{r}^T_1=\mathbf{r}^T_2$.
We also write $\mathbf{r}^T_J:=\mathbf{r}^T_3=\mathbf{r}^{T,(p)}_4$, at the $(\Sigma^{\ast})$-triples with $p\ge\lambda$, in the notation of the unit rule of Definition~\ref{4.26:def-clauses-units}. At $y$ the thresholds satisfy
\begin{equation}\label{4.26:eq-near-pair-dictionary-6}
\theta_{T,1}=\theta_{T,2},\qquad \theta_{T,3}=\theta_{T,4}.
\end{equation}
\end{enumerate}
\end{definition}

\begin{proof}
We justify the relations recorded in parts (a)--(c) above in turn.

\emph{The bound \eqref{4.26:eq-near-pair-dictionary-1}.} This is clause (a) of the geometry of a near pair's box, \eqref{4.26:eq-near-box-geometry-b}. It applies because the hypotheses of that statement are assumed at the outset and $y\in X_T\cap\mathfrak{T}$.

\emph{Why entry $3$ is excluded in part (b).} Entry $3$ has radius exponent $3$: its radius scales as the inverse cube of the unit length. The bounds \cite[(3.26)--(3.27), p.~275]{B-CMP109} and \cite[(3.37)]{B-CMP109} involve the configuration $\mathbf{U}$ alone. In \cite[(3.34)]{B-CMP109} only the configurations $\mathbf{U}$ and $\mathbf{A}$ occur, and in \cite{B-CMP109} the current $\mathbf{J}$ is only rotated. The hypothesis on $\mathbf{J}$ required of the input pair is met clause by clause by the space $\mathfrak{N}^{(n)}(Y)$ of the current stage. The composite current at $y$ (clause entry $4$) is read where \cite[(3.42)]{B-CMP109} reads the current. Hence only entries $5,1$ and $6,7$ feed $\mathfrak{a}_0$ and $\mathfrak{a}_1$.

\emph{The bound \eqref{4.26:eq-near-pair-dictionary-3}.} With the reading in the unit of the box fixed in part (b), \eqref{4.26:eq-near-pair-dictionary-3} is clause (c) of \eqref{4.26:eq-near-box-geometry-b}. That clause gives the following bound for the constants $\mathfrak{a}^{(i)}_0$ (entries $5,1$) and $\mathfrak{a}^{(i)}_1$ (entries $6,7$) formed at the level $i$ of the shell of $\square$:
\begin{equation*}
\mathfrak{a}^{(i)}_{\cdot}\le(1+2\beta)(1+\beta_0)L^2\,\alpha_{\cdot,i}\le\mathsf{a}\,\alpha_{\cdot,i}.
\end{equation*}
The statement is local: neither the weight $w_\ast$ nor entry $3$ enters it.

The entries $5,1,6,7$ have radius exponents at most $2$. Converting their bounds at a cube of $\tilde{\square}^4$ one level finer than $i$ into the level-$i$ unit therefore costs at most the factor $(1+2\beta)(1+\beta_0)L^2$. Since $L\ge1$ and $\beta_0\ge0$, this factor is at most $\mathsf{a}$ of \eqref{4.26:eq-near-pair-dictionary-2}.

The constant $\mathsf{a}$ also satisfies
\begin{equation*}
\mathsf{a}\ge(1+2\beta)(1+\beta_0)L^3 .
\end{equation*}
This is the factor required, in the core-freezing step, for the conversion of entry $3$, whose radius exponent is $3$, at boxes meeting the lifted part of $Z^{\mathrm{on}}$. The constant $\mathsf{a}$ is a power of $L$ times a function of $\beta,\beta_0$. It is fixed after $L$ and before $\gamma$, and imposes no cap in the sense of Definition~\ref{2.3:def-cap}.

\emph{The relations \eqref{4.26:eq-near-pair-dictionary-5}.} The first relation, $s_T\in[1-\beta,1+\beta]$, holds by the choice of the thresholds $\theta_{T,a}$ made for the terms. The second is the window condition of Definition~\ref{4.26:def-radii-window}: it gives $\mathrm{room}_a\ge\mathrm{room}_T$ and $\mathrm{room}_T\ge\beta\hat{c}$.

\emph{The equalities for a one-scale term.} The equalities $\mathbf{r}^T_1=\mathbf{r}^T_2$ and $\mathbf{r}^T_3=\mathbf{r}^{T,(p)}_4$ at triples with $p\ge\lambda$ are the unit rule of Definition~\ref{4.26:def-clauses-units}.

The equalities \eqref{4.26:eq-near-pair-dictionary-6} at $y$ follow from the choice of the shell radii. That choice has $\mathbf{r}_{1,m}=\mathbf{r}_{2,m}$, and $\mathbf{r}_{3,m}=\mathbf{r}_{4,m}$ at triples with $p=m$. It is consistent with the fact that conditions (iii) and (iv) of \cite{B-CMP109} share the radius $\alpha_0$.
\end{proof}

\begin{lemma}[Geometry of a near pair's box]\label{4.26:lem-near-box-geometry}
Work at a fixed level $k$ in the setting of Definition~\ref{4.26:def-stages-regions}, so that
$h=k-N$. Assume, among the conditions \eqref{4.26:eq-inherited-conditions-a}, the two following:
$R_{\min}\ge 5L$, where $R_{\min}$ is the infimum of $R_j$ over the levels $j$ in
play, and the adjacency condition on the radii $\alpha_{\cdot,j}$. Let $n\le N$ be a stage. Let
$(T,\square)$ be a near pair as in Definition~\ref{4.26:def-near-pair-dictionary}. Thus:
\begin{itemize}
\item $\square$ is a cube of the cover of the shell (stratum) $i$ in the sense of \cite{B-CMP119};
\item $X_T\subset\tilde\square^2$;
\item $\square_0=\tilde\square^5\subset Y$, where $Y$ is the union of cubes of
Definition~\ref{4.26:def-regularity-space} on which $\mathfrak{N}^{(n)}(Y)$ is taken.
\end{itemize}
Let $y\in X_T\cap\mathfrak{T}$, and put $m:=m(y)$ and $\ell:=\ell_T$.

Assume moreover that $\square_0$ meets no lifted point, that is, no point at which the weight
$\omega$ exceeds $1$, of any live arrest $\mathfrak{a}$ with $k_a<k$. These are the
points of the thin strata $\Gamma''_{m'}$ of $\mathfrak{a}$ with $m'>k_0^{a}$, where $k_0^{a}$
is the letter of the live piece of $\mathfrak{a}$ corresponding to the index $k_0$ of
Definition~\ref{4.26:def-stages-regions}, together with the points of the lifted stratum of
$\mathfrak{a}$, which we denote by $\Gamma^{\mathrm{lift}}_{\mathfrak{a}}$:
\begin{equation}\label{4.26:eq-near-box-geometry-a}
\begin{gathered}
\square_0\cap\Bigl(\Gamma^{\mathrm{lift}}_{\mathfrak{a}}
\cup\bigcup_{m'>k_0^{a}}\Gamma''_{m'}\Bigr)=\emptyset\\
\text{for every live arrest }\mathfrak{a}\text{ with }k_a<k.
\end{gathered}
\end{equation}
Then:
\begin{enumerate}
\item[(a)] $\square_0$ meets at most two shells. If it meets two, they are adjacent. The shells
$i$ and $m$ are among those it meets. Further $\square_0\subset\Omega_\ell$, that is,
$\ell\le m(y')$ for every $y'\in\square_0$; moreover $m(y')\le p_{\square}$ there, and $\ell\le i$.
\item[(b)] Suppose $\tilde\square^4\cap Z^{\mathrm{on}}\neq\emptyset$. Then $m=k$,
$y\in\Omega_k\setminus\Lambda_k$ and $\ell\le k-1$. Also
$\tilde\square^4\cap Z^{\mathrm{on}}\subset\Omega_{k-1}\setminus\Omega_k$. On this set the
clauses $\mathfrak{L}_n$, read with coefficient $c\in\{1,3\}$ and weight $w$ in the sense of
\cite[(1.66)--(1.68)]{B-CMP122a}, carry the coefficient $c=1$, and their level and weight are
as follows:
\begin{equation*}
\begin{cases}
\text{level }k-1,\ w=1, & n\le N-1,\\
\text{level }k,\ w=L_0^2, & n=N.
\end{cases}
\end{equation*}
Thus the class of weight larger than $1$ is non-empty exactly at the last stage. The table is
read through the index $h+n$ of the stage, the last stage $n=N$ corresponding to $h+n=k$; on a
frozen component $W$ the same table holds with $h+n$ replaced by the frozen level $j^+$ of $W$.
The cubes of $\tilde\square^4\cap Z^{\mathrm{on}}$ at the last stage $n=N$, whose clause is read
at level $k$ with weight $L_0^2$, are called the \emph{raised cubes}.
\item[(c)] Work in the unit of the box, that is, of level $i$. Let $\mathfrak{a}^{(i)}_0$ be the
size constant of the entries $5$ and $1$, and $\mathfrak{a}^{(i)}_1$ that of the entries $6$ and
$7$, of the clauses of Definition~\ref{4.26:def-regularity-space}. Each is defined as
$\alpha_{0,i}$, respectively $\alpha_{1,i}$, times the supremum over the cubes of
$\tilde\square^4$ of the quotient
\begin{equation*}
\frac{\text{the input's bound on the entry at the cube}}{\text{the level-}i\text{ radius of the entry}}.
\end{equation*}
Then
\begin{equation*}
\mathfrak{a}^{(i)}_{\cdot}\le(1+2\beta)(1+\beta_0)L^2\alpha_{\cdot,i}\le\mathsf{a}\,\alpha_{\cdot,i}.
\end{equation*}
Here $\mathsf{a}$ is the constant of Definition~\ref{4.26:def-near-pair-dictionary}. Of it only
the one-sided bound $\mathsf{a}\ge(1+2\beta)(1+\beta_0)L^3$ is used.
\item[(d)] On the raised cubes, $\mathbf{r}_{a,k}/\mathbf{r}^T_a\le(1+\beta_0)/L$ for $a=6,7$,
where $\mathbf{r}_{a,k}$ is the level-$k$ shell radius of entry $a$. Their share of the
zero-layer share $\varkappa_{\square}$ of the near pair, the quantity bounded through condition
\eqref{4.26:eq-near-pairs-outside-a}, is at most $L_0^2(1+\beta_0)L^{-1}\le 0.39$ times the
unraised value, where $L_0$ is the base of the weights of \cite[(1.66)]{B-CMP122a}. Two
facts enter: the factor $L^{-1}$, which is Ba{\l}aban's $L^{j-1}\eta$ of \cite[(3.37)]{B-CMP109},
and $L_0^2\le L/3$. So the raised cubes require nothing further of condition
\eqref{4.26:eq-near-pairs-outside-a}.
\end{enumerate}
Collecting the inequalities of (a)--(d):
\begin{equation}\label{4.26:eq-near-box-geometry-b}
\begin{gathered}
\ell\le m(y')\le p_{\square}\quad(y'\in\square_0),\qquad \ell\le i,\\
m=k,\quad \ell\le k-1\qquad\text{if }\tilde\square^4\cap Z^{\mathrm{on}}\neq\emptyset,\\
\mathfrak{a}^{(i)}_{\cdot}\le(1+2\beta)(1+\beta_0)L^2\alpha_{\cdot,i}\le\mathsf{a}\,\alpha_{\cdot,i},
\qquad \frac{\mathbf{r}_{a,k}}{\mathbf{r}^T_a}\le\frac{1+\beta_0}{L}\quad(a=6,7).
\end{gathered}
\end{equation}
\end{lemma}

\begin{proof}
Lengths are measured in $\pi_i$-sides, the sides of the cubes of Ba{\l}aban's partition $\pi_i$
of level $i$.

\emph{Counting sides.} The cube $\square$ of the cover of shell $i$ has side $2$. Hence
$\square_0=\tilde\square^5$ has side $12$, $\tilde\square^4$ has side $10$, and $\tilde\square^2$
has side $6$. From a point of $\tilde\square^2$, every point of $\square_0$ lies within $9$
$\pi_i$-sides, and every point of $\tilde\square^4$ within $8$. A $\pi_{m'}$-side equals
$L^{m'-i}$ $\pi_i$-sides. The cover of \cite{B-CMP119} is taken stratum by stratum.

We show $\ell\le i$. The localisation domain $X_T$ is a union of $\pi_\ell$-cubes and lies inside
$\tilde\square^2$, whose side is six $\pi_i$-sides. If $\ell>i$, a $\pi_\ell$-cube would have side
at least $L$ $\pi_i$-sides, and $L>6$; so $\ell\le i$. With Ba{\l}aban's $j$ and $n$ read as
$\ell$ and $i$, this is also the inequality $j\le n\le k$ of \cite{B-CMP119}, which in particular
gives $i\le k$.

\emph{Widths of shells.} Consider a shell $m'$ created by the operation $\mathbf{T}$, either a
shell of the nested regions of $Z$ or one of an untouched component $W'$. It contains
$\Lambda_{m'}\setminus\Omega_{m'+1}$, which is at least eight layers of $MR_{m'+1}$-cubes of the
lattice of level $m'+1$ wide. Indeed, in \cite[(3.2)--(3.5)]{B-CMP119} the set $Z'_k$ is the union
of all $LMR_{k+1}$-cubes meeting $Z_k$. The successive regions are then obtained by adding four
layers, one layer, one layer and two layers of such cubes. These are cubes of the lattice of
spacing $\eta$, that is, $MR_{k+1}$-cubes of the lattice of spacing $L^{-1}\eta$, which gives
eight layers. The count of at least ten layers
in \cite{B-CMP122a}, and $S(Z)=Z'^{\sim10}$ in \cite{B-CMP122b}, also include the two further
layers of the complement of $\Lambda$ in \cite[(3.20)]{B-CMP119}. Hence such a shell $m'$ is at
least $8R_{\min}$ $\pi_{m'+1}$-sides wide.

Now take a live older piece. By \cite[(1.11)]{B-CMP122a},
$Z''_j=(\Omega_j^{\sim5})^c\cap Z$. Its un-lifted shells $\Gamma''_{m'}$ with $m'\le k_0^{a}$ are
created by $\mathbf{R}$. Each is at least one layer of $MR_{m'+1}$-cubes wide: the eight layers
above, less five. For $\Gamma''_{k_0^{a}}$, the innermost un-lifted stratum of the piece, the count
is eight layers less seven, so that again at least one layer remains. Hence these shells are at least
$R_{\min}\ge 5L>12$ $\pi_{m'+1}$-sides wide, and $\Omega''_\ell\supset\Omega_\ell\supset\Lambda_\ell$.

Finally, \cite{B-CMP119} gives that the distance between $\partial\Omega_j$ and $\partial\Lambda_j$
is at least $2R_j$ $\pi_j$-sides. This holds for regions determined by the $j$-th renormalisation
transformation and not by the operation $\mathbf{R}$. The proviso is met at $j=k$; inside $Z$,
and on the last $N$ shells of a live piece, the same separation holds by the class~(ii) property
of \cite{B-CMP122a} (Definition~\ref{4.26:def-stages-regions}).

\emph{Regions healed by an earlier $\mathbf{R}$.} Consider a region healed by an earlier
operation $\mathbf{R}$, that is, one on which $\Lambda_j=\Omega_j$ as a result of that operation,
as in \cite{B-CMP119}. By \cite{B-CMP122b},
the new large-field region equals $\tilde Z_k\cup\bigcup_i Y_i$, where $\tilde Z_k$ and the $Y_i$
are the regions so denoted there; the $Y_i$ are unrelated to the set $Y$ of the statement.
Outside $Z_k$ the action is the
usual small-field action, together with additional boundary terms whose localisation domains
contain one of the domains $Y_i$. Hence a term whose localisation domain meets a healed $Y_i$ is a
boundary term of the family $\mathbb{B}$, and so lies outside
$\mathcal{T}_1\cap\mathcal{T}^{\mathbb{E}\mathbb{R}}$.

Moreover, after a healing at step $\ell'$ every structure has level at least $\ell'$ on
$Z\setminus\Lambda$, by \cite[(1.12)]{B-CMP122a}. The same lower bound, level at least $\ell'$
after a healing at step $\ell'$, holds on $\Lambda$, by \cite{B-CMP122a}, and it holds there also
for the structures of an arrested component. Consequently
$\square_0$ meets no shell $\ell-1$ beside such a region.

\emph{Proof of (a).} By the hypothesis \eqref{4.26:eq-near-box-geometry-a}, every shell met by
$\square_0$ is either a shell created by $\mathbf{T}$ or an un-lifted shell $\Gamma''_{m'}$,
$m'\le k_0^{a}$, of a live older piece; lower bounds for the widths of both kinds were obtained
above.
Suppose $\square_0$ met three shells. Take consecutive shells
$m''-1,m'',m''+1$ met by $\square_0$, with $i$ among them, so that $i\le m''+1$. Then $\square_0$
would cross shell $m''$ within $12$ $\pi_i$-sides. Since $i\le m''+1$, a $\pi_{m''+1}$-side is
$L^{m''+1-i}\ge1$ $\pi_i$-sides. Therefore
\begin{equation*}
12\ \pi_i\text{-sides}\le 12\ \pi_{m''+1}\text{-sides}<R_{\min}\ \pi_{m''+1}\text{-sides},
\end{equation*}
because $R_{\min}\ge5L>12$. This contradicts the width of shell $m''$ found above. Hence
$\square_0$ meets at most two shells, and they are adjacent. It meets shell $i$ because
$\square\subset\square_0$. It meets shell $m$ because
\begin{equation*}
y\in X_T\subset\tilde\square^2\subset\square_0.
\end{equation*}

Next, the localisation domain of a term of level $\ell$ is contained in $\Lambda_\ell$
\cite{B-CMP119}, for the terms $\mathbb{R}_\ell$ as well. So
\begin{equation*}
y\in X_T\subset\Lambda_\ell\subset\Omega_\ell,
\end{equation*}
whence $\ell\le m$. We have shown $\ell\le i$.

Suppose $\square_0$ met shell $\ell-1$. Beside a healed region this was excluded above. Otherwise
$\square_0$ would meet shell $\ell-1$ and the shells $m$ and $i$, both at least $\ell$. Since it
meets at most two adjacent shells, $m=i=\ell$. But $y\in\Lambda_\ell$ lies at distance at least
$2R_\ell$ $\pi_\ell$-sides from $\Omega_\ell^c$, and
\begin{equation*}
2R_\ell\ge2R_{\min}\ge10L>9.
\end{equation*}
Every point of $\square_0$ lies within $9$ $\pi_i$-sides, that is, $9$ $\pi_\ell$-sides, of $y$.
Hence $\square_0\subset\Omega_\ell$, and $\square_0$ does not meet shell $\ell-1$, a
contradiction. Therefore $\square_0\subset\Omega_\ell$, which gives $\ell\le m(y')$ at every
$y'\in\square_0$; and $m(y')\le p_{\square}$ holds by the definition of $p_{\square}$ as the top
level on $\square_0$ (Definition~\ref{4.26:def-near-pair-dictionary}).

\emph{Proof of (b).} Let $\tilde\square^4\cap Z^{\mathrm{on}}\neq\emptyset$. Recall that
$\mathfrak{T}=\Omega_k\cup Z^c$ and $Z^{\mathrm{on}}=Z\cap\Omega_k^c$.

First, $y\in Z^c\setminus\Omega_k$ is impossible. The rim $Z\cap\Omega_k$ holds two layers of
$MR_k$-cubes of the lattice of level $k$, namely the layer
$(\Omega_k^c)^{\sim}\setminus\Omega_k^c$ and one further layer, by
\cite[(3.20)]{B-CMP119}, read for the rim of $Z$ at level $k$. Their width is at
least $2R_k>8$ $\pi_k$-sides. A point of $\tilde\square^4$ lies within $8$ $\pi_i$-sides of $y$,
which is at most $8$ $\pi_k$-sides, since $i\le k$. So no such point can reach $Z^{\mathrm{on}}$
from $Z^c\setminus\Omega_k$.

Hence $y\in\Omega_k$, so $m=k$; recall $i\le k$. Then
\begin{equation*}
\operatorname{dist}(y,\Omega_k^c)\le 8\ \pi_k\text{-sides}<2R_k\ \pi_k\text{-sides},
\end{equation*}
so $y\notin\Lambda_k$. Since $y\in X_T\subset\Lambda_\ell$, it follows that $\ell\ne k$. As
$\ell\le i\le k$, we get $\ell\le k-1$. By (a), $\square_0\supset\tilde\square^4$ meets at most
two adjacent shells, one of which is $m=k$. Since $Z^{\mathrm{on}}\subset\Omega_k^c$, the cubes of
$Z^{\mathrm{on}}$ in $\tilde\square^4$ lie in shell $k-1$, that is, in
$\Omega_{k-1}\setminus\Omega_k$.

We now read the clauses $\mathfrak{L}_n$ on this set from the bounds
\cite[(1.66)--(1.68)]{B-CMP122a}, taken at $j:=h+n$. Thus $n=j-h$, which is the index used in
\cite{B-CMP122a}.
\begin{itemize}
\item For $j\le k-2$, that is $n\le N-2$, the clause is \cite[(1.68)]{B-CMP122a}, with level
$k-1$ and $w=c=1$.
\item For $j=k-1$, that is $n=N-1$, it is \cite[(1.66)]{B-CMP122a} read at Ba{\l}aban's shell
index equal to $j$, with weight $L_0^0=1$.
\item For $j=k$, that is $n=N$, it is \cite[(1.66)]{B-CMP122a} read at Ba{\l}aban's shell index
equal to $k-1$, with the radii $\mathbf{r}_{a,k}$ and weight $L_0^2$.
\end{itemize}
These are the bounds of \cite[(1.66)--(1.67), p.~190]{B-CMP122a}, imposed on the intersection of
$X$ with the shells of Ba{\l}aban's shell index $k_0+1,\dots,j-1,j$, with Ba{\l}aban's $k_0$. The
coefficient $c=3$ occurs only at shell index $k$ \cite{B-CMP122a}. Since
$Z^{\mathrm{on}}\cap\Omega_k=\emptyset$, we have $c=1$ on $\tilde\square^4\cap Z^{\mathrm{on}}$.
This is the table of (b).

\emph{Proof of (c).} By (a), each cube of $\tilde\square^4\subset\square_0$ has level $i-1$, $i$
or $i+1$. For an entry $e$, let $p_e$ be the exponent for which the radius of entry $e$ scales as
the $(-p_e)$-th power of the length unit of its level. We measure each cube against the level-$i$
radius of the entry.
\begin{itemize}
\item At its own level, the quotient is at most $1+2\beta$. This follows from the definition of
$\mathfrak{N}^{(n)}$ (Definition~\ref{4.26:def-regularity-space}); the lettered clauses
$\mathfrak{L}_n$ contribute a factor at most $1$.
\item At one level coarser, the quotient is smaller than this.
\item At one level finer, the quotient is at most $(1+2\beta)(1+\beta_0)L^{p_e}$.
\item At a raised cube, which is at level $k$, the quotient is at most $L_0^2\le L/3$ if $i=k$,
and smaller if $i=k-1$.
\end{itemize}
For the entries converted here, that is, the entries re-expressed here in the level-$i$ unit, the
third quotient is at most $(1+2\beta)(1+\beta_0)L^{p}$, where $p$ is the largest exponent $p_e$
among them; here $p=2$, so it is at most $(1+2\beta)(1+\beta_0)L^2$. Entries $5$, $1$ and
$7$ have $p_e=2$. Entry $6$ has $p_e=1$ and gives $(1+2\beta)(1+\beta_0)L$. Entry $3$, with
$p_e=3$, is not converted here.

Each of the four quotients is therefore at most $(1+2\beta)(1+\beta_0)L^2$: for the first two
because $1+2\beta\le(1+2\beta)(1+\beta_0)L^2$, and for the fourth because
$L/3\le L^2\le(1+2\beta)(1+\beta_0)L^2$. Taking the supremum over the cubes of $\tilde\square^4$
and multiplying by $\alpha_{\cdot,i}$ gives
$\mathfrak{a}^{(i)}_{\cdot}\le(1+2\beta)(1+\beta_0)L^2\alpha_{\cdot,i}$. The second inequality
follows because
\begin{equation*}
\mathsf{a}\ge(1+2\beta)(1+\beta_0)L^3\ge(1+2\beta)(1+\beta_0)L^2.
\end{equation*}
The lower bound on $\mathsf{a}$ carries $L^3$ so that the same constant also covers the
conversion of entry $3$ (exponent $3$) at lifted boxes of $Z^{\mathrm{on}}$ in the core-freezing
step. The constant $\mathsf{a}$ is a power of $L$ times factors in $\beta$ and $\beta_0$; it is
fixed after $L$ and before $\gamma$, and it is subject to no cap in the sense of
Definition~\ref{2.3:def-cap}.

\emph{Proof of (d).} Raised cubes occur only when $\tilde\square^4$ meets $Z^{\mathrm{on}}$, and
there $\ell\le k-1$ by (b). For $a=6,7$ the ratio $\mathbf{r}_{a,k}/\mathbf{r}^T_a$, of the
level-$k$ shell radius to the term's level-$\ell$ radius, is at most
$(\alpha_{1,k}/\alpha_{1,\ell})L^{-(k-\ell)}$. In the
present units the power of $L$ is the factor $L^{j-1}\eta$ of \cite[(3.37)]{B-CMP109}. The
adjacency condition on the radii among \eqref{4.26:eq-inherited-conditions-a}, applied to each of
the $k-\ell\ge1$ pairs of consecutive levels from $\ell$ to $k$, bounds this by
$((1+\beta_0)/L)^{k-\ell}$; since $k-\ell\ge1$ and $(1+\beta_0)/L\le1$,
\begin{equation*}
\frac{\alpha_{1,k}}{\alpha_{1,\ell}}\,L^{-(k-\ell)}\le\frac{1+\beta_0}{L}.
\end{equation*}
The raising multiplies by $L_0^2$. Since $\varkappa_{\square}$ is linear in this ratio, the share
of the raised cubes in $\varkappa_{\square}$, relative to the unraised value, is at most
\begin{equation*}
L_0^2(1+\beta_0)L^{-1}\le\frac{1+\beta_0}{3}\le 0.39,
\end{equation*}
using $L_0^2\le L/3$ for the first inequality and $\beta_0\le 0.17$ for the second.
\end{proof}

\begin{remark}[The near-pair budget in the box's unit]\label{4.26:rem-box-unit-budget}
Let $(T,\square)$ be a near pair as in Lemma~\ref{4.26:lem-near-box-geometry}, with $\ell:=\ell_T$.
Write $i$ for the level of the box $\square$, so that $\ell\le i$, and let $r^{[i]}_a$ denote the
level-$i$ radius. For the cubes of the local shell, \cite{B-CMP119} measures every bound attached to
a cube in the unit of the cube's own level (the sentence quoted in
Definition~\ref{4.26:def-near-pair-dictionary}); accordingly the cube side $O(1)M$ of
\cite{B-CMP119} is measured in the unit of level $i$. Consequently the following hold.
\begin{enumerate}
\item In $\varkappa_{\square}$ (defined in \eqref{4.26:eq-box-unit-budget-1} below), in
$G_{\square}$, $\mathfrak{q}_1$, $\mathfrak{q}_2$, $\mathfrak{q}_1'$, $\varepsilon_{\square}$,
$G'_{\square}$ (defined in \eqref{4.26:eq-near-pair-chain-a}), in
\eqref{4.26:eq-near-pair-chain-b}, and in the implication from the ceiling condition, the first of
the conditions collected in \eqref{4.26:eq-inherited-conditions-a}, to its localised form
\eqref{4.26:eq-near-pair-chain-b}, the letters
$\mathfrak{a}_0$, $\mathfrak{a}_1$, $\alpha_{\cdot,m}$ stand for $\mathfrak{a}^{(i)}_0$,
$\mathfrak{a}^{(i)}_1$, $\alpha_{\cdot,i}$. With these letters the zero-layer share of the near pair
is defined by
\begin{equation}\label{4.26:eq-box-unit-budget-1}
\varkappa_{\square}:=B_3^2\,O(1)M\,\frac{\mathfrak{a}^{(i)}_0}{\alpha_{1,i}}\cdot
\frac{r^{[i]}_a}{\mathbf{r}^T_a}.
\end{equation}
Here $O(1)$ is the number in the cube side $O(1)M$, and $B_3$ is the constant of that name entering
\eqref{4.26:eq-near-pair-chain-a}; the constants $\beta$, $\beta_0$ below are likewise those
entering \eqref{4.26:eq-near-pair-chain-a}.
\item The quantities $\mathfrak{q}_1,\mathfrak{q}_2,\mathfrak{q}_1'$ are expressed in the unit
$\alpha_{0,i}$ times the level-$i$ plaquette unit, and this unit is at most
$\mathbf{r}^T_{\partial}$, the radius in whose unit they are compared in
\eqref{4.26:eq-near-pair-chain-a}.
\item No constant of the chain \eqref{4.26:eq-near-pair-chain-a}--\eqref{4.26:eq-near-pair-chain-b}
changes: $\varkappa_{\square}\le\tfrac12(1-\beta)$, and every majorant collected in the ceiling
condition is linear in $\alpha_{\cdot,i}$, where $\alpha_{\cdot,i}\le\bar{\alpha}$, with
$\bar{\alpha}$ the supremum of the radii $\alpha_{\cdot,m}$ over the levels in play.
\end{enumerate}
\end{remark}

\begin{proof}
In \cite[(3.37)]{B-CMP109} the radius carries the factor $L^{j-1}\eta$, with $j$ the level and
$\eta$ the lattice unit of that display. Under the dictionary of
Definition~\ref{4.26:def-near-pair-dictionary} the level $j$ of \cite{B-CMP109} is read as $\ell$,
and the unit $\eta$ is replaced, cube by cube, by the unit of the cube's own level, here the level
$i$ of $\square$. Read in this way, the factor compares the level-$i$ radius with
$\mathbf{r}^T_a$ and gives
\begin{equation}\label{4.26:eq-box-unit-budget-2}
\frac{r^{[i]}_a}{\mathbf{r}^T_a}\le\Bigl(\frac{1+\beta_0}{L}\Bigr)^{i-\ell}\le 1 .
\end{equation}

For the unit of the $\mathfrak{q}$'s: $\alpha_{0,i}$ times the level-$i$ plaquette unit is at most
$\mathbf{r}^T_{\partial}$, the radius in whose unit the $\mathfrak{q}$'s are compared in
\eqref{4.26:eq-near-pair-chain-a}, since $\ell\le i$.

We bound the second factor in \eqref{4.26:eq-box-unit-budget-1} by one, using
\eqref{4.26:eq-box-unit-budget-2}, and we bound the ratio of the input-size constant to the radius
at level $i$ by
\begin{equation*}
\frac{\mathfrak{a}^{(i)}_0}{\alpha_{1,i}}\le(1+2\beta)(1+\beta_0)L^2\,\frac{C_0}{C_1}.
\end{equation*}
Together these give
\begin{equation}\label{4.26:eq-box-unit-budget-3}
\varkappa_{\square}\le B_3^2\,O(1)M\,(1+2\beta)(1+\beta_0)L^2\,\frac{C_0}{C_1}
\le\tfrac12(1-\beta).
\end{equation}
The last inequality is \eqref{4.26:eq-near-pairs-outside-a}. The middle member of
\eqref{4.26:eq-box-unit-budget-3} does not involve the level $i$.

Finally, every majorant collected in the ceiling condition is linear in $\alpha_{\cdot,i}$, and
$\alpha_{\cdot,i}\le\bar{\alpha}$. Hence that condition, and through it
\eqref{4.26:eq-near-pair-chain-b}, holds with unchanged constants.
\end{proof}

\begin{proposition}[Near pairs at locations outside the run]\label{4.26:prop-near-pairs-outside}
Fix a stage $n$, a region $Y$ as in Definition~\ref{4.26:def-inherited-domain}, and a pair
$(\mathbb{U},\mathbb{J})\in\mathfrak{D}^{\flat}_n(Y)$ of the inherited domain of that definition.
Let $(T,\square)$ be a near pair in the
sense of Definition~\ref{4.26:def-radii-window}, with
$T\in\mathcal{T}_1\cap\mathcal{T}^{\mathbb{E}\mathbb{R}}$ a term of level $\ell$, $X_T\subset
\tilde{\square}^2$, $\square_0=\tilde{\square}^5\subset Y$ and $\square\in\mathcal{C}^1$.
Let $(y,a)\in\operatorname{Cl}(T)$ with $y\in X_T\cap\mathfrak{T}$. Assume:
\begin{itemize}
\item the conditions on the constants under which Lemma~\ref{4.26:lem-family-consumption} is
stated for the proof family, among them the window
condition~\eqref{4.26:eq-radii-window-c} of Definition~\ref{4.26:def-radii-window};
\item the ceiling condition on the input sizes under which the bounds of \cite{B-CMP109} are
applied in this section;
\item $R_{\min}\ge 5L$ and the adjacency condition on the radii, both as assumed in
Lemma~\ref{4.26:lem-near-box-geometry};
\item the inequality
\begin{equation}\label{4.26:eq-near-pairs-outside-a}
\frac{C_1}{C_0}\;\ge\;\frac{2\,\mathsf{a}\,B_3^2\,O(1)\,M}{1-\beta},
\end{equation}
where $\mathsf{a}$ is the constant of Definition~\ref{4.26:def-near-pair-dictionary}, $\beta$ is
the schedule constant of Definition~\ref{4.4:def-post-step-space}, and $B_3$
and $O(1)$ are the constants of \cite[(3.37)]{B-CMP109}, with $B_3=B_3(d,L)$ as in
\cite{B-CMP102b}.
\end{itemize}
Then every member $P=P(\varsigma;t,\tau;B')$ of the family~(c) of
Definition~\ref{4.26:def-radii-window} with
\begin{equation*}
\varsigma\in[0,1],\qquad t\in[0,1],\qquad |\tau|\le\mathsf{w}_0\,r^{\mathrm{nr}}_{\square},
\qquad |B'|<\alpha_3(T),
\end{equation*}
satisfies
\begin{equation}\label{4.26:eq-near-pairs-outside-1}
Q_a(P)(y)<\theta_{T,a}(y),
\end{equation}
and so do the two remainder configurations at the remainder radius of
Definition~\ref{4.26:def-radii-window}. Hence these members and the two remainder configurations
satisfy the clauses of $T$ at locations of $\mathfrak{T}$. Together with the clauses of $T$ at
locations of $Z^{\mathrm{on}}$, which they satisfy by the statement on the clauses of terms at
locations of $Z^{\mathrm{on}}$ for the proof family, this gives that they all lie in
$\mathfrak{U}_T$, which is the conclusion of item~(c) of the hypothesis on the proof family.
\end{proposition}

\begin{proof}
If $\square_0$ does not satisfy condition~\eqref{4.26:eq-near-box-geometry-a} of
Lemma~\ref{4.26:lem-near-box-geometry}, that is, if $\square_0$ meets a lifted point in the sense
of that condition, the assertion is that of the proposition on near pairs whose box meets
lifted points.
From now on $\square_0$ satisfies~\eqref{4.26:eq-near-box-geometry-a}.

Throughout, \cite{B-CMP109} is read at the shell scale as in
Definition~\ref{4.26:def-near-pair-dictionary}: in particular the level $j$ of \cite{B-CMP109} is
$\ell=\ell_T$, its $\mathbf{H}_{k+1}$ is read as $\mathbf{H}^{\square}$, its cubes $\square_0$ and
$\tilde{\square}^3$ and its region $X$ are $\square_0=\tilde{\square}^5$, $\square_0^{\sim2}\supset
X_T$ and $X_T$, its lattice scale $\xi$ is the level-$\ell$ unit of $T$, and its length
$L^{j-1}\eta$ is read, in every bound attached to a cube of the local shell, as the unit of that
shell, cube by cube.
Ba{\l}aban's Taylor parameter, written $\tau$ in \cite[(3.37), (3.50)]{B-CMP109} and kept as
$\tau$ inside the two statements of that paper reproduced below, is the parameter $\varsigma$ of
the member $P$; outside those two statements $\tau$ is the stage-layer parameter of $P$.

\emph{The representative; the entries $a=5$ and $a=8$.}
\cite[(3.37)]{B-CMP109}, on p.~277, states that
\begin{align*}
U_j\bigl(\square_0,\exp i\tau Q(L^{-1}\eta\mathbf{H}_{k+1})\bigr)
&=\Bigl(\exp i\xi\,\mathbf{H}_j\bigl(\square_0,\tau Q(L^{-1}\eta\mathbf{H}_{k+1})\bigr)
\Bigr)^{u_j},
\end{align*}
with
\begin{align*}
&\bigl|\mathbf{H}_j\bigl(\square_0,\tau Q(L^{-1}\eta\mathbf{H}_{k+1})\bigr)\bigr|,\\
&\bigl|\nabla^{\xi}\mathbf{H}_j\bigl(\square_0,\tau Q(L^{-1}\eta\mathbf{H}_{k+1})\bigr)\bigr|
<B_3^2\,O(1)\,M\alpha_0L^{j-1}\eta
\end{align*}
on $\tilde{\square}^3$, and $|u_j-1|<B_3^2\,O(1)\,M\alpha_0$. \cite[(3.50)]{B-CMP109}, on p.~279,
states that
\begin{align*}
\mathbf{H}_j\bigl(\square_0,\tau Q(L^{-1}\eta\mathbf{H}_{k+1})+B'\bigr)
&=\mathbf{H}_j\bigl(\square_0,\tau Q(L^{-1}\eta\mathbf{H}_{k+1})\bigr)\\
&\quad+\int_0^1dt_2\,\Bigl\langle\Bigl(\frac{\delta}{\delta B}\mathbf{H}_j\Bigr)
\bigl(\square_0,\tau Q(L^{-1}\eta\mathbf{H}_{k+1})+t_2B'\bigr),B'\Bigr\rangle\\
&=\mathbf{H}_j\bigl(\square_0,\tau Q(L^{-1}\eta\mathbf{H}_{k+1})\bigr)+\mathbf{A}_2,
\end{align*}
and on p.~280 of \cite{B-CMP109} the remainder is bounded by
$|\mathbf{A}_2|,|\nabla^{\xi}\mathbf{A}_2|\le B_3|B'|<B_3\alpha_3$. Applying these two identities
at $j=\ell$, to the argument $\varsigma B(t,\tau)$ of the member $P$, where $B(t,\tau)$ is
decomposed below into the input part and the change produced by the stage layer, and to its
increment $B'$, one finds that the gauge-field component of $P$ is
\begin{equation}\label{4.26:eq-near-pairs-outside-2}
\Bigl(\exp i\xi\bigl[\mathbf{H}_\ell\bigl(\square_0,\varsigma B(t,\tau)\bigr)
+\mathbf{A}_2\bigr]\Bigr)^{u_\ell}.
\end{equation}
For $T\in\mathcal{T}^{\mathbb{E}\mathbb{R}}$ the space $\mathfrak{U}_T$ is a union of
$G^c$-orbits, by its definition and by \cite[(1.10)]{B-CMP109}; and $u_\ell\in G^c$.
Membership in $\mathfrak{U}_T$ therefore depends
only on the orbit of~\eqref{4.26:eq-near-pairs-outside-2}, and the entries $Q_a(P)$ are read on the
representative
\begin{equation*}
\mathbb{U}_P:=\exp i\xi\bigl[\mathbf{H}_\ell\bigl(\square_0,\varsigma B(t,\tau)\bigr)
+\mathbf{A}_2\bigr],
\end{equation*}
whose $G$-valued part is identically $1$. Consequently $Q_5(P)=0$, and the cube gauges of
$\mathbb{U}_P$ vanish identically. The entries $a=5$ and $a=8$ of $P$ therefore lie below every
positive threshold, in particular below $\theta_{T,a}(y)$, which is positive by the lower
bound~\eqref{4.26:eq-near-pairs-outside-4} below.

\emph{The entries $a=6$ and $a=7$.} These entries are read on the part
$\xi(\mathbf{H}_\ell+\mathbf{A}_2)$ of $\mathbb{U}_P$. Write the argument as
$B(t,\tau)=Q(\mathbf{H}^{\square})+Q'$, where $Q(\mathbf{H}^{\square})$ is the input part,
Ba{\l}aban's $Q(L^{-1}\eta\mathbf{H}_{k+1})$ read at the shell scale, and $Q'$ is the change
produced by the stage layer $\mathbf{A}(t,\tau)$. Here $Q(\cdot)$ and $Q'$ are Ba{\l}aban's
averaging map and its stage-layer correction of \cite[(3.49)]{B-CMP109}, not the entries $Q_a$.
There are three contributions, each measured in
the unit $\mathbf{r}^T_a(y)$, $a\in\{6,7\}$.

\emph{First contribution: the input part.} The size $B_3^2\,O(1)\,M\alpha_0$ in
\cite[(3.37)]{B-CMP109} comes from \cite[(3.26)--(3.27)]{B-CMP109}: \cite[(3.26)]{B-CMP109} gives
$|V(b)-1|<O(1)M\alpha_0$ and $|v-1|<O(1)M\alpha_1$, and \cite[(3.27)]{B-CMP109} bounds the
quantities it estimates on $\tilde{\square}^4$ by $B_3O(1)M\alpha_0$ and gives
$|u_{k+1}-1|<B_3O(1)M\alpha_0$ on $\square_0$.
These are bounds in terms of the direct entries of the input on $\tilde{\square}^4$, whose sizes
are measured by the constant $\mathfrak{a}_0$ of Definition~\ref{4.26:def-near-pair-dictionary}.
Hence the sizes of $\xi\mathbf{H}_\ell(\square_0,\varsigma Q(\mathbf{H}^{\square}))$ and of its
$\nabla^{\xi}$-derivative are at most $B_3^2\,O(1)\,M\mathfrak{a}_0$ times the length unit of the
shell, for every $\varsigma\in[0,1]$: the bound depends on $\varsigma$ only through the size of
the datum $\varsigma Q(\mathbf{H}^{\square})$, which is nondecreasing in $\varsigma$. Measured
against $\alpha_{1,\ell}$ times the length unit of $T$, that is, in the unit $\mathbf{r}^T_a(y)$,
this contribution is at most
\begin{equation}\label{4.26:eq-near-pairs-outside-3}
\varkappa_{\square}:=\frac{B_3^2\,O(1)\,M\mathfrak{a}_0}{\alpha_{1,m}}\cdot
\frac{r_a(y)}{\mathbf{r}^T_a(y)}\;\le\;\frac{B_3^2\,O(1)\,M\mathfrak{a}_0}{\alpha_{1,m}},
\end{equation}
where $m=m(y)$. Indeed, the passage from the radius $\alpha_{1,m}$ and the shell's length unit to
the radius $\alpha_{1,\ell}$ and the length unit of $T$ is exactly the ratio
$r_a(y)/\mathbf{r}^T_a(y)$, and $r_a(y)\le\mathbf{r}^T_a(y)$.

\emph{Second contribution: the stage layer.} \cite[(3.49)]{B-CMP109} states that
\begin{equation*}
|Q'|\le O(1)L^j\eta|\mathbf{A}|<O(1)\alpha_2L^j\eta\quad\text{on }\square_0,
\end{equation*}
and the text following it on p.~279 of \cite{B-CMP109} states that $Q'$ is an analytic and almost
local function of $\mathbf{A}$ and $\mathbf{H}_{k+1}$; through the dictionary,
$|Q'|\le O(1)L^\ell\eta|\mathbf{A}|$. By the bound for near pairs in clause~(v) of
Lemma~\ref{4.26:lem-family-consumption}, which carries the constant $K_{\mathrm{cv}}$ and uses
\cite[(3.48)--(3.52)]{B-CMP109} as stated there, every entry of $P$ moves by at most
$\iota^1_n(y)$, in the unit $\mathbf{r}^T_a(y)$, under the stage layer $\mathbf{A}(t,\tau)$, for
all $|\tau|\le\mathsf{w}_0\,r^{\mathrm{nr}}_{\square}$, the near radius of
Definition~\ref{4.26:def-radii-window}.

\emph{Third contribution: the increment $B'$.} By the bound on p.~280 of \cite{B-CMP109} quoted
above and by the definition of $\alpha_3(T)$ in~\eqref{4.26:eq-radii-window-c},
\begin{equation*}
|\mathbf{A}_2|,\ |\nabla^{\xi}\mathbf{A}_2|\;\le\;B_3|B'|\;<\;B_3\,\alpha_3(T)\;\le\;
\tfrac{1}{12}\,\mathrm{room}_T\,\alpha_{\cdot,\ell},
\end{equation*}
where $\mathrm{room}_T$ is the room of $T$ of~\eqref{4.26:eq-radii-window-c} and
$\alpha_{\cdot,\ell}$ is the level-$\ell$ radius entering $\alpha_3(T)$ in
\eqref{4.26:eq-radii-window-c}, one of $\alpha_{0,\ell}$, $\alpha_{1,\ell}$. In either case
$\alpha_{\cdot,\ell}\le\alpha_{1,\ell}$, so this contribution is at most
$\mathrm{room}_T/12$ in the unit $\mathbf{r}^T_a(y)$.

Adding the three contributions, for $a\in\{6,7\}$,
\begin{equation}\label{4.26:eq-near-pairs-outside-5}
Q_a(P)(y)\;\le\;\bigl[\varkappa_{\square}+\iota^1_n(y)+\tfrac{1}{12}\,\mathrm{room}_T\bigr]
\,\mathbf{r}^T_a(y).
\end{equation}
This is to be compared with the threshold. By the table~\eqref{4.26:eq-room-table-a} of rooms
over the source,
\begin{align*}
\theta_{T,a}&\ge\theta_{\mathrm{src},a}=s^{(k)}_m r_a\ge(1-\beta)\mathbf{r}^T_a
&&\text{if }\ell=m,\\
\theta_{T,a}&\ge(1-\beta)\mathbf{r}^T_a&&\text{if }\ell<m,
\end{align*}
the second line by part~(b) of that table. With $s_T$ as in
Definition~\ref{4.26:def-near-pair-dictionary}, in both cases
\begin{equation}\label{4.26:eq-near-pairs-outside-4}
\theta_{T,a}(y)=s_T\,\mathbf{r}^T_a(y)\;\ge\;(1-\beta)\,\mathbf{r}^T_a(y)\;>\;0.
\end{equation}

The passage from $U_j$ to the configuration $\exp i\xi\mathbf{H}_j$ in \cite{B-CMP109} is made
under a proviso: \cite[p.~277]{B-CMP109} states that if $B_3^2\,O(1)\,M\alpha_0<\frac12\alpha_1$,
then the orbit of the configuration $U_j(\square_0,\dots)$ contains the configuration
$\exp i\xi\mathbf{H}_j(\dots)$ satisfying the conditions (i), (ii) of that paper on the cube
$\tilde{\square}^3$, hence on $X$. In the units of this section the proviso reads
$B_3^2\,O(1)\,M\mathfrak{a}_0\le\frac12(1-\beta)\alpha_{1,m}$. It follows from
\eqref{4.26:eq-near-pairs-outside-a}. Indeed, part~(c) of Lemma~\ref{4.26:lem-near-box-geometry}
(see~\eqref{4.26:eq-near-box-geometry-b}) gives $\mathfrak{a}_0\le\mathsf{a}\,\alpha_{0,m}$, as
stated in Definition~\ref{4.26:def-near-pair-dictionary}. The radii $\alpha_{0,\cdot}$ and
$\alpha_{1,\cdot}$ of \cite[(2.28)]{B-CMP119} (see Lemma~\ref{4.4:lem-twist-stability}) carry the
constants $C_0$ and $C_1$, respectively, and the same power of $\log g_m^{-2}$, so that
$\alpha_{0,m}/\alpha_{1,m}=C_0/C_1$. Hence
\begin{equation}\label{4.26:eq-near-pairs-outside-6}
\begin{split}
\varkappa_{\square}&\le\frac{B_3^2\,O(1)\,M\mathfrak{a}_0}{\alpha_{1,m}}
\le\mathsf{a}\,B_3^2\,O(1)\,M\,\frac{\alpha_{0,m}}{\alpha_{1,m}}\\
&=\mathsf{a}\,B_3^2\,O(1)\,M\,\frac{C_0}{C_1}\le\tfrac12(1-\beta),
\end{split}
\end{equation}
the first inequality being~\eqref{4.26:eq-near-pairs-outside-3} and the last
being~\eqref{4.26:eq-near-pairs-outside-a}. The bound~\eqref{4.26:eq-near-pairs-outside-6} holds at
every clause of $T$ at a location of $\mathfrak{T}$ and for every box of $\mathcal{C}^1$, the
boxes at the rim in the last stage included. For the boxes of this proof this uses parts (b)
and (d) of Lemma~\ref{4.26:lem-near-box-geometry}: by part~(b) the boxes at the rim are those
with $\tilde{\square}^4\cap Z^{\mathrm{on}}\ne\emptyset$, and by part~(d) the raised cubes
contribute at most the fraction of the unraised value of $\varkappa_{\square}$ stated there, so
that the raise requires nothing beyond~\eqref{4.26:eq-near-pairs-outside-a}. For the remaining
boxes of
$\mathcal{C}^1$, which meet lifted points, the same bound is part~(c) of the lemma on the geometry
of a near pair's box meeting lifted points, as used in the proposition on near pairs whose box
meets lifted points.

The inequality~\eqref{4.26:eq-near-pairs-outside-a} is the proviso of \cite[p.~277]{B-CMP109}
at the input sizes of this section. Since
$B_3=B_3(d,L)$ depends on $L$ \cite{B-CMP102b}, it is not absorbed into an absolute constant.
The inequality is a one-sided lower bound on the ratio $C_1/C_0$, which is free in
\cite{B-CMP119}, and it is imposed after $L$ and $M$ and before $\gamma$. No upper bound on
$C_1/C_0$ is imposed, so~\eqref{4.26:eq-near-pairs-outside-a} is not a cap in the sense of
Definition~\ref{2.3:def-cap}.

It remains to bound the bracket in~\eqref{4.26:eq-near-pairs-outside-5}. We use the value
$\beta=\frac15$ of the schedule constant $\beta$,
and $\mathsf{w}_0=1/24$ of Definition~\ref{4.26:def-radii-window}. By
\eqref{4.26:eq-family-consumption-a}, $\iota^1_n\le\frac23\mathsf{w}_0\beta<0.006$. Further,
$\mathrm{room}_T\le\max\{2\beta,1+\beta\}$: with $\mathrm{room}_a$ and $s_T$ as in
Definition~\ref{4.26:def-near-pair-dictionary},
\begin{equation*}
\mathrm{room}_T\le\mathrm{room}_a\le s_T\le1+\beta,
\end{equation*}
the middle inequality because $\theta_{\mathrm{src},a}(y)\ge0$. Hence
$\mathrm{room}_T/12\le\frac{1}{12}(1+\beta)=0.1$; moreover $\frac12(1-\beta)=0.4$ and
$1-\beta=0.8$.
Together with
\eqref{4.26:eq-near-pairs-outside-6} and~\eqref{4.26:eq-near-pairs-outside-4} these give
\begin{equation}\label{4.26:eq-near-pairs-outside-b}
\begin{split}
\varkappa_{\square}+\iota^1_n(y)+\tfrac{1}{12}\,\mathrm{room}_T
&\le\tfrac12(1-\beta)+\iota^1_n(y)+\tfrac{1}{12}\,\mathrm{room}_T\\
&\le0.4+0.006+0.1=0.506<0.8\le s_T.
\end{split}
\end{equation}
Inserting~\eqref{4.26:eq-near-pairs-outside-b} into~\eqref{4.26:eq-near-pairs-outside-5} and
using $\theta_{T,a}(y)=s_T\,\mathbf{r}^T_a(y)$ gives $Q_a(P)(y)<\theta_{T,a}(y)$ for
$a\in\{6,7\}$.

\emph{Conclusion.} The bound~\eqref{4.26:eq-near-pairs-outside-1} has thus been obtained for the
members $P$ of family~(c) at the entries $a\in\{5,8\}$ and $a\in\{6,7\}$. At a location
$y\in X_T\cap\mathfrak{T}$ the clauses of $T$ are the bounds~\eqref{4.26:eq-near-pairs-outside-1}
at the pairs $(y,a)\in\operatorname{Cl}(T)$; a configuration satisfying these bounds at all such
pairs therefore satisfies the clauses of $T$ at locations of $\mathfrak{T}$. The locations of $X_T$
not in $\mathfrak{T}$ lie in $Z^{\mathrm{on}}$, and there the clauses of $T$ hold by the statement
on the clauses of terms at locations of $Z^{\mathrm{on}}$ for the proof family. A configuration
satisfying the clauses of $T$ at all locations lies in $\mathfrak{U}_T$, and this is the conclusion
of item~(c) of the hypothesis on the proof family.
\end{proof}

\begin{lemma}[Proof for near pairs: the chain from the box-level entry]\label{4.26:prf-near-pair-chain}
Work at stage $n$ and assume all hypotheses of Proposition~\ref{4.26:prop-near-pairs-outside}, in
particular the conditions of \eqref{4.26:eq-inherited-conditions-a}. Thus $(\mathbb{U},\mathbb{J})\in
\mathfrak{D}^{\flat}_n(Y)$; $(T,\square)$ is a near pair with
$T\in\mathcal{T}_1\cap\mathcal{T}^{\mathbb{E}\mathbb{R}}$ of level $\ell=\ell_T$, $X_T\subset\tilde{\square}^2$
and $\square_0=\tilde{\square}^5\subset Y$; $\square\in\mathcal{C}^1$ satisfies the scale condition of
Proposition~\ref{4.26:prop-near-pairs-outside}, the case in which that condition fails being treated
in the proposition on near pairs at graded-cover boxes; and
$y\in X_T\cap\mathfrak{T}$. Let
$P=P(\varsigma;t,\tau;B')$ be a member of the family (c) of that proposition, with
$\varsigma,t\in[0,1]$, $|\tau|\le\mathsf{w}_0r^{\mathrm{nr}}_{\square}$ and $|B'|<\alpha_3(T)$. Then, for
every $a$ with $(y,a)\in\operatorname{Cl}(T)$: if $a\in\{1,3\}$, the entry $Q_a(P)$ of $P$ (the
curvature entry for $a=1$, the current entry for $a=3$) satisfies
\begin{equation*}
Q_a(P)(y)<\theta_{T,a}(y);
\end{equation*}
and if $a\in\{2,4\}$, assume moreover that at every triple $(p',X',\mathfrak{s}_T)$ of $T$ with
$p'\le\ell$ and $X'\subseteq X_T$ the conjugator $g=\bar{u}_\ell^{-1}u_\ell$ of
\cite[(3.38)]{B-CMP109} at the full datum of $P$ satisfies $\|g(x)\|\,\|g(x)^{-1}\|\le G'_{\square}$
at every site $x$ of $X'^{\sim2}$, the letters $u_\ell$, $\bar{u}_\ell$, the full datum, the set
$X'^{\sim2}$ and the constant $G'_{\square}$ being those of the proof below; then at every such
triple the composite entry satisfies
\begin{equation*}
Q^{(p')}_a(P)(y)<\theta_{T,a}(y;p'),
\end{equation*}
where $\theta_{T,a}(y;p')=s_T\,\mathbf{r}^{T,(p')}_a(y)$ is the threshold of the clause $a$ of $T$ at
that triple.
\end{lemma}

\begin{proof}
Throughout, the analysis of \cite[\S3]{B-CMP109} is read at the shell scale through
Definition~\ref{4.26:def-near-pair-dictionary}: Ba{\l}aban's levels $j$ and $k+1$ become $\ell$ and
$p_{\square}$, his cube $\square_0$ becomes $\square_0=\tilde{\square}^5$, his unit $\xi$ of the level
$j$ becomes the level-$\ell$ unit of $T$, with $\partial^{\xi}$, $\nabla^{\xi}$ the corresponding
difference operators, and his unit $\eta$ is
replaced, cube by cube, by the unit of the local shell at that cube, following the prescription of
\cite{B-CMP119} for all bounds connected with such a cube. Here
$\mathfrak{s}:=\mathbb{B}^{(n)}_k\restriction\square_0$, and on a live older piece of $\mathfrak{T}$
instead $\mathfrak{s}:=\mathfrak{B}^{\mathfrak{c}}\restriction\square_0$; $p_{\square}$ is the top level
of $\mathfrak{s}$ on $\square_0$. Writing $m(y')$ for the level of the local shell at $y'$, we have
$\ell\le m(y')\le p_{\square}$ at every $y'\in\square_0$ by \eqref{4.26:eq-near-box-geometry-b}. The
constants $\mathfrak{a}_0,\mathfrak{a}_1$ are the size constants of the input on $\tilde{\square}^4$
fixed in Definition~\ref{4.26:def-near-pair-dictionary}; $B_3$, $O(1)$ and $\beta$ are the constants
of \cite[\S3]{B-CMP109}; $\alpha_{0,j'},\alpha_{1,j'}$ are the radii at a level $j'$, used below at
$j'=m:=m(y)$, the level of the local shell at $y$, and at $j'=\ell$; $\alpha_{\cdot,j'}$
stands for either of them; $\alpha_3(T)$ and $\mathrm{room}_T$ are as in
\eqref{4.26:eq-radii-window-c}; $\iota^1_l$ is the stage increment of
\eqref{4.26:eq-family-consumption-a}. We write $\mathbf{H}_\ell(\square_0,\cdot)$, $H_{1,\ell}$,
$\mathbf{A}_{1,\ell}$, $H_\ell$, $D_\ell$, $\mathbf{A}_2$ for Ba{\l}aban's functions and operators
$\mathbf{H}_j(\square_0,\cdot)$, $H_{1,j}$, $\mathbf{A}_{1,j}$, $H_j$, $D_j$, $\mathbf{A}_2$ of
\cite[\S3]{B-CMP109} at $j=\ell$, and $u_\ell$,
$\bar{u}_\ell$ for the gauge transformations $u_j$, $\bar{u}_j$ of \cite[(3.37), (3.38)]{B-CMP109} at
$j=\ell$; we write $Q$ for his datum $Q(L^{-1}\eta\mathbf{H}_{k+1})$ of \cite[(3.39)]{B-CMP109}, with
his level $k+1$ read as $p_{\square}$ and his unit $\eta$ read cube by cube as above (the bare letter
$Q$ always denotes this datum, while
the subscripted $Q_a$ are the clause entries of Definition~\ref{4.26:def-inherited-domain}), $Q'$ for
the stage-layer datum that enters at \cite[(3.48)--(3.49)]{B-CMP109} so read, and
$\varsigma Q+\varsigma Q'+B'$ for the full datum of
$P$. We first treat the curvature entry $a=1$; the entry $a=3$ and the composite entries $a=2,4$
follow.

\emph{The entry at the box level.} Write $U_\ell(\square_0,\cdot)$ for the level-$\ell$ minimiser on
$\square_0$ as a function of its constraint and $M_\ell$ for the level-$\ell$ block averages (the
unsubscripted letter $M$ denotes throughout the constant $M=L^{m'}$), and
$\mathbb{U}^{\mathfrak{s}}_{\square}$ for the box minimiser for the structure $\mathfrak{s}$. At zero
stage layer and $\varsigma=1$, the identity \cite[(3.39)]{B-CMP109}, which is derived there from
\cite[(3.26), (3.27)]{B-CMP109} and an identity of an earlier paper of Ba{\l}aban cited there, read
at the
shell scale, states that the $U$-component of $P$, namely
$U_\ell\bigl(\square_0,M_\ell(\mathbb{U}^{\mathfrak{s}}_{\square})\bigr)$, coincides with
$\mathbb{U}^{\mathfrak{s}}_{\square}$ up to the gauge transformations carried in
\cite[(3.39)]{B-CMP109}.
In words: the level-$\ell$ minimiser constrained by the level-$\ell$ averages of the
$\mathfrak{s}$-constrained minimiser is that minimiser, up to these gauge transformations;
this is the composition of minimisers at nested constraints, since $\ell$ does not exceed the levels
of $\mathfrak{s}$ on $\square_0$. The identity is used as stated in \cite[(3.39)]{B-CMP109}.
Consequently the identity of \cite[(3.41)]{B-CMP109} expresses the curvature entry of the zero-layer
pair at $y$ as the image of the input's composite entry at the box-level triple,
$Q_2^{(p_{\square},\square_0,\mathfrak{s})}(y)$, under an operator whose norm is bounded by the
exponential factors of that display, that is, by
\begin{equation}\label{4.26:eq-near-pair-chain-1}
G_{\square}:=\exp\bigl[(B_3^2+B_3)\,O(1)M\mathfrak{a}_0+O(1)M\mathfrak{a}_1\bigr].
\end{equation}
The current entry is obtained in the same way from $Q_4^{(p_{\square},\square_0,\mathfrak{s})}(y)$ by
the parallel identity and bound \cite[(3.42)]{B-CMP109} for the expression replacing the variable
$J$.

Where \cite[(3.40)]{B-CMP109} uses Ba{\l}aban's condition (iv) on the input, taken at the box, we use
the clause
\eqref{4.26:eq-inherited-domain-a} at the triple $(p_{\square},\square_0,\mathfrak{s})$. It
applies, because $(\mathbb{U},\mathbb{J})\in\mathfrak{D}^{\flat}_n(Y)$, the clause $\mathfrak{I}_n$
of Definition~\ref{4.26:def-inherited-domain} is taken with $X'=\square_0\subset Y$, we have
$X_T\subset\tilde{\square}^2\subset\tilde{\square}^3$, the set on which \cite[(3.41)]{B-CMP109} is
stated, and the structure is $(p,\mathfrak{s})=(p_{\square},\mathbb{B}^{(n)}_k)$, respectively
$(p_{\square},\mathfrak{B}^{\mathfrak{c}})$, so that the level $p$ of the triple is the top level of
$\mathfrak{s}$ on $\square_0$. This gives, for $a'\in\{2,4\}$,
\begin{equation}\label{4.26:eq-near-pair-chain-2}
Q_{a'}^{(p_{\square},\square_0,\mathfrak{s})}(y)<\theta^{(n)}_{T,a'}(y)
=\Bigl[s_T-\tfrac13\mathrm{room}_{a'}-\sum_{l'\le n}\iota^1_{l'}(y)\Bigr]\mathbf{r}^T_{a'}(y).
\end{equation}
Here $\theta_{T,a'}(y)=s_T\mathbf{r}^T_{a'}(y)$, and $\mathrm{room}_{a'}$ is defined by
\begin{equation*}
\mathrm{room}_{a'}\,\mathbf{r}^T_{a'}(y)=\theta_{T,a'}(y)-\theta_{\mathrm{src},a'}(y),
\end{equation*}
so that the right side of \eqref{4.26:eq-near-pair-chain-2} is $\theta^{(n)}_{T,a'}$ of
Definition~\ref{4.26:def-inherited-domain} for $T\in\mathcal{T}^{\mathbb{E}\mathbb{R}}$; we write
$\mathrm{room}_{\partial}:=\mathrm{room}_2$ and $\mathrm{room}_J:=\mathrm{room}_4$. The entries $1$ and
$2$ of $T$ are measured in the unit $\mathbf{r}^T_{\partial}$, the entries $3$ and $4$ in the unit
$\mathbf{r}^T_J$; thus $\mathbf{r}^T_2=\mathbf{r}^T_{\partial}$, $\mathbf{r}^T_4=\mathbf{r}^T_J$,
$\theta_{T,1}=s_T\mathbf{r}^T_{\partial}$ and $\theta_{T,3}=s_T\mathbf{r}^T_J$. Write
\begin{equation*}
\rho_0:=Q^{(p_{\square},\square_0,\mathfrak{s})}_{a'}(y)/\mathbf{r}^T_{a'}(y)
\end{equation*}
for the input's composite entry at the box-level triple, expressed in units of
$\mathbf{r}^T_{a'}$, that is, of $\mathbf{r}^T_{\partial}$ for $a'=2$ and of $\mathbf{r}^T_J$ for $a'=4$,
and call it the start ratio; it is taken with $a'=2$ in the chain for the curvature entry and with
$a'=4$ in the chain for the current entry. By \eqref{4.26:eq-near-pair-chain-2}, $\rho_0$ is strictly
less than the
bracket there; moreover $\rho_0\le1+\beta$. The start ratio is taken at the box level $p_{\square}$
and is expressed in the unit of $T$ whatever the value of $\ell\le m(y)$, with no factor $L^{-2}$; a
start at the level $\ell_T<p_{\square}$ would reach the box level only through
\cite[Proposition~9]{B-CMP109}, applied with the constant $2B_3$ in place of $B_3$.

\emph{The chain for the curvature entry.} We follow \cite[pp.~278--280]{B-CMP109}, where each step
raises the constant by $\beta$; here each such increase is replaced by the quantity it bounds. The
losses $\mathfrak{q}_1,\mathfrak{q}_2,\mathfrak{q}_1'$ below are expressed in units of
$\alpha_{0,m}$ times the plaquette unit of the local shell, and this product is at most
$\mathbf{r}^T_{\partial}$.
\begin{enumerate}
\item Zero layer, $\varsigma=1$ (\cite[(3.41)--(3.42)]{B-CMP109}). By the preceding paragraph the
entry is at most $G_{\square}\cdot\rho_0=\rho_0+(G_{\square}-1)\rho_0$, and
$(G_{\square}-1)\rho_0\le(G_{\square}-1)(1+\beta)$, since $\rho_0\le1+\beta$.
\item Passage to the linear size $|\xi\partial^{\xi}\mathbf{H}_\ell|$
(\cite[(3.43)--(3.44)]{B-CMP109}). The bound \cite[(3.43)]{B-CMP109} is obtained there from the
bounds of \cite[(3.37)]{B-CMP109} and the elementary inequality
$|e^{iz}-1-iz|\le\frac12|z|^2e^{|z|}$, valid for every complex $z$; combined with
\cite[(3.41)]{B-CMP109} it yields
\cite[(3.44)]{B-CMP109}. The loss of this step is
\begin{equation}\label{4.26:eq-near-pair-chain-3}
\mathfrak{q}_1:=8\bigl(B_3^2O(1)M\bigr)^2\mathfrak{a}_0^2\,
e^{4B_3^2O(1)M\mathfrak{a}_0}/\alpha_{0,m}.
\end{equation}
\item Passage to the linear functional $H_{1,\ell}$ (\cite[(3.45)--(3.46)]{B-CMP109}). Ba{\l}aban
uses the decomposition
\begin{equation*}
\mathbf{H}_\ell(\square_0,B)=H_{1,\ell}B+\mathbf{A}_{1,\ell}(B)
-H_\ell D_\ell\bigl(H_{1,\ell}B+\mathbf{A}_{1,\ell}(B)\bigr),
\end{equation*}
in which all terms on the right except the first are at least of second order; this gives
\cite[(3.45)]{B-CMP109} and then \cite[(3.46)]{B-CMP109}. The loss of this step is
\begin{equation}\label{4.26:eq-near-pair-chain-4}
\mathfrak{q}_2:=B_3^3\,O(1)^2M^2\mathfrak{a}_0^2/\alpha_{0,m}.
\end{equation}
\item The parameter $\varsigma\in[0,1]$, which plays the role of Ba{\l}aban's $\tau\in[0,1]$ at this
point. The inequality \cite[(3.46)]{B-CMP109} is linear in $Q$, hence remains valid with $\varsigma Q$
in place of $Q$; the value of the linear functional scales by $\varsigma\le1$, and this step carries
no loss. Reversing \cite[(3.45)]{B-CMP109}, as is done there to obtain \cite[(3.44)]{B-CMP109} with
$\varsigma Q$ (Ba{\l}aban's $\tau Q$), carries
the loss $\mathfrak{q}_2$; reversing \cite[(3.43)]{B-CMP109}, which gives \cite[(3.47)]{B-CMP109},
carries the loss $\mathfrak{q}_1$.
\item The stage layer (\cite[(3.48)--(3.49)]{B-CMP109}) carries the loss $\iota^1_n(y)$, the stage
increment of \eqref{4.26:eq-family-consumption-a}, by the bound on the stage layer of a near pair
established in the part of the proof of Proposition~\ref{4.26:prop-near-pairs-outside} that treats
the stage layer.
\item The variable $B'$ (\cite[(3.51)]{B-CMP109}). There
\begin{equation*}
\begin{gathered}
|\partial^{\xi}\mathbf{H}_\ell(\square_0,\varsigma Q+\varsigma Q'+B')|
\le|\partial^{\xi}\mathbf{H}_\ell(\square_0,\varsigma Q+\varsigma Q')|+|\partial^{\xi}\mathbf{A}_2|,\\
|\mathbf{A}_2|,\ |\nabla^{\xi}\mathbf{A}_2|\le B_3|B'|.
\end{gathered}
\end{equation*}
The loss is $2B_3\alpha_3(T)/\alpha_{\cdot,\ell}$ times a ratio of units that is at most
$1$, and $2B_3\alpha_3(T)/\alpha_{\cdot,\ell}\le\frac16\mathrm{room}_T$ by
\eqref{4.26:eq-radii-window-c}.
\item The exponential (\cite[(3.52)]{B-CMP109}). Ba{\l}aban obtains \cite[(3.52)]{B-CMP109} from
\cite[(3.51)]{B-CMP109} and the inequality \cite[(3.43)]{B-CMP109} modified for the present
situation. The loss of this step is
\begin{equation}\label{4.26:eq-near-pair-chain-5}
\mathfrak{q}_1':=8\varepsilon_{\square}^2e^{4\varepsilon_{\square}}/\alpha_{0,m},\qquad
\varepsilon_{\square}:=B_3^2O(1)M\mathfrak{a}_0+B_3\bigl(\alpha_3(T)+\sup|\varsigma Q'|\bigr),
\end{equation}
that is, $\mathfrak{q}_1$ taken at the size which carries $\mathbf{A}_2$ and the stage layer. Indeed
$B_3\alpha_3(T)\le\alpha_{1,m}/3$, so that the member $B_3\alpha_3(T)$ of $\varepsilon_{\square}$ is of
the size $\alpha_{1,m}$ which carries the stage layer; and by the first member of the definition of
$\alpha_3(T)$ in \eqref{4.26:eq-radii-window-c}, $B_3\alpha_3(T)\le0.1\,\alpha_{0,\ell}$. The function
$\mathbf{A}_2$ is measured in the level-$\ell$ unit of $T$. Put $\varepsilon_0:=B_3^2O(1)M\mathfrak{a}_0$
and $\varpi:=\alpha_{0,m}/\alpha_{0,\ell}\ge1$. The loss of \cite[(3.52)]{B-CMP109} in this reading is
$8\bigl(\varepsilon_0L^{-(m-\ell)}+B_3\alpha_3(T)\bigr)^2/\alpha_{0,\ell}$; after conversion to the unit
of level $m$ by the factor $\varpi$ it is still majorised by $\mathfrak{q}_1'$, since
$\varpi L^{-(m-\ell)}\le1$ and, as $\varepsilon_0\ge\alpha_{0,m}$,
\begin{equation*}
\varpi\bigl(B_3\alpha_3(T)\bigr)^2\le0.01\,\alpha_{0,m}\alpha_{0,\ell}
\le0.01\,\alpha_{0,m}^2\le0.01\,\varepsilon_0^2 .
\end{equation*}
Finally $B_3|\varsigma Q'|\le B_3O(1)\alpha_{\cdot,m}$ by \cite[(3.49)]{B-CMP109}.
\end{enumerate}

\emph{Summation.} Adding the contributions of the steps above, with $\rho_0$ in units of
$\mathbf{r}^T_{\partial}$,
\begin{equation}\label{4.26:eq-near-pair-chain-a}
Q_1(P)(y)\le\Bigl\{\rho_0+(G_{\square}-1)(1+\beta)+2\mathfrak{q}_1+2\mathfrak{q}_2
+\mathfrak{q}_1'+\iota^1_n(y)+\tfrac16\mathrm{room}_T\Bigr\}\,\mathbf{r}^T_{\partial}.
\end{equation}
With $G'_{\square}$ defined in the second line, consider the inequality
\begin{equation}\label{4.26:eq-near-pair-chain-b}
\begin{gathered}
(G_{\square}G'_{\square}-1)(1+\beta)+G'_{\square}\bigl(2\mathfrak{q}_1+2\mathfrak{q}_2
+\mathfrak{q}_1'\bigr)+(G'_{\square}-1)\le\tfrac16\beta\hat{c},\\
G'_{\square}:=\Bigl(\frac{1+c_{\mathrm{aux}}\varepsilon_{\square}}
{1-c_{\mathrm{aux}}\varepsilon_{\square}}\Bigr)^2 ,
\end{gathered}
\end{equation}
where $c_{\mathrm{aux}}\ge1$ is the constant of the proposition on the conjugator of
\cite[(3.38)]{B-CMP109} recalled below; it depends on $(d,N,L,M,B_3)$ only and is fixed after
$(N,L,M)$ and before $\gamma$.
Here $G'_{\square}\ge e^{4c_{\mathrm{aux}}\varepsilon_{\square}}\ge e^{4\varepsilon_{\square}}$, and
$G'_{\square}$ is the bound on $\|g(x)\|\,\|g(x)^{-1}\|$, hence on the operator norm of
$\operatorname{Ad}_{g(x)}$, for that conjugator $g$, used for the composite entries below;
the exponential $e^{4c_{\mathrm{aux}}\varepsilon_{\square}}$ alone does not give that bound, because
the bound for the
inverse carries the whole tail of $1/(1-c_{\mathrm{aux}}\varepsilon_{\square})$. The last member of the
left side of
\eqref{4.26:eq-near-pair-chain-b} covers the action of $G'_{\square}$ on
$\iota^1_n+\frac16\mathrm{room}_T\le1$. Since $\hat{c}=\min\{\frac12,\mathsf{c}_e\}$ and
$\mathsf{c}_e=\theta_S$, one has $\frac16\beta\hat{c}=\beta/12$ when $\theta_S\ge\frac12$. The
inequality \eqref{4.26:eq-near-pair-chain-b} is not a further condition: it is implied, at every cube
$\square$, every level $m$ and every $k$, by the ceiling on $\gamma$ that stands first among the
conditions \eqref{4.26:eq-inherited-conditions-a}; this implication is part of the lemma
establishing that ceiling.

Since $G'_{\square}\ge1$, \eqref{4.26:eq-near-pair-chain-b} gives in particular
\begin{equation*}
(G_{\square}-1)(1+\beta)+2\mathfrak{q}_1+2\mathfrak{q}_2+\mathfrak{q}_1'\le\tfrac16\beta\hat{c}.
\end{equation*}
Inserting this and the strict bound \eqref{4.26:eq-near-pair-chain-2} on the start ratio into
\eqref{4.26:eq-near-pair-chain-a}, and then using $\beta\hat{c}\le\mathrm{room}_T\le\mathrm{room}_{\partial}$
(by \eqref{4.26:eq-radii-window-c}) and $\iota^1_n\le\sum_{l'\le n}\iota^1_{l'}$, we obtain
\begin{equation}\label{4.26:eq-near-pair-chain-6}
\begin{aligned}
Q_1(P)(y)&<\Bigl[s_T-\tfrac13\mathrm{room}_{\partial}-\sum_{l'\le n}\iota^1_{l'}
+\tfrac16\beta\hat{c}+\iota^1_n+\tfrac16\mathrm{room}_T\Bigr]\mathbf{r}^T_{\partial}\\
&\le\Bigl[s_T-\tfrac13\mathrm{room}_{\partial}+\tfrac16\mathrm{room}_{\partial}
+\tfrac16\mathrm{room}_{\partial}\Bigr]\mathbf{r}^T_{\partial}
=\theta_{T,1}(y).
\end{aligned}
\end{equation}
The inequality is strict because \eqref{4.26:eq-inherited-domain-a} is a strict inequality.

\emph{The current entry.} The entry $a=3$ is obtained from the start ratio of
\eqref{4.26:eq-near-pair-chain-2} with $a'=4$ by the parallel chain of \cite[pp.~279--280]{B-CMP109}:
starting from \cite[(3.42)]{B-CMP109}, formulas and bounds of an earlier paper of Ba{\l}aban cited
there give a bound of the type
\cite[(3.44)]{B-CMP109} for the second-order operator $\partial^{\xi\ast}\partial^{\xi}$, and the
reasoning of \cite[(3.44)--(3.47)]{B-CMP109} then gives the second inequality of Ba{\l}aban's
condition (iii), that is, the bound for the current entry; the step with $B'$ is
treated in the same way in \cite[p.~280]{B-CMP109}. This chain is used as stated in \cite{B-CMP109},
with the same losses as above, and yields
$Q_3(P)(y)<\theta_{T,3}(y)$.

\emph{The composite entries.} Let $(p',X',\mathfrak{s}_T)$ be a triple of $T$ with $p'\le\ell$ and
$X'\subseteq X_T$. By the identity \cite[(3.38)]{B-CMP109}, the composite pair of $P$ is $P$ itself up
to the gauge transformation $u_\ell\bar{u}_\ell^{-1}$, where $\bar{u}_\ell$ is constant on the blocks
connected with the averaging at the triple and equal to $u_\ell$ at their centres; the function
$\mathbf{H}_\ell$ is the same throughout, and no chain is run at the level $p'$. Let
$g:=\bar{u}_\ell^{-1}u_\ell$ be the conjugator of \cite[(3.38)]{B-CMP109} at the full datum
$\varsigma Q+\varsigma Q'+B'$; it is related to that gauge transformation by
$g=\bar{u}_\ell^{-1}(u_\ell\bar{u}_\ell^{-1})\bar{u}_\ell$, and $g=1$ at the block centres. We use the
following proposition,
whose proof is incomplete at one step, and whose conclusion is the additional hypothesis of the
lemma for $a\in\{2,4\}$: at every triple $(p',X',\mathfrak{s}_T)$ of $T$ with
$p'\le\ell$ and $X'\subseteq X_T$, for every member of the family (c) and for the two remainder
configurations (c$'$) of Proposition~\ref{4.26:prop-near-pairs-outside}, and at every site $x$ of the
set $X'^{\sim2}$, obtained from $X'$ as $\tilde{\square}^3=\square_0^{\sim2}$ is obtained from
$\square_0=\tilde{\square}^5$ in Definition~\ref{4.26:def-near-pair-dictionary},
\begin{equation*}
\|g(x)\|\,\|g(x)^{-1}\|\le G'_{\square}.
\end{equation*}
Since $\|\operatorname{Ad}_h\|\le\|h\|\,\|h^{-1}\|$ for every invertible matrix $h$, in absolute units,
site by site on $X'^{\sim2}$,
\begin{equation*}
Q^{(p')}_{2}(P)(x)\le\bigl\|\operatorname{Ad}_{g(x)}\bigr\|\,Q_1(P)(x)
\le\|g(x)\|\,\|g(x)^{-1}\|\,Q_1(P)(x)\le G'_{\square}\,Q_1(P)(x),
\end{equation*}
and likewise $Q^{(p')}_4(P)(x)\le G'_{\square}Q_3(P)(x)$.
\cite[(3.37)]{B-CMP109} states the bound
$|u_j-1|<B_3^2O(1)M\alpha_0$ in the case $\mathbf{A}=0$ and for the datum $\tau Q$ of
\cite[(3.37)]{B-CMP109}, Ba{\l}aban's $\tau$ being the present $\varsigma$; no bound on $u_\ell$
itself at the full datum is used in this proof.

The passage from $Q_1,Q_3$ to $Q^{(p')}_2,Q^{(p')}_4$ is a further contribution, the one behind the
better constant with which \cite[p.~278]{B-CMP109} deduces condition (iv) from condition (iii)
through \cite[(3.38)]{B-CMP109} and the bound on $u_j$ in \cite[(3.37)]{B-CMP109}; in
\cite[(3.52)]{B-CMP109} it is absorbed through $(1+8\beta)L^{-2}\le1$, and here through
\eqref{4.26:eq-near-pair-chain-b}. Indeed, multiply the bracket of \eqref{4.26:eq-near-pair-chain-a}
(for $a=1$, and its counterpart for $a=3$) by $G'_{\square}$. Since $\rho_0\le1+\beta$ and
$\iota^1_n+\frac16\mathrm{room}_T\le1$,
\begin{align*}
(G'_{\square}-1)\rho_0+G'_{\square}(G_{\square}-1)(1+\beta)
&\le(G_{\square}G'_{\square}-1)(1+\beta),\\
G'_{\square}\bigl(\iota^1_n+\tfrac16\mathrm{room}_T\bigr)
&\le(G'_{\square}-1)+\iota^1_n+\tfrac16\mathrm{room}_T.
\end{align*}
Therefore
\begin{equation}\label{4.26:eq-near-pair-chain-7}
\begin{aligned}
G'_{\square}\cdot\{\cdots\}
&\le\rho_0+(G_{\square}G'_{\square}-1)(1+\beta)
+G'_{\square}\bigl(2\mathfrak{q}_1+2\mathfrak{q}_2+\mathfrak{q}_1'\bigr)\\
&\qquad+(G'_{\square}-1)+\iota^1_n+\tfrac16\mathrm{room}_T,
\end{aligned}
\end{equation}
where $\{\cdots\}$ is the bracket of \eqref{4.26:eq-near-pair-chain-a}. By
\eqref{4.26:eq-near-pair-chain-b} and the estimate leading to \eqref{4.26:eq-near-pair-chain-6}, the
right side, multiplied
by $\mathbf{r}^T_{\partial}$ (respectively $\mathbf{r}^T_J$), is strictly less than
$\theta_{T,1}(y)=s_T\mathbf{r}^T_{\partial}(y)$ (respectively $\theta_{T,3}(y)=s_T\mathbf{r}^T_J(y)$).
By the unit rule \eqref{4.26:eq-clauses-units-b}, $\mathbf{r}^{T,(p')}_2\ge\mathbf{r}^T_{\partial}$ and
$\mathbf{r}^{T,(p')}_4\ge\mathbf{r}^T_J$, so that
$s_T\mathbf{r}^T_{\partial}\le\theta_{T,2}(y;p')$ and $s_T\mathbf{r}^T_J\le\theta_{T,4}(y;p')$. This
proves $Q^{(p')}_a(P)(y)<\theta_{T,a}(y;p')$ for $a\in\{2,4\}$.

\emph{Localisation of the inputs.} The clause \eqref{4.26:eq-inherited-domain-a} is taken at
$y\in X_T$, whereas
\cite[(3.40)]{B-CMP109} is stated on $\tilde{\square}^3$. The bound \cite[(3.41)]{B-CMP109} follows
from \cite[(3.39), (3.40)]{B-CMP109} plaquette by plaquette; the bounds \cite[(3.43), (3.45)]{B-CMP109}
use the sizes of
\cite[(3.37)]{B-CMP109}, that is, the input on $\tilde{\square}^4$ through $\mathfrak{a}_0,
\mathfrak{a}_1$; and the clause \eqref{4.26:eq-inherited-domain-a}, read at the shell scale, holds at
every point of
$\tilde{\square}^3\cap\mathfrak{T}$, as
part of the hypothesis $(\mathbb{U},\mathbb{J})\in\mathfrak{D}^{\flat}_n(Y)$.

\emph{The statements used from \cite{B-CMP109}.} The identities \cite[(3.37)--(3.39), (3.48),
(3.50)]{B-CMP109}, the operator bounds behind \cite[(3.26)--(3.27)]{B-CMP109}, the bounds
\cite[(3.43), (3.45), (3.49)]{B-CMP109}, the bound $|\mathbf{A}_2|,|\nabla^{\xi}\mathbf{A}_2|\le
B_3|B'|$, and the chain for
the current variable are used as stated in
\cite[\S3]{B-CMP109}, read at the shell scale as prescribed in \cite{B-CMP119}; the unit-conversion
constant $K_{\mathrm{cv}}$ of \eqref{4.26:eq-graded-box-geometry-a} is used as fixed there. No bound on
$u_\ell$ or $\bar{u}_\ell$ themselves at the full datum is among these statements, and none is used
in this proof: the composite entries $a=2,4$, and only they, use the bound on the conjugator
$g=\bar{u}_\ell^{-1}u_\ell$ at the full datum recalled above, whose proof is incomplete at one step.
What the argument
establishes is the balance of the chain: every increase by $\beta$ in \cite[pp.~278--280]{B-CMP109}
is replaced by the quantity it bounds, and these quantities, together with the constant
$G'_{\square}$ of \eqref{4.26:eq-near-pair-chain-b}, which accompanies the use of
\cite[(3.38)]{B-CMP109}, are collected in the ceiling on $\gamma$, the first of the
conditions \eqref{4.26:eq-inherited-conditions-a}, through
\eqref{4.26:eq-near-pair-chain-b} and in the two shares $\frac16\mathrm{room}_T$, one from the
definition of $\alpha_3(T)$ in \eqref{4.26:eq-radii-window-c} and one from the reserve
$\frac13\mathrm{room}_{a}$ in the thresholds $\theta^{(n)}_{T,a}$ of the clause $\mathfrak{I}_n$. No
factor $L^{-2}$ is used, so the estimates
hold alike when $\ell_T=m(y)$ and when $\ell_T<m(y)$.
\end{proof}

\begin{lemma}[The remainder pieces of a near pair]\label{4.26:lem-near-remainders}
Let the stage $n$, the region $Y$, the pair $(\mathbb{U},\mathbb{J})\in\mathfrak{D}^{\flat}_n(Y)$, the
near pair $(T,\square)$ and the hypotheses be those of
Proposition~\ref{4.26:prop-near-pairs-outside}. Thus $T\in\mathcal{T}_1\cap\mathcal{T}^{\mathbb{E}\mathbb{R}}$
has level $\ell$, $X_T\subset\tilde{\square}^2$ and $\square_0=\tilde{\square}^5\subset Y$. Let
$P(\varsigma;t,\tau;B')$ be the members of family (c) of Definition~\ref{4.26:def-proof-family}. Let
$\mathsf{D}_{t_\square}$ be the far datum of family (c$'$) there, with
$\mathsf{D}_0$ its value at $t_\square=0$, and put
$\delta\mathbf{H}_\ell:=\partial_{t_\square}|_0\mathbf{H}_\ell(\square_0,\mathsf{D}_{t_\square})$. Let
$D=d_\ell(X_T,\operatorname{supp}(1-\zeta'_\square))$; then $D\ge wL^{(i-\ell)^+}$, where $i$, $w$ and
$K^{\mathrm{rem}}$ are as in \eqref{4.26:eq-radii-window-a}.

The \emph{remainder term} of the near pair is $T$ at the near configuration composed with one more
own-level box layer $e^{i\xi\mathbf{H}_\ell(\square_0,\mathsf{D}_{t_\square})}$. It is written, as in
\cite[(3.21)]{B-CMP109}, as the integral of a derivative, and so it is of first order in the far datum:
\begin{equation}\label{4.26:eq-near-remainders-1}
\int_0^1 d\tau_1\,\partial_{\tau_1}
T\bigl(X_T;e^{i\xi\tau_1\mathbf{H}_\ell(\square_0,\mathsf{D}_{t_\square})}\cdot
P(\varsigma;t,t_\square;B')\bigr);
\end{equation}
here the parameter $\tau$ of $P$, the coefficient of $\zeta_\square$ in the stage layer of family (c), is
set equal to $t_\square$. The \emph{remainder piece} is the derivative of
\eqref{4.26:eq-near-remainders-1} in $t_\square$ at $t_\square=0$. Let $\|T\|_\infty$ denote the supremum
of $|T(X_T;\cdot)|$ on $\mathfrak{U}_T$. Then the following hold for all
$\varsigma,t\in[0,1]$ and $|B'|<\alpha_3(T)$.
\begin{enumerate}
\item[(i)] The remainder piece is the sum of two terms:
\begin{equation}\label{4.26:eq-near-remainders-2}
\begin{split}
&\partial_{t_\square}|_0\int_0^1 d\tau_1\,\partial_{\tau_1}
T\bigl(X_T;e^{i\xi\tau_1\mathbf{H}_\ell(\square_0,\mathsf{D}_{t_\square})}\cdot
P(\varsigma;t,t_\square;B')\bigr)\\
&\quad=\partial_\tau|_0\,T\bigl(X_T;e^{i\xi[\mathbf{H}_\ell(\square_0,\mathsf{D}_0)+\tau\,\delta\mathbf{H}_\ell]}
\cdot P(\varsigma;t,0;B')\bigr)\\
&\qquad+\partial_{\tau_2}|_0\int_0^1 d\tau_1\,\partial_{\tau_1}
T\bigl(X_T;e^{i\xi\tau_1\mathbf{H}_\ell(\square_0,\mathsf{D}_0)}\cdot P(\varsigma;t,\tau_2;B')\bigr).
\end{split}
\end{equation}
\item[(ii)] Let the complex parameters range over the radii \eqref{4.26:eq-radii-window-a}:
\begin{equation}\label{4.26:eq-near-remainders-3}
\begin{split}
&(\mathrm{c}'_1)\quad |\tau|\le \frac{\mathsf{w}_0e^{\frac14\delta D}}{14K^{\mathrm{rem}}\vartheta(\square)};\\
&(\mathrm{c}'_2)\quad |\tau_1|\le\frac{\mathsf{w}_0e^{\frac14\delta D}}{28K^{\mathrm{rem}}\vartheta},
\qquad |\tau_2|\le\mathsf{w}_0r^{\mathrm{nr}}_\square .
\end{split}
\end{equation}
Take any clause $(y,a)\in\operatorname{Cl}(T)$ with $y\in X_T\cap\mathfrak{T}$. Then each of the two
configurations of \eqref{4.26:eq-near-remainders-2} moves the entry of that clause by less than
$\iota^1_n(y)\mathbf{r}^T_a$. Both configurations lie in $\mathfrak{U}_T$.
\item[(iii)] The absolute value of the remainder piece is at most
\begin{equation}\label{4.26:eq-near-remainders-4}
\|T\|_\infty\cdot 28K^{\mathrm{rem}}\vartheta^{\natural}(\square)
e^{-\frac14\delta D}\bigl[1+4K^{\mathrm{nr}}\vartheta^{\natural}\bigr].
\end{equation}
\item[(iv)] Sum the bounds \eqref{4.26:eq-near-remainders-4}, at fixed $\square$, over the terms $T$ that
form a near pair with $\square$. The sum is at most a constant times
$K^{\mathrm{rem}}\vartheta^{\natural}(\square)(E_0+1)$. The constant does not depend on $\square$ or on
the depth $i-\ell$. The increments $\iota^1$, the spend $\mathsf{w}_0=1/24$ and every degree, such as
$\mathsf{m}$ and $\mathsf{m}_0$, are unchanged.
\end{enumerate}
\end{lemma}

\begin{proof}
\cite[(3.21)]{B-CMP109} takes a function $\zeta'_\square$ with $\zeta'_\square=1$ on
$\tilde{\square}^3$ and $\zeta'_\square=0$ outside $\tilde{\square}^4$. It then removes from the near term
the part of the own-level layer localised outside $\tilde{\square}^3$. The derivative of
$\mathbf{H}_j$ with respect to the field decays exponentially. In addition, the localisation domain $X$
satisfies
\[
\operatorname{dist}^{(\xi)}(X,\operatorname{supp}(1-\zeta'_\square))\ge M(L^j\eta)^{-1}.
\]
From these two facts
the removed term obeys a bound of the type \cite[(3.9)]{B-CMP109} with a small factor, and it is
treated in the same way as the terms of \cite[(3.7)]{B-CMP109}. The expression \cite[(3.22)]{B-CMP109}
is localised in $\tilde{\square}^4\setminus\tilde{\square}^3$ and is treated by the same considerations.

In the present setting these remainders carry own-level layers $\mathbf{H}_\ell(\square_0,\cdot)$. The
far datum $\mathsf{D}_{t_\square}$ is supported at $d_\ell$-distance at least $D$ from $X_T$, and
$D\ge wL^{(i-\ell)^+}$. The own-level layer estimate applies to the remainder data without change; its
constant on these data is the constant $K^{\mathrm{rem}}$ of \eqref{4.26:eq-radii-window-a}.

We prove (i). Write
\[
F(\tau',\tau''):=\int_0^1 d\tau_1\,\partial_{\tau_1}
T\bigl(X_T;e^{i\xi\tau_1\mathbf{H}_\ell(\square_0,\mathsf{D}_{\tau'})}\cdot
P(\varsigma;t,\tau'';B')\bigr).
\]
Then the remainder term \eqref{4.26:eq-near-remainders-1} is $F(t_\square,t_\square)$. By the chain rule,
the remainder piece is $\partial_{\tau'}F(0,0)+\partial_{\tau''}F(0,0)$.

The second summand is the second line of \eqref{4.26:eq-near-remainders-2}; this is the definition of
$F$ with $\tau''=\tau_2$.

For the first summand, the fundamental theorem of calculus in $\tau_1$ gives
\[
F(\tau',0)=T\bigl(X_T;e^{i\xi\mathbf{H}_\ell(\square_0,\mathsf{D}_{\tau'})}\cdot P(\varsigma;t,0;B')\bigr)
-T\bigl(X_T;P(\varsigma;t,0;B')\bigr).
\]
The subtracted term does not depend on $\tau'$. Two curves in the layer variable agree to first order at
$0$:
\[
\tau'\mapsto\mathbf{H}_\ell(\square_0,\mathsf{D}_{\tau'})
\quad\text{and}\quad
\tau\mapsto\mathbf{H}_\ell(\square_0,\mathsf{D}_0)+\tau\,\delta\mathbf{H}_\ell .
\]
Both take the value $\mathbf{H}_\ell(\square_0,\mathsf{D}_0)$ at $0$ and have the derivative
$\delta\mathbf{H}_\ell$ there. By the chain rule, $\partial_{\tau'}F(0,0)$ is therefore the first line of
\eqref{4.26:eq-near-remainders-2}.

So the $t_\square$-derivative falls either on the far datum (first line) or on the near pair
(second line). The second line keeps the first-order structure in the far datum.

We prove (ii). Fix a clause $(y,a)\in\operatorname{Cl}(T)$ with $y\in X_T\cap\mathfrak{T}$ and put
$m=m(y)$. Two bounds on the remainder data are used. Both come from the far-layer estimate for
family (b), that is, the estimate for the members
$\mathcal{W}_T(t,\tau)$ of family (b) of Definition~\ref{4.26:def-proof-family}: it has one bound for the
profile of a far own-level layer and one for the profile of its variation, each read with
$d_\ell(y,\cdot)\ge D$. Together with the own-level layer estimate, they give:
\begin{itemize}
\item the profile of the far layer at $y$ is at most $5.5K^{\mathrm{rem}}\vartheta\rho_m\alpha_{\cdot,\ell}$;
\item the profile of its $t_\square$-variation is at most
$4.4K^{\mathrm{rem}}\vartheta(\square)e^{-\frac14\delta D}\rho_m\alpha_{\cdot,\ell}$.
\end{itemize}

We measure the entry moves in the units $\rho_m\mathbf{r}^T_a$ of the clause. Two further inputs enter:
\begin{itemize}
\item the factor $1.1$ of \eqref{4.26:eq-family-consumption-d};
\item the near shares of clause (v) of Lemma~\ref{4.26:lem-family-consumption}, namely
$\frac13\mathsf{w}_0$ for the dilated stage layer and $\frac1{18}\mathsf{w}_0$ for the undilated one.
\end{itemize}

In configuration $(\mathrm{c}'_1)$ the near pair is taken at $\tau=0$, so only its undilated stage
layer contributes. In configuration $(\mathrm{c}'_2)$ it is taken at $\tau_2$, and both its dilated and
its undilated stage layers contribute. The entry moves are therefore bounded by the left sides of the
following two lines:
\begin{equation}\label{4.26:eq-near-remainders-a}
\begin{split}
&1.1\bigl[5.5K^{\mathrm{rem}}\vartheta+4.4|\tau|K^{\mathrm{rem}}\vartheta(\square)e^{-\frac14\delta D}\bigr]
+\tfrac1{18}\mathsf{w}_0\\
&\qquad\le1.1\Bigl(\tfrac{5.5}{28}+\tfrac{4.4}{14}\Bigr)\mathsf{w}_0+\tfrac1{18}\mathsf{w}_0
=0.6173\ldots\,\mathsf{w}_0<\tfrac23\mathsf{w}_0,\\[2pt]
&1.1\cdot5.5K^{\mathrm{rem}}|\tau_1|\vartheta e^{-\frac14\delta D}
+\tfrac13\mathsf{w}_0+\tfrac1{18}\mathsf{w}_0\\
&\qquad\le1.1\cdot\tfrac{5.5}{28}\mathsf{w}_0+\tfrac13\mathsf{w}_0+\tfrac1{18}\mathsf{w}_0
=0.6049\ldots\,\mathsf{w}_0<\tfrac23\mathsf{w}_0.
\end{split}
\end{equation}

The first line uses two estimates.
\begin{itemize}
\item The inequality $28K^{\mathrm{rem}}\vartheta\le\mathsf{w}_0$, contained in the window condition
\eqref{4.26:eq-radii-window-b}, gives $5.5K^{\mathrm{rem}}\vartheta\le\frac{5.5}{28}\mathsf{w}_0$.
\item The radius $(\mathrm{c}'_1)$ gives
$4.4|\tau|K^{\mathrm{rem}}\vartheta(\square)e^{-\frac14\delta D}\le\frac{4.4}{14}\mathsf{w}_0$.
\end{itemize}
The second line uses the radius of $\tau_1$ in $(\mathrm{c}'_2)$, which gives
$K^{\mathrm{rem}}|\tau_1|\vartheta e^{-\frac14\delta D}\le\frac1{28}\mathsf{w}_0$.

The numerical values are as follows. In the first line,
$1.1\bigl(\frac{5.5}{28}+\frac{4.4}{14}\bigr)=\frac{15.73}{28}=0.5617\ldots$ and
$\frac1{18}=0.0555\ldots$. In the second line, $1.1\cdot\frac{5.5}{28}=0.2160\ldots$ and
$\frac13+\frac1{18}=\frac7{18}=0.3888\ldots$.

By \eqref{4.26:eq-near-remainders-a}, both configurations move each entry by less than
$\iota^1_n(y)\mathbf{r}^T_a$. The concluding total of the proof of
Proposition~\ref{4.26:prop-near-pairs-outside} therefore applies to them without change. That total
consists of its starting entry, its $\gamma$-ceiling hypothesis and its two shares of $\frac16$. It
gives $Q_a<\theta_{T,a}(y)$ at both configurations, so both lie in $\mathfrak{U}_T$. This proves (ii).

We prove (iii). By (ii), the functions of $\tau$, $\tau_1$ and $\tau_2$ in
\eqref{4.26:eq-near-remainders-2} are defined on the discs \eqref{4.26:eq-near-remainders-3}. On these
discs they are bounded by $\|T\|_\infty$.

In $\tau$ and in $\tau_2$ the derivative is taken at the centre $0$ of the disc, where Cauchy's
inequality bounds it by the supremum divided by the radius, and hence by twice the supremum divided by
the radius. In $\tau_1$ the bound used is twice the supremum divided by the radius of $\tau_1$ in
$(\mathrm{c}'_2)$, which is the bound Cauchy's inequality gives at the points of the disc of half that
radius.

Cauchy's inequality in $\tau$ on the disc $(\mathrm{c}'_1)$, with the factor $2$, bounds the first line of
\eqref{4.26:eq-near-remainders-2} by
$\|T\|_\infty\cdot28K^{\mathrm{rem}}\vartheta^{\natural}(\square)e^{-\frac14\delta D}$.

For the second line we use Cauchy's inequality in $\tau_2$ and in $\tau_1$, each with twice the
reciprocal of the corresponding radius in $(\mathrm{c}'_2)$. This gives the bound
\[
\|T\|_\infty\cdot\frac{2}{\mathsf{w}_0r^{\mathrm{nr}}_\square}\cdot
\frac{56K^{\mathrm{rem}}\vartheta}{\mathsf{w}_0}e^{-\frac14\delta D}.
\]
Definition~\ref{4.26:def-radii-window} gives $r^{\mathrm{nr}}_\square=1/(K^{\mathrm{nr}}\vartheta(\square))$,
and hence $2/(\mathsf{w}_0r^{\mathrm{nr}}_\square)=2K^{\mathrm{nr}}\vartheta(\square)/\mathsf{w}_0$. With
$\vartheta^{\natural}=\vartheta/\mathsf{w}_0$ and $\vartheta^{\natural}(\square)=\vartheta(\square)/\mathsf{w}_0$,
the bound equals
\[
\|T\|_\infty\cdot28K^{\mathrm{rem}}\vartheta^{\natural}(\square)e^{-\frac14\delta D}\cdot
4K^{\mathrm{nr}}\vartheta^{\natural}.
\]
Adding the bounds for the two lines gives \eqref{4.26:eq-near-remainders-4}.

We prove (iv). Fix $\square$. For each level $\ell$, at most $C_d(5L^{i-\ell})^d$ terms $T$ form a near
pair with $\square$, and for each of them $D\ge wL^{(i-\ell)^+}$. Inserting this into
\eqref{4.26:eq-near-remainders-4}, the sum over $T$ is dominated by the prefactor
$28K^{\mathrm{rem}}\vartheta^{\natural}(\square)\bigl[1+4K^{\mathrm{nr}}\vartheta^{\natural}\bigr]$ times the
series
\[
\sum_\ell C_d(5L^{i-\ell})^dE_0e^{-\frac14\delta wL^{i-\ell}}\le K_S(E_0+1).
\]
This is the series of the summation clause of the far-layer estimate for family (b), with its
constants $C_d$, $E_0$ and $K_S$.

No quantity in this bound grows with the depth $i-\ell$, so no restriction on that depth is needed.
The remainder pieces therefore add at most a constant times
$K^{\mathrm{rem}}\vartheta^{\natural}(\square)(E_0+1)$ to the total of the near-pair contributions at
$\square$. This is of the form of the near-pair member of that total. The increments $\iota^1$, the
spend $\mathsf{w}_0=1/24$ and every degree, such as $\mathsf{m}$ and $\mathsf{m}_0$, are unchanged.
\end{proof}

\begin{definition}[The frozen structure on an older live piece and its weights]
\label{4.26:def-frozen-structure}
Fix the level $k$ of the site and a stage $l$, and let
$\mathfrak{a}=(W_{\mathfrak{a}},k_a,n^+_a,\text{branch})$ be an arrest which is live at $k$, with
$k_a<k$ and $W_{\mathfrak{a}}\subset Z^c$, in the sense of
Definition~\ref{4.26:def-stages-regions}. The following letters belong to the piece itself.
\begin{itemize}
\item $N_a$ is the number written $N=R_{k_a}$ in \cite{B-CMP122a}, and $h_a:=k_a-N_a$; here $R_i$
is the large-field radius of \cite{B-CMP122a} at the level $i$, as in
Lemma~\ref{4.4:lem-collar-avoidance}, and not the ball radius of the observables.
\item $N_0^a$ is the number $N_0$ of \cite{B-CMP122a} at the level $k_a$, namely the smallest
positive integer $N_0$ such that $L^{-N_0+1}MR_{k_a-N_0+1}=M$.
\item We put
\begin{equation}\label{4.26:eq-frozen-structure-1}
k_0:=k_0^a=k_a-N_0^a,\qquad j^+:=h_a+n^+_a\le k_a .
\end{equation}
\end{itemize}
The sets $\Omega_i$, $\Lambda_i$, $Z''_i$ and $\Omega''_i:=(Z''_i)^c$ used below are those of the
piece's own attempt at the level $k_a$, intersected with $W_{\mathfrak{a}}$. For them we use the
description of \cite{B-CMP122a}, written at the level $k_a$ (so that the paper's $h$, $N$, $N_0$
become $h_a$, $N_a$, $N_0^a$, and $k_a-N_0^a=k_0$; and $Z$ in \eqref{4.26:eq-frozen-structure-2}
and in the display below it denotes the large-field region of the piece's own attempt at the site
$k_a$, not the region $Z$ of Definition~\ref{4.26:def-stages-regions} at the site $k$); in it the
cubes are those of the lattice of spacing $L^{-k_a}$
on which the piece's attempt is run. The sets $Z''_{k_a}$
and $Z''_{k_a}\setminus Z''_{k_a-1}$ are unions of $M$-cubes of this lattice, and $Z''_{k_a}$,
$Z''_{k_a-1}$ are separated by one layer of $M$-cubes. Likewise $Z''_{k_a-1}$ and
$Z''_{k_a-1}\setminus Z''_{k_a-2}$ are unions of $L^{-1}M$-cubes of this lattice, and $Z''_{k_a-1}$,
$Z''_{k_a-2}$ are separated by one layer of $L^{-1}M$-cubes, and so on down the levels. For the
remaining levels,
\begin{equation}\label{4.26:eq-frozen-structure-2}
Z''_j=(\Omega_j^{\sim 5})^c\cap Z,\qquad j=k_0,\,k_0-1,\,\dots,\,h_a+1,
\end{equation}
where $\Omega^{\sim 5}_j$ is the enlargement of $\Omega_j$ denoted so in \cite{B-CMP122a}. The top
set is
\begin{equation*}
Z''_{k_a}=(\Omega^{\sim 5}_{k_0+1})^c\cap Z .
\end{equation*}

\emph{The frozen structure.} By the structure freeze on the components excluded at the arrest's own
site $k_a$, as stated for excluded components in Definition~\ref{4.26:def-stages-regions}, the
structure on $W_{\mathfrak{a}}$ is
$\mathbb{B}^{(n^+_a)}_{k_a}\restriction W_{\mathfrak{a}}$. This is the structure of
\cite[(1.16)]{B-CMP122a} at $j=j^+$. In \cite{B-CMP122a}, in passing to
$\mathbb{B}^{(2)}_{k_a}$ the set $(\Omega^c_{h_a+2}\setminus Z''_{h_a+1})^{(h_a+1)}$ is replaced
by the two sets
\begin{equation*}
(\Omega^c_{h_a+2}\setminus Z''_{h_a+2})^{(h_a+2)}
\quad\text{and}\quad
(Z''_{h_a+2}\setminus Z''_{h_a+1})^{(h_a+1)}=\Gamma''_{h_a+1},
\end{equation*}
where the superscript is the level carried by the set. In general $\mathbb{B}^{(j-h_a)}_{k_a}$ is
defined by \cite[(1.16)]{B-CMP122a}. On $W_{\mathfrak{a}}$ the structure therefore carries:
\begin{itemize}
\item level $m$ on $\Gamma''_m=Z''_{m+1}\setminus Z''_m$, for $h_a<m<j^+$;
\item level $j^+$ on $\Omega^c_{j^++1}\setminus Z''_{j^+}$;
\item the old sets $\Gamma_m$ of \cite[(1.16)]{B-CMP122a} above $j^+$, that is, the old shells
$\Omega_m\setminus\Omega_{m+1}$ with $j^+<m<k_a$, each at the level $m$;
\item below these, the core of the structure of \cite[(1.16)]{B-CMP122a}, at its own levels.
\end{itemize}
At $y\in W_{\mathfrak{a}}$ we write $\ell_{\mathfrak{a}}(y)$ for the level and $w_{\mathfrak{a}}(y)$ for
the weight that this frozen structure assigns to $y$. The pair
$(\ell_{\mathfrak{a}}(y),w_{\mathfrak{a}}(y))$ is the \emph{frozen datum} at $y$. Recall from
Definition~\ref{4.26:def-stages-regions} that $\omega(y)=w_{\mathfrak{a}}(y)$ and
\begin{equation}\label{4.26:eq-frozen-structure-3}
\mathfrak{c}(y)=\ell_{\mathfrak{a}}(y)-\tfrac12\log_{L_0}\omega(y).
\end{equation}
Here $L_0$ is the base of the weights of \cite[(1.64), (1.66)]{B-CMP122a}.

\emph{The case $j^+\le k_0$.} Here $\omega\equiv1$ on $W_{\mathfrak{a}}$: since $j^+\le k_0$, the
weights of \cite{B-CMP122a} carry no powers of $L_0$, and the second clause of its weight
prescription is empty. Hence $\mathfrak{c}=\ell_{\mathfrak{a}}$ on $W_{\mathfrak{a}}$.
Moreover every stratum is at least one layer of cubes of side $MR_\cdot$ wide, $R_\cdot$ being the
large-field radius of \cite{B-CMP122a} introduced with $N_a$ above, so the geometry of a near
pair's box, \eqref{4.26:eq-near-box-geometry-b}, applies on $W_{\mathfrak{a}}$ without change.

\emph{The case $j^+>k_0$.} The frozen levels, the weights and the values of $\mathfrak{c}$ on the parts
of $W_{\mathfrak{a}}$ are those of the following table, which is read from \cite[(1.64), (1.66),
(1.68), pp.~190--191]{B-CMP122a}. Every set in the first column is intersected with
$W_{\mathfrak{a}}$.
\begin{equation}\label{4.26:eq-frozen-structure-a}
\begin{array}{l|c|c|c}
\text{set}\ (\cap W_{\mathfrak{a}}) & \ell_{\mathfrak{a}} & \omega & \mathfrak{c}\\ \hline
\text{core};\ \Gamma''_m,\ m\le k_0 & m & 1 & m\\
\Gamma''_m,\ k_0<m<j^+ & m & L_0^{2(m-k_0)} & k_0\\
\Omega''_{j^+}\setminus\Omega_{k_0+1} & j^+ & L_0^{2(j^+-k_0)} & k_0\\
\Omega_m\setminus\Omega_{m+1},\ k_0<m\le j^+ & j^+ & L_0^{2(j^+-m)} & m\\
\Omega_m\setminus\Omega_{m+1},\ j^+<m<k_a & m & 1 & m
\end{array}
\end{equation}
In each row $m$ denotes the index of the stratum $\Gamma''_m$ or of the shell
$\Omega_m\setminus\Omega_{m+1}$ concerned, and on the core its own level in the structure of
\cite[(1.16)]{B-CMP122a}. The strata $\Gamma''_m$ with
$k_0<m<j^+$ are the \emph{thin strata}, and the sets $\Omega_m\setminus\Omega_{m+1}$ are the
\emph{old shells}.

In both cases,
\begin{equation}\label{4.26:eq-frozen-structure-4}
\omega=L_0^{2(\ell_{\mathfrak{a}}-\mathfrak{c})}\qquad\text{everywhere on }W_{\mathfrak{a}} .
\end{equation}
By \eqref{4.26:eq-frozen-structure-3}, $\mathfrak{c}$ is a pointwise function of the frozen datum
$(\ell_{\mathfrak{a}},w_{\mathfrak{a}})$ alone: it reads none of $W_{\mathfrak{a}}$, $k_0$, $j^+$
or $k_a$. For nested arrests, that is, when several live arrests hold the point $x$, the datum at
$x$ is that of the one arrest that lifted $x$,
that is, the arrest for which $\ell_{\mathfrak{a}}(x)-\mathfrak{c}(x)\ge1$ in the sense of
Definition~\ref{4.26:def-stages-regions}. This follows from the structure freeze on excluded
components (Definition~\ref{4.26:def-stages-regions}) and from the lemma, recalled in
Definition~\ref{4.26:def-stages-regions}, by which a live piece that lies in $Z$ lies in the core,
the site's $h=k-N$ of Definition~\ref{4.26:def-stages-regions} satisfying $h\ge k_a$ there.
\end{definition}

\begin{proof}[Verification of the table and of \eqref{4.26:eq-frozen-structure-4}]
It is to be shown that the column $\mathfrak{c}$ of \eqref{4.26:eq-frozen-structure-a} agrees with
\eqref{4.26:eq-frozen-structure-3}, and that \eqref{4.26:eq-frozen-structure-4} holds.
Solving \eqref{4.26:eq-frozen-structure-3} for $\omega$ gives
$\omega=L_0^{2(\ell_{\mathfrak{a}}-\mathfrak{c})}$. Hence \eqref{4.26:eq-frozen-structure-4} holds
at $y$ exactly when $\ell_{\mathfrak{a}}(y)-\tfrac12\log_{L_0}\omega(y)$ equals the value of
$\mathfrak{c}$ assigned to $y$.

If $j^+\le k_0$, then $\omega\equiv1$. Hence $\log_{L_0}\omega=0$ and $\mathfrak{c}=\ell_{\mathfrak{a}}$,
so that
\begin{equation*}
L_0^{2(\ell_{\mathfrak{a}}-\mathfrak{c})}=L_0^0=1=\omega .
\end{equation*}

Let $j^+>k_0$. In the first three rows of \eqref{4.26:eq-frozen-structure-a} the weight is that of
\cite[(1.64)]{B-CMP122a}, namely $L_0^{2\max\{0,m-k_0\}}$ at the level $m$ on the strata
$\Gamma''_m$ and on the set $\Omega''_{j^+}\setminus\Omega_{k_0+1}$,
evaluated at the frozen level. In the fourth row it is that of \cite[(1.66)]{B-CMP122a}, namely level
$j^+$ and weight $L_0^{2(j^+-m)}$ on the old shell $m$. We compute
$\ell_{\mathfrak{a}}-\tfrac12\log_{L_0}\omega$ row by row.
\begin{itemize}
\item Core and $\Gamma''_m$ with $m\le k_0$: $\ell_{\mathfrak{a}}=m$ and
$\max\{0,m-k_0\}=0$, so $\omega=1$. The value is $m-0=m$.
\item Thin strata $\Gamma''_m$ with $k_0<m<j^+$: $\ell_{\mathfrak{a}}=m$ and
$\omega=L_0^{2(m-k_0)}$. The value is $m-(m-k_0)=k_0$.
\item The set $\Omega''_{j^+}\setminus\Omega_{k_0+1}$: $\ell_{\mathfrak{a}}=j^+$ and
$\omega=L_0^{2(j^+-k_0)}$. The value is $j^+-(j^+-k_0)=k_0$.
\item Old shells $\Omega_m\setminus\Omega_{m+1}$ with $k_0<m\le j^+$: $\ell_{\mathfrak{a}}=j^+$ and
$\omega=L_0^{2(j^+-m)}$. The value is $j^+-(j^+-m)=m$.
\item Old shells $\Omega_m\setminus\Omega_{m+1}$ with $j^+<m<k_a$: $\ell_{\mathfrak{a}}=m$ and
$\omega=1$. The value is $m$.
\end{itemize}
In every row the result is the entry of the column $\mathfrak{c}$. Hence
\eqref{4.26:eq-frozen-structure-4} holds on each part of $W_{\mathfrak{a}}$ listed in
\eqref{4.26:eq-frozen-structure-a}, and so on all of $W_{\mathfrak{a}}$.

That $\mathfrak{c}$ is a pointwise function of the frozen datum is immediate from
\eqref{4.26:eq-frozen-structure-3}; the assertion on nested arrests rests on the two statements
appealed to for it in the definition itself.
\end{proof}

\begin{lemma}[The strata of the original-level function are thick]\label{4.26:lem-original-level-strata}
Work in the setting of Definition~\ref{4.26:def-frozen-structure}, with $j^+>k_0$. Thus $\mathfrak{a}$ is a
live piece with $k_a<k$ and $W_{\mathfrak{a}}\subset Z^c$, carrying the letters $N_a$,
$h_a=k_a-N_a$, $N_0^a$, $k_0=k_a-N_0^a$ and $j^+=h_a+n^+_a\le k_a$. The sets $\Omega_i$,
$\Lambda_i$, $Z''_i$ and $\Omega''_i=(Z''_i)^c$ are those of the piece's attempt, cut to
$W_{\mathfrak{a}}$. Let $\ell_{\mathfrak{a}}$ and $\omega$ be the frozen level and the weight given
by the table~\eqref{4.26:eq-frozen-structure-a}. Let
$\mathfrak{c}=\ell_{\mathfrak{a}}-\tfrac12\log_{L_0}\omega$ be the original-level function, so that
$\omega=L_0^{2(\ell_{\mathfrak{a}}-\mathfrak{c})}$ on $W_{\mathfrak{a}}$. Write
$\Gamma''_m=Z''_{m+1}\setminus Z''_m$ and $\Gamma_m=\Omega_m\setminus\Omega_{m+1}$. A superscript
$\sim n$ on a set $X$ denotes Ba{\l}aban's enlargement $X^{\sim n}$ of $X$ by $n$ layers of cubes of
the level in question, as in \cite[(3.2)--(3.5)]{B-CMP119}; in particular $X\subset X^{\sim n}$.
Then the following hold.
\begin{enumerate}
\item[(i)] The sets $\{\mathfrak{c}=i\}\cap W_{\mathfrak{a}}$ are the strata of
$\mathbb{B}^{(k_0-h_a)}_{k_a}\restriction W_{\mathfrak{a}}$. This is the determining set
\cite[(1.16)]{B-CMP122a} at $j=k_0$, the last stage of the first case of
\cite{B-CMP122a}, the case in which the integration regions are disjoint. Its strata are:
\begin{itemize}
\item level $k_0$ on $\Omega^c_{k_0+1}\setminus Z''_{k_0}$;
\item level $m$ on $\Gamma''_m$ for $h_a<m<k_0$, together with the core;
\item the old shells $\Gamma_m=\Omega_m\setminus\Omega_{m+1}$ for $k_0<m<k_a$.
\end{itemize}
\item[(ii)] $\mathfrak{c}\le\ell_{\mathfrak{a}}$ on $W_{\mathfrak{a}}$.
\item[(iii)] Let $i\ge h_a$. The shell $\{\mathfrak{c}=i\}\cap W_{\mathfrak{a}}$ is at least three
layers of $MR_{i+1}$-cubes of the $(i+1)$-lattice wide, and at least eight such layers wide if
$i\ge k_0$; hence it is at least $3R_{\min}$ sides of the cubes of $\pi_{i+1}$ wide. Moreover, for
$h_a<\ell\le k_a$ the distance from $\Lambda_\ell$ to $\{\mathfrak{c}<\ell\}$ is at least $2R_\ell$
sides of the cubes of $\pi_\ell$.
\end{enumerate}
In short:
\begin{equation}\label{4.26:eq-original-level-strata-a}
\begin{aligned}
&\text{(i)}\quad \{\mathfrak{c}=i\}\cap W_{\mathfrak{a}}\ \text{are the strata of}\
 \mathbb{B}^{(k_0-h_a)}_{k_a}\restriction W_{\mathfrak{a}};\qquad
 \text{(ii)}\quad \mathfrak{c}\le\ell_{\mathfrak{a}};\\
&\text{(iii)}\quad \{\mathfrak{c}=i\}\cap W_{\mathfrak{a}}\ \text{is}\ \ge 3\
 (\text{resp.}\ \ge 8\ \text{if}\ i\ge k_0)\ \text{layers of}\ MR_{i+1}\text{-cubes wide},\\
&\phantom{\text{(iii)}}\quad \text{hence}\ \ge 3R_{\min}\
 \text{sides of }\pi_{i+1}\quad(i\ge h_a),\\
&\phantom{\text{(iii)}}\quad \operatorname{dist}\bigl(\Lambda_\ell,\{\mathfrak{c}<\ell\}\bigr)\ge
 2R_\ell\ \text{sides of }\pi_\ell\quad(h_a<\ell\le k_a).
\end{aligned}
\end{equation}
\end{lemma}

\begin{proof}
\emph{Clause (i).} The value of $\mathfrak{c}$ on each set of the table
\eqref{4.26:eq-frozen-structure-a} is given in its last column. We consider the second and third rows
of the table together with the stratum $\Gamma''_{k_0}$ of the first row. On all of these sets
$\mathfrak{c}=k_0$.

The second row consists of the thin strata $\Gamma''_m$ with $k_0<m<j^+$. Their union with
$\Gamma''_{k_0}$ telescopes to $Z''_{j^+}\setminus Z''_{k_0}$.

The third row is $\Omega''_{j^+}\setminus\Omega_{k_0+1}$. Since $\Omega''_{j^+}=(Z''_{j^+})^c$, this set equals
$\Omega^c_{k_0+1}\setminus Z''_{j^+}$.

We next show that $Z''_{j^+}\subset\Omega^c_{k_0+1}$. Since $j^+\le k_a$, the sets satisfy
$Z''_{j^+}\subset Z''_{k_a}$. By the definition of the sets $Z''_j$ recalled in
Definition~\ref{4.26:def-frozen-structure}, $Z''_{k_a}\subset(\Omega^{\sim 5}_{k_0+1})^c$. The operation
$\sim$ enlarges $\Omega_{k_0+1}$, so
$(\Omega^{\sim 5}_{k_0+1})^c\subset\Omega^c_{k_0+1}$. Hence $Z''_{j^+}\subset\Omega^c_{k_0+1}$.

It follows that the three pieces together make up
\begin{equation*}
\bigl(Z''_{j^+}\setminus Z''_{k_0}\bigr)\cup\bigl(\Omega^c_{k_0+1}\setminus Z''_{j^+}\bigr)
=\Omega^c_{k_0+1}\setminus Z''_{k_0}.
\end{equation*}
This is the level-$k_0$ set of \cite[(1.16)]{B-CMP122a} at $j=k_0$.

The remaining rows give the other strata of that determining set. The first row without
$\Gamma''_{k_0}$ gives the core and the strata $\Gamma''_m$ with $h_a<m<k_0$, where $\mathfrak{c}=m$.
The fourth and fifth rows give the old shells $\Gamma_m=\Omega_m\setminus\Omega_{m+1}$ with
$k_0<m<k_a$, where $\mathfrak{c}=m$. This proves (i).

\emph{Clause (ii).} We read the claim set by set from the table
\eqref{4.26:eq-frozen-structure-a}.
\begin{itemize}
\item On the core, on the strata $\Gamma''_m$ with $m\le k_0$, and on the old shells
$\Omega_m\setminus\Omega_{m+1}$ with $j^+<m<k_a$, we have $\mathfrak{c}=\ell_{\mathfrak{a}}=m$.
\item On the thin strata $\Gamma''_m$ with $k_0<m<j^+$, we have $\mathfrak{c}=k_0<m=\ell_{\mathfrak{a}}$.
\item On $\Omega''_{j^+}\setminus\Omega_{k_0+1}$, we have $\mathfrak{c}=k_0<j^+=\ell_{\mathfrak{a}}$.
\item On the old shells $\Omega_m\setminus\Omega_{m+1}$ with $k_0<m\le j^+$, we have
$\mathfrak{c}=m\le j^+=\ell_{\mathfrak{a}}$.
\end{itemize}

\emph{Clause (iii).} At level $k_a$ the piece $\mathfrak{a}$ arose from a component of the class of
large-field components defined in \cite{B-CMP122a}. Condition (ii) of that class reads: ``in the
preceding $N$ renormalization steps no new large field regions were created inside this component'',
where $N=R_{k_a}=N_a$. Hence the regions of the piece with indices in $\,]h_a,k_a]\,$ are produced by
the operation $\mathbf{T}$ alone: the sets $P$ and $Q$ of \cite[(3.2)--(3.5)]{B-CMP119} do not occur
for these indices, and those relations reduce, as in the width statement for the case $j^+\le k_0$
recorded in Definition~\ref{4.26:def-frozen-structure}, to the following: for $h_a\le i<k_a$,
\begin{equation}\label{4.26:eq-original-level-strata-1}
\Omega^c_{i+1}=Z'^{\,\sim 8}_i,\qquad Z'_i\supset\Lambda^c_i\supset\Omega^c_i .
\end{equation}
We take the shells $\{\mathfrak{c}=i\}\cap W_{\mathfrak{a}}$, $i\ge h_a$, in three cases.

\emph{Case $i>k_0$.} By (i) the shell is $\Gamma_i=\Omega^c_{i+1}\setminus\Omega^c_i$. By
\eqref{4.26:eq-original-level-strata-1}, $\Omega^c_{i+1}=Z'^{\,\sim 8}_i$ and
$\Omega^c_i\subset Z'_i$. Hence the shell contains $Z'^{\,\sim 8}_i\setminus Z'_i$, which is
eight layers of $MR_{i+1}$-cubes of the $(i+1)$-lattice wide.

\emph{Case $i=k_0$.} By (i) the shell is $\Omega^c_{k_0+1}\setminus Z''_{k_0}$. We have
\begin{equation*}
Z''_{k_0}\subset(\Omega^{\sim 5}_{k_0})^c\subset\Omega^c_{k_0}\subset Z'_{k_0},
\end{equation*}
the last inclusion by \eqref{4.26:eq-original-level-strata-1}. Hence the shell contains
$Z'^{\,\sim 8}_{k_0}\setminus Z'_{k_0}$, which is again eight layers of $MR_{k_0+1}$-cubes of the
$(k_0+1)$-lattice wide.

\emph{Case $h_a\le i<k_0$.} By (i) the shell is $\Gamma''_i=Z''_{i+1}\setminus Z''_i$, with
$Z''_{i+1}=(\Omega^{\sim 5}_{i+1})^c$. Since $\Omega^c_{i+1}=Z'^{\,\sim 8}_i$ by
\eqref{4.26:eq-original-level-strata-1}, we get
\begin{equation*}
(\Omega^{\sim 5}_{i+1})^c\setminus Z''_i=Z'^{\,\sim 3}_i\setminus Z''_i
\supset Z'^{\,\sim 3}_i\setminus Z'_i .
\end{equation*}
The inclusion uses $Z''_i\subset Z'_i$. For $i>h_a$ this holds because
$Z''_i\subset(\Omega^{\sim 5}_i)^c\subset\Omega^c_i\subset Z'_i$. For $i=h_a$ it holds because
$Z''_{h_a}=Z_{h_a}=\Lambda^c_{h_a}\subset Z'_{h_a}$. The set $Z'^{\,\sim 3}_i\setminus Z'_i$ is
three layers of $MR_{i+1}$-cubes of the $(i+1)$-lattice wide.

In every case the shell has the number of layers asserted in (iii). Since each layer of
$MR_{i+1}$-cubes of the $(i+1)$-lattice is at least $R_{\min}$ sides of the cubes of $\pi_{i+1}$
wide, the shell is in every case at least $3R_{\min}$ such sides wide.

For the last assertion of (iii), let $h_a<\ell\le k_a$. By (i), the set $\{\mathfrak{c}<\ell\}$ is
determined by the position of $\ell$ relative to $k_0$:
\begin{itemize}
\item if $\ell>k_0$, then $\{\mathfrak{c}<\ell\}=\Omega^c_\ell$;
\item if $\ell\le k_0$, then $\{\mathfrak{c}<\ell\}=Z''_\ell$, and $Z''_\ell\subset\Omega^c_\ell$.
\end{itemize}
In both cases $\{\mathfrak{c}<\ell\}\subset\Omega^c_\ell$. The lower bound of \cite{B-CMP119} on the
distance from $\Lambda_\ell$ to $\Omega^c_\ell$ applies, because its proviso is condition (ii) of the
class quoted above. It gives
$\operatorname{dist}(\Lambda_\ell,\{\mathfrak{c}<\ell\})\ge 2R_\ell$ sides of the cubes of $\pi_\ell$.
\end{proof}

\begin{remark}[Why the stage's cover fails on thin strata]\label{4.26:rem-thin-strata}
In the setting of Definition~\ref{4.26:def-frozen-structure} (a site $k$, a stage, and a live
piece $\mathfrak{a}$ carrying the frozen structure, with $k_0=k_0^a$ and $N_0^a$ as fixed there),
let $k_0<i\le j^+$ and let $\Gamma''_i$ be the thin stratum of level $i$ on $W_{\mathfrak{a}}$
(Definition~\ref{4.26:def-frozen-structure}).
Measure lengths in
\emph{$\pi_i$-sides}, that is, in units of the side of a cube of Ba{\l}aban's partition
$\pi_i$. For $n'\ge k_0$ write $\ell_{n'}$ for the side of a block of level $n'$. Let $(R_j)_j$
be the sequence of enlargement parameters of Lemma~\ref{4.4:lem-collar-avoidance}, so that
$R_{k_a}=N_a$ (Definition~\ref{4.26:def-frozen-structure}); $R_j$ is not the ball radius $R$ of
the observables.
Then:
\begin{enumerate}
\item on $\Gamma''_i$ a box of twelve $\pi_i$-sides placed at a cube of the stage's cover
  meets cubes of the frozen structure whose weights, in the unit of the box, retain the
  factor $L^{i-k_0-1}$, and no constant independent of the piece's coupling bounds this
  factor;
\item the frozen set on such a box has no localisation to a smaller box either: for the
  blocks of the highest levels $n'$ of the frozen structure the factor $L^{n'-k_0}$ in the
  side $\ell_{n'}$ reaches $L^{N_0^a}=L\,R_{k_0+1}$, and this exceeds $12M$, the side of a box
  of scale $k_0$ in the units of that scale, as $\gamma\to0$.
\end{enumerate}
\end{remark}

\begin{proof}
For the first assertion we use the geometry of the strata. The stratum $\Gamma''_i$ is
$L$ $\pi_i$-sides thick, whereas all the finer thin strata together are less than two
$\pi_i$-sides thick. Consequently a box of twelve $\pi_i$-sides placed at a cube
$\square\subset\Gamma''_i$ meets every level down to $k_0$. In the unit of that box, a cube
of level $k_0$ contributes the factor $L^{2(i-k_0)}$, up to the point weights $\omega$.
The condition that the stage imposes on the boxes of its cover keeps the factor
$L^{i-k_0-1}$.
The strata run over the levels $k_0<i\le j^+$ (Definition~\ref{4.26:def-frozen-structure}).
On the top strata the retained factor $L^{i-k_0-1}$ is a polylogarithm of the piece's
coupling; it is not a constant, and it cannot be replaced by a constant independent of that
coupling. This proves the first assertion.

For the second assertion, a block of level $n'$ of the frozen structure has side
\begin{equation*}
  \ell_{n'}=L^{n'-k_0}\,\ell_{k_0}.
\end{equation*}
The factor $L^{n'-k_0}$ runs up to
\begin{equation*}
  L^{N_0^a}=L\,R_{k_0+1}.
\end{equation*}
Since $M$ is fixed before $\gamma$ in the order of choice, as $\gamma\to0$ the factor
$L\,R_{k_0+1}$ exceeds $12M$, so a single block of the highest level is larger than a box of
scale $k_0$, and no smaller box localises it.
\end{proof}

\begin{definition}[The graded cover and the structures on a live piece]\label{4.26:def-graded-cover}
Work at the site $k$, with the stage, the regions $Z$, $Z^{\mathrm{on}}$, $\mathfrak{T}$ and the
live arrests of Definition~\ref{4.26:def-stages-regions}. Let $\mathfrak{B}$ be the current
determining-set structure of the stage. For a structure $\mathfrak{s}$, write
$\lambda_{\mathfrak{s}}(y)$ for the level of the block of $\mathfrak{s}$ that contains the point
$y$. Let $\mathfrak{c}$ be the original-scale level function of
Lemma~\ref{4.26:lem-original-level-strata}.

\emph{The graded structure.} The structure $\mathfrak{B}^{\mathfrak{c}}$ is obtained from
$\mathfrak{B}$ as follows. For every arrest $\mathfrak{a}$ live at the site $k$ with $k_a<k$ and
$W_{\mathfrak{a}}\subset Z^c$, the restriction $\mathfrak{B}\restriction W_{\mathfrak{a}}$ is
replaced by the determining set whose strata are the level sets of $\mathfrak{c}$, as described
in Lemma~\ref{4.26:lem-original-level-strata}\,(i) and
display~\eqref{4.26:eq-original-level-strata-a}. On the lifted data of $Z$ (the core pieces
$W_{\mathfrak{a}}\subset Z$, the lifted set $\mathsf{L}_{\mathfrak{d}}$ of the run's own arrest
$\mathfrak{d}$, and the excluded components $W$) the corresponding replacement is made by the
definition of the graded cover on $Z^{\mathrm{on}}$. Elsewhere $\mathfrak{B}$ is left unchanged.
Thus
\begin{equation}\label{4.26:eq-graded-cover-a}
\begin{gathered}
  \mathfrak{B}^{\mathfrak{c}}\restriction W_{\mathfrak{a}}
  := \text{the determining set with strata }
  \{\mathfrak{c}=i\}\cap W_{\mathfrak{a}}
  \quad (\mathfrak{a}\ \text{live},\ k_a<k,\ W_{\mathfrak{a}}\subset Z^c),\\
  \mathfrak{B}^{\mathfrak{c}}\restriction(\text{lifted data of } Z)
  := \text{the graded structure on } Z^{\mathrm{on}},\\
  \mathfrak{B}^{\mathfrak{c}} := \mathfrak{B}
  \ \text{ elsewhere}.
\end{gathered}
\end{equation}
The structure $\mathfrak{B}^{\mathfrak{c}}$ is admissible: on each $W_{\mathfrak{a}}$ it is a
determining set of Ba{\l}aban's own sequence, and it coincides with $\mathfrak{B}$ off the set
$\{\omega>1\}$. It refines $\mathfrak{B}$, that is,
$\lambda_{\mathfrak{B}^{\mathfrak{c}}}\le\lambda_{\mathfrak{B}}$ at every point; and on each
$W_{\mathfrak{a}}$ it is a predecessor of $\mathfrak{B}$ in the construction.

\emph{The two covers.} The cover $\mathcal{C}^1$ is the cover $\{\square\}$ with its partition
of unity $\{\zeta_\square\}$ that is used for the stage (the cover of \cite{B-CMP119}), built on
$\mathfrak{B}^{\mathfrak{c}}$
instead of $\mathfrak{B}$: on the set $\{\mathfrak{c}=i\}$ its cubes are the unions of $2^d$
neighbouring cubes of the partition $\pi_i$, and the enlargement $\tilde{\square}$ of a cube of
scale $i$ is taken in $\pi_i$. The cover $\mathcal{C}$ is the stage's cover, built on
$\mathfrak{B}$. Off the lifted sets
$\mathsf{L}_{\mathfrak{a}}=(\Omega^c_{j^+}\setminus Z''_{k_0^a+1})\cap W_{\mathfrak{a}}$ of the
live pieces (with $j^+$ and $k_0^a$ the piece's own indices) the two sequences of strata
coincide, and hence so do the cubes of $\mathcal{C}^1$
and of $\mathcal{C}$.

\emph{Which cover is used where.} The objects attached to the class $\mathcal{T}_1$ are written
with the cover $\mathcal{C}^1$:
the first localisation of every term $T\in\mathcal{T}^{\mathbb{E}\mathbb{R}}$
together with its counterterm and its pieces on $W$; the split of the terms into the class
$\mathcal{T}_1$ and the tails; the estimates on these localised terms; the interpolation
$\mathcal{W}_T(t,\tau)$ of Definition~\ref{4.26:def-proof-family}\,(b); the radii of the first
localisation, of the higher localisations and
of the remainders; the families (b), (c), (c$'$) of
Definition~\ref{4.26:def-proof-family}; and the fourth of the five concluding estimates of the
stage. The following keep the cover $\mathcal{C}$: family (a) of
Definition~\ref{4.26:def-proof-family}, with the function $\chi_t$ and the bounds
$R^{\natural}_\square$ on the block shifts; the element sets $\mathcal{A}_c$; one further item
of the stage's defining data; and the first, second, third and fifth of the
concluding estimates of the stage.

\emph{Data at a cube of $\mathcal{C}^1$.} At $\square\in\mathcal{C}^1$ put
\begin{equation}\label{4.26:eq-graded-cover-1}
\begin{gathered}
  \square_0 := \tilde{\square}^5, \qquad
  \mathfrak{s} := \mathfrak{B}^{\mathfrak{c}}\restriction\square_0, \qquad
  p_\square := \text{the top level of } \mathfrak{s} \text{ on } \square_0,\\
  \mathbb{U}^{\mathfrak{s}}_\square = \mathbb{U}^w_\square
  := U_{p_\square}\bigl(\square_0, M^{\mathfrak{s}}(\mathbb{U})\bigr).
\end{gathered}
\end{equation}
Here $M^{\mathfrak{s}}$ is the box minimiser for the structure $\mathfrak{s}$; and for each
level $p\le p_\square$ the triple $(p,\square_0,\mathfrak{s})$ uses Ba{\l}aban's localisation
$\mathbb{B}_p(\square_0)\cup\{\Gamma^{\mathfrak{s}}_i\}_{i<p}$ of the structure to $\square_0$,
that is, the level-$p$ boundary-term set $\mathbb{B}_p(\square_0)$ of Ba{\l}aban's construction
on $\square_0$ together with the strata $\Gamma^{\mathfrak{s}}_i$, $i<p$, of $\mathfrak{s}$ below
the level $p$.

\emph{The reading $(\Sigma^{\ast})^{\mathfrak{a}}$ on a live piece.} We write $(\Sigma^{\ast})$
for the reading of the family of derived entries of a triple $(p,X',\mathfrak{s})$ under which
the second and the fourth derived entries are imposed at a prescribed family of structures
$\mathfrak{s}$, and $(\Sigma^{\ast})^{\mathfrak{a}}$ for the description of that family on a
live piece $\mathfrak{a}$ given in clauses (a)--(c) below. Let $\mathfrak{a}$ be live
at the site $k$, with $k_a<k$ and either $W_{\mathfrak{a}}\subset Z^c$ or
$W_{\mathfrak{a}}\subset Z$. For every localisation domain $X'$ meeting $W_{\mathfrak{a}}$, the
structures $\mathfrak{s}$ at which the reading $(\Sigma^{\ast})$ of the family of derived
entries of a triple imposes the second and the fourth derived entries of the triple
$(p,X',\mathfrak{s})$ include the following.
\begin{enumerate}
\item[(a)] On $X'\cap W_{\mathfrak{a}}$: the frozen structure
$\mathbb{B}^{(n^+_a)}_{k_a}\restriction(X'\cap W_{\mathfrak{a}})$ and each predecessor in the
piece's own run, $\mathbb{B}^{(n')}_{k_a}\restriction(X'\cap W_{\mathfrak{a}})$ with
$0\le n'\le n^+_a$, where
$\mathbb{B}^{(0)}_{k_a}:=\mathbb{B}_{k_a}$, each both before and after the operation
$\mathbf{R}$, as the reading $(\Sigma^{\ast})$ lists them at density $k_a$. These are continued on
$X'\setminus W_{\mathfrak{a}}$ by the structures of $\mathfrak{B}$ itself, namely the present
site's current structure and its predecessors.
\item[(b)] At a $\mathfrak{T}$-location (locations being assigned as in
Definition~\ref{4.26:def-regularity-space}) the clauses are read in the frozen weighted unit
$r_e^{(p)}(y')$ of Definition~\ref{4.26:def-stages-regions}: its level and weight are the
point's frozen pair $(m(y'),\omega(y'))$, and its block side is the block side
\begin{equation*}
  l_{y'}=\min\{p,\lambda_{\mathfrak{s}}(y')\}
\end{equation*}
of the shell construction $S^{\mathrm{bump}}$, where $\lambda_{\mathfrak{s}}(y')$ is the level of
$\mathfrak{s}$ at $y'$. Every space of a later
site reads these clauses in this way: the space $\mathfrak{N}^{(n)}$ of
Definition~\ref{4.26:def-regularity-space} at the factor
\begin{equation*}
  1+2\beta-\sum_{l\le n}\iota^1_l
\end{equation*}
fixed there;
the arrest source space of a later site $k'$ at the factor $F^{(k')}(y')$; and an attempt space
at its own bracket at the frozen level $m=m(y')$.
\item[(c)] At a $Z^{\mathrm{on}}$-location the composite clauses are those of the attempt space
$\mathfrak{L}_l$ of Definition~\ref{4.26:def-regularity-space}, never those of $\mathfrak{N}$
with the factor $5\mathsf{w}^{\mathrm{on}}$. At a retained stage of the run's own, they are read
at its birth radius with the bracket $\theta^{(l)}(\lambda_{\mathfrak{s}})$ of
$\mathfrak{L}_l$. At a retained
structure of an older piece held in the core, that is, contained in $Z$, they are read in the
same frozen weighted unit as at a $\mathfrak{T}$-location,
\begin{equation*}
  \omega(y')\, r_e\bigl(\ell_{\mathfrak{a}}(y'),\min\{p,\lambda_{\mathfrak{s}}(y')\}\bigr),
  \qquad\text{with the bracket } \theta^{(l)}\bigl(\ell_{\mathfrak{a}}(y')\bigr),
\end{equation*}
as every space of a later site reads it, the source space
$\tilde{U}^{\mathrm{arrest}}_k=\mathfrak{L}_0$ included.
\end{enumerate}
In particular, for every $\square\in\mathcal{C}^1$ whose enlargement
$\square_0=\tilde{\square}^5$ meets $W_{\mathfrak{a}}$, including those with
$\square_0\not\subset W_{\mathfrak{a}}$,
\begin{equation}\label{4.26:eq-graded-cover-2}
  \bigl(p,\square_0,\mathfrak{B}^{\mathfrak{c}}\restriction\square_0\bigr),
  \qquad p\le p_\square,
\end{equation}
is a $(\Sigma^{\ast})$-triple.
\end{definition}

\begin{proof}[Justification of the properties stated in
Definition~\ref{4.26:def-graded-cover}]
\emph{The graded structure.} By Lemma~\ref{4.26:lem-original-level-strata}\,(i), the sets
$\{\mathfrak{c}=i\}\cap W_{\mathfrak{a}}$ are the strata of the stage structure
$\mathbb{B}^{(k_0^a-h_a)}_{k_a}\restriction W_{\mathfrak{a}}$, with $k_0^a$ and $h_a$ the
piece's own indices. This is Ba{\l}aban's determining
set at the last stage of the case of disjoint integration regions, a stage structure of the
piece's own run.
Therefore the replacement \eqref{4.26:eq-graded-cover-a} puts on $W_{\mathfrak{a}}$ a
determining set of Ba{\l}aban's own sequence. On $W_{\mathfrak{a}}$ the structure
$\mathfrak{B}$ is the frozen structure of the arrest, whose level function is
$\ell_{\mathfrak{a}}$ (the frozen reading of Definition~\ref{4.26:def-stages-regions}). By the
definition of $\mathfrak{c}$ in Lemma~\ref{4.26:lem-original-level-strata} one has
$\mathfrak{c}=\ell_{\mathfrak{a}}-\tfrac12\log_{L_0}\omega$, and by
Lemma~\ref{4.26:lem-original-level-strata}\,(ii) one has $\mathfrak{c}\le\ell_{\mathfrak{a}}$;
hence $\omega\ge1$ on $W_{\mathfrak{a}}$. Therefore off $\{\omega>1\}$ one has $\omega=1$ and
\begin{equation*}
  \mathfrak{c}=\ell_{\mathfrak{a}}-\tfrac12\log_{L_0}\omega=\ell_{\mathfrak{a}},
\end{equation*}
so the replacement leaves $\mathfrak{B}$ unchanged off $\{\omega>1\}$; this is the
admissibility stated in Definition~\ref{4.26:def-graded-cover}. Since that determining set is a
structure of the run of $\mathfrak{a}$, it precedes $\mathfrak{B}$ in the construction on
$W_{\mathfrak{a}}$. Refinement is meant in the sense of the order $\preccurlyeq$ on structures:
$\mathfrak{B}^{\mathfrak{c}}$ refines
$\mathfrak{B}$ means the pointwise inequality
$\lambda_{\mathfrak{B}^{\mathfrak{c}}}\le\lambda_{\mathfrak{B}}$ between the level functions.
On $W_{\mathfrak{a}}$ the levels of $\mathfrak{B}^{\mathfrak{c}}$ are the values of
$\mathfrak{c}$, and these satisfy $\mathfrak{c}\le\ell_{\mathfrak{a}}$ by
Lemma~\ref{4.26:lem-original-level-strata}\,(ii); since the levels of $\mathfrak{B}$ on
$W_{\mathfrak{a}}$ are the values of $\ell_{\mathfrak{a}}$, this gives
$\lambda_{\mathfrak{B}^{\mathfrak{c}}}\le\lambda_{\mathfrak{B}}$ there. On the lifted data of $Z$
the structure $\mathfrak{B}^{\mathfrak{c}}$ is, by \eqref{4.26:eq-graded-cover-a}, the graded
structure on $Z^{\mathrm{on}}$, and admissibility, refinement and the predecessor property hold
there as part of the definition of that structure. Off the pieces
$W_{\mathfrak{a}}\subset Z^c$ and off the lifted data of $Z$ the two structures coincide.
Off the lifted sets the sequences of strata of
$\mathfrak{B}^{\mathfrak{c}}$ and of $\mathfrak{B}$ coincide. The cubes of $\mathcal{C}^1$ and
$\mathcal{C}$ are built from these sequences by the same rule, so they agree there as well.

\emph{The reading $(\Sigma^{\ast})^{\mathfrak{a}}$.} In \cite{B-CMP122a} the new terms
$\mathbb{B}^{(n)}_k$ are kept as a part of all boundary terms. It is stated there that the
definition of the spaces generalises to this case for the next steps. Clauses (a)--(c) fix that
generalisation on a live piece.

Two rules on the spaces written at a density $k'$ after an arrest are used: the frozen-reading
rule and the cube-clause rule. The frozen-reading rule is the following.
Every such space reads at $x\in W_{\mathfrak{a}}$ the frozen triple
$(\ell,w,\mathcal{Q})(x)$ of the innermost live arrest holding $x$, namely its frozen level
$\ell(x)$, its frozen weight $w(x)=\omega(x)$ (Definition~\ref{4.26:def-regularity-space}) and
its frozen cube clauses $\mathcal{Q}(x)$; where this arrest is $\mathfrak{a}$, one has
$\ell(x)=\ell_{\mathfrak{a}}(x)$. This applies to the standard
spaces $\tilde{U}_{k'}(X,\tilde{\alpha})$, the spaces after the operation $\mathbf{T}$, the
$\mathbf{R}$-sources $\tilde{U}^c_k(Y_4,\tilde{\alpha})$, the attempt spaces of sibling attempts
and of enclosing attempts, the spaces $\mathfrak{L}_l$, the tubes, and the spaces of newborn
pieces.
The cube-clause rule has three parts. The cube clauses of such a space are those of
$\mathcal{Q}(x)$ at weight $w(x)$. Off live pieces the spaces are those of
\cite[(2.34)--(2.39)]{B-CMP119} at level $k'$, without change. The readings $(\Sigma^{\flat})$
and $(\Sigma^{\ast})$ are unchanged; here $(\Sigma^{\flat})$ is the reading in which an attempt
space of a structure's own run reads the structure at its birth radius.

The frozen clauses were written at density $k_a$. At that density the reading $(\Sigma^{\ast})$
names exactly the structures $\mathbb{B}^{(n')}_{k_a}$, $0\le n'\le n^+_a$, before and after
$\mathbf{R}$. Since the cube-clause rule leaves the readings unchanged, this family of structures
is carried into every later space on $W_{\mathfrak{a}}$. This is clause (a). Off
$W_{\mathfrak{a}}$, the same rules leave the structures of $\mathfrak{B}$ in place.

That the units and brackets must be those of clauses (b) and (c) is shown by the inclusion of
the source space in an attempt space,
$\mathfrak{L}_l\supseteq\tilde{U}^{\mathrm{arrest}}_{k'}$, in the extended inclusion lemma for the
attempt spaces. Part (a) of that lemma proves this inclusion as follows. Both sides carry the
same frozen $(\ell,w,\mathcal{Q})$, so their cube clauses coincide, and the inclusion reduces to
the inequality between the factors. This argument applies only if the source space and the
attempt space impose the derived entries at the same triples. We examine the two kinds of
triple in turn.

At the triples of the frozen structure itself, the two spaces impose the derived entries at the
same triples, and the argument is as just described.

At a retained structure, where $\lambda_{\mathfrak{s}}(x)<\ell(x)$, the two spaces may differ.
An attempt space of the structure's own run reads the structure at its birth radius; this is
the reading $(\Sigma^{\flat})$. The source space, and every other space written at a later
density, reads instead the radius of the frozen shell, with the block side of
$S^{\mathrm{bump}}$. Against an attempt space of the structure's own run, the inclusion is
therefore the ratio of radii of part (a$'$) of the extended inclusion lemma for the attempt
spaces, of the kind treated in its part (b). Between two spaces written at later densities, it
is the inequality between the factors in one and the same unit. These are the two readings of
clause (c): at a retained stage of the run's own, the birth radius with the bracket
$\theta^{(l)}(\lambda_{\mathfrak{s}})$; at a retained structure of an older piece, the frozen
weighted unit of clause (b) with the bracket $\theta^{(l)}(\ell_{\mathfrak{a}}(y'))$, for every
later space, $\tilde{U}^{\mathrm{arrest}}_k=\mathfrak{L}_0$ included.

For an excluded component $W$ of the present site, the frozen reading of
Definition~\ref{4.26:def-stages-regions} is already of this form. The reading
$(\Sigma^{\ast})^{\mathfrak{a}}$ is the same statement for a piece arrested at an earlier
density $k_a<k$ and read at the present site. Finally, \eqref{4.26:eq-graded-cover-2} is the
instance of clauses (a)--(c)
with $X'=\square_0$ and $\mathfrak{s}=\mathfrak{B}^{\mathfrak{c}}\restriction\square_0$. This
structure is, on $\square_0\cap W_{\mathfrak{a}}$, a structure of the run of $\mathfrak{a}$ by
the first paragraph of this proof. The levels
$p\le p_\square$ do not exceed its top level on $\square_0$.
\end{proof}

\begin{lemma}[Junction boxes and the cost of the added structures]\label{4.26:lem-junction-boxes}
Fix the site $k$ and a stage. Let $\mathfrak{a}$ be live at the site $k$, with $k_a<k$ and either
$W_{\mathfrak{a}}\subset Z^c$ or $W_{\mathfrak{a}}\subset Z$. Use the notation $N_a$, $h_a=k_a-N_a$, $N_0^a$,
$k_0=k_a-N_0^a$, $j^+=h_a+n^+_a\le k_a$, $\Omega_i$, $Z''_i$ of
Definition~\ref{4.26:def-frozen-structure}, the lifted set $\mathsf{L}_{\mathfrak{a}}$, the weight $\omega$, the
lift $g$, the frozen level $\ell_{\mathfrak{a}}$ and the original-level function $\mathfrak{c}$ fixed there,
and the graded structure $\mathfrak{B}^{\mathfrak{c}}$ and graded cover $\mathcal{C}^1$ of
Definition~\ref{4.26:def-graded-cover}, with $\square_0=\tilde{\square}^5$ and $p_{\square}$ the top level of
$\mathfrak{B}^{\mathfrak{c}}\restriction\square_0$ on $\square_0$, for $\square\in\mathcal{C}^1$. We call
$\square\in\mathcal{C}^1$ a \emph{junction box} of $\mathfrak{a}$ if $\square_0$ meets $W_{\mathfrak{a}}$ and
$\square_0\not\subset W_{\mathfrak{a}}$. Then the following hold.
\begin{enumerate}
\item[(a)] $\mathfrak{c}\neq\ell_{\mathfrak{a}}$ only on $\mathsf{L}_{\mathfrak{a}}$, and
$\mathsf{L}_{\mathfrak{a}}\subset\Omega^c_{j^+}\cap W_{\mathfrak{a}}$.
\item[(b)] If $j^+<k_a$, every box of $\mathcal{C}^1$ meeting $\mathsf{L}_{\mathfrak{a}}$ has
$\square_0\subset W_{\mathfrak{a}}$; a box of $\mathcal{C}^1$ whose $\square_0$ reaches $\partial W_{\mathfrak{a}}$
meets no point of $\mathsf{L}_{\mathfrak{a}}$, and on it $\mathfrak{B}^{\mathfrak{c}}$ coincides with $\mathfrak{B}$.
\item[(c)] If $j^+=k_a$, then for every junction box $\square$ of scale $k_a-1$ or of scale $k_a$ the set
$\square_0$ meets neither $\Omega_{k_a+1}$ nor a stratum of a later run, and
\begin{align}
\mathfrak{B}\restriction(\square_0\setminus W_{\mathfrak{a}})
&=\mathbb{B}^{(n')}_{k_a}\restriction(\square_0\setminus W_{\mathfrak{a}})\qquad\text{for every } n',
\notag\\
\mathfrak{B}^{\mathfrak{c}}\restriction\square_0&=\mathbb{B}^{(k_0-h_a)}_{k_a}\restriction\square_0 .
\label{4.26:eq-junction-boxes-1}
\end{align}
The right side is one stage structure of density $k_a$ of the piece's own run: $k_0>h_a$, and
$\mathbb{B}^{(k_0-h_a)}_{k_a}$ is a strict predecessor of the frozen structure
$\mathbb{B}^{(n^+_a)}_{k_a}$ if and only if $j^+>k_0$.
\item[(d)] If $j^+\le k_0$, then $\omega\equiv1$ on $W_{\mathfrak{a}}$ and $\mathfrak{B}^{\mathfrak{c}}=\mathfrak{B}$
there, so that no structure is added on $W_{\mathfrak{a}}$.
\item[(e)] In the case $j^+=k_a$ the levels of $\mathfrak{B}^{\mathfrak{c}}$ on the two sides of
$\partial W_{\mathfrak{a}}$ are $k_a-1$ inside and $k_a$ outside, and $\mathfrak{B}^{\mathfrak{c}}$ is admissible
there; the inequality $\mathsf{m}^{\circ}=\mathfrak{c}\le\mathsf{m}$ between the parent level and the frame
level persists on junction boxes; and no set $\square_0$ with $\square\in\mathcal{C}^1$ meets two disjoint live
pieces.
\item[(f)] Every triple $(p,X',\mathfrak{s})$ of the family $(\Sigma^{\ast})^{\mathfrak{a}}$ of
Definition~\ref{4.26:def-graded-cover}, in particular the triple
$(p,\square_0,\mathfrak{B}^{\mathfrak{c}}\restriction\square_0)$ of every junction box $\square$ and every
$p\le p_{\square}$, imposes no condition beyond those already in force: the real-point bound, the increment and
tube bounds, and the clauses of the arrested source space $\tilde{U}^{\mathrm{arrest}}_k$ hold at it.
\end{enumerate}
\end{lemma}

\begin{proof}
\emph{Step 1: where the lifted set lies.} By the admissibility property of the frozen structure on a
live piece, the original-level function $\mathfrak{c}$ differs from the frozen level $\ell_{\mathfrak{a}}$ only on
the lifted set $\mathsf{L}_{\mathfrak{a}}$, and $\mathsf{L}_{\mathfrak{a}}\subset\Omega^c_{j^+}\cap W_{\mathfrak{a}}$.
This is~(a). Since $j^+\le k_a$, exactly one of the cases $j^+<k_a$ and $j^+=k_a$ occurs.

\emph{Step 2: the case $j^+<k_a$.} Here $\mathsf{L}_{\mathfrak{a}}\subset\Omega^c_{j^+}$ is separated from
$\partial W_{\mathfrak{a}}$ by the old shells $j^+,\dots,k_a-1$. On these shells $\omega=1$. If $j^+\le k_0$, then
$\omega\equiv1$ on $W_{\mathfrak{a}}$ and $\mathfrak{B}^{\mathfrak{c}}=\mathfrak{B}$ there by~(d) (Step~7 below), so
nothing is lifted; let therefore $j^+>k_0$. Then each shell $i$ with $j^+\le i\le k_a-1$ has $i\ge k_0$, and is
therefore at least eight layers of cubes of side $MR_{i+1}$ of the $(i+1)$-lattice wide
(Lemma~\ref{4.26:lem-original-level-strata}(iii)). A box of $\mathcal{C}^1$ meeting $\mathsf{L}_{\mathfrak{a}}$
has scale at most
$j^+$, and its $\square_0=\tilde{\square}^5$ lies within seven sides, of the box's own scale, of the box. Seven
sides of
scale at most $j^+$ do not cross the eight layers of the old shell adjacent to $\mathsf{L}_{\mathfrak{a}}$;
hence $\square_0\subset W_{\mathfrak{a}}$. A box whose $\square_0$ reaches
$\partial W_{\mathfrak{a}}$ meets no lifted point, by the same count of layers. On such a box
$\mathfrak{c}=\ell_{\mathfrak{a}}$ by Step~1, so
the determining
set of the strata of $\mathfrak{c}$ coincides there with the current structure, and by
Definition~\ref{4.26:def-graded-cover} $\mathfrak{B}^{\mathfrak{c}}=\mathfrak{B}$ on it. This is~(b).

\emph{Step 3: the case $j^+=k_a$, boxes of scale $k_a-1$.} Since $j^+=h_a+n^+_a$ and $h_a=k_a-N_a$, the case
$j^+=k_a$ is the case $n^+_a=N_a$. Here $W_{\mathfrak{a}}=X_{\mathfrak{a}}\cap\Omega^c_{k_a}$, where
$X_{\mathfrak{a}}$ denotes the large-field region of the arrest $\mathfrak{a}$, and the outermost
old shell $k_a-1$ of $W_{\mathfrak{a}}$ is lifted up to $\partial\Omega_{k_a}$: on it, with $g$ the lift,
\begin{equation}\label{4.26:eq-junction-boxes-2}
\ell_{\mathfrak{a}}=k_a,\qquad \omega=L_0^2,\qquad \mathfrak{c}=k_a-1,\qquad g=1,
\end{equation}
by the weight table \eqref{4.26:eq-frozen-structure-a}. The weight $L_0^2$ is the factor
$L_0^{2(j-m)}$, read at $j-m=1$, in the bounds \cite[(1.67)]{B-CMP122a}, which in the notation of that paper
read
\begin{equation*}
L^j\eta\,|A|,\ \ (L^j\eta)^2\,|\nabla^{\eta}A| < L_0^{2(j-m)}BCM\alpha_{0,j} .
\end{equation*}
A box of $\mathcal{C}^1$ of scale $k_a-1$ on this shell has $\square_0$ within seven
sides of its own scale of the box, so $\square_0\setminus W_{\mathfrak{a}}$ lies within seven sides of the cubes
of $\pi_{k_a-1}$ of $\partial\Omega_{k_a}$. It therefore lies inside the shell
$\Omega_{k_a}\setminus\Omega_{k_a+1}$. This shell is made by the operation $\mathbf{T}_{k_a+1}$, and by the
width statement for the shells made by the operations $\mathbf{T}$ it is at least eight layers of
$MR_{k_a+1}$-cubes of the $(k_a+1)$-lattice wide.

Two further situations must be considered: a later arrest $\mathfrak{b}$ nested over $\mathfrak{a}$,
$W_{\mathfrak{b}}\supset W_{\mathfrak{a}}$, and the current run when $W_{\mathfrak{a}}\subset Z$, the run's
starting level $h$ then satisfying $h\ge k_a$. In both, the shell $\Omega_{k_a}\setminus\Omega_{k_a+1}$ still
keeps at least three layers of $MR_{k_a+1}$-cubes below the regions of that run. This follows from the lemma on
the position of a later run's regions, in its clause bounding the distance of these regions from the earlier
shells, together with
\begin{equation}\label{4.26:eq-junction-boxes-3}
Z''_{h+1}=\bigl(\Omega^{\sim5}_{h+1}\bigr)^c\cap Z
\end{equation}
from \cite[(1.11)]{B-CMP122a}, where $\Omega^{\sim5}_{h+1}$ is the enlargement of $\Omega_{h+1}$ written so in
\cite[(1.11)]{B-CMP122a}. Seven sides of scale $k_a-1$ stay inside these layers.
Hence $\square_0$ reaches neither $\Omega_{k_a+1}$ nor a stratum of a later run.

\emph{Step 4: the case $j^+=k_a$, boxes of scale $k_a$.} Let $\square\in\mathcal{C}^1$ be of scale $k_a$, lie in
the shell $k_a$, and be such that $\square_0$ meets $W_{\mathfrak{a}}$. Then
$\square_0\setminus W_{\mathfrak{a}}$ lies within twelve sides of the cubes of $\pi_{k_a}$ of
$\partial\Omega_{k_a}$, and $12<3LR_{k_a+1}$. So, by the three layers of Step~3, $\square_0$ again reaches neither
$\Omega_{k_a+1}$ nor a stratum of a later run. On the other side, $\square_0$ reaches at most seven sides of the
cubes of $\pi_{k_a}$ into the old shell $k_a-1$, and $7<8R_{k_a}$. Thus the junction hypothesis of the
real-point lemma holds for $\square$.

\emph{Step 5: the levels on the shell $k_a$.} On the shell $\Omega_{k_a}\setminus\Omega_{k_a+1}$ every structure
since density $k_a$ has level $k_a$ and weight~$1$. Indeed, \cite{B-CMP122a} states that no changes are made in
the operation $\mathbf{T}_k(Z_k\cap Z^c)$ and omits it from the formulas; the stage formulas
\cite[(1.14)--(1.17)]{B-CMP122a} therefore act inside $Z\cap\Omega^c_{k_a}$ at density $k_a$. Consequently no
stage at density $k_a$ changes the determining set on $\Omega_{k_a}$. The later $\mathbf{T}$-steps change levels
inside
$\Omega_{k_a+1}$ only, and $\Omega_{k_a+1}$ is not reached, by Steps~3 and~4. Since
$\square_0\setminus W_{\mathfrak{a}}$ lies in this shell for the boxes of Steps~3 and~4,
\begin{equation}\label{4.26:eq-junction-boxes-4}
\mathfrak{B}\restriction(\square_0\setminus W_{\mathfrak{a}})
=\mathbb{B}^{(n')}_{k_a}\restriction(\square_0\setminus W_{\mathfrak{a}})\qquad\text{for every } n'.
\end{equation}

\emph{Step 6: the identity \eqref{4.26:eq-junction-boxes-1}.} On $\square_0\setminus W_{\mathfrak{a}}$ the
structure $\mathfrak{B}^{\mathfrak{c}}$ is $\mathfrak{B}$, because by Definition~\ref{4.26:def-graded-cover} only
$\mathfrak{B}\restriction W_{\mathfrak{a}}$ is replaced. A second live piece does not meet $\square_0$, by
Step~8 below. By \eqref{4.26:eq-junction-boxes-4} with $n'=k_0-h_a$, on $\square_0\setminus W_{\mathfrak{a}}$ the
structure $\mathfrak{B}^{\mathfrak{c}}$ is $\mathbb{B}^{(k_0-h_a)}_{k_a}$. On $\square_0\cap W_{\mathfrak{a}}$,
$\mathfrak{B}^{\mathfrak{c}}$ is the determining set of the strata of $\mathfrak{c}$, which by
Lemma~\ref{4.26:lem-original-level-strata}(i) is $\mathbb{B}^{(k_0-h_a)}_{k_a}\restriction W_{\mathfrak{a}}$.
The two pieces together give \eqref{4.26:eq-junction-boxes-1} as an identity of structures on $\square_0$, and
its right side is one historical stage structure of density $k_a$.

This stage structure exists. Since $k_0=k_a-N_0^a$ and $h_a=k_a-N_a$, the inequality $k_0>h_a$ is equivalent to
$N_a>N_0^a$, which is the inequality between these two numbers in \cite{B-CMP122a}. Since $j^+=h_a+n^+_a$, we
have $k_0-h_a<n^+_a$ if and only if
$k_0<j^+$. So $\mathbb{B}^{(k_0-h_a)}_{k_a}$ is a strict predecessor of the frozen structure
$\mathbb{B}^{(n^+_a)}_{k_a}$ exactly when $j^+>k_0$. This completes~(c).

\emph{Step 7: the case $j^+\le k_0$.} If $j^+\le k_0$, then the weight $\omega$ of
Definition~\ref{4.26:def-frozen-structure} is identically $1$ on $W_{\mathfrak{a}}$; this value of the weight in
the case $j^+\le k_0$ is given in \cite{B-CMP122a}. Then $\{\omega>1\}$ is empty, and by
Definition~\ref{4.26:def-graded-cover}, under
which
$\mathfrak{B}$ is unchanged off $\{\omega>1\}$, $\mathfrak{B}^{\mathfrak{c}}=\mathfrak{B}$ on $W_{\mathfrak{a}}$.
No structure is added. This is~(d).

\emph{Step 8: continuity, admissibility, separation.} In the case $j^+=k_a$, by \eqref{4.26:eq-junction-boxes-2}
the level of $\mathfrak{B}^{\mathfrak{c}}$ just inside $\partial W_{\mathfrak{a}}$ is $\mathfrak{c}=k_a-1$. By
Step~5 the level just outside is $k_a$. The levels are thus continuous across $\partial W_{\mathfrak{a}}$, and
admissibility follows from the admissibility property of the frozen structure on a live piece. The parent level
of the frame used for the increments is $\mathsf{m}^{\circ}=\mathfrak{c}$. Since
$\mathfrak{c}\le\ell_{\mathfrak{a}}$ by Lemma~\ref{4.26:lem-original-level-strata}(ii), the inequality
$\mathsf{m}^{\circ}=\mathfrak{c}\le\mathsf{m}$ persists. By the separation property of live pieces, disjoint
pieces are at least $R_k$ cubes of side $M$ apart, and $\square_0$ lies within seven sides of its own scale of
$\square$; so no set $\square_0$ with $\square\in\mathcal{C}^1$ meets two of them. This is~(e).

\emph{Step 9: the real point.} We apply the real-point lemma. Every added structure refines the structure
$\mathfrak{B}_{\ast}$ for which the real background $U_{\ast}$ is critical. This $\mathfrak{B}_{\ast}$ is the frozen
structure or, for the arrest branches in which $U_{\ast}$ is critical for the preceding stage structure,
$\mathbb{B}^{(n^+_a-1)}_{k_a}\restriction W_{\mathfrak{a}}$. Every added structure is a predecessor of it by
Definition~\ref{4.26:def-graded-cover} and Step~6.

Ba{\l}aban's nested averaging \cite{B-CMP98}, \cite[(2.11)--(2.12)]{B-CMP119}, \cite{B-CMP122a}, together with the
identification statements for backgrounds on localization domains, states the identity
\begin{equation}\label{4.26:eq-junction-boxes-5}
U_p\bigl(X',M^{\mathfrak{s}}U_{\ast}\bigr)=\bigl(U_{\ast}\restriction X'\bigr)^u ;
\end{equation}
we use it as stated there. Derived entry~$2$ of a triple follows from the plaquettes of $U_{\ast}$ itself, by the
direct
estimate of the derived-entry lemma. Derived entry~$4$ follows from that lemma's bound on entries~$2$ and~$4$,
by which they are multiples, depending on $d$ and $L$ only, of the small-field bound $\varepsilon_{\ell}$ of
\cite{B-CMP109} at the point's level $\ell$. At box triples and at junction triples the structures have two
levels, and there entry~$4$ is given by the two-level form of that bound. At the structures of
$(\Sigma^{\ast})^{\mathfrak{a}}$ with more than two levels entry~$4$ is this bound by multiples of
$\varepsilon_{\ell}$, used here as stated in \cite{B-CMP109},
exactly as in the arrest theorem.

The real-point bound then holds with margin $\tfrac12F^{(k)}$, in the frozen weighted unit, with the block side of
the bump statement. At junction boxes the real-point lemma applies through its junction hypothesis: for those of
scale $k_a$ it is verified in Step~4, and for those of scale $k_a-1$ the geometry is that of Step~3,
$\square_0\setminus W_{\mathfrak{a}}$ lying within seven sides of the cubes of $\pi_{k_a-1}$ of
$\partial\Omega_{k_a}$ inside the shell $\Omega_{k_a}\setminus\Omega_{k_a+1}$. In that hypothesis there are two
cells, the junction structure
satisfies
$\mathfrak{s}_1\preccurlyeq\mathfrak{B}_{\ast}$, and at $p=k_a$ the second of the three units at a retained
structure is at most $0.032$ times the third. The remaining conditions are the half-size $\gamma$-bound of the
derived-entry
lemma and $C_0^{\mathrm{rad}}\ge2.16A_0/\varphi_{\mathrm{arrest}}$ on the radius constant $C_0^{\mathrm{rad}}$ of
the arrested source space. Both are conditions already imposed for the
arrested source space.

\emph{Step 10: increments and tube.} The bump statement and its rescaled companion hold in every frame and at any
$(\Sigma^{\ast})$-triple. The frame $(X',\mathfrak{B},\mathfrak{B}^{\mathfrak{c}}\restriction X')$ falls under the
general clause of the frame statement, because $\mathsf{m}^{\circ}=\mathfrak{c}\le\mathsf{m}=\ell_{\mathfrak{a}}$
(Step~8).

For the tube, the proposition on the tube around the source configuration treats derived entries
$(\Sigma^{\ast})$ for every $(p,X',\mathfrak{s})$ of its space. For each such triple it works around the derived
configuration of the source, and it is generic in $\mathfrak{s}$. The scale-drop room of the arrest
construction is likewise generic in $\mathfrak{s}$. The corollary and the claim of the arrested construction on
the source factor $F^{(k')}$ do not depend on $\mathfrak{s}$.

\emph{Step 11: the source space.} The second condition of the arrested source space puts the same clauses into
$\tilde{U}^{\mathrm{arrest}}_k$. In that condition the block side at a retained structure is
$r_e\bigl(\ell(x),\min\{p,\lambda_{\mathfrak{s}}(x)\}\bigr)$, with $\ell(x)$ the level at $x$. The first clause of
the frozen weighted-unit display
is written clause by clause in the same unit, since the source factor satisfies $F^{(k)}\le1$, below the bound
$1.383$ of that first clause.

The clauses of the term classes, at their birth radii and with the derived entries taken in the form
$(\Sigma^{\flat})$, contain those of the
source. At $\lambda_{\mathfrak{s}}=\ell$ this is clause (a) of the radius comparison lemma. At
$\lambda_{\mathfrak{s}}<\ell$ it is clause (a$'$) of that lemma, by the ratio
\begin{equation}\label{4.26:eq-junction-boxes-6}
(2.625L)^{\ell-\lambda_{\mathfrak{s}}}\ge31.5\ge3/\varphi_{\mathrm{arrest}} .
\end{equation}
Steps~9--11 give~(f).

The inclusion of these triples in $(\Sigma^{\ast})^{\mathfrak{a}}$ is needed. Without it, the fourth clause of
the hypothesis of the construction of the stage spaces, read at lift $g=1$, would require the inequality
$2B_3\cdot0.381\le1+2\beta$, where $\beta$ is the constant in the factor $1+2\beta$ of the stage spaces and $B_3$
is the constant of the bound of \cite{B-CMP109} on which that clause rests; this inequality fails.
\end{proof}

\begin{remark}[Relation to the original-scale rule of an ordinary step]
\label{4.26:rem-original-scale-rule}
The rule by which the graded cover of Definition~\ref{4.26:def-graded-cover} reads structures
on a live piece is, written in the shell form of Ba{\l}aban's construction, the original-scale
rule already used for the shells of an abandoned run of an ordinary step. The goodness test that
accompanies that rule, which yields outright a lower bound on the native index over the enlarged
cube $\square_0$, is not used here. Accordingly the weight bounds, the ratio condition on the
radii and the ceiling condition on the analyticity size enter this chapter as stated, and the
lower bounds $L\ge 13$ and $M\ge L^4$ of that construction are not imposed.
\end{remark}

\begin{proof}
Consider an ordinary step whose large-field run is abandoned, and let $W$, $k_0$ and $j_W$ be
the data that the ordinary step attaches to this run, $j_W$ indexing its lifted stratum. The
shells it leaves are the thin shells $\Gamma''_n$ with $k_0<n<j_W$, together with the lifted
stratum $j_W$. On these shells the current index is $n$. By the original-scale rule of the
ordinary step, the structures read there are, however, the $\mathfrak{s}\in\{l-1,l\}$ with
$l\in[k_0,n]$. Thus $l$ is the original scale, and none of $W$, $k_0$, $j_W$ is ever read.

For the graded cover, the level function $\mathfrak{c}$ on a live piece $\mathfrak{a}$ is
\[
  \mathfrak{c}(y)=\ell_{\mathfrak{a}}(y)-\tfrac12\log_{L_0}\omega(y).
\]
By the comparison of indices for the ordinary step, $\mathfrak{c}$ agrees with the native index
$l$ of the rule just described wherever $\omega>1$. In
Definition~\ref{4.26:def-graded-cover}, $\mathfrak{B}^{\mathfrak{c}}$ replaces
$\mathfrak{B}\restriction W_{\mathfrak{a}}$ by the determining set of the strata of
$\mathfrak{c}$, and leaves it unchanged off $\{\omega>1\}$. Hence on $\{\omega>1\}$ the graded
cover reads the structure of the original scale, exactly as in the abandoned run. This proves
the first assertion.

For the second assertion, the goodness test in six layers is replaced by a counting remark. A
box of $\mathcal{C}^1$ meets at most one finer shell of $\mathfrak{c}$, and that shell is thick.
The constant $\mathsf{a}$ of the ceiling condition contains the factor $L^2$, and by the estimate
on $\mathsf{a}$ that accompanies the ceiling condition this factor already pays for this one
shell. No lower bound on the native index over $\square_0$ is therefore needed. For this reason
the weight bounds, the ratio condition and the ceiling condition on the analyticity size are
applied unchanged, and neither $L\ge 13$ nor $M\ge L^4$ enters.
\end{proof}

\begin{lemma}[Re-grouping the small-term pieces leaves every total unchanged]
\label{4.26:lem-regrouping-exact}
Let $T$ be any one of the terms of the representation of $\mathbb{C}^{(n)}_k(X,\mathcal{S})$, that
is, a term of class $\mathcal{T}$, with localization domain $X_T$. Let $U_\ell$,
$Q_\ell$, $\bar{U}^\ell$ and $\eta$ be as in \cite[(3.4)]{B-CMP109}, and let $\mathbb{H}$ be the
stage layer. For $t\in[0,1]$ put the configuration
\begin{equation*}
  \Phi^t:=U_\ell\big(\exp\big(iQ_\ell(t\eta\mathbb{H})\big)\bar{U}^\ell\big)
\end{equation*}
(here $\Phi^t$ denotes a configuration only and carries no current slot), so that $\Phi^0$ is the
real point. For a finite cover by cubes $\square$ of the region on which $\mathbb{H}$ is
supported, with a partition of unity $\{\zeta_{\square}\}$ on it, so that
$\sum_{\square}\zeta_{\square}=1$ there, define the piece of $T$ on $\square$ by
\begin{equation}\label{4.26:eq-regrouping-exact-1}
  \mathsf{P}(T,\square):=\int_0^1 dt\,
  \frac{\partial}{\partial t_{\square}}\Big|_{t_{\square}=0}
  T\Big(X_T;\,U_\ell\big(\exp\big(iQ_\ell(\eta(t+t_{\square}\zeta_{\square})\mathbb{H})\big)
  \bar{U}^\ell\big)\Big).
\end{equation}
Then, at the real point, for every such cover and every partition of unity on it,
\begin{equation}\label{4.26:eq-regrouping-exact-2}
  \sum_{\square}\mathsf{P}(T,\square)=T(\Phi^1)-T(\Phi^0).
\end{equation}
In particular, for the stage's cover $\mathcal{C}$ and the graded cover $\mathcal{C}^1$ of
Definition~\ref{4.26:def-graded-cover}, each with its partition of unity,
\begin{equation*}
  \sum_{\square\in\mathcal{C}^1}\mathsf{P}(T,\square)
  =\sum_{\square\in\mathcal{C}}\mathsf{P}(T,\square)
  =T(\Phi^1)-T(\Phi^0),
\end{equation*}
and, summing over the terms $T$, the total $\mathbb{C}^{(n)}_k(X,\mathcal{S})$ is the same
for both covers: passing from $\mathcal{C}$ to $\mathcal{C}^1$ changes only the grouping of
the small terms' pieces on the lifted sets $\mathsf{L}_{\mathfrak{a}}$. The identity is asserted
at the real point $\Phi^0$ only; the analytic extensions of the pieces are fixed separately, in
the remark following the proof.
\end{lemma}

\begin{proof}
The first-order contribution in \cite[(3.4)]{B-CMP109} is a sum over cubes of the
integrals \eqref{4.26:eq-regrouping-exact-1}: there is one global parameter $t$, and the
partition of unity enters linearly, through the factor $t+t_{\square}\zeta_{\square}$
multiplying $\eta\mathbb{H}$. Summed over all cubes and over all terms $T$, the pieces reproduce
the left side of \cite[(3.3)]{B-CMP109}, by the statement of this chapter on the sum of all
pieces of the first-order contribution; only the grouping of its far-field part depends on the
cover.

Fix $T$. The pieces \eqref{4.26:eq-regrouping-exact-1} are taken at the real point, in the form
fixed by the statement on the pieces of the first-order contribution at the real point. Their
definition through derivatives in $t_{\square}$ and in $t$ presupposes that the map
\begin{equation*}
  \phi_T(\mathsf{h}):=T\big(X_T;\,U_\ell(\exp(iQ_\ell(\mathsf{h}))\bar{U}^\ell)\big)
\end{equation*}
is continuously differentiable for $\mathsf{h}$ in a neighbourhood, in the space of layers, of
the segment $\{t\eta\mathbb{H}:t\in[0,1]\}$ and of its translates by
$t_{\square}\,\eta\zeta_{\square}\mathbb{H}$ for $|t_{\square}|$ small; we use it in this form.
Write $D\phi_T(\mathsf{h})[\mathsf{v}]$ for its derivative at
$\mathsf{h}$ in the direction $\mathsf{v}$, which is linear in $\mathsf{v}$. For fixed $t$ and
fixed $\square$,
\begin{equation*}
  \eta(t+t_{\square}\zeta_{\square})\mathbb{H}
  =t\eta\mathbb{H}+t_{\square}\,\eta\zeta_{\square}\mathbb{H},
\end{equation*}
so the map $t_{\square}\mapsto\eta(t+t_{\square}\zeta_{\square})\mathbb{H}$ is affine with
derivative $\eta\zeta_{\square}\mathbb{H}$, and the chain rule gives
\begin{equation*}
  \frac{\partial}{\partial t_{\square}}\Big|_{t_{\square}=0}
  \phi_T\big(\eta(t+t_{\square}\zeta_{\square})\mathbb{H}\big)
  =D\phi_T(t\eta\mathbb{H})[\eta\zeta_{\square}\mathbb{H}].
\end{equation*}
The cover is finite, so the sum over $\square$ may be taken inside the integral in $t$.
By linearity of $D\phi_T(t\eta\mathbb{H})[\,\cdot\,]$ and because
$\sum_{\square}\zeta_{\square}=1$, so that
$\sum_{\square}\eta\zeta_{\square}\mathbb{H}=\eta\mathbb{H}$, we obtain
\begin{equation*}
  \sum_{\square}D\phi_T(t\eta\mathbb{H})[\eta\zeta_{\square}\mathbb{H}]
  =D\phi_T(t\eta\mathbb{H})[\eta\mathbb{H}]
  =\frac{d}{dt}\phi_T(t\eta\mathbb{H}).
\end{equation*}
Thus $\sum_{\square}\partial_{t_{\square}}|_{0}=d/dt$ term by term in $T$, for any partition
of unity. Since $\phi_T(t\eta\mathbb{H})=T(\Phi^t)$ by definition, the fundamental theorem of
calculus gives
\begin{equation*}
  \sum_{\square}\mathsf{P}(T,\square)
  =\int_0^1\frac{d}{dt}T(\Phi^t)\,dt
  =T(\Phi^1)-T(\Phi^0),
\end{equation*}
which is \eqref{4.26:eq-regrouping-exact-2}. The right-hand side does not depend on the
cover or on the partition of unity, so the sums over $\mathcal{C}^1$ and over $\mathcal{C}$
coincide. Summing over the terms $T$ shows that the total $\mathbb{C}^{(n)}_k(X,\mathcal{S})$
is unchanged. Since $\mathcal{C}^1$ differs from $\mathcal{C}$ only by the graded
replacement of Definition~\ref{4.26:def-graded-cover}, what changes is the grouping of the
pieces on the lifted sets, and nothing else.
\end{proof}

\begin{remark}
The analytic extension of each piece is defined by the formulas of the proof family
(Definition~\ref{4.26:def-proof-family}) at the cube used; on the lifted sets
$\mathsf{L}_{\mathfrak{a}}$ these formulas are taken as the definition of the pieces' analytic
extensions, the function of the proof family being read with the graded cover and the
structures on a live piece (Definition~\ref{4.26:def-graded-cover} and
display~\eqref{4.26:eq-graded-cover-a}). The interpolation parameters are introduced term by
term, so that the families attached to two different terms never share a parameter; in
particular, the family attached to a term of the $M$-class, a class of terms of the
representation of $\mathbb{C}^{(n)}_k(X,\mathcal{S})$ other than $\mathcal{T}_1$, and the family
attached to a term of class $\mathcal{T}_1$ never share a parameter.
\end{remark}

Throughout, lengths at level $m$ are measured in $\pi_m$-sides, the sides of the cubes of Ba{\l}aban's level-$m$ partition; a $\pi_{m+1}$-side is $L$ $\pi_m$-sides. The length unit $\ell_m$ of level $m$ satisfies $\ell_{m+1}=L\ell_m$.

Recall that the shell radius $\mathbf{r}_{e,m}$ of an entry $e\in\{1,\dots,7\}$ at a level $m$, which enters the clauses of Definition~\ref{4.26:def-regularity-space}, is proportional to $\alpha_{\cdot,m}\ell_m^{-p_e}$, with lengths taken as absolute lengths. Here $\alpha_{\cdot,m}$ stands for $\alpha_{0,m}$ or $\alpha_{1,m}$, as the entry requires, and
\begin{equation*}
(p_1,\dots,p_7)=(2,2,3,2,2,1,2).
\end{equation*}
For $e=4$ the exponent is raised by one through a factor $\ell_{\min\{p,m\}}^{-1}$, where $p$ is the level of the triple at which the radius is read. At a point $y'$ we put $\mathbf{r}^{\mathfrak{c}}_e(y'):=\mathbf{r}_{e,\mathfrak{c}(y')}$. This radius is unweighted, and for $e\in\{2,4\}$ it is taken at the same $p$.

For a live arrest $\mathfrak{a}$ with $k_a<k$ and a point $y'\in W_{\mathfrak{a}}$, write
\begin{equation*}
g(y'):=\ell_{\mathfrak{a}}(y')-\mathfrak{c}(y').
\end{equation*}
By the table \eqref{4.26:eq-frozen-structure-a} of Definition~\ref{4.26:def-frozen-structure}, $\omega(y')=L_0^{2g(y')}$, and a point with $\omega(y')>1$ is called \emph{lifted}. By Lemma~\ref{4.26:lem-original-level-strata}(ii), $g\ge 0$. A \emph{$\mathfrak{c}$-shell} $m$ is the set $\{\mathfrak{c}=m\}$.

\begin{lemma}[Geometry and unit conversion for graded-cover boxes]\label{4.26:lem-graded-box-geometry}
Assume $R_{\min}\ge 5L$ and the adjacency condition $(\alpha\text{-adj})$ of \eqref{4.26:eq-inherited-conditions-a}. Assume further $L_0^2\le L/3$ and $L\ge 12$.

Let $\square\in\mathcal{C}^1$ be a cube of the graded cover (Definition~\ref{4.26:def-graded-cover}) of scale $i$, with $\square_0=\tilde{\square}^5$ and $[\square_0]_{\mathfrak{B}}\subset Y$. Suppose $\square_0$ meets a live older piece $W_{\mathfrak{a}}$ of $\mathfrak{T}$, where $\mathfrak{a}$ is a live arrest with $k_a<k$ (Definition~\ref{4.26:def-frozen-structure}); $Y$ and the stage $n\le N$ are those of $\mathfrak{N}^{(n)}(Y)$ (Definition~\ref{4.26:def-regularity-space}). Let $p_{\square}$ be the top level of $\mathfrak{B}^{\mathfrak{c}}\restriction\square_0$.
The clauses are lettered to match Lemma~\ref{4.26:lem-near-box-geometry}: clauses (a) and (c) below are the counterparts of clauses (a) and (c) there, and there are no clauses (b) and (d).
\begin{itemize}
\item[(a)] $\square_0$ meets at most two $\mathfrak{c}$-shells. They are adjacent, and $i$ is among them:
\begin{equation*}
i-1\le\mathfrak{c}\le p_{\square}\le i+1\quad\text{on }\square_0 .
\end{equation*}
If $(T,\square)$ is a near pair, with $X_T\subset\tilde{\square}^2$ and $T$ of level $\ell$, then $\ell\le i$ and $\ell\le\mathfrak{c}$ on $\square_0$.
\item[(a$'$)] The frozen level $\ell_{\mathfrak{a}}$ takes at most two values on $\square_0$, and they are adjacent.
\item[(c$_0$)] \textup{(Conversion.)} At every $y'\in\square_0\cap W_{\mathfrak{a}}$ and every entry $e\le 7$,
\begin{equation}\label{4.26:eq-graded-box-geometry-a}
r_e(y')=\omega(y')\,\mathbf{r}_{e,\ell_{\mathfrak{a}}(y')}\le 0.381^{\,g(y')}\,\mathbf{r}^{\mathfrak{c}}_e(y').
\end{equation}
For $e\in\{2,4\}$ this is to be read for $r_e^{(p)}$, at every $p$. Hence on $\square_0$ the fat clauses of $\mathfrak{N}^{(n)}(Y)$ in \eqref{4.26:eq-regularity-space-a} imply the unweighted fat clauses of the structure $\mathfrak{B}^{\mathfrak{c}}$ at the same factor $<1+2\beta-\sum\iota^1$. The fat clauses meant here are the direct ones, the clause $(\mathrm{box})^{\mathrm{sh}}$ at every $(\Sigma^{\ast})$-triple $X'\subset\square_0$, and the clause $(\mathrm{L}\text{-}\mathrm{b})$.
\item[(c)] With $\mathfrak{a}^{(i)}_0$ (entries $5,1$) and $\mathfrak{a}^{(i)}_1$ (entries $6,7$) as in clause (c) of Lemma~\ref{4.26:lem-near-box-geometry}, \eqref{4.26:eq-near-box-geometry-b}, all with $p_e\le 2$,
\begin{equation*}
\mathfrak{a}^{(i)}_{\cdot}\le(1+2\beta)(1+\beta_0)L^2\alpha_{\cdot,i}\le\mathsf{a}\,\alpha_{\cdot,i},
\qquad \mathsf{a}\ge(1+2\beta)(1+\beta_0)L^3 .
\end{equation*}
The same lower bound on $\mathsf{a}$ covers the conversion of entry $3$ ($p_3=3$) at the factor $(1+2\beta)(1+\beta_0)L^3$, where the core-freezing step makes it.
\item[(e)] \textup{(Cube clause.)} Let $q'\subset\square_0$ be a cube of Ba{\l}aban's family of cubes for $\mathfrak{B}^{\mathfrak{c}}$, of side $CM\ell_\nu$ with $\nu:=\min_{q'}\mathfrak{c}$ and $C\le 2$. Then $q'$ carries a $G$-valued $u$ with $U^u=e^{\mathrm{i}\eta A}$ and
\begin{equation}\label{4.26:eq-graded-box-geometry-1}
\ell_\nu|A|,\ \ \ell_\nu^2|\nabla^\eta A|<(1+2\beta)BCM\alpha_{0,\nu}.
\end{equation}
Here $U$, $\eta$, $A$, $\nabla^\eta$ and $B$ are as in the cube clause (clause $8$) of $\mathfrak{N}^{(n)}(Y)$ in \eqref{4.26:eq-regularity-space-a}.
\end{itemize}
\end{lemma}

\begin{proof}
Two counts from the proof of Lemma~\ref{4.26:lem-near-box-geometry} are used throughout. First, the box $\square_0$ of a cover cube of scale $i$ has side twelve $\pi_i$-sides. Second, every point of $\square_0$ lies within nine $\pi_i$-sides of every point of $\tilde{\square}^2$.

\emph{(a).} We run the argument of clause (a) of Lemma~\ref{4.26:lem-near-box-geometry} with $\mathfrak{c}$ in place of the shell level.

First we bound the widths of the $\mathfrak{c}$-shells. On $W_{\mathfrak{a}}$, for shells of index at least $h_a$, we use Lemma~\ref{4.26:lem-original-level-strata}(iii), \eqref{4.26:eq-original-level-strata-a}. Off $W_{\mathfrak{a}}$ and in its core we use the widths of the strata established in the proof of Lemma~\ref{4.26:lem-near-box-geometry}. Together these show that every $\mathfrak{c}$-shell $m''$ is at least $R_{\min}$ $\pi_{m''+1}$-sides wide. Since $R_{\min}\ge 5L\ge 60$, this exceeds twelve $\pi_{m''+1}$-sides.

Now suppose $\square_0$ met three $\mathfrak{c}$-shells. It would then meet three consecutive ones, $m''-1,m'',m''+1$, with $i$ among them, so that $i\le m''+1$. The box would cross the shell $m''$ within its side of twelve $\pi_i$-sides. That side is at most twelve $\pi_{m''+1}$-sides, which is less than the width of the shell $m''$; this is impossible. Hence $\square_0$ meets at most two $\mathfrak{c}$-shells, they are adjacent, and one of them is $i$, so $i-1\le\mathfrak{c}\le i+1$ on $\square_0$.

The level of $\mathfrak{B}^{\mathfrak{c}}$ at a point of $\square_0$ is the value of $\mathfrak{c}$ there, by Definition~\ref{4.26:def-graded-cover} and Lemma~\ref{4.26:lem-original-level-strata}(i). Its top level $p_{\square}$ on $\square_0$ therefore satisfies $\mathfrak{c}\le p_{\square}\le i+1$.

Next we show $\ell\le i$. The set $X_T$ is a union of $\pi_\ell$-cubes contained in $\tilde{\square}^2$, whose side is six $\pi_i$-sides. A $\pi_\ell$-cube with $\ell>i$ has side at least $L>6$ $\pi_i$-sides. Hence $\ell\le i$.

Finally we show $\ell\le\mathfrak{c}$ on $\square_0$. Suppose some point of $\square_0$ had $\mathfrak{c}<\ell$. Since $\mathfrak{c}\ge i-1\ge\ell-1$ there, that point lies in $\{\mathfrak{c}=\ell-1\}$, and the two shells met are $\ell-1$ and $i$, whence $i=\ell$. Take $y\in X_T$. Then $y\in\Lambda_\ell$ (as in the proof of Lemma~\ref{4.26:lem-near-box-geometry}(a)), and $\mathfrak{c}(y)=\ell$. The point of $\{\mathfrak{c}=\ell-1\}\cap\square_0$ lies within nine $\pi_\ell$-sides of $y$. We derive a contradiction, in two cases according to $\ell$.
\begin{itemize}
\item Let $h_a<\ell\le k_a$. Then $\Lambda_\ell$ lies at least $2R_\ell>9$ $\pi_\ell$-sides from $\{\mathfrak{c}<\ell\}$, by the last clause of Lemma~\ref{4.26:lem-original-level-strata}(iii); this contradicts the previous sentence.
\item Let $\ell\le h_a$. If $\square_0$ meets a lifted point, note that lifted points have $\mathfrak{c}\ge k_0$, by the table \eqref{4.26:eq-frozen-structure-a}: they lie in the thin strata $\Gamma''_m$ with $k_0<m<j^+$ or in $\Omega''_{j^+}\setminus\Omega_{k_0+1}$, where $\mathfrak{c}=k_0$, or in the old shells $\Omega_m\setminus\Omega_{m+1}$ with $k_0<m<j^+$, where $\mathfrak{c}=m$. By what was just proved, $\mathfrak{c}\ge k_0-1$ on $\square_0$. Moreover $k_0-1>h_a$, that is $N_a>N_0^a+1$; this follows from the condition $(N_0^a+1)(N_0^a+3)\le N_a$ on the integers $N_a$ and $N_0^a$ of the piece, which expresses the assumption stated in \cite{B-CMP122a}: ``We assume that $N>N_0$, in fact \dots\ $N$ is much greater than $N_0$''. So $\mathfrak{c}>h_a\ge\ell$ on $\square_0$, and $\{\mathfrak{c}=\ell-1\}$ is not met. If $\square_0$ meets no lifted point, it satisfies condition $(\mathrm{w}\text{-}\mathrm{sc})$ of \eqref{4.26:eq-near-box-geometry-a}, and Lemma~\ref{4.26:lem-near-box-geometry}(a) applies to it; off $\{\omega>1\}$ the structure $\mathfrak{B}^{\mathfrak{c}}$ coincides with $\mathfrak{B}$ (Definition~\ref{4.26:def-graded-cover}) and its level is $\mathfrak{c}$, so that lemma gives $\ell\le\mathfrak{c}$ on $\square_0$.
\end{itemize}

\emph{(a$'$).} We distinguish where the lifted points lie.
\begin{itemize}
\item Off $\{\omega>1\}$ we have $g=0$, so $\ell_{\mathfrak{a}}=\mathfrak{c}$. If $\square_0\cap W_{\mathfrak{a}}$ contains no lifted point, the claim is therefore (a).
\item In an old shell $\Omega_m\setminus\Omega_{m+1}$ with $k_0<m<j^+$, the table \eqref{4.26:eq-frozen-structure-a} gives $\ell_{\mathfrak{a}}\equiv j^+$.
\item The remaining lifted points lie in $\{\mathfrak{c}=k_0\}$.
\end{itemize}
The frozen strata in $\{\mathfrak{c}=k_0\}$ are the following.
\begin{itemize}
\item The stratum $\Gamma''_{k_0}$. It is at least one layer of $MR_{k_0+1}$-cubes of the $(k_0+1)$-lattice wide, by the widths of the strata in the proof of Lemma~\ref{4.26:lem-near-box-geometry}. That is at least $LR_{k_0+1}$ $\pi_{k_0}$-sides.
\item The thin strata $\Gamma''_m$ with $k_0<m<j^+$. Each is one $M$-cube of scale $m+1$ thick, that is $L^{m+1-k_0}\ge L^2$ $\pi_{k_0}$-sides.
\item The set $\Omega''_{j^+}\setminus\Omega_{k_0+1}$. It is five $M$-cubes of scale $k_a$ thick, because $Z''_{j^+}\subset(\Omega^{\sim5}_{k_0+1})^c$, and by the definition of $N_0^a$ an $MR_{k_0+1}$-cube of scale $k_0+1$ is an $M$-cube of scale $k_a$.
\end{itemize}
By the table, the frozen level is $m$ on $\Gamma''_m$ and $j^+$ on $\Omega''_{j^+}\setminus\Omega_{k_0+1}$. The latter set borders $Z''_{j^+}$, and so borders $\Gamma''_{j^+-1}$. Consecutive strata in this list therefore carry adjacent frozen levels. We now treat the boxes by scale.
\begin{itemize}
\item A box of scale $k_0$ has side $12<L^2$ $\pi_{k_0}$-sides. It therefore meets at most two of these strata, and if two then consecutive ones.
\item A box of scale $k_0+1$ sits in the old shell $k_0+1$. It reaches at most seven $\pi_{k_0+1}$-sides into $\{\mathfrak{c}=k_0\}$. This is less than the thickness $5L^{N_0^a-1}$ $\pi_{k_0+1}$-sides of $\Omega''_{j^+}\setminus\Omega_{k_0+1}$, and $5L^{N_0^a-1}\ge 5L$. Inside $\{\mathfrak{c}=k_0\}$ such a box therefore sees the level $j^+$ only.
\item A box of scale $k_0-1$ reaches at most $7/L$ $\pi_{k_0}$-sides into $\Gamma''_{k_0}$, which is less than the width $LR_{k_0+1}$ of $\Gamma''_{k_0}$; so it sees the frozen levels $k_0-1$ and $k_0$ only.
\end{itemize}
In every case $\ell_{\mathfrak{a}}$ takes at most two adjacent values on $\square_0$.

\emph{(c$_0$).} Put $\mu:=\ell_{\mathfrak{a}}(y')$ and $g=\mu-\mathfrak{c}(y')$. We first show
\begin{equation}\label{4.26:eq-graded-box-geometry-2}
\frac{\mathbf{r}_{e,\mu}}{\mathbf{r}_{e,\mathfrak{c}}}
\le\frac{\alpha_{\cdot,\mu}}{\alpha_{\cdot,\mathfrak{c}}}\,L^{-p_eg}
\le\Bigl(\frac{1+\beta_0}{L}\Bigr)^{g}.
\end{equation}
For $e\ne4$ the first inequality is the radius convention, since $\ell_\mu=L^g\ell_{\mathfrak{c}}$.

For $e=4$, at a triple of level $p$ with structure $\mathfrak{s}$, the length factors give
\begin{equation*}
\frac{\ell_\mu^{-2}\,\ell_{\min\{p,\lambda_{\mathfrak{s}}\}}^{-1}}
{\ell_{\mathfrak{c}}^{-2}\,\ell_{\min\{p,\mathfrak{c}\}}^{-1}}
=L^{-2g}\,\frac{\ell_{\min\{p,\mathfrak{c}\}}}{\ell_{\min\{p,\lambda_{\mathfrak{s}}\}}}\le L^{-2g}.
\end{equation*}
The last step holds because $\lambda_{\mathfrak{s}}\ge\mathfrak{c}$ at lifted points, so the final quotient is at most $1$. At the box triple $\lambda_{\mathfrak{s}}=\mathfrak{c}$, and the last lengths cancel. If the radius is read with $\min\{p,\mu\}$ in place of $\min\{p,\lambda_{\mathfrak{s}}\}$, the same holds, with $\ell_{\min\{p,\mathfrak{c}\}}/\ell_{\min\{p,\mu\}}\le1$ since $\mu\ge\mathfrak{c}$.

For the second inequality in \eqref{4.26:eq-graded-box-geometry-2}, first use $p_e\ge1$, so that $L^{-p_eg}\le L^{-g}$. The ratio $\alpha_{\cdot,\mu}/\alpha_{\cdot,\mathfrak{c}}$ is the product of the $g$ adjacent ratios $\alpha_{\cdot,m+1}/\alpha_{\cdot,m}$ with $\mathfrak{c}\le m<\mu$. The condition $(\alpha\text{-adj})$ bounds this product, against the power $L^{-g}$, by $((1+\beta_0)/L)^g$.

Next, $\omega=L_0^{2g}\le(L/3)^g$ by $L_0^2\le L/3$. Hence
\begin{equation*}
\omega\,\mathbf{r}_{e,\mu}\le\Bigl(\frac{1+\beta_0}{3}\Bigr)^{g}\mathbf{r}_{e,\mathfrak{c}}
=\Bigl(\frac{8}{21}\Bigr)^{g}\mathbf{r}_{e,\mathfrak{c}}\le 0.381^{\,g}\,\mathbf{r}_{e,\mathfrak{c}},
\end{equation*}
where $\beta_0=1/7$ and $8/21<0.381$. This is \eqref{4.26:eq-graded-box-geometry-a}. Read at the single point $y'$, it is the property stated in \cite{B-CMP122a}: ``They form a descending sequence for increasing $n$, if \dots\ $\beta\le1/4$ and $L_0^2\le(1/3)L$. Notice that we get the stronger restriction on $L_0$ because of the bounds for $A'$''.

For the second sentence of (c$_0$): each fat clause of $\mathfrak{N}^{(n)}(Y)$ on $\square_0$ --- direct, $(\mathrm{box})^{\mathrm{sh}}$ at a $(\Sigma^{\ast})$-triple $X'\subset\square_0$ with $r_e^{(p)}$ at every $p$, or $(\mathrm{L}\text{-}\mathrm{b})$ --- bounds its entry by the factor times $r_e(y')$.
By \eqref{4.26:eq-graded-box-geometry-a} and $0.381^{g}\le1$, each of these bounds implies the bound by the same factor times $\mathbf{r}^{\mathfrak{c}}_e$, which is the unweighted clause of $\mathfrak{B}^{\mathfrak{c}}$. The increments $\iota^1$ in the factor are the same numbers on both sides.

\emph{(c).} By (a), the level $\mathfrak{c}$ of every cube of $\tilde{\square}^4\subset\square_0$ lies in $\{i-1,i,i+1\}$. By (c$_0$), every such cube obeys the unweighted clause of its level $\mathfrak{c}$ at factor at most $1+2\beta$.

These are the three cases treated in the proof of clause (c) of Lemma~\ref{4.26:lem-near-box-geometry}, with the bounds measured in the level-$i$ unit.
\begin{itemize}
\item Cubes of level $i$ contribute at most $1+2\beta$.
\item Cubes of level $i+1$ contribute less.
\item Cubes of level $i-1$ contribute at most $(1+2\beta)(1+\beta_0)L^{p_e}\le(1+2\beta)(1+\beta_0)L^2$, because the entries $5,1,6,7$ have $p_e\le2$.
\end{itemize}
This gives the first inequality. The second follows from $\mathsf{a}\ge(1+2\beta)(1+\beta_0)L^3\ge(1+2\beta)(1+\beta_0)L^2$.

\emph{(e).} Put $n':=\min_{q'}\ell_{\mathfrak{a}}$. Since $\ell_{\mathfrak{a}}\ge\mathfrak{c}$ (Lemma~\ref{4.26:lem-original-level-strata}(ii)), we have $n'\ge\nu$.

\emph{Case $n'=\nu$.} At a point of $q'$ with $\ell_{\mathfrak{a}}=\nu$ we have $\nu\le\mathfrak{c}\le\ell_{\mathfrak{a}}=\nu$. Hence $\mathfrak{c}=\nu$ and $\omega=1$ there. The cube $q'$ lies in $\{\ell_{\mathfrak{a}}\ge\nu\}$ and meets $\{\ell_{\mathfrak{a}}=\nu\}$. By (a$'$) it misses $\{\ell_{\mathfrak{a}}\ge\nu+2\}$. It is therefore a cube of the frozen family at level $\nu$, and its weight is $1$: by \cite[(1.65)]{B-CMP122a} if $\nu\le k_0$, by \cite[(1.67)]{B-CMP122a} if $\nu=j^+$, and by \cite[(1.69)]{B-CMP122a} otherwise. The cube clause of $\mathfrak{N}^{(n)}(Y)$ at $q'$ gives the assertion.

\emph{Case $n'>\nu$.} Then $2M\ell_\nu<M\ell_{n'}$. Since $\{\ell_{\mathfrak{a}}\ge n'\}$ contains $q'$ and is a union of $\pi_{n'}$-cubes, every $\pi_{n'}$-cube meeting $q'$ lies in it; these cubes form a box $P_{q'}\subset\{\ell_{\mathfrak{a}}\ge n'\}$. Choose a cube $Q$ of side $M\ell_{n'}$ with $q'\subset Q\subset P_{q'}$. This $Q$ meets $\{\ell_{\mathfrak{a}}=n'\}$. Hence $Q$ is a cube of the frozen family at level $n'$ with $C=1$, since no alignment of the cubes is required in \cite[(2.38)]{B-CMP119}. Its weight $w_Q$ is one of the following.
\begin{itemize}
\item $w_Q=L_0^{2(n'-k_0)}$, by \cite[(1.65)]{B-CMP122a}. In this case $q'$ holds a point with $\mathfrak{c}=k_0$, so $k_0\ge\nu$.
\item $w_Q=L_0^{2(j^+-\min_Q\mathfrak{c})}$, by \cite[(1.67)]{B-CMP122a}, with $j^+=n'$. In this case each $\pi_{j^+}$-cube of $P_{q'}$ lies in one old shell. Indeed, $\partial\Omega_m$ for $m>k_0$ runs along cubes of side $MR_m\ell_m\ge M\ell_{k_a}\ge M\ell_{j^+}$, by \eqref{4.26:eq-lifted-counts-a}. Each such cube also meets $q'$, so $\mathfrak{c}$ on it equals a value of $\mathfrak{c}$ on $q'$. Hence $\min_Q\mathfrak{c}\ge\nu$.
\end{itemize}
In both cases
\begin{equation*}
w_Q\le L_0^{2(n'-\nu)}.
\end{equation*}
On $Q$ the cube clause of $\mathfrak{N}^{(n)}(Y)$ supplies a $G$-valued $u$ with $U^u=e^{\mathrm{i}\eta A}$, at level $n'$ and weight $w_Q$. We restrict $u$ to $q'$. Since $\ell_\nu=L^{-(n'-\nu)}\ell_{n'}$, the weight $w_Q\le(L/3)^{n'-\nu}$, the factors $L^{-(n'-\nu)}$ and $L^{-2(n'-\nu)}$, both at most $L^{-(n'-\nu)}$, and the ratio $\alpha_{0,n'}/\alpha_{0,\nu}$ combine as in the proof of (c$_0$). This gives
\begin{align*}
\ell_\nu|A|&=L^{-(n'-\nu)}\ell_{n'}|A|
<(1+2\beta)\,0.381^{\,n'-\nu}\,BM\alpha_{0,\nu},\\
\ell_\nu^{2}|\nabla^\eta A|&=L^{-2(n'-\nu)}\ell_{n'}^{2}|\nabla^\eta A|
<(1+2\beta)\,0.381^{\,n'-\nu}\,BM\alpha_{0,\nu}.
\end{align*}
This is the bound \eqref{4.26:eq-graded-box-geometry-1} for the cube $q'$.
\end{proof}

\begin{proposition}[Near pairs on the graded cover at untouched clauses]\label{4.26:prop-near-pair-graded}
Let $n$ be a stage, let $Y$ be a union of cubes and let $(\mathbb{U},\mathbb{J})\in\mathfrak{D}^{\flat}_n(Y)$.
Let $(T,\square)$ be a near pair in which:
\begin{itemize}
\item $\square\in\mathcal{C}^1$ is a cube of the graded cover of Definition~\ref{4.26:def-graded-cover},
carrying the top level $p_{\square}$ and the structure $\mathfrak{s}$ fixed in \eqref{4.26:eq-graded-cover-a};
\item $T\in\mathcal{T}_1\cap\mathcal{T}^{\mathbb{E}\mathbb{R}}$ has level $\ell$, $X_T\subset\tilde{\square}^2$,
and $[\square_0]_{\mathfrak{B}}\subset Y$, where $\square_0=\tilde{\square}^5$.
\end{itemize}
Let $(y,a)\in\operatorname{Cl}(T)$ with $y\in X_T\cap\mathfrak{T}$. Assume the standing conditions of
Proposition~\ref{4.26:prop-near-pairs-outside}, among them the window condition of
Definition~\ref{4.26:def-radii-window}, the conditions \eqref{4.26:eq-inherited-conditions-a} and the
ordering \eqref{4.26:eq-near-pairs-outside-a}; the condition \eqref{4.26:eq-near-box-geometry-a} on the
geometry of the box is not assumed. Then (H) below holds, and the conclusion of
Proposition~\ref{4.26:prop-near-pairs-outside} holds in this setting, in the form (M) below:
\begin{itemize}
\item[(H)] the input conditions (i)--(iv) of Ba{\l}aban's treatment of the near term in
\cite{B-CMP109}, condition (iv) being the one displayed in \cite[(1.16)]{B-CMP109}, hold on
$\square_0\cap\mathfrak{T}$ literally in the form stated there, with a factor smaller than
$1+2\beta$, in the unit of the shell of $\mathfrak{c}$;
\item[(M)] every member $P(\varsigma;t,\tau;B')$ of the family (c) of
Proposition~\ref{4.26:prop-near-pairs-outside}, and each of the two remainder configurations (c$'$)
of Lemma~\ref{4.26:lem-near-remainders}, satisfies $Q_a(\cdot)(y)<\theta_{T,a}(y)$.
\end{itemize}
\end{proposition}

\begin{proof}
If $\square_0$ meets no live older piece of $\mathfrak{T}$, then $\square\in\mathcal{C}$,
$\mathfrak{s}=\mathbb{B}^{(n)}_k\restriction\square_0$, and condition
\eqref{4.26:eq-near-box-geometry-a} holds; the assertion is then
Proposition~\ref{4.26:prop-near-pairs-outside} itself.

Assume from now on that $\square_0$ meets a live older piece of $\mathfrak{T}$. Since moreover
$[\square_0]_{\mathfrak{B}}\subset Y$, Lemma~\ref{4.26:lem-graded-box-geometry}, the geometry lemma
for graded-cover boxes, applies to $\square$: its conditions $R_{\min}\ge5L$ and the condition on the
radii $\alpha_{\cdot,m}$ at adjacent levels are both part of \eqref{4.26:eq-inherited-conditions-a},
and its lower bounds on $L$ are among the floors on $L$ fixed for this chapter. Its
clauses (a), (a$'$), (c), (e) and its conversion (c$_0$) are used below. We go through the proof of
Proposition~\ref{4.26:prop-near-pairs-outside} step by step, in its order. When we say that a step
applies unchanged, we mean that it depends on no quantity other than those named in it.

\smallskip\noindent\textit{Step (H).}
Condition (i) holds for entry $5$ by the conversion (c$_0$) of
Lemma~\ref{4.26:lem-graded-box-geometry}, and for the cube clause by clause (e) of that lemma.
Condition (ii) concerns entries $6$ and $7$, and condition (iii) entries $1$ and $3$; for these the
clauses of $\mathfrak{N}^{(n)}(Y)$ of Definition~\ref{4.26:def-regularity-space}, converted by
(c$_0$), give the required bounds.

Condition (iv) is required \emph{on} $\tilde{\square}^3$, which in the notation $X'^{\sim2}$ of the
box clause \eqref{4.26:eq-inherited-domain-a} is $\square_0^{\sim2}$: in
\cite{B-CMP109} it is imposed on the function $U_j$ constructed for the cube
$\square_0=\tilde{\square}^5$ as in condition (iv), whose form is \cite[(1.16)]{B-CMP109}. We first
show that for every $p\le p_{\square}$ the triple $(p,\square_0,\mathfrak{s})$ is a
$(\Sigma^{\ast})$-triple. Recall that a $(\Sigma^{\ast})$-triple is a triple $(p,X',\mathfrak{s})$
in which $p$ does not exceed the top level, $X'\subseteq Y$ is a localization domain
($X'\in\bigcup_{p'}\mathbf{D}_{p'}$), and $\mathfrak{s}$ is a determining-set structure on $X'$ among the
current one and all its predecessors in the construction (the stage structures of stages $n'\le n$,
before and after the operation $\mathbf{R}$), localized to $X'$ as Ba{\l}aban localizes it, namely
as $\mathbb{B}_p(X')\cup\{\Gamma^{\mathfrak{s}}_{i'}\}_{i'<p}$. Here $\square_0\in\mathbf{D}_i$, $i$ being the
scale of $\square$, and $\mathfrak{s}=\mathfrak{B}^{\mathfrak{c}}\restriction\square_0$ belongs to the
list $(\Sigma^{\ast})^{\mathfrak{a}}$ of \eqref{4.26:eq-graded-cover-a}, by the description of that
list for a live piece: it is a stage structure of
the live piece's own run continued by $\mathfrak{B}$, junction boxes included, by clause (i) of
\eqref{4.26:eq-original-level-strata-a}. Hence the composite shell clause of
\eqref{4.26:eq-regularity-space-a} is a clause of $\mathfrak{N}^{(n)}(Y)$ at this triple, since that
clause is imposed at every $(\Sigma^{\ast})$-triple $(p,X',\mathfrak{s})$ with $X'\subset Y$.

This clause is propagated along the proof family by the frame statement for determining sets, because
the triple is a frame in its sense: there $\mathfrak{A}$ is a determining set of the construction, the
unit set, of level $\mathsf{m}$, and
$\mathfrak{B}^{\ast}=\mathbb{B}_p(X')\cup\{\Gamma^{\circ}_{i'}\cap X'\}_{i'<p}$ is
built over a parent sequence $\mathfrak{A}^{\circ}$ of level $\mathsf{m}^{\circ}\le\mathsf{m}$ on
$X'$, $\Gamma^{\circ}_{i'}$ denoting the strata of $\mathfrak{A}^{\circ}$. We take
$\mathfrak{A}=\mathfrak{B}$, $\mathfrak{A}^{\circ}=\mathfrak{B}^{\mathfrak{c}}$,
$\mathsf{m}^{\circ}=\mathfrak{c}\le\mathsf{m}=\ell_{\mathfrak{a}}$. The ball $B(x)$ of radius
$\varrho_0\ell_{\ell_{\mathfrak{a}}(x)}$, $\varrho_0$ being the ball constant of the frame statement
and $\ell_m$ the block side at level $m$,
may leave $\square_0$; the condition of the frame statement
on the balls $B(x)$ then
holds by the convention \eqref{4.26:eq-regularity-space-b}, as at drops of scale. The clause is
stated in the frozen weighted unit $r^{(p)}_e(y')$ at each $y'\in\square_0\cap W_{\mathfrak{a}}$,
whose block side is
$\min\{p,\lambda_{\mathfrak{s}}\}=\min\{p,\mathfrak{c}\}$ here, by the rule fixing the block side of
the frozen weighted unit.

By the conversion (c$_0$), the clause implies \cite[(1.16)]{B-CMP109} in the unit of the shell of
$\mathfrak{c}$. Indeed, for $e\in\{2,4\}$ and $y'\in\square_0\cap W_{\mathfrak{a}}$, with
$m=\ell_{\mathfrak{a}}(y')$ the frozen level and
$g=g(y')$, the lengths of the block sides cancel and
\begin{equation*}
\frac{r^{(p)}_e(y')}{\mathbf{r}^{\mathfrak{c},(p)}_e}
=\omega(y')\,\frac{\alpha_{0,m}}{\alpha_{0,\mathfrak{c}}}\Bigl(\frac{\ell_{\mathfrak{c}}}{\ell_m}\Bigr)^2
\le\Bigl(\frac{(1+\beta_0)L_0^2}{L^2}\Bigr)^{g}
\le\Bigl(\frac{0.381}{L}\Bigr)^{g}.
\end{equation*}
The resulting bounds are
\begin{align*}
|\partial U_p(M^{\mathfrak{s}}\mathbb{U})-1|
&<(1+2\beta)\,\alpha_{0,\mathfrak{c}}\,\eta^2\ell_{\mathfrak{c}}^{-2},\\
|J_p(M^{\mathfrak{s}}\mathbb{U})|
&<(1+2\beta)\,\alpha_{0,\mathfrak{c}}\,\ell_{\mathfrak{c}}^{-2}\ell_{p\wedge\mathfrak{c}}^{-1}
\qquad\text{on }\tilde{\square}^3,
\end{align*}
where $\eta$ is the spacing parameter of \cite{B-CMP109}, replaced by $L^{-i}$ at the scale of
$\square$ as in step (n4);
and the same bounds hold for the configuration $\mathbb{U}$ itself, by the corresponding clause of
\eqref{4.26:eq-regularity-space-a}. By clause (a) of Lemma~\ref{4.26:lem-graded-box-geometry},
$\square_0$ meets at most two shells of $\mathfrak{c}$, and these are adjacent; so these are
hypotheses on at most two adjacent shells, thick by \eqref{4.26:eq-original-level-strata-a}, of an
admissible set. This is exactly what clause
(c) of the stability hypothesis \eqref{4.26:eq-stage-systems-b} supplies on a pair of shells made by
the operation $\mathbf{T}$. This proves (H).

\smallskip\noindent\textit{Step (n1): entries $5$ and $8$.}
This step of the proof of Proposition~\ref{4.26:prop-near-pairs-outside} uses only the structure of
$\mathfrak{U}_T$ as a union of $G^c$-orbits and the identities \cite[(3.37)]{B-CMP109} and
\cite[(3.50)]{B-CMP109}. It applies unchanged: entries $5$ and $8$ of $P$ lie inside every positive
threshold.

\smallskip\noindent\textit{Step (n2): entries $6$ and $7$.}
Every $T\in\mathcal{T}^{\mathbb{E}\mathbb{R}}$ keeps Ba{\l}aban's one-scale space
$U^c_{\ell}(X_T,\cdot)$. Indeed $X_T\subset\Lambda_{\ell}$. If $\ell>k_a$, then $X_T$ misses
$W_{\mathfrak{a}}\subset\Omega^c_{k_a}$, so the modification made by the arrest on $W_{\mathfrak{a}}$
changes nothing for $T$. If $\ell\le k_a$, the clauses of $T$ are fixed before the arrest. At a
lifted $y$ one has $\ell\le\mathfrak{c}(y)$. Consequently $\mathbf{r}^T_a=\mathbf{r}_{a,\ell}$,
unweighted, and
\begin{equation*}
\theta_{T,a}(y)=s_T\,\mathbf{r}^T_a\ge(1-\beta)\,\mathbf{r}^T_a .
\end{equation*}

$(\alpha)$ \emph{The zero stage layer.} As with every cost of the near pair, the input-size constant
and the radius entering $\varkappa_{\square}$ are taken at the scale $i$ of $\square$. With clause (c) of
Lemma~\ref{4.26:lem-graded-box-geometry}, this gives
\begin{align*}
\varkappa_{\square}
&=B_3^2\,O(1)\,M\,\frac{\mathfrak{a}^{(i)}_0}{\alpha_{1,i}}\,\frac{r^{[i]}_a}{\mathbf{r}^T_a}\\
&\le B_3^2\,O(1)\,M\,(1+2\beta)(1+\beta_0)L^2\,\frac{C_0}{C_1}
\Bigl(\frac{1+\beta_0}{L}\Bigr)^{i-\ell}
\le\tfrac12(1-\beta),
\end{align*}
the last inequality by the ordering \eqref{4.26:eq-near-pairs-outside-a} in the form displayed
there. Here $B_3$ and $O(1)$ are the constants appearing in \cite[(3.26)--(3.27)]{B-CMP109} and
\cite[(3.37)]{B-CMP109}; $r^{[i]}_a$ denotes the radius of entry $a$ evaluated at the scale $i$; the
exponents $q_0$
and $q_1$ of the radii $\alpha_{0,\cdot}$ and $\alpha_{1,\cdot}$ coincide, $q_0=q_1$; $\ell\le i$;
and clause (c) is used with its exponent $p_e$ equal to $2$, the value belonging to entries $5$
and $1$. Only the power $L^2$ is needed in the middle line; the ordering displayed next, whose
constant $\mathsf{a}$ carries $L^3$, implies the last inequality a fortiori.

The ordering \eqref{4.26:eq-near-pairs-outside-a} is one-sided:
\begin{equation*}
\frac{C_1}{C_0}\ge\frac{2\,\mathsf{a}\,B_3^2\,O(1)\,M}{1-\beta},
\qquad \mathsf{a}=L^3(1+2\beta)(1+\beta_0)^4 .
\end{equation*}
The ratio $C_1/C_0$ of the two radius constants is chosen after $L$ and $M$ (and after
$B_3=B_3(d,L)$ and $\mathsf{a}=\mathsf{a}(L)$, which its right-hand side reads) and before
$\gamma$, at the place of the radius exponents $q_0,q_1$ and constants $C_0,C_1$ in the fixed order
in which the auxiliary parameters are chosen. Read in this way, it bounds $C_0/C_1$ from above only
and places no condition on $L$ or
$M$; in particular it is not a cap in the sense of Definition~\ref{2.3:def-cap}.

$(\beta)$ \emph{The stage layer, in the units of the shell of $\mathfrak{c}$, cube by cube.}
The functions $\zeta_{\square}$ and $\zeta'_{\square}$ are the functions of the cube $\square$ at its
own scale, for the cover $\mathcal{C}^1$ built on the admissible set $\mathfrak{B}^{\mathfrak{c}}$.
At most five cubes of this cover overlap at any point, by the overlap bound for the covers of
admissible sets; this is the factor $5$ below. Fix $y'\in\square_0$ and put
$n':=\ell_{\mathfrak{a}}(y')$ and $g:=g(y')$. A field has first-order size at level $\mathfrak{c}$ at
most $L^{-g}$ times its first-order size at level $n'$. Let $\mathbb{H}^w$ be the stage layer taken
at the localized pair $(\mathbb{U}^w_{\square},\mathbb{J})$. Therefore
$N^{\mathfrak{c}}_{y'}(\zeta_{\square}\mathbb{H}^w)$, the first-order size of
$\zeta_{\square}\mathbb{H}^w$ at $y'$ measured in the unit of level $\mathfrak{c}$, is at most
$2\cdot5\cdot L^{-g}B_{\sharp}$ times
the first-order majorant of the stage layer at $y'$ at level $n'$, and
\begin{align*}
N^{\mathfrak{c}}_{y'}(\zeta_{\square}\mathbb{H}^w)
&\le 11\,B_{\sharp}\,\vartheta(\square)\,\rho_{n'}\,L^{-g}\alpha_{\cdot,n'}\\
&\le 11\,B_{\sharp}\,\vartheta(\square)\,\rho_{\ell_{\mathfrak{a}}(y')}\,\alpha_{\cdot,\mathfrak{c}(y')}.
\end{align*}
The first inequality uses the first-order bound on the stage layer at the localized pair
$(\mathbb{U}^w_{\square},\mathbb{J})$, the absolute bound taken at the frozen level (which holds for
every admissible $\mathfrak{B}$), and $d^-_{\square}\le d^{\wedge}(\square^{\zeta},\mathcal{C}_j)$,
where $j$ and the set $\mathcal{C}_j$ of level-$j$ bonds of the collar are as in
Lemma~\ref{4.26:lem-family-consumption}.
The second uses
\begin{equation*}
L^{-g}\alpha_{\cdot,n'}\le\Bigl(\frac{1+\beta_0}{L}\Bigr)^{g}\alpha_{\cdot,\mathfrak{c}} .
\end{equation*}
This is the bound on the stage layer used for cubes of the cover $\mathcal{C}$, with
$\alpha_{\cdot,m}$ replaced by $\alpha_{\cdot,\mathfrak{c}}$. By clause (a$'$) of
Lemma~\ref{4.26:lem-graded-box-geometry}, the frozen level takes at most two adjacent values on
$\square_0$, so $\rho_{\ell_{\mathfrak{a}}(\cdot)}$ varies over $\square_0$ as it does over a pair of
shells made by $\mathbf{T}$.

Two ingredients are now used. The first is the formula $K^{\mathrm{nr}}=66B_{\sharp}K_{\mathrm{cv}}$
of Definition~\ref{4.26:def-radii-window}, together with the bound $\bar{\rho}_{\square}\le2\rho_m$
over the two frozen levels of clause (a$'$). The second is the constant $K_{\mathrm{cv}}$, taken at
the box $\square_0$ of $\mathfrak{B}^{\mathfrak{c}}$. Together they give the near line of clause (v) of
Lemma~\ref{4.26:lem-family-consumption}. With $m=\ell_{\mathfrak{a}}(y)$, the dilated member moves
each entry by at most $\tfrac13\mathsf{w}_0\rho_m\mathbf{r}^T_a$, and the undilated member by at most
$\tfrac{1}{18}\mathsf{w}_0\rho_m\mathbf{r}^T_a$. Together the two moves are at most
$\iota^1_n(y)\,\mathbf{r}^T_a$, where
$\iota^1_n(y)=\tfrac23\mathsf{w}_0\rho^{(n)}_{\ell_{\mathfrak{a}}(y)}$ is the increment of
\eqref{4.26:eq-older-pieces-domain-a}.

$(\gamma)$ \emph{The perturbation $B'$.} Its contribution is at most $\mathrm{room}_T/12$,
$\mathrm{room}_T$ being the room of \eqref{4.26:eq-radii-window-c}. This
bound depends only on quantities attached to $T$ itself and applies unchanged.

The closing inequality \eqref{4.26:eq-near-pairs-outside-b}, $0.4+0.006+0.1<0.8\le s_T$, applies
unchanged.

\smallskip\noindent\textit{Step (n3): entries $1$ to $4$.}
\emph{Start.} We use \cite[(3.39)]{B-CMP109} at zero stage layer and $\varsigma=1$:
\begin{equation*}
U_{\ell}\bigl(\square_0,M_{\ell}(\mathbb{U}^{\mathfrak{s}}_{\square})\bigr)
=(\mathbb{U}^{\mathfrak{s}}_{\square})^{(vu)} ,
\end{equation*}
where the superscript $(vu)$ denotes the action of the gauge transformations $v$ and $u$ of
\cite[(3.26)--(3.27)]{B-CMP109}, as in \cite[(3.39)]{B-CMP109}. Here $M_{\ell}$ is the level-$\ell$
averaging, and $U_{\ell}(\square_0,\cdot)$ is the level-$\ell$ minimiser on $\square_0$ constrained
by the given level-$\ell$ averages, as in \cite[(3.39)]{B-CMP109}: the identity says that the
minimiser constrained by the $\ell$-averages of the $\mathfrak{s}$-constrained minimiser
$\mathbb{U}^{\mathfrak{s}}_{\square}$ is that minimiser.
The constraints are nested, because $\ell$ does not exceed the levels of $\mathfrak{s}$ on
$\square_0$ by clause (a) of Lemma~\ref{4.26:lem-graded-box-geometry}. Ba{\l}aban's identity is used
as stated in \cite[(3.39)]{B-CMP109}, at a structure $\mathfrak{s}$ with one or two levels, as on
shells made by $\mathbf{T}$.

The box clause \eqref{4.26:eq-inherited-domain-a} at $(p_{\square},\square_0,\mathfrak{s})$ is a
clause of $\mathfrak{I}_n$. Indeed $(p_{\square},\mathfrak{s})\in(\Sigma^{\ast})^{\mathfrak{a}}$ by
\eqref{4.26:eq-graded-cover-a} and the description of the list $(\Sigma^{\ast})^{\mathfrak{a}}$ for a
live piece, and $\mathfrak{I}_n$ imposes the box clause at every localization
domain $X'\subset Y$ with $X_T\subset X'^{\sim2}$ and every $(p,\mathfrak{s})$ of $(\Sigma^{\ast})$;
here $X'=\square_0$, which lies in $Y$, and $X_T\subset\tilde{\square}^2$, while
$\tilde{\square}^2\subset\tilde{\square}^3$ with $\tilde{\square}^3=\square_0^{\sim2}$.
We have $p_{\square}\ge i$ and $i\ge\ell$, and the unit rule
\eqref{4.26:eq-clauses-units-b} assigns the level $\lambda=\ell$ and hence the unit
$\mathbf{r}^T_{a'}$ to the clause. For $a'\in\{2,4\}$ this gives
\begin{align*}
Q^{(p_{\square},\square_0,\mathfrak{s})}_{a'}(y)
&<s^{(n)}_{T,a'}(y)\,\mathbf{r}^{T,(p_{\square})}_{a'}(y)=\theta^{(n)}_{T,a'}(y)\\
&=\Bigl[s_T-\tfrac13\mathrm{room}_{a'}-\sum_{l'\le n}\iota^1_{l'}(y)\Bigr]\mathbf{r}^T_{a'}(y),
\end{align*}
where we put
\begin{equation*}
\mathrm{room}_{a'}:=\frac{\theta_{T,a'}-\theta_{\mathrm{src},a'}}{\mathbf{r}^T_{a'}} .
\end{equation*}
We claim that $\mathrm{room}_{a'}\ge\mathrm{room}_T\ge\beta\hat{c}$, the last inequality being
\eqref{4.26:eq-radii-window-c}.
\begin{itemize}
\item At a lifted $y$, the conversion (c$_0$) and $\ell\le\mathfrak{c}$ give
$\theta_{\mathrm{src},a'}\le r_{a'}(y)\le0.381\,\mathbf{r}^T_{a'}$. Since
$\theta_{T,a'}=s_T\mathbf{r}^T_{a'}$ with $s_T\ge0.8$ by \eqref{4.26:eq-near-pairs-outside-b}, one
gets $\mathrm{room}_{a'}\ge0.8-0.381>0.4$, as in the second case of
Definition~\ref{4.26:def-radii-window}.
\item At an un-lifted $y$, the claim is clause (5) of \eqref{4.26:eq-older-pieces-domain-a}.
\end{itemize}

\emph{Chain.} The costs $G_{\square}$, $\mathfrak{q}_1$, $\mathfrak{q}_2$, $\mathfrak{q}_1'$,
$\varepsilon_{\square}$, $G'_{\square}$ of \eqref{4.26:eq-near-pair-chain-a} depend only on
$\mathfrak{a}^{(i)}$, $\alpha_{\cdot,i}$ and constants. By clause (c) of
Lemma~\ref{4.26:lem-graded-box-geometry},
$B_3^2\,O(1)\,M\,\mathfrak{a}^{(i)}_0\le\mathsf{S}\,\alpha_{0,i}$. The five majorants of the ceiling
\eqref{4.26:eq-inherited-conditions-a} are linear in $\alpha_{\cdot,i}\le\bar{\alpha}$. Hence the
ceiling \eqref{4.26:eq-inherited-conditions-a} implies the local ceiling
\eqref{4.26:eq-near-pair-chain-b} unchanged; it contains no member involving $w_{\ast}$. The costs
$\mathfrak{q}$ are measured in units $\alpha_{0,i}$ times the level-$i$ plaquette unit, and this unit
is at most $\mathbf{r}^T_{\partial}$, where $\mathbf{r}^T_{\partial}=\mathbf{r}^T_1$ denotes the unit
of the curvature entry of $T$ (and $\mathbf{r}^T_J=\mathbf{r}^T_3$ that of the current entry). The
further contributions are:
\begin{itemize}
\item the parameter $\varsigma$ costs nothing;
\item the stage layer adds $\iota^1_n(y)$, by part $(\beta)$ of step (n2);
\item the perturbation $B'$ adds $\tfrac16\mathrm{room}_T$.
\end{itemize}
The total of the proof of Proposition~\ref{4.26:prop-near-pairs-outside} and its closing inequality
therefore apply unchanged. The same holds for entry $3$, by Ba{\l}aban's parallel chain. It holds
also for entries $2$ and $4$ at the own triples $(p',X',\mathfrak{s}_T)$ of $T$, where $p'\le\ell$,
$X'\subseteq X_T$ and $\mathfrak{s}_T$ is the structure of $T$, by \cite[(3.38)]{B-CMP109} with the
constant $G'_{\square}$ of \eqref{4.26:eq-near-pair-chain-a}. What this uses is a bound
$\|g(x)\|\,\|g(x)^{-1}\|\le G'_{\square}$ on the gauge transformation
$g:=\bar{u}_{\ell}^{-1}u_{\ell}$ that conjugates in \cite[(3.38)]{B-CMP109}, taken at the full
argument of a member of the family: the conditioning bound for the conjugating gauge
transformation, whose proof is incomplete at one step. No bound on $u_{\ell}$ or $\bar{u}_{\ell}$
themselves is used.

\smallskip\noindent\textit{Step (c$'$): the remainders.}
Put $D:=d_{\ell}(X_T,\operatorname{supp}(1-\zeta'_{\square}))$, the distance at level $\ell$ used in
\eqref{4.26:eq-radii-window-a}. Then
$D\ge wL^{(i-\ell)^+}$, where $w$ is the width constant of the remainder distance (not to be
confused with the superscript $w$ of $\mathbb{U}^w_{\square}$ and $\mathbb{H}^w$). Indeed,
$X_T\subset\tilde{\square}^2$ and $\zeta'_{\square}=1$ on
$\tilde{\square}^3$, so a contour from $X_T$ to $\operatorname{supp}(1-\zeta'_{\square})$ crosses one
cube of the partition $\pi_i$. It does so in units of the block length at level $\ell$ of
$\mathfrak{B}$, which is at most the block length $\ell_{\ell}$ of level $\ell$, whatever
the frozen level. The bound on the remainder data used in Lemma~\ref{4.26:lem-near-remainders} holds
at every admissible $\mathfrak{B}$. The following depend only on $i$ and $\ell$, and therefore apply
unchanged:
\begin{itemize}
\item the remainder radii of \eqref{4.26:eq-radii-window-a};
\item the inequality \eqref{4.26:eq-near-remainders-a};
\item the Cauchy bounds;
\item the sum over $T$,
\begin{equation*}
\sum_{\ell}C_d\,(5L^{i-\ell})^d\,E_0\,e^{-\frac14\delta wL^{i-\ell}} ,
\end{equation*}
where $C_d(5L^{i-\ell})^d$ bounds the number of near terms $T$ of level $\ell$ and $E_0$ is the size
constant of the sum in Lemma~\ref{4.26:lem-near-remainders}.
\end{itemize}

\smallskip\noindent\textit{Step (n4): the statements of Ba{\l}aban used as stated.}
The list is that of the proof of Proposition~\ref{4.26:prop-near-pairs-outside}, now taken at the
scale $i$ of a cube of $\mathcal{C}^1$. It consists of the identities \cite[(3.37)--(3.39)]{B-CMP109},
\cite[(3.48)]{B-CMP109} and \cite[(3.50)]{B-CMP109}; the operator bounds behind
\cite[(3.26)--(3.27)]{B-CMP109}, \cite[(3.43)]{B-CMP109}, \cite[(3.45)]{B-CMP109} and
\cite[(3.49)]{B-CMP109}; the chain of bounds for the current in \cite{B-CMP109}; and the constant
$K_{\mathrm{cv}}$. All are taken at the scale $i$ of $\square$, with $\eta$ replaced by $L^{-i}$,
through the transfer of the analysis of \cite[\S3]{B-CMP109} to the cover, printed without proof in
\cite[p.~276]{B-CMP119}; cited as printed. No bound on the gauge transformations
$u_{\ell}$,
$\bar{u}_{\ell}$ themselves at the full argument of a member of the family is among these
statements; entries $2$ and $4$ use the conditioning bound named in step (n3), whose proof is
incomplete at one step.

Steps (n1)--(n3), together with step (c$'$), give (M) for every entry $a$; with step (H) this
completes the proof.
\end{proof}

\begin{lemma}[Further uses of the cover on lifted sets]\label{4.26:lem-cover-lifted-uses}
Work at the site $k$ and at stage $n$ of the run, with the determining-set structure $\mathfrak{B}$,
its graded refinement $\mathfrak{B}^{\mathfrak{c}}$, the stage's cover $\mathcal{C}$ (built on
$\mathfrak{B}$) and the graded cover $\mathcal{C}^1$ (built on $\mathfrak{B}^{\mathfrak{c}}$) of
Definition~\ref{4.26:def-graded-cover}. Let $\mathfrak{a}$ be a live arrest with $k_a<k$, with
piece $W_{\mathfrak{a}}\subset Z^c$, as in Definition~\ref{4.26:def-graded-cover}. Let $Y$ be a union
of cubes,
and let $(\mathbb{U},\mathbb{J})\in\mathfrak{N}^{(n)}(Y)$, the enlarged regularity space of
\eqref{4.26:eq-regularity-space-a}. For $\square\in\mathcal{C}^1$
let, as there,
$\square_0:=\tilde{\square}^5$, let $\mathfrak{s}:=\mathfrak{B}^{\mathfrak{c}}\restriction\square_0$,
let $p_{\square}$ be the top level of $\mathfrak{s}$ on $\square_0$, and
\begin{equation*}
\mathbb{U}^{\mathfrak{s}}_{\square}=\mathbb{U}^w_{\square}
:=U_{p_{\square}}\bigl(\square_0,M^{\mathfrak{s}}(\mathbb{U})\bigr).
\end{equation*}
Write $\Sigma_i$ for the stratum of $\mathfrak{B}$ of scale $i$, and call the set
$\{y:\mathfrak{c}(y)=i\}$ the $\mathfrak{c}$-shell $i$. Each statement of the expansion of the terms
of class $\mathcal{T}_1$ used below is formulated for a cube $\square\subset\Sigma_i$. It holds for
every $\square\in\mathcal{C}^1$, with $i$ the scale of $\square$, in the following form.
\begin{itemize}
\item[(Rd1)] The first localisation of the terms of $\mathcal{T}^{\mathbb{E}\mathbb{R}}$, written
with $\mathcal{C}^1$, and the real-point clause of the stage expansion hold as stated in
Lemma~\ref{4.26:lem-regrouping-exact}.
\item[(Rd2)] In the localised-average clause of the stage expansion, the configuration
$\mathbb{U}^w_{\square}:=U(\mathbb{B}^{(n+1)}_k(\square_0),M^{\cdot}(\mathbb{U}))$, extended as equal
to $\mathbb{U}$ outside $\square_0$, is read with $\mathbb{B}^{(n+1)}_k(\square_0)$ replaced by
$\mathfrak{B}^{\mathfrak{c}}(\square_0)$; here $M^{\cdot}(\mathbb{U})$ denotes the averages of
$\mathbb{U}$ for the structure in the first argument. Thus on $\square_0$
\begin{equation*}
\mathbb{U}^w_{\square}=U\bigl(\mathfrak{B}^{\mathfrak{c}}(\square_0),M^{\cdot}(\mathbb{U})\bigr)
=\mathbb{U}^{\mathfrak{s}}_{\square},
\end{equation*}
and $\mathbb{U}^w_{\square}$ is extended by $\mathbb{U}$ outside $\square_0$.
The stage background $U_{\ast}:=U^{(n+1)}_k$ is critical for $\mathfrak{B}^{\mathfrak{c}}$ at its
own averages. Moreover
$U_{\ast}\restriction\square_0=U(\mathfrak{B}^{\mathfrak{c}}(\square_0),M^{\cdot}(U_{\ast}))$
modulo $\mathcal{G}$. Finally,
$(\mathbb{U}^w_{\square},\mathbb{J})\in\mathfrak{K}_P=\mathfrak{K}_{109}(\mathfrak{B};\gamma_0)$.
\item[(Rd3)] Clauses (i)--(v) of the bounds on the small terms away from their cube
(Lemma~\ref{4.26:prf-expansion-small-terms}) hold at
$\square\in\mathcal{C}^1$, with $i$ its scale. The following hold with the same proofs: the first
display of clause (v) of Lemma~\ref{4.26:lem-family-consumption}; clause (b) of the term bound
required of the terms of $\mathcal{T}$ by the diminished term clauses $\mathfrak{I}_n$; and the
first two estimates of the
near-pair bounds. The cube $q\in\mathsf{Q}$ meeting
$\tilde{X}_T^4$ that the third estimate of the near-pair bounds requires does not occur on
$W_{\mathfrak{a}}$.
\item[(Rd4)] A term $T\in\mathcal{T}^{\mathbb{E}\mathbb{R}}$ is a \emph{tail} if $X_T$ lies in no
$\tilde{\square}^2$ with $\square\in\mathcal{C}^1$. A tail is an element of the $M$-class whose two
set components $Q$ and $\mathcal{S}$, in the notation of the definition of the elements of the
$M$-class, are empty, and whose blocks lie in $\mathcal{C}$. It satisfies
clause (a) of the term bound required of the terms of $\mathcal{T}$ by the diminished term clauses
$\mathfrak{I}_n$.
\item[(Rd5)] Let $e=(T,\square,\mathsf{y})$ be an element, or let $e=(\square,\mathsf{y})$ be a near
group, with $\square\in\mathcal{C}^1$, and let $v_e$ denote the activity of $e$. For a set $S$ let
$[S]_{\mathfrak{B}}$ be the union of the cubes of the local (frozen) scale of $\mathfrak{B}$ whose
interiors meet $S$:
\begin{equation}\label{4.26:eq-cover-lifted-uses-a}
\begin{split}
[S]_{\mathfrak{B}}:=\bigcup\bigl\{\Delta:\ &\Delta\text{ a cube of the frozen scale of }\mathfrak{B},\\
&\operatorname{int}\Delta\cap S\neq\emptyset\bigr\}.
\end{split}
\end{equation}
The domain of $e$ is
\begin{equation*}
\mathsf{d}(e):=[\tilde{X}_T^4\cup\tilde{\square}^4]_{\mathfrak{B}}\cup\mathsf{y},
\qquad\text{resp.}\qquad
\mathsf{d}(e):=[\tilde{\square}^5]_{\mathfrak{B}}\cup\mathsf{y},
\end{equation*}
with $\mathsf{y}\subset\sigma_0$, where $\sigma_0$ is the family of the cubes of the local scale of
$\mathfrak{B}$ disjoint with the interior of $X_0$, and
$X_0:=[\tilde{X}_T^4\cup\tilde{\square}^4]_{\mathfrak{B}}$, resp.
$X_0:=[\tilde{\square}^5]_{\mathfrak{B}}$, is the rounded hull. Then the following hold:
\begin{itemize}
\item $\mathsf{d}(e)$ is connected;
\item $v_e$ depends on $(\mathbb{U},\mathbb{J})$ on $\mathsf{d}(e)$ only;
\item for the localisation domain $X:=[\bigcup\mathsf{d}(o)]_m$ into which the elements $o$ are
filed, $[\,\cdot\,]_m$ being the rounding of that filing, the rounding replaces $\tilde{\square}^5$
by at most $16^d$ cubes of the frozen scale;
\item on $W_{\mathfrak{a}}$, the block set satisfies $|\mathsf{b}(e)|\le|\mathsf{y}|$.
\end{itemize}
\item[(Rd6)] The near group at $\square$ is a finite sum localised in $\tilde{\square}^4$. Its total
is at most $c_{\alpha}C^{\natural}_R(E_0+1)e^{C_2\kappa_1}\vartheta^{\natural}(\square)$,
multiplied by one plus the remainder terms. It carries no size factor: $s(\square)=0$. Here
$c_{\alpha}$, $C^{\natural}_R$, $E_0$, $C_2$ are the constants and $s(\square)$ the size factor of
the localisation statement for the near group at a cube. Moreover, at a point of the
$\mathfrak{c}$-shell $i$ only terms of level $\ell\le i$ are present, as on a shell made by the
operation $\mathbf{T}$, and the counterterm pieces and the $W$-pieces are cut by the same
$\zeta_{\square}$, as in Definition~\ref{4.26:def-graded-cover}.
\end{itemize}
\end{lemma}

\begin{proof}
We take the six clauses in order.

\emph{(Rd1).} This is Lemma~\ref{4.26:lem-regrouping-exact}.

\emph{(Rd2), the localised average.} In the localised-average clause of the stage expansion,
$\mathbb{U}^w_{\square}$ is defined as $U(\mathbb{B}^{(n+1)}_k(\square_0),M^{\cdot}(\mathbb{U}))$
and extended as equal to $\mathbb{U}$ outside $\square_0$. At $\square\in\mathcal{C}^1$ this
expression is taken with the structure $\mathfrak{B}^{\mathfrak{c}}$ in place of
$\mathbb{B}^{(n+1)}_k$. The frozen structure carries no localisation to $\square_0$, so the
localised structure on $\square_0$ is $\mathfrak{s}=\mathfrak{B}^{\mathfrak{c}}\restriction\square_0$.
The expression is therefore
$U(\mathfrak{B}^{\mathfrak{c}}(\square_0),M^{\cdot}(\mathbb{U}))=\mathbb{U}^{\mathfrak{s}}_{\square}$.

\emph{(Rd2), the real point: criticality.} The stage background $U_{\ast}=U^{(n+1)}_k$ is critical
for its own determining set. On $W_{\mathfrak{a}}$ the constraints of that set are functions of the
constraints of the refinement $\mathfrak{B}^{\mathfrak{c}}$, because the $n'$-fold averaging map
factors, $M^{n'}=M^{n'-\mathfrak{c}}\circ M^{\mathfrak{c}}$.

Consider a live piece of the branch of the classification of arrests whose band lies at
level $j+1$ with datum $V_{j+1}$, in the notation of that classification.
There the true background is critical for $\mathbb{B}^{(n^+_a-1)}_{k_a}$. This structure is also
refined by $\mathfrak{B}^{\mathfrak{c}}$. Indeed, the stage of $\mathfrak{B}^{\mathfrak{c}}$ there is
$k_0^a-h_a\le n^+_a-1$ when $j^+_a>k_0^a$, and stages only coarsen.

In both cases we have the same nesting of classes. By the statement on nested classes of
determining sets, $U_{\ast}$ is therefore critical for $\mathfrak{B}^{\mathfrak{c}}$ at its own
averages. The same statement is the one used, for $\mathfrak{B}_\ell$ (in the notation of
Lemma~\ref{4.26:prf-expansion-small-terms}), in clause (v) of
Lemma~\ref{4.26:prf-expansion-small-terms}. At a window on an arrested
component, this criticality is one of the clauses
of the criticality statement for such windows.

\emph{(Rd2), the real point: localisation.} We have
$U_{\ast}\restriction\square_0=U(\mathfrak{B}^{\mathfrak{c}}(\square_0),M^{\cdot}(U_{\ast}))$
modulo $\mathcal{G}$. This is the localisation statement for critical configurations,
$U(B',M_{B'}(U_{\ast}))=(U_{\ast}\restriction_D)^u$ for a gauge transformation $u$, in its
notation, applied with $D=\square_0$.

\emph{(Rd2), the real point: uniqueness.} The configuration critical for
$\mathfrak{B}^{\mathfrak{c}}$ at the averages of $U_{\ast}$ is unique: the real point lies in every
attempt space $\mathfrak{S}^{(k),\mathrm{arrest}}_l$ with margin $\tfrac12$, by the statement on the
margin of the real point in the attempt spaces.

Hence the first localisation and the real-point clause are untouched.

\emph{(Rd2), the kernel.} The triple $(p_{\square},\square_0,\mathfrak{s})$ is a triple of the
derived-entry family on a live piece, \eqref{4.26:eq-graded-cover-a}. The size $\gamma_0$ is
$\gamma_0^{(k)}$ of \eqref{4.26:eq-regularity-radius-a}.

Away from the two outer layers of $\square_0$, the direct entries of $\mathbb{U}^w_{\square}$ are
the composite entries of
$\mathbb{U}$ at $(p_{\square},\square_0,\mathfrak{s})$. In the frozen unit these are smaller than
$1.4\,r_e^{(p_{\square})}(y')$ at every point $y'$ of $\square_0$ away from its two outer layers,
by the composite shell clause of
the enlarged regularity space,
\eqref{4.26:eq-regularity-space-a}. They are thus small at the levels of $\mathfrak{B}$ itself. This
is all that the definition of $\mathfrak{K}_{109}(\mathfrak{B};\gamma_0)$ requires, with
$\gamma_0$ as above.

On the two outer layers we use the following sentence of \cite{B-CMP109}:
\begin{quote}
$U_{k+1}(\square_0,M^{\cdot}(\mathbf{U}))$ has approximately the same regularity properties as
$\mathbf{U}$, hence all necessary theorems are valid for operations with this configuration
\end{quote}
It is used here as stated in
\cite{B-CMP109}, at the level at which the localised-average clause is applied. This regularity of
the outer layers of a localised configuration is proved, level $n$ for level $n$, by the
outer-layer regularity statement for localised configurations.

At a point $y'$ of $\square_0$, the plaquette entries (entry $2$, which involves no block side) are
of size
$1.4\,\omega(y')\alpha_{0,m(y')}$, and
\begin{equation*}
1.4\,\omega(y')\alpha_{0,m(y')}\le 1.4\,w_{\ast}(k_a)\bar{\alpha}_{k_a}
\le\tfrac{1.4}{5}\gamma_0^{(k)}
\end{equation*}
at the levels of $\mathfrak{B}$. This is all that $\mathfrak{K}_{109}$ requires: the regularity of
$\mathbb{U}$ at the levels of $\mathfrak{B}$, and the bound on the current in the definition of
$\mathfrak{K}_{109}(\mathfrak{B};\gamma_0)$.

The block current is not a clause of $\mathfrak{K}_{109}$. Suppose a smallness condition on it is
read at the block side of $\mathfrak{B}$. Its size there at a point $y'$ of $\square_0$, with
$\omega=\omega(y')$, $m=m(y')$ and $\mathfrak{c}=\mathfrak{c}(y')$, is
\begin{equation*}
1.4\,\omega\alpha L^{(p_{\square}\wedge m)-(p_{\square}\wedge\mathfrak{c})}
\le 1.4\,L\omega\alpha\qquad\text{at }p_{\square}=\mathfrak{c}(y')+1 .
\end{equation*}
As on the two outer layers, it is controlled by the same sentence of \cite{B-CMP109}, starting from
the direct entries.

\emph{(Rd3).} The proof of Lemma~\ref{4.26:prf-expansion-small-terms} depends on $\square$ at four
places. In what follows $\delta\mathsf{B}$, $\hat{p}_{\square}$, $\mathfrak{B}_\ell$,
$\ell_{\mathfrak{B}_\ell}$, $\ell_\ell$, $s_i$, $R_n$, $w$, $\Lambda_\ell$, $\mathsf{d}_\ell(X_T)$ and
$Q(\mathbb{U};\eta A)(c)$ are as in the notation of Lemma~\ref{4.26:prf-expansion-small-terms}.
\begin{itemize}
\item[(ii)] The variation $\delta\mathsf{B}$ lives on bonds whose blocks meet
$\operatorname{supp}\zeta_{\square}$, and $\|\delta\mathsf{B}\|_{\hat{p}_{\square}}\le4B_{\sharp}$.
This uses only $0\le\zeta_{\square}\le1$ and the support of $\zeta_{\square}$, through the bound
$|Q(\mathbb{U};\eta A)(c)|\le 2\ell_{\mathfrak{B}_\ell}\sup|A|$ on the averaging operator
$Q(\mathbb{U};\eta A)$. No derivative of $\zeta_{\square}$
enters.
\item[(iii)] The third display is
$d^{\wedge}(\square^{\zeta},\mathcal{C}_j)\ge d^{\wedge}(\square,R_n)$. It follows from the
definition of $d^-_{\square}$, which applies to any set $\square$.
\item[(iv)] In the summation, two facts are used. First, a pair with
$X_T\cap\tilde{\square}\neq\emptyset$ has $\operatorname{diam}X_T\ge 2Ms_i$. Second, if
$X_T\cap\tilde{\square}=\emptyset$, every contour from $X_T$ to $\operatorname{supp}\zeta_{\square}$
crosses a $\pi_i$-cube; measuring lengths at level $\ell$ in units at most $\ell_\ell$, this gives
\begin{equation*}
\mathsf{d}_\ell(X_T)\ge wL^{(i-\ell)^+},\qquad\text{and}\qquad
\mathsf{d}_\ell(X_T)\ge w(\ell-i-1)^+\ \text{ for }\ell>i .
\end{equation*}
Here $i$ is the scale of the cube, and $\ell_{\mathfrak{B}_\ell}\le\ell_\ell$. For $\ell>i$ we have
$X_T\subset\Lambda_\ell$, and $\square$ lies in the $\mathfrak{c}$-shell $i$. The contour therefore
crosses the $\mathfrak{c}$-shells $i+1,\dots,\ell-1$. By \eqref{4.26:eq-lifted-counts-a} of
Lemma~\ref{4.26:lem-lifted-counts}, each of these shells is at least one $M$-cube of the local
scale of $\mathfrak{B}$ wide.
\end{itemize}
Clause (i), the step $(\star)$ and clause (v) do not involve $\square$. Hence clauses (i)--(v) hold
at $\square\in\mathcal{C}^1$, with $i$ its scale.

Three consequences follow with the same proofs: the first display of clause (v) of
Lemma~\ref{4.26:lem-family-consumption}; clause (b) of the term bound required of the terms of
$\mathcal{T}$ by the diminished term clauses $\mathfrak{I}_n$; and the first two estimates of the
near-pair bounds. The third estimate of
the near-pair bounds requires a cube $q\in\mathsf{Q}$
meeting $\tilde{X}_T^4$. No such cube meets $W_{\mathfrak{a}}$,
since $\mathsf{Q}\subset R^{(l)}\subset Z$ and $W_{\mathfrak{a}}\subset Z^c$; so this cube does not
occur on $W_{\mathfrak{a}}$.

\emph{(Rd4).} With the cover $\mathcal{C}^1$, a lifted set carries more tails than with
$\mathcal{C}$: the threshold on the diameter becomes $\operatorname{diam}X_T\ge 2M\ell_{\mathfrak{c}}$
in place of $2M\ell_{\ell_{\mathfrak{a}}}$.

A tail is an element of the $M$-class whose set components $Q$ and $\mathcal{S}$ are empty,
and whose blocks lie in $\mathcal{C}$. For a tail, clause (a) of the term bound required of the
terms of $\mathcal{T}$ by the diminished term clauses $\mathfrak{I}_n$ holds, since that clause
holds for every
$T\in\mathcal{T}$, because the diminished term clauses $\mathfrak{I}_n$ quantify over
$\mathcal{T}$ by the kind of term ($\mathbb{E}$, $\mathbb{R}$, $\mathbb{B}$) and not by class. The
sums over tails are estimated with the frozen-level
sums of the expansion.

\emph{(Rd5), the domain.} The family $\sigma_0$ of the statement is the family of
\cite{B-CMP116} and \cite{B-CMP119}. In \cite{B-CMP116}, $\sigma_0$ is
\begin{quote}
the family of all cubes $\Delta$ disjoint with the interior of $X_0$, where $X_0$ is the smallest
localization domain from $\mathbf{D}_k$ containing $X$
\end{quote}

The set $\mathsf{d}(e)$ is connected, since $\mathsf{y}$ is a connected family of cubes of
$\sigma_0$ adjoining the rounded hull $X_0$.

\emph{(Rd5), the support of $v_e$.} The activity $v_e$ depends on $(\mathbb{U},\mathbb{J})$ on
$\mathsf{d}(e)$ only. The rounded hull cannot be shrunk, since the layer $\mathbb{H}^w(\sigma)$
depends on $\mathbb{U}$ on all of it.
Where the lift $g\ge1$, the enlargement $\tilde{\square}^5$ is not a union of cubes of the frozen
scale. The un-rounded set
$\tilde{\square}^5\cup\mathsf{y}$ is then in general neither connected nor the support of $v_e$;
for this reason the rounding $[\,\cdot\,]_{\mathfrak{B}}$ is used.

\emph{(Rd5), the count.} The localisation domain $X:=[\bigcup\mathsf{d}(o)]_m$ rounds
$\tilde{\square}^5$ up to cubes of the frozen scale. The cube $\tilde{\square}^5$ has side twelve
$\pi_i$-sides, so it meets at most $13^d$ cubes of $\pi_i$ and, since $12/L\le1$, at most $3^d$
cubes of a coarser partition $\pi_{n'}$. In either case at most $16^d$ cubes of the frozen scale
arise. This count enters the hull constant
$e^{(a_k+1+c_J)16^d}$, which is absorbed into the constant $K_{\flat}'$. The same hull is the set
that the $\sigma$-expansion of \cite{B-CMP116} never decouples.

\emph{(Rd5), the block set.} On $W_{\mathfrak{a}}$ we have
$\mathsf{b}(e)=\{q\in\mathsf{Q}:q\text{ meets }\mathsf{y}\}$, so $|\mathsf{b}(e)|\le|\mathsf{y}|$.

\emph{(Rd6).} We apply the conclusion of the localisation statement for the near group at a cube,
without reproducing its proof, at a cube of $\mathcal{C}^1$ of scale $i$. Its three hypotheses,
read at scale $i$ and with $U^c_\ell$ and $\alpha_2$ as in the notation of that statement, are the
following.
\begin{itemize}
\item $T$ is analytic and $G^c$-gauge invariant on $U^c_\ell$, with room of the type $1+2\beta$.
\item The perturbation is smaller than $\alpha_2$ on $\square_0$, in scale-$i$ units.
\item The local minimisers on which that statement rests are taken on the cube sequence built
from $\pi_i$, and their composite entries are the entries of the derived-entry family on a live
piece, \eqref{4.26:eq-graded-cover-a}.
\end{itemize}
These are supplied at scale $i$, respectively, by the analyticity hypothesis of the expansion, by
its perturbation-size hypothesis, and by its local-minimiser hypothesis.

The conclusion is that the near group at $\square$ is a finite sum localised in
$\tilde{\square}^4$. Its total is at most
$c_{\alpha}C^{\natural}_R(E_0+1)e^{C_2\kappa_1}\vartheta^{\natural}(\square)$, multiplied by one
plus the remainder terms of that statement. There is no size factor, $s(\square)=0$: no weight
survives in the box's unit.

At a point of the $\mathfrak{c}$-shell $i$, the levels $\ell$ of the terms present satisfy
$\ell\le i$, since
$X_T\subset\Lambda_\ell$; this is as on a shell made by the operation $\mathbf{T}$. The counterterm
pieces and the $W$-pieces are cut by the same $\zeta_{\square}$, as in
Definition~\ref{4.26:def-graded-cover}.
\end{proof}

\begin{theorem}[Near pairs meeting the lifted points of a live arrest]\label{4.26:thm-near-pair-lifted}
Let $k$ be the current level and $l$ a stage, and take as the cover used in the estimates of the
small terms of types $\mathbb{E}$ and $\mathbb{R}$ the graded cover $\mathcal{C}^1$ of
Definition~\ref{4.26:def-graded-cover}. Let $\mathfrak{a}$ be a
live arrest with piece $W_{\mathfrak{a}}$ and level $k_a$, where
\begin{equation*}
k_a<k, \qquad W_{\mathfrak{a}}\subset Z^c ,
\end{equation*}
and let $\mathsf{L}_{\mathfrak{a}}\subset W_{\mathfrak{a}}$ be its lifted set. Let $Y$ be a union
of cubes and $(\mathbb{U},\mathbb{J})\in\mathfrak{D}^{\flat}_l(Y)$. Let $(T,\square)$,
$\square\in\mathcal{C}^1$, be a near pair of a clause of $T$ at a point of $\mathfrak{T}$, where
$T\in\mathcal{T}_1\cap\mathcal{T}^{\mathbb{E}\mathbb{R}}$ has level $\ell$, $X_T\subset
\tilde{\square}^2$ and $[\square_0]_{\mathfrak{B}}\subset Y$, and with
\begin{equation*}
\square\cap\mathsf{L}_{\mathfrak{a}}\neq\emptyset .
\end{equation*}
Assume for these data the hypotheses of Proposition~\ref{4.26:prop-near-pair-graded} at stage $l$,
and the conditions on the constants of Lemma~\ref{4.26:lem-graded-box-geometry} and of
Lemma~\ref{4.26:lem-lifted-counts}.
Then:
\begin{itemize}
\item[(a)] clause (c) of the stability hypothesis $(\mathrm{H}^{\natural})_l$, display
\eqref{4.26:eq-stage-systems-b}, holds at $Y$, at $(\mathbb{U},\mathbb{J})$ and at the pair
$(T,\square)$: the input hypotheses that this clause
imposes on $\square_0$ are satisfied, and every member of family (c), together with both
configurations (c$'$), lies in $\mathfrak{U}_T$;
\item[(b)] the sums over all such pairs $(T,\square)$ that are taken in deducing the bounds of the
step from clause (c) of \eqref{4.26:eq-stage-systems-b} satisfy the bounds of clauses (F2) and (F3)
of Lemma~\ref{4.26:lem-lifted-counts} and of clauses (fr-a), (fr-$\mu$) and (fr-$\mathsf{M}$) of
Corollary~\ref{4.26:cor-summation-lifted}.
\end{itemize}
\end{theorem}

Such pairs occur because a large-field region $X$ at which the step is estimated may reach
$W_{\mathfrak{a}}$, as \cite[(1.63)]{B-CMP122a} permits.

\begin{proof}
We put together the statements of this section on the graded cover, in the following order.

\emph{The cover.} By Lemma~\ref{4.26:lem-regrouping-exact}, the first localisation of the small
terms is linear in the partition of unity, with one global interpolation parameter. Consequently,
for each term $T$, the sum of the pieces of $T$ localised to the cubes $\square\in\mathcal{C}^1$
equals the sum of its pieces localised to the cubes $\square\in\mathcal{C}$. Hence every total
$\mathbb{C}^{(l)}_k$ is the same for the two covers, and only the grouping of the pieces of the
small terms on the lifted sets changes. It is therefore legitimate to read clause (c) of
\eqref{4.26:eq-stage-systems-b}, and the sums in (b), with the cover $\mathcal{C}^1$, and this is
how the theorem is stated.

On $W_{\mathfrak{a}}$, where $k_a<k$ and $W_{\mathfrak{a}}\subset Z^c$,
Definition~\ref{4.26:def-graded-cover} (display \eqref{4.26:eq-graded-cover-a}) replaces
$\mathfrak{B}\restriction W_{\mathfrak{a}}$ by the determining set of the strata of the
original-level function $\mathfrak{c}$. By Lemma~\ref{4.26:lem-original-level-strata}(i), these
strata are the strata of a determining set of Ba{\l}aban's construction on $W_{\mathfrak{a}}$. The
cover $\mathcal{C}^1$ is therefore the ordinary cover built on $\mathfrak{c}$ in place of the
frozen level.

\emph{Geometry of the box.} By Lemma~\ref{4.26:lem-original-level-strata}(iii), every
$\mathfrak{c}$-shell $i\ge h_a$ of $W_{\mathfrak{a}}$ is at least three layers of
$MR_{i+1}$-cubes of the $(i+1)$-lattice wide (eight for $i\ge k_0$). Since $\square\subset\square_0$
and $\mathsf{L}_{\mathfrak{a}}\subset W_{\mathfrak{a}}$, the
box $\square_0$ meets the piece $W_{\mathfrak{a}}$, which is older than the current level because
$k_a<k$. Lemma~\ref{4.26:lem-graded-box-geometry} therefore applies to $\square$: its conditions on
the constants are assumed in the theorem, and $[\square_0]_{\mathfrak{B}}\subset Y$ holds by
hypothesis.

Let $i$ be the scale of $\square$. Clause (a) of that lemma gives three facts: $\square_0$ meets at
most two $\mathfrak{c}$-shells, these are adjacent and $i$ is among them; the level $\ell$ of the
near pair satisfies $\ell\le i$; and $\ell\le\mathfrak{c}$ on $\square_0$. Its clause (c$_0$)
converts the weighted radii at every point of $\square_0\cap W_{\mathfrak{a}}$ into the unit of the
$\mathfrak{c}$-shell.
These are the geometric inputs of Proposition~\ref{4.26:prop-near-pair-graded} at cubes of
$\mathcal{C}^1$: the shells met by $\square_0$ are thick, so that the two-shell argument for an
ordinary shell applies without change, and by (c$_0$) a weighted clause frozen at a level
$n\ge\mathfrak{c}$ implies the unweighted clause in the unit of the $\mathfrak{c}$-shell.

\emph{Clause (c) of $(\mathrm{H}^{\natural})_l$.} The pair $(T,\square)$ is a near pair on the
graded cover at a clause of $T$ at a point of $\mathfrak{T}$. We apply
Proposition~\ref{4.26:prop-near-pair-graded} to it at stage $l$; its hypotheses are those assumed
in the theorem. Its conclusion (H) states that the input hypotheses (i)--(iv) imposed by clause (c)
of \eqref{4.26:eq-stage-systems-b} hold on $\square_0\cap\mathfrak{T}$ with factor $<1+2\beta$, in
the unit of the $\mathfrak{c}$-shell.
Its conclusion (M) states that every member $P(\varsigma;t,\tau;B')$ of family (c), and both
configurations (c$'$), satisfy, in the notation of Proposition~\ref{4.26:prop-near-pair-graded},
$Q_a(\cdot)(y)<\theta_{T,a}(y)$ at the clauses
$(y,a)\in\operatorname{Cl}(T)$; that is, they lie in $\mathfrak{U}_T$. These are the two parts of
clause (c) of \eqref{4.26:eq-stage-systems-b}, which proves (a).

\emph{The other uses of the cover.} Lemma~\ref{4.26:lem-cover-lifted-uses}, display
\eqref{4.26:eq-cover-lifted-uses-a}, shows that every other place in the estimates of the small
terms of types $\mathbb{E}$ and $\mathbb{R}$ where the cover is used remains valid on the lifted
sets with $\mathcal{C}^1$ in place of $\mathcal{C}$. So no other step of those estimates is
affected by the passage to $\mathcal{C}^1$.

\emph{The sums.} The cost of the graded cover is a count: one frozen cube may contain many finer
cubes of $\mathcal{C}^1$. Lemma~\ref{4.26:lem-lifted-counts}, whose conditions on the constants are
assumed in the theorem, bounds this count at every
$\square\in\mathcal{C}\cup\mathcal{C}^1$ meeting $\mathsf{L}_{\mathfrak{a}}$ by the decay across
the lifted shells, through the lower bound on $D_F(\square)$ in its clause (F1) and the inequalities
(F2) and (F3). Clause (fr-a) of Corollary~\ref{4.26:cor-summation-lifted} bounds the leaf sum at
every cube of $\mathcal{C}$ meeting the lifted set of the live older piece $W_{\mathfrak{a}}$, and
states that the other summation bound at such a cube holds unchanged. The same corollary also
bounds the $T$-sums and the totals of the
near group at the cube $\square\in\mathcal{C}^1$ itself, which lies on a lifted set because
$\square\cap\mathsf{L}_{\mathfrak{a}}\neq\emptyset$ by hypothesis; its clause (fr-$\mu$) shows that
the sum over boxes weighted by the smallness parameter $\vartheta^{\natural}$ is unmoved; and its
clause (fr-$\mathsf{M}$) bounds the anchor sums. These are the sums taken over the pairs of the
theorem, which proves (b).

\emph{Cases not covered above.} A box meeting a healed region, rather than the lifted set of a live
arrest, is not within the hypotheses of the theorem. Pieces with
$W_{\mathfrak{a}}\subset Z$ (core pieces) are excluded by the hypothesis
$W_{\mathfrak{a}}\subset Z^c$, and are the subject of the counterpart of
Lemma~\ref{4.26:lem-lifted-counts} for pieces contained in $Z$.
\end{proof}

\begin{lemma}[Room over the source and the schedule for newborn terms]\label{4.26:lem-room-table}
Let $k$ be the level, and let $\mathcal{T}$, the clause sets $\operatorname{Cl}(T)$, the entries
$Q_a$, the thresholds $\theta_{T,a}$ and the radii $\mathbf{r}^T_a$ be as in
Definition~\ref{4.26:def-clauses-units}. Let $r_a(y)$ be the radius of entry $a$ at $y$, and let
$\theta_{\mathrm{src},a}$ be the threshold of entry $a$ in the source space at level $k$. For
$y\in\mathfrak{T}$ let $m=m(y)$ be the layer of the layered spaces of
Definition~\ref{4.4:def-post-step-space} containing $y$, the top layer being $\Omega_k$. For
$j\ge m$ put
\begin{equation}\label{4.26:eq-room-table-1}
s^{(j)}_m:=1-\beta\bigl(1-2^{-(j-m)}\bigr),\qquad\text{so that } s^{(m)}_m=1 .
\end{equation}
For $T\in\mathcal{T}$ let $\ell_T$ be the level at which $T$ is created (the superscript $j$ of
$\mathbb{B}^{(j)}$, $\mathbb{E}^{(j)}$, $\mathbb{R}^{(j)}$). For $(y,a)\in\operatorname{Cl}(T)$ with
$y\in\mathfrak{T}$ define the \emph{room}
\begin{equation}\label{4.26:eq-room-table-2}
r_{T,a}(y):=\frac{\theta_{T,a}(y)-\theta_{\mathrm{src},a}(y)}{r_a(y)} .
\end{equation}
The clauses listed in (a)--(c) below are those with $a\notin\{5,8\}$.
Let $\theta_S$ be the number fixed once, in the order of choice of the constants, in the interval
$[\tfrac12,\tfrac89]$, and write $\mathsf{c}_e:=\theta_S$. Then the following hold, where
$(\mathrm{s})$ lists the factors $F_m$ of the spaces of the terms created at
level $k$ after the operation $\mathbf{T}$, and $r_{\min}(y)$ is the minimum of $r_{T,a}(y)$ over
the clauses listed in (a) and (c) with $a\notin\{5,8\}$; by definition, the clauses of (b), for
which only lower bounds on $\theta_{T,a}$ are listed, are not included in $r_{\min}$:
\begin{equation}\label{4.26:eq-room-table-a}
\begin{aligned}
&\text{(a)}\ \ r_{T,a}(y)=s^{(j)}_m-s^{(k)}_m\ge\beta2^{-(k-m)}
  &&T=\mathbb{B}^{(j)},\ m\le j\le k-1,\\
&\phantom{\text{(a)}}\ \ r_{T,a}(y)=1-s^{(k)}_m=\beta\bigl(1-2^{-(k-m)}\bigr)\ge\beta/2
  &&T\in\{\mathbb{E}^{(m)},\mathbb{R}^{(m)}\},\ m<k,\\
&\text{(b)}\ \ \theta_{T,a}(y)\ge2r_a(y),\quad \theta_{T,a}(y)\ge(1-\beta)\mathbf{r}^T_a
  &&\ell_T<m,\\
&(\mathrm{s})\ \ F_k=1+\beta,\quad F_m=s^{(k)}_m+\theta_S\beta2^{-(k-m)}
  &&m<k,\\
&\text{(c)}\ \ r_{T,a}(y)=\beta
  &&T\in\{\mathbb{E}^{(k)},\mathbb{R}^{(k)}\},\ m=k,\\
&\phantom{\text{(c)}}\ \ r_{T,a}(y)=\beta e_\ast(m)
  &&T=\mathbb{B}^{(k)},\\
&\phantom{\text{(c)}}\ \ r_{\min}(y)\ge\beta e_\ast(m),\quad e_\ast(k)=1,\ \
  e_\ast(m)=\theta_S2^{-(k-m)}\ (m<k).
\end{aligned}
\end{equation}
For $a\in\{5,8\}$ the room satisfies only $r_{T,a}(y)\ge0$, and these two entries are never
incremented. Every clause of a term $\mathbb{E}^{(\ell)}$ or $\mathbb{R}^{(\ell)}$ with
$\ell_T=\ell=m$ has room at least $\beta/2$, whatever the value of $\theta_S$.
\end{lemma}

\begin{proof}
\emph{Part (a), terms $\mathbb{B}^{(j)}$.} Let $T=\mathbb{B}^{(j)}$ with $m\le j\le k-1$. By the
statement of this chapter that tabulates the thresholds of the inherited terms on the layers, the
room~\eqref{4.26:eq-room-table-2} of a clause of $T$ on layer $m$ is the difference
$s^{(j)}_m-s^{(k)}_m$ of the factors of~\eqref{4.26:eq-room-table-1}. By~\eqref{4.26:eq-room-table-1},
\begin{equation*}
s^{(j)}_m-s^{(k)}_m=\beta\bigl(1-2^{-(k-m)}\bigr)-\beta\bigl(1-2^{-(j-m)}\bigr)
=\beta\bigl(2^{-(j-m)}-2^{-(k-m)}\bigr).
\end{equation*}
Since $j\le k-1$ we have $2^{-(j-m)}\ge2^{-(k-1-m)}=2\cdot2^{-(k-m)}$, so the right side is at
least $\beta2^{-(k-m)}$.

\emph{Part (a), terms $\mathbb{E}^{(m)}$, $\mathbb{R}^{(m)}$ with $m<k$.} By the same statement,
the room of a clause of such a term on layer $m$ is $1-s^{(k)}_m$, and
$1-s^{(k)}_m=\beta(1-2^{-(k-m)})$ by~\eqref{4.26:eq-room-table-1}. Since $k-m\ge1$, we have
$2^{-(k-m)}\le\tfrac12$, whence $1-s^{(k)}_m\ge\beta/2$.

\emph{Part (b).} For a term with $\ell_T<m$ (a drop of scale between the level of the term and
the layer of $y$), the two inequalities $\theta_{T,a}\ge2r_a$ and
$\theta_{T,a}\ge(1-\beta)\mathbf{r}^T_a$ are clause (c) of the theorem in which the
thresholds $\theta_{T,a}$ of the terms are fixed; they are used here as stated there.

\emph{Part (c) and the schedule $(\mathrm{s})$.} Ba{\l}aban~\cite{B-CMP119} states that the terms
created at level $k$ are ``defined on slightly larger spaces, with the coefficients in their
definition bigger by $\beta$ multiplied by a corresponding number''. For this number we use the
schedule proved in the schedule lemma for the layered spaces after $\mathbf{T}$: the factor is
$1+\beta$ on the top layer $\Omega_k$ and $F_m=s^{(k)}_m+\theta_S\beta2^{-(k-m)}$ on every layer
$m<k$, with a fixed $\theta_S\in\,]0,1[$. This is the line $(\mathrm{s})$
of~\eqref{4.26:eq-room-table-a}.

The factor on the layers $m<k$ is not the uniform factor $1+\beta$. Indeed,
by~\eqref{4.26:eq-room-table-1}, with $x=2^{-(k+1-n)}$,
\begin{equation*}
(1+\beta)s^{(k+1)}_n-s^{(k)}_n
=\beta\bigl[1-(1+\beta)(1-x)+1-2x\bigr]
=\beta(1-\beta)\bigl(1-2^{-(k+1-n)}\bigr),
\end{equation*}
which is positive for $n\le k$ because $0<\beta<1$
(Definition~\ref{4.4:def-post-step-space}) and $2^{-(k+1-n)}\le\tfrac12$; this positivity is
claim (i) of the schedule lemma for the layered spaces after $\mathbf{T}$, by which a uniform
factor $1+\beta$ on all layers is incompatible with the level-$k$ space. Moreover,
by~\eqref{4.26:eq-room-table-1},
\begin{equation*}
s^{(k-1)}_m-s^{(k)}_m=\beta\bigl(2^{-(k-1-m)}-2^{-(k-m)}\bigr)=\beta2^{-(k-m)}
\end{equation*}
exactly. Hence the value $\theta_S=1$ would give $F_m=s^{(k-1)}_m$ on every layer $m<k$, and
would leave zero room for the response of the step through $\mathbf{T}$ itself; a schedule with
$\theta_S=1$ therefore cannot be reached by composing the steps, and $\theta_S<1$ is kept fixed.
The part of $\mathbb{B}^{(k)}$ depending on the variable $A_k$ of Ba{\l}aban's boundary terms at
level $k$ is defined on the space obtained from the first factor of Ba{\l}aban's space for these
$A_k$-dependent boundary terms, read layer by layer with the factors $F_m$ of the
schedule lemma for the layered spaces after $\mathbf{T}$.

As in part (a), by the statement of this chapter that tabulates the thresholds of the inherited
terms, a term whose space carries the factor $F$ on layer $m$ has $\theta_{T,a}(y)=F\,r_a(y)$
there, while $\theta_{\mathrm{src},a}(y)=s^{(k)}_m\,r_a(y)$; hence
by~\eqref{4.26:eq-room-table-2} its room is $F-s^{(k)}_m$.
We now compute the rooms of the newborn terms. For $\mathbb{E}^{(k)}$ and $\mathbb{R}^{(k)}$ on
the top layer $m=k$, the factor of their space is the factor $1+\beta$ of
\cite[(3.16)]{B-CMP109} and of the statement of this chapter that tabulates the thresholds of the
inherited terms, and the source factor is $s^{(k)}_k=1$, so the room is
$(1+\beta)-1=\beta$. For $\mathbb{B}^{(k)}$ the room is, on the top layer and on a layer $m<k$
respectively,
\begin{equation*}
F_k-s^{(k)}_k=(1+\beta)-1=\beta e_\ast(k),\qquad
F_m-s^{(k)}_m=\theta_S\beta2^{-(k-m)}=\beta e_\ast(m).
\end{equation*}
This proves the first two lines of (c), with $\mathsf{c}_e=\theta_S$ the newborn number.

\emph{The entries $a\in\{5,8\}$.} For these two entries the only room asserted is
$r_{T,a}(y)\ge0$, and they are never incremented. They are not counted in $r_{\min}$, which is
taken over entries $a\notin\{5,8\}$.

\emph{The bound on $r_{\min}$.} Fix $y$ with layer $m$. If $m=k$, the terms $\mathbb{B}^{(j)}$
with $m\le j\le k-1$ and $\mathbb{E}^{(m)},\mathbb{R}^{(m)}$ with $m<k$ do not occur, and the
clauses listed in (c) all have room $\beta=\beta e_\ast(k)$. If $m<k$, then by (a) a term
$\mathbb{B}^{(j)}$ has room at least $\beta2^{-(k-m)}\ge\beta\theta_S2^{-(k-m)}$ because
$\theta_S<1$; a term $\mathbb{E}^{(m)}$ or $\mathbb{R}^{(m)}$ has room at least
$\beta/2\ge\beta\theta_S2^{-(k-m)}$ because $2^{-(k-m)}\le\tfrac12$ and $\theta_S<1$; and by (c)
the term $\mathbb{B}^{(k)}$ has room exactly $\beta\theta_S2^{-(k-m)}$. In every case the room is
at least $\beta e_\ast(m)$, which is the last line of~\eqref{4.26:eq-room-table-a}. The clauses of
(b) are not counted in $r_{\min}$: for them the table lists only the two lower bounds on
$\theta_{T,a}$, not a room.

\emph{Clauses of $\mathbb{E}$- and $\mathbb{R}$-terms with $\ell_T=m$.} If $m<k$ the room is
$\beta(1-2^{-(k-m)})\ge\beta/2$ by (a); if $m=k$ it is $\beta\ge\beta/2$ by (c). Neither bound
involves $\theta_S$.
\end{proof}

\begin{remark}[The constant $\theta_S$ and the scope of the room table]
The endpoints $\tfrac12$ and $\tfrac89$ of the interval for $\theta_S$ are explained in the remark
on the endpoints of the window for $\theta_S$. Since $\theta_S\le\tfrac89<1$ and
$\theta_S\ge\tfrac12>0$, the quantities $1/(1-\theta_S)$, $1/((1-\theta_S)\beta)$ and
$\tfrac14\min\{\theta_S,1-\theta_S\}$, on which constants chosen later in the order depend, are
finite and positive.

The table of Lemma~\ref{4.26:lem-room-table} concerns the terms of \cite{B-CMP119}. The arrested
left-overs of the large-field operation of \cite{B-CMP122a} enter the diminished term clauses
$\mathfrak{I}_n$ of Definition~\ref{4.26:def-inherited-domain} like any clause of a term of
$\mathcal{T}^{\mathbb{B}}$, so part (a) of the theorem on the diminished term clauses applies to
them. Their rooms, however, are not those of~\eqref{4.26:eq-room-table-a} but those given for
arrested regions in the theorem in which the thresholds $\theta_{T,a}$ are fixed. The theorem on
arrested regions supplies these rooms. On the piece of an arrest it does so by its part
concerning the piece, together with the bound on the sum of the stage increments; off the piece
it does so by its part concerning the complement of the piece.
\end{remark}

\begin{corollary}[Positive room of the inherited domain over the source]\label{4.26:cor-inherited-room}
Assume the hypotheses of Theorem~\ref{4.26:thm-inherited-expansion}, and let $\theta_{\mathrm{src},a}$
be the source thresholds entering Definition~\ref{4.26:def-inherited-domain}. For $y\in\mathfrak{T}$
write $m=m(y)$ for the level of Definition~\ref{4.26:def-regularity-space}, and let
$r_{\min}(y)$ and $e_\ast$ be as in Lemma~\ref{4.26:lem-room-table} and
\eqref{4.26:eq-room-table-a}.

Let $n$ be a stage, $0\le n\le N$, with $N$ the number of stages of
Theorem~\ref{4.26:thm-inherited-expansion}, let $(X,\mathcal{S})$ be as in the
representation in that theorem, let $T\in\mathcal{T}$ with $X_T\subset X$, and let
$(y,a)\in\operatorname{Cl}(T)$ with $y\in X\cap\mathfrak{T}$ and $a\notin\{5,8\}$; write
$r_T=r_{T,a}(y)$ as in Lemma~\ref{4.26:lem-room-table}. Consider the threshold
$\theta^{(n)}_{T,a}(y)$ of $\mathfrak{I}_n$, both for $T$'s own clauses in $\mathfrak{I}_n$ and for
the box clause \eqref{4.26:eq-inherited-domain-a}. If $T\in\mathcal{T}^{\mathbb{B}}$ with $\ell_T\ge m$,
or $T\in\mathcal{T}^{\mathbb{E}\mathbb{R}}$ with $\ell_T=m$, the excess
$\theta^{(n)}_{T,a}(y)-\theta_{\mathrm{src},a}(y)$, in units $r_a(y)$, is the first expression of the
first, respectively second, line below and is bounded below as shown there; in the scale-drop case
$\ell_T<m$ of Lemma~\ref{4.26:lem-room-table} the threshold itself obeys the third line, and at a point
$y$ of the lifted set $\mathsf{L}_{\mathfrak{a}}$ of an older live piece $\mathfrak{a}$ (item~(5) of
\eqref{4.26:eq-older-pieces-domain-a}) the fourth line takes the place of the third:
\begin{equation}\label{4.26:eq-inherited-room-a}
\begin{aligned}
&r_T-\textstyle\sum_{l\le n}\iota_l
  \ge r_T-\tfrac12\mathsf{w}_0\beta e^{-n_W(y)}\ge\tfrac12 r_T
  &&(T\in\mathcal{T}^{\mathbb{B}},\ \ell_T\ge m),\\
&\tfrac23 r_T-\textstyle\sum_{l\le n}\iota^1_l
  \ge\tfrac23 r_T-2\mathsf{w}_0\beta\ge\tfrac12 r_T
  &&(T\in\mathcal{T}^{\mathbb{E}\mathbb{R}},\ \ell_T=m),\\
&\theta^{(n)}_{T,a}(y)\ge 1.5\,r_a(y)
  &&(\ell_T<m,\ y\ \text{in no lifted set}),\\
&\theta^{(n)}_{T,a}(y)\ge 1.35\,r_a(y)
  &&(\ell_T<m,\ y\in\mathsf{L}_{\mathfrak{a}}\ \text{for an older live piece}\ \mathfrak{a}),
\end{aligned}
\end{equation}
where $\iota_l,\iota^1_l$ are the stage increments of \eqref{4.26:eq-family-consumption-a}. Moreover,
at the same locations the space $\mathfrak{N}^{(n)}$ of \eqref{4.26:eq-regularity-space-a} exceeds the
source threshold, in the same units, by at least $2\beta(1-\mathsf{w}_0)\ge\beta$ on the direct entries,
and bounds the composite entries at every $(\Sigma^{\ast})$-triple $X'\subset X$ by
\begin{equation*}
(1+2\beta-2\mathsf{w}_0\beta)\,r_e^{(p)}=1.383\,r_e^{(p)}\ge(1+\beta)\,r_e^{(p)}+0.18\,r_e^{(p)}.
\end{equation*}
Hence the inherited domain $\mathfrak{D}^{\flat}_n(X)$ of Definition~\ref{4.26:def-inherited-domain}
and \eqref{4.26:eq-inherited-domain-b} contains every pair that
\begin{enumerate}
\item is the restriction of a $\mathfrak{K}_P$-pair on the whole domain of $\mathfrak{B}$, fat on the
enlargement $X^+$ of $X$ entering \eqref{4.26:eq-regularity-space-b}, in the sense of that display
read with its existential quantifier;
\item obeys $\mathfrak{L}_n$ on $X\cap Z^{\mathrm{on}}$;
\item has entries $5$ and $8$ at $\mathfrak{T}$-locations obeying the clauses of the source;
\item has its other entries at $\mathfrak{T}$-locations, the composite entries at every
$(\Sigma^{\ast})$-triple $X'\subset X$ included (the box clauses \eqref{4.26:eq-inherited-domain-a}
and \eqref{4.26:eq-regularity-space-a}), exceeding the source thresholds at $y$ by less than
$\mathrm{room}^{\flat}(y)$, in units $r_a(y)$ for the direct entries and in units $r_e^{(p)}(y)$ for
the composite entries at a $(\Sigma^{\ast})$-triple, where
\begin{equation}\label{4.26:eq-inherited-room-1}
\mathrm{room}^{\flat}(y):=\min\{\tfrac12 r_{\min}(y),\beta\}\ge\tfrac12\beta\,e_\ast(m(y)).
\end{equation}
\end{enumerate}
On the shells $m<k$ that the stage leaves untouched, $\mathrm{room}^{\flat}(y)\ge\frac12\beta\theta_S
2^{-(k-m)}$, and on $\Omega_k\cap\mathfrak{T}=(\Omega_k\setminus Z)\cup(Z\cap\Omega_k)$,
$\mathrm{room}^{\flat}(y)\ge\frac12\beta$; these two bounds are the lower bounds on the schedule
$e_\ast$ of \eqref{4.26:eq-room-table-a}, taken at $\mathsf{c}_e=\theta_S$. For $a\in\{5,8\}$ the
threshold is at least
$\theta_{\mathrm{src},a}$: the room is non-negative, and these thresholds are never incremented.
\end{corollary}

\begin{proof}
The excesses are read off from the bounds on the stage increments in part~(c) of
Theorem~\ref{4.26:thm-inherited-expansion}, \eqref{4.26:eq-inherited-expansion-a}, and from the room
table of Lemma~\ref{4.26:lem-room-table}, \eqref{4.26:eq-room-table-a}. The numerical lower bounds
rest on the following checks.

\emph{Untouched shells and $\Omega_k$.} On a shell $m<k$ write $n_W=n_W(y)=k-1-m$. The inequality
$\frac12\mathsf{w}_0\beta e^{-n_W}\le\frac12\beta\theta_S2^{-(n_W+1)}$ is equivalent to
$\mathsf{w}_0\le\frac12\theta_S(e/2)^{n_W}$. Since $(e/2)^{n_W}\ge1$, it is implied by
$\mathsf{w}_0\le\theta_S/12$, which holds because $\theta_S\ge\frac12$ and $\mathsf{w}_0=\frac1{24}$
(fixed in the next step). Since $2^{-(n_W+1)}=2^{-(k-m)}$ and $r_T\ge\beta\theta_S2^{-(k-m)}$ by
\eqref{4.26:eq-room-table-a}, this gives $\frac12\mathsf{w}_0\beta e^{-n_W}\le\frac12 r_T$, which is
the first line of \eqref{4.26:eq-inherited-room-a}. On $\Omega_k$ the check for the first line is
$\frac13\mathsf{w}_0\beta\le\frac12\beta$, with $\frac13\mathsf{w}_0\beta$ the bound that part~(c) of
Theorem~\ref{4.26:thm-inherited-expansion} gives there for the consumed increments; it holds since
$\mathsf{w}_0\le1$.

\emph{The case $T\in\mathcal{T}^{\mathbb{E}\mathbb{R}}$, $\ell_T=m$.} The inequality
$2\mathsf{w}_0\beta\le r_T/6$ is implied, through the room table, by
\begin{equation*}
\mathsf{w}_0\le\tfrac1{12}\min\{\tfrac12,\theta_S\}=\tfrac1{24};
\end{equation*}
this is the condition that fixes $\mathsf{w}_0=\frac1{24}$. With it,
$\frac23 r_T-2\mathsf{w}_0\beta\ge\frac23 r_T-\frac16 r_T=\frac12 r_T$, the second line of
\eqref{4.26:eq-inherited-room-a}. The bound is tight for the majorants at $m=k-1$, where
$\frac23\cdot\frac12\beta-\frac1{12}\beta=\frac14\beta$.

\emph{The scale-drop case $\ell_T<m$.} Recall $\beta=\frac15$. Part~(b) of the room table gives
$\mathbf{r}^T_a\le\theta_{T,a}/(1-\beta)$ and $\theta_{T,a}\ge2r_a$; moreover
$\theta_{\mathrm{src},a}\ge(1-\beta)r_a$. For $T\in\mathcal{T}^{\mathbb{E}\mathbb{R}}$ the threshold is
$\frac23\theta_{T,a}+\frac13\theta_{\mathrm{src},a}-\sum_l\iota^1_l\mathbf{r}^T_a$; part~(c) of
Theorem~\ref{4.26:thm-inherited-expansion}, \eqref{4.26:eq-inherited-expansion-a}, gives
$\sum_l\iota^1_l\le2\mathsf{w}_0\beta$, and
$2\mathsf{w}_0\beta\,\mathbf{r}^T_a\le\frac12\mathsf{w}_0\theta_{T,a}$; hence
\begin{equation*}
\begin{split}
\theta^{(n)}_{T,a}(y)
&\ge\tfrac23\theta_{T,a}+\tfrac13\theta_{\mathrm{src},a}-2\mathsf{w}_0\beta\,\mathbf{r}^T_a\\
&\ge(\tfrac23-\tfrac12\mathsf{w}_0)\cdot2r_a+\tfrac13\cdot0.8\,r_a=(1.6-\mathsf{w}_0)r_a\\
&\ge1.55\,r_a\ge1.5\,r_a,
\end{split}
\end{equation*}
which is the third line of \eqref{4.26:eq-inherited-room-a} for $T\in\mathcal{T}^{\mathbb{E}\mathbb{R}}$.
For $T\in\mathcal{T}^{\mathbb{B}}$ with $\ell_T<m$, the bound
$\sum_l\iota_l\le\frac12\mathsf{w}_0\beta e^{-n_W(y)}\le\frac12\mathsf{w}_0\beta$ of the first line of
\eqref{4.26:eq-inherited-room-a} and $\beta/(1-\beta)=\frac14$ give
$\sum_l\iota_l\mathbf{r}^T_a\le\frac18\mathsf{w}_0\theta_{T,a}$, whence
$\theta^{(n)}_{T,a}\ge\theta_{T,a}(1-\frac18\mathsf{w}_0)\ge1.9\,r_a$.

\emph{A lifted point of an older live piece.} At $y\in\mathsf{L}_{\mathfrak{a}}$, by item~(5) of
\eqref{4.26:eq-older-pieces-domain-a}, the source threshold is bounded below only by
$\theta_{\mathrm{src},a}\ge\varphi_{\mathrm{arrest}}r_a$, where $\varphi_{\mathrm{arrest}}=1-4.51\beta=0.098$
is the floor for the arrested source factor entering item~(5) of
\eqref{4.26:eq-older-pieces-domain-a}, while
$\theta_{T,a}\ge2.1\,r_a$. Hence, using again $\sum_l\iota^1_l\le2\mathsf{w}_0\beta$, the part
$\frac23\theta_{T,a}-\sum_l\iota^1_l\mathbf{r}^T_a$ of the threshold satisfies
\begin{equation*}
\begin{split}
\tfrac23\theta_{T,a}-\textstyle\sum_l\iota^1_l\mathbf{r}^T_a
&\ge\tfrac23\theta_{T,a}-2\mathsf{w}_0\beta\,\mathbf{r}^T_a\\
&\ge(\tfrac23-\tfrac12\mathsf{w}_0)\cdot2.1\,r_a\\
&=1.356\,r_a\ge1.35\,r_a,
\end{split}
\end{equation*}
and adding the third $\frac13\theta_{\mathrm{src},a}\ge\frac13\varphi_{\mathrm{arrest}}r_a$ of the source
threshold gives at least $1.38\,r_a$. For $T\in\mathcal{T}^{\mathbb{B}}$ whose clauses were fixed before
the arrest $\mathfrak{a}$,
$\sum_l\iota_l\le\frac12\mathsf{w}_0\beta$ gives $\sum_l\iota_l\mathbf{r}^T_a\le\frac18\mathsf{w}_0
\theta_{T,a}$, so the threshold is at least $2.1(1-\frac18\mathsf{w}_0)r_a\ge2.08\,r_a$. The quantities
tested against these thresholds stay below them: the source together with the room is at most
$(F^{(k)}(y)+\beta)r_a\le1.2\,r_a$, since $F^{(k)}(y)\le1$ by item~(1) of
Lemma~\ref{4.26:lem-older-pieces-domain}; the tube's image $(\mathcal{U}^{\iota},\mathcal{J}^{\iota})$
is bounded by $1.1\,r_a$; and the quantity controlled by the unit rule \eqref{4.26:eq-clauses-units-b}
is at most $1.12\,r_a$. All three lie below $1.35\,r_a$, hence below the thresholds obtained in this
paragraph.

\emph{The unit rule, $a=4$, every $p$.} Write $s^{(n)}_{T,4}$ for the factor with
$\theta^{(n)}_{T,4}(y;p)=s^{(n)}_{T,4}\mathbf{r}^{T,(p)}_4$. By the rule
\eqref{4.26:eq-clauses-units-b}, $s^{(n)}_{T,4}\ge\frac23s_T-2\mathsf{w}_0\beta$, and this is at least
$0.516$ since $s_T\ge1-\beta$ by part~(b) of the room table; by \eqref{4.26:eq-clauses-units-a}
\begin{equation*}
\mathbf{r}^{T,(p)}_4/r_4^{(p)}\ge\bigl(L^2/(1+\beta_0)\bigr)^{m-\lambda}\ge56,
\end{equation*}
with $\lambda$ as in \eqref{4.26:eq-clauses-units-a}; the last inequality holds since $m-\lambda\ge1$
and $L\ge L_0$, the constant $L_0$ of the quantifier prefix (Notation~\ref{2.1:not-prefix}) being
taken with $L_0^2\ge56(1+\beta_0)$. Hence
$\theta^{(n)}_{T,4}(y;p)\ge28\,r_4^{(p)}\ge1.5\,r_4^{(p)}$.

\emph{Comparison with the tested quantities.} Against the thresholds above, the source together
with $\mathrm{room}^{\flat}$ is at most $(1+\beta)r_a=1.2\,r_a$. For the composite entries at
any triple, the source is at most $1\cdot r_e^{(p)}$, because every $(\Sigma^{\ast})$-triple carries the
factor $s^{(k)}_m\le1$;
adding $\mathrm{room}^{\flat}\le\beta$ gives at most $1.2\,r_e^{(p)}<1.383\,r_e^{(p)}$.

\emph{Direct entries and $a\in\{5,8\}$.} On the direct entries,
$(1+2\beta-2\mathsf{w}_0\beta)-1=0.383\ge\beta$. For $a\in\{5,8\}$ no stage increment is subtracted,
so the two formulas of Definition~\ref{4.26:def-inherited-domain} give $\theta_{T,a}$ and
$\theta_{T,a}-\frac13(\theta_{T,a}-\theta_{\mathrm{src},a})$; since
$\theta_{T,a}\ge\theta_{\mathrm{src},a}$, both are at least $\theta_{\mathrm{src},a}$.
\end{proof}

\begin{remark}
At the same locations the thresholds of the spaces $\tilde{U}^{(n)c,\Sigma\flat}_k$ of
\cite{B-CMP119} lie below the source thresholds by $\sum_{l\le n}\rho^{(l)}_m$, by the identity in
part~(b) of Theorem~\ref{4.26:thm-inherited-expansion}. By Corollary~\ref{4.26:cor-inherited-room},
the thresholds defining $\mathfrak{D}^{\flat}_n$ instead lie above the source thresholds by a positive
room.
\end{remark}

\begin{remark}[Per-species thresholds and the window for $\theta_S$]\label{4.26:rem-theta-window}
Let $\mathsf{w}_0=1/24$, let $\theta_S$ be the newborn number of
Lemma~\ref{4.26:lem-room-table}, let $y\in\mathfrak{T}$, $m:=m(y)$, $a\notin\{5,8\}$, and put
\begin{equation}\label{4.26:eq-theta-window-1}
s:=k-m,\qquad \mathsf{u}_s:=\tfrac{9}{16}\,\theta_S\,2^{-s}.
\end{equation}
The uniqueness theorem uses the inclusion
$\mathfrak{U}[\mathsf{T}^{\flat}_i]\subset\mathfrak{D}^{\flat}_i(X)$ of the clause domain
$\mathfrak{U}[\mathsf{T}^{\flat}_i]$ of its level-$i$ object $\mathsf{T}^{\flat}_i$ (notation of that
theorem) into the inherited domain \eqref{4.26:eq-inherited-domain-b}, and its tables bound the
thresholds required for this inclusion by
\begin{equation}\label{4.26:eq-theta-window-2}
\bigl(F^{(k')}+\tfrac{9}{16}\,\theta_S\,\beta\,2^{-(k'-\ell)}\bigr)\,r_a(y);
\end{equation}
here $k'$ and $\ell$ are the two levels of those tables, and read at $y$ in the present notation
$k'-\ell=k-m=s$. In units $\beta r_a(y)$ the rise of these tables over the source threshold
$\theta_{\mathrm{src},a}(y)$ is therefore $\mathsf{u}_s$. Then:
\begin{enumerate}
\item[(i)] the room $\mathrm{room}^{\flat}$ of \eqref{4.26:eq-inherited-room-a} on the
shells of $W$, as it stands after the newborn schedule of Lemma~\ref{4.26:lem-room-table}(c), does
not give this inclusion, since $\tfrac12\theta_S2^{-s}<\mathsf{u}_s$ for every $\theta_S>0$;
\item[(ii)] for every $\theta_S$ in the window
\begin{equation}\label{4.26:eq-theta-window-3}
\tfrac12\le\theta_S\le\tfrac89,
\end{equation}
in the setting of Corollary~\ref{4.26:cor-inherited-room} (every $n\le N$, every
$(X,\mathcal{S})$, every term $T$ with $(y,a)\in\operatorname{Cl}(T)$, $y\in X\cap\mathfrak{T}$,
$X_T\subset X$), the threshold $\theta^{(n)}_{T,a}(y)$ of the diminished term clauses
$\mathfrak{I}_n$ exceeds the source threshold by at least the rise $\mathsf{u}_s$ (in units
$\beta r_a(y)$) in each of the families: the newborn terms $\mathbb{B}^{(k)}$ with $s\ge1$; the
old terms $\mathbf{B}^{(j)}$, $j\le k-1$, with $s\ge1$; the terms
$\mathbb{E}^{(m)},\mathbb{R}^{(m)}$ with $s=1$ and with $s\ge2$; and the top layer $s=0$
($\mathbb{B}^{(k)}$, $\mathbb{E}^{(k)}$, $\mathbb{R}^{(k)}$); for scale drops and for
$\mathfrak{N}^{(n)}$ the threshold lies above the tables' bound; the upper end $8/9$ is the exact
condition for the families $\mathbb{E}^{(k-1)},\mathbb{R}^{(k-1)}$ at $s=1$;
\item[(iii)] the clauses $(\mathrm{box})$ and $(\mathrm{box})^{\mathrm{sh}}$ carry the same
thresholds, and are covered if and only if the tables of the uniqueness theorem bound their second
and their fourth entry at every $(\Sigma^{\ast})$-triple;
\item[(iv)] the two bounds on the stage increments used below, $\sum_l\iota_l$ and
$\sum_l\iota^1_l$, are absolute inequalities, not among the conditions collected under
$\gamma_U$.
\end{enumerate}
\end{remark}

\begin{proof}
Item~(i) is the inequality $\tfrac12<\tfrac{9}{16}$ multiplied by $\theta_S2^{-s}>0$.

For item~(ii) we take the clause families in turn; all excesses are in units $\beta r_a(y)$, and
they are those of Corollary~\ref{4.26:cor-inherited-room}, evaluated with the rooms of
Lemma~\ref{4.26:lem-room-table}.

\emph{Newborn terms $\mathbb{B}^{(k)}$, $s\ge1$.} The excess is at least
$\theta_S2^{-s}-\tfrac12\mathsf{w}_0e^{-(s-1)}$. This is at least $\mathsf{u}_s$ if and only if
$\tfrac{7}{16}\theta_S2^{-s}\ge\tfrac12\mathsf{w}_0e^{-(s-1)}$, that is,
\begin{equation}\label{4.26:eq-theta-window-4}
\mathsf{w}_0\le\tfrac{7}{16}\,\theta_S\,(e/2)^{s-1},
\end{equation}
which holds under \eqref{4.26:eq-theta-window-3} because $(e/2)^{s-1}\ge1$ and
$\tfrac{7}{16}\cdot\tfrac12>\tfrac1{24}$. If the margin $\tfrac{3}{16}$ of the tables of the
uniqueness theorem is kept, so that the coefficient $\tfrac7{16}$ in
\eqref{4.26:eq-theta-window-4} is replaced by $\tfrac3{16}$, then with $\mathsf{w}_0=1/24$ the
condition reads
\begin{equation*}
\theta_S\ge\tfrac{16}{3}\cdot\tfrac1{24}\,(2/e)^{s-1}=\tfrac29\,(2/e)^{s-1},
\end{equation*}
which is implied by $\theta_S\ge\tfrac29$, since $2/e<1$ and $s\ge1$; and $\tfrac29<\tfrac12$.

\emph{Old terms $\mathbf{B}^{(j)}$, $j\le k-1$, $s\ge1$.} By
Lemma~\ref{4.26:lem-room-table}(a) the room is at least $\beta2^{-s}$, so the excess is at least
$2^{-s}-\tfrac12\mathsf{w}_0e^{-(s-1)}$. Since $\theta_S\le\tfrac89<1$, this is at least the
excess of the previous family, and the condition holds a fortiori.

\emph{Terms $\mathbb{E}^{(m)},\mathbb{R}^{(m)}$, $s=1$.} By
Lemma~\ref{4.26:lem-room-table}(a) the room is $\beta(1-2^{-s})=\beta/2$; two thirds of it are
available and the hereditary spend is at most $2\mathsf{w}_0\beta$, so the excess is at least
\begin{equation*}
\tfrac23\cdot\tfrac12-2\mathsf{w}_0=\tfrac13-\tfrac1{12}=\tfrac14 .
\end{equation*}
This is at least $\mathsf{u}_1=\tfrac{9}{32}\theta_S$ if and only if $\theta_S\le\tfrac89$. This is
the upper end of \eqref{4.26:eq-theta-window-3}; it comes from the reserve of one third of the
room under $\mathbb{E}^{(k-1)}$ at $s=1$, where the excess $\tfrac14\beta$ coincides with
$\tfrac12\beta2^{-1}$.

\emph{Terms $\mathbb{E}^{(m)},\mathbb{R}^{(m)}$, $s\ge2$.} Now $1-2^{-s}\ge\tfrac34$, so the
excess is at least $\tfrac23\cdot\tfrac34-\tfrac1{12}=\tfrac5{12}$, while
\begin{equation*}
\mathsf{u}_s\le\mathsf{u}_2=\tfrac{9}{64}\theta_S\le\tfrac{9}{64}<\tfrac{5}{12}.
\end{equation*}

\emph{Top layer $s=0$.} For $\mathbb{B}^{(k)}$ the excess is at least
$1-\tfrac13\mathsf{w}_0=\tfrac{71}{72}$; for $\mathbb{E}^{(k)},\mathbb{R}^{(k)}$ it is at least
\begin{equation*}
\tfrac23-\tfrac43\mathsf{w}_0=\tfrac23-\tfrac1{18}=\tfrac{11}{18}.
\end{equation*}
Both exceed $\mathsf{u}_0$: indeed
\begin{equation*}
\mathsf{u}_0=\tfrac{9}{16}\theta_S\le\tfrac{9}{16}=\tfrac{81}{144}<\tfrac{88}{144}
=\tfrac{11}{18}<\tfrac{71}{72}.
\end{equation*}

\emph{Scale drops and $\mathfrak{N}^{(n)}$.} Here the thresholds are bounded absolutely: for
$\ell_T<m$ Corollary~\ref{4.26:cor-inherited-room} gives $\theta^{(n)}_{T,a}(y)\ge1.5\,r_a(y)$,
and the threshold of $\mathfrak{N}^{(n)}$ is at least $1.383\,r_a(y)$. The tables' bound is
$(F^{(k')}+\mathsf{u}_s\beta)\,r_a(y)\le1.12\,r_a(y)$, which lies below both.

Item~(iii): the clauses $(\mathrm{box})$ and $(\mathrm{box})^{\mathrm{sh}}$ of $\mathfrak{I}_n$ carry
the thresholds just treated, so the comparisons above apply to them exactly when the tables of the
uniqueness theorem bound their second and their fourth entry at every $(\Sigma^{\ast})$-triple;
this is a condition on those tables.

Item~(iv): the two inequalities $\sum_l\iota_l\le\tfrac12\mathsf{w}_0\beta e^{-n_W(y)}$ and
$\sum_l\iota^1_l\le2\mathsf{w}_0\beta$ of Corollary~\ref{4.26:cor-inherited-room}, used above with
$\mathsf{w}_0=1/24$, hold unconditionally; none of the steps above adds a condition to those
collected under $\gamma_U$.

Collecting the families: every condition of item~(ii) holds for $\theta_S\in[\tfrac12,\tfrac89]$;
from below the families require only $\theta_S\ge\tfrac29$ (with the margin $\tfrac3{16}$ kept),
which the lower end $\tfrac12$ of the window satisfies. This proves the remark.
\end{proof}

\begin{corollary}[Bound on the summed term at real exterior data]\label{4.26:cor-summed-term-bound}
Assume the hypotheses of Theorem~\ref{4.26:thm-inherited-expansion}. Let
$\mathbb{C}^{(n)}_k(X,\mathcal{S};\cdot,\mathbf{A})$ be the terms of the representation of that
theorem, the lengths of the domain $X$ being measured at a level $m\le k$, and put
\begin{equation}\label{4.26:eq-summed-term-bound-1}
\mathbb{C}^{(n)}_k(X;\cdot,\mathbf{A})
:=\sum_{\mathcal{S}}\mathbb{C}^{(n)}_k(X,\mathcal{S};\cdot,\mathbf{A}),
\end{equation}
a finite sum. Let $n\le N$ and let $\mathbf{A}\in\mathfrak{X}^{\sharp}$ be a real exterior datum.
Then the map
\begin{equation*}
(\mathbb{U},\mathbb{J})\longmapsto\mathbb{C}^{(n)}_k\bigl(X;(\mathbb{U},\mathbb{J}),\mathbf{A}\bigr)
\end{equation*}
is holomorphic on the inherited domain $\mathfrak{D}^{\flat}_n(X)$ of
\eqref{4.26:eq-inherited-domain-b}, is local in $X$ and is jointly $\mathcal{G}$-invariant, and
\begin{equation}\label{4.26:eq-summed-term-bound-2}
\bigl|\mathbb{C}^{(n)}_k\bigl(X;(\mathbb{U},\mathbb{J}),\mathbf{A}\bigr)\bigr|
\le C_0^{\sharp}\,e^{-a_m d_m(X)}
\qquad\text{for }(\mathbb{U},\mathbb{J})\in\mathfrak{D}^{\flat}_n(X),
\end{equation}
uniformly in $\mathbf{A}$; for $X$ in the class $\mathcal{K}_{\mathsf{d}}$ the length $d_m(X)$ in
\eqref{4.26:eq-summed-term-bound-2} is replaced by the relative length $\mathsf{d}_m(X)$, by the
relative form of the decay bound on the class $\mathcal{K}_{\mathsf{d}}$. Here
$C_0^{\sharp}$ and $a_m$ are as in Notation~\ref{4.26:not-slot-notation}. At $m=k$,
where $a_k=(1+3\beta)\kappa$, the bound \eqref{4.26:eq-summed-term-bound-2} is the bound
\cite[(1.70)]{B-CMP122a}; for $m<k$ the rate is the constant $a_m$ of
Notation~\ref{4.26:not-slot-notation}.
\end{corollary}

\begin{proof}
We first show that $\mathfrak{X}^{\sharp}\subset\mathfrak{D}_{\mathcal{S}}(X)$ for every
$\mathcal{S}$ occurring in \eqref{4.26:eq-summed-term-bound-1}. The strong coordinates of a datum
in $\mathfrak{X}^{\sharp}$ lie in the slot domain $\mathfrak{D}_{\mathcal{S}}(X)$ by the definition
of the slot data in Notation~\ref{4.26:not-slot-notation}. For the weak coordinates, let
$\mathfrak{b}^{\sharp}$ be their radius in $\mathfrak{X}^{\sharp}$, $\mathfrak{b}$ the
corresponding radius of the slot set, and $\varrho$ the radius of the weak coordinates of
$\mathfrak{D}_{\mathcal{S}}(X)$; then
\begin{equation*}
\mathfrak{b}^{\sharp}\le C_{\mathrm{far}}\,\mathfrak{b}<\varrho ,
\end{equation*}
the first inequality being the comparison of the radii of the weak coordinates that accompanies
the slot data, with $C_{\mathrm{far}}$ the constant of that comparison, and the second holding
because $r_A>C_{\mathrm{far}}$. Hence the weak coordinates
lie in $\mathfrak{D}_{\mathcal{S}}(X)$ as well.

By Theorem~\ref{4.26:thm-inherited-expansion}\,(a), each $\mathbb{C}^{(n)}_k(X,\mathcal{S})$ is a
stored term with slot datum on $\mathfrak{D}^{\flat}_n(X)\times\mathfrak{D}_{\mathcal{S}}(X)$.
Since $\mathbf{A}$ lies in $\mathfrak{D}_{\mathcal{S}}(X)$ for every $\mathcal{S}$, each summand of
\eqref{4.26:eq-summed-term-bound-1}, evaluated at $\mathbf{A}$, is defined on
$\mathfrak{D}^{\flat}_n(X)$ and is holomorphic there, local in $X$ and jointly
$\mathcal{G}$-invariant; a finite sum of such functions has the same three properties.

For the bound, let $(\mathbb{U},\mathbb{J})\in\mathfrak{D}^{\flat}_n(X)$ and fix a point
$x\in X$. With $\kappa_A$ as in Notation~\ref{4.26:not-slot-notation}, by the triangle
inequality, and since $e^{\kappa_A|\mathcal{S}|}\ge1$,
\begin{align*}
\Bigl|\sum_{\mathcal{S}}\mathbb{C}^{(n)}_k\bigl(X,\mathcal{S};(\mathbb{U},\mathbb{J}),
\mathbf{A}\bigr)\Bigr|
&\le\sum_{\mathcal{S}}e^{\kappa_A|\mathcal{S}|}\,
\sup\bigl|\mathbb{C}^{(n)}_k(X,\mathcal{S})\bigr| \\
&\le\bigl\|\mathbb{C}^{(n)}\bigr\|^{A,\flat}_{(n)}\,e^{-a_m d_m(X)} ,
\end{align*}
where the supremum is taken over
$\mathfrak{D}^{\flat}_n(X)\times\mathfrak{D}_{\mathcal{S}}(X)$, which contains the point
$\bigl((\mathbb{U},\mathbb{J}),\mathbf{A}\bigr)$ by the first paragraph, and the second inequality
is the definition of the weighted norm $\|\cdot\|^{A,\flat}_{(n)}$ of
Notation~\ref{4.26:not-slot-notation}, read at the point $x$, which the domain $X$ contains;
for $X\in\mathcal{K}_{\mathsf{d}}$ the same chain holds
with $\mathsf{d}_m(X)$ in place of $d_m(X)$, by the relative form of the decay bound on the class
$\mathcal{K}_{\mathsf{d}}$. The right-hand side does not depend on
$\mathbf{A}$, and by Theorem~\ref{4.26:thm-inherited-expansion} the norm
$\|\mathbb{C}^{(n)}\|^{A,\flat}_{(n)}$ of its term family is at most $C_0^{\sharp}$; this gives
\eqref{4.26:eq-summed-term-bound-2} uniformly in $\mathbf{A}$.
\end{proof}

\begin{corollary}[A further holomorphic parameter under two hypotheses]
\label{4.26:cor-extra-parameter}
For every region $X$ let $\mathcal{W}(X)$ be a product of complex polydiscs, with
$\mathcal{W}(X){\restriction}_{X'}\subset\mathcal{W}(X')$ for every region $X'\subset X$; in
particular $\mathcal{W}(X_T)$ is defined for every localisation domain $X_T$. Let $w$ range over
$\mathcal{W}(X)$.
Assume the following two hypotheses.
\begin{itemize}
\item[(w1)] Each of the following objects, with its localisation domain $X_T$, depends
holomorphically on $w\in\mathcal{W}(X_T)$, jointly with its other complex variables:
\begin{itemize}
\item every term of the input family assumed in
Theorem~\ref{4.26:thm-inherited-expansion};
\item every piece of the types (S3) and (S4) named in the hypothesis of the lemma on holomorphic
parameters of the step's representation;
\item every piece of the Wilson-action difference $\Delta W$.
\end{itemize}
This dependence holds on the same spaces $\mathfrak{U}_T\times\mathfrak{D}_{\mathcal{S}}$ as without the
parameter $w$, under the same majorants. The thresholds $\theta_{T,a}$ do not depend on $w$.
\item[(w2)] The following objects do not depend on $w$: the Gaussian integrals of the
renormalisation transformation of the step; the
products $\chi^{(j)\sharp}\chi'^{(j)}$ of characteristic functions; the stage layer $\mathbb{H}$;
the local current slot
$\mathcal{J}^{\sharp}[1]$; the partition of unity $\zeta_{\square}$; the parameters $\sigma$ and
$t$ of the proof family $\mathfrak{W}^{(l),\natural}(Y)$, and the parameter $u$; and the operators
$\mathbb{C}_G$ of the representation of the step.
\end{itemize}
Then clauses (a)--(c) of Theorem~\ref{4.26:thm-inherited-expansion},
Corollary~\ref{4.26:cor-inherited-room} and Corollary~\ref{4.26:cor-summed-term-bound} hold for the
terms $\mathbb{C}^{(n)}_k(X,\mathcal{S};w)$ on
\begin{equation*}
\mathfrak{D}^{\flat}_n(X)\times\mathfrak{D}_{\mathcal{S}}(X)\times\mathcal{W}(X),
\end{equation*}
holomorphically in $w$ and with the same constants. Moreover, the inherited domain
$\mathfrak{D}^{\flat}_n(X)$ does not depend on $w$.
\end{corollary}

\begin{proof}
We first treat the domain. By \eqref{4.26:eq-inherited-domain-b}, the inherited domain
$\mathfrak{D}^{\flat}_n$ is cut out by three kinds of conditions: the thresholds
$\theta^{(n)}_{T,a}$, the stage increments $\iota_l$, and the lettered clauses $\mathfrak{L}_n$.
The thresholds $\theta^{(n)}_{T,a}$ are formed from the thresholds $\theta_{T,a}$, which do not
depend on $w$ by (w1), and from the stage increments $\iota_l$, which do not involve $w$.
None of these conditions involves $w$. Hence $\mathfrak{D}^{\flat}_n(X)$ does not depend
on $w$.

By (w2), the following ingredients of the construction do not depend on $w$: the
Gaussian integrals, the products $\chi^{(j)\sharp}\chi'^{(j)}$, $\mathbb{H}$,
$\mathcal{J}^{\sharp}[1]$, $\zeta_{\square}$, $\sigma$, $t$, $u$ and the operators
$\mathbb{C}_G$; so in the bounds of Lemma~\ref{4.26:lem-stage-expansion} the dependence on
$w$ comes from the objects of (w1).

Lemma~\ref{4.26:lem-family-consumption} involves no term of the representation. Its statement
concerns only the proof family and the background pair $(\mathbb{U},\mathbb{J})$, and it is used
unchanged.

In Lemma~\ref{4.26:lem-stage-expansion}, every bound is built from the quantities
$\nu^{\natural}(e)$. These are suprema, and in the present setting they are taken over $w$ as well.
To these bounds we apply the lemma on holomorphic parameters of the step's representation, with
\begin{equation*}
\lambda:=(\text{exterior slot variables},\,w).
\end{equation*}
The hypothesis of that lemma is that the pieces of types (S3) and (S4) do not depend on $\lambda$.
Here we read this hypothesis in the weaker form that these pieces are holomorphic in $\lambda$ under
majorants independent of $\lambda$. This weaker form is all that the proof of that lemma uses. That
proof bounds each term by a majorant independent of $\lambda$, and it obtains holomorphy in
$\lambda$ of the resulting sums from Weierstrass's theorem: a locally uniform limit of holomorphic
functions is holomorphic. This is the form in which the lemma is used in the proof of the
correlation theorem. Hypothesis (w1) supplies exactly the weaker form for
$\lambda=(\text{exterior slot variables},w)$. The majorants are the same as without $w$, so the
constants are unchanged.

Therefore the conclusions of Theorem~\ref{4.26:thm-inherited-expansion},
Corollary~\ref{4.26:cor-inherited-room} and Corollary~\ref{4.26:cor-summed-term-bound} hold for
$\mathbb{C}^{(n)}_k(X,\mathcal{S};w)$, holomorphically in $w$ and with the same constants.
\end{proof}

Hypotheses (w1) and (w2) are assumptions of the corollary and are not derived in this section. The
pieces of $\Delta W$ appear in (w1) because, in the application of the corollary in the correlation
theorem, the vertex of the observable sits in the Wilson-action difference.

\begin{proposition}[The image of the source lies in the inherited domain]\label{4.26:prop-source-image}
Let $X\subset Y_4$. At every real configuration $\Phi=\varpi\in\operatorname{supp}\Pi$ put
\[
\mathfrak{E}:=\mathbb{C}^{(n)}_k(X;\cdot,\Phi_A),\qquad
\mathfrak{D}_{\mathfrak{E}}:=\mathfrak{D}^{\flat}_n(X).
\]
Here $Y_4$ is the region of the site within which the localisation domains $X$ of the summed term
are taken. The symbol $\Pi$ denotes the measure on the configurations $\Phi$ of the site; its support
supplies the real points $\Phi=\varpi$ at which $\mathfrak{E}$ is read.
Further, $\mathbb{C}^{(n)}_k(X;\cdot,\cdot)$ is the summed term of
Corollary~\ref{4.26:cor-summed-term-bound}, the exterior (slot) datum
$\Phi_A\in\mathfrak{X}^{\sharp}$ is the one at the point $\Phi$, and
$\mathfrak{D}^{\flat}_n(X)$ is the inherited domain of Definition~\ref{4.26:def-inherited-domain}.
For every $(\mathbb{U},\mathbb{J})\in\mathfrak{S}_{\mathrm{src}}$, the restriction to $X$ of its
image under the tube satisfies
\[
(\mathcal{U}^{\iota}(\mathbb{U}),\mathcal{J}^{\iota})\restriction X\in\mathfrak{D}_{\mathfrak{E}} .
\]
On $\mathfrak{D}_{\mathfrak{E}}$ the functional $\mathfrak{E}$ is holomorphic in the pair and obeys
the bound of Corollary~\ref{4.26:cor-summed-term-bound}, at every real exterior datum.
\end{proposition}

\begin{proof}
We begin with the second assertion. By definition $\mathfrak{D}_{\mathfrak{E}}=\mathfrak{D}^{\flat}_n(X)$.
Corollary~\ref{4.26:cor-summed-term-bound} asserts, for every real exterior datum $\mathbf{A}$, that
$(\mathbb{U},\mathbb{J})\mapsto\mathbb{C}^{(n)}_k(X;(\mathbb{U},\mathbb{J}),\mathbf{A})$ is holomorphic
on $\mathfrak{D}^{\flat}_n(X)$. It also asserts that the bound stated there holds uniformly in
$\mathbf{A}$. Applied to the real exterior datum $\mathbf{A}=\Phi_A$, this is the second assertion.
It remains to prove the
first assertion, namely that the restricted image of a pair $(\mathbb{U},\mathbb{J})\in\mathfrak{S}_{\mathrm{src}}$
satisfies every clause of $\mathfrak{D}^{\flat}_n(X)$ in \eqref{4.26:eq-inherited-domain-b}.
Throughout, the pairs of $\mathfrak{S}_{\mathrm{src}}$ are those of the arrested source
$\tilde{U}^{\mathrm{arrest}}_k$ together with all its $(\Sigma^{\ast})$-triples. We take
these clauses in the order of that definition. Points of $X\cap Z^{\mathrm{on}}$ are treated in part
$(\alpha)$, points of $X\cap\mathfrak{T}$ in part $(\beta)$, and the clause requiring a margin of one
half in part $(\gamma)$. The argument rests on the smallness condition and the numerical inequality
\begin{equation}\label{4.26:eq-source-image-a}
16\max\{P,P',P''\}\le\beta\theta_S,
\qquad
1.1<1+2\beta-2\mathsf{w}_0\beta=1.383 .
\end{equation}

$(\alpha)$ On $X\cap Z^{\mathrm{on}}$, membership in $\mathfrak{D}^{\flat}_n(X)$ is the condition
$\mathfrak{L}_n$ of Definition~\ref{4.26:def-inherited-domain}. On the components of $Z^{\mathrm{on}}$
that survive in the run, the verification that the image satisfies the lettered clauses making up
$\mathfrak{L}_n$ is that of the theorem on inheritance of the domain, unchanged. On a component
$W$ excluded from the run, $\mathfrak{L}_n$ is by Definition~\ref{4.26:def-inherited-domain} the arrested
clause. That clause is one of the five conjuncts of the site hypothesis of the summed term. By the
arrest theorem, the image under the tube of the arrested
source $\tilde{U}^{\mathrm{arrest}}_k$ lies in the arrested clauses with positive room. Hence the clause
holds on $W$.

$(\beta)$ On $X\cap\mathfrak{T}$, membership is the condition $\mathfrak{I}_n$ of
Definition~\ref{4.26:def-inherited-domain}. We distinguish the clause indices $a$.

For $a\in\{5,8\}$, the entries $Q_5,Q_8$ of the image are those of the source. Therefore
$Q_a<\theta_{\mathrm{src},a}$, and $\theta_{\mathrm{src},a}$ does not exceed the threshold of
$\mathfrak{I}_n$ for these two indices. Moreover $Q_a<(1+2\beta)r_a$.

For $a\notin\{5,8\}$, it suffices that the increments added by the layer of the operation $\mathbf{R}$ are
at most one quarter of the room $\operatorname{room}^{\flat}$ in \eqref{4.26:eq-inherited-room-a}.
Indeed, Corollary~\ref{4.26:cor-inherited-room} shows that the threshold of $\mathfrak{I}_n$ exceeds the
source's by this room. The bound on the increments follows from the first inequality in
\eqref{4.26:eq-source-image-a}.

In that inequality, $P$ is the constant of the proposition bounding the increments of the layer:
\[
P:=K_W\,B_H\,K_\Lambda(M)\,(1+\varrho_k)(1+\beta_0)^2\,C_{\beta_0}\,e^{-5\delta_1R_k}.
\]
Here $\varrho_k=\alpha_{1,k}/\alpha_{0,k}$ is the ratio of the two radii at level $k$, $R_k$ is the
size parameter of Ba{\l}aban's construction entering that proposition, $\delta_1$ is the decay rate in
its exponential factor, and $K_W$, $B_H$, $K_\Lambda(M)$ and $C_{\beta_0}$ are the constants of its
estimate.
The constants $P'$ and $P''$ are those of the further increment bound for the layer, which carries
the same smallness condition. In $P$, $P'$ and
$P''$ the constant $C_{\beta_0}$ is replaced by $C_{1/2}=\sqrt{2}$, and the factor $(2+k-m)^{\beta_0}$
by $(1+k-m)^{1/2}$. With these replacements the first inequality in \eqref{4.26:eq-source-image-a} keeps the
form of the smallness condition of those bounds. It is one of the conditions met through the choice of
$\gamma$: its left side carries the factor $e^{-5\delta_1R_k}$, which tends to zero.

At a point $y$ of an older live piece, the increments are at most $P\,2^{-n_W(y)}$ in unweighted units.
This holds stratum by stratum and at every age, by clause (6) of
\eqref{4.26:eq-older-pieces-domain-a} in
Lemma~\ref{4.26:lem-older-pieces-domain}, under the same ceiling
\eqref{4.26:eq-source-image-a}.

The source in all these comparisons is the arrested source together with all its
$(\Sigma^{\ast})$-triples. On a live piece $\mathfrak{a}$ these are the triples of the sub-family of
$(\Sigma^{\ast})^{\mathfrak{a}}$ retained at live pieces in the proof of
Lemma~\ref{4.26:lem-older-pieces-domain}. They are taken at the factor $F^{(k)}\omega$, in the unit of
\eqref{4.26:eq-older-pieces-domain-a}, with the block side used in the proof of
Lemma~\ref{4.26:lem-older-pieces-domain}. This is legitimate because the arrest theorem
leaves the $(\Sigma^{\flat})$ and $(\Sigma^{\ast})$ readings unchanged.

Where the class clauses of a term are read at the radii at which the term was born, the required
inclusion is given by parts (a) and (a$'$) of the extended radius lemma.

We turn to the clauses \eqref{4.26:eq-inherited-domain-a} and the shell clause of
\eqref{4.26:eq-regularity-space-a}. By the construction of $\mathfrak{S}_{\mathrm{src}}$ that accompanies
the proposition bounding the increments of the layer, the space $\mathfrak{S}_{\mathrm{src}}$ carries
every $(\Sigma^{\ast})$-triple $X'\subseteq Y_4$ at $s^{(k)}_m$. By the same proposition, the
increment bounds above hold for all these triples simultaneously. Triple by triple, the
increments are of the form treated in the statement bounding the increments one triple at a time.
Hence the argument of the preceding paragraphs applies to each triple, and the clauses
\eqref{4.26:eq-inherited-domain-a} and the shell clause of \eqref{4.26:eq-regularity-space-a} hold for
the image.

Where the scale drops, the entries of the image at points of $\mathfrak{T}$ are at most $1.1r_a$, by the
site hypothesis of the summed term.
Corollary~\ref{4.26:cor-inherited-room} gives a threshold of at least $1.5r_a$ for $\ell_T<m$, so
\[
Q_a\le 1.1\,r_a<1.5\,r_a\le\theta^{(n)}_{T,a},
\]
and the strict clause of $\mathfrak{I}_n$ holds.
At a lifted point of an older live piece the entries are likewise at most $1.1r_a$. By the same
corollary the threshold there is at least $1.35r_a$, so
\[
Q_a\le 1.1\,r_a<1.35\,r_a\le\theta^{(n)}_{T,a}.
\]

The regularity clause of $\mathfrak{N}^{(n)}$ in \eqref{4.26:eq-regularity-space-a} carries the factor
$1+2\beta-2\mathsf{w}_0\beta$, that is, the threshold $(1+2\beta-2\mathsf{w}_0\beta)\,r_a$.
By the second inequality in \eqref{4.26:eq-source-image-a}, the entries, which are at most
$1.1\,r_a$, lie strictly below this threshold. This holds for direct and composite entries
alike.

Finally, the image under the tube is the restriction of a global field, so the ambient clause
\eqref{4.26:eq-regularity-space-b} holds. On an excluded component $W$ the ambient clause is taken in
the space $\mathfrak{G}(\mathfrak{B}_n)$ of the stage-$n$ determining-set structure $\mathfrak{B}_n$.

$(\gamma)$ The clause of $\mathfrak{D}^{\flat}_n(X)$ requiring a margin of one half on
$X\cap\Lambda^{\natural}$, where $\Lambda^{\natural}$ is the set of the site on which
$\mathfrak{D}^{\flat}_n(X)$ imposes that clause, holds by the argument of part $(\alpha)$.

Parts $(\alpha)$--$(\gamma)$ verify every clause of $\mathfrak{D}^{\flat}_n(X)$ for
$(\mathcal{U}^{\iota}(\mathbb{U}),\mathcal{J}^{\iota})\restriction X$. This proves the first assertion.
\end{proof}

\begin{remark}[Scope of the image statement: real slot data only]\label{4.26:rem-real-slot-scope}
Write $\Phi$ for the datum carrying the slot argument of the functional $\mathfrak{E}$, the argument on which $\mathfrak{E}$ depends besides the background pair $(\mathbb{U},\mathbb{J})$. The slot datum is real when $\Phi$ is real, and in part~(c) the slot argument is the image $A^{[h]}$ of the slot under the gauge map. The following three points fix the scope of Theorem~\ref{4.26:thm-inherited-expansion}, Corollary~\ref{4.26:cor-inherited-room} and Proposition~\ref{4.26:prop-source-image}.
\begin{enumerate}
\item[(a)] Each of these is a statement on the pair $(\mathbb{U},\mathbb{J})$ at real slot data $\Phi$. In this form the conclusion of Theorem~\ref{4.26:thm-inherited-expansion} holds for every term $\mathbb{C}^{(n)}_k(X,\mathcal{S})$ of the representation. For Ba{\l}aban's own terms at the sites treated by the statement on frozen terms of Ba{\l}aban's construction, it is implied by that statement. The level-$k$ collars close for the pair.
\item[(b)] Let $\eta$ be the complex parameter multiplying the shift piece in the correlation theorem, and let $\mathsf{w}_L$ be the radius in that parameter of the polydisc lemma of the correlation theorem, a radius at least $1/\theta'$, with $\theta'$ the constant fixed in that lemma; only an upper bound on $\mathsf{w}_L$ is used in what follows. For each cube $\Delta$ let $\sigma(\Delta)$ be the coordinate of that polydisc indexed by $\Delta$. On the polydisc
\begin{equation*}
\{|\eta|\le \mathsf{w}_L\}\times\{|\sigma(\Delta)|\le e^{\kappa_1}\}
\end{equation*}
the part of the polydisc lemma of the correlation theorem that concerns the pair is supplied by clause~(a) of Theorem~\ref{4.26:thm-inherited-expansion}, by Corollary~\ref{4.26:cor-inherited-room} and by Proposition~\ref{4.26:prop-source-image}. That lemma asks two things for the terms with $X\cap\mathfrak{T}\neq\emptyset$:
\begin{itemize}
\item that the set $\mathsf{R}_p$ of data at which the lemma evaluates the terms on the polydisc be contained in its domain $\mathfrak{D}_v$, and
\item a bound of each such term by its supremum over $\mathfrak{D}_v$.
\end{itemize}
Let $P$, $P'$, $P''$ be the three majorants of the contribution of the stage layer used in the proof of Theorem~\ref{4.26:thm-inherited-expansion}: $P$ on the shells at level $m$ of the large-field region $W$, and $P'$, $P''$ on $\Omega_k$. The two requirements are met provided that, at every $y\in\mathfrak{T}$, each of the quantities
\begin{equation}\label{4.26:eq-real-slot-scope-1}
\mathsf{w}_L\,P\,2^{-(k-1-m)}\ \ \text{(on the shells of $W$ at level $m$)},\qquad
\mathsf{w}_L\,P',\quad \mathsf{w}_L\,P''\ \ \text{(on $\Omega_k$)},
\end{equation}
with $P,P',P''$ evaluated at $|\sigma(\Delta)|\le e^{\kappa_1}$, is at most one quarter of the room by which the thresholds of the inherited domain exceed the thresholds $\theta_{\mathrm{src},a}$ of the source space $\mathfrak{S}_{\mathrm{src}}$. That room is the one given by Corollary~\ref{4.26:cor-inherited-room}.

The condition \eqref{4.26:eq-real-slot-scope-1} is a ceiling on $\gamma$. It is met for $\gamma$ small enough, uniformly in the number of stages of Theorem~\ref{4.26:thm-inherited-expansion} and in the depth $k-m$. It enters the correlation theorem as a hypothesis in supremum form, with $P,P',P''$ taken on $|\sigma(\Delta)|\le e^{\kappa_1}$ and at the halved rate at which that theorem states it. When hypotheses (w1) and (w2) of Corollary~\ref{4.26:cor-extra-parameter} hold, the same conclusions hold jointly in the further parameter $w\in\mathcal{W}(X)$ of Corollary~\ref{4.26:cor-extra-parameter}.
\item[(c)] Analyticity of $\mathfrak{E}$ in the slot at complex values of the source $\zeta$ is not a conclusion of these statements. It is the content of a separate statement, the slot-analyticity theorem at complex sources, which is an input of the correlation theorem, together with the disc condition that accompanies it there. Nothing proved in this section replaces that input.
\end{enumerate}
\end{remark}

\begin{proof}
(a) We first note what each statement is about. Proposition~\ref{4.26:prop-source-image} asserts membership in $\mathfrak{D}_{\mathfrak{E}}$, and holomorphy in the pair, at every real slot datum. Theorem~\ref{4.26:thm-inherited-expansion} and Corollary~\ref{4.26:cor-inherited-room} concern the terms $\mathbb{C}^{(n)}_k(X,\mathcal{S})$ on their spaces, with the slot datum real.

The level-$k$ collars close for the pair by the argument that closes them in the proof of Theorem~\ref{4.26:thm-inherited-expansion}, with the same inputs as there.

In Ba{\l}aban's construction \cite{B-CMP122b} the slot is kept a real parameter, and the gauge map is retained.

(b) The gauge presentation of the pair is as follows. Let $\mathfrak{b}$ be a later arrest, and let $h=h(\zeta)$ be the element of $G^c$ by which the holomorphic site map $\mathfrak{Q}^{\mathfrak{b}}(\zeta)$ of part~(c) below presents the carrier pair at the value $\zeta$ of the source. Proposition~\ref{4.26:prop-source-image} gives
\begin{equation*}
(\mathcal{U}^{\iota},\mathcal{J}^{\iota})\restriction X\in\mathfrak{D}_{\mathfrak{E}}.
\end{equation*}
Its proof keeps the gauge-field part and the cube gauges. Write $(\hat{\mathcal{U}}^{\mathfrak{b}},\hat{\mathcal{J}}^{\mathfrak{b}})$ for the carrier pair of the arrest $\mathfrak{b}$ before gauge presentation. For real $\zeta$ the element $h$ is $G$-valued, so that the presented slots are real, and the statement that places the gauge-presented carrier pair gives
\begin{equation*}
\bigl((\hat{\mathcal{U}}^{\mathfrak{b}})^h,\operatorname{Ad}_h\hat{\mathcal{J}}^{\mathfrak{b}}\bigr)\restriction X\in\mathfrak{D}_{\mathfrak{E}};
\end{equation*}
there the domain $\mathfrak{D}_{\mathfrak{E}}$ and the un-presented pair are those supplied by Proposition~\ref{4.26:prop-source-image}.

Now let $y\in\mathfrak{T}$. Dilating the shift piece by $\eta$, with $|\eta|\le\mathsf{w}_L$, moves a coordinate of the data in the inherited domain by at most $\mathsf{w}_L$ times the contribution of the stage layer to that coordinate, since the three bounds $P\,2^{-(k-1-m)}$, $P'$, $P''$ on that contribution listed in part~(b) of the remark and used in the proof of Theorem~\ref{4.26:thm-inherited-expansion} are linear in the layer. The displacement is therefore bounded by the quantities \eqref{4.26:eq-real-slot-scope-1}, which under the stated condition stay within one quarter of the room, positive by Corollary~\ref{4.26:cor-inherited-room}. Consequently, on the polydisc the moved data remain in the inherited domain, and the bound by the supremum follows from clause~(a) of Theorem~\ref{4.26:thm-inherited-expansion}.

We next show that the condition holds for small $\gamma$. Write $\mathsf{t}=\log g_k^{-2}$. The following bounds hold:
\begin{itemize}
\item $P$ is bounded by a constant times $e^{-3\delta_1R_k}$, where $\delta_1$ is the decay constant carried by the majorant $P$ in the proof of Theorem~\ref{4.26:thm-inherited-expansion};
\item $R_k\ge\mathsf{t}^{r}$;
\item $\mathsf{w}_L$, as chosen in the correlation theorem, is bounded by a constant times $\mathsf{t}^{\mathsf{m}}$.
\end{itemize}
The product $\mathsf{w}_L\,P$ therefore tends to zero as $\mathsf{t}$ grows. Since $\mathsf{t}$ is bounded below by the threshold $\mathsf{t}_\gamma=\log\gamma^{-2}$, which tends to infinity as $\gamma\to0$, the condition is met for $\gamma$ small enough.

On the shells the room carries the factor $2^{-(k-m)}$. The layer carries the factor
\begin{equation*}
2^{-(k-1-m)}=2\cdot2^{-(k-m)}.
\end{equation*}
The ratio of the two is thus independent of $k$ and $m$, which gives the uniformity in the depth.

Nothing in \eqref{4.26:eq-real-slot-scope-1} refers to the number of stages: $P$, $R_k$ and $\mathsf{w}_L$ enter only through the bounds just listed, which involve the level $k$ and not the stage, and the lower bounds on the room in Corollary~\ref{4.26:cor-inherited-room} hold for every stage $n$ with the same constants. This gives the uniformity in the number of stages.

The statement about the further parameter $w\in\mathcal{W}(X)$ is Corollary~\ref{4.26:cor-extra-parameter}.

(c) At complex $\zeta$ the correlation theorem evaluates $\mathfrak{E}$ at the pair and the live slots rotated by the holomorphic site map $\mathfrak{Q}^{\mathfrak{b}}(\zeta)$. Write $A^{[h]}$ for the rotated slot. The quantity used is
\begin{equation*}
F^{\mathfrak{b}}_{\mathfrak{E}}(\zeta;\Phi)
=\mathfrak{E}\bigl(X;(\hat{\mathcal{U}}^{\mathfrak{b}})^h,\operatorname{Ad}_h\hat{\mathcal{J}}^{\mathfrak{b}};A^{[h]}\bigr).
\end{equation*}
Here $h$ lies in the region of $G^c$ near $G$ on which $\mathfrak{E}$ is analytic, $|\log h|\le c_h$, where $c_h$ is the radius of that region in \cite{B-CMP109}.

Three facts combine at this point:
\begin{itemize}
\item the element $h$ is $G$-valued on the real slice only;
\item a term produced by the summation of boundary terms at finer zones is invariant only jointly, under the gauge map acting on pair and slot together;
\item the gauge map is retained in Ba{\l}aban's construction \cite{B-CMP122b}.
\end{itemize}
Hence at complex $\zeta$ the slot $A^{[h]}$ is complex.

None of the available analyticity statements covers a complex slot:
\begin{itemize}
\item the corollary bounding $\mathfrak{E}$ on the tube around the image of the source space $\mathfrak{S}_{\mathrm{src}}$, whose bound Proposition~\ref{4.26:prop-source-image} carries, is stated at real $\mathbf{A}$;
\item the definition of the flat analyticity class sorts the slot coordinates into weak, strong and silent ones (in the sense of that definition) and gives analyticity in the weak coordinates only, while strong and silent slots are real;
\item the two lemmas on the functional $\mathfrak{E}$ used for the stage estimates are stated for an $\mathfrak{E}$ with no slot argument, and so give no slot analyticity;
\item for terms of type $\mathbf{B}$, the analyticity statements of \cite[(3.29)]{B-CMP109}, \cite{B-CMP116} and \cite{B-CMP119} are used here with the slot held real.
\end{itemize}

The lemma on slot analyticity for the inherited expansion gives less than is needed. It gives analyticity, at slot value zero, in the slots that the stages of the inherited expansion themselves create. It also gives analyticity in a coordinate that is weak for every parent term, under the condition $(1+\varepsilon)\mathfrak{b}^{\sharp}<\varrho$ on the slot radius $\mathfrak{b}^{\sharp}$ and the polydisc radius $\varrho$ fixed in that lemma, with $\varepsilon>0$ as there. It gives nothing for strong or silent coordinates.

Therefore slot analyticity at complex $\zeta$ enters the correlation theorem only through the slot-analyticity theorem at complex sources.
\end{proof}

\begin{definition}[Conditions of the inherited-domain theorem]\label{4.26:def-inherited-conditions}
For a level $j$ write $x_j:=g_j^{-2}$ and
$\mathsf{t}_j:=\log x_j$, and put $\mathsf{t}_\gamma:=\log\gamma^{-2}$. The levels in play are the
levels $k$ at which the stages of this section are taken, together with, for each arrest
$\mathfrak{a}$ of the scope below, the levels of its range $I(\mathfrak{a})$. It is part of the
setting that every level $j$ in play satisfies $\mathsf{t}_j\ge\mathsf{t}_\gamma$, that is,
$g_j\le\gamma$.

\emph{Notation.} Let $\alpha_{0,j},\alpha_{1,j}$ be the radii of \cite[(2.28)]{B-CMP119}. Their
constants are written $C_0$ and $C_1$ in \cite[(2.28)]{B-CMP119}; the second is denoted $C_1'$
here, to distinguish it from the polydisc constant $C_1$ of
Definition~\ref{4.26:def-regularity-radius}. The powers of the logarithm in the two radii are
$q_0,q_1$. We put
$q:=\min\{q_0,q_1\}$. Under the condition $q_0=q_1$ imposed below, $q=q_0=q_1$, and the radii are
\begin{equation}\label{4.26:eq-inherited-conditions-1}
\begin{aligned}
\alpha_{0,j}&=C_0\,x_j^{-1/2}(\log x_j)^{q}=C_0\,e^{-\mathsf{t}_j/2}\,\mathsf{t}_j^{\,q},\\
\alpha_{1,j}&=C_1'\,x_j^{-1/2}(\log x_j)^{q}=C_1'\,e^{-\mathsf{t}_j/2}\,\mathsf{t}_j^{\,q}.
\end{aligned}
\end{equation}

We put $\beta_0:=1/7$ and $c_\beta:=2\log(1+\beta_0)$. Let $L_{\mathrm{w}}$ be the base of the
weight $w_\ast$ of Lemma~\ref{4.26:lem-summation}, so that $w_\ast=L_{\mathrm{w}}^{2N_0}$. Further,
\begin{gather*}
\mathsf{a}:=L^3(1+2\beta)(1+\beta_0)^4,\\
\bar{\alpha}_0:=\sup_{\mathsf{t}\ge\mathsf{t}_\gamma}C_0\,e^{-\mathsf{t}/2}\,\mathsf{t}^{\,q},
\qquad
\bar{\alpha}_1:=\sup_{\mathsf{t}\ge\mathsf{t}_\gamma}C_1'\,e^{-\mathsf{t}/2}\,\mathsf{t}^{\,q}.
\end{gather*}
Let $B_3=B_3(d,L)$ be the constant of the input sizes in part (c) of
Lemma~\ref{4.26:lem-near-box-geometry}; $O(1)$ denotes, at each occurrence, the constant so
written at the corresponding place of that part of the lemma or of the chain
\eqref{4.26:eq-near-pair-chain-a}. We set
\begin{gather*}
\mathsf{S}:=B_3^2\,O(1)\,M\,\mathsf{a},\qquad
\mathsf{S}':=\mathsf{S}+\frac{C_1'}{3C_0}+B_3\,O(1),\\
\bar{G}:=\exp\bigl[(1+2B_3^{-1})\,\mathsf{S}\,\bar{\alpha}_0+O(1)\,M\,\mathsf{a}\,\bar{\alpha}_1\bigr],
\qquad
\bar{\varepsilon}:=\mathsf{S}'\,\bar{\alpha}_0,\qquad
\bar{G}':=\Bigl(\frac{1+c_{\mathrm{aux}}\bar{\varepsilon}}{1-c_{\mathrm{aux}}\bar{\varepsilon}}\Bigr)^{2}.
\end{gather*}
Here $c_{\mathrm{aux}}>0$ is the constant of the $\kappa$-bound.
Let $B$ and $c_5$ be the constants of the two-sided bound on the increments of the inverse
couplings along the coupling flow, clause (0) of Theorem~0 (Theorem~\ref{4.1:thm-clause-zero}).
We put $B_{\pm}:=\max\{B,c_5\}$; by that bound,
$|x_j-x_{j+1}|\le B_{\pm}$ for adjacent levels. Let $R_{\min}$ be the infimum of $R_k$ over the
levels in play, and let $\Psi$ be the function of $\mathsf{t}$ in the run-free majorant
$\bar{\gamma}_0(\mathsf{t})$ of the regularity radius, \eqref{4.26:eq-regularity-radius-a}.

The \emph{conditions of the inherited-domain theorem} are listed below. They fall into floors,
constants fixed before $\gamma$, and ceilings on $\gamma$, and they are imposed in the scope
described at the end. Every inequality among them is one-sided; the equalities fix constants.
The floors and the constants are fixed before $\gamma$, in an order compatible with the order of
choice of the auxiliary parameters, and $\gamma$ is chosen last.

\emph{The five displayed conditions.}
\begin{equation}\label{4.26:eq-inherited-conditions-a}
\begin{gathered}
\text{(i)}\quad q=\min\{q_0,q_1\}\ge p_0+7r+1,\\
\text{(ii)}\quad (1+B_{\pm}\gamma^2)^{q+1/2}\le 1+\beta_0,\\
\text{(iii)}\quad R_{\min}\ge 5L,\\
\text{(iv)}\quad \mathsf{t}_\gamma\ge 2(r+q)+c_\beta
\ \text{ and }\ 600(1+\beta_0)^2\,\Psi(\mathsf{t}_\gamma)\le 1,\\
\text{(v)}\quad
\begin{aligned}[t]
&(\bar{G}\bar{G}'-1)(1+\beta)+(\bar{G}'-1)\\
&\quad+\bar{G}'\bigl[16\mathsf{S}^2e^{4\mathsf{S}\bar{\alpha}_0}+2B_3^{-1}\mathsf{S}^2
+8\mathsf{S}'^2e^{4\bar{\varepsilon}}\bigr]\bar{\alpha}_0\le\frac{\beta}{12}.
\end{aligned}
\end{gathered}
\end{equation}
In (i) and (iv), $r$ is the exponent for which the envelope $\check{R}_k$ of $R_k$ is of order
$\mathsf{t}_k^{\,r}$, and $p_0$ is the fixed exponent of the degree function $p_0(\cdot)$.
The second clause of (iv) is stated at $\mathsf{t}_\gamma$. By
\eqref{4.26:eq-regularity-radius-a}, $\Psi(\mathsf{t})$ is a constant multiple of
$\mathsf{t}^{r+q}e^{-\mathsf{t}/2}$, whose derivative
\begin{equation*}
\frac{d}{d\mathsf{t}}\bigl(\mathsf{t}^{r+q}e^{-\mathsf{t}/2}\bigr)
=\mathsf{t}^{r+q-1}e^{-\mathsf{t}/2}\bigl(r+q-\tfrac12\mathsf{t}\bigr)
\end{equation*}
is $\le 0$ for $\mathsf{t}\ge 2(r+q)$. Since $c_\beta\ge 0$, the first clause of (iv) gives
$\mathsf{t}_\gamma\ge 2(r+q)$, hence $\Psi(\mathsf{t}_\gamma)=\sup_{\mathsf{t}\ge\mathsf{t}_\gamma}\Psi(\mathsf{t})$,
and the second clause holds at every $\mathsf{t}\ge\mathsf{t}_\gamma$. Through the
majorant \eqref{4.26:eq-regularity-radius-a} it contains $120\,\gamma_0^{(k)}\le 1$ at every level $k$.

\emph{Floors.}
\begin{itemize}
\item The floors on $\kappa$:
\begin{equation*}
\beta\kappa\ge c_d+2\sqrt{d}\,(2+4\beta),\qquad
\frac{M}{M_1}\,\min\{\delta_0,\delta\}\ge 8\kappa\sqrt{d}.
\end{equation*}
Together with these go the restrictions of \cite{B-CMP116}, taken at $\varpi:=\beta/40$, and
$p_0\ge 5r$.
\item The floor on $\kappa_1$:
\begin{equation*}
\kappa_1\ge\max\{16\kappa+1,\;4L\kappa+1\}.
\end{equation*}
Its second member is
equivalent to $\tfrac12(\kappa_1-1)\ge 2L\kappa$, the restriction assumed in \cite{B-CMP116}.
\item $L\ge L_0$, the floor of the quantifier prefix of Theorem~0
(Notation~\ref{2.1:not-prefix}), under which clause (i) of Theorem~0 and the correlation theorem
are stated.
\item $L_{\mathrm{w}}^2\le L/3$ and $\beta=1/5$.
\item The four floors imposed with the stage spaces of this section: a floor on $\kappa_1$ of
their own, distinct from the floor displayed above; the
level floor $\delta w\ge 16\log L+4$ of Definition~\ref{4.26:def-regularity-radius}; the flat
floor of the stage spaces; and the floor on $M_1$.
\item $M\ge 3L$.
\item The flat floor on $\kappa$ of the cut-shell device, together with its companion floor, the
latter imposed at the cores of that device only.
\item Condition (i) of \eqref{4.26:eq-inherited-conditions-a}, and $q_0=q_1$.
\item The ratio condition \eqref{4.26:eq-near-pairs-outside-a},
\begin{equation*}
\frac{C_1'}{C_0}\ge\frac{2\,\mathsf{a}\,B_3^2\,O(1)\,M}{1-\beta}.
\end{equation*}
The ratio $C_1'/C_0$ is chosen after $L$, $M$, $B_3=B_3(d,L)$ and $\mathsf{a}=\mathsf{a}(L)$,
and before $\gamma$. Condition \eqref{4.26:eq-near-pairs-outside-a} is thus a one-sided
condition on that choice. It places no condition on $L$ or $M$ and caps no parameter.

At the input sizes of part (c) of Lemma~\ref{4.26:lem-near-box-geometry}, it is the
restriction stated in \cite[p.~277]{B-CMP109} for \cite[Lemma~4]{B-CMP109}. The floors on
$C_0$ and on $C_1'$ imposed with the polydisc of constant $C_1$ (the polydisc of
Definition~\ref{4.26:def-regularity-radius}) are proportional to $C_1$ and are lower bounds. None of
them bounds $C_1'/C_0$ from above. They are therefore compatible with
\eqref{4.26:eq-near-pairs-outside-a} and give no cap on $M$.
\end{itemize}
The floor $\delta_1M/M_1\ge 2$ enters through the scope below.

\emph{Constants fixed before $\gamma$.}
\begin{itemize}
\item The absolute constant $\beta_0=1/7$ fixed above.
\item $\theta_S$ is chosen in $[\tfrac12,\tfrac89]$; hence $\mathsf{c}_e:=\theta_S$ and
$\hat{c}:=\min\{\tfrac12,\mathsf{c}_e\}=\tfrac12$.
\item $\mathsf{w}_0=1/24$ and $\theta=\tfrac12$.
\item The constants $K_{\mathrm{cv}}$, $K^{\mathrm{nr}}$, $K^{\mathrm{rem}}$, $K_{\min}$,
$K_{\flat}'$ and the near constants of $C_R^{\natural}$, as in
Definition~\ref{4.26:def-radii-window}.
\item The counting constant $C_V$.
\end{itemize}

\emph{Ceilings on $\gamma$.} In each of them the quantities $R_k$, $N$, $N_0$,
$\mathsf{l}_k$, $\alpha_{i,k}$ of Ba{\l}aban's construction are functions of the one variable
$\mathsf{t}=\mathsf{t}_k$. Each ceiling is required at every $\mathsf{t}\ge\mathsf{t}_\gamma$,
with all its quantities evaluated at one and the same $\mathsf{t}$. In the ceilings $N$ denotes
the envelope $\check{R}_k$ of $R_k$ of the scope below.
\begin{itemize}
\item The window condition $K_{\min}\vartheta\le\mathsf{w}_0$, \eqref{4.26:eq-radii-window-b}.
\item The smallness condition in summed form, \eqref{4.26:eq-summed-smallness-a}.
\item Conditions (v) and (ii) of \eqref{4.26:eq-inherited-conditions-a}.
\item Together with (v), the one-sided smallness conditions
\begin{equation*}
c_{\mathrm{aux}}\bar{\varepsilon}<1,\qquad c_N\bar{\varepsilon}\le\tfrac{1}{25},
\end{equation*}
where $c_N>0$ is a fixed constant, and the smallness of $\alpha_0,\alpha_1$ required for
\cite[(3.24)--(3.27)]{B-CMP109} and the smallness of $\alpha_0'$ required in \cite{B-CMP109},
each read at parameters proportional to $\bar{\varepsilon}$. The quantity bounded in each of
these conditions is non-decreasing in $\bar{\alpha}_0$ and tends to zero with it.
\item With $N=\check{R}_k$, the envelope of $R_k$, and with $\mathsf{l}$ and $R$ read at the
same $\mathsf{t}$, as for every ceiling:
\begin{equation*}
N\,(2L)^d\,\mathsf{l}_k^2\,e^{-4a_\mu R_k/L}\le 1.
\end{equation*}
\item Condition (iii) of \eqref{4.26:eq-inherited-conditions-a}. It contains $R_k\ge L$, used
in the summation lemma, Lemma~\ref{4.26:lem-summation}, and $R_k\ge 2$.
\item $(N_0+1)(N_0+3)\le N$. This governs the degrees in Lemma~\ref{4.26:lem-summation}.
\end{itemize}

\emph{Further conditions.} The following conditions are also imposed; those that
involve a level are read at every $\mathsf{t}\ge\mathsf{t}_\gamma$, as above.
\begin{itemize}
\item $r_A\ge 2$ and $r_A>C_{\mathrm{far}}$.
\item The floor on $\kappa_A$:
\begin{equation*}
\kappa_A(k)\ge\max\bigl\{c_\flat+c_J+1,\ \log 8\max\{c_T^{\mathrm{sz}},C_0^{\sharp}\}\bigr\}.
\end{equation*}
\item The three further conditions of the slot construction on the slot sets
$\mathfrak{X}^{\sharp}_j$ and the slot radii $r_A$, $\kappa_A$.
\item $N\,e^{-\sigma_\ast R_\ast/L}\le 1$.
\item $K_\Delta\gamma_0\le\tfrac12$ and the two further ceilings of the stage spaces that
involve $\gamma_0$. In all three, $\gamma_0=\gamma_0^{(k)}$ at every level $k$, and each is
imposed at the majorant $\bar{\gamma}_0(\mathsf{t}_\gamma)$.
\item Condition (iv) of \eqref{4.26:eq-inherited-conditions-a}.
\item $2.2\,B_\sharp C_1\delta_j'\le 1/10$.
\item $e^{-\delta wR_k/8}\le(4^dC_dR_k)^{-1}$.
\end{itemize}

\emph{Scope.} The conditions are imposed in the following setting.
\begin{itemize}
\item The geometric setting of the stages, with its nesting property.
\item Large-field regions $Z$ of the whole class. Components excluded by an arrest are treated
under the hypotheses of the arrest theorem taken in full:
  \begin{itemize}
  \item its five standing conditions, in force from each arrest on;
  \item the conditions $C_0\ge 2.5A_0/\varphi_{\mathrm{arrest}}$, where
  $\varphi_{\mathrm{arrest}}$ is the lower bound on the arrested source factor of that theorem,
  and $q_0>p_0$, and $\gamma$ below
  a threshold depending on $(d,L,M)$;
  \item the conditions $q_1=q_0$ and $\beta=1/5$; its count, there $N=R_k$ (the envelope
  $\check{R}_k$ of $R_k$ is among the items of this scope, below); $\delta_1M/M_1\ge 2$; its floor
  on $L$;
  and $\gamma$ below its threshold; together with three further conditions of its list, one of
  them a floor on $q$ and one taken at the parameter $\theta_S$ fixed above;
  \item the two further hypotheses of that theorem, namely two statements fixed with the stage
  spaces of this section, each taken together with its own hypotheses: for the first,
  $\beta<\tfrac14$, $L_{\mathrm{w}}^2\le L/3$, $\delta M/M_1\ge 3$ and a floor on $q$; for the
  second, $M\ge M_\ast(L)$, $\gamma\le\gamma_\ast(L)$, its third condition, and its first
  condition taken only on the data for which that statement states it; and $\delta w\ge 8$.
  \end{itemize}
\item The range of levels $I(\mathfrak{a})=\{h_a+1,\dots,k_a+1\}$ of an arrest
$\mathfrak{a}$, with the lemma on the range of an arrest in its extended form on that range.
\item The conclusion of the arrest theorem, taken on its source space
$\tilde{U}^{\mathrm{arrest}}_k$.
\item The frozen terms $\mathfrak{F}(\mathfrak{a})$ of each arrest; their members of the fourth
type are supplied by the frozen-term statement for that type.
\item Older live pieces, under the conditions fixed for them with the stage spaces of this
section.
\item The envelope $\check{R}_k$ of $R_k$.
\item The flat-term convention fixed with the stage spaces of this section.
\end{itemize}
The inequality \eqref{4.26:eq-source-image-a} is not among these conditions. It is a
conclusion, established where the image of the source space is placed in the inherited domain.

\emph{Facts about the conditions,} proved below:
\begin{itemize}
\item[(a)] Under (i), (ii) and the first clause of (iv), for $i=0,1$ and every pair of adjacent
levels $j,j+1$ in play (compare Lemma~\ref{4.4:lem-radii-ratio}),
\begin{equation}\label{4.26:eq-inherited-conditions-2}
\frac{\alpha_{i,j+1}}{\alpha_{i,j}}\le 1+\beta_0,\qquad
\frac{\alpha_{i,j}}{\alpha_{i,j+1}}\le 1+\beta_0 .
\end{equation}
\item[(b)] The first clause of (iv) is equivalent to
$\gamma\le\tfrac78\,e^{-(r+q)}$.
\item[(c)] The left side of (v) depends only on $(L,M,B_3,C_1'/C_0,q)$ and on
$(\bar{\alpha}_0,\bar{\alpha}_1)$. No cube, level, $N_0$ or $w_\ast$ enters it. Each of its
members is non-decreasing in $(\bar{\alpha}_0,\bar{\alpha}_1)$ as long as
$c_{\mathrm{aux}}\bar{\varepsilon}<1$,
and tends to zero with $(\bar{\alpha}_0,\bar{\alpha}_1)$. Hence (v) is one inequality on
$\gamma$, satisfied by every sufficiently small $\gamma$. The same holds for (ii).
\item[(d)] Under (ii) and (iii), condition (v), together with $q_0=q_1$, implies, through
part (c) of Lemma~\ref{4.26:lem-near-box-geometry}, the local condition
\eqref{4.26:eq-near-pair-chain-b} at every near pair.
\end{itemize}
\end{definition}

\begin{proof}[Proof of the facts (a)--(d)]
\emph{(a).} Let $j,j+1$ be adjacent levels in play. Let $x$ be the smaller and $x'$ the larger
of $x_j,x_{j+1}$. Since $\mathsf{t}_j,\mathsf{t}_{j+1}\ge\mathsf{t}_\gamma$, we have
$x\ge\gamma^{-2}$. The bound of Theorem~\ref{4.1:thm-clause-zero} gives $0\le x'-x\le B_{\pm}$. Hence
\begin{equation*}
\frac{x'}{x}\le 1+\frac{B_{\pm}}{x}\le 1+B_{\pm}\gamma^2 .
\end{equation*}
By (i), $q\ge 1$, since $p_0\ge 0$ and $r\ge 0$. By the first clause of (iv),
$\log x\ge\mathsf{t}_\gamma\ge 2(r+q)\ge 1$. Since $\log(1+u)\le u$ for $u\ge 0$,
\begin{equation*}
\frac{\log x'}{\log x}
=1+\frac{\log(x'/x)}{\log x}
\le 1+\frac{B_{\pm}\gamma^2}{\log x}
\le 1+B_{\pm}\gamma^2 .
\end{equation*}
Insert these two bounds into \eqref{4.26:eq-inherited-conditions-1}. Let $C$ be $C_0$ for $i=0$
and $C_1'$ for $i=1$. For $i=0,1$ this gives
\begin{align*}
\frac{Cx^{-1/2}(\log x)^q}{Cx'^{-1/2}(\log x')^q}
&=\Bigl(\frac{x'}{x}\Bigr)^{1/2}\Bigl(\frac{\log x}{\log x'}\Bigr)^{q}
\le(1+B_{\pm}\gamma^2)^{1/2},\\
\frac{Cx'^{-1/2}(\log x')^q}{Cx^{-1/2}(\log x)^q}
&=\Bigl(\frac{x}{x'}\Bigr)^{1/2}\Bigl(\frac{\log x'}{\log x}\Bigr)^{q}
\le(1+B_{\pm}\gamma^2)^{q}.
\end{align*}
Here we used $x/x'\le 1$ and $\log x/\log x'\le 1$. Both right sides are at most
$(1+B_{\pm}\gamma^2)^{q+1/2}$, which is at most $1+\beta_0$ by (ii). Both directions therefore
follow from the same two bounds, and \eqref{4.26:eq-inherited-conditions-2} holds.

Condition (ii) involves no object of the run, only $B_{\pm}$, $q$ and $\beta_0$. At adjacent
levels, \eqref{4.26:eq-inherited-conditions-2} implies the bound
$(1+\beta_0)^2(n-m)^{\beta_0}$ for the ratio of radii at levels $m<n$ stated in
\cite[(2.8)]{B-CMP119}. In \cite{B-CMP119}, $\beta_0>0$ may
be chosen arbitrarily small when $g$ is sufficiently small. Here $\beta_0=1/7$ is fixed, and
the smallness is carried by the ceiling (ii) on $\gamma$. In this section only ratios at
adjacent levels, and products of such ratios with powers of $L$, are used.

\emph{(b).} By definition $\mathsf{t}_\gamma=-2\log\gamma$. Therefore
$\mathsf{t}_\gamma\ge 2(r+q)+c_\beta$ holds if and only if
\begin{equation*}
\log\gamma\le-(r+q)-\tfrac12c_\beta=-(r+q)-\log(1+\beta_0).
\end{equation*}
This is the condition $\gamma\le e^{-(r+q)}/(1+\beta_0)$. Since $\beta_0=1/7$, we have
$1/(1+\beta_0)=7/8$. The condition is one-sided.

\emph{(c).} The quantities $\mathsf{a}$, $\mathsf{S}$ and $\mathsf{S}'$ depend only on $L$,
$M$, $B_3$, $C_1'/C_0$ and the fixed constants $\beta,\beta_0,O(1)$. Then $\bar{G}$,
$\bar{\varepsilon}$ and $\bar{G}'$ depend on these and on $(\bar{\alpha}_0,\bar{\alpha}_1)$
alone, and $\bar{\alpha}_0$ depends on $C_0$, $q$ and $\mathsf{t}_\gamma$, and $\bar{\alpha}_1$
on $C_1'$, $q$ and $\mathsf{t}_\gamma$.

Let $\bar{\alpha}_0,\bar{\alpha}_1\ge 0$ with $c_{\mathrm{aux}}\bar{\varepsilon}<1$. Then:
\begin{itemize}
\item $\bar{G}\ge 1$ is non-decreasing in $(\bar{\alpha}_0,\bar{\alpha}_1)$, being the
exponential of a combination of them with non-negative coefficients;
\item $\bar{\varepsilon}=\mathsf{S}'\bar{\alpha}_0$ is non-decreasing;
\item $\bar{G}'\ge 1$ is non-decreasing, since $u\mapsto(1+u)/(1-u)$ increases on $[0,1)$ and
$u=c_{\mathrm{aux}}\bar{\varepsilon}$ lies in $[0,1)$;
\item the bracket multiplying $\bar{\alpha}_0$ is non-decreasing.
\end{itemize}
Hence $(\bar{G}\bar{G}'-1)(1+\beta)$, $(\bar{G}'-1)$ and the third member of (v) are
non-decreasing. As $(\bar{\alpha}_0,\bar{\alpha}_1)\to 0$ we have $\bar{G}\to1$,
$\bar{\varepsilon}\to 0$ and $\bar{G}'\to 1$, so each member tends to zero.

Next, $C_0e^{-\mathsf{t}/2}\mathsf{t}^q$ and $C_1'e^{-\mathsf{t}/2}\mathsf{t}^q$ tend to $0$ as
$\mathsf{t}\to\infty$. Hence $\bar{\alpha}_i$
tends to zero as $\mathsf{t}_\gamma\to\infty$, that is, as $\gamma\to 0$. Moreover
$\bar{\alpha}_i$ is non-decreasing in $\gamma$, being a supremum over a set that grows with
$\gamma$. So the left side of (v) is non-decreasing in $\gamma$ and tends to zero with it.
Condition (v) therefore holds for all sufficiently small $\gamma$, the threshold depending on
$(L,M,B_3,C_0,C_1',q)$ only.

For (ii): the left side $(1+B_{\pm}\gamma^2)^{q+1/2}$ increases in $\gamma$ and tends to
$1<1+\beta_0$ as $\gamma\to 0$.

\emph{(d).} Fix a near pair, and let $m$ be the level at which its local condition
\eqref{4.26:eq-near-pair-chain-b} is formed. Since $q_0=q_1$,
\eqref{4.26:eq-inherited-conditions-1} gives $\alpha_{1,m}=(C_1'/C_0)\,\alpha_{0,m}$. Thus both
radii at level $m$ are multiples of the single number $\alpha_{0,m}$. Moreover
$\mathsf{t}_m\ge\mathsf{t}_\gamma$, so by \eqref{4.26:eq-inherited-conditions-1} and the
definition of $\bar{\alpha}_0$, $\bar{\alpha}_1$,
\begin{equation*}
\alpha_{0,m}=C_0e^{-\mathsf{t}_m/2}\mathsf{t}_m^{\,q}\le\bar{\alpha}_0,\qquad
\alpha_{1,m}=C_1'e^{-\mathsf{t}_m/2}\mathsf{t}_m^{\,q}\le\bar{\alpha}_1.
\end{equation*}

The hypotheses of Lemma~\ref{4.26:lem-near-box-geometry}, namely $R_{\min}\ge 5L$ and the bound
on ratios of radii at adjacent levels, are conditions (iii) and (ii). Part (c) of that lemma, in
the form \eqref{4.26:eq-near-box-geometry-b}, therefore bounds the input size:
\begin{equation*}
B_3^2\,O(1)\,M\,\mathfrak{a}_0\le\mathsf{S}\,\alpha_{0,m}.
\end{equation*}
Insert this bound and $\alpha_{1,m}=(C_1'/C_0)\alpha_{0,m}$ into the definitions of
$\mathfrak{q}_1$, $\mathfrak{q}_2$, $\varepsilon_{\square}$ and $G_{\square}$ at
\eqref{4.26:eq-near-pair-chain-a}. This gives
\begin{align*}
\mathfrak{q}_1&\le 8\mathsf{S}^2e^{4\mathsf{S}\alpha_{0,m}}\,\alpha_{0,m},\\
\mathfrak{q}_2&=B_3^{-1}\,\frac{\bigl(B_3^2\,O(1)\,M\,\mathfrak{a}_0\bigr)^2}{\alpha_{0,m}}
\le B_3^{-1}\,\frac{\mathsf{S}^2\alpha_{0,m}^2}{\alpha_{0,m}}=B_3^{-1}\mathsf{S}^2\,\alpha_{0,m},\\
\varepsilon_{\square}&\le\mathsf{S}'\,\alpha_{0,m},\\
\log G_{\square}&\le(1+B_3^{-1})\,\mathsf{S}\,\alpha_{0,m}+O(1)\,M\,\mathsf{a}\,\alpha_{1,m}.
\end{align*}
Each right side is linear in the radii $\alpha_{0,m},\alpha_{1,m}$ at the one level $m$. The
only exception is the factor $e^{4\mathsf{S}\alpha_{0,m}}$, which is non-decreasing in
$\alpha_{0,m}$. No quotient of radii at two different levels occurs.

Replace $\alpha_{i,m}$ by $\bar{\alpha}_i\ge\alpha_{i,m}$ and use $1+B_3^{-1}\le 1+2B_3^{-1}$.
Then
\begin{gather*}
\mathfrak{q}_1\le 8\mathsf{S}^2e^{4\mathsf{S}\bar{\alpha}_0}\bar{\alpha}_0,\qquad
\mathfrak{q}_2\le B_3^{-1}\mathsf{S}^2\bar{\alpha}_0,\\
\varepsilon_{\square}\le\bar{\varepsilon},\qquad
G_{\square}\le\bar{G}.
\end{gather*}
These are the quantities through which the local condition \eqref{4.26:eq-near-pair-chain-b}
is formed at the near pair. The left side of \eqref{4.26:eq-near-pair-chain-b} is
non-decreasing in $\mathfrak{q}_1$, $\mathfrak{q}_2$, $\varepsilon_{\square}$ and $G_{\square}$,
the quantity $G'_{\square}$ of \eqref{4.26:eq-near-pair-chain-a} being non-decreasing in
$\varepsilon_{\square}<1$. Replacing these quantities by the bounds just obtained bounds the left
side of \eqref{4.26:eq-near-pair-chain-b} by the left side of (v). Hence (v) implies
\eqref{4.26:eq-near-pair-chain-b} at every near pair. No cube, level or count of the run enters
(v) itself.
\end{proof}

\begin{definition}[The smallness condition in summed form]\label{4.26:def-summed-smallness}
We use the letters of Notation~\ref{4.26:not-slot-notation}, in particular the constant
$K_{\flat}'$ fixed there, $\mathsf{r}$, $V$, $K_2$, $K_Q$, $r_A$, $\varepsilon_S$, $c_{\flat}$ and
$C_0^{\sharp}=3c_G+1$. We also use those of the summation lemma, near groups included
(Lemma~\ref{4.26:lem-summation}), in particular the tile number $\mathsf{l}_j^2$. We put
\begin{equation*}
s^{\natural}:=\frac{\vartheta}{\mathsf{w}_0}+\frac{K_Q}{r_A}+\varepsilon_S .
\end{equation*}
The letters $N$, $N_0$, $E_0$, $B_0$, $\check{R}$, $L_\circ$ and $c_L$ denote the numbers
fixed in Lemma~\ref{4.26:lem-summation}, where they enter the members of
\eqref{4.26:eq-summation-a} and \eqref{4.26:eq-summation-b}. The
\emph{smallness condition in summed form} is the pair of inequalities
\begin{equation}\label{4.26:eq-summed-smallness-a}
\begin{gathered}
K_{\flat}'\,\mathsf{r}\,\Bigl[NV+4^d(100L\check{R})^d(N_0+1)(N_0+3)
+(2^d+L_\circ^6)(100L\check{R})^d\\
+\bigl(N_0+3+2\cdot 52^d+c_L\bigr)\mathsf{l}_j^2\Bigr]
\bigl(C_0^{\sharp}+E_0+B_0+2\bigr)\,s^{\natural}\le\frac14,\\
6NV\Bigl(\frac{K_2\vartheta}{\mathsf{w}_0}+\frac{K_Q}{r_A}\Bigr)\le 1 .
\end{gathered}
\end{equation}
Under \eqref{4.26:eq-summed-smallness-a}, the first member of \eqref{4.26:eq-summation-a} is at
most $1/16$, and that of \eqref{4.26:eq-summation-b} is at most $1/4$. Both bounds hold term by
term, and the members of the near groups are included.

The condition also enters in the following one-variable form. Here $\mathsf{t}=\log g_k^{-2}$,
the threshold is $\mathsf{t}_\gamma=\log\gamma^{-2}$, and $\mathsf{l}^2$ is the size parameter with
$L^{4N_0}\le\mathsf{l}^2$; the exponents $p_0$ and $q$ are among Ba{\l}aban's auxiliary
parameters $\mathfrak{p}$, $p_0$ being the exponent of the degree function $p_0(\cdot)$; and
$\bar{V}$, $C_{\mathrm{far}}$, $\varphi$, $\mathsf{t}_\ast$ and the exponent $p_1$ are parameters
of this form, on which the condition below is imposed. Further, $K'$ is a
constant that majorises $K_{\flat}'$ and the middle summands of the bracket in
\eqref{4.26:eq-summed-smallness-a}, and we put $\mathsf{m}:=q-p_1$ and
$\mathsf{m}_0:=q-p_0$. The form reads
\begin{equation}\label{4.26:eq-summed-smallness-1}
\begin{gathered}
\sup_{\mathsf{t}\ge\mathsf{t}_\gamma}K'\bigl[N\mathsf{r}\bar{V}+(N_0+3)\mathsf{l}^2\bigr]
\Bigl[\mathsf{w}_0^{-1}\mathsf{t}^{-\mathsf{m}}+C_{\mathrm{far}}\,\mathsf{t}^{-\mathsf{m}_0}
+e^{-\varphi(\mathsf{t}-1)+c_{\flat}}\Bigr]\le 1,\\
\mathsf{t}_\gamma-1\ge\mathsf{t}_\ast .
\end{gathered}
\end{equation}
\end{definition}

\begin{proof}[Proof of the bounds $1/16$ and $1/4$]
Put $\mathsf{C}:=C_0^{\sharp}+E_0+B_0+2$. This is the last factor but one on the left of the
first inequality in \eqref{4.26:eq-summed-smallness-a}.

The bounds on the members of \eqref{4.26:eq-summation-a} and \eqref{4.26:eq-summation-b} are
checked term by term. The comparisons are the following inequalities; each bounds a
coefficient occurring in those members by $\mathsf{C}$ times a summand of the bracket in the
first inequality of \eqref{4.26:eq-summed-smallness-a}:
\begin{equation}\label{4.26:eq-summed-smallness-2}
\begin{aligned}
V\bigl(NC_0^{\sharp}+E_0+B_0+1\bigr)&\le NV\mathsf{C},\\
(N_0+1)\bigl[(N_0+3)C_0^{\sharp}+E_0+B_0+1\bigr]&\le(N_0+1)(N_0+3)\mathsf{C},\\
E_0+1&\le\mathsf{C},\\
\mathsf{l}_j^2\bigl(E_0+B_0+1+(N_0+3)C_0^{\sharp}\bigr)&\le(N_0+3)\mathsf{l}_j^2\mathsf{C},\\
NC_0^{\sharp}&\le NV\mathsf{C}.
\end{aligned}
\end{equation}
The third inequality is used on both near members.

Each of these inequalities follows from elementary arithmetic, as follows. Here $N\ge1$,
$V\ge1$ and $N_0\ge0$, and the constants $C_0^{\sharp}$, $E_0$ and $B_0$ are non-negative, as
fixed in Notation~\ref{4.26:not-slot-notation} and Lemma~\ref{4.26:lem-summation}.
\begin{itemize}
\item In the first, the term $NC_0^{\sharp}$ occurs on both sides. Moreover, since $N\ge1$,
\begin{equation*}
E_0+B_0+1\le N(E_0+B_0+2).
\end{equation*}
Adding $NC_0^{\sharp}$ to both sides and multiplying by $V$ gives the first inequality.
\item In the second and in the fourth, $N_0+3\ge1$. Hence the quantity
$(N_0+3)C_0^{\sharp}+E_0+B_0+1$ is at most $(N_0+3)\mathsf{C}$. Multiplying by $N_0+1$, or by
$\mathsf{l}_j^2$, gives the two inequalities.
\item The third holds because $C_0^{\sharp}+B_0+1\ge0$.
\item The fifth holds because $V\ge1$ and $C_0^{\sharp}\le\mathsf{C}$.
\end{itemize}

The last summand of the bracket in \eqref{4.26:eq-summed-smallness-a} carries the tile number
$\mathsf{l}_j^2$ of Lemma~\ref{4.26:lem-summation}, and not the size parameter $\mathsf{l}^2$ of
\eqref{4.26:eq-summed-smallness-1}; where $\mathsf{l}_j^2=16\,\mathsf{l}^2$, the two differ by
the factor $16$.

With these comparisons, every group of the first member of \eqref{4.26:eq-summation-a} and of
the member of \eqref{4.26:eq-summation-b}, near groups included, is bounded by the
corresponding summand of the left side of \eqref{4.26:eq-summed-smallness-a}. The two
inequalities of \eqref{4.26:eq-summed-smallness-a} then give the bound $1/16$ for the first
member of \eqref{4.26:eq-summation-a}, and the bound $1/4$ for that of
\eqref{4.26:eq-summation-b}.

The groups of the overlap window of Lemma~\ref{4.26:lem-summation} are added, never multiplied,
as that lemma does with \eqref{4.26:eq-summation-c}. The termwise comparison therefore applies
to each of them separately.
\end{proof}

\begin{lemma}[Summed smallness from the earlier condition at a constant factor]
\label{4.26:lem-smallness-reduction}
The spend $\mathsf{w}_0=1/24$, the radii and $\theta=\tfrac12$ are those of
Definition~\ref{4.26:def-radii-window}, and the counts $N$, $\mathsf{r}$, $V$, $\bar V$, $N_0$,
$\mathsf{l}_j^2$, $\mathsf{l}^2$, the generic prefactor $K'$, the constants $K_\flat'$, $K_Q$,
$\varepsilon_S$, $C_{\mathrm{far}}$, $\varphi$, $c_\flat$, the threshold $\mathsf{t}_\ast$ and the
radius $r_A$ are those of the
smallness condition in summed form, Definition~\ref{4.26:def-summed-smallness}. Write
$\mathsf{t}=\log g_k^{-2}$, $\mathsf{t}_\gamma=\log\gamma^{-2}$, $\mathsf{m}=q-p_1$ and
$\mathsf{m}_0=q-p_0$.
\begin{itemize}
\item[(a)] Compared with the smallness condition of the analysis of the inherited-domain theorem,
whose conditions are collected in \eqref{4.26:eq-inherited-conditions-a} (that smallness condition
is written out in the form \eqref{4.26:eq-smallness-reduction-2} in (c) below), the condition
\eqref{4.26:eq-summed-smallness-a} differs in exactly one member of the smallness bracket: its
first member, $\vartheta$, is divided by $\mathsf{w}_0$. In addition the constant prefactor
contains further numbers depending on $L$ only, all inside the generic constant $K'$: the
constant $K_\flat$ of that smallness condition is replaced by $K_\flat'$; the constants
$2\cdot 52^d+c_L$ and $L_\circ^6$ of the near groups enter; so do the factor $(2L)^d$ of part (b)
of the summation lemma (Lemma~\ref{4.26:lem-summation}), the constant of the hull estimate, and
the ratio $16=\mathsf{l}_j^2/\mathsf{l}^2$. No other member changes, and no exponent of $r_A$ or
of $\kappa_A$ changes. The count $V$, the radii $r^\sharp_\square$, $r_T(\square)$,
$r^{\mathrm{nr}}_\square$, $r^{\mathrm{rem}}_T(\square)$ of Definition~\ref{4.26:def-radii-window}
before multiplication by $\mathsf{w}_0$, and every sum are those of the analysis of the
inherited-domain theorem (with
$\theta=\tfrac12$), the members of the summation lemma being written out; only the factor
$\mathsf{w}_0$ is new.
\item[(b)] Say that a positive function of $\mathsf{t}$ has degree $a$ if it is bounded by a
constant independent of $\mathsf{t}$ times $\mathsf{t}^{a}$ on $\mathsf{t}\ge 1$ (constants such
as $\mathsf{w}_0$
are absorbed). Degree $0^+$ means degree $\epsilon$ for every $\epsilon>0$, and degree
$-\infty$ means degree $a$ for every real $a$. The degrees of the counts are as follows:
$N$ has degree $r$; $\bar V$ has degree $d(1+\beta_0)r\le 6r$, where
$\beta_0$ is the constant, in Ba{\l}aban's notation, entering the exponent of the count $\bar V$;
$\mathsf{l}^2$ has degree $dr$;
$N_0$ has degree $0^+$; and $\mathsf{r}$ has degree $0$. Consequently the members of the
condition have the following degrees (the third column anticipates the degree condition among
\eqref{4.26:eq-inherited-conditions-a}, see (c)):
\begin{equation}\label{4.26:eq-smallness-reduction-a}
\begin{array}{lll}
\text{member} & \text{degree} & \text{under the degree condition}\\[3pt]
N\mathsf{r}\bar V\cdot\vartheta/\mathsf{w}_0 & 7r-\mathsf{m} & \le -2\\
(N_0+3)\mathsf{l}^2\cdot\vartheta/\mathsf{w}_0 & 4r+0^+-\mathsf{m} & <-2\\
N\mathsf{r}\bar V\cdot K_Q/r_A & 7r-\mathsf{m}_0 & \le -1\\
(N_0+3)\mathsf{l}^2\cdot K_Q/r_A & 4r+0^+-\mathsf{m}_0 & <-1\\
\text{members of \eqref{4.26:eq-summation-c} with }(100L\check R)^d\text{ or }\mathsf{l}^2
 & \le 6r-\mathsf{m},\ 4r-\mathsf{m} & \le -1\\
{}[\,\cdot\,]\cdot\varepsilon_S & -\infty & \text{---}
\end{array}
\end{equation}
In the fifth row, $\mathsf{m}_0$ replaces $\mathsf{m}$ for the members carrying $K_Q/r_A$. In
the last row, $[\,\cdot\,]$ stands for any of the count factors. In the third row the member
$K_Q/r_A$ enters at its full power, the first power of $r_A^{-1}$, so that its degree is
$7r-\mathsf{m}_0$ and not $7r-\mathsf{m}_0/2$.
\item[(c)] Assume the degree condition among the conditions \eqref{4.26:eq-inherited-conditions-a},
under which $\mathsf{m}_0\ge 7r+1$, and $\mathsf{m}\ge 7r+2$ since $p_1<p_0$
\cite{B-CMP122a}. Then every degree in \eqref{4.26:eq-smallness-reduction-a} is at most $-1$.
The set of $\gamma$ for which \eqref{4.26:eq-summed-smallness-a} holds is therefore non-empty,
and no lower bound on $\mathsf{m}$ or $\mathsf{m}_0$ beyond these two is needed. By
Definition~\ref{4.26:def-summed-smallness}, condition \eqref{4.26:eq-summed-smallness-a} is
implied by its form with a generic prefactor $K'$,
\begin{equation}\label{4.26:eq-smallness-reduction-1}
\begin{gathered}
\sup_{\mathsf{t}\ge\mathsf{t}_\gamma} K'\,\bigl[N\mathsf{r}\bar V+(N_0+3)\mathsf{l}^2\bigr]
\bigl[\mathsf{w}_0^{-1}\mathsf{t}^{-\mathsf{m}}+C_{\mathrm{far}}\mathsf{t}^{-\mathsf{m}_0}
+e^{-\varphi(\mathsf{t}-1)+c_\flat}\bigr]\le 1,\\
\mathsf{t}_\gamma-1\ge\mathsf{t}_\ast .
\end{gathered}
\end{equation}
By (a), the smallness condition of the analysis of the inherited-domain theorem, written in the
same form (generic prefactor times the
count bracket), is \eqref{4.26:eq-smallness-reduction-1} with the factor $\mathsf{w}_0^{-1}$
deleted:
\begin{equation}\label{4.26:eq-smallness-reduction-2}
\begin{gathered}
\sup_{\mathsf{t}\ge\mathsf{t}_\gamma} K'\,\bigl[N\mathsf{r}\bar V+(N_0+3)\mathsf{l}^2\bigr]
\bigl[\mathsf{t}^{-\mathsf{m}}+C_{\mathrm{far}}\mathsf{t}^{-\mathsf{m}_0}
+e^{-\varphi(\mathsf{t}-1)+c_\flat}\bigr]\le 1,\\
\mathsf{t}_\gamma-1\ge\mathsf{t}_\ast .
\end{gathered}
\end{equation}
Condition \eqref{4.26:eq-smallness-reduction-1}, and hence \eqref{4.26:eq-summed-smallness-a},
is implied by \eqref{4.26:eq-smallness-reduction-2} with $K'$ replaced by
$K'/\mathsf{w}_0=24K'$. This holds whichever of its forms the prefactor of
\eqref{4.26:eq-smallness-reduction-2} is written in: $K'\mathsf{t}^{7r}$, or a power of degree
$(1+d+d\beta_0)r$ with the degree condition read in the corresponding form. Here the prefactor,
in each of these forms, is read as a majorant of $K_\flat'[N\mathsf{r}\bar V+(N_0+3)\mathsf{l}^2]$,
the constant $K'$ being generic. The count $V$ enters only in the form fixed in (a); no variant
of $V$ counted over a run is used.
\end{itemize}
\end{lemma}

\begin{proof}
(a) The run reads the radii of the discs in the complex parameters $t$, $\tau$ of the proof
family $\mathfrak{W}^{(l),\natural}(Y)$ of Lemma~\ref{4.26:lem-family-consumption} only through
the four quantities
\begin{equation*}
\frac{2}{R^\natural_\square}=K_A\vartheta^\natural(\square),\qquad
\frac{2}{\mathsf{w}_0 r_T(\square)},\qquad
\frac{1}{\mathsf{w}_0 r^{\mathrm{nr}}_\square},\qquad
\frac{2}{\mathsf{w}_0 r^{\mathrm{rem}}_T(\square)}
\end{equation*}
of \eqref{4.26:eq-radii-window-a}, which are constant multiples of the reciprocals of these
radii. These are the quantities that enter the expansion on a stage
system. They are summed by Definition~\ref{4.26:def-radii-window},
by the sums over the cover, and by the summation lemma, Lemma~\ref{4.26:lem-summation}, in the
forms \eqref{4.26:eq-summation-a}, \eqref{4.26:eq-summation-b} and \eqref{4.26:eq-summation-c}.
They are collected in the step of the proof of the expansion on a stage system at which the
summed condition \eqref{4.26:eq-summed-smallness-a} is used. By part (iii) of
Lemma~\ref{4.26:lem-family-consumption}, the circle in the auxiliary complex parameter $u$ of
that lemma is the same as in the analysis of the inherited-domain theorem. Hence the only member
of the bracket in which
the spend appears is $\vartheta^\natural=\vartheta/\mathsf{w}_0$. The numbers listed in (a) enter
only the constant prefactor.

(b) The two remaining factors are $\vartheta$, proportional to $\mathsf{t}^{-\mathsf{m}}$, and
$r_A$, proportional to $\mathsf{t}^{\mathsf{m}_0}$. Thus $\vartheta/\mathsf{w}_0$ has degree $-\mathsf{m}$
and $K_Q/r_A$ has degree $-\mathsf{m}_0$. The degree of a product is the sum of the degrees of
its factors. With $d=4$, the count $\mathsf{l}^2$ has degree $4r$. The count $N\mathsf{r}\bar V$
has degree at most $r+0+6r=7r$, and $(N_0+3)\mathsf{l}^2$ has degree $4r+0^+$. This gives the
second column of \eqref{4.26:eq-smallness-reduction-a}. The last row has degree $-\infty$
because $\varepsilon_S$ has degree $-\infty$ and every count factor has finite degree.

(c) Insert $\mathsf{m}\ge 7r+2$ and $\mathsf{m}_0\ge 7r+1$ into the second column. For the first
and third rows this gives $7r-\mathsf{m}\le -2$ and $7r-\mathsf{m}_0\le -1$. For the second and
fourth rows it gives
\begin{equation*}
4r+0^+-\mathsf{m}\le -3r-2+0^+<-2,\qquad 4r+0^+-\mathsf{m}_0\le -3r-1+0^+<-1,
\end{equation*}
the positive degree $0^+$ being chosen smaller than $3r$. For the fifth row it gives
\begin{equation*}
6r-\mathsf{m}\le -r-2,\qquad 6r-\mathsf{m}_0\le -r-1,\qquad
4r-\mathsf{m}\le -3r-2,\qquad 4r-\mathsf{m}_0\le -3r-1,
\end{equation*}
and each of these is at most $-1$.

Every member of \eqref{4.26:eq-summed-smallness-a} has the degree listed for it in
\eqref{4.26:eq-smallness-reduction-a}. So for $\mathsf{t}\ge\mathsf{t}_\gamma\ge 1$ each member
is bounded by a constant, independent of $\mathsf{t}$ and $\gamma$, times
$\mathsf{t}^{-1}\le\mathsf{t}_\gamma^{-1}$; in the form \eqref{4.26:eq-smallness-reduction-1} of
the condition, $\gamma$ enters only through the lower limit $\mathsf{t}_\gamma$ of the supremum.
The supremum over $\mathsf{t}\ge\mathsf{t}_\gamma$ is
then at most the required bound once $\mathsf{t}_\gamma$ is large enough, that is, once $\gamma$
is small enough. The remaining requirement $\mathsf{t}_\gamma-1\ge\mathsf{t}_\ast$ is also a lower
bound on $\mathsf{t}_\gamma$. The set of admissible $\gamma$ is therefore non-empty, and only the
degree condition on $\mathsf{m}$ and $\mathsf{m}_0$ has been used.

For the implication, note that $\mathsf{w}_0^{-1}=24\ge 1$ and all three members of the brackets
are non-negative. Hence
\begin{equation*}
\mathsf{w}_0^{-1}\mathsf{t}^{-\mathsf{m}}+C_{\mathrm{far}}\mathsf{t}^{-\mathsf{m}_0}
+e^{-\varphi(\mathsf{t}-1)+c_\flat}
\le\mathsf{w}_0^{-1}\bigl[\mathsf{t}^{-\mathsf{m}}+C_{\mathrm{far}}\mathsf{t}^{-\mathsf{m}_0}
+e^{-\varphi(\mathsf{t}-1)+c_\flat}\bigr].
\end{equation*}
Multiply by $K'[N\mathsf{r}\bar V+(N_0+3)\mathsf{l}^2]$ and take the supremum over
$\mathsf{t}\ge\mathsf{t}_\gamma$. The left side of \eqref{4.26:eq-smallness-reduction-1} is then
at most the left side of \eqref{4.26:eq-smallness-reduction-2} with $K'$ replaced by
$K'/\mathsf{w}_0=24K'$. The conditions on $\mathsf{t}_\gamma$ coincide. Since
\eqref{4.26:eq-smallness-reduction-1} implies \eqref{4.26:eq-summed-smallness-a} by
Definition~\ref{4.26:def-summed-smallness}, the latter is implied by
\eqref{4.26:eq-smallness-reduction-2} at $24K'$ as well. This step uses only the
factorisation of the bracket, so it holds in whichever form the prefactor is written.

For example, at $(r,p_1,p_0,q)=(2,40,50,65)$ one has $\mathsf{m}=25$ and $\mathsf{m}_0=15$ (these
values serve only to illustrate the arithmetic of \eqref{4.26:eq-smallness-reduction-a}; they are
not asserted to satisfy the remaining conditions \eqref{4.26:eq-inherited-conditions-a}). The
first four rows of \eqref{4.26:eq-smallness-reduction-a} then have degrees
\begin{equation*}
-11,\qquad -17^+,\qquad -1,\qquad -7^+,
\end{equation*}
and the $\mathsf{m}$-members of the fifth row have degrees $-13$ and $-17$.
\end{proof}

\begin{remark}
Since $R^\natural_\square=\mathsf{w}_0 r^\sharp_\square$ by \eqref{4.26:eq-radii-window-a} and
$\mathsf{w}_0=1/24$, one has $R^\natural_\square=r^\sharp_\square/24$. The members
$\vartheta/\mathsf{w}_0$ and $K_Q/r_A$ enter \eqref{4.26:eq-smallness-reduction-a} at their full
power, not through their square roots. No condition in which these members enter through their
square roots is used, and in particular $q-p_1\ge 14r+1$ is not a hypothesis of
Lemma~\ref{4.26:lem-smallness-reduction}.
\end{remark}

\begin{remark}[The order in which the constants are chosen]\label{4.26:rem-constant-order}
The constants of this chapter are chosen in the following order. The symbols $\Theta_1$,
$\varpi_T$, $C_{\mathrm{far}}$, $\gamma_2$, $\delta_0$, $B_3$ and $C_\ast$ below denote constants
of this chapter, each fixed where it is introduced; this remark uses only their places in the
order and the dependences listed below.
\begin{enumerate}
\item[(a)] The newborn number $\theta_S$ is a number with
\begin{equation*}
\theta_S\in\bigl[\tfrac12,\tfrac89\bigr],
\end{equation*}
fixed before every other constant and depending on none of them. It is fixed once, in this
remark, and the same value serves clause~(i) of Theorem~0, the correlation theorem and the
uniqueness theorem. The following depend on $\theta_S$: the lower bound on $C_0$; the constant
$K_R$ of the correlation theorem and, through it, the constants $c_V$ and $c_A$; and one of the
conditions of the uniqueness theorem.
\item[(b)] Next the blocking factor $L$ is chosen, and then the parameter $\delta=\delta(L)$ of
\cite{B-CMP116}, in the window of the correlation theorem:
\begin{equation}\label{4.26:eq-constant-order-1}
0<\delta\le\delta_J(L);
\end{equation}
here $\delta_J(L)>0$ is the right end of that window.
\item[(c)] After $\delta$ come, in this order: the numbers depending on $(d,L)$ only; the
constants $\Theta_1$ and $\varpi_T$; $\kappa$; $\kappa_1$; $M$; $A_0$ and $A_1$. The triple
$(q,C_0,C_1)$ is chosen together with the condition imposed on it. After the constants just
listed and after this triple come, in this order: the constant $C_{\mathrm{far}}$; the triple
$(\gamma_2,\alpha,c_A)$, where $\alpha$ stands for the radii and $c_A$ is the constant of
item~(a); the spend $\mathsf{w}_0$ of Definition~\ref{4.26:def-radii-window}; and finally
$\gamma$.
\item[(d)] The spend $\mathsf{w}_0$ depends on $\theta_S$ only. The constant $K_{\min}$ depends
on $d$, $L$, $\delta_0$, $\kappa_1$, the constant $\beta$ of
Definition~\ref{4.26:def-inherited-conditions}, $B_3$, $C_0$ and $C_\ast$, and the degree
$q-p_1$, through the constants $\mathsf{K}_r$, $K^{+}$, $K^{\mathrm{nr}}$ and
$K^{\mathrm{rem}}$ of the radii and sums of Definition~\ref{4.26:def-radii-window}, and it is
fixed before $\gamma$.
\item[(e)] At the step~(b), and only there, the window \eqref{4.26:eq-constant-order-1} is used,
as follows. One of the conditions of the inherited-domain theorem
(Definition~\ref{4.26:def-inherited-conditions}) requires, at level $\ell=k$, that the decay
rate $\kappa_T$ of a term $T$ satisfy $\kappa_T\ge(1+4\beta)\kappa$; the decay rate
$\kappa_{\mathrm{out}}$ with which the term $T$ is delivered satisfies
\begin{equation}\label{4.26:eq-constant-order-2}
\kappa_{\mathrm{out}}\ \ge\ (1+4\beta_s)\kappa\ \ge\ (1+4\beta)\kappa,
\qquad \beta_s=\tfrac14,\quad \beta=\tfrac15,
\end{equation}
where $\beta_s$ is the constant in the bound by which $T$ is delivered and $\beta$ is the
constant of Definition~\ref{4.26:def-inherited-conditions},
whenever \eqref{4.26:eq-constant-order-1} holds. Both relations in
\eqref{4.26:eq-constant-order-2} are inequalities in the direction displayed, not equalities.
\end{enumerate}
No constant chosen after $L$ returns to $L$, and no constant chosen after $M$ returns to $M$.
\end{remark}

\begin{proof}
Items~(a), (b) and~(d) list only dependences of a constant on constants placed before it in the
order: $\theta_S$ depends on no other constant; $\delta$ depends on $L$ alone, and the window
\eqref{4.26:eq-constant-order-1} is a one-sided upper bound on $\delta$ in terms of $L$,
imposed once; $\mathsf{w}_0$ depends on $\theta_S$ only; and $K_{\min}$ is fixed before
$\gamma$, so that $\gamma$ may depend on $K_{\min}$ while $K_{\min}$ does not depend on $\gamma$.

For item~(e), the first inequality in \eqref{4.26:eq-constant-order-2},
$\kappa_{\mathrm{out}}\ge(1+4\beta_s)\kappa$, is the output bound by which the term $T$ is
delivered, the statement of this chapter that gives its decay rate $\kappa_{\mathrm{out}}$;
it holds whenever \eqref{4.26:eq-constant-order-1} holds,
that is, for $\delta\le\delta_J(L)$. For the second inequality, $\beta_s=\frac14$ and
$\beta=\frac15$ give $\beta_s\ge\beta$, hence
\begin{equation*}
1+4\beta_s=2\ \ge\ \tfrac95=1+4\beta ;
\end{equation*}
multiplying by the decay rate $\kappa>0$ gives $(1+4\beta_s)\kappa\ge(1+4\beta)\kappa$.
Combining the two inequalities yields $\kappa_{\mathrm{out}}\ge(1+4\beta)\kappa$. Since the
decay rate $\kappa_T$ of the term $T$ so delivered is $\kappa_{\mathrm{out}}$, this is
the requirement $\kappa_T\ge(1+4\beta)\kappa$ of
Definition~\ref{4.26:def-inherited-conditions} at level $\ell=k$. This condition is therefore met
once $\delta$ has been chosen in the window \eqref{4.26:eq-constant-order-1}, and it imposes
no further condition on a constant chosen after $\delta$.

Among the conditions named in this remark, the only one in which $L$ appears is the
window \eqref{4.26:eq-constant-order-1}, which is a condition on $\delta$ with $L$ already
fixed, and $M$ appears in none of them. For the constants of~(c), the last sentence of the
remark, that no constant chosen after $L$ returns to $L$ and none chosen after $M$ returns to
$M$, is a property of the conditions by which each of them is fixed where it is introduced in
this chapter.
\end{proof}

\begin{definition}[Stage data and the frozen data of live older pieces]\label{4.26:not-stage-frozen-data}
The following conventions restate explicitly the stage conventions and the conventions for live
older pieces under which the large-field operation is analysed; they are in force in this section.

\emph{Stage and regions.} Fix the level $k$ and a site of Ba{\l}aban's large-field operation
$\mathbf{R}$ at level $k$. Let $R_k$ be the number of steps of the run of $\mathbf{R}$ at this
site, in the notation of \cite{B-CMP122b} (the run of $\mathbf{R}$ there is written with $N = R_j$).
Put
\begin{equation*}
  \eta = L^{-k}, \qquad h = k - R_k, \qquad k_0 = k - N_0 ,
\end{equation*}
where $N_0 = N_0(k)$ is Ba{\l}aban's integer of \cite{B-CMP122a} at level $k$, the highest
exponent of $L_0^2$ among the weights of \cite{B-CMP122a} at that level. Here and throughout this
section $L_0$ denotes Ba{\l}aban's base of the position-dependent weights of \cite{B-CMP122a}, a
fixed number $L_0>1$; it is not the constant $L_0(N)$ of Notation~\ref{0.2:not-glossary}. For
$l\in\{1,\dots,R_k\}$, stage $l$ integrates the level
\begin{equation}\label{4.26:eq-stage-frozen-data-1}
  j = \ell_l := h + l - 1 .
\end{equation}
Let $Z$ be the union of all connected components of Ba{\l}aban's classes (i) and (ii)
\cite{B-CMP122a} that enter stage $1$ at this site; we assume throughout the nesting condition that
during the $R_k$ steps of the run no components are joined inside $Z$. Let $Z^{[l]}$ be
the union of the components still in the run at stage $l$; let
$\mathfrak{B} := \mathbb{B}^{(l-1)}_k$ be the determining-set structure after $l-1$ stages; let
$R^{(l)}$ and $\mathcal{C}^{(l)}$ denote the collar bond sets of stage $l$. Put
\begin{equation}\label{4.26:eq-stage-frozen-data-2}
  \mathfrak{T} := \Omega_k \cup Z^c, \qquad Z^{\mathrm{on}} := Z \cap \Omega_k^c ,
\end{equation}
with $\Omega_k$ Ba{\l}aban's small-field domain at level $k$ \cite{B-CMP122a}. On $\mathfrak{T}$
no stage changes the determining set, in accordance with \cite{B-CMP122a}. Put $\beta := 1/5$ and
\begin{equation}\label{4.26:eq-stage-frozen-data-3}
  \rho^{(l)}_m := \beta\, 2^{-|m-\ell_l-1|} .
\end{equation}
Let $\mathfrak{L}_l$ be the set of clauses defining the space
$\tilde{U}^{(l)c,\Sigma\flat}_k$ of stage $l$; on an excluded component $W$ (below) it is the
arrest-stage set $\mathfrak{S}^{(k),\mathrm{arrest}}_l$, by which we mean the set of clauses held
at stage $l$ of the run of $\mathbf{R}$ at level $k$ on a component at which that run is arrested.

\emph{Excluded components.} A component $W_{\mathfrak{a}}$ of $Z\cap\Omega_k^c$ that takes a
large-field branch after $n^+$ stages, where $n^+\le R_k$ since the run has $R_k$ stages, is called
\emph{live}; it is out of the run for $l > n^+$, and
its structure is frozen:
\begin{equation}\label{4.26:eq-stage-frozen-data-4}
  \mathfrak{B}\restriction W_{\mathfrak{a}} := \mathbb{B}^{(n^+)}_k\restriction W_{\mathfrak{a}} .
\end{equation}
Such a component, with its data, is written as an \emph{arrest}
$\mathfrak{a} = (W_{\mathfrak{a}}, k_a, n^+_a, \text{branch})$, where $k_a$ is the level of its
site and $n^+_a$ is the number $n^+$ of stages after which it took the branch.

\emph{Live older pieces.} Let $\mathfrak{a} = (W_{\mathfrak{a}}, k_a, n^+_a, \text{branch})$ be an
arrest live at $k$, with $k_a < k$. Put
\begin{gather*}
  N_a := R_{k_a}, \qquad h_a := k_a - N_a, \qquad N_0^a := N_0(k_a), \\
  k_0 := k_0^a := k_a - N_0^a, \qquad j^+ := h_a + n^+_a \le k_a, \\
  I(\mathfrak{a}) := \{h_a+1,\dots,k_a+1\}, \qquad
  \sigma_{\mathfrak{a}}(\nu) := \sum_{i\in I(\mathfrak{a})} 2^{-|\nu-i|} .
\end{gather*}
From here on, and in the table below, $k_0$ abbreviates $k_0^a$.
At $y\in W_{\mathfrak{a}}$ let $\ell_{\mathfrak{a}}(y)$ be the frozen level and $w_{\mathfrak{a}}(y)$
the frozen weight, and put $m(y) := \ell_{\mathfrak{a}}(y)$, $\omega(y) := w_{\mathfrak{a}}(y)$;
the constant $c_{\mathfrak{a}}$ attached to the piece is $c_{\mathfrak{a}} \equiv 1$; in the table
below the index $m$ of a stratum or shell is a running
index, distinct from the function $m(y)$. Put further
\begin{gather*}
  r_a(y) := \omega(y)\,\mathbf{r}_{a,\ell_{\mathfrak{a}}(y)}, \qquad
  r_e^{(p)}(y) := \omega(y)\, r_e\bigl(\ell_{\mathfrak{a}}(y), \min\{p,\ell_{\mathfrak{a}}(y)\}\bigr),\\
  \theta_{\mathrm{src},a}(y) := F^{(k)}(y)\, r_a(y), \qquad
  n_W(y) := k - 1 - \ell_{\mathfrak{a}}(y),
\end{gather*}
where $F^{(k)}(y)$ is the source factor at level $k$ and $r_e(m,p)$ is the unweighted radius at
level $m$ and block side $p$; at a retained structure $\mathfrak{s}$ the block side
$\min\{p,\ell_{\mathfrak{a}}(y)\}$ is replaced by $\min\{p,\lambda_{\mathfrak{s}}(y)\}$. Finally
\begin{equation}\label{4.26:eq-stage-frozen-data-5}
  \mathfrak{c}(y) := \ell_{\mathfrak{a}}(y) - \tfrac12 \log_{L_0}\omega(y),
  \qquad\text{so that}\qquad
  \omega = L_0^{2(\ell_{\mathfrak{a}}-\mathfrak{c})},
\end{equation}
and $g(y) := \ell_{\mathfrak{a}}(y) - \mathfrak{c}(y)$, the lift of $y$ (not the coupling $g$). A
point $y\in W_{\mathfrak{a}}$ is called \emph{lifted} if $g(y)\ge 1$. The lifted set is
$\mathsf{L}_{\mathfrak{a}} := (\Omega^c_{j^+}\setminus Z''_{k_0+1})\cap W_{\mathfrak{a}}$; that it
consists of the lifted points is taken from Ba{\l}aban's strata
\cite[(1.64)--(1.68), p.~190]{B-CMP122a}, in which the table below is written.

\emph{The frozen data.} In Ba{\l}aban's strata $\Gamma''_m$, $\Omega''_j$, $Z''_j$
\cite[(1.64)--(1.68), p.~190]{B-CMP122a}, the triple
$(\ell_{\mathfrak{a}}, \omega, \mathfrak{c})$ takes the following values:
\begin{itemize}
\item on the core and on $\Gamma''_m$ with $m\le k_0$: $(m, 1, m)$;
\item on the thin strata $\Gamma''_m$ with $k_0 < m < j^+$: $(m, L_0^{2(m-k_0)}, k_0)$;
\item on the band $\Omega''_{j^+}\setminus\Omega_{k_0+1}$: $(j^+, L_0^{2(j^+-k_0)}, k_0)$;
\item on the old shells of index $m\in\,]k_0,j^+]$: $(j^+, L_0^{2(j^+-m)}, m)$;
\item on the old shells of index $m$ with $j^+ < m < k_a$: $(m, 1, m)$.
\end{itemize}
Consequently
\begin{equation}\label{4.26:eq-stage-frozen-data-6}
  \ell_{\mathfrak{a}} \le k_a \quad\text{at every } y\in W_{\mathfrak{a}},
  \qquad
  \omega \le L_0^{2N_0^a},
\end{equation}
the second in accordance with \cite{B-CMP122a} read at level $k_a$, where the highest power of
$L_0^2$ appearing is $k_a-k_0 = N_0^a$. When $N_0^a\ge 1$, equality in the second bound holds only on the
band and only when $j^+ = k_a$. If $W_{\mathfrak{a}}\subset Z^c$, then
$W_{\mathfrak{a}}\subset\mathfrak{T}$. If $W_{\mathfrak{a}}\subset Z$, which is the case when
$W_{\mathfrak{a}}$ lies in the core of $Z$ and $k_a\le h$, then $W_{\mathfrak{a}}\subset\Omega_k^c$,
and so $W_{\mathfrak{a}}\subset Z^{\mathrm{on}}$. Finally, for every $m$,
\begin{equation}\label{4.26:eq-stage-frozen-data-7}
  \sum_{l} \rho^{(l)}_m \le 3\beta .
\end{equation}
\end{definition}

\begin{proof}
We first prove \eqref{4.26:eq-stage-frozen-data-7}. The levels $\ell_l = h+l-1$ are pairwise
distinct as $l$ runs over $\{1,\dots,R_k\}$. Hence the integers $m-\ell_l-1$ are pairwise
distinct, and so
\begin{equation*}
  \sum_l \rho^{(l)}_m \le \beta\sum_{n\in\mathbb{Z}} 2^{-|n|} = \beta\,(1+2\cdot 1) = 3\beta .
\end{equation*}

Next we check the bound $j^+\le k_a$. The stage count satisfies $n^+_a\le R_{k_a} = N_a$, since the
run at level $k_a$ has $N_a$ stages, and therefore
\begin{equation*}
  j^+ = h_a + n^+_a \le h_a + N_a = k_a .
\end{equation*}

The identity $\omega = L_0^{2(\ell_{\mathfrak{a}}-\mathfrak{c})}$ in
\eqref{4.26:eq-stage-frozen-data-5} is the definition of $\mathfrak{c}$ solved for $\omega$. Each
entry of the table satisfies it: $2(m-m)=0$; $2(m-k_0)$; $2(j^+-k_0)$; $2(j^+-m)$; $0$.

Reading $g = \ell_{\mathfrak{a}}-\mathfrak{c}$ off the table: $g = 0$ on the core, on
$\Gamma''_m$ with $m\le k_0$ and on the old shells with $j^+<m<k_a$; $g = m-k_0\ge 1$ on the thin
strata; $g = j^+-k_0$ on the band, which is $\ge 1$ exactly when $j^+>k_0$; and $g = j^+-m$ on the
old shells with $k_0<m\le j^+$, which is $\ge 1$ exactly when $m<j^+$. This determines, stratum by
stratum, the lifted points of $W_{\mathfrak{a}}$.

For the first bound in \eqref{4.26:eq-stage-frozen-data-6} we read off $\ell_{\mathfrak{a}}$ from
the table in each case, using $j^+\le k_a$ and, in the first line, $N_0^a\ge 0$:
\begin{align*}
  &\text{core and } \Gamma''_m,\ m\le k_0: && \ell_{\mathfrak{a}} = m \le k_0 = k_a - N_0^a \le k_a,\\
  &\text{thin strata:} && \ell_{\mathfrak{a}} = m < j^+ \le k_a,\\
  &\text{band and old shells with } m\in\,]k_0,j^+]: && \ell_{\mathfrak{a}} = j^+ \le k_a,\\
  &\text{old shells with } j^+<m<k_a: && \ell_{\mathfrak{a}} = m < k_a .
\end{align*}

For the second bound we compare the exponent $2g = 2(\ell_{\mathfrak{a}}-\mathfrak{c})$ of $L_0$ in
$\omega$ with $2N_0^a = 2(k_a-k_0)$:
\begin{align*}
  &\text{core, } \Gamma''_m\ (m\le k_0), \text{ old shells with } j^+<m<k_a: && 2g = 0,\\
  &\text{thin strata:} && 2g = 2(m-k_0) < 2(j^+-k_0) \le 2(k_a-k_0),\\
  &\text{band:} && 2g = 2(j^+-k_0) \le 2(k_a-k_0),\\
  &\text{old shells with } k_0<m\le j^+: && 2g = 2(j^+-m) < 2(j^+-k_0) \le 2(k_a-k_0),
\end{align*}
where on the band equality holds exactly when $j^+ = k_a$.
Thus $2g\le 2N_0^a$ everywhere, and $\omega\le L_0^{2N_0^a}$ follows since $L_0>1$. When
$N_0^a\ge 1$, the exponent $0$ is strictly smaller than $2N_0^a$. In that case equality can occur
only on the band and only when $j^+ = k_a$.

Finally, the two inclusions for $W_{\mathfrak{a}}$. If $W_{\mathfrak{a}}\subset Z^c$, then
$W_{\mathfrak{a}}\subset\Omega_k\cup Z^c = \mathfrak{T}$ by \eqref{4.26:eq-stage-frozen-data-2}. If
$W_{\mathfrak{a}}\subset Z$, the inclusion $W_{\mathfrak{a}}\subset\Omega_k^c$ is the one
established for such pieces of the core in the treatment of $Z^{\mathrm{on}}$ at the sites of
$\mathbf{R}$; together with $W_{\mathfrak{a}}\subset Z$ it gives
$W_{\mathfrak{a}}\subset Z\cap\Omega_k^c = Z^{\mathrm{on}}$ by
\eqref{4.26:eq-stage-frozen-data-2}.
\end{proof}

\begin{definition}[Conventions after an arrest and the four arrest lemmas]
\label{4.26:not-arrest-conventions}
The conventions of Notation~\ref{4.26:not-stage-frozen-data} are in force; in particular $N$ is the
number of steps of a run, and the frozen data of live older pieces are those fixed there. An arrest is
written $\mathfrak{a}$. It carries its region $W_{\mathfrak{a}}$, its level $k_{a}$, its number
$n^+_a$ and its branch. The arrests that enter below are those live at the current density.
Throughout,
$k'$ is a level at which a space is written after the arrest, and $x$ is a point. At a point $x$
of the region of a live arrest, $\ell(x)$, $w(x)$ and $\mathcal{Q}(x)$ denote the frozen level, the
frozen weight and the frozen cube clauses read at $x$ according to clause~(a) below.
\begin{enumerate}
\item[(a)] \emph{Reading at a point of a live arrest.} Every space written after the arrest reads,
at $x\in W_{\mathfrak{a}}$, the frozen triple $(\ell,w,\mathcal{Q})(x)$ of the innermost live
arrest whose region contains $x$, with the quantity $c$ of this reading set identically equal
to $1$.
\item[(b)] \emph{The source factor.} For a level $\nu$ put
\begin{equation}\label{4.26:eq-arrest-conventions-1}
s^{(k')}_{\nu}=1-\beta\bigl(1-2^{-(k'-\nu)}\bigr),
\end{equation}
where $\beta$ is the constant of the one-step rooms
$\rho^{(l)}_m=\beta\,2^{-|m-\ell_l-1|}$, $\ell_l$ being the level integrated at stage $l$
(Notation~\ref{4.26:not-stage-frozen-data}).
The source factor at $x$ is
\begin{equation}\label{4.26:eq-arrest-conventions-2}
F^{(k')}(x)=s^{(k')}_{\ell(x)}
-\beta\sum_{\mathfrak{a}\ \text{live},\ x\in W_{\mathfrak{a}}}\sigma_{\mathfrak{a}}(\ell(x)).
\end{equation}
\item[(c)] \emph{The source space.} For a set $X$, the space
$\tilde{U}^{\mathrm{arrest}}_{k'}(X)$ is defined, at every point $x\in X$ lying in a live piece
(Notation~\ref{4.26:not-stage-frozen-data}), by upper bounds on the
quantities $Q_a(x)$, $a\le 7$, measured at $x$, namely by the conditions
\begin{equation}\label{4.26:eq-arrest-conventions-3}
Q_a(x)<F^{(k')}(x)\,w(x)\,r_a(\ell(x)),\qquad a\le 7,
\end{equation}
together with the cube clauses read at weight $w(x)$. At the points of $X$ off the live pieces,
the conditions of
$\tilde{U}^{\mathrm{arrest}}_{k'}(X)$ are those of \cite[(2.34)--(2.39)]{B-CMP119}, as stated there. The
two readings of the derived-entry family, as fixed where that family is defined, are unchanged.
\item[(d)] \emph{Attempt spaces and dissolution.} The attempt spaces used after the arrest, and
the dissolution of an arrest, are fixed by two further conventions, stated together with
conventions (a)--(c) where arrests are defined, and not restated here.
\end{enumerate}
The following four statements are used together with these conventions; of each we recall only the
part used below.
\begin{itemize}
\item \emph{The first lemma} consists of five clauses, which are stated in that lemma and are not
restated here.
\item \emph{The floor lemma for the source factor.} Put
\begin{equation}\label{4.26:eq-arrest-conventions-4}
\varphi_{\mathrm{arrest}}=1-4.51\beta=0.098.
\end{equation}
The lemma states that $F^{(k')}(x)\ge\varphi_{\mathrm{arrest}}$ at every point $x$ of the
region of a live arrest and at every level $k'$, with no dependence on the level. The bound holds
uniformly in the level because the only properties of the levels on which it depends are the
following two. Number the live arrests whose regions contain $x$ in order of increasing level,
and write $k_{a_1},k_{a_2},\dots$ for their levels. The first property is
$\ell(x)\le k_{a_1}$. The second is the gap condition
$k_{a_{s+1}}\ge k_{a_s}+N$ for every~$s$.
\item \emph{The extended room comparison} has two parts.
Part~(a) states that the room $\rho_n(x;k')$ satisfies
\begin{equation}\label{4.26:eq-arrest-conventions-5}
\rho_n(x;k')\ge\beta\,2^{-(k_a+1-\ell)},
\end{equation}
where $n$ indexes the stage, $k_a$ is the level of a live arrest $\mathfrak{a}$ whose region
contains $x$, $\ell=\ell(x)$ is the frozen level at $x$, and both sides of the comparison read
the same frozen triple $(\ell,w,\mathcal{Q})$.
Part~(b) states that at raised points the source is at most $0.74$ times the threshold
$\theta_{T,a}$ of a term $T$.
\item \emph{The radius-degree lemma} contains a condition indexed by the degree
$\mathsf{m}$, stated in that lemma, and, in addition, the lower bound
\begin{equation}\label{4.26:eq-arrest-conventions-6}
C_0^{\mathrm{rad}}\ge\frac{2.16\,A_0}{\varphi_{\mathrm{arrest}}}
\end{equation}
on its constant $C_0^{\mathrm{rad}}$.
Finally, among the entries of the list stated in the radius-degree lemma, the second, third,
fourth and eighth are multiples of the quantity $\varepsilon_{\ell}$ at the frozen level
$\ell=\ell(x)$, with factors depending on $(d,L)$ only, by \cite[Theorem~1]{B-CMP102b} and
\cite{B-CMP109}.
\end{itemize}
\end{definition}

\begin{lemma}[Levels and weights of retained structures]\label{4.26:lem-retained-levels}
Let the site $k$, the region $Z$, the small-field regions $\Omega_j$, $j\le k$, the regions
$Z''_j$ and strata
$\Gamma''_m$, and the current structure $\mathfrak{B}=\mathbb{B}^{(l-1)}_k$ be as in
Notation~\ref{4.26:not-stage-frozen-data}. Let $\mathfrak{a}$ be an arrest live at the site $k$,
with $k_a<k$, region $W_{\mathfrak{a}}$, run data $h_a$, $n^+_a$, $k_0=k_0^a$, $j^+$, frozen level
$\ell_{\mathfrak{a}}(y)$ and weight $\omega(y)$, $y\in W_{\mathfrak{a}}$, as there. Let
$\mathfrak{c}(y)=\ell_{\mathfrak{a}}(y)-\tfrac12\log_{L_0}\omega(y)$ be the original-scale level
of the piece, and let
\begin{equation*}
\mathsf{L}_{\mathfrak{a}}=\bigl(\Omega^c_{j^+}\setminus Z''_{k_0+1}\bigr)\cap W_{\mathfrak{a}}
\end{equation*}
be its lifted set. For $y'\in W_{\mathfrak{a}}$ let $m(y')$ be the level at $y'$ of the frozen
structure $\mathbb{B}^{(n^+_a)}_{k_a}$; this is the frozen level, $m(y')=\ell_{\mathfrak{a}}(y')$.

Recall, from the definition of the spaces of the construction and from its extension to live
pieces, the two families of structures at which the derived entries are imposed.
\begin{itemize}
\item[$(\Sigma^{\ast})$] In every space the derived entries~$2$ and~$4$ are imposed, with the
factor that Ba{\l}aban's construction attaches to the shell, for every triple
$(p,X',\mathfrak{s})$ in which $p$ is at most the top level, $X'\subseteq Y$ is a localization
domain, $Y$ being the domain on which the space is written, and $\mathfrak{s}$ is a
determining-set structure on $X'$ among the current one and all its predecessors in the
construction (the stage structures $\mathbb{B}^{(n')}_k$, $n'\le l-1$, $l-1$ being the index of
the current structure $\mathfrak{B}$, before and after the operation $\mathbf{R}$), localized as in
Ba{\l}aban's construction; the accompanying condition on $U$ that the clause list of the space
attaches to the derived entries is imposed as well.
\item[$(\Sigma^{\ast})^{\mathfrak{a}}$] On a live piece the structures of $(\Sigma^{\ast})$
include the piece's own structures $\mathbb{B}^{(n')}_{k_a}$ restricted to
$X'\cap W_{\mathfrak{a}}$, $0\le n'\le n^+_a$, continued on $X'\setminus W_{\mathfrak{a}}$ by
the structures of $\mathfrak{B}$. In particular
$(p,\square_0,\mathfrak{B}^{\mathfrak{c}}\restriction\square_0)$,
$p\le p_{\square}$, is a $(\Sigma^{\ast})$-triple for every $\square\in\mathcal{C}^1$ whose
$\square_0=\tilde{\square}^5$ meets $W_{\mathfrak{a}}$, junction boxes
($\square_0\not\subset W_{\mathfrak{a}}$) included, and
$\mathfrak{B}^{\mathfrak{c}}\restriction\square_0
=\mathbb{B}^{(k_0-h_a)}_{k_a}\restriction\square_0$ as an identity of structures, by the
description of the junction boxes of a live piece.
\end{itemize}
For a structure $\mathfrak{s}$ of these families and a point $y'$ let $\lambda_{\mathfrak{s}}(y')$
be the level of the block of $\mathfrak{s}$ at $y'$, and put
\begin{equation}\label{4.26:eq-retained-levels-1}
w_{\mathfrak{s}}(y')=L_0^{\,2(\lambda_{\mathfrak{s}}(y')-\mathfrak{c}(y'))^+}.
\end{equation}
This is the weight that $\mathfrak{s}$ carried at its own stage
(cf.\ \cite[(1.65), (1.67)]{B-CMP122a}). Then the following hold.
\begin{enumerate}
\item[(i)] Let $y'\in W_{\mathfrak{a}}$. If $\mathfrak{s}$ is the frozen structure or any
structure of the present site, then $\lambda_{\mathfrak{s}}(y')=m(y')$. The strict inequality
$\lambda_{\mathfrak{s}}(y')<m(y')$ occurs only when $\mathfrak{s}$ is a retained predecessor
$\mathbb{B}^{(n')}_{k_a}$ of the frozen structure in the piece's run, restricted as in
$(\Sigma^{\ast})^{\mathfrak{a}}$, whose level at $y'$ a later stage of the piece's run raised.
\item[(ii)] Let $y'\in W_{\mathfrak{a}}$, let $\mathfrak{s}=\mathbb{B}^{(n')}_{k_a}$ with
$0\le n'\le n^+_a$, and let $j'=h_a+n'$. Then
\begin{equation}\label{4.26:eq-retained-levels-a}
\begin{gathered}
\lambda_{\mathfrak{s}}(y')=
\begin{cases}
\max\{k_0,\min\{j',m'\}\} & y'\in\Gamma''_{m'}\ \text{a thin stratum},\\
\max\{k_0,j'\} & y'\ \text{in the band},\\
\max\{m',j'\} & y'\ \text{in the old shell of index } m',\ k_0<m'\le j^+,\\
m''-1 & y'\in\Gamma''_{m''}\cap\Omega^c_{m''},\ \Gamma''_{m''}\ \text{thick},\ j'<m'',\\
\ell_{\mathfrak{a}}(y') & \text{otherwise};
\end{cases}\\[4pt]
\lambda_{\mathfrak{s}}\ge\mathfrak{c}\quad\text{on }\mathsf{L}_{\mathfrak{a}}.
\end{gathered}
\end{equation}
Here the thin and thick strata, the band and the old shells of the run of $\mathfrak{a}$ are as
in Notation~\ref{4.26:not-stage-frozen-data}.
\item[(iii)] For every $y'\in W_{\mathfrak{a}}$ and every $\mathfrak{s}$ as in~(ii),
$\mathfrak{c}(y')-1\le\lambda_{\mathfrak{s}}(y')\le m(y')$, and
$\mathfrak{c}\in[k_0,j^+-1]$ on $\mathsf{L}_{\mathfrak{a}}$.
\item[(iv)] For $\mathfrak{s}=\mathfrak{B}^{\mathfrak{c}}$, that is $j'=k_0$: at every
$y'\in\square_0\cap W_{\mathfrak{a}}$,
$\lambda_{\mathfrak{s}}(y')=\mathfrak{c}(y')$ and $w_{\mathfrak{s}}(y')=1$.
\end{enumerate}
\end{lemma}

\begin{proof}
\emph{Part \textup{(i)}.} Let $y'\in W_{\mathfrak{a}}$. For the frozen structure the equality
$\lambda_{\mathfrak{s}}(y')=m(y')$ holds by the definition of $m(y')$. Consider next a structure
of the present site. The stages of the present site act only inside $Z\cap\Omega^c_k$, as in
\cite{B-CMP122a}, where the operation $\mathbf{T}_k$ on $Z_k\cap Z^c$, $Z_k$ being
Ba{\l}aban's large-field region of level $k$, is left unchanged; and on
$W_{\mathfrak{a}}$ every space written after the arrest reads the frozen structure of the
innermost live arrest holding the point, by the conventions after an arrest and the arrest
lemmas of Notation~\ref{4.26:not-arrest-conventions}. Hence at $y'$ such a structure has the
block of the frozen structure, and $\lambda_{\mathfrak{s}}(y')=m(y')$.

The remaining structures of the family $(\Sigma^{\ast})^{\mathfrak{a}}$ at $y'$ are the retained
predecessors $\mathbb{B}^{(n')}_{k_a}$, $n'<n^+_a$, of the frozen structure in the piece's own
run, restricted to $X'\cap W_{\mathfrak{a}}$. The frozen structure arises from such a
predecessor through the later stages of that run. If $\lambda_{\mathfrak{s}}(y')<m(y')$ for such
a predecessor, then the level at $y'$ was changed by one of those later stages, and since the
level $m(y')$ reached at the end of the run exceeds $\lambda_{\mathfrak{s}}(y')$, one of those
stages raised the level at $y'$. Hence $\lambda_{\mathfrak{s}}(y')<m(y')$ occurs only at a
retained predecessor of the piece's run that a later stage of the run raised at $y'$.

\emph{Part \textup{(ii)}.} The table \eqref{4.26:eq-retained-levels-a} is read from the
recursion \cite[(1.16), p.~179]{B-CMP122a} that defines the structures
$\mathbb{B}^{(n')}_{k_a}$. The structure $\mathbb{B}^{(n')}_{k_a}$ is the one in force at the
stage of the run of $\mathfrak{a}$ that integrates the level $j'=h_a+n'$ (stage $n'+1$ in the
numbering of Notation~\ref{4.26:not-stage-frozen-data}). On the set
\begin{equation}\label{4.26:eq-retained-levels-2}
S_{j'}:=\Omega^c_{j'+1}\setminus Z''_{j'}
\end{equation}
the recursion gives blocks of level $j'$, and on a stratum of $\mathfrak{s}$ it gives blocks of
the level of that stratum.

We carry out the first row; let $y'\in\Gamma''_{m'}$ with $\Gamma''_{m'}$ a thin stratum. The
thin strata and the band lie in the old shell $\Omega_{k_0}\setminus\Omega_{k_0+1}$
\cite{B-CMP122a}. There are three cases.
\begin{itemize}
\item $j'<k_0$. The point $y'$ lies in the old shell $\Omega_{k_0}\setminus\Omega_{k_0+1}$,
where the level is $k_0$. Since $\min\{j',m'\}\le j'<k_0$, the tabulated value
$\max\{k_0,\min\{j',m'\}\}$ is also $k_0$.
\item $k_0\le j'\le m'$. The regions $Z''_j$ increase with $j$, so
$Z''_{j'}\subset Z''_{m'}$. The stratum $\Gamma''_{m'}$ does not meet $Z''_{m'}$, so
$y'\notin Z''_{j'}$. Moreover $y'\in\Omega^c_{k_0+1}$ by the location of the thin strata. Since
the sets $\Omega_j$ decrease with $j$ and $j'+1\ge k_0+1$, we have
$\Omega^c_{k_0+1}\subset\Omega^c_{j'+1}$. Hence $y'\in S_{j'}$, and the level at $y'$ is $j'$,
which equals $\max\{k_0,\min\{j',m'\}\}$ in this case.
\item $j'>m'$. The set $\Gamma''_{m'}$ is a stratum of $\mathfrak{s}$, and the level at $y'$ is
$m'$.
\end{itemize}
This gives the first row. The other rows of \eqref{4.26:eq-retained-levels-a} are obtained from
the same recursion in the same manner: one determines the position of $y'$ relative to
$\Omega_{j'+1}$, $Z''_{j'}$ and the strata of $\mathfrak{s}$ at the stage integrating $j'$, and
reads off the level that the recursion assigns there.

\emph{Parts \textup{(ii)}, second line, \textup{(iii)} and \textup{(iv)}.} The table
\eqref{4.26:eq-retained-levels-a}, together with the values of the original-scale level
$\mathfrak{c}$ on the same sets given by the lemma on the original-scale level $\mathfrak{c}$,
yields $\lambda_{\mathfrak{s}}\ge\mathfrak{c}$ at every point of $\mathsf{L}_{\mathfrak{a}}$, which
is the second line of \eqref{4.26:eq-retained-levels-a}, and
$\mathfrak{c}-1\le\lambda_{\mathfrak{s}}\le m$ at every point of $W_{\mathfrak{a}}$. That
$\mathfrak{c}\in[k_0,j^+-1]$ on $\mathsf{L}_{\mathfrak{a}}$ is part of that lemma.

For $\mathfrak{s}=\mathfrak{B}^{\mathfrak{c}}$ we have
$\mathfrak{B}^{\mathfrak{c}}\restriction\square_0=\mathbb{B}^{(k_0-h_a)}_{k_a}\restriction\square_0$,
so $n'=k_0-h_a$ and $j'=h_a+n'=k_0$. That $\lambda_{\mathfrak{s}}=\mathfrak{c}$ for this structure
is clause~(i) of the lemma on the original-scale level $\mathfrak{c}$. Then
$(\lambda_{\mathfrak{s}}-\mathfrak{c})^+=0$, and \eqref{4.26:eq-retained-levels-1} gives
$w_{\mathfrak{s}}=L_0^{0}=1$.
\end{proof}

\begin{remark}[Weights of retained structures and the present run]
The weight $\omega$ of the piece and the weights \eqref{4.26:eq-retained-levels-1} are related
by $\omega\le L_0^{2(m-\lambda_{\mathfrak{s}})}w_{\mathfrak{s}}$, with equality at the points of
$\mathsf{L}_{\mathfrak{a}}$, by the comparison of the weight of the piece with the weights of its
retained structures. For the present run the fact analogous to part~(i) of
Lemma~\ref{4.26:lem-retained-levels} reads as follows: on the points raised by the stages of
the present run, its retained structures $\mathbb{B}^{(n')}_k$ have
$\lambda_{\mathfrak{s}}(y')<m(y')$, with $m(y')$ there denoting the level of the current
structure $\mathfrak{B}$ at $y'$. It belongs to the description of the stage structures of the
present run and is not used in the proof of Lemma~\ref{4.26:lem-retained-levels}.
\end{remark}

\begin{definition}[The space of the current stage and its shell clause]\label{4.26:not-stage-space}
Fix the stage $n$ of the site. The determining set $\mathfrak{B}$, the region $\mathfrak{T}$ and the set
$Z^{\mathrm{on}}$ are those of Notation~\ref{4.26:not-stage-frozen-data}. The triples $(p,X',\mathfrak{s})$
of the family $(\Sigma^{\ast})$ are those of Lemma~\ref{4.26:lem-retained-levels}; for such a triple,
$\lambda_{\mathfrak{s}}(x)$ is the level of the block of $\mathfrak{s}$ containing the point $x$.
For a point $x$ we write $m(x)$ for the shell of the site containing $x$, that is, its level.
The quantities $Q_a$, $a\in\{1,\dots,8\}$, are the entries of the regularity family of a pair, and
$r_a(y)$ is the radius of entry $a$ at $y$. For a triple $(p,X',\mathfrak{s})$, $Q_e^{(p,X',\mathfrak{s})}$,
$e\in\{2,4\}$, is the derived entry $e$ read relative to that triple.

We put
\begin{equation}\label{4.26:eq-stage-space-1}
  \mathsf{w}_0=\tfrac{1}{24},\qquad \iota^1_l(y)=\tfrac23\,\mathsf{w}_0\,\rho^{(l)}_{m(y)},
\end{equation}
where $\rho^{(l)}_m=\beta\,2^{-|m-\ell_l-1|}$ is the one-step room of stage $l$ at the shell $m$, and
$\ell_l$ is the level integrated at stage $l$ (Notation~\ref{4.26:not-stage-frozen-data}).
With these values,
\begin{equation}\label{4.26:eq-stage-space-2}
  1+2\beta-2\mathsf{w}_0\beta=1.383 .
\end{equation}

Let $Y$ be a localization domain. Write $B(x)$ for the block of $\mathfrak{B}$ containing $x$, and put
\begin{equation}\label{4.26:eq-stage-space-10}
  Y^{+}:=Y\cup\bigcup_{x\in Y}B(x).
\end{equation}
The \emph{space of the current stage} $\mathfrak{N}^{(n)}(Y)$ is the set of
pairs satisfying the following clauses.
\begin{enumerate}
\item At the points of $Y^{+}\cap Z^{\mathrm{on}}$, for $a\in\{1,3,5,6,7,8\}$,
  \begin{equation}\label{4.26:eq-stage-space-3}
    Q_a<5\,w_\ast\, r_a .
  \end{equation}
  At a point $x$ the right side is read as $5\,\mathsf{w}^{\mathrm{on}}(x)\,\mathbf{r}_{a,m(x)}$.
\item At the points of $Y^{+}\cap\mathfrak{T}$, the running clauses hold:
  \begin{align}
    Q_a(y)&<\Bigl(1+2\beta-\sum_{l\le n}\iota^1_l(y)\Bigr)\,r_a(y),
      &&a\in\{1,3,6,7\},\label{4.26:eq-stage-space-4}\\
    Q_a&<(1+2\beta)\,r_a, &&a\in\{5,8\}.\label{4.26:eq-stage-space-5}
  \end{align}
\item The \emph{shell clause} holds. Take any triple $(p,X',\mathfrak{s})$ of the family $(\Sigma^{\ast})$
  with $X'\subset Y$. Take any point $y'$ of $X'^{\sim2}\cap\mathfrak{T}$, where $X'^{\sim2}$ is the
  twofold enlargement of $X'$, and any $e\in\{2,4\}$. Then
  \begin{equation}\label{4.26:eq-stage-space-6}
    Q_e^{(p,X',\mathfrak{s})}(y')
      <\Bigl(1+2\beta-\sum_{l\le n}\iota^1_l(y')\Bigr)\,r_e^{(p)}(y').
  \end{equation}
\end{enumerate}

\emph{The radius of the shell clause.} The radius $r_e^{(p)}(y')$ in \eqref{4.26:eq-stage-space-6} is the
shell radius of entry $e$ at $y'$ for the level $p$ of the triple. In the bump-function form of the shell
device it is written $r_e(m,l_{y'})$, with $m=m(y')$ and with the reference level $l_{y'}$ of
\eqref{4.26:eq-stage-space-8} below. Wherever $\lambda_{\mathfrak{s}}(y')=m$, it is given by
\begin{equation}\label{4.26:eq-stage-space-7}
  \begin{aligned}
    r_2&=\alpha_{0,m}\,L^{-2p}\,\bigl(L^{m}L^{-p}\bigr)^{-2},\\
    r_4&=\alpha_{0,m}\,L^{m-\min\{p,m\}}\,\bigl(L^{m}L^{-p}\bigr)^{-3}.
  \end{aligned}
\end{equation}
Among the unweighted shell radii, the same radius appears as $\mathbf{r}_{e,m(b)}$, read at a bond $b$,
with $m(b)$ the shell of the bond $b$.
This is the unit of the box clause itself, the unit of \cite[(1.16)]{B-CMP122a}. It is never the clause
radius $\mathbf{r}^{T,(p)}_e$ of a term $T$.

In the bump-function form of the shell device, the reference level at a point $x$ is
\begin{equation}\label{4.26:eq-stage-space-8}
  l_x:=\min\{p,\mathsf{m}^{\circ}(x)\},
\end{equation}
where $\mathsf{m}^{\circ}(x)$ is the level at $x$ of the parent structure. At a point where the parent
structure's block is coarser than the shell, that is $\mathsf{m}^{\circ}(x)<m(x)$, one has
$l_x=\mathsf{m}^{\circ}(x)<m$. Equivalently,
\begin{equation}\label{4.26:eq-stage-space-9}
  l_x=\min\{p,\lambda_{\mathfrak{s}}(x)\}.
\end{equation}
The description through $r_e(m,l_{y'})$ and the formulas \eqref{4.26:eq-stage-space-7} agree wherever
$\lambda_{\mathfrak{s}}=m$. They differ at retained structures with $\lambda_{\mathfrak{s}}<m$. There the
radius in \eqref{4.26:eq-stage-space-6} is $r_e(m,l_{y'})$, with $l_{y'}$ given by
\eqref{4.26:eq-stage-space-8}, as stated above.

\emph{Ambient pairs.} A member of $\mathfrak{N}^{(n)}(Y)$ is the restriction to $Y$ of a pair in the
polydisc $\mathfrak{K}_P=\mathfrak{K}_{109}(\mathfrak{B};\gamma_0)$ on the whole domain of $\mathfrak{B}$.
The direct clauses \eqref{4.26:eq-stage-space-3}--\eqref{4.26:eq-stage-space-5} of the space of the
current stage are imposed on $Y^{+}$, as in items 1 and 2.

\emph{Clauses that are not part of $\mathfrak{N}^{(n)}(Y)$.}
\begin{itemize}
\item The box clause relative to a term $T$ belongs to the diminished term clauses $\mathfrak{I}_n$ and
  is stated with them.
\item The composite clauses at $Z^{\mathrm{on}}$ belong to the attempt space $\mathfrak{L}_n$ of the stage,
  never to $\mathfrak{N}^{(n)}(Y)$. The clause family $\mathfrak{L}_n$ imposes the derived-entry family in
  the reading $(\Sigma^{\flat})$. In that reading the entries of the retained stages are imposed, each with
  its own shell, radius and gain, under the running threshold $\theta^{(n)}$. The radii are those fixed at
  the stage at which the structure was created, taken at the reference level
  $m_{\mathrm{ref}}:=\mathsf{m}^{\circ}(x)$.
\item The source space $\mathfrak{S}_{\mathrm{src}}$ carries every triple of the family $(\Sigma^{\ast})$
  with $X'\subseteq Y_4$, where $Y_4$ is the domain on which the source space is built, at the factor
  $s^{(k)}_m$. Its composite clauses are read at the shell $m$ of the point.
\end{itemize}
\end{definition}

\begin{definition}[Couplings, the integer $N_0(j)$ and the density-indexed maximum]
\label{4.26:not-density-maximum}
The following symbols are used throughout. The logarithm is the natural one; it is written
$\log$ or $\ln$. The level $k$ (the current density), the structure $\mathfrak{B}$, the level
$k_0$ and the live older pieces with their levels $k_a$ are those of
Notation~\ref{4.26:not-stage-frozen-data}.

\emph{Couplings and radii.} Let $g_j$ be the running couplings of
Definition~\ref{1.2:def-couplings} (Ba{\l}aban's running couplings of \cite{B-CMP119}), with
$x_j=g_j^{-2}$ as there, and put
\begin{equation*}
\mathsf{t}_j:=\log g_j^{-2}=\log x_j .
\end{equation*}
With the constants $C_0,C_1$ and the exponents $q_0=q_1=:q$ of Ba{\l}aban's radii in
\cite{B-CMP119}, put
\begin{equation*}
\alpha_{i,m}:=C_i\,x_m^{-1/2}(\log x_m)^{q}\quad(i=0,1),\qquad
\alpha_{\cdot,m}:=\min_i\alpha_{i,m}
\end{equation*}
(in this section the dot denotes the minimum over $i$, not, as in
Lemma~\ref{4.4:lem-radii-ratio}, either of the two radii),
and, as functions of one variable $\mathsf{t}>0$,
\begin{equation*}
\alpha_i(\mathsf{t}):=C_i\,\mathsf{t}^{q}e^{-\mathsf{t}/2},\qquad
\alpha(\mathsf{t}):=C\,\mathsf{t}^{q}e^{-\mathsf{t}/2},\qquad C:=\max\{C_0,C_1\}.
\end{equation*}

\emph{The numbers $R_j$ and the coupling ratios.} By \cite[(2.5)]{B-CMP119}, $R_j$ is the
smallest integer power $L^{\iota}$ of $L$ with $L^{\iota}\ge(\log g_j^{-2})^{r}$, where $r>0$
is the exponent fixed in \cite[(2.5)]{B-CMP119}. Thus
$R_j=R(\mathsf{t}_j)$, where for $\mathsf{t}>0$ we let $R(\mathsf{t})$ be the least integer power
of $L$ that is $\ge\mathsf{t}^{r}$. In this section $R(\cdot)$ denotes this function of
$\mathsf{t}=\log g^{-2}$; it is not the ball radius $R$ of Notation~\ref{0.2:not-glossary}.
\cite[(2.6), pp.~255--256]{B-CMP119} states, for $n>m$,
$g_n\le(1+\beta_0)(n-m)^{1/2}g_m$ and $g_m\le(1+\beta_0)g_n$, where $\beta_0>0$ can be taken
arbitrarily small provided the couplings are sufficiently small. We use these inequalities with
$\beta_0=1/7$; the smallness of the couplings that this requires is one of the ceilings on
$\gamma$.

\emph{The integer $N_0(j)$.} At a density $j$ at which the operation $\mathbf{R}$ is applied,
$N_0(j)$ is the integer that \cite{B-CMP122a} attaches to such a density: the smallest
positive integer $N_0$ such that $L^{-N_0+1}MR_{k-N_0+1}=M$, read at $k=j$; equivalently,
$N_0(j)$ is the smallest positive integer $n$ with $L^{n-1}=R_{j-n+1}$, so that
\begin{equation}\label{4.26:eq-density-maximum-1}
L^{N_0(j)-1}=R_{j-N_0(j)+1}.
\end{equation}
It is a function of the sequence $(R_i)$ along the trajectory, through $R_{j-N_0(j)+1}$, and
not a function of $\mathsf{t}_j$. Its existence at every density of a run at which
$\mathbf{R}$ is applied is taken from \cite{B-CMP122a}, where it rests on the rule stated there
that in some steps the number $R_k$ decreases by the factor $L^{-1}$. None of the hypotheses
\textup{(H1)}--\textup{(H7)} of Notation~\ref{2.2:not-hypotheses} and no cap of
Definition~\ref{2.3:def-cap} bounds $N_0(j)$. We put
\begin{equation*}
w_\ast(j):=L_0^{2N_0(j)},
\end{equation*}
and write $w_\ast:=w_\ast(k)$ where the weight of the current run is meant.

\emph{The density-indexed maximum.} Put
\begin{equation*}
\bar{\alpha}_{i,j}:=\max_{1\le m\le j}\alpha_{i,m},\qquad
\bar{\alpha}_j:=\max_i\bar{\alpha}_{i,j},
\end{equation*}
the maximum over the levels $m\le j$; the supremum $\sup_{\mathsf{t}'\ge\mathsf{t}_j}\alpha(\mathsf{t}')$
is a different quantity and is not used in place of $\bar{\alpha}_j$. Put
\begin{equation*}
\gamma_0:=\gamma_0^{(k)}:=5\max_{j\le k}w_\ast(j)\,\bar{\alpha}_j,
\end{equation*}
where $j$ ranges over the densities $j\le k$ at which $N_0(j)$ is defined, that is, over the
densities $j\le k$ at which $\mathbf{R}$ is applied; these include $k$ itself and the level
$k_a$ of every live older piece. We put
$\mathfrak{K}_P:=\mathfrak{K}_{109}(\mathfrak{B};\gamma_0^{(k)})$, the analyticity polydisc of
\cite{B-CMP109} over $\mathfrak{B}$ of size $\gamma_0^{(k)}$.

\emph{The one-variable majorant and the conditions on $\gamma$.} Put
\begin{equation}\label{4.26:eq-density-maximum-a}
\begin{gathered}
\Psi(\mathsf{t}):=L^2(L+1)C\,\mathsf{t}^{r+q}e^{-\mathsf{t}/2}
=L^2(L+1)\,\mathsf{t}^{r}\alpha(\mathsf{t}),\\
\bar{\gamma}_0(\mathsf{t}):=5(1+\beta_0)^2\Psi(\mathsf{t}),\qquad
c_\beta:=2\ln(1+\beta_0),\\
\mathsf{t}_j\ge\mathsf{t}_\gamma:=\log\gamma^{-2}
\quad\text{for every density } j\le k \text{ occurring below},\\
\mathsf{t}_\gamma\ge 2(r+q)+c_\beta ,
\end{gathered}
\end{equation}
and let $\Psi_i$ denote $\Psi$ with $C_i$ in place of $C$. The densities occurring below are the
levels $m\le j$ entering $\bar{\alpha}_j$ and the densities $j$ over which $\gamma_0^{(k)}$ is
maximised, among them $k$ and the levels $k_a$ of the live older pieces. Here $\beta_0$ is the
constant of \cite[(2.6)]{B-CMP119}, fixed at
$\beta_0=1/7$, so that $e^{c_\beta/2}=8/7$ and $e^{-c_\beta}=49/64$. The third line of
\eqref{4.26:eq-density-maximum-a} is the coupling ceiling; it is the upper bound
$g_j\le\gamma$ on the couplings and nothing else. Were the coupling ceiling derived from the
bound $g_K=g\le\gamma$ on the last
coupling alone, the second inequality of \eqref{4.26:eq-density-maximum-2} below would give
$\mathsf{t}_j\ge\mathsf{t}_\gamma-c_\beta$ instead, and every threshold in $\mathsf{t}_\gamma$
would shift by $c_\beta$ once. The fourth line is the range condition; it is imposed as part of
the ceiling on $\gamma_0^{(k)}$, one of the ceilings listed at the end of this definition, and is
not a separate hypothesis. The following reading
is a gloss and not part of this definition. Let
\begin{equation*}
\Psi^{\downarrow}(\mathsf{t}):=\sup_{\mathsf{t}'\ge\mathsf{t}}\Psi(\mathsf{t}'),\qquad
\alpha_i^{\downarrow}(\mathsf{t}):=\sup_{\mathsf{t}'\ge\mathsf{t}}\alpha_i(\mathsf{t}')
\end{equation*}
be the non-increasing envelopes of $\Psi$ and $\alpha_i$; by the second line of
\eqref{4.26:eq-density-maximum-b} below they coincide with $\Psi$ and $\alpha_i$ on
$[2(r+q),\infty)$. If $\Psi$ and $\alpha_i$ are read as these envelopes, the range condition is
weakened to $\mathsf{t}_\gamma\ge\max\{2r,c_\beta\}$.

\emph{Elementary functions.} For $s\ge0$ and $\mathsf{t}\ge0$ put
$\varphi_s(\mathsf{t}):=\mathsf{t}^{s}e^{-\mathsf{t}/2}$ (with $0^0=1$).

\emph{The smallness parameter, Ba{\l}aban's $N$, and constants.} Let $A_1$ and $p_1$ be
Ba{\l}aban's auxiliary parameters of these names (among the parameters $\mathfrak{p}$ of
Notation~\ref{2.1:not-prefix}), and let $C_\ast$ be the constant fixed together with them; put
\begin{equation*}
\lambda:=\frac{A_1}{C_\ast}\Bigl(\tfrac12\log g_k^{-2}\Bigr)^{p_1-q}
=\frac{A_1}{C_\ast}\,2^{q-p_1}\,\mathsf{t}_k^{\,p_1-q}.
\end{equation*}
The smallness parameter $\vartheta$ of
the expansion is a function $\vartheta=\vartheta(\mathsf{t}_k)$ of order
$\mathsf{t}_k^{\,p_1-q}$; since $q>p_1$, $\vartheta$ decreases to $0$ as $\mathsf{t}_k$
increases. The number $N$ of Notation~\ref{4.26:not-stage-frozen-data} is Ba{\l}aban's
$N=R_j$ of \cite{B-CMP122b} read at $j=k$, so $N=R_k=R(\mathsf{t}_k)$ depends on
$\mathsf{t}_k$ alone. Its one-variable majorant is
\begin{equation}\label{4.26:eq-density-maximum-c}
\bar{N}(\mathsf{t}):=L\,\mathsf{t}^{r}.
\end{equation}
Put $\mathsf{l}_j:=1$ for $j\le k_0$ and $\mathsf{l}_j:=4\,\mathsf{l}(x_k)$ for $j>k_0$, where
$\mathsf{l}(x_k)$ is the size parameter whose square is the tile number, chosen subject to the
condition $L^{4N_0(k)}\le\mathsf{l}(x_k)^2$. Put
$K^{\mathrm{nr}}:=66\,B_{\mathbb{H}}K_{\mathrm{cv}}$, where
$K_{\mathrm{cv}}=K_{\mathrm{cv}}(d,L,B_3)$
is the constant of \cite[(3.48)--(3.52)]{B-CMP109}, used as stated there, with Ba{\l}aban's
constant $B_3$ of \cite{B-CMP109}, and $B_{\mathbb{H}}$ denotes the
Lipschitz constant of the theorem constructing the stage layer $\mathbb{H}$. The constant
$K_{\min}$ is the maximum of a list of constants that includes $6K^{\mathrm{nr}}$, so that
$K_{\min}\ge6K^{\mathrm{nr}}$; it is subject to the window condition
$K_{\min}\vartheta(\mathsf{t}_k)\le\mathsf{w}_0$.
The three increment constants $P,P',P''$ of the near-pair chain are formed with
$C_{1/2}:=\sqrt2$. The conditions on the constants
of this chapter are grouped as follows: floors; the definitions fixed before $\gamma$, namely
$\beta_0=1/7$, $\theta_S$, $\mathsf{w}_0=1/24$, $\theta=\tfrac12$, $K_{\mathrm{cv}}$,
$K^{\mathrm{nr}}$, $K^{\mathrm{rem}}$, $K_{\min}$, $K_{\flat}'$, $C_V$ and $C_F$; conditions of
scope; ceilings, each a supremum over
$\mathsf{t}\ge\mathsf{t}_\gamma$, taken jointly in one variable $\mathsf{t}$; and the inequality
$16\max\{P,P',P''\}\le\beta\theta_S$.

With these definitions the following hold.
\begin{enumerate}
\item[(a)] $\alpha_{i,m}=\alpha_i(\mathsf{t}_m)$; $\gamma_0^{(k)}$ is non-decreasing in $k$.
\item[(b)] $R$ is non-decreasing on $(0,\infty)$; if $R(\mathsf{t})=L^{\iota}$ with
$\iota\ge1$, then $L^{\iota-1}<\mathsf{t}^{r}\le L^{\iota}$, in particular
$R(\mathsf{t})<L\mathsf{t}^{r}$.
\item[(c)] For $n>m$,
\begin{equation}\label{4.26:eq-density-maximum-2}
\mathsf{t}_m\le\mathsf{t}_n+c_\beta+\ln(n-m),\qquad \mathsf{t}_m\ge\mathsf{t}_n-c_\beta .
\end{equation}
\item[(d)] The elementary inequalities
\begin{equation}\label{4.26:eq-density-maximum-b}
\begin{aligned}
&\varphi_s(\mathsf{t}-c)\le e^{c/2}\varphi_s(\mathsf{t})
&&\text{for } 0\le c\le\mathsf{t},\ s\ge0,\\
&\varphi_s \text{ is non-increasing on } [2s,\infty) &&\text{for } s\ge0,\\
&L^{u}-1\ge u\ln L &&\text{for real } u .
\end{aligned}
\end{equation}
\item[(e)] If $R_k\ge L$, then $N<\bar{N}(\mathsf{t}_k)$.
\end{enumerate}
\end{definition}

\begin{proof}
The form \eqref{4.26:eq-density-maximum-1} of the defining condition of $N_0(j)$ follows from
$L^{-N_0+1}MR_{j-N_0+1}=M$ on dividing by $M>0$ and multiplying by $L^{N_0-1}$.

(a) Since $\mathsf{t}_m=\log x_m$, we have $x_m^{-1/2}=e^{-\mathsf{t}_m/2}$ and
$(\log x_m)^q=\mathsf{t}_m^{\,q}$, so $\alpha_{i,m}=C_i\mathsf{t}_m^{\,q}e^{-\mathsf{t}_m/2}
=\alpha_i(\mathsf{t}_m)$. For the monotonicity, the term $w_\ast(j)\bar{\alpha}_j$ depends on
$j$ only: $w_\ast(j)$ through $N_0(j)$, which is determined by the trajectory up to $j$, and
$\bar{\alpha}_j$ through the $\alpha_{i,m}$ with $m\le j$. If $k\le k'$, the set of densities
$j\le k$ at which $N_0(j)$ is defined is contained in the corresponding set for $k'$, and a
maximum over a larger set of the same terms is not smaller; hence
$\gamma_0^{(k)}\le\gamma_0^{(k')}$. The identity in the first line of
\eqref{4.26:eq-density-maximum-a} is
\begin{equation*}
L^2(L+1)\mathsf{t}^{r}\cdot C\mathsf{t}^{q}e^{-\mathsf{t}/2}
=L^2(L+1)C\mathsf{t}^{r+q}e^{-\mathsf{t}/2}.
\end{equation*}
With $\beta_0=1/7$ we get
$e^{c_\beta/2}=1+\beta_0=8/7$ and $e^{-c_\beta}=(1+\beta_0)^{-2}=49/64$. The coupling ceiling
reads
$\log g_j^{-2}\ge\log\gamma^{-2}$; since the logarithm is increasing and $g_j,\gamma>0$, this is
$g_j^{-2}\ge\gamma^{-2}$, that is, $g_j\le\gamma$.

(b) Let $0<\mathsf{t}\le\mathsf{t}'$. Since $r>0$, $R(\mathsf{t}')\ge\mathsf{t}'^{\,r}\ge
\mathsf{t}^{r}$, and $R(\mathsf{t}')$ is an integer power of $L$; by the minimality in the
definition of $R(\mathsf{t})$, $R(\mathsf{t})\le R(\mathsf{t}')$. Now let
$R(\mathsf{t})=L^{\iota}$ with $\iota\ge1$. Then $\mathsf{t}^{r}\le L^{\iota}$ by definition,
and $L^{\iota-1}$ is an integer power of $L$ smaller than $L^{\iota}$, so by minimality
$L^{\iota-1}\ge\mathsf{t}^{r}$ fails, that is, $L^{\iota-1}<\mathsf{t}^{r}$. Multiplying by $L$
gives $R(\mathsf{t})=L^{\iota}<L\mathsf{t}^{r}$.

(c) Let $n>m$. Squaring $g_n\le(1+\beta_0)(n-m)^{1/2}g_m$ and inverting gives
$g_m^{-2}\le(1+\beta_0)^2(n-m)g_n^{-2}$; taking logarithms,
\begin{equation*}
\mathsf{t}_m\le\mathsf{t}_n+2\ln(1+\beta_0)+\ln(n-m)=\mathsf{t}_n+c_\beta+\ln(n-m).
\end{equation*}
Squaring $g_m\le(1+\beta_0)g_n$ and inverting gives $g_m^{-2}\ge(1+\beta_0)^{-2}g_n^{-2}$;
taking logarithms, $\mathsf{t}_m\ge\mathsf{t}_n-c_\beta$.

(d) For the first inequality, $0\le\mathsf{t}-c\le\mathsf{t}$ and $s\ge0$ give
$(\mathsf{t}-c)^{s}\le\mathsf{t}^{s}$, and $e^{-(\mathsf{t}-c)/2}=e^{c/2}e^{-\mathsf{t}/2}$;
multiplying, $\varphi_s(\mathsf{t}-c)\le e^{c/2}\varphi_s(\mathsf{t})$. For the second, if
$s=0$ then $\varphi_0(\mathsf{t})=e^{-\mathsf{t}/2}$ is non-increasing. If $s>0$, then on
$[2s,\infty)$ we have $\mathsf{t}>0$ and $\varphi_s>0$, and
$(\ln\varphi_s)'(\mathsf{t})=s/\mathsf{t}-\tfrac12\le0$ there, so $\ln\varphi_s$, and hence
$\varphi_s$, is non-increasing on $[2s,\infty)$. For the third, $e^{y}\ge1+y$ for every real $y$
by convexity of the exponential; with $y=u\ln L$ this is $L^{u}-1\ge u\ln L$.

(e) By \cite[(2.5)]{B-CMP119}, $N=R_k=R(\mathsf{t}_k)$ is an integer power $L^{\iota}$ of $L$,
and $R_k\ge L$ forces $\iota\ge1$. By (b),
$N=R(\mathsf{t}_k)<L\mathsf{t}_k^{\,r}=\bar{N}(\mathsf{t}_k)$.
\end{proof}

If the envelope $\check{R}_k$ of $R_k$ is given by $\check{R}_k=R(\log\check{x}_k)$, where
$\check{x}_k=\min_{i\le k}x_i$ and $\log\check{x}_k>0$, then $\check{R}_k\le R(\mathsf{t}_k)$:
indeed $\check{x}_k\le x_k$, so $0<\log\check{x}_k\le\log x_k=\mathsf{t}_k$, and $R$ is
non-decreasing by item (b) of Notation~\ref{4.26:not-density-maximum}.

\begin{definition}[Three units at a retained structure]\label{4.26:def-three-units}
Let $\mathfrak{a}$ be a live arrest, with region $W_{\mathfrak{a}}$, let $(p,X',\mathfrak{s})$ be a
triple of the kind fixed in Lemma~\ref{4.26:lem-retained-levels}, and let
$y'\in X'^{\sim 2}\cap W_{\mathfrak{a}}$ be a point of the live piece of $\mathfrak{a}$. We put
\begin{equation*}
m:=\ell_{\mathfrak{a}}(y'),\qquad \omega:=w_{\mathfrak{a}}(y'),\qquad
\lambda:=\lambda_{\mathfrak{s}}(y'),\qquad w_b:=w_{\mathfrak{s}}(y'),
\end{equation*}
with $\lambda\le m$. Here $\ell_{\mathfrak{a}}(y')$ and $w_{\mathfrak{a}}(y')$ are the frozen level
and the frozen weight of the piece at $y'$, $\lambda_{\mathfrak{s}}(y')$ is the level of the block of
$\mathfrak{s}$ containing $y'$, and $w_{\mathfrak{s}}(y')$ is the weight that $\mathfrak{s}$ carried at
its own stage. Lengths are written $\ell_n:=L^n\eta$, where $\eta>0$ is the length unit in which the
block sides $\ell_n$ are measured, and $p\wedge n:=\min\{p,n\}$; the letter $\ell$ with a point
argument denotes a level, and with an integer subscript a length. The radii
$\alpha_{0,n}$ are those of Lemma~\ref{4.4:lem-twist-stability}, and
$\mathfrak{c}=\mathfrak{c}(y')$ is the original-scale level
$m-\tfrac12\log_{L_0}\omega$, the logarithm being taken to the base $L_0$ of the weights.

For the entry $e=4$ the \emph{three units at the retained structure $\mathfrak{s}$}, and the
unweighted unit of the $\mathfrak{c}$-shell, are
\begin{equation}\label{4.26:eq-three-units-a}
\begin{aligned}
\mathrm{U1}&:=\omega\,\alpha_{0,m}\,\ell_m^{-2}\,\ell_{p\wedge m}^{-1},\\
\mathrm{U2}&:=\omega\,\alpha_{0,m}\,\ell_m^{-2}\,\ell_{p\wedge\lambda}^{-1},\\
\mathrm{U3}&:=w_b\,\alpha_{0,\lambda}\,\ell_\lambda^{-2}\,\ell_{p\wedge\lambda}^{-1},\\
\mathbf{r}^{\mathfrak{c},(p)}_4&:=\alpha_{0,\mathfrak{c}}\,\ell_{\mathfrak{c}}^{-2}\,
\ell_{p\wedge\mathfrak{c}}^{-1}.
\end{aligned}
\end{equation}
For the entry $e=2$ the units $\mathrm{U1}_2,\mathrm{U2}_2,\mathrm{U3}_2$ are obtained from the first
three lines of \eqref{4.26:eq-three-units-a} by dropping the last length; in particular
\begin{equation*}
\mathrm{U1}_2=\mathrm{U2}_2=\omega\,\alpha_{0,m}\,\ell_m^{-2}.
\end{equation*}

Thus $\mathrm{U1}$ is the weighted shell unit taken entirely at the frozen level $m$, in the form in
which the shell clause of Notation~\ref{4.26:not-stage-space} is written; $\mathrm{U2}$ keeps the
frozen level $m$ and the frozen weight $\omega$ and takes the block side at the level $\lambda$ of
$\mathfrak{s}$, that is, $\mathrm{U2}=\omega\,r_4(m,l_x)$ at $l_x=\lambda$, where, for the level $p$
of the triple and a level $l_x\le m$, we put
\begin{equation*}
r_4(m,l_x):=\alpha_{0,m}\,\ell_m^{-2}\,\ell_{p\wedge l_x}^{-1},
\end{equation*}
the radius of entry $4$ at level $m$ with block side at level $l_x$; $\mathrm{U3}$ is the radius of
$\mathfrak{s}$ at its birth, with level $\lambda$ and weight $w_b$ throughout.
\end{definition}

\begin{lemma}[The stage space lies in the analyticity domain at every age]\label{4.26:lem-analyticity-membership}
Let $k$ be the current density, an $\mathbf{R}$-site, and use the stage data, the frozen data of live older pieces
and the space $\mathfrak{N}^{(n)}(Y)$ of the current stage, as fixed in
Notation~\ref{4.26:not-stage-frozen-data} and Notation~\ref{4.26:not-stage-space}. Use the radii $\alpha_{i,m}$,
$\alpha_{\cdot,m}=\min_i\alpha_{i,m}$, and the weights $w_\ast(j)=L_0^{2N_0(j)}$ of
Notation~\ref{4.26:not-density-maximum}. Put
\begin{equation}\label{4.26:eq-analyticity-membership-1}
\bar{\alpha}_{i,j}:=\max_{1\le m\le j}\alpha_{i,m},\qquad
\bar{\alpha}_j:=\max_i\bar{\alpha}_{i,j},\qquad
\gamma_0^{(k)}:=5\max_{j\le k}w_\ast(j)\,\bar{\alpha}_j ,
\end{equation}
the last maximum being the density-indexed maximum of Notation~\ref{4.26:not-density-maximum}, and put
$\mathfrak{K}_P:=\mathfrak{K}_{109}(\mathfrak{B};\gamma_0^{(k)})$. Let $\beta$ be the schedule constant of the
layered analyticity spaces (Definition~\ref{4.4:def-post-step-space}), for which $1+2\beta\le 1.4$. Certain live
pieces carry an attached density $j^+\le k$ (Notation~\ref{4.26:not-stage-frozen-data}); on these the frozen
weights and levels are read at $j^+$ rather than at the creation density. Then the following hold.
\begin{itemize}
\item[(a)] Let $\mathfrak{a}$ be a live arrest with $k_a<k$, let $y\in W_{\mathfrak{a}}\cap\mathfrak{T}$, and let
$\ell:=\ell_{\mathfrak{a}}(y)$ be its frozen level, so that $\ell\le k_a$ by the lemma on the frozen data of a
live piece. Then
\begin{equation}\label{4.26:eq-analyticity-membership-2}
(1+2\beta)\,\omega(y)\,\alpha_{\cdot,\ell}\le 1.4\,w_\ast(k_a)\,\bar{\alpha}_{k_a}
\le \frac{1.4}{5}\,\gamma_0^{(k)}<\gamma_0^{(k)} ,
\end{equation}
so that $\mathfrak{N}^{(n)}\subset\mathfrak{K}_P$ at $y$, whatever the age $k-k_a$ of the piece.
\item[(b)] At a point of $\Omega_k$, or of a component $W'\subset Z_k\setminus Z$ untouched by this site, the
weight is $1$ and the level $m$ satisfies $m\le k$, and $(1+2\beta)\,\alpha_{\cdot,m}\le(1.4/5)\,\gamma_0^{(k)}$.
\item[(c)] At every $x\in Z^{\mathrm{on}}$, in each of the two branches
$\mathsf{w}^{\mathrm{on}}(x)=w_\ast(k)$ and $\mathsf{w}^{\mathrm{on}}(x)=\omega(x)$,
\begin{equation}\label{4.26:eq-analyticity-membership-3}
5\,\mathsf{w}^{\mathrm{on}}(x)\,\alpha_{\cdot,m(x)}\le\gamma_0^{(k)} .
\end{equation}
\item[(d)] The bounds of (a)--(c) hold at a point lying in several nested live pieces, with the datum of the
innermost one, and on the pieces carrying an attached density $j^+$.
\end{itemize}
In every case the weight and the level are read at one density $j\le k$, and the bound is the membership of the
term $w_\ast(j)\bar{\alpha}_j$ in the finite maximum \eqref{4.26:eq-analyticity-membership-1}.
\end{lemma}

\begin{proof}
All quantities involved are positive: the weights are powers of $L_0$ and the radii are
$\alpha_{i,m}=C_i x_m^{-1/2}(\log x_m)^q$. Two facts are used throughout, both definitional unwindings of
\eqref{4.26:eq-analyticity-membership-1}. First, if $1\le m\le j$, then for every $i$
\begin{equation}\label{4.26:eq-analyticity-membership-4}
\alpha_{\cdot,m}\le\alpha_{i,m}\le\bar{\alpha}_{i,j}\le\bar{\alpha}_j ,
\end{equation}
because $\alpha_{i,m}$ is one of the terms of the maximum defining $\bar{\alpha}_{i,j}$. Second, if $j\le k$ is a
density at which $N_0(j)$ is defined, then $w_\ast(j)\bar{\alpha}_j$ is one of the terms of the maximum in
\eqref{4.26:eq-analyticity-membership-1}, whence $5\,w_\ast(j)\bar{\alpha}_j\le\gamma_0^{(k)}$.

\emph{Points of $W_{\mathfrak{a}}\cap\mathfrak{T}$.} Let $\mathfrak{a}$, $y$ and $\ell$ be as in (a). The frozen
weight at $y$ is $\omega(y)=L_0^{2(\ell-\mathfrak{c}(y))}$, and by the table of frozen weights and levels on a live
piece (Notation~\ref{4.26:not-stage-frozen-data}), together with the bound $L_0^{2N_0}$ on the weights in
Ba{\l}aban's construction \cite{B-CMP122a}, read at the density $k_a$,
\begin{equation*}
\omega(y)=L_0^{2(\ell-\mathfrak{c}(y))}\le L_0^{2N_0^a}=w_\ast(k_a).
\end{equation*}
Here $N_0^a=N_0(k_a)$ is the integer of Notation~\ref{4.26:not-density-maximum} at the density $k_a$, that is,
the smallest positive integer $N_0$ with $L^{-N_0+1}MR_{k_a-N_0+1}=M$ in \cite{B-CMP122a}, with $R_i$
Ba{\l}aban's radii as in Notation~\ref{4.26:not-stage-frozen-data}; it exists by the
construction of the piece $W_{\mathfrak{a}}$ itself. Since $\ell\le k_a$, the bound
\eqref{4.26:eq-analyticity-membership-4} with $m=\ell$, $j=k_a$ gives
$\alpha_{\cdot,\ell}\le\bar{\alpha}_{k_a}$: the level $\ell$ is one of the levels $m\le k_a$ of the maximum. Since
$k_a<k$ and $k_a$ is an $\mathbf{R}$-site, the term $w_\ast(k_a)\bar{\alpha}_{k_a}$ belongs to the maximum over
densities $j\le k$, so $w_\ast(k_a)\bar{\alpha}_{k_a}\le\gamma_0^{(k)}/5$. Multiplying the two positive bounds
and using $1+2\beta\le1.4$ gives \eqref{4.26:eq-analyticity-membership-2}; the last inequality there holds
because $\gamma_0^{(k)}>0$. This holds for every live arrest with $k_a<k$, hence at every age, and it is the
inclusion $\mathfrak{N}^{(n)}\subset\mathfrak{K}_P$ at $y$.

\emph{Untouched regions.} On the components $W'\subset Z_k\setminus Z$ untouched by this site and on
$\Omega_k$ the weights equal $1$ and the levels $m$ satisfy $m\le k$. Since $w_\ast(k)\ge1$,
\eqref{4.26:eq-analyticity-membership-4} with $j=k$ gives
\begin{equation*}
(1+2\beta)\cdot1\cdot\alpha_{\cdot,m}\le1.4\,w_\ast(k)\,\bar{\alpha}_k\le\frac{1.4}{5}\,\gamma_0^{(k)},
\end{equation*}
the second inequality being the membership of the term $j=k$. This is (b).

\emph{Points of $Z^{\mathrm{on}}$, branch $\mathsf{w}^{\mathrm{on}}(x)=w_\ast(k)$.} The current structure
$\mathbb{B}^{(l-1)}_k$, and any component $W$ excluded at this site (frozen at density $k$, as in
Notation~\ref{4.26:not-stage-frozen-data}), have levels $m(x)\le k$ and weights at most
$L_0^{2N_0(k)}=w_\ast(k)$. By \eqref{4.26:eq-analyticity-membership-4} with $j=k$ and the membership of the
term $j=k$,
\begin{equation*}
5\,w_\ast(k)\,\alpha_{\cdot,m(x)}\le5\,w_\ast(k)\,\bar{\alpha}_k\le\gamma_0^{(k)} .
\end{equation*}
The thin strata and the band of the run itself (Notation~\ref{4.26:not-stage-frozen-data}) fall under this branch.

\emph{Points of $Z^{\mathrm{on}}$, branch $\mathsf{w}^{\mathrm{on}}(x)=\omega(x)$.} Here $x$ lies on a live piece
$W_{\mathfrak{b}}\subset Z$ created at a density $k_b$ with $k_b\le h<k$, by the placement of live pieces in
Notation~\ref{4.26:not-stage-frozen-data}, and $m(x)=\ell_{\mathfrak{b}}(x)\le k_b$.
By the table of frozen weights read at the density $k_b$, as in the first case, $\omega(x)\le w_\ast(k_b)$;
by \eqref{4.26:eq-analyticity-membership-4} with $j=k_b$, $\alpha_{\cdot,m(x)}\le\bar{\alpha}_{k_b}$. Hence
\begin{equation*}
5\,\omega(x)\,\alpha_{\cdot,m(x)}\le5\,w_\ast(k_b)\,\bar{\alpha}_{k_b}\le\gamma_0^{(k)},
\end{equation*}
the last step being the membership of the term $j=k_b$, $k_b$ being the creation density of a live piece, at
which $N_0(k_b)$ is defined. The two branches together give
\eqref{4.26:eq-analyticity-membership-3}, which is (c).

\emph{Nested arrests.} Let $\mathfrak{a}_1\subset\mathfrak{a}_2\subset\cdots$ be live arrests whose pieces all
contain $x$. By the rule of the arrest construction that reads, at a point held by several nested live arrests,
the datum of the innermost one (on $Z^{\mathrm{on}}$ this is the definition of the frozen weight $\omega(x)$ and
of the level $m(x)$ in Notation~\ref{4.26:not-stage-space}), the datum read at $x$ is
$(\ell_{\mathfrak{a}_1}(x),\omega(x))$, with $\omega(x)\le w_\ast(k_{a_1})$ and
$\ell_{\mathfrak{a}_1}(x)\le k_{a_1}\le k$. The argument of (a), or of the branch
$\mathsf{w}^{\mathrm{on}}(x)=\omega(x)$ of (c), according as $x\in\mathfrak{T}$ or $x\in Z^{\mathrm{on}}$, therefore
applies with $j=k_{a_1}$: \eqref{4.26:eq-analyticity-membership-4} with $m=\ell_{\mathfrak{a}_1}(x)$, $j=k_{a_1}$ and
the membership of the term $j=k_{a_1}$.

\emph{Pieces carrying an attached density $j^+$.} On these pieces the same table of frozen weights and levels
applies at the density $j^+$. By the lemma comparing the real weights of these pieces with the table weight,
the frozen weights actually carried by these pieces lie one power below the table weight
$w_\ast(j^+)=L_0^{2N_0(j^+)}$; in particular they are at most $w_\ast(j^+)$. Since $j^+\le k$ is a density at which
$N_0(j^+)$ is defined (Notation~\ref{4.26:not-stage-frozen-data}), the term $w_\ast(j^+)\bar{\alpha}_{j^+}$
belongs to the maximum \eqref{4.26:eq-analyticity-membership-1}, and the argument of (a), or of the branch
$\mathsf{w}^{\mathrm{on}}(x)=\omega(x)$ of (c), according as the point lies in $\mathfrak{T}$ or in
$Z^{\mathrm{on}}$, applies with $j=j^+$. This proves (d).

In each case above the bound was obtained from a weight and a level read at a single density $j\le k$; no
inequality between the couplings $g_j$ and $g_{j'}$ at two different densities $j\ne j'$ has been used.
\end{proof}

\begin{remark}\label{4.26:thm-density-majorant}
The statement of this theorem is not written out here; parts of its proof are printed below.
\end{remark}

\begin{lemma}[A trajectory-free bound on the weight at one density]\label{4.26:lem-weight-bound}
Let the logarithmic couplings $\mathsf{t}_i=\ln g_i^{-2}$, the integers $R_i$ and $N_0(i)$, the weight
$w_{\ast}(i)=L_0^{2N_0(i)}$, the number $\beta_0$, the constant $c_\beta=2\ln(1+\beta_0)$ and the exponent
$r$ of \cite[(2.5)]{B-CMP119} be as in
Notation~\ref{4.26:not-density-maximum} and \eqref{4.26:eq-density-maximum-a}. Let $j$ be a density at
which $N_0:=N_0(j)$ is defined, and assume the following.
\begin{itemize}
\item $L\ge 3$ and $L_0^2\le L/3$.
\item $R_j\ge L$ (which holds whenever $R_{\min}\ge 5L$).
\item $\mathsf{t}_j\ge 2r$.
\item The running couplings satisfy the first inequality of \cite[(2.6), p.~255]{B-CMP119} at the one
pair $(m,n)=(j-N_0+1,j)$, that is,
\begin{equation}\label{4.26:eq-weight-bound-1}
g_n\le(1+\beta_0)(n-m)^{1/2}g_m,\qquad (m,n)=(j-N_0+1,j).
\end{equation}
\end{itemize}
Then
\begin{equation}\label{4.26:eq-weight-bound-a}
w_{\ast}(j)<0.37\,L^2(L+1)\,\mathsf{t}_j^{\,r}.
\end{equation}
No assumption is made on how $R_i$ varies with the density $i$.
\end{lemma}

\begin{proof}
We recall two definitions from Notation~\ref{4.26:not-density-maximum}. By \cite[(2.5)]{B-CMP119},
$R_i$ is the smallest integer power of $L$ that is $\ge\mathsf{t}_i^{\,r}$. Consequently, if $R_i=L^{a}$
with an integer $a\ge1$, then $\mathsf{t}_i^{\,r}\le L^{a}$, and $L^{a-1}<\mathsf{t}_i^{\,r}$,
since otherwise the integer power $L^{a-1}$ would be $\ge\mathsf{t}_i^{\,r}$ and smaller than $R_i$.
We refer to this as minimality at the density $i$. The integer $N_0(j)$ satisfies the equality
$L^{N_0(j)-1}=R_{j-N_0(j)+1}$ of \cite{B-CMP122a}; we call it the defining equality of $N_0$.

Write $R_j=L^{\iota}$. The hypothesis $R_j\ge L$ gives $\iota\ge1$, and minimality at the density $j$ gives
\begin{equation}\label{4.26:eq-weight-bound-2}
L^{\iota-1}<\mathsf{t}_j^{\,r}\le L^{\iota}.
\end{equation}
Define the integer $d$ by the defining equality of $N_0$:
\begin{equation*}
R_{j-N_0+1}=L^{N_0-1}=:L^{\iota+d}.
\end{equation*}
Then $\iota+d=N_0-1\ge0$. The integer $d$ is the net number of decreases of $R$ by the factor $L^{-1}$
across the densities $j-N_0+1<i\le j$, and it may be $\le0$; nothing about how $R$ varies is assumed.

Since $w_{\ast}(j)=(L_0^2)^{N_0}$, $N_0=\iota+d+1$ and $L_0^2\le L/3$, we have
\begin{equation}\label{4.26:eq-weight-bound-3}
w_{\ast}(j)\le(L/3)^{\iota+d+1}.
\end{equation}
By the left inequality of \eqref{4.26:eq-weight-bound-2},
\begin{equation}\label{4.26:eq-weight-bound-4}
0.37\,L^2(L+1)\,\mathsf{t}_j^{\,r}>0.37\,L^2(L+1)L^{\iota-1}=0.37\,(L+1)L^{\iota+1}.
\end{equation}
By \eqref{4.26:eq-weight-bound-3} and the strict inequality \eqref{4.26:eq-weight-bound-4}, it suffices to
prove
\begin{equation}\label{4.26:eq-weight-bound-5}
(L/3)^{\iota+d+1}\le 0.37\,(L+1)L^{\iota+1}.
\end{equation}

\emph{Case $d\le1$.} Since $L\ge3$, we have $L/3\ge1$, so $(L/3)^a$ is non-decreasing in the
integer $a$, and $\iota+d+1\le\iota+2$. Since $\iota\ge1$, we have $3^{\iota+2}\ge27$; moreover $L\le L+1$.
Hence
\begin{align*}
(L/3)^{\iota+d+1}\le(L/3)^{\iota+2}
&=L^{\iota+1}\cdot\frac{L}{3^{\iota+2}}
\le\frac{L^{\iota+1}(L+1)}{27}
=\frac{0.37\,(L+1)L^{\iota+1}}{0.37\cdot27}.
\end{align*}
Since $0.37\cdot27=9.99$ and $1/9.99<0.1002$, this is \eqref{4.26:eq-weight-bound-5} with a factor
$0.1002$ to spare. No coupling enters this case.

\emph{Case $d\ge2$.} Then $N_0-1=\iota+d\ge3$, so the density $m:=j-N_0+1$ satisfies $m<j$ and is a
density of the trajectory, with $R_m=L^{\iota+d}$ and $\iota+d\ge1$. Minimality at the density $m$ gives
$\mathsf{t}_m^{\,r}>L^{\iota+d-1}$, and with the right inequality of \eqref{4.26:eq-weight-bound-2},
\begin{equation}\label{4.26:eq-weight-bound-6}
\mathsf{t}_m^{\,r}>L^{\iota+d-1}=L^{d-1}L^{\iota}\ge L^{d-1}\mathsf{t}_j^{\,r}.
\end{equation}
Since $\mathsf{t}_j\ge2r>0$, taking $r$-th roots in \eqref{4.26:eq-weight-bound-6} gives
$\mathsf{t}_m>L^{(d-1)/r}\mathsf{t}_j$.

We now use the hypothesis \eqref{4.26:eq-weight-bound-1} at the pair $(m,n)=(j-N_0+1,j)$, for which
$n-m=N_0-1=\iota+d\ge1$; it bounds $\mathsf{t}_m$, the logarithmic coupling at the earlier density,
from above. Raising
\eqref{4.26:eq-weight-bound-1} to the power $-2$ gives
$g_n^{-2}\ge(1+\beta_0)^{-2}(n-m)^{-1}g_m^{-2}$, and taking logarithms, with
$c_\beta=2\ln(1+\beta_0)$,
\begin{equation}\label{4.26:eq-weight-bound-7}
\mathsf{t}_m\le\mathsf{t}_j+c_\beta+\ln(\iota+d).
\end{equation}
Combining $\mathsf{t}_m>L^{(d-1)/r}\mathsf{t}_j$ with \eqref{4.26:eq-weight-bound-7},
\begin{equation*}
\mathsf{t}_j\bigl(L^{(d-1)/r}-1\bigr)<c_\beta+\ln(\iota+d).
\end{equation*}
By the elementary inequality $L^u-1\ge u\ln L$ of \eqref{4.26:eq-density-maximum-b}, applied with
$u=(d-1)/r>0$, the factor $L^{(d-1)/r}-1$ is non-negative and at least $((d-1)/r)\ln L$; together with
$\mathsf{t}_j\ge2r$ this gives
\begin{equation*}
\mathsf{t}_j\bigl(L^{(d-1)/r}-1\bigr)\ge2r\cdot\frac{d-1}{r}\ln L=2(d-1)\ln L.
\end{equation*}
Hence $2(d-1)\ln L<c_\beta+\ln(\iota+d)$; exponentiating, $L^{2(d-1)}<e^{c_\beta}(\iota+d)$. Since
$e^{-c_\beta}=49/64$ by \eqref{4.26:eq-density-maximum-a},
\begin{equation}\label{4.26:eq-weight-bound-8}
\iota+d>e^{-c_\beta}L^{2(d-1)}=\tfrac{49}{64}\,L^{2(d-1)}.
\end{equation}
The logarithm $\ln(n-m)=\ln(\iota+d)$ of \eqref{4.26:eq-weight-bound-7} is absorbed in
\eqref{4.26:eq-weight-bound-8} and is not carried further.

Dividing \eqref{4.26:eq-weight-bound-5} by $L^{\iota+1}>0$ and multiplying by $3^{\iota+d+1}$ shows that
\eqref{4.26:eq-weight-bound-5} is equivalent to
\begin{equation*}
L^{d}\le0.37\,(L+1)\,3^{\iota+d+1}.
\end{equation*}
By \eqref{4.26:eq-weight-bound-8}, and since $\tfrac{49}{64}\ln3>0.84$,
\begin{equation*}
\ln\bigl(3^{\iota+d}\bigr)=(\iota+d)\ln3>\tfrac{49}{64}(\ln3)\,L^{2(d-1)}>0.84\,L^{2(d-1)}.
\end{equation*}
We claim that $0.84\,L^{2(d-1)}\ge d\ln L$ for all $L\ge3$ and $d\ge2$. Indeed,
$L^{2(d-1)}=(L^2)^{d-1}\ge(d-1)L^2$, because $(L^2)^{d-1}=L^2(L^2)^{d-2}$ and
$(L^2)^{d-2}\ge2^{d-2}\ge d-1$;
moreover $d\le2(d-1)$ since $d\ge2$, and $\ln L<L$. Therefore
\begin{equation*}
d\ln L<2(d-1)L\le0.84\,(d-1)L^2\le0.84\,L^{2(d-1)},
\end{equation*}
where the middle inequality holds as soon as $0.84\,L\ge2$, that is, once $L\ge2.39$, and in particular
for $L\ge3$. This proves the claim. Consequently $\ln(3^{\iota+d})>d\ln L$, that is,
$3^{\iota+d}>L^d$. Since $3\cdot0.37\,(L+1)>1$, we obtain
\begin{equation*}
L^d<3^{\iota+d}<0.37\,(L+1)\,3^{\iota+d+1},
\end{equation*}
which is \eqref{4.26:eq-weight-bound-5}.

In both cases \eqref{4.26:eq-weight-bound-5} holds. Combining it with \eqref{4.26:eq-weight-bound-3} and
the strict inequality \eqref{4.26:eq-weight-bound-4},
\begin{equation*}
w_{\ast}(j)\le(L/3)^{\iota+d+1}\le0.37\,(L+1)L^{\iota+1}<0.37\,L^2(L+1)\,\mathsf{t}_j^{\,r},
\end{equation*}
which is \eqref{4.26:eq-weight-bound-a}.

The argument uses only the defining equality of $N_0$, minimality in \cite[(2.5)]{B-CMP119} at the two
densities $j$ and $j-N_0+1$, the inequalities $L_0^2\le L/3$ and $\mathsf{t}_j\ge2r$, and the hypothesis
\eqref{4.26:eq-weight-bound-1} at one pair, the last only in the case $d\ge2$. It uses no monotonicity of
$i\mapsto R_i$, $x_i$ or $\mathsf{t}_i$, not the second inequality of \cite[(2.6)]{B-CMP119}, and not the
minimality in the definition of $N_0(j)$; neither \cite[(2.8), (2.9)]{B-CMP119} nor the bound by
$(L+1)(N_0+1)^{\beta_0}R_k$ in \cite{B-CMP122a} enters.
Since $\Psi(\mathsf{t})=L^2(L+1)\,\mathsf{t}^{\,r}\alpha(\mathsf{t})$ with $\alpha(\mathsf{t})>0$, as in
\eqref{4.26:eq-density-maximum-a}, the bound \eqref{4.26:eq-weight-bound-a} is equivalent to
$w_{\ast}(j)\,\alpha(\mathsf{t}_j)<0.37\,\Psi(\mathsf{t}_j)$.
\end{proof}

\begin{remark}[Existence of $N_0(j)$ and sufficiency of the maximum]
\label{4.26:rem-n0-existence}
Let $j\ge 1$ be a density, and let $R_m$ and $N_0(j)$ be as in
Notation~\ref{4.26:not-density-maximum}, so that $N_0(j)$, when it exists, is the smallest
positive integer $N$ with $L^{N-1}=R_{j-N+1}$. For $1\le N\le j$ put
\begin{equation}\label{4.26:eq-n0-existence-1}
e_N:=\log_L R_{j-N+1}-(N-1).
\end{equation}
As $N$ increases from $1$ to $j$, the density $j-N+1$ runs backwards from $j$ to $1$; we call
$e_N$ the \emph{gap} at the density $j-N+1$, and we call this passage through the densities the
\emph{backward scan}. For every density $m$ let
\begin{equation}\label{4.26:eq-n0-existence-2}
s_m:=\log_L R_m-\log_L R_{m+1}
\end{equation}
be the forward drop of $\log_L R_m$ between the densities $m$ and $m+1$.
\begin{enumerate}
\item[(a)] $N_0(j)$ exists if and only if $e_N=0$ for some $N$ with $1\le N\le j$. Moreover
$e_1=\iota\ge 0$, where $R_j=L^{\iota}$, and $e_{N+1}-e_N=s_{j-N}-1$.
\item[(b)] Suppose that the trajectory obeys the first and the second inequality of
\cite[(2.6)]{B-CMP119} at every pair of adjacent densities, and that $\mathsf{t}\ge 2r$ at the
densities involved. Then $|s_m|\le 1$. Consequently $N\mapsto e_N$ is non-increasing with steps in
$\{0,-1,-2\}$, and it passes from a positive to a negative value without taking the value $0$
exactly when, for some $N$, the gap at the density $j-N+1$ equals $1$ and
$R_{j-N+1}=L\,R_{j-N}$, that is, when the sequence $(R_m)$ rose by the factor $L$, going
forward, into the density at which the backward scan stands at gap $1$. Such a rise is allowed
by the second inequality of \cite[(2.6)]{B-CMP119} and by the first inequality of
\cite[(2.9)]{B-CMP119}; under the rule fixing $R_{m+1}$ from $R_m$ in \cite{B-CMP122a} one has
$s_m\in\{0,1\}$, and such a rise does not occur.
\item[(c)] Nothing is lost by taking the density-indexed maximum
$\gamma_0^{(k)}=5\max_{j\le k}w_\ast(j)\bar{\alpha}_j$ of
Notation~\ref{4.26:not-density-maximum} only over the densities $j\le k$ at which $N_0(j)$, and
hence $w_\ast(j)=L_0^{2N_0(j)}$, is defined.
\item[(d)] The integer
\begin{equation}\label{4.26:eq-n0-existence-3}
N_0'(j):=\min\{N:\ 1\le N\le j,\ L^{N-1}\ge R_{j-N+1}\},
\end{equation}
defined at every density at which this set is non-empty, has the following properties. Under
the hypotheses of (b) it coincides with $N_0(j)$ wherever the equality $L^{N-1}=R_{j-N+1}$ is
solvable; and Lemma~\ref{4.26:lem-weight-bound} holds verbatim with $N_0'(j)$ in place of
$N_0(j)$, the case $d\ge 2$ of its proof using the first inequality of \cite[(2.6)]{B-CMP119} at
the pair $(j-N_0'(j)+2,\,j)$.
\end{enumerate}
\end{remark}

\begin{proof}
(a) By \eqref{4.26:eq-n0-existence-1}, $e_N=0$ if and only if $\log_L R_{j-N+1}=N-1$, that is,
if and only if $L^{N-1}=R_{j-N+1}$. Since \eqref{4.26:eq-n0-existence-1} defines $e_N$ for
$1\le N\le j$, the range in which $N_0(j)$ is sought, the integer $N_0(j)$ exists if and only if
some $e_N$ with $1\le N\le j$ vanishes, and then $N_0(j)$ is the least such $N$. For $N=1$ the
definition gives $e_1=\log_L R_j=\iota$, and $\iota\ge 0$ because $R_j$ is a positive integer,
so that $L^{\iota}=R_j\ge 1$. Finally
\begin{equation*}
e_{N+1}-e_N=\log_L R_{j-N}-\log_L R_{j-N+1}-1=s_{j-N}-1
\end{equation*}
by \eqref{4.26:eq-n0-existence-2} with $m=j-N$.

(b) Let $m$ and $m+1$ be adjacent densities. The first and second inequalities of
\cite[(2.6)]{B-CMP119} at this pair, with $\mathsf{t}\ge 2r$, give
\begin{equation*}
\mathsf{t}_{m+1}^r\ge\Bigl(1-\frac{c_\beta}{2}\Bigr)\mathsf{t}_m^r,
\qquad
\mathsf{t}_{m+1}^r\le\frac{8}{7}\,\mathsf{t}_m^r,
\end{equation*}
where $c_\beta=2\ln(1+\beta_0)$; since $1-c_\beta/2>1/L$, the first of these yields
$\mathsf{t}_{m+1}^r>\mathsf{t}_m^r/L$. Each $R_n$ is an integral power of $L$ fixed by minimality,
$R_n/L<\mathsf{t}_n^r\le R_n$, exactly as at the beginning of the proof of
Lemma~\ref{4.26:lem-weight-bound}. Write $R_m=L^a$ and $R_{m+1}=L^b$, so that $s_m=a-b$. Then
\begin{equation*}
L^b\ge\mathsf{t}_{m+1}^r>\frac{\mathsf{t}_m^r}{L}>L^{a-2},
\end{equation*}
so $b\ge a-1$ and $s_m\le 1$. Also, since $8/7<L$,
\begin{equation*}
L^{b-1}<\mathsf{t}_{m+1}^r\le\frac{8}{7}\,\mathsf{t}_m^r\le\frac{8}{7}\,L^a<L^{a+1},
\end{equation*}
so $b\le a+1$ and $s_m\ge -1$. Hence $s_m\in\{-1,0,1\}$, and by (a) the increments
$e_{N+1}-e_N=s_{j-N}-1$ lie in $\{-2,-1,0\}$: the sequence $(e_N)$ is integer-valued and
non-increasing with steps in $\{0,-1,-2\}$.

Suppose that $e_N\ge 1$ and $e_{N+1}\le -1$ for some $N$. As the step is at least $-2$, this
forces $e_N=1$, $e_{N+1}=-1$ and $s_{j-N}=-1$, that is, $R_{j-N+1}=L\,R_{j-N}$: going forward
from the density $j-N$ to $j-N+1$, the sequence $(R_m)$ rose by the factor $L$, and it did so
into the density $j-N+1$, at which the backward scan stands at gap $e_N=1$. Conversely, if
$e_N=1$ and $s_{j-N}=-1$, then $e_{N'}\ge e_N=1$ for $N'\le N$ by monotonicity, $e_{N+1}=-1$, and
$e_{N'}\le -1$ for all $N'>N$; so the value $0$ is not taken. A forward rise $s_m=-1$ is
compatible with the second inequality of \cite[(2.6)]{B-CMP119} and with the first inequality of
\cite[(2.9)]{B-CMP119}; under the rule fixing $R_{m+1}$ from $R_m$ in \cite{B-CMP122a},
$s_m\in\{0,1\}$, so all increments of $e$ lie in $\{-1,0\}$, the value $0$
cannot be skipped, and such a rise is excluded.

(c) Clause (i) of Lemma~\ref{4.26:lem-analyticity-membership} reads the weight only at two kinds of
densities. The first is the current density $k$, at which Ba{\l}aban's construction
\cite[(1.10)--(1.16)]{B-CMP122a} for the operation $\mathbf{R}$ is being performed; the existence
of $N_0(k)$ is part of the premise of that construction, so $w_\ast(k)$ is defined. The second is
the density $k_a$ of a live arrest $\mathfrak{a}$, at which the same construction was performed
and its data were then frozen; there $N_0^a=N_0(k_a)$ exists by the construction of the piece
itself, so $w_\ast(k_a)$ is defined. Hence every term that clause (i) compares with
$\gamma_0^{(k)}/5$ is a term of the maximum taken over the densities at which $N_0(j)$ exists.
For clause (ii) of the same lemma, a maximum over a smaller index set is no larger than the
maximum over a larger one; so restricting the index set can only help there.

(d) By \eqref{4.26:eq-n0-existence-1}, $L^{N-1}\ge R_{j-N+1}$ holds if and only if $e_N\le 0$.
Thus $N_0'(j)$ is the least $N$ with $e_N\le 0$, while $N_0(j)$ is the least $N$ with $e_N=0$.
Suppose the hypotheses of (b) hold and the equality $L^{N-1}=R_{j-N+1}$ is solvable, so that
$N_0(j)$ exists. Since
$e_{N_0(j)}=0\le 0$, we have $N_0'(j)\le N_0(j)$. If $N_0'(j)<N_0(j)$, then
$e_{N_0'(j)}\ne 0$ by minimality of $N_0(j)$, hence $e_{N_0'(j)}<0$; by the monotonicity of (b),
$e_N<0$ for all $N\ge N_0'(j)$, in particular $e_{N_0(j)}<0$, a contradiction. Therefore
$N_0'(j)=N_0(j)$. The assertion about Lemma~\ref{4.26:lem-weight-bound} is the statement that its
argument runs verbatim with $N_0'(j)$ in place of $N_0(j)$, the case $d\ge 2$ of that argument
invoking the first inequality of \cite[(2.6)]{B-CMP119} at the pair $(j-N_0'(j)+2,\,j)$.
\end{proof}

\begin{lemma}[Earlier levels against one density]\label{4.26:prf-level-comparison}
Let the couplings $g_m$, the numbers $x_m=g_m^{-2}$ and $\mathsf{t}_m=\log g_m^{-2}$, the radii
$\alpha_{i,m}$ ($i\in\{0,1\}$), the functions $\alpha_i(\mathsf{t})=C_i\mathsf{t}^qe^{-\mathsf{t}/2}$ and
$\alpha(\mathsf{t})=C\mathsf{t}^qe^{-\mathsf{t}/2}$ with $C=\max\{C_0,C_1\}$, the maxima $\bar{\alpha}_{i,j}$
and $\bar{\alpha}_j$ over the levels $m\le j$, the number $c_\beta=2\log(1+\beta_0)$ and the
exponents $q$ and $r$ be as in Notation~\ref{4.26:not-density-maximum}. Let $j$ be a density of the
run, and assume:
\begin{itemize}
\item[(a)] the range clause of \eqref{4.26:eq-density-maximum-a}: $\mathsf{t}_m\ge\mathsf{t}_\gamma$ for
every level $m\le j$, where $\mathsf{t}_\gamma=\log\gamma^{-2}$ satisfies
$\mathsf{t}_\gamma\ge 2(r+q)+c_\beta$;
\item[(b)] for every level $m<j$, the running couplings satisfy at the pair $(m,j)$ the second inequality of
\cite[(2.6), pp.~255--256]{B-CMP119}, namely $g_m\le(1+\beta_0)g_n$ for $n>m$, taken with $n=j$.
\end{itemize}
Then, for $i\in\{0,1\}$,
\begin{equation}\label{4.26:eq-level-comparison-1}
\bar{\alpha}_{i,j}\le(1+\beta_0)\,\alpha_i(\mathsf{t}_j),
\qquad\text{and hence}\qquad
\bar{\alpha}_j\le(1+\beta_0)\,\alpha(\mathsf{t}_j).
\end{equation}
The factor $1+\beta_0$ does not depend on $j-m$.
\end{lemma}

\begin{proof}
Fix $i\in\{0,1\}$ and a level $m\le j$. If $m=j$, then $\alpha_{i,m}=\alpha_i(\mathsf{t}_j)$, and this is at most
$(1+\beta_0)\alpha_i(\mathsf{t}_j)$ because $\beta_0>0$ and $\alpha_i(\mathsf{t}_j)>0$.

Let $m<j$. Hypothesis (b) at the pair $(m,j)$ gives $g_m\le(1+\beta_0)g_j$. Squaring and inverting the
positive numbers on both sides,
\begin{equation*}
x_m=g_m^{-2}\ge\frac{g_j^{-2}}{(1+\beta_0)^2}=\frac{x_j}{(1+\beta_0)^2},
\end{equation*}
and taking logarithms,
\begin{equation}\label{4.26:eq-level-comparison-2}
\mathsf{t}_m\ge\mathsf{t}_j-2\log(1+\beta_0)=\mathsf{t}_j-c_\beta .
\end{equation}
This is a lower bound on the earlier quantity $\mathsf{t}_m$, with no term depending on $j-m$. A lower bound is
what is needed: on the range used below $\alpha_i$ is non-increasing, so a lower bound on $\mathsf{t}_m$ bounds
$\alpha_i(\mathsf{t}_m)$ from above. The first inequality of \cite[(2.6)]{B-CMP119}, which bounds $\mathsf{t}_m$
from above and carries a term $\log(j-m)$, is not used.

By (F2) of \eqref{4.26:eq-density-maximum-b}, the function
$\alpha_i(\mathsf{t})=C_i\mathsf{t}^qe^{-\mathsf{t}/2}$ is non-increasing on $[2q,\infty)$. Both arguments
that we compare lie in this interval. Indeed, hypothesis (a) gives
\begin{equation*}
\mathsf{t}_m\ge\mathsf{t}_\gamma\ge2(r+q)+c_\beta\ge 2q .
\end{equation*}
Hypothesis (a) at the level $j$ gives
\begin{equation*}
\mathsf{t}_j-c_\beta\ge\mathsf{t}_\gamma-c_\beta\ge 2(r+q)\ge 2q .
\end{equation*}
Since $\mathsf{t}_m\ge\mathsf{t}_j-c_\beta$ by \eqref{4.26:eq-level-comparison-2}, monotonicity gives
$\alpha_i(\mathsf{t}_m)\le\alpha_i(\mathsf{t}_j-c_\beta)$.

Next we apply (F1) of \eqref{4.26:eq-density-maximum-b} with shift $c_\beta$; this is admissible because
$\mathsf{t}_j\ge c_\beta\ge0$. It gives $\alpha_i(\mathsf{t}_j-c_\beta)\le e^{c_\beta/2}\alpha_i(\mathsf{t}_j)$;
for the function $\alpha_i$ this is the inequality $(\mathsf{t}_j-c_\beta)^q\le\mathsf{t}_j^q$ multiplied by
$C_ie^{-(\mathsf{t}_j-c_\beta)/2}$. Since $e^{c_\beta/2}=1+\beta_0$, the three inequalities combine to
\begin{equation*}
\alpha_{i,m}=\alpha_i(\mathsf{t}_m)\le\alpha_i(\mathsf{t}_j-c_\beta)
\le e^{c_\beta/2}\alpha_i(\mathsf{t}_j)=(1+\beta_0)\,\alpha_i(\mathsf{t}_j).
\end{equation*}

Both cases give $\alpha_{i,m}\le(1+\beta_0)\alpha_i(\mathsf{t}_j)$ for every level $m\le j$. Taking the maximum
over these levels yields the first bound of \eqref{4.26:eq-level-comparison-1}. Since $C_i\le C$, we have
$\alpha_i(\mathsf{t}_j)\le\alpha(\mathsf{t}_j)$ for $i=0,1$. Hence the bound on $\bar{\alpha}_j$ follows from the
bounds on $\bar{\alpha}_{0,j}$ and $\bar{\alpha}_{1,j}$ and the definition of $\bar{\alpha}_j$ in
Notation~\ref{4.26:not-density-maximum}.

The factor $1+\beta_0$ does not depend on $j-m$. The inequality $g_m\le(1+\beta_0)g_j$ is applied once, to the
pair $(m,j)$ itself. It is never iterated along intermediate levels, so every level $m$ is compared with $j$
directly.
\end{proof}

\begin{proof}[Proof of Theorem~\ref{4.26:thm-density-majorant}]
Recall $\Psi(\mathsf{t})=L^2(L+1)C\mathsf{t}^{r+q}e^{-\mathsf{t}/2}$ and
$\bar{\gamma}_0(\mathsf{t})=5(1+\beta_0)^2\Psi(\mathsf{t})$ from \eqref{4.26:eq-density-maximum-a}. The quantity
$\gamma_0^{(k)}$ is the maximum, over the densities $j\le k$ of the run, of the products
$w_\ast(j)\bar{\alpha}_j$.

Fix such a density $j$. Lemma~\ref{4.26:lem-weight-bound}, applied at the density $j$, gives
$w_\ast(j)<0.37\,L^2(L+1)\mathsf{t}_j^r$, which is \eqref{4.26:eq-weight-bound-a}. Its hypothesis
$\mathsf{t}_j\ge2r$ holds by the range clause, since
\begin{equation*}
\mathsf{t}_j\ge\mathsf{t}_\gamma\ge2(r+q)+c_\beta\ge2r .
\end{equation*}
The remaining hypotheses of Lemma~\ref{4.26:lem-weight-bound} at the density $j$, and the second
inequality of \cite[(2.6)]{B-CMP119} at the pairs of levels used below, are hypotheses of
Theorem~\ref{4.26:thm-density-majorant}.
Lemma~\ref{4.26:prf-level-comparison} gives $\bar{\alpha}_j\le(1+\beta_0)C\mathsf{t}_j^qe^{-\mathsf{t}_j/2}$.
Both factors are positive, so
\begin{equation}\label{4.26:eq-level-comparison-3}
w_\ast(j)\,\bar{\alpha}_j<0.37\,(1+\beta_0)\,L^2(L+1)\,C\,\mathsf{t}_j^{r+q}e^{-\mathsf{t}_j/2}
=0.37\,(1+\beta_0)\,\Psi(\mathsf{t}_j).
\end{equation}

By (F2) of \eqref{4.26:eq-density-maximum-b}, $\Psi$ is non-increasing on $[2(r+q),\infty)$. We now compare
$\Psi(\mathsf{t}_j)$ with $\Psi(\mathsf{t}_k)$.

If $j=k$, then \eqref{4.26:eq-level-comparison-3} gives
\begin{equation*}
w_\ast(k)\bar{\alpha}_k<0.37(1+\beta_0)\Psi(\mathsf{t}_k)\le0.37(1+\beta_0)^2\Psi(\mathsf{t}_k).
\end{equation*}

If $j<k$, we apply the second inequality of \cite[(2.6), pp.~255--256]{B-CMP119} once, to the pair $(j,k)$.
This pair is different from the pairs $(m,j)$ used in Lemma~\ref{4.26:prf-level-comparison}. As in the passage
to \eqref{4.26:eq-level-comparison-2}, it gives $\mathsf{t}_j\ge\mathsf{t}_k-c_\beta$. By the range clause at the
level $k$, $\mathsf{t}_k-c_\beta\ge\mathsf{t}_\gamma-c_\beta\ge2(r+q)$, so both $\mathsf{t}_j$ and
$\mathsf{t}_k-c_\beta$ lie in $[2(r+q),\infty)$. A lower bound on the argument of a non-increasing function is an
upper bound on its value, so $\Psi(\mathsf{t}_j)\le\Psi(\mathsf{t}_k-c_\beta)$. Then (F1) of
\eqref{4.26:eq-density-maximum-b} with shift $c_\beta\le\mathsf{t}_k$ is an elementary inequality, the
inequality $(\mathsf{t}_k-c_\beta)^{r+q}\le\mathsf{t}_k^{r+q}$ multiplied by the remaining factors. It gives
\begin{equation*}
\Psi(\mathsf{t}_j)\le\Psi(\mathsf{t}_k-c_\beta)\le e^{c_\beta/2}\Psi(\mathsf{t}_k)=(1+\beta_0)\Psi(\mathsf{t}_k),
\end{equation*}
and with \eqref{4.26:eq-level-comparison-3},
$w_\ast(j)\bar{\alpha}_j<0.37(1+\beta_0)^2\Psi(\mathsf{t}_k)$.

Every term of the finite maximum defining $\gamma_0^{(k)}$ is therefore less than
$0.37(1+\beta_0)^2\Psi(\mathsf{t}_k)$. Hence
\begin{equation*}
\gamma_0^{(k)}<0.37(1+\beta_0)^2\Psi(\mathsf{t}_k)\le0.37\,\bar{\gamma}_0(\mathsf{t}_k)
<5(1+\beta_0)^2\Psi(\mathsf{t}_k)=\bar{\gamma}_0(\mathsf{t}_k),
\end{equation*}
which is the corresponding display.

The maximum must be bounded as a maximum of products, not factor by factor. The majorant
$0.37\,L^2(L+1)\mathsf{t}_j^r$ of $w_\ast(j)$ increases in $\mathsf{t}_j$. Bounding $\max_jw_\ast(j)$ alone by a
function of $\mathsf{t}_k$ would therefore need an upper bound on $\mathsf{t}_j$ in terms of $\mathsf{t}_k$, that is,
the first inequality of \cite[(2.6)]{B-CMP119} in the form
$\mathsf{t}_j\le\mathsf{t}_k+c_\beta+\log(k-j)$, which depends on $k-j$. The product majorant
$0.37(1+\beta_0)\Psi(\mathsf{t}_j)$ is non-increasing in $\mathsf{t}_j$, so only the lower bound
$\mathsf{t}_j\ge\mathsf{t}_k-c_\beta$ is needed. Of the two factors $1+\beta_0$ in $\bar{\gamma}_0$, one comes from
comparing the levels $m\le j$ with $\mathsf{t}_j$ (Lemma~\ref{4.26:prf-level-comparison}). The other comes from
comparing $\mathsf{t}_j$ with $\mathsf{t}_k$.

\emph{The range clause.} Both monotonicity ranges used above follow from the one condition
$\mathsf{t}_j\ge\mathsf{t}_\gamma$ for every density in play, with $\mathsf{t}_\gamma\ge2(r+q)+c_\beta$, of
\eqref{4.26:eq-density-maximum-a}. These are the ranges $[2q,\infty)$ for $\alpha_i$ and $[2(r+q),\infty)$ for
$\Psi$, each required at the argument shifted by $c_\beta$. The same condition gives the hypothesis
$\mathsf{t}_j\ge2r$ of Lemma~\ref{4.26:lem-weight-bound}, since $2(r+q)+c_\beta$ dominates $2q+c_\beta$ and $2r$.

The condition is one-sided in $\gamma$, where $\mathsf{t}_\gamma=\log\gamma^{-2}$; equivalently,
$\gamma\le e^{-(r+q)}/(1+\beta_0)$. It involves only $r$, $q$ and $\beta_0$, it is fixed by these
before $L$ is chosen and involves no degree function, and it is an upper bound on the couplings.
Under it, $\mathsf{t}_k\ge\mathsf{t}_\gamma\ge2(r+q)$ for
every $k$. Since $\bar{\gamma}_0=5(1+\beta_0)^2\Psi$ is non-increasing on $[2(r+q),\infty)$, it follows that
$\bar{\gamma}_0(\mathsf{t}_k)\le\bar{\gamma}_0(\mathsf{t}_\gamma)$ for every $k$.

\emph{Non-increasing envelopes.} Put
\begin{equation*}
\alpha_i^{\downarrow}(\mathsf{t}):=\sup_{\mathsf{t}'\ge\mathsf{t}}\alpha_i(\mathsf{t}'),
\qquad
\alpha^{\downarrow}(\mathsf{t}):=\sup_{\mathsf{t}'\ge\mathsf{t}}\alpha(\mathsf{t}'),
\qquad
\Psi^{\downarrow}(\mathsf{t}):=\sup_{\mathsf{t}'\ge\mathsf{t}}\Psi(\mathsf{t}').
\end{equation*}
These functions are non-increasing by construction, and they coincide with $\alpha_i$, $\alpha$ and $\Psi$ on
the ranges above. For $\mathsf{t}\ge c\ge0$ they satisfy the shift inequality
\begin{equation}\label{4.26:eq-level-comparison-4}
\Psi^{\downarrow}(\mathsf{t}-c)\le e^{c/2}\,\Psi^{\downarrow}(\mathsf{t}),
\end{equation}
and likewise for $\alpha_i^{\downarrow}$.

To prove \eqref{4.26:eq-level-comparison-4}, let $\mathsf{t}'\ge\mathsf{t}-c$ and put
$\mathsf{t}''=\max\{\mathsf{t}',\mathsf{t}\}$. Then $0\le\mathsf{t}'\le\mathsf{t}''$, so
$\mathsf{t}'^{\,r+q}\le\mathsf{t}''^{\,r+q}$. Moreover $\mathsf{t}''-\mathsf{t}'\le c$ and $\mathsf{t}''\ge\mathsf{t}$.
Therefore
\begin{equation*}
\Psi(\mathsf{t}')\le e^{(\mathsf{t}''-\mathsf{t}')/2}\Psi(\mathsf{t}'')\le e^{c/2}\Psi^{\downarrow}(\mathsf{t}),
\end{equation*}
and taking the supremum over $\mathsf{t}'\ge\mathsf{t}-c$ gives \eqref{4.26:eq-level-comparison-4}. The same
argument with the exponent $q$ in place of $r+q$ gives the inequality for $\alpha_i^{\downarrow}$.

With the envelopes, the two proofs above run without the range clause, as follows.
\begin{itemize}
\item For $m<j$: $\alpha_{i,m}\le\alpha_i^{\downarrow}(\mathsf{t}_m)\le\alpha_i^{\downarrow}(\mathsf{t}_j-c_\beta)$,
and then $\alpha_{i,m}\le(1+\beta_0)\alpha_i^{\downarrow}(\mathsf{t}_j)$ by
\eqref{4.26:eq-level-comparison-4} for $\alpha_i^{\downarrow}$, used when $\mathsf{t}_j\ge c_\beta$.
\item Since $C_i\le C$, we have $\alpha_i\le\alpha$ and hence $\alpha_i^{\downarrow}\le\alpha^{\downarrow}$.
Since $\Psi(\mathsf{t}')=L^2(L+1)\mathsf{t}'^{\,r}\alpha(\mathsf{t}')$ and $\mathsf{t}_j^r\le\mathsf{t}'^{\,r}$
for $\mathsf{t}'\ge\mathsf{t}_j$, it follows that
\begin{equation*}
\mathsf{t}_j^r\,\alpha^{\downarrow}(\mathsf{t}_j)\le\frac{\Psi^{\downarrow}(\mathsf{t}_j)}{L^2(L+1)}.
\end{equation*}
\item For $j<k$: $\Psi^{\downarrow}(\mathsf{t}_j)\le\Psi^{\downarrow}(\mathsf{t}_k-c_\beta)$, and then
$\Psi^{\downarrow}(\mathsf{t}_j)\le(1+\beta_0)\Psi^{\downarrow}(\mathsf{t}_k)$ by
\eqref{4.26:eq-level-comparison-4}, used when $\mathsf{t}_k\ge c_\beta$.
\end{itemize}
The result is $\gamma_0^{(k)}<5(1+\beta_0)^2\Psi^{\downarrow}(\mathsf{t}_k)$, with the threshold of the range
clause weakened to $\mathsf{t}_\gamma\ge\max\{2r,c_\beta\}$ and all other hypotheses unchanged. The part
$\mathsf{t}_\gamma\ge2r$ cannot be dropped, since it supplies the
hypothesis $\mathsf{t}_j\ge2r$ of Lemma~\ref{4.26:lem-weight-bound}. In this chapter the bound is used in the
corresponding form, with $\Psi$ itself and the range clause
$\mathsf{t}_\gamma\ge2(r+q)+c_\beta$.
\end{proof}

\begin{remark}[The coupling flow in place of Ba{\l}aban's running-coupling inequalities]
\label{4.26:rem-flow-substitute}
Let $x_j=g_j^{-2}$, $\mathsf{t}_j=\log g_j^{-2}$, $\beta_0=1/7$ and $c_\beta=2\log(1+\beta_0)$ be as
in Notation~\ref{4.26:not-density-maximum} and \eqref{4.26:eq-density-maximum-a}. Here $\log$ is the
natural logarithm, so that $e^{c_\beta}=64/49$. Let $\mathsf{t}_\gamma=\log\gamma^{-2}$, and let (g)
be the lower bound $\mathsf{t}_j\ge\mathsf{t}_\gamma$ of \eqref{4.26:eq-density-maximum-a}. Clause~(0)
of Theorem~0 (Theorem~\ref{4.1:thm-clause-zero}) gives, for all levels $0\le m<n\le K$, the two-sided
flow bounds
\begin{equation}\label{4.26:eq-flow-substitute-1}
\lambda_\sharp(n-m)-c_5\;\le\;x_m-x_n\;\le\;B_\sharp(n-m).
\end{equation}
Assume that $\gamma$ satisfies the two one-sided ceilings
\begin{equation}\label{4.26:eq-flow-substitute-2}
\gamma^2\le\frac{15}{64c_5},\qquad B_\sharp\gamma^2\le\frac{15}{49}.
\end{equation}
Then the following hold for all levels $m<n$ such that (g) holds at $n$:
\begin{itemize}
\item[(a)] $\mathsf{t}_m\ge\mathsf{t}_n-c_\beta$;
\item[(b)] $\mathsf{t}_m\le\mathsf{t}_n+c_\beta+\log(n-m)$.
\end{itemize}
Inequality (a) is the logarithmic form of the second inequality of \cite[(2.6)]{B-CMP119},
$g_m\le(1+\beta_0)g_n$ for $n>m$. Inequality (b) is the logarithmic form of the first,
$g_n\le(1+\beta_0)(n-m)^{1/2}g_m$.

These inequalities have two consequences. First, the trajectory-free bound on the weight
(Lemma~\ref{4.26:lem-weight-bound}) uses the first inequality of \cite[(2.6)]{B-CMP119} exactly once,
and (b) supplies it. Second, the comparison of earlier levels with one density
(Lemma~\ref{4.26:prf-level-comparison}), and the ensuing bound on the density-indexed maximum
$\gamma_0^{(k)}$ of
Notation~\ref{4.26:not-density-maximum}, use the second inequality of \cite[(2.6)]{B-CMP119}, and (a)
supplies that inequality in the same form. Hence item~(ii) of the bounds on the density-indexed
maximum $\gamma_0^{(k)}$, whose items are (i)--(v), holds as an implication from clause~(0) alone
under the two ceilings
\eqref{4.26:eq-flow-substitute-2}, and \cite[(2.6)]{B-CMP119} becomes a consistency check.

For $\mathsf{t}\ge c_\beta$ the inequalities
\begin{equation}\label{4.26:eq-flow-substitute-3}
\Psi(\mathsf{t}-c_\beta)\le(1+\beta_0)\Psi(\mathsf{t}),\qquad
\alpha_i(\mathsf{t}-c_\beta)\le(1+\beta_0)\alpha_i(\mathsf{t})
\end{equation}
are properties of the functions $\Psi$ of \eqref{4.26:eq-density-maximum-a} and $\alpha_i$ of
Notation~\ref{4.26:not-density-maximum}. What uses (a) is something else: the substitution of lower
bounds on the $\mathsf{t}_j$ into these functions where they are non-increasing, which bounds their
values from above. The inequality $\mathsf{t}_m\ge\mathsf{t}_j$ would be the same substitution with
$c_\beta=0$; it is not supplied by \cite[(2.6)]{B-CMP119} and is not used.

Suppose that (g) is not assumed at every level and only $g_K=g\le\gamma$ is assumed. Then
$\mathsf{t}_j\ge\mathsf{t}_\gamma^{\flat}$ for all $j\le K$, where
$\mathsf{t}_\gamma^{\flat}:=\mathsf{t}_\gamma-c_\beta$. This is a shift of the threshold, not a
factor. The range clause of \eqref{4.26:eq-density-maximum-a} is then replaced by
$\mathsf{t}_\gamma\ge 2(r+q)+2c_\beta$. The ceiling on $\gamma$ that bounds $600$ times a multiple
of $\Psi(\mathsf{t}_\gamma)$ by $1$, and within which the range clause is counted, is replaced by
$600(1+\beta_0)^3\Psi(\mathsf{t}_\gamma)\le 1$. Alternatively, $\gamma$ may be
replaced by $\gamma/(1+\beta_0)$.
\end{remark}

\begin{proof}
\emph{Inequality} (a). By the left-hand inequality of \eqref{4.26:eq-flow-substitute-1} and
$\lambda_\sharp(n-m)\ge0$,
\begin{equation*}
x_m\;\ge\;x_n-c_5\;=\;x_n\Bigl(1-\frac{c_5}{x_n}\Bigr).
\end{equation*}
By (g) at $n$ we have $x_n\ge\gamma^{-2}$, so that $c_5/x_n\le c_5\gamma^2$. The first ceiling in
\eqref{4.26:eq-flow-substitute-2} therefore gives $c_5/x_n\le 15/64$. Since $\beta_0=1/7$,
\begin{equation*}
x_m\;\ge\;\frac{49}{64}\,x_n\;=\;(1+\beta_0)^{-2}x_n .
\end{equation*}
Taking logarithms gives $\mathsf{t}_m\ge\mathsf{t}_n-2\log(1+\beta_0)=\mathsf{t}_n-c_\beta$. This is
the inequality that the second inequality of \cite[(2.6)]{B-CMP119} yields at the pair $(m,n)$.
Lemma~\ref{4.26:prf-level-comparison} and the ensuing bound on $\gamma_0^{(k)}$ use that
inequality only in this logarithmic form. Their arguments therefore go through unchanged with (a) in
its place.

\emph{Inequality} (b). By the right-hand inequality of \eqref{4.26:eq-flow-substitute-1},
\begin{equation*}
x_m\;\le\;x_n+B_\sharp(n-m)\;=\;x_n\Bigl(1+\frac{B_\sharp(n-m)}{x_n}\Bigr).
\end{equation*}
Hence $\mathsf{t}_m\le\mathsf{t}_n+\log\bigl(1+B_\sharp(n-m)/x_n\bigr)$. By (g) at $n$ we have
$B_\sharp(n-m)/x_n\le B_\sharp\gamma^2(n-m)$. For $a\ge0$ and an integer $n-m\ge1$,
\begin{equation*}
(1+a)(n-m)-\bigl(1+a(n-m)\bigr)=n-m-1\ge0 .
\end{equation*}
Applying this with $a=B_\sharp\gamma^2$ gives
\begin{equation*}
\mathsf{t}_m\;\le\;\mathsf{t}_n+\log\bigl(1+B_\sharp\gamma^2\bigr)+\log(n-m)
\;\le\;\mathsf{t}_n+c_\beta+\log(n-m).
\end{equation*}
The last inequality is the second ceiling in \eqref{4.26:eq-flow-substitute-2}. Indeed,
$B_\sharp\gamma^2\le 15/49$ gives $1+B_\sharp\gamma^2\le 64/49=e^{c_\beta}$. Inequality (b) is the
inequality that the first inequality of \cite[(2.6)]{B-CMP119} yields at the pair $(m,n)$. The one use
of that inequality in Lemma~\ref{4.26:lem-weight-bound} is therefore supplied by clause~(0). Together
with the previous paragraph, this proves the assertion that item~(ii) holds as an implication from
clause~(0) under \eqref{4.26:eq-flow-substitute-2}.

\emph{The properties of $\Psi$ and $\alpha_i$.} For $\mathsf{t}\ge c_\beta$ and $s\ge0$ one has
$(\mathsf{t}-c_\beta)^s\le\mathsf{t}^s$ and
$e^{-(\mathsf{t}-c_\beta)/2}=e^{c_\beta/2}e^{-\mathsf{t}/2}=(1+\beta_0)e^{-\mathsf{t}/2}$; with
$s=r+q$ and $s=q$ this gives \eqref{4.26:eq-flow-substitute-3}. Moreover
$(\log\alpha_i)'(\mathsf{t})=q/\mathsf{t}-1/2\le0$ on $[2q,\infty)$, and by the same computation with
exponent $r+q$ the function $\Psi$ is non-increasing on $[2(r+q),\infty)$.

\emph{The standing of} (g). The condition (g) is the lower bound $x_j=g_j^{-2}\ge\gamma^{-2}$, that is
$\mathsf{t}_j\ge\mathsf{t}_\gamma$, at every density $j$ under consideration. It is an upper bound on
the couplings only. Suppose only $g_K=g\le\gamma$ is assumed.
Then the second inequality of \cite[(2.6)]{B-CMP119} at the pair $(j,K)$ gives
$\mathsf{t}_j\ge\mathsf{t}_\gamma-c_\beta=\mathsf{t}_\gamma^{\flat}$ for all $j\le K$. The same
conclusion follows from (a) at $n=K$, whose ceiling condition uses only $x_K\ge\gamma^{-2}$. Every
threshold then shifts once by $c_\beta$.

The range clause, imposed at $\mathsf{t}_\gamma^{\flat}$, reads
$\mathsf{t}_\gamma-c_\beta\ge 2(r+q)+c_\beta$, that is, $\mathsf{t}_\gamma\ge 2(r+q)+2c_\beta$. The
ceiling with the factor $600$ described in the remark is read at $\mathsf{t}_\gamma^{\flat}$; the
shift of the argument of
$\Psi$ by $c_\beta$ is paid for by the first inequality of \eqref{4.26:eq-flow-substitute-3},
$\Psi(\mathsf{t}_\gamma-c_\beta)\le(1+\beta_0)\Psi(\mathsf{t}_\gamma)$, and the condition becomes
\begin{equation*}
600(1+\beta_0)^3\Psi(\mathsf{t}_\gamma)\le 1 .
\end{equation*}
Alternatively, assume $g\le\gamma/(1+\beta_0)$. Then
$\mathsf{t}_K\ge\log\bigl((1+\beta_0)^2\gamma^{-2}\bigr)=\mathsf{t}_\gamma+c_\beta$, and the same
argument gives $\mathsf{t}_j\ge\mathsf{t}_\gamma$ for all $j\le K$, which is (g) itself.

In items (i)--(v) of the bounds on the density-indexed maximum $\gamma_0^{(k)}$, (g) enters only as
a lower threshold on the
$\mathsf{t}_j$ and as the left endpoint of suprema over the single variable $\mathsf{t}$; item~(i) does
not use (g) at all. Both uses are preserved under a shift of the threshold, so nothing else changes.
\end{proof}

\begin{proposition}[Every use of the maximum is a membership or a one-variable supremum]
\label{4.26:prop-maximum-uses}
Let $\gamma>0$ and $\mathsf{t}_\gamma:=\log\gamma^{-2}$, and assume $\mathsf{t}_\gamma\ge 2(r+q)+c_\beta$.
Let $k$ be a level of the run to which Theorem~\ref{4.26:thm-density-majorant} applies and at which
$\mathsf{t}_k\ge\mathsf{t}_\gamma$. Let $\gamma_0^{(k)}$, $\Psi$, $\bar{\gamma}_0$, $c_\beta$ and
$\bar{N}(\mathsf{t})=L\mathsf{t}^r$ be as in Notation~\ref{4.26:not-density-maximum}
(displays \eqref{4.26:eq-density-maximum-a} and \eqref{4.26:eq-density-maximum-c}).
Let $\vartheta(\mathsf{t})$ denote the smallness parameter of the expansion at a coupling $g'$ with
$\log g'^{-2}=\mathsf{t}$, and assume that there are $\eta>0$ and an exponent $p_1<q$, fixed with
the auxiliary parameters, such that $\vartheta(\mathsf{t})\ge \eta\,\mathsf{t}^{p_1-q}$ for all
$\mathsf{t}\ge\mathsf{t}_\gamma$.
\begin{enumerate}
\item[(a)] Let $P$ be a condition on a positive number $\gamma_0$ such that $P(\gamma_0)$ implies
$P(\gamma_0')$ for every $\gamma_0'\le\gamma_0$. If $P(\bar{\gamma}_0(\mathsf{t}_\gamma))$ holds, then
$P(\gamma_0^{(k)})$ holds.
\item[(b)] Let $A_1,A_2>0$ be constants independent of $\mathsf{t}$, of $k$ and of the sequence
$(R_j)_{j\le k}$. If
\begin{equation}\label{4.26:eq-maximum-uses-1}
\sup_{\mathsf{t}\ge\mathsf{t}_\gamma}\frac{A_1\,\bar{\gamma}_0(\mathsf{t})}{\vartheta(\mathsf{t})}\le 1
\qquad\text{and}\qquad
\sup_{\mathsf{t}\ge\mathsf{t}_\gamma}
\frac{A_2\,\bar{N}(\mathsf{t})\,\bar{\gamma}_0(\mathsf{t})^2}{\vartheta(\mathsf{t})}\le 1,
\end{equation}
then $A_1\gamma_0^{(k)}\le\vartheta(\mathsf{t}_k)$ and $A_2R_k(\gamma_0^{(k)})^2\le\vartheta(\mathsf{t}_k)$,
where $R_k$ is the integer of \cite[(2.5)]{B-CMP119} at level $k$.
Both suprema in \eqref{4.26:eq-maximum-uses-1} tend to $0$ as $\mathsf{t}_\gamma\to\infty$ faster than
every power of $\mathsf{t}_\gamma^{-1}$. We say that a condition on $\gamma$ has degree $-\infty$ if it is
implied by an inequality $f(\mathsf{t}_\gamma)\le 1$ with $f(\mathsf{t})\to0$ as $\mathsf{t}\to\infty$
faster than every power of $\mathsf{t}^{-1}$; in particular the two conditions
\eqref{4.26:eq-maximum-uses-1} have degree $-\infty$.
\item[(c)] Every condition on the size $\gamma_0$ of the analyticity polydisc
$\mathfrak{K}_{109}(\mathfrak{B};\gamma_0)$ imposed in the construction of this chapter is either a
membership or a smallness condition monotone in $\gamma_0$, which holds at $\gamma_0=\gamma_0^{(k)}$ by
Lemma~\ref{4.26:lem-analyticity-membership} or, by \textup{(a)}, once it is imposed at
$\gamma_0=\bar{\gamma}_0(\mathsf{t}_\gamma)$, a condition of degree $-\infty$ on $\gamma$; or a condition
in which $\gamma_0$ is
compared with $\vartheta(\mathsf{t}_k)$ and which is implied by a supremum over the one variable
$\mathsf{t}$ as in \textup{(b)}. None of them imposes a lower bound on the couplings $g_k$.
\end{enumerate}
\end{proposition}

\begin{proof}
\emph{Monotonicity of $\bar{\gamma}_0$.} By \eqref{4.26:eq-density-maximum-a},
$\bar{\gamma}_0(\mathsf{t})=5(1+\beta_0)^2\Psi(\mathsf{t})$ with
$\Psi(\mathsf{t})=L^2(L+1)C\mathsf{t}^{r+q}e^{-\mathsf{t}/2}$. Hence
\begin{equation*}
\frac{d}{d\mathsf{t}}\log\bar{\gamma}_0(\mathsf{t})=\frac{r+q}{\mathsf{t}}-\frac12\le 0
\qquad\text{for }\mathsf{t}\ge 2(r+q),
\end{equation*}
so, since $\mathsf{t}_\gamma\ge 2(r+q)+c_\beta\ge 2(r+q)$, $\bar{\gamma}_0$ is non-increasing on
$[\mathsf{t}_\gamma,\infty)$.

\emph{Proof of \textup{(a)}.} By Theorem~\ref{4.26:thm-density-majorant}, in the corresponding form, and by the monotonicity just shown, applied with
$\mathsf{t}_k\ge\mathsf{t}_\gamma$,
\begin{equation}\label{4.26:eq-maximum-uses-2}
\gamma_0^{(k)}\le\bar{\gamma}_0(\mathsf{t}_k)\le\bar{\gamma}_0(\mathsf{t}_\gamma).
\end{equation}
Since $P$ passes from a value of $\gamma_0$ to every smaller value, $P(\bar{\gamma}_0(\mathsf{t}_\gamma))$
implies $P(\gamma_0^{(k)})$.

\emph{Proof of \textup{(b)}.} By \eqref{4.26:eq-maximum-uses-2} and the first inequality of
\eqref{4.26:eq-maximum-uses-1} evaluated at $\mathsf{t}=\mathsf{t}_k\ge\mathsf{t}_\gamma$,
$A_1\gamma_0^{(k)}\le A_1\bar{\gamma}_0(\mathsf{t}_k)\le\vartheta(\mathsf{t}_k)$. For the second, write
$R_k=L^{\iota}$; by the minimality of $R_k$ in \cite[(2.5)]{B-CMP119} one has
$L^{\iota-1}<\mathsf{t}_k^r$, hence $R_k<L\mathsf{t}_k^r=\bar{N}(\mathsf{t}_k)$. With
\eqref{4.26:eq-maximum-uses-2} and the second inequality of \eqref{4.26:eq-maximum-uses-1} at
$\mathsf{t}=\mathsf{t}_k$,
\begin{equation*}
A_2R_k\bigl(\gamma_0^{(k)}\bigr)^2\le A_2\bar{N}(\mathsf{t}_k)\,\bar{\gamma}_0(\mathsf{t}_k)^2
\le\vartheta(\mathsf{t}_k).
\end{equation*}
Every factor is here evaluated at the same $\mathsf{t}_k$. Since, by hypothesis,
$\vartheta(\mathsf{t})\ge \eta\,\mathsf{t}^{p_1-q}$ for $\mathsf{t}\ge\mathsf{t}_\gamma$, the two
profiles in \eqref{4.26:eq-maximum-uses-1} are
bounded by constants times
\begin{equation*}
\mathsf{t}^{r+2q-p_1}e^{-\mathsf{t}/2}
\qquad\text{and}\qquad
\mathsf{t}^{3r+3q-p_1}e^{-\mathsf{t}},
\end{equation*}
respectively, the exponents $r+2q-p_1$ and $3r+3q-p_1$ being positive because $q>p_1$. The first
function is decreasing for $\mathsf{t}\ge 2(r+2q-p_1)$, the second for $\mathsf{t}\ge 3r+3q-p_1$; these
thresholds depend on $r,q,p_1$ only and are fixed before $\gamma$. Beyond them each supremum over
$\mathsf{t}\ge\mathsf{t}_\gamma$ is at most a constant times the value of the function at
$\mathsf{t}_\gamma$, which tends to $0$ as $\mathsf{t}_\gamma\to\infty$ faster than every power of
$\mathsf{t}_\gamma^{-1}$. Hence both inequalities \eqref{4.26:eq-maximum-uses-1} hold for all
sufficiently small $\gamma$, a condition on $\gamma$ given the parameters chosen before it.

\emph{Proof of \textup{(c)}.} The conditions on $\gamma_0$ imposed in the construction of this chapter
are of the following kinds.
First, memberships of a background pair $(\mathbb{U},\mathbb{J})$ in the polydisc
$\mathfrak{K}_{109}(\mathfrak{B};\gamma_0)$, required at $\gamma_0=\gamma_0^{(k)}$; these are proved at
$\gamma_0^{(k)}$ by definition in Lemma~\ref{4.26:lem-analyticity-membership}. Second, estimates valid on
$\mathfrak{K}_{109}(\mathfrak{B};\gamma_0)$ whose right sides are non-decreasing in $\gamma_0$; they are
used at $\gamma_0^{(k)}$ and impose nothing. Third, smallness conditions $c\gamma_0\le c'$ with constants
$c,c'$ independent of $\mathsf{t}$, of $k$ and of the trajectory; these are of the kind in \textup{(a)}
and are required at $\gamma_0=\bar{\gamma}_0(\mathsf{t}_\gamma)$, that is at the single value
$\mathsf{t}=\mathsf{t}_\gamma$, where, by \eqref{4.26:eq-density-maximum-a}, their left sides are
bounded by a constant times $\mathsf{t}_\gamma^{r+q}e^{-\mathsf{t}_\gamma/2}$. Fourth, the two
conditions in which $\gamma_0$ is compared with $\vartheta(\mathsf{t}_k)$: one linear in $\gamma_0$, of the
form $A_1\gamma_0\le\vartheta(\mathsf{t}_k)$, and one quadratic in $\gamma_0$ carrying the factor
$R_k$, of the form $A_2R_k\gamma_0^2\le\vartheta(\mathsf{t}_k)$; by \textup{(b)} both hold at
$\gamma_0=\gamma_0^{(k)}$ once \eqref{4.26:eq-maximum-uses-1} holds. These two are the conditions of
the construction of this chapter in which $\gamma_0$ meets $\vartheta(\mathsf{t}_k)$, and
$A_1$, $A_2$ denote the constants that multiply $\gamma_0$ and $R_k\gamma_0^2$ in them. In none of
these kinds is a value of
$\mathsf{t}_k$ bounded from above, so no lower bound on $g_k$ arises.
\end{proof}

If $\gamma_0$ were a single number fixed for the whole run, independent of $k$, the fourth kind of
condition would compare it with $\vartheta(\mathsf{t}_k)$, which, for $\vartheta(\mathsf{t})$ of the
order of $\mathsf{t}^{p_1-q}$, tends to $0$ as $\mathsf{t}_k\to\infty$ because $q>p_1$, and, in the
quadratic condition, would multiply it by $R_k$, which grows like $\mathsf{t}_k^r$; such a condition
fails once $\mathsf{t}_k^{q-p_1}$ is of the order of $1/\gamma_0$, and would be a lower bound on the
couplings $g_k$, that is a bound on the number $K$ of steps. With $\gamma_0=\gamma_0^{(k)}$ part~(b)
applies, no such bound arises, and no condition on the parameters enters beyond the two inequalities
\eqref{4.26:eq-maximum-uses-1}.

The weight $w_\ast(j)=L_0^{2N_0(j)}$ is defined at the density $j$, through the integer $N_0(j)$ of
Notation~\ref{4.26:not-density-maximum}; it depends on the
sequence $(R_j)$ and not on the number $\mathsf{t}_j$ alone. $N_0(j)$ is not determined by
$\mathsf{t}_j$: densities with equal values of $\mathsf{t}$ can carry values of $N_0$ differing by one,
as on the trajectory along which the density-indexed
maximum is attained at every density of an interval below the current level. Accordingly no
expression $w_\ast(\mathsf{t})$, and no supremum over $\mathsf{t}$ of $w_\ast(\mathsf{t})$ times a radius,
is used in this section. Where a function of $\mathsf{t}$ is needed, Lemma~\ref{4.26:lem-weight-bound}
supplies, under its hypotheses, the bound \eqref{4.26:eq-weight-bound-a},
$w_\ast(j)<0.37\,L^2(L+1)\mathsf{t}_j^r$. The integer $N_0(j)$ enters only through upper bounds derived
from its defining equality; no bound on it is assumed. In particular the condition
$(N_0+1)(N_0+3)\le N$ at level $k$, with $N_0=N_0(k)$, is read with $N=R_k$, the integer of
\cite[(2.5)]{B-CMP119} determined by $\mathsf{t}_k$ (in particular $R_k\ge\mathsf{t}_k^r$), and $N_0$
majorised from its defining equality as in the proof of
Lemma~\ref{4.26:lem-weight-bound}; it is a condition in the one variable $\mathsf{t}_k$.

\begin{remark}[No new condition, floor or ceiling]\label{4.26:rem-no-new-condition}
Let $\gamma_0^{(k)}$ be the density-indexed maximum and $\mathsf{t}_k=\log g_k^{-2}$, as in
Notation~\ref{4.26:not-density-maximum}. Let $\bar{\gamma}_0$, $\Psi$, $c_\beta$ and the exponents
$r$, $q$ of $\Psi$ be as in
\eqref{4.26:eq-density-maximum-a}, with $\beta_0$ as in Theorem~\ref{4.26:thm-density-majorant}
and $C=\max\{C_0,C_1\}$ as in Notation~\ref{4.26:not-density-maximum}, and let
$\mathsf{t}_\gamma=\log\gamma^{-2}$. Throughout, the size $\gamma_0$ of the analyticity polydisc
$\mathfrak{K}_{109}(\mathfrak{B};\gamma_0)$, built on the determining-set structure $\mathfrak{B}$
of the current stage, is taken to be $\gamma_0^{(k)}$, and, for a density $j\le k$ of the run,
$N_0$ is taken to be the integer $N_0(j)$ of Notation~\ref{4.26:not-density-maximum}. Compared
with the formulation in which $\gamma_0$ is a fixed
constant, the list of conditions on the parameters under which the constructions of this chapter
are carried out stands as follows.
\begin{enumerate}
\item[(a)] The list has the same number of members as in the formulation with fixed $\gamma_0$;
no member is added. The member which there bounds
$600\,\sup_{\mathsf{t}}w_\ast(\mathsf{t})\bar{\alpha}(\mathsf{t})$ by $1$, for the weight
$w_\ast$ and the maximal radius $\bar{\alpha}$ of Notation~\ref{4.26:not-density-maximum}, takes,
at its place in the list, the form of the single member
\begin{equation}\label{4.26:eq-no-new-condition-1}
\mathsf{t}_\gamma\ge 2(r+q)+c_\beta
\quad\text{and}\quad
600\,(1+\beta_0)^2\,\Psi(\mathsf{t}_\gamma)\le 1 ,
\end{equation}
which is a conjunction of two conditions. Two further members of the list are restated in a
changed form, each at its own place in the list; this neither adds nor removes a member.
\item[(b)] Three smallness conditions of the list are stated in terms of $\gamma_0$: the condition
written $K_{\Delta}\gamma_0\le 1/2$ in the list, with $K_{\Delta}$ the constant by which
$\gamma_0$ is multiplied there, and two further ceilings of the list in which $\gamma_0$ enters.
Each of the three is imposed with $\gamma_0$ replaced by $\bar{\gamma}_0(\mathsf{t}_\gamma)$; in
particular the first is imposed as $K_{\Delta}\bar{\gamma}_0(\mathsf{t}_\gamma)\le 1/2$.
\item[(c)] The member $(N_0+1)(N_0+3)\le R_k$, where $R_k$ is the integer at level $k$ that enters
the defining equality of $N_0(j)$ in Notation~\ref{4.26:not-density-maximum}, denoted $N$ in
Ba{\l}aban's construction, belongs to the list in either formulation. It depends on $N_0$ only
through a majorant obtained from the defining equality of $N_0(j)$. Indexing $N_0$ by the
density, as in Notation~\ref{4.26:not-density-maximum}, neither creates nor worsens this member.
\item[(d)] With $\gamma_0=\gamma_0^{(k)}$ and $N_0=N_0(j)$ no new floor, that is, no new lower
bound on a parameter, and no cap in the sense of Definition~\ref{2.3:def-cap} arises. Moreover,
the member \eqref{4.26:eq-no-new-condition-1} and the three conditions of (b):
\begin{itemize}
\item involve neither the number $K$ of steps, nor the difference $k-k_a$ between the current
level and the level $k_a$ of an arrest, nor the integers $N_0$, $N_0^a$ (the integer $N_0$
attached to an arrest), nor the number of live pieces of an arrest;
\item impose no lower bound on any running coupling $g_k$;
\item contain $(1+\beta_0)$ only to the fixed power $2$, and no power of $(1+\beta_0)$ whose
exponent grows with $k-k_a$.
\end{itemize}
\item[(e)] The constant $L_0$, bound existentially once, is unaffected. The parameters are chosen
in the order
\begin{equation}\label{4.26:eq-no-new-condition-2}
L \;\to\; \mathfrak{M}\ (\text{subject to } R_1\ge L^2) \;\to\; \kappa \;\to\; \kappa_1
\;\to\; M_1 \;\to\; M \;\to\; (q,C) \;\to\; \gamma .
\end{equation}
Here $\mathfrak{M}$ is chosen after $L$ subject to the condition $R_1\ge L^2$, $R_1$ being the
quantity so denoted in the choice of $\mathfrak{M}$.
The function $\Psi$ is formed after $(q,C)$ and before $\gamma$.
\end{enumerate}
\end{remark}

\begin{proof}
(a) By Theorem~\ref{4.26:thm-density-majorant}, in the corresponding form,
$\gamma_0^{(k)}\le\bar{\gamma}_0(\mathsf{t}_k)$.
Proposition~\ref{4.26:prop-maximum-uses} gives
$\bar{\gamma}_0(\mathsf{t}_k)\le\bar{\gamma}_0(\mathsf{t}_\gamma)$. The definition of
$\bar{\gamma}_0$ gives $120\,\bar{\gamma}_0(\mathsf{t})=600(1+\beta_0)^2\Psi(\mathsf{t})$.
Combining these three facts,
\begin{equation*}
120\,\gamma_0^{(k)}\le 120\,\bar{\gamma}_0(\mathsf{t}_k)\le 600\,(1+\beta_0)^2\,\Psi(\mathsf{t}_\gamma).
\end{equation*}
The second condition in \eqref{4.26:eq-no-new-condition-1} bounds the right-hand side by $1$.
Hence $120\,\gamma_0^{(k)}\le 1$ at every level $k$, and \eqref{4.26:eq-no-new-condition-1}
is a single member of the list. The two further members named in (a) keep their places in the
list and change only in form, so the number of members is unchanged.

Since $c_\beta=2\log(1+\beta_0)>0$, the first condition in \eqref{4.26:eq-no-new-condition-1}
gives $\mathsf{t}_\gamma\ge 2(r+q)$. By \eqref{4.26:eq-density-maximum-a}, $\Psi(\mathsf{t})$ is a
positive constant times $\mathsf{t}^{r+q}e^{-\mathsf{t}/2}$, which is non-increasing for
$\mathsf{t}\ge 2(r+q)$; so the conjunction \eqref{4.26:eq-no-new-condition-1} yields
\begin{equation*}
\sup_{\mathsf{t}\ge\mathsf{t}_\gamma}600\,(1+\beta_0)^2\,\Psi(\mathsf{t})\le 1 .
\end{equation*}

(b) Each of the three conditions of (b) is a smallness condition monotone in $\gamma_0$, of the
first of the two kinds described in Proposition~\ref{4.26:prop-maximum-uses}. Imposed with
$\gamma_0$ replaced by $\bar{\gamma}_0(\mathsf{t}_\gamma)$, each of them holds at
$\gamma_0^{(k)}$, since
\begin{equation*}
\gamma_0^{(k)}\le\bar{\gamma}_0(\mathsf{t}_k)\le\bar{\gamma}_0(\mathsf{t}_\gamma).
\end{equation*}
No new member arises.

(c) Notation~\ref{4.26:not-density-maximum} defines $N_0(j)$ by the equality
$L^{N_0(j)-1}=R_{j-N_0(j)+1}$. The member $(N_0+1)(N_0+3)\le R_k$ is imposed with $N_0$
replaced by the majorant furnished by this defining equality. Indexing $N_0$ by
the density $j$ therefore neither introduces this member nor strengthens it.

(d) By (a)--(c), the conditions of the list differ from those of the formulation with fixed
$\gamma_0$ in the member \eqref{4.26:eq-no-new-condition-1}, in the three conditions of (b), in
which $\gamma_0$ stands replaced by the quantity
\begin{equation*}
\bar{\gamma}_0(\mathsf{t}_\gamma)=5(1+\beta_0)^2\Psi(\mathsf{t}_\gamma),
\end{equation*}
and in the form of the two further members named in (a);
the member of (c) has the same content in both formulations. The member
\eqref{4.26:eq-no-new-condition-1} and the quantity $\bar{\gamma}_0(\mathsf{t}_\gamma)$ are built
only from $\mathsf{t}_\gamma=\log\gamma^{-2}$, $r$, $q$, $c_\beta$,
$\beta_0$, $L$ and $C$, and the conditions of (b) bound this quantity from above as they bound
$\gamma_0$ in the formulation with fixed $\gamma_0$. So they restrict, among the parameters, only
the threshold $\gamma$,
impose no lower bound on any running coupling $g_k$, involve none of $K$, $k-k_a$, $N_0$, $N_0^a$
or the number of live pieces of an arrest, and bring in $(1+\beta_0)$ only to the fixed power $2$.
Hence no new floor and no cap arise, and the member \eqref{4.26:eq-no-new-condition-1} and the
three conditions of (b) have the three properties listed in (d).

(e) By (d), no lower bound on $L$ is added or raised, so $L_0$ is unaffected. The member
\eqref{4.26:eq-no-new-condition-1} bounds only $\gamma$ and involves $L$ and $(q,C)$, which are
chosen before $\gamma$ in \eqref{4.26:eq-no-new-condition-2}; so $\Psi$ can be formed after
$(q,C)$ and before $\gamma$, and the order \eqref{4.26:eq-no-new-condition-2} is unchanged.
\end{proof}

\begin{proposition}[The density-indexed maximum on a falling trajectory]\label{4.26:prop-falling-trajectory}
Let the couplings $x_j=g_j^{-2}$, $\mathsf{t}_j=\log g_j^{-2}$, the radii
$\alpha_{e,m}=C_e\mathsf{t}_m^q e^{-\mathsf{t}_m/2}$, $e\in\{0,1\}$ (written $\alpha_{i,m}$ in
Notation~\ref{4.26:not-density-maximum}), the functions
$\alpha(\mathsf{t})=C\mathsf{t}^q e^{-\mathsf{t}/2}$ with $C=\max\{C_0,C_1\}$, the integer $N_0(j)$, the
weights $w_\ast(j)=L_0^{2N_0(j)}$, the maxima $\bar{\alpha}_j$ and the density-indexed maximum
\begin{equation*}
\gamma_0^{(k)}=5\max_{j\le k}w_\ast(j)\,\bar{\alpha}_j
\end{equation*}
(the maximum over the densities $j\le k$ at which $N_0(j)$ is defined) be as in
Notation~\ref{4.26:not-density-maximum}, and $\Psi$, $\bar{\gamma}_0$ as in
\eqref{4.26:eq-density-maximum-a}. Let $m^{\ast}$, $i$, $\mathsf{t}^{\circ}$, $\delta$ and
$\mathsf{t}_\gamma$ be as in the admissible-falling-trajectory proposition, and consider the
trajectory of that proposition: $k=m^{\ast}+i+1$, $k_a=k-1$, and
\begin{align*}
R_m&=L^{i+1}\quad(m\le m^{\ast}), & R_m&=L^{i}\quad(m^{\ast}<m\le k),\\
\mathsf{t}_m&\in(\mathsf{t}^{\circ},\mathsf{t}^{\circ}+\delta]\quad(m\le m^{\ast}), &
\mathsf{t}_m&=\mathsf{t}_\gamma\in[\mathsf{t}^{\circ}-\delta,\mathsf{t}^{\circ}]\quad(m^{\ast}<m\le k),
\end{align*}
which respects \cite[(2.6)]{B-CMP119}, condition (g) of \eqref{4.26:eq-density-maximum-a} and the
second relation of \cite[(2.9)]{B-CMP119}. Assume moreover that $\mathsf{t}_\gamma\ge 2q$ and that
the hypotheses of Lemma~\ref{4.26:lem-weight-bound} hold at the density $j=k_a$ (in particular
$L\ge3$ and $L_0^2\le L/3$). Let $\omega$ be the weight and $\alpha_{e,k_a}$ the radius of that
proposition, for the index $e\in\{0,1\}$ fixed there. Then:
\begin{enumerate}
\item[(a)] $N_0(j)=i+2$ for every $j\le k_a$ at which $N_0(j)$ is defined, and $N_0(k)=i+1$;
\item[(b)] $\gamma_0^{(k)}=5\,w_\ast(k_a)\,\alpha(\mathsf{t}_\gamma)$; the maximum defining
$\gamma_0^{(k)}$ is attained at every density $j$ with $m^{\ast}<j\le k_a$, and not at $j=k$,
whose term is smaller by the factor $L_0^{-2}\le\frac14$;
\item[(c)] the trajectory is majorised by $\gamma_0^{(k)}$: $5\,\omega\,\alpha_{e,k_a}\le\gamma_0^{(k)}$,
with equality when $C_e=C$;
\item[(d)] $\gamma_0^{(k)}<5\cdot 0.37\,\Psi(\mathsf{t}_k)<\bar{\gamma}_0(\mathsf{t}_k)$.
\end{enumerate}
\end{proposition}

\begin{proof}
\emph{The integers $N_0(j)$.} By Notation~\ref{4.26:not-density-maximum}, $N_0(j)$ is the smallest
positive integer $N$ with
\begin{equation}\label{4.26:eq-falling-trajectory-1}
L^{N-1}=R_{j-N+1}.
\end{equation}
We scan the candidates $N$ at three kinds of density.

At $j=k$: for $N\le i$ the index $k-N+1\ge k-i+1=m^{\ast}+2$ exceeds $m^{\ast}$, so
$R_{k-N+1}=L^i$, while $L^{N-1}\le L^{i-1}$; no such $N$ satisfies
\eqref{4.26:eq-falling-trajectory-1}. For $N=i+1$ the index is $k-i=m^{\ast}+1$, and
$R_{m^{\ast}+1}=L^i=L^{N-1}$. Hence $N_0(k)=i+1$.

At $j=m^{\ast}+s$ with $1\le s\le i$, that is $m^{\ast}<j\le k_a$: for $N\le s$ the index
$j-N+1\ge m^{\ast}+1$, so $R_{j-N+1}=L^i$, while $L^{N-1}\le L^{s-1}\le L^{i-1}$; none of these fits.
For $s<N\le i+1$ the index $j-N+1\le m^{\ast}$, so $R_{j-N+1}=L^{i+1}$, while $L^{N-1}\le L^{i}$; none
of these fits either, and in particular $N=i+1$ fails because $R_{j-i}=L^{i+1}$. For $N=i+2$ the
index is $j-i-1\le m^{\ast}$, and $R_{j-i-1}=L^{i+1}=L^{N-1}$, so $N=i+2$ fits wherever the index
$j-i-1$ exists.

At $j\le m^{\ast}$: every index $j-N+1$ with $N\ge1$ is at most $m^{\ast}$, so $R_{j-N+1}=L^{i+1}$,
and \eqref{4.26:eq-falling-trajectory-1} holds exactly for $N=i+2$, wherever the index $j-i-1$
exists. This proves (a). Consequently
\begin{equation}\label{4.26:eq-falling-trajectory-2}
w_\ast(j)=L_0^{2(i+2)}=w_\ast(k_a)\quad(j\le k_a,\ N_0(j)\ \text{defined}),\qquad
w_\ast(k)=L_0^{2(i+1)}.
\end{equation}

\emph{The terms of the maximum.} Since $\alpha_{e,m}=C_e\mathsf{t}_m^q e^{-\mathsf{t}_m/2}$ for
$e\in\{0,1\}$ and $C=\max\{C_0,C_1\}$, the maximum over $e\in\{0,1\}$ of $\alpha_{e,m}$ is
$\alpha(\mathsf{t}_m)$, and therefore
$\bar{\alpha}_j=\max_{m\le j}\alpha(\mathsf{t}_m)$. The derivative of $\log\alpha(\mathsf{t})$ is
$q/\mathsf{t}-\frac12$, which is negative for $\mathsf{t}>2q$; so $\alpha$ is strictly decreasing on
$[2q,\infty)$. Every $\mathsf{t}_m$ of the trajectory is at least $\mathsf{t}_\gamma$, because
$\mathsf{t}_m>\mathsf{t}^{\circ}\ge\mathsf{t}_\gamma$ for $m\le m^{\ast}$ and
$\mathsf{t}_m=\mathsf{t}_\gamma$ for $m>m^{\ast}$; and $\mathsf{t}_\gamma\ge2q$ by hypothesis. Hence
\begin{equation*}
\bar{\alpha}_j=\alpha\Bigl(\min_{m\le j}\mathsf{t}_m\Bigr).
\end{equation*}
For $j\le m^{\ast}$ the minimum is taken over finitely many values in
$(\mathsf{t}^{\circ},\mathsf{t}^{\circ}+\delta]$, so it exceeds $\mathsf{t}^{\circ}\ge\mathsf{t}_\gamma$,
and by \eqref{4.26:eq-falling-trajectory-2} the term at such a density satisfies
\begin{equation*}
w_\ast(j)\,\bar{\alpha}_j<L_0^{2(i+2)}\alpha(\mathsf{t}_\gamma).
\end{equation*}
For $m^{\ast}<j\le k_a$ the minimum includes $\mathsf{t}_{m^{\ast}+1}=\mathsf{t}_\gamma$ and equals
$\mathsf{t}_\gamma$, so the term is
\begin{equation*}
w_\ast(j)\,\bar{\alpha}_j=L_0^{2(i+2)}\alpha(\mathsf{t}_\gamma)=w_\ast(k_a)\,\alpha(\mathsf{t}_\gamma).
\end{equation*}
At $j=k$ the minimum is again $\mathsf{t}_\gamma$, and the term is
$L_0^{2(i+1)}\alpha(\mathsf{t}_\gamma)$, which is $L_0^{-2}$ times the previous one; by the
condition $1.4L_0^2\ge5.6$ of the admissible-falling-trajectory proposition, $L_0^2\ge4$ and
$L_0^{-2}\le\frac14$. Thus the largest term is $w_\ast(k_a)\alpha(\mathsf{t}_\gamma)$, it is attained
at every density $j$ with $m^{\ast}<j\le k_a$ and at no other density $j\le k$, and
\begin{equation}\label{4.26:eq-falling-trajectory-3}
\gamma_0^{(k)}=5\,w_\ast(k_a)\,\alpha(\mathsf{t}_\gamma).
\end{equation}
This proves (b).

\emph{Majorisation.} By the clause of the admissible-falling-trajectory proposition fixing the
weight of the band, $\omega=L_0^{2N_0(k_a)}$, and this is $w_\ast(k_a)$. Since
$\mathsf{t}_{k_a}=\mathsf{t}_\gamma$,
\begin{equation*}
\alpha_{e,k_a}=C_e\mathsf{t}_\gamma^q e^{-\mathsf{t}_\gamma/2}\le C\mathsf{t}_\gamma^q e^{-\mathsf{t}_\gamma/2}
=\alpha(\mathsf{t}_\gamma),
\end{equation*}
with equality when $C_e=C$. Multiplying by $5\omega=5w_\ast(k_a)$ and using
\eqref{4.26:eq-falling-trajectory-3} gives $5\omega\alpha_{e,k_a}\le\gamma_0^{(k)}$, with equality
when $C_e=C$. This proves (c).

The maximum over densities cannot be replaced here by the term at the current density $k$.
Replacing it by that term would require
$1.4\,\omega\,\alpha(\mathsf{t}_\gamma)\le5w_\ast(k)\alpha(\mathsf{t}_\gamma)$, that is
$\omega\le5w_\ast(k)/1.4$, the constant $1.4$ being the one in the condition $1.4L_0^2\ge5.6$ of the
admissible-falling-trajectory proposition; this requirement fails on the trajectory. Indeed,
by \eqref{4.26:eq-falling-trajectory-2}, $5w_\ast(k)\alpha(\mathsf{t}_\gamma)
=5L_0^{2(i+1)}\alpha(\mathsf{t}_\gamma)$, whereas $\omega=L_0^{2}\,w_\ast(k)$ and, by that condition,
\begin{align*}
1.4\,\omega\,\alpha(\mathsf{t}_\gamma)&=1.4L_0^2\cdot L_0^{2(i+1)}\alpha(\mathsf{t}_\gamma)\\
&\ge5.6\,L_0^{2(i+1)}\alpha(\mathsf{t}_\gamma)
=1.12\cdot5\,w_\ast(k)\,\alpha(\mathsf{t}_\gamma).
\end{align*}
So $5w_\ast(k)\alpha(\mathsf{t}_\gamma)$ falls short of $1.4\,\omega\,\alpha(\mathsf{t}_\gamma)$ by a
factor at least $1.12$. This happens already at $k-k_a=1$, after a single drop of $R$ by the factor
$L^{-1}$ (from $R_{m^{\ast}}=L^{i+1}$ to $R_{m^{\ast}+1}=L^{i}$).

\emph{The bound by the current coupling.} By hypothesis, Lemma~\ref{4.26:lem-weight-bound} applies
at the density $k_a$; we record which case of its proof is met. Write $R_j=L^{\iota}$ and define
$d$ by $L^{N_0(j)-1}=L^{\iota+d}$, as in the proof of Lemma~\ref{4.26:lem-weight-bound}. At
$j=k_a$: $R_{k_a}=L^i$, so $\iota=i$, and $N_0(k_a)-1=i+1=\iota+1$ by (a), so $d=1$. At $j=k$:
$\iota=i$ and $N_0(k)-1=i$, so $d=0$. At $j\le m^{\ast}$: $\iota=i+1$ and $N_0(j)-1=i+1$, so $d=0$.
All these densities therefore fall in the case $d\le1$ of the proof of
Lemma~\ref{4.26:lem-weight-bound}, in which no property of the couplings is used. By
\eqref{4.26:eq-weight-bound-a} at the density $k_a$, together with the hypothesis $L_0^2\le L/3$
and \eqref{4.26:eq-falling-trajectory-2},
\begin{equation*}
w_\ast(k_a)=L_0^{2(i+2)}\le(L/3)^{i+2}<0.37\,L^2(L+1)\,\mathsf{t}_\gamma^r .
\end{equation*}
Since $k_a$ and $k$ exceed $m^{\ast}$, $\mathsf{t}_{k_a}=\mathsf{t}_k=\mathsf{t}_\gamma$. Inserting the
last bound into \eqref{4.26:eq-falling-trajectory-3} and recalling
$\Psi(\mathsf{t})=L^2(L+1)C\mathsf{t}^{r+q}e^{-\mathsf{t}/2}$ and
$\bar{\gamma}_0(\mathsf{t})=5(1+\beta_0)^2\Psi(\mathsf{t})$ from
\eqref{4.26:eq-density-maximum-a}, we obtain
\begin{align*}
\gamma_0^{(k)}=5\,w_\ast(k_a)\,\alpha(\mathsf{t}_\gamma)
&<5\cdot0.37\,L^2(L+1)\,\mathsf{t}_\gamma^r\cdot C\mathsf{t}_\gamma^q e^{-\mathsf{t}_\gamma/2}\\
&=5\cdot0.37\,\Psi(\mathsf{t}_k)<5(1+\beta_0)^2\,\Psi(\mathsf{t}_k)=\bar{\gamma}_0(\mathsf{t}_k),
\end{align*}
the last inequality because $0.37<1\le(1+\beta_0)^2$. The factor $(1+\beta_0)^2$ is not needed on
this trajectory. This proves (d), and shows that the corresponding bound of
Theorem~\ref{4.26:thm-density-majorant} holds on the trajectory with room to spare.
\end{proof}

\begin{remark}[Other readings of the density-indexed maximum]\label{4.26:rem-other-readings}
In Notation~\ref{4.26:not-density-maximum} the quantity $\bar{\alpha}_j$ is the maximum of the radii over
the levels $m\le j$; we call this reading (A):
\begin{equation*}
\bar{\alpha}_{i,j}=\max_{1\le m\le j}\alpha_{i,m},\qquad \bar{\alpha}_j=\max_i\bar{\alpha}_{i,j},\qquad
\gamma_0^{(k)}=5\max_{j\le k}w_{\ast}(j)\,\bar{\alpha}_j ,
\end{equation*}
the last maximum being taken over the densities $j\le k$ at which $N_0(j)$ is defined.
Two other readings of $\bar{\alpha}_j$ can be considered:
\begin{align}
&\text{(B)}\qquad \bar{\alpha}_j\ \text{replaced by}\ \bar{\alpha}(\mathsf{t}_j),\qquad
\bar{\alpha}(\mathsf{t}):=\sup_{\mathsf{t}'\ge\mathsf{t}}\alpha(\mathsf{t}'),
\label{4.26:eq-other-readings-1}\\
&\text{(A}'\text{)}\qquad \bar{\alpha}'_j:=\sup_{\mathsf{t}'\ge\mathsf{t}_j-c_\beta}\alpha(\mathsf{t}').
\label{4.26:eq-other-readings-2}
\end{align}
The readings compare as follows; reading (B) is not used in this chapter.
\begin{itemize}
\item[(a)] In the step of Lemma~\ref{4.26:lem-analyticity-membership} in which the weight
$\omega(x)$ of a point $x$ enters, a radius $\alpha_{\cdot,m}$ at a level $m\le k^{\circ}$ is bounded by
the maximum read at a density $k^{\circ}\le k$. Under (B) this requires
$\mathsf{t}_m\ge\mathsf{t}_{k^{\circ}}$,
a monotonicity of $\mathsf{t}$ along the trajectory that is not assumed. The second inequality of
\cite[(2.6)]{B-CMP119} restores the bound only up to a factor $1+\beta_0$, that is, with
$5(1+\beta_0)>5$ in place of the factor $5$ in $\gamma_0^{(k)}$, and the membership on
$Z^{\mathrm{on}}$ requires the factor $5$ exactly. The only way to keep the factor $5$ is to use the
term $j=m(x)$ of the maximum defining $\gamma_0^{(k)}$, provided $m(x)$ is a density at which $N_0$ is
defined; here $m(x)$ is the level $m$ of the radius in question and $w_{\ast}(m(x))\ge\omega(x)$.
This requires that
\begin{equation}\label{4.26:eq-other-readings-3}
k_0(j):=j-N_0(j)\quad\text{be non-decreasing in the density } j .
\end{equation}
\item[(b)] Under (B) the majorant $\bar{\gamma}_0$ would carry only one factor $1+\beta_0$, not the
factor $(1+\beta_0)^2$ of Notation~\ref{4.26:not-density-maximum}. Moreover, the comparison of
the earlier levels $m<j$ against the density $j$ in Lemma~\ref{4.26:prf-level-comparison}, which
holds uniformly in $j-m$, would then compare nothing.
\item[(c)] Reading (A$'$) is equivalent to reading (A) for the purposes of this chapter. It gives the
memberships of Lemma~\ref{4.26:lem-analyticity-membership}, and it gives the bound of
$\gamma_0^{(k)}$ by $\bar{\gamma}_0(\mathsf{t}_k)$ with the factor $(1+\beta_0)^2$. Of the three
readings, (A) requires the least: it needs no inequality on the trajectory for the memberships. It is
the reading used.
\item[(d)] Two further devices are compatible with reading (A). First, a bound on the weights
$w_{\ast}(j)$ at one density that does not depend on the trajectory
is an admissible alternative way to bound these weights, and it is the mechanism behind the bound
\eqref{4.26:eq-weight-bound-a} of Lemma~\ref{4.26:lem-weight-bound}. Second, a constant independent
of the trajectory could stand in place of $\gamma_0^{(k)}$ as an upper value of the size of the
polydisc; being one number for all densities, it would impose on the couplings a threshold uniform in
the density unless it is read inside suprema over a single value of $\mathsf{t}$. So read, it is
subsumed by $\bar{\gamma}_0(\mathsf{t}_k)$.
\end{itemize}
\end{remark}

\begin{proof}
Throughout, $\mathsf{t}_j\ge\mathsf{t}_\gamma$ at every density $j\le k$ that occurs, and
$\mathsf{t}_\gamma\ge 2(r+q)+c_\beta$, as fixed in Notation~\ref{4.26:not-density-maximum}. We also
have $c_\beta=2\ln(1+\beta_0)$.

We use three facts, all recorded in Notation~\ref{4.26:not-density-maximum}; we recall their proofs.
Put $\varphi_q(\mathsf{t})=\mathsf{t}^qe^{-\mathsf{t}/2}$, so that $\alpha=C\varphi_q$.

First, for $0\le c\le s$ we have $(s-c)^q\le s^q$ and $e^{-(s-c)/2}=e^{c/2}e^{-s/2}$. Hence
\begin{equation}\label{4.26:eq-other-readings-4}
\varphi_q(s-c)\le e^{c/2}\varphi_q(s)\qquad (0\le c\le s).
\end{equation}
With $c=c_\beta$ the factor is $e^{c_\beta/2}=1+\beta_0$; this is the shift used in
Lemma~\ref{4.26:prf-level-comparison} and in the bound of $\gamma_0^{(k)}$ by
$\bar{\gamma}_0(\mathsf{t}_k)$.
Second, $(\ln\varphi_q)'(\mathsf{t})=q/\mathsf{t}-\tfrac12\le 0$ for $\mathsf{t}\ge 2q$, so $\varphi_q$
is non-increasing on $[2q,\infty)$.

Finally, the second inequality of \cite[(2.6)]{B-CMP119} reads $g_m\le(1+\beta_0)g_n$ for $n>m$.
Squaring it and inverting gives $x_m\ge x_n/(1+\beta_0)^2$. Taking logarithms,
\begin{equation}\label{4.26:eq-other-readings-5}
\mathsf{t}_m\ge\mathsf{t}_n-c_\beta\qquad(n>m).
\end{equation}
This is a lower bound on the earlier variable.

(a) Since $C_i\le C$, we have $\alpha_{\cdot,m}\le\alpha_{i,m}\le\alpha(\mathsf{t}_m)$. Under (B), the
value $\alpha(\mathsf{t}_m)$ is a term of the supremum $\bar{\alpha}(\mathsf{t}_{k^{\circ}})$ exactly
when $\mathsf{t}_m\ge\mathsf{t}_{k^{\circ}}$. Nothing in the hypotheses orders $\mathsf{t}_m$ and
$\mathsf{t}_{k^{\circ}}$ in this way.

For $m=k^{\circ}$ there is nothing to prove. For $m<k^{\circ}$, what
\eqref{4.26:eq-other-readings-5} gives, applied to the pair $(m,k^{\circ})$, is
$\mathsf{t}_m\ge\mathsf{t}_{k^{\circ}}-c_\beta$. Put $s:=\mathsf{t}_m+c_\beta$, so that
$s\ge\mathsf{t}_{k^{\circ}}$. Moreover
\begin{equation*}
\mathsf{t}_m\ge\mathsf{t}_\gamma-c_\beta\ge 2(r+q)\ge 0,
\end{equation*}
so that $0\le c_\beta\le s$, and \eqref{4.26:eq-other-readings-4} applies. It gives
\begin{equation*}
\alpha(\mathsf{t}_m)=C\varphi_q(s-c_\beta)\le e^{c_\beta/2}\,C\varphi_q(s)
=(1+\beta_0)\,\alpha(s)\le(1+\beta_0)\,\bar{\alpha}(\mathsf{t}_{k^{\circ}}).
\end{equation*}
The membership therefore holds under (B) only if the factor $5$ in $\gamma_0^{(k)}$ is replaced by
$5(1+\beta_0)>5$. The membership on $Z^{\mathrm{on}}$ requires the factor $5$ exactly.

To keep the factor $5$, one bounds $\omega(x)\,\alpha_{\cdot,m}$ by the single term $j=m(x)$ of the
maximum over densities, provided $N_0(m(x))$ is defined, using $w_{\ast}(m(x))\ge\omega(x)$ and
$\alpha_{\cdot,m}\le\alpha(\mathsf{t}_m)\le\bar{\alpha}(\mathsf{t}_m)$. This needs the monotonicity
\eqref{4.26:eq-other-readings-3}. That monotonicity follows from the rule
$R_{j+1}\in\{R_j,\,L^{-1}R_j\}$ stated in \cite{B-CMP122a}; the proofs of this chapter do not use
this rule. It also follows if the numbers $R_j$ are replaced by a monotone envelope
$\check{R}_j$ of them. It follows neither from \cite[(2.6)]{B-CMP119} nor from the two-sided bounds
on the running couplings.

Nor does the first inequality of \cite[(2.9)]{B-CMP119}, $R_n\le LR_m$, give
\eqref{4.26:eq-other-readings-3}. That inequality is compatible with $R$ rising by one power of $L$
at a step where the coupling decreases. Such a step increases $\mathsf{t}$, and with it the value
$R(\mathsf{t})$ of the non-decreasing function $R(\cdot)$ of Notation~\ref{4.26:not-density-maximum},
the least integer power of $L$ that is at least $\mathsf{t}^r$ \cite[(2.5)]{B-CMP119}.

(b) Let $j\le k$ be a density that occurs. Then $\mathsf{t}_j\ge\mathsf{t}_\gamma$ and
$\mathsf{t}_\gamma\ge 2q$. Since $\varphi_q$ is
non-increasing on $[2q,\infty)$, the supremum in \eqref{4.26:eq-other-readings-1} is attained at
$\mathsf{t}'=\mathsf{t}_j$:
\begin{equation*}
\bar{\alpha}(\mathsf{t}_j)=\alpha(\mathsf{t}_j).
\end{equation*}
Under reading (A), by contrast, the comparison of the levels $m<j$ against the density $j$ in
Lemma~\ref{4.26:prf-level-comparison} costs one application of \eqref{4.26:eq-other-readings-5}
for each pair $(m,j)$, uniformly in $j-m$. Its conclusion
yields the second factor $1+\beta_0$ in
$\bar{\gamma}_0=5(1+\beta_0)^2\Psi$. Under (B) no level $m<j$ enters $\bar{\alpha}(\mathsf{t}_j)$,
so that comparison has nothing to compare, and only one factor $1+\beta_0$ arises in $\bar{\gamma}_0$.

(c) Memberships under (A$'$). Let $m\le j$. If $m<j$, then
$\mathsf{t}_m\ge\mathsf{t}_j-c_\beta$ by \eqref{4.26:eq-other-readings-5}, applied once to the pair
$(m,j)$; if $m=j$ this is trivial. Hence
$\alpha_{i,m}\le\alpha(\mathsf{t}_m)\le\bar{\alpha}'_j$, since $\alpha(\mathsf{t}_m)$ is a term of
the supremum \eqref{4.26:eq-other-readings-2}. Therefore $\bar{\alpha}_j\le\bar{\alpha}'_j$. Every
membership of Lemma~\ref{4.26:lem-analyticity-membership} that holds with $\bar{\alpha}_j$ therefore
holds with $\bar{\alpha}'_j$.

The bound under (A$'$). Let $\mathsf{t}'\ge\mathsf{t}_j-c_\beta$ and put $s:=\mathsf{t}'+c_\beta$, so
that $s\ge\mathsf{t}_j$. Since also $\mathsf{t}_j\ge\mathsf{t}_\gamma$ and
$\mathsf{t}_\gamma\ge c_\beta$, we have $s\ge c_\beta$, so
\eqref{4.26:eq-other-readings-4} applies. Since also $\mathsf{t}_j\ge 2q$, we have $s\ge 2q$, and the
monotonicity of $\varphi_q$ applies. Together they give
\begin{equation*}
\alpha(\mathsf{t}')\le(1+\beta_0)\,\alpha(s)\le(1+\beta_0)\,\alpha(\mathsf{t}_j).
\end{equation*}
Hence $\bar{\alpha}'_j\le(1+\beta_0)\alpha(\mathsf{t}_j)$. This is the bound of
Lemma~\ref{4.26:prf-level-comparison} under reading (A). The bound of $\gamma_0^{(k)}$ by
$\bar{\gamma}_0(\mathsf{t}_k)$
therefore holds under (A$'$) with the same factor $(1+\beta_0)^2$.

Comparison of (A) and (A$'$). Under reading (A) the memberships hold by definition: each
$\alpha_{\cdot,m}$ with $m\le j$ is a term of the maximum $\bar{\alpha}_j$, and no inequality on the
trajectory is used. Reading (A$'$) uses \eqref{4.26:eq-other-readings-5} already for the memberships.
This is why (A) is the reading used.
\end{proof}

\begin{lemma}[Ordering of the three units and conversion to the shell unit]
\label{4.26:lem-unit-ordering}
Let $\mathfrak{a}$ be a live arrest with $W_{\mathfrak{a}}\subset Z^c$ or $W_{\mathfrak{a}}\subset Z$.
Let $(p,X',\mathfrak{s})$ be a $(\Sigma^{\ast})^{\mathfrak{a}}$-triple in the sense of
Lemma~\ref{4.26:lem-retained-levels}, and let $y'\in X'^{\sim 2}\cap W_{\mathfrak{a}}$, with
$X'^{\sim 2}$ as in Lemma~\ref{4.26:lem-retained-levels}. Put
\begin{equation*}
m:=\ell_{\mathfrak{a}}(y'),\qquad \omega:=w_{\mathfrak{a}}(y'),\qquad
\mathfrak{c}:=\mathfrak{c}(y')=m-\tfrac12\log_{L_0}\omega,
\end{equation*}
\begin{equation*}
\lambda:=\lambda_{\mathfrak{s}}(y')\le m,\qquad
w_b:=w_{\mathfrak{s}}(y')=L_0^{2(\lambda-\mathfrak{c})^+},
\end{equation*}
the levels and weights being those of Lemma~\ref{4.26:lem-retained-levels}. The weights satisfy
\begin{equation}\label{4.26:eq-unit-ordering-1}
\omega\le L_0^{2(m-\lambda)}\,w_b ,
\end{equation}
with equality whenever $\lambda\ge\mathfrak{c}$, in particular at the lifted points of
Lemma~\ref{4.26:lem-retained-levels}, at which $\lambda\ge\mathfrak{c}$ and $m-\mathfrak{c}\ge1$;
and strictly at the un-lifted inner collar of Lemma~\ref{4.26:lem-retained-levels}, where
$\lambda=\mathfrak{c}-1$ and $\omega=w_b=1$. For structures $\mathfrak{s},\mathfrak{t}$
we write $\mathfrak{s}\preccurlyeq\mathfrak{t}$ if $\mathfrak{s}$ refines $\mathfrak{t}$, and
$\varphi_{\mathrm{arrest}}=0.098$ denotes the floor of the arrested source factor. Let
$\mathrm{U1},\mathrm{U2},\mathrm{U3}$ and $\mathbf{r}^{\mathfrak{c},(p)}_4$ be the units of
Definition~\ref{4.26:def-three-units}, displayed in \eqref{4.26:eq-three-units-a}, with
$\ell_n=L^n\eta$. Let $\beta_0=1/7$ be the constant of Lemma~\ref{4.4:lem-radii-ratio}, by which
$\alpha_{0,j+1}\le(1+\beta_0)\alpha_{0,j}$ for adjacent levels $j,j+1$, and assume
$L_0^2\le L/3$ and $L\ge 12$.
Then the following hold.
\begin{enumerate}
\item[(a)] The units are ordered, $\mathrm{U1}\le\mathrm{U2}\le\mathrm{U3}$, with
\begin{equation}\label{4.26:eq-unit-ordering-a}
\begin{gathered}
\frac{\mathrm{U2}}{\mathrm{U1}}=\frac{\ell_{p\wedge m}}{\ell_{p\wedge\lambda}}
=L^{(p\wedge m)-(p\wedge\lambda)}\in\bigl[1,L^{m-\lambda}\bigr],\\
\frac{\mathrm{U3}}{\mathrm{U2}}
=\frac{w_b}{\omega}\,\frac{\alpha_{0,\lambda}}{\alpha_{0,m}}\Bigl(\frac{\ell_m}{\ell_\lambda}\Bigr)^{2}
\ge\Bigl[\frac{3}{L}\,(1+\beta_0)^{-1}L^2\Bigr]^{m-\lambda}=(2.625\,L)^{m-\lambda}.
\end{gathered}
\end{equation}
The exponent $(p\wedge m)-(p\wedge\lambda)$ equals $0$ for $p\le\lambda$, $p-\lambda$ for
$\lambda<p\le m$, and $m-\lambda$ for $p>m$. For $m>\lambda$ one has
$\mathrm{U3}/\mathrm{U2}\ge 31.5$. At $\lambda=m$ one has $\mathrm{U1}=\mathrm{U2}=\mathrm{U3}$.
\item[(b)] At a lifted point of Lemma~\ref{4.26:lem-retained-levels}, where
$\lambda\ge\mathfrak{c}$, put $g:=m-\mathfrak{c}\ge1$. Then
\begin{equation}\label{4.26:eq-unit-ordering-b}
\frac{\mathrm{U2}}{\mathbf{r}^{\mathfrak{c},(p)}_4}
=\omega\,\frac{\alpha_{0,m}}{\alpha_{0,\mathfrak{c}}}
\Bigl(\frac{\ell_{\mathfrak{c}}}{\ell_m}\Bigr)^{2}
\frac{\ell_{p\wedge\mathfrak{c}}}{\ell_{p\wedge\lambda}}
\le\Bigl(\frac{(1+\beta_0)L_0^2}{L^2}\Bigr)^{g}\le\Bigl(\frac{0.381}{L}\Bigr)^{g}.
\end{equation}
At the box triple of Lemma~\ref{4.26:lem-retained-levels}, where $\lambda=\mathfrak{c}$, the last
length factor equals one, so that there
\begin{equation*}
\frac{\mathrm{U2}}{\mathbf{r}^{\mathfrak{c},(p)}_4}
=\omega\,\frac{\alpha_{0,m}}{\alpha_{0,\mathfrak{c}}}
\Bigl(\frac{\ell_{\mathfrak{c}}}{\ell_m}\Bigr)^{2}\quad\text{exactly.}
\end{equation*}
\item[(c)] At a lifted point $\mathrm{U1}/\mathbf{r}^{\mathfrak{c},(p)}_4$ obeys the same bound.
Consequently the bound $0.381^{g}$ of clause $(c_0)$ of the frozen-weight lemma holds, for
$e\in\{2,4\}$, whether the block's unit is taken to be $\mathrm{U1}$ or $\mathrm{U2}$, and with
one power of $L$ to spare per level.
\item[(d)] In every case $\mathrm{U2}\le(0.381/L)^{m-\lambda}\,\mathrm{U3}$.
\end{enumerate}
\end{lemma}

\begin{proof}
\emph{The inequality \eqref{4.26:eq-unit-ordering-1} and its two cases.}
By the definition of $\mathfrak{c}$ we have $\omega=L_0^{2(m-\mathfrak{c})}$. Since
$(\lambda-\mathfrak{c})^+\ge\lambda-\mathfrak{c}$ and $L_0>1$,
\begin{equation*}
L_0^{2(m-\lambda)}w_b=L_0^{2(m-\lambda)+2(\lambda-\mathfrak{c})^+}
\ge L_0^{2(m-\lambda)+2(\lambda-\mathfrak{c})}=L_0^{2(m-\mathfrak{c})}=\omega ,
\end{equation*}
with equality exactly when $(\lambda-\mathfrak{c})^+=\lambda-\mathfrak{c}$, that is, when
$\lambda\ge\mathfrak{c}$. In that case
\begin{equation*}
L_0^{2(m-\lambda)}w_b=L_0^{2(m-\lambda)}L_0^{2(\lambda-\mathfrak{c})}
=L_0^{2(m-\mathfrak{c})}=\omega .
\end{equation*}
This is equality in \eqref{4.26:eq-unit-ordering-1}. At the un-lifted inner collar of
Lemma~\ref{4.26:lem-retained-levels} one has $\omega=1$, so $\mathfrak{c}=m$, and
$\lambda=\mathfrak{c}-1=m-1$. Then $(\lambda-\mathfrak{c})^+=0$, so $w_b=1$, and the right side of
\eqref{4.26:eq-unit-ordering-1} equals $L_0^{2}>1=\omega$. This is the strict case.

\emph{The first ratio in \eqref{4.26:eq-unit-ordering-a}.}
By Definition~\ref{4.26:def-three-units}, $\mathrm{U1}$ and $\mathrm{U2}$ share the factor
$\omega\,\alpha_{0,m}\ell_m^{-2}$. They differ only in the last length:
$\ell_{p\wedge m}^{-1}$ in $\mathrm{U1}$ and $\ell_{p\wedge\lambda}^{-1}$ in $\mathrm{U2}$. Hence
$\mathrm{U2}/\mathrm{U1}=\ell_{p\wedge m}/\ell_{p\wedge\lambda}$. Since $\ell_n=L^n\eta$, this
equals $L^{(p\wedge m)-(p\wedge\lambda)}$.

Because $\lambda\le m$, the exponent is computed in three cases. For $p\le\lambda$ both minima
equal $p$, and the exponent is $0$. For $\lambda<p\le m$ they are $p$ and $\lambda$, and the
exponent is $p-\lambda$. For $p>m$ they are $m$ and $\lambda$, and the exponent is $m-\lambda$. In
each case the exponent lies in $[0,m-\lambda]$, so the ratio lies in $[1,L^{m-\lambda}]$.

\emph{The second ratio in \eqref{4.26:eq-unit-ordering-a}.}
By Definition~\ref{4.26:def-three-units}, $\mathrm{U3}=w_b\,\alpha_{0,\lambda}\ell_\lambda^{-2}
\ell_{p\wedge\lambda}^{-1}$ and $\mathrm{U2}=\omega\,\alpha_{0,m}\ell_m^{-2}\ell_{p\wedge\lambda}^{-1}$.
The factor $\ell_{p\wedge\lambda}^{-1}$ cancels, which gives the displayed identity. We bound
the three factors from below.
\begin{itemize}
\item By \eqref{4.26:eq-unit-ordering-1}, $w_b/\omega\ge L_0^{-2(m-\lambda)}$. By
$L_0^2\le L/3$ this is at least $(3/L)^{m-\lambda}$.
\item By Lemma~\ref{4.4:lem-radii-ratio}, $\alpha_{0,j+1}\le(1+\beta_0)\alpha_{0,j}$ for each of
the $m-\lambda$ adjacent pairs $(j,j+1)$ with $\lambda\le j<m$; multiplying these bounds gives
$\alpha_{0,m}\le(1+\beta_0)^{m-\lambda}\alpha_{0,\lambda}$, that is,
$\alpha_{0,\lambda}/\alpha_{0,m}\ge(1+\beta_0)^{-(m-\lambda)}$.
\item Since $\ell_n=L^n\eta$, we have $(\ell_m/\ell_\lambda)^2=L^{2(m-\lambda)}$.
\end{itemize}
Multiplying the three bounds gives
\begin{equation*}
\frac{\mathrm{U3}}{\mathrm{U2}}
\ge\Bigl[\frac{3}{L}(1+\beta_0)^{-1}L^2\Bigr]^{m-\lambda}
=\Bigl(\frac{3L}{1+\beta_0}\Bigr)^{m-\lambda}.
\end{equation*}
With $\beta_0=1/7$ we have $3/(1+\beta_0)=21/8=2.625$, so the bound is $(2.625\,L)^{m-\lambda}$.

If $m>\lambda$, the exponent is at least one. Using $L\ge12$ we get
\begin{equation*}
(2.625\,L)^{m-\lambda}\ge 2.625\,L\ge \tfrac{21\cdot 12}{8}=31.5 .
\end{equation*}
Consider moreover the case $\omega=w_b=1$ and $\lambda=m-1$. There the first factor equals one,
and the last two bounds alone give
\begin{equation*}
\frac{\mathrm{U3}}{\mathrm{U2}}\ge\frac{L^2}{1+\beta_0}=\frac{7L^2}{8}.
\end{equation*}
The right side equals $126$ at $L=12$.

\emph{The case $\lambda=m$.}
Then $p\wedge\lambda=p\wedge m$, so the first ratio equals one and $\mathrm{U1}=\mathrm{U2}$. The
second ratio reduces to $w_b/\omega$, because $\alpha_{0,\lambda}=\alpha_{0,m}$ and
$\ell_\lambda=\ell_m$. Here $\omega=L_0^{2(m-\mathfrak{c})}$ with $\mathfrak{c}\le m$, since
$\omega\ge1$ (Lemma~\ref{4.26:lem-retained-levels}). Hence $\lambda=m\ge\mathfrak{c}$, and
equality holds in \eqref{4.26:eq-unit-ordering-1} with $m-\lambda=0$. This gives $w_b=\omega$ and
$\mathrm{U2}=\mathrm{U3}$.

\emph{The ordering.}
Both ratios in \eqref{4.26:eq-unit-ordering-a} are at least one. For the second ratio this holds
because $(2.625\,L)^{m-\lambda}\ge1$. Hence $\mathrm{U1}\le\mathrm{U2}\le\mathrm{U3}$, and (a) is
proved.

\emph{Proof of (b).}
Take $\mathbf{r}^{\mathfrak{c},(p)}_4=\alpha_{0,\mathfrak{c}}\ell_{\mathfrak{c}}^{-2}
\ell_{p\wedge\mathfrak{c}}^{-1}$ from Definition~\ref{4.26:def-three-units}, and divide it into
$\mathrm{U2}$. This gives the identity in \eqref{4.26:eq-unit-ordering-b}. Let
$\lambda\ge\mathfrak{c}$ and $g=m-\mathfrak{c}\ge1$. We bound the four factors from above.
\begin{itemize}
\item By the definition of $\mathfrak{c}$, $\omega=L_0^{2g}$.
\item Apply Lemma~\ref{4.4:lem-radii-ratio} to the $g$ adjacent pairs between $\mathfrak{c}$
and $m$. This gives $\alpha_{0,m}/\alpha_{0,\mathfrak{c}}\le(1+\beta_0)^g$.
\item We have $(\ell_{\mathfrak{c}}/\ell_m)^2=L^{-2g}$.
\item Since $\mathfrak{c}\le\lambda$, we have $p\wedge\mathfrak{c}\le p\wedge\lambda$, so
$\ell_{p\wedge\mathfrak{c}}/\ell_{p\wedge\lambda}\le1$. At the box triple $\lambda=\mathfrak{c}$
this factor equals one, which gives the exact identity stated in (b).
\end{itemize}
The product of the four bounds is $((1+\beta_0)L_0^2/L^2)^g$. By $L_0^2\le L/3$ and
$\beta_0=1/7$,
\begin{equation*}
\frac{(1+\beta_0)L_0^2}{L^2}\le\frac{8}{7}\cdot\frac{L}{3}\cdot\frac{1}{L^2}
=\frac{8}{21\,L}\le\frac{0.381}{L},
\end{equation*}
because $8/21=0.38095\ldots$. This proves (b).

\emph{Proof of (c).}
Divide $\mathrm{U1}$ by $\mathbf{r}^{\mathfrak{c},(p)}_4$ in the same way. The result is the
product of the same first three factors and the length factor
$\ell_{p\wedge\mathfrak{c}}/\ell_{p\wedge m}$. This length factor is at most one, since
$\mathfrak{c}\le m$. Hence $\mathrm{U1}/\mathbf{r}^{\mathfrak{c},(p)}_4$ obeys the bound of
\eqref{4.26:eq-unit-ordering-b}.

For $e=2$ the last length is dropped from the units, and $\mathrm{U1}_2=\mathrm{U2}_2$
(Definition~\ref{4.26:def-three-units}). Dropping it likewise from the shell unit, put
$\mathbf{r}^{\mathfrak{c}}_2:=\alpha_{0,\mathfrak{c}}\ell_{\mathfrak{c}}^{-2}$. The ratio
$\mathrm{U2}_2/\mathbf{r}^{\mathfrak{c}}_2$ is then the product of the first three factors, which
is bounded in the same way.

Thus, whether the block's unit is taken to be $\mathrm{U1}$ or $\mathrm{U2}$, the ratio to the
shell unit is at most $(0.381/L)^g$. This is the bound $0.381^g$ of clause $(c_0)$ of the
frozen-weight lemma, improved by the factor $L^{-g}$.

\emph{Proof of (d).}
By the second line of \eqref{4.26:eq-unit-ordering-a},
\begin{equation*}
\mathrm{U2}\le(2.625\,L)^{-(m-\lambda)}\,\mathrm{U3}
=\Bigl(\frac{8}{21\,L}\Bigr)^{m-\lambda}\mathrm{U3}
\le\Bigl(\frac{0.381}{L}\Bigr)^{m-\lambda}\mathrm{U3}.
\end{equation*}
Here $1/2.625=8/21\le0.381$. The inequality holds for every $\lambda\le m$, including
$\lambda=m$, where both sides are equal.
\end{proof}

\begin{remark}[The literal block side does not propagate]\label{4.26:rem-literal-block-side}
Let $(p,X',\mathfrak{s})$ be a triple and $y'$ a point of a live piece, as in
Definition~\ref{4.26:def-three-units}. Write
\[
m=\ell_{\mathfrak{a}}(y'),\qquad \omega=w_{\mathfrak{a}}(y'),\qquad
\lambda=\lambda_{\mathfrak{s}}(y')\le m,
\]
and let $\mathrm{U1}$, $\mathrm{U2}$, $\mathrm{U3}$ be the units of entry $e=4$ displayed in
\eqref{4.26:eq-three-units-a}; for $e=2$ the last length is dropped and $\mathrm{U1}_2=\mathrm{U2}_2$.
Call the exponent $(p\wedge m)-(p\wedge\lambda)$ the \emph{level gap} of the triple at $y'$.
\begin{enumerate}
\item[(a)] $\mathrm{U2}/\mathrm{U1}=L^{(p\wedge m)-(p\wedge\lambda)}$. Hence
$\mathrm{U1}=\mathrm{U2}=\mathrm{U3}$ wherever $\lambda=m$. This
holds at every triple of the present site's structures on $\mathfrak{T}$ and at the triples of the
frozen structure. $\mathrm{U1}$ and $\mathrm{U2}$ can differ only at retained structures with
$\lambda<m$.
\item[(b)] At such retained structures a shell clause written in $\mathrm{U1}$ is not, in general,
propagated from one stage to the next by the increment lemma for the bump shells; propagation fails
on the thin stratum described in the proof.
\item[(c)] For the region $Y$ and the stage $n$ of Notation~\ref{4.26:not-stage-space}, the inclusion
of the source space $\mathfrak{S}_{\mathrm{src}}$ in $\mathfrak{N}^{(n)}(Y)$ fails
clause by clause if the shell clause of $\mathfrak{N}^{(n)}(Y)$ is read in $\mathrm{U1}$.
\end{enumerate}
Accordingly no clause in what follows is read in the unit $\mathrm{U1}$.
\end{remark}

\begin{proof}
(a) By \eqref{4.26:eq-three-units-a}, $\mathrm{U1}$ and $\mathrm{U2}$ share the factor
$\omega\,\alpha_{0,m}\ell_m^{-2}$. They differ only in the last length, $\ell_{p\wedge m}^{-1}$ against
$\ell_{p\wedge\lambda}^{-1}$. Since $\ell_n=L^n\eta$,
\begin{equation}\label{4.26:eq-literal-block-side-1}
\frac{\mathrm{U2}}{\mathrm{U1}}=\frac{\ell_{p\wedge m}}{\ell_{p\wedge\lambda}}
=L^{(p\wedge m)-(p\wedge\lambda)},
\end{equation}
which is the ordering relation of Lemma~\ref{4.26:lem-unit-ordering}. If $\lambda=m$ then the level
gap is $0$, so $\mathrm{U1}=\mathrm{U2}$.

On $\mathfrak{T}$ no stage changes the determining set, by the freezing statement for the region
excluded from the stages. Hence every structure of the present site on $\mathfrak{T}$ has
$\lambda_{\mathfrak{s}}=m$. At the frozen structure, the structure $\mathfrak{s}$ of
every triple is the frozen structure itself, whose level at $y'$ is $\ell_{\mathfrak{a}}(y')=m$; hence
$\lambda=m$ there as well.

Suppose $\lambda=m$. Since $\omega\ge1$, the definition $\mathfrak{c}(y')=m-\tfrac12\log_{L_0}\omega$
gives $\mathfrak{c}(y')\le m=\lambda$, so the weight relation of Lemma~\ref{4.26:lem-unit-ordering}
holds with equality, $\omega=L_0^{2(m-\lambda)}w_b=w_b$; hence
$\mathrm{U3}=\mathrm{U1}=\mathrm{U2}$.

A difference $\mathrm{U1}\ne\mathrm{U2}$ requires a positive level gap, hence $\lambda<m$. By the two
preceding paragraphs, this occurs only at the retained structures of the family
$(\Sigma^{\ast})^{\mathfrak{a}}$ on a live piece.

(b) A shell clause is carried from one stage to the next by the increment lemma for the bump shells.
That lemma adds to the radius of the clause at $y'$ the bump-shell radius of entry $e$ at frame level
$m$ and block level $\min\{p,\lambda\}$, times $\pi$. By Definition~\ref{4.26:def-three-units} this
radius equals $\mathrm{U2}/\omega$, so the added term is $\pi\,\mathrm{U2}/\omega$.

Measured in the unit of a clause written in $\mathrm{U1}$, this additive term is, by
\eqref{4.26:eq-literal-block-side-1},
\begin{equation}\label{4.26:eq-literal-block-side-2}
\frac{\pi\,\mathrm{U2}}{\omega\,\mathrm{U1}}=\frac{\pi L^{(p\wedge m)-(p\wedge\lambda)}}{\omega}.
\end{equation}

Consider a point $y'$ of a thin stratum of a live piece, that is, a point of the stratum read at the
stage structure $\mathbb{B}_{k_a}$ at which $\mathfrak{c}(y')=k_0$ and $\omega=L_0^{2g}$, with
$g=g(y')\ge1$ the lift, together with a triple whose structure has $\lambda=\lambda_{\mathfrak{s}}(y')=k_0$.
The level $\mathfrak{c}$ is
$\mathfrak{c}(y')=m-\tfrac12\log_{L_0}\omega$, so $m-\lambda=\tfrac12\log_{L_0}\omega=g$. It suffices
to exhibit one triple at which propagation fails; take $p\ge m$, so that the level gap is
$m-\lambda=g$. The weight base satisfies $L_0^{2}\le L/3$, so $\omega=L_0^{2g}\le(L/3)^g$, and the
quantity \eqref{4.26:eq-literal-block-side-2} then satisfies
\[
\frac{\pi L^{g}}{\omega}\;\ge\;\pi L^{g}\Bigl(\frac{3}{L}\Bigr)^{g}=\pi\cdot3^{g}.
\]

To propagate the clause, this term must be absorbed by the stage increment $\iota_n$. The stage
increments are fixed at the value $\iota_n=3.88\pi$, and the bound established for them is only
the inequality $2.4\pi\le\iota_n$. For $g\ge1$ we have $\pi\cdot3^{g}\ge3\pi$, and $3\pi>2.4\pi$, so
this bound does not cover the term. Read in the unit $\mathrm{U1}$, absorption would require, even
with the value $\iota_n=3.88\pi$, that $\pi\cdot3^{g}\le3.88\pi$.
\begin{itemize}
\item For $g\ge2$ absorption is impossible, since $\pi\cdot3^g\ge9\pi$, which exceeds
$\iota_n=3.88\pi$.
\item For $g=1$ the term is $\pi L/L_0^2$. The requirement $\pi L/L_0^2\le3.88\pi$ reads
$L_0^2/L\ge1/3.88$. The inequality $L_0^2\le L/3$ reads
$L_0^2/L\le1/3$. Propagation at $g=1$
would thus need a condition bounding $L_0^2/L$ from both sides. Its lower half, $L\le3.88\,L_0^2$,
bounds $L$ from above, that is, it is a cap in the sense of Definition~\ref{2.3:def-cap}; the
conditions on $L$ in the quantifier prefix of Notation~\ref{2.1:not-prefix} are one-sided lower
bounds, and no cap is among them. So propagation is not available at $g=1$ either.
\end{itemize}

(c) By (b), the unit the increment lemma propagates at a retained structure is $\mathrm{U2}$. The source
space $\mathfrak{S}_{\mathrm{src}}$ therefore writes its composite clause at $y'$ with radius
$F^{(k)}(y')\cdot\mathrm{U2}$, where $F^{(k)}(y')$ is the source factor.

Read literally in $\mathrm{U1}$, the shell clause of $\mathfrak{N}^{(n)}(Y)$ of
Notation~\ref{4.26:not-stage-space} has radius
\[
\Bigl(1+2\beta-\sum_{l\le n}\iota^1_l(y')\Bigr)\mathrm{U1}.
\]
Notation~\ref{4.26:not-stage-space} records the value
\[
1+2\beta-2\mathsf{w}_0\beta=1.383 ,
\]
which the bracket takes when the increments $\iota^1_l(y')$ total $2\mathsf{w}_0\beta$; since the
increments $\iota^1_l=\tfrac23\mathsf{w}_0\rho^{(l)}$ are non-negative, the bracket never exceeds
$1+2\beta$. We read the clause at radius $1.383\cdot\mathrm{U1}$; the computation below can be carried
in the same way with $1+2\beta$ in place of $1.383$, the resulting inequalities being those below
with $1.383$ replaced by $1+2\beta$.
The inclusion of $\mathfrak{S}_{\mathrm{src}}$ in $\mathfrak{N}^{(n)}(Y)$, clause by clause, would
then require $F^{(k)}(y')\,\mathrm{U2}\le1.383\cdot\mathrm{U1}$. By \eqref{4.26:eq-literal-block-side-1}
this is
\begin{equation}\label{4.26:eq-literal-block-side-3}
F^{(k)}(y')\,L^{(p\wedge m)-(p\wedge\lambda)}\le1.383 .
\end{equation}

At level gap $1$, inequality \eqref{4.26:eq-literal-block-side-3} reads $F^{(k)}(y')\le1.383/L$. It
fails for every value of $F^{(k)}(y')$ in the interval $\,]1.383/L,\tfrac12]$.

At level gap at least $2$ and $F^{(k)}(y')\ge\varphi_{\mathrm{arrest}}=0.098$, the left side of
\eqref{4.26:eq-literal-block-side-3} satisfies
\[
F^{(k)}(y')\,L^{(p\wedge m)-(p\wedge\lambda)}\;\ge\;\varphi_{\mathrm{arrest}}L^{2}=0.098\,L^{2}.
\]
So \eqref{4.26:eq-literal-block-side-3} fails for every $F^{(k)}(y')\ge\varphi_{\mathrm{arrest}}$ as soon
as $0.098\,L^{2}>1.383$. In any case the failure at level gap $1$ already witnesses (c).

These failures compare the unit $\mathrm{U2}$, in which $\mathfrak{S}_{\mathrm{src}}$ is written, with
the unit $\mathrm{U1}$ in which $\mathfrak{N}^{(n)}(Y)$ would be literally read. They are not a bound
violated by the real background. Items (b) and (c) are clauses (ii) and (iii) of the proposition
comparing the units $\mathrm{U1}$ and $\mathrm{U2}$ at retained structures, and the failure in (c) is
the example given in its clause (iii).

For these reasons no clause in what follows is read in $\mathrm{U1}$. In particular, no bound of the
form $Q_e^{(p,X',\mathfrak{s})}(y')\le\tfrac12\varphi_{\mathrm{arrest}}\,\mathrm{U1}$ is used. Such a bound
would require the
normalisation of the current in \cite[(1.16)]{B-CMP109} together with the averaging weights of
\cite{B-CMP98}.
\end{proof}

\begin{definition}[Ownership rule for composite clauses at retained structures]
\label{4.26:def-ownership-rule}
Let $\mathfrak{a}$ be a live arrest, with region $W_{\mathfrak{a}}$ and class $k_a$. Let
$(p,X',\mathfrak{s})$ be a $(\Sigma^{\ast})^{\mathfrak{a}}$-triple whose structure $\mathfrak{s}$
is retained from the run of $\mathfrak{a}$ in the sense of Lemma~\ref{4.26:lem-retained-levels},
and let $y'\in X'^{\sim 2}\cap W_{\mathfrak{a}}$, a point of $\mathfrak{T}$ or of
$Z^{\mathrm{on}}$. Let $m$, $\omega$, $\lambda=\lambda_{\mathfrak{s}}(y')\le m$, $w_b$,
$\mathfrak{c}=\mathfrak{c}(y')$ and the units $\mathrm{U1}$, $\mathrm{U2}$, $\mathrm{U3}$ be those
of Definition~\ref{4.26:def-three-units}, display~\eqref{4.26:eq-three-units-a}. Let $g(y')$
denote the lift of $y'$, the exponent of the conversion~\eqref{4.26:eq-unit-ordering-b} of
Lemma~\ref{4.26:lem-unit-ordering}. For
$e\in\{2,4\}$ we write $\mathrm{U2}_e$, $\mathrm{U3}_e$ for these units at entry $e$ (for $e=4$
they are $\mathrm{U2}$, $\mathrm{U3}$; for $e=2$ the last length is dropped). We put
$r_e(m,l):=\alpha_{0,m}\ell_m^{-2}\ell_{l}^{-1}$ for $e=4$ and $r_e(m,l):=\alpha_{0,m}\ell_m^{-2}$
for $e=2$, so that $\mathrm{U2}_e=\omega\,r_e(m,l_{y'})$ with $l_{y'}=\min\{p,\lambda\}$, the
re-blocked shell radius at the frozen level with the block side of $\mathfrak{s}$. A
\emph{composite clause} at the triple and at $y'$ is a bound
$Q^{(p,X',\mathfrak{s})}_e(y')<\theta\,\mathfrak{U}$, $e\in\{2,4\}$, with a factor $\theta$ and a
unit $\mathfrak{U}$; its \emph{typing} is the pair formed by the level at which $\theta$ is
evaluated and the unit $\mathfrak{U}$.

\smallskip
\noindent(R-foreign) The following spaces and letters, written after the arrest of $\mathfrak{a}$
by a run other than the run of $\mathfrak{a}$, are called \emph{foreign} to the retained
structures of $\mathfrak{a}$; they form the list of spaces written after an arrest in
Notation~\ref{4.26:not-arrest-conventions}, as it is quoted in
Lemma~\ref{4.26:lem-retained-levels}: the space $\mathfrak{N}^{(n)}$ of the current stage;
the sources $\tilde{U}^{\mathrm{arrest}}_{k'}$, at points of $\mathfrak{T}$ and of $Z^{\mathrm{on}}$
alike; the spaces written after the operation $\mathbf{T}$ and the letters $\mathbf{B}^{(k')}$ of
Notation~\ref{4.26:not-arrest-conventions}, of every class $k'>k_a$; the inherited domains
$\mathfrak{D}^{\flat}$ of foreign later sites; the attempt spaces $\mathfrak{L}^{(k_b)}_l$ of an
enclosing run ($\mathfrak{b}\supset\mathfrak{a}$) at points of $W_{\mathfrak{a}}$ and, once they
are stored, its class-$k_b$ letters $\mathfrak{F}(\mathfrak{b})$; the spaces of later core sites;
and the tubes. A foreign space reads a composite clause at a retained structure of
$\mathfrak{a}$, at a point of $\mathfrak{T}$ or of $Z^{\mathrm{on}}$, at the frozen level and
weight of the point, with the block side of the re-blocked shell radius: in the unit
$\mathrm{U2}_e$, with its own factor evaluated at the frozen level $m$. These factors are:
for $\mathfrak{N}^{(n)}$, $1+2\beta-\sum\iota^1$ at points of $\mathfrak{T}$, while at
$Z^{\mathrm{on}}$ the space $\mathfrak{N}^{(n)}$ imposes no composite clause; for the source
$\tilde{U}^{\mathrm{arrest}}_{k'}$, $F^{(k')}_m$; for $\mathbf{B}^{(k')}$, $F^{(k')}_m$; for the
attempt space $\mathfrak{L}_l$ of a later site, $\theta^{(l),(k_b)}(m)$; for
$\mathfrak{F}(\mathfrak{b})$, $\theta^{(i),(k_b)}(m)$, with the index $i$ of the stored letter.
Two foreign spaces meeting at one retained triple carry the same frozen $(\ell,w,\mathcal{Q})$ and
are compared by their factors at $m$.

\smallskip
\noindent(R-own) The spaces and letters of the run that built the retained structure, namely the
class-$k_a$ letters ($\mathfrak{F}(\mathfrak{a})$, the letters $\mathfrak{F}(\mathfrak{b}')$ of
sibling arrests $\mathfrak{b}'$, the post-$\mathbf{R}$ stage structure $\mathbb{B}''_{k_a}(W')$ of
class $k_a$ on a region $W'$, and the remaining class-$k_a$ letters) and the attempt letters
$\mathfrak{L}^{(k_a)}_l=\mathfrak{S}^{(k_a),\mathrm{arrest}}_l$ of site $k_a$ on its excluded region
$W$, that is, each run's own retained stages inside its own attempt letter, are called
\emph{own} to $\mathfrak{s}$. An own space reads a composite clause at $\mathfrak{s}$ in the birth
radius $\mathrm{U3}_e$ of $\mathfrak{s}$, with the running threshold of its run evaluated at
$\lambda$; this is the reading $(\Sigma^{\flat})$, in which the entries of retained stages carry
their own shells, radii and gains with the running threshold, the reference level being the
level $\mathsf{m}^{\circ}(y')=\lambda$ of the parent structure.

\smallskip
\noindent\emph{Ownership rule.} The typing of a composite clause at a retained triple is
determined by the pair (class of the imposing space, run in which the structure was born), and
not by the location of the point: for an imposing space $\mathfrak{V}$,
\begin{equation}\label{4.26:eq-ownership-rule-a}
Q^{(p,X',\mathfrak{s})}_e(y')<
\begin{cases}
\theta(\lambda)\,\mathrm{U3}_e, &
\text{if }\mathfrak{V}\text{ belongs to the run that built }\mathfrak{s},\\[2pt]
\theta_{\mathfrak{V}}(m)\,\mathrm{U2}_e, & \text{otherwise},
\end{cases}
\end{equation}
where $\theta(\lambda)$ is the running threshold of the run at $\lambda$ and
$\theta_{\mathfrak{V}}(m)$ is the factor at $m$ of the imposing space $\mathfrak{V}$, listed in
(R-foreign).
A stored space keeps the typing with which it was written.

\smallskip
\noindent The rule has the following consequences, proved below.
\begin{itemize}
\item[(c1)] At points of $\mathfrak{T}$ the source is contained in $\mathfrak{N}^{(n)}$, clause by
clause, in the one unit $\mathrm{U2}$. In particular the source $\tilde{U}^{\mathrm{arrest}}_{k'}$
implies the shell clause $(\mathrm{box})^{\mathrm{sh}}$ of $\mathfrak{N}^{(n)}$ in the shell unit.
\item[(c2)] At retained triples the source is contained in every class letter
($\mathfrak{F}(\mathfrak{a})$, the remaining class-$k_a$ letters, the letters
$\mathfrak{F}(\mathfrak{b}')$ of siblings, $\mathbb{B}''_{k_a}$: the own attempt spaces of the
piece as stored letters), and at $\lambda<m$ the source is at most one third of the letter's
threshold. The attempt letter $\mathfrak{L}^{(k_b)}_l$ of an enclosing run at a core piece is
foreign to the structures of $\mathfrak{a}$, is typed $(m,\mathrm{U2})$ with factor
$\theta^{(l),(k_b)}(m)$, and
lies inside $\tilde{U}^{\mathrm{arrest}}_{k_b}$, the class-$k_a$ letters, the letters
$\mathbf{B}^{(k')}$ and the inherited domains.
\item[(c3)] A composite entry with $Q_e<1.4\,\mathrm{U2}$ propagates under one stage move with
increment at most $\iota_n\mathrm{U2}$.
\item[(c4)] The clause $(\mathrm{box})^{\mathrm{sh}}$ at the box triple
$(p,\square_0,\mathfrak{B}^{\mathfrak{c}}\restriction\square_0)$ is read in $\mathrm{U2}$, and it
implies the shell clause of Ba{\l}aban's construction in the unit of the $\mathfrak{c}$-shell with
factor $1.4(0.381/L)^{g(y')}\le 0.0445<1+2\beta$.
\item[(c5)] At the box triple, for every term $T$ read there whose level satisfies
$\ell_T\le\mathfrak{c}$, the source threshold satisfies
$\theta_{\mathrm{src}}\le 0.381^{g(y')}\,\mathbf{r}^{T,(p)}$, and the room at the start of the box
step is at least $0.8-0.381>0.4$.
\item[(c6)] At points of $\mathfrak{T}$ the composite entry $2$ at the box triple is bounded by
$1.4\,\mathrm{U2}_2$, a unit with no block side, with $\omega\,\alpha_{0,m}\le\gamma_0^{(k)}/5$;
the kernel clauses of the box step hold at the box triple without any block side entering;
where a block current of $\mathbb{U}^w_{\square}$ is read at the block side of $\mathfrak{B}$,
the one factor $L$ at $p_{\square}=\mathfrak{c}(y')+1\le m$ is supplied from the direct entries
of $\mathbb{U}$.
\item[(c7)] Each comparison of the source's $\mathrm{U2}$ with the $\mathrm{U3}$ of a letter, and
each comparison with an unweighted $\mathbf{r}^T$ at a level $\ell_T\le\lambda$, satisfies
respectively
\begin{equation*}
\frac{\mathrm{U2}}{\mathrm{U3}}\le\Bigl(\frac{(1+\beta_0)L_0^2}{L^2}\Bigr)^{m-\lambda}\le 1,
\qquad
\frac{\mathbf{r}^T}{\mathrm{U2}}\ge(2.625L)^{m-\ell_T};
\end{equation*}
at $\lambda<m$ the room is at least $2.087\,\mathrm{U2}$, and tube increments are at most a quarter
of the room.
\item[(c8)] The unit $\mathrm{U2}$ is not the term-level unit $\mathbf{r}^{(p)}$ excluded in
Notation~\ref{4.26:not-stage-space}.
\item[(c9)] At the run's own raised structures the enlarged composite clauses have the form stated
for them, in $\mathrm{U3}$ with threshold $\theta^{(l)}(\lambda)$. At a lifted point of an older
piece in the core the composite clause is $\theta^{(l),(k)}(m)\,\mathrm{U2}$, the one-step tube
closes in $\mathrm{U2}$, the attempt space $\mathfrak{L}_l$ enlarged by the accumulated stage
increments lies in $\tilde{U}^{\mathrm{arrest}}_k$, and $\tilde{U}^{\mathrm{arrest}}_k$ lies in every
older letter. For a localization domain $\hat X\subset Z$ on which the stage spaces are compared
with the kernel-regularity space $\mathfrak{K}_P$, the chain
$\mathfrak{K}_P\supset\mathfrak{N}(\hat X)\supset\tilde{U}^{(n+1)c,\Sigma\flat}_k(\hat X)$ holds,
and so does the inclusion in $\mathfrak{N}$ of the attempt spaces
$\mathfrak{S}^{(k),\mathrm{arrest}}_l$ written after the arrest.
\end{itemize}
\end{definition}

\begin{proof}[Proof of the consistency of the rule and of (c1)--(c9)]
\emph{Comparison of two foreign spaces.} By the convention of
Notation~\ref{4.26:not-arrest-conventions}, every space written after the arrest reads at
$x\in W_{\mathfrak{a}}$ the frozen triple $(\ell,w,\mathcal{Q})(x)$ of the innermost live arrest
holding $x$. Two foreign spaces meeting at one retained triple therefore carry the same frozen
level, weight and cube clauses, and by \eqref{4.26:eq-ownership-rule-a} both read the composite
clause in the same unit $\mathrm{U2}_e$; they are compared by comparing their factors at $m$, as
in the extended source-ratio lemma.

\emph{Consistency of the rule.} Every class-$k_a$ space precedes every later space. An enclosing
run $\mathfrak{b}$ never changes the level of a point of $W_{\mathfrak{a}}$, because
$W_{\mathfrak{a}}\subset Z''_{h_b+1}$. Hence the two cases of
\eqref{4.26:eq-ownership-rule-a} never apply to one space at one pair consisting of a point and a
structure. In particular, at the retained triples of a core piece the attempt space
$\mathfrak{L}^{(k_b)}_0$ coincides with $\tilde{U}^{\mathrm{arrest}}_{k_b}$ clause by clause: the
stage-zero attempt space of the enclosing run is the source of its class by its definition, and
both are foreign, hence both are typed at $(m,\mathrm{U2})$. Likewise the inclusion of the image of
$\mathfrak{L}_{l+1}$ in $\mathfrak{L}_l$ there is the inclusion of clause (ii) of the nesting lemma
for attempt spaces in the form it takes on an excluded region, read in $\mathrm{U2}$ with room
$\rho^{(l+1)}_m\mathrm{U2}$; both sides are in the same unit, so the conversion factor is $1$.

\emph{(c1).} The source $\tilde{U}^{\mathrm{arrest}}_{k'}$ and $\mathfrak{N}^{(n)}$ are both foreign,
so at a point of $\mathfrak{T}$ both composite clauses are read in $\mathrm{U2}$ and are compared
by factor. The source factor $F^{(k')}_m$ of Notation~\ref{4.26:not-arrest-conventions} satisfies
$F^{(k')}_m\le 1$, and the factor of $\mathfrak{N}^{(n)}$ at $\mathfrak{T}$ is
$1+2\beta-\sum\iota^1$, which by Notation~\ref{4.26:not-stage-space} is at least
$1+2\beta-2\mathsf{w}_0\beta=1.383$. Since $1<1.383$, each composite clause of the source implies
the corresponding clause of $\mathfrak{N}^{(n)}$; this is the inclusion of the source in
$\mathfrak{N}^{(n)}$ in the list of inclusions of the source. Since the unit of
$(\mathrm{box})^{\mathrm{sh}}$ is the shell radius $r_e^{(p)}(y')$ of
Notation~\ref{4.26:not-stage-space}, that is, the re-blocked radius whose weighted form is
$\mathrm{U2}_e$, the source implies $(\mathrm{box})^{\mathrm{sh}}$ in the shell unit.

\emph{(c2).} The class letters listed are own to the retained structures, so by
\eqref{4.26:eq-ownership-rule-a} they are read in $\mathrm{U3}$ with the running threshold
$\theta(\lambda)$, while the source is read in $\mathrm{U2}$ with factor $F=F^{(k')}_m$. The
comparison is therefore by ratio, of the same kind as clause (b) of the extended source-ratio
lemma, by which at raised points the source is at most $0.74$ of a term's threshold. Let
$\lambda<m$. We use $F\le 1$; the inequality $\mathrm{U2}\le\mathrm{U3}/31.5$ at $\lambda<m$ from
the comparison of the three units (Lemma~\ref{4.26:lem-unit-ordering}); the identity
$31.5\,\varphi_{\mathrm{arrest}}=3.087$, with $\varphi_{\mathrm{arrest}}=0.098$; and the bound
$\theta(\lambda)\ge\varphi_{\mathrm{arrest}}$ for every threshold, which is the floor lemma for the
arrested factors and does not depend on the level. These give
\begin{equation}\label{4.26:eq-ownership-rule-1}
F\,\mathrm{U2}\le\mathrm{U2}\le\frac{\mathrm{U3}}{31.5}
=\frac{\varphi_{\mathrm{arrest}}}{3.087}\,\mathrm{U3}
\le\frac{\theta(\lambda)}{3.087}\,\mathrm{U3}<\frac13\,\theta(\lambda)\,\mathrm{U3},
\end{equation}
so the source is at most one third of the letter's threshold. At $\lambda=m$ the inclusion is
clause (a) of the extended source-ratio lemma as stated there, together with the complementary
sentence (a$'$) that accompanies that clause.

The attempt letter $\mathfrak{L}^{(k_b)}_l$ of an enclosing run at a core piece belongs to a
run other than that of $\mathfrak{a}$, so by (R-foreign) it is foreign to the structures of
$\mathfrak{a}$ and is typed $(m,\mathrm{U2})$ with factor $\theta^{(l),(k_b)}(m)$. It lies inside
$\tilde{U}^{\mathrm{arrest}}_{k_b}$, the class-$k_a$ letters, the letters $\mathbf{B}^{(k')}$ and
the inherited domains by the inclusions established in (c9) below. The real point lies inside it,
since the composite entries of the real point are bounded by
$\tfrac12\varphi_{\mathrm{arrest}}\,\mathrm{U2}\le\tfrac12\theta^{(l)}(m)\,\mathrm{U2}$ (see the real
background, the corresponding display).

\emph{(c3).} Work in the frame $(X',\mathfrak{B},\mathfrak{s}^{(p)})$ of the general clause of the
re-blocked shell radius; here $\mathfrak{s}^{(p)}$ is as in that general clause, the frame
structure $\mathfrak{B}$ has level $\mathsf{m}=m$ and the parent structure $\mathfrak{s}$ has level
$\mathsf{m}^{\circ}=\lambda\le m$ (the re-blocked case). The propagation estimate for
composite entries under one stage move states: if $Q\le 1.4\,\omega r_e$ then the move is at most
$(1.4\pi_{\times}\omega+\pi)r_e\le 2.4\pi\,\omega r_e$. We apply it with $r_e=r_e(m,l_{y'})$, so
that $\omega r_e=\mathrm{U2}_e$. If $Q_e<1.4\,\mathrm{U2}$, then the updated entry $Q^+_e$ satisfies
\begin{equation}\label{4.26:eq-ownership-rule-2}
\begin{split}
Q^+_e&\le Q_e(1+\pi_{\times})+\pi\,\frac{\mathrm{U2}}{\omega}
\le Q_e+(1.4\pi_{\times}\omega+\pi)\,r_e\\
&\le Q_e+2.4\pi\,\mathrm{U2}\le Q_e+\iota_n\mathrm{U2},
\end{split}
\end{equation}
where the second inequality uses $Q_e\pi_{\times}<1.4\pi_{\times}\mathrm{U2}=1.4\pi_{\times}\omega
r_e$ and $\mathrm{U2}/\omega=r_e$, the third is the quoted estimate, and the last uses
$2.4\pi\le\iota_n$, which holds since the stage increment has the value $\iota_n=3.88\pi$ (see also
Remark~\ref{4.26:rem-literal-block-side}).

\emph{(c4).} The clause (iv) of the hypothesis of the box construction is the clause
$(\mathrm{box})^{\mathrm{sh}}$ at the box triple
$(p,\square_0,\mathfrak{B}^{\mathfrak{c}}\restriction\square_0)$ in the frozen weighted unit,
which by \eqref{4.26:eq-ownership-rule-a} is $\mathrm{U2}$, the clause being foreign. At the
box triple $\lambda=\mathfrak{c}$, and $p\le p_{\square}\le\mathfrak{c}(y')+1$. By
Definition~\ref{4.26:def-three-units},
$\mathrm{U2}/\mathrm{U1}=\ell_{p\wedge m}/\ell_{p\wedge\lambda}$. For $p\le\mathfrak{c}(y')$ both
minima equal $p$, so $\mathrm{U2}=\mathrm{U1}$; for $p=\mathfrak{c}(y')+1\le m$ the ratio is
$\ell_{\mathfrak{c}+1}/\ell_{\mathfrak{c}}=L$, so $\mathrm{U2}=L\,\mathrm{U1}$. The clause (c$_0$)
of that hypothesis states that the box clause implies the shell clause of Ba{\l}aban's
construction in the unit of the
$\mathfrak{c}$-shell. By the conversion \eqref{4.26:eq-unit-ordering-b},
$\mathrm{U2}/\mathbf{r}^{\mathfrak{c},(p)}_4\le(0.381/L)^{g(y')}$, so the factor $1.4$ in
$\mathrm{U2}$ becomes the factor $1.4(0.381/L)^{g(y')}\le 0.0445$ in
$\mathbf{r}^{\mathfrak{c},(p)}_4$, and $0.0445<1+2\beta$.

\emph{(c5).} At the box triple the source threshold is $\theta_{\mathrm{src}}=F\,\mathrm{U2}$. By
$F\le1$ and \eqref{4.26:eq-unit-ordering-b},
$F\,\mathrm{U2}\le(0.381/L)^{g(y')}\,\mathbf{r}^{\mathfrak{c},(p)}_4$. Let $T$ be a term read at the
box triple with $\ell_T\le\mathfrak{c}$; the unit-comparison rule for radii gives
\begin{equation}\label{4.26:eq-ownership-rule-3}
\frac{\mathbf{r}^{\mathfrak{c},(p)}_4}{\mathbf{r}^{T,(p)}}
\le\Bigl(\frac{1+\beta_0}{L^2}\Bigr)^{\mathfrak{c}-\ell_T}\le 1 ,
\end{equation}
and hence
\begin{equation*}
\theta_{\mathrm{src}}\le\Bigl(\frac{0.381}{L}\Bigr)^{g(y')}\mathbf{r}^{\mathfrak{c},(p)}_4
\le 0.381^{g(y')}\,\mathbf{r}^{T,(p)}.
\end{equation*}
This is the bound at the start of the box step, which gives the room
$0.8-0.381>0.4$ in the unit $\mathbf{r}^{T,(p)}$.

\emph{(c6).} Let $y'\in\mathfrak{T}$. The plaquette variables of $\mathbb{U}^w_{\square}$ on
$\square_0^{\sim2}$ are the composite entry $2$ of $\mathbb{U}$ at the box triple. By
\eqref{4.26:eq-ownership-rule-a} this entry is bounded in $\mathrm{U2}_2$, which for $e=2$ carries
no block side:
\begin{equation*}
Q^{(p,\square_0,\mathfrak{B}^{\mathfrak{c}}\restriction\square_0)}_2(y')
<1.4\,\mathrm{U2}_2=1.4\,\omega\,\alpha_{0,m}\ell_m^{-2}
\quad\text{(in the unit of a plaquette)}.
\end{equation*}
It is small at the own levels of $\mathfrak{B}$, since
$\omega\,\alpha_{0,m}\le w_\ast(k_a)\bar{\alpha}_{k_a}\le\gamma_0^{(k)}/5$ by the bound of letter
weights against the analyticity size, where $\gamma_0^{(k)}$ is the size of the analyticity
polydisc $\mathfrak{K}_P$ at the current density. The kernel clauses of the box step are the
plaquette and $U$-regularity clauses of $\mathbb{U}$ at fixed radius and the bound
$\sup\ell^3|\mathbb{J}|\le\gamma_0^{(k')}$, for a space of class $k'$; neither reads a block
current of the input, so the
block side does not enter. Where a kernel clause of the box step bounds a block current of
$\mathbb{U}^w_{\square}$ at the block side of $\mathfrak{B}$, at
$p_{\square}=\mathfrak{c}(y')+1\le m$ the single factor $L$, that is, the bound
$1.4L\,\omega\,\alpha_{0,m}$, follows from the direct entries of $\mathbb{U}$ by the bound of
\cite{B-CMP109}, used as stated there, as on the two outer layers of $\square_0$ and at the core
pieces.

\emph{(c7).} The unit-comparison rule for radii, the comparison of the source with a class letter
at lifted points together with its corollary for the source factor, and clauses (4) and (5) of the
list of inclusions of the source each compare the source's $\mathrm{U2}$ with the $\mathrm{U3}$ of
a letter or with an unweighted $\mathbf{r}^T$ at a level $\ell_T\le\lambda$. In the first case the
comparison display of the unit-comparison rule with its exponent parameter $p_a=2$ bounds the
ratio $\mathrm{U2}/\mathrm{U3}$ by
\begin{equation*}
\Bigl(\frac{(1+\beta_0)L_0^2}{L^2}\Bigr)^{m-\lambda}\le 1;
\end{equation*}
in the second, clause (5) bounds the ratio $\mathbf{r}^T/\mathrm{U2}$ from below by
\begin{equation*}
(2.625L)^{m-\ell_T}\ge 2.625^{g(y')}.
\end{equation*}
Each of these comparisons therefore holds exactly in the
form stated there or a fortiori. By the tube-increment proposition ($P$ its constant), which
holds for every structure $\mathfrak{s}$, the tube increments are at most $P\,2^{-n_W}$ times the
unweighted unit, and the unweighted unit is at most $\mathrm{U2}$; hence the increment is at most
$P\,2^{-n_W}\mathrm{U2}\le P\,\mathrm{U2}$. At $\lambda<m$ the room is, by
$\theta(\lambda)\ge\varphi_{\mathrm{arrest}}$,
$\mathrm{U3}\ge31.5\,\mathrm{U2}$ and $F\le1$,
\begin{equation*}
\theta(\lambda)\,\mathrm{U3}-F\,\mathrm{U2}\ge(31.5\,\varphi_{\mathrm{arrest}}-1)\,\mathrm{U2}
=2.087\,\mathrm{U2}.
\end{equation*}
So the increment is at most a quarter of the room $2.087\,\mathrm{U2}$ provided
$16P\le\beta\theta_S$, which is one of the conditions fixed with the source space
$\mathfrak{S}_{\mathrm{src}}$; the bound holds a fortiori.

\emph{(c8).} The unit excluded in Notation~\ref{4.26:not-stage-space}, the term-level unit
$\mathbf{r}^{(p)}$, places the factor $\ell^{-2}$ at the level $p$ of the triple. The unit
$\mathrm{U2}$ keeps level and weight at the frozen level $m$ of the point and changes only the
block side, to the level $\lambda_{\mathfrak{s}}$ of the structure, which is exactly the
$l_{y'}=\min\{p,\mathsf{m}^{\circ}(y')\}$ of the re-blocked shell radius. Hence $\mathrm{U2}$ is
not $\mathbf{r}^{(p)}$, and the exclusion in Notation~\ref{4.26:not-stage-space} remains in force.

\emph{(c9).} At the run's own raised structures the clauses
are own, so by \eqref{4.26:eq-ownership-rule-a} the enlarged composite clauses have the
form stated for them: the composite clause in $\mathrm{U3}$ with threshold $\theta^{(l)}(\lambda)$,
the start bounds with the constants $0.3$ and $0.75$, and the one-step tube in $\mathrm{U3}$.

At a lifted point of an older piece in the core the clauses are foreign, and
\eqref{4.26:eq-ownership-rule-a} gives the following. The composite clause is
$\theta^{(l),(k)}(m)\,\mathrm{U2}$, and $\mathfrak{L}_0=\tilde{U}^{\mathrm{arrest}}_k$ there by the
definition of $\mathfrak{L}_0$, as shown in the consistency paragraph. The start follows from the
unit conversion of (c5), with room $0.8-0.381=0.419$. The kernel clauses follow from the direct
entries, as in (c6). The one-step tube, the inclusion of the image of $\mathfrak{L}_{l+1}$ in
$\mathfrak{L}_l$, compares like with like in $\mathrm{U2}$: the room is
$\rho^{(l+1)}_m\mathrm{U2}$, because $s^{(k)}_m$ and every $\sigma_{\mathfrak{a}}(m)$ entering
$F^{(k)}_m$ (Notation~\ref{4.26:not-arrest-conventions}) are independent of $l$; the increment is
at most $2.4\pi\,\mathrm{U2}\le\iota_{l+1}\mathrm{U2}$ by the propagation estimate of (c3) with
$r_e=r_e(m,l_{y'})$; and the inclusion closes provided $\iota_{l+1}\le\tfrac12\rho^{(l+1)}_m$,
which is clause (ii) of the nesting lemma for attempt spaces. Writing $\mathfrak{L}_l\oplus\sum\iota$
for $\mathfrak{L}_l$ enlarged by the accumulated stage increments $\sum\iota$, we
obtain
\begin{equation}\label{4.26:eq-ownership-rule-4}
\mathfrak{L}_l\oplus\textstyle\sum\iota\subset\tilde{U}^{\mathrm{arrest}}_k,
\end{equation}
and $\tilde{U}^{\mathrm{arrest}}_k$ is contained in every older letter; both inclusions follow
from (c1), (c2) and clauses (4)(b), (4)(d) and (d$'$) of the list of inclusions of the source.

Finally, for a localization domain $\hat X\subset Z$ as in (c9), the chain
$\mathfrak{K}_P\supset\mathfrak{N}(\hat X)\supset\tilde{U}^{(n+1)c,\Sigma\flat}_k(\hat X)$ and the
inclusion in $\mathfrak{N}$ of the attempt spaces $\mathfrak{S}^{(k),\mathrm{arrest}}_l$ written after
the arrest hold. Indeed
$\hat X\subset Z$; at $Z^{\mathrm{on}}$ the space $\mathfrak{N}$ imposes no composite clause,
because $(\mathrm{box})^{\mathrm{sh}}$ is quantified over points of $\mathfrak{T}$ only
(Notation~\ref{4.26:not-stage-space}), so only direct entries are compared, by the comparison of
direct entries at $Z^{\mathrm{on}}$; and the own retained stages of the present site have
$\lambda=m$ on $\mathfrak{T}$. On pieces in $\mathfrak{T}$ the inclusion used is (c1).
\end{proof}

\begin{lemma}[The bracket at a lower level]\label{4.26:lem-lower-bracket}
Work at the site $k$, with the stage data of Notation~\ref{4.26:not-stage-frozen-data} and the
conventions after an arrest of Notation~\ref{4.26:not-arrest-conventions}. In particular
$s^{(k')}_\nu=1-\beta\bigl(1-2^{-(k'-\nu)}\bigr)$ and $\varphi_{\mathrm{arrest}}=1-4.51\beta$,
where $\beta=\tfrac15$, so that $\varphi_{\mathrm{arrest}}=0.098$.

Let $\mathfrak{a}$ be a live arrest and $x\in W_{\mathfrak{a}}$. Let
$k_{a_1}<k_{a_2}<\cdots$ be the levels of the live arrests holding $x$ (the live chain), put
$\ell:=\ell_{\mathfrak{a}}(x)$, and assume $\ell\le k_{a_1}$.
For a live arrest $\mathfrak{d}\ni x$, let $I(\mathfrak{d})$ be its subtracted range, so that
$\sigma_{\mathfrak{d}}(\nu)=\sum_{i\in I(\mathfrak{d})}2^{-|\nu-i|}$ is the sum entering the
definition of the source factor $F^{(k')}(x)$ in
Notation~\ref{4.26:not-arrest-conventions}. The subtracted range $I(\mathfrak{a})$ of the arrest
$\mathfrak{a}$ itself contains every integer of $[h_a+1,k_a+1]$, with $h_a$ as in (c) below
(see the definition of $\sigma_{\mathfrak{a}}$ accompanying the source factor in
Notation~\ref{4.26:not-arrest-conventions}).
For an integer level $\lambda$ put
\begin{equation}\label{4.26:eq-lower-bracket-1}
\begin{gathered}
\varphi^{(l),(k)}_\lambda:=s^{(k)}_\lambda-\beta\sum_{i=h+1}^{h+l}2^{-|\lambda-i|},\qquad
\theta^{(l)}_{\mathfrak{s}}(\lambda):=\varphi^{(l),(k)}_\lambda
-\beta\!\!\sum_{\mathfrak{d}\text{ live}\ni x,\ k_d<k}\!\!\sigma_{\mathfrak{d}}(\lambda),\\
F^{(k')}_\lambda:=s^{(k')}_\lambda-\beta\sum_{\mathfrak{d}\text{ live}\ni x}\sigma_{\mathfrak{d}}(\lambda).
\end{gathered}
\end{equation}
The attempt-space brackets are the $\theta^{(l)}_{\mathfrak{s}}(\lambda)$; for a stage $n$ and a
term class $T$, the class-letter bracket $\theta^{(n)}_T(\lambda)$ is the threshold of the class
letter of $T$ at stage $n$, taken at level $\lambda$. When $\mathfrak{a}$ is the innermost live
arrest holding $x$, so that $\ell$ is the frozen level read at $x$ under the convention of
Notation~\ref{4.26:not-arrest-conventions} that every space written after an arrest reads at $x$
the frozen data of the innermost live arrest holding $x$, the definition of the source factor in
the same notation gives $F^{(k')}(x)=F^{(k')}_{\ell}$.

Assume the following.
\begin{itemize}
\item[(a)] The floor $F^{(k')}(x)\ge\varphi_{\mathrm{arrest}}$ among the arrest lemmas of
Notation~\ref{4.26:not-arrest-conventions} rests on the following inputs only, each of which
holds at every integer $\lambda\le k_{a_1}$:
\begin{itemize}
\item the leading terms satisfy $s^{(k)}_\lambda,s^{(k')}_\lambda\ge1-\beta$;
\item the subtracted ranges of the live chain and an attempt's own range $\{h+1,\dots,h+l\}$
are pairwise disjoint except for common endpoints;
\item the sum of $2^{-|\lambda-i|}$ over all these ranges, counted with multiplicity, is at most
$3+0.508$.
\end{itemize}
\item[(b)] The proof of the room bound
\begin{equation*}
\rho_n(x;k')=\theta^{(n)}_T(\ell)-F^{(k')}(x)\ \ge\ \beta2^{-(k_a+1-\ell)}
\end{equation*}
among the arrest lemmas of Notation~\ref{4.26:not-arrest-conventions} reads the level $\ell$ only
through the index ranges of the subtracted sums in $\theta^{(n)}_T$ and $F^{(k')}$, and these
ranges do not depend on $\ell$.
\item[(c)] At density $k_a$ the stage counts satisfy $(N_0^a+1)(N_0^a+3)\le N_a$, where
$h_a=k_a-N_a$ and $k_0=k_a-N_0^a$.
\end{itemize}
Then the following hold, with $\lambda\le\ell$ in (i) and $\nu\le k_a$ in (iii):
\begin{equation}\label{4.26:eq-lower-bracket-a}
\begin{gathered}
\text{(i)}\quad \theta^{(l)}_{\mathfrak{s}}(\lambda),\ \theta^{(n)}_T(\lambda),\
F^{(k')}_\lambda\ \ge\ \varphi_{\mathrm{arrest}};\\
\text{(ii)}\quad \text{bracket at }\lambda\ \le\ 1-\tfrac{\beta}{2}-2\beta
=1-\tfrac52\beta=\tfrac12;\\
\text{(iii)}\quad \rho_n(x;k';\nu):=\theta^{(n)}_T(\nu)-F^{(k')}_\nu\ \ge\
\beta2^{-(k_a+1-\nu)}>0.
\end{gathered}
\end{equation}
Clause (ii) holds for each bracket of \eqref{4.26:eq-lower-bracket-1} whose upper index
($k$, respectively $k'$) exceeds $\lambda$, and for which $\lambda-1,\lambda,\lambda+1$ lie in one
of its subtracted ranges. In particular:
\begin{itemize}
\item $F^{(k')}_\nu\le\tfrac12$ for $\nu\in[h_a+2,k_a]$ and $k'>\nu$.
\item Suppose that $x$ is a point of kind (A), in the classification of the points of the lifted set
$\mathsf{L}_{\mathfrak{a}}=(\Omega^c_{j^+}\setminus Z''_{k_0+1})\cap W_{\mathfrak{a}}$ into the
kinds (A), (B), (B$'$), and let
$\mathfrak{c}(x)$, $\lambda_{\mathfrak{s}}(x)$, $m(x)$ be the levels at $x$ of
\eqref{4.26:eq-retained-levels-a}. Suppose the levels
$\nu\in\{\mathfrak{c}(x),\lambda_{\mathfrak{s}}(x),m(x)\}$ lie in $[k_0,j^+]$, and $j^+\le k_a$.
Then $F^{(k')}_\nu\le\tfrac12$ for $k'>\nu$. If moreover every live arrest $\mathfrak{d}$ holding
$x$ satisfies $k_d\le h$, then also
$\theta^{(l)}_{\mathfrak{s}}(\nu)\le F^{(k')}_\nu\le\tfrac12$. These bounds at $x$ are asserted
only as upper bounds and only at lifted points.
\item Suppose that $x$ is a lifted point of kind (B) or (B$'$), at which $\mathfrak{c}(x)$ lies
among the layers of the run's own stages. If $\mathfrak{c}(x)-1,\mathfrak{c}(x),\mathfrak{c}(x)+1$
lie in one of the subtracted ranges of $\theta^{(l)}_{\mathfrak{s}}$ and $k>\mathfrak{c}(x)$, then
$\theta^{(l)}_{\mathfrak{s}}(\mathfrak{c}(x))\le\tfrac12$.
\end{itemize}
\end{lemma}

\begin{proof}
\emph{(i).} Let $\lambda\le\ell$ be an integer. Then $\lambda\le k_{a_1}$.

Each bracket in \eqref{4.26:eq-lower-bracket-1} has the form $s-\beta S$. Here $s$ is
$s^{(k)}_\lambda$ or $s^{(k')}_\lambda$, and $S$ is a sum of the nonnegative numbers
$2^{-|\lambda-i|}$. For $\theta^{(l)}_{\mathfrak{s}}(\lambda)$ the indices $i$ run over the
attempt's own range and the ranges $I(\mathfrak{d})$ with $k_d<k$. For $F^{(k')}_\lambda$ they run
over the ranges of the whole live chain. In both cases $S$ is a sub-sum of the sum in the third
input of (a), and so $S\le3+0.508$ at $\lambda$.

Since $2^{-(k'-\lambda)}>0$, we have $s^{(k')}_\lambda=1-\beta+\beta2^{-(k'-\lambda)}\ge1-\beta$,
and likewise $s^{(k)}_\lambda\ge1-\beta$. Hence
\begin{equation*}
s-\beta S\ \ge\ 1-\beta-3.508\,\beta\ =\ 1-4.508\,\beta\ \ge\ 1-4.51\,\beta=\varphi_{\mathrm{arrest}}.
\end{equation*}
For the class-letter brackets, $\lambda\le\ell\le k_{a_1}\le k_a$, the last inequality because
$\mathfrak{a}$ is a live arrest holding $x$, so that $k_a$ is one of the levels of the live chain.
The proof of (iii) below uses hypothesis (b) only, and not (i). By (iii) at $\nu=\lambda$ and the
bound just proved for $F^{(k')}_\lambda$,
\begin{equation*}
\theta^{(n)}_T(\lambda)=F^{(k')}_\lambda+\rho_n(x;k';\lambda)\ \ge\ F^{(k')}_\lambda
\ \ge\ \varphi_{\mathrm{arrest}}.
\end{equation*}

\emph{(ii).} Let the upper index exceed $\lambda$, say $k>\lambda$. Then
$2^{-(k-\lambda)}\le\frac12$, so $s^{(k)}_\lambda\le1-\frac{\beta}{2}$. The same holds for $k'$.

Suppose $\lambda-1,\lambda,\lambda+1$ lie in one subtracted range. The terms $i=\lambda$ and
$i=\lambda\pm1$ of $S$ contribute $1+\frac12+\frac12=2$, and all other terms are nonnegative.
Since $\beta>0$, the bracket is at most
$1-\frac{\beta}{2}-2\beta=1-\frac52\beta=\frac12$.

For $\nu\in[h_a+2,k_a]$ we have $h_a+1\le\nu-1$ and $\nu+1\le k_a+1$, so
$\nu,\nu\pm1\in I(\mathfrak{a})$. The arrest $\mathfrak{a}$ is live and holds $x$, so
$\sigma_{\mathfrak{a}}(\nu)$ is subtracted in $F^{(k')}_\nu$. With $k'>\nu$ the preceding
paragraph gives $F^{(k')}_\nu\le\frac12$.
More precisely, the terms of $\sigma_{\mathfrak{a}}(\nu)$ with $i\in[h_a+1,k_a+1]$ alone give
\begin{align*}
\sum_{i=h_a+1}^{k_a+1}2^{-|\nu-i|}
&=1+\sum_{i=h_a+1}^{\nu-1}2^{-(\nu-i)}+\sum_{i=\nu+1}^{k_a+1}2^{-(i-\nu)}\\
&=3-2^{-(\nu-h_a-1)}-2^{-(k_a+1-\nu)},
\end{align*}
and all other subtracted terms are nonnegative, so that
\begin{equation*}
F^{(k')}_\nu\ \le\ 1-\tfrac{\beta}{2}-\beta\bigl(3-2^{-(\nu-h_a-1)}-2^{-(k_a+1-\nu)}\bigr).
\end{equation*}
Both powers of $2$ here are at most $\frac12$, since $\nu-h_a-1\ge1$ and $k_a+1-\nu\ge1$; the
right-hand side is therefore at most $1-\frac{\beta}{2}-2\beta=\frac12$.

At a lifted point of kind (A), the three levels lie in $[k_0,j^+]$. We have
$k_0-h_a=N_a-N_0^a$. By (c),
\begin{equation*}
N_a-N_0^a\ \ge\ (N_0^a)^2+4N_0^a+3-N_0^a\ =\ (N_0^a)^2+3N_0^a+3\ \ge\ 3\ \ge\ 2,
\end{equation*}
so $k_0\ge h_a+2$. With $j^+\le k_a$ this gives $[k_0,j^+]\subset[h_a+2,k_a]$, and the previous
paragraph gives $F^{(k')}_\nu\le\frac12$.

Suppose moreover that every live arrest $\mathfrak{d}$ holding $x$ satisfies $k_d\le h$. Since
$h=k-N<k$, every such arrest has $k_d<k$, so the sum
$\sum_{\mathfrak{d}\text{ live}\ni x,\,k_d<k}\sigma_{\mathfrak{d}}(\nu)$ subtracted in
$\theta^{(l)}_{\mathfrak{s}}(\nu)$ coincides with the sum subtracted in $F^{(k')}_\nu$. Further
$\nu\le j^+\le k_a\le h$, and $l\ge1$, the stages being numbered $1,\dots,N$ in
Notation~\ref{4.26:not-stage-frozen-data}; hence the term $i=h+1$ of the attempt's own sum in
$\varphi^{(l),(k)}_\nu$ equals $2^{-(h+1-\nu)}$, and its other terms are nonnegative. Since
$k\ge h+1$, we have $2^{-(k-\nu)}\le2^{-(h+1-\nu)}$, and therefore
\begin{align*}
\varphi^{(l),(k)}_\nu&\le s^{(k)}_\nu-\beta2^{-(h+1-\nu)}
=1-\beta+\beta\bigl(2^{-(k-\nu)}-2^{-(h+1-\nu)}\bigr)\\
&\le1-\beta\ \le\ s^{(k')}_\nu.
\end{align*}
Subtracting from both sides $\beta$ times the common sum of the $\sigma_{\mathfrak{d}}(\nu)$ gives
$\theta^{(l)}_{\mathfrak{s}}(\nu)\le F^{(k')}_\nu$, and with $F^{(k')}_\nu\le\frac12$ from the
previous paragraph, $\theta^{(l)}_{\mathfrak{s}}(\nu)\le F^{(k')}_\nu\le\frac12$.

At a lifted point of kind (B) or (B$'$), suppose that
$\mathfrak{c}(x)-1,\mathfrak{c}(x),\mathfrak{c}(x)+1$ lie in one of the subtracted ranges of
$\theta^{(l)}_{\mathfrak{s}}$ and that $k>\mathfrak{c}(x)$. The first assertion of (ii), applied to
the bracket $\theta^{(l)}_{\mathfrak{s}}$ at $\lambda=\mathfrak{c}(x)$, gives
$\theta^{(l)}_{\mathfrak{s}}(\mathfrak{c}(x))\le\frac12$.

\emph{(iii).} When $\mathfrak{a}$ is the innermost live arrest holding $x$, by its definition and
the definition of the source factor in Notation~\ref{4.26:not-arrest-conventions},
$\rho_n(x;k')$ is $\theta^{(n)}_T(\ell)-F^{(k')}_\ell$, that is, $\rho_n(x;k';\ell)$ in the notation of
\eqref{4.26:eq-lower-bracket-a}.

By (b), the proof of the room bound reads the level only through the index ranges of the
subtracted sums in $\theta^{(n)}_T$ and $F^{(k')}$, and these ranges do not depend on the level at
which the sums are evaluated. The same proof therefore applies with $\ell$ replaced by any level
$\nu\le k_a$, and gives $\rho_n(x;k';\nu)\ge\beta2^{-(k_a+1-\nu)}$. The right-hand side is positive
because $\beta>0$.
\end{proof}

\begin{lemma}[One-step tube in the frozen unit for an enclosing run]\label{4.26:lem-enclosing-tube}
Use the stage data of Notation~\ref{4.26:not-stage-frozen-data}: the site $k$, the level $h=k-N$,
the region $Z$ and the stages $l$ of the run at this site. Let $\mathfrak{a}$ be a live arrest with
$W_{\mathfrak{a}}\subset Z$ and $k_a\le h$, so that the run at site $k$ encloses $W_{\mathfrak{a}}$;
the attempt spaces $\mathfrak{L}_l$ below are those of this enclosing run, read at points of
$W_{\mathfrak{a}}$. Let $l$ and $l+1$ be consecutive stages of the run.
Let $(p,X',\mathfrak{s})$ be a retained triple of $\mathfrak{a}$ at a point $y'$. Let
$m=\ell_{\mathfrak{a}}(y')$ be its frozen level, $\omega=w_{\mathfrak{a}}(y')$ its frozen weight and
$\lambda=\lambda_{\mathfrak{s}}(y')\le m$ the level of the block of $\mathfrak{s}$ at $y'$. Fix an
entry $e\in\{2,4\}$ of the triple, and let $\mathrm{U2}=\mathrm{U2}_e$ be its unit,
as in Definition~\ref{4.26:def-three-units} and \eqref{4.26:eq-three-units-a}. Let $Q_e$ denote a
composite entry of the triple at $y'$.

Write $\theta^{(l)}(m):=\theta^{(l),(k)}(m)$ for the attempt-space bracket at level $m$ of
Lemma~\ref{4.26:lem-lower-bracket}. Recall from \eqref{4.26:eq-lower-bracket-a} that
\begin{equation}\label{4.26:eq-enclosing-tube-1}
\theta^{(l)}(m)=s^{(k)}_m-\beta\sum_{i=h+1}^{h+l}2^{-|m-i|}-\beta\sum_{\mathfrak{d}}\sigma_{\mathfrak{d}}(m),
\end{equation}
where the last sum runs over the live arrests $\mathfrak{d}$ with $y'\in W_{\mathfrak{d}}$ and
arrest level $k_d<k$. Put $\rho^{(l+1)}_m:=\beta\,2^{-|m-h-l-1|}$.

Suppose that at the triple the composite clause of the attempt space of each stage is written in
the unit $\mathrm{U2}$:
\begin{equation}\label{4.26:eq-enclosing-tube-2}
\mathfrak{L}_l=\bigl\{Q_e<\theta^{(l)}(m)\,\mathrm{U2}\bigr\},\qquad
\mathfrak{L}_{l+1}=\bigl\{Q_e<\theta^{(l+1)}(m)\,\mathrm{U2}\bigr\}.
\end{equation}
Let $r_e=r_e(m,p\wedge\lambda)$ be the shell radius of entry $e$ at level $m$ and block level
$p\wedge\lambda$, so that $\mathrm{U2}=\omega\,r_e$ by Definition~\ref{4.26:def-three-units}. Let
$\pi,\pi_\times$ be the nonnegative constants of the composite-move bound, and let
$\iota_{l+1}(y')$ be the stage increment. The amount by which $Q_e$ changes under the
stage-$(l+1)$ step, in absolute value, is
called the stage-$(l+1)$ move of the composite entry, and is written $\operatorname{move}$ in
displays.
Assume the following.
\begin{enumerate}
\item[(1)] (Stage independence of the source factor.) The quantities $s^{(k)}_m$ and
$\sigma_{\mathfrak{d}}(m)$ in \eqref{4.26:eq-enclosing-tube-1} carry no stage index.
\item[(2)] (The composite-move bound in the frame of the triple.) Work in the frame
$(X',\mathfrak{B},\mathfrak{s}^{(p)})$ of the general clause of the composite-move bound (the
bound on the move of a composite entry under one stage); here $\mathfrak{s}^{(p)}$ denotes the
structure $\mathfrak{s}$ of the triple at the
triple's level $p$, the parent level is $\mathsf{m}^{\circ}=\lambda$ and the frame level is
$\mathsf{m}=m$, so that $\lambda\le m$. For every real $\theta$: if a member has
$Q_e<\theta\,\mathrm{U2}$, then the stage-$(l+1)$ move of its composite entry is at most
$\pi_\times\theta\,\mathrm{U2}$ plus the additive term $\pi\,r_e(m,p\wedge\lambda)$.
\item[(3)] (The composite-move line at stage $l+1$ and at the point $y'$.) Put
$r_e=r_e(m,p\wedge\lambda)$. If
$Q_e\le 1.4\,\omega r_e$, then the stage-$(l+1)$ move of the composite entry satisfies
\begin{equation*}
\operatorname{move}\le(1.4\,\pi_\times\omega+\pi)\,r_e\le 2.4\,\pi\,\omega r_e
\le\iota_{l+1}(y')\,\omega r_e .
\end{equation*}
\item[(4)] (The live-piece inequality, read on $W_{\mathfrak{a}}$ as on an excluded region.)
The function $\sigma_{\mathfrak{a}}$ does not depend on the stage, and
\begin{equation*}
\iota_{l+1}(y')\le\frac12\,\rho^{(l+1)}_{m}.
\end{equation*}
\item[(5)] (Ceiling of the bracket.) $\theta^{(l+1)}(m)\le 1.4$.
\end{enumerate}
Then:
\begin{enumerate}
\item[(a)] the brackets of consecutive stages differ exactly by the one-step room:
\begin{equation}\label{4.26:eq-enclosing-tube-3}
\theta^{(l)}(m)-\theta^{(l+1)}(m)=\rho^{(l+1)}_m=\beta\,2^{-(h+l+1-m)};
\end{equation}
\item[(b)] the stage-$(l+1)$ move of a member's composite entry satisfies
\begin{equation}\label{4.26:eq-enclosing-tube-4}
\operatorname{move}\le\Bigl(\theta^{(l+1)}(m)\,\pi_\times+\frac{\pi}{\omega}\Bigr)\mathrm{U2}
\le 2.4\,\pi\,\mathrm{U2}\le\iota_{l+1}(y')\,\mathrm{U2};
\end{equation}
\item[(c)] at the triple, the image of $\mathfrak{L}_{l+1}$ under the stage-$(l+1)$ step lies in
$\mathfrak{L}_l$ with room at least $\tfrac12\rho^{(l+1)}_m\mathrm{U2}$. That is, the composite
entry of every point of the image is below
$\theta^{(l)}(m)\mathrm{U2}-\tfrac12\rho^{(l+1)}_m\mathrm{U2}$.
\end{enumerate}
\end{lemma}

\begin{proof}
(a) By hypotheses (1) and (4), neither $s^{(k)}_m$ nor any $\sigma_{\mathfrak{d}}(m)$ depends on the
stage. Hence in the difference of \eqref{4.26:eq-enclosing-tube-1} at stages $l$ and $l+1$ these
terms cancel. The two sums over $i$ differ only by the term $i=h+l+1$, so
\begin{equation*}
\theta^{(l)}(m)-\theta^{(l+1)}(m)=\beta\,2^{-|m-h-l-1|}=\rho^{(l+1)}_m .
\end{equation*}
Since $m\le k_a\le h<h+l+1$, we have $|m-h-l-1|=h+l+1-m$. This gives
\eqref{4.26:eq-enclosing-tube-3}.

(b) We apply hypothesis (2) in the frame $(X',\mathfrak{B},\mathfrak{s}^{(p)})$ of the general
clause of the composite-move bound. There the parent level is $\mathsf{m}^{\circ}=\lambda\le\mathsf{m}=m$;
this is the case of that clause in which the parent level does not exceed the frame level. A member
of $\mathfrak{L}_{l+1}$ has $Q_e<\theta^{(l+1)}(m)\,\mathrm{U2}$ by
\eqref{4.26:eq-enclosing-tube-2}. So hypothesis (2), with $\theta=\theta^{(l+1)}(m)$, bounds the
move by $\pi_\times\theta^{(l+1)}(m)\,\mathrm{U2}+\pi\,r_e(m,p\wedge\lambda)$.

By Definition~\ref{4.26:def-three-units}, $\mathrm{U2}$ is the shell radius $r_e(m,p\wedge\lambda)$
times the frozen weight $\omega$. Hence
\begin{equation*}
\pi\,r_e(m,p\wedge\lambda)=\pi\,\mathrm{U2}/\omega .
\end{equation*}
The additive term is thus itself a multiple of $\mathrm{U2}$, with no change of unit, and the
first inequality of \eqref{4.26:eq-enclosing-tube-4} follows.

The second and third inequalities come from hypothesis (3), read with $r_e=r_e(m,p\wedge\lambda)$
and $\omega r_e=\mathrm{U2}$. A member of $\mathfrak{L}_{l+1}$ has
$Q_e<\theta^{(l+1)}(m)\,\mathrm{U2}$. By hypothesis (5), $\theta^{(l+1)}(m)\le 1.4$, so the member
has $Q_e\le 1.4\,\mathrm{U2}=1.4\,\omega r_e$, and the premise of hypothesis (3) holds. Moreover,
since $\pi_\times\ge0$ and $\theta^{(l+1)}(m)\le 1.4$,
\begin{equation*}
\Bigl(\theta^{(l+1)}(m)\,\pi_\times+\frac{\pi}{\omega}\Bigr)\mathrm{U2}
\le\Bigl(1.4\,\pi_\times+\frac{\pi}{\omega}\Bigr)\mathrm{U2}.
\end{equation*}
Substituting $\omega r_e=\mathrm{U2}$, that is
$r_e=\mathrm{U2}/\omega$, the chain of hypothesis (3) reads
\begin{equation*}
\Bigl(1.4\,\pi_\times+\frac{\pi}{\omega}\Bigr)\mathrm{U2}\le 2.4\,\pi\,\mathrm{U2}
\le\iota_{l+1}(y')\,\mathrm{U2}.
\end{equation*}

(c) By (a), $\theta^{(l+1)}(m)=\theta^{(l)}(m)-\rho^{(l+1)}_m$. By hypothesis (4),
$\iota_{l+1}(y')\le\tfrac12\rho^{(l+1)}_m$. Therefore
\begin{align*}
\theta^{(l+1)}(m)\,\mathrm{U2}+\iota_{l+1}(y')\,\mathrm{U2}
&=\theta^{(l)}(m)\,\mathrm{U2}-\rho^{(l+1)}_m\mathrm{U2}+\iota_{l+1}(y')\,\mathrm{U2}\\
&\le\theta^{(l)}(m)\,\mathrm{U2}-\frac12\rho^{(l+1)}_m\mathrm{U2}.
\end{align*}
Take a point of $\mathfrak{L}_{l+1}$. Its composite entry is below
$\theta^{(l+1)}(m)\,\mathrm{U2}$. By (b), the stage-$(l+1)$ step moves this entry by at most
$\iota_{l+1}(y')\,\mathrm{U2}$. So the entry of its image is below the right-hand side above. By
\eqref{4.26:eq-enclosing-tube-2}, the image therefore lies in $\mathfrak{L}_l$ with room at least
$\tfrac12\rho^{(l+1)}_m\mathrm{U2}$.

Finally, consider the unit. At the triples of the frozen structure itself one has $\lambda=m$.
There $\mathrm{U2}=\mathrm{U3}$, so (a)--(c) are statements in the unit $\mathrm{U3}$. If the
composite clause were written in $\mathrm{U3}$ at a triple with $\lambda<m$, the comparison of the
increment $\iota_{l+1}(y')$ with the room $\tfrac12\rho^{(l+1)}_m$ would have to absorb the
conversion factor from $\mathrm{U2}$ to $\mathrm{U3}$ in the rule for composite clauses at
retained structures. With the clause written in $\mathrm{U2}$, as in
\eqref{4.26:eq-enclosing-tube-2}, the inequality $\iota_{l+1}(y')\le\tfrac12\rho^{(l+1)}_m$ of
hypothesis (4) is used as it stands, with conversion factor $1$.
\end{proof}

\begin{remark}[The uniform alternative at birth radii]\label{4.26:rem-birth-radii-alternative}
Work at a retained structure $\mathfrak{s}$ of a live arrest $\mathfrak{a}$ at a point $y'$, in the
setting of Lemma~\ref{4.26:lem-unit-ordering}. Thus $m=\ell_{\mathfrak{a}}(y')$ is the frozen level,
$\lambda=\lambda_{\mathfrak{s}}(y')\le m$, $\mathfrak{c}=\mathfrak{c}(y')$, and
$\mathrm{U1}\le\mathrm{U2}\le\mathrm{U3}$ are the three units ordered in
\eqref{4.26:eq-unit-ordering-a}. The assignment rule for composite clauses
(Definition~\ref{4.26:def-ownership-rule}) reads a composite clause at $(\lambda,\mathrm{U3})$ if the
space or letter belongs to the run that built $\mathfrak{s}$, and at $(m,\mathrm{U2})$, with the factor
of the run to which it belongs, otherwise. A choice,
for each space, of the level at which its factor is read and of the unit in which its composite
clauses are read is called a typing. The following typing, the uniform typing, is also admissible;
it is not the one used in this
chapter. In the uniform typing every space reads its composite clauses in the unit $\mathrm{U3}$, with
its factor evaluated at the level $\lambda$:
\begin{equation}\label{4.26:eq-birth-radii-alternative-1}
\mathfrak{N}^{(n)}:\ 1+2\beta-{\textstyle\sum}\iota^1,\qquad
\tilde{U}^{\mathrm{arrest}}_{k'}:\ F^{(k')}_{\lambda},\qquad
\text{class letters}:\ \theta^{(n)}_T(\lambda).
\end{equation}
Below, $\theta^{(n)}(\lambda)$ stands for any of the brackets $\theta^{(n)}_T(\lambda)$ at $\lambda$.
\begin{enumerate}
\item[(a)] In the uniform typing the required inclusions hold:
\begin{itemize}
\item the source space $\tilde{U}^{\mathrm{arrest}}_{k'}$ lies in every class letter. The room is
$\rho_n(\lambda;k')$, where $\rho_n(\lambda;k')$ denotes the difference
$\theta^{(n)}_T(\lambda)-F^{(k')}_{\lambda}$, and it satisfies
$\rho_n(\lambda;k')\ge\beta2^{-(k_a+1-\lambda)}$;
\item the source space lies in $\mathfrak{N}^{(n)}$;
\item the real point $\Phi^0$ lies in every space with margin $\tfrac12$;
\item the box-shell clause of $\mathfrak{N}^{(n)}$ at the box triple is \cite[(1.16)]{B-CMP122a}
literally, without the unit conversion \eqref{4.26:eq-unit-ordering-b};
\item the starting inclusion (clause (c5) of \eqref{4.26:eq-ownership-rule-a}) holds with the rooms
$0.3$ and $0.75$.
\end{itemize}
\item[(b)] In the uniform typing the following would have to be restated:
\begin{itemize}
\item the frozen weighted unit in which the $(\Sigma^{\ast})^{\mathfrak{a}}$-triples are read;
\item the reading of composite clauses at the frozen level $m=m(y')$;
\item the one-step tube statement read in the frozen weighted unit, whose proof uses the
$l$-independence of $\sigma_{\mathfrak{a}}$;
\item the box-shell clause, the starting inclusion and the kernel clause of $\mathfrak{N}^{(n)}$ at
$\mathfrak{T}$, that is, clauses (c4)--(c6) of \eqref{4.26:eq-ownership-rule-a};
\item the comparison, among those of clause (c7) of \eqref{4.26:eq-ownership-rule-a}, of the unit
$\mathrm{U2}$ of the source space with the unit $\mathrm{U3}$ of a letter;
\item the shell unit of the propagation in clause (c3) of \eqref{4.26:eq-ownership-rule-a};
\item the factor clause of the source space;
\item the kernel at $\mathfrak{T}$, which would then be obtained from the direct entries of
$\mathbb{U}$, as at $Z^{\mathrm{on}}$.
\end{itemize}
\item[(c)] Consider the mixed typing, in which the source space reads its composite clauses at
$F^{(k')}_m\cdot\mathrm{U3}$ and the class letters at $\theta^{(n)}(\lambda)\cdot\mathrm{U3}$. Under it
the inclusion of the source space in the class letters does not follow from the comparison of their
clauses. The
factor and the unit must therefore be read at the same place: either at $(m,\mathrm{U2})$ or at
$(\lambda,\mathrm{U3})$.
\end{enumerate}

The assignment rule alters only the block side of the frozen weighted unit; every other statement
entering clauses (c3)--(c6) of \eqref{4.26:eq-ownership-rule-a} is used unchanged. The uniform typing
is not forced: the statement that each run keeps its own retained stages inside its own attempt spaces
concerns only the attempt letters, which the assignment rule already reads at $(\lambda,\mathrm{U3})$
for the run that built $\mathfrak{s}$, while the source space is defined only through the frozen data
$(\ell_{\mathfrak{a}},w_{\mathfrak{a}},\mathcal{Q})(y')$ of the point $y'$ (its frozen level, weight and
cube clauses), and its composite clauses are read at the factors $s^{(k)}_m$. Nor is the
assignment rule forced: if the statement that $(\Sigma^{\ast})$ is the sub-domain twin of
$(\Sigma^{\flat})$ is read as prescribing birth radii for every space of the family $(\Sigma^{\ast})$,
then the uniform typing is the one to be used, at the price of the restatements listed in (b).
\end{remark}

\begin{proof}
\emph{Source inside the class letters.} We apply Lemma~\ref{4.26:lem-lower-bracket}(iii) at
$\nu=\lambda$; its hypotheses hold since $\lambda\le m\le k_a$. Part (iii) states that the room of the
comparison between the source factor and
a term's threshold, evaluated at any level $\nu\le k_a$, equals $\theta^{(n)}_T(\nu)-F^{(k')}_{\nu}$.
It also states that this room is at least $\beta2^{-(k_a+1-\nu)}>0$.

Take $\nu=\lambda$, so that the comparison is carried out at the level $\lambda$ in place of the frozen
level. A member of the source space has composite entries $Q$ at the triple satisfying
\begin{equation*}
Q<F^{(k')}_{\lambda}\,\mathrm{U3}=\theta^{(n)}_T(\lambda)\,\mathrm{U3}-\rho_n(\lambda;k')\,\mathrm{U3}.
\end{equation*}
It therefore satisfies the clause of the class letter with room $\rho_n(\lambda;k')\,\mathrm{U3}$, and
this room is at least $\beta2^{-(k_a+1-\lambda)}\,\mathrm{U3}$.

\emph{Source inside $\mathfrak{N}^{(n)}$.} The
source factor satisfies
\begin{equation*}
F^{(k')}_{\lambda}\le1<1.383\le1+2\beta-{\textstyle\sum}\iota^1,
\end{equation*}
the last inequality being the lower bound on the factor of $\mathfrak{N}^{(n)}$ used in clause (c1)
of \eqref{4.26:eq-ownership-rule-a}. Since both clauses are read in $\mathrm{U3}$, the inclusion holds
clause by clause.

\emph{The real point.} Under the assignment rule, the composite entries of $\Phi^0$ at the triple are at
most $\tfrac12\varphi_{\mathrm{arrest}}\,\mathrm{U2}$ (clause (c2) of
\eqref{4.26:eq-ownership-rule-a}). By \eqref{4.26:eq-unit-ordering-a} we have
$\mathrm{U2}\le\mathrm{U3}$. At $\lambda<m$ the sharper bound $\mathrm{U2}\le\mathrm{U3}/31.5$ holds,
and at $\lambda=m$ we have $\mathrm{U2}=\mathrm{U3}$. Lemma~\ref{4.26:lem-lower-bracket}(i) gives
$\theta^{(n)}(\lambda)\ge\varphi_{\mathrm{arrest}}$ and $F^{(k')}_{\lambda}\ge\varphi_{\mathrm{arrest}}$ for
every bracket at $\lambda$. Hence
\begin{equation*}
\tfrac12\varphi_{\mathrm{arrest}}\,\mathrm{U2}\le\tfrac12\varphi_{\mathrm{arrest}}\,\mathrm{U3}
\le\tfrac12\theta^{(n)}(\lambda)\,\mathrm{U3},
\end{equation*}
and likewise with $F^{(k')}_{\lambda}$ in place of $\theta^{(n)}(\lambda)$, and, for
$\mathfrak{N}^{(n)}$, with $1+2\beta-\sum\iota^1\ge1.383\ge\varphi_{\mathrm{arrest}}$ in place of it. So
$\Phi^0$ lies in each
space with margin $\tfrac12$. Where $\lambda<m$, the margin holds with a further factor $1/31.5$.

\emph{The box-shell clause.} At the box triple we have $\lambda=\mathfrak{c}$ (clause (c4) of
\eqref{4.26:eq-ownership-rule-a}). Under the assignment rule, clause (c4) reads the box-shell
clause in the frozen weighted unit $\mathrm{U2}$. It then reaches \cite[(1.16)]{B-CMP122a} in the unit
of the $\mathfrak{c}$-shell only through the conversion \eqref{4.26:eq-unit-ordering-b}. In the
uniform typing the clause is read in $\mathrm{U3}$ with its factor at $\mathfrak{c}$. It is then
\cite[(1.16)]{B-CMP122a} literally, and the conversion is not used.

\emph{The starting inclusion.} At the box triple the source factor is $F^{(k')}_{\mathfrak{c}}$. The
box triple has $\lambda=\mathfrak{c}$, and at a lifted point Lemma~\ref{4.26:lem-lower-bracket}(ii),
in the case $\nu=\mathfrak{c}$, gives $F^{(k')}_{\mathfrak{c}}\le\tfrac12$; it is read there only as an
upper bound. The starting inclusion at birth radii with the rooms $0.3$ and $0.75$, which clause (c9)
of \eqref{4.26:eq-ownership-rule-a} keeps at the run's own raised structures, uses of the source factor
at $\mathfrak{c}$ only the bound $F^{(k')}_{\mathfrak{c}}\le\tfrac12$. Since this bound holds here, the
starting inclusion holds in the uniform typing with the rooms $0.3$ and $0.75$.

\emph{Restatements.} Each statement in the list of (b) reads a composite clause at the frozen level $m$
or in the frozen weighted unit. In the uniform typing the level is replaced by $\lambda$ and the unit by
$\mathrm{U3}$, so each must be restated. At $\mathfrak{T}$ the kernel clause is no longer the composite
entry read at $\mathfrak{B}$'s own levels. It is then obtained from the direct entries of $\mathbb{U}$,
as at $Z^{\mathrm{on}}$.

\emph{The mixed typing.} Consider a piece of top stage, that is, one with $j^+=k_a$ and $n^+_a=N_a$,
where $N_a$ is the index of the last stage of the run of $\mathfrak{a}$.
Take a point of it at which $m=k_a$ and $\mathfrak{c}=k_0$, and take $\mathfrak{s}=\mathbb{B}_{k_a}$, so
that $\lambda=k_0$. Compare the class letter stored from the attempt space of stage $N_a$ with the
source of class $k'=k_a+1$. At this point the brackets are
\begin{align*}
\theta^{(N_a)}(k_0)&=1-\beta\bigl[4-2^{1-N_0^a}-2^{-(k_0-h_a-1)}\bigr],\\
F^{(k_a+1)}_{k_a}&=1-\beta\bigl[3-2^{1-N_a}\bigr].
\end{align*}
Subtracting, we obtain
\begin{equation}\label{4.26:eq-birth-radii-alternative-2}
\theta^{(N_a)}(k_0)-F^{(k_a+1)}_{k_a}
=-\beta\bigl[1-2^{1-N_0^a}-2^{-(k_0-h_a-1)}+2^{1-N_a}\bigr].
\end{equation}
Two conditions bound the terms in the bracket. The first is $N_0^a\ge3$, which is implied by
$R_{\min}\ge5L$; it gives $2^{1-N_0^a}\le\tfrac14$. The second is $k_0-h_a\ge4$, which is implied by
$(N_0+1)(N_0+3)\le N$ at density $k_a$ (here $N$ is Ba{\l}aban's number $R_{k_a}$, not the rank of the
gauge group); it gives $2^{-(k_0-h_a-1)}\le\tfrac18$. Since also $2^{1-N_a}>0$, the bracket in
\eqref{4.26:eq-birth-radii-alternative-2} is at least $1-\tfrac14-\tfrac18=0.625$. Hence the difference
in \eqref{4.26:eq-birth-radii-alternative-2} is at most $-0.625\beta<-\beta/2$.

In the common unit $\mathrm{U3}$, the threshold of that letter is therefore smaller than the
threshold of the source space by more than $\tfrac{\beta}{2}\,\mathrm{U3}$. The inclusion of the
source space in that letter cannot be obtained by comparing the two clauses. When the unit is
$\mathrm{U3}$, the factor must be read at $\lambda$ for the source space as well; when the factor is
read at $m$, the unit must be $\mathrm{U2}$.
\end{proof}

\begin{remark}[What the enlarged family of triples adds]\label{4.26:rem-enlarged-triples}
Let $\mathfrak{a}$ be an arrest that is live at the present site $k$, at the current stage $n$,
with region $W_{\mathfrak{a}}$, density $k_a<k$ and final stage $n^+_a$, where
$W_{\mathfrak{a}}\subset Z^c$ or $W_{\mathfrak{a}}\subset Z$. Let $j^+$ and $k_0$ be the
levels among the data of the piece for which its lifted set is
\begin{equation*}
\mathsf{L}_{\mathfrak{a}}=(\Omega^c_{j^+}\setminus Z''_{k_0+1})\cap W_{\mathfrak{a}}.
\end{equation*}
Let $\mathfrak{B}$ be the current
determining-set structure of the site, and $\mathfrak{B}^{\mathfrak{c}}$ its graded refinement.
For a localisation domain $X'$ write $\mathfrak{F}_{\mathrm{cur}}(X')$ for the restrictions to $X'$
of $\mathfrak{B}$ and of its predecessors at the present site. Write $\mathfrak{F}_{\mathrm{hyb}}(X')$ for the set
of structures on $X'$ that coincide on $X'\cap W_{\mathfrak{a}}$ with
$\mathbb{B}^{(n')}_{k_a}\restriction(X'\cap W_{\mathfrak{a}})$ for some $0\le n'\le n^+_a$, and
on $X'\setminus W_{\mathfrak{a}}$ with the restriction to $X'\setminus W_{\mathfrak{a}}$ of a member
of $\mathfrak{F}_{\mathrm{cur}}(X')$.

Let $X'$ be a localisation domain meeting $W_{\mathfrak{a}}$.
The structures $\mathfrak{s}$ at which the family $(\Sigma^{\ast})$ of
Lemma~\ref{4.26:lem-retained-levels} imposes the derived entries $2$
and $4$ of a triple $(p,X',\mathfrak{s})$ include the following. On $X'\cap W_{\mathfrak{a}}$ they
include the frozen structure
$\mathbb{B}^{(n^+_a)}_{k_a}\restriction(X'\cap W_{\mathfrak{a}})$ and each predecessor
$\mathbb{B}^{(n')}_{k_a}\restriction(X'\cap W_{\mathfrak{a}})$, $0\le n'\le n^+_a$, in the piece's
own run. On
$X'\setminus W_{\mathfrak{a}}$ these are continued by the members of $\mathfrak{F}_{\mathrm{cur}}(X')$,
that is, by $\mathfrak{B}$ and its predecessors at the present site. By the
proposition describing the triples carried by the attempt spaces, every $(\Sigma^{\ast})$-triple
$(p,X',\mathfrak{s}_{\mathfrak{L}})$ that an attempt space $\mathfrak{L}$ written at a density
$k'\in[k_a,k]$ carries, with $(\Sigma^{\ast})$ read at that density, satisfies
$\mathfrak{s}_{\mathfrak{L}}=\mathfrak{s}$ after
truncation at $p$ for some
$\mathfrak{s}\in\mathfrak{F}_{\mathrm{cur}}(X')\cup\mathfrak{F}_{\mathrm{hyb}}(X')$. This is consistent with the
family just described: the spaces of class $k_a$ (in the sense of
Definition~\ref{4.26:def-ownership-rule})
written at densities $k'\in[k_a,k]$ already carry exactly this hybrid family.

For the spaces of the present site, the added triples fall into three classes.
\begin{itemize}
\item[$(\alpha)$] At junction boxes (cover cubes $\square\in\mathcal{C}^1$ whose enlargement
$\square_0$ meets $W_{\mathfrak{a}}$ without being contained in it), in the case $j^+=k_a$, the
structure
$\mathfrak{s}_1:=\mathfrak{B}^{\mathfrak{c}}\restriction\square_0$ is added. It has two cells.
On the lifted old shell of level $k_a-1$ inside $\square_0$ it has level $\mathfrak{c}=k_a-1$, with
$\ell_{\mathfrak{a}}=k_a$, $\omega=L_0^2$, lift $g(y)=1$ at every point $y$ of that shell and
$w_{\mathfrak{s}}=1$. On
$\square_0\setminus W_{\mathfrak{a}}$ it has level $k_a$ and weight $1$.
\item[$(\beta)$] In the interior of $W_{\mathfrak{a}}$, hybrid triples at lifted points are added.
These are, first, the triples $(p,\square_0,\mathfrak{B}^{\mathfrak{c}}\restriction\square_0)$, with
$\lambda_{\mathfrak{s}}=\mathfrak{c}$. On $\square_0$ such a triple meets at most two adjacent values
of $\mathfrak{c}$ and at most two adjacent frozen levels. Second, they are the triples with a larger
$X'$ carrying a retained $\mathbb{B}^{(n')}_{k_a}\restriction(X'\cap W_{\mathfrak{a}})$,
$n'<n^+_a$. Triples at un-lifted
inner collars are also added; there $\lambda_{\mathfrak{s}}=\ell_{\mathfrak{a}}-1$ and
$\omega=w_{\mathfrak{s}}=1$.
\item[$(\gamma)$] Nothing is added at $\partial W_{\mathfrak{a}}$ when $j^+<k_a$.
\end{itemize}

As clauses, nothing is new relative to $(\Sigma^{\ast})$ read at density $k_a$, together with the
reading convention that leaves $(\Sigma^{\flat})$ and $(\Sigma^{\ast})$ unchanged in spaces written
after the arrest. What is new is that these triples are named by the inherited domain
$\mathfrak{D}^{\flat}_n$ of the present site. Three properties must therefore be checked at every
named triple:
\begin{itemize}
\item[(1)] the real point lies inside the site's spaces at the triple, with a margin;
\item[(2)] the arrested source $\tilde{U}^{\mathrm{arrest}}_{k'}$ imposes the
clauses in the same wording as
$\mathfrak{N}^{(n)}$ and is contained in each attempt space $\mathfrak{L}_l$ and each space
$\mathbf{B}^{(k')}$ written after the arrest;
\item[(3)] the clauses on the block side at the triple are those that the rule for composite
clauses at retained structures, Definition~\ref{4.26:def-ownership-rule}, assigns to it.
\end{itemize}
Property (3) holds by the convention fixed in Definition~\ref{4.26:def-ownership-rule}. Properties
(1) and (2) are
supplied by the lemma placing the real
point at the triples named by $\mathfrak{D}^{\flat}_n$.
\end{remark}

\begin{proof}
The first assertion is clause $(\Sigma^{\ast})^{\mathfrak{a}}$ of
Lemma~\ref{4.26:lem-retained-levels}, applied to a localisation domain $X'$ meeting
$W_{\mathfrak{a}}$.

For the consistency, take a space of class $k_a$ written at density
$k'$ with $k_a\le k'\le k$. On $X'\cap W_{\mathfrak{a}}$ it carries, by $(\Sigma^{\ast})$ read at
density $k_a$ (Lemma~\ref{4.26:lem-retained-levels}) and the reading convention that spaces
written after the arrest read $(\Sigma^{\ast})$ unchanged, the frozen
$\mathbb{B}^{(n^+_a)}_{k_a}\restriction(X'\cap W_{\mathfrak{a}})$ and each retained
$\mathbb{B}^{(n')}_{k_a}\restriction(X'\cap W_{\mathfrak{a}})$ with $n'<n^+_a$. These are the
structures that the
first assertion lists on $X'\cap W_{\mathfrak{a}}$, so such a space carries them. Conversely, by the
proposition describing the triples carried by the attempt spaces, every structure it carries
coincides after truncation at $p$ with a member of
$\mathfrak{F}_{\mathrm{cur}}(X')\cup\mathfrak{F}_{\mathrm{hyb}}(X')$. Hence it carries no structure
outside $\mathfrak{F}_{\mathrm{cur}}(X')\cup\mathfrak{F}_{\mathrm{hyb}}(X')$, and, by the preceding
sentences and the first assertion, it carries every member of $\mathfrak{F}_{\mathrm{hyb}}(X')$.

\emph{Class $(\alpha)$.} By Lemma~\ref{4.26:lem-retained-levels}, the triple
$(p,\square_0,\mathfrak{B}^{\mathfrak{c}}\restriction\square_0)$, $p\le p_{\square}$, is a
$(\Sigma^{\ast})$-triple for every cover cube $\square$ whose enlargement $\square_0$ meets
$W_{\mathfrak{a}}$. This includes the case $\square_0\not\subset W_{\mathfrak{a}}$. When
$j^+=k_a$, the description of junction boxes in Lemma~\ref{4.26:lem-retained-levels} gives the
levels and weights of
$\mathfrak{B}^{\mathfrak{c}}\restriction\square_0$ on the two cells listed in item $(\alpha)$
above; this is $\mathfrak{s}_1$.

\emph{Class $(\beta)$.} At a lifted point in the interior of $W_{\mathfrak{a}}$, the triple
$(p,\square_0,\mathfrak{B}^{\mathfrak{c}}\restriction\square_0)$ has
$\lambda_{\mathfrak{s}}=\mathfrak{c}$. Clauses (a) and $(\mathrm{a}')$ of the lemma on the frozen
weights of a live piece, which bound the frozen weights on the enlargement $\square_0$ of a cover
cube, show
that on $\square_0$ at most two adjacent values of $\mathfrak{c}$ occur, and at most two adjacent
frozen levels. A larger $X'$ meeting $W_{\mathfrak{a}}$ carries, by the first assertion, each
retained $\mathbb{B}^{(n')}_{k_a}\restriction(X'\cap W_{\mathfrak{a}})$ with $n'<n^+_a$; these
triples are also hybrid. At an
un-lifted inner collar one has $\omega=1$, and the structure carries $w_{\mathfrak{s}}=1$ at level
$\lambda_{\mathfrak{s}}=\ell_{\mathfrak{a}}-1$.

\emph{Class $(\gamma)$.} Suppose $j^+<k_a$. By the description of junction boxes in
Lemma~\ref{4.26:lem-retained-levels}, a box that
reaches $\partial W_{\mathfrak{a}}$ meets no lifted point. At such a box
$\mathfrak{B}^{\mathfrak{c}}=\mathfrak{B}$, so the structure is the current one and no triple is
added.

\emph{No new clause.} The piece's own run, at density $k_a$, imposed the derived entries $2$ and $4$
for all triples at once; this is $(\Sigma^{\ast})$ read at density $k_a$. Spaces written after the
arrest read $(\Sigma^{\flat})$ and $(\Sigma^{\ast})$ unchanged, by the reading convention for such
spaces. Hence they impose the same clauses at the same triples, and none of the triples of classes
$(\alpha)$ and $(\beta)$ imposes a clause that was not already imposed. Moreover, the lemma on the
real point in the spaces written after the arrest, applied at every density $k'\ge k_a$, places the
real point in every such space.

\emph{What is new.} The inherited domain $\mathfrak{D}^{\flat}_n$ of the present site names these
triples, so each of the properties (1)--(3) of the remark must hold at each named triple; the
statements that supply them are named above.
\end{proof}

\begin{remark}\label{4.26:lem-real-background}
The statement of this lemma is not written out here; parts of its proof are printed below.
\end{remark}

\begin{lemma}[Nested criticality under composite block averages]\label{4.26:prf-nested-criticality}
This lemma is the second step of the proof of Lemma~\ref{4.26:lem-real-background}; we work in its
setting. The site $k$, the stage $n$, the
live piece $\mathfrak{a}$, the real background $U_{\ast}=U^{(n+1)}_k$ and the structure
$\mathfrak{B}_{\ast}$ for which $U_{\ast}$ is critical are the ones fixed there.

Let $(p,X',\mathfrak{s})$ be a triple of the family $(\Sigma^{\ast})^{\mathfrak{a}}$ of
Lemma~\ref{4.26:lem-retained-levels} with $\mathfrak{s}\preccurlyeq\mathfrak{B}_{\ast}$ on $X'$. Write
\[
\mathfrak{s}=\operatorname{loc}_{X'}\bigl(\operatorname{trunc}_p(\tilde{\mathfrak{s}})\bigr),
\]
where:
\begin{itemize}
\item $\tilde{\mathfrak{s}}$ is the global hybrid structure from which $\mathfrak{s}$ is obtained, and we
assume $\tilde{\mathfrak{s}}\preccurlyeq\mathfrak{B}_{\ast}$ on $X'$ (the refinement step of the proof
of Lemma~\ref{4.26:lem-real-background});
\item $\operatorname{trunc}_p$ replaces every level $>p$ by $p$ inside the same $L$-adic blocks;
\item $\operatorname{loc}_{X'}$ is the localisation to $X'$ used in
Lemma~\ref{4.26:lem-retained-levels}.
\end{itemize}

\emph{Block averages.} Let $M$ be Ba{\l}aban's block-averaging map
(Definition~\ref{1.5:def-block-average}). Let $M^i=M\circ M^{i-1}$ be its $i$-fold iterate
\cite{B-CMP98}. By \cite[(2.11)]{B-CMP119}, the average attached to a structure equals the $j$-fold
iterate $M^j$ on the region where the structure has level $j$; this is what it means for a structure
to have level $j$ on a region.

\emph{Neighbourhoods of bonds.} For a level-$i$ bond $c$ with endpoints $c_-,c_+$, let
$B^i(c_\pm)$ be the level-$i$ blocks containing $c_\pm$, and put
$R^i(c):=B^i(c_-)\cup B^i(c_+)$.

\emph{Feasible sets.} Let $E$ be the set of bonds of the finest lattice, of spacing $\xi$
(Definition~\ref{4.4:def-regular-backgrounds}), contained in $X'$; in the feasible sets below these are
the only bonds on which $U'$ may differ from $U_{\ast}$. For a family $\mathcal{C}$ of
functionals of configurations put
\begin{equation}\label{4.26:eq-nested-criticality-1}
\operatorname{Feas}(\mathcal{C}):=\bigl\{U':\ U'=U_{\ast}\ \text{off }E,\ \
f(U')=f(U_{\ast})\ \text{for all } f\in\mathcal{C}\bigr\}.
\end{equation}

\emph{Constraint families of a structure.} For a structure $\mathfrak{t}$, let
$\mathcal{C}(\mathfrak{t})$ be the family of functionals $U'\mapsto M^{\lambda}U'(b)$, where $\lambda$
runs over the levels of $\mathfrak{t}$ and $b$ over the level-$\lambda$ bonds that $\mathfrak{t}$
retains. Write
$M^{\mathfrak{t}}U'$ for the collection of their values, and put
$\operatorname{Feas}(\mathfrak{t}):=\operatorname{Feas}(\mathcal{C}(\mathfrak{t}))$. We call the members
of $\mathcal{C}(\mathfrak{t})$ the constraints, or conditioning functionals, of $\mathfrak{t}$; those
of level $\lambda$ are its retained level-$\lambda$ functionals. A \emph{cell} of
$\mathfrak{t}$ is a region on which $\mathfrak{t}$ has constant level; an \emph{interface} is the
common boundary of two cells.

Assume the following four statements.
\begin{itemize}
\item[(I)] \emph{Nesting of conditioning functionals along a run.} Each stage of a run integrates
the old variables subject to the condition that the new block field is the one-step average of the
old one \cite{B-CMP122a}. Consequently every
conditioning functional of a later stage structure is a one-step average of conditioning
functionals of the previous stage structure of the same run. In the form of Ba{\l}aban's bond
bookkeeping \cite[(2.11)--(2.12)]{B-CMP119}, for the structure
$\tilde{\mathfrak{s}}$, which refines $\mathfrak{B}_{\ast}$ on $X'$, this gives
\begin{equation}\label{4.26:eq-nested-criticality-2}
\bigl\{U':\ M^{\tilde{\mathfrak{s}}}U'=M^{\tilde{\mathfrak{s}}}U_{\ast}\bigr\}
\subset\bigl\{U':\ M^{\mathfrak{B}_{\ast}}U'=M^{\mathfrak{B}_{\ast}}U_{\ast}\bigr\}\ni U_{\ast}.
\end{equation}
It also gives the following: if $c$ is a level-$j$ bond of $\mathfrak{B}_{\ast}$ such that $R^j(c)$
crosses an interface of cells of $\tilde{\mathfrak{s}}$ that is not an interface of
$\mathfrak{B}_{\ast}$, then $M^jU'(c)$ is a function of finitely many members of
$\mathcal{C}(\tilde{\mathfrak{s}})$ and of $U'\restriction E^c$.
\item[(II)] \emph{The localised problem.} The localised structure on $X'$ is Ba{\l}aban's
localised structure of top level $p$ on $X'$, written $\mathbb{B}_p(X')$, together with the regions
$\{\Gamma^{\mathfrak{s}}_i\cap X'\}_{i<p}$, where
$\Gamma^{\mathfrak{s}}_i$ is the level-$i$ region of $\operatorname{trunc}_p(\tilde{\mathfrak{s}})$,
which for $i<p$ is that of $\tilde{\mathfrak{s}}$. It carries a
staircase near $\partial X'$; we write $X'^{\flat}\subset X'$ for the complement of the staircase in
$X'$, so that the staircase is $X'\setminus X'^{\flat}$. It also carries a boundary gauge layer of
width at most $2M$ at the corresponding scale \cite{B-CMP122b}, where here $M=L^{m'}$ is the
auxiliary parameter of the glossary (Notation~\ref{0.2:not-glossary}) and not the averaging map. The
configuration is extended as equal to the background $\mathbf{U}$ (here $\mathbf{U}=U_{\ast}$) outside
$X'$ \cite{B-CMP109}. The constraints of this structure refine those of
$\operatorname{trunc}_p(\tilde{\mathfrak{s}})$ on $X'$, and the exterior is fixed. The exact
constraint set of the localised problem is the one
of \cite{B-CMP109}. Off the staircase and the boundary layer the localised structure is
$\operatorname{trunc}_p(\tilde{\mathfrak{s}})$, and $X'^{\flat}\supset X'^{\sim2}$, where $X'^{\sim2}$
is the twice-shrunk domain of Lemma~\ref{4.26:lem-real-background}.
\item[(III)] \emph{Identification of the localised minimiser.} Suppose $\mathbf{U}=U_{\ast}$
satisfies the three direct regularity conditions on the background that the identification of
localised minimisers requires, with Ba{\l}aban's $\alpha_0'$ sufficiently small \cite{B-CMP109}.
Then the localised problem has, in Ba{\l}aban's analyticity domain, a critical configuration that is
unique modulo gauge, and there is a gauge transformation $u$ on $X'$ such that the solution
$U_p(X',M^{\mathfrak{s}}U_{\ast})$ of the localised problem on $X'$ at top level $p$ with data
$M^{\mathfrak{s}}U_{\ast}$ equals $(U_{\ast}\restriction X')^u$ off the boundary layer.
\item[(IV)] \emph{Direct bound on the stage background.} At every point $y$ at which the direct
regularity conditions of hypothesis~\textup{(III)} are imposed, with $m=\ell(y)$ and $\omega=w(y)$ the
position-dependent data of Lemma~\ref{4.26:lem-real-background}, $U_{\ast}$ satisfies these
conditions; each direct size is at most
$\tfrac12\varphi_{\mathrm{arrest}}$ times the corresponding unit built on $\omega\alpha_{\cdot,m}$, and
this product is bounded by
\begin{equation}\label{4.26:eq-nested-criticality-3}
1.4\,\omega\alpha_{\cdot,m}\ \le\ \tfrac{1.4}{5}\,\gamma_0^{(k)} .
\end{equation}
\end{itemize}

Then the following hold.
\begin{itemize}
\item[(a)] Let $\mathcal{C}_F,\mathcal{C}_G$ be families of functionals such that every
$g\in\mathcal{C}_G$ is a function of finitely many members of $\mathcal{C}_F$ and of
$U'\restriction E^c$. Then $\operatorname{Feas}(\mathcal{C}_F)\subset\operatorname{Feas}(\mathcal{C}_G)$.
Moreover, a configuration lying in $\operatorname{Feas}(\mathcal{C}_F)$ that is critical (resp.\
minimal) for the Wilson action on $\operatorname{Feas}(\mathcal{C}_G)$ is critical (resp.\ minimal)
on $\operatorname{Feas}(\mathcal{C}_F)$.
\item[(b)] For every level-$i$ bond $c$ and every $\lambda\le i$ there is a fixed function
$\Phi_{c,\lambda}$ with
\begin{equation}\label{4.26:eq-nested-criticality-a}
M^iU'(c)=\Phi_{c,\lambda}\bigl(M^{\lambda}U'(b):\ b\ \text{a level-}\lambda\text{ bond with }
R^{\lambda}(b)\subset R^i(c)\bigr).
\end{equation}
In particular, $\Phi_{c,\lambda}$ depends on the variables of $U'$ in $c\cap X'$ only through
finer averages.
\item[(c)] $\operatorname{Feas}(\tilde{\mathfrak{s}})\subset\operatorname{Feas}(\mathfrak{B}_{\ast})$.
\item[(d)] We have
$U_{\ast}\in\operatorname{Feas}(\mathfrak{s})\subset\operatorname{Feas}(\mathfrak{B}_{\ast})$, and
$U_{\ast}$ is critical for the Wilson action on $\operatorname{Feas}(\mathfrak{s})$. Moreover, for the
gauge transformation $u$ of hypothesis~\textup{(III)},
$U_p(X',M^{\mathfrak{s}}U_{\ast})=(U_{\ast}\restriction X')^u$ off the boundary layer.
\item[(e)] At every $y'\in X'^{\sim2}$ the localised structure is
$\operatorname{trunc}_p(\tilde{\mathfrak{s}})$ and the configuration is $U_{\ast}^{u}$.
\end{itemize}
\end{lemma}

\begin{proof}
\emph{Clause (a).} Let $U'\in\operatorname{Feas}(\mathcal{C}_F)$ and $g\in\mathcal{C}_G$. By
hypothesis, there are $f_1,\dots,f_r\in\mathcal{C}_F$ such that $g(U')$ depends on $U'$ only through
$f_1(U'),\dots,f_r(U')$ and $U'\restriction E^c$. Since $f_l(U')=f_l(U_{\ast})$ for every $l$
and $U'\restriction E^c=U_{\ast}\restriction E^c$, we get $g(U')=g(U_{\ast})$. Thus every $g$ is
constant on $\operatorname{Feas}(\mathcal{C}_F)$, equal to $g(U_{\ast})$. Together with $U'=U_{\ast}$
off $E$, this gives $U'\in\operatorname{Feas}(\mathcal{C}_G)$.

Now let $U$ lie in $\operatorname{Feas}(\mathcal{C}_F)$. The sets are nested, so the tangent space
of $\operatorname{Feas}(\mathcal{C}_F)$ at $U$ is contained in that of
$\operatorname{Feas}(\mathcal{C}_G)$. If $U$ is critical for the Wilson action on
$\operatorname{Feas}(\mathcal{C}_G)$, the differential of the action therefore vanishes on the
smaller tangent space. If $U$ is minimal on $\operatorname{Feas}(\mathcal{C}_G)$, it is minimal on the
subset $\operatorname{Feas}(\mathcal{C}_F)$.

\emph{Clause (b).} By the iteration $M^i=M\circ M^{i-1}$ recalled above,
$M^{i}=M^{i-\lambda}\circ M^{\lambda}$ for $\lambda\le i$.

By the definition of the block average (Definition~\ref{1.5:def-block-average}, \cite{B-CMP98}),
$M^iU'(c)$, for a level-$i$ bond $c$, is a fixed function, equivariant under the gauge group, of the
level-$(i-1)$ averaged configuration $W=M^{i-1}U'$ on the contours of level-$(i-1)$ bonds inside
$R^i(c)$. The parallel transporters of the $i$-th average are
contours of level-$(i-1)$ averaged variables, not of the bond variables of the finest lattice.

We prove \eqref{4.26:eq-nested-criticality-a} by induction on $i-\lambda$.

\emph{Base case.} For $\lambda=i$ the bond $c$ itself satisfies $R^i(c)\subset R^i(c)$, and
$\Phi_{c,i}$ is evaluation at $b=c$.

\emph{First step.} Every level-$i$ block is a union of level-$(i-1)$ blocks
(Definition~\ref{1.5:def-block-average}). Hence the level-$(i-1)$ blocks containing the endpoints
of a level-$(i-1)$ bond $b'$ inside $R^i(c)$ lie in $R^i(c)$, that is,
$R^{i-1}(b')\subset R^i(c)$. The previous paragraph is therefore
\eqref{4.26:eq-nested-criticality-a} with $\lambda=i-1$.

\emph{Inductive step.} Let $\lambda<i-1$. By the inductive hypothesis, each $M^{i-1}U'(b')$ is
$\Phi_{b',\lambda}$ of the values $M^{\lambda}U'(b)$ over bonds $b$ with
$R^{\lambda}(b)\subset R^{i-1}(b')\subset R^i(c)$. Substituting these representations gives
\eqref{4.26:eq-nested-criticality-a} for $\lambda$.

At no step of this induction is a variable of a finest-lattice bond used except through an
average. This is the last assertion of (b).

\emph{Clause (c).} Intersecting both sides of the nesting \eqref{4.26:eq-nested-criticality-2} of
hypothesis~(I) with $\{U':\ U'=U_{\ast}\ \text{off }E\}$ gives (c) directly; the cases below give its
elementary content, Case 3 being exactly the second half of hypothesis~(I).
Take a bond $c$ of $\mathfrak{B}_{\ast}$ of level $j$. We show that
$U'\mapsto M^jU'(c)$ is a function of finitely many members of $\mathcal{C}(\tilde{\mathfrak{s}})$
and of $U'\restriction E^c$, and then apply (a) with $\mathcal{C}_F=\mathcal{C}(\tilde{\mathfrak{s}})$
and $\mathcal{C}_G=\mathcal{C}(\mathfrak{B}_{\ast})$. We distinguish the following cases.

\emph{Case 1: $R^j(c)$ inside one cell of $\tilde{\mathfrak{s}}$.} Suppose $R^j(c)$ lies in one cell
$A$ of $\tilde{\mathfrak{s}}$, of level
$\lambda\le j$. By \eqref{4.26:eq-nested-criticality-a} with $i=j$, $M^jU'(c)$ is a function of the
level-$\lambda$ averages of the level-$\lambda$ bonds $b$ with $R^{\lambda}(b)\subset R^j(c)$. These
are bonds interior to $A\cap R^j(c)$. Their averages are constraints of $\tilde{\mathfrak{s}}$ by
the meaning of \emph{level $\lambda$ on $A$} \cite[(2.11)]{B-CMP119}. This case covers every
interface of $\mathfrak{B}_{\ast}$ interior to a cell of $\tilde{\mathfrak{s}}$.

\emph{Case 2: $R^j(c)$ meets no bond of $E$.} Suppose $R^j(c)$ contains no bond of $E$. By
\eqref{4.26:eq-nested-criticality-a} with
$\lambda=0$, so that $M^0U'=U'$, $M^jU'(c)$ is a function of the values of $U'$ on the bonds inside
$R^j(c)$, all of which lie in $E^c$. It is therefore a function of $U'\restriction E^c$, fixed at
its value for $U'=U_{\ast}$.

\emph{Case 3: $R^j(c)$ crosses a seam.} Suppose $R^j(c)$ crosses an interface of cells of
$\tilde{\mathfrak{s}}$ that is not
an interface of $\mathfrak{B}_{\ast}$. These are the seams left by lifting: the inner collars, the
seams of the band and of the old shells, and, at a junction box, the seam $\partial\Omega_{k_a}$.
Which finer averages near such an interface $\tilde{\mathfrak{s}}$ retains is Ba{\l}aban's bond
bookkeeping. That $M^jU'(c)$ is a composite of retained ones follows from his construction: by
induction over the stages of the run, the one-step averaging of hypothesis~(I) makes every
conditioning functional of a later structure a composite of those of any earlier structure of the
same run. This is the second half of hypothesis~(I), with the nesting
\eqref{4.26:eq-nested-criticality-2}.

Cases 1 and 2 are elementary; Case 3 rests on hypothesis~(I). In each of these cases the hypothesis
of (a) holds for the bond $c$; together with the nesting \eqref{4.26:eq-nested-criticality-2}, this
gives
$\operatorname{Feas}(\tilde{\mathfrak{s}})\subset\operatorname{Feas}(\mathfrak{B}_{\ast})$.

\emph{Clause (d): truncation.} Consider a cell of $\tilde{\mathfrak{s}}$ of level $\lambda>p$. By
\eqref{4.26:eq-nested-criticality-a}, applied with $i=\lambda$ and with $p$ in place of $\lambda$,
every retained level-$\lambda$ functional there is a composite of the level-$p$ averages interior
to the cell.

At a cell interface it is a composite of the level-$p$ averages of the cell and of the neighbouring
cell. By admissibility the neighbour's level is $\lambda\pm1$, and this is at least $p$ because
$\lambda\ge p+1$.

On cells of level at most $p$ the truncation changes nothing. The level-$p$ averages in question
are constraints of $\operatorname{trunc}_p(\tilde{\mathfrak{s}})$. By (a),
$\operatorname{Feas}(\operatorname{trunc}_p(\tilde{\mathfrak{s}}))\subset\operatorname{Feas}(\tilde{\mathfrak{s}})$.

\emph{Clause (d): localisation.} By hypothesis~(II), localisation to $X'$ refines
$\operatorname{trunc}_p(\tilde{\mathfrak{s}})$ on $X'$ and keeps the exterior fixed. The localised
problem has the constraint set of \cite{B-CMP109}. Thus every constraint of
$\operatorname{trunc}_p(\tilde{\mathfrak{s}})$ is a function of finitely many constraints of
$\mathfrak{s}$ and of $U'\restriction E^c$. By (a),
$\operatorname{Feas}(\mathfrak{s})\subset\operatorname{Feas}(\operatorname{trunc}_p(\tilde{\mathfrak{s}}))$.

\emph{Clause (d): criticality.} Altogether,
\begin{equation}\label{4.26:eq-nested-criticality-4}
U_{\ast}\in\operatorname{Feas}(\mathfrak{s})
\subset\operatorname{Feas}(\operatorname{trunc}_p(\tilde{\mathfrak{s}}))
\subset\operatorname{Feas}(\tilde{\mathfrak{s}})
\subset\operatorname{Feas}(\mathfrak{B}_{\ast}).
\end{equation}
The membership $U_{\ast}\in\operatorname{Feas}(\mathfrak{s})$ holds because $U_{\ast}$ satisfies
every defining condition in \eqref{4.26:eq-nested-criticality-1} trivially.

By the definition of $\mathfrak{B}_{\ast}$, the configuration $U_{\ast}$ is critical for the Wilson
action under the constraints of $\mathfrak{B}_{\ast}$. The set $\operatorname{Feas}(\mathfrak{B}_{\ast})$
is contained
in that constraint set, so $U_{\ast}$ is critical on $\operatorname{Feas}(\mathfrak{B}_{\ast})$, by
the nesting of tangent spaces used in (a). By (a) and
\eqref{4.26:eq-nested-criticality-4}, $U_{\ast}$ is critical for the localised problem on
$\operatorname{Feas}(\mathfrak{s})$.

\emph{Clause (d): identification.} Hypothesis~(III) gives uniqueness modulo gauge in Ba{\l}aban's
analyticity domain, in which, as stated in \cite{B-CMP109}, the localised minimiser
has approximately the same regularity properties as its background; hence, for the gauge
transformation $u$ of hypothesis~(III),
$U_p(X',M^{\mathfrak{s}}U_{\ast})=(U_{\ast}\restriction X')^u$ off the boundary layer.

The hypothesis of (III) is that $\mathbf{U}=U_{\ast}$ satisfies the direct regularity conditions
with $\alpha_0'$ sufficiently small. We derive it from hypothesis~(IV) and
Theorem~\ref{4.26:thm-density-majorant}. By hypothesis~(IV), $U_{\ast}$ satisfies the direct
regularity conditions, each direct size is at most $\tfrac12\varphi_{\mathrm{arrest}}$ times its unit
built on $\omega\alpha_{\cdot,m}$, and this product is bounded as in
\eqref{4.26:eq-nested-criticality-3}. By the bound on the density-indexed maximum by the current
coupling (Theorem~\ref{4.26:thm-density-majorant}, the corresponding display),
\begin{equation*}
1.4\,\omega\alpha_{\cdot,m}\ \le\ \tfrac{1.4}{5}\,\gamma_0^{(k)}
\ \le\ \tfrac{1.4}{5}\,\bar{\gamma}_0(\mathsf{t}_k).
\end{equation*}
The majorant $\bar{\gamma}_0(\mathsf{t}_k)$ is a function of the single variable $\mathsf{t}_k$, of
degree $-\infty$ in it; bounded in this way, the direct sizes meet the smallness of $\alpha_0'$
required in (III). No composite clause of Lemma~\ref{4.26:lem-real-background} enters this
bound, so the argument is not circular.

\emph{Clause (e).} Every clause of Lemma~\ref{4.26:lem-real-background} concerns points
$y'\in X'^{\sim2}$ only, as the corresponding bounds of \cite[(1.16)]{B-CMP109} are stated on the
twice-shrunk domain, written $X^{\sim-2}$ in \cite[(1.16)]{B-CMP109} and $X'^{\sim2}$ here. Such a
point lies off the boundary layer of width at most $2M$, with $M=L^{m'}$ as in hypothesis~(II). It
also lies off the
staircase $X'\setminus X'^{\flat}$, since $X'^{\flat}\supset X'^{\sim2}$ by hypothesis~(II).

By hypothesis~(II), at such a point the localised structure is
$\operatorname{trunc}_p(\tilde{\mathfrak{s}})$. By (d), the configuration there is $U_{\ast}^{u}$.
The exact form of the localised constraints at $\partial X'$, which enters only through
hypothesis~(II), does not affect any clause evaluated on $X'^{\sim2}$.
\end{proof}

\begin{lemma}[Plaquette and current entries of the frozen background]
\label{4.26:prf-plaquette-current-entries}
Work at the site $k$, the stage $n$ and the live piece $\mathfrak{a}$ of
Lemma~\ref{4.26:lem-real-background}, with the real background $U_{\ast}=U^{(n+1)}_k$, the real
point $\Phi^0=(U_{\ast},\mathcal{J}(U_{\ast}))$, and the structure $\mathfrak{B}_{\ast}$ for which
$U_{\ast}$ is critical, all as there. Let $(p,X',\mathfrak{s})$ be a triple of the family
$(\Sigma^{\ast})^{\mathfrak{a}}$ with $\mathfrak{s}\preccurlyeq\mathfrak{B}_{\ast}$ on $X'$, that is,
$\mathfrak{s}$ refines $\mathfrak{B}_{\ast}$ there. Let $y'\in X'^{\sim2}$. Let $m$ and $\omega$ be
the level and the weight at $y'$: the frozen data on the live piece and the data of the run
elsewhere, as in Lemma~\ref{4.26:lem-real-background}. Let $\lambda=\lambda_{\mathfrak{s}}(y')$ be
the level of $\mathfrak{s}$ at $y'$, as in Lemma~\ref{4.26:lem-unit-ordering}. Let
$\varphi_{\mathrm{arrest}}$ be the constant fixed there. The units of
\eqref{4.26:eq-three-units-a} at $y'$ are
\begin{equation*}
\mathrm{U2}_2(y')=\omega\,\alpha_{0,m}\,\ell_m^{-2},\qquad
\mathrm{U2}_4(y')=\omega\,\alpha_{0,m}\,\ell_m^{-2}\,\ell_{p\wedge\lambda}^{-1},
\end{equation*}
where $\ell_\nu$ is the block side at level $\nu$. Write $\mathsf{t}_m=\log g_m^{-2}$.

By nested criticality under composite block averages, \eqref{4.26:eq-nested-criticality-a}, the
derived configuration of the triple at $\Phi^0$ is $U_{\ast}^{u}$. Here $u$ is a gauge
transformation, and $U_{\ast}^{u}$ is the solution of the variational problem constrained by the
block averages of the structure $\mathfrak{s}^{(p),\mathrm{loc}}$, the structure $\mathfrak{s}$
localised for $p$. We write $J$ for the current of this constrained problem at $U_{\ast}^{u}$.
Let $Q_2=Q_2^{(p,X',\mathfrak{s})}(y')$ and $Q_4=Q_4^{(p,X',\mathfrak{s})}(y')$ be the plaquette
entry and the current entry at $y'$ of the composite clause of the triple, evaluated at $\Phi^0$.
Thus $Q_2$ is the size at $y'$ of the plaquettes of $U_{\ast}^{u}$ on the lattice of spacing
$\eta$ on which $X'$ lives, and $Q_4$ is the size of $J$ at $y'$.

Assume the following.
\begin{itemize}
\item[(a)] The direct plaquette bound of the arrest lemmas of
Notation~\ref{4.26:not-arrest-conventions}: at every point $x$ of its domain, with frozen data
$(\ell,w)$ at $x$,
\begin{equation}\label{4.26:eq-plaquette-current-entries-1}
|\partial U_{\ast}-1|(x)\le 1.08\,w\,\varepsilon_\ell\,L^{-2\ell},
\end{equation}
together with the floor $C_0^{\mathrm{rad}}\ge 2.16A_0/\varphi_{\mathrm{arrest}}$ on the radius
constant carried by the same arrest lemmas.
\item[(b)] The current bound of the same arrest lemmas, by which the current entries are
$(d,L)$-multiples of $\varepsilon_\ell$, read at the admissible structure $\mathfrak{s}$: there is a
constant $K_J=K_J(d,L)$ such that
\begin{equation}\label{4.26:eq-plaquette-current-entries-2}
|J|\le K_J(d,L)\,\omega\,\varepsilon_m\,\ell_m^{-2}\,\ell_{p\wedge\lambda}^{-1}
\qquad\text{on }X'^{\sim2}.
\end{equation}
The constant $K_J$ is independent of the number of strata of $W_{\mathfrak{a}}$ that $\mathfrak{s}$
meets inside $X'$. At a box triple or a junction triple, (b) is required only in the form in which
$\mathfrak{s}$ is a one- or two-level determining set on the box, with no further uniformity
condition.
\item[(c)] The standing condition of Notation~\ref{4.26:not-arrest-conventions} on the radius
constants: for every $\mathsf{t}\ge\mathsf{t}_\gamma$,
\begin{equation}\label{4.26:eq-plaquette-current-entries-3}
K_J(d,L)\,\frac{A_0}{C_0^{\mathrm{rad}}}\,\mathsf{t}^{\,p_0-q}\le\tfrac12\varphi_{\mathrm{arrest}}.
\end{equation}
\end{itemize}
Then, at $U_{\ast}$,
\begin{equation}\label{4.26:eq-plaquette-current-entries-4}
Q_e^{(p,X',\mathfrak{s})}(y')\le\tfrac12\varphi_{\mathrm{arrest}}\,\mathrm{U2}_e(y'),
\qquad e\in\{2,4\}.
\end{equation}
The case $e=2$ uses only (a), and it holds for every $\mathfrak{s}\preccurlyeq\mathfrak{B}_{\ast}$
and every $p$, with no block side entering. Moreover, the derived configuration is unitary at the
real point. Hence the clause of Lemma~\ref{4.26:lem-real-background} which is read on the
derived configuration itself holds with the same numbers.
\end{lemma}

\begin{proof}
\emph{The entry $e=2$.} On $X'^{\sim2}$, the plaquettes of the derived configuration on the lattice
of spacing $\eta$ are
\begin{equation*}
\partial(U_{\ast}^{u})=u\,(\partial U_{\ast})\,u^{-1}.
\end{equation*}
Conjugation by the $G$-valued $u$ does not change the size of $\partial U-1$. Hence
$Q_2^{(p,X',\mathfrak{s})}(y')$ is the plaquette size of $U_{\ast}$ itself at $y'$, whatever the
structure $\mathfrak{s}$ and the level $p$.

We now bound this size by hypothesis~(a), applied at $x=y'$ with $(\ell,w)=(m,\omega)$, the data
of $y'$ fixed above. The point $y'$ lies in the domain of that bound, since in the setting of
Lemmas~\ref{4.26:lem-unit-ordering} and~\ref{4.26:lem-real-background} it satisfies the
following:
\begin{itemize}
\item the characteristic function $\chi^{(n^+_a)}$ of the small-field conditions at the stage
$n^+_a$ of the arrest $\mathfrak{a}$ equals $1$ at $y'$;
\item the condition \cite[(1.24)]{B-CMP122a} holds at $y'$;
\item the powers of $L_0$ actually carried by the weight at $y'$ are at most those of the
clause's weight $\omega$.
\end{itemize}
The constant $1.08$ in \eqref{4.26:eq-plaquette-current-entries-1} already absorbs the
corrections coming from the $2M_1$-layer and from the layer of the point's own class. Hence
$|\partial U_{\ast}-1|(y')\le 1.08\,\omega\,\varepsilon_m\,L^{-2m}$.

We compare this bound with the unit
$\mathrm{U2}_2(y')=\omega\alpha_{0,m}\ell_m^{-2}$. Here $L^{-2m}$ and $\ell_m^{-2}$ are the same
factor, and the weights $\omega$ cancel, so the ratio is $1.08\,\varepsilon_m/\alpha_{0,m}$. By
the definitions of $\varepsilon_m$ through the degree function $p_0(\cdot)$ and of the radius
$\alpha_{0,m}$ through the radius constant $C_0^{\mathrm{rad}}$,
\begin{equation}\label{4.26:eq-plaquette-current-entries-5}
\frac{1.08\,\varepsilon_m}{\alpha_{0,m}}
=1.08\,\frac{A_0}{C_0^{\mathrm{rad}}}\,(\log g_m^{-2})^{p_0-q}.
\end{equation}
The floor in (a) gives $C_0^{\mathrm{rad}}\ge 2.16A_0/\varphi_{\mathrm{arrest}}$. This is
equivalent to
\begin{equation*}
1.08\,\frac{A_0}{C_0^{\mathrm{rad}}}\le\tfrac12\varphi_{\mathrm{arrest}}.
\end{equation*}
The exponent $p_0-q$ is negative; this is the decrease in $\mathsf{t}$ of the profile in~(c).
In the setting of Lemma~\ref{4.26:lem-real-background} the coupling $g_m$ at the level $m$ of $y'$
is at most $\gamma$, so $\mathsf{t}_m\ge\mathsf{t}_\gamma$; and $\mathsf{t}_\gamma\ge1$ by the
choice of $\gamma$ in that setting. Hence the factor
$(\log g_m^{-2})^{p_0-q}=\mathsf{t}_m^{\,p_0-q}$ is at most one,
the right side of \eqref{4.26:eq-plaquette-current-entries-5} is at most
$\frac12\varphi_{\mathrm{arrest}}$, and
\begin{equation*}
Q_2^{(p,X',\mathfrak{s})}(y')\le\tfrac12\varphi_{\mathrm{arrest}}\,\mathrm{U2}_2(y').
\end{equation*}
Neither $\mathfrak{s}$ nor $p$ entered this argument, and no block side occurred. So the bound
holds for every $\mathfrak{s}\preccurlyeq\mathfrak{B}_{\ast}$ and every $p$.

\emph{The entry $e=4$.} Here $Q_4$ is the derived current of the problem constrained by
$\mathfrak{s}^{(p),\mathrm{loc}}$, at its solution $U_{\ast}^{u}$. This current is the current $J$
of hypothesis~(b), so \eqref{4.26:eq-plaquette-current-entries-2} bounds it on $X'^{\sim2}$.

We compare with $\mathrm{U2}_4(y')=\omega\alpha_{0,m}\ell_m^{-2}\ell_{p\wedge\lambda}^{-1}$. The
factors $\omega$, $\ell_m^{-2}$ and $\ell_{p\wedge\lambda}^{-1}$ cancel, so the ratio is
$K_J(d,L)\,\varepsilon_m/\alpha_{0,m}$. By the same relation between $\varepsilon_m$ and
$\alpha_{0,m}$ as in \eqref{4.26:eq-plaquette-current-entries-5}, this ratio equals
\begin{equation*}
K_J(d,L)\,\frac{A_0}{C_0^{\mathrm{rad}}}\,\mathsf{t}_m^{\,p_0-q}.
\end{equation*}
The left side of \eqref{4.26:eq-plaquette-current-entries-3} is a decreasing function of the one
variable $\mathsf{t}$. Hence hypothesis~(c) is a single condition: a bound on the supremum of this
profile over $\mathsf{t}\ge\mathsf{t}_\gamma$. We read it at $\mathsf{t}=\mathsf{t}_m$, which lies
in the range $\mathsf{t}\ge\mathsf{t}_\gamma$ as noted under the entry $e=2$, and obtain
\begin{equation*}
Q_4^{(p,X',\mathfrak{s})}(y')\le\tfrac12\varphi_{\mathrm{arrest}}\,\mathrm{U2}_4(y').
\end{equation*}

\emph{How hypothesis~(b) reads Ba{\l}aban's current bound.}
In \cite[\S1, (1.16)]{B-CMP109}, Ba{\l}aban considers the functions $U_n(M^{\cdot}(U))$ built from
the $n$-fold block average. In the notation of \cite{B-CMP109}, he bounds their current by
\begin{equation*}
|J_n(M^{\cdot}(U))|<2B_3\,\alpha_1'\,(L^n\xi)^2\qquad\text{on }X^{\sim2}.
\end{equation*}
In words: the current of the constrained minimiser is bounded by $B_3$ times the direct size of
$U$, in the normalisation of the structure, namely the block side of the structure's level at the
point. Together with \cite[Theorem~1]{B-CMP102b}, this is the bound that exhibits the current
entry as a $(d,L)$-multiple of $\varepsilon_\ell$.

Hypothesis~(b) uses this bound at the admissible structure $\mathfrak{s}$, with the output in the
same normalisation. At $y'$ the direct size of $U_{\ast}$ is $\omega'\varepsilon_m$ with
$\omega'\le\omega$, and the block side of the level of $\mathfrak{s}$ at $y'$ is
$\ell_{p\wedge\lambda}$. These are the two quantities that enter
\eqref{4.26:eq-plaquette-current-entries-2}.

The statement of \cite[(1.16)]{B-CMP109}, with \cite[Theorem~1]{B-CMP102b}, is made for every $n$,
the $n$-fold average being the average of a one-level structure. Multi-level determining sets carry
the same averaging condition, \cite[(1.24), condition~(iv)]{B-CMP122a}, also
\cite[(2.34)--(2.39)]{B-CMP119} and \cite[(1.64)--(1.69)]{B-CMP122a}. What is used here is the
bound at the multi-level structure $\mathfrak{s}$, with a constant uniform in the number of strata
that $\mathfrak{s}$ meets inside $X'$. This form, uniformity included, is hypothesis~(b); the lemma
is proved as an implication from it and does not derive it from the one-level statement of
\cite{B-CMP109}.

\emph{Box triples and junction triples.} For these triples the uniformity condition in~(b) plays no
role, for the following reason.

Consider first a box triple, $\mathfrak{s}=\mathfrak{B}^{\mathfrak{c}}\restriction\square_0$ with
$\square_0=\tilde{\square}^5$. The box $\square_0$ meets at most two adjacent
$\mathfrak{c}$-shells and at most two adjacent frozen levels. This follows from clauses (a) and
(a$'$) of the lemma bounding the frozen weight and of the lemma bounding the weight on a box. At a
junction box, $\square_0$ meets two cells, by the junction condition on junction boxes.

In either case the structure on $\square_0$ is a one- or two-level determining set. The transfer of
the current bound to it is therefore the ordinary transfer to two adjacent shells, the one used on
every pair of adjacent shells produced by the operation $\mathbf{T}$. For such a pair,
\cite{B-CMP119} states the following. On the domain $\Omega_j\setminus\Omega_{j+1}$ the
configuration $U_{k+1}$ coincides with a configuration of $U_j$ type. It has the same regularity
properties and bounds, with $\eta=L^{-k}$ replaced by $L^{-n}$.

The identity of \cite{B-CMP119} is thus used as stated, at a one- or two-level structure, exactly
as on shells produced by $\mathbf{T}$. At box and junction triples, hypothesis~(b) is therefore
required only in this two-level form, and no condition enters beyond the standing assumptions
already in force.

\emph{The derived configuration.} At the real point, $U_{\ast}$ is $G$-valued and $u$ is a gauge
transformation. Hence the derived configuration $U_{\ast}^{u}$ is unitary. The clause of
Lemma~\ref{4.26:lem-real-background} which is read on the derived configuration itself therefore
holds with the same numbers as above.

Combining the two entries gives \eqref{4.26:eq-plaquette-current-entries-4}. These are the cases
$e=2$ and $e=4$ of the first inequality of the corresponding display.
\end{proof}

\begin{lemma}[Inclusion of the source in the stage spaces, and the increments]
\label{4.26:prf-source-inclusions}
This lemma gives the inclusions stated after the corresponding display in
Lemma~\ref{4.26:lem-real-background}. Work in the setting and under the hypotheses of
Lemma~\ref{4.26:lem-real-background}: site $k$ and stage $n$, a live piece $\mathfrak{a}$, a
triple $(p,X',\mathfrak{s})\in(\Sigma^{\ast})^{\mathfrak{a}}$ with
$\mathfrak{s}\preccurlyeq\mathfrak{B}_{\ast}$ on $X'$, and a point $y'\in X'^{\sim2}$. Let
$m:=\ell_{\mathfrak{a}}(y')$, $\omega:=w_{\mathfrak{a}}(y')$ and
$\lambda:=\lambda_{\mathfrak{s}}(y')\le m$, and let $e\in\{2,4\}$. Let $\mathrm{U2}_e$ and
$\mathrm{U3}$ be the units of Definition~\ref{4.26:def-three-units}. Write
$Q_e=Q_e^{(p,X',\mathfrak{s})}(y')$ for the composite entry of the real background $U_{\ast}$
of that lemma.

For a block side $l'\le m$ let $r_e(m,l')$ denote the radius of entry $e$ at frozen level $m$
and block side $l'$, as in Definition~\ref{4.26:def-three-units}:
\[
r_4(m,l'):=\alpha_{0,m}\ell_m^{-2}\ell_{l'}^{-1},\qquad r_2(m,l'):=\alpha_{0,m}\ell_m^{-2}.
\]
With these radii $\omega\, r_e(m,p\wedge\lambda)=\mathrm{U2}_e$
(Definition~\ref{4.26:def-three-units}). Put
\[
l_{y'}:=\min\{p,\mathsf{m}^{\circ}(y')\}.
\]
In the frame of (IV) below $\mathsf{m}^{\circ}(y')=\lambda$, so $l_{y'}=p\wedge\lambda$ and
$\omega\,r_e(m,l_{y'})=\mathrm{U2}_e$.
Assume the following.
\begin{enumerate}
\item[(I)] The source factor satisfies $\varphi_{\mathrm{arrest}}\le F^{(k)}\le 1$. We write
$F:=F^{(k)}(y')$.
\item[(II)] The extended comparison of the source $\tilde{U}^{\mathrm{arrest}}_{k'}$ of a class
$k'>k_a$ (in particular of the site's own class $k'=k$) with the class letters of
\eqref{4.26:eq-ownership-rule-a} holds in its clauses (a) and (a$'$):
\begin{itemize}
\item at $\lambda=m$, the source lies inside the class letters;
\item at $\lambda<m$, the ratio of a class letter's threshold to the source's threshold is at
least $3$.
\end{itemize}
\item[(III)] At points of $Z^{\mathrm{on}}$, the image under the one-step tube of a member of
$\mathfrak{L}_{l+1}$ lies in $\mathfrak{L}_l$, as a consequence of the following three
statements:
\begin{itemize}
\item the attempt spaces $\mathfrak{S}^{(k),\mathrm{arrest}}_l$ lie in $\mathfrak{N}^{(n)}$;
\item clause (c) of the extended comparison;
\item the first clause of the arrest theorem concerning $Z^{\mathrm{on}}$.
\end{itemize}
\item[(IV)] Two bounds hold generically.
\begin{itemize}
\item The one-stage bound on the move of a composite entry holds in every frame and for every
triple of the family $(\Sigma^{\ast})$: for an entry with $Q_e<1.4\,\mathrm{U2}_e$, the move of
$Q_e$ over one stage is at most $\pi_{\times}Q_e+\pi\,r_e(m,l_{y'})$. This includes the general
clause of that bound, in which the frame structure $\mathfrak{B}$ has level $\mathsf{m}=m$ and
the parent structure $\mathfrak{s}$ has level $\mathsf{m}^{\circ}=\lambda\le m$.
\item The bound on the derived entries holds for the triples $(p,X',\mathfrak{s})$ of
$\mathfrak{E}$, with its third regularity condition taken around the source's derived
configuration. It is stated generically in the structure $\mathfrak{s}$.
\end{itemize}
\end{enumerate}
Then the following hold.
\begin{enumerate}
\item[(i)] Let $\theta\ge\varphi_{\mathrm{arrest}}$ and $\mathfrak{U}\ge\mathrm{U2}_e$. For every
composite clause of the form $Q_e<\theta\cdot\mathfrak{U}$,
\begin{equation}\label{4.26:eq-source-inclusions-1}
Q_e\le\tfrac12\varphi_{\mathrm{arrest}}\,\mathrm{U2}_e\le\tfrac12\theta\,\mathfrak{U}.
\end{equation}
This applies in particular to three kinds of clause:
\begin{itemize}
\item the source's clause $Q_e<F\cdot\omega\,r_e(m,l')$ whenever its block side satisfies
$l'\le p\wedge\lambda$, and so at $l'=l_{y'}$;
\item the shell clause of $\mathfrak{N}^{(n)}$;
\item every birth-radius letter $\theta\cdot\mathrm{U3}$ with $\theta\ge\varphi_{\mathrm{arrest}}$
(the typing of \eqref{4.26:eq-ownership-rule-a} by the writer's class and the structure's
birth run).
\end{itemize}
\item[(ii)] The source lies in the stage spaces and letters in the following sense.
\begin{itemize}
\item At points of $\mathfrak{T}$ the source lies in $\mathfrak{N}^{(n)}$, clause by clause,
in the one unit $\mathrm{U2}$: each clause of the source implies the corresponding clause of
$\mathfrak{N}^{(n)}$.
\item The source lies inside the class letters.
\item At points of $Z^{\mathrm{on}}$, $\mathfrak{L}_l$ contains the image under the tube
of a member of $\mathfrak{L}_{l+1}$. At the run's own retained stages this holds with the
same unit $\mathrm{U3}$ on both sides, with room $\rho^{(l+1)}_{\lambda}$. At the retained
structures of an older core piece it holds with the same unit $\mathrm{U2}$ on both sides,
with room $\rho^{(l+1)}_m$.
\item The shell clause of $\mathfrak{N}^{(n)}$ is not read at points of $Z^{\mathrm{on}}$.
\end{itemize}
\item[(iii)] At a retained triple, the increment of a composite entry over one stage is
bounded by the one-stage bound of (IV); its additive term is
$\pi\,r_e(m,l_{y'})=\pi\,\mathrm{U2}/\omega$, so the bound is native to the unit
$\mathrm{U2}$. The derived-entry bound of (IV) applies with $\mathfrak{s}$ a retained
structure.
\end{enumerate}
\end{lemma}

\begin{proof}
\emph{Step 1: the bound \eqref{4.26:eq-source-inclusions-1}.} The hypotheses of
Lemma~\ref{4.26:lem-real-background} are in force, so by that lemma the real background
satisfies the corresponding display. That is,
$Q_e\le\frac12\varphi_{\mathrm{arrest}}\,\mathrm{U2}_e$, which is the first inequality in
\eqref{4.26:eq-source-inclusions-1}. For the second, let $\theta\ge\varphi_{\mathrm{arrest}}$ and
$\mathfrak{U}\ge\mathrm{U2}_e$. All of these quantities are nonnegative, so
\[
\tfrac12\varphi_{\mathrm{arrest}}\,\mathrm{U2}_e\le\tfrac12\theta\,\mathrm{U2}_e
\le\tfrac12\theta\,\mathfrak{U}.
\]

\emph{Step 2: the three kinds of clause in (i).} Each of them has the form
$Q_e<\theta\cdot\mathfrak{U}$ with $\theta\ge\varphi_{\mathrm{arrest}}$ and
$\mathfrak{U}\ge\mathrm{U2}_e$.
\begin{itemize}
\item \emph{The source's clause.} The threshold is $\theta=F$, and $F\ge\varphi_{\mathrm{arrest}}$
by (I). The unit is $\mathfrak{U}=\omega r_e(m,l')$. For $e=2$ the radius does not depend on
$l'$, so $\mathfrak{U}=\mathrm{U2}_2$. For $e=4$, the length $\ell_{l'}=L^{l'}\eta$ increases
with $l'$. Hence $l'\le p\wedge\lambda$ gives $\ell_{l'}^{-1}\ge\ell_{p\wedge\lambda}^{-1}$, and
so $\omega r_4(m,l')\ge\omega r_4(m,p\wedge\lambda)=\mathrm{U2}_4$. The side
$l_{y'}=\min\{p,\mathsf{m}^{\circ}(y')\}$ is such a side, because $\mathsf{m}^{\circ}(y')$ is
the parent level $\lambda$ of the frame in (IV).
\item \emph{The shell clause of $\mathfrak{N}^{(n)}$.} By Notation~\ref{4.26:not-stage-space}
and the rule (R-foreign) in \eqref{4.26:eq-ownership-rule-a}, at a retained structure of
$\mathfrak{a}$ this clause is read with unit $\omega r_e(m,l_{y'})=\mathrm{U2}_e$. Its factor
is $1+2\beta-\sum\iota^1\ge1.383$, and $1.383\ge\varphi_{\mathrm{arrest}}$.
\item \emph{A birth-radius letter $\theta\cdot\mathrm{U3}$ with $\theta\ge\varphi_{\mathrm{arrest}}$.}
By Lemma~\ref{4.26:lem-unit-ordering}, \eqref{4.26:eq-unit-ordering-a} gives
$\mathrm{U3}\ge\mathrm{U2}$.
\end{itemize}

\emph{Step 3: the source at $\mathfrak{T}$.} This is consequence (c1) of the rule for typing
composite clauses at retained structures, \eqref{4.26:eq-ownership-rule-a}. At a point of
$\mathfrak{T}$ the clauses of the source and
those of $\mathfrak{N}^{(n)}$ are read in the one unit $\mathrm{U2}$. The source's factor is
$F\le1$ by (I), and the factor of $\mathfrak{N}^{(n)}$ is at least $1.383$. Since
$F\le1<1.383$, every clause of the source implies the corresponding clause of
$\mathfrak{N}^{(n)}$.

\emph{Step 4: the source inside the class letters.} This is consequence (c2) of
\eqref{4.26:eq-ownership-rule-a}.
\begin{itemize}
\item At $\lambda=m$ it is clause (a) of the extended comparison, by (II).
\item At $\lambda<m$ it is clause (a$'$), by the ratio at least $3$ in (II). In the form
given in (c2) of \eqref{4.26:eq-ownership-rule-a}, this reads as follows. At $\lambda<m$ the
ordering of the units in Lemma~\ref{4.26:lem-unit-ordering}, \eqref{4.26:eq-unit-ordering-a},
gives $\mathrm{U2}\le\mathrm{U3}/31.5$, and $31.5\,\varphi_{\mathrm{arrest}}=3.087$ since
$\varphi_{\mathrm{arrest}}=0.098$. Hence, for every bracket
$\theta(\lambda)\ge\varphi_{\mathrm{arrest}}$,
\begin{align*}
F\cdot\mathrm{U2}\le\mathrm{U2}\le\frac{\mathrm{U3}}{31.5}
=\frac{\varphi_{\mathrm{arrest}}}{3.087}\,\mathrm{U3}
\le\frac{\theta(\lambda)}{3.087}\,\mathrm{U3}<\tfrac13\,\theta(\lambda)\,\mathrm{U3},
\end{align*}
so the source's composite threshold is less than one third of the letter's. The lower bound
$\theta(\lambda)\ge\varphi_{\mathrm{arrest}}$ holds at every level $\lambda$ by
Lemma~\ref{4.26:lem-lower-bracket}(i).
\end{itemize}

\emph{Step 5: the tube at $Z^{\mathrm{on}}$.} By (III), $\mathfrak{L}_l$ contains the image
under the tube of a member of $\mathfrak{L}_{l+1}$. The two cases are read as follows.
\begin{itemize}
\item \emph{The run's own retained stages.} Under the rule (R-own) of
\eqref{4.26:eq-ownership-rule-a}, the composite clauses are read in the birth radius
$\mathrm{U3}$ with the running bracket at $\lambda$. $\mathfrak{L}_{l+1}$ and
$\mathfrak{L}_l$ are therefore typed in the same unit $\mathrm{U3}$, and the
room is $\rho^{(l+1)}_{\lambda}$.
\item \emph{The retained structures of an older core piece.} The enclosing run's spaces are
foreign to the structures of $\mathfrak{a}$ under the rule (R-foreign) of
\eqref{4.26:eq-ownership-rule-a}, so both spaces are typed at $\mathrm{U2}$ with brackets
$\theta^{(l)}(m)$ and $\theta^{(l+1)}(m)$. By Lemma~\ref{4.26:lem-enclosing-tube}(a), the
difference of these brackets is exactly $\rho^{(l+1)}_m$, because $\sigma_{\mathfrak{a}}$ does
not depend on $l$; this difference is the room $\rho^{(l+1)}_m$ of (ii), in the unit
$\mathrm{U2}$. By clause (c) of the same lemma, under its condition
$\iota_{l+1}(y')\le\frac12\rho^{(l+1)}_m$, the stage increment consumes at most half of this
room.
\end{itemize}
The shell clause of $\mathfrak{N}^{(n)}$ is quantified over points $y'\in X'^{\sim2}\cap\mathfrak{T}$
only (Notation~\ref{4.26:not-stage-space}). It is therefore never the clause compared at a point
of $Z^{\mathrm{on}}$.

\emph{Step 6: the increments.} By (IV), the one-stage bound holds in every frame and for
every triple of $(\Sigma^{\ast})$. In its general clause the frame structure $\mathfrak{B}$ has
level $\mathsf{m}=m$ and the parent structure $\mathfrak{s}$ has level
$\mathsf{m}^{\circ}=\lambda\le m$. The re-blocked case of that clause,
in which $\mathsf{m}^{\circ}(y')<m$ and the block side at the point is $l_{y'}$, is exactly
the case of a retained structure. The additive term of the bound is
\[
\pi\,r_e(m,l_{y'})=\pi\,\mathrm{U2}/\omega,
\]
so the bound is stated in the unit $\mathrm{U2}$, as in consequence (c3) of
\eqref{4.26:eq-ownership-rule-a}.

The derived-entry bound of (IV) is stated for the triples $(p,X',\mathfrak{s})$ of
$\mathfrak{E}$, around the source's derived configuration, with no condition on the structure
$\mathfrak{s}$. It therefore applies when $\mathfrak{s}$ is a retained structure. The derived
entries so bounded satisfy the clauses of the family $(\Sigma^{\flat})$ of derived entries,
which is the last space in the chain of inclusions, by the first clause of the lemma
introducing that family.
\end{proof}

\begin{lemma}[Nesting of the retained structures at junction boxes]\label{4.26:prf-junction-nesting}
This is the case of junction boxes in the proof of Lemma~\ref{4.26:lem-real-background}, whose
setting and notation are used. Let $\mathfrak{a}$ be a live piece with $j^+_a=k_a$, so that
$n^+_a=N_a$ and, by Lemma~\ref{4.26:lem-retained-levels}, the branch of $\mathfrak{a}$ in the
definition of arrests is one of B5--B7, and
$\mathfrak{B}_{\ast}$ is the frozen structure on $W_{\mathfrak{a}}$. Let $\square\in\mathcal{C}^1$ be
such that $\square_0=\tilde{\square}^5$ meets $W_{\mathfrak{a}}$ and
$\square_0\not\subset W_{\mathfrak{a}}$,
and put $\mathfrak{s}_1:=\mathfrak{B}^{\mathfrak{c}}\restriction\square_0$. Then
$\mathfrak{s}_1\preccurlyeq\mathfrak{B}_{\ast}$ on $\square_0$, and every structure of
$(\Sigma^{\ast})^{\mathfrak{a}}$ on $\square_0$ refines $\mathfrak{B}_{\ast}$. For every level
$p\le k_a$ the triple carried by $\square_0$ at level $p$ satisfies the conclusion of
Lemma~\ref{4.26:lem-real-background}; at $p=k_a$, on $\square_0\cap W_{\mathfrak{a}}$, the ratio of
the units satisfies $\mathrm{U2}/\mathrm{U3}<0.032$ under the standing caps $L_0^2\le L/3$,
$1+\beta_0\le 8/7$ and $L\ge 12$. The same conclusions hold when $\mathfrak{a}$ is
nested in a second live piece and when $W_{\mathfrak{a}}\subset Z$.
\end{lemma}

\begin{proof}
Throughout, a pair $(i,w)$ attached to a structure on a set means that the structure has level
$i$ and weight $w$ at every point of that set. We write $(i,w\mid i',w')$ on $(A\mid B)$ when
it is $(i,w)$ on $A$ and $(i',w')$ on $B$. For a structure $\mathfrak{s}$ and a level $p$,
$\mathfrak{s}^{(p)}$ is its truncation at $p$: levels above $p$ are replaced by $p$ inside the same
$L$-adic blocks. We call $\square_0\cap W_{\mathfrak{a}}$ the \emph{inner cell} and
$\square_0\setminus W_{\mathfrak{a}}$ the \emph{outer cell} of the box.

\emph{Geometry of the box.} We use the geometric description of junction boxes in the construction of
$\mathfrak{B}^{\mathfrak{c}}$. The outer cell lies in the old shell of index $k_a$. There it lies
within seven side lengths of the cubes of $\pi_{k_a-1}$ from $\partial\Omega_{k_a}$ (twelve side
lengths of the cubes of $\pi_{k_a}$ for a box of scale $k_a$). It also lies at least three layers of
cubes of side $MR_{k_a+1}$ below the regions of any later run. On that set every structure from
density $k_a$ on is $(k_a,1)$. The reasons are the following. The stages at density $k_a$ act inside
$Z^{(a)}\cap\Omega^c_{k_a}$, where $Z^{(a)}$ denotes the region $Z$ of
Notation~\ref{4.26:not-stage-frozen-data} taken at the site of density $k_a$, and make no change in
the operation $\mathbf{T}_{k_a}(Z_{k_a}\cap (Z^{(a)})^c)$, as stated in \cite{B-CMP122a}. The
determining sets on that
shell are those of \cite[(1.14)--(1.17)]{B-CMP122a}, of level $k_a$ and weight $1$. Later steps do
not reach the outer cell, because of the three separating layers.

The inner cell lies in the old shell of index $k_a-1$, which is contained in
$\Omega^c_{k_a}\setminus Z''_{k_a-1}$. This shell is lifted by the last stage only. For every
$n'<N_a$ the stage structure $\mathbb{B}^{(n')}_{k_a}$ is $(k_a-1,1)$ there. This is
\cite[(1.67)]{B-CMP122a} at $m=j'=k_a-1$, where the weight factor is $L_0^{0}=1$; here and below
$L_0$ denotes the base of the weights in \cite[(1.67)]{B-CMP122a}. The frozen
structure $\mathbb{B}^{(N_a)}_{k_a}$ is $(k_a,L_0^2)$ there, by the entry for $j^+_a=k_a$ of the
table \eqref{4.26:eq-retained-levels-a} of levels and weights together with
\cite[(1.67)]{B-CMP122a}. Consequently, for every $n'<N_a$,
\begin{equation*}
\mathfrak{s}_1=\mathfrak{B}^{\mathfrak{c}}\restriction\square_0
=\mathbb{B}^{(n')}_{k_a}\restriction\square_0
=(k_a-1,1\mid k_a,1)\ \text{ on }\ (\square_0\cap W_{\mathfrak{a}}\mid\square_0\setminus W_{\mathfrak{a}}),
\end{equation*}
while
\begin{equation*}
\mathfrak{B}_{\ast}\restriction\square_0=(k_a,L_0^2\mid k_a,1).
\end{equation*}
In particular, $\mathfrak{B}_{\ast}$ is flat of level $k_a$ on $\square_0$.

\emph{Refinement.} The set $\Omega_{k_a}$ is a union of cubes of side $MR_{k_a}$ of the lattice of
level $k_a$ \cite{B-CMP119}. Hence $\partial\Omega_{k_a}$ runs along faces of blocks of level $k_a$.
The blocks of level $k_a-1$ of $\mathfrak{s}_1$ on the inner cell are therefore the $L$-adic
sub-blocks of the blocks of level $k_a$ of $\mathfrak{B}_{\ast}$. On the outer cell the two
structures coincide. Thus $\mathfrak{s}_1\preccurlyeq\mathfrak{B}_{\ast}$ on $\square_0$. The
interface clause of Lemma~\ref{4.26:prf-nested-criticality} imposes a condition at the interfaces of
the structure. In the present situation that condition is the last stage's own identification
$\bar{V}_{k_a-1}=V_{k_a}$ on the inner cell, in the notation of the interface clause of
Lemma~\ref{4.26:prf-nested-criticality}.

\emph{Levels $p\le k_a-1$.} Both structures truncate to the flat structure of level $p$ on
$\square_0$. Hence $\mathfrak{s}_1^{(p)}=\mathfrak{B}_{\ast}^{(p)}$, and the triple is a sub-triple of
the frozen structure itself. At every point $\lambda\ge p$, so $p\wedge\lambda=p=p\wedge m$. The
ordering \eqref{4.26:eq-unit-ordering-a} then gives $\mathrm{U2}=\mathrm{U1}$ as well. The triple is
thus one of $\mathfrak{B}_{\ast}$'s own, and Lemma~\ref{4.26:prf-nested-criticality} and
Lemma~\ref{4.26:prf-plaquette-current-entries} apply to it, giving the conclusion of
Lemma~\ref{4.26:lem-real-background}.

\emph{Level $p=k_a$.} Here the structure is $\mathfrak{s}_1$ itself. At a point of the inner cell
the data are
\begin{equation*}
(m,\omega,\mathfrak{c},\lambda,w_{\mathfrak{s}})=(k_a,\,L_0^2,\,k_a-1,\,k_a-1,\,1),
\end{equation*}
so that $p\wedge\lambda=k_a-1$. Since $\mathfrak{s}_1\preccurlyeq\mathfrak{B}_{\ast}$ on $\square_0$,
Lemma~\ref{4.26:prf-nested-criticality} and Lemma~\ref{4.26:prf-plaquette-current-entries} apply to
this triple. They give, for the entries $Q_e^{(k_a,\square_0,\mathfrak{s}_1)}(y')$, $e\in\{2,4\}$, at
every point $y'$ of the inner cell,
\begin{align*}
Q_2^{(k_a,\square_0,\mathfrak{s}_1)}(y')&\le\tfrac12\,\varphi_{\mathrm{arrest}}\,L_0^2\,
\alpha_{0,k_a}\,\ell_{k_a}^{-2},\\
Q_4^{(k_a,\square_0,\mathfrak{s}_1)}(y')&\le\tfrac12\,\varphi_{\mathrm{arrest}}\,L_0^2\,
\alpha_{0,k_a}\,\ell_{k_a}^{-2}\,\ell_{k_a-1}^{-1}.
\end{align*}
With $\omega=L_0^2$, $m=k_a$ and $p\wedge\lambda=k_a-1$, the right-hand sides are
$\tfrac12\varphi_{\mathrm{arrest}}\mathrm{U2}_2$ and $\tfrac12\varphi_{\mathrm{arrest}}\mathrm{U2}_4$.

For the ratio of the second to the third of the units \eqref{4.26:eq-three-units-a} at such a point
we use three facts. The first is the bound $\alpha_{0,k_a}\le(1+\beta_0)\,\alpha_{0,k_a-1}$ on the
ratio of consecutive radii (Lemma~\ref{4.4:lem-radii-ratio}). The second and the third are the
standing caps $L_0^2\le L/3$, and $1+\beta_0\le 8/7$ together with $L\ge 12$, caps in the sense of
Definition~\ref{2.3:def-cap}. Then
\begin{align*}
\frac{\mathrm{U2}}{\mathrm{U3}}
&=L_0^2\,\frac{\alpha_{0,k_a}}{\alpha_{0,k_a-1}}\,L^{-2}
\le\frac{L}{3}\,(1+\beta_0)\,L^{-2}\\
&=\frac{1+\beta_0}{3L}\le\frac{8/7}{36}<0.032 .
\end{align*}
So the triple lies inside the composite clause of Lemma~\ref{4.26:lem-real-background} with margin
$\tfrac12$, in either of the units $\mathrm{U2}$ and $\mathrm{U3}$. On the outer cell $m=\lambda=k_a$
and $\omega=1$, and the triple is the run's own flat triple.

\emph{Nested pieces.} Suppose $\mathfrak{a}\subset\mathfrak{b}$ with $\mathfrak{b}$ a second live
piece. The outer cell $\square_0\setminus W_{\mathfrak{a}}$ lies in the collar
$X_{\mathfrak{a}}\cap\Omega_{k_a}$ of $\mathfrak{a}$. By the lemma on nested arrests, $k_a\le h_b$
and
\begin{equation*}
X_{\mathfrak{a}}\cap\Omega_{k_a}\subset Z_{k_a}\subset Z_{h_b},
\end{equation*}
and $Z_{h_b}$ is the core of $\mathfrak{b}$. By the same lemma, every structure of the run of
$\mathfrak{b}$ keeps $(k_a,1)$ there. The description of the outer cell above is therefore unchanged.

\emph{Pieces in the core.} Suppose $W_{\mathfrak{a}}\subset Z$, so that $k_a\le h$. The run of the
present site re-levels only outside $Z''_{h+1}$, and $Z''_{h+1}\supseteq Z'^{\sim 3}_h$; more
precisely, by the placement of the stage cells of a run, its level-$h$ cell lies in the outer five
layers $Z'^{\sim 8}_h\setminus Z'^{\sim 3}_h$. Moreover the set $\Gamma_h\cap Z''_{h+1}$ of
\cite[(1.15)]{B-CMP122a} is never re-levelled. Junction boxes lie in the collar, and
\begin{equation*}
X_{\mathfrak{a}}\cap\Omega_{k_a}\subset Z_h\subset Z'_h\subseteq Z'^{\sim 3}_h\subseteq Z''_{h+1}.
\end{equation*}
Hence every stage structure $\mathbb{B}^{(n)}_k$ of the run is $(k_a,1)$ on them. Other components are
at least $R_k$ layers of cubes of side $M$ away, by the separation of distinct live pieces in
Notation~\ref{4.26:not-stage-frozen-data}, so no junction box meets them. A healing on the component
of $W_{\mathfrak{a}}$ dissolves $\mathfrak{a}$, by the dissolution rule for arrests, and
$\mathfrak{a}$ is then no longer a live piece.

\emph{No structure fails to refine $\mathfrak{B}_{\ast}$.} On the outer cell every structure of
$(\Sigma^{\ast})^{\mathfrak{a}}$ is flat of level $k_a$, as shown above. On the inner cell these
structures are the $\mathbb{B}^{(n')}_{k_a}\restriction\square_0$ with $n'\le N_a$. Each of them is
either $(k_a-1,1)$ or the frozen $(k_a,L_0^2)$, which is $\mathfrak{B}_{\ast}$. All of them refine
$\mathfrak{B}_{\ast}$, by the refinement argument above.

\emph{Branch B4.} By Lemma~\ref{4.26:lem-retained-levels}, $n^+_a=N_a$ occurs only at branches
B5--B7. For completeness: if an arrest of branch B4 occurred at the last stage (which
Lemma~\ref{4.26:lem-retained-levels} excludes), then
$\mathfrak{B}_{\ast}=\mathbb{B}^{(N_a-1)}_{k_a}\restriction W_{\mathfrak{a}}$, which equals
$\mathfrak{s}_1$ on $\square_0$, and the box triple is one of $\mathfrak{B}_{\ast}$'s own, so it
lies inside the composite clause of Lemma~\ref{4.26:lem-real-background} directly. The frozen
structure's own triples on the inner cell are then
covered by the direct bound on the frozen background that enters
Lemma~\ref{4.26:prf-plaquette-current-entries}, applied at
one level. This is the treatment of branch B4 in the part of the proof of
Lemma~\ref{4.26:lem-real-background} on refinement and scope.

This is the last case, and the proof of Lemma~\ref{4.26:lem-real-background} is complete.
\end{proof}

\begin{definition}[Lower bounds on the parameters and the coupling ceilings]\label{4.26:def-parameter-bounds}
The parameters are chosen in the order stated in \textup{(O)} below. A \emph{floor} is a one-sided
condition among the parameters chosen before $\gamma$. A \emph{$\gamma$-ceiling} is a one-sided
condition on $\gamma$, which is chosen last.

The following letters are those of Notation~\ref{4.26:not-density-maximum}:
\begin{itemize}
\item $\mathsf{t}_j$ and $\mathsf{t}_\gamma$;
\item the constants $C_0,C_1$ and $C=\max\{C_0,C_1\}$;
\item $q=q_0=q_1$ and the exponent $r$;
\item the integer $N_0(j)$;
\item $\Psi$, $\bar{\gamma}_0$ and $c_\beta$, condition \textup{(g)} and the range clause, all as in
\eqref{4.26:eq-density-maximum-a};
\item $\bar{N}(\mathsf{t})=L\mathsf{t}^r$, as in \eqref{4.26:eq-density-maximum-c}.
\end{itemize}
The density-indexed maximum $\gamma_0^{(k)}$ is that of
Notation~\ref{4.26:not-density-maximum}; its bound $\gamma_0^{(k)}\le\bar{\gamma}_0(\mathsf{t}_k)$ is
the corresponding display. We write $R_j=R(\mathsf{t}_j)$ for the radius at density $j$ of
Lemma~\ref{4.26:lem-weight-bound}, so that $R_k=R(\mathsf{t}_k)$, and $R_{\min}$ for the smallest of
the radii $R_j$ at the densities $j$ of the run; this $R$ is not the ball radius of the observables.
The letters $L$, $d$, $M$, $M_1$, $\kappa$, $\kappa_1$, $A_0$, $p_0$, $B_\sharp$, $\bar{B}$, $c_5$, $g$,
$\gamma$ and $\vartheta$ are as in the glossary (Notation~\ref{0.2:not-glossary}). The letters
$\delta w$, $\delta_1$, $\beta_0'$, $B_3$, $\mathsf{a}$, $C_0^{\mathrm{rad}}$, $\varphi_{\mathrm{arrest}}$,
$K_{\mathrm{cv}}$, $K_{\min}$, $K_\flat'$, $C_R^{\natural}$, $K_\Delta$, $C^{b}$, $a_\flat$,
$\kappa_{\max}$, $c_G'$, $\beta_\star$, $p_1$, $a_\mu$, $\mathsf{l}_j$, $\sigma_\ast$, $R_\ast$,
$\delta'_j$, $C_d$, $r_A$ and $\kappa_A(k)$ denote the constants and
quantities that carry these letters in the list of parameter conditions of the large-field expansion,
and each factor written $O(1)$ is the constant so written there; $P,P',P''$ are the three quantities
bounded in \textup{(c9)}. In \textup{(c5)}, $B\,C\,M$ denotes the product of the constant $B$ of that
list with the constants $C$ and $M$.

\smallskip
\noindent\textup{(A)} \emph{Absolute numbers.} These are:
\begin{equation*}
\beta=\tfrac15,\qquad \beta_0=\tfrac17,\qquad \theta=\tfrac12,\qquad
\mathsf{w}_0=\tfrac1{24},\qquad C_{1/2}=\sqrt2 .
\end{equation*}
In addition, $\theta_S$ is a number chosen in the interval $[\tfrac12,\tfrac89]$.

\smallskip
\noindent\textup{(F)} \emph{Floors.} Each floor is one-sided.
\begin{itemize}
\item[(f1)] $L\ge 12$.
\item[(f2)] The auxiliary integer $L_0$ of Ba{\l}aban's large-field construction
\cite{B-CMP122a} satisfies
\begin{equation*}
2\le L_0 \qquad\text{and}\qquad L_0^2\le L/3 .
\end{equation*}
This integer is not the threshold $L_0$ of the quantifier prefix
(Notation~\ref{2.1:not-prefix}).
\item[(f3)] $\delta w\ge 16\ln L+4$.
\item[(f4)] $M\ge 3L$ and $\delta_1 M/M_1\ge 2$.
\item[(f5)] The ratio of the radius constants satisfies
\begin{equation*}
\frac{C_1}{C_0}\ \ge\ \frac{2(1+2\beta)(1+\beta_0')^4B_3^2\,O(1)\,M\mathsf{a}}{1-\beta},
\end{equation*}
where the constant $\mathsf{a}$ is taken with $\mathsf{a}\ge(1+2\beta)(1+\beta_0)L^3$.
\item[(f6)] $C_0^{\mathrm{rad}}\ge 2.16\,A_0/\varphi_{\mathrm{arrest}}$.
\item[(f7)] $\kappa_1\ge\max\{16\kappa+1,\ 4L\kappa+1\}$.
\item[(f8)] $q\ge p_0+7r+1$, and $q_0=q_1$.
\item[(f9)] The remaining floors of the list of parameter conditions of the large-field expansion are
carried over unchanged, each member in the form in which it stands there: a group of floors taken
member by member, the floors involving $\kappa_1$, $M_1$ and $\kappa$, and two further floors.
\end{itemize}

\smallskip
\noindent\textup{(D)} \emph{Numbers fixed after $L$ and before $\gamma$.} These are:
\begin{itemize}
\item $K_{\mathrm{cv}}$;
\item $K^{\mathrm{nr}}=66B_\sharp K_{\mathrm{cv}}$;
\item $K_{\min}$;
\item $K_\flat'$, which carries two absolute numbers depending only on $d$ as factors;
\item $C_F=12(2L^2)^d$;
\item the constant $C_R^{\natural}$, into which the factor $\mathsf{a}^5$ enters.
\end{itemize}

\smallskip
\noindent\textup{(G)} \emph{The $\gamma$-ceilings.} Each ceiling is a condition in one variable.
In \textup{(G)} the letter $N$ denotes the integer $N=R_k=R(\mathsf{t}_k)$ of Ba{\l}aban's large-field
construction, not the rank of the gauge group.
Wherever a supremum over $\mathsf{t}\ge\mathsf{t}_\gamma$ is written, every letter in it that depends
on the coupling, $\vartheta$ and $N$ included, is read at one and the same $\mathsf{t}$.

For several ceilings two further data are given. The first is the \emph{profile}: the dependence
of the left side on $\mathsf{t}$, up to factors independent of $\mathsf{t}$. The second is the
\emph{degree}. Degree $-\infty$ means that the profile decays faster than every power of $\mathsf{t}$.
\begin{itemize}
\item[(c1)] The two conditions
\begin{equation*}
\mathsf{t}_\gamma\ge 2(r+q)+c_\beta
\qquad\text{and}\qquad
600(1+\beta_0)^2\Psi(\mathsf{t}_\gamma)\le 1
\end{equation*}
hold. The profile is $\mathsf{t}^{r+q}e^{-\mathsf{t}/2}$ and the degree is $-\infty$. By
Theorem~\ref{4.26:thm-density-majorant} together with \textup{(c15)}, this ceiling gives
$120\gamma_0^{(k)}\le1$ at every level $k$
(Remark~\ref{4.26:rem-no-new-condition}).
\item[(c2)] $K_\Delta\,\bar{\gamma}_0(\mathsf{t}_\gamma)\le\tfrac12$. The degree is $-\infty$.
\item[(c3)] The further smallness condition on $\gamma_0$ in the list of parameter conditions of the
large-field expansion holds, with $\gamma_0$ read as $\bar{\gamma}_0(\mathsf{t}_\gamma)$. The degree
is $-\infty$.
\item[(c4)] $C^{b}(a_\flat+\kappa_{\max})\le1$, where $a_\flat=O(\bar{\gamma}_0/5)$. The degree is
$-\infty$.
\item[(c5)] The supremum bound
\begin{equation*}
\sup_{\mathsf{t}\ge\mathsf{t}_\gamma}\
c_G'\,O(1)\,(B\,C\,M+1+K_\Delta)\,\frac{\bar{\gamma}_0(\mathsf{t})}{\vartheta}\ \le\ 1
\end{equation*}
holds. The profile is $\mathsf{t}^{r+2q-p_1}e^{-\mathsf{t}/2}$ and the degree is $-\infty$.
\item[(c6)] The supremum bound
\begin{equation*}
\sup_{\mathsf{t}\ge\mathsf{t}_\gamma}\
\beta_\star\,\bar{N}(\mathsf{t})\,(LM)^4\cdot 36\,\frac{\bar{\gamma}_0(\mathsf{t})^2}{\vartheta}
\ \le\ 1
\end{equation*}
holds, with $\bar{N}(\mathsf{t})=L\mathsf{t}^r$. The profile is
$\mathsf{t}^{3r+3q-p_1}e^{-\mathsf{t}}$ and the degree is $-\infty$.
\item[(c7)] The smallness condition of the list of parameter conditions on the constants
$\bar{G},\bar{G}',\mathsf{S},\mathsf{S}',\bar{\varepsilon}$ and on $\bar{\alpha}_0$ holds, in the form
\begin{equation*}
\begin{split}
(\bar{G}\bar{G}'-1)(1+\beta)
&+\bar{G}'\bigl[16\mathsf{S}^2e^{4\mathsf{S}\bar{\alpha}_0}+2B_3^{-1}\mathsf{S}^2
+8\mathsf{S}'^2e^{4\bar{\varepsilon}}\bigr]\bar{\alpha}_0\\
&+(\bar{G}'-1)\ \le\ \frac{\beta}{12}.
\end{split}
\end{equation*}
Here, for $i=0,1$ and within this condition only,
\begin{equation*}
\bar{\alpha}_i:=\sup_{\mathsf{t}\ge\mathsf{t}_\gamma}C_ie^{-\mathsf{t}/2}\mathsf{t}^q .
\end{equation*}
The degree is $-\infty$.
\item[(c8)] $K_{\min}\vartheta\le\mathsf{w}_0$. One of the members of which $K_{\min}$ is formed is
\begin{equation*}
6K^{\mathrm{nr}}=396B_\sharp K_{\mathrm{cv}} .
\end{equation*}
\item[(c9)] $16\max\{P,P',P''\}\le\beta\theta_S$, where $P,P',P''$ are formed with the constant
$C_{1/2}$ and with the factor $(1+k-m)^{1/2}$. The profile is $e^{-5\delta_1R(\mathsf{t})}$ times a
polynomial in $\mathsf{t}$, and the degree is $-\infty$.
\item[(c10)] $(1+\bar{B}\gamma^2)^{q+1/2}\le1+\beta_0$.
\item[(c11)] The $\gamma$-ceiling of the radius condition at parameter $\tfrac12$ of the list of
parameter conditions of the large-field expansion holds.
\item[(c12)] $(N_0+1)(N_0+3)\le N$, read at one $\mathsf{t}$. Here $N_0=N_0(j)$ is not bounded by
hypothesis. It is majorised from its defining equality in
Notation~\ref{4.26:not-density-maximum}, by means of the bound that this equality yields together with
Lemma~\ref{4.26:lem-weight-bound}. Moreover
\begin{equation*}
N=R(\mathsf{t}_k)<\bar{N}(\mathsf{t}_k).
\end{equation*}
\item[(c13)] The range threshold $\mathsf{t}_\gamma\ge 2(r+q)+c_\beta$ holds. It is the first
conjunct of \textup{(c1)}.
\item[(c14)] The inequalities \cite[(2.6), (2.9)]{B-CMP119} are used at $\beta_0=\tfrac17$. About
these inequalities \cite{B-CMP119} states:
\begin{quotation}
$\beta_0>0$ can be chosen arbitrarily small, if $g$ is sufficiently small.
\end{quotation}

What is used here at $\beta_0=\tfrac17$ is as follows. Both inequalities of
\cite[(2.6)]{B-CMP119} are taken from the two-sided flow bounds of
Theorem~\ref{4.1:thm-clause-zero}, under the ceilings
\begin{equation*}
\gamma^2\le\frac{15}{64c_5}\qquad\text{and}\qquad B_\sharp\gamma^2\le\frac{15}{49}.
\end{equation*}
The bound on the weight used here is Lemma~\ref{4.26:lem-weight-bound}; the second inequality of
\cite[(2.9)]{B-CMP119} is not invoked.
\item[(c15)] Condition \textup{(g)} of \eqref{4.26:eq-density-maximum-a} holds:
$\mathsf{t}_j\ge\mathsf{t}_\gamma$ for every density $j\le k$ of the run.
\end{itemize}
The remaining ceilings of the list of parameter conditions of the large-field expansion stand
unchanged. These are the conditions on $r_A$ and on $\kappa_A(k)$, four further members of that list,
and the following:
\begin{gather*}
N(2L)^d\mathsf{l}_j^2e^{-4a_\mu R_k/L}\le1,\qquad Ne^{-\sigma_\ast R_\ast/L}\le1,\qquad
2.2B_\sharp C_1\delta'_j\le\tfrac1{10},\\
e^{-\delta wR_k/8}\le(4^dC_dR_k)^{-1},\qquad R_{\min}\ge 5L .
\end{gather*}
The ceiling $R_{\min}\ge5L$ is implied by $\mathsf{t}_\gamma\ge(5L)^{1/r}$. Indeed, $R(\mathsf{t})$ is
non-decreasing in $\mathsf{t}$, and $R_j=R(\mathsf{t}_j)\ge R(\mathsf{t}_\gamma)$ by \textup{(g)}.

\smallskip
\noindent\textup{(O)} \emph{Order of choice.} The parameters are chosen in the order
\begin{equation*}
L\ \to\ \mathfrak{M}\ (\text{with } R_1\ge L^2)\ \to\ \kappa\ \to\ \kappa_1\ \to\ M_1\ \to\ M
\ \to\ (q,C)\ \to\ \gamma .
\end{equation*}
Here $\mathfrak{M}$ is the entry so written in the list of parameter conditions of the large-field
expansion; it is fixed immediately after $L$, and its choice is made with $R_1\ge L^2$.
The number $\beta_0$ is among the absolute numbers of \textup{(A)} and precedes $L$. The numbers
$K^{\mathrm{nr}}$, $C_{1/2}$ and $C_F$ are fixed after $L$, and $\gamma$ is chosen last. This order has
no cycle.
\end{definition}

\begin{definition}[Blocks, roots and gauge trees]\label{4.26:not-gauge-trees}
Throughout this section $\ell$ is the level of the near pair and $i$ is the index of the shell of
its cover cube, both fixed in the statement on near pairs for which the present notation is set
up; $\ell\le i$, as stated there.
Lengths are measured in the unit of level $\ell$. We put $\xi:=L^{-\ell}$. This is the
spacing of the fine lattice on which the background configuration of level $\ell$ is
defined. In Ba{\l}aban's notation it is the spacing $\xi$ with the level index $j$ replaced
by $\ell$, and his length $L^{j-1}\eta$ reads $L^{\ell-1-i}$ in the unit of level $i$
(compare \cite{B-CMP119}, where Ba{\l}aban's $\eta=L^{-k}$ is replaced by $L^{-n}$; there
$\eta$, $k$ and $n$ are Ba{\l}aban's letters, not those of this chapter).

\emph{Lattices.} For every integer $m\ge0$ put $\varepsilon_m:=L^m\xi$. Thus $\varepsilon_0=\xi$
and $\varepsilon_\ell=1$, and $\varepsilon_m\le1$ for $m\le\ell$. Let $\Lambda_m$ be the set of
sites of spacing $\varepsilon_m$, so that $\Lambda_{m+1}\subset\Lambda_m$. A \emph{site} is a
point of $\Lambda_0$, and a \emph{fine bond} is a bond of $\Lambda_0$.

\emph{Blocks and roots.} For $y\in\Lambda_m$ let $B_m(y)$ be the standard $L$-adic
$m$-block rooted at $y$. The point $y$ is its \emph{root}, and distinct points of
$\Lambda_m$ are roots of distinct $m$-blocks. For a site $x$ let $c_m(x)$ be the root of the
$m$-block containing $x$.

For all integers $m,j,n\ge0$, all $y\in\Lambda_m$, all $y'\in\Lambda_{m+1}$ and all sites
$x$, we use the block axioms (B1)--(B4) below. Their consequences are (c1)--(c5):
\begin{equation}\label{4.26:eq-gauge-trees-a}
\begin{aligned}
&\text{(B1)}\quad\text{every site lies in exactly one $m$-block;}\\
&\text{(B2)}\quad B_{m+1}(y')=\textstyle\bigcup_{y\in B_{m+1}(y')\cap\Lambda_m}B_m(y),
  \ \text{a union of $L^d$ $m$-blocks;}\\
&\text{(B3)}\quad y\in B_m(y);\\
&\text{(B4)}\quad\text{if a fine bond joins sites of two different $m$-blocks,}\\
&\qquad\qquad\text{their roots are joined by a bond of $\Lambda_m$;}\\
&\text{(c1)}\quad c_0(x)=x;\\
&\text{(c2)}\quad c_j\circ c_m=c_j\quad\text{for }j\ge m;\\
&\text{(c3)}\quad z\in\Lambda_m,\ j\le m\ \Longrightarrow\ c_j(z)=z;\\
&\text{(c4)}\quad\text{the only point of $\Lambda_{m+1}$ in an $(m+1)$-block is its root;}\\
&\text{(c5)}\quad R\text{ a union of $n$-blocks},\ x\in R,\ m+1\le n\\
&\qquad\qquad\Longrightarrow\ B_{m+1}(c_{m+1}(x))\subset B_n(c_n(x))\subset R.
\end{aligned}
\end{equation}
Two models satisfy (B1)--(B4). The first is the corner-rooted model
$B_m(y)=y+\varepsilon_m[0,1)^d$. The second, for odd $L$, is the centre-rooted model
$B_m(y)=y+\varepsilon_m[-\tfrac12,\tfrac12)^d$. In \cite[near~(3.38)]{B-CMP109} the blocks are
rooted at their centres. Nothing in what follows depends on the choice of model, except
for the counts of sites in the averaging lemma and one factor $2$ in a corollary of that
lemma.

\emph{Gauge trees.} Let $B$ be an $(m+1)$-block with root $c$. A \emph{gauge tree} of $B$ is
a spanning tree of $B\cap\Lambda_m$ with the following two properties. First, its edges are
bonds of $\Lambda_m$ inside $B$. Second, for every $x\in B\cap\Lambda_m$, its tree path
$\Gamma^{\mathrm{ax}}_{c,x}$ from $c$ to $x$ has at most $d(L-1)$ bonds.

For every $(m+1)$-block $B$ a gauge tree $\mathcal{T}_m(B)$ is fixed. The model is the
\emph{coordinate tree}. It is the union, over $x\in B\cap\Lambda_m$, of the coordinate-ordered
monotone lattice paths from $c$ to $x$. In either model the points of $B\cap\Lambda_m$ differ
from $c$ by at most $(L-1)\varepsilon_m$ in each of the $d$ coordinates. Hence such a path has
at most $d(L-1)$ bonds.

For a site $x$ put
\begin{equation}\label{4.26:eq-gauge-trees-1}
\tau_m(x):=\Gamma^{\mathrm{ax}}_{c_{m+1}(x),\,c_m(x)} .
\end{equation}
This is the path in the gauge tree of $B_{m+1}(c_{m+1}(x))$ from $c_{m+1}(x)$ to $c_m(x)$. It
has at most $d(L-1)$ bonds of $\Lambda_m$, and all of them lie inside
$B_{m+1}(c_{m+1}(x))$. It is trivial if and only if $c_m(x)=c_{m+1}(x)$.

The block average uses its own comb contours $\Gamma_{y,x}$. In Ba{\l}aban's construction
the axial gauge-fixing condition and the average are both built on the contours
$\Gamma_{y,x}$, so there the two families of paths coincide. Nothing in what follows
requires this.

\emph{The forest of gauge-tree edges.} Let $n\ge0$ and let $R$ be a union of $n$-blocks. Put
\begin{equation}\label{4.26:eq-gauge-trees-2}
\mathcal{V}_{n,R}:=(\Lambda_0\cap R)\setminus\Lambda_n .
\end{equation}
Let $v\in\mathcal{V}_{n,R}$. Since $v\in\Lambda_0$, $v\notin\Lambda_n$ and the lattices
$\Lambda_m$ decrease in $m$, there is exactly one integer $m(v)$ with $0\le m(v)<n$ and
$v\in\Lambda_{m(v)}\setminus\Lambda_{m(v)+1}$. By (c3), $c_{m(v)}(v)=v$. Since
$c_{m(v)+1}(v)\in\Lambda_{m(v)+1}$ and $v\notin\Lambda_{m(v)+1}$, we have
$v\ne c_{m(v)+1}(v)$.

The point $v$ lies in $\Lambda_{m(v)}$ and in its own $(m(v)+1)$-block
$B_{m(v)+1}(c_{m(v)+1}(v))$. So the tree path
$\Gamma^{\mathrm{ax}}_{c_{m(v)+1}(v),v}$ in the gauge tree of that block is defined. It has at
least one bond, because its endpoints are distinct. Let $p(v)$ be the predecessor of $v$ on
this path, and let $e_v:=\{p(v),v\}$. This is an edge of the gauge tree.

These objects have the following two properties.
\begin{itemize}
\item[(F1)] The map $v\mapsto e_v$ is a bijection of $\mathcal{V}_{n,R}$ onto the set of all
tree edges of all gauge trees $\mathcal{T}_m(B')$, where $m<n$ and $B'\subset R$ is an
$(m+1)$-block. Hence the number of these tree edges is
\begin{equation}\label{4.26:eq-gauge-trees-3}
|\mathcal{V}_{n,R}|
=\sum_{m<n}\#\{(m+1)\text{-blocks}\subset R\}\cdot(L^d-1)
=|\Lambda_0\cap R|-|\Lambda_n\cap R| .
\end{equation}
If $R$ is a single $n$-block, this number is $L^{dn}-1$.
\item[(F2)] Iterating $p$ from $v\in\mathcal{V}_{n,R}$ terminates at $c_n(v)$. For
$x\in\Lambda_0\cap R$, the chain of edges from $c_n(x)$ down to $x$ is the concatenation
\begin{equation}\label{4.26:eq-gauge-trees-4}
\tau_{n-1}(x)\circ\cdots\circ\tau_0(x).
\end{equation}
Consequently $\mathcal{V}_{n,R}\cup(\Lambda_n\cap R)$, with roots $\Lambda_n\cap R$ and parent
map $p$, is a rooted forest.
\end{itemize}
\end{definition}

\begin{proof}[Proof of \textup{(c1)--(c5)} in \eqref{4.26:eq-gauge-trees-a} and of the properties of $\tau_m(x)$]
We first record a nesting property. Let $m\le j$ and let $x$ be a site. Then the $m$-block
containing $x$ is contained in the $j$-block containing $x$.

To see this, apply (B2) successively at the levels $m,m+1,\dots,j-1$. This writes the
$j$-block containing $x$ as a union of $m$-blocks. The site $x$ lies in one of these
$m$-blocks. By (B1) that block is the $m$-block containing $x$.

Next we observe which block is rooted at $c_m(x)$. By definition, $c_m(x)$ is the root of the
$m$-block containing $x$. So that block is $B_m(c_m(x))$. Since $c_m(x)\in\Lambda_m$, the
point $c_m(x)$ is itself a site.

(c3). Let $z\in\Lambda_m$ and $j\le m$. Iterating $\Lambda_{m'+1}\subset\Lambda_{m'}$ gives
$\Lambda_m\subset\Lambda_j$, so $z\in\Lambda_j$ and the block $B_j(z)$ is defined. By (B3)
it contains $z$. By (B1) it is therefore the $j$-block containing $z$. Its root is $z$, so
$c_j(z)=z$.

(c1). Every site lies in $\Lambda_0$. Hence (c1) is the case $j=m=0$ of (c3).

(c2). Let $j\ge m$. By (B3) the root $c_m(x)$ lies in $B_m(c_m(x))$, which is the $m$-block
containing $x$. By the nesting property, this block lies in the $j$-block containing $x$.
So $c_m(x)$ and $x$ lie in the same $j$-block. By (B1), $c_m(x)$ lies in no other $j$-block.
Hence $c_j(c_m(x))=c_j(x)$.

(c4). Let $B$ be an $(m+1)$-block with root $y$. By (B3), $y\in B$, and $y\in\Lambda_{m+1}$.
Now let $z\in\Lambda_{m+1}\cap B$. By (B1), $B$ is the $(m+1)$-block containing $z$, so
$c_{m+1}(z)=y$. By (c3) with $j=m+1$, $c_{m+1}(z)=z$. Hence $z=y$.

(c5). Let $R$ be a union of $n$-blocks, let $x\in R$ and let $m+1\le n$. The block
$B_{m+1}(c_{m+1}(x))$ is the $(m+1)$-block containing $x$. By the nesting property it lies in
the $n$-block containing $x$, which is $B_n(c_n(x))$. The site $x$ lies in one of the
$n$-blocks whose union is $R$. By (B1) that block is $B_n(c_n(x))$, so
$B_n(c_n(x))\subset R$.

Properties of $\tau_m(x)$. Put $B:=B_{m+1}(c_{m+1}(x))$; its root is $c_{m+1}(x)$. We check
that $c_m(x)\in B\cap\Lambda_m$. First, $c_m(x)\in\Lambda_m$. Second, by (B3),
$c_m(x)\in B_m(c_m(x))$, which is the $m$-block containing $x$. By the nesting property this
block lies in the $(m+1)$-block containing $x$, which is $B$. So $c_m(x)\in B\cap\Lambda_m$,
and the tree path in \eqref{4.26:eq-gauge-trees-1} is defined.

By the definition of a gauge tree, this path has at most $d(L-1)$ bonds. Its bonds are
edges of $\mathcal{T}_m(B)$, hence bonds of $\Lambda_m$ inside $B$.

It remains to check the triviality criterion. A tree path from a vertex to itself has no
bonds. A tree path between distinct vertices has at least one bond. So $\tau_m(x)$ is
trivial if and only if $c_m(x)=c_{m+1}(x)$.
\end{proof}

\begin{proof}[Proof of \textup{(F1)} and \textup{(F2)}]
We use two facts about a tree with a distinguished vertex $c$. First, if $w$ lies on the
tree path from $c$ to $v$, then the tree path from $c$ to $w$ is the initial segment of that
path ending at $w$. This holds because tree paths are unique. Second, let $e=\{a,b\}$ be an
edge. Removing $e$ splits the tree into two components, one of which contains $c$. Exactly
one endpoint of $e$ lies in the other component. We call it the endpoint of $e$
\emph{farther from the root}. The tree path from $c$ to this endpoint ends with $e$. The tree
path from $c$ to the other endpoint does not contain $e$.

(F1), the map lands in the stated set. Let $v\in\mathcal{V}_{n,R}$, put $m:=m(v)$ and
$B':=B_{m+1}(c_{m+1}(v))$. Then $m<n$, and $B'\subset R$ by (c5), since $m+1\le n$. The edge
$e_v$ joins consecutive vertices of a tree path in $\mathcal{T}_m(B')$. Hence it is an edge
of $\mathcal{T}_m(B')$.

(F1), the trees have disjoint edge sets. An edge of $\mathcal{T}_m(B')$ is a bond of
$\Lambda_m$, so it has length $\varepsilon_m$. These lengths are distinct for distinct $m$.
Hence a tree edge determines its level $m$. Both endpoints of the edge lie in $B'$. By (B1)
the $(m+1)$-blocks are disjoint, so the edge also determines $B'$.

(F1), surjectivity. Let $e$ be an edge of $\mathcal{T}_m(B')$, where $m<n$ and $B'\subset R$ is
an $(m+1)$-block with root $c$. Let $w$ be the endpoint of $e$ farther from the root. Then
$w\ne c$ and $w\in B'\cap\Lambda_m$. By (c4), $c$ is the only point of $\Lambda_{m+1}$ in
$B'$, so
\begin{equation*}
w\in(B'\cap\Lambda_m)\setminus\{c\}=(B'\cap\Lambda_m)\setminus\Lambda_{m+1}.
\end{equation*}
We have $w\in\Lambda_0\cap R$. Since $n\ge m+1$ we have $\Lambda_n\subset\Lambda_{m+1}$, so
$w\notin\Lambda_n$. Thus $w\in\mathcal{V}_{n,R}$ and $m(w)=m$. By (B1), $B'$ is the
$(m+1)$-block containing $w$, so $c_{m+1}(w)=c$. The tree path
$\Gamma^{\mathrm{ax}}_{c,w}$ ends with $e$. Hence $p(w)$ is the other endpoint of $e$, and
$e_w=e$.

(F1), injectivity. Suppose $e_v=e_{v'}=e$. By the disjointness just shown, $e$ determines
$m$ and $B'$, and therefore the tree containing it. In that tree, $v$ is an endpoint of $e$,
and the tree path from the root to $v$ ends with $e$. So $v$ is the endpoint of $e$ farther
from the root. The same holds for $v'$, and therefore $v=v'$.

(F1), the count. Let $B'$ be an $(m+1)$-block. By (B2), $B'$ is the union of the $m$-blocks
$B_m(y)$ with $y\in B'\cap\Lambda_m$. There are $L^d$ of these blocks, and distinct $y$ root
distinct blocks. Hence $|B'\cap\Lambda_m|=L^d$. A spanning tree on $L^d$ vertices has
$L^d-1$ edges. By the bijection and the disjointness of the edge sets, we obtain the first
equality in \eqref{4.26:eq-gauge-trees-3}. The second equality holds because
$\Lambda_n\cap R\subset\Lambda_0\cap R$.

Now let $R$ be a single $n$-block. Applying (B2) at levels $n-1,n-2,\dots,m+1$ writes $R$ as a
union of $L^{d(n-m-1)}$ distinct $(m+1)$-blocks. By (B1), every $(m+1)$-block contained in $R$
is one of them. The sum in \eqref{4.26:eq-gauge-trees-3} therefore telescopes:
\begin{equation*}
\sum_{m=0}^{n-1}L^{d(n-m-1)}(L^d-1)
=\sum_{m=0}^{n-1}\bigl(L^{d(n-m)}-L^{d(n-m-1)}\bigr)=L^{dn}-1 .
\end{equation*}

(F2), the chain. Let $x\in\Lambda_0\cap R$ and let $0\le m<n$. In the proof of the properties
of $\tau_m(x)$ above, we showed that $c_m(x)\in B_{m+1}(c_{m+1}(x))\cap\Lambda_m$. Let $w$ be a
vertex of $\tau_m(x)$ other than $c_{m+1}(x)$. Then $w\in B_{m+1}(c_{m+1}(x))\cap\Lambda_m$ and
$w$ is not the root. By (c4), $w\notin\Lambda_{m+1}$, and hence $w\notin\Lambda_n$. We also
have $w\in R$, by (c5). Thus $w\in\mathcal{V}_{n,R}$, with $m(w)=m$ and
$c_{m+1}(w)=c_{m+1}(x)$. By the first tree fact, $\Gamma^{\mathrm{ax}}_{c_{m+1}(x),w}$ is the
initial segment of $\tau_m(x)$ ending at $w$. Hence $p(w)$ is the vertex preceding $w$ on
$\tau_m(x)$. So iterating $p$ from $c_m(x)$ runs backwards along $\tau_m(x)$ and reaches
$c_{m+1}(x)$. If $c_m(x)=c_{m+1}(x)$, then $\tau_m(x)$ is trivial and there is nothing to run.

Start from $c_0(x)=x$, which holds by (c1), and apply this for $m=0,1,\dots,n-1$. The points
visited are the vertices of \eqref{4.26:eq-gauge-trees-4}, read from $x$ up to $c_n(x)$. Let
$u$ be such a point with $u\ne c_n(x)$. Then $u=c_j(x)$ or $u$ is an interior vertex of some
$\tau_j(x)$. If $u$ is an interior vertex of $\tau_j(x)$, put $j':=j$; then $\tau_{j'}(x)$ is
non-trivial and $u$ is a vertex of it other than $c_{j'+1}(x)$. If $u=c_j(x)$, let $j'$ be the
least index with $j'\ge j$ such that $\tau_{j'}(x)$ is non-trivial; such an index exists
because $u\ne c_n(x)$. Then $u=c_{j'}(x)$ is a vertex of $\tau_{j'}(x)$ other than
$c_{j'+1}(x)$. In either case $u$ is a vertex of a non-trivial $\tau_{j'}(x)$ other than its
endpoint $c_{j'+1}(x)$. By the preceding paragraph, $u\in\mathcal{V}_{n,R}$, and $p(u)$ is the
next point of the chain.

The last point is $c_n(x)$, which lies in $\Lambda_n$, so $p$ is not defined there. Thus the
chain of edges from $c_n(x)$ down to $x$ is \eqref{4.26:eq-gauge-trees-4}. Taking $x=v$ for
$v\in\mathcal{V}_{n,R}$ shows that iterating $p$ from $v$ terminates at $c_n(v)$.

(F2), the forest. By (F1), $p$ maps $\mathcal{V}_{n,R}$ into
$\mathcal{V}_{n,R}\cup(\Lambda_n\cap R)$. Indeed, $p(v)$ lies in $\Lambda_0\cap R$, and a point
of $\Lambda_0\cap R$ lies either in $\Lambda_n$ or in $\mathcal{V}_{n,R}$. Consider the graph on
$\mathcal{V}_{n,R}\cup(\Lambda_n\cap R)$ with edges $\{v,p(v)\}$ for $v\in\mathcal{V}_{n,R}$.
Orient each such edge from $v$ to $p(v)$. Then every vertex has at most one outgoing edge,
and a root has none.

Suppose the graph had a cycle with $q$ vertices. The cycle has $q$ edges, and each vertex
has at most one outgoing edge. Hence each vertex of the cycle has exactly one outgoing edge,
and that edge lies in the cycle. So no vertex of the cycle is a root, and $p$ maps the cycle
into itself. Iterating $p$ from a vertex of the cycle would then never reach a root. This
contradicts the termination just proved. The case $q=2$, that is $p(p(v))=v$, is excluded in
the same way. Hence the graph is a forest.

Every component contains a root: iterate $p$ from any of its vertices. Suppose a component
contained two roots $r\ne r'$, and let $r=u_0,u_1,\dots,u_q=r'$ be the tree path between
them. The root $u_0$ has no outgoing edge, so $p(u_1)=u_0$. Then $u_1$ has already used its
outgoing edge, so $p(u_2)=u_1$. Continuing in this way gives $p(u_q)=u_{q-1}$. This
contradicts the fact that $u_q=r'$ is a root. Hence each component contains exactly one root,
and the triple with roots $\Lambda_n\cap R$ and parent map $p$ is a rooted forest.
\end{proof}

\begin{definition}[Gauge action and path transport]\label{4.26:not-gauge-action}
Let $G=SU(N)$ and let $G^{c}=SL(N,\mathbb{C})$ be its complexification. For a bond $b$ write $b_-$ for its initial
and $b_+$ for its final site. For a function $h$ on sites with values in $G^{c}$ near $G$ (as in \cite{B-CMP109}),
a bond field $U$ and a
$\operatorname{Mat}_N(\mathbb{C})$-valued current $J$ on bonds, we put, as in \cite[(1.10), p.~262]{B-CMP109},
\begin{equation}\label{4.26:eq-gauge-action-1}
U^{h}(b)=h(b_-)\,U(b)\,h(b_+)^{-1},\qquad (R(h)J)(b)=\operatorname{Ad}_{h(b_-)}J(b),
\end{equation}
where $\operatorname{Ad}_{g}X=gXg^{-1}$ on $\operatorname{Mat}_N(\mathbb{C})$. For a path
$\gamma=(z_0,\dots,z_r)$ of nearest-neighbour sites, $U(\gamma)$ is the
ordered product $U(z_0,z_1)\cdots U(z_{r-1},z_r)$, where $U(z_{i-1},z_i)=U(b)$ if $b_-=z_{i-1}$, $b_+=z_i$, and
$U(z_{i-1},z_i)=U(b)^{-1}$ if the bond $b$ joining the two sites is traversed in reverse. Then, for $G^{c}$-valued
$a,c$ (with $(ca)(x)=c(x)a(x)$) and every path $\gamma$,
\begin{equation}\label{4.26:eq-gauge-action-a}
\begin{gathered}
(U^{a})^{c}=U^{ca},\qquad (U^{a})^{c}(b)=c(b_-)a(b_-)\,U(b)\,a(b_+)^{-1}c(b_+)^{-1},\\
U^{h}(\gamma)=h(z_0)\,U(\gamma)\,h(z_r)^{-1},
\end{gathered}
\end{equation}
and, for $G$-valued $h$, the map $U\mapsto U^{h}$ preserves the product Haar measure on $G$-valued bond fields.

\emph{Norms.} $\|\cdot\|$ is the operator norm on $\operatorname{Mat}_N(\mathbb{C})$. The matrix norm $|\cdot|$ of
\cite{B-CMP109} is read as a unitarily invariant norm (operator or Hilbert--Schmidt); since $\|X\|\le|X|$,
every upper bound on $|X|$ is an upper bound on $\|X\|$ with the same number. For $|\cdot|$ either of these two
norms and all $A,B,X\in\operatorname{Mat}_N(\mathbb{C})$,
\begin{equation}\label{4.26:eq-gauge-action-3}
|AXB|\le\|A\|\,\|B\|\,|X|.
\end{equation}
\end{definition}

\begin{proof}
\emph{Composition.} Applying \eqref{4.26:eq-gauge-action-1} with $h=c$ to $U^{a}$, and then with $h=a$ to $U$,
\begin{equation*}
(U^{a})^{c}(b)=c(b_-)\,U^{a}(b)\,c(b_+)^{-1}=c(b_-)a(b_-)\,U(b)\,a(b_+)^{-1}c(b_+)^{-1}.
\end{equation*}
Since $c(b_-)a(b_-)=(ca)(b_-)$ and $a(b_+)^{-1}c(b_+)^{-1}=(ca)(b_+)^{-1}$, the right side is $U^{ca}(b)$; this
is the first line of \eqref{4.26:eq-gauge-action-a}.

\emph{Path transport.} We show $U^{h}(z_{i-1},z_i)=h(z_{i-1})U(z_{i-1},z_i)h(z_i)^{-1}$ for each step. If the bond
$b$ is traversed forwards, $b_-=z_{i-1}$, $b_+=z_i$, and this is \eqref{4.26:eq-gauge-action-1}. If it is traversed
in reverse, $b_-=z_i$, $b_+=z_{i-1}$, and
\begin{equation*}
U^{h}(b)^{-1}=\bigl(h(z_i)U(b)h(z_{i-1})^{-1}\bigr)^{-1}=h(z_{i-1})\,U(b)^{-1}\,h(z_i)^{-1}.
\end{equation*}
Multiplying these identities for $i=1,\dots,r$ in order, each factor $h(z_i)^{-1}h(z_i)$ with $1\le i\le r-1$
cancels, leaving the second line of \eqref{4.26:eq-gauge-action-a}.

\emph{Haar measure.} For $G$-valued $h$, the map $U\mapsto U^{h}$ acts coordinatewise: the coordinate $U(b)$ is
sent to $h(b_-)U(b)h(b_+)^{-1}$, a left translation by $h(b_-)\in G$ followed by a right translation by
$h(b_+)^{-1}\in G$. Haar measure on the compact group $G$ is invariant under both, so each factor of the product
measure, and hence the product, is preserved.

\emph{Norms.} With $s_1(X)\ge s_2(X)\ge\dots$ the singular values of $X$, one has $\|X\|=s_1(X)$, while the
Hilbert--Schmidt norm is $(\sum_i s_i(X)^2)^{1/2}\ge s_1(X)$; so both choices of $|\cdot|$ dominate $\|\cdot\|$.
The singular values satisfy $s_i(AXB)\le\|A\|\,\|B\|\,s_i(X)$ for every $i$. For the operator norm this gives
$s_1(AXB)\le\|A\|\,\|B\|\,s_1(X)$, and for the Hilbert--Schmidt norm
\begin{equation*}
\Bigl(\sum_i s_i(AXB)^2\Bigr)^{1/2}\le\|A\|\,\|B\|\,\Bigl(\sum_i s_i(X)^2\Bigr)^{1/2};
\end{equation*}
this is \eqref{4.26:eq-gauge-action-3} for both choices of $|\cdot|$. In particular
$|\operatorname{Ad}_{g}X|\le\|g\|\,\|g^{-1}\|\,|X|$ in either norm, and no factor $\sqrt{N}$ enters the bounds
on $\operatorname{Ad}$.
\end{proof}

\begin{definition}[Axioms for the block average and its iterates]\label{4.26:def-block-average-axioms}
Let $\xi=L^{-\ell}$ be the spacing of the fine lattice on which the background field lives. For
$m\ge 0$ let $\varepsilon_m=L^m\xi$ (so that $\varepsilon_\ell=1$), let $\Lambda_m$ be the set of
sites of spacing $\varepsilon_m$, and let $B_m(y)$ be the $L$-adic $m$-block rooted at
$y\in\Lambda_m$, as in Notation~\ref{4.26:not-gauge-trees}. The gauge action $U\mapsto U^{h}$ on
bond variables is that of Notation~\ref{4.26:not-gauge-action}; $h|_{\Lambda_m}$ is the
restriction of a gauge transformation $h$ to $\Lambda_m$. We write $G=SU(N)$,
$G^{c}=SL(N,\mathbb{C})$ for its complexification, $\|\cdot\|$ for the operator norm on
$\operatorname{Mat}_N(\mathbb{C})$, $\mathfrak{sl}(N,\mathbb{C})$ for the traceless matrices, and
$\operatorname{Tr}=N\operatorname{tr}$ for the unnormalised matrix trace.

The conditions below are a reading of the averaging operation of
\cite[\S1, (11), p.~19]{B-CMP98}; they are formulated so that no particular normalisation
formula is among them. Fix a level $m$, put $\varepsilon:=\varepsilon_m$ and
$\varepsilon':=\varepsilon_{m+1}=L\varepsilon$, and let $M_{m+1}$ be the one-step block average,
which sends a $G^{c}$-valued bond field $W$ on $\Lambda_m$ to a bond field $M_{m+1}W$ on
$\Lambda_{m+1}$; when the level is clear from the context we write $M$ for it. For a
$G^{c}$-valued bond field $W$ and a path $\Gamma$ we write $W(\Gamma)$ for the ordered product
of the bond variables along $\Gamma$, a bond traversed forwards contributing $W(b)$ and a bond
traversed backwards contributing $W(b)^{-1}$. The block average and its iterates are subject
to the following conditions.
\begin{itemize}
\item[(M0)] The iterates are
\begin{equation}\label{4.26:eq-block-average-axioms-a}
M^{m}:=M_m\circ\cdots\circ M_1\quad(m\ge 1),\qquad M^{0}:=\mathrm{id},
\end{equation}
and each $M_{m}$ is given by the same recipe, on $L$-adic blocks, at every level.
\item[(M1)] For a coarse bond $c=\langle y,y'\rangle$ with $y\in\Lambda_{m+1}$ and
$y'=y+\varepsilon' e_\mu$,
\begin{equation}\label{4.26:eq-block-average-axioms-1}
M(W)(c):=\mathcal{N}\bigl((P_x)_{x\in B_{m+1}(y)\cap\Lambda_m}\bigr),\qquad
P_x:=W(\Gamma_{y,x})\,W([x,x'])\,W(\Gamma_{y',x'})^{-1},
\end{equation}
where $x':=x+\varepsilon' e_\mu$ and $[x,x']$ is the straight path of $L$ bonds of level $m$
from $x$ to $x'$. In consequence, for a $G^{c}$-valued bond field $U'$ on $\Lambda_0$ and a bond
$b=\langle b_-,b_+\rangle$ of $\Lambda_m$, the averaged variable $[M^{m}U'](b)$ depends on $U'$
only through its values on bonds with both endpoints in
\begin{equation*}
R_m(b):=B_m(b_-)\cup B_m(b_+).
\end{equation*}
\item[(M2)] $\Gamma_{y,x}$ is the coordinate-ordered monotone path in $\Lambda_m$, inside
$B_{m+1}(y)$, from $y$ to $x$. It is the same for all directions $\mu$, and it is translation
covariant: $\Gamma_{y+v,x+v}=\Gamma_{y,x}+v$ for every lattice vector $v$ of $\Lambda_{m+1}$
(and the block offsets $x-y$ are the same for every block). Its length satisfies
$|\Gamma_{y,x}|\le d(L-1)$, and $n_{\perp}(x;\mu)$ denotes the number of its bonds not parallel
to $e_\mu$.
\item[(M3)] All $L^d$ sites $x$ of the block enter symmetrically, with the uniform weight
$L^{-d}$ in the linearisation, and in the same way for every $\mu$.
\item[(M4)] There are numbers $r_{\mathcal{N}}>0$, $C_{\mathcal{N}}$ and $c_{\Pi}$ such that the
normalisation $\mathcal{N}$ is a holomorphic map from the polydisc
\begin{equation*}
\bigl\{(P_x)_x\in\operatorname{Mat}_N(\mathbb{C})^{L^d}:\ \|P_x-1\|<r_{\mathcal{N}}
\text{ for all }x\bigr\}
\end{equation*}
to $\operatorname{Mat}_N(\mathbb{C})$, which takes values in $G^{c}$ at families of elements of
$G^{c}$ and satisfies $\mathcal{N}(g,\dots,g)=g$ for every $g\in G^{c}$ with $(g,\dots,g)$ in the
polydisc. Its differential at the identity family $\mathbf{1}=(1,\dots,1)$ is the arithmetic
mean on traceless families:
\begin{equation}\label{4.26:eq-block-average-axioms-2}
\Pi[(H_x)_x]:=D\mathcal{N}(\mathbf{1})[(H_x)_x]=L^{-d}\sum_x H_x
\quad\text{whenever every }H_x\in\mathfrak{sl}(N,\mathbb{C}).
\end{equation}
Families are normed by $\|H\|:=\max_x\|H_x\|$, and
\begin{equation*}
\|\Pi\|\le c_{\Pi},\qquad
\|D^2\mathcal{N}(P)[H,K]\|\le 2C_{\mathcal{N}}\|H\|\,\|K\|\quad\text{on the polydisc.}
\end{equation*}
\item[(M5)] Covariance: for $h,h'\in G^{c}$ near $1$,
\begin{equation}\label{4.26:eq-block-average-axioms-3}
\mathcal{N}\bigl((hP_xh'^{-1})_x\bigr)=h\,\mathcal{N}\bigl((P_x)_x\bigr)\,h'^{-1}.
\end{equation}
(For holomorphic $\mathcal{N}$ this follows from the case $h,h'\in G$ by the identity theorem.)
In consequence $M(W^{h})=(MW)^{h|_{\Lambda_{m+1}}}$, and, for the iterates,
\begin{equation}\label{4.26:eq-block-average-axioms-b}
M^{m}(U^{h})=(M^{m}U)^{h|_{\Lambda_m}} .
\end{equation}
\end{itemize}
\end{definition}

\begin{remark}[The contour, admissible normalisations, and the letter $M$]
\emph{The contour.} In \cite[\S1]{B-CMP98} the average is built on comb contours
$\Gamma_{y,x}\cup[x,x']\cup\Gamma_{y',x'}^{-1}$ of the kind described in (M1)--(M2). The contour
$\Gamma_{y,x}$ may, but need not, coincide with the path of the gauge tree of
Notation~\ref{4.26:not-gauge-trees}. That a single contour serves all directions is dictated by
the axial gauge: the axial gauge-fixing $\delta$-functions $\delta(U(\Gamma_{y,x}))$ of
Ba{\l}aban's renormalisation transformation carry no direction index; in that gauge the average
is the plain straight-line mean, and extending it by covariance reproduces the comb with the
gauge's own tree, which does not depend on $\mu$.

\emph{Admissible normalisations.} The mean property \eqref{4.26:eq-block-average-axioms-2} holds
for every $\mathcal{N}$ that is symmetric under permutations of its arguments and satisfies
$\mathcal{N}(g,\dots,g)=g$; this is shown in the proof below. Each of the following
normalisations is holomorphic near $\mathbf{1}$, covariant as in (M5), and $G^{c}$-valued on
families of elements of $G^{c}$, with $(r_{\mathcal{N}},C_{\mathcal{N}},c_{\Pi})$ absolute
numbers; this rests on the operator-norm calculus of $X\mapsto X^{-1/2}$ and of
\begin{equation*}
\det(Y)^{-1/N}=\exp\bigl(-N^{-1}\operatorname{Tr}\operatorname{Log}Y\bigr),\qquad
\bigl|N^{-1}\operatorname{Tr}\operatorname{Log}Y\bigr|\le\|\operatorname{Log}Y\|,
\end{equation*}
where $\operatorname{Log}$ is the principal matrix logarithm.
\begin{itemize}
\item[(a)] $\mathcal{N}=\mathcal{N}_0\bigl(L^{-d}\sum_xP_x\bigr)$ with
$\mathcal{N}_0(X):=X(\det X)^{-1/N}$.
\item[(b)] The holomorphic polar form $\mathcal{N}=Y(\det Y)^{-1/N}$, where
\begin{equation*}
Y:=X(\tilde X X)^{-1/2},\qquad X:=L^{-d}\sum_xP_x,\qquad \tilde X:=L^{-d}\sum_xP_x^{-1}.
\end{equation*}
At real configurations, that is, for $P_x\in G$, one has $\tilde X=X^{*}$, and $Y$ is the
unitary polar part of $X$. For $N=2$ the adjugate $\operatorname{adj}$ is linear on
$2\times2$ matrices and equals the inverse on $SL(2,\mathbb{C})$; hence on families of elements
of $SL(2,\mathbb{C})$ one has $\tilde X=\operatorname{adj}X$, so $\tilde XX=\det X\cdot 1$
and $\mathcal{N}=X/\sqrt{\det X}$, the normalisation used for $SU(2)$ in \cite[\S1]{B-CMP98}.
No identity special to $N=2$ is used in this section.
\item[(c)] Lie-algebra means: $\exp\bigl(L^{-d}\sum_x\operatorname{Log}P_x\bigr)$; the implicit
mean $W$ defined by $\sum_x\operatorname{Log}(P_xW^{-1})=0$; and their variants with a reference
transporter (the last are not symmetric, but have the same linearisation).
\end{itemize}
Which normalisation is used in \cite[\S1]{B-CMP98} plays no role here: of the normalisation,
only the properties listed in (M4) and (M5) are used.

\emph{Two meanings of the letter $M$.} The letter $M$ applied to a configuration (as in
$MW$, $M^{m}U'$, $[M^{m}U'](b)$) always denotes the block average of
Definition~\ref{4.26:def-block-average-axioms} or one of its iterates. A bare $M$ multiplying a
number, as in the cube side of Ba{\l}aban's auxiliary parameters or in the phrase
``$M$ sufficiently large'', is the cube-size integer $M=L^{m'}$ of the glossary. The estimates
for iterated block averages proved in this section are uniform in that integer, and the two
meanings are never interchanged.
\end{remark}

\begin{proof}[Proof of the consequences stated in (M1), (M4) and (M5)]
\emph{Locality in (M1).} We first treat one step. Let $c=\langle y,y'\rangle$ be a bond of
$\Lambda_{m+1}$ with $y'=y+\varepsilon' e_\mu$, and let $x\in B_{m+1}(y)\cap\Lambda_m$. By (M2),
$\Gamma_{y,x}$ lies inside $B_{m+1}(y)$. Since $\varepsilon' e_\mu$ is a lattice vector of
$\Lambda_{m+1}$ and the block offsets are the same for every block, $x'=x+\varepsilon'e_\mu$
lies in $B_{m+1}(y')$ and $\Gamma_{y',x'}=\Gamma_{y,x}+\varepsilon'e_\mu$ lies inside
$B_{m+1}(y')$. The straight path $[x,x']$ has the sites $x+j\varepsilon e_\mu$,
$0\le j\le L$, which differ from $x$ only in the $\mu$-th coordinate. The $L$-adic blocks
$B_{m+1}(y)$ and $B_{m+1}(y')=B_{m+1}(y)+\varepsilon'e_\mu$ are cubes of $L^d$ sites of
$\Lambda_m$; in the $\mu$-th coordinate the sites of the first occupy $L$ consecutive multiples
of $\varepsilon$ and those of the second the next $L$. Since the $\mu$-th coordinate of $x$ is
among the first $L$ and that of $x'$ among the second $L$, every site of $[x,x']$ lies in
$B_{m+1}(y)\cup B_{m+1}(y')$. Hence, by \eqref{4.26:eq-block-average-axioms-1}, $M(W)(c)$
depends on $W$ only through its values on bonds of $\Lambda_m$ with both endpoints in
$B_{m+1}(y)\cup B_{m+1}(y')$.

We now prove by induction on $m\ge 0$ that $[M^{m}U'](b)$ depends on $U'$ only through its values
on bonds with both endpoints in $R_m(b)$. For $m=0$, $M^{0}=\mathrm{id}$ and
$R_0(b)=\{b_-,b_+\}$, so the claim holds. Let $m\ge1$ and assume the claim at $m-1$. By
\eqref{4.26:eq-block-average-axioms-a}, $[M^{m}U'](b)=[M_m(M^{m-1}U')](b)$, and by the one-step
statement just proved (applied at level $m-1$) this depends only on the values
$[M^{m-1}U'](b'')$ at bonds $b''$ of $\Lambda_{m-1}$ with both endpoints in
$B_m(b_-)\cup B_m(b_+)=R_m(b)$. If $z$ is an endpoint of such a bond $b''$, it lies in
$B_m(b_-)$ or in $B_m(b_+)$; since $L$-adic blocks are nested, the $(m-1)$-block $B_{m-1}(z)$ is
contained in that $m$-block. Hence $R_{m-1}(b'')\subset R_m(b)$, and by the induction hypothesis
$[M^{m-1}U'](b'')$ depends on $U'$ only through its values on bonds with both endpoints in
$R_{m-1}(b'')\subset R_m(b)$. This completes the induction.

\emph{The mean property in (M4).} Let $\mathcal{N}$ be symmetric under permutations of its
arguments and satisfy $\mathcal{N}(g,\dots,g)=g$ for $g\in G^{c}$ in the polydisc. The
differential $\Pi=D\mathcal{N}(\mathbf{1})$ is linear, so $\Pi[(H_x)_x]=\sum_x\Phi_x(H_x)$, where
$\Phi_x(H)$ is the value of $\Pi$ on the family with $H$ at the site $x$ and $0$ elsewhere. For a
permutation $\sigma$ of the sites of the block, $\mathcal{N}((P_{\sigma(x)})_x)=\mathcal{N}((P_x)_x)$;
the identity family is fixed by $\sigma$, so differentiating at $\mathbf{1}$ gives
$\Pi[(H_{\sigma(x)})_x]=\Pi[(H_x)_x]$, and therefore all $\Phi_x$ coincide with one linear map
$\Phi$. Let $H\in\mathfrak{sl}(N,\mathbb{C})$. For real $t$ near $0$, $\|e^{tH}-1\|<r_{\mathcal{N}}$
and $\det e^{tH}=e^{t\operatorname{Tr}H}=1$, so $e^{tH}\in G^{c}$ and
$\mathcal{N}(e^{tH},\dots,e^{tH})=e^{tH}$. Differentiating at $t=0$ gives
$\Pi[(H,\dots,H)]=H$, that is, $L^d\Phi(H)=H$. Consequently, if every $H_x$ is traceless,
\begin{equation*}
\Pi[(H_x)_x]=\sum_x\Phi(H_x)=L^{-d}\sum_xH_x,
\end{equation*}
which is \eqref{4.26:eq-block-average-axioms-2}.

\emph{Covariance in (M5).} Under the gauge action of Notation~\ref{4.26:not-gauge-action}, the
ordered product along a path $\Gamma$ from $u$ to $v$ transforms as
$W^{h}(\Gamma)=h(u)\,W(\Gamma)\,h(v)^{-1}$. Applying this to the three factors of $P_x$ in
\eqref{4.26:eq-block-average-axioms-1}, the paths $\Gamma_{y,x}$, $[x,x']$ and
$\Gamma_{y',x'}^{-1}$ running from $y$ to $x$, from $x$ to $x'$ and from $x'$ to $y'$, the
interior factors $h(x)^{-1}h(x)$ and $h(x')^{-1}h(x')$ cancel along the comb, and
\begin{equation*}
P_x(W^{h})=h(y)\,P_x(W)\,h(y')^{-1}\qquad\text{for every }x\in B_{m+1}(y)\cap\Lambda_m .
\end{equation*}
For gauge transformations $h$ whose values $h(y),h(y')$ are as in (M5) and for which the families
involved lie in the polydisc of (M4), \eqref{4.26:eq-block-average-axioms-3} with $h=h(y)$,
$h'=h(y')$ gives
\begin{equation*}
M(W^{h})(c)=h(y)\,M(W)(c)\,h(y')^{-1}=(MW)^{h|_{\Lambda_{m+1}}}(c),
\end{equation*}
since $y,y'\in\Lambda_{m+1}$; thus $M_m(W^{h})=(M_mW)^{h|_{\Lambda_m}}$ at every level $m$.
Finally, \eqref{4.26:eq-block-average-axioms-b} follows by induction on $m$: for $m=0$ both
sides equal $U^{h}$, and if it holds at $m-1$, then, using
\eqref{4.26:eq-block-average-axioms-a}, the one-step covariance at level $m$ and
$\Lambda_m\subset\Lambda_{m-1}$,
\begin{align*}
M^{m}(U^{h})&=M_m\bigl(M^{m-1}(U^{h})\bigr)=M_m\bigl((M^{m-1}U)^{h|_{\Lambda_{m-1}}}\bigr)\\
&=\bigl(M_mM^{m-1}U\bigr)^{h|_{\Lambda_m}}=(M^{m}U)^{h|_{\Lambda_m}}. \qedhere
\end{align*}
\end{proof}

\begin{definition}[Hierarchically block-axial fields]\label{4.26:def-block-axial}
Let $n\ge 0$ be an integer, and let $R$ be a union of $n$-blocks. Blocks, their roots and the gauge trees
$\mathcal{T}_m(B')$ of $(m+1)$-blocks $B'$ on $\Lambda_m$ are those of
Notation~\ref{4.26:not-gauge-trees}. The $m$-fold iterate $M^{m}$ of the one-step block average $M$
is as in Definition~\ref{4.26:def-block-average-axioms}, and $[M^{m}U](b)$ denotes the averaged
bond variable on a bond $b$ of $\Lambda_m$.

A bond field $U$ on the fine lattice of spacing $\xi$ is \emph{hierarchically block-axial to
depth $n$ on $R$}, and we write $U\in\mathrm{HBA}(n,R)$, if for every $m<n$ and every
$(m+1)$-block $B'\subset R$
\begin{equation}\label{4.26:eq-block-axial-a}
  [M^{m}U](b)=1
  \qquad\text{for every bond $b$ of the gauge tree $\mathcal{T}_m(B')$ in $\Lambda_m$.}
\end{equation}

The same condition may be written on the forest of Notation~\ref{4.26:not-gauge-trees}. Let
$\mathcal{V}_{n,R}$ be the vertex set of that forest for depth $n$ on $R$, and for
$v\in\mathcal{V}_{n,R}$ let the forest edge $e_v$, its parent endpoint $p(v)$ and its level
$m(v)$ be as in that notation. We put
\begin{equation*}
  a_v(U):=[M^{m(v)}U]\bigl(p(v)\to v\bigr),
\end{equation*}
the averaged variable on $e_v$ read from $p(v)$ to $v$. Then
$U\in\mathrm{HBA}(n,R)$ if and only if $a_v(U)=1$ for all $v\in\mathcal{V}_{n,R}$.

More generally, for a structure with blocks of varying level, such as the maximal sequence of
domains of \cite[(1.3)--(1.6)]{B-CMP99a}, condition~\eqref{4.26:eq-block-axial-a} is required,
inside each structure block of level $\lambda$, for every $m<\lambda$ and every $(m+1)$-block
$B'$ contained in that structure block.

Of the gauge trees, nothing is used here beyond their definition in
Notation~\ref{4.26:not-gauge-trees}.
\end{definition}

\begin{lemma}[The relative gauge element as a tree product]\label{4.26:lem-tree-product}
Let $R$ be a union of $n$-blocks. Let $U'$ be a $G^{c}$-valued bond field on the fine bonds of
$R$, lying in the domain of the iterated averages $M^{m}$ for $m<n$. Let $u$ be a $G^{c}$-valued
site field. Suppose that $U'^{u}$ is hierarchically block-axial to depth $n$ on $R$, that is,
$U'^{u}\in\mathrm{HBA}(n,R)$ in the sense of Definition~\ref{4.26:def-block-axial}. Then for
every fine site $x\in R$
\begin{equation}\label{4.26:eq-tree-product-a}
\begin{gathered}
u(x)=u(c_n(x))\,P_n(x),\\
P_n(x):=T_{n-1}(x)\,T_{n-2}(x)\cdots T_0(x),\qquad T_m(x):=[M^{m}U'](\tau_m(x)).
\end{gathered}
\end{equation}
Here $[M^{m}U'](\tau_m(x))$ is the ordered product of the factors $[M^{m}U'](b)^{\pm1}$ along the
path $\tau_m(x)$ from $c_{m+1}(x)$ to $c_m(x)$. The exponent is $+1$ for a bond traversed from
$b_-$ to $b_+$ and $-1$ otherwise.

In particular $u(c_n(x))^{-1}u(x)=P_n(x)$ depends on $U'$ alone. For $z\in\Lambda_m\cap R$ we
have $P_n(z)=T_{n-1}(z)\cdots T_m(z)$, and in particular $P_n\equiv 1$ on $\Lambda_n$.

For $m<n$ call the bonds of the gauge trees $\mathcal{T}_m(B')$ of the $(m+1)$-blocks
$B'\subset R$ the tree bonds of level $m$, and put
\begin{equation*}
\delta_m:=\sup_{b}\max\bigl\{\|[M^{m}U'](b)-1\|,\ \|[M^{m}U'](b)^{-1}-1\|\bigr\},
\end{equation*}
the supremum being taken over the tree bonds $b$ of level $m$. Then
\begin{equation}\label{4.26:eq-tree-product-b}
\|P_n(x)-1\|,\ \ \|P_n(x)^{-1}-1\|
\;\le\;\exp\Bigl(d(L-1)\sum_{m<n}\delta_m\Bigr)-1 .
\end{equation}
\end{lemma}

\begin{proof}
\emph{The relation along one tree bond.} Let $m<n$. The covariance of the iterated average,
\eqref{4.26:eq-block-average-axioms-b}, gives
\begin{equation*}
[M^{m}(U'^{u})](b)=u(b_-)\,[M^{m}U'](b)\,u(b_+)^{-1}
\end{equation*}
for every bond $b$ of $\Lambda_m$. Here $U'^{u}$ is the gauge action of
Notation~\ref{4.26:not-gauge-action}.

Let $b=\langle z,z'\rangle$ be a bond of a gauge tree $\mathcal{T}_m(B')$ of an $(m+1)$-block
$B'\subset R$, with $b_-=z$ and $b_+=z'$. By hypothesis $U'^{u}$ satisfies the tree condition
\eqref{4.26:eq-block-axial-a} at $b$. By the covariance identity above, this condition reads
\begin{equation*}
1=u(z)\,[M^{m}U'](b)\,u(z')^{-1}.
\end{equation*}
Equivalently, $u(z')=u(z)\,[M^{m}U'](b)$. Read from the other end, it says
$u(z)=u(z')\,[M^{m}U'](b)^{-1}$.

So, whichever way a tree bond is traversed, the value of $u$ at the far end equals the value of
$u$ at the near end multiplied on the right by the factor $[M^{m}U'](b)^{\pm1}$. The exponent
follows the rule stated in the lemma.

\emph{One level.} Fix a fine site $x\in R$ and $m<n$. The path $\tau_m(x)$ is the gauge-tree path
of the block $B_{m+1}(c_{m+1}(x))$ from its root to $c_m(x)$. By property (c5) in
\eqref{4.26:eq-gauge-trees-a}, this block is contained in $R$. Since $m<n$, the tree
condition holds on all bonds of $\mathcal{T}_m(B_{m+1}(c_{m+1}(x)))$.

Write the path as
\begin{equation*}
\tau_m(x)=(z_0=c_{m+1}(x),\,z_1,\,\dots,\,z_r=c_m(x)),\qquad r\le d(L-1),
\end{equation*}
the bound on $r$ being the length bound for gauge-tree paths of
Notation~\ref{4.26:not-gauge-trees}, and let $b_i$ be the bond joining $z_{i-1}$ and $z_i$. By the
relation along one tree bond,
\begin{equation*}
u(z_i)=u(z_{i-1})\,[M^{m}U'](b_i)^{\pm1},\qquad i=1,\dots,r .
\end{equation*}
Chaining these $r$ relations gives $u(c_m(x))=u(c_{m+1}(x))\,T_m(x)$, by the definition of
$T_m(x)$ as the ordered product along $\tau_m(x)$.

\emph{All levels.} We apply the one-level relation successively for
$m=0,1,\dots,n-1$, using $c_0(x)=x$. This gives
\begin{align*}
u(x)&=u(c_1(x))\,T_0(x)=u(c_2(x))\,T_1(x)\,T_0(x)=\cdots\\
&=u(c_n(x))\,T_{n-1}(x)\cdots T_0(x)=u(c_n(x))\,P_n(x),
\end{align*}
which is \eqref{4.26:eq-tree-product-a}.

Multiplying on the left by $u(c_n(x))^{-1}$ gives $u(c_n(x))^{-1}u(x)=P_n(x)$. The right-hand
side is built from the averaged variables $[M^{m}U'](b)$ alone, so it depends on $U'$ alone.

For $z\in\Lambda_m\cap R$, property (c3) in \eqref{4.26:eq-gauge-trees-a} says that the
paths $\tau_j(z)$ with $j<m$ are trivial. The corresponding factors are therefore $T_j(z)=1$,
and $P_n(z)=T_{n-1}(z)\cdots T_m(z)$. For $z\in\Lambda_n$ no factor remains, so
$P_n(z)=1$.

\emph{The bound.} Let $A,B\in\operatorname{Mat}_N(\mathbb{C})$. From the identity
\begin{equation*}
AB-1=(A-1)(B-1)+(A-1)+(B-1)
\end{equation*}
and the submultiplicativity of the operator norm we obtain
\begin{equation*}
\|AB-1\|\le(1+\|A-1\|)(1+\|B-1\|)-1 .
\end{equation*}
Induction on the number of factors then gives, for $A_1,\dots,A_p\in\operatorname{Mat}_N(\mathbb{C})$,
\begin{equation*}
\Bigl\|\prod_{i=1}^{p}A_i-1\Bigr\|\le\prod_{i=1}^{p}\bigl(1+\|A_i-1\|\bigr)-1 .
\end{equation*}
In the induction step we apply the two-factor inequality with $A$ the product of the first $p-1$
factors, and use the induction hypothesis in the form
$1+\|A-1\|\le\prod_{i<p}(1+\|A_i-1\|)$.

We apply this to $P_n(x)$, written as the product of all the factors of
$T_{n-1}(x),\dots,T_0(x)$. By the description of $\tau_m(x)$ above, each $T_m(x)$ has at most
$d(L-1)$ factors. Each of them is $[M^{m}U'](b)^{\pm1}$ for a tree bond $b$ of level $m$, so it is
within $\delta_m$ of $1$ in either orientation. Hence
\begin{equation*}
\|P_n(x)-1\|\le\prod_{m<n}(1+\delta_m)^{d(L-1)}-1 .
\end{equation*}

The inverse $P_n(x)^{-1}=T_0(x)^{-1}\cdots T_{n-1}(x)^{-1}$ is the reversed product of the
inverses of the same factors. Each inverse factor is $[M^{m}U'](b)^{\mp1}$, again within
$\delta_m$ of $1$, because $\delta_m$ bounds both orientations. The same estimate
therefore holds for $\|P_n(x)^{-1}-1\|$.

Finally, $1+a\le e^{a}$ for $a\ge0$ gives
\begin{equation*}
\prod_{m<n}(1+\delta_m)^{d(L-1)}\le\exp\Bigl(d(L-1)\sum_{m<n}\delta_m\Bigr),
\end{equation*}
which proves \eqref{4.26:eq-tree-product-b}.
\end{proof}

\begin{remark}[Independence of the action convention]
Suppose the gauge action is taken in the mirrored form $U^{h}(b)=h(b_-)^{-1}U(b)h(b_+)$. Then
the argument above, with the roles of the two ends exchanged, gives $u(x)=P'_n(x)\,u(c_n(x))$.
Here $P'_n(x)$ is again an ordered product of averaged tree-bond variables $[M^{m}U'](b)^{\pm1}$.

For the mirrored action the composition rule is $(U^{a})^{h}=U^{ah}$, and the relative element
of $u$ with respect to its root values is $u(x)\,u(c_n(x))^{-1}=P'_n(x)$ in place of
$u(c_n(x))^{-1}u(x)=P_n(x)$. It is again a tree product without conjugation, equal to $1$ at the
roots, that is, on $\Lambda_n\cap R$.

The element that acts in \cite[(3.38)]{B-CMP109} is the one fixed by the chain of equalities
stated there; under either convention it is a conjugation-free tree product equal to the identity
at roots.
\end{remark}

\begin{definition}[Combs, the linearised average and the lattice curl]\label{4.26:not-combs-linearised-average}
We use the blocks, roots and gauge trees of Notation~\ref{4.26:not-gauge-trees} and the axioms
for the block average of Definition~\ref{4.26:def-block-average-axioms}, in particular the
comb transporter $P_x$, the contours $\Gamma_{y,x}$ and the locality regions
$R_m(b)=B_m(b_-)\cup B_m(b_+)$ of \eqref{4.26:eq-block-average-axioms-a}.
Fix a level $m$, and let $A$ be an $\mathfrak{sl}(N,\mathbb{C})$-valued function on the
level-$m$ bonds. We put $W=\exp(i\varepsilon_m A)$, and for a site $z$ and a direction $\mu$
we write $A_\mu(z)$ for the value of $A$ on the bond $\langle z,z+\varepsilon_m e_\mu\rangle$.

Let $c=\langle y,y'\rangle$ be a level-$(m+1)$ bond, $y\in\Lambda_{m+1}$,
$y'=y+\varepsilon_{m+1}e_\mu$, and for $x\in B_{m+1}(y)\cap\Lambda_m$ put
$x'=x+\varepsilon_{m+1}e_\mu$. Write the contour as
$\Gamma_{y,x}=(z_0=y,\dots,z_r=x)$ with $z_i=z_{i-1}+s_i\varepsilon_m e_{\nu_i}$, $s_i=\pm1$,
and let $\beta_i=\langle w_i,w_i+\varepsilon_m e_{\nu_i}\rangle$ be its $i$-th bond, whose
lower endpoint $w_i$ is $z_{i-1}$ if $s_i=+1$ and $z_i$ if $s_i=-1$. The \emph{comb} of $x$
is the path from $y$ to $y'$
\begin{equation*}
\mathrm{comb}_x=\Gamma_{y,x}\cdot[x,x']\cdot(\Gamma_{y',x'})^{-1},
\end{equation*}
where $[x,x']$ is the straight path of $L$ level-$m$ bonds from $x$ to $x'$. For a function
$F$ on level-$m$ bonds we put
\begin{equation*}
\langle F,\mathrm{comb}_x\rangle=\sum_{\beta\in\mathrm{comb}_x}\sigma_\beta F(\beta),
\end{equation*}
where $\sigma_\beta=\pm1$ is the orientation in which $\mathrm{comb}_x$ traverses $\beta$;
$\langle F,\mathrm{comb}_x\rangle$ is traceless if $F$ is. The \emph{linearised comb average}
$\mathcal{Q}A$, the \emph{straight-line average} $SA$ and the
\emph{lattice curl} $d_mA$ at spacing $\varepsilon_m$ are
\begin{equation}\label{4.26:eq-combs-linearised-average-a}
\begin{gathered}
\mathcal{Q}A(c)=L^{-1-d}\sum_{x\in B_{m+1}(y)\cap\Lambda_m}\langle A,\mathrm{comb}_x\rangle,\\
(SA)_\mu(y)=L^{-1}\sum_{u=0}^{L-1}A_\mu(y+u\varepsilon_m e_\mu),\\
(d_mA)_{\mu\nu}(z)=\varepsilon_m^{-1}\bigl[A_\nu(z+\varepsilon_m e_\mu)-A_\nu(z)
-A_\mu(z+\varepsilon_m e_\nu)+A_\mu(z)\bigr].
\end{gathered}
\end{equation}

\emph{Domains.} Let $\Omega$ be a box of the fine lattice. We put
\begin{equation*}
\Omega_m=\{\,b \text{ a level-}m\text{ bond}: R_m(b)\subset\Omega\,\},
\end{equation*}
and we say that a level-$m$ plaquette is in $\Omega_m$ if its four bonds are. If
$c\in\Omega_{m+1}$ and $b$ is a level-$m$ bond in $R_{m+1}(c)$, then
$R_m(b)\subset R_{m+1}(c)\subset\Omega$, so that $b\in\Omega_m$.
\end{definition}

\begin{proof}
Tracelessness. $\langle F,\mathrm{comb}_x\rangle$ is a finite sum of the matrices
$\pm F(\beta)$, and the trace is linear; if every $F(\beta)$ has trace zero, so has the sum.

Nesting. The $L$-adic blocks are nested (Notation~\ref{4.26:not-gauge-trees}): every
$(m+1)$-block is the union of the $L^d$ $m$-blocks rooted at the sites of $\Lambda_m$ that it
contains. Let $b$ be a level-$m$ bond whose endpoints $b_-,b_+$ lie in
$R_{m+1}(c)=B_{m+1}(c_-)\cup B_{m+1}(c_+)$. The endpoint $b_-$ is a site of $\Lambda_m$ lying
in one of the blocks $B_{m+1}(c_\pm)$, hence the $m$-block $B_m(b_-)$ rooted at it is one of the
$m$-blocks of which that $(m+1)$-block is the union, and $B_m(b_-)\subset R_{m+1}(c)$. The same
holds for $b_+$, so
\begin{equation*}
R_m(b)=B_m(b_-)\cup B_m(b_+)\subset R_{m+1}(c).
\end{equation*}
Since $c\in\Omega_{m+1}$ means $R_{m+1}(c)\subset\Omega$, we get $R_m(b)\subset\Omega$, that
is, $b\in\Omega_m$.
\end{proof}

\begin{lemma}[Majorant bounds for products of exponentials]\label{4.26:lem-exponential-products}
Let $X_1,\dots,X_n$ and $\tilde X_1,\dots,\tilde X_n$ be elements of
$\operatorname{Mat}_N(\mathbb{C})$, with the operator norm $\|\cdot\|$, and write $1$ for the
identity matrix. For a word $w=(j_1,\dots,j_r)$ in the letters $1,\dots,n$ put
$X^w=X_{j_1}\cdots X_{j_r}$ and, for $x=(x_1,\dots,x_n)\in[0,\infty)^n$,
$x^w=x_{j_1}\cdots x_{j_r}$. Let
\begin{equation*}
\Phi(X)=\sum_w c_w X^w
\end{equation*}
be a power series in the non-commuting variables $X_1,\dots,X_n$, and let
$\varphi(x)=\sum_w a_w x^w$ be a majorant of it, that is, a series with nonnegative
coefficients $a_w\ge|c_w|$. Let $x_j\ge0$ be numbers with $\|X_j\|\le x_j$ and
$\|\tilde X_j\|\le x_j$ for every $j$, and assume $\varphi(x)<\infty$. Then $\Phi(X)$
converges absolutely,
\begin{equation*}
\|\Phi(X)\|\le\varphi(x),\qquad
\|\Phi(X)-\Phi(\tilde X)\|\le\sum_k\delta_k\,\partial_k\varphi(x),\qquad
\delta_k:=\|X_k-\tilde X_k\|,
\end{equation*}
where $\partial_k\varphi$ is the termwise derivative. Let $s$ be a number with
$\sum_j x_j\le s$ (so $s\ge0$), and put $\delta:=\sum_k\delta_k$. Products
$\prod_j e^{X_j}$ are taken in the fixed order $j=1,\dots,n$. For $Y\in\operatorname{Mat}_N(\mathbb{C})$
with $\|Y\|<1$ let $R_{\log}(Y):=\operatorname{Log}(1+Y)-Y$, with $\operatorname{Log}$ the principal
matrix logarithm. Then, with (C1) using only $\|X_j\|\le x_j$, (C2) using the common
per-factor majorants $x_j$ of $X_j$ and $\tilde X_j$, (C3) holding for $\|Y\|,\|\tilde Y\|\le\eta\le\tfrac12$,
and (C4) holding for $s\le\tfrac14$,
\begin{equation}\label{4.26:eq-exponential-products-a}
\begin{aligned}
&\text{(C1)}\quad \Bigl\|\prod_j e^{X_j}-1\Bigr\|\le e^s-1,\qquad
\Bigl\|\prod_j e^{X_j}-1-\sum_j X_j\Bigr\|\le e^s-1-s;\\
&\text{(C2)}\quad \Bigl\|\prod_j e^{X_j}-\prod_j e^{\tilde X_j}\Bigr\|\le\delta e^s,\\
&\phantom{\text{(C2)}}\quad
\Bigl\|\Bigl[\prod_j e^{X_j}-1-\sum_jX_j\Bigr]
-\Bigl[\prod_j e^{\tilde X_j}-1-\sum_j\tilde X_j\Bigr]\Bigr\|\le\delta(e^s-1);\\
&\text{(C3)}\quad \|R_{\log}(Y)\|\le\frac{\eta^2}{2(1-\eta)}\le\eta^2,\\
&\phantom{\text{(C3)}}\quad
\|R_{\log}(Y)-R_{\log}(\tilde Y)\|\le\frac{\eta}{1-\eta}\,\|Y-\tilde Y\|\le2\eta\,\|Y-\tilde Y\|;\\
&\text{(C4)}\quad e^s-1\le1.14\,s,\qquad e^s-1-s\le0.65\,s^2,\qquad e^s\le1.29.
\end{aligned}
\end{equation}
Both inequalities of (C1) hold also for the reversed product of the $e^{-X_j}$, that is,
with $\prod_j e^{X_j}$ replaced by $e^{-X_n}\cdots e^{-X_1}$ and each $X_j$ by $-X_j$.
\end{lemma}

\begin{proof}
\emph{The majorant bounds.} We argue monomial by monomial. For each word $w$,
submultiplicativity of the operator norm gives $\|c_wX^w\|\le|c_w|\,x^w\le a_wx^w$;
summing over $w$ shows that $\Phi(X)$ converges absolutely and $\|\Phi(X)\|\le\varphi(x)$.
For the difference, let $w=(j_1,\dots,j_r)$ and telescope:
\begin{equation*}
X^w-\tilde X^w=\sum_{i=1}^r X_{j_1}\cdots X_{j_{i-1}}\,(X_{j_i}-\tilde X_{j_i})\,
\tilde X_{j_{i+1}}\cdots\tilde X_{j_r}.
\end{equation*}
Since $X_j$ and $\tilde X_j$ are both bounded by $x_j$, the $i$-th term has norm at most
$\delta_{j_i}\prod_{l\ne i}x_{j_l}$, and summing over $i$ gives
$\|X^w-\tilde X^w\|\le\sum_k\delta_k\,\partial_k(x^w)$. Multiplying by $|c_w|\le a_w$
and summing over $w$ gives the second bound.

\emph{(C1).} Expanding each exponential and multiplying the absolutely convergent series,
\begin{equation*}
\prod_j e^{X_j}-1=\sum_{(m_1,\dots,m_n)\ne0}\frac{X_1^{m_1}\cdots X_n^{m_n}}{m_1!\cdots m_n!},
\end{equation*}
a series with nonnegative coefficients; evaluated commutatively at $x$ it is
$\prod_je^{x_j}-1=e^{\sum_jx_j}-1$, which is therefore a majorant. Hence
$\|\prod_je^{X_j}-1\|\le e^{\sum_jx_j}-1\le e^s-1$, since the exponential is increasing.
The terms of degree one of this series are exactly the $X_j$, with coefficient $1$;
removing them leaves the words of degree at least two with nonnegative coefficients,
whose majorant is $e^{\sum_jx_j}-1-\sum_jx_j$. As $t\mapsto e^t-1-t$ has derivative
$e^t-1\ge0$ for $t\ge0$, this is at most $e^s-1-s$. For $e^{-X_n}\cdots e^{-X_1}$ the
expansion has coefficients $(-1)^{m_1+\dots+m_n}/(m_1!\cdots m_n!)$, of the same absolute
values, on the reversed words in the variables $X_j$, equivalently with nonnegative
coefficients in the variables $-X_j$, whose norms are also bounded by $x_j$; the same
majorants and the same two bounds follow.

\emph{(C2).} Apply the difference bound to $\Phi(X)=\prod_je^{X_j}-1$ with majorant
$\varphi(x)=e^{\sum_jx_j}-1$. Here $\partial_k\varphi(x)=e^{\sum_jx_j}\le e^s$, so
the difference is at most $\sum_k\delta_ke^s=\delta e^s$; the constant $1$ cancels in the
difference. For the second bound take $\Phi(X)=\prod_je^{X_j}-1-\sum_jX_j$ with majorant
$\varphi(x)=e^{\sum_jx_j}-1-\sum_jx_j$; then
$\partial_k\varphi(x)=e^{\sum_jx_j}-1\le e^s-1$, and the bound is $\delta(e^s-1)$.

\emph{(C3).} Since $\|Y\|\le\eta\le\tfrac12<1$, the spectrum of $1+Y$ lies in the disc
$|z-1|<1$ and the principal logarithm is given by the logarithmic series, so
$R_{\log}(Y)=\sum_{k\ge2}(-1)^{k+1}Y^k/k$. Hence
\begin{equation*}
\|R_{\log}(Y)\|\le\sum_{k\ge2}\frac{\eta^k}{k}\le\frac12\sum_{k\ge2}\eta^k
=\frac{\eta^2}{2(1-\eta)}\le\eta^2,
\end{equation*}
the last step because $1-\eta\ge\tfrac12$. The one-variable case of the telescoping identity
gives $\|Y^k-\tilde Y^k\|\le k\eta^{k-1}\|Y-\tilde Y\|$, and therefore
\begin{equation*}
\|R_{\log}(Y)-R_{\log}(\tilde Y)\|\le\sum_{k\ge2}\eta^{k-1}\|Y-\tilde Y\|
=\frac{\eta}{1-\eta}\,\|Y-\tilde Y\|\le2\eta\,\|Y-\tilde Y\|.
\end{equation*}

\emph{(C4).} For $s=0$ the three inequalities hold trivially. For $0<s\le\tfrac14$ the
functions
\begin{equation*}
\frac{e^s-1}{s}=\sum_{m\ge0}\frac{s^m}{(m+1)!},\qquad
\frac{e^s-1-s}{s^2}=\sum_{m\ge0}\frac{s^m}{(m+2)!}
\end{equation*}
are increasing, being power series with positive coefficients. At $s=\tfrac14$ they equal
$1.136\ldots$ and $0.544\ldots$ respectively, since $e^{1/4}=1.284\ldots$; hence
$e^s-1\le1.14\,s$ and $e^s-1-s\le0.65\,s^2$. Finally $e^s\le e^{1/4}=1.284\ldots\le1.29$.
\end{proof}

\begin{lemma}[The comb pairing as straight part plus flux]\label{4.26:lem-straight-part-flux}
Fix a level $m$ and write $\varepsilon=\varepsilon_m$ and $\varepsilon'=\varepsilon_{m+1}=L\varepsilon$.
Let $y$ be the root of an $(m+1)$-block $B_{m+1}(y)$, let $\mu\in\{1,\dots,d\}$, and let
$c=\langle y,y+\varepsilon' e_\mu\rangle$ be the corresponding level-$(m+1)$ bond. We write
$(\mathcal{Q}A)_\mu(y)$ for $\mathcal{Q}A(c)$.

For $x\in B_{m+1}(y)\cap\Lambda_m$ the following are as in
Notation~\ref{4.26:not-combs-linearised-average}:
\begin{itemize}
\item the contour $\Gamma_{y,x}=(z_0=y,\dots,z_r=x)$ and its bonds $\beta_i$, with lower endpoints $w_i$,
orientation signs $s_i$ and directions $\nu_i$;
\item the comb loop $\mathrm{comb}_x$ and the pairing $\langle\cdot,\mathrm{comb}_x\rangle$;
\item the linearised average $\mathcal{Q}A$, the straight-line average $SA$ and the lattice curl $d_mA$.
\end{itemize}
We write $n_\perp(x;\mu)$ for the number of bonds of $\Gamma_{y,x}$ not parallel to $e_\mu$.

Let $F=(F_{\mu\nu}(z))$ be a family of matrices in $\operatorname{Mat}_N(\mathbb{C})$ indexed by the sites
$z\in\Lambda_m$ and pairs of directions. The flux operator $T$ is defined by the third line of the
display below, where, for each $x$, the data $w_i,s_i,\nu_i$ are those of $\Gamma_{y,x}$. We write
$\|TF\|_\infty=\sup_{y,\mu}\|(TF)_\mu(y)\|$, $\sup\|F\|=\sup_{z,\mu,\nu}\|F_{\mu\nu}(z)\|$ and
$\|A\|_\infty=\sup_b\|A(b)\|$ over level-$m$ bonds $b$.

Let $A$ be a bond field, and let $t_\perp$ satisfy $t_\perp\ge L^{-1-d}\sum_{x}n_\perp(x;\mu)$, the sum
running over $x\in B_{m+1}(y)\cap\Lambda_m$, for all blocks and directions. Then for every
$x\in B_{m+1}(y)\cap\Lambda_m$ the first line of the following display holds, the terms with
$\nu_i=\mu$ in it vanishing; consequently the second line holds; the first bound of the last line holds
for every family $F$ as above, and the second bound of the last line holds without any condition on
$t_\perp$:
\begin{equation}\label{4.26:eq-straight-part-flux-a}
\begin{aligned}
&\langle A,\mathrm{comb}_x\rangle=\sum_{k<L}A_\mu(y+k\varepsilon e_\mu)
-\varepsilon\sum_i s_i\sum_{k<L}(d_mA)_{\mu\nu_i}(w_i+k\varepsilon e_\mu),\\
&\mathcal{Q}A=SA+\varepsilon'\,T(d_mA),\\
&(TF)_\mu(y):=-L^{-2-d}\sum_{x\in B_{m+1}(y)\cap\Lambda_m}\ \sum_{i:\,\nu_i\neq\mu}s_i
\sum_{k<L}F_{\mu\nu_i}(w_i+k\varepsilon e_\mu),\\
&\|TF\|_\infty\le t_\perp\sup\|F\|,\qquad \|SA\|_\infty\le\|A\|_\infty .
\end{aligned}
\end{equation}
One may take $t_\perp:=(d-1)(L-1)/(2L)$.
\end{lemma}

\begin{proof}
Write $x'=x+\varepsilon'e_\mu$ and $y'=y+\varepsilon'e_\mu$, so that
$\mathrm{comb}_x=\Gamma_{y,x}\cdot[x,x']\cdot(\Gamma_{y',x'})^{-1}$. By the translation covariance of the
contours among the axioms \eqref{4.26:eq-block-average-axioms-a}
(Definition~\ref{4.26:def-block-average-axioms}), $\Gamma_{y',x'}=\Gamma_{y,x}+\varepsilon'e_\mu$.

We read off the three pieces of the comb. The contour $\Gamma_{y,x}$ traverses $\beta_i$ with orientation
$s_i$ and contributes $\sum_i s_iA_{\nu_i}(w_i)$. The straight segment $[x,x']$ consists of the $L$ bonds
$\langle x+k\varepsilon e_\mu,x+(k+1)\varepsilon e_\mu\rangle$, $k<L$. The reversed translated contour
traverses $\beta_i+\varepsilon'e_\mu$ with orientation $-s_i$. Since $\varepsilon'=L\varepsilon$, this gives
\begin{equation}\label{4.26:eq-straight-part-flux-2}
\langle A,\mathrm{comb}_x\rangle=\sum_i s_i\bigl[A_{\nu_i}(w_i)-A_{\nu_i}(w_i+L\varepsilon e_\mu)\bigr]
+\sum_{k<L}A_\mu(x+k\varepsilon e_\mu).
\end{equation}

By the definition of the lattice curl in \eqref{4.26:eq-combs-linearised-average-a}, for every site $z$ and
all directions $\mu,\nu$,
\begin{equation*}
A_\nu(z+\varepsilon e_\mu)-A_\nu(z)=\varepsilon(d_mA)_{\mu\nu}(z)+A_\mu(z+\varepsilon e_\nu)-A_\mu(z).
\end{equation*}
Apply this at $z=w+k\varepsilon e_\mu$ for $k=0,\dots,L-1$ and sum over $k$. The left side telescopes, and
\begin{equation}\label{4.26:eq-straight-part-flux-3}
\begin{split}
A_\nu(w+L\varepsilon e_\mu)-A_\nu(w)=\sum_{k<L}\bigl[&\varepsilon(d_mA)_{\mu\nu}(w+k\varepsilon e_\mu)\\
&+A_\mu(w+\varepsilon e_\nu+k\varepsilon e_\mu)-A_\mu(w+k\varepsilon e_\mu)\bigr].
\end{split}
\end{equation}

Fix $i$ and $k$. We claim that
\begin{equation*}
s_i\bigl[A_\mu(w_i+\varepsilon e_{\nu_i}+k\varepsilon e_\mu)-A_\mu(w_i+k\varepsilon e_\mu)\bigr]
=A_\mu(z_i+k\varepsilon e_\mu)-A_\mu(z_{i-1}+k\varepsilon e_\mu)
\end{equation*}
in both cases $s_i=\pm1$. Since $z_i=z_{i-1}+s_i\varepsilon e_{\nu_i}$, the lower endpoint of $\beta_i$ is
$w_i=z_{i-1}$ if $s_i=+1$ and $w_i=z_i$ if $s_i=-1$. In the first case $w_i+\varepsilon e_{\nu_i}=z_i$ and
the claim is immediate. In the second case $w_i+\varepsilon e_{\nu_i}=z_{i-1}$, and the factor $-1$ reverses
the difference.

The $i$-th bracket in \eqref{4.26:eq-straight-part-flux-2} is minus the left side of
\eqref{4.26:eq-straight-part-flux-3} with $w=w_i$ and $\nu=\nu_i$. We substitute
\eqref{4.26:eq-straight-part-flux-3} for it, and use the claim summed over $i$: for each $k$ this sum
telescopes along the contour, from $z_0=y$ to $z_r=x$, to
$A_\mu(x+k\varepsilon e_\mu)-A_\mu(y+k\varepsilon e_\mu)$. This gives
\begin{equation*}
\begin{split}
\langle A,\mathrm{comb}_x\rangle={}&\sum_{k<L}A_\mu(x+k\varepsilon e_\mu)
-\varepsilon\sum_i s_i\sum_{k<L}(d_mA)_{\mu\nu_i}(w_i+k\varepsilon e_\mu)\\
&-\sum_{k<L}\bigl[A_\mu(x+k\varepsilon e_\mu)-A_\mu(y+k\varepsilon e_\mu)\bigr].
\end{split}
\end{equation*}
The sums over $A_\mu(x+k\varepsilon e_\mu)$ cancel, and this is the first line of
\eqref{4.26:eq-straight-part-flux-a}. The terms with $\nu_i=\mu$ vanish because $(d_mA)_{\mu\mu}=0$ by the
definition of the curl: the four terms
$A_\mu(z+\varepsilon e_\mu)-A_\mu(z)-A_\mu(z+\varepsilon e_\mu)+A_\mu(z)$ cancel.

By \eqref{4.26:eq-combs-linearised-average-a},
$(\mathcal{Q}A)_\mu(y)=L^{-1-d}\sum_x\langle A,\mathrm{comb}_x\rangle$, the sum running over the $L^d$ sites
of $B_{m+1}(y)\cap\Lambda_m$. We average the first line of \eqref{4.26:eq-straight-part-flux-a} term by term.
The first term does not depend on $x$, so its average is
\begin{equation*}
L^{-1}\sum_{k<L}A_\mu(y+k\varepsilon e_\mu)=(SA)_\mu(y).
\end{equation*}
In the second term we drop the vanishing summands with $\nu_i=\mu$. Since
$\varepsilon L^{-1-d}=\varepsilon' L^{-2-d}$ (that is, $\varepsilon/L=\varepsilon'/L^2$), the average of the
second term is $\varepsilon'\,(T(d_mA))_\mu(y)$ by the third line of \eqref{4.26:eq-straight-part-flux-a}.
This is the second line of \eqref{4.26:eq-straight-part-flux-a}.

For the bound on $T$, fix $x$. In the third line of \eqref{4.26:eq-straight-part-flux-a} the site $x$
contributes at most $n_\perp(x;\mu)\cdot L$ terms, each of norm at most $\sup\|F\|$, with prefactor
$L^{-2-d}$. Hence
\begin{equation*}
\|(TF)_\mu(y)\|\le L^{-1-d}\sum_x n_\perp(x;\mu)\,\sup\|F\|\le t_\perp\sup\|F\|,
\end{equation*}
and taking the supremum over $y$ and $\mu$ gives the first bound of the last line.

It remains to show that $t_\perp=(d-1)(L-1)/(2L)$ is admissible. For the contour $\Gamma_{y,x}$ of
Notation~\ref{4.26:not-combs-linearised-average},
$n_\perp(x;\mu)=\sum_{\nu\neq\mu}|x_\nu-y_\nu|/\varepsilon$. Let $x^{0}$ be the corner of the block with
smallest coordinates, write the sites of $B_{m+1}(y)\cap\Lambda_m$ as
$x^{0}+\varepsilon\sum_\nu j_\nu e_\nu$ with $j_\nu\in\{0,\dots,L-1\}$, and write the root as
$y=x^{0}+\varepsilon\sum_\nu r_\nu e_\nu$. Then $|x_\nu-y_\nu|/\varepsilon=|j_\nu-r_\nu|$. The mean of
$n_\perp(x;\mu)$ over the block is therefore $\sum_{\nu\neq\mu}$ of the mean over $j\in\{0,\dots,L-1\}$ of
$|j-r_\nu|$. For $r\in\{0,\dots,L-1\}$, summing separately over $j\le r$ and $j\ge r$ gives
\begin{equation*}
\frac1L\sum_{j=0}^{L-1}|j-r|=\frac{r(r+1)+(L-1-r)(L-r)}{2L}.
\end{equation*}
The numerator equals $2r^2-2(L-1)r+L(L-1)$. This is convex in $r$, so its maximum over
$r\in\{0,\dots,L-1\}$ is attained at $r\in\{0,L-1\}$, where it equals $L(L-1)$. Hence the mean of $|j-r|$
is at most $(L-1)/2$ for any root position $r$, and the mean of $n_\perp(x;\mu)$ over the block is at most
$(d-1)(L-1)/2$. Since $L^{-1-d}\sum_x n_\perp(x;\mu)$ is $L^{-1}$ times this mean, it is at most
$(d-1)(L-1)/(2L)$.

Finally, $(SA)_\mu(y)$ is the mean of the $L$ values $A_\mu(y+k\varepsilon e_\mu)$, $k<L$. By the triangle
inequality $\|(SA)_\mu(y)\|\le\|A\|_\infty$, which is the second bound of the last line.
\end{proof}

\begin{lemma}[The coarse curl of the linearised average]\label{4.26:lem-coarse-curl}
Assume conditions \textup{(M2)} and \textup{(M3)} of the axioms
\eqref{4.26:eq-block-average-axioms-a} for the block average
(Definition~\ref{4.26:def-block-average-axioms}): one comb contour for all directions,
coarse translation covariance, block offsets that are the same for every block, and
weights that do not depend on the direction. Let $A$ be a bond field, let
$\mathcal{Q}A$ be its linearised average and $d_m$, $d_{m+1}$ the lattice curls at
spacings $\varepsilon_m$, $\varepsilon_{m+1}$
(Notation~\ref{4.26:not-combs-linearised-average}). Let $P=(y;\mu,\lambda)$ be a coarse
plaquette, $y\in\Lambda_{m+1}$, whose four bonds lie in $\Omega_{m+1}$. Then
\begin{equation}\label{4.26:eq-coarse-curl-a}
(d_{m+1}\mathcal{Q}A)_{\mu\lambda}(y)
=L^{-2-d}\sum_{\rho}\ \sum_{p\subset P_\rho}(d_mA)_{\mu\lambda}(p),
\end{equation}
where $\rho$ ranges over the $L^d$ offsets $x-y$ of the sites
$x\in B_{m+1}(y)\cap\Lambda_m$, and $P_\rho$ is the $L\times L$ square of fine plaquettes
in the $\mu\lambda$-plane with corner $y+\rho$. In particular
\begin{equation}\label{4.26:eq-coarse-curl-1}
\|(d_{m+1}\mathcal{Q}A)_{\mu\lambda}(y)\|
\le\sup_{p}\|(d_mA)_{\mu\lambda}(p)\|,
\end{equation}
the supremum taken over the fine plaquettes $p\subset P_\rho$ occurring in
\eqref{4.26:eq-coarse-curl-a}: the right side of \eqref{4.26:eq-coarse-curl-a} is an
average with positive weights, and the factor is exactly $1$.
\end{lemma}

\begin{proof}
Write $\varepsilon=\varepsilon_m$ and $\varepsilon'=\varepsilon_{m+1}=L\varepsilon$. For an
oriented path $\Gamma$ of fine bonds $\beta_1,\dots,\beta_n$, where $\beta_i$ has lower
endpoint $w_i$, direction $\nu_i$ and orientation sign $s_i$ relative to $\Gamma$, we write
\begin{equation*}
\langle A,\Gamma\rangle=\sum_{i=1}^{n}s_i\,A_{\nu_i}(w_i)
\end{equation*}
for the signed sum of $A$ along $\Gamma$; it is additive under concatenation of paths and
changes sign under reversal. For $z\in\Lambda_m$ and a direction $\nu$ let $[z,z+\varepsilon'e_\nu]$
denote the straight segment of $L$ fine bonds from $z$ to $z+\varepsilon'e_\nu$.

By \eqref{4.26:eq-combs-linearised-average-a}
(Notation~\ref{4.26:not-combs-linearised-average}), the linearised average of the coarse
bond $(y',\nu)$, $y'\in\Lambda_{m+1}$, is given by
$\varepsilon'(\mathcal{Q}A)_\nu(y')=L^{-d}\sum_{x}\varepsilon\langle A,\mathrm{comb}_x\rangle$,
the sum running over the sites $x\in B_{m+1}(y')\cap\Lambda_m$, where the comb loop
$\mathrm{comb}_x$ consists of the contour $\Gamma_{y',x}$, the straight segment
$[x,x+\varepsilon'e_\nu]$ and the contour $\Gamma_{y'+\varepsilon'e_\nu,\,x+\varepsilon'e_\nu}$
traversed backwards, from $x+\varepsilon'e_\nu$ to $y'+\varepsilon'e_\nu$; by \textup{(M3)} all
weights equal $L^{-d}$, for every $\nu$. Writing
$x=y'+\rho$, so that by \textup{(M2)} the offsets $\rho$ are the same for every block, and
$\Gamma_{y'+\varepsilon'e_\nu,\,y'+\varepsilon'e_\nu+\rho}=\Gamma_{y',y'+\rho}+\varepsilon'e_\nu$,
this reads
\begin{equation}\label{4.26:eq-coarse-curl-2}
\varepsilon'(\mathcal{Q}A)_\nu(y')=L^{-d}\,\varepsilon\sum_{\rho}
\Bigl(\phi_\rho(y')+\langle A,[y'+\rho,\,y'+\rho+\varepsilon'e_\nu]\rangle
-\phi_\rho(y'+\varepsilon'e_\nu)\Bigr),
\end{equation}
where, for $z\in\Lambda_{m+1}$, $\phi_\rho(z)=\langle A,\Gamma_{z,z+\rho}\rangle$, and, by
\textup{(M2)}, $\Gamma_{z,z+\rho}$
is the translate by $z-y$ of $\Gamma_{y,y+\rho}$ and is the same contour for $\nu=\mu$ and
$\nu=\lambda$.

By the definition of the lattice curl at spacing $\varepsilon'$,
\begin{equation}\label{4.26:eq-coarse-curl-3}
\varepsilon'(d_{m+1}\mathcal{Q}A)_{\mu\lambda}(y)
=(\mathcal{Q}A)_\lambda(y+\varepsilon'e_\mu)-(\mathcal{Q}A)_\lambda(y)
-(\mathcal{Q}A)_\mu(y+\varepsilon'e_\lambda)+(\mathcal{Q}A)_\mu(y).
\end{equation}
The blocks rooted at the neighbouring corners $y+\varepsilon'e_\mu$, $y+\varepsilon'e_\lambda$,
$y+\varepsilon'e_\mu+\varepsilon'e_\lambda$ are translates of $B_{m+1}(y)$, so the same offsets
$\rho$ and the same weights $L^{-d}$ occur in all four terms, and we may insert
\eqref{4.26:eq-coarse-curl-2}, divided by $\varepsilon'$, in \eqref{4.26:eq-coarse-curl-3}:
the right side of \eqref{4.26:eq-coarse-curl-3} becomes $L^{-d}\varepsilon/\varepsilon'$ times
a sum over $\rho$, and we collect, for each fixed $\rho$, the contour terms in it. Put
$v_\mu=\varepsilon'e_\mu$, $v_\lambda=\varepsilon'e_\lambda$. For fixed $\rho$ these contour
terms are
\begin{equation*}
\begin{split}
&\bigl[\phi_\rho(y+v_\mu)-\phi_\rho(y+v_\mu+v_\lambda)\bigr]
-\bigl[\phi_\rho(y)-\phi_\rho(y+v_\lambda)\bigr]\\
&\qquad-\bigl[\phi_\rho(y+v_\lambda)-\phi_\rho(y+v_\mu+v_\lambda)\bigr]
+\bigl[\phi_\rho(y)-\phi_\rho(y+v_\mu)\bigr]=0:
\end{split}
\end{equation*}
the translated contours rooted at the four corners $y$, $y+v_\mu$, $y+v_\lambda$,
$y+v_\mu+v_\lambda$ of the coarse
square each occur once with the sign $+$ and once with the sign $-$, and cancel. This uses
that the contour is the same for the directions $\mu$ and $\lambda$. What remains are the four
straight segments,
\begin{equation*}
\begin{split}
&\langle A,[y+\rho,\,y+\rho+v_\mu]\rangle
+\langle A,[y+\rho+v_\mu,\,y+\rho+v_\mu+v_\lambda]\rangle\\
&\qquad-\langle A,[y+\rho+v_\lambda,\,y+\rho+v_\mu+v_\lambda]\rangle
-\langle A,[y+\rho,\,y+\rho+v_\lambda]\rangle
=\langle A,\partial P_\rho\rangle,
\end{split}
\end{equation*}
the boundary of $P_\rho$ traversed positively in the $\mu\lambda$-plane.

We now apply the lattice Stokes identity on $P_\rho$: every fine bond interior to the square
$P_\rho$ lies in exactly two of its
$L^2$ fine plaquettes, which traverse it in opposite directions, and every bond of
$\partial P_\rho$ lies in exactly one, traversed in the orientation of $\partial P_\rho$.
Hence $\langle A,\partial P_\rho\rangle=\sum_{p\subset P_\rho}\langle A,\partial p\rangle$,
and by the definition of the lattice curl at spacing $\varepsilon$,
$\langle A,\partial p\rangle=\varepsilon(d_mA)_{\mu\lambda}(p)$. Therefore
\begin{equation*}
\varepsilon'(d_{m+1}\mathcal{Q}A)_{\mu\lambda}(y)
=\frac{L^{-d}\varepsilon}{\varepsilon'}\sum_\rho\varepsilon\sum_{p\subset P_\rho}
(d_mA)_{\mu\lambda}(p).
\end{equation*}
Dividing by $\varepsilon'=L\varepsilon$ gives the coefficient
$L^{-d}(\varepsilon/\varepsilon')^2=L^{-2-d}$, which is \eqref{4.26:eq-coarse-curl-a}.

The right side of \eqref{4.26:eq-coarse-curl-a} has $L^d$ offsets $\rho$, each with $L^2$
plaquettes $p\subset P_\rho$, all with the positive coefficient $L^{-2-d}$; the total mass is
$L^{-2-d}\cdot L^d\cdot L^2=1$. The triangle inequality for the operator norm then gives
\begin{equation*}
\|(d_{m+1}\mathcal{Q}A)_{\mu\lambda}(y)\|
\le L^{-2-d}\sum_\rho\sum_{p\subset P_\rho}\|(d_mA)_{\mu\lambda}(p)\|
\le\sup_p\|(d_mA)_{\mu\lambda}(p)\|,
\end{equation*}
which is \eqref{4.26:eq-coarse-curl-1}.
\end{proof}

\begin{proposition}[The one-step averaging estimate]\label{4.26:prop-one-step-averaging}
Work in the setting of Definition~\ref{4.26:def-block-average-axioms} and
Notation~\ref{4.26:not-combs-linearised-average}. A level $m$ is fixed, with
$\varepsilon=\varepsilon_m$ and $\varepsilon'=\varepsilon_{m+1}=L\varepsilon$, and $\Omega$ is a box of the
fine lattice. The numbers $r_{\mathcal{N}}$, $C_{\mathcal{N}}$, $c_{\Pi}$ are the constants of the
normalisation $\mathcal{N}$ in axiom (M4) of Definition~\ref{4.26:def-block-average-axioms}. The number
$t_{\perp}$ is the constant of
Lemma~\ref{4.26:lem-straight-part-flux}. We write $1$ for the identity matrix and $\operatorname{Log}$ for
the principal matrix logarithm.

Let $d\ge2$ and $c=\langle y,y+\varepsilon'e_\mu\rangle\in\Omega_{m+1}$. On the level-$m$ bonds of
$R_{m+1}(c)$ let $A$ be
$\mathfrak{sl}(N,\mathbb{C})$-valued with $\|A\|\le\sigma$, and assume $\|d_mA\|\le f$ on the level-$m$
plaquettes of $R_{m+1}(c)$.

For \textup{(P3)} we assume moreover that $c+\varepsilon'e_\lambda\in\Omega_{m+1}$, that the same
hypotheses hold on $R_{m+1}(c+\varepsilon'e_\lambda)$, and that
\begin{equation*}
\|A(\beta+\varepsilon'e_\lambda)-A(\beta)\|\le\varepsilon' l
\quad\text{for the level-$m$ bonds $\beta$ of $R_{m+1}(c)$.}
\end{equation*}

For \textup{(P4)} let $P=(y;\mu,\lambda)$ be a coarse plaquette with its four bonds in
$\Omega_{m+1}$. Its bonds at $y$ are the bond $c$ above and
$c'=\langle y,y+\varepsilon'e_\lambda\rangle$, and their translates are $c+\varepsilon'e_\lambda$ and
$c'+\varepsilon'e_\mu$. Let $R_{m+1}(P)$ be the union of the four $(m+1)$-blocks at the corners of $P$.
We assume the hypotheses on $\|A\|$ and $\|d_mA\|$ on $R_{m+1}(P)$, together with the hypothesis with
$l$ for the two translation pairs of $\partial P$. That is,
\begin{equation*}
\begin{aligned}
&\|A(\beta+\varepsilon'e_\lambda)-A(\beta)\|\le\varepsilon'l
&&\text{for the level-$m$ bonds $\beta$ of $R_{m+1}(c)$,}\\
&\|A(\beta+\varepsilon'e_\mu)-A(\beta)\|\le\varepsilon'l
&&\text{for the level-$m$ bonds $\beta$ of $R_{m+1}(c')$.}
\end{aligned}
\end{equation*}

Put $s:=(2d+1)\varepsilon'\sigma$ and
\begin{align*}
\mathsf{C}^{\mathrm{av}}_3&:=1.3\,C_{\mathcal{N}}+0.65\,c_{\Pi},&
\mathsf{C}^{\mathrm{av}}_4&:=3\,C_{\mathcal{N}}+1.2\,c_{\Pi},\\
\mathsf{C}^{\mathrm{av}}_1&:=(9\,\mathsf{C}^{\mathrm{av}}_3+4)\,d^2,&
\mathsf{C}^{\mathrm{av}}_2&:=(9\,\mathsf{C}^{\mathrm{av}}_4+24)\,d^2,
\end{align*}
and assume
\begin{equation}\label{4.26:eq-one-step-averaging-a}
s\le\tfrac14,\qquad 1.14\,s<r_{\mathcal{N}},\qquad (2d+1)\,\mathsf{C}^{\mathrm{av}}_3\,s\le1,
\qquad \mathsf{C}^{\mathrm{av}}_4\,s\le1.
\end{equation}

Let $W=\exp(i\varepsilon A)$ on the level-$m$ bonds. Let $P_x$, for $x\in B_{m+1}(y)\cap\Lambda_m$,
be the comb transporters of axiom (M1) at the bond $c$. At every bond
at which the hypotheses hold, define $\mathcal{E}$, $A'$ and $E$ by the identities below. Then:
\begin{equation}\label{4.26:eq-one-step-averaging-b}
\begin{aligned}
&\text{(P1)}\quad
 \|P_x-1\|\le 1.14\,s,\quad \|P_x^{-1}-1\|\le 1.14\,s\quad\text{for every }x,\\
&\phantom{\text{(P1)}\quad}
 M(W)(c)=1+i\varepsilon'\mathcal{Q}A(c)+\mathcal{E}(c),\qquad
 \|\mathcal{E}(c)\|\le\mathsf{C}^{\mathrm{av}}_3\,s^2;\\
&\text{(P2)}\quad
 \|M(W)(c)-1\|\le 0.8\,s\le 0.2,\qquad
 A'(c):=(i\varepsilon')^{-1}\operatorname{Log}M(W)(c),\\
&\phantom{\text{(P2)}\quad}
 A'(c)=\mathcal{Q}A(c)+E(c)=SA(c)+\varepsilon'\,T(d_mA)(c)+E(c),\qquad
 \|E(c)\|\le\mathsf{C}^{\mathrm{av}}_1\,\varepsilon'\sigma^2;\\
&\phantom{\text{(P2)}\quad}
 \text{hence}\quad
 \|A'(c)\|\le\sigma+\varepsilon'\bigl(t_{\perp}f+\mathsf{C}^{\mathrm{av}}_1\sigma^2\bigr);\\
&\text{(P3)}\quad
 \|E(c+\varepsilon'e_\lambda)-E(c)\|\le\mathsf{C}^{\mathrm{av}}_2\,(\varepsilon')^2\sigma l;\\
&\text{(P4)}\quad
 \|d_{m+1}A'(P)\|\le\sup_{R_{m+1}(P)}\|d_mA\|+2\,\mathsf{C}^{\mathrm{av}}_2\,\varepsilon'\sigma l.
\end{aligned}
\end{equation}
\end{proposition}

\begin{proof}
Throughout, $c=\langle y,y'\rangle$ with $y'=y+\varepsilon'e_\mu$. The site $x$ ranges over
$B_{m+1}(y)\cap\Lambda_m$, and $x'=x+\varepsilon'e_\mu$. Families $H=(H_x)_x$ of matrices indexed by
these sites are normed by $\|H\|:=\max_x\|H_x\|$, as in axiom (M4).

\emph{Proof of} (P1).
By axiom (M1) of Definition~\ref{4.26:def-block-average-axioms} the comb of $x$ is
$\mathrm{comb}_x=\Gamma_{y,x}\cdot[x,x']\cdot(\Gamma_{y',x'})^{-1}$. By axiom (M2),
$\Gamma_{y',x'}=\Gamma_{y,x}+\varepsilon'e_\mu$, so both contours have $|\Gamma_{y,x}|\le d(L-1)$
bonds. The straight path $[x,x']$ has $L$ bonds. Hence the number $n_x$ of bonds of $\mathrm{comb}_x$
satisfies
\begin{equation*}
n_x=2|\Gamma_{y,x}|+L\le(2d+1)L .
\end{equation*}
Enumerate the bonds of $\mathrm{comb}_x$ in the order of traversal as
$\beta_{x,1},\dots,\beta_{x,n_x}$. Let $o_{x,j}\in\{\pm1\}$ be the traversal orientation of
$\beta_{x,j}$. Forward bonds contribute $W(b)$,
backward bonds contribute $W(b)^{-1}=\exp(-i\varepsilon A(b))$. Hence
\begin{equation*}
P_x=\prod_{j=1}^{n_x}e^{X_{x,j}},\qquad X_{x,j}:=i\varepsilon\,o_{x,j}A(\beta_{x,j}),\qquad
\sum_{j}X_{x,j}=i\varepsilon\langle A,\mathrm{comb}_x\rangle .
\end{equation*}
The bonds of $\mathrm{comb}_x$ are level-$m$ bonds of $R_{m+1}(c)$, so $\|X_{x,j}\|\le\varepsilon\sigma$
for every $j$. The sum of these majorants is at most
\begin{equation*}
(2d+1)L\varepsilon\sigma=(2d+1)\varepsilon'\sigma=s .
\end{equation*}
By (C1) of Lemma~\ref{4.26:lem-exponential-products},
\begin{equation*}
\|P_x-1\|\le e^s-1
\quad\text{and}\quad
\bigl\|P_x-1-i\varepsilon\langle A,\mathrm{comb}_x\rangle\bigr\|\le e^s-1-s .
\end{equation*}
The same bounds hold for $P_x^{-1}$, which is the reversed product of the $e^{-X_{x,j}}$. With
$s\le\tfrac14$ from \eqref{4.26:eq-one-step-averaging-a}, the numerical bounds (C4) of
Lemma~\ref{4.26:lem-exponential-products} turn these into
\begin{equation*}
P_x=1+i\varepsilon\langle A,\mathrm{comb}_x\rangle+R_x,\qquad \|R_x\|\le0.65\,s^2,\qquad
\|P_x-1\|\le1.14\,s,
\end{equation*}
and likewise $\|P_x^{-1}-1\|\le1.14\,s$. This is the first line of (P1).

Put $V:=(P_x-1)_x$. Then $\|V\|\le1.14\,s<r_{\mathcal{N}}$ by \eqref{4.26:eq-one-step-averaging-a}, so
the family $\mathbf{1}+V=(P_x)_x$ lies in the polydisc on which axiom (M4) makes $\mathcal{N}$
holomorphic. By (M4), $\mathcal{N}(g,\dots,g)=g$; at $g=1$ this gives $\mathcal{N}(\mathbf{1})=1$.

Write $\Pi=D\mathcal{N}(\mathbf{1})$, a linear map on all families, and put
\begin{equation*}
R_{\mathcal{N}}(V):=\mathcal{N}(\mathbf{1}+V)-1-\Pi[V].
\end{equation*}
The segment $\theta\mapsto\mathbf{1}+\theta V$, $0\le\theta\le1$, lies in the polydisc. Taylor's formula
with integral remainder along it, together with the bound
$\|D^2\mathcal{N}[H,K]\|\le2C_{\mathcal{N}}\|H\|\|K\|$ on the polydisc from (M4), gives
\begin{equation*}
\|R_{\mathcal{N}}(V)\|\le C_{\mathcal{N}}\|V\|^2 .
\end{equation*}
By linearity of $\Pi$,
\begin{equation*}
\Pi[V]=\Pi\bigl[(i\varepsilon\langle A,\mathrm{comb}_x\rangle)_x\bigr]+\Pi\bigl[(R_x)_x\bigr].
\end{equation*}
Every $\langle A,\mathrm{comb}_x\rangle$ is traceless, since $A$ is
$\mathfrak{sl}(N,\mathbb{C})$-valued. On traceless families $\Pi$ is the arithmetic mean, by (M4). Using
$\varepsilon=\varepsilon'/L$ and the definition of $\mathcal{Q}A$ in
\eqref{4.26:eq-combs-linearised-average-a},
\begin{equation*}
\Pi\bigl[(i\varepsilon\langle A,\mathrm{comb}_x\rangle)_x\bigr]
=i\varepsilon L^{-d}\sum_x\langle A,\mathrm{comb}_x\rangle
=i\varepsilon' L^{-1-d}\sum_x\langle A,\mathrm{comb}_x\rangle
=i\varepsilon'\mathcal{Q}A(c).
\end{equation*}
By (M1), $M(W)(c)=\mathcal{N}((P_x)_x)=\mathcal{N}(\mathbf{1}+V)$. Hence
$M(W)(c)=1+i\varepsilon'\mathcal{Q}A(c)+\mathcal{E}(c)$ with
$\mathcal{E}(c):=\Pi[(R_x)_x]+R_{\mathcal{N}}(V)$. Since $\|\Pi\|\le c_\Pi$ and $1.14^2\le1.3$,
\begin{equation*}
\|\mathcal{E}(c)\|\le0.65\,c_{\Pi}s^2+C_{\mathcal{N}}(1.14\,s)^2\le\mathsf{C}^{\mathrm{av}}_3s^2 .
\end{equation*}
This completes (P1).

\emph{Proof of} (P2).
We have
\begin{equation*}
\|i\varepsilon'\mathcal{Q}A(c)\|
\le\varepsilon' L^{-1-d}\sum_x\|\langle A,\mathrm{comb}_x\rangle\|
\le\varepsilon L^{-d}\sum_x n_x\sigma .
\end{equation*}
The mean over the $L^d$ sites $x$ of the block of $n_x=2|\Gamma_{y,x}|+L$ is at most $L+d(L-1)$, the
mean of $|\Gamma_{y,x}|$ being at most $d(L-1)/2$. Since $L+d(L-1)\le(d+1)L$ and $d\ge2$,
\begin{equation*}
\|i\varepsilon'\mathcal{Q}A(c)\|\le\varepsilon\bigl(L+d(L-1)\bigr)\sigma\le(d+1)\varepsilon'\sigma
=\frac{d+1}{2d+1}\,s\le0.6\,s .
\end{equation*}
By (P1),
\begin{equation*}
\|\mathcal{E}(c)\|\le\mathsf{C}^{\mathrm{av}}_3s^2\le\frac{s}{2d+1}\le0.2\,s,
\end{equation*}
where the third condition of \eqref{4.26:eq-one-step-averaging-a} and $d\ge2$ were used.

Put $Z:=i\varepsilon'\mathcal{Q}A(c)+\mathcal{E}(c)$, so that $M(W)(c)=1+Z$ by (P1). Then
$\|Z\|\le0.8\,s$, and $0.8\,s\le0.2$ because $s\le\tfrac14$. This is the first assertion of (P2).

Since $\|Z\|<1$, the principal logarithm of $1+Z$ is defined. Write
$\operatorname{Log}(1+Z)=Z+R_{\log}(Z)$. With $\|Z\|\le0.8\,s$ and $\|Z\|\le0.2$, (C3) of
Lemma~\ref{4.26:lem-exponential-products} gives
\begin{equation*}
\|R_{\log}(Z)\|\le\frac{0.64\,s^2}{1.6}=0.4\,s^2 .
\end{equation*}
Consequently
\begin{equation*}
A'(c)=(i\varepsilon')^{-1}\bigl(Z+R_{\log}(Z)\bigr)=\mathcal{Q}A(c)+E(c),\qquad
E(c)=(i\varepsilon')^{-1}\bigl(\mathcal{E}(c)+R_{\log}(Z)\bigr).
\end{equation*}
Using $s^2/\varepsilon'=(2d+1)^2\varepsilon'\sigma^2$ and $(2d+1)^2\le6.25\,d^2$ for $d\ge2$,
\begin{align*}
\|E(c)\|&\le(\mathsf{C}^{\mathrm{av}}_3+0.4)\,\frac{s^2}{\varepsilon'}
=(\mathsf{C}^{\mathrm{av}}_3+0.4)(2d+1)^2\varepsilon'\sigma^2\\
&\le(6.25\,\mathsf{C}^{\mathrm{av}}_3+2.5)\,d^2\varepsilon'\sigma^2
\le\mathsf{C}^{\mathrm{av}}_1\varepsilon'\sigma^2 .
\end{align*}
The decomposition $\mathcal{Q}A(c)=SA(c)+\varepsilon'\,T(d_mA)(c)$ is the identity
\eqref{4.26:eq-straight-part-flux-a} of Lemma~\ref{4.26:lem-straight-part-flux}.

That lemma also gives $\|SA\|_\infty\le\|A\|_\infty$ and $\|TF\|_\infty\le t_\perp\sup\|F\|$. With the
hypotheses $\|A\|\le\sigma$ and $\|d_mA\|\le f$, they give $\|SA(c)\|\le\sigma$ and
$\|\varepsilon'T(d_mA)(c)\|\le\varepsilon't_\perp f$. Together with the bound on $E(c)$,
\begin{equation*}
\|A'(c)\|\le\sigma+\varepsilon'\bigl(t_\perp f+\mathsf{C}^{\mathrm{av}}_1\sigma^2\bigr).
\end{equation*}

\emph{Proof of} (P3).
Put $v:=\varepsilon'e_\lambda$ and $\tilde c:=c+v$. By axioms (M2) and (M3), the contours are
translation covariant under $v$, and the block offsets and the symmetric entry of the sites are the same
for every block. Hence
\begin{equation*}
B_{m+1}(y+v)=B_{m+1}(y)+v,\qquad \mathrm{comb}_{x+v}(\tilde c)=\mathrm{comb}_x(c)+v .
\end{equation*}
So $M(W)(\tilde c)=\mathcal{N}((\tilde P_x)_x)$, where $\tilde P_x$ is the ordered product over the
translated comb. It has the same length $n_x$ and the same orientation signs, and its factors are
$e^{\tilde X_{x,j}}$ with $\tilde X_{x,j}:=i\varepsilon\,o_{x,j}A(\beta_{x,j}+v)$.

The per-factor majorants $\varepsilon\sigma$ are common to $X_{x,j}$ and $\tilde X_{x,j}$, because the
hypothesis $\|A\|\le\sigma$ holds on both $R_{m+1}(c)$ and $R_{m+1}(\tilde c)$. Their sum over $j$ is at
most $s$, as in the proof of (P1). The bonds $\beta_{x,j}$ lie in $R_{m+1}(c)$, so by the hypothesis
with $l$,
\begin{equation*}
\|X_{x,j}-\tilde X_{x,j}\|
=\varepsilon\|A(\beta_{x,j})-A(\beta_{x,j}+v)\|\le\varepsilon\varepsilon' l .
\end{equation*}
Summing over $j$ gives at most $(2d+1)L\varepsilon\varepsilon'l=\delta_0$, where
$\delta_0:=(2d+1)(\varepsilon')^2l$.

First, the comb transporters. Put
$\tilde R_x:=\tilde P_x-1-\sum_j\tilde X_{x,j}$ and $\tilde V:=(\tilde P_x-1)_x$. By (C2) of
Lemma~\ref{4.26:lem-exponential-products} and the numerical bounds (C4) for $s\le\tfrac14$,
\begin{equation*}
\|P_x-\tilde P_x\|\le\delta_0e^s\le1.29\,\delta_0,\qquad
\|R_x-\tilde R_x\|\le\delta_0(e^s-1)\le1.14\,s\,\delta_0
\end{equation*}
for every $x$. Hence the families satisfy
$\|V-\tilde V\|\le1.29\,\delta_0$ and $\|R-\tilde R\|\le1.14\,s\,\delta_0$, with $R:=(R_x)_x$ and
$\tilde R:=(\tilde R_x)_x$.

Second, the remainder $\mathcal{E}$. By the proof of (P1) at $c$ and at $\tilde c$,
\begin{equation*}
\mathcal{E}(c)=\Pi[R]+R_{\mathcal{N}}(V),\qquad
\mathcal{E}(\tilde c)=\Pi[\tilde R]+R_{\mathcal{N}}(\tilde V).
\end{equation*}
For any two families $V_1$ and $V_2$ with $\|V_1\|,\|V_2\|\le\rho:=1.14\,s$, put
$V_\theta:=V_2+\theta(V_1-V_2)$ for $0\le\theta\le1$; these lie in the ball of radius $\rho$ by
convexity, hence inside the polydisc. Then
\begin{equation*}
R_{\mathcal{N}}(V_1)-R_{\mathcal{N}}(V_2)
=\int_0^1\bigl\{D\mathcal{N}(\mathbf{1}+V_\theta)-D\mathcal{N}(\mathbf{1})\bigr\}[V_1-V_2]\,d\theta .
\end{equation*}
Integrating the bound $\|D^2\mathcal{N}\|\le2C_{\mathcal{N}}$ of (M4) along the segment from
$\mathbf{1}$ to $\mathbf{1}+V_\theta$ gives
$\|D\mathcal{N}(\mathbf{1}+V_\theta)-D\mathcal{N}(\mathbf{1})\|\le2C_{\mathcal{N}}\rho$. By (P1),
$\|V\|,\|\tilde V\|\le\rho$, so, taking $V_1=V$ and $V_2=\tilde V$,
\begin{equation*}
\|R_{\mathcal{N}}(V)-R_{\mathcal{N}}(\tilde V)\|\le2C_{\mathcal{N}}\cdot1.14\,s\cdot1.29\,\delta_0
\le2.95\,C_{\mathcal{N}}\,s\,\delta_0 .
\end{equation*}
Since $\Pi[R]-\Pi[\tilde R]=\Pi[R-\tilde R]$ and $\|\Pi\|\le c_\Pi$, it follows that
\begin{equation*}
\|\mathcal{E}(c)-\mathcal{E}(\tilde c)\|\le(1.14\,c_\Pi+2.95\,C_{\mathcal{N}})\,s\,\delta_0
\le\mathsf{C}^{\mathrm{av}}_4\,s\,\delta_0
=\mathsf{C}^{\mathrm{av}}_4(2d+1)^2(\varepsilon')^3\sigma l .
\end{equation*}

Third, the linear part. By the same translation,
\begin{equation*}
\mathcal{Q}A(\tilde c)=L^{-1-d}\sum_x\sum_j o_{x,j}A(\beta_{x,j}+v).
\end{equation*}
The hypothesis with $l$ then gives
\begin{equation*}
\|\mathcal{Q}A(c)-\mathcal{Q}A(\tilde c)\|\le L^{-1-d}\cdot L^d\cdot(2d+1)L\cdot\varepsilon' l
=(2d+1)\varepsilon' l .
\end{equation*}
Put $Z:=M(W)(c)-1$ and $\tilde Z:=M(W)(\tilde c)-1$. By (P1) at both bonds,
\begin{equation*}
Z-\tilde Z=i\varepsilon'\bigl(\mathcal{Q}A(c)-\mathcal{Q}A(\tilde c)\bigr)
+\mathcal{E}(c)-\mathcal{E}(\tilde c).
\end{equation*}
Using $s=(2d+1)\varepsilon'\sigma$ and $\mathsf{C}^{\mathrm{av}}_4s\le1$ from
\eqref{4.26:eq-one-step-averaging-a},
\begin{align*}
\|Z-\tilde Z\|&\le(2d+1)(\varepsilon')^2l+\mathsf{C}^{\mathrm{av}}_4(2d+1)^2(\varepsilon')^3\sigma l\\
&=(2d+1)(\varepsilon')^2l\,\bigl(1+\mathsf{C}^{\mathrm{av}}_4s\bigr)\le2(2d+1)(\varepsilon')^2l .
\end{align*}

Fourth, the logarithm. The hypotheses hold at $c$ and at $\tilde c$, so (P2) applies at both bonds and
gives $\|Z\|,\|\tilde Z\|\le0.8\,s$, and $0.8\,s\le0.2$. With this common bound, (C3) of
Lemma~\ref{4.26:lem-exponential-products} gives
\begin{equation*}
\|R_{\log}(Z)-R_{\log}(\tilde Z)\|\le1.6\,s\cdot2(2d+1)(\varepsilon')^2l
=3.2\,(2d+1)^2(\varepsilon')^3\sigma l .
\end{equation*}
By the proof of (P2), $E(c)=(i\varepsilon')^{-1}[\mathcal{E}(c)+R_{\log}(Z)]$, and likewise at
$\tilde c$ with $\tilde Z$. Hence, using $(2d+1)^2\le6.25\,d^2$ and $3.2\cdot6.25=20$,
\begin{align*}
\|E(c)-E(\tilde c)\|
&\le(\varepsilon')^{-1}\bigl(\mathsf{C}^{\mathrm{av}}_4+3.2\bigr)(2d+1)^2(\varepsilon')^3\sigma l\\
&\le\bigl(6.25\,\mathsf{C}^{\mathrm{av}}_4+20\bigr)d^2(\varepsilon')^2\sigma l
\le\mathsf{C}^{\mathrm{av}}_2(\varepsilon')^2\sigma l .
\end{align*}
This is (P3).

\emph{Proof of} (P4).
For each of the four bonds of $\partial P$, the region $R_{m+1}$ of the bond is the union of the two
$(m+1)$-blocks at its endpoints, hence contained in $R_{m+1}(P)$. So the hypotheses on $\|A\|$ and
$\|d_mA\|$ hold on each of them, and $A'$ and $E$ are defined on all four bonds by (P2).

Apply the lattice curl \eqref{4.26:eq-combs-linearised-average-a} at spacing $\varepsilon'$ to $A'$, and
identify the coarse bonds:
\begin{equation*}
\begin{aligned}
A'_\mu(y)&=A'(c),&\qquad A'_\mu(y+\varepsilon'e_\lambda)&=A'(c+\varepsilon'e_\lambda),\\
A'_\lambda(y)&=A'(c'),&\qquad A'_\lambda(y+\varepsilon'e_\mu)&=A'(c'+\varepsilon'e_\mu).
\end{aligned}
\end{equation*}
Substituting $A'=\mathcal{Q}A+E$ gives
\begin{equation*}
\begin{split}
(d_{m+1}A')_{\mu\lambda}(y)={}&(d_{m+1}\mathcal{Q}A)_{\mu\lambda}(y)\\
&+(\varepsilon')^{-1}\Bigl\{\bigl[E(c'+\varepsilon'e_\mu)-E(c')\bigr]
-\bigl[E(c+\varepsilon'e_\lambda)-E(c)\bigr]\Bigr\}.
\end{split}
\end{equation*}

By Lemma~\ref{4.26:lem-coarse-curl}, identity \eqref{4.26:eq-coarse-curl-a}, the first term equals
\begin{equation*}
L^{-2-d}\sum_\rho\sum_{p\subset P_\rho}(d_mA)_{\mu\lambda}(p).
\end{equation*}
There are $L^d$ offsets $\rho$ and $L^2$ plaquettes in each square $P_\rho$, so this is an average with
weights $L^{-2-d}$ of total mass one. It is taken over fine curls at level-$m$ plaquettes in
$R_{m+1}(P)$, so its norm is at most $\sup_{R_{m+1}(P)}\|d_mA\|$.

Each bracket is a difference of the kind estimated in (P3). The first is taken at the bond $c'$ in the
direction $\mu$, under the hypothesis with $l$ for the pair $(c',c'+\varepsilon'e_\mu)$. The second is
taken at the bond $c$ in the direction $\lambda$, under the hypothesis with $l$ for the pair
$(c,c+\varepsilon'e_\lambda)$. Each is therefore at most $\mathsf{C}^{\mathrm{av}}_2(\varepsilon')^2\sigma l$.
Multiplying by $(\varepsilon')^{-1}$ and adding the two brackets, we obtain
\begin{equation*}
\|d_{m+1}A'(P)\|\le\sup_{R_{m+1}(P)}\|d_mA\|+2\,\mathsf{C}^{\mathrm{av}}_2\varepsilon'\sigma l .
\qedhere
\end{equation*}
\end{proof}

\begin{proof}[Proof of the averaging lemma, continued: the induction step]
\label{4.26:prf-averaging-induction}
We keep the majorants $\sigma_m$, $f_*$, $l_*$, $f^{\sharp}_m$, $\hat{g}$, $s_m$ of
the corresponding display, together with the consequences of their definitions derived
at the beginning of this proof. For every $n\le\ell$ these are
\begin{equation}\label{4.26:eq-averaging-induction-1}
\sum_{j=1}^{n}\varepsilon_j\le\frac{L}{L-1}\,\varepsilon_n\le 2\varepsilon_n,
\qquad f^{\sharp}_n\le f_*,
\qquad \sigma_n\le l_*=(d+\tfrac12)\hat{s}.
\end{equation}
Moreover, for $n\le\ell-1$ the number $s=s_n$ satisfies the hypothesis
\eqref{4.26:eq-one-step-averaging-a} of Proposition~\ref{4.26:prop-one-step-averaging}.
For a function $F$ on the bonds of $\Omega_j$ we write $\|F\|_{\Omega_j}$ for the supremum of
$\|F(b)\|$ over the bonds $b$ of $\Omega_j$.

The base case $m=0$ has been established above. We now let $m\ge0$ with $m+1\le\ell$, assume
the induction hypothesis at level $m$, with its three bounds $(\sigma_m)$, $(f_m)$, $(l_m)$,
and prove it at level $m+1$.

\smallskip\noindent\textit{Application of the one-step estimate.} Fix $t\in[0,1]$ and
$c\in\Omega_{m+1}$. We apply Proposition~\ref{4.26:prop-one-step-averaging} at the step
$m\to m+1$. The data are $A:=A^{(m)}_t$, $(\sigma,f,l):=(\sigma_m,f^{\sharp}_m,l_*)$,
$\varepsilon=\varepsilon_m$ and $\varepsilon'=\varepsilon_{m+1}$. We first check its
hypotheses. By the domain statements established at the beginning of this proof, every
level-$m$ bond used by the average at $c$ and every level-$m$ bond of $R_{m+1}(c)$ lies in
$\Omega_m$. Likewise, every level-$m$ plaquette of $R_{m+1}(c)$, and of $R_{m+1}(P)$ for a
coarse plaquette $P$ with bonds in $\Omega_{m+1}$, has its four bonds in $\Omega_m$. On these
bonds and plaquettes the induction hypothesis at level $m$ gives three facts: that
$A^{(m)}_t$ is $\mathfrak{sl}(N,\mathbb{C})$-valued, the bound $\|A^{(m)}_t\|\le\sigma_m$ by
$(\sigma_m)$, and the bound $\|d_mA^{(m)}_t\|\le f^{\sharp}_m$ by $(f_m)$. The translation
hypothesis with $\varepsilon'=\varepsilon_{m+1}$ is $(l_m)$. This holds for translation pairs
$c,c+\varepsilon_{m+1}e_\lambda\in\Omega_{m+1}$ and for the two translation pairs of the
boundary of a coarse plaquette, since all the bonds involved lie in $\Omega_m$ by the domain
statements. The one-step parameter of Proposition~\ref{4.26:prop-one-step-averaging} is
$s=(2d+1)\varepsilon_{m+1}\sigma_m=s_m$. Since $m\le\ell-1$, it satisfies
\eqref{4.26:eq-one-step-averaging-a}.

By the induction hypothesis, $\exp(i\varepsilon_mA^{(m)}_t)=[M^{m}U'_t]$ on $\Omega_m$.
Moreover, $[M^{m+1}U'_t](c)$ is the one-step average at $c$ of the level-$m$ variables
$[M^{m}U'_t](\beta)$, $\beta$ ranging over the level-$m$ bonds of $R_{m+1}(c)$. Hence
\begin{equation*}
W_t(c):=[M^{m+1}U'_t](c)=M\bigl(\exp(i\varepsilon_mA^{(m)}_t)\bigr)(c).
\end{equation*}
By (P1) and (P2) of \eqref{4.26:eq-one-step-averaging-b} we have
$\|W_t(c)-1\|\le 0.8s_m\le 0.2$. Therefore
$A^{(m+1)}_t(c):=(i\varepsilon_{m+1})^{-1}\operatorname{Log}W_t(c)$ is defined, and it is the
field $A'(c)$ of (P2). Since $\exp\circ\operatorname{Log}$ is the identity on the ball
$\|W-1\|<1$, we get $W_t(c)=\exp(i\varepsilon_{m+1}A^{(m+1)}_t(c))$.

The same application of Proposition~\ref{4.26:prop-one-step-averaging} can be made at every
step $j-1\to j$ with $1\le j\le m+1$, because its hypotheses at that step are supplied by the
induction hypothesis at level $j-1\le m$. For $1\le j\le m+1$ put
$G_j:=A^{(j)}_t-S_jA^{(j-1)}_t$. The decomposition in (P2) at the step $j-1\to j$ writes $G_j$
as the flux term plus the remainder $E_j$ of that step, with $T_j$ the flux operator $T$ of
(P2) at spacing $\varepsilon_j$. By (P2),
$\|T_j(d_{j-1}A^{(j-1)}_t)\|\le t_\perp f^{\sharp}_{j-1}$, so the flux term has norm at most
$\varepsilon_jt_\perp f^{\sharp}_{j-1}$, and
$\|E_j\|\le\mathsf{C}^{\mathrm{av}}_1\varepsilon_j\sigma_{j-1}^2$. Using
\eqref{4.26:eq-averaging-induction-1} in the form $f^{\sharp}_{j-1}\le f_*$ and
$\sigma_{j-1}\le(d+\frac12)\hat{s}$, together with the definition of $\hat{g}$, we obtain the
bound on $G_j$ in the display below. Recall that
$S^{[n,j]}=S_n\circ S_{n-1}\circ\cdots\circ S_{j+1}$, with $S^{[n,n]}$ the identity, denotes the
composition of the straight-line averages. Substituting the identity
$A^{(j)}_t=S_jA^{(j-1)}_t+G_j$ successively for $j=m+1,m,\dots,1$, and using
$A^{(0)}_t=ta$, gives the representation of $A^{(m+1)}_t$ on $\Omega_{m+1}$ in the display
below. By the domain statements, $S^{[m+1,j]}$ maps functions on $\Omega_j$ to functions on
$\Omega_{m+1}$ without increasing sup norms, which gives the bound on $S^{[m+1,j]}G_j$.
Finally, by \eqref{4.26:eq-averaging-induction-1} with $n=m+1\le\ell$,
$\sum_{j=1}^{m+1}\varepsilon_j\le 2\varepsilon_{m+1}$. Altogether,
\begin{equation}\label{4.26:eq-averaging-induction-a}
\begin{gathered}
A^{(m+1)}_t=S_{m+1}A^{(m)}_t+G_{m+1},\qquad
G_j=\varepsilon_jT_j\bigl(d_{j-1}A^{(j-1)}_t\bigr)+E_j\quad(1\le j\le m+1),\\
\|G_j\|_{\Omega_j}\le\varepsilon_j\bigl(t_\perp f^{\sharp}_{j-1}
+\mathsf{C}^{\mathrm{av}}_1\sigma_{j-1}^2\bigr)\le\varepsilon_j\hat{g},\\
A^{(m+1)}_t=S^{[m+1,0]}(ta)+\sum_{j=1}^{m+1}S^{[m+1,j]}G_j\quad\text{on }\Omega_{m+1},\\
\|S^{[m+1,j]}G_j\|_{\Omega_{m+1}}\le\|G_j\|_{\Omega_j},\qquad
\sum_{j=1}^{m+1}\varepsilon_j\le 2\varepsilon_{m+1}.
\end{gathered}
\end{equation}

\smallskip\noindent\textit{Tracelessness.} Let $t'\in[0,1]$. By the induction hypothesis at
every $t'$, each $A^{(j)}_{t'}$ with $j\le m$ is $\mathfrak{sl}(N,\mathbb{C})$-valued. Hence
every comb transporter entering $W_{t'}(c)$ is a product of exponentials of traceless
matrices and lies in $\mathrm{SL}(N,\mathbb{C})$. By the normalisation axiom (M4) of
Definition~\ref{4.26:def-block-average-axioms}, $\mathcal{N}$ is $G^{c}$-valued at families of
elements of $G^{c}$. Therefore $W_{t'}(c)\in\mathrm{SL}(N,\mathbb{C})$.

The map $t'\mapsto W_{t'}(c)$ is continuous. Indeed, it is obtained from
$U'_{t'}=\exp(i\xi t'a)$ by finitely many products of exponentials, inverses and applications
of the holomorphic map $\mathcal{N}$. Each application of $\mathcal{N}$ takes place inside its
polydisc: at every level $j\le m$ and every $t'$, the induction hypothesis and (P1) give
$\|P_x-1\|\le 1.14s_j<r_{\mathcal{N}}$. By the induction hypothesis at every $t'$ and (P2),
$\|W_{t'}(c)-1\|\le 0.2$ for every $t'\in[0,1]$. Hence $t'\mapsto\operatorname{Log}W_{t'}(c)$
is continuous. Since $\det e^{X}=e^{N\operatorname{tr}X}$ for the normalised trace
$\operatorname{tr}$ (with $\operatorname{tr}\mathbf{1}=1$),
\begin{equation*}
\exp\bigl(N\operatorname{tr}\operatorname{Log}W_{t'}(c)\bigr)=\det W_{t'}(c)=1,
\end{equation*}
so $N\operatorname{tr}\operatorname{Log}W_{t'}(c)\in 2\pi i\mathbb{Z}$. A continuous function on
$[0,1]$ with values in $2\pi i\mathbb{Z}$ is constant. At $t'=0$ we have $U'_0=1$ and
$W_0(c)=[M^{m+1}(1)](c)=1$, so the constant is $N\operatorname{tr}\operatorname{Log}1=0$. Thus
$\operatorname{tr}A^{(m+1)}_t(c)=0$, and $A^{(m+1)}_t$ is $\mathfrak{sl}(N,\mathbb{C})$-valued.
No smallness depending on $N$ beyond $\|W-1\|\le0.2$ enters this argument; in particular no
condition of the form $\|\operatorname{Log}W\|<2\pi/N$ is used.

\smallskip\noindent\textit{The sup bound.} By the first hypothesis in
the corresponding display, $\|ta\|\le\hat{s}$, and the straight-line averages do not
increase sup norms. Hence $\|S^{[m+1,0]}(ta)\|_{\Omega_{m+1}}\le\hat{s}$. By
\eqref{4.26:eq-averaging-induction-a} and the bound on
$\sum_j\varepsilon_j$,
\begin{equation}\label{4.26:eq-averaging-induction-2}
\begin{aligned}
\bigl\|A^{(m+1)}_t-S^{[m+1,0]}(ta)\bigr\|_{\Omega_{m+1}}
&\le\hat{g}\sum_{j\le m+1}\varepsilon_j\le\varepsilon_{m+1}\frac{L}{L-1}\,\hat{g}\\
&\le\varepsilon_{m+1}\hat{s}\Bigl[(d-1)\bigl(1+\tfrac{1}{4d}\bigr)
+2\mathsf{C}^{\mathrm{av}}_1(d+\tfrac12)^2\hat{s}\Bigr]\\
&\le\varepsilon_{m+1}\hat{s}\bigl[(d-\tfrac34)+\tfrac14\bigr]=(d-\tfrac12)\,\varepsilon_{m+1}\hat{s}.
\end{aligned}
\end{equation}
The steps are as follows. The second line uses
$\hat{g}=t_\perp f_*+\mathsf{C}^{\mathrm{av}}_1(d+\frac12)^2\hat{s}^2$ and
$f_*=(2+\frac1{2d})\hat{s}$. Since $t_\perp=(d-1)(L-1)/(2L)$, we have
$\frac{L}{L-1}t_\perp=\frac{d-1}{2}$ exactly, and no power of $L$ survives in this term. Then
\begin{equation*}
\tfrac{d-1}{2}\bigl(2+\tfrac1{2d}\bigr)=(d-1)\bigl(1+\tfrac1{4d}\bigr),
\end{equation*}
and $\frac{L}{L-1}\le2$ for $L\ge2$. In the third line,
$(d-1)(1+\frac1{4d})=d-1+\frac{d-1}{4d}\le d-\frac34$. The bound
$2\mathsf{C}^{\mathrm{av}}_1(d+\frac12)^2\hat{s}\le\frac14$ is inequality (8e) of
the corresponding display.

Inequality \eqref{4.26:eq-averaging-induction-2} is assertion (i) of
the corresponding display at level $m+1$. Together with
$\|S^{[m+1,0]}(ta)\|\le\hat{s}$ it gives
$\|A^{(m+1)}_t\|\le\hat{s}(1+(d-\frac12)\varepsilon_{m+1})=\sigma_{m+1}$. This is
$(\sigma_{m+1})$, which is assertion (ii) of the corresponding display.

\smallskip\noindent\textit{The curl bound.} Let $P$ be a level-$(m+1)$ plaquette with its
bonds in $\Omega_{m+1}$. We apply (P4) of \eqref{4.26:eq-one-step-averaging-b}, whose
hypotheses on $R_{m+1}(P)$ are $(f_m)$, $(\sigma_m)$ and $(l_m)$, as checked above. This gives
\begin{equation*}
\begin{aligned}
\bigl\|(d_{m+1}A^{(m+1)}_t)(P)\bigr\|
&\le f^{\sharp}_m+2\mathsf{C}^{\mathrm{av}}_2\varepsilon_{m+1}\sigma_ml_*\\
&\le f^{\sharp}_m+2\mathsf{C}^{\mathrm{av}}_2\varepsilon_{m+1}(d+\tfrac12)^2\hat{s}^2
=f^{\sharp}_{m+1}\le f_*.
\end{aligned}
\end{equation*}
Here the second inequality uses $\sigma_m\le l_*=(d+\frac12)\hat{s}$ from
\eqref{4.26:eq-averaging-induction-1}. The equality is the definition of $f^{\sharp}_{m+1}$ in
the corresponding display. The last inequality is
\eqref{4.26:eq-averaging-induction-1} with $n=m+1\le\ell$. This proves $(f_{m+1})$ and the
first clause of assertion (iv) of the corresponding display.

\smallskip\noindent\textit{The translation bound.} Let $v:=\varepsilon_{m+2}e_\lambda$, and
let $b$ and $b+v$ lie in $\Omega_{m+1}$. Write $b=\langle y,y+\varepsilon_{m+1}e_\mu\rangle$.
By the representation in \eqref{4.26:eq-averaging-induction-a}, where $S^{[m+1,0]}$ is the
straight-line average over the $L^{m+1}$ fine bonds of $b$,
\begin{equation*}
\begin{aligned}
A^{(m+1)}_t(b+v)-A^{(m+1)}_t(b)
&=L^{-(m+1)}\sum_{k<L^{m+1}}t\bigl[a_\mu(y+v+k\xi e_\mu)-a_\mu(y+k\xi e_\mu)\bigr]\\
&\quad+\sum_{j=1}^{m+1}\bigl[S^{[m+1,j]}G_j(b+v)-S^{[m+1,j]}G_j(b)\bigr].
\end{aligned}
\end{equation*}
Consider first the fine differences. Since $v=L^{m+2}\xi e_\lambda$, each difference
$a_\mu(z+v)-a_\mu(z)$, $z=y+k\xi e_\mu$, telescopes into $L^{m+2}$ differences over fine
translates by $\xi e_\lambda$. Since $\Omega$ is a box, all the intermediate fine bonds lie in
$\Omega$. By the second hypothesis in the corresponding display, each of these
differences is at most $\xi\hat{s}$ in norm, and $t\le1$. Hence each fine difference is at
most $L^{m+2}\xi\hat{s}=\varepsilon_{m+2}\hat{s}$, and so is their average.

Each term of the second sum is at most
$2\|S^{[m+1,j]}G_j\|_{\Omega_{m+1}}\le2\|G_j\|_{\Omega_j}\le2\varepsilon_j\hat{g}$, by
\eqref{4.26:eq-averaging-induction-a}. Summing over $j$ and using
$\sum_{j\le m+1}\varepsilon_j\le\frac{L}{L-1}\varepsilon_{m+1}$ and
$\varepsilon_{m+1}L=\varepsilon_{m+2}$, we obtain
\begin{equation}\label{4.26:eq-averaging-induction-3}
\begin{aligned}
\bigl\|A^{(m+1)}_t(b+v)-A^{(m+1)}_t(b)\bigr\|
&\le\varepsilon_{m+2}\Bigl[\hat{s}+\tfrac{2}{L-1}\,\hat{g}\Bigr]\\
&\le\varepsilon_{m+2}\hat{s}\Bigl[1+\frac{(d-1)(2+\frac1{2d})}{L}
+2\mathsf{C}^{\mathrm{av}}_1(d+\tfrac12)^2\hat{s}\Bigr]\\
&\le\varepsilon_{m+2}\hat{s}\bigl[1+(d-\tfrac34)+\tfrac14\bigr]
=(d+\tfrac12)\,\hat{s}\,\varepsilon_{m+2}.
\end{aligned}
\end{equation}
The second line uses the definitions of $\hat{g}$ and $f_*$. It also uses
$\frac{2}{L-1}t_\perp=\frac{d-1}{L}$, which follows from $t_\perp=(d-1)(L-1)/(2L)$, and
$\frac{2}{L-1}\le2$ for $L\ge2$. In the third line, for $L\ge2$,
\begin{equation*}
\frac{(d-1)(2+\frac1{2d})}{L}\le(d-1)\bigl(1+\tfrac1{4d}\bigr)\le d-\tfrac34,
\end{equation*}
and again $2\mathsf{C}^{\mathrm{av}}_1(d+\frac12)^2\hat{s}\le\frac14$ by inequality (8e) of
the corresponding display. Since $(d+\frac12)\hat{s}=l_*$, inequality
\eqref{4.26:eq-averaging-induction-3} is $(l_{m+1})$.

\smallskip\noindent\textit{Conclusion.} We have shown that the induction hypothesis holds at
level $m+1$. By induction it holds at level $\ell$. It contains assertions (i), (ii) and
(iv) of the corresponding display, and the $\mathfrak{sl}(N,\mathbb{C})$-valuedness of
$A^{(m)}_t$, at every $m\le\ell$ and every $t\in[0,1]$.

It remains to prove assertion (iii). Take $t=1$, so that $U'_1=U'$, and write
$A^{(m)}=A^{(m)}_1$. Then $[M^{m}U'](b)=\exp(i\varepsilon_mA^{(m)}(b))$ and
$[M^{m}U'](b)^{-1}=\exp(-i\varepsilon_mA^{(m)}(b))$. We use $\|e^{\pm X}-1\|\le\|X\|e^{\|X\|}$,
the bound $(\sigma_m)$, and
$\varepsilon_m\sigma_m\le\sigma_m\le(d+\frac12)\hat{s}$, which holds since
$\varepsilon_m\le1$. This gives
\begin{equation*}
\bigl\|[M^{m}U'](b)^{\pm1}-1\bigr\|
\le\varepsilon_m\sigma_m e^{\varepsilon_m\sigma_m}
\le\varepsilon_m\hat{s}\bigl(1+(d-\tfrac12)\varepsilon_m\bigr)e^{(d+\frac12)\hat{s}}
=c_m\varepsilon_m\hat{s},
\end{equation*}
where $c_m=(1+(d-\frac12)\varepsilon_m)e^{(d+\frac12)\hat{s}}$ is the closeness constant of
assertion (iii). The bound $c_m\le d+1$ is inequality (8g) of
the corresponding display. This is assertion (iii), and the proof of
the averaging lemma is complete.
\end{proof}

\begin{corollary}[The tree products near the identity]\label{4.26:cor-tree-product-bound}
Assume the hypotheses of the averaging lemma. Let
$1\le n\le\ell$ and let $x\in\Lambda_0$ with $B_n(x)\subset\Omega$. Then every bond of every
path $\tau_m(x)$, $m<n$, lies in $\Omega_m$: for such a bond $b$ the region $R_m(b)$ lies in the
$(m+1)$-block of $x$ inside $B_n(x)$, since the gauge-tree path $\tau_m(x)$ lies in that block.
Let
\begin{equation*}
P_n(x)=\prod_{m=n-1}^{0}\,[M^{m}U'](\tau_m(x))
\end{equation*}
be the tree product of Lemma~\ref{4.26:lem-tree-product}, as in \eqref{4.26:eq-tree-product-a}, and put
\begin{equation*}
Y(x):=\sum_{m<n}\ \sum_{b\in\tau_m(x)}\varepsilon_m\,\|A^{(m)}(b)\|,
\qquad
K(d,L):=d\Bigl[1+\frac{d-\frac12}{L+1}\Bigr]\le\frac{d(2d+5)}{6}\le d(d+1).
\end{equation*}
Then for every such $x$ and every $z\in\Lambda_0\cap B_n(x)$
\begin{equation}\label{4.26:eq-tree-product-bound-a}
\begin{aligned}
\bigl\|P_n(x)^{\pm1}-1\bigr\|
&\le e^{Y(x)}-1\le\exp\bigl(K(d,L)\,\hat{s}\,L^{n-\ell}\bigr)-1\\
&\le\exp\bigl(d(d+1)\,\hat{s}\,L^{n-\ell}\bigr)-1,\\
\bigl\|\bigl(P_n(z)^{-1}P_n(x)\bigr)^{\pm1}-1\bigr\|
&\le e^{Y(z)+Y(x)}-1\le\exp\bigl(2d(d+1)\,\hat{s}\,L^{n-\ell}\bigr)-1.
\end{aligned}
\end{equation}
For $n=0$ one has $P_0\equiv1$.
\end{corollary}

\begin{proof}
By its definition in \eqref{4.26:eq-tree-product-a}, $P_n(x)$ is the ordered product, over
$m=n-1,\dots,0$ and along each path $\tau_m(x)$, of the factors $[M^{m}U'](b)^{\pm1}$,
$b\in\tau_m(x)$; each such factor is $e^{\pm i\varepsilon_m A^{(m)}(b)}$, with $A^{(m)}$ the
level-$m$ field of the averaging lemma. Hence $P_n(x)$ is an ordered product of factors
$e^{X_j}$ with $X_j=\pm i\varepsilon_mA^{(m)}(b)$, $b\in\tau_m(x)$, $m<n$, and $P_n(x)^{-1}$ is
the reversed product of the factors $e^{-X_j}$. Since
$\|X_j\|=\varepsilon_m\|A^{(m)}(b)\|$, the sum $\sum_j\|X_j\|$ equals $Y(x)$. The first bound of
\eqref{4.26:eq-exponential-products-a} in Lemma~\ref{4.26:lem-exponential-products}, applied with
$s=Y(x)$ (it holds for the product of the $e^{X_j}$ and for the reversed product of the
$e^{-X_j}$), gives
\begin{equation*}
\bigl\|P_n(x)^{\pm1}-1\bigr\|\le e^{Y(x)}-1 .
\end{equation*}

We bound $Y(x)$. Every bond of $\tau_m(x)$, $m<n$, lies in $\Omega_m$, as stated in the corollary,
so the bound (ii) of the corresponding display gives $\|A^{(m)}(b)\|\le\sigma_m$ for these
bonds; and by the corresponding display the path $\tau_m(x)$ has $|\tau_m(x)|\le d(L-1)$
bonds. With $\varepsilon_m=L^{m-\ell}$ and the value of $\sigma_m$ in (ii) this yields
\begin{equation*}
\begin{aligned}
Y(x)&\le\sum_{m<n}d(L-1)\,\varepsilon_m\sigma_m\\
&=d\hat{s}\Bigl[(L-1)\sum_{m<n}L^{m-\ell}
+\Bigl(d-\tfrac12\Bigr)(L-1)\sum_{m<n}L^{2(m-\ell)}\Bigr].
\end{aligned}
\end{equation*}
The two geometric sums are
\begin{equation*}
(L-1)\sum_{m=0}^{n-1}L^{m-\ell}=L^{-\ell}(L^{n}-1)\le L^{n-\ell},
\end{equation*}
and, using $L^{2}-1=(L-1)(L+1)$,
\begin{equation*}
(L-1)\sum_{m=0}^{n-1}L^{2(m-\ell)}=\frac{L^{-2\ell}(L^{2n}-1)}{L+1}\le\frac{L^{2(n-\ell)}}{L+1}.
\end{equation*}
Since $n\le\ell$, we have $L^{n-\ell}\le1$, hence $L^{2(n-\ell)}\le L^{n-\ell}$, and therefore
\begin{equation*}
Y(x)\le d\hat{s}L^{n-\ell}\Bigl[1+\frac{d-\frac12}{L+1}\Bigr]=K(d,L)\,\hat{s}\,L^{n-\ell}.
\end{equation*}
Since $L\ge2$, $L+1\ge3$, so
\begin{equation*}
K(d,L)\le d\Bigl[1+\frac{d-\frac12}{3}\Bigr]=\frac{d(2d+5)}{6}\le d(d+1),
\end{equation*}
the last inequality because $2d+5\le6d+6$. As $t\mapsto e^{t}-1$ is increasing, the first line of
\eqref{4.26:eq-tree-product-bound-a} follows.

For the second line let $z\in\Lambda_0\cap B_n(x)$. Then $B_n(z)=B_n(x)\subset\Omega$, so the first
line holds at $z$ as well. For matrices $G,H$ the identity $GH-1=(G-1)(H-1)+(G-1)+(H-1)$ and the
submultiplicativity of the operator norm give
\begin{equation*}
\|GH-1\|\le\bigl(1+\|G-1\|\bigr)\bigl(1+\|H-1\|\bigr)-1 .
\end{equation*}
We apply this with $G=P_n(z)^{-1}$, $H=P_n(x)$, and, for the inverse
$\bigl(P_n(z)^{-1}P_n(x)\bigr)^{-1}=P_n(x)^{-1}P_n(z)$, with $G=P_n(x)^{-1}$, $H=P_n(z)$.
Inserting the first line at $z$ and at $x$, in both cases the right side is at most
$e^{Y(z)}e^{Y(x)}-1=e^{Y(z)+Y(x)}-1$. Finally, by the bounds just proved,
\begin{equation*}
Y(z)+Y(x)\le2\,d(d+1)\,\hat{s}\,L^{n-\ell},
\end{equation*}
which completes the proof of \eqref{4.26:eq-tree-product-bound-a}.
\end{proof}

\begin{remark}[Ba{\l}aban's spaces for partially prepared large-field components
{\cite[(1.63)--(1.69), pp.~190--191]{B-CMP122a}}]
\label{4.27:rem-partial-preparation-spaces}
The spaces described in this remark are those of
\cite[(1.63)--(1.69), pp.~190--191]{B-CMP122a}. They are restated in the present notation, with
Ba{\l}aban's weight base written $\mathsf{L}_0$ to avoid a clash with the constant
$L_0=L_0(N)$ of Theorem~0.

\emph{Indices and constants.} Fix a level $k$ and a component of the large-field region $Z$ at
level $k$ in which no new large-field regions were created during the preceding renormalisation
steps down to a level $h$. Thus $h$ is $k$ minus the number of these steps. The stage $n$ of the
integrations in the component corresponds to the level $j=h+n$, so that $n=j-h$. The number
$\eta$ is the lattice spacing of \cite[(1.63)]{B-CMP122a}, and $L^m\eta$ is the spacing at
level $m$. The integer $p$ appearing in $L^{-p}$ and in $L^{m-\min\{p,m\}}$ is Ba{\l}aban's scale
index of \cite[(1.63)]{B-CMP122a}; it is not a plaquette. The level $k_0$ is the one introduced
in \cite[\S1]{B-CMP122a}; powers of $\mathsf{L}_0$ enter the bounds below only at levels
$m>k_0$. The constant $\beta$ is the one of \cite[(1.63)]{B-CMP122a} that governs the factors
\eqref{4.27:eq-partial-preparation-spaces-3} below. Further, $\mathfrak{D}_m$ is Ba{\l}aban's
class of domains at level $m$, and $\alpha_{0,m}$, $\alpha_{1,m}$ are the radii at level $m$. The
sets $\Omega_m$, $\Omega''_m$, $Z_{j+1}$, $Z''_{j+1}$ are the domains of \cite[\S1]{B-CMP122a}.

\emph{Boundary terms and domains.} For the stage $n$ of the integrations in a large-field
component at level $k$, the boundary terms accumulated up to stage $n$ are
\begin{equation}\label{4.27:eq-partial-preparation-spaces-1}
  \mathbb{B}^{(n)}_k=\sum_{i\le n}\mathbb{C}^{(i)}_k ,
\end{equation}
where $\mathbb{C}^{(i)}_k$ denotes the new boundary terms produced at stage $i$. The terms
considered in \cite[(1.63)]{B-CMP122a} are those whose domains $X$ satisfy
\begin{equation*}
  X\in\mathfrak{D}_m,\qquad X\cap\bigl(Z_{j+1}\setminus Z''_{j+1}\bigr)\neq\emptyset .
\end{equation*}
Let $N_0$ be the smallest positive integer with $L^{-N_0+1}MR_{k-N_0+1}=M$
\cite[\S1]{B-CMP122a}, where the $R_{k'}$ are Ba{\l}aban's scales of \cite[(2.5)]{B-CMP119}.
Ba{\l}aban notes \cite[pp.~190--191]{B-CMP122a} that there are at most $N_0$ such terms.

\emph{The seven quantities and their bounds.} Write
\begin{align*}
  Q_1&=|\partial\mathbb{U}-1|, &
  Q_2&=|\partial U_{p,X}(M_{\cdot}(\mathbb{U}))-1|,\\
  Q_3&=|\mathbb{J}|, &
  Q_4&=|\mathbb{J}_{p,X}(M_{\cdot}(\mathbb{U}))|,\\
  Q_5&=|\partial U-1|, &
  Q_6&=|A'|, &
  Q_7&=|\nabla^{\eta}_{U}A'|,
\end{align*}
with all symbols as in \cite[(1.63)]{B-CMP122a}. For a level $m$, define seven numbers
$\varpi_1(m),\dots,\varpi_7(m)$ by
\begin{equation}\label{4.27:eq-partial-preparation-spaces-2}
\begin{gathered}
  \varpi_1(m)=\alpha_{0,m}\,\eta^2(L^m\eta)^{-2},\qquad
  \varpi_2(m)=\alpha_{0,m}\,L^{-2p}(L^mL^{-p})^{-2},\\
  \varpi_3(m)=\alpha_{0,m}\,(L^m\eta)^{-3},\qquad
  \varpi_4(m)=\alpha_{0,m}\,L^{m-\min\{p,m\}}(L^mL^{-p})^{-3},\\
  \varpi_5(m)=\alpha_{0,m}\,\eta^2(L^m\eta)^{-2},\qquad
  \varpi_6(m)=\alpha_{1,m}\,(L^m\eta)^{-1},\qquad
  \varpi_7(m)=\alpha_{1,m}\,(L^m\eta)^{-2}.
\end{gathered}
\end{equation}
For a level $l$, put
\begin{equation}\label{4.27:eq-partial-preparation-spaces-3}
  \Phi(l)=1-\beta\sum_{i=h+1}^{j}2^{-|l-i|}-\beta\sum_{q=l+1}^{k}2^{-(q-l)} .
\end{equation}
Here $l$ is the level at which the sums are centred; only the values $\Phi(m)$ and $\Phi(j)$
occur below. The space at stage $n=j-h$ consists of the configurations on which, for each
$a\in\{1,\dots,7\}$, the following three groups of strict bounds hold.
\begin{itemize}
\item[(i)] For $m=1,\dots,j-1$ on $(\Omega''_m\setminus\Omega''_{m+1})\cap X$, and for $m=j$ on
$(\Omega''_j\setminus\Omega_{k_0+1})\cap X$,
\begin{equation*}
  Q_a<\Phi(m)\,\mathsf{L}_0^{2\max\{0,m-k_0\}}\,\varpi_a(m).
\end{equation*}
\item[(ii)] For $m=k_0+1,\dots,j-1,j$ on $(\Omega_m\setminus\Omega_{m+1})\cap X$, as in
\cite[(1.66)]{B-CMP122a},
\begin{equation*}
  Q_a<\Phi(j)\,\mathsf{L}_0^{2(j-m)}\,\varpi_a(j).
\end{equation*}
That is, the bracket is the one at level $j$.
\item[(iii)] For $m=j+1,\dots,k-1$ on $(\Omega_m\setminus\Omega_{m+1})\cap X$, and for $m=k$ on
$\Omega_k\cap X$, as in \cite[(1.68)]{B-CMP122a},
\begin{equation*}
  Q_a<\Phi(m)\,c\,\varpi_a(m),
\end{equation*}
where $c=1$ for $m<k$ and $c=3$ for $m=k$.
\end{itemize}
The companion conditions on cubes, \cite[(1.65), (1.67), (1.69)]{B-CMP122a}, carry respectively
the weights $\mathsf{L}_0^{2\max\{0,m-k_0\}}$, $\mathsf{L}_0^{2(j-m)}$ and $1$, and no factor of
the kind \eqref{4.27:eq-partial-preparation-spaces-3}; we do not restate them.

\emph{Degenerate cases and monotonicity.} For $j\le k_0$ the bounds in (i) carry no powers of
$\mathsf{L}_0$, since then $\max\{0,m-k_0\}=0$ for $m\le j$. In the same case the range of (ii) is
empty. For $n=j-h=0$ the space coincides with the space
$\tilde{U}^{c}_k(X,\tilde\alpha_0,\tilde\alpha_1)$ of \cite[pp.~190--191]{B-CMP122a}. The spaces
form a descending sequence for increasing $n$, under conditions stated in
\cite[pp.~190--191]{B-CMP122a} which include
\begin{equation*}
  \beta\le\tfrac14,\qquad \mathsf{L}_0^{2}\le\tfrac13 L .
\end{equation*}

\emph{Components in which the integrations stop.} The case of a component in which the stage
integrations end early, because a new large field appears, is treated in \cite[p.~191]{B-CMP122a}
by the following sentence:
\begin{quotation}
Finally let us remark that if in a component of $Z$ we finish the integrations for some $n$,
because some new large field has been introduced, then we leave the new terms
$\mathbb{B}^{(n)}_k$ as a part of all boundary terms, and it is easy to generalize the definition
of the spaces in such a case for the next steps.
\end{quotation}
Beyond this sentence, the spaces for the subsequent steps in such a component are not written out
in \cite[pp.~190--191]{B-CMP122a}.
\end{remark}

\begin{definition}[Factors, radii and the logarithmic coupling variable]\label{4.27:not-factors-radii-coupling}
The following notation is fixed for the rest of this section. The level of the step is $k$; the letters $m$, $\mu$, $j$, $q'$ denote levels. The summation index $i$ in \eqref{4.27:eq-factors-radii-coupling-3} runs over levels.

\emph{Factors.} Ba{\l}aban's constant $\beta$ in the factors of the spaces of
Remark~\ref{4.27:rem-partial-preparation-spaces} is fixed at $\beta:=1/5$. For $m\le k$ put
\begin{equation}\label{4.27:eq-factors-radii-coupling-1}
s^{(k)}_m:=1-\beta\bigl(1-2^{-(k-m)}\bigr),
\end{equation}
as in \cite[(2.34)--(2.39)]{B-CMP119}. Let $N_k:=\check{R}_k$ be the declared number of stages
of the step at level $k$. Put $h:=k-N_k$ (the integer $h$ of the factors in
\cite[(1.64)--(1.69)]{B-CMP122a}) and $k_0:=k-N_0(k)$, where $N_0(k)$ is
the value at level $k$ of Ba{\l}aban's number $N_0$ of \cite[\S1]{B-CMP122a}. For $n\ge 0$ set
\begin{equation}\label{4.27:eq-factors-radii-coupling-3}
\varphi^{(n),(k)}_m:=s^{(k)}_m-\beta\sum_{i=h+1}^{h+n}2^{-|m-i|}.
\end{equation}
This is the factor of the spaces of Remark~\ref{4.27:rem-partial-preparation-spaces} in each of their three regions, with $n=j-h$, $j$ being the index of \cite[(1.64)--(1.69)]{B-CMP122a}, and $m$ being the level at which the factor is evaluated. Stage $l$ of the step integrates level $\ell_l:=h+l-1$. The one-stage room at level $m$ is
\begin{equation}\label{4.27:eq-factors-radii-coupling-4}
\rho^{(l)}_m:=\varphi^{(l-1),(k)}_m-\varphi^{(l),(k)}_m=\beta\,2^{-|m-\ell_l-1|}.
\end{equation}

\emph{Entries and radii.} The indices $a=1,\dots,7$ number the seven quantities bounded in the spaces of Remark~\ref{4.27:rem-partial-preparation-spaces}, in the order listed there. The index $a=8$ stands for the cube clauses \cite[(1.65), (1.67), (1.69)]{B-CMP122a}. For $a=1,\dots,7$ the radius $r_a(m)$ of entry $a$ at level $m$ is the $a$-th of the seven bounds in the bracket of those spaces at level $m$, with $\eta$ and the integer $p$ as there:
\begin{align*}
r_1(m)&=\alpha_{0,m}\eta^2(L^m\eta)^{-2}, & r_2(m)&=\alpha_{0,m}L^{-2p}(L^mL^{-p})^{-2},\\
r_3(m)&=\alpha_{0,m}(L^m\eta)^{-3}, & r_4(m)&=\alpha_{0,m}L^{m-\min\{p,m\}}(L^mL^{-p})^{-3},\\
r_5(m)&=\alpha_{0,m}\eta^2(L^m\eta)^{-2}, & r_6(m)&=\alpha_{1,m}(L^m\eta)^{-1},\\
r_7(m)&=\alpha_{1,m}(L^m\eta)^{-2}. & &
\end{align*}
Put $(e_1,\dots,e_7):=(2,2,3,2,2,1,2)$, so that $e_a\in\{1,2,3\}$ and $e_6=1$. For a given $a$ write $\alpha_{\cdot,m}$ for $\alpha_{0,m}$ if $a\le 5$ and for $\alpha_{1,m}$ if $a\in\{6,7\}$. Multiplying out the powers of $L$ and $\eta$, for $a\ne 4$
\begin{equation*}
r_a(m)=\eta^{-f_a}\,\alpha_{\cdot,m}L^{-e_am},
\end{equation*}
with the exponents, independent of $m$,
\begin{equation*}
(f_1,f_2,f_3,f_5,f_6,f_7):=(0,0,3,0,1,2),
\end{equation*}
while for $a=4$
\begin{equation*}
r_4(m)=\alpha_{0,m}L^{3p-\min\{p,m\}-2m},
\end{equation*}
which is $\alpha_{0,m}L^{2p}L^{-2m}$ for $m\ge p$ and $\alpha_{0,m}L^{3p}L^{-3m}$ for $m\le p$.

\emph{The logarithmic coupling variable.} Put
\begin{equation*}
\mathsf{x}_j:=g_j^{-2},\qquad \mathsf{t}_j:=\log\mathsf{x}_j,\qquad
\mathsf{t}_\gamma:=\log\gamma^{-2}.
\end{equation*}
The radii are $\alpha_{\iota,m}=\alpha_\iota(\mathsf{t}_m)$ for $\iota=0,1$. For $\mathsf{t}>0$ the functions $\alpha_\iota$ are given by
\begin{equation}\label{4.27:eq-factors-radii-coupling-5}
\alpha_\iota(\mathsf{t}):=C^{\mathrm{rad}}_\iota e^{-\mathsf{t}/2}\mathsf{t}^{q},\qquad
\mathsf{A}(\mathsf{t}):=\max_\iota\alpha_\iota(\mathsf{t})=C^{\mathrm{rad}}e^{-\mathsf{t}/2}\mathsf{t}^{q}.
\end{equation}
Here $C^{\mathrm{rad}}_0,C^{\mathrm{rad}}_1>0$ are the constant prefactors of Ba{\l}aban's radii $\alpha_{0,m},\alpha_{1,m}$ (the radii of Lemma~\ref{4.4:lem-twist-stability}), $C^{\mathrm{rad}}:=\max_\iota C^{\mathrm{rad}}_\iota$, and $q$ is the auxiliary parameter of that name. With $\beta_0>0$ the constant in the first relation of \cite[(2.6)]{B-CMP119} and in the first relation of \cite[(2.7)]{B-CMP119}, put $c_\beta:=2\ln(1+\beta_0)$.

\emph{Ratio of radii.} Let $\mu<m$ be levels, and write $g:=m-\mu\ge 1$ for their difference; the letter $g$ in this paragraph and in the proof of the ratio bound below, and the subscript of $\chi_g$ wherever it occurs, denote a difference of levels, never a coupling. Then for $a=1,\dots,7$
\begin{equation}\label{4.27:eq-factors-radii-coupling-2}
\frac{r_a(\mu)}{r_a(m)}\ge\chi_g\,L^{e_a g},
\end{equation}
where
\begin{equation}\label{4.27:eq-factors-radii-coupling-a}
\chi_g:=(1+\beta_0)^{-2}g^{-1/2}.
\end{equation}

\emph{Stage spaces.} $\mathfrak{S}^{(k)}_n(X)$ denotes the space of Remark~\ref{4.27:rem-partial-preparation-spaces}, that is, \cite[(1.64)--(1.69)]{B-CMP122a} at stage $n$, with the following two modifications, both in force.
\begin{itemize}
\item[$(\Sigma^{\flat})$, $(\Sigma^{*})$] The bounds on the second and fourth entries are required for every triple $(p,X',\mathfrak{s})$ in which $\mathfrak{s}$ is a determining-set structure on $X'$, namely the current one or one of its predecessors in the construction; here $p$ and $X'$ range over the indices of the maps $\partial U_{p,X'}$ and $\mathbb{J}_{p,X'}$ bounded in the second and fourth entries of the spaces of Remark~\ref{4.27:rem-partial-preparation-spaces}.
\item[$(c\mapsto 1)$] The constant $c$ of the third region is replaced by $1$.
\end{itemize}
Later statements of this section refer to the first modification as the readings $(\Sigma^{\flat})$ and $(\Sigma^{*})$ of the stage spaces, and to the second under the name $(c\mapsto 1)$.

\emph{Localisation quantities.} $d^{\wedge}$ denotes the scaled distance, and $\mathsf{w}_M:=M/M_1$ the width ratio, of the localisation estimates of this chapter. Put
\begin{equation*}
\bar{R}(\mathsf{x}):=L(2\log\mathsf{x})^{r},\qquad
\mathsf{l}(\mathsf{x}):=\bigl(L^2\bar{R}(\mathsf{x})\bigr)^2,\qquad
\varepsilon^{\mathsf{l}}:=\mathsf{l}\,\varepsilon,
\end{equation*}
where $\varepsilon$ is the small parameter of the localisation bounds of this chapter, and $r$, one of the auxiliary parameters $\mathfrak{p}$, is the exponent in Ba{\l}aban's length parameter $R_m$ at level $m$. The function $\bar{R}$ dominates Ba{\l}aban's length parameter: $\bar{R}(\mathsf{x}_m)\ge R_m$.

\emph{Regions.} At the site of the step where the operation $\mathbf{R}$ is applied, $Z$ denotes the union of all components of the first two classes of this chapter's domain convention present there. Let $\Omega_k$ be the small-field region at level $k$, and put
\begin{equation*}
\mathfrak{T}_k:=\Omega_k\cup Z^{c},\qquad Z^{\mathrm{on}}_k:=Z\cap\Omega_k^{c}.
\end{equation*}
A stored term is a term of the expansion frozen when a run of the operation $\mathbf{R}$ is arrested; such runs are described under \emph{Attempts} below. The domain of a stored term of stage $n$ on $X$ is
\begin{equation}\label{4.27:eq-factors-radii-coupling-6}
\mathfrak{D}^{\flat}_n(X):=\mathfrak{N}^{(n)}(X)\cap\mathfrak{L}_n\bigl(X\cap Z^{\mathrm{on}}\bigr)
\cap\mathfrak{I}_n\bigl(X\cap\mathfrak{T}\bigr)
\end{equation}
(the level subscript $k$ on $Z^{\mathrm{on}}$ and $\mathfrak{T}$ is suppressed).
Here $\mathfrak{N}^{(n)}$ is the ambient space of pairs at stage $n$, and $\mathfrak{L}_n$ and $\mathfrak{I}_n$ are the class part and the foreign part of the stage-$n$ domain.

\emph{Units.} At a point $x$ the unit for entry $a$ is $\mathsf{w}(x)\,r_a(\ell(x))$. Here $\ell(x)$ is the level at which the spaces of Remark~\ref{4.27:rem-partial-preparation-spaces} read their bounds at the point $x$, and $\mathsf{w}(x)$ is the weight attached to $x$ by the same spaces.

\emph{Attempts.} An attempt $\mathfrak{b}$ at level $k_{\mathfrak{b}}$ is the run of the operation $\mathbf{R}$, at the site of that step, on a component $X_{\mathfrak{b}}$ of one of the two classes above, whether the run is arrested or completed. We set $W_{\mathfrak{b}}:=X_{\mathfrak{b}}\cap\Omega^{c}_{k_{\mathfrak{b}}}$.
\end{definition}

\begin{proof}
\emph{The factor.} For $m\le k$ the geometric sum gives
\begin{equation*}
\sum_{q'=m+1}^{k}2^{-(q'-m)}=1-2^{-(k-m)}.
\end{equation*}
Inserting \eqref{4.27:eq-factors-radii-coupling-1} into \eqref{4.27:eq-factors-radii-coupling-3} gives
\begin{equation*}
\varphi^{(n),(k)}_m=1-\beta\sum_{i=h+1}^{h+n}2^{-|m-i|}-\beta\sum_{q'=m+1}^{k}2^{-(q'-m)}.
\end{equation*}
The right-hand side is the factor written in the spaces of Remark~\ref{4.27:rem-partial-preparation-spaces}, with $j=h+n$, evaluated at level $m$.

\emph{The one-stage room.} By \eqref{4.27:eq-factors-radii-coupling-3}, the sums defining $\varphi^{(l-1),(k)}_m$ and $\varphi^{(l),(k)}_m$ run over $i=h+1,\dots,h+l-1$ and over $i=h+1,\dots,h+l$ respectively. The term $s^{(k)}_m$ is common to both. The difference is therefore the single term $\beta\,2^{-|m-(h+l)|}$. Since $h+l=\ell_l+1$, this term is the right-hand side of \eqref{4.27:eq-factors-radii-coupling-4}.

\emph{The maximal radius.} For $\mathsf{t}>0$ the factor $e^{-\mathsf{t}/2}\mathsf{t}^{q}$ in \eqref{4.27:eq-factors-radii-coupling-5} is positive and does not depend on $\iota$. The maximum over $\iota$ is therefore taken in the constants, which gives the second equality there.

\emph{The ratio of radii.} Two inputs are used: the ratio of radii along the coupling flow (Lemma~\ref{4.4:lem-radii-ratio}), and the first relation of \cite[(2.6)]{B-CMP119} together with the first relation of \cite[(2.7)]{B-CMP119}, with the two indices of \cite[(2.6), (2.7)]{B-CMP119} set equal. Multiplying these two relations and combining the product with Lemma~\ref{4.4:lem-radii-ratio} bounds the radius at the later level $m$ by the radius at the earlier level $\mu=m-g$:
\begin{equation}\label{4.27:eq-factors-radii-coupling-7}
\alpha_{\cdot,m}\le(1+\beta_0)^2g^{1/2}\alpha_{\cdot,\mu}.
\end{equation}
Fix $a\in\{1,2,3,5,6,7\}$. By the formulas for the radii in the statement, with the exponent $f_a$ given there, which does not depend on the level,
\begin{equation*}
r_a(m)=\eta^{-f_a}\alpha_{\cdot,m}L^{-e_am}\quad\text{and}\quad r_a(\mu)=\eta^{-f_a}\alpha_{\cdot,\mu}L^{-e_a\mu}.
\end{equation*}
Both formulas involve the same radius $\alpha_{\cdot}$. Dividing, and using $m-\mu=g$ together with \eqref{4.27:eq-factors-radii-coupling-7},
\begin{equation*}
\frac{r_a(\mu)}{r_a(m)}=\frac{\alpha_{\cdot,\mu}}{\alpha_{\cdot,m}}\,L^{e_a(m-\mu)}
\ge(1+\beta_0)^{-2}g^{-1/2}L^{e_ag}.
\end{equation*}
Now let $a=4$, so that $e_4=2$ and $\alpha_{\cdot,m}=\alpha_{0,m}$. By the formula for $r_4$ in the statement,
\begin{equation*}
\frac{r_4(\mu)}{r_4(m)}=\frac{\alpha_{0,\mu}}{\alpha_{0,m}}\,
L^{2(m-\mu)+\min\{p,m\}-\min\{p,\mu\}}.
\end{equation*}
Since $\mu<m$, we have $\min\{p,m\}\ge\min\{p,\mu\}$, and since $L\ge1$ the factor $L^{\min\{p,m\}-\min\{p,\mu\}}$ is at least $1$. Together with $m-\mu=g$ and \eqref{4.27:eq-factors-radii-coupling-7},
\begin{equation*}
\frac{r_4(\mu)}{r_4(m)}\ge\frac{\alpha_{0,\mu}}{\alpha_{0,m}}\,L^{2g}
\ge(1+\beta_0)^{-2}g^{-1/2}L^{e_4g}.
\end{equation*}
By \eqref{4.27:eq-factors-radii-coupling-a}, in every case this is \eqref{4.27:eq-factors-radii-coupling-2}.
\end{proof}

\begin{definition}[Blocking data and one class per step]\label{4.27:def-one-class-per-step}
Let $\mathsf{t}_j=\log\mathsf{x}_j$, $\mathsf{x}_j=g_j^{-2}$, be the logarithmic coupling variable of
Definition~\ref{4.27:not-factors-radii-coupling}. Let $r$ be the auxiliary parameter of
$\mathfrak{p}$ that enters the one-sided condition $p_0\ge 5r$. The following data are declared,
and are in force for the remainder of this chapter.
\begin{enumerate}
\item[(a)] \emph{Number of stages.} A sequence $(\check{R}_j)_j$ is fixed which is non-increasing
in $j$ and satisfies
\begin{equation*}
\check{R}_{j+1}\le\check{R}_j,\qquad \check{R}_j<L\,\mathsf{t}_j^{\,r}\quad\text{for every }j.
\end{equation*}
At the step leading to level $k$, the number of stages is declared to be
\begin{equation}\label{4.27:eq-one-class-per-step-1}
N_k:=\check{R}_k,\qquad h:=k-N_k .
\end{equation}
Here $h$ is the level integrated by the first stage of the step; it is not a history.
This is a declaration and not a consequence of Ba{\l}aban's construction: there the number
of stages is the number $R_j$ attached to a large-field region in \cite{B-CMP122a}, with $j$ the
index of the step at which the region was created,
whereas \eqref{4.27:eq-one-class-per-step-1} fixes it by the step alone. Consequently every
step treats exactly one class of large-field components, namely the components handled with
the single pair $(h,k)$ of \eqref{4.27:eq-one-class-per-step-1}.
\item[(b)] \emph{Scope.} A step whose stages would have to be organised in two windows, that is,
with two distinct pairs $(h,k)$, is not covered by the present declarations.
\item[(c)] \emph{Schedule constant.} The schedule constant of the layered analyticity spaces
(Definition~\ref{4.4:def-post-step-space}) is $\beta=1/5$.
\item[(d)] \emph{Weight base.} The base $\mathsf{L}_0$ (distinct from the constant
$L_0=L_0(N)$) of the stage weights $\mathsf{L}_0^{2(j-m)}$ attached to level $m$ at step $j$
satisfies
\begin{equation*}
2\le\mathsf{L}_0,\qquad \mathsf{L}_0^{2}\le L/3 .
\end{equation*}
The correlation theorem requires in addition $\mathsf{L}_0\ge 3$, which is compatible with the
two conditions just declared.
\item[(e)] \emph{Two numerical constants.} The input constant $\theta_S$ lies in the interval
$[1/2,\,8/9]$, and the input constant $\mathsf{w}_0$ is $\mathsf{w}_0:=1/24$.
\item[(f)] \emph{Exponents of the radii.} The powers $q_0,q_1$ of $\mathsf{t}_m$ in the radii
$\alpha_{0,m},\alpha_{1,m}$ of Definition~\ref{4.27:not-factors-radii-coupling} are both taken
equal to the auxiliary parameter $q$: $q_0=q_1=q$.
With this choice, the strengthened condition on the radii is in force; it is taken by statement.
\item[(g)] \emph{Nesting.} A geometric declaration, together with its nesting condition, is taken
by statement and is in force.
\end{enumerate}
\end{definition}

The following statements are used as inputs. Throughout this section the step from level $k$ to level $k+1$ is fixed. The blocking data and the constants $\beta=\tfrac15$, $\theta_S$ and $\mathsf{w}_0=\tfrac1{24}$ are those of Definition~\ref{4.27:def-one-class-per-step}. The following are those of Definition~\ref{4.27:not-factors-radii-coupling}: the factors $s^{(k)}_m$; the radii $\alpha_{i,m}=\alpha_i(\mathsf{t}_m)$, with the exponent $q$ and the constants $C^{\mathrm{rad}}_i$; the constant $C=\max_iC^{\mathrm{rad}}_i$; and the logarithmic coupling variable $\mathsf{t}_j=\log g_j^{-2}$. The direct entries of a pair at a point $y$ are the quantities $Q_a(y)$, $a=1,\dots,7$. The $\mathbb{J}$-variable is measured in lattice units as $\ell^3|\mathbb{J}(y)|$, with $\ell=\eta$ at index zero.

\begin{enumerate}
\item[1.] \emph{The statement on presented pairs.} It holds at every determining set, and is used at the $\mathbb{J}$-radius $\gamma_0=\gamma_0^{(k)}$ of the lemma on weights, the dated size parameter and heredity (see \eqref{4.27:eq-weights-heredity-a}). The following parts of it are used.
\begin{itemize}
\item[(a)] A contour pays at least $\mathsf{w}_M$ for each stratum of $\mathfrak{B}$ it crosses, and for each $M$-cube of the local scale of $\mathfrak{B}$ it crosses. Hence, for points $y,y'$ of levels $m,m'$,
\begin{equation*}
d^{\wedge}(y,y')\ge \mathsf{w}_M\,\bigl(|m-m'|-1\bigr)^{+}.
\end{equation*}
\item[(b)] For every $n$ the stage map $\Phi_n\colon\mathfrak{S}_{n+1}(X)\to\mathfrak{S}_n(X)$ is holomorphic and leaves the $G$-valued part fixed.
\item[(c)] Suppose the datum $b$ of a stage layer satisfies $\|b\|_\infty\le r\,C_1\delta'_j$ with $1\le r\le r_*^{\sharp}$. Then $\Phi^{\chi}\in\mathfrak{S}_n(\hat X)$, and at most $\tfrac12\rho_m$ of the one-stage room is used.
\item[(d)] Clause (b), together with a layer bound of the form of the first inequality in \eqref{4.27:eq-t-step-layer-a}, applies on the bonds of $W_{\mathfrak{a}}$. By the clause of the arrest theorem that fixes a single determining-set structure for the source and the target of the map, every entry is compared with the same entry. The range sum $\sigma_{\mathfrak{a}}$ does not depend on $n$ and cancels. Hence the room of one step is $\rho_\ell$, of which clause (b) uses at most $\tfrac12\rho_\ell$.
\item[(e)] \emph{The rescaling bound.} Let $\mu$ be a majorant in the class $\mathfrak{M}^{\wedge}_{\delta/2}(\mathfrak{B})$ of decaying weights. Let $\mathfrak{h}$ be any layer with $N_{\cdot}(\mathfrak{h})\le\mu$, and assume $C^{b}(a_\flat+\kappa_{\max})\le1$. Take any admissible triple with domain $X'$, any bond $b'\in X'^{\flat}$ with location $x$ and level $m=m(b')$, and $e\in\{2,4\}$. Then
\begin{equation*}
Q^{+}_e(b')\le Q_e(b')\,(1+\pi_\times)+\pi\,\mathbf{r}_{e,m(b')},\qquad
\pi:=\frac{K^{b}_S\,\mu(x)}{\tilde\alpha_{0,m}},\qquad \pi_\times\le\pi .
\end{equation*}
Here $Q^+_e$ is the entry after the move. The constant $K^{b}_S$ contains the factor $c_B=e^{2\delta L(\varrho_0+2L)}$, which is one of the numbers fixed with $L$.
\end{itemize}
\item[2.] \emph{Descent of the stage spaces.} For every $n$, $\mathfrak{S}_{n+1}(X)\subset\mathfrak{S}_n(X)$. At a point raised at stage $n$, the ratio of the new threshold to the old one is at most $0.74$.
\item[3.] \emph{Inheritance.} Two clauses of the inheritance statement are used; they carry the letters of that statement.
\begin{itemize}
\item Clause (a): $\mathbb{C}^{(n)}_k(X,\mathcal{S})$ is a stored term on $\mathfrak{D}^{\flat}_n(X)\times\mathfrak{D}_{\mathcal{S}}(X)$, of norm at most $C_0^{\sharp}$.
\item Clause (c): with $n_W(y):=k-1-m(y)$,
\begin{equation*}
\sum_{l\le n}\iota_l(y)\le\tfrac12\,\mathsf{w}_0\,\beta\,e^{-n_W(y)},\qquad
\sum_{l}\iota^{1}_l\le 2\,\mathsf{w}_0\,\beta .
\end{equation*}
\end{itemize}
The room satisfies $r_T\ge\beta\theta_S2^{-(k-m)}$ for $m<k$, and $r_T\ge\beta$ on $\Omega_k$. If $\ell_T<m$, then $\theta_{T,a}\ge 2r_a$. Moreover, by the corollary on target thresholds,
\begin{equation*}
\theta^{(n)}_{T,a}-\theta_{\mathrm{src},a}\ge\tfrac12\,r_T\,r_a,
\end{equation*}
and $\theta^{(n)}_{T,a}\ge 1.5\,r_a$ at drops of scale.
\item[4.] \emph{Layers of the operation $\mathbf{R}$ off the zone.} The presented pair is the identity times layers on the $U$-part, with the cube gauges kept. The increments are linear in the layer. They are bounded by:
\begin{itemize}
\item $P\,2^{-(k-1-m)}$ on foreign shells of level $m\le k-1$;
\item $P'$ on $\Omega_k\setminus Z$;
\item $P''$ on collars.
\end{itemize}
Put $\bar P:=\max\{P,P',P''\}$. The scaled distance of this statement is $d^{\wedge}$. By Lemma~\ref{4.4:lem-radii-ratio}, the coarse-over-fine ratio is used in the form
\begin{equation*}
(1+\beta_0)^2(1+k-m)^{1/2}=(1+\beta_0)^2(2+t)^{1/2},\qquad t:=k-1-m .
\end{equation*}
Accordingly the constant $C_{\beta_0}$ of that statement is replaced by
\begin{equation*}
C_{1/2}:=\sup_{t\ge0}\,(2+t)^{1/2}e^{-3t/5}=\sqrt2 .
\end{equation*}
Indeed, the logarithmic derivative of $(2+t)^{1/2}e^{-3t/5}$ is
\begin{equation*}
\frac{1}{2(2+t)}-\frac35\le\frac14-\frac35<0\qquad(t\ge0),
\end{equation*}
so the supremum is the value $\sqrt2$ at $t=0$. The maximal increment satisfies $16\,\bar P\le\beta\theta_S$.
\item[5.] \emph{The $\mathbf{T}$-step from level $k$ to level $k+1$.} By clause (ii) of the stage-factor lemma, the source factor is $s^{(k+1)}_m+\theta_S\beta2^{-(k+1-m)}$ on the layer of level $m\le k$, and $1+\beta$ on $\Omega_{k+1}$. The layer of the step is controlled as follows. Let $C_H$ be the constant of the background-field layer theorem. Let $K_\chi$ be the constant of clause (iv) of the recursive presentation bounds. It is given by the first formula below, and we put
\begin{equation*}
K_\chi=\bigl(15d+K_H+5K_{J'}+2K_P\bigr)B_\sharp,\qquad
C_H^{\sharp}:=\max\{C_H,\,2K_\chi\}.
\end{equation*}
Here $B_\sharp$ is the constant of the background-field layer theorem; it is not the flow constant $B_\sharp$ of Definition~\ref{1.2:def-flow-constants}. The constants $K_H$, $K_{J'}$ and $K_P$ in this formula are constants of the recursive presentation bounds. Throughout input~5 we assume
\begin{equation}\label{4.27:eq-t-step-layer-2}
C_0^{\mathrm{rad}}\ge\frac{16(1+\beta_0)^2\,\|C\|\,C_1^{\mathrm{pd}}A_1\,C_H^{\sharp}}{(1-\theta_S)\beta},
\end{equation}
and the same lower bound for $C_1^{\mathrm{rad}}$. Here $\|C\|$ denotes the norm that enters the lower bound on $C_0^{\mathrm{rad}}$ in clause (ii) of the stage-factor lemma; it is not the constant $C=\max_iC^{\mathrm{rad}}_i$ of Definition~\ref{4.27:not-factors-radii-coupling}.
\end{enumerate}

\begin{lemma}[The $\mathbf{T}$-step layer in a background field. Stated; proof incomplete at the step marked $(\ast)$]\label{4.27:lem-t-step-layer}
Assume \eqref{4.27:eq-t-step-layer-2}. Let $\mathfrak{B}$ be the determining set of the step, frozen on the live pieces. Let $b$ be the datum of the operation $\mathbf{T}$ on its bonds of level $k$, and let $\mathcal{C}$ be the set of these bonds. Then $\mathcal{C}\subset\Omega_{k+1}$, since the set $S_{k+1}$ on which the operation acts satisfies $S_{k+1}\subset\Omega_{k+1}\cap\Lambda^{c}_{k+1}$ \cite{B-CMP119}. Assume $\sup|b|\le\varpi_k$, where $\varpi_k$ is the size prescribed for the datum of $\mathbf{T}$ at step $k$ in the stage-factor lemma, and let the datum be $\mathsf{d}=(0,b)$. Put
\begin{equation*}
p(y):=\varpi_k\,e^{-\delta d^{\wedge}(y,\mathcal{C})}.
\end{equation*}
Then $p\in\mathfrak{M}_\delta(\mathfrak{B})$, and the map of $\mathbf{T}$ on pairs is
\begin{equation*}
(\mathbb{U},\mathbb{J})\longmapsto\bigl(e^{i\eta\mathfrak{h}}\,\mathbb{U},\ \mathbb{J}+\mathcal{J}^{\sharp}[1]\bigr).
\end{equation*}
Here $\mathfrak{h}$ is the layer of the background-field layer theorem. It is taken in its Landau-gauge representative, as for the layers of the stage maps of input~1(b); Ba{\l}aban's axial representative differs from it by a gauge transformation. The layer and the source increment satisfy
\begin{equation}\label{4.27:eq-t-step-layer-a}
N^{\Delta}_y(\mathfrak{h})\le B_\sharp\,p(y),\qquad
\ell^{3}\,\bigl|\mathcal{J}^{\sharp}[1](y)\bigr|\le 2K_\chi\,p(y).\end{equation}
Every direct entry moves by at most $C_H^{\sharp}p(y)$, in its own units, with $C_H^{\sharp}$ as in input~5 above. In this lemma and in the lemma that bounds the operation $O3$, the constant $C_H$ of the background-field layer theorem is replaced by $C_H^{\sharp}$. Accordingly, the lower bound \eqref{4.27:eq-t-step-layer-2} is the lower bound on $C_0^{\mathrm{rad}}$ of clause (ii) of the stage-factor lemma with its constant $C_H$ replaced by $C_H^{\sharp}$.
\end{lemma}

\begin{proof}
We first check the hypotheses of the background-field layer theorem.

\emph{Membership.} The pair on which $\mathbf{T}$ acts is the pair of the source enlarged over $\mathbb{B}_{k+1}$. Since $\mathfrak{B}\preceq\mathbb{B}_{k+1}$, clause (e3) of the lemma on weights, the dated size parameter and heredity places it in $\mathfrak{K}_{109}(\mathfrak{B};\gamma_0^{(k)})$. The stratum $\Omega_0(\mathfrak{B})\setminus\Omega_1(\mathfrak{B})$ of raw points is treated separately. There the clause of the same lemma on points of index zero gives
\begin{equation*}
\mathfrak{K}_{109}\bigl(\mathfrak{B};\gamma_0^{(k+1)}\bigr)\subset\mathfrak{K}_{109}\bigl(\mathfrak{B};\gamma^{\top}\bigr),
\end{equation*}
and the pair lies in the smaller space. On the larger space the background-field layer theorem holds with the same constants, which do not depend on $\gamma$.

\emph{Size of the datum.} We need
\begin{equation*}
\|\mathsf{d}\|_1=\sup|b|\le\varpi_k\le\varepsilon_\sharp ,
\end{equation*}
where $\varepsilon_\sharp$ is the smallness threshold on the datum in the hypotheses of the background-field layer theorem. The lower bound \eqref{4.27:eq-t-step-layer-2} on $C_0^{\mathrm{rad}}$ gives
\begin{equation}\label{4.27:eq-t-step-layer-1}
C_H^{\sharp}\,\varpi_k\le\frac{(1-\theta_S)\beta}{8}\,\alpha_{0,k}\le 0.0125\,\alpha_{0,k}.
\end{equation}
The second inequality holds because $\theta_S\ge\tfrac12$ and $\beta=\tfrac15$, so that $(1-\theta_S)\beta/8\le\tfrac1{80}$. By the convergence condition on the radii, $\alpha_{0,k}\le\bar\alpha^{\mathrm{fat}}_k\le\varepsilon_\sharp$, where the radius $\bar\alpha^{\mathrm{fat}}_k$ is defined in input~10 below.

$(\ast)$ Passing from \eqref{4.27:eq-t-step-layer-1} to $\varpi_k\le\alpha_{0,k}\le\varepsilon_\sharp$ means dividing by $C_H^{\sharp}$, and this requires $C_H^{\sharp}\ge\tfrac1{80}$. The definitions give, in dimension four,
\begin{equation*}
C_H^{\sharp}\ge 2K_\chi\ge 120\,B_\sharp ,
\end{equation*}
but $B_\sharp$ is only known to be positive. The bound $C_H^{\sharp}\ge\tfrac1{80}$ is not established, and the proof is incomplete at this step. Either of the following one-sided conditions would supply the step. Neither is among the conditions of this chapter, and neither is a cap in the sense of Definition~\ref{2.3:def-cap}.
\begin{itemize}
\item[(i)] Impose \eqref{4.27:eq-t-step-layer-2} with $\max\{C_H^{\sharp},1\}$ in place of $C_H^{\sharp}$. This raises a constant that is fixed together with $K_\chi$, after $\kappa_1$ and before the exponent $q$ and the constant $C$ of the radii determine $\gamma$. The order of choice is therefore kept, nothing chosen earlier changes, and every bound proved at $C_H^{\sharp}$ persists. Then $\varpi_k\le\alpha_{0,k}/80\le\varepsilon_\sharp$.
\item[(ii)] Impose $\sup|b|\le\varepsilon_\sharp$ through the stage-factor lemma, as a ceiling on $\gamma$ of the same kind as the convergence condition on the radii.
\end{itemize}

\emph{The weight.} By input~1(a), $p\in\mathfrak{M}_\delta(\mathfrak{B})$. On $\mathcal{C}$ the distance $d^{\wedge}(\cdot,\mathcal{C})$ vanishes, so $p=\varpi_k$ there. Since $\sup|b|\le\varpi_k$, this gives $\|\mathsf{d}\|_p\le1$.

\emph{The layer bound.} Apply clause (b) of the background-field layer theorem to the data $\mathsf{d}$ and $\mathsf{d}_2:=0$. At $\mathsf{d}_2$ the layer vanishes by clause (a). Together with $\|\mathsf{d}\|_p\le1$ this yields $N^{\Delta}_y(\mathfrak{h})\le B_\sharp p(y)$, the first bound in \eqref{4.27:eq-t-step-layer-a}.

\emph{The source increment.} Apply clause (iv) of the recursive presentation bounds with its cut-off $\chi$ identically equal to $1$. Its proof uses only the first bound in \eqref{4.27:eq-t-step-layer-a} and two further estimates, clause (a) of the sharp estimates and the fourth of the sharp lemmas. None of these depends on the particular determining set $\mathfrak{B}$. Its smallness condition is $2B_\sharp\varpi_k\le\tfrac1{10}$. Here \eqref{4.27:eq-t-step-layer-1} and $C_H^{\sharp}\ge120B_\sharp$ give $120B_\sharp\varpi_k\le0.0125\,\alpha_{0,k}$, that is, $2B_\sharp\varpi_k\le\alpha_{0,k}/4800$, and the smallness condition holds. This gives the second bound in \eqref{4.27:eq-t-step-layer-a}, $\ell^3|\mathcal{J}^{\sharp}[1](y)|\le 2K_\chi p(y)$.

\emph{The direct entries.} The entries $Q_6$, $Q_7$ and $Q_1$ move by at most $3B_\sharp p(y)$. This follows from the two bounds on the movement of these entries under a layer that underlie input~1(b) and input~1(c), in the form in which those clauses use them. By the formula for $K_\chi$ in input~5, $K_\chi\ge 3B_\sharp$.
The entries moved by the layer obey the bounds of the background-field layer theorem, with its constant $C_H$. The source entry moves by at most $2K_\chi p(y)$, and the entries $Q_6$, $Q_7$, $Q_1$ by at most $3B_\sharp p(y)\le K_\chi p(y)$. Therefore every direct entry moves by at most
\begin{equation*}
\max\{C_H,2K_\chi\}\,p(y)=C_H^{\sharp}\,p(y).
\end{equation*}

\emph{Real points.} At a real point, the inequalities of clause (iii) of the recursive presentation bounds apply to the pair of determining sets $(\mathbb{B}_{k+1},\mathbb{B}_k)$. The averaging map to level $k+1$ is the one-step averaging composed with the averaging map to level $k$. Hence the $\mathbb{B}_{k+1}$-class of $U_{k+1}$ contains, through $U_{k+1}$, the $\mathbb{B}_k$-class of $U_{k+1}$. Since $U_{k+1}$ minimises over its $\mathbb{B}_{k+1}$-class, it is critical for $\mathbb{B}_k$ at its own averages. Therefore the current $\mathfrak{j}=\tilde G_0\mathcal{J}(U)$ vanishes. Clause (d) of the background-field layer theorem then returns Ba{\l}aban's $\exp i\eta\mathcal{H}(B)$ of \cite[Proposition~9]{B-CMP102b}.

This identification is the reason clause (d) is used. \cite[Proposition~9]{B-CMP102b} requires the conditions \cite[(3.35)--(3.38)]{B-CMP99b}, which the spaces of pairs of this chapter do not supply, so that proposition is not applied here. The bounds \eqref{4.27:eq-t-step-layer-a} and the bound on the direct entries are those obtained above from input~1 through the background-field layer theorem. At a real point, clause (d) identifies the layer so bounded with Ba{\l}aban's.

\emph{Order of choice.} The constant $K_\chi$ is fixed before $(q,C)$. Replacing $C_H$ by $C_H^{\sharp}$ in the lower bound on $C^{\mathrm{rad}}$ therefore only replaces that lower bound by a larger one, fixed at the same place in the order of choice, and no earlier choice changes.
\end{proof}

The bound \eqref{4.27:eq-t-step-layer-a} is derived from input~1. It is not taken from the hypothesis of the same shape in the stage-factor lemma, which is posed on $\partial U_k$ with the factors $\prod_m e^{-2\delta_0MR_m}$ of the thick strata.

\noindent Input~5, continued. Assume \eqref{4.27:eq-t-step-layer-2} and the same lower bound for $C_1^{\mathrm{rad}}$. Consider the coefficient $\pi$ of the rescaling bound, input~1(e), for the layer of the $\mathbf{T}$-step, written $\pi_m$ on the layer of level $m\le k$ and $\pi_{\mathrm{top}}$ on $\Omega_{k+1}$. It satisfies
\begin{equation*}
\pi_m\le\tfrac14(1-\theta_S)\beta\,2^{-(k+1-m)}\quad(m\le k),\qquad \pi_{\mathrm{top}}\le\tfrac{\beta}{8}.
\end{equation*}
In the stage-factor lemma, the coarse-over-fine ratio $(1+\beta_0)^2(k-n)^{\beta_0}$ is, by Lemma~\ref{4.4:lem-radii-ratio}, used in the form $(1+\beta_0)^2(k-n)^{1/2}$. The bracket of that lemma then becomes $\nu^{1/2}(2u)^{\nu}$ with $\nu\ge1$ and $u\le\tfrac14$, and $\nu^{1/2}(2u)^{\nu}\le\tfrac12$. Indeed, $(2u)^{\nu}\le2^{-\nu}$, equality $\nu^{1/2}2^{-\nu}=\tfrac12$ holds at $\nu=1$, and
\begin{equation*}
\frac{(\nu+1)^{1/2}2^{-\nu-1}}{\nu^{1/2}2^{-\nu}}\le\frac{\sqrt2}{2}<1 .
\end{equation*}
So that lemma holds with its statement unchanged.

\begin{enumerate}
\item[6.] \emph{The keystone bound.} Let $\mathsf{Z}:=\Lambda^{\natural}\cup\mathfrak{N}^{\mathrm{rec}}$. On $\mathsf{Z}$ each entry of the presented pair is at most $\tfrac1{16}$ of the threshold of its consumer term $\mathfrak{E}$, for every consumer term $\mathfrak{E}$. It rests on the core statement of the keystone bound.
\item[7.] \emph{Real points.} Let $U$ be a real point.
\begin{itemize}
\item On determining sets holding excluded attempts, the assignment lemma for excluded attempts applies under its setting. Its geometric conditions, its data clause, its flow conditions and its assignment condition for excluded attempts are in force. By its clause (c), $U$ lies, with three quarters of the radius to spare, in every space whose clause at $x$ has all of the following: the current level of $x$; a weight at least $1$; a multiplicity at least $1$; a factor at least $\varphi_\flat$. Derived entries are included. The remaining quarter is the fraction in the last member of that lemma's lower bound,
\begin{equation*}
K^{\mathrm{own}}\varepsilon^{\mathsf{l}}\le\tfrac14\,\varphi_\flat\,\alpha_0 .
\end{equation*}
\item Off the pieces, the membership of a real point is given by two statements: the real-point clause of the stage-space statements, and the proposition on real points of the stage spaces.
\end{itemize}
\item[8.] \emph{The declaration on the couplings.} The declaration on the couplings $g_j$ is made in clause (ii) of the lemma fixing the weights of an arrest. The bounds of input~10 are stated under it at real points. It is used together with the arrest condition.
\item[9.] \emph{The post-arrest input.} Two parts of the statement on operations after an arrest are used; they carry the letters of that statement.
\begin{itemize}
\item Clauses (a) and (b) describe, exit by exit, the exponent and the objects depending on $(\mathbb{U},\mathbb{J})$ on an excluded region $X$.
\item Clause (d): every operation after an arrest that composes a stored term with a map of pairs $(\mathbb{U},\mathbb{J})$ is one of $O1,\dots,O5$, $O4'$. Every term it creates is assigned its domain by the third clause of the arrest theorem, the clause that fixes the domains of newly created terms.
\end{itemize}
\item[10.] \emph{The localised layer.} The background-field layer theorem holds at the determining set $\mathfrak{B}_Z$ localised to $Z$ as in \cite[(2.14)]{B-CMP122b}, at $\gamma_0^{(k)}$, with the same constants. The function $H_\sigma(b')$ is holomorphic on
\begin{equation*}
\mathfrak{D}_H(Y)\supset\mathfrak{D}^{\flat}_{n+1}(\mathcal{Y})\restriction Y ,
\end{equation*}
where the stage domain $\mathfrak{D}^{\flat}_{n+1}(\mathcal{Y})$ is taken with the following ambient condition.
\begin{itemize}
\item Let $Y$ meet $Z_W^{+}$. Then the continuation lies in $\mathfrak{K}_{109}(\mathfrak{B}^{\circ};\gamma_0^{(k_{\mathfrak{a}})})$ on $Z_W=X_{\mathfrak{a}}$, with $\mathfrak{B}^{\circ}=\mathbb{B}^{(n+1)}_{k_{\mathfrak{a}}}\restriction X_{\mathfrak{a}}$. This membership is supplied by the fifth clause of the arrest theorem, as a membership contained in the reading of the domain of the stored term.
\item For later sources and images, the arrest theorem supplies the ambient condition as $\mathfrak{G}_1(\mathfrak{B}_\tau)$ together with its list of strata:
\begin{itemize}
\item on $X_{\mathfrak{a}}\cap\Omega_1(\cdot)$, at the cubes;
\item at the raw points of $X_{\mathfrak{a}}$ off $Y'$: the index-zero $\mathbb{U}$-members, and the dated $\mathbb{J}$-member with threshold at most $\gamma_0^{(k_{\mathfrak{a}})}$, fixed at birth;
\item on $Y'$, where the map is a layer into $U'$: the gauge member $(3.35)_0$, the index-zero instance of \cite[(3.35)]{B-CMP99b}.
\end{itemize}
\item The members that a map moves on $X_{\mathfrak{a}}\cap Y'$ at index zero constitute the second of the two residual cases of the dating fixed at birth; that case is treated in the statement on the residual cases of the dating fixed at birth.
\end{itemize}
At a real point, where $\sigma\equiv1$, $H_\sigma$ is the function $\mathbb{H}^{(j)}_Z$ of \cite[(1.58)]{B-CMP122a}. For every bond $b'$,
\begin{equation}\label{4.27:eq-t-step-layer-b}
|H_\sigma(b')|\le H_{\max}:=L\,C_F\,B_\sharp\,e^{4\kappa_1}K_{\mathfrak{f}}\,
\bar\alpha^{\mathrm{fat}}_k\,e^{-4R_k},
\end{equation}
where
\begin{equation*}
\bar\alpha^{\mathrm{fat}}_k:=1.4\,\bigl(BCM\,\alpha_{0,k}+\alpha_{1,k}\bigr),
\end{equation*}
with the constants $B$, $C$ and $M$ of this definition of $\bar\alpha^{\mathrm{fat}}_k$. The shift obeys $|\delta_{b'}|\le 2H_{\max}/g_j$. At real points this is stated under the declaration on the couplings $g_j$ of input~8; at complex points it is the bound on the lift.
\end{enumerate}

\begin{definition}[Hypotheses of the arrest theorem]\label{4.27:def-arrest-hypotheses}
Throughout, $k$ is the level of the step, $N_k=\check{R}_k$ is the declared number of stages,
$h=k-N_k$, and the factors, the radii $\alpha_i(\mathsf{t})$, the constant
$c_\beta=2\ln(1+\beta_0)$ and the logarithmic variables
$\mathsf{x}_j=g_j^{-2}$, $\mathsf{t}_j=\log\mathsf{x}_j$, $\mathsf{t}_\gamma=\log\gamma^{-2}$ are those of
Convention~\ref{4.27:not-factors-radii-coupling}. For a level $j$ we put
$\delta'_j=g_jA_1\mathsf{t}_j^{\,p_1}$, with the auxiliary parameter $A_1$ and a fixed exponent
$p_1$, and we write $r$ for the exponent in the declaration $\check{R}_j<L\mathsf{t}_j^{\,r}$ of
Definition~\ref{4.27:def-one-class-per-step}. We write $\varphi_{\mathrm{arrest}}$ for the number below which
no stage factor $\varphi^{(n),(k)}_m$ of Convention~\ref{4.27:not-factors-radii-coupling} falls;
$\varphi_\flat$ for the factor threshold of the real-point lemma on determining sets holding
excluded attempts, that is, the least factor admitted in the spaces in which that lemma places
the real point; and $\varphi_{\min}$ for the lower bound on the stage factors with which the
floors collected in (Fl11) and (Fl12) below are stated in their own setting. The
\emph{hypotheses of the arrest theorem} are the
conjunction of three groups of conditions: the declarations of
Definition~\ref{4.27:def-one-class-per-step}, the floors listed below, and the ceilings listed
below.

\smallskip
\emph{Floors.} The following one-sided conditions hold.
\begin{enumerate}
\item[(Fl1)] $L\ge 8$; in the window $2\le\mathsf{L}_0$, $\mathsf{L}_0^2\le L/3$ of the
declarations this becomes $L\ge 12$, since $\mathsf{L}_0\ge 2$ gives $4\le\mathsf{L}_0^2\le L/3$.
\item[(Fl2)] $\delta\,\mathsf{w}_M\ge 8$, where $\delta$ is the decay constant of the class of
decaying weights and $\mathsf{w}_M=M/M_1$; this is implied by the floor on the levels in the
order of choice of the parameters.
\item[(Fl3)] $\delta M/M_1\ge 3$, in the form in which the descent lemma for the stage spaces
uses it.
\item[(Fl4)] $\delta(M/M_1)\ge 2$, as in \cite[(1.47)]{B-CMP122a}.
\item[(Fl5)] $M\ge M_*(L)$, the floor on $M$ of the descent proposition for the stage spaces.
\item[(Fl6)] The floors on $M$ in the order of choice of the theorem on the stage maps at an
arbitrary determining set.
\item[(Fl7)] The floor on the parameter block $\mathfrak{M}$ of the lemmas on the localised
operators $H_\sigma(b)$, the parameters being chosen in the order
$L,\ \mathfrak{M},\ \kappa,\ \kappa_1,\ M_1,\ M,\ \gamma$.
\item[(Fl8)] The radius condition
\begin{equation*}
C^{\mathrm{rad}}_0\ge 2.5\,A_0/\varphi_{\mathrm{arrest}},\qquad q>p_0 .
\end{equation*}
\item[(Fl9)] The floor on the radius constants of the $\mathbf{T}$-step layer in a background field
(Lemma~\ref{4.27:lem-t-step-layer}, whose proof is incomplete at one step), with the constant
$C_H^\sharp$ of that lemma; it is the second row of \eqref{4.27:eq-arrest-hypotheses-a}.
\item[(Fl10)] The flow conditions of the coupling on the declared stage numbers $\check{R}$ of
Definition~\ref{4.27:def-one-class-per-step}: the reference comparison over at most $N_0+1$
levels ($N_0=N_0(k)$ being the threshold for the number of stages that enters the room
condition in (Fl12) and the condition (Ce2)); the two flow-line conditions on $\check{R}$; and the
first of the decrease conditions on $\check{R}$ --- all in the form, carrying the parameter
$\varepsilon$, in which the theorem on the change of variables $\mathbb{H}$ invokes them.
\item[(Fl11)] The one-sided floors on $\mathsf{x}$ of the real-point lemma on determining sets
holding excluded attempts and of the bending floor for excluded attempts, both taken at the
factor $\min\{\varphi_{\min},\varphi_\flat\}$.
\item[(Fl12)] The room condition on $N_0$ and the accompanying constant condition, both of the
statement fixing the room in the number of stages, the latter with its constant $\Gamma$
replaced by $\Gamma_\flat=(\varphi_\flat/\varphi_{\min})\,\Gamma$.
\item[(Fl13)] The floors on the constants $K_W^\sharp$ and $B_\sharp$ under which the theorem on
the stage maps at an arbitrary determining set bounds each entry of the presented pair by one
sixteenth of the threshold of the term that consumes it.
\end{enumerate}

\smallskip
\emph{Ceilings.} Each of the following conditions that bounds $\gamma$ is an upper bound on
$\gamma$ by a number depending only on constants fixed before $\gamma$, and is understood as
an inequality
\begin{equation*}
\sup_{\mathsf{t}\ge\mathsf{t}_\gamma}F(\mathsf{t})\le 1,
\end{equation*}
where $F$ is built from the scaffold functions $R(\mathsf{t})$ (the block radius),
$N(\mathsf{t})$ (the number of stages), $\mathsf{l}(\mathsf{t})$, $\alpha_i(\mathsf{t})$, the
small parameter $\varepsilon(\mathsf{t})$, $\delta'(\mathsf{t})$,
a further scaffold function $\theta^{\mathrm{sc}}(\mathsf{t})$ and $\bar\gamma_0(\mathsf{t})$, all
functions of one and the same variable
$\mathsf{t}$, and every constant entering $F$ is fixed before $\gamma$. The three ceilings on
$\gamma$ that are explicit inequalities, (Ce1), (Ce11) and (Ce14), are collected, together with
the floor (Fl9) and the conditions (Ce10) and (Ce13), which are named in words, in the display
\begin{equation}\label{4.27:eq-arrest-hypotheses-a}
\begin{gathered}
\mathsf{t}_\gamma\ge 2(r+q)+c_\beta,\qquad
4H_{\max}\le 0.3\,\delta'_j,\qquad
2.2\,K^{b}_{S}\,\mathsf{t}_\gamma^{\,p_0-q}\le 1,\\
\text{the floor (Fl9) on $C^{\mathrm{rad}}_0$ and $C^{\mathrm{rad}}_1$},\\
\text{the hypothesis of the lemmas on $H_\sigma(b)$, less its clause (iv)},\\
\text{the ceiling on $\gamma$ of the lemma fixing the margin $\mathsf{m}$},
\end{gathered}
\end{equation}
where $H_{\max}$ is the quantity \eqref{4.27:eq-t-step-layer-b} of the $\mathbf{T}$-step layer in a
background field (Lemma~\ref{4.27:lem-t-step-layer}, whose proof is incomplete at one step), and
$K^{b}_{S}$ is the constant in the bound of the theorem on the stage maps at an arbitrary
determining set on the increments of its entries two and four.
The full list is as follows.
\begin{enumerate}
\item[(Ce1)] The first inequality of the first row of \eqref{4.27:eq-arrest-hypotheses-a}. It
follows from the inequality $\log\mathsf{x}\ge 4(p_0+2r+q)$, which is a member of the floor of
the real-point lemma on determining sets holding excluded attempts in (Fl11).
\item[(Ce2)] $N_k>N_0(k)$.
\item[(Ce3)] $N_k\ge 6$; this holds whenever $\check{R}_k\ge L$, since $N_k=\check{R}_k$ and
$L\ge 8$ by (Fl1).
\item[(Ce4)] $\beta_0\le 1/7$.
\item[(Ce5)] The ceiling on $\gamma$ of the descent lemma for the stage spaces.
\item[(Ce6)] The ceiling on $\gamma$ of the descent proposition for the stage spaces.
\item[(Ce7)] The list of ceilings on $\gamma$ in the order of choice of the theorem on the stage
maps at an arbitrary determining set; and the following four, each taken at
$\gamma_0=\bar\gamma_0$ (the run-free ceiling of the radius $\gamma_0^{(k)}$): the sharp
$\gamma$-condition of that theorem, the inequality $K_\Delta\gamma_0\le\tfrac12$, the condition
$C^{b}(a_\flat+\kappa_{\max})\le 1$ taken at $a_\flat$, and $\kappa_{\max}\le\gamma_0$. Here
$C^{b}$ is the constant and $a_\flat$, $\kappa_{\max}$ are the two parameters of the smallness
condition $C^{b}(a_\flat+\kappa_{\max})\le 1$ of that theorem, under which its bound on the
increments of the entries two and four holds, and $K_\Delta$ is the constant of its ceiling
$K_\Delta\gamma_0\le\tfrac12$ on the radius $\gamma_0$; the constants entering these four
conditions are fixed before $\gamma$.
\item[(Ce8)] The window condition of the inheritance statement for stored terms.
\item[(Ce9)] The arrest condition, with the constant $r_*^\sharp$ satisfying
$6.6\le r_*^\sharp$; here $r_*^\sharp$ is the largest factor admitted in the bound on the datum $b$
in the theorem on the stage maps at an arbitrary determining set, which applies when
$\|b\|_\infty$ is at most $C_1\delta'_j$ times a factor between $1$ and $r_*^\sharp$.
\item[(Ce10)] The hypothesis of the lemmas on the localised operators $H_\sigma(b)$ in the change
of variables, with its clause (iv) omitted and with its clause (i) in the primed form in which
the relative small-field machine uses it; it is the third row of
\eqref{4.27:eq-arrest-hypotheses-a}. Its clause (vi) contains the requirement that
\begin{equation}\label{4.27:eq-arrest-hypotheses-1}
\mathsf{x}_j\in[\mathsf{x}_k-c_5,\ 2\mathsf{x}_k]\qquad\text{for } h\le j<k,
\end{equation}
which is supplied by the two-sided flow bounds of the coupling together with
$N_k<L\,\mathsf{t}_k^{\,r}$ from the declaration on $\check{R}$. The lower member
$\mathsf{x}_k\le\mathsf{x}_j$ of the original clause (vi), which is a monotonicity of the
coupling, is not part of the hypotheses. In the hypotheses, clause (vi) is taken as the
conjunction of \eqref{4.27:eq-arrest-hypotheses-1} and
$\mathsf{t}_j\in[\mathsf{t}_k-c_\beta,\ \mathsf{t}_k+\ln 2]$ for $h\le j<k$, so that all levels
$h\le j<k$ are governed by the single variable $\mathsf{t}=\mathsf{t}_k$. The upper endpoints of
the two members agree, since $\mathsf{x}_j\le 2\mathsf{x}_k$ is equivalent to
$\mathsf{t}_j\le\mathsf{t}_k+\ln 2$; the lower endpoint $\mathsf{t}_j\ge\mathsf{t}_k-c_\beta$ is an
additional member.
\item[(Ce11)] The second inequality of the first row of \eqref{4.27:eq-arrest-hypotheses-a},
understood as
\begin{equation}\label{4.27:eq-arrest-hypotheses-2}
\sup_{\mathsf{t}\ge\mathsf{t}_\gamma}
\frac{4H_{\max}(\mathsf{t})}{0.3\cdot 2^{-1/2}A_1e^{-\mathsf{t}/2}(\mathsf{t}-c_\beta)^{p_1}}
\le 1 ,
\end{equation}
the supremum being taken over all $\mathsf{t}\ge\mathsf{t}_\gamma$; it is applied below at
$\mathsf{t}=\mathsf{t}_k$, for $\mathsf{t}_k\ge\mathsf{t}_\gamma$, that is, for $g_k\le\gamma$.
The inequality $H_{\max}\le\delta'_j/6$ of clause (iii) of the hypothesis of the lemmas on the
localised operators $H_\sigma(b)$ in (Ce10)
is understood in the same way, as
\begin{equation}\label{4.27:eq-arrest-hypotheses-4}
\sup_{\mathsf{t}\ge\mathsf{t}_\gamma}
\frac{6H_{\max}(\mathsf{t})}{2^{-1/2}A_1e^{-\mathsf{t}/2}(\mathsf{t}-c_\beta)^{p_1}}
\le 1 ;
\end{equation}
in the quantity $\theta_{\mathrm{cv}}$ of that clause the factor
$g_j^{-1}$ enters only through the bound $g_j^{-1}\le(2\mathsf{x}_k)^{1/2}$.
\item[(Ce12)] $16\bar P\le\beta\theta_S$, where $\bar P$ is the maximal increment of the
large-field layers.
\item[(Ce13)] The ceiling on $\gamma$ of the lemma fixing the margin $\mathsf{m}$; it is the
fourth row of \eqref{4.27:eq-arrest-hypotheses-a}.
\item[(Ce14)] The third inequality of the first row of \eqref{4.27:eq-arrest-hypotheses-a},
imposed after $L$ and $M$ are fixed.
\end{enumerate}

\smallskip
\emph{The reading of} (Ce11). By clause (vi) of the hypothesis of the lemmas on the localised
operators $H_\sigma(b)$ in (Ce10), in its member \eqref{4.27:eq-arrest-hypotheses-1}, one has
$\mathsf{x}_j\le 2\mathsf{x}_k$, hence $g_j=\mathsf{x}_j^{-1/2}\ge(2\mathsf{x}_k)^{-1/2}$, and, by
its member in the variable $\mathsf{t}$, $\mathsf{t}_j\ge\mathsf{t}_k-c_\beta$.
Therefore, for $h\le j<k$ and $\mathsf{t}_k\ge\mathsf{t}_\gamma$,
\begin{equation}\label{4.27:eq-arrest-hypotheses-3}
\delta'_j=g_jA_1\mathsf{t}_j^{\,p_1}
\ge(2\mathsf{x}_k)^{-1/2}A_1(\mathsf{t}_k-c_\beta)^{p_1}
=2^{-1/2}A_1e^{-\mathsf{t}_k/2}(\mathsf{t}_k-c_\beta)^{p_1},
\end{equation}
and the second inequality of the first row of \eqref{4.27:eq-arrest-hypotheses-a} holds for
every such $j$, when $\mathsf{t}_k\ge\mathsf{t}_\gamma$, as soon as
\eqref{4.27:eq-arrest-hypotheses-2} holds. In the definition
\eqref{4.27:eq-t-step-layer-b} of $H_{\max}$ (Lemma~\ref{4.27:lem-t-step-layer}, whose proof is
incomplete at one step) the factors depending on $\mathsf{t}$ are the
fattened radius, a fixed linear combination of $\alpha_0(\mathsf{t})$ and $\alpha_1(\mathsf{t})$,
and $e^{-4R(\mathsf{t})}$. Since
$\alpha_i(\mathsf{t})=C_i^{\mathrm{rad}}e^{-\mathsf{t}/2}\mathsf{t}^q$,
the quotient in \eqref{4.27:eq-arrest-hypotheses-2} is a constant multiple of
\begin{equation*}
\mathsf{t}^{\,q}\,(\mathsf{t}-c_\beta)^{-p_1}\,e^{-4R(\mathsf{t})} ,
\end{equation*}
and the constant is fixed before $\gamma$. By the declaration on the stage numbers in
Definition~\ref{4.27:def-one-class-per-step}, $R(\mathsf{t})\ge(\mathsf{t}-c_\beta)^r$. On the
range $\mathsf{t}\ge\mathsf{t}_\gamma$ one has $\mathsf{t}-c_\beta\ge 2(r+q)$ by (Ce1). Hence the
quotient is bounded on this range and tends to $0$ as $\mathsf{t}\to\infty$. Consequently
\eqref{4.27:eq-arrest-hypotheses-2} is a ceiling of the form described above.

\smallskip
None of the conditions above is an upper bound on $L$, $k$, $K$, the number of stages $N_k$,
the ages of the large-field regions present (the number of steps since the step that created
each of them) or their nesting depth, and none is a lower bound on $g$.
\end{definition}

\begin{definition}[Arrests: pieces, collars, bands and frozen terms]\label{4.27:def-arrests-pieces}
We work with the blocking data and the convention of one class per step fixed in
Definition~\ref{4.27:def-one-class-per-step}. Below, $\mathbf{R}$ is Ba{\l}aban's large-field
operation, and the \emph{large-field exits} of the preparation are those of
\cite[\S1]{B-CMP122a}. An exit is a place there at which those components of the class in which
new large fields have been introduced are excluded from the region $Z$; such
components are not changed by the subsequent operations, and are treated in the same way as the
remaining large-field regions. There are nine exits. We label them (e0), (e0$'$), (e1), \dots,
(e7), in the order in which they occur there.

\emph{Arrests.} An \emph{arrest} is a tuple
\begin{equation*}
\mathfrak{a}=(X_{\mathfrak{a}},\,k_{\mathfrak{a}},\,n,\,n^{+},\,e)
\end{equation*}
that records the following event. The event takes place at the application of $\mathbf{R}$ in
step $k_{\mathfrak{a}}$. At that application the declared number of stages
$N:=N_{k_{\mathfrak{a}}}$ and the indices $h_{\mathfrak{a}}$, $k_0^{\mathfrak{a}}$ of the class are
those of Definition~\ref{4.27:def-one-class-per-step}. The class is a set
$Z\subset Z_{k_{\mathfrak{a}}}=\Lambda^{c}_{k_{\mathfrak{a}}}$. The event is that its component
$X_{\mathfrak{a}}$ leaves the preparation through the large-field exit with label $e$, after $n$
performed stages. We put
\begin{equation*}
j:=h_{\mathfrak{a}}+n .
\end{equation*}

\emph{Piece and collar.} The \emph{piece} and the \emph{collar} of $\mathfrak{a}$ are
\begin{equation}\label{4.27:eq-arrests-pieces-1}
W_{\mathfrak{a}}:=X_{\mathfrak{a}}\cap\Omega^{c}_{k_{\mathfrak{a}}},
\qquad
\mathfrak{c}_{\mathfrak{a}}:=X_{\mathfrak{a}}\cap\Omega_{k_{\mathfrak{a}}} .
\end{equation}
The collar lies at level $k_{\mathfrak{a}}$. It consists of two layers of cubes of side
$MR_{k_{\mathfrak{a}}}$.

The split of $X_{\mathfrak{a}}$ into piece and collar is a choice. It is made so that the
statements that use arrests have $W_{\mathfrak{a}}\subset Z^{\mathrm{on}}$ and
$\mathfrak{c}_{\mathfrak{a}}\subset\mathfrak{T}_{k_{\mathfrak{a}}}$ at their disposal. The
constant $c=3$ that \cite[(1.68)]{B-CMP122a} carries at the top level $m=k_{\mathfrak{a}}$ acts on
$\Omega_{k_{\mathfrak{a}}}\cap X_{\mathfrak{a}}$, hence off $W_{\mathfrak{a}}$.

\emph{The exits.} The list below gives, for each exit, four items:
\begin{itemize}
\item the values of $n$ and $n^{+}$;
\item the \emph{frozen boundary sum};
\item the real constraints that are kept;
\item the variables that remain un-integrated besides the block data $V$.
\end{itemize}
In the list $k=k_{\mathfrak{a}}$ and $h=h_{\mathfrak{a}}$. The characteristic functions $\chi_h$,
$\chi^{(0)}$, $\chi^{(n)}$, $\chi^{(n+1)}$, $\chi''_k$, $\chi_{k,\Lambda}$ and the variables
$A_{j'}$ are those of \cite[\S1]{B-CMP122a}, and $B_j$ is the variable introduced at
\cite[(1.53)]{B-CMP122a}. For $n\le N$,
\begin{equation*}
\mathbb{B}^{(n)}_k=\sum_{i\le n}\mathbb{C}^{(i)}_k
\end{equation*}
is the sum of the new boundary terms $\mathbb{C}^{(i)}_k$ of the first $n$ stages, and
$\mathbb{B}_k$ is the sum of the boundary terms of step $k$ present before the first stage, in
the notation of \cite[\S1]{B-CMP122a}.
\begin{itemize}
\item[] \textup{(e0), (e0$'$).} Here $n=0$ and $n^{+}=0$. The frozen boundary sum is
$\mathbb{B}_k$. The constraints kept are \cite[(1.3), (1.4)]{B-CMP122a} together with $\chi_h$;
for (e0$'$) the characteristic function is $\chi^{(0)}$. The un-integrated variables are
$A_{j'}$, $j'\ge h$.
\item[] \textup{(e1).} Here $n^{+}=n$. The frozen boundary sum is $\mathbb{B}^{(n)}_k$. The
constraint kept is $\chi^{(n)}=1$. The un-integrated variables are $A_{j'}$, $j'\ge j$.
\item[] \textup{(e2), (e3).} Here $n^{+}=n$. The frozen boundary sum is $\mathbb{B}^{(n)}_k$. The
constraints kept are $\chi^{(n)}=\chi^{(n+1)}=1$. The un-integrated variables are $A_{j'}$,
$j'\ge j$.
\item[] \textup{(e4).} Here $n\le N-1$ and $n^{+}=n+1$. The frozen boundary sum is
$\mathbb{B}^{(n+1)}_k$, whose band lies at level $j+1$ and whose datum is $V_{j+1}$. The
constraints kept are $\chi^{(n+1)}=1$ and \cite[(1.59)]{B-CMP122a}. The un-integrated variables
are $A_{j'}$, $j'\ge j$, and in addition $B_j$.
\item[] \textup{(e5).} Here $n=N$ and $n^{+}=N$. The frozen boundary sum is
$\mathbb{B}''_k=\mathbb{B}^{(N)}_k$. The constraint kept is $\chi''_k=1$. No variable remains
un-integrated.
\item[] \textup{(e6), (e7).} Here $n=N$ and $n^{+}=N$. The frozen boundary sum is
$\mathbb{B}''_k$. The constraints kept are $\chi''_k=\chi_{k,\Lambda}=1$. No variable remains
un-integrated: the integration over $V_k\restriction\Lambda$ is an operation and leaves no term.
\end{itemize}

\emph{Range and range sum.} Put $j^{+}:=h_{\mathfrak{a}}+n^{+}$. Then $j^{+}\le k_{\mathfrak{a}}$.
For every arrest, whatever its exit, the \emph{range} of $\mathfrak{a}$ and its \emph{range sum}
at an integer level $\ell$ are
\begin{equation}\label{4.27:eq-arrests-pieces-2}
I(\mathfrak{a}):=\{h_{\mathfrak{a}}+1,\dots,k_{\mathfrak{a}}+1\},
\qquad
\sigma_{\mathfrak{a}}(\ell):=\sum_{i\in I(\mathfrak{a})}2^{-|\ell-i|} .
\end{equation}

\emph{Bands.} For a stage $i\le n$, write $\ell_i$ for the level at which stage $i$ works. The
\emph{band} of stage $i$ is
\begin{equation}\label{4.27:eq-arrests-pieces-3}
R^{(i)}_{\mathfrak{a}}
:=\bigl(Z_{\ell_i+1}\setminus Z''_{\ell_i+1}\bigr)\cap X_{\mathfrak{a}} .
\end{equation}

\emph{Frozen structure.} For $x\in W_{\mathfrak{a}}$, let
$(\ell_{\mathfrak{a}},\mathsf{w}_{\mathfrak{a}},\mathcal{Q}_{\mathfrak{a}})(x)$ be the triple
that the stage space $\mathfrak{S}^{(k_{\mathfrak{a}})}_{n^{+}}$ of
\cite[(1.64)--(1.69)]{B-CMP122a} assigns at $x$. Its entries are
the level, the $\mathsf{L}_0$-weight and the cube family.

\emph{Frozen terms.} The \emph{frozen terms} of $\mathfrak{a}$ are
\begin{equation}\label{4.27:eq-arrests-pieces-4}
\mathfrak{F}(\mathfrak{a})
:=\mathfrak{F}^{\mathbb{C}}(\mathfrak{a})\cup\mathfrak{F}^{(\mathrm{e}4)}(\mathfrak{a}).
\end{equation}
The first family is
\begin{equation}\label{4.27:eq-arrests-pieces-5}
\mathfrak{F}^{\mathbb{C}}(\mathfrak{a})
:=\bigl\{\mathbb{C}^{(i)}_{k_{\mathfrak{a}}}(X,\mathcal{S}) :
1\le i\le n,\ X\cap R^{(i)}_{\mathfrak{a}}\neq\emptyset\bigr\},
\end{equation}
where $\mathcal{S}$ is the setting in which the step is taken, in the sense of
Definition~\ref{4.4:def-settings-data-domains}, and $\mathfrak{D}_{\mathcal{S}}(X)$ is the data
domain on $X$ that this setting fixes. The family consists of the new boundary terms of the
performed stages $i\le n$ only, and only of those whose domain meets the band of their stage. The
member with stage index $i$ and domain $X$ is defined on the product
$\mathfrak{D}^{\flat}_i(X)\times\mathfrak{D}_{\mathcal{S}}(X)$. It carries the bound of the first
clause of the stage-term estimates of this chapter.

The second family $\mathfrak{F}^{(\mathrm{e}4)}(\mathfrak{a})$ consists of the terms singled out
by the lemma of this chapter on the exit (e4). It is empty unless $e$ is (e4). When $e$ is (e4),
its members are defined on the domain of stage $n^{+}=n+1$. On the exit (e4), the terms
$\mathbb{C}^{(i)}_{k_{\mathfrak{a}}}$ with $i\le n$ survive only inside clause ($\beta1$) of that
lemma. In $\mathfrak{F}(\mathfrak{a})$ each of them is counted once.

\emph{Frozen variables.} The \emph{frozen variables} $\mathfrak{v}(\mathfrak{a})$ are the
variables listed as un-integrated for the exit $e$ in the list above. The performed stages reach
level $j-1$, and level $j$ is never integrated on $X_{\mathfrak{a}}$. At every later application
of an operation, the frozen variables are exterior coordinates.

\emph{Locality.} The constraints kept are local in $V\restriction X_{\mathfrak{a}}$: they depend
on the block data $V$ only through its restriction to $X_{\mathfrak{a}}$.

\emph{Live and dissolved.} An arrest is \emph{live} from the arrest onwards, until the first later
completed application of $\mathbf{R}$ on the component that holds $X_{\mathfrak{a}}$. An
application of $\mathbf{R}$ is completed when it reaches \cite[(1.70)]{B-CMP122b}, whether it is
of the first or of the second class. From that application on, the arrest is \emph{dissolved}.
\end{definition}

\begin{lemma}[Timing, nesting and bands of arrests]\label{4.27:lem-arrest-timing}
Let $\mathfrak{a}$ be a live arrest in the sense of Definition~\ref{4.27:def-arrests-pieces}, with
component $X_{\mathfrak{a}}$, piece $W_{\mathfrak{a}}$, step $k_a$, range $I(\mathfrak{a})$,
bands $R^{(i)}_{\mathfrak{a}}$ and frozen level $\ell_{\mathfrak{a}}$. Let
$\mathfrak{b}\neq\mathfrak{a}$ be an attempt at the $\mathbf{R}$-site of a step $k_b$ (the
treatment at that site of one component of the class, whether or not it takes a large-field exit),
with component $X_{\mathfrak{b}}$, piece $W_{\mathfrak{b}}$ and
$h_b=k_b-N_{k_b}$ as in Definition~\ref{4.27:def-arrests-pieces}. Assume
$W_{\mathfrak{a}}\cap W_{\mathfrak{b}}\neq\emptyset$ and $k_a\le k_b$. Then the assertions
\textup{(T1)}--\textup{(T5)} below hold; the relations they contain are collected in
\begin{equation}\label{4.27:eq-arrest-timing-a}
\begin{aligned}
&\text{(T1)}\quad k_a<k_b,\qquad X_{\mathfrak{a}}\subset X_{\mathfrak{b}},\qquad
W_{\mathfrak{a}}\subset W_{\mathfrak{b}};\\
&\text{(T2)}\quad k_a\le h_b,\qquad \ell_{\mathfrak{a}}\le k_a\le h_b\ \text{ on }W_{\mathfrak{a}};\\
&\text{(T4)}\quad k_{a_{s+1}}\ge k_{a_s}+N_{k_{a_{s+1}}}\quad\text{for the nested arrests of (T4).}
\end{aligned}
\end{equation}
\begin{itemize}
\item[(T1)] the step increases strictly, and the component and the piece do not decrease, as
stated in line \textup{(T1)} of \eqref{4.27:eq-arrest-timing-a};
\item[(T2)] $k_a\le h_b$; hence $\ell_{\mathfrak{a}}\le k_a\le h_b$ on $W_{\mathfrak{a}}$;
\item[(T3)] every integration region of $\mathfrak{b}$ lies in $\Lambda_{h_b}$, at least three
layers of $MR_{h_b+1}$-cubes away from $X_{\mathfrak{a}}$; consequently no later attempt
assigns a new level to, or integrates the variables at, a point of $X_{\mathfrak{a}}$;
\item[(T4)] if $\mathfrak{b}$ is an arrest, then
$I(\mathfrak{a})\cap I(\mathfrak{b})\subset\{k_a+1\}$, and this intersection is non-void if and
only if $h_b=k_a$; the live arrests whose pieces contain a common point are nested,
$\mathfrak{a}_1\subset\mathfrak{a}_2\subset\cdots$ in the sense of \textup{(T1)}, their steps
satisfy line \textup{(T4)} of \eqref{4.27:eq-arrest-timing-a},
and they are dissolved together;
\item[(T5)] the bands of distinct coexisting chains (a chain being the sequence of stages of one
attempt together with their bands; the coexisting chains are those of the live arrests and those
of the current step) are pairwise disjoint; within one chain the bands of levels
$\le k_0$ are disjoint, and at most $N_0$ bands overlap, where $k_0$ is the datum $k_0^{a}$
that Definition~\ref{4.27:def-arrests-pieces} attaches to the $\mathbf{R}$-site of the chain,
and $N_0$ is the smallest positive integer
with $L^{-N_0+1}MR_{k-N_0+1}=M$, $k$ the step of the chain \cite[p.~179]{B-CMP122a}; every
new boundary term $\mathbb{C}^{(i)}_k(X)$ of a stage $i$ of a chain has its domain $X$ meeting
the band of that stage.
\end{itemize}
If instead $W_{\mathfrak{a}}\cap W_{\mathfrak{b}}=\emptyset$ (with the other hypotheses
unchanged), then $X_{\mathfrak{a}}\cap X_{\mathfrak{b}}=\emptyset$.
\end{lemma}

\begin{proof}
\emph{(T1).} By Definition~\ref{4.27:def-arrests-pieces}, $X_{\mathfrak{a}}$ is a component
excluded from the class at the $\mathbf{R}$-site of step $k_a$ by one of the large-field exits
of \cite[\S1]{B-CMP122a}. Of such a component the paper says: ``They are not changed by the
subsequent operations, and we treat them in the same way as the remaining large field regions''.
Hence a component takes at most one exit in a step, and the site of step $k_a$ performs no
further attempt on $X_{\mathfrak{a}}$. Since $\mathfrak{b}\neq\mathfrak{a}$ and the pieces meet,
$\mathfrak{b}$ cannot be an attempt of step $k_a$; together with $k_a\le k_b$ this gives
$k_a<k_b$.

While $\mathfrak{a}$ is live, $X_{\mathfrak{a}}$ stays in the large-field region, by the same
sentence of \cite[\S1]{B-CMP122a}, and the large-field regions increase,
$Z_{k_a}\subset Z_{k_b}$, by \cite[(2.1)]{B-CMP119}. Thus $X_{\mathfrak{a}}$ is a connected
subset of $Z_{k_b}$, and it meets $X_{\mathfrak{b}}$, because
$\emptyset\neq W_{\mathfrak{a}}\cap W_{\mathfrak{b}}\subset X_{\mathfrak{a}}\cap X_{\mathfrak{b}}$.
The component of $Z_{k_b}$ that meets the connected set $X_{\mathfrak{a}}$ contains it, so
$X_{\mathfrak{a}}\subset X_{\mathfrak{b}}$. Moreover $\Omega^c_{k_a}\subset\Omega^c_{k_b}$, and
therefore
\begin{equation*}
W_{\mathfrak{a}}=X_{\mathfrak{a}}\cap\Omega^c_{k_a}\subset X_{\mathfrak{b}}\cap\Omega^c_{k_b}
=W_{\mathfrak{b}}.
\end{equation*}

\emph{The last sentence.} Let $W_{\mathfrak{a}}\cap W_{\mathfrak{b}}=\emptyset$ and $k_a\le k_b$.
The sets $X_{\mathfrak{a}}$ and $X_{\mathfrak{b}}$ are components of the nested sets
$Z_{k_a}\subset Z_{k_b}$. Components of nested sets are either nested or disjoint, so either
$X_{\mathfrak{a}}\cap X_{\mathfrak{b}}=\emptyset$ or $X_{\mathfrak{a}}\subset X_{\mathfrak{b}}$.
In the second case the inclusion $\Omega^c_{k_a}\subset\Omega^c_{k_b}$ would give
$\emptyset\neq W_{\mathfrak{a}}\subset W_{\mathfrak{b}}$, contradicting
$W_{\mathfrak{a}}\cap W_{\mathfrak{b}}=\emptyset$. Hence
$X_{\mathfrak{a}}\cap X_{\mathfrak{b}}=\emptyset$.

\emph{(T2).} Every exit introduces new large fields: the paper's words are ``In components with
the function $1-\chi^{(0)}_k$ we have introduced new large fields, therefore they do not satisfy
the condition (ii)'' \cite[\S1]{B-CMP122a}. The exit of $\mathfrak{a}$ thus introduces new large
fields at step $k_a$, inside $X_{\mathfrak{a}}$, hence, by (T1), inside $X_{\mathfrak{b}}$.
Condition (ii) of \cite[\S1]{B-CMP122a} for the component $X_{\mathfrak{b}}$ at step $k_b$ reads
``in the preceding $N$ renormalization steps no new large field regions were created inside this
component'', with $N=N_{k_b}$ and $h_b=k_b-N$; the paper describes these as ``the last $N$
one-step operations'', that is, the operations of steps $h_b+1,\dots,k_b$. So step $k_a$ is not
among them, and $k_a\le h_b$. On $W_{\mathfrak{a}}$ the frozen level $\ell_{\mathfrak{a}}$ is a
level assigned at step $k_a$ (Definition~\ref{4.27:def-arrests-pieces}), so
$\ell_{\mathfrak{a}}\le k_a\le h_b$ there.

\emph{(T3).} The integration regions of the stages of $\mathfrak{b}$, that is its bands, are the
sets $(Z_{j+1}\setminus Z''_{j+1})\cap X_{\mathfrak{b}}$, $j\ge h_b$
(Definition~\ref{4.27:def-arrests-pieces}); the final domain $\Lambda$ of a completed
$\mathbf{R}$, which dissolves the arrests of its component, is treated in (T4). The sets
$Z''_j$ of \cite[p.~179]{B-CMP122a} increase
with $j$. Hence, for every
$j\ge h_b$, we have $Z''_{h_b+1}\subset Z''_{j+1}$, and
\begin{equation*}
Z_{j+1}\setminus Z''_{j+1}\subset Z\setminus Z''_{h_b+1}=Z\cap\Omega^{\sim 5}_{h_b+1},
\end{equation*}
the equality being the definition $Z''_{h_b+1}=(\Omega^{\sim 5}_{h_b+1})^c\cap Z$ of
\cite[p.~179]{B-CMP122a}; here and below $A^{\sim n}$ denotes the set $A$ enlarged by $n$ layers
of cubes of the corresponding scale \cite[p.~179]{B-CMP122a}. By \cite[(3.5)]{B-CMP119} one has
$\Omega^c_{h_b+1}\supseteq(Z'_{h_b})^{\sim 8}$;
here $Z'_{h_b}$ is the enlarged large-field region of \cite[p.~195]{B-CMP122a}, which contains
$Z_{h_b}$. The same separation is written in \cite[p.~195]{B-CMP122a} as
\begin{equation*}
\Omega''_{h_b+1}=\Omega^{\sim 5}_{h_b+1}=((Z'_{h_b})^{\sim 3})^c.
\end{equation*}
Therefore $\Omega^{\sim 5}_{h_b+1}$ lies at least three layers of
$MR_{h_b+1}$-cubes away from $Z'_{h_b}$, hence from $Z_{h_b}\subseteq Z'_{h_b}$. By (T2),
$k_a\le h_b$, and since the large-field regions increase \cite[(2.1)]{B-CMP119},
\begin{equation*}
Z_{h_b}\supseteq Z_{k_a}\supseteq X_{\mathfrak{a}}.
\end{equation*}
So every integration region of $\mathfrak{b}$ lies at least three layers of
$MR_{h_b+1}$-cubes away from $Z_{h_b}$; hence it lies in $\Lambda_{h_b}=Z^c_{h_b}$ and, since
$X_{\mathfrak{a}}\subset Z_{h_b}$, at least three layers away from $X_{\mathfrak{a}}$.

It remains to see that the intermediate steps move nothing inside $X_{\mathfrak{b}}$. Consider an
$\mathbf{R}$ completed at a step $k$ with $h_b<k<k_b$. By \cite[(1.33)]{B-CMP122b} it re-sets
the determining set only on its own component. That component is not the one holding
$X_{\mathfrak{a}}$: otherwise $\mathfrak{a}$ would have been dissolved by it
(Definition~\ref{4.27:def-arrests-pieces}). Hence the regions $Z_j$, $Z''_j$ and $\Omega_j$ used
above, intersected with $X_{\mathfrak{b}}$, do not move. Consequently no later attempt assigns a
new level to, or integrates the variables at, a point of $X_{\mathfrak{a}}$.

\emph{(T4).} Let $\mathfrak{b}$ be an arrest. By Definition~\ref{4.27:def-arrests-pieces},
$I(\mathfrak{a})=\{h_a+1,\dots,k_a+1\}$ and $I(\mathfrak{b})=\{h_b+1,\dots,k_b+1\}$. By (T2),
\begin{equation*}
\min I(\mathfrak{b})=h_b+1\ge k_a+1=\max I(\mathfrak{a}).
\end{equation*}
Hence $I(\mathfrak{a})\cap I(\mathfrak{b})\subset\{k_a+1\}$, with equality if and only if
$h_b+1=k_a+1$, that is, $h_b=k_a$.

Let the live arrests whose pieces contain a common point be listed by increasing step. Any two of
them satisfy the hypotheses of the lemma. By (T1) their steps are distinct and their components
and pieces increase, so they are nested, $\mathfrak{a}_1\subset\mathfrak{a}_2\subset\cdots$. By
(T2) applied to the pair $(\mathfrak{a}_s,\mathfrak{a}_{s+1})$,
\begin{equation*}
k_{a_s}\le h_{a_{s+1}}=k_{a_{s+1}}-N_{k_{a_{s+1}}},
\end{equation*}
which is line (T4) of \eqref{4.27:eq-arrest-timing-a}.

A completed $\mathbf{R}$ at a step $k$, with $h=k-N_k$, integrates over the domain
$\Lambda=(\Omega^{\sim 4}_{k_0+1})^c\cap Z$ of \cite[(1.73)]{B-CMP122a}, a domain inside the
class $Z$ that is not to be confused with the small-field regions $\Lambda_j=Z^c_j$. By
\cite[(1.73)]{B-CMP122a}, $\Lambda\supset Z''_k$, and by \cite[p.~179]{B-CMP122a},
$Z''_k\supset Z''_h=Z_h$. By (T2), applied with this
$\mathbf{R}$ in place of $\mathfrak{b}$, every live arrest $\mathfrak{a}'$ whose piece lies in the
component of this $\mathbf{R}$ has $k_{a'}\le h$. Hence
\begin{equation*}
W_{\mathfrak{a}'}\subset X_{\mathfrak{a}'}\subset Z_{k_{a'}}\subset Z_h\subset\Lambda.
\end{equation*}
So $\Lambda$ holds every live piece of that component. The nested arrests
$\mathfrak{a}_1\subset\mathfrak{a}_2\subset\cdots$ all lie in the component holding the largest
of them. The first later completed $\mathbf{R}$ on that component therefore dissolves all of them
together.

\emph{(T5).} By Definition~\ref{4.27:def-arrests-pieces}, every band lies in its own component:
$R^{(i)}_{\mathfrak{a}}\subset X_{\mathfrak{a}}$. Take two distinct coexisting chains.

If their components are nested, disjointness is given by (T3). The bands of the later chain are
integration regions of its stages, and by (T3) these lie at least three layers away from the
earlier component, which contains the bands of the earlier chain.

If their components are not nested, two cases arise. Two chains of the current step belong to
distinct components of the class, which are disjoint, and so are their bands. If one of the two
is a live arrest, then by (T1) their pieces are disjoint. The last sentence of
the lemma then makes the components disjoint, and with them the bands.

Within one chain, the bands of levels $\le k_0$ are disjoint by \cite[\S1]{B-CMP122a}: ``In the
first case the integration regions are disjoint'', the first case being that of levels
$j\le k_0$. That at most $N_0$ bands overlap, with $N_0$ as in the statement, is stated in
\cite[\S1]{B-CMP122a}. That every
new boundary term $\mathbb{C}^{(i)}_k(X)$ of stage $i$ has its domain $X$ meeting the band of
that stage is the localisation \cite[(1.63)]{B-CMP122a}.
\end{proof}

\begin{definition}[Frozen terms, later operations and zones]\label{4.27:def-operations-zones}
Arrests $\mathfrak{a}$, their components $X_{\mathfrak{a}}$, pieces $W_{\mathfrak{a}}$, collars
$\mathfrak{c}_{\mathfrak{a}}$, the numbers $n$, $n^{+}$, $k_a$, the sets of frozen terms
$\mathfrak{F}(\mathfrak{a})$ and the spaces $\mathfrak{D}^{\flat}$ of stored terms are those of
Definition~\ref{4.27:def-arrests-pieces}. For a stored term $\tau$ we write $X_\tau$ for its
domain of dependence. A \emph{chain} is the sequence of stages of Ba{\l}aban's operation
$\mathbf{R}$ on one component; the site \textup{(L)} of a chain is the site at which it is
completed. We also use the rings of a completed chain and the set $\mathsf{Z}$ attached to the
site at which a chain is completed; both are fixed with the description of completion sites in this
chapter.

\emph{The class.} The \emph{class of frozen terms} is
\begin{equation}\label{4.27:eq-operations-zones-1}
\mathcal{C}_{\mathrm{arr}}:=\bigcup_{\mathfrak{a}\ \text{live}}\mathfrak{F}(\mathfrak{a}).
\end{equation}

\emph{Consumers and later operations.} Let $\mathfrak{a}$ be a live arrest and
$\tau\in\mathfrak{F}(\mathfrak{a})$. At every later operation $\mathfrak{O}$ the term $\tau$ is a
\emph{consumer}: its space $\mathfrak{D}^{\flat}$ must contain the image of the source of
$\mathfrak{O}$. The later operations are the following.
\begin{itemize}
\item[(O1)] A stage $l>n^{+}$, performed at step $k_a$ on a sibling of $X_{\mathfrak{a}}$, that
is, on another component of the class of $X_{\mathfrak{a}}$.
\item[(O2)] The site \textup{(L)} of a chain completed at a step $k\ge k_a$ on a component other
than the one holding $X_{\mathfrak{a}}$. The statement applied there is the first part of
\cite[(1.65)]{B-CMP122b} if $X_\tau$ meets that component; otherwise it is the statement of
\cite{B-CMP122b} for terms whose domain of dependence does not meet that component.
\item[(O3)] Every $\mathbf{T}$-step $k'\to k'+1$ with $k'\ge k_a$.
\item[(O4)] The stages of a later attempt $\mathfrak{b}\supset\mathfrak{a}$, that is, one with
$X_{\mathfrak{b}}\supset X_{\mathfrak{a}}$ as in \eqref{4.27:eq-arrest-timing-a}.
\item[(O4$'$)] The stages of a later attempt $\mathfrak{b}$ at a step $k_b>k_a$ on a class that
does not hold $X_{\mathfrak{a}}$.
\item[(O5)] Dissolution: the site \textup{(L)} of the chain completed on the component of a piece
met by $X_\tau$.
\end{itemize}
That every later operation at which $\tau$ is read is one of these is the completeness
statement for the later operations of this chapter.
A term $\tau$ whose domain of dependence $X_\tau$ meets the component of a completed chain is
absorbed there into the term $\mathbf{V}(Y)$, indexed by a domain $Y$, of Ba{\l}aban's
construction \cite{B-CMP122b}; from that point on it ceases to exist as a stored term.

\emph{Zones.} Let $\mathfrak{O}$ be one of the operations above, at a site of step $k$,
and let $x\in X_\tau$. Write $\mu_\tau(x)$ for the level that the condition defining the
space $\mathfrak{D}^{\flat}$ of $\tau$ prescribes at the point $x$. The \emph{source level} at $x$
is the level carried at $x$ by the source of $\mathfrak{O}$. The \emph{zone} of $x$ at
$\mathfrak{O}$ is determined by the following list, read in order.
\begin{itemize}
\item Zone $\mathrm{Z}4$: the point $x$ lies in a piece that is dissolved at $\mathfrak{O}$.
\item Zone $\mathrm{Z}1$: the point $x$ lies in no piece dissolved at $\mathfrak{O}$, and
$x\in Z^{\mathrm{on}}_{k_a}$. Such an $x$ lies either in a live piece of the class of $\tau$, or,
at step $k_a$, in a component of that class which is running (operation (O1)) or completing
(operation (O2)). At such $x$ the condition defining the space of $\tau$ is a class condition,
that is, a condition of the class part of a stage space.
\end{itemize}
Otherwise $x\in\mathfrak{T}_{k_a}$, and the zone is the first of the following that applies.
\begin{itemize}
\item Zone $\mathrm{Z}3\text{-L}$: the point $x$ lies in the set $\mathsf{Z}$ of a chain completing
at $\mathfrak{O}$.
\item Zone $\mathrm{Z}3\mathrm{a}$: the source level at $x$ exceeds $\mu_\tau(x)$, and this excess
was produced by $\mathbf{T}$-steps or by the rings of a completed chain lying off $\mathsf{Z}$.
\item Zone $\mathrm{Z}3\mathrm{b}$: the source level at $x$ exceeds $\mu_\tau(x)$, and this excess
was produced by the stages of a later attempt.
\item Zone $\mathrm{Z}2^{\circ}$: the source level at $x$ equals $\mu_\tau(x)$, the operation
$\mathfrak{O}$ is of type (O4) or (O4$'$), and $x\in Z^{\mathrm{on}}_{k}$.
\item Zone $\mathrm{Z}2$: all remaining cases.
\end{itemize}
Identifying each zone with the set of points of $X_\tau$ that lie in it at $\mathfrak{O}$, we have,
at every later operation $\mathfrak{O}$,
\begin{equation}\label{4.27:eq-operations-zones-a}
\begin{gathered}
\mathfrak{O}\ \text{of one of the types}\ (\mathrm{O}1),\dots,(\mathrm{O}5),(\mathrm{O}4'):\\
X_\tau=\mathrm{Z}4\sqcup\mathrm{Z}1\sqcup\mathrm{Z}3\text{-L}\sqcup\mathrm{Z}3\mathrm{a}
\sqcup\mathrm{Z}3\mathrm{b}\sqcup\mathrm{Z}2^{\circ}\sqcup\mathrm{Z}2 .
\end{gathered}
\end{equation}
The zones form a partition of $X_\tau$, for three reasons. First, $X_\tau$ is contained in
$Z^{\mathrm{on}}_{k_a}\cup\mathfrak{T}_{k_a}$. Second, the level carried at a point of $X_\tau$ by
the successive sources never decreases from one later operation to the next. Third, after step $k_a$
the set $X_\tau\cap Z^{\mathrm{on}}_{k_a}$ consists of live pieces only.

\emph{Pieces of other classes.} At pieces of classes other than that of $\tau$, the arrested
factor takes the place of the corresponding stage variable, in the source of $\mathfrak{O}$
and in the constituents of $\tau$ born later alike. The differences between the source and those
constituents that enter the bounds are therefore unchanged, or room is gained in them.

\emph{Exclusions.} Two kinds of component contribute nothing to $\mathcal{C}_{\mathrm{arr}}$. The
first are the components that fail the first of the conditions, labelled \textup{(i)}, under which
Ba{\l}aban's operation $\mathbf{R}$ is attempted on a component.
The second are the second-class components in the sense of \cite{B-CMP122b}, on
which Ba{\l}aban's operation $\mathbf{R}$ was completed. The second-class components count, for
the timing statements of Lemma~\ref{4.27:lem-arrest-timing}, only as steps at which new large
fields are created.
\end{definition}

\begin{proposition}[The permanent deficit on an arrested piece]\label{4.27:prop-permanent-deficit}
Let $\mathfrak{a}$ be an arrest in the sense of Definition~\ref{4.27:def-arrests-pieces}, taking place
at step $k_a$, with $N=N_{k_a}$ the number of stages as in that definition and $h_a=k_a-N$. Let
$n^+$ be an integer with $1\le n^+<N$, and put
$j^+=h_a+n^+$. Let $\tau\in\mathfrak{F}(\mathfrak{a})$ be a frozen term whose domain is the domain of
stage $n^+$. Let $x\in W_{\mathfrak{a}}$ lie in the first shell that has not been processed, that is,
at level $j^+$, where the $\mathsf{L}_0$-weight is equal to one both in the domain of $\tau$ and in the
standard source domain. With the factors $s^{(k)}_m$, $\varphi^{(n),(k)}_m$ and the constant $\beta$ of
Definition~\ref{4.27:not-factors-radii-coupling}, for every $k'\ge k_a$,
\begin{equation}\label{4.27:eq-permanent-deficit-1}
s^{(k')}_{j^+}-\varphi^{(n^+),(k_a)}_{j^+}\ \ge\ \frac{\beta}{2}.
\end{equation}
\end{proposition}

\begin{proof}
Definition~\ref{4.27:not-factors-radii-coupling} fixes the factors and the stage factors that we use here.
For levels $k$ and $m$,
\begin{equation*}
s^{(k)}_m=1-\beta\bigl(1-2^{-(k-m)}\bigr),\qquad
\varphi^{(n),(k)}_m=s^{(k)}_m-\beta\sum_{i=h+1}^{h+n}2^{-|m-i|},
\end{equation*}
where $h=k-N_k$. We take $k=k_a$, so that $h=h_a$, and we take $n=n^+$ and $m=j^+$.

\emph{The stage sum.} For $h_a+1\le i\le h_a+n^+=j^+$ we have $|j^+-i|=j^+-i$. As $i$ runs over this
range, $r=j^+-i$ runs over $0,1,\dots,n^+-1$. Summing the geometric series gives
\begin{equation}\label{4.27:eq-permanent-deficit-2}
\sum_{i=h_a+1}^{h_a+n^+}2^{-(j^+-i)}=\sum_{r=0}^{n^+-1}2^{-r}=2-2^{1-n^+}.
\end{equation}
By definition of the stage factor, this yields
\begin{equation*}
s^{(k_a)}_{j^+}-\varphi^{(n^+),(k_a)}_{j^+}=\beta\bigl(2-2^{1-n^+}\bigr).
\end{equation*}

\emph{The drift of the standard factor.} We have $j^+\le k_a$. By the formula for $s^{(k)}_m$,
\begin{equation*}
s^{(k_a)}_{j^+}-s^{(k')}_{j^+}=\beta\bigl(2^{-(k_a-j^+)}-2^{-(k'-j^+)}\bigr).
\end{equation*}
The second term in the bracket is positive. Moreover $k_a-j^+=k_a-h_a-n^+=N-n^+$. Hence
\begin{equation}\label{4.27:eq-permanent-deficit-3}
s^{(k_a)}_{j^+}-s^{(k')}_{j^+}\le\beta\,2^{-(N-n^+)}.
\end{equation}

\emph{Conclusion.} We write the left side of \eqref{4.27:eq-permanent-deficit-1} as the sum of
$s^{(k')}_{j^+}-s^{(k_a)}_{j^+}$ and $s^{(k_a)}_{j^+}-\varphi^{(n^+),(k_a)}_{j^+}$. We bound the first
summand from below by \eqref{4.27:eq-permanent-deficit-3} and evaluate the second by
\eqref{4.27:eq-permanent-deficit-2}. This gives
\begin{equation*}
s^{(k')}_{j^+}-\varphi^{(n^+),(k_a)}_{j^+}
\ge\beta\bigl(2-2^{1-n^+}-2^{-(N-n^+)}\bigr).
\end{equation*}
Since $n^+\ge1$, we have $2^{1-n^+}\le1$. Since $n^+<N$, we have $2^{-(N-n^+)}\le\tfrac12$. The right
side is therefore at least $\beta\bigl(2-1-\tfrac12\bigr)=\beta/2$.
\end{proof}

Thus, at $x$, the factor of the standard source domain exceeds the factor of the domain of $\tau$. It does
so permanently, that is, at every step $k'\ge k_a$. At the same time, the map of a later operation
$\mathfrak{O}$, one of the operations of Definition~\ref{4.27:def-operations-zones}, acts at $x$ as the
identity times layers. Consequently the inclusions \cite[(2.41)(ii)]{B-CMP119} and
\cite[(1.65)]{B-CMP122b} cannot be applied, with the domains as assigned, to the terms of the class
$\mathcal{C}_{arr}$ of Definition~\ref{4.27:def-operations-zones}. The inheritance of domains by the
frozen terms under later operations holds on the sets $\mathfrak{T}_k$ only.

\begin{definition}[The arrested reading of the spaces]\label{4.27:def-arrested-spaces}
The reading of the spaces defined here is in force for each arrest from the step at which the arrest
occurs until its dissolution. It is a declaration on the scaffold of the construction. It involves
neither the complex source variables $z$ nor the slot variables. It is the first hypothesis of the
arrest theorem.

\emph{Notation.} The following notation is that of
Notation~\ref{4.27:not-factors-radii-coupling}:
\begin{itemize}
\item the constant $\beta$, the factors $s^{(k)}_m$ and $\varphi^{(n),(k)}_m$ (with $m$ a level), and
the radii $r_a(m)$ of the direct entries $Q_a$, $a=1,\dots,7$;
\item the stage spaces $\mathfrak{S}^{(k)}_n$ with their two readings $(\Sigma^{\flat})$ and
$(\Sigma^{*})$;
\item the sets $Z^{\mathrm{on}}_k=Z\cap\Omega_k^{c}$ and $\mathfrak{T}_k=\Omega_k\cup Z^{c}$;
\item the stage domains
\begin{equation*}
\mathfrak{D}^{\flat}_n(X)=\mathfrak{N}^{(n)}(X)\cap\mathfrak{L}_n(X\cap Z^{\mathrm{on}})
\cap\mathfrak{I}_n(X\cap\mathfrak{T}),
\end{equation*}
with class part $\mathfrak{L}_n$ and foreign part $\mathfrak{I}_n$.
\end{itemize}
In the reading $(\Sigma^{*})$ of that notation, the second and fourth entries of the stage spaces
are required for every triple $(p,X',\mathfrak{s})$ as fixed there, where $\mathfrak{s}$ ranges over the
determining-set structures on $X'$ given by the current one and all its predecessors in the
construction.

The following objects are those of Definition~\ref{4.27:def-arrests-pieces}:
\begin{itemize}
\item for an arrest $\mathfrak{a}$: its component $X_{\mathfrak{a}}$, its piece
$W_{\mathfrak{a}}=X_{\mathfrak{a}}\cap\Omega^{c}_{k_{\mathfrak{a}}}$, its step $k_{\mathfrak{a}}$, its
numbers $n$ and $n^{+}$, and its range $I(\mathfrak{a})$;
\item the range sum
$\sigma_{\mathfrak{a}}(\ell)=\sum_{i\in I(\mathfrak{a})}2^{-|\ell-i|}$;
\item the frozen terms $\mathfrak{F}(\mathfrak{a})$;
\item the frozen structure of $\mathfrak{a}$, that is, the level, the weight (a power of
$\mathsf{L}_0$) and the cube family assigned at $x\in W_{\mathfrak{a}}$ by
$\mathfrak{S}^{(k_{\mathfrak{a}})}_{n^{+}}$.
\end{itemize}
We write the weight as $\mathsf{w}$. Accordingly, $(\ell,\mathsf{w},\mathcal{Q})(x)$ denotes the level,
the weight and the cube family in force at $x$. We write $\mathfrak{B}_k$ for the determining set of
step $k$, and $\mathfrak{B}^{(n)}_k$ for the determining set of stage $n$ of step $k$.

\emph{Forward references.} The parameter $\gamma_0^{(k)}$ is defined below in the lemma on weights, the
dated size parameter and heredity, see~\eqref{4.27:eq-weights-heredity-a}. The rider space
$\mathfrak{G}(\mathfrak{B})$ of a determining set $\mathfrak{B}$ is defined in clause~(e) of that lemma,
and so is the dating rule for the index-$0$ source datum. On the stratum
$\Omega_0\setminus\Omega_1$ of raw points (the points of index $0$), this rule dates that datum at birth
and freezes it.

The arrested reading consists of the following six items.
\begin{enumerate}
\item[(a)] \emph{Structure freeze.} Let $\mathfrak{a}$ be an arrest. Every space written after
$\mathfrak{a}$ reads at $x\in W_{\mathfrak{a}}$ the frozen structure $(\ell,\mathsf{w},\mathcal{Q})(x)$
of the innermost live arrest holding $x$. This applies in particular to the following spaces:
\begin{itemize}
\item the standard spaces $\tilde{U}_{k'}(X,\tilde\alpha)$ and the spaces after an operation
$\mathbf{T}$;
\item the sources $\tilde{U}^{c}_{k}(Y_4,\tilde\alpha)$ of an operation $\mathbf{R}$, taken over the
set $Y_4$ of that operation;
\item the attempt spaces of sibling attempts and of enclosing attempts;
\item the class parts $\mathfrak{L}_l$ and the tubes;
\item the domains of newly born terms.
\end{itemize}
The innermost live arrest holding $x$ is well defined, since by clause~(T4) of
Lemma~\ref{4.27:lem-arrest-timing} the live arrests holding a point are nested. By clause~(T3) of the
same lemma no later attempt re-levels a point of $X_{\mathfrak{a}}$. The frozen level is therefore the
actual level of $x$.

Now let $x\in W_{\mathfrak{a}}\cap(\Omega_0\setminus\Omega_1)$ be a raw point, and let $\mathfrak{c}$ be
the innermost live arrest holding $x$. By the nesting in clause~(T4),
$k_{\mathfrak{c}}\le k_{\mathfrak{a}}$. At such a point the index-$0$ source datum is frozen as
\begin{equation}\label{4.27:eq-arrested-spaces-1}
T(x):=\gamma_0^{(k_{\mathfrak{c}})},
\end{equation}
with $T(x)=\gamma_0^{(k_{\mathfrak{a}})}$ when $\mathfrak{a}$ itself is innermost. The stratum condition
on the source at $x$ is read as $\eta^{3}|\mathbb{J}(x)|<T(x)$. Here $\eta$ is the lattice unit
at index $0$, and $\eta^{3}|\mathbb{J}(x)|$ is the source $1$-form in lattice
units. This reading is the dating rule, in the
form in which it enters the rider space $\mathfrak{G}$ of clause~(e) of the lemma on weights, the dated
size parameter and heredity; see also item~(f) below.

Every raw point of $X_{\mathfrak{a}}$ is a point of $W_{\mathfrak{a}}$. Indeed, write
$\Gamma_0=\Omega_1^{c}$. Since $\Omega_1^{c}\subset\Omega^{c}_{k_{\mathfrak{a}}}$, we have
\begin{equation*}
\Gamma_0\cap X_{\mathfrak{a}}\subset X_{\mathfrak{a}}\cap\Omega^{c}_{k_{\mathfrak{a}}}=W_{\mathfrak{a}}.
\end{equation*}
The constant $c$ entering the stage spaces $\mathfrak{S}^{(k)}_n$ is taken identically equal to $1$.

The family of structures of the reading $(\Sigma^{*})$ is frozen together with the structure. At every
set $X'$ meeting $W_{\mathfrak{a}}$, the structures $\mathfrak{s}$ admitted on $X'\cap W_{\mathfrak{a}}$
include the frozen determining set $\mathfrak{B}^{(n^{+})}_{k_{\mathfrak{a}}}$ and its predecessors in
the piece's own run, namely $\mathfrak{B}^{(n')}_{k_{\mathfrak{a}}}$ for $n'\le n^{+}$, and
$\mathfrak{B}_{k_{\mathfrak{a}}}$. On $X'\setminus W_{\mathfrak{a}}$ they are those of the current
determining set.

\item[(b)] \emph{Factor.} For every step $k'$ set
\begin{equation}\label{4.27:eq-arrested-spaces-2}
F^{(k')}(x):=s^{(k')}_{\ell(x)}
-\beta\sum_{\mathfrak{a}\ \text{live},\ x\in W_{\mathfrak{a}}}\sigma_{\mathfrak{a}}(\ell(x)).
\end{equation}
The arrested attempt spaces $\mathfrak{S}^{(k),\mathrm{arrest}}_n$ are read at locations in
$Z^{\mathrm{on}}_k$ only, as follows.
\begin{itemize}
\item At points of no live piece they are $\mathfrak{S}^{(k)}_n$.
\item Let $x$ lie in a live piece. Such a piece is held there by a component of the class, and it is
either older than the attempt (clause~(T2) of Lemma~\ref{4.27:lem-arrest-timing}) or of the same class.
The attempt spaces then read at $x$ the frozen structure of item~(a) and the factor
\begin{equation}\label{4.27:eq-arrested-spaces-3}
\varphi^{(n),(k)}_{\ell(x)}
-\beta\sum_{\substack{\mathfrak{a}\ \text{live},\ x\in W_{\mathfrak{a}}\\ k_{\mathfrak{a}}<k}}
\sigma_{\mathfrak{a}}(\ell(x)).
\end{equation}
Pieces of the same class contribute no further term inside the run.
\end{itemize}
At a live piece off the class, the stage domain is $\mathfrak{N}^{(n)}\cap\mathfrak{I}_n$. It is
read with item~(a), and with $F^{(k)}$ in place of $s^{(k)}$ in the source threshold
$\theta_{\mathrm{src}}$.

\item[(c)] \emph{Spaces.} For a set $X$ and a step $k'$, the arrested space
$\tilde{U}^{\mathrm{arrest}}_{k'}(X)$ is the set of pairs $(\mathbb{U},\mathbb{J})$ on $X$ that satisfy,
at every point $x$ of $X$ in a live piece,
\begin{equation}\label{4.27:eq-arrested-spaces-a}
Q_a(x)<F^{(k')}(x)\,\mathsf{w}(x)\,r_a(\ell(x)),\qquad a=1,\dots,7,
\end{equation}
together with the cube conditions of $\mathcal{Q}(x)$ at weight $\mathsf{w}(x)$. Off live pieces the
conditions are those of \cite[(2.34)--(2.39)]{B-CMP119} of step $k'$, taken verbatim. The readings
$(\Sigma^{\flat})$ and $(\Sigma^{*})$ are unchanged. The sets $\tilde{U}^{\mathrm{arrest}}_{k'}(X)$ are
open and invariant under gauge transformations.

\item[(d)] \emph{Domains of new terms.} Let a term be born at a later operation $\mathfrak{O}$. It is
assigned a domain --- that of \cite[(2.41)(ii)]{B-CMP119}, that of \cite[(1.68)]{B-CMP122b}, or a
stage domain $\mathfrak{D}^{\flat}_n$ --- and receives here the arrested reading of that domain, as
follows.
\begin{itemize}
\item For the domains of \cite[(2.41)(ii)]{B-CMP119}, the reading is taken with the enlargement of the
factor fixed for newly born terms. On a live piece the enlarged factor is
\begin{equation*}
F^{(k'+1)}(x)+\theta_S\,\beta\,2^{-(k'+1-\ell(x))}.
\end{equation*}
\item The domains of \cite[(1.68)]{B-CMP122b} receive the arrested reading.
\item It applies also to the stage domains $\mathfrak{D}^{\flat}_n$, with
$\mathfrak{L}_n:=\mathfrak{S}^{(k),\mathrm{arrest}}_n$.
\end{itemize}
Frozen terms, that is, the members of $\mathfrak{F}(\mathfrak{a})$ for a live arrest $\mathfrak{a}$,
keep their domains.

\item[(e)] \emph{Dissolution.} Consider the completed operation $\mathbf{R}$ at step $k$ on the
component of $X_{\mathfrak{a}}$. Let $\Lambda$ be the region of that operation and $h=k-N_k$. By
clause~(T4) of Lemma~\ref{4.27:lem-arrest-timing},
$W_{\mathfrak{a}}\subset Z_h\subset\Lambda$. At this operation the arrest $\mathfrak{a}$, and every piece
nested with it, is dropped after the step. On $\Lambda$ the standard structure of the new determining
set $\mathfrak{B}_k$ applies.

\item[(f)] \emph{Ambient rider.} This item is imposed as part of the definition, not derived. It is a
condition of the same kind as the weighted ambient condition.

Every arrested space, source and domain over a determining set $\mathfrak{B}$ is read with the ambient
condition in the following form: its members are restrictions of pairs of $\mathfrak{G}(\mathfrak{B})$
(clause~(e) of the lemma on weights, the dated size parameter and heredity). A stage domain
$\mathfrak{D}^{\flat}_n$ reads $\mathfrak{G}(\mathfrak{B}^{(n)}_k)$, where $\mathfrak{B}^{(n)}_k$ is the
structure of its own class part $\mathfrak{L}_n$.

On the domain itself nothing new is asked, since the direct factors are at most $1+2\beta<5$. The
domains shrink (no supremum grows). They remain open, invariant under gauge transformations, and stable
under restriction. The condition involves neither the complex source variables $z$ nor the slot
variables.

By clause~(e3) of the same lemma, the condition implies three things:
\begin{itemize}
\item the ambient condition at the step $k$ at which the domain was born;
\item the bound $a_{\flat}\le\gamma_0^{(k)}$ on every ball $B(x)$, where $a_{\flat}$ and the balls
$B(x)$ are those of that clause;
\item the weighted ambient condition.
\end{itemize}
On the stratum $\Omega_0\setminus\Omega_1$, the index-$0$ source datum of every arrested space, source,
domain and rider is read by the dating rule. It is frozen at the value $\gamma_0$ of the innermost live
piece, as in item~(a) above and in clause~(e) of that lemma.

Of a rider read in this way, what later sources and images are proved to meet is the following:
\begin{itemize}
\item the sub-rider $\mathfrak{G}^{-J_0}$ of clause~(e3) of that lemma;
\item $\mathfrak{G}_1$ together with the list of stratum conditions in the rider clause of the arrest
theorem, the weighted ambient condition included.
\end{itemize}
The two residual cases of the dating rule are excepted.
\end{enumerate}
\end{definition}

\begin{lemma}[Positivity of the arrested factors]\label{4.27:lem-factor-positivity}
Let $F^{(k')}(x)$ be the arrested factor of item (b) of the arrested reading of the spaces
(Definition~\ref{4.27:def-arrested-spaces}). The factor
$F^{(k')}(x)$, and every factor of an attempt space read at a location of $Z^{\mathrm{on}}$ (both as
in item (b) of Definition~\ref{4.27:def-arrested-spaces}), are
bounded below by
\begin{equation}\label{4.27:eq-factor-positivity-a}
\varphi_{\mathrm{arrest}}:=1-4.51\,\beta=0.098 .
\end{equation}
\end{lemma}

\begin{proof}
Fix $x$ and write $\ell=\ell(x)$. By item (b) of Definition~\ref{4.27:def-arrested-spaces},
\begin{equation*}
F^{(k')}(x)=s^{(k')}_{\ell}-\beta\sum_{\mathfrak{a}\ \mathrm{live},\ x\in\mathfrak{a}}
\sigma_{\mathfrak{a}}(\ell),
\qquad
\sigma_{\mathfrak{a}}(\ell)=\sum_{i\in I(\mathfrak{a})}2^{-|\ell-i|},
\end{equation*}
where $I(\mathfrak{a})=\{h_a+1,\dots,k_a+1\}$ is the range of the arrest $\mathfrak{a}$ and $\beta$ is
the constant multiplying the range sums in item (b) of Definition~\ref{4.27:def-arrested-spaces}.
The factors $s^{(k')}_{\ell}$ fixed in Notation~\ref{4.27:not-factors-radii-coupling} satisfy
$s^{(k')}_{\ell}\ge 1-\beta$.

By clause (T4) of Lemma~\ref{4.27:lem-arrest-timing}, \eqref{4.27:eq-arrest-timing-a}, the live
arrests holding $x$ are nested, $\mathfrak{a}_1\subset\mathfrak{a}_2\subset\cdots$, with
$k_{a_{r+1}}\ge k_{a_r}+N_{k_{a_{r+1}}}$, and $I(\mathfrak{a}_r)\cap I(\mathfrak{a}_{r'})\subset\{k_{a_r}+1\}$
for $r<r'$. Hence an integer lying in three of the ranges, with indices $r<r'<r''$, would equal both
$k_{a_r}+1$ and $k_{a_{r'}}+1$, which is impossible since $k_{a_r}<k_{a_{r'}}$. So the ranges are
disjoint except at the indices $k_{a_r}+1$, each of which lies in at most two of them, and
\begin{equation*}
\sum_r\sigma_{\mathfrak{a}_r}(\ell)
\le\sum_{i\in\mathbb{Z}}2^{-|\ell-i|}+\sum_r 2^{-|\ell-(k_{a_r}+1)|}.
\end{equation*}
The first sum equals $1+2\sum_{j\ge1}2^{-j}=3$. For the second we use $\ell\le k_{a_1}$ (the level
read at $x$ is the frozen level of the innermost live arrest $\mathfrak{a}_1$, by item (a) of
Definition~\ref{4.27:def-arrested-spaces}, and $\ell_{\mathfrak{a}_1}\le k_{a_1}$ on
$W_{\mathfrak{a}_1}$ by clause (T2) of Lemma~\ref{4.27:lem-arrest-timing}), so that
$|\ell-(k_{a_r}+1)|=k_{a_r}+1-\ell$, and the gaps satisfy
$k_{a_{r+1}}-k_{a_r}\ge N_{k_{a_{r+1}}}\ge 6$, so that $k_{a_r}+1-\ell\ge 1+6(r-1)$. Therefore
\begin{equation*}
\sum_r\sigma_{\mathfrak{a}_r}(\ell)\le 3+\tfrac12\sum_{t\ge0}2^{-6t}=3+\tfrac{32}{63}\le 3.51,
\end{equation*}
and $F^{(k')}(x)\ge 1-\beta-3.51\,\beta=\varphi_{\mathrm{arrest}}$.

Now let $x$ be a location of $Z^{\mathrm{on}}$ and consider the factor of an attempt space of an
attempt $\mathfrak{b}$ read there, with range $I(\mathfrak{b})=\{h_b+1,\dots,k_b+1\}$. In this factor
the attempt's own range enters as one more range, besides the ranges of the older live pieces
holding $x$ (possibly none, in which case the factor contains the single range $I(\mathfrak{b})$ and
$\sigma_{\mathfrak{b}}(\ell)\le\sum_{i\in\mathbb{Z}}2^{-|\ell-i|}=3$). By clause (T2) of
Lemma~\ref{4.27:lem-arrest-timing}, $h_b\ge k_a$ for every older
live piece $\mathfrak{a}$ holding $x$, so $I(\mathfrak{b})$ begins at $h_b+1\ge k_a+1$.
It therefore meets the range of an older piece at most at the index $k_{a_{\bar r}}+1$ of the
outermost older live piece $\mathfrak{a}_{\bar r}$, which is one of the indices carried by the second
sum above; so the estimate of the preceding paragraph holds for the chain enlarged by
$I(\mathfrak{b})$. With $s^{(k)}_{\ell}\ge1-\beta$, where $k=k_b$ is the density of the attempt
$\mathfrak{b}$, this shows that the attempt-space factor of item (b) of
Definition~\ref{4.27:def-arrested-spaces} is at least $\varphi_{\mathrm{arrest}}$.
\end{proof}

The positivity $\varphi_{\mathrm{arrest}}>0$ requires $\beta<0.2217$, and the bound does not depend
on the number of
live arrests holding $x$. At a live piece outside $Z^{\mathrm{on}}$ the attempt-space factor can be
$\le-0.175$; this is why
item (b) of Definition~\ref{4.27:def-arrested-spaces} reads the attempt spaces at locations of
$Z^{\mathrm{on}}$ only.

\begin{lemma}[Room of frozen domains over arrested spaces]\label{4.27:lem-frozen-room}
Let $W_{\mathfrak{a}}$ be any live piece of class $k_a$, namely the piece of the live arrest
$\mathfrak{a}$, with range $I(\mathfrak{a})=\{h_a+1,\dots,k_a+1\}$, which fixes the integer $h_a$,
and with range sum defined for every level $\ell$ by
$\sigma_{\mathfrak{a}}(\ell):=\sum_{i\in I(\mathfrak{a})}2^{-|\ell-i|}$; the range sums
$\sigma_{\mathfrak{b}}$ of the other arrests are defined in the same way. For an arrest or attempt
$\mathfrak{b}$ we write $k_b$ for its class, namely the index $k$ of the application of
Ba{\l}aban's operation $\mathbf{R}$ at scale $k$ to which $\mathfrak{b}$ belongs; we also call $k_b$
the \emph{density} of $\mathfrak{b}$ (not to be confused with the effective density $\rho_k$), and
we say that a term is born at density $k$ if it is created
at the step of index $k$. We write $h_b$ for the integer with
$I(\mathfrak{b})=\{h_b+1,\dots,k_b+1\}$; for an arrest $\mathfrak{c}$ we write $k_c$ likewise. Let
$N_{k_a}$ be the declared number of stages at class $k_a$, and let $n^{+}_{\mathfrak{a}}$ be the
stage index of the arrest $\mathfrak{a}$, fixed with the arrest. The frozen
structure of $\mathfrak{a}$ is the level, the weight and the cube family that the space
$\mathfrak{S}^{(k_a)}_{n^{+}_{\mathfrak{a}}}$ assigns at the points of $W_{\mathfrak{a}}$. Throughout,
$\beta$ is the constant multiplying the range sums in the arrested factor $F^{(k)}$ of
Definition~\ref{4.27:def-arrested-spaces}. Let $x\in W_{\mathfrak{a}}$ and
$\ell:=\ell(x)$. Let $E$ be any term whose clause at $x$ is that of the class domain of stage $n$,
where $0\le n\le N_{k_a}$. Examples of such terms are:
\begin{itemize}
\item the frozen terms $\mathfrak{F}(\mathfrak{a})$ of that piece: its new boundary terms of stage
$i$, with $n=i$, and its remaining frozen terms, those not among its new boundary terms, whose clause
at $x$ is that of the class domain of stage $n^{+}_{\mathfrak{a}}$, with $n=n^{+}_{\mathfrak{a}}$;
\item the frozen terms $\mathfrak{F}(\mathfrak{b})$ of sibling arrests, and the terms with
determining set $\mathbb{B}''_{k_a}$;
\item the terms of the convergent expansion of \cite{B-CMP119} born at density at most $k_a$,
with $n=0$.
\end{itemize}
By the \emph{room} of the clause of $E$ over a space at $x$ we mean the excess of the factor by which
the clause of $E$ multiplies the thresholds of the direct entries $Q_a(x)$, $a\le 7$, of
Definition~\ref{4.27:def-arrested-spaces} at $x$ over the factor by which
that space multiplies the same thresholds.
\begin{enumerate}
\item[(a)] Suppose the clause of $E$ reads the frozen structure at $x$. This is always the case
for $n\ge n^{+}_{\mathfrak{a}}$. For $n<n^{+}_{\mathfrak{a}}$ it is the case at the points not
raised in the stages $n+1,\dots,n^{+}_{\mathfrak{a}}$. Then for every $k'\ge k_a$ the room of the
clause of $E$ over
$\tilde{U}^{\mathrm{arrest}}_{k'}$ at $x$, measured in the units of the frozen level $\ell$, is
\begin{equation}\label{4.27:eq-frozen-room-a}
\begin{split}
\rho_n(x;k'):={}&\beta\sum_{i=h_a+n+1}^{k_a+1}2^{-|\ell-i|}
+\beta\bigl(2^{-(k_a-\ell)}-2^{-(k'-\ell)}\bigr)\\
&+\beta\sum_{\mathfrak{b}\ \mathrm{live},\ x\in W_{\mathfrak{b}},\ k_b>k_a}
\sigma_{\mathfrak{b}}(\ell).
\end{split}
\end{equation}
It obeys the following bounds:
\begin{align*}
\rho_n(x;k')&\ge\beta\,2^{-(k_a+1-\ell)},\\
\rho_n(x;k')&\ge\beta\max\bigl\{2^{-|\ell-h_a-n-1|},\,2^{-(k_a-\ell)}\bigr\}
\quad\text{if } n<N_{k_a},\\
\rho_n(x;k')&\ge\beta\,2^{-(k_a-\ell)}\quad\text{if } k'>k_a.
\end{align*}
\item[(b)] At raised points the clause of $E$ contains that of
$\mathfrak{S}^{(k_a),\mathrm{arrest}}_{n^{+}_{\mathfrak{a}}}$. For every $k'\ge k_a$ the room over
$\tilde{U}^{\mathrm{arrest}}_{k'}$, measured in the units of the
frozen level $\ell$, is at least $\rho_{n^{+}_{\mathfrak{a}}}(x;k')$. Moreover, at $x$ the source
member $\mathbb{J}$ of the pairs $(\mathbb{U},\mathbb{J})$ on which these spaces are defined is at
most $0.74$ times the source threshold of the domain of $E$ at $x$.
\item[(c)] The same bounds hold against the space $\mathfrak{S}^{(k_b),\mathrm{arrest}}_{n'}$ of an
enclosing attempt $\mathfrak{b}$. Against a sibling space $\mathfrak{S}^{(k_a),\mathrm{arrest}}_{n'}$
with $n'\ge n$, the room is
\begin{equation*}
\varphi^{(n),(k_a)}_{\ell}-\varphi^{(n'),(k_a)}_{\ell}\ge 0 ,
\end{equation*}
where $\varphi^{(n),(k_a)}_{\ell}$ and $\varphi^{(n'),(k_a)}_{\ell}$ are the stage factors of the
stages $n$ and $n'$ at class $k_a$ and level $\ell$, through which
Definition~\ref{4.27:def-arrested-spaces} reads the attempt spaces at points of live pieces.
\end{enumerate}
\end{lemma}

\begin{proof}
\emph{(a).} Suppose the clause of $E$ reads the frozen structure at $x$. The arrested reading of
the spaces (Definition~\ref{4.27:def-arrested-spaces}) reads $\tilde{U}^{\mathrm{arrest}}_{k'}$ at $x$
in the frozen structure as well. Hence both sides carry the same frozen level, weight and cube family
$\mathcal{Q}(x)$ at $x$.

By the reading of the cube clauses attached to $\mathfrak{a}$, which
Definition~\ref{4.27:def-arrested-spaces} leaves unchanged, every cube of $\mathcal{Q}(x)$ enters
both sides with the same triple, formed by the cube
and its two parameters. The cube clauses therefore coincide. The
inclusion of $\tilde{U}^{\mathrm{arrest}}_{k'}$ in the clause of $E$ at $x$ thus reduces to the
inequality between the factors multiplying the common thresholds of the direct entries. The room is
the difference of these two factors.

By Definition~\ref{4.27:def-arrested-spaces}, the class domain of stage $n$ is read through the
space $\mathfrak{S}^{(k_a),\mathrm{arrest}}_{n}$; its factor at $x$ is the stage factor
$\varphi^{(n),(k_a)}_{\ell}$ less $\beta$ times the sum of $\sigma_{\mathfrak{c}}(\ell)$ over the
live pieces $W_{\mathfrak{c}}\ni x$ with $k_c<k_a$, that is, over the older pieces containing $x$.
Write $s^{(k)}_{\ell}$ for the term of the arrested factor $F^{(k)}$ of
Definition~\ref{4.27:def-arrested-spaces} from which the range sums of the live pieces are
subtracted. By the definitions of the stage factors $\varphi^{(n),(k)}_{\ell}$ of the stage spaces
and of the terms $s^{(k)}_{\ell}$, which Definition~\ref{4.27:def-arrested-spaces} takes as given,
the stage factors and these terms satisfy
\begin{align}
\varphi^{(n),(k_a)}_{\ell}&=s^{(k_a)}_{\ell}-\beta\sum_{i=h_a+1}^{h_a+n}2^{-|\ell-i|}
\quad(0\le n\le N_{k_a}),\label{4.27:eq-frozen-room-3}\\
s^{(k_a)}_{\ell}-s^{(k')}_{\ell}&=\beta\bigl(2^{-(k_a-\ell)}-2^{-(k'-\ell)}\bigr)
\quad(k'\ge k_a).\label{4.27:eq-frozen-room-4}
\end{align}
By \eqref{4.27:eq-frozen-room-3}, the factor of the class domain of stage $n$ at $x$ is
\begin{equation}\label{4.27:eq-frozen-room-1}
s^{(k_a)}_{\ell}-\beta\sum_{i=h_a+1}^{h_a+n}2^{-|\ell-i|}
-\beta\sum_{\mathfrak{c}\ \mathrm{live},\ x\in W_{\mathfrak{c}},\ k_c<k_a}
\sigma_{\mathfrak{c}}(\ell),
\end{equation}
where the last sum runs over the older pieces containing $x$.

The factor of $\tilde{U}^{\mathrm{arrest}}_{k'}$ at $x$ is the arrested factor $F^{(k')}(x)$ of
Definition~\ref{4.27:def-arrested-spaces}. We split its sum over the live pieces containing $x$
into three parts: the piece $W_{\mathfrak{a}}$ itself; the same older pieces as in
\eqref{4.27:eq-frozen-room-1}; and the younger pieces. These parts exhaust the sum: no other live
piece of class $k_a$ contains $x$, since by the first of the relations
\eqref{4.27:eq-arrest-timing-a} of Lemma~\ref{4.27:lem-arrest-timing} a live arrest and a different
attempt whose pieces meet have different classes. The older pieces are still live by the fourth of
the relations \eqref{4.27:eq-arrest-timing-a} of Lemma~\ref{4.27:lem-arrest-timing}, since live
arrests holding a point are nested and dissolve together. This gives
\begin{equation}\label{4.27:eq-frozen-room-2}
\begin{split}
F^{(k')}(x)={}&s^{(k')}_{\ell}-\beta\sigma_{\mathfrak{a}}(\ell)
-\beta\sum_{\mathfrak{c}\ \mathrm{live},\ x\in W_{\mathfrak{c}},\ k_c<k_a}
\sigma_{\mathfrak{c}}(\ell)\\
&-\beta\sum_{\mathfrak{b}\ \mathrm{live},\ x\in W_{\mathfrak{b}},\ k_b>k_a}
\sigma_{\mathfrak{b}}(\ell).
\end{split}
\end{equation}

We subtract \eqref{4.27:eq-frozen-room-2} from \eqref{4.27:eq-frozen-room-1}. The sums over the
older pieces cancel. Since $\sigma_{\mathfrak{a}}(\ell)$ is the sum of $2^{-|\ell-i|}$ over
$i\in I(\mathfrak{a})$, removing the indices $h_a+1,\dots,h_a+n$ leaves the first sum in
\eqref{4.27:eq-frozen-room-a}. By \eqref{4.27:eq-frozen-room-4}, the difference
$s^{(k_a)}_{\ell}-s^{(k')}_{\ell}$ equals the middle term
$\beta(2^{-(k_a-\ell)}-2^{-(k'-\ell)})$. The sum over the younger pieces changes sign and
gives the last term. This proves the identity \eqref{4.27:eq-frozen-room-a}.

For the lower bounds we first note that each term of \eqref{4.27:eq-frozen-room-a} is
non-negative. The sums consist of positive terms. The middle term is non-negative because
$k'\ge k_a$.

Since $n\le N_{k_a}$ and, for the arrest $\mathfrak{a}$, $h_a+N_{k_a}\le k_a$ (the relation between
the range $I(\mathfrak{a})$ and the declared number of stages $N_{k_a}$, fixed with the range), the
lower limit of the
first sum satisfies $h_a+n+1\le k_a+1$; hence the index $i=k_a+1$ lies in its range. Its term is
$\beta\,2^{-|\ell-k_a-1|}$, which equals $\beta\,2^{-(k_a+1-\ell)}$ because the frozen level
satisfies $\ell\le k_a$ on $W_{\mathfrak{a}}$ (cf.\ the second of the
relations \eqref{4.27:eq-arrest-timing-a} of Lemma~\ref{4.27:lem-arrest-timing}, where this
inequality is recorded for the frozen level on $W_{\mathfrak{a}}$). This gives the
first bound.

The index $i=h_a+n+1$ always lies in the range, and its term is
$\beta\,2^{-|\ell-h_a-n-1|}$. If $n<N_{k_a}$, the same relation gives $h_a+n+1\le k_a$, so the
index $i=k_a$ lies in the range as well, with
term $\beta\,2^{-(k_a-\ell)}$. Each of these terms alone bounds $\rho_n(x;k')$ from below, which
gives the second bound.

If $k'>k_a$, then $k'\ge k_a+1$, so that $2^{-(k'-\ell)}\le 2^{-(k_a+1-\ell)}$. Keeping the term
$i=k_a+1$ and the middle term, we obtain
\begin{equation*}
\rho_n(x;k')\ge\beta\bigl(2^{-(k_a+1-\ell)}+2^{-(k_a-\ell)}-2^{-(k'-\ell)}\bigr)
\ge\beta\,2^{-(k_a-\ell)} ,
\end{equation*}
which is the third bound.

\emph{(b).} Let $x$ be a point raised in one of the stages $n+1,\dots,n^{+}_{\mathfrak{a}}$. By
the third of the relations \eqref{4.27:eq-arrest-timing-a} of Lemma~\ref{4.27:lem-arrest-timing},
a raised point lies in no older piece, so the domains there are the published ones, read as in
Definition~\ref{4.27:def-arrested-spaces}. By the descent of the stage spaces, according to which
the space of a later stage lies in that of an earlier one, applied from stage $n$ to stage
$n^{+}_{\mathfrak{a}}$, the clause of $E$ at $x$ contains that of
$\mathfrak{S}^{(k_a),\mathrm{arrest}}_{n^{+}_{\mathfrak{a}}}$. The same descent statement also gives
the second assertion of (b): it bounds the source member $\mathbb{J}$ at $x$ by $0.74$ times the
source threshold of the domain of $E$ at $x$.

The space $\mathfrak{S}^{(k_a),\mathrm{arrest}}_{n^{+}_{\mathfrak{a}}}$ reads the frozen structure at
$x$ by definition, since the frozen structure is the one that
$\mathfrak{S}^{(k_a)}_{n^{+}_{\mathfrak{a}}}$ assigns on $W_{\mathfrak{a}}$. Part (a) therefore
applies to it with $n^{+}_{\mathfrak{a}}$ in place
of $n$. Together with the containment just obtained, this shows that the room of the clause of $E$
is at least $\rho_{n^{+}_{\mathfrak{a}}}(x;k')$.

\emph{(c).} Let $\mathfrak{b}$ be an enclosing attempt. The factor of
$\mathfrak{S}^{(k_b),\mathrm{arrest}}_{n'}$ at $x$ is $F^{(k_b)}(x)$ diminished by $\beta$ times the
terms $2^{-|\ell-i|}$ for $i$ in a part of $I(\mathfrak{b})$. It therefore differs from
$F^{(k_b)}(x)$ only by further subtractions. The room over it is at least the room over
$\tilde{U}^{\mathrm{arrest}}_{k_b}$, to which (a) and (b) apply with $k'=k_b\ge k_a$.

Now let $\mathfrak{S}^{(k_a),\mathrm{arrest}}_{n'}$, $n'\ge n$, be a sibling space. Abbreviate
$\varphi^{(n)}_{\ell}:=\varphi^{(n),(k_a)}_{\ell}$ and
$\varphi^{(n')}_{\ell}:=\varphi^{(n'),(k_a)}_{\ell}$ for the stage factors of
Definition~\ref{4.27:def-arrested-spaces}. By that definition, the factor of the sibling space at
$x$ is $\varphi^{(n')}_{\ell}$ less $\beta$ times the sum of $\sigma_{\mathfrak{c}}(\ell)$ over the
live pieces $W_{\mathfrak{c}}\ni x$ with $k_c<k_a$. The class domain of stage $n$ has the factor
$\varphi^{(n)}_{\ell}$ less the same sum. The room is therefore
$\varphi^{(n)}_{\ell}-\varphi^{(n')}_{\ell}$.

By \eqref{4.27:eq-frozen-room-3}, applied at stage $n$ and at stage $n'$,
\begin{equation*}
\varphi^{(n)}_{\ell}-\varphi^{(n')}_{\ell}=\beta\sum_{i=h_a+n+1}^{h_a+n'}2^{-|\ell-i|}\ge 0 .
\qedhere
\end{equation*}
\end{proof}

\begin{lemma}[Weights, the dated size parameter and heredity; stated; proof incomplete at the
steps marked $(\ast)$]\label{4.27:lem-weights-heredity}
This lemma uses the arrests, pieces and frozen structures of
Definition~\ref{4.27:def-arrests-pieces}; the factors, the direct entries $Q_a$, the radii $r_a$
and the logarithmic coupling variables $\mathsf{x}_j$, $\mathsf{t}_j$, $\mathsf{t}_\gamma$,
$\mathsf{A}(\mathsf{t})$ of Notation~\ref{4.27:not-factors-radii-coupling}; the declarations of
Definition~\ref{4.27:def-one-class-per-step}; and the arrested reading of the spaces of
Definition~\ref{4.27:def-arrested-spaces}. For a step $j$ let $N_0(j)$ be the number $N_0$
attached to density $j$, so that $k_0=j-N_0(j)$ at density $j$, and put
\begin{equation*}
  w_{*}(j):=\mathsf{L}_0^{2N_0(j)},\qquad
  \bar\alpha_j:=\max_{m\le j}\mathsf{A}(\mathsf{t}_m),
\end{equation*}
with $w_{*}(j):=1$ if step $j$ has no $\mathbf{R}$-site. Let $\mathfrak{a}$ be a live arrest, let
$x\in W_{\mathfrak{a}}$, and let $\ell:=\ell_{\mathfrak{a}}(x)$ and $w_{\mathfrak{a}}(x)$ be the
frozen level and the frozen weight at $x$.
\begin{enumerate}
\item[(a)] $w_{\mathfrak{a}}(x)\le w_{*}(k_a)$, and $w_{\mathfrak{a}}(x)=1$ unless
  $k_0^{a}<\ell\le k_a$. A point raised by $g$ levels has weight at most $\mathsf{L}_0^{2g}$; a
  point that is not raised has weight $1$.
\item[(b)] With $R_j$ the number of stages that Ba{\l}aban's large-field operation $\mathbf{R}$
  assigns to a region created at step $j$, and $r$ the exponent of
  Notation~\ref{4.27:not-factors-radii-coupling},
  \begin{equation*}
    w_{*}(j)\le L^2R_j<L^2(L+1)\,\mathsf{t}_j^{\,r},
  \end{equation*}
  and $w_{*}(j)\le\mathsf{l}(\mathsf{x}_n)$ for every $n\le j$.
\item[(c)] Put
  \begin{equation}\label{4.27:eq-weights-heredity-a}
    \gamma_0^{(k)}:=5\max_{j\le k}w_{*}(j)\,\bar\alpha_j
    \;\le\;\bar\gamma_0(\mathsf{t}_k):=5(1+\beta_0)^2\,\Psi(\mathsf{t}_k),
    \qquad
    \Psi(\mathsf{t}):=L^2(L+1)\,C\,\mathsf{t}^{\,r+q}e^{-\mathsf{t}/2}.
  \end{equation}
  Here $C=\max_i C_i^{\mathrm{rad}}$ is the constant of the radius function $\mathsf{A}$ of
  Notation~\ref{4.27:not-factors-radii-coupling}.
  The function $\bar\gamma_0$ is an explicit function of $\mathsf{t}$ and of constants fixed
  before $\gamma$ only; it does not depend on the sequence of couplings of the construction. It
  is decreasing on $\mathsf{t}\ge\mathsf{t}_\gamma$ and tends to $0$ as $\mathsf{t}\to\infty$.
\item[(d)] Put $w^{\vee}(x):=\max\{w_{*}(k),w(x)\}$. Let $\mathfrak{N}^{\vee}$ be the ambient
  space of pairs $\mathfrak{N}$ with its direct clauses on $Z^{\mathrm{on}}$ read as
  \begin{equation*}
    Q_a<5\,w^{\vee}(x)\,r_a(\ell(x)),
  \end{equation*}
  and at density $k$ let $\mathfrak{K}_P:=\mathfrak{K}_{109}(\mathfrak{B};\gamma_0^{(k)})$, where,
  for a determining set $\mathfrak{B}$ of density $k$ and $\gamma>0$,
  $\mathfrak{K}_{109}(\mathfrak{B};\gamma)$ denotes the space of pairs $(\mathbb{U},\mathbb{J})$
  of the type of \cite{B-CMP109} over $\mathfrak{B}$, with radius $\gamma$ in the source
  variable. Every
  arrested space lies in $\mathfrak{N}^{\vee}\subset\mathfrak{K}_P$, and
  $a_\flat\le\gamma_0^{(k)}$ there, by the definition of $\gamma_0^{(k)}$, at every age and every
  coupling; here $a_\flat$ is the quantity bounded in the ceiling
  $C^b(a_\flat+\kappa_{\max})\le1$ of Definition~\ref{4.27:def-arrest-hypotheses}. The
  statement on the stage maps that the arrest theorem takes as an input, and which holds at any
  determining set, holds there under the three ceilings listed together in
  Definition~\ref{4.27:def-arrest-hypotheses} and read there at $\bar\gamma_0$: the ceiling on
  $\gamma$ of that statement on the stage maps, the ceiling $K_\Delta\gamma_0\le\frac12$
  ($K_\Delta$ being its constant), and the ceiling $C^b(a_\flat+\kappa_{\max})\le1$.
\end{enumerate}

\emph{Rule for $\gamma_0$.} Where $\gamma_0$ enters a membership condition it is read as
$\gamma_0^{(k)}$. Where it enters a constant, a floor or a ceiling it is read as
$\bar\gamma_0(\mathsf{t})$, under a single supremum over $\mathsf{t}\ge\mathsf{t}_\gamma$ taken
together with the other functions of $\mathsf{t}$ in that condition, all at the same
$\mathsf{t}$. Since $\mathfrak{K}_{109}(\mathfrak{B};\cdot)$ increases with its argument,
$\mathfrak{K}_P$ may equally be formed at $\bar\gamma_0(\mathsf{t}_k)$.

\emph{Stored domains.} The domain of a stored term $\tau$ is read on the set $\mathfrak{B}_\tau$
of its birth and at its density of birth $k_\tau$. Thus the space $\mathfrak{K}_P$ of the ambient
condition inside the domain of $\tau$ is
$\mathfrak{K}_{109}(\mathfrak{B}_\tau;\gamma_0^{(k_\tau)})$. The bounds on $\tau$ were proved
under ceilings at $\mathsf{t}_{k_\tau}$. Read at a later density $k'$, they would require
$K_\Delta\gamma_0^{(k')}$ to be bounded by the value at $\mathsf{t}_{k_\tau}$ of one of the
functions of $\mathsf{t}$ of which the ceilings of Definition~\ref{4.27:def-arrest-hypotheses}
are formed, which is a condition on the age. That
every later field meets the reading at birth is clause (e).

For stored objects built from the function $\mathbb{H}$ of Ba{\l}aban's change of variables that
carry no clause of their own, $\gamma_0$ may also be left free. Clause (a) of the theorem on the
function $\mathbb{H}$, and the bound on $\mathbb{H}$ on which these stored objects depend, hold at
every $\gamma\le\gamma^{\top}:=\bar\gamma_0(\mathsf{t}_\gamma)$,
because the ceilings are suprema and $\bar\gamma_0$ decreases. Moreover the space
$\mathfrak{K}_{109}(\mathfrak{B}_\tau;\gamma^{\top})$, which contains every later
$\mathfrak{N}^{\vee}$, does not depend on $k$. From now on $\mathfrak{N}$ means
$\mathfrak{N}^{\vee}$.

\begin{enumerate}
\item[(e)] \emph{Heredity.} Order the determining sets of the construction by time, the sets of
  the stages included, each frozen on the pieces live at its date. Let
  $\mathfrak{B}\preccurlyeq\mathfrak{B}'$ be two such sets, $\mathfrak{B}$ of date not later
  than that of $\mathfrak{B}'$, of densities $k\le k'$; the assertions (e1)--(e3) below concern
  this pair.

  A \emph{date} is an operation, the operations being taken in the order in which they are
  performed. At density $k$ the dates are:
  \begin{itemize}
  \item the $\mathbf{T}$-step from $k-1$ to $k$, which writes $\mathbb{B}_k$;
  \item the stages $1<\dots<N_k$ of the single class of the step (one class and one pair $(h,k)$
    per step, by the declarations of Definition~\ref{4.27:def-one-class-per-step}; running
    components share the stage, \cite[(1.63)]{B-CMP122a});
  \item the arrests, at their stages;
  \item the completion site.
  \end{itemize}
  The set of a date is one structure on $\Omega_0$. It is read at $y$ as the structure freeze of
  Definition~\ref{4.27:def-arrested-spaces} reads it, in the following order of precedence:
  \begin{itemize}
  \item the frozen $(\ell,w,\mathcal{Q})$ of the innermost live piece containing $y$;
  \item otherwise, the stage structure on the on-class part $Z^{\mathrm{on}}_k$ of the running
    class;
  \item otherwise, that of $\mathbb{B}_k$ (\cite[(2.2)]{B-CMP119}, and \cite[(1.33)]{B-CMP122b} at
    completed sites).
  \end{itemize}
  A stored domain reads the set of its date of birth. Consecutive dates differ at each $y$ by
  exactly one of the four elementary moves (t1)--(t4) of the lemma on the order of dates, or not
  at all.

  A point $y$ is \emph{held} at a date by a live piece $W_{\mathfrak{c}}$ if
  $y\in W_{\mathfrak{c}}$, where $\mathfrak{c}$ is attached to a date not later than the given one
  and is not dissolved at one. The points of the stratum
  $\Omega_0(\mathfrak{B})\setminus\Omega_1(\mathfrak{B})$ are called \emph{raw}. Put
  \begin{equation*}
    \theta^{\mathfrak{B}}_a(y):=w_{\mathfrak{B}}(y)\,r_a(\ell_{\mathfrak{B}}(y)),
  \end{equation*}
  where $w_{\mathfrak{B}}$ is the actual weight and not $w^{\vee}$.

  Let $\mathfrak{G}(\mathfrak{B})$ be the set of pairs
  $(\mathbb{U},\mathbb{J})=(e^{i\eta A'}U,\mathbb{J})$ on $\Omega_0(\mathfrak{B})$, of the form of
  the members of $\mathfrak{K}_{109}(\mathfrak{B};\cdot)$, that satisfy the following three
  clauses.
  \begin{enumerate}
  \item[(i)] $Q_a(y)<5\theta^{\mathfrak{B}}_a(y)$ for $a\in\{1,3,5,6,7\}$ at every
    $y\in\Omega_1(\mathfrak{B})$.
  \item[(ii)] Every cube of $\mathcal{Q}_{\mathfrak{B}}$ satisfies the cube clause at five times
    the weight of its family. A cube carries the weight of its family, not that of a point: by
    \cite[p.~190, (1.67)]{B-CMP122a}, cubes $\square\subset\Omega_m$ with
    $\square\cap\Omega^{c}_{m+1}\ne\emptyset$ carry the weight $\mathsf{L}_0^{2(j-m)}$. The clause
    for a cube $\square$ of family $(\lambda_\square,w_\square)$
    (\cite[(1.65), (1.67), (1.69)]{B-CMP122a}) is this: there is a gauge transformation $u$ on
    $\square\cap\Omega_0(\mathfrak{B})$ with $U^{u}=e^{i\eta A}$ (as in \cite{B-CMP99b}) and
    \begin{equation*}
      L^{\lambda_\square}\eta\,|A|,\ \ (L^{\lambda_\square}\eta)^2\,|\nabla^{\eta}A|
      \;<\;5\,w_\square\,B\,C\,M\,\alpha_{0,\lambda_\square};
    \end{equation*}
    here $B$ and $C$ denote the constants of Ba{\l}aban's cube clause in
    \cite[(1.65), (1.67), (1.69)]{B-CMP122a}.
  \item[(iii)] On $\Omega_0(\mathfrak{B})\setminus\Omega_1(\mathfrak{B})$, the index-$0$ clause of
    $\mathfrak{K}_{109}(\mathfrak{B};\cdot)$ holds. Ba{\l}aban's spaces start at $m=1$
    \cite[p.~190]{B-CMP122a}, and $\Gamma_0=\Omega_1^{c}$ consists of raw points
    \cite[(2.2)]{B-CMP119}. This clause has two parts.
    \begin{itemize}
    \item Its $U$-members: \cite[(3.35)]{B-CMP99b} on the index-$0$ cubes, which do not depend
      on the date, and \cite[(3.37)]{B-CMP99b} at index $0$; we write $(3.35)_0$ and
      $(3.37)_0$ for these index-$0$ instances. They do not involve $\gamma$, and
      they are taken with the standing they have in the ambient condition on stored terms.
    \item Its member for the source variable, read as
      $\eta^{3}|\mathbb{J}(y)|<T^{\mathfrak{B}}(y)$ (the source variable in lattice units, the
      lattice unit at index $0$ being $\eta$) at every
      $y\in\Omega_0(\mathfrak{B})\setminus\Omega_1(\mathfrak{B})$. Here $T^{\mathfrak{B}}(y)$ is
      defined in \eqref{4.27:eq-weights-heredity-b} below. By clause (T1) of
      Lemma~\ref{4.27:lem-arrest-timing}, the innermost live piece containing $y$ is the oldest
      one.
    \end{itemize}
  \end{enumerate}

  The index-$0$ datum for the source variable is frozen with the structure, as in the structure
  freeze of Definition~\ref{4.27:def-arrested-spaces}. Every arrested space, source, domain and
  ambient clause written while $\mathfrak{c}$ lives reads $T(y)=\gamma_0^{(k_{\mathfrak{c}})}$ at
  the raw points of $W_{\mathfrak{c}}$ that it holds innermost. Here
  $\Gamma_0\cap X_{\mathfrak{c}}\subset W_{\mathfrak{c}}$, since
  $\Omega_1^{c}\subset\Omega^{c}_{k_{\mathfrak{c}}}$ by the nesting \cite[(2.1)]{B-CMP119}.
  Moreover
  $T^{\mathfrak{B}}\le\gamma_0^{(k)}$, since $k_{\mathfrak{c}}\le k$ and $\gamma_0^{(\cdot)}$ is
  non-decreasing, being a maximum over $j\le\cdot$.

  Let $\mathfrak{G}_1(\mathfrak{B})$ be the set of pairs on $\Omega_0(\mathfrak{B})$ that obey the
  clauses of $\mathfrak{G}(\mathfrak{B})$ at every $y\in\Omega_1(\mathfrak{B})$ and at every cube
  of $\mathcal{Q}_{\mathfrak{B}}$; it has no member on the stratum. Let
  $\mathfrak{G}^{-J_0}(\mathfrak{B})$ be the set of pairs of $\mathfrak{G}_1(\mathfrak{B})$ that
  obey in addition the $U$-members of the stratum: $(3.35)_0$ on the index-$0$ cubes and
  $(3.37)_0$ at every $y\in\Omega_0(\mathfrak{B})\setminus\Omega_1(\mathfrak{B})$. Thus
  $\mathfrak{G}^{-J_0}(\mathfrak{B})$ is $\mathfrak{G}(\mathfrak{B})$ less its member for the
  source variable on the stratum, and
  \begin{equation}\label{4.27:eq-weights-heredity-b}
  \begin{gathered}
    T^{\mathfrak{B}}(y):=
    \begin{cases}
      \gamma_0^{(k_{\mathfrak{c}})}, & y \text{ held at the date of } \mathfrak{B}
        \text{ by a live piece,}\\
      & \quad W_{\mathfrak{c}} \text{ the innermost such},\\
      \gamma_0^{(k)}, & y \text{ held at the date of } \mathfrak{B}\text{ by no live piece},
    \end{cases}\\
    \mathfrak{G}(\mathfrak{B})
    =\mathfrak{G}^{-J_0}(\mathfrak{B})\cap
    \bigl\{\eta^{3}|\mathbb{J}(y)|<T^{\mathfrak{B}}(y)\ \text{for all }
    y\in\Omega_0(\mathfrak{B})\setminus\Omega_1(\mathfrak{B})\bigr\}\\
    \subset\mathfrak{G}^{-J_0}(\mathfrak{B})\subset\mathfrak{G}_1(\mathfrak{B}).
  \end{gathered}
  \end{equation}

  The sets $\mathfrak{G}_1$ and $\mathfrak{G}^{-J_0}$ involve no $\gamma_0$. They are membership
  conditions, and they enter no constant, floor, ceiling or room. The full
  $\mathfrak{G}(\mathfrak{B})$, with $T^{\mathfrak{B}}$, is the reading of every arrested space,
  source and domain under the ambient clause of Definition~\ref{4.27:def-arrested-spaces}.

  As sets, $\mathfrak{G}_1(\mathfrak{B})\not\subset\mathfrak{K}_{109}(\mathfrak{B};\gamma_0^{(k)})$,
  since $\mathfrak{G}_1$ has no clause on the stratum. The proof of (d), on the other hand, places
  the clauses of $\mathfrak{G}_1(\mathfrak{B})$ inside those of
  $\mathfrak{K}_{109}(\mathfrak{B};\gamma_0^{(k)})$ at every point of $\Omega_1(\mathfrak{B})$ and
  at every cube.

  At raw points held by no live piece, the threshold $\gamma_0^{(k)}$ plays two different roles.
  It is kept as a membership condition at the density $k$ of $\mathfrak{B}$, so that
  $\mathfrak{G}(\mathfrak{B})\subset\mathfrak{K}_{109}(\mathfrak{B};\gamma_0^{(k)})$ holds on all
  of $\Omega_0(\mathfrak{B})$. It is suspended as a hereditary clause: the member for the source
  variable in the first inclusion of (e3) is not asserted there.

  With these definitions, the following hold.
  \begin{enumerate}
  \item[(e1)] $\ell_{\mathfrak{B}'}\ge\ell_{\mathfrak{B}}$ and
    $w_{\mathfrak{B}'}\le\mathsf{L}_0^{2(\ell_{\mathfrak{B}'}-\ell_{\mathfrak{B}})}\,w_{\mathfrak{B}}$.
  \item[(e2)] $\theta^{\mathfrak{B}'}_a\le\theta^{\mathfrak{B}}_a$, and
    $\theta^{\mathfrak{B}'}_a\le0.436\,\theta^{\mathfrak{B}}_a$ where the level rose. Thresholds
    never rise in time.
  \item[(e3)] As sets on all of $\Omega_0(\mathfrak{B})$,
    \begin{equation*}
      \mathfrak{G}(\mathfrak{B}')\subset\mathfrak{G}^{-J_0}(\mathfrak{B}),\qquad
      \mathfrak{G}^{-J_0}(\mathfrak{B}')\subset\mathfrak{G}^{-J_0}(\mathfrak{B}),\qquad
      \mathfrak{G}_1(\mathfrak{B}')\subset\mathfrak{G}_1(\mathfrak{B}).
    \end{equation*}
    These inclusions rest on three ingredients.
    \begin{itemize}
    \item On $\Omega_1(\mathfrak{B})$: (e2).
    \item On the cubes of $\mathcal{Q}_{\mathfrak{B}}$: the nesting condition of the declarations
      of Definition~\ref{4.27:def-one-class-per-step}, read at one locus.
    \item For the $U$-members of the stratum: at points raw at both dates, the clause does not
      involve $\gamma$ and is the same at the two dates. At points entering
      $\Omega_1(\mathfrak{B}')$, they follow from the level clauses of
      $\mathfrak{G}_1(\mathfrak{B}')$ through the elementary move (t4) of the lemma on the order
      of dates, the nesting condition, and the statement that yields the index-$0$ gauge members
      $(3.35)_0$ from these level clauses, used by its statement.
    \end{itemize}
    Further:
    \begin{enumerate}
    \item[(e3.i)] Every pair of $\mathfrak{G}_1(\mathfrak{B}')$ obeys every stratum member of
      $\mathfrak{G}(\mathfrak{B})$, the member for the source variable included, at every point
      of $(\Omega_0\setminus\Omega_1)(\mathfrak{B})\cap\Omega_1(\mathfrak{B}')$.
    \item[(e3.ii)] Every pair of $\mathfrak{G}(\mathfrak{B}')$ obeys the stratum member for the
      source variable of $\mathfrak{G}(\mathfrak{B})$ at every point raw at both dates that is
      held at the date of $\mathfrak{B}$ by a live piece, since there
      $T^{\mathfrak{B}'}(y)=T^{\mathfrak{B}}(y)$.
    \item[(e3.iii)] Every pair of $\mathfrak{G}(\mathfrak{B}')$ also obeys it at every point raw
      at both dates, held or not, when $\mathfrak{B}'$ is of the density $k$ of $\mathfrak{B}$ or
      $y$ is held by a live piece at some date of density $k$. The datum of an unheld point
      cannot change within a density.
    \end{enumerate}
    Hence $\mathfrak{G}(\mathfrak{B}')\subset\mathfrak{G}(\mathfrak{B})$ holds on all of
    $\Omega_0(\mathfrak{B})$, with one exception. The exception is the member for the source
    variable at points raw at both dates that are held by no live piece at any date of density
    $k$, when $\mathfrak{B}'$ is of a density $k'>k$. There the inclusion of that member is not
    asserted $(\ast)$: the threshold $\gamma_0^{(k)}$ is kept as a membership condition and
    suspended as a hereditary clause across an increase of density.

    Finally, $\mathfrak{G}(\mathfrak{B})\subset\mathfrak{K}_{109}(\mathfrak{B};\gamma_0^{(k)})$ as
    sets on all of $\Omega_0(\mathfrak{B})$, because $T^{\mathfrak{B}}\le\gamma_0^{(k)}$; and
    $a_\flat\le\gamma_0^{(k)}$ on $\mathfrak{G}(\mathfrak{B})$.
  \item[(e4)] The real point $U$ of the lemma on the real critical configuration, and every pair
    that obeys at every $y$
    the direct clauses of an arrested space, source or domain over $\mathfrak{B}$, lie in
    $\mathfrak{G}_1(\mathfrak{B})$. They do so with $3.6$ units to spare, the unit at $y$ being
    $w_{\mathfrak{B}}(y)r_a(\ell_{\mathfrak{B}}(y))$. This is read through the ambient clause of
    Definition~\ref{4.27:def-arrested-spaces}. Such pairs obey the clauses of
    $\mathfrak{G}_1(\mathfrak{B})$ wherever the space imposes a clause, that is, at the points of
    its own domain $X\cap\Omega_1(\mathfrak{B})$ and at its cubes, and there they are
    restrictions of pairs of $\mathfrak{G}_1(\mathfrak{B})$.

    On the stratum $\Omega_0\setminus\Omega_1$ no direct clause is imposed. The attempt spaces of
    \cite[p.~190]{B-CMP122a} start at $m=1$, and the index range of
    \cite[(2.34)--(2.39)]{B-CMP119} is fixed by declaration. There a member of an arrested space,
    source or domain lies in $\mathfrak{G}(\mathfrak{B})$ by the reading of the ambient clause,
    not by proof. Of the real point, only $A'=0$, and hence $(3.37)_0$, is immediate on the
    stratum. Its index-$0$ gauge $(3.35)_0$ and the inequality
    $\eta^{3}|\mathcal{J}(U)(y)|<T^{\mathfrak{B}}(y)$ are not proved here; they form part of the
    residual of (e4) marked $(\ast)$ below.

    Now let $\mathfrak{O}$ be a later operation, with source set $\mathfrak{B}'$ and target set
    $\mathfrak{B}^t\preccurlyeq\mathfrak{B}'$. Here $\mathfrak{B}^t$ is the set in force
    immediately before the map of $\mathfrak{O}$, and $\mathfrak{B}'$ is the set in force
    immediately after it, over which the arrested source of $\mathfrak{O}$ is written. These are:
    \begin{itemize}
    \item for a stage $l$ of an attempt at $k$: $\mathbb{B}^{(l-1)}_k$ and $\mathbb{B}^{(l)}_k$;
    \item for $\mathbf{T}$ from $k'$ to $k'+1$: $\mathbb{B}_{k'}$ and $\mathbb{B}_{k'+1}$;
    \item for the completion site of a chain completed at $k$ (operations $\mathrm{O}2$,
      $\mathrm{O}5$ of the enumeration $\mathrm{O}1$--$\mathrm{O}5$, $\mathrm{O}4'$ of the
      operations that, after an arrest, compose a stored term with a map of the pair):
      $\mathbb{B}''_k=\mathbb{B}^{(N)}_k$ dated during the step, and its image under
      \cite[(1.33)]{B-CMP122b}.
    \end{itemize}
    In the last case, $\mathbb{B}''_k$ is frozen on every piece live during the step, dissolving
    ones included. On the core of the completed chain it consists of the levels $o(x)$ that the
    chain assigns there and the frozen weights; on $\mathsf{Z}$ and on the rings around the core
    it is the stage $N$ of the chain. The re-levelled
    $\mathbb{B}_k$ of \cite[(1.33)]{B-CMP122b} dates after the step (the dissolution clause of
    Definition~\ref{4.27:def-arrested-spaces}). It is not a target of $\mathrm{O}2$ or
    $\mathrm{O}5$: it is the set $\mathfrak{B}'$ there, and the target of $\mathrm{O}3$ from $k$
    to $k+1$.

    Let $\varpi$ be a member of the arrested source of $\mathfrak{O}$ over $Y'$ with continuation
    $\tilde\varpi\in\mathfrak{G}(\mathfrak{B}')$. Then the image of $\varpi$ under $\mathfrak{O}$
    has the continuation $\tilde\varsigma$ defined as follows.
    \begin{itemize}
    \item On $Y'$, $\tilde\varsigma$ is the image.
    \item Off $Y'$, the field of $\tilde\varsigma$ is that of $\tilde\varpi$.
    \item Off $Y'$, the source variable of $\tilde\varsigma$ is that of $\tilde\varpi$ on
      $\Omega_1(\mathfrak{B}^t)$ and at the raw points of
      $(\Omega_0\setminus\Omega_1)(\mathfrak{B}^t)$ held at the date of $\mathfrak{B}^t$ by a live
      piece.
    \item Off $Y'$, the source variable of $\tilde\varsigma$ is $0$ at the raw points held at
      that date by no live piece (the \emph{cut points}). The source variable is free in the
      members of $\mathfrak{K}_{109}$, and the ambient clause is existential.
    \end{itemize}
    With this definition, $\tilde\varsigma\in\mathfrak{G}_1(\mathfrak{B}^t)$.

    Moreover, on $(\Omega_0\setminus\Omega_1)(\mathfrak{B}^t)$ the following hold.
    \begin{itemize}
    \item Off $Y'$, $\tilde\varsigma$ obeys every index-$0$ $U$-member of
      $\mathfrak{G}(\mathfrak{B}^t)$. Its field is that of the input, and
      $\mathfrak{G}(\mathfrak{B}')\subset\mathfrak{G}^{-J_0}(\mathfrak{B}^t)$ by (e3).
    \item Off $Y'$, $\tilde\varsigma$ also obeys the member for the source variable of
      $\mathfrak{G}(\mathfrak{B}^t)$. At held raw points $T^{\mathfrak{B}'}=T^{\mathfrak{B}^t}$;
      if the point enters $\Omega_1(\mathfrak{B}')$, (e3.i) applies. At cut points
      $0<T^{\mathfrak{B}^t}$.
    \item On $Y'$, $\tilde\varsigma$ obeys the index-$0$ cube gauges $(3.35)_0$ of the
      group-valued part wherever the map of $\mathfrak{O}$ is a layer into the group-valued part
      $U'$ of the pair. This is the case for $\mathrm{O}1$, $\mathrm{O}3$, $\mathrm{O}4$ and
      $\mathrm{O}4'$. At the completion site it holds at the points served by the
      layer representation of the presented pair, provided that representation carries the
      shell of index $0$ among its layers.
    \end{itemize}
    Nothing else is asserted on $Y'\cap(\Omega_0\setminus\Omega_1)(\mathfrak{B}^t)$ $(\ast)$. In
    particular, the following are not derived: the members $(3.37)_0$ of $A'$ there; the member
    for the source variable there; and, at the completion site, every index-$0$ member at the
    raw points of $\mathsf{Z}$ and of the core, where the pair is re-presented.
  \end{enumerate}
\end{enumerate}
\end{lemma}

\begin{lemma}[The weight and size-parameter bounds]\label{4.27:prf-weight-bounds}
Clauses (a)--(d) of Lemma~\ref{4.27:lem-weights-heredity} hold (the proof of that lemma is
incomplete at one step, which lies in its clause (e); clauses (a)--(d) are proved in full below).
We use the notation of that lemma. In particular $w_{*}(j)=\mathsf{L}_0^{2N_0(j)}$, where $N_0(j)$
is the number $N_0$ of density $j$, and $w_{*}(j)=1$ if step $j$ has no site of the operation
$\mathbf{R}$. We also use $\bar\alpha_j=\max_{m\le j}\mathsf{A}(\mathsf{t}_m)$ and
$k_0=k-N_0(k)$. For an arrest $\mathfrak{a}$ of density $k_a$ we put $k_0^a=k_a-N_0(k_a)$. As in
clause (d) of that lemma, $w^{\vee}(x)=\max\{w_{*}(k),w(x)\}$, $\mathfrak{N}^{\vee}$ is the ambient
space of pairs $\mathfrak{N}$ with the clauses $Q_a<5w^{\vee}(x)r_a(\ell(x))$ at the points of
$Z^{\mathrm{on}}$, and $\mathfrak{K}_P=\mathfrak{K}_{109}(\mathfrak{B};\gamma_0^{(k)})$ at density $k$.
\end{lemma}

\begin{proof}
\emph{(a).} The level and the weight of a point $x\in W_{\mathfrak{a}}$ are those of the structure
of the stage of index $n$ of the attempt of step $k_a$ at which $\mathfrak{a}$ was arrested.
That stage works at the level $j=h+n\le k_a$, to which it raises the old shells; here $N_{k_a}$ is
the number of stages of that attempt and $h=k_a-N_{k_a}$. That structure is frozen with
the arrest. By \cite[(1.65), (1.67), (1.69)]{B-CMP122a} the structure of that stage
consists of three regions.
\begin{itemize}
\item[(i)] The new shells. On the new shell of level $m\le j$ the level is $m$ and the weight is
$\mathsf{L}_0^{2\max\{0,m-k_0^a\}}$. These shells lie in $\Omega^{c}_{k_0^a+1}$, so the level of
their points before the attempt is at most $k_0^a$.
\item[(ii)] The old shells. On the old shell of level $m$, with $k_0^a<m\le j$, the level is $j$ and
the weight is $\mathsf{L}_0^{2(j-m)}$.
\item[(iii)] The points of level above $j$. These keep their level and have weight $1$.
\end{itemize}
We bound the weight region by region.
\begin{itemize}
\item In region (i) the exponent is $2\max\{0,m-k_0^a\}\le2(k_a-k_0^a)=2N_0(k_a)$.
\item In region (ii) the exponent is $2(j-m)<2(j-k_0^a)\le2N_0(k_a)$.
\item In region (iii) the weight is $1$.
\end{itemize}
Hence $w_{\mathfrak{a}}(x)\le\mathsf{L}_0^{2N_0(k_a)}=w_{*}(k_a)$.

Next we show that the weight is $1$ outside the band $k_0^a<\ell\le k_a$. A weight different from
$1$ occurs only in two cases: in (i) with $\ell=m>k_0^a$, or in (ii), where $\ell=j\ge m>k_0^a$. In
both cases $\ell\le j\le k_a$.

Finally we treat raised and unraised points. Let $\mu$ be the level of $x$ before the attempt, and
let $g=\ell-\mu$.
\begin{itemize}
\item In region (i) we have $\mu\le k_0^a$. Moreover $m\ge\mu$: on the component $X$ of the attempt
one has $\Omega_\mu\cap X\subset\Omega''_\mu$ (this inclusion, valid at every level, is read off
the definition of the intermediate domains $\Omega''_m$ in \cite[p.~191]{B-CMP122a}, and is
established with the closed form of the stage structure in the proof of clause (e) of
Lemma~\ref{4.27:lem-weights-heredity}, whose proof is incomplete at one step, a step other than
this inclusion), so a point of level $\mu$ lies in no new shell of
level below $\mu$. Hence $g\ge m-k_0^a$ and $g\ge0$. Since
$\mathsf{L}_0\ge2$, the weight is at most $\mathsf{L}_0^{2g}$. If $g=0$, then
$m=\mu\le k_0^a$ and the weight is $1$.
\item In region (ii) we have $g=j-m$, and the weight is $\mathsf{L}_0^{2g}$; it equals $1$ when
$g=0$.
\item In region (iii) we have $g=0$ and the weight is $1$.
\end{itemize}

\emph{(b).} Since $\mathsf{L}_0^2\le L/3<L$, we have $w_{*}(j)=\mathsf{L}_0^{2N_0(j)}\le L^{N_0(j)}$.
By the depth bound for the number $N_0$, we have $L^{N_0(j)}\le L^2R_j$. For comparison, the chain of
inequalities in \cite{B-CMP122a} gives
\begin{equation*}
\mathsf{L}_0^{2N_0(j)}\le3^{-N_0(j)}L(L+1)(N_0(j)+1)^{\beta_0}R_j .
\end{equation*}
Since $3^{-N_0(j)}(N_0(j)+1)^{\beta_0}\le1$, this yields $w_{*}(j)\le L(L+1)R_j$.

By \cite[(2.5)]{B-CMP119}, $R_j$ is the least power of $L$ that is at least $\mathsf{t}_j^{\,r}$;
hence $R_j<L\mathsf{t}_j^{\,r}$, a bound also stated in the declaration on the stage counts.
Therefore $L^2R_j<L^3\mathsf{t}_j^{\,r}\le L^2(L+1)\mathsf{t}_j^{\,r}$.

For $n\le j$, the monotonicity line of the declaration on the stage counts gives
\begin{equation*}
L^2R_j\le(L^2R_n)^2\le\mathsf{l}(\mathsf{x}_n).
\end{equation*}
Hence $w_{*}(j)\le\mathsf{l}(\mathsf{x}_n)$.

\emph{(c).} Recall from Notation~\ref{4.27:not-factors-radii-coupling} the maximal radius function
$\mathsf{A}(\mathsf{t})=C\mathsf{t}^{\,q}e^{-\mathsf{t}/2}$ and $c_\beta=2\ln(1+\beta_0)$, so that
$e^{c_\beta/2}=1+\beta_0$.

The functions $\mathsf{A}$ and $\Psi$ are used below only on the ranges $\mathsf{t}>2q$ and
$\mathsf{t}>2(r+q)$ respectively, on which, as shown in (2) and (4), they decrease. The ceiling on
$\gamma$ among the hypotheses of the arrest theorem, \eqref{4.27:eq-arrest-hypotheses-a}, together
with the inequality $\mathsf{t}_m\ge\mathsf{t}_\gamma$ for every level $m$, places every argument
$\mathsf{t}_m$, $\mathsf{t}_j-c_\beta$, $\mathsf{t}_k-c_\beta$ at which they are evaluated below in
these ranges.

(1) Let $m\le j$. Then
\begin{equation}\label{4.27:eq-weight-bounds-1}
\mathsf{t}_m\ge\mathsf{t}_j-c_\beta ,
\end{equation}
uniformly in $j-m$. This is the second inequality of \cite[(2.6)]{B-CMP119}.

The same inequality also follows from the lower bound $\mathsf{x}_m\ge\mathsf{x}_j-c_5$ on the
coupling flow, under the condition $c_5\gamma^2\le1-(1+\beta_0)^{-2}$. Indeed,
$\mathsf{x}_j\ge\gamma^{-2}$, which is $\mathsf{t}_j\ge\mathsf{t}_\gamma$, gives
\begin{equation*}
\frac{\mathsf{x}_m}{\mathsf{x}_j}\ge1-c_5\gamma^2\ge(1+\beta_0)^{-2},
\end{equation*}
and taking logarithms gives \eqref{4.27:eq-weight-bounds-1}. No monotonicity of
$j\mapsto\mathsf{t}_j$ is used.

(2) We have $(\log\mathsf{A})'(\mathsf{t})=-\tfrac12+q/\mathsf{t}<0$ for $\mathsf{t}>2q$, so
$\mathsf{A}$ decreases there. For $m\le j$, \eqref{4.27:eq-weight-bounds-1} therefore gives
\begin{align*}
\mathsf{A}(\mathsf{t}_m)\le\mathsf{A}(\mathsf{t}_j-c_\beta)
&=e^{c_\beta/2}Ce^{-\mathsf{t}_j/2}(\mathsf{t}_j-c_\beta)^{q}\\
&\le(1+\beta_0)Ce^{-\mathsf{t}_j/2}\mathsf{t}_j^{\,q}=(1+\beta_0)\mathsf{A}(\mathsf{t}_j).
\end{align*}
The last inequality uses $0<\mathsf{t}_j-c_\beta\le\mathsf{t}_j$ and $q\ge0$. Taking the maximum over
$m\le j$ gives $\bar\alpha_j\le(1+\beta_0)\mathsf{A}(\mathsf{t}_j)$.

(3) Combine (2) with the bound $w_{*}(j)<L^2(L+1)\mathsf{t}_j^{\,r}$ of (b). Then
\begin{equation}\label{4.27:eq-weight-bounds-2}
w_{*}(j)\bar\alpha_j\le(1+\beta_0)L^2(L+1)C\mathsf{t}_j^{\,r+q}e^{-\mathsf{t}_j/2}
=(1+\beta_0)\Psi(\mathsf{t}_j).
\end{equation}

(4) We have $(\log\Psi)'(\mathsf{t})=-\tfrac12+(r+q)/\mathsf{t}<0$ for $\mathsf{t}>2(r+q)$, so $\Psi$
decreases there. Let $j\le k$. By \eqref{4.27:eq-weight-bounds-1} with $(j,k)$ in place of $(m,j)$,
we have $\mathsf{t}_j\ge\mathsf{t}_k-c_\beta$. Hence
\begin{equation*}
\Psi(\mathsf{t}_j)\le\Psi(\mathsf{t}_k-c_\beta)
=L^2(L+1)Ce^{c_\beta/2}e^{-\mathsf{t}_k/2}(\mathsf{t}_k-c_\beta)^{r+q}
\le(1+\beta_0)\Psi(\mathsf{t}_k).
\end{equation*}
With \eqref{4.27:eq-weight-bounds-2} this gives
\begin{equation*}
\gamma_0^{(k)}=5\max_{j\le k}w_{*}(j)\bar\alpha_j\le5(1+\beta_0)^2\Psi(\mathsf{t}_k)
=\bar\gamma_0(\mathsf{t}_k).
\end{equation*}
By its formula,
\begin{equation*}
\bar\gamma_0(\mathsf{t})=5(1+\beta_0)^2L^2(L+1)C\mathsf{t}^{\,r+q}e^{-\mathsf{t}/2}
\end{equation*}
depends on $\mathsf{t}$ and on $L$, $C$, $r$, $q$, $\beta_0$ only, all fixed before $\gamma$, and on
no datum of the step.
Since $\Psi$ decreases on $\mathsf{t}>2(r+q)$, which contains $\mathsf{t}\ge\mathsf{t}_\gamma$, the
function $\bar\gamma_0$ decreases on $\mathsf{t}\ge\mathsf{t}_\gamma$. It tends to $0$ because
$\mathsf{t}^{\,r+q}e^{-\mathsf{t}/2}\to0$.

\emph{(d).} The factors multiplying the thresholds in the clauses of the spaces of the arrest theorem
are at most $1+\beta<5$. At the points of live pieces lying in the off-class set $\mathfrak{T}_k$,
the units of $\mathfrak{N}$ are the frozen ones, by the freezing rule for live pieces (at a point of
a live piece the level, the weight and the cube family are those frozen at its arrest). The actual
weight $w(x)$ is at most $w^{\vee}(x)$. Hence every such space lies in $\mathfrak{N}^{\vee}$.

By the comparison of the clauses of spaces of the type of \cite{B-CMP109} with weighted radii, each
clause of $\mathfrak{K}_{109}(\mathfrak{B};\gamma)$ is an absolute threshold on the product of a
weight and a radius $\alpha_{i,\ell}$. Here $\ell=\ell(x)$ is the level of the point $x$ itself. It
therefore suffices to bound $5w^{\vee}(x)\alpha_{i,\ell(x)}$ by $\gamma_0^{(k)}$. There are two
cases.
\begin{itemize}
\item If $w(x)\le w_{*}(k)$, then $w^{\vee}(x)=w_{*}(k)$; moreover $\ell(x)\le k$, and
$\alpha_{i,\ell}\le\mathsf{A}(\mathsf{t}_\ell)\le\bar\alpha_k$. Hence
$5w_{*}(k)\alpha_{i,\ell}\le5w_{*}(k)\bar\alpha_k$.
\item If $w(x)>w_{*}(k)$, then $w^{\vee}(x)=w(x)$, and the structure at $x$ is that of the
attempt of some step $k_a\le k$. The
region-by-region bound in the proof of (a) applies to the stage structure of any attempt and gives
$w(x)\le w_{*}(k_a)$, and, since $w(x)>1$, also $\ell\le k_a\le k$. Hence the product is at most
$5w_{*}(k_a)\bar\alpha_{k_a}$.
\end{itemize}
Both bounds are at most $\gamma_0^{(k)}$, by its definition as a maximum over $j\le k$. This holds
whatever the densities $k_a\le k$ of the arrests present and whatever the values of the couplings.
Therefore $\mathfrak{N}^{\vee}\subset\mathfrak{K}_P$.

The size $a_\flat$ of the stage presentation is built from the entries with indices $5,1,6,7$, taken
in the units of the point. It is at most $\sup_x(1+\beta)w^{\vee}(x)\alpha_{i,\ell(x)}$, the
supremum over the points $x$ of the domain and over the radii $\alpha_{i,\ell(x)}$ entering these
four entries. Since $1+\beta<5$, the two cases above give $a_\flat\le\gamma_0^{(k)}$.

The change-of-variables theorem for the operation $\mathbf{T}$, whose hypothesis on the pair is
membership in a space $\mathfrak{K}_{109}(\mathfrak{B};\gamma_0)$, is stated with $\gamma_0$ a free
parameter, subject to its ceiling condition on the radius.
Since $\gamma_0$ is free there, that theorem applies on $\mathfrak{K}_P$ with
$\gamma_0=\gamma_0^{(k)}$. Indeed, its ceiling condition, the condition $K_\Delta\gamma_0\le\tfrac12$
on the constant $K_\Delta$ introduced below, and the smallness hypothesis of the bump estimate on the
held set $W_{\mathfrak{a}}$ are upper bounds on $\gamma_0$, imposed at $\bar\gamma_0(\mathsf{t}_k)$;
by (c), $\gamma_0^{(k)}\le\bar\gamma_0(\mathsf{t}_k)$, so all three hold at $\gamma_0^{(k)}$. This
is the last assertion of clause (d).

We now list where $\gamma_0$ and $w_{*}$ enter the estimates of the stage presentation to which this
clause is applied. Every occurrence of $\gamma_0$ or $w_{*}$ in those estimates is a membership or a
smallness condition; two of them call for comment.
\begin{itemize}
\item The size $a_\flat$ feeds the bump estimate on the held set $W_{\mathfrak{a}}$, whose
hypothesis is the smallness condition on $a_\flat$.
\item The difference bound of the stage presentation, whose constant we denote $K_\Delta$, hands the
product $K_\Delta\gamma_0$ to two estimates of the stored terms.
\end{itemize}
In those two estimates the product meets a threshold that tends to zero and a count that tends to
infinity as $\mathsf{t}_k$ grows, at the current density $k$. There, any single number $\gamma_0$,
for instance a supremum over all densities, would act as a lower bound on the coupling, and hence as
an upper bound on $K$. For this reason clause (c) bounds $\gamma_0^{(k)}$ by a function of
$\mathsf{t}_k$, and the rule stated after \eqref{4.27:eq-weights-heredity-a} applies: memberships
are read at $\gamma_0^{(k)}$, while constants, floors and ceilings are read at
$\bar\gamma_0(\mathsf{t})$, under one supremum over $\mathsf{t}\ge\mathsf{t}_\gamma$ with the other
functions of $\mathsf{t}$ at the same $\mathsf{t}$. In
particular, at the radius $\gamma_0^{(k)}$ the bump estimate on $W_{\mathfrak{a}}$ receives its
smallness hypothesis from the argument just given.
\end{proof}

\begin{lemma}[The time order of determining sets]\label{4.27:lem-determining-set-order}
Order the determining sets of the run by the dates of the operations that write them, as in
clause (e) of Lemma~\ref{4.27:lem-weights-heredity}, whose proof is incomplete at one step. Each
set $\mathfrak{B}$ is one structure on the whole domain $\Omega_0$: a level $\ell_{\mathfrak{B}}(y)$, a
weight $w_{\mathfrak{B}}(y)$ and a cube family $\mathcal{Q}_{\mathfrak{B}}(y)$ at every point $y$, frozen
on the pieces live at its date. We write $\mathfrak{B}\preceq\mathfrak{B}'$ when the date of
$\mathfrak{B}$ precedes or equals that of $\mathfrak{B}'$. For a step $k$ that carries an
$\mathbf{R}$-site we use the following notation:
\begin{itemize}
\item $N=N_k$ is its number of stages and $h=k-N$;
\item $k_0=k-N_0(k)$, where $N_0(k)$ is the integer fixed in \cite[\S1]{B-CMP122a};
\item $j=h+n$ at stage $n$.
\end{itemize}
For a point $y$ put
\begin{equation*}
\mu(y):=\max\{m\le k:\ y\in\Omega_m\},\qquad \mu''(y):=\max\{m\le k:\ y\in\Omega''_m\}.
\end{equation*}
For an arrest $\mathfrak{a}$ of the branch of Definition~\ref{4.27:def-arrests-pieces} with
$n^{+}=n+1$, that is, whose exit is taken during stage $n+1$, the \emph{exit-born terms} of
$\mathfrak{a}$ are the terms of $\mathfrak{F}(\mathfrak{a})$ born at that exit. The symbol
$\mathrm{L}$ in the date $(k;\mathrm{L})$ labels the completion site; it is not the blocking factor.
\begin{enumerate}
\item[(a)] The dates of density $k$, that is, the dates belonging to step $k$, are the following:
\begin{itemize}
\item $(k;\top)$, the $\mathbf{T}$-step from $k-1$ to $k$;
\item $(k;l)$ for $1\le l\le N_k$, stage $l$, present if the step has an $\mathbf{R}$-site and
stage $l$ is performed on at least one component;
\item $(k;\mathrm{L})$, the completion site, present if at least one component of the class
completes.
\end{itemize}
We put $(k;0):=(k;\top)$. The dates are ordered, and the stage structure on a class component $X$
running through stage $n$ is given in closed form at every $y\in X\cap\Omega_1$, by
\begin{equation}\label{4.27:eq-determining-set-order-a}
\begin{gathered}
(k;\top)\prec(k;1)\prec\dots\prec(k;N_k)\prec(k;\mathrm{L})\prec(k+1;\top),\\[2pt]
\lambda_n(y)=\min\{\mu''(y),\max\{\mu(y),j\}\},\qquad
w_n(y)=\mathsf{L}_0^{2(\lambda_n(y)-\max\{\mu(y),k_0\})^{+}} .
\end{gathered}
\end{equation}
The first line of \eqref{4.27:eq-determining-set-order-a} is a finite total order on the dates;
the second gives the level $\lambda_n(y)$ and the weight $w_n(y)$ of the stage structure at $y$.
\item[(b)] The sets of two consecutive dates differ at each $y\in\Omega_0$ by one of the following
elementary moves, or not at all; the same holds at each cube of $\mathcal{Q}_{\mathfrak{B}}$,
through the cube moves that accompany the nesting lemma for lattice cubes in the proof of clause
(e3) of Lemma~\ref{4.27:lem-weights-heredity}, whose proof is incomplete at one step.
\begin{itemize}
\item[(t1)] A $\mathbf{T}$-step from $k$ to $k+1$ re-levels $\Omega_{k+1}$ alone, from level $k$ to
level $k+1$, with weight $1$ before and after.
\item[(t2)] A stage $n\to n+1$ at density $k$ changes the structure on the running class
components according to \eqref{4.27:eq-determining-set-order-a}. The level rises by $0$ or $1$,
the weight is multiplied by at most $\mathsf{L}_0^{2}$, and older live pieces do not move.
\item[(t3)] An arrest writes no value: its component keeps the stage set
$\mathbb{B}^{(n^{+})}_{k_a}$.
\item[(t4)] A completion, with the dissolution it entails, puts level $k$, which is at least the
level $\ell$ previously in force there, and weight $1$ on the completed components.
\end{itemize}
In each move the new level $\ell'$ and weight $w'$ satisfy $\ell'\ge\ell$ and
$w'\le\mathsf{L}_0^{2(\ell'-\ell)}w$. Consequently, for any two sets $\mathfrak{B}\preceq\mathfrak{B}'$
and every $y$,
\begin{equation*}
\ell_{\mathfrak{B}'}(y)\ge\ell_{\mathfrak{B}}(y),\qquad
w_{\mathfrak{B}'}(y)\le\mathsf{L}_0^{2(\ell_{\mathfrak{B}'}(y)-\ell_{\mathfrak{B}}(y))}\,w_{\mathfrak{B}}(y),
\end{equation*}
which is clause (e1) of Lemma~\ref{4.27:lem-weights-heredity}, whose proof is incomplete at one
step.
\item[(c)] Let $\mathfrak{O}$ be one of the later operations $O1,\dots,O5$, $O4'$ of
\eqref{4.27:eq-operations-zones-a}. Let $\mathfrak{B}^t$ be the set in force immediately before
$\mathfrak{O}$'s map and $\mathfrak{B}'$ the one immediately after it. Then
$\mathfrak{B}^t\preceq\mathfrak{B}'$ by one move, as follows.
\begin{itemize}
\item For $O1$, $O4$, $O4'$ the two sets are the class-wide stage sets $l-1$ and $l$ of the run.
The move is (t2) on its running components and the identity on live pieces and off the class.
\item For $O3$ at $k'$ the two sets are $\mathbb{B}_{k'}$ and $\mathbb{B}_{k'+1}$, and the move is
(t1) on $\Omega_{k'+1}\subset\Lambda_{k'}$.
\item For $O2$ and $O5$ the two sets are $\mathbb{B}''_k$, frozen on every piece then live
(dissolving ones included, since dissolution takes effect after the step), and its image under
\cite[(1.33)]{B-CMP122b}. The move is (t4) on the completed components, including the
integration domain $\Lambda$ of \cite[(1.73)]{B-CMP122a} and the dissolved pieces.
\end{itemize}
\item[(d)] A term is composed at $\mathfrak{O}$ only if it was stored before $\mathfrak{O}$: by the
enumeration of the operations of Ba{\l}aban's step that compose stored terms, these operations are
exactly $O1,\dots,O5$, $O4'$, each taking place at a date of the order in (a), and ``before''
refers to that order, by the definition of a stored term. The space in which a stored term $\tau$
is estimated reads the set $\mathfrak{B}_\tau$ of its birth date, by the reading of stored terms in
Definition~\ref{4.27:def-arrested-spaces} and the convention of clause (e) of
Lemma~\ref{4.27:lem-weights-heredity}, whose proof is incomplete at one step; and
$\mathfrak{B}_\tau\preceq\mathfrak{B}^t$ by composition of the moves between the intermediate
dates. In particular:
\begin{itemize}
\item for $O1$, stage $l$ acts on terms born at stages $i\le n^{+}\le l-1$; for the exit taken
during stage $n+1$ one has $n^{+}=n+1$, and the first later sibling stage is $l=n+2$;
\item for $O2$ at $k=k_a$, the intermediate dates are the stages $i,\dots,N$ performed on the
completing sibling.
\end{itemize}
\end{enumerate}
\end{lemma}

\begin{proof}
\emph{Step 0: the elementary moves.} It suffices to prove the inequalities of (b) for one move,
because they compose: the exponents add. Indeed, suppose
\begin{equation*}
\ell''\ge\ell'\ge\ell,\qquad w''\le\mathsf{L}_0^{2(\ell''-\ell')}w',\qquad
w'\le\mathsf{L}_0^{2(\ell'-\ell)}w .
\end{equation*}
Then
\begin{equation*}
w''\le\mathsf{L}_0^{2(\ell''-\ell')}\,\mathsf{L}_0^{2(\ell'-\ell)}\,w=\mathsf{L}_0^{2(\ell''-\ell)}\,w .
\end{equation*}

(t1) A $\mathbf{T}$-step from $k$ to $k+1$ re-levels $\Omega_{k+1}$ alone, from $k$ to $k+1$. By
\cite[(2.1), (2.3)]{B-CMP119} we have $\Omega_{k+1}\subset\Lambda_k=Z_k^{c}$. Hence
$\Omega_{k+1}$ holds no large-field region, so neither a piece nor the structure of an attempt,
and the weight there is $1$ before and after the move. In any case $w'=1\le\mathsf{L}_0^{2}w$.

(t2) Consider the stage $n\to n+1$ at density $k$, with $j=h+n$. Recall the stage spaces of
\cite[pp.~190--191]{B-CMP122a} on a class component $X$. They consist of three families of
regions.
\begin{itemize}
\item Region (i)$_m$ is $(\Omega''_m\setminus\Omega''_{m+1})\cap X$ for $m=1,\dots,j-1$ and, for
$j>k_0$, $(\Omega''_j\setminus\Omega_{k_0+1})\cap X$ for $m=j$ (for $j\le k_0$ see the third fact
in Step~3 below). It carries level $m$ and weight $\mathsf{L}_0^{2\max\{0,m-k_0\}}$.
\item Region (ii)$_m$ is $(\Omega_m\setminus\Omega_{m+1})\cap X$ for $k_0<m\le j$. It carries level
$j$ and weight $\mathsf{L}_0^{2(j-m)}$.
\item Region (iii)$_m$ is $(\Omega_m\setminus\Omega_{m+1})\cap X$ for $j<m<k$, together with
$\Omega_k\cap X$ for $m=k$. It carries level $m$ and weight $1$.
\end{itemize}
For $j\le k_0$ there are no powers of $\mathsf{L}_0$, and the regions (ii)$_m$ are empty.

By this description, the weight at a point $y$ whose level in $\mathbb{B}_k$ is $\mu(y)$ is a
function of its current level $\lambda$ alone:
\begin{equation*}
w=\mathsf{L}_0^{2(\lambda-\max\{\mu(y),k_0\})^{+}} .
\end{equation*}
We check this region by region.
\begin{itemize}
\item Every region (i)$_m$ lies in $\Omega^{c}_{k_0+1}$: for $m<j$ because
$(\Omega''_{m+1})^{c}\subset(\Omega''_k)^{c}\subset\Omega^{c}_{k_0+1}$, the sets $\Omega''$ being
decreasing and $\Omega''_k=\Omega^{\sim5}_{k_0+1}\supseteq\Omega_{k_0+1}$ \cite[\S1]{B-CMP122a},
and for $m=j$ by the definition of (i)$_j$. Hence $\mu\le k_0$, and the weight is
$\mathsf{L}_0^{2\max\{0,\lambda-k_0\}}$.
\item On (ii)$_m$ one has $\mu=m>k_0$ and $\lambda=j$, so the weight is $\mathsf{L}_0^{2(j-m)}$.
\item On the regions (iii)$_m$ one has $\lambda=\mu$, so the weight is $1$.
\item For $j\le k_0$ no powers occur.
\end{itemize}

The sets $\Omega''_m$ and $\Omega_m$ do not depend on the stage \cite[\S1]{B-CMP122a}. Hence
$\lambda$ moves by $0$ or by $1$. To see this, decompose
\begin{equation*}
\Omega''_j\setminus\Omega_{k_0+1}
=(\Omega''_j\setminus\Omega''_{j+1})\cup(\Omega''_{j+1}\setminus\Omega_{k_0+1}).
\end{equation*}
The points of the first set keep their level; those of the second go from $j$ to $j+1$. The
shells of the regions (ii)$_m$ go from $j$ to $j+1$. The set $\Omega_{j+1}\setminus\Omega_{j+2}$
enters (ii)$_{j+1}$ at its own level with weight $1$. So the exponent grows by at most $2$ per
level. Older live pieces do not move: by Lemma~\ref{4.27:lem-arrest-timing} (the clause of
\eqref{4.27:eq-arrest-timing-a} on integration regions), no later attempt re-levels a point of
their components.

(t3) An arrest leaves its component $X_{\mathfrak{a}}$ with the stage set
$\mathbb{B}^{(n^{+})}_{k_a}$. Components leaving the class at an exit are, in the words of
\cite[\S1]{B-CMP122a}, ``not changed by the subsequent operations''. For the exit taken during
stage $n+1$, the passage from $n$ to $n+1$ on $X_{\mathfrak{a}}$ is itself a move (t2).

(t4) At a completion, and at the dissolution it entails, every point of the completed components
receives level $k\ge\lambda$, $\lambda$ its previous level, and weight $1$. This follows from
\cite[(1.33)]{B-CMP122b} and the dissolution rule of Definition~\ref{4.27:def-arrested-spaces}.
With $g=k-\lambda\ge0$ and $w\ge1$,
\begin{equation*}
w'=1\le\mathsf{L}_0^{2g}w .
\end{equation*}

No other operation writes a structure. By the enumeration of the operations of Ba{\l}aban's step
that compose stored terms, the determining sets are written only by the following:
\begin{itemize}
\item the $\mathbf{T}$-step, \cite[(2.2), (3.5)]{B-CMP119};
\item the choice of the regions and the stages, \cite[pp.~190--191]{B-CMP122a};
\item the exits of \cite[\S1]{B-CMP122a};
\item the completion \cite[(1.33)]{B-CMP122b} with the dissolution rule.
\end{itemize}
By the same enumeration, the operations that compose stored terms are exactly those listed in
\eqref{4.27:eq-operations-zones-a}, and each of them takes place at a date of the order established
in Step~1.

\emph{Step 1: dates.} By the scheduling clause of the hypothesis under which the arrested reading
is made (at most one $\mathbf{R}$-site, with one class, per step), a density $k$ has at most one
$\mathbf{R}$-site, with one class. It
therefore has one $N=N_k$, $h=k-N$ and $k_0=k-N_0(k)$. Moreover $h<k_0$, since $N>N_0$ is assumed
in \cite[\S1]{B-CMP122a}. The stage clock is class-wide: \cite[(1.63)]{B-CMP122a} carries one
$j=h+n$ for the whole class. Each component exits at most once per step, by the exits of
\cite[\S1]{B-CMP122a} and clause (T1) of \eqref{4.27:eq-arrest-timing-a} in
Lemma~\ref{4.27:lem-arrest-timing}.

The dates of density $k$ are therefore the following.
\begin{itemize}
\item $(k;\top)$: the $\mathbf{T}$-step from $k-1$ to $k$, writing $\mathbb{B}_k$, by
\cite[(2.2), (3.5)]{B-CMP119}.
\item $(k;l)$ for $1\le l\le N_k$: stage $l$, present if the step has an $\mathbf{R}$-site and
stage $l$ is performed on at least one component.
\item $(k;\mathrm{L})$: present if at least one component of the class completes, through
\cite[(1.65), (1.68), (1.70)]{B-CMP122b} followed by \cite[(1.33)]{B-CMP122b}.
\end{itemize}
They are ordered as in the first line of \eqref{4.27:eq-determining-set-order-a}, which is a finite
total order, and $(k;0):=(k;\top)$.

An arrest with frozen index $n^{+}$ is attached to the date $(k_a;n^{+})$. In particular, an exit
taken before any stage is attached to $(k_a;\top)$, and an exit taken after the last stage to
$(k_a;N_{k_a})$. By (t3) an arrest writes no value.

\emph{The set of a date.} Let $d$ be a date of density $k$ and $y\in\Omega_0$. The set of $d$
at $y$ is given by the first applicable clause below.
\begin{itemize}
\item[(F)] Suppose $y$ lies in a piece of an arrest attached to a date $\preceq d$ and not
dissolved at a date $\preceq d$. Then the set is the frozen $(\ell,w,\mathcal{Q})$ of the
innermost such piece, by the structure freeze of Definition~\ref{4.27:def-arrested-spaces}. This
is well defined: the live pieces holding $y$ are nested, by (T1) and (T4) of
\eqref{4.27:eq-arrest-timing-a} in Lemma~\ref{4.27:lem-arrest-timing}, and pieces of non-nested
arrests are disjoint, by the last assertion of that lemma.
\item[(S)] Otherwise, suppose $d=(k;l)$ and $y$ lies in a class component running through stage
$l$. Then the set is the stage-$l$ structure of \cite[pp.~190--191]{B-CMP122a}, with the factor
$c$ there taken equal to $1$.
\item[(B)] Otherwise the set is that of $\mathbb{B}_k$: level $\mu(y)$, with $\Omega_0$ the whole
domain, weight $1$ and the standard cubes. From $(k;\mathrm{L})$ on, $\mathbb{B}_k$ means the
re-levelled set of \cite[(1.33)]{B-CMP122b}.
\end{itemize}

A stored term reads the set of its birth date, by the reading of stored terms in
Definition~\ref{4.27:def-arrested-spaces} and the convention of clause (e) of
Lemma~\ref{4.27:lem-weights-heredity}, whose proof is incomplete at one step:
\begin{itemize}
\item $\mathbb{C}^{(i)}_k$ reads $(k;i)$;
\item an exit-born term of $\mathfrak{a}$ reads $(k_a;n+1)$;
\item a term born at a $\mathbf{T}$-step reads $(k;\top)$;
\item a term born at the completion site reads $(k;\mathrm{L})$.
\end{itemize}
Live pieces of other classes inside $\mathfrak{T}_k$ fall under (F). They meet no class
component, since pieces and components are nested or disjoint. Second-class components, and
components failing one of the two conditions defining the class in \cite[\S1]{B-CMP122a}, fall
under (B).

\emph{Step 2: one operation per consecutive pair.} By Step~0, each pair of consecutive dates is
separated by exactly one operation that writes a structure.
\begin{itemize}
\item From $(k;l)$ to $(k;l+1)$, $0\le l<N$, the operation is stage $l+1$. This is move (t2); the
arrests attached there make move (t3).
\item From $(k;N)$ to $(k;\mathrm{L})$ the operation is \cite[(1.33)]{B-CMP122b}. This is move
(t4); the arrests attached to $(k;N)$ write nothing.
\item From the last date of $k$ to $(k+1;\top)$ the operation is the $\mathbf{T}$-step. This is
move (t1).
\end{itemize}

\emph{Step 3: the stage structure in closed form.} We use three facts from
\cite[\S1, pp.~190--191]{B-CMP122a}; each is stated below together with the passages of that paper
from which it follows.

\emph{First fact.} On a class component $X$,
\begin{equation*}
\Omega''_m\cap X=\Omega_m\cap X\qquad\text{for } m\le h .
\end{equation*}
Indeed \cite[p.~191]{B-CMP122a} states: ``For $n=j-h=0$ the above space coincides with the space
$\tilde U^c_k(X,\tilde\alpha_0,\tilde\alpha_1)$''. At $n=0$ the regions
$(\Omega''_m\setminus\Omega''_{m+1})\cap X$ for $m<h$, and $(\Omega''_h\setminus\Omega_{h+1})\cap X$,
must therefore be the strata of $\mathbb{B}_k$. Two strata at distinct levels cannot be matched,
because their radii differ by a positive power of $L$. The sets $\Omega''$ are fixed at the site
before the stages begin \cite[\S1]{B-CMP122a}. The identity is also stated in
\cite[\S1]{B-CMP122a}.

\emph{Second fact.} The completion \cite[(1.33)]{B-CMP122b} acts on the completed components only.
Indeed \cite[\S1]{B-CMP122a} states that ``all the characteristic functions introduced above
depend on the field variables local\-ized in the corresponding components''. The stage spaces read
$\Omega''_m$ only through $\Omega''_m\cap X$. Components that exit are ``not changed by the
subsequent operations'' \cite[\S1]{B-CMP122a}. Finally, no later attempt re-levels a point of an
arrested component (Lemma~\ref{4.27:lem-arrest-timing}).

\emph{Third fact.} For $j\le k_0$, region (i)$_{m=j}$ is $(\Omega''_j\setminus\Omega_{j+1})\cap X$,
with cubes $\square\subset\Omega''_j$, $\square\cap\Omega^{c}_{j+1}\neq\emptyset$. The regions
(ii)$_m$ are void, and no power of $\mathsf{L}_0$ occurs. This is \cite[p.~191]{B-CMP122a}: ``For
$j\leqq k_0$ it is simpler\dots''. The definition for $j>k_0$, taken at $j\le k_0$, would overlap
the regions (iii)$_m$ on $\Omega_{j+1}\setminus\Omega_{k_0+1}$.

\emph{Comparison of $\mu$ and $\mu''$.} We have $\mu''\ge\mu$, since
$\Omega_m\cap X\subset\Omega''_m$ for every $m$:
\begin{itemize}
\item for $m\le h$ by the first fact;
\item for $m=h+1$ because $\Omega''_{h+1}=\Omega^{\sim5}_{h+1}\supseteq\Omega_{h+1}$
\cite[\S1]{B-CMP122a};
\item for $h<m\le k_0$ because $\Omega''_m=\Omega^{\sim5}_m\supseteq\Omega_m$
\cite[\S1]{B-CMP122a};
\item for $m>k_0$ because, by \cite[\S1]{B-CMP122a},
\begin{equation*}
\Omega''_m\supseteq\Omega''_k=\Omega^{\sim5}_{k_0+1}\supseteq\Omega_{k_0+1}\supseteq\Omega_m .
\end{equation*}
\end{itemize}
By the last inclusion, moreover, $\mu\ge k_0+1$ implies $\mu''=k$.

\emph{The closed form.} At stage $n$, with $j=h+n$, the level and the weight are given by the
second line of \eqref{4.27:eq-determining-set-order-a}. We verify this case by case.

Let $j>k_0$.
\begin{itemize}
\item $y\in$ (i)$_m$ with $m<j$ if and only if $\mu''=m<j$, and then $\mu\le k_0$. The structure
is $(m,\mathsf{L}_0^{2(m-k_0)^{+}})$. The formula gives the same: $\max\{\mu,j\}\ge j>m$, so
$\lambda_n=m$, and $\max\{\mu,k_0\}=k_0$.
\item $y\in$ (i)$_j$ if and only if $\mu''\ge j$ and $\mu\le k_0$. The structure is
$(j,\mathsf{L}_0^{2(j-k_0)})$, and the formula gives $\lambda_n=\min\{\mu'',j\}=j$ with the same
weight.
\item $y\in$ (ii)$_m$ if and only if $\mu=m\in\,]k_0,j]$. The structure is
$(j,\mathsf{L}_0^{2(j-m)})$. The formula agrees: $\mu''=k\ge j$ gives $\lambda_n=j$, and the
exponent is $2(j-m)$.
\item $y\in$ (iii)$_m$ if and only if $\mu=m>j$. The structure is $(m,1)$, and the formula gives
$\lambda_n=\min\{k,m\}=m$ with exponent $0$.
\end{itemize}
These four cases form a partition of $X\cap\Omega_1$.

Let $j\le k_0$, and use the third fact.
\begin{itemize}
\item $y$ lies in a region (iii)$_m$ if and only if $\mu\ge j+1$. The structure is $(\mu,1)$, and
$\lambda_n=\min\{\mu'',\mu\}=\mu$ because $\mu''\ge\mu$.
\item $y\in$ (i)$_j$ if and only if $\mu\le j\le\mu''$. The structure is $(j,1)$, and
$\lambda_n=j$.
\item $y\in$ (i)$_m$ with $m<j$ if and only if $\mu''=m<j$. The structure is $(m,1)$, and
$\lambda_n=m$.
\end{itemize}
In all three cases the exponent $2(\lambda_n-\max\{\mu,k_0\})^{+}$ vanishes, because
$\lambda_n\le\max\{\mu,j\}\le\max\{\mu,k_0\}$.

At $n=0$, where $j=h<k_0$ and so every weight is $1$, the structure is that of $\mathbb{B}_k$
exactly when the first fact holds.

\emph{The movers.} From the closed form, $\lambda_{n+1}-\lambda_n\in\{0,1\}$. The difference
$\max\{\mu,j+1\}-\max\{\mu,j\}$ equals $1$ exactly when $\mu\le j$, and the minimum with $\mu''$
then rises exactly when $\mu''\ge j+1$. So the difference equals $1$ exactly on the set of movers
\begin{equation*}
M_j:=(\Omega''_{j+1}\setminus\Omega_{j+1})\cap X=\{y\in X:\ \mu(y)\le j,\ \mu''(y)\ge j+1\}.
\end{equation*}
A mover has $\lambda_n=j\ge\mu$. The exponent $e(\lambda):=2(\lambda-\max\{\mu,k_0\})^{+}$ rises,
when $\lambda$ rises by one, by $2$ times the indicator of $\lambda\ge\max\{\mu,k_0\}$. Hence on
$M_j$:
\begin{itemize}
\item if $j\ge k_0$, the move is $(\lambda,w)=(j,w)\mapsto(j+1,\mathsf{L}_0^{2}w)$;
\item if $j<k_0$, the move is $(j,1)\mapsto(j+1,1)$.
\end{itemize}
Every other point is fixed.

The set $Z_h\cap X$ never moves. Indeed, let $y\in Z_h$, and let $Z'_h\supseteq Z_h$ be the region
of \cite[\S1]{B-CMP122a} through which $\Omega''_{h+1}$ is defined. Then
$y\notin\Omega''_{h+1}=(Z'^{\sim3}_h)^{c}$ \cite[\S1]{B-CMP122a}, so $\mu''\le h$. The first fact
then gives $\mu''=\mu$, and $\lambda_n\equiv\mu$. Every older live piece inside $X$ lies in this
set: by (T2) of \eqref{4.27:eq-arrest-timing-a} its density $k_c$ satisfies $k_c\le h$, so
$X_{\mathfrak{c}}\subset Z_{k_c}\subset Z_h$. This is the table of move (t2) stated in Step~0.

\emph{Step 4: the moves (t1) and (t4) pointwise.}

(t1) By \cite[(2.2)]{B-CMP119} the $\mathbf{T}$-step keeps $\Omega_1,\dots,\Omega_k$ and adds
$\Omega_{k+1}\subset\Lambda_k$ \cite[(3.5)]{B-CMP119}. Every live piece, and every class component
of a step $\le k$, lies in $Z_k=\Lambda_k^{c}$. At the last date of $k$ no component is running.
Hence on $\Omega_{k+1}$ clause (B) holds at both dates, and the move is $(k,1)\mapsto(k+1,1)$. Off
$\Omega_{k+1}$ nothing moves. The regions created by the step lie in $\Omega_{k+1}$, are born at
level $k+1$ with weight $1$, and hold no piece.

(t4) Let $Z_{\mathrm{L}}$ denote the union of the completed components, and write
$\Omega_m(\mathbb{B}''_k)$ for the set of points at which the set $\mathbb{B}''_k$ in force at
$(k;N)$ has level at least $m$. Let $X$ be a class component running through stage $N$. A point
of $\Omega_k\cap X$ lies in no piece, since every piece satisfies
$W_{\mathfrak{c}}\subset\Omega^{c}_{k_c}\subset\Omega^{c}_k$; so clause (S) applies there, and at
stage $N$, where $j=k$, the closed form gives $\lambda_N=\min\{\mu'',\max\{\mu,k\}\}=k$, because
$\mu''\ge\mu=k$. Hence $\Omega_k\cap X\subset\Omega_k(\mathbb{B}''_k)$. Off these components
$\Omega_k\subset\Omega_k(\mathbb{B}''_k)$: a point of
$\Omega_k$ there lies in no piece, since every piece lies in
$\Omega^{c}_{k_c}\subset\Omega^{c}_k$, and it is read by (B) at level $k$. By
\cite[(1.33)]{B-CMP122b}, which adds the completed components to the top region of the set in
force and leaves the lower regions unchanged off them, and by the second fact, the new regions are
\begin{align*}
\Omega^{\mathrm{new}}_m&=(\Omega_m(\mathbb{B}''_k)\cap Z_{\mathrm{L}}^{c})\cup Z_{\mathrm{L}}
=\Omega_m(\mathbb{B}''_k)\cup Z_{\mathrm{L}}\qquad(m\le k),\\
\Omega^{\mathrm{new}}_k&=\Omega_k(\mathbb{B}''_k)\cup Z_{\mathrm{L}}\supseteq\Omega_k\cup Z_{\mathrm{L}},
\end{align*}
the last inclusion by what has just been said.
On $Z_{\mathrm{L}}$, every structure in force at $(k;N)$ has $\lambda\le k$: levels of
$\mathbb{B}_k$ are $\le k$, stage levels are $\le j\le k$, and frozen levels are $\le k_c<k$. It
also has $w\ge1$, since weights are of the form $\mathsf{L}_0^{2e}$ with $e\ge0$. This structure is
replaced by $(k,1)$. The frozen structures of the dissolved pieces are dropped after the step, by
the dissolution rule of Definition~\ref{4.27:def-arrested-spaces}.

Pieces meeting $Z_{\mathrm{L}}$ lie in it and dissolve there. Indeed, by (T4) of
\eqref{4.27:eq-arrest-timing-a} and by \cite[\S1]{B-CMP122a},
\begin{equation*}
W_{\mathfrak{c}}\subset Z_h\subset Z''_k\subset\Lambda\subset Z_{\mathrm{L}},
\end{equation*}
where $\Lambda$ is the integration domain of \cite[(1.73)]{B-CMP122a}. Off $Z_{\mathrm{L}}$ nothing
moves.

\emph{Conclusion of Steps 0--4.} Consecutive dates differ pointwise by one of (t1)--(t4), or not
at all, with $\ell'\ge\ell$ and $w'\le\mathsf{L}_0^{2(\ell'-\ell)}w$. By composition the same holds
for all $d\preceq d'$: from $\ell''\ge\ell'\ge\ell$ we get
\begin{equation*}
w''\le\mathsf{L}_0^{2(\ell''-\ell')}w'\le\mathsf{L}_0^{2(\ell''-\ell)}w .
\end{equation*}
Cubewise, consider the matchings $\square\mapsto\square'$ given by the cube moves that accompany the
nesting lemma for lattice cubes in the proof of clause (e3) of
Lemma~\ref{4.27:lem-weights-heredity}, whose proof is incomplete at one step. There is one
matching per consecutive pair, for the same four moves, and they compose: from
$\square\subset\square'\subset\square''$ with level gains $g$ and $g'$,
\begin{equation*}
w''\le\mathsf{L}_0^{2g'}w'\le\mathsf{L}_0^{2(g+g')}w .
\end{equation*}

\emph{Step 5: per operation.}
\begin{itemize}
\item For $O1$ (a stage $l>n^{+}$ of the siblings at $k_a$),
$(\mathfrak{B}^t,\mathfrak{B}')=(\mathfrak{B}_{(k_a;l-1)},\mathfrak{B}_{(k_a;l)})$. The move is (t2)
on the components running through $l$ and the identity on $W_{\mathfrak{a}}$ and on older pieces.
\item For $O4$ and $O4'$ (stage $l$ at $k_b>k_a$),
$(\mathfrak{B}^t,\mathfrak{B}')=(\mathfrak{B}_{(k_b;l-1)},\mathfrak{B}_{(k_b;l)})$. The move is again
(t2) on the components running through $l$; on $W_{\mathfrak{a}}$ and on every other live piece the
set is read frozen by (F), by the structure freeze of Definition~\ref{4.27:def-arrested-spaces},
and is not re-levelled, by Lemma~\ref{4.27:lem-arrest-timing}; off the class nothing moves.
\item For $O3$ at $k'$, $\mathfrak{B}^t$ is the set of the last date of $k'$ and $\mathfrak{B}'$ that
of $(k'+1;\top)$. The move is (t1).
\item For $O2$ and $O5$ at $k$,
$(\mathfrak{B}^t,\mathfrak{B}')=(\mathfrak{B}_{(k;N)},\mathfrak{B}_{(k;\mathrm{L})})$, and the move is
(t4). Here $\mathfrak{B}_{(k;N)}=\mathbb{B}''_k$, read by (F) on every piece then live, dissolving
ones included. During the map at the completion site, the spaces of the target terms still read
the frozen structures, while the $\mathbf{R}$-source is written over the re-levelled
$\mathbb{B}_k=\mathfrak{B}'$ of \cite[(1.33)]{B-CMP122b}. The same re-levelled $\mathbb{B}_k$ is the
set $\mathfrak{B}^t$ of $O3$ from $k$ to $k+1$. These are two consecutive dates.
\end{itemize}

\emph{Births.} A term $\tau=\mathbb{C}^{(i)}_{k_a}$ has birth date $d_\tau=(k_a;i)$, with
$1\le i\le n\le n^{+}$. An exit-born term $\tau$ of $\mathfrak{a}$ has birth date
$d_\tau=(k_a;n+1)=(k_a;n^{+})$. In every case $d_\tau$ precedes or equals the date of
$\mathfrak{B}^t$.
\begin{itemize}
\item For $O1$: $i\le n^{+}\le l-1$, by Definitions~\ref{4.27:def-operations-zones}
and~\ref{4.27:def-arrests-pieces}.
\item For the exit taken during stage $n+1$: stage $n+1$ is one class-wide operation. During it
$X_{\mathfrak{a}}$ runs and composes the terms $\mathbb{C}^{(i)}$, $i\le n$, into the run's own
stage-$(n+1)$ expression, before the arrest is in force. Since the exit-born terms of
$\mathfrak{a}$ are born at $(k_a;n+1)$, the first operation acting on them is stage $n+2$, the
completion site, or the $\mathbf{T}$-step. Even if $\mathfrak{O}$ were taken to be stage $n+1$ on
a sibling, we would have $\mathfrak{B}^t=\mathfrak{B}_{(k_a;n)}\succeq\mathfrak{B}_{(k_a;i)}$.
\item For $O2$ at $k=k_a$: $\mathfrak{O}$ is at the same site, at the later date
$(k_a;N)\succeq(k_a;n^{+})\succeq(k_a;i)$. Equality is possible; the chain then has length $0$,
and $\preceq$ is reflexive.
\item For $O2$ at $k>k_a$, and for $O3$, $O4$, $O4'$, $O5$: these operations are at later
densities. They are reached through the remaining dates of density $k_a$ and then, for each
$k''$ with $k_a<k''\le k$, through the date $(k'';\top)$ (move (t1)), the stages of $k''$ (moves
(t2)), and its completion (move (t4)) if step $k''$ has an $\mathbf{R}$-site.
\end{itemize}
Hence $\mathfrak{B}_\tau\preceq\mathfrak{B}^t$ by composition of the moves between consecutive
intermediate dates.
\end{proof}

\begin{lemma}[Heredity of the direct entries]\label{4.27:prf-direct-entries}
This lemma supplies clause (e2) of \eqref{4.27:eq-weights-heredity-b}, and the part of clause (e3)
that concerns the points of $\Omega_1(\mathfrak{B})$ and the cubes of $\mathcal{Q}_{\mathfrak{B}}$.
Let $\mathfrak{B}\preccurlyeq\mathfrak{B}'$ be two determining sets, ordered by the date order
$\preccurlyeq$ of Lemma~\ref{4.27:lem-determining-set-order}. Let $y\in\Omega_1(\mathfrak{B})$, and put
$\lambda:=\ell_{\mathfrak{B}}(y)$ and $g:=\ell_{\mathfrak{B}'}(y)-\ell_{\mathfrak{B}}(y)$, so that $g\ge 0$
by clause (e1) of \eqref{4.27:eq-weights-heredity-b}.
\begin{enumerate}
\item[(i)] For every direct entry $a\in\{1,\dots,7\}$ we have
$\theta^{\mathfrak{B}'}_a(y)\le\theta^{\mathfrak{B}}_a(y)$, and, if $g\ge 1$,
\begin{equation}\label{4.27:eq-direct-entries-1}
\theta^{\mathfrak{B}'}_a(y)\le\tfrac{64}{147}\,\theta^{\mathfrak{B}}_a(y).
\end{equation}
\item[(ii)] For $a\in\{1,3,5,6,7\}$, every pair obeying the clause
$Q_a(y)<5\theta^{\mathfrak{B}'}_a(y)$ of $\mathfrak{G}(\mathfrak{B}')$ at $y$ obeys the clause
$Q_a(y)<5\theta^{\mathfrak{B}}_a(y)$ of $\mathfrak{G}(\mathfrak{B})$ at $y$.
\item[(iii)] Let $\square$ be a cube of $\mathcal{Q}_{\mathfrak{B}}$ of family $(\lambda_\square,w_\square)$
(level, weight, in the sense of \cite[(1.65), (1.67), (1.69)]{B-CMP122a}). Let $g_\square\ge 0$ be an
integer, and suppose that $\square\subset\square'$ for a cube $\square'$ of
$\mathcal{Q}_{\mathfrak{B}'}$ of family $(\lambda_\square+g_\square,w')$ with
$w'\le\mathsf{L}_0^{2g_\square}w_\square$. If a gauge transformation $u$ on
$\square'\cap\Omega_0(\mathfrak{B}')$ witnesses the cube clause of $\mathfrak{G}(\mathfrak{B}')$ on
$\square'$, then its restriction to $\square\cap\Omega_0(\mathfrak{B})$ witnesses the cube clause of
$\mathfrak{G}(\mathfrak{B})$ on $\square$.
\end{enumerate}
Here $B$, $C$ and $M$ denote the constants of the cube clause in the definition of
$\mathfrak{G}(\mathfrak{B})$ in \eqref{4.27:eq-weights-heredity-b}, whose threshold on a cube of family
$(\lambda_\square,w_\square)$ is $5w_\square BCM\alpha_{0,\lambda_\square}$.
\end{lemma}

\begin{proof}
\emph{Item (i).} Let $g=0$. Clause (e1) gives
$w_{\mathfrak{B}'}(y)\le\mathsf{L}_0^{0}w_{\mathfrak{B}}(y)=w_{\mathfrak{B}}(y)$, and the level is the same at
both dates, so $\theta^{\mathfrak{B}'}_a(y)=w_{\mathfrak{B}'}(y)r_a(\lambda)\le\theta^{\mathfrak{B}}_a(y)$.

Let $g\ge 1$. We claim
\begin{equation}\label{4.27:eq-direct-entries-2}
\begin{split}
\theta^{\mathfrak{B}'}_a(y)&=w_{\mathfrak{B}'}(y)\,r_a(\lambda+g)
\le\mathsf{L}_0^{2g}\,w_{\mathfrak{B}}(y)\,r_a(\lambda+g)\\
&\le w_{\mathfrak{B}}(y)\,(L/3)^g\,\chi_g^{-1}L^{-e_ag}\,r_a(\lambda)
\le 3^{-g}(1+\beta_0)^2g^{1/2}\,\theta^{\mathfrak{B}}_a(y)
\le\tfrac{64}{147}\,\theta^{\mathfrak{B}}_a(y).
\end{split}
\end{equation}
The first inequality is clause (e1), $w_{\mathfrak{B}'}(y)\le\mathsf{L}_0^{2g}w_{\mathfrak{B}}(y)$.

For the second we use the condition $\mathsf{L}_0^2\le L/3$ on the weight base, and we apply the radius
ratio \eqref{4.27:eq-factors-radii-coupling-a} with $\mu=\lambda$ and $m=\lambda+g$. It gives
$r_a(\lambda)/r_a(\lambda+g)\ge\chi_gL^{e_ag}$, that is,
$r_a(\lambda+g)\le\chi_g^{-1}L^{-e_ag}r_a(\lambda)$.

For the third, $(L/3)^gL^{-e_ag}=3^{-g}L^{-(e_a-1)g}\le 3^{-g}$, since $e_a\ge 1$; the exponents
$e_a\in\{1,2,3\}$ of the radii of the direct entries are those of
Definition~\ref{4.27:not-factors-radii-coupling}. Moreover
$\chi_g^{-1}=(1+\beta_0)^2g^{1/2}$, and $w_{\mathfrak{B}}(y)r_a(\lambda)=\theta^{\mathfrak{B}}_a(y)$.

For the fourth, the condition $\beta_0\le 1/7$ gives $(1+\beta_0)^2\le 64/49$. The function
$g\mapsto g^{1/2}3^{-g}$ decreases on $g\ge 1$, since the ratio of consecutive values is
$((g+1)/g)^{1/2}/3\le 2^{1/2}/3<1$. Its value at $g=1$ is $1/3$, and
$(64/49)\cdot(1/3)=64/147$.

This proves \eqref{4.27:eq-direct-entries-1}; note $64/147<0.436$. The factor $\chi_g$ comes from
\cite[(2.6)$_1$, (2.7)$_1$]{B-CMP119}, applied to the ratio of the radius at the later, higher level
$\lambda+g$ to the radius at the earlier level $\lambda$. No monotonicity of
$j\mapsto\mathsf{t}_j$ is used.

\emph{Item (ii).} The entries $Q_a(y)$, $a\in\{1,3,5,6,7\}$, are quantities of the pair at $y$ and do
not depend on the determining set. Since $\ell_{\mathfrak{B}'}(y)\ge\ell_{\mathfrak{B}}(y)\ge 1$, the point
$y$ lies in $\Omega_1(\mathfrak{B}')$, where $\mathfrak{G}(\mathfrak{B}')$ asks
$Q_a(y)<5\theta^{\mathfrak{B}'}_a(y)$. By (i) this gives $Q_a(y)<5\theta^{\mathfrak{B}}_a(y)$.

\emph{Item (iii).} The domain $\Omega_0(\cdot)$ never shrinks along $\preccurlyeq$: this is the
monotonicity of the domains $\Omega_0(\cdot)$ and $\Omega_1(\cdot)$ along the date order, proved below
in this section with the treatment of the points of index zero. Hence
$\square\cap\Omega_0(\mathfrak{B})\subseteq\square'\cap\Omega_0(\mathfrak{B}')$, and $u$ restricts to
$\square\cap\Omega_0(\mathfrak{B})$. The relation $U^u=e^{i\eta A}$ holds there pointwise, with the
same $A$. By the clause on $\square'$,
\begin{equation}\label{4.27:eq-direct-entries-3}
\begin{split}
L^{\lambda_\square}\eta|A|&=L^{-g_\square}\,L^{\lambda_\square+g_\square}\eta|A|\\
&<L^{-g_\square}\cdot 5w'BCM\alpha_{0,\lambda_\square+g_\square},
\end{split}
\end{equation}
and likewise
\begin{equation*}
(L^{\lambda_\square}\eta)^2|\nabla^{\eta}A|
<L^{-2g_\square}\cdot 5w'BCM\alpha_{0,\lambda_\square+g_\square}.
\end{equation*}
This is the arithmetic of \eqref{4.27:eq-direct-entries-2} with exponent $e=1$, respectively $e=2$,
in place of $e_a$.

If $g_\square=0$, then $w'\le w_\square$ and both right-hand sides are at most
$5w_\square BCM\alpha_{0,\lambda_\square}$.

If $g_\square\ge 1$, we use $w'\le\mathsf{L}_0^{2g_\square}w_\square\le(L/3)^{g_\square}w_\square$,
together with the bound
\begin{equation*}
\alpha_{0,\lambda_\square+g_\square}\le(1+\beta_0)^2g_\square^{1/2}\alpha_{0,\lambda_\square}
=\chi_{g_\square}^{-1}\alpha_{0,\lambda_\square}
\end{equation*}
from \cite[(2.6)$_1$, (2.7)$_1$]{B-CMP119}, applied to the ratio of the radius at the later level
$\lambda_\square+g_\square$ to the radius at the earlier level $\lambda_\square$, on which
\eqref{4.27:eq-factors-radii-coupling-a} rests. For $e\in\{1,2\}$ this gives
\begin{equation*}
\begin{split}
L^{-eg_\square}\cdot 5w'BCM\alpha_{0,\lambda_\square+g_\square}
&\le(\mathsf{L}_0^2/L^{e})^{g_\square}\chi_{g_\square}^{-1}\cdot 5w_\square BCM\alpha_{0,\lambda_\square}\\
&\le 3^{-g_\square}(1+\beta_0)^2g_\square^{1/2}\cdot 5w_\square BCM\alpha_{0,\lambda_\square}\\
&\le\tfrac{64}{147}\cdot 5w_\square BCM\alpha_{0,\lambda_\square},
\end{split}
\end{equation*}
by the same four facts as in (i). So the restricted gauge obeys the cube clause of
$\mathfrak{G}(\mathfrak{B})$ on $\square$, with threshold $5w_\square BCM\alpha_{0,\lambda_\square}$.

The hypothesis of (iii) is a geometric statement. It says that each cube of $\mathcal{Q}_{\mathfrak{B}}$
of family $(\lambda_\square,w_\square)$ is, or lies in, a cube of $\mathcal{Q}_{\mathfrak{B}'}$ of family
$(\lambda_\square+g_\square,w')$ with $g_\square\ge 0$ and $w'\le\mathsf{L}_0^{2g_\square}w_\square$, so
that the shell is matched; here the shell of a family is the level stratum $\Omega_m\setminus\Omega_{m+1}$
(respectively $\Omega''_m\setminus\Omega''_{m+1}$) that its cubes meet. This matching is proved cube by
cube, for each of the four elementary moves between consecutive dates of
Lemma~\ref{4.27:lem-determining-set-order}, in the nesting lemma below.

The pointwise clause (e1) does not give it. A cube carries its family's weight
\cite[(1.67)]{B-CMP122a}, not the weight of one of its points. Suppose $\square'$ were taken in the
family one shell below that of $\square$. Then $w'/w_\square=\mathsf{L}_0^4$, and the ratio at $e=1$
becomes $\mathsf{L}_0^4L^{-1}\chi_1^{-1}$. By the bounds used above this ratio is at most
$(L/9)(64/49)$, and it could reach that value, which is $>1$.
\end{proof}

\begin{lemma}[The nested-cube lemma]\label{4.27:lem-nested-cubes}
We work in $\mathbb{R}^d$ with the lattice $\eta\mathbb{Z}^d$ of spacing $\eta$. Any $d$ is
allowed; in this book $d=4$. The following conventions are in force.
\begin{enumerate}
\item A site $x\in\eta\mathbb{Z}^d$ is identified with its closed cell $x+[0,\eta]^d$. A region
is a union of cells, hence a closed set. The complement region $A^{\complement}$ of a region $A$
is the union of the cells not contained in $A$. Thus $A\cup A^{\complement}=\mathbb{R}^d$, and
$A\cap A^{\complement}$ lies in cell boundaries. For a lattice cube $\square$, a condition of the
form $\square\cap(\mathbb{R}^d\setminus A)\neq\emptyset$, as it occurs in the definition of the
cube families in \cite[(1.67)]{B-CMP122a}, is read as $\square\not\subset A$, that is, $\square$
contains a cell of $A^{\complement}$.
\item A lattice box is a set $\prod_{l}[p_l,q_l]$ with $p_l<q_l$ in $\eta\mathbb{Z}$. A cube of
side $s\in\eta\mathbb{N}$ is a lattice box with $q_l-p_l=s$ for every $l$. When a cube is said to
be of a given size, its side is exactly that number, not merely comparable to it; likewise a block
of a partition is a cube whose side is exactly the block side of the partition.
\item Distances $\operatorname{dist}$ and diameters $\operatorname{diam}$ between closed sets are
taken in the sup-norm. Two regions share a cell if and only if their interiors meet.
\item For $S\in\eta\mathbb{N}$ and $v\in\eta\mathbb{Z}^d$, the aligned $S$-blocks
$Q_n:=v+nS+[0,S]^d$, $n\in\mathbb{Z}^d$, tile $\mathbb{R}^d$ with disjoint interiors.
\item The components of a region made of aligned blocks are the connected components of that
closed set. In particular, two blocks sharing a boundary point lie in the same component.
\end{enumerate}
Let $A$ be a region and let $\square\subset A$ be a cube of side $s$. Let $s'\in\eta\mathbb{N}$
with $s'\ge s$. Put
\begin{equation*}
\mathcal{N}:=\{z\in\mathbb{R}^d:\ \operatorname{dist}(z,\square)\le s'\},
\end{equation*}
which is a lattice cube of side $s+2s'$. Assume one of the following.
\begin{itemize}
\item[(a)] Within distance $s'$ of $\square$, the set $A^{\complement}$ is one axis-parallel
box, in the following cell form: there is a lattice box $B$, a union of cells of
$A^{\complement}$ ($B=\emptyset$ allowed), such that every cell of $A^{\complement}$ contained
in $\mathcal{N}$ is a cell of $B$.
\item[(b)] $A$ is a union of aligned $S$-blocks with $S\ge s'$.
\end{itemize}
Then some cube $\square'$ of side $s'$ satisfies
\begin{equation}\label{4.27:eq-nested-cubes-a}
\square\subset\square'\subset A .
\end{equation}
\end{lemma}

\begin{proof}
Write $\square=\prod_{l\le d}I_l$ with $I_l:=[a_l,a_l+s]$.

\emph{Case} (a). First we construct a cube $\square'$ of side $s'$ containing $\square$.

If $B=\emptyset$, put $I'_l:=[a_l,a_l+s']$ for every $l$.

Otherwise write $B=\prod_l J_l$ with $J_l=[b_l,b'_l]$ and $b_l<b'_l$. The cells of $\square$ are
cells of $A$, because $\square\subset A$. The cells of $B$ are cells of $A^{\complement}$. Hence
$\square$ and $B$ share no cell, that is,
$\operatorname{int}\square\cap\operatorname{int}B=\emptyset$. Now
\begin{equation*}
\operatorname{int}\square=\prod_l\,]a_l,a_l+s[,\qquad
\operatorname{int}B=\prod_l\,]b_l,b'_l[ ,
\end{equation*}
so their intersection is the product of the intersections of the factors. A product of sets is
empty only if one of its factors is empty. Hence there is a coordinate $i$ with
$]a_i,a_i+s[\,\cap\,]b_i,b'_i[\,=\emptyset$. Both intervals have positive length, so this means
\begin{equation*}
a_i+s\le b_i \qquad\text{or}\qquad b'_i\le a_i .
\end{equation*}
In the first case put $I'_i:=[a_i+s-s',a_i+s]$, and in the second case put $I'_i:=[a_i,a_i+s']$.
For $l\neq i$ put $I'_l:=[a_l,a_l+s']$.

In both cases $\square':=\prod_l I'_l$ is a lattice cube of side $s'$: its endpoints lie in
$\eta\mathbb{Z}$, since $a_l$, $s$ and $s'$ do. Moreover $I_l\subset I'_l$ for every $l$, so
$\square\subset\square'$.

It remains to show that $\square'\subset A$. We proceed in three steps.

(1) Each $I'_l$ contains $I_l$ and has length $s'$. Hence every $z\in\square'$ has a
coordinatewise nearest point $z^*\in\square$ with $|z_l-z^*_l|\le s'-s$ for every $l$. It follows
that
\begin{equation*}
\operatorname{dist}(z,\square)\le s'-s\le s' .
\end{equation*}
Therefore $\square'\subset\mathcal{N}$, and every cell of $\square'$ is contained in
$\mathcal{N}$.

(2) Suppose $B\neq\emptyset$. In the first case $\sup I'_i=a_i+s\le b_i$; in the second case
$\inf I'_i=a_i\ge b'_i$. Either way
$\operatorname{int}I'_i\cap\operatorname{int}J_i=\emptyset$. Hence
$\operatorname{int}\square'\cap\operatorname{int}B=\emptyset$, so $\square'$ and $B$ share no
cell.

(3) Let $\Delta$ be a cell of $\square'$. By (1), $\Delta\subset\mathcal{N}$. Suppose $\Delta$
were a cell of $A^{\complement}$. By hypothesis (a) it would then be a cell of $B$. This
contradicts (2) when $B\neq\emptyset$, and it contradicts $B=\emptyset$ otherwise. Since the
cells tile $\mathbb{R}^d$, $\Delta$ is a cell of $A$.

Hence $\square'\subset A$, which is \eqref{4.27:eq-nested-cubes-a}. By (1), every point of
$\square'$ lies within distance $s'-s$ of $\square$.

The cell form of hypothesis (a) is needed, and plain disjointness of closed sets would not do.
With closed blocks, a cube $\square\subset A$ and a box $B\subset A^{\complement}$ may share a
face. For example, for any $S\in\eta\mathbb{N}$ with $S\ge s$, take
\begin{equation*}
\square=[0,s]^d,\qquad B=[s,s+S]\times[-S,2S]^{d-1}.
\end{equation*}
Then $I_1\cap J_1=\{s\}$, and $I_l\subset J_l$ for $l\neq1$. A configuration of this kind occurs
in the applications, where $\square$ is a cube of one of the cube families attached to the levels
and $B$ is an adjacent block, or the inner box that remains when a box component of a large-field
region is stripped of its collar. For this reason the argument uses ``no common cell'' in
place of ``disjoint''. Only the separation of $\square$ from $B$ uses the cell form; the
construction of $\square'$ does not.

\emph{Case} (b). Let the aligned lattice be $v+S\mathbb{Z}^d$, with $S$-blocks
\begin{equation*}
Q_n:=\prod_l\,[v_l+n_lS,\ v_l+(n_l+1)S],\qquad n\in\mathbb{Z}^d .
\end{equation*}
Write $A=\bigcup\mathcal{F}$ for a family $\mathcal{F}$ of these blocks. Two distinct blocks meet
only in their boundaries. Put
\begin{equation*}
\mathcal{P}:=\{\text{blocks } Q:\ \operatorname{int}Q\cap\square\neq\emptyset\}.
\end{equation*}

\emph{The index sets.} Fix an axis $l$ and consider the open intervals
$]v_l+nS,v_l+(n+1)S[$ that meet $I_l$. We claim their indices are consecutive and there are at
most two of them. Suppose an index $n$ and some index $n''\ge n+2$ both occurred. Then $I_l$
would contain a point $<v_l+(n+1)S$ and a point $>v_l+(n+2)S$. This would give $s>S$,
contradicting $s\le s'\le S$. Also, at least one such interval meets $I_l$, because $s>0$.

For each $n$ we have
\begin{equation*}
\operatorname{int}Q_n\cap\square=\prod_l\bigl(\,]v_l+n_lS,v_l+(n_l+1)S[\,\cap I_l\bigr).
\end{equation*}
This set is non-empty exactly when every factor is non-empty. Let $\mathcal{I}_l$ be the set of
indices $n$ for which $]v_l+nS,v_l+(n+1)S[$ meets $I_l$; it consists of one or two consecutive
integers. Then
\begin{equation*}
\mathcal{P}=\{Q_n:\ n_l\in\mathcal{I}_l\ \text{for every } l\}.
\end{equation*}

\emph{The box $P$.} The union of these blocks is a box:
\begin{equation*}
P:=\bigcup\mathcal{P}=\prod_l P_l,\qquad
P_l:=\bigl[v_l+(\min\mathcal{I}_l)S,\ v_l+(\max\mathcal{I}_l+1)S\bigr].
\end{equation*}
Every side of $P$ equals $S$ or $2S$, hence is at least $S\ge s'$.

\emph{$\square\subset P$.} Let $t\in I_l$. If $t$ lies inside some open interval
$]v_l+nS,v_l+(n+1)S[$, then $n\in\mathcal{I}_l$ and so $t\in P_l$. If instead $t=v_l+nS$ is a
lattice point, then $t$ is an endpoint of an open interval meeting $I_l$, since $s>0$. The
closure of that interval lies in $P_l$ and contains $t$. Hence $I_l\subset P_l$ for every $l$.

\emph{$P\subset A$.} Let $Q\in\mathcal{P}$ and pick $z\in\operatorname{int}Q\cap\square$. Since
$\square\subset A=\bigcup\mathcal{F}$, there is a block $Q''\in\mathcal{F}$ containing $z$, so
$Q''$ meets $\operatorname{int}Q$. A block other than $Q$ meets $Q$ only in $\partial Q$, because
interiors are disjoint. Hence $Q''=Q$, so $Q\in\mathcal{F}$. Thus $P\subset A$.

\emph{The cube $\square'$.} Inside $P$ choose $\square':=\prod_l[c_l,c_l+s']$, where
\begin{equation*}
c_l:=\max\{v_l+(\min\mathcal{I}_l)S,\ a_l+s-s'\}.
\end{equation*}
Since $v_l$, $S$, $a_l$, $s$ and $s'$ lie in $\eta\mathbb{Z}$, so do $c_l$ and $c_l+s'$; thus
$\square'$ is a lattice cube of side $s'$.
We check $I_l\subset[c_l,c_l+s']$. First, $c_l\le a_l$: the left end of $P_l$ is at most $a_l$
because $a_l\in P_l$, and $a_l+s-s'\le a_l$ because $s\le s'$. Second,
$c_l+s'\ge a_l+s$, directly from the definition of $c_l$.

We check $[c_l,c_l+s']\subset P_l$. By definition $c_l$ is at least the left end of $P_l$. For
the right end,
\begin{equation*}
c_l+s'=\max\{v_l+(\min\mathcal{I}_l)S+s',\ a_l+s\}\le v_l+(\max\mathcal{I}_l+1)S .
\end{equation*}
The first entry satisfies this bound because $s'\le S$ and $P_l$ has length at least $S$. The
second satisfies it because $a_l+s\in P_l$.

Thus
\begin{equation*}
\square\subset\square'\subset P\subset A,
\end{equation*}
with $\square'$ of side $s'$, which is \eqref{4.27:eq-nested-cubes-a}. Each factor of $\square'$
has length $s'$ and contains $I_l$. As in step (1) of case (a), every point of $\square'$
therefore lies within distance $s'-s$ of $\square$.
\end{proof}

\begin{lemma}[Cube families under a stage move. Stated; proof incomplete at the step marked $(\ast)$]\label{4.27:prf-stage-move-cubes}
Fix a step $k$ and work in the setting of Definition~\ref{4.27:def-one-class-per-step} (one class per
step) and under the hypotheses of the arrest theorem, Definition~\ref{4.27:def-arrest-hypotheses}.
Write $N_k=\check{R}_k$ for the number of stages (Definition~\ref{4.27:def-one-class-per-step}),
$h=k-N_k$ and $k_0=k-N_0(k)$; since $N_k>N_0(k)$ by the hypotheses of the arrest theorem, $h<k_0$.
Let $\Omega_m$, $\Omega''_m$, $\Lambda_m$ and
$Z_m=\Lambda_m^{\complement}$ be Ba{\l}aban's regions, $R_m$ the numbers of \cite[(2.5)]{B-CMP119},
$\eta$ the lattice spacing, and $C\le 2$ the constant in the cube sides of \cite[p.~191]{B-CMP122a}.
Put
\begin{equation*}
\xi_m:=ML^m\eta,\qquad \mathsf{S}_m:=MR_mL^m\eta=R_m\xi_m.
\end{equation*}
\begin{enumerate}
\item[(a)] Distinct components of $Z_k$ are at sup-norm distance at least $\mathsf{S}_k$ from one
another, and distinct components of $\Omega_m^{\complement}$, $1\le m\le k$, at least
$\mathsf{S}_m$. Moreover $\mathsf{S}_{i+1}\ge\mathsf{S}_i$ for $1\le i<k$,
$\mathsf{S}_{k_0+1}=\xi_k$, and
\begin{equation}\label{4.27:eq-stage-move-cubes-a}
\mathsf{S}_k\ge 8\xi_k,\qquad
\mathsf{S}_m\ge\xi_{k+1}=L\xi_k\quad\text{for } k_0+2\le m\le k .
\end{equation}
\item[(b)] Let $h\le j\le k-1$, $j=h+n$, put $j_\flat:=\min\{j,k_0\}$, let $X$ be a component of
the class that runs through
stage $n+1$, and let $\mathfrak{B}$, $\mathfrak{B}'$ be the determining sets in force at the dates
$(k;n)$ and $(k;n+1)$ of Lemma~\ref{4.27:lem-determining-set-order}. A cube $\square$ carries a
\emph{family} $(\lambda,w_\square)$: its level $\lambda$ and its weight $w_\square$. A family other
than the top one
is defined by two conditions $\square\subset A$ and $\square\not\subset A'$; we call
$A'^{\complement}$ its \emph{outer set}, and a cube of such a family is \emph{at $X$} when the cells
in which it meets its outer set lie in $X$; the top cubes, $\square\subset\Omega_k$, carry no outer
set. With $(\cdot)^{+}$ denoting the positive part, the cube families of $\mathcal{Q}_{\mathfrak{B}}$
at $X$ at stage $j$ are:
\begin{itemize}
\item $(\mathrm{i})_m$, $1\le m\le j-1$: $\square\subset\Omega''_m$,
$\square\not\subset\Omega''_{m+1}$, side $C\xi_m$, family
$(m,\mathsf{L}_0^{2(m-k_0)^{+}})$;
\item $(\mathrm{i})_{m=j}$: $\square\subset\Omega''_j$, $\square\not\subset\Omega_{j_\flat+1}$,
side $C\xi_j$, family $(j,\mathsf{L}_0^{2(j-k_0)^{+}})$;
\item $(\mathrm{ii})_m$, $k_0<m\le j$: $\square\subset\Omega_m$, $\square\not\subset\Omega_{m+1}$,
side $C\xi_j$, family $(j,\mathsf{L}_0^{2(j-m)})$;
\item $(\mathrm{iii})_m$, $j<m\le k$: $\square\subset\Omega_m$ and, for $m<k$,
$\square\not\subset\Omega_{m+1}$, side $C\xi_m$, family $(m,1)$.
\end{itemize}
Then every cube $\square$ of $\mathcal{Q}_{\mathfrak{B}}$ at $X$, of family $(\lambda,w_\square)$,
lies in a cube $\square'$ of $\mathcal{Q}_{\mathfrak{B}'}$ of family $(\lambda+g,w_{\square'})$ with
$g\in\{0,1\}$ and $w_{\square'}\le\mathsf{L}_0^{2g}w_\square$. Namely, $\square'=\square$, $g=0$ and
$w_{\square'}=w_\square$ for the families
$(\mathrm{i})_m$ with $m<j$, $(\mathrm{iii})_m$ with $m>j$, and for the cubes of
$(\mathrm{i})_{m=j}$ not contained in $\Omega''_{j+1}$; $g=1$ and
$w_{\square'}=\mathsf{L}_0^{2(j+1-k_0)^{+}}\le\mathsf{L}_0^{2}w_\square$ for the cubes of
$(\mathrm{i})_{m=j}$ contained in $\Omega''_{j+1}$; and $g=1$, $w_{\square'}=\mathsf{L}_0^{2}w_\square$,
with $\square'$ in the
stage-$(j+1)$ family $(\mathrm{ii})_m$, for the cubes of $(\mathrm{ii})_m$. Cubes of families at
other components, and the frozen families of the older live pieces inside $X$, are not moved.
\end{enumerate}
\end{lemma}

\begin{proof}
Throughout, the conventions of the nested-cube lemma (Lemma~\ref{4.27:lem-nested-cubes},
display \eqref{4.27:eq-nested-cubes-a}) are in force: a site is its closed cell, a region is a union
of cells, $A^{\complement}$ is the union of the cells not in $A$, the statement
$\square\not\subset A$ means that $\square$ holds a cell of $A^{\complement}$, distances and
diameters are taken in the sup-norm between closed sets, and two regions share a cell when their
interiors meet. For a cube $\square$ of side $s$ and a number $s'\ge s$ we write $\mathcal{N}$ for the
closed sup-norm $s'$-neighbourhood of $\square$, a lattice cube of side $s+2s'$; it is the
neighbourhood that enters hypothesis (a) of the nested-cube lemma.

\emph{Separation.} By \cite[(3.5)]{B-CMP119} and the construction of $\Lambda_k$ in
\cite[\S3]{B-CMP119}, the set $Z_k=\Lambda_k^{\complement}$ is a union of aligned closed cubes of side
$LMR_k$ of the lattice of spacing $L^{k-1}\eta$, that is, of aligned blocks of side
$LMR_kL^{k-1}\eta=\mathsf{S}_k$, and its components are the components of this closed union. Two
blocks belonging to distinct components do not meet; since the blocks are aligned, their indices
then differ by at least $2$ in some coordinate, so that the two blocks are at distance at least
$\mathsf{S}_k$ in that coordinate. Hence distinct components of $Z_k$ are at distance at least
$\mathsf{S}_k$. The same argument applies to $\Omega_m^{\complement}$, which by
\cite[(3.5)]{B-CMP119} is a union of aligned blocks of side $\mathsf{S}_m$: its distinct components
are at distance at least $\mathsf{S}_m$.

Next, $R_k\ge L$. Indeed $R_k$ is the least power of $L$ that is $\ge\mathsf{t}_k^{\,r}$
\cite[(2.5)]{B-CMP119}, where $\mathsf{t}_k=\log g_k^{-2}$ and $r$ is the exponent of that
definition. The hypotheses of the arrest theorem contain $\check{R}_k\ge L$ (it is the condition
from which $N_k\ge 6$ is derived there), while Definition~\ref{4.27:def-one-class-per-step} gives
$\check{R}_k<L\mathsf{t}_k^{\,r}$; hence $L<L\mathsf{t}_k^{\,r}$, that is $\mathsf{t}_k^{\,r}>1$, so
$R_k\ne L^0$ and therefore $R_k\ge L$. Equally, $\mathsf{t}_k\ge\mathsf{t}_\gamma$, and
$\mathsf{t}_\gamma>1$, that is $\gamma<e^{-1/2}$, is contained in the condition
$\mathsf{t}_\gamma\ge 2(r+q)+c_\beta$ of the hypotheses of the arrest theorem as soon as
$2(r+q)+c_\beta>1$; no condition beyond those hypotheses is used. Since
$\mathsf{t}_i\ge\mathsf{t}_\gamma>1$ at every level $i\le k$, the same minimality gives $R_i\ge L$
for every $i\le k$. With the floor $L\ge 8$ of the
hypotheses of the arrest theorem,
\begin{equation*}
\mathsf{S}_k=R_k\xi_k\ge L\xi_k\ge 8\xi_k ,
\end{equation*}
which is the first inequality of \eqref{4.27:eq-stage-move-cubes-a}.

\emph{Scales.} First, $R_{i+1}\ge R_i/L$ for $1\le i<k$. Since $R_i\ge L$,
$R_i/L$ is a power of $L$ smaller than $R_i$, so by the minimality in \cite[(2.5)]{B-CMP119} it is
$<\mathsf{t}_i^{\,r}$; the second inequality of \cite[(2.7)]{B-CMP119}, taken at $n-m=1$, and the
definition of $R_{i+1}$ give
\begin{equation}\label{4.27:eq-stage-move-cubes-1}
L^{-1}R_i<\mathsf{t}_i^{\,r}\le(1+\beta_0)\,\mathsf{t}_{i+1}^{\,r}\le(1+\beta_0)R_{i+1}<LR_{i+1},
\end{equation}
so $R_i/R_{i+1}<L^2$; since this ratio is a power of $L$, it is at most $L$, that is
$R_{i+1}\ge R_i/L$. Consequently
\begin{equation*}
\mathsf{S}_{i+1}=R_{i+1}ML^{i+1}\eta\ge (R_i/L)ML^{i+1}\eta=\mathsf{S}_i .
\end{equation*}

By \cite[p.~179]{B-CMP122a}, $N_0$ is the smallest positive integer with
$L^{-N_0+1}MR_{k-N_0+1}=M$. Since $k_0+1=k-N_0+1$, this equation reads
$R_{k_0+1}=L^{N_0-1}$, and then
\begin{equation*}
\mathsf{S}_{k_0+1}=ML^{N_0-1}L^{k-N_0+1}\eta=ML^k\eta=\xi_k .
\end{equation*}
The minimality of $N_0$ gives, for $k_0+2\le m\le k$, the bound $R_m\ge L^{k+1-m}$. To see this,
let $e(R)$ denote the exponent of a power $R$ of $L$, and put
$\nu_n:=e(R_{k-n+1})-(n-1)$ for $n\ge 1$. By $R_{i+1}\ge R_i/L$ we have
$e(R_{k-n})\le e(R_{k-n+1})+1$, so $\nu_{n+1}\le \nu_n$: the integers $\nu_n$ are non-increasing
in $n$.
Further $\nu_1=e(R_k)\ge 1$ because $R_k\ge L$, and $N_0=\min\{n:\ \nu_n=0\}$. Since $\nu$ is
non-increasing and $\nu_{N_0}=0$, every $n<N_0$ has $\nu_n\ge 0$, and $\nu_n\ne 0$ by the minimality
of $N_0$; hence $\nu_n\ge 1$ for $n<N_0$. For $k_0+2\le m\le k$ put $n:=k+1-m$; then $n\le N_0-1$,
and $\nu_n\ge 1$ reads $e(R_m)\ge n=k+1-m$. Therefore
\begin{equation*}
\mathsf{S}_m=MR_mL^m\eta\ge ML^{k+1}\eta=\xi_{k+1}=L\xi_k\qquad(k_0+2\le m\le k),
\end{equation*}
which is the second inequality of \eqref{4.27:eq-stage-move-cubes-a}. This proves (a).

\emph{Cubes at $X$, and cubes at other components.} The families are those of
\cite[pp.~190--191]{B-CMP122a}, with the cube clauses \cite[(1.65), (1.67), (1.69)]{B-CMP122a},
read at the component $X\subset Z$ of (b), with outer sets as in (b). The cells in which a cube at
$X$ meets its outer set lie in $Z_k$ (see the next paragraph), hence, by the separation in (a), in
one component of $Z_k$, so that a cube at $X$ shares a cell with $X$. Cubes at another component are
governed by that component's own families: the structures are chosen component by component, and a
stage set differs from $\mathbb{B}_k$ on the class only, by the convention of
\cite[(1.33)]{B-CMP122b} and because a stage space splits into its class part, defined on the
on-class region $Z^{\mathrm{on}}_k$, and its foreign part, defined on
$\mathfrak{T}_k=\Omega_k\cup Z^{\complement}$. The stages of $X$ do not move these cubes: the sets
$\Omega''_m$, $\Omega_m$ are
global and do not depend on the stage, the stage index is attached to each component separately,
and, by the next paragraph, no family cube at $X$ lies within $4\xi_k$ of another component.

\emph{The enlarged set $X^{+}$.} Put
\begin{equation*}
X^{+}:=\{z:\ \operatorname{dist}(z,X)\le 4\xi_k\}.
\end{equation*}
Since distinct components of $Z_k$ are at distance at least $\mathsf{S}_k\ge 8\xi_k$ by (a),
$Z_k\cap X^{+}=X$. A cube of a family at $X$ has side at most $C\xi_k\le 2\xi_k$ and shares a cell
with $X$, so every point of it is within $2\xi_k$ of $X$. With $s'\le C\xi_k\le 2\xi_k$, every point
of its $s'$-neighbourhood $\mathcal{N}$, and every cube $\square'$ within $s'-s$ of it, is within
$4\xi_k$ of $X$: all of them lie in $X^{+}$, and in particular no family cube at $X$ lies within
$4\xi_k$ of another component. Every complement met below lies in $Z_k$: for $m\le k$, by
\cite[(2.1), (2.3)]{B-CMP119},
\begin{equation*}
\Omega_m^{\complement}\subset\Lambda_m^{\complement}=Z_m\subset Z_k ;
\end{equation*}
and
\begin{equation*}
\Omega''^{\complement}_{j+1}\subset\Omega^{\complement}_{j_\flat+1}\subset Z_k ,
\end{equation*}
because, by \cite[p.~179]{B-CMP122a} and the definition of $Z''_k$ preceding
\cite[(1.73)]{B-CMP122a},
\begin{align*}
\Omega''_{j+1}&=\Omega^{\sim 5}_{j+1}\supseteq\Omega_{j+1} &&(j<k_0),\\
\Omega''_{j+1}&\supseteq\Omega''_k=\Omega^{\sim 5}_{k_0+1}\supseteq\Omega_{k_0+1} &&(j\ge k_0).
\end{align*}
Hence, inside
$\mathcal{N}$, each of these complements coincides with its part in $X$.

\emph{Diameter.} For $s=C\xi_j$ and $s'=C\xi_{j+1}$,
\begin{equation}\label{4.27:eq-stage-move-cubes-2}
\operatorname{diam}\mathcal{N}=s+2s'=(2+L^{-1})\,C\xi_{j+1}\le\bigl(2+\tfrac18\bigr)\cdot
2\xi_{j+1}=4.25\,\xi_{j+1},
\end{equation}
using $C\le 2$ and the floor $L\ge 8$; and two sets that share cells with $\mathcal{N}$ are at
distance at most $\operatorname{diam}\mathcal{N}$ from each other.

\emph{The family $(\mathrm{i})_{m=j}$ on either side of $k_0$.} For $j>k_0$ the family
$(\mathrm{i})_{m=j}$ is that of \cite[(1.65)]{B-CMP122a}, with outer set
$\Omega^{\complement}_{k_0+1}=\Omega^{\complement}_{j_\flat+1}$. For $j\le k_0$,
\cite[p.~191]{B-CMP122a} states (``For $j\le k_0$ it is simpler'') that there are no powers of
$\mathsf{L}_0$ in the family $(\mathrm{i})$ and that the family $(\mathrm{ii})$ is empty; that
$(\mathrm{iii})$ then starts at $\Omega_{j+1}$ is the closed form of the stage structure in
Lemma~\ref{4.27:lem-determining-set-order}, display \eqref{4.27:eq-determining-set-order-a}. So
for $j\le k_0$ the family $(\mathrm{i})_{m=j}$ sits on $\Omega''_j\setminus\Omega_{j+1}$, with
outer set $\Omega^{\complement}_{j+1}=\Omega^{\complement}_{j_\flat+1}$ and weight $1$. On either
side of $k_0$ this is the level-$j$ stratum $\Omega''_j\setminus\Omega_{j+1}$ of
\cite[(1.15)--(1.16)]{B-CMP122a}. For $j\le k_0$ the outer set must be
$\Omega^{\complement}_{j+1}$, since the families $(\mathrm{i})_{m=j}$ and $(\mathrm{iii})$ partition
$\Omega''_j\cap X$ with one level at each point, and the passage from $(\mathrm{iii})_{m=j+1}$ to
$(\mathrm{i})_{m=j+1}$ at $j+1\le k_0$ uses the outer set $\Omega^{\complement}_{j+2}$. With outer
set $\Omega''^{\complement}_{j+1}$ the early family would be a subfamily of the family above, and for
$j\le k_0$ its move would be the identity, so nothing below changes.

\emph{The move from stage $j$ to stage $j+1$ at $X$.} We go through the families of (b).

\emph{The families $(\mathrm{i})_m$, $m<j$.} The family ``$\square\subset\Omega''_m$,
$\square\not\subset\Omega''_{m+1}$, side $C\xi_m$, weight $\mathsf{L}_0^{2(m-k_0)^{+}}$'' is the
same family at stage $j+1$, since $m\le (j+1)-1$. So $\square'=\square$, $g=0$,
$w_{\square'}=w_\square$.

\emph{The families $(\mathrm{iii})_m$, $m>j$.} For $m\ge j+2$ the family is the same at stage
$j+1$. For $m=j+1$ the family ``$\square\subset\Omega_{j+1}$, $\square\not\subset\Omega_{j+2}$,
side $C\xi_{j+1}$, weight $1$'' is, when $j+1>k_0$, the stage-$(j+1)$ family
$(\mathrm{ii})_{m=j+1}$, whose weight is $\mathsf{L}_0^{2((j+1)-(j+1))}=1$; when $j+1=k$ both
families read $\square\subset\Omega_k$, as there is no $\Omega_{k+1}$ at step $k$
\cite[(1.69)]{B-CMP122a}. When $j+1\le k_0$ it is a cube of the stage-$(j+1)$ family
$(\mathrm{i})_{m=j+1}$, of weight $1$: indeed
$\square\subset\Omega_{j+1}\subset\Omega''_{j+1}$ and $\square$ meets
$\Omega^{\complement}_{j+2}$. In all cases $\square'=\square$, $g=0$, $w_{\square'}=w_\square$: the
same cubes, the weights not raised.

\emph{The family $(\mathrm{i})_{m=j}$.} If $\square\not\subset\Omega''_{j+1}$, then $\square$ is a
cube of the stage-$(j+1)$ family $(\mathrm{i})_{m=j}$, namely ``$\square\subset\Omega''_j$,
$\square\not\subset\Omega''_{j+1}$, side $C\xi_j$'', of weight
$\mathsf{L}_0^{2(j-k_0)^{+}}$: so $g=0$ and $w_{\square'}=w_\square$.

Otherwise $\square\subset\Omega''_{j+1}=:A$, and we apply Lemma~\ref{4.27:lem-nested-cubes}(a) with
$s=C\xi_j$ and $s'=C\xi_{j+1}$. We check its hypothesis in cell form: that there is a lattice box
$B$, a union of cells of $A^{\complement}$ (possibly $B=\emptyset$), such that every cell of
$A^{\complement}$ contained in $\mathcal{N}$ is a cell of $B$. Every cell of
$A^{\complement}=\Omega''^{\complement}_{j+1}$ contained in $\mathcal{N}$ lies in
$Z_k\cap X^{+}=X$, hence in $\Omega''^{\complement}_{j+1}\cap X=Z''_{j+1}\cap X$, hence in a
component of $Z''_{j+1}\cap X$ that shares a cell with $\mathcal{N}$. We show that at most one
component does, and that it is a box of cells of $A^{\complement}$.

Put $\mathsf{S}:=\mathsf{S}_{j_\flat+1}$. For $j<k_0$, $\mathsf{S}=R_{j+1}\xi_{j+1}$, and
$R_{j+1}\xi_{j+1}\ge 8\xi_{j+1}$; for $j\ge k_0$, $\mathsf{S}=\mathsf{S}_{k_0+1}=\xi_k$, and
$\xi_k\ge\xi_{j+1}$. So in either case
$\mathsf{S}\ge\xi_{j+1}$. The set $\Omega^{\complement}_{j_\flat+1}$ is a union of aligned
$\mathsf{S}$-blocks \cite[(3.5)]{B-CMP119}, whose components inside $X$ are boxes:
\cite[\S1]{B-CMP122a} states that ``all the regions connected with the last $N$ steps are
rectangular parallelepipeds'' (here $N=N_k$), used here as stated there; it applies to
$\Omega^{\complement}_{j_\flat+1}$ because $j_\flat+1>h$, so that this region belongs to the last
$N_k$ steps. By (a)
these boxes are pairwise at distance at
least $\mathsf{S}$, and $Z''_{j+1}\cap X$ lies in their union, because
$\Omega''_{j+1}\supseteq\Omega_{j_\flat+1}$. Let
$\mathcal{B}=\prod_l[p_l,p_l+n_l\mathsf{S}]$ be one of these boxes; no other component touches
it. The set $\Omega''_{j+1}\cap\mathcal{B}$ is the collar of $\mathcal{B}$: the five layers of
$\mathsf{S}$-blocks of $\Omega^{\sim 5}_{j_\flat+1}$ at $\partial\mathcal{B}$ (for
$h\le j<k_0$, $\Omega''_{j+1}=\Omega^{\sim 5}_{j+1}$ by \cite[p.~179, p.~195]{B-CMP122a}; for
$j\ge k_0$, $\Omega''_{j+1}\supseteq\Omega''_k=\Omega^{\sim 5}_{k_0+1}$
\cite[the definition of $Z''_k$ preceding (1.73)]{B-CMP122a}),
together with, for $j\ge k_0$, the aligned thin rims laid round $\Omega''_k$ layer by layer, in
Ba{\l}aban's words ``separated by one layer of $M$-cubes. Similarly \dots\ one layer of
$L^{-1}M$-cubes, and so on'' \cite[p.~179]{B-CMP122a}, these rims being taken as in the
component-by-component choice of the structures. Therefore $Z''_{j+1}\cap\mathcal{B}$, which is
$\mathcal{B}$ less its
collar, is one box or empty (this is the statement of \cite[\S1]{B-CMP122a}, used here as stated
there, that ``all
components of the domains $Z''_j$ are also rectangular parallelepipeds for $j=h,\dots,k$''), and it
lies at depth at least $5\mathsf{S}$ inside $\mathcal{B}$; the rims only deepen it. Now let
$I_1\subset\mathcal{B}_1$ and $I_2\subset\mathcal{B}_2$ be these cores in two distinct boxes
$\mathcal{B}_1\ne\mathcal{B}_2$, and let $x\in I_1$, $y\in I_2$. The segment $[x,y]$ leaves the
convex set $\mathcal{B}_1$ at a point $z_1$ and enters $\mathcal{B}_2$ at a point $z_2$, in this
order; since a norm is additive along a segment,
\begin{equation}\label{4.27:eq-stage-move-cubes-3}
|x-y|=|x-z_1|+|z_1-z_2|+|z_2-y|\ge 5\mathsf{S}+\mathsf{S}+5\mathsf{S}=11\mathsf{S}.
\end{equation}
By \eqref{4.27:eq-stage-move-cubes-2},
\begin{equation*}
11\mathsf{S}\ge 11\xi_{j+1}>4.25\,\xi_{j+1}\ge\operatorname{diam}\mathcal{N},
\end{equation*}
so $\mathcal{N}$ shares cells with at most one core, and this core is a box $B$ of cells of
$A^{\complement}$ (or there is none, and $B=\emptyset$). The hypothesis of part (a) of the
nested-cube lemma is met; in this case the nesting condition of
Definition~\ref{4.27:def-one-class-per-step} is not used. The older live
pieces inside $X$ lie inside $Z''_{j+1}$: their components satisfy
\begin{equation*}
X_{\mathfrak{c}}\subset Z_{k_{\mathfrak{c}}}\subset Z_h\subset\Omega''^{\complement}_{h+1}
\subset\Omega''^{\complement}_{j+1},
\end{equation*}
by clauses (T2) and (T3) of
Lemma~\ref{4.27:lem-arrest-timing} (display \eqref{4.27:eq-arrest-timing-a}) and by
$\Omega''_{h+1}=\Omega^{\sim 5}_{h+1}=(Z_h'^{\sim 3})^{\complement}$ \cite[p.~195]{B-CMP122a},
where $Z'_h\supseteq Z_h$ is the region surrounded by layers of cubes in the construction of
$\Omega_{h+1}$ in \cite[\S3]{B-CMP119}; in particular $Z_h\subset Z_h'^{\sim 3}$, which gives the
third inclusion above.
No moved cube touches them: the moved cubes lie in $\Omega''_{j+1}$ or in $\Omega_m$ with
$m>k_0>h$, while
$W_{\mathfrak{c}}\subset\Omega^{\complement}_{k_{\mathfrak{c}}}\subset\Omega^{\complement}_h$.

Part (a) of the nested-cube lemma now gives a cube $\square'\supset\square$ of side $C\xi_{j+1}$
with $\square'\subset\Omega''_{j+1}$. The cube $\square'$ still holds the cell of $\square$ in the
outer set $\Omega^{\complement}_{j_\flat+1}$, and
$\Omega^{\complement}_{j_\flat+1}\subset\Omega^{\complement}_{(j+1)_\flat+1}$: the two sets are
equal for $j\ge k_0$, and for $j<k_0$ this is $\Omega_{j+2}\subset\Omega_{j+1}$. So $\square'$ is a
cube of the stage-$(j+1)$ family $(\mathrm{i})_{m=j+1}$, with $g=1$ and
\begin{equation*}
w_{\square'}=\mathsf{L}_0^{2(j+1-k_0)^{+}}\le\mathsf{L}_0^{2}w_\square,
\end{equation*}
with equality $w_{\square'}=\mathsf{L}_0^2w_\square$ for $j\ge k_0$ and $w_{\square'}=1=w_\square$
for $j<k_0$.

\emph{The families $(\mathrm{ii})_m$, $k_0<m\le j$.} The family is ``$\square\subset\Omega_m$,
$\square\not\subset\Omega_{m+1}$, side $C\xi_j$, weight $\mathsf{L}_0^{2(j-m)}$''. The target is
the stage-$(j+1)$ family $(\mathrm{ii})_m$: ``$\square'\subset\Omega_m$,
$\square'\not\subset\Omega_{m+1}$, side $C\xi_{j+1}$, weight
$\mathsf{L}_0^{2(j+1-m)}=\mathsf{L}_0^2w_\square$''. Any $\square'\supset\square$ with
$\square'\subset\Omega_m$ and side $C\xi_{j+1}$ serves, since $\square'$ inherits the cell of
$\square$ in $\Omega^{\complement}_{m+1}$.

We apply Lemma~\ref{4.27:lem-nested-cubes}(a) at $A=\Omega_m$, with $s=C\xi_j$ and $s'=C\xi_{j+1}$.
In cell form: every cell of $\Omega^{\complement}_m$ contained in $\mathcal{N}$ lies in
$Z_m\subset Z_k$ and in $X^{+}$, hence in $\Omega^{\complement}_m\cap X$. Each component of
$\Omega^{\complement}_m\cap X$ is a box, by the statement of \cite[\S1]{B-CMP122a} quoted above
(here $m>k_0\ge h$, as $N_k>N_0(k)$), and these components are pairwise at distance at least
$\mathsf{S}_m$
\cite[(3.5)]{B-CMP119}, with $\mathsf{S}_{k_0+1}=\xi_k$ and $\mathsf{S}_m\ge\xi_{k+1}$ for
$m\ge k_0+2$ by (a). Two components sharing cells with $\mathcal{N}$ are at distance at most
$\operatorname{diam}\mathcal{N}=(2+L^{-1})C\xi_{j+1}$; by \eqref{4.27:eq-stage-move-cubes-2}
this is $<\xi_k$ if $j+1\le k-1$ (because $4.25\,\xi_{j+1}\le 4.25\,\xi_k/L$ and $4.25<L$), and it is
$<\xi_{k+1}$ always. Hence $\mathcal{N}$ meets at most one component and part (a) of the
nested-cube lemma applies, except when $j+1=k$ and $m=k_0+1$, a case that occurs only if
$N_0\ge 2$.

In that exceptional case, for $C=1$, Lemma~\ref{4.27:lem-nested-cubes}(b) at $A=\Omega_{k_0+1}$
replaces part (a): $\Omega_{k_0+1}$ is a union of aligned blocks of side $\xi_k=s'$, where
$s'\ge s$ and $s=\xi_{k-1}$;
the cube $\square$ need not be aligned, and $\square'\not\subset\Omega_{k_0+2}$ because $\square'$
contains the cell of $\square$ lying in $\Omega^{\complement}_{k_0+2}$.
$(\ast)$ For $C=2$ we use the following property, which is the form in which the nesting condition
on the blocking data of Definition~\ref{4.27:def-one-class-per-step}, among the hypotheses of the
arrest theorem, enters the proof, and which is not derived here: if $\mathcal{N}$ is the sup-norm
$C\xi_k$-neighbourhood of a cube of the stage-$(k-1)$ family $(\mathrm{ii})_{k_0+1}$ at $X$, a set
of sup-norm diameter at most $(2+L^{-1})C\xi_k$, then
\begin{equation}\label{4.27:eq-stage-move-cubes-4}
\mathcal{N}\ \text{shares cells with at most one component of}\ \Omega^{\complement}_{k_0+1}\cap X .
\end{equation}
With \eqref{4.27:eq-stage-move-cubes-4}, $\mathcal{N}$ meets at most one box component, and part
(a) of the nested-cube lemma applies as in the non-exceptional case. This is the one place in the
heredity of the cube clause where the nesting condition is used, and that the nesting condition
yields \eqref{4.27:eq-stage-move-cubes-4} is the step at which the proof is incomplete. Property
\eqref{4.27:eq-stage-move-cubes-4} is implied by the statement that
$\Omega^{\complement}_{k_0+1}\cap X$ is one box, and also by the statement that its components are
more than $4.25\,\xi_k$ apart (for instance at least five blocks apart). Without
\eqref{4.27:eq-stage-move-cubes-4}, two box components across a corridor of width $\xi_k$ would
leave $\square$ inside a cube of the family $(\mathrm{i})_{m=k}$ only, with
$w_{\square'}/w_\square=\mathsf{L}_0^4$ at
$g=1$, beyond the bound $w_{\square'}\le\mathsf{L}_0^{2g}w_\square$; and part (b) of the nested-cube
lemma fails there
too, since there $S=\xi_k$ while $s'=2\xi_k$, so $S<s'$.

Equally, one may use Lemma~\ref{4.27:lem-nested-cubes}(b) at $A=\Omega_m$ with $S=\mathsf{S}_m$. The
condition $S\ge s'=C\xi_{j+1}$ holds for $m\ge k_0+2$, because
$\mathsf{S}_m\ge L\xi_k\ge 8\xi_{j+1}$ by \eqref{4.27:eq-stage-move-cubes-a}; and for $m=k_0+1$ it
holds if and only if $j+1\le k-1$ or $C=1$. The exception is the same one, and on this route the
box property of the components of $\Omega^{\complement}_m$ is not needed at all.

In every case we obtain $\square'\supset\square$ with $\square'\subset\Omega_m$, holding the cell of
$\square$ in $\Omega^{\complement}_{m+1}$, of side $C\xi_{j+1}$: a cube of the stage-$(j+1)$
family $(\mathrm{ii})_m$, the same shell, with $g=1$ and
$w_{\square'}=\mathsf{L}_0^{2(j+1-m)}=\mathsf{L}_0^2w_\square$.

\emph{Frozen families.} The frozen families of the older live pieces inside $X$ are untouched, by
clause (T3) of Lemma~\ref{4.27:lem-arrest-timing}.

This exhausts the families of $\mathcal{Q}_{\mathfrak{B}}$ at $X$. In every case the move is the
cubewise form of the movers in the closed form of the stage structure,
Lemma~\ref{4.27:lem-determining-set-order}, display \eqref{4.27:eq-determining-set-order-a}.
\end{proof}

\begin{lemma}[The cube clause under $\mathbf{T}$-steps and completion]
\label{4.27:prf-cube-clause-heredity}
Write $\xi_m=ML^m\eta$ and $\mathsf{S}_m=MR_mL^m\eta=R_m\xi_m$, where $R_m$ is the block parameter of
the large-field regions at level $m$, the least power of $L$ not smaller than $\mathsf{t}_m^{\,r}$.
Let $C$ be the constant in the side
$C\xi_\lambda$ of the cubes of the families of \cite[(1.65), (1.67), (1.69)]{B-CMP122a}, so that
$C\le 2$ by \cite[p.~191]{B-CMP122a}, and let $B$ be the constant multiplying the radii in the cube
conditions of those families \cite[(1.65)]{B-CMP122a}. For a determining set $\mathfrak{B}$, every
cube $\square$ of
its cube family $\mathcal{Q}_{\mathfrak{B}}$ belongs to a family with a level $\lambda$ and a weight
$w\ge 1$; we say that $\square$ is of family $(\lambda,w)$, and its side is $C\xi_\lambda$. On such a
cube the cube clause of the space $\mathfrak{G}(\mathfrak{B})$ of
\eqref{4.27:eq-weights-heredity-b} asks for a gauge transformation $u$ on
$\square\cap\Omega_0(\mathfrak{B})$ with
$U^u=e^{i\eta A}$, where $U$ is the $G$-valued part of $\mathbb{U}$, and
\begin{equation}\label{4.27:eq-cube-clause-heredity-1}
L^{\lambda}\eta\,|A|<5wBCM\alpha_{0,\lambda},\qquad
(L^{\lambda}\eta)^2\,|\nabla^{\eta}A|<5wBCM\alpha_{0,\lambda}\qquad
\text{on }\square\cap\Omega_0(\mathfrak{B}) .
\end{equation}
\begin{enumerate}
\item[(i)] Let $\mathfrak{B}$ and $\mathfrak{B}'$ be the determining sets of two consecutive dates of
the time order \eqref{4.27:eq-determining-set-order-a} which differ by a $\mathbf{T}$-step
$k\to k+1$, by an arrest, or by a completion at density $k$ (the moves (t1), (t3), (t4) of
\eqref{4.27:eq-weights-heredity-b}). Then every cube $\square\in\mathcal{Q}_{\mathfrak{B}}$ of family
$(\lambda,w)$ is contained in a cube $\square'\in\mathcal{Q}_{\mathfrak{B}'}$ of family
$(\lambda+g,w')$ with $g\ge 0$ and $w'\le\mathsf{L}_0^{2g}w$.
\item[(ii)] Let $\mathfrak{B}\preccurlyeq\mathfrak{B}'$, let $\square\in\mathcal{Q}_{\mathfrak{B}}$
be of family $(\lambda,w)$ and $\square'\in\mathcal{Q}_{\mathfrak{B}'}$ of family $(\lambda+g,w')$,
with $\square\subset\square'$, $g\ge0$ and $w'\le\mathsf{L}_0^{2g}w$. Every pair
$(\mathbb{U},\mathbb{J})$ that satisfies the cube clause of $\mathfrak{G}(\mathfrak{B}')$ on
$\square'$ satisfies the cube clause of $\mathfrak{G}(\mathfrak{B})$ on $\square$.
\end{enumerate}
\end{lemma}

The remaining elementary move, a stage move, is treated in Lemma~\ref{4.27:prf-stage-move-cubes},
whose proof is incomplete at one step.

\begin{proof}
Regions are unions of closed cells of the lattice of spacing $\eta$; two regions share a cell when
their interiors meet; $\square\not\subset A$ means that $\square$ contains a cell of the complement
region of $A$. These are the conventions of the nested-cube lemma
(Lemma~\ref{4.27:lem-nested-cubes}), displayed in \eqref{4.27:eq-nested-cubes-a}. Of that lemma we
use part (b): if $\square\subset A$ is a cube of side $s\le s'$ and $A$ is a union of aligned cubes
of side $S\ge s'$, then some cube $\square'$ of side $s'$ satisfies
$\square\subset\square'\subset A$. We also use that $R_i\ge L$ at every level $i$, the fact recorded
with the separation statement \eqref{4.27:eq-stage-move-cubes-a} of
Lemma~\ref{4.27:prf-stage-move-cubes}, whose proof is incomplete at one step. Finally, the weight
base satisfies $2\le\mathsf{L}_0$ \eqref{4.27:eq-weights-heredity-b} and $\mathsf{L}_0^2\le L/3$
\cite[p.~191]{B-CMP122a}; in particular $L\ge 3\mathsf{L}_0^2$, whence $L>2$ and $L/2>4$.

\emph{(i), the $\mathbf{T}$-step $k\to k+1$.} The step keeps $\Omega_1,\dots,\Omega_k$ and adds
$\Omega_{k+1}\subset\Lambda_k=Z_k^{\,c}$ \cite[(2.2), (3.5)]{B-CMP119}, the small-field region at
level $k$; the set
$\Omega_{k+1}$ enters at level $k+1$. It creates no piece and touches no
frozen family: the held set $W_{\mathfrak{a}}$ of a live arrest $\mathfrak{a}$, of density
$k_a\le k$, satisfies $W_{\mathfrak{a}}\subset\Omega^{c}_{k_a}\subset\Omega^{c}_k$, and so shares no
cell with a top cube, whether the top domain in force is $\Omega_k$ or its enlargement after a
completion (the pieces inside the completed components having been dissolved). The families of
levels $m<k$ are defined through
$\Omega_m$ and $\Omega_{m+1}$ (through $\Omega''_m$ and $\Omega''_{m+1}$ for the stage families),
with $m+1\le k$. These sets are unchanged, so each cube of such a family is a cube of the same family
for $\mathfrak{B}'$: $\square'=\square$, $g=0$, $w'=w$.

There remains the top family. It consists of the cubes $\square\subset\Omega_k(\mathfrak{B})$ of side
$C\xi_k$ and weight $1$, and its thresholds carry no reduction factor of the form
$1-\beta\sum 2^{-|\cdot|}$ \cite[(1.69)]{B-CMP122a}. Here
$\Omega_k(\mathfrak{B})$ is the top domain in force at $\mathfrak{B}$: it equals $\Omega_k$, or
$\Omega''_k\cup Z_L\supseteq\Omega_k$ after a completion at density $k$ (see the completion case
below). The argument is the same in both cases.

If $\square\not\subset\Omega_{k+1}$, then $\square$ is a cube of the level-$k$ family of
$\mathbb{B}_{k+1}$, which consists of the cubes $\square\subset\Omega_k(\mathfrak{B})$ with
$\square\not\subset\Omega_{k+1}$ and side $C\xi_k$. It is kept: $\square'=\square$, $g=0$,
$w'=w=1$.

Otherwise $\square\subset\Omega_{k+1}=:A$. By \cite[(3.5)]{B-CMP119} the set $\Omega_{k+1}$, like its
complement, is a union of aligned $LMR_{k+1}$-blocks of the lattice of spacing $L^k\eta$. These
blocks have side $S=\mathsf{S}_{k+1}=R_{k+1}\xi_{k+1}$. We apply part (b) of the nested-cube lemma
with $s=C\xi_k\le s':=C\xi_{k+1}\le S$. The last inequality is $R_{k+1}\ge C$, and it holds because
$R_{k+1}\ge L$, $L>2$ and $C\le 2$. Part (b) gives a cube $\square'\supset\square$ of side
$C\xi_{k+1}$ with
$\square'\subset\Omega_{k+1}$, that is, a top cube of $\mathbb{B}_{k+1}$
\cite[(1.69)]{B-CMP122a}. Then $g=1$ and $w'=1=w\le\mathsf{L}_0^{2}w$.

\emph{(i), an arrest.} An arrest changes no value and no level, weight or cube family. Thus
$\mathcal{Q}_{\mathfrak{B}'}=\mathcal{Q}_{\mathfrak{B}}$ with the same families, and
$\square'=\square$, $g=0$, $w'=w$.

\emph{(i), a completion at density $k$.} This is \cite[(1.33)]{B-CMP122b}, together with the
dissolution of the pieces after the step. We write $Z_L$ for the union of the completed components.
Each of them is a whole component of $Z_k$, hence a union of blocks of side $\mathsf{S}_k$. On
$Z_L$ the completion sets the level to $k$ and the weight to $1$. On the complement of $Z_L$ the set
$\mathbb{B}''_k$ stands, since \cite[(1.33)]{B-CMP122b} acts on the completed components only, as in
the time order of determining sets \eqref{4.27:eq-determining-set-order-a}. Read with this
convention, \cite[(1.33)]{B-CMP122b} gives: for $m<k$ the new $m$-th region is
$\Omega^{\mathrm{new}}_m=\Omega_m\cup Z_L$, and the new top region is
$\Omega^{\mathrm{new}}_k=\Omega''_k\cup Z_L$, which contains $\Omega_k\cup Z_L$.

Let $\square\in\mathcal{Q}_{\mathfrak{B}}$ be of family $(\lambda,w)$, with $\lambda\le k$; the
frozen families of the dissolving pieces are included. Its side is $s=C\xi_\lambda\le s':=C\xi_k$.

Suppose first that $\square$ shares no cell with $Z_L$. The sets defining its family are untouched
off $Z_L$: \cite[(1.33)]{B-CMP122b} keeps $\mathbb{B}''_k$ there, and the families at other
components, the frozen families of pieces elsewhere and the strata off the class are unchanged. A
top cube not meeting $Z_L$ stays a top cube, since
$\Omega_k(\mathfrak{B})\subset\Omega^{\mathrm{new}}_k$. Hence $\square'=\square$, $g=0$, $w'=w$.

Suppose next that $\square$ shares a cell with $Z_L$, hence with some completed component $X_c$. Put
$X_c^{+}=\{z:\operatorname{dist}(z,X_c)\le 4\xi_k\}$. Since the side of $\square$ is
$C\xi_\lambda\le 2\xi_k\le 4\xi_k$, we have $\square\subset X_c^{+}$. By the separation statement
\eqref{4.27:eq-stage-move-cubes-a} of
Lemma~\ref{4.27:prf-stage-move-cubes}, whose proof is incomplete at one step, we have
$Z_k\cap X_c^{+}=X_c$. Hence the cube
$\square$ meets no other block of $Z_k$, and, as $\Lambda_k=Z_k^{\,c}$,
\begin{equation*}
\square\subset X_c\cup\Lambda_k\subset Z_L\cup\Omega_k=:A,
\end{equation*}
because $X_c\subset Z_L$ and $\Lambda_k\subset\Omega_k$. The set $A$ is a union of blocks of one
aligned partition: the blocks of $Z_k$ and of $\Omega^c_k$ are the $LMR_k$-cubes of the lattice of
spacing $L^{k-1}\eta$ of the step $k-1\to k$ \cite[(3.5)]{B-CMP119}. These blocks have side
$S=\mathsf{S}_k=R_k\xi_k$, and $S\ge s'$ because $R_k\ge C$. Part (b) of the nested-cube lemma gives
a cube $\square'\supset\square$ of side $C\xi_k$ with
$\square'\subset\Omega_k\cup Z_L\subset\Omega^{\mathrm{new}}_k$. Since the top family of the new
$\mathbb{B}_k$ is the family of cubes of side $C\xi_k$ contained in the top domain in force,
$\Omega^{\mathrm{new}}_k$ \cite[(1.69), $m=k$]{B-CMP122a}, the cube $\square'$ is a top cube.
Then $g=k-\lambda\ge0$ and $w'=1\le\mathsf{L}_0^{2g}w$, because $w\ge1$.

The second case covers every family at $X_c$. In the spaces of \cite[pp.~190--191]{B-CMP122a} at
$X_c$, where the current stage index is $k$, a cube of the stage shell of level $m$ below the
current stage contains a cell of $\Omega''^{\,c}_{m+1}$, a cube of the shell of level $m$ with
$k_0<m<k$ contains a cell of $\Omega^{c}_{m+1}$, and a cube of the current-stage shell contains a
cell of $\Omega^{c}_{k_0+1}$; each such cell lies in $X_c$. The second case also covers every
frozen family of the pieces nested in $X_c$: these
lie inside $X_c$ and dissolve now.

\emph{The constants in the two parts of the nested-cube lemma.} We have $C\le2$
\cite[p.~191]{B-CMP122a}, and the cubes need not be aligned. Part (a) of the nested-cube lemma
requires neither alignment of $A$ nor any width of the shells of $A$. In particular it admits the
rims, one $M$-cube wide, laid round the regions. Part (b) is used at the ratios
$S/s'=R_{k+1}/C$ and $R_k/C$, both at least $L/2>4$. For $C=1$ it is used in the stage move
at $\Omega_{k_0+1}$ with $S=s'$.

\emph{Composition.} No operation other than the four elementary moves changes the level, weight or
cube family at any point
\eqref{4.27:eq-weights-heredity-b}. Let one matching give $\square\subset\square'$ with gain $g$,
and a following one give $\square'\subset\square''$ with gain $g'$. Then
$\square\subset\square''$ with gain $g+g'$, and
\begin{equation*}
w''\le\mathsf{L}_0^{2g'}w'\le\mathsf{L}_0^{2(g+g')}w .
\end{equation*}
Together with the stage move of Lemma~\ref{4.27:prf-stage-move-cubes}, whose proof is incomplete at
one step, the matchings of (i) therefore compose along the time order
\eqref{4.27:eq-determining-set-order-a} and give the hypothesis of (ii) for every
$\mathfrak{B}\preccurlyeq\mathfrak{B}'$.

\emph{(ii).} Let $(\mathbb{U},\mathbb{J})$ satisfy the cube clause of $\mathfrak{G}(\mathfrak{B}')$
on $\square'$, and let $u$ be the gauge transformation on $\square'\cap\Omega_0(\mathfrak{B}')$
that this clause provides, with $U^u=e^{i\eta A}$. Its bounds are
\eqref{4.27:eq-cube-clause-heredity-1} with $(\lambda+g,w')$, $\square'$ and $\mathfrak{B}'$ in place
of $(\lambda,w)$, $\square$ and $\mathfrak{B}$. Since $\Omega_0(\cdot)$ does not shrink along the
time order (no elementary move lowers a level, and the moves only add points to $\Omega_1(\cdot)$
and to $\Omega_0(\cdot)$, \eqref{4.27:eq-weights-heredity-b}), we have
$\square\cap\Omega_0(\mathfrak{B})\subseteq\square'\cap\Omega_0(\mathfrak{B}')$; we restrict $u$ to
$\square\cap\Omega_0(\mathfrak{B})$, and the same $A$, restricted there, serves.

For the ratio of the radii we use \eqref{4.27:eq-factors-radii-coupling-a}: it gives
$\alpha_{0,\lambda+g}\le\chi_g^{-1}\alpha_{0,\lambda}$, with $\chi_g^{-1}\le(1+\beta_0)^2g^{1/2}$ for
$g\ge1$. We also use the conditions $\mathsf{L}_0^2\le L/3$ \cite[p.~191]{B-CMP122a} and
$\beta_0\le 1/7$, fixed with the radii in \eqref{4.27:eq-factors-radii-coupling-a}, and the fact that
$g^{1/2}3^{-g}$ decreases for $g\ge1$, so that $g^{1/2}3^{-g}\le 1/3$ there.

For $g\ge1$ these give, on $\square\cap\Omega_0(\mathfrak{B})$,
\begin{align*}
L^{\lambda}\eta\,|A|
&=L^{-g}\,L^{\lambda+g}\eta\,|A|
<L^{-g}\cdot 5w'BCM\alpha_{0,\lambda+g}
\le\Bigl(\frac{\mathsf{L}_0^{2}}{L}\Bigr)^{g}\chi_g^{-1}\cdot 5wBCM\alpha_{0,\lambda}\\
&\le 3^{-g}(1+\beta_0)^2g^{1/2}\cdot5wBCM\alpha_{0,\lambda}
\le\frac{64}{147}\cdot 5wBCM\alpha_{0,\lambda},
\end{align*}
and $64/147<0.436$. This is the computation for the direct entries in
Lemma~\ref{4.27:prf-direct-entries}, at exponent one.

For $g=0$ we have $w'\le w$ and the same level, so
\begin{equation*}
L^{\lambda}\eta\,|A|<5w'BCM\alpha_{0,\lambda}\le 5wBCM\alpha_{0,\lambda}.
\end{equation*}

For the gradient we write
$(L^\lambda\eta)^2|\nabla^\eta A|=L^{-2g}(L^{\lambda+g}\eta)^2|\nabla^\eta A|$. The same chain then
holds with the factor $(\mathsf{L}_0^2/L^2)^g\le(\mathsf{L}_0^2/L)^g$ in place of
$(\mathsf{L}_0^2/L)^g$. So both bounds of \eqref{4.27:eq-cube-clause-heredity-1} hold on
$\square\cap\Omega_0(\mathfrak{B})$ for the family $(\lambda,w)$: every pair satisfying the cube
clause of $\mathfrak{G}(\mathfrak{B}')$ on $\square'$ satisfies the cube clause of
$\mathfrak{G}(\mathfrak{B})$ on $\square$.
\end{proof}

\begin{lemma}[Heredity on the raw stratum]\label{4.27:prf-raw-stratum}
Let $\mathfrak{B}\preceq\mathfrak{B}'$ be the determining sets of two dates, of densities $k\le k'$, in the time
order of Lemma~\ref{4.27:lem-determining-set-order}. A point $y\in\Omega_0(\mathfrak{B})\setminus\Omega_1(\mathfrak{B})$ is
called \emph{raw} at $\mathfrak{B}$, and $(\Omega_0\setminus\Omega_1)(\mathfrak{B})$ is the \emph{raw stratum}. By its
definition in \eqref{4.27:eq-weights-heredity-b}, on the raw stratum $\mathfrak{G}(\mathfrak{B})$ carries the index-$0$
clause of $\mathfrak{K}_{109}(\mathfrak{B};\cdot)$ (for $\gamma>0$, $\mathfrak{K}_{109}(\mathfrak{B};\gamma)$ is the
space of pairs $(\mathbb{U},\mathbb{J})$ on $\Omega_0(\mathfrak{B})$ of the type of the small-field spaces of
\cite{B-CMP109}, with $J$-radius $\gamma$; the subscript refers to that paper): its $U$-members, which do not involve
$\gamma$, namely a gauge obeying
the bounds \cite[(3.35)]{B-CMP99b} at $j=0$ on every cube of the index-$0$ family (the aligned family of scale $0$,
independent of the date, with sides at most $CML\eta$, $C\le2$ being the side constant of the cube clause of
$\mathfrak{G}$) and the bounds \cite[(3.37)]{B-CMP99b} at $j=0$ on $A'$ at every raw point; and its $J$-member, read
pointwise, $\ell^{3}|\mathbb{J}(y)|<T^{\mathfrak{B}}(y)$, with $\ell=\eta$ at index $0$ and the dated threshold
$T^{\mathfrak{B}}$ fixed there.
Then the following hold.
\begin{enumerate}
\item[(a)] $T^{\mathfrak{B}}(y)\le\gamma_0^{(k)}$ at every raw point $y$ of $\mathfrak{B}$, and every pair of
$\mathfrak{G}(\mathfrak{B})$ obeys at every raw point the index-$0$ clause of $\mathfrak{K}_{109}(\mathfrak{B};\gamma)$,
for every $\gamma\ge\gamma_0^{(k)}$.
\item[(b)] $\Omega_1(\mathfrak{B})\subseteq\Omega_1(\mathfrak{B}')$ and
$\Omega_0(\mathfrak{B})\subseteq\Omega_0(\mathfrak{B}')$; no elementary move between consecutive dates lowers a level;
and $(\Omega_0\setminus\Omega_1)(\mathfrak{B}')\cap\Omega_0(\mathfrak{B})\subseteq(\Omega_0\setminus\Omega_1)(\mathfrak{B})$.
\item[(c)] At a point raw at both $\mathfrak{B}$ and $\mathfrak{B}'$ the $U$-members of the stratum clauses of
$\mathfrak{G}(\mathfrak{B}')$ and of $\mathfrak{G}(\mathfrak{B})$ are the same clause on the same cubes. Let
$y\in(\Omega_0\setminus\Omega_1)(\mathfrak{B})\cap\Omega_1(\mathfrak{B}')$, and let $\mathfrak{B}_e$, with
$\mathfrak{B}\prec\mathfrak{B}_e\preceq\mathfrak{B}'$, be the set of the date at which $y$ entered $\Omega_1(\cdot)$,
of density $k_e\ge k$. Then every pair of $\mathfrak{G}_1(\mathfrak{B}')$ obeys at $y$ the $U$-members of the stratum
clause of $\mathfrak{G}(\mathfrak{B})$, and
\begin{equation}\label{4.27:eq-raw-stratum-a}
\begin{split}
\ell^{3}|\mathbb{J}(y)|
&<5\alpha_{0,k_e}\,(L^{k_e}\eta)^{-3}\eta^{3}
=5\alpha_{0,k_e}L^{-3k_e}\\
&\le5\bar\alpha_k\le\gamma_0^{(k)} .
\end{split}
\end{equation}
In particular, if $y$ is held at the date of $\mathfrak{B}$ by no live piece, the pair obeys the stratum $J$-member of
$\mathfrak{G}(\mathfrak{B})$ at $y$.
\end{enumerate}
\end{lemma}

\begin{proof}
The bounds \eqref{4.27:eq-weights-heredity-b} on levels and thresholds, the first inclusion stated there, and the
membership of the real point and of every pair obeying the direct clauses in $\mathfrak{G}_1(\mathfrak{B})$
(Lemma~\ref{4.27:lem-weights-heredity}, whose proof is incomplete at one step) concern the points of
$\Omega_1(\cdot)$ and the cubes of $\mathcal{Q}_{\mathfrak{B}}$. What follows supplies the parts of those statements
that concern the raw stratum, on which $\mathfrak{G}(\mathfrak{B})$ carries, by its definition, the index-$0$ clause
recalled in the statement.

\emph{Proof of} (a). Let $y$ be raw at $\mathfrak{B}$. If $y$ lies at the date of $\mathfrak{B}$ in a live piece, then by
the definition of $T^{\mathfrak{B}}$ in \eqref{4.27:eq-weights-heredity-b} we have
$T^{\mathfrak{B}}(y)=\gamma_0^{(k_{\mathfrak{c}})}$, where $W_{\mathfrak{c}}$ is the innermost live piece holding $y$ and
$k_{\mathfrak{c}}$ the density of its arrest; this piece is live at a date of density $k$, so its arrest took place at a
step $k_{\mathfrak{c}}\le k$. By \eqref{4.27:eq-weights-heredity-a}, $\gamma_0^{(\cdot)}$ is a maximum over
$j\le\cdot$ and hence non-decreasing, so $T^{\mathfrak{B}}(y)\le\gamma_0^{(k)}$. If $y$ lies in no live piece,
$T^{\mathfrak{B}}(y)=\gamma_0^{(k)}$ by the same definition. Now let $\gamma\ge\gamma_0^{(k)}$. The $U$-members of the
index-$0$ clause of
$\mathfrak{K}_{109}(\mathfrak{B};\gamma)$ are the $\gamma$-free members recalled in the statement, which
$\mathfrak{G}(\mathfrak{B})$ carries by definition; its $J$-member, $\sup\ell^{3}|\mathbb{J}|\le\gamma$, read pointwise,
is implied by
\begin{equation*}
\ell^{3}|\mathbb{J}(y)|<T^{\mathfrak{B}}(y)\le\gamma_0^{(k)}\le\gamma .
\end{equation*}
Since
$\mathfrak{K}_{109}(\mathfrak{B};\cdot)$ increases with its radius, this gives (a). Together with the second inclusion
at the points of $\Omega_1(\mathfrak{B})$ and on the cubes of $\mathcal{Q}_{\mathfrak{B}}$
(Lemma~\ref{4.27:lem-weights-heredity}, whose proof is incomplete at one
step), it follows that $\mathfrak{G}(\mathfrak{B})\subset\mathfrak{K}_{109}(\mathfrak{B};\gamma)$ as sets on all of
$\Omega_0(\mathfrak{B})$ for every $\gamma\ge\gamma_0^{(k)}$, the stratum entering by definition.

\emph{Proof of} (b). Throughout, $\Omega_m$, $\Lambda_m$ and $Z_m=\Lambda_m^{c}$ are Ba{\l}aban's domains
\cite[(2.1)--(2.3)]{B-CMP119}, $\Gamma_0=\Omega_1^{c}$, and at a step with an $\mathbf{R}$-site at density $k''$ we write
$N_{k''}$ for the number of stages of that step and $h=k''-N_{k''}$, as in
Lemma~\ref{4.27:lem-determining-set-order}, so that stage $n$ of the step acts at level $j=h+n$; we write $\Omega''_m$
for the intermediate domains of \cite{B-CMP122a}, $Z''_j$ for the corresponding intermediate large-field regions of
\cite{B-CMP122a}, $Z'_h$ for the enlarged large-field region of \cite{B-CMP122a}, and a superscript ${}^{\sim n}$ for the
enlargement of a union of blocks by $n$ layers of blocks as in \cite{B-CMP122a}. By
Lemma~\ref{4.27:lem-determining-set-order}, consecutive dates differ
at each point by one of four elementary moves or not at all: a $\mathbf{T}$-step, a stage, an arrest, a completion.
We examine each.

A $\mathbf{T}$-step $k''\to k''+1$ adds $\Omega_{k''+1}$ at level $k''+1$; here
$\Omega_{k''+1}\subset\Lambda_{k''}\subset\Omega_1$.
A stage raises levels by $0$ or $1$ (Lemma~\ref{4.27:lem-determining-set-order}), and only inside
\begin{equation*}
\Omega''_{h+1}=\Omega^{\sim5}_{h+1}=\bigl((Z'_h)^{\sim3}\bigr)^{c}\subset\Lambda_h\subset\Omega_1
\qquad(h\ge1),
\end{equation*}
by \cite{B-CMP122a} and because older live pieces are not moved by the stages; the last inclusion holds since
$\Lambda_h\subseteq\Lambda_1\subseteq\Omega_1$ for $h\ge1$ \cite[(2.1)--(2.3)]{B-CMP119}. That $h\ge1$ is the form in
which the construction of \cite{B-CMP122a} is written: there $Z''_j=Z_j$ for $j=h,\dots,1$, the factor
$\chi_h(Z_h\cap\Omega_h)$ appears, and the domains of \cite[(2.1),(2.2)]{B-CMP119} start from $\Omega_1$. It is also a
consequence of the fact that every component of $Z_{k''}$ holds a large-field region created at some step not later than
$h$, the functions on $\Omega_1^{c}$ being defined in the first step \cite[p.~255]{B-CMP119}. Were $h=0$ admitted, the
stage family at $m=1$, which lies in $\Omega''_1=\Omega_1^{\sim5}$, would raise raw points to level $1\le k_0$, where
$k_0=k''-N_0(k'')$, with $N_0(k'')$ the number defining the maximal weight $w_{*}(k'')$ in
\eqref{4.27:eq-weights-heredity-a}, at weight $\mathsf{L}_0^{2(1-k_0)^{+}}=1$; that is the entering case treated in (c)
below.

An arrest changes no level, weight or cube family. A completion at density $k_e$ re-levels all of the union $Z_L$ of the
components completed at that date to the top level $k_e$, at weight $1$ \cite[(1.33)]{B-CMP122b}, and it is the only
move that meets
\begin{equation*}
\Gamma_0=\Omega_1^{c}\subseteq\Lambda_1^{c}=Z_1\subseteq Z_{k_e}
\end{equation*}
\cite[(2.1)--(2.3)]{B-CMP119}. The re-inclusion into the large-field domain in \cite{B-CMP122b} only assigns the
component to the large-field domain for the later steps and assigns no level.

No move lowers a level, and none removes a point from the domain. Composing the moves between consecutive dates from
the date of $\mathfrak{B}$ to that of $\mathfrak{B}'$ gives $\Omega_1(\mathfrak{B})\subseteq\Omega_1(\mathfrak{B}')$.
Further $\Omega_0(\mathfrak{B})\subseteq\Omega_0(\mathfrak{B}')$, whether $\Omega_0$ is read as independent of the date
or as the support of $\Omega_1$ of fixed width, a neighbourhood of $\Omega_1$ containing a layer of $M_1$-cubes
\cite[p.~255]{B-CMP119}, which grows with $\Omega_1(\cdot)$. Consequently
$\square\cap\Omega_0(\mathfrak{B})\subseteq\square'\cap\Omega_0(\mathfrak{B}')$ whenever $\square\subset\square'$. This
inclusion is what the restriction of a gauge from a cube of a later family to a cube of an earlier family requires.
Finally, a point of $\Omega_0(\mathfrak{B})$ raw at $\mathfrak{B}'$ is not in $\Omega_1(\mathfrak{B}')$, hence not in
$\Omega_1(\mathfrak{B})$: it is raw at $\mathfrak{B}$.

\emph{Proof of} (c). Let $y$ be raw at both dates. The $U$-clause at $y$ is the same $\gamma$-free clause on the same
index-$0$ cubes, the index-$0$ family being independent of the date. A completed $Z_L$ is a union of coarser aligned
blocks, so an index-$0$ cube lies in $Z_L$ or shares no cell with it. The clause is therefore carried over identically.

Now let $y\in(\Omega_0\setminus\Omega_1)(\mathfrak{B})\cap\Omega_1(\mathfrak{B}')$. By the proof of (b), $y$ entered
$\Omega_1(\cdot)$ through a completion at a date with set $\mathfrak{B}_e$, where
$\mathfrak{B}\prec\mathfrak{B}_e\preceq\mathfrak{B}'$, of density $k_e\ge k$, and there it received level $k_e$ and
weight $1$. Since the level of a point of $\Omega_1(\cdot)$ is at least $1$, $k_e\ge1$. Write $\xi_m=ML^{m}\eta$ and
$\mathsf{S}_m=R_m\xi_m$, with $R_m$ as in Lemma~\ref{4.4:lem-collar-avoidance}, so that $\mathsf{S}_m$ is the side of
Ba{\l}aban's aligned blocks at level $m$ \cite[(3.5)]{B-CMP119}. The index-$0$ cube $\square_0$ of $y$ has side at most
$CML\eta\le C\xi_{k_e}=:s'$. It shares the cell of $y\in Z_L$ and, $Z_L$ being a union of coarser aligned blocks, lies in
$Z_L$. The region $Z_L\cup\Omega_{k_e}$ is a union of blocks of one aligned partition of side $\mathsf{S}_{k_e}$: $Z_L$
is a union of blocks of $Z_{k_e}$, and the blocks of $Z_{k_e}$ and of $\Omega_{k_e}^{c}$ are the cubes of one partition
\cite[(3.5)]{B-CMP119}. Moreover $\mathsf{S}_{k_e}\ge s'$, since $R_{k_e}\ge C$ as in the separation bound
\eqref{4.27:eq-stage-move-cubes-a} (Lemma~\ref{4.27:prf-stage-move-cubes}, on the cube families under a stage move,
whose proof is incomplete at one step). By
Lemma~\ref{4.27:lem-nested-cubes}\,(b), applied to the region $Z_L\cup\Omega_{k_e}$, $\square_0$
lies in a cube $\square'\subset Z_L\cup\Omega_{k_e}$ of side $s'$. After the completion the top family at
$\mathfrak{B}_e$ consists of the cubes of side $C\xi_{k_e}$ inside the new top region, which contains
$\Omega_{k_e}\cup Z_L$, at weight $1$ (the completion move in Lemma~\ref{4.27:prf-cube-clause-heredity}, on the cube
clause under $\mathbf{T}$-steps and completion); so $\square'\subset Z_L\cup\Omega_{k_e}$ is a top cube of
$\mathcal{Q}_{\mathfrak{B}_e}$, whose family weight is $1$.

From $\mathfrak{B}_e$ to $\mathfrak{B}'$ the point $y$ stays in $\Omega_1(\cdot)$ by (b). Therefore
\eqref{4.27:eq-weights-heredity-b} (Lemma~\ref{4.27:lem-weights-heredity}, whose proof is incomplete at one step),
together with Lemma~\ref{4.27:lem-nested-cubes}, applies between these dates. It gives that the clauses of
$\mathfrak{G}_1(\mathfrak{B}')$ at $y$ and on the cubes imply those of $\mathfrak{G}(\mathfrak{B}_e)$ there. At $y$ these
read $Q_a(y)<5r_a(k_e)$ for $a\in\{1,3,5,6,7\}$, with $r_a(m)$ the threshold of entry $a$ at level $m$. On
$\square'\supset\square_0$ they give a gauge transformation $u$ on $\square'$ with $U^{u}=e^{i\eta A}$ and
\begin{equation*}
L^{k_e}\eta|A|<5BCM\alpha_{0,k_e},\qquad (L^{k_e}\eta)^{2}|\nabla^{\eta}A|<5BCM\alpha_{0,k_e},
\end{equation*}
where $B$ is the constant of the cube clause of $\mathfrak{G}$ in \eqref{4.27:eq-weights-heredity-b}.
Transported to index-$0$ units, these thresholds are at most $5\alpha_{\cdot,k_e}L^{-e_ak_e}\le5\alpha_{\cdot,k_e}/L$,
where $e_a\ge1$ is the scaling exponent of entry $a$, as in \eqref{4.27:eq-weights-heredity-b}, and $k_e\ge1$.

The index-$0$ $U$-radii of $\mathfrak{K}_{109}$ majorise every threshold of the form weight$\times$radius whose value at
that index is at most $\gamma_0^{(k_e)}$: the entries $5$ and $8$ give \cite[(3.35)]{B-CMP99b} and the entries $6$ and $7$ give
\cite[(3.37)]{B-CMP99b}. This is the comparison of the clauses of $\mathfrak{K}_{109}$ with the radii of
\cite[(3.35),(3.37)]{B-CMP99b}, used here by statement. It is the same statement through which the second inclusion is
obtained at every index at least $1$. Here
$5\alpha_{\cdot,k_e}\le5\bar\alpha_{k_e}\le\gamma_0^{(k_e)}$, since $\bar\alpha_{k_e}$ is the running maximum over
levels $m\le k_e$ of the maximal radius function, and by \eqref{4.27:eq-weights-heredity-a}. The $U$-members of the
stratum clause of $\mathfrak{G}(\mathfrak{B})$ at $y$ follow. By \eqref{4.27:eq-weights-heredity-b}
(Lemma~\ref{4.27:lem-weights-heredity}, whose proof is incomplete at one step) the thresholds do not rise after
$\mathfrak{B}_e$.

The $J$-part at such an entering point is the clause of entry $3$ at level $k_e$ and weight $1$, namely
$(L^{k_e}\eta)^{3}|\mathbb{J}(y)|<5\alpha_{0,k_e}$. Multiplying by $\eta^{3}(L^{k_e}\eta)^{-3}$, with $\ell=\eta$,
gives the first two members of \eqref{4.27:eq-raw-stratum-a}. For the third, consider first $k_e=k$. Then
$\alpha_{0,k}L^{-3k}\le\alpha_{0,k}\le\bar\alpha_k$. Next let $k_e>k$. The radius-ratio factor $\chi_g$ of
\eqref{4.27:eq-factors-radii-coupling-a}, read as later radius over earlier radius, gives
$\alpha_{0,k_e}\le(1+\beta_0)^{2}(k_e-k)^{1/2}\alpha_{0,k}$. With $\beta_0\le1/7$, as in
\eqref{4.27:eq-weights-heredity-b}, this yields
\begin{equation*}
\begin{split}
5\alpha_{0,k_e}L^{-3k_e}
&\le5\cdot\tfrac{64}{49}(k_e-k)^{1/2}L^{-3k_e}\,\alpha_{0,k}\\
&\le5\alpha_{0,k}\le5\bar\alpha_k .
\end{split}
\end{equation*}
Here $\tfrac{64}{49}(k_e-k)^{1/2}L^{-3k_e}\le1$ for the following reason: $L\ge8$, as stated with the separation
bound \eqref{4.27:eq-stage-move-cubes-a} (Lemma~\ref{4.27:prf-stage-move-cubes}, on the cube families under a stage
move, whose proof is incomplete at one step),
and $k_e-k\le k_e$, and $k_e\ge1$, and $v\mapsto v^{1/2}8^{-3v}$ decreases on $v\ge1$. Hence
\begin{equation*}
\tfrac{64}{49}(k_e-k)^{1/2}L^{-3k_e}\le\tfrac{64}{49}\,k_e^{1/2}8^{-3k_e}\le\tfrac{64}{49}\cdot\tfrac1{512}<1 .
\end{equation*}
The last inequality of \eqref{4.27:eq-raw-stratum-a} holds because, by \eqref{4.27:eq-weights-heredity-a},
$\gamma_0^{(k)}\ge5w_{*}(k)\bar\alpha_k$, and $w_{*}(k)\ge1$. If $y$ is held at the
date of $\mathfrak{B}$ by no live piece, then $T^{\mathfrak{B}}(y)=\gamma_0^{(k)}$, and
\eqref{4.27:eq-raw-stratum-a} gives the stratum $J$-member $\ell^{3}|\mathbb{J}(y)|<T^{\mathfrak{B}}(y)$ of
$\mathfrak{G}(\mathfrak{B})$ at $y$.

\emph{The first inclusion on the stratum, in terms of $\mathfrak{G}_1$ and $\mathfrak{G}^{-J_0}$.} It consists of the
following parts.
\begin{itemize}
\item $\mathfrak{G}_1(\mathfrak{B}')\subset\mathfrak{G}_1(\mathfrak{B})$. This comes from the level and threshold
bounds of \eqref{4.27:eq-weights-heredity-b} (Lemma~\ref{4.27:lem-weights-heredity}, whose proof is incomplete at one
step) and Lemma~\ref{4.27:lem-nested-cubes}, given
$\Omega_1(\mathfrak{B})\subseteq\Omega_1(\mathfrak{B}')$ from (b).
\item $\mathfrak{G}^{-J_0}(\mathfrak{B}')\subset\mathfrak{G}^{-J_0}(\mathfrak{B})$ and
$\mathfrak{G}(\mathfrak{B}')\subset\mathfrak{G}^{-J_0}(\mathfrak{B})$. These follow from the first item together with the
$U$-members in (c) --- the identity at points raw at both dates, the entering case otherwise --- and from
$\mathfrak{G}(\mathfrak{B}')\subset\mathfrak{G}^{-J_0}(\mathfrak{B}')$.
\item Every pair of $\mathfrak{G}_1(\mathfrak{B}')$ obeys every stratum member of $\mathfrak{G}(\mathfrak{B})$ at the
points entering $\Omega_1(\mathfrak{B}')$: the $U$-members by (c), and the $J$-member, at points held by no live piece,
by \eqref{4.27:eq-raw-stratum-a}, derived from the level clauses of $\mathfrak{G}_1(\mathfrak{B}')$. At points held by a
live piece the $J$-member is obeyed by the heredity of the stratum $J$-member at held raw points, a statement
established after the present lemma and not proved in this proof.
\item Every pair of $\mathfrak{G}(\mathfrak{B}')$ obeys the stratum $J$-member of $\mathfrak{G}(\mathfrak{B})$ at held
points raw at both dates, by that same heredity statement, established after the present lemma.
\end{itemize}
The second inclusion is that of the full $\mathfrak{G}$, as in (a).
\end{proof}

Throughout this lemma an operation $\mathbf{R}$ on a component is called \emph{completed} if it reaches
\cite[(1.70)]{B-CMP122b}, whether of the first or of the second class. An arrest lives until the first
later completed $\mathbf{R}$ on the component holding its $X$ (Definition~\ref{4.27:def-arrests-pieces}).
The re-levelling
\cite[(1.33)]{B-CMP122b} is applied to every component that reaches \cite[(1.70)]{B-CMP122b},
components of the second class included. The re-inclusion of second-class components into the
large-field domain in \cite{B-CMP122b} writes membership in the large-field domains of the
$\mathbf{R}$ operations of the later steps and assigns no level. On this reading the completion move
(t4) of Lemma~\ref{4.27:lem-determining-set-order}
lifts a completed component, whatever its class, into $\Omega_1(\cdot)$ at its completion date.

If a completion of the second class instead left the points of its component at index $0$, then a
point $y$ held by the innermost piece $\mathfrak{c}_1$ of the lemma below could be raw and held by no
live piece from the dissolution of $\mathfrak{c}_1$ on. Case (1) of the proof would then end at that
date, and the $\mathbb{J}$-inclusion at $y$ after it would fall under the case of raw points held by
no live piece, where heredity across an increase of density is not asserted. The frozen terms of the
arrest would be unaffected: a frozen term whose domain meets the component of the dissolving nest
ceases at that operation, by the definition of the operations at which frozen terms are consumers.

\begin{lemma}[The current bound at held raw points]\label{4.27:prf-held-raw-points}
Let $\mathfrak{B}$ be the determining set of a date $d$ of density $k$. Let
$y\in(\Omega_0\setminus\Omega_1)(\mathfrak{B})$ be held at $d$ by a live piece, that is,
$y\in W_{\mathfrak{c}}$ for an arrest $\mathfrak{c}$ attached to a date $\preccurlyeq d$ and not
dissolved at a date $\preccurlyeq d$. Then the following hold.
\begin{itemize}
\item The live arrests $\mathfrak{c}$ at $d$ with $y\in X_{\mathfrak{c}}$ form a nest
$\mathfrak{c}_1\subset\mathfrak{c}_2\subset\cdots$. The innermost arrest $\mathfrak{c}_1$ is the one of
smallest density, and
\begin{equation}\label{4.27:eq-held-raw-points-1}
T^{\mathfrak{B}}(y)=\gamma_0^{(k_{\mathfrak{c}_1})}\le\gamma_0^{(k)} .
\end{equation}
\item For every determining set $\mathfrak{B}'\succcurlyeq\mathfrak{B}$ we have
\begin{equation}\label{4.27:eq-held-raw-points-a}
\begin{aligned}
&T^{\mathfrak{B}'}(y)=T^{\mathfrak{B}}(y)
&&\text{if } y\in(\Omega_0\setminus\Omega_1)(\mathfrak{B}'),\\
&\ell^{3}|\mathbb{J}(y)|<T^{\mathfrak{B}}(y)\ \text{ for every pair of }\mathfrak{G}_1(\mathfrak{B}')
&&\text{if } y\in\Omega_1(\mathfrak{B}').
\end{aligned}
\end{equation}
\end{itemize}
In particular, every pair of $\mathfrak{G}(\mathfrak{B}')$ obeys the $\mathbb{J}$-member
$\ell^3|\mathbb{J}(y)|<T^{\mathfrak{B}}(y)$ of $\mathfrak{G}(\mathfrak{B})$ at $y$.
\end{lemma}

\begin{proof}
We recall two facts. First, in the definition \eqref{4.27:eq-weights-heredity-b} of
$\mathfrak{G}(\mathfrak{B})$, the threshold $T^{\mathfrak{B}}(y)$ at a raw point $y$ is
$\gamma_0^{(k_{\mathfrak{c}})}$ for the innermost live piece $W_{\mathfrak{c}}\ni y$ at the date of
$\mathfrak{B}$, and $\gamma_0^{(k)}$ if there is no such piece. Second,
\begin{equation*}
\gamma_0^{(k)}=5\max_{j\le k}w_*(j)\,\bar\alpha_j ,
\end{equation*}
which is non-decreasing in $k$.

\emph{The nest at $d$.} Let $\mathfrak{c}$ be a live arrest at $d$ with $y\in X_{\mathfrak{c}}$. We have
$\Omega_1\supseteq\Omega_{k_{\mathfrak{c}}}$ by \cite[(2.1)]{B-CMP119}. We also have
$\Omega_1(\mathfrak{B})\supseteq\Omega_1$, because the domains $\Omega_0(\cdot)$ and $\Omega_1(\cdot)$
never shrink along the order of dates (this is shown in the proof of
Lemma~\ref{4.27:prf-raw-stratum}). Since $y\notin\Omega_1(\mathfrak{B})$, it follows that
$y\in\Omega^{c}_{k_{\mathfrak{c}}}$, and therefore $y\in X_{\mathfrak{c}}\cap\Omega^c_{k_{\mathfrak{c}}}
=W_{\mathfrak{c}}$.

Any two such arrests are therefore live and their pieces both contain $y$. Hence
Lemma~\ref{4.27:lem-arrest-timing} applies to the pair, ordered so that the first has density at most
that of the second. By (T1) of \eqref{4.27:eq-arrest-timing-a} their densities differ, and the one of
smaller density has the smaller component. By (T4) they are nested,
$\mathfrak{c}_1\subset\mathfrak{c}_2\subset\cdots$. The innermost arrest $\mathfrak{c}_1$ is therefore
the oldest one. Since $\gamma_0^{(\cdot)}$ is non-decreasing,
\begin{equation*}
T^{\mathfrak{B}}(y)=\gamma_0^{(k_{\mathfrak{c}_1})}=\min_s\gamma_0^{(k_{\mathfrak{c}_s})} .
\end{equation*}
The arrest $\mathfrak{c}_1$ is attached to a date $\preccurlyeq d$. Densities do not decrease along the
order of dates (Lemma~\ref{4.27:lem-determining-set-order}, \eqref{4.27:eq-determining-set-order-a}),
so $k_{\mathfrak{c}_1}\le k$. Hence $\gamma_0^{(k_{\mathfrak{c}_1})}\le\gamma_0^{(k)}$, which proves
\eqref{4.27:eq-held-raw-points-1}.

\emph{The interval of liveness.} Let $\mathfrak{B}'\succcurlyeq\mathfrak{B}$ have date
$d'\succcurlyeq d$. By Definition~\ref{4.27:def-arrests-pieces} and the dissolution rule in
\eqref{4.27:eq-arrested-spaces-a}, an arrest is live on an interval of dates. This interval runs from
its attachment date to the date $(k_e;\mathrm{L})$ of the first later completed $\mathbf{R}$ on the
component then holding its $X$. Let $(k_e;\mathrm{L})$ denote this date for $\mathfrak{c}_1$.

At that date the completion move (t4) re-levels the completed components $Z_{\mathrm{L}}$ to level
$k_e\ge1$ at weight $1$, by \cite[(1.33)]{B-CMP122b} and the pointwise description of the completion
move in the proof of Lemma~\ref{4.27:lem-determining-set-order}. These completed components contain
$X_{\mathfrak{c}_1}\ni y$. Indeed, $X_{\mathfrak{c}_1}$ is connected and lies in
$Z_{k_{\mathfrak{c}_1}}\subset Z_{k_e}$ \cite[(2.1)]{B-CMP119}. Components of nested sets are nested
or disjoint, as in the proof of the last assertion of Lemma~\ref{4.27:lem-arrest-timing}. Hence
$X_{\mathfrak{c}_1}$ lies in the component of $Z_{k_e}$ that contains $y$. Since $\Omega_1(\cdot)$
never shrinks, $y\in\Omega_1(\cdot)$ from $(k_e;\mathrm{L})$ on.

Before $(k_e;\mathrm{L})$ the point $y$ is raw. By the proof of Lemma~\ref{4.27:prf-raw-stratum}, the
completion move (t4) is the only move that meets $\Gamma_0=\Omega_1^{c}$, and no move lowers a level.
Moreover, a completed component containing $y$ contains $X_{\mathfrak{c}_1}$, by the argument just
given. So an earlier lifting of $y$ would be an earlier completed $\mathbf{R}$ on the component holding
$X_{\mathfrak{c}_1}$. This contradicts the choice of $(k_e;\mathrm{L})$.

\emph{Case (1): $y\in(\Omega_0\setminus\Omega_1)(\mathfrak{B}')$.} By the preceding paragraph,
$d'\prec(k_e;\mathrm{L})$, so $\mathfrak{c}_1$ is live at $d'$. We show that $\mathfrak{c}_1$ is still
the innermost live arrest holding $y$ at $d'$.

Let $\mathfrak{b}$ be an arrest live at $d'$ with $y\in X_{\mathfrak{b}}$. Then
$y\in W_{\mathfrak{b}}$, by the argument of the first paragraph applied at $d'$. There are two
possibilities.
\begin{itemize}
\item $\mathfrak{b}$ is attached to a date $\preccurlyeq d$. Since liveness is an interval and
$d\preccurlyeq d'$, the arrest $\mathfrak{b}$ is live at $d$. It is then one of the
$\mathfrak{c}_s$, and it encloses $\mathfrak{c}_1$.
\item $\mathfrak{b}$ is attached to a date in $\mathopen{]}d,d'\mathclose{]}$. Its density then
satisfies $k_{\mathfrak{b}}\ge k\ge k_{\mathfrak{c}_1}$.
\end{itemize}
In the second case the equality $k_{\mathfrak{b}}=k_{\mathfrak{c}_1}$ is impossible. The class
components of one step are disjoint, and each exits at most once: an excluded component is not
changed by the subsequent operations of its site \cite{B-CMP122a}, and (T1) of
\eqref{4.27:eq-arrest-timing-a} applies. The class component holding $y$ at step
$k_{\mathfrak{c}_1}$ is $X_{\mathfrak{c}_1}$, and it was arrested at a date $\preccurlyeq d$. Hence
$k_{\mathfrak{b}}>k_{\mathfrak{c}_1}$. Since $y\in W_{\mathfrak{c}_1}\cap W_{\mathfrak{b}}$, (T1) gives
$X_{\mathfrak{c}_1}\subset X_{\mathfrak{b}}$.

In both cases $\mathfrak{b}$ encloses $\mathfrak{c}_1$. Hence
$T^{\mathfrak{B}'}(y)=\gamma_0^{(k_{\mathfrak{c}_1})}=T^{\mathfrak{B}}(y)$: at $y$ the
$\mathbb{J}$-clause of $\mathfrak{G}(\mathfrak{B}')$ coincides with that of $\mathfrak{G}(\mathfrak{B})$.
This is the first line of \eqref{4.27:eq-held-raw-points-a}.

\emph{Case (2): $y\in\Omega_1(\mathfrak{B}')$.} The point $y$ entered $\Omega_1(\cdot)$ at the date
$(k_e;\mathrm{L})$, with determining set
$\mathfrak{B}_e:=\mathfrak{B}_{(k_e;\mathrm{L})}\in\mathopen{]}\mathfrak{B},\mathfrak{B}'\mathclose{]}$,
at level $k_e$ and weight $1$. Here $k_e\ge k$. Moreover $k_e>k_{\mathfrak{c}_1}$, because an arrested
component is not completed at its own step.

The entering line \eqref{4.27:eq-raw-stratum-a} of Lemma~\ref{4.27:prf-raw-stratum} reads only the
clauses at $y$ and on the cube of $\mathcal{Q}_{\mathfrak{B}'}$ containing the index-$0$ cube of $y$.
It therefore applies to
every pair of $\mathfrak{G}_1(\mathfrak{B}')$, and a fortiori to every pair of
$\mathfrak{G}(\mathfrak{B}')$. We run it with $k_{\mathfrak{c}_1}$ in place of $k$:
\begin{equation}\label{4.27:eq-held-raw-points-2}
\begin{aligned}
\ell^3|\mathbb{J}(y)|
&<5\alpha_{0,k_e}L^{-3k_e}
\le 5(1+\beta_0)^2(k_e-k_{\mathfrak{c}_1})^{1/2}L^{-3k_e}\,\alpha_{0,k_{\mathfrak{c}_1}}\\
&\le 5\alpha_{0,k_{\mathfrak{c}_1}}
\le 5\bar\alpha_{k_{\mathfrak{c}_1}}
\le\gamma_0^{(k_{\mathfrak{c}_1})}
=T^{\mathfrak{B}}(y).
\end{aligned}
\end{equation}
We justify each step in turn.
\begin{itemize}
\item The first inequality is the entering line at the density $k_e$ of the entering date.
\item The second is the ratio of radii \eqref{4.27:eq-factors-radii-coupling-a} at the gap
$k_e-k_{\mathfrak{c}_1}\ge1$, read later over earlier.
\item For the third, $(1+\beta_0)^2\le 64/49$ because $\beta_0\le1/7$. Also
$(k_e-k_{\mathfrak{c}_1})^{1/2}\le k_e^{1/2}$, and $L^{-3k_e}\le 8^{-3k_e}$, the blocking factor
satisfying $L\ge8$, as recorded in the separation estimates for the cube families in the proof of the
heredity statement \eqref{4.27:eq-weights-heredity-b}.
Since $k\mapsto k^{1/2}8^{-3k}$ decreases on $k\ge1$ and $k_e\ge1$,
\begin{equation*}
\tfrac{64}{49}(k_e-k_{\mathfrak{c}_1})^{1/2}L^{-3k_e}
\le\tfrac{64}{49}\,k_e^{1/2}8^{-3k_e}
\le\tfrac{64}{49}\cdot\tfrac{1}{512}<1 .
\end{equation*}
\item The fourth uses
\begin{equation*}
\alpha_{0,j}=\alpha_0(\mathsf{t}_j)\le\mathsf{A}(\mathsf{t}_j)\le\bar\alpha_j .
\end{equation*}
\item The fifth uses $w_*(\cdot)\ge1$ and the formula for $\gamma_0^{(\cdot)}$ recalled above.
\item The final equality is \eqref{4.27:eq-held-raw-points-1}.
\end{itemize}
This is the second line of \eqref{4.27:eq-held-raw-points-a}.

\emph{Conclusion.} In both cases, at every raw point of $\Omega_0(\mathfrak{B})$ held at $\mathfrak{B}$ by
a live piece, every pair of $\mathfrak{G}(\mathfrak{B}')$ obeys the $\mathbb{J}$-clause of
$\mathfrak{G}(\mathfrak{B})$. In case (1) this holds because the $\mathbb{J}$-clause of
$\mathfrak{G}(\mathfrak{B}')$ at $y$ is that of $\mathfrak{G}(\mathfrak{B})$; in case (2) it holds by
\eqref{4.27:eq-held-raw-points-2}, which is valid for every pair of
$\mathfrak{G}_1(\mathfrak{B}')\supset\mathfrak{G}(\mathfrak{B}')$.

Combine this with two results from the proof of Lemma~\ref{4.27:prf-raw-stratum}: the monotonicity
of the domains, and the statement for the stratum's $\mathbb{U}$-members together with the entering
line. Together they give the first inclusion of the heredity statement
\eqref{4.27:eq-weights-heredity-b} on the stratum, at every point where that inclusion is asserted.
\end{proof}

\begin{lemma}[Current thresholds across an increase of density. Stated; proof incomplete at the step
marked $(\ast)$]\label{4.27:prf-density-increase}
Work with the arrests, pieces and frozen terms of Definition~\ref{4.27:def-arrests-pieces}, the
later operations of Definition~\ref{4.27:def-operations-zones}, and the dates and determining sets
of Lemma~\ref{4.27:lem-determining-set-order}. Raw and held points and the dated threshold
$T^{\mathfrak{B}}(y)$ of the stratum $J$-member of the rider $\mathfrak{G}(\mathfrak{B})$ are as in
\eqref{4.27:eq-weights-heredity-b} and in the current bound at held raw points
(Lemma~\ref{4.27:prf-held-raw-points}). We recall that a point $y$ is raw at a determining set
$\mathfrak{B}$ if $y\in(\Omega_0\setminus\Omega_1)(\mathfrak{B})$, and held at $\mathfrak{B}$ if it
lies in the piece of an arrest live at the date of $\mathfrak{B}$; at a raw point $y$ the threshold
$T^{\mathfrak{B}}(y)$ has the following value:
\begin{itemize}
\item if $y$ is held at $\mathfrak{B}$, it is $\gamma_0^{(k_c)}$, where $k_c$ is the density of the
innermost live piece holding $y$;
\item otherwise it is $\gamma_0^{(k)}$, where $k$ is the density of the date of $\mathfrak{B}$.
\end{itemize}
Here
\begin{equation*}
\gamma_0^{(k)}=5\max_{j\le k}w_{*}(j)\bar\alpha_j ,
\end{equation*}
which is non-decreasing in $k$.

Let $\mathfrak{a}$ be an arrest of density $k_a$, with component $X_{\mathfrak{a}}$, piece
$W_{\mathfrak{a}}$, number of performed stages $n$ and frozen index $n^+$. Let
$\tau\in\mathfrak{F}(\mathfrak{a})$ have domain $X_\tau$, and let $\mathfrak{B}_\tau$ be the set at
whose date $\tau$ is born. Thus $\mathfrak{B}_\tau=\mathfrak{B}_{(k_a;i)}$ if $\tau$ is the new
boundary term $\mathbb{C}^{(i)}_{k_a}$ of stage $i$, where $1\le i\le n\le n^+$, and
$\mathfrak{B}_\tau=\mathfrak{B}_{(k_a;n^+)}$ if $\tau$ is one of the terms frozen in the branch of
the arrest in which one further stage is performed, so that $n^+=n+1$
(Definition~\ref{4.27:def-arrests-pieces}).

Let $\mathfrak{O}$ be a later operation at which $\tau$ is a consumer. Let $\mathfrak{B}^t$ be the
determining set in force immediately before the map of $\mathfrak{O}$ (its target set), and let
$\mathfrak{B}'$ be the set in force immediately after it (its source set), so that $\mathfrak{a}$ is
live at the date of $\mathfrak{B}^t$. Let $y\in X_\tau$ be raw at $\mathfrak{B}^t$. Distinguish four
cases:
\begin{enumerate}
\item[(1)] $y\in X_{\mathfrak{a}}$;
\item[(2)] $y\notin X_{\mathfrak{a}}$, and $y$ is held at $\mathfrak{B}_\tau$;
\item[(3)] $y\notin X_{\mathfrak{a}}$, $y$ is held by no live piece at $\mathfrak{B}_\tau$, and $y$
lies in a component of the class $Z\subset Z_{k_a}$ of step $k_a$;
\item[(4)] $y\notin X_{\mathfrak{a}}$, $y$ is held by no live piece at $\mathfrak{B}_\tau$, and $y$
lies in no class component of step $k_a$.
\end{enumerate}
Case (4) is split according to the density of the set read: for
$\mathfrak{B}''\in\{\mathfrak{B}^t,\mathfrak{B}'\}$, case (4a) concerns a set $\mathfrak{B}''$ of
density $k_a$, and case (4b) a set $\mathfrak{B}''$ of density $k'>k_a$. Then, as long as $\tau$ is a
consumer,
\begin{equation}\label{4.27:eq-density-increase-a}
\begin{aligned}
&\text{in cases (1), (2) and (3):}\\
&\qquad T^{\mathfrak{B}^t}(y)=T^{\mathfrak{B}_\tau}(y);\\
&\qquad T^{\mathfrak{B}'}(y)=T^{\mathfrak{B}^t}(y)
\qquad\text{if } y\in(\Omega_0\setminus\Omega_1)(\mathfrak{B}'),\\
&\qquad \ell^{3}|\mathbb{J}(y)|<T^{\mathfrak{B}_\tau}(y)\ \text{ for every pair of }
\mathfrak{G}_1(\mathfrak{B}')
\qquad\text{if } y\in\Omega_1(\mathfrak{B}');\\
&\text{in case (4a):}\\
&\qquad T^{\mathfrak{B}''}(y)=T^{\mathfrak{B}_\tau}(y),\ \text{ $y$ being raw at }\mathfrak{B}'';\\
&\text{in case (4b), marked }(\ast):\\
&\qquad T^{\mathfrak{B}''}(y)\ge\gamma_0^{(k_a+1)}\ge\gamma_0^{(k_a)}=T^{\mathfrak{B}_\tau}(y)
\qquad\text{if } y\in(\Omega_0\setminus\Omega_1)(\mathfrak{B}'').
\end{aligned}
\end{equation}
In case (4b) the agreement $T^{\mathfrak{B}''}(y)=T^{\mathfrak{B}_\tau}(y)$, which in cases (1),
(2) and (3) is given by the first two clauses of the display and in case (4a) by the clause for
that case, is not proved; this is the step marked $(\ast)$. What is proved in case (4b) is the
inequality displayed above and the criterion that $T^{\mathfrak{B}''}(y)=T^{\mathfrak{B}_\tau}(y)$
holds if and only if the maximum
defining $\gamma_0^{(\cdot)}$ has not grown on $\,]k_a,k_b]$, where $k_b$ is the density of the
oldest live holder of $y$ at $\mathfrak{B}''$ attached after the dates of density $k_a$,
respectively on $\,]k_a,k']$ if there is no such holder.
Moreover, $T^{\mathfrak{B}}(y)\le\gamma_0^{(k_a)}$ at every raw point $y\in X_{\mathfrak{a}}$ and
every set $\mathfrak{B}$ at whose date $\mathfrak{a}$ is live.
\end{lemma}

\begin{proof}
\emph{Preliminaries.} By the time order of determining sets
(Lemma~\ref{4.27:lem-determining-set-order}), the date of $\mathfrak{B}_\tau$ has density
$k_\tau=k_a$. The date of $\mathfrak{B}^t$ is not earlier than $(k_a;n^+)$, and $(k_a;n^+)$ is not
earlier than the date of $\mathfrak{B}_\tau$.

Since $\tau$ is a consumer at $\mathfrak{O}$, the arrest $\mathfrak{a}$ is live at $\mathfrak{B}^t$.
When $\mathfrak{O}$ is the dissolution (O5) of the nest containing $\mathfrak{a}$, on the dissolving
piece the structure of $\mathfrak{B}^t$ is still the frozen structure of the arrest (same lemma).

Three facts are used throughout.
\begin{itemize}
\item An arrest is live on an interval of dates. The interval runs from its attachment date to the
date of the first later completed large-field operation on the component of the large-field region
then containing the arrest's own component, at which its nest dissolves
(Definition~\ref{4.27:def-arrests-pieces}).
\item Let two live arrests have pieces sharing a point. Then their densities differ, and the one of
smaller density has the smaller component, by \eqref{4.27:eq-arrest-timing-a}. Hence the live pieces
holding a point are nested, and the innermost one is the oldest.
\item The class components of one step are pairwise disjoint, and each of them takes at most one
large-field exit in that step (Definition~\ref{4.27:def-arrests-pieces},
\eqref{4.27:eq-arrest-timing-a}).
\end{itemize}

\emph{Case (1): $y\in X_{\mathfrak{a}}$.} The complement of $\Omega_1(\mathfrak{B}^t)$ is contained
in $\Omega^c_{k_a}$, so $y\in W_{\mathfrak{a}}$.

The innermost live piece holding $y$ at $\mathfrak{B}^t$ is either $W_{\mathfrak{a}}$ itself,
$\mathfrak{a}$ being attached at $(k_a;n^+)$, which is not later than the date of
$\mathfrak{B}^t$, or the piece $W_{\mathfrak{c}}\subset X_{\mathfrak{a}}$ of an older arrest
$\mathfrak{c}$, with $k_c<k_a$, nested in $W_{\mathfrak{a}}$ by
\eqref{4.27:eq-arrest-timing-a}. No other piece of density $k_a$ can hold $y$, because the class
components of one step are disjoint and each exits once.

\emph{Second alternative.} The arrest $\mathfrak{c}$ is attached at a date of density $k_c<k_a$,
hence before every date of density $k_a$. It is live at $\mathfrak{B}^t$. Since liveness is an
interval, it is therefore live at $\mathfrak{B}_\tau$. By the first assertion of the current
bound at held raw points (Lemma~\ref{4.27:prf-held-raw-points},
\eqref{4.27:eq-held-raw-points-a}), by which thresholds persist at held raw points, it is also the
innermost live piece holding $y$ there. Hence
\begin{equation*}
T^{\mathfrak{B}_\tau}(y)=\gamma_0^{(k_c)}=T^{\mathfrak{B}^t}(y).
\end{equation*}

\emph{First alternative.} No piece older than $\mathfrak{a}$ holds $y$ at $\mathfrak{B}_\tau$
either. Such a piece would be live at $\mathfrak{B}^t$, because the nest of $\mathfrak{a}$ dissolves
together, by \eqref{4.27:eq-arrest-timing-a}, and it would be innermost there. Two sub-cases occur.
\begin{itemize}
\item $\mathfrak{a}$ is attached at the date of $\mathfrak{B}_\tau$ (that is, $i=n^+$). Then $y$ is
held at $\mathfrak{B}_\tau$ by $\mathfrak{a}$, and $T^{\mathfrak{B}_\tau}(y)=\gamma_0^{(k_a)}$.
\item $\mathfrak{a}$ is attached later ($i<n^+$). Then $y$ is held by no live piece at
$\mathfrak{B}_\tau$, and $T^{\mathfrak{B}_\tau}(y)$ is the value $\gamma_0^{(k_a)}$ that an unheld
raw point takes at a set whose date has density $k_a$, here $\mathfrak{B}_\tau$.
\end{itemize}
In both sub-cases $T^{\mathfrak{B}^t}(y)=\gamma_0^{(k_a)}$, so the two thresholds are equal.

Moreover $y$ stays raw exactly as long as $\mathfrak{a}$ lives. By
Lemma~\ref{4.27:prf-held-raw-points}, only the dissolution of the nest of $\mathfrak{a}$ lifts $y$
into $\Omega_1$.

At every raw point $y$ of $X_{\mathfrak{a}}$ and every set $\mathfrak{B}$ at whose date
$\mathfrak{a}$ is live one has $T^{\mathfrak{B}}(y)\le\gamma_0^{(k_a)}$. Indeed
$T^{\mathfrak{B}}(y)=\gamma_0^{(k_c)}$ with $k_c\le k_a$, and $\gamma_0^{(\cdot)}$ is
non-decreasing. Hence the $J$-member
$\ell^{3}|\mathbb{J}|\le\gamma_0^{(k_a)}$ of the ambient condition of the arrested stage on
$X_{\mathfrak{a}}$ is implied there, at the same dating.

\emph{Case (2): $y\notin X_{\mathfrak{a}}$, and $y$ lies in a live piece at $\mathfrak{B}_\tau$.}
This is the current bound at held raw points, Lemma~\ref{4.27:prf-held-raw-points}. Its first
assertion \eqref{4.27:eq-held-raw-points-a} gives, at every later set at which $y$ is raw, a
threshold equal to $T^{\mathfrak{B}_\tau}(y)$; in particular
$T^{\mathfrak{B}^t}(y)=T^{\mathfrak{B}_\tau}(y)$.

\emph{Case (3): $y\notin X_{\mathfrak{a}}$, $y$ lies in no live piece at $\mathfrak{B}_\tau$, but $y$
lies in a component $X'$ of the class $Z\subset Z_{k_a}$ of step $k_a$.} Here
$X'\ne X_{\mathfrak{a}}$, the two are disjoint, and $y\in\Gamma_0\subset Z_{k_a}$.

\emph{Target sets at density $k_a$.} While the dates run through density $k_a$, namely at the stages
$l$ of (O1) and at $(k_a;L)$ for (O2),
\begin{equation*}
T^{\mathfrak{B}^t}(y)=\gamma_0^{(k_a)}=T^{\mathfrak{B}_\tau}(y),
\end{equation*}
whether $y$ is held or not. Indeed an arrest of $X'$ has density $k_a$. An older piece inside $X'$
that holds $y$ and is live at $\mathfrak{B}^t$ was live at $\mathfrak{B}_\tau$, which is case (2).

\emph{Target sets past density $k_a$.} Past the dates of density $k_a$, the last of which is
$(k_a;L)$ when a component completes there, one of two things has happened.
\begin{itemize}
\item $X'$ was arrested at step $k_a$. Then $y$ is held at the frozen threshold $\gamma_0^{(k_a)}$
while that arrest lives. Its dissolution is a completed chain met by $X_\tau\ni y$; this is the (O5)
at which $\tau$ ceases to be a consumer (Definition~\ref{4.27:def-operations-zones}), and on the
piece of that arrest the structure of the target set of that (O5) is still the frozen structure of
the arrest.
\item $X'$ completed at $(k_a;L)$. Since $X_\tau$ meets $X'$, the term $\tau$ was absorbed there by
the completion and ceased to be a consumer (Definition~\ref{4.27:def-operations-zones};
\cite[\S1]{B-CMP122b}).
\end{itemize}
In every case the datings at the target sets agree while $\tau$ is a consumer.

\emph{Source sets.} The same holds at the source set $\mathfrak{B}'$ of any such $\mathfrak{O}$,
which lies one move after $\mathfrak{B}^t$ (Lemma~\ref{4.27:lem-determining-set-order}). This is what
the continuation clause of the arrest theorem reads on $X_\tau\setminus Y'$, where $Y'$ is the
domain over which the source of $\mathfrak{O}$ is written.

During density $k_a$ we have $\mathfrak{B}'=\mathfrak{B}_{(k_a;l)}$ for (O1). For (O2) at $k=k_a$,
$\mathfrak{B}'$ is the image under \cite[(1.33)]{B-CMP122b} of the determining set
$\mathbb{B}''_{k_a}$ of the site. Either way $\mathfrak{B}'$ has density $k_a$. The only pieces that
can hold $y$ at $\mathfrak{B}'$ are arrests of $X'$ attached at dates of density $k_a$: an older
piece live at $\mathfrak{B}'$ was live at $\mathfrak{B}_\tau$, which is excluded. Hence
\begin{equation*}
T^{\mathfrak{B}'}(y)=\gamma_0^{(k_a)}=T^{\mathfrak{B}_\tau}(y)
\end{equation*}
while $y$ is raw at $\mathfrak{B}'$.

If $X'$ completes at $(k_a;L)$, then $y$ enters $\Omega_1(\mathfrak{B}')$. The entering bound of the
heredity on the raw stratum \eqref{4.27:eq-raw-stratum-a} then gives, for every pair of
$\mathfrak{G}_1(\mathfrak{B}')$ at $y$,
\begin{equation*}
\ell^{3}|\mathbb{J}(y)|<5\alpha_{0,k_a}L^{-3k_a}\le\gamma_0^{(k_a)}=T^{\mathfrak{B}_\tau}(y).
\end{equation*}

Now consider the dates past density $k_a$, the last date of density $k_a$ being $(k_a;L)$ when a
component completes there. This includes (O3) from $k_a$ to $k_a+1$, whose target set is the last
date of density
$k_a$. Every class component of a step is either arrested at one of its dates up to $(k;N_k)$, at one
of the large-field exits, or completed at $(k;L)$ (Definition~\ref{4.27:def-arrests-pieces}). In the
second alternative $\tau$ has ceased, so $X'$ is arrested. Then $y$ is held at $\mathfrak{B}^t$ by
the arrest of $X'$, and this arrest is innermost: anything older that holds $y$ and is live at
$\mathfrak{B}^t$ was live at $\mathfrak{B}_\tau$. The first assertion of
\eqref{4.27:eq-held-raw-points-a} gives $T^{\mathfrak{B}'}(y)=T^{\mathfrak{B}^t}(y)$ while $y$ is
raw at $\mathfrak{B}'$. Its second assertion, the bound at points entering $\Omega_1$, serves where
$y$ enters $\Omega_1(\mathfrak{B}')$.

\emph{Cases (1) and (2) at the source sets.} In these cases $y$ is held at $\mathfrak{B}^t$: by
$\mathfrak{a}$ or an older nested piece in case (1), by the innermost live piece holding it in case
(2). Both assertions of Lemma~\ref{4.27:prf-held-raw-points}, applied to the pair
$(\mathfrak{B}^t,\mathfrak{B}')$, serve alike. They give
\begin{equation*}
T^{\mathfrak{B}'}(y)=T^{\mathfrak{B}^t}(y)=T^{\mathfrak{B}_\tau}(y)
\end{equation*}
while $y$ is raw at $\mathfrak{B}'$, and the bound of \eqref{4.27:eq-density-increase-a} where $y$
enters $\Omega_1(\mathfrak{B}')$.

\emph{Case (4): $y\notin X_{\mathfrak{a}}$, $y$ lies in no live piece at $\mathfrak{B}_\tau$, and $y$
lies in no class component of step $k_a$.} Every raw point lies in
$Z_{k_a}\supseteq Z_1\supseteq\Gamma_0$. So $y$ is a raw point of a component $X''$ of $Z_{k_a}$
that is not a class component of step $k_a$: $X''$ fails one of the two conditions defining the
class (Definition~\ref{4.27:def-operations-zones}).

Among such components are those holding a second-class component of an earlier step; such a
second-class component, its large-field operation having been completed, is included again into the
large-field region \cite[\S1]{B-CMP122b}, and its re-inclusion counts as a creation for the timing
condition of the class. A class component of step $k_a$ that completes at $(k_a;L)$, whether as first
or second class, falls under case (3).

\emph{Case (4a): dates of density $k_a$.} No arrest of density $k_a$ ever holds $y$. Its component
would be a class component of step $k_a$ holding $y$, hence the component of $Z_{k_a}$ through
$y$, that is $X''$, which is not a class component.

Therefore the second assertion of Lemma~\ref{4.27:prf-unheld-raw-points} below (unheld raw points
within one density), whose proof does not use the present lemma and is incomplete at one step, is
void at $y$. Its first assertion \eqref{4.27:eq-unheld-raw-points-a} gives
\begin{equation*}
T^{\mathfrak{B}''}(y)=\gamma_0^{(k_a)}=T^{\mathfrak{B}_\tau}(y)
\qquad\text{for every set }\mathfrak{B}''\succcurlyeq\mathfrak{B}_\tau\text{ of density }k_a ,
\end{equation*}
in particular for $\mathfrak{B}''\in\{\mathfrak{B}^t,\mathfrak{B}'\}$ when it has density $k_a$.
These sets are the stages $(k_a;l)$ with $l\ge i$, the date $(k_a;N_{k_a})$, and the date $(k_a;L)$
with its image under \cite[(1.33)]{B-CMP122b}. They are every source and target set of (O1) and of
(O2) at $k=k_a$, and the target set of (O3) from $k_a$ to $k_a+1$, which is the last date of density
$k_a$.

The point $y$ stays raw through all of them. No move between consecutive dates lowers a level, and
the only move that meets the raw stratum is the completion move
(Lemma~\ref{4.27:lem-determining-set-order}, together with the monotonicity of $\Omega_0(\cdot)$,
$\Omega_1(\cdot)$ and of levels established with the heredity on the raw stratum). The stage move
raises
levels only inside the intermediate domain $\Omega''_{h+1}\subset\Omega_1$ of the site, where
$h=k_a-N_{k_a}$. The completion move at $(k_a;L)$ re-levels the completed class components only,
and $X''$ is none of them. The arrest move writes nothing.

Finally, an older arrest ($k_c<k_a$) attached after $\mathfrak{B}_\tau$ is impossible, because the
dates after $\mathfrak{B}_\tau$ have density $\ge k_a$. Nothing else can intervene.

\emph{Case (4b) $(\ast)$: dates past density $k_a$.} Let $\mathfrak{B}''$ be either of
$\mathfrak{B}^t,\mathfrak{B}'$, of density $k'>k_a$, with $y$ raw at $\mathfrak{B}''$. Older holders
are excluded as in case (4a) and in
Lemma~\ref{4.27:prf-unheld-raw-points}, whose proof is incomplete at one step. Hence
$T^{\mathfrak{B}''}(y)$ is one of two values:
\begin{itemize}
\item $\gamma_0^{(k_b)}$, where $k_b$ is the density of the oldest live holder of $y$ at
$\mathfrak{B}''$ attached after the dates of density $k_a$, and $k_a<k_b\le k'$;
\item $\gamma_0^{(k')}$, if there is no such holder.
\end{itemize}
In either case, since $\gamma_0^{(\cdot)}$ is non-decreasing,
\begin{equation*}
T^{\mathfrak{B}''}(y)\ \ge\ \gamma_0^{(k_a+1)}\ \ge\ \gamma_0^{(k_a)}=T^{\mathfrak{B}_\tau}(y).
\end{equation*}
Equality holds if and only if the maximum defining $\gamma_0^{(\cdot)}$ has not grown on
$\,]k_a,k_b]$, respectively on $\,]k_a,k']$.

At the source set $\mathbb{B}_{k_a+1}$ of (O3) from $k_a$ to $k_a+1$ the threshold is exactly
$\gamma_0^{(k_a+1)}$: if $y$ is unheld there, this is the value $\gamma_0^{(k_a+1)}$ that an unheld
raw point takes at a set whose date has density $k_a+1$; a holder attached at $(k_a+1;\top)$ is an
arrest of step $k_a+1$ taken before
its first stage, hence of density $k_a+1$ (Lemma~\ref{4.27:lem-determining-set-order}).

The sets $\mathfrak{B}''$ of density $k'>k_a$ that occur are the target and source sets of
$\mathfrak{O}$ at the dates past density $k_a$; among them are the source sets of (O3) from the step
$k_a\to k_a+1$ onward, of (O2) at $k>k_a$, and of (O4), (O4$'$)
and (O5). This is the case of a raw point held by no live piece at any date of the earlier density
and read across an increase of density, which is the step marked $(\ast)$ of
Lemma~\ref{4.27:prf-unheld-raw-points}, at which its proof is incomplete. At these
dates the equality of cases (1), (2), (3) and (4a) in \eqref{4.27:eq-density-increase-a} is not
proved.

\emph{Conclusion.} For the consumers of the arrest theorem, the dating of the stratum $J$-member of
the rider agrees with that of every later target and source set at all raw points of $X_\tau$
except those of case (4b).

The same holds against the rider $\mathfrak{G}(\mathfrak{B}')$ of every later source set at such
points. This follows from the first assertion of \eqref{4.27:eq-held-raw-points-a} with
$\mathfrak{B}'$ in place of $\mathfrak{B}^t$, and from the treatment of source sets in case (3).
That is, at these points every pair of the rider $\mathfrak{G}(\mathfrak{B}')$ of a later source set
satisfies the stratum $J$-member of $\mathfrak{G}(\mathfrak{B}_\tau)$.

For case (1) read at $\mathfrak{B}'$, let $\mathfrak{b}$ be a live arrest holding $y$ at
$\mathfrak{B}'$.
\begin{itemize}
\item If $k_b>k_a$, then $X_{\mathfrak{a}}\subset X_{\mathfrak{b}}$ by
\eqref{4.27:eq-arrest-timing-a}, and $\mathfrak{b}$ does not change the innermost piece.
\item If $k_b=k_a$, then $\mathfrak{b}=\mathfrak{a}$, because each class component takes one exit per
step.
\item If $k_b<k_a$, then $\mathfrak{b}$ is attached before every date of density $k_a$ and is live at
$\mathfrak{B}'$. Since liveness is an interval, it is live at $\mathfrak{B}_\tau$ too. At both dates
the innermost piece holding $y$ is the oldest such piece, by \eqref{4.27:eq-arrest-timing-a}.
\end{itemize}
In the first two alternatives $\mathfrak{b}$ does not displace the innermost live piece holding $y$
found at $\mathfrak{B}^t$ in case (1), namely $W_{\mathfrak{a}}$ or the older nested piece
$W_{\mathfrak{c}}$; in the third, the innermost piece holding $y$ is the oldest such piece at
$\mathfrak{B}_\tau$ and at $\mathfrak{B}'$. Hence, by the two alternatives of case (1),
\begin{equation*}
T^{\mathfrak{B}'}(y)=T^{\mathfrak{B}^t}(y)=T^{\mathfrak{B}_\tau}(y)
\end{equation*}
while $y$ is raw at $\mathfrak{B}'$.
\end{proof}

\begin{lemma}[Unheld raw points within one density.
Stated; the proof below is incomplete at the step marked $(\ast)$]\label{4.27:prf-unheld-raw-points}
Dates, their total order and the determining sets in force at them are those of
Lemma~\ref{4.27:lem-determining-set-order} (see \eqref{4.27:eq-determining-set-order-a}).

For determining sets $\mathfrak{B},\mathfrak{B}'$ in force at two dates:
\begin{itemize}
\item $\mathfrak{B}\preceq\mathfrak{B}'$ means that the date of $\mathfrak{B}$ is not later than that of
$\mathfrak{B}'$;
\item the density of $\mathfrak{B}$ is the density of its date;
\item a date in $\mathopen{]}\mathfrak{B},\mathfrak{B}'\mathclose{]}$ is a date later than that of
$\mathfrak{B}$ and not later than that of $\mathfrak{B}'$.
\end{itemize}

A point $y$ is \emph{raw} at $\mathfrak{B}$ if
$y\in\Omega_0(\mathfrak{B})\setminus\Omega_1(\mathfrak{B})$.

Let $\mathfrak{c}$ be an arrest, of density $k_{\mathfrak{c}}$, with component $X_{\mathfrak{c}}$ and piece
$W_{\mathfrak{c}}$. It \emph{holds} $y$ at a date $d$ if $y\in W_{\mathfrak{c}}$ and $\mathfrak{c}$ is live
at $d$, that is, attached at a date not later than $d$ and not dissolved at a date not later than $d$.

The rider $\mathfrak{G}(\mathfrak{B})$ has a stratum $J$-member
$\ell^{3}|\mathbb{J}(y)|<T^{\mathfrak{B}}(y)$ at its raw points. Its threshold is
$T^{\mathfrak{B}}(y)=\gamma_0^{(k_{\mathfrak{c}})}$ when $W_{\mathfrak{c}}$ is the innermost live piece
containing $y$ at the date of $\mathfrak{B}$. If there is none, $T^{\mathfrak{B}}(y)=\gamma_0^{(k)}$, where
$k$ is the density of $\mathfrak{B}$. Here
$\gamma_0^{(i)}=5\max_{j\le i}w_{*}(j)\bar\alpha_j$. These objects are defined in
Lemma~\ref{4.27:lem-weights-heredity}, whose proof is incomplete at one step; only its definitions are
used here.

Let $\mathfrak{B}$ be of density $k$. Let $y$ be raw at $\mathfrak{B}$ and held at the date of
$\mathfrak{B}$ by no live piece, so that $T^{\mathfrak{B}}(y)=\gamma_0^{(k)}$. Let
$\mathfrak{B}'\succeq\mathfrak{B}$ be such that $y$ is raw at $\mathfrak{B}'$.
\begin{enumerate}
\item[(1)] If $\mathfrak{B}'$ is of density $k$, every live holder of $y$ at $\mathfrak{B}'$ is an arrest
of density $k$ attached at a date in $\mathopen{]}\mathfrak{B},\mathfrak{B}'\mathclose{]}$.
\item[(2)] Suppose an arrest $\mathfrak{b}$ holds $y$ at some date of density $k$ later than that of
$\mathfrak{B}$. Then $k_{\mathfrak{b}}=k$. Moreover $\mathfrak{b}$ is the innermost live holder of $y$ at
every date from its attachment on at which $y$ is raw.
\item[(3)] Suppose $y$ is held by no live piece at any date of density $k$ from that of $\mathfrak{B}$ on,
and $\mathfrak{B}'$ is of density $k'>k$. Then every live holder of $y$ at $\mathfrak{B}'$ has density in
$]k,k']$. Further, $T^{\mathfrak{B}'}(y)=\gamma_0^{(k_{\mathfrak{b}})}$ for the oldest of them,
$\mathfrak{b}$, with $k<k_{\mathfrak{b}}\le k'$. If there is none, $T^{\mathfrak{B}'}(y)=\gamma_0^{(k')}$.
\end{enumerate}
The thresholds then satisfy
\begin{equation}\label{4.27:eq-unheld-raw-points-a}
\begin{aligned}
&\text{in (1):}\quad && T^{\mathfrak{B}'}(y)=\gamma_0^{(k)}=T^{\mathfrak{B}}(y),\\
&\text{in (2):}\quad && T^{\mathfrak{B}''}(y)=\gamma_0^{(k)}=T^{\mathfrak{B}}(y)
 \quad\text{for every }\mathfrak{B}''\succeq\mathfrak{B}\text{ at which }y\text{ is raw},\\
&\text{in (3):}\quad && T^{\mathfrak{B}'}(y)\ge\gamma_0^{(k+1)}\ge\gamma_0^{(k)}=T^{\mathfrak{B}}(y).
\end{aligned}
\end{equation}
In case (2) this holds whatever the density of $\mathfrak{B}''$.

In case (3), both inequalities in the last line of \eqref{4.27:eq-unheld-raw-points-a} are equalities if
and only if
$\max_{j\le i}w_{*}(j)\bar\alpha_j$ takes the same value at $i=k$ as at $i=k_{\mathfrak{b}}$ (resp.
$i=k'$). Where it is larger, consider a pair with $\mathbb{U}\equiv 1$ and $\mathbb{J}$ supported at $y$,
such that
\begin{equation*}
\gamma_0^{(k)}\le\ell^{3}|\mathbb{J}(y)|<T^{\mathfrak{B}'}(y).
\end{equation*}
Each such pair lies in $\mathfrak{G}(\mathfrak{B}')$ and not in $\mathfrak{G}(\mathfrak{B})$.
\end{lemma}

\begin{proof}
\emph{Holding is an interval.} The piece $W_{\mathfrak{c}}$ of an arrest is a fixed set. An arrest is live
from its attachment date up to its dissolution. That dissolution is the date of the first later completed
operation $\mathbf{R}$ on the component then containing $X_{\mathfrak{c}}$. So the dates at which
$\mathfrak{c}$ holds $y$ form an interval of the date order. This follows from the definition of arrests,
of their attachment and of their dissolution after the step, read with
Lemma~\ref{4.27:lem-determining-set-order}.

By the same lemma:
\begin{itemize}
\item the dates are totally ordered;
\item the dates of one density are consecutive, so densities do not decrease along $\preceq$;
\item an arrest of density $k_{\mathfrak{c}}$ is attached to a date of density $k_{\mathfrak{c}}$.
\end{itemize}

\emph{Proof of (1).} Let $\mathfrak{c}$ hold $y$ at $\mathfrak{B}'$. Suppose $\mathfrak{c}$ were attached
at a date not later than that of $\mathfrak{B}$. Its interval of liveness contains its attachment date and
the date of $\mathfrak{B}'$. It therefore contains the date of $\mathfrak{B}$, which lies between them. So
$\mathfrak{c}$ would hold $y$ at $\mathfrak{B}$, which is excluded.

Hence $\mathfrak{c}$ is attached at a date in $\mathopen{]}\mathfrak{B},\mathfrak{B}'\mathclose{]}$. All
these dates have density
$k$, since the dates of one density are consecutive. Thus $k_{\mathfrak{c}}=k$.

If such holders exist, the innermost one has density $k$. By the definition of $T^{\mathfrak{B}'}$, the
threshold is then $\gamma_0^{(k)}$. If there is none, $T^{\mathfrak{B}'}(y)=\gamma_0^{(k)}$ is the
threshold at unheld raw points at the density $k$ of $\mathfrak{B}'$. In both cases
$T^{\mathfrak{B}'}(y)=\gamma_0^{(k)}=T^{\mathfrak{B}}(y)$.

\emph{Proof of (2).} Suppose $\mathfrak{b}$ were attached at a date not later than that of $\mathfrak{B}$.
Then, as in (1), it would hold $y$ at $\mathfrak{B}$. So it is attached later. Its attachment date lies
between the date of $\mathfrak{B}$ and a date of density $k$ at which it holds $y$. It therefore has
density $k$, and $k_{\mathfrak{b}}=k$.

The component $X_{\mathfrak{b}}$ is the class component of step $k$ containing $y$. It is unique, because
the components of $Z_k$ are disjoint. It exits at most once: an arrested component is not changed by the
subsequent operations of the construction, and
\eqref{4.27:eq-arrest-timing-a} applies. Hence $\mathfrak{b}$ is the only arrest of density $k$ that ever
holds $y$.

Now let $\mathfrak{B}''\succeq\mathfrak{B}$ with $y$ raw at $\mathfrak{B}''$.

\emph{Before attachment.} If the date of $\mathfrak{B}''$ precedes the attachment of $\mathfrak{b}$, then
$\mathfrak{B}''$ has density $k$. Part (1), applied to $\mathfrak{B}''$, gives
$T^{\mathfrak{B}''}(y)=\gamma_0^{(k)}$.

\emph{After attachment: $\mathfrak{b}$ is live.} Otherwise $\mathfrak{b}$ is live at $\mathfrak{B}''$. Its
dissolution is the date $(k_e;L)$ of the first later completed operation $\mathbf{R}$ on the component then
containing $X_{\mathfrak{b}}\ni y$. At that date the completion move lifts that whole component, $y$ with
it, into $\Omega_1(\cdot)$. This is the convention on completed components adopted in
Lemma~\ref{4.27:prf-held-raw-points}. The point stays in $\Omega_1(\cdot)$ for good, because
$\Omega_1(\cdot)$ never shrinks (Lemma~\ref{4.27:prf-raw-stratum}). Since $y$ is raw at
$\mathfrak{B}''$, that date has not been reached.

\emph{After attachment: $\mathfrak{b}$ is innermost.} Let $\mathfrak{b}'$ be another live holder of $y$ at
$\mathfrak{B}''$.
\begin{itemize}
\item As in (1), $\mathfrak{b}'$ is attached later than $\mathfrak{B}$, so $k_{\mathfrak{b}'}\ge k$.
\item $k_{\mathfrak{b}'}=k$ would force $\mathfrak{b}'=\mathfrak{b}$. Hence $k_{\mathfrak{b}'}>k_{\mathfrak{b}}$.
\item Since $y$ is raw, $y\notin\Omega_1(\mathfrak{B}'')$. Ba{\l}aban's domain $\Omega_1$ of
\cite[(2.1)]{B-CMP119} satisfies $\Omega_1\subset\Omega_1(\mathfrak{B}'')$
(Lemma~\ref{4.27:prf-raw-stratum}), and $\Omega_{k_{\mathfrak{b}}},\Omega_{k_{\mathfrak{b}'}}\subset\Omega_1$
\cite[(2.1)]{B-CMP119}. Hence $y\in\Omega^{c}_{k_{\mathfrak{b}}}\cap\Omega^{c}_{k_{\mathfrak{b}'}}$.
\item As $y\in W_{\mathfrak{b}}\cap W_{\mathfrak{b}'}$, the timing relations
\eqref{4.27:eq-arrest-timing-a} give $X_{\mathfrak{b}}\subset X_{\mathfrak{b}'}$.
\end{itemize}
So $\mathfrak{b}$ is innermost, and
\[
T^{\mathfrak{B}''}(y)=\gamma_0^{(k_{\mathfrak{b}})}=\gamma_0^{(k)}=T^{\mathfrak{B}}(y).
\]

\emph{Proof of (3).} By the first step of (1), every live holder of $y$ at $\mathfrak{B}'$ is attached later
than $\mathfrak{B}$. An arrest attached at a date of density $k$ with $y$ in its piece would hold $y$ at that
date. By hypothesis no arrest does so, so no holder is attached at a date of density $k$. Its attachment date
is later than that of $\mathfrak{B}$ and not later than that of $\mathfrak{B}'$. Hence its density lies in
$]k,k']$.

As in (2), the timing relations \eqref{4.27:eq-arrest-timing-a} make the oldest holder $\mathfrak{b}$ the
innermost. So $T^{\mathfrak{B}'}(y)=\gamma_0^{(k_{\mathfrak{b}})}$ with $k<k_{\mathfrak{b}}\le k'$. If there
is no holder, $T^{\mathfrak{B}'}(y)=\gamma_0^{(k')}$ is the threshold at unheld raw points at the density of
$\mathfrak{B}'$.

In both cases the index is at least $k+1$. The map $i\mapsto\gamma_0^{(i)}$ is a running maximum times five,
hence non-decreasing, and this gives the third line of \eqref{4.27:eq-unheld-raw-points-a}. Equality holds
exactly when the running maximum has the same value at $k$ and at $k_{\mathfrak{b}}$ (resp. $k'$).

\emph{The separating pairs.} Assume the maximum has grown. Take $\mathbb{U}\equiv1$, that is $U\equiv1$ and
$A'\equiv0$ in $\mathbb{U}=e^{i\eta A'}U$. Let $\mathbb{J}$ be supported at $y$ with
$\gamma_0^{(k)}\le\ell^{3}|\mathbb{J}(y)|<T^{\mathfrak{B}'}(y)$. The pair obeys every clause of
$\mathfrak{G}(\mathfrak{B}')$:
\begin{itemize}
\item the $U$-clauses, the cube clauses and the stratum $U$-members hold trivially with the gauge
$u\equiv1$, since $A=A'=0$;
\item the direct clause on $\mathbb{J}$ at points of $\Omega_1(\mathfrak{B}')$ holds because $\mathbb{J}$
vanishes there;
\item the stratum $J$-member holds at $y$ by the choice of $\mathbb{J}(y)$, and at every other raw point
because $\mathbb{J}$ vanishes there.
\end{itemize}
No further condition is imposed on $\mathbb{J}$. The space of pairs $\mathfrak{K}_{109}(\mathfrak{B};\gamma)$
entering the second inclusion of clause (e3) of Lemma~\ref{4.27:lem-weights-heredity} admits
any $\mathfrak{g}^{\mathbb{C}}$-valued one-form $J$ with $\sup\ell^{3}|J|\le\gamma$. It asks no condition
$D^{*}J=0$ and no second-order condition.

The pair violates the $J$-member of $\mathfrak{G}(\mathfrak{B})$ at $y$, because
$\ell^{3}|\mathbb{J}(y)|\ge\gamma_0^{(k)}=T^{\mathfrak{B}}(y)$. So
$\mathfrak{G}(\mathfrak{B}')\not\subset\mathfrak{G}(\mathfrak{B})$.

The mechanisms that make thresholds hereditary elsewhere have nothing to act on at $y$:
\begin{itemize}
\item at index $0$ there is no clause of the form $w\,r_a$ whose threshold the monotonicity (e2) of
Lemma~\ref{4.27:lem-weights-heredity} (whose proof is incomplete at one step) could carry;
\item no level, weight or freeze of an arrest acts at $y$;
\item freezing the threshold with the structure of the innermost live piece does not apply, since no piece
holds $y$.
\end{itemize}

$(\ast)$ In case (3), wherever the maximum has grown, the $J$-member of the first inclusion of clause (e3)
of Lemma~\ref{4.27:lem-weights-heredity} (see \eqref{4.27:eq-weights-heredity-b}; that lemma's proof is
incomplete at one step) is not established at $y$. By the pairs just described it fails for these riders.

At points held by no live piece, the threshold $\gamma_0^{(k)}$ is kept as a membership condition of
$\mathfrak{G}(\mathfrak{B})$ at the density of $\mathfrak{B}$. With it the second inclusion
$\mathfrak{G}(\mathfrak{B})\subset\mathfrak{K}_{109}(\mathfrak{B};\gamma_0^{(k)})$ holds. It is not used as
a hereditary condition across a density increase.

Consider a stored term $\tau$ of an arrest, born at density $k_\tau$, with domain $X_\tau$. The affected
points are the raw points of $X_\tau$ with three properties:
\begin{itemize}
\item they lie outside every live piece;
\item they lie in components of $Z_{k_\tau}$ that are not class components of step $k_\tau$, because they
fail one of the two conditions defining class components; among these are the components containing
re-included second-class components of earlier steps;
\item they are read by the source of a later operation of density larger than $k_\tau$.
\end{itemize}
Those later operations are $O3$ (the $\mathbf{T}$-steps from $k_\tau\to k_\tau+1$ on), $O2$ (completions
at densities $k>k_\tau$), $O4$ and $O4'$ (stages at densities larger than $k_\tau$), and $O5$ (completions
at which a live arrest dissolves).

Call a \emph{dating} an assignment of the threshold of the stratum $J$-member at a raw point. Off $X_\tau$
the source one-form $\mathbb{J}$ may be cut to zero, whatever dating is used there, for three reasons: the
rider is existential, the value of a stored term reads
$X_\tau$ only, and $\mathfrak{K}_{109}$ admits any $J$. For a stored term of another part of the
construction whose estimate uses the inclusion of riders, the affected
points are the raw points of its domain that are held by no live piece at any date of its birth density and
are later met by a source of a larger density. At that later date such a point is held by an arrest of larger
density, or by none.

The gap at $(\ast)$ would be closed by any one of the following, none of which is established here:
\begin{enumerate}
\item[(a)] No stored value and no size bound of the kind that reads the product $K_\Delta\gamma_0$
($K_\Delta$ the constant of the smallness condition $K_\Delta\gamma_0\le\tfrac12$) reads $\mathbb{J}$ at such
points. Then $\mathbb{J}$ may be cut to zero there, and the stratum $J$-member leaves the rider at those
points.
\item[(b)] A density-independent dating at index $0$ is used, such as
$\gamma^{\top}:=\bar\gamma_0(\mathsf{t}_\gamma)$, which would serve both inclusions. It must then be
checked that the size bounds involving $K_\Delta\gamma_0$, read at a single value of $\gamma_0$, do not
re-create the lower bound on the coupling, equivalently an upper bound on $K$, that the proof of
clause (d) of Lemma~\ref{4.27:lem-weights-heredity} (incomplete at one step) was arranged to avoid.
\item[(c)] The sources of the operations reading stored terms are restricted so as to avoid such points.
\end{enumerate}

\emph{Consequences.} The stratum $J$-member of the first inclusion of clause (e3) is hereditary at every
raw point, held or not, in two situations:
\begin{itemize}
\item between two dates of one density, by (1);
\item from the first holding of the point within that density on, by (2) together with the first part of
\eqref{4.27:eq-held-raw-points-a}.
\end{itemize}

In the analysis of the stored terms of an arrest that accompanies the arrest theorem, this has two
effects. It settles the raw points of
$X_\tau$ in components of $Z_{k_\tau}$ that are not class components, at the dates of density $k_\tau$:
no holder of density $k_\tau$ ever exists there, so (2) is void
and (1) applies. It also re-proves, for points in class components of step $k_\tau$ other than that of the
arrest, the part of that analysis which concerns the dates of density $k_\tau$.

What remains is case (3) of the present lemma: a density increase, at points unheld throughout the earlier
density. For the stored terms of an arrest these are the raw points of $X_\tau$ outside every live piece, in
components of $Z_{k_\tau}$ that are not class components, read past the dates of density $k_\tau$. For other
stored terms it
consists of the raw points of the domain unheld at every date of the birth density and read by a source of
a later density.

The argument uses only the definition of arrests and their liveness, the timing relations
\eqref{4.27:eq-arrest-timing-a}, the uniqueness of the exit of a component, the definition of the dated
threshold, Lemma~\ref{4.27:lem-determining-set-order}, the fact that $\Omega_1(\cdot)$ never shrinks, the
convention on completed components of Lemma~\ref{4.27:prf-held-raw-points}, and the monotonicity of
$\gamma_0^{(\cdot)}$. No estimate and no
constant enters; $k$, $k'$, $k_{\mathfrak{b}}$ and $k_e$ are densities of dates.
\end{proof}

\begin{lemma}[The one-point current probe]\label{4.27:lem-current-probe}
Let $d$ be a date, of density $k(d)$, in the sense of Lemma~\ref{4.27:lem-weights-heredity}, and let
$\mathfrak{B}_d$ be its determining set. Let $\mathfrak{G}(\mathfrak{B}_d)$ be the space of pairs defined in
the statement of Lemma~\ref{4.27:lem-weights-heredity} (stated; proof incomplete at the step marked
$(\ast)$). Only definitions made in its statement are used here (dates, determining sets, $\mathfrak{G}$, the
dating $T^{\mathfrak{B}}$ and $\gamma_0^{(\cdot)}$), none of its conclusions.

For a date $d'$ the set $(\Omega_0\setminus\Omega_1)(\mathfrak{B}_{d'})$ is called the \emph{stratum} at $d'$,
and its points are called \emph{raw} at $d'$; a raw point is one raw at $d$. The condition of
$\mathfrak{G}(\mathfrak{B}_{d'})$ on $\mathbb{J}$ at the raw points is called the stratum $\mathbb{J}$-member.
For a raw point $y$ and a number $T>0$, the stratum $\mathbb{J}$-member at $y$ \emph{dated by} $T$ is the
condition $\ell^{3}|\mathbb{J}(y)|<T$, all other clauses of $\mathfrak{G}(\mathfrak{B}_d)$ being unchanged;
in $\mathfrak{G}(\mathfrak{B}_d)$ itself $T=T^{\mathfrak{B}_d}(y)$.

Let $y$ be a raw point, let $\mathsf{e}\in\mathfrak{g}$ with
$|\mathsf{e}|=1$, and let $c\ge0$. Let $P_y(c)$ be
the following pair. Its field is $\mathbb{U}\equiv1$ on $\Omega_0(\mathfrak{B}_d)$, that is, $U\equiv1$ and
$A'\equiv0$ in the decomposition $\mathbb{U}=e^{i\eta A'}U$. Its source is supported at $y$, with
$\mathbb{J}(y)$ read as in the stratum $\mathbb{J}$-member of $\mathfrak{G}$:
\begin{equation}\label{4.27:eq-current-probe-1}
\mathbb{J}(y):=c\,\mathsf{e},\qquad \mathbb{J}(y'):=0\quad\text{for } y'\neq y.
\end{equation}
Thus $\mathbb{J}$ is the $\mathfrak{g}^{\mathbb{C}}$-valued $1$-form equal to $c\,\mathsf{e}$ at $y$
and to $0$ elsewhere, and $P_y(c)$ is a pair of the form $(e^{i\eta A'}U,\mathbb{J})$ on which
$\mathfrak{G}$ is defined.
Then the following hold.
\begin{enumerate}
\item[(1)] $P_y(c)$ obeys every clause of $\mathfrak{G}(\mathfrak{B}_d)$ other than the stratum
$\mathbb{J}$-member at $y$. This holds for every structure of $\mathfrak{B}_d$ (levels, weights and cube
families). At $y$ one has $\ell^{3}|\mathbb{J}(y)|=\ell^{3}c$.
\item[(2)] For a number $T>0$, $P_y(c)$ obeys the stratum $\mathbb{J}$-member at $y$ dated by $T$ if and
only if $\ell^{3}c<T$. Let $\gamma>0$, and let $\mathfrak{K}_{109}(\mathfrak{B}_d;\gamma)$ be the space of
pairs of the type of \cite{B-CMP109} with current radius $\gamma$; its current clause is
$\sup\ell^{3}|\mathbb{J}|\le\gamma$, and its $\mathbb{J}$-member at $y$ is the condition
$\ell^{3}|\mathbb{J}(y)|\le\gamma$. Then $P_y(c)$ obeys the $\mathbb{J}$-member at $y$ of
$\mathfrak{K}_{109}(\mathfrak{B}_d;\gamma)$ if and only if $\ell^{3}c\le\gamma$.
\item[(3)] Let $T,T'>0$. Every pair that obeys the clauses of $\mathfrak{G}(\mathfrak{B}_d)$ other than the
stratum $\mathbb{J}$-member at $y$, and obeys that member dated by $T'$, obeys it dated by $T$ if and only
if $T'\le T$. Moreover, for
$\gamma>0$, every pair of $\mathfrak{G}(\mathfrak{B}_d)$, its member at $y$ dated by $T$, obeys the
$\mathbb{J}$-member
$\ell^{3}|\mathbb{J}(y)|\le\gamma$ of $\mathfrak{K}_{109}(\mathfrak{B}_d;\gamma)$ at $y$ if and only if
$T\le\gamma$; in particular, if $T>\gamma$, then $\mathfrak{G}(\mathfrak{B}_d)$, its member at $y$ dated by
$T$, is not contained in $\mathfrak{K}_{109}(\mathfrak{B}_d;\gamma)$.
\end{enumerate}
\end{lemma}

\begin{proof}
\emph{(1).} We go through the clauses in the definition of $\mathfrak{G}(\mathfrak{B}_d)$ one by one. For $U$
and $A'$ we use the decomposition $\mathbb{U}=e^{i\eta A'}U$ with $U\equiv1$, $A'\equiv0$.

\emph{Direct entries.} At every point of $\Omega_1(\mathfrak{B}_d)$ the definition bounds the direct entries
$Q_a$, $a\in\{1,3,5,6,7\}$, by $5\theta^{\mathfrak{B}_d}_a$. For $P_y(c)$ we have
\begin{equation*}
Q_1=|\partial\mathbb{U}-1|\equiv0,\quad Q_5=|\partial U-1|\equiv0,\quad
Q_6=|A'|\equiv0,\quad Q_7=|\nabla^{\eta}_{U}A'|\equiv0 .
\end{equation*}
Also $Q_3=|\mathbb{J}|=0$ at every point of $\Omega_1(\mathfrak{B}_d)$, because
$y\notin\Omega_1(\mathfrak{B}_d)$ and $\mathbb{J}$ vanishes off $y$. Each of these entries therefore lies
below its threshold $5\theta^{\mathfrak{B}_d}_a$, which is positive: at a point $y''$ the threshold
\begin{equation*}
5\theta^{\mathfrak{B}_d}_a(y'')=5w_{\mathfrak{B}_d}(y'')\,r_a(\ell_{\mathfrak{B}_d}(y''))
\end{equation*}
is the product of $5$, of the
weight, a non-negative integer power of $\mathsf{L}_0^{2}$ and hence at least $1$, and of $r_a$, a positive
multiple of a radius. The entries $a=2,4$ do not occur in the definition.

\emph{Cube clauses.} On every cube $\square$ of $\mathcal{Q}_{\mathfrak{B}_d}$, of family
$(\lambda_\square,w_\square)$, take the gauge $u\equiv1$. Then $U^{u}=1=e^{i\eta\cdot0}$, so $A\equiv0$. Hence
$L^{\lambda_\square}\eta|A|$ and $(L^{\lambda_\square}\eta)^2|\nabla^{\eta}A|$ vanish. They therefore lie
below the threshold $5w_\square BCM\alpha_{0,\lambda_\square}$ of the cube clause in the definition of
$\mathfrak{G}$ (with its constants $B$, $C$).

\emph{Index-$0$ $U$-members.} The index-$0$ instance of \cite[(3.35)]{B-CMP99b} on the index-$0$ cubes holds
with $u\equiv1$ and $A\equiv0$. The index-$0$ instance of \cite[(3.37)]{B-CMP99b} holds with $A'\equiv0$.
In both, the quantities bounded vanish, and zero lies inside their radii, which are positive and do not
depend on $\gamma$.

\emph{Stratum $\mathbb{J}$-member away from $y$.} At every raw point $y'\neq y$ we have
$\ell^{3}|\mathbb{J}(y')|=0<T^{\mathfrak{B}_d}(y')$; the datum $T^{\mathfrak{B}_d}(y')$ is a value of
$\gamma_0^{(\cdot)}$, defined in the statement of Lemma~\ref{4.27:lem-weights-heredity} as
$5\max_{j\le\cdot}w_{*}(j)\bar\alpha_j$, and positive because $w_{*}(j)\ge1$ and $\bar\alpha_j>0$.

\emph{At $y$.} Since $|\mathsf{e}|=1$, we have $\ell^{3}|\mathbb{J}(y)|=\ell^{3}c$.

\emph{No further conditions on $\mathbb{J}$.} The definition of $\mathfrak{G}$ contains no clause
$D^{*}\mathbb{J}=0$ and no clause of second order on $\mathbb{J}$. Likewise
$\mathfrak{K}_{109}(\mathfrak{B}_d;\gamma)$ admits as current any $\mathfrak{g}^{\mathbb{C}}$-valued $1$-form
$J$ with $\sup\ell^{3}|J|\le\gamma$, and asks no condition $D^{*}J=0$. This proves (1).

\emph{(2).} By (1), $\ell^{3}|\mathbb{J}(y)|=\ell^{3}c$. The member at $y$ dated by $T$ is the condition
$\ell^{3}|\mathbb{J}(y)|<T$, so it holds if and only if $\ell^{3}c<T$. The $\mathbb{J}$-member of
$\mathfrak{K}_{109}(\mathfrak{B}_d;\gamma)$ at $y$ is the condition $\ell^{3}|\mathbb{J}(y)|\le\gamma$, so it
holds if and only if $\ell^{3}c\le\gamma$.

\emph{(3), first assertion.} Suppose first that $T'>T$, and put $c:=T\,\ell^{-3}$. Then
$T\le\ell^{3}c<T'$. By (1), $P_y(c)$ obeys every clause of $\mathfrak{G}(\mathfrak{B}_d)$ other than the
member at $y$; by (2), it obeys the member at $y$ dated by $T'$ and does not obey it dated by $T$. So not
every pair obeying the member dated by $T'$, all other clauses being unchanged, obeys it dated by $T$.

Suppose now that $T'\le T$. Then every pair obeying the member at $y$ dated by $T'$ satisfies
\begin{equation*}
\ell^{3}|\mathbb{J}(y)|<T'\le T,
\end{equation*}
and so obeys the member at $y$ dated by $T$.

\emph{(3), second assertion.} If $T\le\gamma$, every pair of $\mathfrak{G}(\mathfrak{B}_d)$ satisfies
$\ell^{3}|\mathbb{J}(y)|<T\le\gamma$. Suppose that $T>\gamma$, and choose $c\ge0$ with
$\gamma<\ell^{3}c<T$. By (1) and (2), $P_y(c)$ lies in $\mathfrak{G}(\mathfrak{B}_d)$ with its member at $y$
dated by $T$. By (2), it violates the $\mathbb{J}$-member of $\mathfrak{K}_{109}(\mathfrak{B}_d;\gamma)$ at
$y$. Hence not every pair of $\mathfrak{G}(\mathfrak{B}_d)$ obeys that member, and
$\mathfrak{G}(\mathfrak{B}_d)$ is not contained in $\mathfrak{K}_{109}(\mathfrak{B}_d;\gamma)$.

Here $\mathsf{e}$ is any unit vector of $\mathfrak{g}\subset\mathfrak{g}^{\mathbb{C}}$ and $c\ge0$ is free;
no background configuration enters, no commutativity of the gauge group is used, and no constant of the
construction reads the probe.
\end{proof}

\begin{proposition}[Hereditary current thresholds are the non-increasing ones]
\label{4.27:prop-nonincreasing-thresholds}
Dates $d$, their determining sets $\mathfrak{B}_d$, their densities $k(d)$ and the total
order $\preceq$ of dates are those of Lemma~\ref{4.27:lem-determining-set-order}, as are
arrests $\mathfrak{c}$, their pieces $W_{\mathfrak{c}}$ and the dates at which they are
attached and dissolved.
The sets $\Omega_0(\mathfrak{B}_d)\supseteq\Omega_1(\mathfrak{B}_d)$ are the domain of the
date $d$ and its points of level at least one. A point $y$ is \emph{raw} at $d$ if
$y\in(\Omega_0\setminus\Omega_1)(\mathfrak{B}_d)$; this set of raw points is the
\emph{stratum} of $d$. A live piece \emph{holds} $y$ at $d$ if
$y\in W_{\mathfrak{c}}$ for an arrest $\mathfrak{c}$ that is attached at a date $\preceq d$
and is not dissolved at a date $\preceq d$. The \emph{density} $k_{\mathfrak{c}}$ of an
arrest $\mathfrak{c}$ is the density of the date to which it is attached.

The following objects are those of Lemma~\ref{4.27:lem-weights-heredity}, whose proof is
incomplete at one step:
\begin{itemize}
\item the rider $\mathfrak{G}(\mathfrak{B})$ of \eqref{4.27:eq-weights-heredity-b};
\item its stratum current member $\ell^{3}|\mathbb{J}(y)|<T^{\mathfrak{B}}(y)$ at raw
points, with $\ell=\eta$ at index~$0$;
\item the radius $\gamma_0^{(k)}$ of \eqref{4.27:eq-weights-heredity-a}.
\end{itemize}
The current member of $\mathfrak{K}_{109}(\mathfrak{B};\gamma)$ at a raw point $y$ is
$\ell^{3}|\mathbb{J}(y)|\le\gamma$.

Let $y$ be a point. Put
\begin{equation*}
D_0(y):=\{d:\ y\in\Omega_0(\mathfrak{B}_d)\},\qquad
D_1(y):=\{d:\ y\in\Omega_1(\mathfrak{B}_d)\}\subseteq D_0(y).
\end{equation*}
These are final segments of the date order. Hence
\begin{equation*}
I_y:=D_0(y)\setminus D_1(y)=\{d:\ y\in(\Omega_0\setminus\Omega_1)(\mathfrak{B}_d)\}
\end{equation*}
is an interval of that order, possibly empty. For $I_y\neq\emptyset$ put
$d_{\circ}(y):=\min I_y$, which equals $\min D_0(y)$, and
$k_{\circ}(y):=k(d_{\circ}(y))$.

A \emph{dating} at $y$ is any assignment $d\mapsto T(y,d)>0$, $d\in I_y$, of the
threshold of the stratum current member of $\mathfrak{G}(\mathfrak{B}_d)$ at $y$. Every
other clause of $\mathfrak{G}(\mathfrak{B}_d)$ is kept as in
\eqref{4.27:eq-weights-heredity-b}. Examples of datings are:
\begin{itemize}
\item the frozen datum $T^{\mathfrak{B}_d}(y)$, that is, $\gamma_0^{(k_{\mathfrak{c}})}$, where
$\mathfrak{c}$ is the innermost live piece holding $y$ at $d$, and $\gamma_0^{(k(d))}$ if
there is none;
\item the value $\gamma_0^{(k(d))}$ at every date;
\item a radius independent of the density, such as
$\gamma^{\top}:=\bar\gamma_0(\mathsf{t}_\gamma)$ of \eqref{4.27:eq-weights-heredity-a};
\item the dating $T^{\flat}$ of the second remark below.
\end{itemize}
A \emph{completion of the frozen datum} is a dating with
$T(y,d)=\gamma_0^{(k_{\mathfrak{c}})}$, $\mathfrak{c}$ the innermost live piece holding $y$
at $d$, at every $d\in I_y$ at which some live piece holds $y$, and arbitrary at the other
dates of $I_y$. The frozen datum $T^{\mathfrak{B}_d}(y)$ is the completion taking the value
$\gamma_0^{(k(d))}$ at those other dates.
\begin{enumerate}
\item[(i)] \emph{Characterisation.} Consider the following two assertions, each required
for all $d\preceq d'$ in $I_y$:
\begin{itemize}
\item every pair of $\mathfrak{G}(\mathfrak{B}_{d'})$, restricted to
$\Omega_0(\mathfrak{B}_d)$, obeys the current member of $\mathfrak{G}(\mathfrak{B}_d)$ at $y$;
\item every pair of $\mathfrak{G}(\mathfrak{B}_d)$ obeys the current member of
$\mathfrak{K}_{109}(\mathfrak{B}_d;\gamma_0^{(k(d))})$ at $y$.
\end{itemize}
The first is the current part at $y$ of the \emph{first inclusion}
$\mathfrak{G}(\mathfrak{B}_{d'})\subset\mathfrak{G}(\mathfrak{B}_d)$, pairs on
$\Omega_0(\mathfrak{B}_{d'})$ being restricted to
$\Omega_0(\mathfrak{B}_d)\subseteq\Omega_0(\mathfrak{B}_{d'})$. The second is the
current part at $y$ of the \emph{second inclusion}
$\mathfrak{G}(\mathfrak{B}_d)\subset\mathfrak{K}_{109}(\mathfrak{B}_d;\gamma_0^{(k(d))})$.
Both hold as sets if and only if $T(y,\cdot)$ is non-increasing along $I_y$ and
$T(y,d)\le\gamma_0^{(k(d))}$ for every $d\in I_y$. Hence every such dating satisfies
\begin{equation*}
T(y,d)\le T(y,d_{\circ}(y))\le\gamma_0^{(k_{\circ}(y))}\qquad (d\in I_y).
\end{equation*}
The constant dating $T_{\circ}(y,d):=\gamma_0^{(k_{\circ}(y))}$, $d\in I_y$, meets both
conditions, since $\gamma_0^{(k_{\circ}(y))}\le\gamma_0^{(k(d))}$ for every $d\in I_y$, and
it is the pointwise-maximal dating meeting them.
\item[(ii)] \emph{Impossibility under the frozen datum and the second inclusion.} Let
$d\prec d'$ in $I_y$ and put $k:=k(d)$. Assume that no live piece holds $y$ at any date of
density $k$ from $d$ on, and that some live piece holds $y$ at $d'$. Then:
\begin{itemize}
\item $k(d')>k$;
\item the live pieces holding $y$ at $d'$ are attached after $d$, at dates of density in
$(k,k(d')]$;
\item the oldest of them, $\mathfrak{b}$, of density $k_{\mathfrak{b}}\in(k,k(d')]$, is the
innermost.
\end{itemize}
Let the dating keep the frozen datum at the held date, $T(y,d')=\gamma_0^{(k_{\mathfrak{b}})}$,
and the second inclusion at $d$, $T(y,d)\le\gamma_0^{(k)}$. This allows the value
$\gamma_0^{(k)}$ or any smaller one. Suppose $\gamma_0^{(k_{\mathfrak{b}})}>\gamma_0^{(k)}$.
Then, for every such dating and every $c\ge0$ with
\begin{equation*}
T(y,d)\le\ell^{3}c<\gamma_0^{(k_{\mathfrak{b}})},
\end{equation*}
the probe pair $P_y(c)$ of Lemma~\ref{4.27:lem-current-probe}, built on
$\Omega_0(\mathfrak{B}_{d'})$, lies in
$\mathfrak{G}(\mathfrak{B}_{d'})\setminus\mathfrak{G}(\mathfrak{B}_d)$, its restriction to
$\Omega_0(\mathfrak{B}_d)$ being the probe built at $d$. In particular the
current part at $y$ of the first inclusion fails.
\end{enumerate}
Suppose finally that there is a system of dates, determining sets and arrests obeying the
definitions and standing hypotheses of this chapter in which some point $y$ and dates
$d\prec d'$ in $I_y$ satisfy the hypotheses of~(ii) with
$\gamma_0^{(k_{\mathfrak{b}})}>\gamma_0^{(k)}$. Then, in that system, the
following hold.
\begin{itemize}
\item No re-assignment of the current threshold at unheld raw points alone makes the
current part of the first inclusion hold for every pair of dates, provided the frozen
datum at held raw points and the second inclusion are kept.
\item The separating probe of~(ii) exists for every completion of the frozen datum at
unheld raw points.
\item Consequently the two inclusions of \eqref{4.27:eq-weights-heredity-b} together, the
first as sets on all of $\Omega_0(\mathfrak{B}_d)$ for every $d$ and the second kept at
every date, cannot be obtained by completing the frozen datum at unheld raw points.
\end{itemize}
\end{proposition}

\begin{proof}
\emph{Preliminaries.} We collect four facts used below.

First, by the heredity statement on the raw stratum, Lemma~\ref{4.27:prf-raw-stratum}, the sets
$\Omega_0(\cdot)$ and $\Omega_1(\cdot)$ never shrink along the date order. That is,
$d\preceq d'$ implies $\Omega_i(\mathfrak{B}_d)\subseteq\Omega_i(\mathfrak{B}_{d'})$ for
$i=0,1$. Hence, if $d\in D_i(y)$ and $d\preceq d'$, then $d'\in D_i(y)$, so $D_0(y)$ and
$D_1(y)$ are final segments. The inclusion $D_1(y)\subseteq D_0(y)$ holds because
$\Omega_1(\mathfrak{B}_d)\subseteq\Omega_0(\mathfrak{B}_d)$. By
Lemma~\ref{4.27:lem-determining-set-order} the dates form a finite total order.

The complement of the final segment $D_1(y)$ is an initial segment. So
$I_y=D_0(y)\setminus D_1(y)$, the intersection of the final segment $D_0(y)$ with the
initial segment complementary to $D_1(y)$, is an interval.

Suppose $I_y\neq\emptyset$ and let $d_0:=\min D_0(y)$. If $d_0\in D_1(y)$, then the final
segment $D_1(y)$ contains every date $\succeq d_0$, hence all of $D_0(y)$, and $I_y$ would
be empty. Thus $d_0\in I_y$, and, since $I_y\subseteq D_0(y)$, $d_0=\min I_y=d_{\circ}(y)$.
In words: a point lying in $\Omega_1$ at its first date in $\Omega_0$ is never raw.

Second, by the order of dates \eqref{4.27:eq-determining-set-order-a}, the dates of one
density are consecutive, and the last date of density $k$ precedes $(k+1;\top)$. Hence
$k(\cdot)$ is non-decreasing along $\preceq$.

Third, $\gamma_0^{(k)}$ is defined in \eqref{4.27:eq-weights-heredity-a} as five times a
maximum over $j\le k$ of $w_{*}(j)\bar\alpha_j$. A maximum over a larger index set is not
smaller, so $k\mapsto\gamma_0^{(k)}$ is non-decreasing.

Fourth, we use the one-point current probe (Lemma~\ref{4.27:lem-current-probe}) in the
following form. Let $d$ be a date at which $y$ is raw, let $c\ge0$, and let $X$ be a unit
vector of $\mathfrak{g}$. The pair $P_y(c)$ is given by $\mathbb{U}\equiv1$ on
$\Omega_0(\mathfrak{B}_d)$, $\mathbb{J}(y)=cX$ on the bond $y$, and $\mathbb{J}=0$ on every
other bond. It has three properties.
\begin{itemize}
\item It obeys every clause of $\mathfrak{G}(\mathfrak{B}_d)$ other than the stratum
current member at $y$, whatever the structure of $\mathfrak{B}_d$. This includes the
entry $Q_3=|\mathbb{J}|$, which vanishes on $\Omega_1(\mathfrak{B}_d)$ because
$y\notin\Omega_1(\mathfrak{B}_d)$.
\item It obeys the member at $y$ dated by a number $T>0$ if and only if $\ell^{3}c<T$.
\item It obeys the current member of $\mathfrak{K}_{109}(\mathfrak{B}_d;\gamma)$ at $y$ if
and only if $\ell^{3}c\le\gamma$.
\end{itemize}

\emph{Proof of~(i): sufficiency.} Let $T(y,\cdot)$ be non-increasing along $I_y$ and let
$d\preceq d'$ in $I_y$. Since $y$ is raw at $d'$, a pair of $\mathfrak{G}(\mathfrak{B}_{d'})$
satisfies
\begin{equation*}
\ell^{3}|\mathbb{J}(y)|<T(y,d')\le T(y,d),
\end{equation*}
so it obeys the current member of $\mathfrak{G}(\mathfrak{B}_d)$ at $y$.

Now let $T(y,d)\le\gamma_0^{(k(d))}$. A pair of $\mathfrak{G}(\mathfrak{B}_d)$ satisfies
\begin{equation*}
\ell^{3}|\mathbb{J}(y)|<T(y,d)\le\gamma_0^{(k(d))},
\end{equation*}
so it obeys the current member of $\mathfrak{K}_{109}(\mathfrak{B}_d;\gamma_0^{(k(d))})$
at $y$. The other clauses of the second inclusion do not involve the dating. They are
those placed inside $\mathfrak{K}_{109}(\mathfrak{B}_d;\gamma_0^{(k(d))})$ in
Lemma~\ref{4.27:lem-weights-heredity}, whose proof is incomplete at one step.

\emph{Proof of~(i): necessity.} Suppose $T(y,d')>T(y,d)$ for some $d\preceq d'$ in $I_y$.
Choose $c\ge0$ with $T(y,d)\le\ell^{3}c<T(y,d')$, which is possible since the interval is
non-empty. By the fourth preliminary at the date $d'$ (where $y$ is raw), $P_y(c)$, built on
$\Omega_0(\mathfrak{B}_{d'})$, obeys every clause of $\mathfrak{G}(\mathfrak{B}_{d'})$, the
member at $y$ included. Its restriction to $\Omega_0(\mathfrak{B}_d)$ is the probe built at
$d$; by the same fact at $d$, it violates the current member of $\mathfrak{G}(\mathfrak{B}_d)$
at $y$, since $\ell^{3}c\ge T(y,d)$. So the first assertion fails.

Suppose next $T(y,d)>\gamma_0^{(k(d))}$ for some $d\in I_y$. Choose $c$ with
$\gamma_0^{(k(d))}<\ell^{3}c<T(y,d)$. Then $P_y(c)$ lies in $\mathfrak{G}(\mathfrak{B}_d)$
and violates the current member of $\mathfrak{K}_{109}(\mathfrak{B}_d;\gamma_0^{(k(d))})$
at $y$. So the second assertion fails.

\emph{Proof of~(i): the consequences.} Let $T$ be a non-increasing dating with
$T(y,d)\le\gamma_0^{(k(d))}$ on $I_y$. Since $d_{\circ}(y)=\min I_y$, we get
$T(y,d)\le T(y,d_{\circ}(y))$ for every $d\in I_y$. The second condition at
$d_{\circ}(y)$ gives $T(y,d_{\circ}(y))\le\gamma_0^{(k_{\circ}(y))}$.

The constant dating $T_{\circ}$ is non-increasing. For $d\in I_y$ we have
$d_{\circ}(y)\preceq d$, hence $k(d)\ge k_{\circ}(y)$ by the second preliminary, and
$\gamma_0^{(k_{\circ}(y))}\le\gamma_0^{(k(d))}$ by the third. So $T_{\circ}$ meets both
conditions. By the bound just proved it majorises every dating meeting them.

\emph{Proof of~(ii): the structure at $d'$.} We use the statement on unheld raw points
within one density, Lemma~\ref{4.27:prf-unheld-raw-points}, in the form displayed in
\eqref{4.27:eq-unheld-raw-points-a}; its proof is incomplete at one step. Its first part
gives: every live holder of $y$ at a date of density $k$ after $d$ is an arrest of
density $k$. Were $d'$ of density $k$, its live holders of $y$ would therefore be arrests
of density $k$ holding $y$ at a date of density $k$ after $d$, which the hypothesis
excludes, while $y$ is held at $d'$. Hence $d'$ is not of density $k$. Since $d\prec d'$,
the second preliminary gives $k(d')\ge k$, hence $k(d')>k$.

The third part of the same statement then gives three facts. Every live holder of $y$ at
$d'$ is attached after $d$ and, by hypothesis, at no date of density $k$, so its density
lies in $(k,k(d')]$. The oldest such holder $\mathfrak{b}$ is the innermost. The frozen
datum at $d'$ is therefore $\gamma_0^{(k_{\mathfrak{b}})}$ with $k<k_{\mathfrak{b}}\le k(d')$.

\emph{Proof of~(ii): the separation.} Let $T(y,d')=\gamma_0^{(k_{\mathfrak{b}})}$,
$T(y,d)\le\gamma_0^{(k)}$ and $\gamma_0^{(k_{\mathfrak{b}})}>\gamma_0^{(k)}$. Then
\begin{equation*}
T(y,d')=\gamma_0^{(k_{\mathfrak{b}})}>\gamma_0^{(k)}\ge T(y,d),
\end{equation*}
and $T(y,\cdot)$ is not non-increasing at the pair $(d,d')$. This is the necessity part of
(i) at that pair.

Explicitly, let $c\ge0$ with $T(y,d)\le\ell^{3}c<\gamma_0^{(k_{\mathfrak{b}})}$. The point $y$
is raw at $d'$, so the member at $y$ in $\mathfrak{G}(\mathfrak{B}_{d'})$ is dated
$T(y,d')=\gamma_0^{(k_{\mathfrak{b}})}>\ell^{3}c$. Moreover $Q_3$ vanishes on
$\Omega_1(\mathfrak{B}_{d'})$, which does not contain $y$. By the fourth preliminary,
$P_y(c)$, built on $\Omega_0(\mathfrak{B}_{d'})$, obeys every clause of
$\mathfrak{G}(\mathfrak{B}_{d'})$. Its restriction to $\Omega_0(\mathfrak{B}_d)$, the probe
built at $d$, violates the current member of $\mathfrak{G}(\mathfrak{B}_d)$ at $y$, since
$\ell^{3}c\ge T(y,d)$. Thus
$P_y(c)\in\mathfrak{G}(\mathfrak{B}_{d'})\setminus\mathfrak{G}(\mathfrak{B}_d)$.

\emph{Proof of the final assertion.} Take the system of the final assertion, with $y$ and
$d\prec d'$ as in~(ii) and $\gamma_0^{(k_{\mathfrak{b}})}>\gamma_0^{(k)}$. Consider any
completion of the frozen datum that keeps the second inclusion. Its value
at the held date $d'$ is $\gamma_0^{(k_{\mathfrak{b}})}$. Its value at the unheld date $d$
is some number $\le\gamma_0^{(k)}$, chosen arbitrarily otherwise. By~(ii), for every such
choice the probe $P_y(c)$ separates
$\mathfrak{G}(\mathfrak{B}_{d'})$ from $\mathfrak{G}(\mathfrak{B}_d)$ at $y$. Hence no such
re-assignment makes the current part of the first inclusion hold at the pair $(d,d')$,
and a fortiori not for every pair of dates.
\end{proof}

\begin{remark}
The impossibility in~(ii) uses the second inclusion through one constraint only. This is
the bound $T(y,d)\le\gamma_0^{(k(d))}$ that it imposes at a raw point; see the one-point
current probe (Lemma~\ref{4.27:lem-current-probe}). The value $\gamma_0^{(k)}$ assigned at
unheld raw points in \eqref{4.27:eq-weights-heredity-b} is the extreme instance of this
constraint.

At a raw point that no live piece holds before it is lifted into $\Omega_1$, the constant
dating $T_{\circ}$ serves at all its unheld dates; but which points these are is not
determined at the earlier date, and the obstruction of~(ii) sits exactly at the frozen
value at a held date.

Configurations as in~(ii) are not excluded by the definitions. A raw point satisfies
\begin{equation*}
y\in\Gamma_0\subseteq Z_1\subseteq Z_{k_{\mathfrak{b}}},\qquad \Gamma_0=\Omega_1^{c},
\end{equation*}
by \cite[(2.1)--(2.3)]{B-CMP119}, so it lies in a component of $Z_{k_{\mathfrak{b}}}$. The
definition of an arrest allows this component to be a class component $X_{\mathfrak{b}}$ of
step $k_{\mathfrak{b}}$ (a component of $Z_{k_{\mathfrak{b}}}$ on which the stages of step
$k_{\mathfrak{b}}$ run) that is arrested, whereupon
$y\in\Gamma_0\cap X_{\mathfrak{b}}\subseteq W_{\mathfrak{b}}$, since
$\Omega_1^{c}\subseteq\Omega^{c}_{k_{\mathfrak{b}}}$; and the definition
of $\gamma_0^{(\cdot)}$ bounds no growth of $\max_{j\le\cdot}w_{*}(j)\bar\alpha_j$. Neither
the definitions nor the standing hypotheses of this chapter exclude such a configuration,
and none of them produces one. It is not asserted that such a transition with
$\gamma_0^{(k_{\mathfrak{b}})}>\gamma_0^{(k)}$ occurs in a genuine run; what is proved is the
implication stated in the last assertion of the proposition. The standing hypotheses of this
chapter bound neither $K$ nor the growth of
$\max_{j\le\cdot}w_{*}(j)\bar\alpha_j$.
They contain no cap in the sense of Definition~\ref{2.3:def-cap} and no lower bound on the
coupling.
\end{remark}

\begin{remark}
\emph{Alternatives outside the hypotheses of
Proposition~\ref{4.27:prop-nonincreasing-thresholds}.} None of the following removes the
exception of~(ii) within the hypotheses of this chapter; none is adopted as a hypothesis
here, and nothing in what follows uses any of them.

\emph{Setting the current to zero.} A \emph{stored term} $\tau$ is a term kept for
composition at later operations; it is read over the determining set $\mathfrak{B}_\tau$ of
the date at which it is stored, and the density of that date is its \emph{birth density}.
Suppose that no value of a stored term, and no size
bound of a stored term, depends on the current $\mathbb{J}$ at those raw index-$0$ points
of its domain that are held by no live piece at any date of its birth density. The rider is
existential, and $\mathfrak{K}_{109}$ admits every $\mathfrak{g}^{\mathbb{C}}$-valued
current subject only to its bound, with no divergence condition $D^{*}\mathbb{J}=0$. So
the rider's current
may be set to $0$ at such points. The stratum current member then leaves the rider there,
and both inclusions hold with nothing to date.

\emph{Restricting to the sub-riders.} The conclusions of
Lemma~\ref{4.27:lem-weights-heredity} (whose proof is incomplete at one step) are asserted
for the sub-riders $\mathfrak{G}_1$ and $\mathfrak{G}^{-J_0}$, together with an explicit
list of the stratum members derived. For every statement that uses those conclusions, the
same question remains: whether it depends on the current at such points.

\emph{The maximal dating at every raw point.} One may take the maximal dating of~(i) at
every raw point, held or not, namely $T^{\flat}(y):=\gamma_0^{(k_{\circ}(y))}$. This alters
the frozen datum $\gamma_0^{(k_{\mathfrak{c}})}$ at held raw points, so it is a further
hypothesis, lying outside the pair of conditions of~(ii); it is not adopted, and nothing
here uses it.

\emph{A radius independent of the density.} Dating the member at unheld raw points by
$\gamma^{\top}:=\bar\gamma_0(\mathsf{t}_\gamma)$ of \eqref{4.27:eq-weights-heredity-a}
restores the current part of the first inclusion, since
\begin{equation*}
\gamma_0^{(k_{\mathfrak{b}})}\le\bar\gamma_0(\mathsf{t}_{k_{\mathfrak{b}}})\le\gamma^{\top}.
\end{equation*}
It replaces the second inclusion at such points by
$\mathfrak{G}(\mathfrak{B})\subset\mathfrak{K}_{109}(\mathfrak{B};\gamma^{\top})$, so it lies
outside the pair of conditions of~(ii). It is admissible only under the supposition of the
first alternative, since a size bound read at one radius independent of the density
amounts to a lower bound on the coupling; it is dominated by that alternative and is not
proposed.

\emph{Stored terms.} For a stored term the residual case of~(ii) is that of the raw points
of its domain that lie outside the class components of its birth step, that no live piece
holds at any date of its birth density, and that a later operation, whose source is
written over a determining set of larger density, meets; whether the indexing of domains
in \cite[(1.63)]{B-CMP122a} excludes this case is not decided here.
\end{remark}

\begin{remark}
The proposition and the one-point current probe use only the following:
\begin{itemize}
\item the definitions of $\mathfrak{G}$, of the frozen datum $T^{\mathfrak{B}}$ and of
$\mathfrak{K}_{109}$;
\item the fact that $\Omega_0(\cdot)$ and $\Omega_1(\cdot)$ never shrink;
\item the order of dates and the pointwise description of the moves between consecutive
dates (Lemma~\ref{4.27:lem-determining-set-order});
\item the nesting of the live pieces holding a point, the innermost being the oldest;
\item the statement on unheld raw points within one density, with its separating pair,
whose proof is incomplete at one step;
\item the monotonicity of $\gamma_0^{(\cdot)}$.
\end{itemize}
No numerical constant
enters, and no cap in the sense of Definition~\ref{2.3:def-cap} and no bound on the age of a
large-field region is read or introduced. The quantities $k$, $k(d')$, $k_{\mathfrak{b}}$ and
$k_{\circ}(y)$ are densities of dates, and $k_{\mathfrak{b}}-k$ enters only through the
monotonicity of $\gamma_0^{(\cdot)}$. The arguments concern inclusions of open sets of
pairs invariant under the gauge group $G=\mathrm{SU}(N)$, and use no commutativity of $G$:
the probe is $cX$ with $X$ any unit vector of $\mathfrak{g}$; the field $V$ with values in
$G$ is not evaluated.
\end{remark}

\begin{lemma}[No index-zero margin on source domains. Stated; proof incomplete at the step marked $(\ast)$]\label{4.27:prf-index-zero-margin}
Let $\mathfrak{O}$ be an operation that composes stored terms with a map of the pair
$(\mathbb{U},\mathbb{J})$: a $\mathbf{T}$-step, a stage of an attempt, or the completion site
of a completed chain, dissolutions included. Let $Y'$ be its source domain, and let
$\mathfrak{B}^t$ and $\mathfrak{B}'$ be the determining sets in force immediately before and
immediately after its map. For a stored term $\tau$ with domain $X_\tau$, read $Y'$ as
$X_\tau\cap Y'$: the value of $\tau$ depends on the pair on $X_\tau$ only, and in each case of
the classification of later operations the source is written over a domain containing the part
of $X_\tau$ on which the map acts, so that on $X_\tau\cap Y'$ the continuation must be the image
of the map. A stored term is a \emph{consumer} at $\mathfrak{O}$ if its bound must hold for the
image of the current source. Call $y$ \emph{raw} at $\mathfrak{B}$ if
$y\in(\Omega_0\setminus\Omega_1)(\mathfrak{B})$. At a raw point $y$ let $T^{\mathfrak{B}}(y)$
be the threshold of the stratum $J$-member of $\mathfrak{G}(\mathfrak{B})$: it is
$\gamma_0^{(k_{\mathfrak{c}})}$ if $y$ lies in a live piece, $k_{\mathfrak{c}}$ being the
density of the innermost live piece holding $y$, and $\gamma_0^{(k)}$ otherwise, $k$ being the
density of the date of $\mathfrak{B}$ (the \emph{birth-frozen dating}). Write
$\mathbb{U}=e^{i\eta A'}U$ with $U$ the $G$-valued part; the map of $\mathfrak{O}$ is
\emph{a layer on the $A'$-part} at a point if it leaves $U$ and the cube gauges unchanged there.
At every point $y$ of $Y'$ that is raw at $\mathfrak{B}^t$ the following hold.
\begin{enumerate}
\item[(a)] The clause imposed on the source at $y$ and the clause required of the image by a
consumer coincide at index zero at the stages of an attempt, at completion sites, and at a
$\mathbf{T}$-step on live pieces; at a $\mathbf{T}$-step off live pieces they coincide provided
no factor clause of \cite[(2.34)--(2.39)]{B-CMP119} is taken at level $m=0$, with a factor
$s^{(k)}_0$, on the part of the stratum lying in the support of $\Omega_1$. The common clause
consists of the gauge members
on the index-zero cubes (the index-zero instance of \cite[(3.35)]{B-CMP99b}), the
$A'$-members (the index-zero instance of \cite[(3.37)]{B-CMP99b}), and the stratum
$J$-member
\begin{equation*}
\ell^{3}|\mathbb{J}(y)| < T^{\mathfrak{B}}(y).
\end{equation*}
Under the birth-frozen dating the $J$-thresholds of source and consumer are equal at held raw
points (Lemma~\ref{4.27:prf-held-raw-points}), at unheld raw points between dates of one
density and, from a point's first holding within that density on, at every later date
(Lemma~\ref{4.27:prf-unheld-raw-points}, whose proof is incomplete at one step), and
at the raw points of the domains of an arrest's stored terms, by the statement on the datings
of an arrest's stored terms, except at the points unheld at every date of the birth density
and met by a source of a larger density (the residual case of
Lemma~\ref{4.27:prf-density-increase}, whose proof is incomplete at one step); under any
dating that does not depend on the density they are equal at every raw point.
\item[(b)] Wherever the map of $\mathfrak{O}$ is a layer on the $A'$-part, that is, everywhere
except, when $\mathfrak{O}$ is a completion site, on $\mathsf{Z}$ and on the core of the
completion site, the
$G$-part of the image satisfies the index-zero gauge members of $\mathfrak{G}(\mathfrak{B}^t)$.
\item[(c)] There is room at index zero, in the clause required by a consumer, for the members
that the map moves at $y$, namely the $A'$-members and the stratum $J$-member, and, at a
completion site, for every index-zero member of the pair
$(U_2,\mathcal{J}(U_2))$ in which the configuration is re-presented on $\mathsf{Z}$ and on the
core, at raw points of $\mathsf{Z}$ and of the core.
\end{enumerate}
Moreover, for a real $V$ in the support of the density's characteristic functions and
$\mathfrak{B}$ the current determining set, the real point $(U(\mathfrak{B},V),\mathcal{J}(U))$
satisfies, at raw points, the index-zero gauge member of $U$ and the stratum $J$-member.
\end{lemma}

\begin{proof}
Throughout, $\mathfrak{G}(\mathfrak{B})$ is the rider on $\Omega_0(\mathfrak{B})$, and
$\mathfrak{G}_1(\mathfrak{B})$ is its sub-rider at the points of level at least one.
$\mathfrak{G}^{-J_0}(\mathfrak{B})$ is the same sub-rider together with the index-zero
$U$-members of the stratum. $\mathfrak{K}_{109}(\mathfrak{B};\gamma)$ is the space of pairs of
the type of \cite{B-CMP109} with $J$-radius $\gamma$. $\gamma_0^{(k)}$ is the $J$-radius at
density $k$, namely five times the maximum over $j\le k$ of $w_{*}(j)\bar\alpha_j$; it is
non-decreasing in $k$. A raw point is \emph{held} at a date if it lies in the piece
$W_{\mathfrak{c}}$ of an arrest $\mathfrak{c}$ live at that date.

\emph{Step 1: the same index-zero clause.}
Among the statements on which this section rests, none except possibly the factor clauses
discussed after the following list provides, at a raw point, an index-zero clause with a factor
below the rider. The reasons are as follows.
\begin{itemize}
\item Ba{\l}aban's attempt spaces begin at level one \cite[p.~190]{B-CMP122a}.
\item On a class component $X$ the stage structure satisfies
$\Omega''_1\cap X=\Omega_1\cap X$, by the closed form of the stage structure in the lemma on
the order of dates. Hence no stage creates a clause at a raw point.
\item The rules defining the arrested spaces alter no clause at a raw point.
\end{itemize}
The arrested spaces take the factor clauses of \cite[(2.34)--(2.39)]{B-CMP119}
unchanged off live pieces, and the stage factors $s^{(k)}_m$ are taken from them. If these
factor clauses include the level $m=0$ with factor $s^{(k)}_0$ on the part of the stratum
$(\Omega_0\setminus\Omega_1)(\cdot)$ lying in the support of $\Omega_1$
\cite[p.~255]{B-CMP119}, then the $\mathbf{T}$-step carries an index-zero
$J$-margin off live pieces, and the step $(\ast)$ below narrows there.
Therefore the clause imposed on the source and the clause required by a consumer coincide at
$y$ at the stages, at completion sites and at a $\mathbf{T}$-step on live pieces, and at a
$\mathbf{T}$-step off live pieces under the proviso stated in (a).

Under the birth-frozen dating, the $J$-thresholds of this clause are equal at the following
points:
\begin{itemize}
\item at held raw points, by the current bound at held raw points
(Lemma~\ref{4.27:prf-held-raw-points});
\item at unheld raw points between dates of one density and, from a point's first holding
within that density on, at every later date, by
Lemma~\ref{4.27:prf-unheld-raw-points}, whose proof is incomplete at one step;
\item at the raw points of the domains of an arrest's stored terms, whose clause carries the
thresholds of its birth date, by the statement on the datings of an arrest's stored terms,
except at points unheld
at every date of that density and met by a source of a larger density
(Lemma~\ref{4.27:prf-density-increase}, whose proof is incomplete at one step).
\end{itemize}
Under any dating that does not depend on the density, the thresholds are equal trivially.
This is (a).

\emph{Step 2: the maps move $A'$ and $\mathbb{J}$ at raw points.}
Each map moves $A'$ and $\mathbb{J}$ at $y$ by its layers. We take the three kinds of
operation in turn.

At a $\mathbf{T}$-step $k'\to k'+1$, the map on pairs is
\begin{equation*}
(\mathbb{U},\mathbb{J})\longmapsto
\bigl(e^{i\eta\mathfrak{h}}\mathbb{U},\ \mathbb{J}+\mathcal{J}^{\sharp}[1]\bigr),
\end{equation*}
with the bounds \eqref{4.27:eq-t-step-layer-a} of the $\mathbf{T}$-step layer lemma
(Lemma~\ref{4.27:lem-t-step-layer}, whose proof is incomplete at one step):
\begin{equation*}
N^{\Delta}_y(\mathfrak{h})\le B_\sharp\,p(y),\qquad
\ell^{3}|\mathcal{J}^{\sharp}[1](y)|\le 2K_\chi\,p(y).
\end{equation*}
The profile $p$ is the bound on the $\mathbf{T}$-datum times
$e^{-\delta d^{\wedge}(\cdot,\mathcal{C})}$. Write $\varpi_{k'}$ for that bound at step $k'$.

From a point of index zero to the bond set $\mathcal{C}\subset\Omega_{k'+1}$ of the
$\mathbf{T}$-datum, a contour crosses the strata $1,\dots,k'-1$ and one $M$-cube of its own
scale in $\Omega_{k'}\setminus\Omega_{k'+1}$. That is $k'$ crossings. By the contour-cost
bound among the inputs of Lemma~\ref{4.27:lem-t-step-layer}, whose proof is incomplete at one
step, each crossed stratum or $M$-cube of the local scale costs at least $\mathsf{w}_M$, so
$d^{\wedge}(y,\mathcal{C})\ge\mathsf{w}_M k'$. Hence
\begin{equation*}
p(y)\le \varpi_{k'}\,e^{-\delta\mathsf{w}_M k'}.
\end{equation*}
This bound is small but positive, and nothing in Lemma~\ref{4.27:lem-t-step-layer} shows that
the layer itself vanishes at raw points.

At the stages of an attempt, the map acts by the shift layers of the holomorphic stage maps
$\Phi_n\colon\mathfrak{S}_{n+1}(X)\to\mathfrak{S}_n(X)$ and of the stage maps taken at a
bounded datum $b$ (the two families of the hypotheses of Lemma~\ref{4.27:lem-t-step-layer},
whose proof is incomplete at one step). These use at most half the one-stage room
$\rho^{(l)}_m$ of stage $l$ on the shell of level $m\ge1$. At index zero no unit is provided.

At the completion site of a completed chain, including a dissolution, the map acts by the
layers of the operation $\mathbf{R}$ off the zone of the completion site. Their increments are
measured in the units of shells that carry a level. On $\mathsf{Z}$ and on the core of the
completion site, the pair is re-presented as $(U_2,\mathcal{J}(U_2))$, together with the
accompanying $W$-presentation. The bounds on the presented pair on $\mathsf{Z}$ and on the
core are
stated against thresholds whose level, at a point $x$, is the current level of $x$ or a finer
birth index. At a
raw point of the core the current level is $0$, and no threshold is provided there, neither for
the gauges of $U_2$ nor for $\mathcal{J}(U_2)$.

\emph{Step 3: no margin, by any dating.}
Let $T>0$ and let $v\in\mathfrak{g}^{\mathbb{C}}$ with $v\neq0$. The translate by $v$ of the
open set $\{J:\ell^{3}|J|<T\}$ is not contained in that set. Indeed, if $\ell^{3}|v|\ge T$,
take $J=0$: then $\ell^{3}|J+v|\ge T$. Otherwise choose $s$ with
$T/\ell^{3}-|v|<s<T/\ell^{3}$ and put $J=s\,v/|v|$. Then $\ell^{3}|J|<T$, while
\begin{equation*}
\ell^{3}|J+v|=\ell^{3}(s+|v|)>T.
\end{equation*}
The $A'$-members at index zero are open bounds of the same kind, and the same holds for them.

By Step 1, source and consumer are subject to the same index-zero clause, with the same
threshold (at a $\mathbf{T}$-step off live pieces, under the proviso stated in (a)). By
Step 2, the maps move $A'$ and $\mathbb{J}$ at raw points by layers that none of
the available bounds shows to vanish there. Hence, the thresholds being equal, no margin at
index zero is available from these bounds for these members, whatever the dating.

\emph{Step 4: the $G$-part.}
The stage maps keep the $G$-part fixed. The layers of $\mathbf{R}$ off the zone keep the
$U$-part and the cube gauges. The map of the $\mathbf{T}$-step is a layer on the $A'$-part.

The index-zero gauge members are the index-zero instance of \cite[(3.35)]{B-CMP99b} on the
index-zero cubes. These cubes form a family that does not depend on the date, and the clause
does not involve $\gamma$. The source pair obeys them, because
$\mathfrak{G}(\mathfrak{B}')\subset\mathfrak{G}^{-J_0}(\mathfrak{B}^t)$
(Lemma~\ref{4.27:prf-raw-stratum}). Therefore the
image obeys the same members of $\mathfrak{G}(\mathfrak{B}^t)$ wherever the map is a layer
on the $A'$-part. That is everywhere except, when $\mathfrak{O}$ is a completion site, on
$\mathsf{Z}$ and on the core of the completion site. This is (b).

\emph{Step 5: the real point.}
The statement that the real minimiser lies in arrested spaces
(Lemma~\ref{4.27:lem-real-minimiser}, whose proof is incomplete at one step) applies here, as
does the bound on the real point in the presence of excluded attempts. Together they place
$(U(\mathfrak{B},V),\mathcal{J}(U))$ in every space whose clause at a point $x$ is at the
current level of $x$; these two statements concern clauses at levels $\ge1$ only, and at a raw
point the level is $0$.

At a raw point, $A'=0$ holds, and hence the $A'$-members of index zero hold. The index-zero
gauge member of $U$ and the bound $\ell^{3}|\mathcal{J}(U)(y)|<T^{\mathfrak{B}}(y)$ are
supplied by none of the statements used here.

\emph{Step 6 $(\ast)$: the members that the map moves.}
By Steps 2 and 3, the following are not derived from the statements on which this section
rests:
\begin{itemize}
\item the inclusion in the index-zero clauses of $\mathfrak{G}(\mathfrak{B}^t)$ of the members
that the map moves at raw points of $Y'$, namely the $A'$-members and the stratum $J$-member;
\item at a completion site, the inclusion of every index-zero member of the re-presented pair at
raw points of $\mathsf{Z}$ and of the core;
\item by Step 5, the index-zero gauge member and the stratum $J$-member of the real point at raw
points.
\end{itemize}
This is clause (c) together with the last assertion of the statement, and it is the step at
which the proof is incomplete.

Clause (c) would follow either from the four verifications (1)--(4) below, each of which
settles the corresponding part of the list above, or from (5):
\begin{enumerate}
\item that the layers of all the maps above are supported in $\Omega_1(\cdot)$, so that
nothing moves at raw points of the domains of stored terms;
\item that the bounds on the presented pair at a completion site hold at raw points of
$\mathsf{Z}$ and of the core;
\item that the real base pair obeys the index-zero clause there with room;
\item for the $\mathbf{T}$-step off live pieces, the alternative in the conditional of Step~1:
a factor clause of \cite[(2.34)--(2.39)]{B-CMP119} at level $m=0$, with factor $s^{(k)}_0$, on
the part of the stratum lying in the support of $\Omega_1$;
\item an index-zero margin built into the definition of the rider.
\end{enumerate}

\emph{Consequences for the statements of this section.}
The conclusions of the statement on the direct factors of the rider (both of its clauses), of
the arrest theorem's clauses on the rider of later sources and images, and of the second clause
of the proposition on the rider of sources, together with two further statements on the rider
imported by later sections, assert membership in the sub-rider $\mathfrak{G}_1$, and
where they assert an inclusion of riders, an inclusion into $\mathfrak{G}^{-J_0}$. In both
cases the stratum members that are derived are listed explicitly. None of these statements
asserts membership in $\mathfrak{G}$. The remaining clause of the arrest theorem and the first
clause of that proposition carry the same restriction. The members of clause (c) at raw points
of $Y'$, or of $X_\tau\cap Y'$, occur in the conclusion of no other statement of this section,
and no number, room or constant of the arrested construction depends on them.

\emph{The dependence on the birth-frozen dating.}
The birth-frozen dating requires that no estimate of a stored term involve, at raw points of a
live class component of density $k_{\mathfrak{a}}$, a $J$-radius larger than the value
$\gamma_0^{(k_{\mathfrak{a}})}$ of the innermost live piece.
Among the statements on which this section rests, none does. The membership in
$\mathfrak{K}_{109}$
required by the change-of-variables theorem for the function $\mathbb{H}$, the bound of the
source by $\gamma_0$ in the smallness condition of the shift estimate, and the shift estimate
itself are hypotheses satisfied a fortiori by a smaller $\mathbb{J}$. The other lemmas of this
section involve units only at levels $\ge1$. For the estimates of stored terms formed in the
other constructions of the book that store such terms, the requirement is a hypothesis on those
estimates; this lemma does not verify it, and it is not part of the step $(\ast)$.

\emph{The second inclusion on the stratum.}
At the $\mathbf{T}$-step $k\to k+1$ of Lemma~\ref{4.27:lem-t-step-layer}, with $\mathfrak{B}$
the step's determining set, frozen on live pieces, the pair of the enlarged source over
$\mathbb{B}_{k+1}$, with rider $\mathfrak{G}(\mathbb{B}_{k+1})$, lies in
$\mathfrak{K}_{109}(\mathfrak{B};\gamma_0^{(k)})$ on $\Omega_1(\mathfrak{B})$ and in
$\mathfrak{K}_{109}(\mathfrak{B};\gamma_0^{(k+1)})$ on the stratum: on $\Omega_1(\mathfrak{B})$
by the inclusions of the rider applied with $\mathfrak{B}\preccurlyeq\mathbb{B}_{k+1}$, and on
the stratum by the second inclusion of the rider at the density $k+1$ of $\mathbb{B}_{k+1}$.
Both lie in
$\mathfrak{K}_{109}(\mathfrak{B};\gamma^{\top})$, with
\begin{equation*}
\gamma^{\top}:=\bar\gamma_0(\mathsf{t}_\gamma)\ \ge\ \bar\gamma_0(\mathsf{t}_{k''})
\ \ge\ \gamma_0^{(k'')}
\end{equation*}
at every density $k''$. The first inequality holds because
$\mathsf{t}_{k''}\ge\mathsf{t}_\gamma$ and $\bar\gamma_0$ is decreasing there. The second is
the bound of $\gamma_0^{(k'')}$ by its majorant $\bar\gamma_0(\mathsf{t}_{k''})$.

On $\mathfrak{K}_{109}(\mathfrak{B};\gamma^{\top})$, the three clauses of the
change-of-variables theorem for $\mathbb{H}$ used here (the vanishing of $\mathbb{H}$ on a
component without data, the bound on the difference at two data, and the identification at the
real point with the function of Ba{\l}aban's construction) hold with their constants, which do
not depend on
$\gamma$. The smallness conditions entering those constants include the theorem's ceiling on
the radius, the smallness condition $K_\Delta\gamma_0\le\tfrac12$ on the size constant
$K_\Delta$ of the increments of the stored terms, and the smallness condition of the shift
estimate. The quantities these conditions bound are suprema over $\mathsf{t}\ge\mathsf{t}_\gamma$,
attained at $\mathsf{t}_\gamma$. The
second inclusion therefore enters there as a membership only, and no constant, floor, ceiling
or room depends on it.

The same holds for two further memberships:
\begin{itemize}
\item for an arrest $\mathfrak{a}$ of density $k_{\mathfrak{a}}$, the membership in
$\mathfrak{K}_{109}(\mathfrak{B}^{\circ};\gamma_0^{(k_{\mathfrak{a}})})$ at a
stratum point of $X_{\mathfrak{a}}$, if there is one, where $\mathfrak{B}^{\circ}$ is the
determining set $\mathbb{B}^{(n+1)}_{k_{\mathfrak{a}}}$ of stage $n+1$, the frozen stage index
of $\mathfrak{a}$, restricted to $X_{\mathfrak{a}}$;
\item the memberships of the same type used at the exit of an arrest.
\end{itemize}
No dating independent of the density enters any rider or any size of a stored term. For a
stored term $\tau$, with $k_\tau$ the density of its birth date, the size bound
$K_\Delta\gamma_0$ of the stored terms involves
$\gamma_0^{(k_\tau)}\le\bar\gamma_0(\mathsf{t}_{k_\tau})$.

Under the birth-frozen dating this stands. The stratum $J$-datum $T^{\mathfrak{B}_\tau}(y)$ of
a rider is either $\gamma_0^{(k_{\mathfrak{c}})}\le\gamma_0^{(k_\tau)}$, with
$k_{\mathfrak{c}}$ the density of the innermost live piece holding $y$, or
$\gamma_0^{(k_\tau)}$ itself. It is a membership on which no constant, floor, ceiling or room
depends, and no floor of the kind of $K_\Delta\gamma_0$ is created. The sub-riders
$\mathfrak{G}_1$ and $\mathfrak{G}^{-J_0}$ carry no $\gamma_0$ at all.
\end{proof}

\begin{lemma}[Inclusion of the ambient domain at size $\gamma_0^{(k)}$. Stated; proof incomplete
at the step marked $(\ast)$]
\label{4.27:prf-ambient-inclusion}
The following are the second inclusion of (e3) and clause (e4) of
\eqref{4.27:eq-weights-heredity-b}.
Let $\mathfrak{B}$ be a determining set of density $k$, dated in the time order
\eqref{4.27:eq-determining-set-order-a}. Let $\mathfrak{G}(\mathfrak{B})\subset
\mathfrak{G}^{-J_0}(\mathfrak{B})\subset\mathfrak{G}_1(\mathfrak{B})$ and the dated threshold
$T^{\mathfrak{B}}$ be as in \eqref{4.27:eq-arrested-spaces-a}.

For $y\in\Omega_1(\mathfrak{B})$ and a direct entry $Q_a$, $a\in\{1,\dots,7\}$, write
$\theta^{\mathfrak{B}}_a(y)=w_{\mathfrak{B}}(y)\,r_a(\ell_{\mathfrak{B}}(y))$ as in
\eqref{4.27:eq-weights-heredity-b}, $w_{\mathfrak{B}}$ the actual weight. For the cube clause
($a=8$) on a cube $\square$ of $\mathcal{Q}_{\mathfrak{B}}$ of family $(\lambda_\square,w_\square)$
the threshold is $\theta^{\mathfrak{B}}_8(\square)=w_\square BCM\alpha_{0,\lambda_\square}$, as in
\eqref{4.27:eq-arrested-spaces-a}, where the constants $B$ and $C$ are fixed.

The gauge member, the $A'$-member and the $J$-member $\ell^3|\mathbb{J}(y)|<T^{\mathfrak{B}}(y)$ of
the raw stratum $(\Omega_0\setminus\Omega_1)(\mathfrak{B})$ are those of
\eqref{4.27:eq-arrested-spaces-a}; recall $T^{\mathfrak{B}}(y)\le\gamma_0^{(k)}$ there. Recall that
$\gamma_0^{(k)}=5\max_{j\le k}w_*(j)\bar\alpha_j$, which is non-decreasing in $k$.
A point of $(\Omega_0\setminus\Omega_1)(\mathfrak{B})$ is called raw, and a raw point $y$ is
\emph{held} at a date by a live piece $W_{\mathfrak{c}}$ if $y\in W_{\mathfrak{c}}$ and the arrest
$\mathfrak{c}$ is live at that date.

Let $\mathfrak{K}_{109}(\mathfrak{B};\gamma)$ be the space of pairs $(\mathbb{U},\mathbb{J})$ over
$\mathfrak{B}$ of the type of \cite{B-CMP109} with $J$-radius $\gamma$; it increases with $\gamma$,
and its $J$-member admits every $\mathfrak{g}^{\mathbb{C}}$-valued source with
$\sup\ell^3|\mathbb{J}|\le\gamma$, no condition $D^*\mathbb{J}=0$ being imposed. Let $a_\flat$ be
the size formed from the direct entries $5,1,6,7$ in the units of the point. Then:
\begin{enumerate}
\item[(a)] $\mathfrak{G}(\mathfrak{B})\subset\mathfrak{K}_{109}(\mathfrak{B};\gamma_0^{(k)})$ as sets
on all of $\Omega_0(\mathfrak{B})$, and $a_\flat\le\gamma_0^{(k)}$ on $\mathfrak{G}(\mathfrak{B})$.
\item[(b)] The real point of the real-point lemma, and every pair that obeys at every point the
direct and cube clauses of an arrested space, of an arrested source or of the domain of a stored
term over
$\mathfrak{B}$, lie in $\mathfrak{G}_1(\mathfrak{B})$ with $3.6$ units to spare: such a pair obeys
the clauses of $\mathfrak{G}_1(\mathfrak{B})$ wherever the space carries a clause, that is at the
points of its own domain $X\cap\Omega_1(\mathfrak{B})$ and on its cubes, and is there the
restriction of a pair of $\mathfrak{G}_1(\mathfrak{B})$.
\item[(c)] Let $\mathfrak{O}$ be a later operation in the sense of
Lemma~\ref{4.27:lem-determining-set-order} (a stage of an attempt, a $\mathbf{T}$-step, or the
completion site $(\mathrm{L})$ of a completed chain); let $\mathfrak{B}^t$ be the determining set
in force immediately before the map of $\mathfrak{O}$ and $\mathfrak{B}'$ the one after it, over
which the arrested source of $\mathfrak{O}$ is written, so that $\mathfrak{B}^t\preceq
\mathfrak{B}'$. Let $p$ be a member of that source over the domain $Y'$ on which the map acts,
and let $\tilde p\in\mathfrak{G}(\mathfrak{B}')$ be an extension of $p$ to a pair on
$\Omega_0(\mathfrak{B}')$ (only the existence of such an extension is used). Let $\tilde q$ be the
pair equal on
$Y'$ to the image of $p$ under the map of $\mathfrak{O}$, whose $\mathbb{U}$ off $Y'$ is that of
$\tilde p$, and whose $\mathbb{J}$ off $Y'$ is that of $\tilde p$ on $\Omega_1(\mathfrak{B}^t)$ and
at the raw points of $(\Omega_0\setminus\Omega_1)(\mathfrak{B}^t)$ held at the date of
$\mathfrak{B}^t$ by a live piece, and is $0$ at the raw points held at that date by no live
piece. Then $\tilde q\in\mathfrak{G}_1(\mathfrak{B}^t)$. Moreover, on
$(\Omega_0\setminus\Omega_1)(\mathfrak{B}^t)$: off $Y'$, $\tilde q$ obeys every index-zero
$U$-member of $\mathfrak{G}(\mathfrak{B}^t)$ and its $J$-member; on $Y'$, $\tilde q$ obeys the
index-zero gauge members of the $G$-part wherever the map of $\mathfrak{O}$ is a layer into $U'$,
that is, leaves the $G$-part $U$ of $\mathbb{U}=e^{i\eta A'}U$ and the cube gauges unmoved,
namely at the stages and the $\mathbf{T}$-steps. No assertion is made on
$Y'\cap(\Omega_0\setminus\Omega_1)(\mathfrak{B}^t)$
about the $A'$-member of the form \cite[(3.37)]{B-CMP99b} at index zero and the $J$-member, nor,
at the site $(\mathrm{L})$, about the index-zero gauge members of the $G$-part or about any
index-zero member of the re-presented pair at the raw points
of $\mathsf{Z}$ and of the core of the completion site.
\end{enumerate}
\end{lemma}

\begin{proof}
\emph{The second inclusion.} We follow the bound on the ambient space in
Lemma~\ref{4.27:prf-weight-bounds}, with the actual weight $w_{\mathfrak{B}}$ in the place of the
larger weight used there. Each clause of $\mathfrak{K}_{109}(\mathfrak{B};\gamma)$ at a point of
$\Omega_1(\mathfrak{B})$ is an absolute threshold on the product of a weight and a radius
$\alpha_{i,\lambda}$, $i\in\{0,1\}$, at the point's own level $\lambda$; this is the comparison of
the clauses of $\mathfrak{K}_{109}$ with the direct entries used throughout this chapter, under
which entries $5$ and $8$ give the clauses of the form \cite[(3.35)]{B-CMP99b} and entries $6$ and
$7$ those of the form \cite[(3.37)]{B-CMP99b}. Let $y\in\Omega_1(\mathfrak{B})$ and
$\lambda=\ell_{\mathfrak{B}}(y)\le k$. If $w_{\mathfrak{B}}(y)=1$, then
\begin{equation*}
5\alpha_{i,\lambda}\le5\bar\alpha_k\le5w_*(k)\bar\alpha_k ,
\end{equation*}
since $\bar\alpha_k$ is the running maximum of the radii over the levels $\le k$ and
$w_*(k)\ge1$. If $w_{\mathfrak{B}}(y)>1$, the structure at $y$ is that of an attempt of a step
$k_c\le k$ (a live arrest, or the current run), and the weight bound of
Lemma~\ref{4.27:prf-weight-bounds} gives $w_{\mathfrak{B}}(y)\le w_*(k_c)$ and $\lambda\le k_c$,
whence
\begin{equation*}
5w_{\mathfrak{B}}(y)\,\alpha_{i,\lambda}\le5w_*(k_c)\bar\alpha_{k_c}.
\end{equation*}
The cube families are treated in the same way: a family weight larger than $1$ is the weight
$\mathsf{L}_0^{2(\cdot)}$ of a stage family, which is at most $w_*(k_c)$, and the family's level
is at most $k_c$. Both right-hand sides are at most $\gamma_0^{(k)}$, because $\gamma_0^{(k)}$ is
five times the maximum of $w_*(j)\bar\alpha_j$ over $j\le k$ and $k_c\le k$. Entry $3$ is the
$\gamma_0$-clause of $\mathfrak{K}_{109}$; entries $5,8$ and $6,7$ give its clauses of the forms
\cite[(3.35)]{B-CMP99b} and \cite[(3.37)]{B-CMP99b}; and entries $5,1,6,7$ make up $a_\flat$, which
is therefore at most $\gamma_0^{(k)}$ on $\mathfrak{G}(\mathfrak{B})$. On the raw stratum the
$U$-members of $\mathfrak{G}(\mathfrak{B})$ are, by \eqref{4.27:eq-arrested-spaces-a}, the
index-zero $U$-clauses of $\mathfrak{K}_{109}(\mathfrak{B};\cdot)$, which do not depend on the
radius; and the $J$-member $\ell^3|\mathbb{J}(y)|<T^{\mathfrak{B}}(y)$ implies the $J$-member of
$\mathfrak{K}_{109}(\mathfrak{B};\gamma_0^{(k)})$, since $T^{\mathfrak{B}}(y)$ is
$\gamma_0^{(k_{\mathfrak{c}})}$ for an arrest $\mathfrak{c}$ with $k_{\mathfrak{c}}\le k$, or
$\gamma_0^{(k)}$, and $\gamma_0^{(\cdot)}$ is non-decreasing. This proves (a).

\emph{Direct factors.} The factors multiplying the thresholds in the direct clauses are at most
$1+\beta$ on the arrested spaces and on the enlarged sources; at most $1+2\beta$ in the ambient
space $\mathfrak{N}$ at off-class points (points of $\mathfrak{T}_k=\Omega_k\cup Z^c$ for the
density $k$ in force), where the unit is the one frozen on live pieces by
\eqref{4.27:eq-arrested-spaces-a}; and at most $1$ at points of $Z^{\mathrm{on}}$, where this bound
is part of the definition. At the real point of the real-point lemma they are at most
$\mathsf{m}(1+\beta)$, $\mathsf{m}$ the fraction of that lemma. With the value of $\beta$ fixed in
Lemma~\ref{4.27:lem-factor-positivity} one has $1+2\beta=1.4$, so all these factors are at most
$1.4$ ($\mathsf{m}\le1$ being a fraction); since the thresholds of $\mathfrak{G}(\mathfrak{B})$ are
$5\theta^{\mathfrak{B}}_a$, each entry of such a pair is below $1.4\,\theta^{\mathfrak{B}}_a$, which
leaves $5-1.4=3.6$ units. The cube clauses carry no factor: they are stated at the family's
weight $w$ (see \eqref{4.27:eq-arrested-spaces-a} and \cite[(1.65), (1.67), (1.69)]{B-CMP122a}) and
so lie
inside the threshold $5w$; the gauges of the real point are those of the real-point bound on
determining sets holding excluded attempts, which places the real point in every space whose clause
at $x$ has the current level of $x$, a weight $\ge1$, a multiplicity $\ge1$ and a factor
$\ge\varphi_\flat$, taken here at the fraction $\mathsf{m}$. Hence every pair of (b) obeys the
clauses of $\mathfrak{G}(\mathfrak{B})$ at every point of $\Omega_1(\mathfrak{B})$ in its domain
and on every cube of $\mathcal{Q}_{\mathfrak{B}}$ there, which is membership in
$\mathfrak{G}_1(\mathfrak{B})$. On the raw stratum the real point has $A'=0$, which gives the
$A'$-member; nothing else is derived there. This proves (b).

\emph{Maps: the direct entries on $Y'$.} Let $\mathfrak{O}$, $\mathfrak{B}^t\preceq\mathfrak{B}'$,
$p$, $\tilde p$ and $\tilde q$ be as in (c). By (e2) of \eqref{4.27:eq-weights-heredity-b},
thresholds do not rise along $\preceq$, so $\theta^{\mathfrak{B}'}_a\le\theta^{\mathfrak{B}^t}_a$.
At $y\in Y'$ each direct entry of $p$ is below $1.4\,\theta^{\mathfrak{B}'}_a(y)\le1.4\,
\theta^{\mathfrak{B}^t}_a(y)$, and the map of $\mathfrak{O}$ moves it by at most
$0.1\,\theta^{\mathfrak{B}^t}_a(y)$. Indeed, in units of the threshold, the moves are bounded as
follows (Lemma~\ref{4.27:lem-live-piece-increments}): at a stage, by $\tfrac12\rho^{(l)}\le\beta/2$
at points of $Z^{\mathrm{on}}$, by the stage-map statement (the stage map from
$\mathfrak{S}_{n+1}(X)$ to $\mathfrak{S}_n(X)$ is holomorphic, fixes the $G$-part, and its shifts
use at most half the one-stage room), and by the summed increments
$\sum_l\iota^1_l\le2\mathsf{w}_0\beta$ of the inheritance statements for successor stages at points
of $\mathfrak{T}_k$; through the summed foreign increments $\sum_l\iota_l$ of the stages of an
attempt
at a later density, by $\tfrac12\mathsf{w}_0\beta$; at the completion site, by $\bar P\le
\beta\theta_S/16$; at a $\mathbf{T}$-step, by $\beta/8$. Hence at $y\in Y'$ each direct entry of
the image of $p$ is below
\begin{equation*}
1.4\,\theta^{\mathfrak{B}^t}_a(y)+0.1\,\theta^{\mathfrak{B}^t}_a(y)=1.5\,\theta^{\mathfrak{B}^t}_a(y)
<5\,\theta^{\mathfrak{B}^t}_a(y).
\end{equation*}

\emph{Maps: the cube clause on $Y'$.} Where the map of $\mathfrak{O}$ is a layer into $U'$, the
$G$-part $U$ and the cube gauges do not move: at a stage by the stage-map statement; at the
completion site, at the points served by the layer bound at the completion site (there the
presented pair is the identity times layers, the $U$-part and the cube gauges being kept, with
increments at most $\bar P2^{-(k-1-m)}$ on foreign shells $m\le k-1$ and bounded by $\bar P$ on
$\Omega_k\setminus Z$ and on the collars, linear in the layer, and $16\bar P\le\beta\theta_S$); at
a $\mathbf{T}$-step by the $\mathbf{T}$-step statement, under which the map on pairs is
\begin{equation*}
(\mathbb{U},\mathbb{J})\mapsto\bigl(e^{i\eta\mathfrak{h}}\mathbb{U},\
\mathbb{J}+\mathcal{J}^{\sharp}[1]\bigr).
\end{equation*}
So every cube clause passes from $\mathfrak{G}(\mathfrak{B}')$ to $\mathfrak{G}(\mathfrak{B}^t)$ by
(e3) of \eqref{4.27:eq-weights-heredity-b}, that is by the nested-cube statement
\eqref{4.27:eq-nested-cubes-a}. The same two observations place in $\mathfrak{G}_1(\mathfrak{B}^t)$
every member of the family of shifted stage maps of the stage-map statement, whose moves are at
most the one-stage room $\rho_m$ at the level $m$ of the point, and $\rho_m<3.6$. At the site
$(\mathrm{L})$ these points are the foreign shells,
$\Omega_k\setminus Z$, the collars and the rings off $\mathsf{Z}$; there the pair is the identity
times layers, and the increments of the layer bound are in the units of the structure in force at
the point, which are those of $\mathfrak{B}^t$ (native units on foreign shells; level $k$ and
weight $1$ on $\Omega_k\setminus Z$ and on the collars; on the rings off $\mathsf{Z}$ their raised
level in $\mathbb{B}''_k$). So there each entry is at most
\begin{equation*}
1.4\,\theta^{\mathfrak{B}^t}_a+\bar P\,\theta^{\mathfrak{B}^t}_a<1.5\,\theta^{\mathfrak{B}^t}_a .
\end{equation*}
On $\mathsf{Z}$ and on the core the pair is re-presented; its cube clause is placed below.

\emph{Maps: the completion site on $\mathsf{Z}$.} At the site $(\mathrm{L})$ of a completed chain
the set $\mathfrak{B}^t$ is $\mathbb{B}''_k$ dated during the step, with the levels
$o(x):=\ell_{\mathfrak{B}^t}(x)$, the level of $x$ in $\mathbb{B}''_k$, and
the frozen weights, and $\mathfrak{B}_\tau\preceq\mathbb{B}''_k\preceq\mathfrak{B}'$ for every
stored term $\tau$ served (Lemma~\ref{4.27:lem-determining-set-order}). By the bound on the presented
pair at the completion site, on $\mathsf{Z}$ each entry of the presented pair is at most $1/16$ of
the threshold of its consumer $\mathfrak{E}$, for every $\mathfrak{E}$ in the consumer class
$\mathcal{C}$. At $y\in\mathsf{Z}\cap Z^{\mathrm{on}}$ we take $\mathfrak{E}:=
\mathbb{C}^{(N)}_k(X)$ with $X\ni y$. A consumer of a stage $i<N$ would not serve: its clauses at
$y$ are stated with respect to $\mathbb{B}^{(i)}_k\preceq\mathbb{B}''_k$, and by (e2) of
\eqref{4.27:eq-weights-heredity-b} its threshold $\theta^{\mathbb{B}^{(i)}_k}_a$ is at least
$\theta^{\mathbb{B}''_k}_a$, which is the wrong direction. Such an $\mathfrak{E}$ exists through
every $y\in\mathsf{Z}\cap Z^{\mathrm{on}}$: \cite[(1.63)]{B-CMP122a} indexes the stage-$N$ terms by
every $X\in\mathfrak{D}_m$ with $X\cap(Z_k\setminus Z''_k)\ne\emptyset$; the band
$Z\cap\Omega^{\sim5}_{k_0+1}$ \cite[\S1]{B-CMP122a}, where $k_0$ is the threshold level of the
weight formula in \eqref{4.27:eq-weights-heredity-b}, meets every completed component $X_c$,
namely in its collar $X_c\cap\Omega_k$, which is non-empty \cite[\S3]{B-CMP119}, and
$\Omega_k\subset\Omega_{k_0+1}$; and the class
$\mathfrak{D}_m$, the connected unions of cubes of the scale-$m$ cover, holds a domain through any
given $y$ and any cell of the band, and one containing any given cube of $\mathcal{Q}(y)$ (this
property of $\mathfrak{D}_m$ is the step referred to at $(\ast)$ below). These
terms are those of $\mathbb{B}''_k=\mathbb{B}^{(N)}_k$ \cite[\S1]{B-CMP122a}, are kept across the
reset of the determining set at the completion \cite[(1.33)]{B-CMP122b}, and are presented with the
pair by the first part of \cite[(1.65)]{B-CMP122b}. The class part $\mathfrak{L}_N$ of the domain of
$\mathfrak{E}$ has its thresholds at $y$ taken from $\mathbb{B}''_k=\mathfrak{B}^t$, with factor
\begin{equation*}
\varphi^{(N),(k)}_{\ell_{\mathfrak{B}^t}(y)}\le s^{(k)}_{\ell_{\mathfrak{B}^t}(y)}\le1,
\end{equation*}
$s^{(k)}_m$ being the bound on the stage factors fixed by \eqref{4.27:eq-arrested-spaces-a}; on a
live piece the factor is smaller by the range sums $\sigma_{\mathfrak{a}}$ and still at least
$\varphi_{\mathrm{arrest}}>0$ (Lemma~\ref{4.27:lem-factor-positivity}), and a smaller positive
threshold only strengthens the bound. Hence each entry is at most $\tfrac1{16}\theta^{\mathfrak{B}^t}_a$, for
$a=1,\dots,8$, the range of entries of the consumer's clause. At $a=8$ the clause of
$\mathfrak{L}_N$ is that of \cite[(1.65), (1.67), (1.69)]{B-CMP122a} at stage $N$ on the cubes of
$\mathcal{Q}_{\mathbb{B}''_k}(y)=\mathcal{Q}_{\mathfrak{B}^t}(y)$, at the family's weight and
without factor: on each such cube $\square$, of level $\lambda$ and family weight $w_\square$, there
is a gauge $u$ on $\square\cap X$ with $U^u=e^{i\eta A}$ and
\begin{equation*}
L^\lambda\eta|A|,\ (L^\lambda\eta)^2|\nabla^\eta A|<\tfrac1{16}\,w_\square BCM\alpha_{0,\lambda},
\end{equation*}
$U$ being the $G$-part of the one presented $\mathbb{U}$; this is the gauge that
$\mathfrak{G}(\mathfrak{B}^t)$ asks for, on all of $\square$ once $X\supseteq\square$. So the cube
clause holds with $\tfrac1{16}\theta^{\mathfrak{B}^t}_8$ on every cube of
$\mathcal{Q}_{\mathfrak{B}^t}$ sharing a cell with $\mathsf{Z}\cap Z^{\mathrm{on}}$.

At $y\in\mathsf{Z}\cap\mathfrak{T}_k$, with $\mathfrak{T}_k$ taken at the density of the completing
chain, each entry is at most $1.1\,r_a$ in the units of the structure in force at $y$, by the bound
on the presented pair at the
completion site. These units are those of $\mathfrak{B}^t$: on a live piece of another class
in $Z^c$ they are the frozen ones, by \eqref{4.27:eq-arrested-spaces-a}; the chain raises no point
of
$\mathfrak{T}_k$, by the third family of the stage structure of \cite[(1.65), (1.67),
(1.69)]{B-CMP122a} at $m=k$ and because the clauses of the foreign part $\mathfrak{I}_n$ of a stage
domain on $\mathfrak{T}_k$ are in the native units (the inheritance statements for successor
stages); and no
dissolving piece meets $\mathfrak{T}_k=\Omega_k\cup Z^c$, since $W_{\mathfrak{c}}\subset Z\cap Z_h$,
where $h:=k-N$ with $N=N_k$ the number of stages of the step and $Z_h$ is the large-field region at
density $h$,
and $Z_h\cap\Omega_k=\emptyset$ by \cite[(2.1)]{B-CMP119}. So every entry covered by the bound on
the presented pair at the completion site, the cube clause at $\mathcal{Q}_{\mathfrak{B}^t}(y)$
included, is at most $1.1\,\theta^{\mathfrak{B}^t}_a$.

\emph{Maps: the core.} On the core $D_2\setminus B_S$ of the completion site
($S:=\partial\Omega''^{\sim}_{h+1}$, and $B_S$ the set near $S$ that the base bound on the core
excludes), which holds every
piece dissolved at $\mathfrak{O}$, the presented pair is within $\tfrac1{64}+\tfrac3{64}=\tfrac1{16}$
of every threshold $F\cdot w\cdot c\cdot r_e$ whose level is the current level
$o(x)=\ell_{\mathfrak{B}^t}(x)$ of $x$ or a finer birth index, with $F\ge\varphi_\flat$ and
$w,c\ge1$, the cube clause included: the base bound on the core gives $\tfrac1{64}$, read with its
gauges, and the three increments of the increment bound on the core give $\tfrac3{64}$, each a
layer on the re-presented base field $U_2$ on the
cubes of the base. Taking $F=\varphi_\flat$, $w=w_{\mathfrak{B}^t}(x)$ and $c=1$ is admissible,
since the frozen weight $w_{\mathfrak{B}^t}(x)=\mathsf{L}_0^{2(\lambda-\max\{\mu,k_0\})^+}$ of (e1)
in \eqref{4.27:eq-weights-heredity-b} is at least $1$ (it is bounded above by
Lemma~\ref{4.27:prf-weight-bounds}). As $\varphi_\flat\le1$,
\begin{equation*}
\varphi_\flat\,w_{\mathfrak{B}^t}(x)\,r_e(o(x))\le\theta^{\mathfrak{B}^t}_e(x),
\end{equation*}
and each entry is at most $\tfrac1{16}\theta^{\mathfrak{B}^t}_e$.

The presented pair at $(\mathrm{L})$ is one configuration on $\Omega_0$
\cite[(1.50), (1.53), (1.65)]{B-CMP122b}, and the bounds above control it region by region in one
decomposition $e^{i\eta A'}U$. Hence a cube of $\mathcal{Q}_{\mathfrak{B}^t}$ lying across two of
the regions is evaluated on that one $U$ in a single gauge. A cube with a point
$y\in\square\cap\mathsf{Z}$ is
controlled as a whole by the bound for the consumer $\mathfrak{E}=\mathbb{C}^{(N)}_k(X)$ with
$X\supseteq\square$, at $y$, the cube
clauses being attached pointwise (the cube clauses of $\mathcal{Q}(x)$ in
\eqref{4.27:eq-arrested-spaces-a}, and ``every cube of $\mathcal{Q}_{\mathfrak{B}}$'' in the
definition of $\mathfrak{G}$). A cube with no point in $\mathsf{Z}$ is controlled by the gauge kept
by the layer bound at the completion site, through \eqref{4.27:eq-nested-cubes-a} and (e3) of
\eqref{4.27:eq-weights-heredity-b}. A cube of a frozen family of a dissolved piece, of side
$CML^{\ell_{\mathfrak{c}}}\eta$ with $\ell_{\mathfrak{c}}\le h$, lies two $LMR_{h+1}$-layers
inside $S$, hence whole in $D_2\setminus B_S$, and is controlled by the base bound on the core with
its gauges. Each cube is controlled by a single one of these bounds; none needs two of them. In
every case each entry is below $5\theta^{\mathfrak{B}^t}_a$.

\emph{Maps: off $Y'$.} Off $Y'$ the $\mathbb{U}$ of the image is that of the input: the pair
outside the domain of the map is left as it is, and the change-of-variables function $\mathbb{H}$
vanishes on a component without data (clauses (a) and (f) of the theorem on $\mathbb{H}$). Its
$\mathbb{J}$ is taken to be that of the input: only the existence of an extension is required, and
$\mathfrak{K}_{109}$ admits every source with $\sup\ell^3|\mathbb{J}|$ below the radius, so the tail
term that the recursive presentation carries off $Y'$ is not included in the extension. There the
clauses of $\tilde p$ pass to $\mathfrak{B}^t$ by (e2) of \eqref{4.27:eq-weights-heredity-b}, and
by (e3) for the cubes. Thus $\tilde q$ (the image on $Y'$; the input's $\mathbb{U}$ off $Y'$; the
input's $\mathbb{J}$ on $\Omega_1(\mathfrak{B}^t)$ off $Y'$) lies in $\mathfrak{G}_1(\mathfrak{B}^t)$:
every bound above sits at a point of $\Omega_1(\mathfrak{B}^t)$ or on a cube of
$\mathcal{Q}_{\mathfrak{B}^t}$, at the current level $o(x)\ge1$ on the core, and the raw points of
$\mathsf{Z}$ and of the core are points of the stratum, which enter no clause of $\mathfrak{G}_1$.

\emph{The stratum off $Y'$.} On $(\Omega_0\setminus\Omega_1)(\mathfrak{B}^t)$ off $Y'$, the
$\mathbb{U}$ of $\tilde q$ is that of the input and obeys the index-zero $U$-members of
$\mathfrak{G}(\mathfrak{B}^t)$, because $\mathfrak{G}(\mathfrak{B}')\subset
\mathfrak{G}^{-J_0}(\mathfrak{B}^t)$ by (e3) of \eqref{4.27:eq-weights-heredity-b}: at points raw
at $\mathfrak{B}'$ the clause, which does not depend on the radius, is the same at both dates, and
at the points of $Z_L$, the union of the components completed at $(\mathrm{L})$, entering
$\Omega_1(\mathfrak{B}')$ at $(\mathrm{L})$ it follows from the
level clauses by Lemma~\ref{4.27:prf-raw-stratum}. The $\mathbb{J}$ of $\tilde q$ is that of the
input at the raw points held at the date of $\mathfrak{B}^t$ by a live piece, and $0$ at those held
at that date by none. The choice depends on $\mathfrak{B}^t$, not on $\mathfrak{B}'$: a point first
held by
an arrest attached at the date of $\mathfrak{B}'$ itself (for instance an arrest attached to
$(k+1;\top)$ when $\mathfrak{O}$ is the $\mathbf{T}$-step $k\to k+1$,
Lemma~\ref{4.27:lem-determining-set-order}) is unheld at $\mathfrak{B}^t=\mathbb{B}_k$, where
$T^{\mathfrak{B}^t}(y)=\gamma_0^{(k)}$, while the input's $\ell^3|\mathbb{J}(y)|$ may approach
$T^{\mathfrak{B}'}(y)=\gamma_0^{(k+1)}$. With this choice $\tilde q$ obeys the $J$-member of
$\mathfrak{G}(\mathfrak{B}^t)$. At held points, by Lemma~\ref{4.27:prf-held-raw-points} applied to
$\mathfrak{B}^t\preceq\mathfrak{B}'$, $T^{\mathfrak{B}'}(y)=T^{\mathfrak{B}^t}(y)$ while $y$ is raw
at $\mathfrak{B}'$, so $\ell^3|\mathbb{J}(y)|<T^{\mathfrak{B}^t}(y)$ because
$\tilde p\in\mathfrak{G}(\mathfrak{B}')$; if $y\in\Omega_1(\mathfrak{B}')$, the bound of the same
lemma at points entering $\Omega_1$ gives the same inequality. At the points where $\mathbb{J}$ is
cut, $0<T^{\mathfrak{B}^t}(y)$.

\emph{The stratum on $Y'$.} On $Y'\cap(\Omega_0\setminus\Omega_1)(\mathfrak{B}^t)$ the index-zero
gauges of the $G$-part, of the form \cite[(3.35)]{B-CMP99b} at index zero, are kept wherever the
map is a layer into $U'$, since the map does not move the $G$-part there and the input obeys them,
$\mathfrak{G}(\mathfrak{B}')\subset\mathfrak{G}^{-J_0}(\mathfrak{B}^t)$.

$(\ast)$ The moved members on $Y'\cap(\Omega_0\setminus\Omega_1)(\mathfrak{B}^t)$, namely the
$A'$-member of the form \cite[(3.37)]{B-CMP99b} at index zero and the $J$-member, and, at the site
$(\mathrm{L})$, every index-zero member of the re-presented pair at the raw points of $\mathsf{Z}$
and of the core, are not derived by the argument above. The existence of a consumer
$\mathbb{C}^{(N)}_k(X)$ through every point of $\mathsf{Z}\cap Z^{\mathrm{on}}$ used above rests on
a property of the class $\mathfrak{D}_m$ of \cite{B-CMP122a}, the class over which
\cite[(1.63)]{B-CMP122a} is indexed, which is not established by the argument above.
\end{proof}

\begin{lemma}[The real minimiser lies in arrested spaces. Stated; proof incomplete at the step
marked $(\ast)$]\label{4.27:lem-real-minimiser}
Let $\mathsf{m}\in\,]0,\tfrac12]$. Assume the following hypotheses. The first three are among the
hypotheses of the arrest theorem collected in Definition~\ref{4.27:def-arrest-hypotheses}. The
next two together form the margin condition at $\mathsf{m}$, which is among the members displayed
in \eqref{4.27:eq-arrest-hypotheses-a}. The last item consists of the statements from which the
clauses at points of level $\ge1$ are derived.
\begin{itemize}
\item The radius condition
\begin{equation}\label{4.27:eq-real-minimiser-1}
C_0^{\mathrm{rad}}\ge 2.5\,\frac{A_0}{\varphi_{\mathrm{arrest}}},\qquad q>p_0 .
\end{equation}
\item The flow bounds of Definition~\ref{4.27:def-arrest-hypotheses}.
\item The size condition on $N_k$, namely $N_k>N_0(k)$.
\item The first member of the margin condition:
\begin{equation}\label{4.27:eq-real-minimiser-2}
1.08\,\frac{A_0}{C_0^{\mathrm{rad}}}\,\mathsf{t}_\gamma^{\,p_0-q}\le \mathsf{m}\,\varphi_{\mathrm{arrest}}.
\end{equation}
\item The second member of the margin condition: the floor of the real-point lemma with excluded
attempts, with its last member taken at the margin $\mathsf{m}$:
\begin{equation}\label{4.27:eq-real-minimiser-3}
K^{\mathrm{own}}\,\varepsilon^{\mathsf{l}}\le \mathsf{m}\,\varphi_\flat\,\alpha_0 ,
\end{equation}
where $K^{\mathrm{own}}$ is the constant of that floor, $\alpha_0$ is the radius entering that
floor, and $\varphi_\flat=\min\{1-\beta,1-4.51\beta\}$ is the factor floor of that lemma, which
equals $\varphi_{\mathrm{arrest}}=1-4.51\beta$ of Lemma~\ref{4.27:lem-factor-positivity}, since
$\beta>0$.
\item The statements on the real point: the two statements on the real point away from arrested
pieces, and the real-point lemma with excluded attempts together with its statement on refinement
of the determining set (criticality and uniqueness of the minimiser for a refined problem), whose
conclusion holds in its setting under its conditions on the geometry, the data and the flow (these
conditions are established in the proof below).
\end{itemize}
Here $\varphi_{\mathrm{arrest}}$ is the floor of the arrested factors, \eqref{4.27:eq-factor-positivity-a};
$C_0^{\mathrm{rad}}$ is the radius constant bounded below in the floors of
Definition~\ref{4.27:def-arrest-hypotheses}, and $\mathsf{t}_\gamma=\log\gamma^{-2}$.

Let $V$ be a real configuration in the support of the characteristic functions of the effective
density. Let $\mathfrak{B}$ be the current determining set: either the determining set
$\mathbb{B}_{k'}$ of the current step $k'$, or the stage set of a later attempt. Write
$U=U(\mathfrak{B},V)$ for the real
minimiser determined by $\mathfrak{B}$ and $V$, and $\mathcal{J}(U)$ for the source $1$-form
attached to it.

Then the pair $(U,\mathcal{J}(U))$ lies in every arrested space on $\mathfrak{B}$
(Definition~\ref{4.27:def-arrested-spaces}) in the sense of clauses (a)--(c) below. Moreover, for
every live arrest $\mathfrak{a}$ and
at every triple of the family of triples attached to $\mathfrak{a}$, written
$(\Sigma^{*})^{\mathfrak{a}}$, each entry is at most $\mathsf{m}$ times its threshold. For the
entries $Q_1$ and $Q_5$ this holds with $\mathsf{m}=\tfrac12$ under
\eqref{4.27:eq-real-minimiser-1} alone. Clause by clause, the assertion is the following.
\begin{enumerate}
\item[(a)] At every point $x$ of level $\ge1$, each lettered entry $Q_a(x)$, $a=1,\dots,7$, and
each cube clause of $\mathcal{Q}(x)$ is at most $\mathsf{m}$ times its threshold.
\item[(b)] The arrested spaces carry a rider, by the rider clause of
\eqref{4.27:eq-arrested-spaces-a}. At points of level $\ge1$ the real point meets this rider as
the level-$\ge1$ rider $\mathfrak{G}_1(\mathfrak{B})$.
\item[(c)] At raw points $y$ of index $0$ of $\Omega_0(\mathfrak{B})$ the pair meets the member
$(3.37)_0$ of the rider, the index-$0$ instance of \cite[(3.37)]{B-CMP99b}, that is, the member
$A'=0$. It also meets, at such points, the two remaining members of the rider:
  \begin{itemize}
  \item the index-$0$ gauge condition $(3.35)_0$ of \cite[(3.35)]{B-CMP99b};
  \item the bound $\eta^{3}|\mathcal{J}(U)(y)|<T^{\mathfrak{B}}(y)$ on the source, in lattice
  units.
  \end{itemize}
  The proof of these two members is the step marked $(\ast)$, at which the proof below is
  incomplete.
\end{enumerate}
Clause (a) and the bound at every triple are derived from the hypotheses above together with
\cite[(1.24)]{B-CMP122a}, \cite[Proposition~2]{B-CMP98},
Lemmas~\ref{4.27:lem-arrest-timing} and~\ref{4.27:lem-factor-positivity}, and part (b) of
Lemma~\ref{4.27:lem-weights-heredity}, whose proof is incomplete at one step. Clause (b) rests on
clause (e4) of Lemma~\ref{4.27:lem-weights-heredity}, whose proof is incomplete at one step, and
on clause (B) of the arrest theorem.
\end{lemma}

\begin{proof}
The clauses at points of level $\ge1$ other than the rider are derived from four inputs: the
statements on the real point, \cite[(1.24)]{B-CMP122a}, \cite[Proposition~2]{B-CMP98}, and the flow
bounds. The rider at raw points of index $0$ is treated at the end, at the step marked
$(\ast)$.

\emph{Points off the arrested pieces.} Let a point lie in the held set of no live arrest. There
the assertion is the content of the two statements on the real point away from arrested pieces,
among the hypotheses above.

\emph{Points of a held set.} Let $\mathfrak{a}$ be a live arrest of density
$k_{\mathfrak{a}}$, and let $x\in W_{\mathfrak{a}}$. Put $\ell:=\ell_{\mathfrak{a}}(x)$, the
level at $x$, and let $w_{\mathfrak{a}}(x)$ be the weight at $x$, as in
Lemma~\ref{4.27:lem-weights-heredity}, whose proof is incomplete at one step. Throughout,
$\beta$ is the constant of the floor $\varphi_{\mathrm{arrest}}=1-4.51\beta=0.098$ in
\eqref{4.27:eq-factor-positivity-a}, so that $\beta=0.2$; the same constant enters the factor
$3(1-\beta/2)$ below.

\emph{Entries $6$ and $7$.} These entries are formed from $A'$. At the real point $A'=0$, so both
entries vanish.

\emph{Entries $1$ and $5$.} Let $\chi^{(n^{+})}_{k_{\mathfrak{a}}}$ be the characteristic
functions of Ba{\l}aban's construction at density $k_{\mathfrak{a}}$ that stand at $x$, frozen
with the arrest; $n^{+}$ is their stage.

On the processed shells $\chi^{(n^{+})}$ takes the value $1$ at $V$, since $V$ lies in the
support of the characteristic functions. The bound
\cite[(1.24)]{B-CMP122a} at stage $n^{+}$ then gives, for Ba{\l}aban's background field
computed from the restriction $V\restriction X_{\mathfrak{a}}$,
\begin{equation}\label{4.27:eq-real-minimiser-4}
\bigl|\partial U^{(n^{+})}_{k_{\mathfrak{a}},Z}(V\restriction X_{\mathfrak{a}})(x)-1\bigr|
<w'\,\varepsilon_\ell\,L^{-2\ell},\qquad w'\le w_{\mathfrak{a}}(x).
\end{equation}
The inequality $w'\le w_{\mathfrak{a}}(x)$ holds because the powers of $\mathsf{L}_0$ carried by
\cite[(1.24)]{B-CMP122a} are at most those of the weight $w_{\mathfrak{a}}(x)$. At the last stage
$n^{+}=N_{k_{\mathfrak{a}}}$ and at layer $m=k_{\mathfrak{a}}-1$,
\cite{B-CMP122a} carries the factor $3(1-\beta/2)$, and there
\begin{equation*}
3(1-\beta/2)=2.7\le\mathsf{L}_0^{2}=w_{\mathfrak{a}}(x),
\end{equation*}
because $\beta=0.2$ and $\mathsf{L}_0\ge2$.

Consider next the shells lying in $\Omega_{j^{+}+1}$, on which \cite[(1.24)]{B-CMP122a} at stage
$n^{+}$ carries no clause ($j^{+}$ being the level index of that bound at stage $n^{+}$); these are
the shells not processed by stage $n^{+}$. Consider also the points at which the step keeps the
bounds \cite[(1.3), (1.4)]{B-CMP122a}. On these shells and at these points those bounds
give \eqref{4.27:eq-real-minimiser-4} with $w'=1$.

The real minimiser $U$ has, on $X_{\mathfrak{a}}$, the same determining set and the same data
as this background field. Indeed, by clause (T3) of Lemma~\ref{4.27:lem-arrest-timing},
displayed in \eqref{4.27:eq-arrest-timing-a}, no later attempt re-levels or integrates a point
of $X_{\mathfrak{a}}$. Further, the collar $\mathfrak{c}_{\mathfrak{a}}$ keeps the level
$k_{\mathfrak{a}}$, since $\Omega_{k_{\mathfrak{a}}+1}\subset\Lambda_{k_{\mathfrak{a}}}$.

Near $\partial X_{\mathfrak{a}}$ the real minimiser is regular, by the small-field
characteristic functions kept on the adjacent strata. Therefore, as in
\cite[(1.30), (1.56)]{B-CMP122a}, $U$ is the product of $e^{i\eta\mathbb{H}}$ with the background
field of \eqref{4.27:eq-real-minimiser-4}. The function $\mathbb{H}$ depends on data supported in
the layer of width $2M_1$ at $\partial X_{\mathfrak{a}}$. That layer lies two layers of
$MR_{k_{\mathfrak{a}}}$-cubes away from $W_{\mathfrak{a}}$, and at that distance from its data
the estimate on $\mathbb{H}$ shows that the factor $e^{i\eta\mathbb{H}}$
changes $|\partial U(x)-1|$ by at most
$(\beta/10)\,w_{\mathfrak{a}}(x)\,\varepsilon_\ell\,L^{-2\ell}$. Since $1+\beta/10=1.02$, this gives
\begin{equation*}
|\partial U(x)-1|\le 1.02\,w_{\mathfrak{a}}(x)\,\varepsilon_\ell\,L^{-2\ell}.
\end{equation*}

By Lemma~\ref{4.27:lem-factor-positivity}, the threshold of these entries is at least
$\varphi_{\mathrm{arrest}}\,w_{\mathfrak{a}}(x)\,\alpha_{0,\ell}\,L^{-2\ell}$. The constant $1.08$ of
\eqref{4.27:eq-real-minimiser-2} exceeds $1.02$. Replacing $1.02$ by $1.08$ and using
$\varepsilon_\ell/\alpha_{0,\ell}=(A_0/C_0^{\mathrm{rad}})\,\mathsf{t}_\ell^{\,p_0-q}$, we obtain
\begin{equation}\label{4.27:eq-real-minimiser-5}
\frac{|\partial U(x)-1|}{\varphi_{\mathrm{arrest}}\,w_{\mathfrak{a}}(x)\,\alpha_{0,\ell}\,L^{-2\ell}}
\le\frac{1.08\,\varepsilon_\ell}{\varphi_{\mathrm{arrest}}\,\alpha_{0,\ell}}
=\frac{1.08}{\varphi_{\mathrm{arrest}}}\,\frac{A_0}{C_0^{\mathrm{rad}}}\,\mathsf{t}_\ell^{\,p_0-q}.
\end{equation}
At $\mathsf{m}=\tfrac12$ the right side of \eqref{4.27:eq-real-minimiser-5} is at most
$\tfrac12$ provided $C_0^{\mathrm{rad}}\ge 2.16\,A_0/\varphi_{\mathrm{arrest}}$. Condition
\eqref{4.27:eq-real-minimiser-1} provides this, since $2.5\ge2.16$. At a general margin
$\mathsf{m}$, the right side of \eqref{4.27:eq-real-minimiser-5} is at most $\mathsf{m}$. Indeed,
the hypotheses of Definition~\ref{4.27:def-arrest-hypotheses} are imposed at every value
$\mathsf{t}\ge\mathsf{t}_\gamma$ of the logarithmic coupling variable, and every running coupling
$g_\ell$ of the flow satisfies $g_\ell\le\gamma$ (Definition~\ref{4.27:def-arrest-hypotheses}),
that is, $\mathsf{t}_\ell=\log g_\ell^{-2}\ge\mathsf{t}_\gamma$. Since $q>p_0$ by
\eqref{4.27:eq-real-minimiser-1}, the function $\mathsf{t}\mapsto\mathsf{t}^{\,p_0-q}$ is
decreasing, and \eqref{4.27:eq-real-minimiser-2} at $\mathsf{t}_\gamma$ implies the same bound
at $\mathsf{t}_\ell$.

\emph{Entries $2$, $3$, $4$ and the cube clause.} Apply part (c) of the real-point lemma with
excluded attempts, at the triple $(\mathfrak{B},V,\hat k)$, where $\hat k$ is the top level of
$\mathfrak{B}$. Take the fraction $\mathsf{m}$: this is the fraction in the last member of that
lemma's floor, here \eqref{4.27:eq-real-minimiser-3}. The base statement at completion sites
applies the same part (c) with the fraction $1/64$. The hypotheses of the real-point lemma with
excluded attempts are established at our triple in the order in which they are checked at a
completion site, as follows.

($\alpha$) \emph{The set.} The strata of $\mathfrak{B}$ are of four kinds:
\begin{itemize}
\item strata made by $\mathbf{T}$;
\item strata rebuilt by completed $\mathbf{R}$-operations, at level $k''$ on $X''$, by
\cite[(1.33)]{B-CMP122b};
\item the frozen strata of live arrests;
\item for a stage set, the strata of the treated site.
\end{itemize}
The frozen strata are given by the branches in the definition of an arrest, and these are
exactly the exclusions from the large-field region $Z$ that the setting of the real-point lemma
with excluded attempts admits. The strata of the treated site are admitted by that setting. No
stratum of $\mathfrak{B}$ is of any other kind: by the dissolution clause of
\eqref{4.27:eq-arrested-spaces-a}, pieces of dissolved arrests are rebuilt.

($\beta$) \emph{The geometry.} The nesting condition of the real-point lemma with excluded
attempts is the nesting of the arrested and treated sites. The proof of its three conditions on
thin links (in the terminology of that lemma's geometry conditions: the thin links between nested
sites and the blocks of levels containing them) uses the following:
\begin{itemize}
\item the count $4+1+1+2$ entering Ba{\l}aban's choice of $N_0$ and the estimate on the
function $\mathbb{H}$ used above;
\item the condition of that proof on nested sites which, read literally at nested sites, is
clause (T2) of \eqref{4.27:eq-arrest-timing-a};
\item the size condition on $N_k$, $N_k>N_0(k)$.
\end{itemize}
Thus the thin links lie in one block per arrested or treated site. For each such site, of level
$k''$ say, its block lies at levels in $[k''-N_0(k''),k''-1]$; the blocks are disjoint, and are
wide at the top.

($\gamma$) \emph{The data, on the support of the characteristic functions.}
\begin{itemize}
\item On the stratum made by $\mathbf{T}$ at level $n$, the data are of size $O(L^2)\varepsilon_n$,
by the regularity $|\partial V-1|<O(L^2)\varepsilon$ assumed of the incoming field in
\cite{B-CMP119}.
\item On strata rebuilt by completed $\mathbf{R}$-operations, they are of size
$K_J\varepsilon_{k''}$, where $K_J$ is the constant of the bound at completed exits of
$\mathbf{R}$.
\item On a block raised by $\nu$ levels, at a point of the block of level $m$,
\cite[(1.24)]{B-CMP122a} and \cite[(54)]{B-CMP98} give, on the plaquettes of the block, the bound
\begin{equation*}
|\partial V-1|\le 6\,\mathsf{L}_0^{2\nu}\varepsilon_m\le K^{\mathrm{reg}}\,w_{*}(k_{\mathfrak{a}})\,
\varepsilon_m ,
\end{equation*}
with $2\nu$ the exponent of $\mathsf{L}_0$ that \cite[(1.24)]{B-CMP122a} carries on the block,
and
$K^{\mathrm{reg}}$ the regularity constant of the data clause of the real-point lemma with
excluded attempts. This bound is local, by \cite[Proposition~2]{B-CMP98}. The smallness
hypothesis \cite[(52)]{B-CMP98} of that proposition, taken with radius $3\varepsilon^{\mathsf{l}}$,
requires this radius to be at most $c_2'$; this is the member $3\varepsilon^{\mathsf{l}}\le c_2'$
of the floor of the real-point lemma with excluded attempts. The characteristic function in force
on the block is $\chi^{(n^{+})}_{k_{\mathfrak{a}}}$, and the subsequent operations do not change
it; at a treated site it is that site's own $\chi^{(n')}$.
\item On shells not yet processed, the bounds are \cite[(1.3), (1.4)]{B-CMP122a}.
\end{itemize}

($\delta$) \emph{The size.} By Lemma~\ref{4.27:lem-weights-heredity}(b), whose proof is
incomplete at one step, the data are $K^{\mathrm{reg}}\varepsilon^{\mathsf{l}}_n$-regular on every
stratum.

($\varepsilon$) \emph{The flow.} These are the flow bounds, among the hypotheses above.

The clause at $x$ has three properties. It is evaluated at the current level of $x$, by the
structure freeze of \eqref{4.27:eq-arrested-spaces-a}. Its weight satisfies
$w_{\mathfrak{a}}(x)\ge1$. Its
factor is at least $\varphi_\flat$, which equals the floor $\varphi_{\mathrm{arrest}}$ of
\eqref{4.27:eq-factor-positivity-a}, by Lemma~\ref{4.27:lem-factor-positivity}.
Hence part (c) of the real-point lemma with excluded attempts applies, and entries $2$, $3$, $4$
are at most $\mathsf{m}$ times their thresholds. The cube clause holds with the gauges of part (a)
of that lemma.

\emph{Predecessor triples.} Let a triple of $(\Sigma^{*})^{\mathfrak{a}}$ be a predecessor
triple on a domain $X'$. Let $\mathbb{B}_p(X')$ be the determining set of the predecessor on
$X'$, and $\Gamma^{\mathfrak{s}}_i$, $i<p$, its earlier strata. Stages only coarsen, so the family
$\mathbb{B}_p(X')\cup\{\Gamma^{\mathfrak{s}}_i\}_{i<p}$ refines $\mathfrak{B}$ on $X'$.

By the statement on refinement of the determining set that accompanies the real-point lemma with
excluded attempts, among the hypotheses above, the
restriction $U\restriction X'$ is critical for the refined variational problem at its own
averages, in the space on which that problem is posed, and the critical point there is unique.
Let $U_{p,X'}(M^{\mathfrak{s}}U)$ be the minimiser of the refined problem at the averages
$M^{\mathfrak{s}}U$ of $U$, and $\mathbb{J}_{p,X'}$ its source. Hence
$U_{p,X'}(M^{\mathfrak{s}}U)$ is a gauge copy of $U\restriction X'$, and $\mathbb{J}_{p,X'}$ is
a gauge copy of $\mathcal{J}(U)$.

Therefore entries $2$ and $4$ carry the right members of entries $1$ and $3$ respectively. This
rests on the identities, with the exponents $n$ and $p$ of the two triples,
\begin{equation*}
L^{-2p}\,(L^{n}L^{-p})^{-2}=L^{-2n},\qquad L^{\,n-\min\{p,n\}}\ge1 .
\end{equation*}
No tail terms occur.

The appeals above are not circular: the base statement at completion sites, invoked in the step
for entries $2$, $3$, $4$, appeals, among the clauses of the arrested spaces, only to the first
three clauses of \eqref{4.27:eq-arrested-spaces-a}, beginning with the structure freeze, and to
Lemma~\ref{4.27:lem-factor-positivity}.

\emph{The rider at levels $\ge1$.} By clause (e4) of Lemma~\ref{4.27:lem-weights-heredity}, whose
proof is incomplete at one step, and by clause (B) of the arrest theorem, the real point meets the
rider of the rider clause of \eqref{4.27:eq-arrested-spaces-a} as $\mathfrak{G}_1(\mathfrak{B})$.

$(\ast)$ \emph{Raw points of index $0$.} Let $y$ be a raw point of index $0$ of
$\Omega_0(\mathfrak{B})$. The member $(3.37)_0$ of \cite[(3.37)]{B-CMP99b} holds at $y$, because
$A'=0$ at the real point.
The argument above does not establish the remaining two members of the rider at $y$: the gauge
condition $(3.35)_0$ of \cite[(3.35)]{B-CMP99b}, and the bound
$\eta^{3}|\mathcal{J}(U)(y)|<T^{\mathfrak{B}}(y)$. At this step the proof is incomplete.
\end{proof}

\begin{lemma}[Later increments at a live piece]\label{4.27:lem-live-piece-increments}
Let $\mathfrak{a}$ be a live arrest, let $x\in W_{\mathfrak{a}}$, and let $\ell:=\ell(x)$ be the level that
the spaces read at $x$ under the arrested reading (Definition~\ref{4.27:def-arrested-spaces}).
\begin{enumerate}
\item[(G)] The determining set of the operation under consideration, frozen on live pieces, is
admissible. Its levels are consecutive across its strata, and each stratum is at least one
$M$-cube of its own scale wide. A contour from a cell of level $\lambda$ to $x$ has
\begin{equation*}
d^{\wedge}\ \ge\ \mathsf{w}_M\,(\lambda-\ell-1)^{+}.
\end{equation*}
\end{enumerate}
Under the later operations of Definition~\ref{4.27:def-operations-zones}, entries $5$ and $8$
of the clause read at $x$ do not move. Every other entry moves linearly in the layer and, measured
in the units of the level $\ell$, at every triple, by at most the following amounts.
\begin{itemize}
\item Under \textup{(O1)} and \textup{(O4)}, at stage $l$: $\tfrac12\rho^{(l)}_{\ell}$.
\item Under \textup{(O4$'$)}, where the later attempt is at density $k_b$:
\begin{equation}\label{4.27:eq-live-piece-increments-1}
\sum_{l}\iota_l(x)\ \le\ \tfrac12\,\mathsf{w}_0\,\beta\,e^{-(k_b-1-\ell)}.
\end{equation}
\item Under \textup{(O2)} at step $k$: $\bar P\,2^{-(k-1-\ell)^{+}}$.
\item Under \textup{(O3)}, the $\mathbf{T}$-step $k'\to k'+1$: $\tfrac14(1-\theta_S)\beta\,2^{-(k'+1-\ell)}$.
\end{itemize}
Here $\rho^{(l)}_{\ell}$ is the room of one stage $l$ at level $\ell$, $\iota_l(x)$ is the increment at
$x$ of stage $l$ of the later attempt, $\bar P$ is the largest of the three increment constants
$P,P',P''$ of the bound on the layers of the large-field operation off the zone, and $\theta_S$,
$\beta$ are the constants of the room schedule.
\end{lemma}

\begin{proof}
\emph{Proof of \textup{(G)}.} The structures of the determining set are chosen component by
component, each from an admissible sequence in the sense of \cite[(1.16)]{B-CMP122a}, subject
to the admissibility condition on such sequences. This gives admissibility, consecutive levels
across the strata, and a width of at least one $M$-cube of the own scale for every stratum. The
contour bound is then the contour property of determining sets: a contour pays at least
$\mathsf{w}_M$ for every crossed stratum, or $M$-cube of the local scale, so that
$d^{\wedge}(y,y')\ge \mathsf{w}_M(|m-m'|-1)^{+}$ for points $y,y'$ of levels $m,m'$. Taking $y$ in
a cell of level $\lambda$ and $y'=x$, of level $\ell$, gives the display in \textup{(G)}.

The strata need not be thick. Let $k_0^{\mathfrak{a}}$ be the index of
Lemma~\ref{4.27:lem-weights-heredity}(a), whose proof is incomplete at one step, such that a point
of $W_{\mathfrak{a}}$ whose level is at most $k_0^{\mathfrak{a}}$ carries weight one. Whenever the
intermediate domains $Z''_i$ of Ba{\l}aban's construction extend above this index, their rims
$Z''_{i+1}\setminus Z''_i$ with $i>k_0^{\mathfrak{a}}$ are exactly one $M$-cube wide, by the
admissibility condition on such sequences, and no decay at a rate proportional to the block
parameter $R_i$ of that scale per shell holds
there. The bounds invoked below nevertheless apply, because their proofs use the geometry only
through the payment of at least $2$ in the exponent, that is a factor at most $e^{-2}$, per crossed
stratum, measured in the units of the own level.

\emph{The operations \textup{(O1)} and \textup{(O4)}.} These act by the stage maps of the
large-field operation on the bonds of $W_{\mathfrak{a}}$. By the arrested reading of the spaces
(Definition~\ref{4.27:def-arrested-spaces}), the source and the target of such a map read one and
the same structure at $x$, so that every entry is compared with the same entry. The range sum
$\sigma_{\mathfrak{a}}(\ell)$ of the arrest does not depend on the stage and cancels from the
comparison, the room of one stage at $x$ is $\rho^{(l)}_{\ell}$, and the room statement for the stage
maps on the held set shows that a stage uses at most $\tfrac12\rho^{(l)}_{\ell}$. That statement is
applied at the radius $\gamma_0^{(k)}$ of the density $k$ at which the operation takes place, which
is permitted by Lemma~\ref{4.27:lem-weights-heredity}(d),
whose proof is incomplete at one step. The derived entries $2$ and $4$ are covered by the
bound $\mathsf{K}_r$.

\emph{The operation \textup{(O4$'$)}.} The bound on the new boundary terms of later attempts
states that
\begin{equation*}
\sum_{l\le n}\iota_l(y)\ \le\ \tfrac12\,\mathsf{w}_0\,\beta\,e^{-n_W(y)},
\qquad n_W(y):=k_b-1-m(y),
\end{equation*}
where $m(y)$ is the level read at $y$. Its proof uses the geometry only through the two
inequalities
\begin{equation*}
d^{\wedge}\bigl(y,\mathcal{C}^{(l)}\bigr)\ \ge\ \mathsf{w}_M\,n_W(y),
\qquad \tfrac14\,\delta\,\mathsf{w}_M\ \ge\ 1,
\end{equation*}
where $\mathcal{C}^{(l)}$ is the bond set of stage $l$. The first is a contour bound of the form
supplied by \textup{(G)}; the second follows from $\delta\mathsf{w}_M\ge 8$, which is among the
hypotheses of the arrest theorem (Definition~\ref{4.27:def-arrest-hypotheses}). At $y=x$ we have
$m(x)=\ell$, hence $n_W(x)=k_b-1-\ell$, and the bound becomes \eqref{4.27:eq-live-piece-increments-1}.

\emph{The operation \textup{(O2)}.} We use the bound on the layers of the large-field operation
off the zone: the presented pair is the identity times layers, and the increments are linear in the
layer and at most $P\,2^{-(k-1-m)}$ on foreign shells of level $m\le k-1$, at most $P'$ on
$\Omega_k\setminus Z$, and at most $P''$ on collars. Here $\bar P$ is the largest of $P,P',P''$. In
that bound the decay factor $e^{-10\delta_1 R_{\min}}$ is replaced by $e^{-\delta\mathsf{w}_M}$ per
crossed stratum \cite[(1.47)]{B-CMP122a}, and $e^{-\delta\mathsf{w}_M}\le e^{-2}$, since
$\delta\mathsf{w}_M\ge 8$ is among the hypotheses of the arrest theorem
(Definition~\ref{4.27:def-arrest-hypotheses}). With
$t:=k-1-\ell$, the $t$-dependent factor $(2+t)^{1/2}$ of the coarse-over-fine ratio
$(1+\beta_0)^2(2+t)^{1/2}$ of that bound is then controlled by
\begin{equation}\label{4.27:eq-live-piece-increments-2}
(2+t)^{1/2}e^{-2t}\ \le\ \sqrt2\cdot 2^{-t}\qquad(t\ge 0).
\end{equation}
Indeed, both sides equal $\sqrt2$ at $t=0$, and the logarithmic derivative of the left side,
$1/(2(2+t))-2\le-\tfrac74$, is smaller than that of the right side, $-\ln 2$. The factor $\sqrt2$
on the right of \eqref{4.27:eq-live-piece-increments-2} is the constant $C_{1/2}=\sqrt2$ with which
that bound is stated, so it is already contained in $\bar P$. The derived entries
are covered by the bound $K_D$. This gives $\bar P\,2^{-(k-1-\ell)^{+}}$.

\emph{The operation \textup{(O3)}.} Consider the $\mathbf{T}$-step $k'\to k'+1$, whose datum is
supported on a set $\mathcal{C}$ of bonds of level $k'$ contained in $\Omega_{k'+1}$ and has size
at most $\varpi_{k'}$, in the notation of Lemma~\ref{4.27:lem-t-step-layer}. Put, giving $t$ a new
meaning in this paragraph and below,
\begin{equation*}
t:=k'-\ell,\qquad u:=2^{-(t+1)},\qquad \Lambda_\gamma:=\mathsf{t}_\gamma^{\,p_0-q}.
\end{equation*}
Then $\Lambda_\gamma\le 1$ under the hypotheses of the arrest theorem
(Definition~\ref{4.27:def-arrest-hypotheses}), which include $q>p_0$ and the lower bound
$\mathsf{t}_\gamma\ge 2(r+q)+c_\beta$ of \eqref{4.27:eq-arrest-hypotheses-a}.

The strata of levels $\ell+1,\dots,k'$ separate $x$ from $\Omega_{k'+1}$. More precisely, in
$\mathbb{B}_{k'}$ the point $x$ is separated from $\mathcal{C}$ by the strata of levels
$\ell+1,\dots,k'-1$ and by at least one $M$-cube of the own scale of $\Omega_{k'}\setminus\Omega_{k'+1}$
\cite{B-CMP119}. By the contour property of determining sets, therefore,
\begin{equation}\label{4.27:eq-live-piece-increments-3}
d^{\wedge}(x,\mathcal{C})\ \ge\ \mathsf{w}_M\,t,
\qquad d^{\wedge}(x,\Omega_{k'+1})\ \ge\ \mathsf{w}_M\,t.
\end{equation}

Two ratios enter. The bound on the size of the $\mathbf{T}$-datum gives, with the constants
$\|C\|$, $C_1^{\mathrm{pd}}$, $A_1$ and $C_0^{\mathrm{rad}}$ of the floor for the $\mathbf{T}$-step in
\eqref{4.27:eq-arrest-hypotheses-a},
\begin{equation*}
\frac{\varpi_{k'}}{\alpha_{0,k'}}
\ \le\ \frac{2\,\|C\|\,C_1^{\mathrm{pd}}A_1}{C_0^{\mathrm{rad}}}\,\Lambda_\gamma ;
\end{equation*}
here the factor $\Lambda_\gamma$ is kept. For $t\ge1$ the ratio of radii along the coupling flow
\eqref{4.27:eq-factors-radii-coupling-a} gives
\begin{equation*}
\frac{\alpha_{0,k'}}{\alpha_{0,\ell}}\ \le\ (1+\beta_0)^2\,t^{1/2},
\end{equation*}
and at $t=0$ this ratio equals one.

\emph{Direct entries.} By the $\mathbf{T}$-step layer bound \eqref{4.27:eq-t-step-layer-a} of
Lemma~\ref{4.27:lem-t-step-layer}, whose proof is incomplete at one step, every direct entry moves
by at most $C_H^{\sharp}p(x)$ in the units of the level $\ell$, where
$p(y)=\varpi_{k'}e^{-\delta d^{\wedge}(y,\mathcal{C})}$. Write $\Delta_x$ for the
move of a direct entry at $x$. By
\eqref{4.27:eq-live-piece-increments-3} and $\delta\mathsf{w}_M\ge 8$ we have
$p(x)\le \varpi_{k'}e^{-2t}$. Define
\begin{equation*}
\mathsf{D}_t:=\Lambda_\gamma\,
\frac{2(1+\beta_0)^2\,\|C\|\,C_1^{\mathrm{pd}}A_1C_H^{\sharp}}{C_0^{\mathrm{rad}}}\,
\bigl[t^{1/2}(2e^{-2})^{t}\bigr]\,2^{-t}.
\end{equation*}
Since $(2e^{-2})^t\,2^{-t}=e^{-2t}$, the two ratios give
\begin{align*}
\frac{\Delta_x}{\alpha_{0,\ell}}
&\le C_H^{\sharp}\,\frac{\varpi_{k'}}{\alpha_{0,\ell}}\,e^{-2t}
 = C_H^{\sharp}\,\frac{\varpi_{k'}}{\alpha_{0,k'}}\,
   \frac{\alpha_{0,k'}}{\alpha_{0,\ell}}\,e^{-2t}\\
&\le \mathsf{D}_t .
\end{align*}
At $t=0$, where the ratio of radii is one, the bracket is read as $1$. For $t\ge1$ the bracket is
at most $1$, since $2e^{-2}<\tfrac13$ and $t^{1/2}\le 3^{t}$. The floor for the $\mathbf{T}$-step in
\eqref{4.27:eq-arrest-hypotheses-a} reads
\begin{equation*}
C_0^{\mathrm{rad}}\ \ge\ \frac{16(1+\beta_0)^2\,\|C\|\,C_1^{\mathrm{pd}}A_1C_H^{\sharp}}{(1-\theta_S)\beta},
\end{equation*}
so the constant factor in $\mathsf{D}_t$ is at most $\tfrac18(1-\theta_S)\beta$. Hence
\begin{equation}\label{4.27:eq-live-piece-increments-4}
\mathsf{D}_t\ \le\ \Lambda_\gamma\cdot\tfrac14(1-\theta_S)\beta\,u
\ \le\ \tfrac14(1-\theta_S)\beta\,2^{-(k'+1-\ell)} .
\end{equation}

\emph{Entries $2$ and $4$.} These entries are not derived again from the direct ones. We apply
the relative stability estimate for entries $2$ and $4$. For every layer $\mathfrak{h}$ whose local
norm is majorised by a weight $\hat p$ of the class $\mathfrak{M}^{\wedge}_{\delta/2}(\mathfrak{B})$,
with $\mathfrak{B}$ the determining set of \textup{(G)},
under its smallness condition, at any triple, at every bond $c$ of the domain of the estimate, with
$m=m(c)$ the level read at $c$, and for $e\in\{2,4\}$, it states
\begin{equation*}
Q_e^{+}(c)\ \le\ Q_e(c)\,(1+\pi_\times)+\pi\,\mathbf{r}_{e,m},
\qquad \pi:=K_S^{b}\,\frac{\hat p(x)}{\tilde\alpha_{0,m}},\qquad \pi_\times\le\pi ,
\end{equation*}
where $\tilde\alpha_{0,m}$ is the radius of the arrested space at level $m$ and $\mathbf{r}_{e,m}$ is
the unit attached to entry $e$ at level $m$ in that estimate. The quantity $K_S^{b}$ is the constant
of that estimate, independent of the bond and of the triple; its superscript is part of its name,
and it is the constant entering the condition $2.2\,K_S^{b}\,\mathsf{t}_\gamma^{\,p_0-q}\le1$ of
\eqref{4.27:eq-arrest-hypotheses-a}. At the bonds at $x$ we have
$m(c)=\ell$.
We take
\begin{equation*}
\hat p:=C_H^{\sharp}\,\varpi_{k'}\,
e^{-\frac12\delta d^{\wedge}(\cdot,\Omega_{k'+1})}.
\end{equation*}
This weight majorises the bound $N^{\Delta}_{\cdot}(\mathfrak{h})\le B_\sharp\,p$ of
\eqref{4.27:eq-t-step-layer-a} (Lemma~\ref{4.27:lem-t-step-layer}, whose proof is incomplete at
one step), and it lies in $\mathfrak{M}^{\wedge}_{\delta/2}(\mathfrak{B})$ by the
contour property of determining sets.

The smallness condition bounds a multiple of $a_\flat+\kappa_{\max}$, where $a_\flat$ is the radius
entering that condition and $\kappa_{\max}$ is the largest value of the majorant $\hat p$. Here, with
$\gamma_0=\gamma_0^{(k')}$,
\begin{equation*}
a_\flat\le\gamma_0,\qquad
\kappa_{\max}=C_H^{\sharp}\varpi_{k'}\le\alpha_{0,k'}\le\gamma_0,
\end{equation*}
so the smallness condition follows from its instance with $a_\flat$ and $\kappa_{\max}$ replaced by
$\gamma_0$, since the condition bounds a multiple of their sum by one and therefore persists when
they decrease. That instance in turn follows from the hypothesis of the arrest theorem
(Definition~\ref{4.27:def-arrest-hypotheses}) which imposes the condition with both quantities at
the ceiling $\bar\gamma_0(\mathsf{t}_{k'})$, since $\gamma_0^{(k')}\le\bar\gamma_0(\mathsf{t}_{k'})$ by
Lemma~\ref{4.27:lem-weights-heredity}(c), whose proof is incomplete at one step.

By \eqref{4.27:eq-live-piece-increments-3} and $\tfrac12\delta\mathsf{w}_M\ge4$,
\begin{equation*}
\frac{\hat p(x)}{\alpha_{0,\ell}}
\ \le\ C_H^{\sharp}\,\frac{\varpi_{k'}}{\alpha_{0,\ell}}\,e^{-4t}
\ \le\ \mathsf{D}_t .
\end{equation*}
In units, $Q_e\le 1.2$, $\pi_\times\le\pi$, and the weight at $x$ is at least one. The move of entry
$e$ is therefore at most $2.2\,\pi\le 2.2\,K_S^{b}\,\mathsf{D}_t$. By
\eqref{4.27:eq-live-piece-increments-4},
\begin{equation*}
2.2\,K_S^{b}\,\mathsf{D}_t\ \le\ \bigl(2.2\,K_S^{b}\,\mathsf{t}_\gamma^{\,p_0-q}\bigr)\,
\tfrac14(1-\theta_S)\beta\,u\ \le\ \tfrac14(1-\theta_S)\beta\,u ,
\end{equation*}
where the last step is the condition $2.2\,K_S^{b}\,\mathsf{t}_\gamma^{\,p_0-q}\le1$ of
\eqref{4.27:eq-arrest-hypotheses-a}. This is the bound
$\tfrac14(1-\theta_S)\beta\,2^{-(k'+1-\ell)}$ stated in
Lemma~\ref{4.27:lem-live-piece-increments} for the $\mathbf{T}$-step $k'\to k'+1$.

\emph{Triples and weights.} The estimates used hold for all triples simultaneously. The bound on
the layers off the zone holds for all triples at once, and the relative stability estimate holds at
any triple. The margins $\mathsf{m}^{\circ}$ of the frames of the relative stability estimate satisfy
$\mathsf{m}^{\circ}\le\mathsf{m}$, the margin of \eqref{4.27:eq-arrest-hypotheses-a}.
The argument is therefore
generic in the structure label $\mathfrak{s}$. Finally, the weight at $x$ is at least one, so dividing
by it only lowers the relative sizes, and the bounds above hold at every weight.
\end{proof}

\begin{lemma}[Terms of an interrupted change of variables]\label{4.27:lem-interrupted-terms}
Let $\mathfrak{a}$ be an arrest taking the exit~(e4) of
Definition~\ref{4.27:def-arrests-pieces}, and write $k=k_a$. At the site of the operation
$\mathbf{R}$ of step $k$ the component $X_{\mathfrak{a}}$ takes its large-field exit in stage
$n+1$, at level $j=h_a+n$, where $n\le N_k-1$ and $n^{+}=n+1$. Write $P=(\mathbb{U},\mathbb{J})$
for the pair on which the spaces are defined, and $\mathsf{B}$ for the variable $B_j$, which is
left un-integrated on $X_{\mathfrak{a}}$; the variables $A_{j'}$, $j'\ge j$, are those left
un-integrated on $X_{\mathfrak{a}}$ beside the frozen data, as listed for the exit~(e4)
in Definition~\ref{4.27:def-arrests-pieces}. The map of pairs of the stages is
the one described among the inputs of Lemma~\ref{4.27:lem-t-step-layer}; its datum is
\begin{equation}\label{4.27:eq-interrupted-terms-1}
b(\mathsf{B})=g_jC\mathsf{B}-h\tilde{D}(g_jC\mathsf{B}),
\qquad |b(\mathsf{B})|\le C_1\,g_j\sup|\mathsf{B}|,
\end{equation}
with the operators $C$, $h$, $\tilde{D}$ as used in
\cite[(1.53)--(1.59)]{B-CMP122a} (here $h$ is Ba{\l}aban's operator, not a history), and with
$C_1$ the constant of the bound on the data in that statement on the map of pairs. The new
terms of clause $(\beta3)$ below are indexed by the set of level-$j$ bonds
\begin{equation}\label{4.27:eq-interrupted-terms-2}
S^{P}=\bigl(\Omega^{c}_{j+1}\setminus Z''_{j+1}\bigr)^{(j)*}\cap X_{\mathfrak{a}}.
\end{equation}
The radius of the $\mathsf{B}$-variable is, with $\delta'_j$ the threshold of
\cite[(1.55)]{B-CMP122a},
\begin{equation}\label{4.27:eq-interrupted-terms-3}
\varrho^{B}_j=\frac{6.6\,\delta'_j}{g_j}.
\end{equation}
Let $\Phi_{\chi,\sigma}(\mathsf{b})[P]$ be that map of pairs with datum
$\mathsf{b}$, with cut-off $\chi$ and continuation parameter $\sigma$; in this lemma the cut-off
is $\chi\equiv\mathbf{1}$, and we write $\Phi(\mathsf{b})[P]=\Phi_{\mathbf{1},\sigma}(\mathsf{b})[P]$.
Let
$A(\mathsf{c},\cdot)$ be the Wilson action weighted by $\mathsf{c}$
(Definition~\ref{1.1:def-wilson-action}). Put
\begin{equation}\label{4.27:eq-interrupted-terms-4}
\begin{gathered}
\mathsf{c}=\bigl(g^{(n)}_{k}(\cdot)\bigr)^{-2}-g_j^{-2},\\
\mathsf{W}(P,\mathsf{B})=A\bigl(\mathsf{c},\Phi(b(\mathsf{B}))[P]\bigr)-A(\mathsf{c},P).
\end{gathered}
\end{equation}
The functional $\mathsf{W}$ is decomposed by the partition of unity $\{\zeta_\square\}$ of
\cite[p.~276]{B-CMP119} into pieces, called the $\mathsf{W}$-pieces; we call this decomposition,
with its grouping step and its step on the remaining pieces, the localisation of $\mathsf{W}$.
Then the following hold.
\begin{enumerate}
\item[(a)] The exponent left on $X_{\mathfrak{a}}$, as a function of $(P,A_j,\mathsf{B})$ with
$\mathsf{B}=B_j$, is the sum of three families.
\begin{itemize}
\item[$(\beta1)$] The composites
\begin{equation*}
c\bigl(S;\,\Phi_{\mathbf{1},\sigma}(b(\mathsf{B}))[P];\,\cdot\,\bigr),
\end{equation*}
where $c$ runs over the constituents of Ba{\l}aban's term classes
$\mathbb{E}_k+\mathbb{R}_k+\mathbb{B}_k+\mathbb{B}^{(n)}_k$ of \cite{B-CMP122a} (in this sum
$\mathbb{B}_k$ and $\mathbb{B}^{(n)}_k$ are terms, not determining sets), and $S$ is the first
argument of $c$, which does not depend on the pair and which the composite leaves unchanged.
They are taken in the localised elements $v_e$ into which these terms are
decomposed by the localisation lemma for the stored terms, each with a domain $D_e$ and reading
$\mathsf{B}$ only on a bond set $\mathcal{B}_e$. Elements near one another are grouped as in
the grouping step of the localisation of $\mathsf{W}$, and each group is taken together with
the $\mathsf{W}$-pieces matched to it there.
\item[$(\beta2)$] The functional
\begin{equation*}
-\tfrac12\langle\mathsf{B},C^{*}\Delta^{(j)}C\mathsf{B}\rangle
+\mathbb{P}^{(j)}(g_j,A_j,\mathsf{B}),
\end{equation*}
where $\Delta^{(j)}$ and $\mathbb{P}^{(j)}$ are Ba{\l}aban's operator and term of these names
in \cite{B-CMP122a}, together with the $\mathsf{W}$-pieces not matched to a group in
$(\beta1)$, which the step on the remaining pieces of the localisation of $\mathsf{W}$ leaves
over. Such pieces occur
because the support of $\mathsf{c}$ is not contained in $X_{\mathfrak{a}}$.
\item[$(\beta3)$] The usual new terms $f_b$, $b\in S^{P}$, of the change of variables. Here
$f_b$ is the sum of the trace of the logarithm of the Jacobian of \cite[(1.22)]{B-CMP119} at $b$
and of the $b$-th increment of the sum of $(\beta1)$ and $(\beta2)$, $\mathsf{W}$ included, under
the shift $\delta_b$. The change of variables is made with the function $\mathbb{H}^{(j)}_Z$ of
\cite[(1.58)]{B-CMP122a} in place of the function $\mathbb{H}_{1,\Lambda_1}$ with which it is
made in \cite[(1.22)]{B-CMP119}.
\end{itemize}
Set
$\mathfrak{F}^{(\mathrm{e}4)}(\mathfrak{a})=(\beta1)\cup(\beta2)\cup(\beta3)$. There is no
further family, and no member is a boundary term $\mathbb{C}^{(i)}_k$ of a stage $i$. At every
rim level of every exit, that is, at every level $j'$ at which the variables $A_{j'}$ are left
un-integrated beside the frozen data (listed for each exit in
Definition~\ref{4.27:def-arrests-pieces}), the weight in the variables $A_{j'}$ is that of
\cite[(2.21)]{B-CMP119}. This weight is not a frozen term of the arrest: it is not a member of
the family $\mathfrak{F}(\mathfrak{a})$ of frozen terms of
Definition~\ref{4.27:def-arrests-pieces}.
\item[(b)] The members of $(\beta1)$ and $(\beta2)$ are holomorphic in $(P,\mathsf{B})$ on
$\mathfrak{D}^{\flat}_{n+1}(\cdot)\times\{|\mathsf{B}|<\varrho^{B}_j\}$. The members of
$(\beta3)$ are holomorphic in $P$ at real $\mathsf{B}$ with $|\mathsf{B}|<6.3\,\delta'_j/g_j$,
and in the lifted variables $(\mathsf{B},\mathsf{s},\mathsf{v})$ of the lift of the cut-off
dependence. A composite is bounded by the supremum of the constituent from which it is formed.
Consequently the domain of every member on $W_{\mathfrak{a}}$ is the class domain of stage
$n^{+}=n+1$. Thus Lemma~\ref{4.27:lem-frozen-room}(a) applies to it with $n=n^{+}$. The
domain of the members of $(\beta3)$ is the domain $\widehat{\mathfrak{D}}^{\flat}_{n+1}(\cdot)$
of stage $n+1$ read with the ambient membership condition on $W_{\mathfrak{a}}$ prescribed by
the arrested reading of the spaces (Definition~\ref{4.27:def-arrested-spaces}).
\item[(b$'$)] Let $\mathfrak{K}_{109}(\mathfrak{B};\gamma_0)$ be the space of pairs of the type
of \cite{B-CMP109} at the determining set $\mathfrak{B}$ in force on $X_{\mathfrak{a}}$ and at
radius $\gamma_0=\gamma_0^{(k)}$. Then $\mathsf{W}$ is holomorphic on
$\mathfrak{K}_{109}(\mathfrak{B};\gamma_0)\times\{|\mathsf{B}|<\varrho^{B}_j\}$. It carries no
domain condition of its own, and no domain condition is imposed on it as a frozen term.
\item[(c)] For $b\in S^{P}$,
\begin{equation}\label{4.27:eq-interrupted-terms-5}
|f_b|\le C^{P}_{22}\,(1+G_1)\,\frac{A_0}{A_1}\,\bigl(\varrho^{B}_j\bigr)^{2}e^{-R_j},
\end{equation}
where $C^{P}_{22}$ is the constant and $G_1$ the factor of the estimate on the new terms stated
among the inputs of Lemma~\ref{4.27:lem-t-step-layer}, and $R_j$ is the exponent in
Ba{\l}aban's bound, a constant times $e^{-R_j}\varepsilon_j$ with $\varepsilon_j$ the small
parameter of \cite{B-CMP122a} at level $j$, on the function $\mathbb{H}^{(j)}_Z$ of
\cite[(1.58)]{B-CMP122a}, stated in \cite[after~(1.55)]{B-CMP122a}; and $B_j$
is a weak
chart coordinate at radius $\varrho^{B}_j$, in the sense of the statement on weak chart
coordinates among the same inputs.
\item[(d)] This lemma makes no assertion on how the members of
$\mathfrak{F}^{(\mathrm{e}4)}(\mathfrak{a})$ are filed or on their total norm. These are
treated in the filing lemmas for the large-field terms. Of
$\mathfrak{F}^{(\mathrm{e}4)}(\mathfrak{a})$
the arrest theorem uses only the domains given in clause (b).
\end{enumerate}
\end{lemma}

\begin{proof}
Every appeal in this lemma and its proof to Lemma~\ref{4.27:lem-t-step-layer} is to a statement
among its inputs; that lemma is stated, and its proof is incomplete at the step marked $(\ast)$
there.

\emph{Clause (a).} Among the inputs of Lemma~\ref{4.27:lem-t-step-layer} (proof incomplete) is
the description, exit by exit, of the exponent and of the pair-dependent objects on an excluded
component. We apply that description to the exit of
stage $n+1$. It reads this exit in the form of \cite[(1.53)--(1.59)]{B-CMP122a}. The term with
the coupling difference
$\mathsf{c}$ of \eqref{4.27:eq-interrupted-terms-4} is put into the interaction, which gives
$\mathsf{W}$. The elements $v_e$ of $(\beta1)$ are those of the localisation lemma for the
stored terms; the grouping of near elements with their matching $\mathsf{W}$-pieces is the
grouping step of the localisation of $\mathsf{W}$, and the unmatched pieces of $(\beta2)$ are
those of its step on the remaining pieces. The change of variables of \cite[(1.22)]{B-CMP119},
made with $\mathbb{H}^{(j)}_Z$ in place of $\mathbb{H}_{1,\Lambda_1}$, gives the usual new terms
$(\beta3)$. The same description lists no further term on $X_{\mathfrak{a}}$. The boundary terms
$\mathbb{C}^{(i)}_k$, $i\le n$, occur only as constituents of the composites $(\beta1)$ and are
not members on their own (Definition~\ref{4.27:def-arrests-pieces}).

The rim weight is not a frozen term, by the statement on the rim form and the operator
$\Delta^{(j)}$ among the inputs of
the same lemma (proof incomplete). Its dependence on the background is
the subject of the statement on the background dependence of the rim weight.

\emph{The stage structure.} During stage $n+1$ the component $X_{\mathfrak{a}}$ has not yet
been arrested. It is still in the sequence of stages of step $k$. Hence the structures on it
(levels, weights, cube families) are those of the stages, and so are the domains. In
particular the domain of stage $n+1$ on $X_{\mathfrak{a}}$ is $\mathfrak{D}^{\flat}_{n+1}$.

\emph{The datum.} Let $|\mathsf{B}|<\varrho^{B}_j$. Combine the bound in
\eqref{4.27:eq-interrupted-terms-1} with the definition \eqref{4.27:eq-interrupted-terms-3}:
\begin{equation}\label{4.27:eq-interrupted-terms-6}
\|b(\mathsf{B})\|_\infty\le C_1g_j\sup|\mathsf{B}|\le C_1g_j\varrho^{B}_j=6.6\,C_1\delta'_j.
\end{equation}

\emph{The map on the on-class part.} Among the inputs of Lemma~\ref{4.27:lem-t-step-layer}
(proof incomplete) is the following statement on the descending stage spaces. Let $X$ be a
domain, and suppose
$\|\mathsf{b}\|_\infty\le rC_1\delta'_j$ with $1\le r\le r^{\sharp}_{*}$, where $r^{\sharp}_{*}$
is the admissible multiple of that statement. Then the map $\Phi_{\chi,\sigma}(\mathsf{b})$
of the stage, defined on the stage-$(n+1)$ space $\mathfrak{S}_{n+1}(X)$, takes values in the
stage-$n$ space $\mathfrak{S}_n$ on the enlargement of $X$,
and it uses at most half of the
one-stage room $\rho_m$.

We apply this statement with $r=6.6$, for which $r\le r^{\sharp}_{*}$, and with the cut-off
$\chi\equiv1$. By \eqref{4.27:eq-interrupted-terms-6}, on the on-class part $Z^{\mathrm{on}}$
the map $\Phi(b(\mathsf{B}))$ takes the class part $\mathfrak{L}_{n+1}$ of the domain of
stage $n+1$ into the class part $\mathfrak{L}_n$ of the domain of stage $n$.

\emph{The map on the off-class locations.} On $\mathfrak{T}_{k}$ we use instead the bound for
the stage map at off-class locations, part (ii) of the lemma on that map among the same inputs.
We take it at six times the amplitude. That bound is homogeneous of degree one in the amplitude.
The condition it then imposes is the amplitude condition for arrests, stated with the new-term
estimate among the inputs of Lemma~\ref{4.27:lem-t-step-layer} (proof incomplete).

\emph{The composites $(\beta1)$.} The image lies in the domain of every constituent $c$ in
$(\beta1)$. The constituents are of two kinds: the boundary terms $\mathbb{C}^{(i)}_k$,
$i\le n$, and the terms of the classes of \cite{B-CMP119}. First let $c$ be a boundary term
$\mathbb{C}^{(i)}_k$ with $i\le n$. The claim then
follows from two inputs of Lemma~\ref{4.27:lem-t-step-layer} (proof incomplete). The first
states that
$\mathbb{C}^{(i)}_k(X,\mathcal{S})$ is a stored term on
$\mathfrak{D}^{\flat}_i(X)\times\mathfrak{D}_{\mathcal{S}}(X)$. The second is the descent
$\mathfrak{S}_{i+1}(X)\subset\mathfrak{S}_i(X)$ of the stage spaces. The descent carries the
image, which lies in stage $n$, into stage $i$. Next let $c$ be a term of the classes of
\cite{B-CMP119}. Here the same conclusion is the statement on these terms among the
same inputs of Lemma~\ref{4.27:lem-t-step-layer} (proof incomplete).

Hence every composite is defined on
$\mathfrak{D}^{\flat}_{n+1}(\cdot)\times\{|\mathsf{B}|<\varrho^{B}_j\}$. It is holomorphic
there, as the composition of the holomorphic map of the stages with the holomorphic
constituent $c$. It is bounded there by the supremum of $c$ over its domain, since all its
values are values of $c$. Matching $\mathsf{W}$-pieces are treated under (b$'$) below.

\emph{The family $(\beta2)$.} Its members are holomorphic on
$\mathfrak{K}_{109}(\mathfrak{B};\gamma_0)$. The quadratic term is a polynomial in $\mathsf{B}$
whose coefficients are given by $C$ and $\Delta^{(j)}$, hence entire in $\mathsf{B}$. For the
operator $\Delta^{(j)}$ in that term holomorphy in the pair is the statement on the rim form and
the operator $\Delta^{(j)}$ among the inputs of Lemma~\ref{4.27:lem-t-step-layer} (proof
incomplete).
For $\mathbb{P}^{(j)}$ it is \cite[Lemma~2]{B-CMP116}, taken by statement. For the unmatched
$\mathsf{W}$-pieces it is (b$'$).

\emph{Clause (b$'$).} On the group,
\begin{equation*}
\operatorname{Re}\operatorname{tr}U(\partial p)
=\tfrac12\operatorname{tr}\bigl(U(\partial p)+U(\partial p)^{-1}\bigr),
\end{equation*}
since $U(\partial p)^{-1}=U(\partial p)^{*}$. In this form $A(\mathsf{c},\cdot)$ is a
polynomial in the bond variables and their
inverses. On the complexified group $SL(N,\mathbb{C})$, where the determinant is one, the
inverse of a matrix is
its adjugate, a polynomial in its entries. Hence $A(\mathsf{c},\cdot)$ is entire.

By the inputs of Lemma~\ref{4.27:lem-t-step-layer} (proof incomplete),
$\Phi(\mathsf{b})[P]$ is holomorphic on
$\mathfrak{K}_{109}(\mathfrak{B};\gamma_0)\times\{\|\mathsf{b}\|\le\varepsilon_\sharp\}$, where
$\varepsilon_\sharp$ is the radius of the polydisc of data stated there. Therefore
$A(\mathsf{c},\Phi(b(\mathsf{B}))[P])$, the composition of the entire functional
$A(\mathsf{c},\cdot)$ with this holomorphic map, and with it the difference
\eqref{4.27:eq-interrupted-terms-4}, is holomorphic on
$\mathfrak{K}_{109}(\mathfrak{B};\gamma_0)\times\{|\mathsf{B}|<\varrho^{B}_j\}$. Each
$\mathsf{W}$-piece is obtained from $\mathsf{W}$ by an operation linear in the weights
$\zeta_\square$, which depend only on position \cite[p.~276]{B-CMP119} and not on $P$; it is
therefore holomorphic on the same domain. This
domain is that of the pair variable, so $\mathsf{W}$ imposes no domain condition of its own.

\emph{The family $(\beta3)$: why only real $\mathsf{B}$.} Ba{\l}aban's cut-off function of
\cite[p.~252]{B-CMP119}, a function distinct from the cut-off $\chi$ of the map $\Phi$, is
$C^{\infty}$. It vanishes on a neighbourhood of zero and equals one
away from it, so it has plateaux. A function constant on an open set and not constant is not
analytic. Hence analyticity of the new terms in $\mathsf{B}$ is not available; they are
$C^{\infty}$ in $\mathsf{B}$, and their holomorphy in $P$ is stated at real $\mathsf{B}$ only.
Their holomorphy in the
lifted variables $(\mathsf{B},\mathsf{s},\mathsf{v})$ is the statement on the lift of the
cut-off dependence, taken by statement.

\emph{The family $(\beta3)$: the function $\mathbb{H}^{(j)}_Z$.}
Of the function $\mathbb{H}^{(j)}_Z$ of \cite[(1.58)]{B-CMP122a},
\cite[after~(1.55)]{B-CMP122a} states that it is an analytic function of the configuration of
stage $n+1$ restricted to $Z$ and that it is bounded by a constant times
$e^{-R_j}\varepsilon_j$, with $\varepsilon_j$ the small parameter of that paper at level $j$;
the configurations considered there are real.
Moreover $\mathbb{H}^{(j)}_Z$ lives on the determining set $\mathfrak{B}_Z$, localised as in
\cite{B-CMP122b}. That set lies outside the geometry in which the spaces of
pairs of the inputs of Lemma~\ref{4.27:lem-t-step-layer} are defined.

A domain of holomorphy in $P$ and a bound on that domain are supplied by the part of
Lemma~\ref{4.27:lem-t-step-layer} on the localised layer function (proof incomplete), used here
as stated there. This part
gives the following two facts.
\begin{itemize}
\item[(i)] For a bond $b\in S^{P}$, the continuation $H_\sigma(b)$ of the value
$\mathbb{H}^{(j)}_Z(b)$ is holomorphic in $P$ on a
domain $\mathfrak{D}_H$. This domain contains the restriction of
$\widehat{\mathfrak{D}}^{\flat}_{n+1}$ to the domain $Y$ on which $H_\sigma(b)$ is localised.
At real points with $\sigma\equiv1$ it equals $\mathbb{H}^{(j)}_Z(b)$.
\item[(ii)] $|H_\sigma(b)|\le H_{\max}$, with $H_{\max}$ as in
\eqref{4.27:eq-t-step-layer-b}. This bound is of the same kind as the bound of
\cite{B-CMP122a} recalled above: constants times a radius times an exponentially small factor.
The mechanism is clause (b) of the layer theorem, applied at
$\mathfrak{B}_Z$ with a decay profile proportional to
$e^{-\delta d^{\wedge}(\cdot,\mathfrak{L})}$; here $\delta$ and $d^{\wedge}$ are the decay rate
and the scaled distance of the layer theorem, and $\mathfrak{L}$ is the set from which that
distance is measured there, not a class part $\mathfrak{L}_n$.
\end{itemize}
By the same part of that lemma, the shift obeys $|\delta_b|\le 2H_{\max}/g_j$. The condition
on the layer bound among the hypotheses \eqref{4.27:eq-arrest-hypotheses-a} of the arrest
theorem (Definition~\ref{4.27:def-arrest-hypotheses}) then gives
\begin{equation}\label{4.27:eq-interrupted-terms-7}
|\delta_b|\le\frac{2H_{\max}}{g_j}<\frac{0.3\,\delta'_j}{g_j}.
\end{equation}

\emph{The family $(\beta3)$: the shifted arguments.} Let $\mathsf{B}$ be real with
$|\mathsf{B}|<6.3\,\delta'_j/g_j$. By \eqref{4.27:eq-interrupted-terms-7}, every point
$\mathsf{B}'$ of the segment from $\mathsf{B}$ to $\mathsf{B}+\delta_b$ satisfies
\begin{equation*}
|\mathsf{B}'|<\frac{6.3\,\delta'_j}{g_j}+\frac{0.3\,\delta'_j}{g_j}
=\frac{6.6\,\delta'_j}{g_j}=\varrho^{B}_j .
\end{equation*}
So the increments of $(\beta1)$, $(\beta2)$ and $\mathsf{W}$ in $(\beta3)$ are taken inside
the domain established above for these families. The Jacobian factor is holomorphic in $P$ on
$\mathfrak{D}_H$ by (i), and the increments are holomorphic in $P$ on
$\mathfrak{D}^{\flat}_{n+1}(\cdot)$ and on the space of pairs of (b$'$). All of them are
therefore holomorphic in $P$ on the restriction of $\widehat{\mathfrak{D}}^{\flat}_{n+1}(\cdot)$
to the domain $Y$ of (i),
which lies in $\mathfrak{D}^{\flat}_{n+1}(\cdot)$, in the space of pairs by its ambient
condition, and in $\mathfrak{D}_H$ by (i). This proves the assertions of (b) on $(\beta3)$,
with domain $\widehat{\mathfrak{D}}^{\flat}_{n+1}(\cdot)$.

\emph{The real range.} The actual values satisfy $|g_jB_j|<6\,\delta'_j$
\cite[(1.54), (1.59)]{B-CMP122a}. Hence $|B_j|<6\,\delta'_j/g_j<6.3\,\delta'_j/g_j$, and the
real range of $\mathsf{B}$ lies inside the range covered by (b).

\emph{The consequence in (b).} By the stage structure, the domain of stage $n+1$ is the class
domain of stage $n^{+}=n+1$. By what was just shown, the domain of every member on
$W_{\mathfrak{a}}$ is that class domain. Therefore Lemma~\ref{4.27:lem-frozen-room}(a)
applies to the members with $n=n^{+}$.

\emph{Clause (c).} This is the estimate on the new terms and the statement on the weak chart
coordinate among the inputs of Lemma~\ref{4.27:lem-t-step-layer} (proof incomplete). We take
it by statement, at radius $\varrho^{B}_j$.

\emph{Clause (d).} Clause (d) is a delimitation and needs no proof.
\end{proof}

\begin{theorem}[The arrest theorem. Stated; proof incomplete at the step marked $(\ast)$]
\label{4.27:thm-arrest-theorem}
Let the arrested reading of the spaces (Definition~\ref{4.27:def-arrested-spaces},
\eqref{4.27:eq-arrested-spaces-a}) be in force from each arrest on. Assume the hypotheses of
the arrest theorem (Definition~\ref{4.27:def-arrest-hypotheses}) and the inputs
listed in Lemma~\ref{4.27:lem-t-step-layer}, whose proof is incomplete at one step. Those used
below are named by their content: the stage-map input, the descent input, the stored-term
input (with its room table and the corollary of that table), the $\mathbf{R}$-layer input, the
$\mathbf{T}$-step input, the keystone input, the post-arrest input and the localised
change-of-variables input.

We use the following notation. Arrests, live arrests $\mathfrak{a}$ with their densities
$k_a$, components $X_{\mathfrak{a}}$, pieces $W_{\mathfrak{a}}$ and stored terms
$\mathfrak{F}(\mathfrak{a})$ are those of Definition~\ref{4.27:def-arrests-pieces}. The later
operations (O1)--(O5), (O4$'$), the completion site (L) of a chain and the zones of a point
are those of Definition~\ref{4.27:def-operations-zones} and
\eqref{4.27:eq-operations-zones-a}. Determining sets are ordered in time by $\preccurlyeq$
(Lemma~\ref{4.27:lem-determining-set-order}). A stored term $\tau$ lives on a region
$X_\tau$, is born over its birth set $\mathfrak{B}_\tau$ at density $k_\tau$, and is defined
on its polydisc; the condition this polydisc imposes at a point $x$ is \emph{$\tau$'s clause
at $x$}. For a later operation $\mathfrak{O}$, $\mathfrak{B}^t$ and $\mathfrak{B}'$ are its
target and source sets as in (e4), \eqref{4.27:eq-weights-heredity-b}, of
Lemma~\ref{4.27:lem-weights-heredity}, whose proof is incomplete at one step; the arrested
source of $\mathfrak{O}$ is written over
$\mathfrak{B}'$, and $Y'$ is the region of a member of that source. A point of
$\Omega_0(\mathfrak{B})\setminus\Omega_1(\mathfrak{B})$ is called \emph{raw at
$\mathfrak{B}$}; these points form the \emph{stratum} of $\mathfrak{B}$. The spaces of pairs
$\mathfrak{G}(\mathfrak{B})\subset\mathfrak{G}^{-J_0}(\mathfrak{B})\subset
\mathfrak{G}_1(\mathfrak{B})$ on $\Omega_0(\mathfrak{B})$ are those of clause (e) of that
lemma; their pairs have the form
$(e^{i\eta A'}U,\mathbb{J})$. The \emph{$G$-part} of a pair $(\mathbb{U},\mathbb{J})$ is its
variable $\mathbb{U}$; the map of a later operation is a \emph{layer into $U'$} when it
multiplies the $G$-part by a layer
$e^{i\eta\mathfrak{h}}$, as the $\mathbf{T}$-map of the $\mathbf{T}$-step input does, $U'$
denoting the $G$-part
of the image. At a raw point $y$ the \emph{index-zero members} of
$\mathfrak{G}(\mathfrak{B})$ are its two $U$-members, the cube gauges $(3.35)_0$ on the
index-zero cubes and the member $(3.37)_0$ bearing on $A'$, where $(3.35)_0$ and $(3.37)_0$
denote the index-zero instances of \cite[(3.35)]{B-CMP99b} and \cite[(3.37)]{B-CMP99b}; and
its \emph{current member}
\begin{equation*}
\eta^{3}|\mathbb{J}(y)|<T^{\mathfrak{B}}(y),
\end{equation*}
with $\eta$ the lattice unit at index zero and $T^{\mathfrak{B}}(y)$ the dated threshold fixed
there (the member is written $\ell^3|\mathbb{J}(y)|<T^{\mathfrak{B}}(y)$ in
\eqref{4.27:eq-weights-heredity-b}).

Then for every $K$, every label, every live arrest $\mathfrak{a}$, every
$\tau\in\mathfrak{F}(\mathfrak{a})$ (at points of zone Z1: every term whose polydisc is as in
Lemma~\ref{4.27:lem-frozen-room}), every later operation $\mathfrak{O}$ among (O1)--(O5),
(O4$'$), and every $x\in X_\tau$, the following hold.
\begin{itemize}
\item[(A)] The image under the map of $\mathfrak{O}$ of the arrested source of $\mathfrak{O}$
satisfies $\tau$'s clause at $x$, with the room, spend and ratio of
\eqref{4.27:eq-arrest-theorem-a}.
\item[(A$'$)] Let $\tau\in\mathfrak{F}(\mathfrak{a})$ be a stored term, born over
$\mathfrak{B}_\tau$ at density $k_\tau=k_a$, and let $p$ be a member over $Y'$ of the arrested
source of a later operation $\mathfrak{O}$, with continuation
$\tilde p\in\mathfrak{G}(\mathfrak{B}')$ given by the ambient rider of
Definition~\ref{4.27:def-arrested-spaces}. Then the image of $p$ under $\mathfrak{O}$ has a
continuation $\tilde q$ which equals the image on $Y'$ and the input on
$X_\tau\setminus Y'$ (and the value of $\tau$ depends on its argument on $X_\tau$ only); off
$X_\tau\cup Y'$ its variable $\mathbb{U}$ is that of the input, and its variable $\mathbb{J}$
is that of the input on $\Omega_1(\mathfrak{B}^t)$ and zero at the points raw at
$\mathfrak{B}^t$. This continuation
has the following properties.
\begin{enumerate}
\item[(1)] $\tilde q\in\mathfrak{G}_1(\mathfrak{B}_\tau)$; hence $\tilde q$ satisfies the
clauses of $\mathfrak{K}_{109}(\mathfrak{B}_\tau;\gamma_0^{(k_\tau)})$ at every point of
$\Omega_1(\mathfrak{B}_\tau)$ and at every cube.
\item[(2)] At the points of the stratum of $\mathfrak{B}_\tau$ that lie in
$\Omega_1(\mathfrak{B}^t)$, $\tilde q$ obeys every index-zero member of
$\mathfrak{G}(\mathfrak{B}_\tau)$, the current member included.
\item[(3)] At the points of the stratum of $\mathfrak{B}_\tau$ that are raw at $\mathfrak{B}^t$
and lie off $Y'$, $\tilde q$ obeys every index-zero $U$-member of
$\mathfrak{G}(\mathfrak{B}_\tau)$. It obeys the current member of
$\mathfrak{G}(\mathfrak{B}_\tau)$ at every such point off $X_\tau$, and at every such point of
$X_\tau\setminus Y'$ that is dated by \eqref{4.27:eq-density-increase-a} (a statement whose
proof is incomplete at one step): a point inside
$X_{\mathfrak{a}}$, inside a class component of step $k_a$, or inside a piece live at
$\mathfrak{B}_\tau$; or, when the source set $\mathfrak{B}'$ of $\mathfrak{O}$ is of density
$k_a$ (for (O1), and for (O2) at $k=k_a$), any such point.
\item[(4)] At the points of the stratum of $\mathfrak{B}_\tau$ that are raw at
$\mathfrak{B}^t$ and lie on $Y'$, the $G$-part of $\tilde q$ obeys the index-zero cube gauges
$(3.35)_0$ wherever the map of $\mathfrak{O}$ is a layer into $U'$: for (O1), (O3), (O4),
(O4$'$), and, at the completion site (L), at the points served by the $\mathbf{R}$-layer
input, provided the polydisc of that input carries the shell of index zero; this proviso is
part of the second locus marked $(\ast)$ below and is not asserted.
\end{enumerate}
$(\ast)$ Not asserted in (A$'$) are the following two loci. First, the current member at
raw points of $X_\tau\setminus Y'$ that lie in components of $Z_{k_a}$ which are not class
components of step $k_a$ and lie outside every piece live at $\mathfrak{B}_\tau$, when the
source set of $\mathfrak{O}$ is of a density $>k_a$, that is, for (O3), for (O2) at $k>k_a$,
and for (O4), (O4$'$), (O5). Here the components that are not class components of step $k_a$
are those failing one of the two conditions of a class component in
Definition~\ref{4.27:def-operations-zones}; among them are the components that hold a
second-class component of an earlier step, such second-class components being created for the
timing in the second condition.
Second, on $Y'$ at points raw at $\mathfrak{B}^t$, the members that the map of
$\mathfrak{O}$ moves, namely $(3.37)_0$ for $A'$ and the current member; and, at the
completion site (L), every index-zero member at raw points of $\mathsf{Z}$ and of the core of
the completion site, that is, of the part of the site on which the determining set dated
during the step keeps the frozen levels $o(x)$ and the frozen weights.

For $\tau$ of kind $(\beta 3)$ of Lemma~\ref{4.27:lem-interrupted-terms} one has
$\mathfrak{B}_\tau=\mathbb{B}^{(n+1)}_{k_a}$; put
$\mathfrak{B}^{\circ}:=\mathfrak{B}_\tau\restriction X_{\mathfrak{a}}$. The ambient membership
on $X_{\mathfrak{a}}$ required by the localised change-of-variables input, namely a
continuation in
$\mathfrak{K}_{109}(\mathfrak{B}^{\circ};\gamma_0^{(k_a)})$, is supplied on the same scope: its
clauses on $X_{\mathfrak{a}}\cap\Omega_1(\mathfrak{B}_\tau)$ and at the cubes; and, at the raw
points of $X_{\mathfrak{a}}$, its $U$-members off $Y'$, its gauge $(3.35)_0$ on $Y'$ where the
map is a layer into $U'$, and its current member off $Y'$ (there $T\le\gamma_0^{(k_a)}$ on
$X_{\mathfrak{a}}$). The members moved on
$X_{\mathfrak{a}}\cap Y'$ belong to the second locus above and are not asserted.
\item[(B)] The conclusion of Lemma~\ref{4.27:lem-real-minimiser}, whose proof is incomplete
at one step, holds, its rider included in the form: the real point lies in
$\mathfrak{G}_1(\mathfrak{B})$. At raw points of index zero of $\Omega_0(\mathfrak{B})$ the
real point obeys the member $(3.37)_0$. $(\ast)$ Its index-zero gauge $(3.35)_0$ and its
current member $\eta^{3}|\mathcal{J}(U)|<T^{\mathfrak{B}}$ at these points are not asserted.
\item[(C)] The polydiscs assigned to newborn terms by the arrested reading
(Definition~\ref{4.27:def-arrested-spaces}) are legitimate, as far as (A$'$) reaches. Here
the polydisc of a term born at a later operation $\mathfrak{O}$, which is a stored term
composed with the map of $\mathfrak{O}$, is called \emph{legitimate} when the polydisc of that
stored term, read with its ambient rider, holds the image under the map of $\mathfrak{O}$ of
the arrested source of $\mathfrak{O}$, so that the composite is defined on it.
$(\ast)$ For the riders of old terms at the two loci not asserted in (A$'$), legitimacy is not
asserted; the first clause of the proposition on the riders of newborn terms inherits this
restriction.
\item[(D)] Off live pieces every space is the unarrested one, its rider understood as in the
ambient rider of Definition~\ref{4.27:def-arrested-spaces}. On $W_{\mathfrak{a}}$,
\begin{equation*}
\tilde{U}^{\mathrm{arrest}}_{k'}\subset\mathfrak{S}^{(k_a),\mathrm{arrest}}_{n^+}
\subset\mathfrak{S}^{(k_a),\mathrm{arrest}}_{0}.
\end{equation*}
Every arrested space lies in $\mathfrak{N}^{\vee}\subset\mathfrak{K}_{P}$. Differences of
factors are those of the unarrested construction.
\item[(E)] No constant depends on $k$, $k'-k_a$, the number of stages $N$, the threshold
$N_0$, the nest depth, $K$, $\gamma_0^{(k)}$ or a slot (the arrested reading of
Definition~\ref{4.27:def-arrested-spaces} is declared free of slots). $L$ enters only through
floors, and
the coupling only through ceilings in one variable $\mathsf{t}$.
\end{itemize}
In \eqref{4.27:eq-arrest-theorem-a} all quantities are in units at $x$; $m$ and $\ell$ denote
the level; $u:=2^{-(k'+1-\ell)}$; and $\rho_n$ is read as $\rho_{\max(n,n^+_W)}$, where
$n^+_W$ is the stage $n^+$ of the live piece $W$ holding $x$
(Definition~\ref{4.27:def-arrests-pieces}). The index $a\in\{1,\dots,7\}$ runs over the direct
entries, and $r_a(\cdot)$ is the threshold of the entry $a$ as a function of the level;
$\mu=\mu_\tau(x)$ is the level at which $\tau$'s clause is evaluated at $x$
(Definition~\ref{4.27:def-operations-zones}); $\theta_T$ is the threshold of $\tau$'s own
clause at $x$ ($\tau$ being the term whose clause must hold the image), as in the room table
of the stored-term input; $k_b$ is the density of the later attempt; $k$ is the density of the
step at which the chain of (O2) completes, and $k'$ that of the $\mathbf{T}$-step $k'\to k'+1$
of (O3); in clause (D), $k'\ge k_a$. The constant $\beta$ is the schedule constant of the
layered spaces (Definition~\ref{4.4:def-post-step-space}), $\theta_S$ is the enlargement
fraction of the $\mathbf{T}$-step input, $\bar{P}$ is the maximal layer increment of the
$\mathbf{R}$-layer input, and $\mathsf{w}_0=1/24$.
\begin{equation}\label{4.27:eq-arrest-theorem-a}
\begin{array}{l|l|l|l}
\text{zone, operation} & \text{room}\ \ge & \text{spend}\ \le & \text{ratio}\ \le\\
\hline
\mathrm{Z1},\ (\mathrm{O1}),(\mathrm{O4}) & \rho^{(l)}_\ell & \tfrac12\rho^{(l)}_\ell
 & \tfrac12\\
\mathrm{Z1},\ (\mathrm{O4}') & \beta2^{-(k_a-\ell)}
 & \tfrac12\mathsf{w}_0\beta2^{-(k_a-\ell)} & \tfrac1{48}\\
\mathrm{Z1},\ (\mathrm{O2}) & \beta2^{-(k_a+1-\ell)} & \bar{P}2^{-(k-1-\ell)^+}
 & 4\bar{P}/\beta\le\theta_S/4\\
\mathrm{Z1},\ (\mathrm{O3}) & (2-\theta_S)\beta u & \tfrac14(1-\theta_S)\beta u & \tfrac18\\
\mathrm{Z2},\ (\mathrm{O2}) & \tfrac12\theta_S\beta2^{-(k_a-m)} & \bar{P}2^{-(k-1-m)^+}
 & \tfrac14\\
\mathrm{Z2},\ (\mathrm{O3}) & \beta u & \tfrac14(1-\theta_S)\beta u & \tfrac18\\
\mathrm{Z2},\ (\mathrm{O4}),(\mathrm{O4}') & \tfrac34\beta2^{-(k_a-m)}
 & \tfrac12\mathsf{w}_0\beta2^{-(k_a-m)} & \tfrac1{36}\\
\mathrm{Z2}^{\circ} & \text{that of Z2, untouched} & 0 & 0\\
\mathrm{Z3a} & 1.5-1.2 & \max\{\beta/8,\bar{P}\} & \tfrac1{12}\\
\mathrm{Z3b} & \tau\text{'s},\ \ge0.78\,r_a(\mu) & \le0.436\,r_a(\mu) & \tfrac23\\
\mathrm{Z3\text{-}L} & \tau\text{'s},\ \ge\tfrac12\theta_T & \theta_T/16 & \tfrac18\\
\mathrm{Z4}' & \theta & \theta/64+3\theta/64 & \tfrac1{16}
\end{array}
\end{equation}
In the first row the map takes $\mathfrak{S}^{\mathrm{arrest}}_l$ into
$\mathfrak{S}^{\mathrm{arrest}}_{l-1}$, which lies in $\tau$'s polydisc. In the second row the
room is $\rho_n(x;k_b)$, which is at least $\beta2^{-(k_a-\ell)}$. In the third row
$k\ge k_a$, the same density included. In the rows for (O3) the room is that after the
enlargement. In the row for Z2 with (O4), (O4$'$) the point satisfies
$x\in\mathfrak{T}_{k_b}$. Zone $\mathrm{Z2}^{\circ}$ is the descent, zone Z3a the scale drop,
and zone Z3b the raised points, read clause to clause. The last row, labelled
$\mathrm{Z4}'$, concerns points of zone Z4, that is, points in a piece dissolved at
$\mathfrak{O}$ (Definition~\ref{4.27:def-operations-zones}), on the core of the completion
site; there the spend consists of a
base $\theta/64$ and increments $3\theta/64$, and
\begin{equation*}
\theta:=\theta_{\tau,a}(x)=F_\tau\,\varpi\,r_a(\lambda),\qquad \lambda\le o(x),\quad
\varpi\ge1,\quad F_\tau\ge\varphi_{\mathrm{arrest}},
\end{equation*}
where $\lambda$ is a level, $\varpi$ is the weight at $x$, $o(x)$ the level at $x$ on the core
of the completion site, $F_\tau$ the factor of $\tau$'s clause, and $\varphi_{\mathrm{arrest}}$
the floor of Lemma~\ref{4.27:lem-factor-positivity}.
\end{theorem}

\begin{proof}
We treat the clauses in the order (A), (A$'$), (B), (C), (D), (E).

\emph{Clause (A).} Fix a row of \eqref{4.27:eq-arrest-theorem-a}. Its room is a lower bound,
in units at $x$, either for the amount by which each entry of the arrested source of
$\mathfrak{O}$ lies below the corresponding threshold of $\tau$'s clause at $x$ (zones Z1, Z2,
$\mathrm{Z2}^{\circ}$, Z3a), or, where the table writes $\tau$'s threshold or $\theta$, for the
threshold of $\tau$'s clause itself (zones Z3b, Z3-L and the row $\mathrm{Z4}'$). Its spend is
correspondingly an upper bound either for the amount by which the map of $\mathfrak{O}$ moves
that entry or for the size of the presented entry itself; its ratio bounds the quotient of the
spend by the room. Each ratio of the table is less than one, so the spend is smaller than the
room, and in either reading the image of the arrested source satisfies $\tau$'s clause at $x$.
It remains to identify the rooms and spends and to check the ratios. Some rooms and spends are
taken from two statements proved after this theorem in this section, the statement on the room
in the regular zones and the statement on the room at dissolving pieces.

\emph{The constants $\beta$ and $\theta_S$.} Lemma~\ref{4.27:lem-factor-positivity} states
$\varphi_{\mathrm{arrest}}=1-4.51\beta=0.098$ for the schedule constant $\beta$ of the layered
spaces (Definition~\ref{4.4:def-post-step-space}); hence $4.51\beta=0.902$, that is,
$\beta=0.2$. For the enlargement fraction $\theta_S$ of the $\mathbf{T}$-step input we use
$\tfrac12\le\theta_S\le1$.

\emph{Rooms.} In zone Z1 the clause of $\tau$ at $x$ is a class polydisc of the live piece
holding $x$, read with the frozen structure; its room over the arrested space is
$\rho_n(x;k')$ of Lemma~\ref{4.27:lem-frozen-room}(a), \eqref{4.27:eq-frozen-room-a}, with
$\rho_n$ read as $\rho_{\max(n,n^+_W)}$. For (O4$'$) this is the room $\rho_n(x;k_b)$ of the
second row; it is at least $\beta2^{-(k_a-\ell)}$, as follows. By (T2) of
Lemma~\ref{4.27:lem-arrest-timing}, $\ell\le k_a$, and by
Definition~\ref{4.27:def-operations-zones} the later attempt of (O4$'$) has $k_b\ge k_a+1$.
Every term of \eqref{4.27:eq-frozen-room-a} is non-negative; its first sum runs up to
$i=k_a+1$ and, the stage levels $h_a+n$ of the arrest not exceeding $k_a$
(Definition~\ref{4.27:def-arrests-pieces}), contains the term $i=k_a+1$, which is
$\beta2^{-|\ell-k_a-1|}=\beta2^{-(k_a+1-\ell)}$; its
second term at $k'=k_b$ is at least
$\beta(2^{-(k_a-\ell)}-2^{-(k_a+1-\ell)})=\beta2^{-(k_a+1-\ell)}$. Hence
\begin{equation*}
\rho_n(x;k_b)\ge\beta2^{-(k_a+1-\ell)}+\beta2^{-(k_a+1-\ell)}=\beta2^{-(k_a-\ell)}.
\end{equation*}
For (O1)
and (O4) at stage $l$, source and target read one frozen structure
(Definition~\ref{4.27:def-arrested-spaces}), so the terms $\beta\sigma_{\mathfrak{a}}(\ell)$ of
the arrested factor, which do not depend on the stage, are the same in both and cancel, and the
room of one stage is $\rho^{(l)}_\ell$ (the stage-map input). In the rows for (O3) the room is
taken
after the enlargement $F^{(k'+1)}(x)+\theta_S\beta2^{-(k'+1-\ell)}$ of the source factor on a
live piece, which the arrested reading of the polydiscs of newborn terms prescribes following
the $\mathbf{T}$-step input; its values $(2-\theta_S)\beta u$ in zone Z1 and $\beta u$ in zone
Z2 after the
enlargement are those established in the statement on the room in the regular zones, given
below in this section.
In zone Z2 with (O2) the room is that of the room
table of the stored-term input and its corollary: the threshold of $\tau$ exceeds that of the
source by at
least $\tfrac12 r_T r_a$, where $r_T$ is the room parameter of the consuming term in that room
table, with $r_T\ge\beta\theta_S2^{-(k_a-m)}$ for $m<k_a$, and with $r_T=\beta$
on $\Omega_{k_a}$, where $m=k_a$, so that there, as $\theta_S\le1$,
\begin{equation*}
r_T=\beta\ge\beta\theta_S=\beta\theta_S2^{-(k_a-m)};
\end{equation*}
in both cases this is the room
$\tfrac12\theta_S\beta2^{-(k_a-m)}$ of that row. In zone Z2 with (O4), (O4$'$) the room
$\tfrac34\beta2^{-(k_a-m)}$ at a point of $\mathfrak{T}_{k_b}$ is that established in the
statement on the room in the regular zones. In zone
Z3a, at a scale drop, the threshold of $\tau$ is at least $1.5\,r_a$ by the corollary of the
room table of the stored-term input; the value $1.2$ subtracted from it in the room of that row
is the bound on the
corresponding entry of the source at a scale drop established in the statement on the room in
the regular zones. In zone Z3b the source level has risen by the stages of a
later attempt, and by (e2) of Lemma~\ref{4.27:lem-weights-heredity}, whose proof is incomplete
at one step, thresholds fall by the factor $0.436$ at least where the level rose; this is the
factor of the spend $0.436\,r_a(\mu)$ of that row, while the room $0.78\,r_a(\mu)$ of $\tau$'s
clause at points where the level of $\tau$'s clause lies below the source level is that
established in the statement on the room in the regular zones. In zone Z3-L and in the row
$\mathrm{Z4}'$ the rooms are those of $\tau$'s clause, $\tfrac12\theta_T$ and $\theta$
respectively, as established in the statement on the room at dissolving pieces, given below in
this section.

\emph{Spends.} In zone Z1, whose points lie in live pieces, the spends for (O1)--(O4$'$) are
the increments of Lemma~\ref{4.27:lem-live-piece-increments}: $\tfrac12\rho^{(l)}_\ell$ at stage
$l$ of (O1), (O4); $\sum_l\iota_l(x)\le\tfrac12\mathsf{w}_0\beta e^{-(k_b-1-\ell)}$ for (O4$'$);
$\bar{P}2^{-(k-1-\ell)^+}$ for (O2) at $k$; and
$\tfrac14(1-\theta_S)\beta2^{-(k'+1-\ell)}=\tfrac14(1-\theta_S)\beta u$ for (O3). In zone Z2 the
spends are the increment bounds of the $\mathbf{R}$-layer input, the $\mathbf{T}$-step input
and part (c) of the stored-term input at the level $m$ of $x$: for
(O2) at $k$, the $\mathbf{R}$-layer input bounds the increments by $P2^{-(k-1-m)}$ on foreign
shells $m\le k-1$,
by $P'$ on $\Omega_k\setminus Z$ and by $P''$ on collars, which, with $\bar{P}$ the maximum of
these constants, is the spend $\bar{P}2^{-(k-1-m)^+}$; for (O3) at $k'\to k'+1$, the
$\mathbf{T}$-step input
gives $\pi_m\le\tfrac14(1-\theta_S)\beta2^{-(k'+1-m)}$; and for (O4), (O4$'$), part (c) of the
stored-term input
gives $\tfrac12\mathsf{w}_0\beta e^{-n_W}$ with $n_W=k_b-1-m$. For (O4$'$) in zone Z1,
the level satisfies $\ell\le k_a$ by (T2) of Lemma~\ref{4.27:lem-arrest-timing}, and
$k_b\ge k_a+1$, so the exponent $k_b-1-\ell$ is at least $k_a-\ell\ge0$, and
\begin{equation*}
e^{-(k_b-1-\ell)}\le 2^{-(k_b-1-\ell)}\le 2^{-(k_a-\ell)},
\end{equation*}
which gives the spend $\tfrac12\mathsf{w}_0\beta2^{-(k_a-\ell)}$. In zone Z2, for (O4) and
(O4$'$), the level $m=\mu_\tau(x)$ read by $\tau$'s clause, which is written at the birth density
$k_a$ of $\tau$, satisfies $m\le k_a$, and, as in zone Z1, $k_b\ge k_a+1$; so the same
inequality with $m$ in place of $\ell$ gives the spend $\tfrac12\mathsf{w}_0\beta2^{-(k_a-m)}$.
In zone
$\mathrm{Z2}^{\circ}$ the source level equals $\mu_\tau(x)$ and the map is the descent, which
does not move the entry; the spend is $0$. In zone Z3a the spend is bounded by the bound
$\beta/8$ of the $\mathbf{T}$-step input on the top shell $\Omega_{k'+1}$ and by the increment
bound $\bar{P}$
of the $\mathbf{R}$-layer input.
In zone Z3-L the keystone input bounds each entry of the presented pair on $\mathsf{Z}$ by one
sixteenth of the threshold of $\tau$'s clause, which is the spend $\theta_T/16$. In the row
$\mathrm{Z4}'$, at the points of zone Z4 on the core of the completion site, the spend for the
re-presented pair is a base
$\theta/64$ plus increments totalling $3\theta/64$, as established in the statement on the room
at dissolving pieces.

\emph{Ratios.} For (O1), (O4) in zone Z1 the quotient is
$\tfrac12\rho^{(l)}_\ell/\rho^{(l)}_\ell=\tfrac12$. For (O4$'$) in zone Z1 it is
$\tfrac12\mathsf{w}_0=\tfrac1{48}$, since $\mathsf{w}_0=1/24$. For (O2), let $k\ge k_a$. If
$\ell\le k-1$, then $2^{-(k-1-\ell)}=2\cdot2^{-(k-\ell)}\le2\cdot2^{-(k_a-\ell)}$; if $\ell\ge
k$, then $2^{-(k-1-\ell)^+}=1\le2^{-(k_a-\ell)}$, as $\ell\ge k_a$. In both cases
\begin{equation*}
2^{-(k-1-\ell)^+}\le2\cdot2^{-(k_a-\ell)}=4\cdot2^{-(k_a+1-\ell)},
\end{equation*}
so the quotient in the row Z1 with (O2) is at most $4\bar{P}/\beta$, and this is at most
$\theta_S/4$ by the hypothesis $16\bar{P}\le\beta\theta_S$ of
Definition~\ref{4.27:def-arrest-hypotheses}. With $m$ in place of $\ell$ the same
inequality bounds the spend of the row Z2 with (O2) by $2\bar{P}2^{-(k_a-m)}$, and the quotient
by
\begin{equation*}
\frac{2\bar{P}}{\tfrac12\theta_S\beta}=\frac{4\bar{P}}{\theta_S\beta}\le\frac14.
\end{equation*}
For (O3) in zone Z1 the quotient is $(1-\theta_S)/(4(2-\theta_S))$, where $2-\theta_S\ge1>0$
as $\theta_S\le1$. Since $\theta_S\ge0$, one has
$1-\theta_S\le\tfrac12(2-\theta_S)$, and the quotient is at most $\tfrac18$. For (O3) in zone
Z2 the quotient is $\tfrac14(1-\theta_S)\beta u/(\beta u)=\tfrac14(1-\theta_S)$, which is at
most $\tfrac18$ as $\theta_S\ge\tfrac12$. For
(O4),
(O4$'$) in zone Z2 it is $\tfrac12\mathsf{w}_0/\tfrac34=\tfrac23\mathsf{w}_0=\tfrac1{36}$. In
zone $\mathrm{Z2}^{\circ}$ it is $0$. In zone Z3a it is the quotient of
$\max\{\beta/8,\bar{P}\}$ by $1.5-1.2=0.3$; for the term $\beta/8$ it is, as $\beta=0.2$,
\begin{equation*}
\frac{0.2/8}{0.3}=\frac{0.025}{0.3}=\frac1{12},
\end{equation*}
and the term $\bar{P}$ does not exceed $\beta/8$: since $\theta_S\le1$, the hypothesis
$16\bar{P}\le\beta\theta_S$ gives $\bar{P}\le\beta/16$; hence the quotient is $\tfrac1{12}$.
In zone Z3b it
is $0.436/0.78<0.56<\tfrac23$; in zone Z3-L it is $\tfrac1{16}/\tfrac12=\tfrac18$; in the row
$\mathrm{Z4}'$ it is $(\theta/64+3\theta/64)/\theta=\tfrac1{16}$. This gives (A).

\emph{Clause (A$'$): the continuation.} Let $\tau\in\mathfrak{F}(\mathfrak{a})$ be born over
$\mathfrak{B}_\tau$. By the lemma on the time order of determining sets
(Lemma~\ref{4.27:lem-determining-set-order}), its density is
$k_\tau=k_a$. Let $\mathfrak{O}$ be a later operation and let $p$ be a member of its arrested
source over $Y'$. That source is read with the ambient rider of
Definition~\ref{4.27:def-arrested-spaces}; hence $p$ is the restriction of a pair
$\tilde p\in\mathfrak{G}(\mathfrak{B}')$, its continuation. Define $\tilde q$ as the image of
$p$ on $Y'$ and as $\tilde p$ on $X_\tau\setminus Y'$; off $X_\tau\cup Y'$ let its variable
$\mathbb{U}$ be that of $\tilde p$, and its variable $\mathbb{J}$ be that of $\tilde p$ on
$\Omega_1(\mathfrak{B}^t)$ and zero at the points raw at $\mathfrak{B}^t$. By clauses (a) and
(f) of the theorem on
the function of the change of variables, the image of the map of $\mathfrak{O}$ coincides with
its input off
$Y'$; hence $\tilde q$ is a continuation of the image of $p$, equal to the image on $Y'$ and
to the input on $X_\tau\setminus Y'$. By the localisation clause of the stored-term
representation, the value of $\tau$
depends on its argument on $X_\tau$ only. The rider asks only for the existence of a
continuation, so the choice of $\mathbb{J}$ at the raw points off $X_\tau\cup Y'$ is free, and
zero is chosen there.

\emph{Item (1).} By (e4) of Lemma~\ref{4.27:lem-weights-heredity}, whose proof is incomplete at
one step, applied to $\mathfrak{O}$ with target set $\mathfrak{B}^t$ and source set
$\mathfrak{B}'$ and to the continuation $\tilde p\in\mathfrak{G}(\mathfrak{B}')$, the
continuation so defined satisfies $\tilde q\in\mathfrak{G}_1(\mathfrak{B}^t)$. By the same
lemma on the time order of determining sets, $\mathfrak{B}_\tau\preccurlyeq\mathfrak{B}^t$.
Therefore (e1)--(e3) of Lemma~\ref{4.27:lem-weights-heredity}, whose proof is incomplete at one
step, and the nested-cube lemma (Lemma~\ref{4.27:lem-nested-cubes},
\eqref{4.27:eq-nested-cubes-a}) give
$\mathfrak{G}_1(\mathfrak{B}^t)\subset\mathfrak{G}_1(\mathfrak{B}_\tau)$, whence
$\tilde q\in\mathfrak{G}_1(\mathfrak{B}_\tau)$. By (d) of
Lemma~\ref{4.27:lem-weights-heredity}, whose proof is incomplete at one step, the clauses of
$\mathfrak{G}_1(\mathfrak{B}_\tau)$ lie inside those of
$\mathfrak{K}_{109}(\mathfrak{B}_\tau;\gamma_0^{(k_\tau)})$ at every point of
$\Omega_1(\mathfrak{B}_\tau)$ and at every cube. This is item (1).

\emph{Item (2).} Let $y$ be raw at $\mathfrak{B}_\tau$ and lie in $\Omega_1(\mathfrak{B}^t)$.
Since $\tilde q\in\mathfrak{G}_1(\mathfrak{B}^t)$, the heredity statement on the raw stratum,
Lemma~\ref{4.27:prf-raw-stratum}, gives every index-zero member of
$\mathfrak{G}(\mathfrak{B}_\tau)$ at $y$. For the current member this is the entering line of
that statement together with \eqref{4.27:eq-held-raw-points-a}(2).

\emph{Item (3).} Let $y$ be raw at $\mathfrak{B}_\tau$ and at $\mathfrak{B}^t$, and lie off
$Y'$. There the variable $\mathbb{U}$ of $\tilde q$ is that of the input $\tilde p$. By (e3) of
Lemma~\ref{4.27:lem-weights-heredity}, whose proof is incomplete at one step,
$\mathfrak{G}(\mathfrak{B}')\subset\mathfrak{G}^{-J_0}(\mathfrak{B}_\tau)$; so $\tilde q$
obeys every index-zero $U$-member of $\mathfrak{G}(\mathfrak{B}_\tau)$ at $y$. If $y$ lies off
$X_\tau$, the variable $\mathbb{J}$ of $\tilde q$ is zero at $y$, and since the threshold
$T^{\mathfrak{B}_\tau}(y)$, a value of $\gamma_0^{(\cdot)}$, is positive, the current member
holds there.
Let now $y\in X_\tau\setminus Y'$ be dated by \eqref{4.27:eq-density-increase-a}, that is, $y$
lies inside $X_{\mathfrak{a}}$, inside a class component of step $k_a$, or inside a piece live
at $\mathfrak{B}_\tau$, or $\mathfrak{B}'$ is of density $k_a$. In the last case any such $y$ is
dated, by \eqref{4.27:eq-density-increase-a}(4a) (the statement on current thresholds across an
increase of density, Lemma~\ref{4.27:prf-density-increase}, whose proof is incomplete at one
step) and by \eqref{4.27:eq-unheld-raw-points-a} (the statement on unheld raw points within one
density, Lemma~\ref{4.27:prf-unheld-raw-points}, whose proof is incomplete at one step). At
such a point, as long as $y$ is raw at $\mathfrak{B}'$, the dating of the index-zero threshold
--- frozen at the innermost live piece where one holds $y$
(Definition~\ref{4.27:def-arrested-spaces}, and the definition of
$\mathfrak{G}$ in (e) of Lemma~\ref{4.27:lem-weights-heredity}, whose proof is incomplete at one
step), and unchanged within one
density where none does --- gives, together with
\eqref{4.27:eq-held-raw-points-a}(1), \eqref{4.27:eq-density-increase-a} and
\eqref{4.27:eq-unheld-raw-points-a} (the last two from statements whose proofs are incomplete
at one step),
\begin{equation}\label{4.27:eq-arrest-theorem-1}
T^{\mathfrak{B}'}(y)=T^{\mathfrak{B}^t}(y)=T^{\mathfrak{B}_\tau}(y).
\end{equation}
Since $\tilde p\in\mathfrak{G}(\mathfrak{B}')$ and $y$ is raw at $\mathfrak{B}'$, the input
obeys $\eta^{3}|\mathbb{J}(y)|<T^{\mathfrak{B}'}(y)$, and by \eqref{4.27:eq-arrest-theorem-1}
this is $\eta^{3}|\mathbb{J}(y)|<T^{\mathfrak{B}_\tau}(y)$. Where $y$ enters
$\Omega_1(\mathfrak{B}')$, the entering line of \eqref{4.27:eq-held-raw-points-a}(2) serves.
This is item (3).

\emph{Item (4).} This is the last clause of (e4) of Lemma~\ref{4.27:lem-weights-heredity},
whose proof is incomplete at one step, applied as in item (1): at the points raw at
$\mathfrak{B}^t$ that lie on $Y'$, $\tilde q$ obeys the $G$-part's index-zero cube gauges
$(3.35)_0$ wherever the map of $\mathfrak{O}$ is a layer into $U'$, namely for (O1), (O3),
(O4), (O4$'$), and, at the completion site (L), at the points served by the $\mathbf{R}$-layer
input (one of the inputs of Lemma~\ref{4.27:lem-t-step-layer}, whose proof is incomplete at one
step); the presented pair
of that input is the identity times layers and keeps the cube gauges. The index-zero cubes
carry no date, so these gauges are the
index-zero gauges of $\mathfrak{G}(\mathfrak{B}_\tau)$ at points raw at $\mathfrak{B}_\tau$.
At the completion site this requires that the polydisc of the $\mathbf{R}$-layer input carry
the shell of index
zero; this proviso is part of the second locus marked $(\ast)$ below and is not asserted.

$(\ast)$ The argument does not reach the two loci listed in (A$'$). At the first, the raw
points of $X_\tau\setminus Y'$ in components of $Z_{k_a}$ that are not class components of step
$k_a$ and lie outside every piece live at $\mathfrak{B}_\tau$, with the source of
$\mathfrak{O}$ of a density $>k_a$, none of the datings used for item (3) applies, and across
an increase of density the maximum defining $\gamma_0^{(k)}$ in (c) of
Lemma~\ref{4.27:lem-weights-heredity}, whose proof is incomplete at one step, may grow; so the
threshold identity
\eqref{4.27:eq-arrest-theorem-1} is not available there, and the current member is not
asserted.
At the second, on $Y'$ at points raw at $\mathfrak{B}^t$, the map of $\mathfrak{O}$ moves the
member $(3.37)_0$ for $A'$ and the current member, and at the completion site (L) the pair is
re-presented at the raw points of $\mathsf{Z}$ and of the core of the completion site; none of
these members is asserted.

\emph{Terms of kind $(\beta 3)$.} Let $\tau$ be of kind $(\beta 3)$ of
Lemma~\ref{4.27:lem-interrupted-terms}, so that $\mathfrak{B}_\tau=\mathbb{B}^{(n+1)}_{k_a}$
and $\mathfrak{B}^{\circ}=\mathfrak{B}_\tau\restriction X_{\mathfrak{a}}$. Items (1)--(4)
supply, on the same scope, the continuation required by the localised change-of-variables
input: its clauses on
$X_{\mathfrak{a}}\cap\Omega_1(\mathfrak{B}_\tau)$ and at the cubes come from item (1). The raw
points of $X_{\mathfrak{a}}$ are raw at $\mathfrak{B}^t$ while $\mathfrak{a}$ lives, by
\eqref{4.27:eq-density-increase-a}(1) (from a statement whose proof is incomplete at one step);
there item (3) gives the $U$-members off $Y'$, and item (4) the gauge $(3.35)_0$ on $Y'$ where
the map is a layer into $U'$. For the current member off $Y'$: by the freezing of the
index-zero datum, the threshold on $X_{\mathfrak{a}}$ is $\gamma_0^{(k_c)}$ for the
innermost live arrest $\mathfrak{c}$ holding the point, $k_c$ being the density of
$\mathfrak{c}$, and $k_c\le k_a$; as $\gamma_0^{(k)}$ is a running maximum over $j\le k$
((c) of Lemma~\ref{4.27:lem-weights-heredity}, whose proof is incomplete at one step), it is
non-decreasing in $k$, so $T\le\gamma_0^{(k_a)}$ on $X_{\mathfrak{a}}$; the points of
$X_{\mathfrak{a}}$ are dated, and
item (3) applies. $(\ast)$ The members moved on $X_{\mathfrak{a}}\cap Y'$ belong to the second
locus and are not asserted.

\emph{Clause (B).} The conclusion is that of Lemma~\ref{4.27:lem-real-minimiser}, whose proof
is incomplete at one step; its rider is included as $\mathfrak{G}_1(\mathfrak{B})$ by (e4) of
Lemma~\ref{4.27:lem-weights-heredity}, whose proof is incomplete at one step, which places the
real point in $\mathfrak{G}_1(\mathfrak{B})$. At raw points of index zero of
$\Omega_0(\mathfrak{B})$ the real point has $A'=0$, so its member $(3.37)_0$ holds.
$(\ast)$ Its index-zero gauge $(3.35)_0$ and its current member
$\eta^{3}|\mathcal{J}(U)|<T^{\mathfrak{B}}$ at these points are not asserted.

\emph{Clause (C).} By part (d) of the post-arrest input, one of the inputs listed in
Lemma~\ref{4.27:lem-t-step-layer}, whose proof is incomplete at one step, every operation after
an arrest that composes a stored
term with a map of $(\mathbb{U},\mathbb{J})$ is one of (O1)--(O5), (O4$'$), and every term it
creates receives its polydisc by the arrested reading of
Definition~\ref{4.27:def-arrested-spaces}. By the definition of legitimacy in (C), the
polydisc of such a term is legitimate when the polydisc of the stored term it composes, read
with its ambient rider, holds the image under the map of $\mathfrak{O}$ of the arrested source
of $\mathfrak{O}$; by Definition~\ref{4.27:def-operations-zones} this is the requirement placed
on a stored term at every later operation. At the points of $X_\tau$ this requirement is (A),
and for the ambient rider it is the existence of the continuation $\tilde q$ of (A$'$).
Legitimacy of the newborn polydiscs therefore holds as far as (A$'$) reaches.
$(\ast)$ At the two loci not asserted in (A$'$) it gives nothing for the riders of old terms,
and the first clause of the proposition on the riders of newborn terms inherits this
restriction.

\emph{Clause (D).} Off live pieces the arrested spaces of
Definition~\ref{4.27:def-arrested-spaces} are the spaces \cite[(2.34)--(2.39)]{B-CMP119} at
density $k'$ without change, and their rider is understood as in the ambient rider of that
definition. On $W_{\mathfrak{a}}$, the second inclusion is the descent
$\mathfrak{S}_{n+1}\subset\mathfrak{S}_n$ of the descent input, applied $n^+$ times in the
arrested
reading, in which source and target read one frozen structure. The first inclusion is
Lemma~\ref{4.27:lem-frozen-room}(a) at $n=n^+$: the class polydisc of stage $n^+$ exceeds
$\tilde{U}^{\mathrm{arrest}}_{k'}$ by the room $\rho_{n^+}(x;k')\ge0$,
\eqref{4.27:eq-frozen-room-a}. That every arrested space lies in
$\mathfrak{N}^{\vee}\subset\mathfrak{K}_{P}$ is (d) of Lemma~\ref{4.27:lem-weights-heredity},
whose proof is incomplete at one step. Finally, the arrested factor $F^{(k')}(x)$ is the factor
$s^{(k')}_{\ell(x)}$ of the unarrested construction less
$\beta\sum_{\mathfrak{a}\ \mathrm{live}\ni x}\sigma_{\mathfrak{a}}(\ell(x))$; where source and
target read one structure this sum is the same for both and does not depend on the stage
(the stage-map input), so it cancels in every difference of factors, which is therefore that of
the
unarrested construction.

\emph{Clause (E).} By the rule on $\gamma$ accompanying (d) of
Lemma~\ref{4.27:lem-weights-heredity}, whose proof is incomplete at one step, a membership is
read at $\gamma_0^{(k)}$, while every constant, floor and ceiling is read at
$\bar\gamma_0(\mathsf{t})$, under one supremum over $\mathsf{t}\ge\mathsf{t}_\gamma$ jointly
with the other scaffold functions at the same $\mathsf{t}$. The hypotheses of
Definition~\ref{4.27:def-arrest-hypotheses} are floors on $L$ and on the auxiliary parameters,
and ceilings $\gamma\le\gamma_{\mathrm{arrest}}$, each read as a supremum over
$\mathsf{t}\ge\mathsf{t}_\gamma$ of scaffold functions of one $\mathsf{t}$, with the parameters
fixed before $\gamma$ in the order of choice; nothing in them bounds $L$, $k$, $K$, $N$, the
ages or the nest depth
from above, or the coupling from below. The arrested reading of
Definition~\ref{4.27:def-arrested-spaces} involves no slot. Hence no constant depends on the
quantities listed in (E), $L$ enters only through floors, and the coupling only through
ceilings in one variable $\mathsf{t}$.
\end{proof}

\begin{lemma}[Room, spend and ratio outside dissolving pieces]\label{4.27:prf-room-regular-zones}
Assume the hypotheses of the arrest theorem, Theorem~\ref{4.27:thm-arrest-theorem}, whose proof is
incomplete at one step. Let $\mathfrak{a}$ be a live arrest, $\tau\in\mathfrak{F}(\mathfrak{a})$,
$\mathfrak{O}\in\{\mathrm{O}1,\dots,\mathrm{O}5,\mathrm{O}4'\}$ a later operation and $x\in X_\tau$, as there. Then at every
point $x$ of the zones Z1, Z2 and Z3 of~\eqref{4.27:eq-operations-zones-a} the following rows of the
table~\eqref{4.27:eq-arrest-theorem-a} hold:
\begin{itemize}
\item in zone Z1, the rows for O1 and O4, for O4$'$, for O2 and for O3;
\item in zone Z2, the rows for O2, for O3, for O4 and O4$'$, and the descent row Z2$^{\circ}$;
\item in zone Z3, the scale-drop row Z3a, the raised row Z3b and the consumer row Z3-L.
\end{itemize}
That is, in each of these rows the room of the clause of $\tau$ at $x$ over the image of the source
of $\mathfrak{O}$ is at least, the spend is at most, and the ratio of spend to room is at most, the
entries of the row.
\end{lemma}

\begin{proof}
All quantities are measured in the units at $x$ of the table~\eqref{4.27:eq-arrest-theorem-a}. We
write $\ell$ (zone Z1) or $m$ (zones Z2 and Z3) for the level of $x$, and $\mu$ for the level from
which a raised point is raised; we put
$u:=2^{-(k'+1-\ell)}$, and read $\rho_n$ as $\rho_{\max(n,n^+_W)}$, with $\rho_n(x;k')$ the room of
frozen domains defined in~\eqref{4.27:eq-frozen-room-a}. Here $n$ is the stage whose class domain
is the clause of $\tau$ at $x$ and $N_{k_a}$ is the number of stages of step $k_a$, so that
$0\le n\le N_{k_a}$; $W=W_{\mathfrak{a}}$ and $n^+_W$ are as in Lemma~\ref{4.27:lem-frozen-room};
$k'\ge k_a$ is the density at which $\mathfrak{O}$ reads its source; and $k_b>k_a$ is the density
of a live arrest $\mathfrak{b}\ni x$ enclosing $\mathfrak{a}$. The ratio of a row is the quotient
of its spend by its room.

\emph{Zone Z1, operations O1 and O4.} Let the operation act at stage $l$. By
Lemma~\ref{4.27:lem-live-piece-increments} the entries
of the clause move by at most $\tfrac12\rho^{(l)}_\ell$. We compare this with the one-step room,
that is, the room of the arrested stage space $\mathfrak{S}_l$ inside $\mathfrak{S}_{l-1}$,
which is $\rho^{(l)}_\ell$: the range sums $\sigma$ of the live pieces containing $x$ enter the
factors of both stages, do not depend on $l$, and therefore cancel in the difference. The ratio
is thus at most $\tfrac12$. The target $\mathfrak{S}_{l-1}$ is contained in the
domain of $\tau$ by Lemma~\ref{4.27:lem-frozen-room}(c) and~(b): part~(c) gives a nonnegative room
against the stage spaces
of an enclosing attempt and against sibling spaces of equal or higher stage, and part~(b) covers
the raised points. Two cases occur. On a running sibling the operation is the stage step, which
maps the stage-$l$ space into the stage-$(l-1)$ space, and the inclusion is the descent of the
stage spaces, hypothesis~[I2] of Theorem~\ref{4.27:thm-arrest-theorem}. On a completing sibling
$\tau$ ceases to exist after
the operation; there the completion claim of hypothesis~[I3] of
Theorem~\ref{4.27:thm-arrest-theorem} is
applied by statement, at the floor $\varphi_{\min}=1-4\beta$. Points of older pieces belong to zone
Z4 and are not treated here.

\emph{Zone Z1, operation O4$'$.} Here $W_{\mathfrak{a}}\subset\mathfrak{T}_{k_b}$. On this set
$\tau$ enters the foreign part $\mathfrak{I}_l$ of the stage domain directly as one of its terms,
by the scope clause of hypothesis~[I3] of Theorem~\ref{4.27:thm-arrest-theorem}; hence the room is
$\rho_n(x;k_b)$, and Lemma~\ref{4.27:lem-frozen-room}(a)
with $k'=k_b>k_a$ gives $\rho_n(x;k_b)\ge\beta2^{-(k_a-\ell)}$. By
Lemma~\ref{4.27:lem-live-piece-increments} the spend is
$\sum_l\iota_l(x)\le\tfrac12\mathsf{w}_0\beta e^{-(k_b-1-\ell)}$. We have
$e^{-(k_b-1-\ell)}\le2^{-(k_a-\ell)}$; this uses $k_b-1\ge k_a$ and the inequality
$e^{-t}\le2^{-t}$, valid for $t\ge0$, at $t=k_b-1-\ell$. The spend is therefore at most
$\tfrac12\mathsf{w}_0\beta2^{-(k_a-\ell)}$, and the ratio is at most
$\tfrac12\mathsf{w}_0=\tfrac1{48}$.

\emph{Zone Z1, operation O2 at density $k\ge k_a$, the density $k_a$ included.} By
Lemma~\ref{4.27:lem-frozen-room}(a) the
room is at least $\beta2^{-(k_a+1-\ell)}$, and by Lemma~\ref{4.27:lem-live-piece-increments} the
spend is at most
$\bar P2^{-(k-1-\ell)^+}$. We claim
\begin{equation*}
2^{-(k-1-\ell)^+}\,2^{k_a+1-\ell}\le4 .
\end{equation*}
If $k-1-\ell\ge0$, the left side is $2^{k_a-k+2}\le4$, since $k\ge k_a$. If $k-1-\ell<0$, then
$\ell\ge k\ge k_a$ and the left side is $2^{k_a+1-\ell}\le2$. Hence the ratio is at most
$4\bar P/\beta$, the entry of the row.

\emph{Zone Z1, operation O3, after the enlargement.} Put $v:=2^{-(k_a+1-\ell)}$. Since $k'\ge k_a$,
we have $v\ge u$. In the defining formula~\eqref{4.27:eq-frozen-room-a} with $k'+1$ in place of
$k'$, the middle term is
\begin{equation*}
\beta\bigl(2^{-(k_a-\ell)}-2^{-(k'+1-\ell)}\bigr)=\beta(2v-u),
\end{equation*}
and the first sum contains the index $i=k_a+1$, because $n\le N_{k_a}$; this term contributes
$\beta v$.
Hence $\rho_n(x;k'+1)\ge\beta(3v-u)$. The enlargement uses $\theta_S\beta u$ of this room. What is
left is at least
\begin{equation*}
\beta(3v-u)-\theta_S\beta u\ge\beta(3u-u-\theta_Su)=(2-\theta_S)\beta u .
\end{equation*}
By Lemma~\ref{4.27:lem-live-piece-increments} the spend is at most
$\tfrac14(1-\theta_S)\beta u$. Since
$(1-\theta_S)/(2-\theta_S)\le\tfrac12$ for $\theta_S\ge0$, the ratio is at most $\tfrac18$.

\emph{Zone Z2.} In this zone the clauses of $\tau$ are those of its constituents $T$, with
diminished thresholds, together with the clause of the ambient space $\mathfrak{N}^{(i)}$. For
the latter the comparison is of $1.383$ against a source of at most $1.2$; this is part of
hypothesis~[I3] of Theorem~\ref{4.27:thm-arrest-theorem}. Apply the factor corollary in
hypothesis~[I3] at density $k_a$, together
with the descent of the factors $s^{(k)}_m$ in $k$. The room of the clause of $\tau$ over the
source at time $k'$ is then
\begin{equation}\label{4.27:eq-room-regular-zones-1}
\ge\tfrac12r_T+\beta\bigl(2^{-(k_a-m)}-2^{-(k'-m)}\bigr),
\qquad r_T\ge\beta\theta_S2^{-(k_a-m)} ,
\end{equation}
where $r_T$ is the room that the factor corollary attaches to the constituent $T$. The increments
$\iota_l$ that enter here are those of the step $k_a$, hence fixed numbers, independent of $k'$.

For O1, the required bound is the step of the inheritance statement of hypothesis~[I3] itself.

For O2 at density $k\ge k_a$, the spend is $\bar P2^{-(k-1-m)^+}$, and
$\bar P2^{-(k-1-m)^+}\le2\bar P2^{-(k_a-m)}$. Indeed, if $k-1-m\ge0$ then
$2^{-(k-1-m)}=2\cdot2^{-(k-m)}\le2\cdot2^{-(k_a-m)}$, and if $k-1-m<0$ then $m\ge k\ge k_a$, so
that $2\cdot2^{-(k_a-m)}\ge2>1$. By~\eqref{4.27:eq-room-regular-zones-1} the room is at least
$\tfrac12r_T\ge\tfrac12\theta_S\beta2^{-(k_a-m)}$, the entry of the row. The quotient of spend by
room is therefore at most $4\bar P/(\theta_S\beta)$. This is at most $\tfrac14$ by the bound
$4\bar P/\beta\le\theta_S/4$ entered in the row for O2 in zone Z1.

For O3, after the enlargement, the source is at time $k'+1$, so
$u=2^{-(k'+1-m)}\le\tfrac12\cdot2^{-(k_a-m)}$. The bound~\eqref{4.27:eq-room-regular-zones-1},
with $k'+1$ in place of $k'$, gives a room of at least
\begin{equation*}
\tfrac12r_T+\beta\bigl(2^{-(k_a-m)}-u\bigr)\ge\beta\theta_Su+\beta(2u-u)=\beta(1+\theta_S)u .
\end{equation*}
The enlargement uses $\theta_S\beta u$, which leaves at least $\beta u$. The spend is at most
$\tfrac14(1-\theta_S)\beta u$. This follows from the floor of hypothesis~[I5] of
Theorem~\ref{4.27:thm-arrest-theorem}, with the
entries numbered two and four bounded as in the proof of
Lemma~\ref{4.27:lem-live-piece-increments} for O3. The ratio is at most
$\tfrac14(1-\theta_S)$, and this is at most $\tfrac18$ because $\theta_S\ge\tfrac12$, the
inequality underlying the case O4, O4$'$ below.

For O4 and O4$'$ the point $x$ lies in $\mathfrak{T}_{k_b}$. The source is at time $k_b\ge k_a+1$,
so $2^{-(k_b-m)}\le\tfrac12\cdot2^{-(k_a-m)}$, and~\eqref{4.27:eq-room-regular-zones-1} gives a
room of at least
\begin{equation*}
\bigl(\tfrac12\theta_S+\tfrac12\bigr)\beta2^{-(k_a-m)}\ge\tfrac34\beta2^{-(k_a-m)} .
\end{equation*}
Against the spend $\tfrac12\mathsf{w}_0\beta2^{-(k_a-m)}$ of the row, the ratio is
$\tfrac12\mathsf{w}_0\big/\tfrac34=\tfrac1{36}$.

\emph{Zone Z2$^{\circ}$ (descent).} The point $x$ is unraised, so its weight is $w=1$ by
Lemma~\ref{4.27:lem-weights-heredity}(a), whose proof is incomplete at one step. The stage step
maps $\mathfrak{S}^{(k)}_l$ into $\mathfrak{S}^{(k)}_{l-1}$, and by the
descent of the stage spaces, hypothesis~[I2] of Theorem~\ref{4.27:thm-arrest-theorem},
$\mathfrak{S}^{(k)}_{l-1}\subset\mathfrak{S}^{(k)}_0$. The factor of the
image is at most $s^{(k)}_m$, and $s^{(k)}_m\le s^{(k_a)}_m$. It therefore lies below the clause of
$\tau$, with the room of zone Z2 to spare; on the collar $\mathfrak{c}_{\mathfrak{a}}$ this room is
at least $\beta/2$. The one-step room of the stage sequence of step $k$ itself covers the move.
The room of zone Z2 is
left untouched, and the spend and the ratio are $0$.

\emph{Zone Z3, scale drop (row Z3a).} We have
\begin{equation*}
(1-\beta)\chi_gL^g\ge0.8\cdot\bigl(\tfrac78\bigr)^2\cdot8=4.9\ge2 ,
\end{equation*}
with $\chi_g$ the radius ratio of~\eqref{4.27:eq-factors-radii-coupling-a}; here we use that
$g\mapsto g^{-1/2}L^g$ is increasing, so that the smallest value is the one at $g=1$. Hence the
factor corollary in hypothesis~[I3] of Theorem~\ref{4.27:thm-arrest-theorem} gives the threshold
$\theta^{(i)}_{T,a}\ge1.5\,r_a(m)$. The source is at most $(1+\beta)r_a(m)$, which gives the room
$1.5-1.2$ of the row. The spend is at most $\max\{\beta/8,\bar P\}$: the first bound is given by
the input constant $\pi_{\mathrm{top}}$ of Theorem~\ref{4.27:thm-arrest-theorem}, the second by
hypothesis~[I4] of that theorem. Against the room $0.3\,r_a(m)$ the ratio is at most
$\max\{\beta/8,\bar P\}/0.3$, which is at most $\tfrac1{12}$, the entry of the row; for the first
member this is the inequality $\beta/2.4\le\tfrac1{12}$, that is $\beta\le0.2$, the bound already
used in $1+\beta\le1.2$.

\emph{Zone Z3, raised points (row Z3b).} By the timing and nesting relations (T3)
of~\eqref{4.27:eq-arrest-timing-a}, a point lies in no live piece at the moment it is raised. Let
it be raised by $g$ levels from the level $\mu$. Then the source threshold satisfies
$\theta_{\mathrm{src}}\ge0.8\,r_a(\mu)$, and so
\begin{equation*}
\theta^{(i)}_{T,a}\ge\theta_{\mathrm{src}}-2\mathsf{w}_0\beta r_a\ge0.78\,r_a(\mu).
\end{equation*}
By Lemma~\ref{4.27:lem-weights-heredity}(a), whose proof is incomplete at one step, the raised
weighted clause is bounded by
\begin{equation*}
\mathsf{L}_0^{2g}r_a(\mu+g)\le3^{-g}\chi_g^{-1}r_a(\mu)\le\tfrac13\bigl(\tfrac87\bigr)^2r_a(\mu)
=\tfrac{64}{147}\,r_a(\mu)\le0.436\,r_a(\mu) .
\end{equation*}
The factor $3^{-g}\chi_g^{-1}$ is largest at $g=1$, because $g\mapsto g^{1/2}3^{-g}$ is
decreasing. Consequently the clause of the class part $\mathfrak{L}_{n'}$ of the stage domain
implies the clause of $\tau$. This is the inclusion of hypothesis~[I3] together with
hypothesis~[I7] of Theorem~\ref{4.27:thm-arrest-theorem}, read for the
diminished thresholds. The ratio is at most $0.436/0.78\le\tfrac23$.

\emph{Zone Z3, consumers (row Z3-L).} Apply the entry bound for consumers, hypothesis~[I6] of
Theorem~\ref{4.27:thm-arrest-theorem},
in its plain form, to the old consumer $T\in\mathcal{C}$. Each entry is at most
$\theta_T/16$. Moreover
\begin{equation*}
\theta^{(i)}_T\ge\tfrac23\theta_T-2\mathsf{w}_0\beta\,\mathbf{r}^T\ge\tfrac12\theta_T ,
\end{equation*}
where $\mathbf{r}^T$ denotes the threshold scale of the consumer $T$ entering that bound.
So the room is at least $\tfrac12\theta_T$, the spend is at most $\theta_T/16$, and the ratio is
at most $\tfrac18$.
\end{proof}

\begin{lemma}[Room, spend and ratio at dissolving pieces]\label{4.27:prf-room-dissolving}
This is the line for dissolving pieces of the table \eqref{4.27:eq-arrest-theorem-a} of the
arrest theorem (Theorem~\ref{4.27:thm-arrest-theorem}, whose proof is incomplete at one
step). Let $\mathfrak{c}$ be a live piece of density $k_{\mathfrak{c}}$, read at the step
at which it dissolves. Let $h=h_{\mathfrak{b}}$, in the notation of
Lemma~\ref{4.27:lem-arrest-timing}, for the later attempt $\mathfrak{b}$ at which
$\mathfrak{c}$ dissolves, so that $k_{\mathfrak{c}}\le h$ by clause (T2) of that lemma.
Write $k$ for the level of the step being performed and $\mathbb{B}''_k$ for its
determining set. Put
\begin{equation*}
D_2:=(Z'_h)^{\sim 2}\cap\Lambda,\qquad S:=\partial(\Omega''_{h+1})^{\sim}.
\end{equation*}
Here ${}^{\sim}$ denotes the enlargements of Ba{\l}aban's notation, and $\Lambda$ is the region
of that name in the sealing construction of the dissolution step. The set $B_S$ is the subset
of $S$ fixed by that construction which the pieces considered avoid.
The sealing construction is the construction carried out at the dissolution step on the region
$D_2$ beside $S$; it fixes the configuration $U_2$, a function of the raw configuration $V''$
(see the proof), together with the set $B_S$, the objects $\mathfrak{b}_3$, $\mathfrak{b}_4$ and
the set $\Gamma_\flat$ used in the proof.
Let $x\in W_{\mathfrak{c}}$,
and let $\tau$ be a term whose clause is read at $x$. Write $o(x)$ for the level of $x$ in
$\mathbb{B}''_k$. The room of the clause of $\tau$ at $x$ is its threshold. For an entry $a$,
every threshold of that clause has the form
\begin{equation*}
\theta:=\theta_{\tau,a}(x)=F_\tau\,\mathsf{w}\,r_a(n) ,\qquad \mathsf{w}\ge 1 .
\end{equation*}
Here $r_a$ is the threshold function of the entry $a$. The number $\mathsf{w}$ is the weight
factor of the threshold of that entry at $x$, that is, the weight by which the thresholds
$\theta^{\mathfrak{B}}_a$ multiply the threshold function; it is not one of the constants
$\mathsf{w}_M$, $\mathsf{w}_0$, $\mathsf{w}^{\mathrm{on}}$. The threshold function is evaluated at a
level $n$ with $n\le o(x)$ and $o(x)\le h$, and this level is the current level of $x$ or a
finer birth index. The number $F_\tau$ is the factor of $\tau$. Then the following hold.
\begin{itemize}
\item[(a)] $F_\tau\ge\varphi_{\mathrm{arrest}}$, where
$\varphi_{\mathrm{arrest}}=\varphi_\flat:=\min\{1-\beta,\,1-4.51\beta\}$.
\item[(b)] On every cell of $D_2\setminus B_S$, the
base pair $(U_2,\mathcal{J}(U_2))$ satisfies, for every entry $a$, the entry $a$ of the clause
with bound at most $\theta_{\tau,a}(x)/64$.
\item[(c)] The increments together take at most $3\theta/64$.
\item[(d)] Consequently the spend is at most $\theta/16$, and the ratio is $1/16$.
\end{itemize}
During the stages of the chain at which $\mathfrak{c}$ dissolves, the clause read is that of
the line of the same table for
terms whose domains are assigned as in Lemma~\ref{4.27:lem-frozen-room}, taken at the operation
$O4$.
\end{lemma}

\begin{proof}
\emph{Place.} Clauses (T2) and (T4) of Lemma~\ref{4.27:lem-arrest-timing} (see
\eqref{4.27:eq-arrest-timing-a}) give
\begin{equation*}
X_{\mathfrak{c}}\subset Z_{k_{\mathfrak{c}}}\subset Z_h .
\end{equation*}

By the placement statement of the sealing construction for pieces beside $S$, the set
$W_{\mathfrak{c}}$ lies two layers
of cubes of side $LMR_{h+1}$ inside $S$, and it lies off $B_S$.

At such points the pair $(\mathbb{U},\mathbb{J})$ (field and source one-form) of the
presentation in force at the dissolution step reads the raw configuration $V''$. Indeed, by
\cite[(1.50), (1.53)]{B-CMP122b}, the configuration $U''_k$, and with it the configuration
$U_2$ of the sealing construction, reads the raw configuration $V''$ through its restriction
\begin{equation*}
V''\restriction\bigl((\Omega''_{h+1})^{\sim 2}\bigr)^{\complement}.
\end{equation*}
In addition, the integral over the region $Z_h$ is left unchanged in the subsequent
discussion of \cite{B-CMP122a}. This raw configuration is rough at its own level, so no size
bound of level $k$ holds there. The bounds at $x$ therefore come from the base and the
increments below.

\emph{Threshold.} Four cases occur.
\begin{enumerate}
\item The clause of $\tau$ at $x$ is the class part of a stage domain,
in the sense of Lemma~\ref{4.27:lem-frozen-room}. Its threshold is then the one given by
Lemma~\ref{4.27:lem-frozen-room}.
\item The clause is one on an older piece lying inside $X_{\mathfrak{c}}$. Its factor is
then at least $\varphi_{\mathrm{arrest}}$, by Lemma~\ref{4.27:lem-factor-positivity}.
\item There is a piece $\mathfrak{a}\subsetneq\mathfrak{c}$ and
$x\in W_{\mathfrak{c}}\setminus W_{\mathfrak{a}}$. The clause is then a clause of finer
birth index, with $\mu_\tau(x)\le k_{\mathfrak{a}}$, where $\mu_\tau(x)$ denotes the birth
index of $\tau$ at $x$.
\item The clause is that of a diminished constituent. Its threshold is then at least the
threshold $\theta_{\mathrm{src}}$ of its source term, by the corollary bounding the threshold
of a diminished constituent below by that of its source term. Here $F_\tau$ is the
constituent's own factor.
\end{enumerate}
In each of the four cases $F_\tau\ge\varphi_\flat$. Since $\beta>0$, the minimum defining
$\varphi_\flat$ is attained at its second member. Hence $\varphi_\flat=1-4.51\beta$, which
is $\varphi_{\mathrm{arrest}}$ of \eqref{4.27:eq-factor-positivity-a}. This proves (a).

\emph{Base.} Apply the own-level bound for the configuration $U_2$ of the sealing
construction, in its clause on cells, on every cell of $D_2\setminus B_S$. It states that the
pair $(U_2,\mathcal{J}(U_2))$ satisfies the defining inequalities of the analyticity spaces,
derived entries included. It does so with bound at most $1/64$ of every threshold
$F\,\mathsf{w}\,c\,r_a$ whose level is the current level of $x$ or a finer birth index, for
all weight factors $\mathsf{w}\ge 1$, all multipliers $c\ge 1$ and all $F\ge\varphi_\flat$.
The constant $1/64$ is independent of
$h-n$ and of the numbers of exclusions, blocks and strata.

By the threshold step, the thresholds of the clause of $\tau$ at $x$ are of this form, with
$c=1$ and $F=F_\tau\ge\varphi_\flat$. On cubes, the cube clause of the same own-level bound is
met with the gauges of its first clause, under the last member of the ownership floor condition
of the sealing construction. The member $A'$ vanishes: $A'=0$. This proves (b).

\emph{Increments.} We use two statements.
\begin{itemize}
\item The corollary of the sealing construction that bounds the ratio of the increments
independently of the age applies, because $n\le o(x)$ and $o(x)\le h$.
\item The share statement of the sealing construction says the following for the objects
$\mathfrak{b}_3$ and $\mathfrak{b}_4$ of that construction (distinct from the attempt
$\mathfrak{b}$). On every bond, the three shares of that corollary
add up to at most $3/64$ of the threshold.
\end{itemize}
The first share is controlled under the bend floor condition of the sealing construction. The
second share, that of the source, and the third, that of the map $\mathbb{H}'$, are
controlled under the floors that the statement of the sealing construction on floors imposes
at its set $\Gamma_\flat$. This proves (c).

\emph{Spend and ratio.} By (b) and (c) the spend at $x$ is at most
\begin{equation*}
\frac{\theta}{64}+\frac{3\theta}{64}=\frac{4\theta}{64}=\frac{\theta}{16}.
\end{equation*}
Dividing by the room $\theta$ gives the ratio $1/16$ of (d).

\emph{No circle.} The own-level bound used in the base step and the corollary bounding the
ratio independently of the age used in the increments step use, from the arrest apparatus,
only two things:
\begin{itemize}
\item the definition of the arrested spaces, through the first three of their defining
conditions;
\item Lemma~\ref{4.27:lem-factor-positivity}.
\end{itemize}
None of their inputs cites the hybrid weight statement or the key stage statement. Hence the
appeals made to them in the base step and in the increments step do not return to the
statement being proved.
\end{proof}

\begin{lemma}[Newborn terms, space inclusions and continuation of stored domains.
Stated; proof incomplete at the steps marked $(\ast)$]
\label{4.27:prf-newborn-continuation}
This lemma consists of clauses (C), (D) and (E) of the arrest theorem,
Theorem~\ref{4.27:thm-arrest-theorem}, whose proof is incomplete at the steps marked $(\ast)$ below,
together with its clause (A$'$), on the continuation of stored domains. Assume the hypotheses of that
theorem.
The arrested reading of the spaces \eqref{4.27:eq-arrested-spaces-a} is in force from each arrest on.

Let $\mathfrak{a}$ be a live arrest of density $k_a$, and let $\tau\in\mathfrak{F}(\mathfrak{a})$ be a
stored term with domain $X_\tau$, born over the determining set $\mathfrak{B}_\tau$, of density
$k_\tau=k_a$. Let $\mathfrak{O}\in\{O1,\dots,O5,O4'\}$ be a later operation. The term $\tau$ is either a
new boundary term $\mathbb{C}^{(i)}_{k_a}$ of stage $i$, born at the date $(k_a;i)$ with
$1\le i\le n\le n^{+}$, or a composed term of $\mathfrak{a}$, born at the date $(k_a;n^{+})$; here
$n^{+}$ is the frozen stage index of $\mathfrak{a}$ and $n\le n^{+}$ is the index of the last stage
whose new boundary terms are stored in $\mathfrak{F}(\mathfrak{a})$. Composed terms occur for arrests at
the exit at which the stage $n+1$, during which $X_{\mathfrak{a}}$ runs, composes the terms
$\mathbb{C}^{(i)}_{k_a}$, $i\le n$; for these arrests $n^{+}=n+1$. The
operations $O1$, $O4$ and $O4'$ are stages $l$ of a run; $O2$ and $O5$ act at the completion site of a
density $k$, whose step has $N=N_k$ stages; $O3$ is the $\mathbf{T}$-step from density $k'$ to $k'+1$.
Write $\mathfrak{B}^{t}$
for the target set of $\mathfrak{O}$ (the set in force immediately before $\mathfrak{O}$'s map)
and $\mathfrak{B}'$ for
its source set (the set in force immediately after it), so that $\mathfrak{B}^{t}\preceq\mathfrak{B}'$
in the order of dates \eqref{4.27:eq-determining-set-order-a}. Let $p$ be the member, over
$\mathfrak{O}$'s source domain $Y'$, of the source space that the arrested reading assigns to
$\mathfrak{O}$ (as in the statement of Theorem~\ref{4.27:thm-arrest-theorem}), with continuation
$\tilde p\in\mathfrak{G}(\mathfrak{B}')$. For a pair
$q$ write $\mathbb{U}_q,\mathbb{J}_q$ for its field and source components.

Let $\tilde q$ be the pair that equals the image of $p$ under $\mathfrak{O}$'s map on $Y'$ and that,
off $Y'$, has $\mathbb{U}_{\tilde q}:=\mathbb{U}_{\tilde p}$ and
\begin{equation}\label{4.27:eq-newborn-continuation-1}
\mathbb{J}_{\tilde q}(y):=
\begin{cases}
0, & y\in(\Omega_0\setminus\Omega_1)(\mathfrak{B}^{t})\setminus(X_\tau\cup Y'),\\
\mathbb{J}_{\tilde p}(y), & \text{for all other }y\notin Y'.
\end{cases}
\end{equation}
A point of $(\Omega_0\setminus\Omega_1)(\mathfrak{B})$ is called \emph{raw at $\mathfrak{B}$}. It is
\emph{held at $\mathfrak{B}$} if it lies in a piece that is live at the date of $\mathfrak{B}$. At a raw
point $y$ the rider $\mathfrak{G}(\mathfrak{B})$ has the following index-$0$ members:
\begin{itemize}
\item the $U$-members: the gauge condition $(3.35)_0$, which is the instance at index $0$ of
\cite[(3.35)]{B-CMP99b} on the index-$0$ cubes, and the condition $(3.37)_0$, which is the instance at
index $0$ of \cite[(3.37)]{B-CMP99b} for $A'$;
\item the $\mathbb{J}$-member $\ell^{3}|\mathbb{J}(y)|<T^{\mathfrak{B}}(y)$.
\end{itemize}
We say that $\mathfrak{O}$'s map is a \emph{layer into $U'$} at a point if it leaves the $G$-part $U$ of
the field and the cube gauges unchanged there, moving only the remaining components of the pair.
Then:
\begin{enumerate}
\item[(a)] $\tilde q\in\mathfrak{G}_1(\mathfrak{B}^{t})$ and $\tilde q\in\mathfrak{G}_1(\mathfrak{B}_\tau)$.
At every point of $\Omega_1(\mathfrak{B}_\tau)$ and on every cube, $\tilde q$ obeys the clauses of
$\mathfrak{K}_{109}(\mathfrak{B}_\tau;\gamma_0^{(k_\tau)})$.
\item[(b)] Let $y\in(\Omega_0\setminus\Omega_1)(\mathfrak{B}_\tau)$.
\begin{itemize}
\item If $y\in\Omega_1(\mathfrak{B}^{t})$, then $\tilde q$ obeys at $y$ every index-$0$ member of
$\mathfrak{G}(\mathfrak{B}_\tau)$.
\item If $y$ is raw at $\mathfrak{B}^{t}$ and $y\notin Y'$, then $\tilde q$ obeys at $y$ every
index-$0$ $U$-member of $\mathfrak{G}(\mathfrak{B}_\tau)$. It obeys the $\mathbb{J}$-member at $y$ if
$y\notin X_\tau$, and if $y\in X_\tau\setminus Y'$ is a point of cases (1)--(3) of
\eqref{4.27:eq-density-increase-a} (Lemma~\ref{4.27:prf-density-increase}, whose proof is incomplete
at a marked step), or a point of its case (4) with $\mathfrak{B}'$ of density $k_a$.
\item If $y$ is raw at $\mathfrak{B}^{t}$ and $y\in Y'$, then $\tilde q$ obeys at $y$ the index-$0$
cube gauges $(3.35)_0$ of $\mathfrak{G}(\mathfrak{B}_\tau)$ wherever $\mathfrak{O}$'s map is a layer
into $U'$.
\end{itemize}
Not asserted: the $\mathbb{J}$-member at the points of $X_\tau\setminus Y'$ of case (4) of
\eqref{4.27:eq-density-increase-a} when $\mathfrak{B}'$ has density larger than $k_a$; at points raw at
$\mathfrak{B}^{t}$ on $Y'$, the members moved by the map (the member $(3.37)_0$ of $A'$ and the
$\mathbb{J}$-member), and, at the completion site, every index-$0$ member at the raw points of
$\mathsf{Z}$ and of the core (zone Z4$'$ of
\eqref{4.27:eq-operations-zones-a}).
\item[(c)] Let $\tau$ be a composed term of $\mathfrak{a}$ whose domain is prescribed by the inclusion
of the ambient domain (Lemma~\ref{4.27:prf-ambient-inclusion}, whose proof is incomplete at a marked
step), so that $n^{+}=n+1$ and
$\mathfrak{B}_\tau=\mathbb{B}^{(n+1)}_{k_a}$, and let
$\mathfrak{B}^\circ$ be the restriction of $\mathfrak{B}_\tau$ to $X_{\mathfrak{a}}$. Then the
continuation required by that inclusion is supplied by
$\tilde q$: on $X_{\mathfrak{a}}$, $\tilde q$ obeys the clauses of
$\mathfrak{K}_{109}(\mathfrak{B}^\circ;\gamma_0^{(k_a)})$ at the points of
$X_{\mathfrak{a}}\cap\Omega_1(\mathfrak{B}_\tau)$
and at the cubes; and, at the raw points of $X_{\mathfrak{a}}$, its $U$-members and its
$\mathbb{J}$-member, with threshold $\le\gamma_0^{(k_a)}$, off $Y'$, and its gauge $(3.35)_0$ on
$X_{\mathfrak{a}}\cap Y'$ where the map is a layer into $U'$. The members moved by the map on
$X_{\mathfrak{a}}\cap Y'$ are not asserted.
\item[(d)] Clauses (C), (D) and (E) of Theorem~\ref{4.27:thm-arrest-theorem} hold. In particular, the
domains that the arrested reading assigns to the terms born at $\mathfrak{O}$ are legitimate as far as
parts (a)--(c) reach; at the loci excluded there (the steps marked $(\ast)$ below) nothing is
asserted.
\end{enumerate}
\end{lemma}

\begin{proof}
The proof treats the clauses of Theorem~\ref{4.27:thm-arrest-theorem} in the following order: first
(C), then (D), then clause (A$'$) (parts (a)--(c)), and last (E).

\emph{Clause (C).} A term born at $\mathfrak{O}$ is a cluster functional of two kinds of arguments: old
terms composed with $\mathfrak{O}$'s map, and the operators $H_\sigma(b)$.

Consider the old terms first. The stored terms of arrests are holomorphic on the source space that the
arrested reading assigns to $\mathfrak{O}$ by two statements. One is clause (A) of
Theorem~\ref{4.27:thm-arrest-theorem}, whose proof is
incomplete at the steps marked $(\ast)$ below, used inductively as described at the end of this proof.
The other is clause (A$'$) of that theorem. That clause reaches only as far as parts (a)--(c) below
establish it, namely
membership in $\mathfrak{G}_1$ together with the stratum members derived in steps (3) and (4); at the
two loci marked $(\ast)$ in step (3) nothing is derived.

For the other old terms, those stored by no arrest, holomorphy is given by the bounds of the
construction in force off live pieces.
The domains of such terms on live pieces are arrested domains, by induction. Measured against the
enlarged source of the step from density $k'$ to $k'+1$, the rooms are as follows, with $u$ the unit in
which Lemma~\ref{4.27:lem-frozen-room} measures room:
\begin{itemize}
\item a term born at the $\mathbf{T}$-step of a density $k''\le k'$ has room at least
$(1+\theta_S)\beta u$;
\item a term born at an $\mathbf{R}$-step has room at least $(1-\theta_S)\beta u$, against a requirement
of $\tfrac14(1-\theta_S)\beta u$.
\end{itemize}
The operators $H_\sigma(b)$ are holomorphic on $\mathfrak{K}_P$. All bounds are suprema over the domain.

\emph{Clause (D).} It follows from three statements:
\begin{itemize}
\item part (a) of Lemma~\ref{4.27:lem-frozen-room} at the density $k'$;
\item the statement on the descending stage spaces $\mathfrak{S}_n(X)$;
\item part (d) of Lemma~\ref{4.27:lem-weights-heredity}, whose proof is incomplete at
a marked step.
\end{itemize}

\emph{Clause (A$'$): the pair $\tilde q$.} Let $\tau$, $\mathfrak{O}$, $p$,
$\tilde p\in\mathfrak{G}(\mathfrak{B}')$ and $Y'$ be as in the statement. That
$\mathfrak{B}^{t}\preceq\mathfrak{B}'$ is part of the analysis of the individual operations in the
lemma on the time order of determining sets (Lemma~\ref{4.27:lem-determining-set-order}).

On $Y'$, $\tilde q$ is the image of $p$. Off $Y'$ the image of a pair under $\mathfrak{O}$'s map is the
pair itself, by the representation of the images and by parts (a) and (f) of the change-of-variables
theorem. There $\tilde q$ has the field component of $\tilde p$, and its source component is that of
$\tilde p$ except at the raw points of $(\Omega_0\setminus\Omega_1)(\mathfrak{B}^{t})$ off
$X_\tau\cup Y'$, where it is $0$; see \eqref{4.27:eq-newborn-continuation-1}.

This choice is admissible for three reasons:
\begin{itemize}
\item the rider is existential;
\item the space $\mathfrak{K}_{109}$ admits as source component any $\mathfrak{g}^{\mathbb{C}}$-valued
one-form whose size is bounded in the sense of its $\mathbb{J}$-member, with no condition
$D^{*}\mathbb{J}=0$;
\item the value of $\tau$ depends on the pair on $X_\tau$ only, and there $\tilde q$ is either the image
(on $Y'$) or the input, which equals the image (on $X_\tau\setminus Y'$).
\end{itemize}

\emph{Step (1): $\tilde q\in\mathfrak{G}_1(\mathfrak{B}^{t})$.} We follow the proof of the continuation
statement in clause (e) of Lemma~\ref{4.27:lem-weights-heredity}, whose proof is
incomplete at a marked step.

On $Y'$ we use the bounds on the maps given there. Each direct entry of the input is below
$1.4\,\theta^{\mathfrak{B}'}_a\le1.4\,\theta^{\mathfrak{B}^{t}}_a$, by the threshold comparison in
\eqref{4.27:eq-weights-heredity-b} (proof incomplete at a marked step), and $\mathfrak{O}$'s map moves it
by at most $0.1\,\theta^{\mathfrak{B}^{t}}_a$. The corresponding entry of the image is therefore below
\begin{equation*}
1.4\,\theta^{\mathfrak{B}^{t}}_a+0.1\,\theta^{\mathfrak{B}^{t}}_a=1.5\,\theta^{\mathfrak{B}^{t}}_a
<5\,\theta^{\mathfrak{B}^{t}}_a,
\end{equation*}
the threshold of $\mathfrak{G}(\mathfrak{B}^{t})$.
The cube clause holds wherever $\mathfrak{O}$'s map is a layer into $U'$. There it follows from the
nested-cube statement \eqref{4.27:eq-nested-cubes-a} and the inclusion of the riders in clause (e) of
Lemma~\ref{4.27:lem-weights-heredity}, \eqref{4.27:eq-weights-heredity-b} (proof incomplete at a marked
step).

At the completion site the bounds are as follows.
\begin{itemize}
\item At the points served by the layer statement of the completion site, the increments are in the
units of $\mathfrak{B}^{t}$, and each entry stays below $1.5\,\theta^{\mathfrak{B}^{t}}_a$.
\item On $\mathsf{Z}\cap Z^{\mathrm{on}}$, by the consumer bound of the completion site, each entry of
the pair in which the configuration is presented at the completion site (below called the presented
pair) is at most one sixteenth of the threshold of its consumer; at a point
$y\in\mathsf{Z}\cap Z^{\mathrm{on}}$ the consumer is the stage-$N$ boundary term
$\mathfrak{E}:=\mathbb{C}^{(N)}_k(X)$ of a domain $X\ni y$. Such a domain exists through every point of
$\mathsf{Z}\cap Z^{\mathrm{on}}$, and one exists containing any given cube of
$\mathcal{Q}_{\mathfrak{B}^{t}}(y)$: the stage-$N$ terms are indexed by every domain $X$ of the class
$\mathfrak{D}_m$ of \cite[(1.63)]{B-CMP122a} with $X\cap(Z_k\setminus Z''_k)\neq\emptyset$; the band
$Z_k\setminus Z''_k$ meets every completed component in its collar; and the class $\mathfrak{D}_m$,
consisting of the connected unions of cubes of the cover at scale $m$, contains a domain through any
given point and any given cell of the band, and one containing any given cube of
$\mathcal{Q}_{\mathfrak{B}^{t}}(y)$. The clause of $\mathfrak{E}$ reads
$\mathbb{B}''_k=\mathfrak{B}^{t}$ at $y$ with a stage factor at most $1$, so that its threshold is at
most that of $\mathfrak{B}^{t}$. Hence each entry of the presented pair, the cube clause on the cubes
of $\mathcal{Q}_{\mathfrak{B}^{t}}(y)$ included, is at most $\tfrac1{16}\theta^{\mathfrak{B}^{t}}_a$.
\item On $\mathsf{Z}\cap\mathfrak{T}_k$ each entry of the presented pair is at most
$1.1\,\theta^{\mathfrak{B}^{t}}_a$.
\item On the core each entry of the presented pair is at most $\tfrac1{64}+\tfrac3{64}=\tfrac1{16}$ of
every threshold at
the current level $\ell_{\mathfrak{B}^{t}}(x)$ of the point $x$ or at a finer birth index; such a
threshold is the level threshold $r_a$ multiplied by a factor at least $\varphi_\flat$ and by a weight
and a further factor each at least $1$. Since the frozen weight $w_{\mathfrak{B}^{t}}(x)$ is a
non-negative power of $\mathsf{L}_0^{2}$, it is at least $1$, so
$\varphi_\flat\,w_{\mathfrak{B}^{t}}(x)\,r_a(\ell_{\mathfrak{B}^{t}}(x))$ is such a threshold.
In particular the entry is at most one sixteenth of
\begin{equation*}
\varphi_\flat\,w_{\mathfrak{B}^{t}}(x)\,r_a(\ell_{\mathfrak{B}^{t}}(x))
\le\theta^{\mathfrak{B}^{t}}_a(x),
\end{equation*}
so each entry, the cube clause included, is at most $\tfrac1{16}\theta^{\mathfrak{B}^{t}}_a(x)$.
\end{itemize}
All these bounds hold at points of $\Omega_1(\mathfrak{B}^{t})$ and on cubes of
$\mathcal{Q}_{\mathfrak{B}^{t}}$, and each lies below $5\,\theta^{\mathfrak{B}^{t}}_a$.

Off $Y'$, the clauses obeyed by $\tilde p$ carry over to $\mathfrak{B}^{t}$. This follows from the
threshold comparison and the inclusion of the riders in \eqref{4.27:eq-weights-heredity-b} (proof
incomplete at a marked step), since $\mathfrak{B}^{t}\preceq\mathfrak{B}'$.

Finally, $\tilde q$ differs from the continuation constructed in that proof only in the source
component at raw points off $Y'$. No clause of $\mathfrak{G}_1$ reads the source component at such
points. Hence $\tilde q\in\mathfrak{G}_1(\mathfrak{B}^{t})$.

\emph{Step (2): $\mathfrak{B}_\tau\preceq\mathfrak{B}^{t}$, and part (a).} The term $\tau$ was stored
before $\mathfrak{O}$.
\begin{itemize}
\item For $\mathfrak{O}=O1$ this is the inequality $l-1\ge n^{+}\ge i$.
\item For $O2,\dots,O5$ and $O4'$, the set $\mathfrak{B}^{t}$ is of a later date. For $O2$ with $k=k_a$
it is the same site at the later date $(k_a;N)$.
\end{itemize}
In all cases $\mathfrak{B}_\tau\preceq\mathfrak{B}^{t}$ follows from the analysis of the individual
operations in Lemma~\ref{4.27:lem-determining-set-order}.

Next, $\mathfrak{G}_1(\mathfrak{B}^{t})\subset\mathfrak{G}_1(\mathfrak{B}_\tau)$. This is the
inclusion of the riders in \eqref{4.27:eq-weights-heredity-b} (proof incomplete at a marked step),
composed along $\preceq$ by the composition of moves in Lemma~\ref{4.27:lem-determining-set-order}. Here
$\Omega_1(\mathfrak{B}_\tau)\subseteq\Omega_1(\mathfrak{B}^{t})$, because the domains
$\Omega_0(\cdot)$ and $\Omega_1(\cdot)$ never shrink along the date order. This monotonicity of the
domains is established in the proof of clause (e) of Lemma~\ref{4.27:lem-weights-heredity}, whose
proof is incomplete at a marked step. With step (1) we obtain
$\tilde q\in\mathfrak{G}_1(\mathfrak{B}_\tau)$.

The proof of part (d) of Lemma~\ref{4.27:lem-weights-heredity} (whose proof is
incomplete at a marked step) places the clauses of $\mathfrak{G}_1(\mathfrak{B}_\tau)$ inside those of
$\mathfrak{K}_{109}(\mathfrak{B}_\tau;\gamma_0^{(k_\tau)})$, at every point of
$\Omega_1(\mathfrak{B}_\tau)$ and on every cube. This proves part (a).

\emph{Step (3): the stratum, part (b).} By the monotonicity of the domains,
\begin{equation*}
(\Omega_0\setminus\Omega_1)(\mathfrak{B}_\tau)\subseteq\Omega_0(\mathfrak{B}^{t})
\subseteq\Omega_0(\mathfrak{B}').
\end{equation*}
Let $y$ be a point of $(\Omega_0\setminus\Omega_1)(\mathfrak{B}_\tau)$. We treat three cases.

\emph{(3a) $y\in\Omega_1(\mathfrak{B}^{t})$.} We apply the inclusion at entering points in clause (e)
of Lemma~\ref{4.27:lem-weights-heredity} (whose proof is incomplete at a marked step), with
$(\mathfrak{B}_\tau,\mathfrak{B}^{t})$ in place of $(\mathfrak{B},\mathfrak{B}')$, to
$\tilde q\in\mathfrak{G}_1(\mathfrak{B}^{t})$. It gives every index-$0$ member of
$\mathfrak{G}(\mathfrak{B}_\tau)$ at $y$, both the $U$-members and the $\mathbb{J}$-member.
\begin{itemize}
\item The $U$-members come from heredity on the raw stratum
(Lemma~\ref{4.27:prf-raw-stratum}).
\item At points unheld at $\mathfrak{B}_\tau$, the entering bound \eqref{4.27:eq-raw-stratum-a} gives
$\ell^{3}|\mathbb{J}_{\tilde q}(y)|<\gamma_0^{(k_a)}=T^{\mathfrak{B}_\tau}(y)$.
\item At points held at $\mathfrak{B}_\tau$, with $\mathfrak{c}$ the innermost live piece holding $y$,
the second case of \eqref{4.27:eq-held-raw-points-a}, that is, the case of points entering $\Omega_1$,
gives $\ell^{3}|\mathbb{J}_{\tilde q}(y)|<\gamma_0^{(k_{\mathfrak{c}})}=T^{\mathfrak{B}_\tau}(y)$.
\end{itemize}

\emph{(3b) $y$ raw at $\mathfrak{B}^{t}$ and $y\notin Y'$.}

\emph{The $U$-members.} Here $\mathbb{U}_{\tilde q}=\mathbb{U}_{\tilde p}$. Moreover
$\tilde p\in\mathfrak{G}(\mathfrak{B}')\subset\mathfrak{G}^{-J_0}(\mathfrak{B}_\tau)$, by the inclusion
of the riders \eqref{4.27:eq-weights-heredity-b} (proof incomplete at a marked step) composed along
$\mathfrak{B}_\tau\preceq\mathfrak{B}'$. So $\tilde q$ obeys every index-$0$ $U$-member of
$\mathfrak{G}(\mathfrak{B}_\tau)$ at $y$.

\emph{The $\mathbb{J}$-member off $X_\tau$.} If $y\notin X_\tau\cup Y'$, then
$\mathbb{J}_{\tilde q}(y)=0<T^{\mathfrak{B}_\tau}(y)$.

\emph{The $\mathbb{J}$-member on $X_\tau\setminus Y'$.} Here
$\mathbb{J}_{\tilde q}(y)=\mathbb{J}_{\tilde p}(y)$. Suppose first that $y$ is a point dated by the
statement on current thresholds across an increase of density (Lemma~\ref{4.27:prf-density-increase},
whose proof is incomplete at a marked step) in its cases (1)--(3) of
\eqref{4.27:eq-density-increase-a}. These are the cases (1) $y\in X_{\mathfrak{a}}$; (2)
$y\notin X_{\mathfrak{a}}$ and $y$ inside a piece live at $\mathfrak{B}_\tau$; (3) $y$ inside no such
piece but inside a class component $X'\neq X_{\mathfrak{a}}$ of step $k_a$.

If $y$ is raw at $\mathfrak{B}'$, then
\begin{equation}\label{4.27:eq-newborn-continuation-2}
\ell^{3}|\mathbb{J}_{\tilde p}(y)|<T^{\mathfrak{B}'}(y)=T^{\mathfrak{B}^{t}}(y)
=T^{\mathfrak{B}_\tau}(y).
\end{equation}
The equalities hold for the following reasons.

Lemma~\ref{4.27:prf-density-increase} (proof incomplete at a marked step) gives
$T^{\mathfrak{B}^{t}}(y)=T^{\mathfrak{B}_\tau}(y)$. It also gives that $y$ is held at
$\mathfrak{B}^{t}$ in cases (1) and (2), and in case (3) past the dates of density $k_a$. In those cases
the first case of \eqref{4.27:eq-held-raw-points-a} (points raw at both dates), applied to the pair
$(\mathfrak{B}^{t},\mathfrak{B}')$, gives $T^{\mathfrak{B}'}(y)=T^{\mathfrak{B}^{t}}(y)$.

In case (3) during density $k_a$ (that is, $\mathfrak{O}=O1$, or $O2$ with $k=k_a$), the set
$\mathfrak{B}'$ has density $k_a$. The only pieces that can hold $y$ at $\mathfrak{B}'$ are then arrests
of density $k_a$: an older piece holding $y$ and live at $\mathfrak{B}'$ would have been live at
$\mathfrak{B}_\tau$, which case (3) excludes. Hence
$T^{\mathfrak{B}'}(y)=\gamma_0^{(k_a)}=T^{\mathfrak{B}_\tau}(y)$.

For $O3$ from $k_a$ to $k_a+1$, the source date lies past every date of density $k_a$ (past the
completion date $(k_a;L)$ when a component completes there), and the target date is the last of them.
By that date the sibling component $X'$ containing $y$ has been either arrested or completed:
\begin{itemize}
\item if it has been arrested, $y$ is held at $\mathfrak{B}^{t}$;
\item if it has been completed, $\tau$ has ceased.
\end{itemize}
That every class component of the step is arrested or completed by then follows from the definition
of an arrest, from the description of the attempt at a step as ending with each class component
arrested or completed, and from the list of exits of the attempt. That $\tau$ has ceased in the second
case follows from the definition of the stored terms.

If instead $y$ enters $\Omega_1(\mathfrak{B}')$, then $\mathfrak{O}$ is $O2$ or $O5$, whose completion
move \textup{(t4)} acts on a component holding $y$. Let $k_e$ be
the density of the date at which $y$ enters. The second
case of \eqref{4.27:eq-held-raw-points-a} and the entering bound \eqref{4.27:eq-raw-stratum-a} give
\begin{equation*}
\ell^{3}|\mathbb{J}_{\tilde p}(y)|<5\alpha_{0,k_e}L^{-3k_e}\le\gamma_0^{(k'')}
\qquad\text{for every }k''\le k_e ,
\end{equation*}
and in particular $\ell^{3}|\mathbb{J}_{\tilde p}(y)|<T^{\mathfrak{B}_\tau}(y)$.

Now let $y$ be a point of case (4) of \eqref{4.27:eq-density-increase-a}. We distinguish the density of
$\mathfrak{B}'$.

\emph{$\mathfrak{B}'$ of density $k_a$} ($\mathfrak{O}=O1$, or $O2$ with $k=k_a$). The point $y$, raw at
$\mathfrak{B}^{t}$, is raw at $\mathfrak{B}'$: the moves \textup{(t2)} and \textup{(t4)} do not touch
the raw points of a
component that is not a class component. This is part (4a) of Lemma~\ref{4.27:prf-density-increase}
(proof incomplete at a marked step). By that part, no arrest of density $k_a$ holds $y$ and no older one
is live. The first part of \eqref{4.27:eq-unheld-raw-points-a}, from the statement on unheld raw points
within one density (Lemma~\ref{4.27:prf-unheld-raw-points}, whose proof is incomplete at a marked step),
then gives $T^{\mathfrak{B}'}(y)=\gamma_0^{(k_a)}=T^{\mathfrak{B}_\tau}(y)$. So the input satisfies
$\ell^{3}|\mathbb{J}_{\tilde p}(y)|<T^{\mathfrak{B}_\tau}(y)$.

\emph{$\mathfrak{B}'$ of density larger than $k_a$.} This is part (4b) of
Lemma~\ref{4.27:prf-density-increase} (proof incomplete at a marked step): $\mathfrak{O}$ is $O3$, $O2$
with $k>k_a$, $O4$, $O4'$ or $O5$.
At these points no bound of $\ell^{3}|\mathbb{J}_{\tilde p}(y)|$ by $T^{\mathfrak{B}_\tau}(y)$ is
derived $(\ast)$.

\emph{(3c) $y$ raw at $\mathfrak{B}^{t}$ and $y\in Y'$.} Wherever $\mathfrak{O}$'s map is a layer into
$U'$, the $G$-part $U$ of the image is that of the input. Indeed:
\begin{itemize}
\item the stage layers fix the $G$-part;
\item the layers at the completion site keep the $U$-part and the cube gauges;
\item the map of the $\mathbf{T}$-step multiplies the field by $e^{i\eta\mathfrak{h}}$, a layer on
$A'$ in the decomposition $\mathbb{U}=e^{i\eta A'}U$, which leaves $U$ and the cube gauges unchanged.
\end{itemize}
The input $\tilde p$ obeys $(3.35)_0$ there, since for the $U$-members
\begin{equation*}
\mathfrak{G}(\mathfrak{B}')\subset\mathfrak{G}^{-J_0}(\mathfrak{B}^{t})
\subset\mathfrak{G}^{-J_0}(\mathfrak{B}_\tau)
\end{equation*}
by the inclusion of the riders \eqref{4.27:eq-weights-heredity-b} (proof incomplete at a marked step).
Hence $\tilde q$ obeys the index-$0$ cube gauges of $\mathfrak{G}(\mathfrak{B}_\tau)$ at $y$.

Two groups of members are not derived $(\ast)$:
\begin{itemize}
\item the members that the map moves, namely the member $(3.37)_0$ of $A'$ and the $\mathbb{J}$-member;
\item at the completion site, every index-$0$ member at raw points of $\mathsf{Z}$ and of the core.
\end{itemize}

\emph{Step (4): part (c).} Let $\tau$ be a composed term of $\mathfrak{a}$, as in part (c), so that
$\mathfrak{B}_\tau=\mathbb{B}^{(n+1)}_{k_a}$. While $\mathfrak{a}$
lives, the set in force on $X_{\mathfrak{a}}$ is
$\mathfrak{B}^\circ=\mathfrak{B}_\tau\restriction X_{\mathfrak{a}}$. This is by the move \textup{(t3)},
and because the collar $\mathfrak{c}_{\mathfrak{a}}$ of $\mathfrak{a}$ (not one of the older pieces
$\mathfrak{c}$ considered below) keeps level $k_a$:
\begin{equation*}
X_{\mathfrak{a}}\subset\Lambda^{c}_{k_a}\subset\Omega^{c}_{k_a+1}.
\end{equation*}
Steps (2) and (3), restricted to $X_{\mathfrak{a}}$, give the following clauses of
$\mathfrak{K}_{109}(\mathfrak{B}^\circ;\gamma_0^{(k_a)})$.

\emph{On $X_{\mathfrak{a}}\cap\Omega_1(\mathfrak{B}_\tau)$ and at the cubes.} Here we use the proof of
part (d) of Lemma~\ref{4.27:lem-weights-heredity} (whose proof is incomplete at a marked step), with
the inputs $w_{\mathfrak{B}^\circ}(x)\le w_{*}(k_a)$ and $\ell_{\mathfrak{B}^\circ}(x)\le k_a$ on
$W_{\mathfrak{a}}$ supplied by its part (a). On the collar
$\mathfrak{c}_{\mathfrak{a}}=X_{\mathfrak{a}}\cap\Omega_{k_a}$ the structure of $\mathfrak{B}^\circ$ is
the stage structure of $\mathbb{B}^{(n+1)}_{k_a}$; its level there is $k_a$, by the displayed inclusion,
and its weight is the closed-form weight of \eqref{4.27:eq-weights-heredity-b} (proof incomplete at a
marked step), with level and $\mathbb{B}_{k_a}$-level both equal to $k_a$, namely
$\mathsf{L}_0^{2(k_a-\max\{k_a,k_0\})^{+}}=1$.
So the case of weight $1$ in the same proof applies there and gives, for both radii,
\begin{equation*}
5\alpha_{0,k_a},\ 5\alpha_{1,k_a}\le5\bar\alpha_{k_a}\le\gamma_0^{(k_a)}.
\end{equation*}
At the points of an older piece $\mathfrak{c}$ nested in
$X_{\mathfrak{a}}$ the same part, applied to $\mathfrak{c}$, gives $w\le w_{*}(k_{\mathfrak{c}})$ and
$\ell\le k_{\mathfrak{c}}$ with $k_{\mathfrak{c}}<k_a$, so that
$5w_{*}(k_{\mathfrak{c}})\bar\alpha_{k_{\mathfrak{c}}}$ lies inside the maximum over $j\le k_a$
defining $\gamma_0^{(k_a)}$.

\emph{At the raw points.}
\begin{itemize}
\item The $U$-members hold at the raw points of $X_{\mathfrak{a}}\setminus Y'$.
\item The gauge $(3.35)_0$ holds on $X_{\mathfrak{a}}\cap Y'$ where the map is a layer into $U'$.
\item The $\mathbb{J}$-member holds at the raw points of $X_{\mathfrak{a}}\setminus Y'$ with threshold
$T\le\gamma_0^{(k_a)}$. By case (1) of \eqref{4.27:eq-density-increase-a} (proof incomplete at a marked
step) such a point is held by $\mathfrak{a}$ or by an older nested piece, and
\[
T^{\mathfrak{B}'}(y)=T^{\mathfrak{B}^{t}}(y)=T^{\mathfrak{B}_\tau}(y)\le\gamma_0^{(k_a)}.
\]
\end{itemize}
The moved members on $X_{\mathfrak{a}}\cap Y'$ are not derived, as at the second step marked $(\ast)$.

The proof of clause (A$'$) (steps (1)--(4)) uses only the following inputs:
\begin{itemize}
\item the stage structures;
\item the timing, nesting and band statements for arrests \eqref{4.27:eq-arrest-timing-a};
\item the definitions of an arrest and of the stored terms;
\item clauses (a), (d) and (e) of Lemma~\ref{4.27:lem-weights-heredity} (proof incomplete at a marked
step);
\item the spends.
\end{itemize}
It reads no room of the table of zones \eqref{4.27:eq-operations-zones-a}, alters no constant of the
construction, and contains no circular appeal.

\emph{The date order used above.} $\mathfrak{B}^{t}$ is the set in force immediately before
$\mathfrak{O}$'s map and $\mathfrak{B}'$ the one after it, so $\mathfrak{B}^{t}\preceq\mathfrak{B}'$ by
one move (Lemma~\ref{4.27:lem-determining-set-order}). The relation
$\mathfrak{B}_\tau\preceq\mathfrak{B}^{t}$ is reached as follows, operation by operation.
\begin{itemize}
\item For $O1$, through the stages $i\le n^{+}\le l-1$. The siblings run from stage $i$ to stage
$l-1$. The held set $W_{\mathfrak{a}}$ runs from stage $i$ to stage $n^{+}$ and is then frozen; for an
arrest with $n^{+}=n+1$ the first later sibling stage is $l=n+2$.
\item For $O2$ with $k=k_a$, at the same site at the later date: from stage $i$ to stage $N$ on the
completing sibling, with $\mathfrak{B}^{t}=\mathbb{B}''_{k_a}$.
\item For $O2$ with $k>k_a$, for $O3$, $O4$, $O5$ and for $O4'$, through one $\mathbf{T}$-step move
\textup{(t1)} for each density from $k_a+1$ up to the density of $\mathfrak{B}^{t}$, and the moves
\textup{(t2)}--\textup{(t4)} of the intervening sites on their own classes.
\end{itemize}

\emph{Clause (E).} The clause concerns the constants that occur in (A)--(D) and in the statements cited
for them. It follows by inspection of these constants, together with two inputs:
\begin{itemize}
\item the positivity of the arrested factors (Lemma~\ref{4.27:lem-factor-positivity});
\item the rule by which $\gamma_0^{(k)}$ enters only as the radius of a membership.
\end{itemize}

\emph{Order of the induction.} The proof is an induction along the time-ordered sites. Clause (A) at
$\mathfrak{O}$ uses the inheritance statement for clause (A) at sites not later than $\mathfrak{O}$'s.
That inheritance statement at site $k$ uses only the arithmetic of Lemma~\ref{4.27:lem-frozen-room}
and the table of zones \eqref{4.27:eq-operations-zones-a}, for
pieces arrested before $k$.
\end{proof}

\begin{proposition}[Bounds of later steps on live arrested pieces. Stated; proof incomplete at the step marked $(\ast)$]\label{4.27:prop-live-piece-bounds}
The arrested spaces of \eqref{4.27:eq-arrested-spaces-a} are not contained in the standard spaces
of Ba{\l}aban's construction. The product of an arrested factor with a weight may exceed the
size parameter of the standard spaces, since
\begin{equation*}
\varphi_{\mathrm{arrest}}\,\mathsf{L}_0^{4}\ge 1.5
\end{equation*}
(see \eqref{4.27:eq-factor-positivity-a}). Bounds proved for consumer terms reach
live pieces as follows.
\begin{enumerate}
\item[(i)] A term born after the arrest $\mathfrak{a}$ (of density $k_{\mathfrak{a}}$) is stated on
its arrested analyticity domain, as in
\eqref{4.27:eq-arrested-spaces-a}. The proof of its bound reads that space through two properties
only. The first is that the old terms, composed with the map of a later operation
$\mathfrak{O}$, are holomorphic on the arrested source of $\mathfrak{O}$. This is clauses (A)
and (C) of the arrest theorem, Theorem~\ref{4.27:thm-arrest-theorem}, whose proof is incomplete
at one step, with the scope which
clause (C) inherits from the clause of that theorem fixing the scope of (A) for stored terms.
The second is that the operators $H_\sigma(b)$ act on
$\mathfrak{K}_P=\mathfrak{K}_{109}(\mathfrak{B};\gamma_0^{(k)})$, formed over the determining set
$\mathfrak{B}$ in force, of density $k$, in which every arrested space lies; this inclusion is
clause (D) of the same theorem.
\item[(ii)] A term $\tau$ born before the arrest, with lattice domain $X_\tau$, keeps its
analyticity domain over $X_\tau$ and its rider space at its birth set $\mathfrak{B}_\tau$, of
density $k_\tau=k_{\mathfrak{a}}$; every later source lies in that analyticity domain.
Precisely, let $\mathfrak{B}'$ be the determining set of a later source, so that
$\mathfrak{B}_\tau\preceq\mathfrak{B}'$. The rider space $\mathfrak{G}(\mathfrak{B}')$ in which
that source is read, by \eqref{4.27:eq-arrested-spaces-a}, satisfies the heredity
\eqref{4.27:eq-weights-heredity-b} of Lemma~\ref{4.27:lem-weights-heredity}, whose proof is
incomplete at one step, with the index-zero thresholds $T^{\mathfrak{B}_\tau}(y)$ and
$T^{\mathfrak{B}'}(y)$ frozen, at each raw point $y$ lying in a live piece, with the structure of
the innermost live piece containing $y$. This heredity gives the following three assertions.
\begin{enumerate}
\item[(a)] $\mathfrak{G}(\mathfrak{B}')\subset\mathfrak{G}^{-J_0}(\mathfrak{B}_\tau)$ on all of
$\Omega_0(\mathfrak{B}_\tau)$, namely at the points of $\Omega_1(\mathfrak{B}_\tau)$, at the cubes,
and for the $\mathbb{U}$-members of the stratum
$\Omega_0(\mathfrak{B}_\tau)\setminus\Omega_1(\mathfrak{B}_\tau)$.
\item[(b)] Every pair of $\mathfrak{G}(\mathfrak{B}')$ obeys every stratum member of
$\mathfrak{G}(\mathfrak{B}_\tau)$ at the points of the stratum which have since entered
$\Omega_1(\mathfrak{B}')$.
\item[(c)] Every pair of $\mathfrak{G}(\mathfrak{B}')$ obeys the $\mathbb{J}$-member of
$\mathfrak{G}(\mathfrak{B}_\tau)$, namely
\begin{equation*}
\ell^{3}|\mathbb{J}(y)|<T^{\mathfrak{B}_\tau}(y),
\end{equation*}
where $\ell$ denotes the lattice unit,
at certain points $y$ of $X_\tau$ which are raw at both dates and are dated by
\eqref{4.27:eq-density-increase-a} or \eqref{4.27:eq-unheld-raw-points-a} (statements whose
proofs are incomplete at one step).
These are the points $y$ that lie inside $X_{\mathfrak{a}}$, inside a class component of step
$k_{\mathfrak{a}}$, or inside a piece live at $\mathfrak{B}_\tau$. When $\mathfrak{B}'$ is of
density $k_\tau=k_{\mathfrak{a}}$, they are all the points of $X_\tau$ raw at both dates. At each
of these points
$T^{\mathfrak{B}'}(y)=T^{\mathfrak{B}_\tau}(y)$, and no map of a later operation intervenes.
\end{enumerate}
$(\ast)$ At the raw points of $X_\tau$ in components of $Z_{k_{\mathfrak{a}}}$ that are not class
components of step $k_{\mathfrak{a}}$, outside every live piece, the $\mathbb{J}$-member is not
asserted for sources $\mathfrak{B}'$ of a density $>k_{\mathfrak{a}}$.

A further inclusion holds at every point $x$ of a live piece of class $k_{\mathfrak{a}}$. Let
$\ell(x)$ be the level of $x$, let $n^{+}$ be the stage at which the piece was arrested (the
index $n^{+}_W$ of Lemma~\ref{4.27:lem-frozen-room}), let
$0\le n\le n^{+}$ be any stage, and let $k'\ge k_{\mathfrak{a}}$. Then
\begin{equation}\label{4.27:eq-live-piece-bounds-1}
\tilde{U}^{\mathrm{arrest}}_{k'}\subset\mathfrak{S}^{(k_{\mathfrak{a}}),\mathrm{arrest}}_{n^{+}}
\subset\mathfrak{S}^{(k_{\mathfrak{a}}),\mathrm{arrest}}_{n},
\end{equation}
and each of these spaces is contained in the analyticity domain $\mathfrak{D}^{\flat}_n(X_{\tau'})$
of every term $\tau'$ born before the arrest whose analyticity domain at $x$ is the class domain
of stage $n$. The room between $\tilde{U}^{\mathrm{arrest}}_{k'}$ and the analyticity domain of such a term
is at least
$\beta 2^{-(k_{\mathfrak{a}}+1-\ell(x))}$, in the units of
Lemma~\ref{4.27:lem-frozen-room}, with $\beta$ the constant of that lemma.
\item[(iii)] The hypothesis which the auxiliary lemma used in the proof of the correlation
theorem places on a live piece is exactly clause (A) of
Theorem~\ref{4.27:thm-arrest-theorem}, whose proof is incomplete at one
step.
\item[(iv)] The real domain lies inside the arrested spaces in the sense, and within the scope,
of clause (B) of the same theorem. For every determining set
$\mathfrak{B}$, the real point lies in
$\mathfrak{G}_1(\mathfrak{B})$. At the raw points of $\Omega_0(\mathfrak{B})$ it satisfies the
index-zero instance of \cite[(3.37)]{B-CMP99b}. Two further conditions are not asserted: its
index-zero gauge condition, the instance of \cite[(3.35)]{B-CMP99b}, and its $\mathbb{J}$-member
$\ell^{3}|\mathcal{J}(U)(y)|<T^{\mathfrak{B}}(y)$.
\end{enumerate}
\end{proposition}

\begin{proof}
\emph{Item (i).} This item is established by the argument that proves clause (C) of the arrest
theorem, Theorem~\ref{4.27:thm-arrest-theorem}, whose proof is incomplete at one step; clause (C)
states that the newborn terms may be stated on their arrested analyticity domains. The bound of
such a term
reads its arrested space through two properties only. The first is the holomorphy on the arrested
source of the old terms composed with the map of the later operation, which is stated by
clauses (A) and (C). The second is the action of the operators $H_\sigma(b)$ on
$\mathfrak{K}_P$, in which every arrested space lies by clause (D).
Clause (C) holds only within the scope it inherits from the clause fixing the scope of (A).
Item (i) therefore holds with that scope and no further.

\emph{Item (ii): the analyticity domain of the older term and the chain
\eqref{4.27:eq-live-piece-bounds-1}.}
We apply parts (a) and (b) of Lemma~\ref{4.27:lem-frozen-room} at the point $x$ of the live piece
of class $k_{\mathfrak{a}}$, with the level $\ell(x)$, the stage $n\le n^{+}$ and
$k'\ge k_{\mathfrak{a}}$. Together with the descent of the stage spaces of the arrested class,
\begin{equation*}
\mathfrak{S}^{(k_{\mathfrak{a}}),\mathrm{arrest}}_{n^{+}}\subset\dots\subset
\mathfrak{S}^{(k_{\mathfrak{a}}),\mathrm{arrest}}_{0},
\end{equation*}
which is among the hypotheses of Theorem~\ref{4.27:thm-arrest-theorem}, whose proof is
incomplete at one step, and which clause (D) of that theorem records on the live piece,
they give the inclusions
\eqref{4.27:eq-live-piece-bounds-1} and the stated room. Consequently every later source lies in
the analyticity domain that $\tau$ keeps.

\emph{Item (ii): inclusions (a) and (b).} Take the determining sets
$\mathfrak{B}_\tau\preceq\mathfrak{B}'$. We apply to them heredity (e3) of
Lemma~\ref{4.27:lem-weights-heredity}, whose proof is incomplete at one step; see
\eqref{4.27:eq-weights-heredity-b}. Its first assertion is
$\mathfrak{G}(\mathfrak{B}')\subset\mathfrak{G}^{-J_0}(\mathfrak{B}_\tau)$ as sets on all of
$\Omega_0(\mathfrak{B}_\tau)$, which is (a). Its assertion on the points of the stratum which
enter $\Omega_1(\mathfrak{B}')$ is (b); at points held by a live piece it uses the case of the
current bound at held raw points, \eqref{4.27:eq-held-raw-points-a}, for points that enter
$\Omega_1(\mathfrak{B}')$.

\emph{Item (ii): inclusion (c).} Let $y\in X_\tau$ be raw at both dates. We distinguish three
cases.
\begin{itemize}
\item Suppose $y$ is held at the date of $\mathfrak{B}_\tau$ by a live piece. The current bound
at held raw points, \eqref{4.27:eq-held-raw-points-a},
gives $T^{\mathfrak{B}'}(y)=T^{\mathfrak{B}_\tau}(y)$.
\item Suppose $y$ lies inside $X_{\mathfrak{a}}$ or inside a class component of step
$k_{\mathfrak{a}}$. The statement on current thresholds across an increase of density,
\eqref{4.27:eq-density-increase-a}, whose proof is incomplete at one step,
dates $y$ and gives the same equality.
\item Suppose $\mathfrak{B}'$ is of density $k_\tau=k_{\mathfrak{a}}$. The statement on unheld raw
points within one density, \eqref{4.27:eq-unheld-raw-points-a}, whose proof is incomplete
at one step, gives the same equality at every such $y$.
\end{itemize}
In each of these cases, a pair of $\mathfrak{G}(\mathfrak{B}')$ satisfies
\begin{equation*}
\ell^{3}|\mathbb{J}(y)|<T^{\mathfrak{B}'}(y)=T^{\mathfrak{B}_\tau}(y),
\end{equation*}
which is the $\mathbb{J}$-member of $\mathfrak{G}(\mathfrak{B}_\tau)$ at $y$. This is a
comparison of two memberships at one point, so no map of a later operation enters.

$(\ast)$ Consider the raw points of $X_\tau$ in components of $Z_{k_{\mathfrak{a}}}$ that are not
class components of step $k_{\mathfrak{a}}$, lying outside every live piece, and take
$\mathfrak{B}'$ of a density $>k_{\mathfrak{a}}$. None of the three dating statements above
applies there. Across such an increase of density the radius $\gamma_0^{(k)}$, a maximum over
$j\le k$, may grow, and with it the datum an unheld point would read.
The proof does not derive the $\mathbb{J}$-member at these points, and the proposition does not
assert it.

\emph{Items (iii) and (iv).} These are clauses (A) and (B) of
Theorem~\ref{4.27:thm-arrest-theorem}, whose proof is incomplete at one step. They are read as
those clauses state them, and (iv) carries the scope of clause (B): the index-zero gauge
condition and the $\mathbb{J}$-member of the real point are not asserted.
\end{proof}

\begin{remark}[The theorem concerns the field factor alone]\label{4.27:rem-field-factor}
Write $\mathcal{C}_{\mathrm{arr}}$ for the class of frozen terms $\tau\in\mathfrak{F}(\mathfrak{a})$ of live
arrests $\mathfrak{a}$ (Definition~\ref{4.27:def-arrests-pieces}); for such a $\tau$ let $\mathfrak{a}$
denote its arrest and $\mathfrak{v}(\mathfrak{a})$ the frozen variables of $\mathfrak{a}$.

Every $\tau\in\mathcal{C}_{\mathrm{arr}}$ is a stored term, in the sense of
Definition~\ref{4.27:def-arrests-pieces}, on
$\mathfrak{D}_\tau\times\mathfrak{D}_{\mathcal{S}}$. Here $\mathfrak{D}_\tau$ is the domain of
$\tau$ in the pair $(\mathbb{U},\mathbb{J})$, and $\mathfrak{D}_{\mathcal{S}}$ is the stored
domain of its rim arguments (the arguments that $\tau$ carries at the rim levels of $\mathfrak{a}$).
The arguments of $\tau$ other than $(\mathbb{U},\mathbb{J})$ are of
two kinds: its rim arguments, and, when $\mathfrak{a}$ is of branch B4, the variable
$\mathsf{B}=B_j$ of the interrupted change of variables, supplied by the statement that carries the
stored terms of branch B4 holomorphically to the auxiliary variables of that change of variables
(cf.\ Lemma~\ref{4.27:lem-interrupted-terms}). All of them belong to
$\mathfrak{v}(\mathfrak{a})$.

The arrest theorem (Theorem~\ref{4.27:thm-arrest-theorem}; stated, proof incomplete at the step
marked $(\ast)$) is a statement on the field factor alone, that is, on the factor of each frozen
term that depends on the pair $(\mathbb{U},\mathbb{J})$. Its clauses
(A)--(E) hold with these further arguments as parameters, each fixed in its stored domain: the rim
arguments in $\mathfrak{D}_{\mathcal{S}}$, and $\mathsf{B}$ in its stored domain of
Lemma~\ref{4.27:lem-interrupted-terms}, namely the polydisc $|\mathsf{B}|<\varrho^{B}_j$ for the
families $(\beta1)$, $(\beta2)$ of that lemma, and, for its family $(\beta3)$, real $\mathsf{B}$ with
$|\mathsf{B}|<6.3\,\delta'_j/g_j$; complex values of $\mathsf{B}$ enter that family through its
holomorphic extension, given by the statement named above, to the auxiliary variables of the
interrupted change of variables.

Later operations present the further arguments in rotated form \cite{B-CMP122b}. Where they then range
is stated in the range statement for the frozen arguments; that statement is not
used here. That statement uses Definition~\ref{4.27:def-arrests-pieces},
Lemma~\ref{4.27:lem-arrest-timing} and Lemma~\ref{4.27:lem-interrupted-terms}. Neither these three
statements nor the arrest theorem, read as above, uses the range statement. The dependence
therefore runs in one direction only, and there is no circle.
\end{remark}

\begin{proof}
The description of the frozen terms comes from Definition~\ref{4.27:def-arrests-pieces}. Each
$\tau\in\mathcal{C}_{\mathrm{arr}}$ appears there as a stored term with a domain in
$(\mathbb{U},\mathbb{J})$, together with the stored domain $\mathfrak{D}_{\mathcal{S}}$ of its rim
arguments. The same definition places these arguments among the frozen variables
$\mathfrak{v}(\mathfrak{a})$.

On branch B4 the additional argument $\mathsf{B}=B_j$ is supplied by the statement that carries the
stored terms of branch B4 to the auxiliary variables of the interrupted change of variables, which
places it in $\mathfrak{v}(\mathfrak{a})$. By Lemma~\ref{4.27:lem-interrupted-terms} the exponent
left on $X_{\mathfrak{a}}$ depends on $P=(\mathbb{U},\mathbb{J})$ and on $\mathsf{B}$; the families
$(\beta1)$, $(\beta2)$ are holomorphic in $(P,\mathsf{B})$ on the product of their stored domain in
$P$ with the polydisc $\{|\mathsf{B}|<\varrho^{B}_j\}$, and the family $(\beta3)$ is holomorphic in
$P$ at real $\mathsf{B}$ with $|\mathsf{B}|<6.3\,\delta'_j/g_j$ and in the auxiliary variables.
\end{proof}

\begin{remark}[The forms of the arrest theorem used in later steps. Stated; the proof below is incomplete at the step marked $(\ast)$]\label{4.27:rem-later-step-forms}
Throughout, $\mathfrak{a}$ is a live arrest of density $k_a$, with component $X_{\mathfrak{a}}$, held set
$W_{\mathfrak{a}}$, level $h_a$ and frozen level function $\ell=\ell(x)$. In this remark the letter $N$
denotes the number of stages of the arrest, as in the frozen-room lemma
(Lemma~\ref{4.27:lem-frozen-room}), and not the rank of the gauge group. The letter $W$ denotes a live
piece (in (K1), $W=W_{\mathfrak{a}}$), and $n^{+}_W$ is the stage index of that lemma for the piece $W$.
The index $n$, with $0\le n\le N$, is the stage index of that lemma: a term's clause at a point is the
class letter of stage $n$.
Further:
\begin{itemize}
\item $\mathfrak{B}$ denotes the determining set in force at the date under consideration, and $k$ its
density;
\item $\tau\in\mathfrak{F}(\mathfrak{a})$ is a stored term with domain $X_\tau$, born over the
determining set $\mathfrak{B}_\tau$ at density $k_\tau$, so that $k_\tau=k_a$;
\item $\mathfrak{O}$ is a later operation among $O1,\dots,O5,O4'$, with target set $\mathfrak{B}^t$
(the determining set in force immediately before its map), source set $\mathfrak{B}'$ (the one in force
after it) and source domain $Y'$.
\end{itemize}
The rider spaces, their inclusions
\begin{equation*}
\mathfrak{G}(\mathfrak{B})\subset\mathfrak{G}^{-J_0}(\mathfrak{B})\subset\mathfrak{G}_1(\mathfrak{B}),
\end{equation*}
and the dated threshold $T^{\mathfrak{B}}(y)$ are those of \eqref{4.27:eq-weights-heredity-b}. A point
$y$ is \emph{raw} at $\mathfrak{B}$ if $y\in\Omega_0(\mathfrak{B})\setminus\Omega_1(\mathfrak{B})$; this
set of raw points is the \emph{stratum} of $\mathfrak{B}$, and the index-$0$ members of a rider at its
points are its \emph{stratum members}. At
raw points the rider $\mathfrak{G}(\mathfrak{B})$ of pairs $(e^{i\eta A'}U,\mathbb{J})$ carries its
index-$0$ members, and these are of two kinds. The first are its $U$-members: the index-$0$ cube gauges,
which are the index-$0$ instance of \cite[(3.35)]{B-CMP99b}, and the index-$0$ condition on $A'$, which
is the index-$0$ instance of \cite[(3.37)]{B-CMP99b}. The second is its \emph{current member}
$\ell^{3}|\mathbb{J}(y)|<T^{\mathfrak{B}}(y)$. A raw point is \emph{held} at a date if it lies in the held
set $W_{\mathfrak{c}}$ of an arrest $\mathfrak{c}$ live at that date. With this notation, the arrest
theorem (Theorem~\ref{4.27:thm-arrest-theorem}, whose proof is incomplete at one step) enters the later
steps in the following seven forms. Every appeal in this remark and in its proof to
Theorem~\ref{4.27:thm-arrest-theorem}, to Proposition~\ref{4.27:prop-live-piece-bounds}, to
Lemma~\ref{4.27:lem-real-minimiser} or to Lemma~\ref{4.27:lem-weights-heredity} is an appeal to a
statement whose proof is incomplete at one step.
\begin{enumerate}
\item[(K1)] (The form in which the arrested reading enters the correlation theorem.)
From each arrest $\mathfrak{a}$ to its dissolution the arrested reading
\eqref{4.27:eq-arrested-spaces-a} is in force, so that every space is taken on $W_{\mathfrak{a}}$ in the
following form.
\begin{itemize}
\item The frozen structure $(\ell,w_{\mathfrak{a}},\mathcal{Q})$ and the frozen family
$(\Sigma^{*})^{\mathfrak{a}}$ are used.
\item The index-$0$ current datum at the raw points of $W_{\mathfrak{a}}$ is frozen at the radius
$\gamma_0$ of the innermost live piece.
\item The quantity $c$ of \eqref{4.27:eq-arrested-spaces-a} is identically equal to $1$.
\item The factor is $F^{(k')}$ on standard spaces and on sources of $\mathbf{R}$, and on attempt spaces
it is
\begin{equation*}
\varphi^{(n),(k)}-\beta\sum_{\mathfrak{b}}\sigma_{\mathfrak{b}}(\ell),
\end{equation*}
the sum running over the older live arrests $\mathfrak{b}$, as in the factor clause of
\eqref{4.27:eq-arrested-spaces-a}; each of these factors is
at least $\varphi_{\mathrm{arrest}}$.
\item The spaces depend neither on the complex source parameters $z$ nor on the slot of a term, and they
lie inside $\mathfrak{K}_{109}(\mathfrak{B};\gamma_0^{(k)})$. Memberships depend on $\gamma_0^{(k)}$,
while every constant, floor and ceiling depends on $\bar\gamma_0(\mathsf{t})$, all under one supremum
over $\mathsf{t}\ge\mathsf{t}_\gamma$ taken together with the other majorant functions of $\mathsf{t}$
at the same $\mathsf{t}$, as in clause~(d) of Lemma~\ref{4.27:lem-weights-heredity}.
\item Riders are taken to be $\mathfrak{G}(\mathfrak{B})$.
\end{itemize}
Later sources are proved to lie in $\mathfrak{G}^{-J_0}$, and later fields and images in
$\mathfrak{G}_1$, in each case together with the list of stratum members in the rider part of
clause~(A) of the arrest theorem and in (K7) below. This never includes the current member of
$\mathfrak{G}$ at raw points held by no live piece. It never includes a member that a map moves at
index $0$.

The real minimiser lies in the arrested spaces on $\Omega_1(\cdot)$ and at the cubes, with rider
$\mathfrak{G}_1$, by clause~(B) of the arrest theorem. At raw index-$0$ points of $W_{\mathfrak{a}}$ no
property beyond that list is asserted, whether of later sources, of images or of the real point; in
particular the real point's index-$0$ gauge member and its current member are not asserted.
Bounds of later steps reach live pieces through Proposition~\ref{4.27:prop-live-piece-bounds}, whose
proof is incomplete at one step.

The hypotheses are those under which Theorem~\ref{4.27:thm-arrest-theorem} is stated. Of these, the
following are carried, by name, into the hypotheses of the correlation theorem and of the other later
statements in which (K1) is used:
\begin{itemize}
\item the flow bounds on the running inverse couplings up to the step considered;
\item the size condition on the number of stages, $N>N_0$ (at density $k$, the condition
$N_k>N_0(k)$);
\item the condition on the class $\mathfrak{M}_{\delta}$ of decaying weights;
\item the condition on the change of variables;
\item the condition on the function $\mathbb{H}$ of the change of variables;
\item the condition on the operation $O3$;
\item one of the conditions, evaluated at $\bar\gamma_0(\mathsf{t})$, under which the first input
hypothesis of the arrest theorem holds on $\mathfrak{N}^{\vee}$;
\item the monotonicity condition in the radius $\gamma$;
\item the rider clause of \eqref{4.27:eq-arrested-spaces-a}, which carries the ambient continuation on
$W$.
\end{itemize}
A dilation $\mathsf{w}$ of a layer multiplies the spend of the table \eqref{4.27:eq-arrest-theorem-a};
the condition $16\mathsf{w}_L\bar P\le\beta\theta_S$, in which $\mathsf{w}_L$ denotes the dilation that
the rows of $O2$ use and $\bar P$ is the maximal increment, serves those rows.
\item[(K2)] The stored terms of $\mathfrak{a}$ are
\begin{equation}\label{4.27:eq-later-step-forms-1}
\mathfrak{F}(\mathfrak{a})=\{\mathbb{C}^{(i)}:\ i\le n\}\cup\mathfrak{F}^{B4}(\mathfrak{a}).
\end{equation}
Thus there is one further family for each live arrest of branch B4, and the difference terms
$\mathsf{W}(P,\mathsf{B})$ of Lemma~\ref{4.27:lem-interrupted-terms} are among its members. Its domain is
that of clauses (b) and (b$'$) of
that lemma; for the terms of item $(\beta3)$ of that lemma it comes with the ambient continuation on $W$.
Its size is that of clause (c), and it is entered as in that lemma. No term is counted twice.
\item[(K3)] (The form used in the uniqueness theorem.) Let $i<N$ and let $x$ be a point not raised after
stage $i$. Then the room of $\mathbb{C}^{(i)}_{k_a}$ at $x$, which is $\rho_i(x;k')$ of
\eqref{4.27:eq-frozen-room-a}, satisfies, at every $k'\ge k_a$,
\begin{equation}\label{4.27:eq-later-step-forms-2}
\rho_i(x;k')\ \ge\ \beta\max\bigl\{2^{-|\ell-h_a-i-1|},\,2^{-(k_a-\ell)}\bigr\}.
\end{equation}
At raised points, and for $i=N$, the room is at least $\tfrac12\beta2^{-(k_a-\ell)}$, and at least
$\beta2^{-(k_a-\ell)}$ for $k'>k_a$.
\item[(K4)] Each term of a chain is charged once, at a cell of its own band, as in clause (T5) of
Lemma~\ref{4.27:lem-arrest-timing}.
\item[(K5)] (The form used in the increment statement.) At $k'=k_a$ on $W_{\mathfrak{a}}$ the room is at
least $\tfrac32\beta2^{-(k_a-\ell)}$ if $\max(n,n^{+}_W)<N$, and at least $\tfrac12\beta2^{-(k_a-\ell)}$
otherwise.
\item[(K6)] On live pieces put $r_\tau:=\rho_n(\cdot\,;k)$, the room of
\eqref{4.27:eq-frozen-room-a}. The following hold.
\begin{itemize}
\item The row of $O4'$ in \eqref{4.27:eq-arrest-theorem-a} gives $r_\tau-\sum\iota\ge\tfrac12r_\tau$.
\item The third input hypothesis of Theorem~\ref{4.27:thm-arrest-theorem}, restricted to $W$, is clause (c)
of Lemma~\ref{4.27:lem-frozen-room}.
\item For the frozen terms of $\mathcal{C}_{\mathrm{arrest}}$ off $W$, the row of
\eqref{4.27:eq-arrest-theorem-a} for those terms applies.
\item The frozen family $(\Sigma^{*})^{\mathfrak{a}}$, the space $\mathfrak{N}^{\vee}$ and the radius
$\gamma_0^{(k)}$ are those of (K1).
\item The real-domain lemma (Lemma~\ref{4.27:lem-real-minimiser}, whose proof is incomplete at one step)
is used at margin $\mathsf{m}=1/16$, under that lemma's margin condition taken at $\mathsf{m}=1/16$, and
with the scope of clause~(B) of the arrest theorem. Its rider is $\mathfrak{G}_1(\mathfrak{B})$. At raw
index-$0$ points only the real point's $A'$-member is asserted; its index-$0$ gauge member and its
current member $\ell^{3}|\mathcal{J}(U)(y)|<T^{\mathfrak{B}}(y)$ are not.
\end{itemize}
\item[(K7)] (The form used wherever the domain of a term is stored, the correlation theorem included.)
\begin{enumerate}
\item[(a)] \emph{Stored domains.} The domain of a stored term is formed over its birth set
$\mathfrak{B}_\tau$ at its density $k_\tau$. Its rider is taken to be $\mathfrak{G}(\mathfrak{B}_\tau)$,
and $\mathfrak{G}(\mathfrak{B}_\tau)\subset\mathfrak{K}_{109}(\mathfrak{B}_\tau;\gamma_0^{(k_\tau)})$.
Hence no later $\gamma_0^{(k')}$ enters a stored size.
\item[(b)] \emph{Later sources.} Every later source over $\mathfrak{B}'$ is proved to satisfy
\begin{equation}\label{4.27:eq-later-step-forms-3}
\mathfrak{G}(\mathfrak{B}')\subset\mathfrak{G}^{-J_0}(\mathfrak{B}_\tau)\subset\mathfrak{G}_1(\mathfrak{B}_\tau).
\end{equation}
In addition it satisfies:
\begin{itemize}
\item every stratum member of $\mathfrak{G}(\mathfrak{B}_\tau)$ at points since lifted into
$\Omega_1(\mathfrak{B}')$;
\item the current member of $\mathfrak{G}(\mathfrak{B}_\tau)$ at those raw points of the stored term's
domain that are held by a live piece at $\mathfrak{B}_\tau$ or at some date of density $k_\tau$;
\item when $\mathfrak{B}'$ is of density $k_\tau$, the current member of $\mathfrak{G}(\mathfrak{B}_\tau)$
at every raw point of that domain.
\end{itemize}
\item[(c)] \emph{Later images.} What a later image is proved to have is a continuation in
$\mathfrak{G}_1(\mathfrak{B}_\tau)$ together with the stratum list of the rider part of clause~(A) of the
arrest theorem.
\item[(d)] \emph{Points at which nothing is asserted.} Nothing is asserted of later sources or images in
the following cases:
\begin{itemize}
\item at raw index-$0$ points of the domain held by no live piece at any date of density $k_\tau$, for
the current member, when the source is of a density $>k_\tau$;
\item for the members (those of $A'$ and of $\mathbb{J}$) that a later map moves on the domain at
index~$0$;
\item at the completion site of a chain, for every index-$0$ member at the raw points of the set
$\mathsf{Z}$ of that site and of the core of the completed chain.
\end{itemize}
\end{enumerate}
\end{enumerate}
\end{remark}

\begin{proof}
\emph{(K1).} The first list of (K1) is the arrested reading \eqref{4.27:eq-arrested-spaces-a} itself.
The frozen structure, the frozen family and the freezing of the index-$0$ current datum at the innermost
live piece's $\gamma_0$ are part of that reading. The last of these is the definition of $T^{\mathfrak{B}}(y)$
in \eqref{4.27:eq-weights-heredity-b}. The factors and their floor $\varphi_{\mathrm{arrest}}$ are the
factor clause of the same reading. The inclusion of every arrested space in $\mathfrak{N}^{\vee}$, and
the inclusion
\begin{equation*}
\mathfrak{N}^{\vee}\subset\mathfrak{K}_{109}(\mathfrak{B};\gamma_0^{(k)}),
\end{equation*}
are clause~(d) of Lemma~\ref{4.27:lem-weights-heredity}, whose proof is incomplete at one step. The rule
that memberships depend on $\gamma_0^{(k)}$ while constants, floors and ceilings depend on
$\bar\gamma_0(\mathsf{t})$, under one supremum over $\mathsf{t}\ge\mathsf{t}_\gamma$, is stated in the
same clause. That riders are taken to be $\mathfrak{G}(\mathfrak{B})$ is the rider clause of
\eqref{4.27:eq-arrested-spaces-a}, with $\mathfrak{G}(\mathfrak{B})$ as in clause~(e) of that lemma.

What later sources, fields and images meet of a rider comes from two statements. For sources it is
clause~(e3) of Lemma~\ref{4.27:lem-weights-heredity}. For images it is clause~(e4) of the same lemma,
together with the rider part of clause~(A) of the arrest theorem. Both are spelled out in the proof of
(K7) below, and the exclusions in (K1) are those of (K7)(d).

The real minimiser lies in the arrested spaces by clause~(B) of the arrest theorem. That clause asserts
Lemma~\ref{4.27:lem-real-minimiser} with its rider taken to be $\mathfrak{G}_1(\mathfrak{B})$, through
clause~(e4) of
Lemma~\ref{4.27:lem-weights-heredity}. At raw index-$0$ points clause~(B) asserts the real point's
$A'$-member, because $A'=0$ there. It does not assert the index-$0$ gauge member or the current member
$\ell^{3}|\mathcal{J}(U)(y)|<T^{\mathfrak{B}}(y)$, and so no property beyond the list is asserted at
these points.

Bounds of later steps reach live pieces as follows. Terms born after the arrest are covered by clause~(i)
of Proposition~\ref{4.27:prop-live-piece-bounds}, whose reach is that of the rider part of clause~(A) of
the arrest theorem. Terms born before it are covered by clause~(ii) of that proposition, which rests on
clause~(e3) of Lemma~\ref{4.27:lem-weights-heredity} and does not assert the current member at raw
points held by no live piece when the source is of a density larger than $k_a$. The hypotheses of (K1)
are those under which the arrest theorem is stated.

\emph{(K2).} The stored terms of a live arrest are the new boundary terms $\mathbb{C}^{(i)}$ of its stages
$i\le n$. To these is added, when the arrest is of branch B4, the family of
Lemma~\ref{4.27:lem-interrupted-terms}. The domain and the size of that family are clauses (b), (b$'$)
and (c) of that lemma, and the family is entered in $\mathfrak{F}(\mathfrak{a})$ by the clause of that
lemma that places its terms. For the terms of item $(\beta3)$ of that lemma the domain carries the
ambient continuation on $W$. That clause places each term in one family only, so no term is counted
twice.

\emph{(K3) and (K5).} Both are lower bounds on the room $\rho_n(x;k')$ of the frozen-room lemma,
\eqref{4.27:eq-frozen-room-a}. (K3) is clause~(a) of Lemma~\ref{4.27:lem-frozen-room} for the term
$\mathbb{C}^{(i)}_{k_a}$, that is at stage $n=i$, at points not raised after stage $i$; at raised
points, and for $i=N$, it is the raised-point clause of that lemma. (K5) is clause~(a) of the same lemma
at $k'=k_a$ on $W_{\mathfrak{a}}$.

\emph{(K4).} This is the band clause (T5) of Lemma~\ref{4.27:lem-arrest-timing}.

\emph{(K6).} The entries of (K6) are supplied as follows.
\begin{itemize}
\item The room $r_\tau$ is that of Lemma~\ref{4.27:lem-frozen-room} at $k'=k$.
\item The inequality $r_\tau-\sum\iota\ge\tfrac12r_\tau$ is the spend entry of the row of $O4'$ in
\eqref{4.27:eq-arrest-theorem-a}.
\item The third input hypothesis of Theorem~\ref{4.27:thm-arrest-theorem}, restricted to $W$, is clause (c)
of Lemma~\ref{4.27:lem-frozen-room}.
\item The frozen terms of $\mathcal{C}_{\mathrm{arrest}}$ off $W$ are covered by their row in
\eqref{4.27:eq-arrest-theorem-a}.
\item The frozen family, $\mathfrak{N}^{\vee}$ and $\gamma_0^{(k)}$ are as in (K1).
\item The real-domain lemma at $\mathsf{m}=1/16$ is Lemma~\ref{4.27:lem-real-minimiser} under its margin
condition at that margin.
\item The index-$0$ scope of that lemma is exactly that of clause~(B) of the arrest theorem, as recalled
in the proof of (K1).
\end{itemize}

\emph{(K7)(a).} By clause~(d) of Lemma~\ref{4.27:lem-weights-heredity}, the domain of a stored term is
formed over the determining set and the density of its birth. Its rider is then
$\mathfrak{G}(\mathfrak{B}_\tau)$ by the rider clause of \eqref{4.27:eq-arrested-spaces-a}. Next, by
clause~(e3) of Lemma~\ref{4.27:lem-weights-heredity},
\begin{equation*}
\mathfrak{G}(\mathfrak{B}_\tau)\subset\mathfrak{K}_{109}(\mathfrak{B}_\tau;\gamma_0^{(k_\tau)})
\end{equation*}
as sets on all of $\Omega_0(\mathfrak{B}_\tau)$. The reason is that
$T^{\mathfrak{B}_\tau}(y)\le\gamma_0^{(k_\tau)}$: at raw points held by a live piece this threshold is
$\gamma_0^{(k_{\mathfrak{c}})}$ with $k_{\mathfrak{c}}\le k_\tau$, and $\gamma_0^{(\cdot)}$ is
non-decreasing, being a maximum over $j\le\cdot$; at raw points held by no live piece it equals
$\gamma_0^{(k_\tau)}$. Hence a stored size involves $\gamma_0^{(k_\tau)}$ only.

\emph{(K7)(b).} The chain \eqref{4.27:eq-later-step-forms-3} is the first inclusion of clause~(e3) of
Lemma~\ref{4.27:lem-weights-heredity}, followed by
$\mathfrak{G}^{-J_0}\subset\mathfrak{G}_1$ from \eqref{4.27:eq-weights-heredity-b}. The first inclusion
of (e3) holds on all of $\Omega_0(\mathfrak{B}_\tau)$, and it is obtained in three parts:
\begin{itemize}
\item on $\Omega_1(\cdot)$, because the thresholds $w_{\mathfrak{B}}r_a$ never rise in time (clause (e2));
\item at the cubes of $\mathcal{Q}_{\mathfrak{B}_\tau}$, because the cubes match the shells by the
nested-cube inclusion \eqref{4.27:eq-nested-cubes-a};
\item for the stratum's index-$0$ $U$-members, because these are date-free (clause~(e)).
\end{itemize}

Consider next a point since lifted into $\Omega_1(\mathfrak{B}')$. By clause~(e3), every pair of
$\mathfrak{G}_1(\mathfrak{B}')$ obeys every stratum member of $\mathfrak{G}(\mathfrak{B}_\tau)$, the
current member included, at every point of
$(\Omega_0\setminus\Omega_1)(\mathfrak{B}_\tau)\cap\Omega_1(\mathfrak{B}')$. At held points this is the
entering case, part (2) of \eqref{4.27:eq-held-raw-points-a}.

Consider now a point raw at both dates and held at the date of $\mathfrak{B}_\tau$ by a live piece.
There the threshold is frozen at the innermost piece's
$\gamma_0$, so $T^{\mathfrak{B}'}(y)=T^{\mathfrak{B}_\tau}(y)$ by part (1) of
\eqref{4.27:eq-held-raw-points-a}. Therefore the current member transfers. Within one density the same
transfer holds at held and unheld raw points alike, by \eqref{4.27:eq-unheld-raw-points-a}, because the
datum of an unheld point cannot change within a density. This gives the current member at points held at
some date of density $k_\tau$, and at every raw point when $\mathfrak{B}'$ is of density $k_\tau$.

\emph{(K7)(c).} The member of the later operation's source over $Y'$ has a continuation
$\tilde p\in\mathfrak{G}(\mathfrak{B}')$ by the rider clause of \eqref{4.27:eq-arrested-spaces-a}; from
this, clause~(e4) of Lemma~\ref{4.27:lem-weights-heredity} gives its image a continuation
$\tilde q\in\mathfrak{G}_1(\mathfrak{B}^t)$. Since $\mathfrak{B}_\tau\preceq\mathfrak{B}^t$ by the
date-order lemma (the birth date of a stored term precedes the target date of every later operation in
the total order $\preceq$ of dates), clauses (e1)--(e3) and \eqref{4.27:eq-nested-cubes-a} give
$\mathfrak{G}_1(\mathfrak{B}^t)\subset\mathfrak{G}_1(\mathfrak{B}_\tau)$. The stratum members that
$\tilde q$ carries in addition are those listed in the rider part of clause~(A) of the arrest theorem.

\emph{(K7)(d) and the exclusions of (K1).} $(\ast)$ Consider first the raw points of the domain held by no
live piece at any date of density $k_\tau$, when $\mathfrak{B}'$ is of a density $k'>k_\tau$. There no
inclusion of the current member of $\mathfrak{G}(\mathfrak{B}')$ in that of $\mathfrak{G}(\mathfrak{B}_\tau)$
is derived: across a density increase the maximum defining $\gamma_0^{(\cdot)}$ may grow, and
clause~(e3) keeps the threshold $\gamma_0^{(k)}$ at such points as a membership at the density of the
date, not as a hereditary clause. Consider next the members that a later map moves at index $0$ on the
domain, namely the $A'$-member and the current member, and, at the completion site of a chain, every
index-$0$ member at the raw points of $\mathsf{Z}$ and of the core. For these, neither clause~(e4) nor
clause~(A) of the arrest theorem derives any member. At all these points the rider
$\mathfrak{G}(\mathfrak{B}_\tau)$ that (K1) and (K7)(a) attach to the domain of a stored term is the
form in which the rider is taken, by the rider clause of \eqref{4.27:eq-arrested-spaces-a}, and not a
property proved of later sources and images.
\end{proof}

\begin{remark}[The dated radius $\gamma_0^{(k)}$ and its two coupling ceilings]
\label{4.27:rem-dated-radius}
Fix a density $k$, that is, the effective density $\rho_k$ at level $k$. Write
$\mathsf{x}_j=g_j^{-2}$, $\mathsf{t}_j=\log\mathsf{x}_j$ and $\mathsf{t}_\gamma=\log\gamma^{-2}$.
For $j\le k$ let $N_0(j)$ be the integer $N_0$ of \cite{B-CMP122a} at density $j$, and let
$\alpha_{0,m},\alpha_{1,m}$ be the two radii at level $m$; write $\alpha_{i,m}$, $i\in\{0,1\}$. Put
\begin{equation}\label{4.27:eq-dated-radius-1}
\begin{gathered}
w_{*}(j):=\mathsf{L}_0^{2N_0(j)}\quad\bigl(w_{*}(j):=1\text{ if step $j$ has no site of }
\mathbf{R}\bigr),\qquad
\bar\alpha_j:=\max_{m\le j}\max_i\alpha_{i,m},\\
\gamma_0^{(k)}:=5\max_{j\le k}w_{*}(j)\,\bar\alpha_j,\qquad w_{*}:=w_{*}(k).
\end{gathered}
\end{equation}
This is the dated size parameter of \eqref{4.27:eq-weights-heredity-a}. Parts (b) and (c) of
Lemma~\ref{4.27:lem-weights-heredity}, whose proof is incomplete at one step, give under the
monotonicity floor on $\mathsf{t}_\gamma$ among the hypotheses \eqref{4.27:eq-arrest-hypotheses-a}
the majorant
\begin{equation}\label{4.27:eq-dated-radius-2}
\gamma_0^{(k)}\le\bar\gamma_0(\mathsf{t}_k):=5(1+\beta_0)^2\,\Psi(\mathsf{t}_k),\qquad
\Psi(\mathsf{t}):=L^2(L+1)\,C\,\mathsf{t}^{\,r+q}e^{-\mathsf{t}/2},\qquad C:=\max_iC_i,
\end{equation}
with the constants $C_i$ entering that lemma. The function $\bar\gamma_0$ is explicit, involves
only $\mathsf{t}$ and constants fixed before $\gamma$, does not depend on the large-field history
of the steps, is decreasing on $\mathsf{t}\ge\mathsf{t}_\gamma$, and decays faster than every power
of $\mathsf{t}$.

Radii are taken by the following rule. Every membership in a space
$\mathfrak{K}_{109}(\mathfrak{B};\cdot)$ of pairs of the type of \cite{B-CMP109} is taken at
$\gamma_0^{(k)}$; we write $\mathfrak{K}_P:=\mathfrak{K}_{109}(\mathfrak{B};\gamma_0^{(k)})$ for the
determining set $\mathfrak{B}$ of a date of density $k$. Every constant, floor and ceiling in which
the radius enters is taken at
$\bar\gamma_0(\mathsf{t})$, under one supremum over $\mathsf{t}\ge\mathsf{t}_\gamma$ in which the
numbers $R_j$ of part (b) of Lemma~\ref{4.27:lem-weights-heredity}, the number of stages
$\check{R}_k$, $N_0$, $\mathsf{l}$, the radii $\alpha_i$ and $\vartheta$ are evaluated at the same
$\mathsf{t}$. Let $K_\Delta$ be the constant, independent of $\gamma$, with which the radius enters
the size estimates of the step with arrests. In particular the ceiling
$K_\Delta\bar\gamma_0\le\frac12$, the sharp radius condition and the smallness condition
\eqref{4.27:eq-dated-radius-5} of the bump estimate below are imposed at $\bar\gamma_0$, among the
ceilings of the hypotheses of the arrest theorem
(Definition~\ref{4.27:def-arrest-hypotheses}). The conditions on the radius used in the estimates
below are the ceilings $\sup_{\mathsf{t}\ge\mathsf{t}_\gamma}120\,\bar\gamma_0(\mathsf{t})\le1$ and
\eqref{4.27:eq-dated-radius-7}, and the monotonicity floor
$\mathsf{t}_\gamma\ge2(r+q)+c_\beta$, with $c_\beta=2\ln(1+\beta_0)$.

For a point $x$ let $w_{\mathfrak{B}}(x)$ be the actual weight at $x$ in the determining set
$\mathfrak{B}$ of the date of density $k$ at which it is read, taken as in
Definition~\ref{4.27:def-arrested-spaces}; at a point $x$ of the held set $W_{\mathfrak{a}}$ of a
live arrest $\mathfrak{a}$ holding $x$ innermost it is the frozen weight $w_{\mathfrak{a}}(x)$.
At a point held by no live arrest we assume $w_{\mathfrak{B}}(x)\le w_{*}(k)$.
Let $m(x)\le k$ be the level whose radii are read at
$x$, and put
\begin{equation}\label{4.27:eq-dated-radius-3}
w^{\vee}(x):=\max\{w_{*}(k),\,w_{\mathfrak{B}}(x)\}.
\end{equation}
Let $Q_a$ be the quantity bounded by the direct clause $a$, and $\mathbf{r}_{a,m}$ its threshold
radius at level $m$; these thresholds are the ones written $r_a(\cdot)$ in the clauses of the rider
space of part (e) of Lemma~\ref{4.27:lem-weights-heredity}. Let $\mathfrak{N}^{\vee}$ be the ambient
space of pairs, whose direct clauses at the points $x$ of $Z^{\mathrm{on}}_k$ read
$Q_a(x)<5w^{\vee}(x)\,\mathbf{r}_{a,m(x)}$.

\emph{Hypotheses.} The size estimates of the step with arrests are assumed in the form in which
they are written with the weight $w^{\vee}$; explicitly:
\begin{itemize}
\item[(E1)] the radius $\alpha_5$ of those estimates satisfies
$\alpha_5=O\bigl(5w^{\vee}(\tilde\alpha_0+\tilde\alpha_1)\bigr)+K_\Delta\gamma_0^{(k)}$,
where, at each point $x$ at which the estimate is read, $\tilde\alpha_0,\tilde\alpha_1$ denote
radii with $\tilde\alpha_i\le\alpha_{i,m(x)}$ and $w^{\vee}=w^{\vee}(x)$;
\item[(E2)] those estimates hold provided $\Pi_1\le\vartheta$ and $\Pi_2\le\vartheta$, where
\begin{equation}\label{4.27:eq-dated-radius-4}
\begin{aligned}
\Pi_1&:=c_G'\,O(1)\bigl[5w^{\vee}(BCM\tilde\alpha_0+\tilde\alpha_1)
+K_\Delta\gamma_0^{(k)}\bigr],\\
\Pi_2&:=\beta_\star\check{R}_k(LM)^4\cdot4(5w^{\vee})^2
(2\tilde\alpha_0+\tilde\alpha_1)^2;
\end{aligned}
\end{equation}
here $O(1)$ stands for a fixed numerical constant of those estimates, independent of $\gamma$;
$c_G'$ and $\beta_\star$ are constants of those estimates, independent of $\gamma$; $BCM$ is the
constant factor written so in the cube clauses of the rider space of
Lemma~\ref{4.27:lem-weights-heredity}(e), its letter $C$ being a different constant from the $C$ of
\eqref{4.27:eq-dated-radius-2}; $\check{R}_k$ is the number of stages, with
$\check{R}_k<L\mathsf{t}_k^{\,r}$; and
$\vartheta=\vartheta(\mathsf{t}_k)$, where the smallness function
$\vartheta(\mathsf{t})$ of those estimates is a constant multiple of $(\mathsf{t}/2)^{p_1-q}$, with
$p_1$ a fixed exponent;
\item[(E3)] the bump estimate, applied on a set $Y$ with $Y^+:=Y\cup\bigcup_{x\in Y}B(x)$, where
$B(x)$ is the set attached to $x$ by the bump estimate, so that $Y^+\setminus Y$ is the collar of
$Y$, has a conclusion that is pointwise at $x\in Y$, with a constant $K_S^b$ independent of the
frame, of the level at $x$, of the stage and density indices, of the cube, of the stored term, of
the size of the domain and of the two remaining parameters on which the stage estimates depend;
its only hypothesis on the restriction of the pair to $B(x)$ is the smallness condition
\begin{equation}\label{4.27:eq-dated-radius-5}
C^b\,(a_\flat+\kappa_{\max})\le1,
\end{equation}
with $C^b$ the constant of that condition, $a_\flat$ the amplitude entering it and $\kappa_{\max}$
the quantity added to $a_\flat$ in it; and among the size estimates it is the only one that
uses clauses at points of $Y^+\setminus Y$ (and of $X^+\setminus X$, with $X^+$ formed from $X$ as
$Y^+$ from $Y$).
\end{itemize}

\emph{Conclusions.}
\begin{itemize}
\item[(i)] For every point $x$ and every $i$,
$5w^{\vee}(x)\,\alpha_{i,m(x)}\le\gamma_0^{(k)}$. The factor is $5$ however many
steps separate the density $k$ from the density of a live arrest holding $x$, and no
monotonicity in $\mathsf{t}$ is used. The inclusion $\mathfrak{N}^{\vee}\subset\mathfrak{K}_P$ of
part (d) of Lemma~\ref{4.27:lem-weights-heredity}, whose proof is incomplete at one step, holds by
definition of $\gamma_0^{(k)}$ through this bound.
\item[(ii)] Under (E1) and (E2), the quantities $\alpha_5$, $\Pi_1$, $\Pi_2$ obey
\begin{equation}\label{4.27:eq-dated-radius-6}
\begin{aligned}
\alpha_5&\le\bigl(2\,O(1)+K_\Delta\bigr)\gamma_0^{(k)},\qquad
\Pi_1\le c_G'\,O(1)\,(BCM+1+K_\Delta)\,\gamma_0^{(k)},\\
\Pi_2&\le\beta_\star\check{R}_k(LM)^4\cdot36\bigl(\gamma_0^{(k)}\bigr)^2 ,
\end{aligned}
\end{equation}
with the implied constant of (E1) in the first bound.
\item[(iii)] At every density $k$ with $\mathsf{t}_k\ge\mathsf{t}_\gamma$, the two requirements
$\Pi_1\le\vartheta$ and $\Pi_2\le\vartheta$ of (E2) hold as soon as the two conditions
\begin{equation}\label{4.27:eq-dated-radius-7}
\begin{aligned}
&\sup_{\mathsf{t}\ge\mathsf{t}_\gamma}
\frac{c_G'\,O(1)\,(BCM+1+K_\Delta)\,\bar\gamma_0(\mathsf{t})}{\vartheta(\mathsf{t})}\le1,\\
&\sup_{\mathsf{t}\ge\mathsf{t}_\gamma}
\frac{\beta_\star L\mathsf{t}^{\,r}(LM)^4\cdot36\,\bar\gamma_0(\mathsf{t})^2}{\vartheta(\mathsf{t})}
\le1
\end{aligned}
\end{equation}
hold. Both suprema in \eqref{4.27:eq-dated-radius-7} tend to zero as
$\gamma\to0$. Thus \eqref{4.27:eq-dated-radius-7} is a ceiling on $\gamma$ alone: it puts no lower
bound on the coupling $g$ and no upper bound on $K$.
\item[(iv)] Let $\mathsf{t}_k\ge\mathsf{t}_\gamma$. Let $\mathfrak{D}^{\flat}_n$ be a stage domain
on $Y$ at stage $n$ of density $k$, let $\mathfrak{L}_n$ be its class, and let
$\mathfrak{B}_n:=\mathbb{B}^{(n)}_k$ be the determining set of $\mathfrak{L}_n$. An input of the
estimates on $\mathfrak{D}^{\flat}_n$ is the restriction to $Y$ of a pair of
$\mathfrak{G}(\mathfrak{B}_n)$ on the whole domain that
\begin{itemize}
\item[(a)] obeys the fat direct clauses (the form of the direct clauses taken as input by the
stage estimates) at every point of $Y$, and
\item[(b)] on $Y^+\setminus Y$ is subject only to the clauses
$Q_a(y)<5\theta^{\mathfrak{B}_n}_a(y)$ of $\mathfrak{G}(\mathfrak{B}_n)$.
\end{itemize}
Where an input is described on a set $X$, it is the restriction of a
$\mathfrak{G}(\mathfrak{B}_n)$-pair on the whole domain, fat on $X$. The tube of the stage, one of
the spaces read as in Definition~\ref{4.27:def-arrested-spaces}, is the restriction of a global
field of $\mathfrak{G}$: it uses at most $1.1$ units of room in the direct clauses on
$\mathfrak{T}_k=\Omega_k\cup Z_k^{c}$, obeys the direct clauses at $Z^{\mathrm{on}}_k$, and is at
most $1/16$ on the set $\mathsf{Z}$ of a completion site and on the core of that site. For
inputs so described, every estimate proved for inputs in $\mathfrak{K}_P$ holds with no constant
increased, and the bump estimate applies at every $x\in Y$.
\end{itemize}
\end{remark}

\begin{proof}
\emph{(i).} Let $x$ and $i$ be given. We distinguish the two terms that can realise the maximum in
\eqref{4.27:eq-dated-radius-3}.

First case: $w_{\mathfrak{B}}(x)\le w_{*}(k)$, so that $w^{\vee}(x)=w_{*}(k)$. Since
$m(x)\le k$, the definition of $\bar\alpha_k$ in \eqref{4.27:eq-dated-radius-1} gives
$\alpha_{i,m(x)}\le\bar\alpha_k$. The index $j=k$ occurs in the maximum defining $\gamma_0^{(k)}$,
so
\[
5w^{\vee}(x)\alpha_{i,m(x)}\le5w_{*}(k)\bar\alpha_k\le\gamma_0^{(k)}.
\]

Second case: $w_{\mathfrak{B}}(x)>w_{*}(k)$, so that $w^{\vee}(x)=w_{\mathfrak{B}}(x)$. By the
assumption on the weight at points held by no live arrest, $x$ lies in the held set
$W_{\mathfrak{a}}$ of a live arrest $\mathfrak{a}$ of density $k_{\mathfrak{a}}\le k$, holding $x$
innermost, and $w_{\mathfrak{B}}(x)=w_{\mathfrak{a}}(x)$. By part (a) of
Lemma~\ref{4.27:lem-weights-heredity}, whose proof is incomplete at one step,
$w_{\mathfrak{a}}(x)\le w_{*}(k_{\mathfrak{a}})$. Moreover, by
Definition~\ref{4.27:def-arrested-spaces} the radii at $x$ are read at the frozen level
$\ell_{\mathfrak{a}}(x)$ of the arrest, and $\ell_{\mathfrak{a}}(x)\le k_{\mathfrak{a}}$; so
$m(x)\le k_{\mathfrak{a}}$, hence $\alpha_{i,m(x)}\le\bar\alpha_{k_{\mathfrak{a}}}$. The index
$j=k_{\mathfrak{a}}\le k$ occurs in the maximum defining $\gamma_0^{(k)}$, so
\[
5w_{\mathfrak{B}}(x)\alpha_{i,m(x)}\le5w_{*}(k_{\mathfrak{a}})\bar\alpha_{k_{\mathfrak{a}}}
\le\gamma_0^{(k)}.
\]

In both cases the factor is $5$, whatever the number $k-k_{\mathfrak{a}}$ of steps separating the
density of the arrest from $k$. The comparison is between terms
of one maximum, with no use of monotonicity of any function of $\mathsf{t}$.

\emph{(ii).} By (i) with $i=0$ and $i=1$, and since $\tilde\alpha_i\le\alpha_{i,m(x)}$, we have
$5w^{\vee}\tilde\alpha_0\le5w^{\vee}\alpha_{0,m(x)}\le\gamma_0^{(k)}$
and
$5w^{\vee}\tilde\alpha_1\le5w^{\vee}\alpha_{1,m(x)}\le\gamma_0^{(k)}$.
Hence $5w^{\vee}(\tilde\alpha_0+\tilde\alpha_1)\le2\gamma_0^{(k)}$. By (E1), the
first term of $\alpha_5$ is at most the implied constant of (E1), written $O(1)$, times
$2\gamma_0^{(k)}$; this gives the bound on $\alpha_5$ in \eqref{4.27:eq-dated-radius-6}.

Multiplying the bound for $\tilde\alpha_0$ just obtained by the constant $BCM$ and adding the bound
for $\tilde\alpha_1$ gives
$5w^{\vee}(BCM\tilde\alpha_0+\tilde\alpha_1)\le BCM\,\gamma_0^{(k)}+\gamma_0^{(k)}$.
Adding $K_\Delta\gamma_0^{(k)}$ and multiplying by $c_G'\,O(1)$ yields the bound on $\Pi_1$.

Finally, $5w^{\vee}(2\tilde\alpha_0+\tilde\alpha_1)\le3\gamma_0^{(k)}$. Squaring,
multiplying by $4$ and then by $\beta_\star\check{R}_k(LM)^4$ gives
\[
\begin{split}
\Pi_2&=\beta_\star\check{R}_k(LM)^4\cdot4(5w^{\vee})^2
(2\tilde\alpha_0+\tilde\alpha_1)^2\\
&\le\beta_\star\check{R}_k(LM)^4\cdot36\bigl(\gamma_0^{(k)}\bigr)^2 ,
\end{split}
\]
which is the last bound of (ii).

\emph{(iii).} Let $\mathsf{t}_k\ge\mathsf{t}_\gamma$. By \eqref{4.27:eq-dated-radius-2},
$\gamma_0^{(k)}\le\bar\gamma_0(\mathsf{t}_k)$, and by (E2), $\check{R}_k<L\mathsf{t}_k^{\,r}$.
Evaluating the suprema of \eqref{4.27:eq-dated-radius-7} at $\mathsf{t}=\mathsf{t}_k$ and combining
with (ii), we obtain
\[
\Pi_1\le c_G'\,O(1)(BCM+1+K_\Delta)\,\bar\gamma_0(\mathsf{t}_k)\le\vartheta(\mathsf{t}_k)
\]
and
\[
\Pi_2\le\beta_\star\check{R}_k(LM)^4\cdot36\bigl(\gamma_0^{(k)}\bigr)^2
\le\beta_\star L\mathsf{t}_k^{\,r}(LM)^4\cdot36\,\bar\gamma_0(\mathsf{t}_k)^2
\le\vartheta(\mathsf{t}_k).
\]
These are the two requirements of (E2), with $\vartheta=\vartheta(\mathsf{t}_k)$.

By \eqref{4.27:eq-dated-radius-2}, $\bar\gamma_0(\mathsf{t})$ is a constant multiple of
$\mathsf{t}^{\,r+q}e^{-\mathsf{t}/2}$, and $\vartheta(\mathsf{t})$ is a constant multiple of
$(\mathsf{t}/2)^{p_1-q}$. Hence the quotient in the first supremum is a constant multiple of
$\mathsf{t}^{\,r+2q-p_1}e^{-\mathsf{t}/2}$. The quotient in the second is a constant multiple of
$\mathsf{t}^{\,r}\cdot\mathsf{t}^{\,2r+2q}e^{-\mathsf{t}}\cdot\mathsf{t}^{\,q-p_1}
=\mathsf{t}^{\,3r+3q-p_1}e^{-\mathsf{t}}$.

Each is continuous for $\mathsf{t}\ge\mathsf{t}_\gamma$ and tends to zero as
$\mathsf{t}\to\infty$, so each is bounded on $[\mathsf{t}_\gamma,\infty)$. Its supremum over
$[\mathsf{t}_\gamma,\infty)$ tends to zero as $\mathsf{t}_\gamma\to\infty$, that is, as
$\gamma\to0$, the multiplying constants being fixed before $\gamma$. So
\eqref{4.27:eq-dated-radius-7} is a condition on $\gamma$ only. Once it holds, the two
requirements hold at every density $k$ with $\mathsf{t}_k\ge\mathsf{t}_\gamma$, with no condition
on $g$ from below and none on $K$ from above.

\emph{(iv), clause (a).} By the arrested reading of the spaces
(Definition~\ref{4.27:def-arrested-spaces}) and part (e) of Lemma~\ref{4.27:lem-weights-heredity},
whose proof is incomplete at one step,
$\mathfrak{G}(\mathfrak{B}_n)\subset\mathfrak{G}(\mathfrak{B})\subset\mathfrak{K}_P$ for every
determining set $\mathfrak{B}\preceq\mathfrak{B}_n$ of density $k$, where $\preceq$ is the order of
dates by time. In particular the input is the restriction of a pair of the ambient space
required by the stage estimates. Clause (a) therefore replaces the set of inputs (restrictions of
pairs of $\mathfrak{K}_P$) by a subset.

An estimate proved for every input of the larger set holds for every input of the subset.
Every constant of such an estimate is a supremum over the inputs, and a supremum over a subset is
not larger. So every estimate proved for inputs in $\mathfrak{K}_P$ holds, and no constant grows.

\emph{(iv), clause (b).} Clause (b) enlarges the set of inputs only on $Y^+\setminus Y$ (and on
$X^+\setminus X$ for inputs described on $X$). By (E3), the only estimate that uses clauses there
is the bump estimate. Its conclusion is pointwise at $x\in Y$, with a constant $K_S^b$ that does
not depend on the data listed in (E3). Its only hypothesis on the restriction of the pair to $B(x)$
is the smallness condition \eqref{4.27:eq-dated-radius-5}.

On $\mathfrak{G}$ the amplitude satisfies $a_\flat\le\gamma_0^{(k)}$ by part (e) of
Lemma~\ref{4.27:lem-weights-heredity}, whose proof is incomplete at one step. Moreover
$\gamma_0^{(k)}\le\bar\gamma_0(\mathsf{t}_k)$ by \eqref{4.27:eq-dated-radius-2}. The condition
\eqref{4.27:eq-dated-radius-5} is imposed at $\bar\gamma_0$, under one supremum over
$\mathsf{t}\ge\mathsf{t}_\gamma$, among the ceilings of
Definition~\ref{4.27:def-arrest-hypotheses}; since $\mathsf{t}_k\ge\mathsf{t}_\gamma$, it may be
evaluated at $\mathsf{t}=\mathsf{t}_k$, so
\[
C^b(a_\flat+\kappa_{\max})\le C^b\bigl(\bar\gamma_0(\mathsf{t}_k)+\kappa_{\max}\bigr)\le1 .
\]
Hence the bump estimate applies at every $x\in Y$ to the inputs of (a)--(b), with the same
constant $K_S^b$.
\end{proof}

\begin{remark}[The one-box geometric hypothesis at the last stage]\label{4.27:rem-one-box}
Stated; proof incomplete at the steps marked $(\ast)$.

Throughout, $X$ is a class component of the step of density $k$. Let $N_k$ be the number of stages
of that step, $h=k-N_k$ (in this section $h$ is this integer, not a history), $N_0(k)$ the least
integer $n\ge1$ with $L^{-n+1}MR_{k-n+1}=M$
(Ba{\l}aban's number of last stages \cite{B-CMP122a}), and
$k_0=k-N_0(k)$. We write
$\xi_m=ML^m\eta$, $\mathsf{S}_m=MR_mL^m\eta$ as in the lemma on cube families under a stage move
(Lemma~\ref{4.27:prf-stage-move-cubes}, whose proof is incomplete at one step). $C$ is the constant
in Ba{\l}aban's cube side $C\xi_m$, and $B$ is the constant in the thresholds of his cube clause,
which at level $\lambda$ are multiples of $BCM\alpha_{0,\lambda}$. Cells, lattice boxes, regions,
complement regions
$A^{\complement}$ and the relation ``two regions share a cell'' are those of the conventions of the
nested-cube lemma (Lemma~\ref{4.27:lem-nested-cubes}), with distances in the sup-norm. The stage
families of \cite[(1.65), (1.67), (1.69)]{B-CMP122a} at stage $j$ used below are named in words.
\begin{itemize}
\item The \emph{new-shell family of index $m\le j$}: for $m<j$, cubes $\square\subset\Omega''_m$ with
$\square\not\subset\Omega''_{m+1}$, of side $C\xi_m$ and weight $\mathsf{L}_0^{2(m-k_0)^+}$; for $m=j$,
cubes $\square\subset\Omega''_j$ with $\square\not\subset\Omega_{\min\{j,k_0\}+1}$, of side $C\xi_j$
and weight $\mathsf{L}_0^{2(j-k_0)^+}$.
\item The \emph{old-shell family of index $m$}, for $k_0<m\le j$: cubes $\square$ with
$\square\subset\Omega_m$ and $\square\not\subset\Omega_{m+1}$, of side $C\xi_j$ and weight
$\mathsf{L}_0^{2(j-m)}$.
\item The \emph{top family}: cubes $\square\subset\Omega_k$ of side $C\xi_k$ and weight $1$.
\end{itemize}
Besides what is proved in this chapter and what is quoted from Ba{\l}aban's papers, the proofs of
this section use the following ingredients.
\begin{enumerate}
\item[(i)] The \emph{one-box condition} at the move of the old-shell family of index $k_0+1$ from
stage $k-1$ to stage $k$, in the case $N_0(k)\ge 2$ and $C=2$. It reads as follows. For every cube
$\square$ of that family at $X$, the set
\begin{equation}\label{4.27:eq-one-box-1}
\mathcal{N}(\square):=\{z:\ \operatorname{dist}(z,\square)\le C\xi_k\}
\end{equation}
shares a cell with at most one connected component of $\Omega^{\complement}_{k_0+1}\cap X$. This
condition is read from the geometric declaration on the nesting of the large-field regions, which
is among the standing declarations of the arrest theorem (Theorem~\ref{4.27:thm-arrest-theorem},
whose proof is incomplete at one step).
It is the only place where heredity of the cube clause reads a geometric declaration.
\item[(ii)] The definition of Ba{\l}aban's domain class $\mathfrak{D}_m$
\cite[near (1.61)--(1.63)]{B-CMP122a}. It is used to place a consumer term through every point of
$\mathsf{Z}\cap Z^{\mathrm{on}}$.
\item[(iii)] The inequality $C_H^\sharp\ge 1/80$, tacitly read in the $\mathbf{T}$-step and supplied
by no stated lower bound; see part~(iii) of the proof.
\item[(iv)] Three further points: the range of the bound on the presented pair on $\mathsf{Z}$, with
the cube clause among its entries; the interface at the completion site; and the clauses at
index~$0$.
\item[(v)] The units at points of $\mathsf{Z}\cap\mathfrak{T}_k$.
\item[(vi)] The width of the thin rims.
\item[(vii)] The meaning of ``component''.
\item[(viii)] The collection of the facts above that are read, without quotation, from Ba{\l}aban's
papers, from the nesting declaration, or from statements of this chapter used by statement.
\end{enumerate}
\end{remark}

\begin{proof}
\emph{(i).} In the proof of heredity of the cube clause
(Lemma~\ref{4.27:prf-stage-move-cubes}, whose proof is incomplete at one step), a cube $\square$ of
side $s$ is moved to a cube $\square'\supset\square$ of side $s'$. The move is made by clause~(a)
(the one-box hypothesis) or by clause~(b) (the aligned-cube hypothesis) of the nested-cube lemma,
Lemma~\ref{4.27:lem-nested-cubes}. At every move that enlarges a cube either the hypothesis of
clause~(a) holds by separation or clause~(b) applies, with one exception described after the
following list.
Specifically:
\begin{itemize}
\item At the new-shell family of index $j$ (stage $j\to j+1$), the hypothesis of clause~(a) holds by
the derived separation of Lemma~\ref{4.27:prf-stage-move-cubes} (whose proof is incomplete at the
one step isolated below): distinct components of
$Z''_{j+1}\cap X$ are at least $11\,\mathsf{S}_{\min\{j,k_0\}+1}$ apart, and
\begin{equation*}
11\,\mathsf{S}_{\min\{j,k_0\}+1}\ge11\,\xi_{j+1}>4.25\,\xi_{j+1}\ge(2+L^{-1})C\xi_{j+1},
\end{equation*}
the last quantity being the diameter of the $C\xi_{j+1}$-neighbourhood of $\square$. This uses the
statement of \cite{B-CMP122a} that all components of the domains $Z''_j$ are rectangular
parallelepipeds, together with the collar structure.
\item At every other pair formed by an old-shell family of index $m$ and a stage move $j\to j+1$,
clause~(b) applies with $S=\mathsf{S}_m\ge C\xi_{j+1}$, by the scales
\eqref{4.27:eq-stage-move-cubes-a} of that lemma. Hence the box property of the components of
$\Omega^{\complement}_m$ stated in \cite{B-CMP122a} is not needed there.
\item At the $\mathbf{T}$-step move and at the completion move, clause~(b) applies, since
$R_{k+1}\ge C$ at the $\mathbf{T}$-step and $R_k\ge C$ at the completion.
\end{itemize}
The exception is the move $j+1=k$, $m=k_0+1$ with $N_0(k)\ge2$ and $C=2$. There
$s=C\xi_{k-1}$ and $s'=C\xi_k$, so the set \eqref{4.27:eq-one-box-1} has sup-norm diameter
\begin{equation*}
s+2s'=(2+L^{-1})C\xi_k .
\end{equation*}
Clause~(b) fails: by \eqref{4.27:eq-stage-move-cubes-a} of that lemma the block side
$\mathsf{S}_{k_0+1}$ equals $\xi_k$, while $s'=2\xi_k$. Clause~(a) needs exactly
the one-box condition. That condition is implied by either of two sufficient forms:
\begin{itemize}
\item $\Omega^{\complement}_{k_0+1}\cap X$ is one box;
\item the components of $\Omega^{\complement}_{k_0+1}\cap X$ are more than $4.25\,\xi_k$ apart, for
instance at least five blocks of side $\mathsf{S}_{k_0+1}$ apart.
\end{itemize}
$(\ast)$ The one-box condition is read from the nesting declaration among the standing
declarations of the arrest theorem. That declaration enters here only through this condition, and it
is not established here that it implies one of the two sufficient forms. Should it imply neither,
the one-box condition is, at this one locus, a strengthening of that geometric declaration. In no
case does it bound a parameter: it is neither a lower bound, nor an upper bound, nor a condition on
the coupling, nor a restriction of the window of $\mathsf{L}_0$.

If Ba{\l}aban's constant is $C=1$, clause~(b) at $A=\Omega_{k_0+1}$ replaces clause~(a) at this
locus, and no declaration is read. Indeed $A$ is then a union of aligned blocks of side
$\xi_k=s'\ge s$. The value of $C$ is fixed by Ba{\l}aban's cube geometry, with $C\le 2$
\cite[p.~191]{B-CMP122a}; it is not at our disposal.

Without the declaration, take two box components of $\Omega^{\complement}_{k_0+1}\cap X$ across a
corridor of width $\xi_k$. Then $\square$ lies inside a cube of the new-shell family of index $k$
only. For this move of one level the weight ratio is $w_{\square'}/w_\square=\mathsf{L}_0^4$, and the
restriction ratio for the clause of first order in $L^\lambda\eta$ (the bound on $|A|$) is
$\mathsf{L}_0^4(1+\beta_0)^2L^{-1}$. With $\mathsf{L}_0^2=L/3$ and $\beta_0\le 1/7$ this equals
\begin{equation*}
\mathsf{L}_0^4(1+\beta_0)^2L^{-1}=\frac{L}{9}(1+\beta_0)^2,
\end{equation*}
which reaches $(L/9)(64/49)>1$ for every $L\ge 8$. The ratio for the clause of second order in
$L^\lambda\eta$ (the bound on $|\nabla^\eta A|$) is
$\mathsf{L}_0^4(1+\beta_0)^2L^{-2}\le 64/441$, and this one is harmless. The repair that avoids the
declaration would be the lower bound $L\ge(1+\beta_0)^2\mathsf{L}_0^4$. Equivalently, it is the upper
bound $\mathsf{L}_0^4\le L/(1+\beta_0)^2$, which narrows the admissible window of $\mathsf{L}_0$. A
narrowed window of $\mathsf{L}_0$ is a cap, and it is not imposed. At this locus, therefore, clause~(a)
of the nested-cube lemma rests on the wording of the nesting declaration.

The following observation is not used. \cite[p.~191]{B-CMP122a} states a descending sequence of
spaces under the conditions $\beta\le1/4$ and $\mathsf{L}_0^2\le L/3$, in the notation of this book.
Restriction of gauges alone gives the cube-clause part of that descent only where the corridor
configuration above does not occur, as is the case when ``rectangular parallelepipeds'' in
\cite{B-CMP122a} is read as one box per level inside $X$, which is the reading of the nesting
declaration. Alternatively the cube clause could be regenerated from the plaquette entry by the
gauge lemma invoked at \cite[p.~191]{B-CMP122a}, with $B=B_3$ taken from Theorem~1 of the paper
listed as reference~16 of \cite{B-CMP122a}; on that route heredity of the cube clause would need no
geometry at all. This route is not examined here.

\emph{(ii).} At $y\in\mathsf{Z}\cap Z^{\mathrm{on}}$ the consumer is
$\mathfrak{E}:=\mathbb{C}^{(N_k)}_k(X_{\mathfrak{E}})$ with $X_{\mathfrak{E}}\ni y$. It is a term of
$\mathbb{B}''_k=\mathbb{B}^{(N_k)}_k$, and \cite[(1.63)]{B-CMP122a} indexes these terms over every
domain $X_{\mathfrak{E}}\in\mathfrak{D}_m$ with $X_{\mathfrak{E}}\cap(Z_k\setminus Z''_k)\ne\emptyset$.
The band
$Z\cap\Omega^{\sim5}_{k_0+1}$ meets every completed component in its collar $X\cap\Omega_k$. The term
is kept across the reset of the determining set at the completion \cite[(1.33)]{B-CMP122b}, and it is
presented with the pair by the first part of \cite[(1.65)]{B-CMP122b}. It is therefore a consumer
that reads $\mathbb{B}''_k=\mathfrak{B}^t$ at $y$, with a factor that is at most $1$ and positive.

$(\ast)$ Ba{\l}aban defines $\mathfrak{D}_m$ as the class of connected unions of cubes of the
scale-$m$ cover \cite[near (1.61)--(1.63)]{B-CMP122a}; the step uses the following consequence of
that definition, which is not derived here: $\mathfrak{D}_m$ contains a domain through any given $y$
and any given cell of the
band, and one containing any given cube of $\mathcal{Q}(y)$. This is the one ingredient of
Ba{\l}aban's papers used at this step without quotation. It is not part of the content of the arrest
theorem.

The bounds on $\mathsf{Z}\cap\mathfrak{T}_k$ and on the core region $D_2\setminus B_S$ of the
completion site are absolute in the
units of $\mathfrak{B}^t$, because $w_{\mathfrak{B}^t}\ge1$. On $\mathsf{Z}\cap\mathfrak{T}_k$ the
bound is $\le 1.1\,r_a$; on the core the presented pair lies within $1/64+3/64=1/16$ of every
threshold.

\emph{(iii).} Consider the lower bound on the radius constant $C^{\mathrm{rad}}_0$ in the
$\mathbf{T}$-step lemma
(Lemma~\ref{4.27:lem-t-step-layer}, whose proof is incomplete at one step); the bound on
$C^{\mathrm{rad}}_1$ is treated alike. It is stated at $C_H$ and read at
$C_H^\sharp:=\max\{C_H,2K_\chi\}$. With $w_k$ the bound on the $\mathbf{T}$-datum and
$\varepsilon_\sharp$ the smallness threshold of the theorem on $\mathbb{H}$, the inequality
$w_k\le\varepsilon_\sharp$ is obtained by dividing
\begin{equation*}
C_H^\sharp w_k\le\tfrac{(1-\theta_S)\beta}{8}\,\alpha_{0,k}\le 0.0125\,\alpha_{0,k}
\end{equation*}
by $C_H^\sharp$ and then using $\alpha_{0,k}\le\varepsilon_\sharp$. $(\ast)$ The division yields
$w_k\le\alpha_{0,k}$ only if $C_H^\sharp\ge1/80$, which no statement used here
supplies. At $d=4$ one has $C_H^\sharp\ge2K_\chi\ge120B_\sharp$, but $B_\sharp>0$ is bounded below by
name only, among the lower bounds of the standing hypotheses of the arrest theorem. So
$120B_\sharp\ge1/80$ is not available.

Two one-sided remedies exist, and neither is imposed here.
\begin{itemize}
\item Read the same lower bound at $\max\{C_H^\sharp,1\}$ in place of $C_H$. This is an upward
change, fixed with $K_\chi(d,L,\kappa_1)$ after $\kappa_1$, and before the pair $(q,C)$, which
precedes $\gamma$ in the order of choice. Every bound at $C_H^\sharp$ persists, and no cap arises.
\item Impose $\sup|b|\le w_k\le\varepsilon_\sharp$ as an upper bound on $\gamma$, where $b$ is the
$\mathbf{T}$-step's datum on its level-$k$ bonds. By
$w_k\le 2\|C_{\mathrm{op}}\|C_1^{\mathrm{pd}}A_1e^{-\mathsf{t}_k/2}\mathsf{t}_k^{q}$, with
$\|C_{\mathrm{op}}\|$ a fixed operator norm and $C_1^{\mathrm{pd}}$ a constant, both entering that
bound on $w_k$, this
is the condition
\begin{equation*}
\sup_{\mathsf{t}\ge\mathsf{t}_\gamma}
\frac{2\|C_{\mathrm{op}}\|C_1^{\mathrm{pd}}A_1e^{-\mathsf{t}/2}\mathsf{t}^{q}}{\varepsilon_\sharp}\le1,
\end{equation*}
which has degree $-\infty$ in $\mathsf{t}_\gamma$ (the quantity under the supremum decays
exponentially) and adds one member to the list of upper bounds on $\gamma$.
\end{itemize}
No cap arises either way. Until one of the two is imposed, the hypothesis
$\|\mathsf{d}\|_1\le\varepsilon_\sharp$ of the theorem on $\mathbb{H}$ in the $\mathbf{T}$-step
is used as stated and is not derived here from the lower bound on $C_0^{\mathrm{rad}}$.

\emph{(iv). Entry $8$.} The cube clause on $\mathsf{Z}\cap Z^{\mathrm{on}}$ rests on the bound on the
presented pair on $\mathsf{Z}$, taken by statement. Its entries are $a=1,\dots,7$, the direct
quantities, and $a=8$, the cube clause. This is the same range that clause~(A) of the arrest
theorem uses (each entry $\le\theta_T/16$, where $\theta_T$ is the threshold of the consumer term,
the cube clauses of $\mathcal{Q}(x)$ included) and that
the real-base lemma uses at entries $2,3,4,8$.

At $a=8$ it means the following. For the consumer $\mathfrak{E}$ and $y\in\mathsf{Z}$, on each cube
$\square$ of $\mathcal{Q}_{\mathbb{B}''_k}(y)$ there is a gauge $u$ on $\square\cap X_{\mathfrak{E}}$;
here $\square$ has level $\lambda$ and family weight $w_\square$, and $A$ is defined by
$U^u=e^{i\eta A}$, with $U$ the $G$-part of the one presented $\mathbb{U}$. The gauge satisfies
\begin{equation*}
L^\lambda\eta|A|,\ (L^\lambda\eta)^2|\nabla^\eta A|<\tfrac1{16}\,w_\square BCM\alpha_{0,\lambda}.
\end{equation*}
This is the gauge that $\mathfrak{G}(\mathfrak{B}^t)$ asks for, since $1/16<5$. ``Within a fraction
$\mathsf{m}$ of a cube threshold'' is the notion used by the real-base lemma, here with
$\mathsf{m}=1/16$.

$(\ast)$ Suppose that bound enumerated the seven direct quantities only. Then the cube clause on
$\mathsf{Z}\setminus(D_2\setminus B_S)$ would need one further argument, and there are two
candidates.
\begin{itemize}
\item The real-base route: the real-base lemma at $(\mathbb{B}''_k,V)$, $V$ a real background, with
$\mathsf{m}=1/16$, with
the gauges of the real-point lemma on determining sets holding excluded attempts. This works if the
presentation on $\mathsf{Z}$ is a layer on the real point $U(\mathbb{B}''_k,V)$ with $V$ real in the
support, as
the vanishing $A'=0$ in the core bounds suggests.
\item The gauge lemma invoked at \cite[p.~191]{B-CMP122a} (Theorem~1 of the paper listed as
reference~16 of \cite{B-CMP122a}, with $B=B_3$), applied from the plaquette entry~$5$.
\end{itemize}
Each candidate reads one fact not quoted here. Nothing among the statements used restricts
``entry'' to $a\le7$.

\emph{(iv). Interface.} The presented pair at the completion site is one configuration
$(e^{i\eta A'}U,\mathbb{J})$ on $\Omega_0(\mathfrak{B}^t)$
\cite[(1.50), (1.53), (1.65)]{B-CMP122b}. It is bounded region by region in that one
decomposition:
\begin{itemize}
\item on $\mathsf{Z}\cap Z^{\mathrm{on}}$, by the bound on the presented pair at
$\mathfrak{E}=\mathbb{C}^{(N_k)}_k(X_{\mathfrak{E}})$ with $X_{\mathfrak{E}}\ni y$;
\item on $\mathsf{Z}\cap\mathfrak{T}_k$, by the bound $\le1.1\,r_a$ of (v), and on the core
$D_2\setminus B_S$, by the core bounds;
\item at the points served by the layer bound off the zone (the bound on the layers of the
$\mathbf{R}$-operation at the points it serves: foreign shells, $\Omega_k\setminus Z$, collars, and
the rings off $\mathsf{Z}$), by its kept $U$-part and cube gauges,
followed by the nested-cube lemma and the heredity of the cube clause
(Lemma~\ref{4.27:prf-stage-move-cubes}, whose proof is incomplete at one step);
\item on $\Omega_0\setminus Y'$, by the input.
\end{itemize}
Every one of these bounds states the cube clause per cube of
$\mathcal{Q}_{\mathbb{B}''_k}=\mathcal{Q}_{\mathfrak{B}^t}$, at its family's weight and free of
factors \cite[(1.65), (1.67), (1.69)]{B-CMP122a}, on the whole cube and under its own gauge $u$ on
that cube. A cube astride two regions therefore reads the one $U$ under one gauge, chosen as
follows:
\begin{itemize}
\item the consumer's gauge, with $X_{\mathfrak{E}}\supseteq\square$, if the cube has a point in
$\mathsf{Z}$ (the
bound quantifies the point $y$);
\item else the core's gauge, if the cube lies in $D_2\setminus B_S$, as the cubes of the frozen
families of dissolved pieces do;
\item else the kept gauge of the layer bound off the zone.
\end{itemize}
No cube is assembled from two bounds. $(\ast)$ Two facts are taken as given by the layer bound off
the zone and by the bound on the presented pair, and are not proved here: that $Y'\setminus\mathsf{Z}$
is exhausted by the list of the layer bound, with the rings of the chain off $\mathsf{Z}$; and that
its increments there are in the units of the structure in force. The units are level $k$ on
$\Omega_k\setminus Z$ and on the collars, and the raised level on the rings.

\emph{(iv). Index $0$.} Let $\mathfrak{B}$ be a determining set of density $k$, not necessarily that
of the step of $X$. On the stratum $\Omega_0(\mathfrak{B})\setminus\Omega_1(\mathfrak{B})$ of
raw points, $\mathfrak{G}(\mathfrak{B})$ carries the index-$0$ clause of
$\mathfrak{K}_{109}(\mathfrak{B};\cdot)$. It has two parts.
\begin{itemize}
\item The $U$-members: \cite[(3.35), (3.37)]{B-CMP99b} at index $0$, with $\gamma$-free radii, on the
date-free index-$0$ cubes.
\item The $J$-member $\ell^3|\mathbb{J}(y)|<T^{\mathfrak{B}}(y)$. Here
$T^{\mathfrak{B}}(y)=\gamma_0^{(k_{\mathfrak{c}})}$, with $k_{\mathfrak{c}}$ the density of the
innermost, equivalently the oldest, live piece $\mathfrak{c}$ holding $y$, if one exists, and
$T^{\mathfrak{B}}(y)=\gamma_0^{(k)}$ otherwise. This extends the freezing of the arrest to the
index-$0$ $J$-datum.
\end{itemize}
At raw points held by no live piece, the number $\gamma_0^{(k)}$ is kept as a membership at the
density of $\mathfrak{B}$ and is not asserted to be hereditary.

Let $\mathfrak{G}_1(\mathfrak{B})$ be the space of pairs satisfying those clauses of
$\mathfrak{G}(\mathfrak{B})$ that sit at points of $\Omega_1(\mathfrak{B})$ or on cubes of
$\mathcal{Q}_{\mathfrak{B}}$, with no stratum member, and let $\mathfrak{G}^{-J_0}(\mathfrak{B})$ be
the space of pairs satisfying these clauses and the stratum's $U$-members. Both contain
$\mathfrak{G}(\mathfrak{B})$, and neither involves $\gamma_0$. For determining sets
$\mathfrak{B}\preceq\mathfrak{B}'$ of an earlier and a later date, the first inclusion (heredity) is
$\mathfrak{G}(\mathfrak{B}')\subset\mathfrak{G}(\mathfrak{B})$, and the second inclusion is
$\mathfrak{G}(\mathfrak{B})\subset\mathfrak{K}_{109}(\mathfrak{B};\gamma_0^{(k)})$, with $k$ the
density of $\mathfrak{B}$. The statements of this section assert the following.
\begin{itemize}
\item The inclusions
\begin{equation}\label{4.27:eq-one-box-2}
\mathfrak{G}(\mathfrak{B}')\subset\mathfrak{G}^{-J_0}(\mathfrak{B}),\qquad
\mathfrak{G}^{-J_0}(\mathfrak{B}')\subset\mathfrak{G}^{-J_0}(\mathfrak{B}),\qquad
\mathfrak{G}_1(\mathfrak{B}')\subset\mathfrak{G}_1(\mathfrak{B})
\end{equation}
of Lemma~\ref{4.27:prf-raw-stratum}, with every stratum member at points entering
$\Omega_1(\mathfrak{B}')$ and the $J$-member at held points raw at both dates.
\item The membership statement for continuations, and the part of clause~(A) of the arrest theorem on
continuations: membership in $\mathfrak{G}_1$, together with an explicit list of stratum members.
\item Clause~(B) of the arrest theorem: the real point lies in $\mathfrak{G}_1(\mathfrak{B})$,
together with $(3.37)_0$, which follows from $A'=0$.
\item The second part of the proposition on later sources, and the corresponding item of the summary
statement of this section: later sources lie in $\mathfrak{G}^{-J_0}(\mathfrak{B}_\tau)$, together
with the stratum members at entered points and the dated $J$-member; a further item of the summary
statement carries the same restriction.
\item Clause~(C) of the arrest theorem and the first part of the proposition on later sources inherit
the scope of the part of clause~(A) on continuations.
\end{itemize}
None asserts membership in
$\mathfrak{G}$. The full $\mathfrak{G}(\mathfrak{B})$ remains the space in which every space of the
arrest, every source and every stored term is read, and it is the left member of the second
inclusion.

\emph{What is proved.} The following are proved in Lemma~\ref{4.27:prf-raw-stratum} (heredity on
the raw stratum) and Lemma~\ref{4.27:prf-held-raw-points} (the current bound at held raw points),
together with, for the clauses at points of $\Omega_1(\cdot)$ and on cubes, the heredity of the cube
clause (Lemma~\ref{4.27:prf-stage-move-cubes}, whose proof is incomplete at the one step of (i))
and the nested-cube lemma (Lemma~\ref{4.27:lem-nested-cubes}).
\begin{itemize}
\item The second inclusion holds as a set inclusion on all of $\Omega_0(\mathfrak{B})$, because
$T^{\mathfrak{B}}\le\gamma_0^{(k)}$.
\item The inclusions \eqref{4.27:eq-one-box-2} hold as sets.
\item Every stratum member holds at the points lifted into $\Omega_1(\mathfrak{B}')$ by
\cite[(1.33)]{B-CMP122b}.
\item The $J$-member of the first inclusion holds at every raw point held at $\mathfrak{B}$ by a live
piece. The reason is that the innermost live piece holding the point cannot change before its nest
dissolves, and at that moment \cite[(1.33)]{B-CMP122b} lifts the point into $\Omega_1$ for good; at
points entering $\Omega_1(\mathfrak{B}')$ the bound at entering points is used at the density of the
innermost piece.
This uses the reading that \cite[(1.33)]{B-CMP122b} applies to every completed component, second
class included (the reading questioned in (viii)).
\item The $J$-member also holds at every raw point, held or not, when $\mathfrak{B}'$ has the density
of $\mathfrak{B}$, and from the point's first holding within that density on. The reason is that the
holders attached between two dates of one density $k$ have density $k$, and at most one of them ever
holds a given raw point.
\end{itemize}
Call a \emph{dating} of the stratum $J$-member at a point $y$ an assignment $d\mapsto T(y,d)>0$ of
the threshold in $\ell^3|\mathbb{J}(y)|<T(y,d)$ to the dates $d$ at which $y$ is raw; the dating used
above is $T(y,d)=T^{\mathfrak{B}_d}(y)$, with $\mathfrak{B}_d$ the determining set of the date $d$.
Consequently, for the stored terms $\tau\in\mathfrak{F}(\mathfrak{a})$ of the arrest theorem, the
dating of the stratum $J$-member agrees with that of every later source and target in the following
cases:
\begin{itemize}
\item at the raw points of $X_\tau$ inside $X_{\mathfrak{a}}$;
\item at the raw points of $X_\tau$ inside a class sibling of step $k_{\mathfrak{a}}$;
\item at the raw points of $X_\tau$ inside any live piece at $\mathfrak{B}_\tau$;
\item at the raw points of the components of $Z_{k_{\mathfrak{a}}}$ that are not class components,
against every source and target of density $k_{\mathfrak{a}}$.
\end{itemize}
This includes the dating of the ambient membership on $X_{\mathfrak{a}}$, since the dating there is at
most $\gamma_0^{(k_{\mathfrak{a}})}$. Let the continuation be the pair equal to the image on $Y'$, the
source domain of the later operation; to the input on $X_\tau\setminus Y'$; to the input's
$\mathbb{U}$ and the $\Omega_1$-part of its $\mathbb{J}$ elsewhere; and with $\mathbb{J}$ cut to $0$
at raw points off $X_\tau\cup Y'$.
Hence the continuation lies in $\mathfrak{G}_1(\mathfrak{B}_\tau)$. It obeys, in addition:
\begin{itemize}
\item every stratum member at the points that have entered $\Omega_1$;
\item its stratum $U$-members off $Y'$;
\item its index-$0$ gauges on $Y'$, where the map is a layer into $U'$, that is, where it keeps the
$G$-part $U$ and the cube gauges fixed;
\item its $J$-member off $X_\tau\cup Y'$ and at the dated points of $X_\tau\setminus Y'$.
\end{itemize}
The same holds for the inclusions of later sources, where no map intervenes and the $U$-members pass
by identity.

\emph{What is not proved.} Two points remain open at index $0$. The first concerns the heredity of
the $J$-datum at unheld raw points; the second concerns room at index $0$.

$(\ast)$ First, at raw points held at $\mathfrak{B}$ by no live piece, the $J$-part of
$\mathfrak{G}(\mathfrak{B}')\subset\mathfrak{G}(\mathfrak{B})$ is not proved. There
$T^{\mathfrak{B}}=\gamma_0^{(k)}$, while $T^{\mathfrak{B}'}$, with $k'$ the density of
$\mathfrak{B}'$, equals $\gamma_0^{(k')}$ or
$\gamma_0^{(k_b)}$ of a first later arrest, which is in general larger. The pair $\mathbb{U}\equiv1$
with $\mathbb{J}$ supported at $y$ and
\begin{equation*}
\gamma_0^{(k)}\le\ell^3|\mathbb{J}(y)|<T^{\mathfrak{B}'}(y)
\end{equation*}
lies in $\mathfrak{G}(\mathfrak{B}')\setminus\mathfrak{G}(\mathfrak{B})$. At such points no freeze
acts: there is no piece whose $\gamma_0$ could be frozen, and no clause of the form weight times
threshold for the comparison of thresholds to carry. The failure is confined to
the case in which $\mathfrak{B}'$ has a density $k'>k$ and the point is held by no live piece at any
date of density $k$.

For the arrest theorem this matters only at the raw points of $X_\tau\setminus Y'$ that satisfy all
of the following:
\begin{itemize}
\item they lie in components of $Z_{k_{\mathfrak{a}}}$ that are not class components of step
$k_{\mathfrak{a}}$, among them
the components holding a second-class component of an earlier step that is included again into the
large-field domain \cite{B-CMP122b} and counted as a creation for the timing condition that defines
class components; a class component of step $k_{\mathfrak{a}}$ that completes at the completion
date of that step, as first or as second class, is not among them;
\item they lie outside every live piece;
\item they are read by a source of density $>k_{\mathfrak{a}}$: the $\mathbf{T}$-step from
$k_{\mathfrak{a}}$ to $k_{\mathfrak{a}}+1$ on, a completion at a density $k>k_{\mathfrak{a}}$, the
stages of a later attempt, or a completion at which live pieces dissolve. Against sources of density
$k_{\mathfrak{a}}$ the dating is proved.
\end{itemize}
Off $X_\tau\cup Y'$, cutting $\mathbb{J}$ to $0$ suffices.

This gap cannot be closed by re-assigning the datum at the unheld raw points alone. Keep the frozen
datum at held points and keep the second inclusion. The probe lemma and the characterisation of
hereditary datings show that the datings for which both inclusions hold at a raw point $y$, for every
pair of dates at which $y$ is raw, are exactly the datings that are non-increasing along the interval
of dates at which $y$ is raw and bounded, at each date $d$, by $\gamma_0$ at the density of $d$.
All of them are at most
$\gamma_0^{(k_\circ(y))}$, with $k_\circ(y)$ the density of the first date at which $y$ is raw.
Hence, at every transition of $y$ from unheld to held across which the maximum defining
$\gamma_0^{(\cdot)}$ has grown, the $J$-inclusion fails as a set inclusion for every such
re-assignment, by the probe pair above. This conclusion is conditional on such a transition, which
the hypotheses admit and which no statement used here produces or excludes. The
configuration in which this matters for the arrest theorem is therefore not excluded by the
statements used here; whether Ba{\l}aban's definition of $\mathfrak{D}_m$ excludes it is among the
facts of (ii) and (viii). The maximal such dating, $T_\circ(y,\cdot):\equiv\gamma_0^{(k_\circ(y))}$,
imposed at every raw point, held or not, would lower the held-point datum, and with it the $J$-radius
that the consumers of the continuation statements read. It is not adopted and is used nowhere.

The gap would close in any of three ways:
\begin{itemize}
\item if no stored value and no size bound involving $K_\Delta\gamma_0$ ($K_\Delta$ the constant with
which that size bound multiplies $\gamma_0$) reads $\mathbb{J}$ at such
points, for then $\mathbb{J}$ may be cut to $0$ there;
\item by a $k$-free index-$0$ $J$-datum, provided $K_\Delta\gamma_0$ is checked not to become a lower
bound on the coupling when $\gamma_0$ is read at one number;
\item by restricting the sources of the consumers.
\end{itemize}

$(\ast)$ Second, there is no room at index $0$ on source domains; this is the content of the
statement that there is no index-zero margin on source domains
(Lemma~\ref{4.27:prf-index-zero-margin}, whose proof is incomplete at one step). At raw points of
$Y'$, and for a stored $\tau$ at raw points of $X_\tau\cap Y'$, no statement used here provides an
index-$0$ clause with a factor below that of $\mathfrak{G}$. Three facts stand behind this:
\begin{itemize}
\item Ba{\l}aban's attempt spaces start at $m=1$ \cite[p.~190]{B-CMP122a};
\item on a class component $\Omega''_1\cap X=\Omega_1\cap X$;
\item the freezing conventions move nothing there.
\end{itemize}
Whether \cite[(2.34)--(2.39)]{B-CMP119} carry an $m=0$ clause with factor $s^{(k)}_0$ (the stage
factor $s^{(k)}_m$ at $m=0$) on
$\Gamma_0$ intersected with the support of $\Omega_1$ in the sense of \cite[p.~255]{B-CMP119} is not
determined here. So source and consumer read the same
index-$0$ clause, while the map moves $A'$ and $\mathbb{J}$ there by layers. At the $\mathbf{T}$-step
the layers satisfy
\begin{equation*}
N^{\Delta}_y(\mathfrak{h})\le B_\sharp p(y),\qquad
\ell^3|\mathcal{J}^\sharp[1]|\le2K_\chi p(y),\qquad
p=w\,e^{-\delta d^{\wedge}(\cdot,\mathcal{C})},
\end{equation*}
with $d^{\wedge}$ at least $\mathsf{w}_M$ times the density of the step, from index $0$; these bounds
are small, and no statement used here shows that the layers vanish at raw points. The other layers
are
stated only in units of levels $\ge 1$:
\begin{itemize}
\item the stage shift layers use at most $\frac12\rho_m$ of the one-stage room $\rho_m$ at level
$m$, a room stated for levels $\ge1$ only;
\item the increments of the layers of the $\mathbf{R}$-operation off the zone are stated in units of
shells carrying a level;
\item on $\mathsf{Z}$ and on the core, the re-presented pair $(U_2,\mathcal{J}(U_2))$, $U_2$ the
re-presented field at the completion site, is bounded at
thresholds whose level is the level of the point in force or a finer birth index.
\end{itemize}
None of these is stated at a raw point, so there is no margin under any dating. The same holds for
the real point: the real-base lemma bounds clauses at levels $\ge1$, and at index $0$ only $A'=0$,
hence $(3.37)_0$, is immediate.

The statements of this section therefore exclude three kinds of index-$0$ members:
\begin{itemize}
\item the members moved on $Y'\cap(\Omega_0\setminus\Omega_1)$;
\item the index-$0$ members of the re-presented pair at raw points of $\mathsf{Z}$ and of the core;
\item the index-$0$ gauge and $J$-member of the real point.
\end{itemize}
The index-$0$ gauges of the $G$-part carry over wherever the map is a layer into $U'$.

This second gap would close in any of the following ways:
\begin{itemize}
\item if the layers of the maps are supported in $\Omega_1(\cdot)$, so that nothing moves at raw
points of the domains of stored terms;
\item if the core bounds and the bound on the presented pair bound the re-presented pair at raw
points of $\mathsf{Z}$ and of the core, and the real base pair obeys the index-$0$ clause there with
room;
\item if \cite[(2.34)--(2.39)]{B-CMP119} carry an $m=0$ clause with a factor;
\item by an index-$0$ margin stated outright.
\end{itemize}

The freezing of the $J$-datum requires one further property: no estimate at raw points of a live
class component may read a $J$-radius larger than $\gamma_0^{(k_{\mathfrak{a}})}$ of the innermost
live piece.
Among the statements used here none does. The hypotheses of the theorem on $\mathbb{H}$ and of the
perturbation bounds of the $\mathbf{T}$-step are met by smaller $\mathbb{J}$, and the threshold
tables read units at levels $\ge1$. $(\ast)$ That no other estimate of a stored term at such points
does so is not established here.

In the $\mathbf{T}$-step the hypothesis on the stratum is read as the membership
\begin{equation*}
\mathfrak{K}_{109}(\mathfrak{B};\gamma_0^{(k+1)})\subset\mathfrak{K}_{109}(\mathfrak{B};\gamma^{\top}).
\end{equation*}
Here $\gamma^{\top}:=\bar\gamma_0(\mathsf{t}_\gamma)\ge\gamma_0^{(k'')}$ for every density $k''$, and
the
theorem on $\mathbb{H}$ holds with its $\gamma$-free constants, because its ceilings are suprema
attained at $\mathsf{t}_\gamma$. No $k$-free dating enters any rider or any stored size, and
$K_\Delta\gamma_0$ reads $\gamma_0^{(k_\tau)}\le\bar\gamma_0(\mathsf{t}_{k_\tau})$, with $k_\tau$ the
density of the birth date of the stored term $\tau$. The ambient membership on $X_{\mathfrak{a}}$ and
the clause-free space of pairs of the type of $\mathfrak{K}_{109}$ used in the arrest are likewise
memberships, so the consumers of the second inclusion are untouched by the passage to
$\mathfrak{G}_1$, $\mathfrak{G}^{-J_0}$. No lower bound on the coupling of the
kind that $K_\Delta\gamma_0$ would create arises, and $\mathfrak{G}_1$, $\mathfrak{G}^{-J_0}$ involve
no $\gamma_0$. Finally, $\mathfrak{K}_{109}$ writes
the $J$-member with ``$\le$'' while $\mathfrak{G}$ writes it open; this makes no difference.

\emph{(v).} $(\ast)$ In the bound $\le1.1\,r_a$ on $\mathsf{Z}\cap\mathfrak{T}_k$, the set
$\mathfrak{T}_k=\Omega_k\cup Z^{c}$ is taken at the density of the completing chain. The own units at
$y$ are those of $\mathfrak{B}^t$, for the following reasons.
\begin{itemize}
\item On $\Omega_k$: in every $\mathbb{B}^{(n)}_k$, $0\le n\le N_k$, the set $\Omega_k\cap X$ has
level
$k$ and weight $1$. The constant $c=3$ of \cite{B-CMP122a} is a factor multiplying radii and is not
part of the unit.
\item On $Z^c\cap\Omega_k^c$: the stage spaces read $\mathfrak{T}_k$ natively, and a live piece of
another class there carries the same frozen units in the statement of the bound $\le1.1\,r_a$ and in
$\mathfrak{B}^t$.
\item No piece dissolving at the operation meets $\mathfrak{T}_k$, since
$W_{\mathfrak{c}}\subset X_{\mathfrak{c}}\subset Z$ and $W_{\mathfrak{c}}\subset Z_h$ with
$Z_h\cap\Omega_k=\emptyset$.
\item The re-levelling of \cite[(1.33)]{B-CMP122b} comes after $\mathfrak{B}^t$ and touches $Z$ only.
\end{itemize}
This reading is needed. Suppose $y$ were read two levels up in $\mathfrak{B}^t$ while the bound is
stated in $r_a(\lambda)$. The comparison of thresholds under a rise of levels would then give
\begin{gather*}
5\theta^{\mathfrak{B}^t}_a\le5\cdot3^{-2}\chi_2^{-1}r_a(\lambda)=1.026\,r_a(\lambda)
<1.1\,r_a(\lambda),\\
\text{since}\quad 9.9\,\chi_2=9.9\cdot\tfrac{49}{64}\cdot2^{-1/2}=5.36>5.
\end{gather*}
This is attained when the weight two levels up is $\mathsf{L}_0^{4}$ times the weight at $y$'s own
level, with $\mathsf{L}_0^2=L/3$ and $e_a=1$, where $L^{-e_a}$ is, up to the ratio $\chi_g$, the
factor by which
the threshold $r_a$ decreases per level; a smaller $\beta_0$ only worsens it. The allowance
$1.1\,r_a$
would then not fit. A misreading by one level would still fit, since
$5(1+\beta_0)^2/3=2.18\ge1.1$; from two levels on it fails. The four reasons above are what secure
the reading.

\emph{(vi).} $(\ast)$ The thin rim $Z''_{i+1}\setminus Z''_i$, $i>k_0$, is described in
\cite[p.~179]{B-CMP122a} as
``separated by one layer of $M$-cubes \dots\ one layer of $L^{-1}M$-cubes, and so on''. From this
wording it cannot be decided whether the rim is one $M$-cube of scale $i$ or of scale $i+1$ wide. No
conclusion depends on it. Each rim is at least one $M$-cube of its scale wide on either reading.
The depth in the nested-cube lemma only grows with the rims. A contour pays at least the stratum's
weight for each stratum it crosses of any determining set, in own-level units.

\emph{(vii).} $(\ast)$ The separations used with the nested-cube lemma read Ba{\l}aban's
``component'' (\cite[p.~255]{B-CMP119}, \cite{B-CMP122a}) as a connected component of the closed union
of blocks, with blocks sharing a boundary point being connected. Under this reading, distinct
components of a region of aligned $\mathsf{S}$-blocks are disjoint closed sets, hence at least
$\mathsf{S}$ apart. If Ba{\l}aban's connectivity were nearest-neighbour on sites, boxes touching at a
corner would be distinct components at distance $0$, and the separation step would need to be
restated for that reading.

\emph{(viii).} $(\ast)$ The facts to be read from Ba{\l}aban's papers, from the nesting declaration,
or from statements of this chapter used by statement, are the following.
\begin{itemize}
\item That the nesting declaration implies the one-box condition of (i), namely that the set
\eqref{4.27:eq-one-box-1} shares a cell with at most one component, at the one
locus of (i). The declaration-free alternative $L\ge(1+\beta_0)^2\mathsf{L}_0^4$ is a cap and is not
adopted, and the route through clause~(b) at $C=1$ reads nothing.
\item The whole of the bound on the presented pair and of the core bounds: entry $8$ among ``each
entry'' on $\mathsf{Z}\cap Z^{\mathrm{on}}$, and the coverage of $Y'\setminus\mathsf{Z}$ by the list
of the layer bound off the zone, as in (iv).
\item The definition of $\mathfrak{D}_m$ \cite[near (1.61)--(1.63)]{B-CMP122a}, as in (ii). It also
bears on whether the domain of a class term can reach raw points of a component that is not a class
component.
\item The sentence of \cite{B-CMP122b} on re-inclusion into the large-field domain, with its
surroundings: does \cite[(1.33)]{B-CMP122b} apply to a component completing as second class?
\item \cite[(2.34)--(2.39)]{B-CMP119}: the index range, and whether an $m=0$ clause with factor
$s^{(k)}_0$ exists on $\Gamma_0$ intersected with the support of $\Omega_1$ in the sense of
\cite[p.~255]{B-CMP119}.
\end{itemize}
None of these is part of the content of Theorem~\ref{4.27:thm-arrest-theorem}, whose proof is
incomplete at one step.
\end{proof}

\begin{remark}[The current radius read at raw points of live pieces. Stated; proof incomplete at the step marked $(\ast)$]\label{4.27:rem-raw-current-radius}
Let $y$ be a raw point, that is, a point of index $0$: $y\in\Omega_0(\mathfrak{B})\setminus\Omega_1(\mathfrak{B})$
for the determining set $\mathfrak{B}$ of the date at which $y$ is read, and suppose that $y$ is held
by at least one live piece. Let $\mathfrak{a}$ be the
innermost live arrest whose held set $W_{\mathfrak{a}}$ contains $y$, and let $k_{\mathfrak{a}}$ be its
density. In the arrested reading of the spaces (Definition~\ref{4.27:def-arrested-spaces}), the source
variable at $y$ is subject to the condition
\begin{equation}\label{4.27:eq-raw-current-radius-1}
\ell^{3}|\mathbb{J}(y)|<\gamma_0^{(k_{\mathfrak{a}})},\qquad \ell=\eta \text{ at index } 0 .
\end{equation}
The question is whether any estimate at $y$ that is made where terms of an arrest are stored or
consumed requires a radius in $\mathbb{J}$ larger than $\gamma_0^{(k_{\mathfrak{a}})}$
there, that is, whether it fails under the condition \eqref{4.27:eq-raw-current-radius-1}. The
estimates concerned are of the following four kinds:
\begin{enumerate}
\item a Cauchy estimate in $\mathbb{J}$ for a newborn term, taken on a polydisc in $\mathbb{J}$ about the real point;
\item the use of a source space, or of a tube, that must contain the pairs with $\ell^{3}|\mathbb{J}(y)|$ up to $\gamma_0^{(k')}$, where $k'$ is the current density;
\item a translation $\mathbb{J}\mapsto\mathbb{J}+\mathcal{J}$ whose size is measured against $\gamma_0^{(k')}$;
\item a reading at $y$ of the size $K_{\Delta}\gamma_0$ prescribed by the third size condition, $K_{\Delta}$ being the constant of that condition.
\end{enumerate}
Inside the present section no step reads such a larger radius (proof below); at the five places
at which stored terms of an arrest are kept and consumed the question is not settled (the step
marked $(\ast)$).

Under an alternative dating that would read at $y$ the radius of the density $k_0(y)$ of the first
raw date of $y$, the radius in question would be
\begin{equation*}
\gamma_0^{(k_0(y))}\le\gamma_0^{(k_{\mathfrak{a}})}
\end{equation*}
in place of $\gamma_0^{(k_{\mathfrak{a}})}$, which makes the question more demanding, and the separate
question concerning density increases at raw points held by no live piece throughout the birth
density would then not arise. That dating is not the one in force here.
\end{remark}

\begin{proof}
Within this section no step reads a radius in $\mathbb{J}$ at index $0$ as a source datum. We go through the places where such a reading could occur.

The change-of-variables theorem of the $\mathbf{T}$-step requires that its data lie in the space
$\mathfrak{K}_{109}(\mathfrak{B};\gamma_0^{(k)})$ of pairs of the type of \cite{B-CMP109}, at the
source radius $\gamma_0^{(k)}$ of the density $k$ at which the step is taken. The smallness
hypothesis on the source and the source bound that accompany it in this section are
hypotheses of the same kind. Each of these is a condition of the form ``$\mathbb{J}$ is small
enough'', so it is met by every $\mathbb{J}$ of smaller size. None of these conditions is a reading
of a larger radius at $y$.

The bounds for the third operation in the list of later operations of the arrest theorem
(Theorem~\ref{4.27:thm-arrest-theorem}, whose proof is incomplete at one step), and that theorem's
table of rooms, spends and ratios, read units only at points of level $\ge1$. They therefore read
nothing at a raw point.

Consequently, inside the arrested reading and its inputs as established in this section, none of the four kinds of estimate listed in the remark reads, at a raw point of $W_{\mathfrak{a}}$, a radius exceeding $\gamma_0^{(k_{\mathfrak{a}})}$.

$(\ast)$ It remains to settle the question at the five statements, outside this section, at which the
stored terms of an arrest are kept and consumed. At each of them it must be shown that none of the
four kinds of estimate listed in the remark reads at $y$ a radius larger than
$\gamma_0^{(k_{\mathfrak{a}})}$. The five places are the following.
\begin{enumerate}
\item The correlation theorem. There the question concerns its domains, units, tubes and recursion,
and it enters through the forms of the arrest theorem used in later steps
(Remark~\ref{4.27:rem-later-step-forms}, whose proof is incomplete at one step).
\item The second statement to which, together with the correlation theorem, the arrested reading of
the spaces (Definition~\ref{4.27:def-arrested-spaces}) is passed. The question
enters there through the same forms of the arrest theorem
(Remark~\ref{4.27:rem-later-step-forms}, whose proof is incomplete at one step).
\item The statement at which stored data are lifted to a later density. The question enters there
through the variant, used for that lifting, of the statement fixing the source radius
$\gamma_0^{(k)}$ at density $k$, through the slots of that statement, and through the lifted data
and the lifting map.
\item A fourth statement at which stored terms of an arrest are kept and consumed.
\item The statement that uses the domain $\mathfrak{D}_H$. The question enters there through that
domain, with its ambient space restricted to the region occupied by the
$\mathsf{W}$-pieces of that statement.
\end{enumerate}
This is not established in this section at any of the five places.
\end{proof}

\begin{remark}[Unheld raw points read at a later density. Stated; proof incomplete at the step
marked $(\ast)$]\label{4.27:rem-unheld-later-density}

We use the dates of this section. They are totally ordered by $\preceq$, and the density $k(d)$ of a date $d$ is non-decreasing along $\preceq$. We write $\mathfrak{B}_d$ for the determining set in force at $d$.

A point $y$ is \emph{raw} at $d$ if $y\in(\Omega_0\setminus\Omega_1)(\mathfrak{B}_d)$. A live arrest $\mathfrak{c}$ \emph{holds} $y$ at $d$ if $y\in W_{\mathfrak{c}}$ and $\mathfrak{c}$ is live at $d$.

At a raw point $y$ the stratum current member of $\mathfrak{G}(\mathfrak{B})$ is the clause $\ell^{3}|\mathbb{J}(y)|<T^{\mathfrak{B}}(y)$. Recall from the arrested reading of the spaces (Definition~\ref{4.27:def-arrested-spaces}) the birth-frozen dating $T^{\mathfrak{B}}(y)$:
\begin{itemize}
\item $T^{\mathfrak{B}}(y)=\gamma_0^{(k_{\mathfrak{c}})}$ if some live arrest holds $y$ at the date of $\mathfrak{B}$, where $\mathfrak{c}$ is the innermost one.
\item $T^{\mathfrak{B}}(y)=\gamma_0^{(k)}$, with $k$ the density of $\mathfrak{B}$, if no live arrest holds $y$ at the date of $\mathfrak{B}$.
\end{itemize}
The function $\gamma_0^{(\cdot)}$ is that of \eqref{4.27:eq-weights-heredity-a}.

Let $\mathfrak{a}$ be a live arrest, and let $\tau\in\mathfrak{F}(\mathfrak{a})$ be a stored term. It is born over $\mathfrak{B}_\tau$ at density $k_\tau=k_{\mathfrak{a}}$ and has domain $X_\tau$.
\begin{enumerate}
\item[(i)] \emph{One density.} Let $d\preceq d'$, and let $y$ be raw at both dates. Then
$T^{\mathfrak{B}_{d'}}(y)=T^{\mathfrak{B}_d}(y)$ in each of the following cases:
 \begin{itemize}
 \item $k(d')=k(d)$, whether $y$ is held at $d$ or not;
 \item whatever $k(d')$, $y$ is held at $d$ by a live arrest;
 \item whatever $k(d')$, $y$ is held by a live arrest at some date of density $k(d)$ following $d$.
 \end{itemize}
 Consequently, at every raw point of $X_\tau$ the stratum current member of $\mathfrak{G}(\mathfrak{B}_\tau)$ coincides with that of the space in force at every source and every target of density $k_\tau$ of a later operation reading $\tau$. These are the later stages of the run at density $k_\tau$, the completion at density $k_\tau$, and the target of the $\mathbf{T}$-step from $k_\tau$ to $k_\tau+1$.
\item[(ii)] \emph{No completion of the dating.} Fix a point $y$ that is raw at some date. Let $I_y$ be the set of dates at which $y$ is raw, and let $k_{\circ}(y)$ be the density of the first of them. A \emph{dating} at $y$ is an assignment $d\mapsto T(y,d)>0$, $d\in I_y$, of the threshold of the stratum current member at $y$, every other clause of $\mathfrak{G}$ being unchanged.
 \begin{itemize}
 \item The current parts at $y$ of both inclusions in \eqref{4.27:eq-weights-heredity-b} hold for all $d\preceq d'$ in $I_y$ if and only if two conditions hold: $T(y,\cdot)$ is non-increasing along $I_y$, and $T(y,d)\le\gamma_0^{(k(d))}$ on $I_y$. Every such dating is bounded by $\gamma_0^{(k_{\circ}(y))}$.
 \item Let $d\prec d'$ in $I_y$. Suppose that $y$ is held by no live arrest at any date of density
 $k(d)$ from $d$ on, and that $y$ is held at $d'$. Suppose further that the oldest live holder of
 $y$ at $d'$ has density $k_{\mathfrak{b}}$ with $\gamma_0^{(k_{\mathfrak{b}})}>\gamma_0^{(k(d))}$.
 Then, for every dating with $T(y,d')=\gamma_0^{(k_{\mathfrak{b}})}$ and
 $T(y,d)\le\gamma_0^{(k(d))}$, the current part at $y$ of
 $\mathfrak{G}(\mathfrak{B}_{d'})\subset\mathfrak{G}(\mathfrak{B}_d)$ fails.
 \end{itemize}
\item[(iii)] \emph{The remaining case.} Let $y$ be a raw point of $X_\tau$ satisfying both of the following:
 \begin{itemize}
 \item $y$ lies in a component of $Z_{k_\tau}$ that is not a class component of step $k_\tau$;
 \item $y$ is held by no live arrest at any date of density $k_\tau$.
 \end{itemize}
 Let $\mathfrak{B}'$, of density $k'>k_\tau$, be the source of a later operation, with $y$ raw at
 $\mathfrak{B}'$. Then $T^{\mathfrak{B}'}(y)$ equals $\gamma_0^{(k_{\mathfrak{b}})}$, where $\mathfrak{b}$ is the oldest
 live holder of $y$ at $\mathfrak{B}'$ and $k_\tau<k_{\mathfrak{b}}\le k'$, or equals
 $\gamma_0^{(k')}$ if $y$ has no live holder at $\mathfrak{B}'$; hence
 \begin{equation}\label{4.27:eq-unheld-later-density-1}
 T^{\mathfrak{B}'}(y)\ \ge\ \gamma_0^{(k_\tau+1)}\ \ge\ \gamma_0^{(k_\tau)}\ =\ T^{\mathfrak{B}_\tau}(y),
 \end{equation}
 and $T^{\mathfrak{B}'}(y)=T^{\mathfrak{B}_\tau}(y)$ if and only if the maximum
 $\max_{j\le\cdot}w_{*}(j)\bar\alpha_j$ has not grown on $\mathopen]k_\tau,k_{\mathfrak{b}}]$,
 respectively on $\mathopen]k_\tau,k']$. At such $y$ the current part of the first inclusion of
 \eqref{4.27:eq-weights-heredity-b} between $\mathfrak{B}_\tau$ and $\mathfrak{B}'$ is not asserted;
 this is the step marked $(\ast)$ in the proof below.
\item[(iv)] \emph{A dating not adopted.} For a point $y$ lying in $\Omega_0(\cdot)$ at some date, let
 $k_{\circ}(y)$ denote the density of the first date at which $y\in\Omega_0(\cdot)$; at every point
 that is ever raw this is the density $k_{\circ}(y)$ of clause (ii), as shown in the proof below. Put
 \begin{equation}\label{4.27:eq-unheld-later-density-2}
 T^{\flat}(y):=\gamma_0^{(k_{\circ}(y))}.
 \end{equation}
 Let $\mathfrak{G}^{\flat}(\mathfrak{B})$ be $\mathfrak{G}(\mathfrak{B})$ with its stratum current member read at every raw point of $\mathfrak{B}$, held or not, as $\ell^3|\mathbb{J}(y)|<T^{\flat}(y)$. Then:
 \begin{itemize}
 \item $T^{\flat}(y)$ does not depend on the date at which it is read;
 \item $T^{\flat}(y)\le\gamma_0^{(k)}$ whenever $y\in\Omega_0(\mathfrak{B})$ and $\mathfrak{B}$ has density $k$.
 \end{itemize}
 Granted the entering line \eqref{4.27:eq-raw-stratum-a} and the remaining clauses of
 Lemma~\ref{4.27:lem-weights-heredity} (whose proof is incomplete at one step), under this dating the
 current parts of both inclusions of \eqref{4.27:eq-weights-heredity-b} would hold at every raw point,
 held or unheld. This dating is not adopted in this book, and no statement of this book uses it.
\end{enumerate}
\end{remark}

\begin{proof}
\emph{Clause \textup{(i)}.} Put $k:=k(d)$.

Suppose first that $y$ is held at $d$ by a live arrest. The part of the heredity on the raw stratum (Lemma~\ref{4.27:prf-raw-stratum}) that treats held raw points then gives $T^{\mathfrak{B}_{d'}}(y)=T^{\mathfrak{B}_d}(y)$ for as long as $y$ stays raw.

Suppose next that $y$ is held at $d$ by no live arrest. We apply the statement on unheld raw points
within one density (Lemma~\ref{4.27:prf-unheld-raw-points}; its proof is incomplete at one step), in
the form of its display \eqref{4.27:eq-unheld-raw-points-a}.
\begin{itemize}
\item Its first clause treats $k(d')=k$. Every live holder of $y$ at $d'$ is an arrest of density $k$
attached at a date in the interval $\mathopen]d,d']$. Hence
$T^{\mathfrak{B}_{d'}}(y)=\gamma_0^{(k)}=T^{\mathfrak{B}_d}(y)$.
\item Its second clause treats the case where an arrest $\mathfrak{b}$ holds $y$ at some date of
density $k$ after $d$. Then $\mathfrak{b}$ is the only arrest of density $k$ ever holding $y$. Hence
$T^{\mathfrak{B}_{d''}}(y)=\gamma_0^{(k)}=T^{\mathfrak{B}_d}(y)$ for every $d''\succeq d$ at which $y$
is raw, whatever its density: for $d''$ before the attachment of $\mathfrak{b}$ by the first clause,
and from the attachment on because $\mathfrak{b}$ is then the innermost live holder of $y$.
\end{itemize}

For $\tau$ we use the statement on current thresholds across an increase of density
(Lemma~\ref{4.27:prf-density-increase}; its proof is incomplete at one step), in the form of its
display \eqref{4.27:eq-density-increase-a}. Its first three cases cover the raw points of $X_\tau$ that
lie in $X_{\mathfrak{a}}$, in a class component of step $k_\tau$, or in a piece live at
$\mathfrak{B}_\tau$, and give $T^{\mathfrak{B}'}(y)=T^{\mathfrak{B}_\tau}(y)$ at every later source
and target $\mathfrak{B}'$ at which $y$ is raw.

It remains to treat a raw point $y$ of $X_\tau$ held by no live arrest at $\mathfrak{B}_\tau$ and lying in a component $C$ of $Z_{k_\tau}$ that is not a class component of step $k_\tau$; then $T^{\mathfrak{B}_\tau}(y)=\gamma_0^{(k_\tau)}$. No arrest of density $k_\tau$ ever holds $y$. Indeed, its component would be a class component of step $k_\tau$ containing $y$, hence equal to $C$, which is not a class component. So the second clause of \eqref{4.27:eq-unheld-raw-points-a} (Lemma~\ref{4.27:prf-unheld-raw-points}, whose proof is incomplete at one step) is void at $y$. Its first clause then gives $T^{\mathfrak{B}'}(y)=\gamma_0^{(k_\tau)}=T^{\mathfrak{B}_\tau}(y)$ for every $\mathfrak{B}'\succeq\mathfrak{B}_\tau$ of density $k_\tau$.

By the fourth case of \eqref{4.27:eq-density-increase-a} (Lemma~\ref{4.27:prf-density-increase}, whose proof is incomplete at one step), at density $k_\tau$, the point $y$ stays raw through all these dates. They comprise:
\begin{itemize}
\item the later stages of the run;
\item the completion and its image under the re-levelling of the completed components;
\item the last date of density $k_\tau$, which is the target of the $\mathbf{T}$-step from $k_\tau$ to $k_\tau+1$.
\end{itemize}
This is clause (i).

\emph{Clause \textup{(ii)}.} The characterisation is the first clause of the statement that hereditary current thresholds are the non-increasing ones (Proposition~\ref{4.27:prop-nonincreasing-thresholds}). Its bound holds because a non-increasing dating satisfies $T(y,d)\le T(y,d_{\circ})\le\gamma_0^{(k_{\circ}(y))}$, where $d_{\circ}$ is the first date of $I_y$.

Now let $d\prec d'$ be as in the second part of clause (ii), and put $k:=k(d)$.
\begin{itemize}
\item Since $y$ is held at $d'$ and at no date of density $k$ from $d$ on, $d'$ is not of density $k$;
densities being non-decreasing along $\preceq$ and $d\prec d'$, we get $k(d')>k$.
\item By the third clause of \eqref{4.27:eq-unheld-raw-points-a}
(Lemma~\ref{4.27:prf-unheld-raw-points}, whose proof is incomplete at one step), the live holders of
$y$ at $d'$ are attached after $d$ at densities in $\mathopen]k,k(d')]$. The oldest of them,
$\mathfrak{b}$, is the innermost.
\item For every dating with $T(y,d')=\gamma_0^{(k_{\mathfrak{b}})}$ and $T(y,d)\le\gamma_0^{(k)}$ we get
\begin{equation*}
T(y,d')=\gamma_0^{(k_{\mathfrak{b}})}>\gamma_0^{(k)}\ge T(y,d),
\end{equation*}
so the dating is not non-increasing.
\item A separating pair is the one-point current probe $P_y(c)$ of
Lemma~\ref{4.27:lem-current-probe}, with $T(y,d)\le\ell^3c<\gamma_0^{(k_{\mathfrak{b}})}$. It obeys
every clause of $\mathfrak{G}(\mathfrak{B}_{d'})$, including the stratum current member at $y$, which
is read there because $y$ is raw at $d'$ and holds since
$\ell^3c<\gamma_0^{(k_{\mathfrak{b}})}=T(y,d')$; the entry $Q_3=|\mathbb{J}|$ vanishes at every point of
$\Omega_1(\mathfrak{B}_{d'})$, which does not contain $y$. It violates the current member of
$\mathfrak{G}(\mathfrak{B}_d)$ at $y$.
\end{itemize}

The growth $\gamma_0^{(k_{\mathfrak{b}})}>\gamma_0^{(k)}$ is admitted by the conditions of this section. The hypotheses of the arrest theorem bound neither $K$ nor the growth of the maximum $\max_{j\le\cdot}w_*(j)\bar\alpha_j$. It is not claimed that such a transition occurs in an actual run of the renormalization steps; the statement quantifies over the datings and histories that these conditions admit.

\emph{Clause \textup{(iii)}.} The bound \eqref{4.27:eq-unheld-later-density-1} is the part of the fourth case of \eqref{4.27:eq-density-increase-a} (Lemma~\ref{4.27:prf-density-increase}, whose proof is incomplete at one step) that lies past density $k_\tau$. There $T^{\mathfrak{B}'}(y)$ equals one of two values:
\begin{itemize}
\item $\gamma_0^{(k_{\mathfrak{b}})}$, where $\mathfrak{b}$ is the oldest live holder attached after the dates of density $k_\tau$, so that $k_\tau<k_{\mathfrak{b}}\le k'$;
\item $\gamma_0^{(k')}$, if there is no such holder.
\end{itemize}
In both cases $T^{\mathfrak{B}'}(y)=T^{\mathfrak{B}_\tau}(y)$ holds if and only if the maximum
$\max_{j\le\cdot}w_{*}(j)\bar\alpha_j$ has not grown on $\mathopen]k_\tau,k_{\mathfrak{b}}]$,
respectively on $\mathopen]k_\tau,k']$. The sets $\mathfrak{B}'$ in question are the sources of the
$\mathbf{T}$-step from $k_\tau$ to $k_\tau+1$ and of all later $\mathbf{T}$-steps, of the stages and
completions at densities larger than $k_\tau$, and of the completions dissolving a chain. At those
$\mathfrak{B}'$ at which $y$ is held and across which the maximum has grown, clause (ii) shows that no
re-assignment of the dating at unheld raw points alone makes the current part at $y$ of the first
inclusion of \eqref{4.27:eq-weights-heredity-b} hold between $\mathfrak{B}_\tau$ and $\mathfrak{B}'$,
as long as two things stand: the birth-frozen dating at held raw points, and the second inclusion of
\eqref{4.27:eq-weights-heredity-b}. Where $y$ is unheld at $\mathfrak{B}'$ as well, the constant
dating $\gamma_0^{(k_{\circ}(y))}$ of clause (ii) would serve; the obstruction sits at the
birth-frozen dating at held raw points. At a raw point never held before it is lifted into
$\Omega_1(\cdot)$ the constant dating serves throughout; but a dating rule cannot tell in advance
which points these are, and the transitions from unheld to held of clause (ii) are admitted by the
conditions of this section.

$(\ast)$ The current part of the first inclusion at the points of clause (iii) is not proved here. It
would be settled by a negative answer to question (a) below, by a negative answer to question (b)
below for the stored terms $\mathfrak{F}(\mathfrak{a})$, or by adopting the dating of clause (iv),
which this book does not do.

\emph{Question \textup{(a)}.} Does the value of any stored term, or any estimate of the step in which
the current radius $\gamma_0^{(k)}$ enters through its product $K_{\Delta}\gamma_0^{(k)}$ with a fixed
constant $K_{\Delta}$, read $\mathbb{J}$ at raw points of its own domain of the following kind?
\begin{itemize}
\item They lie in a component of $Z_{k_\tau}$ that is not a class component of step $k_\tau$. Such
components are those failing at least one of the two conditions defining the class; among them are the
components holding a second-class component of an earlier step that has been included again into the
large-field region.
\item They are held by no live arrest at any date of density $k_\tau$.
\end{itemize}
For $\tau\in\mathfrak{F}(\mathfrak{a})$ these are exactly the points of clause (iii), read only by
sources of density larger than $k_\tau$. For a stored space of another statement that imports the
first inclusion of \eqref{4.27:eq-weights-heredity-b}, they are the raw points of its domain that are
unheld at every date of its birth density and are later met by a source of larger density.

If the answer is no, the current may be cut to zero on the domain at such points. This is allowed
because the space $\mathfrak{K}_{109}(\mathfrak{B};\gamma)$ admits, for every $\gamma$, any
$\mathfrak{g}^{\mathbb{C}}$-valued current one-form subject only to the bound on
$\sup\ell^3|\mathbb{J}|$, with no divergence condition on $\mathbb{J}$ and no second-order condition.
The stratum current member then leaves
$\mathfrak{G}$ at these points. The inclusions \eqref{4.27:eq-weights-heredity-b}, and the clauses of
the arrest theorem and of the statements importing them, then extend to the stratum current member of
$\mathfrak{G}(\mathfrak{B}_\tau)$ everywhere, with nothing left to date.

\emph{Question \textup{(b)}.} Can a term born in the run of the class at step $k_\tau$ have a domain
$X_\tau\in\mathfrak{D}_m$ that reaches raw points of another component of $Z_{k_\tau}$? Here
$\mathfrak{D}_m$ is the class of domains of \cite{B-CMP122a} that enters
\cite[(1.63)]{B-CMP122a}. Two facts bear on this:
\begin{itemize}
\item distinct components of $Z_{k_\tau}$ are at least $\mathsf{S}_{k_\tau}$ apart, by
\eqref{4.27:eq-stage-move-cubes-a};
\item \cite[(1.63)]{B-CMP122a} indexes the terms of stage level $j$ by domains $X$ with
\begin{equation*}
X\in\mathfrak{D}_m,\qquad X\cap(Z_{j+1}\setminus Z''_{j+1})\neq\emptyset;
\end{equation*}
no further condition on $X$ is used in this section.
\end{itemize}
The configuration is not excluded by the conditions used in this section. The set
$\mathfrak{T}_{k_\tau}=\Omega_{k_\tau}\cup Z^{c}$, with $Z$ the region of the class at step $k_\tau$,
contains every component of $Z_{k_\tau}$ that is not
a class component, so that the domains of the stored terms are permitted to meet such components, and
no bound on the diameter of a domain is available here. It is decided by the definition of
$\mathfrak{D}_m$ in \cite{B-CMP122a}; no claim is made that it occurs in a run.
If the answer is no, the points of clause (iii) do not occur for $\tau\in\mathfrak{F}(\mathfrak{a})$.
For the stored spaces of other statements importing the first inclusion, question (b) does not dispose
of these points, and question (a) remains.

Suppose the answer to (a) is yes, and the answer to (b) is yes for $\mathfrak{F}(\mathfrak{a})$. One
could then date the unheld raw points by a $k$-free number, such as
$\gamma^{\top}=\bar\gamma_0(\mathsf{t}_\gamma)$. Dated so, the unheld raw points would satisfy the
current part of the first inclusion, since
\begin{equation*}
\gamma_0^{(k_{\mathfrak{b}})}\le\bar\gamma_0(\mathsf{t}_{k_{\mathfrak{b}}})\le\gamma^{\top}
\end{equation*}
by \eqref{4.27:eq-weights-heredity-a}, $\bar\gamma_0$ being decreasing on the range
$\mathsf{t}\ge\mathsf{t}_\gamma$ in which $\mathsf{t}_{k_{\mathfrak{b}}}$ lies. But such a dating lies
outside the two standing conditions of clause (ii) (the birth-frozen dating at held raw points and the
second inclusion of \eqref{4.27:eq-weights-heredity-b}), because it moves the target of the second
inclusion at these points to $\mathfrak{K}_{109}(\mathfrak{B};\gamma^{\top})$. The alternative is
admissible only if the estimates in which $K_{\Delta}\gamma_0^{(k)}$ enters, read at the single number
$\gamma^{\top}$, do not re-create the floor on the coupling removed in the proof of
Lemma~\ref{4.27:lem-weights-heredity} (whose proof is incomplete at one step). Such a floor would be a
cap on $K$. The alternative stands or falls with question (a), and it is not proposed. Failing such a
dating, the sources of the later operations would have to be shown to avoid
such points. The remaining alternative is the dating of clause (iv).

\emph{Clause \textup{(iv)}.} Let $k_{\circ}(y)$ and $T^{\flat}$ be as in clause (iv),
\eqref{4.27:eq-unheld-later-density-2}. At every point that is ever raw, the first date at which the point lies in $\Omega_0(\cdot)$ is the first date of $I_y$, so that the two descriptions of $k_{\circ}(y)$ in clauses (ii) and (iv) agree. Indeed, $\Omega_1(\cdot)$ never shrinks along $\preceq$, by the heredity on the raw stratum (Lemma~\ref{4.27:prf-raw-stratum}). So a point lying in $\Omega_1$ at its first date in $\Omega_0$ lies in $\Omega_1$ at every later date and is never raw. Hence, at every point $y$ that is ever raw, $T^{\flat}(y)$ read on $I_y$ is the pointwise-maximal hereditary dating $T_{\circ}(y,\cdot)$ of Proposition~\ref{4.27:prop-nonincreasing-thresholds}.

The first assertion of clause (iv) is immediate from \eqref{4.27:eq-unheld-later-density-2}.

For the second assertion of clause (iv), let $y\in\Omega_0(\mathfrak{B})$ with $\mathfrak{B}$ of density $k$. The first date of $y$ in
$\Omega_0(\cdot)$ precedes or equals the date of $\mathfrak{B}$, and densities are non-decreasing
along $\preceq$, so $k_{\circ}(y)\le k$. Since $\gamma_0^{(\cdot)}$ is non-decreasing by
\eqref{4.27:eq-weights-heredity-a}, it follows that $T^{\flat}(y)\le\gamma_0^{(k)}$.

By the first assertion of clause (iv), at a point raw at both dates $\mathfrak{B}\preceq\mathfrak{B}'$ the stratum current members of $\mathfrak{G}^{\flat}(\mathfrak{B}')$ and $\mathfrak{G}^{\flat}(\mathfrak{B})$ are the same inequality.

Now let $y\in(\Omega_0\setminus\Omega_1)(\mathfrak{B})$ enter $\Omega_1(\mathfrak{B}')$ at a date
$\mathfrak{B}_e$ of density $k_e$. As in the proof of the second assertion of clause (iv), $k_{\circ}(y)\le k$, where $k$ is the density of
$\mathfrak{B}$, and $k\le k_e$ since $\mathfrak{B}\preceq\mathfrak{B}_e$; hence $k_e\ge k_{\circ}(y)$. The
entering line \eqref{4.27:eq-raw-stratum-a}, run at $k_{\circ}(y)\le k_e$, gives for every pair of
$\mathfrak{G}_1(\mathfrak{B}')$
\begin{equation}\label{4.27:eq-unheld-later-density-3}
\begin{aligned}
\ell^{3}|\mathbb{J}(y)|&<5\alpha_{0,k_e}L^{-3k_e}\le5\alpha_{0,k_{\circ}(y)}\le5\bar\alpha_{k_{\circ}(y)}\\
&\le\gamma_0^{(k_{\circ}(y))}=T^{\flat}(y).
\end{aligned}
\end{equation}
We justify the second inequality in two cases.
\begin{itemize}
\item If $k_e=k_{\circ}(y)$, it holds because $L^{-3k_{\circ}(y)}\le1$.
\item If $k_e-k_{\circ}(y)\ge1$, then \eqref{4.27:eq-factors-radii-coupling-a} gives
$\alpha_{0,k_e}\le(1+\beta_0)^2(k_e-k_{\circ}(y))^{1/2}\alpha_{0,k_{\circ}(y)}$. The hypotheses of the arrest
theorem include $\beta_0\le1/7$, so $(1+\beta_0)^2\le64/49$; and $L\ge3$, the blocking factor $L$ being
an odd integer greater than one. Since $k_e-k_{\circ}(y)\le k_e$,
\begin{equation*}
\tfrac{64}{49}\,(k_e-k_{\circ}(y))^{1/2}L^{-3k_e}\le\tfrac{64}{49}\,k_e^{1/2}3^{-3k_e}
\le\tfrac{64}{49}\cdot\tfrac1{27}<1 .
\end{equation*}
Here $k^{1/2}3^{-3k}$ is decreasing for $k\ge1$, because the ratio of consecutive values is $((k+1)/k)^{1/2}/27<1$.
\end{itemize}
The third inequality uses $\alpha_{0,j}\le\mathsf{A}(\mathsf{t}_j)\le\bar\alpha_j$. The fourth uses
$w_{*}(\cdot)\ge1$ in \eqref{4.27:eq-weights-heredity-a}.

The second inclusion follows from the second assertion of clause (iv) together with the remaining clauses of \eqref{4.27:eq-weights-heredity-b}. With those remaining clauses as in Lemma~\ref{4.27:lem-weights-heredity}, whose proof is incomplete at one step, the dating of clause (iv) would give
\begin{equation*}
\mathfrak{G}^{\flat}(\mathfrak{B}')\subset\mathfrak{G}^{\flat}(\mathfrak{B})
\subset\mathfrak{K}_{109}(\mathfrak{B};\gamma_0^{(k)})
\end{equation*}
as sets on all of $\Omega_0(\mathfrak{B})$, at held and unheld points alike. The exception at the
points of clause (iii) in the clause of the arrest theorem that places the continuations of the
sources in $\mathfrak{G}(\mathfrak{B}_\tau)$ (the rider clause), and in the statements importing it,
would disappear.

In particular the memberships at the radius $\gamma_0^{(k_\tau)}$ asked of the stored terms at birth
would hold a fortiori at the raw points of $\Omega_0(\mathfrak{B}_\tau)$, since $k_{\circ}(y)\le k_\tau$
there by the second assertion of clause (iv), and no floor on the coupling of the kind produced by $K_{\Delta}\gamma_0^{(k)}$ would
arise. Estimates that use only an upper bound on the current radius are unaffected, since
$T^{\flat}\le\gamma_0^{(k_\tau)}\le\bar\gamma_0(\mathsf{t}_{k_\tau})$; estimates that would use a lower
bound on the current radius are the subject of the first condition below. The power
of $L$ in
\eqref{4.27:eq-unheld-later-density-3} is the change of units of the current variable (exponent three)
between level $0$ and level $k_e$, for a pair held in $\mathfrak{G}_1(\mathfrak{B}')$ at level $k_e$;
it is not a gain per step.

The dating \eqref{4.27:eq-unheld-later-density-2} is a new hypothesis and not a change of notation. The
birth-frozen dating freezes the threshold at the birth density of the holder,
$\gamma_0^{(k_{\mathfrak{c}})}$. The dating \eqref{4.27:eq-unheld-later-density-2} freezes it at the
first density of the point in $\Omega_0(\cdot)$. Let $y$ be a held raw point with innermost holder
$\mathfrak{c}$. Then $T^{\flat}(y)\le T^{\mathfrak{B}}(y)$, with
strict inequality whenever the maximum defining $\gamma_0^{(\cdot)}$ grew on
$\mathopen]k_{\circ}(y),k_{\mathfrak{c}}]$; such growth is admitted, since the hypotheses bound neither $K$
nor that growth. Whether $k_{\circ}(y)$ is, at every raw point, the density of the first date, one
number on the whole raw stratum, depends on how $\Omega_0(\cdot)$ is read: if it is read date-free, it
is; if $\Omega_0(\cdot)$ is read as the support of $\Omega_1(\cdot)$ of width $M_1\eta$
\cite[p.~255]{B-CMP119}, it is granted the inequality $M_1\le M$ between the auxiliary parameters of
\cite{B-CMP119}, and otherwise under a condition $M\ge M_1$ that would have to be declared.
This comparison is not used. Where the inequality $T^{\flat}(y)\le T^{\mathfrak{B}}(y)$ holds strictly at a held raw point $y$ of
$X_{\mathfrak{a}}$, the probe $P_y(c)$ of Lemma~\ref{4.27:lem-current-probe} with
$T^{\flat}(y)\le\ell^3c<T^{\mathfrak{B}_\tau}(y)$ lies in
$\mathfrak{G}(\mathfrak{B}_\tau)\setminus\mathfrak{G}^{\flat}(\mathfrak{B}_\tau)$, so that then
$\mathfrak{G}^{\flat}(\mathfrak{B}_\tau)\subsetneq\mathfrak{G}(\mathfrak{B}_\tau)$.

Adopting the dating would therefore change the following objects:
\begin{itemize}
\item the threshold of the arrested reading of the spaces (Definition~\ref{4.27:def-arrested-spaces});
\item the space $\mathfrak{G}$;
\item the heredity and membership clauses of Lemma~\ref{4.27:lem-weights-heredity};
\item the rider clauses of the arrest theorem, and the rider and membership clauses of the statements
importing them.
\end{itemize}
Two further conditions would then have to be met:
\begin{itemize}
\item The question whether an estimate of a stored term reads, at raw points of a live class component, a lower bound on the current radius would have to be posed at the smaller radius $\gamma_0^{(k_{\circ}(y))}\le\gamma_0^{(k_{\mathfrak{a}})}$ --- when $\Omega_0(\cdot)$ is read date-free, a single number, the value of $\gamma_0^{(\cdot)}$ at the density of the first date.
\item The index-$0$ current member of the real base pair $(U(\mathfrak{B},V),\mathcal{J}(U))$ of the
expansion, $\ell^3|\mathcal{J}(U)(y)|<\gamma_0^{(k_{\circ}(y))}$,
which is not asserted in this chapter, would have to hold below that radius.
\end{itemize}
The situation is therefore as follows. The remaining case is not closable inside the arrested reading under the two standing conditions of clause (ii). The dating \eqref{4.27:eq-unheld-later-density-2} alters the number fixed at held raw points. It is not adopted, nothing in this book uses it, and the remaining case stands as described at $(\ast)$.
\end{proof}

\begin{remark}[Index-zero members under maps, at the real point and at completion. Stated; proof incomplete at the step marked $(\ast)$]
\label{4.27:rem-index-zero-members}
Let $\mathfrak{B}$ be a determining set of density $k$, with domain $\Omega_0(\mathfrak{B})$ and set
$\Omega_1(\mathfrak{B})$ of points of level at least one. Write $\Omega_1$, without argument, for the
region of \cite[(2.2)]{B-CMP119}, and $\Gamma_0=\Omega_1^{c}$ for the raw stratum. One has
$\Omega_1\subseteq\Omega_1(\mathfrak{B})$, so every point of $(\Omega_0\setminus\Omega_1)(\mathfrak{B})$
lies in $\Gamma_0$. Call the points of $(\Omega_0\setminus\Omega_1)(\mathfrak{B})$ the
\emph{raw points} of $\Omega_0(\mathfrak{B})$. The stratum $\Gamma_0$ is the lowest of the
strata of \cite[(2.2)]{B-CMP119}, on which that construction prescribes no block average.

At a raw point $y$ the rider $\mathfrak{G}(\mathfrak{B})$, that is, the space of pairs
$(\mathbb{U},\mathbb{J})$ on $\Omega_0(\mathfrak{B})$ attached to $\mathfrak{B}$, carries the
index-zero clause of the space $\mathfrak{K}_{109}(\mathfrak{B};\gamma)$ of pairs of the type of
\cite{B-CMP109} with $J$-radius $\gamma$. This clause has two parts.
\begin{itemize}
\item The \emph{$U$-members}: writing $\mathbb{U}=e^{i\eta A'}U$ with $U$ group-valued and
$A'$ valued in $\mathfrak{g}^{\mathbb{C}}$, on the cubes of scale zero a gauge obeying the
index-zero instance of \cite[(3.35)]{B-CMP99b}, and, for $A'$, the index-zero instance of
\cite[(3.37)]{B-CMP99b}.
\item The \emph{$J$-member}: $\ell^{3}|\mathbb{J}(y)|<T^{\mathfrak{B}}(y)$, where $\ell=\eta$
at index zero and $T^{\mathfrak{B}}(y)$ is the dated threshold of the stratum; that is, the
$J$-clause of $\mathfrak{K}_{109}(\mathfrak{B};\gamma)$ with the radius $\gamma$ read at $y$ as
the dated threshold $T^{\mathfrak{B}}(y)$.
\end{itemize}
We write $\mathfrak{G}^{-J_0}(\mathfrak{B})$ for $\mathfrak{G}(\mathfrak{B})$ with the
$J$-member at raw points removed and the $U$-members kept.

For a later operation $\mathfrak{O}$ we use the following notation: $Y'$ is its source
domain, $\mathfrak{B}^t$ is the determining set in force immediately before its map, and
$\mathfrak{B}_\tau$ is the birth set of a stored term $\tau$. We formulate three
properties; none of them is asserted among the statements of this chapter.

\begin{enumerate}
\item[(a)] \emph{Structural vanishing at raw points.} For every configuration $V$ for which the
minimiser $U(\mathfrak{B},V)$ attached to $V$ and $\mathfrak{B}$ is defined, this minimiser
coincides with $V$ on $\Gamma_0$.
Consider data supported inside $\Omega_1(\mathfrak{B})$, namely:
\begin{itemize}
\item the datum $b$ of the $\mathbf{T}$-step taken at density $k'$, on its bond set
$\mathcal{C}\subset\Omega_{k'+1}$;
\item the shift data of the stages, on their bands, which lie in
$\Omega''_{h+1}\subset\Omega_1$, where $h=k-N_k$ and $N_k$ is the number of stages of the
large-field operation at density $k$;
\item the data of the completion site.
\end{itemize}
For such data the following vanish at the raw points of $\Omega_0(\mathfrak{B})$:
\begin{itemize}
\item the layer $\mathfrak{h}=\mathbb{H}(b)$ of the $\mathbf{T}$-step, bounded in
\eqref{4.27:eq-t-step-layer-a};
\item the source increment $\mathcal{J}^{\sharp}[1]$ of the $\mathbf{T}$-map;
\item the shift layers of the stage maps $\Phi_n$ and $\Phi^{\chi}$;
\item the layers at the completion site.
\end{itemize}
Here ``vanish'' is meant up to the range of $U\mapsto\mathcal{J}(U)$ near the boundary of
$\Omega_1$, which a statement of this property would have to specify. Moreover, the
statement on the layers at the completion site treats the shell of index zero as the
identity.
\item[(b)] \emph{The real point at index zero.} The real point
$(U(\mathfrak{B},V),\mathcal{J}(U))$ of Lemma~\ref{4.27:lem-real-minimiser}, whose proof is
incomplete at one step, obeys the
index-zero clause of $\mathfrak{K}_{109}(\mathfrak{B};\gamma)$ at raw points of term
domains, with room. That is, its gauge obeys the index-zero instance of
\cite[(3.35)]{B-CMP99b} on the cubes of scale zero, and
\begin{equation}\label{4.27:eq-index-zero-members-1}
\ell^{3}|\mathcal{J}(U)(y)|<T^{\mathfrak{B}}(y)
\end{equation}
holds with room.
\item[(c)] \emph{The re-presented pair at the completion site.} At the completion site, one
of the following holds:
\begin{itemize}
\item the re-presented pair $(U_2,\mathcal{J}(U_2))$ and the completion layers are bounded by
their thresholds at raw points of the set $\mathsf{Z}$ of the completion site and of the core
$D_2\setminus B_S$;
\item $\mathsf{Z}$ and the re-presented core are disjoint from $\Gamma_0$.
\end{itemize}
\end{enumerate}

The three properties have the following consequences, and their failure the following
alternatives.

\emph{If \textup{(a)} holds,} every map of the later operations is the identity on
$(A',\mathbb{J})$ at raw points of its source domain $Y'$. The map clause of
Lemma~\ref{4.27:prf-ambient-inclusion}, the inclusion of the ambient domain at size
$\gamma_0^{(k)}$, whose proof is incomplete at one step, then extends
there by identity to $\mathfrak{G}^{-J_0}(\mathfrak{B}^t)$. Likewise, the rider-membership
clause of part~(A) of the arrest theorem (Theorem~\ref{4.27:thm-arrest-theorem}), whose proof
is incomplete at one step, extends there by
identity to $\mathfrak{G}^{-J_0}(\mathfrak{B}_\tau)$. The extension also reaches the
$J$-member of $\mathfrak{G}$ at every raw point whose dating is fixed by one of the
following:
\begin{itemize}
\item the dating at raw points held by a live piece, that is, lying in the held set
$W_{\mathfrak{c}}$ of an arrest $\mathfrak{c}$ that is live at the date in question;
\item the dating at raw points of the domains $X_\tau$ of the stored terms $\tau$ of an
arrest;
\item the dating at unheld raw points between dates of one density.
\end{itemize}
On $Y'$ there then remain only the raw points held by no live piece at any date of the
density of $\mathfrak{B}^t$, respectively of $\mathfrak{B}_\tau$, from its date on, and read by
a source of a strictly larger density.

\emph{If \textup{(a)} fails,} that is, if some map moves $(A',\mathbb{J})$ at raw points, then
the map clauses at raw points require an index-zero margin. By this we mean a factor at
level zero of the kind the arrested spaces carry at levels $\ge1$; no such factor is
introduced in this chapter. Alternatively, if among the factors of
\cite[(2.34)--(2.39)]{B-CMP119} there is one of index $m=0$ on the intersection of $\Gamma_0$
with the support of $\Omega_1$ in the sense of \cite[p.~255]{B-CMP119}, then the room of the
$\mathbf{T}$-step off live pieces at raw points is the one that factor gives in
\cite{B-CMP119}.

\emph{If \textup{(a)} holds, then in \textup{(b)}} one has $U=V$ on $\Gamma_0$. Property~(b)
is then a small-field condition of the first renormalisation step on the restriction
$V\restriction(\Omega_0\setminus\Omega_1)$, together with a gauge for it.

\emph{If \textup{(c)} holds,} either because the bounds hold or because the sets are
disjoint, then the rider-membership clause of part~(A) of
Theorem~\ref{4.27:thm-arrest-theorem}, whose proof is incomplete at one step, extends at the
completion site to the stratum $U$-members there. With the datings listed above, it
extends to the $J$-member as well.

\emph{An alternative dating.} Suppose the threshold at a raw point $y$ were dated by the
radius $T^{\flat}(y)=\gamma_0^{(k_{\circ}(y))}$ of its first raw date. Then
\eqref{4.27:eq-index-zero-members-1} would read
$\ell^{3}|\mathcal{J}(U)(y)|<\gamma_0^{(k_{\circ}(y))}$; that is, the member would be bounded
by the smallest of the radii $\gamma_0^{(k(d))}$ over all dates $d$ at which $y$ is raw.
This is not the dating used in this chapter.
\end{remark}

\begin{proof}
We record what the available statements give at raw points, and we mark the step at which
they stop.

\emph{The layers at raw points.} We appeal to Lemma~\ref{4.27:lem-t-step-layer}, whose proof
is incomplete at one step. In the bound \eqref{4.27:eq-t-step-layer-a} of that lemma the
local norm $N^{\Delta}_y(\mathfrak{h})$ of the layer is bounded by a constant times the decay
profile $p(y)$, where $p(y)$ is a positive weight times
$e^{-\delta d^{\wedge}(y,\mathcal{C})}$, and is therefore strictly positive. The contour
bound of the same lemma states that a contour pays at least $\mathsf{w}_M$ for each stratum
and each $M$-cube of own scale it crosses. From a point of index zero to
$\mathcal{C}\subset\Omega_{k'+1}$ a contour crosses the strata $1,\dots,k'-1$ and one $M$-cube
of own scale in $\Omega_{k'}\setminus\Omega_{k'+1}$, and therefore has
$d^{\wedge}\ge\mathsf{w}_M k'$. This makes the bound small at raw
points. It does not make it zero, so \eqref{4.27:eq-t-step-layer-a} does not decide~(a).

The vanishing statement for the layer in the same lemma does not decide (a) either. That
statement holds on the connected components of the layer's domain that carry no datum.
However, the component of the support of $\mathfrak{B}$ containing a raw point also contains
the datum.

\emph{The real point at index zero.} We appeal to Lemma~\ref{4.27:lem-real-minimiser}, whose
proof is incomplete at one step. That lemma, and the bound, among the hypotheses of the arrest
theorem, that places the real point in the arrested spaces at every level $\ge1$, place
$(U(\mathfrak{B},V),\mathcal{J}(U))$ in every space whose clause at $x$ is
taken at the current level of $x$, a level $\ge1$. At index zero only one member follows
at once. The real point has $A'=0$ there, and this is the index-zero instance of
\cite[(3.37)]{B-CMP99b}. The index-zero gauge condition of \cite[(3.35)]{B-CMP99b} and the
$J$-member \eqref{4.27:eq-index-zero-members-1} are not contained in either statement.

If (a) holds, then $U=V$ on $\Gamma_0$. The two missing members are then conditions on
$V\restriction(\Omega_0\setminus\Omega_1)$ alone, that is, a small-field condition of the
first renormalisation step and a gauge for it. Such a condition is not among the inputs of
this chapter.

\emph{The completion site.} The bounds on the re-presented pair $(U_2,\mathcal{J}(U_2))$ and
on the completion layers at the completion site compare them with thresholds taken at the
current level $o(x)$ of $x$ or at a finer birth index. A raw point has level zero, so at a
raw point of $\mathsf{Z}$ or of the core no threshold of these bounds applies.

Nor is the core known to avoid $\Gamma_0$. The core contains the dissolving pieces
$W_{\mathfrak{c}}\subset\Omega^{c}_{k_{\mathfrak{c}}}$, and these pieces may meet $\Gamma_0$.

\emph{The consequences.} Suppose the layers of the maps vanish at raw points of $Y'$. Then
each map is the identity on $(A',\mathbb{J})$ there. The image of a pair therefore coincides
with the input at those points, and the map clauses pass by identity. This is the form in
which the consequence of~(a) is stated in the remark. For~(c), the pair on $\mathsf{Z}$ and on
the core is the re-presented one, and the extension to the stratum $U$-members follows either
from the bounds at the raw points there or from the absence of raw points there. In the
failure of~(a) an index-zero margin on source domains would be needed; that none is supplied
by the statements at hand is the content of Lemma~\ref{4.27:prf-index-zero-margin}, whose
proof is incomplete at one step.

$(\ast)$ The argument stops at the three points above. Properties (a), (b) and (c) are not
derived from the statements available in this
chapter. The bound \eqref{4.27:eq-t-step-layer-a} and the vanishing statement of
Lemma~\ref{4.27:lem-t-step-layer}, whose proof is incomplete at one step, leave~(a) open.
Lemma~\ref{4.27:lem-real-minimiser}, whose proof is incomplete at one step, gives
only the index-zero instance of \cite[(3.37)]{B-CMP99b} towards~(b). The bounds at the
completion site contain no threshold at level zero, which leaves~(c) open.
\end{proof}

\begin{remark}[What the published analyticity spaces say at index zero]
\label{4.27:rem-index-zero-reading}
Let $\Omega_1\supset\Omega_2\supset\dots\supset\Omega_j$ be the small-field regions of the
multi-scale construction of \cite{B-CMP119}. Let
\[
  \Omega_0=\Omega_1^{\sim}
\]
be the enlargement of $\Omega_1$ fixed in \cite[p.~262]{B-CMP119}. We call
$\Omega_0\setminus\Omega_1$ the \emph{collar}. Let $X$ be a domain of a term of the effective
action, and let $(\mathbb{U},\mathbb{J})$ be complex bond variables and a current, the arguments
on which the analyticity spaces are defined.
\begin{enumerate}
\item[(a)] The remark on \cite[p.~191]{B-CMP122a} beginning ``Finally let us remark'' asserts
exactly two things:
\begin{itemize}
\item[($\alpha$)] when the integrations in a component of $Z$ are stopped after $n$ of them
because a new large field has appeared, the terms produced by the $n$
integrations already performed remain in the action, as part of the boundary terms;
\item[($\beta$)] a definition of the analyticity spaces for the subsequent steps exists and is
easy to give.
\end{itemize}
The remark contains no definition of the generalised spaces, no statement that a later
configuration belongs to one of them, and no estimate. Neither assertion refers to the collar.
\item[(b)] In the multi-scale analyticity spaces of \cite[p.~261]{B-CMP119} and
\cite[p.~190]{B-CMP122a}, the conditions on $(\mathbb{U},\mathbb{J})$ carry indices starting at
$1$. Consequently they impose no condition at index zero; for \cite{B-CMP119} this is the set
$X\cap(\Omega_0\setminus\Omega_1)$.
The terms, however, may depend on the fields on the collar.
\item[(c)] The one-scale definitions \cite[(1.11), (1.13), (1.14)]{B-CMP109} impose their
conditions on all of $X$. The displays \cite[(3.35), (3.37)]{B-CMP99b} are written for
$j=0,\dots,k$, so the index $j=0$ is included.
\end{enumerate}
\end{remark}

\begin{proof}
(a) The remark of \cite[p.~191]{B-CMP122a} reads:
\begin{quotation}
Finally let us remark that if in a component of $Z$ we finish the integrations for some $n$,
because some new large field has been introduced, then we leave the new terms
[produced by these $n$ integrations] as a part of all boundary terms, and it is easy to generalize
the definition of the spaces in such a case for the next steps.
\end{quotation}
Its first clause, ``we leave the new terms \dots\ as a part of all boundary terms'', is
($\alpha$). Its second clause, ``it is easy to generalize the definition of the spaces in such a
case for the next steps'', is ($\beta$). The sentence contains no further assertion: it gives no
definition of the generalised spaces, does not assert that any later configuration lies in one of
them, and contains no estimate; neither clause mentions the set $\Omega_0\setminus\Omega_1$. The
spaces for the steps that follow such an interrupted sequence of integrations are defined in the
definition that accompanies the arrest theorem.

(b) In \cite[p.~261]{B-CMP119} the conditions defining the space are imposed
\begin{quote}
on $X\cap(\Omega_n\setminus\Omega_{n+1})$ for $n=1,\dots,j-1$, or on $X\cap\Omega_j$ for $n=j$.
\end{quote}
The union of these sets is $X\cap\Omega_1$.

In \cite[p.~190]{B-CMP122a} the conditions are imposed
\begin{quote}
on $(\Omega''_m\setminus\Omega''_{m+1})\cap X$ for $m=1,\dots,j-1$.
\end{quote}
Here the $\Omega''_m$ are the regions of that construction, entering these conditions as the
$\Omega_n$ enter those of \cite[p.~261]{B-CMP119}. Again the index starts at $1$.

Hence neither definition writes a clause at index zero; in particular no clause of the definition
of \cite[p.~261]{B-CMP119} concerns $X\cap(\Omega_0\setminus\Omega_1)$.

On the other hand, \cite[p.~262]{B-CMP119} sets $\Omega_0=\Omega_1^{\sim}$ and considers
$X\cap\Omega_0$ as a union of components. It states:
\begin{quote}
If a component does not connect to $\Omega_j$ \dots\ then the function does not depend on the
fields $(U,J)$, $A$ restricted to this component
\end{quote}
(the fields $(U,J)$, $A$ in the notation of \cite[p.~262]{B-CMP119}).
This statement excludes dependence only on the components that do not connect to
$\Omega_j$. On the components that do connect, a term may therefore depend on the fields on the
part $X\cap(\Omega_0\setminus\Omega_1)$ of the collar.

(c) The one-scale definition \cite[(1.11), (1.13), (1.14)]{B-CMP109} requires
\[
  |\partial U-1|<\alpha_0\xi^2, \qquad |A'|,\ |\nabla A'|<\alpha_1, \qquad
  |\partial\mathbb{U}-1|<\alpha_0\xi^2, \quad |\mathbb{J}|<\gamma_0 ,
\]
where, in the notation of \cite{B-CMP109}, $\partial U$ denotes the plaquette variables of $U$,
$\xi$ the lattice spacing (as in Definition~\ref{4.4:def-regular-backgrounds}), and
$\alpha_0,\alpha_1,\gamma_0$ the radii of the spaces; $A'$ and $\nabla A'$ are as in
\cite{B-CMP109}. Each of these inequalities is imposed on all of $X$. The displays
\cite[(3.35), (3.37)]{B-CMP99b} are written for $j=0,\dots,k$, which includes the index $j=0$.
\end{proof}

\begin{definition}[Level-zero conventions on the never-averaged collar]\label{4.27:not-collar-conventions}
On matrices, $|\cdot|$ is the operator norm; it is submultiplicative and invariant under unitary
conjugation. We write $\eta$ for the level-zero lattice spacing, so that a cell of level $j$ has size
$L^{j}\eta$. For a determining set $\mathfrak{B}$, the sets $\Omega_0(\mathfrak{B})\supset\Omega_1(\mathfrak{B})$
are the first two members of its decreasing sequence of domains in the sense of
\cite[(3.35)--(3.37)]{B-CMP99b}. The
clock $\check{R}_k$ and the radii $\alpha_{i,m}$ are those fixed earlier in this section.

\emph{The level-zero domain.} Dates $d$, their total order and the set $\mathfrak{B}_d$ of a date are
those of the ordering lemma for dates in the setting of the arrest theorem. At every set
$\mathfrak{B}=\mathfrak{B}_d$ of a date $d$ we set
\begin{equation}\label{4.27:eq-collar-conventions-a}
\Omega_0(\mathfrak{B}) := \Omega_1(\mathfrak{B})^{\sim},
\end{equation}
that is, $\Omega_1(\mathfrak{B})$ enlarged by one layer of cubes of side $LM\check{R}_1$, with the
operation ${}^{\sim}$ as in \cite[p.~262]{B-CMP119}. This is an admissible choice of Ba{\l}aban's
$\Omega_0$: \cite{B-CMP99b} allows $\Omega_0$ to be formed by adding to $\Omega_1$ a thinner layer of the
big blocks surrounding $\Omega_1$. In every operator statement the level-zero cells of $\mathfrak{B}$ are
those of $\Omega_0(\mathfrak{B})\setminus\Omega_1(\mathfrak{B})$. Fields, sources and currents are
considered on $\Omega_0(\mathfrak{B})$ only.

\emph{The collar.} For a date $d$ we put
$\operatorname{St}(d):=\Omega_0(\mathfrak{B}_d)\setminus\Omega_1(\mathfrak{B}_d)$. This set is the
\emph{never-averaged collar} of $d$. When $\mathfrak{B}$ is fixed we write
$\operatorname{St}:=\Omega_0(\mathfrak{B})\setminus\Omega_1(\mathfrak{B})$.

\emph{Level-zero radii.} For $a\in\{1,3,5,6,7\}$ we put
\begin{equation}\label{4.27:eq-collar-conventions-1}
\begin{gathered}
r_1(0)=r_5(0):=\alpha_{0,0},\qquad r_3(0):=\alpha_{0,0}\,\eta^{-3},\\
r_6(0):=\alpha_{1,0}\,\eta^{-1},\qquad r_7(0):=\alpha_{1,0}\,\eta^{-2}.
\end{gathered}
\end{equation}
The entries $a=1$ and $a=5$ refer to $|\partial\mathbb{U}-1|$ and $|\partial U-1|$ respectively.

\emph{Bonds and plaquettes.} A bond is written $b=\langle b_-,b_-+\nu\rangle$. A plaquette $p$ spanned by
two directions $\mu<\nu$ at a point $x$ has the corners $x$, $x+\mu$, $x+\nu$, $x+\mu+\nu$, and $x$ is its
\emph{base point}. A bond or a plaquette \emph{of} $\Omega_0(\mathfrak{B})$ is one all of whose vertices
lie in $\Omega_0(\mathfrak{B})$. As in \cite{B-CMP102b}, a bond with an end-point in
$\Omega_1(\mathfrak{B})$, and a plaquette with a corner in $\Omega_1(\mathfrak{B})$, belong to
$\Omega_1(\mathfrak{B})$. Two further notions are used below:
\begin{itemize}
\item a \emph{stratum bond} is a bond of $\Omega_0(\mathfrak{B})$ with both ends in $\operatorname{St}$;
\item a \emph{pinned plaquette} is a plaquette all four of whose bonds are stratum bonds.
\end{itemize}

\emph{The level-zero entries.} Let $(\mathbb{U},\mathbb{J})$ be a pair with
$\mathbb{U}=e^{i\eta A'}U$, and let $y\in\operatorname{St}$. The quantities $Q_1(y)$ and $Q_5(y)$ are the
maxima of $|\partial\mathbb{U}(p)-1|$ and of $|\partial U(p)-1|$, taken over the pinned plaquettes $p$ with
base point $y$. The quantities $Q_3(y)$ and $Q_6(y)$ are the maxima of $|\mathbb{J}(b)|$ and of
$|A'(b)|$, taken over the stratum bonds $b$ with $b_-=y$. Write $B_\nu(x)=B(\langle x,x+\nu\rangle)$ and
define the covariant difference by
\begin{equation}\label{4.27:eq-collar-conventions-2}
\eta\,(\nabla_{U,\mu}B_\nu)(x):=\operatorname{Ad}_{U(\langle x,x+\mu\rangle)}B_\nu(x+\mu)-B_\nu(x).
\end{equation}
Its stencil at $y$ consists of the bonds $\langle y,y+\mu\rangle$, $\langle y,y+\nu\rangle$ and
$\langle y+\mu,y+\mu+\nu\rangle$. The quantity $Q_7(y)$ is the maximum of
$|(\nabla_{U,\mu}A'_\nu)(y)|$ over the pairs $(\nu,\mu)$ of directions such that
$\langle y,y+\nu\rangle$ is a stratum bond and the stencil at $y$ lies in $\Omega_0(\mathfrak{B})$.
Of the list of entries of a pair, the entries with index $2$ and $4$ are not imposed at level zero, and in
the set $\mathfrak{G}(\mathfrak{B})$ of global pairs obeying the weighted clauses they are imposed at no
level: there only $a\in\{1,3,5,6,7\}$ occurs.

Call a bond or a plaquette \emph{at} $y$ if $y$ is one of its vertices. A stratum bond $b$ is
\emph{attached} to the point $b_-$, and a pinned plaquette to its base point.
These conventions have three consequences:
\begin{enumerate}
\item[(i)] the point to which a stratum bond or a pinned plaquette is attached lies in
$\operatorname{St}$, and it is the only point $y\in\operatorname{St}$ for which the object belongs to a
family over which one of the maxima $Q_a(y)$ is taken; thus every stratum bond and every pinned
plaquette is attached to exactly one $y\in\operatorname{St}$;
\item[(ii)] a bond of $\Omega_0(\mathfrak{B})$ at $y\in\operatorname{St}$ that is not a stratum bond
belongs to $\Omega_1(\mathfrak{B})$, and so does a plaquette of $\Omega_0(\mathfrak{B})$ at $y$ that is
not pinned; conditions on such objects are conditions on $\Omega_1(\mathfrak{B})$;
\item[(iii)] any bound on the maximum of the quantity defining $Q_1(y)$ or $Q_5(y)$ over all pinned
plaquettes at $y$, or of the quantity defining $Q_3(y)$ or $Q_6(y)$ over all bonds at $y$, implies the
same bound for $Q_a(y)$; for $Q_7$ the corresponding maximum is that of
$|(\nabla_{U,\mu}A'_\nu)(y)|$ over the directions $\mu$ and the bonds $\langle y,y+\nu\rangle$ at $y$.
In particular the collar estimate along the presentation chain, which controls the families of all
pinned plaquettes at $y$, of all bonds of the enlarged cell $\tilde{\Delta}(y)$ at $y$, and of all bonds
at $y$, gives the same estimate for $Q_a(y)$.
\end{enumerate}
\end{definition}

\begin{proof}
We prove the three consequences in turn.

(i) Let $b$ be a stratum bond. Both of its ends lie in $\operatorname{St}$, so in particular
$b_-\in\operatorname{St}$. The quantities $Q_3$ and $Q_6$ at a point $y$ are maxima over the stratum
bonds with $b_-=y$, and $Q_7$ at $y$ is a maximum over stratum bonds $\langle y,y+\nu\rangle$ only,
namely those whose stencil lies in $\Omega_0(\mathfrak{B})$; all of these have $b_-=y$. Hence $b$
belongs to one of these families only for the single point $y=b_-$, and for no other point.

Now let $p$ be a pinned plaquette. Each of its four bonds is a stratum bond, so every end of these bonds
lies in $\operatorname{St}$. These ends are the corners of $p$, and the base point of $p$ is one of its
corners, so the base point lies in $\operatorname{St}$. The quantities $Q_1$ and $Q_5$ at $y$ use exactly
the pinned plaquettes with base point $y$. Since the base point of $p$ is unique, $p$ enters them at
exactly one $y\in\operatorname{St}$.

(ii) Let $b$ be a bond of $\Omega_0(\mathfrak{B})$ at a point $y\in\operatorname{St}$, and suppose $b$ is
not a stratum bond. Then some end of $b$ is not in $\operatorname{St}$. That end lies in
$\Omega_0(\mathfrak{B})$, and $\Omega_0(\mathfrak{B})\setminus\operatorname{St}=\Omega_1(\mathfrak{B})$,
so the end lies in $\Omega_1(\mathfrak{B})$. By the membership convention of \cite{B-CMP102b},
$b$ therefore belongs to $\Omega_1(\mathfrak{B})$.

Next let $p$ be a plaquette of $\Omega_0(\mathfrak{B})$ at $y$, and suppose $p$ is not pinned.
Then one of its bonds is not a stratum bond. That bond joins two corners of $p$, so it is a bond of
$\Omega_0(\mathfrak{B})$. By the same argument as for bonds, one of its ends lies in
$\Omega_1(\mathfrak{B})$. This end is a corner of $p$, so by the membership convention $p$ belongs to
$\Omega_1(\mathfrak{B})$.

(iii) Each $Q_a(y)$ is a maximum over a family of plaquettes, of bonds, or, for $Q_7$, of pairs of a bond
and a direction. The family for $Q_1$ and $Q_5$ consists of pinned plaquettes with base point $y$; the
base point is a corner, hence a vertex, so these are pinned plaquettes at $y$. The families for $Q_3$ and
$Q_6$ consist of stratum bonds $b$ with $b_-=y$, and the bonds $\langle y,y+\nu\rangle$ entering $Q_7$
are of the same kind; since $y$ is a vertex of each of them, they are bonds at $y$. For $Q_7$ the pairs
$(\nu,\mu)$ used are among all pairs formed by a direction $\mu$ and a bond $\langle y,y+\nu\rangle$ at
$y$. A maximum over a subfamily does not exceed the maximum over the whole family. Hence a bound on the
maximum over all pinned plaquettes at $y$, respectively over all bonds at $y$, bounds $Q_a(y)$ by the
same quantity. The collar estimate along the presentation chain controls, among its three families, the
pinned plaquettes at $y$ and the bonds at $y$, so it gives the same estimate for $Q_a(y)$.
\end{proof}

\begin{definition}[Dated allowances and the collar clause of the domains]
\label{4.27:def-dated-allowances}
\emph{Scaffold dates.} For a level $k$ let $\check{R}_k$ be the value of the clock at level $k$, let
$h_k=k-\check{R}_k$ be the first level, and let $N_k$ be the number of stages at level $k$. The
\emph{scaffold dates} are the dates of the ordering lemma for sets of dates:
\begin{itemize}
\item $(k;\top)$ for every $k\ge 1$ (the averaging step at level $k$);
\item $(k;l)$ with $1\le l\le N_k$ (stage $l$), and $(k;L)$ (the completion site), for every
density index $k$ with $h_k\ge 1$, that is, at the levels at which the operation $\mathbf{R}$ acts.
\end{itemize}
At the other levels no class of stage or completion-site objects exists, and no such dates occur.
The scaffold dates are ordered by
\begin{equation*}
(k;\top)<(k;1)<\dots<(k;N_k)<(k;L)<(k+1;\top).
\end{equation*}

\emph{Allowances, cumulated allowance, stratum factor.} Let $\beta$ be the schedule constant of
the layered analyticity spaces (Definition~\ref{4.4:def-post-step-space}), and let
$\rho^{(l),(k)}_m$ be the allowance profile of stage $l$ at level $k$. To every scaffold date
$d$ we assign its \emph{allowance} $\mathfrak{c}(d)$, its \emph{cumulated allowance}
$\mathfrak{s}(d)$ and its \emph{stratum factor} $\mathsf{f}_0(d)$:
\begin{equation}\label{4.27:eq-dated-allowances-b}
\begin{gathered}
\mathfrak{c}(1;\top):=0,\qquad \mathfrak{c}(k;\top):=\beta\,2^{-k}\quad(k\ge 2),\\
\mathfrak{c}(k;l):=\rho^{(l),(k)}_0=\beta\,2^{-(h_k+l)},\qquad
\mathfrak{c}(k;L):=\tfrac14\,\beta\,2^{-(k-1)},\\
\mathfrak{s}(d):=\sum_{d''\le d}\mathfrak{c}(d''),\qquad \mathsf{f}_0(d):=4-\mathfrak{s}(d).
\end{gathered}
\end{equation}
The sum defining $\mathfrak{s}(d)$ runs over every scaffold date $d''\le d$, whether or not the
operation of that date is performed on the label under consideration (the datum that determines
at which scaffold dates operations are performed). For a date $d$ at level $k$, $\mathfrak{s}(d)$
is a function of $d$ and of $(\check{R}_i)_{i\le k}$ alone; it depends neither on the complex
sources $z$, nor on the slot at which it is evaluated, nor on the label.

\emph{Idle dates and predecessors.} A scaffold date not performed on a label (a stage that no
component runs, or a completion site at which no component completes) writes nothing: its set
$\mathfrak{B}_d$ is the set of its scaffold predecessor, and its allowance $\mathfrak{c}(d)$ is
spent without use. For an operation $\mathfrak{O}$ performed at a date $d'$ we write $d^t$ for
the scaffold predecessor of $d'$; then $\mathfrak{B}_{d^t}$ is the set in force immediately
before the map of $\mathfrak{O}$, in the sense of the arrest theorem.

\emph{Dating of spaces, sources and stored quantities.} Every space, source and stored quantity
of the construction in the arrest theorem is written at an operation, hence at a date:
for each level $k'$, the sources of $\mathbf{T}$ and the quantities created by $\mathbf{T}$ over
$\mathbb{B}_{k'+1}$ at $(k'+1;\top)$; the attempt spaces
$\mathfrak{S}^{(k),\mathrm{ARR}}_n$, $\mathfrak{L}_n$, $\mathfrak{D}^{\flat}_n$ at $(k;n)$; the
sources of $\mathbf{R}$ and the quantities created at the completion site over the set
$\mathbb{B}_k$ of \cite[(1.33)]{B-CMP122a} at $(k;L)$, by the ordering lemma for sets of dates.
Accordingly, in the ambient-rider
convention of the arrest theorem the words ``over a set $\mathfrak{B}$'' are read as ``over the
set $\mathfrak{B}_d$ of a date $d$''. A stored quantity $\tau$ reads its birth date $d_\tau$:
where the arrest theorem reads the birth set and the birth density, here, by the ordering lemma
for sets of dates, it reads the set $\mathfrak{B}_{d_\tau}$ of the birth date.

\emph{The collar clause.} Let $\operatorname{St}(d)=\Omega_0(\mathfrak{B}_d)\setminus
\Omega_1(\mathfrak{B}_d)$ be the never-averaged collar of the date $d$, that is, the part of the
domain $\Omega_0(\mathfrak{B}_d)$ of $\mathfrak{B}_d$ outside its averaged part
$\Omega_1(\mathfrak{B}_d)$ (Notation~\ref{4.27:not-collar-conventions}). On
$\operatorname{St}(d)$ the rider set
$\mathfrak{G}(\mathfrak{B}_d)$ imposes the following two requirements. First, at every
$y\in\operatorname{St}(d)$,
\begin{equation}\label{4.27:eq-dated-allowances-a}
Q_a(y)<\mathsf{f}_0(d)\,r_a(0),\qquad a\in\{1,3,5,6,7\}.
\end{equation}
Second, on the
index-$0$ cubes it imposes the cube condition \cite[(3.35)]{B-CMP99b} at $j=0$. These two
requirements take the place, on $\operatorname{St}(d)$, of the index-$0$ clause of the kernel class
$\mathfrak{K}_{109}(\mathfrak{B}_d;\gamma_0^{(k)})$.

\emph{Dated thresholds.} For $a\in\{1,3,5,6,7\}$ we put
\begin{equation*}
\Theta^{d}_a(y):=
\begin{cases}
5\,\theta^{\mathfrak{B}_d}_a(y), & y\in\Omega_1(\mathfrak{B}_d),\\[2pt]
\mathsf{f}_0(d)\,r_a(0), & y\in\operatorname{St}(d),
\end{cases}
\end{equation*}
and on $\operatorname{St}(d)$ the level and weight functions of the set $\mathfrak{B}_d$ are put
equal to $(\ell_{\mathfrak{B}_d},w_{\mathfrak{B}_d})(y):=(0,1)$.

\emph{On the constant $4$.} The constant $4$ in $\mathsf{f}_0(d)=4-\mathfrak{s}(d)$ is chosen
below the factor $5$ of the thresholds on $\Omega_1(\mathfrak{B}_d)$ for the following three
reasons. First, the inequality $4(1+\beta_0)<5$ makes the inclusion of the rider set
$\mathfrak{G}(\mathfrak{B})$ in the kernel class $\mathfrak{K}_{109}(\mathfrak{B};\gamma_0^{(k)})$
independent of whether the running maximum $\bar{\alpha}_j$, a maximum over $m\le j$, starts at
$m=0$ or at $m=1$; here $\beta_0$ denotes the constant with which that inclusion is established.
Next, $4<3(1+2\beta)=4.2$, so that the bound \eqref{4.27:eq-dated-allowances-a} lies below the
threshold $3c_\alpha\mathbf{r}_a(0)$ of the correlation theorem at $\pi=0$, should a statement of
that theorem ever read a bound of the form \eqref{4.27:eq-dated-allowances-a}; here $c_\alpha$
and $\pi$ are the constant and the
parameter of that threshold, and $\mathbf{r}_a(0)$ is the variant of the level-zero radius
$r_a(0)$ used there. Finally, no larger constant is needed: the largest object required to
satisfy \eqref{4.27:eq-dated-allowances-a} has $Q_a\le\tfrac12\,r_a(0)$.
\end{definition}

\begin{lemma}[Boundedness of the summed allowances]\label{4.27:lem-allowance-sum}
Let the scaffold dates, the allowances $\mathfrak{c}$ of \eqref{4.27:eq-dated-allowances-b}, the
cumulated allowance $\mathfrak{s}$ and the stratum factor $\mathsf{f}_0(d)=4-\mathfrak{s}(d)$ be as
in Definition~\ref{4.27:def-dated-allowances}. For every scaffold date $d$,
\begin{equation*}
0\le \mathfrak{s}(d)<\tfrac{7}{4}\beta=0.35 ,
\end{equation*}
so that $3.65<\mathsf{f}_0(d)\le 4$. Moreover $\mathsf{f}_0$ is non-increasing in $d$. If
$d^{t}<d'$ are consecutive scaffold dates, that is, $d^{t}$ is the scaffold date immediately
preceding $d'$, then $\mathsf{f}_0(d^{t})-\mathsf{f}_0(d')=\mathfrak{c}(d')$.
\end{lemma}

\begin{proof}
By Definition~\ref{4.27:def-dated-allowances}, $\mathfrak{s}(d)$ is the sum of $\mathfrak{c}(d'')$
over all scaffold dates $d''\le d$. Only finitely many scaffold dates precede $d$, because the
dates of each level $k$ are finite in number and are followed by $(k+1;\top)$. By
\eqref{4.27:eq-dated-allowances-b}, every allowance is a non-negative multiple of $\beta$. Hence
$\mathfrak{s}(d)\ge 0$, and $\mathfrak{s}(d)$ is bounded above by the sum of all allowances of each
of the three kinds of date. We bound the three contributions separately.

\emph{Averaging steps.} We have $\mathfrak{c}(1;\top)=0$ and $\mathfrak{c}(k;\top)=\beta 2^{-k}$ for
$k\ge 2$. The dates $(k;\top)$ with $(k;\top)\le d$ therefore contribute at most
\begin{equation*}
\sum_{k\ge 2}\beta 2^{-k}=\frac{\beta}{2}.
\end{equation*}

\emph{Stages.} The dates $(k;l)$, $1\le l\le N_k$, exist only at the levels $k$ at which
Ba{\l}aban's operation $\mathbf{R}$ acts; we call these levels $\mathbf{R}$-levels. At such
a level $h_k\ge 1$, where $h_k=k-\check{R}_k$. For one $\mathbf{R}$-level $k$ the allowances are
$\mathfrak{c}(k;l)=\beta 2^{-(h_k+l)}$, and
\begin{equation*}
\sum_{l=1}^{N_k}\beta 2^{-(h_k+l)}=\beta 2^{-h_k}\bigl(1-2^{-N_k}\bigr)<\beta 2^{-h_k}.
\end{equation*}
The clock $\check{R}$ is non-increasing in the setting of this section, so
$\check{R}_{k+1}\le\check{R}_k$. Hence
\begin{equation*}
h_{k+1}=k+1-\check{R}_{k+1}\ge k+1-\check{R}_k=h_k+1 .
\end{equation*}
Thus $k\mapsto h_k$ is strictly increasing, and the numbers $h_k$ belonging to distinct levels are
distinct integers, each $\ge 1$ at an $\mathbf{R}$-level. Consequently
\begin{equation*}
\sum_{k}\beta 2^{-h_k}\le\beta\sum_{h\ge 1}2^{-h}=\beta ,
\end{equation*}
where the left-hand sum runs over the $\mathbf{R}$-levels. Combining this with the strict bound for
each level, the stage dates preceding $d$ contribute strictly less than $\beta$. If no stage date
precedes $d$, their contribution is $0<\beta$.

\emph{Completion sites.} Each $(k;L)$ also lies at an $\mathbf{R}$-level $k$. By the standing
condition of the setting of this section at the levels at which an $\mathbf{R}$ acts, such a level
satisfies $k>\check{R}_k\ge 1$, so $k\ge 2$. With $\mathfrak{c}(k;L)=\tfrac14\beta 2^{-(k-1)}$,
these dates contribute at most
\begin{equation*}
\sum_{k\ge 2}\tfrac14\beta 2^{-(k-1)}=\tfrac14\beta\sum_{j\ge 1}2^{-j}=\tfrac14\beta .
\end{equation*}

Adding the three contributions, and using that the second one is strict, gives
\begin{equation*}
\mathfrak{s}(d)<\beta\Bigl(\tfrac12+1+\tfrac14\Bigr)=\tfrac74\beta=0.35 .
\end{equation*}
Since $\mathsf{f}_0(d)=4-\mathfrak{s}(d)$, the bounds $0\le\mathfrak{s}(d)<0.35$ are equivalent to
$3.65<\mathsf{f}_0(d)\le 4$.

For the last clause, let $d^{t}<d'$ be consecutive. The dates $\le d'$ are those $\le d^{t}$
together with $d'$ itself. Hence $\mathfrak{s}(d')=\mathfrak{s}(d^{t})+\mathfrak{c}(d')$, and
\begin{equation*}
\mathsf{f}_0(d^{t})-\mathsf{f}_0(d')=\mathfrak{s}(d')-\mathfrak{s}(d^{t})=\mathfrak{c}(d')\ge 0 .
\end{equation*}
Applying this along each pair of consecutive dates shows that $\mathsf{f}_0$ is non-increasing in
$d$.

The argument uses neither the numbers $N_k$ of stages, nor the threshold function $N_0$ of the
setting, nor the number $K$ of steps, nor any age or depth. In particular, the bound holds although
the number of scaffold dates preceding $d$ is not bounded uniformly in $d$.
\end{proof}

\begin{lemma}[Local algebra of a left exponential factor at unit spacing]
\label{4.27:lem-left-factor-algebra}
Let $y\in\operatorname{St}$, the never-averaged collar of
Notation~\ref{4.27:not-collar-conventions}, and let $\eta$ denote the unit of length at level zero.
Let $\mathbb{U}=e^{i\eta A'}U$ with $U$ $G$-valued, and let $\mathfrak{h}$ be a
$\mathfrak{g}^{c}$-valued $1$-form, all given on the bonds of
$\tilde{\Delta}(y)\cap\Omega_0(\mathfrak{B})$, with $\Omega_0(\mathfrak{B})$ the level-zero region of
Notation~\ref{4.27:not-collar-conventions} (the level-zero families of
Notation~\ref{4.27:not-collar-conventions} read no other bond). Put
\begin{equation*}
\mathsf{n}:=N_y(\mathfrak{h})
=\sup_{\tilde{\Delta}(y)}\max\bigl\{\eta|\mathfrak{h}|,\ \eta^{2}|\nabla_{\mathbb{U}}\mathfrak{h}|\bigr\},
\qquad
\mathsf{a}:=\sup_{\tilde{\Delta}(y)}\eta|A'|,
\end{equation*}
and let $\theta:=\max_p|\partial\mathbb{U}(p)-1|$, the maximum being taken over the plaquettes $p$
based at $y$, where $\partial\mathbb{U}(p)$ denotes the plaquette variable $\mathbb{U}(\partial p)$,
the ordered product of $\mathbb{U}$ around $p$. Assume
\begin{equation*}
\mathsf{n}\le \tfrac{1}{40},\qquad \mathsf{a}\le\tfrac{1}{80},\qquad \theta\le\tfrac{1}{10}.
\end{equation*}
Then, with $A'^{+}$ defined bondwise by
\begin{equation}\label{4.27:eq-left-factor-algebra-1}
A'^{+}:=(i\eta)^{-1}\log\bigl(e^{i\eta\mathfrak{h}}e^{i\eta A'}\bigr),
\end{equation}
one has $\mathbb{U}^{+}:=e^{i\eta\mathfrak{h}}\mathbb{U}=e^{i\eta A'^{+}}U$; the $G$-valued part $U$ is
unchanged, and so is every gauge of $U$; and
\begin{enumerate}
\item[(i)] $|\partial\mathbb{U}^{+}(p)-\partial\mathbb{U}(p)|\le 4\mathsf{n}$ for every plaquette $p$
based at $y$;
\item[(ii)] $\eta|A'^{+}(b)-A'(b)|\le 1.06\,\mathsf{n}$ for every bond $b$ in $\tilde{\Delta}(y)$;
\item[(iii)] $\eta^{2}|\nabla_{U,\mu}(A'^{+}-A')_{\nu}(y)|\le 1.2\,\mathsf{n}$ for all directions
$\mu,\nu$.
\end{enumerate}
\end{lemma}

\begin{proof}
Throughout, $|\cdot|$ denotes the matrix norm on $M_N(\mathbb{C})$, used only through its
submultiplicativity, $|1|=1$, and its invariance under conjugation by unitary matrices, and
$\|\cdot\|$ denotes the norm it induces on linear maps of $M_N(\mathbb{C})$. We put
$\operatorname{Ad}_{Z}C:=ZCZ^{-1}$ and $\operatorname{ad}_{F}C:=[F,C]$, so that
$\|\operatorname{ad}_F\|\le 2|F|$ and $\operatorname{Ad}_{e^{F}}=e^{\operatorname{ad}_F}$.
We write $y+\mu$ for the nearest neighbour of $y$ in the positive
$\mu$-direction, $\mathfrak{h}_\mu(y):=\mathfrak{h}(\langle y,y+\mu\rangle)$, and the covariant
derivative in the form
\begin{equation*}
(\nabla_{V,\mu}C_\nu)(y)=\eta^{-1}\bigl(\operatorname{Ad}_{V(\langle y,y+\mu\rangle)}C_\nu(y+\mu)
-C_\nu(y)\bigr)
\end{equation*}
for a configuration $V$ and a $\mathfrak{g}^{c}$-valued $1$-form $C$. The three clauses are proved in
the order (ii), (iii), (i).

\emph{Proof of} (ii). Fix a bond $b$ of $\tilde{\Delta}(y)$ and put $X:=i\eta\mathfrak{h}(b)$,
$Y:=i\eta A'(b)$, so that $|X|\le\mathsf{n}\le 1/40$ and $|Y|\le\mathsf{a}\le 1/80$. Expanding the
exponentials, for $t\in[0,1]$,
\begin{equation*}
|e^{tX}e^{Y}-1|\le e^{t|X|+|Y|}-1\le e^{3/80}-1<1 .
\end{equation*}
Hence $F(t):=\log(e^{tX}e^{Y})$, with the principal branch
$\log(1+Z)=\sum_{n\ge1}(-1)^{n+1}Z^{n}/n$ for $|Z|<1$, is real-analytic on $[0,1]$, satisfies
$e^{F(t)}=e^{tX}e^{Y}$, and $F(0)=\log e^{Y}=Y$ because $|Y|\le 1/80$. At $t=1$ the definition
\eqref{4.27:eq-left-factor-algebra-1} reads $F(1)=i\eta A'^{+}(b)$, and
\begin{equation*}
e^{i\eta A'^{+}(b)}U(b)=e^{i\eta\mathfrak{h}(b)}e^{i\eta A'(b)}U(b)
=e^{i\eta\mathfrak{h}(b)}\mathbb{U}(b)=\mathbb{U}^{+}(b).
\end{equation*}
This is the asserted representation $\mathbb{U}^{+}=e^{i\eta A'^{+}}U$; the factor $U$ is the same in
both representations, so the $G$-valued part $U$, and every gauge of $U$, is unchanged.

Differentiating $e^{F(t)}=e^{tX}e^{Y}$ gives $\frac{d}{dt}e^{F}=Xe^{F}$. On the other hand
\begin{equation*}
\frac{d}{dt}e^{F}=\bigl(\varphi(\operatorname{ad}_F)F'\bigr)e^{F},
\qquad \varphi(z):=\frac{e^{z}-1}{z},
\end{equation*}
so $\varphi(\operatorname{ad}_F)F'=X$. Let
\begin{equation*}
\psi(z):=\frac{z}{e^{z}-1}=1-\frac{z}{2}+\frac{z^{2}}{12}-\frac{z^{4}}{720}+\cdots ,
\end{equation*}
holomorphic for $|z|<2\pi$, with $\psi\varphi=1$; as long as $\|\operatorname{ad}_F\|<2\pi$ this gives
$\psi(\operatorname{ad}_F)\varphi(\operatorname{ad}_F)=1$, hence $F'=\psi(\operatorname{ad}_F)X$.
Since $\psi(z)+z/2=(z/2)\coth(z/2)$ is even, apart from the term $-z/2$ only even powers of $z$
occur in $\psi$, and their coefficients are bounded in absolute value by $1/12$. Put
$\zeta:=\|\operatorname{ad}_F\|\le 2|F|$. For $\zeta\le 1/4$,
\begin{equation}\label{4.27:eq-left-factor-algebra-2}
\|\psi(\operatorname{ad}_F)-1\|\le \frac{\zeta}{2}+\sum_{n\ge1}\frac{\zeta^{2n}}{12}
=\frac{\zeta}{2}+\frac{\zeta^{2}}{12(1-\zeta^{2})}\le\frac{\zeta}{2}+\frac{\zeta^{2}}{11},
\end{equation}
since $12(1-\zeta^{2})\ge 12\cdot 15/16>11$.

Suppose that $|F(s)|\le 0.06$ for all $s\in[0,t]$. Then $\zeta\le 0.12\le 1/4$ on $[0,t]$, and
\eqref{4.27:eq-left-factor-algebra-2} gives
\begin{equation*}
\|\psi(\operatorname{ad}_F)-1\|\le\zeta\Bigl(\frac12+\frac{0.12}{11}\Bigr)\le 0.511\,\zeta
\le 1.03\,|F| ;
\end{equation*}
hence $|F'(s)|\le(1+1.03\cdot 0.06)|X|\le 1.07\,|X|$ and
\begin{equation*}
|F(s)|\le |Y|+\int_0^{s}|F'|\le \frac{1}{80}+\frac{1.07}{40}<0.04
\qquad(s\in[0,t]).
\end{equation*}
Let $t_*$ be the supremum of those $t\in[0,1]$ with $|F|\le 0.06$ on $[0,t]$. Since
$|F(0)|=|Y|\le 1/80$, continuity gives $t_*>0$ and $|F|\le 0.06$ on $[0,t_*]$, hence $|F|<0.04$
there; if $t_*<1$, continuity would give $|F|\le 0.06$ on a larger interval, contrary to the
definition of $t_*$. Thus $t_*=1$, and $|F|<0.04$ and
$\|\psi(\operatorname{ad}_F)-1\|\le 1.03\,|F|$ hold on all of $[0,1]$. Therefore
\begin{equation*}
|F(1)-Y-X|\le\int_0^{1}|F'(t)-X|\,dt
=\int_0^{1}\bigl|(\psi(\operatorname{ad}_{F(t)})-1)X\bigr|\,dt
\le 1.03\cdot 0.04\,|X|\le 0.042\,|X| .
\end{equation*}
Put $E:=A'^{+}-A'-\mathfrak{h}$ on the bonds of $\tilde{\Delta}(y)$. Then
$i\eta E(b)=F(1)-Y-X$, so
\begin{equation}\label{4.27:eq-left-factor-algebra-3}
\eta|E(b)|\le 0.042\,\eta|\mathfrak{h}(b)|\le 0.042\,\mathsf{n},
\end{equation}
and
\begin{equation*}
\eta|A'^{+}(b)-A'(b)|\le\eta|\mathfrak{h}(b)|+\eta|E(b)|\le 1.042\,\mathsf{n}\le 1.06\,\mathsf{n}.
\end{equation*}
This is (ii).

\emph{Proof of} (iii). Fix $\mu,\nu$ and put $b':=\langle y,y+\mu\rangle$. We have
$A'^{+}-A'=\mathfrak{h}+E$. Since $U(b')$ is unitary, $|\operatorname{Ad}_{U(b')}C|=|C|$, and
\eqref{4.27:eq-left-factor-algebra-3} applied to the bonds carrying $E_\nu(y+\mu)$ and $E_\nu(y)$
gives
\begin{equation*}
\eta^{2}|\nabla_{U,\mu}E_\nu(y)|
=\eta\bigl|\operatorname{Ad}_{U(b')}E_\nu(y+\mu)-E_\nu(y)\bigr|
\le\eta|E_\nu(y+\mu)|+\eta|E_\nu(y)|\le 0.084\,\mathsf{n}.
\end{equation*}
Next, $\mathbb{U}(b')=e^{Y'}U(b')$ with $Y':=i\eta A'(b')$, so
$\operatorname{Ad}_{\mathbb{U}(b')}=\operatorname{Ad}_{e^{Y'}}\operatorname{Ad}_{U(b')}$, and, since
$\operatorname{Ad}_{e^{Y'}}=e^{\operatorname{ad}_{Y'}}$ with $\|\operatorname{ad}_{Y'}\|\le 2|Y'|\le 1/40$,
\begin{equation*}
|e^{Y'}Ce^{-Y'}-C|\le\bigl(e^{2|Y'|}-1\bigr)|C|\le 0.026\,|C| .
\end{equation*}
Consequently
\begin{equation*}
\eta^{2}\bigl|\nabla_{U,\mu}\mathfrak{h}_\nu(y)-\nabla_{\mathbb{U},\mu}\mathfrak{h}_\nu(y)\bigr|
=\eta\bigl|(\operatorname{Ad}_{e^{Y'}}-1)\operatorname{Ad}_{U(b')}\mathfrak{h}_\nu(y+\mu)\bigr|
\le 0.026\,\eta|\mathfrak{h}_\nu(y+\mu)|,
\end{equation*}
and therefore, adding the bound on $E$ obtained above,
\begin{align*}
\eta^{2}|\nabla_{U,\mu}\mathfrak{h}_\nu(y)|
&\le\eta^{2}|\nabla_{\mathbb{U},\mu}\mathfrak{h}_\nu(y)|+0.026\,\eta|\mathfrak{h}_\nu(y+\mu)|
\le 1.026\,\mathsf{n},\\
\eta^{2}|\nabla_{U,\mu}(A'^{+}-A')_\nu(y)|
&\le 0.084\,\mathsf{n}+1.026\,\mathsf{n}\le 1.11\,\mathsf{n}\le 1.2\,\mathsf{n}.
\end{align*}
This is (iii).

\emph{Proof of} (i). Let $p=(y;\mu,\nu)$ be a plaquette based at $y$, with bonds
$b_1=\langle y,y+\mu\rangle$, $b_2=\langle y+\mu,y+\mu+\nu\rangle$,
$b_3=\langle y+\nu,y+\mu+\nu\rangle$, $b_4=\langle y,y+\nu\rangle$, so that
$\partial\mathbb{U}(p)=\mathbb{U}_1\mathbb{U}_2\mathbb{U}_3^{-1}\mathbb{U}_4^{-1}$ with
$\mathbb{U}_j:=\mathbb{U}(b_j)$; write $\mathfrak{h}_j:=\mathfrak{h}(b_j)$ and
$W:=\partial\mathbb{U}(p)$. Using $\mathbb{U}_1\mathbb{U}_2\mathbb{U}_3^{-1}=W\mathbb{U}_4$,
\begin{align*}
\partial\mathbb{U}^{+}(p)
&=e^{i\eta\mathfrak{h}_1}\mathbb{U}_1e^{i\eta\mathfrak{h}_2}\mathbb{U}_2\,
\mathbb{U}_3^{-1}e^{-i\eta\mathfrak{h}_3}\,\mathbb{U}_4^{-1}e^{-i\eta\mathfrak{h}_4}\\
&=e^{i\eta\mathfrak{h}_1}e^{i\eta\operatorname{Ad}_{\mathbb{U}_1}\mathfrak{h}_2}\,
W\,e^{-i\eta\operatorname{Ad}_{\mathbb{U}_4}\mathfrak{h}_3}e^{-i\eta\mathfrak{h}_4}
=e^{A_1}e^{A_2}e^{A_3}e^{A_4}\,W,
\end{align*}
where
\begin{equation}\label{4.27:eq-left-factor-algebra-4}
A_1=i\eta\mathfrak{h}_1,\quad
A_2=i\eta\operatorname{Ad}_{\mathbb{U}_1}\mathfrak{h}_2,\quad
A_3=-i\eta\operatorname{Ad}_{W}\operatorname{Ad}_{\mathbb{U}_4}\mathfrak{h}_3,\quad
A_4=-i\eta\operatorname{Ad}_{W}\mathfrak{h}_4 .
\end{equation}
Since $\mathfrak{h}_1=\mathfrak{h}_\mu(y)$, $\mathfrak{h}_2=\mathfrak{h}_\nu(y+\mu)$,
$\mathfrak{h}_3=\mathfrak{h}_\mu(y+\nu)$, $\mathfrak{h}_4=\mathfrak{h}_\nu(y)$, the covariant
derivative above gives
\begin{equation*}
\eta\nabla_{\mathbb{U},\mu}\mathfrak{h}_\nu(y)-\eta\nabla_{\mathbb{U},\nu}\mathfrak{h}_\mu(y)
=\mathfrak{h}_1+\operatorname{Ad}_{\mathbb{U}_1}\mathfrak{h}_2
-\operatorname{Ad}_{\mathbb{U}_4}\mathfrak{h}_3-\mathfrak{h}_4,
\end{equation*}
and hence
\begin{equation}\label{4.27:eq-left-factor-algebra-5}
S:=\sum_{j=1}^{4}A_j
=i\eta^{2}\bigl[(\nabla_{\mathbb{U},\mu}\mathfrak{h}_\nu)-(\nabla_{\mathbb{U},\nu}\mathfrak{h}_\mu)\bigr](y)
+i\eta(1-\operatorname{Ad}_W)\bigl(\operatorname{Ad}_{\mathbb{U}_4}\mathfrak{h}_3+\mathfrak{h}_4\bigr).
\end{equation}

We bound the adjoint actions. For a bond $b$ of $\tilde{\Delta}(y)$,
$\operatorname{Ad}_{\mathbb{U}(b)}=e^{\operatorname{ad}_{i\eta A'(b)}}\operatorname{Ad}_{U(b)}$, so
$\|\operatorname{Ad}_{\mathbb{U}(b)}\|\le e^{2\mathsf{a}}\le 1.03$. Since $|W-1|\le\theta$, the Neumann
series gives $|W^{-1}|\le(1-\theta)^{-1}$ and $|W|\le 1+\theta$; from
$WCW^{-1}-C=\bigl((W-1)C-C(W-1)\bigr)W^{-1}$,
\begin{equation*}
|(1-\operatorname{Ad}_W)C|\le\frac{2\theta}{1-\theta}|C|\le\frac{2}{9}|C|,
\qquad
\|\operatorname{Ad}_W\|\le\frac{1+\theta}{1-\theta}\le\frac{11}{9}.
\end{equation*}
By \eqref{4.27:eq-left-factor-algebra-5} and the definition of $\mathsf{n}$,
\begin{equation*}
|S|\le 2\mathsf{n}+\frac{2}{9}\,(1.03+1)\,\mathsf{n}\le 2.46\,\mathsf{n},
\end{equation*}
and by \eqref{4.27:eq-left-factor-algebra-4} each $|A_j|$ is at most
$a:=\tfrac{11}{9}\cdot 1.03\,\mathsf{n}<1.26\,\mathsf{n}$, with $4a\le 0.126\le 1/7$. Then
$|e^{A_1}\cdots e^{A_4}-1|\le e^{4a}-1<1$, so $\Xi:=\log(e^{A_1}e^{A_2}e^{A_3}e^{A_4})$ (principal
branch) satisfies $\partial\mathbb{U}^{+}(p)=e^{\Xi}W$.

We bound $\Xi-S$. Put $s:=4a\le 1/7$, $P:=e^{A_1}e^{A_2}e^{A_3}e^{A_4}$ and $Q:=P-1$.
Multiplying out the four exponential series, $P$ is the sum over $(n_1,\dots,n_4)$ of
$A_1^{n_1}A_2^{n_2}A_3^{n_3}A_4^{n_4}/(n_1!\,n_2!\,n_3!\,n_4!)$; by the multinomial identity the
norms of the terms of total degree $n$ add up to at most $s^{n}/n!$. The terms of degree one add up
to $S$, and those of degree two to
\begin{equation*}
P_2:=\frac12\sum_{j}A_j^{2}+\sum_{i<j}A_iA_j=\frac12S^{2}+\frac12\sum_{i<j}[A_i,A_j],
\end{equation*}
as one sees by writing $S^{2}=\sum_jA_j^{2}+\sum_{i<j}(A_iA_j+A_jA_i)$. Let $P_{\ge3}$ be the sum of
the terms of degree at least three, so that $Q=S+P_2+P_{\ge3}$. Using $2/(k+2)!\le 3^{-k}$,
$6/(k+3)!\le 4^{-k}$ for $k\ge0$, and $s\le 1/7$,
\begin{align*}
|Q-S|&\le\sum_{n\ge2}\frac{s^{n}}{n!}\le\frac{s^{2}/2}{1-s/3}\le 0.525\,s^{2},\qquad
|P_{\ge3}|\le\sum_{n\ge3}\frac{s^{n}}{n!}\le\frac{s^{3}/6}{1-s/4}\le 0.173\,s^{3},\\
|Q|&\le s+0.525\,s^{2}\le 1.075\,s\le 0.154 .
\end{align*}
Since $|Q|<1$, $\Xi=\sum_{n\ge1}(-1)^{n+1}Q^{n}/n$. Writing $Q^{2}=S^{2}+S(Q-S)+(Q-S)Q$,
\begin{equation*}
\Xi-S=\frac12\sum_{i<j}[A_i,A_j]+P_{\ge3}-\frac12\bigl(S(Q-S)+(Q-S)Q\bigr)
+\sum_{n\ge3}\frac{(-1)^{n+1}Q^{n}}{n}.
\end{equation*}
Here $|[A_i,A_j]|\le 2a^{2}$ for each of the six pairs, so the first term is at most $6a^{2}$; the
third is at most $\frac12\cdot 0.525\cdot(1+1.075)\,s^{3}\le 0.545\,s^{3}$, since $|S|\le s$; and the
last is at most $|Q|^{3}/(3(1-|Q|))\le 1.075^{3}s^{3}/(3\cdot 0.846)\le 0.49\,s^{3}$. With
$s^{3}=64a^{3}$ and $(0.173+0.545+0.49)\cdot 64<80$, and since $a<1.26/40$,
\begin{equation*}
|\Xi-S|\le 6a^{2}+80a^{3}\le 8.6\,a^{2}\le 13.7\,\mathsf{n}^{2}\le 0.35\,\mathsf{n}.
\end{equation*}
Hence
\begin{align*}
|\Xi|&\le 2.81\,\mathsf{n}\le 0.071,\qquad |e^{\Xi}-1|\le|\Xi|e^{|\Xi|}\le 3.02\,\mathsf{n},\\
|\partial\mathbb{U}^{+}(p)-\partial\mathbb{U}(p)|&=|(e^{\Xi}-1)W|\le 3.02\,\mathsf{n}\,|W|
\le 3.33\,\mathsf{n}\le 4\,\mathsf{n}.
\end{align*}
This is (i).
\end{proof}

\begin{lemma}[Coupling ratios, distances and absorption at level zero]\label{4.27:lem-level-zero-ratios}
We use the notation of Notation~\ref{4.27:not-collar-conventions} and of
Definition~\ref{4.27:def-dated-allowances}. In particular:
\begin{itemize}
\item $\mathsf{t}_m=\log g_m^{-2}$, and $\alpha_{\cdot,m}=\min\{\alpha_{0,m},\alpha_{1,m}\}$;
\item $r_a(m)$ are the level-$m$ radii of the entry $a$, and $\mathsf{A}$ is the maximal radius
function of $\mathsf{t}$;
\item $d^{\wedge}$ is the weighted contour distance, and $\mathsf{w}_M$ is the price a contour pays
per crossed stratum;
\item $\operatorname{St}(d)=\Omega_0(\mathfrak{B}_d)\setminus\Omega_1(\mathfrak{B}_d)$ is the
never-averaged collar of the set $\mathfrak{B}_d$ of a date $d$;
\item $\rho^{(l),(k)}_0$ is the allowance of the date $(k;l)$;
\item $\beta_0$ is the constant of the coupling comparisons \cite[(2.6),(2.7)]{B-CMP119};
\item $c_\beta=2\log(1+\beta_0)$, so that $e^{c_\beta/2}=1+\beta_0$;
\item $p_1$ and $C_*$ are the exponent and the constant of the stage-ratio statement whose content
is hypothesis (h4) below; in this lemma they enter only through (h1) and (h4).
\end{itemize}
The constants $C_1$ and $\delta$, the size $\delta'_j$ of the stage-$j$ datum, its support
$\mathcal{C}_j$ and the closure factor $\vartheta$ are those of the absorption inequality, which
states that, for the stage $l$ of the density $k$, of level $j=h_k+l-1$, and for every cell $y$
of $\mathfrak{B}_{(k;l-1)}$ of level $m(y)$,
\begin{equation*}
C_1\delta'_j\,e^{-\frac14\delta d^{\wedge}(y,\mathcal{C}_j)}
\le\vartheta\,\rho^{(l),(k)}_{m(y)}\,\alpha_{\cdot,m(y)} .
\end{equation*}
Throughout (h4), (h6) and part (c), $k\ge1$ is one fixed density index; the numbers
$\lambda_\vartheta$ and $\vartheta$ defined in (h4) and (h6) depend on it.
Assume the following.
\begin{itemize}
\item[(h0)] \emph{Numerical conditions.} $L\ge8$ and $(1+\beta_0)^2\le 64/49$.
\item[(h1)] \emph{The coupling comparisons of} \cite[(2.6),(2.7)]{B-CMP119} \emph{at the pairs
$(n,0)$.} For every $n\ge 1$ and for $p=q$ and $p=p_1$,
\begin{equation}\label{4.27:eq-level-zero-ratios-1}
g_n\le(1+\beta_0)\,n^{1/2}g_0,\qquad g_0\le(1+\beta_0)\,g_n,\qquad
(\log g_n^{-2})^{p}\le(1+\beta_0)(\log g_0^{-2})^{p}.
\end{equation}
Assume moreover that for every $k\ge1$
\begin{equation*}
\log g_0^{-2}\ge\tfrac12\log g_k^{-2}.
\end{equation*}
\item[(h2)] \emph{The coupling form of the radii.} For $i\in\{0,1\}$ and $\lambda\ge1$,
\begin{equation}\label{4.27:eq-level-zero-ratios-2}
\frac{\alpha_{i,\lambda}}{\alpha_{i,0}}
=\frac{g_\lambda}{g_0}\Bigl(\frac{\mathsf{t}_\lambda}{\mathsf{t}_0}\Bigr)^{q}.
\end{equation}
For every entry $a$ there is an exponent $e_a\ge1$ such that
\begin{equation}\label{4.27:eq-level-zero-ratios-3}
\frac{r_a(m)}{r_a(0)}=\frac{\alpha_{\cdot,m}}{\alpha_{\cdot,0}}\,L^{-e_am}
\qquad\text{for every level } m.
\end{equation}
\item[(h3)] \emph{Properties of the maximal radius function.} The function $\mathsf{A}$ is
decreasing on $\mathsf{t}>2q$. For every $k\ge1$,
\begin{equation*}
\mathsf{A}(\mathsf{t}_k-c_\beta)\le e^{c_\beta/2}\mathsf{A}(\mathsf{t}_k).
\end{equation*}
Moreover $\mathsf{t}_0$ and $\mathsf{t}_k-c_\beta$ lie in the range
$\mathsf{t}>2q$.
\item[(h4)] \emph{The stage-ratio bound at the pair $(j,0)$.} Put
$\lambda_\vartheta:=(A_1/C_*)(\tfrac12\mathsf{t}_k)^{p_1-q}$. For every stage of the density $k$
whose level $j$ satisfies $j\ge1$ the following holds: if
\begin{equation*}
g_j\le(1+\beta_0)j^{1/2}g_0,\qquad
(\log g_j^{-2})^{p_1}\le(1+\beta_0)(\log g_0^{-2})^{p_1},\qquad
\log g_0^{-2}\ge\tfrac12\log g_k^{-2},
\end{equation*}
then
\begin{equation}\label{4.27:eq-level-zero-ratios-4}
\frac{\delta'_j}{\alpha_{\cdot,0}}\le 3(1+j)\,\lambda_\vartheta .
\end{equation}
\item[(h5)] \emph{The contour price and the layering of the set of a date.}
\begin{itemize}
\item The set $\mathfrak{B}$ is divided into strata $S_0,S_1,\dots$, the stratum $S_i$ consisting
of the cells of $\mathfrak{B}$-level $i$; the regions $\Omega_0\supset\Omega_1\supset\Omega_2
\supset\cdots$ carrying the successive levels are nested. A contour pays at least $\mathsf{w}_M$
per crossed stratum of $\mathfrak{B}$ or per crossed $M$-cube of the local scale of $\mathfrak{B}$.
\item For the set of a date, the levels are consecutive across its strata, and each stratum is at
least one $M$-cube of its scale wide.
\end{itemize}
\item[(h6)] \emph{The data of the absorption inequality.} Its closure factor is
$\vartheta=K_\vartheta\lambda_\vartheta$, with $K_\vartheta\ge 24C_1/\beta$; the support
$\mathcal{C}_j$ of the stage-$j$ datum is a set of bonds of $\mathfrak{B}$-level $j$; and the
constants $\delta$, $\mathsf{w}_M$ satisfy
\begin{equation*}
\tfrac14\delta\mathsf{w}_M\ge 4\ln L+1,
\end{equation*}
so that $\tfrac14\delta\mathsf{w}_M>9$ by (h0).
\end{itemize}
Then the following hold.
\begin{enumerate}
\item[(a)] For $\lambda\ge1$ and $i\in\{0,1\}$,
\begin{equation*}
\alpha_{i,\lambda}\le(1+\beta_0)^2\lambda^{1/2}\alpha_{i,0},
\end{equation*}
and therefore
\begin{equation*}
r_a(\lambda)\le(1+\beta_0)^2\lambda^{1/2}L^{-e_a\lambda}r_a(0)\le 0.164\,r_a(0).
\end{equation*}
Moreover $\mathsf{A}(\mathsf{t}_0)\le(1+\beta_0)\mathsf{A}(\mathsf{t}_k)$ for every $k\ge1$.
\item[(b)] In every set $\mathfrak{B}_d$ of a date $d$, let $y\in\operatorname{St}(d)$ and let $y'$
be a cell of level $\lambda\ge1$. Then $d^{\wedge}(y,y')\ge\mathsf{w}_M(\lambda-1)$.
\item[(c)] Consider a stage of level $j=h_k+l-1\ge1$, put
$\mathfrak{B}:=\mathfrak{B}_{(k;l-1)}$, and let $y\in\operatorname{St}((k;l-1))$. Then
\begin{equation}\label{4.27:eq-level-zero-ratios-5}
C_1\delta'_j\,e^{-\frac14\delta d^{\wedge}(y,\mathcal{C}_j)}
\le\vartheta\,\rho^{(l),(k)}_0\,\alpha_{\cdot,0},
\end{equation}
with the same closure factor $\vartheta$ as in the absorption inequality.
\end{enumerate}
\end{lemma}

\begin{proof}
\emph{(a), the ratios of the radii.} Fix $i\in\{0,1\}$ and $\lambda\ge1$. By
\eqref{4.27:eq-level-zero-ratios-2},
\begin{equation*}
\frac{\alpha_{i,\lambda}}{\alpha_{i,0}}=\frac{g_\lambda}{g_0}\Bigl(\frac{\mathsf{t}_\lambda}{\mathsf{t}_0}\Bigr)^{q}.
\end{equation*}
The first inequality of \eqref{4.27:eq-level-zero-ratios-1} at $n=\lambda$ gives
$g_\lambda/g_0\le(1+\beta_0)\lambda^{1/2}$. Its third inequality at $n=\lambda$ with $p=q$ reads
$\mathsf{t}_\lambda^{\,q}\le(1+\beta_0)\mathsf{t}_0^{\,q}$. Since $\mathsf{t}_0>2q>0$ by (h3), this
is the bound
\begin{equation*}
(\mathsf{t}_\lambda/\mathsf{t}_0)^{q}\le 1+\beta_0 .
\end{equation*}
Multiplying the two bounds gives $\alpha_{i,\lambda}\le(1+\beta_0)^2\lambda^{1/2}\alpha_{i,0}$.

This holds for $i=0$ and for $i=1$; and if each of two positive numbers is at most a fixed
multiple of its counterpart among two others, then the minimum of the first two is at most the
same multiple of the minimum of the other two.
Hence $\alpha_{\cdot,\lambda}\le(1+\beta_0)^2\lambda^{1/2}\alpha_{\cdot,0}$.
Inserting this into \eqref{4.27:eq-level-zero-ratios-3} at $m=\lambda$ gives
\begin{equation*}
r_a(\lambda)\le(1+\beta_0)^2\lambda^{1/2}L^{-e_a\lambda}r_a(0).
\end{equation*}

\emph{(a), the numerical bound.} Since $e_a\ge1$ and $L>1$, we have $L^{-e_a\lambda}\le L^{-\lambda}$.
The function $\lambda\mapsto\lambda^{1/2}L^{-\lambda}$ is decreasing on $\lambda\ge1$ for $L\ge8$,
which holds by (h0): its logarithmic derivative is $1/(2\lambda)-\ln L<0$ there. Its maximum on
$\lambda\ge1$ is therefore $L^{-1}$. With $(1+\beta_0)^2\le 64/49$ and $L\ge8$, both from (h0),
this gives
\begin{equation*}
(1+\beta_0)^2\lambda^{1/2}L^{-e_a\lambda}\le\frac{64/49}{8}<0.164,
\end{equation*}
which is the second inequality of (a).

\emph{(a), the last clause.} The second inequality of \eqref{4.27:eq-level-zero-ratios-1} at
$n=k$ reads $g_0\le(1+\beta_0)g_k$. Squaring, inverting and taking logarithms gives
\begin{equation*}
\mathsf{t}_0\ge\mathsf{t}_k-2\log(1+\beta_0)=\mathsf{t}_k-c_\beta ,
\end{equation*}
by the value of $c_\beta$ recalled above. By (h3), both $\mathsf{t}_0$ and $\mathsf{t}_k-c_\beta$ lie
in the range $\mathsf{t}>2q$, on which $\mathsf{A}$ is decreasing. Hence
\begin{equation*}
\mathsf{A}(\mathsf{t}_0)\le\mathsf{A}(\mathsf{t}_k-c_\beta)\le e^{c_\beta/2}\mathsf{A}(\mathsf{t}_k)
=(1+\beta_0)\mathsf{A}(\mathsf{t}_k),
\end{equation*}
where the middle inequality is the second property in (h3).

\emph{(b).} We derive the bound from the two clauses of (h5), for a cell of level $0$ and a cell
of level $\lambda$.
The level-$0$ cells of $\mathfrak{B}_d$ are the cells of the stratum $S_0$, that is, the cells of
$\Omega_0\setminus\Omega_1=\operatorname{St}(d)$. This is a stratum of $\mathfrak{B}_d$ in
the distance of Ba{\l}aban's construction, whose defining sum runs over all levels of
$\mathfrak{B}_d$. So $y$ is a cell of level $0$ and $y'$ a cell of level $\lambda$.

By the second clause of (h5), the levels of $\mathfrak{B}_d$ are consecutive across its
strata, and each stratum is at least one $M$-cube of its scale wide. By the first clause of (h5),
the regions $\Omega_1\supset\Omega_2\supset\cdots$ of the successive levels of $\mathfrak{B}_d$ are
nested. A part of a contour lying outside $\Omega_0$
reaches no level $\ge1$ without crossing stratum $1$. Therefore every contour from $y$ to $y'$
crosses the strata $1,\dots,\lambda-1$, whichever way it runs. By the price clause of (h5) each
crossing costs at least $\mathsf{w}_M$, and so
\begin{equation*}
d^{\wedge}(y,y')\ge\mathsf{w}_M(\lambda-1)^{+}=\mathsf{w}_M(\lambda-1).
\end{equation*}

\emph{(c), the three inputs.} We follow the proof of the absorption inequality at a point $y$ of
level $m(y)=0$. Since $y\in\operatorname{St}((k;l-1))$ is a cell of the stratum $S_0$, its level
is indeed $0$, and its distance $|0-j|$ to the level of the stage is $j\ge1$. We need three inputs.

First, we apply (h4) to this stage, which is a stage of the density $k$ of level $j\ge1$, at the
pair $(j',m)=(j,0)$, and verify its three premises.
The first two are the first and third inequalities of \eqref{4.27:eq-level-zero-ratios-1} at
$n=j$, the latter with $p=p_1$. The third, $\log g_0^{-2}\ge\frac12\log g_k^{-2}$, is the
inequality assumed last in (h1), at the present $k$. Hence
\eqref{4.27:eq-level-zero-ratios-4} holds:
\begin{equation*}
\delta'_j/\alpha_{\cdot,0}\le 3(1+j)\lambda_\vartheta .
\end{equation*}

Second, by (h6), $\mathcal{C}_j$ consists of bonds of level $j$. Part (b) with $\lambda=j\ge1$,
applied to each cell of $\mathcal{C}_j$, gives $d^{\wedge}(y,\mathcal{C}_j)\ge\mathsf{w}_M(j-1)$.

Third, by Definition~\ref{4.27:def-dated-allowances},
\begin{equation*}
\rho^{(l),(k)}_0=\beta2^{-(h_k+l)}=\beta2^{-(j+1)}.
\end{equation*}

\emph{(c), combining the inputs.} The three inputs together give
\begin{equation*}
\frac{C_1\delta'_j\,e^{-\frac14\delta d^{\wedge}(y,\mathcal{C}_j)}}{\rho^{(l),(k)}_0\,\alpha_{\cdot,0}}
\le\frac{C_1\lambda_\vartheta}{\beta}\cdot 3(1+j)\,2^{j+1}\,e^{-\frac14\delta\mathsf{w}_M(j-1)} .
\end{equation*}
We bound the factor $3(1+j)2^{j+1}e^{-\frac14\delta\mathsf{w}_M(j-1)}$ by $24$.
\begin{itemize}
\item If $j=1$, this factor equals $3\cdot2\cdot4=24$.
\item If $j\ge2$, then $\frac14\delta\mathsf{w}_M\ge4\ln L+1>9$ by (h6). The factor is therefore at
most $f(j):=3(1+j)2^{j+1}e^{-9(j-1)}$. Here $f(2)=72e^{-9}<1$, and
$f(j+1)/f(j)=2(j+2)(j+1)^{-1}e^{-9}<1$. Hence $f(j)<1$ for all $j\ge2$.
\end{itemize}
In both cases the factor is at most $24$. Consequently
\begin{equation*}
\frac{C_1\delta'_j\,e^{-\frac14\delta d^{\wedge}(y,\mathcal{C}_j)}}{\rho^{(l),(k)}_0\,\alpha_{\cdot,0}}
\le\frac{24C_1\lambda_\vartheta}{\beta}\le K_\vartheta\lambda_\vartheta=\vartheta,
\end{equation*}
by the bound $K_\vartheta\ge24C_1/\beta$ of (h6). This is \eqref{4.27:eq-level-zero-ratios-5}. The
closure factor $\vartheta=K_\vartheta\lambda_\vartheta$ is the one of the absorption inequality.
\end{proof}

\begin{lemma}[Displacement of the collar quantities by one operation]
\label{4.27:lem-collar-displacement}
We use the conventions of Notation~\ref{4.27:not-collar-conventions}, and the scaffold dates, the
allowances $\mathfrak{c}(d)$ and the stratum factor $\mathsf{f}_0(d)$ of
Definition~\ref{4.27:def-dated-allowances}. Let $\mathfrak{O}$ be an operation with target date
$d^t$ and source date $d'$, two consecutive scaffold dates. Put
$\mathfrak{B}:=\mathfrak{B}_{d^t}$, and let $k$ be the density in the sense of
Definition~\ref{4.27:def-dated-allowances}. Let $\tilde p=(\mathbb{U},\mathbb{J})$ be a pair on
$\Omega_0(\mathfrak{B}_{d'})$ which lies in $\mathfrak{K}_{109}(\mathfrak{B};\gamma_0^{(k)})$ and
obeys \eqref{4.27:eq-dated-allowances-a} at the date $d'$. For a map $\Phi$ of pairs, a point $y$
and an entry $a$ put
\begin{equation}\label{4.27:eq-collar-displacement-1}
\Delta_a(y):=\bigl|Q_a(\Phi\tilde p)(y)-Q_a(\tilde p)(y)\bigr| .
\end{equation}
When the source date $d'=(k;l)$ is a stage we put $j:=h_k+l-1$ and
$\rho_0:=\rho^{(l),(k)}_0$. For a datum $b$ supported on $\mathcal{C}_j$ write
$p(y):=\sup|b|\,e^{-\delta d^{\wedge}(y,\mathcal{C}_j)}$, and let $\ell(y):=L^{j(y)}\eta$ be the
length attached to the cell $\Delta(y)$, $y\in\Lambda_{j(y)}$, of the determining set; $j(y)$ is
the level of that cell and is distinct from $j$.

The maps of pairs in part (S) are the following. Let $\mathbb{H}$ be the response field of the
response-field theorem, let $\{\zeta_{\square}\}$ be the partition of unity subordinate to the
cover by cubes $\square$, of which every cut-off $\chi$ below is a linear combination,
and, for a cut-off $\chi$, let $\mathcal{J}^{\sharp}[\chi]$ be the current increment of the
current-closure statement; put
\begin{equation*}
\Phi^{\chi}(\mathbb{U},\mathbb{J}):=\bigl(e^{i\eta\chi\mathbb{H}}\mathbb{U},\;
\mathbb{J}+\mathcal{J}^{\sharp}[\chi]\bigr).
\end{equation*}
The current-closure statement has the families (C1), (C2)($\alpha$) and (C2)($\beta$) of such
maps. In (C2)($\alpha$), $\chi=\sum_{\square}t_{\square}\zeta_{\square}$ with $|t_{\square}|\le2$,
the datum satisfies $\|b\|_{\infty}\le rC_1\delta'_j$, and $1\le r\le r_*^{\sharp}$; in
(C2)($\beta$), $r\le1.1$ and one further amplitude $\tau$ with $|\tau|\le r_{\square}^{\sharp}$ is
present. The simultaneous-amplitude statement has a family (a) of cut-offs $\chi_t$, whose
amplitudes $t_{\square}$ satisfy $|t_{\square}|\le R^{\natural}_{\square}$ for all cubes $\square$
at once; for such $\chi_t$ put
\begin{equation*}
\Phi_t(\mathbb{U},\mathbb{J}):=\bigl(e^{i\eta\chi_t\mathbb{H}}\mathbb{U},\;
\mathbb{J}+\chi_t\mathcal{J}^{\sharp}[1]\bigr).
\end{equation*}
The number $R_0$ is the bound on the further amplitude in clause (vi) of hypothesis (c) below;
in that clause $\tilde{\square}$ is the enlarged cube of the cube $\square$ and
$d(y,\tilde{\square})$ the distance from $y$ to it, and $d^{-}$ is the distance through which that
clause lets the admissible size of $\tau$ grow. The function $\bar q(y)$ is the cube profile of the
simultaneous-amplitude statement formed with the amplitudes $t_{\square}$ of $\chi_t$; it plays for
the family (a) the role that $q(y)$ plays in clause (v) of hypothesis (c).

Assume the statements (a)--(f$'$) below, each at every point $y$ and at level $0$, together with
(g)--(i).
\begin{enumerate}
\item[(a)] Clause (b) of the response-field theorem: for every profile in the class
$\mathfrak{M}_{\delta}$ and all data $\mathsf{d}_1,\mathsf{d}_2$ admissible for that theorem, in
particular obeying its smallness condition $\|\mathsf{d}_i\|_1\le\varepsilon_{\sharp}$, $i=1,2$,
with $\varepsilon_{\sharp}$ the smallness radius of that theorem, the differences of the response
fields $\mathbb{H}_1-\mathbb{H}_2$ and of the associated fields of that theorem are bounded, in the
norms weighted by that profile, by $B_{\sharp}$ times the norm of $\mathsf{d}_1-\mathsf{d}_2$
weighted by the same profile, for $|\sigma|\le e^{\kappa_1}$ and at any determining set.
\item[(b)] The linear bounds of the stage, which bound $N^{\Delta}_y(A')$, $N^{\Delta}_y(\mathbb{H})$
and the associated quantities of the response-field theorem by $B_{\sharp}\,p(y)$.
\item[(c)] Clauses (iv)--(vi) of the recursive bounds of the stage, with the profiles $a(y)$,
$q(y)$ and the cube profiles $q_{\square}(y)$ of those bounds: (iv) under the hypotheses of that
clause, among which are
$\mathsf{p}(\mathbb{U})\le1$ at the cells of the stratum and $2a(y)B_{\sharp}p(y)\le 1/10$,
\begin{equation}\label{4.27:eq-collar-displacement-2}
\ell(y)^3\,\bigl|\mathcal{J}^{\sharp}[\chi]\bigr|\le K_{\chi}\bigl[a(y)p(y)+q(y)\bigr];
\end{equation}
(v) if $\chi=\sum_{\square}c_{\square}\zeta_{\square}$, then
$q:=10\sum_{\square}|c_{\square}|q_{\square}$ lies in $\mathfrak{M}_{2\delta}$ and $ap\le q$, and
$|c_{\square}|\le2$ for every $\square$ implies $q\le 20\,\mathsf{c}_{\square}\,p$; (vi) if, for a
number $\varsigma$, $|\tau|\le R_0e^{\varsigma\delta d^{-}}$, then at every point $y$
\begin{equation}\label{4.27:eq-collar-displacement-3}
10|\tau|\,q_{\square}(y)\le 10R_0\sup|b|\;
e^{-(1-\varsigma)\delta d^{\wedge}(y,\mathcal{C}_j)}\,e^{-(1+\varsigma)\delta d(y,\tilde{\square})}.
\end{equation}
\item[(d)] The closure conditions of the current-closure statement:
\begin{equation}\label{4.27:eq-collar-displacement-4}
r\vartheta\le\frac{\beta\,\mathsf{c}_{\vartheta}}{32\,\mathsf{K}_r},\qquad
R_0\vartheta\le\frac{1}{2\mathsf{K}_r},\qquad
\mathsf{K}_r:=2^9\mathsf{c}_{\square}\mathsf{c}_{\vartheta}\max\{K_{\chi},K_S^bB_{\sharp}\},
\end{equation}
with $\mathsf{c}_{\square},\mathsf{c}_{\vartheta}\ge1$ and $K_{\chi}\ge 15dB_{\sharp}=60B_{\sharp}$,
where $\vartheta$ is the factor of Lemma~\ref{4.27:lem-level-zero-ratios}(c).
\item[(e)] The response-size corollary: for the maps $\Phi^{\chi}$ of the families (C1),
(C2)($\alpha$) and (C2)($\beta$),
\begin{equation}\label{4.27:eq-collar-displacement-5}
N_y(\chi\mathbb{H})\le 2a(y)N_y(\mathbb{H})\le 2B_{\sharp}q(y)\le\varkappa_0(y)
:=\hat b\,e^{-\frac12\delta d^{\wedge}(y,\mathcal{C}_j)},
\end{equation}
where $\hat b:=2B_{\sharp}(20\mathsf{c}_{\square}r+11R_0)C_1\delta'_j$.
\item[(f)] The simultaneous-amplitude statement: for the members $\chi_t$ of its family (a), and
at every location $y$,
\begin{equation}\label{4.27:eq-collar-displacement-6}
N_y(\chi_t\mathbb{H})\le\varkappa_{\mathsf{w}}(y)
:=22B_{\sharp}\mathsf{c}_{\square}\mathsf{w}_0R_0C_1\delta'_j\,
e^{-\frac12\delta d^{\wedge}(y,\mathcal{C}_j)},\qquad \mathsf{w}_0=\tfrac1{24}.
\end{equation}
\item[(f$'$)] The level-zero bound on the current component for the family (a): for its members
$\chi_t$ and at every location $y$,
\begin{equation}\label{4.27:eq-collar-displacement-14}
\eta^3\bigl|\chi_t\mathcal{J}^{\sharp}[1]\bigr|\le a(y)\cdot2K_{\chi}\,p(y)\le2K_{\chi}\,\bar q(y)
\le\frac{11}{2^9}\,\mathsf{w}_0\rho_0\alpha_{0,0}\,e^{-\frac14\delta d^{\wedge}(y,\mathcal{C}_j)}.
\end{equation}
\item[(g)] Lemma~\ref{4.27:lem-level-zero-ratios}.
\item[(h)] The ceiling $120\,\bar{\gamma}_0(\mathsf{t}_k)\le1$ on the $\gamma$-ceiling function.
\item[(i)] The value $6.6$ lies in the range $1\le r\le r_*^{\sharp}$.
\end{enumerate}
Then part (S) below holds. Parts (T), (L$^{\mathrm{for}}$) and (L$^{\mathrm{Col}}$) hold under the
further hypotheses of the lemma on the displacements by the $\mathbf{T}$-step and at the completion
site, in which they are proved.
\begin{itemize}
\item[(S)] Let $d'=(k;l)$ be a stage, with $j$ and $\rho_0$ as above, and let $\mathfrak{O}$ be
any of the operations with source date $d'$; the three such operations are treated alike. Let
$\Phi$ be a map $\Phi^{\chi}$ of one of the families (C1), (C2)($\alpha$) and (C2)($\beta$) above
(in (C2)($\alpha$) the value $r=6.6$ is included, by (i)), or a map $\Phi_t$ with $\chi_t$ in the
family (a) of the simultaneous-amplitude statement. Let
$|\sigma|\le e^{\kappa_1}$. Then for every $y\in\operatorname{St}(d')$ we have $\Delta_5(y)=0$, the
gauges on the cubes are unchanged, and
\begin{equation}\label{4.27:eq-collar-displacement-b}
\Delta_a(y)\le\tfrac14\,\rho_0\,e^{-\frac14\delta d^{\wedge}(y,\mathcal{C}_j)}\,r_a(0)
\le\tfrac14\,\mathfrak{c}(d')\,r_a(0),\qquad a\in\{1,3,6,7\}.
\end{equation}
For the maps $\Phi_t$ the sharper bound
\begin{equation}\label{4.27:eq-collar-displacement-7}
\Delta_a(y)\le\tfrac16\,\mathsf{w}_0\,\rho_0\,e^{-\frac14\delta d^{\wedge}(y,\mathcal{C}_j)}\,r_a(0),
\qquad a\in\{1,3,6,7\},
\end{equation}
holds; its right-hand side is the value $\iota_l(y)\,r_a(0)$ of the profile $\iota_l$ at a point
$y$ of cell level $j(y)=0$.
\item[(T)] Let $d'=(k'+1;\top)$ with $k'\ge1$, and let $\mathfrak{O}$ be the $\mathbf{T}$-step
from level $k'$ to level $k'+1$. For every datum $b$ on $\mathcal{C}\subset\Omega_{k'+1}$ with
$\sup|b|\le w_{k'}$ and every $y\in\operatorname{St}(d')$,
\begin{equation}\label{4.27:eq-collar-displacement-c}
\Delta_a(y)\le\Lambda_{\gamma}\cdot\tfrac14(1-\theta_S)\,\beta\,2^{-(k'+1)}\,r_a(0)
\le\tfrac18\,\mathfrak{c}(d')\,r_a(0),
\end{equation}
$\Delta_5(y)=0$, and the gauges on the cubes are unchanged.
\item[(L$^{\mathrm{for}}$)] Let $d'=(k;L)$, let $\mathfrak{O}$ be either of the two operations at
the completion site, and let $y\in\operatorname{St}(d')$, that is, $y$ lies in no completing
component. Then the
pair presented at $y$ has the form identity${}\times{}$layers, $\Delta_5(y)=0$, the gauges on the
cubes are unchanged, and
\begin{equation}\label{4.27:eq-collar-displacement-d}
\Delta_a(y)\le\bar{P}^{0}\,2^{-(k-1)}\,r_a(0)\le\frac{\theta_S}{4}\,\mathfrak{c}(d')\,r_a(0)
\end{equation}
under the condition
\begin{equation}\label{4.27:eq-collar-displacement-a}
16\,\bar{P}^{0}\le\beta\,\theta_S,\qquad
\bar{P}^{0}:=\bar{P}\cdot\max\Bigl\{1,\frac{K_W^{0}}{K_W}\Bigr\},\qquad
K_W^{0}:=\max\{4,\;3+18d+2.2K_J\}.
\end{equation}
Here $K_W^{0}$ is a number fixed before $\gamma$, with which the $\gamma$-ceiling $K_W$ of the
completion-site bound is re-maximised; the enlargement is void, $\bar{P}^{0}=\bar{P}$, when
$K_W\ge K_W^{0}$. The bound \eqref{4.27:eq-collar-displacement-d} is
linear in the layer: a dilation $\mathsf{w}$ of a layer multiplies it by $\mathsf{w}$.
\item[(L$^{\mathrm{Col}}$)] Let $d'=(k;L)$ and
$y\in\operatorname{St}(d^t)\setminus\operatorname{St}(d')=\bigcup_C\operatorname{Col}(C)$, the
union over the completing components $C$. For every
$\mathfrak{m}=(\zeta,\varpi,\vec{\mathsf{s}})\in\mathfrak{P}$ (the letter $\zeta$ here has no
relation to the partition of unity $\{\zeta_{\square}\}$), the pairs presented at
$\mathfrak{b}_2$,
$\mathfrak{b}_3$, $\mathfrak{b}_4$ and on the segments between them satisfy, under the level-zero
condition on the response field of the completion site,
\begin{equation}\label{4.27:eq-collar-displacement-e}
Q_a(y)\le\frac{\varphi_{\flat}}{16}\,\mathbf{r}_a(0)\le\frac1{16}\,r_a(0),
\qquad a\in\{1,3,5,6,7\}.
\end{equation}
\end{itemize}
\end{lemma}

We prove here the preliminary smallness bounds at the stratum, which are used in all four parts,
and part (S). Parts (T), (L$^{\mathrm{for}}$) and (L$^{\mathrm{Col}}$) are proved in the lemma on
the displacements by the $\mathbf{T}$-step and at the completion site.

\begin{proof}[Proof of the preliminary bounds and of part (S)]
\emph{Preliminary: smallness at the stratum.} Since $\tilde p$ lies in
$\mathfrak{K}_{109}(\mathfrak{B};\gamma_0^{(k)})$, we have $\mathbb{U}=e^{i\eta A'}U$ with $U$
$G$-valued; let $\theta$ and $\mathsf{a}$ be the quantities of
Lemma~\ref{4.27:lem-left-factor-algebra} formed with $(\mathbb{U},A')$, namely the maximum of
$|\partial\mathbb{U}-1|$ over the plaquettes based at $y$ and the supremum of $\eta|A'|$ over
$\tilde{\Delta}(y)$. The bound \eqref{4.27:eq-dated-allowances-a} at the date $d'$ bounds the
entries at $y$ by the stratum factor $\mathsf{f}_0(d')$ times the radii at level $0$, and
$\mathsf{f}_0(d')\le4$ by Lemma~\ref{4.27:lem-allowance-sum}. Hence, at the bonds of the stratum,
\begin{equation}\label{4.27:eq-collar-displacement-8}
\theta\le4\alpha_{0,0},\qquad Q_5\le4\alpha_{0,0},\qquad \mathsf{a}\le4\alpha_{1,0}.
\end{equation}
At any bonds and plaquettes of level $1$ inside the enlarged cell $\tilde{\Delta}(y)$, the
clauses defining $\Omega_1$ give $\eta|A'|<5\alpha_{1,1}L^{-1}$ and
$|\partial\mathbb{U}-1|<5\alpha_{0,1}L^{-2}$. By Lemma~\ref{4.27:lem-level-zero-ratios}(a) with
$\lambda=1$ we have $\alpha_{i,1}\le(1+\beta_0)^2\alpha_{i,0}$ for $i\in\{0,1\}$, so these bounds
are at most $5(1+\beta_0)^2L^{-1}\alpha_{1,0}$ and $5(1+\beta_0)^2L^{-2}\alpha_{0,0}$. By the bound
$r_a(1)\le0.164\,r_a(0)$ of Lemma~\ref{4.27:lem-level-zero-ratios}(a), these are at most
$5\cdot0.164<4$ times the corresponding level-$0$ radii, hence smaller than the bounds
\eqref{4.27:eq-collar-displacement-8}.

Next,
\begin{equation}\label{4.27:eq-collar-displacement-9}
4\mathsf{A}(\mathsf{t}_0)\le4(1+\beta_0)\mathsf{A}(\mathsf{t}_k)<5\bar{\alpha}_k
\le\gamma_0^{(k)}\le\bar{\gamma}_0(\mathsf{t}_k)\le\frac1{120}.
\end{equation}
The first inequality is the last assertion of Lemma~\ref{4.27:lem-level-zero-ratios}(a). The
second holds because $4(1+\beta_0)<5$, that is $\beta_0<\frac14$, and
$\mathsf{A}(\mathsf{t}_k)\le\bar{\alpha}_k$, the running
maximum $\bar{\alpha}_k$ being the maximum of $\mathsf{A}(\mathsf{t}_m)$ over $m\le k$. The third
holds by the definition $\gamma_0^{(k)}=5\max_{m\le k}w_*(m)\bar{\alpha}_m$ of the dated size
parameter, with $w_*(m)=\mathsf{L}_0^{2N_0(m)}$, taken at $m=k$, since the weight $w_*(k)\ge1$; the
fourth is its bound by the $\gamma$-ceiling function. The last is hypothesis (h). Since
$\mathsf{A}$ is the maximum of the
radius functions, $\alpha_{0,0}\le\mathsf{A}(\mathsf{t}_0)$ and
$\alpha_{1,0}\le\mathsf{A}(\mathsf{t}_0)$. By \eqref{4.27:eq-collar-displacement-8} and
\eqref{4.27:eq-collar-displacement-9},
\begin{equation*}
\theta\le4\alpha_{0,0}\le\frac1{120}\le\frac1{10},\qquad
\mathsf{a}\le4\alpha_{1,0}\le\frac1{120}\le\frac1{80}.
\end{equation*}
The same bounds give $\mathsf{p}(\mathbb{U})\le1$ at the cells of level $0$, where
$\mathsf{p}(\mathbb{U})$ is the regularity size of $\mathbb{U}$ entering clause (iv) of
hypothesis (c). Thus the hypotheses of
Lemma~\ref{4.27:lem-left-factor-algebra} on $(\mathbb{U},A')$ hold at every $y$ of the stratum,
and so does the hypothesis $\mathsf{p}(\mathbb{U})\le1$ of clause (iv) of hypothesis (c).

\emph{Part (S): the setting.} Let $d'=(k;l)$, put $n:=l-1$, so that $j=h_k+n$, and let
$\rho_0$ be as in the statement. The hypotheses under which (b) and (c) are stated hold:
$(\mathbb{U},\mathbb{J})$ lies in $\mathfrak{K}_{109}(\mathfrak{B};\gamma_0^{(k)})$, the class on
which the recursive bounds of the stage are stated, and the set is
$\mathfrak{B}=\mathbb{B}^{(n)}_k$, frozen on the live pieces. This is a determining set, and the
response-field theorem, together with the bounds (b) and (c), holds at any determining set.

\emph{Step 1: the profiles.} No level is read in this step. The two bounds
\eqref{4.27:eq-collar-displacement-5} and \eqref{4.27:eq-collar-displacement-6},
$N_y(\chi\mathbb{H})\le\varkappa_0(y)$ and $N_y(\chi_t\mathbb{H})\le\varkappa_{\mathsf{w}}(y)$, are
consequences of the linear bounds (b), of clauses (v) and (vi) of (c), and of the response-field
theorem (a), each taken at the point $y$. Each of these statements holds at every point $y$.
Their quantifier runs over the cells $\Delta(y)$, $y\in\Lambda_{j(y)}$, of $\mathfrak{B}$, for
every value $j(y)\ge0$ of the cell level. For $j(y)=0$ these are the cells of the stratum
$\Lambda_0=\Omega_0\setminus\Omega_1$, at
which $\ell(y)=\eta$. Hence both bounds hold at every $y$ of the stratum.

\emph{Step 2: the one line that depends on the level.} By
Lemma~\ref{4.27:lem-level-zero-ratios}(c),
\begin{equation*}
C_1\delta'_je^{-\frac14\delta d^{\wedge}(y,\mathcal{C}_j)}\le\vartheta\rho_0\alpha_{\cdot,0}.
\end{equation*}
Writing $e^{-\frac12\delta d^{\wedge}}=e^{-\frac14\delta d^{\wedge}}\cdot e^{-\frac14\delta d^{\wedge}}$
in the definition of $\varkappa_0$ in \eqref{4.27:eq-collar-displacement-5} gives
\begin{equation}\label{4.27:eq-collar-displacement-10}
\varkappa_0(y)\le2B_{\sharp}(20\mathsf{c}_{\square}r+11R_0)\,\vartheta\rho_0\alpha_{\cdot,0}\,
e^{-\frac14\delta d^{\wedge}(y,\mathcal{C}_j)}.
\end{equation}
In the same way \eqref{4.27:eq-collar-displacement-6} and Lemma~\ref{4.27:lem-level-zero-ratios}(c)
give
\begin{equation*}
\varkappa_{\mathsf{w}}(y)\le22B_{\sharp}\mathsf{c}_{\square}\mathsf{w}_0R_0\vartheta\rho_0
\alpha_{\cdot,0}\,e^{-\frac14\delta d^{\wedge}(y,\mathcal{C}_j)}.
\end{equation*}
The second condition in
\eqref{4.27:eq-collar-displacement-4}, $R_0\vartheta\le1/(2\mathsf{K}_r)$, then yields
\begin{equation}\label{4.27:eq-collar-displacement-11}
\varkappa_{\mathsf{w}}(y)\le\frac{11B_{\sharp}\mathsf{c}_{\square}}{\mathsf{K}_r}\,
\mathsf{w}_0\rho_0\alpha_{\cdot,0}\,e^{-\frac14\delta d^{\wedge}(y,\mathcal{C}_j)},
\end{equation}
which is the profile bound of the simultaneous-amplitude statement read at level $0$.

We now evaluate the coefficient in \eqref{4.27:eq-collar-displacement-10} with the closure
conditions (d). Recall $K_{\chi}\ge15dB_{\sharp}=60B_{\sharp}$ and
$\mathsf{c}_{\square},\mathsf{c}_{\vartheta}\ge1$. The definition of $\mathsf{K}_r$ gives
$\mathsf{K}_r\ge2^9\mathsf{c}_{\square}\mathsf{c}_{\vartheta}K_{\chi}$, and $\beta=1/5$ by
Lemma~\ref{4.27:lem-allowance-sum}. For the first term, the first condition in
\eqref{4.27:eq-collar-displacement-4} gives
\begin{equation*}
40B_{\sharp}\mathsf{c}_{\square}r\vartheta
\le\frac{40B_{\sharp}\mathsf{c}_{\square}\beta\mathsf{c}_{\vartheta}}{32\mathsf{K}_r}
\le\frac{40B_{\sharp}\beta}{32\cdot2^9K_{\chi}}
\le\frac{40\beta}{32\cdot2^9\cdot60}<10^{-5}.
\end{equation*}
For the second term, the second condition in \eqref{4.27:eq-collar-displacement-4} gives
\begin{equation*}
22B_{\sharp}R_0\vartheta\le\frac{11B_{\sharp}}{\mathsf{K}_r}
\le\frac{11B_{\sharp}}{2^9K_{\chi}}\le\frac{11}{2^9\cdot60}<4\cdot10^{-4}.
\end{equation*}
The coefficient $2B_{\sharp}(20\mathsf{c}_{\square}r+11R_0)\vartheta$ is the sum of these two
terms, so by \eqref{4.27:eq-collar-displacement-5} and \eqref{4.27:eq-collar-displacement-10}
\begin{equation}\label{4.27:eq-collar-displacement-12}
\mathsf{n}:=N_y(\chi\mathbb{H})\le\varkappa_0(y)\le5\cdot10^{-4}\,\rho_0\alpha_{\cdot,0}\,
e^{-\frac14\delta d^{\wedge}(y,\mathcal{C}_j)}\le\frac1{40}.
\end{equation}
The last inequality holds because $\rho_0=\beta2^{-(h_k+l)}\le1$ by
Definition~\ref{4.27:def-dated-allowances}, because
\begin{equation*}
\alpha_{\cdot,0}\le\alpha_{0,0}\le\mathsf{A}(\mathsf{t}_0)\le\frac1{480}
\end{equation*}
by \eqref{4.27:eq-collar-displacement-9}, and because the exponential is at most $1$. In the same
way, \eqref{4.27:eq-collar-displacement-11} together with
$\mathsf{K}_r\ge2^9\mathsf{c}_{\square}K_{\chi}\ge2^9\mathsf{c}_{\square}B_{\sharp}$ gives
\begin{equation*}
N_y(\chi_t\mathbb{H})\le\varkappa_{\mathsf{w}}(y)
\le\frac{11}{2^9}\,\mathsf{w}_0\rho_0\alpha_{\cdot,0}\le\frac1{40}.
\end{equation*}

\emph{Step 3: the entries $1$, $6$, $7$.} The map $\Phi^{\chi}$ sends $(\mathbb{U},\mathbb{J})$ to
$(e^{i\eta\chi\mathbb{H}}\mathbb{U},\mathbb{J}+\mathcal{J}^{\sharp}[\chi])$. Its first component is
the left multiplication of Lemma~\ref{4.27:lem-left-factor-algebra} with
$\mathfrak{h}=\chi\mathbb{H}$. The hypotheses of that lemma hold: $\mathsf{n}\le1/40$ by
\eqref{4.27:eq-collar-displacement-12}, and $\mathsf{a}\le1/80$, $\theta\le1/10$ by the preliminary.
Each $Q_a$ is a maximum of norms, and for a maximum of norms the triangle inequality gives
$|Q(x^+)-Q(x)|\le\max|x^+-x|$. Hence $\Delta_a(y)$ is bounded by the largest change of the
corresponding quantities, and Lemma~\ref{4.27:lem-left-factor-algebra} gives
\begin{equation*}
\Delta_1(y)\le4\mathsf{n},\qquad \eta\,\Delta_6(y)\le1.06\,\mathsf{n},\qquad
\eta^2\Delta_7(y)\le1.2\,\mathsf{n}.
\end{equation*}
By \eqref{4.27:eq-collar-displacement-12} each right-hand side is at most
$2\cdot10^{-3}\rho_0\alpha_{\cdot,0}e^{-\frac14\delta d^{\wedge}(y,\mathcal{C}_j)}$. Moreover
$\alpha_{\cdot,0}\le\alpha_{0,0}$ and $\alpha_{\cdot,0}\le\alpha_{1,0}$. Since
$2\cdot10^{-3}<\frac14$, this gives the first inequality of
\eqref{4.27:eq-collar-displacement-b} for $a\in\{1,6,7\}$, the radii $r_a(0)$ being expressed
through $\alpha_{0,0}$ and $\alpha_{1,0}$ as in Notation~\ref{4.27:not-collar-conventions}.

\emph{Step 4: the entry $3$.} Let $b_0$ be a bond of the stratum with initial point $y$ and
$b_0\in\Delta(y)$. Since $\ell(y)=\eta$ at such $y$, clause (iv) of (c) applies. Clause (v) gives
$a(y)p(y)\le q(y)$, and then, as in the current-closure statement for the families (C2),
\begin{equation}\label{4.27:eq-collar-displacement-13}
\eta^3\bigl|\mathcal{J}^{\sharp}[\chi](b_0)\bigr|
\le K_{\chi}\bigl[a(y)p(y)+q(y)\bigr]\le2K_{\chi}q(y)
\le40K_{\chi}\mathsf{c}_{\square}\,p(y)+20K_{\chi}|\tau|\,q_{\square}(y).
\end{equation}
The last inequality uses clause (v) for the amplitudes $|t_{\square}|\le2$ and the further amplitude
$\tau$ of the family (C2)($\beta$).

For the first member, $p(y)=\sup|b|\,e^{-\delta d^{\wedge}(y,\mathcal{C}_j)}$ and
$\sup|b|\le rC_1\delta'_j$. Lemma~\ref{4.27:lem-level-zero-ratios}(c) gives
\begin{equation*}
C_1\delta'_je^{-\delta d^{\wedge}}\le\vartheta\rho_0\alpha_{\cdot,0}e^{-\frac34\delta d^{\wedge}}.
\end{equation*}
Then $r\vartheta\le\beta\mathsf{c}_{\vartheta}/(32\mathsf{K}_r)$, from
\eqref{4.27:eq-collar-displacement-4}, together with
$\mathsf{K}_r\ge2^9\mathsf{c}_{\square}\mathsf{c}_{\vartheta}K_{\chi}$,
$\alpha_{\cdot,0}\le\alpha_{0,0}$ and $e^{-\frac34\delta d^{\wedge}}\le e^{-\frac14\delta d^{\wedge}}$,
gives
\begin{equation*}
\begin{split}
40K_{\chi}\mathsf{c}_{\square}\,p(y)
&\le40K_{\chi}\mathsf{c}_{\square}rC_1\delta'_j\,e^{-\delta d^{\wedge}(y,\mathcal{C}_j)}
\le40K_{\chi}\mathsf{c}_{\square}r\vartheta\,\rho_0\alpha_{\cdot,0}\,
e^{-\frac34\delta d^{\wedge}(y,\mathcal{C}_j)}\\
&\le\tfrac18\,\rho_0\alpha_{0,0}\,e^{-\frac14\delta d^{\wedge}(y,\mathcal{C}_j)},
\end{split}
\end{equation*}
since
\begin{equation*}
\frac{40K_{\chi}\mathsf{c}_{\square}\beta\mathsf{c}_{\vartheta}}{32\mathsf{K}_r}
\le\frac{40\beta}{32\cdot2^9}<\frac18.
\end{equation*}

For the second member, clause (vi) of (c), applied with $\varsigma=\frac12$ at the point $y$ (it
holds at every $y$), together with $e^{-\frac32\delta d(y,\tilde{\square})}\le1$ and
$\sup|b|\le1.1\,C_1\delta'_j$ (in the family (C2)($\beta$), $r\le1.1$), gives
\begin{equation*}
20K_{\chi}|\tau|\,q_{\square}(y)\le22K_{\chi}R_0C_1\delta'_j\,
e^{-\frac12\delta d^{\wedge}(y,\mathcal{C}_j)}.
\end{equation*}
Then
Lemma~\ref{4.27:lem-level-zero-ratios}(c), the condition $R_0\vartheta\le1/(2\mathsf{K}_r)$ and
$\mathsf{K}_r\ge2^9K_{\chi}$ give
\begin{equation*}
\begin{split}
20K_{\chi}|\tau|\,q_{\square}(y)
&\le22K_{\chi}R_0C_1\delta'_j\,e^{-\frac12\delta d^{\wedge}(y,\mathcal{C}_j)}
\le22K_{\chi}R_0\vartheta\,\rho_0\alpha_{\cdot,0}\,e^{-\frac14\delta d^{\wedge}(y,\mathcal{C}_j)}\\
&\le\frac{11K_{\chi}}{\mathsf{K}_r}\,\rho_0\alpha_{0,0}\,e^{-\frac14\delta d^{\wedge}(y,\mathcal{C}_j)}
\le\tfrac18\,\rho_0\alpha_{0,0}\,e^{-\frac14\delta d^{\wedge}(y,\mathcal{C}_j)},
\end{split}
\end{equation*}
since $11/2^9<\frac18$.

The second component of $\Phi^{\chi}$ adds $\mathcal{J}^{\sharp}[\chi]$ to $\mathbb{J}$. By the
inequality for a maximum of norms used in Step 3, $\eta^3\Delta_3(y)$ is therefore at most the
maximum over the bonds $b_0$ of the left-hand side of \eqref{4.27:eq-collar-displacement-13}. The
two estimates just obtained give
\begin{equation*}
\eta^3\Delta_3(y)\le\tfrac14\,\rho_0\alpha_{0,0}\,e^{-\frac14\delta d^{\wedge}(y,\mathcal{C}_j)}.
\end{equation*}
This is the first inequality of \eqref{4.27:eq-collar-displacement-b} for $a=3$.

\emph{The maps $\Phi_t$.} For a map $\Phi_t$ the second component adds
$\chi_t\mathcal{J}^{\sharp}[1]$ to $\mathbb{J}$. By the inequality for a maximum of norms used in
Step 3, $\eta^3\Delta_3(y)$ is at most the left-hand side of
\eqref{4.27:eq-collar-displacement-14}, and hypothesis (f$'$) bounds it by
\begin{equation*}
\frac{11}{2^9}\,\mathsf{w}_0\rho_0\alpha_{0,0}\,e^{-\frac14\delta d^{\wedge}(y,\mathcal{C}_j)}.
\end{equation*}
By the end of Step 2, Lemma~\ref{4.27:lem-left-factor-algebra} applies with
$\mathfrak{h}=\chi_t\mathbb{H}$. Using its constant $4$ for all three entries $1$, $6$, $7$, these
entries move by at most $4\varkappa_{\mathsf{w}}(y)/\alpha_{\cdot,0}$ in units of $r_a(0)$. By
\eqref{4.27:eq-collar-displacement-11} and $\mathsf{K}_r\ge2^9\mathsf{c}_{\square}B_{\sharp}$,
\begin{equation*}
\frac{4\varkappa_{\mathsf{w}}(y)}{\alpha_{\cdot,0}}
\le\frac{44}{2^9}\,\mathsf{w}_0\rho_0\,e^{-\frac14\delta d^{\wedge}(y,\mathcal{C}_j)}.
\end{equation*}
Since $44/2^9<\frac16$ and $11/2^9<\frac16$, this gives \eqref{4.27:eq-collar-displacement-7}
for $a\in\{1,3,6,7\}$.

\emph{The hypothesis of clause (iv).} Step 4 used clause (iv) of (c), whose hypotheses include
$2a(y)B_{\sharp}p(y)\le1/10$ (the hypothesis $\mathsf{p}(\mathbb{U})\le1$ was obtained in the
preliminary). This follows from Step 2: $a(y)p(y)\le q(y)$ by clause (v), and
\eqref{4.27:eq-collar-displacement-5} and \eqref{4.27:eq-collar-displacement-12} then give
\begin{equation*}
2a(y)B_{\sharp}p(y)\le2B_{\sharp}q(y)\le\varkappa_0(y)\le\frac1{40}<\frac1{10}.
\end{equation*}

\emph{The entry $5$ and the gauges.} The map $\Phi^{\chi}$ multiplies $\mathbb{U}$ on the left by
$e^{i\eta\chi\mathbb{H}}$, and the map $\Phi_t$ multiplies it on the left by
$e^{i\eta\chi_t\mathbb{H}}$. By Lemma~\ref{4.27:lem-left-factor-algebra}, the $G$-part $U$ and every
gauge of $U$ are unchanged under such a multiplication. Hence $\Delta_5(y)=0$ and the gauges on the
cubes are unchanged.

\emph{Conclusion.} Steps 3 and 4 give the first inequality of
\eqref{4.27:eq-collar-displacement-b} for $a\in\{1,3,6,7\}$. For the second inequality,
$e^{-\frac14\delta d^{\wedge}(y,\mathcal{C}_j)}\le1$, and
$\mathfrak{c}(d')=\mathfrak{c}(k;l)=\rho^{(l),(k)}_0=\rho_0$ by
Definition~\ref{4.27:def-dated-allowances}. This proves part (S).
\end{proof}

\begin{lemma}[Collar displacement at the averaging step and the completion site]
\label{4.27:prf-displacement-averaging}
This lemma is the part of Lemma~\ref{4.27:lem-collar-displacement} that concerns an operation
$\mathfrak{O}$ whose source date $d'$ is either an averaging date $(k'+1;\top)$ with $k'\ge 1$, or a
completion site $(k;L)$. The notation is that of Lemma~\ref{4.27:lem-collar-displacement}:
\begin{itemize}
\item $d^t$ is the target date, consecutive to $d'$, and $k$ is the density;
\item $\tilde p=(\mathbb{U},\mathbb{J})$ is a pair on $\Omega_0(\mathfrak{B}_{d'})$ lying in
$\mathfrak{K}_{109}(\mathfrak{B}_{d^t};\gamma_0^{(k)})$ and obeying
\eqref{4.27:eq-dated-allowances-a} at the date $d'$;
\item $\Phi$ is the map of $\mathfrak{O}$ on pairs, and
$\Delta_a(y)=|Q_a(\Phi\tilde p)(y)-Q_a(\tilde p)(y)|$;
\item $\mathfrak{c}(d')$ is the allowance of the date $d'$ fixed in
\eqref{4.27:eq-dated-allowances-b};
\item $r_a(0)$, for $a\in\{1,3,5,6,7\}$, is the level-zero radius of the entry $a$, built on one
of $\alpha_{0,0}$, $\alpha_{1,0}$: on $\alpha_{0,0}$ for $a\in\{1,3\}$ and on $\alpha_{1,0}$ for
$a\in\{6,7\}$.
\end{itemize}
Assume the following.
\begin{enumerate}
\item[(a)] \emph{The response field of the averaging operation.} Let $k'\ge 1$ and
$d'=(k'+1;\top)$, so that $\mathfrak{O}$ is Ba{\l}aban's operation $\mathbf{T}$ from level $k'$
to level $k'+1$. Let $\mathfrak{B}$ be the set of the last date of density $k'$.

The hypotheses under which the response field $\mathfrak{h}$ of the $\mathbf{T}$-step is
constructed hold, namely:
\begin{itemize}
\item the pair lies in $\mathfrak{K}_{109}(\mathfrak{B};\gamma_0^{(k')})$;
\item the datum $\mathsf{d}$ of the response field lies on a set
$\mathcal{C}\subset\Omega_{k'+1}$ and obeys $\|\mathsf{d}\|_1\le w_{k'}\le\varepsilon_\sharp$,
where $\varepsilon_\sharp$ is the smallness threshold of that construction;
\item $2B_\sharp w_{k'}\le 1/10$.
\end{itemize}
The map of $\mathbf{T}$ on pairs is
\begin{equation*}
(\mathbb{U},\mathbb{J})\longmapsto
\bigl(e^{i\eta\mathfrak{h}}\mathbb{U},\ \mathbb{J}+\mathcal{J}^{\sharp}[1]\bigr).
\end{equation*}
Here $\mathfrak{h}$ is that response field, constructed by the theorem on the response field of
a datum on a determining set (in what follows, the response-field theorem), and
$\mathcal{J}^{\sharp}[1]$ is the current increment at $\chi\equiv1$. For
$y\in\operatorname{St}(d')$ and every bond $b\in\Delta(y)$, the cell at $y$ (of which
$\tilde\Delta(y)$ is the enlargement),
\begin{equation}\label{4.27:eq-displacement-averaging-1}
N^{\Delta}_y(\mathfrak{h})\le B_\sharp\,p(y),\qquad
\eta^3|\mathcal{J}^{\sharp}[1](b)|\le 2K_\chi\,p(y),\qquad
p(y):=w_{k'}\,e^{-\delta d^{\wedge}(y,\mathcal{C})}.
\end{equation}
The constants $B_\sharp$, $K_\chi$, $\delta$ do not depend on the level. The first bound is
part~(b) of the response-field theorem; the second is the bound on the current increment in its
recursive construction. Both are taken at $\chi\equiv1$.
\item[(b)] \emph{The floor of the radii.} Put $C_H^{\sharp}:=\max\{C_H,2K_\chi\}$, where $C_H$ is
a constant of the response-field theorem, and $\Lambda_\gamma:=\mathsf{t}_\gamma^{\,p_0-q}\le 1$,
where $\mathsf{t}_\gamma$ denotes the value, attached to $\gamma$, of the logarithm $\mathsf{t}$
of the inverse squared coupling (written $\mathsf{t}_k$ at level $k$). Then
\begin{equation}\label{4.27:eq-displacement-averaging-2}
\frac{w_{k'}}{\alpha_{0,k'}}\le
\frac{2\|C\|C_1^{\mathrm{pd}}A_1}{C_0^{\mathrm{rad}}}\,\Lambda_\gamma,
\qquad
C_0^{\mathrm{rad}}\ge
\frac{16(1+\beta_0)^2\|C\|C_1^{\mathrm{pd}}A_1C_H^{\sharp}}{(1-\theta_S)\beta}.
\end{equation}
Here $\theta_S$ is a fixed number, independent of the date and of the level, which measures shares
of the allowance: conclusion~(1) below uses the share $\tfrac14(1-\theta_S)$ of it at the averaging
date and conclusion~(2) the share $\tfrac{\theta_S}{4}$ at the completion site. The same two
inequalities hold with $\alpha_{1,k'}$ and $C_1^{\mathrm{rad}}$ in place of
$\alpha_{0,k'}$ and $C_0^{\mathrm{rad}}$. Moreover $\theta_S\ge\tfrac12$ and
$e^{-\delta\mathsf{w}_M}\le e^{-2}$.
\item[(c)] \emph{The increment bound at the completion site, off the completing components.}
Let $d'=(k;L)$, and let $\mathfrak{O}$ be one of the operations of the completion site. Let $Z_L$
be the set on which the completion site changes the determining set (see~(f)), and let
$\mathcal{Z}\subset Z_L$ be the union of the zones of the completing components.
\begin{itemize}
\item Off $\mathcal{Z}$, the presented pair is
\begin{equation*}
\Bigl(\prod_s e^{i\eta\mathbb{H}_s}\mathbb{U},\ \mathbb{J}
+\sum_s\mathfrak{s}^{\sharp}_{\sigma_s}[\mathbb{H}_s]\Bigr),
\end{equation*}
the identity times the layers $\mathbb{H}_s$, with the $G$-part and the cube gauges kept; here
$\sigma_s$ is the parameter of the layer current $\mathfrak{s}^{\sharp}_{\sigma_s}[\cdot]$ of the
layer $\mathbb{H}_s$. The data $\mathrm{arg}_s$ of the layers are supported in $\mathcal{Z}$.
\item For $a\in\{1,3,6,7\}$, the bound on the increment of the entries at a completion site
(in what follows, the increment bound) states that
the increment at a bond $b$ of level $m=m(b)$ is at most
$K_W\Theta(h)(b)/\alpha_{\cdot,m}$ in threshold units, where, in the unit of $b$,
\begin{equation}\label{4.27:eq-displacement-averaging-3}
\Theta(h)(b)\le B_H\sum_s|\mathrm{arg}_s|\,
e^{-\delta_1 d_{\mathrm{sc}}(b,\operatorname{supp}\mathrm{arg}_s)}.
\end{equation}
Here the letter $h$ is that of the increment bound, $B_H$ and $\delta_1>0$ are constants
independent of the level, and $d_{\mathrm{sc}}$ is the distance in the scaled units of that bound.
The layer data obey $|\mathrm{arg}_s|\le 6\delta'_k$, with $\delta'_k$ the quantity $\delta'_j$ of
Lemma~\ref{4.27:lem-level-zero-ratios}(c) at $j=k$, or
$|\mathrm{arg}_s|\le K_\Lambda(M)(\alpha_{0,k}+\alpha_{1,k})$, with $K_\Lambda(M)$ a constant
depending on $M$; in the second case, in units of the threshold $\alpha_{\cdot,m}$, the bound
carries the quotient $\alpha_{0,k}/\alpha_{\cdot,m}$.
This bound is linear in $\mathrm{arg}_s$ up to the Lipschitz constant of the response-field
theorem.
\item For $b$ outside the completing components the far-distance bound
$d_{\mathrm{sc}}\ge 5R_k+10R_{\min}(k-1-m)$ holds, where $R_k$ is the separation, in own-scale
$M$-cubes, of distinct components of the large-field region $Z_k$ at level $k$, and each of the
$k-1-m$ strata crossed contributes $10R_{\min}$; the per-stratum factor obeys
$e^{-10\delta_1R_{\min}}\le e^{-2}$.
\item At levels $m\ge1$ the quotient $\alpha_{0,k}/\alpha_{\cdot,m}$ is bounded by
$(1+\beta_0)^2(2+t)^{1/2}$, $t:=k-1-m$, times the level-free quotient of the constants
$C_i^{\mathrm{rad}}$.
\item The increment bound obtained by multiplying these factors is stated with a number $P$ in
which $K_W$ enters as a factor, and $P\le\bar P$.
\end{itemize}
\item[(d)] \emph{The constants of the completion site.} The level-zero condition
\eqref{4.27:eq-collar-displacement-a} of the completion site holds, that is,
\begin{equation*}
16\bar P^{0}\le\beta\theta_S,\qquad
\bar P^{0}:=\bar P\max\{1,K_W^{0}/K_W\},\qquad
K_W^{0}:=\max\{4,\,3+18d+2.2K_J\},
\end{equation*}
with $d=4$ the dimension, here and below, and $K_J$ the constant of the current bound
\eqref{4.27:eq-displacement-averaging-4} in~(e). Moreover $4\alpha+\bar P^{0}\alpha\le 1/80$, where
$\alpha$ denotes the auxiliary parameter of that name in the tuple $\mathfrak{p}$ of
Ba{\l}aban's auxiliary parameters.
\item[(e)] \emph{The level-zero chain bounds and the collar corollary.} We use the proposition
bounding the entries of the presented chain at level zero (in what follows, the level-zero chain
proposition) and its corollary bounding the chain on the collars (in what follows, the collar
corollary), and we assume the level-zero condition on the response field under which the collar
corollary is stated. Then:
\begin{itemize}
\item Clauses (i)--(iii) of the level-zero chain proposition bound every entry $Q_a$ of the
presented chain by functions increasing in the layer sizes $F_s$. Clause (iv) contains no layer
amplitude.
\item Clause (iii) contains the bound
\begin{equation}\label{4.27:eq-displacement-averaging-4}
\eta^3\bigl|\mathfrak{s}^{\sharp}_{\sigma_s}[\mathbb{H}_s](b)\bigr|
\le(3+18d+2.2K_J)\,F_s(y).
\end{equation}
\item The collar corollary states: for every parameter point $\varsigma\in\mathfrak{P}$,
every $y\in\bigcup_C\operatorname{Col}(C)$, the union over the completing components $C$, and every
$a\in\{1,3,5,6,7\}$,
\begin{equation*}
Q_a(\operatorname{chain}(\varsigma))(y)\le(\varphi_\flat/16)\,\mathbf{r}_a(0).
\end{equation*}
Here $\mathbf{r}_a(0)$ is $\alpha_{*,0}$ times the unit, $\alpha_{*,0}=\min_i\alpha_{i,0}$, and
$\varphi_\flat\le1$.
\end{itemize}
\item[(f)] \emph{Facts at the completion site.} We use the construction of the layers at the
completion site (in what follows, the layer construction). In it, $h$ and $h'$ denote the two
levels at which it reads the large-field region $Z_{h'}$ and the layer width
$LM\check{R}_{h+1}$; $S$ denotes the region inside which the two layers forming $D_2$ below are
placed, and $B_S$ the part of $S$ avoided by the core $D_2\setminus B_S$ on which the pair is
read; $\Lambda$ denotes the region with which $Z_{h'}^{\sim2}$ is intersected to form $D_2$ (it
is not the region $\Lambda_k$ used in the proof, and it is unrelated to $\Lambda_\gamma$); and
$U_2$ is the configuration of the presented pair at the presentation point $\mathfrak{b}_2$, with
current $J_2$. The first fact below and the identity in the second are supplied by the arrest
theorem; the inequality $h_k\ge1$ in the second by a lemma on the completing components; the third
by a corollary on the collars of the completing components; the fourth by the arrest theorem
together with the layer construction; the fifth by the layer construction.
\begin{itemize}
\item The set $\mathfrak{B}_{d'}$ differs from $\mathfrak{B}_{d^t}=\mathbb{B}''_k$ only on $Z_L$,
where the level becomes $k$.
\item For every completing component $C$ one has $h_k\ge1$ and
$\Omega''_1\cap C=\Omega_1\cap C$, where $\Omega''_1:=\Omega_1(\mathbb{B}''_k)$.
\item $\operatorname{Col}(C)\subset Z_{h'}\cap C$.
\item On the core $D_2\setminus B_S$, where $D_2=Z_{h'}^{\sim2}\cap\Lambda$ consists of two layers
of width $LM\check{R}_{h+1}$ inside $S$ and off $B_S$, the presented pair has base
$(U_2,\mathcal{J}(U_2))$ and, as increments, the layers $\mathbb{H}_3$, $\mathbb{H}_4$.
\item The rotation of the presentation is trivial on $D_2$: its factors obey
$h^{\mathfrak{b}}_q=1$ for $\mathfrak{b}\ge\mathfrak{b}_2$, and $\mathfrak{Q}^{\mathfrak{b}_2}=1$.
\end{itemize}
\end{enumerate}
Then the following hold.
\begin{enumerate}
\item[(1)] At an averaging date $d'=(k'+1;\top)$, for every $y\in\operatorname{St}(d')$ and
$a\in\{1,3,6,7\}$, the bound \eqref{4.27:eq-collar-displacement-c} holds:
\begin{equation*}
\Delta_a(y)\le\Lambda_\gamma\cdot\tfrac14(1-\theta_S)\beta\,2^{-(k'+1)}r_a(0)
\le\tfrac18\,\mathfrak{c}(d')\,r_a(0).
\end{equation*}
Moreover $\Delta_5(y)=0$ and the cube gauges are kept.
\item[(2)] At a completion site $d'=(k;L)$, for every $y\in\operatorname{St}(d')$, the presented
pair at $y$ is the identity times the layers, $\Delta_5(y)=0$, the cube gauges are kept, and for
$a\in\{1,3,6,7\}$ the bound \eqref{4.27:eq-collar-displacement-d} holds:
\begin{equation*}
\Delta_a(y)\le\bar P^{0}\,2^{-(k-1)}r_a(0)\le\tfrac{\theta_S}{4}\,\mathfrak{c}(d')\,r_a(0),
\end{equation*}
and the bound is linear in the layer: a dilation of a layer by $\mathsf{w}$ multiplies the bound
by $\mathsf{w}$.
\item[(3)] At a completion site, for every parameter point of $\mathfrak{P}$ and every
$y\in\operatorname{St}(d^t)\setminus\operatorname{St}(d')=\bigcup_C\operatorname{Col}(C)$, the
union over the completing components $C$, the presented pairs at
$\mathfrak{b}_2$, $\mathfrak{b}_3$, $\mathfrak{b}_4$ and on the segments between them obey the
bound \eqref{4.27:eq-collar-displacement-e}:
\begin{equation*}
Q_a(y)\le(\varphi_\flat/16)\,\mathbf{r}_a(0)\le\tfrac1{16}\,r_a(0)
\qquad(a\in\{1,3,5,6,7\}).
\end{equation*}
\end{enumerate}
\end{lemma}

\begin{proof}
\emph{The averaging step.} Let $d'=(k'+1;\top)$ with $k'\ge1$, and let $\mathfrak{B}$ be the set
of the last date of density $k'$. By hypothesis~(a), the conditions of the response-field theorem
hold for the pair and the datum. Hence, at every $y\in\operatorname{St}(d')$, the response field
$\mathfrak{h}$ and the current increment $\mathcal{J}^{\sharp}[1]$ obey the two bounds of
\eqref{4.27:eq-displacement-averaging-1}, with the profile $p(y)=w_{k'}e^{-\delta
d^{\wedge}(y,\mathcal{C})}$ and with constants that do not depend on the level.

We first bound the distance from $y$ to $\mathcal{C}$.
\begin{itemize}
\item Since $\mathcal{C}\subset\Omega_{k'+1}$, a contour from $y\in\operatorname{St}(d')$ to
$\mathcal{C}$ crosses the strata $1,\dots,k'-1$.
\item It also crosses at least one $M$-cube of its own scale in $\Omega_{k'}\setminus
\Omega_{k'+1}$. This follows from the separation of boundaries stated in
\cite[p.~259]{B-CMP119} (the distance between the boundaries is at least $2MR_j$, with $j$ the
level in the notation there), together with
the geometric description of the strata used in the proof of the arrest theorem.
\item Each stratum crossed contributes at least the per-stratum contour price $\mathsf{w}_M$.
\end{itemize}
Therefore
\begin{equation}\label{4.27:eq-displacement-averaging-5}
d^{\wedge}(y,\mathcal{C})\ge\mathsf{w}_M\,k'.
\end{equation}

Next we check the size condition of Lemma~\ref{4.27:lem-left-factor-algebra} for the left factor
$e^{i\eta\mathfrak{h}}$. The size $\mathsf{n}$ there is at most $B_\sharp p(y)$ by the first bound
of \eqref{4.27:eq-displacement-averaging-1}. With \eqref{4.27:eq-displacement-averaging-5} and
the hypothesis $2B_\sharp w_{k'}\le1/10$ of~(a),
\begin{equation*}
\mathsf{n}\le B_\sharp p(y)\le B_\sharp w_{k'}e^{-8k'}
\le\tfrac1{20}\,e^{-8}<\tfrac1{40}.
\end{equation*}
Lemma~\ref{4.27:lem-left-factor-algebra} then writes $e^{i\eta\mathfrak{h}}\mathbb{U}$ again as a
left exponential factor times the unmoved $G$-part $U$, and every gauge of $U$ is unmoved. Hence
the cube gauges are kept and $\Delta_5(y)=0$. The displacements of the entries read from the bond
variables are bounded by
\begin{equation*}
\Delta_1,\ \eta\Delta_6,\ \eta^2\Delta_7\le 4B_\sharp p(y)\le 2K_\chi p(y)\le C_H^{\sharp}p(y).
\end{equation*}
The current moves by $\mathcal{J}^{\sharp}[1]$. By the second bound of
\eqref{4.27:eq-displacement-averaging-1}, and since $C_H^{\sharp}=\max\{C_H,2K_\chi\}$,
\begin{equation*}
\eta^3|\mathcal{J}^{\sharp}[1](b)|\le 2K_\chi p(y)\le C_H^{\sharp}p(y).
\end{equation*}

It remains to measure $C_H^{\sharp}p(y)$ against the level-zero radius. We follow the arithmetic
of the estimate for the operation $\mathbf{T}$ in the proof of the arrest theorem, carried out
here at level zero. Put $t:=k'$ and $u:=2^{-(k'+1)}$. We use three inputs:
\begin{itemize}
\item the first inequality of \eqref{4.27:eq-displacement-averaging-2} for $w_{k'}/\alpha_{0,k'}$;
\item Lemma~\ref{4.27:lem-level-zero-ratios}(a) with $\lambda=k'\ge1$, which gives
$\alpha_{0,k'}/\alpha_{0,0}\le(1+\beta_0)^2t^{1/2}$;
\item \eqref{4.27:eq-displacement-averaging-5} together with $e^{-\delta\mathsf{w}_M}\le e^{-2}$,
which gives $e^{-\delta d^{\wedge}(y,\mathcal{C})}\le e^{-\delta\mathsf{w}_Mt}\le e^{-2t}$.
\end{itemize}
Writing $(1+\beta_0)^2t^{1/2}e^{-2t}=(1+\beta_0)^2\,t^{1/2}(2e^{-2})^t\,2^{-t}$, we obtain
\begin{align*}
\frac{C_H^{\sharp}p(y)}{\alpha_{0,0}}
&\le C_H^{\sharp}\,\frac{w_{k'}}{\alpha_{0,k'}}\cdot\frac{\alpha_{0,k'}}{\alpha_{0,0}}
\cdot e^{-\delta\mathsf{w}_Mt}\\
&\le\Lambda_\gamma\,
\Bigl(\frac{2(1+\beta_0)^2\|C\|C_1^{\mathrm{pd}}A_1C_H^{\sharp}}{C_0^{\mathrm{rad}}}\Bigr)
\cdot\bigl[t^{1/2}(2e^{-2})^t\bigr]\cdot2^{-t}.
\end{align*}

The bracket is at most $1$ for $t\ge1$. Indeed, the logarithmic derivative of $t\mapsto
t^{1/2}(2e^{-2})^t$ is $\tfrac1{2t}+\log(2e^{-2})$. This is negative for $t\ge1$, because
$\log(2e^{-2})=\log2-2<-\tfrac12$. Hence the bracket is largest at $t=1$, where it equals
$2e^{-2}<0.271$.

By the second inequality of \eqref{4.27:eq-displacement-averaging-2}, the constant in parentheses
is at most $(1-\theta_S)\beta/8$. Hence
\begin{equation*}
\frac{C_H^{\sharp}p(y)}{\alpha_{0,0}}
\le\Lambda_\gamma(1-\theta_S)\frac{\beta}{8}\,2^{-t}
=\Lambda_\gamma\cdot\tfrac14(1-\theta_S)\,\beta u .
\end{equation*}
This bounds $\Delta_a(y)/r_a(0)$ for $a\in\{1,3\}$. For $a\in\{6,7\}$ the same chain runs with
$\alpha_{1,\cdot}$ and $C_1^{\mathrm{rad}}$, using the second half of hypothesis~(b). This is the
first inequality of \eqref{4.27:eq-collar-displacement-c}.

For the second inequality, $\Lambda_\gamma\le1$ and $\theta_S\ge\tfrac12$ give
$\tfrac14(1-\theta_S)\le\tfrac18$. Hence the right-hand side is at most $\tfrac18\beta u$. This is
the second inequality of \eqref{4.27:eq-collar-displacement-c}, with $\mathfrak{c}(d')$ as fixed
in \eqref{4.27:eq-dated-allowances-b} at the averaging date. This proves~(1).

\emph{The completion site, off the completing components.} Let $d'=(k;L)$ and
$y\in\operatorname{St}(d')$. By \cite[(2.1)--(2.3)]{B-CMP119}, $y$ lies in
$\Omega_1^{c}\subset Z_1\subset Z_k$, where $Z_1$ and $Z_k$ are the large-field regions at levels
$1$ and $k$. It therefore lies in a component $W$ of $Z_k$; since $y\in\operatorname{St}(d')$
lies in no completing component, $W$ does not complete at $(k;L)$. We collect five facts at level
$m(b)=0$.

\emph{Form.} By hypothesis~(c), off the zone the presented pair is the identity times the layers:
\begin{equation*}
\Bigl(\prod_s e^{i\eta\mathbb{H}_s}\mathbb{U},\ \mathbb{J}
+\sum_s\mathfrak{s}^{\sharp}_{\sigma_s}[\mathbb{H}_s]\Bigr).
\end{equation*}
The $G$-part and the cube gauges are kept. The data $\mathrm{arg}_s$ are supported in the zones
$\mathcal{Z}\subset Z_L$ of the completing components.

\emph{Size.} In the unit of $b$, \eqref{4.27:eq-displacement-averaging-3} holds, with constants
independent of the level. Equivalently, it is part~(b) of the response-field theorem applied with
the profile $p:=e^{-\delta d^{\wedge}(\cdot,Z_L)}$.

\emph{Distance.} The component $W$ and the completing components are distinct components of the
complement $\Lambda_k^{c}$ of the level-$k$ region $\Lambda_k\subset\Omega_k$. Hence a contour
from $Z_L$ to $y$ first passes through $\Lambda_k\subset
\Omega_k$, at level $k$, and then crosses the strata $k-1,\dots,1$ of $W$. This is the
stratum-crossing argument behind Lemma~\ref{4.27:lem-level-zero-ratios}(b). By the far-distance
bound of hypothesis~(c) at $m(b)=0$,
\begin{equation*}
d_{\mathrm{sc}}\ge 5R_k+10R_{\min}(k-1).
\end{equation*}
The per-stratum factor obeys $e^{-10\delta_1R_{\min}}\le e^{-2}$, in the same way as
$e^{-\delta\mathsf{w}_M}\le e^{-2}$. Hence the exponential in
\eqref{4.27:eq-displacement-averaging-3} is at most $e^{-5\delta_1R_k}\,e^{-2(k-1)}$.

\emph{Ratio.} By Lemma~\ref{4.27:lem-level-zero-ratios}(a) with $\lambda=k\ge1$,
\begin{equation*}
\alpha_{i,k}/\alpha_{i,0}\le(1+\beta_0)^2k^{1/2}\le(1+\beta_0)^2(1+k)^{1/2}.
\end{equation*}
So at $m=0$ the quotient $\alpha_{0,k}/\alpha_{\cdot,m}$ of the increment bound is bounded by the
same expression as at $m\ge1$, namely $(1+\beta_0)^2(2+t)^{1/2}$ with $t:=k-1-m$; at $m=0$,
$2+t=1+k$. The level-free quotient of the constants $C_i^{\mathrm{rad}}$ is the same number at
every $m$.

\emph{Local algebra.} Let $\mathsf{n}$ be the sum of the sizes of the layers at $y$. Apply
Lemma~\ref{4.27:lem-left-factor-algebra} layer by layer, with $\mathfrak{h}=\mathbb{H}_s$ and with
$\mathbb{U}$ replaced by $W_s$, the pair obtained after the layers preceding $s$.
\begin{itemize}
\item The size $N_y(\mathbb{H}_s)$ is taken with the covariant derivative $\nabla_{W_s}$. This is
the derivative $\nabla_{\mathbb{U}}$ of Lemma~\ref{4.27:lem-left-factor-algebra} at each
application.
\item The intermediate pairs obey the hypotheses of that lemma: each intermediate configuration
$W_s$ differs from $\mathbb{U}$ by a left factor whose size is at most the total $\mathsf{n}$ of
the layer sizes, and $4\alpha+\bar P^{0}\alpha\le1/80$ by hypothesis~(d), which gives the bound
$\mathsf{a}\le1/80$ of that lemma for each of them.
\item The $G$-part is the same $U$ throughout.
\end{itemize}
Summing over the layers gives
\begin{equation*}
\Delta_1,\ \eta\Delta_6,\ \eta^2\Delta_7\le 4\mathsf{n}.
\end{equation*}
Since the $G$-part is the same $U$ throughout and, at each application of
Lemma~\ref{4.27:lem-left-factor-algebra}, every gauge of $U$ is unmoved, the cube gauges are kept
and $\Delta_5(y)=0$.
The current moves by $\sum_s\mathfrak{s}^{\sharp}_{\sigma_s}[\mathbb{H}_s]$. The bound
\eqref{4.27:eq-displacement-averaging-4} of clause (iii) of the level-zero chain proposition
applies at the spacing $\ell=\eta$; it rests on the definition of the layer currents
$\mathfrak{s}^{\sharp}_{\sigma}[\cdot]$ and on the kernel bounds used there. Thus both
local-algebra constants, $4$ and $3+18d+2.2K_J$, are at most
$K_W^{0}$, and $K_W^{0}$ takes the place of the constant $K_W$ of the increment bound.

\emph{Multiplication.} Put $t:=k-1$. The inequality
\begin{equation*}
(2+t)^{1/2}e^{-2t}\le\sqrt2\cdot2^{-t}
\end{equation*}
holds for $t\ge0$. At $t=0$ it is an equality. The logarithmic derivative of
$(2+t)^{1/2}(2e^{-2})^t$ is $\tfrac1{2(2+t)}+\log2-2$, which is at most $\tfrac14+\log2-2<0$.

Multiplying the factors supplied by the last four facts --- the size factor, the exponential
$e^{-5\delta_1R_k}e^{-2t}$, the radius quotient $(1+\beta_0)^2(2+t)^{1/2}$ and the local-algebra
constant --- and using that $K_W$ enters $P$ as a factor, we obtain
\begin{equation*}
\frac{\Delta_a(y)}{r_a(0)}\le\frac{K_W^{0}}{K_W}\,P\,2^{-(k-1)}
\le\bar P\max\{1,K_W^{0}/K_W\}\,2^{-(k-1)}=\bar P^{0}\,2^{-(k-1)}.
\end{equation*}
Here $P\le\bar P$ by hypothesis~(c). This is the first inequality of
\eqref{4.27:eq-collar-displacement-d}.

By \eqref{4.27:eq-collar-displacement-a}, $16\bar P^{0}\le\beta\theta_S$, hence
\begin{equation*}
\bar P^{0}\,2^{-(k-1)}\le\frac{\theta_S}{4}\cdot\frac14\beta\,2^{-(k-1)}.
\end{equation*}
This is the second inequality of \eqref{4.27:eq-collar-displacement-d}, with $\mathfrak{c}(d')$
as fixed in \eqref{4.27:eq-dated-allowances-b} at the completion site.

\emph{Linearity.} The bound \eqref{4.27:eq-displacement-averaging-3} is linear in
$\mathrm{arg}_s$ up to the Lipschitz constant of the response-field theorem. The estimate for the
operation $\mathbf{T}$ in the proof of the arrest theorem states it in this form, linearly in the
layer. Hence a dilation of a layer by $\mathsf{w}$ multiplies the bound by $\mathsf{w}$.

\emph{An alternative derivation for the layers $s=3,4$.} The same shape of bound also follows from
the construction of the layers at the completion site.
\begin{itemize}
\item The layer data $B_3$, $B_4$ are supported in the completing component.
\item $\mathsf{J}_3$ is built from the current at $\mathfrak{b}_2$ through the cut-offs of that
construction, and is supported there.
\item The amplitude obeys $\mathsf{a}_{\star}(k)\le C_{\mathsf{a}}\alpha_{*,k}$, with a constant
$C_{\mathsf{a}}$ of that construction.
\item $\|\tilde G_0(\sigma_3)\mathsf{J}_3\|^{\Delta}_p\le K_\sharp\|\mathsf{J}_3\|_p$, with a
constant $K_\sharp$ of that construction.
\end{itemize}
With $p:=e^{-\delta d^{\wedge}(\cdot,Z_L)}\in\mathfrak{M}_\delta$, part~(b) of the response-field
theorem gives
\begin{equation*}
N_y(\mathbb{H}_s)\le B_\sharp\,\mathsf{a}_{\star}\,p(y)\le B_\sharp C_{\mathsf{a}}\,
\alpha_{*,k}\,e^{-\delta\mathsf{w}_M(k-1)}\,e^{-\delta\mathsf{w}_MR_k}.
\end{equation*}
The last factor appears because the components of $Z_k$ are at least $R_k$ own-scale $M$-cubes
apart. By Lemma~\ref{4.27:lem-level-zero-ratios}(a),
$\alpha_{*,k}\le(1+\beta_0)^2k^{1/2}\alpha_{*,0}$. The bound thus has the same shape as above,
with the smallness in $\gamma$ carried by $e^{-\delta\mathsf{w}_MR_k}$. This proves~(2).

\emph{The completion site, on the collars.} By hypothesis~(f), $\mathfrak{B}_{d'}$ differs from
$\mathfrak{B}_{d^t}=\mathbb{B}''_k$ only on $Z_L$, where the level becomes $k$. Hence
\begin{equation*}
\operatorname{St}(d^t)\setminus\operatorname{St}(d')=\operatorname{St}(d^t)\cap Z_L.
\end{equation*}
Let $C$ be a completing component. By hypothesis~(f), $h_k\ge1$ and $\Omega''_1\cap
C=\Omega_1\cap C$. Hence
\begin{equation*}
\operatorname{St}(d^t)\cap C=(\Omega''^{\sim}_1\setminus\Omega''_1)\cap C,
\end{equation*}
which is the collar $\operatorname{Col}(C)$ by its definition.

At $\mathfrak{b}_4$, the bound \eqref{4.27:eq-collar-displacement-e} is the collar corollary of
hypothesis~(e).

At $\mathfrak{b}_3$, at $\mathfrak{b}_2$ and on the segments between them, the chain is
\begin{equation*}
e^{i\eta t_4\mathbb{H}_4}e^{i\eta t_3\mathbb{H}_3}U_2,\qquad
\mathcal{J}=J_2+\sum_s\mathfrak{s}^{\sharp}_{\sigma_s}[t_s\mathbb{H}_s],\qquad t_s\in[0,1].
\end{equation*}
The parameters $t_s$ enter through the definition of the layer currents
$\mathfrak{s}^{\sharp}_{\sigma_s}[t_s\mathbb{H}_s]$, which carries them. Clauses (i)--(iii) of the
level-zero chain proposition bound every entry by functions increasing in the layer sizes $F_s$.
For the chain above these sizes become $t_sF_s\le F_s$, so the bounds at $\mathfrak{b}_4$ bound
the chain at every point of the segments. Clause (iv) does not involve the $t_s$. This gives the
first inequality of \eqref{4.27:eq-collar-displacement-e}.

For the second inequality: $\mathbf{r}_a(0)$ is $\alpha_{*,0}$ times the unit, with
$\alpha_{*,0}=\min_i\alpha_{i,0}$. Since $\alpha_{*,0}$ is at most each of $\alpha_{0,0}$ and
$\alpha_{1,0}$, on which the radii $r_a(0)$ are built, $\mathbf{r}_a(0)\le r_a(0)$; with
$\varphi_\flat\le1$ this gives $(\varphi_\flat/16)\mathbf{r}_a(0)\le\tfrac1{16}r_a(0)$.

Finally, the pair bounded here is the one the arrest theorem uses on the core $D_2\setminus B_S$,
and this core contains the collar. Indeed, $\operatorname{Col}(C)\subset Z_{h'}\cap C$, and
$D_2=Z_{h'}^{\sim2}\cap\Lambda$ consists of two layers of width $LM\check{R}_{h+1}$ inside $S$,
off $B_S$. By hypothesis~(f), on this core the base is $(U_2,\mathcal{J}(U_2))$ and the increments
are the layers $\mathbb{H}_3$, $\mathbb{H}_4$. The rotation of the presentation is trivial there:
on $D_2$, $h^{\mathfrak{b}}_q=1$ for $\mathfrak{b}\ge\mathfrak{b}_2$, and
$\mathfrak{Q}^{\mathfrak{b}_2}=1$. This proves~(3).
\end{proof}

\begin{lemma}[The real configuration lies deep inside the collar clause]\label{4.27:lem-real-point-collar}
In this lemma $\eta$ is the spacing of the lattice on which the configurations of level zero live;
for a set $\mathfrak{B}$ of a date, $\Omega_0(\mathfrak{B})$, $\Omega_1(\mathfrak{B})$,
$\Omega_2(\mathfrak{B})$ are its regions of levels $0,1,2$, so that
$\operatorname{St}(d)=\Omega_0(\mathfrak{B}_d)\setminus\Omega_1(\mathfrak{B}_d)$ is the never-averaged
collar; $U(\mathfrak{B},V)$ is the configuration that $\mathfrak{B}$ attaches to the data $V$;
$\mathcal{J}(U)$ is the current of $U$; $Q_a(y)$, $a\in\{1,3,5,6,7\}$, are the entries of a pair at
$y$, and $r_a(0)$ is the level-zero radius of entry $a$. A plaquette is \emph{pinned} if none of its
corners lies in $\Omega_1(\mathfrak{B})$. By the definition of the entries and radii of the spaces of
the arrest theorem, at $y\in\operatorname{St}(d)$ entries one and five of a pair are read on the
plaquette variables of the pinned plaquettes at $y$, entry three on $\eta^3$ times the norm of the
current of the pair at the bonds at $y$, and entries six and seven on the potential $A'$ of the pair;
and $\alpha_{0,0}\le r_1(0)$, $\alpha_{0,0}\le r_3(0)$. In $2(d-1)$ and in $c(d,L)$
the letter $d=4$ is the dimension; elsewhere in this lemma $d$ is a date. For an $N\times N$ complex
matrix $X$ write $\operatorname{Tr}X$ for its ordinary matrix trace and
$\pi(X)=X-N^{-1}(\operatorname{Tr}X)\,1$, with $1$ the identity matrix,
and let $c_\pi$ denote the least constant with $|\pi(X)|\le c_\pi|X|$ for all $X$.
Assume the following.
\begin{enumerate}
\item[(a)] The small-field bound at level zero, taken by statement \cite[p.~248]{B-CMP119}: for real
data $V$ in the support of the characteristic functions of the density, every pinned plaquette $p$
within distance $2LM\check{R}_1$ of $\Omega_1$ satisfies $|V_0(\partial p)-1|<\varepsilon_0$, where
$V_0$ is the level-zero component of the data $V$ and
\begin{equation}\label{4.27:eq-real-point-collar-1}
\varepsilon_0=A_0\,e^{-\mathsf{t}_0/2}\,\mathsf{t}_0^{p_0}.
\end{equation}
\item[(b)] The floor of that bound: $\mathsf{a}:=A_0/A_1\ge c(d,L)$, with $c(d,L)$ depending on $d$
and $L$ only.
\item[(c)] The radius condition with factor $\varphi_{\mathrm{ARR}}$ attached to the real-point bound
of the arrest theorem, in the form in which it is used below: together with $\mathsf{t}_0\ge1$ it
gives $(A_0/C_0^{\mathrm{rad}})\,\mathsf{t}_0^{p_0-q}\le\varphi_{\mathrm{ARR}}/2.5$.
\item[(c$'$)] The factor satisfies $\varphi_{\mathrm{ARR}}<1/10$, so that
$\varphi_{\mathrm{ARR}}/2.5<0.04$.
\item[(d)] The real-point bound of the arrest theorem at a fraction $\mathsf{m}\in\,]0,\tfrac12]$,
under its hypothesis at that fraction, together with the statement that all entries are at most one
half of the radii. At level one these give, for the configuration $U=U(\mathfrak{B},V)$ of the
conclusion, at a level-one plaquette $p$ to which that bound attaches weight one (its weight rule is
recalled in the proof),
$|U(\partial p)-1|\le\mathsf{m}\,(1+\beta)\,\alpha_{0,1}L^{-2}$.
\item[(e)] The current formula at a real configuration: for $W$ real,
$\mathcal{J}(W)=D^*_W\eta^{-2}\pi\operatorname{Im}\partial W$ \cite[(1.8)]{B-CMP109}, where
\begin{equation}\label{4.27:eq-real-point-collar-2}
(D^*_WF)(b)=\eta^{-1}\sum_{p\ni b}\pm\operatorname{Ad}_{\tau_p}F(p),
\end{equation}
the sum running over the $2(d-1)$ plaquettes $p$ containing the bond $b$, $\tau_p$ being the
parallel transporter of $W$ attached to $p$ and $b$ in \cite[(1.8)]{B-CMP109} and the sign being
the relative orientation of $b$ and $\partial p$.
\item[(f)] The axial-gauge bound on cubes, in the form in which it is used below, on cubes of at
most $10M$ sites per side: if $\square$ is a cube with $s$ sites per side,
$s\le10M$, and $\varepsilon>0$ satisfies $\sup_{p\subset\square}|U(\partial p)-1|\le\varepsilon$ and
$2(d-1)s\varepsilon\le1/10$, then there is a gauge transformation on $\square$ with values in the
gauge group $G$ of Definition~\ref{4.4:def-regular-backgrounds}, in which the
Lie-algebra-valued gauge potential $A$ of $U$ (not the action $A(U)$) obeys
$\eta|A|\le25(d-1)M\varepsilon$ and
$\eta^2|\nabla A|\le50(d-1)M\varepsilon$.
\item[(g)] The two one-sided conditions
\begin{equation}\label{4.27:eq-real-point-collar-a}
\sup_{\mathsf{t}\ge\mathsf{t}_\gamma-c_\beta}24M\,\mathsf{A}(\mathsf{t})\le1,
\qquad
\sup_{\mathsf{t}\ge\mathsf{t}_\gamma-c_\beta}6\,\mathsf{A}(\mathsf{t})\le\alpha_0^{(99)},
\end{equation}
where $\mathsf{t}$ is the logarithmic coupling variable and the range
$\mathsf{t}\ge\mathsf{t}_\gamma-c_\beta$ of the suprema is the one attached to the size parameter
$\gamma$, $\mathsf{t}_\gamma$ and
$c_\beta$ being the point and the offset that fix the lower end of that range; $\mathsf{A}$ is the
largest of the radii as a function of $\mathsf{t}$; and $\alpha_0^{(99)}$ is the radius $\alpha_0$ of
\cite{B-CMP99b}. Both conditions are one-sided caps of degree $-\infty$ in the sense of
Definition~\ref{2.3:def-cap}, $M$ and $\alpha_0^{(99)}$ being fixed before $\gamma$.
\end{enumerate}
Let $V$ be real and in the support of the characteristic functions of the density, let $d$ be a
date, and let $\mathfrak{B}=\mathfrak{B}_d$ be any set of that date; put $U=U(\mathfrak{B},V)$. Then
the pair $(U,\mathcal{J}(U))$ has, at every $y\in\operatorname{St}(d)$,
\begin{gather}
Q_6(y)=Q_7(y)=0,\label{4.27:eq-real-point-collar-3}\\
Q_1(y)=Q_5(y)<\varepsilon_0\le(\varphi_{\mathrm{ARR}}/2.5)\,r_1(0)<0.04\,r_1(0),
\label{4.27:eq-real-point-collar-4}\\
Q_3(y)\le c_{\mathcal{J}}\cdot0.04\cdot r_3(0)\le0.48\,r_3(0),\qquad
c_{\mathcal{J}}:=2(d-1)c_\pi\le12,\label{4.27:eq-real-point-collar-5}
\end{gather}
and it obeys \cite[(3.35)]{B-CMP99b} on the cubes of index zero.
\end{lemma}

\begin{proof}
\emph{Entries six and seven.} For the pair $(U,\mathcal{J}(U))$ the potential $A'$ on which entries
six and seven are read vanishes, $A'=0$, and this gives \eqref{4.27:eq-real-point-collar-3}.

\emph{Entries one and five.} At level zero the constraint is the identity. Indeed,
\cite[(2.10)--(2.11)]{B-CMP119} state that $V=V_j$ on $\Gamma_j$ and that
$M_{\mathbf{B}}(U)=M^j(U)$ on $\Gamma_j$, where $\Gamma_j$ is the region of level $j$, $V_j$ the
datum there, $M^j$ the $j$-fold averaging and $M_{\mathbf{B}}$ the averaging of the operation
$\mathbf{B}$; and $M^0$ is the identity map, by the level-zero statement that the averaging map of
level zero is the identity.
Hence $U=V_0$ on the bonds
of the stratum $\operatorname{St}(d)$, and on a pinned plaquette $p$
\begin{equation*}
|U(\partial p)-1|=|V_0(\partial p)-1|<\varepsilon_0
\end{equation*}
by hypothesis (a). Hypothesis (a) applies here for two reasons. First, for the region $\Omega_1$
made by the operation $\mathbf{T}_1$ one has
$\operatorname{St}(d)\subset\tilde{\Omega}_1\setminus\Omega_1$, where $\tilde{\Omega}_1$ is the set of
points within distance $2LM\check{R}_1$ of $\Omega_1$, by the level-zero conventions on the
never-averaged collar (Notation~\ref{4.27:not-collar-conventions}). Second, the characteristic
functions of step zero are factors of
\begin{equation*}
\mathbf{T}_k(X)=\prod_j\mathbf{T}^{(j)}(Z_{j+1}\cap X)
\end{equation*}
at every date before the completion of $X$ \cite[(2.20)--(2.22)]{B-CMP119}, so that the
support condition on $V$ in (a) is the one in force at every such date. Next, the level-zero radius
has the form $\alpha_{0,0}=C_0^{\mathrm{rad}}\,e^{-\mathsf{t}_0/2}\,\mathsf{t}_0^{q}$; hence, by
\eqref{4.27:eq-real-point-collar-1}, by hypotheses (c) and (c$'$) and by $\mathsf{t}_0\ge1$,
\begin{equation}\label{4.27:eq-real-point-collar-6}
\frac{\varepsilon_0}{\alpha_{0,0}}
=\frac{A_0}{C_0^{\mathrm{rad}}}\,\mathsf{t}_0^{p_0-q}
\le\frac{\varphi_{\mathrm{ARR}}}{2.5}<0.04.
\end{equation}
Since $U=V_0$ on the stratum bonds, the plaquette variables of $U$ and of $V_0$ at the pinned
plaquettes at $y$ coincide; hence entries one and five agree, $Q_1(y)=Q_5(y)$, and both are at most
the largest value of $|U(\partial p)-1|$ over the pinned plaquettes $p$ at $y$, which is less than
$\varepsilon_0$. With $\alpha_{0,0}\le r_1(0)$ and \eqref{4.27:eq-real-point-collar-6}, these are the
bounds on the entries one and five in
\eqref{4.27:eq-real-point-collar-4}; the last inequality there is $\varphi_{\mathrm{ARR}}/2.5<0.04$,
by hypothesis (c$'$).

\emph{Entry three: the current.} By hypothesis (e), for the real configuration $W=U$ the current is
$\mathcal{J}(U)=D^*_U\eta^{-2}\pi\operatorname{Im}\partial U$, and by
\eqref{4.27:eq-real-point-collar-2}
\begin{equation*}
\eta^3\,\mathcal{J}(U)(b)=\sum_{p\ni b}\pm\operatorname{Ad}_{\tau_p}\,
\pi\bigl(\operatorname{Im}U(\partial p)\bigr),
\end{equation*}
a sum of $2(d-1)$ terms. Since $U$ is real, each $\operatorname{Ad}_{\tau_p}$ has norm one. For
unitary $X$ one has $|\operatorname{Im}X|\le|X-1|$, since the imaginary part of $X-1$ is
that of $X$. Finally $c_\pi\le2$, because
$|N^{-1}(\operatorname{Tr}X)\,1|\le|X|$, so that $|\pi(X)|\le2|X|$ for all $X$. Therefore
\begin{equation}\label{4.27:eq-real-point-collar-7}
\eta^3|\mathcal{J}(U)(b)|\le c_{\mathcal{J}}\max_{p\ni b}|U(\partial p)-1|,
\qquad c_{\mathcal{J}}=2(d-1)c_\pi\le2\cdot3\cdot2=12.
\end{equation}

It remains to bound $|U(\partial p)-1|$ for the plaquettes $p$ containing a stratum bond $b$. Such a
plaquette is of one of two kinds.

\emph{(i) $p$ is pinned.} Then $|U(\partial p)-1|<\varepsilon_0$, as shown above. If $p$ leaves
$\Omega_0$, its corners lie within distance $\eta$ of $\operatorname{St}(d)$, hence inside the margin
$2LM\check{R}_1$ of hypothesis (a), which therefore applies to $p$.

\emph{(ii) $p$ has a corner in $\Omega_1(\mathfrak{B})$.} Then $p$ has no corner in
$\Omega_2(\mathfrak{B})$. Indeed, $p$ contains the stratum bond $b$, whose endpoints lie in
$\operatorname{St}(d)$ and hence outside $\Omega_1$, so $p$ has a corner in the complement of
$\Omega_1$; by \cite[p.~277, (1)]{B-CMP102b} the distance between the
complement of $\Omega_1$ and $\Omega_2$ exceeds the constant of that display times $M_1L\eta$, and
this exceeds the diameter of $p$; and the levels of the regions $\Omega_j(\mathfrak{B})$ are
consecutive, by the geometric statement that these regions have consecutive levels. Thus $p$
is a plaquette of level one. The weight that the real-point bound of the arrest theorem attaches to
an entry $a$ at a level $\ell$ is one unless $k_0^a<\ell\le k_a$, where $k_0^a$ and $k_a$ are the
ends of the window of levels in which that weight may differ from one; at $\ell=1$ it is one because
$k_0^a>h_a\ge1$, $h_a$ being the first level of the clock attached to entry $a$. Hypothesis (d)
therefore applies to $p$. By (d) with $\mathsf{m}\le\frac12$, by the numerical bound
$1+\beta\le1.2$ on the constant $\beta$ that enters the allowances of
Definition~\ref{4.27:def-dated-allowances}, and by the ratio of radii of
Lemma~\ref{4.27:lem-level-zero-ratios}(a) at $\lambda=1$, which gives
$\alpha_{0,1}\le(1+\beta_0)^2\alpha_{0,0}$, together with the numerical bound
$(1+\beta_0)^2\le\frac{64}{49}$,
\begin{equation}\label{4.27:eq-real-point-collar-8}
|U(\partial p)-1|\le\mathsf{m}\,(1+\beta)\,\alpha_{0,1}L^{-2}
\le\tfrac12\,(1.2)\,\tfrac{64}{49}\,L^{-2}\alpha_{0,0}<0.013\,\alpha_{0,0};
\end{equation}
the last step uses the one-sided lower bound on the blocking factor $L$ under which
$\tfrac12\cdot1.2\cdot\tfrac{64}{49}\,L^{-2}<0.013$, a condition on $L$ that holds for all
sufficiently large $L$.

In both cases, by \eqref{4.27:eq-real-point-collar-6} and \eqref{4.27:eq-real-point-collar-8},
\begin{equation*}
\max_{p\ni b}|U(\partial p)-1|\le\max\{\varepsilon_0,\,0.013\,\alpha_{0,0}\}<0.04\,\alpha_{0,0}.
\end{equation*}
Inserted in \eqref{4.27:eq-real-point-collar-7}, this gives
$\eta^3|\mathcal{J}(U)(b)|\le c_{\mathcal{J}}\cdot0.04\cdot\alpha_{0,0}$ at every stratum bond $b$.
Since entry three is read on $\eta^3|\mathcal{J}(U)(b)|$ at the bonds $b$ at $y$, and
$\alpha_{0,0}\le r_3(0)$, this is the first inequality of \eqref{4.27:eq-real-point-collar-5}. The
second inequality there is
the arithmetic $c_{\mathcal{J}}\cdot0.04\le12\cdot0.04=0.48$.

\emph{The cubes of index zero.} Let $\square$ be a cube of index zero, so that
$\square\subset B^0(\Lambda_0)\cup B^1(\Lambda_1)$, in the notation of \cite[p.~396]{B-CMP99b}. As
stated in \cite[p.~396]{B-CMP99b}, the size of a cube of index $j$ in the lattice of spacing $\eta$
is $O(1)ML^j\eta$, where $O(1)$ denotes a number at most $10$. For $j=0$ this gives that $\square$
has $s\le10M$ sites per side. Every plaquette $p$ of $\square$ is either pinned, with
$|U(\partial p)-1|<\varepsilon_0$ by hypothesis (a), or of level one, with
\eqref{4.27:eq-real-point-collar-8} by hypothesis (d); hence
\begin{equation*}
\sup_{p\subset\square}|U(\partial p)-1|\le\varepsilon_\square
:=\max\{\varepsilon_0,\,0.013\,\alpha_{0,0}\}<0.04\,\alpha_{0,0}.
\end{equation*}
Hypothesis (f), applied with $\varepsilon=\varepsilon_\square$, gives a $G$-valued gauge on $\square$
with
\begin{equation*}
\eta|A|\le25(d-1)M\varepsilon_\square,\qquad
\eta^2|\nabla A|\le50(d-1)M\varepsilon_\square,
\end{equation*}
provided $2(d-1)s\,\varepsilon_\square\le1/10$. Two conditions remain: this proviso, and the
comparison of these bounds with the radius $O(1)M\alpha_0^{(99)}$ of \cite{B-CMP99b}. They are the
two members of \eqref{4.27:eq-real-point-collar-a}. Since $\mathsf{A}$ is the largest of the radii
as a function of $\mathsf{t}$, $\alpha_{0,0}\le\mathsf{A}(\mathsf{t}_0)$, and $\mathsf{t}_0$ lies in
the range $\mathsf{t}\ge\mathsf{t}_\gamma-c_\beta$ of the suprema in
\eqref{4.27:eq-real-point-collar-a}; it therefore suffices to bound the left sides by multiples of
$\mathsf{A}(\mathsf{t}_0)$. For the proviso,
\begin{equation*}
2(d-1)s\,\varepsilon_\square\le60M\,\varepsilon_\square<60M\cdot0.04\,\alpha_{0,0}
=2.4\,M\alpha_{0,0},
\end{equation*}
and $2.4\,M\alpha_{0,0}\le1/10$ follows from the first member, $24M\mathsf{A}\le1$. For the
comparison,
\begin{equation*}
\eta^2|\nabla A|\le50(d-1)M\cdot0.04\,\alpha_{0,0}=6M\alpha_{0,0},
\end{equation*}
and $6M\alpha_{0,0}\le M\alpha_0^{(99)}$ follows from the second member,
$6\mathsf{A}\le\alpha_0^{(99)}$. The
bound on $\eta|A|$ is $3M\alpha_{0,0}$ at most and is dominated in the same way. Hence the pair obeys
\cite[(3.35)]{B-CMP99b} on the cubes of index zero.
\end{proof}

\begin{theorem}[Each operation stays within its collar allowance]\label{4.27:thm-collar-allowance}
Assume the following.
\begin{itemize}
\item[(a)] The statements from which Lemma~\ref{4.27:lem-collar-displacement}, the displacement of
the collar quantities by one operation, is proved as an implication.
\item[(b)] The statements from which Lemma~\ref{4.27:lem-real-point-collar}, on the real
configuration on the stratum, is proved as an implication.
\item[(c)] From the inputs of the arrest theorem: every operation after the arrest that composes a
stored term with a map of $(\mathbb{U},\mathbb{J})$ is one of the operations O1, O2, O3, O4,
O4$'$, O5. Namely, it is one of the following: the operation $\mathbf{T}$ of the averaging step,
which composes with one map; one of the stages, each of which composes with one map; or one of
the four layers of the completion site.
\end{itemize}
Let $\mathfrak{O}$ be one of the operations O1--O5, O4$'$, with target date $d^t$ and source date
$d'$, where $d^t$ and $d'$ are consecutive dates. Let $k$ be the density of $\mathfrak{O}$. Let $\Phi$
be any member of the family of maps that $\mathfrak{O}$ composes with stored terms. Let
$\tilde{p}\in\mathfrak{G}(\mathfrak{B}_{d'})$ satisfy the collar clause
\eqref{4.27:eq-dated-allowances-a} at the date $d'$. Then the following hold.
\begin{itemize}
\item[(i)] \emph{Budget.} For every date $d$ we have $0\le\mathfrak{s}(d)<7\beta/4$, so that
$3.65<\mathsf{f}_0(d)\le 4$.
\item[(ii)] \emph{Room against spend.} For every $y\in\operatorname{St}(d')$ and every
$a\in\{1,3,6,7\}$,
\begin{equation*}
\bigl|Q_a(\Phi\tilde{p})(y)-Q_a(\tilde{p})(y)\bigr|\le\mathfrak{c}(d')\,r_a(0).
\end{equation*}
Moreover, $Q_5$ and the gauges on the cubes of index $0$ do not move under $\Phi$.
\item[(iii)] \emph{Replacement.} Points $y\in\operatorname{St}(d^t)\setminus\operatorname{St}(d')$
occur at a completion site only. At every such $y$ and for every entry $a$,
\begin{equation*}
Q_a(\Phi\tilde{p})(y)\le\tfrac{1}{16}\,r_a(0).
\end{equation*}
\item[(iv)] \emph{Real point.} For every date $d$ and every real $V$ in the support of the
density's characteristic functions, the real point
$(U(\mathfrak{B}_d,V),\mathcal{J}(U))$ of Lemma~\ref{4.27:lem-real-point-collar} satisfies
$Q_a(y)\le\frac12 r_a(0)$ for every entry $a$ and at every $y\in\operatorname{St}(d)$.
\end{itemize}
\end{theorem}

Here $\beta$ is the allowance constant of Definition~\ref{4.27:def-dated-allowances}, and
$\mathfrak{c}$, $\mathfrak{s}$, $\mathsf{f}_0=4-\mathfrak{s}$, the strata
$\operatorname{St}(d)=\Omega_0(\mathfrak{B}_d)\setminus\Omega_1(\mathfrak{B}_d)$ and the level-zero
radii $r_a(0)$ are also those of that definition. The entries $Q_a$,
$a\in\{1,3,5,6,7\}$, of a pair are those that enter the definition of the rider sets
$\mathfrak{G}(\mathfrak{B})$.

\begin{proof}
The four conclusions are proved in turn.

\emph{Budget.} This is Lemma~\ref{4.27:lem-allowance-sum}, the boundedness of the summed
allowances. For every date $d$ it gives $0\le\mathfrak{s}(d)<\frac74\beta=0.35$. Since
$\mathsf{f}_0(d)=4-\mathfrak{s}(d)$, we obtain $3.65<\mathsf{f}_0(d)\le4$.

\emph{Room against spend.} By hypothesis~(c), $\mathfrak{O}$ is one of three kinds of operation:
\begin{itemize}
\item the operation $\mathbf{T}$ of the averaging step, which composes with one map;
\item one of the stages, each of which composes with one map;
\item one of the layers of the completion site.
\end{itemize}
Accordingly, the source date $d'$ is either a stage $(k;l)$, or an averaging step $(k;\top)$, or a
completion site $(k;L)$. These are the three kinds of scaffold date of
Definition~\ref{4.27:def-dated-allowances}.

We apply Lemma~\ref{4.27:lem-collar-displacement} with the consecutive dates $d^t,d'$, with the
density $k$ and with $\mathfrak{B}=\mathfrak{B}_{d^t}$. In each of the three cases of
hypothesis~(c), the case headings of Lemma~\ref{4.27:lem-collar-displacement}, which assign the
operations to its cases, show that the maps with which $\mathfrak{O}$ composes stored terms belong
to the families of maps for which the lemma is stated in that case; so the lemma applies to
$\Phi$. Its hypotheses on the pair are two. The first is membership in the kernel class
$\mathfrak{K}_{109}(\mathfrak{B}_{d^t};\gamma_0^{(k)})$, where $\gamma_0^{(k)}$ is the dated size
parameter at the density $k$. Since $d^t<d'$ are consecutive dates, the order on the sets of dates
gives $\mathfrak{B}_{d^t}\preccurlyeq\mathfrak{B}_{d'}$, and the monotonicity of the rider sets then
gives
\begin{equation*}
\mathfrak{G}(\mathfrak{B}_{d'})\subset\mathfrak{G}(\mathfrak{B}_{d^t})
\subset\mathfrak{K}_{109}(\mathfrak{B}_{d^t};\gamma_0^{(k)});
\end{equation*}
hence $\tilde{p}$ lies in the kernel class. The second is the collar clause
\eqref{4.27:eq-dated-allowances-a} at the date $d'$, which is assumed. Write
$\Delta_a(y)=|Q_a(\Phi\tilde{p})(y)-Q_a(\tilde{p})(y)|$. The lemma bounds $\Delta_a(y)$ at
$y\in\operatorname{St}(d')$ as follows.
\begin{itemize}
\item If $d'=(k;l)$ is a stage, display \eqref{4.27:eq-collar-displacement-b} gives
$\Delta_a(y)\le\frac14\mathfrak{c}(d')\,r_a(0)$. It covers O1, O4 and O4$'$ alike.
\item If $d'=(k;\top)$ is an averaging step, display \eqref{4.27:eq-collar-displacement-c} gives
$\Delta_a(y)\le\frac18\mathfrak{c}(d')\,r_a(0)$.
\item If $d'=(k;L)$ is a completion site, display \eqref{4.27:eq-collar-displacement-d} gives
$\Delta_a(y)\le\frac{\theta_S}{4}\mathfrak{c}(d')\,r_a(0)$ for the points of the stratum of the
source date, where $\theta_S$ is the constant of that display of
Lemma~\ref{4.27:lem-collar-displacement}, which satisfies $\theta_S\le4$ there, so that
$\theta_S/4\le1$.
\end{itemize}
The stages are the operations O1, O4 and O4$'$; by hypothesis~(c) the remaining operations O2, O3
and O5 are therefore the operation $\mathbf{T}$ of the averaging step and the layers of the
completion site, so the three cases exhaust the six operations. Since each of the factors
$\frac14$, $\frac18$ and $\theta_S/4$ is at most $1$, in each case
$\Delta_a(y)\le\mathfrak{c}(d')\,r_a(0)$ for $a\in\{1,3,6,7\}$. The same lemma gives
$\Delta_5=0$, and it shows that the gauges on the cubes of index $0$ are not moved by $\Phi$. This
is the second conclusion.

\emph{Replacement.} By Lemma~\ref{4.27:lem-collar-displacement}, the set
$\operatorname{St}(d^t)\setminus\operatorname{St}(d')$ can be non-empty only when the source date is
a completion site $(k;L)$, and its points then lie in the collars $\operatorname{Col}(C)$ of the
completing components $C$. At such points, display \eqref{4.27:eq-collar-displacement-e} of
Lemma~\ref{4.27:lem-collar-displacement} gives $Q_a(\Phi\tilde{p})(y)\le\frac1{16}r_a(0)$ for every
entry $a$.

\emph{Real point.} Let $d$ be a date, and apply Lemma~\ref{4.27:lem-real-point-collar} with
$\mathfrak{B}=\mathfrak{B}_d$; the statements from which it is proved are assumed in
hypothesis~(b). At every
$y\in\operatorname{St}(d)$, the pair $(U(\mathfrak{B}_d,V),\mathcal{J}(U))$, with $V$ real in the
support of the density's characteristic functions, satisfies the following:
\begin{align*}
&Q_6(y)=Q_7(y)=0,\\
&Q_1(y)=Q_5(y)<0.04\,r_1(0),\\
&Q_3(y)\le 0.48\,r_3(0).
\end{align*}
The radii $r_1(0)$ and $r_5(0)$ are both the level-zero radius $\alpha_{0,0}$
(Definition~\ref{4.27:def-dated-allowances}). Hence $Q_5(y)<0.04\,r_5(0)$ as well. Since
$0<0.04<0.48<\frac12$, every entry obeys $Q_a(y)\le\frac12r_a(0)$ at every
$y\in\operatorname{St}(d)$.
\end{proof}

\begin{theorem}[Rigidity of the arrest theorem on the never-averaged collar]\label{4.27:thm-collar-rigidity}
Let the standing conventions of the arrest theorem be in force. The rider set
$\mathfrak{G}(\mathfrak{B}_d)$ of a date $d$ is read as in
Definition~\ref{4.27:def-dated-allowances}: its members are pairs $(e^{i\eta A'}U,\mathbb{J})$ on
the whole of $\Omega_0(\mathfrak{B}_d)$; on $\Omega_1(\mathfrak{B}_d)$ they obey the weighted
clauses of the arrest theorem, and on the never-averaged collar
\[
\operatorname{St}(d)=\Omega_0(\mathfrak{B}_d)\setminus\Omega_1(\mathfrak{B}_d)
\]
they obey the level-zero clause \eqref{4.27:eq-dated-allowances-a} at the date $d$, together with
the index-zero cube clause. Assume moreover:
\begin{enumerate}
\item the standing hypothesis of the arrest theorem and the inputs listed with it;
\item the level-zero statements on the currents of the completion site, among them the
index-zero cube clause on the collar $\operatorname{Col}(C)$ of a completing component $C$;
\item the level-zero small-field statement, with the floor $\mathsf{a}=A_0/A_1\ge c(4,L)$ of
Lemma~\ref{4.27:lem-real-point-collar}, where $c(4,L)$ depends only on the dimension $4$ and on
$L$;
\item the conditions \eqref{4.27:eq-collar-displacement-a} and
\eqref{4.27:eq-real-point-collar-a}, and the level-zero conditions on the response field
$\mathbb{H}$ and on the collars;
\item from the arrest theorem: its ordering lemma on dates, clauses (e1) and (e2) of its rider
lemma, its nesting statement, its index-zero statements $(\alpha)$ and $(\beta)$, its separation
statement, and its input listing the operations that follow the arrest.
\end{enumerate}
Write $\ell_d$ and $w_d$ for the level and weight functions of $\mathfrak{B}_d$, so that
$\theta^{\mathfrak{B}_d}_a=w_d\,r_a(\ell_d)$. Write $\Theta^d_a$ for the dated threshold of
Definition~\ref{4.27:def-dated-allowances}, equal to $5\theta^{\mathfrak{B}_d}_a$ on
$\Omega_1(\mathfrak{B}_d)$ and to $\mathsf{f}_0(d)\,r_a(0)$ on $\operatorname{St}(d)$, for
$a\in\{1,3,5,6,7\}$. Let $\gamma_0^{(k)}$ be the dated size parameter of the arrest theorem and
$\mathfrak{K}_{109}(\mathfrak{B};\gamma_0)$ the kernel class. Let $a_\flat$ be the size quantity
read on the balls of the small-field bump estimate; at a point of $\operatorname{St}(d)$ its four
members are $Q_5$, $Q_1$, $\eta Q_6$ and $\eta^2Q_7$, each in its own units. Then for all dates
$d\le d'$, of densities $k\le k'$ respectively,
\begin{equation}\label{4.27:eq-collar-rigidity-a}
\begin{aligned}
&\text{(e1)}'\quad \ell_{d'}\ge\ell_d,\quad
 w_{d'}\le\mathsf{L}_0^{2(\ell_{d'}-\ell_d)}\,w_d\ \text{ on }\Omega_0(\mathfrak{B}_d),
 \qquad \operatorname{St}(d')\subset\operatorname{St}(d);\\
&\text{(e2)}'\quad \Theta^{d'}_a\le\Theta^{d}_a\ \text{ on }\Omega_0(\mathfrak{B}_d);\\
&\text{(e3)}'\quad \mathfrak{G}(\mathfrak{B}_{d'})\subset\mathfrak{G}(\mathfrak{B}_d)
 \subset\mathfrak{K}_{109}(\mathfrak{B}_d;\gamma_0^{(k)}),\qquad
 a_\flat\le\gamma_0^{(k)}\ \text{ on }\mathfrak{G}(\mathfrak{B}_d);\\
&\text{(e4)}'\quad \bigl(U(\mathfrak{B}_d,V),\mathcal{J}(U)\bigr)\in\mathfrak{G}(\mathfrak{B}_d),
 \qquad \Phi\,\mathfrak{G}(\mathfrak{B}_{d'})\subset\mathfrak{G}(\mathfrak{B}_{d^t})
 \ \text{ at the stratum},
\end{aligned}
\end{equation}
in the following sense. In (e3)$'$ both inclusions hold as inclusions of sets of pairs on the
whole of $\Omega_0(\mathfrak{B}_d)$, and the bound on $a_\flat$ holds at every point,
$\operatorname{St}(d)$ included. In the first half of (e4)$'$, $V$ is real and lies in the
support of the characteristic functions of the density, as in
Lemma~\ref{4.27:lem-real-point-collar}; the real pair lies in $\mathfrak{G}(\mathfrak{B}_d)$ with a
margin of at least $3.1\,r_a(0)$ below $\Theta^d_a$ at every point of $\operatorname{St}(d)$. In
the second half of (e4)$'$, $\mathfrak{O}$ is an operation with target date $d^t$ and source date
$d'$, consecutive dates with $d^t<d'$, and $\Phi$ is any member of the families of
Lemma~\ref{4.27:lem-collar-displacement}; for every $\tilde p\in\mathfrak{G}(\mathfrak{B}_{d'})$
the global pair $\Phi\tilde p$ obeys \eqref{4.27:eq-dated-allowances-a} at the date $d^t$ at
every $y\in\operatorname{St}(d^t)$, together with the index-zero cube clause.

Consequently, the following four statements hold as stated with the arrest theorem, on the
whole of $\Omega_0$ of the determining sets concerned, points of the never-averaged collar
included:
\begin{itemize}
\item the rider clause of the arrest theorem;
\item clause (ii) of its proposition on terms born earlier, by which a term born earlier keeps
its rider at its birth set and every later source lies in it;
\item its statement that a stored term reads the rider set and the size parameter of its birth
set and birth density only;
\item the ambient kernel-class condition of the flat spaces, read at the determining set over
which that condition is stated.
\end{itemize}
\end{theorem}

\begin{proof}
Throughout, dates are the dates of Definition~\ref{4.27:def-dated-allowances}, in the order
fixed there. For a determining set, $\Omega_m$ denotes its domain of level $m$. We write
$\Lambda_k$ for the small-field region at level $k$ and $Z_k=\Lambda_k^c$; for a subset $S$ of the
lattice, $S^\sim$ denotes its enlargement by one layer of $LM\check{R}_1$-cubes. The four clauses
are proved in the order (e1)$'$, (e2)$'$, (e3)$'$, (e4)$'$, after a geometric step $(\alpha)$.

\emph{Step $(\alpha)$: geometry of the collar.} We use the index-zero statement $(\alpha)$ of the
arrest theorem, read with the level-zero conventions \eqref{4.27:eq-collar-conventions-a}. No
transformation of the determining sets lowers a level. Of the set transformations (t1)--(t4) of
the arrest theorem:
\begin{itemize}
\item (t1) and (t2) act inside $\Omega_1$;
\item (t3) changes no set;
\item (t4) replaces $\Omega_m$ by $\Omega_m\cup Z_L$ for every $m\le k$, where $Z_L$ is a union
of components of $Z_k=\Lambda_k^c$.
\end{itemize}
In particular every domain $\Omega_m$ can only grow along the dates.

The one-layer enlargement $Z_L^\sim$ stays in $Z_L\cup\Lambda_k\subset Z_L\cup\Omega_1$. Indeed,
by the separation statement of the arrest theorem, distinct components of $Z_k$ are at a mutual
distance that exceeds the width of one layer of $LM\check{R}_1$-cubes. Since enlargement commutes
with unions,
$(\Omega_1\cup Z_L)^\sim=\Omega_1^\sim\cup Z_L^\sim$. Removing $\Omega_1\cup Z_L$ therefore
removes all of $Z_L^\sim$. Together with $\Omega_1^\sim\subset(\Omega_1\cup Z_L)^\sim$ this gives
\begin{equation}\label{4.27:eq-collar-rigidity-1}
(\Omega_1\cup Z_L)^\sim\setminus(\Omega_1\cup Z_L)
=\Omega_1^\sim\setminus(\Omega_1\cup Z_L)
=(\Omega_1^\sim\setminus\Omega_1)\setminus Z_L .
\end{equation}
Hence $\operatorname{St}(d')$ is $\operatorname{St}(d)$ with the components completed at the
dates of $]d,d']$ removed. In particular $\operatorname{St}(d')\subset\operatorname{St}(d)$, and
$\operatorname{St}(d)\subset\Omega_1^\sim\setminus\Omega_1$, with $\Omega_1$ the set made by the
first averaging step $\mathbf{T}$. A point of $\operatorname{St}(d)$ that does not lie in
$\operatorname{St}(d')$ lies in a component completed at a date of $]d,d']$. By (t4) it then lies
in $\Omega_1(\mathfrak{B}_{d'})$. Thus
\begin{equation}\label{4.27:eq-collar-rigidity-2}
\operatorname{St}(d)\subset\operatorname{St}(d')\cup\Omega_1(\mathfrak{B}_{d'}),\qquad
\Omega_0(\mathfrak{B}_d)\subset\Omega_0(\mathfrak{B}_{d'}).
\end{equation}

\emph{Proof of (e1)$'$.} By the conventions \eqref{4.27:eq-collar-conventions-a},
$(\ell,w)=(0,1)$ on the stratum. Consider first a point that stays in the stratum up to $d'$.
There $(\ell_{d'},w_{d'})=(\ell_d,w_d)=(0,1)$, and (e1)$'$ holds with equality. By Step
$(\alpha)$, a point leaves the stratum only through (t4), at some date $d_e\in\,]d,d']$. There it
receives $(\ell,w)=(k_e,1)$, and from $d_e$ on clause (e1) of the rider lemma of the arrest theorem
applies; this gives (e1)$'$ at such a point. On $\Omega_1(\mathfrak{B}_d)$, (e1)$'$ is clause (e1)
of the rider lemma. The inclusion $\operatorname{St}(d')\subset\operatorname{St}(d)$ was shown in
Step $(\alpha)$.

\emph{Proof of (e2)$'$.} By \eqref{4.27:eq-collar-rigidity-2} there are three cases.
\begin{itemize}
\item On $\Omega_1(\mathfrak{B}_d)$ the inequality is clause (e2) of the rider lemma, multiplied
by $5$.
\item On $\operatorname{St}(d)\cap\operatorname{St}(d')$ we have
$\Theta^{d'}_a=\mathsf{f}_0(d')r_a(0)$ and $\Theta^d_a=\mathsf{f}_0(d)r_a(0)$. By
Lemma~\ref{4.27:lem-allowance-sum}, $\mathsf{f}_0$ is non-increasing, so
$\mathsf{f}_0(d')\le\mathsf{f}_0(d)$.
\item Let $y\in\operatorname{St}(d)\cap\Omega_1(\mathfrak{B}_{d'})$. Then $y$ entered
$\Omega_1$ at a date $d_e\in\,]d,d']$ with $(\ell,w)=(k_e,1)$ and $k_e\ge2$. Hence
\begin{equation}\label{4.27:eq-collar-rigidity-3}
\begin{aligned}
\Theta^{d'}_a(y)\le\Theta^{d_e}_a(y)=5\,r_a(k_e)
&\le5\cdot0.164\,r_a(0)=0.82\,r_a(0)\\
&<3.65\,r_a(0)<\Theta^d_a(y).
\end{aligned}
\end{equation}
The first inequality is clause (e2) of the rider lemma, applied from $d_e$ on. The equality is
$\Theta^{d_e}_a=5w\,r_a(\ell)$ at $(\ell,w)=(k_e,1)$. The bound $r_a(k_e)\le0.164\,r_a(0)$ is
clause (a) of Lemma~\ref{4.27:lem-level-zero-ratios}, since $k_e\ge1$. The last inequality is
$\Theta^d_a(y)=\mathsf{f}_0(d)r_a(0)$ with $\mathsf{f}_0(d)>3.65$
(Lemma~\ref{4.27:lem-allowance-sum}).
\end{itemize}

\emph{Proof of (e3)$'$, first inclusion.} Let $\tilde p\in\mathfrak{G}(\mathfrak{B}_{d'})$. By
\eqref{4.27:eq-collar-rigidity-2} it is defined on $\Omega_0(\mathfrak{B}_d)$. We check the
clauses of $\mathfrak{G}(\mathfrak{B}_d)$ in three parts.
\begin{itemize}
\item The entries. The $Q_a(y)$ are quantities of the pair at $y$. The families of bonds and
plaquettes read at $y$ at level zero (Notation~\ref{4.27:not-collar-conventions}) are
sub-families of those read at $y$ by a clause of level $\ge1$. So the bound
$Q_a(y)<\Theta^{d'}_a(y)$ given by $\mathfrak{G}(\mathfrak{B}_{d'})$ controls the entries asked
for by $\mathfrak{G}(\mathfrak{B}_d)$ at $y$, and (e2)$'$ turns it into $Q_a(y)<\Theta^d_a(y)$.
\item Cubes of level $\ge1$. The cube clauses follow from the nesting statement of the arrest
theorem.
\item Index-zero cubes. We use the index-zero statement $(\beta)$ of the arrest theorem. The
index-zero clause involves neither $\gamma_0^{(k)}$ nor the date. At cubes of index zero at both
dates it is therefore the same clause at $d$ and at $d'$. At cubes that entered a higher level
between $d$ and $d'$ it follows from clause (b) of the nesting statement and from the
restriction of the gauge of a cube of the top level.
\end{itemize}

\emph{Proof of (e3)$'$, second inclusion.} On $\Omega_1(\mathfrak{B}_d)$ this is clause (e3) of
the rider lemma. On $\operatorname{St}(d)$ we check the three conditions of the kernel class.

\emph{The current.} At a point of $\operatorname{St}(d)$, with $\beta_0$ the constant appearing
in clause (a) of Lemma~\ref{4.27:lem-level-zero-ratios},
\begin{equation}\label{4.27:eq-collar-rigidity-4}
\begin{aligned}
\eta^3|\mathbb{J}|<4\alpha_{0,0}\le4\mathsf{A}(\mathsf{t}_0)
&\le4(1+\beta_0)\mathsf{A}(\mathsf{t}_k)\le\tfrac{32}{7}\,\bar{\alpha}_k\\
&<5\,w_*(k)\,\bar{\alpha}_k\le\gamma_0^{(k)} .
\end{aligned}
\end{equation}
The steps are justified as follows.
\begin{itemize}
\item The first inequality is the current clause of $\mathfrak{G}(\mathfrak{B}_d)$ at a point of
$\operatorname{St}(d)$, whose factor $\mathsf{f}_0(d)$ is at most $4$
(Lemma~\ref{4.27:lem-allowance-sum}).
\item The second holds because $\mathsf{A}$ is the maximum of the radii.
\item The third is clause (a) of Lemma~\ref{4.27:lem-level-zero-ratios}.
\item The fourth uses that $\bar{\alpha}_k=\max_{m\le k}\mathsf{A}(\mathsf{t}_m)$ is a running
maximum, so that $\mathsf{A}(\mathsf{t}_k)\le\bar{\alpha}_k$, together with the bound
$1+\beta_0\le\frac87$ on the constant $\beta_0$.
\item The fifth uses the bound $w_*(k)>\frac{32}{35}$ on the weight
$w_*(k)=\mathsf{L}_0^{2N_0(k)}$.
\item The last is the definition of the dated size parameter,
$\gamma_0^{(k)}=5\max_{j\le k}w_*(j)\bar{\alpha}_j$.
\end{itemize}
If the maximum defining $\bar{\alpha}_k$ includes $m=0$, the third step is not needed, since
then $4\mathsf{A}(\mathsf{t}_0)\le4\bar{\alpha}_k$ directly.

\emph{The condition \cite[(3.37)]{B-CMP99b} at $j=0$.} At a point of $\operatorname{St}(d)$ the
clause of $\mathfrak{G}(\mathfrak{B}_d)$ gives $\eta|A'|<4\alpha_{1,0}$ and
$\eta^2|\nabla_UA'|<4\alpha_{1,0}$. Moreover
$4\alpha_{1,0}\le4\mathsf{A}(\mathsf{t}_0)<\gamma_0^{(k)}$ by the chain
\eqref{4.27:eq-collar-rigidity-4}. In \cite[(3.37)]{B-CMP99b} one radius $\alpha_1$ serves all
indices $j=0,\dots,k$. The inequality $\gamma_0^{(k)}\le\alpha_1$, which the arrest theorem uses
at the levels $\ge1$ to obtain \cite[(3.37)]{B-CMP99b}, is therefore the inequality needed at
$j=0$. The radius in \cite[(3.37)]{B-CMP99b} carries no index $j$.

\emph{The condition \cite[(3.35)]{B-CMP99b} at $j=0$.} It is part of the index-zero cube clause
of $\mathfrak{G}(\mathfrak{B}_d)$.

\emph{The bound on $a_\flat$.} At a point of the stratum the four members of $a_\flat$ are
$Q_5$, $Q_1$, $\eta Q_6$ and $\eta^2Q_7$, each read in its own units. By the clause of
$\mathfrak{G}(\mathfrak{B}_d)$ there they are smaller than $4\mathsf{A}(\mathsf{t}_0)$. By the
chain \eqref{4.27:eq-collar-rigidity-4}, $4\mathsf{A}(\mathsf{t}_0)<\gamma_0^{(k)}$. On
$\Omega_1(\mathfrak{B}_d)$ the bound is clause (e3) of the rider lemma.

\emph{Proof of (e4)$'$: the real point.} By Lemma~\ref{4.27:lem-real-point-collar}, at every
$y\in\operatorname{St}(d)$ the real pair has
\[
Q_6=Q_7=0,\qquad Q_1=Q_5<0.04\,r_1(0),\qquad Q_3\le0.48\,r_3(0),
\]
and it obeys \cite[(3.35)]{B-CMP99b} on the index-zero cubes. Hence $Q_a(y)\le\frac12r_a(0)$.
Since $\frac12<3.65<\mathsf{f}_0(d)$ by Lemma~\ref{4.27:lem-allowance-sum}, the clause
\eqref{4.27:eq-dated-allowances-a} holds with a margin of at least
$(3.65-\frac12)\,r_a(0)>3.1\,r_a(0)$. On $\Omega_1(\mathfrak{B}_d)$ the real point lies in the
rider set by clause (e4) of the rider lemma.

The first clause of (e4) of the rider lemma states that every pair obeying, at every point, the
direct clauses of a space of the arrest theorem over $\mathfrak{B}$ lies in
$\mathfrak{G}(\mathfrak{B})$. This is a statement on $\Omega_1(\mathfrak{B})$. At level zero, a
space over $\mathfrak{B}_d$ has exactly one clause, namely that of $\mathfrak{G}(\mathfrak{B}_d)$
itself, by the ambient rider convention among the standing conventions of the arrest theorem.
So there is nothing to show. Moreover, a global pair with $Q_a\le3.65\,r_a(0)$ on the stratum
meets the level-zero clause of every date, because $\mathsf{f}_0>3.65$ at every date
(Lemma~\ref{4.27:lem-allowance-sum}).

\emph{Proof of (e4)$'$: the images.} Let $\mathfrak{O}$ have target date $d^t$ and source date
$d'$, consecutive dates with $d^t<d'$. Let $\tilde p\in\mathfrak{G}(\mathfrak{B}_{d'})$. By
(e3)$'$,
\[
\tilde p\in\mathfrak{G}(\mathfrak{B}_{d'})\subset\mathfrak{G}(\mathfrak{B}_{d^t})
\subset\mathfrak{K}_{109}(\mathfrak{B}_{d^t};\gamma_0^{(k)}).
\]
Hence $\tilde p$ satisfies the hypothesis of Lemma~\ref{4.27:lem-collar-displacement}, and
$\Delta_a(y)=|Q_a(\Phi\tilde p)(y)-Q_a(\tilde p)(y)|$ is controlled by it; since the proof of
(e3)$'$ above used none of the displacement bounds of that lemma, no circularity arises. By
(e1)$'$, $\operatorname{St}(d')\subset\operatorname{St}(d^t)$; we treat the two parts of
$\operatorname{St}(d^t)$ separately.

\emph{Points of $\operatorname{St}(d')$.} Let $y\in\operatorname{St}(d')$ and $a\in\{1,3,6,7\}$.
Then
\begin{equation}\label{4.27:eq-collar-rigidity-5}
\begin{aligned}
Q_a(\Phi\tilde p)(y)\le Q_a(\tilde p)(y)+\Delta_a(y)
&<\mathsf{f}_0(d')\,r_a(0)+\mathfrak{c}(d')\,r_a(0)\\
&=\mathsf{f}_0(d^t)\,r_a(0).
\end{aligned}
\end{equation}
Here the bound on $Q_a(\tilde p)(y)$ is the clause of $\mathfrak{G}(\mathfrak{B}_{d'})$ at $y$.
The bound $\Delta_a(y)\le\mathfrak{c}(d')r_a(0)$ is, according to the operation,
\eqref{4.27:eq-collar-displacement-b}, \eqref{4.27:eq-collar-displacement-c} or
\eqref{4.27:eq-collar-displacement-d} of Lemma~\ref{4.27:lem-collar-displacement}. The equality
is $\mathsf{f}_0(d^t)-\mathsf{f}_0(d')=\mathfrak{c}(d')$ for consecutive dates
(Lemma~\ref{4.27:lem-allowance-sum}). For $a=5$ the entry is not moved by $\Phi$, and
$\mathsf{f}_0(d')\le\mathsf{f}_0(d^t)$.

\emph{Points of $\operatorname{St}(d^t)\setminus\operatorname{St}(d')$.} By
\eqref{4.27:eq-collar-displacement-e} of Lemma~\ref{4.27:lem-collar-displacement},
$Q_a(\Phi\tilde p)(y)\le\frac1{16}r_a(0)$, which is smaller than $\mathsf{f}_0(d^t)r_a(0)$ by
Lemma~\ref{4.27:lem-allowance-sum}.

\emph{The index-zero cubes.} On $\operatorname{St}(d')$ the group-valued part and its gauges are
kept, by Lemma~\ref{4.27:lem-left-factor-algebra} and the fourth input of the arrest theorem. On
the collar $\operatorname{Col}(C)$ of a completing component the group-valued part is the
configuration $U_2$ of the completion site. There the index-zero cube clause is clause (b) of
the level-zero current proposition, which is among the level-zero statements on the currents of
the completion site assumed in the theorem.

\emph{No other operation occurs.} By the inputs of the arrest theorem, every operation after the
arrest that composes a stored term with a map of $(\mathbb{U},\mathbb{J})$ is the step
$\mathbf{T}$, a stage, or a layer of a completion site. These are the operations treated above.
At the levels $\ge1$, (e4) is clause (e4) of the rider lemma and is not affected.

\emph{The global image is a continuation.} No truncation of the current is needed and no part of
the image is discarded. The bounds of Lemma~\ref{4.27:lem-collar-displacement} hold at every
point of the stratum, on or off the domain $Y'$ of the source of $\mathfrak{O}$, so the global
image is itself a continuation on the stratum. The continuation used by the arrest theorem at
the levels $\ge1$ is built as follows: on $Y'$ it is the image; off $Y'$ it is the input's
$\mathbb{U}$, which there equals the image's, together with the input's current $\mathbb{J}$.
This continuation also obeys \eqref{4.27:eq-dated-allowances-a} at the date $d^t$. Indeed, the
entry $a=3$ is read bond by bond. At each stratum bond it is either that of the image, bounded
by Lemma~\ref{4.27:lem-collar-displacement}, or that of the input, smaller than
$\mathsf{f}_0(d')r_3(0)\le\mathsf{f}_0(d^t)r_3(0)$. Thus one global pair serves both halves.

\emph{The rider clause of the arrest theorem.} Let $\tau$ be a stored term, born over
$\mathfrak{B}_\tau=\mathfrak{B}_{d_\tau}$ at density $k_\tau$. The arrest theorem proves the rider
clause in two steps: a continuation of the image lies in the rider set of the target, and
(e3) is then applied. Both steps now hold on all of $\Omega_0$. By (e4)$'$ the continuation of
the image lies in $\mathfrak{G}(\mathfrak{B}_{d^t})$, at the stratum as well. Since $\tau$ was
stored before $\mathfrak{O}$, clause (5) of the ordering lemma of the arrest theorem gives
$d_\tau\le d^t$. Applying (e3)$'$ with $d=d_\tau$ and $d'=d^t$ gives
\[
\mathfrak{G}(\mathfrak{B}_{d^t})\subset\mathfrak{G}(\mathfrak{B}_{d_\tau})
\subset\mathfrak{K}_{109}(\mathfrak{B}_{d_\tau};\gamma_0^{(k_\tau)})
\]
on the whole of $\Omega_0(\mathfrak{B}_{d_\tau})$.

\emph{The ambient kernel-class condition.} It is the rider clause at the determining set over
which that condition is stated, together with the second inclusion of (e3)$'$, on the whole of
the region on which the condition is read, its stratum points included.

\emph{Clause (ii) of the proposition on terms born earlier, and the statement on stored sizes.}
Both follow from (e3)$'$, whose inclusions now hold on the whole of $\Omega_0$.

\emph{Clause (C) of the arrest theorem at a stratum point.} Clause (C) states that old terms
composed with the map of $\mathfrak{O}$ are holomorphic on the source of $\mathfrak{O}$, by
clause (A) and the rider clause. Let $X_\tau$ be the domain of $\tau$ and
$x\in X_\tau\cap\operatorname{St}(d_\tau)$. At $x$ the only clause of $\tau$ is
\eqref{4.27:eq-dated-allowances-a} at the date $d_\tau$. Hence clause (A) at $x$ is the rider
clause at $x$, which has been established.

\emph{Order of the argument.} The proof is an induction along the dates. Clause (e3)$'$ uses only
arithmetic and the geometry of Step $(\alpha)$. Clause (e4)$'$ at $\mathfrak{O}$ uses only (e3)$'$
and the inputs at $\mathfrak{O}$. No step of the argument uses its own conclusion.
\end{proof}

\begin{remark}[Allowances and displacements on the collar, tabulated]\label{4.27:rem-collar-table}
Let $\mathfrak{O}$ be an operation with target date $d^t$ and source date $d'$, consecutive
scaffold dates in the sense of Definition~\ref{4.27:def-dated-allowances}. Let the pair
$\tilde p=(\mathbb{U},\mathbb{J})$ and the map $\Phi$ be as in
Lemma~\ref{4.27:lem-collar-displacement}, and put
$\Delta_a(y)=|Q_a(\Phi\tilde p)(y)-Q_a(\tilde p)(y)|$. Let $\theta_S$, $\Lambda_{\gamma}$ and
$\bar{P}^{0}$ be the constants of Lemma~\ref{4.27:lem-collar-displacement}.

The \emph{room} of $\mathfrak{O}$ is
$\mathsf{f}_0(d^t)-\mathsf{f}_0(d')\ge\mathfrak{c}(d')$.
In the first three rows of the table below:
\begin{itemize}
\item the spend bounds $\Delta_a(y)$ at $y\in\operatorname{St}(d')$ for $a\in\{1,3,6,7\}$;
\item there $\Delta_5=0$ and the cube gauges are kept;
\item the ratio column gives a number $q$ with
\begin{equation*}
\Delta_a(y)\le q\,\bigl(\mathsf{f}_0(d^t)-\mathsf{f}_0(d')\bigr)\,r_a(0),
\end{equation*}
the displacement measured against the room.
\end{itemize}
In the last row:
\begin{itemize}
\item the point $y$ ranges over $\operatorname{St}(d^t)\setminus\operatorname{St}(d')$, the
union of the collars $\operatorname{Col}(C)$ of the completing components $C$;
\item the room entry is the threshold factor $\mathsf{f}_0(d^t)$;
\item the spend bounds $Q_a(y)$ of the presented pairs of
\eqref{4.27:eq-collar-displacement-e}, $a\in\{1,3,5,6,7\}$, under the condition stated there;
\item the ratio is the spend over the room.
\end{itemize}
All entries are in units of $r_a(0)$. The third row holds under the condition
$16\bar{P}^{0}\le\beta\theta_S$ stated with \eqref{4.27:eq-collar-displacement-d}.
\begin{center}
\begin{tabular}{|l|c|c|c|}
\hline
source date $d'$ of $\mathfrak{O}$ & room $\ge$ & spend $\le$ & ratio $\le$\\
\hline
$(k;l)$, a stage & $\beta 2^{-(h_k+l)}$ & $\frac14\beta 2^{-(h_k+l)}$ & $\frac14$\\
$(k'+1;\top)$, averaging $k'\to k'+1$ & $\beta 2^{-(k'+1)}$
 & $\Lambda_{\gamma}\cdot\frac14(1-\theta_S)\beta 2^{-(k'+1)}$ & $\frac18$\\
$(k;L)$, off the completing components & $\frac14\beta 2^{-(k-1)}$ & $\bar{P}^{0}2^{-(k-1)}$
 & $\frac{\theta_S}{4}$\\
$(k;L)$, on $\operatorname{Col}(C)$ (replacement) & $\mathsf{f}_0(d^t)>3.65$ & $\frac1{16}$
 & $\frac1{58}$\\
\hline
\end{tabular}
\end{center}
\end{remark}

\begin{proof}
By Lemma~\ref{4.27:lem-allowance-sum}, the consecutive dates of $\mathfrak{O}$ satisfy
$\mathsf{f}_0(d^t)-\mathsf{f}_0(d')=\mathfrak{c}(d')$. So in each of the first three rows the
room is the allowance $\mathfrak{c}(d')$ of Definition~\ref{4.27:def-dated-allowances}. We take
the rows in turn.

\emph{Stage.} Here $d'=(k;l)$, and
$\mathfrak{c}(k;l)=\rho^{(l),(k)}_0=\beta 2^{-(h_k+l)}$ by definition. Display
\eqref{4.27:eq-collar-displacement-b} of Lemma~\ref{4.27:lem-collar-displacement} gives
$\Delta_a(y)\le\frac14\mathfrak{c}(d')r_a(0)$ for every $y\in\operatorname{St}(d')$ and
$a\in\{1,3,6,7\}$. It also gives $\Delta_5=0$ and that the cube gauges are kept. This is the
spend $\frac14\beta2^{-(h_k+l)}$ and the ratio $\frac14$.

\emph{Averaging step.} Here $d'=(k'+1;\top)$ with $k'\ge1$. Since $k'+1\ge2$, the allowance
is $\mathfrak{c}(d')=\beta2^{-(k'+1)}$. Display
\eqref{4.27:eq-collar-displacement-c} of Lemma~\ref{4.27:lem-collar-displacement} gives
\begin{equation*}
\Delta_a(y)\le\Lambda_{\gamma}\cdot\tfrac14(1-\theta_S)\beta2^{-(k'+1)}r_a(0)
\le\tfrac18\mathfrak{c}(d')r_a(0).
\end{equation*}
It also gives $\Delta_5=0$ and that the cube gauges are kept. The spend column records this
bound, and the ratio $\frac18$ is its final inequality.

\emph{Completion site, off the completing components.} Here $d'=(k;L)$, and
$\mathfrak{c}(k;L)=\frac14\beta2^{-(k-1)}$. Display
\eqref{4.27:eq-collar-displacement-d} of Lemma~\ref{4.27:lem-collar-displacement} gives
$\Delta_a(y)\le\bar{P}^{0}2^{-(k-1)}r_a(0)$, with
$\Delta_5=0$ and the cube gauges kept. Under $16\bar{P}^{0}\le\beta\theta_S$,
\begin{equation*}
\bar{P}^{0}2^{-(k-1)}\le\tfrac{\beta\theta_S}{16}2^{-(k-1)}
=\tfrac{\theta_S}{4}\cdot\tfrac14\beta2^{-(k-1)}=\tfrac{\theta_S}{4}\mathfrak{c}(d').
\end{equation*}
This gives the ratio $\theta_S/4$.

\emph{Completion site, on the collars.} At
$y\in\operatorname{St}(d^t)\setminus\operatorname{St}(d')$ the value is replaced rather than
displaced. It is therefore compared with the threshold $\mathsf{f}_0(d^t)r_a(0)$, and
$\mathsf{f}_0(d^t)>3.65$ by Lemma~\ref{4.27:lem-allowance-sum}. Display
\eqref{4.27:eq-collar-displacement-e} of Lemma~\ref{4.27:lem-collar-displacement} gives
$Q_a(y)\le\frac1{16}r_a(0)$ for
$a\in\{1,3,5,6,7\}$. This holds at the presented pairs and under the condition stated there.
Since $16\cdot3.65=58.4$,
\begin{equation*}
\frac{1/16}{\mathsf{f}_0(d^t)}<\frac{1}{58.4}<\frac1{58}. \qedhere
\end{equation*}
\end{proof}

\begin{remark}[The one-sided conditions added on the collar and their order]
\label{4.27:rem-collar-conditions}
The conditions on the parameters
that enter Theorem~\ref{4.27:thm-collar-rigidity} and
Lemmas~\ref{4.27:lem-collar-displacement} and~\ref{4.27:lem-real-point-collar} are listed
below; these statements carry further hypotheses, which are stated with them. The conditions in
(b) and (c) are one-sided. Those in (a) are read unchanged, in the form in which they stand
among the standing hypotheses of the arrest theorem.
\begin{enumerate}
\item[(a)] \emph{Read unchanged from the standing hypotheses of the arrest theorem:} $L\ge 8$;
the two inequalities
\begin{equation*}
\delta\mathsf{w}_M\ge 8,\qquad \tfrac14\,\delta\mathsf{w}_M\ge 4\ln L+1,
\end{equation*}
where $\mathsf{w}_M$ is the per-stratum contour price,
both implied by the level condition of those hypotheses; $\beta=\frac15$, where $\beta$ is the
constant of the allowance profiles; $\beta_0\le\frac17$;
$\theta_S\in[\frac12,\frac89]$; the floor at the value $\max\{C_H^{\sharp},1\}$ that those
hypotheses impose, $C_H^{\sharp}$ being a constant of the $\mathbf{T}$-step; the real-point
condition with factor $\varphi_{\mathrm{ARR}}$; the condition with fraction
$\mathsf{m}\in\,]0,\frac12]$;
the monotonicity condition in $\gamma$; the list of $\gamma$-ceilings of the propagator
statements, together with the closures of the second family of maps $\Phi^{\chi}$ of
Lemma~\ref{4.27:lem-collar-displacement}; the window condition; the arrest-range condition;
$\Lambda_{\gamma}\le 1$; and the conditions on the dates, namely a condition on the clock
$\check{R}_k$, the further condition that the definition of the dates of this section imposes
beside the clock condition, and the requirement of one class per step.
\item[(b)] \emph{Made part of the standing hypotheses of the arrest theorem:} the inequality
$120\,\bar{\gamma}_0\le 1$, a $\gamma$-ceiling of degree $-\infty$. It is used only
in the preliminary step of the proof of Lemma~\ref{4.27:lem-collar-displacement}.
\item[(c)] \emph{New to those hypotheses:}
\begin{itemize}
\item condition~\eqref{4.27:eq-collar-displacement-a}, a $\gamma$-ceiling of degree $-\infty$;
\item condition~\eqref{4.27:eq-real-point-collar-a}, consisting of two $\gamma$-ceilings of
degree $-\infty$;
\item the level-zero ceiling on the response field $\mathbb{H}$ and the collar ceiling. These
are two $\gamma$-ceilings in sup form, of degree at most $-(r+2)$ in $\log x$, with $x$ the
inverse squared running coupling, by the statement bounding the degree in $\log x$ of ceilings
in sup form; both belong to the list of ceilings on $\gamma_J$;
\item the floor $\mathsf{a}\ge c(4,L)$ on Ba{\l}aban's free ratio $\mathsf{a}=A_0/A_1$, where
$c(4,L)$, the number written $c(d,L)$ with $d$ the dimension in
Lemma~\ref{4.27:lem-real-point-collar} and Theorem~\ref{4.27:thm-collar-rigidity}, depends only
on the dimension $4$ and on $L$. This
floor is also a hypothesis of the correlation theorem.
\end{itemize}
\item[(d)] \emph{Order of choice.} The order
\begin{equation*}
L\to\mathfrak{M}\to\kappa\to\kappa_1\to M_1\to M\to(q,C)\to\gamma
\end{equation*}
is kept, $\mathfrak{M}$ being the class of exponentially decaying profiles. No parameter is
chosen again after a later one.
\item[(e)] \emph{Memberships.} The stratum factor $\mathsf{f}_0$, the cumulated allowance
$\mathfrak{s}$, the level-zero radii $r_a(0)$ and the condition imposed on the pair
$\tilde{p}$ in Lemma~\ref{4.27:lem-collar-displacement} are memberships in the sense of the
rule governing the choice of $\gamma$: no constant, floor, ceiling or size is chosen as a
function of them.
\item[(f)] \emph{How the cutoff, the dates and the bare coupling enter.} The numbers $K$, $k$,
$N_k$, $N_0$, the ages (the numbers of steps since the creation of a large-field region), the
nest depth, the number of dates and $|X_\tau|$ enter only through sums
$\sum 2^{-h}$, sums $\sum 2^{-k}$ and the factor $\lambda^{1/2}e^{-8\lambda}$. Apart from this,
the bare coupling $g_0$, and with it $K=K(g_0,g)$, enters only through the radii $r_a(0)$.
Wherever these radii
occur in an inequality they occur on both sides of it, and the quotients of radii are bounded
uniformly by part~(a) of the radii-quotient lemma of this section.
\item[(g)] \emph{The group.} Of $SU(N)$ only two facts are used: the local-size lemma of this
section uses a submultiplicative, unitarily invariant norm; and
the constant $c_\pi$ of the projection onto $\mathfrak{sl}(N,\mathbb{C})$ satisfies
$c_\pi\le 2$. No other step depends on $N$.
\item[(h)] \emph{Scope.} Steps with two windows are not covered. The statement of the lemma of
this section on distinct first levels, that the first levels $h_k$ are distinct, applies to
steps carrying exactly one pair $(h,k)$.
\end{enumerate}
\end{remark}

\begin{proof}
Every condition in (b) and (c) is an inequality in one direction.

Condition~\eqref{4.27:eq-collar-displacement-a} is a ceiling on $\gamma$. Its degree is
$-\infty$ because the constant $\bar{P}$ in it is proportional to $e^{-5\delta_1 R_k}$. The
number $K_W^{0}$ in it depends only on the dimension $4$ and on $K_J$, so it is fixed before
$\gamma$. The two
members of~\eqref{4.27:eq-real-point-collar-a} are ceilings of the same kind, each of degree
$-\infty$. The level-zero ceiling on $\mathbb{H}$ and the collar ceiling are ceilings in sup
form, and their degrees in $\log x$ are at most $-(r+2)$ by the statement bounding the degree in
$\log x$ of ceilings in sup form. The remaining new condition is a
floor on the ratio $A_0/A_1$ by a number $c(4,L)$ that depends only on the dimension $4$ and
on $L$.

Hence each new condition bounds either $\gamma$ from above or $A_0/A_1$ from below. None of
them bounds any of $L,\mathfrak{M},\kappa,\kappa_1,M_1,M,q,C$, and the floor on $A_0/A_1$
involves only the dimension and $L$, which is chosen first. Therefore the order in (d) is
unchanged, and no parameter has to be chosen again after a later one.

For item~(e): the cumulated allowance $\mathfrak{s}$ enters only through the stratum factor
$\mathsf{f}_0(d)=4-\mathfrak{s}(d)$; the stratum factor and the radii $r_a(0)$ enter only through
the dated thresholds $\Theta^d_a$, which for points of the collar are $\mathsf{f}_0\,r_a(0)$;
and the condition on $\tilde{p}$ enters only as a hypothesis of
Lemma~\ref{4.27:lem-collar-displacement}. None of these quantities is an argument of any of the
constants, floors, ceilings or sizes chosen in the order of (d).

For (f), first consider the radii $r_a(0)$. Wherever they occur in an inequality they occur on
both sides of it. Dividing such an inequality by the $r_a(0)$ on its right-hand side leaves
only quotients of radii, and these quotients are bounded uniformly by part~(a) of the
radii-quotient lemma of this section. The other quantities
listed in (f) appear only inside the geometric sums $\sum 2^{-h}$ and $\sum 2^{-k}$ and the
factor $\lambda^{1/2}e^{-8\lambda}$.

Item~(g) follows from the local-size lemma of this section, whose
only input from the group is a submultiplicative, unitarily invariant norm, together with the
bound $c_\pi\le 2$ on $\mathfrak{sl}(N,\mathbb{C})$.

Item~(h) records the scope of the hypotheses: the lemma on distinct first levels is applied
with one pair $(h,k)$ per step, and steps with two windows are outside the hypotheses of
Theorem~\ref{4.27:thm-collar-rigidity}.
\end{proof}

\begin{definition}[Notation for an interrupted preparation: levels, stages and block structures]
\label{4.28:not-preparation-notation}
Throughout this section the notation below is in force.
\begin{enumerate}
\item A level $k$ is fixed at which the large-field operation $\mathbf{R}$ acts (an
\emph{$\mathbf{R}$-site}), together with the integers $N$ and $N_0$ attached to it in the
relative small-field machine; $N$ is the number of stages of the preparation at $k$. In this
section the letter $N$ denotes this integer and not the $N$ of the gauge group $SU(N)$. We write
\begin{equation*}
h=k-N,\qquad k_0=k-N_0 .
\end{equation*}
\item A \emph{stage} is an integer $l$ with $l\le N$, the index of the stage-by-stage
construction of the relative small-field machine (written $n$ there). To a stage $l$ we attach
the level
\begin{equation*}
j:=h+l-1 .
\end{equation*}
\item The input's block structure at stage $l$ is $\mathbb{B}^{(l)}_k$. We write
$\mathfrak{B}:=\mathbb{B}^{(l-1)}_k$ for the block structure at stage $l-1$, and
$\ell_{\mathfrak{B}}$ for its level function.
\item $Z$ is the large-field region, and $Z^{[l]}$, $Z^{\mathrm{on}}$, $\mathfrak{T}$ are its
parts during the interrupted preparation; $R^{(l)}=R_{l-1}$ is the integration region of
stage $l$ (a region, not to be confused with the radii $R_i$ of \cite[(2.9)]{B-CMP119});
$\mathcal{C}$ is the cover of the region, $\mathcal{T}_1$ the family of terms of the
expansion, $\mathfrak{L}_l$ the attempt space at stage $l$ and $\mathfrak{D}^{\flat}_l$ its
domain of configurations. All of these are as fixed in the relative small-field machine, in its
opening definitions and in its stage-by-stage construction.
\item For every level $n$, $\ell_n$ denotes the lattice spacing of level $n$, and $\pi_n$
denotes the family of cubes of side $M\ell_n$; this side length is called the
\emph{$\pi_n$-side}. (The subscript of $\ell$ is a block structure when a level function is
meant and a level $n$ when a spacing is meant.)
\item The enlarged regions $Z'_i$ and the operation $S\mapsto S^{\sim}$, which adds one layer
of cubes to a union of cubes, are those of \cite[(3.2)~ff.]{B-CMP119}.
\end{enumerate}
\end{definition}

\begin{definition}[Attempt data and the three kinds of lifted points]\label{4.28:def-attempt-data}
The notation is that of Definition~\ref{4.28:not-preparation-notation}. Let $\mathsf{k}$ be a level at
which Ba{\l}aban's operation $\mathbf{R}$ is performed (an \emph{$\mathbf{R}$-site}), and put
\begin{equation}\label{4.28:eq-attempt-data-1}
\mathsf{h}:=\mathsf{k}-N_{\mathsf{k}},\qquad \mathsf{k}_0:=\mathsf{k}-N_0(\mathsf{k}),
\end{equation}
where $N_{\mathsf{k}}$ and $N_0(\mathsf{k})$ are the numbers $N$ and $N_0$ of
Definition~\ref{4.28:not-preparation-notation} for the site $\mathsf{k}$.
Let $\Omega_i$, $\Lambda_i$, $Z_i=\Lambda_i^c$ and $Z''_i$ be the regions of the site $\mathsf{k}$
\cite[(1.10)--(1.11)]{B-CMP122a}, and put $\Gamma''_i:=Z''_{i+1}\setminus Z''_i$.
An \emph{attempt datum} is a tuple
\begin{equation}\label{4.28:eq-attempt-data-2}
\mathfrak{d}=(C;\mathsf{k},\mathsf{h},\mathsf{k}_0,J),
\end{equation}
where $C$ is a component of a region of the class defined by conditions (i), (ii) of
\cite{B-CMP122a}, subject to the nesting condition on excluded regions of the interrupted
preparation, taken at the $\mathbf{R}$-site $\mathsf{k}$; where $J\in[\mathsf{h},\mathsf{k}]$; and
where $C\cap\Omega_{\mathsf{k}}^c$ carries the levels $\ell$ and the weights $\omega$ of
\cite[(1.64)--(1.69)]{B-CMP122a} at $j=J$. The regions $\Omega_i$, $\Lambda_i$, $Z_i$,
$Z''_i$, $\Gamma''_i$ attached to $\mathfrak{d}$ are those of the site $\mathsf{k}$, intersected
with $C$. The datum $\mathfrak{d}$ is \emph{lifted} if $J>\mathsf{k}_0$.

At the site $k$ and stage $l$ of the run, with $h=k-N$, $k_0=k-N_0$ and $Z$, $Z^{[l]}$,
$\mathfrak{L}_l$ as in Definition~\ref{4.28:not-preparation-notation}, attempt data arise in
three kinds.
\begin{itemize}
\item[(A)] Let $x\in Z^{\mathrm{on}}$, let $\mathfrak{a}$ be the innermost live arrested older
piece at $x$, created at the site $k_a<k$, and suppose that its excluded region
$W_{\mathfrak{a}}$ satisfies $W_{\mathfrak{a}}\subset Z$. Its datum is
\begin{equation*}
\mathfrak{d}=(X_{\mathfrak{a}};k_a,h_a,k_0^a,j^+_a),
\end{equation*}
where $X_{\mathfrak{a}}$ is the region of $\mathfrak{a}$, $h_a:=k_a-N_{k_a}$,
$k_0^a:=k_a-N_0(k_a)$, and $j^+_a$ is the index $J$ of the datum at the arrest. This datum is
frozen at the arrest, by the arrest theorem.
\item[(B)] Let $\bar c$ be a component of $Z^{[l]}$. Its datum is
$\mathfrak{d}=(\bar c;k,h,k_0,h+l)$, the datum with respect to which the attempt space
$\mathfrak{L}_l$ is defined.
\item[(B$'$)] Let $W$ be a region excluded at a stage $\le l-1$ of the run, in the sense of the
nesting condition on excluded regions. For each component $C_W$ of $W$ the datum is
$\mathfrak{d}=(C_W;k,h,k_0,j^+)$ with $j^+\le k-1$, frozen at the exclusion.
\end{itemize}
These data have the following two properties.
\begin{enumerate}
\item[(P)] For a datum of kind (A), with $Z_{k_a}$, $Z_h$, $Z''_h$ the regions
\cite[(1.10)--(1.11)]{B-CMP122a} of the site $k$,
\begin{equation}\label{4.28:eq-attempt-data-3}
k_a\le h,\qquad X_{\mathfrak{a}}\subset Z_{k_a}\subset Z_h=Z''_h .
\end{equation}
\item[(E)] At $(k,l)$ every weight $\omega\ne1$ on $Z^{\mathrm{on}}$ is carried either by an
attempt datum of kind (B) or (B$'$), or by a live arrested older piece $\mathfrak{a}$, whose
datum is of kind (A) when $W_{\mathfrak{a}}\subset Z$; off the run and off live pieces
$\omega=1$; and on $Z^{\mathrm{on}}$ the factor $c$ of \cite[(1.64)--(1.69)]{B-CMP122a}
equals $1$.
\end{enumerate}
\end{definition}

\begin{proof}
\emph{Property (P).} By the clause of the arrest theorem describing the branches of
Ba{\l}aban's step, an arrest of an operation at the site $k_a$ takes
place only on the branch of Ba{\l}aban's step in which new large fields were introduced at the
step $k_a$. Condition (ii) of the class of \cite{B-CMP122a}, imposed at the site $k$, forbids
the introduction of new large fields in the last $N$ steps before $k$. Hence $k_a\le k-N=h$.
By the same clause of the arrest theorem, the piece $X_{\mathfrak{a}}$ lies in the region
$Z_{k_a}$ of the site $k$, and $Z_h=Z''_h$; the
chain \eqref{4.28:eq-attempt-data-3}
follows, the inclusion $Z_{k_a}\subset Z_h$ holding because the regions $\Lambda_i$ of
\cite[(1.10)--(1.11)]{B-CMP122a} decrease as $i$ increases, so that their complements
$Z_i=\Lambda_i^c$ increase, and $k_a\le h$.

Kind (A) is the generic case. In \cite{B-CMP122b} a region $\hat Z=S(\hat Z_0)$ is
considered, $S$ being the enlargement operation of that paper and $\hat Z_0$ a region
satisfying conditions (i), (ii), in the case in which a new large field was introduced in the
preparatory operations; for that case \cite{B-CMP122b} states a containment in a cube of side
$100MR_{j+1}$, where $j+1$ is the level index of that paper and not the index $j$ of
Definition~\ref{4.28:not-preparation-notation}, and deduces that the number of steps attached
to $\hat Z$ there equals $R_{j+1}$. A datum of kind (A) arises in this situation: during that
number of steps the piece left un-joined there is a live arrested older piece $\mathfrak{a}$,
with $X_{\mathfrak{a}}$ that piece, so that by \eqref{4.28:eq-attempt-data-3} it is contained
in $Z_h=Z''_h$; after that number of steps it belongs to the class again and lies in the core
of the region being prepared.

\emph{Property (E).} The weights $\omega$ are defined by \cite[(1.64)]{B-CMP122a} and
\cite[(1.66)]{B-CMP122a} alone. Off
the run and off live pieces every bound in force is one of the bounds of
\cite[(2.34)--(2.39)]{B-CMP119} or of \cite[(1.68)]{B-CMP122a}, and there $\omega=1$. It
remains to dispose of an operation $\mathbf{R}$, at a site $\mathsf{k}$, that has been completed.
For such an operation the small-field region satisfies $\Lambda\supset Z''_{\mathsf{k}}$
\cite{B-CMP122a}. The region $Z''_{\mathsf{k}}$ contains all the thin strata
$\Gamma''_i$ of that operation; the completed operation re-levels $C\setminus\Lambda$ at the
level $\mathsf{k}$; and, by the clause of the arrest theorem on completed operations, it drops its
pieces. Hence a completed operation leaves no weight $\omega\ne1$ behind, and a weight
$\omega\ne1$ at $(k,l)$ belongs either to the run at $k$ itself, through a component of
$Z^{[l]}$ (kind (B)) or through a region excluded at an earlier stage (kind (B$'$)), or to a
live arrested older piece, which is of kind (A) when its excluded region $W_{\mathfrak{a}}$
lies in $Z$. Finally, the factor $c$ of \cite[(1.64)--(1.69)]{B-CMP122a}
equals $1$ on $Z^{\mathrm{on}}$; the value $c=3$, which those
formulae prescribe when their level index equals the site $k$, occurs only on $\Omega_k$.
\end{proof}

\begin{definition}[Levels, weights and the cutting level on a lifted datum]\label{4.28:def-cutting-level}
Let $\mathfrak{d}=(C;\mathsf{k},\mathsf{h},\mathsf{k}_0,J)$ be an attempt datum in the sense of
Definition~\ref{4.28:def-attempt-data}, and assume that it is lifted, that is, $J>\mathsf{k}_0$. Let
$\Omega_i$, $Z_i$, $Z''_i$ and $\Gamma''_i=Z''_{i+1}\setminus Z''_i$ be the regions of the
$\mathbf{R}$-site $\mathsf{k}$, cut to $C$. On $C\cap\Omega^c_{\mathsf{k}}$ the datum carries the levels
$\ell$ and the weights $\omega$ of \cite[(1.64), (1.66), (1.68)]{B-CMP122a}, with the index of
those formulas taken equal to $J$. In \eqref{4.28:eq-cutting-level-a} and
\eqref{4.28:eq-cutting-level-b}, $L_{\mathrm{w}}$ denotes the base of the weights of
\cite[(1.64), (1.66), (1.68)]{B-CMP122a}, written $L_0$ there. On the
sets of the first column, each intersected with
$C\cap\Omega^c_{\mathsf{k}}$, these functions and the function $\mathfrak{c}$ defined below are
\begin{equation}\label{4.28:eq-cutting-level-a}
\begin{array}{llll}
\text{set} & \ell & \omega & \mathfrak{c}\\
\hline
Z_{\mathsf{h}};\ \ \Gamma''_m,\ \mathsf{h}\le m\le\mathsf{k}_0
  & \text{its own};\ m & 1 & \ell\\
\Gamma''_m,\ \mathsf{k}_0<m<J\ \ (\text{thin zones})
  & m & L_{\mathrm{w}}^{2(m-\mathsf{k}_0)} & \mathsf{k}_0\\
\Omega^c_{\mathsf{k}_0+1}\setminus Z''_J\ \ (\text{band})
  & J & L_{\mathrm{w}}^{2(J-\mathsf{k}_0)} & \mathsf{k}_0\\
\Omega_m\setminus\Omega_{m+1},\ \mathsf{k}_0<m\le J\ \ (\text{old shells})
  & J & L_{\mathrm{w}}^{2(J-m)} & m\\
\Omega_m\setminus\Omega_{m+1},\ J<m<\mathsf{k}\ \ (\text{old shells})
  & m & 1 & m
\end{array}
\end{equation}
where on $Z_{\mathsf{h}}$ the level $\ell$ is the one already carried there, unaltered. This is
the table of levels and weights at a live arrested piece $\mathfrak{a}$ (the pieces entering
kind (A) of Definition~\ref{4.28:def-attempt-data}), with $(W_{\mathfrak{a}},k_0^a,j^+_a,k_a)$
replaced by $(C,\mathsf{k}_0,J,\mathsf{k})$. The
\emph{cutting level} $\mathfrak{c}$, the \emph{depth} $\mathsf{e}$ and the set $\mathsf{L}_{\mathfrak{d}}$
are
\begin{equation}\label{4.28:eq-cutting-level-b}
\mathfrak{c}:=\ell-\tfrac12\log_{L_{\mathrm{w}}}\omega,\qquad \mathsf{e}:=\ell-\mathfrak{c},\qquad
\mathsf{L}_{\mathfrak{d}}:=\bigl(\Omega^c_J\setminus Z''_{\mathsf{k}_0+1}\bigr)\cap C .
\end{equation}
A point $y\in C\cap\Omega^c_{\mathsf{k}}$ is \emph{lifted} if $\mathsf{e}(y)\ge1$; the lifted points are
exactly the points of $\mathsf{L}_{\mathfrak{d}}$.

Let $k$, $l$, $j$, $\mathfrak{B}=\mathbb{B}^{(l-1)}_k$ and its level function
$\ell_{\mathfrak{B}}$ be as in Notation~\ref{4.28:not-preparation-notation}. If $\mathfrak{d}$ is of
kind (A) or (B$'$), then $\ell_{\mathfrak{B}}=\ell$. If $\mathfrak{d}$ is of kind (B), then
$j=J-1$, and $\ell_{\mathfrak{B}}$ is the level function $\ell$ of
\eqref{4.28:eq-cutting-level-a} read with $J-1$ in place of $J$ (when $J-1=\mathsf{k}_0$ the second
and fourth rows of \eqref{4.28:eq-cutting-level-a} are empty and the band carries weight $1$). In
every case
\begin{equation}\label{4.28:eq-cutting-level-1}
\mathfrak{c}\le\ell_{\mathfrak{B}}\le\ell\qquad\text{on } C\cap\Omega^c_{\mathsf{k}} .
\end{equation}
\end{definition}

\begin{proof}
\emph{The last column.} Inserting the second and third columns of
\eqref{4.28:eq-cutting-level-a} into \eqref{4.28:eq-cutting-level-b} gives
$\mathfrak{c}=\ell$ on the first row. On the thin zones it gives
$\mathfrak{c}=m-(m-\mathsf{k}_0)=\mathsf{k}_0$, and on the band $\mathfrak{c}=J-(J-\mathsf{k}_0)=\mathsf{k}_0$.
On the old shells with $m\le J$ it gives $\mathfrak{c}=J-(J-m)=m$, and on those with $m>J$ it
gives $\mathfrak{c}=m$. These are the values of the last column.

\emph{The lifted set.} From these values, $\mathsf{e}=0$ on the first row, on the old shells with
$m>J$, and on the old shell $m=J$. Elsewhere the depth is positive: $\mathsf{e}=m-\mathsf{k}_0\ge1$ on the
thin zones, $\mathsf{e}=J-\mathsf{k}_0\ge1$ on the band (because $J>\mathsf{k}_0$), and $\mathsf{e}=J-m\ge1$ on the old
shells with $\mathsf{k}_0<m<J$. In particular $\mathsf{e}\ge0$ everywhere.

Hence the lifted points are the points of the thin zones, of the band, and of the old shells
with $\mathsf{k}_0<m<J$. The thin zones together form $Z''_J\setminus Z''_{\mathsf{k}_0+1}$. The old
shells with $\mathsf{k}_0<m<J$ together form
$\Omega_{\mathsf{k}_0+1}\setminus\Omega_J=\Omega^c_J\setminus\Omega^c_{\mathsf{k}_0+1}$. The regions
of the site $\mathsf{k}$ are nested \cite[(1.10)--(1.11)]{B-CMP122a}, cut to $C$ as in
Definition~\ref{4.28:def-attempt-data}, and the rows of \eqref{4.28:eq-cutting-level-a} are
pairwise disjoint subsets of $C\cap\Omega^c_{\mathsf{k}}$; hence on $C\cap\Omega^c_{\mathsf{k}}$
\begin{equation*}
Z''_{\mathsf{k}_0+1}\subset Z''_J\subset\Omega^c_{\mathsf{k}_0+1}\subset\Omega^c_J .
\end{equation*}
Therefore the thin zones, the band and these old shells together form
$\Omega^c_J\setminus Z''_{\mathsf{k}_0+1}$, cut to $C$, which is $\mathsf{L}_{\mathfrak{d}}$.

\emph{Kind (B).} Here $J-1=j\ge\mathsf{k}_0$, so \eqref{4.28:eq-cutting-level-a} can be read with
$J-1$ in place of $J$. We compare the two readings row by row.
\begin{itemize}
\item The first row does not involve $J$, so both readings agree there.
\item The thin zones $\Gamma''_m$ with $m<J-1$ agree. If $J-1>\mathsf{k}_0$, the zone
$\Gamma''_{J-1}$ is a thin zone at $J$, with $\ell=J-1$ and $\mathfrak{c}=\mathsf{k}_0$; at $J-1$ it
lies in the band $\Omega^c_{\mathsf{k}_0+1}\setminus Z''_{J-1}$, again with $\ell=J-1$ and
$\mathfrak{c}=\mathsf{k}_0$. If $J-1=\mathsf{k}_0$, the zone $\Gamma''_{J-1}=\Gamma''_{\mathsf{k}_0}$ lies in
the first row in both readings, with the same values.
\item On $\Omega^c_{\mathsf{k}_0+1}\setminus Z''_J$ the level drops from $J$ to $J-1$, while
$\mathfrak{c}=\mathsf{k}_0$ in both readings.
\item On the old shells with $\mathsf{k}_0<m\le J-1$ the level drops from $J$ to $J-1$, while
$\mathfrak{c}=m$ in both readings.
\item The old shell $m=J$ passes from the fourth row to the fifth. There $\ell=J$ and
$\mathfrak{c}=J$ in both readings.
\item The old shells with $m>J$ agree.
\end{itemize}
Consequently the last column is the same at $J$ and at $J-1$. Moreover $\ell_{\mathfrak{B}}$
equals $\ell-1$ on
\begin{equation*}
\bigl(\Omega^c_{\mathsf{k}_0+1}\setminus Z''_J\bigr)\cup\bigl(\Omega^c_J\setminus
\Omega^c_{\mathsf{k}_0+1}\bigr)
=\Omega^c_J\setminus Z''_J=\Omega^c_{j+1}\setminus Z''_{j+1},
\end{equation*}
which is the region where $V_j$ is integrated in \cite{B-CMP122a}; and $\ell_{\mathfrak{B}}=\ell$
elsewhere. The argument of the previous paragraph, applied to the reading at $J-1$, shows that
its depth is $\ge0$. The reading at $J-1$ has the same cutting level $\mathfrak{c}$, so
$\mathfrak{c}\le\ell_{\mathfrak{B}}$. Since $\ell_{\mathfrak{B}}\in\{\ell-1,\ell\}$, also
$\ell_{\mathfrak{B}}\le\ell$.

\emph{Kinds (A) and (B$'$).} Here $\ell_{\mathfrak{B}}=\ell$, and $\mathfrak{c}\le\ell$ because
$\mathsf{e}\ge0$. This proves \eqref{4.28:eq-cutting-level-1}.
\end{proof}

\begin{lemma}[Strata and widths of the cutting level]\label{4.28:lem-cutting-strata}
Let $\mathfrak{d}=(C;\kappa,\mathsf{h},\kappa_0,J)$ be a lifted attempt datum in the sense of
Definition~\ref{4.28:def-attempt-data}, with the level $\ell$, the weight $\omega$ and the cutting
level $\mathfrak{c}$ of \eqref{4.28:eq-cutting-level-a} and \eqref{4.28:eq-cutting-level-b}.
Let $R_i$ be the level-$i$ radii of \cite[(2.9)]{B-CMP119} and $R_{\min}$ their minimal value, and
assume
\[
R_{\min}\ge 5L .
\]
Lengths at level $n\le\kappa$ are measured in $\pi_n$-sides, a $\pi_n$-cube being a cube of side
$M\ell_n$ of the level-$n$ lattice of spacing $\ell_n$. Assume moreover the following.
\begin{enumerate}
\item[(a)] The component $C$ satisfies the nesting condition on class components at the site
$\kappa$ of the operation $\mathbf{R}$.
\item[(b)] The radii satisfy the second inequality of \cite[(2.9), p.~256]{B-CMP119} in the form
\[
R_m\le L\bigl(1+g_n^2\beta'(n-m)\bigr)^{\beta_0}R_n\qquad(n>m),
\]
where $\beta'$ is the constant and $\beta_0=1/7$ the exponent of that inequality as stated in
\cite[(2.9)]{B-CMP119}, and, by \cite[(2.5)]{B-CMP119}, every ratio $R_i/R_{i+1}$ is an integral
power of $L$.
\item[(c)] The radii satisfy the second part of the window-envelope bound for $\check{R}$, namely
$R_i\le LR_{i+1}$ for $\kappa_0\le i<\kappa$.
\item[(d)] The one-sided ceiling on the coupling
\begin{equation}\label{4.28:eq-cutting-strata-1}
g_n^2\beta'\le 1
\end{equation}
holds for every $n$ with $\kappa_0<n\le\kappa$.
\end{enumerate}
Then the following hold.
\begin{enumerate}
\item[(i)] The sets $\{\mathfrak{c}=i\}\cap C$ are the strata of the block structure
$\mathbb{B}^{(\kappa_0-\mathsf{h})}_{\kappa}\restriction C$, that is, of the structure
\cite[(1.16)]{B-CMP122a} at $j=\kappa_0$: the set $\Gamma''_i$ for $i<\kappa_0$, the set
$\Omega^c_{\kappa_0+1}\setminus Z''_{\kappa_0}$ for $i=\kappa_0$, and the old shells
$\Omega_i\setminus\Omega_{i+1}$ for $i>\kappa_0$. The function $\mathfrak{c}$ does not depend
on $J$.
\item[(ii)] For $\mathsf{h}\le i<\kappa$ the $\mathfrak{c}$-shell $\{\mathfrak{c}=i\}\cap C$ is at
least three layers wide, and for $\kappa_0\le i<\kappa$ at least eight layers wide, the layers
consisting of $MR_{i+1}$-cubes of the $(i+1)$-lattice; in particular
\begin{equation}\label{4.28:eq-cutting-strata-b}
\{\mathfrak{c}=i\}\cap C\ \text{has width at least}\ 3R_{\min}\ \pi_{i+1}\text{-sides}
\qquad(\mathsf{h}\le i<\kappa).
\end{equation}
\item[(iii)] For $\mathsf{h}<n\le\kappa$,
\[
\operatorname{dist}\bigl(\Lambda_n,\{\mathfrak{c}<n\}\bigr)\ge 2R_n\ \pi_n\text{-sides}.
\]
\item[(iv)] Inside $\{\mathfrak{c}=\kappa_0\}$ the strata of $\ell$ are $\Gamma''_{\kappa_0}$, of
width at least $LR_{\kappa_0+1}$ $\pi_{\kappa_0}$-sides; the thin sets $\Gamma''_m$,
$\kappa_0<m<J$, each of width at least $L^2$ $\pi_{\kappa_0}$-sides; and the band
$\Omega^c_{\kappa_0+1}\setminus Z''_J$, of width at least $5L^{N_0}$ $\pi_{\kappa_0}$-sides.
Moreover $N_0(\kappa)\ge 3$.
\item[(v)] The numbers
\begin{equation}\label{4.28:eq-cutting-strata-c}
\rho_m:=\frac{R_{m+1}\ell_{m+1}}{\ell_J}\quad(\kappa_0\le m<J)\qquad\text{satisfy}\qquad
1\le\rho_{\kappa_0}\le\rho_{\kappa_0+1}\le\dots\le\rho_{J-1}.
\end{equation}
For $\kappa_0<m<J$ the old shell $\Omega_m\setminus\Omega_{m+1}$, on which $\ell=J$, has width at
least $8\rho_m$ $\pi_J$-sides. Every $\pi_n$-cube, $n\le\kappa$, contained in
$\Omega_{\kappa_0+1}$ lies in a single old shell. For $\kappa_0\le m<J$ the set
$\Omega^c_{m+1}\cap C$ lies in a cube of side $100\rho_m$ $\pi_J$-sides.
\end{enumerate}
\end{lemma}

\begin{proof}
All sets below are cut to $C$. We use throughout that $\ell_{n+1}=L\ell_n$, so that an
$LMR_{i+1}$-cube of scale $i$ is an $MR_{i+1}$-cube of the $(i+1)$-lattice, of side
$MR_{i+1}\ell_{i+1}$, that is, of $R_{i+1}$ $\pi_{i+1}$-sides.

\emph{The nested containments.} Fix $i$ with $\mathsf{h}\le i<\kappa$. Since $C$ is a component of
the class of \cite{B-CMP122a} at the site $\kappa$, clause (ii) of that class says that in the
$N_\kappa=\kappa-\mathsf{h}$ renormalisation steps preceding $\kappa$ no new large-field region was
created inside $C$; in particular none of the regions $P$, $Q$, $R$ of the construction
\cite[(3.2)--(3.5)]{B-CMP119} is made inside $C$ at the steps producing the indices $i+1$ with
$\mathsf{h}<i+1\le\kappa$. The construction \cite[(3.2)--(3.5)]{B-CMP119} therefore applies there
without these regions. In it $Z'_i$ is the union of all $LMR_{i+1}$-cubes meeting $Z_i$, taken
from partitions compatible with all the previous ones, and it is enlarged successively by four
layers, a layer, a layer and two layers of such cubes; since $4+1+1+2=8$, this gives
$\Omega^c_{i+1}\supseteq {Z'_i}^{\sim8}$. Next, \cite[(3.20)]{B-CMP119} gives
\[
Z_{i+1}=\Lambda^c_{i+1}=\bigl((\Omega^c_{i+1})^{\sim}\cup R'_{i+1}\bigr)^{\sim}
\supseteq(\Omega^c_{i+1})^{\sim2}\supseteq {Z'_i}^{\sim10},
\]
where $R'_{i+1}$ is the set of \cite[(3.20)]{B-CMP119}; the first containment holds because the
operation ${}^{\sim}$ is monotone, and the second is $\Omega^c_{i+1}\supseteq{Z'_i}^{\sim8}$ with
two further layers added. When none of the regions $P$, $Q$, $R$ is made, which is the case under
clause (ii) of
the class, $R'_{i+1}$ is empty, $\Omega^c_{i+1}={Z'_i}^{\sim8}$, and $\Lambda^c_{i+1}$ is
$(\Omega^c_{i+1})^{\sim}$ enlarged by one further layer, so that equality holds throughout, in
agreement with \cite{B-CMP122b}, where the corresponding set is the ten-layer enlargement
${Z'_i}^{\sim10}$; this equality is not used below, where only the containments enter. We also have
$Z'_i\supset Z_i\supset\Omega^c_i$; the layers in question are layers of $LMR_{i+1}$-cubes of
scale $i$, that is, of $MR_{i+1}$-cubes of the $(i+1)$-lattice; and all the sets involved are
rectangular parallelepipeds, as \cite{B-CMP122a} states for the regions connected with the last
$N_\kappa$ steps.
The operation ${}^{\sim}$ applied to $\Omega_{i+1}$ likewise adds one layer of such cubes
\cite{B-CMP119}.

Let now $\mathsf{h}\le i<\kappa_0$. By \cite[(1.11), (1.10)]{B-CMP122a},
$Z''_{i+1}:=(\Omega^{\sim5}_{i+1})^c\cap Z$. A cube of ${Z'_i}^{\sim3}$ has every cube within five
adjacency steps inside ${Z'_i}^{\sim8}\subseteq\Omega^c_{i+1}$, hence it lies outside
$\Omega^{\sim5}_{i+1}$; and it lies inside ${Z'_i}^{\sim10}\subseteq Z_{i+1}\subset Z$. Hence
$Z''_{i+1}\supseteq{Z'_i}^{\sim3}$. Finally, we use the two equalities for the top indices as
stated in \cite[(1.10), p.~179]{B-CMP122a}:
$Z''_{\kappa_0+1}={Z'_{\kappa_0}}^{\sim}$ and $Z''_{\kappa}={Z'_{\kappa_0}}^{\sim3}$. We collect:
\begin{equation}\label{4.28:eq-cutting-strata-a}
\begin{gathered}
\Omega^c_{i+1}\supseteq{Z'_i}^{\sim8}\qquad(\mathsf{h}\le i<\kappa),\\
Z_{i+1}=\Lambda^c_{i+1}\supseteq(\Omega^c_{i+1})^{\sim2}\supseteq{Z'_i}^{\sim10}
\qquad(\mathsf{h}\le i<\kappa),\\
Z''_{i+1}\supseteq{Z'_i}^{\sim3}\qquad(\mathsf{h}\le i<\kappa_0),\\
Z''_{\kappa_0+1}={Z'_{\kappa_0}}^{\sim},\qquad Z''_{\kappa}={Z'_{\kappa_0}}^{\sim3}.
\end{gathered}
\end{equation}
In addition $Z''_{\mathsf{h}}=Z_{\mathsf{h}}$, and the sets $Z''_i$ increase with $i$
\cite[pp.~178--179]{B-CMP122a}.

\emph{Proof of \textup{(i)}.} By the last column of \eqref{4.28:eq-cutting-level-a},
$\mathfrak{c}=\kappa_0$ on the thin sets $\Gamma''_m$, $\kappa_0<m<J$, and on the band
$\Omega^c_{\kappa_0+1}\setminus Z''_J$, and $\mathfrak{c}=\kappa_0$ on $\Gamma''_{\kappa_0}$ by the
first row. Since $\Gamma''_m=Z''_{m+1}\setminus Z''_m$ and the $Z''_m$ increase, the union of
$\Gamma''_m$ over $\kappa_0\le m<J$ is $Z''_J\setminus Z''_{\kappa_0}$. Moreover, by
\eqref{4.28:eq-cutting-strata-a},
\[
Z''_J\subset Z''_\kappa={Z'_{\kappa_0}}^{\sim3}\subset{Z'_{\kappa_0}}^{\sim8}\subseteq
\Omega^c_{\kappa_0+1}.
\]
Hence
\[
(Z''_J\setminus Z''_{\kappa_0})\cup(\Omega^c_{\kappa_0+1}\setminus Z''_J)
=\Omega^c_{\kappa_0+1}\setminus Z''_{\kappa_0},
\]
which is the level-$\kappa_0$ set of \cite[(1.16)]{B-CMP122a} at $j=\kappa_0$. The remaining rows of
\eqref{4.28:eq-cutting-level-a} are the remaining strata of that structure: $\mathfrak{c}=\ell$ on
$Z_{\mathsf{h}}$, with the levels that the structure assigns there, $\mathfrak{c}=m$ on
$\Gamma''_m$ for $m<\kappa_0$, and $\mathfrak{c}=m$ on the old shell $\Omega_m\setminus\Omega_{m+1}$
for every $m>\kappa_0$, whether $m\le J$ or $m>J$. Since these strata and the values of
$\mathfrak{c}$ on them are those of the structure at $j=\kappa_0$, into which $J$ does not enter,
$\mathfrak{c}$ does not depend on $J$.

\emph{Proof of \textup{(ii)}.} Only the containments in \eqref{4.28:eq-cutting-strata-a} are used.
For $\mathsf{h}\le i<\kappa_0$, the shell is $\Gamma''_i=Z''_{i+1}\setminus Z''_i$, and
\[
\Gamma''_i\supseteq{Z'_i}^{\sim3}\setminus Z''_i\supset{Z'_i}^{\sim3}\setminus Z'_i ,
\]
because $Z''_i\subset\Omega^c_i\subset Z'_i$ for $i>\mathsf{h}$ (as
$Z''_i=(\Omega^{\sim5}_i)^c\cap Z$) and $Z''_{\mathsf{h}}=Z_{\mathsf{h}}\subset Z'_{\mathsf{h}}$. For
$i=\kappa_0$,
\[
\Omega^c_{\kappa_0+1}\setminus Z''_{\kappa_0}\supseteq{Z'_{\kappa_0}}^{\sim8}\setminus Z''_{\kappa_0}
\supset{Z'_{\kappa_0}}^{\sim8}\setminus Z'_{\kappa_0}.
\]
For $\kappa_0<i<\kappa$, the old shell is $\Omega_i\setminus\Omega_{i+1}=\Omega^c_{i+1}\setminus
\Omega^c_i$, and
\[
\Omega^c_{i+1}\setminus\Omega^c_i\supseteq{Z'_i}^{\sim8}\setminus\Omega^c_i\supset
{Z'_i}^{\sim8}\setminus Z'_i .
\]
In each case the right-hand side consists of three, respectively eight, layers of
$MR_{i+1}$-cubes of the $(i+1)$-lattice. One layer is $R_{i+1}\ge R_{\min}$ $\pi_{i+1}$-sides
wide, which gives \eqref{4.28:eq-cutting-strata-b}.

\emph{Proof of \textup{(iii)}.} Let $\mathsf{h}<n\le\kappa$. By (i), $\{\mathfrak{c}<n\}$ is
$\Omega^c_n$ if $n>\kappa_0$, and it is $Z''_n\subset\Omega^c_n$ if $n\le\kappa_0$. By
\eqref{4.28:eq-cutting-strata-a} at $i=n-1$, $\Lambda^c_n\supseteq(\Omega^c_n)^{\sim2}$. Thus
$\Lambda_n$ misses the two-layer neighbourhood of $\Omega^c_n$ formed by $MR_n$-cubes of the
$n$-lattice, whose layers are $R_n$ $\pi_n$-sides wide.

\emph{Proof of \textup{(iv)}.} By the definition of $N_0$ in \cite[pp.~178--179]{B-CMP122a},
$L^{-N_0+1}MR_{\kappa_0+1}=M$, that is, $L^{N_0-1}=R_{\kappa_0+1}$. Since
$R_{\kappa_0+1}\ge R_{\min}\ge 5L$, in particular $L^{N_0-1}>L$, so $N_0-1>1$, that is, $N_0\ge3$.
An $MR_{\kappa_0+1}$-cube
of scale $\kappa_0+1$ is an $M$-cube of scale $\kappa$ \cite[pp.~178--179]{B-CMP122a}; its side is
\[
MR_{\kappa_0+1}\ell_{\kappa_0+1}=ML^{N_0-1}\ell_{\kappa_0+1}=M\ell_\kappa=L^{N_0}M\ell_{\kappa_0},
\]
that is, $L^{N_0}=LR_{\kappa_0+1}$ $\pi_{\kappa_0}$-sides. By (i) and
\eqref{4.28:eq-cutting-level-a}, the strata of $\ell$ inside $\{\mathfrak{c}=\kappa_0\}$ are
$\Gamma''_{\kappa_0}$ (level $\kappa_0$), the thin $\Gamma''_m$ (level $m$) and the band (level
$J$). First,
\[
\Gamma''_{\kappa_0}=Z''_{\kappa_0+1}\setminus Z''_{\kappa_0}
\supset{Z'_{\kappa_0}}^{\sim}\setminus Z'_{\kappa_0},
\]
since $Z''_{\kappa_0}\subset Z'_{\kappa_0}$ as in (ii); this is one layer of such cubes, of
width $LR_{\kappa_0+1}$ $\pi_{\kappa_0}$-sides. Second, since $J\le\kappa$ and the $Z''_i$
increase,
\[
\Omega^c_{\kappa_0+1}\setminus Z''_J\supseteq{Z'_{\kappa_0}}^{\sim8}\setminus Z''_\kappa
={Z'_{\kappa_0}}^{\sim8}\setminus{Z'_{\kappa_0}}^{\sim3},
\]
by \eqref{4.28:eq-cutting-strata-a} and the equality $Z''_\kappa={Z'_{\kappa_0}}^{\sim3}$ of
\cite[(1.10)]{B-CMP122a}. The right-hand side is five
layers of $M$-cubes of scale $\kappa$, of width $5L^{N_0}$ $\pi_{\kappa_0}$-sides. Third, for
$\kappa_0+1<m<J$ the thin set $\Gamma''_m$ is one layer of $M$-cubes of scale $m+1$
\cite[pp.~178--179]{B-CMP122a}, of width $L^{m+1-\kappa_0}\ge L^2$ $\pi_{\kappa_0}$-sides. For
$m=\kappa_0+1$, the set $\Gamma''_{\kappa_0+1}$ is what remains of the two layers of
\cite[(1.10)]{B-CMP122a} between
$Z''_{\kappa_0+1}={Z'_{\kappa_0}}^{\sim}$ and $Z''_\kappa={Z'_{\kappa_0}}^{\sim3}$, and its width is
at least $(2-L/(L-1))L^{N_0}$ $\pi_{\kappa_0}$-sides. Since $L$ is an odd integer greater than $1$,
we have $L\ge3$, so $L^2-3L+1\ge1>0$, that is, $L(L-2)\ge L-1$. With $N_0\ge3$ this gives
\[
\Bigl(2-\frac{L}{L-1}\Bigr)L^{N_0}=\frac{L-2}{L-1}\,L^{N_0}\ge\frac{L(L-2)}{L-1}\,L^2\ge L^2 .
\]

\emph{Proof of \textup{(v)}.} We first show $R_i\le LR_{i+1}$ for $\kappa_0\le i<\kappa$. This is
hypothesis (c); \cite{B-CMP122a} states the same inequality where the existence of $N_0$ is shown.
Independently, it is a consequence of hypotheses (b) and (d): apply hypothesis (b) with $m=i$ and
$n=i+1$, so $n-m=1$. By \cite[(2.5)]{B-CMP119},
$R_i/R_{i+1}=L^a$ with $a\in\mathbb{Z}$. By hypothesis (d),
\[
g_n^2\beta'\le 1,
\]
and since the exponent $\beta_0$ is less than one,
\[
(1+g_n^2\beta')^{\beta_0}\le1+\beta_0<2\le L ,
\]
so that
\[
L^a=\frac{R_i}{R_{i+1}}\le L(1+g_n^2\beta')^{\beta_0}<L^2 .
\]
Hence $a\le1$, that is, $R_i\le LR_{i+1}$; for this conclusion only
$(1+g_n^2\beta')^{\beta_0}<L$ is needed. Since $\ell_{i+1}=L\ell_i$, it follows that
$R_{i+1}\ell_{i+1}=LR_{i+1}\ell_i\ge R_i\ell_i$. Thus $R_i\ell_i$ is non-decreasing in $i$. By the
proof of (iv), $R_{\kappa_0+1}\ell_{\kappa_0+1}=\ell_\kappa\ge\ell_J$, because $J\le\kappa$.
Hence $\rho_{\kappa_0}=\ell_\kappa/\ell_J\ge1$, and $\rho_m$ is non-decreasing in $m$, which is
\eqref{4.28:eq-cutting-strata-c}.

\emph{Width.} For $\kappa_0<m<J$, by (ii) the old shell $\Omega_m\setminus\Omega_{m+1}$ is at least
eight layers of $MR_{m+1}$-cubes of the $(m+1)$-lattice wide. One such layer has width
$MR_{m+1}\ell_{m+1}=\rho_m M\ell_J$, that is, $\rho_m$ $\pi_J$-sides. By the fourth row of
\eqref{4.28:eq-cutting-level-a}, $\ell=J$ on this shell.

\emph{One old shell per cube.} For $m>\kappa_0$ the boundary of $\Omega_m$ runs along cubes of side
$MR_m\ell_m$ of partitions compatible with all the previous ones \cite{B-CMP119}. By the
monotonicity just proved,
\[
MR_m\ell_m\ge MR_{\kappa_0+1}\ell_{\kappa_0+1}=M\ell_\kappa\ge M\ell_n\qquad(n\le\kappa).
\]
A $\pi_n$-cube, $n\le\kappa$, therefore lies on one side of the boundary of each $\Omega_m$,
$m>\kappa_0$. Hence a $\pi_n$-cube contained in $\Omega_{\kappa_0+1}$ lies in a single old shell.

\emph{The container.} For $\kappa_0\le m<J$ we have $\Omega^c_{m+1}\subset Z_{m+1}$. By the second
part of clause (ii) of the class \cite{B-CMP122a}, the previous regions contained in $C$ satisfy
condition (i) of the class on the corresponding scales. Together with the nesting condition at
$\kappa$ (hypothesis (a)), this makes $\Omega^c_{m+1}\cap C$ a single rectangular parallelepiped
contained in a cube of side $100MR_{m+1}$ at scale $m+1$. That side is
\[
100MR_{m+1}\ell_{m+1}=100\rho_m M\ell_J ,
\]
that is, $100\rho_m$ $\pi_J$-sides.
\end{proof}

\begin{definition}[The refined block structure and cover on lifted sets]\label{4.28:def-refined-cover}
Let $\mathfrak{B}:=\mathbb{B}^{(l-1)}_k$ be the block structure at stage $l-1$ of the run at site $k$, with level function $\ell_{\mathfrak{B}}$. Let the attempt data $\mathfrak{d}=(C;\kappa,\mathsf{h},\kappa_0,J)$ be as in the definition of attempt data, and let the cutting level $\mathfrak{c}$, the depth $g=\ell-\mathfrak{c}$, the lifted points and the lifted set $\mathsf{L}_{\mathfrak{d}}$ be as in Definition~\ref{4.28:def-cutting-level}. In the definition of the cover $\mathcal{C}$ and of the restricted block structures at its boxes, which the present definition modifies and which is called the unmodified definition below, the condition $W_{\mathfrak{a}}\subset Z^c$ is dropped, and the following objects are added for the stages of the run.
\begin{enumerate}
\item \emph{The refined block structure} $\mathfrak{B}^{\mathfrak{c}}$ is the block structure that coincides with $\mathfrak{B}$ except that, on the set $C$ of every lifted datum $\mathfrak{d}$, its determining set is that of the strata of $\mathfrak{c}$, that is, of the sets $\{\mathfrak{c}=i\}\cap C$, so that on $C$ the level function of $\mathfrak{B}^{\mathfrak{c}}$ is $\mathfrak{c}$. On such a set $C$ one has $\mathfrak{c}\le\ell_{\mathfrak{B}}$. Outside $\{\mathfrak{c}<\ell_{\mathfrak{B}}\}$, $\mathfrak{B}^{\mathfrak{c}}$ and $\mathfrak{B}$ coincide.
\item \emph{The refined cover} $\mathcal{C}^1$ is Ba{\l}aban's cover $\{\square\}$ together with its partition of unity $\{\zeta_{\square}\}$, as constructed in \cite{B-CMP119}, but built on $\mathfrak{B}^{\mathfrak{c}}$ in place of $\mathfrak{B}$.
\item \emph{Box structures.} For $\square\in\mathcal{C}^1$ put $\square_0:=\tilde{\square}^5$, the enlargement ${}^{\sim}$ being taken five times in the cube family $\pi_i$, where $i$ is the level of the cube $\square$. If $\square_0$ meets a lifted point, set
\begin{equation*}
\mathfrak{s}:=\mathfrak{B}^{\mathfrak{c}}\restriction\square_0 ,\qquad
p_{\square}:=\text{the top level of }\mathfrak{s},
\end{equation*}
and let the box configuration be
\begin{equation}\label{4.28:eq-refined-cover-1}
\mathbb{U}^{w}_{\square}=\mathbb{U}^{\mathfrak{s}}_{\square}
:=U_{p_{\square}}\bigl(\square_0,M^{\mathfrak{s}}(\mathbb{U})\bigr).
\end{equation}
Here $\mathbb{U}$ is the input configuration, and the maps $M^{\mathfrak{s}}(\cdot)$ and $U_{p}(\square_0,\cdot)$ are those of the unmodified definition. For every other $\square\in\mathcal{C}^1$, the structure $\mathfrak{s}$ and the objects built from it are those of the unmodified definition.
\item \emph{Local-scale saturation.} For a set $S$, $[S]_{\sharp}$ is the union of the cubes of the input's local scale $\pi_{\ell(\cdot)}$ that meet $S$. Every condition of the form ``$\square_0\subset Y$'' is read as $[\square_0]_{\sharp}\subset Y$. For an element $e$ attached to a box $\square$, with attached set $\mathsf{y}$ as in the unmodified definition, the domain of $e$ is
\begin{equation*}
\mathsf{d}(e):=[\tilde{\square}^5]_{\sharp}\cup\mathsf{y}.
\end{equation*}
At a pair $(T,\square)$ of a term $T$ and a box $\square$ that is far in the sense of the unmodified definition, the set $[\tilde{X}_T^4]_{\sharp}$ is added to this union; here $X_T$ is the domain of the term $T$, and $\tilde{X}_T^4$ is its enlargement by four applications of ${}^{\sim}$.
\item \emph{Choice of cover.} The terms of the family $\mathcal{T}_1$ are localised with the cover $\mathcal{C}^1$. The terms of the $M$-class, as classified in the unmodified definition, keep the cover $\mathcal{C}$, as there.
\end{enumerate}
With these objects the following hold.
\begin{itemize}
\item[(a)] Let $\square\in\mathcal{C}^1$ be a box with $\square_0$ meeting a lifted point. If the datum is of kind (B) or (B$'$) (kinds as in the definition of attempt data), the triple $(p_{\square},\square_0,\mathfrak{s})$ belongs to the list $(\Sigma^{\ast})$ of retained triples of the sixth-family lemma. If it is of kind (A), it is assumed, as a hypothesis, that the triple $(p_{\square},\square_0,\mathfrak{s})$ belongs to the list $(\Sigma^{\ast})^{\mathfrak{a}}$ of retained triples of the arrested piece $\mathfrak{a}$.
\item[(b)] $\mathfrak{B}^{\mathfrak{c}}$ is admissible, in the sense in which the block structures $\mathbb{B}^{(l)}_k$ of the run are admissible, and refines $\mathfrak{B}$.
\item[(c)] Localisation with $\mathcal{C}^1$ is exact. The total of the stage-$l$ fluctuation terms $\mathbb{C}^{(l)}_k$ is the same as with $\mathcal{C}$. The pieces of the small terms are only regrouped, now at every localisation domain $X$ meeting $\mathsf{L}_{\mathfrak{d}}$, that is, on the run, in the stages whose level $j$ exceeds $\kappa_0$.
\end{itemize}
\end{definition}

\begin{proof}
\emph{The box lies in $C$.} Let $\square\in\mathcal{C}^1$ be such that $\square_0$ meets a lifted point of a datum $\mathfrak{d}$. Then $\square_0\subset C$, by clause (o) of the lemma below on the boxes of the refined cover that meet a lifted point. Hence $\mathfrak{s}=\mathfrak{B}^{\mathfrak{c}}\restriction\square_0$ is determined by the strata of $\mathfrak{c}$ on $C$ alone.

\emph{Proof of (a).} Consider first data of kinds (B) and (B$'$). Such a datum belongs to the present run, so its site is the present site, $\kappa=k$. By Lemma~\ref{4.28:lem-cutting-strata}(i), applied to the lifted datum $\mathfrak{d}$ under that lemma's condition $R_{\min}\ge 5L$, the sets $\{\mathfrak{c}=i\}\cap C$ are the strata of $\mathbb{B}^{(\kappa_0-\mathsf{h})}_{k}\restriction C$. Here the strata of $\mathbb{B}^{(\kappa_0-\mathsf{h})}_{k}\restriction C$ are those given by \cite[(1.16)]{B-CMP122a} at $j=\kappa_0$. Since $\square_0\subset C$, this gives
\begin{equation*}
\mathfrak{s}=\mathbb{B}^{(\kappa_0-\mathsf{h})}_{k}\restriction\square_0 .
\end{equation*}
Thus $\mathfrak{s}$ is the restriction to $\square_0$ of the stage structure of stage $n'=\kappa_0-\mathsf{h}\le l-1$ of the present run. The list $(\Sigma^{\ast})$ contains, by its definition in the sixth-family lemma, the restrictions of the stage structures of all stages $n'$ not exceeding the present stage, taken before and after $\mathbf{R}$. Since $n'=\kappa_0-\mathsf{h}\le l-1$, the triple $(p_{\square},\square_0,\mathfrak{s})$ is therefore a member of $(\Sigma^{\ast})$ as that list is written.

For a datum of kind (A) the membership of $(p_{\square},\square_0,\mathfrak{s})$ in the list $(\Sigma^{\ast})^{\mathfrak{a}}$ of retained triples of the arrested piece $\mathfrak{a}$ is the hypothesis made in (a) for such data.

\emph{Proof of (b).} The admissibility of $\mathfrak{B}^{\mathfrak{c}}$ rests on two facts.

First, on the set $C$ of a lifted datum the strata of $\mathfrak{B}^{\mathfrak{c}}$ are the $\mathfrak{c}$-shells. By Lemma~\ref{4.28:lem-cutting-strata}(ii), for $\mathsf{h}\le i<\kappa$ the $\mathfrak{c}$-shell $i$ has the following width:
\begin{itemize}
\item it is at least three layers of $MR_{i+1}$-cubes of the $(i+1)$-lattice wide, and at least eight layers if $i\ge\kappa_0$;
\item it is at least $3R_{\min}$ sides of $\pi_{i+1}$-cubes wide.
\end{itemize}

Second, two distinct lifted data either lie in distinct components, or are nested, with at least $N_\kappa-N_0(\kappa)$ shells of the width just described lying between them, $N_\kappa$ and $N_0(\kappa)$ being as in the definition of attempt data.

Outside the union of the sets $\{\mathfrak{c}<\ell_{\mathfrak{B}}\}$ nothing is changed. This proves admissibility.

For the refinement, recall that on each set $C$ one has $\mathfrak{c}\le\ell_{\mathfrak{B}}$. So $\mathfrak{B}^{\mathfrak{c}}$ coincides with $\mathfrak{B}$ off $\{\mathfrak{c}<\ell_{\mathfrak{B}}\}$, and on that set it replaces the level $\ell_{\mathfrak{B}}$ by the smaller level $\mathfrak{c}$. Hence $\mathfrak{B}^{\mathfrak{c}}$ refines $\mathfrak{B}$.

\emph{Proof of (c).} For a term $T$ and a box $\square\in\mathcal{C}^1$, write $\operatorname{piece}(T,\square)$ for the piece of $T$ localised at $\square$. The localisation interpolates between the endpoints $\Phi^0$ and $\Phi^1$. It does so with one global interpolation parameter $t$, common to all boxes, and the partition of unity $\{\zeta_{\square}\}$ enters linearly, as in \cite[(3.4)]{B-CMP109}. Because the parameter is common and the dependence on $\{\zeta_{\square}\}$ is linear, summing the pieces over the boxes reconstructs the increment of $T$ between the two endpoints, term by term:
\begin{equation}\label{4.28:eq-refined-cover-2}
\sum_{\square\in\mathcal{C}^1}\operatorname{piece}(T,\square)=T(\Phi^1)-T(\Phi^0).
\end{equation}
The right-hand side of \eqref{4.28:eq-refined-cover-2} does not depend on the cover. Summing over the terms $T$ of $\mathbb{C}^{(l)}_k$, the total of $\mathbb{C}^{(l)}_k$ is therefore the same for $\mathcal{C}^1$ as for $\mathcal{C}$.

The replacement of $\mathcal{C}$ by $\mathcal{C}^1$ only changes how the pieces of the small terms are distributed over the boxes. By construction of $\mathfrak{B}^{\mathfrak{c}}$, this redistribution takes place at every localisation domain $X$ meeting $\mathsf{L}_{\mathfrak{d}}$. On the run this regrouping occurs in the stages whose level $j$ exceeds $\kappa_0$.
\end{proof}

\begin{lemma}[Thresholds at lifted points, with running factor at most one half]
\label{4.28:lem-lifted-thresholds}
Fix the level $k$ and the stage $l$ of the run at the site $k$, and let $h:=k-N_k$ be the lowest
level of that run. Let $\mathfrak{d}=(C;\kappa,\mathsf{h},\kappa_0,J)$ be a lifted attempt datum of
one of the three kinds (A), (B), (B$'$) of Definition~\ref{4.28:def-attempt-data}. Let $\ell$,
$\omega$, $\mathfrak{c}$, $g$ and $\mathsf{L}_{\mathfrak{d}}$ be its levels, weights, cutting level,
depth and lifted set, as in \eqref{4.28:eq-cutting-level-a} and \eqref{4.28:eq-cutting-level-b}.
Let $\mathfrak{B}^{\mathfrak{c}}$ and, for a cube $\square$ of the refined cover $\mathcal{C}^1$, the
box $\square_0=\tilde{\square}^5$ be as in Definition~\ref{4.28:def-refined-cover}. We write
$Q_e(y')$, $e=1,\dots,8$, for the entries of the smallness vector of the attempt space
$\mathfrak{L}_l$ at $y'$, and $Q_e^{(p,X',\mathfrak{s})}(y')$ for the entry $e$ evaluated at a triple
$(p,X',\mathfrak{s})$ of the list $(\Sigma^{\ast})$. Assume the following.
\begin{enumerate}
\item[(a)] \emph{Radii.} For $e=1,\dots,7$ and every level $n$ the radius $\mathbf{r}_{e,n}$ is a fixed
constant multiple of $\alpha_{\cdot,n}\ell_n^{-p_e}$, where
\begin{equation*}
(p_1,\dots,p_7)=(2,2,3,2,2,1,2),
\end{equation*}
and for $e=4$ the radius carries, at the index $p$, one further factor $\ell_{\min\{p,n\}}^{-1}$.
We write $r_e(n,q)$ for the radius at level $n$ in which this further factor is
$\ell_q^{-1}$, and we set $\mathbf{r}^{\mathfrak{c}}_e(y'):=\mathbf{r}_{e,\mathfrak{c}(y')}$ and
$\mathbf{r}^{\mathfrak{c},(p)}_e(y'):=r_e(\mathfrak{c}(y'),p\wedge\mathfrak{c}(y'))$.
\item[(b)] \emph{The bracket.} $\beta=1/5$ and, for $m\le k$,
$s^{(k)}_m:=1-\beta\bigl(1-2^{-(k-m)}\bigr)$. This is the bracket of
\cite[(1.64)--(1.68)]{B-CMP122a}, as it enters the arrest theorem.
\item[(c)] \emph{The running-factor identity.} For all $n$ and $m$,
\begin{equation}\label{4.28:eq-lifted-thresholds-1}
\varphi^{(n),(k)}_m=s^{(k)}_m-\beta\sum_{i=h+1}^{h+n}2^{-|m-i|}.
\end{equation}
This identity is part of the arrest theorem.
\item[(d)] \emph{The floor at the current level.} $\Phi(y')\ge\varphi_{\mathrm{arr}}$ for every
$y'\in Z^{\mathrm{on}}$, where $\Phi$ is defined in \eqref{4.28:eq-lifted-thresholds-3} below and
$\varphi_{\mathrm{arr}}$ denotes the floor of the running factor in the arrest theorem.
This bound is part of the arrest theorem.
\item[(e)] \emph{The level-free floor.} $\theta^{(l)}(n)\ge\varphi_{\mathrm{arr}}$ at every level $n$,
with $\theta^{(l)}(n)$ as in \eqref{4.28:eq-lifted-thresholds-2} below. This bound, in which no
level enters, is part of the arrest theorem.
\item[(f)] \emph{Arrest weights.} Each live piece $\mathfrak{a}$ with $k_a<k$ has
$\sigma_{\mathfrak{a}}(n)\ge0$ at every level $n$. Moreover $[h_a+1,k_a]$ is contained in the
index interval $I(\mathfrak{a})$ of $\sigma_{\mathfrak{a}}$, and $\sigma_{\mathfrak{a}}(n)\ge2$
whenever $n-1,n,n+1\in I(\mathfrak{a})$. These properties of the arrest weights are part of the
arrest theorem.
\item[(g)] \emph{Retained structures.} In kinds (B) and (B$'$) the structure $\mathfrak{s}$ of a triple
of $(\Sigma^{\ast})$ is a retained stage structure $\mathbb{B}^{(n')}_k$ of this run. In kind (A) it
is the restriction of a retained structure $\mathbb{B}^{(n')}_{k_a}$ of the list $(\Sigma^{\ast})$
of $\mathfrak{a}$. In all cases the block level $\lambda_{\mathfrak{s}}(y')$ of $\mathfrak{s}$ at $y'$
satisfies
\begin{equation*}
\mathfrak{c}(y')-1\le\lambda_{\mathfrak{s}}(y')\le\ell(y').
\end{equation*}
\item[(h)] \emph{The kind-(A) reading.} Let $\mathfrak{s}$ be a retained structure of the older piece
$\mathfrak{a}$, and let $\ell_{\mathfrak{a}}$ be the level function of its datum. Then the enclosing
run's attempt space $\mathfrak{L}_l$ reads the composite clause
\begin{equation}\label{4.28:eq-lifted-thresholds-5}
Q_e^{(p,X',\mathfrak{s})}(y')
<\theta^{(l)}(\ell_{\mathfrak{a}}(y'))\,\omega(y')\,
r_e\bigl(\ell_{\mathfrak{a}}(y'),p\wedge\lambda_{\mathfrak{s}}(y')\bigr).
\end{equation}
This is the frozen weighted unit with the bracket of $\mathfrak{L}_l$ at the frozen level. At the box
triple of a lifted point one has $\lambda_{\mathfrak{s}}=\mathfrak{c}<\ell_{\mathfrak{a}}$ and
$g\ge1$. There the conversion of units from the level $\ell_{\mathfrak{a}}$ to the level
$\mathfrak{c}$, whose margin is at least $0.419$, turns \eqref{4.28:eq-lifted-thresholds-5} into
$Q_e<(0.381/L)^g\,\mathbf{r}^{\mathfrak{c},(p)}_e(y')$, and
$(0.381/L)^g\le\varphi_{\mathrm{arr}}$. This is the kind-(A) reading of the list
$(\Sigma^{\ast})$ of the older piece.
\item[(i)] At the site $\kappa$ one has $(N_0(\kappa)+1)(N_0(\kappa)+3)\le N_\kappa$.
\end{enumerate}
Let $y'\in Z^{\mathrm{on}}$. For a level $n$ put
\begin{equation}\label{4.28:eq-lifted-thresholds-2}
\theta^{(l)}(n):=\varphi^{(l),(k)}_n
-\beta\sum_{\substack{\mathfrak{a}\ \mathrm{live},\ \mathfrak{a}\ni y'\\ k_a<k}}
\sigma_{\mathfrak{a}}(n).
\end{equation}
This quantity depends on $y'$ through the set of live pieces containing $y'$. Put further
\begin{equation}\label{4.28:eq-lifted-thresholds-3}
\Phi(y'):=\theta^{(l)}(\ell(y')),\qquad
\theta^{(l)}_{\mathfrak{s}}(y'):=\theta^{(l)}(\lambda_{\mathfrak{s}}(y')),\qquad
\theta^{(l)}(\mathfrak{c}):=\theta^{(l)}(\mathfrak{c}(y')),
\end{equation}
and
\begin{equation*}
w_{\mathfrak{s}}(y'):=L_w^{2(\lambda_{\mathfrak{s}}(y')-\mathfrak{c}(y'))^+}.
\end{equation*}
Here $L_w$ is the weight base of \cite[(1.65), (1.67)]{B-CMP122a}, with $2\le L_w$ and
$L_w^2\le L/3$. The quantity
$w_{\mathfrak{s}}(y')$ is the weight $y'$ carried, by \cite[(1.65), (1.67)]{B-CMP122a} at that
stage, when its level was $\lambda_{\mathfrak{s}}(y')$.

The attempt space $\mathfrak{L}_l$ imposes at $y'$ two families of clauses.

\emph{Direct entries} $e\in\{1,3,5,6,7\}$; the cube clause, entry eight, is read at the weight
$\omega$:
\begin{equation}\label{4.28:eq-lifted-thresholds-4}
Q_e(y')<\theta_{\mathfrak{L},e}(y'):=\Phi(y')\,\omega(y')\,\mathbf{r}_{e,\ell(y')}.
\end{equation}

\emph{Composite entries} $e\in\{2,4\}$, together with the clause bounding the background
configuration $U$. These are imposed at every triple $(p,X',\mathfrak{s})$ of $(\Sigma^{\ast})$,
hybrid triples included, in birth radii, that is, at the block level of $\mathfrak{s}$:
\begin{equation}\label{4.28:eq-lifted-thresholds-6}
Q_e^{(p,X',\mathfrak{s})}(y')
<\theta^{(l)}_{\mathfrak{s}}(y')\,w_{\mathfrak{s}}(y')\,
r_e\bigl(\lambda_{\mathfrak{s}}(y'),p\wedge\lambda_{\mathfrak{s}}(y')\bigr).
\end{equation}
The composite clauses are those of $\mathfrak{L}_l$. The weight $5\mathsf{w}^{\mathrm{on}}$ of the
companion space $\mathfrak{N}^{(l)}$ and the shell box clause $(\mathrm{box})^{\mathrm{sh}}$ enter only
on $\square_0\cap\mathfrak{T}$. This set is the un-lifted collar $\bar{c}\cap\Omega_k$ of kind (B)
at the last stage $l=N_k$, since the box $\square_0$ lies in $C$, as recorded in
Definition~\ref{4.28:def-refined-cover}.

Then the following hold.
\begin{enumerate}
\item[(1)] $\Phi(y')\in[\varphi_{\mathrm{arr}},s^{(k)}_{\ell(y')}]$ and
$\theta^{(l)}_{\mathfrak{s}}(y')\in[\varphi_{\mathrm{arr}},s^{(k)}_{\lambda_{\mathfrak{s}}(y')}]$.
\item[(2)] At the current structure, where $\lambda_{\mathfrak{s}}=\ell$, the threshold of
\eqref{4.28:eq-lifted-thresholds-6} is $\Phi\,\omega\,r_e(\ell,p\wedge\ell)$.
\item[(3)] At the box triple $\mathfrak{s}=\mathfrak{B}^{\mathfrak{c}}\restriction\square_0$ one has
$\lambda_{\mathfrak{s}}=\mathfrak{c}$ and $w_{\mathfrak{s}}=1$. Hence
\eqref{4.28:eq-lifted-thresholds-6} takes the form of the first line of the following display, and
for all three kinds its second line holds:
\begin{equation}\label{4.28:eq-lifted-thresholds-a}
\begin{gathered}
Q_e^{(p,\square_0,\mathfrak{B}^{\mathfrak{c}})}(y')
<\theta^{(l)}(\mathfrak{c}(y'))\,\mathbf{r}^{\mathfrak{c},(p)}_e(y'),\\
\varphi_{\mathrm{arr}}\le\theta^{(l)}(\mathfrak{c})\le1-2.5\beta=\tfrac12
\qquad\text{on }\mathsf{L}_{\mathfrak{d}}.
\end{gathered}
\end{equation}
In the first line the birth radius is the unweighted radius of the $\mathfrak{c}$-shell.
\item[(4)] In kind (A) the clause \eqref{4.28:eq-lifted-thresholds-5} coincides with
\eqref{4.28:eq-lifted-thresholds-6} wherever $\lambda_{\mathfrak{s}}(y')=\ell_{\mathfrak{a}}(y')$.
At a lifted point's box triple it gives
\begin{equation*}
Q_e<(0.381/L)^g\,\mathbf{r}^{\mathfrak{c},(p)}_e(y')
\le\varphi_{\mathrm{arr}}\,\mathbf{r}^{\mathfrak{c},(p)}_e(y')
\le\theta^{(l)}(\mathfrak{c})\,\mathbf{r}^{\mathfrak{c},(p)}_e(y').
\end{equation*}
Thus the first line of \eqref{4.28:eq-lifted-thresholds-a}, and the starting estimate of the box
analysis, hold at kind (A) a fortiori by the conversion of units, not literally.
\end{enumerate}
\end{lemma}

\begin{proof}
\emph{Clause (1).} For $m\le k$ the geometric sum gives
\begin{equation}\label{4.28:eq-lifted-thresholds-8}
s^{(k)}_m=1-\beta\sum_{q=m+1}^{k}2^{-(q-m)}.
\end{equation}
By \eqref{4.28:eq-lifted-thresholds-1} we have $\varphi^{(l),(k)}_n\le s^{(k)}_n$, since the
subtracted sum is non-negative. By (f) every $\sigma_{\mathfrak{a}}(n)$ in
\eqref{4.28:eq-lifted-thresholds-2} is non-negative. Hence $\theta^{(l)}(n)\le s^{(k)}_n$ at every
level $n$. At $n=\ell(y')$ this is the upper end for $\Phi$, and at
$n=\lambda_{\mathfrak{s}}(y')$ it is the upper end for $\theta^{(l)}_{\mathfrak{s}}$. The lower ends
are (d) for $\Phi$ and (e) at $n=\lambda_{\mathfrak{s}}(y')$ for $\theta^{(l)}_{\mathfrak{s}}$.

\emph{Clause (2).} If $\lambda_{\mathfrak{s}}(y')=\ell(y')$, then
$\theta^{(l)}_{\mathfrak{s}}(y')=\Phi(y')$ by \eqref{4.28:eq-lifted-thresholds-3}. By
\eqref{4.28:eq-cutting-level-b} we have $\omega=L_w^{2(\ell-\mathfrak{c})}$ and $g=\ell-\mathfrak{c}$.
Every weight in the table \eqref{4.28:eq-cutting-level-a} is at least one, so $g\ge0$. Therefore
$w_{\mathfrak{s}}=L_w^{2(\ell-\mathfrak{c})^+}=\omega$, and the threshold of
\eqref{4.28:eq-lifted-thresholds-6} is $\Phi\,\omega\,r_e(\ell,p\wedge\ell)$.

\emph{Clause (3): the weight.} At the box triple $\lambda_{\mathfrak{s}}=\mathfrak{c}$, so
$w_{\mathfrak{s}}=L_w^{0}=1$. Moreover
$r_e(\mathfrak{c},p\wedge\mathfrak{c})=\mathbf{r}^{\mathfrak{c},(p)}_e$ by (a), and
$\theta^{(l)}_{\mathfrak{s}}=\theta^{(l)}(\mathfrak{c})$. This is the first line of
\eqref{4.28:eq-lifted-thresholds-a}.

We next check that $w_{\mathfrak{s}}$ is the weight of \eqref{4.28:eq-cutting-level-a} at the
stage of $\mathfrak{s}$, in kind (B). Let $\mathfrak{s}=\mathbb{B}^{(n')}_k$ with top level
$j'=h+n'\in\,]\kappa_0,J[$. The table applies at that stage with $J$ replaced by $j'$; the cutting
level $\mathfrak{c}$ is that of the current stage. We treat the points row by row.
\begin{itemize}
\item A band point lay at stage $n'$ in that stage's top set, at level $\lambda=j'$ with weight
$L_w^{2(j'-\kappa_0)}$. Its cutting level is $\mathfrak{c}=\kappa_0$. Hence this weight is
$L_w^{2(\lambda-\mathfrak{c})}$.
\item A thin-stratum point of $\Gamma''_m$ has $\mathfrak{c}=\kappa_0$. If $\kappa_0<m<j'$, it had
level $\lambda=m$ and weight $L_w^{2(m-\kappa_0)}=L_w^{2(\lambda-\mathfrak{c})}$. If $m\ge j'$, it lies
outside $Z''_{j'}$, since $\Gamma''_m=Z''_{m+1}\setminus Z''_m$ and $Z''_{j'}\subset Z''_m$, and is
read in the band row with $J$ replaced by $j'$: level $j'$ and weight
$L_w^{2(j'-\kappa_0)}$, as for a band point.
\item An old-shell point of $\Omega_m\setminus\Omega_{m+1}$, with $\kappa_0<m\le J$, has
$\mathfrak{c}=m$. If $m\le j'$, it had level $\lambda=j'$ and weight
$L_w^{2(j'-m)}=L_w^{2(\lambda-\mathfrak{c})}$. If $j'<m$, it had level $\lambda=m=\mathfrak{c}$ and
weight $1=L_w^{2(\lambda-\mathfrak{c})^+}$.
\end{itemize}

\emph{Clause (3): the range of $\mathfrak{c}$ on $\mathsf{L}_{\mathfrak{d}}$.} By
\eqref{4.28:eq-cutting-level-b}, $y'$ is lifted if and only if $g(y')\ge1$. We read the rows of
\eqref{4.28:eq-cutting-level-a}.
\begin{itemize}
\item On the first row and on the old shells with $J<m<\kappa$ one has $\mathfrak{c}=\ell$, so
$g=0$.
\item On the thin strata and on the band one has $\mathfrak{c}=\kappa_0$.
\item On an old shell with $\kappa_0<m\le J$ one has $\mathfrak{c}=m$ and $g=J-m$, and $g\ge1$ if
and only if $m\le J-1$.
\end{itemize}
Since $J>\kappa_0$ for a lifted datum, every lifted $y'$ has
\begin{equation*}
c:=\mathfrak{c}(y')\in[\kappa_0,J-1].
\end{equation*}
Further, $\kappa_0-\mathsf{h}=N_\kappa-N_0(\kappa)$. By (i),
$N_\kappa\ge(N_0(\kappa)+1)(N_0(\kappa)+3)\ge N_0(\kappa)+3$, hence
\begin{equation}\label{4.28:eq-lifted-thresholds-9}
\kappa_0\ge\mathsf{h}+2.
\end{equation}
Also $J\le\kappa$ gives $J-1\le\kappa-1$.

\emph{Kind (B), the run.} Here $\kappa=k$, $\mathsf{h}=h$ and $J=h+l$. By the placement of live
older pieces in Definition~\ref{4.28:def-attempt-data}, every such piece lies in $Z_h=Z''_h$. Since
$h\le\kappa_0+1$, we have $Z''_h\subset Z''_{\kappa_0+1}$. As
$\mathsf{L}_{\mathfrak{d}}=(\Omega^c_J\setminus Z''_{\kappa_0+1})\cap C$, no live older piece meets
$\mathsf{L}_{\mathfrak{d}}$, and the sum in \eqref{4.28:eq-lifted-thresholds-2} is empty. By
\eqref{4.28:eq-lifted-thresholds-1} and \eqref{4.28:eq-lifted-thresholds-8},
\begin{equation}\label{4.28:eq-lifted-thresholds-10}
\theta^{(l)}(c)=\varphi^{(l),(k)}_c
=1-\beta\sum_{i=h+1}^{h+l}2^{-|c-i|}-\beta\sum_{q=c+1}^{k}2^{-(q-c)}.
\end{equation}
The first sum contains the terms $i=c-1,c,c+1$. Indeed $c-1\ge\kappa_0-1\ge h+1$ by
\eqref{4.28:eq-lifted-thresholds-9}, and $c+1\le J=h+l$. These terms contribute
$\tfrac12+1+\tfrac12=2$. The second sum contains the term $q=c+1$, since $c+1\le J\le k$, and this
term contributes $\tfrac12$. All terms are non-negative. Hence the subtracted quantity is at least
$2.5\beta=\tfrac12$, and $\theta^{(l)}(c)\le\tfrac12$ for every $c\in[\kappa_0,J-1]$.

\emph{Kind (A), the core piece.} Here $\kappa=k_a\le h$, $\mathsf{h}=h_a$, $\kappa_0=k_0^a$ and
$J=j^+_a$. The attempt space is that of the arrest, so $\mathfrak{a}$ is among the live pieces
containing $y'\in\mathsf{L}_{\mathfrak{d}}\subset C=X_{\mathfrak{a}}$. Drop the other,
non-negative, $\sigma$-terms in \eqref{4.28:eq-lifted-thresholds-2} and use
$\varphi^{(l),(k)}_c\le s^{(k)}_c$. This gives
\begin{equation*}
\theta^{(l)}_{\mathfrak{s}}(y')=\theta^{(l)}(c)\le s^{(k)}_c-\beta\sigma_{\mathfrak{a}}(c).
\end{equation*}
By \eqref{4.28:eq-lifted-thresholds-9} and $J\le\kappa$, the levels $c-1,c,c+1$ lie in
$[k_0^a-1,j^+_a]\subset[h_a+1,k_a]$, which lies in $I(\mathfrak{a})$ by (f). Hence
$\sigma_{\mathfrak{a}}(c)\ge2$ by (f). Next, $c\le j^+_a-1\le k_a-1<k$. So $k-c\ge1$, hence
$1-2^{-(k-c)}\ge\tfrac12$ and $s^{(k)}_c\le1-\beta/2$ by (b). Therefore
\begin{equation*}
\theta^{(l)}(c)\le1-\tfrac12\beta-2\beta=1-2.5\beta=\tfrac12
\qquad\text{on }[k_0^a,j^+_a-1]\subset[h_a+2,k_a].
\end{equation*}

\emph{Kind (B$'$), an excluded region of this site.} Here $\kappa=k$, $\mathsf{h}=h$ and
$J=h+n^+\le h+l-1$, where $n^+$ is the stage of Definition~\ref{4.28:def-attempt-data} attached to
the excluded region. For $l>n^+$ the running factor is that of the sibling attempt space,
$\varphi^{(l),(k)}$. Pieces of the same class add no $\sigma$-term, by the nesting condition of
Definition~\ref{4.28:def-attempt-data}. Live older pieces do not meet $\mathsf{L}_{\mathfrak{d}}$, as
in kind (B). So \eqref{4.28:eq-lifted-thresholds-10} holds, with the first sum running to
$h+l\ge J+1$. It therefore contains $c-1,c,c+1$, using $c-1\ge h+1$ by
\eqref{4.28:eq-lifted-thresholds-9}. The second sum contains $q=c+1$, because $c+1\le J<k$. As in
kind (B), $\theta^{(l)}(c)\le\tfrac12$.

\emph{Lower end.} By (e), $\theta^{(l)}(\mathfrak{c})\ge\varphi_{\mathrm{arr}}$, whatever the level.
This completes the second line of \eqref{4.28:eq-lifted-thresholds-a}.

\emph{Clause (4).} Suppose $\lambda_{\mathfrak{s}}(y')=\ell_{\mathfrak{a}}(y')$. In kind (A) the
level function of the datum is $\ell=\ell_{\mathfrak{a}}$, so clause (2) gives
$\theta^{(l)}_{\mathfrak{s}}=\theta^{(l)}(\ell_{\mathfrak{a}})$ and $w_{\mathfrak{s}}=\omega$. The two
radii $r_e(\ell_{\mathfrak{a}},p\wedge\lambda_{\mathfrak{s}})$ agree. Hence
\eqref{4.28:eq-lifted-thresholds-5} and \eqref{4.28:eq-lifted-thresholds-6} coincide.

At a lifted point's box triple, (h) gives
$Q_e<(0.381/L)^g\,\mathbf{r}^{\mathfrak{c},(p)}_e(y')$ and
$(0.381/L)^g\le\varphi_{\mathrm{arr}}$. The lower end in the second line of
\eqref{4.28:eq-lifted-thresholds-a} gives
$\varphi_{\mathrm{arr}}\le\theta^{(l)}(\mathfrak{c})$. Together these give the displayed chain, and
so the first line of \eqref{4.28:eq-lifted-thresholds-a} at kind (A).
\end{proof}

\begin{lemma}[The box lemma at lifted points: scales, conversion and rooms]\label{4.28:lem-box-lemma}
Let $\mathfrak{d}=(C;\kappa,\mathsf{h},\kappa_0,J)$ be an attempt datum of the step at the site $k$, with
$\mathbf{R}$-site $\kappa$. Let its level
function $\ell$, weight $\omega$, cutting level $\mathfrak{c}$, depth $g=\ell-\mathfrak{c}$ and lifted set
$\mathsf{L}_{\mathfrak{d}}$ be as in Definition~\ref{4.28:def-cutting-level}. Let the refined block
structure $\mathfrak{B}^{\mathfrak{c}}$ and the cover $\mathcal{C}^1$ be as in
Definition~\ref{4.28:def-refined-cover}; for $\square\in\mathcal{C}^1$ this gives the box
$\square_0=\tilde{\square}^5$, the structure
$\mathfrak{s}=\mathfrak{B}^{\mathfrak{c}}\restriction\square_0$ and its top level $p_\square$. Let the
radii $\mathbf{r}_{e,n}$, $\mathbf{r}^{\mathfrak{c}}_e$, $\mathbf{r}^{\mathfrak{c},(p)}_e$ with exponents
$p_e$, the running factor $\Phi$, the factor $\theta^{(l)}(\mathfrak{c})$ and the attempt space
$\mathfrak{L}_l$ be as in Lemma~\ref{4.28:lem-lifted-thresholds}. Lengths are counted in
\emph{$\pi_n$-sides}, the side $M\ell_n$ of a cube of the family $\pi_n$; the spacing $\ell_n$ grows by
the factor $L$ from one level to the next. Assume the following.
\begin{itemize}
\item $R_{\min}\ge 5L$ and $L_0^2\le L/3$, where $L_0$ denotes the weight base of the weight
function $\omega$ of Definition~\ref{4.28:def-cutting-level}; $\beta_0=1/7$, and the blocking factor
satisfies the one-sided condition $(1+\beta_0)/L\le 2/21$. Moreover
$(N_0(\kappa)+1)(N_0(\kappa)+3)\le N_\kappa$ at the $\mathbf{R}$-site $\kappa$ of $\mathfrak{d}$, so that
$\kappa_0=\kappa-N_0(\kappa)$ and $\mathsf{h}=\kappa-N_\kappa$; in particular at $\kappa=k$ for the datum
of the run at the site $k$. Here $N_\kappa$ is the number of stages of the run at the $\mathbf{R}$-site
$\kappa$ (not the rank of the gauge group); below we abbreviate $N_0:=N_0(\kappa)$ and $N:=N_\kappa$.
\item The adjacent-ratio bound for the radii: for every level $n$,
\begin{equation*}
\alpha_{\cdot,n+1}\le(1+\beta_0)\,\alpha_{\cdot,n},\qquad
\alpha_{\cdot,n}\le(1+\beta_0)\,\alpha_{\cdot,n+1}.
\end{equation*}
\item The widths statement for the shells of $\mathfrak{B}^{\mathfrak{c}}$. Off the lifted data, every
$\mathfrak{c}$-shell $m''$ of $\mathfrak{B}^{\mathfrak{c}}$ is at least $R_{\min}$ sides of a
$\pi_{m''+1}$-cube wide. A shell made by $\mathbf{T}$ contains
$(Z'_{m''})^{\sim 8}\setminus Z'_{m''}$ by \cite[(3.5)]{B-CMP119}, for every choice of the sets entering
that display; an unlifted shell made by $\mathbf{R}$ contains one layer of $MR_{m''+1}$-cubes.
\item The box radius constant $\mathsf{a}$ is a power of $L$, fixed after $L$ and before $\gamma$, with
no upper bound. It is subject only to the one-sided condition
\begin{equation*}
\mathsf{a}\ge(1+2\beta)(1+\beta_0)L^3 .
\end{equation*}
\item The rule for the composite radii: at a retained triple the composite entries carry the radii
$r_e(\lambda_{\mathfrak{s}},p\wedge\lambda_{\mathfrak{s}})$; for the terms $T$ near $\square$ in the
sense of the next item this rule gives $p_\square\ge\ell_T$; and it comes with an inequality form,
through which the ratios obtained for the direct entries apply to the composite entries.
\item The threshold form of the clauses of near terms. A term
$T\in\mathcal{T}^{\mathbb{E}\mathbb{R}}$ is \emph{near} $\square$ if $X_T\subset\tilde{\square}^2$. The
thresholds of the clauses of such a $T$ have the form $\theta_{T,a}=s_T\mathbf{r}^T_a$, with
$\mathbf{r}^T_a=\mathbf{r}_{a,\ell_T}$. Here $s_T=1$ at points $y$ with $\ell_T=\ell(y)$, and
$s_T\ge 1-\beta$ everywhere.
\item The quantity $\alpha_3(T)$ is defined through a maximum whose first member is $24/\beta$, the
first member of the maximum in \cite[(3.54)]{B-CMP109}, and in which $\operatorname{room}_T$ enters
through $12/\operatorname{room}_T$.
\item Every cube of Ba{\l}aban's family for $\mathfrak{B}^{\mathfrak{c}}$ contained in $\mathfrak{T}$
satisfies the cube clause, entry $8$, of the companion space $\mathfrak{N}^{(l)}$, with factor at most
$1+2\beta$.
\end{itemize}
Let $\square\in\mathcal{C}^1$ have scale $i$, and suppose that $\square_0$ meets
$\mathsf{L}_{\mathfrak{d}}$. Then the following hold.
\begin{itemize}
\item[(a)] $\square_0$ meets at most two $\mathfrak{c}$-shells; they are adjacent, and in every case the
shell $i$ is one of them. Moreover $i\le J$.
\item[(o)] One has
\begin{equation}\label{4.28:eq-box-lemma-c}
[\square_0]_{\sharp}\subset Z_J\cap C .
\end{equation}
Hence, for a datum of kind (A), $[\square_0]_{\sharp}\subset X_{\mathfrak{a}}\subset Z_h$, where $h$ is
the base level of the run at the site $k$; here the kinds (A), (B), (B$'$) of an attempt datum are those
distinguished in Lemma~\ref{4.28:lem-lifted-thresholds}: the datum of a live arrested older piece
$\mathfrak{a}$, the run at the site $k$ itself, and an excluded region of that site. The box
$\square_0$ meets $\mathfrak{T}$ only for kind (B) at stage $l=N$, and then inside
$\bar{c}\cap\Omega_k$.
\item[(a$''$)] Let $T$ be near $\square$. Then
\begin{equation*}
\ell_T\le i,\qquad \ell_T\le\mathfrak{c}\ \text{ on }\square_0,\qquad \ell_T\le\kappa-1 .
\end{equation*}
\item[(a$'$)] $\ell$ takes at most two values on $\square_0$, and if two, they are adjacent; the same
holds for $\ell_{\mathfrak{B}}$.
\item[(c$_0$)] For every $y'\in\square_0\cap C$ and every entry $e\le 7$ (for $e\in\{2,4\}$, at every
value of $p$),
\begin{equation}\label{4.28:eq-box-lemma-a}
\omega(y')\,\mathbf{r}_{e,\ell(y')}\le 0.381^{\,g(y')}\,\mathbf{r}^{\mathfrak{c}}_e(y') .
\end{equation}
Consequently, on $\square_0\cap Z^{\mathrm{on}}$ the following hold.
\begin{itemize}
\item The direct clauses of $\mathfrak{L}_l$ (entries $1,3,5,6,7$; the cubes are treated in (e)) imply
the unweighted direct clauses of level $\mathfrak{c}$, with factor at most $\Phi\cdot 0.381^{g}\le 1$.
\item The composite clause at a box triple
$(p,\square_0,\mathfrak{B}^{\mathfrak{c}}\restriction\square_0)$ with $p\le p_\square$, together with
the bound on the configuration that accompanies the composite clauses, is already written with the radii
of the $\mathfrak{c}$-shell. This is the birth radius with $\lambda=\mathfrak{c}$ and $w_{\mathfrak{s}}=1$.
Its factor is $\theta^{(l)}(\mathfrak{c})\le\frac12$ where $g\ge1$, and $\Phi\le1$ where $g=0$.
\item For kinds (B) and (B$'$) no conversion is involved. For kind (A) at $g\ge1$ the composite clause
follows by converting the frozen weighted radius, which is at most
$(0.381/L)^{g}\,\mathbf{r}^{\mathfrak{c},(p)}_e$.
\end{itemize}
The bound \eqref{4.28:eq-box-lemma-a} serves the direct entries only.
\item[(c)] The box radius constants $\mathsf{b}^{(i)}_{\cdot}$ at scale $i$ satisfy
\begin{equation}\label{4.28:eq-box-lemma-b}
\mathsf{b}^{(i)}_{\cdot}\le(1+2\beta)(1+\beta_0)L^{\bar p}\,\alpha_{\cdot,i}\le\mathsf{a}\,\alpha_{\cdot,i}.
\end{equation}
Here $\bar p$ is the largest $p_e$ ($\le 3$) among the entries converted. For the converted entries $6$
and $7$, $\bar p\le 2$; for the converted entries $1$ to $4$, $\bar p=3$.
\item[(e)] Let $q'\subset\square_0$ be a cube of Ba{\l}aban's family for $\mathfrak{B}^{\mathfrak{c}}$.
Put $\nu:=\min_{q'}\mathfrak{c}$, and let the side of $q'$ be $C_{q'}M\ell_\nu$ with $C_{q'}\le2$. Then
$q'$ carries a gauge transformation $u$ with values in the gauge group $G$ such that the transformed
configuration satisfies $U^u=e^{i\eta A}$ on $q'$, where $U$ is the configuration on which the cube
clause is imposed and $\eta$ is its lattice spacing, and
\begin{equation*}
\ell_\nu|A|,\ \ell_\nu^2|\nabla^\eta A|<(1+2\beta)BC_{q'}M\alpha_{0,\nu} ,
\end{equation*}
where $B$ is the constant of the cube clause of the attempt space $\mathfrak{L}_l$.
\item[(r)] Let $T\in\mathcal{T}^{\mathbb{E}\mathbb{R}}$ be near $\square$, let
$(y,a)\in\operatorname{Cl}(T)$ with $y\in Z^{\mathrm{on}}$ and $a\notin\{5,8\}$, and put
\begin{equation*}
\operatorname{room}^{\mathrm{on}}_a(y)
:=\frac{(\theta_{T,a}-\theta_{\mathfrak{L},a})(y)}{\mathbf{r}^T_a(y)} .
\end{equation*}
Then
\begin{equation*}
\operatorname{room}^{\mathrm{on}}_a(y)\ge
\begin{cases}
\beta/2 & \text{if } g(y)=0 \text{ and } \ell_T=\ell(y),\\
0.7 & \text{if } g(y)=0 \text{ and } \ell_T<\ell(y),\\
0.419 & \text{if } g(y)\ge1 \text{ and } a \text{ is a direct entry.}
\end{cases}
\end{equation*}
Now consider the composite entries $a'\in\{2,4\}$ of the starting bound at the box triple
$(p_\square,\square_0,\mathfrak{B}^{\mathfrak{c}}\restriction\square_0)$. For these put
\begin{equation*}
\theta_{\mathfrak{L},a'}(y):=\theta^{(l)}(\mathfrak{c}(y))\,\mathbf{r}^{\mathfrak{c},(p_\square)}_{a'}(y).
\end{equation*}
Then
\begin{equation*}
\operatorname{room}^{\mathrm{on}}_{a'}(y)\ge
\begin{cases}
s_T-\tfrac12\ge 0.3 & \text{if } g(y)\ge1 \text{ and } \ell_T=\mathfrak{c}(y),\\
0.75 & \text{if } g(y)\ge1 \text{ and } \ell_T<\mathfrak{c}(y).
\end{cases}
\end{equation*}
No age, and neither $N$ nor $N_0$, enters these bounds. For $T$ with a clause at a point of
$Z^{\mathrm{on}}$, the stage-free number
\begin{equation}\label{4.28:eq-box-lemma-d}
\operatorname{room}^{\dagger}_T:=\frac{\beta}{2}=\beta\hat{c},\qquad \hat{c}=\tfrac12,
\end{equation}
which is at most $\operatorname{room}_T$ and at most every $\operatorname{room}^{\mathrm{on}}_a$,
replaces $\operatorname{room}_T$ in $\alpha_3(T)$.
\end{itemize}
\end{lemma}

\begin{proof}
\emph{Clause (a).} Every $\mathfrak{c}$-shell $m''$ of $\mathfrak{B}^{\mathfrak{c}}$ is at least
$R_{\min}$ sides of a $\pi_{m''+1}$-cube wide. On lifted data this is
Lemma~\ref{4.28:lem-cutting-strata}(ii), which gives at least $3R_{\min}$ such sides. Elsewhere it is the
widths statement: shells made by $\mathbf{T}$ contain $(Z'_{m''})^{\sim8}\setminus Z'_{m''}$ by
\cite[(3.5)]{B-CMP119}, and unlifted shells made by $\mathbf{R}$ contain one layer of
$MR_{m''+1}$-cubes. The condition $(1+\beta_0)/L\le2/21$ gives $5L\ge60$, so $R_{\min}\ge60$, and every
such shell is more than $12$ sides of a $\pi_{m''+1}$-cube wide.

The box $\square_0$ is connected, so the values of $\mathfrak{c}$ on it form an interval of integers,
and this interval contains $i$. Suppose it contained at least three values. Take three consecutive ones
containing $i$, and let $m''$ be the middle one. Then $m''+1\ge i$. The box $\square_0$, which extends
over at most $12$ $\pi_i$-sides, would cross the whole shell $m''$ within $12$ $\pi_i$-sides. Since
$m''+1\ge i$, these are at most $12$ $\pi_{m''+1}$-sides, which contradicts the width just established.
So there are at most two values, and they are adjacent.

At a point $y'$ of $\mathsf{L}_{\mathfrak{d}}$ we have $g(y')\ge1$, and $\ell(y')\le J$ by the rows of
\eqref{4.28:eq-cutting-level-a}. Hence $\mathfrak{c}(y')=\ell(y')-g(y')\le J-1$, so $\square_0$ holds a
value of $\mathfrak{c}$ at most $J-1$. Since the values on $\square_0$ are at most two, adjacent, and
include $i$, it follows that $i\le J$.

\emph{Clause (o).} Since $i\le J$ and $\ell\le J$ there, the set $[\square_0]_{\sharp}$ lies within
$13$ $\pi_J$-sides of $\Omega^c_J\cap C$. By the inclusion $\supseteq$ in
\eqref{4.28:eq-cutting-strata-a}, $Z_J\supseteq(\Omega^c_J)^{\sim2}$; this adds two layers of
$MR_J$-cubes, that is $2R_J\ge 2R_{\min}>13$ $\pi_J$-sides. Hence $[\square_0]_{\sharp}\subset Z_J\cap C$,
which is \eqref{4.28:eq-box-lemma-c}. For $J<\kappa$ one has
\begin{equation*}
Z_J\subset\Omega^c_{J+1}\subset\Omega^c_\kappa .
\end{equation*}
Hence, when $J<\kappa$, $[\square_0]_{\sharp}\subset\Omega^c_\kappa$.
For kind (B$'$) one has $J\le k-1$.

\emph{Clause (a$''$).} The set $X_T$ is a union of $\pi_{\ell_T}$-cubes, and it lies inside
$\tilde{\square}^2$, that is, inside six $\pi_i$-sides. A $\pi_{\ell_T}$-cube with $\ell_T\ge i+1$ has
side at least $L$ $\pi_i$-sides, and $L>6$, since $(1+\beta_0)/L\le2/21$; so
$\ell_T\le i$. Next, $X_T\subset\Lambda_{\ell_T}$ by
\cite{B-CMP119}, and $X_T\subset C\subset\Lambda^c_\kappa$. The regions $\Lambda_m$ are nested, so
$\ell_T<\kappa$. It remains to show $\ell_T\le\mathfrak{c}$ on $\square_0$; there are three cases.
\begin{itemize}
\item If $\ell_T<i$, then $\mathfrak{c}\ge i-1\ge\ell_T$ on $\square_0$ by (a).
\item If $\ell_T=i>\mathsf{h}$, then $\square_0$ reaches at most $9$ $\pi_i$-sides from
$X_T\subset\Lambda_i$. Since $9<2R_i$, Lemma~\ref{4.28:lem-cutting-strata}(iii) shows that $\square_0$
does not meet $\{\mathfrak{c}<i\}$.
\item If $\ell_T\le\mathsf{h}$, then $\mathsf{h}<\kappa_0-1\le\mathfrak{c}$ on $\square_0$. The first
inequality is $\kappa_0=\kappa-N_0\ge\mathsf{h}+2$, which follows from the hypothesis
$(N_0+1)(N_0+3)\le N$, as in the proof of Lemma~\ref{4.28:lem-lifted-thresholds}. The second
inequality holds by (a), because $\square_0$ contains a lifted point, and there $\mathfrak{c}\ge\kappa_0$
by the rows of \eqref{4.28:eq-cutting-level-a}.
\end{itemize}

\emph{Clause (a$'$).} Off $\mathsf{L}_{\mathfrak{d}}$ we have $g=0$, that is $\ell=\mathfrak{c}$, and (a)
applies. On an old shell $m\in(\kappa_0,J]$ we have $\ell\equiv J$. Inside $\{\mathfrak{c}=\kappa_0\}$,
every stratum of $\ell$ is at least $L^2>12$ $\pi_{\kappa_0}$-sides thick, by
Lemma~\ref{4.28:lem-cutting-strata}(iv). A box of scale $\kappa_0$ therefore meets at most two of these
strata, and they are consecutive. A box of scale $\kappa_0+1$ reaches at most $7$ of its own sides into
a band of width $5L^{N_0-1}\ge5L^2$, so it meets level $J$ only. For $\ell_{\mathfrak{B}}$ the same
argument applies with $J-1$ in place of $J$.

\emph{Clause (c$_0$).} Let $y'\in\square_0\cap C$. Write $n:=\ell(y')$ and $g=g(y')$, so that
$n=\mathfrak{c}(y')+g$. Since the radii are proportional to $\alpha_{\cdot,n}\ell_n^{-p_e}$,
\begin{equation}\label{4.28:eq-box-lemma-1}
\frac{\mathbf{r}_{e,n}}{\mathbf{r}_{e,\mathfrak{c}}}
\le\frac{\alpha_{\cdot,n}}{\alpha_{\cdot,\mathfrak{c}}}\,L^{-p_eg}
\le\Bigl(\frac{1+\beta_0}{L}\Bigr)^{g}.
\end{equation}
For $e\ne4$ the first relation follows from the proportionality of the radii; for $e=4$ the
length factor is
\begin{equation*}
\frac{\ell_n^{-2}\ell_{p\wedge n}^{-1}}{\ell_{\mathfrak{c}}^{-2}\ell_{p\wedge\mathfrak{c}}^{-1}}
\le L^{-2g}.
\end{equation*}
The second inequality in \eqref{4.28:eq-box-lemma-1} uses $L^{-p_eg}\le L^{-g}$, since $p_e\ge1$, and
the adjacent-ratio bound applied to the $g$ adjacent ratios between the levels $\mathfrak{c}$ and $n$.

From $\mathfrak{c}=\ell-\frac12\log_{L_0}\omega$ we get $\omega=L_0^{2g}$. By the hypothesis
$L_0^2\le L/3$ this gives $\omega(y')\le(L/3)^g$. Combining with \eqref{4.28:eq-box-lemma-1},
\begin{equation*}
\omega(y')\,\mathbf{r}_{e,n}
\le\Bigl(\frac{1+\beta_0}{3}\Bigr)^{g}\mathbf{r}^{\mathfrak{c}}_e(y'),
\end{equation*}
and
\begin{equation*}
\frac{1+\beta_0}{3}=\frac{8/7}{3}=\frac{8}{21}\le 0.381 .
\end{equation*}
This proves \eqref{4.28:eq-box-lemma-a}; it is the weight bound of \cite{B-CMP122a} read at a single
point.

Finally, $\Phi\le1$ by Lemma~\ref{4.28:lem-lifted-thresholds}. For a direct entry $e$, the clause
$Q_e<\Phi\,\omega\,\mathbf{r}_{e,\ell}$ of $\mathfrak{L}_l$ then gives
$Q_e<\Phi\cdot0.381^g\,\mathbf{r}^{\mathfrak{c}}_e$, with factor at most $1$. The composite clause at the
box triple is the birth-radius clause of Lemma~\ref{4.28:lem-lifted-thresholds} with $\lambda=\mathfrak{c}$
and $w_{\mathfrak{s}}=1$. Its factor is $\theta^{(l)}(\mathfrak{c})\le\frac12$ at $g\ge1$ by
\eqref{4.28:eq-lifted-thresholds-a}, and $\Phi\le1$ at $g=0$.
For kind (A) at $g\ge1$, the frozen weighted radius of the older piece gives at the box triple
\begin{equation*}
Q_e<\Bigl(\frac{0.381}{L}\Bigr)^{g}\mathbf{r}^{\mathfrak{c},(p)}_e
\le\varphi_{\mathrm{arr}}\,\mathbf{r}^{\mathfrak{c},(p)}_e
\le\theta^{(l)}(\mathfrak{c})\,\mathbf{r}^{\mathfrak{c},(p)}_e ,
\end{equation*}
the last inequality by Lemma~\ref{4.28:lem-lifted-thresholds}, so that the composite clause holds at
kind (A) as well.

\emph{Clause (c).} By (a) and (c$_0$), a cube of $\tilde{\square}^4$ obeys the unweighted clause of some
level $\mathfrak{c}\in\{i-1,i,i+1\}$ with factor at most $1$. On $\mathfrak{T}$ the factor is at most
$1+2\beta$, by the cube clause of the companion space $\mathfrak{N}^{(l)}$. Measured in the radii of
level $i$, the factor is as follows.
\begin{itemize}
\item At the cube's own level it is at most $1+2\beta$.
\item At the coarser level it is smaller.
\item At the finer level it is at most $(1+2\beta)(1+\beta_0)L^{p_e}$, with $p_e\le3$.
\end{itemize}
For the finer level two facts are used. First, the radii are proportional to $\ell^{-p_e}$, with the
exponents listed in Lemma~\ref{4.28:lem-lifted-thresholds}: entry $3$ has exponent $3$; entry $4$ has
exponent $2$, with a further power $1$ through $\ell_{p\wedge n}$; entry $6$ has exponent $1$; all other
entries have exponent $2$. Second, the adjacent-ratio bound is used in the direction lower level over
higher level, $\alpha_{\cdot,i-1}/\alpha_{\cdot,i}\le1+\beta_0$. Hence $\bar p=3$ at most, and
$\bar p\le2$ for entries $6$ and $7$. This gives the first inequality in \eqref{4.28:eq-box-lemma-b}.
The second follows from the one-sided condition on $\mathsf{a}$, since $L^{\bar p}\le L^3$.

\emph{Clause (e).} Put $n':=\min_{q'}\ell$. Since $g\ge0$, $n'\ge\nu$. There are two cases.

\emph{Case $n'=\nu$.} A point of $q'$ with $\ell=\nu=\mathfrak{c}$ has $g=0$, hence $\omega=1$. Moreover
$q'\subset\{\ell\ge\nu\}$, and by (a$'$) $q'$ misses $\{\ell\ge\nu+2\}$. So $q'$ is a cube of
Ba{\l}aban's family for the attempt space at level $\nu$, with weight $1$. The weight is read from
\cite[(1.65)]{B-CMP122a} on the rows $\Gamma''_m$, $\mathsf{h}\le m\le\kappa_0$, of
\eqref{4.28:eq-cutting-level-a}, from \cite[(1.67)]{B-CMP122a} on the rows of
\eqref{4.28:eq-cutting-level-a} with $\ell=J$, and from \cite[(1.69)]{B-CMP122a}. The cube clause of
$\mathfrak{L}_l$ at $q'$ has factor $\Phi\le1<1+2\beta$, and it gives the bound of (e).

\emph{Case $n'>\nu$.} Here $2M\ell_\nu<M\ell_{n'}$. The $\pi_{n'}$-cubes meeting $q'$ form a box
$\mathfrak{b}$ with
\begin{equation*}
\mathfrak{b}\subset\{\ell\ge n'\}\cap[\square_0]_{\sharp}\subset Y ,
\end{equation*}
where $Y$ denotes the region on which the attempt space imposes its cube clauses.
Choose a cube $Q$ of side $M\ell_{n'}$ with $q'\subset Q\subset\mathfrak{b}$. Then $Q$ belongs to
Ba{\l}aban's family at level $n'$ with side constant $C_Q=1$, since no alignment is required of these
cubes. Its weight is one of two
kinds.
\begin{itemize}
\item It is $L_0^{2(n'-\kappa_0)}$, by \cite[(1.65)]{B-CMP122a}; in this case $q'$ holds a point with
$\mathfrak{c}=\kappa_0\ge\nu$.
\item It is $L_0^{2(J-\min_Q\mathfrak{c})}$, by \cite[(1.67)]{B-CMP122a}; in this case each cube of
$\mathfrak{b}$ lies in one old shell, by Lemma~\ref{4.28:lem-cutting-strata}(v), and meets $q'$.
\end{itemize}
In both cases the weight $w_Q$ of $Q$ satisfies $w_Q\le L_0^{2(n'-\nu)}$. Restricting to $q'$ the $u$
carried by $Q$ gives, for $j=1,2$,
\begin{equation*}
\ell_\nu^{j}|\nabla^{j-1}A|
=L^{-j(n'-\nu)}\ell_{n'}^{j}|\nabla^{j-1}A|
<0.381^{\,n'-\nu}BM\alpha_{0,\nu}.
\end{equation*}
On $\mathfrak{T}$ the bound comes from the cube clause, entry $8$, of $\mathfrak{N}^{(l)}$.

\emph{Clause (r).} By (a$''$), $\ell_T<\kappa\le k$. So $T$ carries Ba{\l}aban's one-scale small-field
condition on $X_T$ at level $\ell_T$, imposed before the attempt, and the attempt space adds no further
condition to it. By the threshold form of the clauses of near terms,
\begin{equation*}
\theta_{T,a}=s_T\mathbf{r}^T_a,\qquad \mathbf{r}^T_a=\mathbf{r}_{a,\ell_T}.
\end{equation*}
Here $s_T=1$ where $\ell_T=\ell(y)$, and $s_T\ge1-\beta=0.8$ everywhere, since $\beta=1/5$. We treat the
direct entries first, in three cases; write $m:=\ell(y)$.

\emph{Case $g(y)=0$ and $\ell_T=m$.} Here $m=\ell_T<k$, and, since $\omega(y)=1$,
$\theta_{\mathfrak{L}}\le s^{(k)}_m\mathbf{r}_{a,m}$, with $1-s^{(k)}_m\ge\beta/2$. As $s_T=1$, the room
is at least $1-s^{(k)}_m\ge\beta/2$.

\emph{Case $g(y)=0$ and $\ell_T<m$.} By the adjacent-ratio bound and the condition
$(1+\beta_0)/L\le2/21$,
\begin{equation*}
\theta_{\mathfrak{L}}\le\mathbf{r}_{a,m}\le\frac{1+\beta_0}{L}\,\mathbf{r}^T_a\le0.1\,\mathbf{r}^T_a .
\end{equation*}
The room is therefore at least $s_T-0.1\ge0.7$.

\emph{Case $g(y)\ge1$, direct entries.} By (c$_0$) and $\ell_T\le\mathfrak{c}(y)$, which holds by
(a$''$),
\begin{equation*}
\theta_{\mathfrak{L}}\le0.381\,\mathbf{r}^{\mathfrak{c}}_a\le0.381\,\mathbf{r}^T_a .
\end{equation*}
The room is therefore at least $0.8-0.381=0.419$.

We turn to the composite entries $a'\in\{2,4\}$ at the box triple. Their clause in the attempt space is
the birth radius with $\lambda=\mathfrak{c}$ and $w_{\mathfrak{s}}=1$
(Lemma~\ref{4.28:lem-lifted-thresholds}), and
$p_\square\ge\ell_T$ by the rule for the composite radii. For a composite entry the radius
$\mathbf{r}^T_{a'}$ of $T$ depends on the level $p$ of the triple; we write $\mathbf{r}^{T,(p)}_{a'}$ for
it, that is, the radius $\mathbf{r}^{\mathfrak{c},(p)}_{a'}$ with the level $\mathfrak{c}$ replaced by
$\ell_T$, in which $p$ enters through $\ell_{p\wedge\ell_T}$; at the box triple $p=p_\square$.

\emph{Case $g=0$.} Here $\theta_{\mathfrak{L},a'}=\Phi\,\mathbf{r}_{a',\ell}$, and the ratios of the
direct cases apply, by the inequality form of the rule for the composite radii.

\emph{Case $g\ge1$ and $\ell_T=\mathfrak{c}(y)$.} Here
$\mathbf{r}^{\mathfrak{c},(p)}_{a'}=\mathbf{r}^{T,(p)}_{a'}$, $\theta^{(l)}(\mathfrak{c})\le\frac12$ by
\eqref{4.28:eq-lifted-thresholds-a}, and $s_T\ge0.8$. The room is therefore at least
$s_T-\frac12\ge0.3$.

\emph{Case $g\ge1$ and $\ell_T<\mathfrak{c}(y)$.} The adjacent-ratio bound, set against the powers of
$L$, together with $\ell_{p\wedge\ell_T}\le\ell_{p\wedge\mathfrak{c}}$ and the condition
$(1+\beta_0)/L\le2/21$, gives
\begin{equation*}
\frac{\mathbf{r}^{\mathfrak{c},(p)}_{a'}}{\mathbf{r}^{T,(p)}_{a'}}
\le\Bigl(\frac{1+\beta_0}{L^2}\Bigr)^{\mathfrak{c}-\ell_T}
\le\frac{1+\beta_0}{L}\le\frac{2}{21}\le0.0953 .
\end{equation*}
Hence
\begin{equation*}
\theta_{\mathfrak{L},a'}\le\tfrac12\cdot0.0953\,\mathbf{r}^{T,(p)}_{a'}\le0.05\,\mathbf{r}^{T,(p)}_{a'},
\end{equation*}
and the room is at least $0.8-0.05=0.75$.

The values of these bounds depend only on $\beta$, $\beta_0$ and $L$; no age, and neither $N$ nor
$N_0$, enters them.
Since $\beta/2=0.1$ does not exceed any of the lower bounds obtained, $\operatorname{room}^{\dagger}_T$
is at most every $\operatorname{room}^{\mathrm{on}}_a$. Finally,
$12/\operatorname{room}^{\dagger}_T=24/\beta$. This is already the first member of the maximum in
\cite[(3.54)]{B-CMP109}, so that maximum is unchanged.
\end{proof}

\begin{proposition}[The box construction at near pairs meeting lifted points]
\label{4.28:prop-near-pair-box}
Let $l$ be a stage, let $(\mathbb{U},\mathbb{J})\in\mathfrak{D}^{\flat}_l(Y)$, and let $(T,\square)$
be a near pair with the following properties:
\begin{itemize}
\item $\square\in\mathcal{C}^1$ has scale $i$;
\item $T\in\mathcal{T}_1\cap\mathcal{T}^{\mathbb{E}\mathbb{R}}$ and $X_T\subset\tilde{\square}^2$;
\item $[\square_0]_{\sharp}\subset Y$, where $\square_0=\tilde{\square}^5$ is the box of
Definition~\ref{4.28:def-refined-cover};
\item $\square_0$ meets a lifted point of $Z^{\mathrm{on}}$, that is, a point of
$Z^{\mathrm{on}}\cap\mathsf{L}_{\mathfrak{d}}$ in the sense of \eqref{4.28:eq-cutting-level-b}.
\end{itemize}
Assume the hypotheses of Lemma~\ref{4.28:lem-box-lemma} on the constants, assume $M\ge3L$, and
assume the following statements.
\begin{enumerate}
\item[(h1)] The near-pair box construction away from lifted points, together with the statements of
Ba{\l}aban that it invokes by statement, each taken at the scale of a cube of $\mathcal{C}^1$.
\item[(h2)] The ordering of the two radius constants $C_0,C_1$:
\begin{equation}\label{4.28:eq-near-pair-box-a}
\frac{C_1}{C_0}\ \ge\ \frac{2(1+2\beta)(1+\beta_0')^4\,B_3^2\,O(1)\,M\,\mathsf{a}}{1-\beta},
\end{equation}
where $B_3=B_3(d,L)$ is the constant of \cite{B-CMP109}, $O(1)$ is the numerical constant of the
ratio $\varpi_{\square}$ of the near-pair box construction (h1), and $\mathsf{a}=\mathsf{a}(L)$ is the
constant of clause (c) of Lemma~\ref{4.28:lem-box-lemma}. In the fixed order of choice of the
auxiliary parameters, the ratio $C_1/C_0$ is chosen after $L$ and $M$ (hence after $B_3$ and
$\mathsf{a}$, which the right side involves) and before $\gamma$. So read,
\eqref{4.28:eq-near-pair-box-a} is one-sided: it bounds $C_0/C_1$ from above only and places no
condition on $L$ or on $M$; it is not a cap in the sense of Definition~\ref{2.3:def-cap}.
\item[(h3)] The condition of \cite{B-CMP109} on the constants of the chain of operators implies
its localised form.
\item[(h4)] The first-order field-size bound at the pair $(\mathbb{U}^{\mathfrak{s}}_{\square},\mathbb{J})$;
the absolute field-size bound at the level of $\mathfrak{B}$, valid at every point and for every
admissible block structure, the collar included; and the multiplicity bound on the partition of
unity, which gives $a_{\square}\le 5$ when $M\ge 3L$, $\zeta_{\square},\zeta'_{\square}$ are functions
of the cube's own scale $i$ and the norm is read in the unit of a level $\mathfrak{c}\le i+1$.
\item[(h5)] The nested-averaging identity \cite[(3.39)]{B-CMP109}, used by statement.
\item[(h6)] The shell box clause of the companion space
$\mathfrak{N}^{(l)}$ on $\mathfrak{T}$.
\item[(h7)] The list $(\Sigma^{\ast})$ of retained triples $(p,X',\mathfrak{s})$ at which the
attempt space $\mathfrak{L}_l$ imposes its composite clauses. For lifted data of kinds (B) and (B$'$),
its structures include the retained stage structures of the present run at stages at most $l-1$,
among them the structure $\mathbb{B}^{(\kappa_0-\mathsf{h})}_{\kappa}$ of stage
$\kappa_0-\mathsf{h}\le l-1$ restricted to $\square_0$. For
an older piece $\mathfrak{a}$ of kind (A), the list is its list $(\Sigma^{\ast})^{\mathfrak{a}}$.
\item[(h8)] The conditions on the constants of the near-pair analysis, with its near-pair constant
taken equal to $66B_{\sharp}K_{\mathrm{cv}}$, under which, for a clause
$(y,a)\in\operatorname{Cl}(T)$ and $m:=\ell_{\mathfrak{B}}(y)$, the dilated and the undilated
members of family (c) move each entry by at most $\tfrac13\mathsf{w}_0\rho_m\mathbf{r}^T_a$ and
$(\mathsf{w}_0/18)\rho_m\mathbf{r}^T_a$ respectively, where $\mathsf{w}_0=1/24$ is the numerical
constant of that analysis, together by at most
$\iota^1_l(y)\mathbf{r}^T_a$, and the layer cost satisfies $\iota^1_l(y)\le\beta/36$.
\item[(h9)] The treatment of the clauses on $\mathfrak{T}$ in the near-pair box construction away
from lifted points.
\item[(h10)] The remainder estimates for the two configurations of clause (c$'$) of the near-pair
box construction.
\end{enumerate}
Then:
\begin{equation}\label{4.28:eq-near-pair-box-b}
\begin{aligned}
&\textup{(H)}\ \ \text{the input hypotheses (i)--(iv) of \cite[Lemma~4]{B-CMP109} hold on }
\square_0,\\
&\qquad\text{in the unit of the }\mathfrak{c}\text{-shell, literally,}\\
&\qquad\text{at factor }\le 1\text{ on }Z^{\mathrm{on}}\text{ and at factor }<1+2\beta
\text{ on }\mathfrak{T};\\
&\textup{(M)}\ \ Q_a(P)(y)<\theta_{T,a}(y)\quad\text{for every }(y,a)\in\operatorname{Cl}(T),
\end{aligned}
\end{equation}
where in (M) the configuration $P$ ranges over every member $P(\varsigma;t,\tau;B')$ of the family
(c) of the near-pair box construction with $|B'|<\alpha_3(T)$, and over both configurations of its
clause (c$'$). Here $\alpha_3(T)$ is formed with $\operatorname{room}^{\dagger}_T$ of
\eqref{4.28:eq-box-lemma-d} in place of $\operatorname{room}_T$.
\end{proposition}

\begin{proof}
We follow the steps of the proof of the near-pair box construction away from lifted points
(hypothesis (h1)). In doing so we use no statement other than (h1)--(h10),
Lemma~\ref{4.28:lem-box-lemma} and the statements and displays cited below.

The box lemma at lifted points, Lemma~\ref{4.28:lem-box-lemma}, is applied to $\square$, which has
scale $i$ and whose box $\square_0$ meets $\mathsf{L}_{\mathfrak{d}}$. Its clauses (c$_0$), (c),
(a), (o), (a$''$), (a$'$), (e) and (r) are those displayed in \eqref{4.28:eq-box-lemma-a},
\eqref{4.28:eq-box-lemma-b}, \eqref{4.28:eq-box-lemma-c} and \eqref{4.28:eq-box-lemma-d}. Throughout,
$g=\ell-\mathfrak{c}$ is the depth of \eqref{4.28:eq-cutting-level-b}, and
$\mathfrak{s}=\mathfrak{B}^{\mathfrak{c}}\restriction\square_0$, with top level $p_{\square}$
(Definition~\ref{4.28:def-refined-cover}).

\emph{The hypotheses} (H). Hypothesis (i) of \cite[Lemma~4]{B-CMP109} is obtained in two parts.
The first is the entry-$5$ clauses of $\mathfrak{L}_l$, converted into the unit of the
$\mathfrak{c}$-shell by (c$_0$). The second is the cube condition, which clause (e) supplies.
Hypothesis (ii) is obtained from entries $6$ and $7$, and hypothesis (iii) from entries $1$ and $3$.
In both cases the clauses of $\mathfrak{L}_l$ are converted by (c$_0$). On $\square_0\cap\mathfrak{T}$
the clauses of $\mathfrak{N}^{(l)}$ are used in their place. This gives factor at most $1$ on
$Z^{\mathrm{on}}$ and less than $1+2\beta$ on $\mathfrak{T}$.

Hypothesis (iv) is required on the cube written $\tilde{\square}^3$ in \cite[\S3]{B-CMP109}; in the
present notation this is the set $\square_0^{\sim 2}$, obtained from $\square_0$ by adding two layers
of cubes. Its points of $Z^{\mathrm{on}}$ lie in $\square_0^{\flat}$, the set on which
$\mathfrak{L}_l$ imposes the
clauses of the triples $(p,\square_0,\mathfrak{s})$. Let $p\le p_{\square}$. The triple
$(p,\square_0,\mathfrak{s})$ belongs to $(\Sigma^{\ast})$. Indeed, the structure $\mathfrak{s}$ is the
one attached to $\square$ by Definition~\ref{4.28:def-refined-cover}; the box $\square_0$ belongs to
Ba{\l}aban's class $\mathbf{D}_i$ of level-$i$ domains;
and, according to the kind of the datum:
\begin{itemize}
\item for kinds (B) and (B$'$), on a lifted datum the strata of $\mathfrak{c}$ are those of
$\mathbb{B}^{(\kappa_0-\mathsf{h})}_{\kappa}$ (Lemma~\ref{4.28:lem-cutting-strata}), so that
$\mathfrak{s}=\mathbb{B}^{(\kappa_0-\mathsf{h})}_{\kappa}\restriction\square_0$ is
a stage structure of the present run at stage $n'=\kappa_0-\mathsf{h}\le l-1$, retained by (h7);
\item for kind (A), the triple belongs to $(\Sigma^{\ast})^{\mathfrak{a}}$ by (h7).
\end{itemize}
The attempt space $\mathfrak{L}_l$ imposes its composite clauses, at birth radii, at every triple
of $(\Sigma^{\ast})$. Hence, for every $p\le p_{\square}$ and every
$y'\in\square_0^{\sim 2}\cap Z^{\mathrm{on}}$,
\begin{equation}\label{4.28:eq-near-pair-box-1}
Q_e^{(p,\square_0,\mathfrak{s})}(y')\ <\
\begin{cases}
\theta^{(l)}(\mathfrak{c}(y'))\,\mathbf{r}^{\mathfrak{c},(p)}_e(y')
\le\tfrac12\,\mathbf{r}^{\mathfrak{c},(p)}_e(y') & \text{if } g(y')\ge1,\\[3pt]
\Phi(y')\,\mathbf{r}^{\mathfrak{c},(p)}_e(y')\le\mathbf{r}^{\mathfrak{c},(p)}_e(y')
& \text{if } g(y')=0.
\end{cases}
\end{equation}
The bound by $\tfrac12$ is \eqref{4.28:eq-lifted-thresholds-a}, and the bound by $1$ at $g=0$ holds
since $\Phi\le s^{(k)}_{\ell}\le1$ (Lemma~\ref{4.28:lem-lifted-thresholds}). The same holds for the
clause on the configuration $U$ which $\mathfrak{L}_l$ imposes at the same triples, by the definition
of $(\Sigma^{\ast})$.

For kinds (B) and (B$'$), \eqref{4.28:eq-near-pair-box-1} is the condition \cite[(1.16)]{B-CMP122a}
written literally in the unit of the $\mathfrak{c}$-shell, at factor at most $\tfrac12$ on
$\mathsf{L}_{\mathfrak{d}}$; this is what hypothesis (iv) requires, and no conversion of units is
needed. For kind (A) with $g\ge1$ the
condition follows instead by conversion: the frozen weighted unit in which $\mathfrak{L}_l$ writes
the clause is at most $(0.381/L)^g\mathbf{r}^{\mathfrak{c},(p)}_e$.

On $\square_0^{\sim 2}\cap\mathfrak{T}$ the same condition is supplied by the shell box clause (h6),
at factor less than $1+2\beta$. By clause (o) of Lemma~\ref{4.28:lem-box-lemma} this occurs only
for kind (B) at $l=N$, on the un-lifted
collar, where the unit of the shell is the birth unit. This proves (H).

\emph{The orbit structure.} The space $\mathfrak{U}_T$ carries the $G^c$-orbit structure of
\cite[(3.37), (3.50)]{B-CMP109}, and this part of the argument is taken over from (h1) without change.
The entries $5$ and $8$ of a member $P$ lie within every positive threshold.

\emph{Entries $6$ and $7$.} There are three contributions.

$(\alpha)$ The ratio $\varpi_{\square}$ of the near-pair box construction satisfies
\begin{align*}
\varpi_{\square}
&=B_3^2\,O(1)\,M\,\frac{\mathfrak{a}^{(i)}_0}{\alpha_{1,i}}\cdot\frac{r^{[i]}_a}{\mathbf{r}^T_a}\\
&\le B_3^2\,O(1)\,M(1+2\beta)(1+\beta_0)L^2\,\frac{C_0}{C_1}
\Big(\frac{1+\beta_0}{L}\Big)^{i-\ell_T}\ \le\ \tfrac12(1-\beta).
\end{align*}
Here $r^{[i]}_a$ denotes the radius at which entry $a$ is read at the scale $i$. The first inequality
uses clause (c) at $p\le2$, which is the range of the converted entries $6$ and $7$. It also uses
(a$''$), which gives $\ell_T\le i$. The second inequality is the ordering
\eqref{4.28:eq-near-pair-box-a}, used as an upper bound on $C_0/C_1$.

$(\beta)$ \emph{The stage layer, in the units of the $\mathfrak{c}$-shell.} Let $y'\in\square_0$ and put
$n':=\ell_{\mathfrak{B}}(y')\ge\mathfrak{c}(y')$ and $g':=n'-\mathfrak{c}(y')$. A field of
first-order size $\mathsf{N}$ at level $n'$ has size at most $L^{-g'}\mathsf{N}$ at level
$\mathfrak{c}(y')$. The functions $\zeta_{\square},\zeta'_{\square}$ live on the scale of $\square$,
and $\mathfrak{c}\le i+1$ on $\square_0$. Since $M\ge3L$, the multiplicity bound of (h4) gives
$a_{\square}\le5$. Consequently
\begin{align*}
N^{\mathfrak{c}}_{y'}(\zeta_{\square}\mathbb{H}^{w})
&\le 2\cdot5\cdot L^{-g'}B_{\sharp}\cdot\big(\text{first-order size at }y'\big)
\le 11B_{\sharp}\,\vartheta(\square)\,\rho_{n'}L^{-g'}\alpha_{\cdot,n'}\\
&\le 11B_{\sharp}\,\vartheta(\square)\,\rho_{\ell_{\mathfrak{B}}(y')}\,\alpha_{\cdot,\mathfrak{c}(y')},
\end{align*}
where the first-order size at $y'$ is the one that (h4) bounds. The second inequality combines two
parts of (h4): the first-order size bound at $(\mathbb{U}^{\mathfrak{s}}_{\square},\mathbb{J})$, and the
absolute
bound at the level of $\mathfrak{B}$. The latter holds at every point and for every admissible
$\mathfrak{B}$, the collar included. The distance $d^{-}$ to the integration regions that enters these
bounds, taken at $\square$, is at most the distance $d^{\wedge}$ of the support $\square^{\zeta}$ of
$\zeta_{\square}$ from those regions.

Now fix the clause $(y,a)\in\operatorname{Cl}(T)$. By (a$'$),
$\rho_{\ell_{\mathfrak{B}}(\cdot)}\le2\rho_{\ell_{\mathfrak{B}}(y)}$ on $\square_0$. Put
$m:=\ell_{\mathfrak{B}}(y)$. With the near-pair constant equal to $66B_{\sharp}K_{\mathrm{cv}}$ as in
(h8), the dilated
member moves each entry by at most $\tfrac13\mathsf{w}_0\rho_m\mathbf{r}^T_a$, and the undilated
member by at most $(\mathsf{w}_0/18)\rho_m\mathbf{r}^T_a$. Together they move it by at most
$\iota^1_l(y)\mathbf{r}^T_a\le(\beta/36)\mathbf{r}^T_a$, by (h8).

$(\gamma)$ The contribution of $B'$ is at most $\operatorname{room}^{\dagger}_T/12$.

With $\beta=1/5$ the arithmetic check of the construction away from lifted points is unchanged:
\begin{equation*}
0.4+0.006+0.1\ <\ 0.8\ \le\ s_T.
\end{equation*}
This is the required inequality
for entries $6$ and $7$. The inequality that carries the argument on $Z^{\mathrm{on}}$
is \eqref{4.28:eq-near-pair-box-c} below.

\emph{Entries $1$ to $4$ at a clause on $Z^{\mathrm{on}}$: the start.} Apply the identity
\cite[(3.39)]{B-CMP109} of (h5) at the zero layer, with $\varsigma=1$, at the structure $\mathfrak{s}$,
whose levels on $\square_0$ are the one or two adjacent values of $\mathfrak{c}$ there, by (a).
The nested constraints hold because $\ell_T$ is at most the
levels of $\mathfrak{s}$, by (a$''$). The identity gives
\begin{equation*}
U_{\ell_T}\big(\square_0,M(\mathbb{U}^{\mathfrak{s}}_{\square})\big)
=(\mathbb{U}^{\mathfrak{s}}_{\square})^{(\mathrm{vu})},
\end{equation*}
where $(\mathbb{U}^{\mathfrak{s}}_{\square})^{(\mathrm{vu})}$ denotes the right side of
\cite[(3.39)]{B-CMP109} evaluated at $\mathbb{U}^{\mathfrak{s}}_{\square}$.
Hence the entry $a'$ of the zero-layer pair at $y$ is $G_{\square}$ times the entry of the input,
and for the latter
\begin{equation*}
Q_{a'}^{(p_{\square},\square_0,\mathfrak{s})}(y)\ <\ \theta_{\mathfrak{L},a'}(y)
=\theta^{(l)}(\mathfrak{c}(y))\,\mathbf{r}^{\mathfrak{c},(p_{\square})}_{a'}
\le\tfrac12\,\mathbf{r}^{\mathfrak{c},(p_{\square})}_{a'}
\le\tfrac12\,\mathbf{r}^{T,(p_{\square})}_{a'}
\end{equation*}
if $g\ge1$. The last step uses $\ell_T\le\mathfrak{c}$ (by (a$''$)) and the unit rule for the radii.
If $g=0$ the bound is $Q_{a'}^{(p_{\square},\square_0,\mathfrak{s})}(y)<\Phi\,\mathbf{r}_{a',\ell}$.
By the definition of $\operatorname{room}^{\mathrm{on}}_{a'}$ in (r), in both cases the threshold
equals $[s_T-\operatorname{room}^{\mathrm{on}}_{a'}(y)]\mathbf{r}^T_{a'}$. Clause (r) gives the
lower bounds
\begin{equation*}
\operatorname{room}^{\mathrm{on}}_{a'}(y)\ \ge\
\begin{cases}
0.3 & \text{if } g\ge1,\ \ell_T=\mathfrak{c}(y),\\
0.75 & \text{if } g\ge1,\ \ell_T<\mathfrak{c}(y),\\
\beta/2 & \text{if } g=0,\ \ell_T=\ell(y),\\
0.7 & \text{if } g=0,\ \ell_T<\ell(y).
\end{cases}
\end{equation*}
Clause (r) applies because $y\in X_T\subset\tilde{\square}^3$ and $p_{\square}\ge i\ge\ell_T$, so
that by the unit rule the unit is $\mathbf{r}^T_{a'}$.

At this point the argument of \cite{B-CMP109} uses \cite[(3.40)]{B-CMP109}, and on $\mathfrak{T}$ the
construction away from lifted points uses the box clause. On $Z^{\mathrm{on}}$, by contrast, the start
is the clause of $\mathfrak{L}_l$ itself. No diminished threshold and no reserve of one third enter,
because the subtraction in the thresholds of $\mathfrak{L}_l$ steps down with $l$.

\emph{The chain.} The operators and constants $G_{\square},\mathfrak{q}_1,\mathfrak{q}_2,\mathfrak{q}_1'$
of the chain depend only on $\mathfrak{a}^{(i)}$, on $\alpha_{\cdot,i}$ and on constants. Clause (c)
enters here at $p=3$. By (c), the condition of (h3) is read at
$\mathsf{a}\ge(1+2\beta)(1+\beta_0)L^3$, and (h3) then gives the localised condition without change
and without any term involving the weight ceiling $w_{\ast}$. The chain therefore costs at most
$\tfrac16\beta\hat{c}$. The condition involves $L$ only through an absolute power of $L$, which enters
the ceiling on $\gamma$ along a single trajectory $\mathsf{t}$; it is not a cap. The parameter
$\varsigma$ costs nothing. The stage layer adds $\iota^1_l(y)$, which is at most $\beta/36$ by the
estimate for entries $6$ and $7$. The parameter $B'$ adds $\tfrac16\operatorname{room}^{\dagger}_T$.

\emph{Total.} Adding these costs,
\begin{equation}\label{4.28:eq-near-pair-box-2}
\begin{aligned}
Q_a(P)(y)
&<\Big[s_T-\operatorname{room}^{\mathrm{on}}_a(y)+\tfrac16\beta\hat{c}
+\tfrac16\operatorname{room}^{\dagger}_T+\tfrac{\beta}{36}\Big]\mathbf{r}^T_a(y)\\
&<s_T\,\mathbf{r}^T_a(y)=\theta_{T,a}(y).
\end{aligned}
\end{equation}
The second inequality holds because $\hat{c}=\tfrac12$ and
$\operatorname{room}^{\dagger}_T=\beta/2$, so that, with $\beta=1/5$, the three costs satisfy
\begin{equation}\label{4.28:eq-near-pair-box-c}
\frac{\beta}{12}+\frac{\beta}{12}+\frac{\beta}{36}=\frac{7\beta}{36}=\frac{7}{180}
\le0.039<\frac{\beta}{2}\le\operatorname{room}^{\mathrm{on}}_a(y)
\end{equation}
in every case, and $7/180<0.3$ at the composite start at a lifted point. Entry $3$ is treated by the
parallel chain of \cite{B-CMP109}. Entries $2$ and $4$ at the triples of $T$ itself are treated by
\cite[(3.38)]{B-CMP109}, as in the construction away from lifted points.

\emph{The two configurations of clause} (c$'$). The scaled distance satisfies
\begin{equation*}
D:=d_{\ell_T}\big(X_T,\operatorname{supp}(1-\zeta'_{\square})\big)\ \ge\ w\,L^{(i-\ell_T)^+},
\end{equation*}
where $w$ is the cost per crossed $M$-cube entering the scaled distance $d_{\ell_T}$,
because the two sets are separated by at least one $\pi_i$-cube, and the distance $d_{\ell_T}$ is
measured in units at most $\ell_{\ell_T}$.
By the remainder estimates (h10) for these configurations, both
configurations move each entry by less than $\iota^1_l(y)\mathbf{r}^T_a$. Hence
\eqref{4.28:eq-near-pair-box-2} applies to them. The Cauchy bounds and the sum over $T$ depend only on
$i$ and $\ell_T$ and are taken over without change.

\emph{Clauses on $\mathfrak{T}$.} By clause (o) of Lemma~\ref{4.28:lem-box-lemma}, these occur only
for kind (B) at $l=N$, at points
$y\in\bar{c}\cap\Omega_k\setminus\Lambda_k$, which are not lifted. We follow the treatment of these
clauses in the construction away from lifted points (hypothesis (h9)), with
Lemma~\ref{4.28:lem-box-lemma} in place of
the box lemma used there. That treatment reads clauses (a)--(e) and the box clause at the triple
$(p_{\square},\square_0,\mathfrak{s})$. This is a clause imposed at stage $l$, because $\mathfrak{s}$
is a structure of $(\Sigma^{\ast})$ and every pair $(p,\mathfrak{s})$ of $(\Sigma^{\ast})$ carries
the clause. Clauses (b) and (d) of the box lemma used there remain true and are not needed.

\emph{Statements used by statement.} The statements of Ba{\l}aban invoked by statement are those of
the list in (h1), taken at the scale $i$, with the quantity $\eta$ of \cite{B-CMP119} replaced
by $L^{-n}$ at level $n$.
\end{proof}

\begin{proposition}[Kernel membership and anchored sums on the refined cover]
\label{4.28:prop-kernel-membership}
Let $d=4$, $L\ge 12$, $R_{\min}\ge 5L$, $\beta_0=1/7$, $\beta=1/5$, and let $\mathsf{w}_0=1/24$ be the
fixed fraction in the upper bound on the layer costs in (IX) below. Let $L_0$ denote
the weight base of \cite{B-CMP122a}, with $2\le L_0$ and $L_0^2\le L/3$. Fix the site $k$ and a stage $l$
of the preparation, and put $\mathfrak{B}:=\mathbb{B}^{(l-1)}_k$, with level function $\ell_{\mathfrak{B}}$.
Let $\mathfrak{B}^{\mathfrak{c}}$ and $\mathcal{C}^1$ be as in Definition~\ref{4.28:def-refined-cover}.
For a cube $\square\in\mathcal{C}^1$, let $\square_0=\tilde{\square}^5$, $\mathfrak{s}=\mathfrak{B}^{\mathfrak{c}}\restriction\square_0$
and the box configuration $\mathbb{U}^{w}_{\square}$ be as there as well.

Let $\mathfrak{d}=(C;\kappa,\mathsf{h},\kappa_0,J)$ be a lifted datum of kind (A), (B) or (B$'$). Its
\emph{density} $k^{\circ}$ is $k^{\circ}=k_a$ in kind (A) and $k^{\circ}=k$ in kinds (B), (B$'$).
In each kind $k^{\circ}=\kappa$, the site of the datum, and $\ell\le\kappa$ on $C$.
Let $Y$ be a region, and let $(\mathbb{U},\mathbb{J})\in\mathfrak{D}^{\flat}_l(Y)$ be a configuration of the
stage-$l$ domain on $Y$. Let $\square\in\mathcal{C}^1$ have scale $i$, with
$[\square_0]_\sharp\subset Y$ and $\square_0$ meeting a lifted point. Terms of the expansion are denoted
$T$; for a term $T\in\mathcal{T}_1\cap\mathcal{T}^{\mathbb{E}\mathbb{R}}$, $X_T$ is its region,
$\ell_T$ its level, and its clauses $a\in\operatorname{Cl}(T)$ carry the radii $\mathbf{r}^T_a$; such a
$T$ is \emph{near} to $\square$ if $X_T\subset\tilde{\square}^2$.
Let $\mathcal{Q}\subset R^{(l)}$ be the family of anchor cubes. On the run, that is, for a datum
of kind (B) on a component $\bar{c}$ of $Z^{[l]}$, write $J=j+1$ and $k_0:=\kappa_0$, assume $j>k_0$
(for $j\le k_0$ one has $\mathcal{C}^1=\mathcal{C}$ on the run, and the sums below over
$\mathcal{C}^1\setminus\mathcal{C}$ are empty), and note that there $\mathcal{Q}\subset\pi_{j+1}$.
Write $\mathcal{M}^n$ for the $n$-fold composite of
Ba{\l}aban's block-averaging map, so that
\begin{equation*}
\mathcal{M}^n=\mathcal{M}^{n-\mathfrak{c}}\circ \mathcal{M}^{\mathfrak{c}} .
\end{equation*}
Assume the following.
\begin{enumerate}
\item[(I)] \emph{Regularity of the box configuration.} This is the assertion of \cite{B-CMP109} that
the configuration constructed for the cube $\square_0=\tilde{\square}^5$, extended off $\square_0$,
``has approximately the same regularity properties as $\mathbf{U}$, hence all necessary theorems are
valid'', where $\mathbf{U}$ is the background.
\item[(II)] \emph{Criticality under refinement.} Let a configuration be critical for a block
structure whose constraints are functions of the constraints of a refinement. Then it is critical for
that refinement.
\item[(III)] \emph{Restriction of critical configurations.} Let $U$ be critical for
$\mathfrak{B}^{\mathfrak{c}}$. Then $U\restriction\square_0$ coincides, up to a gauge transformation,
with $U(\mathfrak{B}^{\mathfrak{c}}(\square_0),\mathcal{M}^{\cdot}(U))$. The latter is the configuration on
$\square_0$ determined by the restriction of $\mathfrak{B}^{\mathfrak{c}}$ to $\square_0$ and by the
block averages $\mathcal{M}^{n}(U)$.
\item[(IV)] \emph{The weight ceiling.} The weights $\omega$ of the datum and the weight ceiling
$w_{\ast}(k^{\circ})$, the maximum of the weights admitted at density $k^{\circ}$, satisfy
$\omega\le L_0^{2(J-\kappa_0)}\le w_{\ast}(k^{\circ})$. Moreover
$w_{\ast}(k^{\circ})\,\bar{\alpha}_{k^{\circ}}\le\gamma_0^{(k)}/5$ for every density $k^{\circ}\le k$.
Here $\gamma_0^{(k)}$ is the
radius of the kernel class $\mathfrak{K}_{109}(\mathfrak{B};\gamma_0^{(k)})$ at site $k$, and
$\bar{\alpha}_{k^{\circ}}\ge\alpha_{\cdot,n}$ for all $n\le k^{\circ}$.
\item[(V)] \emph{The smallness hypotheses of the stage theorem.} Consider the theorem on the stage-$l$
fluctuation terms $\mathbb{C}^{(l)}_k$. Its three hypotheses on $\gamma_0$ hold with $\gamma_0$ replaced by
$2B_3\bar{\gamma}_0(\mathsf{t})$, with suprema taken over one trajectory $\mathsf{t}$. Here
$B_3=B_3(d,L)$ is the constant of \cite{B-CMP109}, $\bar{\gamma}_0(\mathsf{t})$ is the kernel radius
$\gamma_0$ of those hypotheses associated with the trajectory $\mathsf{t}$, and $B_3$ is chosen before
$\gamma$.
\item[(VI)] \emph{The sum statement for the near terms.} This statement has four parts.
  \begin{itemize}
  \item It has a kept factor $e^{-\frac14\delta d_{\ell_T}(X_T,\square^{\zeta})}$ at the radius it fixes
  for the clauses of a term $T$.
  \item Its clause (iv) carries the series over the terms $T$ at rate $\delta w/8$, with constant
  $K''K_S(E_0+1)\psi^{\natural}$, where $\psi^{\natural}$ is the smallness parameter of that statement.
  \item Its anchor rule attaches a term $T$ to a cube $q$ of the family $\mathcal{Q}$ meeting
  $\tilde{X}_T^4$. The rule
  counts at most $16^d$ blocks with $\mathsf{y}=\emptyset$, inside the animal constant $K_\flat'$.
  \item By its cost bound, every cube of side $M$ in the units of $\mathfrak{B}$'s local scale crossed
  by a contour adds at least $w$ to $d_{\ell_T}$.
  \end{itemize}
\item[(VII)] \emph{The one-cube bound.} This is the bound on the total at one cube of the cover. Its
hypotheses are (M) and (H) of~\eqref{4.28:eq-near-pair-box-b} and the stage-layer bound established
in the proof of Proposition~\ref{4.28:prop-near-pair-box}: the dilated and the undilated member of the
stage layer together move each entry by at most $\iota^1_l(y)\mathbf{r}^T_a\le(\beta/36)\mathbf{r}^T_a$.
\item[(VIII)] \emph{The depth-gain lemma.} Its two parts are as follows.
  \begin{itemize}
  \item Part (i). Let $(p,X',\mathfrak{s})$ be a retained triple at a lifted point $y'$, $\mathfrak{s}$ now
  denoting its structure (for the box triple, $\mathfrak{s}=\mathfrak{B}^{\mathfrak{c}}\restriction\square_0$),
  considered in the frame
  $(X',\mathfrak{B},\mathfrak{B}^{\mathfrak{c}}\restriction X')$. In this frame the triple belongs to the
  list $(\Sigma^{\flat})$ of triples at which the lemma's composite clauses are imposed, and the
  birth level is $m_{\mathrm{ref}}=\mathsf{m}^{\circ}(y')=\mathfrak{c}(y')$, where $\mathsf{m}^{\circ}$ is
  the birth-level function of the frame; the current level is
  $m:=\mathsf{m}=\ell(y')\ge\mathfrak{c}(y')$, and $l_x=\min\{p,\mathfrak{c}(y')\}$. Then the stage-$l$
  layer moves $Q_e^{(p,X',\mathfrak{s})}(y')$, $e\in\{2,4\}$, by at most
  $\iota_l(y')(4.4/L^2)^{\Delta}$ times the birth radius $r_e(\mathfrak{c},p\wedge\mathfrak{c})$, with
  $\Delta=g(y')$. The additive term of the move is $\pi r_e(m,l_x)$; here $\pi$ and
  $\pi_\times\le\pi$, without subscript, denote the two small constants of the bump estimate of the
  stage layer, not a cube family.
  \item Part (ii). In kind (A) the step is stated in the frozen weighted unit at the frozen
  level $m$, which we denote $\mathsf{u}_m$; the one-step room at that level is
  $\nu^{(l+1)}_m\cdot\mathsf{u}_m$, the stage layer moves the
  entry by at most $\iota_{l+1}\cdot\mathsf{u}_m$, and
  $\iota_{l+1}\le\frac12\nu^{(l+1)}_m$, with $\nu^{(l+1)}_m$ as in (IX).
  \end{itemize}
\item[(IX)] \emph{Running factors and layer costs.} Put $\nu_x:=\beta2^{-|x-c_{\ast}|}$, where
$c_\ast:=\ell_{\cdot}+1$ and $\ell_{\cdot}$ is the level entering the running factor at the step in
question; the index of that step is $l$ unless indicated, and we write $\nu^{(l+1)}_x$ when it is
$l+1$. At a
retained structure the terms $\sigma$ of the
running factors do not depend on $l$, and the difference of consecutive running factors is
\begin{equation*}
\bigl[\theta^{(l)}_{\mathfrak{s}}-\theta^{(l+1)}_{\mathfrak{s}}\bigr](y')=\nu_{\mathfrak{c}(y')} .
\end{equation*}
The layer costs satisfy $2.4\pi\le\iota_l(y')\le\frac16\mathsf{w}_0\nu_m$.
\item[(X)] \emph{Decay per crossed cube.} $\delta w\ge 16\ln L+4$.
\item[(XI)] \emph{The statements for the unrefined cover.} These are the statements that the
estimates used for the cover $\mathcal{C}$ make at a cube of $\mathcal{C}$. Among them is the tail
estimate. It bounds the terms of a class fixed there, which we call the $M$-class: the tail terms,
whose regions lie in no $\tilde{\square}^2$ of the cover at hand. Its hypothesis list
$(H^{\natural})$ has a clause (a), which one step of its proof delivers from its input bound; here
that input bound is used restricted to $Z^{\mathrm{on}}$.
\item[(XII)] \emph{The distance and counting bounds on the lifted set.} These are the
bounds~\eqref{4.28:eq-old-region-distance-a} on the distance from an old region's lifted set to the
integration regions and the counts~\eqref{4.28:eq-refined-count-a} of the refined cubes along the
current preparation.
\end{enumerate}
Then the following hold.
\begin{enumerate}
\item[(a)] The configuration $U:=U^{(l)}_k$ of stage $l$ at site $k$, frozen on the arrested pieces, is
critical for
$\mathfrak{B}^{\mathfrak{c}}$. Moreover
$U\restriction\square_0=U(\mathfrak{B}^{\mathfrak{c}}(\square_0),\mathcal{M}^{\cdot}(U))$
up to a gauge transformation. The same holds for the background of a piece frozen at a block structure
of top level $J-1\ge\kappa_0$.
\item[(b)] $(\mathbb{U}^{w}_{\square},\mathbb{J})\in\mathfrak{K}_{109}(\mathfrak{B};\gamma_0^{(k)})$.
\item[(c)] Let $y\in Z^{\mathrm{on}}$, and take a clause $a\in\operatorname{Cl}(T)$ at $y$ at the radius
of the kept factor.
There the stage layer moves the entry by at most $\iota^1_l(y)\mathbf{r}^T_a$, and this is less than
$\operatorname{room}^{\mathrm{on}}_a(y)\mathbf{r}^T_a$.
\item[(d)] In kinds (A) and (B$'$) no anchor cube meets $\tilde{X}_T^4$. In these kinds the sum over the
lifted cubes of $\mathcal{C}^1$ is bounded through the second and third of the
bounds~\eqref{4.28:eq-old-region-distance-a} of (XII), with no factor $\mathsf{l}_j$; here
$\mathsf{l}_j$ is the logarithmic count of the present coupling, which bounds the multiplicity in the
second of the counts~\eqref{4.28:eq-refined-count-a} of (XII). On the run, fix
$q\in\mathcal{Q}$
and consider the near terms anchored at $q$ through cubes $\square\in\mathcal{C}^1\setminus\mathcal{C}$.
Their total is at most
\begin{equation}\label{4.28:eq-kernel-membership-1}
c_d'\,\mathsf{l}_j^2\,K''K_S(E_0+1)\psi^{\natural},
\end{equation}
where $C_d:=d2^d$ and
\begin{equation*}
c_d':=C_d\sum_{r\ge0}(r+3)^{d-1}e^{-\frac18\delta w(r-3)^+}\le 435\,C_d .
\end{equation*}
\item[(e)] At fixed $T$,
\begin{equation}\label{4.28:eq-kernel-membership-2}
\sum_{\square}e^{-\frac14\delta d_{\ell_T}(X_T,\square^{\zeta})}\le 2C_d16^d ,
\end{equation}
where the sum runs over the cubes $\square\in\mathcal{C}^1$ that meet $X_T$ through one coarse cube
only, that is, whose element attached to $X_T$ has empty domain, $\mathsf{y}=\emptyset$.
\item[(f)] Let $X_T$ lie in no $\tilde{\square}^2$ with $\square\in\mathcal{C}^1$. Then $T$ is a tail
term for the cover $\mathcal{C}^1$, hence of $M$-class, and the step of the tail estimate of (XI) that
uses its input bound, with that bound restricted to $Z^{\mathrm{on}}$, delivers clause (a) of
$(H^{\natural})$ for $T$.
\item[(g)] $|\mathsf{b}(e)|\le 48^d+|\mathsf{y}|\le(12L)^d+|\mathsf{y}|$, where $\mathsf{b}(e)$ is the set
of cubes attached to an element $e$ of the expansion (here $e$ denotes an element, not an entry index).
\item[(h)] The one-cube bound applies at $\square$. It gives a total of at most
\begin{equation*}
c_\alpha C_R^{\natural}(E_0+1)e^{C_2\kappa_1}\psi^{\natural}(\square)
\end{equation*}
times one plus the remainder terms of that bound, with $s(\square)=0$. Here $c_\alpha$, $C_R^{\natural}$
and $C_2$ are the constants of the one-cube bound (VII), $\psi^{\natural}(\square)$ is the smallness
parameter $\psi^{\natural}$ of (VI) evaluated at $\square$, and $s(\square)$ is the cube-dependent
parameter of that bound; at the cubes of $\mathcal{C}^1$ it is zero.
\item[(i)] Let $(p,X',\mathfrak{s})$ be a retained triple at a lifted point $y'$, in kind (B) or (B$'$),
$\mathfrak{s}$ now denoting its structure. This
includes the box triple $(p,\square_0,\mathfrak{s})$, with
$\mathfrak{s}=\mathfrak{B}^{\mathfrak{c}}\restriction\square_0$. Then, in its frame, the stage-$l$ layer moves
$Q_e^{(p,X',\mathfrak{s})}(y')$, $e\in\{2,4\}$, by at most
$\frac12\nu_{\mathfrak{c}(y')}\,r_e(\mathfrak{c},p\wedge\mathfrak{c})$. This is half of the one-step
room. In kind (A), where the step is stated in the frozen weighted unit $\mathsf{u}_m$ of (VIII) at the
frozen level $m$, the stage layer moves the entry by at most
$\iota_{l+1}\cdot\mathsf{u}_m\le\frac12\nu^{(l+1)}_m\cdot\mathsf{u}_m$, half of the one-step room
$\nu^{(l+1)}_m\cdot\mathsf{u}_m$.
\end{enumerate}
\end{proposition}

\begin{proof}
Each clause below modifies the corresponding statement of (XI) at a cube of $\mathcal{C}^1$. Where no
modification is stated, (XI) applies verbatim.

\emph{Clause (a).} The configuration $U=U^{(l)}_k$, frozen on the arrested pieces, is critical for its
own block
structure. Since $\mathcal{M}^n=\mathcal{M}^{n-\mathfrak{c}}\circ \mathcal{M}^{\mathfrak{c}}$, the
constraints of that structure are
functions of the constraints of the refinement $\mathfrak{B}^{\mathfrak{c}}$. By~(II), $U$ is critical for
$\mathfrak{B}^{\mathfrak{c}}$. By~(III),
$U\restriction\square_0=U(\mathfrak{B}^{\mathfrak{c}}(\square_0),\mathcal{M}^{\cdot}(U))$
up to a gauge transformation. The background of a frozen piece is critical for the block structure of
top level $J-1\ge\kappa_0$. The structure $\mathfrak{B}^{\mathfrak{c}}$ refines that structure too, so
(II) and (III) apply in the same way. All of this is taken by statement from~(II) and~(III).

\emph{Clause (b).} Membership is derived from the direct entries of the input. By
Lemma~\ref{4.28:lem-lifted-thresholds}, on $[\square_0]_\sharp\subset Y$ the input obeys, for every direct
entry $e$ and every $y'$,
\begin{equation}\label{4.28:eq-kernel-membership-3}
Q_e(y')<\Phi(y')\,\omega(y')\,\mathbf{r}_{e,\ell(y')}\le\omega(y')\,\mathbf{r}_{e,\ell_{\mathfrak{B}}(y')} .
\end{equation}
The second inequality holds because $\ell_{\mathfrak{B}}\in\{\ell-1,\ell\}$ and $(1+\beta_0)/L\le1$. So
the input is small at $\mathfrak{B}$'s own levels, with size $\omega(y')\alpha_{\cdot,\ell(y')}$. Moreover
\begin{equation*}
\omega(y')\,\alpha_{\cdot,\ell(y')}\le w_{\ast}(k^{\circ})\,\bar{\alpha}_{k^{\circ}}\le\gamma_0^{(k)}/5 .
\end{equation*}
Indeed, by (IV), $\omega\le L_0^{2(J-\kappa_0)}\le w_{\ast}(k^{\circ})$. Since $\ell(y')\le\kappa=k^{\circ}$,
we have $\alpha_{\cdot,\ell(y')}\le\bar{\alpha}_{k^{\circ}}$. Finally
$w_{\ast}(k^{\circ})\bar{\alpha}_{k^{\circ}}\le\gamma_0^{(k)}/5$ by (IV), since $k^{\circ}\le k$.
No comparison between different trajectories $\mathsf{t}$ enters.

By (I), taken by statement on the whole of $\square_0$, the configuration $\mathbb{U}^{w}_{\square}$,
extended by $\mathbb{U}$ off $\square_0$, has approximately the same regularity properties as the
background, so the theorems used for $\mathbf{U}$ apply to it. This gives
(b).

Where a number is wanted, \cite{B-CMP109} bounds the direct entries of the pair
$(\mathbb{U}^{w}_{\square},\mathbb{J})$ by $2B_3(d,L)$ times
those of the input. The factor $2B_3$ is carried inside the three hypotheses of (V), with $\gamma_0$
replaced by
$2B_3\bar{\gamma}_0(\mathsf{t})$. These are suprema over one trajectory. They carry no factor of the
degree function $p_0(\cdot)$, since $B_3$ is chosen before $\gamma$. They are thus a ceiling on $\gamma$,
not a cap in the sense of Definition~\ref{2.3:def-cap}.

\emph{Clause (c).} The sum statement (VI) depends on $\square$ only through
$\operatorname{supp}\zeta_\square$, through the distance $d^-_\square$ of $\square$ to the integration
regions, and through the sum over $T$ at the scale of the cube. For $\ell_T>i$ a contour
from $X_T$ crosses the $\mathfrak{c}$-shells $i+1,\dots,\ell_T-1$. By clause (v) of
Lemma~\ref{4.28:lem-cutting-strata},
each of these shells is at least one cube of side $M$ in the units of $\mathfrak{B}$'s local scale wide;
these crossings
enter the distance $d_{\ell_T}$ of the kept factor of (VI). At a clause of $T$ at
$y\in Z^{\mathrm{on}}$, at the radius of the kept factor, the stage layer moves the entry by at most
$\iota^1_l(y)\mathbf{r}^T_a$, by the stage-layer bound recalled in (VII). By~\eqref{4.28:eq-box-lemma-d},
this is less than
$\operatorname{room}^{\mathrm{on}}_a(y)\mathbf{r}^T_a$.

\emph{Clause (d): kinds (A) and (B$'$).} The anchor of (VI) is a cube $q\in\mathcal{Q}$ meeting
$\tilde{X}_T^4$. In kinds (A) and (B$'$) there is none, and the anchor rule of (VI) is void in these
kinds. The reason is that here $\mathcal{Q}\subset R^{(l)}$, while, by the first of the distance
bounds~\eqref{4.28:eq-old-region-distance-a} of (XII), on the distance from an old region's lifted set
to the integration regions, the cubes $\square$ of $\mathcal{C}^1$ meeting the lifted set lie at distance
$D(\square)\ge d^-_{\mathfrak{d}}$ from the integration regions, the lower bound on $d^-_{\mathfrak{d}}$
given there being a positive multiple of $w$ in each of the two kinds.

In these kinds, with $u=u(\square)$ the depth of $\square$ and $D=D(\square)$ as
in~\eqref{4.28:eq-old-region-distance-a} of (XII), the square root $(3L^{u+1})^de^{-\frac14\delta D}\le1$
of the second of the
bounds~\eqref{4.28:eq-old-region-distance-a} absorbs the multiplicity of the lifted cubes of
$\mathcal{C}^1$ under one cube of scale $J$. The other square root, $e^{-\frac14\delta D}$, survives as
decay at rate $\frac14\delta$, summable against
the count in the third of those bounds. No $\mathsf{l}_j$ enters.

\emph{Clause (d): the run, setting up the count.} By~\eqref{4.28:eq-refined-count-a} of (XII), the cubes
$\square\in\mathcal{C}^1\setminus\mathcal{C}$ have scales $\mathfrak{c}\in[k_0,j-1]$. The cubes of
$\mathcal{C}$, scale $J$ included, keep the anchor of (VI) as stated. We sum over $\square$ first, at a
fixed $q\in\mathcal{Q}\subset\pi_{j+1}$.

We claim that the cubes $\square\in\mathcal{C}^1\setminus\mathcal{C}$ at $\pi_{j+1}$-distance $r$ from $q$
number at most $C_d(r+3)^{d-1}\mathsf{l}_j^2$. We apply the second of the
counts~\eqref{4.28:eq-refined-count-a} of (XII)
to the $\pi_{j+1}$-cubes that such a $\square$ touches. Each of them lies in $\{\mathfrak{c}=k_0\}$ or in one
old shell; this is shown in the next paragraph.

\emph{Clause (d): location of the touched cubes.} The cube $\square$ lies within two $\pi_{\mathfrak{c}}$-sides
of the set
\begin{equation*}
\{\mathfrak{c}<\ell_{\mathfrak{B}}\}=(\Omega^c_j\setminus Z''_{k_0+1})\cap\bar{c} .
\end{equation*}
By \cite[(1.10)]{B-CMP122a}, $Z''_{k_0+1}=Z'^{\sim}_{k_0}$. This set is $Z'_{k_0}\supseteq Z''_{k_0}$ with
one whole layer of $\pi_k$-cubes added. Here~\eqref{4.28:eq-cutting-strata-a} and clause (iv) of
Lemma~\ref{4.28:lem-cutting-strata} give
\begin{equation*}
\Gamma''_{k_0}\supseteq Z'^{\sim}_{k_0}\setminus Z'_{k_0} ,
\end{equation*}
and clause (ii) of Lemma~\ref{4.28:lem-cutting-strata} and its proof give
\begin{equation*}
Z''_{k_0}\subset\Omega^c_{k_0}\subset Z'_{k_0} .
\end{equation*}
Since $\mathfrak{c}\le j-1\le k-2$, two $\pi_{\mathfrak{c}}$-sides are at most $2L^{-2}$ $\pi_k$-sides.
Consider the $\pi_k$-cube holding a $\pi_{j+1}$-cube that $\square$ touches. It lies at distance less than
one $\pi_k$-side from the complement of $Z''_{k_0+1}=Z'^{\sim}_{k_0}$, whereas every $\pi_k$-cube of
$Z'_{k_0}$ is separated from that complement by the whole added layer. It is therefore none of the
cubes of $Z'_{k_0}$. So it lies in $\Gamma''_{k_0}\subset\{\mathfrak{c}=k_0\}$, or outside $Z''_{k_0+1}$.
Outside $Z''_{k_0+1}$ the boundaries $\partial\Omega_m$, $m>k_0$, run along $\pi_k$-cubes,
by clause (v) of Lemma~\ref{4.28:lem-cutting-strata} and its proof. Hence the $\pi_{j+1}$-cube lies
either in $\{\mathfrak{c}=k_0\}$ or inside a single old shell, as claimed.

\emph{Clause (d): the count.} By the second of the counts~\eqref{4.28:eq-refined-count-a} of (XII), at most
$\mathsf{l}_j^2$ cubes of $\mathcal{C}^1$ of the families $\ge k_0$ touch each such $\pi_{j+1}$-cube. These
families include all of $\mathcal{C}^1\setminus\mathcal{C}$. The $\pi_{j+1}$-cubes at sup-distance $r$ from
$q$ number at most
\begin{equation*}
2d\cdot2^{d-1}(r+3)^{d-1}=C_d(r+3)^{d-1} .
\end{equation*}
This proves the claim.

\emph{Clause (d): the distance bound.} Let $T$ be a near term anchored at $q$ through the hull of its
element at such a $\square$. We show that
\begin{equation*}
d_{\ell_T}(X_T,\square^{\zeta})\ge w(r-3)^+ .
\end{equation*}
There are two cases.
\begin{itemize}
\item Suppose $q$ meets $[\tilde{\square}^4]_\sharp$. The cube $\square$ has scale at most $j-1$, so
$[\tilde{\square}^4]_\sharp$ reaches less than two $\pi_{j+1}$-sides beyond $\square$. Hence $r\le2$, and
$(r-3)^+=0$.
\item Otherwise $q$ meets $[\tilde{X}_T^4]_\sharp$. Then every contour from $X_T$ to $\square^{\zeta}$
crosses at least $(r-3)^+$ whole $\pi_{j+1}$-cubes. By the cost bound of (VI), each costs at least $w$
in $d_{\ell_T}$.
\end{itemize}

\emph{Clause (d): the total.} The series over $T$ in clause (iv) of (VI) runs at rate $\delta w/8$. So half
of the kept rate $\frac14\delta$ is free. Summing over $\square$ first and then over $T$, the total at $q$
is at most
\begin{equation*}
\sum_{r\ge0}C_d(r+3)^{d-1}\mathsf{l}_j^2\,e^{-\frac18\delta w(r-3)^+}\,K''K_S(E_0+1)\psi^{\natural}
=c_d'\,\mathsf{l}_j^2\,K''K_S(E_0+1)\psi^{\natural} .
\end{equation*}
By (X) and $L\ge12$,
\begin{equation*}
\tfrac18\delta w\ge 2\ln L+\tfrac12\ge 5.469 .
\end{equation*}
At $d=4$, the terms $r\le3$ contribute $\sum_{r\le3}(r+3)^3=432$. For $r=s+3$ with $s\ge1$ we have
$(s+6)^3\le343\cdot2^s$. Hence
\begin{equation*}
\frac{c_d'}{C_d}\le 432+\sum_{s\ge1}343\bigl(2e^{-5.469}\bigr)^s\le 432+3=435 .
\end{equation*}
This is the member of the multiplicity sum that carries $\mathsf{l}_j^2$. It carries exactly one power
$\mathsf{l}_j^2$ and no factor of the degree function $p_0(\cdot)$.

\emph{Clause (e).} At fixed $T$, with elements taken through their hulls as in (d), many cubes of
$\mathcal{C}^1$ meet $X_T$ through one coarse
cube only, that is, with $\mathsf{y}=\emptyset$. They are summed by the kept factor of (VI). By the cost
bound of (VI), a
$\pi_i$-cube costs at least $w$ in $d_{\ell_T}$. Grouping the cubes by the number $r$ of such crossings
gives
\begin{equation*}
\sum_{\square}e^{-\frac14\delta d_{\ell_T}(X_T,\square^{\zeta})}
\le\sum_{r\ge0}C_d(2r+16)^de^{-\frac14\delta wr} .
\end{equation*}
By (X), $e^{-\frac14\delta w}\le L^{-4}e^{-1}$. At $d=4$,
\begin{equation*}
(2r+16)^4=16^4(1+r/8)^4\le16^4e^{r/2}\le16^4\,2^r .
\end{equation*}
Moreover, at $L\ge12$,
\begin{equation*}
\sum_{r\ge0}(2e^{-1}L^{-4})^r\le1.0001 .
\end{equation*}
Together these give~\eqref{4.28:eq-kernel-membership-2}. This constant, depending on $d$ only, replaces
the count of at most $16^d$ blocks with $\mathsf{y}=\emptyset$ in the anchor rule of (VI). It sits inside
the same animal constant $K_\flat'$. With (d), $K_\flat'$ is thus multiplied by two constants depending on
$d$ only.

\emph{Clause (f).} Let $X_T$ lie in no $\tilde{\square}^2$ with $\square\in\mathcal{C}^1$. Such terms $T$
are tail terms for the cover $\mathcal{C}^1$, hence of $M$-class. The tail estimate of (XI) therefore
applies to them; the step of it that uses its input bound, with that bound restricted
to $Z^{\mathrm{on}}$, delivers clause (a) of $(H^{\natural})$ for them.

\emph{Clause (g).} The set $[\tilde{\square}^5]_\sharp$ has at most $16^d$ cubes. Each of them touches at
most $3^d$ cubes of $\mathcal{Q}\subset\pi_{j+1}$. Hence $|\mathsf{b}(e)|\le48^d+|\mathsf{y}|$. Since
$48\le12L$, this is at most $(12L)^d+|\mathsf{y}|$.

\emph{Clause (h).} The bound (VII) is applied by statement at a cube of $\mathcal{C}^1$. Its three
hypotheses are (M), the stage-layer bound and (H). All three are supplied by
Proposition~\ref{4.28:prop-near-pair-box} and its proof. The total is as stated, with $s(\square)=0$.

\emph{Clause (i): the frame and the gain.} This step is taken from the depth-gain lemma (VIII), by
statement: in kinds (B) and (B$'$) from its part (i), in kind (A) from its part (ii). In kinds (B) and
(B$'$) the lemma applies at the retained triples, regarded as triples of $(\Sigma^{\flat})$,
with birth level
$m_{\mathrm{ref}}=\mathsf{m}^{\circ}(y')=\mathfrak{c}(y')$, in the frame
$(X',\mathfrak{B},\mathfrak{B}^{\mathfrak{c}}\restriction X')$. In this frame
\begin{equation*}
\mathsf{m}^{\circ}=\mathfrak{c}\le\mathsf{m}=\ell=m ,\qquad l_x=\min\{p,\mathfrak{c}\},\qquad g=m-\mathfrak{c} .
\end{equation*}
By (VIII)(i) with $\Delta=g$, the stage-$l$ layer moves $Q_e^{(p,X',\mathfrak{s})}(y')$, $e\in\{2,4\}$, by at
most $\iota_l(y')(4.4/L^2)^g$ times the birth radius $r_e(\mathfrak{c},p\wedge\mathfrak{c})$.

\emph{Clause (i): the additive term.} Since $\mathfrak{c}=\ell-\frac12\log_{L_0}\omega$, we have
$\omega=L_0^{2g}\ge1$. The radii ratio satisfies
\begin{equation*}
\frac{\omega\,r_e(m,l_x)}{r_e(\mathfrak{c},p\wedge\mathfrak{c})}
=L_0^{2g}\,\frac{\alpha_{0,m}}{\alpha_{0,\mathfrak{c}}}\,L^{-2g}\le\Bigl(\frac{0.381}{L}\Bigr)^g .
\end{equation*}
This bound is the restriction on $L_0$ imposed in \cite{B-CMP122a}, applied at the single point $y'$.
Hence the additive term alone satisfies
\begin{equation*}
\pi r_e(m,l_x)\le\pi\,\omega\,r_e(m,l_x)\le\pi\Bigl(\frac{0.381}{L}\Bigr)^g r_e(\mathfrak{c},p\wedge\mathfrak{c}) .
\end{equation*}

\emph{Clause (i): comparison with the room.} By (IX), the one-step room is
\begin{equation*}
[\theta^{(l)}_{\mathfrak{s}}-\theta^{(l+1)}_{\mathfrak{s}}](y')\,r_e(\mathfrak{c},p\wedge\mathfrak{c})
=\nu_{\mathfrak{c}}\,r_e(\mathfrak{c},p\wedge\mathfrak{c}) .
\end{equation*}
Since $|\mathfrak{c}-c_\ast|\le|m-c_\ast|+(m-\mathfrak{c})$, we have $\nu_{\mathfrak{c}}\ge2^{-g}\nu_m$.
By (IX), $\pi\le\iota_l\le\frac16\mathsf{w}_0\nu_m$. The depth factors therefore pair, at $L\ge12$:
\begin{equation*}
2^g\Bigl(\frac{4.4}{L^2}\Bigr)^g=\Bigl(\frac{8.8}{L^2}\Bigr)^g\le(0.062)^g ,
\qquad
2^g\Bigl(\frac{0.381}{L}\Bigr)^g\le(0.064)^g .
\end{equation*}
Hence the increment is at most
\begin{equation*}
\tfrac16\mathsf{w}_0(0.064)^g\nu_{\mathfrak{c}}\,r_e(\mathfrak{c},p\wedge\mathfrak{c})
\le\tfrac12\nu_{\mathfrak{c}}\,r_e(\mathfrak{c},p\wedge\mathfrak{c}) .
\end{equation*}
This is half of the one-step room.

\emph{Clause (i): kind (A).} In kind (A) the step is stated in the frozen weighted unit $\mathsf{u}_m$ at
the frozen level $m$. By (VIII)(ii) the one-step room at that level is $\nu^{(l+1)}_m\cdot\mathsf{u}_m$,
the increment is at most $\iota_{l+1}\cdot\mathsf{u}_m$, and
$\iota_{l+1}\cdot\mathsf{u}_m\le\frac12\nu^{(l+1)}_m\cdot\mathsf{u}_m$, that is, the increment
is at most half of the one-step room.
\end{proof}

\begin{remark}
One other route is not used in the proof of (b). At boxes meeting a lifted point, membership is not
derived from the composite clauses at birth radii. Put $g':=\ell_{\mathfrak{B}}-\mathfrak{c}$. At such a
clause these would give only a size of order $\alpha_{0,\mathfrak{c}}L^{2g'}$ in $\mathfrak{B}$'s unit,
and this is not bounded by $\gamma_0^{(k)}/5$.

The depth gain in (i) cannot be replaced by the bump estimate for the stage-$l$ layer alone under the
convention of (IX), in which the one-step room at the triple is $\nu_{\mathfrak{c}(y')}$ times the birth
radius, the running
factor being evaluated at the level $\mathfrak{c}$ of the retained structure. From the bump estimate
the move is at most $\pi_\times E+\pi r_e(m,l_x)$, where $E$ is the size of the entry before the step.
Here
$E<\theta^{(l)}(\mathfrak{c})\cdot r_e(\mathfrak{c},p\wedge\mathfrak{c})$, and this is at most half the
birth radius by~\eqref{4.28:eq-lifted-thresholds-a}. The additive part has the gain, but the
multiplicative part $\frac12\pi_\times$ has none, since only $\pi_\times\le\pi$ is available. In the worst
case $\nu_{\mathfrak{c}}=2^{-g}\nu_m$, it would require $2^g\lesssim1/\mathsf{w}_0$. That is a bound on the
depth, and no such bound is imposed.

Under the alternative convention in which the running factor is evaluated at the point's current shell
$\ell(y')$, the
one-step room at the triple is $\nu_{\ell(y')}$ times the birth radius. The move is then at most
\begin{equation*}
\bigl(\tfrac12\pi_\times+\pi\bigr)r_e(\mathfrak{c},p\wedge\mathfrak{c})
\le1.5\,\pi\,r_e(\mathfrak{c},p\wedge\mathfrak{c})
\le\iota_l\,r_e(\mathfrak{c},p\wedge\mathfrak{c})
\le\tfrac12\nu_\ell\,r_e(\mathfrak{c},p\wedge\mathfrak{c}) ,
\end{equation*}
using $2.4\pi\le\iota_l$. So, under that convention, the increment is again at most half of the
one-step room, and this needs only the bump estimate. Under that convention
the thresholds at the starting point of the chain of bounds in the proof of
Proposition~\ref{4.28:prop-near-pair-box} are at most $0.61$, $0.51$ and $0.55$ of the unit of that
chain. The rooms are at
least $0.19$ (when $\ell_T=\mathfrak{c}$) and at least $0.74$ (when $\ell_T<\mathfrak{c}$). These rooms are
still at least $\beta/2$ and larger than $0.039$, the bound on the three costs
in~\eqref{4.28:eq-near-pair-box-c}.
\end{remark}

\begin{definition}[Notation for sums over the refined cover]\label{4.28:not-cover-sum-notation}
Let $\mathfrak{d}=(C;\kappa,\mathsf{h},\kappa_0,J)$ be a lifted datum. The level function $\ell$,
the weight $\omega$, the cutting level $\mathfrak{c}$ and the lifted set $\mathsf{L}_{\mathfrak{d}}$
are those of Definition~\ref{4.28:def-cutting-level}. We write $\mathcal{C}$ for
Ba{\l}aban's cover of the large-field region, with its partition of unity, as in
\cite{B-CMP119}, and $\mathcal{C}^1$ for the refined cover of
Definition~\ref{4.28:def-refined-cover}, which is built on the refined block structure
$\mathfrak{B}^{\mathfrak{c}}$.

The symbols below are used at a stage at which $\mathfrak{d}$ is not in the run. This covers
data of kind (A), data of kind (B$'$), and the pieces of the construction that lie on
$\mathfrak{T}$; the kinds are those of the classification of lifted data in this chapter.
Throughout, $\square$ ranges over the boxes $\square\in\mathcal{C}\cup\mathcal{C}^1$
that meet $\mathsf{L}_{\mathfrak{d}}$, and $d=4$.

\begin{enumerate}
\item \emph{Depth.} The depth of $\square$ is
\begin{equation*}
u(\square):=J-\min_{\square}\mathfrak{c}\ \ge 1 .
\end{equation*}
\item \emph{Multiplicity weight.} With $u=u(\square)$, we set
\begin{equation*}
\mathsf{m}(\square):=
\begin{cases}
(2L^{u+1})^{d}\,(u+4), & \square\in\mathcal{C},\\
1, & \square\in\mathcal{C}^1 .
\end{cases}
\end{equation*}
\item \emph{Exit set.} The exit set $\partial_{\mathfrak{d}}$ is the set of $\pi_J$-cubes of
$\Omega_J$ that touch $\Omega^c_J\cap C$.
\item \emph{Distance to the integration regions.} Let $n$ be a stage, with integration region
$R^{(n)}$, and let $\bar c'$ be a component of the large-field region $Z$. We set
\begin{equation*}
D(\square):=d_{\mathrm{sc}}\bigl(\square,\,R^{(n)}\cap\bar c'\bigr),
\end{equation*}
where $d_{\mathrm{sc}}=d^{\wedge}_{\mathfrak{B}}$ is the scaled distance of the block structure
$\mathfrak{B}$ of Definition~\ref{4.28:def-refined-cover}, as fixed in the statement on scaled
distances of this chapter.

This definition is made for every stage $n$ and every component $\bar c'$ of $Z$. In each
use, $n$ is the stage whose sum over boxes contains $D(\square)$. One such sum is the one in
the volume constant
\begin{equation*}
V^{\mathrm{fr}}=\sup_{n,\,\bar c'}\ \sum_{\square} e^{-\frac12\delta D(\square)},
\end{equation*}
where $\delta$ is the decay rate fixed in the level estimate of this chapter. Another is the
stage-$n$ summation over
boxes of the summation statement for the refined cover. The first clause of the lemma in which
these sums are bounded is proved for each fixed stage $n$, uniformly.

For a datum of kind (B$'$), let $n^{+}$ be the last stage at which it is in the run; only the
stages $n\ge n^{+}+1$ occur, and at the stages $n\le n^{+}$ the datum is in the run and is of
kind (B).
\item \emph{Distance of the exit set.} With $n$ and $\bar c'$ as in the preceding item, and with
$D$ evaluated on the cubes of $\partial_{\mathfrak{d}}$ by the same formula, we set
\begin{equation*}
d^-_{\mathfrak{d}}:=\min_{\partial_{\mathfrak{d}}} D .
\end{equation*}
\end{enumerate}
\end{definition}

\begin{lemma}[distance from an old region's lifted set to the integration regions]
\label{4.28:lem-old-region-distance}
Let $d=4$, $L\ge 12$ and $R_{\min}\ge 5L$. Let the decay rate $\delta$ and the cost $w$ per crossed
$M$-cube satisfy the level condition $\delta w\ge 16\ln L+4$. Assume the following.
\begin{itemize}
\item[(a)] The per-cube cost of the scaled distance: along every contour, each $M$-cube of the
local scale of $\mathfrak{B}$ that the contour crosses contributes at least $w$ to
$d_{\mathrm{sc}}=d^{\wedge}_{\mathfrak{B}}$.
\item[(b)] The nesting of live arrested data: if live arrested pieces $\mathfrak{a}$ and
$\mathfrak{b}$, with data $(k_a,h_a,k_0^a,j^+_a)$ and $(k_b,h_b,k_0^b,j^+_b)$, are nested (one of
$X_{\mathfrak{a}}$, $X_{\mathfrak{b}}$ contains the other), then, with $\mathfrak{a}$ the piece of
smaller site index,
\begin{equation*}
j^+_a\le k_a\le h_b<j^+_b.
\end{equation*}
\item[(c)] The multiplicity of cube families: for every level $m$, a $\pi_m$-cube lies in at most
$2^d$ cubes of each family of cubes of $\mathcal{C}\cup\mathcal{C}^1$; the cover $\mathcal{C}^1$
consists of two families and $\mathcal{C}$ of at most four.
\item[(d)] The frozen-volume bound: let $\mathcal{F}$ be a family of lifted data not in the run
such that each datum of $\mathcal{F}$ satisfies (F3) below; each cube of $\mathcal{C}^1$ with $u=1$
meeting the lifted set of a datum of $\mathcal{F}$ is assigned to a cube of $\mathcal{C}$ of scale
$J$ containing it, each cube of $\mathcal{C}$ receiving at most $(2L)^d$ cubes and $D$ not
decreasing under the assignment; $\mathsf{m}=5(2L^2)^d$ at $u=1$; and a cube of $\mathcal{C}$ of
scale $J$ serves at most $2^d$ data of $\mathcal{F}$. Then, with the sums defining $V^0$ and
$V^{\mathrm{fr}}$ taken over $\mathcal{F}$, one has $V^{\mathrm{fr}}\le 6(2L^2)^dV^0$ and
$\max\{V^0,V^{\mathrm{fr}}\}\le C_FV^{0,\mathrm{th}}$ with $C_F=12(2L^2)^d$.
\item[(e)] For a lifted datum of kind (B$'$) of Definition~\ref{4.28:def-attempt-data}, not in the
run at the stage $n$ considered: the integration region of stage $n$ is contained in $Z^{[l]}$ and
misses $C$, $J\le k-1$, every contour from $\partial_{\mathfrak{d}}$ to that integration region
crosses the collar $C\cap\Omega_k$, and consequently the bounds (F1)--(F4) below hold with
$d^-_{\mathfrak{d}}\ge 2R_kw$ in place of $d^-_{\mathfrak{d}}\ge 3LR_{h+1}w$.
\end{itemize}
Let $(k,l)$ be the current site and stage, and let $h$ be the base level of the preparation at
$k$ (the level $\mathsf{h}$ of Definition~\ref{4.28:def-attempt-data} at the site $k$). Let
$\mathfrak{d}=(C;\kappa,\mathsf{h},\kappa_0,J)$ be a lifted datum of kind (A) of
Definition~\ref{4.28:def-attempt-data}, so that $C=X_{\mathfrak{a}}$, $\kappa=k_a$,
$\mathsf{h}=h_a$, $\kappa_0=k_0^a$, $J=j^+_a$, and so that, by the placement of kind (A) in that
definition, $k_a\le h$ and
$X_{\mathfrak{a}}\subset Z_{k_a}\subset Z_h$. Let $\mathsf{L}_{\mathfrak{d}}$ be the lifted set
of~\eqref{4.28:eq-cutting-level-b}, let $u(\square)$, $\mathsf{m}(\square)$,
$\partial_{\mathfrak{d}}$, $D(\square)$, $d^-_{\mathfrak{d}}$, $V^0$, $V^{\mathrm{fr}}$ be as in
Notation~\ref{4.28:not-cover-sum-notation}, and let $\rho_m$ be as in
Lemma~\ref{4.28:lem-cutting-strata}(v). Then (F1) and (F2) hold for every
$\square\in\mathcal{C}\cup\mathcal{C}^1$ meeting $\mathsf{L}_{\mathfrak{d}}$, with
$m:=\min_{\square}\mathfrak{c}$ and $u:=u(\square)$, and (F3), (F4) hold:
\begin{equation}\label{4.28:eq-old-region-distance-a}
\begin{aligned}
&\text{(F1)}\quad D(\square)\ge d^-_{\mathfrak{d}}+6w(u-1)\rho_m,
\qquad d^-_{\mathfrak{d}}\ge 3LR_{h+1}w,\\
&\text{(F2)}\quad (3L^{u+1})^{2d}(u+4)\,e^{-\frac12\delta D(\square)}\le 1,\\
&\text{(F3)}\quad \sum_{\substack{\square\in\mathcal{C}\cup\mathcal{C}^1,\
\square\cap\mathsf{L}_{\mathfrak{d}}\ne\emptyset\\ u(\square)\ge 2}}
\mathsf{m}(\square)\,e^{-\frac12\delta D(\square)}
\le L^{-20}e^{-\frac12\delta d^-_{\mathfrak{d}}},\\
&\text{(F4)}\quad V^{\mathrm{fr}}\le 6(2L^2)^dV^0,\qquad
\max\{V^0,V^{\mathrm{fr}}\}\le C_FV^{0,\mathrm{th}},\qquad C_F=12(2L^2)^d,
\end{aligned}
\end{equation}
where in (F4) the sums defining $V^0$ and $V^{\mathrm{fr}}$ are taken over all data not in the run,
that is, over $\mathcal{C}^1$ with the cubes of the run counted in~\eqref{4.28:eq-refined-count-a}
removed, and $C_F$ is the same constant as in (d). These bounds are absolute: their constants depend
on none of $k-k_a$, the numbers of stages $N_k$, $N_{k_a}$ and $N_0(k_a)$ of the preparations at
$k$ and at $k_a$, $J-\kappa_0$, nor on the coupling. For a datum of kind (B$'$) of
Definition~\ref{4.28:def-attempt-data} the bounds (F1)--(F4) hold with $d^-_{\mathfrak{d}}\ge 2R_kw$
in place of $d^-_{\mathfrak{d}}\ge 3LR_{h+1}w$.
\end{lemma}

\begin{proof}
Throughout, the stage $n$ and the component $\bar{c}'$ of $Z$ entering
$D(\square)=d_{\mathrm{sc}}(\square,R^{(n)}\cap\bar{c}')$ are fixed, $R^{(n)}$ denoting the
integration region of stage $n$; every bound below is uniform in $n$ and $\bar{c}'$.

\emph{(F1), the bound on $d^-_{\mathfrak{d}}$.} Let $\bar{c}$ be the component of $Z$ containing
$X_{\mathfrak{a}}$. Each cube of $\partial_{\mathfrak{d}}$ lies within one $\pi_J$-side of
$\Omega^c_J\cap C$, and one $\pi_J$-side is less than $2R_J$ $\pi_J$-sides; by
Lemma~\ref{4.28:lem-cutting-strata}(iii) at $\ell=J$ it therefore does not meet $\Lambda_J$. With
the placement of kind (A) this gives
\begin{equation*}
\partial_{\mathfrak{d}}\subset Z_J\cap C\subset X_{\mathfrak{a}}\subset Z_h\cap\bar{c}.
\end{equation*}
The integration region satisfies $R^{(n)}\subset Z_{j'+1}\setminus Z''_{j'+1}$ with
$j'=h+n\ge h$; since the sets $Z''_i$ increase with $i$, $R^{(n)}$ lies off $Z''_{h+1}$. Every
component $\bar{c}'\ne\bar{c}$ is disjoint from $\bar{c}$ and so also lies off
$Z''_{h+1}\cap\bar{c}$. Hence $R^{(n)}\cap\bar{c}'$ does not meet $Z''_{h+1}\cap\bar{c}$. Write
${}^{\sim}$ for the addition of one layer of cubes and ${}^{\sim 3}$ for its threefold application,
and $Z'_h$ for the enlargement of \cite[(3.2)]{B-CMP119}. Reading only the inclusions
$\supseteq$ of~\eqref{4.28:eq-cutting-strata-a} at the site $k$, we have
\begin{equation*}
Z''_{h+1}\cap\bar{c}\supseteq (Z'_h)^{\sim 3}\supseteq Z'_h\supseteq Z_h\cap\bar{c}.
\end{equation*}
Therefore every contour from a cube $\Delta\in\partial_{\mathfrak{d}}$ to
$R^{(n)}\cap\bar{c}'$ crosses the three layers $(Z'_h)^{\sim 3}\setminus Z'_h$ of
$LMR_{h+1}$-cubes of scale $h$. These layers lie in $\Lambda_h\setminus\Omega_{h+1}$ and in
$Z''_{h+1}$. They carry level $h$ in the block structure $\mathbb{B}_k$ from which the construction
of \cite[(1.15)]{B-CMP122a} starts at the site $k$, and in every $\mathbb{B}^{(n)}_k$: in the
construction of the block structures in \cite[(1.15)]{B-CMP122a} the stratum $\Gamma_h$ of level $h$
is replaced by the two sets $(\Gamma_h\cap Z''^{c}_{h+1})^{(h+1)}$ and $\Gamma_h\cap Z''_{h+1}$,
and the second of these is never given a new level. An $LMR_{h+1}$-cube of scale $h$ has
$LR_{h+1}$ $M$-cubes of scale $h$ along each side. The contour thus crosses at least $3LR_{h+1}$
$M$-cubes of the local scale of $\mathfrak{B}$, and by (a) its scaled length is at least
$3LR_{h+1}w$. Taking the minimum over $\Delta\in\partial_{\mathfrak{d}}$ gives
$d^-_{\mathfrak{d}}\ge 3LR_{h+1}w$.

\emph{(F1), the inside part.} A contour from $\square$ to $R^{(n)}\cap\bar{c}'$ leaves
$\Omega^c_J\cap C$ through some $\Delta\in\partial_{\mathfrak{d}}$. Its scaled length is therefore
at least the length of its part inside plus $D(\Delta)\ge d^-_{\mathfrak{d}}$. The cube $\square$
meets at most the $\mathfrak{c}$-shells $m$ and $m+1$: for $\square\in\mathcal{C}^1$ this is
Lemma~\ref{4.28:lem-box-lemma}(a); for $\square\in\mathcal{C}$ the side of $\square$ is at most
$2M\ell_J<8\rho_mM\ell_J$, because $\rho_m\ge 1$. Moreover $\square$ reaches at most two
$\pi_J$-sides into the shell $m+1$. For $u\ge 2$ the contour crosses the rest of the shell $m+1$ and
the shells $m+2,\dots,J-1$. These are old shells above $\kappa_0$, at level $J$. By
Lemma~\ref{4.28:lem-cutting-strata}(v) each shell $m'$ among them is at least $8\rho_{m'}$ of its own
$M$-cubes wide, and $\rho_{m'}\ge\rho_m$ because $\rho_{m'}$ is non-decreasing in $m'$
(Lemma~\ref{4.28:lem-cutting-strata}(v)). The number of such cubes
crossed is therefore at least
\begin{equation*}
8(u-1)\rho_m-2\ge 6(u-1)\rho_m,
\end{equation*}
where the inequality holds because $2(u-1)\rho_m\ge 2$ for $u\ge 2$. By (a) each of these cubes
contributes at least $w$. This gives $D(\square)\ge d^-_{\mathfrak{d}}+6w(u-1)\rho_m$ for $u\ge 2$.

\emph{(F2).} By (F1), $\rho_m\ge 1$ and the level condition,
\begin{equation*}
\tfrac12\delta D(\square)\ge \tfrac32LR_{h+1}\delta w+3(u-1)\delta w
\ge LR_{h+1}(24\ln L+6)+(u-1)(48\ln L+12),
\end{equation*}
so that
\begin{equation*}
e^{-\frac12\delta D(\square)}\le (L^{-24}e^{-6})^{LR_{h+1}}\,L^{-48(u-1)}e^{-12(u-1)}.
\end{equation*}
Since $d=4$, we have $(3L^{u+1})^{2d}=3^8L^{8u+8}$. For $u\ge 1$ the exponents satisfy
$8u+8\le 48(u-1)+24$. Moreover $LR_{h+1}\ge 1$, so $L^{-24LR_{h+1}}\le L^{-24}$. Hence
\begin{equation*}
(3L^{u+1})^{2d}(u+4)e^{-\frac12\delta D(\square)}
\le 3^8(u+4)e^{-12(u-1)}\,e^{-6LR_{h+1}}.
\end{equation*}
The function $(u+4)e^{-12(u-1)}$ decreases in $u\ge 1$ and equals $5$ at $u=1$. Further
$3^8\cdot 5=32805<e^{11}$, and $e^{11}<e^{6LR_{h+1}}$ because $L\ge 12$ and
$R_{h+1}\ge R_{\min}\ge 1$. This proves (F2).

\emph{(F3).} Fix $m$ with $u=J-m\ge 2$, and let $S_m$ denote the part of the sum in (F3) over the
cubes $\square$ with $\min_\square\mathfrak{c}=m$. Since the partitions are compatible, such a cube
contains a $\pi_m$-cube of $\Omega^c_{m+1}\cap C$. By (c) a $\pi_m$-cube lies in at most
$2^d(2+4)\le 2^{d+3}$ cubes of $\mathcal{C}\cup\mathcal{C}^1$. By
Lemma~\ref{4.28:lem-cutting-strata}(v), $\Omega^c_{m+1}\cap C$ lies in a cube of $100\rho_m$
$\pi_J$-sides. One $\pi_J$-side is $L^u$ $\pi_m$-sides, so there are at most
$2^{d+3}(100\rho_mL^u)^d$ cubes in $S_m$. Each carries $\mathsf{m}(\square)\le(2L^{u+1})^d(u+4)$,
with equality on $\mathcal{C}$ and $\mathsf{m}=1$ on $\mathcal{C}^1$. By (F1) and the level
condition,
\begin{equation*}
e^{-\frac12\delta D(\square)}\le e^{-\frac12\delta d^-_{\mathfrak{d}}}\,
e^{-3\delta w(u-1)\rho_m}
\le e^{-\frac12\delta d^-_{\mathfrak{d}}}\,L^{-48(u-1)}e^{-12(u-1)\rho_m},
\end{equation*}
where $\rho_m\ge 1$ was used in the power of $L$. Multiplying, with $d=4$,
\begin{equation*}
S_m\le 2^{11}10^8\rho_m^4(u+4)e^{-12(u-1)\rho_m}\cdot L^{8u+4-48(u-1)}
e^{-\frac12\delta d^-_{\mathfrak{d}}}.
\end{equation*}
For $u\ge 2$ the function $\rho^4e^{-12(u-1)\rho}$ decreases in $\rho\ge 1$, since its derivative is
$\rho^3e^{-12(u-1)\rho}(4-12(u-1)\rho)<0$. Hence it is at most $e^{-12(u-1)}$. The function
$(u+4)e^{-12(u-1)}$ decreases in $u\ge 2$. Therefore
\begin{equation*}
2^{11}10^8\rho_m^4(u+4)e^{-12(u-1)\rho_m}\le 2^{11}10^8\cdot 6e^{-12}<8\cdot 10^6<12^7\le L^7.
\end{equation*}
The exponent of $L$ becomes $8u+4-48(u-1)+7=-21-40(u-2)$. Distinct $m$ correspond to distinct
$u\ge 2$, and therefore
\begin{equation*}
\sum_{u\ge 2}L^{-21-40(u-2)}=\frac{L^{-21}}{1-L^{-40}}\le L^{-20}.
\end{equation*}
This gives (F3).

\emph{(F4).} We apply (d) to the family of all lifted data not in the run, to which the present
datum belongs. Its inputs concern each datum: the bound (F3); the assignment of each cube of
$\mathcal{C}^1$ with $u=1$ to a cube of $\mathcal{C}$ of scale $J$ containing it; the value of
$\mathsf{m}$ at $u=1$; and the fact that a cube of $\mathcal{C}$ of scale $J$ serves at most $2^d$
data. The environment of the datum enters them only through $d^-_{\mathfrak{d}}$, which appears in
(F3) as a parameter. Each of these inputs holds for the present datum.
\begin{itemize}
\item (F3) has been proved above.
\item In the assignment, each cube of $\mathcal{C}$ of scale $J$ receives at most $(2L)^d$ cubes of
$\mathcal{C}^1$. If $\square^1\subset\square$, then $D(\square^1)\ge D(\square)$.
\item By definition, $\mathsf{m}=(2L^2)^d(1+4)=5(2L^2)^d$ at $u=1$.
\item Nested live data have different top levels, by (b). Disjoint data have disjoint sets
$\Omega^c_J\cap C$, aligned at $MR_J\ell_J>4M\ell_J$, as $R_J\ge R_{\min}\ge 5L$.
\end{itemize}
The last three items use only the definitions, (b) and $R_{\min}\ge 5L$, and not the kind of the
datum. The bound (F3) holds for every lifted datum of kind (A) not in the run, by the argument above
applied to that datum, and for every lifted datum of kind (B$'$) not in the run by (e). Hence (d)
applies to the family of all lifted data not in the run, with its sums taken over that family, and
gives (F4) with the same $C_F=12(2L^2)^d$.

\emph{Absoluteness.} The numerical constants obtained above ($3$, $6$, $1$, $L^{-20}$,
$6(2L^2)^d$ and $C_F$) arise only from $d=4$, $L\ge 12$, $R_{\min}\ge 5L$, the level condition and
$\rho_m\ge 1$. None of $k-k_a$, $N_k$, $N_{k_a}$, $N_0(k_a)$, $J-\kappa_0$ or the coupling enters
them.

\emph{Data of kind (B$'$).} For such a datum the bounds (F1)--(F4) with
$d^-_{\mathfrak{d}}\ge 2R_kw$ are the content of (e).
\end{proof}

In what follows, a length of $a$ \emph{$\pi_n$-sides} is $a\,M\ell_n$, and
$\ell_n/\ell_{n'}=L^{n-n'}$ for levels $n,n'$.

\begin{lemma}[Counting the refined cubes along the current preparation]\label{4.28:lem-refined-count}
Let $\mathfrak{B}=\mathbb{B}^{(l-1)}_k$ be the block structure at stage $l-1$ of the preparation at
the site $k$, $\ell_{\mathfrak{B}}$ its level function, $\mathcal{C}$ Ba{\l}aban's cover and
partition of unity \cite{B-CMP119} built on $\mathfrak{B}$, as in
Definition~\ref{4.28:def-refined-cover}, and $\mathcal{C}^1$ the refined cover of
Definition~\ref{4.28:def-refined-cover}. Let $\mathfrak{d}=(C;\kappa,\mathsf{h},\kappa_0,J)$ be the
lifted datum carried by the preparation, with $\kappa=k$ the site of the preparation and $C$ its
class component, and with cutting level $\mathfrak{c}$ and lifted set $\mathsf{L}_{\mathfrak{d}}$ as in
Definition~\ref{4.28:def-cutting-level}; write $k_0$ for its level $\kappa_0$,
$N_0=N_0(\kappa)=\kappa-\kappa_0$, and $j=J-1$ for the
level below its top level $J$. Assume $R_{\min}\ge 5L$, the hypothesis of
Lemma~\ref{4.28:lem-cutting-strata}, so that that lemma applies to $\mathfrak{d}$. The regions
$\Omega_i$, $Z_i$, $Z''_i$ are those of the site cut to $C$, and $\mathfrak{c}$ is defined on $C$
(Definition~\ref{4.28:def-cutting-level}); intersections of these sets with $\bar{c}$ below are
therefore subsets of $\bar{c}\cap C$.
For $m\ge k_0$ the \emph{$\mathfrak{c}$-shell} $m$ is the set
$\{\mathfrak{c}=m\}$, and the \emph{family} $m$ is the set of cubes of $\mathcal{C}^1$ of scale $m$
determined by the stratum $\{\mathfrak{c}=m\}$ of $\mathfrak{B}^{\mathfrak{c}}$, that is, those
containing a $\pi_m$-cube of $\{\mathfrak{c}=m\}$; for
$m>k_0$ the $\mathfrak{c}$-shell $m$ is the old shell $\Omega_m\setminus\Omega_{m+1}$.

If $j\le k_0$, then $\mathfrak{c}=\ell_{\mathfrak{B}}$ and $\mathcal{C}^1=\mathcal{C}$ along the
preparation. Let $j>k_0$, let $\bar{c}$ be a component of $Z^{[l]}$, and assume:
\begin{itemize}
\item[$(\alpha)$] for every $m$ with $k_0\le m\le j-1$, the set $Z_{m+1}\cap\bar{c}$ is a single
parallelepiped contained in a cube of side $100LR_{m+1}$ $\pi_m$-sides;
\item[$(\beta)$] the window envelope $\check{R}$ of the radii satisfies $R_{m+1}\le\check{R}$ for
every $m$ with $k_0<m+1\le j$;
\item[$(\gamma)$] the logarithmic count of the present coupling satisfies
$(2L^{N_0})^d\le\mathsf{l}_j^2$;
\item[$(\delta)$] in the bound on the overlaps of the cubes of the cover (in what follows, the
overlap bound), for every $\square'\in\mathcal{C}$ whose enlarged cube $\tilde{\square}'$ meets
$\{\mathfrak{c}<\ell_{\mathfrak{B}}\}$, the cube $\tilde{\square}'$ has side at most four
$\pi_j$-sides and is tiled by cubes of $\mathcal{C}^1$ of families $m\ge k_0$, each of side two
$\pi_m$-sides.
\end{itemize}
Then the following hold.
\begin{itemize}
\item[(a)] The number of cubes of $\mathcal{C}^1\setminus\mathcal{C}$ on $\bar{c}$ obeys bound (a)
of \eqref{4.28:eq-refined-count-a}.
\item[(b)] Let $\Delta$ be a $\pi_n$-cube, $n\le j+1$, in $\bar{c}$, lying in
$\{\mathfrak{c}=k_0\}=\Omega^c_{k_0+1}\setminus Z''_{k_0}$ or in one old shell; every $\pi_n$-cube
with $n\le j+1$ in $\bar{c}$ that meets $\mathsf{L}_{\mathfrak{d}}$ is of this kind. The number of
cubes of $\mathcal{C}^1$ of the
families $m\ge k_0$ touching $\Delta$ obeys bound (b) of \eqref{4.28:eq-refined-count-a}. Every
cube of $\mathcal{C}^1\setminus\mathcal{C}$ belongs to these families; the cubes of $\mathcal{C}$
of scale below $k_0$ keep the multiplicity with which they are counted in the multiplicity bound
of the sum over cubes.
\item[(c)] For $\square'\in\mathcal{C}$ whose enlarged cube $\tilde{\square}'$ meets
$\{\mathfrak{c}<\ell_{\mathfrak{B}}\}$, the number of cubes of $\mathcal{C}^1$ in the tiling of
$\tilde{\square}'$ in the overlap bound obeys bound (c) of \eqref{4.28:eq-refined-count-a}. (Off
$\{\mathfrak{c}<\ell_{\mathfrak{B}}\}$ the covers $\mathcal{C}^1$ and $\mathcal{C}$ coincide, and
the overlap bound is unchanged.)
\end{itemize}
\begin{equation}\label{4.28:eq-refined-count-a}
\begin{aligned}
\text{(a)}&\quad \#\{\square\in\mathcal{C}^1\setminus\mathcal{C}\ \text{on}\ \bar{c}\}
 \le 2^d(N_0-1)(100L\check{R})^d,\\
\text{(b)}&\quad \#\{\square\in\mathcal{C}^1\ \text{of a family}\ m\ge k_0:
 \ \square\ \text{touches}\ \Delta\}\le 3(L^{N_0}+3)^d\le\mathsf{l}_j^2,\\
\text{(c)}&\quad \#\{\square\in\mathcal{C}^1:\ \square\ \text{in the tiling of}\
 \tilde{\square}'\}\le\mathsf{l}_j^2.
\end{aligned}
\end{equation}
\end{lemma}

Hypotheses $(\alpha)$--$(\delta)$ are supplied, in this order, by the container statement for the
regions $Z_{m+1}$ of the preparation together with the nesting of
those regions, by the window-envelope bound $R_{m+1}\le\check{R}$ on the radii, by the bound on
the logarithmic count $\mathsf{l}_j$ of the present coupling, and by the tiling of the enlarged
cubes $\tilde{\square}'$ in the overlap bound.

\begin{proof}
\emph{Clause} (a). By Definition~\ref{4.28:def-refined-cover}, the refined block structure
$\mathfrak{B}^{\mathfrak{c}}$ agrees with $\mathfrak{B}$ off the set $\{\mathfrak{c}<\ell_{\mathfrak{B}}\}$,
so $\mathcal{C}^1$ differs from $\mathcal{C}$ only on this set. On $\bar{c}$ it is
\begin{equation*}
\{\mathfrak{c}<\ell_{\mathfrak{B}}\}\cap\bar{c}=(\Omega^c_j\setminus Z''_{k_0+1})\cap\bar{c},
\end{equation*}
so the cubes of $\mathcal{C}^1\setminus\mathcal{C}$ on $\bar{c}$ belong to the families of the
$\mathfrak{c}$-shells $m$ with $k_0\le m\le j-1$.

Fix such an $m$. By Lemma~\ref{4.28:lem-cutting-strata}\,(i), the $\mathfrak{c}$-shell $m$ is
$\Omega^c_{k_0+1}\setminus Z''_{k_0}$ if $m=k_0$ and the old shell $\Omega_m\setminus\Omega_{m+1}$
if $m>k_0$; in both cases it lies in $\Omega^c_{m+1}$. By \eqref{4.28:eq-cutting-strata-a},
$\Omega^c_{m+1}\subset Z_{m+1}$. Hence
\begin{equation*}
\{\mathfrak{c}=m\}\cap\bar{c}\subset\Omega^c_{m+1}\cap\bar{c}\subset Z_{m+1}\cap\bar{c},
\end{equation*}
and by hypothesis $(\alpha)$ the right-hand side is one parallelepiped inside a cube of side
$100LR_{m+1}$ $\pi_m$-sides. This cube contains at most $(100LR_{m+1})^d$ $\pi_m$-cubes, so the
$\mathfrak{c}$-shell $m$ on $\bar{c}$ contains at most that many.

Each cube of the family $m$ contains a $\pi_m$-cube of the shell, and each $\pi_m$-cube lies in at
most $2^d$ cubes of the family. Assigning to each cube of the family a $\pi_m$-cube of the shell
that it contains, every $\pi_m$-cube is assigned at most $2^d$ times. The number of cubes of the
family $m$ on $\bar{c}$ is therefore at most $2^d(100LR_{m+1})^d$, and by hypothesis $(\beta)$,
applicable since $k_0<m+1\le j$,
\begin{equation*}
2^d(100LR_{m+1})^d\le 2^d(100L\check{R})^d .
\end{equation*}
There are $j-k_0$ admissible values of $m$, and, because the top level $J=j+1$ does not exceed the
site $\kappa=k$ (Definition~\ref{4.28:def-cutting-level}),
\begin{equation*}
j-k_0=J-1-\kappa_0\le\kappa-1-\kappa_0=N_0-1 .
\end{equation*}
Summing over $m$ gives bound (a) of \eqref{4.28:eq-refined-count-a}.

\emph{Clause} (b). Let $\Delta$ be a $\pi_n$-cube with $n\le j+1\le k$, in $\bar{c}$, of the
kind described in (b). First, a $\pi_n$-cube meeting
$\mathsf{L}_{\mathfrak{d}}=(\Omega^c_J\setminus Z''_{k_0+1})\cap C$, with $J=j+1$
(Definition~\ref{4.28:def-cutting-level}), is of this kind. Indeed, by
\eqref{4.28:eq-cutting-strata-a},
\cite[(1.10)]{B-CMP122a} and Lemma~\ref{4.28:lem-cutting-strata}\,(v), the boundaries
$\partial Z''_{k_0+1}=\partial Z'^{\sim}_{k_0}$ and $\partial\Omega_m$, $m>k_0$, run along
$\pi_k$-cubes. Hence either $\Delta\subset\Omega^c_{k_0+1}\setminus Z''_{k_0+1}$, which lies in
$\{\mathfrak{c}=k_0\}$ because $Z''_{k_0}\subset Z''_{k_0+1}$, the sets $Z''_i$ defined in
\cite[(1.10)]{B-CMP122a} being increasing in $i$ (see the proof of
Lemma~\ref{4.28:lem-cutting-strata}); or $\Delta$ lies in one old shell
$m$ with $k_0<m\le J-1$.

In either case $\Delta$ lies in one $\mathfrak{c}$-shell $m\ge k_0$. Among the families
$m'\ge k_0$, it is touched only by cubes of the families $m'=m-1$ (if $m>k_0$), $m'=m$ and
$m'=m+1$, and by at most $(L^t+3)^d$ cubes of each such family $m'$, where $L^t$, $t=n-m'$, is the
side of $\Delta$ in $\pi_{m'}$-sides. Since $m'\ge k_0$ and $n\le j+1$,
\begin{equation*}
t\le j+1-k_0\le N_0 ,
\end{equation*}
the last inequality being $j-k_0\le N_0-1$, shown in (a). The number of cubes
of $\mathcal{C}^1$ of the families $m\ge k_0$ touching $\Delta$ is therefore at most
$3(L^{N_0}+3)^d$. By the proof of (a), the cubes of $\mathcal{C}^1\setminus\mathcal{C}$ belong to
families with $k_0\le m\le j-1$, so they are among those counted.

For the second inequality of (b), recall $d=4$. Since $L$ is odd and greater than one and
$N_0\ge 3$ by Lemma~\ref{4.28:lem-cutting-strata}\,(iv), we have $L^{N_0}\ge 27$, so
$3\le\tfrac14L^{N_0}$ and $L^{N_0}+3\le\tfrac54L^{N_0}$. As $3(5/4)^4=1875/256<16$,
\begin{equation*}
3(L^{N_0}+3)^4\le 3\Bigl(\tfrac54\Bigr)^4L^{4N_0}<16\,L^{4N_0}=(2L^{N_0})^4 .
\end{equation*}
By hypothesis $(\gamma)$, $(2L^{N_0})^d\le\mathsf{l}_j^2$. This is bound (b) of
\eqref{4.28:eq-refined-count-a}.

\emph{Clause} (c). Let $\square'\in\mathcal{C}$ be such that $\tilde{\square}'$ meets
$\{\mathfrak{c}<\ell_{\mathfrak{B}}\}$. By hypothesis $(\delta)$ it remains to count the tiling
cubes. By hypothesis $(\delta)$, $\tilde{\square}'$ has
side at most four $\pi_j$-sides, and the tiling cubes have side two $\pi_m$-sides with $m\ge k_0$.
Since
$\ell_j/\ell_m=L^{j-m}\le L^{j-k_0}$, each tiling cube has side at least
$2M\ell_j L^{-(j-k_0)}$. The cubes of a tiling have disjoint interiors, so their number is at most
the ratio of the volumes, and
\begin{equation*}
\frac{\operatorname{vol}(\tilde{\square}')}{\bigl(2M\ell_jL^{-(j-k_0)}\bigr)^d}
\le\Bigl(\frac{4M\ell_j}{2M\ell_jL^{-(j-k_0)}}\Bigr)^d=(2L^{j-k_0})^d .
\end{equation*}
As $j-k_0\le N_0-1$, hypothesis $(\gamma)$ gives
\begin{equation*}
(2L^{j-k_0})^d\le(2L^{N_0})^d\le\mathsf{l}_j^2 .
\end{equation*}
This is bound (c) of \eqref{4.28:eq-refined-count-a}.
\end{proof}

\begin{lemma}[Sums over near pairs at lifted points. Stated; the proof below is incomplete at the step
marked $(\ast)$]\label{4.28:lem-near-pair-sums}
Let $\mathfrak{d}=(C;\kappa,\mathsf{h},\kappa_0,J)$ be an attempt datum, with lifted set
$\mathsf{L}_{\mathfrak{d}}$ and cutting level $\mathfrak{c}$ as in \eqref{4.28:eq-cutting-level-b}.
Let $\mathcal{C}$ be the cover of the region, and let $\mathcal{C}^1$ be the refined cover built on the
refined block structure $\mathfrak{B}^{\mathfrak{c}}$. For the current preparation let $j>k_0$ and let
$\bar{c}$ be a component of $Z^{[l]}$, as in Lemma~\ref{4.28:lem-refined-count}. Assume the
conditions on $L$, $L_0$, $R_{\min}$, $\beta_0$, $\delta w$, $N_0$ and $N$, and on the radii
$\alpha_{\cdot,n}$, imposed in the hypotheses of Lemmas~\ref{4.28:lem-box-lemma},
\ref{4.28:lem-old-region-distance} and \ref{4.28:lem-refined-count}, and the following statements:
\begin{itemize}
\item the far-pair clauses of the sum lemma (its bound, its multiplicity clause and its budget clause
for far pairs) and the near-term bound of its clause (c);
\item clauses (a)--(c), the multiplicity clause and the budget clause with its volume inequality of
the sum lemma, and the three recursive bounds whose proofs use the overlap statement;
\item the single-cube estimate, with its bounds (S3) and (S5), in particular (S5)$(\alpha)$;
\item the overlap statement;
\item the kernel-membership proposition.
\end{itemize}
Then the following hold. In clauses (b)--(f) the datum lies on the current preparation, and the cover
is the refined cover $\mathcal{C}^1$.
\begin{enumerate}
\item[(a)] \emph{Data not in the current preparation.} For data $\mathfrak{d}$ of kinds (A) and
(B$'$), the far-pair clauses of the sum lemma hold on $\mathsf{L}_{\mathfrak{d}}$.
\item[(b)] \emph{The current preparation, clauses (a)--(c) of the sum lemma.} Clauses (a), (b) and
(c) of the sum lemma hold as stated on the current preparation with the cover $\mathcal{C}^1$, the
sums including the new tails, which belong to the class of terms these clauses treat.
Here a \emph{new tail} is a term $T$ whose set
$X_T$ lies in no $\tilde{\square}^2$ with $\square\in\mathcal{C}^1$.
\item[(c)] \emph{The current preparation, the family $\mathcal{T}_1$.} The index $s(\square)$ of
the single-cube estimate vanishes at every $\square\in\mathcal{C}^1$, and
\begin{equation}\label{4.28:eq-near-pair-sums-a}
\sum_{\substack{\square\in\mathcal{C}^1\\ \tilde{\square}^5\cap\bar{c}\neq\emptyset}}
\vartheta^{\natural}(\square)
\le\bigl[V+2^d(N_0-1)(100L\check{R})^d\bigr]\,\vartheta^{\natural};
\end{equation}
this bound serves in place of the near-term bound of clause (c) of the sum lemma.
\item[(d)] \emph{Multiplicity.} The multiplicity clause of the sum lemma holds as stated with the
cover $\mathcal{C}^1$.
\item[(e)] \emph{Budget.} The budget clause of the sum lemma holds with the cover $\mathcal{C}^1$,
with its third bracketed summand
omitted and with the coefficient $4^d$ in its second member replaced by $4^d+2^{d+1}$, where
\begin{equation*}
4^d+2^{d+1}=288\le 625=5^d\qquad (d=4).
\end{equation*}
Its volume inequality holds with the coefficient $5^d$:
\begin{equation}\label{4.28:eq-near-pair-sums-1}
5^d(N_0+1)(N_0+3)(100L\check{R})^d\le\frac{5^d}{C_V}\,N\bar{V},
\end{equation}
where $N$ and $N_0$ are as in Lemma~\ref{4.28:lem-box-lemma}, so that $(N_0+1)(N_0+3)\le N$,
and its right-hand side is absorbed in the budget $K'N\mathsf{r}\bar{V}$, the factor $5^d/C_V$ being
absorbed in $K'\mathsf{r}$, where $K'$ no longer carries the
factors $c_L$, $L_{\circ}^6$ and $2\cdot 52^d$. The degrees $6r$ and $4r$, the floors, and the
condition of the budget clause denoted there $(Q^{+})^{\natural,N}$ are as stated there.
\item[(f)] \emph{Overlap.} The overlap statement, that the total of the bounds (S3) over a tiling
carries the factor $\mathsf{l}_j^2$, holds as a count of the groups of cubes of $\mathcal{C}^1$
summed under one cube of the stage. The three recursive bounds whose proofs use the overlap statement
hold with the overlap factor in their bracketed majorant equal to $\mathsf{l}_j^2$.
\end{enumerate}
\end{lemma}

\begin{proof}
We take the clauses in order.

\emph{Clause} (a). For data of kinds (A) and (B$'$) the proof of the far-pair clauses of the sum
lemma is carried out on $\mathsf{L}_{\mathfrak{d}}$ with the following inputs. The tiling used in
the overlap statement is taken to be the tiling by the cubes of $\mathcal{C}^1$. The bound (S5)
of the single-cube estimate is applied on thick $\mathfrak{c}$-shells, a $\mathfrak{c}$-shell being a
level set $\{\mathfrak{c}=m\}$. At most $u+4$ short stages
occur, $u=u(\square)$ being the depth of the cube. The second assertion of
Lemma~\ref{4.28:lem-old-region-distance}, see \eqref{4.28:eq-old-region-distance-a}, is used; it
leaves a factor $L^{24LR_{h+1}-16}$, with $h$ as in Lemma~\ref{4.28:lem-old-region-distance}, to
spare, and this factor pays for the hulls, the hull of a cube $\square$ of scale $i$ being a set that
exceeds $\square$ by five sides of $\pi_i$ and one cube of the local scale (clause (d) below).
Finally the fourth
assertion of Lemma~\ref{4.28:lem-old-region-distance} supplies the volume bounds, with its sums
taken over all data not in the current preparation.

\emph{Clause} (b). For a cube $\square_0\in\mathcal{C}$, the proof of clause (a) of the sum lemma
tiles the enlarged cube $\tilde{\square}_0$ by the
original cover. By the third count of Lemma~\ref{4.28:lem-refined-count}, see
\eqref{4.28:eq-refined-count-a}, this tiling is the tiling by $\mathcal{C}^1$. The proof then
applies (S5)$(\alpha)$ at each
tile $\square''\in\mathcal{C}^1$. The terms so bounded are exactly the terms $T$ with
\begin{equation*}
X_T\cap\tilde{\square}''\neq\emptyset,\qquad X_T\not\subset\tilde{\square}''^{2},\qquad
\square''\in\mathcal{C}^1 .
\end{equation*}
A new tail has $X_T$ in no $\tilde{\square}''^{2}$. Since the tiles cover $\tilde{\square}_0$, a new
tail with $X_T$ meeting $\tilde{\square}_0$ meets some tile $\square''$, hence meets
$\tilde{\square}''$, and is among these terms; so the new tails are included in the sum. In clause
(c) of the sum lemma these terms
are fine objects, each anchored once. For them the first clause of the overlap statement, which
applies (S5)$(\alpha)$ as stated, carries the factor $\mathsf{l}_j^2$. The corresponding display of
the kernel-membership proposition already carries the factor $\mathsf{l}_j^2$.

\emph{Clause} (c). At every cube $\square\in\mathcal{C}^1$ the bound (S3) of the single-cube estimate
applies as stated, and its conclusion there carries the index $s(\square)=0$; this is the first
assertion of clause (c). For
\eqref{4.28:eq-near-pair-sums-a} we split the sum into the cubes of $\mathcal{C}$ and the cubes of
$\mathcal{C}^1\setminus\mathcal{C}$. The first part is bounded by $V\vartheta^{\natural}$, by the
definition of $V$. For the second part we use two facts. By the first count of
Lemma~\ref{4.28:lem-refined-count}, see \eqref{4.28:eq-refined-count-a}, the cubes of
$\mathcal{C}^1\setminus\mathcal{C}$ on $\bar{c}$ number at most $2^d(N_0-1)(100L\check{R})^d$.
Each summand satisfies $\vartheta^{\natural}(\square)\le\vartheta^{\natural}$. So the second part is
at most $2^d(N_0-1)(100L\check{R})^d\vartheta^{\natural}$. Adding the two parts gives
\eqref{4.28:eq-near-pair-sums-a}.

\emph{Clause} (d). Fix an anchor cube $q_0\in\mathcal{Q}$, or a cube of an animal. We count the cubes
$\square\in\mathcal{C}^1\setminus\mathcal{C}$ whose hull touches it. Such a cube belongs to the
family of one $\mathfrak{c}$-shell $\{\mathfrak{c}=i\}$ with $k_0\le i\le j-1$, and has scale $i$,
since $\mathcal{C}^1$ differs from $\mathcal{C}$ only on these shells, as in the first count of
Lemma~\ref{4.28:lem-refined-count}.

The hull of $\square$ exceeds $\square$ by two pieces. The first is five sides of $\pi_i$, and
$5M\ell_i\le 5L^{-2}M\ell_{j+1}$. The second is one cube of the local scale, of side at most
$M\ell_{j+1}$. Since $L^2>5$,
\begin{equation*}
5M\ell_i+M\ell_{j+1}\le 5L^{-2}M\ell_{j+1}+M\ell_{j+1}<2M\ell_{j+1}.
\end{equation*}
So the hull exceeds $\square$ by fewer than two sides of $\pi_{j+1}$. Therefore, if the hull touches
$q_0$, then $\square$ touches one of the $5^d$ cubes of $\pi_{j+1}$ that form the box of side five
around $q_0$.

Each $\pi_{j+1}$-cube touched by such a $\square$ lies in $\{\mathfrak{c}=k_0\}$ or in one old
shell, since $\square$ lies within two sides of $\pi_i$ of the set $\{\mathfrak{c}<\ell_{\mathfrak{B}}\}$
and the boundaries of the shells run along cubes of $\pi_k$, $k$ being the level of the step at which
the current preparation is made. So the second count of
Lemma~\ref{4.28:lem-refined-count}, see
\eqref{4.28:eq-refined-count-a}, applies to it: each of these $5^d$ cubes is touched by at most
$\mathsf{l}_j^2$ cubes of $\mathcal{C}^1$ of the
families $m\ge k_0$. Every cube of $\mathcal{C}^1\setminus\mathcal{C}$ is among them. This gives at
most $5^d\mathsf{l}_j^2$ groups, that is, summation groups of the multiplicity clause, one for each
such cube. By clause (c) every group has $s(\square)=0$, so the factor $L_{\circ}^{6s(\square)}$ of
the multiplicity clause equals $1$, and its majorant
$\mathsf{l}_j^2$ as stated in that clause pays for the factor $\mathsf{l}_j^2$ of the count; the
total is the one stated in the multiplicity clause times $5^d$. The
remaining factor satisfies $5^d\le 2\cdot 52^d$, and $2\cdot 52^d$ is a factor of $K_{\flat}'$. This
factor of $K_{\flat}'$ is the number used by the near-term bound of clause (c) of the sum lemma. It
plays no role in the multiplicity clause, and it is free once that bound is replaced by
\eqref{4.28:eq-near-pair-sums-a}.

The cubes of $\mathcal{C}$, including those of scale $J$, keep the multiplicity stated in the
multiplicity clause. Alternatively, one may count all cubes of $\mathcal{C}^1$ of the families
$m\ge k_0$, the cubes of $\mathcal{C}$ of scale up to $J$ included; the cubes of scale below $k_0$
keep the multiplicity stated in the multiplicity clause. The hull of a
cube of scale $J$ exceeds it by six sides of $\pi_{j+1}$, so the box around $q_0$ has side
$2\cdot 6+1=13$, and $5^d$ is replaced by $13^d$. At $d=4$,
\begin{equation*}
13^4=28\,561\le 2\cdot 52^4 ,
\end{equation*}
so $13^d\le 2\cdot 52^d$, and $2\cdot 52^d$ is again a factor of $K_{\flat}'$.

\emph{Clause} (e). The budget clause is used with its third bracketed summand omitted.

$(\ast)$ In its second member the coefficient $4^d$ becomes $4^d+2^{d+1}$. The additional $2^{d+1}$
is the anchor sum of the far terms over $\mathcal{C}^1$, placed beside the anchor sum of the
anchoring clause it refines. The bound $4^d+2^{d+1}\le 5^d$ is the arithmetic
$256+32=288\le 625$ at $d=4$.

The volume inequality of the budget clause is used with the coefficient $5^d$, in the form
\eqref{4.28:eq-near-pair-sums-1}. Its right-hand side $(5^d/C_V)N\bar{V}$ is absorbed in the budget
$K'N\mathsf{r}\bar{V}$, the factor $5^d/C_V$ being absorbed in $K'\mathsf{r}$, where $K'$ no longer
carries the factors $c_L$, $L_{\circ}^6$ and $2\cdot 52^d$.
The degrees $6r$ and $4r$, the floors and the condition $(Q^{+})^{\natural,N}$ of the budget clause
enter unchanged.

\emph{Clause} (f). Let $\square_0\in\mathcal{C}$. By the third count of
Lemma~\ref{4.28:lem-refined-count}, see \eqref{4.28:eq-refined-count-a}, $\tilde{\square}_0$ is
covered by at most $\mathsf{l}_j^2$ cubes of $\mathcal{C}^1$. So the groups of $\mathcal{C}^1$
summed under one cube of the stage number at most $\mathsf{l}_j^2$. This is the factor
$\mathsf{l}_j^2$ in the total of the bounds (S3) asserted by the overlap statement.

Consequently the three recursive bounds whose proofs use the overlap
statement keep their form, with the overlap factor in their bracketed majorant equal to
$\mathsf{l}_j^2$.
\end{proof}

\begin{remark}[Thin-zone crossing for refined hulls and coarse blocks]\label{4.28:rem-thin-zone-crossing}
Let $\sigma_0$ be the family of cells of the thin zones. Assume the following two statements.
\begin{enumerate}
\item[(i)] \emph{The thin-zone crossing requirement.} A block $\square$ meets the thin-zone
crossing requirement if a face of its enlargement $\tilde{\square}^4$ touches
\begin{equation*}
\bigl(10L^{N_0-2}\bigr)^{d-1}
\end{equation*}
cells of $\sigma_0$, so that $\tilde{\square}^4$ crosses all thin zones; the requirement asks
for such a block.
\item[(ii)] \emph{The crossing property of the cover $\mathcal{C}$.} Every block of
Ba{\l}aban's cover $\mathcal{C}$ of the region meets the requirement of~(i).
\end{enumerate}
Then for every $\square\in\mathcal{C}^1$ the hull $[\tilde{\square}^5]_{\sharp}$ touches
at most
\begin{equation}\label{4.28:eq-thin-zone-crossing-1}
16^d(L+2)^d
\end{equation}
cells of $\sigma_0$, a number depending on $L$ only; the block asked for in~(i) is
accordingly a block of the cover $\mathcal{C}$, for which requirement~(i) holds by~(ii).
\end{remark}

\begin{proof}
Let $\square\in\mathcal{C}^1$. By clause (a$'$) of the box lemma at lifted points
(Lemma~\ref{4.28:lem-box-lemma}, \eqref{4.28:eq-box-lemma-c}), the level function $\ell$,
and with it $\ell_{\mathfrak{B}}$, takes at most two values on
$\square_0=\tilde{\square}^5$, and these values are adjacent. Hence the hull
$[\tilde{\square}^5]_{\sharp}$ consists of cubes of at most two adjacent levels. Each
stratum (level set) of $\ell$ is at least $L$ of its own cubes thick. Consequently the hull
$[\tilde{\square}^5]_{\sharp}$ touches at most the number \eqref{4.28:eq-thin-zone-crossing-1}
of cells of $\sigma_0$. This bound depends on $L$ and the dimension $d$ only, and not on
$N_0$, in contrast with the count in requirement~(i); the block asked for in
requirement~(i) is therefore taken from the cover $\mathcal{C}$, whose blocks meet that
requirement by hypothesis~(ii).
\end{proof}

\begin{proposition}[The core-freezing step]\label{4.28:prop-core-freezing}
Let $k$ be a site of the operation $\mathbf{R}$ and $l$ a stage of the preparation at $k$ (the
\emph{run} at $k$). Let
$\mathcal{C}^1$ be the cover of Definition~\ref{4.28:def-refined-cover}, and for
$\square\in\mathcal{C}^1$ of scale $i$ let $\square_0=\tilde{\square}^5$, the cube $\square$
enlarged by five layers of cubes of its own family $\pi_i$, be its box. Let $\mathfrak{L}_l$ be the
attempt space at stage $l$ and $\mathfrak{D}^{\flat}_l(Y)$ its domain of configurations on a
region $Y$, with its retained triples $(\Sigma^{\ast})$ and its composite clauses. A
\emph{lifted point} of
$Z^{\mathrm{on}}$ is a point of $Z^{\mathrm{on}}$ at which the weight $\omega$ exceeds one. Near
pairs $(T,\square)$ are those of Proposition~\ref{4.28:prop-near-pair-box}. We say that \emph{the
near-pair hypothesis at
stage $l$ holds at} $(T,\square)$ for a configuration $(\mathbb{U},\mathbb{J})$ if the input
hypotheses of \cite[Lemma~4]{B-CMP109} hold on $\square_0$ in the sense of clause (H) of
\eqref{4.28:eq-near-pair-box-b}, and every member of family (c), with $|B'|<\alpha_3(T)$ at
$\operatorname{room}^{\dagger}_T$ as in clause (M), and both configurations (c$'$) of
Proposition~\ref{4.28:prop-near-pair-box} lie in $\mathfrak{U}_T$ in the sense of clause (M) of
\eqref{4.28:eq-near-pair-box-b}.

Assume the following.
\begin{enumerate}
\item[(i)] The one-sided conditions $R_{\min}\ge 5L$, $L\ge 12$ and $M\ge 3L$ hold; $2\le L_0$
and $L_0^2\le L/3$, where $L_0$ is the base of the weight function $\omega$ of
\cite[(1.64)--(1.69)]{B-CMP122a} (not the constant $L_0(N)$ fixed in
Notation~\ref{0.2:not-glossary}); the decay
rate $\delta$ and the cost $w$ per crossed $M$-cube of Lemma~\ref{4.28:lem-old-region-distance}
satisfy $\delta w\ge 16\ln L+4$; and for every attempt datum
$\mathfrak{d}=(C;\kappa,\mathsf{h},\kappa_0,J)$, in the sense of the definition of attempt data in
this section, $(N_0+1)(N_0+3)\le N$, where $N_0=N_0(\kappa)=\kappa-\kappa_0$, an integer whose
existence is stated in \cite{B-CMP122a}, and $N$ are as in that condition of
Lemma~\ref{4.28:lem-box-lemma}.
Moreover $\beta_0=1/7$; every remaining hypothesis of Lemma~\ref{4.28:lem-box-lemma} holds; and
the ceiling $g_n^2\beta'\le 1$ on the running coupling $g_n$ holds at the levels $n$ at which the
statement on the strata and widths of the cutting level imposes it, $\beta'$ being the constant of
that statement. Here $R_{\min}$ and
$R_n$ are the radii of \cite[(2.6)--(2.9)]{B-CMP119}. These conditions, with the near-row
constant $K^{\mathrm{nr}}=66B_{\sharp}K_{\mathrm{cv}}$, belong to the list of conditions under
which Proposition~\ref{4.28:prop-near-pair-box} is stated; that list is assumed.
\item[(ii)] The statements on arrested pieces hold, among them that a completed operation
$\mathbf{R}$ drops its pieces.
\item[(iii)] The box propositions on $\mathfrak{T}$, which give at every near pair the clauses of
$T$ imposed on $\mathfrak{T}$, for data in the run and for data not in it, hold.
\item[(iv)] The statement giving the clauses on $Z^{\mathrm{on}}$ of a term $T$ at boxes
that meet no lifted point holds.
\item[(v)] The sum lemma over near pairs used in the inheritance argument holds, together with
its clause for data not in the run (its clause at lifted points is the subject of (vii)).
\item[(vi)] The near pairs at which the near-pair step of the inheritance theorem imposes its
hypothesis are the near pairs defined above, and for the region $Y$ on which that step is taken
their boxes satisfy $[\square_0]_{\sharp}\subset Y$; and $R_n\le\check{R}$ for every level $n$
with $k_0<n\le j$ ($k_0$, $j$ as in Lemma~\ref{4.28:lem-refined-count}), where
$\check{R}$ is the window envelope of the radii.
\item[(vii)] The bounds of clauses (c)--(f) of Lemma~\ref{4.28:lem-near-pair-sums}, the first of
which is displayed in \eqref{4.28:eq-near-pair-sums-a}, hold, each as stated there; clause (e) of
that lemma is stated there with its proof incomplete at the step marked $(\ast)$.
\end{enumerate}
Then, for every region $Y$ and every $(\mathbb{U},\mathbb{J})\in\mathfrak{D}^{\flat}_l(Y)$, the
near-pair hypothesis at stage $l$ holds at every near pair $(T,\square)$ with
$[\square_0]_{\sharp}\subset Y$ whose box $\square_0$ meets a lifted point of $Z^{\mathrm{on}}$;
and the two anchored sums over near pairs taken in the inheritance argument, restricted to such
pairs, satisfy the bounds of clauses (c)--(f) of Lemma~\ref{4.28:lem-near-pair-sums}.
\end{proposition}

\begin{proof}
The near-pair step of the inheritance theorem requires the near-pair hypothesis at stage $l$ at
every near pair on $\mathcal{C}^1$, and the two anchored sums of the inheritance argument over all
such pairs; by (vi) these are the near pairs defined above. We sort the clauses of a term $T$
according to where they are imposed and according to whether the box of the pair meets a lifted
point.

\emph{Clauses on $\mathfrak{T}$.} At every near pair $(T,\square)$ the clauses of $T$ on
$\mathfrak{T}$ are provided by the box propositions on $\mathfrak{T}$ of (iii), for data in the run
and for data not in it.

\emph{Pairs whose box meets a lifted point.} Let $Y$ be a region,
$(\mathbb{U},\mathbb{J})\in\mathfrak{D}^{\flat}_l(Y)$, and $(T,\square)$ a near pair with
$[\square_0]_{\sharp}\subset Y$ and $\square_0$
meeting a lifted point of $Z^{\mathrm{on}}$. Proposition~\ref{4.28:prop-near-pair-box} applies to
this pair at stage $l$; the conditions it uses are
\begin{equation*}
\begin{gathered}
R_{\min}\ge 5L,\quad L\ge 12,\quad L_0^2\le L/3,\quad \beta_0=1/7,\\
\delta w\ge 16\ln L+4,\quad M\ge 3L,\quad (N_0+1)(N_0+3)\le N,
\end{gathered}
\end{equation*}
and the remaining hypothesis of Lemma~\ref{4.28:lem-box-lemma}, all of which are listed in (i),
together with the ceiling $g_n^2\beta'\le 1$ on the coupling, which is also in (i); these belong
to the list of conditions of Proposition~\ref{4.28:prop-near-pair-box}, assumed in (i). Clause (H) of
\eqref{4.28:eq-near-pair-box-b} gives the input hypotheses of \cite[Lemma~4]{B-CMP109} on
$\square_0$, in the unit of the $\mathfrak{c}$-shell, at factor at most one on $Z^{\mathrm{on}}$
and less than $1+2\beta$ on $\mathfrak{T}$; clause (M) of \eqref{4.28:eq-near-pair-box-b} gives,
for every member $P(\varsigma;t,\tau;B')$ of family (c), with $|B'|<\alpha_3(T)$ at
$\operatorname{room}^{\dagger}_T$ as in clause (M), and for both configurations (c$'$), the
inequalities $Q_a(\cdot)(y)<\theta_{T,a}(y)$ at every pair $(y,a)$ of the index set of clause
(M), which define membership in $\mathfrak{U}_T$. Proposition~\ref{4.28:prop-near-pair-box} gives
these for every clause of $T$ at such a pair, those imposed on
$Z^{\mathrm{on}}$ and those imposed on $\mathfrak{T}$ alike. Hence the near-pair hypothesis at
stage $l$ holds at $(T,\square)$, which is the first assertion.

\emph{Clauses on $Z^{\mathrm{on}}$ at boxes meeting no lifted point.} These are provided by (iv).

Combining the three cases, and using (vi) for the inclusion $[\square_0]_{\sharp}\subset Y$ at
the pairs whose box meets a lifted point, the near-pair step of the inheritance theorem has
the near-pair hypothesis at stage $l$ at every near pair on $\mathcal{C}^1$.

\emph{The sums.} The near pairs over which the two anchored sums of the inheritance argument run
are divided into those treated by the sum lemma of (v), those treated by its clause for
data not in the run,
and those whose box meets a lifted point of $Z^{\mathrm{on}}$. For the last class the bounds are
those of clauses (c)--(f) of Lemma~\ref{4.28:lem-near-pair-sums}, assumed in (vii); clause (e) of
that lemma is stated with its proof incomplete at the step marked $(\ast)$, and each of these
bounds is used as stated there. The count of refined cubes on which these bounds rest, the count in
Lemma~\ref{4.28:lem-refined-count} of the cubes of $\mathcal{C}^1\setminus\mathcal{C}$ on a
component $\bar{c}$ of $Z^{[l]}$, is
taken with $R_n\le\check{R}$ for $k_0<n\le j$, as in (vi). The third class alone gives the second
assertion; adding the three classes gives the two anchored sums over all near pairs required by
the near-pair step.

\emph{Induction on the site.} The inheritance theorem is proved by induction on the site of the
operation $\mathbf{R}$. Proposition~\ref{4.28:prop-near-pair-box} at the site $k$ uses, of the
earlier sites, only the statements (ii) on arrested pieces, and never clause (a) of the
inheritance theorem at an earlier site; so the argument above may be used inside that induction
without circularity. Applied at each earlier site $k_b<k$, to the thin layers of that site's own
run, the same argument supplies the near-pair hypothesis, at the stages of that run, at the near
pairs there whose
box meets a lifted point, which is the hypothesis that clause (a) of the inheritance theorem at
$k_b$ uses in the corresponding step of its proof.

\emph{Comparison of the covers $\mathcal{C}^1$ and $\mathcal{C}$.} Let $\mathcal{C}$ be the cover
given by the construction of Definition~\ref{4.28:def-refined-cover} on $\mathfrak{B}$ in place of
$\mathfrak{B}^{\mathfrak{c}}$. Passing from $\mathcal{C}$ to $\mathcal{C}^1$ changes only the cover
and box structure of the terms of $\mathcal{T}_1$ on the lifted sets of $Z^{\mathrm{on}}$
(Definition~\ref{4.28:def-refined-cover}) and the room $\operatorname{room}^{\dagger}_T$, and
hence the grouping of near pairs. The following are the same for $\mathcal{C}^1$ and for
$\mathcal{C}$, so that the statements appealed to above apply in the form in which they are
stated.
\begin{itemize}
\item The form of $\mathfrak{D}^{\flat}_l$. On lifted sets its retained triples
$(\Sigma^{\ast})$ are those already named: for data of kinds (B) and (B$'$), in the sense of the
definition of attempt data in this section, the retained stage
structures of the present run, and for data of kind (A) (an older arrested piece $\mathfrak{a}$)
the retained triples of $\mathfrak{a}$.
\item The composite clauses, which are those of $\mathfrak{L}_l$. Every real point meets them
with margin one half. Here $U2$ and $U3$ denote the frozen weighted units of the attempt space;
the composite entries at real points satisfy
\begin{equation*}
Q_e\le \tfrac12\,\varphi_{\mathrm{arr}}\cdot U2,
\end{equation*}
and the right side equals $0.049$ times the clause radius when $\lambda_{\mathfrak{s}}=m$ and is
at most $0.0016$ times the birth radius when $\lambda_{\mathfrak{s}}<m$, because
$U2\le U3/31.5$; here $m$ is the level against which the block level $\lambda_{\mathfrak{s}}$ is
compared in that bound. This bound and the constants entering it are those recorded in the remark
on the margin at boxes which follows this proof in the present section, and they are used as
stated there.
\item Every threshold, the layer costs $\iota_l$ and $\iota^1_l$, and every radius;
$\mathsf{w}_0=1/24$; the condition \eqref{4.28:eq-near-pair-box-a} of
Proposition~\ref{4.28:prop-near-pair-box}, the remaining hypothesis of
Lemma~\ref{4.28:lem-box-lemma} listed in (i), and the further conditions accompanying them in
Proposition~\ref{4.28:prop-near-pair-box}.
\item The conditions under which Proposition~\ref{4.28:prop-near-pair-box} is stated; of these the
argument above uses only those listed in (i).
\item The constants $C_F$, $\bar{V}$ and $K_{\flat}'$, the last up to two absolute constants
depending only on $d$; and $24K'$.
\item Clauses (a)--(c) of the inheritance theorem, its two corollaries, and the auxiliary claim
proved with them.
\end{itemize}
\end{proof}

\begin{remark}[Margin one half at the box triple, with constants]\label{4.28:rem-box-margin}
Let $y'$ be a point of the lifted set $\mathsf{L}_{\mathfrak{d}}$. At $y'$ write $\ell$ for the level, $\omega$ for the
weight, $\mathfrak{c}$ for the cutting level and $g=\ell-\mathfrak{c}$ for the depth (here $g$ is not the
reference coupling). Let $e$ be a composite
entry, read at a retained triple $(p,X',\mathfrak{s})$ with block level $\lambda_{\mathfrak{s}}=\lambda_{\mathfrak{s}}(y')$. Write $Q_e$ for the
corresponding entry of the smallness vector at $y'$. Assume that clause (3) of the lemma bounding the
composite entries at real points by the frozen weighted unit $U2$ holds:
\begin{equation}\label{4.28:eq-box-margin-1}
Q_e\le\tfrac12\,\varphi_{\mathrm{arr}}\,\omega\,r_e(\ell,p\wedge\mathfrak{c})=\tfrac12\,\varphi_{\mathrm{arr}}\,U2 .
\end{equation}
Then the following hold.
\begin{itemize}
\item[(i)] If $\lambda_{\mathfrak{s}}$ is a block level at which the frozen weighted units $U2$ and $U3$ coincide,
$U3$ being the clause's own radius, then $Q_e\le 0.049\,U3$, that is, $0.049$ times the clause's own
radius.
\item[(ii)] If $g\ge1$, then at the box triple $(p,\square_0,\mathfrak{B}^{\mathfrak{c}})$, where $\lambda_{\mathfrak{s}}=\mathfrak{c}$ and
$w_{\mathfrak{s}}=1$, we have
\begin{equation}\label{4.28:eq-box-margin-3}
Q_e\le 0.0016\,\mathbf{r}^{\mathfrak{c},(p)}_e
\qquad\text{and}\qquad
Q_e\le\tfrac12\,\theta^{(l)}(\mathfrak{c})\,\mathbf{r}^{\mathfrak{c},(p)}_e .
\end{equation}
\end{itemize}
Thus the composite entries have margin one half, in the sense of the attempt space $\mathfrak{L}_l$, at every
depth.
\end{remark}

\begin{proof}
Suppose first that $\lambda_{\mathfrak{s}}$ is a block level at which $U2=U3$, $U3$ being the clause's own radius.
Then \eqref{4.28:eq-box-margin-1} and the bound $\varphi_{\mathrm{arr}}\le 0.098$ of the arrest theorem
give
\[
Q_e\le\tfrac12\,\varphi_{\mathrm{arr}}\,U3\le\tfrac12\cdot0.098\,U3=0.049\,U3 .
\]
This is (i). Since the running factor of the clause is at least $\varphi_{\mathrm{arr}}$, the middle term
$\tfrac12\varphi_{\mathrm{arr}}U3$ is also at most one half of the clause's threshold.

Now suppose that $g\ge1$. At depth $g\ge1$ the frozen weighted unit $U2$ satisfies
\begin{equation}\label{4.28:eq-box-margin-2}
U2\le\Bigl(\frac{8}{21L}\Bigr)^{g}\,\mathbf{r}^{\mathfrak{c},(p)}_e .
\end{equation}
Under the standing condition $L\ge 12$ of the construction,
\[
\frac{8}{21L}\le\frac{8}{21\cdot12}=\frac{2}{63}=\frac{1}{31.5}<1 .
\]
Since $g\ge1$, this gives $(8/(21L))^{g}\le 1/31.5$. Hence, by \eqref{4.28:eq-box-margin-2},
\[
U2\le\frac{\mathbf{r}^{\mathfrak{c},(p)}_e}{31.5}.
\]
Inserting this bound into \eqref{4.28:eq-box-margin-1} and using $\varphi_{\mathrm{arr}}\le 0.098$, we obtain
\[
Q_e\le\frac{\tfrac12\,\varphi_{\mathrm{arr}}}{31.5}\,\mathbf{r}^{\mathfrak{c},(p)}_e
\le\frac{\tfrac12\cdot0.098}{31.5}\,\mathbf{r}^{\mathfrak{c},(p)}_e\le 0.0016\,\mathbf{r}^{\mathfrak{c},(p)}_e .
\]
This is the first assertion of \eqref{4.28:eq-box-margin-3}.

By Lemma~\ref{4.28:lem-lifted-thresholds}, the running factor at the cutting level satisfies
$\theta^{(l)}(\mathfrak{c})\ge\varphi_{\mathrm{arr}}$ on $\mathsf{L}_{\mathfrak{d}}$. Together with $1/31.5<1$, this gives
\[
\frac{\tfrac12\,\varphi_{\mathrm{arr}}}{31.5}\le\tfrac12\,\theta^{(l)}(\mathfrak{c}).
\]
Hence $Q_e\le\tfrac12\theta^{(l)}(\mathfrak{c})\mathbf{r}^{\mathfrak{c},(p)}_e$, which is the second assertion of
\eqref{4.28:eq-box-margin-3}.

The right-hand side of the second assertion of \eqref{4.28:eq-box-margin-3} is one half of the threshold
$\theta^{(l)}(\mathfrak{c})\mathbf{r}^{\mathfrak{c},(p)}_e$ of the composite clause at the box triple
(see \eqref{4.28:eq-lifted-thresholds-a}); this is the asserted margin one half.
\end{proof}

\begin{definition}[Scales, cube partitions and the native label sequence]\label{4.29:not-scales-labels}
The notation below serves the $(k+1)$-st step with the operation $\mathbf{T}$ in the induction
of \cite[Theorem~1]{B-CMP119}, read as an induction on Ba{\l}aban's families
$\mathbf{E}^{(k)}$, $\mathbf{R}^{(k)}$, $\mathbf{B}^{(k)}$. That step must produce Ba{\l}aban's families
$\mathbf{E}^{(k+1)}$, $\mathbf{R}^{(k+1)}$, $\mathbf{B}^{(k+1)}$ of \cite{B-CMP119}, with
$\mathbf{E}^{(k+1)}+\mathbf{R}^{(k+1)}+\mathbf{B}^{(k+1)}$ satisfying
\cite[(2.26)--(2.29), (2.31), (2.41)--(2.42)]{B-CMP119} on every label (the labels at the
$\mathbf{T}$-stage are introduced below). The new $\mathbf{B}$-pieces must be analytic and bounded on the domain on which
that theorem requires them. The place where this is decided is the first localisation
\cite[(3.4)]{B-CMP109} of the $\mathbf{E}_k$- and $\mathbf{R}_k$-families in three actions:
\begin{itemize}
\item the action of \cite[(3.27)]{B-CMP119}, in which the $k$-th action $A_k$ of
\cite{B-CMP119} enters at full strength;
\item the action of \cite[(3.29)]{B-CMP119}, in which $A_k$ enters multiplied by the
interpolation parameter;
\item the last logarithm of \cite[(3.28)]{B-CMP119}, which contains only the term denoted $A$
in that display.
\end{itemize}
The term $V^{(k)}(\mathsf{S}_{k+1})$ of \cite{B-CMP119}, where $\mathsf{S}_{k+1}$ denotes the
set written $S_{k+1}$ in \cite[(3.27)]{B-CMP119} (the letter $\mathsf{S}$ distinguishes it from
the boundary strips $S_j$ defined below), belongs to the action of \cite[(3.27)]{B-CMP119}
(the radius $R_k$ is defined below). By
\cite[\S3]{B-CMP119}, it arises as the sum of the expressions discussed there, expanded with
respect to the perturbative terms in the transformations \cite[(1.22)]{B-CMP119}. It therefore
has the same support as that sum, and its size differs from that of the sum by a factor
$1+O(e^{-R_k})$. As stated in
\cite[\S3]{B-CMP119}, the integrands of \cite[(3.26), (3.28)]{B-CMP119} are expanded and
localised in the same way as the corresponding expressions in the actions. The same place is
where the analysis of \cite[\S\S3--4]{B-CMP109} has to run at each cover cube (defined below);
\cite[p.~276]{B-CMP119} states that the analysis of \cite[\S3]{B-CMP109} is repeated using the
cover cubes and their partition of unity.

\emph{Scales and cubes.} We set $\eta=L^{-k}$ and $\ell_n:=L^n\eta$. For each $n$, let
$\pi_n$ denote the partition of the sites of spacing $\eta$ into cubes of side $M\ell_n$, which
we call $M$-cubes of scale $n$. All these partitions are compatible. For a site $x$ we write
$\Delta_s(x)\in\pi_s$ for the cube containing $x$. For $\Delta\in\pi_s$ and $r\ge 0$, the set
$\tilde\Delta^{r}$ is $\Delta$ together with $r$ layers of $\pi_s$-cubes around it: the full
cube of side $(1+2r)M\ell_s$ concentric with $\Delta$, corners included \cite{B-CMP109}. A
\emph{cover cube} of scale $s$ is a union of $2^d$ neighbouring $\pi_s$-cubes
\cite[\S3]{B-CMP109}, and $\tilde\Box^{r}$ is defined for a cover cube $\Box$ in the same way.
We write $|\cdot|_\infty$ for the sup-norm distance between sites of spacing $\eta$, taken with
the periodic identifications of the lattice, and
$\operatorname{dist}_\infty$ for the distance between sets induced by it. With
this notation,
\begin{equation}\label{4.29:eq-scales-labels-a}
x\in\Delta\in\pi_s,\quad y\in\tilde\Delta^{6}\quad\Longrightarrow\quad |y-x|_\infty<7M\ell_s .
\end{equation}

\emph{Labels at the $\mathbf{T}$-stage.} The \emph{current sequence} is $\{\Omega_i\}_{i\le k+1}$.
We put $\Gamma_a:=\Omega_a\setminus\Omega_{a+1}$, and we let $n(x)$ be the index $a$ with
$x\in\Gamma_a$. The \emph{native sequence} $\{\Omega^{\natural}_i\}$ is the sequence of regions
made by the operations $\mathbf{T}$, with those components deleted that have been healed, that
is, whose large-field factor a later step has absorbed back into the small-field format. The
\emph{native label} $l(x):=n^{\natural}(x)$ is the index of $x$ with respect to the native
sequence: $l(x)=a$ if and only if $x\in\Omega^{\natural}_a\setminus\Omega^{\natural}_{a+1}$.

\emph{Action domains.} The action domains are $\{\Lambda_j\}_{j\le k}$, and $\Lambda:=\Lambda_k$.
For $r\ge0$, $\Lambda^{(-r)}$ is the union of the $\pi_k$-cubes $\Delta$ with
$\tilde\Delta^{r}\subset\Lambda$. The boundary strip of $\Lambda_j$ is
\begin{equation*}
S_j:=\bigl\{x\in\Lambda_j:\ \operatorname{dist}_\infty(x,\Lambda_j^{c})\le(MR_j+1)\ell_j\bigr\},
\end{equation*}
with $R_j$ as fixed at the end of this definition. Throughout, the
standing hypotheses of the
$\mathbf{T}$-stage are in force: the nesting of the action domains $\Lambda_i$ in the native
regions $\Omega^{\natural}_i$ and of these in the current regions $\Omega_i$; the cube structure
of these regions ($\Lambda_i$ is a union of cubes of side $MR_i\ell_i$, and $\Omega_i$,
$\Omega^{\natural}_i$ are unions of $\pi_i$-cubes); the separation
conditions among these regions; the arrest rule relating the labels $n(x)$ and $l(x)$; and the
regularity of the configuration $U_{k+1}$.

\emph{Constants and the flow.} The decay exponent of \cite{B-CMP119} used below is
$\kappa'=(1+4\beta)\kappa$, where $\kappa$ is Ba{\l}aban's auxiliary parameter of the glossary
(Definition~\ref{0.2:not-glossary}) and $\beta$ is the schedule constant of
Definition~\ref{4.4:def-post-step-space}. We write $x_j:=g_j^{-2}$. The two-sided flow
bound
\begin{equation}\label{4.29:eq-scales-labels-b}
\lambda_\sharp(j-i)-c_5\ \le\ x_i-x_j\ \le\ B_\sharp(j-i),\qquad i\le j,
\end{equation}
is in force; here $\lambda_\sharp$, $B_\sharp$ and $c_5$ are the flow constants of
Definition~\ref{1.2:def-flow-constants}, and the bound is the coupling-flow bound of clause~(0)
of Theorem~0 (Theorem~\ref{4.1:thm-clause-zero}). Further, $R_j:=\check R_j$, where
$\check R_j$ denotes the level-$j$ radius, so that under the standing hypotheses above
$\Lambda_j$ is a union of cubes of side $MR_j\ell_j$; the
quantities $\varepsilon_j$, $\alpha_{0,j}$, $\alpha_{1,j}$ are those of
\cite[(2.4), (2.28)]{B-CMP119}; and $q:=\min\{q_0,q_1\}$, with $q_0,q_1$ the exponents of
\cite{B-CMP119}.
\end{definition}

\begin{proof}[Proof of \eqref{4.29:eq-scales-labels-a}]
Let $x\in\Delta\in\pi_s$ and $y\in\tilde\Delta^{6}$. Fix a lift of $\Delta$ to a cube of
$(\eta\mathbb{Z})^d$. The cube of side $13M\ell_s$ concentric with this lift, made of the lift
and six layers of cubes of side $M\ell_s$ around it, projects onto $\tilde\Delta^{6}$. Hence
$x$ has a lift in the lift of $\Delta$, and $y$ has a lift in that cube. The periodic distance of
two sites is at most the sup-norm distance in $(\eta\mathbb{Z})^d$ of any two of their lifts, so
it suffices to bound the latter for these lifts.

The cube $\Delta$ belongs to a partition of the sites of spacing $\eta$ into cubes of side
$M\ell_s$. Hence in each coordinate direction $\mu$ it contains exactly $M\ell_s/\eta=ML^{s}$
sites. These have coordinates $a_\mu,\ a_\mu+\eta,\ \dots,\ a_\mu+M\ell_s-\eta$ for some
$a_\mu$. The set $\tilde\Delta^{6}$ is $\Delta$ together with six layers of $\pi_s$-cubes of side
$M\ell_s$ on every side. Therefore the $\mu$-coordinates of its sites lie between
$a_\mu-6M\ell_s$ and $a_\mu+7M\ell_s-\eta$.

Taking the largest admissible $y_\mu$ and the smallest admissible $x_\mu$ gives
\begin{equation*}
y_\mu-x_\mu\le (a_\mu+7M\ell_s-\eta)-a_\mu=7M\ell_s-\eta .
\end{equation*}
Taking the largest admissible $x_\mu$ and the smallest admissible $y_\mu$ gives
\begin{equation*}
x_\mu-y_\mu\le (a_\mu+M\ell_s-\eta)-(a_\mu-6M\ell_s)=7M\ell_s-\eta .
\end{equation*}
Hence $|y_\mu-x_\mu|\le 7M\ell_s-\eta$ for every $\mu$. Since $\eta>0$, this gives
$|y-x|_\infty\le 7M\ell_s-\eta<7M\ell_s$.
\end{proof}

\begin{definition}[Good cubes, the native scale, tiles and the cover]\label{4.29:def-tiles-cover}
The following objects are those of Definition~\ref{4.29:not-scales-labels}:
\begin{itemize}
\item the scales $\ell_n=L^n\eta$ with $\eta=L^{-k}$, and the partitions $\pi_n$ into $M$-cubes of scale $n$;
\item the cube $\Delta_s(x)\in\pi_s$ containing the site $x$, and the enlargements $\tilde\Delta^r$ and
$\tilde\square^r$;
\item the cover cubes of scale $s$, the distance $\operatorname{dist}_\infty$, and the native sequence
$\{\Omega_i^\natural\}$ beside the current sequence $\{\Omega_i\}$;
\item the action domains $\{\Lambda_j\}_{j\le k}$, $\Lambda=\Lambda_k$, the sets $\Lambda^{(-r)}$, and the numbers $R_j$.
\end{itemize}
All the partitions $\pi_n$ are compatible: for $s'\le s$ every cube of $\pi_s$ is a union of cubes of $\pi_{s'}$.

\smallskip
\noindent\emph{Good cubes and the native scale.}
A cube $\Delta\in\pi_s$, $1\le s\le k$, is \emph{good} if $\tilde\Delta^6\subset\Lambda_s$. The
\emph{native scale} of a site $x$ is
\begin{equation}\label{4.29:eq-tiles-cover-1}
\mathfrak{s}(x):=\max\{s\le k:\ \Delta_s(x)\ \text{is good}\},
\end{equation}
and $\mathfrak{s}(x):=0$ if no such $s$ exists.

\smallskip
\noindent\emph{The enlarged domains.}
For a generation $j\le k$ let
\begin{equation}\label{4.29:eq-tiles-cover-2}
\Lambda_j^{\sim}:=\Lambda_j\cup\bigcup\bigl\{\Delta\in\pi_j:\ \tilde\Delta^1\cap\Lambda_j\neq\emptyset\bigr\},
\end{equation}
that is, $\Lambda_j$ together with one layer of $\pi_j$-cubes. We call $\Lambda_j^{\sim}\setminus\Lambda_j$ the
\emph{ring}.

The operation $\sim$ in \eqref{4.29:eq-tiles-cover-2} is taken in $\pi_j$-cubes. Here $\tilde\Delta^n$ is understood
as in \cite{B-CMP109}: for $\Delta\in\pi_j$ it is the full cube of side $(1+2n)M\ell_j$ formed by $\Delta$ and $n$
layers of $\pi_j$-cubes, corners included.

The set $\Lambda_j^{\sim}$ is therefore not the set obtained by the operation $\sim$ that \cite{B-CMP119} applies
to $\Lambda_j$. That operation adjoins a layer of cubes of the family in terms of which the sets are defined,
namely of cubes of side $MR_j\ell_j$; the set written $\Lambda_j^0$ in \cite{B-CMP119} is $\Lambda_j$ with one
such layer removed. Nor is $\Lambda_j^{\sim}$ the set $\Lambda^{\sim}_{k+1}$ of \cite{B-CMP119}, on
which only small-field characteristic functions occur; that set is obtained by adjoining a layer of
$LMR_{k+1}$-cubes.

The set $\Lambda_j^{\sim}$ is a union of $\pi_j$-cubes. Since the partitions are compatible and $M_1$ divides
$M$ \cite{B-CMP119}, it is also a union of cubes of side $M_1\ell_j$.

Moreover
\begin{equation}\label{4.29:eq-tiles-cover-3}
\Lambda_j^{\sim}\subset\Omega_j^\natural\subset\Omega_j .
\end{equation}
To see this, let $x$ be a site of the ring. Then $x\in\Delta$ for some $\Delta\in\pi_j$ with
$\tilde\Delta^1\cap\Lambda_j\neq\emptyset$; pick $z\in\tilde\Delta^1\cap\Lambda_j$. Since $x\in\Delta$ and
$z\in\tilde\Delta^1$, the two sites differ by at most $2M\ell_j$ in each coordinate. Hence
\begin{equation*}
\operatorname{dist}_\infty(x,\Lambda_j)\le 2M\ell_j\le 2MR_j\ell_j,
\end{equation*}
because $R_j\ge1$. The inclusions \eqref{4.29:eq-tiles-cover-3} then follow from the lemma on the native
sequence of domains, which gives that every site within distance $2MR_j\ell_j$ of $\Lambda_j$ lies in
$\Omega_j^\natural$, and that $\Omega_j^\natural\subset\Omega_j$.

For $x\in\Lambda_j^{\sim}$ put
\begin{equation}\label{4.29:eq-tiles-cover-4}
\mathfrak{s}_j(x):=\max\{j,\mathfrak{s}(x)\}.
\end{equation}

\smallskip
\noindent\emph{Tiles.}
The \emph{tiles} of generation $j$ are the cubes
\begin{equation}\label{4.29:eq-tiles-cover-5}
\mathcal{T}_j:=\bigl\{\Delta_{\mathfrak{s}_j(x)}(x):\ x\in\Lambda_j^{\sim}\bigr\}.
\end{equation}
A tile belonging to $\pi_s$ has \emph{scale} $s$. A tile $\Delta$ is \emph{default} if $\mathfrak{s}(y)<j$ for every
$y\in\Delta$.

On the ring every tile is a default $\pi_j$-cube. Indeed, let $x\in\Lambda_j^{\sim}\setminus\Lambda_j$. If
$s=\mathfrak{s}(x)\ge j$, then $\Delta:=\Delta_s(x)$ is good, so
\begin{equation*}
x\in\Delta\subset\tilde\Delta^6\subset\Lambda_s\subset\Lambda_j ,
\end{equation*}
which is impossible. Hence $\mathfrak{s}(x)<j$, $\mathfrak{s}_j(x)=j$, and the tile of $x$ is
$\Delta_j(x)\in\pi_j$.

This tile is default. Suppose $y\in\Delta_j(x)$ had $s=\mathfrak{s}(y)\ge j$. Then the good cube $\Delta_s(y)$ lies
in $\Lambda_s\subset\Lambda_j$. By compatibility it contains $\Delta_j(y)=\Delta_j(x)\ni x$, so again
$x\in\Lambda_j$, which is impossible.

\smallskip
\noindent\emph{The cover.}
For each tile $\Delta\in\mathcal{T}_j$ of scale $s$, take the $2^d$ cover cubes of scale $s$ that contain
$\Delta$; each of them lies in $\tilde\Delta^1$. These cubes are said to be \emph{generated} by $\Delta$. The
\emph{cover} $\mathcal{C}_j$ is the set of all cover cubes generated by the tiles of $\mathcal{T}_j$. A cube
$\square\in\mathcal{C}_j$ is \emph{default} if every tile generating it is default.

\smallskip
\noindent\emph{The partition of unity.}
For $\square\in\mathcal{C}_j$ of scale $s$ with centre $y$, let $\zeta_\square$ be the function of
\cite{B-CMP109} at scale $s$:
\begin{equation}\label{4.29:eq-tiles-cover-6}
\zeta_\square(x)=\prod_{\mu=1}^{d}\zeta\bigl((M\ell_s)^{-1}(x_\mu-y_\mu)\bigr).
\end{equation}
The properties of this standard cut-off are taken as stated in \cite{B-CMP109}:
\begin{itemize}
\item $\zeta(t)=1$ for $|t|\le\tfrac13$, and $\zeta(t)=0$ for $|t|\ge\tfrac23$;
\item the derivatives of $\zeta$ are bounded by $5$, and $0\le\zeta\le1$;
\item $\sum\zeta_\square=1$, the sum running over the complete family of cover cubes of scale $s$.
\end{itemize}
These properties are used in the lemma on the cover below, for the bounds $1\le S$ and
$S\le 1+4\cdot 2^d$ on the set specified there, and for the cost estimates of the normalised functions
$\bar\zeta_\square$.

Put
\begin{equation}\label{4.29:eq-tiles-cover-7}
S:=\sum_{\square\in\mathcal{C}_j}\zeta_\square,\qquad
\bar\zeta_\square:=\frac{\zeta_\square}{S}\quad\text{on }\{S>0\}.
\end{equation}
The set $\{S>0\}$ contains the $\tfrac13M\ell_j$-neighbourhood of $\Lambda_j^{\sim}$; this is proved in the lemma on
the cover that follows this definition.

\smallskip
\noindent\emph{Use in the first localisation.}
The first localisation \cite[(3.4)]{B-CMP109} of the families $\mathbf{E}^{(j)}$ and $\mathbf{R}^{(j)}$ is written
with the pair $(\mathcal{C}_j,\bar\zeta)$, and with the block field $U_j$ specified in what follows. At a cube of
scale $s$ the following conventions apply in those formulas:
\begin{itemize}
\item the operation $\sim$ is taken in $\pi_s$;
\item the enlarged cube is $\square_0:=\tilde\square^5$, called the \emph{window} of $\square$;
\item $\eta$ is replaced by $L^{-s}$.
\end{itemize}

\smallskip
\noindent\emph{Two interior sets.}
Finally,
\begin{equation}\label{4.29:eq-tiles-cover-8}
A:=\Lambda^{(-6)},\qquad P:=\Lambda^{(-12)} .
\end{equation}
\end{definition}

\begin{proposition}[The native-scale cover and its partition of unity]\label{4.29:prop-native-cover}
Use the scales $\ell_n$, the partitions $\pi_n$, the cubes $\Delta_s(x)$, their enlargements
$\tilde\Delta^{r}$ and the label at the $\mathbf{T}$-stage of
Notation~\ref{4.29:not-scales-labels}, and the good cubes, the native scale $\mathfrak{s}$, the
sets $\Lambda_j^{\sim}$, the functions $\mathfrak{s}_j$, the tiles $\mathcal{T}_j$, the cover
$\mathcal{C}_j$, the functions $\zeta_\square$, the sum $S=\sum_{\square\in\mathcal{C}_j}\zeta_\square$,
the quotients $\bar\zeta_\square=\zeta_\square/S$, the cubes $\square_0=\tilde\square^{5}$ and the
sets $A=\Lambda^{(-6)}$, $P=\Lambda^{(-12)}$ of Definition~\ref{4.29:def-tiles-cover}. For a set $E$
and $r>0$ let $B_\infty(E,r)$ denote the set of points whose $\operatorname{dist}_\infty$-distance
to $E$ is at most $r$. We say that a quantity \emph{reads} a datum (one of the sets of the label, or
the block structure $\mathbf{B}'$ of (h$'$)) on a set $D$, or within $r$ of a set $E$, if it depends
on that datum only through its trace on $D$, respectively on $B_\infty(E,r)$. Assume the following.
\begin{enumerate}
\item[(1)] $L\ge 13$; $R_i\ge 7$ for every $i$; $M\ge L^4$; and $M>M_2=L^2M_1$.
\item[(2)] Here and throughout this section $\delta_0>0$ denotes the exponential decay constant of
Ba{\l}aban's propagator estimates (a condition of this form appears in \cite[(2.59)]{B-CMP96}), and
\begin{equation}\label{4.29:eq-native-cover-a}
L^{d-1}e^{-\delta_0M_1/16}\le\tfrac12 .
\end{equation}
\item[(3)] The conclusions of the label lemma at the $\mathbf{T}$-stage: (T0), that
$\Lambda_i\subset\Omega_i^{\natural}\subset\Omega_i$ and $\Omega_{i+1}^{\natural}\subset\Lambda_i$, so
that $\Lambda_{i+1}\subset\Lambda_i$, with $\Lambda_i$ a union of $MR_i$-cubes and $\Omega_i$,
$\Omega_i^{\natural}$ unions of $\pi_i$-cubes; the bounds
\begin{align*}
&\text{(T1)}\quad\operatorname{dist}_\infty(\Omega^{\natural}_{i+1},\Lambda_i^{c})\ge
8MR_{i+1}\ell_{i+1}\ \ (i\le k),\quad\Omega_{k+1}=\Omega^{\natural}_{k+1}\supset\Lambda_{k+1};\\
&\text{(T2)}\quad\operatorname{dist}_\infty(\Lambda_i,(\Omega_i^{\natural})^{c})\ge 2MR_i\ell_i;\\
&\text{(G1)}\quad\operatorname{dist}_\infty(\Omega_{a+1},\Omega_a^{c})\ge M\ell_a;
\end{align*}
(A), which confines the points where the current index $n(x)$ exceeds the native index $l(x)$ to
arrested components; and (REG)$^+$, the regularity of the minimal configuration $U_{k+1}$. And, for
clause (h$''$), conclusion (T1) of the label lemma at the later $\mathbf{T}$-stage that creates
$\Lambda_{k''}$.
\item[(4)] Clause (b) of the lemma on the radii $R_j$, the identification lemma, and the filing
lemma for the $(X,z)$-functions.
\end{enumerate}
Then for every generation $j\le k$ the following hold.
\begin{enumerate}
\item[(a)] At every site $x$ the scales $s$ for which $\Delta_s(x)$ is good form the initial
segment $\{1,\dots,\mathfrak{s}(x)\}$. The function $\mathfrak{s}_j$ is constant on tiles, and
$\mathcal{T}_j$ is a partition of $\Lambda_j^{\sim}$. If $\Delta$ is a tile of scale $s$, then
$\mathfrak{s}_j\in\{s-1,s,s+1\}$ on $\tilde\Delta^{1}\cap\Lambda_j^{\sim}$. Two cubes of
$\mathcal{C}_j$ that overlap have scales differing by at most $2$. On every tile of scale $s$,
\begin{equation*}
\sum_{\square\in\mathcal{C}_j,\ \text{of scale } s}\zeta_\square=1 .
\end{equation*}
One has $1\le S\le 1+4\cdot 2^d$ on $\Lambda_j^{\sim}$, and $S\ge 1$ on the
$\tfrac13M\ell_j$-neighbourhood of $\Lambda_j^{\sim}$ (the points at
$\operatorname{dist}_\infty$-distance less than $\tfrac13M\ell_j$ from $\Lambda_j^{\sim}$); on that
neighbourhood $\bar\zeta=(\bar\zeta_\square)_{\square\in\mathcal{C}_j}$ is a partition of unity
subordinate to $\mathcal{C}_j$, and for $\square\in\mathcal{C}_j$ of scale $s$
\begin{equation*}
\ell_s\,|\nabla\bar\zeta_\square|\le C_\zeta\frac{L^2}{M},\qquad
\ell_s^2\,|\nabla^2\bar\zeta_\square|\le C_\zeta\frac{L^4}{M^2},
\end{equation*}
where $C_\zeta=C_\zeta(d)$ is a constant depending on $d$ only.
\item[(b)] $\mathfrak{s}_j\equiv k$ on $A$, and
\begin{equation*}
P\supset\Lambda^{(-8LR_{k+1})}\supset\Omega_{k+1}\supset\Lambda_{k+1}.
\end{equation*}
The cubes of $\mathcal{C}_j$ that meet $P$ are exactly the scale-$k$ cover cubes of
\cite{B-CMP109} used in \cite[p.~276]{B-CMP119} that meet $P$; every cube of
$\mathcal{C}_j$ overlapping one of them is one of Ba{\l}aban's cover cubes; and $S\equiv 1$ on them.
\item[(c)] If $\square\in\mathcal{C}_j$ has scale $s$ and is not default, then
\begin{equation*}
\tilde\square^{2}\subset\square_0\subset\Lambda_s\subset\Lambda_j\cap\Omega_s^{\natural}.
\end{equation*}
If $\square\in\mathcal{C}_j$ is default, then $s=j$, each of its tiles lies within $7M\ell_j$ of
$\Lambda_j^{c}$ (inside the collar
$\{x\in\Lambda_j:\operatorname{dist}_\infty(x,\Lambda_j^{c})\le(MR_j+1)\ell_j\}$ of
Notation~\ref{4.29:not-scales-labels}) or in the ring
$\Lambda_j^{\sim}\setminus\Lambda_j$, and $\square_0\subset\Omega_j^{\natural}$. In both cases
$l\ge s$ on $\square_0$, and $\square_0\subset\Omega_s$.
\item[(d)] The real form (d) of the hypothesis on this cover, and its complex form (d)$^{c}$: see
the proposition on the complex form of the hypothesis below.
\item[(e$^{\natural}$)] For every $x\in\Lambda_j^{\sim}$,
\begin{equation*}
\max\{j,\,l(x)-1\}\le\mathfrak{s}_j(x)\le\min\{l(x),\,k\}.
\end{equation*}
On the ring $\Lambda_j^{\sim}\setminus\Lambda_j$ one has $\mathfrak{s}_j=j=l(x)$, and the tiles
there are default. A tile is default only where $l=j$. At a site $x$ at most one generation,
namely $j=l(x)$, has a default tile, and all generations $j\le\mathfrak{s}(x)$ share the tile
$\Delta_{\mathfrak{s}(x)}(x)$.
\item[(f)] A cube $\square$ of $\mathcal{C}_j$ that does not meet $P$ satisfies
$\operatorname{dist}_\infty(\square,\Omega_{k+1})\ge 7LMR_{k+1}\ell_k$. By (b), every cube whose
scale, whose $\bar\zeta$ or whose overlaps differ from those of Ba{\l}aban's cover is such a cube;
the localised pieces attached to it have domains not contained in $\Lambda_{k+1}$, that is, they
belong to the terms $\mathbf{B}^{(k+1)}$ of \cite{B-CMP119}.
\item[(g)] The bookkeeping statement for the sites $x\in\Lambda_j^{\sim}\setminus P$: see the
lemma on far-field paths below.
\item[(h)] For $\square\in\mathcal{C}_j$: membership in $\mathcal{C}_j$, the scale, the default
status, and $\bar\zeta_\square$ near $\square\cap\Lambda_j^{\sim}$ depend on the label only through
the traces of $\Lambda_j$ and of $\{\Lambda_{s'}\}_{s'\le k}$ on $B_\infty(\square,9M\ell_k)$; and
$9M\ell_k<M\ell_{k+1}$.
\item[(h$'$)] For $D\subset\Lambda_j^{\sim}$, the trace on $D$ of the block structure
$\mathbf{B}'=\mathbf{B}_j(\Lambda_j^{\sim})$ of the restricted minimiser (which blocks, of which
levels) depends on the label only through
\begin{equation*}
\Lambda_j\cap B_\infty(D,3M_1\ell_j+2M\ell_j)\subset B_\infty(D,3M\ell_j).
\end{equation*}
For a localisation domain $Y$, the localised pieces attached to a cube $\square\in\mathcal{C}_j$
with $\square\subset Y$ read $\mathbf{B}'$ only inside a set $Y^{+}\subset B_\infty(Y,2M\ell_k)$.
The total reach is less than $5M\ell_k$, hence less than
$9M\ell_k$: it lies within the collar of (h).
\item[(h$''$)] This clause bears on the cut \cite[(1.49)]{B-CMP122b}. Let $k''\ge k+1$, let $C$ be a
connected component of $\Lambda_{k''}^{c}$, and let
$Y$ be a localisation domain with $C\cap Y=\emptyset$. Then
\begin{equation*}
\operatorname{dist}_\infty(Y,\Lambda_{s'}^{c}\cap C)\ge 8LMR_{k+1}\ell_k\qquad
\text{for all } s'\le k ,
\end{equation*}
and this distance exceeds the reach of every dependence of a piece $\mathbf{B}^{(k+1)}(Y)$ on label
data of index at most $k$, and of the test of index $k+1$ applied to that piece, which depends on
$\Omega_{k+1}$. These dependences and their reaches are:
\begin{itemize}
\item the cover, reach $9M\ell_k$, by (h);
\item the restricted minimisers, reach $5M\ell_k$, by (h$'$);
\item the native ranges $z\in\Lambda_j^{0}$ and the cut-offs $\varphi_j$, $j\le k$, whose reach is
at most
\begin{equation*}
2MR_j\ell_j\le 2MR_{k+1}\ell_{k+1}=2LMR_{k+1}\ell_k
\end{equation*}
by clause (b) of the lemma on the radii $R_j$;
\item the test of index $k+1$ applied to the piece, which reads $\Omega_{k+1}$ within $5M\ell_k$,
while $\Omega_{k+1}^{c}\cap C$ lies at depth at least $2LMR_{k+1}\ell_k$ inside $C$ by (T2) at level
$k+1$;
\item directly, the sequences $\{\Omega_i\}_{i\le k}$ and $\{S_i\}_{i\le k}$, which are data on the
side of $\Lambda_i^{c}$ and are read within $2M\ell_k$ of $Y$, whereas inside $C$ they are carried
by $\Lambda_i^{c}\cap C$, $i\le k$, which lies at distance at least $8LMR_{k+1}\ell_k$ from $Y$ by
the display above. Here $S_i$ is Ba{\l}aban's
large-field set of \cite[(2.40)--(2.42)]{B-CMP119}, with $S_i\subset\Omega_i\cap\Lambda_i^{c}$ a
union of $LM_2R_i$-cubes \cite[(3.21)]{B-CMP119}; it is not the collar written out
in (c), which Notation~\ref{4.29:not-scales-labels} also denotes by $S_j$; the two sets are
disjoint, and nothing about the one is inferred from the other.
\end{itemize}
Consequently $\mathbf{B}^{(k+1)}(Y)$ reads nothing of $C$ of index at most $k$, nor of
$\Omega_{k+1}^{c}\cap C$. Of index $k+1$ it reads $\Lambda_{k+1}$ and $S_{k+1}$ on $Y$
(\cite[(2.41)--(2.42)(i)]{B-CMP119} and the identification lemma) and, through the filing at
$z\in\Lambda_{k+1}^{0}$ \cite{B-CMP119}, the $(X,z)$-function being unchanged by the
identification lemma and the filing lemma, on the collar of $Y$ of one $MR_{k+1}$-layer. Hence it
reads $C$ only if $C^{\sim}\cap Y\neq\emptyset$, where $C^{\sim}$ is $C$ enlarged by one layer of
$MR_{k+1}$-cubes in Ba{\l}aban's sense. When $k''\ge k+2$ even this reading of $C$ is empty,
whether or not $C^{\sim}\cap Y\neq\emptyset$. Indeed, (T1) at $i=k+1$ is then in force when $C$ is
created, since the label lemma at the $\mathbf{T}$-stage that creates $\Lambda_{k''}$ gives (T1)
for all $i\le k''-1$; so $\Lambda_{k+1}^{c}\cap C$ lies at depth at least $8MR_{k+2}\ell_{k+2}$
inside $C$, and $8MR_{k+2}\ell_{k+2}>MR_{k+1}\ell_{k+1}$ by clause (b) of the lemma on the radii
$R_j$. The collar of $Y$ of one $MR_{k+1}$-layer, where it enters $C$, lies at depth at most
$MR_{k+1}\ell_{k+1}$ inside $C$, because $C\cap Y=\emptyset$; it therefore contains no point of
$\Lambda_{k+1}^{c}$, and the trace of $\Lambda_{k+1}$ on it is the whole of it.
\end{enumerate}
\end{proposition}

\begin{proof}[Proof of (a)]
We use throughout that the partitions $\pi_n$ are compatible and nested, that $M\ell_{n+1}=LM\ell_n$,
and that $\Lambda_{i+1}\subset\Lambda_i$ by (T0).

We first prove heredity of goodness. Let $\Delta_s(x)$ be good, with $s\ge 2$. Six layers of
$\pi_{s-1}$-cubes have total thickness $6M\ell_{s-1}<LM\ell_{s-1}=M\ell_s$, because $6<L$; since all
the sets involved are unions of aligned cubes of compatible partitions, and
$\Delta_{s-1}(x)\subset\Delta_s(x)$, it follows that
\begin{equation*}
\widetilde{\Delta_{s-1}(x)}^{\,6}\subset\widetilde{\Delta_s(x)}^{\,1}
\subset\widetilde{\Delta_s(x)}^{\,6}\subset\Lambda_s\subset\Lambda_{s-1},
\end{equation*}
the third inclusion being the goodness of $\Delta_s(x)$ and the last one (T0). Thus $\Delta_{s-1}(x)$
is good. The set of good scales at $x$ is therefore closed under $s\mapsto s-1$ for $s\ge 2$, and
it is the initial segment $\{1,\dots,\mathfrak{s}(x)\}$, $\mathfrak{s}(x)$ being its maximum.

Next, constancy on tiles. Let $x\in\Lambda_j^{\sim}$, $s^*:=\mathfrak{s}_j(x)$ and
$x'\in\Delta_{s^*}(x)$. Since the partitions are nested, $\Delta_s(x')=\Delta_s(x)$ for all
$s\ge s^*$, so the goodness tests at all scales $s\ge s^*\ge j$ coincide at $x$ and at $x'$. If
$\mathfrak{s}(x)\ge j$, then $s^*=\mathfrak{s}(x)$, the cube $\Delta_{s^*}(x')$ is good and no
larger scale is good at $x'$, so $\mathfrak{s}(x')=s^*$. If $\mathfrak{s}(x)<j$, then $s^*=j$, no
scale $\ge j$ is good at $x$, hence none at $x'$, and $\mathfrak{s}_j(x')=j$. In both cases
$\mathfrak{s}_j(x')=s^*$, and $x'$ is default if and only if $x$ is. For this to make sense we need
$\Delta_{s^*}(x)\subset\Lambda_j^{\sim}$: a good cube $\Delta_{s^*}(x)$ lies in
$\tilde\Delta^{6}\subset\Lambda_{s^*}\subset\Lambda_j$ by (T0), as $s^*\ge j$; a default one is
$\Delta_j(x)$, and $\Lambda_j^{\sim}$, being a union of $\pi_j$-cubes and containing $x$, contains
$\Delta_j(x)$. Now let two tiles
have a common site $z$. By constancy each tile has, at each of its sites, scale equal to
$\mathfrak{s}_j$ there; so both have scale $\mathfrak{s}_j(z)$ and both equal
$\Delta_{\mathfrak{s}_j(z)}(z)$. Since every $x\in\Lambda_j^{\sim}$ lies in its tile
$\Delta_{\mathfrak{s}_j(x)}(x)\subset\Lambda_j^{\sim}$, the tiles partition $\Lambda_j^{\sim}$.

Next, the grading. Lower bound: let $\Delta$ be a tile of scale $s\ge j+1$; then
$\mathfrak{s}=s$ on $\Delta$, so $\Delta$ is good. Let $y\in\tilde\Delta^{1}$. As in the proof of
heredity, $\Delta_{s-1}(y)\subset\Delta_s(y)\subset\tilde\Delta^{1}$, and six layers of
$\pi_{s-1}$-cubes around it stay within one further layer of $\pi_s$-cubes, so
\begin{equation*}
\widetilde{\Delta_{s-1}(y)}^{\,6}\subset\tilde\Delta^{2}\subset\tilde\Delta^{6}
\subset\Lambda_s\subset\Lambda_{s-1}.
\end{equation*}
Hence $y\in\Lambda_s\subset\Lambda_j$, the cube $\Delta_{s-1}(y)$ is good, and
$\mathfrak{s}_j(y)\ge s-1$. For a tile of scale $s=j$ the bound $\mathfrak{s}_j\ge j\ge s-1$ holds
by the definition of $\mathfrak{s}_j$. Upper bound: let $\Delta$ be a tile of scale $s$, and
suppose that some $y\in\tilde\Delta^{1}\cap\Lambda_j^{\sim}$ has $\mathfrak{s}_j(y)=s''\ge s+2$.
The tile $\Delta''\ni y$ has scale $s''>j$, so it is good and $y\in\Lambda_j$. For $x\in\Delta$,
\begin{equation*}
|x-y|_\infty\le 2M\ell_s<M\ell_{s''},
\end{equation*}
so $x\in\tilde\Delta''^{\,1}$, and the lower bound at the tile $\Delta''$ of scale $s''\ge j+1$
gives $\mathfrak{s}_j(x)\ge s''-1>s$. But $\mathfrak{s}_j(x)=s$ by constancy, since
$x\in\Delta\subset\Lambda_j^{\sim}$; this is a contradiction. We have proved
\begin{equation}\label{4.29:eq-native-cover-1}
\mathfrak{s}_j\in\{s-1,s,s+1\}\quad\text{on }\tilde\Delta^{1}\cap\Lambda_j^{\sim},
\qquad\Delta\in\mathcal{T}_j\text{ of scale } s .
\end{equation}

Next, overlaps. Let $\square,\square'\in\mathcal{C}_j$ have scales $s,s'$, let $\Delta,\Delta'$ be
tiles generating them, so that $\square\subset\tilde\Delta^{1}$ and
$\square'\subset\tilde\Delta'^{\,1}$, and let $\square$ and $\square'$ share a site $z$. If
$z\in\Lambda_j^{\sim}$, then \eqref{4.29:eq-native-cover-1} at both tiles gives
$\mathfrak{s}_j(z)\in\{s-1,s,s+1\}\cap\{s'-1,s',s'+1\}$, whence $|s-s'|\le 2$. If
$z\notin\Lambda_j^{\sim}$, then $z\notin\Lambda_j\subset\Lambda_j^{\sim}$, so neither
$\tilde\Delta^{1}$ nor $\tilde\Delta'^{\,1}$ lies in $\Lambda_j$. A good tile $\Delta$ of scale
$s\ge j$ satisfies $\tilde\Delta^{1}\subset\tilde\Delta^{6}\subset\Lambda_s\subset\Lambda_j$; so
neither tile is good, both are default, and $s=s'=j$.

Next, the partition. Let $x$ be a site of a tile $\Delta=\prod_\mu[a_\mu,a_\mu+M\ell_s)$ of scale
$s$. A cube $\square$ of the complete family of cover cubes of scale $s$, with centre $y$, has
$\zeta_\square(x)\neq 0$ only if $|x_\mu-y_\mu|<\tfrac23M\ell_s$ for all $\mu$. Since $y$ is a
vertex of $\pi_s$, that is $y_\mu\in a_\mu+M\ell_s\mathbb{Z}$, and
$x_\mu\in[a_\mu,a_\mu+M\ell_s)$, this forces $y_\mu\in\{a_\mu,a_\mu+M\ell_s\}$: $y$ is a vertex of
$\Delta$, so $\square\supset\Delta$, so $\square\in\mathcal{C}_j$. Hence the sum over the cubes of
$\mathcal{C}_j$ of scale $s$ of $\zeta_\square(x)$ equals the sum over the complete scale-$s$
family, which is $1$. Let now $\square'\in\mathcal{C}_j$ have scale $s'\neq s$ and
$\zeta_{\square'}(x)\neq0$. Since $\zeta_{\square'}$ vanishes outside $\square'$, $x\in\square'$,
so $x$ lies in $\tilde\Delta'^{\,1}\cap\Lambda_j^{\sim}$ for a tile $\Delta'$ generating
$\square'$, and \eqref{4.29:eq-native-cover-1} at $\Delta'$ gives $s=\mathfrak{s}_j(x)\in
\{s'-1,s',s'+1\}$, that is $s'\in\{s-1,s+1\}$. For each such $s'$ the centre of $\square'$ must lie
within $\tfrac23M\ell_{s'}$ of $x$ in every coordinate, which leaves at most two values per
coordinate; so at most $2^d$ such cubes per scale have $\zeta_{\square'}(x)\neq0$. As
$0\le\zeta\le1$ (Definition~\ref{4.29:def-tiles-cover}), we obtain $1\le S(x)\le 1+4\cdot 2^d$. The
argument in fact gives $S(x)\le 1+2\cdot2^d$, the scales other than $s$ present at $x$ being $s-1$
and $s+1$ only; the bound stated in (a) follows a fortiori.

Next, the derivatives. Let $\square\in\mathcal{C}_j$ have scale $s$. The quotient
$\bar\zeta_\square=\zeta_\square/S$ varies only on the support of $\zeta_\square$, which lies in
$\square$; every cube $\square'$ whose function $\zeta_{\square'}$ is non-zero there overlaps
$\square$, hence has scale $s'\ge s-2$ by the overlap statement, and $\ell_{s'}\ge\ell_s/L^2$. Since
$\zeta_{\square'}(x)=\prod_\mu\zeta((M\ell_{s'})^{-1}(x_\mu-y'_\mu))$, with $y'$ the centre of
$\square'$, and the derivatives of $\zeta$ are at most $5$,
\begin{equation*}
|\partial\zeta_{\square'}|\le\frac{5L^2}{M\ell_s},\qquad
|\partial^2\zeta_{\square'}|\le\frac{25L^4}{(M\ell_s)^2}.
\end{equation*}
By the counting of the previous paragraph at most $2^d$ such cubes occur per scale, and by the
overlap statement at most five scales occur, so $|\nabla S|$ and $|\nabla^2S|$ obey the same bounds
multiplied by a factor depending on $d$ only. The quotient rule,
\begin{equation*}
\nabla\frac{\zeta_\square}{S}=\frac{\nabla\zeta_\square}{S}-\frac{\zeta_\square\nabla S}{S^2},
\end{equation*}
and its second-order analogue, whose terms are $\nabla^2\zeta_\square/S$,
$\nabla\zeta_\square\otimes\nabla S/S^2$, $\zeta_\square\nabla^2S/S^2$ and
$\zeta_\square\nabla S\otimes\nabla S/S^3$ with coefficients depending on $d$ only, together with
$S\ge1$ and $0\le\zeta_\square\le 1$, give the two bounds of (a) with $C_\zeta=C_\zeta(d)$.
Finally $\sum_\square\bar\zeta_\square=S/S=1$ wherever $S>0$, and
$\operatorname{supp}\bar\zeta_\square\subset\operatorname{supp}\zeta_\square\subset\square$, so
$\bar\zeta$ is a partition of unity subordinate to $\mathcal{C}_j$ wherever $S\ge1$.

Finally, the neighbourhood. On $\Lambda_j^{\sim}$ itself $S\ge1$ by the partition step. Let $x$
satisfy $\operatorname{dist}_\infty(x,\Lambda_j^{\sim})<\tfrac13M\ell_j$. Then there is a site $z$
of $\Lambda_j^{\sim}$ with $|x-z|_\infty<\tfrac13M\ell_j$; $z$ lies in a tile
$\Delta=\prod_\mu[a_\mu,a_\mu+M\ell_s)$ of some scale $s\ge j$, so
$\operatorname{dist}_\infty(x,\Delta)<\tfrac13M\ell_s$. The tile is half-open, and its vertices are
the $2^d$ points of $\prod_\mu\{a_\mu,a_\mu+M\ell_s\}$. A cube $\square_y$ of the complete scale-$s$
family, with centre $y\in a+M\ell_s\mathbb{Z}^d$, has $\zeta_{\square_y}(x)\neq0$ only if
$|x_\mu-y_\mu|<\tfrac23M\ell_s$ for all $\mu$. Since
$x_\mu\in(a_\mu-\tfrac13M\ell_s,\,a_\mu+\tfrac43M\ell_s)$, this forces
$y_\mu\in\{a_\mu,a_\mu+M\ell_s\}$: $y$ is a vertex of $\Delta$, so $\square_y\supset\Delta$ and
$\square_y\in\mathcal{C}_j$. Put $t_\mu:=(x_\mu-a_\mu)/(M\ell_s)\in(-\tfrac13,\tfrac43)$. Then
\begin{equation*}
\sum_{y\text{ vertex of }\Delta}\zeta_{\square_y}(x)=\prod_\mu\bigl[\zeta(t_\mu)+\zeta(t_\mu-1)\bigr]=1,
\end{equation*}
because the one-dimensional sum $\sum_n\zeta(t_\mu-n)=1$ of the complete family has only the two
terms $n=0$ and $n=1$ non-zero for $t_\mu\in(-\tfrac13,\tfrac43)$. All other terms of $S(x)$ are
non-negative, so $S(x)\ge1$. This completes the proof of (a).
\end{proof}

The remaining clauses (b), (c), (e$^{\natural}$), (f) and (h)--(h$''$) are proved below.

\begin{proof}[Proof of clauses (b), (c), (e)$^{\natural}$ and (f) of Proposition~\ref{4.29:prop-native-cover}]
\label{4.29:prf-window-clauses-proof}
We prove the four clauses for a fixed generation $j\le k$. We write $\Lambda=\Lambda_k$,
$A=\Lambda^{(-6)}$ and $P=\Lambda^{(-12)}$, and we call $\Lambda_j^{\sim}\setminus\Lambda_j$ the
\emph{ring} of generation $j$. Since $\ell_n=L^n\eta$, we have $\ell_{n+1}=L\ell_n$. As in
Proposition~\ref{4.29:prop-native-cover}, $L\ge 13$ and $R_i\ge 7$ for every $i$.

Among the hypotheses of Proposition~\ref{4.29:prop-native-cover} are the following properties of
the action domains $\Lambda_i$, the sequence $\Omega_i$ and the native sequence
$\Omega^{\natural}_i$. For every level $i$,
\begin{equation}\label{4.29:eq-window-clauses-proof-1}
\Lambda_i\subset\Omega^{\natural}_i\subset\Omega_i,\qquad
\Omega^{\natural}_{i+1}\subset\Lambda_i,\qquad\text{hence }\Lambda_{i+1}\subset\Lambda_i .
\end{equation}
For every $i\le k$,
\begin{equation}\label{4.29:eq-window-clauses-proof-2}
\operatorname{dist}_\infty\bigl(\Omega^{\natural}_{i+1},\Lambda_i^{c}\bigr)\ge 8MR_{i+1}\ell_{i+1},
\qquad \Omega_{k+1}=\Omega^{\natural}_{k+1}\supset\Lambda_{k+1}.
\end{equation}
For every level $i$,
\begin{equation}\label{4.29:eq-window-clauses-proof-3}
\operatorname{dist}_\infty\bigl(\Lambda_i,(\Omega^{\natural}_i)^{c}\bigr)\ge 2MR_i\ell_i .
\end{equation}
By \eqref{4.29:eq-window-clauses-proof-1} the native sequence is decreasing. Hence, since the
native index $l(x)$ is the index $a$ with $x\in\Omega^{\natural}_a\setminus\Omega^{\natural}_{a+1}$,
every point of $\Omega^{\natural}_s$ has $l\ge s$.

Two elementary facts about cubes are used repeatedly. For $\Delta\in\pi_s$ the set
$\widetilde{\Delta}^{r}$ is the cube of side $(2r+1)M\ell_s$ concentric with $\Delta$. Hence for
$x\in\Delta$ every site of $\widetilde{\Delta}^{r}$ lies at sup-distance less than $(r+1)M\ell_s$
from $x$; for $r=6$ this is \eqref{4.29:eq-scales-labels-a}. Moreover, any two sites of
$\widetilde{\Delta}^{1}$ are at sup-distance less than $3M\ell_s$. Finally, $\Lambda^{(-r)}$
decreases as $r$ increases.

\smallskip
\noindent\emph{Clause (b).} Let $x\in A$. By the definition of $\Lambda^{(-6)}$ we have
$\widetilde{\Delta_k(x)}^{6}\subset\Lambda=\Lambda_k$, so the cube $\Delta_k(x)$ is good. Since
$\mathfrak{s}(x)\le k$ by definition, $\mathfrak{s}(x)=k$. As $j\le k$, this gives
$\mathfrak{s}_j(x)=\max\{j,k\}=k$.

We apply \eqref{4.29:eq-window-clauses-proof-2} at $i=k$ and use
$\Omega_{k+1}=\Omega^{\natural}_{k+1}$. This gives
\begin{equation*}
\operatorname{dist}_\infty(\Omega_{k+1},\Lambda^{c})\ge 8MR_{k+1}\ell_{k+1}=8LR_{k+1}\cdot M\ell_k .
\end{equation*}
Moreover $8LR_{k+1}\ge 12$, because $L\ge 13$ and $R_{k+1}\ge 7$. Since $\Lambda^{(-r)}$ decreases
in $r$, it follows that $\Lambda^{(-8LR_{k+1})}\subset\Lambda^{(-12)}=P$.

Next, among the same hypotheses, $\Omega_{k+1}$ is a union of $\pi_{k+1}$-cubes and
$\Lambda=\Lambda_k$ is a union of $MR_k$-cubes; since the partitions are compatible, both sets
are unions of $\pi_k$-cubes, and so is $\Lambda^{c}$. Put $r=8LR_{k+1}$, let $\Delta\in\pi_k$
with $\Delta\subset\Omega_{k+1}$, and let $y\in\widetilde{\Delta}^{r}$. The cube $\Delta_k(y)$
lies within $r$ layers of $\pi_k$-cubes of $\Delta$, so it contains a site at sup-distance at most
$(r-1)M\ell_k$ plus one lattice spacing from a site of $\Delta$. As the spacing of the lattice of
sites is smaller than the side $M\ell_k$ of a $\pi_k$-cube, this gives
$\operatorname{dist}_\infty(\Delta_k(y),\Delta)<rM\ell_k$. If $y$ were not in $\Lambda$, the whole
cube $\Delta_k(y)$ would lie in $\Lambda^{c}$, and the bound just displayed would be violated.
Hence $\widetilde{\Delta}^{r}\subset\Lambda$ for every $\pi_k$-cube $\Delta\subset\Omega_{k+1}$,
that is, $\Omega_{k+1}\subset\Lambda^{(-8LR_{k+1})}$. The inclusion
$\Omega_{k+1}\supset\Lambda_{k+1}$ is part of \eqref{4.29:eq-window-clauses-proof-2}.

Now let $\square\in\mathcal{C}_j$ meet $P$ at a site $z$, and let $\Delta$ be a tile generating
$\square$, of scale $s\le k$. Then $\square\subset\widetilde{\Delta}^{1}$, so
$z\in\widetilde{\Delta}^{1}$. Every site $y$ of $\widetilde{\Delta}^{1}$ therefore satisfies
\begin{equation*}
|y-z|_\infty<3M\ell_s\le 3M\ell_k .
\end{equation*}
For such $y$, the cube $\Delta_k(y)$ lies within three layers of $\pi_k$-cubes of $\Delta_k(z)$.
Since $z\in P$, we obtain
\begin{equation*}
\widetilde{\Delta_k(y)}^{9}\subset\widetilde{\Delta_k(z)}^{12}\subset\Lambda .
\end{equation*}
Thus $\widetilde{\Delta}^{1}\subset\Lambda^{(-9)}\subset A$. In particular $\Delta\subset A$.
Since $\Delta=\Delta_{\mathfrak{s}_j(x)}(x)$ for some $x\in\Delta\subset A$, and
$\mathfrak{s}_j=k$ on $A$ by the first paragraph, the tile $\Delta$ has scale $s=k$. Hence
$\square$ is a cover cube of scale $k$, that is, one of the
cubes of Ba{\l}aban's construction at scale $k$.

Let $\square'\in\mathcal{C}_j$ overlap $\square$ at a site $z'$. Since
$\square\subset\Lambda^{(-9)}$, we have $z'\in\Lambda^{(-9)}$. Let $\Delta'$ be a tile
generating $\square'$. Then $z'\in\square'\subset\widetilde{\Delta'}^{1}$, and every site of
$\widetilde{\Delta'}^{1}$ is at sup-distance less than $3M\ell_k$ from $z'$. The argument of the
previous paragraph, with $\Lambda^{(-9)}$ in place of $P$ and six layers in place of nine, gives
$\widetilde{\Delta'}^{1}\subset\Lambda^{(-6)}=A$. Hence the tile $\Delta'$ lies in $A$ and has
scale $k$.

Each $\zeta_{\square'}$ vanishes off $\square'$. Hence only cubes of scale $k$ contribute to
$S=\sum_{\square'\in\mathcal{C}_j}\zeta_{\square'}$ on $\square$, and on $\square$ the sum $S$
equals the sum over the cubes of scale $k$ of $\mathcal{C}_j$. Every site of $\square$ lies in
$A$, hence in a tile of scale $k$. By clause (a) of Proposition~\ref{4.29:prop-native-cover}, the
sum over the scale-$k$ cubes of $\mathcal{C}_j$ equals $1$ on such a tile. Therefore $S=1$ on
$\square$.

Conversely, let a cover cube of scale $k$ meet $P$ at $z$. It is a union of $\pi_k$-cubes, so it
contains $\Delta_k(z)$. Since $P$ is a union of $\pi_k$-cubes, $\Delta_k(z)\subset P\subset A$.
Because $\mathfrak{s}_j=k$ on $A$, the cube $\Delta_k(z)$ is a tile of scale $k$, and every cover
cube of scale $k$ containing it belongs to $\mathcal{C}_j$.

Finally, by \eqref{4.29:eq-window-clauses-proof-1} we have
\begin{equation*}
P\subset A\subset\Lambda_k\subset\Lambda_j .
\end{equation*}
So the ring of generation $j$ lies outside
$\Lambda_j$, hence outside $A$ and $P$.

\smallskip
\noindent\emph{Clause (c).} Let $\square\in\mathcal{C}_j$ have scale $s$, and let $\Delta$ be a
tile generating it. Then $\square\subset\widetilde{\Delta}^{1}$. Taking the enlargements in
$\pi_s$, this gives
\begin{equation*}
\square_0=\widetilde{\square}^{5}\subset\widetilde{\Delta}^{6}.
\end{equation*}

Suppose first that $\square$ is not default. Then it has a generating tile $\Delta$ that is not
default, so $\mathfrak{s}\ge j$ at a point $x$ of $\Delta$. Consequently
$\mathfrak{s}_j(x)=\mathfrak{s}(x)=s$, and $\Delta=\Delta_{\mathfrak{s}(x)}(x)$ is good, that is,
$\widetilde{\Delta}^{6}\subset\Lambda_s$. Since $s\ge j$, \eqref{4.29:eq-window-clauses-proof-1}
gives $\Lambda_s\subset\Lambda_j\cap\Omega^{\natural}_s$. Hence
\begin{equation*}
\square_0\subset\widetilde{\Delta}^{6}\subset\Lambda_s\subset\Lambda_j\cap\Omega^{\natural}_s .
\end{equation*}

Suppose next that $\square$ is default. Then every generating tile is default, and we may take
$\Delta=\Delta_j(x)$ with $x\in\Lambda_j^{\sim}$, $\mathfrak{s}(x)<j$ and $s=j$.

If $x\in\Lambda_j$, the cube $\Delta_j(x)$ is not good, since otherwise
$\mathfrak{s}(x)\ge j$. So $\widetilde{\Delta}^{6}\not\subset\Lambda_j$. By
\eqref{4.29:eq-scales-labels-a} some site outside $\Lambda_j$ lies at sup-distance less than
$7M\ell_j$ from $x$. Since $R_j\ge 7$, this gives
\begin{equation*}
\operatorname{dist}_\infty(x,\Lambda_j^{c})<7M\ell_j\le MR_j\ell_j .
\end{equation*}

If instead $x$ lies in the ring, then $\widetilde{\Delta}^{1}$ meets $\Lambda_j$. Since
$x\in\Delta$, the site $x$ lies within $2M\ell_j$ of $\Lambda_j$.

We write $B_\infty(x,r)$ for the set of sites at sup-distance less than $r$ from $x$, and
$B_\infty(Y,r)$ for the union of these balls over $x\in Y$.
In both cases \eqref{4.29:eq-scales-labels-a} gives the first inclusion below, and the position of
$x$ relative to $\Lambda_j$ gives the second:
\begin{equation*}
\square_0\subset\widetilde{\Delta}^{6}\subset B_\infty(x,7M\ell_j)
\subset B_\infty(\Lambda_j,9M\ell_j)\subset\Omega^{\natural}_j .
\end{equation*}
The last inclusion follows from
\eqref{4.29:eq-window-clauses-proof-3}, because $2MR_j\ell_j\ge 14M\ell_j>9M\ell_j$.

In both cases, therefore, $l\ge s$ on $\square_0$, since every point of $\Omega^{\natural}_s$ has
$l\ge s$. In the non-default case $\square_0\subset\Lambda_s\subset\Omega^{\natural}_s$; in the
default case $s=j$ and $\square_0\subset\Omega^{\natural}_j$. Finally,
$\Omega^{\natural}_s\subset\Omega_s$ by \eqref{4.29:eq-window-clauses-proof-1}, so
$\square_0\subset\Omega_s$.

\smallskip
\noindent\emph{Clause (e)$^{\natural}$.} Let $x\in\Lambda_j^{\sim}$.

\emph{Upper bound.} If $\Delta_s(x)$ is good, then $\Delta_s(x)\subset\Lambda_s\subset\Omega^{\natural}_s$
by \eqref{4.29:eq-window-clauses-proof-1}, which forces $l(x)\ge s$. Hence
$\mathfrak{s}(x)\le l(x)$. By Definition~\ref{4.29:def-tiles-cover},
$\Lambda_j^{\sim}\subset\Omega^{\natural}_j$, so $l(x)\ge j$. Therefore
$\mathfrak{s}_j(x)=\max\{j,\mathfrak{s}(x)\}\le l(x)$. Also $\mathfrak{s}_j(x)\le k$, since
$\mathfrak{s}(x)\le k$ and $j\le k$.

\emph{Lower bound for $x\in\Lambda_j$.} Put $l:=l(x)$, so that $l\ge j$.

Suppose $l>j$. Then $x\in\Omega^{\natural}_l$, and $l\le k+1$ because the native sequence is
$\{\Omega^{\natural}_i\}_{i\le k+1}$. So $i=l-1\le k$, and \eqref{4.29:eq-window-clauses-proof-2}
at $i=l-1$ gives
\begin{equation*}
\operatorname{dist}_\infty(x,\Lambda_{l-1}^{c})\ge 8MR_l\ell_l=8LR_l\,M\ell_{l-1}>7M\ell_{l-1}.
\end{equation*}
By \eqref{4.29:eq-scales-labels-a} every site of $\widetilde{\Delta_{l-1}(x)}^{6}$ is at
sup-distance less than $7M\ell_{l-1}$ from $x$, hence lies in $\Lambda_{l-1}$. So
$\Delta_{l-1}(x)$ is good at the scale $l-1\ge j$, and
$\mathfrak{s}(x)\ge l-1\ge j$.

Suppose $l=j$. Then the bound $\mathfrak{s}_j(x)\ge l-1$ holds by the definition
$\mathfrak{s}_j=\max\{j,\mathfrak{s}\}$.

In both cases $\max\{j,l(x)-1\}\le\mathfrak{s}_j(x)$ for every $x\in\Lambda_j$; the ring is treated
below. If the tile at $x\in\Lambda_j$ is default, then $\mathfrak{s}(x)<j$. This excludes the
case $l>j$, in which $\mathfrak{s}(x)\ge l-1\ge j$. Hence $j=l$: a tile is default only where
$l=j$.

\emph{The ring.} Let $x\in\Lambda_j^{\sim}\setminus\Lambda_j$. Then $x\in\Omega^{\natural}_j$, and
$x\notin\Omega^{\natural}_{j+1}$ because $\Omega^{\natural}_{j+1}\subset\Lambda_j$ by
\eqref{4.29:eq-window-clauses-proof-1}. So $l(x)=j$. A good cube $\Delta_s(x)$ with $s\ge j$
would lie in $\Lambda_s\subset\Lambda_j$, which does not contain $x$. Hence $\mathfrak{s}(x)<j$,
so $\mathfrak{s}_j(x)=j=l(x)$, the tile at $x$ is the default cube $\Delta_j(x)$, and both bounds
hold.

\emph{At most one ring.} A point lies in the ring of at most one generation. Let $x$ lie in the
ring of generation $j$, and let $j'\ne j$.

If $j'>j$, then by Definition~\ref{4.29:def-tiles-cover} and \eqref{4.29:eq-window-clauses-proof-1}
\begin{equation*}
\Lambda_{j'}^{\sim}\subset\Omega^{\natural}_{j'}\subset\Lambda_{j'-1}\subset\Lambda_j,
\end{equation*}
and $\Lambda_j$ does not contain $x$. So $x\notin\Lambda_{j'}^{\sim}$.

If $j'<j$, then by \eqref{4.29:eq-window-clauses-proof-1}
\begin{equation*}
x\in\Omega^{\natural}_j\subset\Lambda_{j-1}\subset\Lambda_{j'},
\end{equation*}
so $x$ is not in the ring of generation $j'$. Moreover $\mathfrak{s}(x)=j-1\ge j'$. Indeed, the
bound \eqref{4.29:eq-window-clauses-proof-2} at $i=j-1$, together with
\eqref{4.29:eq-scales-labels-a} as in the case $l>j$ above, shows that $\Delta_{j-1}(x)$ is a
good tile. On the other hand $\mathfrak{s}(x)<j$ was shown in the ring paragraph.

Consequently a default tile occurs at a point $x$ only in the generation $j=l(x)$, whether
$x\in\Lambda_j$ or $x$ lies in the ring of generation $j$; so at most one generation has a default
tile at $x$. Moreover, for every generation $j\le\mathfrak{s}(x)$ with $x\in\Lambda_j^{\sim}$ we
have $\mathfrak{s}_j(x)=\max\{j,\mathfrak{s}(x)\}=\mathfrak{s}(x)$, so all these generations have
the same tile $\Delta_{\mathfrak{s}(x)}(x)$ at $x$.

\smallskip
\noindent\emph{Clause (f).} Let $\square\in\mathcal{C}_j$ with $\square\cap P=\emptyset$.

Let $y\in\square\cap\Lambda$. Since $y\notin P$, we have
$\widetilde{\Delta_k(y)}^{12}\not\subset\Lambda$. Every site of this cube is at sup-distance less
than $13M\ell_k$ from $y$, so $\operatorname{dist}_\infty(y,\Lambda^{c})<13M\ell_k$. By the bound
$\operatorname{dist}_\infty(\Omega_{k+1},\Lambda^{c})\ge 8LR_{k+1}M\ell_k$ from clause (b) and the
triangle inequality,
\begin{equation*}
\operatorname{dist}_\infty(y,\Omega_{k+1})\ge(8LR_{k+1}-13)M\ell_k\ge 7LR_{k+1}M\ell_k .
\end{equation*}
The last inequality holds because $LR_{k+1}\ge 13$.

Let $y\in\square$ with $y\notin\Lambda$; these are the ring sites when $j=k$, and the sites of
cubes extending beyond $\Lambda$. For such $y$, \eqref{4.29:eq-window-clauses-proof-2} at $i=k$
gives $\operatorname{dist}_\infty(y,\Omega_{k+1})\ge 8LR_{k+1}M\ell_k$ directly.

Combining the two cases,
\begin{equation*}
\operatorname{dist}_\infty(\square,\Omega_{k+1})\ge 7LMR_{k+1}\ell_k .
\end{equation*}

Consider a term of the localised expansion born at $\square$, that is, arising at the cube
$\square$ of the cover. By the localisation of \cite[\S3]{B-CMP109}, taken by statement, a term
born at $\square$ is localised in a domain $Y$ with $Y\supset\square_0$. On the present cover this
localisation is used through the transfer of the analysis of \cite[\S3]{B-CMP109} to the cover,
the sentence printed without proof in \cite[p.~276]{B-CMP119}; cited as printed. Hence
$Y\supset\square_0\supset\square$. Since $\square\cap P=\emptyset$ and $\Lambda_{k+1}\subset P$
by clause (b) (or directly by the distance bound just proved), $\square\cap\Lambda_{k+1}=\emptyset$;
as $Y\supset\square\ne\emptyset$, therefore $Y\not\subset\Lambda_{k+1}$.

Conversely, consider a term with domain $Y\subset\Lambda_{k+1}$. By the same localisation it is
born at a cube $\square\subset Y$, so that
\begin{equation*}
\square\subset Y\subset\Lambda_{k+1}\subset P .
\end{equation*}
By clause (b),
$\square$ is one of Ba{\l}aban's cubes of scale $k$, every cube of $\mathcal{C}_j$ overlapping it
is one of Ba{\l}aban's, and $S=1$ on them. Hence $\bar\zeta_\square=\zeta_\square$, and the term
reads only $\bar\zeta_\square=\zeta_\square$ and the auxiliary function $\zeta'_\square$ attached
to $\square$ in Ba{\l}aban's construction. Therefore the family of terms with domain
$Y\subset\Lambda_{k+1}$ is Ba{\l}aban's family without change. This family is the material of the
localised families $\mathbf{E}^{(k+1)}$ and $\mathbf{R}''^{(k+1)}$, of the coupling increment
$\beta_{k+1}$ and of the term $E_0$ in \cite{B-CMP119}.
\end{proof}

\begin{proof}[Proof of clauses (h), (h$'$) and (h$''$) of Proposition~\ref{4.29:prop-native-cover}]
\label{4.29:prf-label-locality-proof}
We use the scales, cubes and label of Notation~\ref{4.29:not-scales-labels} and the native scale,
tiles and cover of Definition~\ref{4.29:def-tiles-cover}. For a set $A$ of sites and $r>0$ we
write $B_\infty(A,r)$ for the set of sites at $\ell^\infty$-distance less than $r$ from $A$, and
$B_\infty(x,r)=B_\infty(\{x\},r)$. A cube of $\pi_s$ has side $M\ell_s$, and since
$\ell_{s'}=L^{s'}\eta$, we have $\ell_{s'}\le\ell_k$ for $s'\le k$.

From the lemma on the native sequence at the $\mathbf{T}$-stage, which is a hypothesis of
Proposition~\ref{4.29:prop-native-cover}, we use three properties. For
$i\le k$,
\begin{align}
&\Lambda_i\subset\Omega_i^{\natural}\subset\Omega_i,\qquad
\Omega^{\natural}_{i+1}\subset\Lambda_i,\qquad\text{hence }\Lambda_{i+1}\subset\Lambda_i;
\label{4.29:eq-label-locality-proof-1}\\
&\operatorname{dist}_\infty\bigl(\Omega^{\natural}_{i+1},\Lambda_i^{c}\bigr)
\ge 8MR_{i+1}\ell_{i+1},\qquad
\Omega_{k+1}=\Omega^{\natural}_{k+1}\supset\Lambda_{k+1};
\label{4.29:eq-label-locality-proof-2}
\end{align}
and, for $i\le k+1$,
\begin{equation}\label{4.29:eq-label-locality-proof-3}
\operatorname{dist}_\infty\bigl(\Lambda_i,(\Omega_i^{\natural})^{c}\bigr)\ge 2MR_i\ell_i .
\end{equation}
For the radii we use the hypothesis $R_i\ge7$ of Proposition~\ref{4.29:prop-native-cover} and
the inequality $LR_{i+1}\ge R_i$ between consecutive radii, which is part (b) of the statement
fixing the radii $R_i$; since $\ell_{i+1}=L\ell_i$, the latter gives
\begin{equation}\label{4.29:eq-label-locality-proof-4}
R_j\ell_j\le R_{j+1}\ell_{j+1}\le\dots\le R_{k+1}\ell_{k+1}\qquad(j\le k+1).
\end{equation}

\emph{Clause (h).} Let $y\in\Lambda_j$. By Definition~\ref{4.29:def-tiles-cover}, $\mathfrak{s}_j(y)$
--- and with it the tile through $y$, its scale and its default status --- is decided by the
test ``$y\in\Lambda_j$'' together with the tests
``$\widetilde{\Delta_{s'}(y)}^{\,6}\subset\Lambda_{s'}$'' for $s'\le k$. Apply
\eqref{4.29:eq-scales-labels-a} with $s=s'$, $x=y$ and $\Delta=\Delta_{s'}(y)$: every site of
$\widetilde{\Delta_{s'}(y)}^{\,6}$ lies at $\ell^\infty$-distance less than $7M\ell_{s'}$ from $y$.
Hence the test at scale $s'$ reads $\Lambda_{s'}$ only inside
\begin{equation*}
B_\infty(y,7M\ell_{s'})\subset B_\infty(y,7M\ell_k).
\end{equation*}
Next let $y\notin\Lambda_j$. Whether the tile $\Delta_j(y)$ belongs to the ring
$\Lambda_j^{\sim}\setminus\Lambda_j$ is, by the definition of $\Lambda_j^{\sim}$, the test
``$\widetilde{\Delta_j(y)}^{\,1}\cap\Lambda_j\neq\emptyset$''. The site $y$ lies in the central
cube $\Delta_j(y)$, of side $M\ell_j$, and one layer of $\pi_j$-cubes, each of side $M\ell_j$, is
added; so this test reads $\Lambda_j$ only inside
$B_\infty(y,2M\ell_j)\subset B_\infty(y,7M\ell_k)$. Finally, by the construction of the cover,
$\square\in\mathcal{C}_j$ with scale $s$ if and only if $\square$ is a cover cube of scale $s$
containing the tile $\Delta_{\mathfrak{s}_j(y')}(y')$, of scale $s$, of some site $y'\in\square$; and
$\zeta_\square$ depends on nothing but $\square$ once $\square$ is given.

Consider now the sum $S$ of clause (a) of Proposition~\ref{4.29:prop-native-cover} on $\square$. At
a point $x\in\square$ it is the sum of $\zeta_{\square'}$ over the cubes
$\square'\in\mathcal{C}_j$ containing $x$. Such a $\square'$, of scale $s'\le k$, is a union of
$2^d$ neighbouring $\pi_{s'}$-cubes, so it lies, together with its generating tiles, in
$B_\infty(x,2M\ell_{s'})\subset B_\infty(x,2M\ell_k)$. Its membership in $\mathcal{C}_j$ and its
scale are decided by the tests above at points of $B_\infty(x,2M\ell_k)$, each of which reads the
label within $7M\ell_k$ of its point; so they are decided inside $B_\infty(x,(2+7)M\ell_k)$. The
same bound covers the tests deciding $\square$ itself, which are made at points $y\in\square$.
Therefore everything listed in clause (h) reads only the traces of $\Lambda_j$ and of
$\{\Lambda_{s'}\}_{s'\le k}$ on $B_\infty(\square,9M\ell_k)$. Since $L\ge13$,
\begin{equation*}
9M\ell_k<LM\ell_k=M\ell_{k+1}.
\end{equation*}

\emph{Clause (h$'$).} The trace on $D$ of the restricted minimiser
$\mathbf{B}'=\mathbf{B}_j(\Lambda_j^{\sim})$ is read through its layers $D_n$, taken as in
\cite[(2.13)]{B-CMP119} with $D_0=\Lambda_j^{\sim}$. Their maximality is cubewise: $D_n$ is the
union of the cubes $c$ of side $M_1\ell_n$ such that $c\subset D_{n-1}$ and
$\operatorname{dist}(c,D_{n-1}^{c})\ge M_1\ell_n$. If $x\in c$, every site of $c$ lies within
$M_1\ell_n$ of $x$, and the distance test for $c$ involves $D_{n-1}$ only within $M_1\ell_n$ of
$c$. Hence $1_{D_n}(x)$ reads $D_{n-1}$ only on $B_\infty(x,2M_1\ell_n)$. By induction on $n$,
and because $\sum_{n\le j}\ell_n\le\ell_j L/(L-1)$ and $2L/(L-1)\le3$, the indicator $1_{D_n}(x)$
for $n\le j$ reads $D_0=\Lambda_j^{\sim}$ only on
\begin{equation*}
B_\infty\Bigl(x,2M_1\sum_{n\le j}\ell_n\Bigr)\subset
B_\infty\Bigl(x,\tfrac{2M_1\ell_jL}{L-1}\Bigr)\subset B_\infty(x,3M_1\ell_j).
\end{equation*}
The indicator $1_{\Lambda_j^{\sim}}$ at a site reads $\Lambda_j$ only within $2M\ell_j$ of it, by
the ring test in the proof of clause (h). So the trace of $\mathbf{B}'$ on $D$ reads the label
only through $\Lambda_j\cap B_\infty(D,3M_1\ell_j+2M\ell_j)$. Since
$M>M_2=L^2M_1\ge 3M_1$, we have $3M_1+2M<3M$, and this set lies in $B_\infty(D,3M\ell_j)$.

The localized pieces born at a cube $\square\subset Y$ are assembled from operator pieces
attached to cubes $\Delta'$ with $\widetilde{\Delta'}\subset Y^{+}$, where
$Y^{+}\subset B_\infty(Y,2M\ell_k)$ as in clause (h$'$). Each such operator piece is fixed by the
trace on $\widetilde{\Delta'}$ of the determining set of the step (the set whose traces fix the
operator pieces, in the sense of parts (b) and (d) of the identification lemma for localized
pieces), by those parts of that lemma and by \cite[\S3]{B-CMP119}. Hence these pieces read
$\mathbf{B}'$ only inside $Y^{+}$. Applying the preceding paragraph with $D=Y^{+}$, and using
$j\le k$, they read $\Lambda_j$ only within distance $2M\ell_k+3M\ell_j\le5M\ell_k$ of $Y$; and
$5M\ell_k<9M\ell_k$, so this is inside the collar of clause (h).

\emph{Clause (h$''$).} Let $k''\ge k+1$, let $C$ be a component of $\Lambda_{k''}^{c}$, and let
$Y$ be a localization domain with $C\cap Y=\emptyset$.

Suppose first that $\Lambda_k^{c}\cap C=\emptyset$. By \eqref{4.29:eq-label-locality-proof-1},
$\Lambda_k\subset\Lambda_{s'}$ for $s'\le k$, so $\Lambda_{s'}^{c}\subset\Lambda_k^{c}$, and every
$\Lambda_{s'}^{c}\cap C$ with $s'\le k$ is empty.

Otherwise let $p\in\Lambda_k^{c}\cap C$ and put $B:=B_\infty(p,8MR_{k+1}\ell_{k+1})$. By
\eqref{4.29:eq-label-locality-proof-2} at $i=k$, $B$ misses $\Omega^{\natural}_{k+1}$. By
\eqref{4.29:eq-label-locality-proof-1} and \eqref{4.29:eq-label-locality-proof-2},
$\Omega^{\natural}_{k+1}\supset\Lambda_{k+1}\supset\Lambda_{k''}$. For $k''\ge k+2$ the last
inclusion uses the nesting $\Lambda_{i+1}\subset\Lambda_i$ for $k+1\le i\le k''-1$, given by the
lemma on the native sequence at the later $\mathbf{T}$-stages together with the distance bound of
\eqref{4.29:eq-label-locality-proof-2} at those indices, used below. Thus $B\subset\Lambda_{k''}^{c}$. The
ball $B$ is connected in any convention, the wall-connectedness of \cite[\S1]{B-CMP109} included, and
it contains $p\in C$; as $C$ is a component of $\Lambda_{k''}^{c}$, we get $B\subset C$. Since
$Y\cap C=\emptyset$, $Y$ misses $B$, and
\begin{equation*}
\operatorname{dist}_\infty(Y,p)\ge8MR_{k+1}\ell_{k+1}=8LMR_{k+1}\ell_k .
\end{equation*}
As $\Lambda_{s'}^{c}\cap C\subset\Lambda_k^{c}\cap C$ for $s'\le k$, this is the asserted
distance bound.

We compare this distance with the reaches. Since $L\ge13$ and $R_{k+1}\ge7$, we have $9<8LR_{k+1}$
and $5<8LR_{k+1}$. So the cover of clause (h), with reach $9M\ell_k$, and the restricted
minimisers of clause (h$'$), with reach $5M\ell_k$, do not reach $\Lambda_{s'}^{c}\cap C$. By
\eqref{4.29:eq-label-locality-proof-4}, for $j\le k$,
\begin{equation*}
2MR_j\ell_j\le2LMR_{k+1}\ell_k<8LMR_{k+1}\ell_k .
\end{equation*}
A third route by which the level-$j$ functionals read $\Lambda_j$, with reach $2MR_j\ell_j$, is
the range $z\in\Lambda_j^{0}$ of \cite[\S2]{B-CMP119} and the cut-off functions $\varphi_j$. It
belongs to Ba{\l}aban's construction, through the native reading-range hypothesis and the statement
identifying an old boundary term with the new localized term, and not to the cover of
Proposition~\ref{4.29:prop-native-cover}; the last display shows that this reach, too, is smaller
than $8LMR_{k+1}\ell_k$.

Let $q\in\Omega_{k+1}^{c}\cap C$. By \eqref{4.29:eq-label-locality-proof-3} at $i=k+1$ and by
$\Omega_{k+1}=\Omega^{\natural}_{k+1}$, the ball $B_\infty(q,2MR_{k+1}\ell_{k+1})$ misses
$\Lambda_{k+1}\supset\Lambda_{k''}$. By the connectedness argument above, it lies in $C$. Thus
$\Omega_{k+1}^{c}\cap C$ lies at depth at least
\begin{equation*}
2MR_{k+1}\ell_{k+1}=2LMR_{k+1}\ell_k>5M\ell_k
\end{equation*}
inside $C$. This is what the membership test in $\Omega_{k+1}$ performed by the localized piece
requires, since that test reads $\Omega_{k+1}$ only within $5M\ell_k$ of $Y$.

In Ba{\l}aban's construction the big structure $\mathbb{B}$, the walk cubes $\sigma_0$ and the
characteristic functions read $\{\Omega_i\}_{i\le k+1}$ and $\{S_i\}_{i\le k}$ within $2M\ell_k$
of $Y$. Here $S_i$ are Ba{\l}aban's large-field sets \cite[(2.40)--(2.42)]{B-CMP119}, not the
collar sets of the same letter in Notation~\ref{4.29:not-scales-labels}. For $i\le k$,
$\Omega_i^{c}\subset\Lambda_i^{c}$ by \eqref{4.29:eq-label-locality-proof-1}, and
$S_i\subset\Omega_i\cap\Lambda_i^{c}$ by \cite[\S2]{B-CMP119} and \cite[(3.21)]{B-CMP119}.
So their traces on $C$ lie in $\Lambda_i^{c}\cap C$, at distance at least $8LMR_{k+1}\ell_k$
from $Y$. The set $\Omega_{k+1}$ is handled as in the membership test above.

It remains to treat index $k+1$. Ba{\l}aban's functions are taken ``restricted to $X$'' with the
sets $\{S_i\cap X\}$ \cite[(2.42)(i)]{B-CMP119}; for the piece at hand this restriction is to
$Y$, and $Y\cap C=\emptyset$. Further, Ba{\l}aban puts
$E^{(k+1)}(X,z)=E_0^{(k+1)}(\Lambda_{k+1},X,z)$ for $z\in\Lambda^{0}_{k+1}$ and
$X\subset\Lambda_{k+1}$. He includes the remaining terms of the expansions
\cite[(3.47)]{B-CMP119} in the boundary term $B^{(k+1)}$ \cite[\S3]{B-CMP119}. This assignment
reads $\Lambda_{k+1}$ only on the $MR_{k+1}$-cubes around the points $z\in Y$. The functions
themselves are those of the whole lattice. They do not depend on $\Lambda_{k+1}$
\cite[\S3]{B-CMP119}, by the identification lemma for localized pieces. By the statement
identifying an old boundary term with the new localized term, for a point $z'$ and a level $i$ with
$z'\notin\Lambda_i^{0}$, the term of level $i$ at $z'$, which is then a boundary term, is the same
function as the $E$-term of level $k+1$ filed at $z'$.

Finally let $k''\ge k+2$. The lemma on the native sequence at the $\mathbf{T}$-stage creating
$\Lambda_{k''}$ gives the nesting $\Lambda_{i+1}\subset\Lambda_i$ and the distance bound of
\eqref{4.29:eq-label-locality-proof-2} for all $i\le k''-1$, in particular at $i=k+1$. Let
$q\in\Lambda_{k+1}^{c}\cap C$. As above, the ball $B_\infty(q,8MR_{k+2}\ell_{k+2})$ misses
$\Omega^{\natural}_{k+2}\supset\Lambda_{k+2}\supset\Lambda_{k''}$, is connected and contains
$q$, hence lies in $C$. By the inequality $LR_{k+2}\ge R_{k+1}$ of the statement fixing the radii
and $\ell_{k+2}=L\ell_{k+1}$, we have
$8MR_{k+2}\ell_{k+2}\ge8MR_{k+1}\ell_{k+1}$. So $\Lambda_{k+1}^{c}\cap C$ lies at depth at least
$8MR_{k+1}\ell_{k+1}$ inside $C$, beyond the $MR_{k+1}$-collar.
\end{proof}

\begin{remark}[The native scale compared with the standard window cube]\label{4.29:rem-native-scale-window}
Work under the hypotheses of Proposition~\ref{4.29:prop-native-cover}, with the native scale
$\mathfrak{s}$, the generation-$j$ scale $\mathfrak{s}_j=\max\{j,\mathfrak{s}\}$, the sets
$\Lambda_j^{\sim}$ and the tiles of Definition~\ref{4.29:def-tiles-cover}. The native index
$l(x)$, which enters Proposition~\ref{4.29:prop-native-cover}, and the native strata are those of
the geometry of arrested components (Notation~\ref{0.2:not-glossary}); the native stratum of index
$l$ is the set of points $x$
with $l(x)=l$. For $j\le k+1$, $\Omega^{\natural}_j$ denotes the region of generation $j$ of that
geometry; for $j\le k$ it contains $\Lambda_j^{\sim}$, an inclusion recorded with
Definition~\ref{4.29:def-tiles-cover}. The
abandoned onions, with their thin shells $\Gamma''_n$ ($k_0<n<j_W$) and their lifted stratum
$j_W$, the datum $W$, the sets $S_j$ (not the sum $S$ of Definition~\ref{4.29:def-tiles-cover})
and the sets $\Gamma^{\natural}_j$, each contained in the native stratum of index $j$, are those of
the geometry
of arrested components (Notation~\ref{0.2:not-glossary}); the abandoned onions and the data $W$,
$k_0$, $j_W$ and $\Gamma''_n$ are
introduced with the definition of the analyticity spaces. For $1\le l\le k$ let
$\Lambda_l^{(-6)}$ be $\Lambda_l$ with six layers of $\pi_l$-cubes removed, that is, the union of
the cubes $\Delta\in\pi_l$ with $\tilde\Delta^{6}\subset\Lambda_l$; these are exactly the good cubes
of $\pi_l$. Then the following hold.
\begin{enumerate}
\item On $\Lambda_l^{(-6)}\setminus\Omega^{\natural}_{l+1}$ one has $\mathfrak{s}=l$; for every
generation $j\le l$ the tile of $\mathcal{T}_j$ at such a point $x$ is $\Delta_l(x)$, the standard
window cube of scale $l$ of the cover of \cite{B-CMP109}.
\item On the points $x\in\Lambda_1^{\sim}$ of the native stratum of index $l$ that lie outside
$\Lambda_l^{(-6)}$, one has $\mathfrak{s}(x)=l-1$. At such a point the standard window cube
$\Delta_l(x)$ satisfies $\widetilde{\Delta_l(x)}{}^{6}\not\subset\Lambda_l$,
whereas for $l\ge2$ the cube $\Delta_{l-1}(x)$ read by the native scale is good, so
$\widetilde{\Delta_{l-1}(x)}{}^{6}\subset\Lambda_{l-1}$.
\item On an abandoned onion, that is, on its thin shells $\Gamma''_n$, $k_0<n<j_W$, and on its
lifted stratum $j_W$, the index under which a point is entered in the geometry of arrested
components of Notation~\ref{0.2:not-glossary} is the index $n$ of the shell on $\Gamma''_n$ and
$j_W$ on the lifted stratum, whereas
the native scale reads the native index $l(x)$ of the point: $\mathfrak{s}(x)\in\{l(x)-1,l(x)\}$,
with
$k_0\le l(x)\le n$ on $\Gamma''_n$ and $k_0\le l(x)\le j_W$ on the lifted stratum, at every point
$x$ of the onion that lies in $\Lambda_1^{\sim}$.
\item On $S_j$, at the points where $\mathfrak{s}_j$ is defined, that is, on
$S_j\cap\Lambda_j^{\sim}$, one has $\mathfrak{s}_j=j=l(x)$ exactly.
\item The data $W$, $k_0$, $j_W$ do not enter the definition of $\mathfrak{s}$ or of
$\mathfrak{s}_j$.
\item On the ring $\Lambda_j^{\sim}\setminus\Lambda_j$ one has $\mathfrak{s}_j=j=l(x)$, and the
tile is default. Each ring cube, that is, each cube of $\pi_j$ that meets
$\Lambda_j^{\sim}\setminus\Lambda_j$, is a default tile of $\mathcal{T}_j$ and is not a tile of
$\mathcal{T}_{j'}$ for any generation $j'\neq j$.
\end{enumerate}
\end{remark}

\begin{proof}
We use two inclusions recorded with Definition~\ref{4.29:def-tiles-cover}:
$\Lambda_s\subset\Lambda_j$ for $s\ge j$, and
$\Lambda_j\subset\Lambda_j^{\sim}\subset\Omega^{\natural}_j$. We also use two clauses of
Proposition~\ref{4.29:prop-native-cover}: clause (a), by which the good scales at $x$ form the
initial segment $[1,\mathfrak{s}(x)]$, and the clause that bounds the loss to one native scale: for
every $x\in\Lambda_j^{\sim}$,
\begin{equation}\label{4.29:eq-native-scale-window-1}
\max\{j,\,l(x)-1\}\le\mathfrak{s}_j(x)\le\min\{l(x),k\},
\end{equation}
and a tile of generation $j$ is default only at points where $l(x)=j$.
Let $x$ be a point with $\mathfrak{s}(x)\ge1$, and let $j$ be a generation with
$1\le j\le\mathfrak{s}(x)$. The cube $\Delta_{\mathfrak{s}(x)}(x)$ is good, so
\begin{equation*}
x\in\Delta_{\mathfrak{s}(x)}(x)\subset\Lambda_{\mathfrak{s}(x)}\subset\Lambda_j
\subset\Lambda_j^{\sim},
\end{equation*}
and $\mathfrak{s}_j(x)=\max\{j,\mathfrak{s}(x)\}=\mathfrak{s}(x)$. Hence
\eqref{4.29:eq-native-scale-window-1} at generation $j$ gives
$l(x)-1\le\mathfrak{s}(x)\le l(x)$. Now let $x\in\Lambda_1^{\sim}$ with $\mathfrak{s}(x)=0$. Then
$\mathfrak{s}_1(x)=\max\{1,0\}=1$, and the tile of generation $1$ at $x$ is $\Delta_1(x)$. For
$y\in\Delta_1(x)$ one has $\Delta_1(y)=\Delta_1(x)$, which is not good since $\mathfrak{s}(x)=0$;
by clause (a) the good scales at $y$ form the initial segment $[1,\mathfrak{s}(y)]$, which
therefore does not contain $1$, so $\mathfrak{s}(y)=0<1$. The tile $\Delta_1(x)$ is thus default in
the sense of Definition~\ref{4.29:def-tiles-cover}, and by the loss clause $l(x)=1$, so
$\mathfrak{s}(x)=0=l(x)-1$. We have shown
\begin{equation}\label{4.29:eq-native-scale-window-2}
l(x)-1\le\mathfrak{s}(x)\le l(x)\qquad\text{for every }x\in\Lambda_1^{\sim}.
\end{equation}

\emph{First assertion.} Let $x\in\Lambda_l\setminus\Omega^{\natural}_{l+1}$, and
suppose that $\Delta_s(x)$ is good for some $s\ge l+1$ (for $l=k$ there is no such $s$, since
scales do not exceed $k$). Then
\begin{equation*}
x\in\Delta_s(x)\subset\widetilde{\Delta_s(x)}{}^{6}\subset\Lambda_s\subset\Lambda_{l+1}
\subset\Omega^{\natural}_{l+1},
\end{equation*}
which is a contradiction. Hence $\mathfrak{s}(x)\le l$.
If moreover $x\in\Lambda_l^{(-6)}$, then $\Delta_l(x)$ is good by the definition of
$\Lambda_l^{(-6)}$, so $\mathfrak{s}(x)=l$. For every generation $j\le l$ one has
\begin{equation*}
x\in\Lambda_l\subset\Lambda_j\subset\Lambda_j^{\sim}
\end{equation*}
and $\mathfrak{s}_j(x)=\max\{j,l\}=l$, so
the tile of $\mathcal{T}_j$ at $x$ is $\Delta_l(x)$.

\emph{Second assertion.} Let $x\in\Lambda_1^{\sim}$ lie in the native stratum of index $l$, so
that $l(x)=l$, and outside $\Lambda_l^{(-6)}$. Then $\Delta_l(x)$ is not good, so
$\widetilde{\Delta_l(x)}{}^{6}\not\subset\Lambda_l$, and by clause (a) $l\notin[1,\mathfrak{s}(x)]$,
that is, $\mathfrak{s}(x)\le l-1$. Since $x\in\Lambda_1^{\sim}$,
\eqref{4.29:eq-native-scale-window-2} gives $\mathfrak{s}(x)\ge l(x)-1=l-1$. Hence
$\mathfrak{s}(x)=l-1$, and when $l\ge2$ the cube
$\Delta_{l-1}(x)=\Delta_{\mathfrak{s}(x)}(x)$ is good by the definition of $\mathfrak{s}$, so
$\widetilde{\Delta_{l-1}(x)}{}^{6}\subset\Lambda_{l-1}$.

\emph{Third assertion.} On the thin shells $\Gamma''_n$ the native index satisfies
$k_0\le l(x)\le n$, and on the lifted stratum $k_0\le l(x)\le j_W$, by the geometry of arrested
components (Notation~\ref{0.2:not-glossary}). The bounds
\eqref{4.29:eq-native-scale-window-2} involve only the native index $l(x)$. They do not involve the
index $n$ or $j_W$ under which the point is entered in that geometry. Hence
$\mathfrak{s}(x)\in\{l(x)-1,l(x)\}$ at every point of the onion that lies in
$\Lambda_1^{\sim}$.

\emph{Fourth assertion.} By the clause of the geometry of arrested components
(Notation~\ref{0.2:not-glossary}) that places $S_j$
inside $\Gamma_j^{\natural}$, one has $S_j\subset\Gamma_j^{\natural}$, and $\Gamma_j^{\natural}$ is
contained in the native stratum of index $j$, so $l(x)=j$ on $S_j$. Let
$x\in S_j\cap\Lambda_j^{\sim}$. Since $j\le k$,
\eqref{4.29:eq-native-scale-window-1} at $x$ reads
\begin{equation*}
j=\max\{j,\,j-1\}\le\mathfrak{s}_j(x)\le\min\{j,k\}=j .
\end{equation*}

\emph{Fifth assertion.} By Definition~\ref{4.29:def-tiles-cover}, $\mathfrak{s}$ is determined by
the partitions $\pi_s$ and the sets $\Lambda_s$ alone, and $\mathfrak{s}_j$ is determined by
$\mathfrak{s}$ and $j$. Neither definition refers to $W$, $k_0$ or $j_W$.

\emph{Sixth assertion.} Let $x\in\Lambda_j^{\sim}\setminus\Lambda_j$, and suppose that
$\Delta_s(x)$ is good for some $s\ge j$. Then
\begin{equation*}
x\in\Delta_s(x)\subset\Lambda_s\subset\Lambda_j ,
\end{equation*}
which is impossible. Hence $\mathfrak{s}(x)<j$ and $\mathfrak{s}_j(x)=j$. By
Proposition~\ref{4.29:prop-native-cover}, on the ring one has $\mathfrak{s}_j=j=l(x)$ and the tile
is default. By the same proposition, at most one generation, namely $j=l(x)$, has a default tile
at $x$. A ring cube is thus the default tile of its own generation $j$; it is not a tile of any
other generation $j'$, as follows. A tile of generation $j'$ containing a point $y$ has
scale $\mathfrak{s}_{j'}(y)\ge j'$, so $j'>j$ is excluded. If $j'<j$ and
$\mathfrak{s}_{j'}(y)=\max\{j',\mathfrak{s}(y)\}=j$, then $\mathfrak{s}(y)=j$, so the ring cube
$\Delta_j(y)$ is good and hence contained in $\Lambda_j$; this is impossible, since it contains a
point of $\Lambda_j^{\sim}\setminus\Lambda_j$.
\end{proof}

\begin{lemma}[Exactness of the cube-by-cube localisation on the collared domain]
\label{4.29:lem-localisation-exact}
Work at level $k$, with $\eta=L^{-k}$, $\ell_n=L^n\eta$, the partitions $\pi_n$ and a label at the
$\mathbf{T}$-stage as in Notation~\ref{4.29:not-scales-labels}. That notation supplies the current
sequence $\{\Omega_i\}_{i\le k+1}$, the differences $\Gamma_a=\Omega_a\setminus\Omega_{a+1}$, the level
$n(x)$ of a site $x$, the native sequence $\{\Omega_i^\natural\}$ and the small-field regions
$\{\Lambda_s\}_{s\le k}$. Fix a generation $j\le k$. Let $\Lambda_j^{\sim}$, the native scale
$\mathfrak{s}_j$, the cover $\mathcal{C}_j$ and the functions $\bar\zeta_\square$,
$\square\in\mathcal{C}_j$, be those of Definition~\ref{4.29:def-tiles-cover}.
The superscript $\sim$ on $\Lambda_j$ is that of Definition~\ref{4.29:def-tiles-cover}, one layer of
$\pi_j$-cubes; where the operation $\sim$ is taken for a different family of cubes, this is said at
the place.

\emph{The collared minimiser.} Take $\Omega:=\Lambda_j^{\sim}$ in \cite[(2.13)]{B-CMP119}, read on the
lattice of spacing $\eta$ with $\ell_n$ in place of $L^n\xi$. Let
$D_0=\Lambda_j^{\sim}\supset D_1\supset\dots\supset D_j$ be the sequence of maximal domains given
there. Each $D_n$ is a union of $\ell_nM_1$-cubes, with
\[
\operatorname{dist}(D_n,D_{n-1}^c)\ge \ell_nM_1 .
\]
The last domain $D_j$ satisfies the condition of \cite[(2.13)]{B-CMP119}, written there as
$\operatorname{dist}(\Omega_j,\Omega^c)\le 2M_1$, or $\Omega^{\sim-2}\subset\Omega_j$ with $\sim$ taken
for $M_1$-cubes. Let $\mathbf{B}_j(\Lambda_j^{\sim})$ be the set determined by this sequence. It has
level $n$ on $D_n\setminus D_{n+1}$ for $n<j$, and level $j$ on $D_j$. Write $M_{\mathbf{B}}$ for the
averaging map of a set $\mathbf{B}$. We put
\begin{equation}\label{4.29:eq-localisation-exact-a}
\begin{gathered}
\mathbf{B}':=\mathbf{B}_j(\Lambda_j^{\sim}),\qquad M':=M_{\mathbf{B}'},\qquad
U_j:=U_{j,\Lambda_j^{\sim}}=U(\mathbf{B}_j(\Lambda_j^{\sim}),\cdot\,),\\
M'(e^{i\eta\mathbf{a}}U)=e^{iQ'(\eta\mathbf{a})}M'(U).
\end{gathered}
\end{equation}
Here $U(\mathbf{B}',V)=U_{j,\Lambda_j^{\sim}}(V)$ is the minimal configuration of
\cite[(2.13)]{B-CMP119} at the data $V$. The second line defines $Q'$, the averaging operation induced
by $M'$ on Lie-algebra valued fields $\mathbf{a}$, as in \cite[(97)]{B-CMP98}. The relation is exact
and in general non-linear in $\mathbf{a}$.

\emph{The two endpoints.} Let $B'$ be a real field as in \cite[(3.15)]{B-CMP119}, and let
$\mathbf{H}_k=\mathbf{H}_k(B')$ be the field of \cite[(3.3)]{B-CMP109}. Put
\[
U^0:=U_{k+1},\qquad U^1:=U_k(e^{iB'}V^{(k)}),\qquad W_0:=M'(U^0),
\]
the two configurations compared in \cite[(3.3)]{B-CMP109}. Let $F$ be a term of one of the following
two families of localised functionals of \cite{B-CMP119}:
\begin{itemize}
\item $F(U)=\mathbf{E}^{(j)}(X,U,z)$ with $X\subset\Lambda_j$, a term of the family of
\cite[(2.27)]{B-CMP119};
\item $F(U)=\mathbf{R}^{(j)}(X,U,z)$ with $X\subset\Lambda_j^{\sim-1}$, a term of the second family of
\cite{B-CMP119}, for which properties (i)--(iii) of \cite[(2.27)]{B-CMP119} are stated there; the
letter $\mathbf{R}^{(j)}$ is that of \cite{B-CMP119} and does not denote the large-field operation,
and the operation $\sim$ in this exponent is that of \cite{B-CMP119}, taken in the cubes in terms of
which $\Lambda_j$ is defined.
\end{itemize}
In both cases $z$ denotes the remaining arguments, as in \cite[(2.27)]{B-CMP119}. Put
\begin{equation}\label{4.29:eq-localisation-exact-1}
f(t):=F\bigl(U_{j,\Lambda_j^{\sim}}(e^{iQ'(\eta t\mathbf{H}_k)}W_0)\bigr),\qquad t\in[0,1].
\end{equation}

Assume the following.
\begin{enumerate}
\item[(1)] \emph{Restriction of minimisers.} Restriction of minimisers holds for the domain
$\Lambda_j^{\sim}$ and the set $\mathbf{B}'$ in the following form. Let
$U^*$ be a minimiser of \cite[Theorem~1]{B-CMP102b} for a set whose levels on $\Lambda_j^{\sim}$
dominate the levels of $\mathbf{B}'$, with all partitions compatible. Suppose that the minimal orbit of
the variational problem for $\mathbf{B}'$ at the data $M'(U^*)$ is unique. Then
\[
U(\mathbf{B}',M'(U^*))=(U^*\lceil_{\Lambda_j^{\sim}})^{u}
\]
for a gauge transformation $u$, the restriction being to $\Lambda_j^{\sim}$ as a set of bonds, so that
blocks of $\mathbf{B}'$ straddling the boundary of $\Lambda_j^{\sim}$ are included in it.
\item[(2)] \emph{Uniqueness of the minimal orbit.} For $\iota=0,1$, the minimal orbit of the
variational problem for $\mathbf{B}'$ at the data $M'(U^\iota)$ is unique.
\item[(3)] \emph{The cover.} The hypotheses of Proposition~\ref{4.29:prop-native-cover} hold, and hence
clause (a) of Proposition~\ref{4.29:prop-native-cover}.
\item[(4)] \emph{Size of the induced source.} For every bond $c$ of $\mathbf{B}'$ and every
$t\in[0,1]$,
\[
|Q'(\eta t\mathbf{H}_k)(c)|\le c_Q\,\ell_j\sup|\mathbf{H}_k|,
\]
with the constant $c_Q$ of the source bound \eqref{4.29:eq-collar-staircase-b}; and
$c_Q\ell_j\sup|\mathbf{H}_k|$ does not exceed the radius of the neighbourhood of $W_0$, in the
supremum norm over the bonds of $\mathbf{B}'$, on which \cite[\S G, Proposition~9]{B-CMP102b} applies.
\end{enumerate}
Then the following hold.
\begin{enumerate}
\item[(a)] For $\iota=0,1$, $U^\iota\lceil_{\Lambda_j^{\sim}}=U_{j,\Lambda_j^{\sim}}(M'(U^\iota))$ up to
gauge. Moreover $M'(U^1)=e^{iQ'(\eta\mathbf{H}_k)}W_0$ up to gauge.
\item[(b)] $F(U^1)-F(U^0)=f(1)-f(0)$; that is, \cite[(3.3)]{B-CMP109} holds with
$U_j=U_{j,\Lambda_j^{\sim}}$.
\item[(c)] $f$ is of class $C^1$ on $[0,1]$. Moreover
\begin{equation}\label{4.29:eq-localisation-exact-2}
\begin{split}
F(U^1)-F(U^0)
=\sum_{\square\in\mathcal{C}_j}\int_0^1 dt\,\frac{\partial}{\partial t_\square}
F\Bigl(U_{j,\Lambda_j^{\sim}}\bigl(e^{iQ'(\eta(t+t_\square\bar\zeta_\square)\mathbf{H}_k)}W_0\bigr)
\Bigr)\Big|_{t_\square=0}.
\end{split}
\end{equation}
That is, the expansion \cite[(3.4)]{B-CMP109}, written with the cover $\mathcal{C}_j$, the
functions $\bar\zeta_\square$ and the minimiser $U_{j,\Lambda_j^{\sim}}$, is exact.
\item[(d)] Every cube of $\mathcal{C}_j$ has scale $\mathfrak{s}_j\in[j,k]$.
\end{enumerate}
\end{lemma}

\begin{proof}
\emph{The expansion to be made exact.} The expansion \cite[(3.4)]{B-CMP109}, on p.~270 of that paper,
has the form
\[
\sum_{\square}\int_0^1 dt\,\frac{d}{dt_\square}\,
\mathbf{E}^{(j)}\Bigl(U_j\bigl(\exp iQ_j(\eta(t+t_\square\zeta_\square)\mathbf{H}_k(B'))
\,\bar U^{j}_{k+1}\bigr)\Bigr)\Big|_{t_\square=0},
\]
$Q_j$ and $\bar U^{j}_{k+1}$ being the averaging operation and the averaged configuration at level $j$
denoted so in \cite[(3.4)]{B-CMP109}.
It has one global interpolation parameter $t$, and the partition of unity enters linearly through
$t_\square\zeta_\square$. The argument is a minimiser $U_j(\cdot)$ at level $j$, which depends on
$\mathbf{H}_k$ through $Q_j(\eta\mathbf{H}_k)$. The statement replaces $U_j$ by the minimiser
$U_{j,\Lambda_j^{\sim}}$ of \eqref{4.29:eq-localisation-exact-a}, and the family of cubes and functions
by $(\mathcal{C}_j,\bar\zeta_\square)$. We prove (a)--(d) in this setting, in seven steps.

\emph{Step 1: the endpoints as minimisers.} By \cite[\S G, Proposition~9]{B-CMP102b} and
\cite{B-CMP119},
\[
U^1=U_k(e^{iB'}V^{(k)})=\bigl(e^{i\eta\mathbf{H}_k(B')}U^0\bigr)^{u}
\]
for a gauge transformation $u$. Both configurations are minimisers in the sense of
\cite[Theorem~1]{B-CMP102b}, as stated in \cite{B-CMP119}:
\begin{itemize}
\item $U^0$ is the minimiser for the set $\mathbb{B}=B(\{\Omega_i\}_{i\le k+1})$ with data
$\{V_{k+1}\lceil\Gamma_{k+1},V_k\lceil\Gamma_k,\dots\}$;
\item $U^1$ is the minimiser for the set of level $k$ on $\Omega_{k+1}$ with the shifted data.
\end{itemize}
Here $B'$ is real, as in \cite[(3.15)]{B-CMP119}. Both sides of \cite[(3.3)]{B-CMP109} are analytic
in $B'$ by \cite[Proposition~9, (181)]{B-CMP102b}. By Definition~\ref{4.29:def-tiles-cover},
$\Lambda_j^{\sim}\subset\Omega_j^\natural\subset\Omega_j$, and on $\Lambda_j^{\sim}$ both sets have levels
$n(x)\ge j$. The levels of $\mathbf{B}'$ are at most $j$: they are $n<j$ on $D_n\setminus D_{n+1}$ and
$j$ on $D_j$. Hence on $\Lambda_j^{\sim}$ the levels of both sets dominate every level of
$\mathbf{B}'$. All partitions involved are compatible.

\emph{Step 2: restriction to the collared domain; proof of (a).} Let $\iota\in\{0,1\}$. By Step 1,
$U^\iota$ is a minimiser of \cite[Theorem~1]{B-CMP102b} for a set whose levels on $\Lambda_j^{\sim}$
dominate those of $\mathbf{B}'$, with compatible partitions. By hypothesis (2), the minimal orbit of
the variational problem for $\mathbf{B}'$ at the data $M'(U^\iota)$ is unique. Hypothesis (1) therefore
applies with $U^*=U^\iota$, and the definition \eqref{4.29:eq-localisation-exact-a} gives
\[
U^\iota\lceil_{\Lambda_j^{\sim}}=U_{j,\Lambda_j^{\sim}}(M'(U^\iota))\quad\text{up to gauge},\qquad
\iota=0,1.
\]
For the data we use the covariance of averages \cite[(98), (181)]{B-CMP102b}. Applied to the identity
$U^1=(e^{i\eta\mathbf{H}_k}U^0)^u$ of Step 1, it shows that $M'(U^1)$ coincides, up to gauge, with
$M'(e^{i\eta\mathbf{H}_k}U^0)$. By the second line of \eqref{4.29:eq-localisation-exact-a} with
$\mathbf{a}=\mathbf{H}_k$, the latter equals $e^{iQ'(\eta\mathbf{H}_k)}M'(U^0)$, that is,
$e^{iQ'(\eta\mathbf{H}_k)}W_0$. This proves (a).

\emph{Step 3: the difference as an increment of $f$; proof of (b).} Consider first
$F(U)=\mathbf{E}^{(j)}(X,U,z)$ with $X\subset\Lambda_j$. By \cite[(2.27)(i),(iii)]{B-CMP119}, it depends
on its configuration argument only through the restriction to $X$, and it is gauge invariant.
Consider next $F(U)=\mathbf{R}^{(j)}(X,U,z)$ with $X\subset\Lambda_j^{\sim-1}$. The same properties
(i)--(iii) hold by \cite{B-CMP119}, statement following (2.27). In both cases
$X\subset\Lambda_j\subset\Lambda_j^{\sim}$. So $F(U^\iota)$ is unchanged if $U^\iota$ is replaced by any
configuration that agrees with it on $\Lambda_j^{\sim}$ up to gauge.

Apply this with $\iota=0$. By (a), $U^0$ agrees on $\Lambda_j^{\sim}$, up to gauge, with
$U_{j,\Lambda_j^{\sim}}(W_0)$. Hence $F(U^0)=f(0)$. Apply it with $\iota=1$. By (a), $U^1$ agrees on
$\Lambda_j^{\sim}$, up to gauge, with $U_{j,\Lambda_j^{\sim}}(M'(U^1))$. The data $M'(U^1)$ equal
$e^{iQ'(\eta\mathbf{H}_k)}W_0$ up to gauge. Hence $F(U^1)=f(1)$. Therefore
\[
\mathbf{E}^{(j)}(X,U^1,z)-\mathbf{E}^{(j)}(X,U^0,z)=f(1)-f(0)
\]
in the first case, and the same identity holds for $\mathbf{R}^{(j)}$ in the second. This is
\cite[(3.3)]{B-CMP109} with $U_j=U_{j,\Lambda_j^{\sim}}$, which proves (b).

\emph{Step 4: regularity of $f$, and one linear map for both derivatives.} By hypothesis (4),
\[
|Q'(\eta t\mathbf{H}_k)(c)|\le c_Q\ell_j\sup|\mathbf{H}_k|
\]
for every bond $c$ of $\mathbf{B}'$ and every $t\in[0,1]$, and by the same hypothesis this quantity
does not exceed the radius of the neighbourhood of $W_0$ on which \cite[\S G, Proposition~9]{B-CMP102b}
applies. The data $e^{iQ'(\eta t\mathbf{H}_k)}W_0$ therefore stay in that neighbourhood of $W_0$. So
$t\mapsto U_{j,\Lambda_j^{\sim}}(e^{iQ'(\eta t\mathbf{H}_k)}W_0)$
is of class $C^1$, and hence so is $f$.

For a field $\mathbf{b}$ on the bonds of $\mathbf{B}'$ (the space in which $Q'$ takes values), write
\[
G(\mathbf{b}):=F\bigl(U_{j,\Lambda_j^{\sim}}(e^{i\mathbf{b}}W_0)\bigr).
\]
Write $DQ'|_{\eta t\mathbf{H}_k}$ for the differential of $Q'$ at $\eta t\mathbf{H}_k$. Only this
differential is used below, not the non-linear map $Q'$ itself. Then $f(t)=G(Q'(\eta t\mathbf{H}_k))$,
and the chain rule gives
\[
f'(t)=DG|_{Q'(\eta t\mathbf{H}_k)}\bigl[DQ'|_{\eta t\mathbf{H}_k}[\eta\mathbf{H}_k]\bigr].
\]
For $\square\in\mathcal{C}_j$, the insertion in \eqref{4.29:eq-localisation-exact-2} is
$G\bigl(Q'(\eta(t+t_\square\bar\zeta_\square)\mathbf{H}_k)\bigr)$. Differentiating in $t_\square$ at
$t_\square=0$ gives
\[
\frac{\partial}{\partial t_\square}G\bigl(Q'(\eta(t+t_\square\bar\zeta_\square)\mathbf{H}_k)\bigr)
\Big|_{t_\square=0}
=DG|_{Q'(\eta t\mathbf{H}_k)}\bigl[DQ'|_{\eta t\mathbf{H}_k}[\eta\bar\zeta_\square\mathbf{H}_k]\bigr].
\]
Thus $f'(t)$ and each insertion derivative at $t_\square=0$ are values of one and the same linear
map, $DG|_{Q'(\eta t\mathbf{H}_k)}\circ DQ'|_{\eta t\mathbf{H}_k}$. The map is applied to
$\eta\mathbf{H}_k$ in the first case and to $\eta\bar\zeta_\square\mathbf{H}_k$ in the second.

\emph{Step 5: locality of $Q'$ and the partition of unity.} By \cite{B-CMP98}, the average
$(\bar U^k)_c$ over a bond $c$ depends only on the variables $U_b$ with
$b\subset B^k(c_-)\cup B^k(c_+)$, the blocks of the two end-points of $c$. Through the defining
relation \eqref{4.29:eq-localisation-exact-a}, $Q'(\eta\mathbf{a})(c)$ reads $\mathbf{a}$ only on the two
blocks of $\mathbf{B}'$ containing the end-points of $c$. Consequently, for every bond $c$ of
$\mathbf{B}'$, the value $DQ'|_{\eta t\mathbf{H}_k}[\mathbf{a}](c)$ depends only on the restriction of
$\mathbf{a}$ to these two blocks.

The blocks of $\mathbf{B}'$ lie inside $\Lambda_j^{\sim}$ when $\Lambda_j^{\sim}$ is taken as a set of
bonds, namely the bonds at least one end-point of which belongs to $\Lambda_j^{\sim}$
\cite{B-CMP99a}. Every such bond lies within distance $\eta$ of $\Lambda_j^{\sim}$. By
clause (a) of Proposition~\ref{4.29:prop-native-cover}, available by hypothesis (3), $\bar\zeta$ is a
partition of unity on the $\tfrac13 M\ell_j$-neighbourhood of $\Lambda_j^{\sim}$, subordinate to
$\mathcal{C}_j$. In particular
\[
\sum_{\square\in\mathcal{C}_j}\bar\zeta_\square=1
\quad\text{on the $\tfrac13 M\ell_j$-neighbourhood of $\Lambda_j^{\sim}$.}
\]
This neighbourhood contains the $\eta$-neighbourhood of $\Lambda_j^{\sim}$. Indeed, $M\ge L^4>3$ among
the hypotheses of that proposition, and $\ell_j=L^j\eta\ge\eta$, so that
$\eta\le\ell_j<\tfrac13 M\ell_j$.

The fields $\eta\bigl(\sum_{\square\in\mathcal{C}_j}\bar\zeta_\square\bigr)\mathbf{H}_k$ and
$\eta\mathbf{H}_k$ therefore coincide on every block of $\mathbf{B}'$. Now
$DQ'|_{\eta t\mathbf{H}_k}$ is linear, and $\mathcal{C}_j$ is a finite family, since the lattice is
finite. Hence
\[
\sum_{\square\in\mathcal{C}_j}DQ'|_{\eta t\mathbf{H}_k}[\eta\bar\zeta_\square\mathbf{H}_k]
=DQ'|_{\eta t\mathbf{H}_k}\Bigl[\eta\Bigl(\sum_{\square\in\mathcal{C}_j}\bar\zeta_\square\Bigr)
\mathbf{H}_k\Bigr]
=DQ'|_{\eta t\mathbf{H}_k}[\eta\mathbf{H}_k].
\]

\emph{Step 6: exactness; proof of (c).} Apply the linear map
$DG|_{Q'(\eta t\mathbf{H}_k)}$ of Step 4 to the identity of Step 5. Together with the two formulas of
Step 4, this gives, for every $t\in[0,1]$,
\[
\sum_{\square\in\mathcal{C}_j}\frac{\partial}{\partial t_\square}
G\bigl(Q'(\eta(t+t_\square\bar\zeta_\square)\mathbf{H}_k)\bigr)\Big|_{t_\square=0}=f'(t).
\]
Since $f$ is of class $C^1$ by Step 4, $f(1)-f(0)=\int_0^1f'(t)\,dt$. Inserting the last identity and
exchanging the finite sum with the integral, (b) yields \eqref{4.29:eq-localisation-exact-2}. This
proves (c).

The argument shows that the listed conditions suffice for exactness; only this sufficiency is proved
here, and only this sufficiency is used. The sum in \eqref{4.29:eq-localisation-exact-2} runs over
$\mathcal{C}_j$ alone. No cube whose function $\zeta_\square$ is supported far from $\Lambda_j^{\sim}$
occurs, since $\mathcal{C}_j$ contains none. The sum of all the pieces is the left-hand side of
\cite[(3.3)]{B-CMP109}; compared with the expansion \cite[(3.4)]{B-CMP109}, only the grouping of the
far field changes. In the case distinction of \cite[\S3]{B-CMP109}, the first case,
$X\cap\tilde\square=\emptyset$, arises for cubes of $\mathcal{C}_j$ inside $\Lambda_j^{\sim}$.

\emph{Step 7: scales of the cubes; proof of (d), and the reading of the localisation.} Every cube of
$\mathcal{C}_j$ has the scale of a tile generating it. By Definition~\ref{4.29:def-tiles-cover}, this
scale is $\mathfrak{s}_j(x)=\max\{j,\mathfrak{s}(x)\}$ for a site $x$ of the tile, and
$\mathfrak{s}(x)\le k$ while $j\le k$. Hence $\mathfrak{s}_j(x)\in[j,k]$, which is (d).

Consequently, each cube $\square$ of the expansion \eqref{4.29:eq-localisation-exact-2} is a union of
$2^d$ cubes of $\pi_n$ with $j\le n\le k$. For such a cube, the operation $\sim$ is defined using cubes
of $\pi_n$, as in \cite{B-CMP119}, and this holds here exactly. The restriction of the term
$\mathbf{E}^{(j)}$ to the domain $\Lambda_j$, described in \cite{B-CMP119}, holds here exactly with
$\Lambda_j^{\sim}$ in place of $\Lambda_j$. This is the one-layer collar inside which the localised
object is read in this book. In \cite[(1.16)]{B-CMP109} a localised object is read on $X^{\sim-2}$,
the domain obtained from $X$ by taking away two layers of cubes. The one-layer reading adopted here is
that of \cite{B-CMP119}: for $X\subset\square^{\sim n-j}$ one reads through
$\square_0=\square^{\sim n-j+1}$ and $U_k=U_{j,\square_0}(M_{\square_0}(U_k))$, with $M_{\square_0}$ the
averaging map of the set determined by $\square_0$.

By Definition~\ref{4.29:def-tiles-cover}, the cover $\mathcal{C}_j$ and the functions $\bar\zeta_\square$
are determined by the sequence $\{\Lambda_s\}_{s\le k}$ alone.
\end{proof}

\begin{remark}[Why the minimiser is taken on the collared domain]
Suppose the minimiser in \eqref{4.29:eq-localisation-exact-2} were instead one whose induced averaging
operation reads $\mathbf{H}_k$ on the whole lattice. Then the sum of the insertion derivatives would
differ from $f'(t)$ by the first variation of $\mathbf{E}^{(j)}$ with respect to $\mathbf{H}$, paired
with $(1-\sum_{\square}\bar\zeta_\square)\mathbf{H}_k$. This pairing is in general different from zero,
and the expansion would not be exact.
\end{remark}

Throughout, the objects of the generation-$j$ localisation are those of \eqref{4.29:eq-localisation-exact-a}. By the \emph{staircase} of $\mathbf{B}'$ we mean $D_0\setminus D_j$, the part of $\Lambda_j^{\sim}$ on which $\mathbf{B}'$ has levels below $j$. The twist bound used below is stated in the units of \cite[Proposition~6]{B-CMP98}. We write $\alpha_1^{(98)}$ for the radius $\alpha_1$ of that proposition. In \cite[\S3]{B-CMP109} and \cite{B-CMP102b} the fine spacing is written $\xi$; it is the spacing $\eta$ here.

\begin{lemma}[The collar staircase, decay at observation points and induced sources]
\label{4.29:lem-collar-staircase}
Fix the level $k$ and a generation $j\le k$. We use the following objects.
\begin{itemize}
\item From Notation~\ref{4.29:not-scales-labels}: the spacing $\eta=L^{-k}$, the scales $\ell_n=L^n\eta$, the mutually compatible partitions $\pi_n$ into cubes of side $M\ell_n$, and $\mathrm{dist}_\infty$.
\item From Definition~\ref{4.29:def-tiles-cover}: the collared domain $\Lambda_j^{\sim}$ and the good cubes.
\item From Proposition~\ref{4.29:prop-native-cover}: the cover $\mathcal{C}_j$, the partition $(\bar\zeta_\Box)_{\Box\in\mathcal{C}_j}$ and the function $S$.
\item From \eqref{4.29:eq-localisation-exact-a}:
 \begin{itemize}
 \item the set $\mathbf{B}'=\mathbf{B}_j(\Lambda_j^{\sim})$ of \cite[(2.13)]{B-CMP119};
 \item its maximal domains $\Lambda_j^{\sim}=D_0\supset D_1\supset\dots\supset D_j$, where $\mathbf{B}'$ has level $n$ on $D_n\setminus D_{n+1}$ and level $j$ on $D_j$;
 \item the averaging operation $M'=M_{\mathbf{B}'}$ and the induced operation $Q'$, with $M'(e^{i\eta\mathbf{a}}U)=e^{iQ'(\eta\mathbf{a})}M'(U)$;
 \item the configuration $U^0=U_{k+1}$, the data $W_0=M'(U^0)$, and the minimiser $U_{j,\Lambda_j^{\sim}}$.
 \end{itemize}
\end{itemize}
For an element $c$ of $\mathbf{B}'$ of level $m(c)\le j$, with end-points $c_\pm$ and the averaging blocks $B^{m(c)}(c_\pm)$ of \cite{B-CMP98}, of scale $\ell_{m(c)}$, containing them (not to be confused with the field $B$), put
\[
A(c):=B^{m(c)}(c_-)\cup B^{m(c)}(c_+).
\]
Assume that $L\ge3$, that $M>M_2=L^2M_1$ with $M_1\ge1$ the cube size of \cite{B-CMP119}, and assume the following three hypotheses.
\begin{itemize}
\item[(tw)] \emph{Twist bound for averages along an element.} Let $c$ be an element of level $s$ with end-block set $\tilde B(c)$, and let $M_{\mathbf{S}}$ be the averaging operation of a set $\mathbf{S}$ of elements containing $c$. Let $U_0$ satisfy \cite[(52)]{B-CMP98} at scale $\ell_s$ on $\tilde B(c)$. Let $\mathbb{U}=U'U_0$ with $A'=(i\eta)^{-1}\log U'$ and $\ell_s\sup_{\tilde B(c)}|A'|\le\alpha_1^{(98)}/4$. Let $\mathcal{X}$ be a twist with values in the complexified Lie algebra, with $\ell_s\sup_{\tilde B(c)}|\mathcal{X}|\le\alpha_1^{(98)}/4$. Then there is $\mathfrak{h}(c)$ in the complexified algebra with
\[
\begin{gathered}
M_{\mathbf{S}}(e^{i\eta\mathcal{X}}\mathbb{U})(c)=e^{i\mathfrak{h}(c)}M_{\mathbf{S}}(\mathbb{U})(c),\\
|e^{i\mathfrak{h}(c)}-1|\le c_6'\,\ell_s\sup_{\tilde B(c)}|\mathcal{X}|,
\qquad c_6':=16/\alpha_1^{(98)}.
\end{gathered}
\]
The map $\mathcal{X}\mapsto M_{\mathbf{S}}(e^{i\eta\mathcal{X}}\mathbb{U})(c)$ is analytic on the polydisc $\ell_s\sup_{\tilde B(c)}|\mathcal{X}|\le\alpha_1^{(98)}/4$. This is the linear form, in the twist, of the bound \cite[(164)]{B-CMP98}: that bound estimates the deviation of the twisted average from the untwisted one, $|\overline{U'U_0}^k(\bar U_0^k)^{-1}-1|$, by $O(1)\alpha_1$ on the polydisc $|A'|<\alpha_1$, with $A'$ in block units \cite[(162)]{B-CMP98}, and vanishes at zero twist, so that the Schwarz lemma in the twist gives the linear form.
\item[(dec)] \emph{The summed form of \cite[(190)]{B-CMP102b}.} Let $e\in\{0,1,2\}$, where for $e=2$ the operator $\nabla^2$ stands also for the covariant operators $D^*D$ and $\Delta^{\xi}_{\mathbf{U}}$ of \cite[(3.14)]{B-CMP109}, $\mathbf{U}$ being the background variable of \cite[(3.13)]{B-CMP109}. Let $B$ lie in the polydisc of \cite[Proposition~9]{B-CMP102b}. Then the bound \eqref{4.29:eq-collar-staircase-a} below holds at every point $x$ at which $\mathbf{B}'$ has level $j$, $x$ lying in the block $\Delta(y)$ of a point $y$ of the level-$j$ lattice (the lattice written $\Lambda_j$ in \cite[(190)]{B-CMP102b}), with constants $C_{90}$ and $\delta_0$ uniform in $k$. Here $\operatorname{d}$ is the scaled distance of \cite{B-CMP102b}, measured along contours made of bonds of each element's own lattice (cf.\ \cite[(2.46), (2.48)]{B-CMP96}).
\item[(reg)] \emph{Regularity of the reference configuration.} $U^0$ is regular at all levels $n\ge j$, in the form which yields \cite[(52)]{B-CMP98} at scale $\ell_m$ on $A(c)$ for every element $c$ of $\mathbf{B}'$ of level $m\le j$.
\end{itemize}
Then the following hold.
\begin{enumerate}
\item[(i)] \emph{Where the staircase is.}
 \begin{itemize}
 \item $D_j$ contains every cube $C$ of side $M_1\ell_j$ with $C\subset\Lambda_j^{\sim}$ and $\mathrm{dist}_\infty(C,(\Lambda_j^{\sim})^c)\ge2M_1\ell_j$.
 \item $\Lambda_j\subset D_j$.
 \item The elements of $\mathbf{B}'$ of levels $<j$ lie in the outer two layers of $M_1\ell_j$-cubes of the ring $\Lambda_j^{\sim}\setminus\Lambda_j$, at $\mathrm{dist}_\infty\ge(M-2M_1)\ell_j$ from $\Lambda_j$.
 \end{itemize}
\item[(ii)] \emph{Observation points.} Let
\[
H_j:=\mathcal{H}',
\]
where $\mathcal{H}'$ is the function of \cite[Proposition~9]{B-CMP102b} attached to $\mathbf{B}'$, $W_0$ and the external configuration $U^0$ restricted to $\Lambda_j^{\sim}$. Thus, for some gauge transformation $u$,
\begin{equation}\label{4.29:eq-collar-staircase-1}
U_{j,\Lambda_j^{\sim}}(e^{iB}W_0)=\bigl(e^{i\eta\mathcal{H}'(B)}U^0\bigr)^{u}
\quad\text{on }\Lambda_j^{\sim}.
\end{equation}
An \emph{observation point} is any of the following:
 \begin{itemize}
 \item a point of a domain $X$ of \cite[(2.27)]{B-CMP119} or \cite[(2.30)]{B-CMP119}, and the localisation point $z$ of \cite[(2.27)]{B-CMP119}, which lies in $\Lambda_j$;
 \item a plaquette entering the expressions \cite[(3.10)--(3.13)]{B-CMP109};
 \item a point of a block entering the operator $P_1=DPD^*$ of \cite[(3.14)]{B-CMP109};
 \item a point of the cube $\Box_0$ of \cite[p.~274]{B-CMP109} attached to a good cube $\Box$, as in clause~(c) of Proposition~\ref{4.29:prop-native-cover}.
 \end{itemize}
Every observation point lies in, or within $2\ell_j+\eta$ of, $\Lambda_j\subset D_j$. It is therefore a point at which $\mathbf{B}'$ has level $j$, at $\mathrm{dist}_\infty\ge(M-2M_1-3)\ell_j$ from every element of level $<j$. In particular, $U_{j,\Lambda_j^{\sim}}$ is never evaluated on its staircase.
\item[(iii)] \emph{Decay at observation points.} Let $x$ be an observation point, lying in the block $\Delta(y)$ of a point $y$ of the level-$j$ lattice, let $e\in\{0,1,2\}$, and let $B$ lie in the polydisc of \cite[Proposition~9]{B-CMP102b}. Then, uniformly in $k$,
\begin{equation}\label{4.29:eq-collar-staircase-a}
\ell_j^{1+e}\,|\nabla^e\mathcal{H}'(B)(x)|
\le C_{90}\sum_{c\in\mathbf{B}'}e^{-\delta_0 \operatorname{d}(y,c)/8}\,|B(c)|,
\end{equation}
where $\operatorname{d}$ is the scaled distance of \cite{B-CMP102b}. Moreover,
\[
P_1\mathcal{H}'=\Delta^{\xi}_{\mathbf{U}}\mathcal{H}'-D^*D\mathcal{H}'+(\text{local lower order}).
\]
\item[(iv)] \emph{Sources.} Let $c$ be an element of $\mathbf{B}'$ of level $m(c)\le j$. Let the configuration at which $Q'$ is evaluated be $U'U^0$ with $\ell_{m(c)}\sup_{A(c)}|A'|\le\alpha_1^{(98)}/4$. Let $\mathbf{a}$ be a twist with $64\,\ell_{m(c)}\sup_{A(c)}|\mathbf{a}|\le\alpha_1^{(98)}$; at the twists at which $Q'$ and $DQ'$ are evaluated, this is secured by \eqref{4.29:eq-constants-order-b}. Then
\begin{equation}\label{4.29:eq-collar-staircase-b}
|Q'(\eta\mathbf{a})(c)|\le c_Q\,\ell_{m(c)}\sup_{A(c)}|\mathbf{a}|,
\qquad c_Q:=2c_6',
\end{equation}
and, for every twist $\mathbf{b}$,
\begin{equation}\label{4.29:eq-collar-staircase-2}
\bigl|DQ'|_{\eta\mathbf{a}}[\eta\mathbf{b}](c)\bigr|
\le 2c_Q\,\ell_{m(c)}\sup_{A(c)}|\mathbf{b}|.
\end{equation}
Since $\ell_{m(c)}\le\ell_j$, a finer source is a smaller source.
\item[(v)] \emph{The cube insertions.} For $\Box\in\mathcal{C}_j$, the $\Box$-term of the expansion of Lemma~\ref{4.29:lem-localisation-exact} depends on $\bar\zeta_\Box$ only through $DQ'$, and only through its values near $\Box\cap\Lambda_j^{\sim}$. There $S\ge1$ and the gradient bounds of clause~(a) of Proposition~\ref{4.29:prop-native-cover} hold.
\end{enumerate}
\end{lemma}

\begin{proof}
The proof of (i) proceeds in four steps.

\emph{Step 1: the inclusion for cubes.} By \cite[(2.13)]{B-CMP119}, read on the lattice of spacing $\eta$ with $\ell_n$ in place of $L^n\xi$, we have $D_0=\Lambda_j^{\sim}$. For $1\le n\le j$, $D_n$ is the maximal union of cubes of side $M_1\ell_n$ such that
\[
\mathrm{dist}_\infty(D_n,D_{n-1}^c)\ge M_1\ell_n.
\]
By maximality, $D_n$ contains every $M_1\ell_n$-cube $C\subset D_{n-1}$ with $\mathrm{dist}_\infty(C,D_{n-1}^c)\ge M_1\ell_n$.

We prove by induction on $n\in\{0,\dots,j\}$ the following claim.
\begin{quote}
$(\mathrm{P}_n)$: every $M_1\ell_n$-cube $C\subset\Lambda_j^{\sim}$ with $\mathrm{dist}_\infty(C,(\Lambda_j^{\sim})^c)\ge2M_1\ell_n$ lies in $D_n$.
\end{quote}
For $n=0$ the claim holds because $D_0=\Lambda_j^{\sim}$.

Let $n\ge1$ and let $C$ be as in $(\mathrm{P}_n)$. Let $v$ be a point with $\mathrm{dist}_\infty(v,C)<M_1\ell_n$. Then
\[
\mathrm{dist}_\infty(v,(\Lambda_j^{\sim})^c)>M_1\ell_n,
\]
so $v\in\Lambda_j^{\sim}$. Let $C'$ be the cube of side $M_1\ell_{n-1}$ of the compatible partition that contains $v$. Then
\[
\mathrm{dist}_\infty\bigl(C',(\Lambda_j^{\sim})^c\bigr)
\ge 2M_1\ell_n-M_1\ell_n-M_1\ell_{n-1}
=(L-1)M_1\ell_{n-1}\ge2M_1\ell_{n-1},
\]
the last step because $L\ge3$. By $(\mathrm{P}_{n-1})$ we get $C'\subset D_{n-1}$. Hence the $M_1\ell_n$-neighbourhood of $C$ lies in $D_{n-1}$. This means $C\subset D_{n-1}$ and $\mathrm{dist}_\infty(C,D_{n-1}^c)\ge M_1\ell_n$, so $C\subset D_n$ by maximality. For $n=j$, $(\mathrm{P}_j)$ is the first assertion of (i). In the notation of \cite[(2.13)]{B-CMP119} it is the inclusion $\Omega^{\sim-2}\subset\Omega_j$, the operation $\sim$ being taken for $M_1$-cubes.

\emph{Step 2: the ring.} Let $q\in\Lambda_j$ and let $\mathsf{C}\in\pi_j$ be the cube containing $q$. Then $\mathsf{C}\cap\Lambda_j\ne\emptyset$, so $\mathsf{C}\subset\Lambda_j^{\sim}$ by the definition of $\Lambda_j^{\sim}$ (Definition~\ref{4.29:def-tiles-cover}). Every cube $\mathsf{C}'\in\pi_j$ adjacent to $\mathsf{C}$, corners included, has $\mathsf{C}\subset\tilde{\mathsf{C}}'^{1}$, hence $\tilde{\mathsf{C}}'^{1}\cap\Lambda_j\ne\emptyset$ and $\mathsf{C}'\subset\Lambda_j^{\sim}$. Thus $\tilde{\mathsf{C}}^{1}\subset\Lambda_j^{\sim}$. Since $\tilde{\mathsf{C}}^1$ extends one cube of side $M\ell_j$ beyond $\mathsf{C}$ in every direction,
\[
\mathrm{dist}_\infty\bigl(\Lambda_j,(\Lambda_j^{\sim})^c\bigr)\ge M\ell_j .
\]
So the ring $\Lambda_j^{\sim}\setminus\Lambda_j$ is one $\pi_j$-cube, of side $M\ell_j$, thick. In \cite{B-CMP119} the cube sizes satisfy $M_1<M_2<M$ with $M_2=L^2M_1$. By hypothesis $M>L^2M_1$, so the ring is at least $L^2M_1\ell_j$ thick.

\emph{Step 3: $\Lambda_j\subset D_j$.} The $M_1\ell_j$-cube $C$ containing a point $q\in\Lambda_j$ lies in the $\pi_j$-cube $\mathsf{C}$ of $q$, by compatibility of the partitions. By Step~2,
\[
\mathrm{dist}_\infty(C,(\Lambda_j^{\sim})^c)\ge M\ell_j\ge2M_1\ell_j,
\]
and Step~1 gives $C\subset D_j$. Hence $\Lambda_j\subset D_j$.

\emph{Step 4: the staircase.} The elements of levels $<j$ lie in $D_0\setminus D_j$. By Step~1, an $M_1\ell_j$-cube of $\Lambda_j^{\sim}$ not contained in $D_j$ is at $\mathrm{dist}_\infty<2M_1\ell_j$ from $(\Lambda_j^{\sim})^c$. Since $\Lambda_j^{\sim}$ is a union of $\pi_j$-cubes, and hence of $M_1\ell_j$-cubes, these cubes lie in the two outermost layers of $M_1\ell_j$-cubes of $\Lambda_j^{\sim}$, whose union has depth $2M_1\ell_j$. Since $2M_1<M$, they lie in the ring, and by Step~2 they are at $\mathrm{dist}_\infty\ge(M-2M_1)\ell_j$ from $\Lambda_j$. This completes (i).

The proof of (ii) proceeds in four steps.

\emph{Step 1: which readings of $H_j$ enter \cite[\S3]{B-CMP109}.} The estimates \cite[(3.15)--(3.17)]{B-CMP109} use $H_j$ through the two conditions ``$H_j(B(t))$ satisfies (3.14) with $\frac13\alpha_2$'' and ``$\delta H_j$ satisfies (3.9)'' \cite[(3.9)]{B-CMP109}. Both are imposed on $X$ in level-$j$ units. Condition \cite[(3.14)]{B-CMP109} reads
\[
|A|,\ |P_1A|,\ |\nabla^{\xi}_{\mathbf{U}}A|,\ |\Delta^{\xi}_{\mathbf{U}}A|<\alpha_2
\quad\text{on }X,
\]
where $\mathbf{U}$ is the variable that replaces $U_{k+1}$ in \cite[(3.13)]{B-CMP109}. These conditions are required at every cube and in both cases of \cite[(3.5)]{B-CMP109}.

At a good cube where \cite[(3.18)]{B-CMP109} is entered, the term removed in \cite[(3.21)]{B-CMP109} depends on $H_j$ through
\[
(1-\tilde\zeta_\Box)\,Q\bigl(\xi H_j(B_\Box)\bigr)\quad\text{on }\Box_0\setminus\tilde\Box^{3},
\]
where $Q$ and $B_\Box$ are the induced averaging map and the localised field of \cite[(3.18)--(3.21)]{B-CMP109}. Here $\tilde\zeta_\Box$ is the cut-off of \cite[p.~274]{B-CMP109}: it is supported in $\Box_0$, equals $1$ on $\tilde\Box^3$ and vanishes outside $\tilde\Box^4$. Here $H_j=\mathcal{H}'$ as in \eqref{4.29:eq-collar-staircase-1}, by \cite[Proposition~9]{B-CMP102b} applied to $\mathbf{B}'$, $W_0$ and $U^0$ restricted to $\Lambda_j^{\sim}$. The set $\mathbf{B}'$ is admissible for that proposition: \cite{B-CMP119} takes $M_1$ as ``the size of large cubes for which all the theorems of the previous papers \dots\ are valid'', which is at least the cube size required in \cite[p.~75]{B-CMP99a}.

\emph{Step 2: location of the observation points.}
\begin{itemize}
\item The domains $X$ of \cite[(2.27)]{B-CMP119} are ``contained in $\Lambda_j$'', and so is their localisation point $z$.
\item For the terms of \cite[(2.30)]{B-CMP119}, $X$ lies in $\Lambda_j$ with one layer of cubes removed.
\item The plaquettes entering the expressions \cite[(3.10)--(3.13)]{B-CMP109} lie within $\eta$ of $X$.
\item The blocks entering $P_1=DPD^*$ lie within $2\ell_j$ of $X$.
\item For every good cube $\Box$, of scale $s\ge j$, clause~(c) of Proposition~\ref{4.29:prop-native-cover} gives $\Box_0\subset\Lambda_s\subset\Lambda_j$.
\end{itemize}
Hence every observation point lies in $\Lambda_j$ or within $2\ell_j+\eta$ of it.

\emph{Step 3: observation points lie in $D_j$.} Let $v$ be within $2\ell_j+\eta$ of $\Lambda_j$. By Step~2 of the proof of (i), and since $\eta\le\ell_j$,
\[
\mathrm{dist}_\infty(v,(\Lambda_j^{\sim})^c)\ge M\ell_j-2\ell_j-\eta\ge(M-3)\ell_j .
\]
The $M_1\ell_j$-cube $C_v$ containing $v$ therefore satisfies
\[
\mathrm{dist}_\infty(C_v,(\Lambda_j^{\sim})^c)\ge(M-3-M_1)\ell_j\ge2M_1\ell_j,
\]
because $M\ge3M_1+3$. This inequality follows from $M>L^2M_1$, $L\ge3$ and $M_1\ge1$:
\[
M\ge9M_1+1\ge3M_1+3 .
\]
It is the only consequence of $M>L^2M_1$ used in this step. By (i), $C_v\subset D_j$, so $v$ is a point of level $j$ of $\mathbf{B}'$.

\emph{Step 4: distance from the staircase.} By (i), the elements of levels $<j$ are at $\mathrm{dist}_\infty\ge(M-2M_1)\ell_j$ from $\Lambda_j$. Hence they are at $\mathrm{dist}_\infty$ at least
\[
(M-2M_1)\ell_j-2\ell_j-\eta\ge(M-2M_1-3)\ell_j
\]
from $v$. Since every point at which $H_j=\mathcal{H}'$, and hence $U_{j,\Lambda_j^{\sim}}$ through \eqref{4.29:eq-collar-staircase-1}, is evaluated is an observation point, $U_{j,\Lambda_j^{\sim}}$ is never evaluated on $D_0\setminus D_j$. This proves (ii).

\emph{Proof of (iii).} Let $x$ be an observation point, lying in the block $\Delta(y)$ of a point $y$ of the level-$j$ lattice. By (ii), $\mathbf{B}'$ has level $j$ at $x$. This is the level $j$ of the point at which \cite[(190)]{B-CMP102b} is evaluated, that display being stated for ``$x\in\Delta(y)$, $y\in\Lambda_j$, $y'\in\Lambda_{j'}$'', in which $\Lambda_j$ and $\Lambda_{j'}$ are the lattices of levels $j$ and $j'$, with the operators $\nabla^e$, $e\in\{0,1,2\}$, including $D^*D$ and $\Delta^{\xi}_{\mathbf{U}}$, normalised by $(L^j\eta)^{-1-e}=\ell_j^{-1-e}$. Hypothesis (dec) therefore applies at $x$ for each $e\in\{0,1,2\}$ and gives \eqref{4.29:eq-collar-staircase-a}, with $C_{90}$ and $\delta_0$ uniform in $k$ by (dec); conclusion (iii) is hypothesis (dec) at the points supplied by (ii).

By (dec) and \cite[(2.46), (2.48)]{B-CMP96}, the distance $\operatorname{d}$ is measured along contours whose part lying in the level-$j$ blocks consists of bonds of the level-$j$ lattice. The layered instance displayed in \cite[(1.45)--(1.47)]{B-CMP122a}, ``defined in terms of $M_1$-cubes on corresponding scales'', is used here for the layered form of the bound only.

Finally, $\mathcal{H}'$ is in the Landau gauge by \cite[Proposition~9]{B-CMP102b}. In this gauge the identity $P_1\mathcal{H}'=\Delta^{\xi}_{\mathbf{U}}\mathcal{H}'-D^*D\mathcal{H}'+(\text{local lower order})$ is the one \cite[p.~272]{B-CMP109} uses for its $H_j$. Hence $|P_1\mathcal{H}'|$ is controlled by \eqref{4.29:eq-collar-staircase-a} for $e\le2$. This proves (iii).

The proof of (iv) proceeds in four steps. Write $m=m(c)$. All levels of $\mathbf{B}'$ are at most $j$, so $m\le j$.

\emph{Step 1: locality.} By the locality of the averaging operations \cite{B-CMP98}, $Q'(\eta\mathbf{a})(c)$ depends on $\mathbf{a}$ only through its restriction to $A(c)$. Consequently, if $\mathbf{b}$ vanishes on $A(c)$, then $DQ'|_{\eta\mathbf{a}}[\eta\mathbf{b}](c)=0$, and \eqref{4.29:eq-collar-staircase-2} holds trivially. We therefore assume $\sup_{A(c)}|\mathbf{b}|>0$ when proving \eqref{4.29:eq-collar-staircase-2}.

\emph{Step 2: applying (tw).} Take $s=m$, $\tilde B(c)=A(c)$, $\mathbf{S}=\mathbf{B}'$, so that $M_{\mathbf{S}}=M'$, $U_0=U^0$ and $\mathcal{X}$ a twist with $32\,\ell_m\sup_{A(c)}|\mathcal{X}|\le\alpha_1^{(98)}$. We check the hypotheses of (tw).
\begin{itemize}
\item $U^0$ satisfies \cite[(52)]{B-CMP98} at scale $\ell_m$ on $A(c)$ by (reg), since $U^0$ is regular at levels $\ge j\ge m$.
\item $\ell_m\sup_{A(c)}|A'|\le\alpha_1^{(98)}/4$ by assumption.
\item $\ell_m\sup_{A(c)}|\mathcal{X}|\le\alpha_1^{(98)}/32<\alpha_1^{(98)}/4$.
\end{itemize}
So $M'(e^{i\eta\mathcal{X}}\mathbb{U})(c)=e^{i\mathfrak{h}(c)}M'(\mathbb{U})(c)$, and
\[
|e^{i\mathfrak{h}(c)}-1|\le c_6'\,\ell_m\sup_{A(c)}|\mathcal{X}|
\le\frac{16}{\alpha_1^{(98)}}\cdot\frac{\alpha_1^{(98)}}{32}=\frac12 .
\]

\emph{Step 3: the bound \eqref{4.29:eq-collar-staircase-b}.} Put $\varpi:=e^{i\mathfrak{h}(c)}-1$. On the convex set of twists with $32\,\ell_m\sup_{A(c)}|\mathcal{X}|\le\alpha_1^{(98)}$ we have $|\varpi|\le\frac12$. We take $\mathfrak{h}(c)$ to be the determination of $-i\log(1+\varpi)$ that is continuous from $\mathcal{X}=0$, where it vanishes. Comparison with $M'(e^{i\eta\mathcal{X}}\mathbb{U})=e^{iQ'(\eta\mathcal{X})}M'(\mathbb{U})$, with $Q'(0)=0$, identifies this determination with $Q'(\eta\mathcal{X})(c)$.

The logarithmic series gives
\[
|\mathfrak{h}(c)|\le\sum_{n\ge1}\frac{|\varpi|^n}{n}=-\log(1-|\varpi|).
\]
The function $-\log(1-\cdot)$ is convex on $[0,\frac12]$, vanishes at $0$ and takes the value $\log2\le1$ at $\frac12$. Hence it is at most twice its argument on $[0,\frac12]$, and since $|\varpi|\le\frac12$,
\[
|Q'(\eta\mathcal{X})(c)|=|\mathfrak{h}(c)|\le2|\varpi|\le2c_6'\,\ell_m\sup_{A(c)}|\mathcal{X}|
=c_Q\,\ell_m\sup_{A(c)}|\mathcal{X}| .
\]
For $\mathcal{X}=\mathbf{a}$ with $64\,\ell_m\sup_{A(c)}|\mathbf{a}|\le\alpha_1^{(98)}$ this is \eqref{4.29:eq-collar-staircase-b}.

\emph{Step 4: the derivative bound \eqref{4.29:eq-collar-staircase-2}.} Consider the circle
\[
\mathbf{a}+\varsigma\mathbf{b},\qquad
|\varsigma|=\varsigma_0:=\frac{\alpha_1^{(98)}}{64\,\ell_m\sup_{A(c)}|\mathbf{b}|}.
\]
On this circle,
\[
\ell_m\sup_{A(c)}|\mathbf{a}+\varsigma\mathbf{b}|
\le\frac{\alpha_1^{(98)}}{64}+\frac{\alpha_1^{(98)}}{64}=\frac{\alpha_1^{(98)}}{32}.
\]
So the twist $\mathbf{a}+\varsigma\mathbf{b}$ lies in the set of Step~3, and
\[
|Q'(\eta(\mathbf{a}+\varsigma\mathbf{b}))(c)|\le c_Q\,\alpha_1^{(98)}/32 .
\]
The map $\varsigma\mapsto Q'(\eta(\mathbf{a}+\varsigma\mathbf{b}))(c)$ is analytic on the closed disc $|\varsigma|\le\varsigma_0$. Indeed, the averages are analytic in the twist by (tw), and the chosen determination of the logarithm is analytic where $|\varpi|\le\frac12$. Cauchy's formula
\[
DQ'|_{\eta\mathbf{a}}[\eta\mathbf{b}](c)
=\frac{1}{2\pi i}\oint_{|\varsigma|=\varsigma_0}\frac{Q'(\eta(\mathbf{a}+\varsigma\mathbf{b}))(c)}{\varsigma^{2}}\,d\varsigma
\]
then gives
\[
\bigl|DQ'|_{\eta\mathbf{a}}[\eta\mathbf{b}](c)\bigr|
\le\frac{c_Q\,\alpha_1^{(98)}}{32\,\varsigma_0}
=2c_Q\,\ell_m\sup_{A(c)}|\mathbf{b}|,
\]
which is \eqref{4.29:eq-collar-staircase-2}.

\emph{Finer sources are smaller.} Since $m(c)\le j$ we have $\ell_{m(c)}\le\ell_j$, so both bounds are at most their values with $\ell_j$. This is the direction used in \cite[(3.21)]{B-CMP109}: ``from the exponential decay of the derivative $(\delta/\delta B)\mathbf{H}_j$ and from the condition $\mathrm{dist}^{(\xi)}(X,\operatorname{supp}(1-\tilde\zeta_\Box))\ge M(L^j\eta)^{-1}$, we obtain a bound of the type (3.9)'' \cite[p.~274]{B-CMP109}. There the sources lie on $\Box_0\setminus\tilde\Box^4$, which contains the staircase of the determining set attached to $\Box_0$ in the localisation of \cite[\S3]{B-CMP109}, and they are observed on $X\subset\tilde\Box^2$. The norms at the observation point in (dec) are taken at spacing $\eta$. This proves (iv).

\emph{Proof of (v).} As shown in the proof of Lemma~\ref{4.29:lem-localisation-exact}, and in the notation of that lemma, the insertion of $\Box\in\mathcal{C}_j$ enters only through $DQ'|_{\eta t\mathbf{H}_k}[\eta\bar\zeta_\Box\mathbf{H}_k]$. By Step~1 of the proof of (iv), $DQ'$ depends on $\bar\zeta_\Box$ only through its values on the end-blocks $A(c)$ of elements $c$ of $\mathbf{B}'$. These blocks consist of bonds with at least one end-point in $\Lambda_j^{\sim}$ \cite[p.~75]{B-CMP99a}, so they lie within $\eta$ of $\Lambda_j^{\sim}$. Hence only the values of $\bar\zeta_\Box$ near $\Box\cap\Lambda_j^{\sim}$ are used. These points lie in the $\frac13M\ell_j$-neighbourhood of $\Lambda_j^{\sim}$, on which $S\ge1$ and the gradient bounds hold by clause~(a) of Proposition~\ref{4.29:prop-native-cover}.
\end{proof}

By \eqref{4.29:eq-collar-staircase-a} and \eqref{4.29:eq-collar-staircase-b}, the sums of \cite[\S3]{B-CMP109} require, beyond the estimates made there, only a count of the sources of level below $j$.

\begin{lemma}[Summation of finer sources over the collar]\label{4.29:lem-collar-sources}
Let $j$ be a level, and let $\Lambda_j^{\sim}$, the set $\mathbf{B}'=\mathbf{B}_j(\Lambda_j^{\sim})$ of
\cite[(2.13)]{B-CMP119} and its maximal domains $D_n$ be as in
Lemma~\ref{4.29:lem-localisation-exact}. For $n<j$ the elements of $\mathbf{B}'$ on
$D_n\setminus D_{n+1}$ have level $n$, those on $D_j$ have level $j$, and $m(c)$ denotes the level
of an element $c$. Let $\ell_n$ be the side of a block of level $n$, so that $\ell_{n+1}=L\ell_n$.
Let $\delta_0>0$ be the decay constant of \cite[Lemma~2.1]{B-CMP96} and $M_1$ the large-cube size
of \cite{B-CMP119}. Assume
the condition \eqref{4.29:eq-native-cover-a} of the native-scale cover, which, with $M_1$ the
large-cube size of \cite{B-CMP119} as fixed above, reads
\begin{equation*}
L^{d-1}e^{-\delta_0 M_1/16}\le\tfrac12 .
\end{equation*}
Let $y$ be a level-$j$ block of $\mathbf{B}'$, and let $\beta\ge0$ be a function on the elements of
$\mathbf{B}'$. Then
\begin{equation}\label{4.29:eq-collar-sources-a}
\sum_{\substack{c\in\mathbf{B}'\\ m(c)<j}} e^{-\frac18\delta_0 d(y,c)}\,\ell_{m(c)}\,\beta(c)
\;\le\; 2e^{\delta_0/8}C_d L^{d-1}\,\ell_j\sum_{\hat c}
e^{-\frac18\delta_0 d^{(j)}(y,\hat c)}\,\hat\beta(\hat c).
\end{equation}
Here:
\begin{itemize}
\item $\hat c$ runs over the level-$j$ blocks that tile the collar $\Lambda_j^{\sim}\setminus D_j$;
\item an element $c$ with $m(c)<j$ is \emph{inside} the block $\hat c$ that contains its end-point
of its own level; a level-$0$ bond with one end-point outside $\Lambda_j^{\sim}$ is inside the
block of its inner end-point;
\item $\hat\beta(\hat c)$ is the maximum of $\beta$ over the elements inside $\hat c$;
\item $d(y,c)$ is the scaled distance of \cite[(2.46), (2.48)]{B-CMP96} on $\mathbf{B}'$: along a
path, a segment in
$D_n\setminus D_{n+1}$ counts its length in sites of level $n$, and a segment in $D_j$ counts it
in sites of level $j$;
\item $d^{(j)}(y,\cdot)$ is the scaled distance in which all of $\Lambda_j^{\sim}$ is counted at
level $j$, so that $d^{(j)}(y,c)\le d(y,c)$;
\item $C_d$ is a constant that depends on the dimension $d=4$ only.
\end{itemize}
\end{lemma}

\begin{proof}
We write $\iota:=j-m(c)\ge1$ for the number of levels by which an element $c$ of the left side
of \eqref{4.29:eq-collar-sources-a} is finer than $y$.

\emph{Counting.} The partitions into blocks of the different levels are compatible. A level-$j$
block $\hat c$ therefore contains at most $L^{d\iota}$ blocks of level $j-\iota$. An element of
level $j-\iota$ is such a block together with one of the $d$ bond directions. Hence $\hat c$
contains at most $C_dL^{d\iota}$ elements of level $j-\iota$. For $\iota=j$ the elements are
the level-$0$ bonds. Those with one end-point outside $\Lambda_j^{\sim}$ are, by the convention
of the statement, inside the block of their inner end-point. A site of such a block is then the
inner end-point of at most $2d$ level-$0$ bonds instead of $d$; the resulting factor is absorbed
in $C_d$.

\emph{Distance.} Let $c$ have level $m(c)=j-\iota$. Its end-point of its own level lies in
$D_{j-\iota}\setminus D_{j-\iota+1}$. Consider a path from the block $y$, which lies in $D_j$, to
$c$. It
crosses each layer $D_n\setminus D_{n+1}$ with $n=j-1,\dots,j-\iota+1$; there are $\iota-1$ such
layers. By the layer property of \cite[(2.13)]{B-CMP119}, the distance from $D_{n+1}$ to the
complement of $D_n$ is at least $M_1\ell_{n+1}$, so the path crosses such a layer over a physical
length $P\ge M_1\ell_{n+1}$. In $d(y,c)$ this segment is counted at the local level $n$, that is,
as $P/\ell_n$ sites. In $d^{(j)}(y,c)$ it is counted at level $j$, as $P/\ell_j$ sites, and
$P/\ell_j\le P/(L\ell_n)$ because $\ell_j\ge\ell_{n+1}=L\ell_n$ for $n+1\le j$. Each crossed layer
therefore gains
\begin{equation*}
\frac{P}{\ell_n}-\frac{P}{\ell_j}\ge\frac{P(L-1)}{L\ell_n}\ge(L-1)M_1
\end{equation*}
in $d(y,c)$ over $d^{(j)}(y,c)$. Every other segment of the path counts in $d(y,c)$ at least
its level-$j$ length, because its local level is at most $j$. Taking the infimum over paths gives
\begin{equation}\label{4.29:eq-collar-sources-1}
d(y,c)\ge d^{(j)}(y,c)+(\iota-1)(L-1)M_1\ge d^{(j)}(y,\hat c)-1+(\iota-1)(L-1)M_1 ,
\end{equation}
where $\hat c$ is the level-$j$ block inside which $c$ lies.

The blocks $\hat c$ are the level-$j$ blocks tiling the collar $\Lambda_j^{\sim}\setminus D_j$, and
every element $c$ with $m(c)<j$ is inside one of them: its own-level end-point lies in
$D_{j-\iota}\setminus D_{j-\iota+1}$, and this set is contained in $\Lambda_j^{\sim}\setminus D_j$
because $D_j\subset D_{j-\iota+1}$.

The second inequality in \eqref{4.29:eq-collar-sources-1} holds for the following reason. The
length $\ell_{j-\iota}$ of $c$ is less than one level-$j$ site, and distances to $\hat c$ are
counted in level-$j$ blocks in the $\ell^{\infty}$ sense. The offset $-1$ is therefore one
level-$j$ unit. If distances to $\hat c$ were counted bond by bond, the offset would be $-d$, one
unit for each of the $d=4$ coordinate directions, and
the factor $e^{\delta_0/8}$ below would become $e^{d\delta_0/8}$.

\emph{Summation.} Group the left side of \eqref{4.29:eq-collar-sources-a} by the block $\hat c$
inside which $c$ lies and by $\iota$. Bound $\beta(c)$ by $\hat\beta(\hat c)$. Use
$\ell_{m(c)}=L^{-\iota}\ell_j$, the count of at most $C_dL^{d\iota}$ elements of level $j-\iota$
inside $\hat c$, and \eqref{4.29:eq-collar-sources-1}. Extending the sum over $\iota$, which
stops at $\iota=j$, to all $\iota\ge1$, the left side is at most
\begin{align*}
&\sum_{\hat c}\hat\beta(\hat c)\sum_{\iota\ge1}C_dL^{d\iota}L^{-\iota}\ell_j\,
e^{\delta_0/8}e^{-\frac18\delta_0 d^{(j)}(y,\hat c)}e^{-\frac18\delta_0(\iota-1)(L-1)M_1}\\
&\qquad=e^{\delta_0/8}\ell_j\sum_{\hat c}e^{-\frac18\delta_0 d^{(j)}(y,\hat c)}\hat\beta(\hat c)
\cdot C_dL^{d-1}\sum_{\iota\ge1}\bigl(L^{d-1}e^{-\frac18\delta_0(L-1)M_1}\bigr)^{\iota-1}.
\end{align*}

Since $L\ge2$, we have $(L-1)/8\ge1/16$, and the assumed condition gives
\begin{equation*}
L^{d-1}e^{-\frac18\delta_0(L-1)M_1}\le L^{d-1}e^{-\delta_0M_1/16}\le\tfrac12 .
\end{equation*}
If the two boundary bonds of each crossed layer are left out of the gain, $L\ge3$ is needed in
place of $L\ge2$; this holds because the blocking factor $L$ is an odd integer larger than $1$.

The geometric series in $\iota$ is therefore at most $2$. This gives
\eqref{4.29:eq-collar-sources-a} with the constant $2e^{\delta_0/8}C_dL^{d-1}$.

For $\iota=1$ the count is at most $C_dL^{d}$ and the weight is $\ell_{j-1}=L^{-1}\ell_j$, which
together give the factor $C_dL^{d-1}\ell_j$; no decay is used for this layer. The constant
$2e^{\delta_0/8}C_dL^{d-1}$ does not depend on $j$, on the index $k$ of the step in
Lemma~\ref{4.29:lem-localisation-exact}, on the number $K$ of renormalisation steps, or on the
size of $\Lambda_j$.
\end{proof}

\begin{remark}
The condition \eqref{4.29:eq-native-cover-a} assumed in Lemma~\ref{4.29:lem-collar-sources} is a
one-sided condition on $M_1$ in terms of $L$ only; $M_1$ is chosen after $L$ and before
$M_2$, $M$ and $\gamma$. It is of the same kind as \cite[(163)]{B-CMP102b}, \cite[(2.59)]{B-CMP96}
and \cite[(1.4)]{B-CMP99a}.

The proof of Lemma~\ref{4.29:lem-collar-sources} uses only one property of the domain
$\Lambda_j^{\sim}$: the layer property of the maximal domains of \cite[(2.13)]{B-CMP119}, namely
that the distance from $D_{n+1}$ to the complement of $D_n$ is at least $M_1\ell_{n+1}$.

\emph{A cube in place of the collared domain.} The same argument applies verbatim with
$\Lambda_j^{\sim}$ replaced by a cube $\square_0$ at a window and $\mathbf{B}'$ by the determining
set of $\square_0$ of \cite[(2.13)]{B-CMP119}, taken at the level at which the window is read, as
identified in the corresponding clause of the lemma on the determining set of a cube at a window;
the maximal domains of that determining set have the layer property used in the proof by
\cite[(2.13)]{B-CMP119}.

\emph{No length weight.} It also applies without the weight $\ell_{m(c)}$ on the left side, for
data that carry no such weight. The first finer layer then contributes $L^{d}$ in place of
$L^{d-1}$, and the ratio of consecutive layers, $L^{d}e^{-\frac18\delta_0(L-1)M_1}$, is bounded
as follows. By the assumed condition, $e^{-\delta_0M_1/16}\le\frac12L^{-(d-1)}$. Hence
\begin{equation}\label{4.29:eq-collar-sources-2}
\begin{split}
L^{d}e^{-\frac18\delta_0(L-1)M_1}&\le L^{d}e^{-\delta_0(L-1)M_1/16}\\
&\le L^{d}\bigl(\tfrac12L^{-(d-1)}\bigr)^{L-1}\le\tfrac12 .
\end{split}
\end{equation}
The last inequality holds because $L\ge3$ gives $d-(d-1)(L-1)\le2-d\le0$ and
$2^{-(L-1)}\le\frac12$. The geometric series is therefore again at most $2$, and the bound holds
with the constant $2e^{\delta_0/8}C_dL^{d}$ and no factor $\ell_j$, that is, the left side is at
most $2e^{\delta_0/8}C_dL^{d}$ times
$\sum_{\hat c}e^{-\frac18\delta_0 d^{(j)}(y,\hat c)}\hat\beta(\hat c)$.

This is the sum formed in \cite[(3.21), (3.54), (4.1)]{B-CMP109}.
\cite[Lemma~2.1, (2.61)]{B-CMP96} states the bound
$\sup_y\sum_{y'}e^{-a\delta_0 d(y,y')}\le c_1(a)$, with $y'$ over the determining set of that
lemma and $a>0$ its exponent, under the condition \cite[(2.59)]{B-CMP96};
Lemma~\ref{4.29:lem-collar-sources} is its weighted form on the collar. The same bound without
the length weight is the one used in the proof of the lemma in which the product $B_0B_3$ of
Ba{\l}aban's constants, taken uniformly in the structure of the domains, enters, at the step where
that product enters; the finiteness of $B_3$, established in the corollary bounding it, therefore
presupposes the condition \eqref{4.29:eq-native-cover-a}.
\end{remark}

\begin{definition}[The analyticity space of the pieces born at the step]
\label{4.29:def-analyticity-space}
Fix the level $k$ and consider the stage of the step at which Ba{\l}aban's large-field operation
$\mathbf{T}$ acts. We use the scales $\eta=L^{-k}$ and $\ell_n=L^n\eta$, the cubes of scale $n$
and their enlargements, the current label sequence $\{\Omega_i\}_{i\le k+1}$ with its index
$n(x)$, and the native label sequence $\{\Omega^{\natural}_i\}$, all as in
Notation~\ref{4.29:not-scales-labels}. We write $l(x)$ for the native index of a point $x$, that
is, the index $i$ with $x\in\Omega^{\natural}_i\setminus\Omega^{\natural}_{i+1}$. We write
$L_\circ$ for the constant denoted $L_0$ in \cite[(1.64)--(1.69)]{B-CMP122a}; in \cite{B-CMP122a}
it satisfies $L_\circ\ge1$, and it is distinct from
the existential constant $L_0=L_0(N)$ of Theorem~0 (Notation~\ref{0.2:not-glossary}). We write
$\beta$ for the constant of \cite[(2.34)]{B-CMP119}, and $\alpha_{0,n}$, $\alpha_{1,n}$ for the
radii entering the spaces of \cite[(2.34)--(2.39)]{B-CMP119}.

An \emph{arrested component} $W$ is a component of the large-field region left in arrested form
by the arrest construction. It carries data $h<k_0<k_W\le k$ (the letter $h$ as in
\cite[(1.64)]{B-CMP122a}, not a history), and $j_W\le k_W$ denotes its last
level, the datum written $j^{+}$ in the definition of the arrest. Inside $W$ we use the regions
$\Omega''_m$ of \cite[(1.64)]{B-CMP122a}, and we write
$\Gamma''_i=\Omega''_i\setminus\Omega''_{i+1}$. We write $\Lambda_i$ for the small-field region at
level $i$, $Z'_i$ for the regions of the arrest, which satisfy $Z'_i\supset\Lambda_i^{c}$, and
$Z''_i$ for their enlargements.
A superscript $\sim r$ denotes enlargement by $r$ layers of cubes, as for $\tilde{\Delta}^r$ in
Notation~\ref{4.29:not-scales-labels}, and $\sim -r$ denotes removal of $r$ layers. A piece is
\emph{live} if it still carries its large-field factor at the level read.

Let $\square_0$ be the cube at which the hypothesis at a cube is posed, and let $Y$ be a
domain with $Y\supset\square_0$. The \emph{analyticity space of the pieces born at the step},
$\mathfrak{S}(Y)$, is defined as follows.

\begin{enumerate}
\item At every point of $Y$ off live arrested pieces, the defining conditions of
$\mathfrak{S}(Y)$ are those of the space of \cite[(3.43), p.~276]{B-CMP119},
\begin{equation*}
\tilde{U}^{c}_{k+1}\bigl(Y,(1+\beta)\tilde{\alpha}_0,(1+\beta)\tilde{\alpha}_1\bigr),
\end{equation*}
where $\tilde{\alpha}_0$ and $\tilde{\alpha}_1$ are the radii appearing there. In other words,
they are the conditions of \cite[(2.34)--(2.39)]{B-CMP119} at top index $k+1$, with the radii
enlarged as in \cite[(3.43)]{B-CMP119}. On live arrested pieces they are modified as in the next
item.

\item For live arrested pieces we turn the following sentence of \cite{B-CMP122a} into a
definition:
\begin{quotation}
it is easy to generalize the definition of the spaces in such a case for the next steps
\end{quotation}
The definition is as follows.
On an arrested component $W$, the structure of the clauses of \cite[(1.64)--(1.69)]{B-CMP122a}
is frozen. This structure consists of the level, the weight and the family of cubes, and it is
taken with the parameters $(h,j_W,k_W,k_0)$ in place of $(h,j,k,k_0)$. The factor multiplying the
bracket is lowered by the frozen profile of the arrest. The level and the weight are frozen at a
single stage of the arrest; this is why the weight bound~\eqref{4.29:eq-analyticity-space-1}
below holds on every branch. The conventions for the collar of the arrest are those in which the
source constant is $c_{\mathrm{src}}=1$ and the level interval of an arrest $\mathfrak{a}$ is
$I(\mathfrak{a})=[h_{\mathfrak{a}}+1,k_{\mathfrak{a}}+1]$.
\end{enumerate}

All conditions are read with the convention on the localisation domains $X'$ and on their
trimmed domains $X'^{\flat}$, and with the convention that the conditions on the localised
quantities hold also for the gauge-group-valued factor $U$ of $\mathbb{U}$; the content of these
conventions that enters here is stated in the first item after $(\mathfrak{S}1)$ below.

For comparison, we recall the printed conditions. In \cite[(2.34)--(2.37), (2.39)]{B-CMP119} the
seven quantities $\mathcal{Q}_1,\dots,\mathcal{Q}_7$ listed below are bounded by the factor
\begin{equation*}
\bigl(1-\beta(1-2^{-(j-n)})\bigr)
\end{equation*}
times the bracket displayed below, with $n$ and $\xi$ in place of $m$ and $\eta$. These bounds hold
\begin{quotation}
on $X\cap(\Omega_n\setminus\Omega_{n+1})$ for $n = 1,\dots,j-1$, or on $X\cap\Omega_j$ for $n = j$
\end{quotation}
and \cite[(2.38)]{B-CMP119} reads:
\begin{quotation}
(ii) For each cube $\square\subset\Omega_n\setminus\Omega_{n+2}$,
$\square\cap\Omega^{c}_{n+1}\neq\emptyset$, $n = 1,\dots,j-1$, of the size $CML^n\xi$, or
$\square\subset\Omega_j$ and of the size $CM$, there exists a gauge transformation $u$ defined on
$\square\cap X$ and such, that $U^u = \exp i\xi A$,
\begin{equation*}
L^n\xi|A|,\ (L^n\xi)^2|\nabla^{\xi}A| < BCM\alpha_{0,n} \tag*{\cite[(2.38)]{B-CMP119}}
\end{equation*}
\end{quotation}
Here, in the notation of \cite[p.~261]{B-CMP119},
\begin{equation*}
U_{p,X}(M^{\cdot}(\mathbb{U}))=U\bigl(\mathbb{B}_p(X)\cup\{\Gamma_i\}_{i<p},M^{\cdot}(\mathbb{U})\bigr).
\end{equation*}
In \cite[p.~190]{B-CMP122a}, where $L_0$ is the constant we write $L_\circ$, the conditions read:
\begin{quotation}
(i) $|\partial\mathbb{U}-1|$, $|\partial U_{p,X}(M^{\cdot}(\mathbb{U}))-1|$, $|\mathbb{J}|$,
$|J_{p,X}(M^{\cdot}(\mathbb{U}))|$, $|\partial U-1|$, $|A'|$, $|\nabla^{\eta}_{U}A'|$
\begin{multline*}
< \Bigl(1-\beta\sum_{i=h+1}^{j}2^{-|m-i|}-\beta\sum_{q=m+1}^{k}2^{-(q-m)}\Bigr)
L_0^{2\max\{0,m-k_0\}}\cdot\\
\bigl[\alpha_{0,m}\eta^2(L^m\eta)^{-2},\ \alpha_{0,m}L^{-2p}(L^mL^{-p})^{-2},\\
\alpha_{0,m}(L^m\eta)^{-3},\ \alpha_{0,m}L^{m-\min\{p,m\}}(L^mL^{-p})^{-3},\\
\alpha_{0,m}\eta^2(L^m\eta)^{-2},\ \alpha_{1,m}(L^m\eta)^{-1},\ \alpha_{1,m}(L^m\eta)^{-2}\bigr]
\end{multline*}
on $(\Omega''_m\setminus\Omega''_{m+1})\cap X$ for $m = 1,\dots,j-1$, and on
$(\Omega''_j\setminus\Omega_{k_0+1})\cap X$ for $m = j$, $p = 1,\dots,k$.
\begin{equation*}
L^m\eta|A|,(L^m\eta)^2|\nabla^{\eta}A| < L_0^{2\max\{0,m-k_0\}}BCM\alpha_{0,m}
\tag*{\cite[(1.65)]{B-CMP122a}}
\end{equation*}
for $\square\subset\Omega''_m$, $\square\cap\Omega''^{c}_{m+1}\neq\emptyset$ \dots{} and for
$\square\subset\Omega''_j$, $\square\cap\Omega^{c}_{k_0+1}\neq\emptyset$ if $m = j$

(ii)
\begin{equation*}
< (\dots)L_0^{2(j-m)}\cdot[\alpha_{0,j}\eta^2(L^j\eta)^{-2}, \dots,
\alpha_{1,j}(L^j\eta)^{-2}]\tag*{\cite[(1.66)]{B-CMP122a}}
\end{equation*}
on $(\Omega_m\setminus\Omega_{m+1})\cap X$ for $m = k_0+1,\dots,j-1,j$.
\begin{equation*}
L^j\eta|A|,(L^j\eta)^2|\nabla^{\eta}A| < L_0^{2(j-m)}BCM\alpha_{0,j}
\tag*{\cite[(1.67)]{B-CMP122a}}
\end{equation*}
for $\square\subset\Omega_m$, $\square\cap\Omega^{c}_{m+1}\neq\emptyset$, $\square$ is of the size
$CML^j\eta$
\end{quotation}
In \cite[p.~191]{B-CMP122a} the conditions continue:
\begin{quotation}
(iii)
\begin{equation*}
< (\dots)\cdot c[\alpha_{0,m}\eta^2(L^m\eta)^{-2}, \dots]\tag*{\cite[(1.68)]{B-CMP122a}}
\end{equation*}
on $(\Omega_m\setminus\Omega_{m+1})\cap X$ for $m = j+1,\dots,k-1$, and on $\Omega_k\cap X$ for
$m = k$, where $c = 1$ for $m < k$, and $c = 3$ for $m = k$,
\begin{equation*}
L^m\eta|A|,(L^m\eta)^2|\nabla^{\eta}A| < BCM\alpha_{0,m}\tag*{\cite[(1.69)]{B-CMP122a}}
\end{equation*}
\end{quotation}
and
\begin{quotation}
For $j\le k_0$ \dots{} there are no powers of $L_0$ in the point (i), and the point (ii) is empty.
\end{quotation}
The index $p$ in these quotations is the index $p=1,\dots,k$ of \cite{B-CMP122a}, not a plaquette.

\emph{Pointwise form.} Let $\mathbb{U}$ be a complexified configuration and $\mathbb{J}$ a complex
current. For $e=1,\dots,7$ let $\mathcal{Q}_e$ denote, in this order, the quantities
\begin{gather*}
|\partial\mathbb{U}-1|,\quad |\partial U_{p,X'}(M^{\cdot}\mathbb{U})-1|,\quad |\mathbb{J}|,\quad
|J_{p,X'}(M^{\cdot}\mathbb{U})|,\\
|\partial U-1|,\quad |A'|,\quad |\nabla^{\eta}_{U}A'|.
\end{gather*}
Here $M^{\cdot}(\mathbb{U})$, $U_{p,X'}$ and $J_{p,X'}$ are as in \cite[p.~261]{B-CMP119}, with
$X'$ in place of $X$.
Let $\mathfrak{r}_e(n)$ be the $e$-th member of the bracket above at level $n$:
\begin{align*}
\mathfrak{r}_1(n)&=\alpha_{0,n}\eta^2\ell_n^{-2}, &
\mathfrak{r}_2(n)&=\alpha_{0,n}L^{-2p}(L^nL^{-p})^{-2},\\
\mathfrak{r}_3(n)&=\alpha_{0,n}\ell_n^{-3}, &
\mathfrak{r}_4(n)&=\alpha_{0,n}L^{n-\min\{p,n\}}(L^nL^{-p})^{-3},\\
\mathfrak{r}_5(n)&=\alpha_{0,n}\eta^2\ell_n^{-2}, &
\mathfrak{r}_6(n)&=\alpha_{1,n}\ell_n^{-1},\\
\mathfrak{r}_7(n)&=\alpha_{1,n}\ell_n^{-2}.
\end{align*}
Only $\mathfrak{r}_4$ depends on $p$; in $\mathfrak{r}_2$ the powers of $L^{p}$ cancel. The
homogeneity of
$\mathfrak{r}_e(n)$ in $L^n$ is $\nu_e$, where $\nu_1,\dots,\nu_7$ are $2,2,3,\nu_4,2,1,2$ and
$\nu_4$ is $2$ or $3$. That is, for $e\ne4$, with $\alpha_{\cdot,n}$ the radius occurring in
$\mathfrak{r}_e$,
\begin{equation*}
\frac{\mathfrak{r}_e(n)}{\mathfrak{r}_e(m)}
=\frac{\alpha_{\cdot,n}}{\alpha_{\cdot,m}}\,L^{-\nu_e(n-m)}.
\end{equation*}
For $e=4$,
\begin{equation*}
\frac{\mathfrak{r}_4(n)}{\mathfrak{r}_4(m)}
=\frac{\alpha_{0,n}}{\alpha_{0,m}}\,L^{-2(n-m)}\,L^{-(\min\{p,n\}-\min\{p,m\})}.
\end{equation*}
This follows from the exponent
$(n-\min\{p,n\})-(m-\min\{p,m\})-3(n-m)$ of $L$. Hence $\nu_4=3$ when $n,m\le p$, and $\nu_4=2$
when $n,m\ge p$.

For each $e$ let $\Phi_e(x)$ denote the scalar factor other than the weight $\varpi(x)$
of~\eqref{4.29:eq-analyticity-space-b} below that multiplies the bracket member $\mathfrak{r}_e$
in the condition on $\mathcal{Q}_e$ at $x$ (in the quotation from \cite[p.~190]{B-CMP122a}, the
factor in large parentheses containing $\beta$), and let $\Phi^{\mathrm{nat}}_e(n)$
denote the value of this factor in the native conditions at level $n$, that is, in the conditions
before the modification on live arrested pieces described above.
Then $(\mathbb{U},\mathbb{J})\in\mathfrak{S}(Y)$ means that $\mathbb{U}=U'U$, with $U$ valued in
the gauge group and $U'=e^{i\eta A'}$, and that the conditions $(\mathfrak{S}1)$ and
$(\mathfrak{S}2)$ hold, with the weights $\varpi(x)$ and $\varpi_\square$
of~\eqref{4.29:eq-analyticity-space-b} below, the factor $\Phi_e$ having property $(\Phi)$:
\begin{equation}\label{4.29:eq-analyticity-space-a}
\begin{aligned}
&(\mathfrak{S}1)\quad \mathcal{Q}_e(x)<\Phi_e(x)\,\varpi(x)\,\mathfrak{r}_e(n(x)),\\
&(\mathfrak{S}2)\quad \ell_n|A|,\ \ell_n^2|\nabla^{\eta}A|<\varpi_\square\,BCM\alpha_{0,n},\\
&(\Phi)\quad \Phi_e(x)=\Phi^{\mathrm{nat}}_e(n(x)),\quad
1-\beta\le\Phi^{\mathrm{nat}}_e(\cdot)\le1+\beta\quad\text{off live pieces},\\
&\phantom{(\Phi)}\quad \Phi_e(x)\le\Phi^{\mathrm{nat}}_e(n(x))\quad\text{on live pieces}.
\end{aligned}
\end{equation}
The quantifiers and readings are as follows.
\begin{itemize}
\item $(\mathfrak{S}1)$ holds at every $x\in Y$ for $e=1,3,5,6,7$. For $e=2,4$ it holds for every
$p\le k+1$, every localisation domain $X'\subset Y$ and every $x$ in the trimmed domain
$X'^{\flat}$ of $X'$. For $e=2,4$ it holds also with $U$ in place of $\mathbb{U}$. By the lemma
on trimmed localisation domains, $X'^{\flat}\supset X'^{\,\sim-2}$ whenever
$\ell_{m_{X'}}\ge\ell_p$, where $m_{X'}$ is the scale of $X'$.
\item $(\mathfrak{S}2)$ holds for every $n$ and every cube $\square$ of the family at level $n$.
Off $W$ these are the cubes of side $CM\ell_n$ with $\square\subset\Omega_n$ and
$\square\cap\Omega^{c}_{n+1}\neq\emptyset$; on $W$ the family is that of
\cite[(1.65), (1.67), (1.69)]{B-CMP122a}. The condition requires a gauge-group-valued $u$ on
$\square\cap Y$ with $U^u=e^{i\eta A}$ such that the displayed inequality holds. The constants $B$
and $C$ are those of \cite[(2.38)]{B-CMP119}.
\item On live pieces, the inequality in $(\Phi)$ expresses that the frozen profile only subtracts
from the native factor.
\end{itemize}
Of the factor $\Phi_e$ only property $(\Phi)$ is used below. Hence nothing below depends on the
form in which the enlargement of the radii in \cite[(3.43)]{B-CMP119} is written (radii ``bigger
by $\beta$ multiplied by a corresponding number'', in the words of \cite{B-CMP119}), as long as
$(\Phi)$ holds.

\emph{Levels and weights.} Inside $W$ the native indices are those fixed in the construction of
the native label sequence on arrested components, by the relations
\begin{gather*}
Z''_i=Z'^{\,\sim3}_{i-1}\ (h<i\le k_0),\qquad Z''_{k_0+1}=Z'^{\,\sim1}_{k_0},\qquad
Z''_{k_W}=Z'^{\,\sim3}_{k_0},\\
(\Omega^{\natural}_{i+1})^{c}\cap W=Z'^{\,\sim8}_i,\qquad Z'_i\supset\Lambda_i^{c}.
\end{gather*}
In the fourth, fifth and sixth rows of the table below, and in the display for the fourth row
that follows it, the unprimed regions $\Omega_m$ are those of the quotation from
\cite[p.~190]{B-CMP122a}, that is, the label sequence of $W$ before the arrest; on the fifth row
the index $m$ of this sequence is the native index $l$, while $n$ is the index of the current
label sequence of Notation~\ref{4.29:not-scales-labels}.
The level $n$, the native index $l$, and the weights $\varpi(x)$ and $\varpi_\square$ are given by
the following table.
\begin{equation}\label{4.29:eq-analyticity-space-b}
\begin{array}{c|l|c|c|c}
 & \text{region} & n & l & \varpi\ (\text{and }\varpi_\square)\\ \hline
1 & \text{off live pieces} & n & n & 1\\
2 & W:\ \Gamma''_i,\ i\le k_0;\ \text{old core} & i & i\text{ or }i-1 & 1\\
3 & W:\ \Gamma''_i,\ k_0<i<j_W & i & k_0 & L_\circ^{2(i-k_0)}\\
4 & W:\ \Omega''_{j_W}\setminus\Omega_{k_0+1} & j_W & k_0 & L_\circ^{2(j_W-k_0)}\\
5 & W:\ \Omega_m\setminus\Omega_{m+1},\ k_0<m\le j_W & j_W & m & L_\circ^{2(j_W-m)}\\
6 & W:\ \Omega_m\setminus\Omega_{m+1},\ j_W<m\le k_W & m & m & 1
\end{array}
\end{equation}
The third row consists of the thin shells. They satisfy
\begin{equation*}
\Gamma''_i\subset Z'^{\,\sim3}_{k_0}\setminus Z'^{\,\sim1}_{k_0}
\subset\Lambda_{k_0}\setminus\Omega^{\natural}_{k_0+1}.
\end{equation*}
In the fourth row,
\begin{equation*}
\Omega''_{j_W}\setminus\Omega_{k_0+1}\subset Z'^{\,\sim8}_{k_0}\setminus Z'^{\,\sim1}_{k_0}.
\end{equation*}
In the fifth row the index $m$ of the label sequence of $W$ before the arrest is the native index
$l$, and the row is point (ii) of the quotation
above, that is, \cite[(1.66), (1.67)]{B-CMP122a} with $j_W$ in place of $j$. The third to fifth
rows are void if $j_W\le k_0$. Where arrests are nested, the
structure at a point $x$ is that of the one arrest that lifted $x$, by the lemma on the label
sequences at the $\mathbf{T}$-stage and the construction of the arrest. With these weights,
\begin{equation}\label{4.29:eq-analyticity-space-1}
\varpi(x)\le L_\circ^{2(n(x)-l(x))}\quad\text{for every }x,
\end{equation}
with equality in the third, fourth and fifth rows of~\eqref{4.29:eq-analyticity-space-b}. The
corresponding real conditions \cite[(1.24), (1.49), (1.51)]{B-CMP122a} carry, by the construction
of the native label sequence on arrested components, the smaller weight
$L_\circ^{2\max\{0,n-l-1\}}$.
\end{definition}

\begin{proof}[Proof of the weight bound~\eqref{4.29:eq-analyticity-space-1}]
We first assign to each point $x$ of $Y$ a row of~\eqref{4.29:eq-analyticity-space-b}. By the
lemma on the label sequences at the $\mathbf{T}$-stage, $n(x)>l(x)$ only inside arrested
components, and each point is lifted by at most one arrest. Where arrests are nested,
the level and the weight at $x$ are therefore read from~\eqref{4.29:eq-analyticity-space-b} with
the data $(h,j_W,k_W,k_0)$ of the unique arrest that lifted $x$; a point of an arrested
component $W$ with $n(x)=l(x)$ is read with the data of the arrest of $W$. Points off live pieces
are in the first row. By the
construction of the arrest, for every arrest the level and the weight are frozen at a single
stage of the arrest, so the reading of $x$ in~\eqref{4.29:eq-analyticity-space-b}
is the same on every branch. It therefore suffices to compare $\varpi$ with
$L_\circ^{2(n-l)}$ row by row in~\eqref{4.29:eq-analyticity-space-b}.

\emph{First row.} Here $n=l$ and $\varpi=1=L_\circ^{0}$.

\emph{Second row.} Here $n=i$ and $l\in\{i,i-1\}$, so $n-l\in\{0,1\}$. Since $L_\circ\ge1$, we
get $\varpi=1\le L_\circ^{2(n-l)}$.

\emph{Third row.} Here $n=i$ and $l=k_0$, so $L_\circ^{2(n-l)}=L_\circ^{2(i-k_0)}=\varpi$.

\emph{Fourth row.} Here $n=j_W$ and $l=k_0$, so $L_\circ^{2(n-l)}=L_\circ^{2(j_W-k_0)}=\varpi$.

\emph{Fifth row.} Here $n=j_W$ and $l=m$, so $L_\circ^{2(n-l)}=L_\circ^{2(j_W-m)}=\varpi$.

\emph{Sixth row.} Here $n=l=m$ and $\varpi=1=L_\circ^{0}$.

The cube weights $\varpi_\square$ are read from the same rows, so the same comparison applies to
them. This proves~\eqref{4.29:eq-analyticity-space-1}, with equality in the third, fourth and
fifth rows.
\end{proof}

\begin{lemma}[Conversion of weighted radii to a native scale]\label{4.29:lem-radius-conversion}
Let $k\ge 0$ and write $x_i=g_i^{-2}$ for the running inverse couplings. Assume the following.
\begin{itemize}
\item[(a)] The flow bounds \eqref{4.29:eq-scales-labels-b} of Notation~\ref{4.29:not-scales-labels}
hold. That is, with the constants $\lambda_\sharp>0$, $B_\sharp>0$ and $c_5\ge 0$ fixed there, for all indices
$i\le j\le k+1$,
\begin{equation*}
\lambda_\sharp(j-i)-c_5\le x_i-x_j\le B_\sharp(j-i).
\end{equation*}
\item[(b)] $L_\circ$ is a positive integer with $L_\circ^2\le L/3$.
\item[(c)] $x_{k+1}\ge B_\sharp+c_5$.
\item[(d)] Let $\beta_0>0$, and let $\gamma\in(0,1)$ be so small that
\begin{equation*}
\Bigl(\frac{\log(x+c_5)}{\log x}\Bigr)^{q'}\le 1+\beta_0
\quad\text{for every } x\ge\gamma^{-2} \text{ and each } q'\in\{p_0,q_0,q_1\},
\end{equation*}
where $p_0,q_0,q_1$ are the exponents entering the sequences of \cite[(2.4), (2.28)]{B-CMP119}.
Such a smallness threshold exists, depends only on $p_0,q_0,q_1,c_5,\beta_0$, and is written
$\gamma(p_0,q_0,q_1,c_5,\beta_0)$. Assume moreover that the couplings lie in the regime
$x_i\ge\gamma^{-2}$ for all $i\le k+1$.
\item[(e)] $(a_n)$ is one of the three sequences $(\varepsilon_n)$, $(\alpha_{0,n})$,
$(\alpha_{1,n})$ of \cite[(2.4), (2.28)]{B-CMP119}.
\item[(f)] $p\ge 1$ is a real number, and $0\le m\le l\le n\le k+1$ are integers.
\end{itemize}
Then
\begin{equation}\label{4.29:eq-radius-conversion-a}
L_\circ^{2(n-l)}\,a_n\,L^{-p(n-m)}\le (1+\beta_0)\,3^{-(n-m)}\,(1+n-m)^{1/2}\,a_m .
\end{equation}
\end{lemma}

\begin{proof}
\emph{The powers of $L_\circ$ and $L$.}
By (b), $L\ge 3L_\circ^2\ge 3$. In particular $L\ge 1$, so $p\ge 1$ gives
$L^{-p(n-m)}\le L^{-(n-m)}$. Since $n-m=(n-l)+(l-m)$, we obtain
\begin{equation}\label{4.29:eq-radius-conversion-1}
L_\circ^{2(n-l)}L^{-p(n-m)}
\le \Bigl(\frac{L_\circ^2}{L}\Bigr)^{n-l}L^{-(l-m)}
\le 3^{-(n-l)}\,3^{-(l-m)}
=3^{-(n-m)} .
\end{equation}
Here the second inequality uses $L_\circ^2/L\le 1/3$ from (b) and $L\ge 3$, together with
$n-l\ge 0$ and $l-m\ge 0$.

\emph{The ratio of the radii.}
By their definitions in \cite[(2.4), (2.28)]{B-CMP119}, for each of the three sequences in (e) there
is an exponent $q'\in\{p_0,q_0,q_1\}$ such that
\begin{equation}\label{4.29:eq-radius-conversion-2}
\frac{a_n}{a_m}
=\Bigl(\frac{x_m}{x_n}\Bigr)^{1/2}\Bigl(\frac{\log x_n}{\log x_m}\Bigr)^{q'} .
\end{equation}
We bound the two factors separately.

First, we show that $x_n\ge B_\sharp$. Apply the lower flow bound of (a) to the pair $n\le k+1$. It gives
$x_n-x_{k+1}\ge\lambda_\sharp(k+1-n)-c_5\ge -c_5$, because $\lambda_\sharp(k+1-n)\ge 0$. Hence, by (c),
\begin{equation*}
x_n\ge x_{k+1}-c_5\ge B_\sharp .
\end{equation*}
Next, apply the upper flow bound of (a) to the pair $m\le n$. It gives $x_m-x_n\le B_\sharp(n-m)$. Since
$B_\sharp\le x_n$,
\begin{equation*}
x_m\le x_n+B_\sharp(n-m)\le x_n(1+n-m).
\end{equation*}
Therefore $(x_m/x_n)^{1/2}\le(1+n-m)^{1/2}$.

Second, apply the lower flow bound of (a) to the pair $m\le n$. It gives
$x_m-x_n\ge\lambda_\sharp(n-m)-c_5\ge -c_5$, that is, $x_n\le x_m+c_5$. By (d) we have
$x_m,x_n\ge\gamma^{-2}>1$, so $\log x_m>0$ and $\log x_n>0$. Since $q'>0$, the logarithm and the map
$t\mapsto t^{q'}$ are increasing. Hence
\begin{equation*}
(\log x_n)^{q'}\le\bigl(\log(x_m+c_5)\bigr)^{q'}\le(1+\beta_0)(\log x_m)^{q'} .
\end{equation*}
The last inequality is the displayed inequality of (d), applied at
$x=x_m\ge\gamma^{-2}$ with the exponent $q'$ of \eqref{4.29:eq-radius-conversion-2}. Inserting both
bounds into \eqref{4.29:eq-radius-conversion-2} gives
\begin{equation*}
a_n\le(1+\beta_0)(1+n-m)^{1/2}\,a_m .
\end{equation*}

\emph{Conclusion.}
Multiply this inequality by $L_\circ^{2(n-l)}L^{-p(n-m)}>0$ and use \eqref{4.29:eq-radius-conversion-1}.
This yields \eqref{4.29:eq-radius-conversion-a}.

Hypothesis (b) is the condition under which \cite{B-CMP122a} states that its sequences of radii are
descending; it is imposed there, together with $\beta\le 1/4$, on an integer that is written
$L_\circ$ here, and \cite{B-CMP122a} attributes the restriction on that integer to the bounds for
the field written $A'$ there. Inequality \eqref{4.29:eq-radius-conversion-a} is that condition read at
a single point, and the case $p=1$ is the one that enters through the bounds for $A'$.
\end{proof}

\begin{proposition}[The complex hypothesis at a cover cube]\label{4.29:prop-complex-hypothesis}
Assume the hypotheses of Proposition~\ref{4.29:prop-native-cover} and the conditions on $L_\circ$,
$x_{k+1}$ and $\gamma$ among the hypotheses of Lemma~\ref{4.29:lem-radius-conversion}, $L_\circ$ being
the constant of that lemma.
Let the cover $\mathcal{C}_j$, the scale of its cubes, the core $\square_0$ and the enlargements
$\tilde\square^{r}$ of a cover cube $\square$, the regions $\Lambda_j\subset\Lambda_j^{\sim}$, the
native small-field regions $\Omega^{\natural}_n$ and the cube radii $R_i$ be those of
Proposition~\ref{4.29:prop-native-cover}.
Let the space $\mathfrak{S}(Y)$, the quantities $\mathcal{Q}_1,\dots,\mathcal{Q}_7$, the bracket members
$\mathfrak{r}_e(n)$ with
their homogeneities $\nu_e$ in $L^n$, the factors $\Phi_e$ and $\Phi^{\mathrm{nat}}_e$, the level $n(x)$
and the native index $l(x)$ of a point be those of Definition~\ref{4.29:def-analyticity-space}, with the
levels, native indices and weights of the table~\eqref{4.29:eq-analyticity-space-b}; the weight is
written $\varpi(x)$ at a point and $\varpi_Q$ at a cube $Q$ of Ba{\l}aban's family. Let the regions
$\Omega_n$, the shells $\Gamma_i$ and $\Gamma''_i$, the bond sets $\mathbb{B}_p(X)$, the configurations
$U_{p,X}(M^{\cdot}\mathbb{U})$, the factorisation $\mathbb{U}=U'U$ with $U$ $G$-valued and
$U'=e^{i\eta A'}$, and the constant $B$ of \cite[(2.38)]{B-CMP119} also be those of
Definition~\ref{4.29:def-analyticity-space}. For every level
$n$, $\ell_n=L^n\eta$ is the unit of length at level $n$, $\eta$ being the lattice spacing, and $\pi_n$ is
the grid of cubes of side $M\ell_n$. The class $\mathbf{D}_j$ is the class of localisation domains at
level $j$ of Definition~\ref{4.4:def-settings-data-domains}. Assume the following statements.
\begin{itemize}
\item[(I)] \emph{Localisation domains.} For every $p\le k+1$ and every localisation domain
$X'\subset Y$, the clauses $e=2,4$ of \eqref{4.29:eq-analyticity-space-a} for the pair $(p,X')$, that
is, the bounds on $|\partial U_{p,X'}(M^{\cdot}\mathbb{U})-1|$ and on
$|J_{p,X'}(M^{\cdot}\mathbb{U})|$, belong to $\mathfrak{S}(Y)$. This statement is an input of the
reading of $\mathfrak{S}(Y)$ in Definition~\ref{4.29:def-analyticity-space} and is taken here as a
hypothesis.
\item[(II)] \emph{The group-valued factor.} The clauses $e=2,4$ of
\eqref{4.29:eq-analyticity-space-a} hold also with $U$ in place of $\mathbb{U}$. This statement is an
input of the reading of $\mathfrak{S}(Y)$ in Definition~\ref{4.29:def-analyticity-space} and is taken
here as a hypothesis.
\item[(III)] \emph{Siting of the localised clauses.} The clause for $(p,X')$ is imposed at the points of
a set $X'^{\flat}$, and $X'^{\flat}\supset X'^{\sim-2}$ whenever $\ell_{m_{X'}}\ge\ell_p$, where
$m_{X'}$ is the level of the localisation domain $X'$ and $\ell_{m_{X'}}$ its unit. The siting is part
of the reading of $\mathfrak{S}(Y)$ in Definition~\ref{4.29:def-analyticity-space}; the inclusion is
assumed here.
\item[(IV)] \emph{Arrested pieces.} On a live arrested piece $W$ the space $\mathfrak{S}(Y)$ has the
arrested form used in Definition~\ref{4.29:def-analyticity-space}: the structure of the clauses
(level, weight, cube family) of \cite[(1.64)--(1.69)]{B-CMP122a}, with the indices $j_W$, $k_W$, $k_0$
of the piece, is frozen, and the factor is only lowered.
\item[(V)] \emph{Separation of consecutive small-field regions.} For every level $n$,
$\operatorname{dist}_\infty(\Omega_{n+2},\Omega^c_{n+1})\ge M\ell_{n+1}$. This separation is a
clause of the lemma on the small-field regions that is among the hypotheses of
Proposition~\ref{4.29:prop-native-cover}.
\item[(VI)] \emph{The cube radii.} $R_m\ge R_{k_0+1}/L$ for $m>k_0+1$; this lower bound on
the radii $R_i$ is assumed here.
\item[(VII)] \emph{Reading of the gauge condition.} Condition \cite[(1.12)]{B-CMP109} is read in the
form \cite[(2.38)]{B-CMP119}; this is the one reading of \cite[(1.12)]{B-CMP109} adopted in this
chapter.
\end{itemize}
Let $\beta\le\frac14$ and let the constant $\beta_0$ of Lemma~\ref{4.29:lem-radius-conversion} satisfy
$\beta_0\le\frac1{10}$. Let $\square\in\mathcal{C}_j$ have scale $\sigma$
(default or not), let $Y\supset\square_0$, let $(\mathbb{U},\mathbb{J})\in\mathfrak{S}(Y)$, and put
\begin{equation}\label{4.29:eq-complex-hypothesis-1}
\theta_1:=(1+\beta)(1-\beta)^{-1}(1+\beta_0)\frac{\sqrt2}{3}\qquad(\le 0.87).
\end{equation}
Then the following hold.
\begin{itemize}
\item[(K1)] For $e=1,\dots,7$ and every $x\in\square_0$,
\[
\mathcal{Q}_e(x)<\Phi^{\mathrm{nat}}_e(\sigma)\,\mathfrak{r}_e(\sigma),
\]
with the further factor $\theta_1$ on the right at every $x\in\square_0$ with $n(x)>\sigma$. For
$e=2,4$ this holds for every $p\le\sigma$ and every localisation domain $X'\subset\square_0$, in
particular for $X'=\square_0$ on $\tilde\square^{3}$, where
\[
U_{p,\square_0}(M^{\cdot}\mathbb{U})=U\bigl(\mathbb{B}_p(\square_0),M^{\cdot}\mathbb{U}\bigr)
\]
is the function of \cite[(1.15)]{B-CMP109}; and likewise with $U$ in place of $\mathbb{U}$.
\item[(K2)] For $C\in\{1,2\}$, every cube $q'\subset\square_0$ of side $CM\ell_\sigma$, aligned with
$\pi_\sigma$ or not (the cover cube $\square$ itself is one of them), carries a $G$-valued gauge
transformation $u$ such that on $q'$
\[
U^u=e^{i\eta A},\qquad \ell_\sigma|A|<BCM\alpha_{0,\sigma},\qquad
\ell_\sigma^2|\nabla^{\eta}A|<BCM\alpha_{0,\sigma}.
\]
This covers also the cubes $q''$ of Ba{\l}aban's family with $q''\not\subset\square_0$ and
$q''\cap\square_0\neq\emptyset$: translating $q''$ along each coordinate axis inwards until it lies in
$\square_0$ (possible because $\square_0$ has side $12M\ell_\sigma\ge CM\ell_\sigma$) gives a cube
$q'\subset\square_0$ of the same side with $q'\supset q''\cap\square_0$, and the restriction of the
gauge transformation of $q'$ serves.
\item[(K3)] (Membership.) The same holds on $X$ at level $j$, for $X\in\mathbf{D}_j$ with
$X\subset\Lambda_j\cap Y$: (K1) with $\sigma$ replaced by $j$ on $X$, and (K2) for every cube
$q'\subset X$ of side $CM\ell_j$, $C\in\{1,2\}$. Likewise at every point of $\Lambda_j^{\sim}\cap Y$,
pointwise, and at finer levels a fortiori, since \cite{B-CMP99a} states of these regularity
conditions that
\begin{quote}
if these conditions hold for some $j=l$, then they hold for all $j<l$.
\end{quote}
Hence the operators $\mathcal{H}'$ and $Q'$ of the step, which are built on backgrounds
satisfying these regularity conditions, are holomorphic on $\mathfrak{S}(Y)$. As for the cubes
\emph{of a size $O(1)LM$} contained in $X$ of \cite[(1.12)]{B-CMP109}: under the reading (VII), the
cube clause of (K3) is the whole clause, and the cubes of side a fixed multiple of $LM\ell_j$ contained
in $X$ and meeting $\Gamma_j$ of \cite{B-CMP119}, which are not supplied by $C\le2$, lie at points
whose native index equals the level, so that under (VII) no condition on them remains.
\end{itemize}
In other words,
\begin{equation}\label{4.29:eq-complex-hypothesis-a}
\mathfrak{S}(Y)\restriction\square_0\;\subset\;\mathfrak{N}_\sigma(\square_0)\;\subset\;
\mathcal{U}^c_\sigma\bigl(\square_0,(1+2\beta)\alpha_{0,\sigma},(1+2\beta)\alpha_{1,\sigma},\cdot\bigr),
\end{equation}
where $\mathfrak{N}_\sigma(\square_0)$ is the space of Definition~\ref{4.29:def-analyticity-space} for
a label (in the sense of Definition~\ref{4.29:def-analyticity-space}) for which $\square_0$ lies in a
stratum of level $\sigma$ that is never arrested, and $\restriction$
denotes reading on $\square_0$ and restricting to the clauses with $p\le\sigma$. The triples with
$p>\sigma$ coincide with those with $p=\sigma$ on $\square_0\subset\Omega_\sigma$, since $J$ and
$\mathfrak{r}_4$
both scale by $L^{3(p-\sigma)}$. The space on the right is Ba{\l}aban's, defined by the conditions
(i)--(iv) of \cite[(1.11)--(1.16)]{B-CMP109} at $X=\square_0$, with $j$ replaced by $\sigma$ and unit
$\ell_\sigma$; its requirement that (i)--(iii) hold \emph{on the whole lattice} \cite{B-CMP109} is read,
as in the transition from \cite[(6.34)]{B-CMP109} to \cite[(3.43)]{B-CMP119}, as membership in
$\mathfrak{S}(Y)$, and its constant for $\mathbb{J}$ is the native one,
$\Phi^{\mathrm{nat}}_3(\sigma)\alpha_{0,\sigma}$.
\end{proposition}

\begin{proof}
By clause (c) of Proposition~\ref{4.29:prop-native-cover}, $l(x)\ge\sigma$ and $n(x)\ge\sigma$ for
every $x\in\square_0$, and $\square_0\subset\Omega_\sigma$.

\emph{Two elementary facts.} First, for every integer $u\ge1$,
\begin{equation}\label{4.29:eq-complex-hypothesis-2}
3^{-u}(1+u)^{1/2}\le\frac{\sqrt2}{3}:
\end{equation}
equality holds at $u=1$, and passing from $u$ to $u+1$ multiplies the left-hand side by
$\frac13\bigl((u+2)/(u+1)\bigr)^{1/2}<1$. Secondly, by \eqref{4.29:eq-complex-hypothesis-1},
$\theta_1(1-\beta)=(1+\beta)(1+\beta_0)\sqrt2/3$.

\emph{(K1) for $e\neq2,4$.} Let $x\in\square_0$. In every row of the
table~\eqref{4.29:eq-analyticity-space-b} the native index does not exceed the level, and the weight
obeys $\varpi(x)\le L_\circ^{2(n(x)-l(x))}$.

Suppose $n(x)=\sigma$. Then $\sigma\le l(x)\le n(x)=\sigma$, so $l(x)=\sigma$, and $\varpi(x)=1$ by
the table. Clause $(\mathfrak{S}1)$ of \eqref{4.29:eq-analyticity-space-a} gives
$\mathcal{Q}_e(x)<\Phi_e(x)\mathfrak{r}_e(\sigma)$. By $(\Phi)$ of \eqref{4.29:eq-analyticity-space-a},
$\Phi_e(x)\le\Phi^{\mathrm{nat}}_e(\sigma)$, which is the claim.

Suppose $n:=n(x)>\sigma$, and write $l=l(x)$. By $(\mathfrak{S}1)$,
$\mathcal{Q}_e(x)<\Phi_e(x)\varpi(x)\mathfrak{r}_e(n)$. By $(\Phi)$,
$\Phi_e(x)\le\Phi^{\mathrm{nat}}_e(n)\le1+\beta$, and
by the table, $\varpi(x)\le L_\circ^{2(n-l)}$. The bracket member $\mathfrak{r}_e(n)$ is a multiple,
independent of
$n$, of $a_nL^{-\nu_en}$, where $(a_n)$ is one of $(\alpha_{0,n})$, $(\alpha_{1,n})$ and $\nu_e\ge1$ is its
homogeneity. Hence
\[
L_\circ^{2(n-l)}\mathfrak{r}_e(n)=L_\circ^{2(n-l)}a_nL^{-\nu_e(n-\sigma)}\,
\frac{\mathfrak{r}_e(\sigma)}{a_\sigma}.
\]
Lemma~\ref{4.29:lem-radius-conversion} applies with $m:=\sigma$ and $p:=\nu_e$, since
$\sigma\le l\le n$. Its conclusion, multiplied by $\mathfrak{r}_e(\sigma)/a_\sigma$ and combined with the
identity above, gives
\begin{align*}
\mathcal{Q}_e(x)&<(1+\beta)L_\circ^{2(n-l)}\mathfrak{r}_e(n)
\le(1+\beta)(1+\beta_0)3^{-(n-\sigma)}(1+n-\sigma)^{1/2}\mathfrak{r}_e(\sigma)\\
&\le\theta_1(1-\beta)\,\mathfrak{r}_e(\sigma)\le\theta_1\,\Phi^{\mathrm{nat}}_e(\sigma)\,\mathfrak{r}_e(\sigma).
\end{align*}
The third inequality is \eqref{4.29:eq-complex-hypothesis-2} with $u=n-\sigma\ge1$, combined with
the identity for $\theta_1(1-\beta)$. The last one is $1-\beta\le\Phi^{\mathrm{nat}}_e(\sigma)$,
from $(\Phi)$. Since $\theta_1<1$, the bound without the factor $\theta_1$ follows as well.

\emph{(K1) for $e=2,4$.} Let $p\le\sigma$. Let $X'\subset\square_0\subset Y$ be a localisation domain.

By hypothesis (I), the clauses $e=2,4$ of \eqref{4.29:eq-analyticity-space-a} for the pair $(p,X')$
belong to $\mathfrak{S}(Y)$. Since $\square_0\subset\Omega_\sigma$, we have $\Gamma_i\cap X'=\emptyset$
for $i<p\le\sigma$, and so
\[
\mathbb{B}_p(X')\cup\{\Gamma_i\}_{i<p}=\mathbb{B}_p(X').
\]
Thus the function $U_{p,X'}(M^{\cdot}\mathbb{U})$ of the clause is $U(\mathbb{B}_p(X'),M^{\cdot}\mathbb{U})$,
which for $X'=\square_0$ is the function of \cite[(1.15)]{B-CMP109}.

By hypothesis (III), applied with $p\le\sigma$, the only range of $p$ used here, the bound is sited on
$X'^{\flat}\supset X'^{\sim-2}$. It converts as in the case
$e\neq2,4$: for $p\le\sigma\le n$, the members $\mathfrak{r}_2$ and $\mathfrak{r}_4$ have homogeneity
$2$, because the
factor $L^{-(\min\{p,n\}-\min\{p,\sigma\})}$ in $\mathfrak{r}_4(n)/\mathfrak{r}_4(\sigma)$ equals $1$. The
bounds with $U$ in
place of $\mathbb{U}$ follow in the same way from hypothesis (II).

\emph{(K2).} Let $q'\subset\square_0$ have side $CM\ell_\sigma$, $C\in\{1,2\}$, and put
$n':=\min_{q'}n(\cdot)$. By clause (c) of Proposition~\ref{4.29:prop-native-cover}, $n'\ge\sigma$ and
$l\ge\sigma$ on $q'$. Every point of $q'$ has level at least $n'$, so $q'\subset\Omega_{n'}$.

\emph{Case $n'=\sigma$.} Here $q'\subset\Omega_\sigma$, and the point of $q'$ at which the minimum
is attained gives $q'\cap\Omega^c_{\sigma+1}\neq\emptyset$. Moreover $q'\cap\Omega_{\sigma+2}=\emptyset$
by hypothesis (V), the inequality $L>2$ and $C\le2$:
\[
\operatorname{dist}_\infty(\Omega_{\sigma+2},\Omega^c_{\sigma+1})\ge M\ell_{\sigma+1}
>2M\ell_\sigma\ge\operatorname{diam}_\infty q'.
\]
Hence $q'$ is itself a cube of Ba{\l}aban's family at level $\sigma$, with its own $C$.

Its weight is $1$. A weighted clause at level $\sigma$ (rows three to five of the
table~\eqref{4.29:eq-analyticity-space-b}) would need a point of $q'$ with $l<\sigma$, and this is
excluded by clause (c) of Proposition~\ref{4.29:prop-native-cover}. Clause $(\mathfrak{S}2)$ of
\eqref{4.29:eq-analyticity-space-a} on $q'\cap Y=q'$ is then the assertion (K2).

\emph{Case $n'>\sigma$.} Let $\mathfrak{b}$ be the box formed by the cubes of $\pi_{n'}$ that meet $q'$.
Each of them contains a point of $q'\subset\Omega_{n'}$, and $\Omega_{n'}$ is a union of cubes of
$\pi_{n'}$; hence $\mathfrak{b}\subset\Omega_{n'}$. As $2M\ell_\sigma<M\ell_{n'}$, the box
$\mathfrak{b}$ has at most two cubes in each direction. We may therefore choose a cube $Q$ of side
$M\ell_{n'}$, in general not aligned with $\pi_{n'}$, with $q'\subset Q\subset\mathfrak{b}$.

Then $Q\subset\Omega_{n'}$, and $Q\cap\Omega^c_{n'+1}\supset q'\cap\Omega^c_{n'+1}\neq\emptyset$.
By hypothesis (V), since $M\ell_{n'+1}>\operatorname{diam}_\infty Q$, also
$Q\cap\Omega_{n'+2}=\emptyset$. So $Q$ is a cube of Ba{\l}aban's family at level $n'$ with $C=1$.
Clause $(\mathfrak{S}2)$ gives a $G$-valued $u$ on $Q\cap Y\supset q'$ with $U^u=e^{i\eta A}$ and
\[
\ell_{n'}|A|<\varpi_Q\,BM\alpha_{0,n'},\qquad \ell_{n'}^2|\nabla^{\eta}A|<\varpi_Q\,BM\alpha_{0,n'}.
\]

\emph{The weight of $Q$.} We show that $\varpi_Q\le L_\circ^{2(n'-l_Q)}$ for an index $l_Q\ge\sigma$,
examining the rows of the table~\eqref{4.29:eq-analyticity-space-b} under the arrested form of
hypothesis (IV).

\emph{Off $W$, and in rows two and six.} Here $\varpi_Q=1$, and we take $l_Q:=n'$.

\emph{Row three.} Here $\varpi_Q=L_\circ^{2(n'-k_0)}$. The cube $q'$ contains a point of $\Gamma''_{n'}$,
of native index $k_0$, and $l\ge\sigma$ on $q'$; so $l_Q:=k_0\ge\sigma$.

\emph{Rows four and five.} Here the level is $n'=j_W$. For $k_0<m\le k_W$, the region
$\Omega^{\natural}_m$ is a union of cubes of side $MR_m\ell_m\in M\ell_{k_W}L^{\mathbb{N}_0}$, because
$R_{k_0+1}\ell_{k_0+1}=\ell_{k_W}$ and $R_m\ge R_{k_0+1}/L$ for $m>k_0+1$ by hypothesis (VI) (for
the cube structure of $\Omega^{\natural}_m$ and the radii see \cite{B-CMP122a} and \cite{B-CMP119}).
Consequently each cube of $\pi_{j_W}$ in $\mathfrak{b}$ lies in one native
stratum $m\ge k_0$ (the set of points of native index $m$), so that the weight of $Q$ depends on these
cubes only through the smallest native
index $\min_Ql$. Each of these cubes meets $q'$, where $l\ge\sigma$, so
$\min_Ql\ge\sigma$. By \cite[(1.65)]{B-CMP122a} (for $m=k_0$) or \cite[(1.67)]{B-CMP122a},
\[
\varpi_Q=L_\circ^{2(j_W-\max\{k_0,\min_Ql\})}\le L_\circ^{2(j_W-\min_Ql)},
\]
and we take $l_Q:=\min_Ql$.

In all cases, therefore,
\[
\varpi_Q\le L_\circ^{2(n'-l_Q)},\qquad l_Q\ge\sigma.
\]

\emph{Conversion to level $\sigma$.} Restrict $u$ to $q'$. Since $\ell_\sigma=L^{-(n'-\sigma)}\ell_{n'}$,
for $p\in\{1,2\}$
\[
\ell_\sigma^p|(\nabla^{\eta})^{p-1}A|=L^{-p(n'-\sigma)}\ell_{n'}^p|(\nabla^{\eta})^{p-1}A|
<L_\circ^{2(n'-l_Q)}\alpha_{0,n'}L^{-p(n'-\sigma)}BM.
\]
Apply Lemma~\ref{4.29:lem-radius-conversion}
with $n=n'$, $l=l_Q$, $m=\sigma$, $p=1,2$ and $a_n=\alpha_{0,n}$, and then
\eqref{4.29:eq-complex-hypothesis-2} with $u=n'-\sigma\ge1$. This gives
\begin{align*}
\ell_\sigma|A|,\ \ell_\sigma^2|\nabla^{\eta}A|
&<(1+\beta_0)3^{-(n'-\sigma)}(1+n'-\sigma)^{1/2}BM\alpha_{0,\sigma}\\
&\le BM\alpha_{0,\sigma}\le BCM\alpha_{0,\sigma},
\end{align*}
since $(1+\beta_0)\sqrt2/3<1$ and $C\ge1$.

\emph{The cubes $q''\not\subset\square_0$.} Let $q''$ meet $\square_0$. Each coordinate interval of
$q''$ has length $CM\ell_\sigma\le12M\ell_\sigma$, the length of the corresponding interval of
$\square_0$, and meets that interval. If the interval of $q''$ is translated by the least amount that
brings it inside the interval of $\square_0$, the translated interval contains the intersection of the
original one with the interval of $\square_0$. The translated cube $q'$ is contained in $\square_0$, has the
same side, and contains $q''\cap\square_0$; the restriction to $q''\cap\square_0$ of the gauge
transformation given by (K2) for $q'$ has the required properties.

\emph{(K3).} Since $X\subset\Lambda_j\subset\Omega^{\natural}_j$, we have $l\ge j$ on $X$. The proofs
of (K1) and (K2) above then apply verbatim with $\sigma$ replaced by $j$, the cubes $q'\subset X$ of
side $CM\ell_j$ included. On $\Lambda_j^{\sim}\cap Y$ the same argument applies pointwise, since
$l\ge j$ there by clauses (c) and $(e)^{\natural}$ of Proposition~\ref{4.29:prop-native-cover}.

\emph{The inclusions \eqref{4.29:eq-complex-hypothesis-a}.} The first inclusion is the content of
(K1)--(K3), the clauses with $p>\sigma$ being reduced to $p=\sigma$ as stated. For the second, check
the conditions of \cite[(1.11)--(1.16)]{B-CMP109} at $X=\square_0$ and level $\sigma$ one by one.
\begin{itemize}
\item (K1) at $e=5,6,7,1,3$ is \cite[(1.11), (1.13), (1.14)]{B-CMP109}, with
$\Phi\le1+\beta<1+2\beta$.
\item (K1) at $e=2,4$ with $X'=\square_0$ is \cite[(1.16)]{B-CMP109}.
\item No cube of \cite[(1.12)]{B-CMP109} of a size $O(1)LM$, that is, of side $\mathsf{o}LM\ell_\sigma$
with $\mathsf{o}$ the number of Definition~\ref{4.4:def-local-data-classes}, fits in $\square_0$, whose
side satisfies
\[
12M\ell_\sigma<LM\ell_\sigma\le\mathsf{o}LM\ell_\sigma,
\]
because $L\ge13$ by the hypotheses of Proposition~\ref{4.29:prop-native-cover} and $\mathsf{o}\ge1$.
The layered gauge clause is (K2) with its clause on the cubes $q''\not\subset\square_0$.
\end{itemize}
Here \cite[(1.12)]{B-CMP109} is read, by hypothesis (VII), in its successor form
\cite[(2.38)]{B-CMP119}, which \cite[\S2]{B-CMP119} introduces with the words
\begin{quote}
similarly as in (I.1.11)--(I.1.16), but now we must take into account the existence of many different
scales
\end{quote}
(here [I] is \cite{B-CMP109}) and describes, in the same section, as
\begin{quote}
a more detailed and precise form, suitable for our inductive constructions, of the regularity
conditions for background fields, \dots
\end{quote}
In that form the cubes have side $CM\ell_n$ with $C\le2$, the gauge transformation
$u$ is defined on the intersection of the cube with the domain, and the bound is
$\varpi_Q\,BCM\alpha_{0,n}$. Every cube of the family contained in the core $\square_0$ of a cover
cube has level $\sigma$, since a cube at a level $n\ge\sigma+1$ has side at least
$M\ell_{\sigma+1}=LM\ell_\sigma>12M\ell_\sigma$, and it has weight $\varpi_Q=1$, as in the case
$n'=\sigma$ above; so (K2) carries no weight factor.
\end{proof}

\begin{remark}[Ba{\l}aban's family of gauge cubes]
The cube family of (K2) is that of \cite[(2.38)]{B-CMP119}, which reads:
\begin{quote}
(ii) For each cube $\square\subset\Omega_n\setminus\Omega_{n+2}$,
$\square\cap\Omega^c_{n+1}\neq\emptyset$, $n=1,\dots,j-1$, of the size $CML^n\xi$, or
$\square\subset\Omega_j$ and of the size $CM$, there exists a gauge transformation $u$ defined on
$\square\cap X$ \dots
\end{quote}
No alignment of the cubes with a grid is required there. On the constant $C$, \cite{B-CMP122a} states:
\begin{quote}
The constant $C$ above is a positive integer. It is usually small, because of geometric constraints on
the cubes, for example we can assume that $C\le2$.
\end{quote}
Accordingly, $C\in\{1,2\}$ throughout.
\end{remark}

\begin{remark}[The factor $c=3$ at the last level]
In Definition~\ref{4.29:def-analyticity-space} the factor $c$ of \cite[(1.68)]{B-CMP122a} is taken
equal to $1$ on the live arrested pieces, and the proposition uses this value. If the value $c=3$, which
\cite[(1.68)]{B-CMP122a} carries at $m=k$, were kept there, it would sit only on
$W\cap\Omega^{\natural}_{k_W}$, where $n=l=k_W$. The core of a cover cube that is not default, if it
contains such a point, lies in $\Lambda_\sigma$ with $\sigma<k_W$, because $W\subset\Lambda^c_{k_W}$.
Moreover
\[
3(1+\beta)(1-\beta)^{-1}(1+\beta_0)\frac{\sqrt2}{L}\le1\qquad\text{for }L\ge9,
\]
which absorbs the factor $3$ in the conversion of (K1) at such a point, where $l=n$, so that no power
of $L_\circ$ occurs, and $n-\sigma\ge1$. The threshold $L\ge9$ lies below the condition $L\ge13$
among the hypotheses of Proposition~\ref{4.29:prop-native-cover}.
\end{remark}

\begin{proposition}[The minimiser as a native-scale configuration on the window]
\label{4.29:prop-window-minimiser}
Let $j\le k$, and let $M>R_{\mathrm{adm}}M_1$, where $R_{\mathrm{adm}}$ is the constant of the
admissibility condition of \cite{B-CMP99a}. Let $\square\in\mathcal{C}_j$ be a cube of scale $s$ of
the native-scale cover of
Proposition~\ref{4.29:prop-native-cover}, and let $\square_0$ be its window. Write $\mathbb{B}$ for the
determining set of the sequence $\{\Omega_i\}_{i\le k+1}$, and $\tilde{\mathbf{B}}_s$ for the
determining set of the truncated sequence $\{\Omega_i\}_{i\le s}$. Write $M_{\mathbb{B}}$,
$M_{\tilde{\mathbf{B}}_s}$ for the corresponding averaging maps and $U_{\mathbb{B}}(\cdot)$,
$U_{\tilde{\mathbf{B}}_s}(\cdot)$ for the corresponding minimal configurations, $M^i$ for the
$i$-fold block average, and $U_i(\square_0,\cdot)=U_{i,\square_0}(\cdot)$ for the minimiser of
\cite[(2.13)]{B-CMP119} for the domain $\square_0$ at level $i$. Assume the following statements.
\begin{enumerate}
\item[(1)] The criticality-transfer and uniqueness step of the construction of the minimal
configuration, in the form it takes for the step from level $k$ to level $k+1$: if the constraints
of one determining set are functions of those
of a second one, a configuration critical for the first is critical for the second; and under the
regularity of hypothesis~(3) the minimal orbit at given data is unique.
\item[(2)] The restriction lemma for minimisers: for a domain $D$, a determining set $\mathbf{B}'$ with
support in $D$ and a configuration $U^*$ regular as in hypothesis~(3), one has
$U_{\mathbf{B}'}(M_{\mathbf{B}'}(U^*))=(U^*\upharpoonright D)^u$ for a gauge transformation $u$,
blocks straddling the boundary of $D$ included.
\item[(3)] The regularity clause and the stratum-separation clause
of the lemma on the domains in force at the $\mathbf{T}$-stage: $U_{k+1}=U_{\mathbb{B}}(V)$ is the
minimal configuration of hypothesis~(4), with the data $V$, the block fields of the levels of
$\mathbb{B}$, carrying a tower at the parameter
$\varepsilon_1^{g}$ of the tower condition; its plaquette and current
constants are those of its own layer, namely $c_{\mathrm{reg}}\varepsilon_n$ at level $n$ off the
arrested components, and on them the real plaquette and current bounds of the large-field
construction on arrested components, with weights at most
$L_\circ^{2(n-l)}$ and constant at most $6$; in particular $U_{k+1}$ belongs to the class of
regular configurations of the minimiser theorem of hypothesis~(4) at the parameter
$B_3\varepsilon_1^{g}$, both defining conditions of that class holding at every level $\le s$,
and the block averages $M^i(U_{k+1})$, $i\le s$, satisfy the tower condition at the parameter
$2L^2B_3\varepsilon_1^{g}$;
$\operatorname{dist}_\infty(\Omega_{a+1},\Omega_a^{c})\ge M\ell_a$ for every $a$. Here $n=n(x)$ is the level of
$\mathbb{B}$ at the point $x$, $l=l(x)$ is the level function of that lemma, the exponent $l$ in
the weights $L_\circ^{2(n-l)}$, with $l\ge s$ on the window by
Proposition~\ref{4.29:prop-native-cover}, and $L_\circ$ is the base of the weights on the arrested
components.
\item[(4)] The minimiser theorem under a tower: for an admissible sequence of domains with
determining set $\mathbf{B}$, and data whose witnesses, the block averages $M^i$ of the data
configuration, satisfy the tower condition at a parameter
$\varepsilon_1'$, there is a unique critical orbit at radius $a_0$, the radius within which that
theorem asserts uniqueness.
\item[(5)] The real-point lemma for the pieces born at the step: for every domain
$Y\supset\square_0$, the real configuration $(U_{k+1},J_{k+1})$ of the step lies in the space
$\mathfrak{S}(Y)$ of Definition~\ref{4.29:def-analyticity-space} with all radii halved, in the
arrested form of that space on the arrested components and in its native form off them.
\item[(6)] The local form of the restriction identity: the identity of hypothesis~(2) holds at
$D=\square_0$ for every cube scale $m$ with $j\le m$ and $m$ not larger than the index of the
highest stratum $\Omega_i$ meeting $\square_0$.
\item[(7)] The hypotheses of Lemma~\ref{4.29:lem-radius-conversion} hold.
\end{enumerate}
Then the following hold on $\square_0$:
\begin{equation}\label{4.29:eq-window-minimiser-a}
\begin{aligned}
&\text{(d1)}\quad U_{k+1}=U_{\tilde{\mathbf{B}}_s}\bigl(M_{\tilde{\mathbf{B}}_s}U_{k+1}\bigr),\\
&\phantom{\text{(d1)}\quad} U_{k+1}\upharpoonright\square_0=U_i\bigl(\square_0,M^{\cdot}(U_{k+1})\bigr)
\ \text{up to gauge},\ i\le s;\\
&\text{(d2)}\quad \bigl|U_{k+1}(\partial p)-1\bigr|\le C_{\mathrm{reg}}\,\varepsilon_s\,\eta^2\,
\ell_s^{-2}\qquad (p\subset\square_0),
\end{aligned}
\end{equation}
where
\begin{equation*}
C_{\mathrm{reg}}:=(1+\beta_0)\max\{6,\,c_{\mathrm{reg}},\,O(1)B_3B_5M^6\}.
\end{equation*}
Here $\eta$ is the spacing of the fine lattice on which $U_{k+1}$ is defined, $B_3$ is the constant
of the regularity clause of hypothesis~(3), and $B_5$ is a constant of the minimiser construction;
$\beta_0$ and the sequence $(\varepsilon_n)$ are the constant and the radius sequence of
Lemma~\ref{4.29:lem-radius-conversion}, and $\ell_s$ is the native length of level $s$.
In \textup{(d1)} the whole window lies in the top stratum $\Omega_s$ of the truncated sequence, and
$M^{\cdot}(U_{k+1})$ is the sequence of block averages $M^i(U_{k+1})$; thus, up to gauge, the
restriction of $U_{k+1}$ to $\square_0$ is a minimiser of level $s$ on $\square_0$, namely
$U_s(\square_0,M^{\cdot}(U_{k+1}))$. In \textup{(d2)} the constant $C_{\mathrm{reg}}$
carries no factor $LR_n$, $R_n$ being the level-$n$ parameter of
Proposition~\ref{4.29:prop-native-cover}; moreover $(U_{k+1},J_{k+1})\upharpoonright\square_0$,
$J_{k+1}$ the current of the step, lies in the native analyticity space
$\mathfrak{N}_s(\square_0)$ of \eqref{4.29:eq-complex-hypothesis-a} with all radii halved.
\end{proposition}

\begin{proof}
\emph{Reading of \textup{(d1)}.}
The identity (d1) is the form of the representation $U_k=U_{j,\square_0}(M^{\cdot}(U_k))$ of
\cite{B-CMP119}, where, in the notation of that paper, $\square$ is a cube of the level-$j$
partition $\pi_j$ and $\square_0=\square^{\sim n'-j+1}$, $n'$ being the scale index there; in
\cite[(3.18)]{B-CMP109} the window is $\square_0=\tilde{\square}^5$. With
$U_j:=U_{j,\Lambda_j^{\sim}}$, hypothesis~(2) gives, on $\Lambda_j^{\sim}$,
$U_{j,\Lambda_j^{\sim}}(M'(U_{k+1}))=U_{k+1}$ up to gauge, where $M'$ is the averaging map of the
determining set $\mathbf{B}_j(\Lambda_j^{\sim})$; so (d1) is the form in which the analysis of
\cite[\S3]{B-CMP109} reads the minimiser at the cube. The representation \cite[(3.18)]{B-CMP109}
is entered at good cubes, in the sense of Proposition~\ref{4.29:prop-native-cover}, of scale $s>j$
and at cubes with $\square_0\subset\Omega_{k+1}$
(Lemma~\ref{4.29:lem-localisation-exact}); at these $\square_0\subset\Lambda_j\subset D_j$, where
$D_j$ is the domain of level $j$ of $\mathbf{B}_j(\Lambda_j^{\sim})$.

\emph{The first identity of \textup{(d1)}.} We apply hypothesis~(1) to the pair $\mathbb{B}$,
$\tilde{\mathbf{B}}_s$.
For $i>s$ one has $M^i=M^{i-s}\circ M^s$ on $\Gamma_i$, the set on which the constraints of level
$i$ are imposed; hence the constraints of $\mathbb{B}$ are
functions of those of $\tilde{\mathbf{B}}_s$. The configuration $U_{k+1}$, being critical for
$\mathbb{B}$, is therefore critical for $\tilde{\mathbf{B}}_s$. The truncated sequence
$\{\Omega_i\}_{i\le s}$ is admissible, by the separation clause of hypothesis~(3) and the
condition $M>R_{\mathrm{adm}}M_1$ of the statement. By the regularity clause of hypothesis~(3),
$U_{k+1}$ lies in the class of regular
configurations of the minimiser theorem at the parameter $B_3\varepsilon_1^{g}$, with both defining
conditions of that class at every level $\le s$; and, by the same clause, the truncated witnesses
$W_i:=M^i(U_{k+1})$, $i\le s$, satisfy the tower condition at
\begin{equation*}
\varepsilon_1'=2L^2B_3\varepsilon_1^{g}.
\end{equation*}
Hypothesis~(4), applied to $\tilde{\mathbf{B}}_s$ at the parameter $\varepsilon_1'$, gives a unique
critical orbit at radius $a_0$ with data $M_{\tilde{\mathbf{B}}_s}U_{k+1}$. Since $U_{k+1}$ is
critical for $\tilde{\mathbf{B}}_s$ with these data, it lies on this orbit; this is the first
identity of \eqref{4.29:eq-window-minimiser-a}.

\emph{The window lies in the top stratum.} The inclusion $\square_0\subset\Omega_s$ is the window
clause of Proposition~\ref{4.29:prop-native-cover}.

\emph{The local form.} Hypothesis~(2) at $D=\square_0$, together with hypothesis~(6), which covers
every cube scale $m$ with $j\le m$ up to the index of the highest stratum meeting $\square_0$, gives
$U_{k+1}\upharpoonright\square_0=U_i(\square_0,M^{\cdot}(U_{k+1}))$ up to gauge for every $i$ with
$j\le i\le s$, since $\square_0\subset\Omega_s$ makes $s$ not larger than the index of the highest
stratum meeting $\square_0$; for these $i$ this is the second
identity of \eqref{4.29:eq-window-minimiser-a}.

\emph{The bound \textup{(d2)}.} By the regularity clause of hypothesis~(3), at a point of level $n$
the plaquette deviation of $U_{k+1}$ is bounded with the constant of its own layer: $c_{\mathrm{reg}}
\varepsilon_n$ off the arrested components, and on them the real plaquette and current bounds of the
large-field construction on arrested components,
with weights at most $L_\circ^{2(n-l)}$ and constant at most $6$. The real column of the
analyticity space, \eqref{4.29:eq-analyticity-space-b}, reads these bounds at the native length
$\ell_n$ of level $n$, in units of $\eta^2\ell_n^{-2}$; the entries $6$ and $c_{\mathrm{reg}}$ of the
maximum defining $C_{\mathrm{reg}}$ are the constants of the regularity clause, and the third entry
$O(1)B_3B_5M^6$ collects the constant $B_3$ of that clause, the constant $B_5$ of the minimiser
construction and the power $M^6$. To pass to the native scale $s$
of the window we use
Lemma~\ref{4.29:lem-radius-conversion} with $a=\varepsilon$, $p=2$ and $m=s$. On $\square_0$ one has
$l\ge s$ by the window clause of Proposition~\ref{4.29:prop-native-cover}, and $n\ge s$ since
$\square_0\subset\Omega_s$. Wherever the weight $L_\circ^{2(n-l)}$ exceeds $1$ one has $l<n$, hence
$s\le l<n$, and the lemma applies with this $l$; elsewhere the weight is at most $1$, and the lemma
applies with $l=n$. In either case, with $l$ replaced
by $n$ in the second, the lemma gives
\begin{equation*}
L_\circ^{2(n-l)}\varepsilon_n L^{-2(n-s)}\le(1+\beta_0)\,3^{-(n-s)}(1+n-s)^{1/2}\varepsilon_s
\le(1+\beta_0)\,\varepsilon_s,
\end{equation*}
the last step because $3^{r}\ge 1+r\ge(1+r)^{1/2}$ for every integer $r\ge0$. Since
$\ell_n=L^{n-s}\ell_s$, the factor $L^{-2(n-s)}$ is $\ell_s^2\ell_n^{-2}$, so a bound in units of
$\eta^2\ell_n^{-2}$ at level $n$ becomes a bound in units of $\eta^2\ell_s^{-2}$ at level $s$. This
gives \textup{(d2)} with the constant $C_{\mathrm{reg}}$ of the statement.

\emph{Membership.} Let $Y\supset\square_0$. By hypothesis~(5), the real configuration
$(U_{k+1},J_{k+1})$ of the step lies in $\mathfrak{S}(Y)$ with all radii halved, both on the
arrested components and off them. Proposition~\ref{4.29:prop-complex-hypothesis} applies to pairs
in $\mathfrak{S}(Y)$, and its proof is
homogeneous in the radii; carried out with all radii halved and applied to $(U_{k+1},J_{k+1})$, it
places $(U_{k+1},J_{k+1})\upharpoonright\square_0$ in $\mathfrak{N}_s(\square_0)$ with all radii
halved.
\end{proof}

\begin{definition}[Path lengths in local cubes and the decay exponent of the response field]
\label{4.29:def-path-length-decay}
We use the scales, the cube partitions and the current sequence $\{\Omega_i\}_{i\le k+1}$ of
Notation~\ref{4.29:not-scales-labels}: $\eta=L^{-k}$, $\ell_n=L^n\eta$,
$\Gamma_a=\Omega_a\setminus\Omega_{a+1}$, and $n(x)$ is the index $a$ with $x\in\Gamma_a$. Let
$\mathbf{H}_k(B')$ be the response of the level-$k$ structure on $\Omega_{k+1}$, as at the
end-points of the configurations in Lemma~\ref{4.29:lem-localisation-exact}. Its sources are the
bonds of the block structure on $\Gamma_{k+1}$ that carry $B'$, and the determining set of
$\mathbf{H}_k$ is of level $k$ on $\Omega_{k+1}$. The fluctuation fields $A$, $A_k$, the cutoff
$\chi^{(k)}$ of \cite[(3.16)]{B-CMP119}, the cutoff $\chi'_k$, and the region
$\Lambda^{(k)*}_{k+1}$ in which the fluctuation field $A$ is localised in the last exponential of
\cite[(3.28), p.~275]{B-CMP119} are those of \cite[\S3]{B-CMP119}. The regions
$\Lambda_{k+1}\subset\Omega_{k+1}$ and $S_{k+1}$, the parameters $R_k$, the functions $p_1(\cdot)$,
$\varepsilon_k$, $\delta_k$ and the constants $A_0$, $A_1$ are those of \cite{B-CMP119}; $B_3$ is
the constant so denoted in \cite[(1.45)]{B-CMP122a} and in the display of \cite[p.~267]{B-CMP119},
and $\beta_0$ is the constant in that display, recalled in the proof below. For a domain
$\Omega$ that is a union of cubes of a given side, $\Omega^{\sim-1}$ is $\Omega$ with one layer of
these cubes removed and $\Omega^{\sim}$ is $\Omega$ with one layer of these cubes added; for
$\Omega_{k+1}$ the side is $LMR_{k+1}$, and $\Lambda^{\sim}_{k+1}$ is the region so denoted in
\cite[\S3]{B-CMP119}. The set $S_{k+1}$ contains $\Omega_{k+1}\setminus\Omega^{\sim-1}_{k+1}$.

\emph{Hypotheses.} The following five statements are assumed.
\begin{enumerate}
\item[(I1)] \emph{Layered decay of the response field in local units.} $\mathbf{H}_k(B')$ and its
lattice derivatives obey the layered form of Ba{\l}aban's exponential decay bound for the response
field. At a point $y$ the field is weighted by $\ell_a$, its gradient by $\ell_a^2$, and so on, with
\[
a=\min\{n(y),k\},
\]
and the decay is exponential in the number of local $M$-cubes between $y$ and
$\operatorname{supp}B'$; that is, the second line of \eqref{4.29:eq-path-length-decay-a} below
holds for an exponent $\mathsf{d}_{\mathbf{H}}$, with the constant $B_3$.
\item[(I2)] \emph{The constant of the fluctuation bound on $\Omega_{k+1}$.} On $\Omega_{k+1}$, the
bonds of $\partial\Omega_{k+1}$ included, the bound $|A'_k|<O(1)\varepsilon_k$ of
\cite[(3.14)]{B-CMP119} holds with the constant written $O(1)$ equal to $a_\sharp$, with
$a_\sharp$ as in \eqref{4.29:eq-path-length-decay-1} below.
\item[(I3)] \emph{The shift of the centre.} After the change of variables
\cite[(3.21)]{B-CMP119} the centre of the fluctuation integration moves by the configuration
$\mathbf{H}^{(k)}$ of \cite[(3.21)]{B-CMP119}, which satisfies
\[
|\mathbf{H}^{(k)}|\le O(1)\,\varepsilon_k\,e^{-R_k},
\]
with the constant written $O(1)$ satisfying $O(1)e^{-R_k}\le1$.
\item[(I4)] \emph{The field $A_k$ as an external small field.} On
$\Omega_{k+1}\setminus\Lambda_{k+1}$ the field $A_k$ is one more external small field. It satisfies
$|A_k|\le p_1(g_k)$ on $\operatorname{supp}\chi'_k$ and enters the expansions of
\cite[\S3]{B-CMP119} exactly as $A$ does.
\item[(I5)] \emph{The bound on the sources on the enlarged domain.} With $\delta^\sharp_k$ and
$\mathfrak{D}^\sharp$ as in \eqref{4.29:eq-path-length-decay-b} and
\eqref{4.29:eq-path-length-decay-2}, one has
\[
\sup|B'(\tau A_k,A)|\le\delta^\sharp_k
\]
for $\tau\in[0,1]$ on $\mathfrak{D}^\sharp$, and $\mathfrak{D}^\sharp$ contains the supports of
the integrations over $A_k$ and $A$.
\end{enumerate}

\emph{Path lengths and the decay exponent.} For a lattice path $\omega$ and $a\le k$, let
$\lambda_a(\omega)$ be the length of $\omega$ inside $\Gamma_a$. For a site $y$ and a set of sites
$D$, the infimum below runs over the paths $\omega$ from $y$ to $D$. We put
\begin{equation}\label{4.29:eq-path-length-decay-a}
\begin{gathered}
|\omega|_{\mathbb{B}}:=\sum_{a\le k}\frac{\lambda_a(\omega)}{M\ell_a},\qquad
|\omega|_{k+1}:=\sum_{a\le k}\frac{\lambda_a(\omega)}{M\ell_{k+1}},\qquad
d_{\mathbb{B}}(y,D):=\inf_{\omega}|\omega|_{\mathbb{B}},\\[2pt]
\ell_a|\mathbf{H}_k(B')(y)|,\ \ \ell_a^2|\nabla\mathbf{H}_k(B')(y)|,\ \ \ldots\ \le\
B_3\,e^{-\mathsf{d}_{\mathbf{H}}\,d_{\mathbb{B}}(y,\operatorname{supp}B')}\,\sup|B'|,
\qquad a=\min\{n(y),k\}.
\end{gathered}
\end{equation}
Thus $|\omega|_{\mathbb{B}}$ is the length of $\omega$ counted in local $M$-cubes, and
$|\omega|_{k+1}$ counts the same stretches in $M$-cubes of scale $k+1$. The
\emph{decay exponent of the response field} $\mathsf{d}_{\mathbf{H}}$ is the exponent of (I1), that
is, the decay exponent per local $M$-cube of the layered bound in (I1), written as in the second
line of \eqref{4.29:eq-path-length-decay-a}. At a point $y\in\Omega_{k+1}$ one has $a=k$, because the
determining set of $\mathbf{H}_k$ is of level $k$ there.

In \cite{B-CMP109} this bound carries the decay rate $\delta_0$ on the scale of the lattice
spacing, with a reference to the exponential decay bound of an earlier paper in its list of
references; in this normalisation it supplies (I1) with any exponent
$\mathsf{d}_{\mathbf{H}}\le\delta_0M$. \cite[(1.45)]{B-CMP122a} states the bound in layered form.
At a point $y$ of the $i$-th component of the determining set, $L^i\eta|H|$ and
$(L^i\eta)^2|\nabla H|$ are bounded by $B_3e^{-\delta d(y,\cdot)}\,4\delta'_j$, with the rate
$\delta$, the scaled distance $d(y,\cdot)$ and the factor $4\delta'_j$ of that display; the scaled
distance is defined through $M_1$-cubes on the corresponding scales. In that normalisation the
bound supplies (I1) with any exponent $\mathsf{d}_{\mathbf{H}}\le\delta M/M_1$. In this chapter
(I1) is taken at the rate $\delta_0/8$, that is, with $\mathsf{d}_{\mathbf{H}}=\delta_0M/8$. Let
$\kappa'$ be the decay exponent of the kernel estimate
fixed in the glossary (Notation~\ref{0.2:not-glossary}); it is the exponent of the decay factor
$e^{-\kappa'd_{k+1}(Y)}$, for a domain $Y$ of the level-$(k+1)$ expansion, in the inductive bound
on the localised terms used in this chapter. The
lower bound
$\mathsf{d}_{\mathbf{H}}\ge\kappa'+4\log L+1$, imposed among the far-field thresholds of this
section, then reads
\[
\delta_0M\ge 8(\kappa'+4\log L+1).
\]
It is a one-sided condition on $M$, which is chosen after $L$ and $\kappa$.

The restriction $a\le k$ is part of the definitions. The length $|\omega|_{k+1}$ omits the stretch
of $\omega$ inside $\Omega_{k+1}$, which $d_{k+1}(Y)$ counts. There the local unit is $\ell_{k+1}$,
and the lower bound $\mathsf{d}_{\mathbf{H}}\ge\kappa'$ leaves a surplus $\mathsf{d}_{\mathbf{H}}-\kappa'$
per local $M$-cube on that stretch, with no factor $L^{-1}$. This is the form used for the
near-field class at the cubes beside the fluctuation support. The length of $\omega$ inside
$\Omega_{k+1}$ and the lengths $|\omega|^+_{\mathbb{B}}$, $|\omega|^+_{k+1}$ are defined there, in
\eqref{4.29:eq-near-field-scope-b}. Here $a\le k$ is kept, because the inequality
\eqref{4.29:eq-path-length-decay-3} below is what the far-field bounds of this section use, and it
fails on $\Gamma_{k+1}$.

\emph{Constants and the enlarged domain.} Let $C_{\mathrm{II}}$ be the absolute constant of
\cite[(1.20)]{B-CMP116}. The subscript distinguishes it from the constants $C_1$ of
\cite[(2.28)]{B-CMP119} and \cite[(3.43)]{B-CMP119}; we write $C_1^{(3.43)}$ for the constant of
\cite[(3.43)]{B-CMP119}.
We put
\begin{equation}\label{4.29:eq-path-length-decay-1}
a_\sharp:=1+c_{14},\qquad
c_{14}:=\bigl(2+d(L-1)\bigr)\frac{A_1}{A_0}+184\,d^2L^4(1+\beta_0)B_3,
\end{equation}
\begin{equation}\label{4.29:eq-path-length-decay-b}
\delta^\sharp_k:=C_{\mathrm{II}}\max\bigl\{4a_\sharp\varepsilon_k,\ C_1^{(3.43)}\delta_k\bigr\}
=C_{\mathrm{II}}\max\Bigl\{4a_\sharp\frac{A_0}{A_1},\ C_1^{(3.43)}\Bigr\}\,\delta_k,
\end{equation}
\begin{equation}\label{4.29:eq-path-length-decay-2}
\mathfrak{D}^\sharp:=\bigl\{|g_kA_k|\le4a_\sharp\varepsilon_k\bigr\}\times
\bigl\{|A|<C_1^{(3.43)}p_1(g_k)\bigr\}.
\end{equation}
The second factor in \eqref{4.29:eq-path-length-decay-2} is the polydisc of
\cite[(3.43)]{B-CMP119}. The second form of $\delta^\sharp_k$ holds because $\varepsilon_k$ and
$\delta_k$ carry the same power of $\log x_k$, in the ratio $\varepsilon_k/\delta_k=A_0/A_1$.

\emph{Conclusions.} Under (I1)--(I5) the following hold.
\begin{enumerate}
\item[(1)] $|\omega|_{k+1}\le L^{-1}|\omega|_{\mathbb{B}}$ for every path $\omega$.
\item[(2)] $\operatorname{supp}B'\subset\Omega_{k+1}$. The source $A$ lives on $\Lambda_{k+1}$ (in
fact on $\Lambda^{(k)*}_{k+1}$), and the source $A_k$ on $\Omega_{k+1}\setminus\Lambda_{k+1}$.
$A_k$ enters at full strength in the action of \cite[(3.27)]{B-CMP119}, as $tA_k$ in that of
\cite[(3.29)]{B-CMP119}, and in the integrands of \cite[(3.26), (3.28)]{B-CMP119}. It is absent from
the last logarithm of \cite[(3.28)]{B-CMP119}.
\item[(3)] $|g_kA_k|\le2a_\sharp\varepsilon_k$ on $S_{k+1}\cap\Omega_{k+1}$, so that the support
of the $A_k$-integration there lies in the first factor of $\mathfrak{D}^\sharp$. The field-size
condition on the sources of the response is read with $\delta^\sharp_k$, whichever of the two
parts of $B'$ carries the source. The parameter $t$ of \cite[(3.29)]{B-CMP119} satisfies $t\le1$,
so that the bound of (I5) applies to the field $tA_k$ with $\tau=t$.
\item[(4)] $d_{\mathbb{B}}(y,\operatorname{supp}B')\ge d_{\mathbb{B}}(y,\Omega_{k+1})$ for every
site $y$. Consequently a bound written with distances measured from $\Omega_{k+1}$, or along paths
starting in $\Omega_{k+1}$, is a weakening of the decay in \eqref{4.29:eq-path-length-decay-a},
and holds for both kinds of source as stated; this applies to the far-field bounds of this
section, to the two further estimates stated with them and to the two field-size conditions at
the cubes. Only the near-field class at $j=k$, whose cubes may
touch $\operatorname{supp}A_k$, distinguishes them.
\end{enumerate}
\end{definition}

\begin{proof}
\emph{Conclusion} (1). Let $a\le k$. Then
\[
M\ell_{k+1}=L^{k+1-a}M\ell_a\ge LM\ell_a,
\]
so $\lambda_a(\omega)/(M\ell_{k+1})\le L^{-1}\lambda_a(\omega)/(M\ell_a)$ term by term. Summing
over $a\le k$ gives
\begin{equation}\label{4.29:eq-path-length-decay-3}
|\omega|_{k+1}\le L^{-1}|\omega|_{\mathbb{B}}.
\end{equation}
If the sum also ran over $a=k+1$, the term $\lambda_{k+1}/(M\ell_{k+1})$ would appear with
coefficient one on both sides, and \eqref{4.29:eq-path-length-decay-3} would fail. This is why the
restriction $a\le k$ is kept.

\emph{Conclusion} (2). In \cite[\S3]{B-CMP119}, $A_k$ denotes the fluctuation field on
$\Omega_{k+1}\cap\Lambda^{c}_{k+1}$ and $A$ denotes this field on $\Lambda_{k+1}$. By (I4), $A_k$ is
one more external small field and enters exactly as $A$ does. The sources $B'$ are therefore formed
from $A$ on $\Lambda_{k+1}$ and from $A_k$ on $\Omega_{k+1}\setminus\Lambda_{k+1}$. Since
$\Lambda_{k+1}\subset\Omega_{k+1}$, it follows that $\operatorname{supp}B'\subset\Omega_{k+1}$.

For the last exponential of \cite[(3.28), p.~275]{B-CMP119}, that page states that the
fluctuation field $A$ is localised in $\Lambda^{(k)*}_{k+1}$. This is the refinement of the support
of $A$, and the reason $A_k$ does not occur in the last logarithm of \cite[(3.28)]{B-CMP119}.
The action of \cite[(3.27)]{B-CMP119} contains the field $A_k$ at full strength, with coefficient
one; the action of \cite[(3.29)]{B-CMP119} contains it as $tA_k$.
\cite{B-CMP119} states that the action in \cite[(3.29)]{B-CMP119} differs only by the additional
dependence on the field $tA_k$, and that this dependence was already taken into account. The
integrands of \cite[(3.26), (3.28)]{B-CMP119} contain $A_k$ as termwise derivatives of the same
localised expansion \cite{B-CMP119}. At the place where these terms are estimated, \cite{B-CMP119}
states that the fluctuation field is localised in $\Omega_{k+1}$, and that the exponential decay of
propagators and minimising functions yields additional small factors. This gives the four
occurrences listed in (2).

\emph{Conclusion} (3). We bound the sources in three regions.

\emph{First region.} Where the cutoff $\chi^{(k)}$ of \cite[(3.16)]{B-CMP119} acts, namely on
$\Lambda^{\sim}_{k+1}$ and on the cubes of the complement of a second region of
\cite[\S3]{B-CMP119}, denoted there by the same letter as the parameter $R_{k+1}$, the sources
satisfy
$|B'|\le C_{\mathrm{II}}\delta_k$, by \cite[(1.20)]{B-CMP116}. The decomposition of unity
\cite[(3.16)]{B-CMP119} is introduced only on $\Omega^{\sim-1}_{k+1}$, the domain obtained by
taking off one layer of $LMR_{k+1}$-cubes from $\Omega_{k+1}$.

\emph{Second region.} On the polydisc $\{|A|<C_1^{(3.43)}p_1(g_k)\}$ of \cite[(3.43)]{B-CMP119},
the same constant gives
\[
|B'|\le C_{\mathrm{II}}C_1^{(3.43)}\delta_k .
\]

\emph{Third region.} The set $S_{k+1}$ of \cite{B-CMP119} contains
\[
\Omega_{k+1}\setminus\Omega^{\sim-1}_{k+1}=\Omega_{k+1}\cap(\Omega^{c}_{k+1})^{\sim}.
\]
On this set the bounds of \cite{B-CMP119} that apply are \cite[(3.13)--(3.14), p.~267]{B-CMP119}. On the
enlarged cubes $\square'^{\sim}$ of that page one has
\[
2\delta_k+(4dL)^2(1+\beta_0)B_3\varepsilon_k=O(1)\varepsilon_k .
\]
This implies $|V'_k-1|<O(1)\varepsilon_k$ on $\Omega_{k+1}$, where $V'_k$ is the group-valued
variable of \cite[p.~267]{B-CMP119} whose logarithm defines the fluctuation field. With
$A'_k=(1/i)\log V'_k$ it follows that $|A'_k|<O(1)\varepsilon_k$, and $A'_k=g_kA_k$
\cite{B-CMP119}.

The cutoff $\chi'_k$ restricts the approximate fluctuation fields \cite{B-CMP119}. By (I2), the
constant $O(1)$ of \cite[(3.14)]{B-CMP119} on $\Omega_{k+1}$, the bonds of $\partial\Omega_{k+1}$
included, is $a_\sharp$ with $a_\sharp$ from \eqref{4.29:eq-path-length-decay-1}. After the change
of variables \cite[(3.21)]{B-CMP119} the centre moves by $\mathbf{H}^{(k)}$, and by (I3)
$|\mathbf{H}^{(k)}|\le O(1)\varepsilon_ke^{-R_k}$. From the bound of (I2) on the fields restricted
by $\chi'_k$ and from $A'_k=g_kA_k$ one has $|g_kA_k|<a_\sharp\varepsilon_k$ about the original
centre; together with this displacement of the centre,
\[
|g_kA_k|\le a_\sharp\varepsilon_k+O(1)\,\varepsilon_ke^{-R_k}\le2a_\sharp\varepsilon_k
\quad\text{on } S_{k+1}\cap\Omega_{k+1},
\]
where the second inequality holds by (I3) and $a_\sharp\ge1$.
The interior quantity $2\delta_k+(4dL)^2(1+\beta_0)B_3\varepsilon_k$ displayed above is at most
$c_{14}\varepsilon_k$, since $2\delta_k=2(A_1/A_0)\varepsilon_k\le(2+d(L-1))(A_1/A_0)\varepsilon_k$
and $(4dL)^2\le184\,d^2L^4$; the constant of (I2) holds on $\Omega_{k+1}$ with the bonds of
$\partial\Omega_{k+1}$ included.

\emph{The enlarged domain.} Since $2a_\sharp\varepsilon_k\le4a_\sharp\varepsilon_k$, the bound just
obtained places the support of the $A_k$-integration inside the first factor of
$\mathfrak{D}^\sharp$. That the supports of the integrations over $A_k$ and over $A$ lie in
$\mathfrak{D}^\sharp$ is part of (I5).

\emph{The constant $\delta^\sharp_k$.} The two entries of the maximum in
\eqref{4.29:eq-path-length-decay-b} are the source sizes on the two factors of
$\mathfrak{D}^\sharp$: $C_{\mathrm{II}}\cdot4a_\sharp\varepsilon_k$ on the first, and
$C_{\mathrm{II}}C_1^{(3.43)}\delta_k$ on the second, from the second region; the bound itself is
(I5). By (I5),
$\sup|B'(\tau A_k,A)|\le\delta^\sharp_k$ for $\tau\in[0,1]$ on $\mathfrak{D}^\sharp$. The field
$tA_k$ in the action of \cite[(3.29)]{B-CMP119} has $t\le1$, and the bound is applied to it with
$\tau=t$. Because
$\delta^\sharp_k$ dominates the constants of all three regions, the field-size condition on the
sources may be read with $\delta^\sharp_k$ whether the source is carried by $A$ or by $A_k$.

The second form in \eqref{4.29:eq-path-length-decay-b} follows from the first. Substitute
$\varepsilon_k=(A_0/A_1)\delta_k$, which is the common power of $\log x_k$ in $\varepsilon_k$ and
$\delta_k$ written out, and take $\delta_k$ out of the maximum.

\emph{Conclusion} (4). By (2), $\operatorname{supp}B'\subset\Omega_{k+1}$. Every path from $y$ to
$\operatorname{supp}B'$ is therefore a path from $y$ to $\Omega_{k+1}$. The infimum defining
$d_{\mathbb{B}}(y,\Omega_{k+1})$ runs over a set of paths that contains those defining
$d_{\mathbb{B}}(y,\operatorname{supp}B')$, so
\[
d_{\mathbb{B}}(y,\operatorname{supp}B')\ge d_{\mathbb{B}}(y,\Omega_{k+1}).
\]
The decay factor in \eqref{4.29:eq-path-length-decay-a} is therefore bounded by the same factor
with $\operatorname{supp}B'$ replaced by $\Omega_{k+1}$.

A bound written with distances measured from $\Omega_{k+1}$, or along paths starting in
$\Omega_{k+1}$, contains this weakened factor; the far-field bounds of this section, the two
further estimates stated with them and the two field-size conditions at the cubes are written
in this form, their distances being measured from $\Omega_{k+1}$ and the paths of the first
far-field bound starting in $\Omega_{k+1}$. Such a bound does not depend on which part of $B'$
carries the source, and it holds for both kinds of source as stated. The near-field class at $j=k$ is
different: its cubes may touch $\operatorname{supp}A_k$, where $d_{\mathbb{B}}(\cdot,\operatorname{supp}B')$
may vanish. That class is the only place where the two kinds of source are distinguished.
\end{proof}

\begin{lemma}[Path-length bounds in the far field]\label{4.29:lem-far-field-paths}
Fix a level $j\le k$. We use the following objects, each with its definition:
\begin{itemize}
\item from Notation~\ref{4.29:not-scales-labels}: the scales $\ell_a=L^a\eta$; the current sequence
$\{\Omega_i\}_{i\le k+1}$, its strata $\Gamma_a=\Omega_a\setminus\Omega_{a+1}$ and its index $n(\cdot)$;
the native sequence $\{\Omega_i^\natural\}$ and its index $l(\cdot)$; the action domains
$\{\Lambda_i\}_{i\le k}$; the radii $R_i$; the sup-distance $\operatorname{dist}_\infty$; and the
constant $(1+4\beta)\kappa$, which we denote $\kappa_\beta$;
\item from Proposition~\ref{4.29:prop-native-cover}: the enlarged action domain $\Lambda_j^{\sim}$,
the near-field region $P$ and the scale function $\mathfrak{s}_j$;
\item from Definition~\ref{4.29:def-path-length-decay}: the lengths $\lambda_a(\omega)$ of a lattice
path $\omega$ inside the strata $\Gamma_a$ ($a\le k$), its length $|\omega|_{\mathbb{B}}$ in local
$M$-cubes, its length $|\omega|_{k+1}$ in $M$-cubes of scale $k+1$, and the decay exponent
$\mathsf{d}_{\mathbf{H}}$, as displayed in \eqref{4.29:eq-path-length-decay-a}.
\end{itemize}
Assume the hypotheses of Proposition~\ref{4.29:prop-native-cover}, and assume:
\begin{itemize}
\item clauses (T0), (T1), (T2), (G1) and (A) of the lemma on the label sequences at the
$\mathbf{T}$-stage.
In the form used below: (T0) gives $\Lambda_i\subset\Omega_i^\natural\subset\Omega_i$
and $\Omega_{i+1}^\natural\subset\Lambda_i$; (T1) gives
$\operatorname{dist}_\infty(\Omega^\natural_{i+1},\Lambda_i^{c})\ge 8MR_{i+1}\ell_{i+1}$ for $i\le k$,
and $\Omega_{k+1}=\Omega_{k+1}^\natural$; (T2)
gives $\Lambda_j^{\sim}\subset\Omega_j^\natural$; (G1) is
$\operatorname{dist}_\infty(\Omega_{a+1},\Omega_a^{c})\ge M\ell_a$; and (A) says that a point $y$ with
$n(y)>l(y)$ (a \emph{lifted} point) lies in an \emph{arrested component} $W$, to which integers
$k_0<k_W\le k$, $N_0:=k_W-k_0$ with $L^{N_0-1}=R_{k_0+1}$, and $j_W\le k_W$ are attached; each point
is lifted at most once, and a lifted point has $n(y)\le j_W$ and either $n(y)=l(y)+1$ or
$k_0\le l(y)<n(y)$; every point of $W$ has $n\le k_W$ at all later times;
\item the following three properties of the radii $R_i$ fixed in
Notation~\ref{4.29:not-scales-labels}, in the instances used below: the \emph{logarithmic lower
bound} $R_{k+1}\ge\log x_{k+1}$; the \emph{step comparison}: for every arrested component $W$ and
with the exponent $\beta_0$ of that property, $R_{k_W}\le(L+1)(k+1-k_W)^{\beta_0}R_{k+1}$, where
$k+1-k_W\ge1$; and the \emph{arrest comparison} $R_{k_0+1}\le LR_{k_W}$,
the index difference $k_W-k_0-1=\log_LR_{k_0+1}$ being at most $L^2R_{k_0+1}$. In the step
comparison the earlier index carries the larger radius, and the ratio is bounded by a power of the
number of steps, as the flow inequality \eqref{4.29:eq-scales-labels-b} gives;
\item the bound $x_{k+1}\ge\gamma^{-2}$ on the inverse coupling $x_{k+1}=g_{k+1}^{-2}$ at level $k+1$;
\item clauses (b), (e)$^\natural$ and (f) of Proposition~\ref{4.29:prop-native-cover}.
\end{itemize}
Let $x\in\Lambda_j^{\sim}\setminus P$ with $n:=n(x)\le k$, put $s:=\mathfrak{s}_j(x)$ and
$u:=\max\{k-1-n,0\}$, let $\beta_0\le 1/10$, and put $\kappa_L:=\kappa_\beta(1-1/L)$, which is
positive since $\kappa>0$ and $L\ge2$. Then
\begin{equation}\label{4.29:eq-far-field-paths-a}
\begin{gathered}
\text{(g1)}\quad |\omega|_{\mathbb{B}}\ge 7LR_{k+1}+u
\quad\text{for every lattice path $\omega$ from $\Omega_{k+1}$ to $x$;}\\
\text{(g2)}\quad L^{4(n-s)}(1+k-s)^{1/2}\le e\,L^{21}R_{k+1}^{5}e^{u};\\
\text{floor:}\quad \mathsf{d}_{\mathbf{H}}\ge\kappa_\beta+4\log L+1;\\
\text{(g3)}\quad L^{4(k-s)}(1+k-s)^{1/2}e^{-\mathsf{d}_{\mathbf{H}}|\omega|_{\mathbb{B}}}
\le e^{-\kappa_\beta|\omega|_{k+1}}\,e^{-6R_{k+1}},\\
e^{-6R_{k+1}}\le x_{k+1}^{-6}\le\gamma^{12};\\
\text{(g3)}^{+}\quad L^{4(k-s)}(1+k-s)^{1/2}e^{-\mathsf{d}_{\mathbf{H}}|\omega|_{\mathbb{B}}}
\le e^{-\kappa_\beta|\omega|_{k+1}}\,e^{-6R_{k+1}}\,e^{-\kappa_Lu}.
\end{gathered}
\end{equation}
Here (g1) and (g2) hold as stated; (g3) and (g3)$^+$ hold for every path $\omega$ as in (g1)
whenever the floor holds. The floor is one of the one-sided lower bounds collected in
\eqref{4.29:eq-constants-order-c}.
\end{lemma}

\begin{proof}
\emph{Domain.} The argument below uses of the point $x$ only the following: $x\notin P$;
$x\notin\Omega_a$ for $a>n(x)$; clause (e)$^\natural$ in the form $n-s\le n-l+1$, where $l:=l(x)$;
and (A) together with the step comparison and the arrest comparison of the radii, all of which are
pointwise statements in $x$. For $x\in\Lambda_j\setminus P$ these
hold directly: $x\notin P$ by assumption; $x\notin\Omega_a$ for $a>n(x)$ because $x\in\Gamma_{n(x)}$
and the sequence $\{\Omega_a\}$ decreases ($\Omega_{a+1}\subset\Omega_a$ by (G1)); and the lower bound
$l(x)-1\le\mathfrak{s}_j(x)$ of (e)$^\natural$ gives $n-s\le n-l+1$.

Let now $x$ lie in the ring $\Lambda_j^{\sim}\setminus\Lambda_j$. By (T0) and clause (b) we have
\begin{equation*}
\Lambda_j\supset\Lambda_k\supset P\supset\Omega_{k+1};
\end{equation*}
since $x\notin\Lambda_j$, it follows that
$x\notin P$. Moreover $j\le n(x)\le k$: on the one hand, by (T2) and (T0),
\begin{equation*}
x\in\Lambda_j^{\sim}\subset\Omega_j^\natural\subset\Omega_j ;
\end{equation*}
on the other hand, by (T1) and (T0),
\begin{equation*}
\Omega_{k+1}=\Omega_{k+1}^\natural\subset\Lambda_k ,
\end{equation*}
and $x\notin\Lambda_k$, so $x\notin\Omega_{k+1}$. On the ring, clause (e)$^\natural$
gives $s=\mathfrak{s}_j(x)=j=l(x)$, so that $n-s=n-l$, which is the required form of (e)$^\natural$.
Hence the argument below applies to every $x\in\Lambda_j^{\sim}\setminus P$.

\emph{Proof of (g1).} Let $\omega$ be a lattice path from a site of $\Omega_{k+1}$ to $x$. Since
$x\notin P$, after its last exit from $\Omega_{k+1}$ the path $\omega$ leaves $P$. Between its last
exit from $\Omega_{k+1}$ and its first subsequent exit from $P$ it runs in
\begin{equation*}
P\setminus\Omega_{k+1}\subset\Lambda_k\setminus\Omega_{k+1}\subset\Omega_k\setminus\Omega_{k+1}=\Gamma_k,
\end{equation*}
using $P\subset\Lambda_k$ (see above) and $\Lambda_k\subset\Omega_k$ (T0). On this stretch it covers a
sup-distance at least $\operatorname{dist}_\infty(\Omega_{k+1},P^{c})$, and this is at least
$7LR_{k+1}M\ell_k$ by the arithmetic of the proof of clause (f) of
Proposition~\ref{4.29:prop-native-cover}: there a site $y\in\Lambda_k$ of a cube of the cover
$\mathcal{C}_j$ of that proposition not meeting $P$
satisfies $\operatorname{dist}_\infty(y,\Lambda_k^{c})<13M\ell_k$, and since (T1) at $i=k$ together
with $\Omega_{k+1}=\Omega_{k+1}^\natural$ gives
$\operatorname{dist}_\infty(\Omega_{k+1},\Lambda_k^{c})\ge8MR_{k+1}\ell_{k+1}=8LR_{k+1}M\ell_k$,
such a site has
\begin{equation*}
\operatorname{dist}_\infty(y,\Omega_{k+1})\ge(8LR_{k+1}-13)M\ell_k\ge7LR_{k+1}M\ell_k ,
\end{equation*}
because $LR_{k+1}\ge13$; a site $y\notin\Lambda_k$ has
$\operatorname{dist}_\infty(y,\Omega_{k+1})\ge8LR_{k+1}M\ell_k$ by (T1) directly. Since each step
of a lattice path changes the sup-distance by at most its own
length,
\begin{equation*}
\lambda_k(\omega)\ge 7LR_{k+1}M\ell_k .
\end{equation*}
Next let $n<a<k$. The path starts inside $\Omega_{k+1}\subset\Omega_{a+1}$ and ends at $x$, which lies
outside $\Omega_a$, because $\Omega_a\subset\Omega_{n+1}$ and $x\notin\Omega_{n+1}$. Between its last
exit from $\Omega_{a+1}$ and its next entry into $\Omega_a^{c}$ the path runs in
$\Omega_a\setminus\Omega_{a+1}=\Gamma_a$, and by (G1) it covers there a sup-distance at least
$\operatorname{dist}_\infty(\Omega_{a+1},\Omega_a^{c})\ge M\ell_a$; hence $\lambda_a(\omega)\ge M\ell_a$.
The portions so found lie in different strata $\Gamma_k$ and $\Gamma_a$, $n<a<k$, which are pairwise
disjoint, and there are exactly $\max\{k-1-n,0\}=u$ indices $a$ with $n<a<k$. By the definition of
$|\omega|_{\mathbb{B}}$ as $\sum_a\lambda_a(\omega)/(M\ell_a)$,
\begin{equation*}
|\omega|_{\mathbb{B}}\ge\frac{\lambda_k(\omega)}{M\ell_k}+\sum_{n<a<k}\frac{\lambda_a(\omega)}{M\ell_a}
\ge 7LR_{k+1}+u .
\end{equation*}
An estimate of the same kind occurs in \cite[(1.47)]{B-CMP122a}: there, in Ba{\l}aban's indices $i$,
$j$ (unrelated to the level $j$ fixed above) and with his decay rate $\delta$, the product of
$1+(j-i)^{1/2}$ and the exponential factor $\exp(-\delta(M/M_1)(j-i))$ is estimated by
$\exp(-(j-i))$.

\emph{Proof of (g2).} By (e)$^\natural$, $n-s\le n-l+1$ (on the ring $n-s=n-l$). Further
$n\ge l\ge s$: $s\le l$ is the upper bound $\mathfrak{s}_j(x)\le\min\{l(x),k\}$ of (e)$^\natural$, and
$l\le n$ because $x\in\Omega_l^\natural\subset\Omega_l$ by (T0).

Put $t_W:=k+1-k_W$ if $x$ lies in an arrested component $W$, and $t_W:=1$ otherwise; in both cases
$t_W\ge1$.

If $x$ is not lifted ($n=l$) or is lifted by one ($n=l+1$), then $n-s\le n-l+1\le2$, and
$L^{5(n-s)}\le L^{10}\le L^{21}t_W^{5\beta_0}R_{k+1}^{5}$ because $t_W\ge1$ and $R_{k+1}\ge1$.
Otherwise, by
(A), $x$ lies in an arrested component $W$ with $k_0\le l$ and $n\le j_W\le k_W$; then
\begin{equation*}
n-s\le n-l+1\le k_W-k_0+1=N_0+1,
\end{equation*}
and
\begin{equation*}
L^{5(n-s)}\le L^{5(N_0+1)}=L^{10}L^{5(N_0-1)}=L^{10}R_{k_0+1}^{5}.
\end{equation*}
By the arrest comparison of the radii, $R_{k_0+1}\le LR_{k_W}$, and by the step comparison of the
radii, applied to this $W$, for which $t_W=k+1-k_W$, we have
$R_{k_W}\le(L+1)t_W^{\beta_0}R_{k+1}$. Therefore
\begin{equation*}
L^{5(n-s)}\le L^{10}L^{5}(L+1)^{5}t_W^{5\beta_0}R_{k+1}^{5}
=L^{20}\Bigl(1+\frac1L\Bigr)^{5}t_W^{5\beta_0}R_{k+1}^{5}\le L^{21}t_W^{5\beta_0}R_{k+1}^{5},
\end{equation*}
since $(1+1/L)^5\le L$ holds for $L\ge4$, in particular under the lower bound on $L$ among the
hypotheses of Proposition~\ref{4.29:prop-native-cover}. So in all cases
\begin{equation}\label{4.29:eq-far-field-paths-1}
L^{5(n-s)}\le L^{21}t_W^{5\beta_0}R_{k+1}^{5}.
\end{equation}
Next, $k-n\le u+1$ by the definition of $u$, so
\begin{equation*}
1+k-s=1+(k-n)+(n-s)\le(u+2)+(n-s)\le(u+2)(1+n-s),
\end{equation*}
the last step because $u+2\ge1$ and $n-s\ge0$. For every integer $v\ge0$ we have
$(1+v)^{1/2}\le2^{v}\le L^{v}$, so $(1+n-s)^{1/2}\le L^{n-s}$. With \eqref{4.29:eq-far-field-paths-1},
\begin{equation*}
L^{4(n-s)}(1+k-s)^{1/2}\le L^{5(n-s)}(u+2)^{1/2}\le L^{21}R_{k+1}^{5}\,t_W^{5\beta_0}(u+2)^{1/2}.
\end{equation*}
If $x$ lies in an arrested component $W$, then $n\le k_W$, hence
\begin{equation*}
t_W=k+1-k_W\le k+1-n\le u+2 ;
\end{equation*}
if not, $t_W=1\le u+2$. Since $5\beta_0\le1/2$,
\begin{equation*}
t_W^{5\beta_0}(u+2)^{1/2}\le(u+2)^{5\beta_0+1/2}\le u+2\le e^{u+1},
\end{equation*}
and (g2) follows.

\emph{Proof of (g3).} Assume the floor. By definition, $|\omega|_{k+1}=\sum_{a\le k}\lambda_a/(M\ell_{k+1})$,
and $\ell_a/\ell_{k+1}=L^{a-k-1}\le L^{-1}$ for $a\le k$, so $|\omega|_{k+1}\le L^{-1}|\omega|_{\mathbb{B}}$.
Since $\kappa_\beta=(1+4\beta)\kappa>0$, the floor gives
$\mathsf{d}_{\mathbf{H}}-\kappa_\beta/L\ge\mathsf{d}_{\mathbf{H}}-\kappa_\beta\ge4\log L+1$, and with (g1)
\begin{equation*}
\mathsf{d}_{\mathbf{H}}|\omega|_{\mathbb{B}}-\kappa_\beta|\omega|_{k+1}
\ge\Bigl(\mathsf{d}_{\mathbf{H}}-\frac{\kappa_\beta}{L}\Bigr)|\omega|_{\mathbb{B}}
\ge(4\log L+1)(7LR_{k+1}+u).
\end{equation*}
Moreover $L^{4(k-s)}=L^{4(k-n)}L^{4(n-s)}\le L^{4(u+1)}L^{4(n-s)}$. Combining these with (g2), the left
side of (g3) is at most
\begin{equation*}
e^{-\kappa_\beta|\omega|_{k+1}}\cdot L^{4}\,e\,L^{21}R_{k+1}^{5}\,e^{-(4\log L+1)7LR_{k+1}}
\cdot\bigl[L^{4u}e^{u}e^{-(4\log L+1)u}\bigr].
\end{equation*}
The bracket equals $1$, and $e^{-(4\log L+1)7LR_{k+1}}\le e^{-7LR_{k+1}}$; so the left side of (g3)
is at most $e^{-\kappa_\beta|\omega|_{k+1}}\cdot e\,L^{25}R_{k+1}^{5}e^{-7LR_{k+1}}$. It remains to
see that
\begin{equation*}
1+25\log L+5\log R_{k+1}\le(7L-6)R_{k+1},
\end{equation*}
which holds for $L,R_{k+1}\ge7$, as the hypotheses of Proposition~\ref{4.29:prop-native-cover}
provide. Indeed, $\log y\le y/3$ for $y\ge7$, so the difference of the two sides is at least
\begin{equation*}
L\Bigl(7R_{k+1}-\frac{25}{3}\Bigr)-\frac{23R_{k+1}}{3}-1
\ge49R_{k+1}-\frac{175}{3}-\frac{23R_{k+1}}{3}-1=\frac{124R_{k+1}-178}{3}>0,
\end{equation*}
where we used $7R_{k+1}-25/3>0$ and $L\ge7$. Hence $e\,L^{25}R_{k+1}^{5}e^{-7LR_{k+1}}\le e^{-6R_{k+1}}$,
which is (g3). Finally, $e^{-6R_{k+1}}\le x_{k+1}^{-6}$ because $R_{k+1}\ge\log x_{k+1}$ by the
logarithmic lower bound on the radii; the inequality $x_{k+1}^{-6}\le\gamma^{12}$ is the assumed bound
$x_{k+1}\ge\gamma^{-2}$ on the inverse coupling $x_{k+1}=g_{k+1}^{-2}$ at level $k+1$.

\emph{Proof of (g3)$^+$.} Under the same floor,
\begin{equation*}
\mathsf{d}_{\mathbf{H}}-\frac{\kappa_\beta}{L}\ge(4\log L+1)+\kappa_\beta\Bigl(1-\frac1L\Bigr)=(4\log L+1)+\kappa_L .
\end{equation*}
The proof of (g3) used only the part $(4\log L+1)|\omega|_{\mathbb{B}}$ of
$(\mathsf{d}_{\mathbf{H}}-\kappa_\beta/L)|\omega|_{\mathbb{B}}$. The remaining part satisfies, by (g1),
\begin{equation*}
\kappa_L|\omega|_{\mathbb{B}}\ge\kappa_L(7LR_{k+1}+u)\ge\kappa_Lu,
\end{equation*}
so the estimate of (g3) carries the additional factor $e^{-\kappa_L|\omega|_{\mathbb{B}}}\le e^{-\kappa_Lu}$,
which is (g3)$^+$. No condition beyond the floor is used.
\end{proof}

\begin{remark}[Far-field decay pays the cube count uniformly in the number of steps]
\label{4.29:rem-far-field-count}
Fix levels $j\le k$. We use four sets of notation:
\begin{itemize}
\item the cover $\mathcal{C}_j$, its tiles and the native scale $\mathfrak{s}_j$ of
Proposition~\ref{4.29:prop-native-cover};
\item the path lengths $|\omega|_{\mathbb{B}}$, $|\omega|_{k+1}$, the distance $d_{\mathbb{B}}$ and the
decay exponent $\mathsf{d}$ of the response field $H_k=H_k(B')$, whose source satisfies
$\operatorname{supp}B'\subset\Omega_{k+1}$, all from
Definition~\ref{4.29:def-path-length-decay};
\item the cubes of the partition $\pi_k$ of level $k$, called $\pi_k$-cubes;
\item for $x\in\Lambda_j^{\sim}\setminus P$, the quantities $n=n(x)$, $s=\mathfrak{s}_j(x)$,
$u=(k-1-n)_+$ and $R=R_{k+1}$ of Lemma~\ref{4.29:lem-far-field-paths}, and the generation
$l(x)$ of Proposition~\ref{4.29:prop-native-cover}.
\end{itemize}
The number of cubes of $\mathcal{C}_j$ in a $\pi_k$-cube, the ring cubes of scale $j$
included, is at most $2^dL^{4(k-s)}$. Every piece born at a cube $\square$ is a
derivative in the direction $\bar\zeta_\square H_k(B')$. It therefore carries the factor
$e^{-\mathsf{d}\,d_{\mathbb{B}}(\square,\operatorname{supp}B')}$, which is at most
$e^{-\mathsf{d}\,d_{\mathbb{B}}(\square,\Omega_{k+1})}$. The factor $(1+k-s)^{1/2}$, which
satisfies $(1+k-s)^{1/2}\ge g_k/g_s$, is the drift between the size of the fluctuation
and the radii of level $s$.

Assume $L\ge 13$, $R=R_{k+1}\ge7$ (the radii satisfy $R_i\ge7$ in
Proposition~\ref{4.29:prop-native-cover}), the floor $\mathsf{d}\ge\kappa'+4\log L+1$, and
the following.
\begin{itemize}
\item \emph{The far-field path bounds.} These are the bounds of
Lemma~\ref{4.29:lem-far-field-paths} collected in \eqref{4.29:eq-far-field-paths-a}: the
path bound (g1), the bound (g2) together with the radius comparisons used in its
proof, the bound (g3) with $e^{-6R}\le x_{k+1}^{-6}\le\gamma^{12}$, and the refined bound
(g3)$^+$. The last states that every path $\omega$ from $\Omega_{k+1}$ to $x$ satisfies
\begin{equation*}
L^{4(k-s)}(1+k-s)^{1/2}e^{-\mathsf{d}|\omega|_{\mathbb{B}}}
\le e^{-\kappa'|\omega|_{k+1}}\,e^{-6R}\,e^{-\kappa'(1-1/L)u}.
\end{equation*}
\item \emph{Level structure.} $\Omega_{k+1}\subset P$. The regions are nested,
$\Omega_{a+1}\subset\Omega_a$ for every level $a$. The sets $\Omega_k$ and $\Omega_{k+1}$
are unions of $\pi_k$-cubes. A point $y$ of current level $n(y)$ lies in
$\Omega_{n(y)}$ and does not lie in $\Omega_{n(y)+1}$.
\item \emph{Width of the strata.} For every level $a$,
$\operatorname{dist}_\infty(\Omega_{a+1},\Omega_a^{c})\ge M\ell_a$.
\item \emph{Generations in arrested components.} Each arrested component $W$ carries two
levels $k_0(W)<k_W$ and a generation count
\begin{equation*}
N_0(W)=k_W-k_0(W)=1+\log_LR_{k_0(W)+1}.
\end{equation*}
Put $t_W:=k+1-k_W$; it satisfies $t_W\ge1$. Let $x$ be a point of current level $n$.
\begin{itemize}
\item If $x$ lies in an arrested component $W$, then $n\le k_W$, and $l(x)\in[k_0(W),n]$ or
$l(x)\in\{n-1,n\}$.
\item If $x$ lies in no arrested component, then $l(x)=n$.
\end{itemize}
\item \emph{Radius comparisons.} For every arrested component $W$,
\begin{equation*}
R_{k_0(W)+1}\le L\,R_{k_W},\qquad R_{k_W}\le(L+1)\,t_W^{\beta_0}\,R_{k+1},
\end{equation*}
with the exponent $\beta_0\le 1/10$ of Lemma~\ref{4.29:lem-far-field-paths}.
\item \emph{The count lemma for the expansion of \cite{B-CMP119}.} This lemma has two
parts.
\begin{itemize}
\item Pieces at a point of current level $n$ are counted in Ba{\l}aban's estimate with
the factor $L^{4(k-n)}$.
\item At good cubes, the sum over the generations $j\le s$ that share a tile is bounded,
uniformly in $s$ and $k$, by a sum of terms of four kinds: terms $e^{-\delta_0ML^{s-j}}$
per point; terms $L^{-(5-\beta)(s-j)}$ per point $z$, over at most $L^{4(s-j)}$ points $z$
\cite{B-CMP119}; terms bounded by \cite[(3.48)]{B-CMP119}; and the pieces of the
operation $\mathbf{R}$, bounded by the lemma's clause for them. Here $\delta_0$ is the
decay rate on the lattice scale of \cite{B-CMP109} from which $\mathsf{d}$ is obtained,
and $\beta$ is the schedule constant of the layered analyticity spaces of
Definition~\ref{4.4:def-post-step-space}.
\end{itemize}
\end{itemize}

Then the following hold.
\begin{enumerate}
\item[(i)] Consider any connection from $\operatorname{supp}B'\subset\Omega_{k+1}$ to
$\square$. Along it, the part of the decay of $H_k$ beyond the factor
$e^{-\kappa'|\omega|_{k+1}}$, which is retained for the kernel bound of rate $\kappa'$ in
the metric $d_{k+1}$ (the form used in this chapter, which sharpens
\cite[(2.42)]{B-CMP119}), pays three things: Ba{\l}aban's count $L^{4(k-n)}$, the extra
native count $L^{4(n-s)}$, and the drift $(1+k-s)^{1/2}$.
What is left is the factor $e^{-6R_{k+1}}\le\gamma^{12}$. This factor absorbs every
constant depending only on $(L,M,\mathfrak{p})$, among them the loss of one native scale
in the irrelevance factors and the constant $C_\zeta L^4$ of the partition of unity
$\{\zeta_\square\}$ of Proposition~\ref{4.29:prop-native-cover}, as soon as
$\gamma\le\gamma_\natural$, where $\gamma_\natural>0$ denotes the threshold on $\gamma$
into which these constants, and the constant $C_{\mathrm{ff}}$ of (ii), are absorbed.
The factor is drawn once, at the step at which the piece is
created, so nothing compounds over the steps.
\item[(ii)] Take the default pieces and the ring pieces at points of one $\pi_k$-cube
disjoint from $P$. After each
piece has kept its own factor $e^{-\kappa'|\omega|_{k+1}}$, their sum is at most
\begin{equation}\label{4.29:eq-far-field-count-a}
\begin{gathered}
2^d\,C_{\mathrm{ff}}\,e^{-5R},\\
C_{\mathrm{ff}}:=\sum_{u\ge0}\bigl(6+\log(u+2)\bigr)e^{-\kappa'(1-1/L)u}<\infty .
\end{gathered}
\end{equation}
The constant $C_{\mathrm{ff}}$ depends only on $L$, $\kappa$ and $\beta$. It enters
$\gamma_\natural$ through $e^{-5R}\le\gamma^{10}$.
\item[(iii)] At good cubes, the sum over the generations $j\le s$ sharing a tile is
bounded, uniformly in $s$ and $k$, by the second part of the count lemma. The sum over
$s$ is controlled by the factor $e^{-\kappa'(1-1/L)u}$.
\end{enumerate}
No term is summed up to the number of steps, and no cap is used.
\end{remark}

\begin{proof}
\emph{Transfer from tile points to cubes.} The bound (g3) is stated at tile points $x$.
It is carried by the cubes $\square\subset\tilde\Delta^1$, where $\Delta$ is the tile of
$x$ and has scale $s$. Let $y\in\square$.

Then $|x-y|_\infty\le 2M\ell_s<M\ell_{n(x)+1}$, because $n(x)\ge s$ and $L\ge3$. Suppose
$y\in\Omega_{n(x)+2}$. The width of the strata, applied with $a=n(x)+1$, gives
\begin{equation*}
\operatorname{dist}_\infty(\Omega_{n(x)+2},\Omega_{n(x)+1}^{c})\ge M\ell_{n(x)+1}
>|x-y|_\infty .
\end{equation*}
This forces $x\in\Omega_{n(x)+1}$, which the level structure excludes. Hence
$y\notin\Omega_{n(x)+2}$. Since $y\in\Omega_{n(y)}$ and the regions are nested,
$n(y)\ge n(x)+2$ would give $y\in\Omega_{n(x)+2}$; therefore $n(y)\le n(x)+1$, so
$u(y)\ge u(x)-1$.

By (g1) at $y$, every path $\omega$ from $\Omega_{k+1}$ to $y$ therefore has
$|\omega|_{\mathbb{B}}\ge 7LR+u(x)-1$. For such a path, the inequality
$|\omega|_{k+1}\le L^{-1}|\omega|_{\mathbb{B}}$ of
Definition~\ref{4.29:def-path-length-decay} and the floor give
\begin{equation*}
\mathsf{d}|\omega|_{\mathbb{B}}-\kappa'|\omega|_{k+1}
\ge\bigl(4\log L+1+\kappa'(1-1/L)\bigr)|\omega|_{\mathbb{B}} .
\end{equation*}
Let $s$, $n$, $u$ be the quantities of $x$. Since $k-n\le u+1$, the bound (g2) at $x$ gives
\begin{equation*}
L^{4(k-s)}(1+k-s)^{1/2}=L^{4(k-n)}L^{4(n-s)}(1+k-s)^{1/2}
\le L^{4(u+1)}\,eL^{21}R^5e^{u}.
\end{equation*}
Insert $|\omega|_{\mathbb{B}}\ge7LR+u-1$ into the first of these two displays. The part
$(4\log L+1)u$ of the exponent cancels $L^{4u}e^{u}$, and the part $\kappa'(1-1/L)u$
leaves $e^{-\kappa'(1-1/L)u}$; the shift by $-1$ costs the factor $eL^4e^{\kappa'(1-1/L)}$,
whose last factor is compensated by $e^{-\kappa'(1-1/L)\cdot7LR}$ since $7LR\ge1$. Hence
\begin{equation*}
L^{4(k-s)}(1+k-s)^{1/2}e^{-\mathsf{d}|\omega|_{\mathbb{B}}}
\le e^{-\kappa'|\omega|_{k+1}}\,e^{-\kappa'(1-1/L)u}\,e^{2}L^{29}R^5e^{-(4\log L+1)7LR}.
\end{equation*}
Finally $2+29\log L+5\log R\le(28\log L+7)LR-6R$ for $L\ge13$ and $R\ge1$, so the
product $e^{2}L^{29}R^5e^{-(4\log L+1)7LR}$ is at most $e^{-6R}$. Thus (g3)$^+$, and with
it (g3), hold with the quantities of $x$ for every path from $\Omega_{k+1}$ to any point of
$\square$.

\emph{Clause (i).} By (g3), the factor $e^{-\mathsf{d}|\omega|_{\mathbb{B}}}$ along any
such connection pays three things at once: the factor $L^{4(k-s)}$, the drift
$(1+k-s)^{1/2}$, and the factor $e^{-\kappa'|\omega|_{k+1}}$, leaving $e^{-6R}$. Write
\begin{equation*}
L^{4(k-s)}=L^{4(k-n)}L^{4(n-s)} .
\end{equation*}
The first factor is Ba{\l}aban's count in the first part of the count lemma; the second is
the extra native count. The leftover satisfies $e^{-6R}\le x_{k+1}^{-6}\le\gamma^{12}$ by
\eqref{4.29:eq-far-field-paths-a}. Hence a constant depending only on $(L,M,\mathfrak{p})$
is absorbed once $\gamma\le\gamma_\natural$.

The factor is produced by the single derivative in the direction
$\bar\zeta_\square H_k(B')$ at the creation of the piece. No later step reproduces it.

\emph{Number of generations per point.} Consider first a default piece or a ring piece.
By Proposition~\ref{4.29:prop-native-cover}, at most one generation has a default tile
at $x$, namely $j=l(x)$. On the ring one also has $\mathfrak{s}_j=j=l(x)$. So each point
carries one generation.

Now fix a $\pi_k$-cube disjoint from $P$, so that every current level present in it
satisfies $n\le k$, and fix such a level $n$. If no arrested component contains a point of
level $n$ of the cube, the hypothesis on generations in arrested components gives $l=n$ at
all these points, so one generation occurs. Otherwise let $N_0$ be the largest
$N_0(W)$ over the arrested components $W$ that have a point of level $n$ in the cube,
and among these components let $W_*$ be one with the least $k_0(W)$.

By the hypothesis on generations in arrested components, every generation $l$ of a point of
level $n$ in the cube lies in one of three sets:
\begin{equation*}
[k_0(W),n],\qquad \{n-1,n\},\qquad \{n\},
\end{equation*}
where $W$ is one of these components. All of these lie in the interval
$[\min\{k_0(W_*),n-1\},n]$. Since $n\le k_{W_*}$, the number of its elements is
\begin{equation*}
\max\{n-k_0(W_*)+1,\,2\}\le N_0(W_*)+2\le N_0+2 .
\end{equation*}
So at most $N_0+2$ generations occur.

For each such $W$ we have $n\le k_W$, hence $t_W=k+1-k_W\le k+1-n$. If $n\le k-1$
then $k+1-n=u+2$. If $n=k$ then $u=0$ and $t_W\le1\le u+2$. In either case $t_W\le u+2$.

By the radius comparisons,
\begin{equation*}
R_{k_0(W)+1}\le L(L+1)\,t_W^{\beta_0}R\le L(L+1)(u+2)^{\beta_0}R .
\end{equation*}
With $L+1\le L^2$ this gives
\begin{align*}
N_0(W)&=1+\log_LR_{k_0(W)+1}\\
&\le 1+\log_L\bigl(L(L+1)\bigr)+\beta_0\log_L(u+2)+\log_LR\\
&\le 4+\beta_0\log_L(u+2)+\log_LR .
\end{align*}
Since $L\ge3$ we have $\log_L\le\log$ on $[1,\infty)$. Moreover $\beta_0\le1$,
$u+2\ge1$ and $R\ge1$. Therefore
\begin{equation*}
N_0+2\le 6+\log(u+2)+\log R .
\end{equation*}
The single generation of the case in which no arrested component contains a point of level
$n$ of the cube also satisfies this bound.
All of this uses only $L\ge3$, which holds since $L\ge13$.

\emph{Distinct levels in one cube have distinct $u$.} Let a $\pi_k$-cube disjoint from
$P$ be given. By the level structure, $\Omega_k$ and $\Omega_{k+1}$ are unions of
$\pi_k$-cubes, and $\Omega_{k+1}\subset P$. So the cube lies either in
$\Omega_k\setminus\Omega_{k+1}$ or in $\Omega_k^{c}$.
\begin{itemize}
\item In the first case every point of the cube has $n=k$, so $u=0$.
\item In the second case every point of the cube has $n\le k-1$. There the map
$n\mapsto u=k-1-n$ is injective.
\end{itemize}
In particular the levels $k-1$ and $k$ never occur in the same $\pi_k$-cube.

\emph{Proof of \eqref{4.29:eq-far-field-count-a}.} Fix a level $n$ present in the cube
and a generation. By the cube count, the pieces of that generation in the cube number at
most $2^dL^{4(k-s)}$. Each carries the drift $(1+k-s)^{1/2}$ and the factor
$e^{-\mathsf{d}\,d_{\mathbb{B}}(\square,\Omega_{k+1})}$. Since
$d_{\mathbb{B}}=\inf_\omega|\omega|_{\mathbb{B}}$, this factor equals
$\sup_\omega e^{-\mathsf{d}|\omega|_{\mathbb{B}}}$, the supremum running over paths from
$\Omega_{k+1}$ to $\square$.

By (g3)$^+$, transferred to cubes as above, the contribution of this generation is at
most
\begin{equation*}
2^d\sup_\omega e^{-\kappa'|\omega|_{k+1}}\;e^{-6R}\;e^{-\kappa'(1-1/L)u} .
\end{equation*}
The factor $2^d$ comes from the cube count and is absorbed into the constants depending
only on $(L,M,\mathfrak{p})$; the factor $e^{-\kappa'|\omega|_{k+1}}$ is kept by each
piece.

We sum over the at most $6+\log(u+2)+\log R$ generations of each level. We then sum over
the levels present, which have distinct values of $u\ge0$. The sum of the default and
ring pieces in the cube is therefore at most
\begin{equation*}
2^d\sup_\omega e^{-\kappa'|\omega|_{k+1}}\;e^{-6R}
\sum_{u\ge0}\bigl(6+\log(u+2)+\log R\bigr)e^{-\kappa'(1-1/L)u} .
\end{equation*}

Split the last sum as
\begin{equation*}
\sum_{u\ge0}\bigl(6+\log(u+2)+\log R\bigr)e^{-\kappa'(1-1/L)u}
=C_{\mathrm{ff}}+\log R\sum_{u\ge0}e^{-\kappa'(1-1/L)u} .
\end{equation*}
Since $6+\log(u+2)\ge1$ for every $u$, the geometric sum on the right is at most
$C_{\mathrm{ff}}$. The whole expression is therefore at most $C_{\mathrm{ff}}(1+\log R)$.

Finally $1+\log R\le e^R$, so
\begin{equation*}
2^dC_{\mathrm{ff}}(1+\log R)e^{-6R}\le 2^dC_{\mathrm{ff}}e^{-5R},
\end{equation*}
with each piece keeping its own factor $e^{-\kappa'|\omega|_{k+1}}$. This is
\eqref{4.29:eq-far-field-count-a}.

The series defining $C_{\mathrm{ff}}$ converges, because $\kappa'(1-1/L)>0$ and
$\log(u+2)$ grows more slowly than every exponential. From $e^{-6R}\le\gamma^{12}$ we get
$e^{-R}\le\gamma^2$, hence $e^{-5R}\le\gamma^{10}$.

\emph{Good cubes.} At good cubes, all generations $j\le\mathfrak{s}(x)$ share the tile
$\Delta_{\mathfrak{s}(x)}(x)$, by Proposition~\ref{4.29:prop-native-cover}. Their sum is
the one bounded in the second part of the count lemma, whose four kinds of terms are
treated as follows.
\begin{itemize}
\item For the terms of the first kind the factor per point is $e^{-\delta_0ML^{s-j}}$.
\item For the terms of the second kind the factor is $L^{-(5-\beta)(s-j)}$ per point $z$.
Net of the number $L^{4(s-j)}$ of points $z$, this leaves $L^{-(1-\beta)(s-j)}$. This is
summed over $j$ as a geometric series, in the manner of the geometric sum of
\cite[(2.45)]{B-CMP119}:
\begin{equation*}
\sum_{j\le s}L^{-(1-\beta)(s-j)}<\bigl(1-L^{-(1-\beta)}\bigr)^{-1}.
\end{equation*}
\item The terms of the third kind are governed by \cite[(3.48)]{B-CMP119}.
\item The pieces of $\mathbf{R}$ are governed by the count lemma's clause for them.
\end{itemize}
The sum over the scales $s$ is controlled by the factor $e^{-\kappa'(1-1/L)u}$ of
(g3)$^+$. No sum above runs up to the number of steps, and no cap is used.
\end{proof}

\begin{lemma}[Thresholds across scales and the window rim]\label{4.29:lem-scale-thresholds}
Assume the hypotheses of Proposition~\ref{4.29:prop-native-cover}; in particular $L\ge 13$, $M\ge L^4$
and $M>L^2M_1$. Assume the following statements.
\begin{itemize}
\item[(A)] The nesting lemma for the radii across scales: for every integer $v\ge 1$,
\begin{equation*}
(1+7\beta)(1+\beta_0)(1+v)^{1/2}\le L^{v}.
\end{equation*}
\item[(B)] The expansion-ratio lemma, with its inequality $(1+v)^{5/2}L^{-(\beta/2)v}\le C_{\beta,L}$,
its bound on $\sup|B|$, which comes from \cite[(3.32)]{B-CMP109}, and its summation clause for the
factors $\varpi_{js}:=\max\{1,(x_j/x_s)^2\}$ occurring in the terms of the large-field
operation $\mathbf{R}$.
\item[(C)] Clauses (c)--(e) of the rim lemma for the window.
\item[(D)] The bound of a datum by the field: for an element $c$ of level $j$,
$|B(c)|\le c_Q\ell_j\sup|A|$, the supremum being taken over the blocks of $c$, provided
$64\,\ell_j\sup|A|\le\alpha_1^{(98)}$ there, where $\alpha_1^{(98)}$ is the radius of \cite{B-CMP98}
appearing in \eqref{4.29:eq-constants-order-b}.
\item[(E)] The layered decay bound for the first variation of $H_j(\square_0,\cdot)$ in a direction
of the data, taken at the top level $j$ of $\mathbf{B}_j(\square_0)$, with constant $C_{90}$ and rate
$\delta_0/8$.
\item[(F)] The clause of Lemma~\ref{4.29:lem-collar-sources} for the domain $\Omega=\square_0$, without
the weight $\ell_{m(c)}$ and at the rate $\delta_0/16$: under the floor
$L^{d-1}e^{-\delta_0M_1/16}\le\tfrac12$ among the hypotheses of
Proposition~\ref{4.29:prop-native-cover}, for every level-$j$ block $y$ of $\mathbf{B}_j(\square_0)$ and
every function $\phi\ge0$ on elements,
\begin{align*}
&\sum_{c\in\mathbf{B}_j(\square_0),\,m(c)<j}e^{-\delta_0d(y,c)/16}\phi(c)\\
&\qquad\le2e^{\delta_0/16}C_dL^d\sum_{\hat c}e^{-\delta_0d^{(j)}(y,\hat c)/16}\hat\phi(\hat c),
\end{align*}
where $\hat c$ ranges over the level-$j$ blocks tiling the part of $\square_0$ that carries elements of
level $<j$, $\hat\phi(\hat c)$ is the maximum of $\phi$ over the elements inside $\hat c$, and
$d^{(j)}$ is the scaled distance with all of $\square_0$ counted at level $j$.
\item[(G)] The regularity corollary for the hybrid background.
\item[(H)] The arrest theorem, with its clause on newly created regions and its clauses (v) and (vi).
\item[(I)] The lemma on late levels, which carries a bound on $\sup|B|$ of the same kind as that of
the expansion-ratio lemma of hypothesis (B).
\item[(J)] The exponents satisfy $q\ge p_0+1$, and the exponents and constants of the radii
\cite[(2.28)]{B-CMP119} satisfy $q_0=q_1$ and
$C_1\ge C_0$; and $\gamma\le\gamma_\natural$, where $\gamma_\natural$ denotes the coupling ceiling of
this section, whose conditions include the hypotheses of Lemma~\ref{4.29:lem-radius-conversion} on
$\gamma$ and $x_{k+1}$, the ceilings \eqref{4.29:eq-constants-order-b}, the ceiling among
\eqref{4.29:eq-constants-order-c} on the members at the hybrid background, and the condition
$3b_{\mathrm{out}}\le 2C^{(172)}_1\varepsilon_1^{\mathrm{abs}}$ of clause (c) of the rim lemma, with
$b_{\mathrm{out}}$ as in \eqref{4.29:eq-scale-thresholds-b}, $\varepsilon_1^{\mathrm{abs}}$ the
polydisc radius of \cite[(172)]{B-CMP102b} and $C_1^{(172)}$ the absolute constant there.
\end{itemize}
Let $0\le j\le s$ and $v:=s-j$, and write $\alpha_{\cdot,i}$ for either of the radii
$\alpha_{0,i},\alpha_{1,i}$. Here $C(L)$, $C$ and $C'$ denote constants independent of $j$, $s$, $k$
and $v$, $C'$ depending on $L$ and $\mathfrak{p}$ only; $C_\zeta$ is the gradient constant of clause (a)
of Proposition~\ref{4.29:prop-native-cover}, and $C_d$ the constant of
Lemma~\ref{4.29:lem-collar-sources}. Then the following hold.

\textup{($\tau$1)} At the same scale, clause (d2) of Proposition~\ref{4.29:prop-window-minimiser}
applies, and the inequality it requires, the first line of \eqref{4.29:eq-scale-thresholds-a}, holds.

\textup{($\tau$2)} Across scales, for $v\ge1$, the inequalities of the second line of
\eqref{4.29:eq-scale-thresholds-a} hold.

\textup{($\tau$3)} The expansion ratios obey the third line of \eqref{4.29:eq-scale-thresholds-a},
with $\alpha_{3,j}$ as in \eqref{4.29:eq-scale-thresholds-b}; this bound grows with $v$ and is not
bounded by $\tfrac14$.

\textup{($\tau$4)} A field whose layered members are bounded by $a\,\ell_{n(x)}^{-p}$ has level-$s$
members bounded by $2a$ on $\square_0$ (for $H_k$, $n(x)$ is read as $\min\{n(x),k\}$); multiplication
by $\bar\zeta_\square$ costs at most the factor $1+C_\zeta$; and
$(t\zeta'_\square+t_\square\bar\zeta_\square)H_k(B')$ obeys \cite[(3.14), (3.31)]{B-CMP109} at level
$s$ with constant at most the left member of the fourth line of \eqref{4.29:eq-scale-thresholds-a},
itself bounded by the right member, where
$d_{\mathbb{B}}=d_{\mathbb{B}}(y,\operatorname{supp}B')$ at the evaluation point $y$. Moreover $H_k$ and
$Q_j$ are holomorphic on
$\mathfrak{S}(Y)$ with weighted constants bounded by $L(L+1)N_0^{\beta_0}R_{k_W}\alpha_{\cdot,n}$,
$n\le k_W$, where $R_{k_W}$ is the constant $R$ of \cite{B-CMP122a} taken at the arrest time $k_W$ of
the piece $W$, and $N_0:=1+\log_LR_{k_0+1}$, with $k_0$ the index so denoted in the structure clause
of Definition~\ref{4.29:def-analyticity-space}; these constants tend to $0$ uniformly as $\gamma\to0$;
the same holds at the hybrid background and
for $\mathcal{H}'$, $Q'$ on $\Lambda_j^\sim\cap Y$.
\begin{equation}\label{4.29:eq-scale-thresholds-a}
\begin{aligned}
&2B_3C_{\mathrm{reg}}\varepsilon_s\le\tfrac12(1-\beta)\alpha_{0,s},\\
&(1+7\beta)\alpha_{0,s}L^{-2v}\le\alpha_{0,j},\qquad
(1+2\beta)\alpha_{1,s}L^{-v}\le\alpha_{1,j}\qquad(v\ge1),\\
&\alpha_{\cdot,s}/\alpha_{3,j}\le C(L)(1+v)^{1/2},\\
&CB_3\delta_k^\sharp e^{-\mathsf{d}d_{\mathbb{B}}}
\le C'(1+k-s)^{1/2}e^{-\mathsf{d}d_{\mathbb{B}}}\alpha_{\cdot,s}.
\end{aligned}
\end{equation}

\textup{($\tau$5)} (The rim of the window.) Let the representation \cite[(3.18)]{B-CMP109} be entered
at the cover cube
$\square$ of scale $s$, let $b$ be as in clause (g) of the inputs to the transfer of the small-field
analysis ($b=s$, and $b=k+1$ where
$\square_0\subset\Omega_{k+1}$), let $Y\supset\square_0$ and $(\mathbb{U},\mathbb{J})\in\mathfrak{S}(Y)$.
Define $b_{\mathrm{out}},\varepsilon_{\mathrm{w}},\alpha_{2,j},\alpha_{3,j},C_{\mathrm{rim}},C_\Sigma,
C_{\mathrm{w}}$ by \eqref{4.29:eq-scale-thresholds-b}, where $K_{\mathrm{w}}$ is an absolute
constant containing the factor $c_QB_5(1+2B_3)O(1)$, $\beta_s>0$ is a fixed
constant entering the radius $\alpha_{3,j}$ there, independent of $j$, $s$, $k$ and $v$ (its subscript
does not denote the scale $s$), and $c_2>0$ is taken so
small
that $c_2\alpha_{3,j}$ does not exceed the multiple of $\alpha_{\cdot,j}$ admitted in the estimate of
the finer sources entering \cite[(3.14)]{B-CMP109}; $\alpha_{2,j}$ so defined is the one meant in every
use of $\alpha_{2,j}$ in this section. Then:
\begin{itemize}
\item[(a)] the data $B$ of $U_j(\square_0,\cdot)$ on the rim
$\mathfrak{r}_\square:=\square_0\setminus\tilde\square^4$ obey $|B(c)|\le b_{\mathrm{out}}$ at every
determining element $c$ of $U_j(\square_0)$, whatever its level;
\item[(b)] for $E\in\{\tilde\square^2,\tilde\square^3\}$, $E^+$ the union of $E$ and one layer of
$\pi_s$-cubes, $B_t$ in the polydisc of \cite[Proposition~9]{B-CMP102b}, and a direction
$\mathfrak{b}$ supported on the elements of $\mathbf{B}_j(\square_0)$ outside $E^+$, all members
$e\le2$ ($D^*D$ and $\Delta$ included) obey on $E$
\begin{equation}\label{4.29:eq-scale-thresholds-1}
\ell_j^{1+e}\bigl|\nabla^e\bigl\langle(\delta/\delta B)H_j(\square_0,B_t),\mathfrak{b}\bigr\rangle\bigr|
\le C_{90}C_{\mathrm{rim}}e^{-\delta_0(ML^v-2)/16}\sup|\mathfrak{b}|;
\end{equation}
\item[(c)] under the floor $\mathsf{d}\ge(1+4\beta)\kappa+4\log L+1$ of
\eqref{4.29:eq-constants-order-c}, which at the decay rate $\delta_0/8$ reads
$\delta_0M\ge8((1+4\beta)\kappa+4\log L+1)$, and the rim floor on $M$, the inequality
$\delta_0(M-2)/16\ge\log(8C_{\mathrm{w}}M)$ of \eqref{4.29:eq-scale-thresholds-b},
the bound $\varepsilon_{\mathrm{w}}\le\tfrac14L^{-6v}\alpha_{2,j}$ holds for every $v\ge0$:
\end{itemize}
\begin{equation}\label{4.29:eq-scale-thresholds-b}
\begin{gathered}
b_{\mathrm{out}}:=K_{\mathrm{w}}(B_3+B_5)M\alpha_{0,b},\qquad
\varepsilon_{\mathrm{w}}:=2C_{90}C_{\mathrm{rim}}b_{\mathrm{out}}e^{-\delta_0(ML^v-2)/16}
\le\tfrac14L^{-6v}\alpha_{2,j},\\
\delta_0(M-2)/16\ge\log(8C_{\mathrm{w}}M),\\
\alpha_{2,j}:=c_2\alpha_{3,j},\qquad
\alpha_{3,j}:=\frac{\beta_s\min\{\alpha_{0,j},\alpha_{1,j}\}}{2L^2B_3}
=\frac{\beta_s\alpha_{0,j}}{2L^2B_3},\\
C_{\mathrm{rim}}:=C_\Sigma\bigl(1+2e^{\delta_0/16}C_dL^d\bigr),\qquad
C_\Sigma:=2d\sum_{n\in\mathbb{Z}^d}e^{-\delta_0|n|_\infty/16},\\
C_{\mathrm{w}}:=\frac{4C_{90}C_{\mathrm{rim}}K_{\mathrm{w}}B_3(B_3+B_5)L^2(1+\beta_0)}{c_2\beta_s}.
\end{gathered}
\end{equation}
The constant $C_{\mathrm{rim}}$ depends on $d,L,\delta_0$ only; $C_{\mathrm{w}}$ depends only on
$d,L,\delta_0$ (through $C_{\mathrm{rim}}$, $C_{90}$ and $K_{\mathrm{w}}$), on $\beta_0,\beta_s,c_2$ and
on the constants $B_3,B_5$; in particular on none of $M,C_0,C_1,A_0,A_1$. Consequently each of the six
entries of the rim data listed in
clause (g) of the inputs to the transfer of the small-field analysis is at most
$\tfrac14L^{-6v}\alpha_{2,j}$ in all four members of
\cite[(3.14)]{B-CMP109}, at every window of this section, $v=0$ and $v=1$ included.
\end{lemma}

\begin{proof}
Throughout, the flow line of Lemma~\ref{4.29:lem-radius-conversion},
\eqref{4.29:eq-radius-conversion-a}, whose hypotheses are among those of $\gamma\le\gamma_\natural$,
gives for $j\le s$
\begin{equation}\label{4.29:eq-scale-thresholds-2}
\alpha_{\cdot,s}\le(1+\beta_0)(1+v)^{1/2}\alpha_{\cdot,j},
\end{equation}
and the same inequality with $\varepsilon_s,\varepsilon_j$ in place of $\alpha_{\cdot,s},\alpha_{\cdot,j}$.
By hypothesis (J), $q_0=q_1$ and $C_1\ge C_0$, so that $\min\{\alpha_{0,j},\alpha_{1,j}\}=\alpha_{0,j}$;
this is the second equality in the definition of $\alpha_{3,j}$ in \eqref{4.29:eq-scale-thresholds-b}.

\emph{Clause \textup{($\tau$1)}.} At the same scale the threshold is clause (d2) of
Proposition~\ref{4.29:prop-window-minimiser} (see \eqref{4.29:eq-window-minimiser-a}), which bounds
$|U_{k+1}(\partial p)-1|$ by $C_{\mathrm{reg}}\varepsilon_s\eta^2\ell_s^{-2}$. Its use requires the
first line of \eqref{4.29:eq-scale-thresholds-a},
$2B_3C_{\mathrm{reg}}\varepsilon_s\le\tfrac12(1-\beta)\alpha_{0,s}$. This inequality holds because the
exponents satisfy $q\ge p_0+1$ (hypothesis (J)) and $\gamma$ is small: the ratio
$\varepsilon_s/\alpha_{0,s}$ is a constant times $(\log x_s)^{p_0-q}\le(\log x_s)^{-1}$, which is
small in the coupling regime $x_s\ge\gamma^{-2}$; the smallness required is one of the conditions
of $\gamma\le\gamma_\natural$.

\emph{Clause \textup{($\tau$2)}.} Let $v\ge1$. By \eqref{4.29:eq-scale-thresholds-2} and hypothesis (A),
\begin{align*}
(1+7\beta)\alpha_{0,s}L^{-2v}
&\le(1+7\beta)(1+\beta_0)(1+v)^{1/2}L^{-2v}\alpha_{0,j}
\le L^{-v}\alpha_{0,j}\le\alpha_{0,j},\\
(1+2\beta)\alpha_{1,s}L^{-v}
&\le(1+7\beta)(1+\beta_0)(1+v)^{1/2}L^{-v}\alpha_{1,j}\le\alpha_{1,j},
\end{align*}
where in the second line $1+2\beta\le1+7\beta$ was used. At $v=1$ the first inequality is the
counterpart, at coupling-dependent radii, of the condition $(1+7\beta)L^{-2}\le1$ stated in
\cite{B-CMP109}. This is the second line of
\eqref{4.29:eq-scale-thresholds-a}.

\emph{Clause \textup{($\tau$3)}.} With $\alpha_{3,j}$ as in \eqref{4.29:eq-scale-thresholds-b},
\eqref{4.29:eq-scale-thresholds-2} gives $\alpha_{\cdot,s}/\alpha_{3,j}\le C(L)(1+v)^{1/2}$, the third
line of \eqref{4.29:eq-scale-thresholds-a}. This ratio is unbounded in $v$ and is not at most
$\tfrac14$. In the expansion-ratio lemma (hypothesis (B)) these ratios enter in the combination
$(1+v)^{5/2}L^{-5v}$, and they are compensated there by the inequality
\begin{equation*}
(1+v)^{5/2}L^{-(\beta/2)v}\le C_{\beta,L}
\end{equation*}
of that lemma. The bound on $\sup|B|$ carried by that lemma, which comes from
\cite[(3.32)]{B-CMP109}, is taken over $\tilde\square^4$, the domain on which
\cite[(3.32)]{B-CMP109} is stated, and not over the rim $\mathfrak{r}_\square=\square_0\setminus
\tilde\square^4$; the data of the window's own units on the rim are removed at the cost computed in
clause ($\tau$5) (see clause (g) of the inputs to the transfer of the small-field analysis, and
hypothesis (C)). The same restriction
applies to the expansion-ratio lemma at $\mathbf{A}=0$ and to the lemma on late levels (hypothesis
(I)). In the terms of the large-field operation $\mathbf{R}$ there occurs the factor
\begin{equation*}
\varpi_{js}=\max\{1,(x_j/x_s)^2\}\le(1+v)^2,
\end{equation*}
since $x_j\le(1+v)x_s$ by the comparison of the couplings along the flow used in the proof of
Lemma~\ref{4.29:lem-radius-conversion}; and the summation clause of the expansion-ratio lemma gives,
with the exponent $\kappa_0$ and the constant $C_\lambda$ of that clause,
\begin{equation*}
\sum_{j\le s}\varpi_{js}g_j^{\kappa_0}\le4C_\lambda g_s^{\kappa_0-2}.
\end{equation*}

\emph{Clause \textup{($\tau$4)}, fields.} Let a field have layered members bounded by
$a\,\ell_{n(x)}^{-p}$ at $x$ (for $H_k$, with $n(x)$ read as $\min\{n(x),k\}$). On $\square_0$ one has
$n(x)\ge s$, since the whole window lies in the top stratum $\Omega_s$ (clause (d1) of
Proposition~\ref{4.29:prop-window-minimiser}), so $\ell_s^{p}\,a\,\ell_{n(x)}^{-p}\le a$ at every point;
a H\"older quotient between
points of different strata is bounded through this sup bound, which costs the factor $2$. Hence the
level-$s$ members are bounded by $2a$ on $\square_0$. Next, $0\le\zeta\le1$ gives
$0\le\bar\zeta_\square\le1$, and by clause (a) of Proposition~\ref{4.29:prop-native-cover} together
with $M\ge L^4$, multiplication by $\bar\zeta_\square$ costs at most
\begin{equation*}
1+C_\zeta L^4/M\le1+C_\zeta .
\end{equation*}
Apply these two facts to $H_k(B')$, whose layered members at a point $y$ of level
$\min\{n(y),k\}$ are bounded by $B_3e^{-\mathsf{d}d_{\mathbb{B}}(y,\operatorname{supp}B')}\sup|B'|$
(Definition~\ref{4.29:def-path-length-decay}), with $\sup|B'|\le\delta_k^\sharp$
by \eqref{4.29:eq-path-length-decay-b}. It follows that $(t\zeta'_\square+t_\square\bar\zeta_\square)
H_k(B')$ obeys \cite[(3.14), (3.31)]{B-CMP109} at level $s$ with constant at most
$CB_3\delta_k^\sharp e^{-\mathsf{d}d_{\mathbb{B}}}$. Since
$\delta_k^\sharp=C_{II}\max\{4a_\sharp A_0/A_1,C_1^{(3.43)}\}\delta_k$, this is at most
$C'(1+k-s)^{1/2}e^{-\mathsf{d}d_{\mathbb{B}}}\alpha_{\cdot,s}$, the fourth line of
\eqref{4.29:eq-scale-thresholds-a}, the factor $(1+k-s)^{1/2}$ being that of
\eqref{4.29:eq-scale-thresholds-2} applied between the levels $s\le k$; here $C'$ depends on $L$ and
$\mathfrak{p}$ only and contains the factor $C_{II}\max\{4a_\sharp A_0/A_1,C_1^{(3.43)}\}$ once.
In the far field
$\Lambda_j^\sim\setminus P$, with $P$ the set of Lemma~\ref{4.29:lem-far-field-paths}, this constant is
small by clause (g3) of that lemma, the far-field constants being contained in those paying
$e^{-6R}$; on $P$ it is small directly, without recourse to clause (g3).

\emph{Clause \textup{($\tau$4)}, holomorphy.} By \cite[Proposition~9]{B-CMP102b}, $H_k$ is an analytic
function of $B$ and of the external gauge field configuration $U$. \cite{B-CMP109} states that the
function $P^{(k)}(g_k,U,J,B)$, which is built of the same operators, is ``analytic on the space of all
configurations $U$, $J$ satisfying the conditions (i)--(iii) in Sect.~1, for some $\alpha_0'$,
$\alpha_1'$ sufficiently small, but much bigger than $\alpha_0$, $\alpha_1$''. On $\mathfrak{S}(Y)$
(Definition~\ref{4.29:def-analyticity-space}) the weighted constants of $H_k$ and $Q_j$ are at most
\begin{equation*}
L(L+1)N_0^{\beta_0}R_{k_W}\alpha_{\cdot,n},\qquad n\le k_W,
\end{equation*}
where $R_{k_W}$ is the constant $R$ of \cite{B-CMP122a} taken at the arrest time $k_W$ of the piece
$W$ (Definition~\ref{4.29:def-analyticity-space}), and not the radius
$R=R_{k+1}$ of Lemma~\ref{4.29:lem-far-field-paths}; no comparison of $R_{k_W}$ with $R_{k+1}$ is made
or needed. By the arrest theorem (hypothesis (H)),
\begin{equation*}
R_{k_W}\alpha_{\cdot,n}\le LC(1+\beta_0)(\log x_n)^{r+q}x_n^{-1/2},\qquad x_n\ge\gamma^{-2},
\end{equation*}
and
\begin{equation*}
N_0=1+\log_LR_{k_0+1}\le2+r\log_L\log x_{k_0+1},\qquad x_{k_0+1}\le x_n+B_{\mathrm{fl}}N_0 ,
\end{equation*}
where $B_{\mathrm{fl}}$ is the constant written $B$ in the hypothesis $x_{k+1}\ge B+c_5$ of
Lemma~\ref{4.29:lem-radius-conversion}, which bounds the increments of the flow. Hence
the weighted constants tend to $0$ uniformly as $\gamma\to0$, and lie inside the absolute radii,
independent of the coupling, of the analyticity domains of the operators of Ba{\l}aban's earlier
papers on which \cite{B-CMP119} rests, and inside the radii $\alpha_0',\alpha_1'$ of
the statement quoted above. The same holds at the hybrid background
$\mathfrak{V}=[\mathbb{U}\mid U_b(\square_0,M^{\cdot}\mathbb{U})]$ of \cite[(3.24)]{B-CMP109}: by the
regularity corollary for the hybrid background (hypothesis (G)), $\mathfrak{V}^h=e^{i\eta A''}U$ with
the same $G$-valued part $U$, $A''=A'$ off $\square_1$, and the level-$n(x)$ members of $A''$ are at
most $L^2K_{\mathrm{rim}}C_\flat B_{\mathrm{h}}M$ times the bound
$L(L+1)N_0^{\beta_0}R_{k_W}\alpha_{\cdot,n}$ displayed above, $K_{\mathrm{rim}}$,
$C_\flat$ and $B_{\mathrm{h}}$ being the constants of that corollary; the ceiling on the members at the
hybrid background in hypothesis (J) places these members inside the radii. The same holds for
$\mathcal{H}'$ and
$Q'$ on $\Lambda_j^\sim\cap Y$, at the levels occurring there, by the clause on the ring in
\eqref{4.29:eq-complex-hypothesis-a}, with absolute radii.

\emph{Clause \textup{($\tau$5)}, part \textup{(a)}.} Let $c$ be a determining element of
$U_j(\square_0)$ on the rim. If $c$ has level $n<j$, the datum is reproduced:
\begin{equation*}
e^{iB(c)}=u(c_-)V(c)u(c_+)^{-1}
\end{equation*}
(\cite{B-CMP98}; \cite[(20)]{B-CMP102b}), where $c_\mp$ are the end-points of $c$. It is bounded by
\cite[(3.26)]{B-CMP109} and by the half of \cite[(3.27)]{B-CMP109} that bounds $u$; these read $V$ and
$u$ only, and give $|B(c)|\le O(1)M\alpha_{0,b}$. If $c$ has level $j$ and $c\subset\square_j$,
hypothesis (D) gives $|B(c)|\le c_Q\ell_j\sup|A|$. By \cite[(19)]{B-CMP102b},
``$|A|<\varepsilon_2(L^j\eta)^{-1}$ \dots\ on $\Omega_j$, $j=0,1,\dots,k$'', with
$\varepsilon_2=B_5\varepsilon_1$; this bound needs only $n_{\mathrm{w}}\ge j$ on the blocks of $c$,
$n_{\mathrm{w}}$ denoting the level function of the window, so
$\ell_j|A(x)|\le\ell_{n_{\mathrm{w}}(x)}|A(x)|<\varepsilon_2$ there. It is applied with
$\varepsilon_1$ taken equal to the size of the data themselves,
\begin{equation*}
\varepsilon_1:=O(1)(1+2B_3)M\alpha_{0,b},
\end{equation*}
a bound by the data and not by the supremum over the polydisc. The proviso of hypothesis (D) holds
there, since $\ell_j|A|<\varepsilon_2$ and $64\varepsilon_2\le\alpha_1^{(98)}$, with $\varepsilon_2$
taken at this value of $\varepsilon_1$, is one of the ceilings \eqref{4.29:eq-constants-order-b}. Hence
$|B(c)|\le c_QB_5O(1)(1+2B_3)M\alpha_{0,b}$. In neither case is the other half of
\cite[(3.27)]{B-CMP109}, the only part carrying the radius $\alpha_1$, used. The absolute constant
$K_{\mathrm{w}}$ is fixed so that $b_{\mathrm{out}}=K_{\mathrm{w}}(B_3+B_5)M\alpha_{0,b}$ dominates
both bounds; it contains the factor $c_QB_5(1+2B_3)O(1)$. This is clause (c) of the rim lemma
(hypothesis (C)); its condition $3b_{\mathrm{out}}\le2C_1^{(172)}\varepsilon_1^{\mathrm{abs}}$, with
$\varepsilon_1^{\mathrm{abs}}$ the polydisc radius of \cite[(172)]{B-CMP102b}, $C_1^{(172)}$ the
absolute constant there, and $B_5$ independent of $\varepsilon_1$, is among the conditions of
$\gamma\le\gamma_\natural$.

\emph{Clause \textup{($\tau$5)}, part \textup{(b)}.} Apply the layered decay bound of hypothesis (E) at
the top level $j$ of $\mathbf{B}_j(\square_0)$. A contour from $y\in E$ to an element $c\not\subset
E^+$ crosses $E^+\setminus E$, which lies in the pure level $j$ and has width
$M\ell_s=ML^v\ell_j$; hence every such contour has scaled level-$j$ length at least $ML^v-2$, so the
scaled distance $d(y,c)$ of Lemma~\ref{4.29:lem-collar-sources} satisfies $d(y,c)\ge ML^v-2$. Split the
rate as
$\tfrac18=\tfrac1{16}+\tfrac1{16}$:
\begin{equation*}
e^{-\delta_0d(y,c)/8}\le e^{-\delta_0(ML^v-2)/16}\,e^{-\delta_0d(y,c)/16}.
\end{equation*}
The remaining sum of $e^{-\delta_0d(y,c)/16}$ over the elements $c$ is bounded by $C_{\mathrm{rim}}$:
the elements of level $j$ give $C_\Sigma$, and the sum over the finer levels is taken by the clause of
Lemma~\ref{4.29:lem-collar-sources} for $\Omega=\square_0$ without the weight $\ell_{m(c)}$, read at the
rate $\delta_0/16$ (hypothesis (F)), which gives the factor $2e^{\delta_0/16}C_dL^d$; it depends on
neither $j$ nor $k$.
This proves \eqref{4.29:eq-scale-thresholds-1}. The constant $C_{\mathrm{rim}}$ depends on
$d,L,\delta_0$ only and imposes no upper bound on any parameter; where $e^{\delta_0/16}\le2$ it is at
most $C_\Sigma(1+4C_dL^d)$.

\emph{Clause \textup{($\tau$5)}, part \textup{(c)}.} By \eqref{4.29:eq-scale-thresholds-2} with $b$ in
place of $s$, $\alpha_{0,b}\le(1+\beta_0)(1+b-j)^{1/2}\alpha_{0,j}$, and $b-j\le v+1$. Inserting
$b_{\mathrm{out}}$ and $\alpha_{2,j}=c_2\beta_s\alpha_{0,j}/(2L^2B_3)$,
\begin{align*}
\frac{\varepsilon_{\mathrm{w}}}{\alpha_{2,j}}
&=\frac{4C_{90}C_{\mathrm{rim}}K_{\mathrm{w}}B_3(B_3+B_5)L^2}{c_2\beta_s}\,
M\,\frac{\alpha_{0,b}}{\alpha_{0,j}}\,e^{-\delta_0(ML^v-2)/16}\\
&\le C_{\mathrm{w}}M(2+v)^{1/2}e^{-\delta_0(ML^v-2)/16}.
\end{align*}
For $v=0$, the rim floor on $M$ gives $e^{-\delta_0(M-2)/16}\le(8C_{\mathrm{w}}M)^{-1}$, hence
$\varepsilon_{\mathrm{w}}/\alpha_{2,j}\le\sqrt2/8\le\tfrac14$. For $v\ge1$ write
\begin{equation*}
C_{\mathrm{w}}M(2+v)^{1/2}e^{-\delta_0(ML^v-2)/16}
=\Bigl[\sqrt2\,C_{\mathrm{w}}Me^{-\delta_0(M-2)/16}\Bigr]
(1+v/2)^{1/2}e^{-\delta_0M(L^v-1)/16}.
\end{equation*}
The bracket is at most $\sqrt2/8$ by the rim floor on $M$. The floor on $\mathsf{d}$ in
\eqref{4.29:eq-constants-order-c}, read as $\delta_0M\ge8((1+4\beta)\kappa+4\log L+1)$,
gives $\delta_0M/16\ge2\log L$, and $L\ge13$ gives $L^v-1\ge12v$; therefore the factor
\begin{equation*}
(1+v/2)^{1/2}L^{6v}e^{-\delta_0M(L^v-1)/16}\le(1+v/2)^{1/2}L^{6v}L^{-24v}
=(1+v/2)^{1/2}L^{-18v}\le1 .
\end{equation*}
Hence $\varepsilon_{\mathrm{w}}/\alpha_{2,j}\le(\sqrt2/8)L^{-6v}\le\tfrac14L^{-6v}$ for every $v\ge0$.
The constant $C_{\mathrm{w}}$ depends only on $d,L,\delta_0$ (through $C_{\mathrm{rim}}$, $C_{90}$ and
$K_{\mathrm{w}}$), on $\beta_0,\beta_s,c_2$ and on the constants $B_3,B_5$; in particular on none of
$M,C_0,C_1,A_0,A_1$;
therefore the rim floor is a lower bound on $M$, imposed after $L$ and $M_1$ are fixed, and it
holds for all sufficiently large $M$ because $M/\log M\to\infty$. It is among the floors
\eqref{4.29:eq-constants-order-c}. Any smaller $\alpha_2$ is admissible in \cite[(3.14)]{B-CMP109};
$C_{\mathrm{w}}$, and the per-cube constant $C^\sharp$ of the ceilings on $\gamma$, depend on $c_2$.

\emph{Use.} By parts (a)--(c), each of the six entries of the rim data listed in clause (g) of the
inputs to the transfer of the small-field analysis --- each a direction or an increment supported
outside
$\tilde\square^3$ or $\tilde\square^4$ and observed on $\tilde\square^2$ or $\tilde\square^3$ --- is at
most $\varepsilon_{\mathrm{w}}\le\tfrac14L^{-6v}\alpha_{2,j}$ in all four members of
\cite[(3.14)]{B-CMP109}. Thus the estimate ``by an arbitrary power of $L^j\eta$ using the exponential
decay properties'' of \cite{B-CMP109} holds at the coupling-dependent radii of \cite{B-CMP119} at every
window of this section, $v=0$ ($j=k$) and $v=1$ included, and the argument of \cite{B-CMP109} at that
place completes the estimate. None of the constants used depends on $k$, $K$ or $N_0$.
At \cite[(3.21)--(3.22)]{B-CMP109} the direction is
$\mathfrak{b}=(1-\zeta'_\square)Q(\xi\mathcal{H}'(B_\square))$, taken at an external field different
from the identity; its supremum is bounded by clause ($\tau$4), through
Lemma~\ref{4.29:lem-collar-staircase}, while $b_{\mathrm{out}}$ bounds the increments of the window's
own data.
\end{proof}

\begin{lemma}[Membership and decay bounds for the interpolated field]
\label{4.29:lem-interpolated-field-bounds}
Let $j\le k$. We use the native-scale cover and its partition of unity of
Proposition~\ref{4.29:prop-native-cover}, the notation of the collar staircase
(Lemma~\ref{4.29:lem-collar-staircase}): the set $\mathbf{B}'$ of elements on which the variable
$B$ of $\mathcal{H}'$ lives, the domains $\Lambda_j\subset D_j$ and
$\Lambda_j^{\sim}$, the function $H_j=\mathcal{H}'$ of \cite[Proposition~9]{B-CMP102b}, the
constant $C_{90}$ and the scaled distance $d(y,c)$ of \eqref{4.29:eq-collar-staircase-a}, and
the constant $c_Q$, the end-blocks $A(c)$ and the map $Q'$ from fields to data whose bound is
\eqref{4.29:eq-collar-staircase-b}; $\eta$ is the spacing of the finest lattice, so that
$\ell_j=L^j\eta$. We also use
the quantities $d_{\mathbb{B}}$ and $\mathsf{d}$ of \eqref{4.29:eq-path-length-decay-a},
$\delta_k^{\sharp}$ of \eqref{4.29:eq-path-length-decay-b}, and the response field
$\mathbf{H}_k=\mathbf{H}_k(B')$, whose decay exponent is fixed in
\eqref{4.29:eq-path-length-decay-a}, evaluated at the datum $B'$ as in
\eqref{4.29:eq-scale-thresholds-a}, together with the constant $B_3$ of that display. The cubes
$\square$ of the cover, their enclosing cubes $\square_0$ and the functions $\zeta_\square$,
$\bar\zeta_\square$ are those of Proposition~\ref{4.29:prop-native-cover}. Put
\begin{equation}\label{4.29:eq-interpolated-field-bounds-a}
C_{\mathfrak{s}}:=1+4e^{\delta_0/8}C_dL^{d-1},
\end{equation}
where $C_d$ is the absolute constant of Lemma~\ref{4.29:lem-collar-sources} and $\delta_0$ is
the absolute decay rate of \eqref{4.29:eq-collar-staircase-a}. Thus $C_{\mathfrak{s}}$ depends on
$d$, $L$ and the absolute rate $\delta_0$ only. For the interpolation parameter $t$,
$|t|\le 1$, put $B(t):=Q'(\eta t\mathbf{H}_k)$, and assume that $B(t)$ lies in the polydisc of
\cite[Proposition~9]{B-CMP102b}. Assume further the following statements.
\begin{itemize}
\item[(i)] \emph{The derivative bound \cite[(190)]{B-CMP102b} in kernel form.} The layered
display \eqref{4.29:eq-collar-staircase-a} holds for the directional derivative. Let $x\in\Delta(y)$
be a level-$j$ point of $\mathbf{B}'$, let $e\in\{0,1,2\}$ (the members $D^*D$ and $\Delta$ included),
and let $B$ lie in the polydisc of \cite[Proposition~9]{B-CMP102b}. Then, for every direction
$\mathfrak{b}$ on the elements of $\mathbf{B}'$,
\begin{equation*}
\ell_j^{1+e}\Bigl|\nabla^e\Bigl\langle\frac{\delta\mathcal{H}'}{\delta B}(B),\mathfrak{b}\Bigr\rangle(x)\Bigr|
\le C_{90}\sum_{c\in\mathbf{B}'}e^{-\frac18\delta_0 d(y,c)}\,|\mathfrak{b}(c)|.
\end{equation*}
\item[(ii)] Lemma~\ref{4.29:lem-collar-sources} (summation of finer sources over the collar),
with the sum \eqref{4.29:eq-collar-sources-a}.
\item[(iii)] Clause $(\tau4)$ of Lemma~\ref{4.29:lem-scale-thresholds} (thresholds across scales
and the window rim), as displayed in \eqref{4.29:eq-scale-thresholds-a}.
\item[(iv)] Clauses (g1) and (g3) of Lemma~\ref{4.29:lem-far-field-paths}, under the floor on
$\mathsf{d}$ displayed in \eqref{4.29:eq-far-field-paths-a}.
\item[(v)] The far-field clause (f) of Proposition~\ref{4.29:prop-native-cover}.
\item[(vi)] The dictionary between the notation of \cite[\S3]{B-CMP109} and that of the
native-scale cover.
\item[(vii)] The statement that the second-order member of the kernel-form bound is not
transferable to a skin.
\end{itemize}
Let $\gamma$ satisfy the ceiling \eqref{4.29:eq-constants-order-a}. Then the following hold.
\begin{itemize}
\item[(N1)] \emph{Membership.} Let $X\subset\Lambda_j$ be a domain of the expansion. On $X$ the
level-$j$ members of $H_j(B(t))$ are bounded by $C_{\mathfrak{s}}C_{90}c_Q$ times the layered
supremum of $|\mathbf{H}_k|$ that clause $(\tau4)$ converts. Here $n(x)$ is the big-structure
level of $x$, and that supremum is the one in the layered bound
\begin{equation}\label{4.29:eq-interpolated-field-bounds-1}
\ell_j|\mathbf{H}_k(x)|=L^{j-n}\ell_n|\mathbf{H}_k(x)|
\le L^{j-n}B_3e^{-\mathsf{d}\,d_{\mathbb{B}}}\sup|B'|,\qquad n=\min\{n(x),k\}\ge j.
\end{equation}
Consequently $H_j(B(t))$ satisfies \cite[(3.14)]{B-CMP109} on $X$ with
$\tfrac13\alpha_{2,j}$ in place of $\tfrac13\alpha_2$, where $\alpha_{2,j}$ is the radius fixed in
\eqref{4.29:eq-scale-thresholds-b}. Neither a floor nor a cap is imposed beyond the ceiling on
$\gamma$.
\item[(N2)] \emph{The decay bound \cite[(3.9)]{B-CMP109}.} For a cube $\square$ of the cover put
\begin{equation}\label{4.29:eq-interpolated-field-bounds-2}
\delta H_j=\Bigl\langle\frac{\delta\mathcal{H}'}{\delta B}(B(t)),\mathfrak{b}_\square\Bigr\rangle,
\qquad
\mathfrak{b}_\square:=DQ'\big|_{\eta t\mathbf{H}_k}\bigl[\eta\bar\zeta_\square\mathbf{H}_k\bigr].
\end{equation}
Suppose $(\square+2\ell_j)\cap(\Lambda_j^{\sim}\setminus D_j)=\emptyset$; this holds for every
good cube, and in particular for every cube with $\square_0\subset\Omega_{k+1}$. Then $\delta H_j$
satisfies \cite[(3.9)]{B-CMP109} with the constants of \cite[\S3]{B-CMP109}. At every other cube,
$\delta H_j$ satisfies \cite[(3.9)]{B-CMP109} with $B_3$ replaced by $C_{\mathfrak{s}}B_3$ and
with the same exponent. In both cases the same holds for the members $e=1,2$.
\end{itemize}
\end{lemma}

\begin{proof}
\emph{The sources of $B(t)$.} Apply \eqref{4.29:eq-collar-staircase-b} with $\mathbf{a}=t\mathbf{H}_k$
and use $|t|\le1$. For every element $c$ of $\mathbf{B}'$, of level $m(c)$,
\begin{equation*}
|B(t)(c)|\le c_Q\ell_{m(c)}\sup_{A(c)}|t\mathbf{H}_k|\le \ell_{m(c)}\varphi(c),
\qquad \varphi(c):=c_Q\sup_{A(c)}|\mathbf{H}_k|.
\end{equation*}

\emph{Proof of \textup{(N1)}.} Let $X\subset\Lambda_j$ and let $x\in X$, $x\in\Delta(y)$. By
Lemma~\ref{4.29:lem-collar-staircase}, $\Lambda_j\subset D_j$, so $x$ is a level-$j$ point of
$\mathbf{B}'$. We may therefore apply \eqref{4.29:eq-collar-staircase-a} at $B=B(t)$, which lies
in the polydisc. With the bound on the sources this gives
\begin{equation*}
\ell_j^{1+e}|\nabla^eH_j(B(t))(x)|\le C_{90}\sum_{c\in\mathbf{B}'}e^{-\frac18\delta_0d(y,c)}
\ell_{m(c)}\varphi(c).
\end{equation*}
We split the sum according to the level of $c$.

The elements with $m(c)=j$ give a sum of the layered form of \cite[(1.45)]{B-CMP122a}. This is
the sum formed in \cite[\S3]{B-CMP109}.

For the elements with $m(c)<j$, the staircase sources, we apply
Lemma~\ref{4.29:lem-collar-sources} with the function $\varphi$ at the level-$j$ block $y$. It
bounds their sum by its constant times
\begin{equation*}
\ell_j\sum_{\hat c}e^{-\frac18\delta_0d^{(j)}(y,\hat c)}\hat\varphi(\hat c).
\end{equation*}
Here $\hat c$ runs over the level-$j$ blocks tiling the collar $\Lambda_j^{\sim}\setminus D_j$, and
$\hat\varphi(\hat c)$ is the maximum of $\varphi$ over the elements inside $\hat c$. The constant of
\eqref{4.29:eq-collar-sources-a} is at most $4e^{\delta_0/8}C_dL^{d-1}=C_{\mathfrak{s}}-1$, and the
sum just displayed is a level-$j$-type sum.

Adding the two contributions, the level-$j$ members of $H_j(B(t))$ on $X$ are bounded by
$C_{\mathfrak{s}}C_{90}c_Q$ times the layered supremum of $|\mathbf{H}_k|$. That supremum is the
quantity which clause $(\tau4)$ converts. At a point of big-structure level $n\ge j$, the layered
bound of $\mathbf{H}_k$ reads as \eqref{4.29:eq-interpolated-field-bounds-1}.

For $\mathbf{H}_k$ the level $n$ is read as $\min\{n,k\}$. Indeed, the determining set of
$\mathbf{H}_k$ is of level $k$ on $\Omega_{k+1}$, by \eqref{4.29:eq-localisation-exact-a}, so no
factor $L^{-1}$ appears at $n=k+1$.

Clause $(\tau4)$, applied with its constant multiplied by $C_{\mathfrak{s}}C_{90}c_Q$, turns this
layered bound into the statement that the field satisfies \cite[(3.14)]{B-CMP109} at level $j$.
The constant is bounded through $\delta_k^{\sharp}$ of \eqref{4.29:eq-path-length-decay-b} and
carries the factor $e^{-\mathsf{d}\,d_{\mathbb{B}}}$. On $P$, where the cover and its functions are
those of \cite[\S3]{B-CMP109}, this constant is small by the argument given there; off $P$, in the
far field, it is small by (g3). In \cite[\S3]{B-CMP109} the analysis requires that
$H_j(B(t))$ satisfy \cite[(3.14)]{B-CMP109} with $\frac13\alpha_2$ once $\varepsilon_1$ is small
enough. Here the requirement is met after multiplying the constant of clause $(\tau4)$ by
$C_{\mathfrak{s}}C_{90}c_Q$; the one new constant, $C_{\mathfrak{s}}$, depends on $(d,L,\delta_0)$
only and is absorbed into the ceiling \eqref{4.29:eq-constants-order-a}. No floor is added.

\emph{The drift.} The factor $(1+k-j)^{1/2}$ carried by $\delta_k^{\sharp}/\alpha_{\cdot,j}$ is paid
as follows. For $j\le n\le k$ we have
\begin{equation*}
(1+k-n)(1+n-j)=1+(k-j)+(k-n)(n-j)\ge 1+k-j,
\end{equation*}
and hence
\begin{equation*}
(1+k-j)^{1/2}\le(1+k-n)^{1/2}(1+n-j)^{1/2}.
\end{equation*}
The second factor is absorbed by the factor $L^{j-n}$ of \eqref{4.29:eq-interpolated-field-bounds-1}.
Indeed, $1+v\le 2^v\le L^{2v}$ for every integer $v\ge0$, so
\begin{equation*}
(1+n-j)^{1/2}L^{-(n-j)}\le 1.
\end{equation*}

The first factor is absorbed by the decay factor. Put $u:=(k-1-n)_+$. At points of
$P\subset\Lambda_k$ the level $n$, read as $\min\{n,k\}$, equals $k$, so $u=0$ and the claim
below holds trivially. Clause (g1) is stated at
points $x\in\Lambda_j^{\sim}\setminus P$; by the far-field clause (f), the cubes not meeting $P$ lie
at distance at least $7LMR_{k+1}\ell_k$ from $\Omega_{k+1}$. At such points (g1) states that every
path $\omega$ from $\Omega_{k+1}$ to $x$ has $|\omega|_{\mathbb{B}}\ge 7LR_{k+1}+u$, so that
$d_{\mathbb{B}}\ge u$ and the factor $e^{-\mathsf{d}\,d_{\mathbb{B}}}$ is at most $e^{-\mathsf{d}u}$.
We claim
\begin{equation*}
(1+k-n)^{1/2}e^{-\mathsf{d}u}\le\sqrt2.
\end{equation*}
For $n\ge k-1$ we have $u=0$ and $1+k-n\le2$. For $n\le k-2$ we have $1+k-n=2+u$. The floor in
\eqref{4.29:eq-far-field-paths-a} gives $\mathsf{d}\ge1$, so
\begin{equation*}
e^{2\mathsf{d}u}\ge1+2u\ge1+\tfrac u2,
\end{equation*}
that is, $2+u\le 2e^{2\mathsf{d}u}$. This proves the claim.

\emph{The radius.} The radius $\frac13\alpha_2$ of \cite[(3.14)]{B-CMP109} is read as
$\frac13\alpha_{2,j}$, where $\alpha_{2,j}$ is the radius fixed in \eqref{4.29:eq-scale-thresholds-b}.
The constant there is chosen so small that $\alpha_{2,j}$ is at most a fixed multiple of the native
room times $\alpha_{\cdot,j}$; here the native room is the margin, in units of $\alpha_{\cdot,j}$,
available in the layered analyticity spaces of the native-scale cover between the radius at which
the field is controlled and the radius required in \cite[(3.14)]{B-CMP109}. The native room is at
least $\beta/2$, and equals $\beta$ at $j=k$,
with $\beta$ the schedule constant of the layered analyticity spaces.

The native room is the relevant one at the domains $X$ that touch the edge, that is, that contain
points $x$ whose big-structure level $n(x)$ and whose native level $l(x)$ both equal $j$.
There no factor $L^{-2}$ is available, and so the space of half radius of \cite[(3.13)]{B-CMP109}
is not available either. The radius $\alpha_{2,j}$ does not depend on $k$, and its choice imposes
no cap in the sense of Definition~\ref{2.3:def-cap}. This proves (N1).

\emph{Proof of \textup{(N2)}.} Consider first the support of $\mathfrak{b}_\square$. Since $Q'$ is
local, $\mathfrak{b}_\square(c)$ depends only on $\bar\zeta_\square\mathbf{H}_k$ on $A(c)$. Hence
$\mathfrak{b}_\square$ is supported on the elements $c$ with $A(c)\cap\operatorname{supp}\zeta_\square\ne\emptyset$.
Since $\zeta_\square\ne0$ only inside $\square$, and the end-blocks of an element of level at most
$j$ have side at most $\ell_j$, every such $c$ satisfies $c\subset\square+2\ell_j$.

Next, the Cauchy bound in \eqref{4.29:eq-collar-staircase-b} gives
\begin{equation*}
|\mathfrak{b}_\square(c)|\le 2c_Q\ell_{m(c)}\sup|\bar\zeta_\square\mathbf{H}_k|.
\end{equation*}
The variation $\delta H_j$ is bounded by input (i) with $\mathfrak{b}=\mathfrak{b}_\square$,
uniformly for $B(t)$ in the polydisc. This input is used as stated. It is not derived here from the
summed display \eqref{4.29:eq-collar-staircase-a}, and it is the form in which it is read in
Lemma~\ref{4.29:lem-collar-staircase} and in clause $(\tau5)$ of
Lemma~\ref{4.29:lem-scale-thresholds}. At the corresponding place,
\cite[(3.21)]{B-CMP109} obtains its bound from the exponential decay of the derivative
$(\delta/\delta B)H_j$.

\emph{Cubes whose enlargement misses the collar.} Let
$(\square+2\ell_j)\cap(\Lambda_j^{\sim}\setminus D_j)=\emptyset$. By
Lemma~\ref{4.29:lem-collar-staircase}, the elements of $\mathbf{B}'$ of level below $j$ lie in the
collar $\Lambda_j^{\sim}\setminus D_j$. Every element $c$ in the support of $\mathfrak{b}_\square$
lies in $\square+2\ell_j$, so it is of level $j$. All sources are therefore of level $j$, and the
bound of input (i) is the bound \cite[(3.9)]{B-CMP109} term by term, through the dictionary (vi).
The rate $\delta_0$ written in \cite{B-CMP109} appears here as $\frac18\delta_0$; this is a change
of letter only.

Every good cube falls into this case. Indeed, for a good cube
\begin{equation*}
\square+2\ell_j\subset\square_0\subset\Lambda_j\subset D_j.
\end{equation*}
In particular every cube with $\square_0\subset\Omega_{k+1}$ falls into this case.

\emph{The remaining cubes.} Their tiles lie in the ring $\Lambda_j^{\sim}\setminus\Lambda_j$ or in
the outermost layer of $\Lambda_j$ of width $M\ell_j$. By the far-field clause (f), they lie at
distance at least $7LMR_{k+1}\ell_k$ from $\Omega_{k+1}$.

At such a cube, Lemma~\ref{4.29:lem-collar-sources} bounds the staircase part of the sum in (i) by
$C_{\mathfrak{s}}-1$ times a sum of level-$j$ type. Every source $c$ satisfies
\begin{equation*}
d(y,c)\ge d^{(j)}(y,c)\ge\operatorname{dist}^{(\xi)}(X,\square)-2,
\end{equation*}
where $\operatorname{dist}^{(\xi)}$ is the distance in units of the $\xi$-lattice of
\cite[(3.9)]{B-CMP109}.
So \cite[(3.9)]{B-CMP109} holds with $B_3$ replaced by $C_{\mathfrak{s}}B_3$ and with the same
exponent.

Two kinds of pieces occur at these cubes. The pieces whose domain $X$ contains a point $z$ of the
set $\Lambda_j^{0}$ of \cite[\S2]{B-CMP119} satisfy the first condition of \cite[(3.5)]{B-CMP109}.
Counterterm pieces created at these cubes may instead have $X\subset\tilde\square^{2}$, with
$\tilde\square^{2}$ the enlarged cube of \cite[(3.5)]{B-CMP109}; this is the second condition of
\cite[(3.5)]{B-CMP109}. These are covered by \cite[(3.9)]{B-CMP109} without the exponential factor
at level-$j$ points, as in the second case of \cite[\S3]{B-CMP109}, that of a big domain $X$. There
the required bound is derived only from two facts: that $H_j(B(t))$ satisfies
\cite[(3.14)]{B-CMP109} with $\frac13\alpha_2$, which is (N1), and that $\delta H_j$ satisfies
\cite[(3.9)]{B-CMP109}.

The factor $C_{\mathfrak{s}}$ is a constant depending on $(d,L,\delta_0)$ and acts on far-field
material of the large-field operation $\mathbf{B}$. Clauses (g1) and (g3) are stated on
$\Lambda_j^{\sim}\setminus P$. There the
surplus factor $e^{-6R_{k+1}}\le\gamma^{12}$ of (g3) pays every constant depending only on
$(L,M,\mathfrak{p})$, and in particular pays $C_{\mathfrak{s}}$.

\emph{Higher members.} \cite[\S3]{B-CMP109} states that the same holds for derivatives and for the
second-order operators. Here these are the members $e=1,2$ of input (i), estimated in the same way.

\emph{Conclusion.} The analysis of \cite[(3.6)--(3.17)]{B-CMP109} therefore runs at every cube with
the notation of \cite{B-CMP109}, the dictionary (vi) and inputs of the native-scale cover only.
The second-order member, which by (vii) is not transferable to a skin, is never required on a skin.
\end{proof}

\begin{proposition}[The near field and the cubes beside the fluctuation support]
\label{4.29:prop-near-field-scope}
Let $k$ be the level of the step, let $j\le k$, and work under the hypotheses of
Proposition~\ref{4.29:prop-native-cover}. The cover $\mathcal{C}_j$ of $\Lambda_j^\sim$,
its scale function $\mathfrak{s}_j$, its partition of unity $\{\bar\zeta_\square\}$, the near-field
set $P$, the good and the default cubes, the domain $\Lambda^{(-9)}$ containing the good cubes,
and the cube $\square_0$ and the enlargements $\tilde\square^2$, $\tilde\square^4$ attached to a
cube $\square$ are those constructed there; for a cube of scale $k$,
$\square_0=B_\infty(\square,5M\ell_k)$. A domain $Y\supset\square_0$ is a localisation domain of a
piece born in the step, as in Definition~\ref{4.29:def-analyticity-space}, and $Y^\sim$ is $Y$
enlarged by the operation $\sim$ of \cite{B-CMP119}. Assume the following.
\begin{itemize}
\item $L\ge 13$, $R_i\ge 7$ for every $i$, $R_i$ the large-field parameters of
Proposition~\ref{4.29:prop-native-cover}, $\beta\le\frac14$, where $\beta$ is the schedule
constant of Definition~\ref{4.4:def-post-step-space}, $\kappa'=(1+4\beta)\kappa$, and the
exponent $\mathsf{d}$ of \eqref{4.29:eq-path-length-decay-a} obeys the floor
$\mathsf{d}\ge\kappa'+4\log L+1$; see \eqref{4.29:eq-far-field-paths-a} and
\eqref{4.29:eq-constants-order-c}.
\item The geometric relations in force at the step, which are among the hypotheses of
Proposition~\ref{4.29:prop-native-cover} and are used here as stated there: the nesting
relations of the domains $\Lambda_i$, $\Omega_i^\natural$, $\Omega_i$ together with their cube
structure, the distance bound between $\Omega^\natural_{i+1}$ and $\Lambda_i^c$ together with the
equality $\Omega_{k+1}=\Omega^\natural_{k+1}$, and the distance bound between $\Lambda_i$ and
$(\Omega_i^\natural)^c$; and the ratio bound $8LR_{j+1}\ge R_j$ of these parameters.
\item Reproduction of minimal configurations on a domain: if $U^*$ is the minimal
configuration for a determining set $\mathcal B_1$ at its data, regular at the levels of
$\mathcal B_1$, and $\mathcal B_2$ is a determining set on the domain finer than or equal to
$\mathcal B_1$ about the domain, then the minimiser for $\mathcal B_2$ whose data are the
$\mathcal B_2$-averages of $U^*$ coincides with the restriction of $U^*$ to the domain up to a
gauge transformation, blocks straddling its boundary included.
\item The thresholds \eqref{4.29:eq-scale-thresholds-a} and \eqref{4.29:eq-scale-thresholds-b}
of Lemma~\ref{4.29:lem-scale-thresholds}, and the bounds \eqref{4.29:eq-interpolated-field-bounds-a}
of Lemma~\ref{4.29:lem-interpolated-field-bounds}.
\item The fluctuation fields and the bound on the sources. As in \cite[\S3]{B-CMP119}, $A_k$
denotes the fluctuation field on $\Omega_{k+1}\cap\Lambda^c_{k+1}$, and $A$ denotes this field on
$\Lambda_{k+1}$. The source $B'$ of the response $\mathbf{H}_k(B')$ of the level-$k$ structure
on $\Omega_{k+1}$ is carried by these fields, $\operatorname{supp}B'\subset\Omega_{k+1}$, and
satisfies $\sup|B'|\le\delta_k^\sharp$ on the domain of \eqref{4.29:eq-path-length-decay-b}, with
$\delta_k^\sharp$ as in \eqref{4.29:eq-path-length-decay-b}, whichever of the fields $A$, $A_k$
carries it. Let
$Q'$ be the induced averaging operation of Lemma~\ref{4.29:lem-localisation-exact}, $\eta$ the
lattice spacing of \cite[\S3]{B-CMP109}, $B(t):=Q'(\eta t\mathbf H_k)$, and, for $j\le k$,
$H_j(B(t))$ the field of \cite[\S3]{B-CMP109} built on the source $B(t)$, whose membership in the
space \cite[(3.14)]{B-CMP109} is in question. The membership of $H_j(B(t))$ at $j=k$ in that space
with $\tfrac13\alpha_{2,k}$ holds on $P$ whatever the position of the
sources, once $3C_NB_3\delta_k^\sharp\le\alpha_{2,k}$, with $\alpha_{2,k}$ of
\eqref{4.29:eq-scale-thresholds-b} and $C_N:=C_{\mathfrak s}C_{\mathrm{ker}}\cdot 2c_Q\cdot C$,
where $C$ is
the constant of the membership display of \eqref{4.29:eq-scale-thresholds-a}, $C_{\mathrm{ker}}$
that of the
kernel bound below, $C_{\mathfrak s}$ that of \eqref{4.29:eq-interpolated-field-bounds-a}, and
$c_Q$ that of \eqref{4.29:eq-collar-staircase-b}; and $A_k$ enters the integrands of
\cite[(3.26), (3.28)]{B-CMP119} through
termwise derivatives of the same localised expansion.
\item Clauses (T3), (T4), (T5) and (T7) of the transfer lemma for the analysis of
\cite[\S3]{B-CMP109}, which rests on the sentence printed without proof in
\cite[p.~276]{B-CMP119} and cited as printed; its clauses used here are (T3), the gains at cubes
of scale $s>j$; (T4), the treatment within
one layer in its own units, in which the chain of \cite[Lemma~4]{B-CMP109} is not entered; (T5),
the count of layer-$j$ cubes inside $Y_0\in\mathbf D_k$ is compensated by the surplus decay once
$d_k(Y_0)\ge1$ and by the per-cube small factors otherwise; and (T7), the distance in the decay
bound is measured in the local units.
\item The identification lemma, which identifies the terms computed with $\Omega_{k+1}$ read on
$B_\infty(\square,5M\ell_k)\subset Y^\sim$ with those computed in the whole-lattice setting.
\item The bound on the induced averaging operation $Q'$ of
Lemma~\ref{4.29:lem-localisation-exact} and on its differential at twists $\mathbf a$ in the
polydisc of that bound, which contains the twists $\mathbf a=t\mathbf H_k(B')$, $t\in[0,1]$, at
which it is used here:
$\bigl|DQ'\big|_{\eta\mathbf a}[\mathbf b](c)\bigr|\le 2c_Q\ell_{m(c)}\sup|\mathbf b|$, where
$m(c)$ is the
level of the
element $c$ and the supremum is over the union of the two blocks containing the end-points of $c$.
\item The kernel form of the decay bound for the derivative of the response, that is, the
bound of \eqref{4.29:eq-collar-staircase-a} written for the directional derivative, with
constant $C_{\mathrm{ker}}$ and with the
source constant $c_Q$ of \eqref{4.29:eq-collar-staircase-b}.
\item The standing choice $q_0=q_1=:q$ of the exponents of the radii \cite[(2.28)]{B-CMP119},
with $q\ge p_0+1$.
\item The coupling regime $x_i\ge\gamma^{-2}$ at $i=k,k+1$, with $\gamma\le e^{-1/2}$, and the
ceiling \eqref{4.29:eq-constants-order-a}.
\end{itemize}
Then, for $\square\in\mathcal{C}_j$ of scale $s$:
\begin{enumerate}
\item[(1)] \emph{Near field.} If $\square_0\subset\Omega_{k+1}$, then $\square\subset\Omega_{k+1}
\subset P$, $s=k$, the cube, its overlaps and $\bar\zeta_\square=\zeta_\square$ are those of
\cite[\S3]{B-CMP119}, and $\square$ is good, $\square\subset\Lambda^{(-9)}$. Every
$\square\subset Y\subset\Lambda_{k+1}$ is such a cube; conversely, $\square_0\not\subset
\Omega_{k+1}$ forces $\square\not\subset\Lambda_{k+1}$, and the pieces of such a cube are terms
of the boundary part $\mathbf{B}^{(k+1)}$ of the effective action in \cite{B-CMP119}. At a cube
with $\square_0\subset\Omega_{k+1}$ and for $X\subset
\tilde\square^2$, the analysis \cite[(3.18)--(3.28)]{B-CMP109}, \cite[Lemma~4]{B-CMP109} and the
completeness of the local family $\{X\in \mathbf D_j: X\subset\tilde\square^2\}$ hold as in
\cite{B-CMP109} for every $j\le k$, $j=k$ included. For $j<k$ the same holds at every good cube
of scale $s>j$, with the gains (T3) of the transfer lemma and the background of
Proposition~\ref{4.29:prop-complex-hypothesis}. The test $\square_0\subset\Omega_{k+1}$ reads
$\Omega_{k+1}$ only on $B_\infty(\square,5M\ell_k)\subset Y^\sim$, and is always satisfied in the
whole-lattice setting, so that the identification lemma applies.
\item[(2)] \emph{The other cubes of scale $s=j$.} At every cube of scale $s=j$ with
$\square_0\not\subset\Omega_{k+1}$, default or not (for $j<k$: every scale-$j$ cube, none of
which meets $P$), the terms, $X\subset\tilde\square^2$ included, are treated by
\cite[(3.6)--(3.17)]{B-CMP109} alone, with membership from \eqref{4.29:eq-complex-hypothesis-a}
and fields from \eqref{4.29:eq-scale-thresholds-a} and \eqref{4.29:eq-interpolated-field-bounds-a};
no factor $L^{j-s}$, no window and no completeness is used. All these pieces are terms of
$\mathbf{B}^{(k+1)}$. Off $P$ they are compensated by Lemma~\ref{4.29:lem-far-field-paths} and by
the far-field distance bound, by which the default cubes lie at distance at least
$7LMR_{k+1}\ell_k$ from $\Omega_{k+1}$; they meet $P$ only at $j=k$.
\item[(3)] \emph{On $P$ at $j=k$.} At the cubes of scale $k$ meeting $P$ with
$\operatorname{dist}_\infty(\square,\Omega_{k+1}^c)\le 5M\ell_k$ and $\square_0\not\subset
\Omega_{k+1}$, on either side of $\partial\Omega_{k+1}$, what must be compensated is the count
$2^dL^4$ of scale-$k$ cover cubes in a $\pi_{k+1}$-cube and the constants of
\cite[(3.15)--(3.17)]{B-CMP109}, which depend on $L$, $M$, $\mathfrak p$ only and are bounded by
a constant $C(L,M,\mathfrak p)$ containing the factor $e^{16\kappa_1}$ of
\cite[(1.21)--(1.22)]{B-CMP116}; both are compensated source by source, according to the part of
the source field $B'$ that carries the piece:
\begin{enumerate}
\item[(a)] for the part of $B'$ carried by $A$ on $\Lambda_{k+1}$: with $\lambda_{k+1}(\omega)$
the length of a path $\omega$ inside $\Omega_{k+1}$ and
\begin{equation}\label{4.29:eq-near-field-scope-b}
|\omega|^+_{\mathbb B}:=|\omega|_{\mathbb B}+\frac{\lambda_{k+1}(\omega)}{M\ell_{k+1}},\qquad
|\omega|^+_{k+1}:=|\omega|_{k+1}+\frac{\lambda_{k+1}(\omega)}{M\ell_{k+1}}
\le|\omega|^+_{\mathbb B},
\end{equation}
the piece carries, along every path $\omega$ from $\Lambda_{k+1}$ to $\square$, the factor
\begin{equation*}
e^{-\mathsf{d}|\omega|^+_{\mathbb B}}\le e^{-\kappa'|\omega|^+_{k+1}}e^{-6R_{k+1}}L^{-4},
\end{equation*}
so that $L^4$ is compensated by $L^{-4}$, and $2^d$ and the constants by
$e^{-6R_{k+1}}\le\gamma^{12}$, by \eqref{4.29:eq-far-field-paths-a};
\item[(b)] for the part of $B'$ carried by $A_k$ on $\Omega_{k+1}\setminus\Lambda_{k+1}$, no
surplus decay is used, the per-cube factor of \cite[(3.15)--(3.17)]{B-CMP109} is at most
$3C_NB_3\delta_k^\sharp/\alpha_{2,k}$, and
\begin{equation}\label{4.29:eq-near-field-scope-a}
\begin{aligned}
\frac{3C_NB_3\delta_k^\sharp}{\alpha_{2,k}}&\le C^\sharp(\log x_k)^{-1},\\
C^\sharp&:=\frac{6L^2B_3^2C_NC_{\mathrm{II}}\max\{4a_\sharp A_0,\,C_1^{(3.43)}A_1\}}
{c_2\beta_s\min C_\cdot},
\end{aligned}
\end{equation}
with $C_N$ as in the hypotheses, $C_{\mathrm{II}}$ the absolute constant of
\cite[(1.20)]{B-CMP116}, $a_\sharp$ and $C_1^{(3.43)}$ those of
\eqref{4.29:eq-path-length-decay-b}, $\alpha_{2,k}$, $c_2$ and $\beta_s$ those of
\eqref{4.29:eq-scale-thresholds-b}, and $\min C_\cdot=\min\{C_0,C_1\}$ with $C_0$, $C_1$ of
\cite[(2.28)]{B-CMP119}; and
$2^dL^4C(L,M,\mathfrak p)C^\sharp(\log x_k)^{-1}\le 1$.
\end{enumerate}
No new floor on $L$ or $M$ and no condition on $k$ or $K$ arises.
\item[(4)] The analysis \cite[(3.18)]{B-CMP109} is entered exactly at the good cubes with
$s>j$ and at the cubes with $\square_0\subset\Omega_{k+1}$.
\end{enumerate}
\end{proposition}

\begin{proof}
\emph{Clause (1).} Let $\square_0\subset\Omega_{k+1}$. Since $\square\subset\square_0$, we have
$\square\subset\Omega_{k+1}$, and $\Omega_{k+1}\subset P$ by clause (b) of
Proposition~\ref{4.29:prop-native-cover}. By the same clause,
the cubes of $\mathcal{C}_j$ meeting $P$ are exactly the scale-$k$ cubes of \cite[\S3]{B-CMP119}
meeting $P$, every cube of $\mathcal{C}_j$ overlapping one of them is one of those cubes, and the
normalising sum of the partition is identically $1$ on them; since $\square$ meets $P$, $s=k$,
the cube and its overlaps are those of \cite[\S3]{B-CMP119}, and
$\bar\zeta_\square=\zeta_\square$. The cube is good: $\square\subset\Lambda^{(-9)}$.

Let $\square\subset Y\subset\Lambda_{k+1}$. By the distance bound between $\Lambda_{k+1}$ and
$(\Omega^\natural_{k+1})^c$ among the hypotheses of Proposition~\ref{4.29:prop-native-cover}, with
the equality $\Omega^\natural_{k+1}=\Omega_{k+1}$ stated there and $\ell_{k+1}=L\ell_k$,
\begin{equation*}
\operatorname{dist}_\infty\bigl(\Lambda_{k+1},\Omega^c_{k+1}\bigr)\ge 2LMR_{k+1}\ell_k>5M\ell_k,
\end{equation*}
so $\square_0=B_\infty(\square,5M\ell_k)\subset\Omega_{k+1}$ and $\square$ is such a cube. The
same inequality read contrapositively shows that $\square_0\not\subset\Omega_{k+1}$ forces
$\square\not\subset\Lambda_{k+1}$; the pieces of such a cube are then terms of
$\mathbf{B}^{(k+1)}$, as in \cite[\S3]{B-CMP119} and \cite[p.~279]{B-CMP119}.

Now let $X\subset\tilde\square^2$ at a cube with $\square_0\subset\Omega_{k+1}$. The cube has
level $k$ and the background level $k+1$, so the line of \cite[\S3]{B-CMP109} which reads
\emph{$j\le k$, hence $L^j\eta\le1$ and $(1+7\beta)L^{-2}\le1$} is available for every $j\le k$,
$j=k$ included, and $\zeta'_\square=\zeta_{\square_0}$ of \cite[(3.21)]{B-CMP109} is as there. To
prove that the local family is complete we use the chain
\begin{equation*}
\tilde\square^2\subset\Omega_{k+1}\subset\Omega^\natural_{j+1}\subset\Lambda_j^{\sim-1}.
\end{equation*}
The first inclusion holds because $\tilde\square^2\subset\square_0$. For the second: if $j=k$ it is
the equality $\Omega_{k+1}=\Omega^\natural_{k+1}$ among the hypotheses of
Proposition~\ref{4.29:prop-native-cover}; if $j<k$, the nesting relations there give
\begin{equation*}
\Omega_{k+1}=\Omega^\natural_{k+1}\subset\Lambda_k\subset\Lambda_{j+1}\subset\Omega^\natural_{j+1}.
\end{equation*}
For the third, $\Lambda_j^{\sim-1}$ denotes $\Lambda_j$ with one layer of $MR_j$-cubes removed
(the domain written $\Lambda_j^0$ in \cite[\S2]{B-CMP119}; it is not $\Lambda_j^\sim$), and the
distance bound between $\Omega^\natural_{j+1}$ and $\Lambda_j^c$, with the ratio bound, gives
\begin{equation*}
\operatorname{dist}_\infty\bigl(\Omega^\natural_{j+1},\Lambda_j^c\bigr)\ge 8MR_{j+1}\ell_{j+1}
=8LR_{j+1}\,M\ell_j\ge MR_j\ell_j,
\end{equation*}
while the nesting relations place $\Omega^\natural_{j+1}$ inside $\Lambda_j$. Hence the local family
$\{X\in \mathbf D_j: X\subset\tilde\square^2\}$ is complete. Further,
$\square_0\subset\Omega^\natural_{k+1}$ lies at distance at least
$8MR_{k+1}\ell_{k+1}=8LMR_{k+1}\ell_k$ inside $\Lambda_k$ by the distance bound between
$\Omega^\natural_{k+1}$ and $\Lambda_k^c$; $\Lambda_k\subset
\Lambda_j$ by the nesting relations; and, with $D_n$ the maximal domains of the determining set of
Lemma~\ref{4.29:lem-localisation-exact}, $\Lambda_j\subset D_j$ by
Lemma~\ref{4.29:lem-collar-staircase}, which places in $D_j$ every $M_1\ell_j$-cube of
$\Lambda_j^\sim$ at distance at least $2M_1\ell_j$ from the complement of $\Lambda_j^\sim$, the ring
$\Lambda_j^\sim\setminus\Lambda_j$ having width $M\ell_j\ge L^2M_1\ell_j$. Therefore the identity of
\cite[(3.18)]{B-CMP109} expressing $U_j(\exp iB_\square\bar U)$ as $U_j(\square_0,M^\cdot(\dots))$,
invoked in \cite[\S3]{B-CMP119}, is an instance of the reproduction of minimal configurations
among the hypotheses, as follows. Set
$U^*:=U_{j,\Lambda_j^\sim}(e^{iQ'(\eta\mathbf a)}W_0)$, completed by its level-$0$ data off $D_1$,
where $U_{j,\Lambda_j^\sim}$, $Q'$ and $W_0$ are those of Lemma~\ref{4.29:lem-localisation-exact}
and $\mathbf a=(t+t_\square\zeta_\square)\mathbf H_k(B')$, $t\in[0,1]$, as in that lemma. Then
$U^*$ is the minimal configuration for the determining set $\mathbf B_j(\Lambda_j^\sim)$ of
Lemma~\ref{4.29:lem-localisation-exact} at the data $e^{iQ'(\eta\mathbf a)}W_0$, and it is regular
at the levels of that set; this is the regularity under which
Lemma~\ref{4.29:lem-localisation-exact} applies the reproduction hypothesis. Apply the reproduction
hypothesis on the domain $\square_0$, with the determining set of \cite[(3.18)]{B-CMP109}, the cube
family $\{\square_n\}$ with top level $j$ on $\square_j\supset\tilde\square^4$, which is finer than
or equal to the pure level $j$ of $\mathbf{B}_j(\Lambda_j^\sim)$ about $\square_0$; the result is
the identity of \cite[(3.18)]{B-CMP109}. Polymers $X$
satisfying the first of the two conditions of \cite[(3.5)]{B-CMP109} are treated at these cubes as
at every cube, by Lemma~\ref{4.29:lem-collar-staircase}.

For $j<k$ let $\square$ be good of scale $s>j$. By clause (c) of
Proposition~\ref{4.29:prop-native-cover}, $\square_0\subset\Lambda_s$, and, since $s\ge j+1$,
the nesting relations give
\begin{equation*}
\square_0\subset\Lambda_s\subset\Lambda_{j+1}\subset\Omega^\natural_{j+1}.
\end{equation*}
By the distance bound between $\Omega^\natural_{j+1}$ and $\Lambda_j^c$, $\square_0$ lies at
distance at least $8MR_{j+1}\ell_{j+1}$ inside $\Lambda_j$,
and the chain
\begin{equation*}
\tilde\square^2\subset\square_0\subset\Omega^\natural_{j+1}\subset\Lambda_j^{\sim-1}
\end{equation*}
gives completeness as above; since $\square_0\subset\Lambda_j\subset D_j$, the identity of
\cite[(3.18)]{B-CMP109} is again an instance of the reproduction hypothesis. The gains (T3) of the
transfer lemma and the background of Proposition~\ref{4.29:prop-complex-hypothesis} apply at
these cubes.

The test $\square_0\subset\Omega_{k+1}$ involves $\Omega_{k+1}$ only on
$B_\infty(\square,5M\ell_k)$, which lies in $Y^\sim$; it is therefore a condition local to the
piece. In the whole-lattice setting it is always satisfied, so that the identification lemma
compares terms of the same form.

\emph{Clause (2).} Let $\square$ have scale $s=j$ and $\square_0\not\subset\Omega_{k+1}$, default
or not. For $j<k$ every scale-$j$ cube is of this kind, and none meets $P$, because by clause (b)
of Proposition~\ref{4.29:prop-native-cover} the cubes of $\mathcal C_j$ meeting $P$ are of scale
$k$. The terms of such a cube, $X\subset\tilde\square^2$ included, are treated by
\cite[(3.6)--(3.17)]{B-CMP109} alone: membership is \eqref{4.29:eq-complex-hypothesis-a}, the
fields are controlled by \eqref{4.29:eq-scale-thresholds-a} together with
\eqref{4.29:eq-interpolated-field-bounds-a}, and, since $s=j$, no factor $L^{j-s}$ is asked;
neither the window nor completeness is used. The chain of \cite[Lemma~4]{B-CMP109}, closed by the
factor $L^{-2}$ of the line \emph{$j\le k$, hence $L^j\eta\le1$ and $(1+7\beta)L^{-2}\le1$}, is
not entered within a layer; the treatment is that of clause (T4) of the transfer lemma. Such a
cube is not contained in $\Lambda_{k+1}$: for $j<k$ it does not meet $P\supset\Lambda_{k+1}$, and
for $j=k$ this is the converse in clause (1); so these pieces are terms of
$\mathbf B^{(k+1)}$, as in \cite[\S3]{B-CMP119} and \cite[p.~279]{B-CMP119}. Off $P$ they are
compensated by the far-field distance bound, by which the default cubes lie at distance at least
$7LMR_{k+1}\ell_k$ from $\Omega_{k+1}$, and by Lemma~\ref{4.29:lem-far-field-paths}, which is
stated on $\Lambda_j^\sim\setminus P$. On $P$ they occur only at $j=k$: they are the scale-$k$
cubes of \cite[\S3]{B-CMP119} with $\operatorname{dist}_\infty(\square,\Omega^c_{k+1})\le5M\ell_k$,
on either side of $\partial\Omega_{k+1}$. There the cube and the scale are those of
\cite[\S3]{B-CMP119}, so the factors $L^{4(k-s)}$ and $(1+k-s)^{1/2}$ of
\eqref{4.29:eq-far-field-paths-a} in Lemma~\ref{4.29:lem-far-field-paths} equal $1$ ($s=k$),
and nothing is due to them; membership is given by \eqref{4.29:eq-scale-thresholds-a}
on $P$ in its own units. What remains to be compensated there is treated in clause (3).

\emph{Clause (3).} Let $j=k$ and let $\square$ be a cube of scale $k$ meeting $P$ with
$\square_0\not\subset\Omega_{k+1}$ and $\operatorname{dist}_\infty(\square,\Omega^c_{k+1})\le
5M\ell_k$. By clause (2) its cube and scale are those of \cite[\S3]{B-CMP119}, and what must be
compensated at it is the count of \cite{B-CMP119} itself: $2^dL^4$, the number of scale-$k$ cover
cubes per $\pi_{k+1}$-cube, the $\pi_{k+1}$-cube being the unit of $d_{k+1}$ (level $k+1$ against
scale $k$); it is the analogue one level up of the count in (T5), which counts layer-$j$ cubes in
$Y_0\in\mathbf D_k$. Also to be compensated are the constants $E_0$, $B_0B_3$ and
$\sum_{X\cap\tilde\square^2\ne\emptyset}E_0e^{-\kappa d_k(X)}$ of
\cite[(3.15)--(3.17)]{B-CMP109}, which depend on $L$, $M$, $\mathfrak p$ only and are bounded by
$C(L,M,\mathfrak p)$.

The decay needs no compensation: the decay $e^{-\kappa'd_{k+1}(Y)}$, of the form of
\cite[(2.42)]{B-CMP119} used in this chapter with exponent $\kappa'$, is that of the tree inside
$Y$, assembled from three parts. First, $X$ carries
$e^{-\kappa d_k(X)}$. By \cite[\S1]{B-CMP109}, $d_j(X)$ is the length of a shortest tree graph
contained in $X$ and intersecting all the cubes of $X$, divided by $M$, in the units in which the
cubes of $\pi_j$ are unit cubes. Let $\bar X$ be the $\pi_{k+1}$-hull of $X$; the same tree,
measured in the coarser unit, gives $d_{k+1}(\bar X)\le d_k(X)/L$. Since $\kappa'=(1+4\beta)\kappa$
and $L\ge13>1+4\beta$,
\begin{equation*}
\kappa'd_{k+1}(\bar X)\le\frac{(1+4\beta)\kappa\,d_k(X)}{L}\le\kappa\,d_k(X).
\end{equation*}
Second, the walk from the source bond, which lies in $Y$, to $\square$ decays at the rate
$\mathsf d$ per local $M$-cube, which is at least $\kappa'$ per $M\ell_{k+1}$-cube on each
stratum it crosses (a stratum $\Gamma_a$, $a\le k$, of \eqref{4.29:eq-path-length-decay-a}, or
$\Omega_{k+1}$), since the local units $\ell_a$ do not exceed $\ell_{k+1}$ and $\mathsf d\ge\kappa'$
by the floor. Third, from $\square$ to $X$ the decay is the exponent
$\tfrac12\delta_0\operatorname{dist}^{(\xi)}(X,\square)$ of \cite[(3.17)]{B-CMP109}, with the
rate read as $\tfrac18\delta_0$ as in the kernel bound, that is $\delta_0/16$ per $\ell_k$-site;
this is at least $\kappa'/(LM)$ per site as soon as $\delta_0M\ge8\kappa'$ (using $L\ge2$), and
this holds because $\mathsf d=\delta_0M/8$ at the rate of \eqref{4.29:eq-path-length-decay-a},
so that the floor gives $\delta_0M\ge8(\kappa'+4\log L+1)$.

Hence only the count and the constants are to be compensated, and they are compensated source by
source, according to the part of $B'$ that carries the piece. The piece is linear in its direction,
and
\begin{equation*}
\mathbf H_k(B')=\int_0^1d\tau\,\Bigl\langle\frac{\delta\mathbf H_k}{\delta B'}(\tau B'),B'\Bigr\rangle
\end{equation*}
is additive in the last slot, the kernel bound holding on the polydisc; since
$\operatorname{supp}B'\subset\Omega_{k+1}$, the splitting
$B'=B'1_{\Lambda_{k+1}}+B'1_{\Omega_{k+1}\setminus\Lambda_{k+1}}$ splits the piece into a part
whose sources are carried by $A$ and a part whose sources are carried by $A_k$. In
\cite[(3.27), (3.29)]{B-CMP119} every piece of the class has both parts, and the Cauchy factor of
part (b) below compensates the whole.

\emph{Part (a).} The field $A$ lives on $\Lambda_{k+1}$; it is the whole of $B'$ in the last
logarithm of \cite[(3.28)]{B-CMP119}, where, as stated in \cite[p.~275]{B-CMP119}, the fluctuation
field $A$ is localised in $\Lambda^{(k)*}_{k+1}$. Since $\square_0=B_\infty(\square,5M\ell_k)\not
\subset\Omega_{k+1}$, there is a point of $\Omega^c_{k+1}$ within $5M\ell_k$ of $\square$, and
$\square$, a union of $2^d$ cubes of $\pi_k$ (\cite[\S3]{B-CMP119}), has side $2M\ell_k$; so every
point of $\square$ lies within $7M\ell_k$ of $\Omega^c_{k+1}$, and the distance bound between
$\Lambda_{k+1}$ and $(\Omega^\natural_{k+1})^c=\Omega^c_{k+1}$ gives
\begin{equation*}
\operatorname{dist}_\infty(\square,\Lambda_{k+1})\ge(2LR_{k+1}-7)M\ell_k.
\end{equation*}
The path lengths of \eqref{4.29:eq-near-field-scope-b} add to those of
\eqref{4.29:eq-path-length-decay-a} the stretch of $\omega$ inside $\Omega_{k+1}$:
$|\omega|^+_{k+1}$ is what $d_{k+1}(Y)$ counts, while $|\omega|_{\mathbb B}$, $|\omega|_{k+1}$ and
$d_{\mathbb B}$ of \eqref{4.29:eq-path-length-decay-a} keep the strata of levels $a\le k$, as
Lemma~\ref{4.29:lem-far-field-paths} requires; the inequality
$|\omega|_{k+1}\le L^{-1}|\omega|_{\mathbb B}$ used there is not claimed on $\Omega_{k+1}$. Let
$\omega$ be a path from $\Lambda_{k+1}$ to $\square$. Its length inside $\Omega_{k+1}$ counts
$1/(M\ell_{k+1})$ per unit in $|\omega|^+_{\mathbb B}$, and its length inside a stratum
$\Gamma_a$, $a\le k$, counts $1/(M\ell_a)\ge1/(M\ell_{k+1})$
per unit; hence, with $\ell_{k+1}=L\ell_k$ and $L\ge7$,
\begin{equation*}
|\omega|^+_{\mathbb B}\ge\frac{\text{length of }\omega}{M\ell_{k+1}}
\ge\frac{(2LR_{k+1}-7)M\ell_k}{LM\ell_k}=2R_{k+1}-\frac7L\ge2R_{k+1}-1.
\end{equation*}
The response $\mathbf H_k(B')$ decays on $\Omega_{k+1}$ as well: its determining set is of level $k$
there
(Lemma~\ref{4.29:lem-localisation-exact}), and by (T7) the distance is measured in the local
units, the members and sites being of level $k$; so the layered rate of the kernel bound per
$M\ell_{k+1}$-cube on that stretch is at least $\mathsf d$. (The local cubes there have side
$M\ell_k$, so the true rate per $M\ell_{k+1}$-cube is $L\mathsf d$; only $\mathsf d$ is used.)
Hence the piece carries $e^{-\mathsf d|\omega|^+_{\mathbb B}}$, and, writing
$R=R_{k+1}$ and using $|\omega|^+_{k+1}\le|\omega|^+_{\mathbb B}$ and $\mathsf d\ge\kappa'$,
\begin{equation*}
\mathsf d|\omega|^+_{\mathbb B}=\kappa'|\omega|^+_{\mathbb B}
+(\mathsf d-\kappa')|\omega|^+_{\mathbb B}
\ge\kappa'|\omega|^+_{k+1}+(\mathsf d-\kappa')(2R-1);
\end{equation*}
no factor $1/L$ enters on $\Omega_{k+1}$. The floor $\mathsf d\ge\kappa'+4\log L+1$ gives
\begin{equation*}
(\mathsf d-\kappa')(2R-1)\ge(4\log L+1)(2R-1)=6R+4\log L+\bigl(8(R-1)\log L-4R-1\bigr),
\end{equation*}
and $8(R-1)\log L\ge4R+1$ for $L\ge13$, $R\ge7$, since then $\log L>2$ and
$16(R-1)\ge4R+1$. Therefore
\begin{equation*}
e^{-\mathsf d|\omega|^+_{\mathbb B}}\le e^{-\kappa'|\omega|^+_{k+1}}e^{-6R_{k+1}}L^{-4}:
\end{equation*}
$L^4$ is compensated by $L^{-4}$, and $2^d$ and the constants by
$e^{-6R_{k+1}}\le\gamma^{12}$ by \eqref{4.29:eq-far-field-paths-a}.

\emph{Part (b).} The field $A_k$ lives on $\Omega_{k+1}\setminus\Lambda_{k+1}$; it enters at full
strength in the action of \cite[(3.27)]{B-CMP119}, as $tA_k$ in that of
\cite[(3.29)]{B-CMP119}, and in the integrands of \cite[(3.26), (3.28)]{B-CMP119}, through
termwise derivatives of the same localised expansion. As stated in \cite[\S3]{B-CMP119}, the
action in \cite[(3.29)]{B-CMP119} differs only by the additional dependence on the field $tA_k$,
which was already
taken into account there, and at this site the fluctuation field is localised in $\Omega_{k+1}$.
Here $\square$ may lie on $\operatorname{supp}A_k$ or next to it, with
$\operatorname{dist}_\infty(\square,\Omega^c_{k+1})\le5M\ell_k$ on either side of
$\partial\Omega_{k+1}$, and $d_{\mathbb B}(\square,\operatorname{supp}B')$ may vanish; no surplus
decay is claimed. The count $L^4$, the factor $2^d$ and the constants are compensated by the
per-cube factor of \cite[(3.15)--(3.17)]{B-CMP109}. There, \cite[p.~273]{B-CMP109}, the
$t_\square$-derivative is the Cauchy integral over $|t_\square|=r$, with the radius $r$ given by
\begin{equation*}
r\max\bigl\{|\delta H_j|_X,\ |P_1\delta H_j|_X,\ |\nabla^\xi_{\mathbf U}\delta H_j|_X,\
|\Delta^\xi_{\mathbf U}\delta H_j|_X\bigr\}=\tfrac13\alpha_2
\end{equation*}
in the notation of \cite[(3.14)]{B-CMP109}; so the piece is at most $3/\alpha_{2,k}$ times the
members of \cite[(3.14)]{B-CMP109} on $X$ of the direction field, times the bound of $\mathbf E$
of \cite{B-CMP109}, which, as stated in \cite[\S3]{B-CMP109}, follows only from the fact that
$H_j(B(t))$ satisfies \cite[(3.14)]{B-CMP109} with $\tfrac13\alpha_2$ and $\delta H_j$ satisfies
\cite[(3.9)]{B-CMP109}. The
mechanism is named in \cite[\S3]{B-CMP109}: the constant $E_0$ of the inductive assumption
\cite[(1.18)]{B-CMP109} is recovered by the differentiation with respect to $t_\square$ at
$t_\square=0$, which yields a factor $O(1)\alpha_2^{-1}\varepsilon_1$, made small by taking
$\varepsilon_1$ small; for the pieces of $\mathbf E_k$ it is the factor
$6\alpha_2^{-1}L^j\eta B_0B_3|\zeta_\square\mathbf H_k(B')|$ displayed in
\cite[(3.17), p.~273]{B-CMP109}.

The members in question are those of
$\delta H_j=\langle(\delta\mathcal H'/\delta B)(B(t)),\mathfrak b_\square\rangle$, with
$\mathfrak b_\square=DQ'\big|_{\eta t\mathbf H_k}[\eta\zeta_\square\mathbf H_k]$ and $\mathcal H'$ the
response on the determining set of Lemma~\ref{4.29:lem-localisation-exact}, as in
\eqref{4.29:eq-interpolated-field-bounds-a}; here $\bar\zeta_\square=\zeta_\square$ by clause (b)
of Proposition~\ref{4.29:prop-native-cover}. At level $s=k$ they are bounded by the kernel bound,
by the bound on $DQ'$ among the hypotheses and by the membership display of
\eqref{4.29:eq-scale-thresholds-a} read with
$\delta_k^\sharp$:
\begin{equation*}
C_{\mathfrak s}C_{\mathrm{ker}}\cdot2c_Q\cdot CB_3\delta_k^\sharp=C_NB_3\delta_k^\sharp,
\end{equation*}
with or without decay, all sources being of level $k$ on $P$ (the first case of
\eqref{4.29:eq-interpolated-field-bounds-a}); the constant $C$ already contains $C_{\mathrm{II}}$
once, so $C_{\mathrm{II}}$ enters $C^\sharp$ twice; this only enlarges the bound. Take
$\alpha_{2,k}=c_2\alpha_{3,k}=c_2\beta_s(2L^2B_3)^{-1}\min\alpha_{\cdot,k}$ of
\eqref{4.29:eq-scale-thresholds-b}, the $\alpha_2$ meant throughout, that of
\eqref{4.29:eq-interpolated-field-bounds-a} included; any smaller $\alpha_2$ is admissible in
\cite[(3.14)]{B-CMP109}. Then the factor is at most $3C_NB_3\delta_k^\sharp/\alpha_{2,k}$. With
$C_\flat(L):=6C_NL^2B_3/(c_2\beta_s)$, with $\delta_k^\sharp$ of
\eqref{4.29:eq-path-length-decay-b}, and with $\min\alpha_{\cdot,k}=(\min C_\cdot)g_k(\log x_k)^q$,
which is \cite[(2.28)]{B-CMP119} at $q_0=q_1=q$,
\begin{align*}
\frac{3C_NB_3\delta_k^\sharp}{\alpha_{2,k}}
&=\frac{C_\flat(L)B_3\delta_k^\sharp}{\min\alpha_{\cdot,k}}\\
&\le C_\flat(L)B_3C_{\mathrm{II}}\max\{4a_\sharp A_0,C_1^{(3.43)}A_1\}(\min C_\cdot)^{-1}
(\log x_k)^{p_0-q}\\
&\le C^\sharp(\log x_k)^{-1},
\end{align*}
the last step because $q\ge p_0+1$ and $\log x_k\ge\log\gamma^{-2}\ge1$ ($\gamma\le e^{-1/2}$),
and $C^\sharp=C_\flat(L)B_3C_{\mathrm{II}}\max\{4a_\sharp A_0,C_1^{(3.43)}A_1\}/\min C_\cdot$ is the
constant of \eqref{4.29:eq-near-field-scope-a}. This proves \eqref{4.29:eq-near-field-scope-a}.

The condition $q_0=q_1\ge p_0+1$ is a choice of the free exponents of \cite[(2.28)]{B-CMP119}:
there $q_0$, $q_1$ are integers greater than $1$ and $C_0$, $C_1$ sufficiently large positive
numbers, and $p_0$ is the positive integer of \cite[(2.4)]{B-CMP119}; the choice lies in that
range because $p_0+1\ge2$. It is one-sided, a floor on $q$ with no ceiling; it involves none of
$L$, $M$, $k$, $K$, and $\gamma$ is chosen after it. The lower bound on $q$ fixed, together with
$q_0=q_1$, in the setting of Theorem~0 implies $q\ge p_0+1$; no other condition on $q$ is used
here. The same mechanism is stated in
\cite[\S3]{B-CMP119} for the boundary terms $\mathbf B^{(k)}$ of the effective action, treated by
the same
\cite[(3.6), (3.7)]{B-CMP109}: the terms of the expansion can be estimated as in
\cite[(2.42)]{B-CMP119}, with the additional factor $O(\varepsilon_k)$ arising from the bound of
the fluctuation field; it is also the per-cube small factor of (T5).

Since $x_k\ge\gamma^{-2}$, $(\log x_k)^{-1}\le(\log\gamma^{-2})^{-1}$, and the ceiling
\eqref{4.29:eq-constants-order-a} gives
\begin{equation*}
2^dL^4C(L,M,\mathfrak p)C^\sharp(\log x_k)^{-1}\le
\frac{2^dL^4C(L,M,\mathfrak p)C^\sharp}{\log\gamma^{-2}}\le1;
\end{equation*}
this ceiling is one-sided and free of $k$, and $\gamma$ is chosen after $L$ and $M$. In
particular $3C_NB_3\delta_k^\sharp\le\alpha_{2,k}$: by \eqref{4.29:eq-near-field-scope-a} the
quotient $3C_NB_3\delta_k^\sharp/\alpha_{2,k}$ is at most $C^\sharp/\log\gamma^{-2}$, which is at
most $1$ by the ceiling because $2^dL^4C(L,M,\mathfrak p)\ge1$ (the constant $C(L,M,\mathfrak p)$
is an upper bound and may be taken at least $1$); so the membership condition of the hypotheses
is in force. The same
factor is present for the part carried by $A$, which is thus compensated twice. The far-field
statements, Lemma~\ref{4.29:lem-far-field-paths} and the far-field distance bound, by which the
default cubes lie at distance at least $7LMR_{k+1}\ell_k$ from $\Omega_{k+1}$, measure from
$\Omega_{k+1}\supset\operatorname{supp}B'$ and are unaffected. No new floor on $L$ or $M$ arises,
and nothing in $k$ or $K$.

\emph{Clause (4).} In the proof of clause (1), \cite[(3.18)]{B-CMP109} is entered at the cubes
with $\square_0\subset\Omega_{k+1}$ and, for $j<k$, at the good cubes of scale $s>j$; in the proof
of clause (2) it is not entered, the terms being treated by \cite[(3.6)--(3.17)]{B-CMP109} alone.
\end{proof}

\begin{remark}[Counterterm pieces and pieces inside the small-field region]
\label{4.29:rem-counterterm-pieces}
We use the notation of Lemma~\ref{4.29:lem-localisation-exact}: $\mathcal{C}_j$ is the cover of
generation $j$ with the partition of unity $\{\bar\zeta_\Box\}$, $U_{j,\Lambda_j^{\sim}}$ is the
minimiser attached to the determining set $\mathbf{B}'=\mathbf{B}_j(\Lambda_j^{\sim})$, and
$D_j$ is the maximal domain of $\mathbf{B}'$ on which it has level $j$. For a set $X$ and $r>0$,
$B_\infty(X,r)$ denotes the set of points at $\operatorname{dist}_\infty$-distance at most $r$
from $X$. Assume:
\begin{itemize}
\item[(1)] the exactness of the cube-by-cube localisation on the collared domain
(Lemma~\ref{4.29:lem-localisation-exact});
\item[(2)] clauses (a), (b), (f) and (h$'$) of the native-scale cover and its partition of unity
(Proposition~\ref{4.29:prop-native-cover});
\item[(3)] clauses (b) and (d) of the lemma identifying the label-local expansion with the
whole-lattice one;
\item[(4)] the near field and the cubes beside the fluctuation support
(Proposition~\ref{4.29:prop-near-field-scope}).
\end{itemize}
Then the following hold for every $j\le k$.
\begin{itemize}
\item[(i)] The counterterm carried by $\mathbf{E}^{(j)}$ is localised by the cover
$\mathcal{C}_j$, the partition $\{\bar\zeta_\Box\}$ and the configuration
$U_{j,\Lambda_j^{\sim}}$ exactly as the other terms of $\mathbf{E}^{(j)}$ are; its pieces read
$U_{j,\Lambda_j^{\sim}}$ at points of level $j$ only, and the cube-by-cube localisation of the
counterterm is exact.
\item[(ii)] The pieces with domain $Y\subset\Lambda_{k+1}$ coincide with those of
the expansion of \cite[\S3]{B-CMP119}.
\end{itemize}
\end{remark}

\begin{proof}
(i) The generation-$j$ part $\mathbf{E}^{(j)}$ of the effective action contains the counterterm
\begin{equation}\label{4.29:eq-counterterm-pieces-1}
-\beta_j(g_{j-1})\,A(\varphi_j,\cdot\,),
\end{equation}
where, in the notation of \cite[\S2]{B-CMP119}, $\varphi_j$ is a smooth function supported in
$\Lambda_j$ and equal to $1$ on $\Lambda_j^{\sim-1}$, the region $\Lambda_j$ with one layer of
$MR_j$-cubes removed, and $A(\varphi_j,\cdot\,)$ is the weighted Wilson action of
Definition~\ref{1.1:def-wilson-action}.
In \cite[(4.38)]{B-CMP109} the difference of this term at the two configurations compared in
\eqref{4.29:eq-localisation-exact-a} is localised
``by the partition of unity connected with the partition $\pi_j$, and we apply to it the whole
procedure of these two sections'' (here $\pi_j$ is the partition of the lattice into the cubes
of side $M\ell_j$ of \cite[\S1]{B-CMP109}). The same procedure applies to it here. Indeed,
$\operatorname{supp}\varphi_j\subset\Lambda_j$, and $\Lambda_j\subset D_j$; hence every piece of
the counterterm depends on $U_{j,\Lambda_j^{\sim}}$ only at points of $D_j$, where
$\mathbf{B}'$ has level $j$. Moreover $D_j$ is one of the maximal domains of $\mathbf{B}'$, so
$D_j\subset\Lambda_j^{\sim}$. By clause (a) of
Proposition~\ref{4.29:prop-native-cover}, $\sum_{\Box\in\mathcal{C}_j}\bar\zeta_\Box=1$ on the
$\tfrac13 M\ell_j$-neighbourhood of $\Lambda_j^{\sim}$, and in particular on $\Lambda_j$. The
exactness argument of Lemma~\ref{4.29:lem-localisation-exact} uses of the partition only this
identity and the linearity of each insertion in $\bar\zeta_\Box$; both are available for the
counterterm, whose pieces, like the other terms treated there, are gauge invariant (a gauge
transformation conjugates each $U(\partial p)$ and so leaves each plaquette term
$1-\operatorname{Re}\operatorname{tr}U(\partial p)$ unchanged) and depend on
$U_{j,\Lambda_j^{\sim}}$ on $\Lambda_j^{\sim}$ only. Hence the argument of
that lemma applies to the counterterm verbatim: the sum over $\Box\in\mathcal{C}_j$ of the
$t_\Box$-derivatives at $t_\Box=0$ equals the derivative in the global parameter, and the
localisation of the counterterm is exact.

(ii) Let $Y\subset\Lambda_{k+1}$. Near $Y$ the cover is forced by the two separation
conditions between the small-field regions of successive levels; this is the content of
clauses (b) and (f) of Proposition~\ref{4.29:prop-native-cover}, which we now apply.
By clause (b), the near-field region of Proposition~\ref{4.29:prop-native-cover} contains
the region $\Omega_{k+1}$ on which the fluctuation field of the step is localised, and
$\Omega_{k+1}\supset\Lambda_{k+1}$; the cubes of $\mathcal{C}_j$ meeting the near-field region
are exactly Ba{\l}aban's scale-$k$ cubes meeting it, and every cube of $\mathcal{C}_j$ overlapping
one of them is Ba{\l}aban's, with $\bar\zeta_\Box=\zeta_\Box$. By clause (f), every
cube whose scale, partition function or overlaps differ from Ba{\l}aban's has
$\operatorname{dist}_\infty(\Box,\Omega_{k+1})\ge 7LMR_{k+1}\ell_k$, and its pieces have domains
not contained in $\Lambda_{k+1}$. Consequently no
cube of the latter kind contributes to a piece with domain $Y$.

Next, $\mathbf{B}'$ has pure level $j$ on $D_j$. Moreover,
\begin{equation*}
D_j\supset B_\infty\bigl(\Lambda_{k+1},\,8LMR_{k+1}\ell_k\bigr)
\end{equation*}
by the separation conditions between the small-field regions of successive
levels, recalled in Proposition~\ref{4.29:prop-near-field-scope}. On this set the whole-lattice
determining set of generation $j$ has pure level $j$ as well. The random-walk expansions
and the operators decoupled in the cluster-expansion parameters that enter a piece with
domain $Y\subset\Lambda_{k+1}$ read nothing but this pure-level-$j$ structure, on which
$\mathbf{B}'$ and the whole-lattice determining set of generation $j$ agree, by clause (h$'$) of
Proposition~\ref{4.29:prop-native-cover}, clauses (b) and (d) of the lemma identifying the
label-local expansion with the whole-lattice one, and
Proposition~\ref{4.29:prop-near-field-scope}. Hence each piece with $Y\subset\Lambda_{k+1}$ is
built from the same cubes, the same partition functions and the same level-$j$ data as in
the expansion of \cite[\S3]{B-CMP119}, and is therefore unchanged.
\end{proof}

\begin{corollary}[Label dependence of boundary pieces born at a cover cube]
\label{4.29:cor-boundary-label}
Let $j\le k$. Dependence on the label is understood as in clause (h) of
Proposition~\ref{4.29:prop-native-cover}. Assume the following statements.
\begin{enumerate}
\item[(1)] Clauses (h), (h$'$) and (h$''$) of Proposition~\ref{4.29:prop-native-cover}.
\item[(2)] Independence of the localised terms from the small-field region
(Corollary~\ref{4.4:cor-localised-independence}), in the form: for every level $i$, the localised
functions of $(X,z)$ at level $i$ are the functions of the whole lattice and do not depend on
$\Lambda_i$.
\item[(3)] The statement that, at every level $i$, an old term of $(X,z')$ with
$z'\notin\Lambda_i^0$, filed as a boundary term, is the same function as the corresponding new
term of $\mathbf{E}$, together with the statement giving the collar of pieces lifted to later
levels.
\item[(4)] The clause of the remark on the localised terms $V(Y,U_k)$ that governs their
dependence on the large-field data of $Z$, and whose hypotheses include a collar of $Z$ at the
index $k''$ of $Z$ itself, with $k''$ and $Z$ as in (iv) below.
\item[(5)] The statement fixing how the region $Z$ of the summation over localisation domains in
\cite{B-CMP122b} is to be read (as a union of components of $\Lambda_{k''}^c$, with $k''$ as in
(iv) below).
\item[(6)] The statement that healing deletes whole components of the large-field region.
\item[(7)] The statements that the cut \cite[(1.49)]{B-CMP122b} uses only the functions of
$(X,z)$ and their total, and that the regrouped term $\mathbf{R}'$ is label-free given
clauses (h)--(h$''$) of Proposition~\ref{4.29:prop-native-cover}.
\end{enumerate}
Let $Y$ be a localisation domain and let $\mathbf{B}^{(k+1)}(Y)$ be a boundary piece born at a
cube $\square\subset Y$ of the cover of Proposition~\ref{4.29:prop-native-cover}. With
$\pi_{k+1}$ the partition into cubes of the lattice of level $k+1$, put
\begin{equation}\label{4.29:eq-boundary-label-a}
Y^{\sim}:=Y\ \text{together with one layer of}\ \pi_{k+1}\text{-cubes}.
\end{equation}
Then the following hold.
\begin{enumerate}
\item[(i)] Through the cover, $\mathbf{B}^{(k+1)}(Y)$ reads the label only on
$B_\infty(\square,9M\ell_k)$, and $B_\infty(\square,9M\ell_k)\subset Y^{\sim}$.
\item[(ii)] Through the restricted minimisers $\mathbf{B}_j(\Lambda_j^{\sim})$ it reads the label
only within distance $5M\ell_k$ of $Y$, and hence inside $Y^{\sim}$, the set of (i).
\item[(iii)] Through Ba{\l}aban's own construction, namely the ranges $z\in\Lambda_j^0$ and the
cut-offs $\varphi_j$, it reads the label only within distance $2LMR_{k+1}\ell_k$ of $Y$.
\item[(iv)] Let $k''\ge k+1$ and let $Z$ be a union of components of $\Lambda_{k''}^c$, as in
\cite{B-CMP122a}. For families $\mathfrak{L}$ and $\mathfrak{F}$ of components of
$\Lambda_{k''}^c$, write $\Phi(X;\mathfrak{L})$ for a function of $(X,z)$ in the new action
localised in $X$ and computed with the label data carried by $\mathfrak{L}$, and write
$Z(\mathfrak{F})$ for the union of the components in $\mathfrak{F}$. For label data of index
$<k''$,
\begin{equation}\label{4.29:eq-boundary-label-1}
\Phi(X;\mathfrak{L}\cup\mathfrak{F})=\Phi(X;\mathfrak{L})
\quad\text{whenever}\quad X\cap Z(\mathfrak{F})=\emptyset .
\end{equation}
At the index $k''$ of $Z$ itself the functions of $(X,z)$ are unchanged, and only the filing
between $\mathbf{E}$ and $\mathbf{B}$ of those $(X,z)$ with $z$ within one $MR_{k''}$-layer of
$Z$ (the filing by $z\in\Lambda^0_{k''}$) changes.
\item[(v)] Compared with the step performed with the cover and the partition of unity of
\cite[p.~276]{B-CMP119}, the parts $\mathbf{E}^{(k+1)}$ and
$\mathbf{R}''^{(k+1)}$ of the effective density produced by the step, the coefficient
$\beta_{k+1}$ and the quantity $E_0$ of the effective density are the same. The part
$\mathbf{R}'$ of that density differs only by a regrouping of the same density, so its total is
the same; and it is label-free.
\end{enumerate}
\end{corollary}

\begin{proof}
\emph{The cover.} The piece $\mathbf{B}^{(k+1)}(Y)$ is born at $\square\subset Y$.
Clause (h) of Proposition~\ref{4.29:prop-native-cover} states that membership of $\square$ in the
cover, its scale, its default status and the function $\bar\zeta_\square$ near
$\square\cap\Lambda_j^{\sim}$ depend on the label only through the traces of $\Lambda_j$ and of
$\{\Lambda_{s'}\}_{s'\le k}$ on $B_\infty(\square,9M\ell_k)$. The same clause gives
$9M\ell_k<M\ell_{k+1}$; since $\square\subset Y$, the ball $B_\infty(\square,9M\ell_k)$ lies
within distance $9M\ell_k<M\ell_{k+1}$ of $Y$, and so, the cubes of $\pi_{k+1}$ having side at
least $M\ell_{k+1}$, inside $Y^{\sim}$. This is (i).

\emph{The restricted minimisers.} By clause (h$'$) of
Proposition~\ref{4.29:prop-native-cover}, the trace of $\mathbf{B}_j(\Lambda_j^{\sim})$ on a set
$D\subset\Lambda_j^{\sim}$ depends on the label only through $\Lambda_j$ on
$B_\infty(D,3M\ell_j)$. The same clause shows that the localised pieces born at
$\square\subset Y$ read $\mathbf{B}_j(\Lambda_j^{\sim})$ only inside a set contained in
$B_\infty(Y,2M\ell_k)$. Since $j\le k$, the total reach is $2M\ell_k+3M\ell_j\le5M\ell_k$.
Because $5M\ell_k<9M\ell_k<M\ell_{k+1}$, this route of dependence lies in $Y^{\sim}$, as in (i).
This is (ii).

\emph{Ba{\l}aban's own construction.} Clause (h$''$) of Proposition~\ref{4.29:prop-native-cover}
gives, for $j\le k$,
\begin{equation*}
2MR_j\ell_j\le 2MR_{k+1}\ell_{k+1}=2LMR_{k+1}\ell_k .
\end{equation*}
Here $2MR_j\ell_j$ is the distance within which the ranges $z\in\Lambda_j^0$ and the cut-offs
$\varphi_j$ read $\Lambda_j$. This is (iii).

\emph{Where the dependence is used.} The dependence in (i)--(iii) enters through three
statements of \cite{B-CMP122b} and through hypothesis (4).

First, \cite{B-CMP122b} states that ``the summation is over domains $Y\in\mathbf{D}_k$ satisfying
the property that the intersection $Y\cap Z^{\sim}$ is a union of components of $Z^{\sim}$''.
Here $Z$ is read as in hypothesis (5), and the superscript $\sim$ on $Z$ is that of
\cite{B-CMP122b}: the enlargement of $Z$ by one layer of cubes of side $MR_{k''}\ell_{k''}$ at the
index $k''$ of $Z$; for $k''=k+1$ this is the reading of clause (h$''$) of
Proposition~\ref{4.29:prop-native-cover}. It is to be
distinguished from the enlargement~\eqref{4.29:eq-boundary-label-a} of localisation domains.

Second, \cite{B-CMP122b} states that ``the terms $V(Y,U_k)$ depend on the sequence
$\{\Omega_j,Z_j\}$ restricted to the components of $Z$ contained in $Y$''.

Third, there is the clause of the remark on the terms $V(Y,U_k)$ in hypothesis (4).

Finally, there is the cut \cite[(1.49)]{B-CMP122b}, which is made by $Z$ itself, without a
collar: ``the second part includes the remaining terms, i.e., the terms with localization domains
disjoint with $Z$ $\ldots$ we obtain all the terms of the new action with localization domains
disjoint with $Z$''. What this cut uses is exactly the identity~\eqref{4.29:eq-boundary-label-1}.

\emph{Label data of index $<k''$.} The set $Z$ is a union of components of
$\Lambda_{k''}^c$, and it remains one after healing, since by hypothesis (6) healing deletes
whole components; so $\mathfrak{F}$ ranges over families of components $C$ of
$\Lambda_{k''}^c$, and $X\cap Z(\mathfrak{F})=\emptyset$ says, by the definition of
$Z(\mathfrak{F})$, that $X\cap C=\emptyset$ for each $C\in\mathfrak{F}$. For such $X$ and $C$,
clause (h$''$) of Proposition~\ref{4.29:prop-native-cover} states that a piece localised in $X$
reads nothing of $C$ of index $\le k$.

In the case $k''=k+1$ these are all the indices $<k''$. For $k''\ge k+2$, the index $k+1$ is the
case of clause (h$''$) in which the reading of $C$ at index $k+1$ is void. A piece born at
step $k+1$ reads no label data of index greater than $k+1$; hence the indices $\le k$ and $k+1$
exhaust the data of index $<k''$. So all label data of index $<k''$ are covered by
clause (h$''$), and~\eqref{4.29:eq-boundary-label-1} holds for them.

\emph{The index $k''$ of $Z$ itself.} By hypothesis (2), applied at the level $i=k''$, and
\cite[\S3]{B-CMP119}, the functions of $(X,z)$ are
unchanged. What changes is only the filing between $\mathbf{E}$ and $\mathbf{B}$ of those $(X,z)$
with $z$ within one $MR_{k''}$-layer of $Z$; this is \cite[\S3, following (3.47)]{B-CMP119}
together with hypothesis (3).
This proves (iv).

\emph{No statement that uses the dependence is affected.}
\begin{itemize}
\item By hypothesis (7), the cut \cite[(1.49)]{B-CMP122b} uses the functions of $(X,z)$ and their
total, and neither of these changes.
\item The domains of the terms $V(Y,U_k)$ meet $Z$, so these terms are boundary terms whichever
filing is used.
\item The clause of the remark in hypothesis (4) carries the collar at the index $k''$ of $Z$
itself in its own hypothesis.
\end{itemize}

\emph{Later steps and lifts.} Later applications of the operation $\mathbf{T}$ inherit the radii of
(i)--(iii). At each application the radius is measured from the localisation domain containing the
piece, so there is no sum over steps. For lifts of pieces the collar is the one given by the
collar statement for lifted pieces in hypothesis (3).

\emph{Pieces with $Y\subset\Lambda_{k+1}$.} By Proposition~\ref{4.29:prop-near-field-scope},
for these pieces the cube, its overlaps and the partition-of-unity function
$\bar\zeta_\square=\zeta_\square$ are those of Ba{\l}aban's
construction, and there is nothing to add, since there $\mathbf{B}_j(\Lambda_j^{\sim})$ is purely
of level $j$. Pieces born at cubes $\square\not\subset\Lambda_{k+1}$ have
$Y\not\subset\Lambda_{k+1}$, since $\square\subset Y$, and are therefore terms of
$\mathbf{B}^{(k+1)}$ by the filing of \cite[\S3, following (3.47)]{B-CMP119} (compare
Proposition~\ref{4.29:prop-near-field-scope}); their dependence on the label is the one
treated in (i)--(iv).

\emph{Conclusion.} Hence, compared with the step performed with the cover and the partition of
unity of \cite[p.~276]{B-CMP119}, the quantities $\mathbf{E}^{(k+1)}$, $\mathbf{R}''^{(k+1)}$,
$\beta_{k+1}$ and $E_0$ are the same. The term $\mathbf{R}'$ is only regrouped: its totals in the
grouping with the cover of Proposition~\ref{4.29:prop-native-cover} and in the grouping of
\cite[p.~276]{B-CMP119} are equal, because both groupings are exact rewritings of one density.
By hypothesis (7), $\mathbf{R}'$ is label-free given
clauses (h)--(h$''$) of Proposition~\ref{4.29:prop-native-cover}. This is (v).
\end{proof}

\begin{remark}[The sharpened locality clause for boundary terms]\label{4.29:rem-sharpened-locality}
Throughout, $j$ is an index of the inductive construction. In the step from level $k$ to level
$k+1$ the statements below are used at $j=k+1$. The radii $R_j$, $\check R_j$ and the unit
$\ell_j$ are those of Proposition~\ref{4.29:prop-native-cover}. For a set $Y$ and $r>0$,
$B_\infty(Y,r)$ is the set of points at $\sup$-norm distance at most $r$ from $Y$. For a component
$C$ of $\Lambda_j^{c}$, $C^{\sim}$ is $C$ enlarged by one layer of width $MR_j\ell_j$, as in
Ba{\l}aban's construction \cite{B-CMP119}. We write $\mathfrak{l}$ for the label, the large-field
datum at which the terms of the step are evaluated. Its entries of index $i$ include the sets
$\Lambda_i$, $\Omega_i$ and Ba{\l}aban's large-field set $S_i$ of \cite[(2.40)--(2.42)]{B-CMP119}.
We write $\sigma$ for the remaining argument of a term in the inductive statement of
hypothesis~(b) below. The localisation domain $D$ of a term is $D=X$ for $\mathbf{B}^{(j)}(X)$,
and $D=X^{\sim}$, the enlarged domain of \cite[(1.71)]{B-CMP122b} in the definition of the operation
$\mathbf{T}'_j$, for $\mathbf{T}'_j(X)$. The index $j$ of a term is called its own index.
The clause sharpened here is the label-locality clause recalled in~(b): clause~(1) below
replaces its wording, the trace of the label on $X$, by a statement that names the ball on which
the label is read and, separately for the data of index below $j$ and for the own index $j$, the
components of $\Lambda_j^c$ that enter.

Assume the following.
\begin{enumerate}
\item[(a)] Proposition~\ref{4.29:prop-native-cover}, and in particular its clause (h$''$). Let
$k''\ge k+1$, let $C$ be a component of $\Lambda^c_{k''}$, and let $Y$ be a localisation domain
with $C\cap Y=\emptyset$. Then a term of the step read on the localisation domain $Y$, namely a
piece $\mathbf{B}^{(k+1)}(Y)$, or $\mathbf{T}'_{k+1}(X)$ with $Y=X^{\sim}$, depends on no datum of
index at most $k$ carried by $C$, and it does not depend on $\Omega^c_{k+1}\cap C$. Of index $k+1$ it
depends on $\Lambda_{k+1}$ and $S_{k+1}$ only on $Y$ and, through the assignment of the pairs
with $z\in\Lambda^0_{k+1}$ described in~(d), on the collar of width $MR_{k+1}\ell_{k+1}$ about
$Y$; hence it depends on $C$ only if $C^\sim\cap Y\ne\emptyset$.
\item[(b)] The locality clauses of the inductive statement of this chapter on the terms of
index $j$: clause (iii), for the terms $\mathbf{B}^{(j)}(X)$ and $\mathbf{T}'_j(X)$, and clause
(ii), for the terms of the operation $\mathbf{R}'$ of the inductive statement. Both read as
follows: the term depends on
$(\sigma,\mathfrak{l})$ only through the trace of $\mathfrak{l}$ on $X$; a term with this
property is called label-local. A further clause of the same inductive statement allots a collar
of radius $60M\check R_j\ell_j$ to the re-expression of a term of index $j$ at a later index (its
lift).
\item[(c)] Ba{\l}aban's cut-off functions satisfy $\varphi_i\in C_0^\infty(\Lambda_i)$ and
$\varphi_i=1$ on $\Lambda_i^{\sim-1}$ \cite{B-CMP119}, where $\Lambda_i^{\sim-1}\subset\Lambda_i$ is
the set of his construction obtained by removing from $\Lambda_i$ one layer along
$\Lambda_i^c$, so that $\Lambda_i\setminus\Lambda_i^{\sim-1}\subset\bigcup_C(C^\sim\setminus C)$,
$C$ running over the components of $\Lambda_i^c$. In his step, at $j=k+1$, the term
$\beta_{k+1}(g_k)A(\varphi_{k+1},U_{k+1})$ is subtracted from the difference
$E^{(k+1)}-E^{(k+1)}(1)$ of the effective action $E^{(k+1)}$ and its value $E^{(k+1)}(1)$ at the
configuration $1$, and the
function $g_k^{-2}(\cdot)$ is replaced by $g_{k+1}^{-2}(\cdot)$ \cite{B-CMP119}. Moreover, with
the weighted forms of the Wilson action of Definition~\ref{1.1:def-wilson-action}, the following
identity of the inductive statement of this chapter and of the lemma of this chapter on the
coupling weights holds plaquette by plaquette:
\begin{equation}\label{4.29:eq-sharpened-locality-1}
-A\bigl(g_j^{-2}(\cdot),U\bigr)-\sum_{i\le j}\beta_i\,A(\varphi_i,U)=-g_0^{-2}A(U).
\end{equation}
\item[(d)] At the own index $j$, Ba{\l}aban's step attaches to each localisation domain $X$ and
each point $z$ of $\Lambda_j^0$, the native range of the points $z$ in his step, a function of
the configuration, and files it
among the small-field terms $\mathbf{E}$ or among the boundary terms $\mathbf{B}$. These
functions are those of Ba{\l}aban's step; only their assignment to the terms $\mathbf{E}$ or to
the terms $\mathbf{B}$ depends on $\Lambda_j$ outside the localisation domain $D$. Each of the
following own-index dependences involves only components $C$ of $\Lambda_j^c$ with
$C^\sim\cap D\ne\emptyset$:
\begin{itemize}
\item the assignment of the pairs with $z\in\Lambda_j^0$ \cite{B-CMP119};
\item the dependence through the enlargement $Z_j^{\sim}$ of the large-field region $Z_j$
\cite{B-CMP122b};
\item the dependence through touching regions, in the sense of clause (u1) of the remark
in~(f);
\item the terms of $\mathbf{R}''$ at the rim, those with $X\subset\Lambda_j^{\sim-1}$, in the form
of the statement of this chapter on that operation, the points $z$ running over the range fixed
among the hypotheses of the remark in~(f).
\end{itemize}
\item[(e)] The radii satisfy $R_j\ell_j<2\check R_j\ell_j$ and $R_i\ell_i\le R_{i+1}\ell_{i+1}$,
as fixed in the display of this chapter comparing the radii.
The large-field regions also have the following separation property, which is among the
hypotheses of Proposition~\ref{4.29:prop-native-cover}. For every index $k''$, every component
$C$ of $\Lambda^c_{k''}$ and every $i\le k''-1$, when $C$ is created, $\Lambda_i^c\cap C$ lies at
$\sup$-norm distance at least $8MR_{i+1}\ell_{i+1}$ inside $C$.
\item[(f)] The remark on the dependence of the polymer terms on the large-field data. Its
clause (u1) concerns touching regions. Its clause (u2) states that the terms are label-local,
in the sense of~(b), by induction. Its clause (u3) concerns terms of an own index, which we
write $k''$ (the index written $k$ in that remark); its hypothesis is that $X$ is disjoint from
$Z_{k''}^{\sim}\cup\bigcup_i Y_i^{\sim}$, where $Z_{k''}$ is the large-field region and the $Y_i$
are the domains of that remark, and ${}^{\sim}$ denotes their enlargements by one layer.
\end{enumerate}
Then the following hold.
\begin{enumerate}
\item[(1)] With $D$ the localisation domain fixed above, each
term $\mathbf{B}^{(j)}(X)$ and $\mathbf{T}'_j(X)$ depends on $(\sigma,\mathfrak{l})$ only through
the trace of $\mathfrak{l}$ on $B_\infty(D,2M\check R_j\ell_j)$. More precisely:
\begin{itemize}
\item of the data of index $<j$ and of $\Omega_j$, only what lies in components of
$\Lambda_j^c$ meeting $D$ enters;
\item $\Lambda_j$ and $S_j$ enter on $D$;
\item through the own-index dependences of~(d), $\Lambda_j$ also enters from the components $C$
of $\Lambda_j^c$ with $C^\sim\cap D\ne\emptyset$.
\end{itemize}
The radius $2M\check R_j\ell_j$ is smaller than $60M\check R_j\ell_j$. Clause (1) is stable
under the induction: no collar accumulates over the steps.
\item[(2)] Clause (u2) holds with \emph{local in the sense of} (1), by induction, in place of
label-local. Clause (u3) holds as stated.
\item[(3)] Clause (ii) of~(b), for the terms of $\mathbf{R}'$, holds in the wording of~(1).
\end{enumerate}
\end{remark}

\begin{proof}
We prove (1) first, in the order of its clauses.

\emph{Data of index $<j$ and $\Omega_j$.} Take $k''=j=k+1$ in clause (h$''$) of
Proposition~\ref{4.29:prop-native-cover}, quoted in~(a). We apply it with $Y=D$: for
$\mathbf{B}^{(j)}(X)$ this is $Y=X$, and for $\mathbf{T}'_j(X)$ it is $Y=X^{\sim}$, itself a
localisation domain at level $j$. Then a term with localisation domain
$D$ depends on no datum of index $\le k=j-1$ carried by a component $C$ of $\Lambda_j^c$ with
$C\cap D=\emptyset$, and it does not depend on $\Omega_j$ there. Hence only the components
meeting $D$ enter.

\emph{The set $\Lambda_j$.} By the own-index clause of~(a), applied with $Y=D$ and $k+1=j$, the
sets $\Lambda_j$ and $S_j$ enter on $D$. Beyond $D$, $\Lambda_j$
enters only through the own-index dependences listed in~(d). By~(d), each of them involves only
components $C$ with $C^\sim\cap D\ne\emptyset$. What is read there lies in one layer of width
$MR_j\ell_j$ about $D$. By~(e) we have $MR_j\ell_j<2M\check R_j\ell_j$. This layer therefore
lies inside $B_\infty(D,2M\check R_j\ell_j)$, the radius of the first assertion of~(1). Since
$2M\check R_j\ell_j<60M\check R_j\ell_j$, this radius stays inside the collar of radius
$60M\check R_j\ell_j$ used for lifts in~(b).

\emph{The cut-off functions.} By~(c), $\varphi_j$ is supported in $\Lambda_j$ and equals $1$ on
$\Lambda_j^{\sim-1}$, and $\Lambda_j\setminus\Lambda_j^{\sim-1}\subset\bigcup_C(C^\sim\setminus C)$.
Hence
\begin{equation*}
\{\varphi_j\ne1\}\cap\Lambda_j\subset\bigcup_C\bigl(C^\sim\setminus C\bigr),
\end{equation*}
where $C$ runs over the components of $\Lambda_j^c$. The value of $\varphi_j$ therefore enters
only on $C^\sim\setminus C$. It enters the terms $\mathbf{E}$ and not the terms $\mathbf{B}$.
Indeed, at $j=k+1$ the subtraction of $\beta_{k+1}(g_k)A(\varphi_{k+1},U_{k+1})$ and the
replacement of $g_k^{-2}(\cdot)$ by $g_{k+1}^{-2}(\cdot)$ in~(c) are made in
$E^{(k+1)}-E^{(k+1)}(1)$.

Add and subtract $\sum_{i\le j}\beta_iA(\varphi_i,U)$, and combine the subtracted copy with the
weighted action term $-A(g_j^{-2}(\cdot),U)$. By \eqref{4.29:eq-sharpened-locality-1}, this pair
equals $-g_0^{-2}A(U)$ plaquette by plaquette. The right-hand side contains no cut-off function,
so the pair so grouped is free of the cut-off functions. The total attached to each $X$ is therefore the same
whether the terms carrying the cut-off functions are kept with this pair or counted separately.
Moreover, the dependence on $\varphi_j$ lives in components $C$ with
$C^\sim\cap D\ne\emptyset$, which is the criterion of the third item of~(1). Hence~(1) remains
true. Since $\varphi_{k+1}$ enters the terms $\mathbf{E}$ only, a piece $\mathbf{B}^{(k+1)}(Y)$
produced by the operation $\mathbf{T}$ does not depend on $\varphi_{k+1}$.

\emph{Stability under the induction.} By the second inequality of~(e), $R_i\ell_i\le
R_{i+1}\ell_{i+1}$, and therefore
\begin{equation*}
MR_i\ell_i<8MR_{i+1}\ell_{i+1}.
\end{equation*}
Let $C$ be a component created at a later index $k''>i$. By the separation property in~(e),
$\Lambda_i^c\cap C$ lies at distance at least $8MR_{i+1}\ell_{i+1}$ inside $C$. The own-index
layer at index $i$ has width $MR_i\ell_i$, which is smaller than this distance. So that layer
is contained in the component at every later index. At that index, the first item of~(1) covers
it. Consequently no collar accumulates over the steps.

\emph{Proof of (2).} The induction in clause (u2) runs with clause~(1) at each index. By the
preceding paragraph, (1) is stable under that induction, so (u2) holds with ``local in the
sense of~(1)''.

Clause (u3) keeps its statement; we now give its justification. Let $k''$ be the own index of
the terms in (u3), as in~(f), so that the hypothesis of (u3) reads
$X\cap\bigl(Z_{k''}^{\sim}\cup\bigcup_iY_i^{\sim}\bigr)=\emptyset$. For the indices $<k''$ we use
clause (h$''$) of Proposition~\ref{4.29:prop-native-cover}. By it, no polymer contained in $X$
depends on a component of the large-field data of (u3) that $X$ does not meet, whatever the
distance. For the own index $k''$ we use this disjointness. At the own index the terms depend on
$Z_{k''}$ only through the own-index dependences of~(d), which involve only components $C$ with
$C^{\sim}\cap X\ne\emptyset$; the hypothesis $X\cap Z_{k''}^{\sim}=\emptyset$ excludes every such
component. The enlargement $Z_{k''}^{\sim}$
contains the collar of width $MR_{k''}\ell_{k''}$. A
separation of order $MR_{k''}\ell_{k''}$ obtained from distance alone would not be enough here,
because $MR_{k''}\ell_{k''}<2M\check R_{k''}\ell_{k''}$, the radius of~(1). It is therefore the
disjointness from the collar that is used.

\emph{Proof of (3).} For the terms of $\mathbf{R}'$, clause (ii) of~(b) is read in the same
words as~(1), and the steps above carry over to it.
\end{proof}

\begin{remark}[The inputs to the transfer of the small-field analysis]
\label{4.29:rem-transfer-inputs}
Fix levels $j\le k$. We use the notation of Lemma~\ref{4.29:lem-localisation-exact}, recalled here.
\begin{itemize}
\item $\ell_n=L^n\eta$, and $\pi_n$ is the partition into cubes of side $M\ell_n$. $\Lambda_j^{\sim}$
is the collared domain, obtained from $\Lambda_j$ by adding one layer of cubes of $\pi_j$, and
$\mathbf{B}'=\mathbf{B}_j(\Lambda_j^{\sim})$ is the determining set of \cite[(2.13)]{B-CMP119} with
$\Omega=\Lambda_j^{\sim}$, read on the lattice of spacing $\eta$. For a configuration $U$ and a
domain $D$ we write $U\restriction D$ for the restriction of $U$ to $D$.
\item $D_n$ are the maximal domains of $\mathbf{B}'$, which has level $n$ on $D_n\setminus D_{n+1}$
and level $j$ on $D_j$. A point at which $\mathbf{B}'$ has level $j$ is called a level-$j$ point
of $\mathbf{B}'$.
\item $M'=M_{\mathbf{B}'}$, and $Q'$ is the induced averaging operation.
\item $U^0=U_{k+1}$ and $U^1=(e^{i\eta\mathbf{H}_k(B')}U^0)^u$, $u$ a gauge transformation, are the
two endpoint configurations of \cite[(3.3)]{B-CMP109}; here $B'$ is the real fluctuation variable
of \cite[(3.15)]{B-CMP119}.
\item $\mathbf{H}_k=\mathbf{H}_k(B')$.
\item $\mathcal{C}_j$ is the cover, with partition of unity $\{\bar\zeta_\square\}$. Each
$\square\in\mathcal{C}_j$ has a scale $s\in[j,k]$, a window $\square_0$, the nested domains
$\tilde\square^2\subset\tilde\square^3\subset\tilde\square^4\subset\tilde\square^5=\square_0$ of
\cite[\S3]{B-CMP109}, and the cut-off $\zeta'_\square=\zeta_{\square_0}$ of \cite[(3.21)]{B-CMP109},
equal to $1$ on $\tilde\square^3$ and vanishing outside $\tilde\square^4$.
\item For a determining set $\mathfrak{B}$ and data $V$, $U(\mathfrak{B},V)$ is the corresponding
minimiser, so that $U_{j,\Lambda_j^{\sim}}(V)=U(\mathbf{B}',V)$. For a window $\square_0$ and a
level $i$, $\mathbf{B}_i(\square_0)$ is the determining set of \cite[(2.13)]{B-CMP119} with
$\Omega=\square_0$, truncated at level $i$; $M^{\cdot}U$ is the family of averages of $U$ at the
levels of this set, and $U_i(\square_0,M^{\cdot}(U))=U(\mathbf{B}_i(\square_0),M^{\cdot}U)$ is the
function of \cite[(1.15)]{B-CMP109}.
\end{itemize}
Good cubes are the cubes of $\mathcal{C}_j$ so called in the construction of the cover, as used in
Proposition~\ref{4.29:prop-near-field-scope}. In this remark
$X\subset\Lambda_j$ denotes the localisation domain of the term of \cite[\S3]{B-CMP109} under
consideration, $P_1$ the operator entering \cite[(3.14)]{B-CMP109}, and $B_\square$ the field of
the cube $\square$ entering \cite[(3.23)--(3.25)]{B-CMP109}. $\mathcal{H}'$ is the
function of \cite[Proposition~9]{B-CMP102b} attached to $\mathbf{B}'$, to $W_0:=M'(U^0)$ and to
the external field $U^0\restriction\Lambda_j^{\sim}$, as in
Lemma~\ref{4.29:lem-collar-staircase}. We write $\sigma_0$ for the family of walk cubes of the
random-walk expansion which \cite{B-CMP119} builds for the sequence $\{\Omega_n\}$; in
$\Omega_n\setminus\Omega_{n+1}$ these cubes have side $M\ell_n$. For a small real field
$\mathbf{a}$ we put
\begin{equation}\label{4.29:eq-transfer-inputs-1}
W[\mathbf{a}]:=U\bigl(\mathbf{B}',M'(e^{i\eta\mathbf{a}}U^0)\bigr).
\end{equation}
For $\square\in\mathcal{C}_j$ of scale $s$ we put $v:=s-j$ and call
$\mathfrak{r}_\square:=\square_0\setminus\tilde\square^4$ the rim of the window. For such a cube let
$Y\supset\square_0$ and $(\mathbb{U},\mathbb{J})\in\mathfrak{S}(Y)$ be as in
Proposition~\ref{4.29:prop-complex-hypothesis}, and let $\mathbf{A}$ denote the field of
\cite[(3.31)]{B-CMP109}.

Assume the following statements.
\begin{enumerate}
\item[(1)] Lemma~\ref{4.29:lem-localisation-exact}.
\item[(2)] The local reproduction lemma for minimisers. Let $D$ be a domain and let $\mathfrak{B}$
be a determining set with support in $D$. Let $U^*$ be a minimiser for a determining set whose
levels on $D$ are not below the local levels of $\mathfrak{B}$ and whose partitions are compatible
with those of $\mathfrak{B}$, and let $U^*$ be regular at all levels not below the local levels of
$\mathfrak{B}$. Then
$U(\mathfrak{B},M_{\mathfrak{B}}(U^*))=(U^*\restriction D)^u$ for a
gauge transformation $u$, blocks straddling $\partial D$ included.
\item[(3)] The constraint identity underlying \cite[(3.20)]{B-CMP109}.
\item[(4)] Lemma~\ref{4.29:lem-collar-staircase} and
Proposition~\ref{4.29:prop-near-field-scope}.
\item[(5)] Lemma~\ref{4.29:lem-interpolated-field-bounds}, under the floor of
Lemma~\ref{4.29:lem-collar-sources}.
\item[(6)] Clause (K1) of Proposition~\ref{4.29:prop-complex-hypothesis}.
\item[(7)] The rim lemma for the window minimiser, clauses (a)--(e).
\item[(8)] Lemma~\ref{4.29:lem-scale-thresholds}, with the rim bound
\eqref{4.29:eq-scale-thresholds-b}.
\item[(9)] The regularity theorem for the hybrid background, with its corollary.
\item[(10)] The walk-cube lemma on the collared domain.
\item[(11)] The transfer lemma, by which the analysis of \cite[\S\S3--4]{B-CMP109} is read at the
cubes of $\mathcal{C}_j$.
\item[(12)] The transfer-inputs lemma on a general domain $\Omega$, clauses (a)--(f) together with
the supplementary clause stated with them, of which the present dictionary is the case
$\Omega=\Lambda_j^{\sim}$.
\end{enumerate}
The transfer lemma leaves the choice of the minimiser $U_j$ open. Under hypotheses (1)--(12),
together with \cite[Theorems~3.7--3.10]{B-CMP99b}, the following dictionary holds for the
transfer lemma:
\begin{equation}\label{4.29:eq-transfer-inputs-a}
U_j:=U_{j,\Lambda_j^{\sim}},\qquad H_j\mapsto\mathcal{H}',\qquad Q_j\mapsto Q',\qquad
\sigma_0\restriction\Lambda_j^{\sim}:=\pi_j .
\end{equation}
Its clauses (a)--(g), referred to together with \eqref{4.29:eq-transfer-inputs-a}, are the
following.
\begin{enumerate}
\item[(a)] $U_j:=U_{j,\Lambda_j^{\sim}}$, with the collar $\Lambda_j^{\sim}$ formed in $\pi_j$ as
above. The
configuration $W[\mathbf{a}]$ depends on $\mathbf{a}$ only through its values within $\eta$ of
$\Lambda_j^{\sim}$. Up to gauge,
$W[0]=U^0\restriction\Lambda_j^{\sim}$ and $W[\mathbf{H}_k(B')]=U^1\restriction\Lambda_j^{\sim}$.
Moreover the expansion \cite[(3.4)]{B-CMP109}, written with the cover $\mathcal{C}_j$, the
partition $\{\bar\zeta_\square\}$ and this $U_j$, holds exactly, at all scales $j\le s\le k$.
\item[(b)] $H_j$ is read as $\mathcal{H}'$ and $Q_j$ as $Q'$. The bound \cite[(190)]{B-CMP102b}
is applied at the level of the evaluation element. Every evaluation point of
\cite[\S3]{B-CMP109} is a level-$j$ point of $\mathbf{B}'$. These points are the domain $X$, the
plaquettes entering \cite[(3.10)--(3.13)]{B-CMP109}, the blocks read by $P_1$, and $\square_0$
at good cubes. For the window's own cube set $\{\square_n\}$ this holds on $\square_j$ only;
see (g).
\item[(c)] At good cubes of scale $s>j$, and at cubes with $\square_0\subset\Omega_{k+1}$, the
relations \cite[(3.18)--(3.20)]{B-CMP109} hold for $W[(t+t_\square\bar\zeta_\square)\mathbf{H}_k]$.
After \cite[(3.21)--(3.22)]{B-CMP109} the main term is
\begin{equation}\label{4.29:eq-transfer-inputs-2}
U_j\bigl(\square_0,M^{\cdot}\bigl(e^{i\eta(t\zeta'_\square+t_\square\bar\zeta_\square)
\mathbf{H}_k}U^0\bigr)\bigr).
\end{equation}
The two terms removed at \cite[(3.21)--(3.22)]{B-CMP109} evaluate $\mathcal{H}'$ at level-$j$
points of $\square_0$ only. The main term \cite[(3.23)--(3.25)]{B-CMP109} is built from the
objects of $\square_0$ alone, namely $U_j(\square_0,\cdot)$, $\zeta_{\square_0}$ and $B_\square$.
It contains no $\mathbf{B}'$ and coincides symbol for symbol with that of \cite[\S3]{B-CMP109}.
The clauses of the transfer lemma, read with the transfer lemma's cube-level index $n$ replaced
by $s$, are unaffected.
\item[(d)] Put $B(t)=Q'(\eta t\mathbf{H}_k)$. The inputs \cite[(3.14)]{B-CMP109} for $H_j(B(t))$
and \cite[(3.9)]{B-CMP109} for $\delta H_j$ on $X$ are supplied by
Lemma~\ref{4.29:lem-interpolated-field-bounds}.
\item[(e)] For $\mathcal{H}'$ the walk cubes $\sigma_0\restriction\Lambda_j^{\sim}$ are those of
$\pi_j$, the ring cubes being taken whole. Under this choice
\cite[Theorems~3.7--3.10]{B-CMP99b} apply to $\mathbf{B}'$.
\item[(f)] The one-level offset treated within the transfer lemma holds in the cube's own scale
and units, as there, and is unaffected by the choice of
$U_j$. The configuration $e^{i\eta\mathcal{H}'(B(t))}U^0$ differs from $U^0$ by the term of
\cite[\S3]{B-CMP109} built from $H_j$, observed at level $j$.
\item[(g)] Where \cite[(3.18)]{B-CMP109} is entered, its cube set is
$\{\square_n\}=\mathbf{B}_b(\square_0)$, with $b=s$, and $b=k+1$ if
$\square_0\subset\Omega_{k+1}$. Then $\tilde\square^4\subset\square_b\subset\square_j$, and the
window's own staircase, formed by the elements of level $<b$ of $\mathbf{B}_b(\square_0)$ and of
level $<j$ of $\mathbf{B}_j(\square_0)$, lies within $2LM_1\ell_s$ of $\partial\square_0$, outside
$\tilde\square^4$, at $\ell^\infty$-distance at least $(M-2LM_1)\ell_s>0$ from it. Every
dimensionful bound of \cite[\S\S3--4]{B-CMP109} on an object of the window keeps the siting it
has in \cite{B-CMP109}, on $\square_0$, $\tilde\square^3$ or $\tilde\square^4$. Of the space of
configurations only conditions (i)--(iii) of \cite[\S1]{B-CMP109} for $(\mathbb{U},\mathbb{J})$
at level $b$ and condition \cite[(3.31)]{B-CMP109} for $\mathbf{A}$ are required on
$\mathfrak{r}_\square$. On the rim $\mathfrak{r}_\square$ the datum of $U_j(\square_0,\cdot)$ enters
\cite[\S\S3--4]{B-CMP109} in three ways and in no others:
\begin{itemize}
\item as data, six times, each removed at the cost bounded in \eqref{4.29:eq-scale-thresholds-b};
\item as structure, once, without reference to the configuration;
\item as background, once, in the form covered by the regularity theorem for the hybrid
background.
\end{itemize}
No term, counterterm piece, gauge cube or derived entry is read on the window's staircase. The
cost of the single background use depends on $v$ only.
\end{enumerate}
\end{remark}

\begin{proof}
\emph{Clause (a).} Lemma~\ref{4.29:lem-localisation-exact} takes $U_j:=U_{j,\Lambda_j^{\sim}}$. It
proves that \cite[(3.4)]{B-CMP109}, written with $(\mathcal{C}_j,\bar\zeta,U_{j,\Lambda_j^{\sim}})$,
is exact; every cube of $\mathcal{C}_j$ has a scale in $[j,k]$, so this covers the scales
$j\le s\le k$. The same lemma shows that $Q'(\eta\mathbf{a})(c)$ depends on $\mathbf{a}$ only
through its values on the blocks of the end-points of $c$. These blocks lie inside
$\Lambda_j^{\sim}$ as a set of bonds, that is, within $\eta$ of $\Lambda_j^{\sim}$. Since
$M'(e^{i\eta\mathbf{a}}U^0)=e^{iQ'(\eta\mathbf{a})}M'(U^0)$
(Lemma~\ref{4.29:lem-localisation-exact}), the configuration $W[\mathbf{a}]$ likewise depends on
$\mathbf{a}$ only through its values within $\eta$ of $\Lambda_j^{\sim}$.

For $\mathbf{a}=0$ we have $W[0]=U(\mathbf{B}',M'(U^0))$. The local reproduction lemma,
hypothesis (2),
applied with $D=\Lambda_j^{\sim}$, $\mathfrak{B}=\mathbf{B}'$ and $U^*=U^0$, gives
$W[0]=U^0\restriction\Lambda_j^{\sim}$ up to gauge. Its hypotheses hold: $U^0$ is the minimiser for
the determining set of \cite[(3.3)]{B-CMP109}, whose levels on $\Lambda_j^{\sim}$ are at least $j$
and so dominate every level of $\mathbf{B}'$, the partitions of the two sets are compatible, and
$U^0$ is regular at all levels $\ge j$ on $\Lambda_j^{\sim}$
(Lemma~\ref{4.29:lem-localisation-exact}).

For $\mathbf{a}=\mathbf{H}_k(B')$ we use $U^1=(e^{i\eta\mathbf{H}_k(B')}U^0)^u$. By the gauge
covariance of the averaging operation $M'$ and the definition of $Q'$
(Lemma~\ref{4.29:lem-localisation-exact}),
\[
M'(U^1)=e^{iQ'(\eta\mathbf{H}_k)}M'(U^0)=M'(e^{i\eta\mathbf{H}_k(B')}U^0)
\]
up to gauge. Since the minimiser $U(\mathbf{B}',\cdot)$ is covariant under gauge transformations
of its data \cite[(181)]{B-CMP102b}, it follows that
$W[\mathbf{H}_k(B')]=U(\mathbf{B}',M'(U^1))$ up to gauge. The configuration $U^1$ is likewise a
minimiser for a determining set whose levels on $\Lambda_j^{\sim}$ are at least $j$, with partitions
compatible with those of $\mathbf{B}'$, and it is regular at all levels $\ge j$ on
$\Lambda_j^{\sim}$ (Lemma~\ref{4.29:lem-localisation-exact}). The local reproduction lemma,
hypothesis (2), applied
with $U^*=U^1$, gives $U(\mathbf{B}',M'(U^1))=U^1\restriction\Lambda_j^{\sim}$ up to gauge, hence
$W[\mathbf{H}_k(B')]=U^1\restriction\Lambda_j^{\sim}$ up to gauge.

\emph{Clause (b).} Lemma~\ref{4.29:lem-collar-staircase} gives, on $\Lambda_j^{\sim}$, the identity
\[
U_{j,\Lambda_j^{\sim}}(e^{iB}W_0)=(e^{i\eta\mathcal{H}'(B)}U^0)^u .
\]
Hence the function $H_j$ of \cite[\S3]{B-CMP109} is $\mathcal{H}'$, and the averaging operation
$Q_j$ is $Q'$. The same lemma shows that each evaluation point listed in (b) lies in $\Lambda_j$
or within $2\ell_j+\eta$ of it, and that $\Lambda_j\subset D_j$ lies at $\ell^\infty$-distance at
least $(M-2M_1)\ell_j$ from the elements of $\mathbf{B}'$ of level $<j$. Such a point is therefore
a level-$j$ point of $\mathbf{B}'$, at $\ell^\infty$-distance at least $(M-2M_1-3)\ell_j$ from
every finer element of $\mathbf{B}'$. Consequently the bound \cite[(190)]{B-CMP102b} is used with
its level index equal to the level of the evaluation element. The exception for the window's own
cube set is treated in (g).

\emph{Clause (c).} By Proposition~\ref{4.29:prop-near-field-scope}, the cubes named in (c) are
those at which \cite[(3.18)]{B-CMP109} is entered, and at these cubes the argument of the
relations \cite[(3.18)--(3.28)]{B-CMP109} may be $W[\mathbf{a}]$ with any small real
$\mathbf{a}$. In particular $\mathbf{a}=(t+t_\square\bar\zeta_\square)\mathbf{H}_k$ is admitted.

At these cubes the identity \cite[(3.18)]{B-CMP109}, which writes $U_j$ at the cube as
$U_j(\square_0,M^\cdot(\cdot))$,
is the local reproduction lemma, hypothesis (2), with $D=\square_0$, with determining set
$\mathbf{B}_j(\square_0)$, the truncation at level $j$ of the cube set $\{\square_n\}$ of (g),
whose top stratum $\square_j$ contains $\tilde\square^4$, and with $U^*=W[\mathbf{a}]$ completed by
its level-$0$ data off
$D_1$. Its hypotheses hold: $W[\mathbf{a}]$ is the minimiser for $\mathbf{B}'$, and on
$\square_0\subset\Lambda_j\subset D_j$ the set $\mathbf{B}'$ has pure level $j$, which dominates
every level of $\mathbf{B}_j(\square_0)$, the partitions being compatible; the regularity of
$W[\mathbf{a}]$ required there is supplied by Proposition~\ref{4.29:prop-near-field-scope}. The
relation \cite[(3.20)]{B-CMP109} follows from its constraint identity, hypothesis (3).

The splitting \cite[(3.21)--(3.22)]{B-CMP109} replaces $t$ by $t\zeta'_\square$ in the argument
and produces the main term \eqref{4.29:eq-transfer-inputs-2}. The two removed terms contain
$\mathcal{H}'$ evaluated only at points of $\square_0$. By Proposition~\ref{4.29:prop-near-field-scope}
we have $\square_0\subset\Lambda_j\subset D_j$ at these cubes. By (b), these points are therefore
level-$j$ points of $\mathbf{B}'$.

The main term \cite[(3.23)--(3.25)]{B-CMP109} is formed from $U_j(\square_0,\cdot)$,
$\zeta_{\square_0}$ and $B_\square$ only. It contains no $\mathbf{B}'$, and so coincides with
that of \cite[\S3]{B-CMP109}. By clause (e) of the transfer-inputs lemma on a general domain,
hypothesis (12),
the clauses of the transfer lemma, read with $n$ replaced by $s$, are unaffected.

\emph{Clause (d).} The analysis \cite[(3.15)--(3.17)]{B-CMP109} rests on two inputs on $X$, in
level-$j$ units. The first is that $H_j(B(t))$ satisfies \cite[(3.14)]{B-CMP109} with
$\tfrac13\alpha_2$. The second is that $\delta H_j$ satisfies \cite[(3.9)]{B-CMP109}. These are
the membership clause and the clause on \cite[(3.9)]{B-CMP109} of
Lemma~\ref{4.29:lem-interpolated-field-bounds}, with the constant
\eqref{4.29:eq-interpolated-field-bounds-a}. That lemma holds under the floor of
Lemma~\ref{4.29:lem-collar-sources}, which is assumed in (5).

\emph{Clause (e).} \cite{B-CMP119} builds the random walks for the sequence $\{\Omega_n\}$,
with walk cubes $\sigma_0$ of side $M\ell_n$ in $\Omega_n\setminus\Omega_{n+1}$. For
$\mathcal{H}'$ the walk-cube lemma, hypothesis (10), takes
$\sigma_0\restriction\Lambda_j^{\sim}:=\pi_j$, the ring
cubes being taken whole. This choice is made because the layers of $\mathbf{B}'$ are at
most $2LM_1\ell_n$ thick and do not accommodate cubes of side $M\ell_n$.

Each $D_n$ is a union of $M_1\ell_n$-cubes. Hence the walk cubes of \cite{B-CMP99b} nest,
and a crossing between $\pi_j$-cubes costs at least the level-$j$ rate. The determining set
$\mathbf{B}'$ is admissible in the sense of \cite{B-CMP99a}, the block size $M_1$ of
\cite{B-CMP119} being at least the one required there (Lemma~\ref{4.29:lem-collar-staircase}).
That \cite[Theorems~3.7--3.10]{B-CMP99b} apply to $\mathbf{B}'$ with this choice of walk cubes
is part of the transfer lemma, hypothesis (11), which reads these theorems as \cite{B-CMP119}
invokes them, ``in their full generality''. This applicability statement is used here as stated
in the transfer lemma, hypothesis (11); the theorems themselves are used as stated in
\cite{B-CMP99b}.

\emph{Clause (f).} By clause (f) of the transfer-inputs lemma on a general domain, hypothesis
(12), the
one-level offset handled inside the transfer lemma, hypothesis (11), holds in the cube's own scale
and units and is unaffected by the choice
of the minimiser $U_j$ made in (a). By the identity recalled in the
proof of (b), the configuration $e^{i\eta\mathcal{H}'(B(t))}U^0$ differs from $U^0$ by the term
built from $H_j=\mathcal{H}'$. By (b) this term is observed at level-$j$ points, that is, at
level $j$.

\emph{Clause (g): the cube set and the staircase.} By Proposition~\ref{4.29:prop-near-field-scope},
\cite[(3.18)]{B-CMP109} is entered exactly at two kinds of cubes: good cubes with $s>j$, and
cubes with $\square_0\subset\Omega_{k+1}$. There the cube set of \cite[\S3]{B-CMP109} is
$\{\square_n\}=\mathbf{B}_b(\square_0)$. Here $b=s$ at good cubes, and $b=k+1$ where
$\square_0\subset\Omega_{k+1}$ (there $s=k$ by Proposition~\ref{4.29:prop-near-field-scope}, so
that $b\le s+1$ in both cases); in the latter case \cite[(3.18)]{B-CMP109} is accompanied by the
statement ``hence
$\tilde\square^4\subset\square_{k+1}$''. By the layers of \cite[(2.13)]{B-CMP119},
\begin{equation}\label{4.29:eq-transfer-inputs-3}
\square_n\supset\{x:\operatorname{dist}_\infty(x,\square_0^c)>2M_1\ell_n\},\qquad n\le b .
\end{equation}
Hence $\tilde\square^4\subset\square_b\subset\square_j$.

The window's own staircase consists of the elements of level $<b$ of $\mathbf{B}_b(\square_0)$
and of level $<j$ of $\mathbf{B}_j(\square_0)$. An element of level $n<b$ lies outside
$\square_{n+1}$. By \eqref{4.29:eq-transfer-inputs-3} it therefore lies within
$2M_1\ell_{n+1}\le 2M_1\ell_b\le 2LM_1\ell_s$ of $\partial\square_0$. Since
$\square_0=\tilde\square^5$ is obtained from $\tilde\square^4$ by adding a layer of width
$M\ell_s$ (\cite[\S3]{B-CMP109}), the staircase thus lies
outside $\tilde\square^4$, at $\ell^\infty$-distance at least $(M-2LM_1)\ell_s$ from it. This
distance is positive because $M>L^2M_1$ (clause (a) of the rim lemma, hypothesis (7)).

\emph{Clause (g): siting of the dimensionful bounds.} Every dimensionful bound of
\cite[\S\S3--4]{B-CMP109} on an object of the window keeps the siting it has in
\cite{B-CMP109}:
\begin{itemize}
\item \cite[(1.16)]{B-CMP109} at $X'=\square_0$;
\item \cite[(3.37)--(3.52)]{B-CMP109} on $\tilde\square^3$;
\item \cite[(3.20), (3.27), (3.32)]{B-CMP109}, the support of $\zeta'_\square$, and the product
$\zeta'B$ of \cite[(4.1)]{B-CMP109}, all on $\tilde\square^4$.
\end{itemize}
At \cite[(3.54)]{B-CMP109} the text reads: ``this bound holds on $\tilde\square^4$ only. We
should localize the inner product in domains $\tilde\square^4$, $(\tilde\square^4)^c$, and if at
least one factor has a localization in $(\tilde\square^4)^c$, then the expression can be
estimated by an arbitrary power of $L^j\eta$ using the exponential decay properties''. At
\cite[(4.1)]{B-CMP109} it reads: ``Terms with $m<n$ are exponentially small in $L^j\eta$''.

\emph{Clause (g): what is required on the rim.} Of the space of configurations,
\cite{B-CMP109} requires on $\mathfrak{r}_\square$ only two things:
\begin{itemize}
\item conditions (i)--(iii) of \cite[\S1]{B-CMP109} for $(\mathbb{U},\mathbb{J})$ at level $b$;
\item condition \cite[(3.31)]{B-CMP109} for $\mathbf{A}$.
\end{itemize}
The first is supplied as follows. In the case $s>j$ it is clause (K1) of
Proposition~\ref{4.29:prop-complex-hypothesis}. When $\square_0\subset\Omega_{k+1}$ it is the
clause of the space $\mathfrak{S}(Y)$ of that proposition at its top level. For the second, we
have
$\mathbf{A}=(t\zeta'_\square+t_\square\bar\zeta_\square)\mathbf{H}_k\equiv0$ on
$\mathfrak{r}_\square$, because both cut-offs $\zeta'_\square=\zeta_{\square_0}$ and
$\bar\zeta_\square$ vanish outside $\tilde\square^4$.

\emph{Clause (g): the rim datum as data.} On $\mathfrak{r}_\square$ the datum $B(c)$ of
$U_j(\square_0,\cdot)$ at a determining element $c$ of $U_j(\square_0,\cdot)$ is bounded in its own
units only. By clause (c) of the rim lemma, hypothesis (7),
$|B(c)|\le b_{\mathrm{out}}$, with $b_{\mathrm{out}}$ as in \eqref{4.29:eq-scale-thresholds-b}.
It enters \cite{B-CMP109} as data six times: at \cite[(3.21)]{B-CMP109}, at
\cite[(3.22)]{B-CMP109}, at \cite[(3.37)]{B-CMP109} and \cite[(3.45)]{B-CMP109}, at
\cite[(3.54)]{B-CMP109}, and at \cite[(4.1)]{B-CMP109}.

Each time it enters as a direction or an increment outside $\tilde\square^3$ or
$\tilde\square^4$, observed on $\tilde\square^2$ or $\tilde\square^3$. Each such entry is removed
by the device of \cite[(3.21)]{B-CMP109} and \cite[(4.1)]{B-CMP109}. By clauses (d)--(e) of the
rim lemma, hypothesis (7), and Lemma~\ref{4.29:lem-scale-thresholds}, each such entry is at most
$\tfrac14L^{-6v}\alpha_{2,j}$ in all four members of \cite[(3.14)]{B-CMP109}; this is the rim
bound \eqref{4.29:eq-scale-thresholds-b}.

\emph{Clause (g): the rim datum as structure.} It enters once as structure, in the passage of
\cite{B-CMP109} beginning ``Replacing the functions $H_j(\square_0)$ by $H_j$''. That
passage is independent of the configuration.

\emph{Clause (g): the rim datum as background.} It enters once as background:
$U_b(\square_0,M^{\cdot}\mathbb{U})$ under the operators of $\mathbf{H}_k$ on the whole lattice.
Of this configuration \cite[\S3]{B-CMP109} states: ``$U_{k+1}(\square_0,M^{\cdot}(U))$ has
approximately the same regularity properties as $U$, hence all necessary theorems are valid for
operations with this configuration''. This use is covered by the regularity theorem for the
hybrid background and its corollary, hypothesis (9). Let $\mathfrak{V}$ denote the hybrid
configuration of \cite[(3.24)]{B-CMP109} formed from $\mathbb{U}$ and the window minimiser
$U_b(\square_0,M^{\cdot}\mathbb{U})$. That theorem gives a gauge
transformation $u''$, equal to $1$ on the layer on which the data are pinned, and an
algebra-valued field $A''$ such that $\mathfrak{V}^{u''}=e^{i\eta A''}U$, where $U$ is the
$G$-valued part of $\mathbb{U}$ itself, the comparison being made level $n$ for level $n$.

\emph{Clause (g): conclusion.} The list above exhausts what \cite[\S\S3--4]{B-CMP109} reads on the
rim. It contains no term, counterterm piece, gauge cube or derived entry evaluated on the
window's staircase. The single background use is the one just described, and its cost depends
on $v$ only.
\end{proof}

\begin{corollary}[The small-field analysis at each cover cube at its native scale]
\label{4.29:cor-cubewise-analysis}
Let $j\le k$. Let $\mathcal{C}_j$, with the partition of unity $\{\bar\zeta_\square\}$, be the
cover of Proposition~\ref{4.29:prop-native-cover}. We use the notation of that proposition:
scales, good and default cubes, windows $\square_0$, the domains $\Lambda_j$, $\Lambda_j^{\sim}$,
$\Omega_i$, $\Omega^{\natural}_i$ and the pavings $\pi_s$. We call the sets
$\Omega_i\setminus\Omega_{i+1}$ the strata of the sequence $\{\Omega_i\}$. We also use the class
$\mathbf{D}_j$ of localisation domains of \cite{B-CMP119}, and the domain $D_j$ of level $j$ of
the determining set of the minimiser $U_{j,\Lambda_j^{\sim}}$
(Lemma~\ref{4.29:lem-localisation-exact}).
We use further, in the notation of Proposition~\ref{4.29:prop-native-cover}, the sets
$\Lambda_j^{\sim-1}$ and $\Lambda_j^0$, the units $\ell_s$ and the enlargements
$\tilde\square^{2},\tilde\square^{3},\dots$ of a cube. For
$\square\in\mathcal{C}_j$ write $s=s(\square)$ for its scale and $v:=s-j$. Assume the following
statements.
\begin{enumerate}
\item Proposition~\ref{4.29:prop-native-cover}, clauses (a)--(h$''$).
\item Lemma~\ref{4.29:lem-localisation-exact}, with the definitions
\eqref{4.29:eq-localisation-exact-a}.
\item Proposition~\ref{4.29:prop-near-field-scope}.
\item Proposition~\ref{4.29:prop-window-minimiser}, that is, (d1) and (d2) of
\eqref{4.29:eq-window-minimiser-a}.
\item Proposition~\ref{4.29:prop-complex-hypothesis}, that is, (K1)--(K3) and the inclusion
\eqref{4.29:eq-complex-hypothesis-a}.
\item Lemma~\ref{4.29:lem-scale-thresholds}: the thresholds across scales
\eqref{4.29:eq-scale-thresholds-a}, in particular its bounds on fields, and the thresholds at
the window's rim \eqref{4.29:eq-scale-thresholds-b}.
\item The inputs \eqref{4.29:eq-transfer-inputs-a} of Remark~\ref{4.29:rem-transfer-inputs}.
\item The transfer of the analysis of \cite[\S\S3--4]{B-CMP109} to the cover and its partition
of unity (in what follows, the transfer to the cover), printed without proof in
\cite[p.~276]{B-CMP119} and cited as printed. At a cube of
scale $s$ we read it with the level $n$ of the cubes of \cite{B-CMP119} replaced by $s$ and with
the substitutions of the inputs \eqref{4.29:eq-transfer-inputs-a}; we call this the reading with
$n\mapsto s$.
\item The conditions on the sequences $\{\Omega_i\}$, $\{\Lambda_i\}$ and on the radii $R_i$
under which Proposition~\ref{4.29:prop-native-cover} is stated, and the restriction lemma for
the minimiser on a subdomain, on which Lemma~\ref{4.29:lem-localisation-exact} rests.
\end{enumerate}
Fix $\square\in\mathcal{C}_j$ of scale $s$, let $Y\supset\square_0$ be a localisation domain and
$(\mathbb{U},\mathbb{J})\in\mathfrak{S}(Y)$, the space of
Proposition~\ref{4.29:prop-complex-hypothesis}. Consider the terms of \cite[(3.4)]{B-CMP109}
connected with $\square$. For these terms \cite[(3.5)]{B-CMP109} and the enlargement $\sim$ are
read in $\pi_s$, and distances are read in the unit $\ell_s$:
\begin{equation*}
\operatorname{dist}(X,\square)>M\ell_s,\qquad\text{hence}\qquad
\operatorname{dist}^{(\xi)}(X,\square)>ML^{s-j},
\end{equation*}
where $\operatorname{dist}^{(\xi)}$ is the distance of \cite[(3.5)]{B-CMP109} at level $j$, in
the unit $\ell_j$. The
minimiser $U_j$ of the expansion is $U_{j,\Lambda_j^{\sim}}$. Moreover the following hold.
\begin{enumerate}
\item[(A)] Suppose $\square$ is good with $s>j$, or $\square_0\subset\Omega_{k+1}$. Then:
\begin{itemize}
\item the local family $\{X\in\mathbf{D}_j: X\subset\tilde\square^{2}\}$ is complete, that is,
every $X\in\mathbf{D}_j$ with $X\subset\tilde\square^{2}$ occurs among the domains summed over
$z\in\Lambda_j^0$ in \cite[(2.26)--(2.27)]{B-CMP119};
\item $\square_0\subset\Lambda_j\subset D_j$, so \cite[(3.18)]{B-CMP109} is the instance at the
domain $\square_0$ of the restriction lemma for the minimiser;
\item the background is, on $\square_0$, a $U_s$-type configuration in the sense of (d1) of
\eqref{4.29:eq-window-minimiser-a};
\item every $(\mathbb{U},\mathbb{J})\in\mathfrak{S}(Y)$ lies on $\square_0$ in the level-$s$
cube space $\mathfrak{N}_s(\square_0)$, and on each $X\subset\Lambda_j$ in the term's own space,
with the margin in the radii that Proposition~\ref{4.29:prop-complex-hypothesis} provides;
\item the fields induced by the fluctuation obey the level-$s$ bounds of
\eqref{4.29:eq-scale-thresholds-a}.
\end{itemize}
Consequently all the formulas and bounds of \cite[\S\S3--4]{B-CMP109} hold with their parameter
$\eta$ replaced by $L^{-s}$, as bounds on $\mathfrak{S}(Y)$.
\item[(B)] Let $\square\in\mathcal{C}_j$ be any other cube of scale $s=j$, with
$\square_0\not\subset\Omega_{k+1}$, default or not. Its terms are
treated by \cite[(3.6)--(3.17)]{B-CMP109} alone. These read $(\mathbb{U},\mathbb{J})$ on $X$ and
the field $\mathcal{H}'$, which replaces the field $H_j$ of \cite[\S3]{B-CMP109} in the inputs
\eqref{4.29:eq-transfer-inputs-a}, at level-$j$ points of $X$. No factor $L^{j-s}$ is extracted,
and neither the completeness of the local family nor the chain of \cite[Lemma~4]{B-CMP109} is
used.
\end{enumerate}
\end{corollary}

\begin{proof}
\emph{The reading at scale $s$.} The cube $\square$ has scale $s$. The paving underlying
\cite[(3.5)]{B-CMP109} and the enlargement $\sim$ are therefore those of $\pi_s$, and lengths are
measured in $\ell_s$. Since $\ell_s=L^{s-j}\ell_j$, a distance exceeding $M\ell_s$ is a distance
exceeding $ML^{s-j}$ in units of $\ell_j$, the unit in which \cite[(3.5)]{B-CMP109} measures
distances at level $j$ and in which the domains $X\in\mathbf{D}_j$ are written. This is the
displayed reading.

\emph{The minimiser.} By Lemma~\ref{4.29:lem-localisation-exact}, with the definitions
\eqref{4.29:eq-localisation-exact-a}, the minimiser $U_j$ entering the expansion
\cite[(3.4)]{B-CMP109} is $U_{j,\Lambda_j^{\sim}}$.

\emph{Case (A): completeness.} Let $\square$ be good with $s>j$, or let
$\square_0\subset\Omega_{k+1}$. In either case
\begin{equation*}
\tilde\square^{2}\subset\Omega^{\natural}_{j+1}\subset\Lambda_j^{\sim-1}.
\end{equation*}
For a good cube with $s>j$ this follows from clause (c) of
Proposition~\ref{4.29:prop-native-cover} together with the separation conditions on the
sequences $\{\Omega_i\}$, $\{\Lambda_i\}$ and the bound on the radii $R_i$ under which that
proposition is stated. For a cube with $\square_0\subset\Omega_{k+1}$ it follows from
Proposition~\ref{4.29:prop-near-field-scope}, under the same conditions. The sums of
\cite[(2.26)--(2.27)]{B-CMP119} run over $z\in\Lambda_j^0$. By the displayed inclusion, every
$X\in\mathbf{D}_j$ with $X\subset\tilde\square^{2}$ occurs in them, so the local family
$\{X\in\mathbf{D}_j:X\subset\tilde\square^{2}\}$ is complete.

\emph{Case (A): the restriction identity.} A good cube with $s>j$ is not default, because a
default cube has scale $j$ by clause (c) of Proposition~\ref{4.29:prop-native-cover}. The same
clause then gives $\square_0\subset\Lambda_s\subset\Lambda_j$. For a cube with
$\square_0\subset\Omega_{k+1}$, Proposition~\ref{4.29:prop-near-field-scope} gives
$\square_0\subset\Lambda_j$. Moreover $\Lambda_j\subset D_j$ by
Lemma~\ref{4.29:lem-localisation-exact}, with the definitions
\eqref{4.29:eq-localisation-exact-a}. Hence at these cubes $\square_0\subset\Lambda_j\subset D_j$.

Identity \cite[(3.18)]{B-CMP109} is therefore the instance at the domain $\square_0$ of the
restriction lemma for the minimiser on which Lemma~\ref{4.29:lem-localisation-exact} rests.
By (d1) of \eqref{4.29:eq-window-minimiser-a}, the background is a $U_s$-type
configuration on $\square_0$.

\emph{Case (A): the space and the fields.} By Proposition~\ref{4.29:prop-complex-hypothesis}, the
inclusion \eqref{4.29:eq-complex-hypothesis-a} holds: every
$(\mathbb{U},\mathbb{J})\in\mathfrak{S}(Y)$ lies, on $\square_0$, in the level-$s$ cube space
$\mathfrak{N}_s(\square_0)$. By clause (K3) of the same proposition, on each $X\in\mathbf{D}_j$
with $X\subset\Lambda_j\cap Y$ the pair lies in the space of the term $X$ at level $j$, with the
margin in the radii that Proposition~\ref{4.29:prop-complex-hypothesis} provides. By the bounds
on fields in \eqref{4.29:eq-scale-thresholds-a} of Lemma~\ref{4.29:lem-scale-thresholds}, the
fields induced by the fluctuation obey the bounds at level $s$ on $\square_0$.

\emph{Case (A): conclusion.} The statements just established, together with the inputs
\eqref{4.29:eq-transfer-inputs-a} of Remark~\ref{4.29:rem-transfer-inputs}, which use the
thresholds \eqref{4.29:eq-scale-thresholds-b} at the window's rim, are the hypotheses of the
transfer to the cover. Read with $n\mapsto s$, it passes from these inputs to the formulas and
bounds of \cite[\S\S3--4]{B-CMP109}, with their parameter $\eta$ replaced by $L^{-s}$. They hold
as bounds on
$\mathfrak{S}(Y)$, since the inputs hold for every $(\mathbb{U},\mathbb{J})\in\mathfrak{S}(Y)$.

\emph{Case (B).} Let $\square$ be a cube of $\mathcal{C}_j$ not covered by (A), of scale $s=j$
and with $\square_0\not\subset\Omega_{k+1}$, default or not. Its terms are treated by
\cite[(3.6)--(3.17)]{B-CMP109} alone. These formulas read $(\mathbb{U},\mathbb{J})$ on $X$, which
is supplied by clause (K3) of Proposition~\ref{4.29:prop-complex-hypothesis}. They also read
$\mathcal{H}'$ at level-$j$ points of $X$, which is supplied by the inputs
\eqref{4.29:eq-transfer-inputs-a}.

Because $s=j$, no factor $L^{j-s}$ is extracted. Completeness of the local family is not used.
Nor is the chain of \cite[Lemma~4]{B-CMP109}; that chain is closed in \cite{B-CMP109} by the
inequality $(1+7\beta)L^{-2}\le1$, a gain that is absent within one layer. This is the collar case
as it is treated in \cite{B-CMP119}.

\emph{The level of the background.} In case (A) the passage from the inputs to the bounds is the
transfer to the cover, read with $n\mapsto s$. Two points of this reading concern the level of
the background relative to the level of the cube.
\begin{itemize}
\item No further term enters the reading: by the inputs \eqref{4.29:eq-transfer-inputs-a}, the
configuration at which the cube is read differs from the background only by the term of the field
$H_j$ of \cite[\S3]{B-CMP109} itself, observed at level $j$.
\item The one-level offset is the one of \cite{B-CMP119} and is unchanged.
There the cube has level $n$ and the background level $n$; in \cite[\S3]{B-CMP109} the cube has
level $k$ and the background level $k+1$. This offset is felt only at $v=0$, where no gain is
asked.
\end{itemize}
In case (B) the one-level offset is likewise the one of \cite{B-CMP119}.
\end{proof}

\begin{remark}[Families not treated through the window]
The following families are not treated through the window $\square_0$ of
Corollary~\ref{4.29:cor-cubewise-analysis}.
\begin{itemize}
\item The families of $\mathbf{B}$-terms: the old terms $\mathbf{B}^{(j)}$, the terms
$\mathbb{B}^{(n)}_{k_W}$ of level $n$ frozen at the arrest level $k_W$, and the Wilson terms with
variable coupling. These are
treated by \cite[(3.6)--(3.7)]{B-CMP109} with the plaquette-wise partition of unity $\zeta$, as
in \cite{B-CMP119}, with no window; their treatment is that of the arrest theorem.
\item The function $\mathbf{P}^{(k)}$ and the quadratic form. On $\Omega_{k+1}$, which lies in
the near-field set $P$ of clause (b) of Proposition~\ref{4.29:prop-native-cover}, they
enter, in the notation of \cite[(3.27), (3.29)]{B-CMP119}, as
$\mathbf{P}^{(k)}(g_k,(tA_k,A))$ and $\langle A,C^{*}\Delta^{(k)}CA_k\rangle$ in those displays.
On $\Lambda_{k+1}$ they enter the last logarithm of
\cite[(3.28)]{B-CMP119}. As \cite{B-CMP119} states, ``the quadratic form in the exponential
couples only the fields $A_j$ in the same component of $Z_{j+1}\cap\Omega_{j+1}$''.
\item The background-field operators built for the sequence $\{\Omega_i\}$ that \cite{B-CMP119}
uses. Their windows consist of twelve $M_1$-blocks, with $M$ larger than $M_1$ times the radius
parameter of the admissibility condition of \cite{B-CMP122a} on the sequence $\{\Omega_i\}$.
On nested sequences
of domains they reach one stratum beyond their own, as in \cite{B-CMP119}; such a nested
sequence is admissible in the sense of \cite{B-CMP122a}.
\end{itemize}
\end{remark}

Every condition in the remark below is one-sided; the order in which they are chosen is fixed in \eqref{4.29:eq-constants-order-2}.

\begin{remark}[Constants, one-sided lower bounds and their order of choice]
\label{4.29:rem-constants-order}
In this remark the following notation is used.
\begin{itemize}
\item $\delta_0$ is the decay rate of \cite[Proposition~9, (190)]{B-CMP102b}.
\item $\alpha_1^{(98)}$ is the radius $\alpha_1$ of \cite[Proposition~6]{B-CMP98}.
\item $C_{\mathrm{II}}$ is the absolute constant of \cite[(1.20)]{B-CMP116}.
\item $C_1^{(3.43)}$ is the constant $C_1$ of \cite[(3.43)]{B-CMP119}.
\item $C_0,C_1$ are the constants of \cite[(2.28)]{B-CMP119}.
\item $\varepsilon^{\mathrm{abs}}_1$ is the radius of the polydisc of
\cite[Proposition~9, (172)]{B-CMP102b}.
\item $\varepsilon^{\mathrm{I}}_1$ is the number $\varepsilon_1$ in the hypothesis
``$g_k|B|<\varepsilon_1$'' of \cite[(1.20)]{B-CMP116}; it is fixed after $M$.
\item $a_1^{(99)}$ is the constant $a_1$ of \cite{B-CMP99b}.
\item $\kappa'=(1+4\beta)\kappa$ and $\kappa_L:=\kappa'(1-1/L)=(1+4\beta)\kappa(1-1/L)$.
\end{itemize}
All other letters are those of the statements of this section to which we refer.

\emph{(a) Constants.} The following constants are independent of the labels, of $k$, of $K$ and of the torus
$\mathbb{T}_m$.
\begin{itemize}
\item The numerical constant $\theta\le 0.87$ that enters the estimates of this section.
\item The constants $C_\zeta$ and $C_{\mathrm{reg}}$ of \eqref{4.29:eq-scale-thresholds-a}.
\item The factor $eL^{25}$ of the proof of Lemma~\ref{4.29:lem-far-field-paths}.
\item The constant $C_{\mathfrak{s}}=1+4e^{\delta_0/8}C_dL^{d-1}$ of
\eqref{4.29:eq-interpolated-field-bounds-a}; under $e^{\delta_0/8}\le2$ it is at most $1+8C_dL^{d-1}$.
\item The constant $C_{90}$ of the layered kernel bound used in Lemma~\ref{4.29:lem-collar-staircase}.
\item The constant $c_Q=2c_6'$ of \eqref{4.29:eq-collar-staircase-b}.
\end{itemize}
The last three enter the threshold $\bar\gamma=\bar\gamma(L,M,\mathfrak{p})$ on $\gamma$, fixed in (c) below,
through \eqref{4.29:eq-scale-thresholds-a} and through $e^{-6R_{k+1}}\le\gamma^{12}$.

The constant $C_{\mathrm{ff}}$ of \eqref{4.29:eq-far-field-count-a} is built from the rate $\kappa_L$ and depends
on $(L,\kappa,\beta)$ only. It enters
$\bar\gamma$, together with the factor $2^d$ of that display, through $e^{-5R_{k+1}}\le\gamma^{10}$.

The constants $C_\Sigma$, $K_w$ and $C_w$ of \eqref{4.29:eq-scale-thresholds-b}, and
$C_{\mathrm{rim}}=C_\Sigma(1+2e^{\delta_0/16}C_dL^d)$, belong to this list. In $C_{\mathrm{rim}}$ the factor
$e^{\delta_0/16}$ is kept in place of its bound $2$, which would give $C_\Sigma(1+4C_dL^d)$, exactly as
$C_{\mathfrak{s}}$ and the factor $4C_dL^{d-1}$ carry $e^{\delta_0/8}$ in place of $2$ in
Lemma~\ref{4.29:lem-collar-sources}. The constants $K_w$ and $C_w$ depend only on $d$, $L$, $\delta_0$,
$\beta_0$, $\beta_s$, $c_2$, $B_3$, $B_5$ and absolute constants, and are independent of $M$, $C_0$, $C_1$,
$A_0$ and $A_1$.

Finally, the constant of the per-cube factor of \eqref{4.29:eq-near-field-scope-a} is
\begin{equation}\label{4.29:eq-constants-order-1}
C^\sharp:=\frac{6L^2B_3^2\,C_N\,C_{\mathrm{II}}\,
\max\{4a_\sharp A_0,\;C_1^{(3.43)}A_1\}}{c_2\,\beta_s\,\min\{C_0,C_1\}} .
\end{equation}
Here $C_N:=C_{\mathfrak{s}}\,C_{90}\cdot 2c_Q\cdot C$, where $C$ is the constant of
\eqref{4.29:eq-interpolated-field-bounds-a} and \eqref{4.29:eq-scale-thresholds-a};
$a_\sharp=1+c_{14}$ is as in \eqref{4.29:eq-path-length-decay-b}, and $c_2,\beta_s$ are those of
\eqref{4.29:eq-scale-thresholds-b}. The constant $C^\sharp$ is fixed after $L$, $M_1$, $M$, $C_0$, $C_1$, $A_0$,
$A_1$ and before $\gamma$.

\emph{(b) Floors.} The following one-sided conditions are imposed. All of them are lower bounds on $L$, $M_1$ or
$M$, except $R_i\ge7$, which follows from the coupling regime, and the last four, which are a smallness condition
on $\beta$ and one-sided choices of the free exponents and constants of \cite[(2.28)]{B-CMP119}.
\begin{itemize}
\item $L\ge 13$. This floor is imposed as it stands. It holds in particular whenever the integer $L_\circ$ of
Lemma~\ref{4.29:lem-radius-conversion}, which satisfies $L_\circ^2\le L/3$, is at least $2$, since $L$ is odd.
\item $R_i\ge 7$ for every $i$.
\item $M\ge L^4$.
\item $M\ge C_\zeta L^2M_{\mathrm{I}}$ wherever \cite[\S\S3--4]{B-CMP109} uses smallness of the gradients of the
partition of unity as a threshold. Here $M_{\mathrm{I}}$ is the lower bound on $M$ of that analysis. This
condition is one-sided and is imposed after $L$. It enters through the transfer of the analysis of
\cite[\S3]{B-CMP109} to the cover, printed without proof in \cite[p.~276]{B-CMP119} and cited as printed.
\item $M>R^{\mathrm{abs}}M_1$, where $R^{\mathrm{abs}}$ is an absolute constant of Ba{\l}aban's construction.
\item $M>M_2=L^2M_1$, as in \cite{B-CMP119}.
\item The second line of \eqref{4.29:eq-constants-order-c}. It is a floor on $M$ chosen after $L$ and $\kappa$.
\item The native-cover floor \eqref{4.29:eq-native-cover-a}, $L^{d-1}e^{-\delta_0M_1/16}\le\frac12$, for the
parameter $M_1$ of \cite{B-CMP119}. It involves $L$ and $M_1$ only, and $M_1$ is chosen after $L$ and before
$M_2=L^2M_1<M$.
\item The first line of \eqref{4.29:eq-constants-order-c}. In it, $\mathfrak{M}^\natural$ is the least value of
$M_1$ admitted by the native-cover floor \eqref{4.29:eq-native-cover-a}, and $R_\star^\natural$ is the least value
admitted by the floor on the parameter $R_\star$ at $M_1=\mathfrak{M}^\natural$. The condition depends on $L$ only.
\item The third line of \eqref{4.29:eq-constants-order-c}. It is a floor on $M$ chosen after $L$ and $M_1$ and
before $\gamma$.
\item $\beta\le\frac14$; the standing value $\beta=\frac15$ of this chapter satisfies it.
\item $q\ge p_0+1$, where $q:=\min\{q_0,q_1\}$ and $q_0,q_1$ are the exponents of \cite[(2.28)]{B-CMP119}.
\item $q_0=q_1$ wherever \eqref{4.29:eq-scale-thresholds-b} is invoked.
\item $C_1\ge C_0$ wherever \eqref{4.29:eq-scale-thresholds-b} is invoked.
\end{itemize}
The displayed conditions are
\begin{equation}\label{4.29:eq-constants-order-c}
\begin{aligned}
&M_1\ \text{is a power of}\ L,\qquad M_1\ge R_\star^\natural\,\mathfrak{M}^\natural;\\
&\mathsf{d}\ge(1+4\beta)\kappa+4\log L+1;\\
&\delta_0(M-2)/16\ge\log(8C_wM);\\
&\max\{4a_\sharp\varepsilon_k,\;C_1^{(3.43)}\delta_k\}<\varepsilon^{\mathrm{I}}_1
\quad\text{and}\quad
\max\{4a_\sharp\varepsilon_k,\;C_1^{(3.43)}\delta_k\}\le\varepsilon^{\mathrm{abs}}_1;\\
&L^2K_{\mathrm{rim}}C_\flat\,B_{\mathrm{rim}}\,M\,\bar\varpi\le
\min\{a_1^{(99)},\;\alpha_0',\;\alpha_1',\;\alpha_1^{(98)}/4\}.
\end{aligned}
\end{equation}
The first three lines of \eqref{4.29:eq-constants-order-c} are floors; the last two are ceilings on $\gamma$ and
belong to (c) below.

In \eqref{4.29:eq-constants-order-c}:
\begin{itemize}
\item $C_w$ is the constant of \eqref{4.29:eq-scale-thresholds-b}; its dependence is as stated in (a).
\item $\bar\varpi$ is the quantity of \eqref{4.29:eq-scale-thresholds-a}. It equals
$L(L+1)N_0^{\beta_0}R_{k_W}(\alpha_{0,n}+\alpha_{1,n})$, $n\le k_W$, on the region $W$ of that display, and
$\alpha_{0,n}+\alpha_{1,n}$ off $W$.
\item $\alpha_0',\alpha_1'$ are the radii of the analyticity domain in \cite{B-CMP109}, taken sufficiently small
but much bigger than $\alpha_0,\alpha_1$.
\item $K_{\mathrm{rim}}$, $C_\flat$ and the factor $B_{\mathrm{rim}}$ are the constants of the rim regularity bound.
\end{itemize}
The order of choice is a single chain:
\begin{equation}\label{4.29:eq-constants-order-2}
L\to\mathfrak{M}^\natural\to R_\star^\natural\to M_1\to M_2=L^2M_1\to R_\star(M_1)\to\kappa\to M\to\gamma .
\end{equation}
Here $R_\star(M_1)$ is the floor on $R_\star$ evaluated at the chosen value of the parameter $M_1$ of
\cite{B-CMP119}.

\emph{(c) The ceiling on $\gamma$.} The threshold $\bar\gamma(L,M,\mathfrak{p})$ is chosen with
$\bar\gamma(L,M,\mathfrak{p})\le\gamma_9$, where $\gamma_9$ is the threshold on $\gamma$ of the statements used in
this section. The following conditions are
imposed. The first item is a smallness condition on $\beta_0$ together with the lower bound $x_{k+1}\ge B+c_5$; the
second item is the coupling regime, a lower bound on the $x_i$, that is, an upper bound on the couplings; the fourth
item is a smallness condition on the primed radius. Every other item is met by every
$\gamma\le\bar\gamma(L,M,\mathfrak{p})$.
\begin{enumerate}
\item $\beta_0\le 1/10$, and $x_{k+1}\ge B+c_5$, with $B$ the flow constant of the hypotheses of
Lemma~\ref{4.29:lem-radius-conversion}; this number $B$ is not the field $B$ of the sources.
\item $x_i\ge\gamma^{-2}$ at every flow index $i\le k+1$ at which $x_i$ enters: $i=k$ in the per-cube factor of
\eqref{4.29:eq-near-field-scope-a}, $i=k+1$ in \eqref{4.29:eq-far-field-paths-a}, and the indices entering
\eqref{4.29:eq-scale-thresholds-a}, \eqref{4.29:eq-scale-thresholds-b} and
Lemma~\ref{4.29:lem-radius-conversion}, all of them at most $k+1$. We also require $\gamma\le e^{-1/2}$.
\item $\gamma\le\gamma(p_0,q_0,q_1,c_5,\beta_0)$, so small that
\[
\Bigl(\frac{\log(x+c_5)}{\log x}\Bigr)^{q'}\le 1+\beta_0
\qquad\text{for every } x\ge\gamma^{-2} \text{ and every } q'\in\{p_0,q_0,q_1\}.
\]
\item $\varepsilon_1'\le a_1$ and $B_3\varepsilon_1'\le a_0$. Here $\varepsilon_1'$ is the primed radius and
$a_0$, $a_1$ are the constants of the ceilings on the primed radii,
Definition~\ref{4.4:def-radius-ceilings}; they are not identified with $a_1^{(99)}$.
\item The same-scale threshold of \eqref{4.29:eq-scale-thresholds-a} holds.
\item The quantities $L(L+1)N_0^{\beta_0}R_{k_W}\alpha_{\cdot,n}$, $n\le k_W$, lie inside the absolute radii of the
analyticity domains used in \cite{B-CMP109}.
\item The condition $\ell_m\sup|A'|\le\alpha_1^{(98)}/4$ of the source bound \eqref{4.29:eq-collar-staircase-b}
holds. Moreover, at every twist $\mathbf{a}$ at which $Q'$ or $DQ'$ is evaluated, we require
$64\,\ell_{m(c)}\sup_{A(c)}|\mathbf{a}|\le\alpha_1^{(98)}$. The twists in question are:
\begin{itemize}
\item in Lemma~\ref{4.29:lem-collar-staircase}, Lemma~\ref{4.29:lem-interpolated-field-bounds} and
Proposition~\ref{4.29:prop-near-field-scope}, the twist $\mathbf{a}=t\mathbf{H}_k(B')$ with $t\le1$;
\item in \eqref{4.29:eq-scale-thresholds-b}, the twist $\mathbf{a}=A$.
\end{itemize}
These requirements amount to
\begin{equation}\label{4.29:eq-constants-order-b}
64\,C B_3\,\delta_k^\sharp\le\alpha_1^{(98)},\qquad
64\,\varepsilon_2=64\,B_5\cdot O(1)\,(1+2B_3)\,M\,\alpha_{0,b}\le\alpha_1^{(98)} .
\end{equation}
Here $\varepsilon_2$ is the constant of the bound \cite[(19)]{B-CMP102b} on the fluctuation field, which is a
hypothesis of \cite[Proposition~9]{B-CMP102b}, taken at the data's own size, as in
\eqref{4.29:eq-scale-thresholds-b}; it is not the fixed radius $B_5\varepsilon^{\mathrm{abs}}_1$. The index $b$
is the level whose radius $\alpha_{0,b}$ bounds the rim data in \eqref{4.29:eq-scale-thresholds-b}; in the
notation there, $b-j\le v+1$.
\item Let $C(L,M,\mathfrak{p})$ denote a constant depending on $(L,M,\mathfrak{p})$ only which bounds the product of
the constants depending on $(L,M,\mathfrak{p})$ only that the surplus factors of
\eqref{4.29:eq-far-field-paths-a} and \eqref{4.29:eq-far-field-count-a} and the per-cube factor of
\eqref{4.29:eq-near-field-scope-a} have to absorb, among them the constants of \cite[(3.15)--(3.17)]{B-CMP109};
it is free upward. Then $C(L,M,\mathfrak{p})\gamma^{12}\le1$ and
$2^dC_{\mathrm{ff}}\,C(L,M,\mathfrak{p})\gamma^{10}\le1$.
\item The product ceiling
\begin{equation}\label{4.29:eq-constants-order-a}
\frac{2^dL^4\cdot C(L,M,\mathfrak{p})\cdot C^\sharp}{\log\gamma^{-2}}\le 1 .
\end{equation}
Among the factors of $C(L,M,\mathfrak{p})$ is $e^{16\kappa_1}$, which comes from the factor $4B_0C_1e^{16\kappa_1}$
of \cite[(1.21)--(1.22)]{B-CMP116} (in the notation there).
\item The fourth line of \eqref{4.29:eq-constants-order-c} holds. It is the proviso on the domain of
\eqref{4.29:eq-path-length-decay-b}.
\item The condition $3b_{\mathrm{out}}\le 2C_1^{\mathrm{abs}}\varepsilon^{\mathrm{abs}}_1$ on the rim data holds.
Here $b_{\mathrm{out}}$ is as in \eqref{4.29:eq-scale-thresholds-b}, and $C_1^{\mathrm{abs}}$ is the absolute
constant $C_1$ of \cite[Proposition~9, (172)]{B-CMP102b}. This condition is an input of the bound on the rim data.
\item The fifth line of \eqref{4.29:eq-constants-order-c} holds, with $a_\square\le c_{\mathrm{rim}}$: the
amplitude $a_\square$ attached to the cube $\square$ in the hypothesis of the rim regularity bound is at most
the threshold $c_{\mathrm{rim}}$ of that hypothesis.
\end{enumerate}
None of these conditions bounds $L$, $M$, $N$, $N_0$, $K$, $k$, $|W|$ or the arrest levels from above. No window
is imposed. The condition $L_\circ^2\le L/3$ is Ba{\l}aban's choice of an integer $L_\circ$ below a bound that
grows with $L$. No condition compares a power of $L$ with $L_0$.
\end{remark}

\begin{proof}
We justify the assertions of the remark that are not definitions, in the order in which they appear.

\emph{Dependence of the constants.} Each constant listed in (a) is built from constants that are independent of
the labels, of $k$, of $K$ and of $\mathbb{T}_m$. The constant $C^\sharp$ of \eqref{4.29:eq-constants-order-1}
depends on $L$, $B_3$, $C_N$, $C_{\mathrm{II}}$, $a_\sharp$, $A_0$, $A_1$, $C_1^{(3.43)}$, $c_2$, $\beta_s$, $C_0$
and $C_1$. All of these are fixed once $L$, $M_1$, $M$, $C_0$, $C_1$, $A_0$, $A_1$ are, so $C^\sharp$ can be fixed
before $\gamma$.

\emph{The floor on $L$.} The floor $L\ge13$ is a lower bound, hence one-sided. It holds in particular whenever
$L_\circ\ge 2$: then $L_\circ^2\le L/3$ gives $L\ge 3L_\circ^2\ge 12$, and since $L$ is odd, $L\ge 13$. The
integer $L_\circ$ is chosen below the bound $(L/3)^{1/2}$, which grows with $L$, so this imposes
no upper bound on $L$.

\emph{The floor on $\mathsf{d}$.} By Definition~\ref{4.29:def-path-length-decay}, $\mathsf{d}$ is the decay exponent
per local $M$-cube of the layered bound for $H_k$. At the rate $\delta_0/8$ per site, it equals
$\mathsf{d}=\delta_0M/8$. Since $\kappa'=(1+4\beta)\kappa$, the second line of
\eqref{4.29:eq-constants-order-c} is the floor $\mathsf{d}\ge\kappa'+4\log L+1$ of
\eqref{4.29:eq-far-field-paths-a}. Equivalently,
\[
\delta_0M\ge 8(\kappa'+4\log L+1).
\]
Once $L$ and $\kappa$ are fixed, this is a lower bound on $M$, so $M$ is chosen after $L$ and $\kappa$.

\emph{The third line of \eqref{4.29:eq-constants-order-c}.} The constant $C_w$ does not depend on $M$. The
difference
\[
\delta_0(M-2)/16-\log(8C_w)-\log M
\]
is therefore a function of $M$ that tends to $+\infty$ as $M\to\infty$, since $M/\log M\to\infty$. So the condition
is non-empty and holds for all $M$ beyond a value depending on $L$ and on the quantities entering $C_w$. It is
thus a condition of the form $M\ge M_*(L)$. As $C_w$ is fixed after $L$, and $M$ is in any case chosen after $M_1$
(since $M>L^2M_1$), this is a floor on $M$ chosen after $L$ and $M_1$ and before $\gamma$. It is one-sided.

\emph{The floors on $M_1$.} The left side of \eqref{4.29:eq-native-cover-a} involves only $L$ and $M_1$ and
decreases as $M_1$ increases. So it is a lower bound on $M_1$ chosen after $L$. The first line of
\eqref{4.29:eq-constants-order-c} asks $M_1$ to be a power of $L$ not smaller than $R_\star^\natural\mathfrak{M}^\natural$.
Once $L$, $\mathfrak{M}^\natural$ and $R_\star^\natural$ are fixed, this asks only that $M_1$ be large. Both are
one-sided. Together with $M_2=L^2M_1<M$, these floors produce the chain \eqref{4.29:eq-constants-order-2}.

\emph{Classification of the floors.}
\begin{itemize}
\item $R_i\ge 7$ is a consequence of the coupling regime $x_i\ge\gamma^{-2}$, through the bound
$R_i\ge\log x_i$ on Ba{\l}aban's length parameters $R_i$.
\item $\beta\le\frac14$ and $\beta_0\le 1/10$ are one-sided smallness conditions on auxiliary constants of the
estimates.
\item $q\ge p_0+1$, $q_0=q_1$ and $C_1\ge C_0$ are one-sided choices of the free exponents and constants of
\cite[(2.28)]{B-CMP119}, free upward. There $q_0,q_1$ are integers greater than $1$, and $C_0,C_1$ are taken
sufficiently large.
\item $C_1\ge C_0$ and $q_0=q_1$ give $\min\alpha_{\cdot,j}=\alpha_{0,j}$ in $\alpha_{3,j}$ of
\eqref{4.29:eq-scale-thresholds-b}, since
\[
\alpha_{0,j}=g_jC_0(\log g_j^{-2})^{q_0},\qquad\alpha_{1,j}=g_jC_1(\log g_j^{-2})^{q_1}.
\]
\item All the remaining floors are lower bounds on $L$, $M_1$ or $M$, in the order
\eqref{4.29:eq-constants-order-2}.
\end{itemize}
None of them bounds $L$, $M$, $K$, $k$, $N_0$ or the coupling from above. None bounds the coupling from below.

\emph{The coupling regime.} The coupling-regime condition of (c) bounds the coupling from above only. Since
$\gamma\le e^{-1/2}$, we have $x_i\ge\gamma^{-2}\ge e$, hence $\log x_i\ge1$. This is used twice.
\begin{itemize}
\item Let $q:=\min\{q_0,q_1\}$. Since $\log x_k\ge1$, for $\iota=0,1$,
\[
\alpha_{\iota,k}=g_kC_\iota(\log x_k)^{q_\iota}\ge(\min\{C_0,C_1\})\,g_k(\log x_k)^q ,
\]
so
\[
\min\alpha_{\cdot,k}\ge(\min\{C_0,C_1\})\,g_k(\log x_k)^q .
\]
\item Since $q\ge p_0+1$, the exponent $p_0-q$ is at most $-1$. With $\log x_k\ge1$ this gives
\[
(\log x_k)^{p_0-q}\le(\log x_k)^{-1} .
\]
\end{itemize}
These are the two bounds used in the per-cube factor of \eqref{4.29:eq-near-field-scope-a}.

\emph{The ceilings \eqref{4.29:eq-constants-order-b}.} The quantities $\delta_k^\sharp$ of
\eqref{4.29:eq-path-length-decay-b} and $\varepsilon_2$ tend to zero with $\gamma$, uniformly in $k$. So both
inequalities in \eqref{4.29:eq-constants-order-b} are one-sided ceilings on $\gamma$, chosen after $L$ and $M$.

At the twist $\mathbf{a}=t\mathbf{H}_k(B')$ with $t\le1$, \eqref{4.29:eq-scale-thresholds-a} gives
$\ell_{m(c)}\sup_{A(c)}|\mathbf{a}|\le CB_3\delta_k^\sharp$. The first inequality then yields
$64\,\ell_{m(c)}\sup_{A(c)}|\mathbf{a}|\le\alpha_1^{(98)}$.

At the twist $\mathbf{a}=A$ of \eqref{4.29:eq-scale-thresholds-b}, the quantity
$\ell_{m(c)}\sup_{A(c)}|\mathbf{a}|$ is smaller than $\varepsilon_2$. The second inequality yields the same
conclusion.

On the circle $\mathbf{a}+\varsigma\mathbf{b}$ with
$|\varsigma|=\alpha_1^{(98)}/(64\,\ell_{m(c)}\sup_{A(c)}|\mathbf{b}|)$, we therefore have
\[
32\,\ell_{m(c)}\sup_{A(c)}|\mathbf{a}+\varsigma\mathbf{b}|
\le 32\Bigl(\frac{\alpha_1^{(98)}}{64}+\frac{\alpha_1^{(98)}}{64}\Bigr)=\alpha_1^{(98)} .
\]
Both the small-twist branch of Lemma~\ref{4.29:lem-collar-staircase} and this Cauchy circle therefore lie in the
region $32\,\ell_{m(c)}\sup_{A(c)}|\cdot|\le\alpha_1^{(98)}$. That region is contained in the polydisc on which the
source bound \eqref{4.29:eq-collar-staircase-b} holds. This is what makes the constant $c_Q=2c_6'$ and the
factor $2$ in $|DQ'|\le 2c_Q\,\ell_{m(c)}\sup|\mathbf{b}|$ available.

\emph{The ceiling \eqref{4.29:eq-constants-order-a}.} By \eqref{4.29:eq-near-field-scope-a}, the per-cube factor
at the cubes of Proposition~\ref{4.29:prop-near-field-scope} is bounded as follows:
\[
\frac{3C_NB_3\delta_k^\sharp}{\alpha_{2,k}}\le C^\sharp(\log x_k)^{-1}\le\frac{C^\sharp}{\log\gamma^{-2}} .
\]
These are the cubes that meet the region $P$ there and carry sources of the fluctuation field $A_k$ on
$\Omega_{k+1}\setminus\Lambda_{k+1}$. The last inequality holds because $x_k\ge\gamma^{-2}$.

Hence \eqref{4.29:eq-constants-order-a} gives $2^dL^4\cdot C(L,M,\mathfrak{p})\cdot C^\sharp(\log x_k)^{-1}\le1$.
So the per-cube factor absorbs, at these cubes, Ba{\l}aban's native count $2^dL^4$ and the constants that depend
on $(L,M,\mathfrak{p})$ only.

Moreover, $C(L,M,\mathfrak{p})$ is free upward, so we may assume without loss that $2^dL^4C(L,M,\mathfrak{p})\ge1$.
Then \eqref{4.29:eq-constants-order-a} gives $C^\sharp/\log\gamma^{-2}\le1$, and therefore
\[
3C_NB_3\delta_k^\sharp\le\alpha_{2,k}.
\]
The level-$k$ members of $H_k(B(t))$ on $P$ are at most $C_NB_3\delta_k^\sharp$, by
\eqref{4.29:eq-interpolated-field-bounds-a} together with \eqref{4.29:eq-scale-thresholds-a}. Hence they are at most
$\frac13\alpha_{2,k}$. That is, $H_k(B(t))$ satisfies \cite[(3.14)]{B-CMP109} with $\frac13\alpha_{2,k}$ on $P$,
whatever the position of the sources. For $j<k$ the bound with $\frac13\alpha_{2,j}$ follows from the pair bound in
\eqref{4.29:eq-interpolated-field-bounds-a}.

So the single condition \eqref{4.29:eq-constants-order-a} serves both purposes. It is one-sided, independent of
$k$, and chosen after $L$, $M_1$, $M$ and $\mathfrak{p}$.

\emph{The conditions with $\gamma^{12}$ and $\gamma^{10}$.} By \eqref{4.29:eq-far-field-paths-a}, the far field
carries the surplus $e^{-6R_{k+1}}\le\gamma^{12}$. By \eqref{4.29:eq-far-field-count-a}, the sum over generations
carries $2^dC_{\mathrm{ff}}e^{-5R_{k+1}}$ with $e^{-5R_{k+1}}\le\gamma^{10}$. The two conditions of (c) with
$\gamma^{12}$ and $\gamma^{10}$ make both surpluses absorb $C(L,M,\mathfrak{p})$.

\emph{The remaining ceilings.} The fourth line of \eqref{4.29:eq-constants-order-c} compares the source sizes of
\eqref{4.29:eq-path-length-decay-b} with $\varepsilon^{\mathrm{I}}_1$ and with $\varepsilon^{\mathrm{abs}}_1$. Since
$\varepsilon_k$ and $\delta_k$ tend to zero with $\gamma$ uniformly in $k$ under the coupling regime, it is a
ceiling on $\gamma$, uniform in $k$, chosen after $L$ and $M$.

In the fifth line, $\bar\varpi$ tends to zero with $\gamma$, uniformly in $k$. Indeed, by
\eqref{4.29:eq-scale-thresholds-a},
\[
R_{k_W}\alpha_{\iota,n}\le L\,C'(1+\beta_0)(\log x_n)^{r+q_\iota}x_n^{-1/2},\qquad\iota=0,1,\qquad x_n\ge\gamma^{-2},
\]
where $r$ is the exponent of \cite[(2.5)]{B-CMP119} and $C'$ is a constant independent of $k$ and of $\gamma$.
Moreover, by \eqref{4.29:eq-scale-thresholds-a}, with $k_0$ the first level of the region $W$, so that
$N_0=1+\log_LR_{k_0+1}$,
\[
N_0\le 2+r\log_L\log x_{k_0+1},\qquad x_{k_0+1}\le x_n+BN_0 ,
\]
so that $N_0^{\beta_0}(\log x_n)^{r+q_\iota}x_n^{-1/2}$ tends to $0$ as $\gamma\to0$, since $x_n\ge\gamma^{-2}$,
uniformly in $k$.
So that line is also a one-sided ceiling on $\gamma$.

Every condition of (c) other than the first, second and fourth items is an upper bound on $\gamma$. The first item
is a smallness condition on $\beta_0$ together with the lower bound $x_{k+1}\ge B+c_5$ on the inverse coupling, the
second bounds the couplings from above, and the fourth is a smallness condition on the primed radius. None of them
imposes an upper bound on $L$, $M$, $N$, $N_0$, $K$, $k$, $|W|$ or the arrest levels, and none involves the torus.
This proves the final sentence of the remark.
\end{proof}

\begin{definition}[Stratified regions and the stratified distance over free contours]
\label{4.30:def-stratified-distance}
Let $d\ge 2$ and $\eta>0$, and let $T$ be the lattice $\eta\mathbb{Z}^d$ or a torus of
$\eta\mathbb{Z}^d$, with the $\ell^1$-norm $|\cdot|$. Let $L>1$ be an odd integer, and put
$\ell_i:=L^i\eta$ for $i\in\mathbb{Z}_{\ge0}$; thus $\ell_i$ is the side of a cell at level $i$.

\emph{Strata and cells.} A \emph{stratification} $\mathfrak{S}$ of a region $T'\subset T$ is a
finite family of closed sets $S$, the \emph{strata}, with pairwise disjoint interiors and union
$T'$, such that each stratum
$S$ carries a level $i(S)\in\mathbb{Z}_{\ge0}$ and is a union of closed cubes of side $\ell_{i(S)}$ with
corners in $\ell_{i(S)}\mathbb{Z}^d$. These cubes are the \emph{cells} of $S$; we write
$\mathcal{C}(S)$ for the set of cells of $S$ and $\mathcal{C}:=\bigcup_{S}\mathcal{C}(S)$ for the
set of all cells. Each cell $c$ has a centre $\hat{c}$. Two strata are \emph{adjacent} if their
closures meet in a $(d-1)$-dimensional set, which is then called their \emph{interface}.
Interfaces are indexed by the pair of levels they separate, and not by connected component.
We write $T'^{\circ}$ for the interior of
$T'$.

\emph{Contours.} A \emph{contour} $\Gamma$ is a nearest-neighbour path of $T$; its length is
$\eta$ times its number of bonds. Each bond of $\Gamma$ lying in $T'$ is charged to exactly one
stratum whose closure contains it, chosen of highest level among these; for a stratum $S$ we write
$|\Gamma\cap\overline{S}|$ for $\eta$ times the number of bonds of $\Gamma$ charged to $S$.

\emph{Hypotheses.} Throughout, the following four hypotheses are in force:
\begin{equation}\label{4.30:eq-stratified-distance-a}
\mathfrak{S}\ \text{satisfies}\ (\mathrm{S}0),\ (\mathrm{S}1),\ (\mathrm{S}2),\ (\mathrm{S}3),
\end{equation}
where the conditions are these:
\begin{itemize}
\item[(S0)] (No pinches.) For every lattice point $p$ of $T'$, the closed $\eta$-cubes contained
in $T'$ that contain $p$ are connected to one another through chains of such cubes in which
consecutive cubes share a $(d-1)$-dimensional face, or else they all belong to one stratum.
\item[(S1)] (Unit step.) Adjacent strata have levels differing by exactly $1$.
\item[(S2)] (Width.) For a stratum $S$, let $\partial^{+}S$ be the part of $\partial S$ shared
with strata of level $i(S)+1$, and $\partial^{-}S$ the part of $\partial S$ shared with strata of
level $i(S)-1$; by (S1) and the cube-adjacency argument of the remark below,
\begin{equation*}
\partial S\cap T'^{\circ}\subset\partial^{+}S\cup\partial^{-}S .
\end{equation*}
A contour portion lying in $\overline{S}$ and running from a point of $\partial^{+}S$ to a point
of $\partial^{-}S$ is said to \emph{cross} $S$. Each stratum $S$ carries a number $W_S\ge 1$, the
\emph{width} of $S$ in its own units, and every contour portion $\Gamma$ crossing $S$ has length at
least $W_S\,\ell_{i(S)}$.
\item[(S3)] (Interfaces are coarse-cube faces.) The interface between adjacent strata $S_1$ of
level $i$ and $S_2$ of level $i-1$ is a union of faces of cubes of side $\ell_i$ with corners in
$\ell_i\mathbb{Z}^d$.
\end{itemize}

\emph{The stratified length and the stratified distance.} For a contour $\Gamma$, and for
lattice points $x,x'$ of $T'$, we put
\begin{equation}\label{4.30:eq-stratified-distance-b}
\begin{split}
\operatorname{len}_{\mathfrak{S}}(\Gamma)
&:=\sum_{S\in\mathfrak{S}}\ell_{i(S)}^{-1}\,|\Gamma\cap\overline{S}|,\\
d^{\wedge}(x,x')
&:=\inf\bigl\{\operatorname{len}_{\mathfrak{S}}(\Gamma):\ \Gamma\subset T'
\text{ a contour from } x \text{ to } x'\bigr\}.
\end{split}
\end{equation}
Thus $\operatorname{len}_{\mathfrak{S}}(\Gamma)$ is the length of $\Gamma$ measured stratum by
stratum in the unit of the stratum, with the convention that a bond lying on an interface is
charged at the coarser of the two units; which of the two units is used for such bonds is
immaterial in the uses of $d^{\wedge}$. The infimum defining $d^{\wedge}(x,x')$ is taken over all
contours in $T'$ from $x$ to $x'$, with no further restriction; we call them \emph{free
contours}, and $d^{\wedge}$ the \emph{stratified distance} over free contours.

Contours run between lattice points, and a cell centre is in general not a lattice point (for a
cell of level $i$ each coordinate of $\hat{c}$ differs from a corner coordinate by
$L^i\eta/2$, which is not a multiple of $\eta$ since $L$ is odd). Distances to cell centres are
therefore defined as follows: for a lattice point $x$ of $T'$ and a cell $c$, owned by the stratum
$S_c$,
\begin{equation*}
d^{\wedge}(x,\hat{c}):=\inf\Bigl\{\operatorname{len}_{\mathfrak{S}}(\Gamma)
+\frac{|p-\hat{c}|}{\ell_{i(S_c)}}:\ p \text{ a lattice point of } c,\ \Gamma\subset T'
\text{ a contour from } x \text{ to } p\Bigr\}.
\end{equation*}
Accordingly, a contour from $x$ to $\hat{c}$ means such a contour $\Gamma$ followed by the
segment $[p,\hat{c}]$, whose last stratum is $S_c$.

\emph{Rate.} A decay rate $a>0$ is fixed. For $b>0$ we put
\begin{equation*}
c_0(b):=\sum_{z\in\mathbb{Z}}e^{-b|z|}<2\bigl(1-e^{-b}\bigr)^{-1},
\end{equation*}
the bound being the one of \cite[p.~233]{B-CMP96}; indeed
$c_0(b)=(1+e^{-b})(1-e^{-b})^{-1}$ by summing the two geometric series, and $1+e^{-b}<2$.
\end{definition}

\begin{remark}[On the no-pinch condition]
Condition (S0) is automatic at every interior lattice point $p$ of $T'$. At such a point all
$2^d$ closed $\eta$-cubes containing $p$, one in each orthant at $p$, lie in $T'$. Label each of
them by the vector in $\{\pm1\}^d$ of the signs of its orthant; two of these cubes share a
$(d-1)$-dimensional face exactly when their sign vectors differ in one coordinate, so the
face-adjacency graph of the cubes at $p$ is the hypercube graph on $\{\pm1\}^d$, which is
connected. Hence (S0) imposes nothing when $T'=T$ or when $T'$ is a whole torus of
$\eta\mathbb{Z}^d$, and in general it is a condition on the lattice points of $\partial T'$ only.

Condition (S0) cannot be dropped. Let $d\ge3$, let $J\ge1$, and write $e_1,\dots,e_d$ for the
unit coordinate vectors of $\mathbb{R}^d$. Consider the two strata
\begin{equation*}
S':=[-\ell_J,\ell_J]\times[-\ell_J,0]^{d-1}\quad\text{(level $J$)},\qquad
S'':=[0,\ell_J]\times[0,\eta]^{d-1}\quad\text{(level $0$)},
\end{equation*}
and $T':=S'\cup S''$. Their closures meet in the segment $[0,\ell_J]\times\{0\}^{d-1}$, which for
$d\ge3$ is not $(d-1)$-dimensional; so $S'$ and $S''$ are not adjacent, the sets
$\partial^{\pm}S'$ and $\partial^{\pm}S''$ are empty, and (S1)--(S3), as well as the width
condition imposed in the cell-counting bound for sums over cells, hold vacuously. Yet,
for every $a>0$ and every odd $L>1$,
\begin{equation}\label{4.30:eq-stratified-distance-1}
\sum_{c\in\mathcal{C}}e^{-a\,d^{\wedge}(0,\hat{c})}\ \ge\ (L^J-1)\,e^{-a(2+d/2)}
\ \longrightarrow\ \infty\qquad(J\to\infty),
\end{equation}
the $L^J$ cells of $S''$ being reached from $0$ through the pinch at cost less than $2$ inside
$S'$. Indeed, for $1\le l<L^J$ consider the contour
\begin{equation*}
0\ \to\ -\eta e_2\ \to\ \eta e_1-\eta e_2\ \to\ \cdots\ \to\ l\eta e_1-\eta e_2\ \to\ l\eta e_1 .
\end{equation*}
It has $l+2$ bonds: one from $0$ to $-\eta e_2$, then $l$ in the direction $e_1$, then one from
$l\eta e_1-\eta e_2$ to $l\eta e_1$. Along all of them the first coordinate lies in
$[0,l\eta]\subset[-\ell_J,\ell_J]$, the second in $[-\eta,0]\subset[-\ell_J,0]$, and the others
vanish; so all bonds lie in $S'$, and none lies in $S''$, whose second coordinate ranges over
$[0,\eta]$ (the first and the last bond meet $S''$ only in an endpoint, $0$ and $l\eta e_1$
respectively). Its stratified length is therefore
computed in the unit $\ell_J$ of $S'$ alone:
\begin{equation*}
\operatorname{len}_{\mathfrak{S}}(\Gamma)\ \le\ \frac{(l+2)\eta}{\ell_J}
=\frac{l+2}{L^J}\ <\ 2,
\end{equation*}
since $l+2\le L^J+1<2L^J$. The contour ends at $l\eta e_1$, a lattice point of the $l$-th
$\eta$-cell of $S''$ (the cell $[(l-1)\eta,l\eta]\times[0,\eta]^{d-1}$), and the centre of that
cell lies at $\ell^1$-distance $d\eta/2$ from each of its corners, that is at distance $d/2$ in the
unit $\eta=\ell_0$ of $S''$. By the definition of $d^{\wedge}(0,\hat{c})$, applied with
$p=l\eta e_1$, this gives $d^{\wedge}(0,\hat{c})<2+d/2$ for each of these cells. Summing over the
$L^J-1$ cells so reached gives \eqref{4.30:eq-stratified-distance-1}.

Condition (S0) fails at every lattice point $p$ of $S'\cap S''=[0,\ell_J]\times\{0\}^{d-1}$. The
$\eta$-cubes of $T'$ present at $p$ are the cubes of $S'$, whose coordinates $2,\dots,d$ range
over $[-\eta,0]$, and the cube or cubes of $S''$ (one at the two ends of the segment, two at its
interior lattice points), whose coordinates $2,\dots,d$ range over $[0,\eta]$.
In the labelling by sign vectors these differ in the $d-1\ge2$ signs of coordinates
$2,\dots,d$, while every $\eta$-cube at $p$ whose signs in these coordinates are mixed lies
neither in $S'$ nor in $S''$. A chain of cubes in which consecutive cubes share a face changes one
sign at a time, so it cannot pass from a cube of $S'$ to a cube of $S''$ inside $T'$; and the
cubes at $p$ do not all belong to one stratum.
\end{remark}

\begin{remark}[Ba{\l}aban's counting lemma for admissible contours {\cite[Lemma~2.1]{B-CMP96}}]\label{4.30:rem-admissible-counting}
We state, in the notation of \cite[\S2]{B-CMP96}, the counting lemma proved there, together with
the structure of its proof. Let $\eta$ be the lattice spacing and $L$ the blocking factor. Let
$\Lambda_j$, $0\le j\le k$, be the lattices of \cite[\S2]{B-CMP96}, $B^j(\Lambda_j)$ the
region of \cite[\S2]{B-CMP96} in which lengths are measured in the unit $L^{j}\eta$, $\Sigma_j$
the interfaces between these regions, and $\mathfrak{B}$ the point set of \cite[(2.46)]{B-CMP96}
on which the distance is defined. For $y,y'\in\mathfrak{B}$ let $\Gamma_{y,y'}$ range over the
admissible contours from $y$ to $y'$ described in \cite[p.~231]{B-CMP96}. Here $|\cdot|$ denotes the
length of a contour, or of its intersection with a region, and, between two points, the distance
$|y-y'|$ of \cite[\S2]{B-CMP96}. By \cite[(2.46)]{B-CMP96},
\begin{equation}\label{4.30:eq-admissible-counting-1}
  d(y,y')=\inf_{\Gamma_{y,y'}}\sum_{j=0}^{k}(L^{j}\eta)^{-1}\,
  \bigl|\Gamma_{y,y'}\cap B^{j}(\Lambda_j)\bigr|,\qquad y,y'\in\mathfrak{B},
\end{equation}
the infimum being taken over all admissible contours. In this remark $d(\cdot,\cdot)$ is this
distance, and the letter $d$ standing alone is the dimension. Let $\delta_0$ be the constant of
\cite[\S2]{B-CMP96} and let $\alpha>0$; let $R$ and $M$ be the parameters entering
\cite[(2.57)]{B-CMP96} (here $R$ is the parameter of \cite{B-CMP96}, not the ball radius $R$ of
the glossary, Notation~\ref{0.2:not-glossary}), and put $c_0(b)=\sum_{z\in\mathbb{Z}}e^{-b|z|}$.

\cite[Lemma~2.1, (2.61)]{B-CMP96} states
\begin{equation}\label{4.30:eq-admissible-counting-2}
  \sup_{y\in\mathfrak{B}}\sum_{y'\in\mathfrak{B}}e^{-\alpha\delta_0d(y,y')}\le c_1(\alpha),
  \qquad c_1(\alpha)=12\,c_0^{d}(\tfrac12\alpha);
\end{equation}
its proof is carried out under the condition \cite[(2.59)]{B-CMP96},
\begin{equation*}
  \tfrac14\alpha\delta_0RM>2d\log c_0(\tfrac12\alpha)+1.
\end{equation*}
The proof is given in \cite[pp.~232--233]{B-CMP96}; we describe its structure. Write $j$ and $j'$
for the levels of the regions $B^{j}(\Lambda_j)$ and $B^{j'}(\Lambda_{j'})$ at $y$ and $y'$. A
contour $\Gamma_{y,y'}$ is decomposed at its successive first and last meetings with the
interfaces \cite[(2.47)]{B-CMP96}. This gives points $y_1,\dots,y_m$ and
$y'_1,\dots,y'_m$, where $y_l,y'_l$ are the entry and exit points at the $l$-th interface met;
$j_l$ is the index of the interface $\Sigma_{j_l}$ met there, so that
$y_l,y'_l\in\Sigma_{j_l}$; $j_{l,l+1}$ is the level of the region run through by the contour
between $y'_l$ and $y_{l+1}$; and $y_{m+1}=y'$. This decomposition yields the lower bound
\cite[(2.48)]{B-CMP96}:
\begin{align*}
  d(y,y')&\ge(L^{j}\eta)^{-1}|\Gamma_{y,y_1}|+\sum_{l=1}^{m}(L^{j_l}\eta)^{-1}|\Gamma_{y_l,y'_l}|
  +\sum_{l=1}^{m}(L^{j_{l,l+1}}\eta)^{-1}|\Gamma_{y'_l,y_{l+1}}|\\
  &\ge(L^{j}\eta)^{-1}|y-y_1|+\sum_{l}(L^{j_l}\eta)^{-1}|y_l-y'_l|
  +\sum_{l}(L^{j_{l,l+1}}\eta)^{-1}|y'_l-y_{l+1}|.
\end{align*}
For the segments between consecutive interfaces, \cite[(2.57)]{B-CMP96} gives
\begin{equation*}
  (L^{j_{l,l+1}}\eta)^{-1}|y'_l-y_{l+1}|>RM.
\end{equation*}
The sum over the crossing points is then carried out in \cite[(2.58)]{B-CMP96}. Under
\cite[(2.59)]{B-CMP96} the sum over $m$ is estimated in \cite[p.~233]{B-CMP96} by
\begin{equation*}
  c_0^{d}(\tfrac12\alpha)\,e^{-\frac14\alpha\delta_0RM\max\{|j-j'|-1,0\}}\sum_{m\ge0}e^{-m}
  \le6\,c_0^{d}(\tfrac12\alpha)\,e^{-|j-j'|}.
\end{equation*}
The summation over $j'$ then gives the constant $12\,c_0^{d}(\frac12\alpha)=c_1(\alpha)$.

In this argument the decomposition \cite[(2.47)]{B-CMP96}, the lower bound \cite[(2.48)]{B-CMP96} and
the bound \cite[(2.57)]{B-CMP96} use only that $\Gamma_{y,y'}$ is a contour from $y$ to $y'$. They
therefore hold for any contour. Admissibility is used only in the counting \cite[(2.58)]{B-CMP96},
where the crossing points $y_l,y'_l$ are points of the lattice $\Lambda_{j_l}$ of the interface they
lie on. The distance of Definition~\ref{4.30:def-stratified-distance} is instead an infimum over free
contours.
\end{remark}

\begin{definition}[Cubes and strata at a lattice point]\label{4.30:not-cubes-at-point}
Throughout this section $T$ is the lattice $\eta\mathbb{Z}^d$ (or a torus of it), $T'\subset T$ a
region, and $\mathfrak{S}$ a stratification of $T'$ as in
Definition~\ref{4.30:def-stratified-distance}, satisfying the hypotheses
\eqref{4.30:eq-stratified-distance-a}. From here on we make the standing assumptions $W_S\ge1$ for
every stratum $S$, levels $i(S)\ge0$, $L\ge3$ and $d\ge2$. An \emph{$\eta$-cube} is a closed cube
of side $\eta$ with corners in $\eta\mathbb{Z}^d$.

For a lattice point $p$ of $T'$ let $Q(p)$ be the set of $\eta$-cubes $q\subset T'$ with $p\in q$.
Each $q\in Q(p)$ lies in exactly one stratum, its \emph{home stratum}. Let $\operatorname{Str}(p)$
be the set of strata containing $p$; it coincides with the set of home strata of the cubes of
$Q(p)$.

We record two observations. Here $A$ is a finite union of $\eta$-cubes.
\begin{equation}\label{4.30:eq-cubes-at-point-a}
\begin{gathered}
\text{(o1)}\quad
\begin{aligned}[t]
&p\in S\cap S',\ S\neq S'\ \Longrightarrow\ p\in\partial S\cap\partial S';\\
&q,q'\ \eta\text{-cubes},\ q\subset S,\ q'\subset S',\ S\neq S',\\
&\qquad q,q'\ \text{share a }(d-1)\text{-dimensional face}\ f\ \Longrightarrow\
 S,S'\text{ adjacent},\\
&\qquad |i(S)-i(S')|=1,\ f\subset\partial^{+}S\ \text{if}\ i(S')=i(S)+1,\\
&\qquad f\subset\partial^{-}S\ \text{if}\ i(S')=i(S)-1;
\end{aligned}\\[4pt]
\text{(o2)}\quad
\begin{aligned}[t]
&[u,v]\ \text{a bond of}\ T,\ |u-v|=\eta,\ [u,v]\subset A\ \Longrightarrow\\
&\qquad [u,v]\subset q\ \text{for some}\ \eta\text{-cube}\ q\ \text{of}\ A;\\
&\text{hence, if}\ [u,v]\subset T':\ [u,v]\subset q\ \text{for some}\ q\in Q(u)\cap Q(v),\\
&\qquad\text{and every stratum containing}\ [u,v]\ \text{belongs to}\
 \operatorname{Str}(u)\cap\operatorname{Str}(v).
\end{aligned}
\end{gathered}
\end{equation}
\end{definition}

\begin{proof}
\emph{Home strata.} A cell of a stratum $S$ is a closed cube of side $\ell_{i(S)}=L^{i(S)}\eta$ with
corners in $\ell_{i(S)}\mathbb{Z}^d$. Since $i(S)\ge0$, $L^{i(S)}$ is an integer, so the cell is
the union of the $\eta$-cubes it contains. Hence every stratum, being the union of its cells, is a
union of $\eta$-cubes; distinct strata have disjoint interiors by
Definition~\ref{4.30:def-stratified-distance}.

Let $q\subset T'$ be an $\eta$-cube with centre $z$. The coordinates of $z$ lie in
$\eta(\mathbb{Z}+\tfrac12)$, and $q$ is the only $\eta$-cube containing $z$. Since $z\in T'$, $z$
lies in some stratum $S$, hence in an $\eta$-cube of $S$, which must be $q$; thus $q\subset S$. If
also $q\subset S'$ with $S'\neq S$, then the nonempty open set $\operatorname{int}q$ lies in
$\operatorname{int}S\cap\operatorname{int}S'$, which is impossible. So the home stratum exists and
is unique.

If $q\in Q(p)$, its home stratum contains $q\ni p$. Conversely, if $p\in S$, then $p$ lies in an
$\eta$-cube $q\subset S\subset T'$. This cube belongs to $Q(p)$, and its home stratum is $S$. Hence
$\operatorname{Str}(p)$ is the set of home strata of the cubes of $Q(p)$.

\emph{(o1).} Let $p\in S\cap S'$ with $S\neq S'$. The stratum $S'$ is a union of closed cubes, so
it is the closure of its interior, and every neighbourhood of $p$ meets $\operatorname{int}S'$.
Suppose some open neighbourhood $U$ of $p$ were contained in $S$. Then $U\subset\operatorname{int}S$,
and $U\cap\operatorname{int}S'$ would be a nonempty subset of
$\operatorname{int}S\cap\operatorname{int}S'$, a contradiction. Hence $p\notin\operatorname{int}S$.
As $S$ is closed, this gives $p\in\partial S$. Exchanging $S$ and $S'$ gives $p\in\partial S'$.

Now let $q\subset S$, $q'\subset S'$ with $S\neq S'$, and let $f$ be a common $(d-1)$-dimensional
face of $q$ and $q'$. The closed sets $S$ and $S'$ then meet in a set containing $f$, so they are
adjacent. By (S1) in \eqref{4.30:eq-stratified-distance-a}, $|i(S)-i(S')|=1$. The argument of the
first part uses no property of lattice points and so applies to every point of $S\cap S'$; hence
every point of $f\subset S\cap S'$ lies in $\partial S$, so $f\subset\partial S\cap S'$. By the
definition of $\partial^{\pm}S$ as the parts of $\partial S$ shared with strata of level
$i(S)\pm1$, we get $f\subset\partial^{+}S$ if $i(S')=i(S)+1$ and $f\subset\partial^{-}S$ if
$i(S')=i(S)-1$.

\emph{(o2).} After exchanging $u$ and $v$ if necessary, write $u=\eta n$ and $v=\eta(n+e_\mu)$ with
$n=(n_\nu)\in\mathbb{Z}^d$ and $e_\mu$ the $\mu$-th coordinate vector. An arbitrary $\eta$-cube is
\begin{equation*}
q=\prod_{\nu}[\eta m_\nu,\eta(m_\nu+1)],\qquad m=(m_\nu)\in\mathbb{Z}^d.
\end{equation*}
It meets the bond only if $m_\nu\in\{n_\nu-1,n_\nu\}$ for all $\nu\neq\mu$. In that case the
intersection is determined by the $\mu$-th coordinate. It is the whole bond if $m_\mu=n_\mu$,
$\{u\}$ if $m_\mu=n_\mu-1$, $\{v\}$ if $m_\mu=n_\mu+1$, and empty otherwise. Thus an $\eta$-cube
meets the bond in $\emptyset$, $\{u\}$, $\{v\}$ or the whole bond.

If $[u,v]\subset A$, the cubes of $A$ cover the midpoint of the bond. A cube of $A$ containing the
midpoint meets the bond in none of $\emptyset$, $\{u\}$, $\{v\}$, so it contains the whole bond.

Apply this with $A=T'$. Some $\eta$-cube $q\subset T'$ contains $[u,v]$, hence contains $u$ and
$v$, so $q\in Q(u)\cap Q(v)$.

Apply it next with $A$ a stratum $S$ containing the bond. Some $\eta$-cube $q\subset S$ contains
$[u,v]$, so $q\in Q(u)\cap Q(v)$ and its home stratum is $S$. Since $\operatorname{Str}(u)$ and
$\operatorname{Str}(v)$ are the sets of home strata of the cubes of $Q(u)$ and of $Q(v)$, we
conclude $S\in\operatorname{Str}(u)\cap\operatorname{Str}(v)$.
\end{proof}

\begin{lemma}[No lattice point on both boundaries of a stratum. Stated; proof incomplete
at the step marked $(\ast)$]\label{4.30:lem-no-two-sided-point}
Let $\mathfrak{S}$ be a stratification of the region $T'\subset T$ as in
Definition~\ref{4.30:def-stratified-distance}, and use the cubes and strata at a lattice point of
Definition~\ref{4.30:not-cubes-at-point}, with the standing conditions $d\ge 2$ and $L\ge 3$ recorded
there; let $\ell_i=L^{i}\eta$ be the
cell side at level $i$. There is no stratum $S\in\mathfrak{S}$ and no lattice point $p$
of $T'$ with both of the following properties:
\begin{itemize}
\item $p$ lies on an $\eta$-face $f^{+}$ shared by an $\eta$-cube of $S$ and an $\eta$-cube of a stratum
of level $i(S)+1$;
\item $p$ lies on an $\eta$-face $f^{-}$ shared by an $\eta$-cube of $S$ and an $\eta$-cube of a stratum
of level $i(S)-1$.
\end{itemize}
\end{lemma}

\begin{proof}
Suppose that $S$ and $p$ have both properties, and let $f^{+}$ and $f^{-}$ be the two faces. By the first
observation of Definition~\ref{4.30:not-cubes-at-point}, a $(d-1)$-dimensional face shared by an
$\eta$-cube of $S$ and an $\eta$-cube of a stratum $S'\neq S$ lies in $\partial^{+}S$ if
$i(S')=i(S)+1$ and in $\partial^{-}S$ if $i(S')=i(S)-1$. Hence
\begin{equation*}
f^{+}\subset\partial^{+}S,\qquad f^{-}\subset\partial^{-}S .
\end{equation*}

The lattice point $p$ lies on the $\eta$-face $f^{-}$, so it is a vertex of $f^{-}$. Since $d\ge 2$, the
face $f^{-}$ has dimension $d-1\ge 1$, and therefore it contains an edge $[p,p']$ issuing from $p$, with
$|p-p'|=\eta$. This edge satisfies
\begin{equation*}
[p,p']\subset f^{-}\subset q\subset S,
\end{equation*}
where $q$ is the $\eta$-cube of $S$ of which $f^{-}$ is a face. Thus the one-bond path $(p,p')$ is a
contour portion inside $S$ running from the point $p\in f^{+}\subset\partial^{+}S$ to the point
$p'\in f^{-}\subset\partial^{-}S$; that is, it crosses $S$. Its length is $\eta$. The width condition
in \eqref{4.30:eq-stratified-distance-a} of Definition~\ref{4.30:def-stratified-distance}, together
with $W_S\ge 1$, gives
\begin{equation*}
\eta\ \ge\ W_S\,\ell_{i(S)}\ \ge\ \ell_{i(S)} .
\end{equation*}

$(\ast)$ On the other hand, $S$ has a neighbour of level $i(S)-1$, namely the stratum whose $\eta$-cube
shares the face $f^{-}$ with an $\eta$-cube of $S$, and this level satisfies $i(S)-1\ge 0$. Hence
$i(S)\ge 1$, and since $\ell_i=L^{i}\eta$ and $L\ge 3$,
\begin{equation*}
\ell_{i(S)}\ \ge\ L\eta\ \ge\ 3\eta .
\end{equation*}

Combining the two displays gives $\eta\ge 3\eta$, which is impossible because $\eta>0$. This
contradiction proves the lemma.
\end{proof}

\begin{lemma}[Two consecutive levels at each interface point]\label{4.30:lem-interface-levels}
Let $p$ be a lattice point of $T'$, with $Q(p)$ and $\operatorname{Str}(p)$ as in
Definition~\ref{4.30:not-cubes-at-point}; the \emph{home stratum} of a cube of $Q(p)$ is the
unique stratum containing it. Put
\begin{equation*}
j:=\min\{\,i(S):\ S\in\operatorname{Str}(p)\,\}.
\end{equation*}
\begin{enumerate}
\item[(a)] Every $S\in\operatorname{Str}(p)$ has $i(S)\in\{j,j+1\}$.
\item[(b)] Suppose that $p$ lies in two or more strata; such a point is called an
\emph{interface point}. Then
\begin{equation*}
\{\,i(S):\ S\in\operatorname{Str}(p)\,\}=\{j,j+1\}=:K(p),
\end{equation*}
and this pair is called the \emph{kind} of $p$. Moreover, for every
$S\in\operatorname{Str}(p)$ there is an $\eta$-face $f\ni p$ shared by a cube of $Q(p)$
with home stratum $S$ and a cube of $Q(p)$ whose home stratum is some stratum $S'\neq S$,
where $S$ and $S'$ are adjacent and $\{i(S),i(S')\}=K(p)$.
\item[(c)] Consequently, let $\Gamma$ be a contour in $T'$ and $u$ a vertex of $\Gamma$
with $\operatorname{Str}(u)=\{S\}$. Then the first vertex $v$ of $\Gamma$ after $u$ with
$\operatorname{Str}(v)\neq\{S\}$, if there is one, is an interface point with
$S\in\operatorname{Str}(v)$, so that $i(S)\in K(v)$. In words: a contour leaves a stratum
only at lattice points where that stratum meets an adjacent one, and the level of the
stratum it leaves belongs to the kind of such a point.
\end{enumerate}
\end{lemma}

\begin{proof}
We call a \emph{chain} a finite sequence $q_0,\dots,q_n$ of cubes of $Q(p)$ in which any
two consecutive cubes $q_{s-1},q_s$ share a $(d-1)$-dimensional face. Since two distinct
closed $\eta$-cubes that share a $(d-1)$-dimensional face intersect exactly in that face,
and both contain $p$, the shared face contains $p$. By condition (S0) of
\eqref{4.30:eq-stratified-distance-a}, any two cubes of $Q(p)$ are joined by a chain.
Recall that $\operatorname{Str}(p)$ is the set of home strata of the cubes of $Q(p)$. If all
cubes of $Q(p)$ have one and the same home stratum, then $\operatorname{Str}(p)$ is a
singleton; in that case (a) holds trivially and the hypothesis of (b) fails, so there is
nothing to prove. We therefore use chains only where the cubes of $Q(p)$ have at least two
home strata.

Along a chain $q_0,\dots,q_n$ put $\lambda_s:=i(\text{home stratum of }q_s)$. We claim that
\begin{equation}\label{4.30:eq-interface-levels-1}
|\lambda_s-\lambda_{s-1}|\le 1\quad\text{for }1\le s\le n,\qquad
\lambda_s=\lambda_{s-1}\iff q_s,q_{s-1}\text{ have the same home stratum}.
\end{equation}
Indeed, if $q_{s-1}$ and $q_s$ have the same home stratum, then $\lambda_s=\lambda_{s-1}$.
If they have different home strata $S\neq S''$, then $q_{s-1}\subset S$ and
$q_s\subset S''$ share a
$(d-1)$-dimensional face, so by observation (o1) of \eqref{4.30:eq-cubes-at-point-a} the
strata $S$ and $S''$ are adjacent, and condition (S1) of
\eqref{4.30:eq-stratified-distance-a} gives $|i(S)-i(S'')|=1$; thus in this case the level
moves by exactly $1$. Both assertions of \eqref{4.30:eq-interface-levels-1} follow.

(a) Suppose, to the contrary, that some stratum $S_2\in\operatorname{Str}(p)$ has
$i(S_2)\ge j+2$. Take a chain $q_0,\dots,q_n$ from a cube $q_0$ whose home stratum has
level $j$ to a cube $q_n$ with home stratum $S_2$. Then $\lambda_0=j$ and
$\lambda_n\ge j+2$; since by
\eqref{4.30:eq-interface-levels-1} the levels change by at most $1$ per step, there is a
first index $s$ with $\lambda_s=j+2$. Since $\lambda_0=j$ and $0<s$, there is a last index
$r<s$ with $\lambda_r=j$. For every $u$ with $r<u<s$ we have $\lambda_u=j+1$: indeed
$\lambda_u\ge j$ by the minimality of $j$ (the home stratum of $q_u$ belongs to
$\operatorname{Str}(p)$), $\lambda_u\neq j$ by the choice of $r$, and $\lambda_u\le j+1$ by
the choice of $s$. This range of indices is nonempty, because the level passes from
$\lambda_r=j$ to $\lambda_s=j+2$ in steps of size at most $1$, so $s\ge r+2$. By
\eqref{4.30:eq-interface-levels-1}, consecutive equal levels mean equal home strata; hence
the cubes $q_{r+1},\dots,q_{s-1}$ all have one home stratum $S_1$, with $i(S_1)=j+1$.

Now $p$ lies on the face shared by $q_r$ and $q_{r+1}$; here $q_{r+1}$ is an $\eta$-cube of
$S_1$ and $q_r$ is an $\eta$-cube of a stratum of level $j=i(S_1)-1$, so by (o1) this face
lies in $\partial^{-}S_1$. Likewise $p$ lies on the face shared by $q_{s-1}$ and $q_s$,
where $q_{s-1}$ is an $\eta$-cube of $S_1$ and $q_s$ is an $\eta$-cube of a stratum of
level $j+2=i(S_1)+1$, so this face lies in $\partial^{+}S_1$. Thus $p$ lies both on an
$\eta$-face shared by an $\eta$-cube of $S_1$ and an $\eta$-cube of a stratum of level
$i(S_1)+1$, and on an $\eta$-face shared by an $\eta$-cube of $S_1$ and an $\eta$-cube of
a stratum of level $i(S_1)-1$. This contradicts
Lemma~\ref{4.30:lem-no-two-sided-point}, whose proof is incomplete at one step. Hence
every $S\in\operatorname{Str}(p)$ has $i(S)\in\{j,j+1\}$.

(b) Let $p$ be an interface point and $S\in\operatorname{Str}(p)$. Since $S$ is the home
stratum of a cube of $Q(p)$, there is a cube $q\in Q(p)$ with home stratum $S$; since
$|\operatorname{Str}(p)|\ge 2$, there is a cube $q''\in Q(p)$ whose home stratum is not $S$.
Take a chain $q=q_0,\dots,q_n=q''$, and let $q'=q_t$ be its first cube whose home stratum
is not $S$; it exists because the home stratum of $q''$ is not $S$, and $t\ge1$ because
the home stratum of $q_0$ is $S$.
The cube $q_{t-1}$ has home stratum $S$ and shares with $q'$ a $(d-1)$-dimensional face
$f$, which contains $p$ by the remark at the beginning of the proof. Let $S'\neq S$ be the
home stratum of $q'$. By (o1), $S$ and $S'$ are adjacent, and $|i(S)-i(S')|=1$. By (a), the
only levels present at $p$ are $j$ and $j+1$; two levels in $\{j,j+1\}$ differing by $1$
are $j$ and $j+1$, so $\{i(S),i(S')\}=\{j,j+1\}$. In particular both levels $j$ and $j+1$
occur among
$\{i(S''):S''\in\operatorname{Str}(p)\}$, which by (a) is contained in $\{j,j+1\}$; hence
this set equals $\{j,j+1\}=K(p)$, and the face $f$ and the stratum $S'$ have the properties
asserted in (b).

(c) Let $v$ be the first vertex of $\Gamma$ after $u$ with $\operatorname{Str}(v)\neq\{S\}$,
and let $v^{-}$ be the vertex of $\Gamma$ immediately before $v$. Either $v^{-}=u$, or
$v^{-}$ is a vertex after $u$ and before $v$; in both cases
$\operatorname{Str}(v^{-})=\{S\}$, in the second by the choice of $v$. The bond
$[v^{-},v]$ of $\Gamma$ is contained in $T'$, so by observation (o2) of
\eqref{4.30:eq-cubes-at-point-a} some stratum containing this bond belongs to
\begin{equation*}
\operatorname{Str}(v^{-})\cap\operatorname{Str}(v)=\{S\}\cap\operatorname{Str}(v).
\end{equation*}
This intersection is therefore nonempty, that is, $S\in\operatorname{Str}(v)$. Together
with $\operatorname{Str}(v)\neq\{S\}$ this forces $|\operatorname{Str}(v)|\ge 2$, so $v$ is
an interface point. Applying (b) at $v$ to the stratum $S\in\operatorname{Str}(v)$ gives
$i(S)\in K(v)$, together with an $\eta$-face through $v$ at which $S$ meets an adjacent
stratum $S'$ with $\{i(S),i(S')\}=K(v)$.
\end{proof}

\begin{lemma}[Interface points lie on coarse faces]\label{4.30:lem-coarse-faces}
Let $\mathfrak{S}$ be a stratification of the stratified region $T'\subset T$ as in
Definition~\ref{4.30:def-stratified-distance}, with cell sides $\ell_i=L^i\eta$. Let $p$ be an
interface point, in the sense of Lemma~\ref{4.30:lem-interface-levels}\,(b), of kind
$K(p)=\{i,i+1\}$. Then $p$ lies on a face $F$ of a cube of side $\ell_{i+1}$ with corners in
$\ell_{i+1}\mathbb{Z}^d$, and $F$ is contained in the interface of two adjacent strata of levels
$i+1$ and $i$. Consequently, for each $i'\in\{i,i+1\}$ there is a point
$z\in F\cap\ell_{i'}\mathbb{Z}^d$ with
\begin{equation}\label{4.30:eq-coarse-faces-1}
|p-z|\le \frac{(d-1)\,\ell_{i'}}{2}.
\end{equation}
The coordinate of $p$ normal to $F$ already lies in
$\ell_{i+1}\mathbb{Z}\subset\ell_{i'}\mathbb{Z}$; only the $d-1$ tangential coordinates move,
each by at most $\ell_{i'}/2$.
\end{lemma}

\begin{proof}
Throughout, an \emph{$\eta$-face} is a $(d-1)$-dimensional face of a closed cube of side $\eta$
of the lattice $T$, and measure on an affine hyperplane means $(d-1)$-dimensional Lebesgue
measure.

\emph{The $\eta$-face through $p$.} By clause (b) of Lemma~\ref{4.30:lem-interface-levels},
applied to the interface point $p$ of kind $\{i,i+1\}$, there is an $\eta$-face $f$ with
$p\in f$ which is shared by an $\eta$-cube owned by a stratum $S$ of level $i+1$ and an
$\eta$-cube owned by a stratum $S'$ of level $i$, the strata $S$ and $S'$ being adjacent.
The face $f$ therefore lies in the interface of $S$ and $S'$. By condition (S3) of the
stratification (see \eqref{4.30:eq-stratified-distance-a}), applied to the adjacent strata $S$
of level $i+1$ and $S'$ of level $i$, this interface is a union
\begin{equation*}
\bigcup_{G\in\mathcal{F}} G
\end{equation*}
of a family $\mathcal{F}$ of faces $G$ of cubes of side $\ell_{i+1}$ with corners in
$\ell_{i+1}\mathbb{Z}^d$. The faces of this grid form a countable, locally finite family.

\emph{The face $f$ is covered by the members of $\mathcal{F}$ in its own hyperplane.} Let $\Pi$
be the affine hyperplane containing $f$, and put
\begin{equation*}
C_\Pi:=\bigcup\{G\in\mathcal{F}:\ G\subset\Pi\},
\end{equation*}
a closed set, being the union of a locally finite family of closed faces. A face
$G\in\mathcal{F}$ not contained in $\Pi$ lies in an affine hyperplane different from $\Pi$, and
two distinct affine hyperplanes are disjoint or meet in an affine subspace of dimension $d-2$;
hence $G\cap\Pi$ has measure zero. Since $f\subset\Pi$ and $f$ is contained in the union of the
members of $\mathcal{F}$, the set $f\setminus C_\Pi$ is contained in the countable union of the
sets $G\cap\Pi$ with $G\not\subset\Pi$, and therefore has measure zero. On the other hand
$f\setminus C_\Pi$ is relatively open in $f$, $C_\Pi$ being closed; a non-empty relatively open
subset of the $(d-1)$-dimensional cube $f$ has positive measure. Hence $f\setminus C_\Pi$ is
empty, that is, $f\subset C_\Pi$.

\emph{Identification of the coarse face $F$.} Since $f$ is non-empty and $f\subset C_\Pi$, some
$G\in\mathcal{F}$ is contained in $\Pi$. Every face of a cube of side $\ell_{i+1}$ with corners
in $\ell_{i+1}\mathbb{Z}^d$ lies in a coordinate hyperplane of the form $\{x_\nu=c\}$ with
$c\in\ell_{i+1}\mathbb{Z}$; hence $\Pi=\{x_\nu=c\}$ for some coordinate direction $\nu$ and some
$c\in\ell_{i+1}\mathbb{Z}$. Inside $\Pi$ the members $G\in\mathcal{F}$ with $G\subset\Pi$ are
cubes of the $(d-1)$-dimensional grid of side $\ell_{i+1}$, while $f$ is a cube of the
$(d-1)$-dimensional grid of side $\eta$, which is finer, since $\ell_{i+1}=L^{i+1}\eta$ is an
integer multiple of $\eta$. Therefore $f$ lies in exactly one cube $F$ of the $\ell_{i+1}$-grid
of $\Pi$, and meets every other cube of that grid in a set of measure zero. If $F$ were not a
member of $\mathcal{F}$, then $f\subset C_\Pi$ would be covered by countably many cubes each
meeting $f$ in a null set, so $f$ would have measure zero; since $f$ has positive measure, $F$
is one of the members of $\mathcal{F}$. Hence $F$ is a face of a cube of side $\ell_{i+1}$ with
corners in $\ell_{i+1}\mathbb{Z}^d$, it is contained in the interface of the adjacent strata
$S$ (level $i+1$) and $S'$ (level $i$), and $p\in f\subset F$.

\emph{Rounding.} Write
\begin{equation*}
F=\{x_\nu=\sigma\,\ell_{i+1}\}\times\prod_{\mu\ne\nu}
\bigl[\sigma_\mu\ell_{i+1},\,(\sigma_\mu+1)\ell_{i+1}\bigr]
\end{equation*}
with integers $\sigma$ and $\sigma_\mu$ ($\mu\ne\nu$). Fix $i'\in\{i,i+1\}$; then
$\ell_{i+1}/\ell_{i'}\in\{1,L\}$, so $\ell_{i'}$ divides $\ell_{i+1}$. The normal coordinate
$p_\nu=\sigma\ell_{i+1}$ therefore lies in $\ell_{i'}\mathbb{Z}$. For $\mu\ne\nu$ consider the
points
\begin{equation*}
\sigma_\mu\ell_{i+1}+e\,\ell_{i'},\qquad 0\le e\le \ell_{i+1}/\ell_{i'},\quad e\in\mathbb{Z}.
\end{equation*}
They lie in $\ell_{i'}\mathbb{Z}$ and in the $\mu$-th interval
$[\sigma_\mu\ell_{i+1},(\sigma_\mu+1)\ell_{i+1}]$ of $F$, they have spacing $\ell_{i'}$, and they
include both end points of that interval ($e=0$ and $e=\ell_{i+1}/\ell_{i'}$). Since
$p_\mu$ lies in this interval, one of these points, call it $z_\mu$, satisfies
$|p_\mu-z_\mu|\le\ell_{i'}/2$. Let $z$ be the point with $z_\nu:=p_\nu$ and with the coordinates
$z_\mu$ just chosen for $\mu\ne\nu$. Then every coordinate of $z$ lies in $\ell_{i'}\mathbb{Z}$,
and $z_\nu=\sigma\ell_{i+1}$ together with $z_\mu\in[\sigma_\mu\ell_{i+1},(\sigma_\mu+1)\ell_{i+1}]$
gives $z\in F$; thus $z\in F\cap\ell_{i'}\mathbb{Z}^d$. Finally, by the triangle inequality along
the coordinate directions, the coordinate unit vectors having length one,
\begin{equation*}
|p-z|\le\sum_{\mu\ne\nu}|p_\mu-z_\mu|\le\frac{(d-1)\,\ell_{i'}}{2},
\end{equation*}
which is \eqref{4.30:eq-coarse-faces-1}.
\end{proof}

\begin{lemma}[The hyperbolic cotangent comparison]\label{4.30:lem-coth-comparison}
For $b>0$ put
\begin{equation*}
c_0(b):=\sum_{z\in\mathbb{Z}}e^{-b|z|},\qquad t:=\tanh(b/4),\qquad
\theta_b:=\frac{1+t^2}{2}\in\Big(\frac12,1\Big).
\end{equation*}
Let $d\ge 1$ be an integer and let $L\ge 2$ be real. Then
\begin{equation}\label{4.30:eq-coth-comparison-a}
\begin{aligned}
&\text{(T1)}\quad c_0(b)=\theta_b\,c_0(b/2);\\
&\text{(T2)}\quad c_0(b/L)\le \frac{L}{2}\,c_0(b/2);\\
&\text{(T3)}\quad d\,c_0(b)\,c_0(b/L)^{d-1}
\le \frac{d}{2^{d-1}}\,\theta_b\,L^{d-1}c_0(b/2)^d\le L^{d-1}c_0(b/2)^d .
\end{aligned}
\end{equation}
\end{lemma}

\begin{proof}
Since $b>0$, we have $t\in(0,1)$, hence $t^2\in(0,1)$ and $\theta_b\in(\frac12,1)$.

We first identify $c_0$. Splitting the sum into $z=0$ and $\pm z$ with $z\ge1$ and summing the
geometric series, we obtain
\begin{equation*}
c_0(b)=1+2\sum_{z\ge1}e^{-bz}=1+\frac{2e^{-b}}{1-e^{-b}}
=\frac{1+e^{-b}}{1-e^{-b}}=\coth(b/2),
\end{equation*}
the last equality after multiplying numerator and denominator by $e^{b/2}$. In particular
$c_0>0$ on $(0,\infty)$.

\emph{Proof of} (T1). Put $u:=b/4$, so that $t=\tanh u$. By the identity for $c_0$,
$c_0(b)=\coth 2u$ and $c_0(b/2)=\coth u$. The double-angle formula gives
$\tanh 2u=2t/(1+t^2)$. Therefore
\begin{equation*}
\frac{c_0(b)}{c_0(b/2)}=\frac{\coth 2u}{\coth u}=\frac{\tanh u}{\tanh 2u}
=t\cdot\frac{1+t^2}{2t}=\frac{1+t^2}{2}=\theta_b .
\end{equation*}

\emph{Proof of} (T2). The function $\tanh$ is concave on $[0,\infty)$, since its second
derivative $-2\tanh(v)\operatorname{sech}^2(v)$ is non-positive there, and $\tanh 0=0$.
Hence, for $v\ge0$ and $\lambda\in[0,1]$, concavity applied to the points $0$ and $v$ gives
\begin{equation*}
\tanh(\lambda v)\ge \lambda\tanh v+(1-\lambda)\tanh 0=\lambda\tanh v .
\end{equation*}
Take $v:=b/4$ and $\lambda:=2/L$, which lies in $(0,1]$ because $L\ge2$. Then
$\tanh(b/(2L))\ge (2/L)\tanh(b/4)>0$, and passing to reciprocals,
\begin{equation*}
\coth\big(b/(2L)\big)\le \frac{L}{2}\coth(b/4).
\end{equation*}
By the identity for $c_0$, the left side is $c_0(b/L)$ and $\coth(b/4)=c_0(b/2)$; this is (T2).

\emph{Proof of} (T3). All quantities involved are positive, so (T2) may be raised to the
power $d-1\ge0$ and multiplied by (T1):
\begin{equation*}
c_0(b)\,c_0(b/L)^{d-1}\le \theta_b\,c_0(b/2)\Big(\frac{L}{2}\Big)^{d-1}c_0(b/2)^{d-1}
=\frac{1}{2^{d-1}}\,\theta_b\,L^{d-1}c_0(b/2)^d .
\end{equation*}
Multiplying by $d$ gives the first inequality of (T3). For the second, $d\le 2^{d-1}$ for every
integer $d\ge1$ (by induction: equality at $d=1$, and $d+1\le 2d\le 2^{d}$ from $d\le 2^{d-1}$),
and $\theta_b<1$; hence $(d/2^{d-1})\,\theta_b\le1$, and the second inequality follows.
\end{proof}

For the decay rate $a>0$ fixed in this section we write $\theta:=\theta_a$; by (T1) of
\eqref{4.30:eq-coth-comparison-a}, $\theta=c_0(a)/c_0(a/2)$.

\begin{remark}\label{4.30:prop-cell-counting}
The statement of this proposition is not written out here; parts of its proof are printed below.
\end{remark}

\begin{proof}[Proof of Proposition~\ref{4.30:prop-cell-counting}, continued]
\label{4.30:prf-crossing-rounding}
\emph{Rounding the crossing points to the appropriate lattices.} We keep the notation of the
preceding part of the proof: $\Gamma$ is a free contour from $x$ to the centre $\hat{c}$ of a
cell, $\Sigma^{1},\dots,\Sigma^{m}$ are the successive interface kinds met by $\Gamma$, and
$y_l,y'_l$ are its entry and exit points at $\Sigma^{l}$. The stratum $S_{l-1}$ is the one left at
$\Sigma^{l}$, and $S_l$ is the one run through next. We write $\ell_S:=\ell_{i(S)}$, and
$\ell(\Sigma^{l})^{+}$ for the coarser of the two units at $\Sigma^{l}$. The lower bound for $\operatorname{len}_{\mathfrak{S}}(\Gamma)$ is in force.

\emph{Entry points.} For $2\le l\le m$ let $z_l$ be a nearest point to $y_l$ of
$\Sigma^{l}\cap\ell_{S_{l-1}}\mathbb{Z}^d$, the lattice of the unit of the stratum just left.
Exactly two terms of the corresponding display contain $y_l$:
\begin{equation*}
  \tfrac12\,|y'_{l-1}-y_l|/\ell_{S_{l-1}}
  \quad\text{and}\quad
  |y_l-y'_l|/\ell(\Sigma^{l})^{+} .
\end{equation*}
Both are measured in units at least $\ell_{S_{l-1}}$, because $S_{l-1}$ is one of the two strata
adjacent at $\Sigma^{l}$; moreover, by conditions (S3) and (S1) of
\eqref{4.30:eq-stratified-distance-a}, $\Sigma^{l}$ is a union of faces of cubes of side
$\ell_{S_{l-1}}$ or coarser. These two facts are the reason for rounding $y_l$ to the lattice
$\ell_{S_{l-1}}\mathbb{Z}^d$.

\emph{Exit points.} For $1\le l\le m$ let $z'_l$ be a nearest point to $y'_l$ of
$\Sigma^{l}\cap\ell_{S_l}\mathbb{Z}^d$, the lattice of the stratum about to be run through. The
terms of the corresponding display containing $y'_l$ are the following. The first is
$|y_l-y'_l|/\ell(\Sigma^{l})^{+}$; for $l=1$ it is the merged term
$|x-y'_1|/\ell(\Sigma^{1})^{+}$. The second is $\tfrac12|y'_l-y_{l+1}|/\ell_{S_l}$ if $l\le m-1$,
and $|y'_m-\hat{c}|/\ell_{S_m}$ if $l=m$. All of these are measured in units at least
$\ell_{S_l}$.

\emph{Size of each displacement.} Each of the points $y_l,y'_l$ lies on a face of a cube of the
coarser unit at its interface, by (S3). Hence its coordinate normal to that face already lies on
both lattices, and only the $d-1$ tangential coordinates move. By Lemma~\ref{4.30:lem-coarse-faces},
applied with $i'$ equal to the level of $S_{l-1}$, resp.\ of $S_l$, there is a point of the
relevant lattice on that face at distance at most $(d-1)\ell_{i'}/2$, that is, at most $(d-1)/2$
in the relevant unit. The nearest
point chosen above is no farther away. Measured in the unit of any term that contains it, the
displacement is therefore at most $(d-1)/2$. Each rounded point enters at most two terms, each
with coefficient at most one. By the triangle inequality, replacing the point by its rounding
changes the right side of the corresponding display by at most $d-1$.

\emph{Counting the rounded points.} The rounded points are $y'_1$, and $y_l$ and $y'_l$ for
$2\le l\le m$: $2m-1$ points in all. The point $y_1$ does not occur, since the first term of
the corresponding display is based at $x$ itself. Hence
\begin{equation}\label{4.30:eq-crossing-rounding-1}
\begin{split}
  \operatorname{len}_{\mathfrak{S}}(\Gamma) \ge{}&
  \frac{|x-z'_1|}{\ell(\Sigma^{1})^{+}}
  +\sum_{l=2}^{m}\frac{|z_l-z'_l|}{\ell(\Sigma^{l})^{+}}
  +\sum_{l=1}^{m-1}\Bigl[\tfrac12 W_{S_l}
  +\tfrac12\,\frac{|z'_l-z_{l+1}|}{\ell_{S_l}}\Bigr]\\
  &+\frac{|z'_m-\hat{c}|}{\ell_{S_m}}-(d-1)(2m-1).
\end{split}
\end{equation}
Since $(d-1)(2m-1)\le 2dm$, passing from $e^{-a\operatorname{len}_{\mathfrak{S}}(\Gamma)}$ to the
exponential of $-a$ times the rounded right side costs a factor at most
\begin{equation}\label{4.30:eq-crossing-rounding-2}
  e^{a(d-1)(2m-1)}\le \bigl(e^{2ad}\bigr)^{m}.
\end{equation}
Each rounded point costs a factor at most $e^{a(d-1)}\le e^{ad}$, and a crossing carries at most
two rounded points; so in the comparison sum $R_m$ of \eqref{4.30:eq-crossing-histories-b} this
factor is counted as $e^{2ad}$ per crossing. The exact count $(d-1)(2m-1)$ is used only once, in
the comparison $T_m\le R_m/3$ of the summation that follows, where $T_m$ denotes the contribution
of the contours that meet $m$ interface kinds.

\emph{Consequences of the rounding rule.} With this rule there is no entry sum at all, because the
first term of the corresponding display is based at $x$ itself. Every crossing sum of the
summation that follows, that is, every sum over a rounded entry point $z_l$, runs over the lattice
$\ell_{S_{l-1}}\mathbb{Z}^d$, which is no finer than the unit $\ell_{S_{l-1}}$ of its cost, and so
carries no lateral factor. The lateral factor $L^{d-1}$ appears only in
the exit sums over the points $z'_l$ that lie on the lattice of a finer stratum. There is one such
factor for each passage to a finer side, and the summation that follows uses exactly this.
\end{proof}

\begin{proof}[Proof of Proposition~\ref{4.30:prop-cell-counting}, continued: counting crossing
histories against a geometric sum]
\label{4.30:prf-crossing-histories}
The proof of Proposition~\ref{4.30:prop-cell-counting} is incomplete at the step marked
$(\ast)$, which is not among the steps below. These steps carry out the count of
\cite[(2.58), p.~233]{B-CMP96}, adapted to free contours.

\emph{Notation.} We put
\begin{equation*}
\theta:=\frac{c_0(a)}{c_0(a/2)} .
\end{equation*}
Since $c_0(b)=\sum_{z\in\mathbb{Z}}e^{-b|z|}=\coth(b/2)$, clause (T1) of the hyperbolic cotangent
comparison (Lemma~\ref{4.30:lem-coth-comparison}, display~\eqref{4.30:eq-coth-comparison-a}),
applied with $b=a$, gives $\theta=\theta_a=(1+\tanh^2(a/4))/2$, and in particular
$\theta\in(\tfrac12,1)$.

For a cell $c$ we use the objects attached to the pair $(x,\hat{c})$ in the preceding parts of
this proof:
\begin{itemize}
\item the free contour $\Gamma$ in $T'$ from $x$ to $\hat{c}$ and the tolerance $\varepsilon>0$
with which it was chosen;
\item the number $m\ge0$ of interface kinds that $\Gamma$ meets, and these kinds
$\Sigma^{1},\dots,\Sigma^{m}$ in the order in which they are met;
\item the stratum $S_0$ containing $x$, of level $j$, and, for $1\le l\le m$, the stratum $S_l$
into which $\Gamma$ runs at $\Sigma^{l}$, of level $i_l$;
\item the entry and exit points $y_l,y'_l$ at $\Sigma^{l}$ and their roundings $z_l,z'_l$.
\end{itemize}
In this part of the proof the letter $m$ is used only in this sense. For a stratum $S$ we write
$\ell_S:=\ell_{i(S)}$ for its cell side, and for an interface kind $\Sigma$ we write
$\ell(\Sigma)^{+}$ for the coarser of the two units present at $\Sigma$.

\emph{The summation.} Let $\Phi_{\Gamma}$ denote the right-hand side
of the corresponding display for the contour $\Gamma$. By the
corresponding bound,
\begin{equation*}
e^{-a\,d^{\wedge}(x,\hat{c})}\le e^{a\varepsilon}\,e^{-a\,\Phi_{\Gamma}} .
\end{equation*}
We sum this over
the cells $c$ by summing, in turn, over the following data:
\begin{itemize}
\item the number $m\ge0$;
\item the first kind $\Sigma^{1}$, for which there are at most two choices given $x$, namely the
kinds containing the level $j$ of $S_0$;
\item for each $l\ge1$, the side $S_l$ of $\Sigma^{l}$ into which $\Gamma$ runs.
\end{itemize}
For the last item there are two choices, namely its level $i_l\in\Sigma^{l}$. This choice fixes
$\Sigma^{l+1}$ as the other kind containing $i_l$, by the construction of the crossing points in
the first part of this proof, which is used here as stated there. Only the level of $S_l$ and its
width factor enter below, never
the identity of $S_l$. We then sum over the rounded points $z'_1$ and $z_l,z'_l$
($2\le l\le m$), and finally over $c\in\mathcal{C}(S_m)$.

The right-hand side of the corresponding display is a sum of nonnegative terms, and we use
$e^{-a(u+v)}=e^{-au}e^{-av}$ term by term. Each stratum $S_l$ that is crossed in full
contributes the width factor $e^{-(a/2)W_{S_l}}$, which is bounded by
\begin{equation*}
\omega:=\max_{S}e^{-(a/2)W_S}.
\end{equation*}

\emph{The point sums.} Each of the point sums below is bounded uniformly in its base point.

(i) There is no entry sum. The approach leg was merged into the first lateral term in the first
part of this proof, and the exit sum (ii) at $l=1$ has base point $z_1:=x$.

(ii) For $l\ge1$, the lateral sum is
\begin{equation*}
\sum_{z'_l\in\Sigma^{l}\cap\ell_{S_l}\mathbb{Z}^d}e^{-a|z_l-z'_l|/\ell(\Sigma^{l})^{+}} .
\end{equation*}
There are two cases.
\begin{itemize}
\item If $S_l$ is the coarse side of $\Sigma^{l}$, then $\ell_{S_l}=\ell(\Sigma^{l})^{+}$, and
the sum is at most $c_0(a)^d$.
\item If $S_l$ is the fine side, then $\ell_{S_l}=\ell(\Sigma^{l})^{+}/L$, and the sum runs over
the finer lattice at the coarse cost.
\end{itemize}
In the second case we argue as follows. By condition (S3) of~\eqref{4.30:eq-stratified-distance-a},
the points of $\Sigma^{l}\cap\ell_{S_l}\mathbb{Z}^d$ lie on faces of cubes of side
$\ell(\Sigma^{l})^{+}$, and these faces have $d$ orientations. For the points on faces of one
orientation the sum factorises:
\begin{itemize}
\item the coordinate normal to the face runs over the coarse lattice at the coarse cost, and
contributes at most $c_0(a)$;
\item each of the $d-1$ tangential coordinates runs over the fine lattice at the coarse cost,
and contributes at most
\begin{equation*}
\sum_{n\in\mathbb{Z}}e^{-(a/L)|t-n|}\le c_0(a/L),
\end{equation*}
where $t$ is the corresponding coordinate of the base point in units of $\ell_{S_l}$.
\end{itemize}
Summing over the $d$ orientations, the lateral sum is bounded by
\begin{equation}\label{4.30:eq-crossing-histories-1}
d\,c_0(a)\,c_0(a/L)^{d-1}\le\frac{d}{2^{d-1}}\,\theta\,L^{d-1}c_0(a/2)^d\le L^{d-1}c_0(a/2)^d .
\end{equation}
This is clause (T3) of Lemma~\ref{4.30:lem-coth-comparison} with $b=a$. It rests on three facts:
\begin{itemize}
\item $c_0(a)=\theta\,c_0(a/2)$, which is (T1);
\item the bound
\begin{equation*}
c_0(a/L)=\coth\bigl(a/(2L)\bigr)\le\frac{L}{2}\coth(a/4)=\frac{L}{2}\,c_0(a/2),
\end{equation*}
which is (T2); it holds for $L\ge2$ because $\tanh$ is concave on $\mathbb{R}_+$ with
$\tanh0=0$, so that $\tanh(2v/L)\ge(2/L)\tanh v$;
\item $d\le2^{d-1}$ and $\theta<1$.
\end{itemize}
The factor $d$ in~\eqref{4.30:eq-crossing-histories-1} cannot be dropped, since all $d$
orientations of faces occur at an interface that is not flat; the comparison (T3) absorbs it.

Thus the bound $L^{d-1}c_0(a/2)^d$ carries the lateral factor $L^{d-1}$. It arises because the
contour ran along
$\Sigma^{l}$ on the coarse side, at cost $1/L$ per fine unit, and may enter the finer stratum at
any point of $\Sigma^{l}$. Summed over the two sides of $\Sigma^{l}$, the lateral sum is at most
\begin{equation*}
B:=c_0(a)^d+\frac{d}{2^{d-1}}L^{d-1}c_0(a/2)^d\le\frac43\,L^{d-1}c_0(a/2)^d .
\end{equation*}
For the last inequality, note first that $L\ge3$, since the blocking factor $L$ is an odd integer
greater than one (Definition~\ref{1.1:def-lattice}); in particular (T2) applies above. Next,
$c_0(a)\le c_0(a/2)$ since $\theta<1$, so
\begin{equation*}
B\le c_0(a/2)^d\Bigl[1+\frac{d}{2^{d-1}}L^{d-1}\Bigr].
\end{equation*}
For $d=2$ the bracket is $1+L$, and $1+L\le 4L/3$ holds exactly when $L\ge3$. For $d\ge3$ and
$L\ge3$ we divide by $L^{d-1}$ and use
\begin{equation*}
L^{1-d}+\frac{d}{2^{d-1}}\le\frac19+\frac34<\frac43 .
\end{equation*}
At $l=1$ we include the at most two choices of $\Sigma^{1}$, and the first step contributes at
most
\begin{equation*}
2B\le\frac83\,L^{d-1}c_0(a/2)^d .
\end{equation*}

(iii) For $1\le l\le m-1$, the crossing sum is
\begin{equation*}
\sum_{z_{l+1}}e^{-(a/2)|z'_l-z_{l+1}|/\ell_{S_l}}\le c_0(a/2)^d .
\end{equation*}
Here $z_{l+1}$ runs over $\Sigma^{l+1}\cap\ell_{S_l}\mathbb{Z}^d$, a lattice no finer than the
unit $\ell_{S_l}$ of the cost. This is the lattice of the rounding of the crossing points, and it
is available on $\Sigma^{l+1}$ because both interfaces of $S_l$ are unions of faces of
$\ell_{S_l}$-cubes or of coarser cubes, by conditions (S1) and (S3)
of~\eqref{4.30:eq-stratified-distance-a}.

(iv) The final cell sum satisfies
\begin{equation*}
\sum_{c\in\mathcal{C}(S_m)}e^{-a|z'_m-\hat{c}|/\ell_{S_m}}\le c_0(a)^d\,e^{ad/2}
=\theta^d c_0(a/2)^d\,e^{ad/2} .
\end{equation*}
Indeed, the centres $\hat{c}$ are half-integer points in the unit $\ell_{S_m}$. In each
coordinate the triangle inequality therefore costs at most a factor $e^{a/2}$ when the sum is
compared with a sum over the integer lattice.

\emph{Collecting.} We sum in the following order:
\begin{equation*}
\hat{c},\ z'_m,\ z_m,\ \text{side of }\Sigma^{m},\ z'_{m-1},\ \dots,\ z_2,\
\text{side of }\Sigma^{2},\ z'_1,\ \text{side of }\Sigma^{1},\ \Sigma^{1}.
\end{equation*}
We include the rounding cost $e^{a(d-1)(2m-1)}$ from the rounding of the crossing points to the
appropriate lattices earlier in this proof, used here as stated there. Let $T_0$ be the term in
which no
interface is met, so that $c\in\mathcal{C}(S_0)$, and for $m\ge1$ let $T_m$ collect the cells
reached after $m$ interface kinds have been met. These terms obey
\begin{equation}\label{4.30:eq-crossing-histories-a}
\begin{aligned}
T_0&\le e^{ad/2}c_0(a)^d\le e^{ad/2}c_0(a/2)^d,\\
T_m&\le 2B\,\bigl[c_0(a/2)^d\,B\,\omega\bigr]^{m-1}\,e^{ad/2}\theta^d c_0(a/2)^d\,
e^{a(d-1)(2m-1)}\qquad(m\ge1).
\end{aligned}
\end{equation}
In the bound for $T_m$, each later step $l\ge2$ contributes three factors:
\begin{itemize}
\item one crossing sum, at most $c_0(a/2)^d$;
\item one two-sided exit bound $B$;
\item one width factor $\omega$.
\end{itemize}
The width factor is that of $S_{l-1}$, which $\Gamma$ crossed in full between $\Sigma^{l-1}$ and
$\Sigma^{l}$; by the construction of the crossing points in the first part of this proof, these
two kinds are the two sides of $S_{l-1}$. The first step carries the bare factor $2B$, which is
at most $\tfrac83L^{d-1}c_0(a/2)^d$.

This is the content of the following observation.

(${\alpha}$) At most one passage to a finer side can occur without an intervening full crossing,
namely the first one ($l=1$). Every later passage to a finer side at an interface $\Sigma^{l}$
($l\ge2$) is preceded by the full crossing of $S_{l-1}$. The width factor
$e^{-(a/2)W_{S_{l-1}}}$ of that crossing is therefore available, and it pays for the factor
$L^{d-1}$ of that step through $(\mathrm{W}^{b})$, the corresponding display.

(${\beta}$) Precisely, put
\begin{equation}\label{4.30:eq-crossing-histories-b}
R_m:=2L^{d-1}c_0(a/2)^{2d}e^{2ad}\cdot
\bigl[2L^{d-1}c_0(a/2)^{2d}e^{2ad}\,\omega\bigr]^{m-1}\qquad(m\ge1).
\end{equation}
Each step carries the entropy factor $2c_0(a/2)^{2d}e^{2ad}$ and the lateral factor $L^{d-1}$,
and each step after the first carries in addition one width factor.

We compare $T_m$ with $R_m$. Using $B\le\tfrac43L^{d-1}c_0(a/2)^d$ in the $m$ factors $B$ of
\eqref{4.30:eq-crossing-histories-a}, the powers of $L$ and of $c_0(a/2)$ and the factor
$\omega^{m-1}$ cancel against those of $R_m$. Since $2(4/3)^m/2^m=\tfrac43(2/3)^{m-1}$, this
gives
\begin{equation}\label{4.30:eq-crossing-histories-2}
\begin{aligned}
\frac{T_m}{R_m}
&\le\frac43\Bigl(\frac23\Bigr)^{m-1}\theta^d\exp\Bigl\{a\Bigl[\frac d2+(2m-1)(d-1)-2dm\Bigr]\Bigr\}\\
&=\frac43\Bigl(\frac23\Bigr)^{m-1}\theta^d\,e^{-a(2m+d/2-1)}\\
&\le\frac43\,\theta^2e^{-2a}\le\frac13 .
\end{aligned}
\end{equation}
The equality holds because $(2m-1)(d-1)=2md-2m-d+1$. The first inequality in the last line holds
for $m\ge1$ and $d\ge2$, for three reasons: $(2/3)^{m-1}\le1$; $2m+d/2-1\ge2$; and
$\theta^d\le\theta^2$, since $\theta<1$. For the last inequality, note the equivalences
\begin{equation*}
\frac43\,\theta^2e^{-2a}\le\frac13
\iff 2\theta e^{-a}\le1
\iff \bigl(1+\tanh^2(a/4)\bigr)e^{-a}\le1 .
\end{equation*}
The last condition holds, because $\tanh u\le u$ for $u\ge0$ gives
\begin{equation*}
1+\tanh^2(a/4)\le1+\frac{a^2}{16}\le e^{a}.
\end{equation*}

\emph{The geometric sum.} Taking the maximum over the strata $S$ in $(\mathrm{W}^{b})$,
the corresponding display, gives
\begin{equation*}
2L^{d-1}c_0(a/2)^{2d}e^{2ad}\,\omega\le e^{-2}.
\end{equation*}
Inserting this into~\eqref{4.30:eq-crossing-histories-b} yields
\begin{equation*}
R_m\le2L^{d-1}c_0(a/2)^{2d}e^{2ad}\,e^{-2(m-1)} .
\end{equation*}
By~\eqref{4.30:eq-crossing-histories-2} we have $T_m\le R_m/3\le R_m$ for every $m\ge1$, and
hence
\begin{align*}
\sum_{m\ge1}T_m\le\sum_{m\ge1}R_m
&\le2L^{d-1}c_0(a/2)^{2d}e^{2ad}\sum_{m\ge1}e^{-2(m-1)}\\
&\le2L^{d-1}c_0(a/2)^{2d}e^{2ad}\,(1-e^{-2})^{-1}. \qedhere
\end{align*}
\end{proof}

\begin{proof}[Proof of Proposition~\ref{4.30:prop-cell-counting}, continued]
\label{4.30:prf-counting-constant}
\emph{The constant.} It remains to add the term with no interface kind met, and then to remove the
margin $\varepsilon>0$. Here $\varepsilon$ and the factor $e^{a\varepsilon}$ are those of the
earlier part of the proof, and $d^{\wedge}$ is the stratified distance of
\eqref{4.30:eq-stratified-distance-b}. Adding the term $T_0$ (the case $m=0$) to the terms
$T_m$ with $m\ge 1$, the sum over all cells $c\in\mathcal{C}$ is bounded by
\begin{equation}\label{4.30:eq-counting-constant-1}
\sum_{c\in\mathcal{C}} e^{-a\,d^{\wedge}(x,\hat{c})}
\le e^{a\varepsilon}\Bigl[T_0+\sum_{m\ge 1}T_m\Bigr].
\end{equation}
We bound the two contributions by the count of crossing histories in
\eqref{4.30:eq-crossing-histories-a} and the comparison that follows it in the proof of
Proposition~\ref{4.30:prop-cell-counting}; that proof is incomplete at the step marked $(\ast)$.
The bound on $T_0$
in \eqref{4.30:eq-crossing-histories-a}, together with $c_0(a)\le c_0(a/2)$, gives
\begin{equation*}
T_0\le e^{ad/2}c_0(a/2)^d .
\end{equation*}
For $m\ge 1$ the same counting gives $T_m\le R_m$, with $R_m$ the comparison terms defined
immediately after \eqref{4.30:eq-crossing-histories-a}. By the width condition
$(\mathrm{W}^{b})$ their sum is dominated by a geometric series
of ratio $e^{-2}$:
\begin{equation*}
\sum_{m\ge 1}T_m\le 2e^{2ad}L^{d-1}c_0(a/2)^{2d}\bigl(1-e^{-2}\bigr)^{-1}.
\end{equation*}
Inserting both bounds into \eqref{4.30:eq-counting-constant-1} yields
\begin{align*}
\sum_{c\in\mathcal{C}} e^{-a\,d^{\wedge}(x,\hat{c})}
&\le e^{a\varepsilon}\Bigl[e^{ad/2}c_0(a/2)^d
+2e^{2ad}L^{d-1}c_0(a/2)^{2d}\bigl(1-e^{-2}\bigr)^{-1}\Bigr]\\
&= e^{a\varepsilon}\,C_{\mathrm{cc}}(a,d,L).
\end{align*}
The equality is the definition of $C_{\mathrm{cc}}(a,d,L)$ in the corresponding display.
The left-hand side does not depend on $\varepsilon$, and $\varepsilon>0$ was arbitrary.
Letting $\varepsilon\downarrow 0$ therefore gives the bound of
the corresponding display.
\end{proof}

\begin{remark}
For a stratum $S$ write $\ell_S:=\ell_{i(S)}$ for its unit.

Suppose one kept the approach leg $|x-y_1|/\ell_{S_0}$ of the contour as a separate piece,
with an entry sum of its own bounded by a $d$-th power of $c_0$. The count would still be
correct. However,
the whole sum over $m$ would then carry the prefactor $e^{ad/2}c_0(a/2)^d$, and the argument
would prove only the constant
\begin{align*}
C'&:=e^{ad/2}c_0(a/2)^d\Bigl[1+2e^{2ad}L^{d-1}c_0(a/2)^{2d}\bigl(1-e^{-2}\bigr)^{-1}\Bigr]\\
&\le 4e^{5ad/2}L^{d-1}c_0(a/2)^{3d}.
\end{align*}
This constant has the same form as $C_{\mathrm{cc}}(a,d,L)$ and lies one power $c_0(a/2)^d$ above it.
The argument above merges the approach leg into the first crossing. This removes the separate
entry sum, and with it the prefactor, so that the constant $C_{\mathrm{cc}}(a,d,L)$ of
the corresponding display is obtained. On the part $m\ge 1$ there is a factor $3$ to
spare, since the comparison following \eqref{4.30:eq-crossing-histories-a} in the proof of
Proposition~\ref{4.30:prop-cell-counting}, which is incomplete at the step marked $(\ast)$,
gives $T_m\le R_m/3$ for $m\ge 1$.

\cite[p.~233]{B-CMP96} states the constant $c_1(a)=12c_0(a/2)^d$ for admissible contours.
In \cite{B-CMP96} the entropy of the first crossing is contained in the bound
\begin{equation*}
6\sum_{j'}e^{-|j-j'|}\le 12,
\end{equation*}
in which $j$ and $j'$ are indices in the notation of that paper,
because for admissible contours the first lateral sum is bounded by a $d$-th power of $c_0$
without the factor $L^{d-1}$. For the free contours considered here, the first lateral sum
carries the factor $L^{d-1}$, and the constant obtained is $C_{\mathrm{cc}}(a,d,L)$. It depends on
$a$, $d$ and $L$ only.
\end{remark}

\begin{corollary}[The bound for bonds and plaquettes attached to cells]\label{4.30:cor-attached-objects}
Let $\mathfrak{S}$ be a stratification of $T'$, let $a>0$ and $x\in T'$, and assume that for these
data the conclusion of the cell-counting bound for free contours
(Proposition~\ref{4.30:prop-cell-counting}) holds, namely
\begin{equation*}
\sum_{c\in\mathcal{C}} e^{-a\,d^{\wedge}(x,\hat{c})}\le C_{\mathrm{cc}}(a,d,L),
\end{equation*}
with $d^{\wedge}$ the stratified distance of \eqref{4.30:eq-stratified-distance-b} and
$C_{\mathrm{cc}}(a,d,L)$ the constant of the corresponding display. Consider a family of objects $o$,
each a bond or a plaquette, each attached to a cell $c(o)\in\mathcal{C}$ with centre
$\hat{c}(o)$, and each carrying a number $\tau_o$. Suppose that
\begin{enumerate}
\item[(i)] at most $2d$ objects of the family are attached to each cell $c\in\mathcal{C}$;
\item[(ii)] $\tau_o\ge d^{\wedge}(x,\hat{c}(o))-d$ for every object $o$ of the family.
\end{enumerate}
Then
\begin{equation}\label{4.30:eq-attached-objects-1}
\sum_{o} e^{-a\tau_o}\le 2d\, e^{ad}\, C_{\mathrm{cc}}(a,d,L).
\end{equation}
In particular, at $a=\tfrac14\delta$, where $\delta=\tfrac18\delta_0$ is the decay rate per block
unit of the stratified setting, for which the inequality $e^{d\delta}\le 2$ holds, one has
$e^{ad}\le 2^{1/4}$ and $e^{2ad}\le 2^{1/2}$.
\end{corollary}

\begin{proof}
Hypothesis (ii) says that the datum $\tau_o$ falls short of the stratified distance from $x$ to
the centre of the cell of $o$ by at most $d$. It is not restrictive: it holds with
$\tau_o=d^{\wedge}(x,o)$, the stratified distance from $x$ to the object $o$, in the sense fixed
with \eqref{4.30:eq-stratified-distance-b}. Indeed, by the triangle inequality for the stratified
distance, together with the half-diagonal of the cell, measured in the unit of
its stratum and counted twice, one has
\begin{equation*}
\bigl|d^{\wedge}(x,o)-d^{\wedge}(x,\hat{c}(o))\bigr|\le d ,
\end{equation*}
and in particular $d^{\wedge}(x,o)\ge d^{\wedge}(x,\hat{c}(o))-d$. In the proof of
\eqref{4.30:eq-attached-objects-1} below only hypotheses (i) and (ii) are used.

Now let $o$ be any object of the family. Since $a>0$, hypothesis (ii) gives
\begin{equation*}
e^{-a\tau_o}\le e^{-a\,(d^{\wedge}(x,\hat{c}(o))-d)}=e^{ad}\,e^{-a\,d^{\wedge}(x,\hat{c}(o))}.
\end{equation*}
All terms are non-negative, so we may sum over the family by grouping the objects according to
the cell to which they are attached. For $c\in\mathcal{C}$ let $n(c)$ be the number of objects
attached to $c$; then $n(c)\le 2d$ by hypothesis (i). Hence
\begin{align*}
\sum_{o} e^{-a\tau_o}
&\le e^{ad}\sum_{c\in\mathcal{C}} n(c)\, e^{-a\,d^{\wedge}(x,\hat{c})}
 \le 2d\, e^{ad}\sum_{c\in\mathcal{C}} e^{-a\,d^{\wedge}(x,\hat{c})}\\
&\le 2d\, e^{ad}\, C_{\mathrm{cc}}(a,d,L).
\end{align*}
The last inequality is the assumed bound. It is the conclusion of
Proposition~\ref{4.30:prop-cell-counting}, whose proof is incomplete at one step. This proves
\eqref{4.30:eq-attached-objects-1}.

Finally, let $a=\tfrac14\delta$, and recall $e^{d\delta}\le 2$. Then $ad=\tfrac14 d\delta$ and
$2ad=\tfrac12 d\delta$. Since $t\mapsto t^{1/4}$ and $t\mapsto t^{1/2}$ are increasing on
$[0,\infty)$, we obtain
\begin{equation*}
e^{ad}=\bigl(e^{d\delta}\bigr)^{1/4}\le 2^{1/4},
\qquad
e^{2ad}=\bigl(e^{d\delta}\bigr)^{1/2}\le 2^{1/2}. \qedhere
\end{equation*}
\end{proof}

\begin{proposition}[The cell-counting bound on the stratified frames of the induction]
\label{4.30:prop-frames-instance}
Let $(X',\mathfrak{A},\mathfrak{B}^{*})$ be a frame of the induction: $X'$ is its region, $\mathfrak{A}$ its
parent structure, and $\mathfrak{B}^{*}$ the structure localised to $X'$, made of the part
$\mathbb{B}_{\bar\imath}(X')$ at the top level $\bar\imath$ and of the pieces $\Gamma^{\circ}_i\cap X'$ at the
levels $i<\bar\imath$. Assume that $\mathfrak{B}^{*}$ induces a stratification of $X'$ as in hypothesis~(g)
below, whose strata are of two kinds: a \emph{parent stratum} is a stratum that contains a piece of the parent
structure, and the other strata are made only of grades of the staircase along the boundary $\partial X'$; we
call such a stratum a \emph{staircase stratum}. Let
\[
  \delta := \tfrac18\,\delta_0
\]
be the decay rate per block unit, where $\delta_0$ is the decay constant of Ba{\l}aban's estimate
\cite[(2.61)]{B-CMP96}, and put
\begin{equation}\label{4.30:eq-frames-instance-1}
  a := \tfrac14\,\delta = \tfrac{1}{32}\,\delta_0 = \alpha\,\delta_0,
  \qquad \alpha = \tfrac{1}{32}.
\end{equation}
Let $d_{*} := d^{\wedge}_{\mathfrak{B}^{*}}$ be the stratified distance of
Definition~\ref{4.30:def-stratified-distance}, taken over contours in $T' = X'$. The cells are the cells of
$\mathfrak{B}^{*}$ at every level; no level and no cell of $\mathfrak{B}^{*}$ is left out. Suppose that to each
cell of $\mathfrak{B}^{*}$ at most $2d$ bonds are attached. Put $\mathsf{w} := M/M_1$, the floor for the width,
in own units, of a crossed parent stratum, and let $\mathfrak{M} := R_1 M_1^{[4]}$ be the width parameter of the
staircase grades, $R_1$ and $M_1^{[4]}$ being constants entering the floor $(\mathrm{F}\text{-}M_1)^{b}$.
Assume the following.
\begin{itemize}
  \item[(a)] The floors $(\lambda 6)$, $(\mathrm{F}^{b})$ and $(\mathrm{F}\text{-lev})$ hold; they are used for
  the parent strata, with their constants unchanged.
  \item[(b)] The floor $(\mathrm{F}\text{-}M_1)^{b}$ holds; it is used for the staircase grades, with its
  constants unchanged, and it contains the inequality $\delta_0 R_1 M_1 > 128$.
  \item[(c)] $e^{d\delta} \le 2$.
  \item[(d)] Hypothesis $(\mathrm{S0})^{\mathfrak{B}^{*}}$, the form taken at $\partial X'$ by condition (S0)
  of~\eqref{4.30:eq-stratified-distance-a}: let $y$ be a lattice point of $\partial X'$; then the closed
  $\eta$-cubes of $X'$ containing $y$ are connected to one another through shared faces, or they all belong to
  one stratum of $\mathfrak{B}^{*}$. In other words, $X'$ has no boundary pinch between different strata.
  \item[(e)] Hypothesis $(\mathrm{S1})^{\mathfrak{B}^{*}}$: adjacent strata of $\mathfrak{B}^{*}$ differ by one
  level.
  \item[(f)] Hypothesis $(\mathrm{S2})^{\mathrm{stair}}$: each grade of the boundary staircase of $X'$ is crossed
  at a width, in its own unit, at least the width $L\mathfrak{M}$ entering the floor
  $(\mathrm{F}\text{-}M_1)^{b}$, and every staircase stratum $S$ satisfies condition (S2)
  of~\eqref{4.30:eq-stratified-distance-a} with $W_S \ge L\mathfrak{M}$.
  \item[(g)] The structure $\mathfrak{B}^{*}$ induces a stratification of $X'$ in the sense of
  Definition~\ref{4.30:def-stratified-distance}: its strata are unions of whole aligned blocks of their level,
  each stratum being of one level, and a grade of the boundary staircase and a parent piece of the same level
  that meet along a $(d-1)$-dimensional set belong to one stratum. This stratification satisfies condition (S3)
  of~\eqref{4.30:eq-stratified-distance-a}, and condition (S2) of~\eqref{4.30:eq-stratified-distance-a} with
  $W_S \ge \mathsf{w}$ for every stratum $S$ containing a piece of the parent structure.
  \item[(h)] The conclusion of Proposition~\ref{4.30:prop-cell-counting} (stated; proof incomplete at the step
  marked $(\ast)$) holds: for every stratification and every rate $a>0$ satisfying the hypotheses of that
  proposition, its cell-counting bound holds.
\end{itemize}
Then the following hold.
\begin{enumerate}
  \item The stratification of $X'$ by $\mathfrak{B}^{*}$ satisfies, at the rate $a$
  of~\eqref{4.30:eq-frames-instance-1}, the hypotheses of Proposition~\ref{4.30:prop-cell-counting}. Hence, by
  hypothesis~(h), the conclusion of that proposition holds with $d^{\wedge} = d_{*}$ and with the constant
  $C_{\mathrm{cc}}(\tfrac14\delta, d, L)$.
  \item If, moreover, $x$ is a point of $X'$ and the distance datum $\tau_o$ of each bond $o$ attached at a
  cell $c$ satisfies $\tau_o \ge d_{*}(x,\hat{c}) - d$, then the conclusion of
  Corollary~\ref{4.30:cor-attached-objects} holds at $a = \tfrac14\delta$ for the bonds attached to the cells.
  \item In the constants of those statements,
  \begin{equation}\label{4.30:eq-frames-instance-2}
    e^{ad} \le 2^{1/4}, \qquad e^{2ad} \le 2^{1/2};
  \end{equation}
  this is the evaluation that accompanies Corollary~\ref{4.30:cor-attached-objects}.
\end{enumerate}
\end{proposition}

\begin{proof}
We check the hypotheses of Proposition~\ref{4.30:prop-cell-counting} for the stratification of $T' = X'$ by
$\mathfrak{B}^{*}$ at the rate~\eqref{4.30:eq-frames-instance-1}, one at a time. We then draw the three
assertions. Proposition~\ref{4.30:prop-cell-counting} (stated; proof incomplete at the step marked $(\ast)$)
enters only through hypothesis~(h), which takes its conclusion by statement.

\emph{The data.} In a frame the bound is applied to the structure $\mathfrak{B}^{*}$ and to the distance
$d_{*} = d^{\wedge}_{\mathfrak{B}^{*}}$, with contours in $X'$. The cells entering the sum are the cells of
$\mathfrak{B}^{*}$ at every level, no level and no cell being left out, and at most $2d$ bonds are attached
to a cell. The rate is $a = \frac14\delta$ with $\delta = \frac18\delta_0$. Hence
\[
  a = \tfrac{1}{32}\,\delta_0 = \alpha\,\delta_0, \qquad \alpha = \tfrac{1}{32}.
\]
This is the value of Ba{\l}aban's parameter $\alpha$ at which his estimate \cite[(2.61)]{B-CMP96} is read.

\emph{Hypothesis} (S0). For comparison, we first consider Ba{\l}aban's nested regions $B^i(\Lambda_i)$ on the
torus, \cite[p.~231]{B-CMP96}. There every lattice point is an interior point, and (S0) is empty. For those sets
nothing beyond the setting of \cite[Lemma~2.1]{B-CMP96} is required.

For the structures $\mathfrak{B}^{*}$ localised to $X'$, the contours are taken in $X'$, so $T' = X'$ has a
boundary. At an interior lattice point of $X'$, (S0) holds automatically. At a lattice point $y$ of $\partial X'$,
(S0) asks the following: the closed $\eta$-cubes of $X'$ containing $y$ are connected to one another through
shared faces, or they all belong to one stratum of $\mathfrak{B}^{*}$. This is a property of the construction of
$X'$. It is not contained in hypothesis~(e), which speaks only of contacts that skip
levels. It is hypothesis~(d), and (S0) holds by it.

\emph{Hypotheses} (S1) \emph{and} (S3). For comparison, the strata of Ba{\l}aban's nested regions on the torus
are his regions $S_i = B^i(\Lambda_i)$ of \cite[p.~231]{B-CMP96}. Each is a union of blocks of level $i$, and the
strata are nested level by level. Ba{\l}aban states that ``The surface $\Sigma_j$ separates the sets
$B^j(\Lambda_j)$ and $B^{j-1}(\Lambda_{j-1})$'' \cite[p.~231]{B-CMP96}. The compatible partitions of
\cite{B-CMP119} have their corners on the coarser lattice. So for those regions (S1) and (S3) hold: adjacent
regions differ by one level, and their interfaces lie on the coarser lattice.

Consider next the structures $\mathfrak{B}^{*}$ localised to $X'$. For them the level-separation bound for the
frames of the induction is stated in unit-step form; applied to $\mathfrak{B}^{*}$, it reads
\begin{equation}\label{4.30:eq-frames-instance-3}
  d^{\wedge}_{\mathfrak{B}^{*}}(y,y') \;\ge\; \mathsf{w}\,\bigl(|i-i'|-1\bigr)^{+}
\end{equation}
for points $y, y'$ of cells $c, c'$ of $\mathfrak{B}^{*}$ at levels $i, i'$. The
form~\eqref{4.30:eq-frames-instance-3} presumes that every intermediate level between $i$ and $i'$ is crossed as
a stratum, and (S1) presumes the same.

Suppose a structure had an interface with a level gap $n \ge 2$. At that interface, the lateral factor in the
proof of Proposition~\ref{4.30:prop-cell-counting}, at the passage to a finer side, would be $L^{n(d-1)}$, and
$(\mathrm{W}^{b})$ would have to carry this factor. Both~\eqref{4.30:eq-frames-instance-3} and the lateral count
in the bond-sum constant of the frames are stated for $n = 1$. In that count a single factor $L^{d-1}$ is taken at
the first adjacent finer layer. The statement that adjacent strata of $\mathfrak{B}^{*}$ differ by one level is
hypothesis~(e), and (S1) holds by it. Condition (S3) holds by hypothesis~(g), and so
does the fact that $\mathfrak{B}^{*}$ induces a stratification of $X'$ in the sense of
Definition~\ref{4.30:def-stratified-distance}.

\emph{Hypothesis} (S2) \emph{for parent strata.} For parent strata the width floor is $\mathsf{w} = M/M_1$, the
floor, in own units, for a crossed stratum. It is the floor in which the thickness bound for the strata of the
parent structure, that every stratum of $\mathfrak{A}$ is at least $\mathsf{w}$ own units thick, and the
level-separation bound~\eqref{4.30:eq-frames-instance-3} are written. A stratum of $\mathfrak{B}^{*}$ containing a
parent piece may, by hypothesis~(g), also contain grades of the boundary staircase of the same level; the bound
$W_S \ge \mathsf{w}$ for every stratum $S$ containing a parent piece is part of hypothesis~(g), and (S2) holds for
parent strata by that hypothesis.

\emph{Hypothesis} (S2) \emph{for staircase strata.} The width of a grade of the boundary staircase of $X'$ enters
in the form in which it is used. Let $W$ be the width of a further crossing, as in
\cite[(2.58)]{B-CMP96}. That crossing leaves the factor $e^{-\frac12\alpha\delta_0 W}$. For a grade of the
staircase, with $W = L\mathfrak{M}$, the half of the exponent of this factor that is set against the lateral factor
$L^{d-1}$ gives
\begin{equation}\label{4.30:eq-frames-instance-4}
  e^{-\delta_0 L\mathfrak{M}/128}\,L^{d-1} \;\le\; e^{-L}\,L^{d-1}
\end{equation}
by the floor $(\mathrm{F}\text{-}M_1)^{b}$, which contains $\delta_0 R_1 M_1 > 128$. Thus a grade of the staircase
is crossed at a cost of at least its width $L\mathfrak{M}$, counted in the sites of the grade, times the rate per
site; here $\mathfrak{M} = R_1 M_1^{[4]}$. The bound $W_S \ge L\mathfrak{M}$ for every staircase stratum $S$ is
part of hypothesis~(f), and (S2) holds for staircase strata by that hypothesis.

\emph{The width condition} $(\mathrm{W}^{b})$. Beyond the factors of \cite[(2.47)--(2.59)]{B-CMP96}, the proof of
Proposition~\ref{4.30:prop-cell-counting} for free contours carries two further factors: the rounding factor
$e^{2ad}$ for each crossing, and the lateral factor $L^{d-1}$ for each passage to a finer side. All of these after
the first crossing are paid by the width of the crossed stratum, which is what the width condition
$(\mathrm{W}^{b})$ of the corresponding display expresses; it is required for every stratum $S$ at the rate
$a = \frac14\delta$.

With $W_S \ge \mathsf{w}$ for parent strata and $W_S \ge L\mathfrak{M}$ for staircase strata, as obtained above,
$(\mathrm{W}^{b})$ follows from two sets of floors:
\begin{itemize}
  \item $(\lambda 6)$ together with $(\mathrm{F}^{b})$ for the parent strata, by
  hypothesis~(a);
  \item $(\mathrm{F}\text{-}M_1)^{b}$ for the staircase grades, by hypothesis~(b) and
  \eqref{4.30:eq-frames-instance-4}.
\end{itemize}
That it follows, at $a = \frac14\delta$ and with these widths, is the content of the statement of the width
condition and the bond-sum constant on the frames, whose constant is displayed
in~\eqref{4.30:eq-bond-constant-a}. The floors enter there with their constants unchanged and in one-sided form.

\emph{Conclusion.} All hypotheses of Proposition~\ref{4.30:prop-cell-counting} now hold for the stratification of
$X'$ by $\mathfrak{B}^{*}$ at $a = \frac14\delta$: the stratification property of hypothesis~(g), (S0)--(S3) as
shown above, and $(\mathrm{W}^{b})$. By hypothesis~(h), which takes the conclusion of that proposition (stated;
proof incomplete at the step marked $(\ast)$) by statement, its cell-counting bound holds for $d^{\wedge} = d_{*}$
with the constant $C_{\mathrm{cc}}(\frac14\delta, d, L)$. This is assertion~(1).

Suppose now, as in assertion~(2), that $x$ is a point of $X'$ and that $\tau_o \ge d_{*}(x,\hat{c}) - d$ for each
bond $o$ attached at a cell $c$. To each cell of $\mathfrak{B}^{*}$ at most $2d$ bonds are attached, by
assumption. Both hypotheses of Corollary~\ref{4.30:cor-attached-objects} therefore hold with $d^{\wedge} = d_{*}$,
and the corollary, applied at $a = \frac14\delta$ together with assertion~(1), gives its bound for the bonds
attached to the cells. This is assertion~(2).

Finally, $ad = \frac14 d\delta$. Hypothesis~(c) gives $e^{d\delta} \le 2$, so
\[
  e^{ad} = \bigl(e^{d\delta}\bigr)^{1/4} \le 2^{1/4},
  \qquad
  e^{2ad} = \bigl(e^{d\delta}\bigr)^{1/2} \le 2^{1/2}.
\]
This is~\eqref{4.30:eq-frames-instance-2}, and assertion~(3) follows.
\end{proof}

\begin{lemma}[The width condition and the bond-sum constant on the frames]\label{4.30:prf-bond-constant}
Let $d\ge 2$. Let $\delta_0$ be the constant of \cite{B-CMP96}, and let $\alpha$ be the parameter of
\cite{B-CMP96}, taken here at the value $\alpha=1/32$. Let $\delta:=\tfrac18\delta_0$
be the decay rate per block unit of the frames, and put
\begin{equation*}
a:=\frac{\delta_0}{32}=\frac{\delta}{4},\qquad\text{so that}\qquad \frac a2=\frac{\delta_0}{64},\qquad
\frac a4=\frac{\delta_0}{128}=\frac{\delta}{16}.
\end{equation*}
Recall that $c_0(b)=\sum_{z\in\mathbb{Z}}e^{-b|z|}=\coth(b/2)$ (Lemma~\ref{4.30:lem-coth-comparison}).
Let $\mathfrak{B}^{*}$ be the stratification of the frame's region $X'$, with stratified distance
$d_{*}:=d^{\wedge}_{\mathfrak{B}^{*}}$ (contours in $X'$) and widths $W_S$. Let $w:=M/M_1$ be the width
parameter of the parent strata, let $\varrho_0:=\lceil 8\delta^{-1}\ln L\rceil$, and let
$\mathfrak{M}$ be the width parameter of the strata of staircase grade, fixed with the frames. Assume:
\begin{enumerate}
\item[] the inequality
\begin{equation}\label{4.30:eq-bond-constant-1}
e^{d\delta}\le 2;
\end{equation}
\item[$(\lambda 6)$] $\dfrac{\delta_0}{128}\,w>2d\ln c_0(\delta_0/64)+1$;
\item[$(\mathrm{F}^{b})$] $w\ge 3L(\varrho_0+2L)$;
\item[$(\mathrm{F}\text{-}M_1)^{b}$] $\dfrac{\delta_0\mathfrak{M}}{128}>2d\ln c_0(\delta_0/64)+1$;
\item[] the width dichotomy: every stratum $S$ of $\mathfrak{B}^{*}$ is either a parent stratum, with
$W_S\ge w$, or a stratum of staircase grade, with $W_S\ge L\mathfrak{M}$.
\end{enumerate}
Then the following hold.
\begin{enumerate}
\item[(a)] The width condition $(\mathrm{W}^{b})$ of the corresponding display holds at $a=\delta_0/32$
for every stratum $S$ of $\mathfrak{B}^{*}$.
\item[(b)] The constant $C_{\mathrm{cc}}$ of the corresponding display satisfies
\begin{equation}\label{4.30:eq-bond-constant-2}
C_{\mathrm{cc}}(\delta_0/32,d,L)\le 4\cdot 2^{1/2}L^{d-1}c_0(\delta_0/64)^{2d}.
\end{equation}
If, moreover, $\mathfrak{B}^{*}$ satisfies conditions (S0)--(S3) of
Proposition~\ref{4.30:prop-cell-counting}, and bonds $o$ are attached to the cells of $\mathfrak{B}^{*}$,
at most $2d$ to each cell, with distance data $\tau_o$ relative to a point $x$ of $X'$ satisfying the
condition on the distance data in Corollary~\ref{4.30:cor-attached-objects} at $a=\delta_0/32$, then
\begin{equation}\label{4.30:eq-bond-constant-a}
\sum_{o}e^{-a\tau_o}\le c_1^{b'}:=2d\cdot 2^{3/4}\cdot 4L^{d-1}c_0(\delta_0/64)^{2d}.
\end{equation}
\item[(c)] For a value $\alpha'$ of the parameter of \cite{B-CMP96} let
$c_1(\alpha'):=12c_0(\tfrac12\alpha'\delta_0)^d$ be the constant of \cite{B-CMP96}, and put
\begin{equation*}
c_1^{b}:=2d\cdot 2L^{d-1}c_1(1/32)=2d\cdot 2L^{d-1}\cdot 12c_0(\delta_0/64)^d .
\end{equation*}
Then
\begin{equation}\label{4.30:eq-bond-constant-3}
c_1^{b'}=\frac{4\cdot 2^{3/4}}{24}\,c_0(\delta_0/64)^d\,c_1^{b}.
\end{equation}
\end{enumerate}
\end{lemma}

This lemma is used in the proof of Proposition~\ref{4.30:prop-frames-instance}.

\begin{proof}
\emph{Step 1: reduction of $(\mathrm{W}^{b})$ to two quarters.} For a stratum $S$ consider the two
inequalities
\begin{align}
c_0(a/2)^{2d}\cdot e\cdot e^{-(a/4)W_S}&\le 1, \label{4.30:eq-bond-constant-4}\\
2L^{d-1}e^{2ad}\cdot e\cdot e^{-(a/4)W_S}&\le 1. \label{4.30:eq-bond-constant-5}
\end{align}
Their product is
\begin{equation*}
2L^{d-1}e^{2ad}c_0(a/2)^{2d}\,e^{2}\,e^{-(a/2)W_S}\le 1,
\end{equation*}
which is $(\mathrm{W}^{b})$ of the corresponding display, the condition of the cell-counting bound for
free contours (Proposition~\ref{4.30:prop-cell-counting}, whose proof is incomplete at one step). Only this
direction is used: the two quarters imply $(\mathrm{W}^{b})$. At $a=\delta_0/32$ we have $a/2=\delta_0/64$
and $a/4=\delta_0/128=\delta/16$, so \eqref{4.30:eq-bond-constant-4} and \eqref{4.30:eq-bond-constant-5}
read
\begin{equation}\label{4.30:eq-bond-constant-6}
\frac{\delta_0}{128}W_S\ge 2d\ln c_0(\delta_0/64)+1,\qquad
\frac{\delta}{16}W_S\ge (d-1)\ln L+1+\ln 2+2ad.
\end{equation}
Both right-hand sides are independent of $S$ and both left-hand sides increase with $W_S$, so by
the width dichotomy it suffices to prove \eqref{4.30:eq-bond-constant-6} with $W_S=w$ for parent
strata and with $W_S=L\mathfrak{M}$ for strata of staircase grade.

Two consequences of \eqref{4.30:eq-bond-constant-1} are used repeatedly. Since $a=\delta/4$,
\begin{equation}\label{4.30:eq-bond-constant-7}
2ad=\tfrac12 d\delta\le\tfrac12\ln 2,\qquad e^{ad}\le 2^{1/4},\qquad e^{2ad}\le 2^{1/2}.
\end{equation}
Moreover, $c_0(\delta_0/64)=\coth(\delta_0/128)$ by Lemma~\ref{4.30:lem-coth-comparison}, and
$\coth u>1/u$ for $u>0$ because $\tanh u<u$; since $\delta_0=8\delta\le 8\ln 2/d$ by
\eqref{4.30:eq-bond-constant-1},
\begin{equation}\label{4.30:eq-bond-constant-8}
c_0(\delta_0/64)>\frac{128}{\delta_0}\ge\frac{16d}{\ln 2}.
\end{equation}
In particular, as $d\ge 2$, $c_0(\delta_0/64)>32/\ln 2>46$.

\emph{Step 2: parent strata, first quarter.} With $W_S=w$, the first inequality of
\eqref{4.30:eq-bond-constant-6} is $(\delta_0/128)w\ge 2d\ln c_0(\delta_0/64)+1$, which follows from
$(\lambda 6)$. Hypothesis $(\lambda 6)$ is the inequality \cite[(2.59), p.~233]{B-CMP96}, namely
$\tfrac14\alpha\delta_0RM>2d\log c_0(\tfrac12\alpha)+1$, taken at $\alpha=1/32$ with $RM$ replaced by $w$;
there $c_0(\tfrac12\alpha)$ denotes $c_0$ at the rate $\tfrac12\alpha\delta_0=\delta_0/64$, and
$\tfrac14\alpha\delta_0 w=\delta_0 w/128$.

\emph{Step 3: parent strata, second quarter.} We show
\begin{equation}\label{4.30:eq-bond-constant-9}
\frac{\delta}{16}w\ge(d-1)\ln L+1+\ln 2+2ad .
\end{equation}
Two lower bounds for $\delta w/16$ are available. First, by $(\mathrm{F}^{b})$, $w\ge 3L\varrho_0$, and
$\delta\varrho_0\ge 8\ln L$ by the definition of $\varrho_0$, hence
\begin{equation*}
\frac{\delta}{16}w\ge\frac{3L}{16}\,\delta\varrho_0\ge\frac{3L}{16}\cdot 8\ln L=\frac{3L}{2}\ln L .
\end{equation*}
Second, since $\delta/16=\delta_0/128$, hypothesis $(\lambda 6)$ gives $\delta w/16>2d\ln c_0(\delta_0/64)+1$.
Therefore
\begin{equation*}
\frac{\delta}{16}w\ge\max\Bigl\{\frac{3L}{2}\ln L,\ 2d\ln c_0(\delta_0/64)+1\Bigr\}.
\end{equation*}
By \eqref{4.30:eq-bond-constant-7} the right-hand side of \eqref{4.30:eq-bond-constant-9} is at most
$(d-1)\ln L+1+\tfrac32\ln 2$. Since $L$ is odd and $L>1$, we have $L\ge 3$. We distinguish two cases.

\emph{Case $L\ge d$.} Then $\tfrac{3L}{2}-d+1\ge\tfrac L2+1\ge\tfrac52$ and $\ln L\ge\ln 3$, so
\begin{equation*}
\frac{3L}{2}\ln L-(d-1)\ln L=\Bigl(\frac{3L}{2}-d+1\Bigr)\ln L\ge\frac52\ln 3 .
\end{equation*}
As $\ln 3>1.09$ and $\ln 2<0.7$,
\begin{equation*}
\frac52\ln 3>2.7>2.05>1+\frac32\ln 2\ge 1+\ln 2+2ad,
\end{equation*}
the last step by \eqref{4.30:eq-bond-constant-7}. Hence the first lower bound yields
\eqref{4.30:eq-bond-constant-9}.

\emph{Case $L<d$.} Here $(d-1)\ln L<(d-1)\ln d$. By \eqref{4.30:eq-bond-constant-8},
$\ln c_0(\delta_0/64)>\ln d+\ln(16/\ln 2)$, and $16/\ln 2>22>e^3$, so
$\ln c_0(\delta_0/64)>\ln d+3$. Consequently
\begin{equation*}
2d\ln c_0(\delta_0/64)+1>2d\ln d+6d+1>(d-1)\ln d+1+\frac32\ln 2,
\end{equation*}
because $2d\ln d\ge(d-1)\ln d$ and $6d>\tfrac32\ln 2$. Hence the second lower bound yields
\eqref{4.30:eq-bond-constant-9}.

In either case \eqref{4.30:eq-bond-constant-9} is obtained as a maximum of the two lower bounds, and no
condition beyond the hypotheses is used. The first lower bound alone does not suffice in the second case: at
$L=3$, $d=4$ one has $\tfrac92\ln 3<3\ln 3+1+\ln 2$, since
\begin{equation*}
\frac92\ln 3-3\ln 3=\frac32\ln 3<1.65<1.69<1+\ln 2 ;
\end{equation*}
there the second lower bound, which rests on $(\lambda 6)$ and \eqref{4.30:eq-bond-constant-1}, is the one
used.

\emph{Step 4: strata of staircase grade.} Let $W_S=L\mathfrak{M}$. Both quarters are obtained from
$(\mathrm{F}\text{-}M_1)^{b}$ as stated. For the first inequality of \eqref{4.30:eq-bond-constant-6}: since
$L\ge 1$,
\begin{equation*}
\frac{\delta_0}{128}L\mathfrak{M}\ge\frac{\delta_0\mathfrak{M}}{128}>2d\ln c_0(\delta_0/64)+1 .
\end{equation*}
For the second quarter \eqref{4.30:eq-bond-constant-5}, exponentiating $(\mathrm{F}\text{-}M_1)^{b}$ gives
$e^{-\delta_0\mathfrak{M}/128}<\bigl(c_0(\delta_0/64)^{2d}e\bigr)^{-1}$, hence
\begin{equation*}
e^{-\delta_0L\mathfrak{M}/128}=\bigl(e^{-\delta_0\mathfrak{M}/128}\bigr)^{L}
<\bigl(c_0(\delta_0/64)^{2d}e\bigr)^{-L}.
\end{equation*}
Together with $e^{2ad}\le 2^{1/2}$ from \eqref{4.30:eq-bond-constant-7}, the left-hand side of
\eqref{4.30:eq-bond-constant-5} is less than $2\sqrt2\,e\,L^{d-1}\bigl(c_0(\delta_0/64)^{2d}e\bigr)^{-L}$. It
therefore suffices that
\begin{equation}\label{4.30:eq-bond-constant-10}
2\sqrt2\,e\,L^{d-1}\le\bigl(c_0(\delta_0/64)^{2d}e\bigr)^{L}.
\end{equation}
Taking logarithms and writing $c_0=c_0(\delta_0/64)$, \eqref{4.30:eq-bond-constant-10} is $g(L,d)\ge 0$, where
\begin{equation*}
g(L,d):=L\bigl(2d\ln c_0+1\bigr)-(d-1)\ln L-1-\frac32\ln 2 .
\end{equation*}
We show that $g(L,d)>0$ for all $L\ge 3$ and $d\ge 2$ as soon as $c_0\ge 1.7$; this holds amply, since
$c_0>46$ by Step 1. Let $c_0\ge 1.7$, so $\ln c_0>0.53$. Regard $g$ as a function of real $L\ge 3$ and real
$d\ge 2$. Then
\begin{align*}
\partial_L g&=2d\ln c_0+1-\frac{d-1}{L}>1.06\,d+1-\frac{d-1}{3}>0,\\
\partial_d g&=2L\ln c_0-\ln L>1.06\,L-\ln L>0,
\end{align*}
so $g$ is increasing in $L$ and in $d$. It is also increasing in $c_0$, the coefficient $2dL$ of $\ln c_0$
being positive. Hence, using $\ln 3<1.1$ and $\tfrac32\ln 2<1.05$,
\begin{equation*}
g(L,d)\ge g(3,2)\big|_{c_0=1.7}>3(4\cdot 0.53+1)-1.1-1-1.05>6>0 .
\end{equation*}
This proves \eqref{4.30:eq-bond-constant-10}, hence \eqref{4.30:eq-bond-constant-5} for strata of staircase
grade. No lower bound on $L$ other than $L\ge 3$ enters.

By Steps 1--4, \eqref{4.30:eq-bond-constant-4} and \eqref{4.30:eq-bond-constant-5} hold for every stratum of
$\mathfrak{B}^{*}$, and Step 1 gives (a).

\emph{Step 5: the constant.} By the defining display of the constant of the
cell-counting bound for free contours (Proposition~\ref{4.30:prop-cell-counting}, whose proof is incomplete at
one step),
\begin{equation*}
C_{\mathrm{cc}}(a,d,L)\le 4e^{2ad}L^{d-1}c_0(a/2)^{2d}.
\end{equation*}
At $a=\delta_0/32$, so that $a/2=\delta_0/64$, and with $e^{2ad}\le 2^{1/2}$ from
\eqref{4.30:eq-bond-constant-7}, this is \eqref{4.30:eq-bond-constant-2}. Under the further hypotheses stated
in (b), $\mathfrak{B}^{*}$ satisfies (S0)--(S3), and by (a) it satisfies $(\mathrm{W}^{b})$ at
$a=\delta_0/32$; so the cell-counting bound for free contours
(Proposition~\ref{4.30:prop-cell-counting}, whose proof is incomplete at one step), and hence
Corollary~\ref{4.30:cor-attached-objects}, applies to $\mathfrak{B}^{*}$ at $a=\delta_0/32$.
Corollary~\ref{4.30:cor-attached-objects} applied with the bonds as attached objects gives
$\sum_o e^{-a\tau_o}\le 2d\,e^{ad}C_{\mathrm{cc}}(a,d,L)$. With $e^{ad}\le 2^{1/4}$ from
\eqref{4.30:eq-bond-constant-7} and \eqref{4.30:eq-bond-constant-2},
\begin{equation*}
\sum_o e^{-a\tau_o}\le 2d\cdot 2^{1/4}\cdot 4\cdot 2^{1/2}L^{d-1}c_0(\delta_0/64)^{2d}
=2d\cdot 2^{3/4}\cdot 4L^{d-1}c_0(\delta_0/64)^{2d},
\end{equation*}
which is \eqref{4.30:eq-bond-constant-a}. This proves (b).

\emph{Step 6: comparison.} Dividing,
\begin{equation*}
\frac{c_1^{b'}}{c_1^{b}}=\frac{2d\cdot 2^{3/4}\cdot 4L^{d-1}c_0(\delta_0/64)^{2d}}
{2d\cdot 2L^{d-1}\cdot 12c_0(\delta_0/64)^{d}}
=\frac{4\cdot 2^{3/4}}{24}\,c_0(\delta_0/64)^{d},
\end{equation*}
which is \eqref{4.30:eq-bond-constant-3}. Both constants depend on $(d,L,\delta_0)$ only. The additional factor
$c_0(\delta_0/64)^{d}$ in $c_1^{b'}$ is the entropy of the first crossing in the proof of
Proposition~\ref{4.30:prop-cell-counting}. This proves (c).
\end{proof}

\begin{definition}[The complexified group, operator norms and the gauge action on pairs]
\label{4.31:not-complexified-group}
Throughout, $d=4$ and $G=SU(N)$ with $N\ge 2$. The following conventions are in force.
\begin{enumerate}
\item \emph{Algebras and groups.} The Lie algebra $\mathfrak{g}$ is the space of traceless Hermitian
$N\times N$ matrices, so that a bond variable is written $U=e^{iA}$ with $A\in\mathfrak{g}$, as in
Ba{\l}aban's papers. Further, $\mathfrak{g}^{\mathbb{C}}=\mathfrak{sl}(N,\mathbb{C})$ and
$G^{\mathbb{C}}=SL(N,\mathbb{C})$.
\item \emph{Norms.} The symbol $|\cdot|$ denotes the operator norm of an $N\times N$ matrix; for a
configuration it denotes the supremum over bonds of the operator norms of its values. The norm
written $|\cdot|$ in \cite{B-CMP99b} is the Hilbert--Schmidt norm
$|Y|_{\mathrm{HS}}=\bigl(\sum_{a,b}|Y_{ab}|^2\bigr)^{1/2}$, and
\begin{equation}\label{4.31:eq-complexified-group-1}
|Y|\le |Y|_{\mathrm{HS}}\le \sqrt{N}\,|Y| .
\end{equation}
The factor $\sqrt{N}$ in \eqref{4.31:eq-complexified-group-1} is absorbed into the thresholds on
the free radii and enters nowhere else.
\item \emph{Trace, real and imaginary parts, projection.} The trace is normalised,
$\operatorname{tr}\mathbf{1}=1$. For an invertible complex $N\times N$ matrix $U$, in particular
for $U\in G^{\mathbb{C}}$, we put, as in \cite{B-CMP99b},
\begin{equation*}
\operatorname{Re}U=\tfrac12\bigl(U+U^{-1}\bigr),\qquad
\operatorname{Im}U=\tfrac{1}{2i}\bigl(U-U^{-1}\bigr).
\end{equation*}
The map $\pi$ is the traceless projection $\pi Y=Y-(\operatorname{tr}Y)\mathbf{1}$ on $N\times N$
matrices.
\item \emph{Gauge action.} A gauge transformation $g$ is a $G$- or $G^{\mathbb{C}}$-valued function
on the sites of the lattice. It acts on bond variables, and on pairs $(U,J)$ with $J$ a function on
bonds with values in $\mathfrak{g}^{\mathbb{C}}$ (a current), by
\begin{align*}
U^{g}(\langle x,y\rangle)&=g(x)\,U(\langle x,y\rangle)\,g(y)^{-1},\qquad
(U,J)^{g}=\bigl(U^{g},\operatorname{Ad}(g)J\bigr),\\
\bigl(\operatorname{Ad}(g)J\bigr)(b)&=g(b_-)\,J(b)\,g(b_-)^{-1},
\end{align*}
where $b_-$ is the initial point of the bond $b$. We write $(\partial U)(p)=U(\partial p)$ for the
plaquette variable of Definition~\ref{1.1:def-wilson-action}, and $x_p$ for the corner of $p$ at
which that product starts and ends.
\item \emph{Fixed parameters.} The blocking factor $L$ is odd with $L\ge L_0$; the tuple
$\mathfrak{p}$ of auxiliary parameters, $\gamma=\gamma_H$ and $m\ge m_{\mathbb{T}}$ are fixed as in
the principal theorem of this section. We write $x_k=g_k^{-2}$ and, following
\cite[(2.4), (2.28)]{B-CMP119},
\begin{equation*}
\varepsilon(x)=x^{-1/2}A_0(\log x)^{p_0},\qquad
\alpha_i(x)=x^{-1/2}C_i(\log x)^{q_i},\quad q_i\ge 0,
\end{equation*}
with $\varepsilon_k=\varepsilon(x_k)$ and $\alpha_{i,k}=\alpha_i(x_k)$.
\item \emph{Smallness ceiling.} We assume $\gamma^2\le 1/e$. This is a smallness condition on the
coupling, imposed last among the choices, within the range of couplings of \cite{B-CMP119}.
\item \emph{Constants.} The constant $B_3$ is Ba{\l}aban's constant so denoted, and $B$ is the
parameter of the space $\mathcal{H}(B;U_0)$ of complex background fields used below; both enter
only upper bounds. Since constants in upper bounds may be enlarged, we take
$B_3\ge 1$ and $B\ge 1$. We put
\begin{equation}\label{4.31:eq-complexified-group-2}
A_{*}=\sup_{x\ge\gamma^{-2}}\frac{\alpha_0(x)}{2B_3}.
\end{equation}
\end{enumerate}
The assertions contained in these conventions are the following.
\begin{enumerate}
\item[(a)] The inequalities \eqref{4.31:eq-complexified-group-1} hold.
\item[(b)] $\pi$ is a projection onto traceless matrices and satisfies
$\pi(hYh^{-1})=h(\pi Y)h^{-1}$ for every $h\in G^{\mathbb{C}}$.
\item[(c)] For every gauge transformation $g$, whether or not it is constant on blocks, and every
plaquette $p$,
\begin{equation}\label{4.31:eq-complexified-group-3}
\partial(U^{g})(p)=g(x_p)\,(\partial U)(p)\,g(x_p)^{-1},
\end{equation}
and hence
\begin{equation}\label{4.31:eq-complexified-group-4}
\bigl|\partial(U^{g})(p)-\mathbf{1}\bigr|
\le |g(x_p)|\,|g(x_p)^{-1}|\,\bigl|(\partial U)(p)-\mathbf{1}\bigr| .
\end{equation}
\item[(d)] For $x\ge\gamma^{-2}$ one has $\log x\ge 1$ and $\alpha_i(x)\ge C_i x^{-1/2}$.
\item[(e)] $A_{*}<\infty$.
\end{enumerate}
\end{definition}

\begin{proof}[Proof of the assertions]
(a) Let $\lambda_1\ge\dots\ge\lambda_N\ge 0$ be the eigenvalues of the positive semidefinite matrix
$Y^{*}Y$. Then $|Y|^2=\lambda_1$ and $|Y|_{\mathrm{HS}}^2=\sum_{a,b}|Y_{ab}|^2=\lambda_1+\dots+\lambda_N$.
Since every $\lambda_j$ is non-negative and at most $\lambda_1$, this sum lies between $\lambda_1$
and $N\lambda_1$. Taking square roots gives \eqref{4.31:eq-complexified-group-1}.

(b) Since $\operatorname{tr}\mathbf{1}=1$, we have
$\operatorname{tr}\pi Y=\operatorname{tr}Y-\operatorname{tr}Y=0$. Hence $\pi$ maps into traceless
matrices, and $\pi(\pi Y)=\pi Y$. For $h\in G^{\mathbb{C}}$ we have
$\operatorname{tr}(hYh^{-1})=\operatorname{tr}Y$. Therefore
$\pi(hYh^{-1})=hYh^{-1}-(\operatorname{tr}Y)\mathbf{1}=h(\pi Y)h^{-1}$.

(c) Let $x_p=y_0,y_1,y_2,y_3,y_4=x_p$ be the consecutive corners of $\partial p$, and let $U_j$
denote the variable of the step from $y_{j-1}$ to $y_j$. This is $U(\langle y_{j-1},y_j\rangle)$,
or the inverse of the variable of the oppositely oriented bond. In either case the transformed
variable is $g(y_{j-1})U_jg(y_j)^{-1}$, because
\begin{equation*}
\bigl(g(y_j)U(\langle y_j,y_{j-1}\rangle)g(y_{j-1})^{-1}\bigr)^{-1}
=g(y_{j-1})U(\langle y_j,y_{j-1}\rangle)^{-1}g(y_j)^{-1}.
\end{equation*}
In the product of the four transformed variables the interior factors $g(y_j)^{-1}g(y_j)$,
$j=1,2,3$, cancel. Since $y_4=y_0=x_p$, this proves \eqref{4.31:eq-complexified-group-3}. The
argument uses only the definition of $U^{g}$ on bonds, so it holds for every gauge transformation
on the sites of the lattice. By \eqref{4.31:eq-complexified-group-3},
\begin{equation*}
\partial(U^{g})(p)-\mathbf{1}
=g(x_p)\bigl((\partial U)(p)-\mathbf{1}\bigr)g(x_p)^{-1}.
\end{equation*}
Submultiplicativity of the operator norm then gives \eqref{4.31:eq-complexified-group-4}.

(d) Let $x\ge\gamma^{-2}$. Since $\gamma^{2}\le 1/e$ we have $\gamma^{-2}\ge e$, hence $x\ge e$ and
$\log x\ge 1$. Since $q_i\ge 0$, it follows that $(\log x)^{q_i}\ge 1$, and so
$\alpha_i(x)\ge C_ix^{-1/2}$.

(e) On $[\gamma^{-2},\infty)$ the function
$\alpha_0(x)=C_0x^{-1/2}(\log x)^{q_0}$ is continuous. It tends to $0$ as $x\to\infty$, because
$(\log x)^{q_0}x^{-1/2}\to 0$. A continuous function on a closed half-line with limit $0$ at infinity
is bounded. Since $B_3\ge 1$, the supremum in \eqref{4.31:eq-complexified-group-2} is finite.
\end{proof}

\begin{definition}[Cutoffs, the depth-$s$ torus, plaquette windows and complex tubes]
\label{4.31:def-windows-tubes}
We use the complexified group $G^{\mathbb{C}}=SL(N,\mathbb{C})$, the real form $\mathfrak{g}$ of
traceless Hermitian $N\times N$ matrices, the operator norm $|\cdot|$ (for fields, the supremum
over bonds) and the gauge action of Definition~\ref{4.31:not-complexified-group}.

\emph{Cutoffs.} Let $\mathcal{K}\subset\mathbb{N}$ be an infinite set. For every $K\in\mathcal{K}$
a bare coupling $g_0=g_0(K)$ is given whose running couplings $g_0,\dots,g_K$
(Definition~\ref{1.2:def-couplings}) lie in $(0,\gamma]$, and which is chosen so that, with
$x_K=g_K^{-2}$ and for some numbers $x_-,x_+$ independent of $K$,
\begin{equation}\label{4.31:eq-windows-tubes-a}
0<x_-\le x_+<\infty,\qquad x_K\in[x_-,x_+]\quad\text{for every }K\in\mathcal{K}.
\end{equation}
That such a choice is possible is the following statement on the running couplings, used here by
statement: for every $K$ there is a bare coupling $g_0$ for which $|x_K-g^{-2}|\le b_+$, with a
number $b_+$ independent of $K$.
Condition \eqref{4.31:eq-windows-tubes-a} is the form in which the constants of the statements
below depend on the fixed reference coupling $g$; every bound below in which $g^{-2}$ enters holds
with $g^{-2}$ replaced by $x_K$. If the bare coupling is tuned so that $x_K=g^{-2}$ exactly, then
one may take $x_-=x_+=g^{-2}$. Condition \eqref{4.31:eq-windows-tubes-a} neither imposes nor
narrows a window on any coupling, and no lower bound on any coupling is used.

\emph{The depth-$s$ torus.} Let $0\le s\le K$ and put $k:=K-s$. The torus of
\cite[(0.1)]{B-CMP109} determines a sequence of tori: the $k$-th is obtained by multiplying the
lattice spacing by $L^k$. At level $k$ it is
\begin{equation}\label{4.31:eq-windows-tubes-1}
T^{(k)}=\Bigl\{x\in L^{-s}\mathbb{Z}^4+\tfrac12 L^{-s}(1,1,1,1):\ |x_\mu|<L^m,\ \mu=1,\dots,4\Bigr\}.
\end{equation}
This is a subset of the torus $\mathbb{T}_m$ of Definition~\ref{1.1:def-lattice} and does not
depend on $K$. Rescaled to unit spacing it is
denoted $\mathbb{T}_s$, and it has $(2L^{m+s})^4$ sites. The fine lattice $T_\eta$ is the bare
lattice; in the units of $\mathbb{T}_s$ its spacing is $\eta=L^{-k}$. Put
\begin{equation*}
\mathcal{G}_s:=SU(N)^{\operatorname{bonds}(\mathbb{T}_s)},\qquad
\mathcal{G}_s^{\mathbb{C}}:=(G^{\mathbb{C}})^{\operatorname{bonds}(\mathbb{T}_s)},
\end{equation*}
and let $dV$ be the Haar probability measure on $\mathcal{G}_s$.

The effective density at level $K-s$ of the theory with cutoff $K$ is
$\rho_{K,s}:=\rho^{(K)}_{K-s}=(\mathbf{R}\mathbf{T})^{K-s}\rho_0$, with $\rho_0$ the bare density
of Definition~\ref{1.5:def-densities}. It is always read as the
measure $\rho_{K,s}\,dV$. Put $Z_K:=\int\rho_0\,dV$, the integral being over the bond
configurations of the bare lattice (depth $s=K$), and $\nu^{B}_{K,s}:=\rho_{K,s}\,dV/Z_K$.

\emph{Scales.} For $0\le t\le k-1$ let $j:=k-t$. The $j$-lattice $T^{(j)}_1$ has spacing $L^{-t}$
in the units of $\mathbb{T}_s$, so $|T^{(j)}_1|=L^{4t}|\mathbb{T}_s|$. The fine spacing in
scale-$j$ units is $\xi:=L^{-j}=L^{t}\eta$. Ba{\l}aban's partition $\pi_j$ of $T^{(j)}_1$
consists of cubes of side $M$ in scale-$j$ units, that is, of side $ML^{-t}$ in the units of
$\mathbb{T}_s$. Hence $\#\pi_j=M^{-4}L^{4t}|\mathbb{T}_s|$.

\emph{Windows and tubes.} For $\varpi>0$ and $r>0$ the \emph{plaquette window} is
\begin{equation}\label{4.31:eq-windows-tubes-2}
\mathcal{V}_s(\varpi):=\bigl\{V\in\mathcal{G}_s:\ |V(\partial p)-\mathbf{1}|<\varpi
\text{ for all plaquettes }p\bigr\},
\end{equation}
and the \emph{complex tube} over it is
\begin{equation}\label{4.31:eq-windows-tubes-3}
\mathfrak{D}_s(\varpi,r):=\bigl\{e^{-Y}V_0:\ V_0\in\mathcal{V}_s(\varpi),\
Y\in\mathfrak{g}^{\operatorname{bonds}(\mathbb{T}_s)},\ |Y|<r\bigr\}.
\end{equation}
Here $(e^{-Y}V_0)(b)=e^{-Y(b)}V_0(b)$ for every bond $b$.

The tube has three properties:
\begin{itemize}
\item[(a)] $\mathfrak{D}_s(\varpi,r)$ is open in $\mathcal{G}_s^{\mathbb{C}}$;
\item[(b)] its real slice $\mathfrak{D}_s(\varpi,r)\cap\mathcal{G}_s$ equals
$\mathcal{V}_s(\varpi)$, which contains $\mathbf{1}$;
\item[(c)] $\mathfrak{D}_s(\varpi,r)$ is invariant under $SU(N)$-valued gauge transformations.
\end{itemize}
\end{definition}

\begin{proof}
Define
\begin{equation*}
\Phi:\mathcal{G}_s\times\mathfrak{g}^{\operatorname{bonds}(\mathbb{T}_s)}\to
\mathcal{G}_s^{\mathbb{C}},\qquad
\Phi(V,Y)(b)=e^{-Y(b)}V(b).
\end{equation*}
The map acts bond by bond, so it suffices to consider one bond. By the polar decomposition,
every $A\in SL(N,\mathbb{C})$ can be written uniquely as $A=PW$ with $P$ positive definite
Hermitian and $W$ unitary. Since $\det P>0$ and $|\det W|=1$, the relation
$\det P\,\det W=1$ forces $\det P=\det W=1$. Moreover $P=e^{-Y}$ for a unique traceless
Hermitian $Y$. The factors $P$ and $W$ depend smoothly on $A$, and $(W,Y)\mapsto e^{-Y}W$ is
smooth. Hence the single-bond map is a diffeomorphism $SU(N)\times\mathfrak{g}\to SL(N,\mathbb{C})$.
Taking the product over the finitely many bonds, $\Phi$ is a diffeomorphism.

(a) For each plaquette $p$, the variable $V(\partial p)$ is a product of four bond variables or
their inverses. Hence $V\mapsto|V(\partial p)-\mathbf{1}|$ is continuous on $\mathcal{G}_s$. There
are finitely many plaquettes and the inequalities in \eqref{4.31:eq-windows-tubes-2} are strict,
so $\mathcal{V}_s(\varpi)$ is open in $\mathcal{G}_s$. The set $\{Y:|Y|<r\}$ is open in
$\mathfrak{g}^{\operatorname{bonds}(\mathbb{T}_s)}$, because $|\cdot|$ is a supremum over finitely
many bonds of a norm. By \eqref{4.31:eq-windows-tubes-3},
\begin{equation*}
\mathfrak{D}_s(\varpi,r)=\Phi\bigl(\mathcal{V}_s(\varpi)\times\{|Y|<r\}\bigr).
\end{equation*}
The right-hand side is the image of an open set under a homeomorphism, hence open.

(b) Let $\mathbf{V}=\Phi(V_0,Y)\in\mathfrak{D}_s(\varpi,r)$ and suppose $\mathbf{V}\in\mathcal{G}_s$.
Then also $\mathbf{V}=\Phi(\mathbf{V},0)$. Injectivity of $\Phi$ gives $Y=0$, hence
$\mathbf{V}=V_0\in\mathcal{V}_s(\varpi)$.
Conversely, for $V_0\in\mathcal{V}_s(\varpi)$ we have $V_0=\Phi(V_0,0)$ with $|0|<r$, so
$V_0\in\mathfrak{D}_s(\varpi,r)$. Therefore
$\mathfrak{D}_s(\varpi,r)\cap\mathcal{G}_s=\mathcal{V}_s(\varpi)$. Finally
$\mathbf{1}(\partial p)=\mathbf{1}$ for every $p$, and $0<\varpi$, so
$\mathbf{1}\in\mathcal{V}_s(\varpi)$.

(c) Let $v$ be an $SU(N)$-valued function on the sites. Take $b=\langle x,y\rangle$ and
$\mathbf{V}=e^{-Y}V_0$ in the tube. By the gauge action of
Definition~\ref{4.31:not-complexified-group},
\begin{equation*}
\mathbf{V}^{v}(b)=v(x)e^{-Y(b)}V_0(b)v(y)^{-1}=e^{-Y'(b)}\,V_0^{v}(b),
\qquad Y'(b):=v(x)Y(b)v(x)^{-1}.
\end{equation*}
Each $Y'(b)$ is traceless Hermitian. Since $v(x)$ is unitary, $\operatorname{Ad}(v(x))$ is an
isometry for the operator norm, so $|Y'|=|Y|<r$.

The plaquette variables of $V_0^{v}$ are conjugates of those of $V_0$:
\begin{equation*}
V_0^{v}(\partial p)=v(x_p)V_0(\partial p)v(x_p)^{-1},
\end{equation*}
where $x_p$ is the base point of $p$. Hence
\begin{equation*}
|V_0^{v}(\partial p)-\mathbf{1}|=|v(x_p)(V_0(\partial p)-\mathbf{1})v(x_p)^{-1}|
=|V_0(\partial p)-\mathbf{1}|<\varpi,
\end{equation*}
so $V_0^{v}\in\mathcal{V}_s(\varpi)$ and $\mathbf{V}^{v}\in\mathfrak{D}_s(\varpi,r)$. The same
argument applies to $v^{-1}$, so the gauge transformation by $v$ maps $\mathfrak{D}_s(\varpi,r)$
onto itself.
\end{proof}

\begin{definition}[The exponent of the never-large term at depth $s$]
\label{4.31:def-never-large-exponent}
Let the following be as in Definition~\ref{4.31:def-windows-tubes}:
\begin{itemize}
\item the cutoff $K$, the depth $s$ and the level $k=K-s$;
\item the depth-$s$ torus $\mathbb{T}_s$ and the fine lattice $T_{\eta}$ of spacing $\eta=L^{-k}$;
\item the configuration space $\mathcal{G}_s$ of $SU(N)$-valued bond variables on $\mathbb{T}_s$;
\item for $0\le t\le k-1$, the scale $j=k-t$, the $j$-lattice $T^{(j)}_1$ and Ba{\l}aban's
partition $\pi_j$ of $T^{(j)}_1$ into $M$-cubes.
\end{itemize}
Write $x_k=g_k^{-2}$. For $V\in\mathcal{G}_s$, let $U_k(V)$ be the background configuration
determined by $V$ (Definition~\ref{1.5:def-domains}). Let
$\mathbf{E}^{(j)}(X,\cdot,z)$ and $\mathbf{R}^{(j)}(X,\cdot)$ be the local effective-action pieces
of scale $j$ of \cite[(2.23), (2.25)--(2.27), (2.30)]{B-CMP119}, let $\beta_j(\cdot)$ be
Ba{\l}aban's coefficient function of scale $j$ in those displays, whose value $\beta_j(g_{j-1})$
at the coupling $g_{j-1}$ is the coefficient of the term proportional to the action
$A(U_k(V))$ that accompanies the scale-$j$ pieces, and let $\mathbf{D}_j$ be
the family of connected unions of closed cubes of $\pi_j$. We put
\begin{equation}\label{4.31:eq-never-large-exponent-1}
\mathcal{A}_{K,s}(V):=A(U_k(V))=\sum_{p\subset T_{\eta}}\bigl[1-\operatorname{Re}\operatorname{tr}
U_k(V)(\partial p)\bigr],
\end{equation}
and, for $0\le t\le k-1$ and $j=k-t$,
\begin{equation}\label{4.31:eq-never-large-exponent-2}
\begin{aligned}
\mathsf{E}_{K,s,t}(V)&:=\sum_{z\in T^{(j)}_1}\ \sum_{X\in\mathbf{D}_j,\,X\ni z}
\bigl[\mathbf{E}^{(j)}(X,U_k(V),z)-\mathbf{E}^{(j)}(X,\mathbf{1},z)\bigr]\\
&\qquad-\beta_j(g_{j-1})\,\mathcal{A}_{K,s}(V),\\
\mathsf{R}_{K,s,t}(V)&:=\sum_{X\in\mathbf{D}_j}\bigl[\mathbf{R}^{(j)}(X,U_k(V))
-\mathbf{R}^{(j)}(X,\mathbf{1})\bigr].
\end{aligned}
\end{equation}
The \emph{exponent of the never-large term at depth $s$} is
\begin{equation}\label{4.31:eq-never-large-exponent-3}
\mathsf{S}_{K,s}(V):=-x_k\,\mathcal{A}_{K,s}(V)
+\sum_{t=0}^{k-1}\bigl(\mathsf{E}_{K,s,t}(V)+\mathsf{R}_{K,s,t}(V)\bigr).
\end{equation}
Let $h_{\emptyset}$ be the history in the expansion \cite[(2.18)]{B-CMP119} in which
$\Omega_j=\Lambda_j=T$ for every $j$, where $T$ is the whole torus on which that expansion is
written. The term of $h_{\emptyset}$ in \cite[(2.18)]{B-CMP119}, the \emph{never-large term}, is
\begin{equation}\label{4.31:eq-never-large-exponent-4}
\chi_k(V)\,\exp\bigl[\mathsf{S}_{K,s}(V)-E^{\emptyset}_{K,s}\bigr].
\end{equation}
Here $E^{\emptyset}_{K,s}$ is a constant, independent of $V$. The subtractions at $\mathbf{1}$ in
\eqref{4.31:eq-never-large-exponent-2} normalise the exponent so that
$\mathsf{S}_{K,s}(\mathbf{1})=0$.

\emph{Index set.} In the units of $\mathbb{T}_s$, the family $\mathbf{D}_j$ is the family
$\mathbf{D}_s^{[t]}$ of connected unions of closed cubes of side $ML^{-t}$ of $\pi_j$. The
variable $z$ runs over $T^{(j)}_1$, the lattice of spacing $L^{-t}$ in the units of
$\mathbb{T}_s$. This index set is countable and does not
depend on $K$. The piece indexed by $(t,X,z)$ is present exactly when $K\ge s+t+1$. This is a
condition for the piece to exist and places no upper bound on $K$. An absent piece is set equal
to $0$. Finally, $\mathbf{R}^{(j)}(X):=0$ unless $X$ is a union of cubes of side $MR_j$ in
scale-$j$ units, that is, of side $R_j$ times the side $M$ of the cubes of $\pi_j$, where $R_j$
is a parameter of scale $j$ of \cite{B-CMP119}.
\end{definition}

\begin{proof}
We verify the three assertions made above: that the term of $h_{\emptyset}$ in
\cite[(2.18)]{B-CMP119} is \eqref{4.31:eq-never-large-exponent-4}, that
$\mathsf{S}_{K,s}(\mathbf{1})=0$, and that the index set does not depend on $K$.

For the history $h_{\emptyset}$ every $\Omega_j$ and every $\Lambda_j$ is the whole torus. The
expansion of \cite{B-CMP119} admits the possibility $\Lambda_j=\Omega_j$, so $h_{\emptyset}$ is
one of its histories. For this history each large-field region $Z_k$ is empty.

For this history the product in \cite[(2.19)]{B-CMP119} defining $\mathbf{T}_k$ is empty, and
$\mathbf{T}_k=1$.

The term $\mathbf{B}_k$ of \cite[(2.41)]{B-CMP119} is a sum over sets $X$ with
$X\cap Z_j^{\sim}\neq\emptyset$, where $Z_j^{\sim}$ is the enlargement of $Z_j$ used there.
Since $Z_j=\emptyset$, this sum has no terms, and $\mathbf{B}_k=0$.

For this history the functions $\phi_j$ are identically $1$. By \cite[(2.24)]{B-CMP119}, the
position-dependent coupling $g_k(\cdot)$ is then the constant $g_k$.
Moreover $\Lambda_j^{0}=T^{(j)}_1$.

We insert these values into \cite[(2.23), (2.25)--(2.27), (2.30)]{B-CMP119}. The term of
$h_{\emptyset}$ becomes the product of $\chi_k(V)$ with the exponential of the following sum:
\begin{itemize}
\item the action $-x_k A(U_k(V))$;
\item for each scale $j=k-t$, $0\le t\le k-1$, the localised pieces
$\mathbf{E}^{(j)}(X,U_k(V),z)$, summed over $z\in T^{(j)}_1$ and $X\in\mathbf{D}_j$ with $X\ni z$,
together with the coupling term $-\beta_j(g_{j-1})A(U_k(V))$;
\item for the same scales, the pieces $\mathbf{R}^{(j)}(X,U_k(V))$, summed over
$X\in\mathbf{D}_j$;
\item a constant.
\end{itemize}
We subtract from each localised piece its value at the identity configuration, as in
\eqref{4.31:eq-never-large-exponent-2}, and collect the subtracted values together with the
constant into $-E^{\emptyset}_{K,s}$. The result is \eqref{4.31:eq-never-large-exponent-4} with
$\mathsf{S}_{K,s}$ given by \eqref{4.31:eq-never-large-exponent-3}, and $E^{\emptyset}_{K,s}$ does
not depend on $V$.

For the normalisation we evaluate at $V=\mathbf{1}$. The identity configuration has action
$A(\mathbf{1})=0$, the least value of $A$, and its block average is $\mathbf{1}$. The background
configuration determined by $V=\mathbf{1}$ is the minimiser of $A$ under this block-average
constraint, and this constrained minimiser is unique in the gauge fixed in
Definition~\ref{1.5:def-domains}; hence $U_k(\mathbf{1})=\mathbf{1}$.
Then every bracket in \eqref{4.31:eq-never-large-exponent-2} is the difference of two equal
numbers and vanishes. Also $\mathcal{A}_{K,s}(\mathbf{1})=0$, because
$\operatorname{tr}\mathbf{1}=1$ makes each summand of \eqref{4.31:eq-never-large-exponent-1}
zero. Hence $\mathsf{S}_{K,s}(\mathbf{1})=0$.

For the index set, the scale $j=k-t$ ranges over $j\ge1$, that is, $t\le k-1$. Since $k=K-s$,
the inequality $k-t\ge1$ is equivalent to $K\ge s+t+1$. For fixed $t$, the lattice of spacing
$L^{-t}$ and the cubes of side $ML^{-t}$ in $\mathbb{T}_s$ units are determined by $t$, $M$, $L$,
$m$ and $s$ alone, as in Definition~\ref{4.31:def-windows-tubes}. Hence the index set does not
depend on $K$.
\end{proof}

\begin{definition}[Standing inputs: coupling flow, irrelevance and stability]
\label{4.31:def-standing-inputs}
In the remainder of this section the statements (a)--(e) below are the
\emph{standing inputs}. Each of them is taken by statement from the place indicated with it,
except that the all-pairs growth bound within (a) is an additional hypothesis, and the local
boundedness of $E_{\pm}$ in (d) is an assumption of this section. The
consequences stated inside (b), (c) and (e) are derived from those statements in the proof
below. Throughout, $G=SU(N)$ with $N\ge 2$ arbitrary, $x_l=g_l^{-2}$, $\xi=L^{-j}$ is the fine
spacing measured in units of scale $j$, and $\eta=L^{-k}$ is the spacing of the fine lattice
$T_{\eta}$ at level $k$.

\begin{itemize}
\item[(a)] \emph{Coupling flow.} There are constants $\lambda_{\sharp},B_{\sharp},c_2,C_{\beta}>0$,
independent of $K$, such that for every sequence $g_0,\dots,g_K$ produced by the flow with all
$g_l\in\,]0,\gamma]$
\begin{equation}\label{4.31:eq-standing-inputs-a}
\begin{aligned}
&u\lambda_{\sharp}-2c_2\le x_{K-u}-x_K\le uB_{\sharp}+2c_2,\\
&(j'-i)\lambda_{\sharp}-2c_2\le x_i-x_{j'}\le (j'-i)B_{\sharp}+2c_2,\\
&x_{l-1}-x_l=\beta_l(g_{l-1}),\qquad |\beta_l|\le C_{\beta}.
\end{aligned}
\end{equation}
The first line, the \emph{linear-growth bound}, holds for $0\le u\le K$; the second line,
the \emph{all-pairs growth bound}, holds for $0\le i<j'\le K$; the third line is the
\emph{beta-function bound}.
The linear-growth bound and the beta-function bound are the content of clause (0) of Theorem~0
(Theorem~\ref{4.1:thm-clause-zero}), with the constants $\lambda_{\sharp}$, $B_{\sharp}$ and
$c_5=2c_2$ of that theorem. The all-pairs growth bound, the same two-sided bound for all pairs
of levels, is used in this section only for the inequalities comparing different depths, and
those inequalities are asserted under the all-pairs growth bound; no other conclusion of this
section depends on it.
No positive lower bound (lower slope) on the single-step increment $x_{l-1}-x_l$ is assumed or
used.

\item[(b)] \emph{Form of the local effective-action pieces}
\cite[(1.8)--(1.19)]{B-CMP109}, \cite[(2.31)]{B-CMP119}. For every scale $j$ and every region
$X$, the pieces $\mathbf{E}^{(j)}(X,(\mathbf{U},\mathbf{J}),z)$ and
$\mathbf{R}^{(j)}(X,(\mathbf{U},\mathbf{J}))$ are analytic on the space
$\mathfrak{U}^{c}_j(X,\alpha_{0,j},\alpha_{1,j})$ of \cite[(1.10)--(1.16)]{B-CMP109}; they
depend on the pair only through its restriction $(\mathbf{U},\mathbf{J})|_X$; they are invariant
under all $G^{\mathbb{C}}$-valued gauge transformations of pairs \cite[(1.10),
(1.19)]{B-CMP109}, the space $\mathfrak{U}^{c}_j(X,\alpha_{0,j},\alpha_{1,j})$ being a union
of orbits; and
\begin{equation}\label{4.31:eq-standing-inputs-b}
\begin{aligned}
|\mathbf{E}^{(j)}(X,(\mathbf{U},\mathbf{J}),z)|&\le E_0\,e^{-\kappa d_j(X)},\\
|\mathbf{R}^{(j)}(X,(\mathbf{U},\mathbf{J}))|&\le g_j^{\kappa_0}\,e^{-\kappa d_j(X)},
\qquad \kappa_0\ge 7,
\end{aligned}
\end{equation}
the first by \cite[(1.18)]{B-CMP109} with $E_0,\alpha_0,\alpha_1$ absolute constants
independent of $X$ and $j$, the second by \cite[(2.31)]{B-CMP119}. For a configuration $U$ we
write $\mathbf{E}^{(j)}(X,U,z):=\mathbf{E}^{(j)}(X,(U,J^{\xi}(U)),z)$, where the current of $U$
is \cite[(1.8)--(1.9), (3.10)]{B-CMP109}
\begin{equation}\label{4.31:eq-standing-inputs-1}
J^{\xi}(U):=D_U^{\xi*}\,\xi^{-2}\,\pi\operatorname{Im}\partial U,
\end{equation}
with $D^{\xi*}_U$ the adjoint of the covariant lattice derivative of
\cite[(1.8)--(1.9)]{B-CMP109}, $\pi Y=Y-(\operatorname{tr}Y)\mathbf{1}$ the traceless
projection, and $\operatorname{Im}W=(W-W^{*})/(2i)$ on $G$, extended holomorphically to
$G^{\mathbb{C}}$ as $(W-W^{-1})/(2i)$. For every $G^{\mathbb{C}}$-valued gauge transformation
$g$,
\begin{equation}\label{4.31:eq-standing-inputs-2}
J^{\xi}(U^{g})=\operatorname{Ad}(g)\,J^{\xi}(U).
\end{equation}

\item[(c)] \emph{Irrelevance on real configurations}
\cite[Theorem~2, (2.43)]{B-CMP119}. For $\beta<1$ and for sufficiently regular configurations
$U_k=U_k(V)$, for instance for $V$ restricted by the characteristic functions in
\cite[(2.18)]{B-CMP119},
\begin{equation}\label{4.31:eq-standing-inputs-c}
\begin{split}
\Bigl|\sum_{z\in\Lambda_j^{0}\cap\Omega}\bigl[\mathbf{E}^{(j)}(\Lambda_j,U_k,z)
-\mathbf{E}^{(j)}(\Lambda_j,\mathbf{1},z)\bigr]&-\beta_j(g_{j-1})\,A(\phi,U_k)\Bigr|\\
&\le E_1\sum_{n=j}^{k}\bigl(L^{j-n}\bigr)^{\beta}\,|\Gamma_n\cap\Omega|,
\end{split}
\end{equation}
with $E_1$ independent of $j$, $k$, $\Omega$, $\{\Omega_j\}$, $\{\Lambda_j\}$ and the torus $T$
of \cite[Theorem~2]{B-CMP119}; here
$A(\phi,U_k)$ denotes the action of $U_k$ weighted by the function $\phi$, in the notation of
\cite[Theorem~2]{B-CMP119}. The bound \eqref{4.31:eq-standing-inputs-c} is used with
$\Omega=B^{j}(\Lambda_j^{0})$, the region determined by $\Lambda_j^{0}$ in
\cite[Theorem~2]{B-CMP119}, and with the weight $\phi=\phi_j$ there. The
exponent $\beta\in\,]0,1[$ is fixed once and for all; it is the exponent of
\cite[(2.43)]{B-CMP119}, not to be confused with the functions $\beta_l$ of
\eqref{4.31:eq-standing-inputs-a} nor with the schedule constant of
Definition~\ref{4.4:def-post-step-space}. Consequently, for the empty history
$h_{\emptyset}$ (in which $\Gamma_n=\emptyset$ for $n<k$ and $\Gamma_k=T^{(k)}_1$), for
$k=K-s$, $j=k-t$, and for every $V$ with $\chi_k(V)=1$,
\begin{equation}\label{4.31:eq-standing-inputs-3}
|\mathsf{E}_{K,s,t}(V)|\le E_1\,L^{-\beta t}\,|\mathbb{T}_s|,
\end{equation}
where $\mathbb{T}_s$ is the depth-$s$ unit lattice of Definition~\ref{4.31:def-windows-tubes}
and $\mathsf{E}_{K,s,t}$ is the piece of the exponent of
Definition~\ref{4.31:def-never-large-exponent}.

\item[(d)] \emph{Stability and equality of the integrals}
\cite[Theorem~1, (0.1)]{B-CMP122b}, used only in the following form: the lower member as an
inequality of measures, the upper member integrated at $k=K$, and the equality of the
integrals,
\begin{equation}\label{4.31:eq-standing-inputs-d}
\begin{gathered}
\rho_k\,dV\ \ge\ \chi_k\,e^{-x_kA(U_k)-E_{-}(x_k)|T_{\eta}|}\,dV\ \ge\ 0,\\
Z_K=\int\rho_K\,dV\le e^{E_{+}(x_K)|\mathbb{T}_0|},\qquad
\int\rho_{K,s}\,dV=Z_K\quad\text{for every }s;
\end{gathered}
\end{equation}
the first line rests on the lower member of \cite[Theorem~1, (0.1)]{B-CMP122b} together with
the positivity statement for the terms of the expansion \cite[(2.18)]{B-CMP119} (every term
is non-negative), which is taken by statement from elsewhere in this book; the upper member
enters only integrated, at $k=K$. The
constants $E_{\pm}=E_{\pm}(x_k)$ are independent of $\eta$ and of the torus but depend on $g_k$,
through logarithmic terms \cite{B-CMP119}. As part of (d) we assume that
$x\mapsto E_{\pm}(x)$ is locally bounded on $[\gamma^{-2},\infty[$; this is consistent with
the dependence of $E_{\pm}$ on $g_k$ through logarithmic terms just stated, and we put
\begin{equation*}
E_{\pm}^{(s)}:=\sup_{x\in I_s}E_{\pm}(x)<\infty,
\end{equation*}
where $I_s$ is the compact window, independent of $K$, for $x_{K-s}$ fixed earlier in this
section; the constants
$E_{\pm}^{(s)}$ are independent of $K$. The volume $|T_{\eta}|$ in
\eqref{4.31:eq-standing-inputs-d} is the number of unit cubes,
$|T_{\eta}|=|T^{(k)}_1|=|\mathbb{T}_s|$ at level $k=K-s$; it does not depend on the lattice
spacing \cite{B-CMP109}. A pointwise upper bound on $\rho_k$ is not part of (d); it is used in
this section only in a statement that assumes it explicitly.

\item[(e)] \emph{Real background fields} \cite[Theorem~1, (7)--(10)]{B-CMP102b}, with
$\Omega_0=\dots=\Omega_k=T_{\eta}$. If $|\partial V_0-1|<\varepsilon_1$ and
$\varepsilon_1\le a_1$, then the minimal orbit $U_k(V_0)$ exists and
\begin{equation}\label{4.31:eq-standing-inputs-e}
U_k(V_0)\in\mathfrak{U}_k(B_3\varepsilon_1),\qquad
|\partial U_k(V_0)-1|<B_3\varepsilon_1\eta^{2}.
\end{equation}
Moreover \cite[(9), p.~279]{B-CMP102b}, in which we write $M'$ for the cube parameter denoted
$M$ there and which we apply to $U=U_k(V_0)$, states the following: for a cube of side
$2M'L^{j}\eta$ whose
correspondingly enlarged cube lies in the domains, with $M'$ a multiple of $R_1M_1$ and
$M'\le M(\varepsilon_1)=R_1M_1a_1/\varepsilon_1$, there is a $G$-valued gauge
transformation $u_0$, written $u^{-1}$ in \cite[(9)]{B-CMP102b}, with
$U^{u_0}=e^{i\eta\mathsf{a}}$ on the cube, where $\mathsf{a}$ is a
$\mathfrak{g}$-valued field on the bonds of the cube, written $A$ in \cite[(9)]{B-CMP102b}, and
\begin{equation}\label{4.31:eq-standing-inputs-4}
|\mathsf{a}|<B_3M'\varepsilon_1\,(L^{j}\eta)^{-1},\qquad
|\nabla^{\eta}\mathsf{a}|<B_3M'\varepsilon_1\,(L^{j}\eta)^{-2}.
\end{equation}
Consequently, every cube $\square\subset\mathbb{T}_s$ of side $2M'$ with corners on the grid
of mesh $R_1M_1$, where $M'$ is a multiple of $R_1M_1$ and $M'\le M(\varepsilon_1)$, admits a
$G$-valued $u_0$ and a $\mathfrak{g}$-valued field $\mathsf{a}_{\square}$ on the bonds of
$T_{\eta}$ lying in $\square$ with
\begin{equation}\label{4.31:eq-standing-inputs-5}
U_k(V_0)^{u_0}=e^{i\eta\mathsf{a}_{\square}}\ \text{on }\square,\qquad
|\mathsf{a}_{\square}|<B_3M'\varepsilon_1,\quad
|\nabla^{\eta}\mathsf{a}_{\square}|<B_3M'\varepsilon_1.
\end{equation}
\end{itemize}
\end{definition}

\begin{proof}[Derivation of the consequences stated in (b), (c) and (e)]
\emph{The covariance \eqref{4.31:eq-standing-inputs-2}.} Let $g$ be $G^{\mathbb{C}}$-valued.
The plaquette variables transform by conjugation, $\partial(U^{g})=\operatorname{Ad}(g)\,\partial U$
at the base point of each plaquette. The holomorphic extension of $\operatorname{Im}$ on the
plaquette variables is $\operatorname{Im}W=(W-W^{-1})/(2i)$, and
\begin{equation*}
\frac{gWg^{-1}-gW^{-1}g^{-1}}{2i}=g\,\Bigl(\frac{W-W^{-1}}{2i}\Bigr)\,g^{-1},
\end{equation*}
so it commutes with
$\operatorname{Ad}(g)$. The projection $\pi$ is $\operatorname{Ad}$-equivariant, because
$\operatorname{tr}(gYg^{-1})=\operatorname{tr}Y$ gives
$\pi(gYg^{-1})=gYg^{-1}-(\operatorname{tr}Y)\mathbf{1}=g(\pi Y)g^{-1}$. Finally, for a bond
$b=\langle x,y\rangle$ the bond variable transforms as $U^{g}(b)=g(x)U(b)g(y)^{-1}$, so that the
product $U(\ell)$ of bond variables along a lattice path $\ell$ from $x$ to $y$ transforms as
$U^{g}(\ell)=g(x)U(\ell)g(y)^{-1}$. The covariant derivative of \cite[(1.8)--(1.9)]{B-CMP109}
takes fields on bonds to fields on plaquettes, so its adjoint $D^{\xi*}_{U}$ takes a field $Y$
on plaquettes, such as $\xi^{-2}\pi\operatorname{Im}\partial U$, to a field on bonds. Its value
at a bond $b$ with base point $x_b$ is a signed sum, with a factor depending only on $\xi$, over
the plaquettes $p$ containing $b$, of $U(\ell_{b,p})\,Y(p)\,U(\ell_{b,p})^{-1}$, where
$\ell_{b,p}$ is the path along $\partial p$ from $x_b$ to the base point $x_p$ of $p$ at which
$Y(p)$ is based. Replacing $U$ by $U^{g}$ and $Y$ by $\operatorname{Ad}(g)Y$, that is, $Y(p)$
by $g(x_p)Y(p)g(x_p)^{-1}$, turns each such term into
\begin{equation*}
g(x_b)\,U(\ell_{b,p})\,Y(p)\,U(\ell_{b,p})^{-1}\,g(x_b)^{-1},
\end{equation*}
since the factors $g(x_p)$ and $g(x_p)^{-1}$ cancel. Hence
$D^{\xi*}_{U^{g}}\operatorname{Ad}(g)=\operatorname{Ad}(g)D^{\xi*}_{U}$. Applying these
facts in turn to \eqref{4.31:eq-standing-inputs-1} --- first
$\partial(U^{g})=\operatorname{Ad}(g)\,\partial U$, then the three commutation relations for
$\operatorname{Im}$, $\pi$ and $D^{\xi*}$ --- gives \eqref{4.31:eq-standing-inputs-2}. No
step uses a particular value of $N$.

\emph{The bound \eqref{4.31:eq-standing-inputs-3}.} Let $k=K-s$, $j=k-t$, and let $V$ satisfy
$\chi_k(V)=1$. Then \eqref{4.31:eq-standing-inputs-c} applies with the empty history
$h_{\emptyset}$, in which $\Gamma_n=\emptyset$ for $j\le n<k$ and $\Gamma_k=T^{(k)}_1$, at the
pair $\Omega=B^{j}(\Lambda_j^{0})$, $\phi=\phi_j$; the quantity inside the absolute value on
its left-hand side is then $\mathsf{E}_{K,s,t}(V)$, by
Definition~\ref{4.31:def-never-large-exponent}. On the right-hand side every term with
$n<k$ vanishes, and the term $n=k$ is
\begin{equation*}
E_1\bigl(L^{j-k}\bigr)^{\beta}|\Gamma_k\cap\Omega|=E_1L^{-\beta t}\,|T^{(k)}_1\cap\Omega|
\le E_1L^{-\beta t}\,|T^{(k)}_1|,
\end{equation*}
since $k-j=t$. By the volume identification in (d), $|T^{(k)}_1|=|\mathbb{T}_s|$, which gives
\eqref{4.31:eq-standing-inputs-3}.

\emph{The class cubes \eqref{4.31:eq-standing-inputs-5}.} Take $j=k$ in
\eqref{4.31:eq-standing-inputs-4}. Since $\eta=L^{-k}$, we have $L^{k}\eta=1$, so the cube has
side $2M'$ in units of the unit lattice, and both bounds of \eqref{4.31:eq-standing-inputs-4}
become $B_3M'\varepsilon_1$; with the same $u_0$ and with $\mathsf{a}_{\square}=\mathsf{a}$
this is
\eqref{4.31:eq-standing-inputs-5} for the cubes to which \cite[(9)]{B-CMP102b} applies. Since
all domains equal the whole torus, $\Omega_j=T_{\eta}$ for every $j$, the condition that the
enlarged cube lie in the domains holds for every cube, and the statement made for the cubes of
one fixed partition applies equally to the partition translated by any vector of the grid of
mesh
$R_1M_1$: the construction of \cite[Theorem~1]{B-CMP102b} can be carried out with the
translated partition, and by uniqueness of the minimal orbit it yields the same orbit
$U_k(V_0)$. Every cube $\square$ of side $2M'$ with corners on that grid belongs to one such
translated partition, and \eqref{4.31:eq-standing-inputs-5} follows for it.
\end{proof}

\begin{definition}[The complex background field, global and local forms]\label{4.31:def-complex-background}
Throughout, $G=SU(N)$. The space $\mathfrak{g}$ consists of the traceless Hermitian $N\times N$ matrices, and
$\mathfrak{g}^{\mathbb{C}}=\mathfrak{g}\oplus i\mathfrak{g}=\mathfrak{sl}(N,\mathbb{C})$. Sup norms of matrix-valued
fields are taken bondwise. For a constant $v\in G$, a configuration $U$ and a
$\mathfrak{g}^{\mathbb{C}}$-valued field $D$ we write $U^{v}(b)=vU(b)v^{-1}$ and
$(\operatorname{Ad}(v)D)(b)=vD(b)v^{-1}$ for every bond $b$. The lattices $T_{\eta}$, $\eta=L^{-k}$, and
$\mathbb{T}_s$ are those of Definition~\ref{4.31:def-windows-tubes}.

\emph{Global form.} Let $\alpha_0=\alpha_1>0$, and let $V_0$ be a real top-level datum satisfying the
small-field conditions of the real background field with parameter $\varepsilon_1=\alpha_0$; put
$U_0:=U_k(V_0)$, the real minimal configuration (real background field) of the datum $V_0$. Let $\beta>0$ and
$B\in(\mathfrak{g}^{\mathbb{C}})^{\mathrm{bonds}(\mathbb{T}_s)}$ with $\sup|B|\le\beta$, and put
$\rho:=\alpha_0+\alpha_1+2B_0\beta$. Assume $\rho\le\rho_W$ and the side condition of the construction of the
complex background field
\begin{equation*}
MB_3\alpha_0+C_{19}\rho\le\min\{a_{\mathrm{op}},a_J,\hat{a}_1\},
\end{equation*}
where $a_{\mathrm{op}}$, $a_J$ and $\hat{a}_1$ are the three side thresholds of that construction and $B_3$ is
the constant of Ba{\l}aban's small-field bounds that enters it.
Consider the $k$-step minimisation problem on $T_{\eta}\supset\mathbb{T}_s$ with top-level datum $e^{iB}V_0$.
By the construction of the complex background field (its existence theorem, its gauge-equivalence statement
for real data, and its bounds on the exponent, on the plaquette variables and on the current), there is a
configuration $W=e^{i\eta\mathcal{H}(B;U_0)}U_0$, the \emph{complex background field}, holomorphic in $B$ and
lying in the complex small-field space $\mathfrak{U}^{\mathbb{C}}_k$ of the construction with radii
$B_3\alpha_0$ and $C_{19}\rho$; writing $\nabla^{\eta}_{U_0}$ for the covariant lattice derivative at spacing
$\eta$ in the background $U_0$, it satisfies
\begin{equation}\label{4.31:eq-complex-background-a}
\begin{gathered}
|\mathcal{H}(B;U_0)|,\ |\nabla^{\eta}_{U_0}\mathcal{H}(B;U_0)|\le C_{19}\rho,\qquad
|W(\partial p)-1|\le 6(B_3\alpha_0+C_{19}\rho)\eta^{2},\\
W\in\mathfrak{U}^{\mathbb{C}}_k(B_3\alpha_0,C_{19}\rho),\qquad
|J^{\eta}(W)|,\ |J^{\eta}(U_0)|\le 4(B_3\alpha_0+C_{19}\rho).
\end{gathered}
\end{equation}
For $\mathfrak{g}$-valued $B$, $W$ is equivalent under a $G$-valued gauge transformation to $U_k(e^{iB}V_0)$, the
datum of the latter problem being $e^{iB}V_0$ itself, with no seam bonds. The constants $\rho_W$, $C_{19}$, $B_0$
depend only on $(d,L,M)$. They are independent of $k$; neither the terrace levels, nor the torus, nor positions
enter them.

\emph{Local form.} Let $\square_0\subset T^{(j)}_1$ be an aligned scale-$j$ cube, that is, a union of cubes of
the partition $\pi_j$, and let $\mathfrak{B}(\square_0)$ be the determining set based on $\square_0$; $\xi$ is the
fine spacing of the scale-$j$ problem. For a union $Y$ of such cubes and $n\ge0$, $Y^{\sim n}$ denotes $Y$
enlarged by $n$ layers of cubes and $Y^{\sim -n}$ denotes $Y$ with $n$ layers of cubes removed. Let $\beta_W>0$
depend only on $(d,L,M)$, chosen so that $2B_0\beta_W\le\rho_W$ and so that the side thresholds of the global
form and of the nesting statements for small complex minimisers stated with
\eqref{4.31:eq-window-counting-nesting-e} hold. Let $D$ be a
$\mathfrak{g}^{\mathbb{C}}$-valued function on the bonds of $\mathfrak{B}(\square_0)$ on which the datum is
prescribed (its terrace bonds), with $\sup|D|\le\beta\le\beta_W$. The local minimiser and its exponent satisfy
\begin{equation}\label{4.31:eq-complex-background-b}
\begin{gathered}
\mathbf{U}[D]=e^{i\xi\mathcal{H}_j(\square_0;D,\mathbf{1})}\ \text{on}\ \mathfrak{B}(\square_0),\qquad
D\mapsto\mathcal{H}_j(\square_0;D,\mathbf{1})\ \text{holomorphic on}\ \{\sup|D|<\beta_W\},\\
\mathcal{H}_j(\square_0;\operatorname{Ad}(v)D,\mathbf{1})
=\operatorname{Ad}(v)\,\mathcal{H}_j(\square_0;D,\mathbf{1})\qquad(v\in G\ \text{constant}),
\end{gathered}
\end{equation}
in the following precise sense.
\begin{enumerate}
\item[(a)] The $\square_0$-local $j$-step minimisation problem with terrace datum $e^{iD}$ and base point the
identity configuration $\mathbf{1}$ has a small complex minimiser
$\mathbf{U}[D]=e^{i\xi\mathcal{H}_j(\square_0;D,\mathbf{1})}$ on $\mathfrak{B}(\square_0)$. It is unique among small
configurations satisfying the normalisation of the global construction relative to the base point $\mathbf{1}$.
In the top-level form, that is, with the letters of \eqref{4.31:eq-complex-background-a} specialised as in
\eqref{4.31:eq-complex-background-1} below, the following bounds hold on $\square_0^{\sim-2}$, which contains
$X^{\sim3}$ for every union $X$ of cubes of $\pi_j$ with $X^{\sim5}\subset\square_0$:
\begin{equation}\label{4.31:eq-complex-background-1}
\begin{gathered}
|\mathcal{H}_j|,\ |\nabla^{\xi}\mathcal{H}_j|\le 2C_{19}B_0\beta,\qquad
|\partial\mathbf{U}[D]-1|\le 12C_{19}B_0\beta\,\xi^{2},\\
|J^{\xi}(\mathbf{U}[D])|\le 8C_{19}B_0\beta.
\end{gathered}
\end{equation}
These are the bounds of \eqref{4.31:eq-complex-background-a} on the exponent, on the plaquette variables and on
$|J^{\eta}(W)|$, read at spacing $\xi$ and base point $\mathbf{1}$, with $\alpha_0=\alpha_1=0$ and
$\rho=2B_0\beta$, since $C_{19}\rho=2C_{19}B_0\beta$, $6C_{19}\rho=12C_{19}B_0\beta$ and
$4C_{19}\rho=8C_{19}B_0\beta$; the local form lists no counterpart of the membership clause or of the bound on
$|J^{\eta}(U_0)|$. Moreover
$\mathcal{H}_j(\square_0;0,\mathbf{1})=0$. For $\mathfrak{g}$-valued $D$, $\mathbf{U}[D]$ is unitary and represents
the real minimal orbit $U_j(\square_0,e^{iD})$. The constants depend only on $(d,L,M)$; they are independent of
the position and size of $\square_0$, of $j$ and of $K$.
\item[(b)] The map $D\mapsto\mathcal{H}_j(\square_0;D,\mathbf{1})$ is holomorphic on the open ball
$\{\sup|D|<\beta_W\}$.
\item[(c)] For every constant $v\in G$ and every $D$ in that ball, the covariance identity in the second line
of \eqref{4.31:eq-complex-background-b} holds.
\end{enumerate}
Items (a) and (b) are the local form of the construction of the complex background field. Item (c) is proved
below from (a), (b), the identity principle and the gauge covariance of the minimisers.

\emph{Identity principle.} In this paragraph and in the proof below, $V$, $\Omega$, $E$, $F_1$, $F_2$ are local
letters. Let $V$ be a finite-dimensional real vector space, $V^{\mathbb{C}}=V\oplus iV$, and $E$
a finite-dimensional complex vector space. The identity principle used below is
\begin{equation}\label{4.31:eq-complex-background-c}
\begin{gathered}
\Omega\subset V^{\mathbb{C}}\ \text{open, convex},\ \Omega\cap V\neq\emptyset,\
F_1,F_2\colon\Omega\to E\ \text{holomorphic},\\
F_1=F_2\ \text{on}\ \Omega\cap V\ \Longrightarrow\ F_1=F_2\ \text{on}\ \Omega.
\end{gathered}
\end{equation}
\end{definition}

\begin{proof}[Proof of item \textup{(c)}]
We begin by justifying \eqref{4.31:eq-complex-background-c}. At a point $x\in\Omega\cap V$, the set $\Omega\cap V$ is
open in $V$, so all real partial derivatives of all orders of $F_1-F_2$ along $V$ vanish at $x$. The complex
derivatives of all orders at $x$ are complex multilinear, and $V$ spans $V^{\mathbb{C}}$ over $\mathbb{C}$, so
they vanish too. The power series of $F_1-F_2$ at $x$ is therefore zero, so $F_1-F_2$ vanishes near $x$. Since
the convex set $\Omega$ is connected, the identity theorem gives $F_1-F_2=0$ on $\Omega$.

Fix a constant $v\in G$ and consider first $\mathfrak{g}$-valued $D$ with $\sup|D|<\beta_W$.

Since $\operatorname{Ad}(v)$ is conjugation by a unitary matrix, $\operatorname{Ad}(v)D$ is again
$\mathfrak{g}$-valued. Also $\sup|\operatorname{Ad}(v)D|=\sup|D|<\beta_W$, so $\mathbf{U}[\operatorname{Ad}(v)D]$
is defined by item (a). On every bond,
\begin{equation*}
e^{i\operatorname{Ad}(v)D}=v\,e^{iD}v^{-1}.
\end{equation*}
Hence the datum $e^{i\operatorname{Ad}(v)D}$ is the constant gauge transform $(e^{iD})^{v}$ of the datum
$e^{iD}$. The base point is unchanged, since $\mathbf{1}^{v}=v\mathbf{1}v^{-1}=\mathbf{1}$.

By the gauge covariance of the minimisers \eqref{4.31:eq-window-counting-nesting-e} (\cite[(181)]{B-CMP102b}),
for every $G$-valued gauge transformation $v$ of the datum lattice, with natural extension $\tilde{v}$, the
minimal configurations for the transformed datum are the $\tilde{v}$-transforms of those for the original
datum. This holds likewise for
the local problems, and $\tilde{v}=v$ for constant $v$. Applied to the $\square_0$-local $j$-step problem with
datum $e^{iD}$, it shows that the minimal configurations for $(e^{iD})^{v}$ are the $v$-conjugates of those for
$e^{iD}$. In particular $\mathbf{U}[D]^{v}$ is a minimal configuration for the datum
$e^{i\operatorname{Ad}(v)D}$, and
\begin{equation*}
\mathbf{U}[D]^{v}=v\,e^{i\xi\mathcal{H}_j(\square_0;D,\mathbf{1})}v^{-1}
=e^{i\xi\operatorname{Ad}(v)\mathcal{H}_j(\square_0;D,\mathbf{1})}.
\end{equation*}

The configuration $\mathbf{U}[D]$ is small and normalised relative to the base point $\mathbf{1}$ by item (a). We
read the uniqueness in item (a) as uniqueness within a class of small configurations, normalised relative to the
base point $\mathbf{1}$, that is carried into itself by conjugation with any constant $v\in G$: the smallness
conditions are sup-norm bounds, which $\operatorname{Ad}(v)$ preserves, and a linear gauge-covariant normalising
condition is covariant under $\operatorname{Ad}(v)$. With this reading $\mathbf{U}[D]^{v}$ is again small and
normalised relative to $\mathbf{1}$. The uniqueness in item (a), applied to the datum
$e^{i\operatorname{Ad}(v)D}$, therefore gives
\begin{equation}\label{4.31:eq-complex-background-2}
\mathcal{H}_j(\square_0;\operatorname{Ad}(v)D,\mathbf{1})
=\operatorname{Ad}(v)\,\mathcal{H}_j(\square_0;D,\mathbf{1})
\qquad\text{for every $\mathfrak{g}$-valued $D$ with $\sup|D|<\beta_W$.}
\end{equation}

We now pass to complex $D$. Let $\Omega=\{\sup|D|<\beta_W\}$, an open convex ball in the complex vector space
$V^{\mathbb{C}}$ of $\mathfrak{g}^{\mathbb{C}}$-valued functions on the terrace bonds of $\mathfrak{B}(\square_0)$.
Here $V$ is the real subspace of $\mathfrak{g}$-valued functions, and $V^{\mathbb{C}}=V\oplus iV$ because
$\mathfrak{g}^{\mathbb{C}}=\mathfrak{g}\oplus i\mathfrak{g}$. Thus $\Omega\cap V$ is the maximal totally real
slice of $\Omega$ consisting of the $\mathfrak{g}$-valued $D$, and it contains $D=0$. The map
$D\mapsto\operatorname{Ad}(v)D$ is complex linear, and it preserves the bondwise norm because $v$ is unitary, so
it is a linear isometry of $\Omega$ onto itself. Put
\begin{equation*}
F_1(D):=\mathcal{H}_j(\square_0;\operatorname{Ad}(v)D,\mathbf{1}),\qquad
F_2(D):=\operatorname{Ad}(v)\,\mathcal{H}_j(\square_0;D,\mathbf{1}),\qquad D\in\Omega.
\end{equation*}
By item (b), $F_1$ is holomorphic on $\Omega$ as the composition of the holomorphic map
$D\mapsto\mathcal{H}_j(\square_0;D,\mathbf{1})$ with a linear map of $\Omega$ into itself. Also by item (b),
$F_2$ is holomorphic on $\Omega$ as the image of that holomorphic map under the linear map $\operatorname{Ad}(v)$.
By \eqref{4.31:eq-complex-background-2}, $F_1=F_2$ on $\Omega\cap V$. The identity principle
\eqref{4.31:eq-complex-background-c} gives $F_1=F_2$ on $\Omega$, which is item (c).
\end{proof}

\begin{definition}[Window, counting and nesting of small complex minimisers]
\label{4.31:def-window-counting-nesting}
Throughout this section the following statements and conventions are in force, in addition to
those of Definitions~\ref{4.31:def-windows-tubes}, \ref{4.31:def-standing-inputs}
and~\ref{4.31:def-complex-background}. The depth-$s$ torus $\mathbb{T}_s$, the plaquette windows
$\mathcal{V}_s(w)$, the scale $j$ and the spacings $\xi$, $\eta$ are those of
Definition~\ref{4.31:def-windows-tubes}.

\begin{enumerate}
\item[(i)] \emph{Window.} Let $\varepsilon_k=\varepsilon(x_k)$ be the small-field function of
\cite[(2.4)]{B-CMP119}, let $B_3$ be the constant of the same name in \cite{B-CMP119}, and let $a_1$ be
the threshold of the same name in the choice of radii of this section. Then, for every $w>0$,
\begin{equation}\label{4.31:eq-window-counting-nesting-a}
B_3\,w\le\varepsilon_k\quad\text{and}\quad w\le a_1
\qquad\Longrightarrow\qquad
\chi_k\equiv 1\ \text{ on }\ \mathcal{V}_s(w).
\end{equation}
This is part~(a) of the window proposition for the small-field characteristic function $\chi_k$,
used here as stated there. The smallness condition $B_3w\le\varepsilon_k$
has the form of the threshold $B_3^{-1}\varepsilon_j$ of \cite[p.~380]{B-CMP122b}.

\item[(ii)] \emph{Counting.} Let $\mathbf{D}_j$ be the family of connected unions of the cubes of
Ba{\l}aban's partition $\pi_j$, let $d_j(X)$ be the tree length of $X\in\mathbf{D}_j$, and let $\kappa$
be the auxiliary parameter of that name. We put
\begin{equation}\label{4.31:eq-window-counting-nesting-b}
K_0:=\sup_{j}\ \sup_{\square}\ \sum_{X\in\mathbf{D}_j,\ X\supset\square}
e^{-(\kappa-1)\,d_j(X)}\ <\ \infty ,
\end{equation}
where $\square$ runs over the cubes of $\pi_j$. The finiteness of $K_0$ is the bound
\cite[(0.26)]{B-CMP109}. The factor $\exp(-(\kappa-1)d_j(X))$ is the one by which the sum over $X$
is controlled in \cite{B-CMP119}.

\item[(iii)] \emph{Nesting of small complex minimisers.} The following three statements are those
of the lemma on nesting of small complex minimisers, used here as stated there.

In them:
\begin{itemize}
\item $U_n(X,\cdot)$ and $J_n(X,\cdot)$ are the local minimisers and their currents of
\cite{B-CMP109};
\item $M^{n}$ is the $n$-fold average of \cite{B-CMP109}, and $M^{\cdot}(\cdot)$ denotes the averages
of a local problem;
\item $n$ ranges over the values allowed in \cite[(1.16)]{B-CMP109};
\item $J^{\xi}(U)$ is the current of $U$;
\item $X^{\sim5}$ is $X$ enlarged by five layers of cubes;
\item $X^{\sim-2}$ is the set obtained from $X$ by removing two layers of cubes. It is the domain on
which condition (iv) in the definition of the space $\mathfrak{U}^{c}_j(X,\alpha_0,\alpha_1)$ of
\cite[(1.10)--(1.16)]{B-CMP109} is imposed.
\end{itemize}
Each $\bar{u}$ below is a genuine $G^{\mathbb{C}}$-valued lattice gauge transformation
$x\mapsto\bar{u}(x)$ which is constant on each block. The letter $\bar{u}_1$ in (b)
denotes a $G^{\mathbb{C}}$-valued gauge transformation on $\square_0$, of which only the bound
$|\bar{u}_1^{\pm1}|\le2$ is asserted. The constant $c_u$ depends
only on $(d,L,M)$. The identities of the three statements read:
\begin{equation}\label{4.31:eq-window-counting-nesting-c}
\begin{aligned}
&\text{(a)}\quad U_n(X,M^{n}\mathbf{U}[D])=\mathbf{U}[D]^{\bar{u}^{-1}},\qquad
J_n(X,M^{n}\mathbf{U}[D])=(L^{n}\xi)^{3}\operatorname{Ad}(\bar{u}^{-1})\,J^{\xi}(\mathbf{U}[D]);\\
&\text{(b)}\quad W^{g}=\mathbf{U}[D]^{\bar{u}_1};\\
&\text{(c)}\quad U_n(X,M^{n}W)=W^{\bar{u}^{-1}},\qquad
J_n(X,M^{n}W)=(L^{n}\xi)^{3}\operatorname{Ad}(\bar{u}^{-1})\,J^{\xi}(W).
\end{aligned}
\end{equation}
\begin{itemize}
\item[(a)] \emph{Local form.} Let $\square_0$, $D$, $\mathbf{U}[D]$ and $\beta$ be as in
\eqref{4.31:eq-complex-background-b}. Let $X\in\mathbf{D}_j$ with $X^{\sim5}\subset\square_0$, and
let $n$ be in the range above. Then there is $\bar{u}$ with
$|\bar{u}^{\pm1}|\le 1+c_u\beta$ for which the first line of
\eqref{4.31:eq-window-counting-nesting-c} holds on $X^{\sim-2}$.

\item[(b)] \emph{From the global field to the local problem.} Let $W=W_{V_0}(B)$ be the complex
background field of Definition~\ref{4.31:def-complex-background}, taken with the parameters
$(w_s,r_s,\rho_s)$ of this section, $V_0\in\mathcal{V}_s(w_s)$ and $|B|<4r_s$. Let $\square_0$ be an
aligned scale-$j$ cube lying in one of the class cubes of
Definition~\ref{4.31:def-standing-inputs}. Let $g$ be a
$G^{\mathbb{C}}$-valued gauge transformation on $\square_0$ with $|g^{\pm1}|\le2$. Put
\begin{equation*}
D:=-i\log M^{\cdot}\bigl((W^{g})|_{\square_0}\bigr),
\end{equation*}
that is, $D$ is built from the averages of $W^{g}$ on the determining set
$\mathfrak{B}(\square_0)$ that the local problem itself uses. Assume that $\sup|D|\le\beta_W$.

Then there is $\bar{u}_1$ with $|\bar{u}_1^{\pm1}|\le2$ for which the second line of
\eqref{4.31:eq-window-counting-nesting-c} holds on the part of $\square_0$ that carries the
conditions and the determining sets of the sets $X\in\mathbf{D}_j$ with $X^{\sim5}\subset\square_0$.
Equivalently, for every $X\in\mathbf{D}_j$ with $X^{\sim5}\subset\square_0$ and every function $f$
of the local class of Definition~\ref{4.31:def-standing-inputs},
\begin{equation*}
f\bigl((W^{g},\operatorname{Ad}(g)J^{\xi}(W))|_{X}\bigr)
=f\bigl((\mathbf{U}[D],J^{\xi}(\mathbf{U}[D]))|_{X}\bigr).
\end{equation*}

\item[(c)] \emph{From the global field to all $X$.} Let $W$ be as in (b). Let $X\in\mathbf{D}_j$ be
arbitrary; large sets and sets that wrap around the torus are included. Let $n$ be in the range
above. Then there is $\bar{u}$ with
\begin{equation*}
|\bar{u}^{\pm1}|\le 1+c_u\bigl(B_3w_s+C_{19}\rho_s\bigr)
\end{equation*}
for which the third line of \eqref{4.31:eq-window-counting-nesting-c} holds on $X^{\sim-2}$.
\end{itemize}

Statement (b) is part of the lemma on nesting of small complex minimisers; the following reduction
of it is not used in what
follows. For $\mathfrak{g}$-valued $B$, the field $W$ and the gauge transformations of its
construction are unitary, and $D$ is $\mathfrak{g}$-valued. In that case (b) reduces to the locality
and nesting of real minimal configurations \cite[(1.2)--(1.3)]{B-CMP109}, \cite{B-CMP102b}: the global
minimiser, restricted, is the local minimiser of its own averages. The general case then follows from
three facts: both sides of the identity between values of $f$ are holomorphic in $B$ on the convex
set $\{|B|<4r_s\}$ (by the holomorphy statements of Definition~\ref{4.31:def-complex-background},
among them \eqref{4.31:eq-complex-background-b}, of the averaging lemma, and of the local class of
Definition~\ref{4.31:def-standing-inputs}); they agree for $\mathfrak{g}$-valued $B$; and
\eqref{4.31:eq-complex-background-c} holds.

For the family of local minimisers considered in \cite{B-CMP109}, the identity corresponding to (a)
is \cite[(3.38)]{B-CMP109}. The text there continues:
\begin{quote}
where $\bar{u}$ is a gauge transformation constant on blocks naturally connected with the function
$U_n(X,\cdot)$, and equal to $U_j$ at centers of the blocks. This identity and the bound on $U_j$ in
(3.37) imply that the condition (iv) is a consequence of the condition (iii), with a bit better
constant.
\end{quote}
Here (3.37) is \cite[(3.37)]{B-CMP109}, and (iii), (iv) are the conditions of those names in
\cite[(1.10)--(1.16)]{B-CMP109}.

\item[(iv)] \emph{Positivity and the never-large term.} Every term of \cite[(2.18)]{B-CMP119} is a
non-negative density. For the contribution $E^{\emptyset}_{K,s}$ of the empty large-field history
$h_{\emptyset}$, evaluated at $V=\mathbf{1}$,
\begin{equation}\label{4.31:eq-window-counting-nesting-d}
\bigl|E^{\emptyset}_{K,s}\bigr|\ \le\ E_{-}(x_k)\,|\mathbb{T}_s| ,
\end{equation}
where $|\mathbb{T}_s|$ is the number of sites of $\mathbb{T}_s$. Both statements are taken from the
positivity statement for the terms of \cite[(2.18)]{B-CMP119}, which includes a bound on the
never-large term; \eqref{4.31:eq-window-counting-nesting-d} is that bound read at $V=\mathbf{1}$,
where both the exponent $\mathsf{S}_{K,s}$ of the never-large term and the Wilson action vanish.

\item[(v)] \emph{Gauge covariance of the minimisers.} For every $G$-valued gauge transformation $v$
of the lattice on which the datum $V$ lives,
\begin{equation}\label{4.31:eq-window-counting-nesting-e}
U_k(V^{v})=U_k(V)^{\tilde{v}} .
\end{equation}
Here $\tilde{v}$ is the natural extension of $v$, and $\tilde{v}=v$ when $v$ is constant. The same
holds for the local problems $U_n(Y,\cdot)$. This is \cite[(181)]{B-CMP102b}.

\item[(vi)] \emph{Size of the torus.} The torus exponent $m$ satisfies
\begin{equation}\label{4.31:eq-window-counting-nesting-f}
\begin{aligned}
2L^{m}\ \ge\ &\text{the side of every class cube used in this section,}\\
&\text{enlarged as in \cite{B-CMP102b}.}
\end{aligned}
\end{equation}
These sides are at most $4\max\{M',c_{\square}LM\}$. Hence the class cubes fit into $\mathbb{T}_0$,
which is contained in every $\mathbb{T}_s$. Condition \eqref{4.31:eq-window-counting-nesting-f} does
not involve $K$. It is a one-sided lower bound on $m$, of the same kind as the floor
$m\ge m_{\mathbb{T}}$ on the torus exponent. Where it is not already implied by that floor, it is
imposed in addition, as a lower bound on $m$ alone. It is not a cap in the sense of
Definition~\ref{2.3:def-cap}, and it does not restrict the couplings. It is used once in this
section.

\item[(vii)] \emph{Classical facts.} We use the following classical facts:
\begin{itemize}
\item the Banach--Alaoglu theorem for a dual with separable predual;
\item the Riesz representation theorem;
\item the theorems of Montel and Weierstrass on complex manifolds;
\item the Cauchy estimates;
\item the statement \eqref{4.31:eq-complex-background-c};
\item the two-constants theorem, proved in this section as a lemma;
\item the polar decomposition;
\item Tannery's theorem;
\item the following consequence of Schur's lemma: for $N\ge2$, with $dv$ the Haar probability
measure on $SU(N)$,
\begin{equation}\label{4.31:eq-window-counting-nesting-g}
\int_{SU(N)}\operatorname{Ad}(v)\,Y\,dv=0\qquad\text{for every } Y\in\mathfrak{sl}(N,\mathbb{C}).
\end{equation}
Equivalently, the adjoint representation of the simple Lie algebra $\mathfrak{sl}(N,\mathbb{C})$ has
no non-zero invariant vectors.
\end{itemize}
\end{enumerate}
\end{definition}

\begin{proof}[Proof of \eqref{4.31:eq-window-counting-nesting-g}]
Let $Y\in\mathfrak{sl}(N,\mathbb{C})$ and put
\begin{equation*}
\bar{Y}:=\int_{SU(N)}\operatorname{Ad}(v)\,Y\,dv=\int_{SU(N)}vYv^{-1}\,dv .
\end{equation*}
The integrand is a continuous matrix-valued function on the compact group $SU(N)$, so $\bar{Y}$ is a
well-defined $N\times N$ matrix.

Fix $u\in SU(N)$. By linearity of the integral,
\begin{equation*}
u\bar{Y}u^{-1}=\int_{SU(N)}(uv)Y(uv)^{-1}\,dv .
\end{equation*}
By left invariance of the Haar measure under $v\mapsto uv$, the right-hand side equals $\bar{Y}$.
Hence $\bar{Y}$ commutes with every element of $SU(N)$, that is, with every operator of the defining
representation of $SU(N)$ on $\mathbb{C}^{N}$.

This representation is irreducible. By Schur's lemma, $\bar{Y}=c\,\mathbf{1}$ for some $c\in\mathbb{C}$.

The trace is linear and invariant under conjugation, so $\operatorname{tr}(vYv^{-1})=\operatorname{tr}Y=0$
for every $v$. Integrating against the probability measure $dv$ gives $\operatorname{tr}\bar{Y}=0$. On
the other hand $\operatorname{tr}\bar{Y}=c\operatorname{tr}\mathbf{1}=c$, since
$\operatorname{tr}\mathbf{1}=1$. Hence $c=0$ and $\bar{Y}=0$.

The reformulation follows because the average $\bar{Y}$ is, for any representation of the compact
group, the projection of $Y$ onto the vectors fixed by the group. By connectedness, these are the
vectors $Z$ with $[H,Z]=0$ for every $H\in\mathfrak{g}$; since $[H,Z]$ is complex-linear in $H$ and
$\mathfrak{sl}(N,\mathbb{C})=\mathfrak{g}\oplus i\,\mathfrak{g}$, these are the vectors $Z$ with
$[H,Z]=0$ for every $H\in\mathfrak{sl}(N,\mathbb{C})$. Hence $\bar{Y}=0$ for all $Y$
says exactly that $\mathfrak{sl}(N,\mathbb{C})$ has no non-zero vector invariant under the adjoint
action.
\end{proof}

\begin{definition}[The analyticity space, determining sets and averaging in closed form]
\label{4.31:def-analyticity-space}
Fix a scale $j$ and measure lengths in scale-$j$ units, so that the fine lattice has spacing
$\xi=L^{-j}$. The complexified group $G^{\mathbb{C}}$, the norms $|\cdot|$ and the gauge action on
pairs are those of Notation~\ref{4.31:not-complexified-group}, and the set $\mathcal{K}$ of
cutoffs is that of Definition~\ref{4.31:def-windows-tubes}. We write $\pi_j$ for Ba{\l}aban's
partition of the scale-$j$ lattice into cubes of side $M$. For a union $X$ of $\pi_j$-cubes,
$X^{\sim-2}$ denotes the set obtained from $X$ by taking away two layers of cubes, as in
\cite[\S1]{B-CMP109}. Items (a)--(c) below state the definitions of \cite{B-CMP109},
\cite{B-CMP119} and \cite{B-CMP98} in the form used below; each is part of the present
definition.

\smallskip
\noindent\textup{(a)} \emph{The analyticity space.} Let $X$ be a union of $\pi_j$-cubes, and let
$\alpha_0,\alpha_1>0$; in \cite[\S1]{B-CMP109} the radius $\alpha_0$ is also written $y_0$. Let
$c_{\square}$ be the absolute number written $O(1)$ in \cite[(1.12)]{B-CMP109}, and $B\ge1$ the
constant of \cite[(1.12)]{B-CMP109}. Let $\mathcal{N}_j\subset\{0,\dots,j\}$ be the range of levels of
\cite[(1.15)--(1.16)]{B-CMP109}; for $n=0$ the condition below indexed by $n$ is vacuous, and along
$\mathcal{K}$ the gap $j-n$ is not bounded. For $n\in\mathcal{N}_j$ and a configuration $\mathbf{U}$
on the bonds of $X$, let $M^{\cdot}(\mathbf{U})=(M^i(\mathbf{U}))_{i\le n}$ be the family of averages
of $\mathbf{U}$ at the levels $i\le n$ (item~(c); a superscripted $M^{i}$, $M^{n}$, $M^{\cdot}$
always denotes an averaging map of item~(c), never a power of the cube side $M$); let
$U_n(X,M^{n}\mathbf{U})$ be the $n$-step local
minimiser based on $X$ with this datum (a terrace in the sense of item~(b)), and let
\begin{equation}\label{4.31:eq-analyticity-space-1}
J_n(X,M^{n}\mathbf{U}):=(L^n\xi)^3\,J^{\xi}\bigl(U_n(X,M^{n}\mathbf{U})\bigr)
\end{equation}
be its current, written in the units of the $n$-th lattice (the current has homogeneity $3$); here
$J^{\xi}(U)$ is the current of a configuration $U$. In the conditions below, $u_{\square}$ denotes a
gauge transformation on the cube $\square$ whose existence is part of condition (i), $\nabla^{\xi}$ is
the lattice gradient of spacing $\xi$, and $\nabla^{\xi}_{U}$ is the lattice covariant derivative in
the background $U$. Then $\mathfrak{U}^{c}_j(X,\alpha_0,\alpha_1)$
\cite[(1.10)--(1.16)]{B-CMP109} is the set of pairs $(\mathbf{U},\mathbf{J})$ on $X$ such that
$\mathbf{U}=U'U$ with configurations $U$, $U'$ on the bonds of $X$ for which the following hold:
\begin{equation}\label{4.31:eq-analyticity-space-a}
\begin{aligned}
&\text{(i)}\ \ |\partial U-\mathbf{1}|<\alpha_0\xi^2,\quad\text{and on every cube }\square
  \text{ of scale } j \text{ and side } c_{\square}LM:\\
&\qquad U^{u_{\square}}=e^{i\xi A},\quad |A|<B\,c_{\square}LM\,\alpha_0,
  \quad |\nabla^{\xi}A|<B\,c_{\square}LM\,\alpha_0;\\
&\text{(ii)}\ \ U'=e^{i\xi A'},\quad |A'|<\alpha_1,\quad |\nabla^{\xi}_{U}A'|<\alpha_1;\\
&\text{(iii)}\ \ |\partial\mathbf{U}-\mathbf{1}|<\alpha_0\xi^2,\quad |\mathbf{J}|<\alpha_0,
  \quad |J^{\xi}(U)|<\alpha_0;\\
&\text{(iv)}\ \ \text{for every } n\in\mathcal{N}_j,\ \text{on } X^{\sim-2}:\\
&\qquad |\partial U_n(X,M^{n}\mathbf{U})-\mathbf{1}|<\alpha_0\xi^2,
  \quad |J_n(X,M^{n}\mathbf{U})|<\alpha_0(L^n\xi)^2 .
\end{aligned}
\end{equation}
Conditions (i), (ii), (iii), (iv) are \cite[(1.11)--(1.12)]{B-CMP109},
\cite[(1.13)]{B-CMP109}, \cite[(1.14)]{B-CMP109}, \cite[(1.15)--(1.16)]{B-CMP109} respectively.

The normalisation in which (iv) is written is the one of the display of \cite[\S1]{B-CMP109} for the
minimisers built from the averages of a configuration $U$,
\begin{equation}\label{4.31:eq-analyticity-space-2}
\begin{aligned}
|\partial U_n(M^{\cdot}(U))-\mathbf{1}|&<2B_3\alpha_0'L^{-2n}(L^n\xi)^2=2B_3\alpha_0'\xi^2,\\
|J_n(M^{\cdot}(U))|&<2B_3\alpha_0'(L^n\xi)^2\qquad\text{on } X^{\sim-2},
\end{aligned}
\end{equation}
with Ba{\l}aban's constant $B_3$ and primed radius $\alpha_0'$. This display is quoted only to
fix the units in which (iv) is written; condition (iv) is a defining condition of the space and is not
derived from \eqref{4.31:eq-analyticity-space-2}.

Every condition in \eqref{4.31:eq-analyticity-space-a} is a strict upper bound whose right-hand side
is non-decreasing in $\alpha_0$ (conditions (i), (iii), (iv)) or in $\alpha_1$ (condition (ii)).
Hence if $\alpha_0\le\alpha_0''$ and $\alpha_1\le\alpha_1''$, a pair satisfying the conditions with
$(\alpha_0,\alpha_1)$ satisfies them with $(\alpha_0'',\alpha_1'')$, the same factorisation
$\mathbf{U}=U'U$ and the same gauge transformations $u_{\square}$ serving, so that
\begin{equation}\label{4.31:eq-analyticity-space-3}
\mathfrak{U}^{c}_j(X,\alpha_0,\alpha_1)\subset\mathfrak{U}^{c}_j(X,\alpha_0'',\alpha_1'').
\end{equation}
This is the elementary half of the remark of \cite[\S1]{B-CMP109} on the dependence of the space on
its constants.

\smallskip
\noindent\textup{(b)} \emph{Determining sets} \cite[(1.2)--(1.3)]{B-CMP109},
\cite[(1.2)--(1.3)]{B-CMP119}. Let $Y$ be a union of $\pi_j$-cubes. The determining set
$\mathfrak{B}(Y)$ based on $Y$ is a terrace: a family $(\mathfrak{B}_i(Y))_i$ of sets of bonds,
$\mathfrak{B}_i(Y)$ consisting of bonds of level $i$. Its coarsest level lies on $Y^{\sim-2}$, in
accordance with \cite{B-CMP119}, where the coarsest level of the terrace is placed on the region with
two layers of cubes removed, and its finer levels lie towards $\partial Y$. In display form,
\begin{equation}\label{4.31:eq-analyticity-space-b}
\begin{gathered}
\mathfrak{B}_i(Y)\subset Y\ \text{for every level } i,\\
\text{the coarsest level of } \mathfrak{B}(Y) \text{ lies on } Y^{\sim-2},\\
U_n(Y,\cdot)\ \text{depends on data inside } Y \text{ only},\\
\text{the blocks entering } M^i \text{ on } \mathfrak{B}_i(Y) \text{ lie inside } Y .
\end{gathered}
\end{equation}
The last clause means: for every level $i$, the blocks entering the level-$i$ averages of item~(c) at
the bonds of $\mathfrak{B}_i(Y)$ lie inside $Y$, and hence so do the fine bonds on which these
averages depend.

\smallskip
\noindent\textup{(c)} \emph{Averaging in closed form} (the averaging $M^{\cdot}(U)$ of
\cite{B-CMP98}, \cite[\S1]{B-CMP109}). Let $i\le j$ and let $c$ be a bond of level $i$, of length
$L^i\xi\le1$ in scale-$j$ units, with endpoints $c_-$, $c_+$. For a configuration $U$ on the fine
bonds,
\begin{equation}\label{4.31:eq-analyticity-space-c}
M^i(U)(c)=P\Bigl(L^{-4i}\sum_{x}U(\sigma_{c,x})\Bigr),
\end{equation}
where $x$ runs over the base points of $c$ and the factor $L^{-4i}$ normalises the sum;
$\sigma_{c,x}$ is the centre-to-centre contour attached to the base point $x$, consisting of at most
$c_dL^i$ fine bonds and lying inside the double $i$-block of $c$; and $U(\sigma)$ is the ordered
product of the fine bond variables along $\sigma$. The map $P$ is Ba{\l}aban's projection onto $G$,
extended holomorphically to the set $\{V:|V-\mathbf{1}|\le r_P\}$, which contains $\mathbf{1}$; it
satisfies
\begin{equation}\label{4.31:eq-analyticity-space-4}
P(\mathbf{1})=\mathbf{1},\qquad |P(V)-\mathbf{1}|\le c_P|V-\mathbf{1}|,\qquad
P(gVh^{-1})=g\,P(V)\,h^{-1}\quad(g,h\in G),
\end{equation}
and, by the identity principle, the last identity persists for $g,h\in G^{\mathbb{C}}$ near $G$.
The numbers $c_d$, $c_P$, $r_P$ are Ba{\l}aban's absolute numbers; for $P$ given by the polar part
with the holomorphic adjoint one may take $c_P=2$ on $r_P=\tfrac14$. Consequently
$M^i(\mathbf{1})=\mathbf{1}$, and the averaging maps $M^i$ are gauge covariant: for a gauge
transformation $g$,
\begin{equation}\label{4.31:eq-analyticity-space-5}
M^i(U^{g})(c)=g(c_-)\,M^i(U)(c)\,g(c_+)^{-1}.
\end{equation}

If Ba{\l}aban's level-$i$ average is instead the $i$-fold iterate of one-step averages with a
projection at each intermediate step, then \eqref{4.31:eq-analyticity-space-c} is not a literal
transcription of it. In that reading the estimate by telescoping which, starting from
\eqref{4.31:eq-analyticity-space-c}, bounds the level-$i$ averages of a configuration
$e^{i\xi\mathbf{A}}$ with a small $\mathfrak{g}^{\mathbb{C}}$-valued potential $\mathbf{A}$ loses a
factor $(2d+1)^i$ and does not apply; a proof under that reading has to use the cancellation between
forward and backward transports, at the price of a one-sided lower bound on $L$. Under that reading
the following statement is then assumed: there is an absolute number $c_{\mathrm{av}}$ (the
averaging constant) such that for every box $Y$ of scale $j$ and every
$\mathfrak{g}^{\mathbb{C}}$-valued function $\mathbf{A}$ on the fine bonds of $Y$ with
$c_{\mathrm{av}}\sup|\mathbf{A}|\le1$, putting $U=e^{i\xi\mathbf{A}}$, one has for every level
$i\le j$ and every level-$i$ bond $c$ whose double $i$-block lies in $Y$ that $M^i(U)(c)=e^{iD(c)}$
with $D(c)\in\mathfrak{g}^{\mathbb{C}}$ and $|D(c)|\le c_{\mathrm{av}}L^i\xi\sup|\mathbf{A}|$; $D$
depends holomorphically on $\mathbf{A}$, and $D$ is $\mathfrak{g}$-valued when $\mathbf{A}$ is. This
statement is the conclusion of the lemma on terrace averages of a small field, stated below in this
section and proved there from \eqref{4.31:eq-analyticity-space-c}; if Ba{\l}aban's average is the
iterated one, that proof does not apply and the conclusion is assumed as just stated.
\end{definition}

For $u\ge 0$ we define the intervals \eqref{4.31:eq-fixed-depth-convergence-a} below, where $x_{-}$, $x_{+}$ are the endpoints of the $K$-free window of \eqref{4.31:eq-windows-tubes-a}, $\lambda$, $B$, $c_2$ are the constants of the coupling flow \eqref{4.31:eq-standing-inputs-a}, and $\gamma$ is the coupling-smallness parameter of the glossary (Notation~\ref{0.2:not-glossary}).
\begin{equation}\label{4.31:eq-fixed-depth-convergence-a}
\begin{gathered}
I_u := \bigl[x_{-}^{(u)},\,x_{+}^{(u)}\bigr],\\
x_{-}^{(u)} := \max\bigl\{\gamma^{-2},\; x_{-}+\lambda u-2c_2\bigr\},
\qquad
x_{+}^{(u)} := x_{+}+Bu+2c_2 .
\end{gathered}
\end{equation}
These intervals do not depend on $K$.

For $s\ge 0$ we write:
\begin{itemize}
\item $\mathbb{T}_s$ for the depth-$s$ torus of Definition~\ref{4.31:def-windows-tubes};
\item $\mathcal{G}_s$ for the product of copies of $SU(N)$ over the bonds of $\mathbb{T}_s$;
\item $\mathcal{G}_s^{\mathbb{C}}$ for its complexification;
\item $dV$ for the Haar probability measure on $\mathcal{G}_s$;
\item $\rho_{K,s}$ for the effective density $\rho_{K-s}$ at level $K-s$, read as a measure $\rho_{K,s}\,dV$ on $\mathcal{G}_s$.
\end{itemize}
We further put
\begin{equation*}
E_{\pm}^{(s)}:=\sup_{x\in I_s}E_{\pm}(x),
\end{equation*}
where $E_{\pm}$ are the functions of the stability input \eqref{4.31:eq-standing-inputs-d}.

\begin{proposition}[Convergence at fixed depth with uniform analyticity on a tube]
\label{4.31:prop-fixed-depth-convergence}
Let $\mathcal{K}$ be an infinite set of cutoffs $K$, and assume the following.
\begin{itemize}
\item The standing inputs \textup{(I1)}--\textup{(I4)} of Definition~\ref{4.31:def-standing-inputs}
hold along $\mathcal{K}$. These include the coupling flow \eqref{4.31:eq-standing-inputs-a},
the irrelevance bound \eqref{4.31:eq-standing-inputs-c} and the stability bound
\eqref{4.31:eq-standing-inputs-d}. Moreover, for every $K\in\mathcal{K}$ the window condition
\eqref{4.31:eq-windows-tubes-a} of Definition~\ref{4.31:def-windows-tubes} holds.
\item Input \textup{(I5)} holds: the hypothesis \cite[(9)]{B-CMP102b} of
\cite[Theorem~1]{B-CMP102b} is satisfied on every aligned cube $\square_0$ of the local form of
the complex background field (Definition~\ref{4.31:def-complex-background}),
so that the conclusion of that theorem holds there.
\item Inputs \textup{(I6)} and \textup{(I6$'$)} hold; that is, the statements on the complex
background field of Definition~\ref{4.31:def-complex-background} hold in the global form and in
the local form \textup{(W1)}--\textup{(W3)}.
\item Inputs \textup{(I7)}, \textup{(I8)} and \textup{(I9)} of
Definition~\ref{4.31:def-window-counting-nesting} hold. These are the window, the counting bound
\cite[(0.26)]{B-CMP109}, and the nesting identities \textup{(N1)}--\textup{(N3)} for small complex
minimisers. The statement \eqref{4.31:eq-window-counting-nesting-d} of the same definition also holds.
\item The natural extension $\tilde{v}$ of a gauge transformation of the datum lattice
satisfies \cite[(181)]{B-CMP102b}.
\item The torus exponent obeys the floor $m\ge m_{\mathbb{T}}$, with $m_{\mathbb{T}}$ as in the
glossary (Notation~\ref{0.2:not-glossary}).
\item The averaging lemma, an estimate on Ba{\l}aban's averaging operations, holds, these
operations being taken in the closed form of Definition~\ref{4.31:def-analyticity-space}.
\end{itemize}
Then there are numbers $w_s,r_s>0$ and $C_s<\infty$ $(s=0,1,2,\dots)$, depending only on
$d$, $L$, $M$, $\mathfrak{p}$, $x_{-}$, $x_{+}$, $m$ and $s$, and an infinite subset
$\mathcal{K}'\subset\mathcal{K}$, with the following property.

Put $\mathcal{V}_s:=\mathcal{V}_s(w_s)$ and $\mathfrak{D}_s:=\mathfrak{D}_s(w_s,r_s)$, the plaquette window and the complex tube of Definition~\ref{4.31:def-windows-tubes}. Let $\mathcal{A}_{K,s}$, $\mathsf{E}_{K,s,t}$, $\mathsf{R}_{K,s,t}$, $\mathsf{S}_{K,s}$ and $E^{\emptyset}_{K,s}$ be as in Definition~\ref{4.31:def-never-large-exponent}. Then the statements \textup{(D1)}--\textup{(D5)} below hold for every $s\ge 0$ simultaneously, all limits being taken as $K\to\infty$ in $\mathcal{K}'$.

\smallskip
\noindent\textup{(D1)} \emph{Couplings.} With $\beta_{K-s}(g_{K-s-1})=x_{K-s-1}-x_{K-s}$,
\begin{equation}\label{4.31:eq-fixed-depth-convergence-b}
x_{K-s}\to x_s^{\infty}\in I_s,
\qquad
\beta_{K-s}(g_{K-s-1})\to\beta_s^{\infty}:=x^{\infty}_{s+1}-x^{\infty}_s,
\qquad
|\beta_s^{\infty}|\le C_{\beta}.
\end{equation}
If, in addition, hypothesis \textup{(I1b)} of Definition~\ref{4.31:def-standing-inputs} holds (the all-pairs form of \eqref{4.31:eq-standing-inputs-a}), then for $s<s'$
\begin{equation*}
(s'-s)\lambda-2c_2\;\le\;x^{\infty}_{s'}-x^{\infty}_s\;\le\;(s'-s)B+2c_2 .
\end{equation*}

\smallskip
\noindent\textup{(D2)} \emph{Densities, as measures.} For all $K\in\mathcal{K}$, $\rho_{K,s}\,dV$ is a positive Borel measure on $\mathcal{G}_s$ of mass $Z_K$, with
\begin{equation*}
z_{*}\le Z_K\le e^{E_{+}^{(0)}|\mathbb{T}_0|},
\end{equation*}
where $z_{*}>0$ is independent of $K$ and $s$. Moreover $Z_K\to Z^{\infty}$. The normalised measures $\nu^{B}_{K,s}:=Z_K^{-1}\rho_{K,s}\,dV$ converge weakly, that is, against every continuous function $F$ on $\mathcal{G}_s$, to a measure $\nu_s^{\infty}$, and $\nu_s^{\infty}$ is a probability measure.

\smallskip
\noindent\textup{(D2$^{+}$)} \emph{Bounded densities; a conditional clause.} This proposition does not assert the hypothesis of this clause. Suppose that, for every $K\in\mathcal{K}$, $\rho_{K,s}$ is a function on $\mathcal{G}_s$ satisfying, $dV$-almost everywhere, the pointwise form
\begin{equation*}
\rho_{K,s}\le e^{E_{+}^{(s)}|\mathbb{T}_s|}
\end{equation*}
of the upper bound in \cite[(0.1)]{B-CMP122b}. Then $\rho_{K,s}\to\rho_s^{\infty}$ weak-$*$ in $L^{\infty}(\mathcal{G}_s,dV)$, with
\begin{equation*}
0\le\rho_s^{\infty}\le e^{E_{+}^{(s)}|\mathbb{T}_s|}
\qquad\text{and}\qquad
\nu_s^{\infty}=\frac{\rho_s^{\infty}\,dV}{Z^{\infty}}\ll dV .
\end{equation*}
Moreover $\nu^{B}_{K,s}(F)\to\nu_s^{\infty}(F)$ for every bounded Borel function $F$ on $\mathcal{G}_s$.

\smallskip
\noindent\textup{(D3)} \emph{Uniform analyticity on a fixed domain.} For every $K\in\mathcal{K}$ with $K>s$, put $k=K-s$. Then $\chi_{K-s}\equiv 1$ on $\mathcal{V}_s(w_s+16r_s)$. Define
\begin{equation*}
e^{K}_{t,X,z}(V):=\mathbf{E}^{(K-s-t)}\bigl(X,U_k(V),z\bigr),
\qquad
r^{K}_{t,X}(V):=\mathbf{R}^{(K-s-t)}\bigl(X,U_k(V)\bigr),
\end{equation*}
where $\mathbf{E}^{(j)}$ and $\mathbf{R}^{(j)}$ are Ba{\l}aban's local terms of scale $j$ in \cite[(2.18)]{B-CMP119}. The functions
\begin{equation*}
\mathcal{A}_{K,s},\quad e^{K}_{t,X,z},\quad r^{K}_{t,X},\quad \mathsf{E}_{K,s,t},\quad
\mathsf{R}_{K,s,t},\quad \mathsf{S}_{K,s}
\end{equation*}
are initially defined on $\mathcal{V}_s$; here $t$ ranges over the scale gaps $0\le t\le K-s-1$, and $X$ and $z\in X$ range over the domains and points that enter \cite[(2.18)]{B-CMP119} at scale $K-s-t$. These functions extend to holomorphic functions on $\mathfrak{D}_s$ which satisfy the bounds \textup{(a)}--\textup{(e)} below. The constants $C_{\mathcal{A}}$, $C_E$, $C_R$, $E_0$ and $C_s$ in these bounds are independent of $K$. In these bounds:
\begin{itemize}
\item $d(X)$ is the tree length of $X$ in units $ML^{-t}$, which does not depend on $K$;
\item $E_0$, $\kappa$, $\kappa_0$ are the constants of \cite[(2.18)--(2.43)]{B-CMP119};
\item $\beta\in(0,1)$ is the exponent of \eqref{4.31:eq-standing-inputs-c}; the letter $\beta$ without
subscript denotes this exponent, never a coupling increment $\beta_j$.
\end{itemize}
\begin{itemize}
\item[(a)] $|\mathcal{A}_{K,s}|\le C_{\mathcal{A}}(w_s+r_s)^2|\mathbb{T}_s|$;
\item[(b)] the local terms obey
\begin{equation*}
|e^{K}_{t,X,z}|\le E_0\,e^{-\kappa d(X)},
\qquad
|r^{K}_{t,X}|\le \bigl(x_{-}^{(s+t)}\bigr)^{-\kappa_0/2}e^{-\kappa d(X)} ;
\end{equation*}
\item[(c)] $|\mathsf{E}_{K,s,t}|\le C_E\bigl(x_{+}^{(s+t)}\bigr)^{1/2}L^{-\beta t/2}|\mathbb{T}_s|$;
\item[(d)] $|\mathsf{R}_{K,s,t}|\le C_R\bigl(x_{-}^{(s+t)}\bigr)^{1-\kappa_0/2}|\mathbb{T}_s|$;
\item[(e)] $|\mathsf{S}_{K,s}|\le C_s$.
\end{itemize}

\smallskip
\noindent\textup{(D4)} \emph{Convergence of the holomorphic extensions.} The following convergences hold uniformly on compact subsets of $\mathfrak{D}_s$, together with all derivatives:
\begin{gather*}
\mathcal{A}_{K,s}\to\mathcal{A}_s^{\infty},\qquad
e^{K}_{t,X,z}\to e^{\infty}_{t,X,z},\qquad
r^{K}_{t,X}\to r^{\infty}_{t,X}\quad\text{(every $(t,X,z)$)},\\
\mathsf{E}_{K,s,t}\to\mathsf{E}^{\infty}_{s,t},\qquad
\mathsf{R}_{K,s,t}\to\mathsf{R}^{\infty}_{s,t},\qquad
\mathsf{S}_{K,s}\to\mathsf{S}_s^{\infty}.
\end{gather*}
The limits have the following properties.
\begin{itemize}
\item They are holomorphic on $\mathfrak{D}_s$ and invariant under $SU(N)$-valued gauge transformations.
\item They obey the bounds \textup{(a)}--\textup{(e)}.
\item The limits $\mathcal{A}_s^{\infty}$, $\mathsf{E}^{\infty}_{s,t}$, $\mathsf{R}^{\infty}_{s,t}$ and $\mathsf{S}_s^{\infty}$ vanish at $\mathbf{1}$.
\item They have the form of \cite[(0.22)--(0.23)]{B-CMP109}:
\begin{align*}
\mathsf{S}_s^{\infty}&=-x_s^{\infty}\mathcal{A}_s^{\infty}
+\sum_{t\ge0}\bigl(\mathsf{E}^{\infty}_{s,t}+\mathsf{R}^{\infty}_{s,t}\bigr),\\
\mathsf{E}^{\infty}_{s,t}&=\sum_{z}\sum_{X\ni z}
\bigl[e^{\infty}_{t,X,z}-e^{\infty}_{t,X,z}(\mathbf{1})\bigr]
-\beta^{\infty}_{s+t}\mathcal{A}_s^{\infty},\\
\mathsf{R}^{\infty}_{s,t}&=\sum_{X}\bigl[r^{\infty}_{t,X}-r^{\infty}_{t,X}(\mathbf{1})\bigr].
\end{align*}
The series in $t$ converges uniformly on $\mathfrak{D}_s$.
\end{itemize}

\smallskip
\noindent\textup{(D5)} \emph{Lower bound on the limit measure.} As measures on $\mathcal{V}_s$,
\begin{equation*}
Z^{\infty}\nu_s^{\infty}\;\ge\;
\exp\bigl[-x_s^{\infty}\mathcal{A}_s^{\infty}-E_{-}^{(s)}|\mathbb{T}_s|\bigr]\,dV .
\end{equation*}
Moreover (this part rests on \eqref{4.31:eq-window-counting-nesting-d}), $e^{-E^{\emptyset}_{K,s}}\to c_s^{\infty}$ with
\begin{equation*}
e^{-E_{-}^{(s)}|\mathbb{T}_s|}\;\le\;c_s^{\infty}\;\le\;e^{E_{-}^{(s)}|\mathbb{T}_s|},
\end{equation*}
and, as measures on $\mathcal{V}_s$,
\begin{equation*}
Z^{\infty}\nu_s^{\infty}\;\ge\;c_s^{\infty}\,e^{\mathsf{S}_s^{\infty}}\,dV .
\end{equation*}
\end{proposition}

The proof occupies the sections that follow.

\begin{remark}[What fixed-depth convergence does not assert]\label{4.31:rem-fixed-depth-not-asserted}
Proposition~\ref{4.31:prop-fixed-depth-convergence} asserts the convergence statements
\eqref{4.31:eq-fixed-depth-convergence-b}, with the exception of the bounded-density clause
(D2$^{+}$) (item~(8) below). It does not assert any of the following.
\begin{enumerate}
\item Convergence of $\rho_{K,s}$ in any sense stronger than the one stated there. Analyticity
of $\rho_{K,s}$ is not asserted either: $\rho_{K,s}$ contains large-field characteristic functions
and, through the components of the second class, $\delta$-functions (item~(8)).
\item The identity
\begin{equation*}
Z^{\infty}\nu_s^{\infty}=c_s^{\infty}\,e^{\mathsf{S}_s^{\infty}}\,dV
\qquad\text{on } \mathcal{V}_s .
\end{equation*}
Large-field histories at finer scales contribute on $\mathcal{V}_s$. For their smallness the
bounds \cite[(1.79)--(1.89)]{B-CMP122b} are used as stated there and are not derived again.
\item Any relation between $\nu_s^{\infty}$ and $\nu_{s-1}^{\infty}$ beyond the equality of their
masses. Compare the words ``equivalent, but not equal'' in \cite{B-CMP122b}.
\item Uniqueness of the limits. Uniformity in $s$ is not asserted either. Write $x_{+}^{(s)}$ for
the upper end of the window $I_s$ of (D1). The numbers $w_s$ and $r_s$ decrease to $0$ like
$(x_{+}^{(s)})^{-1/2}$, and $C_s$ is proportional to $|\mathbb{T}_s|$.
\item Anything about $T^{k}\rho_0$ or about $\mu_K$. Here $T^{k}\rho_0$ denotes the density
obtained from the bare density $\rho_0$ by $k$ renormalisation steps, and $\mu_K$ the lattice
measure at cutoff $K$.
\item Invariance of $\mathsf{S}_s^{\infty}$ under lattice translations and lattice isometries
beyond what clause~(i), assumed in Proposition~\ref{4.31:prop-fixed-depth-convergence}, asserts.
Compare \cite[(2.32)]{B-CMP119}.
\item The remark \cite[(1.17)]{B-CMP109} in its complex form, taken in general. By this we mean
the implication from conditions (i)--(iii) there, read for complex configurations, to condition
(iv) there, for configurations that are not minimisers. It is neither asserted nor used.
\item The bounded-density clause (D2$^{+}$). It is conditional on the pointwise bound
\begin{equation*}
\rho_{K,s}\le e^{E_{+}^{(s)}|\mathbb{T}_s|}\qquad dV\text{-almost everywhere}.
\end{equation*}
After the large-field operation $\mathbf{R}$, the components of the second class carry
$\delta$-functions in the new field variables $V_k$ \cite{B-CMP122b}. Hence $\rho_k\,dV_k$ has a
part singular with respect to Haar measure. The estimate in the proof of
\cite[(2.50)]{B-CMP119} is a bound on the integral $\int dV_k\,\rho_k$.
\end{enumerate}
\end{remark}

\begin{lemma}[Cutoff-free coupling windows and summable scale majorants]
\label{4.31:lem-cutoff-free-windows}
Let $\mathcal{K}$ be the infinite set of cutoffs along which the standing inputs
\eqref{4.31:eq-standing-inputs-a} and the window \eqref{4.31:eq-windows-tubes-a} hold, and for
$u\ge0$ let $I_u=[x_{-}^{(u)},x_{+}^{(u)}]$ be the interval of \eqref{4.31:eq-fixed-depth-convergence-a}.
Let $\varepsilon(x)=x^{-1/2}A_0(\log x)^{p_0}$ and $\alpha_i(x)=x^{-1/2}C_i(\log x)^{q_i}$, $i=0,1$,
$q_i\ge0$, be the functions of \cite[(2.4), (2.28)]{B-CMP119}, and let $B_3\ge1$ be the constant,
as fixed in Notation~\ref{4.31:not-complexified-group}. Throughout, $s,t,u$ are non-negative integers.
\begin{enumerate}
\item[(a)] For every $K\in\mathcal{K}$ and every $0\le u\le K$ one has $x_{K-u}\in I_u$.
\item[(b)] The numbers
\begin{equation}\label{4.31:eq-cutoff-free-windows-a}
\begin{gathered}
\varepsilon_u^{-}:=\min_{x\in I_u}\varepsilon(x),\qquad
\alpha_u^{\flat}:=\min_{x\in I_u}\min\Bigl\{\frac{\alpha_0(x)}{2B_3},\frac{\alpha_1(x)}{2}\Bigr\},\\
c_{\flat}:=\min\Bigl\{\frac{C_0}{2B_3},\frac{C_1}{2}\Bigr\},\qquad
\mathfrak{a}_s:=\inf_{t\ge0}L^{t}\alpha^{\flat}_{s+t}
\end{gathered}
\end{equation}
are well defined and positive, and $\alpha_u^{\flat}\ge c_{\flat}\,(x_{+}^{(u)})^{-1/2}$ for every $u$.
\item[(c)] Let $\kappa_0$ be the exponent of the irrelevance input among the standing inputs; it
satisfies $\kappa_0\ge7$. For every $s\ge0$ and every $\beta>0$,
\begin{equation}\label{4.31:eq-cutoff-free-windows-1}
\sum_{t\ge0}\bigl(x_{-}^{(s+t)}\bigr)^{1-\kappa_0/2}<\infty,\qquad
\sum_{t\ge0}\bigl(x_{+}^{(s+t)}\bigr)^{1/2}L^{-\beta t/2}<\infty .
\end{equation}
\end{enumerate}
\end{lemma}

\begin{proof}
(a) Let $K\in\mathcal{K}$ and $0\le u\le K$. The anchored flow bound of
\eqref{4.31:eq-standing-inputs-a}, applied at this $u$, gives
$u\lambda-2c_2\le x_{K-u}-x_K\le uB+2c_2$, with the constants $\lambda,B,c_2>0$ free of $K$.
Adding the window $x_{-}\le x_K\le x_{+}$ of \eqref{4.31:eq-windows-tubes-a} yields
\begin{equation*}
x_{-}+\lambda u-2c_2\le x_{K-u}\le x_{+}+Bu+2c_2 .
\end{equation*}
Moreover $x_{K-u}\ge\gamma^{-2}$, because $0<g_{K-u}\le\gamma$ along the flow of
\eqref{4.31:eq-standing-inputs-a}. Hence
\begin{equation*}
\max\{\gamma^{-2},x_{-}+\lambda u-2c_2\}\le x_{K-u}\le x_{+}+Bu+2c_2,
\end{equation*}
which is $x_{K-u}\in I_u$ by the definition of $x_{\pm}^{(u)}$.

(b) Since $\mathcal{K}$ is infinite, for each $u$ there is $K\in\mathcal{K}$ with $K\ge u$, and then
(a) shows that $I_u$ is non-empty; it is compact, and $I_u\subset[\gamma^{-2},\infty)$ since
$x_{-}^{(u)}\ge\gamma^{-2}$. By the ceiling $\gamma^2\le 1/e$ imposed in
Notation~\ref{4.31:not-complexified-group} we have $\log x\ge1$ for
$x\ge\gamma^{-2}$; hence $\varepsilon$, $\alpha_0$, $\alpha_1$ are continuous and positive on
$[\gamma^{-2},\infty)$, and $(\log x)^{q_i}\ge1$ gives $\alpha_i(x)\ge C_ix^{-1/2}$ there. A continuous
positive function attains a positive minimum on a non-empty compact set, so $\varepsilon_u^{-}>0$ and
$\alpha_u^{\flat}>0$. For $x\in I_u$ one has $x^{-1/2}\ge(x_{+}^{(u)})^{-1/2}$, hence
$\alpha_0(x)/2B_3\ge (C_0/2B_3)(x_{+}^{(u)})^{-1/2}$ and $\alpha_1(x)/2\ge(C_1/2)(x_{+}^{(u)})^{-1/2}$;
taking minima gives $\alpha_u^{\flat}\ge c_{\flat}(x_{+}^{(u)})^{-1/2}$, and $c_{\flat}>0$.
For $\mathfrak{a}_s$ this bound gives
\begin{equation*}
L^{t}\alpha^{\flat}_{s+t}\ge c_{\flat}\,L^{t}\bigl(x_{+}+B(s+t)+2c_2\bigr)^{-1/2},
\end{equation*}
and the right side tends to $\infty$ as $t\to\infty$ because $L>1$. Hence only finitely many $t$
satisfy $L^{t}\alpha^{\flat}_{s+t}\le\alpha^{\flat}_{s}$, the value at $t=0$; the infimum is
therefore a minimum over finitely many positive numbers, and $\mathfrak{a}_s>0$.

(c) For the first series:
\begin{equation*}
x_{-}^{(u)}=\max\{\gamma^{-2},x_{-}+\lambda u-2c_2\}\ge\gamma^{-2}\ge e>1,
\end{equation*}
and $1-\kappa_0/2\le-5/2$ because $\kappa_0\ge7$. Therefore
$(x_{-}^{(s+t)})^{1-\kappa_0/2}\le(x_{-}^{(s+t)})^{-5/2}\le\gamma^{5}$ for every $t$, and, as soon as
$t\ge 2(2c_2-x_{-}-\lambda s)/\lambda$, also $x_{-}^{(s+t)}\ge x_{-}+\lambda(s+t)-2c_2\ge\lambda t/2$, so
that the summand is at most $(\lambda t/2)^{-5/2}$; since $\sum_{t\ge1}t^{-5/2}<\infty$ the series
converges. For the second series, $x_{+}^{(s)}\ge x_{-}^{(s)}$ since $I_s\ne\emptyset$, and
$x_{-}^{(s)}\ge\gamma^{-2}>0$; hence
\begin{equation*}
x_{+}^{(s+t)}=x_{+}^{(s)}+Bt\le(1+t)\bigl(x_{+}^{(s)}+B\bigr),
\end{equation*}
so the series is dominated by $(x_{+}^{(s)}+B)^{1/2}\sum_{t\ge0}(1+t)^{1/2}L^{-\beta t/2}$, which is
finite for every $\beta>0$ because $L>1$.

Of the flow inputs in \eqref{4.31:eq-standing-inputs-a}, only the anchored cumulative bound has been
used; the bound for all pairs of levels and the per-step statement do not enter.
\end{proof}

\begin{definition}[Choice of the two free radii]\label{4.31:def-radii-choice}
Fix the depth $s$. Two lower bounds on $M$ are in force throughout this section and are recalled
here, not imposed; they are among the lower bounds on $M$ assumed in the main theorem of this
section:
\begin{equation*}
R_1M_1\le M,\qquad 2M_1\le M,
\end{equation*}
where $R_1$, $M_1$ are Ba{\l}aban's constants. Let $c_{\square}$ be the constant written $O(1)$
in \cite[(1.12)]{B-CMP109}. Let $M'$ be the least multiple of $R_1M_1$ not smaller than $M$, and
let $M''$ be the least multiple of $R_1M_1$ not smaller than $\max\{M',c_{\square}LM\}$ (the first
two rows of the display below). The two radii
$w_s,r_s$ are positive numbers, fixed so small that, with
\begin{equation}\label{4.31:eq-radii-choice-a}
\begin{aligned}
&M':=\min\{nR_1M_1:\ n\in\mathbb{N},\ nR_1M_1\ge M\},\\
&M'':=\min\bigl\{nR_1M_1:\ n\in\mathbb{N},\ nR_1M_1\ge\max\{M',c_{\square}LM\}\bigr\},\\
&\rho_s:=2B_3w_s+8B_0r_s,\qquad h_s:=8\bigl(B_3M'w_s+C_{19}\rho_s\bigr),\\
&C_m:=6\bigl(B_3+8C_{19}(B_3+B_0)\bigr),\\
&\text{(S1)}\quad r_s\le\tfrac1{20},\qquad
w_s+16r_s\le\min\{a_1,\ \varepsilon_s^{-}/B_3\},\qquad \rho_s\le\rho_W,\\
&\phantom{\text{(S1)}}\quad M''B_3w_s+C_{19}\rho_s\le\min\{a_{\mathrm{op}},a_J,\hat{a}_1\},
\qquad M''\le M(w_s),\\
&\phantom{\text{(S1)}}\quad c_u\bigl(B_3w_s+C_{19}\rho_s\bigr)\le\tfrac1{10},\\
&\text{(S2)}\quad 2C_m(w_s+r_s)\le\mathfrak{a}_s,\\
&\text{(S3)}\quad 32M'h_s\le1,\\
&\text{(S4)}\quad 2c_{\mathrm{av}}C_gMh_s\le\min\{1,\beta_W\}.
\end{aligned}
\end{equation}
Here $B_0$, $C_{19}$, $\rho_W$ and the side thresholds $a_{\mathrm{op}}$, $a_J$,
$\hat{a}_1$ are the constants of the complex background field
(Definition~\ref{4.31:def-complex-background}), $B_3$ is Ba{\l}aban's constant, $\beta_W$ is
the threshold on the datum in its
local form (same definition), $c_u$ is the constant of the nesting identities
(Definition~\ref{4.31:def-window-counting-nesting}), and $\varepsilon_s^{-}$, $\mathfrak{a}_s$
are the cutoff-free quantities of \eqref{4.31:eq-cutoff-free-windows-a}; $a_1$ is Ba{\l}aban's
constant and $M(\cdot)$ his function, for which the condition $M''\le M(w_s)$ is equivalent to
\begin{equation*}
w_s\le R_1M_1a_1/M''.
\end{equation*}
Further, $c_{\mathrm{av}}:=\max\{4c_dc_P,\ 2c_d/r_P,\ 8c_d\}$, where $c_d$, $c_P$, $r_P$ are
the absolute numbers of the bound on averages of a small field established later in this section,
and $C_g$ is the absolute constant in the bound, established later in this section, on a potential
from which its value at a base site has been removed by a $G^{\mathbb{C}}$-valued gauge
transformation; neither constant is given a numerical value. All conditions
\textup{(S1)}--\textup{(S4)} are upper bounds on $(w_s,r_s)$.

Such a choice exists, and it can be made independently of $K$: the set of pairs
$(w_s,r_s)$ satisfying \textup{(S1)}--\textup{(S4)} is non-empty and depends only on
$(d,L,M,\mathfrak{p},x_{\pm},s)$. With this choice,
the charts used below have radius $4r_s$ and the complex tube $\mathfrak{D}_s(w_s,r_s)$ has
radius $r_s$; the cubes of the covering class are accommodated in the torus by
\eqref{4.31:eq-window-counting-nesting-f}.
\end{definition}

\begin{proof}
We first record the elementary bounds on $M'$ and $M''$. The least multiple of $R_1M_1$ that is
$\ge M$ is smaller than $M+R_1M_1$; since $R_1M_1\le M$, this gives $M\le M'\le 2M$. Likewise
$M''$ is smaller than $\max\{M',c_{\square}LM\}+R_1M_1$, and $R_1M_1\le M\le M'$, so that
\begin{equation*}
\max\{M',c_{\square}LM\}\le M''\le 2\max\{M',c_{\square}LM\}.
\end{equation*}

Next, $\rho_s$ and $h_s$ are linear forms in $(w_s,r_s)$ with non-negative coefficients, built
from $B_0$, $B_3$, $C_{19}$ and $M'$, and $C_m$ does not depend on $(w_s,r_s)$. Hence every
condition in \textup{(S1)}--\textup{(S4)} has the following shape: a linear form in
$(w_s,r_s)$ with non-negative coefficients is bounded above by a positive number. Indeed, the
right-hand sides are $\tfrac1{20}$; $\min\{a_1,\varepsilon_s^{-}/B_3\}$, positive since
$\varepsilon_s^{-}>0$ by Lemma~\ref{4.31:lem-cutoff-free-windows}; $\rho_W$;
$\min\{a_{\mathrm{op}},a_J,\hat{a}_1\}$; $R_1M_1a_1/M''$ (for the condition $M''\le M(w_s)$, in
its equivalent form above); $\tfrac1{10}$ (divided by $c_u$, the condition being void if
$c_u=0$); $\mathfrak{a}_s/(2C_m)$, positive since $\mathfrak{a}_s>0$ by
Lemma~\ref{4.31:lem-cutoff-free-windows}; $1/(32M')$; and
$\min\{1,\beta_W\}/(2c_{\mathrm{av}}C_gM)$. A linear form with non-negative coefficients is at
most its coefficient sum times $\max\{w_s,r_s\}$, so each condition holds as soon as
$\max\{w_s,r_s\}$ lies below a positive number; since there are finitely many conditions, all of
them hold as soon as $\max\{w_s,r_s\}$ lies below the smallest of these positive numbers. The set
of pairs satisfying \textup{(S1)}--\textup{(S4)} is therefore non-empty.

Finally, none of these thresholds depends on $K$. The constants $B_0$, $C_{19}$, $\rho_W$ and
$\beta_W$ depend on $(d,L,M)$ only, by
Definition~\ref{4.31:def-complex-background}; $c_u$ is independent of $K$ by
Definition~\ref{4.31:def-window-counting-nesting}; $\varepsilon_s^{-}$ and $\mathfrak{a}_s$ are
free of $K$ by Lemma~\ref{4.31:lem-cutoff-free-windows} and depend on the window endpoints
$x_{\pm}$ and on $s$; $M'$ and $M''$ depend on $R_1$, $M_1$, $c_{\square}$, $L$ and $M$;
$B_3$, $a_1$, $R_1$, $M_1$, $c_{\square}$ are Ba{\l}aban's constants and, together with the side
thresholds $a_{\mathrm{op}}$, $a_J$, $\hat{a}_1$, are determined by $d$, $L$, $M$ and the auxiliary
parameters $\mathfrak{p}$, independently of $K$; $c_{\mathrm{av}}$ is a maximum of absolute
numbers and $C_g$ is absolute.
Every threshold is thus a positive number depending on $(d,L,M,\mathfrak{p},x_{\pm},s)$ only, and
the choice of $(w_s,r_s)$ is free of $K$.
\end{proof}

\begin{lemma}[The small-field characteristic functions equal one on the window]
\label{4.31:lem-window-characteristic}
Assume the following two inputs, with the characteristic function $\chi_k$ of the small-field
region at level $k=K-s$ read as in \cite[(2.16)--(2.17)]{B-CMP119}:
\begin{itemize}
\item the window statement \eqref{4.31:eq-window-counting-nesting-a} of
Definition~\ref{4.31:def-window-counting-nesting};
\item the real background statement \eqref{4.31:eq-standing-inputs-e}, that is,
\cite[Theorem~1, (7)--(8)]{B-CMP102b}.
\end{itemize}
If $w\le a_1$ and $B_3w\le\varepsilon_k$, then $\chi_k\equiv1$ on $\mathcal{V}_s(w)$. In particular,
for every $K$, $\chi_k\equiv1$ on $\mathcal{V}_s(w_s+16r_s)$, with $w_s,r_s$ the radii of
Definition~\ref{4.31:def-radii-choice}.
\end{lemma}

\begin{proof}
The first assertion is the content of \eqref{4.31:eq-window-counting-nesting-a}. We show how it is
read off from \eqref{4.31:eq-standing-inputs-e}. Write $\eta=L^{-k}$. By \cite[(2.17)]{B-CMP119},
$\chi_k$ is a product over the cubes $\square$ of the partition at level $k$,
\begin{equation}\label{4.31:eq-window-characteristic-1}
\chi_k(V)=\prod_{\square}\chi\Bigl(\Bigl\{\sup_{p\subset\square^{\sim}}
\bigl|U_{k,\square}(V,\partial p)-1\bigr|<\varepsilon_k\eta^2\Bigr\}\Bigr),
\end{equation}
where $\square^{\sim}$ is the enlargement of $\square$ used there, and where, by
\cite[(2.16)]{B-CMP119} together with \cite[(1.2)--(1.3)]{B-CMP119}, $U_{k,\square}(V)$ is the
minimal configuration on the minimal determining set based on $\square^{\sim4}$ whose block datum
is that of $V$.

Let $V\in\mathcal{V}_s(w)$ and put $\varepsilon_1=\sup_p|V(\partial p)-1|$, so that
$\varepsilon_1<w\le a_1$. The block datum of $V$ then satisfies hypothesis (7) of
\cite[Theorem~1]{B-CMP102b} with this $\varepsilon_1$, and conclusion (8) of the same theorem, in
the form \eqref{4.31:eq-standing-inputs-e}, gives
\begin{equation*}
\bigl|U_{k,\square}(V,\partial p)-1\bigr|<B_3\varepsilon_1\eta^2<B_3w\eta^2
\end{equation*}
for every plaquette $p$ of the top domain of the determining set. This top domain contains
$\square^{\sim}$, by the inclusion $\Omega^{\sim-2}\subset\Omega_j$ of \cite{B-CMP119}, in which
$\Omega^{\sim-2}$ denotes $\Omega$ with two layers of cubes removed. Since $B_3w\le\varepsilon_k$,
every event in \eqref{4.31:eq-window-characteristic-1} occurs, every factor equals one, and
$\chi_k(V)=1$. The same threshold $B_3^{-1}\varepsilon_j$ for the datum, obtained by an
application of \cite[Theorem~1]{B-CMP102b}, is the one used in \cite[p.~380]{B-CMP122b} and in
\cite{B-CMP109}.

For the second assertion take $w=w_s+16r_s$. By Lemma~\ref{4.31:lem-cutoff-free-windows},
applied with $u=s$, we have $x_k=x_{K-s}\in I_s$, hence
\begin{equation*}
\varepsilon_k=\varepsilon(x_k)\ge\min_{I_s}\varepsilon=\varepsilon_s^{-}
\end{equation*}
(see \eqref{4.31:eq-cutoff-free-windows-a}), and this holds for every $K$. Condition (S1) in
\eqref{4.31:eq-radii-choice-a} gives $B_3(w_s+16r_s)\le\varepsilon_s^{-}$ and
$w_s+16r_s\le a_1$. Hence $w=w_s+16r_s$ satisfies $w\le a_1$ and
$B_3w\le\varepsilon_s^{-}\le\varepsilon_k$, and the first assertion applies.

If the datum is taken to satisfy (7) of \cite[Theorem~1]{B-CMP102b} only with
$\varepsilon_1=C_Q\sup_p|V(\partial p)-1|$, where $C_Q=C_Q(d,L)$ is a constant depending on $d$
and $L$ only, the argument above applies unchanged once $B_3$ is replaced by $C_QB_3$ in
condition (S1).
\end{proof}

\begin{lemma}[Holomorphic charts on the tube and their gluing]\label{4.31:lem-charts-gluing}
Let $K>s$, with $k$ the level and $\eta=L^{-k}$ the spacing of the fine lattice
$T_{\eta}\supset\mathbb{T}_s$ of the $k$-step problem of Definition~\ref{4.31:def-complex-background};
let $w_s,r_s>0$ be the radii fixed in \eqref{4.31:eq-radii-choice-a}, with $\rho_s$ as
defined there, and let $V_0\in\mathcal{V}_s(w_s)$, the plaquette window of
Definition~\ref{4.31:def-windows-tubes}. Let $U_0:=U_k(V_0)$ be the configuration of
\eqref{4.31:eq-standing-inputs-e} with $\varepsilon_1=w_s$. For
$B\in(\mathfrak{g}^{\mathbb{C}})^{\mathrm{bonds}(\mathbb{T}_s)}$ with $|B|<4r_s$, where
$|B|=\sup_b|B(b)|$ and $|\cdot|$ on matrices is the operator norm, put
\begin{equation*}
W_{V_0}(B):=e^{i\eta\mathcal{H}(B;U_0)}\,U_0 ,
\end{equation*}
the complex background field of Definition~\ref{4.31:def-complex-background} in its global form
\eqref{4.31:eq-complex-background-a}, taken at $\alpha_0=\alpha_1=B_3w_s$, $\beta=4r_s$ and
$\rho=\rho_s\le\rho_W$. Only one level of that construction enters, because the determining set
$\mathfrak{B}_k$ is the whole lattice $T^{(k)}_1$.
\begin{enumerate}
\item[(a)] The map $B\mapsto W_{V_0}(B)$ is holomorphic on $\{|B|<4r_s\}$, and for every such $B$
and every plaquette $p$
\begin{align*}
|\mathcal{H}(B;U_0)|,\ |\nabla^{\eta}_{U_0}\mathcal{H}(B;U_0)| &\le C_{19}\rho_s,\\
|W_{V_0}(B)(\partial p)-\mathbf{1}| &\le 6(B_3w_s+C_{19}\rho_s)\,\eta^2,\\
|J^{\eta}(W_{V_0}(B))|,\ |J^{\eta}(U_0)| &\le 4(B_3w_s+C_{19}\rho_s).
\end{align*}
\item[(b)] If $B$ is $\mathfrak{g}$-valued, then $W_{V_0}(B)$ is equivalent under a $G$-valued gauge
transformation to $U_k(e^{iB}V_0)$, and
\begin{equation*}
e^{iB}V_0\in\mathcal{V}_s(w_s+16r_s).
\end{equation*}
\item[(c)] Let $f^{\mathbb{R}}$ be a function on $\mathcal{V}_s(w_s+16r_s)$ and, for each
$V_0\in\mathcal{V}_s(w_s)$, let $\Phi_{V_0}$ be holomorphic on $\{|B|<4r_s\}$ with
$\Phi_{V_0}(B)=f^{\mathbb{R}}(e^{iB}V_0)$ for every $\mathfrak{g}$-valued $B$ with $|B|<4r_s$. Put
\begin{equation*}
\mathfrak{T}_s:=\bigcup_{V_0\in\mathcal{V}_s(w_s)}\{e^{iB}V_0:\ |B|<r_s\},
\qquad F(e^{iB}V_0):=\Phi_{V_0}(B)\quad(|B|<r_s).
\end{equation*}
Then $\mathfrak{T}_s$ is open, $\mathfrak{T}_s\supset\mathfrak{D}_s(w_s,r_s)$, $F$ is well defined
and holomorphic on $\mathfrak{T}_s$, and $F=f^{\mathbb{R}}$ on $\mathcal{V}_s(w_s)$.
\end{enumerate}
\end{lemma}

\begin{proof}
\emph{Part (a) and the gauge equivalence in (b).} We apply the global form
\eqref{4.31:eq-complex-background-a} of the complex background field to the $k$-step problem with
top-level datum $e^{iB}V_0$, at $\alpha_0=\alpha_1=B_3w_s$ and $\beta=4r_s$. With these values
\begin{equation*}
\alpha_0+\alpha_1+2B_0\beta=2B_3w_s+8B_0r_s=\rho_s ,
\end{equation*}
and the condition $\rho_s\le\rho_W$, together with the side thresholds of the complex background
field, is part of (S1) in \eqref{4.31:eq-radii-choice-a}. The global form then provides the
configuration $W_{V_0}(B)=e^{i\eta\mathcal{H}(B;U_0)}U_0$, holomorphic in $B$, with the bound
$C_{19}\rho_s$ on $\mathcal{H}$ and on its covariant derivative $\nabla^{\eta}_{U_0}\mathcal{H}$, the
bound $6(B_3w_s+C_{19}\rho_s)\eta^2$ on $|W_{V_0}(B)(\partial p)-\mathbf{1}|$, and the bound
$4(B_3w_s+C_{19}\rho_s)$ on the currents $J^{\eta}(W_{V_0}(B))$ and $J^{\eta}(U_0)$; this is (a).
For $\mathfrak{g}$-valued $B$ the same statement gives that $W_{V_0}(B)$ is equivalent under a
$G$-valued gauge transformation to $U_k(e^{iB}V_0)$, which is the first assertion of (b).

\emph{The window in (b).} Let $B$ be $\mathfrak{g}$-valued with $|B|<4r_s$. For each bond $b$ the
matrix $B(b)$ is Hermitian and traceless, so $e^{iB(b)}$ is unitary, and it has determinant one
because $B(b)$ is traceless; hence $e^{iB}V_0\in\mathcal{G}_s$. The eigenvalues of $B(b)$ are real
numbers $\theta$, and $|e^{i\theta}-1|\le|\theta|$, so
\begin{equation*}
|e^{\pm iB(b)}-\mathbf{1}|\le|B(b)|<4r_s .
\end{equation*}
Fix a plaquette $p$. The variable $V_0(\partial p)$ is a product of four factors, each of the form
$V_0(b)$ or $V_0(b)^{-1}$ for a bond $b$ of $\partial p$. In $(e^{iB}V_0)(\partial p)$ each factor
$V_0(b)$ is replaced by $e^{iB(b)}V_0(b)$ and each factor $V_0(b)^{-1}$ by $V_0(b)^{-1}e^{-iB(b)}$;
the distance between the old and the new factor is $|e^{iB(b)}-\mathbf{1}|$, respectively
$|e^{-iB(b)}-\mathbf{1}|$, since the norm is invariant under multiplication by unitary matrices on
either side. For the same reason, if one factor of a product of unitary matrices is replaced by
another unitary matrix, the product moves by exactly the distance between the two factors. Replacing
the four factors one at a time, each replacement moves the product by less than $4r_s$, and
therefore, the sum running over the four bonds of $\partial p$ with the sign of the corresponding
factor,
\begin{align*}
|(e^{iB}V_0)(\partial p)-\mathbf{1}|
&\le|V_0(\partial p)-\mathbf{1}|+\sum_{b\in\partial p}|e^{\pm iB(b)}-\mathbf{1}|\\
&<w_s+4\cdot 4r_s=w_s+16r_s .
\end{align*}
As $p$ was arbitrary, $e^{iB}V_0\in\mathcal{V}_s(w_s+16r_s)$. In particular the hypothesis of (c)
relating $\Phi_{V_0}$ and $f^{\mathbb{R}}$ refers only to points where $f^{\mathbb{R}}$ is defined.

\emph{Part (c): well-definedness.} Let $V_0,V_0'\in\mathcal{V}_s(w_s)$ and
$B,B'\in(\mathfrak{g}^{\mathbb{C}})^{\mathrm{bonds}(\mathbb{T}_s)}$ with $|B|,|B'|<r_s$ and
$e^{iB}V_0=e^{iB'}V_0'$. On every bond $V_0'V_0^{-1}=e^{-iB'}e^{iB}$. The left-hand side is
unitary, as a product of elements of $SU(N)$. For the right-hand side we use
$|e^{X}-\mathbf{1}|\le|X|e^{|X|}$ and $|e^{X}|\le e^{|X|}$; this gives, on every bond,
\begin{align*}
|V_0'V_0^{-1}-\mathbf{1}|
&=|e^{-iB'}(e^{iB}-\mathbf{1})+(e^{-iB'}-\mathbf{1})|\\
&\le r_se^{2r_s}+r_se^{r_s}=r_se^{r_s}(e^{r_s}+1)\le 2.2\,r_s ,
\end{align*}
since $r_s\le 1/20$ by condition (S1) of \eqref{4.31:eq-radii-choice-a} and
$e^{0.05}(e^{0.05}+1)=2.157\ldots\le 2.2$. A unitary matrix $u$ with
$|u-\mathbf{1}|\le 2.2r_s$ has eigenvalues $e^{i\theta}$ with $2|\sin(\theta/2)|\le 2.2r_s$; its
principal logarithm $-i\log u$ is Hermitian with eigenvalues $\theta$,
$|\theta|\le 2\arcsin(1.1r_s)$. We put $C:=-i\log(V_0'V_0^{-1})$ bondwise; $C$ is Hermitian,
traceless because $\det(V_0'V_0^{-1})=1$, so $C\in\mathfrak{g}^{\mathrm{bonds}(\mathbb{T}_s)}$,
$V_0'=e^{iC}V_0$, and
\begin{equation*}
|C|\le 2\arcsin(1.1\,r_s)<2.4\,r_s ,
\end{equation*}
because $y\mapsto\arcsin(y)/y$ is increasing on $(0,1]$ and $2\arcsin(0.055)/0.05\le 2.4$.

On the convex ball $\{|B''|<1.2r_s\}$ in $(\mathfrak{g}^{\mathbb{C}})^{\mathrm{bonds}(\mathbb{T}_s)}$
define, bondwise, $b(B''):=-i\log(e^{iB''}e^{iC})$. By the Baker--Campbell--Hausdorff series,
$\log(e^{iB''}e^{iC})$ is $iB''+iC$ plus nested commutators, hence traceless, and
\begin{equation*}
|b(B'')|\le|B''|+|C|+2|B''||C|<3.6\,r_s+5.76\,r_s^2\le 3.888\,r_s<4r_s ,
\end{equation*}
using $r_s\le 1/20$. Thus $b$ is holomorphic on the ball with values in $\{|B|<4r_s\}$. Put
\begin{equation*}
\Psi(B''):=\Phi_{V_0}(b(B''))-\Phi_{V_0'}(B''),\qquad |B''|<1.2\,r_s ;
\end{equation*}
both terms are defined, since $|b(B'')|<4r_s$ and $|B''|<1.2r_s<4r_s$, and $\Psi$ is holomorphic.
If $B''$ is $\mathfrak{g}$-valued, then $e^{iB''}e^{iC}$ is unitary of determinant one and near
$\mathbf{1}$, so $b(B'')$ is $\mathfrak{g}$-valued, and
$e^{ib(B'')}V_0=e^{iB''}e^{iC}V_0=e^{iB''}V_0'$; by the hypothesis of (c) both terms of $\Psi(B'')$
equal $f^{\mathbb{R}}(e^{iB''}V_0')$, so $\Psi(B'')=0$. The real points of the ball form a maximal
totally real slice of the convex ball, and the identity principle
\eqref{4.31:eq-complex-background-c} gives $\Psi\equiv0$.

Now take $B''=B'$, which lies in the ball since $|B'|<r_s$. Then
\begin{equation*}
e^{ib(B')}=e^{iB'}e^{iC}=e^{iB'}V_0'V_0^{-1}=e^{iB'}e^{-iB'}e^{iB}=e^{iB}.
\end{equation*}
Both $b(B')$ and $B$ have norm below $4r_s\le1/5$, and the exponential map is injective near $0$;
hence $b(B')=B$. From $\Psi(B')=0$ we obtain $\Phi_{V_0}(B)=\Phi_{V_0'}(B')$, so $F$ is well defined.

\emph{Holomorphy and openness.} Fix $V_0\in\mathcal{V}_s(w_s)$. Near $0$ the exponential map is a
biholomorphism onto a neighbourhood of $\mathbf{1}$ with inverse the principal logarithm, so the set
$\{e^{iB}V_0:|B|<r_s\}$ is open in $\mathcal{G}_s^{\mathbb{C}}$, and $\mathfrak{T}_s$, a union of such
sets, is open. On that set, by well-definedness with the same $V_0$,
\begin{equation*}
F(V)=\Phi_{V_0}\bigl(-i\log(VV_0^{-1})\bigr),
\end{equation*}
a composition of holomorphic maps. Hence $F$ is holomorphic near every point of $\mathfrak{T}_s$.

\emph{The tube and the real slice.} If $V=e^{-Y}V_0\in\mathfrak{D}_s(w_s,r_s)$ with
$V_0\in\mathcal{V}_s(w_s)$, $Y\in\mathfrak{g}^{\mathrm{bonds}(\mathbb{T}_s)}$, $|Y|<r_s$, take $B=iY$:
then $e^{iB}=e^{-Y}$ and $|B|=|Y|<r_s$, so $V\in\mathfrak{T}_s$. Finally, for
$V\in\mathcal{V}_s(w_s)$ take $V_0=V$ and $B=0$; then $F(V)=\Phi_V(0)=f^{\mathbb{R}}(V)$.
\end{proof}

\begin{lemma}[Membership of the complex background field in the analyticity space]
\label{4.31:lem-background-membership}
Assume the following statements of this section:
\begin{itemize}
\item the form, analyticity, locality, gauge invariance and bounds of Ba{\l}aban's local
effective-action pieces, \eqref{4.31:eq-standing-inputs-b}, and the real background field,
\eqref{4.31:eq-standing-inputs-e} (Definition~\ref{4.31:def-standing-inputs});
\item the complex background field in its global form, \eqref{4.31:eq-complex-background-a}
(Definition~\ref{4.31:def-complex-background});
\item the nesting identity for small complex minimisers in its global form,
\eqref{4.31:eq-window-counting-nesting-c} (Definition~\ref{4.31:def-window-counting-nesting}).
\end{itemize}
Let $s\ge 0$, $t\ge 0$ and $K$ be integers with $K\ge s+t+1$, and let $w_s,r_s$ be the radii of
Definition~\ref{4.31:def-radii-choice}. Let $k=K-s$ and $\eta=L^{-k}$ be the level and
the fine spacing of the torus $\mathbb{T}_s$ as in Lemma~\ref{4.31:lem-charts-gluing}, and put
$j=k-t$ and $\xi=L^{-j}$. Let $X\in\mathbf{D}_j$, $V_0\in\mathcal{V}_s(w_s)$, let
$B\in(\mathfrak{g}^{\mathbb{C}})^{\mathrm{bonds}(\mathbb{T}_s)}$ with $|B|<4r_s$, and let
$W=W_{V_0}(B)$ be the complex background field of Lemma~\ref{4.31:lem-charts-gluing}. Then
\begin{equation}\label{4.31:eq-background-membership-1}
\bigl(W,J^{\xi}(W)\bigr)\big|_X\in\mathfrak{U}^{c}_j(X,\alpha_{0,j},\alpha_{1,j}).
\end{equation}
Consequently the maps
\begin{gather*}
B\longmapsto\mathbf{E}^{(j)}\bigl(X,(W_{V_0}(B),J^{\xi}(W_{V_0}(B))),z\bigr),\\
B\longmapsto\mathbf{R}^{(j)}\bigl(X,(W_{V_0}(B),J^{\xi}(W_{V_0}(B)))\bigr)
\end{gather*}
are holomorphic on $\{|B|<4r_s\}$, obey the bounds of \cite[(1.18)]{B-CMP109} and
\cite[(2.31)]{B-CMP119} respectively, and for $B$ with values in $\mathfrak{g}$ they are equal to
Ba{\l}aban's $\mathbf{E}^{(j)}(X,U_k(e^{iB}V_0),z)$ and $\mathbf{R}^{(j)}(X,U_k(e^{iB}V_0))$.
\end{lemma}

\begin{proof}
\emph{Step 1: units.} Throughout, one and the same fine configuration on $\mathbb{T}_s$ is read
at two levels: at level $k$, in the units of $\mathbb{T}_s$, where the fine spacing is $\eta$,
and at scale $j$, where the fine spacing is $\xi=L^{t}\eta$. The regularity of the configurations
we use is known at level $k$, by Lemma~\ref{4.31:lem-charts-gluing}(a) and by the real background
field \eqref{4.31:eq-standing-inputs-e}, that is, \cite[Theorem~1, (9)]{B-CMP102b}. If a bond
variable is written as $e^{i\eta a}=e^{i\xi a'}$, then the potential, which is homogeneous of
degree one, its difference quotient (degree two), a current (degree three) and the plaquette
threshold (degree two) transform as
\begin{equation}\label{4.31:eq-background-membership-2}
a'=L^{-t}a,\qquad \nabla^{\xi}a'=L^{-2t}\nabla^{\eta}a,\qquad J^{\xi}=L^{-3t}J^{\eta},
\qquad \eta^2=L^{-2t}\xi^2 ;
\end{equation}
compare \cite[(3.12)]{B-CMP109}. The same rescaling holds for covariant difference quotients.
A length measured in scale-$j$ units is multiplied by $L^{-t}\le1$ when it is measured in the
units of $\mathbb{T}_s$.

\emph{Step 2: constants.} Put
\begin{equation*}
\alpha_0':=\frac{\alpha_{0,j}}{2B_3},\qquad \alpha_1':=\frac{\alpha_{1,j}}{2}.
\end{equation*}
Since $j=K-(s+t)$, Lemma~\ref{4.31:lem-cutoff-free-windows} gives $x_j\in I_{s+t}$; since
$\alpha_{i,j}=\alpha_i(x_j)$ for $i=0,1$, and $\alpha^{\flat}_{s+t}$ is the minimum of
$\min\{\alpha_0(x)/2B_3,\alpha_1(x)/2\}$ over $x\in I_{s+t}$, it follows that
$\min\{\alpha_0',\alpha_1'\}\ge\alpha^{\flat}_{s+t}$. By the definition of $\mathfrak{a}_s$ as an
infimum over $t\ge0$ we have $L^{t}\alpha^{\flat}_{s+t}\ge\mathfrak{a}_s$, and the second of the
conditions \eqref{4.31:eq-radii-choice-a} reads $\mathfrak{a}_s\ge 2C_m(w_s+r_s)$. Therefore
\begin{equation}\label{4.31:eq-background-membership-3}
\min\{\alpha_0',\alpha_1'\}\ge\alpha^{\flat}_{s+t}\ge L^{-t}\mathfrak{a}_s
\ge 2C_m(w_s+r_s)L^{-t}.
\end{equation}
Moreover $\rho_s=2B_3w_s+8B_0r_s\le 8(B_3+B_0)(w_s+r_s)$, and therefore, since
$C_m=6B_3+48C_{19}(B_3+B_0)$,
\begin{equation}\label{4.31:eq-background-membership-4}
6(B_3w_s+C_{19}\rho_s)\le 6B_3(w_s+r_s)+48C_{19}(B_3+B_0)(w_s+r_s)=C_m(w_s+r_s).
\end{equation}

\emph{Step 3: conditions (i)--(iii) of \eqref{4.31:eq-analyticity-space-a}.} We verify them with
the constants $(\alpha_0',\alpha_1')$ for $\mathbf{U}=U'U$, where
\begin{equation*}
U:=U_0=U_k(V_0),\qquad U':=e^{i\eta\mathcal{H}}=e^{i\xi A'},\qquad A':=L^{-t}\mathcal{H},
\end{equation*}
with $\mathcal{H}=\mathcal{H}(B;U_0)$, so that $\mathbf{U}=W$ by the definition of $W_{V_0}(B)$ in
Lemma~\ref{4.31:lem-charts-gluing}; the current of the pair is $\mathbf{J}=J^{\xi}(W)$. All bounds
below hold on the whole of $\mathbb{T}_s$, hence on $X$.

Condition (i), first part, \cite[(1.11)]{B-CMP109}. By the definition of $\mathcal{V}_s(w_s)$ we
have $|\partial V_0-1|<w_s$, and $w_s\le a_1$ by the first of the conditions
\eqref{4.31:eq-radii-choice-a}; so the real background field \eqref{4.31:eq-standing-inputs-e}
applies with $\varepsilon_1=w_s$ and gives $|\partial U_0-1|<B_3w_s\eta^2$. By
\eqref{4.31:eq-background-membership-2} and \eqref{4.31:eq-background-membership-3}, and since
$C_m\ge 6B_3$ and $L^{-2t}\le L^{-t}$,
\begin{equation*}
|\partial U_0-1|<B_3w_s\eta^2=B_3w_sL^{-2t}\xi^2<2C_m(w_s+r_s)L^{-t}\xi^2\le\alpha_0'\xi^2 .
\end{equation*}

Condition (i), second part, \cite[(1.12)]{B-CMP109}. Let $B^{109}\ge1$ be the constant written
$B$ in \cite[(1.12)]{B-CMP109}, as in Definition~\ref{4.4:def-local-data-classes}, where the
number $c_{\square}$ is written $\mathsf{o}$. Let $\square$ be a scale-$j$ cube of side
$c_{\square}LM$. By Step~1 its
side in the units of $\mathbb{T}_s$ is at most $c_{\square}LM$. Slide the corner of $\square$ onto
the grid of mesh $R_1M_1$: the cube of side $2M''$ whose corner is the grid point at distance less
than $R_1M_1$ below that corner in each coordinate contains $\square$, because
$c_{\square}LM+R_1M_1\le 2M''$; indeed $M''\ge c_{\square}LM$ and
$M''\ge M'\ge M\ge R_1M_1$ by the standing floor $R_1M_1\le M$ of
Definition~\ref{4.31:def-radii-choice}. This cube is a class cube of the real background field
\eqref{4.31:eq-standing-inputs-e} with half-side $M''$: $M''$ is a multiple of $R_1M_1$, and
$M''\le M(w_s)$ by the first of the conditions \eqref{4.31:eq-radii-choice-a}; it fits into the
torus by the floor \eqref{4.31:eq-window-counting-nesting-f} on the torus exponent. Hence there is
a $G$-valued $u_0$ with $U_0^{u_0}=e^{i\eta A_0}$ on it and $|A_0|,|\nabla^{\eta}A_0|<B_3M''w_s$.
With $A:=L^{-t}A_0$ we have $U_0^{u_0}=e^{i\xi A}$, and by \eqref{4.31:eq-background-membership-2}
\begin{equation*}
|A|\le L^{-t}B_3M''w_s,\qquad |\nabla^{\xi}A|\le L^{-2t}B_3M''w_s .
\end{equation*}
Now
\begin{equation*}
M''\le 2\max\{M',c_{\square}LM\}\le 4c_{\square}LM\,(M'/M)\le 8c_{\square}LM,
\end{equation*}
and, since $C_m\ge 6B_3+48C_{19}B_3$, we have $16B_3\le 2C_m$, so that by
\eqref{4.31:eq-background-membership-3}
\begin{equation*}
16B_3w_sL^{-t}\le 2C_m(w_s+r_s)L^{-t}\le\alpha_0'.
\end{equation*}
Hence, using $L^{-2t}\le L^{-t}$ and $B^{109}\ge1$, both $|A|$ and $|\nabla^{\xi}A|$ are bounded by
\begin{equation*}
8c_{\square}LM\,B_3w_sL^{-t}\le\tfrac12 c_{\square}LM\,\alpha_0'<c_{\square}LM\,B^{109}\alpha_0'.
\end{equation*}

Condition (ii), \cite[(1.13)]{B-CMP109}. Lemma~\ref{4.31:lem-charts-gluing}(a) gives
$|\mathcal{H}|,|\nabla^{\eta}_{U_0}\mathcal{H}|\le C_{19}\rho_s$, so by
\eqref{4.31:eq-background-membership-2}
$|A'|\le L^{-t}C_{19}\rho_s$ and $|\nabla^{\xi}_{U}A'|\le L^{-2t}C_{19}\rho_s$. By
\eqref{4.31:eq-background-membership-4}, $6C_{19}\rho_s\le C_m(w_s+r_s)$, and with
\eqref{4.31:eq-background-membership-3} both quantities are at most
$(C_m/6)(w_s+r_s)L^{-t}<\alpha_1'$.

Condition (iii), \cite[(1.14)]{B-CMP109}. By Lemma~\ref{4.31:lem-charts-gluing}(a),
\eqref{4.31:eq-background-membership-2}, \eqref{4.31:eq-background-membership-4} and
\eqref{4.31:eq-background-membership-3},
\begin{equation}\label{4.31:eq-background-membership-5}
|\partial\mathbf{U}-1|\le 6(B_3w_s+C_{19}\rho_s)L^{-2t}\xi^2\le\tfrac12\alpha_0'L^{-t}\xi^2
<\alpha_0'\xi^2 .
\end{equation}
By Lemma~\ref{4.31:lem-charts-gluing}(a) and Step~1, and by
\eqref{4.31:eq-background-membership-4} and \eqref{4.31:eq-background-membership-3},
\begin{equation*}
|J^{\xi}(\mathbf{U})|,\ |J^{\xi}(U)|\le 4(B_3w_s+C_{19}\rho_s)L^{-3t}
\le\tfrac13\alpha_0'L^{-2t}<\alpha_0' .
\end{equation*}
In particular $|\mathbf{J}|=|J^{\xi}(W)|<\alpha_0'$.

Since $B_3\ge1$ we have $\alpha_0'\le\alpha_{0,j}$, and $\alpha_1'\le\alpha_{1,j}$. All conditions
are strict inequalities with constants on the right, so they remain true when the constants are
enlarged; hence (i)--(iii) hold with the constants $(\alpha_{0,j},\alpha_{1,j})$.

\emph{Step 4: condition (iv) of \eqref{4.31:eq-analyticity-space-a}.} Write $X^{\sim -2}$ for the
set obtained from $X$ by removing two layers of cubes. Condition (iv),
\cite[(1.15)--(1.16)]{B-CMP109}, with $\alpha_0=\alpha_{0,j}$ asks that for every $n$ in the range
of \cite[(1.16)]{B-CMP109}, which is contained in $\{0,\dots,j\}$, the local minimiser
$U_n(X,M^{n}W)$ and its current $J_n(X,M^{n}W)$ satisfy
$|\partial U_n(X,M^{n}W)-1|<\alpha_{0,j}\xi^2$ and $|J_n(X,M^{n}W)|<\alpha_{0,j}(L^n\xi)^2$ on
$X^{\sim -2}$. We derive it from (iii) through the nesting identity
\eqref{4.31:eq-window-counting-nesting-c}, which applies to $W=W_{V_0}(B)$ with
$V_0\in\mathcal{V}_s(w_s)$, $|B|<4r_s$, to every $X\in\mathbf{D}_j$ and to every $n$ in the range:
on $X^{\sim -2}$,
\begin{equation*}
U_n(X,M^{n}W)=W^{\bar{u}^{-1}},\qquad J_n(X,M^{n}W)=(L^n\xi)^3\operatorname{Ad}(\bar{u}^{-1})
J^{\xi}(W),
\end{equation*}
where $\bar{u}$ is a block-constant $G^{\mathbb{C}}$-valued lattice gauge transformation with
$|\bar{u}^{\pm1}|\le 1+c_u(B_3w_s+C_{19}\rho_s)$. The first of the conditions
\eqref{4.31:eq-radii-choice-a} gives $c_u(B_3w_s+C_{19}\rho_s)\le1/10$, hence
$|\bar{u}^{\pm1}|\le 1.1$. A plaquette variable of a gauge transform is conjugate, by the value of
the gauge transformation at the base point of the plaquette, to the plaquette variable of the
original configuration: for a gauge transformation $h$ and a plaquette $p$ based at $x$ one has
$U^{h}(\partial p)=h(x)U(\partial p)h(x)^{-1}$, since in the ordered product of the bond
variables around $p$ the values of $h$ at the other three corners cancel. So by
\eqref{4.31:eq-background-membership-5}
\begin{equation*}
|\partial U_n(X,M^{n}W)-1|\le|\bar{u}^{-1}|\,|\bar{u}|\,|\partial W-1|
\le 1.21\cdot\tfrac12\alpha_0'\xi^2<\alpha_0'\xi^2\le\alpha_{0,j}\xi^2 .
\end{equation*}
For the current, $|\operatorname{Ad}(\bar{u}^{-1})Y|\le|\bar{u}^{-1}|\,|\bar{u}|\,|Y|\le 1.21|Y|$
for every $Y\in\mathfrak{sl}(N,\mathbb{C})$,
and the bound on $J^{\xi}(W)$ of Step~3 gives
\begin{align*}
|J_n(X,M^{n}W)|&\le(L^n\xi)^3\cdot 1.21\cdot 4(B_3w_s+C_{19}\rho_s)L^{-3t}
<(L^n\xi)^2\cdot 1.21\,\alpha_0'\\
&\le 2B_3\alpha_0'(L^n\xi)^2=\alpha_{0,j}(L^n\xi)^2 .
\end{align*}
Here we used $(L^n\xi)^3\le(L^n\xi)^2$, which holds because $L^n\xi=L^{n-j}\le1$ for $n\le j$
(the current is homogeneous of degree three, the threshold of (iv) of degree two), and
$B_3\ge1$. The bound is uniform in $n$: the scales $n$ and $j$ enter only through
$L^n\xi\le1$. Thus (iv) holds with $\alpha_0=\alpha_{0,j}$, and
\eqref{4.31:eq-background-membership-1} follows from the definition
\eqref{4.31:eq-analyticity-space-a} of $\mathfrak{U}^{c}_j(X,\alpha_{0,j},\alpha_{1,j})$.

\emph{Step 5: holomorphy and bounds.} The map $B\mapsto W_{V_0}(B)$ is holomorphic on
$\{|B|<4r_s\}$ by Lemma~\ref{4.31:lem-charts-gluing}(a); the map $U\mapsto J^{\xi}(U)$ is rational
in the entries of $U$; and by \eqref{4.31:eq-standing-inputs-b} the functionals
$\mathbf{E}^{(j)}(X,\cdot,z)$ and $\mathbf{R}^{(j)}(X,\cdot)$ are analytic on
$\mathfrak{U}^{c}_j(X,\alpha_{0,j},\alpha_{1,j})$ and depend only on the restriction of the pair to
$X$. By \eqref{4.31:eq-background-membership-1} the composition is defined for every $B$ with
$|B|<4r_s$, hence holomorphic there, and the bounds of \cite[(1.18)]{B-CMP109} and
\cite[(2.31)]{B-CMP119}, which \eqref{4.31:eq-standing-inputs-b} states on that space, hold at
every such $B$.

\emph{Step 6: real $B$.} Let $B$ take values in $\mathfrak{g}$. By
Lemma~\ref{4.31:lem-charts-gluing}(b) there is a $G$-valued gauge transformation $g$ with
$W=U_k(e^{iB}V_0)^{g}$. By the equivariance $J^{\xi}(U^{g})=\operatorname{Ad}(g)J^{\xi}(U)$ stated
in \eqref{4.31:eq-standing-inputs-b},
\begin{equation*}
\bigl(W,J^{\xi}(W)\bigr)=\bigl(U_k(e^{iB}V_0),J^{\xi}(U_k(e^{iB}V_0))\bigr)^{g}.
\end{equation*}
By the gauge invariance of $\mathbf{E}^{(j)}$ and $\mathbf{R}^{(j)}$,
\cite[(1.19)]{B-CMP109}, as stated in \eqref{4.31:eq-standing-inputs-b}, the values at
$(W,J^{\xi}(W))$ equal those at $(U_k(e^{iB}V_0),J^{\xi}(U_k(e^{iB}V_0)))$; and by the reading of
$\mathbf{E}^{(j)}(X,U,z)$ as $\mathbf{E}^{(j)}(X,(U,J^{\xi}(U)),z)$ fixed in
\eqref{4.31:eq-standing-inputs-b}, and likewise for $\mathbf{R}^{(j)}$, these are Ba{\l}aban's
$\mathbf{E}^{(j)}(X,U_k(e^{iB}V_0),z)$ and $\mathbf{R}^{(j)}(X,U_k(e^{iB}V_0))$.
\end{proof}

\begin{remark}[Condition (iv) is not cutoff-uniform for non-minimiser data]
\label{4.31:rem-non-minimiser-data}
Let $W$ be the complex background field of \eqref{4.31:eq-complex-background-a}, at depth $s$ and
level $k=K-s$, with fine spacing $\eta=L^{-k}$. Let $t\ge 0$ with $K\ge s+t+1$, let $j=k-t$ and
$\xi=L^{-j}$, and let $X\in\mathbf{D}_j$. Let $n\ge 1$ lie in the range of condition (iv) of
Definition~\ref{4.31:def-analyticity-space}, and let $U_n(X,M^{n}W)$ be the $n$-step local
minimiser of that condition (see \eqref{4.31:eq-analyticity-space-a}).
When the complex datum is measured by the supremum of its Lie-algebra-valued exponent, as in
\eqref{4.31:eq-complex-background-a}, the contribution of the complex datum to the bound on
$|\partial U_n(X,M^{n}W)-1|$ is only linear in the size of the scale-$n$ complex
datum. This bound does not give condition (iv) uniformly along the set of cutoffs
$\mathcal{K}$, already at $n=1$. Condition (iv) for $W$ follows from the nesting identity
\eqref{4.31:eq-window-counting-nesting-c}, which compares $U_n$ with $W$ itself.
\end{remark}

\begin{proof}
Write $\rho_s$ for the radius combination of \eqref{4.31:eq-complex-background-a} at depth $s$.
Write $\beta_n$ (in this remark only) for the supremum norm of the scale-$n$ complex datum
$M^{n}W$, measured as the datum is measured in \eqref{4.31:eq-complex-background-a}.

A bond of the $n$-th lattice is a product of $L^n$ fine bonds. Each fine bond carries the complex
increment $\eta\mathcal{H}$. By \eqref{4.31:eq-complex-background-a},
$|\mathcal{H}|\le C_{19}\rho_s$, so each increment has size at most $\eta C_{19}\rho_s$. In
general this accumulation cannot be removed by a gauge transformation. This happens, for example,
when $\mathcal{H}$ is of the order of a constant field $Y$ whose components do not commute. Hence
$\beta_n$ is of the order $L^n\eta C_{19}\rho_s$; write $\beta_n = c\,L^n\eta C_{19}\rho_s$,
with $c$ a constant of order one.

Read on the $n$-th lattice, the linear bound of \eqref{4.31:eq-complex-background-a} bounds the
contribution of the complex datum to $|\partial U_n(X,M^{n}W)-1|$ by at most a quantity of order
\begin{equation*}
12\,C_{19}B_0\,\beta_n L^{-2n} = 12\,c\,C_{19}^2B_0\,\rho_s\,L^{-n-k},
\end{equation*}
since $\eta=L^{-k}$. Condition (iv) of Definition~\ref{4.31:def-analyticity-space} asks for
\begin{equation*}
|\partial U_n(X,M^{n}W)-1| < \alpha_{0,j}\xi^2 = \alpha_{0,j}L^{2t-2k}.
\end{equation*}
The ratio of the first quantity to the second is
\begin{equation}\label{4.31:eq-non-minimiser-data-1}
\frac{12\,c\,C_{19}^2B_0\,\rho_s\,L^{-n-k}}{\alpha_{0,j}L^{2t-2k}}
= \frac{12\,c\,C_{19}^2B_0\,\rho_s\,L^{K-s-2t-n}}{\alpha_{0,j}},
\end{equation}
where $k=K-s$ has been used. For fixed $s$ and $t$, and $n=1$, the ratio
\eqref{4.31:eq-non-minimiser-data-1} is unbounded along $\mathcal{K}$.

Consequently, with the datum controlled only in this supremum norm, condition (iv) is not
obtained here from the linear bound. The remark following \cite[(1.16)]{B-CMP109} states that,
for the radius $\alpha_0'$ of conditions (i)--(iii) in \cite[(1.11)--(1.14)]{B-CMP109}
sufficiently small, conditions (i)--(iii) imply condition (iv). When (i)--(iii) for the datum
$M^{n}W$ are available only through the bound linear in $\beta_n$, the smallness so required is
smallness of the ratio \eqref{4.31:eq-non-minimiser-data-1}, which is not uniform along
$\mathcal{K}$. The same holds for a derivation of condition (iv) from
\cite[Proposition~9]{B-CMP102b} in which the datum enters only through $\beta_n$.

The nesting identity \eqref{4.31:eq-window-counting-nesting-c} compares $U_n(X,M^{n}W)$ with $W$
itself, not with a configuration controlled only by $\beta_n$. This is the route taken in the proof
of Lemma~\ref{4.31:lem-background-membership}.
\end{proof}

\begin{lemma}[Terrace averages of a small complex field]\label{4.31:lem-terrace-averages}
Assume that at every level $i\ge0$ the averages $M^{i}$ are given in the closed form
\eqref{4.31:eq-analyticity-space-c} of Definition~\ref{4.31:def-analyticity-space}, with the
absolute numbers $c_d$, $c_P$, $r_P$ fixed there, and put
\begin{equation*}
  c_{\mathrm{av}} := \max\Bigl\{4c_dc_P,\ \frac{2c_d}{r_P},\ 8c_d\Bigr\}.
\end{equation*}
Let $Y$ be a box of scale $j$, in whose units the fine lattice has spacing $\xi=L^{-j}$. Let
$\mathbf{A}$ be a $\mathfrak{g}^{\mathbb{C}}$-valued function on the fine bonds of $Y$. Put
$N_0:=\sup|\mathbf{A}|$, assume $c_{\mathrm{av}}N_0\le 1$, and let $U:=e^{i\xi\mathbf{A}}$,
bond by bond. Then for every level $i\le j$ and every level-$i$ bond $c$ whose double $i$-block
lies in $Y$ there is $D(c)\in\mathfrak{g}^{\mathbb{C}}$ with
\begin{equation}\label{4.31:eq-terrace-averages-1}
  M^{i}(U)(c)=e^{iD(c)},\qquad |D(c)|\le c_{\mathrm{av}}L^{i}\xi N_0\le c_{\mathrm{av}}N_0 .
\end{equation}
The datum $D$ depends holomorphically on $\mathbf{A}$, and if $\mathbf{A}$ is
$\mathfrak{g}$-valued, then $D$ is $\mathfrak{g}$-valued.
\end{lemma}

\begin{proof}
Throughout, $|\cdot|$ is the operator norm on $N\times N$ complex matrices, which is
submultiplicative and satisfies $|\mathbf{1}|=1$. By the definition of $c_{\mathrm{av}}$ and the
hypothesis $c_{\mathrm{av}}N_0\le1$ we have
\begin{equation*}
  4c_dc_PN_0\le 1,\qquad 2c_dN_0\le r_P,\qquad c_dN_0\le \tfrac18 .
\end{equation*}
Moreover $\xi=L^{-j}\le 1$, and a level-$i$ bond with $i\le j$ has length $L^{i}\xi\le 1$ in the
units of scale $j$. Every contour of \eqref{4.31:eq-analyticity-space-c} along $c$ joins the
centres of two adjacent $i$-blocks, which are $L^{i}$ fine spacings apart, so it has at least
$L^{i}$ fine bonds; since it has at most $c_dL^{i}$ of them, $c_d\ge 1$. Hence $N_0\le c_dN_0$,
and so $N_0\le\frac18$.

\emph{One fine factor.} For every fine bond $b$ of $Y$ the power series of the exponential gives
$|e^{i\xi\mathbf{A}(b)}-\mathbf{1}|\le \xi N_0e^{\xi N_0}$. Since $\xi N_0\le N_0\le\frac18\le\frac14$
and $e^{1/4}\le 1.3$, this yields
\begin{equation*}
  |e^{i\xi\mathbf{A}(b)}-\mathbf{1}|\le 1.3\,\xi N_0=:\delta .
\end{equation*}

\emph{A product of fine factors.} Let $\Pi=u_1u_2\cdots u_{n_\Pi}$ be a product of
$n_\Pi\le c_dL^{i}$ such factors, taken in any order; the factors need neither commute nor be
unitary. Writing $u_l=\mathbf{1}+a_l$ and expanding, $\Pi-\mathbf{1}$ is the sum, over the
non-empty subsets of $\{1,\dots,n_\Pi\}$, of the ordered products of the corresponding $a_l$.
Submultiplicativity of the norm and $|a_l|\le\delta$ give
\begin{equation*}
  |\Pi-\mathbf{1}|\le \prod_{l=1}^{n_\Pi}\bigl(1+|u_l-\mathbf{1}|\bigr)-1
  \le e^{n_\Pi\delta}-1\le n_\Pi\delta\, e^{n_\Pi\delta}.
\end{equation*}
With $n_\Pi\le c_dL^{i}$, $L^{i}\xi\le 1$ and $c_dN_0\le\frac18$ this becomes
\begin{equation}\label{4.31:eq-terrace-averages-2}
  |\Pi-\mathbf{1}|\le 1.3\,c_dL^{i}\xi N_0\,e^{1.3c_dN_0}\le 2c_dL^{i}\xi N_0 ,
\end{equation}
because $1.3\,e^{1.3/8}\le 2$.

\emph{The average.} By \eqref{4.31:eq-analyticity-space-c}, $M^{i}(U)(c)$ is $P$ applied to a
normalised sum, with weights $L^{-4i}$ over base points, of ordered products of fine bond
variables along centre-to-centre contours of at most $c_dL^{i}$ fine bonds inside the double
$i$-block of $c$; these fine bonds lie in $Y$, so each product satisfies
\eqref{4.31:eq-terrace-averages-2}. Since the weights are non-negative and sum to one, the
normalised sum $S$ satisfies
\begin{equation*}
  |S-\mathbf{1}|\le 2c_dL^{i}\xi N_0\le 2c_dN_0\le r_P,
\end{equation*}
so $S$ lies in the domain of $P$. The bound $|P(V)-\mathbf{1}|\le c_P|V-\mathbf{1}|$ of
\eqref{4.31:eq-analyticity-space-c} then gives
\begin{equation*}
  |M^{i}(U)(c)-\mathbf{1}|\le 2c_dc_PL^{i}\xi N_0\le 2c_dc_PN_0\le \tfrac12 .
\end{equation*}

\emph{The logarithm.} On $\{|V-\mathbf{1}|\le\frac12\}$ the logarithm is defined by its power
series, and summing that series gives $|\log V|\le |V-\mathbf{1}|/(1-|V-\mathbf{1}|)\le
2|V-\mathbf{1}|$. Put $D(c):=-i\log M^{i}(U)(c)$. Then $M^{i}(U)(c)=e^{iD(c)}$ and
\begin{equation*}
  |D(c)|\le 2\,|M^{i}(U)(c)-\mathbf{1}|\le 4c_dc_PL^{i}\xi N_0\le c_{\mathrm{av}}L^{i}\xi N_0 ,
\end{equation*}
and the last quantity is at most $c_{\mathrm{av}}N_0$ because $L^{i}\xi\le1$. This is
\eqref{4.31:eq-terrace-averages-1}.

\emph{Holomorphy.} The map $\mathbf{A}\mapsto D(c)$ is the composition of the entire maps
$\mathbf{A}(b)\mapsto e^{i\xi\mathbf{A}(b)}$, of the polynomial map forming the normalised sum of
products, of $P$, which is holomorphic on its domain, and of $-i\log$, which is holomorphic on
$\{|V-\mathbf{1}|<1\}$ and hence on a neighbourhood of the closed ball of radius $\frac12$ in
which the values lie; the estimates above keep each intermediate value in the domain of
the next map. A composition of holomorphic maps is holomorphic.

\emph{Reality.} Let $\mathbf{A}$ be $\mathfrak{g}$-valued. Every fine factor
$e^{i\xi\mathbf{A}(b)}$ is unitary and lies in $G$. The map $P$ is the holomorphic extension of
the projection onto $G$, and at the $G$-valued configurations arising here it is that
projection, so $V:=M^{i}(U)(c)\in G$, and $|V-\mathbf{1}|\le\frac12$ by the estimate of the
average. Since $V$ is unitary, it is normal, and its eigenvalues have modulus one and lie within
distance $|V-\mathbf{1}|\le\frac12$ of $1$. The power series $\log V$ is then the matrix with the
same orthonormal eigenvectors whose eigenvalues are the principal logarithms of those of $V$;
these are purely imaginary, because the eigenvalues of $V$ have modulus one. Hence $\log V$ is
normal with purely imaginary spectrum, that is, anti-Hermitian, and $D(c)=-i\log V$ is Hermitian.
Next, $\det V=1$ and, since $\operatorname{tr}$ is normalised by $\operatorname{tr}\mathbf{1}=1$,
$\det e^{iD(c)}=e^{iN\operatorname{tr}D(c)}$, so
$\operatorname{tr}D(c)\in\frac{2\pi}{N}\mathbb{Z}$. The
set of $\mathfrak{g}$-valued $\mathbf{A}$ with $c_{\mathrm{av}}N_0\le1$ is a ball in a real
vector space, hence connected, and on it $\mathbf{A}\mapsto\operatorname{tr}D(c)$ is continuous,
as a composition of continuous maps (holomorphic in the interior); being valued in the discrete
set $\frac{2\pi}{N}\mathbb{Z}$, it is constant
there. At $\mathbf{A}=0$ every fine factor equals $\mathbf{1}$, so every product and the
normalised sum equal $\mathbf{1}$, and $P(\mathbf{1})=\mathbf{1}$ by
\eqref{4.31:eq-analyticity-space-c}; thus $M^{i}(U)(c)=\mathbf{1}$ and $D(c)=0$. Hence
$\operatorname{tr}D(c)=0$ for every $\mathfrak{g}$-valued $\mathbf{A}$ with
$c_{\mathrm{av}}N_0\le1$, and $D(c)$, being Hermitian and traceless, lies in $\mathfrak{g}$.

\emph{Tracelessness for complex $\mathbf{A}$.} The map $\mathbf{A}\mapsto\operatorname{tr}D(c)$
is holomorphic on the open convex set of $\mathfrak{g}^{\mathbb{C}}$-valued $\mathbf{A}$ with
$c_{\mathrm{av}}N_0<1$. Since $\mathfrak{g}^{\mathbb{C}}=\mathfrak{g}\oplus i\mathfrak{g}$, the
$\mathfrak{g}$-valued points of this set form a non-empty open subset of a maximal totally real
subspace, and there the map vanishes by the preceding paragraph. By the identity principle it
vanishes on the whole open set, and, being continuous, also when $c_{\mathrm{av}}N_0=1$. Hence
$D(c)$ is traceless, that is $D(c)\in\mathfrak{g}^{\mathbb{C}}$, for every
$\mathfrak{g}^{\mathbb{C}}$-valued $\mathbf{A}$ with $c_{\mathrm{av}}N_0\le1$.

The argument uses the closed form \eqref{4.31:eq-analyticity-space-c} only through the
representation of $M^{i}(U)(c)$ as $P$ applied to a normalised sum of products of at most
$c_dL^{i}$ fine factors. Suppose instead that $M^{i}$ is read as the $i$-fold iterate of one-step
averages with a projection after each step. Then this representation fails, and with it the
estimate of the product above. In that reading the conclusion of the lemma enters the
subsequent arguments as a hypothesis.
\end{proof}

\begin{remark}\label{4.31:lem-complex-marginal-bound}
The statement of this lemma is not written out here; parts of its proof are printed below.
\end{remark}

\begin{lemma}[Complex marginal bound: cube gauge and constant subtraction]
\label{4.31:prf-marginal-cube-gauge}
This lemma carries out the first two steps of the proof of
Lemma~\ref{4.31:lem-complex-marginal-bound} in the case there called Case~A; of the statements
from which Lemma~\ref{4.31:lem-complex-marginal-bound} is proved as an implication, only the
class-cube form of the real background statement \eqref{4.31:eq-standing-inputs-e} and
Lemma~\ref{4.31:lem-charts-gluing} enter these two steps. Assume the
following.
\begin{enumerate}
\item The conclusions of Lemma~\ref{4.31:lem-charts-gluing}\,(a): $K>s$,
$V_0\in\mathcal{V}_s(w_s)$, $U_0:=U_k(V_0)$, and for
$B\in(\mathfrak{g}^{\mathbb{C}})^{\mathrm{bonds}(\mathbb{T}_s)}$ with $|B|<4r_s$ the
configuration $W=W_{V_0}(B)=e^{i\eta\mathcal{H}(B;U_0)}U_0$ depends holomorphically on $B$,
and $|\mathcal{H}|,\ |\nabla^{\eta}_{U_0}\mathcal{H}|\le C_{19}\rho_s$.
\item The class-cube form of the real background statement, the standing input
\eqref{4.31:eq-standing-inputs-e}, taken from \cite[Theorem~1, (9)]{B-CMP102b}, at
$\varepsilon_1=w_s$, together with the conditions of (S1) in
\eqref{4.31:eq-radii-choice-a} under which it applies at $\varepsilon_1=w_s$.
\item Condition (S3) of \eqref{4.31:eq-radii-choice-a}: $32M'h_s\le 1$.
\item The Case~A geometry of the corresponding display: the polymer $X$ is
contained in the aligned cube $\square_0$ of $\mathbb{T}_s$-side $\delta_0\le M$, and
$\square_0$ lies in a class cube $\square$ of side $2M'$; in particular $\square_0$ does not
wrap around $\mathbb{T}_s$ and is a coordinate box.
\end{enumerate}
Write $h:=h_s$. Then the following hold.
\begin{enumerate}
\item[(a)] There are a $G$-valued gauge transformation $u_0$ on $\square$ and a field $H$ on
the bonds of $\square$, holomorphic in $B$, with $W^{u_0}=e^{i\eta H}$ and
$|H|,\ |\nabla^{\eta}H|\le h$ on $\square$; for real $B$, $H$ is Hermitian.
\item[(b)] Fix a fine site $x_0\in X$, put
$c_\mu:=H(\langle x_0,x_0+\eta e_\mu\rangle)$, $\varphi(x):=\sum_\mu c_\mu(x-x_0)_\mu$,
$u(x):=e^{i\varphi(x)}\in G^{\mathbb{C}}$ and $g:=uu_0$. Then $g$ is holomorphic in $B$,
$|g^{\pm1}|<2$ on $\square_0$, $g$ is unitary for real $B$, and the field $H^{1}$ on the
bonds of $\square_0$ defined by $e^{i\eta H^{1}}=W^{g}$ satisfies, with an absolute
constant $C_g$,
\begin{equation}\label{4.31:eq-marginal-cube-gauge-a}
|H^{1}|\le C_g(\delta_0+\eta)\,h,\qquad |\nabla^{\eta}H^{1}|\le C_g\,h
\qquad\text{on }\square_0 .
\end{equation}
\end{enumerate}
\end{lemma}

\begin{proof}
Gauge transformations act on bonds $b=\langle x,x+\eta e_\mu\rangle$ by
$U^{u}(b)=u(x)U(b)u(x+\eta e_\mu)^{-1}$; this is the action used in the definition of
$H^{1}$ below. We use repeatedly that $h=h_s=8(h_0+h_1)$ with $h_0$, $h_1$ as defined in the
next paragraph, that $\eta=L^{-k}\le 1$, and that $h\le 1$ by (S3).

\emph{The cube gauge.} By the class-cube statement \eqref{4.31:eq-standing-inputs-e}, which
follows from \cite[Theorem~1, (9)]{B-CMP102b} at $j=k$ and $\varepsilon_1=w_s$, applied on
the class cube $\square$, there is a $G$-valued $u_0$ with
\begin{equation*}
U_0^{u_0}=e^{i\eta A_0}\ \text{on }\square,\qquad
|A_0|,\ |\nabla^{\eta}A_0|<h_0:=B_3M'w_s .
\end{equation*}
For $b=\langle x,x+\eta e_\mu\rangle$ put $\mathcal{H}'(b):=\operatorname{Ad}(u_0(x))\mathcal{H}(b)$
and $h_1:=C_{19}\rho_s$. Since $u_0$ is unitary, hypothesis~1 gives $|\mathcal{H}'|\le h_1$.
Since $u_0$ transports $U_0$ to $e^{i\eta A_0}$ on $\square$, the gauge transform by $u_0$ of
the covariant difference $\nabla^{\eta}_{U_0}\mathcal{H}$ is the covariant difference of
$\mathcal{H}'$ along $e^{i\eta A_0}$, and this differs from the ordinary difference
$\nabla^{\eta}\mathcal{H}'$ by $\eta^{-1}(\operatorname{Ad}(e^{i\eta A_0})-1)$ applied to
$\mathcal{H}'$. Now $|\operatorname{Ad}(e^{i\eta A_0})-1|\le e^{2\eta|A_0|}-1\le 4\eta h_0$,
because $2\eta h_0\le h/4\le 1$ and $e^{y}-1\le 2y$ for $0\le y\le1$. Hence
\begin{equation*}
|\nabla^{\eta}\mathcal{H}'|
\le|\nabla^{\eta}_{U_0}\mathcal{H}|+\eta^{-1}|\operatorname{Ad}(e^{i\eta A_0})-1|\,
|\mathcal{H}'|
\le h_1+4h_0h_1\le 2h_1,
\end{equation*}
where the last inequality uses $4h_0\le\tfrac12 h_s\le 1$, which follows from
$h_s=8(h_0+h_1)$ and (S3).

Since $W=e^{i\eta\mathcal{H}}U_0$, on each bond $b=\langle x,x+\eta e_\mu\rangle$ of $\square$
\begin{align*}
W^{u_0}(b)&=u_0(x)e^{i\eta\mathcal{H}(b)}u_0(x)^{-1}\,u_0(x)U_0(b)u_0(x+\eta e_\mu)^{-1}\\
&=e^{i\eta\mathcal{H}'(b)}e^{i\eta A_0(b)}=:e^{i\eta H(b)} .
\end{align*}
The map $(a,a')\mapsto -i\log(e^{ia}e^{ia'})$ is analytic near $0$, vanishes at $(0,0)$ and,
by the Baker--Campbell--Hausdorff series, has Lipschitz constant at most $2$ in each argument
on $\{|a|,|a'|\le\tfrac12\}$. The arguments $a=\eta\mathcal{H}'(b)$ and $a'=\eta A_0(b)$ lie in
this set, since $\eta h_1,\ \eta h_0\le h/8\le\tfrac18$. Hence $|H|\le 2(h_1+h_0)$, and,
comparing the values at the bonds $b$ and $b+\eta e_\nu$,
\begin{equation*}
|\nabla^{\eta}H|\le 2(|\nabla^{\eta}\mathcal{H}'|+|\nabla^{\eta}A_0|)\le 2(2h_1+h_0).
\end{equation*}
Thus
\begin{equation*}
|H|,\ |\nabla^{\eta}H|\le 2(h_0+2h_1)\le 8(h_0+h_1)=h\qquad\text{on }\square .
\end{equation*}
The field $H$ is holomorphic in $B$, because $W$ is (hypothesis~1) and $u_0$ does not depend
on $B$.
For real $B$ the configuration $W$ is unitary, hence so is $W^{u_0}=e^{i\eta H}$, and $H$ is
Hermitian. This proves~(a).

\emph{Subtracting the constant by a complex gauge.} With $c_\mu$, $\varphi$, $u$ as in~(b),
$|c_\mu|\le h$ by~(a). By hypothesis~4 the cube $\square_0$ is a coordinate box of side
$\delta_0$, so
$|(x-x_0)_\mu|\le\delta_0$ for $x\in\square_0$, and by $\delta_0\le M\le M'$ and (S3)
\begin{equation*}
|\varphi|\le 4h\delta_0\le 4hM\le 4hM'\le\tfrac18\qquad\text{on }\square_0 .
\end{equation*}
On a bond $b=\langle x,x+\eta e_\mu\rangle$ of $\square_0$ the definition of $H^{1}$ reads
$e^{i\eta H^{1}(b)}=u(x)e^{i\eta H(b)}u(x+\eta e_\mu)^{-1}$, that is,
$e^{i\eta H^{1}}=(W^{u_0})^{u}=W^{g}$ with $g=uu_0$. The field $u$ is holomorphic in $B$
because the $c_\mu$ are, and $|u^{\pm1}|\le e^{|\varphi|}\le e^{1/8}$; as $u_0$ is unitary,
$|g^{\pm1}|\le e^{1/8}<2$. For real $B$ the $c_\mu$, hence $\varphi$, are Hermitian, so $u$
and $g$ are unitary.

For matrices $p$ with $|p|\le\tfrac12$ put
\begin{equation*}
\mathfrak{c}_\mu(p):=\frac{i}{\eta}\log\bigl(e^{-i(p+\eta c_\mu)}e^{ip}\bigr).
\end{equation*}
Since $\varphi(x+\eta e_\mu)=\varphi(x)+\eta c_\mu$, we have
$e^{-i\eta\mathfrak{c}_\mu(\varphi(x))}=e^{-i(\varphi(x)+\eta c_\mu)}e^{i\varphi(x)}$, hence
\begin{equation}\label{4.31:eq-marginal-cube-gauge-1}
u(x+\eta e_\mu)^{-1}=e^{-i\eta\mathfrak{c}_\mu(\varphi(x))}u(x)^{-1},\qquad
u(x+\eta e_\mu)=u(x)e^{i\eta\mathfrak{c}_\mu(\varphi(x))} .
\end{equation}
Differentiating $\sigma\mapsto e^{-i\sigma(p+\eta c)}e^{i\sigma p}$ gives the Duhamel formula
\begin{equation*}
e^{-i(p+\eta c)}e^{ip}=1-i\eta\int_0^1 e^{-i\sigma(p+\eta c)}\,c\,e^{i\sigma p}\,d\sigma ,
\end{equation*}
and
\begin{equation*}
e^{-i\sigma p}ce^{i\sigma p}-c=\sum_{n\ge1}(-i\sigma\operatorname{ad}_p)^n c/n!,
\end{equation*}
where every
term contains $[p,c]$; since $|\operatorname{ad}_p^{\,n}c|\le(2|p|)^{n-1}|[p,c]|$, for
$2|p|\le1$ the sum has norm at most $(e-1)|[p,c]|\le 1.72\,|[p,c]|$.

($\alpha$) The function $\mathfrak{c}_\mu$ is holomorphic on $|p|\le\tfrac12$, with
$|\mathfrak{c}_\mu|\le 4h$ there: the Duhamel integrand has norm at most $e^{1.02}h$, up to a
factor $1+O(\eta h)$, and the logarithm of $1-i\eta X$ equals $-i\eta X$ up to the same kind
of factor. By the Cauchy estimate on complex lines (a ball of radius $\tfrac14$ about a point
with $|p|\le\tfrac14$ lies in $|p|\le\tfrac12$), $\mathfrak{c}_\mu$ is $16h$-Lipschitz on
$|p|\le\tfrac14$.

($\beta$) For $|p|\le\tfrac12$,
\begin{equation*}
|\mathfrak{c}_\mu(p)-c_\mu|\le 2|[p,c_\mu]|+8\eta h^2 ;
\end{equation*}
here the constant $1.72$ of the commutator bound is enlarged to $2$ to absorb a factor
$1+O(\eta h)$, and the term $8\eta h^2$ comes from
$\log(1-i\eta X)=-i\eta X+O(\eta^2|X|^2)$. Moreover
$[\varphi(x),c_\mu]=\sum_{\nu\ne\mu}(x-x_0)_\nu[c_\nu,c_\mu]$, so that
\begin{equation*}
|[\varphi(x),c_\mu]|\le 3\delta_0\cdot 2h^2\le 8\delta_0h^2\qquad(x\in\square_0);
\end{equation*}
the constants $c_\nu$, $c_\mu$ are matrices and need not commute.

Let $b=\langle x,x+\eta e_\mu\rangle$ be a bond of $\square_0$ and define $H^{1/2}(b)$ by
$e^{i\eta H^{1/2}(b)}:=e^{i\eta H(b)}e^{-i\eta\mathfrak{c}_\mu(\varphi(x))}$. By
\eqref{4.31:eq-marginal-cube-gauge-1},
\begin{equation*}
e^{i\eta H^{1}(b)}=u(x)e^{i\eta H(b)}e^{-i\eta\mathfrak{c}_\mu(\varphi(x))}u(x)^{-1}
=u(x)e^{i\eta H^{1/2}(b)}u(x)^{-1},
\end{equation*}
so $H^{1}(b)=\operatorname{Ad}(u(x))H^{1/2}(b)$. By the Baker--Campbell--Hausdorff series,
whose first correction is the commutator of two terms of sizes $\eta h$ and $4\eta h$,
\begin{equation*}
\bigl|H^{1/2}(b)-\bigl(H(b)-\mathfrak{c}_\mu(\varphi(x))\bigr)\bigr|\le 8\eta h^2 .
\end{equation*}
A coordinate path inside $\square_0$ from $x_0$ to $x$ has length at most $4\delta_0$; moving
the bond of direction $\mu$ along it one lattice step at a time changes $H$ by at most
$\eta\sup_{\square}|\nabla^{\eta}H|$ per step; hence, by~(a) and by $|\varphi|\le\tfrac18$,
\begin{equation*}
|H(b)-c_\mu|\le 4\delta_0\sup_{\square}|\nabla^{\eta}H|\le 4\delta_0h,\qquad
|\operatorname{Ad}(u)|\le e^{2|\varphi|}\le e^{1/4}\le 2 .
\end{equation*}
Combining these bounds with
($\beta$) at $p=\varphi(x)$ (admissible since $|\varphi|\le\tfrac18$),
\begin{align*}
|H^{1}(b)|
&\le 2\bigl(|H(b)-c_\mu|+|c_\mu-\mathfrak{c}_\mu(\varphi(x))|+8\eta h^2\bigr)\\
&\le 2\bigl(4\delta_0h+16\delta_0h^2+16\eta h^2\bigr)\le C_g(\delta_0+\eta)h ,
\end{align*}
using $h\le1$. This is the first bound of \eqref{4.31:eq-marginal-cube-gauge-a}.

For the second bound let $b$ and $b+\eta e_\nu$ be bonds of $\square_0$, with
$b=\langle x,x+\eta e_\mu\rangle$. Since
\begin{equation*}
H^{1}(b+\eta e_\nu)=\operatorname{Ad}(u(x+\eta e_\nu))H^{1/2}(b+\eta e_\nu),
\end{equation*}
we have
\begin{align*}
\nabla^{\eta}_{\nu}H^{1}(b)
&=\operatorname{Ad}(u(x+\eta e_\nu))\nabla^{\eta}_{\nu}H^{1/2}(b)\\
&\quad+\eta^{-1}\bigl[\operatorname{Ad}(u(x+\eta e_\nu))-\operatorname{Ad}(u(x))\bigr]H^{1/2}(b).
\end{align*}
In the first term, $\eta H^{1/2}$ is the Baker--Campbell--Hausdorff combination of
$\eta H(b)$ and $-\eta\mathfrak{c}_\mu(\varphi(x))$, with Lipschitz constant at most $2$ in
each argument. Passing from $b$ to $b+\eta e_\nu$ changes $H$ by at most
$\eta|\nabla^{\eta}_{\nu}H|$, and changes $\varphi(x)$ by $\eta c_\nu$, hence, by the
$16h$-Lipschitz bound of ($\alpha$) on $|p|\le\tfrac14$, changes
$\mathfrak{c}_\mu(\varphi(x))$ by at most $16h\,\eta|c_\nu|$. Therefore, by~(a) and
$h\le\tfrac1{16}$ (from (S3)),
\begin{equation*}
|\nabla^{\eta}_{\nu}H^{1/2}|\le 2\bigl(|\nabla^{\eta}_{\nu}H|+16h|c_\nu|\bigr)
\le 2(h+16h^2)\le 4h ,
\end{equation*}
and $|\operatorname{Ad}(u(x+\eta e_\nu))|\le 2$. In the second term, by
\eqref{4.31:eq-marginal-cube-gauge-1} in direction $\nu$ the bracket equals
$\operatorname{Ad}(u(x))\bigl(\operatorname{Ad}(e^{i\eta\mathfrak{c}_\nu(\varphi(x))})-1\bigr)$,
and since $|\eta\mathfrak{c}_\nu|\le 4\eta h$ by ($\alpha$), its norm is at most
$2\,(e^{8\eta h}-1)\le 18\eta h$. Further
$|H^{1/2}|\le|H|+|\mathfrak{c}_\mu|+8\eta h^2\le 6h$. Altogether
\begin{equation*}
|\nabla^{\eta}H^{1}|\le 2\cdot 4h+18h\cdot 6h\le C_g\,h\qquad\text{on }\square_0 ,
\end{equation*}
which is the second bound of \eqref{4.31:eq-marginal-cube-gauge-a}. The constant $C_g$ is
absolute: it depends on none of $L$, $M$, $K$ or the couplings.
\end{proof}

\begin{lemma}[Complex marginal bound: minimiser path and Cauchy remainder]
\label{4.31:prf-marginal-minimiser-path}
This lemma is the concluding step, in Case~A, of the proof of
Lemma~\ref{4.31:lem-complex-marginal-bound}. Assume the following statements:
\begin{itemize}
\item Lemma~\ref{4.31:lem-background-membership} (membership of the complex background field
in the analyticity space);
\item the standing input \eqref{4.31:eq-standing-inputs-b} (holomorphy and gauge invariance of
the local pieces);
\item the standing input on class cubes, which provides the class cube containing $\square_0$
required by \textup{(N2)};
\item the local form \textup{(W1)}--\textup{(W3)} of the complex background field,
\eqref{4.31:eq-complex-background-b} (Definition~\ref{4.31:def-complex-background});
\item the nesting identities \textup{(N1)} and \textup{(N2)} for small complex minimisers,
\eqref{4.31:eq-window-counting-nesting-c} (Definition~\ref{4.31:def-window-counting-nesting});
\item Lemma~\ref{4.31:lem-terrace-averages} (terrace averages of a small complex field).
\end{itemize}
Let $K,s,t$, $k$, $j=k-t$, $\xi=L^{-j}$, $\eta=L^{-k}$, $X\in\mathbf{D}_j$,
$V_0\in\mathcal{V}_s(w_s)$, $B$ with
$|B|<4r_s$, $W=W_{V_0}(B)$, the function $f$ and the number $b$ be as in
Lemma~\ref{4.31:lem-complex-marginal-bound}; let $h=h_s$, and let $c_R$ and $C_5'$ be the
constants of the corresponding display, so that
$C_5'=1800\,c_{\mathrm{av}}^2C_g^2c_R^{-2}c_{\flat}^{-2}$, where $C_g$ is the absolute constant
of \eqref{4.31:eq-marginal-cube-gauge-a}, $c_{\mathrm{av}}$ that of
Lemma~\ref{4.31:lem-terrace-averages} and $c_{\flat}$ that of
Lemma~\ref{4.31:lem-cutoff-free-windows}. Suppose that Case~A of
the corresponding display holds, with the number $\delta_0$ and the aligned
scale-$j$ cube $\square_0$ fixed there. Let $g=uu_0$ be the $G^{\mathbb{C}}$-valued gauge
transformation on $\square_0$ of \eqref{4.31:eq-marginal-cube-gauge-a}, with the field
$\varphi$ of \eqref{4.31:eq-marginal-cube-gauge-a}, $|\varphi|\le 1/8$, and let $H^{1}$ be the
potential on $\square_0$ given there,
so that $W^{g}=e^{i\eta H^{1}}$ on $\square_0$. Write $d_j=d_j(X)$. Then
\begin{equation}\label{4.31:eq-marginal-minimiser-path-1}
\bigl|f\bigl((W,J^{\xi}(W))\lceil X\bigr)-f(\mathbf{1},0)\bigr|
\le C_5'\,M^2h^2x_j\,L^{-4t}\,b\,e^{-(\kappa-1)d_j(X)} .
\end{equation}
\end{lemma}

\begin{proof}
Throughout, for $n\ge 1$ we write $X^{\sim n}$ for the set obtained from $X$ by adding $n$
layers of cubes, and $X^{\sim -2}$ for the set obtained from $X$ by removing two layers of
cubes; $X^{\sim -2}$ is the domain of condition (iv) of \eqref{4.31:eq-analyticity-space-a}.
For a pair $(\mathbf{U},\mathbf{J})$ defined near $X$, $(\mathbf{U},\mathbf{J})\lceil X$
denotes its restriction to $X$.

\emph{Units and the datum.} We work in scale-$j$ units and put $\mathbf{A}^{1}:=L^{-t}H^{1}$
on $\square_0$. Since
\begin{equation*}
\xi L^{-t}=L^{-j-t}=L^{-k}=\eta,
\end{equation*}
we have
$e^{i\eta H^{1}}=e^{i\xi\mathbf{A}^{1}}$: the potential has homogeneity exponent one under the
change of units between the levels $k$ and $j$, and both exponentials describe the same
configuration $W^{g}$, whose regularity is the content of
\eqref{4.31:eq-marginal-cube-gauge-a}. The lattice gradient carries one further factor, so
that $|\nabla^{\xi}\mathbf{A}^{1}|=L^{-2t}|\nabla^{\eta}H^{1}|$ (exponent two). Further,
$\eta\le L^{-t}$ since $\eta=L^{-k}$ and $k\ge t$, and $\delta_0\le L^{-t}M(d_j+14)$ by the
choice of $\delta_0$ in Case~A of the corresponding display; in the second bound
the factor $L^{-t}$ is
the conversion of the scale-$j$ side of $\square_0$ into the units of $\mathbb{T}_s$, a
geometric factor and not a gain of the renormalisation. The bounds of
\eqref{4.31:eq-marginal-cube-gauge-a} on $H^{1}$ and on $\nabla^{\eta}H^{1}$, read through
the two identities just stated, therefore give
\begin{equation}\label{4.31:eq-marginal-minimiser-path-2}
\begin{aligned}
N_0:={}&\sup_{\square_0}\max\bigl\{|\mathbf{A}^{1}|,|\nabla^{\xi}\mathbf{A}^{1}|\bigr\}
\le\max\bigl\{L^{-t}C_g(\delta_0+\eta)h,\;L^{-2t}C_gh\bigr\}\\
\le{}&C_gL^{-2t}\bigl(M(d_j+14)+1\bigr)h\le C_gL^{-2t}M(d_j+15)h ,
\end{aligned}
\end{equation}
where the second inequality inserts $\delta_0\le L^{-t}M(d_j+14)$ and $\eta\le L^{-t}$ into
the first member of the maximum and uses $M(d_j+14)+1\ge1$ for the second member, and the
third inequality uses $M\ge1$.

In Case~A one has $L^{-t}(d_j+14)\le 1$, and hence
\begin{equation*}
L^{-2t}(d_j+15)=L^{-t}\cdot L^{-t}(d_j+14)+L^{-2t}\le 2L^{-t}\le 2 .
\end{equation*}
Consequently $N_0\le 2C_gMh$, and by condition (S4) of \eqref{4.31:eq-radii-choice-a},
\begin{equation*}
c_{\mathrm{av}}N_0\le 2c_{\mathrm{av}}C_gMh_s\le\min\{1,\beta_W\}.
\end{equation*}
Define
\begin{equation*}
D:=-i\log M^{\cdot}\bigl((W^{g})\lceil\square_0\bigr)=-i\log M^{\cdot}(e^{i\xi\mathbf{A}^{1}})
\end{equation*}
on the terrace bonds of the determining set $\mathfrak{B}(\square_0)$; these are bonds of
levels $i\le j$, and by \eqref{4.31:eq-analyticity-space-b} their blocks lie in $\square_0$.
We apply Lemma~\ref{4.31:lem-terrace-averages} with $Y=\square_0$ and
$\mathbf{A}=\mathbf{A}^{1}$; its hypothesis $c_{\mathrm{av}}\sup|\mathbf{A}^{1}|\le1$ holds
because $\sup|\mathbf{A}^{1}|\le N_0$ and $c_{\mathrm{av}}N_0\le 1$. The lemma gives that
$D$ is defined, depends holomorphically on $B$, is $\mathfrak{g}$-valued for real $B$, and
satisfies $|D(c)|\le c_{\mathrm{av}}L^{i}\xi\sup|\mathbf{A}^{1}|\le c_{\mathrm{av}}N_0$ at a
bond $c$ of level $i$, since $L^{i}\xi\le1$ for $i\le j$ (the top terrace level $i=j$ gives the
largest factor, equal to one). By \eqref{4.31:eq-marginal-minimiser-path-2} and
$c_{\mathrm{av}}N_0\le\beta_W$,
\begin{equation}\label{4.31:eq-marginal-minimiser-path-3}
N':=\sup|D|\le c_{\mathrm{av}}N_0\le c_{\mathrm{av}}C_gL^{-2t}M(d_j+15)h,
\qquad N'\le\beta_W ,
\end{equation}
for every $|B|<4r_s$, with constants independent of $K$. No bound on a gradient of $D$ is
needed, because \textup{(W1)} requires of its datum only a bound on $\sup|D|$.

\emph{The minimiser path and its membership in the analyticity space.} Let
\begin{equation*}
\alpha':=\min\{\alpha_{0,j}/2B_3,\;\alpha_{1,j}/2\}.
\end{equation*}
By Lemma~\ref{4.31:lem-cutoff-free-windows} and the definition
$A_{*}=\sup_{x\ge\gamma^{-2}}\alpha_0(x)/2B_3$, one has
$c_{\flat}x_j^{-1/2}\le\alpha'\le A_{*}$. Put
\begin{equation}\label{4.31:eq-marginal-minimiser-path-a}
R'_X:=N'^{-1}\min\Bigl\{\frac{\alpha'}{12B_0C_{19}},\;\beta_W,\;\frac{1}{10c_u}\Bigr\}
\ge\frac{c_R\,\alpha'}{N'} .
\end{equation}
The inequality holds because each member of the minimum is at least $c_R\alpha'$, where
$c_R=\min\{1/(12B_0C_{19}),\beta_W/A_{*},1/(10c_uA_{*})\}$: this is immediate for the first
member, and for the second and third members it follows from $\alpha'\le A_{*}$, which gives
$\beta_W\ge(\beta_W/A_{*})\alpha'$ and $1/(10c_u)\ge\alpha'/(10c_uA_{*})$. The second and third
members thus only change the constant.

Let $|\tau|<R'_X$. The datum $\tau D$ satisfies
\begin{equation*}
\sup|\tau D|\le|\tau|N'<R'_XN'\le\beta_W.
\end{equation*}
Hence, by \textup{(W1)} and \textup{(W2)}, the configuration and current
\begin{equation*}
\mathbf{U}_{\tau}:=\mathbf{U}[\tau D]=e^{i\xi\mathcal{H}_j(\square_0;\tau D,1)},
\qquad\mathbf{J}_{\tau}:=J^{\xi}(\mathbf{U}_{\tau}),
\end{equation*}
are defined on $\mathfrak{B}(\square_0)\supseteq X^{\sim3}$ and depend holomorphically on
$\tau$, the map $\tau\mapsto\tau D$ being linear with values in the ball of \textup{(W2)}. We
check the conditions (i)--(iv) of \eqref{4.31:eq-analyticity-space-a} (that is,
\cite[(1.11)--(1.16)]{B-CMP109}) for the pair $(\mathbf{U}_{\tau},\mathbf{J}_{\tau})\lceil X$,
written as $\mathbf{U}_{\tau}=U'U$ with $U\equiv\mathbf{1}$ and $U'=\mathbf{U}_{\tau}$. The
bounds below are those of \textup{(W1)} with $\beta=|\tau|N'$, followed by
$|\tau|N'<\alpha'/(12B_0C_{19})$, $\alpha'\le\alpha_{1,j}/2$, and
$\alpha'\le\alpha_{0,j}/2B_3\le\alpha_{0,j}$ (as $B_3\ge1$).
\begin{itemize}
\item[(i)] Condition \cite[(1.11)]{B-CMP109} holds because $\partial\mathbf{1}=\mathbf{1}$, and
condition \cite[(1.12)]{B-CMP109} holds with $A=0$.
\item[(ii)] With base point $\mathbf{1}$ the covariant gradient is the ordinary gradient, and
\begin{equation*}
|\mathcal{H}_j(\tau D)|,\ |\nabla^{\xi}\mathcal{H}_j(\tau D)|\le 2C_{19}B_0|\tau|N'
<\alpha'/6<\alpha_{1,j}.
\end{equation*}
\item[(iii)] The plaquette and current bounds are
\begin{equation*}
\begin{aligned}
|\partial\mathbf{U}_{\tau}-1|&\le 12C_{19}B_0|\tau|N'\xi^2<\alpha'\xi^2\le\alpha_{0,j}\xi^2,\\
|\mathbf{J}_{\tau}|&\le 8C_{19}B_0|\tau|N'<\tfrac23\alpha'<\alpha_{0,j},
\end{aligned}
\end{equation*}
and $J^{\xi}(\mathbf{1})=0$.
\item[(iv)] We apply \textup{(N1)} with $\beta=|\tau|N'$; since $|\tau|N'<1/(10c_u)$, the
third member of the minimum in \eqref{4.31:eq-marginal-minimiser-path-a} being used here,
the gauge transformation $\bar{u}$ of \textup{(N1)} is a block-constant lattice gauge
transformation with $|\bar{u}^{\pm1}|\le1+c_u|\tau|N'\le1.1$, and \textup{(N1)} gives, for
every $n$ in the range of \cite[(1.16)]{B-CMP109}, on $X^{\sim-2}$,
\begin{equation*}
U_n(X,M^{n}\mathbf{U}_{\tau})=\mathbf{U}_{\tau}^{\bar{u}^{-1}},\qquad
J_n(X,M^{n}\mathbf{U}_{\tau})=(L^{n}\xi)^3\operatorname{Ad}(\bar{u}^{-1})J^{\xi}(\mathbf{U}_{\tau}).
\end{equation*}
Hence, for every $n$ in this range, on $X^{\sim-2}$,
\begin{equation}\label{4.31:eq-marginal-minimiser-path-4}
|\partial U_n(X,M^{n}\mathbf{U}_{\tau})-1|\le|\bar{u}^{-1}|\,|\bar{u}|\,
|\partial\mathbf{U}_{\tau}-1|\le1.21\,\alpha'\xi^2\le\frac{1.21\,\alpha_{0,j}\xi^2}{2B_3}
<\alpha_{0,j}\xi^2,
\end{equation}
using (iii), $\alpha'\le\alpha_{0,j}/2B_3$ and $1.21<2\le2B_3$; and
\begin{equation}\label{4.31:eq-marginal-minimiser-path-5}
\begin{aligned}
|J_n(X,M^{n}\mathbf{U}_{\tau})|&=(L^{n}\xi)^3\,
|\operatorname{Ad}(\bar{u}^{-1})J^{\xi}(\mathbf{U}_{\tau})|
\le(L^{n}\xi)^2\cdot1.21\cdot\tfrac23\alpha'\\
&<2B_3\alpha'(L^{n}\xi)^2\le\alpha_{0,j}(L^{n}\xi)^2 ,
\end{aligned}
\end{equation}
where the current carries the exponent three against the exponent two of the threshold, and
$L^{n}\xi\le1$ and $B_3\ge1$ are used. Both bounds are uniform in $n$.
\end{itemize}
Therefore $(\mathbf{U}_{\tau},\mathbf{J}_{\tau})\lceil X\in
\mathfrak{U}^{c}_j(X,\alpha_{0,j},\alpha_{1,j})$ for every $|\tau|<R'_X$, uniformly in $K$.

\emph{The function $\psi$.} Put $\psi(\tau):=f\bigl((\mathbf{U}_{\tau},\mathbf{J}_{\tau})\lceil
X\bigr)$ for $|\tau|<R'_X$. It is holomorphic there: $\tau\mapsto\mathcal{H}_j(\square_0;\tau
D,1)$ is holomorphic by \textup{(W2)}, $J^{\xi}$ is a rational function of the bond variables,
and $f$ is holomorphic on the analyticity space by \eqref{4.31:eq-standing-inputs-b}. Since the
pair lies in $\mathfrak{U}^{c}_j(X,\alpha_{0,j},\alpha_{1,j})$, the hypothesis on $f$ gives
$|\psi|\le b\,e^{-\kappa d_j}$ on the whole disc. Finally $\psi(0)=f(\mathbf{1},0)$, because
$\mathcal{H}_j(\square_0;0,1)=0$ by \textup{(W1)} and $J^{\xi}(\mathbf{1})=0$.

\emph{The endpoint.} This paragraph is used only when $R'_X\ge2$, so that $\tau=1$ lies in
the disc; then also $N'\le\beta_W/2$, since $N'R'_X\le\beta_W$. By
\cite[(1.19)]{B-CMP109}, which applies through \eqref{4.31:eq-standing-inputs-b} to the
$G^{\mathbb{C}}$-valued gauge transformation $g=uu_0$ (with $|\varphi|\le1/8$), and then by
\textup{(N2)}, whose datum bound $\sup|D|\le\beta_W$ is \eqref{4.31:eq-marginal-minimiser-path-3},
\begin{equation*}
f\bigl((W,J^{\xi}(W))\lceil X\bigr)
=f\bigl((W^{g},\operatorname{Ad}(g)J^{\xi}(W))\lceil X\bigr)
=f\bigl((\mathbf{U}[D],J^{\xi}\mathbf{U}[D])\lceil X\bigr)=\psi(1).
\end{equation*}
Equivalently, the second equality may be read as follows: by \textup{(N2)},
$W^{g}=\mathbf{U}[D]^{\bar{u}_1}$ near $X$, and
$\operatorname{Ad}(g)J^{\xi}(W)=J^{\xi}(W^{g})=\operatorname{Ad}(\bar{u}_1)J^{\xi}\mathbf{U}[D]$;
one then applies \cite[(1.19)]{B-CMP109} with the gauge transformation $\bar{u}_1$.

\emph{Absence of the linear term.} Let
$\Psi(D'):=f\bigl((\mathbf{U}[D'],J^{\xi}\mathbf{U}[D'])\lceil X\bigr)$ on the ball
$\{\sup|D'|<N'R'_X\}$ of $\mathfrak{g}^{\mathbb{C}}$-valued functions on the terrace bonds of
$\mathfrak{B}(\square_0)$. On this ball the estimates of the membership paragraph, with
$|\tau|N'$ replaced by $\sup|D'|$, place the pair in
$\mathfrak{U}^{c}_j(X,\alpha_{0,j},\alpha_{1,j})$; $\Psi$ is holomorphic there by
\textup{(W2)} and \eqref{4.31:eq-standing-inputs-b}, and $\psi(\tau)=\Psi(\tau D)$. For a
constant $v\in G$, $\operatorname{Ad}(v)$ is a linear isometry of this ball, and
\begin{equation*}
\begin{aligned}
\Psi(\operatorname{Ad}(v)D')
&=f\bigl((\mathbf{U}[D']^{v},J^{\xi}(\mathbf{U}[D']^{v}))\lceil X\bigr)\\
&=f\bigl((\mathbf{U}[D']^{v},\operatorname{Ad}(v)J^{\xi}\mathbf{U}[D'])\lceil X\bigr)
=\Psi(D').
\end{aligned}
\end{equation*}
The first equality is \textup{(W3)}: it gives
\begin{equation*}
\mathbf{U}[\operatorname{Ad}(v)D']=e^{i\xi\operatorname{Ad}(v)\mathcal{H}_j(D')}
=\mathbf{U}[D']^{v},
\end{equation*}
the transform of $\mathbf{U}[D']$ by the constant gauge transformation
$v$. The second is the equivariance of $J^{\xi}$ under constant gauge transformations, and the
third is \cite[(1.19)]{B-CMP109}. Since $\Psi\circ\operatorname{Ad}(v)=\Psi$ on the ball and
$\operatorname{Ad}(v)$ is linear, differentiation at $0$ gives
$D\Psi_0\circ\operatorname{Ad}(v)=D\Psi_0$: the $\mathbb{C}$-linear functional $D\Psi_0$ is
$\operatorname{Ad}(G)$-invariant. It acts on a finite-dimensional space and so commutes with
integration over the Haar probability measure $dv$ of $G$; therefore
\begin{equation*}
D\Psi_0(D')=\int_G D\Psi_0(\operatorname{Ad}(v)D')\,dv
=D\Psi_0\Bigl(\int_G\operatorname{Ad}(v)D'\,dv\Bigr)=D\Psi_0(0)=0,
\end{equation*}
because, bond by bond, $\int_{SU(N)}\operatorname{Ad}(v)Y\,dv=0$ for every
$Y\in\mathfrak{sl}(N,\mathbb{C})$ and every $N\ge2$, by Schur's lemma in the form recorded
in \eqref{4.31:eq-window-counting-nesting-g}. By the chain rule,
$\psi'(0)=D\Psi_0(D)=0$. This vanishing is the marginal gain, and it is here that the
non-abelian character of $G$ enters: for an abelian group $\operatorname{Ad}(v)$ is the
identity and the argument gives nothing.

\emph{The Cauchy remainder.} Suppose first $R'_X\ge2$. Since $1<R'_X$ and $\psi'(0)=0$, the
Taylor series of $\psi$ at $0$ converges at $\tau=1$ and
$\psi(1)-\psi(0)=\sum_{n\ge2}\psi^{(n)}(0)/n!$. The Cauchy estimates on the circle of radius
$r<R'_X$, together with $|\psi|\le b\,e^{-\kappa d_j}$, give
$|\psi^{(n)}(0)/n!|\le b\,e^{-\kappa d_j}r^{-n}$, and letting $r\to R'_X$ gives
$|\psi^{(n)}(0)/n!|\le b\,e^{-\kappa d_j}R_X'^{-n}$. Hence, using $R_X'^{-1}\le1/2$,
\begin{equation*}
|\psi(1)-\psi(0)|\le b\,e^{-\kappa d_j}\sum_{n\ge2}R_X'^{-n}
=\frac{b\,e^{-\kappa d_j}R_X'^{-2}}{1-R_X'^{-1}}\le2\,b\,e^{-\kappa d_j}R_X'^{-2},
\end{equation*}
and by the paragraphs on $\psi$ and on the endpoint the left side equals
$|f((W,J^{\xi}(W))\lceil X)-f(\mathbf{1},0)|$. Suppose next $R'_X<2$. By
Lemma~\ref{4.31:lem-background-membership}, $(W,J^{\xi}(W))\lceil X\in
\mathfrak{U}^{c}_j(X,\alpha_{0,j},\alpha_{1,j})$, so
$|f((W,J^{\xi}(W))\lceil X)|\le b\,e^{-\kappa d_j}$; and
$|f(\mathbf{1},0)|=|\psi(0)|\le b\,e^{-\kappa d_j}$. Therefore, as $R_X'^{-2}>1/4$,
\begin{equation*}
\bigl|f\bigl((W,J^{\xi}(W))\lceil X\bigr)-f(\mathbf{1},0)\bigr|\le2\,b\,e^{-\kappa d_j}
\le8\,b\,e^{-\kappa d_j}R_X'^{-2}.
\end{equation*}
In both cases
\begin{equation}\label{4.31:eq-marginal-minimiser-path-6}
\bigl|f\bigl((W,J^{\xi}(W))\lceil X\bigr)-f(\mathbf{1},0)\bigr|\le8\,b\,e^{-\kappa d_j}R_X'^{-2}.
\end{equation}

\emph{Collection of the constants.} By \eqref{4.31:eq-marginal-minimiser-path-a},
$\alpha'\ge c_{\flat}x_j^{-1/2}$ and \eqref{4.31:eq-marginal-minimiser-path-3},
\begin{equation*}
R_X'^{-2}\le c_R^{-2}\alpha'^{-2}N'^2
\le c_R^{-2}c_{\flat}^{-2}x_j\cdot c_{\mathrm{av}}^2C_g^2L^{-4t}M^2(d_j+15)^2h^2 .
\end{equation*}
For $d\ge0$, using $1+y\le e^{y}$,
\begin{equation*}
(d+15)^2=225\,(1+d/15)^2\le225\,e^{2d/15}\le225\,e^{d}.
\end{equation*}
Inserting these two bounds into \eqref{4.31:eq-marginal-minimiser-path-6}, with $d=d_j$,
gives
\begin{equation*}
\begin{aligned}
\bigl|f\bigl((W,J^{\xi}(W))\lceil X\bigr)-f(\mathbf{1},0)\bigr|
&\le8\cdot225\cdot c_{\mathrm{av}}^2C_g^2c_R^{-2}c_{\flat}^{-2}\cdot M^2h^2x_j\cdot
L^{-4t}\,b\,e^{-(\kappa-1)d_j}\\
&=C_5'\,M^2h^2x_j\,L^{-4t}\,b\,e^{-(\kappa-1)d_j(X)},
\end{aligned}
\end{equation*}
since $8\cdot225=1800$. This is \eqref{4.31:eq-marginal-minimiser-path-1}. Combined with the
treatment of Case~B of the corresponding display in the proof of
Lemma~\ref{4.31:lem-complex-marginal-bound}, this gives the bound of that lemma.
\end{proof}

\begin{remark}[Cutoff uniformity of the marginal bound and the factor $L^{-4t}$]
\label{4.31:rem-marginal-uniformity}
Fix the depth $s$ and the gap $t$. The constants $c_R$, $C_5'$ and $48e^{14}$ of the complex
marginal bound for local pieces (Lemma~\ref{4.31:lem-complex-marginal-bound}), and the
constants of~\eqref{4.31:eq-complex-background-b} and~\eqref{4.31:eq-window-counting-nesting-c}
used in its proof, depend on none of the following: the cutoff $K$; the scale $j$; the
level $n$ of the nesting identities; the age of a region; the volume; the torus; the
position of $X$. The auxiliary quantities of the proof are the radius $\alpha'$, the sizes
$N_0$ and $N'$ of the local datum, and the radius $R'_X$. They are bounded in terms of these
constants and of $M$, $L^{-t}$, $h_s$, $x_j\in I_{s+t}$ and $d_j(X)$ alone.

The factor $L^{-4t}$ in the bound is not a gain per renormalisation step. It is the
product of two contributions.
\begin{itemize}
\item The first is marginality: there is no linear term, and what remains is the Cauchy
remainder $R_X'^{-2}$.
\item The second is a change of units: $N'^2$ carries
$L^{-4t}=(L^{-t})^2\cdot(L^{-t})^2$. This comes from the homogeneity of exponent one of
the potential between the levels $k$ and $j$, and from the size
$L^{-t}M(d_j+14)$ of the cube $\square_0$ in the units of $\mathbb{T}_s$.
\end{itemize}
\end{remark}

\begin{proof}
We go through every place at which $K$, $j$ or $n$ could enter the constants of
Lemma~\ref{4.31:lem-complex-marginal-bound}. We use $j=K-s-t$ throughout.

\emph{The radius $\alpha'$.} Recall that
\begin{equation*}
\alpha'=\min\{\alpha_{0,j}/2B_3,\ \alpha_{1,j}/2\}.
\end{equation*}
It enters the proof only through the two bounds
\begin{equation*}
c_{\flat}\,x_j^{-1/2}\le\alpha'\le A_{*}.
\end{equation*}
By Lemma~\ref{4.31:lem-cutoff-free-windows}, applied with $u=s+t$ and resting on the anchored
lower slope of the coupling flow, we have
$x_j=x_{K-(s+t)}\in I_{s+t}$. The window $I_{s+t}$ and the constant $c_{\flat}$ do not
depend on $K$. The constant $A_{*}$ is the supremum of $\alpha_0(x)/2B_3$ over
$x\ge\gamma^{-2}$, so it does not depend on $K$ either.

\emph{The sizes $N_0$ and $N'$.} These enter through three things.
\begin{itemize}
\item The first is the bound $\delta_0\le L^{-t}M(d_j+14)$ on the side of
$\square_0$ in the units of $\mathbb{T}_s$
(see the corresponding display).
\item The second is the geometry of $X$ at scale $j$, measured in the units of
$\mathbb{T}_s$.
\item The third consists of the constants $h=h_s$, $C_g$ and $c_{\mathrm{av}}$. Of these, $C_g$ is
the absolute constant of~\eqref{4.31:eq-marginal-cube-gauge-a}, $c_{\mathrm{av}}$ is the
averaging constant, and $h_s$ does not depend on $K$.
\end{itemize}
None of these depends on $K$. Consequently the bound~\eqref{4.31:eq-marginal-minimiser-path-a}
on $N'$ does not depend on $K$.

\emph{The complex background field.} The constants of its global bound
in~\eqref{4.31:eq-complex-background-b} depend on $(d,L,M)$ only. By the statement of
that bound, they do not depend on the position or the size of $\square_0$, on $j$, or on
$K$.

\emph{The nesting statements.} The bound in~\eqref{4.31:eq-window-counting-nesting-c} on
the block-constant gauge transformations $\bar{u}$ is uniform in $n$. In that bound the
scales $n$ and $j$ are compared only through the inequality $L^{n}\xi\le1$. The nesting
identity in~\eqref{4.31:eq-window-counting-nesting-c} is an identity and carries no
constants.

\emph{The second case.} Consider the case $L^{-t}(d_j(X)+14)>1$
of the corresponding display. There the only constant is the absolute
number $48e^{14}$.

\emph{The constant $c_R$.} It is given by
\begin{equation*}
c_R=\min\{1/(12B_0C_{19}),\ \beta_W/A_{*},\ 1/(10c_uA_{*})\}.
\end{equation*}
It therefore depends only on $B_0$, $C_{19}$, $\beta_W$, $c_u$ and $A_{*}$, and none of
these depends on $K$. Hence the constant
\begin{equation*}
C_5'=1800\,c_{\mathrm{av}}^2C_g^2c_R^{-2}c_{\flat}^{-2}
\end{equation*}
does not depend on $K$ either. Hence none of the constants listed depends on $K$, $j$, $n$,
an age, the volume, the torus or the position of $X$, and the auxiliary quantities
$\alpha'$, $N_0$, $N'$, $R'_X$ are bounded in terms of them and of $M$, $L^{-t}$, $h_s$,
$x_j$ and $d_j(X)$ alone. This proves the first assertion.

We turn to the factor $L^{-4t}$. By the absence of a linear term in the proof of
Lemma~\ref{4.31:lem-complex-marginal-bound}, the Cauchy estimate leaves the remainder
$8be^{-\kappa d_j}R_X'^{-2}$, and by~\eqref{4.31:eq-marginal-minimiser-path-a} we have
$R_X'^{-2}\le c_R^{-2}\alpha'^{-2}N'^2$. One factor $L^{-t}$ in $N'$ is the homogeneity of
exponent one of the potential between the levels $k$ and $j$, namely
$\mathbf{A}^{1}=L^{-t}H^{1}$ for the same configuration; the gradient term carries
$L^{-2t}$ directly, since $|\nabla^{\xi}\mathbf{A}^{1}|=L^{-2t}|\nabla^{\eta}H^{1}|$. The
other factor $L^{-t}$ is the bound $\delta_0\le L^{-t}M(d_j+14)$, together with
$\eta\le L^{-t}$; it is the conversion of the scale-$j$ side of $\square_0$ into the units
of $\mathbb{T}_s$, a matter of geometry, not a renormalisation gain.
By~\eqref{4.31:eq-marginal-minimiser-path-a},
\begin{equation*}
N'\le c_{\mathrm{av}}C_gL^{-2t}M(d_j+15)h ,
\end{equation*}
so $N'^2$ carries $L^{-4t}=(L^{-t})^2\cdot(L^{-t})^2$, as stated.

Apart from the factor $L^{-4t}$, the proof of Lemma~\ref{4.31:lem-complex-marginal-bound}
produces no further power of $L^{-t}$; in particular no gain of the form $L^{-2}$ per step
is asserted or used. The geometric decay
in $t$ of the sum of the pieces $\mathsf{E}_{K,s,t}$ comes only from the real bound
\eqref{4.31:eq-standing-inputs-c}. That bound is \cite[Theorem~2, (2.43)]{B-CMP119}. It
is interpolated by the two-constants interpolation lemma of this section, which gives the
exponent $L^{-\beta t/2}$. The sum of the pieces $\mathsf{R}_{K,s,t}$ decays through the
factor $x_j^{1-\kappa_0/2}$ and through the anchored lower slope of the coupling flow.
\end{proof}

\begin{lemma}[The two-constants estimate on the disc]\label{4.31:lem-two-constants}
Let $\psi$ be holomorphic on the open unit disc $\{\zeta\in\mathbb{C}:|\zeta|<1\}$. Let
$\mathsf{M}$ and $\mathsf{m}$ be non-negative numbers such that $|\psi|\le\mathsf{M}$ on the
disc and $|\psi|\le\mathsf{m}$ on the open interval $]-1,1[$. Then
\begin{equation*}
|\psi(\zeta)|\le(\mathsf{m}\mathsf{M})^{1/2}\qquad\text{for }|\zeta|\le\tfrac14 .
\end{equation*}
\end{lemma}

\begin{proof}
If $\mathsf{m}\ge\mathsf{M}$, then $\mathsf{M}^2\le\mathsf{m}\mathsf{M}$, and the assertion follows
from $|\psi|\le\mathsf{M}$. If $\mathsf{m}=0$, then $\psi$ vanishes on $]-1,1[$, so
$\psi\equiv0$ by the identity theorem. We may therefore assume $0<\mathsf{m}<\mathsf{M}$.
If $\zeta$ is real, then $|\psi(\zeta)|\le\mathsf{m}<(\mathsf{m}\mathsf{M})^{1/2}$. The
function $\zeta\mapsto\overline{\psi(\bar\zeta)}$ is holomorphic on the disc, satisfies the same
two bounds, and has modulus $|\psi(\bar\zeta)|$ at $\zeta$. It therefore suffices, by this
symmetry, to treat $\zeta$ with $\operatorname{Im}\zeta>0$.

Let $D^{+}=\{\zeta:|\zeta|<1,\ \operatorname{Im}\zeta>0\}$ be the upper half-disc. For
$\zeta=e^{i\theta}$ with $0<\theta<\pi$ one has
$(1+\zeta)/(1-\zeta)=i\cot(\theta/2)$. Hence the M\"obius map $\zeta\mapsto(1+\zeta)/(1-\zeta)$
maps $D^{+}$ onto the open first quadrant. It sends the diameter $]-1,1[$ onto the positive real
axis and the open upper arc onto the positive imaginary axis. Put
\begin{equation*}
\omega(\zeta)=1-\frac{2}{\pi}\arg\frac{1+\zeta}{1-\zeta},\qquad\zeta\in D^{+},
\end{equation*}
with the argument taken in $]0,\pi/2[$. Then $\omega$ is harmonic on $D^{+}$, being an affine
function of the imaginary part of a holomorphic branch of $\log((1+\zeta)/(1-\zeta))$. It takes
values in $]0,1[$, tends to $1$ at points of the diameter and to $0$ at points of the arc.

Consider
\begin{equation*}
u=\log|\psi|-\omega\log\mathsf{m}-(1-\omega)\log\mathsf{M}\quad\text{on }D^{+}.
\end{equation*}
It is subharmonic, since $\log|\psi|$ is subharmonic and the other two terms are harmonic. Since
$0\le\omega\le1$ and $\log\mathsf{m}<\log\mathsf{M}$, it is bounded above by
$\log\mathsf{M}-\log\mathsf{m}$.

Consider first a point $x\in\,]-1,1[$. By upper semicontinuity of $\log|\psi|$ and continuity of
$\omega$ up to $x$, we have $\limsup_{\zeta\to x}u(\zeta)\le\log\mathsf{m}-\log\mathsf{m}=0$.
Next consider a point of the open arc. There $\log|\psi|\le\log\mathsf{M}$ and $\omega\to0$, so
the $\limsup$ is at most $0$. Thus $\limsup u\le0$ at every boundary point of $D^{+}$ except
$\pm1$.

The extended maximum principle applies to this situation: a subharmonic function bounded above on
a bounded domain, whose $\limsup$ is at most $0$ at every boundary point outside a finite set, is
at most $0$ on the domain. Hence $u\le0$ on $D^{+}$, that is,
\begin{equation}\label{4.31:eq-two-constants-1}
|\psi(\zeta)|\le\mathsf{m}^{\omega(\zeta)}\mathsf{M}^{1-\omega(\zeta)}
=\mathsf{M}\Big(\frac{\mathsf{m}}{\mathsf{M}}\Big)^{\omega(\zeta)},\qquad\zeta\in D^{+}.
\end{equation}

Now fix $\zeta\in D^{+}$ with $|\zeta|\le\frac14$ and put $\varrho=|\zeta|$. The map
$\zeta\mapsto(1+\zeta)/(1-\zeta)$ sends the circle $|\zeta|=\varrho$, which is symmetric about
the real axis and does not pass through the pole $1$, onto a circle symmetric about the real axis.
This image circle meets the real axis in the images $(1+\varrho)/(1-\varrho)$ and
$(1-\varrho)/(1+\varrho)$ of $\pm\varrho$. Its centre and radius are therefore
\begin{equation*}
\frac{1+\varrho^2}{1-\varrho^2}
\quad\text{and}\quad
\frac{2\varrho}{1-\varrho^2}.
\end{equation*}
Every point of a circle with centre $c>0$ and radius $r<c$ has argument at most
$\arcsin(r/c)$. Here $r/c=2\varrho/(1+\varrho^2)$. Moreover
$\sin(2\arctan\varrho)=2\varrho/(1+\varrho^2)$ with $2\arctan\varrho\in[0,\pi/2]$. Hence
\begin{equation*}
\arg\frac{1+\zeta}{1-\zeta}\le\arcsin\frac{2\varrho}{1+\varrho^2}=2\arctan\varrho
\le2\arctan\tfrac14 .
\end{equation*}
Consequently
\begin{equation*}
\omega(\zeta)\ge1-\frac{4}{\pi}\arctan\tfrac14=0.688\ldots\ge\tfrac12 .
\end{equation*}
Since $0<\mathsf{m}/\mathsf{M}<1$, the function $t\mapsto(\mathsf{m}/\mathsf{M})^{t}$ is
decreasing. Inserting this into \eqref{4.31:eq-two-constants-1} gives
\begin{equation*}
|\psi(\zeta)|\le\mathsf{M}\Big(\frac{\mathsf{m}}{\mathsf{M}}\Big)^{\omega(\zeta)}
\le\mathsf{M}\Big(\frac{\mathsf{m}}{\mathsf{M}}\Big)^{1/2}=(\mathsf{m}\mathsf{M})^{1/2}.
\end{equation*}
This proves the lemma.
\end{proof}

\begin{lemma}[Fixed-depth convergence: uniform analyticity on the tube]\label{4.31:prf-tube-analyticity}
This is clause (D3) of Proposition~\ref{4.31:prop-fixed-depth-convergence}. Assume the following
statements:
\begin{itemize}
\item the standing inputs of Definition~\ref{4.31:def-standing-inputs}: the bound
$|\beta_j(g_{j-1})|\le C_\beta$ of \eqref{4.31:eq-standing-inputs-a} and the real irrelevance
bound \eqref{4.31:eq-standing-inputs-c};
\item the counting bound \eqref{4.31:eq-window-counting-nesting-b} of
Definition~\ref{4.31:def-window-counting-nesting}, and gauge invariance of the real functionals,
\cite[(1.19)]{B-CMP109} together with \eqref{4.31:eq-window-counting-nesting-e};
\item Lemmas~\ref{4.31:lem-cutoff-free-windows}, \ref{4.31:lem-window-characteristic},
\ref{4.31:lem-charts-gluing}, \ref{4.31:lem-background-membership},
\ref{4.31:lem-complex-marginal-bound} and \ref{4.31:lem-two-constants}.
\end{itemize}
Then the following holds. Let $w_s,r_s$ be the radii of Definition~\ref{4.31:def-radii-choice}, write
$\mathcal{V}_s=\mathcal{V}_s(w_s)$, $\mathfrak{D}_s=\mathfrak{D}_s(w_s,r_s)$, and for $K\in\mathcal{K}$
with $K>s$ put $k=K-s$, and, for $0\le t<k$, $j=k-t$, $\xi=L^{-j}$, $\eta=L^{-k}$. For
$V\in\mathcal{V}_s$ put
\[
e^{K}_{t,X,z}(V):=\mathbf{E}^{(j)}(X,U_k(V),z),\qquad r^{K}_{t,X}(V):=\mathbf{R}^{(j)}(X,U_k(V)).
\]
Then $\chi_{k}\equiv 1$ on $\mathcal{V}_s(w_s+16r_s)$, and the functions $\mathcal{A}_{K,s}$,
$\mathsf{E}_{K,s,t}$, $\mathsf{R}_{K,s,t}$ and $\mathsf{S}_{K,s}$ of
Definition~\ref{4.31:def-never-large-exponent} and the functions $e^{K}_{t,X,z}$, $r^{K}_{t,X}$
extend from $\mathcal{V}_s$ to holomorphic functions
on $\mathfrak{D}_s$ which satisfy, with $d(X)=d_j(X)$ the tree length of $X$ in units of the cubes
of side $ML^{-t}$ (so measured, $d(X)$ does not depend on $K$),
\begin{align*}
&\text{(a)}\quad |\mathcal{A}_{K,s}|\le C_{\mathcal{A}}(w_s+r_s)^2|\mathbb{T}_s|;\\
&\text{(b)}\quad |e^{K}_{t,X,z}|\le E_0e^{-\kappa d(X)},\qquad
|r^{K}_{t,X}|\le (x_{-}^{(s+t)})^{-\kappa_0/2}e^{-\kappa d(X)};\\
&\text{(c)}\quad |\mathsf{E}_{K,s,t}|\le C_E(x_{+}^{(s+t)})^{1/2}L^{-\beta t/2}|\mathbb{T}_s|;\\
&\text{(d)}\quad |\mathsf{R}_{K,s,t}|\le C_R(x_{-}^{(s+t)})^{1-\kappa_0/2}|\mathbb{T}_s|;\\
&\text{(e)}\quad |\mathsf{S}_{K,s}|\le C_s,
\end{align*}
where $\beta\in\,]0,1[$ is the exponent of \eqref{4.31:eq-standing-inputs-c}, the constants
$C_{\mathcal{A}},C_R,C_E$ are those of \eqref{4.31:eq-tube-analyticity-a} and $C_s$ that of
\eqref{4.31:eq-tube-analyticity-1} below; none of them depends on $K$.
\end{lemma}

\begin{proof}
Throughout, $C_m=6(B_3+8C_{19}(B_3+B_0))$, $\rho_s$ and $h_s$ are as in
Definition~\ref{4.31:def-radii-choice}, $C_5'$ is the constant of
Lemma~\ref{4.31:lem-complex-marginal-bound}, $K_0$ that of \eqref{4.31:eq-window-counting-nesting-b},
$E_0,\kappa,\kappa_0$ are the constants of Ba{\l}aban's bounds entering
Lemma~\ref{4.31:lem-background-membership}, and $E_1,\beta$ those of \eqref{4.31:eq-standing-inputs-c}.
We put
\begin{equation}\label{4.31:eq-tube-analyticity-a}
\begin{gathered}
C_{\mathcal{A}}:=6C_m^2,\qquad C_R:=(C_5'M^2h_s^2+48e^{14})K_0M^{-4},\\
C_E':=(C_5'M^2h_s^2+48e^{14})E_0K_0+C_\beta C_{\mathcal{A}}(w_s+r_s)^2,\qquad
C_E:=(E_1C_E')^{1/2},\\
\mathsf{M}_t:=C_E'\,x_{+}^{(s+t)}|\mathbb{T}_s|,\qquad \mathsf{m}_t:=E_1L^{-\beta t}|\mathbb{T}_s|.
\end{gathered}
\end{equation}

\emph{Characteristic functions.} By (S1) of Definition~\ref{4.31:def-radii-choice},
$w_s+16r_s\le\min\{a_1,\varepsilon_s^{-}/B_3\}$, and Lemma~\ref{4.31:lem-window-characteristic} gives
$\chi_k\equiv 1$ on $\mathcal{V}_s(w_s+16r_s)$ for every $K$.

\emph{The chart scheme.} Fix $V_0\in\mathcal{V}_s$ and let $W_{V_0}(B)$, $|B|<4r_s$, be the complex
background field of Lemma~\ref{4.31:lem-charts-gluing}. In each item below we define a function
$\Phi_{V_0}$ of $B$ on the chart $\{|B|<4r_s\}$ which is holomorphic there and which, at real $B$,
equals the named real function evaluated at $e^{iB}V_0$. The second property holds because, by
Lemma~\ref{4.31:lem-charts-gluing}(b), for real $B$ the field $W_{V_0}(B)$ is gauge equivalent, by a
$G$-valued gauge transformation, to $U_k(e^{iB}V_0)$; because the real functions are invariant under
such transformations (\cite[(1.19)]{B-CMP109} and \eqref{4.31:eq-window-counting-nesting-e}); and,
for the pieces $\mathbf{E}^{(j)},\mathbf{R}^{(j)}$, by the clause of
Lemma~\ref{4.31:lem-background-membership} on real $B$. Lemma~\ref{4.31:lem-charts-gluing}(c) then
extends the function holomorphically to the union $\mathfrak{T}_s$ of the charts, which contains
$\mathfrak{D}_s$, and the extension obeys every supremum bound that holds on all the charts. It
therefore suffices, for (a), (b), (d) and the first half of (c), to bound $\Phi_{V_0}$ on the chart
uniformly in $V_0$.

\emph{(a)} Put $\Phi_{V_0}(B):=A(W_{V_0}(B))$, a polynomial in the entries of $W$ and $W^{-1}$,
which at real $B$ equals $\mathcal{A}_{K,s}(e^{iB}V_0)$ by the chart scheme; here
$1-\operatorname{Re}\operatorname{tr}W(\partial p)$ stands for its holomorphic extension
$1-\tfrac12\operatorname{tr}(W(\partial p)+W(\partial p)^{-1})$. Let
$\delta:=|W(\partial p)-\mathbf{1}|$. By Lemma~\ref{4.31:lem-charts-gluing}(a),
\[
\delta\le 6(B_3w_s+C_{19}\rho_s)\eta^2,
\]
and this is at most $\tfrac12$, since $\eta\le1$, $M'\ge1$, and
$8(B_3w_s+C_{19}\rho_s)\le h_s\le 1/(32M')$ by the definition of $h_s$ and (S3) of
Definition~\ref{4.31:def-radii-choice}.
Hence $|W(\partial p)^{-1}|\le 2$ by the Neumann series. Since
$W+W^{-1}-2\cdot\mathbf{1}=W^{-1}(W-\mathbf{1})^2$ and $\operatorname{tr}\mathbf{1}=1$,
\[
|1-\operatorname{Re}\operatorname{tr}W(\partial p)|
=\bigl|\tfrac12\operatorname{tr}\bigl(W(\partial p)^{-1}(W(\partial p)-\mathbf{1})^2\bigr)\bigr|
\le\delta^2 ,
\]
using $|\operatorname{tr}Y|\le|Y|$ for the normalised trace, for every $N$. In $d=4$ there are six
plaquettes per site, hence $6\eta^{-4}|\mathbb{T}_s|$ plaquettes, and
\[
|\Phi_{V_0}|\le 6\eta^{-4}|\mathbb{T}_s|\cdot 36(B_3w_s+C_{19}\rho_s)^2\eta^4
=216(B_3w_s+C_{19}\rho_s)^2|\mathbb{T}_s|.
\]
Since $\rho_s=2B_3w_s+8B_0r_s$, we have
\begin{align*}
B_3w_s+C_{19}\rho_s&=(B_3+2C_{19}B_3)w_s+8C_{19}B_0r_s\le\tfrac16C_m(w_s+r_s),\\
216(B_3w_s+C_{19}\rho_s)^2&=6\bigl(6(B_3w_s+C_{19}\rho_s)\bigr)^2\le 6C_m^2(w_s+r_s)^2,
\end{align*}
which is (a) with $C_{\mathcal{A}}=6C_m^2$.

\emph{(b)} Lemma~\ref{4.31:lem-background-membership} shows that
$B\mapsto\mathbf{E}^{(j)}(X,(W,J^{\xi}(W)),z)$ and $B\mapsto\mathbf{R}^{(j)}(X,(W,J^{\xi}(W)))$ are
holomorphic on the chart and bounded by $b\,e^{-\kappa d_j(X)}$, with $b=E_0$, respectively
$b=x_j^{-\kappa_0/2}$. By Lemma~\ref{4.31:lem-cutoff-free-windows}, $x_j=x_{K-(s+t)}\in I_{s+t}$, so
$x_j\ge x_{-}^{(s+t)}$ and $x_j^{-\kappa_0/2}\le(x_{-}^{(s+t)})^{-\kappa_0/2}$. This gives (b).

\emph{(d)} Let $\Phi^{\mathsf{R}}_{V_0}(B)$ be the sum over $X$ of
$\mathbf{R}^{(j)}(X,(W,J^{\xi}(W)))-\mathbf{R}^{(j)}(X,(\mathbf{1},0))$, which at real $B$ equals
$\mathsf{R}_{K,s,t}(e^{iB}V_0)$. Every $X$ contains a cube of the partition $\pi_j$ into scale-$j$
cubes, and $\#\pi_j=M^{-4}L^{4t}|\mathbb{T}_s|$. Bounding each summand by
Lemma~\ref{4.31:lem-complex-marginal-bound} with $b=x_j^{-\kappa_0/2}$, attaching each $X$ to every
cube of $\pi_j$ it contains, and using \eqref{4.31:eq-window-counting-nesting-b}, we get on the whole
chart
\begin{align*}
|\Phi^{\mathsf{R}}_{V_0}|
&\le\sum_{\square\in\pi_j}\sum_{X\supset\square}(C_5'M^2h_s^2x_j+48e^{14})L^{-4t}
x_j^{-\kappa_0/2}e^{-(\kappa-1)d_j(X)}\\
&\le(C_5'M^2h_s^2x_j+48e^{14})x_j^{-\kappa_0/2}K_0M^{-4}L^{-4t}\cdot L^{4t}|\mathbb{T}_s|\\
&\le C_Rx_j^{1-\kappa_0/2}|\mathbb{T}_s|\le C_R(x_{-}^{(s+t)})^{1-\kappa_0/2}|\mathbb{T}_s|,
\end{align*}
where the third inequality uses $x_j\ge1$, so that
$C_5'M^2h_s^2x_j+48e^{14}\le(C_5'M^2h_s^2+48e^{14})x_j$, and the last uses that the exponent
$1-\kappa_0/2$ is negative and $x_j\ge x_{-}^{(s+t)}$. The powers $L^{4t}$ and $L^{-4t}$ cancel
exactly: the first is the number of scale-$j$ cubes in $\mathbb{T}_s$, the second is the factor in
the bound of Lemma~\ref{4.31:lem-complex-marginal-bound}. This is (d).

\emph{(c)} Let $\Phi^{\mathsf{E}}_{V_0}(B)$ be the sum over $z\in T^{(j)}_1$ and $X\ni z$ of
$\mathbf{E}^{(j)}(X,(W,J^{\xi}(W)),z)-\mathbf{E}^{(j)}(X,(\mathbf{1},0),z)$, minus
$\beta_j(g_{j-1})\Phi_{V_0}(B)$ with $\Phi_{V_0}$ from (a); at real $B$ it equals
$\mathsf{E}_{K,s,t}(e^{iB}V_0)$. Now $|T^{(j)}_1|=L^{4t}|\mathbb{T}_s|$, and each $X\ni z$ contains the
cube of $\pi_j$ through $z$. Lemma~\ref{4.31:lem-complex-marginal-bound} with $b=E_0$,
\eqref{4.31:eq-window-counting-nesting-b}, item (a) and the bound $|\beta_j(g_{j-1})|\le C_\beta$
of \eqref{4.31:eq-standing-inputs-a} give on the whole chart
\begin{align*}
|\Phi^{\mathsf{E}}_{V_0}|
&\le\bigl[(C_5'M^2h_s^2x_j+48e^{14})E_0K_0+C_\beta C_{\mathcal{A}}(w_s+r_s)^2\bigr]|\mathbb{T}_s|
\le\mathsf{M}_t,
\end{align*}
again with an exact cancellation of $L^{4t}$ against $L^{-4t}$; the last inequality uses
$1\le x_j\le x_{+}^{(s+t)}$. For real $B$, $e^{iB}V_0\in\mathcal{V}_s(w_s+16r_s)$ by
Lemma~\ref{4.31:lem-charts-gluing}(b), and $\chi_k=1$ there by
Lemma~\ref{4.31:lem-window-characteristic}; hence \eqref{4.31:eq-standing-inputs-c} gives
$|\Phi^{\mathsf{E}}_{V_0}(B)|\le\mathsf{m}_t$.

Now let $V=e^{-Y}V_0\in\mathfrak{D}_s$ with $|Y|<r_s$, and put
$\psi(\zeta):=\Phi^{\mathsf{E}}_{V_0}(4\zeta Y)$. Since $|4\zeta Y|<4r_s$ for $|\zeta|<1$, $\psi$ is
holomorphic on the unit disc and $|\psi|\le\mathsf{M}_t$ there; for real $\zeta$ the argument
$4\zeta Y$ is $\mathfrak{g}$-valued, so $|\psi|\le\mathsf{m}_t$ on $\,]{-1},1[\,$. Moreover
$\psi(i/4)=\Phi^{\mathsf{E}}_{V_0}(iY)=\mathsf{E}_{K,s,t}(V)$ by
Lemma~\ref{4.31:lem-charts-gluing}(c), since $e^{i(iY)}V_0=V$. Lemma~\ref{4.31:lem-two-constants} at
$|\zeta|=\tfrac14$ gives
\[
|\mathsf{E}_{K,s,t}(V)|\le(\mathsf{m}_t\mathsf{M}_t)^{1/2}
=C_E(x_{+}^{(s+t)})^{1/2}L^{-\beta t/2}|\mathbb{T}_s|,
\]
which is (c). The decay $L^{-\beta t/2}$ is the real decay of \eqref{4.31:eq-standing-inputs-c},
interpolated; no other decay enters. Any $\beta>0$ suffices for what follows, because
$\mathsf{M}_t$ contains no positive power of $L^{t}$, by the cancellation above.

\emph{(e)} $\mathsf{S}_{K,s}$ is a finite sum of the functions treated in (a), (c) and (d), the
function of (a) entering with the factor $x_{K-s}\le x_{+}^{(s)}$
(Lemma~\ref{4.31:lem-cutoff-free-windows}); it is therefore holomorphic on $\mathfrak{D}_s$ and
$|\mathsf{S}_{K,s}|\le C_s$ with
\begin{equation}\label{4.31:eq-tube-analyticity-1}
\begin{split}
C_s:=\Bigl[&x_{+}^{(s)}C_{\mathcal{A}}(w_s+r_s)^2\\
&+\sum_{t\ge0}\Bigl(C_E(x_{+}^{(s+t)})^{1/2}L^{-\beta t/2}
+C_R(x_{-}^{(s+t)})^{1-\kappa_0/2}\Bigr)\Bigr]|\mathbb{T}_s|,
\end{split}
\end{equation}
which is finite because both series converge by Lemma~\ref{4.31:lem-cutoff-free-windows}.
\end{proof}

\begin{lemma}[Fixed-depth convergence: couplings and densities as measures]
\label{4.31:prf-couplings-densities}
Assume the hypotheses of Proposition~\ref{4.31:prop-fixed-depth-convergence}, and let
$w_s,r_s$ be its radii, $\mathcal{V}_s=\mathcal{V}_s(w_s)$, and $I_u=[x_{-}^{(u)},x_{+}^{(u)}]$
the windows of \eqref{4.31:eq-fixed-depth-convergence-a}. The proof below uses the following
statements:
\begin{itemize}
\item[(a)] Lemma~\ref{4.31:lem-cutoff-free-windows}: $x_{K-u}\in I_u$ for $K\in\mathcal{K}$
and $0\le u\le K$;
\item[(b)] the increment relation $x_{l-1}-x_l=\beta_l(g_{l-1})$, $|\beta_l|\le C_{\beta}$,
of \eqref{4.31:eq-standing-inputs-a};
\item[(c)] the stability and equal-integral statements \eqref{4.31:eq-standing-inputs-d},
taken from \cite[Theorem~1, (0.1)]{B-CMP122b}, with
$E_{\pm}^{(s)}=\sup_{x\in I_s}E_{\pm}(x)<\infty$, finite because, in the reading of
\eqref{4.31:eq-standing-inputs-d}, $x\mapsto E_{\pm}(x)$ is locally bounded on
$[\gamma^{-2},\infty)$; together with the non-negativity
statement \eqref{4.31:eq-window-counting-nesting-d};
\item[(d)] Lemma~\ref{4.31:lem-window-characteristic};
\item[(e)] on $\mathcal{V}_0$ one has $\mathcal{A}_{K,0}\le C_{\mathcal{A}}(w_0+r_0)^2|\mathbb{T}_0|$,
with the constant $C_{\mathcal{A}}$ of \eqref{4.31:eq-tube-analyticity-a}.
\end{itemize}
Put $\nu^{B}_{K,s}=Z_K^{-1}\rho_{K,s}\,dV$ and, with $w_0,r_0$ the radii $w_s,r_s$ of
Proposition~\ref{4.31:prop-fixed-depth-convergence} at $s=0$,
\begin{equation*}
z_{*}=\exp\Bigl[-\bigl(E_{-}^{(0)}+x_{+}^{(0)}C_{\mathcal{A}}(w_0+r_0)^2\bigr)
|\mathbb{T}_0|\Bigr]\int_{\mathcal{V}_0}dV .
\end{equation*}
Then $z_{*}>0$, and $z_{*}$ depends on neither $K$ nor $s$; moreover there are infinite sets
$\mathcal{K}_2\subset\mathcal{K}_1\subset\mathcal{K}$ such that the following hold, all limits
being taken as $K\to\infty$ in the set indicated.
\begin{enumerate}
\item[(D1)] Along $\mathcal{K}_1$, for every $s\ge0$, $x_{K-s}\to x_s^{\infty}\in I_s$ and
\begin{equation*}
\beta_{K-s}(g_{K-s-1})=x_{K-s-1}-x_{K-s}\longrightarrow\beta_s^{\infty}
:=x_{s+1}^{\infty}-x_s^{\infty},\qquad|\beta_s^{\infty}|\le C_{\beta}.
\end{equation*}
Moreover, if in addition for every $K\in\mathcal{K}$ the all-pairs inequalities of
\eqref{4.31:eq-standing-inputs-a} hold, with the $K$-free constants $\lambda$, $B$, $c_2$
fixed there,
\begin{equation*}
(j'-i)\lambda-2c_2\le x_i-x_{j'}\le(j'-i)B+2c_2\qquad(0\le i<j'\le K),
\end{equation*}
then for all $0\le s<s'$
\begin{equation*}
(s'-s)\lambda-2c_2\le x_{s'}^{\infty}-x_s^{\infty}\le(s'-s)B+2c_2 .
\end{equation*}
\item[(D2)] For all $K\in\mathcal{K}$ and $0\le s\le K$, $\rho_{K,s}\,dV$ is a positive Borel
measure on $\mathcal{G}_s$ of mass $Z_K\in[z_{*},e^{E_{+}^{(0)}|\mathbb{T}_0|}]$. Along
$\mathcal{K}_2$, $Z_K\to Z^{\infty}$, and for every $s\ge0$ the measures $\nu^{B}_{K,s}$
converge weakly (against continuous $F$ on $\mathcal{G}_s$) to a probability measure
$\nu_s^{\infty}$.
\item[(D2$^{+}$)] If, in addition, for some $s$ and every $K\in\mathcal{K}$ with $K\ge s$ the
measure $\rho_{K,s}\,dV$ has a density $\rho_{K,s}$ with
$\rho_{K,s}\le e^{E_{+}^{(s)}|\mathbb{T}_s|}$ $dV$-almost everywhere, then along a further
infinite subset of $\mathcal{K}_2$ one has $\rho_{K,s}\to\rho_s^{\infty}$ weak-$*$ in
$L^{\infty}(\mathcal{G}_s,dV)$, $0\le\rho_s^{\infty}\le e^{E_{+}^{(s)}|\mathbb{T}_s|}$,
$\nu_s^{\infty}=\rho_s^{\infty}\,dV/Z^{\infty}$, in particular $\nu_s^{\infty}\ll dV$, and
$\nu^{B}_{K,s}(F)\to\nu_s^{\infty}(F)$ for every bounded Borel $F$.
\end{enumerate}
The additional hypotheses of the conditional part of (D1) and of (D2$^{+}$) are not among the
statements (a)--(e): input (c) uses the upper member of \eqref{4.31:eq-standing-inputs-d} only
integrated at level $K$, and the lower member only as an inequality of measures. These additional
hypotheses are not proved here; the conditional part of (D1) and (D2$^{+}$) are proved only as
implications from them, and nothing in what follows uses either.
\end{lemma}

\begin{proof}
\emph{Clause (D1).} For $K\in\mathcal{K}$ define the sequence $y^K=(y^K_u)_{u\ge0}$ by
$y^K_u=x_{K-u}$ for $0\le u\le K$ and $y^K_u=x_{-}^{(u)}$ for $u>K$. By (a),
$y^K\in\prod_{u\ge0}I_u$. Each $I_u$ is a compact interval, so the product is compact by
Tychonoff's theorem, and it is metrisable as a countable product of metric spaces; hence it is
sequentially compact. There is therefore an infinite $\mathcal{K}_1\subset\mathcal{K}$ along
which $y^K$ converges in the product topology, that is, coordinatewise: $x_{K-u}\to x_u^{\infty}$
for every $u$ (for fixed $u$ only finitely many $K$ have $K<u$), and $x_u^{\infty}\in I_u$
because $I_u$ is closed. For fixed $s$ and $K\in\mathcal{K}_1$ with $K>s$, relation (b) at
$l=K-s\ge1$ gives
\begin{equation*}
\beta_{K-s}(g_{K-s-1})=x_{K-s-1}-x_{K-s}=y^K_{s+1}-y^K_s\longrightarrow
x_{s+1}^{\infty}-x_s^{\infty}=\beta_s^{\infty},
\end{equation*}
and $|\beta_s^{\infty}|\le C_{\beta}$ since every $\beta_{K-s}(g_{K-s-1})$ lies in the closed
interval $[-C_{\beta},C_{\beta}]$.

\emph{Clause (D1): the conditional part.} For $0\le s<s'$ and $K\ge s'$, the assumed
inequalities at $(i,j')=(K-s',K-s)$ read
\begin{equation*}
(s'-s)\lambda-2c_2\le x_{K-s'}-x_{K-s}\le(s'-s)B+2c_2 .
\end{equation*}
These are closed conditions on $(x_{K-s'},x_{K-s})$, so they pass to the limit along
$\mathcal{K}_1$. They are not used in what follows.

\emph{Clause (D2): positivity, masses and the upper bound.} At every level $k=K-s$ the lower
member of (c), read as an inequality of measures, and the non-negativity statement in (c) give
$\rho_{K-s}\,dV\ge0$, so $\rho_{K,s}\,dV$ is a positive Borel measure on $\mathcal{G}_s$. At
$k=K$, where $\rho_{K,0}\,dV=\rho_K\,dV$, the lower member reads
\begin{equation*}
\rho_{K,0}\,dV\ \ge\ \chi_K\,e^{-x_K\mathcal{A}_{K,0}-E_{-}(x_K)|\mathbb{T}_0|}\,dV\ \ge\ 0 .
\end{equation*}
By the equal-integral
statement in (c), $\int\rho_{K,s}\,dV=Z_K$ for every $s$. By the integrated upper member of (c)
at depth $0$, and since $x_K\in I_0$ by (a) at $u=0$,
\begin{equation*}
Z_K\le e^{E_{+}(x_K)|\mathbb{T}_0|}\le e^{E_{+}^{(0)}|\mathbb{T}_0|}.
\end{equation*}

\emph{Clause (D2): the lower bound.} By Lemma~\ref{4.31:lem-window-characteristic},
$\chi_K\equiv1$ on $\mathcal{V}_0(w_0+16r_0)$, which contains $\mathcal{V}_0=\mathcal{V}_0(w_0)$
because the windows \eqref{4.31:eq-windows-tubes-a} increase with their radius; everywhere else
$\chi_K\ge0$. Integrating the displayed inequality of measures and discarding the complement of
$\mathcal{V}_0$, on which the integrand is non-negative, we get, with $E_{-}(x_K)\le E_{-}^{(0)}$
because $x_K\in I_0$,
\begin{equation*}
Z_K=\int\rho_{K,0}\,dV\ \ge\ e^{-E_{-}^{(0)}|\mathbb{T}_0|}
\int_{\mathcal{V}_0}e^{-x_K\mathcal{A}_{K,0}}\,dV .
\end{equation*}
On $\mathcal{V}_0$, by (e) and $0<x_K\le x_{+}^{(0)}$,
\begin{equation*}
e^{-x_K\mathcal{A}_{K,0}}\ \ge\ e^{-x_{+}^{(0)}C_{\mathcal{A}}(w_0+r_0)^2|\mathbb{T}_0|},
\qquad\text{and therefore}\qquad Z_K\ \ge\ z_{*}.
\end{equation*}
The set $\mathcal{V}_0$ is open, non-empty and independent of $K$, and
Haar measure charges every non-empty open set, so $\int_{\mathcal{V}_0}dV>0$; since moreover
$E_{-}^{(0)}$ and $C_{\mathcal{A}}$ are finite, $z_{*}>0$, and none
of its factors depends on $K$ or $s$.

\emph{Clause (D2): convergence.} The group $\mathcal{G}_s$ is compact and metrisable, so
$C(\mathcal{G}_s)$ is separable. By the Riesz representation theorem the probability measures on
$\mathcal{G}_s$ are the positive functionals of norm one on $C(\mathcal{G}_s)$, a weak-$*$
closed subset of the unit ball of $C(\mathcal{G}_s)^{*}$; by the Banach--Alaoglu theorem and
separability that ball is weak-$*$ sequentially compact. Hence the probability measures on
$\mathcal{G}_s$ are weakly sequentially compact. The numbers $Z_K$ lie in the compact interval
$[z_{*},e^{E_{+}^{(0)}|\mathbb{T}_0|}]$. Extracting successively, first for $Z_K$ and then for
$s=0,1,2,\dots$, a further infinite subset along which $\nu^{B}_{K,s}$ converges weakly, and
taking the diagonal sequence (its $n$-th element is the $n$-th element of the $n$-th extraction,
so that it is eventually contained in each extraction), we obtain an infinite
$\mathcal{K}_2\subset\mathcal{K}_1$ along which $Z_K\to Z^{\infty}\in
[z_{*},e^{E_{+}^{(0)}|\mathbb{T}_0|}]$ and, for every $s$, $\nu^{B}_{K,s}\to\nu_s^{\infty}$
weakly, with $\nu_s^{\infty}$ a probability measure.

\emph{Clause (D2$^{+}$).} Since $\mathcal{G}_s$ is compact metrisable, $L^{1}(\mathcal{G}_s,dV)$
is separable, so the balls of $L^{\infty}(\mathcal{G}_s,dV)=L^{1}(\mathcal{G}_s,dV)^{*}$ are
weak-$*$ sequentially compact. With $c=e^{E_{+}^{(s)}|\mathbb{T}_s|}$, the set
$\{0\le\rho\le c\}$ is weak-$*$ closed, being the set of $\rho$ with $\int\rho\varphi\,dV\ge0$
and $\int(c-\rho)\varphi\,dV\ge0$ for all non-negative $\varphi\in L^{1}$. One more extraction
from $\mathcal{K}_2$ (diagonal over the depths $s$ for which the additional hypothesis is made)
gives $\rho_{K,s}\to\rho_s^{\infty}$ weak-$*$ with $0\le\rho_s^{\infty}\le c$. For bounded Borel
$F$ we have $F\in L^{1}(\mathcal{G}_s,dV)$, and since $Z_K\to Z^{\infty}\ge z_{*}>0$,
\begin{equation*}
\nu^{B}_{K,s}(F)=\frac{1}{Z_K}\int F\rho_{K,s}\,dV\longrightarrow
\frac{1}{Z^{\infty}}\int F\rho_s^{\infty}\,dV .
\end{equation*}
For continuous $F$ the left side also converges to $\nu_s^{\infty}(F)$ by (D2); the two limits
agree on $C(\mathcal{G}_s)$, so by the uniqueness part of the Riesz representation theorem
$\nu_s^{\infty}=\rho_s^{\infty}\,dV/Z^{\infty}$, and the displayed convergence is
$\nu^{B}_{K,s}(F)\to\nu_s^{\infty}(F)$ for every bounded Borel $F$.
\end{proof}

\begin{lemma}[Fixed-depth convergence: normal families and the limit lower bound]
\label{4.31:prf-normal-families}
Assume the hypotheses of Proposition~\ref{4.31:prop-fixed-depth-convergence}; of these the proof
below uses the following:
\begin{itemize}
\item the stability bounds and the equality of integrals \eqref{4.31:eq-standing-inputs-d};
\item the positivity of the terms of the effective density and the bound on the term of the empty
large-field history, \eqref{4.31:eq-window-counting-nesting-d};
\item the covariance identity \eqref{4.31:eq-window-counting-nesting-e};
\item the identity principle in the form \eqref{4.31:eq-complex-background-c}, and the classical
facts collected in \eqref{4.31:eq-window-counting-nesting-g};
\item Lemma~\ref{4.31:lem-window-characteristic}.
\end{itemize}
Assume further that there is an infinite set $\mathcal{K}_2\subset\mathcal{K}$ with the following
properties. Here clauses (D1)--(D5) are those of
Proposition~\ref{4.31:prop-fixed-depth-convergence}, displayed in
\eqref{4.31:eq-fixed-depth-convergence-b} with the constants of
\eqref{4.31:eq-tube-analyticity-a}.
\begin{itemize}
\item For every $u\ge0$, $x_{K-u}\to x_u^{\infty}\in I_u$ and
$\beta_{K-u}(g_{K-u-1})\to\beta_u^{\infty}$ as $K\to\infty$ in $\mathcal{K}_2$ (these are the
convergence statements of clause (D1)), and clause (D2) holds along $\mathcal{K}_2$ for every
$s\ge0$; both are established in Lemma~\ref{4.31:prf-couplings-densities}.
\item For every $s\ge0$, clause (D3) holds for all $K\in\mathcal{K}$ with $K>s$, as established in
Lemma~\ref{4.31:prf-tube-analyticity}.
\end{itemize}
Then there is an infinite set $\mathcal{K}'\subset\mathcal{K}_2$ such that, as $K\to\infty$ in
$\mathcal{K}'$, the convergence statements of (D1) and (D2) assumed above continue to hold and
clauses (D4) and (D5) hold, for every $s\ge 0$ simultaneously.
\end{lemma}

\begin{proof}
Throughout, $s\ge0$ is the depth, $k=K-s$, and $\mathcal{K}_2\subset\mathcal{K}$ is the infinite
set of the hypotheses, reached after the extractions made in the proofs of clauses (D1)--(D3)
(Lemmas~\ref{4.31:prf-tube-analyticity} and~\ref{4.31:prf-couplings-densities}). We write
$\mathcal{V}_s=\mathcal{V}_s(w_s)$ and $\mathfrak{D}_s=\mathfrak{D}_s(w_s,r_s)$.

\emph{Step 1: clause} (D4). Fix $s$. Consider the families, indexed by $K\in\mathcal{K}_2$,
\begin{equation*}
\{\mathcal{A}_{K,s}\},\quad \{e^{K}_{t,X,z}\},\quad \{r^{K}_{t,X}\},\quad
\{\mathsf{E}_{K,s,t}\},\quad \{\mathsf{R}_{K,s,t}\},\quad \{\mathsf{S}_{K,s}\},
\end{equation*}
one family for each admissible value of the indices $t$, $X$, $z$; there are countably many of them.
By clause (D3), each of these functions extends from $\mathcal{V}_s$ to a holomorphic function on
the complex manifold $\mathfrak{D}_s$, and the bounds (a)--(e) of (D3) bound each family on
$\mathfrak{D}_s$ uniformly in $K$, the right-hand sides of (a)--(e) being independent of $K$. By
Montel's theorem on complex manifolds (one of the classical facts of
\eqref{4.31:eq-window-counting-nesting-g}), each family is therefore normal: every infinite subset
of $\mathcal{K}_2$ contains a further infinite subset along which the family converges uniformly on
compact subsets of $\mathfrak{D}_s$. The set of all pairs consisting of a depth $s$ and one of the
families above is countable. A diagonal extraction over this set, performed after the extractions
of the earlier steps and after the additional extraction made in Step~2 below, yields an infinite
set $\mathcal{K}'\subset\mathcal{K}_2$ along which every one of these families converges uniformly
on compact subsets of the corresponding $\mathfrak{D}_s$. By the Weierstrass theorem (again from
\eqref{4.31:eq-window-counting-nesting-g}), the limits are holomorphic on $\mathfrak{D}_s$ and all
derivatives converge uniformly on compact subsets as well. Since $\mathcal{K}'\subset\mathcal{K}_2$,
the convergence statements of (D1) and (D2) assumed above hold along $\mathcal{K}'$.

Gauge invariance. Let $v$ be a $G$-valued gauge transformation of $\mathbb{T}_s$. On $\mathcal{V}_s$
each of the functions above is invariant under $V\mapsto V^{v}$: by
\eqref{4.31:eq-window-counting-nesting-e} (\cite[(181)]{B-CMP102b}) one has
$U_k(V^{v})=U_k(V)^{\tilde v}$, with $\tilde v$ the
natural extension of $v$, and the invariance follows from this together with
\cite[(1.19)]{B-CMP109}. The map $V\mapsto V^{v}$ is holomorphic and preserves $\mathfrak{D}_s$.
Hence, for each function $f$ of the list, $V\mapsto f(V^{v})-f(V)$ is holomorphic on
$\mathfrak{D}_s$ and vanishes on $\mathcal{V}_s$. Every connected component of $\mathfrak{D}_s$
retracts onto its real slice, the retraction sending the complex direction $Y$ of the tube to $0$,
and this real slice is part of $\mathcal{V}_s$, which is a maximal totally real submanifold. By the
identity principle \eqref{4.31:eq-complex-background-c}, $f(V^{v})=f(V)$ on every component, that
is, on all of $\mathfrak{D}_s$.

Properties of the limits. The bounds (a)--(e), the vanishing at $\mathbf{1}$ of
$\mathcal{A}_{K,s}$, $\mathsf{E}_{K,s,t}$, $\mathsf{R}_{K,s,t}$, $\mathsf{S}_{K,s}$, and the gauge
invariance just proved are each preserved under pointwise limits, their constants being independent
of $K$. The limits $\mathcal{A}_s^{\infty}$, $e^{\infty}_{t,X,z}$, $r^{\infty}_{t,X}$,
$\mathsf{E}^{\infty}_{s,t}$, $\mathsf{R}^{\infty}_{s,t}$, $\mathsf{S}_s^{\infty}$ therefore have all
these properties.

The formulas for $\mathsf{E}^{\infty}_{s,t}$ and $\mathsf{R}^{\infty}_{s,t}$. At finite $K$ the
functions $\mathsf{E}_{K,s,t}$ and $\mathsf{R}_{K,s,t}$ are given by Ba{\l}aban's formulas
\cite[(0.22)--(0.23)]{B-CMP109}, with $e^{K}_{t,X,z}$, $r^{K}_{t,X}$, $\mathcal{A}_{K,s}$ in place of
the limits and with the coupling increment at the level $K-s-t$ as the coefficient of
$\mathcal{A}_{K,s}$. At fixed $t$ the sums over $(X,z)$ are finite, because $\mathbb{T}_s$ is
finite. The coefficient converges to $\beta^{\infty}_{s+t}$ by clause (D1) read at depth $s+t$, and
each of the finitely many remaining terms converges by what was proved above. Taking the limit term
by term gives the displayed formulas for $\mathsf{E}^{\infty}_{s,t}$ and $\mathsf{R}^{\infty}_{s,t}$
in \eqref{4.31:eq-fixed-depth-convergence-b}.

The $t$-series. At finite $K$,
\begin{equation*}
\mathsf{S}_{K,s}=-x_k\mathcal{A}_{K,s}+\sum_{t\ge0}\bigl(\mathsf{E}_{K,s,t}+\mathsf{R}_{K,s,t}\bigr).
\end{equation*}
The terms with $t\ge K-s$ are zero. By (c) and (d) of
(D3), on $\mathfrak{D}_s$ and for every $K$,
\begin{equation*}
|\mathsf{E}_{K,s,t}|+|\mathsf{R}_{K,s,t}|
\le \Bigl(C_E\bigl(x_{+}^{(s+t)}\bigr)^{1/2}L^{-\beta t/2}
+C_R\bigl(x_{-}^{(s+t)}\bigr)^{1-\kappa_0/2}\Bigr)|\mathbb{T}_s| ,
\end{equation*}
and the right-hand side is independent of $K$ and summable in $t$. By Tannery's theorem (from
\eqref{4.31:eq-window-counting-nesting-g}), the limit $K\to\infty$ in $\mathcal{K}'$ may be taken
under the sum over $t$, uniformly on compact subsets of $\mathfrak{D}_s$; with $x_k\to x_s^{\infty}$
(clause (D1)) this gives the displayed formula for $\mathsf{S}_s^{\infty}$. The same majorant shows
that the limit series converges uniformly on $\mathfrak{D}_s$. This proves (D4).

\emph{Step 2: clause} (D5). Let $F\ge0$ be continuous with compact support in $\mathcal{V}_s$.
Since $\nu^{B}_{K,s}$ is the normalised measure, $Z_K\nu^{B}_{K,s}=\rho_{K,s}\,dV$. The lower member
of \eqref{4.31:eq-standing-inputs-d} is an inequality of measures; read at level $k$ on the depth-$s$
lattice, where the action of the block field is $\mathcal{A}_{K,s}$ and the volume is
$|\mathbb{T}_s|$, it is
\begin{equation*}
\rho_{K,s}\,dV\ \ge\ \chi_k\,e^{-x_k\mathcal{A}_{K,s}-E_{-}(x_k)|\mathbb{T}_s|}\,dV ,
\end{equation*}
valid because every term of the effective density is non-negative, by
\eqref{4.31:eq-window-counting-nesting-d}. By Lemma~\ref{4.31:lem-window-characteristic},
$\chi_k\equiv1$ on $\mathcal{V}_s(w_s+16r_s)\supset\mathcal{V}_s$. Since $x_k\in I_s$ and
$E_{-}^{(s)}$ is the supremum of $E_{-}$ over $I_s$, we have $E_{-}(x_k)\le E_{-}^{(s)}$. Hence
\begin{equation}\label{4.31:eq-normal-families-1}
\begin{split}
Z_K\nu^{B}_{K,s}(F)
&\ge \int F\,\exp\bigl[-x_k\mathcal{A}_{K,s}-E_{-}(x_k)|\mathbb{T}_s|\bigr]\,dV\\
&\ge \int F\,\exp\bigl[-x_k\mathcal{A}_{K,s}-E_{-}^{(s)}|\mathbb{T}_s|\bigr]\,dV .
\end{split}
\end{equation}
As $K\to\infty$ in $\mathcal{K}'$, the left side of \eqref{4.31:eq-normal-families-1} tends to
$Z^{\infty}\nu_s^{\infty}(F)$ by clause (D2). On the right, $x_k\to x_s^{\infty}$ by clause (D1), and
$\mathcal{A}_{K,s}\to\mathcal{A}_s^{\infty}$ uniformly on the compact set $\operatorname{supp}F$ by
(D4); the right side therefore tends to
\begin{equation*}
\int F\,\exp\bigl[-x_s^{\infty}\mathcal{A}_s^{\infty}-E_{-}^{(s)}|\mathbb{T}_s|\bigr]\,dV .
\end{equation*}
Functions $F$ of this kind determine a measure on the open set $\mathcal{V}_s$ (Riesz, from
\eqref{4.31:eq-window-counting-nesting-g}), so the limiting inequality is the inequality of measures
\begin{equation*}
Z^{\infty}\nu_s^{\infty}\ \ge\ \exp\bigl[-x_s^{\infty}\mathcal{A}_s^{\infty}
-E_{-}^{(s)}|\mathbb{T}_s|\bigr]\,dV
\end{equation*}
on $\mathcal{V}_s$ asserted in (D5).

The additional statement of (D5). By \eqref{4.31:eq-window-counting-nesting-d} at $V=\mathbf{1}$,
\begin{equation*}
|E^{\emptyset}_{K,s}|\le E_{-}(x_k)|\mathbb{T}_s|\le E_{-}^{(s)}|\mathbb{T}_s| ,
\end{equation*}
so $e^{-E^{\emptyset}_{K,s}}$ lies in the compact interval
$[e^{-E_{-}^{(s)}|\mathbb{T}_s|},e^{E_{-}^{(s)}|\mathbb{T}_s|}]$ for every $K$. Extracting once more,
for all $s$ at once by a diagonal argument over $s$, and doing so before the final diagonal
extraction of Step~1, we obtain $e^{-E^{\emptyset}_{K,s}}\to c_s^{\infty}$ with $c_s^{\infty}$ in
that interval. Next, since all the other terms of the effective density are non-negative, by
\eqref{4.31:eq-window-counting-nesting-d}, keeping only the term of the empty large-field history
gives
\begin{equation*}
\rho_{K,s}\,dV\ \ge\ \chi_k\,e^{\mathsf{S}_{K,s}-E^{\emptyset}_{K,s}}\,dV .
\end{equation*}
The argument of the preceding paragraph, with $\chi_k\equiv1$ on $\mathcal{V}_s$ by
Lemma~\ref{4.31:lem-window-characteristic}, with $\mathsf{S}_{K,s}\to\mathsf{S}_s^{\infty}$
uniformly on compact subsets by (D4), and with $e^{-E^{\emptyset}_{K,s}}\to c_s^{\infty}$, gives
$Z^{\infty}\nu_s^{\infty}\ge c_s^{\infty}e^{\mathsf{S}_s^{\infty}}\,dV$ on $\mathcal{V}_s$. This
proves (D5).
\end{proof}

\begin{proposition}[Positivity dichotomy of the large-field operations. Proof outstanding]
\label{4.32:prop-positivity-dichotomy}
The following holds inside the one induction on the level $k$ that carries clauses (0) and (i)
of Theorem~0 (clause (0) is Theorem~\ref{4.1:thm-clause-zero}).

\emph{Notation.} Let $k\le K$ and let $X$ be a component on which Ba{\l}aban's operation
$\mathbf{T}'_k(X)$ of \cite[(1.71)]{B-CMP122b} acts, inside a term of \cite[(2.18)]{B-CMP119},
equivalently of \cite[(1.72)]{B-CMP122b}, taken with its label, that is, with its index in that
expansion.

\begin{itemize}
\item A \emph{history} $\omega$ of $X$ is one choice of everything summed over in
\cite[(1.71)]{B-CMP122b}, and $\mathcal{A}(X)$ is the finite set of histories.
\item $\Xi_\omega$ is the integration space of $\omega$ and $\mu_\omega$ its product of Haar and
Lebesgue measures.
\item $S_\omega\subset\Xi_\omega$ is the set on which every characteristic function of $\omega$
equals $1$.
\item $e^{\Phi_\omega}$ is the product of all factors of $\omega$ other than its characteristic
functions, $\delta$-functions and integration measures (the weight factors).
\item We put $\mathfrak{m}_\omega=\mu_\omega(S_\omega)$ and
$\mathfrak{m}(X)=\sum_{\omega\in\mathcal{A}(X)}\mathfrak{m}_\omega$; the component $X$, and the
label of the term carrying $\mathbf{T}'_k(X)$, are called \emph{non-null} if $\mathfrak{m}(X)>0$.
\item $\mathfrak{D}_k(X)$ is the complex space $U^c_k(X^{\sim},\alpha_{0,k},\alpha_{1,k})$ of
pairs $(\mathbf{U},\mathbf{J})$, with the radii $\alpha_{0,k},\alpha_{1,k}$ as in
Lemma~\ref{4.4:lem-twist-stability}, and $\operatorname{Reg}_k(X)$ is its set of real pairs.
\end{itemize}

Then:
\begin{enumerate}
\item[(a)] \emph{Reality.} For every history $\omega\in\mathcal{A}(X)$, every factor of
\cite[(1.71)]{B-CMP122b} that is not a characteristic function, a $\delta$-function or a measure
is a strictly positive real number at every point of $S_\omega\times\operatorname{Reg}_k(X)$, in
particular at the real pairs $(U,0)$ and at the diagonal pair $(U_k,J(U_k))$, where $J(U_k)$ is
the current determined by $U_k$. Every expression in \cite[(1.71)]{B-CMP122b} is real on
the real integration variables. This clause is carried in the induction: it is assumed at every
level $j<k$ and is part of what is established at level $k$.
\item[(b)] \emph{Positivity.} For every real pair $(U,J)$, $\mathbf{T}'_k(X,(U,J))$ is integration
against a non-negative measure $\nu_{X,(U,J)}$. Thus $F\ge0$ implies
$\mathbf{T}'_k(X,(U,J))F\ge0$, and
\begin{equation*}
|\mathbf{T}'_k(X,(U,J))F|\le\sup|F|\cdot\mathbf{T}'_k(X,(U,J))1 .
\end{equation*}
Every term of \cite[(2.18)]{B-CMP119}, equivalently of \cite[(1.72)]{B-CMP122b}, is a non-negative
measure on real fields.
\item[(c)] \emph{Reference normalisations.} Let $X_i$ run over the components indexed by $i$
on \cite[p.~378]{B-CMP122b} (those of the first class carry the gauge-fixing $\delta$-functions
$\delta_{G_0\cap X_i}$ and the characteristic functions $\chi_k(X_i^{\sim-6})$), and let
$\Lambda_i$ be the domains of \cite{B-CMP122a} to which the block field $V_k$ is extended, with
$V_{\Lambda_i}$ the extended field and $U_{k,X_i}(V'V_{\Lambda_i})$ the configuration built from
it. Let $\chi_{k,\Lambda_i}$ be the
characteristic function of \cite{B-CMP122a} which restricts the field $V_k$ to have an extension
satisfying the regularity condition \cite[(1.75)]{B-CMP122a}. By an \emph{exterior} we mean a
configuration of the block field $V_k$ outside $X_i$. The following are defined and strictly
positive at every exterior in $\operatorname{supp}\chi_k(\Omega_k^{\sim4})\chi_{k,\Lambda_i}$,
where $\Omega_k$ is the old $k$-th domain of \cite[p.~378]{B-CMP122b}, that is, the $k$-th domain
of the sequence of domains as it stands before the step:
\begin{itemize}
\item the denominators of \cite[(0.3)]{B-CMP122a}, which are those of \cite[(1.100)]{B-CMP122a};
\item the normalisations of the reference laws \cite[(1.100)--(1.102)]{B-CMP122a}.
\end{itemize}
Hence $\mathbf{R}$ is well defined on every term it meets. Moreover \cite[(0.4)]{B-CMP122a} is
Fubini's theorem for finite positive measures, combined with the reversed Faddeev--Popov
equivalence (an equivalence, not an equality) for the $\delta$-functions $\delta_{G_0\cap X_i}$ on
the components carrying the first (small-field) function \cite[p.~378]{B-CMP122b}.
\item[(d)] \emph{Exterior-uniform dichotomy.} The following holds for every label, including labels
whose factors contain the gauge-invariant characteristic functions of $\mathbf{R}$ that are not
products: the functions $\chi_j,\chi^c_j$ of \cite[(2.17), (3.2)]{B-CMP119}, the functions
$1-\chi_{j,\Lambda}$ of \cite{B-CMP122a}, and the functions $\chi^c_j(Y^{\sim-6})$ of
\cite{B-CMP122b}.
\begin{enumerate}
\item[(d1)] $S_\omega$ is a Borel subset of $\Xi_\omega$. The triple
$(\Xi_\omega,\mu_\omega,S_\omega)$ is \emph{exterior-free}, that is, it is determined by
$(X,\omega)$, $k$, $g_0,\dots,g_k$ and $\mathfrak{p}$. It does not depend on
$(\mathbf{U},\mathbf{J})$, on $V_k$, on anything outside $X$, on $K$, or on the torus.
\item[(d2)] For $(\mathbf{U},\mathbf{J})\in\mathfrak{D}_k(X)$ and $F$ bounded Borel on
$\bigsqcup_\omega\Xi_\omega$,
\begin{equation}\label{4.32:eq-positivity-dichotomy-1}
\mathbf{T}'_k(X,(\mathbf{U},\mathbf{J}))F=(2^dd!)^{-1}\sum_{\omega\in\mathcal{A}(X)}
\int_{S_\omega}F(\omega,\xi)\,e^{\Phi_\omega(\xi;(\mathbf{U},\mathbf{J}))}\,\mu_\omega(d\xi).
\end{equation}
Here $\Phi_\omega$ is finite on $S_\omega\times\mathfrak{D}_k(X)$ and real on
$S_\omega\times\operatorname{Reg}_k(X)$. In particular, at a real pair the measure of (b) is
\begin{equation*}
\nu_{X,(U,J)}=(2^dd!)^{-1}\sum_\omega\mathbf{1}_{S_\omega}e^{\Phi_\omega}\mu_\omega .
\end{equation*}
\item[(d3)] The following are equivalent:
\begin{itemize}
\item $\mathfrak{m}(X)>0$;
\item $\mathbf{T}'_k(X,(U,J))1>0$ for some $(U,J)\in\operatorname{Reg}_k(X)$;
\item $\mathbf{T}'_k(X,(U,J))1>0$ for every $(U,J)\in\operatorname{Reg}_k(X)$, in particular at
every $(U,0)$.
\end{itemize}
The same equivalence holds history by history, with $\mathfrak{m}_\omega$ in place of
$\mathfrak{m}(X)$.
\item[(d4)] If $\mathfrak{m}(X)=0$, then $\mathbf{T}'_k(X,(\mathbf{U},\mathbf{J}))F=0$ for every
$(\mathbf{U},\mathbf{J})\in\mathfrak{D}_k(X)$ and every $F$. Consequently each of the following is
the zero function on its whole real and complex domain:
\begin{itemize}
\item each term of \cite[(1.72)]{B-CMP122b} that contains the factor $\mathbf{T}'_k(X)$;
\item each such summand of an activity \cite[(1.91)]{B-CMP122b};
\item each such product \cite[(1.103)]{B-CMP122b}.
\end{itemize}
Omitting these terms changes none of the following: the effective density $\rho_k$, the activities
$F(X')$ of \cite[(1.91)]{B-CMP122b}, the function $R'^{(k)}$ and the new boundary terms
$\mathbf{B}'^{(k)}$ of \cite{B-CMP122b}; hence it changes none of their covariance, locality or
bounds. Wherever a
quotient by $\mathbf{T}'_k(X)1$ occurs in \cite{B-CMP122b}, the label is non-null or the term is
absent.
\item[(d5)] The operations $\mathbf{T}_k$ produced at step $k$ again have the following property.
For $Y$ a union of components of $Z_k$, every characteristic function, $\delta$-function and
integration measure of $\mathbf{T}_k(Y)$ is a Borel object in the variables that
$\mathbf{T}_k(Y)$ integrates and in $V_k\restriction_Y$ alone.
\end{enumerate}
\end{enumerate}
\end{proposition}

Proof outstanding.

\smallskip\noindent\emph{Route.}
The proof would run at level $k$, with the clauses at every level $j<k$ as inductive hypotheses.
These are reality (a) at level $j$, the statement that every factor of \cite[(1.71)]{B-CMP122b}
at level $j$ is defined, finite and analytic (the second step below), and, for $\mathbf{T}_j$, the
Borel and localisation property stated in (d5). The steps would be taken in the following order.

First, (c) would be established directly from the construction of $\mathbf{R}$, with no
inductive input. By \cite[(1.89)]{B-CMP122a} the factor $\chi''_k$ may be dropped where the
denominators of \cite[(1.100)]{B-CMP122a} stand. The exterior is therefore restricted there by
$\chi_k(\Omega_k^{\sim4})\chi_{k,\Lambda_i}$. On its support, \cite[p.~193]{B-CMP122a} shows that
$V_k$ has an extension to the domain denoted $\Lambda$ there, satisfying the regularity bound
displayed there on $Z$, so
that the configuration $U_{k,Z}$ built from this extension is defined and the condition in
\cite[(1.75)]{B-CMP122a} has a meaning. The extended block field $V_\Lambda$ and the configuration
$U_{k,X_i}(V'V_{\Lambda_i})$ of \cite{B-CMP122a} exist there by \cite[Proposition~1]{B-CMP122a},
whose proof is given in \cite[p.~359]{B-CMP122b}.

Second, one would use the definedness part of clause (i-a) at level $k$, in the same induction:
every factor of \cite[(1.71)]{B-CMP122b} is defined, finite and
analytic in $(\mathbf{U},\mathbf{J})$ on $S_\omega\times\mathfrak{D}_k(X)$. For the factors depending
on $(\mathbf{U},\mathbf{J})$ this comes from \cite[(1.65)]{B-CMP122b} and the
convergence of the resummations \cite[(1.66)--(1.69)]{B-CMP122b} built on it. For the factors that
do not depend on $(\mathbf{U},\mathbf{J})$ it comes from definedness at $j<k$ and
\cite[(1.95)--(1.96)]{B-CMP122a}.

Third, (a) would be established at level $k$ from (a) at all $j<k$. In \cite{B-CMP109} the
extension of \cite[(1.7)]{B-CMP109} to pairs $(U,J)$ is introduced through the assumption that
there are analytic functions $E^{(j)}(X,\dots,U,J)$, the extension by the same formulas being
stated only for the explicit terms \cite[(1.4)]{B-CMP109}, while for $E^{(j)}(X,\dots,U,J)$ only
analyticity is assumed; reality on the diagonal $J=J(U)$ does not by itself propagate off the
diagonal, and this is why (a) is carried as a clause of the induction.

Fourth, (d1) would be derived from the property of (d5) at all $j<k$.

Fifth, (d2) would be obtained by writing \cite[(1.71)]{B-CMP122b} history by history. Finiteness
of $\Phi_\omega$ comes from the second step and its reality from the third; (b) follows from (d2).

Sixth, (d3) and (d4) would follow from \eqref{4.32:eq-positivity-dichotomy-1}, the strict
positivity of $e^{\Phi_\omega}$ at real pairs, and the independence of $\mu_\omega(S_\omega)$ from
$(\mathbf{U},\mathbf{J})$ given by (d1).

Last, (d5) would be read off the construction of $\mathbf{T}_k$ at step $k$, in which every
characteristic function, $\delta$-function and integration measure is introduced in the field
variables localised in the corresponding components of $Z_k$ \cite{B-CMP119}, \cite{B-CMP122a}.

Throughout, $k\le K$, the setting is that of Theorem~0, and $\mathfrak{p}$ denotes Ba{\l}aban's auxiliary parameters.

We write $N_k$ for the number of unit cubes of the scale-$k$ lattice on $\mathbb{T}_m$. Thus $N_j=L^{-4(j-k)}N_k$ for $j\ge k$. For a union $Y$ of scale-$j$ unit cubes, $|Y|_j$ is the number of such cubes in $Y$. We put $\ell_j:=\log g_j^{-2}$.

By Theorem~\ref{4.20:thm-representation}, in the notation of Notation~\ref{4.20:not-representation}, the measure $\rho_k\,dV_k$ is a finite sum $\sum_{\mathcal{S}}\tau_{\mathcal{S}}$ of terms of Ba{\l}aban's representation \cite[(2.18)]{B-CMP119}. Each term carries a label $\mathcal{S}=(\Omega_j,\Lambda_j)_{1\le j\le k}$, which is Ba{\l}aban's sequence of small-field regions. To a label we attach the following objects:
\begin{itemize}
\item $Z_j:=\mathbb{T}_m\setminus\Lambda_j$, the large-field region of scale $j$, a union of cubes of side $MR_j$, where $R_j$ denotes the large-field length of scale $j$ in Ba{\l}aban's construction;
\item $Z_j^+$, which is $Z_j$ together with the cubes of side $MR_j$ of $\Lambda_j$ that touch $Z_j$;
\item $\Gamma_n:=\Omega_n\setminus\Omega_{n+1}$, with $\Omega_0:=\mathbb{T}_m$ and $\Omega_{k+1}:=\emptyset$;
\item $|\Gamma_n|:=|\Gamma_n|_n$, the number of scale-$n$ unit cubes in $\Gamma_n$.
\end{itemize}
We write $E_k[\mathcal{S}]$ for the additive constant of the term $\tau_{\mathcal{S}}$: the constant $E_k$ of Notation~\ref{4.20:not-representation} at the label $\mathcal{S}$, which depends on the label through the regions $\Omega_{j+1}$, $\Lambda_{j+1}$, $j<k$.

The never-large label $\mathcal{S}_\emptyset$ has $\Omega_j=\Lambda_j=\mathbb{T}_m$ for all $j$. Its term is
\begin{equation}\label{4.33:eq-stability-1}
\tau_{\mathcal{S}_\emptyset}=\chi_k\exp\bigl[-g_k^{-2}A(U_k)+\mathbf{E}_k(U_k)+\mathbf{R}_k(U_k)
-E_k[\mathcal{S}_\emptyset]\bigr]\,dV_k ,
\end{equation}
where $U_k=U_k(V_k)$ is the background field and $\chi_k$ is the product of the small-field characteristic functions.

For each label $\mathcal{S}$, let $Y(\mathcal{S})$ be the union of all large-field components of $\tau_{\mathcal{S}}$. These are the components on which $\mathbf{R}$ has acted, of the first and of the second class, together with those on which it has not. By the representation of Theorem~\ref{4.20:thm-representation}, every operation composing $\mathbf{T}'_k(\mathcal{S})$ acts inside $Y(\mathcal{S})$, together with its delta-function constraints between consecutive block fields and its gauge-fixing factors. Outside $Y(\mathcal{S})$ the term $\tau_{\mathcal{S}}$ is the product of three factors:
\begin{itemize}
\item small-field characteristic functions;
\item the exponential of a bounded action;
\item $dV^{\mathrm{ext}}$.
\end{itemize}
Here $V^{\mathrm{ext}}$ is the block field outside $Y(\mathcal{S})$ and $dV^{\mathrm{ext}}$ is the product Haar probability. Hence the marginal of $\tau_{\mathcal{S}}$ on $V^{\mathrm{ext}}$ is absolutely continuous with respect to $dV^{\mathrm{ext}}$. We write $\tau_{\mathcal{S}}=dV^{\mathrm{ext}}\otimes\tau^{V^{\mathrm{ext}}}_{\mathcal{S}}$ for the disintegration, and $\|\cdot\|$ for total variation.

\begin{theorem}[Integrated stability (clause (i-b)) and equality of integrals (clause (i-c))]
\label{4.33:thm-stability}
Let $k\le K$, and assume the following.
\begin{itemize}
\item[(h1)] For every $j\le K$, $\rho_j\,dV_j$ has the representation of Theorem~\ref{4.20:thm-representation}, with the inductive bounds and termwise majorants stated there. Every term is a non-negative measure on real block fields, by the positivity dichotomy of the large-field operations. Every term has finite total mass, by the majorants \cite[(1.73), (1.79)]{B-CMP122b} carried by Theorem~\ref{4.20:thm-representation}.
\item[(h2)] On the support of $\chi_k$, at real fields, $\mathbf{E}_k(U_k)$ and $\mathbf{R}_k(U_k)$ are real and
\[
|\mathbf{E}_k(U_k)|+|\mathbf{R}_k(U_k)|\le C_EN_k,\qquad C_E:=E_1^\sharp(1-L^{-\beta_T})^{-1}+1 .
\]
Here $E_1^\sharp$ and $\beta_T$ are the constant and the rate of the bounds \cite[(2.45)--(2.46)]{B-CMP119} in the form carried by Theorem~\ref{4.20:thm-representation}. They depend only on $L$, $N$ and $\mathfrak{p}$, and are independent of $g_k$.
\item[(h3)] $\ell_j>0$ for all $j\le k$. This is supplied by the flow bounds of clause~(0) (Proposition~\ref{4.19:prop-clause-zero}).
\item[(h4)] The vacuum-energy theorem applies. Its hypotheses are supplied by clause~(0) (Proposition~\ref{4.19:prop-clause-zero}) and by the representation of Theorem~\ref{4.20:thm-representation}; the hypothesis on the localised terms is supplied by the independence of the localised terms from the small-field region (Corollary~\ref{4.4:cor-localised-independence}) together with the independence of the values at the unit configuration of the terms made by $\mathbf{R}$ from the label. Two of its conclusions are used, parts (a) and (c).
\begin{itemize}
\item[(a)] We have $-E_k[\mathcal{S}_\emptyset]=\mathfrak{l}_k-\mathfrak{r}_k$. The logarithmic part $\mathfrak{l}_k$ is non-negative whenever $\ell_j>0$ for all $j\le k$. The remaining part satisfies $|\mathfrak{r}_k|\le C_{\mathfrak{r}}N_k$, where $C_{\mathfrak{r}}\ge0$ depends only on $L$, $N$ and $\mathfrak{p}$.
\item[(c)] In the form used here, part (c) is the following upper bound: for every label $\mathcal{S}$,
\[
-E_k[\mathcal{S}]\le C_\Gamma(g_k)|\Gamma_k|
+\sum_{j=1}^k\bigl(C_{\log}\ell_{j-1}+C_{\mathrm{loc}}\bigr)|Z_j^+|_j ,
\]
where $C_{\log}$ and $C_{\mathrm{loc}}$ are constants of the vacuum-energy theorem, which enter (h6) below, and
\[
C_\Gamma(g_k):=\tfrac32(N^2-1)(\ell_k+B_\sharp\gamma^2)+C_{\mathfrak{r}} .
\]
\end{itemize}
\item[(h5)] For every label $\mathcal{S}$ and every real $V^{\mathrm{ext}}$,
\[
\|\tau^{V^{\mathrm{ext}}}_{\mathcal{S}}\|\le\mathfrak{M}_k(\mathcal{S})
\exp\Bigl[-E_k[\mathcal{S}]+C_{\mathrm{maj}}\sum_{n=0}^k|\Gamma_n|\Bigr].
\]
Here $\mathfrak{M}_k(\mathcal{S})\ge0$ is independent of $V^{\mathrm{ext}}$, and $C_{\mathrm{maj}}\ge0$ depends only on $L$, $N$ and $\mathfrak{p}$. The bound is supplied as follows.
\begin{itemize}
\item On the components of $Y(\mathcal{S})$ on which $\mathbf{R}$ has acted, by the majorant \cite[(1.79)]{B-CMP122b}.
\item On the components on which it has not acted, by the majorant of the operation $\mathbf{T}''_k(Z_k)$ of \cite[(1.72)]{B-CMP122b}. This is applied after the integral in $V_k$ there has removed the constraint between $V_{k-1}$ and $V_k$ by~(h7). On these components \cite[(2.48)]{B-CMP119} is used, in the form carried by Theorem~\ref{4.20:thm-representation}, for the terms frozen on components not renormalised by later steps.
\item Outside $Y(\mathcal{S})$, by \cite[(2.43)--(2.44)]{B-CMP119}, in the same form, for the small-field terms.
\end{itemize}
The Wilson part of the exponent is non-positive.
\item[(h6)] Put
\begin{equation}\label{4.33:eq-stability-2}
q_k(\mathcal{S}):=C_{\mathrm{maj}}\sum_{n<k}|\Gamma_n|
+\sum_{j=1}^k\bigl(C_{\log}\ell_{j-1}+C_{\mathrm{loc}}\bigr)|Z_j^+|_j .
\end{equation}
Then $\sum_{\mathcal{S}}\mathfrak{M}_k(\mathcal{S})\,e^{q_k(\mathcal{S})}\le e^{c_MN_k}$, where $c_M\ge0$ depends only on $L$, $N$ and $\mathfrak{p}$.
\item[(h7)] The block-averaging map $\mathcal{M}$ of Definition~\ref{1.5:def-block-average} is defined on all configurations as a Borel map with values in the block configurations. It has unit fibre mass: $\int dV\,\delta(\mathcal{M}(U)V^{-1})=1$ for every $U$.
\item[(h8)] Every normalising integral of the reference densities \cite[(1.100)--(1.102)]{B-CMP122a} lies in $(0,\infty)$ wherever it is used; this is the positivity of the reference normalisations. Consequently, on each component on which it acts, $\mathbf{R}$ replaces a term $\tau$ by $m_\tau\,\mathfrak{s}$, the extension by the normalised reference law. Here:
\begin{itemize}
\item $m_\tau$ is the integral of $\tau$ over the block field on the component, a function of the block field exterior to that component, which we denote $\widetilde{V}$;
\item $\mathfrak{s}(\cdot;\widetilde{V})$ is a probability measure on the block field of the component. On second-class components it lives on the slice fixed by the gauge-fixing factors, and is singular with respect to Haar measure.
\end{itemize}
\item[(h9)] On first-class components, the gauge-fixing factors are removed by the reversed Faddeev--Popov step. This step preserves integrals of gauge-invariant functions. It is an equivalence of densities, not an equality \cite[p.~378]{B-CMP122b}.
\item[(h10)] Dropping the characteristic function $\chi''_k$ under the integral, as in \cite[(1.89)]{B-CMP122a}, is an equality of integrals, in the form of the proposition on dropping $\chi''_k$ that accompanies the treatment of the creating step.
\item[(h11)] The creating step is an exact finite decomposition of unity inside $\mathbf{T}$; this is the exactness statement proved with the creating step.
\end{itemize}
Then the following hold.
\begin{itemize}
\item[(i-b)] Stability.
\begin{itemize}
\item[(1)] \emph{Lower member.} Under (h1), (h2), (h3) and (h4)(a), with $E_-:=C_E+C_{\mathfrak{r}}$,
\begin{equation}\label{4.33:eq-stability-3}
\chi_k\,e^{-g_k^{-2}A(U_k)-E_-N_k}\,dV_k\le\rho_k\,dV_k
\quad\text{as measures on real block fields}.
\end{equation}
The constant $E_-=E_-(L,\mathfrak{p},N)$ is independent of $g_k$, $k$, $K$, $m$ and $g_0$.
\item[(2)] \emph{Integrated upper member.} Under (h1), (h3), (h4)(a), (h4)(c), (h5), (h6) and (h7),
\begin{equation}\label{4.33:eq-stability-4}
\int\rho_k\,dV_k\le e^{E_+N_k},\qquad
E_+:=\tfrac32(N^2-1)(\log g_k^{-2}+B_\sharp\gamma^2)+C_{\mathfrak{r}}+C_{\mathrm{maj}}+c_M .
\end{equation}
The bound \eqref{4.33:eq-stability-4} holds after integration only; the pointwise bound \cite[(0.1)]{B-CMP122b} is not asserted by this theorem. Moreover $E_+\le C(L,\mathfrak{p},\gamma,N)(1+\log g_k^{-2})$.
\item[(3)] \emph{Integrated upper member with the constant of the last level.} Assume the hypotheses of (2) with $k=K$, and the hypotheses (h1) and (h7)--(h11) of (i-c). Let $E_+(g_K)$ denote the constant $E_+$ of \eqref{4.33:eq-stability-4} at level $K$, that is,
\[
E_+(g_K)=\tfrac32(N^2-1)(\log g_K^{-2}+B_\sharp\gamma^2)+C_{\mathfrak{r}}+C_{\mathrm{maj}}+c_M .
\]
Then $E_+(g_K)\ge0$, and
\[
\int\rho_j\,dV_j\le e^{E_+(g_K)N_j}\qquad\text{for all }j\le K.
\]
If moreover $g_K\ge g_-(g)$, then
\[
E_+(g_K)\le\tfrac32(N^2-1)\bigl(\log(g^{-2}+B_\sharp)+B_\sharp\gamma^2\bigr)
+C_{\mathfrak{r}}+C_{\mathrm{maj}}+c_M .
\]
\end{itemize}
\item[(i-c)] Under (h1) and (h7)--(h11), $\int\rho_k\,dV_k=\int\rho_0\,dU$ for every $k\le K$.
\end{itemize}
\end{theorem}

\begin{proof}
\emph{Lower member.} By (h1) every term $\tau_{\mathcal{S}}$ is non-negative at real fields. Hence $\rho_k\,dV_k\ge\tau_{\mathcal{S}_\emptyset}$.

By (h3) we have $\ell_j>0$ for $j\le k$, so $\mathfrak{l}_k\ge0$. Part (a) of (h4) then gives
\[
-E_k[\mathcal{S}_\emptyset]=\mathfrak{l}_k-\mathfrak{r}_k\ge-C_{\mathfrak{r}}N_k ,
\]
and the non-negative logarithmic part $\mathfrak{l}_k$ is discarded. By (h2), on the support of $\chi_k$ at real fields the exponent in \eqref{4.33:eq-stability-1} is real and satisfies
\[
\mathbf{E}_k(U_k)+\mathbf{R}_k(U_k)\ge-C_EN_k .
\]
Combining the two bounds in the exponent of \eqref{4.33:eq-stability-1} gives
\[
\tau_{\mathcal{S}_\emptyset}\ge\chi_k\,e^{-g_k^{-2}A(U_k)-(C_E+C_{\mathfrak{r}})N_k}\,dV_k ,
\]
which is \eqref{4.33:eq-stability-3}. Now $E_-=C_E+C_{\mathfrak{r}}$ is built from $E_1^\sharp$, $\beta_T$, $L$ and $C_{\mathfrak{r}}$. By (h2) and (h4)(a) these depend only on $L$, $N$ and $\mathfrak{p}$, which gives the stated independence.

\emph{Integrated upper member.} At real fields $\int\rho_k\,dV_k$ is real and equals the finite sum $\sum_{\mathcal{S}}\int\tau_{\mathcal{S}}$. By the disintegration above and $\int dV^{\mathrm{ext}}=1$,
\[
\Bigl|\int\tau_{\mathcal{S}}\Bigr|\le\int dV^{\mathrm{ext}}\,\|\tau^{V^{\mathrm{ext}}}_{\mathcal{S}}\|
\le\sup_{V^{\mathrm{ext}}}\|\tau^{V^{\mathrm{ext}}}_{\mathcal{S}}\| .
\]
We apply (h5) to the right-hand side. This gives
\[
\Bigl|\int\tau_{\mathcal{S}}\Bigr|\le\mathfrak{M}_k(\mathcal{S})
\exp\Bigl[-E_k[\mathcal{S}]+C_{\mathrm{maj}}\sum_{n=0}^k|\Gamma_n|\Bigr].
\]
The exponent on the right is
\[
-E_k[\mathcal{S}]+C_{\mathrm{maj}}\sum_{n=0}^k|\Gamma_n| .
\]
We split the sum over $n$ as
\[
C_{\mathrm{maj}}\sum_{n=0}^k|\Gamma_n|=C_{\mathrm{maj}}\sum_{n<k}|\Gamma_n|+C_{\mathrm{maj}}|\Gamma_k| ,
\]
so that the first summand on the right is the first term of $q_k(\mathcal{S})$ in
\eqref{4.33:eq-stability-2}.

We bound $-E_k[\mathcal{S}]$ by part (c) of the vacuum-energy theorem, that is, by (h4)(c). Its sum
over $j$ is the second term of $q_k(\mathcal{S})$ in \eqref{4.33:eq-stability-2}. Together with the
splitting above, this bounds the exponent by
\[
q_k(\mathcal{S})+\bigl(C_\Gamma(g_k)+C_{\mathrm{maj}}\bigr)|\Gamma_k| .
\]

We also use $|\Gamma_k|\le N_k$. Indeed, $\Omega_{k+1}=\emptyset$, so $\Gamma_k=\Omega_k$. Hence
$\Gamma_k\subset\mathbb{T}_m$, and $|\Gamma_k|=|\Gamma_k|_k$ counts scale-$k$ unit cubes of
$\mathbb{T}_m$, of which there are $N_k$. The coefficient of $|\Gamma_k|$ is non-negative:
$C_\Gamma(g_k)\ge0$ because $\ell_k>0$ by (h3), $B_\sharp>0$
(Definition~\ref{1.2:def-flow-constants}) and $C_{\mathfrak{r}}\ge0$ by (h4)(a), and
$C_{\mathrm{maj}}\ge0$ by (h5). So we may replace $|\Gamma_k|$ by $N_k$ in this upper bound.

This gives, for every label $\mathcal{S}$,
\begin{equation}\label{4.33:eq-stability-5}
\Bigl|\int\tau_{\mathcal{S}}\Bigr|\le\mathfrak{M}_k(\mathcal{S})\,e^{q_k(\mathcal{S})}\,
e^{(C_\Gamma(g_k)+C_{\mathrm{maj}})N_k} .
\end{equation}

We now sum \eqref{4.33:eq-stability-5} over the finitely many labels $\mathcal{S}$. Its last factor
does not depend on $\mathcal{S}$; taking it outside the sum gives
\[
\int\rho_k\,dV_k\le\sum_{\mathcal{S}}\Bigl|\int\tau_{\mathcal{S}}\Bigr|
\le e^{(C_\Gamma(g_k)+C_{\mathrm{maj}})N_k}\sum_{\mathcal{S}}\mathfrak{M}_k(\mathcal{S})\,
e^{q_k(\mathcal{S})} .
\]
By (h6) the remaining sum is at most $e^{c_MN_k}$, so
\[
\int\rho_k\,dV_k\le e^{(C_\Gamma(g_k)+C_{\mathrm{maj}}+c_M)N_k} .
\]
By the definition of $C_\Gamma(g_k)$ in (h4)(c), and since $\ell_k=\log g_k^{-2}$,
\begin{align*}
C_\Gamma(g_k)+C_{\mathrm{maj}}+c_M
&=\tfrac32(N^2-1)(\log g_k^{-2}+B_\sharp\gamma^2)+C_{\mathfrak{r}}+C_{\mathrm{maj}}+c_M\\
&=E_+ .
\end{align*}
This is \eqref{4.33:eq-stability-4}.

For the bound $E_+\le C(L,\mathfrak{p},\gamma,N)(1+\log g_k^{-2})$, write
$E_+=\tfrac32(N^2-1)\ell_k+C'$. Here $C'$ is the following non-negative number:
\begin{align*}
C'&=\tfrac32(N^2-1)B_\sharp\gamma^2+C_{\mathfrak{r}}+C_{\mathrm{maj}}+c_M\\
&\le\tfrac32(N^2-1)(C_\beta+1)\gamma^2+C_{\mathfrak{r}}+C_{\mathrm{maj}}+c_M ,
\end{align*}
where we used $B_\sharp\le C_\beta+1$ (Definition~\ref{1.2:def-flow-constants}). Since $\ell_k>0$ by (h3), it follows that $E_+\le C(1+\ell_k)$, with
\[
C:=\max\bigl\{\tfrac32(N^2-1),\ \tfrac32(N^2-1)(C_\beta+1)\gamma^2+C_{\mathfrak{r}}
+C_{\mathrm{maj}}+c_M\bigr\}.
\]
The constant $C$ depends only on $L$, $\mathfrak{p}$, $\gamma$ and $N$: $C_\beta$ is one of Ba{\l}aban's auxiliary parameters, and $C_{\mathfrak{r}}$, $C_{\mathrm{maj}}$ and $c_M$ depend only on $L$, $N$ and $\mathfrak{p}$ by (h4)(a), (h5) and (h6).

\emph{Equality of integrals.} We argue by induction on $k$; the case $k=0$ is trivial. We show that each operation producing $\rho_{k+1}$ from $\rho_k$ preserves the total mass.

First, $\mathbf{T}$. By (h7), $\int dV\,\delta(\mathcal{M}(U)V^{-1})=1$ for every $U$. By (h1), $\rho_k\,dV_k$ is a finite non-negative measure, so Fubini's theorem applies:
\begin{align*}
\int(\mathbf{T}\rho_k)\,dV_{k+1}
&=\int dV_k\,\rho_k(V_k)\int dV_{k+1}\,\delta\bigl(\mathcal{M}(V_k)V_{k+1}^{-1}\bigr)\\
&=\int\rho_k\,dV_k .
\end{align*}

Second, the decomposition of unity of (h11). It multiplies the integrand by finitely many functions whose pointwise sum is $1$. So the total masses of the pieces add up to the original mass.

Third, $\mathbf{R}$. By (h8) and Fubini's theorem,
\[
\int m_\tau\,\mathfrak{s}=\int d\widetilde{V}\,m_\tau(\widetilde{V})\int\mathfrak{s}(\cdot;\widetilde{V})
=\int\tau ,
\]
where only $\int\mathfrak{s}(\cdot;\widetilde{V})=1$ has been used. The remaining operations inside $\mathbf{R}$ preserve integrals: the reversed Faddeev--Popov step by (h9), applied to the gauge-invariant function $1$, and the dropping of $\chi''_k$ by (h10). Together these give $\int\mathbf{R}\rho=\int\rho$, which is \cite[(0.4)]{B-CMP122a}.

Hence $\int\rho_{k+1}\,dV_{k+1}=\int\rho_k\,dV_k$ for every $k<K$, and by induction $\int\rho_k\,dV_k=\int\rho_0\,dU$ for every $k\le K$.

\emph{Integrated upper member with the constant of the last level.} We first show that $E_+(g_K)$ is non-negative. By (h3) at level $K$ we have $\ell_K>0$, and $B_\sharp>0$ (Definition~\ref{1.2:def-flow-constants}), so the
term $\tfrac32(N^2-1)(\ell_K+B_\sharp\gamma^2)$ is non-negative. The constant $C_{\mathfrak{r}}$
is non-negative, since it bounds $|\mathfrak{r}_k|$ in (h4)(a). The constants $C_{\mathrm{maj}}$ and
$c_M$ of the upper bounds in (h5) and (h6) are non-negative, as stated there. Hence
$E_+(g_K)\ge0$.

Let $j\le K$. By (i-c), applied at level $j$ and at
level $K$, both $\int\rho_j\,dV_j$ and $\int\rho_K\,dV_K$ equal $\int\rho_0\,dU$; in particular they
are equal to each other. By the integrated upper member \eqref{4.33:eq-stability-4} at level $K$,
\[
\int\rho_j\,dV_j=\int\rho_K\,dV_K\le e^{E_+(g_K)N_K}\le e^{E_+(g_K)N_j}.
\]
The last inequality holds because $N_K=L^{-4(K-j)}N_j\le N_j$ and $E_+(g_K)\ge0$.

Suppose now that $g_K\ge g_-(g)$. Since $g_-(g)=(g^{-2}+B_\sharp)^{-1/2}$, this gives
$g_K^{-2}\le g^{-2}+B_\sharp$, hence $\log g_K^{-2}\le\log(g^{-2}+B_\sharp)$ because the logarithm
is increasing. Inserting this into the expression for $E_+(g_K)$ gives the last bound.
\end{proof}

Hypothesis (h6) is the summed form in which the inductive bounds of Theorem~\ref{4.20:thm-representation} enter here. It holds under Ba{\l}aban's conditions relating the exponent $p_0$ of the degree function $p_0(\cdot)$ to the exponent $r$ in $R_j$. Three contributions meet in it.
\begin{itemize}
\item The charge $q_k$ is $O(1)\log g_j^{-2}$ per unit cube of $Z_j^+$. Since $\log g_j^{-2}\le R_j$, this is at most a constant times $M^4R_j^5$ per cube of side $MR_j$ meeting $Z_j$. It is the contribution $O(1)\log g_j^{-2}|Z_j|$ of each large-field domain $Z_j$ described in \cite[p.~380]{B-CMP122b}.
\item The choice of the sequences of large-field regions costs factors $\exp O(1)(MR_j)^{-4}|Z_j|$.
\item The majorants carry per-cube factors $e^{-p_0(g_j)}$, which dominate these costs under the conditions just named.
\end{itemize}
What remains after these factors are paired is bounded by $e^{c_MN_k}$, with $c_M$ free of $g_k$.

No later clause of Theorem~0 evaluates $\rho_k$ at a point. Clauses (ii)--(iv) use only the following:
\begin{itemize}
\item total integrals and masses of terms;
\item the bound $|\langle F\rangle_K|\le\sup|F|$;
\item the explicit exponent of the never-large term \eqref{4.33:eq-stability-1};
\item the lower member \eqref{4.33:eq-stability-3};
\item termwise integrated majorants.
\end{itemize}
The object estimated in the proof of \cite[Corollary~3]{B-CMP119} is the integral $\int dV_k\,\rho_k$, estimated there by a sum of terms. The upper member \eqref{4.33:eq-stability-4} is that integrated statement.

The pointwise bound \cite[(0.1)]{B-CMP122b} is not asserted; it fails at every scale carrying a live second-class component. For the measures $\rho_k\,dV_k$ at such a scale, a bound $\rho_k\,dV_k\le e^{EN_k}\,dV_k$ fails for every $E$; this is the content of the proposition on the failure of the pointwise upper bound.

\begin{lemma}[Exactness of the step under complex sources]\label{4.34:lem-sourced-exactness}
Let $w=\sum_i z_i f_i$ be a source with complex coefficients $z=(z_i)_i$. Let $r$ be the number,
fixed in the setting of the correlation theorem, in terms of which the running shifts are bounded,
and suppose that the running shifts satisfy
\begin{equation*}
\sup|\zeta_j|\le\frac{8}{3}\,r .
\end{equation*}
Then at every level $k$ at which the induction hypothesis of the correlation theorem holds at
source $w$,
\begin{equation}\label{4.34:eq-sourced-exactness-1}
\int\rho^w_k\,dV_k=\int\rho^w_0\,dU ,
\end{equation}
and both members of \eqref{4.34:eq-sourced-exactness-1} are analytic in $w$. Here $\rho^w_k$ is
the effective density at level $k$ (Definition~\ref{1.5:def-densities}) computed at source $w$,
$dV_k$ is the product Haar measure on the block fields $V_k$ at level $k$, and $\rho^w_0$ is the
sourced bare density
\begin{equation*}
\rho^w_0(U)=\exp\bigl[-A(x_0-w,U)-E\bigr],
\end{equation*}
where $A(\cdot,U)$ is the weighted Wilson action of Definition~\ref{1.1:def-wilson-action} and $E$
is the additive constant of the bare density; its integral over the bare configurations is $Z(z)$.
\end{lemma}

\begin{proof}
Under the induction hypothesis at level $k$ the density $\rho^w_k$ is obtained from $\rho^w_0$ by
the renormalisation steps up to level $k$. Each step is composed of moves of the following five
kinds. Each move is an identity at fixed $w$: it
replaces the integrand by another one with the same integral, and it is linear in the integrand.
We check this for each kind in turn.

\emph{(a) Block averaging.} For every configuration, the block average $\bar U$ is given by
Ba{\l}aban's block-averaging map, extended to all configurations as in
Definition~\ref{1.5:def-block-average}. Inserting the new block variable $V$ uses
\begin{equation*}
\int dV\,\delta(\bar U V^{-1})=1 ,
\end{equation*}
which holds for every $\bar U$, where $\delta$ is the delta function of Haar measure. The
integrations over $U$ and over $V$ are then interchanged by Fubini's theorem. The total integral
is unchanged, and the move is linear in the integrand.

\emph{(b) Gauge fixing.} Let $\varepsilon_k$ denote the small-field bound on the bond variables at
level $k$. Fix a bond $(y,x)$ along which the gauge is fixed, and let $u$ be a gauge
transformation with $u(y)=1$; the transformed bond variable is
\begin{equation*}
U^u(y,x)=u(y)\,U(y,x)\,u(x)^{-1}=U(y,x)\,u(x)^{-1}.
\end{equation*}
For complex $a$ put
\begin{equation*}
\mathfrak{z}(a):=\int dv\,\chi\bigl(|v-1|<\varepsilon_k\bigr)\,e^{-a(1-\operatorname{Re}\operatorname{tr}v)} ,
\end{equation*}
the integral being over $SU(N)$ with Haar measure. The substitution $v=U(y,x)u(x)^{-1}$ preserves
Haar measure, so
\begin{equation*}
\int du(x)\,\chi\bigl(|U^u(y,x)-1|<\varepsilon_k\bigr)\,
e^{-a(1-\operatorname{Re}\operatorname{tr}U^u(y,x))}=\mathfrak{z}(a)
\end{equation*}
for every configuration $U$. Once $\mathfrak{z}(a)\ne0$ is known, dividing by $\mathfrak{z}(a)$
gives a resolution of $1$ valid for every $U$.

We now show $\mathfrak{z}(a)\ne0$ for the value of $a$ that occurs in the move, namely the complex
weight with which the density $1-\operatorname{Re}\operatorname{tr}U(y,x)$ of the gauge-fixed bond
enters the integrand. The imaginary part of this weight comes from the running shifts, so
$|\operatorname{Im}a|\le\sup|\zeta_j|\le\tfrac83 r$. Since $v$ is unitary and
$\operatorname{tr}\mathbb{1}=1$,
\begin{equation*}
1-\operatorname{Re}\operatorname{tr}v=\tfrac12\operatorname{tr}\bigl((v-1)^*(v-1)\bigr)
\le\tfrac12\,|v-1|^2 ,
\end{equation*}
with $|\cdot|$ the operator norm; hence $0\le1-\operatorname{Re}\operatorname{tr}v\le\varepsilon_k^2$
on the support of the indicator. On that support one therefore has
\begin{equation*}
|\operatorname{Im}a|\,\bigl(1-\operatorname{Re}\operatorname{tr}v\bigr)\le\frac{8}{3}\,r\,\varepsilon_k^2\le1 ,
\end{equation*}
the second inequality being the relation between $r$ and $\varepsilon_k$ in the setting of the
correlation theorem. Hence
$\cos\bigl((\operatorname{Im}a)(1-\operatorname{Re}\operatorname{tr}v)\bigr)\ge\cos1$ on the
support, and
\begin{align*}
|\mathfrak{z}(a)|\ \ge\ \operatorname{Re}\mathfrak{z}(a)
&=\int dv\,\chi\bigl(|v-1|<\varepsilon_k\bigr)\,
e^{-(\operatorname{Re}a)(1-\operatorname{Re}\operatorname{tr}v)}
\cos\bigl((\operatorname{Im}a)(1-\operatorname{Re}\operatorname{tr}v)\bigr)\\
&\ge\cos1\cdot\mathfrak{z}(\operatorname{Re}a)\ >\ 0 .
\end{align*}
The last inequality holds because the integrand of $\mathfrak{z}(\operatorname{Re}a)$ is positive
on the neighbourhood $\{|v-1|<\varepsilon_k\}$ of the identity, which has positive Haar measure.

The gauge-fixing move then interchanges the integration over $u$ with the integration over the
configuration. This interchange, as carried out in \cite{B-CMP109}, uses only the gauge invariance
of the remaining integrand. The weighted action $A(c,\cdot)$ is gauge invariant for every weight
$c$, complex weights included: each plaquette density $1-\operatorname{Re}\operatorname{tr}U(\partial p)$
is unchanged when $U(\partial p)$ is conjugated by a gauge transformation, and $c(p)$ does not
depend on $U$. The interchange therefore holds at complex source $w$ exactly as at real coupling.

\emph{(c) Representation moves.} These are decompositions of unity, changes of variables and Mayer
identities. Each of them is an identity for every fixed integrand and is linear in it. The
refinement of the large-field operation $\mathbf{T}$ used in this chapter is exact as well: by the
lemma of this chapter that establishes this refinement, the refined decomposition is exact, and
for each configuration exactly one of its terms is non-zero. The logarithms that occur are taken
of quantities represented by the
convergent expansions used for the correlation theorem. Each such logarithm is the sum of its
expansion, and exponentiating it returns the quantity, so passing to it is again an identity.

\emph{(d) Large-field normalisation.} Let $\Lambda$ be the region on which the operation
$\mathbf{R}$ integrates, and write $V|_{\Lambda}$ for the restriction of $V$ to $\Lambda$. The move
replaces an integrand $\tau^w$ by $\mathsf{s}^0N^w/D^0$, where
\begin{equation*}
N^w:=\int dV|_{\Lambda}\,\tau^w ,\qquad D^0:=\int dV|_{\Lambda}\,\mathsf{s}^0 ,
\end{equation*}
and $\mathsf{s}^0\ge0$ is the small-field density, which does not depend on $z$. In the form
stated in \cite[pp.~176--177]{B-CMP122a}, the denominators are positive and are determined by
small-field effective actions only. Thus $D^0$ is a $z$-independent positive number, and $N^w$ does
not depend on $V|_{\Lambda}$. Therefore
\begin{equation*}
\int dV|_{\Lambda}\,\frac{\mathsf{s}^0\,N^w}{D^0}=\frac{N^w}{D^0}\int dV|_{\Lambda}\,\mathsf{s}^0=N^w ,
\end{equation*}
and the integral is preserved. The move is linear in $\tau^w$, since $N^w$ is.

\emph{(e) Forward and reversed transformations.} The forward and reversed transformations of
\cite{B-CMP122b} are changes of variables. They preserve integrals by gauge invariance (the
weighted action being invariant at every complex weight, as in (b)) and by the invariance of Haar
measure.

Composing the moves of all steps up to level $k$ gives \eqref{4.34:eq-sourced-exactness-1}.

For analyticity, note first that $A(x_0-w,U)$ is linear in the coefficients $z_i$. Hence
$\rho^w_0$ is continuous in $U$ and entire in the $z_i$. Its integral over the compact
configuration space of the bare lattice, with respect to the product Haar measure, is therefore
analytic in $w$. By \eqref{4.34:eq-sourced-exactness-1} the left member is analytic as well.

In this argument no denominator depending on the source occurs, and no logarithm of an integral
over a large-field region is taken.
\end{proof}

A renormalisation step $\rho_{k+1}=\mathbf{R}\mathbf{T}\rho_k$ (Definition~\ref{1.5:def-densities})
uses two gauge-fixing devices inside $\mathbf{T}$, and $\mathbf{R}$ reverses a gauge fixing on part of
the large-field region. The total integrals depend on none of these three devices. In particular
they do not depend on the form $\mathsf{G}$ of device (2); this is part (ii) of
Proposition~\ref{4.35:prop-gauge-fixing} below. Two facts give this independence. First, the
representations of devices (1) and (2) are exact rewritings of the push-forward
$\mathcal{M}_*\tau$ (Definition~\ref{1.5:def-block-average}); here device (2) is the exponential
gauge-fixing term of \cite[(1.5), (2.11)]{B-CMP109}, with its normalisation included. Second,
device (3) preserves total mass.

The expectations of gauge-invariant observables at cutoff $K$ are defined through the bare density
$\rho_0$ alone, and so they depend on no device. The representative of a density after a reversed
Faddeev--Popov step does depend on device (3). No statement is made on how the densities, or their
individual localised terms, depend on $\mathsf{G}$. Throughout, $G=SU(N)$ is the gauge group.

\emph{The devices.} The three devices are the following.

(1) Inside $\mathbf{T}$ the fluctuation integral is expressed in the axial gauge on the trees of the
blocks \cite[(1.5)]{B-CMP119}. This is the gauge of \cite[(0.16)]{B-CMP109}.

(2) The quadratic form of the fluctuation field $B$ is that of \cite[(1.5), (2.11)]{B-CMP109}. It is
the form of \cite[(3.156)]{B-CMP99b}, with the $\delta$-function gauge-fixing term replaced by an
exponential gauge-fixing term. The gauge-fixing term itself is not quadratic in $B$. Its
contribution to the quadratic form $\Delta^{(k)}$ of the fluctuation field
(Definition~\ref{4.4:def-quadratic-form}) is $\mathsf{G}(B,B)$,
where $\mathsf{G}$ is the following bilinear form. $\mathsf{G}(B,B')$ is a sum over the blocks
$\square_y$ with centre $y\in\Lambda'=\Omega_t\cap T^{(t)}$, and over sites $x\in\square_y$ with
$x\neq y$. Each summand is the trace of the product of the part linear in $B$ and the part linear
in $B'$ of
\begin{equation*}
\frac{1}{i}\log\bigl[(e^{iB}V)\langle y,x\rangle\,V\langle y,x\rangle^{-1}\bigr].
\end{equation*}
Here $V$, $W\langle y,x\rangle$ (for a configuration $W$), the blocks $\square_y$, the index $t$, the
set $\Omega_t$ and hence $\Lambda'$ are those fixed with the quadratic form of the fluctuation field
in Definition~\ref{4.4:def-quadratic-form}, and $T^{(t)}$ is as in
Definition~\ref{1.5:def-block-average}. The logarithm is the intrinsic logarithm
of $G$ near the identity.

(3) Inside $\mathbf{R}$, on first-class components of the large-field region (the division of the
components of the large-field region into a first and a second class is that of Ba{\l}aban's
large-field operation \cite{B-CMP122a}), the gauge-fixing
$\delta$-functions are removed by reversing the Faddeev--Popov procedure. Ba{\l}aban notes that the
result is ``equivalent, but not equal'' to the density before removal \cite[p.~378]{B-CMP122b}.
Likewise he states that the output of his operation on the renormalised densities ``may not be
equal to $\rho_k$, but is equivalent to it'' \cite{B-CMP122a}. After a reversed Faddeev--Popov
step, $\rho_k\,dV_k$ keeps $\delta$-function factors on second-class components. It is therefore a
positive measure, which in general is not a function of $V_k$ times $dV_k$. We give Ba{\l}aban's
phrase the following meaning for measures.

Two positive finite Borel measures $\tau,\tau'$ on the space of block fields $V$ at one scale are
\emph{equivalent}, written $\tau\sim\tau'$, if
\begin{equation*}
\int F\,d\tau=\int F\,d\tau'
\end{equation*}
for every bounded Borel function $F$ of $V$ that is invariant under $G$-valued gauge
transformations. A density $\sigma$ is identified with the measure $\sigma\,dV$.

The construction of the step rests on two theorems of Ba{\l}aban on gauge fixing, which are cited by
statement in parts (A) and (B) below.

\begin{proposition}[Gauge fixing inside the step: the three devices, Ba{\l}aban's gauge-fixing
theorem and the axial-gauge minimal configuration cited by statement, and independence of the
total integrals from the devices]\label{4.35:prop-gauge-fixing}
\leavevmode
\begin{itemize}
\item[(A)] \emph{Ba{\l}aban's gauge-fixing theorem} \cite[Theorem~2]{B-CMP99a}. There exist
constants $B_1$, $B_2(\beta_0)$ and $c_1$ with the following property. Let $U_0$ and $U'$ be
configurations satisfying the conditions \cite[(1.33)--(1.35)]{B-CMP99a}, with parameters
$\alpha_0,\alpha_1$ such that $\alpha_0+\alpha_1\le c_1$. Then there exists exactly one gauge
transformation $u$ with the following two properties: $u$ satisfies \cite[(1.29)]{B-CMP99a}, and
the conditions \cite[(1.36)--(1.39)]{B-CMP99a} hold for the configuration $U_1$ obtained from $U'$
by the gauge transformation $u$, in the convention of \cite{B-CMP99a}. The constants $B_1$,
$B_2(\beta_0)$ and the parameter $\beta_0$ are those of \cite[Theorem~2]{B-CMP99a}.
\item[(B)] \emph{The minimal orbit and its regularity} \cite[Theorem~1]{B-CMP102b}. In the notation
of \cite{B-CMP102b}, there exist positive constants $a_0$, $a_1$, $B_3$ and $M(\varepsilon_1)$,
together with the further constants fixed there, with the following properties for every
configuration $V$ that satisfies \cite[(7)]{B-CMP102b} with $\varepsilon_1\le a_1$.
\begin{itemize}
\item[(B1)] There exists a minimal orbit in the space \cite[(8)]{B-CMP102b}.
\item[(B2)] If $B_3\varepsilon_1\le\varepsilon_0$ and $\varepsilon_0\le a_0$, this orbit is the
unique critical orbit in the space \cite[(6)]{B-CMP102b}.
\item[(B3)] The minimal configurations have the following regularity property. For every cube of
the class described in \cite{B-CMP102b}, of size $2ML^j\eta$ with $M\le M(\varepsilon_1)$, there
exists a gauge transformation $u$ defined on a neighbourhood of the cube such that, on the cube,
the transformed minimal configuration satisfies the bounds \cite[(9)]{B-CMP102b}.
\end{itemize}
\item[(C)] \emph{Independence of the total integrals from the devices.} Assume the following.
\begin{itemize}
\item[(a)] The block averaging map $\mathcal{M}$ is Borel. Moreover, for every $G$-valued gauge
function $u$ on the finer lattice there is a $G$-valued gauge function $\bar u$ on the coarser
lattice such that $\mathcal{M}(V^u)=(\mathcal{M}V)^{\bar u}$ at every configuration $V$.
\item[(b)] For every measure $\tau$ to which $\mathbf{T}$ is applied, $\mathbf{T}\tau$ is the
push-forward $\mathcal{M}_*\tau$. Moreover, for every bounded Borel gauge-invariant $F$,
Ba{\l}aban's representation of $\mathbf{T}\tau$ integrates $F$ to $\int(F\circ\mathcal{M})\,d\tau$.
This representation has the following features:
\begin{itemize}
\item it is written in the axial gauge \cite[(1.5)]{B-CMP119};
\item it uses the exponential gauge-fixing term of \cite[(1.5), (2.11)]{B-CMP109} together with its
normalisation;
\item it uses the decompositions of unity of \cite[\S3]{B-CMP119}, refined by the additional
thresholds at the creation of large-field regions;
\item the integrand of each piece is averaged over gauge transformations wherever a
representative after a reversed Faddeev--Popov step is read.
\end{itemize}
\item[(c)] $\mathbf{R}$ is well defined on every measure $\tau$ to which it is applied in the
construction of $\rho_1,\dots,\rho_K$. Moreover, with the normalisation \cite[(1.102)]{B-CMP122a},
the total mass of $\mathbf{R}\tau$ equals that of $\tau$ \cite[(0.4)]{B-CMP122a}. This holds for
every representative $\mathbf{R}\tau$ of the equivalence class produced by the reversed
Faddeev--Popov step.
\end{itemize}
Then the following hold.
\begin{itemize}
\item[(i)] If $\tau\sim\tau'$, then $\mathbf{T}\tau\sim\mathbf{T}\tau'$.
\item[(ii)] For every $k\le K$,
\begin{equation*}
\int\rho_k\,dV_k=\int\rho_0\,dU.
\end{equation*}
In particular this integral does not depend on the axial gauge inside $\mathbf{T}$, on the form
$\mathsf{G}$ of the exponential gauge fixing, or on the representatives chosen at the reversed
Faddeev--Popov steps.
\end{itemize}
\end{itemize}
\end{proposition}

Part (A) is \cite[Theorem~2]{B-CMP99a}, and part (B) is \cite[Theorem~1]{B-CMP102b}. Their proofs
are Ba{\l}aban's and are not reproduced. For the densities of the construction, part (ii) of (C) is
clause (i-c), the equality of integrals of Theorem~\ref{4.33:thm-stability}.

\begin{proof}[Proof of \textup{(C)}]
(i) The measures $\mathbf{T}\tau$ and $\mathbf{T}\tau'$ are positive and finite, since by (b) they
are push-forwards of positive finite measures. Let $F$ be a bounded Borel gauge-invariant function
of the field at the coarser scale. By (b),
\begin{equation*}
\int F\,d(\mathbf{T}\tau)=\int(F\circ\mathcal{M})\,d\tau,
\end{equation*}
and likewise for $\tau'$. By (a), $\mathcal{M}$ is Borel. Hence $F\circ\mathcal{M}$ is Borel, and
it is bounded because $F$ is. It is also gauge invariant. Indeed, let $u$ be a $G$-valued gauge
function on the finer lattice, and let $\bar u$ be the function that (a) provides. Then
\begin{equation*}
(F\circ\mathcal{M})(V^u)=F\bigl((\mathcal{M}V)^{\bar u}\bigr)=F(\mathcal{M}V),
\end{equation*}
where the second equality is the gauge invariance of $F$. Since $\tau\sim\tau'$, the definition of
equivalence gives
\begin{equation*}
\int(F\circ\mathcal{M})\,d\tau=\int(F\circ\mathcal{M})\,d\tau'.
\end{equation*}
Therefore $\int F\,d(\mathbf{T}\tau)=\int F\,d(\mathbf{T}\tau')$ for every such $F$, that is,
$\mathbf{T}\tau\sim\mathbf{T}\tau'$.

(ii) We argue by induction on $k$. For $k=0$ there is nothing to prove. Assume the identity for
$k<K$, and recall that $\rho_{k+1}=\mathbf{R}(\mathbf{T}\rho_k)$. By (c), applied to the measure
$\mathbf{T}\rho_k$, the total mass of $\mathbf{R}(\mathbf{T}\rho_k)$ equals that of
$\mathbf{T}\rho_k$. The constant function $F=1$ is bounded, Borel and gauge invariant. So (b) with
$F=1$ gives $\int d(\mathbf{T}\rho_k)=\int(1\circ\mathcal{M})\rho_k\,dV_k$. Together,
\begin{equation*}
\int\rho_{k+1}\,dV_{k+1}=\int d\,\mathbf{R}(\mathbf{T}\rho_k)=\int d\,(\mathbf{T}\rho_k)
=\int\rho_k\,dV_k .
\end{equation*}
By the induction hypothesis the right side equals $\int\rho_0\,dU$. The integral
$\int\rho_0\,dU$ is defined by the Wilson action alone, so it involves no choice made inside the
steps.
\end{proof}

Hypotheses (a)--(c) hold in the construction for the following reasons.

Hypothesis (a) holds in the construction. The averaging map fixed in the setting
(Definition~\ref{1.5:def-block-average}) is Borel and gauge covariant at every configuration,
including configurations off the small-field sets.

Hypothesis (b) also holds in the construction, for three reasons. First, $\mathbf{T}$ is the
push-forward \cite[(0.1)]{B-CMP119} under $\mathcal{M}$. Second, the decompositions of unity in (b)
are exact decompositions of unity, so $\mathbf{T}$ preserves total mass. Third, for gauge-invariant
$F$, Ba{\l}aban's display evaluates $\int(F\circ\mathcal{M})\,d\tau$. This is because the
Faddeev--Popov Jacobian of the axial gauge on trees is $1$, and because the exponential
gauge-fixing term of \cite[(1.5), (2.11)]{B-CMP109}, with its normalisation, is an exact rewriting
of the same integral.

The well-definedness of $\mathbf{R}$ in hypothesis (c) rests on the positivity dichotomy of the
large-field operations. This well-definedness comprises two positivity statements: the positivity
of the normalising integrals, which is a premise of \cite[(0.4)]{B-CMP122a}, and the dichotomy of
\cite[Proposition~1]{B-CMP122b}. Once $\mathbf{R}$ is well defined, the mass identity in (c)
follows from four facts:
\begin{itemize}
\item Fubini's theorem;
\item the unit mass of the fibres of $\mathcal{M}$;
\item the normalisation \cite[(1.102)]{B-CMP122a};
\item the Faddeev--Popov identity on first-class components (axial gauge on trees, Jacobian $1$),
by which all representatives of the reversed step have equal total mass.
\end{itemize}

The expectation of a gauge-invariant observable at cutoff $K$ is defined through the bare density
and depends on no choice made inside the steps. Its representation through the histories of
$\rho_{K-s}$ evaluates the reversed step only against gauge-invariant functions.

No statement that $\mathbf{R}$ itself respects the equivalence $\sim$ defined above is made. The
tests and decompositions of $\mathbf{R}$ are evaluated on the representative Ba{\l}aban displays.

\begin{definition}[The group $SU(N)$, its Lie algebra and two norms]\label{4.36:not-group-norms}
The dimension is $d=4$; the letter $d$ is kept in formulas where no confusion can arise. The
integer $N\ge 2$ is fixed before every other parameter, and $G=SU(N)$. We write
$\mathfrak{g}$ for the real vector space of traceless Hermitian $N\times N$ matrices, and group
elements are written $U=e^{iA}$ with $A\in\mathfrak{g}$, as in \cite{B-CMP109}. The
complexifications are $\mathfrak{g}^{c}=\mathfrak{sl}(N,\mathbb{C})$ and
$G^{c}=SL(N,\mathbb{C})$. We put
\[
D=N^2-1,\qquad \mathsf{n}=(N-1)^{1/2}.
\]
\begin{enumerate}
\item \emph{Trace.} $\operatorname{tr}=N^{-1}\operatorname{Tr}$ is the normalised trace, so
that $\operatorname{tr}\mathbf{1}=1$, as in \cite{B-CMP109} and \cite{B-CMP99b}.
\item \emph{Operator norm.} For a complex $N\times N$ matrix $Y$, $|Y|$ is the operator norm of
$Y$ on $\mathbb{C}^N$ with its standard Hermitian inner product. This is the norm used in
\cite[(19)]{B-CMP98}, and it satisfies
\begin{equation}\label{4.36:eq-group-norms-1}
|XY|\le |X|\,|Y|,
\end{equation}
as in \cite[(20)]{B-CMP98}.
\item \emph{Trace norm.} For a complex $N\times N$ matrix $Y$ we put
$\|Y\|^2=\operatorname{tr} Y^{*}Y$, with the normalised trace. This is the norm $\|\cdot\|$
used in \cite{B-CMP98}.
\item \emph{Norm on $\mathfrak{g}$.} For $X\in\mathfrak{g}$ we put
$|X|_{\mathfrak{g}}=(\operatorname{tr}X^2)^{1/2}$.
\item \emph{Eigen-angles.} A unitary matrix $u$ has eigenvalues $e^{i\theta_p}$,
$p=1,\dots,N$, with $\theta_p\in\,]-\pi,\pi]$; the $\theta_p$ are the eigen-angles of $u$, as
in \cite{B-CMP98}.
\end{enumerate}
For every Hermitian $X$ with eigenvalues $\theta_1,\dots,\theta_N$ (so that $e^{iX}$ has
eigenvalues $e^{i\theta_p}$),
\begin{equation}\label{4.36:eq-group-norms-2}
|X|=\max_p|\theta_p|.
\end{equation}
For every unitary $u$ with eigen-angles $\theta_1,\dots,\theta_N$,
\begin{equation}\label{4.36:eq-group-norms-3}
|u-\mathbf{1}|=\max_p|e^{i\theta_p}-1|=\max_p 2\,|\sin(\theta_p/2)|.
\end{equation}
For every $X\in\mathfrak{g}$ with eigenvalues $\theta_1,\dots,\theta_N$,
\begin{equation}\label{4.36:eq-group-norms-4}
|X|_{\mathfrak{g}}=\Bigl(N^{-1}\sum_{p=1}^N\theta_p^2\Bigr)^{1/2}=\|X\|.
\end{equation}
Moreover $D=\dim_{\mathbb{R}}\mathfrak{g}$. For $N=2$, let $\sigma_1,\sigma_2,\sigma_3$ be the
Pauli matrices, which are traceless, Hermitian and satisfy
$\sigma_a\sigma_b+\sigma_b\sigma_a=2\delta_{ab}\mathbf{1}$, where $\delta_{ab}$ is the
Kronecker symbol; let $n\in\mathbb{R}^3$ be a unit vector, $n\cdot\sigma=\sum_{a=1}^3 n_a\sigma_a$,
$\theta\ge 0$ and $X=\theta\, n\cdot\sigma$. Then $X\in\mathfrak{g}$ and
$|X|=|X|_{\mathfrak{g}}=\theta$.
\end{definition}

\begin{proof}
A Hermitian matrix $X$ and a unitary matrix $u$ are normal, hence unitarily diagonalisable. For
a normal matrix $Y=W\Delta W^{*}$ with $W$ unitary and $\Delta$ diagonal, unitary invariance of
the Euclidean norm on $\mathbb{C}^N$ gives $|Y|=|\Delta|$. This is the maximal modulus of the
diagonal entries, that is, of the eigenvalues of $Y$. Applied to a Hermitian $X$ with
eigenvalues $\theta_1,\dots,\theta_N$, this gives \eqref{4.36:eq-group-norms-2}.

Let $u$ be unitary with eigen-angles $\theta_1,\dots,\theta_N$. The matrix $u-\mathbf{1}$ is
diagonalised by the same unitary as $u$, so it is normal with
eigenvalues $e^{i\theta_p}-1$. This gives the first equality in \eqref{4.36:eq-group-norms-3}.
The second follows from the identity
\[
e^{i\theta}-1=e^{i\theta/2}\bigl(e^{i\theta/2}-e^{-i\theta/2}\bigr)
=2i\,e^{i\theta/2}\sin(\theta/2),
\]
whose modulus is $2|\sin(\theta/2)|$.

Let $X\in\mathfrak{g}$ have eigenvalues $\theta_1,\dots,\theta_N$. The matrix $X^2$ has
eigenvalues $\theta_p^2$, so
$\operatorname{tr}X^2=N^{-1}\sum_p\theta_p^2$. Since $X^{*}=X$, we also have
$\|X\|^2=\operatorname{tr}X^{*}X=\operatorname{tr}X^2$. This proves
\eqref{4.36:eq-group-norms-4}.

The Hermitian $N\times N$ matrices form a real vector space of dimension $N^2$: there are $N$
real diagonal entries and $N(N-1)/2$ complex entries above the diagonal. The condition
$\operatorname{Tr}X=0$ is one nonzero real linear condition on this space. Hence
$\dim_{\mathbb{R}}\mathfrak{g}=N^2-1=D$.

For $N=2$, each $\sigma_a$ is Hermitian and traceless, so $X=\theta\,n\cdot\sigma$ lies in
$\mathfrak{g}$. From $\sigma_a\sigma_b+\sigma_b\sigma_a=2\delta_{ab}\mathbf{1}$ and
$|n|=1$ we get $(n\cdot\sigma)^2=\mathbf{1}$. Hence $X^2=\theta^2\mathbf{1}$, and the
eigenvalues of $X$ are $\pm\theta$, because $\operatorname{Tr}X=0$. By
\eqref{4.36:eq-group-norms-2}, $|X|=\theta$. By \eqref{4.36:eq-group-norms-4},
\[
|X|_{\mathfrak{g}}=\bigl(\tfrac12(\theta^2+\theta^2)\bigr)^{1/2}=\theta.
\]
\end{proof}

\begin{lemma}[Elementary inequalities for eigen-angles, traces and norms]
\label{4.36:lem-eigenangle-inequalities}
We use the group $SU(N)$, the Lie algebra $\mathfrak{g}$ of traceless Hermitian $N\times N$ matrices,
its complexification $\mathfrak{g}^{c}=sl(N,\mathbb{C})$, the normalised trace
$\operatorname{tr}=N^{-1}\operatorname{Tr}$, the operator norm $|\cdot|$, the norm
$\|Y\|=(\operatorname{tr}Y^{*}Y)^{1/2}$, the norm $|X|_{\mathfrak{g}}=(\operatorname{tr}X^{2})^{1/2}$ on
$\mathfrak{g}$ and the number $\mathsf{n}=(N-1)^{1/2}$, all as in
Notation~\ref{4.36:not-group-norms}. We write $M_N(\mathbb{C})$ for the algebra of complex $N\times N$
matrices and $U(N)$ for the unitary group. For a Hermitian $X$ we write $\theta_1,\dots,\theta_N$ for
its eigenvalues, so that $|X|=\max_p|\theta_p|$ and, for $X\in\mathfrak{g}$,
$|X|_{\mathfrak{g}}^2=N^{-1}\sum_p\theta_p^2$ and $\sum_p\theta_p=0$. Then the following hold.
\begin{enumerate}
\item[(F1)] For $X\in\mathfrak{g}$,
\begin{equation}\label{4.36:eq-eigenangle-inequalities-a}
|X|_{\mathfrak{g}}\le|X|\le\mathsf{n}\,|X|_{\mathfrak{g}},
\end{equation}
and the second inequality is an equality at $X=\operatorname{diag}(1,\dots,1,1-N)$.
\item[(F2)] For Hermitian $X$ and all $p\neq q$,
\begin{equation*}
(\theta_p-\theta_q)^2\le2(\theta_p^2+\theta_q^2)\le2\sum_r\theta_r^2=2N\operatorname{tr}X^2,
\end{equation*}
which is $2N|X|_{\mathfrak{g}}^2$ for $X\in\mathfrak{g}$; and
\begin{equation*}
\sum_{p<q}(\theta_p-\theta_q)^2=N\sum_r\theta_r^2-\Bigl(\sum_r\theta_r\Bigr)^2,
\end{equation*}
which for traceless $X$ equals $N^2|X|_{\mathfrak{g}}^2$.
\item[(F3)] Let $\operatorname{sinc}y=(\sin y)/y$ for $y\neq0$ and $\operatorname{sinc}0=1$. Then
$\operatorname{sinc}$ is positive on $(0,\pi)$ and strictly decreasing on $(0,\pi]$. Hence
$u\mapsto(1-\cos u)/u^2=\frac12\operatorname{sinc}^2(u/2)$ is decreasing on $(0,2\pi)$, and for every
$t\in(0,2\pi)$
\begin{equation*}
1-\cos\theta\ge\frac{1-\cos t}{t^2}\,\theta^2\qquad(|\theta|\le t).
\end{equation*}
At $t=\pi$ this gives
\begin{equation*}
\frac{2}{\pi^2}\,\theta^2\le1-\cos\theta\le\frac12\,\theta^2\qquad(|\theta|\le\pi),
\end{equation*}
and the upper bound holds for all real $\theta$. Moreover $\sin x\ge(2/\pi)x$ for $x\in[0,\pi/2]$,
equivalently $\arcsin y\le(\pi/2)y$ for $y\in[0,1]$.
\item[(F4)] $\operatorname{sinc}^2y\ge1-y^2/3$ for all real $y$.
\item[(F5)] For $a_1,\dots,a_n\in[0,1]$,
\begin{equation*}
\prod_{i=1}^n(1-a_i)\ge1-\sum_{i=1}^na_i.
\end{equation*}
\item[(F6)] For $X\in\mathfrak{g}$: $|e^{iX}-\mathbf{1}|\le|X|\le\mathsf{n}|X|_{\mathfrak{g}}$. For
$Z\in M_N(\mathbb{C})$: $|e^{Z}-\mathbf{1}|\le e^{|Z|}-1$. For all $X,Y\in M_N(\mathbb{C})$:
$|\operatorname{tr}Y|\le|Y|$; $|\operatorname{tr}(XY)|\le\|X\|\,\|Y\|$; $\|XY\|\le|X|\,\|Y\|$ and
$\|XY\|\le\|X\|\,|Y|$; $\|Y\|\le|Y|\le\sqrt N\,\|Y\|$. If $E\in M_N(\mathbb{C})$ and
$|E|\le\tau<1$, then $\mathbf{1}+E$ is invertible and $|(\mathbf{1}+E)^{-1}|\le(1-\tau)^{-1}$.
\item[(F7)] For $W\in U(N)$,
\begin{equation*}
N^{-1}|W-\mathbf{1}|^2\le2(1-\operatorname{Re}\operatorname{tr}W)=\|W-\mathbf{1}\|^2
\le|W-\mathbf{1}|^2.
\end{equation*}
\item[(F8)] For $A,B,Y\in M_N(\mathbb{C})$ with $|A|,|B|\le1$ and $T,T'\in\mathfrak{g}$:
\begin{equation*}
|\operatorname{tr}(AT)|\le|T|_{\mathfrak{g}},\qquad
|\operatorname{tr}(ATBT')|\le|T|_{\mathfrak{g}}|T'|_{\mathfrak{g}},\qquad
|\operatorname{tr}(T^2Y)|\le|T|_{\mathfrak{g}}^2|Y|.
\end{equation*}
\item[(F9)] For $W\in SU(N)$: $|W-\mathbf{1}|\le\mathsf{n}\,\|W-\mathbf{1}\|$.
\item[(F10)] For $Z\in\mathfrak{g}^{c}$ the inequality $|Z|\le\sqrt N\|Z\|$ of (F6) holds, and the
constant $\sqrt N$ is attained at $Z=E_{12}$, the matrix with entry $1$ in position $(1,2)$ and $0$
elsewhere: $|E_{12}|=1$, $\|E_{12}\|=N^{-1/2}$. In particular the constant $\mathsf{n}$ of (F1)
does not extend from $\mathfrak{g}$ to $\mathfrak{g}^{c}$.
\end{enumerate}
\end{lemma}

\begin{proof}
(F1). Since $N^{-1}\sum_q\theta_q^2\le\max_q\theta_q^2$, we have $|X|_{\mathfrak{g}}^2\le|X|^2$.
For the second inequality fix $p$. Since $\sum_q\theta_q=0$, $\theta_p=-\sum_{q\neq p}\theta_q$, and
the Cauchy--Schwarz inequality for the $N-1$ terms gives
$\theta_p^2\le(N-1)\sum_{q\neq p}\theta_q^2$. Adding $(N-1)\theta_p^2$ to both sides,
\begin{equation*}
N\theta_p^2\le(N-1)\sum_q\theta_q^2=(N-1)N|X|_{\mathfrak{g}}^2,
\end{equation*}
so $\theta_p^2\le\mathsf{n}^2|X|_{\mathfrak{g}}^2$ for every $p$, that is
$|X|\le\mathsf{n}|X|_{\mathfrak{g}}$. At $X=\operatorname{diag}(1,\dots,1,1-N)$, which is traceless,
$|X|=N-1$ because $N\ge2$, and
\begin{equation*}
\operatorname{tr}X^2=\frac{(N-1)+(N-1)^2}{N}=N-1,
\end{equation*}
so $\mathsf{n}|X|_{\mathfrak{g}}=(N-1)^{1/2}(N-1)^{1/2}=N-1=|X|$.

(F2). The first inequality is $2(\theta_p^2+\theta_q^2)-(\theta_p-\theta_q)^2=(\theta_p+\theta_q)^2\ge0$;
the second holds because for $p\neq q$ the sum $\sum_r\theta_r^2$ contains both $\theta_p^2$ and
$\theta_q^2$; and $\sum_r\theta_r^2=\operatorname{Tr}X^2
=N\operatorname{tr}X^2$. For the identity, expanding the squares,
\begin{equation*}
\sum_{p,q}(\theta_p-\theta_q)^2=2N\sum_r\theta_r^2-2\Bigl(\sum_r\theta_r\Bigr)^2,
\end{equation*}
and the left-hand side is twice $\sum_{p<q}(\theta_p-\theta_q)^2$, the diagonal terms being zero.
If $X$ is traceless, $\sum_r\theta_r=0$ and $N\sum_r\theta_r^2=N^2|X|_{\mathfrak{g}}^2$.

(F3). On $(0,\pi)$, $\sin y>0$, so $\operatorname{sinc}y>0$. For $y\neq0$,
$(\operatorname{sinc})'(y)=h(y)/y^2$ with $h(y)=y\cos y-\sin y$. Here $h(0)=0$ and
$h'(y)=\cos y-y\sin y-\cos y=-y\sin y<0$ on $(0,\pi)$; hence $h$ is strictly decreasing on
$[0,\pi]$ and $h<0$ on $(0,\pi]$, so $\operatorname{sinc}$ is strictly decreasing on $(0,\pi]$.
Since $1-\cos u=2\sin^2(u/2)$, we have $(1-\cos u)/u^2=\frac12\operatorname{sinc}^2(u/2)$ for
$u\neq0$; for $u\in(0,2\pi)$ the point $u/2$ lies in $(0,\pi)$, where $\operatorname{sinc}$ is positive
and decreasing, so $\operatorname{sinc}^2(u/2)$ is decreasing in $u$ there. Now let $t\in(0,2\pi)$ and
$0<|\theta|\le t$. The function $\theta\mapsto(1-\cos\theta)/\theta^2$ is even, so by monotonicity
$(1-\cos\theta)/\theta^2=(1-\cos|\theta|)/|\theta|^2\ge(1-\cos t)/t^2$; at $\theta=0$ both sides of
the claimed inequality vanish. At $t=\pi$, $(1-\cos\pi)/\pi^2=2/\pi^2$. The upper bound follows from
$1-\cos\theta=2\sin^2(\theta/2)$ and $|\sin x|\le|x|$, which give
$1-\cos\theta\le2(\theta/2)^2=\frac12\theta^2$ for all real $\theta$. Finally, $\sin$ is concave on
$[0,\pi/2]$, so it lies above its chord from $(0,0)$ to $(\pi/2,1)$: $\sin x\ge(2/\pi)x$ there.
Putting $x=\arcsin y\in[0,\pi/2]$ for $y\in[0,1]$, this reads $y\ge(2/\pi)\arcsin y$, that is
$\arcsin y\le(\pi/2)y$; the converse substitution $y=\sin x$ gives the equivalence.

(F4). Let $f(u)=1-u^2/2+u^4/24-\cos u$. Then $f(0)=0$, $f'(u)=-u+u^3/6+\sin u$ gives $f'(0)=0$,
and $f''(u)=\cos u-1+u^2/2$. The derivative of $f''$ is $u-\sin u$, which is $\ge0$ for $u\ge0$;
since $f''(0)=0$, $f''\ge0$ on $[0,\infty)$, and $f''$ is even, so $f''\ge0$ on $\mathbb{R}$. A function
with $f(0)=f'(0)=0$ and $f''\ge0$ is nonnegative, so $1-\cos u\ge u^2/2-u^4/24$. For $u\neq0$ this
gives
\begin{equation*}
\operatorname{sinc}^2(u/2)=\frac{2(1-\cos u)}{u^2}\ge1-\frac{u^2}{12},
\end{equation*}
the equality being the identity of (F3). Putting $u=2y$ gives $\operatorname{sinc}^2y\ge1-y^2/3$
for $y\neq0$; at $y=0$ both sides equal $1$.

(F5). By induction on $n$; for $n\le1$ there is equality. Let $n\ge2$. If
$1-\sum_{i<n}a_i<0$, then $1-\sum_{i\le n}a_i<0$, while the product is $\ge0$ because every factor
is $\ge0$. Otherwise, since $1-a_n\ge0$ and by the induction hypothesis,
\begin{align*}
\prod_{i\le n}(1-a_i)&\ge(1-a_n)\Bigl(1-\sum_{i<n}a_i\Bigr)\\
&=1-\sum_{i\le n}a_i+a_n\sum_{i<n}a_i\ge1-\sum_{i\le n}a_i.
\end{align*}

(F6). For $X\in\mathfrak{g}$, $e^{iX}-\mathbf{1}$ is normal with eigenvalues $e^{i\theta_p}-1$, so
$|e^{iX}-\mathbf{1}|=\max_p|e^{i\theta_p}-1|$; and $|e^{i\theta}-1|=2|\sin(\theta/2)|\le|\theta|$.
Hence $|e^{iX}-\mathbf{1}|\le|X|$, and $|X|\le\mathsf{n}|X|_{\mathfrak{g}}$ is (F1). For
$Z\in M_N(\mathbb{C})$, $e^{Z}-\mathbf{1}=\sum_{n\ge1}Z^n/n!$, and the submultiplicativity
$|XY|\le|X||Y|$ of the operator norm gives $|Z^n|\le|Z|^n$, hence
$|e^{Z}-\mathbf{1}|\le\sum_{n\ge1}|Z|^n/n!=e^{|Z|}-1$. For an orthonormal basis $(e_p)$ of
$\mathbb{C}^N$, $\operatorname{tr}Y=N^{-1}\sum_p\langle e_p,Ye_p\rangle$ and
$|\langle e_p,Ye_p\rangle|\le|Y|$, so $|\operatorname{tr}Y|\le|Y|$. The form
$\langle A,B\rangle=\operatorname{tr}A^{*}B$ is an inner product on $M_N(\mathbb{C})$ with norm
$\|\cdot\|$, and $\operatorname{tr}(XY)=\langle X^{*},Y\rangle$; by the Cauchy--Schwarz inequality
$|\operatorname{tr}(XY)|\le\|X^{*}\|\,\|Y\|$, and $\|X^{*}\|=\|X\|$, since
$\operatorname{tr}(X\,X^{*})=\operatorname{tr}(X^{*}X)$ by cyclicity of the trace. Next,
$X^{*}X\le|X|^2\mathbf{1}$ as Hermitian forms, so
$Y^{*}X^{*}XY\le|X|^2Y^{*}Y$, and taking the trace (which is monotone on Hermitian forms)
$\|XY\|^2\le|X|^2\|Y\|^2$. Applying this to $(XY)^{*}=Y^{*}X^{*}$ and using $\|Z^{*}\|=\|Z\|$ and
$|Y^{*}|=|Y|$ gives $\|XY\|\le|Y|\,\|X\|$. Further,
$\|Y\|^2=N^{-1}\sum_j|Ye_j|^2\le|Y|^2$. If $\sigma_1,\dots,\sigma_N$ are the singular values of
$Y$, then
\begin{equation*}
|Y|^2=\max_i\sigma_i^2\le\sum_i\sigma_i^2=\operatorname{Tr}Y^{*}Y=N\|Y\|^2.
\end{equation*}
Finally, if
$|E|\le\tau<1$, the Neumann series $\sum_{n\ge0}(-E)^n$ converges absolutely in operator norm, its
sum is the inverse of $\mathbf{1}+E$, and its norm is at most $\sum_{n\ge0}\tau^n=(1-\tau)^{-1}$.

(F7). Let $e^{i\theta_p}$, $\theta_p\in\mathopen{]}-\pi,\pi]$, be the eigenvalues of $W$. Since $W$ is
normal,
$W-\mathbf{1}$ is unitarily diagonalisable with eigenvalues $e^{i\theta_p}-1$, and
$|e^{i\theta_p}-1|^2=(\cos\theta_p-1)^2+\sin^2\theta_p=2-2\cos\theta_p$. Hence
\begin{equation*}
\|W-\mathbf{1}\|^2=N^{-1}\sum_p|e^{i\theta_p}-1|^2=N^{-1}\sum_p(2-2\cos\theta_p)
=2(1-\operatorname{Re}\operatorname{tr}W),
\end{equation*}
because $\operatorname{Re}\operatorname{tr}W=N^{-1}\sum_p\cos\theta_p$, while
$|W-\mathbf{1}|^2=\max_p|e^{i\theta_p}-1|^2$. For nonnegative numbers
$\max\le\sum\le N\cdot\max$; dividing by $N$ gives both inequalities.

(F8). Let $v_1,\dots,v_N$ be an orthonormal eigenbasis of $T$, $Tv_p=\theta_pv_p$, and put
$\|T\|_1=\sum_p|\theta_p|$. By the Cauchy--Schwarz inequality
$\|T\|_1\le\sqrt N(\sum_p\theta_p^2)^{1/2}=N|T|_{\mathfrak{g}}$. Since
$\operatorname{Tr}(AT)=\sum_p\theta_p\langle v_p,Av_p\rangle$, we get
$|\operatorname{Tr}(AT)|\le|A|\,\|T\|_1\le N|T|_{\mathfrak{g}}$, and dividing by $N$ gives the first
bound. For the second, write $\|Z\|_{\mathrm{HS}}=(\operatorname{Tr}Z^{*}Z)^{1/2}=\sqrt N\|Z\|$.
The Cauchy--Schwarz inequality of (F6), multiplied by $N$, gives
$|\operatorname{Tr}((AT)(BT'))|\le\|AT\|_{\mathrm{HS}}\|BT'\|_{\mathrm{HS}}$, and by (F6),
$\|AT\|_{\mathrm{HS}}\le|A|\,\|T\|_{\mathrm{HS}}\le\sqrt N|T|_{\mathfrak{g}}$ and likewise
$\|BT'\|_{\mathrm{HS}}\le\sqrt N|T'|_{\mathfrak{g}}$. The product is at most
$N|T|_{\mathfrak{g}}|T'|_{\mathfrak{g}}$; divide by $N$. For the third,
$\operatorname{Tr}(T^2Y)=\operatorname{Tr}(YT^2)=\sum_p\theta_p^2\langle v_p,Yv_p\rangle$, so
$|\operatorname{Tr}(T^2Y)|\le|Y|\operatorname{Tr}T^2$; dividing by $N$ gives
$|\operatorname{tr}(T^2Y)|\le|Y|\operatorname{tr}T^2=|T|_{\mathfrak{g}}^2|Y|$.

(F9). Let $e^{i\theta_p}$, $\theta_p\in\mathopen{]}-\pi,\pi]$, be the eigenvalues of $W$; since
$\det W=1$, $\sum_p\theta_p\in2\pi\mathbb{Z}$. Put $s_p=\sin^2(\theta_p/2)$. Then
$\theta_p/2\equiv-\sum_{q\neq p}\theta_q/2$ modulo $\pi$, and since $|\sin|$ is even and
$\pi$-periodic, $|\sin(\theta_p/2)|=|\sin(\sum_{q\neq p}\theta_q/2)|$. From
\begin{equation*}
|\sin(a+b)|\le|\sin a||\cos b|+|\cos a||\sin b|\le|\sin a|+|\sin b|
\end{equation*}
and induction on the number of terms,
\begin{equation*}
|\sin(\theta_p/2)|\le\sum_{q\neq p}|\sin(\theta_q/2)|
\le(N-1)^{1/2}\Bigl(\sum_{q\neq p}s_q\Bigr)^{1/2},
\end{equation*}
the last step by the Cauchy--Schwarz inequality. Squaring, $s_p\le(N-1)\sum_{q\neq p}s_q$, and
adding $(N-1)s_p$ gives $Ns_p\le(N-1)\sum_qs_q$. As in the proof of (F7),
$|e^{i\theta}-1|^2=4\sin^2(\theta/2)$, so $|W-\mathbf{1}|^2=4\max_ps_p$ and
$\|W-\mathbf{1}\|^2=N^{-1}\sum_q4s_q$. Hence
\begin{equation*}
|W-\mathbf{1}|^2=4\max_ps_p\le(N-1)\,N^{-1}\sum_q4s_q=(N-1)\|W-\mathbf{1}\|^2.
\end{equation*}

(F10). The inequality $|Z|\le\sqrt N\|Z\|$ is the last-but-one assertion of (F6), valid for all
$Z\in M_N(\mathbb{C})$. The matrix $E_{12}$ is traceless, so $E_{12}\in\mathfrak{g}^{c}$. Since
$E_{12}^{*}E_{12}=E_{22}$, the matrix with entry $1$ in position $(2,2)$ and $0$ elsewhere, the only
nonzero singular value of $E_{12}$ is $1$, so $|E_{12}|=1$, and
$\|E_{12}\|^2=\operatorname{tr}E_{22}=N^{-1}$. Thus $|E_{12}|=\sqrt N\|E_{12}\|$. On $\mathfrak{g}$
one has $\|X\|=|X|_{\mathfrak{g}}$, so (F1) reads $|X|\le\mathsf{n}\|X\|$ there; since
$\sqrt N>\mathsf{n}$, this inequality fails at $Z=E_{12}\in\mathfrak{g}^{c}$.
\end{proof}

The right-hand inequality of (F7) gives $1-\operatorname{Re}\operatorname{tr}W\le\frac12|W-\mathbf{1}|^2$
for $W\in U(N)$ with a constant independent of $N$, whereas the left-hand inequality carries the
factor $N^{-1}$; where that lower bound is used together with \cite[(70)]{B-CMP102a} in the
normalisation $\operatorname{tr}\mathbf{1}=1$, the constant $2$ of that display enters as $2N$ and the
constant $\gamma_0$ there as $\gamma_0/N$. The bounds of (F8) carry no constant depending on $N$;
in particular, for $W\in U(N)$, a real $x\ge0$ and $T\in\mathfrak{g}$ with $|T|_{\mathfrak{g}}=1$, the
third bound of (F8) and $|W|=1$ give $x|\operatorname{tr}(T^2W)|\le x$.

\begin{definition}[Haar measure in exponential coordinates and two families of balls]
\label{4.36:def-haar-balls}
Let $G=SU(N)$, $\mathfrak{g}$, $\operatorname{tr}$, the operator norm $|\cdot|$, the norm
$|\cdot|_{\mathfrak{g}}$, $D=N^2-1$ and $\mathsf{n}=(N-1)^{1/2}$ be as in
Notation~\ref{4.36:not-group-norms}. For $v\in\mathfrak{g}$ we write $v_1,\dots,v_N$ for its
eigenvalues, the eigen-angles of $e^{iv}$. Thus $|v|=\max_p|v_p|$,
$|v|_{\mathfrak{g}}^2=N^{-1}\sum_p v_p^2$ and $\sum_p v_p=0$.
\begin{enumerate}
\item[(1)] \emph{Spread and cut locus.} The \emph{spread} of $v\in\mathfrak{g}$ is
$\operatorname{spr}(v):=\max_p v_p-\min_p v_p$. We put
\begin{equation*}
\mathfrak{D}:=\{v\in\mathfrak{g}:\operatorname{spr}(v)<2\pi\},\qquad
\mathrm{Cut}:=\{e^{iv}:v\in\mathfrak{g},\ \operatorname{spr}(v)=2\pi\},
\end{equation*}
and $\exp i\,S:=\{e^{iv}:v\in S\}$ for $S\subset\mathfrak{g}$.
\item[(2)] \emph{The chart.} We use the following statement about $G$. The set $\mathrm{Cut}$
is compact
and Haar-null, and $v\mapsto e^{iv}$ is an analytic diffeomorphism of $\mathfrak{D}$ onto
$G\setminus\mathrm{Cut}$. Its inverse is denoted $\Lambda_G$.

The normaliser $\sigma_N$ below integrates over all of $\mathfrak{D}$ and therefore rests on this
statement: for $N\ge3$, the points of $\mathfrak{D}$ satisfy
only $|v_p|<2\pi(N-1)/N$, and this bound exceeds $\pi$. On a set all of whose points satisfy
$|v|<\pi$, injectivity of $v\mapsto e^{iv}$ and analyticity of its inverse hold directly: there
the principal logarithm \cite[(22)]{B-CMP98} inverts $v\mapsto e^{iv}$ analytically.
\item[(3)] \emph{The density.} For $y\in\mathbb{C}$ put
\begin{equation*}
\Psi_{+}(y):=\sum_{n\ge0}\frac{(iy)^n}{(n+1)!},
\end{equation*}
so that $\Psi_{+}(y)=(e^{iy}-1)/(iy)$ for $y\ne0$ and $\Psi_{+}(0)=1$.
Write $\operatorname{ad}_vY:=[v,Y]$. We put
\begin{equation*}
J(v):=\prod_{p<q}\operatorname{sinc}^2\Big(\frac{v_p-v_q}{2}\Big),\qquad
\operatorname{sinc}y:=\frac{\sin y}{y},\ \operatorname{sinc}0:=1 .
\end{equation*}
We write $dv$ for the Lebesgue measure of the Euclidean space $(\mathfrak{g},\langle X,Y\rangle)$,
where $\langle X,Y\rangle=\operatorname{tr}XY$, identified with $\mathbb{R}^D$. Finally
$\sigma_N:=\big(\int_{\mathfrak{D}}J\,dv\big)^{-1}$.
\item[(4)] \emph{Haar measure.} We use the following statement about the Haar measure of $G$.
The Haar probability measure
$\operatorname{Haar}$ of $G$, restricted to $G\setminus\mathrm{Cut}$, is the image of
$\sigma_N J(v)\,dv$ on $\mathfrak{D}$ under $v\mapsto e^{iv}$. In short,
$\operatorname{Haar}=\sigma_N J(v)\,dv$ on $\mathfrak{D}$.
\item[(5)] \emph{Balls.} For $u\in G$ and $r>0$ put
\begin{equation*}
B^{\mathfrak{g}}_r(u):=\{ue^{iv}:v\in\mathfrak{g},\ |v|_{\mathfrak{g}}<r\},\qquad
B^{\mathrm{op}}_r(u):=\{ue^{iv}:v\in\mathfrak{g},\ |v|<r\}.
\end{equation*}
Let $\omega_D$ be the $dv$-volume of $\{|v|_{\mathfrak{g}}<1\}$, and $\omega^{\mathrm{op}}_N$ the
$dv$-volume of $\{|v|<1\}$. The latter is a definite number once $N$ is fixed.
\end{enumerate}
These objects have the following properties.
\begin{enumerate}
\item[(a)] $J(v)=|\det\Psi_{+}(\operatorname{ad}_v)|$. Moreover $0\le J\le1$, $J(0)=1$,
$J(-v)=J(v)$, and $J$ is continuous.
\item[(b)] $\sigma_N\in(0,\infty)$. Furthermore
$\sigma_2=1/(2\pi^2)$.
\item[(c)] For every $u\in G$, every Borel set $S\subset\mathfrak{D}$ and every Borel function
$\phi\ge0$ on $G$,
\begin{equation}\label{4.36:eq-haar-balls-a}
\int_{u\exp i\,S}\phi\,d\operatorname{Haar}=\sigma_N\int_S\phi(ue^{iv})\,J(v)\,dv .
\end{equation}
\item[(d)] For all $u\in G$ and $r>0$,
\begin{equation*}
B^{\mathfrak{g}}_{r/\mathsf{n}}(u)\subset B^{\mathrm{op}}_r(u)\subset B^{\mathfrak{g}}_r(u).
\end{equation*}
\item[(e)] $\omega_D=\pi^{D/2}/\Gamma(D/2+1)$, and
$\omega^{\mathrm{op}}_N\in[(N-1)^{-D/2}\omega_D,\ \omega_D]$. Furthermore
$\omega^{\mathrm{op}}_2=\omega_3=4\pi/3$.
\item[(f)] Both unit balls $\{|v|_{\mathfrak{g}}<1\}$ and $\{|v|<1\}$ are symmetric under
$v\mapsto-v$. The measure $\operatorname{Haar}$ is both left- and right-invariant.
\end{enumerate}
\end{definition}

\begin{proof}
\emph{(a).} The map $i\operatorname{ad}_v$ sends $\mathfrak{g}$ into $\mathfrak{g}$: for Hermitian
$v,Y$ the matrix $i[v,Y]$ is Hermitian and traceless. Hence $\Psi_{+}(\operatorname{ad}_v)$, a
power series in $i\operatorname{ad}_v$, is a real-linear map of $\mathfrak{g}$. It is the
right-trivialised differential of $v\mapsto e^{iv}$, divided by $i$:
\begin{equation*}
\frac{d}{d\epsilon}\Big|_{\epsilon=0}e^{i(v+\epsilon Y)}e^{-iv}
=i\int_0^1e^{itv}Ye^{-itv}\,dt
=i\int_0^1e^{it\operatorname{ad}_v}Y\,dt=i\,\Psi_{+}(\operatorname{ad}_v)Y .
\end{equation*}
Choose an orthonormal basis of $\mathbb{C}^N$ consisting of eigenvectors of $v$, the $p$-th
with eigenvalue $v_p$, and let $e_{pq}$ be the matrix whose entry in that basis is $1$ at
$(p,q)$ and $0$ elsewhere. Then
$\operatorname{ad}_v e_{pq}=(v_p-v_q)\,e_{pq}$ for $p\ne q$, on the $N^2-N$ root
directions. Moreover $\operatorname{ad}_v$ vanishes on the traceless matrices that are diagonal
in that basis.
These span $\mathfrak{g}^{c}=\mathfrak{sl}(N,\mathbb{C})$. The real determinant of
$\Psi_{+}(\operatorname{ad}_v)$ on $\mathfrak{g}$ equals the complex determinant of its
complex-linear extension to $\mathfrak{g}^{c}$, which is the product of the eigenvalues
$\Psi_{+}(v_p-v_q)$, $p\ne q$, and $\Psi_{+}(0)=1$. For real $x$,
\begin{equation*}
|\Psi_{+}(x)|=\frac{|e^{ix}-1|}{|x|}=\frac{2|\sin(x/2)|}{|x|}=\Big|\operatorname{sinc}\frac{x}{2}\Big|,
\end{equation*}
and both sides equal $1$ at $x=0$. Since $\operatorname{sinc}$ is even, the factors for $(p,q)$
and $(q,p)$ coincide, so
\begin{equation*}
|\det\Psi_{+}(\operatorname{ad}_v)|
=\prod_{p\ne q}\Big|\operatorname{sinc}\frac{v_p-v_q}{2}\Big|=J(v).
\end{equation*}
In particular $J$ does not depend on the numbering of the eigenvalues. The Haar density of
$SU(N)$ in these coordinates is written $\sigma(\cdot)$ in \cite{B-CMP102a}, which states that it can
be calculated explicitly for all classical groups and is analytic, positive, even and invariant.

Since $|\sin y|\le|y|$, every factor of $J$ lies in $[0,1]$, so $0\le J\le1$. At $v=0$ all
$v_p$ vanish and $\operatorname{sinc}0=1$, so $J(0)=1$. The eigenvalues of $-v$ are $-v_p$;
this changes the sign of every difference $v_p-v_q$, and $\operatorname{sinc}^2$ is even, so
$J(-v)=J(v)$. The map $v\mapsto\operatorname{ad}_v$ is linear. The series of $\Psi_{+}$ has
infinite radius of convergence, since $\sum_n R^n/(n+1)!<\infty$ for every $R>0$; hence it
converges in operator norm uniformly on bounded sets of operators. Hence
$v\mapsto\Psi_{+}(\operatorname{ad}_v)$ is continuous, and so is
$J=|\det\Psi_{+}(\operatorname{ad}_v)|$.

\emph{(b).} The largest and the smallest eigenvalue of a Hermitian matrix are $1$-Lipschitz
functions of the matrix in the operator norm, by their variational characterisation. Hence
$\operatorname{spr}$ is continuous and
$\mathfrak{D}$ is open. Since $\sum_qv_q=0$,
\begin{equation*}
N\max_pv_p=\sum_q\big(\max_pv_p-v_q\big)\le(N-1)\operatorname{spr}(v),
\end{equation*}
because the term with $v_q$ maximal vanishes. In the same way
$-N\min_pv_p\le(N-1)\operatorname{spr}(v)$. So every $v\in\mathfrak{D}$ satisfies
$|v|<2\pi(N-1)/N<2\pi$, and $\mathfrak{D}$ is bounded. With $J\le1$ this gives
$\int_{\mathfrak{D}}J\,dv<\infty$.

Conversely, $|v|<\pi$ gives $\operatorname{spr}(v)\le2|v|<2\pi$, so $\{|v|<\pi\}\subset\mathfrak{D}$.
By continuity and $J(0)=1$, we have $J\ge1/2$ on a neighbourhood of $0$ of positive volume, so
$\int_{\mathfrak{D}}J\,dv>0$. Hence $\sigma_N\in(0,\infty)$.

For $N=2$ the eigenvalues of $v$ are $\pm a$ with $a\ge0$, and
$|v|_{\mathfrak{g}}^2=(a^2+a^2)/2=a^2$. So $a=|v|_{\mathfrak{g}}=|v|$,
$\operatorname{spr}(v)=2|v|_{\mathfrak{g}}$, and $J(v)=\operatorname{sinc}^2|v|_{\mathfrak{g}}$.
Therefore $\mathfrak{D}=\{|v|_{\mathfrak{g}}<\pi\}$ in the three-dimensional Euclidean space
$(\mathfrak{g},\langle\cdot,\cdot\rangle)$. Polar coordinates give
\begin{equation*}
\int_{\mathfrak{D}}J\,dv=4\pi\int_0^\pi\frac{\sin^2t}{t^2}\,t^2\,dt
=4\pi\int_0^\pi\sin^2t\,dt=2\pi^2,
\end{equation*}
that is, $\sigma_2=1/(2\pi^2)$. This is the value $\sigma(0)$ of the density of
\cite{B-CMP102a}.

\emph{(c).} By item (2), $v\mapsto e^{iv}$ is a homeomorphism of $\mathfrak{D}$ onto the open set
$G\setminus\mathrm{Cut}$. Hence $\exp i\,S$ is Borel in $G$, and
$\exp i\,S\subset G\setminus\mathrm{Cut}$.
The image of $\operatorname{Haar}$ under $y\mapsto uy$ is $\operatorname{Haar}$ (left invariance).
Therefore
\begin{equation*}
\int_{u\exp i\,S}\phi\,d\operatorname{Haar}
=\int_G\mathbf{1}_{\exp i\,S}(y)\,\phi(uy)\,d\operatorname{Haar}(y).
\end{equation*}
The integrand vanishes off $G\setminus\mathrm{Cut}$, so item (4) applies to it. Since
$v\mapsto e^{iv}$ is injective on $\mathfrak{D}$, we have
$\mathbf{1}_{\exp i\,S}(e^{iv})=\mathbf{1}_S(v)$
for $v\in\mathfrak{D}$. The right-hand side thus equals
$\sigma_N\int_S\phi(ue^{iv})J(v)\,dv$, which is \eqref{4.36:eq-haar-balls-a}.

\emph{(d).} By inequality (F1) of Lemma~\ref{4.36:lem-eigenangle-inequalities}, see
\eqref{4.36:eq-eigenangle-inequalities-a}, every $X\in\mathfrak{g}$ satisfies
$|X|_{\mathfrak{g}}\le|X|\le\mathsf{n}|X|_{\mathfrak{g}}$. If $|v|_{\mathfrak{g}}<r/\mathsf{n}$,
then $|v|\le\mathsf{n}|v|_{\mathfrak{g}}<r$. If $|v|<r$, then $|v|_{\mathfrak{g}}\le|v|<r$. Hence
\begin{equation*}
\{|v|_{\mathfrak{g}}<r/\mathsf{n}\}\subset\{|v|<r\}\subset\{|v|_{\mathfrak{g}}<r\},
\end{equation*}
and applying $v\mapsto ue^{iv}$ gives the asserted inclusions.

\emph{(e).} Since $dv$ is the Lebesgue measure of a $D$-dimensional Euclidean space and
$\{|v|_{\mathfrak{g}}<1\}$ is its unit ball, $\omega_D=\pi^{D/2}/\Gamma(D/2+1)$. The inclusions
displayed in (d) with $r=1$, together with the scaling of Lebesgue measure, give
$\mathsf{n}^{-D}\omega_D\le\omega^{\mathrm{op}}_N\le\omega_D$. Moreover
$\mathsf{n}^{-D}=(N-1)^{-D/2}$. For $N=2$ we have $\mathsf{n}=1$ and $D=3$; indeed
$|v|=|v|_{\mathfrak{g}}$ for $N=2$, as computed in the proof of (b). So
\begin{equation*}
\omega^{\mathrm{op}}_2=\omega_3=\frac{\pi^{3/2}}{\Gamma(5/2)}
=\frac{\pi^{3/2}}{\tfrac34\sqrt{\pi}}=\frac{4\pi}{3}.
\end{equation*}

\emph{(f).} Both norms satisfy $|-v|_{\mathfrak{g}}=|v|_{\mathfrak{g}}$ and $|-v|=|v|$, so both
unit balls are symmetric under $v\mapsto-v$. Finally, $G$ is compact, hence unimodular, so its
Haar measure is left- and right-invariant.
\end{proof}

\begin{definition}[Scales, running couplings, flow inequalities and thresholds]
\label{4.36:not-flow-thresholds}
Fix a level $k$ with $0\le k\le K$. We write $\eta:=L^{-k}$ and, for $n\le k$,
\[
\ell_n:=L^{n}\eta=L^{-(k-n)},
\]
the unit of scale $n$, so that $\eta\le\ell_n$. For every level $j$ we write $x_j:=g_j^{-2}$ and
$\ell\mathrm{og}_j:=\log x_j$; the abbreviation $\ell_j$ for $\ell\mathrm{og}_j$ is used only where
no scale unit $\ell_n$ occurs.

\emph{Functions of the coupling.} With the auxiliary parameters $A_0>0$, $A_1$, $C_0$, $C_1$, $q_0$,
$q_1$ of $\mathfrak{p}$, the positive integer $p_0$ and the exponent $r$, we use the degree
function $p_0(g):=A_0(\log g^{-2})^{p_0}$ and, as in \cite[(2.4), (2.5), (2.28)]{B-CMP119},
\[
\varepsilon_j:=g_j\,p_0(g_j),\qquad \delta_j:=\frac{A_1}{A_0}\,\varepsilon_j,\qquad
\alpha_{0,j}:=g_jC_0(\log x_j)^{q_0},\qquad \alpha_{1,j}:=g_jC_1(\log x_j)^{q_1},
\]
and $R_j$ denotes the least power of $L$ not smaller than $(\log x_j)^{r}$. As in the hypothesis
of \cite[Theorem~1]{B-CMP122b}, the running couplings lie in an interval $(0,\gamma]$ with
$\gamma<1$, so $x_j>1$ and $\log x_j>0$ for every $j$.

\emph{The flow inequalities.} By Theorem~\ref{4.1:thm-clause-zero}, which is clause (0) of
Theorem~0, for $K=K(g_0,g)$ the inverse couplings satisfy
\begin{equation}\label{4.36:eq-flow-thresholds-1}
\lambda(j-i)-c_5\le x_i-x_j\le B(j-i),\qquad 0\le i\le j\le K,
\end{equation}
with $B=B^{\sharp}(N,L,\mathfrak{p})$, $c_5=c_5(L,N)$ and
$\lambda\ge\lambda^{\sharp}>0$, none of these constants depending on $K$. Here
$\lambda^{\sharp}=c_0/2-\bar{\varsigma}$ and $B^{\sharp}=C_\beta+\bar{\varsigma}$, where
$c_0(N)=(11N^2/12\pi^2)\log L$ is the one-loop constant of Definition~\ref{1.2:def-flow-constants}
and $\bar{\varsigma}$, the margin by which $\lambda^{\sharp}$ falls below $c_0/2$ and $B^{\sharp}$
exceeds $C_\beta$, is a number with $\bar{\varsigma}\le c_0/4$ and $\bar{\varsigma}\le 1$;
hence $\lambda^{\sharp}\ge c_0/4$ and $B^{\sharp}\le C_\beta+1$. The first-passage index
$K=K(g_0,g)$ satisfies $g^{-2}<x_K\le g^{-2}+B^{\sharp}$, that is $g_K\in[g_-(g),g)$ (the anchor of
the flow, Theorem~\ref{4.1:thm-clause-zero}). For $1\le j\le K$, $\beta_j=x_{j-1}-x_j$ are the
coupling increments of Definition~\ref{1.2:def-couplings}. Let $\bar{B}=C_\beta+1$ be the constant
of Definition~\ref{1.2:def-flow-constants}, and put
\[
\bar{B}^{+}:=\max\{\bar{B},\,c_5\}=\max\{C_\beta+1,\,c_5\},
\]
a number fixed once $(N,L,\mathfrak{p})$ are fixed, independent of $\gamma$ and of $K$. The
following hold (proved below):
\begin{enumerate}
\item[(a)] $\beta_j\in[\lambda-c_5,B]$ for $1\le j\le K$, and
\[
\bar{B}^{+}\ge\max\{B,c_5\}\ge\sup_j|\beta_j|\ge0;
\]
\item[(b)] $K(g_0,g)$ is unbounded as $g_0\downarrow0$, $B\ge\lambda$, and
\begin{equation}\label{4.36:eq-flow-thresholds-2}
\bar{B}^{+}\ge B\ge\lambda\ge\lambda^{\sharp}\ge\frac{c_0(N)}{4}=\frac{11N^2}{48\pi^2}\log L;
\end{equation}
\item[(c)] with $s:=K-k$ the depth of the level $k$, and all indices in $\{0,\dots,K\}$,
\begin{equation}\label{4.36:eq-flow-thresholds-a}
\begin{aligned}
&(\mathrm{Fl}1)\quad g^{-2}+\lambda s-c_5\le x_k\le x_K+Bs\le g^{-2}+B^{\sharp}(s+1)\\
&\phantom{(\mathrm{Fl}1)}\quad\ \text{at }K=K(g_0,g);\\
&(\mathrm{Fl}2)\quad x_n\le x_k+B(k-n)\ \ (n\le k),\qquad x_n\le x_j+c_5\ \ (j\le n);\\
&(\mathrm{Fl}3)\quad g_n\le g_j\bigl(1+B(n-j)\bigr)^{1/2}\ \ (j\le n),\\
&\phantom{(\mathrm{Fl}3)}\quad\ \delta_n\ge\tfrac12\,\delta_k\bigl(1+B(k-n)\bigr)^{-1/2}
\ \ (n\le k);\\
&(\mathrm{Fl}4)\quad \sum_{0\le j\le k}g_j^{\kappa_0}\le C_\lambda\,x_k^{1-\kappa_0/2},
\end{aligned}
\end{equation}
where $C_\lambda$ depends only on $\lambda$, $c_5$ and $\kappa_0$. The second inequality of
$(\mathrm{Fl}3)$ is proved below at every level $k$ at which
\[
x_k>1+c_5\quad\text{and}\quad\bigl(\log(x_k-c_5)\bigr)^{p_0}\ge\tfrac12\,(\log x_k)^{p_0}.
\]
Since $\log(x-c_5)/\log x\to1$ as $x\to\infty$, there is a number
$x_{\delta}=x_{\delta}(c_5,p_0)\ge1+c_5$ such that this floor holds whenever $x_k\ge x_{\delta}$;
the threshold $x_{\mathrm{th}}$ is taken with $x_{\mathrm{th}}\ge x_{\delta}$, so that
under the threshold $x_k\ge x_{\mathrm{th}}$ the second inequality of $(\mathrm{Fl}3)$ holds at
every level $k\le K$;
\item[(d)] let $\gamma'>0$ be any number with $\gamma'^{-2}\ge x_{\mathrm{th}}+c_5$; if
$g\le\gamma'$ and $K=K(g_0,g)$, then $x_k\ge x_{\mathrm{th}}(N,L,M,\mathfrak{p})$ for all $k\le K$.
\end{enumerate}
The threshold $x_{\mathrm{th}}$ and the ceiling $\gamma_N$ are fixed with the transport bounds for
every $L$ under a coupling ceiling, \eqref{4.36:eq-transport-every-l-a}, and there
$x_{\mathrm{th}}\ge\bar{B}^{+}$ and
$\gamma_N^{-2}\ge x_{\mathrm{th}}+c_5$, so that (d) applies with $\gamma'=\gamma_N$. No lower bound
on $L$ of the form $L^2\ge1+\bar{B}^{+}$ is imposed among the thresholds below; by
\eqref{4.36:eq-flow-thresholds-2} such a condition would read
\[
L^2\ge1+\bar{B}^{+}\ge1+\frac{11N^2}{48\pi^2}\log L
\]
and would force a lower bound on $L$ that grows with $N$. Under the ceiling, where
$x_{\mathrm{th}}\ge\bar{B}^{+}$, every estimate proved in this chapter into which such a condition
would enter holds with its constant unchanged, at every $N$ and for every value of $L$; this is
proved together with the bound \eqref{4.36:eq-transport-every-l-a}. One further estimate into
which such a condition would enter is not proved in this chapter but belongs to the inputs that
the chapter takes by statement; the corresponding assertion, printed beside the bound
\eqref{4.36:eq-transport-every-l-a}, is stated; proof incomplete at the step marked $(\ast)$ there.
Where couplings are
transported between levels under the ceiling, the bound \eqref{4.36:eq-transport-every-l-a} is used
instead: if $x_n\ge\bar{B}^{+}$ for all $n\le K$, it gives $g_n\le(1+(n-j))^{1/2}g_j$ for $j\le n$,
with no dependence on $N$ or on $B$, for every value of $L$.

\emph{Conditions on the exponents.} Put $q:=\min\{q_0,q_1\}$. We use the conditions
\begin{equation}\label{4.36:eq-flow-thresholds-3}
q\ge p_0+3r+2,\qquad p_0\ge5r,\qquad r\ge2,\qquad
\kappa\ge\max\{10,8c_d+8\},\qquad \kappa_0\ge7;
\end{equation}
these are imposed together with the conditions $\kappa\ge10\log L$ and $\kappa_0\ge17$ on the
auxiliary parameters fixed with the quantifier prefix (Notation~\ref{2.1:not-prefix}). We also use
the lower torus window, the one-sided lower bound on the torus exponent $m$ by the function
$m_{\mathbb{T}}$ of Definition~\ref{1.2:def-flow-constants}, in the form fixed there.

\emph{Thresholds.} With the constants listed further below, and with $\mathsf{r}_0>0$ a radius
subject to $(\mathrm{t}8)$, we use
\begin{equation}\label{4.36:eq-flow-thresholds-b}
\begin{aligned}
&(\mathrm{t}1)\quad p_0(g_k)\ge8C_{\mathrm{loc}};\\
&(\mathrm{t}2)\quad \mathsf{R}:=\min\Bigl\{\frac{c_{\star}\delta_k}{4C'_{\chi}C_H},\,
\frac{\varepsilon_k}{16C_{\mathrm{loc}}},\,\frac{\mathsf{r}_0}{2}\Bigr\},\qquad
r_i:=\min\Bigl\{\mathsf{r}_0,\,\frac{c_B\delta_i}{8C_H}\Bigr\};\\
&(\mathrm{t}3)',\ (\mathrm{t}4)'\quad\text{the conditions stated after this display};\\
&(\mathrm{t}5)\quad \kappa\ge\max\{10,8c_d+8\};\\
&(\mathrm{t}6)\quad L^2\ge20(1+\beta^{\circ}_0);\\
&(\mathrm{t}8)\quad \mathsf{r}_0\le\min\bigl\{\tfrac12,\,C^{(102)}_1a'_1/4,\,(4C_H)^{-1},
\,(2C_P)^{-1}\bigr\};\\
&(\mathrm{t}9)\quad R_j\ge L^2;\\
&x_k\ge x_{\mathrm{th}}(N,L,M,\mathfrak{p})\ \ (0\le k\le K);\qquad s\ge s_*.
\end{aligned}
\end{equation}
Here $(\mathrm{t}3)'$ comprises the absorption conditions $|\partial^2F^{E}|\le x_k/8$, the bound
$\tfrac12$ for each of the two sums controlled by $\mathsf{R}$, the bound $1$ for the corresponding
sum of the $F^{B}$ estimate, and $\eta_{\mathrm{far}}\le1$, each required when $\mathsf{R}$ and
$r_i$ are replaced by any one of the members of the minima in $(\mathrm{t}2)$; and $(\mathrm{t}4)'$
is the floor on $M$ under which \eqref{4.36:eq-flow-thresholds-c} below holds. In $(\mathrm{t}2)$,
$r_i$ is defined for every level $i$ and is used at the working scale $i=m(j,n)$ defined below. In
$(\mathrm{t}6)$, $\beta^{\circ}_0$ is a constant fixed where it first occurs, distinct from the
increments $\beta_j$. In $(\mathrm{t}8)$, $C^{(102)}_1$ is the constant of \cite[(172)]{B-CMP102b}
and $a_0,a_1,a'_1$ are the constants of \cite[Theorem~1]{B-CMP102b}.

\emph{Objects at a marked plaquette.} Let $p$ be a plaquette of the lattice $\mathbb{T}^{(k)}$ of
level $k$ (Definition~\ref{1.5:def-block-average}), $b_1\subset\partial p$ a marked bond, and
$u_1:=V_k(b_1)$ the free link; we write $V=(u_1,V_{\mathrm{off}})$, where $V_{\mathrm{off}}$ collects
the values of $V_k$ on the other bonds. Put
\[
R_*:=20L^2M\max\{R_k,R_{k+1}\}.
\]
A history $h$ is \emph{clean} if $\operatorname{dist}_k(p,\Lambda_k(h)^{c})\ge R_*$, where
$\Lambda_k(h)$ is the region that $h$ carries at level $k$ and $\operatorname{dist}_k$ is the
distance in units of level $k$. Let $\mathcal{N}$ be the set of cubes $\square$ of side $LM_2R_k$
with $\partial p\cap\tilde{\square}^4\neq\emptyset$, where $\tilde{\square}$, $\tilde{\square}^4$,
$\tilde{\square}^5$ are the enlargements of $\square$ used in Ba{\l}aban's construction of the local
problems \cite[(2.16)]{B-CMP119}; $\mathcal{N}$
has at most $9^4$ members. For $\square\in\mathcal{N}$, $U_{k,\square}$ is the minimiser, in the
sense of \cite{B-CMP102b}, of the local variational problem of \cite[(2.16)]{B-CMP119} on
$\tilde{\square}^4$; its data
$\mathfrak{d}_{\square}(V)$ are an explicit map, built from finitely many products, of the
restriction of $V$ to $\tilde{\square}^4$. We put
\[
\chi_{\square}(V):=\mathbf{1}\Bigl\{\sup_{p''\subset\tilde{\square}}
\bigl|U_{k,\square}(V,\partial p'')-\mathbf{1}\bigr|<\varepsilon_k\eta^2\Bigr\},
\]
the \emph{window} of the free link
\[
W(V_{\mathrm{off}}):=\Bigl\{u\in G:\ \prod_{\square\in\mathcal{N}}
\chi_{\square}(u,V_{\mathrm{off}})=1\Bigr\},
\]
an open subset of $G$, and the \emph{deep set}
\[
\mathcal{D}:=\Bigl\{V:\ \sup_{\square\in\mathcal{N},\ p''\subset\tilde{\square}}
\bigl|U_{k,\square}(V,\partial p'')-\mathbf{1}\bigr|\le\tfrac12\,\varepsilon_k\eta^2\Bigr\}.
\]
For a clean history $h$, $\mathbf{1}_{G(h)}$ is its good-set factor, a Borel function of $V$ with
values in $[0,1]$. $C_{\mathrm{loc}}$ is the constant of the local lemma on the minimisers
$U_{k,\square}$; by \cite{B-CMP102b} it depends on $d$ and $L$ only, at the
generality $G\subset U(N)$, and it is used at its value for $G=SU(N)$. We put
$\varrho:=\varepsilon_k/(4C_{\mathrm{loc}})$, a local letter used only for the window of the free
link.

\emph{The conditional exponent and the block geometry.} $F^{W}$, $F^{E}$, $F^{R}$, $F^{B}$,
$F^{\mathrm{far}}$ are the members of the conditional exponent, with
$F^{\mathrm{far}}:=Q_h+\sum_Y\Psi_Y$; for a plaquette $q$, $w_q:=1-\operatorname{Re}\operatorname{tr}
U_k(\partial q)$. Let $(\Omega_n)_{n\le k}$ and $(\Lambda_j)_{j\le k}$ be the two families of
regions carried by the history $h$ in Ba{\l}aban's construction, at the levels $n$ and $j$
respectively; at $j=k$, $\Lambda_k=\Lambda_k(h)$ is the region from whose complement the distance
in the definition of a clean history is measured. We put
$\Gamma_n:=\Omega_n\setminus\Omega_{n+1}$ for $n<k$, $\Gamma_k:=\Omega_k$, and
$Z_j:=\Lambda_j^{c}$. For
$y\in\Gamma_n$, $\Delta(y)$ is the block of $y$, a cube of side $\ell_n$; $\pi_n$ denotes the
$M$-cubes of scale $n$. The working scale is $m(j,n):=\max\{j,n-2\}$, always written with its
arguments to distinguish it from the torus exponent $m$, and $\square_0:=\tilde{\square}^5$;
$d_j(X)$ is as in Notation~\ref{0.2:not-glossary}. We write $d(\cdot,\cdot)$
for the multiscale distance of \cite[(2.46)]{B-CMP96}, $\delta_0>0$ for the decay constant of
Ba{\l}aban's estimates read at $G=SU(N)$, and $e_y:=e^{-\delta_0d(y,b_1)/8}\le1$. For a set
$\mathsf{N}$ of bonds, $\mathfrak{b}(\mathsf{N})$ is the set of blocks meeting $\mathsf{N}$,
\[
\hat{e}_{\mathsf{N}}:=\sum_{y\in\mathfrak{b}(\mathsf{N})}e_y^{1/2},\qquad
\check{e}_{\mathsf{N}}:=\max_{y\in\mathfrak{b}(\mathsf{N})}e_y,
\]
and $\mathsf{N}^{+}$ is the set of blocks at adjacency at most $2$ from $\mathfrak{b}(\mathsf{N})$.
Under the floor $(\mathrm{t}4)'$ on $M$, the stratum-sum bound, a geometric consequence of
\cite[Lemma~2.1]{B-CMP96}, reads, with a constant $C_\sigma$,
\begin{equation}\label{4.36:eq-flow-thresholds-c}
\sigma^{(\alpha)}_n:=\sum_{y\in\Gamma_n}e^{-\alpha\delta_0d(b_1,y)}\le C_\sigma e^{-(k-n)},
\qquad \alpha\in\bigl[\tfrac{1}{288},\tfrac18\bigr].
\end{equation}
This bound is used in this chapter by statement; its proof is given with the estimates on the
strata $\Gamma_n$.

\emph{Constants.} The constants $C_H$, $C_P$, $c_{\star}$, $E'_1$ and $R'_1$ are constants of
Ba{\l}aban's papers, among those listed in the next sentence; the constants $C'_W$, $C_T$,
$C'_{\chi}$, $C_B$, $c_B$, $C_q$, $\beta^{\circ}_0$, $c_d$ and $\eta_{\mathrm{far}}$ are fixed in the
statements where they first occur. A statement holds \emph{at the $SU(N)$ values of its inputs} if
every constant of Ba{\l}aban's papers entering it ($E_0$, $E'_1$, $R'_1$, $B_0$, $a_0$, $a_1$,
$a'_1$, $B_3$, $B_5$,
$C^{(102)}_1$, $C_H$, $C_P$, $C_{\mathrm{loc}}$, $c_{\star}$, $\delta_0$, $c^{(96)}_0$,
$C_{\mathrm{av}}$,
\dots) and every flow constant ($\bar{B}^{+}$, $c_5$, $C_\lambda$) is read at $G=SU(N)$; here
$c^{(96)}_0$ is the absolute constant of \cite{B-CMP96} entering $(\mathrm{t}4)'$, distinct from the
one-loop constant $c_0(N)$. That these constants exist at $G=SU(N)$ with the dependences stated is
the content of clause (i-a) of Theorem~0 in its form for $G=SU(N)$, together with
\eqref{4.36:eq-flow-thresholds-1}.
\end{definition}

\begin{proof}[Proof of \textup{(a)}--\textup{(d)} in Notation~\ref{4.36:not-flow-thresholds}]
All indices below lie in $\{0,\dots,K\}$, and we use throughout that $x_j>1$ for every $j$.

\emph{(a).} Applying \eqref{4.36:eq-flow-thresholds-1} with $(i,j)=(j-1,j)$ gives
\[
\lambda-c_5\le x_{j-1}-x_j=\beta_j\le B.
\]
If $\beta_j\ge0$ then $|\beta_j|=\beta_j\le B$; if $\beta_j<0$ then, since $\lambda>0$,
\[
|\beta_j|=-\beta_j\le c_5-\lambda<c_5.
\]
Hence $\sup_j|\beta_j|\le\max\{B,c_5\}$. Further
\[
B=B^{\sharp}=C_\beta+\bar{\varsigma}\le C_\beta+1=\bar{B}\le\bar{B}^{+}
\]
and $c_5\le\bar{B}^{+}$ by the definition of $\bar{B}^{+}$, so $\bar{B}^{+}\ge\max\{B,c_5\}$.

\emph{(b).} Fix $g$ and let $K=K(g_0,g)$. Since $x_0=g_0^{-2}$ and $x_K\le g^{-2}+B^{\sharp}$ by the
anchor, the upper inequality of \eqref{4.36:eq-flow-thresholds-1} with $(i,j)=(0,K)$ gives
\[
g_0^{-2}-g^{-2}-B^{\sharp}\le x_0-x_K\le BK.
\]
As $g_0\downarrow0$ the left side tends to $+\infty$ while $B$ is fixed, so $K(g_0,g)$ is unbounded.
The lower inequality with $(i,j)=(0,K)$ gives
\[
\lambda K-c_5\le x_0-x_K\le BK,
\]
that is $B\ge\lambda-c_5/K$. Since $\lambda$, $B$ and $c_5$ do not depend on $K$ and $K$ takes
arbitrarily large values, $B\ge\lambda$. Combining this with the bounds $\lambda\ge\lambda^{\sharp}$
and $\lambda^{\sharp}\ge c_0(N)/4$, the identity $c_0(N)/4=(11N^2/48\pi^2)\log L$, and
$\bar{B}^{+}\ge B$ from (a) gives \eqref{4.36:eq-flow-thresholds-2}.

\emph{(Fl1).} Let $K=K(g_0,g)$. The lower inequality of \eqref{4.36:eq-flow-thresholds-1} with
$(i,j)=(k,K)$ gives
\[
x_k\ge x_K+\lambda(K-k)-c_5=x_K+\lambda s-c_5,
\]
and $x_K>g^{-2}$ by the anchor; this is the lower
bound. The upper inequality with $(i,j)=(k,K)$ gives $x_k\le x_K+Bs$. Finally
$x_K\le g^{-2}+B^{\sharp}$ by the anchor and $B=B^{\sharp}$, so
\[
x_K+Bs\le g^{-2}+B^{\sharp}+B^{\sharp}s=g^{-2}+B^{\sharp}(s+1).
\]

\emph{(Fl2).} For $n\le k$ the upper inequality with $(i,j)=(n,k)$ reads $x_n-x_k\le B(k-n)$. For
$j\le n$ the lower inequality with $(i,j)=(j,n)$ reads $x_j-x_n\ge\lambda(n-j)-c_5$, and
$\lambda(n-j)\ge0$, so $x_n\le x_j+c_5$.

\emph{(Fl3).} Let $j\le n$. By the first inequality of (Fl2), applied with $(n,k)$ replaced by
$(j,n)$, and since $x_n>1$ and $B\ge\lambda>0$ by (b),
\[
x_j\le x_n+B(n-j)\le x_n\bigl(1+B(n-j)\bigr).
\]
Taking this to the power $-1/2$ gives $g_j\ge g_n(1+B(n-j))^{-1/2}$, which is the first inequality.
For the second, let $n\le k$ and let $x_k$ satisfy the floor displayed in (c), as it does under the
threshold $x_k\ge x_{\mathrm{th}}$. By definition
\[
\delta_i=\frac{A_1}{A_0}\,g_iA_0(\log g_i^{-2})^{p_0}=A_1g_i(\log x_i)^{p_0},
\qquad\text{so}\qquad
\frac{\delta_n}{\delta_k}=\frac{g_n}{g_k}\Bigl(\frac{\log x_n}{\log x_k}\Bigr)^{p_0}.
\]
By the first half of (Fl3) with $(j,n)$ replaced by $(n,k)$, $g_n/g_k\ge(1+B(k-n))^{-1/2}$. By the
lower inequality of \eqref{4.36:eq-flow-thresholds-1} with $(i,j)=(n,k)$, by $\lambda(k-n)\ge0$ and
by the floor,
\[
x_n\ge x_k+\lambda(k-n)-c_5\ge x_k-c_5>1,
\]
so $\log x_n\ge\log(x_k-c_5)>0$, and by the floor the last factor is at least $\tfrac12$. Hence
$\delta_n\ge\tfrac12\,\delta_k(1+B(k-n))^{-1/2}$.

\emph{(Fl4).} Let $j\le k$ and put $a:=x_k+\lambda(k-j)$. The lower inequality of
\eqref{4.36:eq-flow-thresholds-1} with $(i,j)$ replaced by $(j,k)$ gives $x_j\ge a-c_5$, and
$x_j>1$; hence
\[
x_j\ge\max\{1,a-c_5\}\ge\frac{a}{1+c_5}.
\]
Indeed, if $a\le1+c_5$ then $1\ge a/(1+c_5)$, and if $a>1+c_5$ then
\[
a-c_5-\frac{a}{1+c_5}=\frac{c_5\,(a-1-c_5)}{1+c_5}\ge0.
\]
Therefore
\[
g_j^{\kappa_0}=x_j^{-\kappa_0/2}\le(1+c_5)^{\kappa_0/2}\bigl(x_k+\lambda(k-j)\bigr)^{-\kappa_0/2}.
\]
The function $t\mapsto(x_k+\lambda t)^{-\kappa_0/2}$ is decreasing on $[0,\infty)$ and, since
$\kappa_0\ge7$ by \eqref{4.36:eq-flow-thresholds-3}, so that $\kappa_0>2$, integrable there;
bounding the term with $t=k-j\ge1$ by the integral over $[t-1,t]$ gives
\begin{align*}
\sum_{0\le j\le k}g_j^{\kappa_0}
&\le(1+c_5)^{\kappa_0/2}\Bigl(x_k^{-\kappa_0/2}
+\int_0^\infty(x_k+\lambda t)^{-\kappa_0/2}\,dt\Bigr)\\
&=(1+c_5)^{\kappa_0/2}\Bigl(x_k^{-\kappa_0/2}+\frac{2\,x_k^{1-\kappa_0/2}}{\lambda(\kappa_0-2)}\Bigr)
\le C_\lambda\,x_k^{1-\kappa_0/2},
\end{align*}
with $C_\lambda:=(1+c_5)^{\kappa_0/2}\bigl(1+2/(\lambda(\kappa_0-2))\bigr)$, where in the last step
$x_k^{-\kappa_0/2}\le x_k^{1-\kappa_0/2}$ because $x_k>1$. This constant depends only on $\lambda$,
$c_5$ and $\kappa_0$.

\emph{(d).} Let $g\le\gamma'$, $K=K(g_0,g)$ and $k\le K$. The lower bound in (Fl1), the inequality
$\lambda s\ge0$ and the assumption $\gamma'^{-2}\ge x_{\mathrm{th}}+c_5$ give
\[
x_k\ge g^{-2}-c_5\ge\gamma'^{-2}-c_5\ge x_{\mathrm{th}}.
\]
\end{proof}

\begin{lemma}[Arithmetic of the transport factors]\label{4.36:lem-transport-arithmetic}
Let $\bar{B}\ge 0$ be real, and let $L\ge 2$, $m\ge 1$ and $u\ge 1$ be integers. Then:
\begin{enumerate}
\item[(i)] $1+\bar{B}m\le(1+\bar{B})m$;
\item[(ii)] $m^{1/2}\le 2^{m-1}\le L^{m-1}$;
\item[(iii)] the chain
\begin{equation*}
(1+\bar{B}m)^{1/2}L^{-m}\le\frac{(1+\bar{B})^{1/2}}{L}\le(1+\bar{B})^{1/2}
\end{equation*}
holds;
\item[(iv)] $(1+\bar{B}u)^{1/2}L^{-2u}\le(1+\bar{B})^{1/2}/L^{2}$;
\item[(v)] for every integer $m\ge 0$, $(1+m)^{1/2}L^{-m}\le 1$ and $(1+m)^{1/2}L^{-2m}\le 1$.
\end{enumerate}
\end{lemma}

\begin{proof}
(i) Expanding, $(1+\bar{B})m-(1+\bar{B}m)=m-1$, and $m-1\ge 0$ because $m\ge 1$.

(ii) We show $m\le 4^{m-1}$ for all integers $m\ge1$ by induction on $m$. For $m=1$ both sides equal $1$. If $m\le 4^{m-1}$ for some $m\ge 1$, then, since $1\le m$,
\begin{equation*}
m+1\le 2m\le 2\cdot 4^{m-1}\le 4^{m},
\end{equation*}
which is the claim for $m+1$. Taking square roots gives $m^{1/2}\le 2^{m-1}$. Since $2\le L$ and $m-1\ge 0$, we have $2^{m-1}\le L^{m-1}$.

(iii) Since $\bar{B}\ge 0$, both sides of (i) are non-negative, and its square root together with (ii) gives
\begin{equation*}
(1+\bar{B}m)^{1/2}\le(1+\bar{B})^{1/2}m^{1/2}\le(1+\bar{B})^{1/2}L^{m-1}.
\end{equation*}
Dividing by $L^{m}>0$ yields the first inequality of (iii); the second holds because $L\ge 1$.

(iv) By (ii) applied with $u$ in place of $m$, and since $u-1\ge 0$ and $L\ge 1$, we have $u^{1/2}\le L^{u-1}\le L^{2u-2}$. As in (iii), (i) with $u$ in place of $m$ then gives
\begin{equation*}
(1+\bar{B}u)^{1/2}\le(1+\bar{B})^{1/2}u^{1/2}\le(1+\bar{B})^{1/2}L^{2u-2},
\end{equation*}
and dividing by $L^{2u}$ yields (iv).

(v) For $m=0$ both left-hand sides equal $1$. Let $m\ge 1$. Then $1+m\le 2^{m}$: for $m=1$ both sides equal $2$, and $1+m\le 2^{m}$ implies $m+2\le 2(m+1)\le 2^{m+1}$. Hence, since $2^{m}\ge 1$ and $2\le L$,
\begin{equation*}
(1+m)^{1/2}\le 2^{m/2}\le 2^{m}\le L^{m},
\end{equation*}
which is the first inequality of (v). The second follows from the first, because $L^{-2m}\le L^{-m}$ for $L\ge 1$ and $m\ge 0$.
\end{proof}

By (iii), the bound $(1+\bar{B}m)^{1/2}L^{-m}\le 1$ holds for every integer $m\ge 1$ as soon as it holds for $m=1$, that is, as soon as $L^{2}\ge 1+\bar{B}$; and by (v) its case $\bar{B}=1$ holds for every integer $m\ge 0$ and every integer $L\ge 2$ without any further condition.

\begin{lemma}[Comparison of powers of logarithms under a bounded shift]\label{4.36:lem-log-comparison}
Let $c_5\ge 0$, let $p_0\ge 1$ be an integer and let $\beta_0\in(0,1]$. Put
\begin{equation*}
x_{\delta}:=\max\{e,\,2c_5,\,4p_0c_5\},\qquad x_{\beta}:=\max\{e,\,2p_0c_5/\beta_0\}.
\end{equation*}
Let $x$ and $y$ be real numbers.
\begin{enumerate}
\item[(i)] If $x\ge x_{\delta}$ and $y\ge x-c_5$, then
\begin{equation}\label{4.36:eq-log-comparison-1}
(\log y)^{p_0}\ge \tfrac12(\log x)^{p_0}.
\end{equation}
\item[(ii)] If $x\ge x_{\beta}$ and $1\le y\le x+c_5$, then
\begin{equation}\label{4.36:eq-log-comparison-2}
(\log y)^{p_0}\le (1+\beta_0)(\log x)^{p_0}.
\end{equation}
\end{enumerate}
\end{lemma}

\begin{proof}
In both parts $x\ge e$, hence $\log x\ge 1$; and $t\mapsto t^{p_0}$ is non-decreasing on $[0,\infty)$.

(i) If $y\ge x$, then $\log y\ge\log x\ge 1$, so
$(\log y)^{p_0}\ge(\log x)^{p_0}\ge\tfrac12(\log x)^{p_0}$.
Suppose now $y<x$. Since $x\ge 2c_5$ we have $x-c_5\ge x/2$, and since $x\ge e>2$ we have $x/2>1$;
hence $y\ge x-c_5\ge x/2>1$, and in particular $\log y>0$ and $x-c_5>0$. Using $y\ge x-c_5$ and the
inequality $\log(1+t)\le t$ for $t>-1$, we obtain
\begin{equation*}
\log x-\log y\le\log\frac{x}{x-c_5}\le\frac{x}{x-c_5}-1=\frac{c_5}{x-c_5}\le\frac{2c_5}{x},
\end{equation*}
the last step by $x-c_5\ge x/2$. Therefore $\log y\ge(1-u)\log x$ with
\begin{equation*}
u:=\frac{2c_5}{x\log x}\le\frac{2c_5}{x}\le 1,
\end{equation*}
where the first inequality uses $\log x\ge1$ and the second uses $x\ge 2c_5$. Thus $0\le 1-u$, and
since both sides of $\log y\ge(1-u)\log x$ are non-negative,
$(\log y)^{p_0}\ge(1-u)^{p_0}(\log x)^{p_0}$. By Bernoulli's inequality for the integer $p_0\ge1$ and
$u\in[0,1]$, together with $u\le 2c_5/x$ and $x\ge 4p_0c_5$,
\begin{equation*}
(1-u)^{p_0}\ge 1-p_0u\ge 1-\frac{2p_0c_5}{x}\ge \frac12 .
\end{equation*}
This proves \eqref{4.36:eq-log-comparison-1}.

(ii) If $y\le x$, then $1\le y\le x$ gives $0\le\log y\le\log x$, so
$(\log y)^{p_0}\le(\log x)^{p_0}\le(1+\beta_0)(\log x)^{p_0}$.
Suppose now $y>x$. Then $\log y>\log x\ge 1$, and by $y\le x+c_5$, the inequality
$\log(1+t)\le t$ for $t\ge 0$, and $\log x\ge 1$,
\begin{equation*}
\log y\le\log x+\log\Bigl(1+\frac{c_5}{x}\Bigr)\le\log x+\frac{c_5}{x}
\le\Bigl(1+\frac{c_5}{x}\Bigr)\log x .
\end{equation*}
Both sides being positive, $(\log y)^{p_0}\le(1+c_5/x)^{p_0}(\log x)^{p_0}$. Using $1+t\le e^{t}$ and
$x\ge 2p_0c_5/\beta_0$, that is $p_0c_5/x\le\beta_0/2$,
\begin{equation*}
\Bigl(1+\frac{c_5}{x}\Bigr)^{p_0}\le e^{p_0c_5/x}\le e^{\beta_0/2}\le 1+\beta_0 .
\end{equation*}
The last inequality is $e^{t/2}\le 1+t$ for $t\in[0,1]$, applied with $t=\beta_0$: the function
$t\mapsto e^{t/2}$ is convex, so on $[0,1]$ it lies below its chord,
$e^{t/2}\le(1-t)+te^{1/2}$, and $e^{1/2}\le 2$ because $e\le 4$; hence $e^{t/2}\le 1+t$. This proves
\eqref{4.36:eq-log-comparison-2}.
\end{proof}

In this chapter the lemma is applied with $x$ and $y$ among the inverse couplings $x_j=g_j^{-2}$ of
Notation~\ref{4.36:not-flow-thresholds} and with $c_5$ the constant of the flow inequalities recorded
there; for instance, part~(ii) turns $x_{n'}\le x_j+c_5$ into
$(\log x_{n'})^{p_0}\le(1+\beta_0)(\log x_j)^{p_0}$.

\begin{proposition}[Transport bounds for every $L$ under a coupling ceiling. Stated; proof
incomplete at the step marked $(\ast)$]\label{4.36:prop-transport-every-l}
Let the levels, the inverse couplings $x_j=g_j^{-2}$, the degree function $p_0(\cdot)$, the radii
$\varepsilon_j=g_jp_0(g_j)$, $\delta_j=(A_1/A_0)\varepsilon_j$, $\alpha_{0,j}$, $\alpha_{1,j}$ and the
lengths $R_j$ be those of Notation~\ref{4.36:not-flow-thresholds}. Let $x_\delta$, $x_\beta$ and
$\beta_0\in(0,1]$ be as in Lemma~\ref{4.36:lem-log-comparison} (comparison of powers of logarithms
under a bounded shift), and put $\hat\varepsilon_n:=c_{\mathrm{reg}}LR_n\varepsilon_n$, where
$c_{\mathrm{reg}}$ is a fixed constant not depending on the level. Assume:
\begin{itemize}
\item the lower and upper flow inequalities at $SU(N)$ of Notation~\ref{4.36:not-flow-thresholds}, with
lower constant $\lambda^{\sharp}$, upper constant $B^{\sharp}$, shift constant $c_5$, and the fixed
number $\bar B\ge B^{\sharp}$, together with their
consequences $(\mathrm{Fl}1)$--$(\mathrm{Fl}4)$ in \eqref{4.36:eq-flow-thresholds-a};
\item the anchoring of the flow at the last level, in the form $x_K\ge g^{-2}$;
\item the coupling ceiling
\begin{equation}\label{4.36:eq-transport-every-l-1}
\gamma\le\gamma_{\bar B}:=\bigl(\max\{\bar B,x_\delta,x_\beta\}+c_5\bigr)^{-1/2}.
\end{equation}
\end{itemize}
Then the following hold.
\begin{enumerate}
\item[(a)] $x_n\ge\max\{\bar B,x_\delta,x_\beta\}$ for every level $n\le K$.
\item[(b)] For all levels $j\le n\le k\le K$,
\begin{equation}\label{4.36:eq-transport-every-l-a}
\begin{gathered}
x_j\le(1+(n-j))\,x_n,\qquad g_n\le(1+(n-j))^{1/2}g_j,\\
\varepsilon_n\le(1+\beta_0)(1+(n-j))^{1/2}\varepsilon_j,\qquad
\delta_n\le(1+\beta_0)(1+(n-j))^{1/2}\delta_j,\\
\hat\varepsilon_n\le(1+\beta_0)(1+(n-j))^{1/2}\hat\varepsilon_j,\\
\delta_n\ge\tfrac12(1+(k-n))^{-1/2}\delta_k,\qquad p_0(g_n)\ge\tfrac12\,p_0(g_k).
\end{gathered}
\end{equation}
These bounds contain neither $N$ nor $\bar B$; they replace the first half of $(\mathrm{Fl}2)$ and
the inequality $(\mathrm{Fl}3)$ of \eqref{4.36:eq-flow-thresholds-a}. Moreover, for every integer
$L\ge2$ and all levels $j\le n\le K$,
\begin{align}
\delta_nL^{-(n-j)}&\le(1+\beta_0)\,\delta_j,\label{4.36:eq-transport-every-l-2}\\
\varepsilon_nL^{-(n-j)}&\le(1+\beta_0)\,\varepsilon_j,\label{4.36:eq-transport-every-l-3}\\
\varepsilon_nL^{-2(n-j)}&\le(1+\beta_0)\,\varepsilon_j,\label{4.36:eq-transport-every-l-4}\\
\hat\varepsilon_nL^{-2(n-j)}&\le(1+\beta_0)\,\hat\varepsilon_j.\label{4.36:eq-transport-every-l-5}
\end{align}
\item[(c)] Each of the following inequalities of this chapter, whose proofs from
$(\mathrm{Fl}2)$ and $(\mathrm{Fl}3)$ invoke the condition $L^2\ge1+\bar B$ among the threshold
conditions \eqref{4.36:eq-flow-thresholds-b}, holds with the constant with which it is stated,
for every integer $L\ge2$ (levels $j\le n'$ and $j\le m\le n$):
\begin{itemize}
\item the unit transport line $\delta_{n'}L^{j-n'}\le(1+\beta_0)\delta_j\le2\delta_j$, after
which $2\delta_j\le\frac14\alpha_{1,j}$ follows from the radius comparison above the
threshold $x_{\mathrm{th}}$, unchanged;
\item the per-plaquette curvature bound by $2(1+\beta_0)B_3\varepsilon_j$, where $B_3$ is the
constant of that bound;
\item the bound $c^{\mathrm{re}}_m\le(1+\beta_0)C_{\mathrm{re}}\varepsilon_m$, with
$C_{\mathrm{re}}$ the constant of that bound, which lies within
the bound $2(1+\beta_0)C_{\mathrm{re}}\varepsilon_m$ with which it is used;
\item the transport inequalities
\begin{equation*}
\varepsilon_mL^{j-m}\le(1+\beta_0)\varepsilon_j,\qquad
c^{\mathrm{im}}_mL^{j-m}\le c_B(1+\beta_0)\delta_j,
\end{equation*}
which occur in the bound on $c^{\mathrm{im}}_m$ (with $c_B$ the constant of that bound) and in the
opening transport step of the transported data-norm lemma; they give, above the threshold
$x_{\mathrm{th}}$ and by the radius comparison,
\begin{equation*}
C_Q\cdot2\bigl(c^{\mathrm{re}}_m+c^{\mathrm{im}}_m\bigr)L^{j-m}\le\tfrac12\,\alpha_{3,j},
\end{equation*}
where $C_Q$ is the constant with which these coefficients enter the transported data-norm lemma
and $\alpha_{3,j}$ is a further radius at level $j$;
\item the inequality $\hat\varepsilon_nL^{-2(n-m)}\le(1+\beta_0)\hat\varepsilon_m$, and the
thresholds with which the transported data-norm lemma is applied in this chapter.
\end{itemize}
Consequently, on each of the inequalities just listed, the condition $L^2\ge1+\bar B$ of
\eqref{4.36:eq-flow-thresholds-b} is needed only in the form $L^2\ge2$, which every integer
$L\ge2$ satisfies; on these inequalities it is therefore superfluous. The hypothesis
$L^2\ge1+\bar B$ of the transported data-norm lemma is not needed. Likewise, the hypothesis
$L^2\ge1+\bar B$ of the data-norm lemma, with which that lemma is used in this chapter, is not
needed; this is the assertion whose proof is incomplete at the step marked $(\ast)$. The
threshold $x_{\mathrm{th}}$ of
\eqref{4.36:eq-flow-thresholds-b} is kept, enlarged only to
$\max\{x_{\mathrm{th}},\bar B,x_\delta,x_\beta\}$.
\item[(d)] In the proofs of this chapter that use the flow inequalities, write $B$ for either of
$B^{\sharp}$ and $\bar B$. Every occurrence of $B$ other than those in (c) is admissible for every
$B\ge0$ and every $L>1$. It enters either through the value of a constant that is finite for all
such $B$ and $L$, or through a lower bound on the inverse couplings (an $x$-floor).
\end{enumerate}
\end{proposition}

\begin{proof}
\emph{Proof of} (a). Let $n\le K$. The lower flow inequality between the levels $n$ and $K$ gives
$x_n\ge x_K+\lambda^{\sharp}(K-n)-c_5$. Here $\lambda^{\sharp}(K-n)\ge0$ and $x_K\ge g^{-2}$ by the
anchoring of the flow, so $x_n\ge g^{-2}-c_5$. The bound used below is the one given by
$(\mathrm{Fl}1)$ in \eqref{4.36:eq-flow-thresholds-a} at $k=n$, namely $x_n\ge\gamma^{-2}-c_5$. Hence
\begin{equation*}
x_n\ge\gamma^{-2}-c_5\ge\max\{\bar B,x_\delta,x_\beta\}.
\end{equation*}
The last inequality holds because \eqref{4.36:eq-transport-every-l-1} is equivalent to
$\gamma^{-2}\ge\max\{\bar B,x_\delta,x_\beta\}+c_5$.

\emph{Proof of} (b). Let $j\le n\le k\le K$. The upper flow inequality gives
$x_j-x_n\le B^{\sharp}(n-j)$. Since $B^{\sharp}\le\bar B$, and $\bar B\le x_n$ by (a), we get
\begin{equation*}
x_j-x_n\le B^{\sharp}(n-j)\le\bar B(n-j)\le x_n(n-j).
\end{equation*}
Hence $x_j\le(1+(n-j))x_n$. Equivalently
\begin{equation*}
g_n^2=\frac1{x_n}\le\frac{1+(n-j)}{x_j}=(1+(n-j))\,g_j^2,
\end{equation*}
which is the bound on $g_n$.

Next, $p_0(g')=A_0(\log g'^{-2})^{p_0}$ for $g'>0$, so $\varepsilon_j=A_0g_j(\log x_j)^{p_0}$. By (a),
$x_n,x_j\ge x_\beta\ge e$, so all logarithms involved are at least $1$ and
\begin{equation*}
\frac{\varepsilon_n}{\varepsilon_j}
=\frac{g_n}{g_j}\Bigl(\frac{\log x_n}{\log x_j}\Bigr)^{p_0}.
\end{equation*}
The first factor is at most $(1+(n-j))^{1/2}$ by what was just proved. For the second factor,
apply Lemma~\ref{4.36:lem-log-comparison}(ii) with $x:=x_j$ and $y:=x_n$. Its hypotheses hold:
$x_j\ge x_\beta$ by (a); $x_n\ge e\ge1$ by (a); and $x_n\le x_j+c_5$ by $(\mathrm{Fl}2)$. The
lemma gives $(\log x_n)^{p_0}\le(1+\beta_0)(\log x_j)^{p_0}$. Together,
$\varepsilon_n\le(1+\beta_0)(1+(n-j))^{1/2}\varepsilon_j$.

Since $\delta_j=(A_1/A_0)\varepsilon_j$ at every level, the same bound holds for $\delta$. For
$\hat\varepsilon$ we use $R_n\le R_j$, from clause (a) of the lemma on the lengths $R_j$:
\begin{equation*}
\hat\varepsilon_n=c_{\mathrm{reg}}LR_n\varepsilon_n
\le c_{\mathrm{reg}}LR_j(1+\beta_0)(1+(n-j))^{1/2}\varepsilon_j
=(1+\beta_0)(1+(n-j))^{1/2}\hat\varepsilon_j.
\end{equation*}

For the lower bounds, apply the first inequality of \eqref{4.36:eq-transport-every-l-a} to the
pair of levels $n\le k$. It gives $x_n\le(1+(k-n))x_k$, hence
\begin{equation*}
g_n/g_k=(x_k/x_n)^{1/2}\ge(1+(k-n))^{-1/2}.
\end{equation*}
Apply Lemma~\ref{4.36:lem-log-comparison}(i) with $x:=x_k$ and $y:=x_n$. Its hypotheses hold:
$x_k\ge x_\delta$ by (a), and, by the lower flow inequality between the levels $n$ and $k$,
\begin{equation*}
x_n\ge x_k+\lambda^{\sharp}(k-n)-c_5\ge x_k-c_5.
\end{equation*}
The
lemma gives $(\log x_n)^{p_0}\ge\frac12(\log x_k)^{p_0}$, that is,
$p_0(g_n)\ge\frac12p_0(g_k)$. Multiplying,
\begin{equation*}
\frac{\delta_n}{\delta_k}=\frac{g_n}{g_k}\Bigl(\frac{\log x_n}{\log x_k}\Bigr)^{p_0}
\ge\tfrac12(1+(k-n))^{-1/2}.
\end{equation*}
None of these bounds contains $N$ or $\bar B$.

For the transport bounds, let $L\ge2$ be an integer. By
Lemma~\ref{4.36:lem-transport-arithmetic}(v) (arithmetic of the transport factors), with $n-j\ge0$
in place of its $m$,
\begin{equation*}
(1+(n-j))^{1/2}L^{-(n-j)}\le1,\qquad(1+(n-j))^{1/2}L^{-2(n-j)}\le1.
\end{equation*}
Multiplying the bounds on $\delta_n$,
$\varepsilon_n$ and $\hat\varepsilon_n$ in \eqref{4.36:eq-transport-every-l-a} by $L^{-(n-j)}$ or by
$L^{-2(n-j)}$ gives \eqref{4.36:eq-transport-every-l-2}--\eqref{4.36:eq-transport-every-l-5}.

\emph{Proof of} (c). Each inequality listed in (c) is a bound of a smallness parameter at a higher
level by one at a lower level, multiplied by $L^{-\mathrm{gap}}$ or by
$L^{-2\,\mathrm{gap}}$, where the gap is the difference of the two levels. Its logarithmic part is
Lemma~\ref{4.36:lem-log-comparison}(ii), applied with $x$ the inverse coupling at the lower level
$j$ and $y$ that at the higher level $n'$, using $x_{n'}\le x_j+c_5$ from $(\mathrm{Fl}2)$. The
$x$-floor needed there is $x_j\ge x_\beta$. This is the same floor that already enters these
proofs, where $x_{n'}\le x_j+c_5$ is used to give
$(\log x_{n'})^{p_0}\le(1+\beta_0)(\log x_j)^{p_0}$; by (a) it holds at every level.

Replacing $(\mathrm{Fl}3)$ by \eqref{4.36:eq-transport-every-l-a} and applying
Lemma~\ref{4.36:lem-transport-arithmetic}(v) gives the transport bounds
\eqref{4.36:eq-transport-every-l-2}--\eqref{4.36:eq-transport-every-l-5} for every integer
$L\ge2$, with no further condition on $L$. The individual lines are obtained as below.
\begin{itemize}
\item The unit transport line is \eqref{4.36:eq-transport-every-l-2} with $n=n'$. Then
$(1+\beta_0)\delta_j\le2\delta_j$ because $\beta_0\le1$.
\item The per-plaquette curvature bound is obtained from \eqref{4.36:eq-transport-every-l-4}.
\item The bound $c^{\mathrm{re}}_m\le(1+\beta_0)C_{\mathrm{re}}\varepsilon_m$ is covered by the
argument of the preceding two paragraphs: its factor $(1+\beta_0)$ is the logarithmic comparison of
Lemma~\ref{4.36:lem-log-comparison}(ii) at the $x$-floor $x_\beta$, which holds at every level
by (a), and its flow comparison is taken from \eqref{4.36:eq-transport-every-l-a} in place of
$(\mathrm{Fl}3)$. It therefore holds with the constant $(1+\beta_0)C_{\mathrm{re}}$ for every
integer $L\ge2$, within the bound $2(1+\beta_0)C_{\mathrm{re}}\varepsilon_m$ with which it is used.
\item The first inequality of the opening transport step of the transported data-norm lemma is
\eqref{4.36:eq-transport-every-l-3} with $n=m$. The inequality
$c^{\mathrm{im}}_mL^{j-m}\le c_B(1+\beta_0)\delta_j$ carries a quantity at level $m$ to level $j$
by the factor $L^{j-m}$, with the constant $c_B$ and the factor $(1+\beta_0)$ of
\eqref{4.36:eq-transport-every-l-2} applied with $n=m$.
\item The inequality $\hat\varepsilon_nL^{-2(n-m)}\le(1+\beta_0)\hat\varepsilon_m$ is
\eqref{4.36:eq-transport-every-l-5} with the pair of levels $(j,n)$ replaced by $(m,n)$.
\end{itemize}
The remaining inequalities $2\delta_j\le\frac14\alpha_{1,j}$ and
$C_Q\cdot2(c^{\mathrm{re}}_m+c^{\mathrm{im}}_m)L^{j-m}\le\frac12\alpha_{3,j}$ follow from the
radius comparison above the threshold $x_{\mathrm{th}}$, which the present argument leaves
unchanged. By (a), replacing
$x_{\mathrm{th}}$ by $\max\{x_{\mathrm{th}},\bar B,x_\delta,x_\beta\}$ costs nothing further.

The number $\bar B$ does not occur in \eqref{4.36:eq-transport-every-l-a}. The factor
$(1+(n-j))^{1/2}$ there is the factor of Lemma~\ref{4.36:lem-transport-arithmetic}(iii) with
$\bar B$ replaced by $1$ and $n-j$ in place of its $m$. Accordingly, the
condition $L^2\ge1+\bar B$ is needed on the inequalities listed in (c) only with $\bar B$
replaced by $1$, that is, as $L^2\ge2$.

It remains to treat the assertion of (c) on the data-norm lemma; this is a different statement
from the transported data-norm lemma, and the step marked $(\ast)$ below concerns this assertion
alone. The data-norm lemma is to be used in this chapter with its hypothesis
$L^2\ge1+\bar B$ discharged by replacing $(\mathrm{Fl}3)$ with
\eqref{4.36:eq-transport-every-l-a}; the occurrences of the scale index in its proof for which
this is carried out here are the following:
\begin{itemize}
\item the bounds on $c^{\mathrm{re}}$ and $c^{\mathrm{im}}$;
\item the transport in its opening step;
\item the quantities $\vartheta_{jn}$ of that lemma;
\item the constant $C_{B,\beta,L}$ of \eqref{4.36:eq-transport-every-l-6}, which enters its proof
and does not involve the condition $L^2\ge1+\bar B$.
\end{itemize}
The first three have the shape
\begin{equation*}
(\text{smallness at level }n)\cdot L^{-(n-j)}\le C\cdot(\text{quantity at level }j),
\end{equation*}
or the same with $L^{-2(n-j)}$. This shape is covered by
\eqref{4.36:eq-transport-every-l-2}--\eqref{4.36:eq-transport-every-l-5}. It is the only shape in
which $B$ enters when a smallness at a level $n'$ is carried to the level $j$, since such a
transport always comes with the geometric gain of the change of scale in \cite{B-CMP119}.
$(\ast)$ That the proof of the data-norm lemma uses the hypothesis $L^2\ge1+\bar B$ in no other
way, that is, that it nowhere compares $\bar B$ with $L^2$ except in the shapes just listed, is
not proved here. If such a comparison occurred, this assertion would hold only under an
additional one-sided lower bound on $L$ depending on $N$ alone; this would be a lower bound on
$L$, not a cap in the sense of Definition~\ref{2.3:def-cap} and not a condition on the coupling,
and the inequalities listed in (c) would hold for every integer $L\ge2$ as proved above.

\emph{Proof of} (d). The remaining occurrences of $B$ are the following.
\begin{itemize}
\item The constant of the curvature lemma,
\begin{equation*}
C_W=36\,C_q\,c_0^d\sum_{\ell\ge0}\bigl(1+B(\ell+2)\bigr)e^{-\ell}<\infty .
\end{equation*}
\item The constant $C'_W$ of the curvature lemma. Its defining sum converges for every
$B,\bar B\ge0$ and $L>1$, because
\begin{equation*}
\log_L\bigl(1+\log(1+Bu)\bigr)\le\log(1+Bu)/\log L\le Bu/\log L.
\end{equation*}
The constants $C_{B,r}$ and $C_1(B,r)$ that accompany it are finite. The condition
$\bar BP_1(x_k)\le x_k$, with $P_1$ a fixed function of the inverse coupling accompanying these
constants, is a ceiling on the coupling, equivalently an $x$-floor.
\item The factor $(1+\log(1+B(k-j)))^{2p_0+4r}$ and the polylogarithmic factors in $k-j$ of the
two bounds proved with it. They stand against weights $e^{-(k-j)}$, so their sums are finite for
every $B\ge0$.
\item The companion bound for the decay weights $e_y$, where $\nu_0$ denotes the decay rate of
these weights and $R_*$ the length entering them,
\begin{equation*}
\frac{e_y^{1/2}}{\delta_n}
\le2\bigl(1+B(k-n)\bigr)^{1/2}e^{-\nu_0M(k-1-n)/16}\,e^{-\nu_0R_*/16}\,\frac1{\delta_k}
\le\frac1{\delta_k}.
\end{equation*}
It does not use the condition $L^2\ge1+\bar B$. It uses
$(1+B(\ell+1))^{1/2}e^{-\ell}\le(1+B)^{1/2}$ for integers $\ell\ge0$, which follows from
Lemma~\ref{4.36:lem-transport-arithmetic}(i), and the $x$-floor
$e^{\nu_0R_*/16}\ge2(1+B)^{1/2}$. Under \eqref{4.36:eq-transport-every-l-a} it becomes
independent of $N$, with the $x$-floor
\begin{equation*}
e^{\nu_0R_*/16}\ge2\sup_{\ell\ge0}(2+\ell)^{1/2}e^{-\ell}=2\sqrt2.
\end{equation*}
\item The factor $1+B(k-m)$ in the quantity $\Theta_m$ of these proofs. Summed over the level,
this factor appears inside
$\sum_ne^{-(k-n)}(1+B(k-n))<\infty$, so it enters as a value. It can be made free of $B$ under
\eqref{4.36:eq-transport-every-l-a} if preferred.
\item In the data-norm lemma, the bound $\vartheta_{jn}\le C(L,M)(1+B(n-j)/x_n)^{1/2}$. Under (a),
$B\le\bar B\le x_n$, so it is at most $C(L,M)(1+(n-j))^{1/2}$. Another step of its proof uses the
constant
\begin{equation}\label{4.36:eq-transport-every-l-6}
C_{B,\beta,L}:=\sup_{u\ge0}(1+Bu)^{5/2}L^{-\beta u/2},
\end{equation}
which is finite for every $B\ge0$, $\beta>0$ and $L>1$.
\item In the second part of the transported data-norm lemma, the bound
$\varpi_{jm}\le4\varpi_{jk}$ on the weights $\varpi$ of that lemma. It follows from $x_m\ge c_5$,
which is automatic under (a), since $x_\delta\ge2c_5$.
\item The uses of $B$ in the proof of the energy theorem. They enter as values, and none of them
is compared with a power of $L$.
\item Two further clauses of the lemma on the lengths $R_j$. There the powers of $L$ are those
of Ba{\l}aban's construction.
\end{itemize}
\end{proof}

\begin{remark}[The order of choice of the ceiling]
The number $\bar B$ of the flow inequalities is fixed; it depends neither on $\gamma$ nor on $K$.
Hence $\gamma_{\bar B}$ in \eqref{4.36:eq-transport-every-l-1} is a number determined by $N$, $L$
and $\mathfrak{p}$, and the ceiling $\gamma\le\gamma_{\bar B}$ is imposed after these, as the last
choice. For every $L$ the interval $(0,\min\{\gamma_N,\gamma_{\bar B}\}]$ is non-empty, where
$\gamma_N$ is the coupling ceiling that corresponds to the thresholds
\eqref{4.36:eq-flow-thresholds-b}. The same mechanism appears in \cite[(3.8)]{B-CMP119}: there the
factor $2(1+\beta_0)L^{-2}$ controls $g_{k+1}/g_k$ through the condition $B/x\le\beta_0$ on the
constant $B$ and the inverse coupling $x$ of that display, which is a ceiling on the coupling.
\end{remark}

\begin{remark}[An alternative route without the ceiling $\gamma_{\bar B}$]
By Lemma~\ref{4.36:lem-transport-arithmetic}(iii), the bounds $(\mathrm{Fl}2)$ and
$(\mathrm{Fl}3)$ of \eqref{4.36:eq-flow-thresholds-a} give, for every integer $L\ge2$,
\begin{equation*}
\delta_nL^{-(n-j)}\le(1+\beta_0)(1+\bar B)^{1/2}\delta_j,
\end{equation*}
and the analogous bound for $\varepsilon_n$. By Lemma~\ref{4.36:lem-transport-arithmetic}(iv), the
same holds with $L^{-2(n-j)}$ for $\varepsilon_n$ and $\hat\varepsilon_n$. The inequalities of
clause (c) of Proposition~\ref{4.36:prop-transport-every-l}, grouped as below, then hold under
ceilings on $\gamma$ that absorb the factor $(1+\bar B)^{1/2}$ into the spare
powers of the logarithm in the radius comparison:
\begin{enumerate}
\item the unit transport line holds provided
\begin{equation*}
(\log x_j)^{q_1-p_0}\ge8(1+\bar B)^{1/2}A_1/C_1;
\end{equation*}
\item the per-plaquette curvature bound and the bound on $c^{\mathrm{re}}_m$ hold provided
\begin{equation*}
(\log x_j)^{q_0-p_0}\ge8(1+\beta_0)(1+\bar B)^{1/2}B_3A_0/C_0,
\end{equation*}
together with the smallness condition used with these bounds;
\item the transport of $c^{\mathrm{im}}_m$ holds provided
$(\log x_j)^{q-p_0}\ge C(L,M)(1+\bar B)$;
\item the inequality $\hat\varepsilon_nL^{-2(n-m)}\le(1+\beta_0)\hat\varepsilon_m$ and the
thresholds with which the transported data-norm lemma is applied in this chapter hold provided
\begin{equation*}
(\log x_j)^{q_0-p_0-r}\ge4(1+\beta_0)^2(1+\bar B)^{1/2}c_{\mathrm{reg}}L^2A_0/C_0.
\end{equation*}
\end{enumerate}
The exponents on the left are at least $2r+2>0$ by the radius comparison. The conclusion is the
same as in (c). Proposition~\ref{4.36:prop-transport-every-l} is preferred to this route because
it keeps every constant of the inequalities in (c) unchanged.
\end{remark}

\begin{lemma}[The one-link partition number at $SU(N)$]\label{4.36:lem-one-link-partition}
Let $G=SU(N)$ with $N\ge 2$, its Lie algebra $\mathfrak{g}$ of traceless Hermitian $N\times N$
matrices, $D=N^2-1=\dim\mathfrak{g}$, $\mathsf{n}=(N-1)^{1/2}$, the operator norm $|\cdot|$, the norm
$|\cdot|_{\mathfrak{g}}$ and the eigen-angles be as in Notation~\ref{4.36:not-group-norms}. Let $dX$
denote the Lebesgue measure of the Euclidean space
$(\mathfrak{g},\langle X,Y\rangle=\operatorname{tr}XY)$, written $dv$ in
Definition~\ref{4.36:def-haar-balls}, and let the domain $\mathfrak{D}$, the Haar density $J$ in
exponential coordinates, its normaliser $\sigma_N$ and the volume $\omega_D$ of the unit
$|\cdot|_{\mathfrak{g}}$-ball be as in that definition. Let the degree function $p_0(\cdot)$, with $A_0>0$ and $p_0$ a
positive integer, and the running couplings $g_j$ be as in Notation~\ref{4.36:not-flow-thresholds}.
Put
\begin{equation*}
r_G:=\frac{\sqrt6}{N},\qquad
\mathfrak{B}_z:=\{X\in\mathfrak{g}:\ |X|_{\mathfrak{g}}<r_G\},\qquad
\varepsilon_G:=2\sin\Bigl(\frac{\sqrt6}{2N}\Bigr),\qquad
v_2:=\Bigl(\frac{\pi^3}{2}\Bigr)^{D/2},
\end{equation*}
\begin{equation}\label{4.36:eq-one-link-partition-1}
z(g):=\int_G\chi\bigl(\{|u-\mathbf{1}|<g\,p_0(g)\}\bigr)\,
e^{-g^{-2}(1-\operatorname{Re}\operatorname{tr}u)}\,du,\qquad z_j:=z(g_j),
\end{equation}
where $\chi(S)$ denotes the characteristic function of a set $S\subset G$ and $du$ is the Haar
probability measure on $G$, and
\begin{equation}\label{4.36:eq-one-link-partition-2}
\gamma_z=\gamma_z(N,A_0,p_0):=\sup\Bigl\{g_1\in\bigl(0,\min\{e^{-p_0},r_G/2\}\bigr]:\
g_1p_0(g_1)\le\varepsilon_G\ \text{and}\ p_0(g_1)\ge2\mathsf{n}\Bigr\}.
\end{equation}
The number $z(g)$ is the one-link integral of \cite[(1.5)]{B-CMP119}, with the coupling $g$ in
place of $g_0$ and $g\,p_0(g)$ in place of $\varepsilon_0$; the measure $du$ is normalised as in
\cite[(0.15)--(0.16)]{B-CMP109}, so that the integral over $u$ of the constant $1$ equals $1$. Then
the following hold.
\begin{enumerate}
\item[(o)] $\gamma_z>0$, and for every $0<g\le\gamma_z$,
\begin{equation*}
g\le\min\{e^{-p_0},r_G/2\},\qquad g\,p_0(g)\le\varepsilon_G,\qquad p_0(g)\ge2\mathsf{n}.
\end{equation*}
Each of these three conditions is an upper bound on $g$; the number $\gamma_z$ is fixed after
$N$, $A_0$, $p_0$ and depends on them only, and no condition is placed on $A_0$.
\item[(i)] The map $X\mapsto e^{iX}$ restricted to $\mathfrak{B}_z$ is an analytic diffeomorphism
onto an open set $O_G\subset G$ with $O_G\supseteq\{u\in G:|u-\mathbf{1}|<\varepsilon_G\}$. On
$\mathfrak{B}_z$ one has $\tfrac12<J\le1$, so that there the Haar probability is $\sigma\,dX$ with
$\sigma:=\sigma_NJ$ and
\begin{equation*}
\tfrac12\sigma_N<\sigma\le\sigma_N=\sigma(0).
\end{equation*}
\item[(ii)] For $0<g\le\gamma_z$,
\begin{equation}\label{4.36:eq-one-link-partition-3}
\tfrac12\sigma_Ne^{-1/2}\omega_D\,g^D\le z(g)\le\sigma_Nv_2\,g^D .
\end{equation}
Hence, whenever $0<g_j\le\gamma_z$,
\begin{equation}\label{4.36:eq-one-link-partition-4}
\begin{split}
&\log z_j=D\log g_j+\zeta_j,\\
&|\zeta_j|\le c_G(N):=\max\bigl\{\bigl|\log\bigl(\tfrac12\sigma_Ne^{-1/2}\omega_D\bigr)\bigr|,\
|\log(\sigma_Nv_2)|\bigr\}.
\end{split}
\end{equation}
\item[(iii)] The number $\nu_{\mathfrak{g}}$ defined below satisfies
\begin{equation*}
\nu_{\mathfrak{g}}:=\log\int_{\mathfrak{g}}e^{-\frac12|X|^2_{\mathfrak{g}}}\,dX=\frac{D}{2}\log 2\pi.
\end{equation*}
\end{enumerate}
\end{lemma}

\begin{proof}
\emph{(o).} Put $\phi(g):=g(\log g^{-2})^{p_0}$ and $\psi(g):=(\log g^{-2})^{p_0}$ for $0<g<1$, so that
$g\,p_0(g)=A_0\phi(g)$ and $p_0(g)=A_0\psi(g)$. Differentiating,
\begin{equation*}
\phi'(g)=(\log g^{-2})^{p_0}-2p_0(\log g^{-2})^{p_0-1}
=(\log g^{-2})^{p_0-1}\bigl(\log g^{-2}-2p_0\bigr),
\end{equation*}
which is positive for $0<g<e^{-p_0}$, since there $\log g^{-2}>2p_0>0$; moreover $\phi(g)\to0$ as
$g\downarrow0$. The function $\psi$ is strictly decreasing on $(0,1)$ and $\psi(g)\to\infty$ as
$g\downarrow0$. Let $I_z$ be the set whose supremum is taken in
\eqref{4.36:eq-one-link-partition-2}, a subset of $(0,g_*]$ with
$g_*:=\min\{e^{-p_0},r_G/2\}<1$. By the two limits, $I_z$ contains every sufficiently small
$g_1>0$, so it is non-empty and $\gamma_z>0$. If $g_1\in I_z$ and $0<g<g_1$, then
$\phi(g)\le\phi(g_1)$ and $\psi(g)\ge\psi(g_1)$ because $\phi$ is increasing and $\psi$ decreasing on
$(0,g_*]$; hence $g\in I_z$, and $I_z$ is an initial segment of $(0,g_*]$. Since
$\phi$ and $\psi$ are continuous, the two defining inequalities are preserved under limits within
$(0,g_*]$, so $\gamma_z\in I_z$. Therefore $I_z=(0,\gamma_z]$, which is the assertion
that the three bounds hold for every $0<g\le\gamma_z$. Finally, since $\log g^{-2}>0$ on $(0,1)$, the
condition $p_0(g)\ge2\mathsf{n}$, that is $A_0(\log g^{-2})^{p_0}\ge2\mathsf{n}$, is equivalent to
$\log g^{-2}\ge(2\mathsf{n}/A_0)^{1/p_0}$, that is to
\begin{equation*}
g\le\exp\Bigl(-\tfrac12\bigl(2\mathsf{n}/A_0\bigr)^{1/p_0}\Bigr),
\end{equation*}
an upper bound on $g$ for every $A_0>0$. The other two conditions are upper bounds on $g$ by the
initial-segment property just shown.

\emph{(i), the window lies in the image of the chart.} Let $u\in G$ with $|u-\mathbf{1}|<\varepsilon_G$,
and write its spectral decomposition $u=\sum_pe^{i\theta_p}P_p$ with eigen-angles
$\theta_p\in\,]-\pi,\pi]$ and mutually orthogonal rank-one projections $P_p$. Since $u-\mathbf{1}$ is
normal, $|u-\mathbf{1}|=\max_p|e^{i\theta_p}-1|$, and $|e^{i\theta}-1|=2|\sin(\theta/2)|$; hence for
every $p$
\begin{equation*}
2|\sin(\theta_p/2)|\le|u-\mathbf{1}|<2\sin\Bigl(\frac{\sqrt6}{2N}\Bigr).
\end{equation*}
Here $|\theta_p|/2\in[0,\pi/2]$ and $\sqrt6/(2N)\le\sqrt6/4<\pi/2$ because $N\ge2$; since $\sin$ is
increasing on $[0,\pi/2]$, it follows that $|\theta_p|<\sqrt6/N$. The condition $\det u=1$ gives
$e^{i\sum_p\theta_p}=1$, so $\sum_p\theta_p\in2\pi\mathbb{Z}$, while
$|\sum_p\theta_p|<N\cdot\sqrt6/N=\sqrt6<2\pi$; therefore $\sum_p\theta_p=0$. Consequently the
principal logarithm $X:=\sum_p\theta_pP_p$ (\cite[(22)]{B-CMP98}) is Hermitian and traceless, so
$X\in\mathfrak{g}$; by the first inequality of (F1) in
Lemma~\ref{4.36:lem-eigenangle-inequalities}, $|X|_{\mathfrak{g}}\le|X|=\max_p|\theta_p|<r_G$; and
$u=e^{iX}$. Thus $u\in\exp i(\mathfrak{B}_z)$. We note in passing that the central elements
$e^{2\pi ik/N}\mathbf{1}$, $k=1,\dots,N-1$, other than the identity satisfy
$|e^{2\pi ik/N}\mathbf{1}-\mathbf{1}|=2|\sin(\pi k/N)|\ge2\sin(\pi/N)>\varepsilon_G$, since
$\sqrt6/(2N)<\pi/N\le\pi/2$; they lie outside the window, as they must.

\emph{(i), the chart.} Let $X\in\mathfrak{B}_z$ have eigen-angles $\theta_1,\dots,\theta_N$. By the
second inequality of (F1) in Lemma~\ref{4.36:lem-eigenangle-inequalities},
\begin{equation*}
|\theta_p|\le|X|\le\mathsf{n}|X|_{\mathfrak{g}}<\mathsf{n}\,\frac{\sqrt6}{N}
=\frac{(6(N-1))^{1/2}}{N}\le\frac{\sqrt6}{2}<\pi,
\end{equation*}
where $6(N-1)/N^2\le3/2$ is equivalent to $N^2-4N+4\ge0$, that is to $(N-2)^2\ge0$. By (F2) of
Lemma~\ref{4.36:lem-eigenangle-inequalities}, $(\theta_p-\theta_q)^2\le2N|X|^2_{\mathfrak{g}}$ for all
$p,q$, so the spread $s(X)=\theta_{\max}-\theta_{\min}$ satisfies
\begin{equation*}
s(X)\le(2N)^{1/2}|X|_{\mathfrak{g}}<(2N)^{1/2}\frac{\sqrt6}{N}=\Bigl(\frac{12}{N}\Bigr)^{1/2}
\le\sqrt6<2\pi .
\end{equation*}
Hence $\mathfrak{B}_z\subset\mathfrak{D}$. By Definition~\ref{4.36:def-haar-balls}, $X\mapsto e^{iX}$ is
an analytic diffeomorphism of $\mathfrak{D}$ onto $G\setminus\mathrm{Cut}$; its restriction to the
open set $\mathfrak{B}_z$ is therefore an analytic diffeomorphism onto an open image, which we call
$O_G$. This can also be seen directly: since every $|\theta_p|<\pi$, the principal logarithm inverts
$X\mapsto e^{iX}$ analytically on the image, and the differential has $|\det|=J>\frac12>0$ by the
bound proved next. By the previous step, $O_G\supseteq\{u:|u-\mathbf{1}|<\varepsilon_G\}$.

\emph{(i), the Jacobian.} By Definition~\ref{4.36:def-haar-balls},
$J(X)=\prod_{p<q}\operatorname{sinc}^2((\theta_p-\theta_q)/2)\le1$. For $p<q$ put
$a_{pq}:=(\theta_p-\theta_q)^2/12$. By (F2) of Lemma~\ref{4.36:lem-eigenangle-inequalities},
\begin{equation*}
a_{pq}\le\frac{2N|X|^2_{\mathfrak{g}}}{12}<\frac{2N\cdot6/N^2}{12}=\frac1N\le\frac12,
\qquad
\sum_{p<q}a_{pq}=\frac{N^2|X|^2_{\mathfrak{g}}}{12}<\frac{N^2\cdot6/N^2}{12}=\frac12,
\end{equation*}
the identity $\sum_{p<q}(\theta_p-\theta_q)^2=N^2|X|^2_{\mathfrak{g}}$ for traceless $X$ being the
second part of (F2). By (F4) of Lemma~\ref{4.36:lem-eigenangle-inequalities}, each factor satisfies
$\operatorname{sinc}^2((\theta_p-\theta_q)/2)\ge1-a_{pq}\ge0$; by (F5) of the same lemma,
\begin{equation*}
J(X)\ge\prod_{p<q}(1-a_{pq})\ge1-\sum_{p<q}a_{pq}>\frac12 .
\end{equation*}
Since $\mathfrak{B}_z\subset\mathfrak{D}$, the identity \eqref{4.36:eq-haar-balls-a} with $u=\mathbf{1}$
and Borel sets $S\subset\mathfrak{B}_z$ shows that on $\mathfrak{B}_z$ the Haar probability is
$\sigma\,dX$ with $\sigma=\sigma_NJ$, and the bounds on $J$ give
$\frac12\sigma_N<\sigma\le\sigma_N$; finally $\sigma(0)=\sigma_NJ(0)=\sigma_N$ because $J(0)=1$.

\emph{(ii).} Let $0<g\le\gamma_z$ and put $\varepsilon:=g\,p_0(g)$; by (o),
$0<\varepsilon\le\varepsilon_G$. By (i), $\{u:|u-\mathbf{1}|<\varepsilon\}\subset O_G$, so
\begin{equation*}
\{u\in G:|u-\mathbf{1}|<\varepsilon\}=\exp i(\mathcal{V}_g),\qquad
\mathcal{V}_g:=\{X\in\mathfrak{B}_z:\ |e^{iX}-\mathbf{1}|<\varepsilon\},
\end{equation*}
and $\mathcal{V}_g$ is a Borel (indeed open) subset of $\mathfrak{B}_z\subset\mathfrak{D}$. Applying
\eqref{4.36:eq-haar-balls-a} with $u=\mathbf{1}$, $S=\mathcal{V}_g$ and
$h(u')=e^{-g^{-2}(1-\operatorname{Re}\operatorname{tr}u')}$ for $u'\in G$, so that
$h(e^{iX})=e^{-\Phi(X)/g^2}$ with $\Phi$ as below, to \eqref{4.36:eq-one-link-partition-1} gives
\begin{equation*}
z(g)=\sigma_N\int_{\mathcal{V}_g}J(X)\,e^{-\Phi(X)/g^2}\,dX,\qquad
\Phi(X):=1-\operatorname{Re}\operatorname{tr}e^{iX}=\frac1N\sum_p(1-\cos\theta_p),
\end{equation*}
the last identity because $e^{iX}$ has eigenvalues $e^{i\theta_p}$ and $\operatorname{tr}$ is the
normalised trace. On $\mathfrak{B}_z$ every $|\theta_p|<\pi$, so (F3) of
Lemma~\ref{4.36:lem-eigenangle-inequalities} applies to each $\theta_p$; averaging over $p$ and using
$|X|^2_{\mathfrak{g}}=N^{-1}\sum_p\theta_p^2$ gives
\begin{equation*}
\frac{2}{\pi^2}|X|^2_{\mathfrak{g}}\le\Phi(X)\le\frac12|X|^2_{\mathfrak{g}}
\qquad(X\in\mathfrak{B}_z).
\end{equation*}

\emph{Upper bound.} Using $J\le1$, the lower bound on $\Phi$, and enlarging $\mathcal{V}_g$ to
$\mathfrak{g}$, and then the Gaussian integral
$\int_{\mathbb{R}^D}e^{-a|x|^2}\,dx=(\pi/a)^{D/2}$ with $a=2/(\pi^2g^2)$ on
$(\mathfrak{g},\operatorname{tr})\cong\mathbb{R}^D$,
\begin{equation*}
z(g)\le\sigma_N\int_{\mathfrak{g}}e^{-\frac{2}{\pi^2}|X|^2_{\mathfrak{g}}/g^2}\,dX
=\sigma_N\Bigl(\frac{\pi\cdot\pi^2g^2}{2}\Bigr)^{D/2}=\sigma_Nv_2\,g^D .
\end{equation*}

\emph{Lower bound.} We claim $\{X\in\mathfrak{g}:|X|_{\mathfrak{g}}\le g\}\subset\mathcal{V}_g$.
Indeed, such $X$ lie in $\mathfrak{B}_z$ because $g\le r_G/2<r_G$ by (o); and by (F6) and (F1) of
Lemma~\ref{4.36:lem-eigenangle-inequalities} together with $p_0(g)\ge2\mathsf{n}$ from (o),
\begin{equation*}
|e^{iX}-\mathbf{1}|\le|X|\le\mathsf{n}|X|_{\mathfrak{g}}\le\mathsf{n}g\le\tfrac12g\,p_0(g)
=\frac\varepsilon2<\varepsilon .
\end{equation*}
This is the only point at which a ball for $|\cdot|_{\mathfrak{g}}$ has to be contained in the window
defined by the operator norm; it costs the factor $\mathsf{n}$, which is absorbed in the condition
$p_0(g)\ge2\mathsf{n}$ in the definition of $\gamma_z$. On this ball
$\Phi(X)/g^2\le\frac12|X|^2_{\mathfrak{g}}/g^2\le\frac12$ and $J(X)>\frac12$ by (i); its
$dX$-volume is $\omega_Dg^D$. Restricting the integral to the ball therefore gives
$z(g)\ge\frac12\sigma_Ne^{-1/2}\omega_Dg^D$, which completes \eqref{4.36:eq-one-link-partition-3}.

\emph{Logarithmic form.} If $0<g_j\le\gamma_z$, taking logarithms in
\eqref{4.36:eq-one-link-partition-3} at $g=g_j$ and putting $\zeta_j:=\log z_j-D\log g_j$ gives
\begin{equation*}
\log\bigl(\tfrac12\sigma_Ne^{-1/2}\omega_D\bigr)\le\zeta_j\le\log(\sigma_Nv_2),
\end{equation*}
so $|\zeta_j|$ is at most the larger of the absolute values of the two ends, which is $c_G(N)$;
this is \eqref{4.36:eq-one-link-partition-4}.

\emph{(iii).} Under the isometry $(\mathfrak{g},\operatorname{tr})\cong\mathbb{R}^D$, the same Gaussian
integral with $a=\frac12$ gives $\int_{\mathfrak{g}}e^{-\frac12|X|^2_{\mathfrak{g}}}\,dX
=(\pi/\tfrac12)^{D/2}=(2\pi)^{D/2}$, whence $\nu_{\mathfrak{g}}=\frac D2\log2\pi$.

Finally, if the Lebesgue measure $dX$ on $\mathfrak{g}$ is replaced by $c\,dX$ with $c>0$, then
$\sigma_N$ is replaced by $\sigma_N/c$, the volume $\omega_D$ and the Gaussian integral $v_2g^D$ of
the upper bound are multiplied by $c$, and $\nu_{\mathfrak{g}}$ is replaced by
$\nu_{\mathfrak{g}}+\log c$; consequently $\sigma_N\omega_D$, $\sigma_Nv_2$, $c_G(N)$, $\zeta_j$ and
$\log\sigma_N+\nu_{\mathfrak{g}}$ do not depend on the normalisation of $dX$.
\end{proof}

\begin{definition}[Cell and bond counts and the one-step normalisation numbers]
\label{4.36:def-step-normalisation}
Write $\mathbb{T}^{(i)}$ for the block lattice of level $i$ (Definition~\ref{1.5:def-block-average}),
so that $\mathbb{T}^{(0)}$ is the lattice on which the steps act; an \emph{$i$-cell} is a cell of
$\mathbb{T}^{(i)}$ and an \emph{$i$-bond} is a bond of $\mathbb{T}^{(i)}$. Recall $d=4$ and
$D=N^2-1=\dim\mathfrak{g}$.

\emph{Cell and bond counts.} Let $Y$ be a union of $(j+1)$-cells. We put
\begin{equation*}
|Y|_j:=\#\bigl(Y\cap \mathbb{T}^{(j)}\bigr),
\end{equation*}
the number of points of $\mathbb{T}^{(j)}$ in $Y$. Each $(j+1)$-cell is the union of $L^{d}$ cells of
$\mathbb{T}^{(j)}$, so that $|\mathbb{T}^{(0)}|_{j+1}=L^{-d}|\mathbb{T}^{(0)}|_j$. For every level $i$
we let $\mathfrak{b}_i(Y)$ be the number of bonds of $\mathbb{T}^{(i)}$ with at least one end-point
in $Y$. Let $Y^{(j)*}$ be the set of $j$-bonds of $Y$ from which, for every $(j+1)$-bond $c$ of $Y$,
the one $j$-bond that \cite[(1.9)]{B-CMP119} associates with $c$ has been removed. In this notation
\begin{equation*}
\bigl|Y^{(j)*}\bigr|=\mathfrak{b}_j(Y)-\mathfrak{b}_{j+1}(Y).
\end{equation*}

\emph{Labels.} A \emph{label} at level $k$ is a family $\mathcal{S}=(\Omega_j,\Lambda_j)_{1\le j\le k}$
of nested unions of cubes. The label in which no region is ever large is written $\emptyset$; it has
$\Omega_j=\Lambda_j=\mathbb{T}^{(0)}$ for every $j$. For a label we put $Z_j:=\Lambda_j^{c}$. With
$R_j$ the level-$j$ radius of Lemma~\ref{4.4:lem-collar-avoidance}, we let $\Lambda^{0}_j$ be
$\Lambda_j$ with one layer of $MR_j$-cubes removed, with the convention that nothing is removed from
the whole lattice, $(\mathbb{T}^{(0)})^{0}=\mathbb{T}^{(0)}$, and we put
$Z^{+}_j:=\mathbb{T}^{(0)}\setminus\Lambda^{0}_j$. For a region $Y\subset\mathbb{T}^{(0)}$ we write
$Y^{(j)}:=Y\cap\mathbb{T}^{(j)}$. Further,
\begin{equation*}
\Gamma_0=\Omega_1^{c},\qquad
\Gamma_j=\bigl(\Omega_j\setminus\Omega_{j+1}\bigr)^{(j)}\quad(1\le j\le k-1),\qquad
\Gamma_k=\Omega_k^{(k)}.
\end{equation*}
A label is \emph{geometric} if every $Z_j$ is a union of $MR_j$-cubes.

\emph{The one-step normalisation numbers.} Let $z(g)$ be the one-link partition number of
Lemma~\ref{4.36:lem-one-link-partition} and
$\nu_{\mathfrak{g}}=\log\int_{\mathfrak{g}}e^{-|X|^2_{\mathfrak{g}}/2}\,dX=(D/2)\log 2\pi$ the Gaussian
constant of clause (iii) of that lemma, let $\sigma_N$ be the normaliser of the
Haar density in exponential coordinates, and put $z_j:=z(g_j)$. Let $0\le j<K$, and let
$\Lambda\subset\Omega$ be two unions of cubes of $\pi_{j+1}$, the partition of the lattice into the
cubes of level $j+1$ of \cite[\S3]{B-CMP119}. We define
\begin{equation}\label{4.36:eq-step-normalisation-1}
\begin{split}
\mathsf{N}^{(j)}(\Omega,\Lambda)
:={}&\bigl(D\log g_j+\log\sigma_N\bigr)\,\bigl|\Omega^{(j)*}\bigr|
-(L^{d}-1)\log z_j\cdot|\Omega|_{j+1}\\
&+\nu_{\mathfrak{g}}\,\bigl|\Lambda^{(j)*}\bigr|,
\end{split}
\end{equation}
and
\begin{equation}\label{4.36:eq-step-normalisation-2}
E^{(j)}(\Omega,\Lambda):=\mathsf{N}^{(j)}(\Omega,\Lambda)
+\mathbf{E}^{(j+1)}(\Lambda,1)+\mathbf{R}^{(j+1)}(\Lambda,1),
\end{equation}
where $\mathbf{E}^{(j+1)}(\Lambda,1)$ and $\mathbf{R}^{(j+1)}(\Lambda,1)$ are the terms so denoted in
\cite[\S3]{B-CMP119}, with the arguments as written there.
\end{definition}

The numbers \eqref{4.36:eq-step-normalisation-1} and \eqref{4.36:eq-step-normalisation-2} are
Ba{\l}aban's one-step numbers, taken as printed in the first exponential of
\cite[(3.15), p.~267]{B-CMP119} and in \cite[(3.33), p.~273]{B-CMP119} and written in the notation
of this book; they are not re-derived here. The gauge group enters them only through $D$,
$\sigma_N$, $\nu_{\mathfrak{g}}$ and $z_j$.

\begin{lemma}[The energy constant of a label; polymer-sum and bond-count bounds]
\label{4.36:lem-label-energy}
We use the cell counts $|Y|_j$, the bond counts $\mathfrak{b}_i(Y)$, the bond sets $Y^{(j)*}$,
the labels $\mathcal{S}=(\Omega_j,\Lambda_j)_{1\le j\le k}$, the whole torus $T$ and the
one-step normalisation numbers $\mathsf{N}^{(j)}(\Omega,\Lambda)$ and $E^{(j)}(\Omega,\Lambda)$
of Definition~\ref{4.36:def-step-normalisation}. We write $T^{(j)}$ for the lattice of level $j$,
$\pi_j$ for the partition of $T^{(j)}$ into the cubes of level $j$ of \cite{B-CMP119},
$\mathbf{D}_j$ for the polymers of level $j$, which are the face-connected unions of cubes of
$\pi_j$ (as in \cite{B-CMP119}), $d_j(X)$ for
Ba{\l}aban's size of a polymer $X$, and $|X|_{\square}$ for the number of cubes of $\pi_j$
contained in $X$ (the subscript $\square$ is notational: $|X|_{\square}$ does not depend on a
choice of cube).
\begin{enumerate}
\item[(a)] For $\Lambda$ a union of cubes of $\pi_j$ which is either $T$ or the region
$\Lambda_j$ of a label we put
\begin{equation}\label{4.36:eq-label-energy-1}
\begin{aligned}
\mathbf{E}^{(j)}(\Lambda,1)&:=\sum_{z\in\Lambda^{0}}\ \sum_{X\in\mathbf{D}_j,\ z\in X\subset\Lambda}
\mathbf{E}^{(j)}(X,1,z),\\
\mathbf{R}^{(j)}(\Lambda,1)&:=\sum_{X\in\mathcal{R}(\Lambda)}\mathbf{R}^{(j)}(X,1),
\end{aligned}
\end{equation}
where $\mathbf{E}^{(j)}(X,1,z)$ and $\mathbf{R}^{(j)}(X,1)$ are Ba{\l}aban's terms of
\cite{B-CMP119} evaluated at the configuration written $1$, all of whose bond variables are the
identity, $\mathcal{R}(\Lambda)$ is the family of subsets
$X\subset\Lambda$ over which the sum of the terms $\mathbf{R}^{(j)}$ ranges in
\cite{B-CMP119}, and $\Lambda^{0}$ is $\Lambda^{0}_j$, the set obtained from $\Lambda_j$ by
removing one layer of $MR_j$-cubes, when $\Lambda=\Lambda_j$, while $T^{0}=T$; here the
$MR_j$-cubes are the cubes of side $MR_j$ of which the complement $T\setminus\Lambda_j$ of the
region $\Lambda_j$ of a label is a union, $R_j$ being the level-$j$ factor in that side.
The \emph{energy constant} $E$ and the \emph{energy constant of the
label} $\mathcal{S}=(\Omega_j,\Lambda_j)_{1\le j\le k}$ are
\begin{equation}\label{4.36:eq-label-energy-a}
\begin{aligned}
E&:=\sum_{j=0}^{K-1}E^{(j)}(T,T),\\
E_k[\mathcal{S}]&:=E-\sum_{j=0}^{k-1}E^{(j)}(\Omega_{j+1},\Lambda_{j+1}).
\end{aligned}
\end{equation}
\item[(b)] Put
\begin{equation}\label{4.36:eq-label-energy-2}
\kappa(d):=3\cdot 2^{d+1}\,(2\log(2d)+1),
\qquad
K_{*}:=1+\sum_{n\le 2^{d+1}}n\,(2d)^{2(n-1)}.
\end{equation}
If $\kappa\ge\kappa(d)$, then for every $\square\in\pi_j$ and every $z\in T^{(j)}$
\begin{equation}\label{4.36:eq-label-energy-3}
\sum_{X\in\mathbf{D}_j,\ X\supset\square}|X|_{\square}\,e^{-\kappa d_j(X)}\le K_{*},
\qquad
\sum_{X\in\mathbf{D}_j,\ X\ni z}|X|_{\square}\,e^{-\kappa d_j(X)}\le K_{*}.
\end{equation}
\item[(c)] Let $Y\subset Y'$ be unions of $(j+1)$-cells. Then
\begin{equation}\label{4.36:eq-label-energy-4}
\bigl|\,|Y'^{(j)*}|-|Y^{(j)*}|\,\bigr|\le 2dL^{d}\,|Y'\setminus Y|_{j+1},
\qquad
|Y^{(j)*}|\le 2dL^{d}\,|Y|_{j+1},
\end{equation}
and
\begin{equation}\label{4.36:eq-label-energy-5}
|T^{(j)*}|=d\,(L^{d}-1)\,|T|_{j+1}.
\end{equation}
\end{enumerate}
\end{lemma}

\begin{proof}
Part (a) consists of definitions. That the number $E_k[\mathcal{S}]$ of
\eqref{4.36:eq-label-energy-a} is, at every step $k$ and for every label $\mathcal{S}$, the
constant of \cite[(2.23)]{B-CMP119} is the content of the lemma of this chapter identifying the
energy constant with that constant; that lemma also shows, by an
argument of homogeneity, that the identification is not circular.

(b) Fix $\square\in\pi_j$. For $n\ge1$ let $s_n$ be the sum of $|X|_{\square}e^{-\kappa d_j(X)}$
over those $X\in\mathbf{D}_j$ with $X\supset\square$ and $|X|_{\square}=n$; the left-hand side of
the first inequality in \eqref{4.36:eq-label-energy-3} is $\sum_{n\ge1}s_n$.

Every polymer entering $s_n$ is a face-connected family of $n$ cubes of $\pi_j$ through
$\square$; we bound the number of such families by $(2d)^{2(n-1)}$. Indeed, each cube of
$\pi_j$ has at most $2d$ face neighbours; fix an enumeration of the $2d$ face directions. Given a
face-connected family $F$ of $n$ cubes containing $\square$, choose a spanning tree of the
face-adjacency graph on $F$; its depth-first traversal from $\square$ crosses each of the $n-1$
edges of the tree twice, hence is a sequence of $2(n-1)$ moves, each in one of $2d$ directions,
and the family $F$ is the set of cubes visited. Fixing for each such family $F$ one spanning
tree and one depth-first traversal, the map assigning to $F$ the resulting sequence of $2(n-1)$
directions is injective, since $F$ is the set of cubes visited; hence there are at most
$(2d)^{2(n-1)}$ such families.

Next, by the definition of Ba{\l}aban's size $d_j(X)$ in \cite{B-CMP119}, a polymer
$X\in\mathbf{D}_j$ consisting
of $n$ cubes satisfies $n\le 2^{d}(3d_j(X)+1)$,
that is,
\begin{equation}\label{4.36:eq-label-energy-6}
d_j(X)\ge \tfrac13\bigl(n2^{-d}-1\bigr).
\end{equation}

If $n\le 2^{d+1}$ we use $e^{-\kappa d_j(X)}\le 1$ (since $d_j(X)\ge0$) and obtain
$s_n\le n(2d)^{2(n-1)}$. If $n>2^{d+1}$, then $1<n2^{-d-1}$, so
$n2^{-d}-1>n2^{-d}-n2^{-d-1}=n2^{-d-1}$, and by \eqref{4.36:eq-label-energy-6} and
$\kappa\ge\kappa(d)$,
\begin{equation*}
\kappa d_j(X)\ge\frac{\kappa\,n}{3\cdot2^{d+1}}\ge\frac{\kappa(d)\,n}{3\cdot 2^{d+1}}
=n\,(2\log(2d)+1).
\end{equation*}
Hence, for $n>2^{d+1}$,
\begin{equation*}
s_n\le n\,(2d)^{2(n-1)}e^{-n(2\log(2d)+1)}=n\,(2d)^{-2}e^{-n}\le n\,e^{-n}.
\end{equation*}
Summing,
\begin{equation*}
\sum_{n\ge1}s_n\le\sum_{n\le2^{d+1}}n\,(2d)^{2(n-1)}+\sum_{n\ge1}n\,e^{-n}.
\end{equation*}
Finally $\sum_{n\ge1}n\,e^{-n}=e/(e-1)^2<1$, because $(e-1)^2-e=e^2-3e+1>0$, the larger root of
$x^2-3x+1$ being $(3+\sqrt5)/2<e$. This gives the first inequality in
\eqref{4.36:eq-label-energy-3}. For the second, let $\square_z$ be the cube of $\pi_j$ containing
$z$; since $\pi_j$ partitions $T^{(j)}$ and every $X\in\mathbf{D}_j$ is a union of cubes of
$\pi_j$, we have $z\in X$ if and only if $X\supset\square_z$, and the second sum is the first sum
for $\square=\square_z$.

(c) By Definition~\ref{4.36:def-step-normalisation},
$|Y^{(j)*}|=\mathfrak{b}_j(Y)-\mathfrak{b}_{j+1}(Y)$ for every union $Y$ of $(j+1)$-cells. We first
show that for $i\in\{j,j+1\}$
\begin{equation}\label{4.36:eq-label-energy-7}
0\le\mathfrak{b}_i(Y')-\mathfrak{b}_i(Y)\le 2d\,|Y'\setminus Y|_i .
\end{equation}
A bond of $T^{(i)}$ with an end-point in $Y$ has an end-point in $Y'\supset Y$; this gives the
lower bound. A bond counted in $\mathfrak{b}_i(Y')$ but not in $\mathfrak{b}_i(Y)$ has an
end-point in $Y'$ and none in $Y$, hence an end-point in $(Y'\setminus Y)\cap T^{(i)}$; each
point of $T^{(i)}$ is an end-point of $2d$ bonds, which gives the upper bound.

Now
\begin{equation*}
|Y'^{(j)*}|-|Y^{(j)*}|=\bigl(\mathfrak{b}_j(Y')-\mathfrak{b}_j(Y)\bigr)
-\bigl(\mathfrak{b}_{j+1}(Y')-\mathfrak{b}_{j+1}(Y)\bigr).
\end{equation*}
By \eqref{4.36:eq-label-energy-7} the first bracket lies in $[0,2d|Y'\setminus Y|_j]$ and the
second in $[0,2d|Y'\setminus Y|_{j+1}]$, so the difference lies in
$[-2d|Y'\setminus Y|_{j+1},\,2d|Y'\setminus Y|_j]$. Since $Y'\setminus Y$ is a union of
$(j+1)$-cells, $|Y'\setminus Y|_j=L^{d}|Y'\setminus Y|_{j+1}$, and as $L^d\ge1$ both ends are
bounded in absolute value by $2dL^{d}|Y'\setminus Y|_{j+1}$. This is the first inequality in
\eqref{4.36:eq-label-energy-4}. The second follows from the first with $\emptyset$ in place of
$Y$ and $Y$ in place of $Y'$, because $\mathfrak{b}_i(\emptyset)=0$ gives
$|\emptyset^{(j)*}|=0$.

For \eqref{4.36:eq-label-energy-5}: on the periodic lattice $T^{(i)}$ the number of bonds is $d$
times the number of points, so $\mathfrak{b}_i(T)=d|T|_i$; with $|T|_j=L^{d}|T|_{j+1}$,
\begin{equation*}
|T^{(j)*}|=\mathfrak{b}_j(T)-\mathfrak{b}_{j+1}(T)=d|T|_j-d|T|_{j+1}=d\,(L^{d}-1)\,|T|_{j+1}.
\qedhere
\end{equation*}
\end{proof}

\begin{theorem}[The exponent of the never-large term and label corrections at $SU(N)$]
\label{4.36:thm-never-large-exponent}
Let $d=4$ and $D=N^2-1=\dim\mathfrak{g}$. Let $T$ be the torus, $T^{(j)}$ its lattice at level $j$,
$|Y|_j$ and $Y^{(j)*}$ the cell and bond counts, and $\mathsf{N}^{(j)}(\Omega,\Lambda)$,
$E^{(j)}(\Omega,\Lambda)$ the one-step normalisation numbers and energies of
Definition~\ref{4.36:def-step-normalisation}. Let $\sigma_N$, $\nu_{\mathfrak{g}}$, $\zeta_j$, $c_G(N)$
and $\gamma_z=\gamma_z(N,A_0,p_0)$ be as in Lemma~\ref{4.36:lem-one-link-partition}, and let
$K_*$, $\kappa(d)$, the energy constant $E$ and its truncations $E_k[\mathcal{S}]$ at a label
$\mathcal{S}=(\Omega_j,\Lambda_j)_{1\le j\le k}$ be as in Lemma~\ref{4.36:lem-label-energy} and
\eqref{4.36:eq-label-energy-a}; $\emptyset$ denotes the never-large label, $\Omega_j=\Lambda_j=T$
for all $j$. For a label, $Z_j:=\Lambda_j^{c}$; $\Lambda^{0}_j$ is $\Lambda_j$ with one layer of
cubes of side $MR_j$ removed, so that $T^{0}=T$; $Z^{+}_j:=T\setminus\Lambda^{0}_j$; and
$\Gamma_0:=\Omega_1^{c}$, $\Gamma_n:=(\Omega_n\setminus\Omega_{n+1})^{(n)}$ for $0<n<k$,
$\Gamma_k:=\Omega_k^{(k)}$. Let $E_0=E_0(N,L,\mathfrak{p})$ be a constant, let $\kappa\ge\kappa(d)$
and $\kappa_0>0$, and consider, for $0\le j\le K$ ($0\le j<K$ in the increment condition) and at
$(\mathbf{U},\mathbf{J})=(1,0)$, the conditions
\begin{equation}\label{4.36:eq-never-large-exponent-a}
\begin{gathered}
|\mathbf{E}^{(j)}(X,1,z)|\le E_0\,e^{-\kappa d_j(X)},\qquad
|\mathbf{R}^{(j)}(X,1)|\le g_j^{\kappa_0}\,e^{-\kappa d_j(X)},\\
0<g_j\le\gamma<1,\qquad |x_{j+1}-x_j|=|\beta_{j+1}|\le\bar{B},\qquad
\gamma\le\gamma_z(N,A_0,p_0).
\end{gathered}
\end{equation}
Assume:
\begin{itemize}
\item[(h1)] the effective density $\rho_k$ has the format of
\cite[(2.18), (2.23), (2.25)--(2.27), (2.30)]{B-CMP119}, with the energy constant
\eqref{4.36:eq-label-energy-a};
\item[(h2)] the two bounds in the first line of \eqref{4.36:eq-never-large-exponent-a} hold, for all
polymers $X$ and points $z$; these are \cite[(1.18)]{B-CMP109}, from
\cite[Theorem~1]{B-CMP109}, which is stated for an arbitrary semisimple compact Lie group
$G\subset U(N)$, and \cite[(2.31)]{B-CMP119};
\item[(h3)] the label-locality of the terms $\mathbf{E}^{(j)}$, $\mathbf{R}^{(j)}$ of
\cite{B-CMP119}, and, for the primed remainder $\mathbf{R}'$, the locality remark attached to the
theorem on the exponent of the never-large term and label corrections at $SU(2)$ (in what
follows, the $SU(2)$ theorem);
\item[(h4)] the three conditions in the second line of \eqref{4.36:eq-never-large-exponent-a} hold.
\end{itemize}
Put
\begin{align*}
c^{\sharp}_G(N)&:=d\,|\log\sigma_N+\nu_{\mathfrak{g}}|+c_G(N),\qquad
C_{\mathfrak{r}}:=c^{\sharp}_G+\frac{K_*(E_0+1)}{L^d-1},\\
C_{\log}(N)&:=\tfrac12(2d+1)DL^d=\tfrac92(N^2-1)L^4,\\
C^{\mathrm{lab}}_{\mathrm{loc}}(N)&:=L^d\bigl(2d|\log\sigma_N|+2d|\nu_{\mathfrak{g}}|+c_G(N)\bigr)+K_*(2E_0+1),\\
C_{\mathrm{lab}}&:=C_{\mathfrak{r}}+2\cdot 3^d\,C^{\mathrm{lab}}_{\mathrm{loc}}.
\end{align*}
Then the following hold.
\begin{itemize}
\item[(a)] (Under (h1), (h2), (h4) only.) For $0\le k\le K$ put
$\mathfrak{E}_k:=E_k[\emptyset]$ and
\begin{equation*}
\mathfrak{l}_k:=\tfrac12(d-1)D(L^d-1)\sum_{j=k}^{K-1}\log x_j\,|T|_{j+1}.
\end{equation*}
Then $\mathfrak{E}_k=\sum_{j=k}^{K-1}E^{(j)}(T,T)=-\mathfrak{l}_k+\mathfrak{r}_k$ with
$|\mathfrak{r}_k|\le C_{\mathfrak{r}}|T|_k$; $\mathfrak{E}_K=\mathfrak{l}_K=\mathfrak{r}_K=0$; and for
$0\le k<K$
\begin{equation}\label{4.36:eq-never-large-exponent-b}
\begin{split}
\tfrac12(d-1)(N^2-1)(1-L^{-d})\log x_k\,|T|_k
&\le\mathfrak{l}_k\\
&\le\tfrac12(d-1)(N^2-1)\Bigl(\log x_k+\frac{\bar{B}\gamma^2}{L^d-1}\Bigr)|T|_k .
\end{split}
\end{equation}
At $d=4$, $\tfrac12(d-1)D=\tfrac32(N^2-1)$.
\item[(b)] For the logarithmic and regular parts $\mathfrak{e}^{\log}_j$, $\mathfrak{e}^{\mathrm{reg}}_j$
of the label corrections of the $SU(2)$ theorem, taken at $SU(N)$,
\begin{equation*}
|\mathfrak{e}^{\log}_j|\le C_{\log}(N)\log x_{j-1}\,|\Omega_j^{c}|_j,\qquad
|\mathfrak{e}^{\mathrm{reg}}_j|\le C^{\mathrm{lab}}_{\mathrm{loc}}(N)\,|Z^{+}_j|_j,
\end{equation*}
and $|Z^{+}_j|_j\le 3^d|Z_j|_j$ when $Z_j$ is a union of cubes of side $MR_j$.
\item[(c)] With $\mathfrak{L}_k[\mathcal{S}]:=-\mathfrak{l}_k+\sum_j\mathfrak{e}^{\log}_j$,
\begin{equation*}
\bigl|E_k[\mathcal{S}]-\mathfrak{L}_k[\mathcal{S}]\bigr|\le C_{\mathrm{lab}}\sum_{n=0}^{k}|\Gamma_n|.
\end{equation*}
\item[(d)] The two-label bound of part (d) of the $SU(2)$ theorem holds at $SU(N)$, with
$D$, $\sigma_N$, $\nu_{\mathfrak{g}}$ and $c_G(N)$ in place of that theorem's group quantities.
\end{itemize}
Parts (b)--(d) use (h3), that is, the locality remark for $\mathbf{R}'$.
\end{theorem}

Each of (h1)--(h4) is a clause, or part of the hypotheses, of the statement of the format of the
effective densities at $SU(N)$. The condition $|x_{j+1}-x_j|\le\bar{B}$ in (h4) is implied one step
at a time by the flow inequalities of Notation~\ref{4.36:not-flow-thresholds}: the lower and the
upper one give, respectively,
\begin{equation*}
\begin{gathered}
\beta_j\ge\lambda^{\sharp}-c_5\ge-c_5\ge-\bar{B},\\
\beta_j\le B^{\sharp}\le\bar{B}.
\end{gathered}
\end{equation*}
The positivity of the drift, $\lambda^{\sharp}>0$ (asymptotic freedom), is not used. The flow
inequalities enter only through this per-step bound: in the upper estimate of
\eqref{4.36:eq-never-large-exponent-b}, and in the inequality
$\log x_{j-1}\le\log x_j+\bar{B}\gamma^2$ of the corollary at $SU(N)$ in which it is stated.

\begin{proof}
\emph{Step 1: the normalisation number of the never-large pair.} Fix $0\le j<K$. In
Definition~\ref{4.36:def-step-normalisation} take $\Omega=\Lambda=T$. Then
$|\Omega^{(j)*}|=|\Lambda^{(j)*}|=|T^{(j)*}|$, and by the bond-count bounds of
Lemma~\ref{4.36:lem-label-energy}, $|T^{(j)*}|=d(L^d-1)|T|_{j+1}$. Since $g_j\le\gamma\le\gamma_z$ by
(h4), Lemma~\ref{4.36:lem-one-link-partition}(ii) applies at $g=g_j$ and gives
$\log z_j=D\log g_j+\zeta_j$ with $|\zeta_j|\le c_G(N)$. Substituting both facts into the definition
of $\mathsf{N}^{(j)}(T,T)$, collecting the terms in $|T|_{j+1}$, and using $\log g_j=-\tfrac12\log x_j$
(which is the definition $x_j=g_j^{-2}$ of Notation~\ref{4.36:not-flow-thresholds}), we obtain
\begin{equation}\label{4.36:eq-never-large-exponent-1}
\begin{split}
\mathsf{N}^{(j)}(T,T)
&=(D\log g_j+\log\sigma_N+\nu_{\mathfrak{g}})\,d(L^d-1)|T|_{j+1}
-(L^d-1)(D\log g_j+\zeta_j)|T|_{j+1}\\
&=(L^d-1)|T|_{j+1}\bigl\{(d-1)D\log g_j+d(\log\sigma_N+\nu_{\mathfrak{g}})-\zeta_j\bigr\}\\
&=(L^d-1)|T|_{j+1}\bigl\{-\tfrac12(d-1)D\log x_j+d(\log\sigma_N+\nu_{\mathfrak{g}})-\zeta_j\bigr\}.
\end{split}
\end{equation}
The coefficient $(d-1)D$ of $\log g_j$ records that, of the $d\cdot D$ Gaussian directions per site
carried by the starred bonds, $(d-1)\cdot D$ survive gauge fixing. This is the only place in the proof
of (a) where the group enters.

\emph{Step 2: the terms $\mathbf{E}^{(j+1)}(T,1)$ and $\mathbf{R}^{(j+1)}(T,1)$.} These are the sums of
Lemma~\ref{4.36:lem-label-energy} at level $j+1$ with $\Lambda=T$; here $T^{0}=T$, since $T$ has no
boundary layer to remove.

By the first bound of (h2) at level $j+1$, the sum over polymers $X\ni z$ of
$|\mathbf{E}^{(j+1)}(X,1,z)|$ is at most $E_0\sum_{X\ni z}e^{-\kappa d_{j+1}(X)}$. Since
$\kappa\ge\kappa(d)$, the polymer-sum bound of Lemma~\ref{4.36:lem-label-energy} (the form with the sum
over $X\ni z$) bounds this by $E_0K_*$. Summing over the $|T|_{j+1}$ points $z$ gives
$|\mathbf{E}^{(j+1)}(T,1)|\le K_*E_0|T|_{j+1}$.

For the second term, by the second bound of (h2) and $g_{j+1}^{\kappa_0}\le1$ (because
$g_{j+1}\le\gamma<1$ by (h4)), each summand is at most $e^{-\kappa d_{j+1}(X)}$. Every polymer $X$
contains a cube $\square$ of the cover $\pi_{j+1}$ of Lemma~\ref{4.36:lem-label-energy}, and its
cube count satisfies $|X|_{\square}\ge1$. Hence the sum over $X$ is at most
the sum over the cubes $\square\subset T$ of
$\sum_{X\supset\square}|X|_{\square}e^{-\kappa d_{j+1}(X)}$. By the polymer-sum bound each of these
inner sums is at most $K_*$. The cubes of $\pi_{j+1}$ in $T$ number $M^{-d}|T|_{j+1}$, so
\begin{equation*}
|\mathbf{R}^{(j+1)}(T,1)|\le K_*M^{-d}|T|_{j+1}\le K_*|T|_{j+1}.
\end{equation*}

\emph{Step 3: the one-step energy of the never-large pair.} By
Definition~\ref{4.36:def-step-normalisation},
$E^{(j)}(T,T)=\mathsf{N}^{(j)}(T,T)+\mathbf{E}^{(j+1)}(T,1)+\mathbf{R}^{(j+1)}(T,1)$. Inserting
\eqref{4.36:eq-never-large-exponent-1} gives
\begin{equation}\label{4.36:eq-never-large-exponent-2}
E^{(j)}(T,T)=-\tfrac12(d-1)D(L^d-1)\log x_j\,|T|_{j+1}+r_j,
\end{equation}
where
\begin{equation*}
r_j:=(L^d-1)|T|_{j+1}\bigl\{d(\log\sigma_N+\nu_{\mathfrak{g}})-\zeta_j\bigr\}
+\mathbf{E}^{(j+1)}(T,1)+\mathbf{R}^{(j+1)}(T,1).
\end{equation*}
By $|\zeta_j|\le c_G(N)$, the definition of $c^{\sharp}_G$ and Step~2,
\begin{equation}\label{4.36:eq-never-large-exponent-3}
|r_j|\le\bigl[(L^d-1)c^{\sharp}_G+K_*(E_0+1)\bigr]|T|_{j+1}.
\end{equation}

\emph{Step 4: telescoping and the remainder.} By \eqref{4.36:eq-label-energy-a} at the never-large label,
$E_k[\emptyset]=E-\sum_{j=0}^{k-1}E^{(j)}(T,T)$ with $E=\sum_{j=0}^{K-1}E^{(j)}(T,T)$. Hence
$\mathfrak{E}_k=\sum_{j=k}^{K-1}E^{(j)}(T,T)$. By \eqref{4.36:eq-never-large-exponent-2} this equals
$-\mathfrak{l}_k+\mathfrak{r}_k$ with $\mathfrak{r}_k:=\sum_{j=k}^{K-1}r_j$. At $k=K$ all three sums are
empty, so $\mathfrak{E}_K=\mathfrak{l}_K=\mathfrak{r}_K=0$.

Iterating the relation $|Y|_i=L^d|Y|_{i+1}$ of the cell counts gives
$|T|_{j+1}=L^{-d(j+1-k)}|T|_k$. Put $q:=L^d>1$ and $j=k+t$. By
\eqref{4.36:eq-never-large-exponent-3},
\begin{equation*}
|\mathfrak{r}_k|\le\bigl[(L^d-1)c^{\sharp}_G+K_*(E_0+1)\bigr]|T|_k\sum_{t\ge0}q^{-(t+1)}
=\frac{(L^d-1)c^{\sharp}_G+K_*(E_0+1)}{L^d-1}\,|T|_k,
\end{equation*}
because $\sum_{t\ge0}q^{-(t+1)}=1/(q-1)$. The right-hand side is $C_{\mathfrak{r}}|T|_k$. The bound is
independent of $K$: only the geometric series $\sum_tL^{-dt}$ was used.

\emph{Step 5: two series.} For $q>1$,
\begin{equation*}
\sum_{t\ge0}(q-1)q^{-(t+1)}=(q-1)\,\frac{1/q}{1-1/q}=1,
\end{equation*}
and
\begin{equation*}
\sum_{t\ge0}t\,(q-1)q^{-(t+1)}=(q-1)q^{-1}\sum_{t\ge0}t\,q^{-t}
=(q-1)q^{-1}\,\frac{1/q}{(1-1/q)^2}=\frac{1}{q-1}.
\end{equation*}

\emph{Step 6: the lower estimate in \eqref{4.36:eq-never-large-exponent-b}.} By (h4),
$g_j\le\gamma<1$, so $x_j=g_j^{-2}>1$ and $\log x_j>0$ for every $j$. Every term of $\mathfrak{l}_k$ is
therefore non-negative, and $\mathfrak{l}_k$ is at least its term $j=k$:
\begin{equation*}
\mathfrak{l}_k\ge\tfrac12(d-1)D(L^d-1)L^{-d}\log x_k\,|T|_k=\tfrac12(d-1)D(1-L^{-d})\log x_k\,|T|_k.
\end{equation*}
With $D=N^2-1$ this is the left inequality.

\emph{Step 7: the upper estimate in \eqref{4.36:eq-never-large-exponent-b}.} Let $t\ge0$. By (h4),
$x_{k+t}-x_k\le\sum_{i=k}^{k+t-1}|x_{i+1}-x_i|\le\bar{B}t$. If $x_{k+t}\ge x_k$, then
$\log y\le y-1$ at $y=x_{k+t}/x_k$ gives
\begin{equation*}
\log x_{k+t}-\log x_k\le\frac{x_{k+t}-x_k}{x_k}\le\bar{B}t\,g_k^2\le\bar{B}\gamma^2t,
\end{equation*}
using $1/x_k=g_k^2\le\gamma^2$. If $x_{k+t}<x_k$, the left side is negative and the same inequality
holds. By $|T|_{j+1}=L^{-d(j+1-k)}|T|_k$,
\begin{equation*}
\mathfrak{l}_k=\tfrac12(d-1)D\,|T|_k\sum_{t=0}^{K-1-k}(q-1)q^{-(t+1)}\log x_{k+t}.
\end{equation*}
The weights $(q-1)q^{-(t+1)}$ are positive. Insert $\log x_{k+t}\le\log x_k+\bar{B}\gamma^2t$, and
extend both resulting sums to all $t\ge0$; this is allowed because $\log x_k>0$ and $\bar{B}\gamma^2t\ge0$.
The two series of Step~5 then give
\begin{equation*}
\mathfrak{l}_k\le\tfrac12(d-1)D\,|T|_k\Bigl[\log x_k\cdot1+\frac{\bar{B}\gamma^2}{q-1}\Bigr].
\end{equation*}
With $q=L^d$ and $D=N^2-1$ this is the right inequality. At $d=4$, $\tfrac12(d-1)D=\tfrac32(N^2-1)$.
This proves (a). Hypothesis (h3) was not used.

\emph{Step 8: parts (b)--(d).} These are the conclusions of the fourth, fifth and sixth steps of the
proof of the $SU(2)$ theorem, that is, of the steps of that proof which bound the logarithmic and
regular parts of the label corrections, the difference $E_k[\mathcal{S}]-\mathfrak{L}_k[\mathcal{S}]$,
and the two labels of the two-label bound. They are used here as stated there, with
$\dim\mathfrak{g}=D$, with $\sigma_N$ in place of the $SU(2)$ normaliser, and with
$\nu_{\mathfrak{g}}$ and $c_G(N)$ in place of their $SU(2)$ values. The group
enters those steps only in the fourth, and only through the difference of normalisation numbers.
Put $n_{j-1}(Y):=|Y^{(j-1)*}|-(L^d-1)|Y|_j$, and let $(\Omega,W)$ be the pair of arguments at which
the fourth step evaluates the normalisation number. By Definition~\ref{4.36:def-step-normalisation},
the relation $|T|_j-|\Omega|_j=|\Omega^{c}|_j$, and $\log z_{j-1}=D\log g_{j-1}+\zeta_{j-1}$, this
difference is
\begin{equation*}
\begin{split}
\mathsf{N}^{(j-1)}(T,T)-\mathsf{N}^{(j-1)}(\Omega,W)
&=D\log g_{j-1}\bigl[n_{j-1}(T)-n_{j-1}(\Omega)\bigr]
+\log\sigma_N\bigl(|T^{(j-1)*}|-|\Omega^{(j-1)*}|\bigr)\\
&\quad-(L^d-1)\zeta_{j-1}|\Omega^{c}|_j
+\nu_{\mathfrak{g}}\bigl(|T^{(j-1)*}|-|W^{(j-1)*}|\bigr).
\end{split}
\end{equation*}
It is estimated by the bond-count bounds of Lemma~\ref{4.36:lem-label-energy} and by
$|\zeta_{j-1}|\le c_G(N)$, which holds by Lemma~\ref{4.36:lem-one-link-partition}(ii) since
$g_{j-1}\le\gamma\le\gamma_z$. With these substitutions the three steps yield the bounds (b) and (c)
with the constants $C_{\log}(N)$, $C^{\mathrm{lab}}_{\mathrm{loc}}(N)$ and $C_{\mathrm{lab}}$ defined
above, and the two-label bound (d). These steps use hypothesis (h3), that is, the locality remark for
$\mathbf{R}'$.
\end{proof}

\begin{corollary}[The density constants $E_{-}$ and $E_{+}$ at $SU(N)$]\label{4.36:cor-density-constants}
Let $d=4$. Assume the following statements:
\begin{itemize}
\item[(A)] the theorem on the exponent of the never-large term and the label corrections at $SU(N)$,
Theorem~\ref{4.36:thm-never-large-exponent}, together with its hypotheses
\eqref{4.36:eq-never-large-exponent-a};
\item[(B)] the lower-bound lemma for the effective density at $SU(N)$, with its density constant
$E^{\sharp}_1$ and its exponent $\beta_T$;
\item[(C)] the uniformity lemma for the energy constants;
\item[(D)] the integration lemma, the density-format lemma and the step lemma of this chapter.
\end{itemize}
Let $C_{\mathfrak{r}}$, $C_{\log}(N)$ and $C^{\mathrm{lab}}_{\mathrm{loc}}(N)$ be the constants of
Theorem~\ref{4.36:thm-never-large-exponent}, let $\mathfrak{E}_k=E_k[\emptyset]$ be the energy of the
never-large label and $\mathfrak{l}_k$ the exponent of the never-large term defined there, let the region
$T$, the labels $\mathcal{S}$ at level $k$ with their energies $E_k[\mathcal{S}]$, the large-field regions
$Z_j$ of a label with their enlargements $Z^{+}_j$, and the cell counts $|\cdot|_j$ be as in that theorem,
and let $0<g_j\le\gamma<1$ and $\bar{B}$ be as in \eqref{4.36:eq-never-large-exponent-a}. Then the
following hold.
\begin{enumerate}
\item[(1)] Put
\begin{equation}\label{4.36:eq-density-constants-1}
E_{-}=E_{-}(N,L,\mathfrak{p}):=E^{\sharp}_1\bigl(1-L^{-\beta_T}\bigr)^{-1}+1+C_{\mathfrak{r}} .
\end{equation}
The constant $E_{-}$ does not depend on $g_k$, $k$, $K$ or $m$, and
\begin{equation}\label{4.36:eq-density-constants-2}
\chi_k\exp\bigl[-g_k^{-2}A(U_k)-E_{-}N_k\bigr]\,dV_k\le\rho_k\,dV_k .
\end{equation}
\item[(2)] Put
\begin{equation}\label{4.36:eq-density-constants-3}
C_{\Gamma}(g_k):=\tfrac32(N^2-1)\bigl(\log x_k+\bar{B}\gamma^2\bigr)+C_{\mathfrak{r}} .
\end{equation}
For every label $\mathcal{S}$ at level $k$,
\begin{equation}\label{4.36:eq-density-constants-4}
-E_k[\mathcal{S}]\le C_{\Gamma}(g_k)|T|_k
+\sum_{j=1}^{k}\bigl(C_{\log}(N)\log x_{j-1}+C^{\mathrm{lab}}_{\mathrm{loc}}(N)\bigr)|Z^{+}_j|_j .
\end{equation}
Hence the constant $E_{+}$ of \cite[Corollary~3]{B-CMP119} may be taken as
$E_{+}=E_{+}(N,L,\mathfrak{p};g_k)=O(1+\log g_k^{-2})$, with leading member
\begin{equation*}
\tfrac12(d-1)\cdot\dim SU(N)\cdot\log g_k^{-2}=\tfrac32(N^2-1)\log g_k^{-2};
\end{equation*}
by clause (i-c) of Theorem~0 and statement (C), $E_{+}$ may be taken at $g_K$. The share of level $j$
in \eqref{4.36:eq-density-constants-4} is at most
\begin{equation}\label{4.36:eq-density-constants-5}
3^d(MR_j)^d\bigl(C_{\log}(N)(\log x_j+\bar{B}\gamma^2)+C^{\mathrm{lab}}_{\mathrm{loc}}(N)\bigr)
\end{equation}
per cube of side $MR_j$ meeting $Z_j$, where $M$ is the auxiliary parameter of
Notation~\ref{0.2:not-glossary} and $R_j$ is as in Notation~\ref{4.36:not-flow-thresholds}.
\item[(3)] For $0\le s\le K$,
\begin{equation}\label{4.36:eq-density-constants-6}
|\mathfrak{E}_{K-s}|\le C_{\Gamma}(g_{K-s})|\mathbb{T}_s| ,
\end{equation}
and at fixed $(g,s)$ the right-hand side is bounded independently of $K$.
\end{enumerate}
\end{corollary}

\begin{proof}
Throughout, $x_j=g_j^{-2}$. Since $0<g_j\le\gamma<1$ by \eqref{4.36:eq-never-large-exponent-a},
$x_j>1$ and therefore $\log x_j>0$ for every $j$.

\emph{Step 1: an upper bound on the never-large energy.}
By clause (a) of Theorem~\ref{4.36:thm-never-large-exponent}, $\mathfrak{E}_k=-\mathfrak{l}_k+\mathfrak{r}_k$
with $|\mathfrak{r}_k|\le C_{\mathfrak{r}}|T|_k$. For $0\le k<K$ the lower half of the two-sided estimate
\eqref{4.36:eq-never-large-exponent-b} gives
\begin{equation*}
\mathfrak{l}_k\ge\tfrac12(d-1)(N^2-1)(1-L^{-d})\log x_k|T|_k\ge0 ,
\end{equation*}
because $1-L^{-d}>0$ and $\log x_k>0$; for $k=K$, $\mathfrak{l}_K=0$. Discarding $-\mathfrak{l}_k\le0$
therefore yields
\begin{equation}\label{4.36:eq-density-constants-7}
\mathfrak{E}_k\le\mathfrak{r}_k\le C_{\mathfrak{r}}|T|_k\qquad(0\le k\le K).
\end{equation}

\emph{Step 2: proof of (1).}
The bound \eqref{4.36:eq-density-constants-7} is the only information on $\mathfrak{E}_k$ that we
feed into the lower-bound lemma (B), with its constant $E^{\sharp}_1$. That lemma, applied together with
the integration lemma, the density-format lemma and the step lemma of (D), then gives
\eqref{4.36:eq-density-constants-2} with $E_{-}$ as in \eqref{4.36:eq-density-constants-1}: the
contribution of the never-large energy to the constant is the $C_{\mathfrak{r}}$ of
\eqref{4.36:eq-density-constants-7}. The three members $E^{\sharp}_1(1-L^{-\beta_T})^{-1}$, $1$ and
$C_{\mathfrak{r}}$ depend on $N$, $L$ and $\mathfrak{p}$ only. In particular $C_{\mathfrak{r}}$ is
independent of $K$, since clause (a) of Theorem~\ref{4.36:thm-never-large-exponent} obtains it from the
geometric sum $\sum_{t\ge0}L^{-dt}$ alone. Hence $E_{-}$ is free of $g_k$, $k$, $K$ and $m$.

\emph{Step 3: proof of \eqref{4.36:eq-density-constants-4}.}
The upper half of \eqref{4.36:eq-never-large-exponent-b} at $d=4$ reads
\begin{equation*}
\mathfrak{l}_k\le\tfrac32(N^2-1)\Bigl(\log x_k+\frac{\bar{B}\gamma^2}{L^d-1}\Bigr)|T|_k .
\end{equation*}
Since $L^d-1\ge1$ and $\bar{B}\gamma^2\ge0$, we may replace $\bar{B}\gamma^2/(L^d-1)$ by
$\bar{B}\gamma^2$. Adding the bound $C_{\mathfrak{r}}|T|_k$ on $|\mathfrak{r}_k|$ gives
\begin{equation*}
\mathfrak{l}_k+|\mathfrak{r}_k|\le C_{\Gamma}(g_k)|T|_k ,
\end{equation*}
with $C_{\Gamma}(g_k)$ as in \eqref{4.36:eq-density-constants-3}.

It remains to treat a general label $\mathcal{S}$. Let $\mathfrak{e}^{\log}_j$ and $\mathfrak{e}^{\mathrm{reg}}_j$
be the logarithmic and the regular corrections of level $j$, $\Omega^{c}_j$, $\Gamma_n$ and $C_{\mathrm{lab}}$ the sets
and the constant, and $\mathfrak{L}_k[\mathcal{S}]=-\mathfrak{l}_k+\sum_{j}\mathfrak{e}^{\log}_j$ the
leading part, all as in Theorem~\ref{4.36:thm-never-large-exponent}. Writing
$E_k[\mathcal{S}]=\mathfrak{L}_k[\mathcal{S}]+(E_k[\mathcal{S}]-\mathfrak{L}_k[\mathcal{S}])$ gives
\begin{equation}\label{4.36:eq-density-constants-8}
-E_k[\mathcal{S}]\le\mathfrak{l}_k+\sum_{j}|\mathfrak{e}^{\log}_j|
+\bigl|E_k[\mathcal{S}]-\mathfrak{L}_k[\mathcal{S}]\bigr| .
\end{equation}
Clauses (b)--(d) of Theorem~\ref{4.36:thm-never-large-exponent} provide, for the corrections of level $j$
and for the remainder, the bounds
\begin{align*}
|\mathfrak{e}^{\log}_j|&\le C_{\log}(N)\log x_{j-1}|\Omega^{c}_j|_j,\\
|\mathfrak{e}^{\mathrm{reg}}_j|&\le C^{\mathrm{lab}}_{\mathrm{loc}}(N)|Z^{+}_j|_j,\\
\bigl|E_k[\mathcal{S}]-\mathfrak{L}_k[\mathcal{S}]\bigr|&\le C_{\mathrm{lab}}\sum_{n=0}^{k}|\Gamma_n| .
\end{align*}
Theorem~\ref{4.36:thm-never-large-exponent}, through these clauses (b)--(d) and with the never-large
members estimated by the display preceding \eqref{4.36:eq-density-constants-8}, bounds the right-hand side
of \eqref{4.36:eq-density-constants-8} by the right-hand side of \eqref{4.36:eq-density-constants-4}: the
never-large part carries the coefficient $C_{\Gamma}(g_k)$ of $|T|_k$, and the corrections of level $j$
carry the coefficient $C_{\log}(N)\log x_{j-1}+C^{\mathrm{lab}}_{\mathrm{loc}}(N)$ of $|Z^{+}_j|_j$. This is
\eqref{4.36:eq-density-constants-4} for every label $\mathcal{S}$.

\emph{Step 4: the size of $E_{+}$.}
Since $\log x_k=\log g_k^{-2}$,
\begin{equation*}
C_{\Gamma}(g_k)=\tfrac32(N^2-1)\log g_k^{-2}+\Bigl(\tfrac32(N^2-1)\bar{B}\gamma^2+C_{\mathfrak{r}}\Bigr).
\end{equation*}
The bracket depends on $N$, $L$ and $\mathfrak{p}$ only, and is at most
$\tfrac32(N^2-1)\bar{B}+C_{\mathfrak{r}}$ because $\gamma<1$. Hence $E_{+}=O(1+\log g_k^{-2})$, and its
leading member is $\tfrac12(d-1)\dim SU(N)\log g_k^{-2}$. Since $\dim SU(N)=N^2-1$ and $d=4$, this
equals $\tfrac32(N^2-1)\log g_k^{-2}$. By clause (i-c) of Theorem~0 and the uniformity lemma (C), the
coupling at which $E_{+}$ is evaluated may be taken to be $g_K$.

\emph{Step 5: the shares of the labels, proof of \eqref{4.36:eq-density-constants-5}.}
First,
\begin{equation*}
\log x_{j-1}-\log x_j\le\frac{x_{j-1}-x_j}{x_j}\le\bar{B}g_j^2\le\bar{B}\gamma^2 .
\end{equation*}
The first inequality is $\log y\le y-1$ at $y=x_{j-1}/x_j$. The second uses $|x_j-x_{j-1}|\le\bar{B}$
and $1/x_j=g_j^2$, and the third uses $g_j\le\gamma$; both bounds are contained in
\eqref{4.36:eq-never-large-exponent-a}. Thus $\log x_{j-1}\le\log x_j+\bar{B}\gamma^2$.

Second, under the geometric hypothesis of Theorem~\ref{4.36:thm-never-large-exponent},
$|Z^{+}_j|_j\le3^d|Z_j|_j$. Partition the lattice of level $j$ into cubes of side $MR_j$, measured in
units of level $j$; each of them contains $(MR_j)^d$ cells of level $j$. Hence $|Z_j|_j$ is at most
$(MR_j)^d$ times the number of cubes of the partition meeting $Z_j$.

Combining the two facts, the level-$j$ term in \eqref{4.36:eq-density-constants-4} is at most
\eqref{4.36:eq-density-constants-5} per cube of side $MR_j$ meeting $Z_j$. \cite{B-CMP122b} states that
every large field domain $Z_j$ contributes the constant $O(1)\log g_j^{-2}|Z_j|$. The shares
\eqref{4.36:eq-density-constants-5} are of this form. Their dependence on $N$, through $C_{\log}(N)$ and
$C^{\mathrm{lab}}_{\mathrm{loc}}(N)$, is compensated by the ceiling $\gamma$ on the couplings, which is chosen
after $N$.

\emph{Step 6: proof of (3).}
Let $0\le s\le K$ and $k=K-s$. By clause (a) of Theorem~\ref{4.36:thm-never-large-exponent} and
Step~1,
\begin{equation*}
|\mathfrak{E}_k|\le\mathfrak{l}_k+|\mathfrak{r}_k|\le C_{\Gamma}(g_k)|T|_k ,
\end{equation*}
where the last inequality is Step~3. For $k=K-s$ the cell count $|T|_{K-s}$ is written
$|\mathbb{T}_s|$, and this is \eqref{4.36:eq-density-constants-6}.

For the independence from $K$ we use two inputs. The first is the flow inequality $(\mathrm{Fl}1)$ of
\eqref{4.36:eq-flow-thresholds-a}. The second is the anchoring $g_K\in[g_{-}(g),g)$ of the flow at the
last step (Notation~\ref{4.36:not-flow-thresholds}; the first passage $K=K(g_0,g)$ is that of
Definition~\ref{1.2:def-flow-constants}). Together they give
\begin{equation*}
x_{K-s}\le g^{-2}+B^{\sharp}(s+1).
\end{equation*}
Since $C_{\Gamma}$ is increasing in $\log x_{K-s}$, we obtain
\begin{equation*}
C_{\Gamma}(g_{K-s})\le\tfrac32(N^2-1)\Bigl(\log\bigl(g^{-2}+B^{\sharp}(s+1)\bigr)+\bar{B}\gamma^2\Bigr)
+C_{\mathfrak{r}} .
\end{equation*}
At fixed $(g,s)$ the right-hand side does not depend on $K$.
\end{proof}

\begin{remark}[Necessity of the coupling dependence of $E_{+}$]\label{4.36:rem-coupling-dependence}
Let $k<K$. Let $\rho_k$ be the effective density at level $k$, a density with respect to the product
Haar measure $dV_k$ on the level-$k$ variables $V_k$. Let $|T|_k$ be the number of level-$k$ cells of
the lattice $T$ on which $\rho_k$ is defined, as in
Theorem~\ref{4.36:thm-never-large-exponent}, and write $x_k=g_k^{-2}$. There is a neighbourhood
$\mathcal{U}$ of the unit configuration $V_k=\mathbf{1}$ on which
\begin{equation}\label{4.36:eq-coupling-dependence-1}
\rho_k\,dV_k\ \ge\ \exp\Bigl\{\Bigl[\tfrac32(N^2-1)\bigl(1-L^{-4}\bigr)\log x_k
-C_{\mathfrak{r}}\Bigr]|T|_k-1\Bigr\}\,dV_k .
\end{equation}
Consequently every constant $E_{+}$ for which a density bound $\rho_k\le e^{E_{+}|T|_k}$ holds satisfies
\begin{equation}\label{4.36:eq-coupling-dependence-2}
E_{+}\ \ge\ \tfrac32(N^2-1)\bigl(1-L^{-4}\bigr)\log g_k^{-2}-C_{\mathfrak{r}}-1 .
\end{equation}
Since $d=4$, the coefficient of $\log g_k^{-2}$ on the right equals
$\tfrac12(d-1)\dim SU(N)\,(1-L^{-d})$. The dependence of $E_{+}$ on $g_k$ is therefore necessary, and
its coefficient is proportional to $\dim SU(N)$. In \cite[\S1]{B-CMP122b} the constants $E_{-}$,
$E_{+}$ of the bounds on the effective densities are stated to be independent of $k$, of the lattice
and of $U_k$.
In this chapter the densities $\rho_k$ are used in the normalisation of
Corollary~\ref{4.36:cor-density-constants}, in which the upper constant depends on $g_k$; for these
densities \eqref{4.36:eq-coupling-dependence-2} shows that an upper constant uniform in $g_k$ is not
available.

The size of \eqref{4.36:eq-coupling-dependence-1} is consistent with the equality
$\int\rho_k\,dV_k=\int\rho_K\,dV_K$ of the integrals of the effective densities of
Definition~\ref{1.5:def-densities}. To leading order in $g_k$, the integral
$\int dV_k\,e^{-g_k^{-2}A(U_k)}$ is of size $g_k^{(d-1)(N^2-1)|T|_k}$. This factor cancels
$\mathfrak{l}_k$ to within the correction of order $\bar{B}\gamma^2$ in
\eqref{4.36:eq-never-large-exponent-b}. Thus the integral of the term of $\rho_k$ without large-field
regions is consistent with that equality.
\end{remark}

\begin{proof}
The argument has five steps.

\emph{The never-large term at the unit configuration.} Let $V_k=\mathbf{1}$. Then the block
configuration is $U_k=\mathbf{1}$. The configuration satisfies the small-field conditions, so
$\chi_k=1$. The action vanishes, $A(\mathbf{1})=0$.

In the format of $\rho_k$ fixed in Theorem~\ref{4.36:thm-never-large-exponent}, the contributions
$\mathbf{E}_k$ and $\mathbf{R}_k$ enter as differences that
vanish at the unit configuration, so $\mathbf{E}_k(\mathbf{1})=\mathbf{R}_k(\mathbf{1})=0$. The
contribution $\mathbf{B}_k$ is carried by the large-field region. That region is empty for the term of
$\rho_k$ in which no large-field region occurs, so $\mathbf{B}_k=0$ for that term.

Hence this term, the never-large term of Theorem~\ref{4.36:thm-never-large-exponent}, equals
$e^{-\mathfrak{E}_k}$ at $V_k=\mathbf{1}$.

\emph{A neighbourhood of the unit configuration.} The unit configuration satisfies the small-field
conditions strictly, with a margin. These conditions therefore continue to hold, and $\chi_k=1$ remains
true, for all real configurations in a neighbourhood of $\mathbf{1}$. On that neighbourhood the
never-large term depends continuously on $V_k$.

At $V_k=\mathbf{1}$ the term takes the value $e^{-\mathfrak{E}_k}$, which is strictly larger than
$e^{-\mathfrak{E}_k-1}$. By continuity there is a neighbourhood $\mathcal{U}$ of $\mathbf{1}$ on which
the never-large term is at least $e^{-\mathfrak{E}_k-1}$.

\emph{The remaining terms.} Every other term of $\rho_k$, that is, every term carrying a non-empty
large-field region, is non-negative on real configurations. This is the positivity of these terms on
real configurations on which the lower density bound of Corollary~\ref{4.36:cor-density-constants}
rests, in the form stated in \cite{B-CMP122b}. Since $V_k$ ranges over
$SU(N)$-valued, hence real, configurations, $\rho_k$ is bounded below by its never-large term. On
$\mathcal{U}$ this gives
\[
\rho_k\ \ge\ e^{-\mathfrak{E}_k-1}.
\]

\emph{The lower bound on the exponent.} In the notation of
Theorem~\ref{4.36:thm-never-large-exponent}, $\mathfrak{E}_k$ is the sum of $-\mathfrak{l}_k$ and a
part bounded above by $C_{\mathfrak{r}}|T|_k$; discarding $-\mathfrak{l}_k\le0$ from this gives the
bound $\mathfrak{E}_k\le C_{\mathfrak{r}}|T|_k$ used in Corollary~\ref{4.36:cor-density-constants}.
For $k<K$ the
lower side of the two-sided estimate \eqref{4.36:eq-never-large-exponent-b} of $\mathfrak{l}_k$ in
Theorem~\ref{4.36:thm-never-large-exponent} bounds $\mathfrak{l}_k$ below by
$\tfrac32(N^2-1)(1-L^{-4})\log x_k\,|T|_k$. Together the two bounds give
\[
-\mathfrak{E}_k\ \ge\ \Bigl[\tfrac32(N^2-1)\bigl(1-L^{-4}\bigr)\log x_k-C_{\mathfrak{r}}\Bigr]|T|_k .
\]
Inserting this into the bound of the previous step yields \eqref{4.36:eq-coupling-dependence-1} on
$\mathcal{U}$.

\emph{The consequence for $E_{+}$.} Suppose that $\rho_k\le e^{E_{+}|T|_k}$. The set $\mathcal{U}$ is
a non-empty open set, so it has positive Haar measure. Hence there is a point of $\mathcal{U}$ at which
both this bound and \eqref{4.36:eq-coupling-dependence-1} hold. At that point
\[
E_{+}|T|_k\ \ge\ \Bigl[\tfrac32(N^2-1)\bigl(1-L^{-4}\bigr)\log x_k-C_{\mathfrak{r}}\Bigr]|T|_k-1 .
\]
Divide by $|T|_k$ and use $|T|_k\ge1$, which gives $-1/|T|_k\ge-1$, together with $x_k=g_k^{-2}$. The
result is \eqref{4.36:eq-coupling-dependence-2}.
\end{proof}

\begin{lemma}[Haar measure of small balls and small-action sets]\label{4.36:lem-haar-small-sets}
Let $N\ge 2$, $G=SU(N)$, $D=N^2-1=\dim\mathfrak{g}$ and $\mathsf{n}=(N-1)^{1/2}$. The norms
$|\cdot|$ and $|\cdot|_{\mathfrak{g}}$ on $\mathfrak{g}$ are those of
Notation~\ref{4.36:not-group-norms}. The set $\mathfrak{D}$, the density $J$, the normaliser
$\sigma_N$, the volumes $\omega_D$ and $\omega^{\mathrm{op}}_N$ and the balls
$B^{\mathfrak{g}}_r(u)$, $B^{\mathrm{op}}_r(u)$ are those of
Definition~\ref{4.36:def-haar-balls}. Put
\begin{equation}\label{4.36:eq-haar-small-sets-1}
\begin{gathered}
\kappa_N:=\frac{N^2\sin^2(\pi/N)}{2\pi^2(N-1)^2},\qquad
c_3(N):=\sigma_N\,\omega_D\Big(\frac{2}{\pi}\Big)^{N^2-N},\qquad
c^{\flat}_3(N):=(N-1)^{-D/2}c_3(N),\\
c^{\mathrm{op}}_3(N):=\sigma_N\,\omega^{\mathrm{op}}_N\Big(\frac{2}{\pi}\Big)^{N^2-N},\qquad
C_5(N):=\sigma_N\,\omega_D\,\kappa_N^{-D/2}.
\end{gathered}
\end{equation}
Then the following hold.
\begin{itemize}
\item[(a)] For every $u\in G$ and $0<r\le\pi/\sqrt{2N}$,
$\operatorname{Haar}(B^{\mathfrak{g}}_r(u))\ge c_3(N)\,r^D$.
\item[(a$^{\mathrm{op}}$)] For every $u\in G$ and $0<r\le\pi/2$, a radius that does not depend
on $N$,
\begin{equation*}
\operatorname{Haar}(B^{\mathrm{op}}_r(u))\ge c^{\mathrm{op}}_3(N)\,r^D\ge c^{\flat}_3(N)\,r^D,
\end{equation*}
and also, directly from \textup{(a)},
\begin{equation*}
\operatorname{Haar}(B^{\mathrm{op}}_r(u))\ge\operatorname{Haar}(B^{\mathfrak{g}}_{r/\mathsf{n}}(u))
\ge c^{\flat}_3(N)\,r^D .
\end{equation*}
\item[(b)] For every $\tau>0$,
\begin{equation*}
\operatorname{Haar}\{u\in G:\ 1-\operatorname{Re}\operatorname{tr}u<\tau\}\le C_5(N)\,\tau^{D/2}.
\end{equation*}
\item[(c)] For every $u'\in G$ the maps $u\mapsto uu'$ and $u\mapsto u'u$ preserve
$\operatorname{Haar}$.
\item[(d)] For every $r>0$, as subsets of $\mathfrak{g}$,
\begin{equation*}
\{v:\ |v|_{\mathfrak{g}}<r/\mathsf{n}\}\subset\{v:\ |v|<r\}
\subset\{v:\ |e^{iv}-\mathbf{1}|<r\}.
\end{equation*}
\end{itemize}
Moreover
\begin{equation}\label{4.36:eq-haar-small-sets-2}
\sigma_N\ge\big(\omega_D(2\pi)^D\big)^{-1},\qquad
\frac{2}{\pi^2(N-1)^2}\le\kappa_N\le\frac{1}{2(N-1)^2}\le\frac12 .
\end{equation}
\end{lemma}

\begin{proof}
For $v\in\mathfrak{g}$ we write $v_1,\dots,v_N$ for its eigenvalues, so that
$\sum_p v_p=0$, $|v|=\max_p|v_p|$ and $|v|_{\mathfrak{g}}^2=N^{-1}\sum_p v_p^2$.

\emph{Reduction to $u=\mathbf{1}$.} For a Borel set $S\subset\mathfrak{D}$, the identity
\eqref{4.36:eq-haar-balls-a} of Definition~\ref{4.36:def-haar-balls}, applied with $h=1$, gives
\begin{equation*}
\operatorname{Haar}(u\exp iS)=\sigma_N\int_S J(v)\,dv ,
\end{equation*}
and the right side does not depend on $u$. Since $B^{\mathfrak{g}}_r(u)=u\exp iS$ with
$S=\{|v|_{\mathfrak{g}}<r\}$ and $B^{\mathrm{op}}_r(u)=u\exp iS$ with $S=\{|v|<r\}$, it suffices in
(a) and (a$^{\mathrm{op}}$), once these sets $S$ are shown to lie in $\mathfrak{D}$, to evaluate
$\sigma_N\int_S J\,dv$.

\emph{(a).} Let $|v|_{\mathfrak{g}}<r\le\pi/\sqrt{2N}$. By Lemma~\ref{4.36:lem-eigenangle-inequalities}
(in \eqref{4.36:eq-eigenangle-inequalities-a}), $(v_p-v_q)^2\le 2(v_p^2+v_q^2)\le
2N|v|_{\mathfrak{g}}^2$, hence
\begin{equation*}
|v_p-v_q|\le\sqrt{2N}\,|v|_{\mathfrak{g}}<\pi\qquad\text{for all }p,q .
\end{equation*}
The spread of $v$ is therefore less than $\pi<2\pi$, so $v\in\mathfrak{D}$, on which $\exp i$ is
injective by Definition~\ref{4.36:def-haar-balls}. Each of the $N(N-1)/2$ factors of
$J(v)=\prod_{p<q}\operatorname{sinc}^2((v_p-v_q)/2)$ has argument of modulus less than $\pi/2$;
since $\operatorname{sinc}$ is even, positive and decreasing on $[0,\pi/2]$
(Lemma~\ref{4.36:lem-eigenangle-inequalities}), each factor is at least
$\operatorname{sinc}^2(\pi/2)=(2/\pi)^2$, so $J(v)\ge(2/\pi)^{N^2-N}$. The $dv$-volume of
$\{|v|_{\mathfrak{g}}<r\}$ is $\omega_D r^D$, because $\mathfrak{g}$ with the inner product
$\operatorname{tr}XY$ is a Euclidean space of dimension $D$. Hence
\begin{equation*}
\operatorname{Haar}(B^{\mathfrak{g}}_r(u))=\sigma_N\int_{|v|_{\mathfrak{g}}<r}J(v)\,dv
\ge\sigma_N\Big(\frac{2}{\pi}\Big)^{N^2-N}\omega_D r^D=c_3(N)\,r^D .
\end{equation*}

\emph{(a$^{\mathrm{op}}$).} Let $|v|<r\le\pi/2$. Then
\begin{equation*}
|v_p-v_q|\le|v_p|+|v_q|\le 2|v|<\pi ,
\end{equation*}
and the
argument of (a) applies word for word, with the $dv$-volume of $\{|v|<r\}$ equal to
$\omega^{\mathrm{op}}_N r^D$ by homogeneity of the operator norm. This gives the first inequality.
The second, $c^{\mathrm{op}}_3(N)\ge c^{\flat}_3(N)$, is the inequality
$\omega^{\mathrm{op}}_N\ge(N-1)^{-D/2}\omega_D$ of Definition~\ref{4.36:def-haar-balls}. For the
second route, (d) at the same $r$ gives $B^{\mathfrak{g}}_{r/\mathsf{n}}(u)\subset B^{\mathrm{op}}_r(u)$.
Moreover $r/\mathsf{n}\le\pi/\sqrt{2N}$ is equivalent to $r\le\pi\sqrt{(N-1)/(2N)}$, and
$(N-1)/(2N)\ge 1/4$ is equivalent to $N\ge2$; hence $r\le\pi/2$ implies
$r/\mathsf{n}\le\pi/\sqrt{2N}$, and (a) at radius $r/\mathsf{n}$ gives
\begin{equation*}
\operatorname{Haar}(B^{\mathfrak{g}}_{r/\mathsf{n}}(u))\ge c_3(N)\,(r/\mathsf{n})^D
=c_3(N)(N-1)^{-D/2}r^D=c^{\flat}_3(N)\,r^D .
\end{equation*}

\emph{(b).} Let $v\in\mathfrak{D}$ and let $v_{\max}$, $v_{\min}$ be its largest and smallest
eigenvalues. Since the spread is less than $2\pi$, every $v_q>v_{\max}-2\pi$. Summing over the
$N-1$ indices other than one index at which the maximum is attained,
\begin{equation*}
0=\sum_p v_p>v_{\max}+(N-1)(v_{\max}-2\pi)=Nv_{\max}-2\pi(N-1),
\end{equation*}
that is, $v_{\max}<t_N:=2\pi(N-1)/N$. The matrix $-v$ is traceless with the same spread, so
$-v\in\mathfrak{D}$, and the same argument applied to $-v$ gives $v_{\min}>-t_N$. Thus
\begin{equation*}
|v_p|<t_N<2\pi\qquad\text{for every }p\text{ and every }v\in\mathfrak{D};
\end{equation*}
in particular $|v|<2\pi$ on $\mathfrak{D}$. By the sinc inequalities of
Lemma~\ref{4.36:lem-eigenangle-inequalities}, used in the form
$1-\cos y\ge((1-\cos y_0)/y_0^2)\,y^2$ for $|y|\le y_0$ and $0<y_0<2\pi$ (which follows from
$1-\cos y=\tfrac12y^2\operatorname{sinc}^2(y/2)$ and the monotonicity of $\operatorname{sinc}$
on $[0,\pi)$), taken at $y_0=t_N$, we get $1-\cos v_p\ge\kappa_N v_p^2$, because
\begin{equation*}
\frac{1-\cos t_N}{t_N^2}
=\frac{2\sin^2(\pi-\pi/N)\,N^2}{4\pi^2(N-1)^2}=\kappa_N .
\end{equation*}
Since $\operatorname{Re}\operatorname{tr}e^{iv}=N^{-1}\sum_p\cos v_p$, averaging over $p$ gives
\begin{equation}\label{4.36:eq-haar-small-sets-3}
1-\operatorname{Re}\operatorname{tr}e^{iv}\ge\kappa_N|v|_{\mathfrak{g}}^2\qquad(v\in\mathfrak{D}).
\end{equation}
By Definition~\ref{4.36:def-haar-balls}, $\mathrm{Cut}$ is Haar-null and $\exp i$ maps $\mathfrak{D}$
bijectively onto $G\setminus\mathrm{Cut}$. Every $u\notin\mathrm{Cut}$ with
$1-\operatorname{Re}\operatorname{tr}u<\tau$ is therefore $e^{iv}$ with $v\in\mathfrak{D}$, and by
\eqref{4.36:eq-haar-small-sets-3} $|v|_{\mathfrak{g}}<(\tau/\kappa_N)^{1/2}$. Hence, up to a null set,
\begin{equation*}
\{u:\ 1-\operatorname{Re}\operatorname{tr}u<\tau\}\subset
\exp i\,\{v\in\mathfrak{D}:\ |v|_{\mathfrak{g}}<(\tau/\kappa_N)^{1/2}\},
\end{equation*}
and by \eqref{4.36:eq-haar-balls-a} with $u=\mathbf{1}$, $h=1$, and $J\le1$, the right side has
Haar measure at most $\sigma_N\omega_D(\tau/\kappa_N)^{D/2}=C_5(N)\tau^{D/2}$. No restriction on
$\tau$ was used; for $\tau>2$ the set on the left is all of $G$, since
$\operatorname{Re}\operatorname{tr}u\ge-1$, and the bound then reads $1\le C_5(N)\tau^{D/2}$,
which carries no further information.

\emph{(c).} $G$ is compact, hence unimodular, and $\operatorname{Haar}$ is both left- and
right-invariant (Definition~\ref{4.36:def-haar-balls}).

\emph{(d).} If $|v|_{\mathfrak{g}}<r/\mathsf{n}$, then $|v|\le\mathsf{n}|v|_{\mathfrak{g}}<r$ by the
inequality $|X|\le\mathsf{n}|X|_{\mathfrak{g}}$ of Lemma~\ref{4.36:lem-eigenangle-inequalities}. If
$|v|<r$, then $e^{iv}-\mathbf{1}$ is normal with eigenvalues $e^{iv_p}-1$, of modulus
$2|\sin(v_p/2)|\le|v_p|$, so $|e^{iv}-\mathbf{1}|\le|v|<r$
(Lemma~\ref{4.36:lem-eigenangle-inequalities}).

\emph{The bound on $\sigma_N$.} By the definition of $\sigma_N$ and $J\le1$,
\begin{equation*}
1=\sigma_N\int_{\mathfrak{D}}J(v)\,dv\le\sigma_N\operatorname{vol}(\mathfrak{D}).
\end{equation*}
By the proof of (b), $\mathfrak{D}\subset\{|v|<2\pi\}$, and since $|v|_{\mathfrak{g}}\le|v|$
(Lemma~\ref{4.36:lem-eigenangle-inequalities}), $\{|v|<2\pi\}\subset\{|v|_{\mathfrak{g}}<2\pi\}$,
whose volume is $\omega_D(2\pi)^D$. This gives the first bound in
\eqref{4.36:eq-haar-small-sets-2}.

\emph{The bounds on $\kappa_N$.} Since $N\ge2$, $y=\pi/N\in(0,\pi/2]$, and the sine bounds
$2y/\pi\le\sin y\le y$ on $(0,\pi/2]$ of Lemma~\ref{4.36:lem-eigenangle-inequalities} give
$2/N\le\sin(\pi/N)\le\pi/N$. Inserting these in the definition of $\kappa_N$,
\begin{equation*}
\frac{N^2\cdot 4N^{-2}}{2\pi^2(N-1)^2}\le\kappa_N\le\frac{N^2\cdot\pi^2N^{-2}}{2\pi^2(N-1)^2},
\end{equation*}
which is $2/(\pi^2(N-1)^2)\le\kappa_N\le1/(2(N-1)^2)$; the last bound is at most $1/2$ because
$N-1\ge1$.

For $N=2$ one has $D=3$, $\mathsf{n}=1$, $\sigma_2=1/(2\pi^2)$ and
$\omega^{\mathrm{op}}_2=\omega_3=4\pi/3$ (Definition~\ref{4.36:def-haar-balls}); the definitions
\eqref{4.36:eq-haar-small-sets-1} then give $\kappa_2=2/\pi^2$,
\begin{equation*}
c^{\mathrm{op}}_3(2)=c^{\flat}_3(2)=c_3(2)=\frac{1}{2\pi^2}\cdot\frac{4\pi}{3}\cdot\frac{4}{\pi^2}
=\frac{8}{3\pi^3},\qquad
C_5(2)=\frac{1}{2\pi^2}\cdot\frac{4\pi}{3}\Big(\frac{\pi^2}{2}\Big)^{3/2}=\frac{\pi^2}{3\sqrt2}. \qedhere
\end{equation*}
\end{proof}

\begin{lemma}[The mode bound for a log-density along one link]\label{4.36:lem-mode-bound}
Let $G=SU(N)$, let $W\subset G$ be open, and let $(\mathcal{W},\mu)$ be a $\sigma$-finite measure space.
Let $F=F_1+F_2\colon W\times\mathcal{W}\to\mathbb{R}$ be jointly measurable, and put
\begin{equation*}
Z(u):=\int_{\mathcal{W}}e^{F(u,y)}\,\mu(dy)\in[0,\infty],\qquad \Phi:=\log Z ;
\end{equation*}
by Tonelli's theorem $Z$ is Borel on $W$. Assume
\begin{equation*}
0<\int_W e^{\Phi}\,d\operatorname{Haar}<\infty ,
\end{equation*}
and let $f:=\mathbf{1}_W e^{\Phi}\big/\int_W e^{\Phi}\,d\operatorname{Haar}$, a probability density with
respect to Haar measure. Let $u^0\in W$ with $Z(u^0)<\infty$. Fix one of the following two cases:
\begin{itemize}
\item[(op)] $|\cdot|_{\bullet}=|\cdot|$, the operator norm, and $0<r\le\pi/2$;
\item[($\mathfrak{g}$)] $|\cdot|_{\bullet}=|\cdot|_{\mathfrak{g}}$, and $0<r\le\pi/\sqrt{2N}$.
\end{itemize}
Let $S:=\{v\in\mathfrak{g}:|v|_{\bullet}<r\}$; then $S\subset\mathfrak{D}$, as shown in the proof of
Lemma~\ref{4.36:lem-haar-small-sets}. Let
\begin{equation*}
B:=u^0\exp(iS)=\{u^0e^{iv}:v\in S\},
\end{equation*}
which is the ball $B^{\mathrm{op}}_r(u^0)$ in case (op) and the ball $B^{\mathfrak{g}}_r(u^0)$ in case
($\mathfrak{g}$) of Definition~\ref{4.36:def-haar-balls}, and assume $B\subset W$. Let $\Lambda,\eta_2\ge0$
be such that for all $y\in\mathcal{W}$, all $u\in B$ and all $T\in\mathfrak{g}$ with $|T|_{\bullet}=1$:
\begin{equation}\label{4.36:eq-mode-bound-a}
\begin{aligned}
&\text{(m1)}\quad t\mapsto F_1(ue^{itT},y)\ \text{is of class } C^2 \text{ on a neighbourhood of } t=0,\\
&\phantom{\text{(m1)}\quad}\text{and}\ \Bigl|\partial_t^2F_1(ue^{itT},y)\big|_{t=0}\Bigr|\le\Lambda;\\
&\text{(m2)}\quad |F_2(u,y)-F_2(u^0,y)|\le\eta_2 .
\end{aligned}
\end{equation}
Then
\begin{equation}\label{4.36:eq-mode-bound-1}
f(u^0)\le\frac{e^{\Lambda r^2/2+\eta_2}}{\operatorname{Haar}(B)} ,
\end{equation}
and hence, in case (op),
\begin{equation}\label{4.36:eq-mode-bound-2}
f(u^0)\le c^{\mathrm{op}}_3(N)^{-1}e^{\Lambda r^2/2+\eta_2}r^{-D}
\le c^{\flat}_3(N)^{-1}e^{\Lambda r^2/2+\eta_2}r^{-D},
\end{equation}
and, in case ($\mathfrak{g}$),
\begin{equation}\label{4.36:eq-mode-bound-3}
f(u^0)\le c_3(N)^{-1}e^{\Lambda r^2/2+\eta_2}r^{-D},
\end{equation}
with the constants $c_3(N)$, $c^{\flat}_3(N)$, $c^{\mathrm{op}}_3(N)$ of
Lemma~\ref{4.36:lem-haar-small-sets}.
\end{lemma}

\begin{proof}
Since $u^0\in W$, we have $f(u^0)=Z(u^0)\big/\int_W e^{\Phi}\,d\operatorname{Haar}$. If $Z(u^0)=0$, then
$f(u^0)=0$ and there is nothing to prove. We assume from now on that $Z(u^0)\in(0,\infty)$, so that
$\Phi(u^0)$ is a real number.

\emph{Step 1: one line through $u^0$, both signs.} Fix $y\in\mathcal{W}$ and $v\in S\setminus\{0\}$, and put
$\hat v:=v/|v|_{\bullet}$, so that $|\hat v|_{\bullet}=1$. For $\tau\in(-r,r)$ we have
$|\tau\hat v|_{\bullet}=|\tau|<r$, hence $u^0e^{i\tau\hat v}\in B\subset W$, and we may define
\begin{equation*}
\varphi_y(\tau):=F_1(u^0e^{i\tau\hat v},y),\qquad \tau\in(-r,r).
\end{equation*}
Let $\tau_0\in(-r,r)$ and $u_1:=u^0e^{i\tau_0\hat v}$; as just observed, $u_1\in B$. The matrices
$e^{i\tau_0\hat v}$ and $e^{it\hat v}$ commute, so $e^{i(\tau_0+t)\hat v}=e^{i\tau_0\hat v}e^{it\hat v}$ and
\begin{equation*}
\varphi_y(\tau_0+t)=F_1(u_1e^{it\hat v},y).
\end{equation*}
By hypothesis (m1) of \eqref{4.36:eq-mode-bound-a}, applied with $u=u_1$ and $T=\hat v$, the right-hand
side is of class $C^2$ in $t$ on a neighbourhood of $t=0$, with second derivative at $t=0$ bounded in
absolute value by $\Lambda$; that is, $\varphi_y$ is $C^2$ near $\tau_0$ and $|\varphi_y''(\tau_0)|\le\Lambda$.
Since being of class $C^2$ is a local property and $\tau_0\in(-r,r)$ was arbitrary,
$\varphi_y\in C^2((-r,r))$ with $|\varphi_y''|\le\Lambda$ on $(-r,r)$.

Apply Taylor's formula with Lagrange remainder to $\varphi_y$ at $0$, at the two points
$\tau=\pm|v|_{\bullet}$, both of which lie in $(-r,r)$: with $a:=|v|_{\bullet}\varphi_y'(0)\in\mathbb{R}$
and for some $\theta_\pm$ between $0$ and $\pm|v|_{\bullet}$,
\begin{equation*}
\begin{aligned}
\varphi_y(\pm|v|_{\bullet})&=\varphi_y(0)\pm a+\tfrac12\varphi_y''(\theta_\pm)|v|_{\bullet}^2\\
&\ge\varphi_y(0)\pm a-\tfrac12\Lambda|v|_{\bullet}^2.
\end{aligned}
\end{equation*}
Since $u^0e^{i(\pm|v|_{\bullet})\hat v}=u^0e^{\pm iv}$, one and the same function $\varphi_y$ serves for
both signs, and the last inequality reads
\begin{equation*}
F_1(u^0e^{\pm iv},y)-F_1(u^0,y)\ge\pm a-\tfrac12\Lambda|v|_{\bullet}^2 .
\end{equation*}
Both points $u^0e^{\pm iv}$ lie in $B$, so hypothesis (m2) of \eqref{4.36:eq-mode-bound-a} at
$u=u^0e^{\pm iv}$ gives $F_2(u^0e^{\pm iv},y)-F_2(u^0,y)\ge-\eta_2$. For $v\in S$ put
\begin{equation*}
\delta_y(v):=F(u^0e^{iv},y)-F(u^0,y)\in\mathbb{R}.
\end{equation*}
Adding the two bounds, $\delta_y(\pm v)\ge\pm a-\tfrac12\Lambda|v|_{\bullet}^2-\eta_2$. Hence, using
$\tfrac12(e^{a}+e^{-a})\ge1$ and $|v|_{\bullet}<r$,
\begin{equation}\label{4.36:eq-mode-bound-4}
\tfrac12\bigl(e^{\delta_y(v)}+e^{\delta_y(-v)}\bigr)
\ge e^{-\frac12\Lambda|v|_{\bullet}^2-\eta_2}\cdot\tfrac12\bigl(e^{a}+e^{-a}\bigr)
\ge e^{-\Lambda r^2/2-\eta_2},
\end{equation}
pointwise in $y\in\mathcal{W}$. At $v=0$ we have $\delta_y(0)=0$, the left-hand side of
\eqref{4.36:eq-mode-bound-4} equals $1$, and \eqref{4.36:eq-mode-bound-4} holds as well. Thus
\eqref{4.36:eq-mode-bound-4} holds for every $y\in\mathcal{W}$ and every $v\in S$.

\emph{Step 2: the reference probability.} Since $Z(u^0)\in(0,\infty)$,
\begin{equation*}
\nu(dy):=\frac{e^{F(u^0,y)}\,\mu(dy)}{Z(u^0)}
\end{equation*}
is a probability measure on $\mathcal{W}$. For $v\in S$, writing
$e^{F(u^0e^{iv},y)}=e^{\delta_y(v)}e^{F(u^0,y)}$ in the integral defining $Z(u^0e^{iv})$, we obtain
\begin{equation*}
\frac{Z(u^0e^{iv})}{Z(u^0)}=\int_{\mathcal{W}}e^{\delta_y(v)}\,\nu(dy)\in(0,\infty],
\end{equation*}
the value being positive because the integrand is positive and $\nu$ is a probability measure.

\emph{Step 3: the lower bound on the normalisation.} Since $B\subset W$ and the integrand is
non-negative,
\begin{equation*}
\int_W e^{\Phi-\Phi(u^0)}\,d\operatorname{Haar}\ge\int_B\frac{Z(u)}{Z(u^0)}\,\operatorname{Haar}(du).
\end{equation*}
By the formula \eqref{4.36:eq-haar-balls-a} for Haar measure in exponential coordinates, which applies
because $S\subset\mathfrak{D}$, together with the invariance of Haar measure under left translation by
$u^0$ (Lemma~\ref{4.36:lem-haar-small-sets}(c)), the right-hand side equals
\begin{equation*}
\sigma_N\int_S\frac{Z(u^0e^{iv})}{Z(u^0)}\,J(v)\,dv
=\sigma_N\int_{\mathcal{W}}\nu(dy)\int_S e^{\delta_y(v)}J(v)\,dv ,
\end{equation*}
where we inserted Step~2 and exchanged the order of integration by Tonelli's theorem; the integrand
$(v,y)\mapsto e^{\delta_y(v)}J(v)$ is non-negative and jointly measurable, since $F$ is jointly
measurable and $v\mapsto u^0e^{iv}$ is continuous.

Fix $y$. The map $v\mapsto-v$ maps $S$ onto $S$ (the set $S$ is defined by a norm), preserves Lebesgue
measure $dv$, and $J(-v)=J(v)$. Therefore, in $[0,\infty]$,
\begin{equation*}
\begin{aligned}
\int_S e^{\delta_y(v)}J(v)\,dv&=\int_S e^{\delta_y(-v)}J(v)\,dv\\
&=\int_S\tfrac12\bigl(e^{\delta_y(v)}+e^{\delta_y(-v)}\bigr)J(v)\,dv\\
&\ge e^{-\Lambda r^2/2-\eta_2}\int_S J(v)\,dv,
\end{aligned}
\end{equation*}
the inequality by \eqref{4.36:eq-mode-bound-4} and $J\ge0$. Integrating this bound against the
probability measure $\nu$ and using \eqref{4.36:eq-haar-balls-a} once more, now for the indicator of
$B=u^0\exp(iS)$, we arrive at
\begin{equation*}
\int_W e^{\Phi-\Phi(u^0)}\,d\operatorname{Haar}
\ge e^{-\Lambda r^2/2-\eta_2}\,\sigma_N\int_S J(v)\,dv
=e^{-\Lambda r^2/2-\eta_2}\operatorname{Haar}(B).
\end{equation*}

\emph{Step 4: conclusion.} Since $u^0\in W$ and $\Phi(u^0)$ is finite,
\begin{equation*}
f(u^0)=\frac{e^{\Phi(u^0)}}{\int_W e^{\Phi}\,d\operatorname{Haar}}
=\frac{1}{\int_W e^{\Phi-\Phi(u^0)}\,d\operatorname{Haar}}
\le\frac{e^{\Lambda r^2/2+\eta_2}}{\operatorname{Haar}(B)},
\end{equation*}
which is \eqref{4.36:eq-mode-bound-1}. In case (op), $B=B^{\mathrm{op}}_r(u^0)$ with $0<r\le\pi/2$, and
Lemma~\ref{4.36:lem-haar-small-sets}(a$^{\mathrm{op}}$) gives
$\operatorname{Haar}(B)\ge c^{\mathrm{op}}_3(N)r^D\ge c^{\flat}_3(N)r^D$; inserting this into
\eqref{4.36:eq-mode-bound-1} yields \eqref{4.36:eq-mode-bound-2}. In case ($\mathfrak{g}$),
$B=B^{\mathfrak{g}}_r(u^0)$ with $0<r\le\pi/\sqrt{2N}$, and
Lemma~\ref{4.36:lem-haar-small-sets}(a) gives $\operatorname{Haar}(B)\ge c_3(N)r^D$, whence
\eqref{4.36:eq-mode-bound-3}.

The proof uses only the upper bound (m1) on second derivatives along one-parameter curves through points
of $B$, the symmetry of $S$ under $v\mapsto-v$, the evenness of $J$ and
Lemma~\ref{4.36:lem-haar-small-sets}; it uses no hypothesis on first derivatives of $F_1$ and no
integrability of $\varphi_y'(0)$ in $y$, since $e^{a}+e^{-a}\ge2$ holds pointwise in $y$.
\end{proof}

\begin{definition}[One-link complexification; operator-unit and $\mathfrak{g}$-unit directions]
\label{4.36:def-unit-directions}
Let $G=SU(N)$, let $\mathfrak{g}$ be the space of traceless Hermitian $N\times N$ matrices and
$\mathfrak{g}^{c}=\mathfrak{sl}(N,\mathbb{C})$. Write $|\cdot|$ for the operator norm and
$|X|_{\mathfrak{g}}=(\operatorname{tr}X^2)^{1/2}$ for $X\in\mathfrak{g}$, with
$\operatorname{tr}\mathbf{1}=1$.

\emph{Setting.} Fix the level $k$, a plaquette $p$ of the lattice at level $k$, a bond
$b_1\subset\partial p$, and write the configuration as $V=(u_1,V_{\mathrm{off}})$, where
$u_1=V_k(b_1)$ is the free link. Put $R_*:=20L^2M\max\{R_k,R_{k+1}\}$, with $R_k$ the integers of
Notation~\ref{4.36:not-flow-thresholds}. Let $\rho^{B}$ be an effective density at level $k$ in the
sense of Definition~\ref{1.5:def-densities}, written as a sum of terms $h$, each with a region
$\Lambda_k(h)$. A term $h$ is \emph{clean} at $p$ if $\Lambda_k(h)$ contains $p$ with
$\operatorname{dist}_k(p,\Lambda_k(h)^{c})\ge R_*$, where $\operatorname{dist}_k$ is the distance
at level $k$. Let $h$ be a clean term. Its conditional exponent splits as
\begin{equation}\label{4.36:eq-unit-directions-1}
\begin{gathered}
F_h=F^{\mathrm{near}}+F^{\mathrm{far}},\qquad
F^{\mathrm{near}}:=A_k(h)=F^{W}+F^{E}+F^{R}+F^{B}+\mathrm{const},\\
F^{\mathrm{far}}:=Q_h+\sum_{Y}\Psi_Y ,
\end{gathered}
\end{equation}
where the near exponent $A_k(h)$ of $h$ with its four members, and the far members $Q_h$ and
$\Psi_Y$, are the members of the expansion of $\rho^{B}$ at level $k$.
The exponent is considered at a real pair $(V,w)$, both the configuration $V$ and the second
argument $w$ of $F_h$ being real, with the free link in the window
$W(V_{\mathrm{off}})$, the open
set of $u\in G$ at which every local small-field indicator $\chi_{\square}$ of a cube
$\square\in\mathcal{N}$ near $p$ equals one.

\emph{One-link complexification.} For $u\in G$, $T\in\mathfrak{g}$ and $\zeta\in\mathbb{C}$ put
\begin{equation*}
T':=uTu^{-1},\qquad V^{\zeta,T}:=(u\,e^{i\zeta T},V_{\mathrm{off}}),
\end{equation*}
and let $V'$ be the configuration equal to $e^{i\zeta T'}$ on $b_1$ and to $\mathbf{1}$ on every
other bond. Then:
\begin{itemize}
\item[(a)] $T'\in\mathfrak{g}$, $|T'|=|T|$, $|T'|_{\mathfrak{g}}=|T|_{\mathfrak{g}}$, and
$u\,e^{i\zeta T}=e^{i\zeta T'}u$; hence $V^{\zeta,T}=V'\cdot(u,V_{\mathrm{off}})$ bondwise. This is
the notation of \cite[Proposition~9]{B-CMP102b}, which concerns an analytic function of
small configurations $V'$ with values in the complexification $G^{c}=SL(N,\mathbb{C})$ of $G$,
with logarithmic variable $B=\frac1i\log V'$, under the
hypothesis $|V'-\mathbf{1}|<C^{(102)}_1\varepsilon_1$ of \cite[(172)]{B-CMP102b}, where
$C^{(102)}_1,\varepsilon_1$ are the constants of that display; the superscript separates
$C^{(102)}_1$ from the radius constant $C_1$ of Notation~\ref{4.36:not-flow-thresholds}.
\item[(b)] $|V'-\mathbf{1}|\le e^{|\zeta||T|}-1$.
\item[(c)] If $|\operatorname{Re}\zeta|\,|T|<\pi$, the principal logarithm gives $B=\zeta T'$ on
$b_1$ and $B=0$ elsewhere. This $B$ is traceless, $\partial_\zeta B(b_1)=T'$ and
$|\partial_\zeta B(b_1)|=|T|$.
\end{itemize}
Let $\mathcal{H}$ be the analytic function of \cite[Proposition~9]{B-CMP102b}, as a function of the
logarithmic variable. For every $\zeta$ with $|\operatorname{Re}\zeta|\,|T|<\pi$ at which $V'$
satisfies the hypothesis of \cite[(172)]{B-CMP102b}, the logarithmic variable of $V'$ is
$\zeta T'$ by (c), and we put $\mathcal{H}_{\zeta}:=\mathcal{H}(\zeta T')$; by (b) that hypothesis
holds whenever $e^{|\zeta||T|}-1<C^{(102)}_1\varepsilon_1$. It is
$\mathfrak{g}^{c}$-valued. At real $\zeta$ it is $\mathfrak{g}$-valued, and there the continuation
coincides with the minimal configuration \cite{B-CMP102b}.

In the curvature lemma for the conditional exponent at $SU(N)$, the norm of the direction $T$
enters only through three quantities: $|V'-\mathbf{1}|$, $\partial_\zeta B(b_1)=T'$ with
$|\partial_\zeta B(b_1)|=|T|$, and the tracelessness of $T'$.
Everything else in that lemma is a function of the field $\mathcal{H}_{\zeta}$ and of scalar data.
This covers its analyticity bounds, the mode bound, the radii $2R$ and $r_m$ of threshold
$(\mathrm{t}2)$ in \eqref{4.36:eq-flow-thresholds-b}, the Cauchy bounds, the four near members,
and the far part.

\emph{Unit directions.} A direction $T\in\mathfrak{g}$ is \emph{operator-unit} if $|T|=1$, and
\emph{$\mathfrak{g}$-unit} if $|T|_{\mathfrak{g}}=1$. At $N=2$ the two notions coincide; for
$N\ge3$ they differ. The deviation of a link variable from its reference value is measured in the
operator norm $|\cdot|$ in \cite[(19)]{B-CMP98} and in \cite[(172)]{B-CMP102b}, where it is
$|V'-\mathbf{1}|$. Accordingly, a unit direction in the curvature lemma means, at every $N$, an
operator-unit direction.

\emph{Constants.} Every constant of the curvature lemma is read at the values its inputs take
for $G=SU(N)$, in the sense of Notation~\ref{4.36:not-flow-thresholds}. In particular, with
$\varpi$ the polylogarithmic weight of the arrest step, a function of the integer $j\ge0$,
we put
\begin{align}
C_q(N)&:=2C_P\bigl(B_3a'_1/r_0+C_P\bigr),\label{4.36:eq-unit-directions-2}\\
C'_W(N)&:=36\,C_q(N)\,C_\sigma\,S_{\mathfrak{n}},\qquad
S_{\mathfrak{n}}:=\sum_{j\ge0}e^{-j}\bigl[2+\bar B\bigl(j+\varpi(j)\bigr)\bigr].
\label{4.36:eq-unit-directions-3}
\end{align}
Here $a'_1$ is the constant of \cite[Theorem~1]{B-CMP102b}, $C_P$ is the constant that enters the
threshold $(\mathrm{t}8)$ of \eqref{4.36:eq-flow-thresholds-b} through $r_0\le(2C_P)^{-1}$, and
$B_3$ is one of the constants of Ba{\l}aban's construction; $a'_1$, $B_3$ and $C_P$ are all read
at $G=SU(N)$, and
$r_0$ is the radius constrained by the threshold $(\mathrm{t}8)$ of
\eqref{4.36:eq-flow-thresholds-b}.
The constant $C_\sigma$ is the one in the bound
$\sigma^{(\alpha)}_n\le C_\sigma e^{-(k-n)}$ on the stratum weight sums, and $\bar B$ is the
frozen flow constant of Notation~\ref{4.36:not-flow-thresholds}.

Thus $C'_W(N)$ is a number depending on $(N,L,\mathfrak{p})$. It carries $\bar B$, which satisfies
\begin{equation*}
\bar B\ \ge\ c_0(N)/4=\frac{11N^2}{48\pi^2}\log L .
\end{equation*}
The constant $C'_W(N)$ is never compared with a power of $L$, and no bound on it uniform in $N$ is
used.

\emph{Thresholds.} The curvature lemma is used under the thresholds $(\mathrm{t}1)$,
$(\mathrm{t}2)$, $(\mathrm{t}3)'$, $(\mathrm{t}4)'$,
$(\mathrm{t}5)$, $(\mathrm{t}6)$, $(\mathrm{t}8)$, $(\mathrm{t}9)$ of
\eqref{4.36:eq-flow-thresholds-b}, under
$x_k\ge x_{\mathrm{th}}(N,L,\mathfrak{p})$, and under
\begin{equation}\label{4.36:eq-unit-directions-4}
(\mathrm{t}8)^{+}\qquad g_k<r_0 .
\end{equation}
Condition \eqref{4.36:eq-unit-directions-4} ensures that the real segment $[-2g_k,2g_k]$ lies in
the open disc $|\zeta|<2r_0$, on which the analyticity bound of the curvature lemma holds.

No condition of the form $L^2\ge1+\bar B$ is imposed, at any $N$. Each inequality in the estimates
accompanying the curvature lemma that would call for such a condition is obtained instead from the
transport bounds under the coupling ceiling of Proposition~\ref{4.36:prop-transport-every-l}, whose
proof is incomplete at one step. For the one further statement applied together with the curvature
lemma, the absence of such a condition from its hypotheses is stated, with a proof that is
incomplete.
In this way no condition on $L$ depending on $N$ enters the hypotheses under which the curvature
lemma is used.
\end{definition}

\begin{proof}
\emph{Proof of (a).} Since $u$ is unitary, $u^{-1}=u^{*}$. Hence $(T')^{*}=uT^{*}u^{*}=T'$.
Moreover $\operatorname{tr}T'=\operatorname{tr}T=0$ by cyclicity of the trace, so $T'\in\mathfrak{g}$.
The operator norm is invariant under unitary conjugation, so $|T'|=|T|$. Since
$T'^2=uT^2u^{-1}$, cyclicity gives $\operatorname{tr}T'^2=\operatorname{tr}T^2$, that is,
$|T'|_{\mathfrak{g}}=|T|_{\mathfrak{g}}$.

Summing the exponential series term by term and using $u(i\zeta T)^nu^{-1}=(i\zeta T')^n$, we get
$u\,e^{i\zeta T}u^{-1}=e^{i\zeta T'}$, that is, $u\,e^{i\zeta T}=e^{i\zeta T'}u$. Therefore
$V^{\zeta,T}$ agrees with $V'\cdot(u,V_{\mathrm{off}})$ on $b_1$. On every other bond $V'$ equals
$\mathbf{1}$, so the two configurations agree there as well.

\emph{Proof of (b).} This is the exponential-series bound among the elementary inequalities of
Lemma~\ref{4.36:lem-eigenangle-inequalities}; we verify it directly. Put $X=i\zeta T'$. By (a),
$|X|=|\zeta||T|$. By the triangle inequality and submultiplicativity of the operator norm,
\begin{equation*}
|V'(b_1)-\mathbf{1}|=\Bigl|\sum_{n\ge1}\frac{X^n}{n!}\Bigr|
\le\sum_{n\ge1}\frac{|X|^n}{n!}=e^{|\zeta||T|}-1 .
\end{equation*}
On the other bonds $V'-\mathbf{1}=0$.

\emph{Proof of (c).} Since $T'$ is Hermitian, there is a unitary $v$ with
$T'=v\operatorname{diag}(\theta_1,\dots,\theta_N)v^{*}$, where the eigen-angles $\theta_q$ are real,
$|\theta_q|\le|T'|=|T|$ and $\sum_q\theta_q=0$. Then
\begin{equation*}
V'(b_1)=v\operatorname{diag}\bigl(e^{i\zeta\theta_1},\dots,e^{i\zeta\theta_N}\bigr)v^{*},
\end{equation*}
so $V'(b_1)$ is normal. Each eigenvalue can be written
$e^{i\zeta\theta_q}=e^{-(\operatorname{Im}\zeta)\theta_q}\,e^{i(\operatorname{Re}\zeta)\theta_q}$.
Its argument is $(\operatorname{Re}\zeta)\theta_q$, which satisfies
\begin{equation*}
|(\operatorname{Re}\zeta)\theta_q|\le|\operatorname{Re}\zeta|\,|T|<\pi .
\end{equation*}
So every eigenvalue lies off $(-\infty,0]$. The principal logarithm of $e^{i\zeta\theta_q}$ is
$i\zeta\theta_q$, because its imaginary part $(\operatorname{Re}\zeta)\theta_q$ lies in
$(-\pi,\pi)$.

The principal logarithm of a normal matrix with spectrum off $(-\infty,0]$ acts on the eigenvalues
in the same unitary frame. Hence
\begin{equation*}
\log V'(b_1)=v\operatorname{diag}(i\zeta\theta_q)v^{*}=i\zeta T',
\qquad\text{so}\qquad B(b_1)=\tfrac1i\log V'(b_1)=\zeta T'.
\end{equation*}
On the other bonds $\log\mathbf{1}=0$. Since $\operatorname{tr}T'=0$, the variable $B$ is traceless.
The map $\zeta\mapsto\zeta T'$ is linear, so $\partial_\zeta B(b_1)=T'$, and
$|\partial_\zeta B(b_1)|=|T'|=|T|$ by (a).

\emph{The field $\mathcal{H}_{\zeta}$.} If $e^{|\zeta||T|}-1<C^{(102)}_1\varepsilon_1$, then (b)
gives $|V'-\mathbf{1}|<C^{(102)}_1\varepsilon_1$, which is the hypothesis of
\cite[(172)]{B-CMP102b}. Now let $\zeta$ be as in the definition of
$\mathcal{H}_{\zeta}$. Since $T'$ is traceless by (a), $\zeta T'\in\mathfrak{g}^{c}$
for every such $\zeta$; the configuration $V'$ satisfies the hypothesis of
\cite[(172)]{B-CMP102b} by the choice of $\zeta$, and by (c) its logarithmic variable is
$\zeta T'$. Hence
$\mathcal{H}_{\zeta}=\mathcal{H}(\zeta T')$ is the value of the continuation of
\cite[Proposition~9]{B-CMP102b} at the $\mathfrak{g}^{c}$-valued argument $\zeta T'$, and by that
proposition it is $\mathfrak{g}^{c}$-valued. At real $\zeta$ the argument $\zeta T'$ lies in
$\mathfrak{g}$, and
$V'(b_1)=e^{i\zeta T'}\in G$. There the continuation coincides with the minimal configuration
\cite{B-CMP102b}, and it is $\mathfrak{g}$-valued.

\emph{How the direction enters.} The four near members, the far part and the radii and bounds
listed above are functions of $\mathcal{H}_{\zeta}$ and of scalar data. By definition
$\mathcal{H}_{\zeta}=\mathcal{H}(\zeta T')$. The admissibility of the argument under
\cite[(172)]{B-CMP102b} is controlled by $|V'-\mathbf{1}|$, through (b). The derivative in $\zeta$
is governed by $\partial_\zeta B=T'$, whose norm is given by (c), and the tracelessness of $T'$
comes from (a). Hence the norm of $T$ enters only through these three quantities.

\emph{Coincidence at $N=2$ and difference for $N\ge3$.} At $N=2$ a nonzero $X\in\mathfrak{g}$ has
eigen-angles $\pm\theta$ with $\theta>0$. Then $|X|=\theta$ and
\begin{equation*}
|X|_{\mathfrak{g}}=\bigl((\theta^2+\theta^2)/2\bigr)^{1/2}=\theta ,
\end{equation*}
so the two unit conditions coincide. For $N\ge3$, Lemma~\ref{4.36:lem-eigenangle-inequalities} gives
$|X|_{\mathfrak{g}}\le|X|\le\mathsf{n}|X|_{\mathfrak{g}}$, with $\mathsf{n}=(N-1)^{1/2}>1$ and
equality in the second inequality at $X=\operatorname{diag}(1,\dots,1,1-N)$. At that $X$ the two
norms differ, so the two unit conditions differ.

\emph{The lower bound on $\bar B$.} Notation~\ref{4.36:not-flow-thresholds} gives
$\bar B\ge c_0(N)/4$ with $c_0(N)=\frac{11N^2}{12\pi^2}\log L$. Dividing by four gives the displayed
value $\frac{11N^2}{48\pi^2}\log L$.

\emph{The threshold \eqref{4.36:eq-unit-directions-4}.} If $g_k<r_0$, then every real $\zeta$ with
$|\zeta|\le2g_k$ satisfies $|\zeta|\le 2g_k<2r_0$, so it lies in the open disc $|\zeta|<2r_0$.
\end{proof}

\begin{lemma}[Second directional derivatives scale with the squared norm]
\label{4.36:lem-derivative-rescaling}
Let $u\in SU(N)$, let $T\in\mathfrak{g}$ with $T\neq0$, and put $\hat{T}:=T/|T|$. Let $F$ be any
function defined on a neighbourhood of $u$: a neighbourhood in $SU(N)$ when only real parameters
are used, and in $SL(N,\mathbb{C})$ when complex parameters are used. Then, for real or complex $t$,
\begin{equation}\label{4.36:eq-derivative-rescaling-1}
F\bigl(u e^{itT}\bigr)=F\bigl(u e^{i(t|T|)\hat{T}}\bigr)
\end{equation}
identically, wherever either side is defined. Consequently:
\begin{enumerate}
\item the function $t\mapsto F(ue^{itT})$ is $C^2$ (respectively analytic) near $0$ if and only if
$s\mapsto F(ue^{is\hat{T}})$ is;
\item when these functions are $C^2$ near $0$, the directional derivatives
\begin{equation*}
X_TF:=\frac{d}{dt}\Big|_{t=0}F\bigl(ue^{itT}\bigr),\qquad
X_T^2F:=\frac{d^2}{dt^2}\Big|_{t=0}F\bigl(ue^{itT}\bigr),
\end{equation*}
and the corresponding $X_{\hat{T}}F$, $X^2_{\hat{T}}F$, satisfy
\begin{equation}\label{4.36:eq-derivative-rescaling-2}
X_TF=|T|\,X_{\hat{T}}F,\qquad X_T^2F=|T|^2\,X^2_{\hat{T}}F;
\end{equation}
\item for every $S>0$ and every property $P$ of matrices, the assertion
\emph{$P(ue^{is\hat{T}})$ holds for all real $s$ with $|s|\le S$} is literally the assertion
\emph{$P(ue^{itT})$ holds for all real $t$ with $|t|\,|T|\le S$}.
\end{enumerate}
Moreover, if $|T|_{\mathfrak{g}}=1$ then $|T|^2\le N-1$, with equality for the direction
$T=\operatorname{diag}(1,\dots,1,1-N)/\mathsf{n}$.
\end{lemma}

\begin{proof}
Since $T\neq0$, its operator norm $|T|$ is positive, so $\hat{T}$ is defined. It is a real multiple
of a traceless Hermitian matrix, hence $\hat{T}\in\mathfrak{g}$, and $|\hat{T}|=1$. For every real or
complex $t$ we have the matrix identity $tT=(t|T|)\hat{T}$. Hence
$e^{itT}=e^{i(t|T|)\hat{T}}$, so the two arguments $ue^{itT}$ and $ue^{i(t|T|)\hat{T}}$ are the same
matrix. This gives \eqref{4.36:eq-derivative-rescaling-1}.

Put $\varphi(t):=F(ue^{itT})$ and $\psi(s):=F(ue^{is\hat{T}})$. By
\eqref{4.36:eq-derivative-rescaling-1}, $\varphi(t)=\psi(|T|t)$, and equivalently
$\psi(s)=\varphi(s/|T|)$. Multiplication by $|T|$ is a linear bijection of $\mathbb{R}$, and also of
$\mathbb{C}$, that maps neighbourhoods of $0$ onto neighbourhoods of $0$. Each of $\varphi$, $\psi$
is therefore the composition of the other with a linear map. Such a composition is $C^2$,
respectively analytic, near $0$ if and only if the original function is, which proves item~1.

By the chain rule applied to $\varphi(t)=\psi(|T|t)$ we obtain $\varphi'(0)=|T|\psi'(0)$ and
$\varphi''(0)=|T|^2\psi''(0)$. These are the two identities in
\eqref{4.36:eq-derivative-rescaling-2}, which proves item~2.

For item~3, the map $t\mapsto s=t|T|$ is a bijection from $\{t\in\mathbb{R}:|t|\,|T|\le S\}$ onto
$\{s\in\mathbb{R}:|s|\le S\}$. By \eqref{4.36:eq-derivative-rescaling-1} it carries the point
$ue^{itT}$ to the same point $ue^{is\hat{T}}$. The two families of points therefore coincide,
and so do the two assertions.

Finally, let $|T|_{\mathfrak{g}}=1$. The second inequality of
Lemma~\ref{4.36:lem-eigenangle-inequalities}, in \eqref{4.36:eq-eigenangle-inequalities-a}, gives
$|T|\le\mathsf{n}|T|_{\mathfrak{g}}=\mathsf{n}$, hence $|T|^2\le\mathsf{n}^2=N-1$.

Now let $X:=\operatorname{diag}(1,\dots,1,1-N)$. This matrix is Hermitian, and its trace
$(N-1)+(1-N)$ vanishes, so $X\in\mathfrak{g}$. Its operator norm is $|X|=N-1$. With the normalised
trace,
\begin{equation*}
\operatorname{tr}X^2=\frac{(N-1)+(N-1)^2}{N}=N-1,
\end{equation*}
so $|X|_{\mathfrak{g}}=\mathsf{n}$. Put $T_0:=X/\mathsf{n}$. Then $T_0\in\mathfrak{g}$ and
$|T_0|_{\mathfrak{g}}=1$. Moreover $|T_0|=(N-1)/\mathsf{n}=\mathsf{n}$, so $|T_0|^2=N-1$, and the
bound is attained in the direction $T_0$.
\end{proof}

\begin{lemma}[The complexified plaquette action: holomorphy, increment and size]
\label{4.36:lem-plaquette-action}
Let $|\cdot|$ denote the operator norm and $\|\cdot\|$ the normalised trace norm,
$\|Y\|^2=\operatorname{tr}Y^*Y$, as in Notation~\ref{4.36:not-group-norms}, with
$\operatorname{tr}=N^{-1}\operatorname{Tr}$ and $\operatorname{tr}\mathbf{1}=1$. On $GL(N,\mathbb{C})$ put
\begin{equation}\label{4.36:eq-plaquette-action-1}
\mathfrak{w}(W):=1-\tfrac12\bigl(\operatorname{tr}W+\operatorname{tr}W^{-1}\bigr).
\end{equation}
\begin{itemize}
\item[(0)] $\mathfrak{w}$ is holomorphic on $GL(N,\mathbb{C})$, and
\begin{equation}\label{4.36:eq-plaquette-action-2}
\mathfrak{w}(W)=-\tfrac12\operatorname{tr}\bigl[W^{-1}(W-\mathbf{1})^2\bigr].
\end{equation}
On $U(N)$ one has $\mathfrak{w}(W)=1-\operatorname{Re}\operatorname{tr}W$. On $SL(N,\mathbb{C})$,
$\mathfrak{w}$ is Ba{\l}aban's prescription for $1-\operatorname{Re}\operatorname{tr}U(\partial p)$ at
$SL(N,\mathbb{C})$-valued configurations in \cite{B-CMP99b}, and hence the holomorphic continuation of
$1-\operatorname{Re}\operatorname{tr}$ from $SU(N)$.
\item[(a$^\flat$)] (Mixed form.) Let $\mathsf{k}\ge0$ and $\tau\ge0$. If $W_0\in U(N)$,
$\|W_0-\mathbf{1}\|\le\mathsf{k}$, and $W=QW_0$ with $\|Q-\mathbf{1}\|\le\tau$ and
$|Q-\mathbf{1}|\le\frac12$, then
\begin{equation*}
|\mathfrak{w}(W)-\mathfrak{w}(W_0)|\le(\mathsf{k}+\tau)\tau .
\end{equation*}
\item[(a)] (Operator norm.) Let $\mathsf{k}\ge0$ and $\tau\ge0$. If $W_0\in U(N)$,
$|W_0-\mathbf{1}|\le\mathsf{k}$, and $W=QW_0$ with $|Q-\mathbf{1}|\le\tau\le\frac12$, then
\begin{equation}\label{4.36:eq-plaquette-action-3}
|\mathfrak{w}(W)-\mathfrak{w}(W_0)|\le(\mathsf{k}+\tau)\tau .
\end{equation}
This bound holds with constant one, for every $N$, and without any restriction such as
$\mathsf{k}\le1$.
\item[(b)] If $W\in GL(N,\mathbb{C})$ and $|W-\mathbf{1}|\le\frac12$, then
\begin{equation}\label{4.36:eq-plaquette-action-4}
|\mathfrak{w}(W)|\le\tfrac12|W^{-1}|\,\|W-\mathbf{1}\|^2\le\|W-\mathbf{1}\|^2\le|W-\mathbf{1}|^2 .
\end{equation}
\end{itemize}
\end{lemma}

\begin{proof}
We use four elementary facts about the two norms, all contained in
\eqref{4.36:eq-eigenangle-inequalities-a} of Lemma~\ref{4.36:lem-eigenangle-inequalities}. First,
$\|Y\|\le|Y|$. Second, $\|XY\|\le|X|\,\|Y\|$ and $\|XY\|\le\|X\|\,|Y|$. Third,
$|\operatorname{tr}(XY)|\le\|X\|\,\|Y\|$. Fourth, if $|Y-\mathbf{1}|\le\frac12$, then $Y$ is
invertible and $|Y^{-1}|\le2$.

\emph{Proof of} (0). On $GL(N,\mathbb{C})$ we have $W^{-1}=(\det W)^{-1}\operatorname{adj}W$; the
entries of $\operatorname{adj}W$ are polynomials in the entries of $W$, and $\det W$ is a
polynomial that does not vanish there, so the entries of $W^{-1}$ are holomorphic. Since
$\operatorname{tr}$ is linear, $\mathfrak{w}$ is holomorphic. Expanding the square,
\begin{equation*}
W^{-1}(W-\mathbf{1})^2=W^{-1}(W^2-2W+\mathbf{1})=W-2\cdot\mathbf{1}+W^{-1},
\end{equation*}
and taking the normalised trace with $\operatorname{tr}\mathbf{1}=1$ gives
$\operatorname{tr}[W^{-1}(W-\mathbf{1})^2]=\operatorname{tr}W+\operatorname{tr}W^{-1}-2$. Multiplying by
$-\frac12$ yields \eqref{4.36:eq-plaquette-action-2}. For $W\in U(N)$ we have $W^{-1}=W^*$, so
$\operatorname{tr}W^{-1}=\operatorname{tr}W^*=\overline{\operatorname{tr}W}$, and
$\frac12(\operatorname{tr}W+\operatorname{tr}W^{-1})=\operatorname{Re}\operatorname{tr}W$, hence
$\mathfrak{w}(W)=1-\operatorname{Re}\operatorname{tr}W$. At a configuration $U_0$ with values in
$SL(N,\mathbb{C})$, \cite{B-CMP99b} defines the real part of the plaquette variable
$U_0(\partial p)$ by
$\operatorname{Re}U_0(\partial p)=\frac12\bigl(U_0(\partial p)+U_0(-\partial p)\bigr)$ with
$U_0(-\partial p)=(U_0(\partial p))^{-1}$. Taking $\operatorname{tr}$ and subtracting from $1$ gives
exactly $\mathfrak{w}(U_0(\partial p))$. Thus $\mathfrak{w}$ is this prescription; it is holomorphic on
$SL(N,\mathbb{C})$ by the first assertion and coincides with $1-\operatorname{Re}\operatorname{tr}$ on
$SU(N)\subset U(N)$ by the second, so it is the holomorphic continuation of
$1-\operatorname{Re}\operatorname{tr}$ from $SU(N)$.

\emph{Proof of} (a$^\flat$) \emph{and} (a). Put $E:=Q-\mathbf{1}$. In both cases $|E|\le\frac12$:
in (a$^\flat$) this is a hypothesis, and in (a) it follows from $|E|\le\tau\le\frac12$. Hence $Q$ is
invertible with $|Q^{-1}|\le2$, and so is $W=QW_0$. The matrices $E$ and $Q=\mathbf{1}+E$ commute.
Therefore $Q^{-1}-\mathbf{1}=Q^{-1}(\mathbf{1}-Q)=-EQ^{-1}$. Since $Q-E=\mathbf{1}$, we also have
$-E+E^2Q^{-1}=-E(Q-E)Q^{-1}=-EQ^{-1}$, so
\begin{equation*}
Q^{-1}-\mathbf{1}=-EQ^{-1}=-E+E^2Q^{-1}.
\end{equation*}
Next, $W-W_0=(Q-\mathbf{1})W_0=EW_0$ and $W^{-1}=W_0^{-1}Q^{-1}$. Consequently
\begin{equation*}
W^{-1}-W_0^{-1}=W_0^{-1}(Q^{-1}-\mathbf{1})=-W_0^{-1}E+W_0^{-1}E^2Q^{-1}.
\end{equation*}
Inserting these two differences into \eqref{4.36:eq-plaquette-action-1} gives
\begin{align*}
\mathfrak{w}(W)-\mathfrak{w}(W_0)
&=-\tfrac12\bigl[\operatorname{tr}(EW_0)-\operatorname{tr}(W_0^{-1}E)
+\operatorname{tr}(W_0^{-1}E^2Q^{-1})\bigr]\\
&=-\tfrac12\bigl[\operatorname{tr}\bigl(E(W_0-W_0^{-1})\bigr)
+\operatorname{tr}\bigl((W_0^{-1}E)(EQ^{-1})\bigr)\bigr].
\end{align*}
In the second line we used cyclicity of the trace,
$\operatorname{tr}(W_0^{-1}E)=\operatorname{tr}(EW_0^{-1})$.

We estimate the two traces. By the third fact,
$|\operatorname{tr}(E(W_0-W_0^{-1}))|\le\|E\|\,\|W_0-W_0^{-1}\|$. Since
$W_0^{-1}(W_0-\mathbf{1})=\mathbf{1}-W_0^{-1}$, we can write
\begin{equation*}
W_0-W_0^{-1}=(W_0-\mathbf{1})+W_0^{-1}(W_0-\mathbf{1}).
\end{equation*}
Because $W_0$ is unitary, $|W_0^{-1}|=1$. The second fact then gives
\begin{equation*}
\|W_0-W_0^{-1}\|\le\|W_0-\mathbf{1}\|+|W_0^{-1}|\,\|W_0-\mathbf{1}\|=2\|W_0-\mathbf{1}\|.
\end{equation*}
By the third and second facts, together with $|Q^{-1}|\le2$,
\begin{equation*}
|\operatorname{tr}((W_0^{-1}E)(EQ^{-1}))|\le\|W_0^{-1}E\|\,\|EQ^{-1}\|
\le|W_0^{-1}|\,\|E\|\cdot\|E\|\,|Q^{-1}|\le2\|E\|^2 .
\end{equation*}
Altogether,
\begin{equation*}
|\mathfrak{w}(W)-\mathfrak{w}(W_0)|\le\tfrac12\bigl(2\|E\|\,\|W_0-\mathbf{1}\|+2\|E\|^2\bigr)
\le\tfrac12(2\mathsf{k}\tau+2\tau^2)=(\mathsf{k}+\tau)\tau .
\end{equation*}
The last inequality holds provided $\mathsf{k}\ge\|W_0-\mathbf{1}\|$ and $\tau\ge\|E\|$. Both sets of
hypotheses give this. In (a$^\flat$) the two inequalities are assumed. In (a) they follow from
$|W_0-\mathbf{1}|\le\mathsf{k}$ and $|E|\le\tau$ by the first fact, $\|\cdot\|\le|\cdot|$. This proves
(a$^\flat$) and (a).

\emph{Proof of} (b). Since $|W-\mathbf{1}|\le\frac12$, the fourth fact gives $|W^{-1}|\le2$. By
\eqref{4.36:eq-plaquette-action-2} and the third and second facts,
\begin{equation*}
|\mathfrak{w}(W)|=\tfrac12\bigl|\operatorname{tr}\bigl[(W^{-1}(W-\mathbf{1}))(W-\mathbf{1})\bigr]\bigr|
\le\tfrac12\|W^{-1}(W-\mathbf{1})\|\,\|W-\mathbf{1}\|\le\tfrac12|W^{-1}|\,\|W-\mathbf{1}\|^2 .
\end{equation*}
Using $|W^{-1}|\le2$, and then $\|W-\mathbf{1}\|\le|W-\mathbf{1}|$, we obtain the remaining two
inequalities of \eqref{4.36:eq-plaquette-action-4}.
\end{proof}

The smallness hypotheses $|Q-\mathbf{1}|\le\frac12$ in (a$^\flat$) and (a), and
$|W-\mathbf{1}|\le\frac12$ in (b), are in the operator norm, and they cannot be replaced by the
same bounds in $\|\cdot\|$. Two examples show this.

First, take $N=3$ and $W=\operatorname{diag}(1/7,1,1)$. Then
\begin{equation*}
\|W-\mathbf{1}\|^2=\tfrac13\bigl(\tfrac67\bigr)^2=\tfrac{12}{49}<\tfrac14,\qquad
\operatorname{tr}W=\tfrac57,\qquad \operatorname{tr}W^{-1}=3,
\end{equation*}
so
\begin{equation*}
\mathfrak{w}(W)=1-\tfrac12\cdot\tfrac{26}{7}=-\tfrac67,\qquad
|\mathfrak{w}(W)|=\tfrac{42}{49}>\tfrac{12}{49}.
\end{equation*}
Hence (b) with constant one fails under the hypothesis $\|W-\mathbf{1}\|\le\frac12$.

Second, for $N\ge4$, take $W=Q=\operatorname{diag}(\delta,1,\dots,1)$ and $W_0=\mathbf{1}$ with
$0<\delta<1$. Then $\|W-\mathbf{1}\|=(1-\delta)N^{-1/2}\le\frac12$, while
\begin{equation*}
\mathfrak{w}(W)=N^{-1}-\frac{\delta+\delta^{-1}}{2N}\longrightarrow-\infty
\qquad(\delta\downarrow0).
\end{equation*}
Since $\mathfrak{w}(W_0)=0$, no constant multiplying the right-hand sides of (a) or (b) rescues
these bounds under smallness hypotheses in $\|\cdot\|$. The conclusions, on the other hand, may be
taken in $\|\cdot\|^2$, as in (a$^\flat$) and in the middle member of
\eqref{4.36:eq-plaquette-action-4}.

Part (a) is applied to the plaquette variables of a moved configuration, with $W_0$ the
$SU(N)$-valued plaquette variable of the unmoved configuration and $Q$ the ratio of the moved to
the unmoved plaquette variable. Part (b) is applied to the complex plaquette variable
$\mathfrak{U}(\partial q)$ under the bound $|\mathfrak{U}(\partial q)-\mathbf{1}|\le\frac12$, where it
gives $|\mathfrak{w}(\mathfrak{U}(\partial q))|\le\|\mathfrak{U}(\partial q)-\mathbf{1}\|^2$ for every $N$.

\begin{lemma}[The Cauchy bound for a second derivative]\label{4.36:lem-cauchy-bound}
Let $\rho>0$. In this lemma $\zeta$ denotes a complex variable and $g$ a complex-valued function of
$\zeta$, not the coupling. Suppose that $g$ is continuous on the closed disc
$\{\zeta\in\mathbb{C}:|\zeta|\le\rho\}$ and holomorphic in the open disc $\{|\zeta|<\rho\}$. Then
\begin{equation}\label{4.36:eq-cauchy-bound-1}
|g''(0)|\le 2\rho^{-2}\sup_{|\zeta|=\rho}|g(\zeta)-g(0)| .
\end{equation}
\end{lemma}

\begin{proof}
Put $h(\zeta):=g(\zeta)-g(0)$. Then $h$ is continuous on the closed disc of radius $\rho$ and
holomorphic in the open disc, and $h''(0)=g''(0)$, because $h$ and $g$ differ by a constant.

Fix $\rho'$ with $0<\rho'<\rho$. The circle $\{|\zeta|=\rho'\}$ lies in the open disc, where $h$ is
holomorphic, so Cauchy's integral formula for the second derivative gives
\begin{equation*}
h''(0)=\frac{2}{2\pi i}\oint_{|\zeta|=\rho'}\frac{h(\zeta)}{\zeta^{3}}\,d\zeta .
\end{equation*}
The circle has length $2\pi\rho'$ and $|\zeta^{3}|=\rho'^{3}$ on it; hence
\begin{equation*}
|h''(0)|\le\frac{2}{2\pi}\cdot 2\pi\rho'\cdot\frac{1}{\rho'^{3}}\sup_{|\zeta|=\rho'}|h(\zeta)|
=2\rho'^{-2}\sup_{|\zeta|=\rho'}|h(\zeta)| .
\end{equation*}
Since $h$ is continuous on the closed disc of radius $\rho$ and holomorphic in its interior, the
maximum principle shows that $|h|$ attains its maximum over this closed disc on the boundary circle
$\{|\zeta|=\rho\}$. The circle $\{|\zeta|=\rho'\}$ is contained in the closed disc, so
\begin{equation*}
\sup_{|\zeta|=\rho'}|h(\zeta)|\le\sup_{|\zeta|=\rho}|h(\zeta)| .
\end{equation*}
Combining the last two estimates,
\begin{equation*}
|g''(0)|=|h''(0)|\le 2\rho'^{-2}\sup_{|\zeta|=\rho}|g(\zeta)-g(0)|
\end{equation*}
for every $\rho'\in(0,\rho)$. The supremum on the right does not depend on $\rho'$ and is finite,
because $h$ is continuous on the compact circle $\{|\zeta|=\rho\}$. Letting $\rho'$ increase to
$\rho$ gives \eqref{4.36:eq-cauchy-bound-1}.
\end{proof}

\begin{lemma}[Short broken paths stay in the small-field window]\label{4.36:lem-broken-paths}
Let $G=SU(N)$, let $\mathfrak{g}$ be the traceless Hermitian $N\times N$ matrices, let $|\cdot|$ be
the operator norm, and let $\eta=L^{-k}$ and $\varepsilon_k=g_k\,p_0(g_k)$, with $p_0(\cdot)$ the
degree function, be as in
Notation~\ref{4.36:not-flow-thresholds}. Let $h$ be a clean term, in the setting of
Definition~\ref{4.36:def-unit-directions}, and let $\mathbf{1}_{G(h)}$ be the indicator of its
good set (the letter $G$ in $\mathbf{1}_{G(h)}$ names the good set of $h$, not the group). Let $p$
be the plaquette and $b_1\subset\partial p$ the marked bond of that setting. We
write a configuration as $V=(u,V_{\mathrm{off}})$, where $u=V(b_1)$ is the free link and
$V_{\mathrm{off}}$ is the collection of the remaining links. For a cube $\square$ of the set
$\mathcal{N}$ of cubes near $p$, let $\mathfrak{d}_{\square}(V)$ be the data of the local
variational problem \cite[(2.16)]{B-CMP119} associated with $\square$ and posed on an enlargement
of $\square$. These data are a finite product map of the links of $V$ on that enlargement. Let
$U_{k,\square}(V)$ be the minimiser of that problem
given by \cite[Theorem~1]{B-CMP102b}. The letter $p''$ ranges over the fine plaquettes contained in
$\tilde{\square}$, the enlargement of $\square$ over which the window is read; the local problem
itself is posed on a larger enlargement of $\square$. We recall the window characteristic
functions, the window of the free link and the deep set:
\begin{gather*}
\chi_{\square}(V)=\begin{cases}
 1&\text{if }\sup_{p''\subset\tilde{\square}}
 |U_{k,\square}(V,\partial p'')-\mathbf{1}|<\varepsilon_k\eta^2,\\
 0&\text{otherwise,}
 \end{cases}\\
W(V_{\mathrm{off}})=\Bigl\{u\in G:\ \prod_{\square\in\mathcal{N}}\chi_{\square}(u,V_{\mathrm{off}})=1\Bigr\},\\
\mathcal{D}=\Bigl\{V:\ \sup_{\square\in\mathcal{N},\,p''\subset\tilde{\square}}
 |U_{k,\square}(V,\partial p'')-\mathbf{1}|\le\tfrac12\varepsilon_k\eta^2\Bigr\}.
\end{gather*}
We recall that for $u^0\in G$ and $r>0$ the operator ball is
$B^{\mathrm{op}}_r(u^0)=\{u^0e^{iv}:\ v\in\mathfrak{g},\ |v|<r\}$. Let $a_1'$ be the constant of
\cite[Theorem~1]{B-CMP102b}, and let $C_1$ be Ba{\l}aban's constant of the hypothesis
\cite[(172)]{B-CMP102b} (in this lemma $C_1$ denotes this constant only, not the constant of the
radius $\alpha_{1,j}$ of
Notation~\ref{4.36:not-flow-thresholds}). Let $r_0$ be the radius entering the thresholds
\eqref{4.36:eq-flow-thresholds-b} (not the depth offset $r_0$ of the witness construction,
Notation~\ref{0.2:not-glossary}). The thresholds
give in particular $r_0\le\tfrac12$ and $r_0\le C_1a_1'/4$.

Assume the following three hypotheses.
\begin{itemize}
\item[$(\mathrm{D\text{-}i})$] If $\mathbf{1}_{G(h)}(V)>0$, then $V\in\mathcal{D}$.
\item[$(\mathrm{P}9_{\mathrm{loc}})$] The free link enters each product map
$\mathfrak{d}_{\square}$ in at most $C'(L)$ items, $C'(L)$ a constant depending on $L$ only. Each
such item is one factor $d=A\,u\,B'$, where
$A$ and $B'$ are words in $V_{\mathrm{off}}$. Hence, for $v\in\mathfrak{g}$,
\begin{equation*}
d(ue^{iv})=d(u)\,e^{iv_d},\qquad v_d:=\operatorname{Ad}(B'^{-1})v,\qquad |v_d|=|v|.
\end{equation*}
Let $C_{\mathrm{loc}}\ge1$ be the constant entering the thresholds
\eqref{4.36:eq-flow-thresholds-b}; it depends on the dimension and on $L$ only. The hypothesis
asserts the following. Suppose that the data
$\mathfrak{d}_{\square}((u',V_{\mathrm{off}}))$ satisfy \cite[(7)]{B-CMP102b} with
$\varepsilon_1=a_1'$, a constant independent of the coupling, and suppose that
$v'\in\mathfrak{g}$ satisfies $|v'|\le r_0$. Then
\begin{equation}\label{4.36:eq-broken-paths-a}
\begin{split}
\sup_{p''\subset\tilde{\square}}
 |U_{k,\square}((u'e^{iv'},V_{\mathrm{off}}),\partial p'')-\mathbf{1}|
 &\le\sup_{p''\subset\tilde{\square}}
 |U_{k,\square}((u',V_{\mathrm{off}}),\partial p'')-\mathbf{1}|\\
 &\quad+C_{\mathrm{loc}}|v'|\eta^2 .
\end{split}
\end{equation}
\item[$(7_{\mathrm{loc}})$] At the base $(u^0,V_{\mathrm{off}})$ the data
$\mathfrak{d}_{\square}(u^0,V_{\mathrm{off}})$, $\square\in\mathcal{N}$, satisfy
\cite[(7)]{B-CMP102b} with $\varepsilon_1=a_1'$. The same holds for all data obtained from them by
multiplying the items containing the free link by unitary factors at operator distance at most
$\varepsilon_k$ from $\mathbf{1}$.
\end{itemize}
Fix $V_{\mathrm{off}}$, and let $(u^0,V_{\mathrm{off}})\in\mathcal{D}$ satisfy
$(7_{\mathrm{loc}})$. Put $\varrho:=\varepsilon_k/(4C_{\mathrm{loc}})$. Then
$\varrho\le\varepsilon_k/4$, and the following hold.
\begin{itemize}
\item[(a)] Let $u=u^0e^{iv_1}e^{iv_2}\cdots e^{iv_s}$ with $v_i\in\mathfrak{g}$ and
$\sigma:=\sum_{i\le s}|v_i|\le\varrho$. Then for every $\square\in\mathcal{N}$
\begin{equation}\label{4.36:eq-broken-paths-1}
\sup_{p''\subset\tilde{\square}}|U_{k,\square}((u,V_{\mathrm{off}}),\partial p'')-\mathbf{1}|
\le\bigl(\tfrac12\varepsilon_k+C_{\mathrm{loc}}\sigma\bigr)\eta^2
\le\tfrac34\varepsilon_k\eta^2<\varepsilon_k\eta^2 .
\end{equation}
In particular $u\in W(V_{\mathrm{off}})$. The same holds for every partial product
$u^0e^{iv_1}\cdots e^{iv_j}$, $j\le s$.
\item[(b)] $B^{\mathrm{op}}_{\varrho}(u^0)\subset W(V_{\mathrm{off}})$.
\item[(c)] Assume the thresholds (t1) and (t2) of \eqref{4.36:eq-flow-thresholds-b}. Threshold
(t1) contains the member $\varrho\ge2g_k$, equivalently $p_0(g_k)\ge8C_{\mathrm{loc}}$. Threshold
(t2) contains the members $R\le\varepsilon_k/(16C_{\mathrm{loc}})=\varrho/4$ and $2R\le r_0$. Here
$R$ is the radius of the one-link moves of Definition~\ref{4.36:def-unit-directions}, the quantity
bounded in threshold (t2); it is not the ball radius $R$ of the observables. Let
$u\in B^{\mathrm{op}}_{g_k}(u^0)$, let
$\hat{T}\in\mathfrak{g}$ with $|\hat{T}|=1$, and let $t$ be real with $|t|\le2R$. Then
$ue^{it\hat{T}}\in W(V_{\mathrm{off}})$. Moreover, for every $\square\in\mathcal{N}$, the supremum
in \eqref{4.36:eq-broken-paths-1} taken at $(ue^{it\hat{T}},V_{\mathrm{off}})$ is at most
$\tfrac34\varepsilon_k\eta^2$.
\end{itemize}
\end{lemma}

\begin{proof}
Note that $\varrho\le\varepsilon_k/4$ because $C_{\mathrm{loc}}\ge1$. We use two elementary facts
about unitary matrices. (i) If $A$ and $B$ are unitary, then
\begin{equation*}
AB-\mathbf{1}=(A-\mathbf{1})B+(B-\mathbf{1}),
\end{equation*}
and therefore $|AB-\mathbf{1}|\le|A-\mathbf{1}|+|B-\mathbf{1}|$. (ii) For $X\in\mathfrak{g}$,
item (F6) of Lemma~\ref{4.36:lem-eigenangle-inequalities} gives $|e^{iX}-\mathbf{1}|\le|X|$.

(a) \emph{Refinement.} Suppose $|v_i|>r_0$ for some $i$. Put $m_i:=\lceil|v_i|/r_0\rceil$. Then
$|v_i|/m_i\le r_0$, and $e^{iv_i}=(e^{iv_i/m_i})^{m_i}$ as matrices. Replace the factor
$e^{iv_i}$ by the $m_i$ factors $e^{iv_i/m_i}$. This leaves the endpoint $u$ unchanged. It also
leaves unchanged the original partial products, which occur among the new ones. The total length
is preserved, since $m_i\cdot|v_i|/m_i=|v_i|$, so $\sigma$ is unchanged. Every partial sum of the
refined lengths is at most $\sigma$. It therefore suffices to treat paths with $|v_i|\le r_0$ for
all $i$ and to prove, for each partial product, the bound with $\sigma$ replaced by the
corresponding partial sum, which is at most $\sigma$; this is stronger. We assume this from now on.

\emph{Induction on the number of factors.} For $0\le j\le s$ put
$u_j:=u^0e^{iv_1}\cdots e^{iv_j}$ and $\sigma_j:=\sum_{i\le j}|v_i|$. Let $P(j)$ be the
statement that for every $\square\in\mathcal{N}$
\begin{equation}\label{4.36:eq-broken-paths-2}
\sup_{p''\subset\tilde{\square}}|U_{k,\square}((u_j,V_{\mathrm{off}}),\partial p'')-\mathbf{1}|
\le\bigl(\tfrac12\varepsilon_k+C_{\mathrm{loc}}\sigma_j\bigr)\eta^2 .
\end{equation}
Since $u_0=u^0$ and $\sigma_0=0$, the statement $P(0)$ is the hypothesis
$(u^0,V_{\mathrm{off}})\in\mathcal{D}$, by the definition of $\mathcal{D}$.

Let $1\le j\le s$ and assume $P(j-1)$. Put $u':=u_{j-1}$ and
$E:=e^{iv_1}\cdots e^{iv_{j-1}}$, so that $u'=u^0E$. Fix $\square\in\mathcal{N}$.

If the free link does not enter the data $\mathfrak{d}_{\square}$, then these data, and with them
the minimiser $U_{k,\square}$, do not depend on the free link. The left side of
\eqref{4.36:eq-broken-paths-2} is then the same at $u_j$ as at $u^0$, hence at most
$\tfrac12\varepsilon_k\eta^2$, and $P(j)$ holds for this $\square$.

Otherwise, by $(\mathrm{P}9_{\mathrm{loc}})$ the free link enters at most $C'(L)$ items of
$\mathfrak{d}_{\square}$, each in the form $d=A\,u\,B'$. For each such item,
\begin{equation*}
d(u^0E)=A\,u^0B'\cdot B'^{-1}EB'=d(u^0)\cdot\operatorname{Ad}(B'^{-1})(E).
\end{equation*}
So the data $\mathfrak{d}_{\square}(u',V_{\mathrm{off}})$ differ from
$\mathfrak{d}_{\square}(u^0,V_{\mathrm{off}})$ only in these items. Each of them is multiplied on
the right by the unitary factor $\operatorname{Ad}(B'^{-1})(E)$. The word $B'$ is unitary, so
conjugation by it preserves the operator norm. Applying the telescoping inequality repeatedly (for
$j=1$ there is nothing to do, since $E=\mathbf{1}$) and then (F6) to each factor, we get
\begin{equation*}
|\operatorname{Ad}(B'^{-1})(E)-\mathbf{1}|=|E-\mathbf{1}|
\le\sum_{i<j}|e^{iv_i}-\mathbf{1}|\le\sum_{i<j}|v_i|\le\sigma\le\varrho\le\varepsilon_k .
\end{equation*}
By $(7_{\mathrm{loc}})$, the data $\mathfrak{d}_{\square}(u',V_{\mathrm{off}})$ therefore satisfy
\cite[(7)]{B-CMP102b} with $\varepsilon_1=a_1'$. The hypotheses of
\cite[Proposition~9]{B-CMP102b} that enter through $(\mathrm{P}9_{\mathrm{loc}})$ bear on two
things only: on the base data, through \cite[(7)]{B-CMP102b}, and on the increment, through
\cite[(172)]{B-CMP102b}. The regularity of the base minimiser $U_{k,\square}((u',V_{\mathrm{off}}))$
is a conclusion of \cite[Theorem~1]{B-CMP102b}, namely the bound \cite[(8)]{B-CMP102b}. It is not
an additional hypothesis.

The move $u'\mapsto u'e^{iv_j}=u_j$ multiplies each item containing the free link by
$e^{iv'_d}$, where $v'_d=\operatorname{Ad}(B'^{-1})v_j$. This is a unitary conjugate of
$e^{iv_j}$, and $|v'_d|=|v_j|\le r_0$. So hypothesis $(\mathrm{P}9_{\mathrm{loc}})$ applies at the
base $(u',V_{\mathrm{off}})$ with $v'=v_j$. Inequality \eqref{4.36:eq-broken-paths-a} and $P(j-1)$
give
\begin{align*}
\sup_{p''\subset\tilde{\square}}|U_{k,\square}((u_j,V_{\mathrm{off}}),\partial p'')-\mathbf{1}|
&\le\sup_{p''\subset\tilde{\square}}|U_{k,\square}((u',V_{\mathrm{off}}),\partial p'')-\mathbf{1}|
 +C_{\mathrm{loc}}|v_j|\eta^2\\
&\le\bigl(\tfrac12\varepsilon_k+C_{\mathrm{loc}}\sigma_{j-1}\bigr)\eta^2
 +C_{\mathrm{loc}}|v_j|\eta^2
=\bigl(\tfrac12\varepsilon_k+C_{\mathrm{loc}}\sigma_j\bigr)\eta^2 .
\end{align*}
This is $P(j)$ for this $\square$. Since $\square\in\mathcal{N}$ was arbitrary, $P(j)$ holds, and
by induction $P(j)$ holds for $0\le j\le s$.

Finally, $\sigma_j\le\sigma$ and $\sigma\le\varrho$, and $\varrho=\varepsilon_k/(4C_{\mathrm{loc}})$;
hence $C_{\mathrm{loc}}\sigma_j\le\varepsilon_k/4$. Since $\varepsilon_k\eta^2>0$, it follows that
\begin{equation*}
\bigl(\tfrac12\varepsilon_k+C_{\mathrm{loc}}\sigma_j\bigr)\eta^2
\le\tfrac34\varepsilon_k\eta^2<\varepsilon_k\eta^2 .
\end{equation*}
For $j=s$ this is \eqref{4.36:eq-broken-paths-1}. Moreover, for every $j$ and every
$\square\in\mathcal{N}$ we get $\chi_{\square}(u_j,V_{\mathrm{off}})=1$, so the product over
$\mathcal{N}$ equals $1$. Thus $u_j\in W(V_{\mathrm{off}})$ for every partial product, and in
particular $u=u_s\in W(V_{\mathrm{off}})$.

(b) An element of $B^{\mathrm{op}}_{\varrho}(u^0)$ has the form $u^0e^{iv}$ with $v\in\mathfrak{g}$
and $|v|<\varrho$. This is the case $s=1$, $v_1=v$, $\sigma=|v|\le\varrho$ of (a). Hence it lies
in $W(V_{\mathrm{off}})$.

(c) By $\varepsilon_k=g_k\,p_0(g_k)$ and (t1), $p_0(g_k)\ge8C_{\mathrm{loc}}$. Hence
\begin{equation*}
g_k\le\frac{g_k\,p_0(g_k)}{8C_{\mathrm{loc}}}=\frac{\varepsilon_k}{8C_{\mathrm{loc}}}
=\frac{\varrho}{2}.
\end{equation*}
By (t2), $2R\le\varrho/2$. Write $u=u^0e^{iv}$ with $v\in\mathfrak{g}$ and $|v|<g_k$. Then
$ue^{it\hat{T}}=u^0e^{iv}e^{it\hat{T}}$, with $t\hat{T}\in\mathfrak{g}$ and $|t\hat{T}|=|t|$. This
is a broken path with two legs, $v$ and $t\hat{T}$, of total operator length
\begin{equation*}
\sigma=|v|+|t|<g_k+2R\le\frac{\varrho}{2}+\frac{\varrho}{2}=\varrho .
\end{equation*}
If moreover (t8)$^{+}$ of \eqref{4.36:eq-flow-thresholds-b} holds, both legs have length at most
$r_0$: $|v|<g_k$ and $g_k<r_0$ by (t8)$^{+}$, and $|t|\le2R$, which is at most $r_0$ by the member
$2R\le r_0$ of (t2); so no refinement is needed. In any case part (a) applies, since its
refinement step removes every restriction on the lengths of the legs. Part (a) with $s=2$, $v_1=v$
and $v_2=t\hat{T}$ gives $ue^{it\hat{T}}\in W(V_{\mathrm{off}})$. It also gives, for every
$\square\in\mathcal{N}$, that the supremum in \eqref{4.36:eq-broken-paths-1} at
$ue^{it\hat{T}}$ is at most $\tfrac34\varepsilon_k\eta^2$.
\end{proof}

\begin{remark}[On the hypotheses and the form of Lemma~\ref{4.36:lem-broken-paths}]
Hypothesis $(\mathrm{D\text{-}i})$ is obtained from \cite[(3.6)--(3.8)]{B-CMP119} together with
the threshold (t6)
of \eqref{4.36:eq-flow-thresholds-b}. The lemma is applied at bases with
$\mathbf{1}_{G(h)}>0$. By $(\mathrm{D\text{-}i})$, these bases lie in $\mathcal{D}$, and this is
the hypothesis
$(u^0,V_{\mathrm{off}})\in\mathcal{D}$ of the lemma.

Hypothesis $(\mathrm{P}9_{\mathrm{loc}})$ is the form in which \cite[Proposition~9]{B-CMP102b},
with its hypothesis \cite[(172)]{B-CMP102b} and its bound \cite[(190)]{B-CMP102b}, is used for the
local problem \cite[(2.16)]{B-CMP119}. It is combined with the plaquette formula for that problem,
which expresses its plaquette variables through the Baker--Campbell--Hausdorff series in $|\cdot|$.
The steps are these.
\begin{itemize}
\item The increment satisfies \cite[(172)]{B-CMP102b}. By (F6) and the threshold (t8),
\begin{equation*}
|e^{iv'_d}-\mathbf{1}|\le|v'|\le r_0\le C_1a_1'/4<C_1a_1'.
\end{equation*}
Moreover $v'_d=\tfrac{1}{i}\log e^{iv'_d}$ with the principal branch of the logarithm, since
$|v'_d|\le r_0$ and $r_0\le\tfrac12$, so that $|v'_d|<\pi$.
\item \cite[Proposition~9]{B-CMP102b} then represents
$U_{k,\square}((u'e^{iv'},V_{\mathrm{off}}),\cdot)$ as $[e^{i\eta\mathcal{H}}U_0]^{\mathrm{g}}$.
Here $U_0$ is the base minimiser, $\mathrm{g}$ is a $G$-valued gauge transformation, and
$\mathcal{H}$ is analytic in $v'$, vanishes at $v'=0$, and obeys \cite[(190)]{B-CMP102b} in
$|\cdot|$.
\item The plaquette formula uses the Baker--Campbell--Hausdorff series in $|\cdot|$, with
constants independent of $N$. It applies under
\begin{equation*}
\eta|\mathcal{H}|\le C_H|v'|\le C_Hr_0\le\tfrac14,
\end{equation*}
where $C_H$ denotes the constant in the bound $\eta|\mathcal{H}|\le C_H|v'|$, and the last
inequality is the member
$r_0\le(4C_H)^{-1}$ of (t8). Together with the bound \cite[(8)]{B-CMP102b} for $U_0$ from
\cite[Theorem~1]{B-CMP102b}, it gives, for every fine plaquette $p''\subset\tilde{\square}$,
\begin{equation*}
|[e^{i\eta\mathcal{H}}U_0](\partial p'')\,U_0(\partial p'')^{-1}-\mathbf{1}|
\le C_{\mathrm{loc}}|v'|\eta^2 .
\end{equation*}
The constant $C_{\mathrm{loc}}$ is enlarged to be at least $1$ if necessary. This only
strengthens (t1) and lowers the upper bound on $R$ in (t2).
\item Inequality \eqref{4.36:eq-broken-paths-a} then follows from two facts. The quantity
$|U(\partial p'')-\mathbf{1}|$ is gauge invariant. For unitaries $\tilde{U}$ and $U_0$, the
telescoping inequality gives
\begin{equation*}
|\tilde{U}U_0^{-1}\cdot U_0-\mathbf{1}|\le|\tilde{U}U_0^{-1}-\mathbf{1}|+|U_0-\mathbf{1}|.
\end{equation*}
\end{itemize}

Hypothesis $(7_{\mathrm{loc}})$ requires that Ba{\l}aban's condition \cite[(7)]{B-CMP102b} on the
data be kept along the moves of the free link.

Finally, membership in $W(V_{\mathrm{off}})$ is a property of the endpoint configuration alone.
Part (a) establishes it directly along broken paths. For this reason no triangle inequality for a
distance on $SU(N)$ built from the chart $\Lambda_G$ is used in part (c). The same holds for
parts (a) and (b).
\end{remark}

\begin{remark}[Where the direction's normalisation enters the curvature inputs]
\label{4.36:rem-normalisation-inputs}
Let $G=SU(N)$ and work in the setting of Definition~\ref{4.36:def-unit-directions}: $K$, $k$, a
clean term $h$ with conditional exponent $F_h=F^{\mathrm{near}}+F^{\mathrm{far}}$, and a real windowed
pair $(V,w)$ whose free bond $b_1$ carries the link variable $u$. Write
$\mathsf{n}=(N-1)^{1/2}$. Let $\hat{T}\in\mathfrak{g}$ be an operator-unit direction,
$|\hat{T}|=1$, and $T\in\mathfrak{g}$ a $\mathfrak{g}$-unit direction, $|T|_{\mathfrak{g}}=1$. Call a
bound \emph{scalar} if it bounds a complex-valued function on a domain defined by conditions in the
operator norm $|\cdot|$. The inputs of the bounds on $F^{\mathrm{near}}$ and $F^{\mathrm{far}}$
along these directions are exactly those taken one by one in the proof below; for each of them the
following holds.
\begin{itemize}
\item[(a)] It is stated in $|\cdot|$, or it is a scalar bound, or it involves no group variable at all,
as for the coefficients of the plaquette terms. A scalar bound depends on the direction
only through the radius of the disc in $\zeta$, or of the data ball, on which it holds, and this radius
is where the rescaling acts.
\item[(b)] Along $\hat{T}$ its use produces no factor from the normalisation of the direction. In
particular, the one-link factor $V'=e^{i\zeta\hat{T}'}$, $\hat{T}'=u\hat{T}u^{-1}$, satisfies
$|V'-\mathbf{1}|\le e^{|\zeta|}-1$ in either of the norms $|\cdot|$, $\|\cdot\|$. Hence no norm on the
input side costs a power of $\mathsf{n}$.
\item[(c)] Along $T$ its use costs the power of $\mathsf{n}$ stated for it in the proof, which is at most
$\mathsf{n}^2$. The cost is none for the coefficients of the plaquette terms and at most $\mathsf{n}$ for
the inputs linear in the direction. Equivalently, the radius on which the input holds shrinks by the
factor $|T|\le\mathsf{n}$. The window of Lemma~\ref{4.36:lem-broken-paths} and the one-bond move bound
for the second-class regions are used only at operator reach below $g_k$, where they cover both kinds
of direction with no power of $\mathsf{n}$.
\end{itemize}
Here \emph{no power of $\mathsf{n}$} means no factor coming from the normalisation of the direction. The
constants involved, such as $C_q(N)$ below, depend on $N$ through the values of their own inputs.
\end{remark}

\begin{proof}
\emph{How the direction enters.} By Definition~\ref{4.36:def-unit-directions}, for $u\in G$,
$S\in\mathfrak{g}$ and $\zeta\in\mathbb{C}$ we have $ue^{i\zeta S}=e^{i\zeta S'}u$ with
$S'=uSu^{-1}\in\mathfrak{g}$. Moreover $|S'|=|S|$ and $|S'|_{\mathfrak{g}}=|S|_{\mathfrak{g}}$, since
$u$ is unitary. The moved configuration is $V'\cdot(u,V_{\mathrm{off}})$, where $V'=e^{i\zeta S'}$ sits
at $b_1$ and $\mathbf{1}$ sits elsewhere.

Summing the exponential series term by term gives
\begin{equation}\label{4.36:eq-normalisation-inputs-1}
|V'-\mathbf{1}|\le\sum_{j\ge1}\frac{|\zeta|^j|S'|^j}{j!}=e^{|\zeta||S|}-1 .
\end{equation}
This is the inequality (F6) of \eqref{4.36:eq-eigenangle-inequalities-a}. Since the normalised trace
norm is bounded by the operator norm, $\|\cdot\|\le|\cdot|$, the bound
\eqref{4.36:eq-normalisation-inputs-1}
holds in both norms. For $S=\hat{T}$ it reads $|V'-\mathbf{1}|\le e^{|\zeta|}-1$, which proves the second
sentence of (b).

For $|\operatorname{Re}\zeta|\,|S|<\pi$ the principal logarithm gives $B=(1/i)\log V'=\zeta S'$ for the
quantity $B$ of Ba{\l}aban's continuation: the matrix $V'$ is normal with eigenvalues $e^{i\zeta\mu_l}$,
where $\mu_1,\dots,\mu_N$ are the eigenvalues of $S'$, and their arguments $(\operatorname{Re}\zeta)\mu_l$
lie in $]-\pi,\pi[$ because $|\mu_l|\le|S'|=|S|$. This field is traceless, since $S'\in\mathfrak{g}$. Its
derivative is $\partial_\zeta B(b_1)=S'$, with $|\partial_\zeta B(b_1)|=|S|$. In particular this applies
on the discs $|\zeta|\,|S|<2r_0$ of the items below, where $2r_0\le1$.

Next compare the two kinds of direction. By Lemma~\ref{4.36:lem-derivative-rescaling}, with
$\hat{T}=T/|T|$,
\[
F(ue^{i\zeta T})=F(ue^{i(\zeta|T|)\hat{T}})
\]
for every function $F$ near $u$. Hence a bound valid for $|\zeta|<\rho$ along $\hat{T}$ is valid for
$|\zeta||T|<\rho$ along $T$, that is, on a disc of radius $\rho/|T|$. Derivatives of order one and two
along $T$ are $|T|$ and $|T|^2$ times those along $\hat{T}$. For $|T|_{\mathfrak{g}}=1$ we have
$|T|^2\le N-1=\mathsf{n}^2$ by (F1) of \eqref{4.36:eq-eigenangle-inequalities-a}, as recorded in
Lemma~\ref{4.36:lem-derivative-rescaling}. So every radius shrinks by at most $\mathsf{n}$, every first
variation grows by at most $\mathsf{n}$, and every second derivative obtained by a Cauchy bound at such
a radius grows by at most $\mathsf{n}^2$.

We now take the inputs one by one. In each item $r_0$ and $R$ are the radii fixed in the thresholds
$(\mathrm{t}8)$ and $(\mathrm{t}2)$ of \eqref{4.36:eq-flow-thresholds-b}; they are the radii of the discs
$|\zeta|<2r_0$ and $|\zeta|<2R$ below, distinct from the depth offset $r_0$ of the witness
construction and from the ball radius $R$ of the observables. Throughout, $\eta$ is the lattice
spacing, $e_y$ is the decay weight at the point $y$, $n=n(y)$ is the stratum of the point $y$ at which
a bound is read, and $\ell_n=L^n\eta$ is the unit of that scale.

\emph{The one-link complexification.} By \cite[Proposition~9]{B-CMP102b}, the background field is an
analytic function of small configurations $V'$ with values in the complexification
$G^c=SL(N,\mathbb{C})$ of $G$, with $B=(1/i)\log V'$, under the hypothesis
\cite[(172)]{B-CMP102b}, $|V'-\mathbf{1}|<C_1\varepsilon_1$, with Ba{\l}aban's constants $C_1$ and
$\varepsilon_1$ of that hypothesis. This hypothesis is stated in $|\cdot|$, in
the convention of \cite[(19)]{B-CMP98}. Here $V'=e^{i\zeta S'}$, and \eqref{4.36:eq-normalisation-inputs-1}
bounds $|V'-\mathbf{1}|$ by $e^{|\zeta||S|}-1$. On the discs used here the field $B=\zeta S'$ is
traceless.

Let $S=\hat{T}$ and $|\zeta|<2r_0$, where $2r_0\le\min\{1,C_1a_1'/2\}$ and $a_1'$ is the constant of the
small-field bounds that also enters the bound on $\kappa_q$ in the plaquette item below. The function
$s\mapsto e^s-1$ is
convex and vanishes at $0$, and $|\zeta|\le 1$, so $e^{|\zeta|}-1\le(e-1)|\zeta|$. Since $e-1<2$ and
$|\zeta|<C_1a_1'/2$, we obtain
\begin{equation}\label{4.36:eq-normalisation-inputs-2}
|V'-\mathbf{1}|\le e^{|\zeta|}-1\le(e-1)|\zeta|<C_1a_1' .
\end{equation}
Cost along $\hat{T}$: none; the condition is $|\zeta|<2r_0$, exactly as in $(\mathrm{t}8)$ of
\eqref{4.36:eq-flow-thresholds-b}. Cost along $T$: the
disc is $|\zeta|<2r_0/|T|$, of radius at least $2r_0/\mathsf{n}$.

\emph{The first variation of the background field.} Let $\mathcal{H}_\zeta$ be the field of
\cite[Proposition~9]{B-CMP102b} at $B=\zeta S'$. It is $\mathfrak{g}^c$-valued because $S'$ is traceless.
Contract the bound \cite[(190), p.~308]{B-CMP102b} with $\partial_\zeta B(b_1)=S'$. For
$\iota\in\{0,1\}$ this gives
\begin{equation}\label{4.36:eq-normalisation-inputs-3}
\ell_n^{1+\iota}\,|\nabla^\iota\partial_\zeta\mathcal{H}_\zeta(y)|\le C_H\,e_y\,|S'| ,
\end{equation}
with a constant $C_H$ supplied by that bound.
The members of that bound involving $D^*D$, the Laplacian and the H\"older norm are treated in the same
way, and the same bound holds for $\mathcal{H}_\zeta/\zeta$.

The bound \cite[(190)]{B-CMP102b} is a bound, in the operator norm of \cite[(19)--(20)]{B-CMP98}, for
the linear map $D_{B(y')}\mathcal{H}$, the derivative of $\mathcal{H}$ with respect to the value of $B$
at the bond $y'$. Its data are supplied in $|\cdot|$ by
\cite[(3.6)--(3.8)]{B-CMP119}, and the field is bounded with no factor of the dimension $D=N^2-1$ of
$\mathfrak{g}$.

Cost along $\hat{T}$: none, since $|S'|=1$ for $S=\hat{T}$. Cost along $T$: for $S=T$ one has
$|S'|=|T|\le\mathsf{n}$, so the right side of \eqref{4.36:eq-normalisation-inputs-3} is multiplied by
$|T|\le\mathsf{n}$, or equivalently the radius is divided by $\mathsf{n}$.

\emph{The plaquette formula.} The Baker--Campbell--Hausdorff formula is applied in $|\cdot|$ to a
product of four exponentials whose exponents have operator norm at most $\tfrac14$. Its constants do
not depend on $N$ in $|\cdot|$, by
\cite[(20)]{B-CMP98}. Together with \cite[(8)]{B-CMP102b} and the plaquette formula of
\cite{B-CMP119}, it gives the increment size
\[
\theta_q=C_P\,\eta^2\ell_n^{-2}e_y ,
\]
where, under the threshold $(\mathrm{t}8)$ above, $C_P$ is a numerical multiple of $C_H$. Put
$\kappa_q=|U_k(V)(\partial q)-\mathbf{1}|$.
This quantity is computed at the real minimiser and does not depend on the direction, and
\[
\kappa_q\le B_3a_1'\eta^2\ell_n^{-2},
\]
with the constant $B_3$ of the small-field bounds.
Cost along $\hat{T}$: none. Cost along $T$: a factor $\mathsf{n}$.

\emph{The window.} The window characteristic functions $\chi_{\square}$ are defined in $|\cdot|$, in
Ba{\l}aban's construction of the small-field cubes.
Lemma~\ref{4.36:lem-broken-paths} takes the inputs $(\mathrm{P}9_{\mathrm{loc}})$ and
$(7_{\mathrm{loc}})$, and its window ball is the operator ball $B^{\mathrm{op}}_{\varrho}(u^0)$ of
radius $\varrho$ about the base point $u^0$ of that lemma.

Cost along $\hat{T}$: none; the lemma is used under $(\mathrm{t}1)$ of \eqref{4.36:eq-flow-thresholds-b}
exactly as stated.

Cost along $T$: none. Since $|v|\le\mathsf{n}|v|_{\mathfrak{g}}$ by (F1) of
\eqref{4.36:eq-eigenangle-inequalities-a}, we have
$B^{\mathfrak{g}}_{g_k/\mathsf{n}}(u^0)\subset B^{\mathrm{op}}_{g_k}(u^0)$. By contrast, the ball
$B^{\mathfrak{g}}_{g_k}(u^0)$ is only contained in $B^{\mathrm{op}}_{\mathsf{n}g_k}(u^0)$. Working on it
would require the condition $p_0(g_k)\ge8C_{\mathrm{loc}}\mathsf{n}$ on the degree function
$p_0(\cdot)$, with $C_{\mathrm{loc}}$ the constant
of Lemma~\ref{4.36:lem-broken-paths}, which is not among the thresholds and is not used.

\emph{Membership in Ba{\l}aban's domains, and the norm $\|\cdot\|_{\star}$.} Put
$\mathfrak{h}=\chi_{\mathsf{N}}\mathcal{H}_\zeta$, where $\chi_{\mathsf{N}}$ is the smooth cutoff to the
bond set $\mathsf{N}$ described in the cutoff item below. The configuration $e^{i\eta\mathfrak{h}}U_k(V)$
lies in
Ba{\l}aban's complex domains $U^c_j(X,\alpha_{0,j},\alpha_{1,j})$ (of radii $\alpha_{0,j}$,
$\alpha_{1,j}$) and $\tilde{U}_j$. At half radius it also lies in the
domain of the propagator bounds \cite[Theorems~3.4, 3.10]{B-CMP99b}.

The defining clauses \cite[(1.11)--(1.16)]{B-CMP109} and \cite[(2.34)--(2.39)]{B-CMP119} are smallness
conditions in $|\cdot|$ on $\mathfrak{g}^c$-valued fields. The configuration norm $\|\cdot\|_{\star}$ is
likewise built on $|\cdot|$ of $\mathfrak{g}^c$-valued fields. The field $\mathfrak{h}$ inherits its bounds
from \eqref{4.36:eq-normalisation-inputs-3}, and $\mathcal{H}_\zeta\in\mathfrak{g}^c$ because $S'$ is
traceless. The remaining conditions entering the membership concern only the coupling flow.

Cost along $\hat{T}$: none. Cost along $T$: the radius is divided by $\mathsf{n}$.

\emph{The cutoff bound.} The smooth cutoffs $\chi_{\mathsf{N}}$ are those of \cite{B-CMP119}, and
\[
\|\chi_{\mathsf{N}}\mathcal{H}_\zeta\|_{\star}
\le C'_{\chi}C_H|\zeta|\min\{1,\hat{e}_{\mathsf{N}^{+}}\}/\delta_k ,
\]
with the constant $C'_{\chi}$ of the cutoff construction and the level-$k$ decay constant $\delta_k$ of
that construction.
A coupling ceiling ensures $e^{\delta_0R_*/16}\ge2(1+B)^{1/2}$, a condition free of $N$ under
\eqref{4.36:eq-transport-every-l-a}; here $R_*$ is the distance from $p$ in the cleanness condition on
$h$ used below, $\delta_0$ is the constant of that ceiling, and $B$ is the constant entering the
polylogarithmic factor of the flow bounds. This bound is a scalar
cutoff times \eqref{4.36:eq-normalisation-inputs-3}.

Cost along $\hat{T}$: none. Cost along $T$: a factor $\mathsf{n}$.

\emph{The Cauchy bounds.} These are Lemma~\ref{4.36:lem-cauchy-bound}, the directional Cauchy bound,
the Cauchy bound on $|\zeta|<2R$, the two-constants bound and its product form. All are one-variable or
polydisc bounds on $|\zeta|<2R$, or on the $\star$-ball of radius $c_{\star}/2$, with $c_{\star}$ the
radius constant of the norm $\|\cdot\|_{\star}$. Once the disc is given they involve no norm.

Cost along $\hat{T}$: none. Cost along $T$: a second derivative is taken at a radius divided by
$|T|$, so by Lemma~\ref{4.36:lem-cauchy-bound} it gains a factor $|T|^2\le\mathsf{n}^2$.

\emph{Cauchy bounds in the data norms.} The data norms $\|\cdot\|_{(m)}$ and $\|\cdot\|'$ of a data
field $\mathbf{A}$ are built on
$|\mathbf{A}|$, $|\nabla^{\eta}\mathbf{A}|$ and H\"older quotients, as in
\cite[(3.31)--(3.32)]{B-CMP109}. By \eqref{4.36:eq-normalisation-inputs-3}, the velocity satisfies
\[
\|\partial_\zeta\mathbf{A}_\zeta\|'\le2\max_y v_y ,\qquad v_y=C_He_y .
\]
The constant of \cite[Proposition~5]{B-CMP98} enters only inside the constant $C_B$ of these bounds,
and is absorbed by a threshold, which gives $|\partial^2F^{R}|\le1$. The radius constant is
\[
\Theta_m=\frac{4C_H}{c_B\delta_mr_m},
\]
where $c_B$ is a further constant of these bounds, $\delta_m$ is the decay constant $\delta_k$ of the
cutoff item taken at level $m$, and $r_m$
the data radius at scale $m$. Here and in the items below, $m$ denotes the level at which the data
norms are taken, not the torus index of $\mathbb{T}_m$.
All of this is stated in $|\cdot|$. This input is used together with the bound on the localised
effective-action terms in the data norm $\|\cdot\|'$, taken with both of its level parameters equal to
$m$ and with the parameters fixed there, and with the first and second variants of that bound and the
corollary of the second variant.

Cost along $\hat{T}$: none. Cost along $T$: a factor $\mathsf{n}^2$, namely $\mathsf{n}$ from the velocity
and $\mathsf{n}$ from $r_m$ divided by $\mathsf{n}$.

\emph{The Wilson part, plaquette by plaquette.} Write $\mathfrak{w}_q$ for the value of $\mathfrak{w}$
at the complexified plaquette variable of $q$. Apply Lemma~\ref{4.36:lem-plaquette-action}(a) with
$\mathsf{k}=\kappa_q$ and $\tau=|\zeta|\theta_q$; here $\tau\le\tfrac12$ on $|\zeta|\le r_0$. Combine it
with Lemma~\ref{4.36:lem-cauchy-bound} at radius $r_0$. This gives
\begin{equation}\label{4.36:eq-normalisation-inputs-4}
|\partial_t^2\mathfrak{w}_q|\le C_q\,\eta^4\ell_n^{-4}e_y,\qquad
C_q(N)=2C_P\Bigl(\frac{B_3a_1'}{r_0}+C_P\Bigr).
\end{equation}
The block $\Delta(y)$ of side $\ell_n$ contains $(\ell_n/\eta)^4$ sites at spacing $\eta$. In $d=4$ each
site carries $\binom{4}{2}=6$ plaquettes, so
\[
\sum_{q\subset\Delta(y)}\eta^4\ell_n^{-4}=6 .
\]
The trace is the normalised one, as in \cite{B-CMP109} and \cite{B-CMP99b}, and smallness is measured in
$|\cdot|$. Hence the constant relating the plaquette term to its trace is $1$ and does not depend on
$N$.

Cost along $\hat{T}$: none. Cost along $T$: a factor $\mathsf{n}^2$.

\emph{The coefficients of the plaquette terms.} By \cite[(2.24)--(2.26)]{B-CMP119}, the coefficient of
the plaquette $q$ is
\[
c(q)=x_k+\sum_{j\le k}\beta_j\bigl(1-H_j(q)\bigr) ,
\]
where $H_j$ and $\varphi_j$ are the functions of those displays, the terms $\varphi_j$ cancelling. The
count $\mathfrak{n}(q)$ attached to $q$ satisfies
\[
\mathfrak{n}(q)\le(k-n)+8+3\log_LR_n ,
\]
where $R_n$ is the scale-$n$ parameter of those displays.
Moreover \eqref{4.36:eq-flow-thresholds-c} is used at $\alpha=\tfrac18$, together with the coupling
ceiling $\bar{B}P_1(x_k)\le x_k$, where $P_1$ is the polylogarithmic factor of the flow bounds. These
are statements about the lattice and the flow; no group element enters.

Cost along either direction: none. The constant is the value $\bar{B}(N,L,\mathfrak{p})$.

\emph{The member $F^{E}$: main part.} We use the bound \cite[(1.18)]{B-CMP109} on the ball of the
membership statement above. There the function $\Psi$ bounded in \cite[(1.18)]{B-CMP109} satisfies
\[
|\Psi|\le A_{\Psi}=2E_0K_0+\bar{B}\mathfrak{A},\qquad \mathfrak{A}=96C_{\mathrm{pl}}^2 ,
\]
with constants $E_0$ and $K_0$, and with the plaquette constant $C_{\mathrm{pl}}$ of
Lemma~\ref{4.36:lem-plaquette-action}(b).
The value of $\mathfrak{A}$ comes from Lemma~\ref{4.36:lem-plaquette-action}(b) under the operator
smallness $\kappa_q+2R\theta_q\le\tfrac12$. To this we add clause (I) of the first variant of the
bound on the localised effective-action terms in the data norm, at $\mathbf{A}=0$, and the
two-constants bound.

Alternatively, one uses clauses (c) and (f) of the second variant of that bound together with the
corollary of the second variant, which carries the factor
\[
E_1'L^{-(5-\beta)(m-j)}\check{e}_{\square_0}\Theta_m ,
\]
where $E_1'$ and $\beta$ are the constant and the exponent of that corollary, and $\square_0$ is the
base cube of that corollary.
In either route the bounds are scalar on domains defined in $|\cdot|$, and $E_0$ and $E_1'$ are values
depending on $N$.

Cost along $\hat{T}$: none. Cost along $T$: a factor $\mathsf{n}^2$.

\emph{The member $F^{E}$: tail.} For the tail regions $d_j(X)\ge2(L^{m-j}-1)$, and
\cite[(1.18)]{B-CMP109} together with \cite[(3.48)]{B-CMP119} gives, on the ball of the membership
statement,
\[
|\mathbf{E}^{(j)}(X,\cdot,z)|\le e^2L^8E_0L^{-4(n-j)}e^{-(n-j)}e^{-\kappa d_j(X)/2} ,
\]
where $z$ is the remaining argument of the terms $\mathbf{E}^{(j)}$ in \cite[(1.18)]{B-CMP109}.
These bounds are combined with the Cauchy bound on $|\zeta|<2R$ and with the counting bound for the
regions. All of them are scalar.

Cost along $\hat{T}$: none. Cost along $T$: a factor $\mathsf{n}^2$.

\emph{The member $F^{R}$: small part.} We use clause (II) of the first variant of the bound on the
localised effective-action terms in the data norm, or clause (d) of its second variant. It bounds each
small term by
\[
R_1'\varpi_{jm}L^{-4(m-j)}g_j^{\kappa_0}e^{-\kappa d_j} ,
\]
with the constant $R_1'$, the weights $\varpi_{jm}$ and the exponent $\kappa_0$ of that clause.
This is combined with the Cauchy bounds in the data norms, with \cite[(2.31)]{B-CMP119} and
\cite[(3.29)]{B-CMP109}, with the counts $7^4L^{4(m-j)}$ and $13^4$, and with the flow inequality
$(\mathrm{Fl}4')$ among the consequences of the coupling flow. The bound is scalar on a data ball
defined in $|\cdot|$.

Cost along $\hat{T}$: none. Cost along $T$: a factor $\mathsf{n}^2$.

\emph{The member $F^{R}$: tail.} We use \cite[(2.31)]{B-CMP119} on the ball of the
membership statement, together with \cite[(3.48)]{B-CMP119}, the Cauchy bound on $|\zeta|<2R$ and the
counting bound. All of these are scalar.

Cost along $\hat{T}$: none. Cost along $T$: a factor $\mathsf{n}^2$.

\emph{The large-field terms.} On $\tilde{U}_j(X)$, \cite[(2.42)]{B-CMP119} gives
\[
|\mathbf{B}^{(j)}(X)|<B_0e^{-\kappa d_j(X)} ,
\]
with Ba{\l}aban's constant $B_0$.
We also use the frozen terms of \cite{B-CMP122a}, through at most $\mathfrak{N}\le R_m$
stages of the operation $\mathbf{R}$ as counted in \cite{B-CMP122b}, $R_m$ being the bound on the number
of stages. To these we add the bound $(\mathrm{P}3)^{\natural}$, stated for $N=2$ and read here at the
$SU(N)$ values of its inputs,
the arrest bound on the window $W(V_{\mathrm{off}})$ of Lemma~\ref{4.36:lem-broken-paths}, and the Cauchy
bound on $|\zeta|<2R$, in both of the classes (A) and (B) of the classification used in this chapter
for these terms. All of these are scalar bounds on domains defined in $|\cdot|$.

Cost along $\hat{T}$: none. Cost along $T$: a factor $\mathsf{n}^2$.

\emph{The far part: propagators in the moved background.} We use the propagator theorems
\cite[Theorems~3.4, 3.10]{B-CMP99b}, whose propagators are analytic in a background perturbation of the
size \cite[(3.37)]{B-CMP99b} and have random-walk decay. We also use the bound
\[
|A_j|\le C_AL^2p_0(g_j)\ell_j^{-1}
\]
of \cite[(1.8)]{B-CMP119}, for the field $A_j$ of that display, with its constant $C_A$, the degree
function $p_0(\cdot)$ and
$\ell_j=L^j\eta$. The perturbation $e^{i\eta\mathcal{H}_t}$ is controlled by
\eqref{4.36:eq-normalisation-inputs-3}. Covering the real reach $2g_k$ inside the open disc
$|\zeta|<2r_0$ requires $g_k<r_0$, which is the threshold $(\mathrm{t}8)^{+}$ of
\eqref{4.36:eq-flow-thresholds-b}.

The $\ell^2$ inner product of \cite{B-CMP99b}, written with the trace, enters only through the bound of
the absolute value of the normalised trace of a matrix by its operator norm, or through
$\|\cdot\|\le|\cdot|$. No sum over a basis of $\mathfrak{g}$ occurs, and
no factor $D$ appears. The resulting bound is linear in the reach, and a threshold absorbs it into
$\eta_{\mathrm{far}}\le1$.

Cost along $\hat{T}$: none, with $\eta_{\mathrm{far}}$ taken at reach $2g_k$. Cost along $T$: the reach
becomes $2g_k/|T|$. The main argument does not need this, because it uses operator reach below $g_k$.

\emph{The far part: one-bond moves of second-class regions.} The input is the one-bond variation bound
for large-field regions of the second class, whose input from \cite{B-CMP122b} is used as stated there;
we use it in the following form. Let $Y$ be a large-field region of the second class in the
classification of the large-field regions of the clean term $h$,
created at step $j_Y\le k$. Fix the earlier steps arbitrarily, and fix arbitrarily the carrier of the
term. Consider a real move
$u_1\mapsto u_1e^{it\hat{T}}$ of the link variable $u_1$ of one bond, with $\hat{T}\in\mathfrak{g}$,
$|\hat{T}|=1$ and $|t|\le2g_k$,
and write $V^t$ for the moved configuration. Then the activities $\mathbf{V}_{\mathsf{i}}$ of $Y$, with
their index $\mathsf{i}$ and their further argument $\mathsf{a}$ held fixed, satisfy
\begin{equation}\label{4.36:eq-normalisation-inputs-5}
\begin{split}
&\bigl|\mathbf{V}_{\mathsf{i}}(Y^{\sim};\mathsf{a};U_k(V^t))
-\mathbf{V}_{\mathsf{i}}(Y^{\sim};\mathsf{a};U_k(V))\bigr|\\
&\qquad\le C^{\mathsf{i}}_{\sigma}\,x_k^{A_\sigma}(2+k-j_Y)^{A_\sigma}
\max_{y\in\mathfrak{b}((Y^{\sim2})^{+})}e_y ,
\end{split}
\end{equation}
where $Y^{\sim}$ and $Y^{\sim2}$ are the enlarged regions attached to $Y$ in that bound,
$C^{\mathsf{i}}_{\sigma}$ is its constant, and $A_\sigma$ is its exponent increased by one.
In the summation of these bounds over the regions, the decay weight $\hat{e}_{\mathsf{N}}$ of a bond set
$\mathsf{N}$ is replaced by
\[
\hat{e}^{(3r)}_{\mathsf{N}}=\sum_{y\in\mathfrak{b}(\mathsf{N})}\bigl(1+k-n(y)\bigr)^{3r}e_y^{1/2} ,
\]
with $r$ the exponent of the flow polylogarithm in the last of the four facts below. Since
$\hat{e}_{\mathsf{N}}\le\hat{e}^{(3r)}_{\mathsf{N}}$, a bound stated with $\hat{e}_{\mathsf{N}}$ implies
the same bound stated with $\hat{e}^{(3r)}_{\mathsf{N}}$. Accordingly the exponent of the constant
$C_{\Sigma}$ of the summed bound is $A_\sigma+1+3r$.
Summed over the regions, these bounds total at most $1$. This uses four facts:
\begin{itemize}
\item cleanness of $h$, in class (A) of the classification above: the large-field regions of the term
lie in the region $Z_k(h)$ of the clean term, at distance at least $R_*$ from $p$;
\item near a block of stratum $n'$, the regions with $j_Y\ge n'-2$ have per-level multiplicity at most
$11^4$;
\item the inequality \eqref{4.36:eq-flow-thresholds-c}, used at $\alpha=\tfrac1{16}$;
\item the coupling ceiling $e^{-\delta_0R_*/64}\le e^{-5L^2(\log x_k)^r}$ with the exponent $r\ge2$
of the flow polylogarithm.
\end{itemize}
The bound is scalar and concerns a real move of one bond.

It is used at operator reach below $g_k$, and this reach lies inside
\eqref{4.36:eq-normalisation-inputs-5} whichever kind of unit direction is used to state it. Indeed,
apply the bound with $u_1=u^0$, and
let $v\ne0$ with $|v|<g_k$. Taking $\hat{T}=v/|v|$ and $t=|v|$ gives $u^0e^{iv}=u^0e^{it\hat{T}}$ with
$|t|<g_k$. Taking $T=v/|v|_{\mathfrak{g}}$ and $t'=|v|_{\mathfrak{g}}$ gives
$u^0e^{iv}=u^0e^{it'T}$ with $|T|_{\mathfrak{g}}=1$; here $|v|_{\mathfrak{g}}\le|v|$ because the trace is
normalised, so $|t'|<g_k$. Hence
\[
\{u^0e^{iv}:|v|<g_k\}\subset\{u^0e^{it\hat{T}}:|\hat{T}|=1,\ |t|<g_k\}
\]
and
\[
\{u^0e^{iv}:|v|<g_k\}\subset\{u^0e^{it'T}:|T|_{\mathfrak{g}}=1,\ |t'|<g_k\} .
\]
The ball $B^{\mathfrak{g}}_{g_k}(u^0)$ would instead require reach $\mathsf{n}g_k$. This is at most
$2g_k$ exactly when $N-1\le4$, that is, when $N\le5$; that ball is not used.

\emph{First derivatives and analyticity.} The first derivatives are bounded by the directional Cauchy
bound, the first claim of the corollary of the second variant of the bound on the localised
effective-action terms in the data norm, and
\[
r_0^{-1}\sup_{|\zeta|\le r_0}|\mathfrak{w}_q(\zeta)-\mathfrak{w}_q(0)| ,
\]
with the same sums as above. These bounds are uniform in
the source $w$ of the windowed pair: the bounds \cite[(1.18)]{B-CMP109} and \cite[(2.31)]{B-CMP119} do
not involve $w$, and \cite[(2.42)]{B-CMP119} is uniform on $\operatorname{supp}m^{V_{\mathrm{off}}}_h$.

Analyticity of $u\mapsto F^{\mathrm{near}}$ follows from three ingredients:
\begin{itemize}
\item the background field of \cite[Proposition~9]{B-CMP102b}, which is analytic in
$B\in\mathfrak{g}^c$ at $b_1$;
\item the analyticity of the pieces of $F^{\mathrm{near}}$;
\item the chart $u\mapsto(1/i)\log(uu_1^{-1})$.
\end{itemize}
These derivatives are finite along either direction.

These are the inputs used; each is stated in $|\cdot|$, is a scalar bound, or involves no
group variable, which gives (a). Along $\hat{T}$ none
of them carries a factor from the normalisation, which together with \eqref{4.36:eq-normalisation-inputs-1}
gives (b). Along $T$ each carries at most the factor $\mathsf{n}^2$ stated in its item, which gives (c).
\end{proof}

\begin{lemma}[Curvature of the conditional exponent along one link at $SU(N)$]
\label{4.36:lem-one-link-curvature}
Let $G=SU(N)$ and $\mathsf{n}=(N-1)^{1/2}$. We work in the setting of
Definition~\ref{4.36:def-unit-directions}:
\begin{itemize}
\item $h$ is a clean term of the effective density;
\item $F_h=F^{\mathrm{near}}+F^{\mathrm{far}}$, where
$F^{\mathrm{near}}=F^{W}+F^{E}+F^{R}+F^{B}+\mathrm{const}$ and
$F^{\mathrm{far}}=Q_h+\sum_Y\Psi_Y$;
\item $(V,w)$ is real windowed, $V=(u,V_{\mathrm{off}})$, with $u$ the variable on the free
link $b_1$;
\item $V^{\zeta,T}=(ue^{i\zeta T},V_{\mathrm{off}})$.
\end{itemize}
The scales, running couplings, flow inequalities and thresholds are those of
Notation~\ref{4.36:not-flow-thresholds}. For a function $\varphi$ of a real variable $t$ we
write $|\partial_t^j\varphi|_0$ for the absolute value of its $j$-th derivative at $t=0$.
Traces are normalised, $\operatorname{tr}\mathbf{1}=1$, and smallness is measured in the
operator norm $|\cdot|$. All constants depend on $(N,L,M,\mathfrak{p})$ only.

The assertions below are made for all $K$, $k$, every clean term $h$, every real windowed
$(V,w)$ and every $u^0\in W(V_{\mathrm{off}})$ with $(u^0,V_{\mathrm{off}})\in\mathcal{D}$; such
data are fixed in what follows. For $\hat T\in\mathfrak{g}$ with $|\hat T|=1$, a point
$u\in B^{\mathrm{op}}_{g_k}(u^0)$ and $V=(u,V_{\mathrm{off}})$, and for a plaquette $q$ contained
in a block $\Delta(y)$, $y\in\Gamma_n$, $n\le k$, put
\begin{equation*}
\mathcal{P}_q(\zeta):=\bigl(e^{i\eta\mathcal{H}_\zeta}U_k(V)\bigr)(\partial q),\qquad
\theta_q:=C_P\,\eta^2\ell_n^{-2}e_y .
\end{equation*}
Here $\mathcal{H}_\zeta$ is the field of Ba{\l}aban's continuation \cite[Proposition~9]{B-CMP102b}
along the source $e^{i\zeta\hat T'}$ at $b_1$, and $\hat T'=u\hat Tu^{-1}$. The constants $C_P$,
$B_3$, $a'_1$ and the radius $r_0$ are those entering the constant $C_q(N)$ of
Definition~\ref{4.36:def-unit-directions}. For a far term $\Psi_Y$, a point $u$ and real $t$ write
\begin{equation*}
\Psi_Y[t]:=\Psi_Y\bigl((ue^{it\hat T},V_{\mathrm{off}}),w\bigr).
\end{equation*}
Assume the following.
\begin{itemize}
\item[(I)] \emph{Flow.} The upper and lower flow inequalities of
Notation~\ref{4.36:not-flow-thresholds} hold at $SU(N)$ with $\bar B$ fixed. The thresholds
$(\mathrm{t}1)$, $(\mathrm{t}2)$, $(\mathrm{t}3)'$, $(\mathrm{t}4)'$, $(\mathrm{t}5)$, $(\mathrm{t}6)$,
$(\mathrm{t}8)$, $(\mathrm{t}8)^{+}$ and $(\mathrm{t}9)$ of \eqref{4.36:eq-flow-thresholds-b} hold,
among them $g_k<r_0$, together with the coupling ceilings
$x_k\ge x_{\mathrm{th}}(N,L,\mathfrak{p})$, which include $\bar BP_1(x_k)\le x_k$ and $x_k\ge1$,
with $P_1(x_k)$ the polylogarithmic factor in $x_k$ of the coefficient bound of
hypothesis (VII). Moreover the
coupling ceiling $C_Tg_kp_0(g_k)^2\le1$ holds.
\item[(II)] \emph{Complexification at $SU(N)$, in the operator norm.}
\begin{itemize}
\item[--] The map $\zeta\mapsto\mathcal{P}_q(\zeta)$ is holomorphic on $|\zeta|<2r_0$, with values
in $SL(N,\mathbb{C})$. A change of the representative in \cite[Proposition~9]{B-CMP102b}
conjugates $\mathcal{P}_q(\zeta)$ by the value of a $G^c$-valued gauge function at the base
point of $\partial q$. At real $\zeta=t$ one has
$\mathcal{P}_q(t)=U_k(V^{t,\hat T})(\partial q)\in SU(N)$.
\item[--] On each block $\Delta(y)$, $y\in\Gamma_n$, and for $|\zeta|<2r_0$, the field obeys
$\ell_n\sup_{\Delta(y)}|\mathcal{H}_\zeta|\le C_He_y|\zeta|$.
\item[--] The plaquette formula \cite[(8)]{B-CMP102b} holds in the form: whenever
$\eta|\mathcal{H}_\zeta|\le\frac14$,
\begin{equation}\label{4.36:eq-one-link-curvature-1}
\bigl|\mathcal{P}_q(\zeta)\mathcal{P}_q(0)^{-1}-\mathbf{1}\bigr|\le|\zeta|\theta_q,\qquad
\mathsf{k}_q:=|\mathcal{P}_q(0)-\mathbf{1}|\le B_3a'_1\eta^2\ell_n^{-2}.
\end{equation}
\item[--] The configurations $e^{i\eta\chi_{\mathsf{N}}\mathcal{H}_\zeta}U_k(V)$ lie in the
analyticity domains \cite[(1.11)--(1.16)]{B-CMP109}, \cite[(2.34)--(2.39)]{B-CMP119}, at half
radius, whenever $\|\chi_{\mathsf{N}}\mathcal{H}_\zeta\|_\star\le c_\star/2$; this holds for
$|\zeta|\le2R$, where $R$ is the radius of the threshold $(\mathrm{t}2)$ in
\eqref{4.36:eq-flow-thresholds-b}. The four near members $F^{W}$, $F^{E}$, $F^{R}$, $F^{B}$ are
analytic functions of the configuration on these domains.
\item[--] For $|\zeta|\le2R$ the cut-off fields obey
\begin{equation*}
\|\chi_{\mathsf{N}}\mathcal{H}_\zeta\|_\star
\le C'_\chi C_H|\zeta|\min\{1,\hat e_{\mathsf{N}^{+}}\}/\delta_k .
\end{equation*}
\item[--] Derivatives in $\zeta$ at $0$ are estimated by the one-variable Cauchy inequalities
on the disc $|\zeta|<2R$, and on the ball $\|\cdot\|_\star\le c_\star/2$ by the corresponding
Cauchy inequalities in each complex direction.
\end{itemize}
\item[(III)] \emph{Data-ball bounds.} Let $l$ denote a level. For data $\mathfrak{d}$ in the ball
$\|\mathfrak{d}\|'<c_B\delta_l$, the bounds \cite[(3.29), (3.31)--(3.32)]{B-CMP109} and
\cite[(2.31)]{B-CMP119} for the localised terms hold at $G=SU(N)$ in the form of the data-ball
lemma of this chapter with both of its level indices equal to $l$, together with its chain
forms, its corollary bounding derivatives along a data path, and the nesting condition for the
data domains. Along the data path $\zeta\mapsto\mathfrak{d}_\zeta$ of the move,
$\|\partial_\zeta\mathfrak{d}_\zeta\|'\le2\max_yC_He_y$.
\item[(IV)] \emph{The three remaining near members at $SU(N)$.} For every $\hat T\in\mathfrak{g}$
with $|\hat T|=1$ and every $u\in B^{\mathrm{op}}_{g_k}(u^0)$, the second derivatives at $t=0$ of
the functions $t\mapsto F^{E},F^{R},F^{B}$ at $(V^{t,\hat T},w)$ satisfy
\begin{equation}\label{4.36:eq-one-link-curvature-2}
|\partial_t^2F^{E}|_0\le x_k/8,\qquad|\partial_t^2F^{R}|_0\le1,\qquad|\partial_t^2F^{B}|_0\le1 .
\end{equation}
These three inequalities hold under the coupling ceiling $(\mathrm{t}3)'$ of
\eqref{4.36:eq-flow-thresholds-b}, read at the $SU(N)$ values of the constants.
\item[(IV$'$)] \emph{Analyticity and first derivatives of the three remaining near members.}
For the same $\hat T$ and $u$, the functions $t\mapsto F^{E},F^{R},F^{B}$ at
$(V^{t,\hat T},w)$ are real-analytic near $0$, and their first derivatives at $t=0$ are bounded
uniformly in $w$.
\item[(V)] \emph{The background part of the far bound at $SU(N)$.} For real $|t|\le2g_k$, the
background perturbation produced by the move $u^0\mapsto u^0e^{it\hat T}$ has size at most
$2C_Hg_ke_y$ in the operator norm, which is small compared with $a_1$, the size of background
perturbation admitted in \cite[Theorems~3.4, 3.10]{B-CMP99b}. From a perturbation of that size,
together with the bound $\sum_Y|\Psi_Y[t]-\Psi_Y[0]|\le g_kp_0(g_k)^2$ at the base point
$u=u^0$ on the changes of the far terms, the bound in clause~(b) below follows with the
constant $\eta_{\mathrm{far}}$.
\item[(VI)] \emph{One-bond response of the far terms at $SU(N)$.} Let $\Psi_Y$ be a term of the
second class in $h$, created at step $j_Y\le k$, and consider a real one-bond move
$u\mapsto ue^{it\hat T}$ with $|\hat T|=1$ and $|t|\le2g_k$. Then
\begin{equation*}
|\Psi_Y[t]-\Psi_Y[0]|\le C^{\mathfrak{p}}_\sigma\,x_k^{A_\sigma}(2+k-j_Y)^{A_\sigma}\max_y e_y ,
\end{equation*}
where the maximum runs over the blocks $y$ of the doubly enlarged support of $Y$, and where
$C^{\mathfrak{p}}_\sigma$ and $A_\sigma$ are the constant and the exponent of the one-bond
response bound, depending on $(N,L,M,\mathfrak{p})$ only.
\item[(VII)] \emph{Coefficients and sum of the Wilson member at $SU(N)$.} For every
$\hat T\in\mathfrak{g}$ with $|\hat T|=1$ and every $u\in B^{\mathrm{op}}_{g_k}(u^0)$, the member
$F^{W}$ at $(V^{t,\hat T},w)$ is the sum, over the fine plaquettes $q$ contained in the region
$B(\Omega_1(h))$ of the clean term, of
$c(q)\bigl(1-\operatorname{Re}\operatorname{tr}U_k(V^{t,\hat T})(\partial q)\bigr)$, plus terms
independent of $u$; the coefficients $c(q)$ are lattice and flow quantities independent of $t$.
If $\Xi\ge0$ and numbers $b_q\ge0$ satisfy $b_q\le\Xi\,\eta^4\ell_n^{-4}e_y$ for each such $q$
contained in a block $\Delta(y)$, $y\in\Gamma_n$, $n\le k$, then
\begin{equation*}
\begin{split}
\sum_q|c(q)|\,b_q
&\le6\Xi\sum_{n\le k}\sigma^{(1/8)}_n\sup_{\Gamma_n}|c|\\
&\le6\Xi C_\sigma\sum_{a\ge0}e^{-a}x_k\bigl[2+\bar B(a+\pi_{\mathfrak{n}}(a))\bigr].
\end{split}
\end{equation*}
Here $\sigma^{(1/8)}_n$ are the stratum weight sums of the weights $e_y$, and the weighted sum
\eqref{4.36:eq-flow-thresholds-c} at $\alpha=\frac18$ gives
$\sigma^{(1/8)}_n\le C_\sigma e^{-(k-n)}$; $C_\sigma$ and the polylogarithmic function
$\pi_{\mathfrak{n}}$ are those entering the constants $C'_W(N)$ and $S_{\mathfrak{n}}$ of
Definition~\ref{4.36:def-unit-directions}. The coefficient bound contained in the second
inequality uses the ceilings $\bar BP_1(x_k)\le x_k$ and $x_k\ge1$ of hypothesis (I).
\end{itemize}
Put
\begin{equation}\label{4.36:eq-one-link-curvature-3}
C'^{\mathrm{op}}_F(N):=C'_W(N)+3,\qquad
C'^{\mathfrak{g}}_F(N):=(N-1)\bigl(C'_W(N)+3\bigr)=\mathsf{n}^2C'^{\mathrm{op}}_F(N),
\end{equation}
with $C'_W(N)$ the constant of Definition~\ref{4.36:def-unit-directions}, and
$\eta_{\mathrm{far}}:=C_Tg_kp_0(g_k)^2$.

Then the following hold for these data.

\smallskip\noindent
\emph{The one-parameter bounds.} Let $\hat T\in\mathfrak{g}$ with $|\hat T|=1$ and
$u\in B^{\mathrm{op}}_{g_k}(u^0)$, and put $V=(u,V_{\mathrm{off}})$.
\begin{itemize}
\item[(a)] The map $t\mapsto F^{\mathrm{near}}(V^{t,\hat T},w)$ is real-analytic near $0$, and
\begin{equation}\label{4.36:eq-one-link-curvature-4}
\bigl|\partial_t^2F^{\mathrm{near}}(V^{t,\hat T},w)\bigr|_0\le C'^{\mathrm{op}}_F(N)\,x_k .
\end{equation}
\item[(b)] For real $|t|\le2g_k$,
\begin{equation*}
\bigl|F^{\mathrm{far}}((u^0e^{it\hat T},V_{\mathrm{off}}),w)
-F^{\mathrm{far}}((u^0,V_{\mathrm{off}}),w)\bigr|\le\eta_{\mathrm{far}}\le1 .
\end{equation*}
\item[(c)] One has
$\sup_w\bigl|\partial_tF^{\mathrm{near}}(V^{t,\hat T},w)\bigr|_0<\infty$.
Moreover $u\mapsto F^{\mathrm{near}}((u,V_{\mathrm{off}}),w)$ is real-analytic near every point
of $W(V_{\mathrm{off}})$.
\end{itemize}
\noindent
($\mathfrak{g}$) Let $T\in\mathfrak{g}$ with $|T|_{\mathfrak{g}}=1$.
\begin{itemize}
\item[(a)] Clause (a) of the one-parameter bounds holds along $T$ with the constant
$C'^{\mathfrak{g}}_F(N)$.
\item[(b)] Clause (b) of the one-parameter bounds holds along $T$ for real $|t||T|\le2g_k$, in
particular for $|t|\le2g_k/\mathsf{n}$.
\item[(c)] Clause (c) of the one-parameter bounds holds along $T$, the first derivative being
$|T|$ times the one along $\hat T=T/|T|$.
\end{itemize}
All of this holds at every point of $B^{\mathrm{op}}_{g_k}(u^0)$, a ball that contains
$B^{\mathfrak{g}}_{g_k/\mathsf{n}}(u^0)$.

\smallskip
At $N=2$ one has $\mathsf{n}=1$, and the two constants in
\eqref{4.36:eq-one-link-curvature-3} coincide, both being equal to $C'_W(2)+3$.
\end{lemma}

\begin{proof}
We first prove the one-parameter bounds. Fix $\hat T$ with $|\hat T|=1$, a point
$u\in B^{\mathrm{op}}_{g_k}(u^0)$, and $V=(u,V_{\mathrm{off}})$.

\emph{The Wilson member.} Only the fine plaquettes $q$ contained in the region $B(\Omega_1(h))$
of the clean term respond to the variable $u$ on $b_1$. Outside this region $U_k$ coincides
with the frozen configuration $V_0$ \cite{B-CMP119}. Each responding plaquette is assigned to
one block $\Delta(y)$, $y\in\Gamma_n$, $n\le k$.

Put $\mathfrak{w}_q(\zeta):=\mathfrak{w}(\mathcal{P}_q(\zeta))$, where $\mathfrak{w}$ is the
complexified plaquette action of Lemma~\ref{4.36:lem-plaquette-action}. This quantity does not
depend on the representative chosen in \cite[Proposition~9]{B-CMP102b}: a change of
representative conjugates $\mathcal{P}_q(\zeta)$ by the value of a $G^c$-valued gauge function
at the base point of $\partial q$, and $\operatorname{tr}\mathcal{P}$ and
$\operatorname{tr}\mathcal{P}^{-1}$ are invariant under conjugation of
$\mathcal{P}\in GL(N,\mathbb{C})$, so $\mathfrak{w}(\mathcal{P}_q(\zeta))$ is unchanged.

The function $\mathfrak{w}_q$ is holomorphic on $|\zeta|<2r_0$. Indeed, by hypothesis (II) the
map $\zeta\mapsto\mathcal{P}_q(\zeta)$ is holomorphic there with values in
$SL(N,\mathbb{C})\subset GL(N,\mathbb{C})$. By part (0) of Lemma~\ref{4.36:lem-plaquette-action},
$\mathfrak{w}$ is holomorphic on $GL(N,\mathbb{C})$.

At real $t$ we have $\mathcal{P}_q(t)=U_k(V^{t,\hat T})(\partial q)\in SU(N)$, and on $U(N)$ the
function $\mathfrak{w}$ equals $1-\operatorname{Re}\operatorname{tr}$
(Lemma~\ref{4.36:lem-plaquette-action}). Hence
\begin{equation*}
\mathfrak{w}_q(t)=1-\operatorname{Re}\operatorname{tr}U_k(V^{t,\hat T})(\partial q)
\qquad(t\in\mathbb{R}),
\end{equation*}
and $t\mapsto\mathfrak{w}_q(t)$ is real-analytic near $0$.

Let $\mathcal{P}_q(0)=U_k(V)(\partial q)\in SU(N)$ and
$\mathcal{M}_q(\zeta):=\mathcal{P}_q(\zeta)\mathcal{P}_q(0)^{-1}$. We apply the plaquette formula
\eqref{4.36:eq-one-link-curvature-1} on $|\zeta|\le r_0$; its precondition holds there. By
hypothesis (II) and $\eta\le\ell_n$, $e_y\le1$,
\begin{equation*}
\eta|\mathcal{H}_\zeta|\le\ell_n|\mathcal{H}_\zeta|\le C_H|\zeta|\le C_Hr_0\le\tfrac14 .
\end{equation*}
The last inequality is the member $(4C_H)^{-1}$ of the threshold $(\mathrm{t}8)$ in
\eqref{4.36:eq-flow-thresholds-b}. Using in addition $e_y\le1$, $\eta\le\ell_n$ and the member
$(2C_P)^{-1}$ of $(\mathrm{t}8)$, we obtain on $|\zeta|\le r_0$
\begin{equation*}
|\mathcal{M}_q(\zeta)-\mathbf{1}|\le|\zeta|\theta_q\le r_0\theta_q\le r_0C_P\le\tfrac12 .
\end{equation*}
Part (a) of Lemma~\ref{4.36:lem-plaquette-action} applies with base size $\mathsf{k}=\mathsf{k}_q$
and increment size $|\zeta|\theta_q\le\frac12$. Together with
$\mathsf{k}_q\le B_3a'_1\eta^2\ell_n^{-2}$ and $e_y\le1$ it gives, for $|\zeta|\le r_0$,
\begin{equation*}
|\mathfrak{w}_q(\zeta)-\mathfrak{w}_q(0)|\le(\mathsf{k}_q+r_0\theta_q)\,r_0\theta_q
\le C'r_0\,\eta^4\ell_n^{-4}e_y ,\qquad C':=(B_3a'_1+r_0C_P)C_P .
\end{equation*}
Indeed,
$\mathsf{k}_q+r_0\theta_q\le(B_3a'_1+r_0C_P)\eta^2\ell_n^{-2}$ and
$r_0\theta_q=r_0C_P\eta^2\ell_n^{-2}e_y$.

The closed disc $|\zeta|\le r_0$ lies inside the disc of holomorphy, so the Cauchy bound of
Lemma~\ref{4.36:lem-cauchy-bound} applies at radius $r_0$:
\begin{equation}\label{4.36:eq-one-link-curvature-5}
|\partial_t^2\mathfrak{w}_q|_0=|\mathfrak{w}''_q(0)|\le2r_0^{-2}\,C'r_0\,\eta^4\ell_n^{-4}e_y
=C_q(N)\,\eta^4\ell_n^{-4}e_y ,
\end{equation}
with $C_q(N)=2C'/r_0=2C_P(B_3a'_1/r_0+C_P)$. This bound carries no power of $x_k$ and no factor
from the normalisation of the direction. It uses only the normalised trace and smallness in the
operator norm.

By hypothesis (VII), the Wilson member is the sum over the responding plaquettes $q$ of
$c(q)\,\mathfrak{w}_q(t)$, plus terms independent of $u$, with coefficients $c(q)$ independent
of $t$. Hence it is real-analytic in $t$ near $0$.

We now sum the per-plaquette bounds by hypothesis (VII), applied with
$b_q=|\partial_t^2\mathfrak{w}_q|_0$ and $\Xi=C_q(N)$, which is admissible by
\eqref{4.36:eq-one-link-curvature-5}. In that hypothesis the factor $6$ is a count of
four-dimensional lattice geometry: a block $\Delta(y)$ of the stratum $n$ carries
$(\ell_n/\eta)^4$ sites and six plaquettes per site, so
$\sum_{q\subset\Delta(y)}\eta^4\ell_n^{-4}=6$. The weights $e_y$ of the blocks of the stratum
$n$ are collected into the stratum weight sums $\sigma^{(1/8)}_n\le C_\sigma e^{-a}$, $a=k-n$,
and on the stratum $n$ the coefficients obey
\begin{equation*}
\sup_{\Gamma_n}|c|\le x_k\bigl[2+\bar B(a+\pi_{\mathfrak{n}}(a))\bigr].
\end{equation*}
We obtain
\begin{equation*}
\begin{split}
|\partial_t^2F^{W}|_0
&\le\sum_q|c(q)|\,|\partial_t^2\mathfrak{w}_q|_0
\le6C_q(N)\sum_{n\le k}\sigma^{(1/8)}_n\sup_{\Gamma_n}|c|\\
&\le6C_q(N)C_\sigma\sum_{a\ge0}e^{-a}x_k\bigl[2+\bar B(a+\pi_{\mathfrak{n}}(a))\bigr]
\le C'_W(N)\,x_k .
\end{split}
\end{equation*}
In the last step we used $S_{\mathfrak{n}}=\sum_{a\ge0}e^{-a}[2+\bar B(a+\pi_{\mathfrak{n}}(a))]$,
the definition $C'_W(N)=36C_q(N)C_\sigma S_{\mathfrak{n}}$
(Definition~\ref{4.36:def-unit-directions}) and $6\le36$.

\emph{The members $F^{E}$, $F^{R}$, $F^{B}$.} By hypotheses (IV$'$) and (IV) these three members
are real-analytic in $t$ near $0$ and satisfy \eqref{4.36:eq-one-link-curvature-2}. The three
absorptions hold under the coupling ceiling $(\mathrm{t}3)'$ of
\eqref{4.36:eq-flow-thresholds-b} at the $SU(N)$ constants.

We record the structure of these bounds insofar as it concerns the direction; the inputs are
listed in Remark~\ref{4.36:rem-normalisation-inputs}. Every input is a scalar bound on a
configuration-space or data-space domain defined in the operator norm, such as
\cite[(1.18)]{B-CMP109}. Each is reached through one of two balls:
\begin{itemize}
\item the ball $\|\cdot\|_\star\le c_\star/2$ for $|\zeta|\le2R$ of hypothesis (II);
\item the data ball $\|\cdot\|'<c_B\delta_l$ of hypothesis (III).
\end{itemize}
The velocities come from the field bound of hypothesis (II).

Every count is a count of four-dimensional lattice geometry. These counts are:
\begin{itemize}
\item six plaquettes per site, and $96=2^4\cdot6$;
\item $(18M)^4/(26M)^4$, $13^4$, $11^4$, $9^4$, and $7^4L^{4(l-j)}$, with $l$ as in
hypothesis (III) and $j$ the level of the term counted;
\item $2^4/2^5$, and $16(1+d_j)M^4$;
\item the lattice-animal factor $e^{c_d\ell}$, and $K_0$, $C_\sigma$.
\end{itemize}
The same holds for the working cube, the counting of levels, the frozen terms, the
natural-sign statement entering the bound of $F^{B}$, the split of the large-field terms into
two classes, and the sum weighted by \eqref{4.36:eq-flow-thresholds-c}. The group never indexes
a sum.

No Gaussian or Haar normaliser enters an oscillation bound. The normalisers $\log z_j$ and
$z^{(k)}$, which carry the dimension $D=N^2-1$, enter only in two places. They sit in the
constant of $F^{\mathrm{near}}$, which does not depend on $u$ and is annihilated by $\partial_t$.
They also sit in the measure $m^{V_{\mathrm{off}}}_h$, which does not depend on $u$.

The flow factors are the flow consequences $(\mathrm{Fl}2)$--$(\mathrm{Fl}4)$ of
Notation~\ref{4.36:not-flow-thresholds}, with $(\mathrm{Fl}3)$ replaced, under the flow anchor
$x_K\ge g^{-2}$ and the coupling ceiling of Proposition~\ref{4.36:prop-transport-every-l} (whose
proof is incomplete at one step), by the transport bound \eqref{4.36:eq-transport-every-l-a}.

No input of these three members carries a norm of the direction. Along $\hat T$ they therefore
enter with no factor from its normalisation, and they enter the collection below as
$x_k/8+1+1$.

\emph{The far part.} Consider real $|t|\le2g_k$. By the threshold $(\mathrm{t}8)^{+}$ of
\eqref{4.36:eq-flow-thresholds-b}, $g_k<r_0$, so $2g_k<2r_0$. Hypothesis (V) gives that the
background perturbation along $\hat T$ has size at most $2C_Hg_ke_y$ in the operator norm,
small compared with $a_1$, the size of background perturbation admitted in
\cite[Theorems~3.4, 3.10]{B-CMP99b}. Since $\|\cdot\|\le|\cdot|$,
the same bound holds a fortiori in the normalised trace norm.

For the far terms $\Psi_Y$ we use hypothesis (VI) at the base point $u=u^0$. By cleanness of
$h$, the second-class terms that respond lie at distance at least $R_*$ from the plaquette $p$
of $h$; they lie in the region $Z_k(h)$ attached to $h$. This produces the factor
$e^{-\delta_0R_*/64}$.

Consider a block of the stratum $n'$, and put $a=k-n'$. At most $k-n'+3=a+3$ levels reach it.
At each level at most $11^4$ terms with $j_Y\ge n'-2$ lie near it, and for these
$2+k-j_Y\le a+4$. The stratum weight obeys $\sigma^{(1/16)}_{n'}\le C_\sigma e^{-(k-n')}$ by
\eqref{4.36:eq-flow-thresholds-c} at $\alpha=\frac1{16}$. Hence
\begin{equation*}
\begin{split}
\sum_Y|\Psi_Y[t]-\Psi_Y[0]|
&\le C\,x_k^{A_\sigma}\cdot11^4C_\sigma
\Bigl[\sum_{a\ge0}e^{-a}(a+3)(a+4)^{A_\sigma}\Bigr]e^{-\delta_0R_*/64}\\
&\le g_kp_0(g_k)^2 .
\end{split}
\end{equation*}
Here $C$ is a constant depending on $(N,L,M,\mathfrak{p})$ only.
The last inequality holds for $x_k\ge x_{\mathrm{th}}$, by the following argument. Since
$\delta_0M\ge16$ and $R_*\ge20L^2M(\log x_k)^r$,
\begin{equation*}
\frac{\delta_0R_*}{64}\ge\frac{16\cdot20\,L^2(\log x_k)^r}{64}=5L^2(\log x_k)^r .
\end{equation*}
With $r\ge2$, the factor $e^{-5L^2(\log x_k)^r}$ decreases faster than every power of $x_k$.
This is a coupling ceiling imposed after $N$. Because the number of levels reaching a stratum-$n'$
block is paid for by the decay of $\sigma^{(1/16)}_{n'}$, the number $K$ enters nowhere.

By hypothesis (V), this bound on the far terms together with the size of the background
perturbation gives clause (b) with $\eta_{\mathrm{far}}=C_T(N,L,M)\,g_kp_0(g_k)^2$. The inequality
$\eta_{\mathrm{far}}\le1$ is the coupling ceiling of hypothesis (I), imposed after $N$.

\emph{The ball.} Let $u\in B^{\mathrm{op}}_{g_k}(u^0)$ with $(u^0,V_{\mathrm{off}})\in\mathcal{D}$.
By part (c) of Lemma~\ref{4.36:lem-broken-paths}, $(u,V_{\mathrm{off}})$ is real windowed: the
window condition holds at every real point of the support. The same lemma keeps windowed every
real segment of length $2R$ from $u$ along $\hat T$.

Real segments are needed only in two places: the minimiser identity used in the data
comparison, and the data-ball bounds of hypothesis (III) at real points. Clause (b), taken from
the base point $u^0$, uses only $|t|\le2g_k<2r_0$, by $(\mathrm{t}8)^{+}$.

\emph{Collection.} The constant in $F^{\mathrm{near}}$ is annihilated by $\partial_t$. Adding
the Wilson bound and \eqref{4.36:eq-one-link-curvature-2} gives
\begin{equation*}
|\partial_t^2F^{\mathrm{near}}|_0\le C'_W(N)\,x_k+\frac{x_k}8+1+1\le\bigl(C'_W(N)+3\bigr)x_k .
\end{equation*}
Here we used $x/8+2\le3x$ for $x\ge16/23$ and $x_k\ge1$. This is
\eqref{4.36:eq-one-link-curvature-4}. The real-analyticity in (a) holds because each of the
four members is real-analytic in $t$ near $0$. Along $\hat T$ no factor arises from the
normalisation of the direction; this does not say that $C'_W(N)$ is independent of $N$.

\emph{Clause (c).} The function $\mathfrak{w}_q$ is holomorphic on $|\zeta|<2r_0$. Cauchy's
formula for $\mathfrak{w}_q-\mathfrak{w}_q(0)$ on the circle $|\zeta|=r_0$, which is the
first-order form of the argument of Lemma~\ref{4.36:lem-cauchy-bound}, together with the bound
on $|\mathfrak{w}_q(\zeta)-\mathfrak{w}_q(0)|$ obtained above, gives
\begin{equation*}
|\partial_t\mathfrak{w}_q|_0=|\mathfrak{w}'_q(0)|
\le r_0^{-1}\sup_{|\zeta|=r_0}|\mathfrak{w}_q(\zeta)-\mathfrak{w}_q(0)|
\le C'\eta^4\ell_n^{-4}e_y .
\end{equation*}
By hypothesis (VII), applied with $b_q=|\partial_t\mathfrak{w}_q|_0$ and $\Xi=C'$, we obtain
\begin{equation*}
\begin{split}
|\partial_tF^{W}|_0
&\le\sum_q|c(q)|\,|\partial_t\mathfrak{w}_q|_0
\le6C'\sum_{n\le k}\sigma^{(1/8)}_n\sup_{\Gamma_n}|c|\\
&\le6C'C_\sigma\sum_{a\ge0}e^{-a}x_k\bigl[2+\bar B(a+\pi_{\mathfrak{n}}(a))\bigr]
=6C'C_\sigma S_{\mathfrak{n}}\,x_k .
\end{split}
\end{equation*}
This bound does not depend on $w$. For $F^{E}$, $F^{R}$, $F^{B}$ the first derivatives at $t=0$
are bounded uniformly in $w$ by hypothesis (IV$'$). Hence
$\sup_w|\partial_tF^{\mathrm{near}}(V^{t,\hat T},w)|_0<\infty$.

The real-analyticity of $u\mapsto F^{\mathrm{near}}((u,V_{\mathrm{off}}),w)$ near every point of
$W(V_{\mathrm{off}})$ comes from the analyticity of the continuation
\cite[Proposition~9]{B-CMP102b} in its $\mathfrak{g}^c$-valued variable at $b_1$. This is
composed with the analyticity of the four near members on the domains of hypothesis (II) and
with the chart $u\mapsto\frac1i\log(u\tilde u^{-1})$ near a point $\tilde u$. This proves the
one-parameter bounds.

\emph{Proof of ($\mathfrak{g}$).} Let $T\in\mathfrak{g}$ with $|T|_{\mathfrak{g}}=1$, and put
$\hat T:=T/|T|$, so that $|\hat T|=1$. By Lemma~\ref{4.36:lem-derivative-rescaling},
\begin{equation*}
F^{\mathrm{near}}(V^{t,T},w)=F^{\mathrm{near}}(V^{t|T|,\hat T},w)
\end{equation*}
identically. Hence $t\mapsto F^{\mathrm{near}}(V^{t,T},w)$ is real-analytic near $0$ exactly when
the map along $\hat T$ is. By clause (a) of the one-parameter bounds and $|T|^2\le N-1$ (the
same lemma),
\begin{equation*}
\bigl|\partial_t^2F^{\mathrm{near}}(V^{t,T},w)\bigr|_0
=|T|^2\bigl|\partial_s^2F^{\mathrm{near}}(V^{s,\hat T},w)\bigr|_0
\le(N-1)\bigl(C'_W(N)+3\bigr)x_k .
\end{equation*}
The factor $|T|^2\le\mathsf{n}^2$ multiplies the second derivative of the whole sum
$F^{W}+F^{E}+F^{R}+F^{B}$. It is the rescaling factor, common to all four members, and no member
carries it separately.

By the reparametrisation clause of Lemma~\ref{4.36:lem-derivative-rescaling}, clause (b) along
$T$ for real
$|t||T|\le2g_k$ is literally clause (b) along $\hat T$ for $|s|\le2g_k$. Since $|T|\le\mathsf{n}$,
the condition $|t|\le2g_k/\mathsf{n}$ implies $|t||T|\le2g_k$.

By the same lemma, $\partial_tF^{\mathrm{near}}(V^{t,T},w)$ at $t=0$ equals $|T|$ times the
corresponding derivative along $\hat T$. Hence clause (c) holds along $T$, and the analyticity
in $u$ is unchanged.

All of this holds at the points of $B^{\mathrm{op}}_{g_k}(u^0)$. This ball contains
$B^{\mathfrak{g}}_{g_k/\mathsf{n}}(u^0)$, because for $v\in\mathfrak{g}$ one has
$|v|\le\mathsf{n}|v|_{\mathfrak{g}}$ (Lemma~\ref{4.36:lem-derivative-rescaling}).
\end{proof}

\begin{remark}[The curvature constant under two readings of the background-field response]
\label{4.36:rem-curvature-readings}
Write $\mathsf{n}=(N-1)^{1/2}$, let $|\cdot|$ be the operator norm and $\|X\|=(\operatorname{tr}X^{*}X)^{1/2}$
the normalised trace norm (with $\operatorname{tr}\mathbf{1}=1$). Let $C'_W(N)$ be the constant that bounds the
second derivative of the Wilson member $F^{W}$ of the conditional exponent by $C'_W(N)x_k$, as in the
proof of Lemma~\ref{4.36:lem-one-link-curvature}.

Along operator-unit directions $\hat{T}\in\mathfrak{g}$, $|\hat{T}|=1$, parts (a), (b) and (c) of
Lemma~\ref{4.36:lem-one-link-curvature} hold at $SU(N)$ with $C'^{\mathrm{op}}_{F}(N)=C'_W(N)+3$. No
factor comes from the normalisation of the direction. Along $|\cdot|_{\mathfrak{g}}$-unit directions part (a)
holds with $C'^{\mathfrak{g}}_{F}(N)=(N-1)(C'_W(N)+3)$, that is, with the factor $\mathsf{n}^{2}$; parts (b)
and (c) hold in the form stated in Lemma~\ref{4.36:lem-one-link-curvature}.

Two bounds of Ba{\l}aban enter the member of $F^{\mathrm{near}}$ that carries $x_k$:
\begin{itemize}
\item \cite[(190), p.~308]{B-CMP102b};
\item \cite[(8)]{B-CMP102b}.
\end{itemize}
We read both as operator-norm bounds, and call this the \emph{operator-norm reading}. Read instead
as bounds in the normalised trace norm with the same numbers, they give what we call the \emph{normalised
trace-norm reading}. In the display below, $C'^{\mathrm{op}}_{F}(N)$ denotes, on each line, the constant of
part (a) of Lemma~\ref{4.36:lem-one-link-curvature} along operator-unit directions under the reading named
on that line, whereas $C'_W(N)$ always denotes its value under the operator-norm reading:
\begin{equation}\label{4.36:eq-curvature-readings-a}
\begin{aligned}
C'^{\mathrm{op}}_{F}(N)&=C'_W(N)+3
&&\text{(operator-norm reading)},\\
C'^{\mathrm{op}}_{F}(N)&\le N(N-1)^{1/2}\,C'_W(N)+3
&&\text{(normalised trace-norm reading)},\\
C'^{\mathrm{op}}_{F}(N)&\le\sqrt{N}\,C'_W(N)+3
&&\text{(normalised trace-norm reading, sharper route)},\\
C'^{\mathfrak{g}}_{F}(N)&=(N-1)\,C'^{\mathrm{op}}_{F}(N)
&&\text{(either reading)}.
\end{aligned}
\end{equation}
The sharper route is described in the proof.

Under the operator-norm reading nothing else changes. Under the normalised trace-norm reading the
following change in value:
\begin{itemize}
\item the radius $\mathsf{r}_0$ of the disc in $\zeta$ on which the moved plaquette variables are
controlled;
\item the radius of the disc in $\zeta$ on which the mode bound applies;
\item the window parameters $r_m$, $\varrho$, $\Theta_m$;
\item the threshold $x_{\mathrm{th}}$ and the ceiling $\gamma_N$.
\end{itemize}
In addition, the condition on the degree function $p_0(\cdot)$ among the thresholds
\eqref{4.36:eq-flow-thresholds-b} becomes $p_0(g_k)\ge 8\mathsf{n}C_{\mathrm{loc}}$.

In the derivation of the constant $a_0(N)=(3/(8C_4C_5))^{2/D}$ from the constant
$C'^{\mathrm{op}}_F(N)$ of part (a) of Lemma~\ref{4.36:lem-one-link-curvature}, $a_0(N)$ changes by
the factor $e^{-\Delta C'^{\mathrm{op}}_F/D}$, where $\Delta C'^{\mathrm{op}}_F$ is the change of
$C'^{\mathrm{op}}_F(N)$. If the ball entering that derivation is read literally as a
$|\cdot|_{\mathfrak{g}}$-ball, the factor is
\begin{equation}\label{4.36:eq-curvature-readings-2}
e^{-\Delta C'^{\mathfrak{g}}_F/D}=e^{-\Delta C'^{\mathrm{op}}_F/(N+1)} .
\end{equation}
In either case $c_1=a_0/2$.

Under every reading the order of choice is unchanged: the constants are chosen after $N$, and the
ceilings on $\gamma$ last. No condition on $L$ arises, and no lower bound on $g$ is needed.
\end{remark}

\begin{proof}
\emph{The two direction norms.} The values of $C'^{\mathrm{op}}_F(N)$ and $C'^{\mathfrak{g}}_F(N)$ are
those of Lemma~\ref{4.36:lem-one-link-curvature}, which is proved as an implication from the statements it
names. The lemma gives $C'^{\mathfrak{g}}_F(N)=\mathsf{n}^{2}C'^{\mathrm{op}}_F(N)$ and
$\mathsf{n}^{2}=N-1$. This ratio does not depend on which reading fixes $C'_W(N)$. That gives the last line
of \eqref{4.36:eq-curvature-readings-a}.

The norms of Ba{\l}aban's papers enter in two ways, and we treat them separately. We use two comparisons.
For every complex $N\times N$ matrix $X$ with singular values $\sigma_1,\dots,\sigma_N$,
\begin{equation}\label{4.36:eq-curvature-readings-3}
\|X\|^{2}=\frac1N\sum_{i}\sigma_i^{2}\le\max_i\sigma_i^{2}=|X|^{2},
\qquad
|X|^{2}=\max_i\sigma_i^{2}\le\sum_i\sigma_i^{2}=N\|X\|^{2}.
\end{equation}

\emph{Norms on the input side.} These are the norms that define the following sets:
\begin{itemize}
\item the domains in \cite[(1.11)--(1.16)]{B-CMP109} and \cite[(2.34)--(2.39)]{B-CMP119};
\item the window in \cite[(172)]{B-CMP102b};
\item the data radii in \cite[(3.31)]{B-CMP109};
\item the size condition \cite[(3.37)]{B-CMP99b}.
\end{itemize}
Every quantity placed in these sets is bounded in $|\cdot|$. By the first inequality of
\eqref{4.36:eq-curvature-readings-3}, $\|\cdot\|\le|\cdot|$. Hence membership holds a fortiori if any of
these norms is the normalised trace norm instead of $|\cdot|$, and nothing changes.

\emph{Bounds on matrix-valued outputs that feed only absorbed members.} Two results are of this kind:
\begin{itemize}
\item the first variation of the averaging operation in \cite[Proposition~5]{B-CMP98};
\item \cite[Theorems~3.4 and 3.10]{B-CMP99b}.
\end{itemize}
Both enter only members that are absorbed by thresholds, namely $|\partial_t^{2}F^{R}|\le1$ and
$\eta_{\mathrm{far}}\le1$. A dimensional factor in either therefore changes $x_{\mathrm{th}}(N)$, not
$C'^{\mathrm{op}}_F(N)$.

\emph{The two bounds that feed the member carrying $x_k$.} The first is \cite[(190)]{B-CMP102b}.
\begin{itemize}
\item It yields the constant $C_H$ of the velocity bound $\eta|\mathcal{H}_{\zeta}|\le C_H|\zeta|$.
\item From $C_H$ comes the constant $C_P$ with $\theta_q\le C_P$. Here $\theta_q$ is the rate in
$|Q(\zeta)-\mathbf{1}|\le|\zeta|\theta_q$.
\item In that bound, $W(\zeta)\in SL(N,\mathbb{C})$ is the plaquette variable of the moved configuration,
$W_0=W(0)\in SU(N)$, and $Q(\zeta)=W(\zeta)W_0^{-1}$.
\item Through the upper bounds $\mathsf{r}_0\le(4C_H)^{-1}$ and $\mathsf{r}_0\le(2C_P)^{-1}$ among the
thresholds \eqref{4.36:eq-flow-thresholds-b}, it also fixes the radius $\mathsf{r}_0$.
\item Its local instance in the window construction is the constant $C_{\mathrm{loc}}$.
\end{itemize}
The second is \cite[(8)]{B-CMP102b}. It gives $\kappa_q:=|W_0-\mathbf{1}|\le B_3a_1'\eta^{2}\ell_n^{-2}$,
where $\ell_n=L^{n}\eta$.

These two bounds determine
\begin{equation}\label{4.36:eq-curvature-readings-4}
C_q(N)=\frac{2C_PB_3a_1'}{\mathsf{r}_0}+2C_P^{2},
\end{equation}
which is the constant of the bound
\begin{equation*}
\bigl|\partial_t^{2}w_q(t)|_{t=0}\bigr|\le C_q(N)\eta^{4}\ell_n^{-4}e_y
\end{equation*}
for the complexified action $w_q(\zeta)=\mathfrak{w}(W(\zeta))$ of one plaquette $q$, with $\mathfrak{w}$ as in
Lemma~\ref{4.36:lem-plaquette-action}. Summing over the plaquettes with
the lattice and flow coefficients yields $C'_W(N)$ as a fixed multiple of $C_q(N)$.

\emph{The operator-norm reading.} This is the norm declared in \cite[(19)--(20)]{B-CMP98}. It is also the
norm used in \cite[(3.6)--(3.8)]{B-CMP119}: there the operator norm of the field $\mathbf{H}$ is read off
\cite[(190)]{B-CMP102b}, and the small-field condition
\cite[(2.17)]{B-CMP119}, stated in the operator norm, is inferred from the inequality
\begin{equation*}
2(1+\beta_0)L^{-2}+4/10<1 ,
\end{equation*}
where $\beta_0$ is the constant of \cite{B-CMP119} which, together with $L$, controls the small-field
condition \cite[(2.17)]{B-CMP119} through this inequality.
In this inequality $N$ does not occur. The operator-norm reading also agrees with the small-field
conditions used in this chapter, which are Ba{\l}aban's and are stated in $|\cdot|$. Under this reading
\cite[(190)]{B-CMP102b} and \cite[(8)]{B-CMP102b}
are bounds in $|\cdot|$, exactly as they are used in the proof of
Lemma~\ref{4.36:lem-one-link-curvature}. Nothing changes, and the first line of
\eqref{4.36:eq-curvature-readings-a} is that lemma's value.

\emph{The normalised trace-norm reading.} Suppose instead that the two bounds are normalised trace-norm
bounds with the same numbers. Converting each of them to $|\cdot|$ costs the following.
\begin{itemize}
\item On $\mathfrak{g}^{c}$-valued outputs over the complex disc the cost is at most $\sqrt{N}$, by items
(F6) and (F10) of \eqref{4.36:eq-eigenangle-inequalities-a} in
Lemma~\ref{4.36:lem-eigenangle-inequalities} (this is also the second inequality of
\eqref{4.36:eq-curvature-readings-3}). The factor $\mathsf{n}$ is not available there. Indeed, let
$E_{12}\in\mathfrak{g}^{c}$ be the matrix whose only non-zero entry is a $1$ in place $(1,2)$. Then
$|E_{12}|=1$ and $\|E_{12}\|=N^{-1/2}$, so the factor $\sqrt N$ is attained. Hence $C_H$ and $C_P$ are
multiplied by at most $\sqrt{N}$.
\item On the $SU(N)$-valued real plaquette variables the cost is at most $\mathsf{n}$, by item (F9) of
\eqref{4.36:eq-eigenangle-inequalities-a} in Lemma~\ref{4.36:lem-eigenangle-inequalities}. Hence
$B_3a_1'$ is multiplied by at most $\mathsf{n}$.
\item On the $\mathfrak{g}$-valued real field $\mathcal{H}_t$ of the window construction the cost is at
most $\mathsf{n}$: for $X\in\mathfrak{g}$ one has $\|X\|=(\operatorname{tr}X^{2})^{1/2}=|X|_{\mathfrak{g}}$,
and the second inequality $|X|\le\mathsf{n}|X|_{\mathfrak{g}}$ of item (F1) of
\eqref{4.36:eq-eigenangle-inequalities-a} in Lemma~\ref{4.36:lem-eigenangle-inequalities} applies.
Hence $C_{\mathrm{loc}}$ is multiplied by at most
$\mathsf{n}$. This is why the condition on $p_0(g_k)$ takes the form $p_0(g_k)\ge8\mathsf{n}C_{\mathrm{loc}}$.
\end{itemize}
Through the two upper bounds $(4C_H)^{-1}$ and $(2C_P)^{-1}$, the radius $\mathsf{r}_0$ shrinks by a
factor of at most $\sqrt{N}$.

Insert these factors in \eqref{4.36:eq-curvature-readings-4}.
\begin{itemize}
\item The first summand is multiplied by at most $\sqrt{N}\cdot\mathsf{n}\cdot\sqrt{N}=N\mathsf{n}$.
\item The second summand is multiplied by at most $(\sqrt{N})^{2}=N$.
\end{itemize}
A sum of non-negative terms is multiplied by at most the larger factor. Since $\mathsf{n}\ge1$ for $N\ge2$,
\begin{equation*}
\max\{N\mathsf{n},N\}=N\mathsf{n} .
\end{equation*}
So $C_q(N)$ grows by at most the factor $N\mathsf{n}=N(N-1)^{1/2}$. Because $C'_W(N)$ is a fixed multiple of
$C_q(N)$, it grows by the same factor.

The second derivatives of the three remaining members $F^{E}$, $F^{R}$, $F^{B}$ of $F^{\mathrm{near}}$ are
bounded by $x_k/8$, $1$ and $1$. These bounds are absorbed at the moved thresholds, so the summand $3$ is
unchanged. This proves the second line of
\eqref{4.36:eq-curvature-readings-a}.

In the same way the radius of the mode-bound disc, $r_m$, $\varrho$, $\Theta_m$, the further constants of
the window construction, $x_{\mathrm{th}}$ and $\gamma_N$ change in value. The conditions of the step that
require room in the coupling remain ceilings on $\gamma$, because the exponent $q_0$ occurring in them
exceeds the exponent $p_0$ of the degree function $p_0(\cdot)$.

\emph{The sharper route.} Keep the numbers of the two bounds as normalised trace-norm sizes. Secure only
the operator-norm margins
\begin{equation*}
|Q(\zeta)-\mathbf{1}|\le\tfrac12,\qquad|\eta\mathcal{H}_{\zeta}|\le\tfrac14 ,
\end{equation*}
by working at the radius $\mathsf{r}_0/\sqrt{N}$. Then part~(a$^{\flat}$) of
Lemma~\ref{4.36:lem-plaquette-action} gives that $C_q(N)$ grows by at most the factor $\sqrt{N}$. This
gives the third line of \eqref{4.36:eq-curvature-readings-a}. The sharper bound is not used elsewhere.

\emph{The unnormalised Frobenius norm.} Suppose the bounds are read in the unnormalised Frobenius norm
$\sqrt{N}\|\cdot\|$. By the second inequality of \eqref{4.36:eq-curvature-readings-3} this norm dominates
$|\cdot|$. The conclusions of the operator-norm reading then hold verbatim.

\emph{The constant $a_0(N)$.} Along the derivation of $a_0(N)=(3/(8C_4C_5))^{2/D}$ from
$C'^{\mathrm{op}}_F(N)$ named in the remark, a
change $\Delta C'^{\mathrm{op}}_F$ of the curvature constant changes $a_0(N)$ by the factor
$e^{-\Delta C'^{\mathrm{op}}_F/D}$. If the ball entering that derivation is read literally as a
$|\cdot|_{\mathfrak{g}}$-ball, the constant is
$C'^{\mathfrak{g}}_F$ instead. By the last line of \eqref{4.36:eq-curvature-readings-a},
$\Delta C'^{\mathfrak{g}}_F=(N-1)\Delta C'^{\mathrm{op}}_F$. Since $D=N^{2}-1=(N-1)(N+1)$,
\begin{equation*}
\frac{\Delta C'^{\mathfrak{g}}_F}{D}=\frac{\Delta C'^{\mathrm{op}}_F}{N+1},
\end{equation*}
which is \eqref{4.36:eq-curvature-readings-2}. The constant $c_1=a_0/2$ follows $a_0$.

\emph{Order of choice.} In every case above the only quantities that move are constants depending on $N$,
chosen after $N$, and ceilings on $\gamma$, chosen last. No condition on $L$ and no lower bound on $g$
appear.
\end{proof}

\begin{remark}[Frozen-background curvature, minimiser response and operator-norm smallness]
\label{4.36:rem-frozen-background}
Throughout, $|\cdot|$ is the operator norm, $\|\cdot\|$ the normalised trace norm
$\|Y\|^2=\operatorname{tr}(Y^*Y)$ for an $N\times N$ matrix $Y$, $|\cdot|_{\mathfrak{g}}$ the norm
$(\operatorname{tr}X^2)^{1/2}$ on $\mathfrak{g}$, and $\mathsf{n}=(N-1)^{1/2}$.
\begin{enumerate}
\item[(i)] \emph{Two different curvatures.} Let $b_1$ be the free bond and $u_1\in SU(N)$ its
variable. Let $p$ be a plaquette containing $b_1$, and let $\Pi_1\in SU(N)$ be the ordered product
along $\partial p$ of the three other bond variables, held fixed, so that $U(\partial p)$ is
conjugate to $u_1\Pi_1$. For $T\in\mathfrak{g}$ the frozen-background plaquette function
$t\mapsto1-\operatorname{Re}\operatorname{tr}(e^{itT}u_1\Pi_1)$ satisfies
\begin{equation}\label{4.36:eq-frozen-background-1}
\begin{gathered}
\frac{d^2}{dt^2}\Bigl(1-\operatorname{Re}\operatorname{tr}\bigl(e^{itT}u_1\Pi_1\bigr)\Bigr)
=\operatorname{Re}\operatorname{tr}\bigl(T^2e^{itT}u_1\Pi_1\bigr),\\
\bigl|\operatorname{Re}\operatorname{tr}\bigl(T^2e^{itT}u_1\Pi_1\bigr)\bigr|
\le|T|_{\mathfrak{g}}^2 ,
\end{gathered}
\end{equation}
a bound in which $N$ does not appear. The Wilson member
$F^{W}=-\sum_q c(q)\,\mathfrak{w}\bigl(U_k(V^{t,T})(\partial q)\bigr)$ of the
conditional exponent, with the coefficients $c(q)$ appearing in the proof of
Lemma~\ref{4.36:lem-one-link-curvature}, which depend on the lattice geometry and the flow only,
is a different object: it sums over the minimiser's fine plaquettes $U_k(V^{t,T})(\partial q)$,
where $V^{t,T}$ is the configuration moved by $t$ along $T$ at the free bond, as in
Lemma~\ref{4.36:lem-one-link-curvature}; at real $t$ each summand equals
$1-\operatorname{Re}\operatorname{tr}U_k(V^{t,T})(\partial q)$, since
$\mathfrak{w}=1-\operatorname{Re}\operatorname{tr}$ on $U(N)$
(Lemma~\ref{4.36:lem-plaquette-action}, part (0)). Its dependence
on $u_1$ passes through the field $\mathcal{H}_\zeta$ of the minimiser's response, where $\zeta$
is the complex parameter extending $t$ (the moved configuration being $V^{\zeta,T}$), and its
second derivative is bounded by Cauchy's inequality on a disc in $\zeta$ of radius
$r_{\mathcal{H}}$ (of radius $r_{\mathcal{H}}/|T|$ along a general $T\in\mathfrak{g}$), contained
in the domain of holomorphy in $\zeta$ of the minimiser's response; $r_{\mathcal{H}}$ is an
operator-norm quantity, measured in $|\cdot|$.
Along $\hat T$ with $|\hat T|=1$ this bound, the one entering
clause (a) of the one-parameter bounds of Lemma~\ref{4.36:lem-one-link-curvature}, carries no
power of $\mathsf{n}$.
Along $T$ with $|T|_{\mathfrak{g}}=1$ the same bound is multiplied by $|T|^2\le N-1$. No
bound on the second derivative of $F^{W}$ along $\mathfrak{g}$-unit directions
without the factor $|T|^2$ is asserted. Such a bound would require two ingredients: the estimate
\cite[(190)]{B-CMP102b} in a form in which both its input and its output are measured in the
trace norm $\|\cdot\|$, and a second-variation formula for the
minimiser's plaquettes.
\item[(ii)] \emph{Smallness in the operator norm.} In
Lemma~\ref{4.36:lem-plaquette-action}, parts (a) and (b), the smallness hypotheses are read in
$|\cdot|$:
\begin{equation*}
|W_0-\mathbf{1}|\le\mathsf{k},\qquad |Q-\mathbf{1}|\le\tau\le\tfrac12,\qquad |W-\mathbf{1}|\le\tfrac12 .
\end{equation*}
This is the norm in which their inputs are given: $\kappa_q:=|W_0-\mathbf{1}|$, the
operator-norm size of the base plaquette $W_0=U_k(V)(\partial q)$, plays the role of
$\mathsf{k}$ and is bounded by means of \cite[(8)]{B-CMP102b}, and $\theta_q$ is the
operator-norm response rate of the plaquette $q$: the ratio $Q(\zeta)=W(\zeta)W_0^{-1}$ of the
fine plaquette $W(\zeta)$ of the moved configuration to the base plaquette satisfies
$|Q(\zeta)-\mathbf{1}|\le|\zeta|\theta_q$ for $|\zeta|\le r_{\mathcal{H}}$, and the same rate
$\theta_q$ enters the smallness hypothesis for the complex plaquette variable below.
When the smallness hypothesis for the complex plaquette variable $\mathfrak{U}(\partial q)$,
read on the disc in $\zeta$ of radius $2R_{\mathfrak{U}}$, is taken in the operator norm,
\begin{equation*}
|\mathfrak{U}(\partial q)-\mathbf{1}|\le\kappa_q+2R_{\mathfrak{U}}\theta_q\le\tfrac12 ,
\end{equation*}
both summands being operator-norm quantities, one has
\begin{equation}\label{4.36:eq-frozen-background-2}
|\mathfrak{w}(\mathfrak{U}(\partial q))|\le\|\mathfrak{U}(\partial q)-\mathbf{1}\|^2
\qquad\text{at every }N,
\end{equation}
with constant $1$: the hypothesis is in the operator norm, the conclusion in the square of the
trace norm. If the hypotheses are
instead stated with the normalised $\|\cdot\|$, both \eqref{4.36:eq-frozen-background-2} and the
increment bound of Lemma~\ref{4.36:lem-plaquette-action}(a) fail for $N\ge3$: with constant $1$
at $N=3$, with every constant for $N\ge4$. At $N=2$
the hypotheses in $\|\cdot\|$ cost only the factor $1+1/\sqrt2$ in the conclusions. If
$\|\cdot\|$ is replaced by the unnormalised Frobenius norm, they cost nothing.
\end{enumerate}
\end{remark}

\begin{proof}
(i) Differentiating $t\mapsto\operatorname{tr}(e^{itT}u_1\Pi_1)$ twice brings down $(iT)^2=-T^2$.
This gives the identity in \eqref{4.36:eq-frozen-background-1}, since $\operatorname{Re}$ and
$\operatorname{tr}$ are linear.

For the bound, let $e_1,\dots,e_N$ be the standard basis of $\mathbb{C}^N$. The matrix
$e^{itT}u_1\Pi_1$ is unitary, so its operator norm is $1$. By cyclicity of the trace and since
$T$ is Hermitian,
\begin{equation*}
\operatorname{tr}\bigl(T^2e^{itT}u_1\Pi_1\bigr)=\operatorname{tr}\bigl(Te^{itT}u_1\Pi_1T\bigr)
=N^{-1}\sum_{j=1}^N\bigl\langle Te_j,\,e^{itT}u_1\Pi_1\,Te_j\bigr\rangle ,
\end{equation*}
and each summand has modulus at most the squared Euclidean length $|Te_j|^2$ of the vector
$Te_j$, since the unitary factor has operator norm $1$. Hence
\begin{equation*}
\bigl|\operatorname{Re}\operatorname{tr}\bigl(T^2e^{itT}u_1\Pi_1\bigr)\bigr|
\le N^{-1}\sum_{j=1}^N|Te_j|^2=\operatorname{tr}T^2=|T|_{\mathfrak{g}}^2 .
\end{equation*}
This is the $N$-free curvature inequality among the elementary inequalities
\eqref{4.36:eq-eigenangle-inequalities-a} of Lemma~\ref{4.36:lem-eigenangle-inequalities}.

For the Wilson member we combine three facts.
\begin{itemize}
\item In the proof of clause (a) of the one-parameter bounds of
Lemma~\ref{4.36:lem-one-link-curvature}, the Wilson member
is bounded along $\hat T$ with $|\hat T|=1$ by
$|\partial_t^2F^{W}(V^{t,\hat T})|_{t=0}|\le C'_W(N)x_k$. This bound is obtained from
Lemma~\ref{4.36:lem-plaquette-action}(a) and Cauchy's inequality at the operator-norm radius
$r_{\mathcal{H}}$, with no direction norm entering; hence it carries no factor $\mathsf{n}$.
\item Put $\hat T:=T/|T|$ and $s:=t|T|$. Then $V^{t,T}=V^{s,\hat T}$ by definition. The chain
rule gives
\begin{equation*}
\partial_t^2F^{W}(V^{t,T})=|T|^2\,\partial_s^2F^{W}(V^{s,\hat T}).
\end{equation*}
\item The inequality $|T|\le\mathsf{n}|T|_{\mathfrak{g}}$ of
Lemma~\ref{4.36:lem-eigenangle-inequalities} gives $|T|^2\le N-1$ when $|T|_{\mathfrak{g}}=1$.
\end{itemize}
Hence along $\hat T$ the bound on the second derivative of $F^{W}$ carries no factor
$\mathsf{n}$, and along $T$ with $|T|_{\mathfrak{g}}=1$ it reads
\begin{equation*}
\bigl|\partial_t^2F^{W}(V^{t,T})|_{t=0}\bigr|\le(N-1)\,C'_W(N)\,x_k ;
\end{equation*}
this proves the assertions of (i).

(ii) Let $|W-\mathbf{1}|\le\frac12$. Then \eqref{4.36:eq-frozen-background-2} is
Lemma~\ref{4.36:lem-plaquette-action}(b), applied to $W=\mathfrak{U}(\partial q)$.

Failure for $N=3$ with constant $1$. Take $W=\operatorname{diag}(1/7,1,1)$. Then
\begin{equation*}
\|W-\mathbf{1}\|^2=\tfrac13\bigl(\tfrac67\bigr)^2=\tfrac{12}{49}<\tfrac14 ,\qquad
\operatorname{tr}W=\tfrac57,\qquad \operatorname{tr}W^{-1}=3 .
\end{equation*}
Hence $\mathfrak{w}(W)=-\frac67$ and $|\mathfrak{w}(W)|=\frac{42}{49}>\frac{12}{49}$, so
\eqref{4.36:eq-frozen-background-2} fails under $\|W-\mathbf{1}\|\le\frac12$. With
$W_0=\mathbf{1}$, $\mathsf{k}=0$, $Q=W$ and $\tau=\|Q-\mathbf{1}\|$, one has
$\mathfrak{w}(W_0)=0$, and the increment bound of Lemma~\ref{4.36:lem-plaquette-action}(a) with
constant $1$ reads $|\mathfrak{w}(W)|\le\tau^2=\frac{12}{49}$; it fails by the same computation.
For $N\ge4$ take $W=\operatorname{diag}(\delta,1,\dots,1)$ with $0<\delta<1$. Then
\begin{equation*}
\|W-\mathbf{1}\|=(1-\delta)N^{-1/2}\le\tfrac12 ,\qquad
\mathfrak{w}(W)=N^{-1}-\frac{\delta+\delta^{-1}}{2N}\longrightarrow-\infty
\quad(\delta\downarrow0).
\end{equation*}
With $W_0=\mathbf{1}$ and $Q=W$, the right-hand sides of both bounds with a constant $C$ are at
most $C/4$, while $|\mathfrak{w}(W)|$ is unbounded; so no constant suffices.

At $N=2$, for a $2\times2$ matrix $Y$, $|Y|^2$ is the larger eigenvalue of $Y^*Y$, and
$2\|Y\|^2=\operatorname{Tr}(Y^*Y)$ is the sum of both eigenvalues; hence
$|Y|\le\sqrt2\|Y\|$. A hypothesis
$\|W-\mathbf{1}\|\le\frac12$ therefore gives $|W-\mathbf{1}|\le1/\sqrt2$. The Neumann series
then gives
\begin{equation*}
|W^{-1}|\le\bigl(1-1/\sqrt2\bigr)^{-1}=2+\sqrt2 .
\end{equation*}
The first inequality of Lemma~\ref{4.36:lem-plaquette-action}(b),
$|\mathfrak{w}(W)|\le\frac12|W^{-1}|\,\|W-\mathbf{1}\|^2$, uses only invertibility. It
therefore yields $|\mathfrak{w}(W)|\le(1+1/\sqrt2)\|W-\mathbf{1}\|^2$.

In the same way, in the proof of Lemma~\ref{4.36:lem-plaquette-action}(a) the bound $|Q^{-1}|\le2$
is replaced by $|Q^{-1}|\le2+\sqrt2$, so that
\begin{equation*}
|\mathfrak{w}(W)-\mathfrak{w}(W_0)|\le
\mathsf{k}\tau+\bigl(1+1/\sqrt2\bigr)\tau^2\le\bigl(1+1/\sqrt2\bigr)(\mathsf{k}+\tau)\tau .
\end{equation*}
Finally, for the unnormalised Frobenius norm $(\operatorname{Tr}Y^*Y)^{1/2}$ one has
$|Y|\le(\operatorname{Tr}Y^*Y)^{1/2}$. A hypothesis stated in that norm is therefore already
an operator-norm hypothesis.
\end{proof}

\begin{lemma}[Anti-concentration: cancellation of the coupling powers at every $N$]
\label{4.36:lem-anti-concentration}
Let $G=SU(N)$ with normalised Haar measure $du$, and put $D=N^2-1$. Let $M_0\in G$, $x>0$ and
$C_4\ge 0$. Let $S\subset G$ be a Borel set, and let $f\ge 0$ be a Borel function on $G$ with
$f\le C_4x^{D/2}$ on $S$. Then for every $a>0$
\begin{equation}\label{4.36:eq-anti-concentration-1}
\int_S \mathbf{1}\bigl\{1-\operatorname{Re}\operatorname{tr}(uM_0)<a/x\bigr\}\,f(u)\,du
\;\le\; C_4\,C_5(N)\,a^{D/2},
\end{equation}
where $C_5(N)$ is the constant of Lemma~\ref{4.36:lem-haar-small-sets}.
\end{lemma}

\begin{proof}
Fix $a>0$. The function $u\mapsto 1-\operatorname{Re}\operatorname{tr}(uM_0)$ is continuous on
$G$, so the set
\begin{equation*}
E=\bigl\{u\in G:\ 1-\operatorname{Re}\operatorname{tr}(uM_0)<a/x\bigr\}
\end{equation*}
is open, hence Borel, and the integrand in \eqref{4.36:eq-anti-concentration-1} is a
non-negative Borel function. On $S$ we have $0\le f\le C_4x^{D/2}$, so the integrand is bounded
on $S$ by $C_4x^{D/2}\mathbf{1}_E$. Integrating over $S$ and enlarging $S\cap E$ to $E$ gives
\begin{equation*}
\int_S \mathbf{1}_E(u)\,f(u)\,du\;\le\;C_4x^{D/2}\operatorname{Haar}(S\cap E)
\;\le\;C_4x^{D/2}\operatorname{Haar}(E).
\end{equation*}
The set $E$ is the preimage of
\begin{equation*}
E'=\bigl\{g'\in G:\ 1-\operatorname{Re}\operatorname{tr}g'<a/x\bigr\}
\end{equation*}
under the right translation $u\mapsto uM_0$. By Lemma~\ref{4.36:lem-haar-small-sets}(c) this
translation preserves Haar measure, so $\operatorname{Haar}(E)=\operatorname{Haar}(E')$. By
Lemma~\ref{4.36:lem-haar-small-sets}(b), which holds for every $\tau>0$, applied with
$\tau=a/x>0$, we have $\operatorname{Haar}(E')\le C_5(N)(a/x)^{D/2}$. Combining these,
\begin{align*}
\int_S \mathbf{1}_E(u)\,f(u)\,du
&\le C_4x^{D/2}\operatorname{Haar}(E)
 = C_4x^{D/2}\operatorname{Haar}(E')\\
&\le C_4x^{D/2}\,C_5(N)\,(a/x)^{D/2}
 = C_4\,C_5(N)\,a^{D/2},
\end{align*}
the last equality being $x^{D/2}(a/x)^{D/2}=a^{D/2}$. This is
\eqref{4.36:eq-anti-concentration-1}.
\end{proof}

Where the lemma is applied in this chapter, the constant $C_4x^{D/2}$ is a pointwise density bound
of the form
\begin{equation*}
C_4x^{D/2}=c^{-1}e^{\Lambda r^2/2+1}r^{-D},
\end{equation*}
in which $c>0$ and $\Lambda$ are fixed constants and the bound is obtained on a ball of radius $r$
proportional to $x^{-1/2}$. The two exponents that meet in the proof are therefore the same number
$D/2$, half the dimension of the one-link fibre $G$. It enters once through the volume of the ball
of radius $r\propto x^{-1/2}$, which gives the factor $x^{D/2}$, and once through the sublevel set
of height $a/x$, which gives the factor $x^{-D/2}$. Consequently the power $x^{(N^2-1)/2}$ cancels
identically at every $N$, and no relation between $a$ and $x$ is required.

\begin{theorem}[lower and anti-concentration bounds for the averaged plaquette at $\mathrm{SU}(N)$]
\label{4.36:thm-plaquette-lower-bound}
Let $G=\mathrm{SU}(N)$, $D=N^2-1=\dim\mathfrak{g}$ and $\mathsf{n}=(N-1)^{1/2}$. Let $V_k$ be the
block field at level $k$ of Definition~\ref{1.5:def-block-average}. Let $g_k$, $x_k=g_k^{-2}$ and
$R_k$ be as in Notation~\ref{4.36:not-flow-thresholds}. For a plaquette $p$ of the block lattice
$\mathbb{T}^{(k)}$ put
\[
O:=1-\operatorname{Re}\operatorname{tr}V_k(\partial p).
\]
For a history $h$ let $\mathrm{term}_h\ge0$ be its term in the effective density at level $k$. Put
\[
\nu^{B}(F)=\mathcal{Z}^{-1}\sum_h\int F\,\mathrm{term}_h\,dV_k ,
\]
with $\mathcal{Z}$ chosen so that $\nu^{B}$ is a probability measure. Let $\operatorname{Av}_p$ be
the average over $\tau\in\mathfrak{P}$ of the measures $\nu^{B}\circ\tau$. Here $\mathfrak{P}$ is
the set of translations of $\mathbb{T}^{(k)}$ provided by hypothesis (h) below.

Let $b_1$ be the first bond of $\partial p$. Write $V_k=(u,V_{\mathrm{off}})$, where $u\in G$ is
the variable on $b_1$ and $V_{\mathrm{off}}$ collects the variables on the other bonds. Clean
histories, the good sets $G(h)$ (the letter $G$ with an argument always denotes a good set, never
the group), the window $W(V_{\mathrm{off}})$, the deep set $\mathcal{D}$ and
the members $F^{\mathrm{near}}$, $F^{\mathrm{far}}$ of the conditional exponent are those of
Lemma~\ref{4.36:lem-broken-paths} and Lemma~\ref{4.36:lem-one-link-curvature}.

Assume the following statements at $G=\mathrm{SU}(N)$.
\begin{enumerate}
\item[(a)] \emph{Factorisation of a clean term.} For every clean $h$ and every
$V_{\mathrm{off}}$,
\[
\mathrm{term}_h(u,V_{\mathrm{off}})
=\mathbf{1}_{W(V_{\mathrm{off}})}(u)\,\Psi_h(V_{\mathrm{off}})\,Z_h(u),
\qquad
Z_h(u):=\int e^{F_h(u,V_{\mathrm{off}};\varpi)}\,m^{V_{\mathrm{off}}}_h(d\varpi).
\]
Here $F_h=F^{\mathrm{near}}+F^{\mathrm{far}}$ exactly, and $\Psi_h\ge0$ is Borel. The integration
variable is written $\varpi$; it is the variable written $w$ in Lemma~\ref{4.36:lem-mode-bound} and
Lemma~\ref{4.36:lem-one-link-curvature}. The measure
$m^{V_{\mathrm{off}}}_h\ge0$ is $\sigma$-finite on a measure space $\mathcal{W}$. It carries only
characteristic functions and $\delta$-functions that do not depend on $u$, the $\delta$-functions
being resolved in Ba{\l}aban's own variables \cite{B-CMP119}. Every exponent that depends on $u$
through $U_k$, including those of the large-field operations $\mathbf{T}'_k(Y)$, $Y$ a large-field
region of the history $h$, is a member of
$F^{\mathrm{near}}+F^{\mathrm{far}}$. All integrands are non-negative and Borel.
\item[(b)] \emph{The large-field discard bound.} There are
$\gamma^{B}=\gamma^{B}(N,L,M,\mathfrak{p})>0$ and $s^{B}_*=s^{B}_*(N,L,M)$ such that
for $0<g\le\gamma^{B}$, every admissible $K$, every $s$
with $s^{B}_*\le s\le K-1$ and $k=K-s$ the following holds. After the average
$\operatorname{Av}_p$, the total mass of the terms of non-clean histories and of the factors
$1-\mathbf{1}_{G(h)}$ of clean ones is at most $\vartheta^{B}(s)$, and
$\vartheta^{B}(s)\le\frac18$. In
this bound the constant $\gamma_0$ is replaced by $\gamma_0/N$, and $c^{B}=\gamma_0(A_1/A_0)^2$ by
$c^{B}/N$.
\item[(c)] \emph{The window inputs.} These are the three inherited inputs of
Lemma~\ref{4.36:lem-broken-paths} (the statement placing the good points of a clean term in the
deep set $\mathcal{D}$, the localised field bound, and the localised estimate that the lemma
assumes at its base point), together with the starting bounds of the window analysis, the bounds in
the configuration norm $\|\cdot\|_{\star}$ and the remaining estimates of the same group, all taken
in the operator norm $|\cdot|$. In
particular, if $u^0\in W(V_{\mathrm{off}})$ and $\mathbf{1}_{G(h)}((u^0,V_{\mathrm{off}}))>0$,
then $(u^0,V_{\mathrm{off}})\in\mathcal{D}$ and the localised base-point estimate of
Lemma~\ref{4.36:lem-broken-paths} holds at $(u^0,V_{\mathrm{off}})$.
\item[(d)] \emph{The data-norm estimate and the nesting statement.} This is the data-norm
estimate in $\|\cdot\|'$ at coinciding levels, with its three constants held fixed, together with
the nesting statement.
\item[(e)] \emph{The bound on the far member $F^{\mathrm{far}}$ of the conditional exponent along one
link.} It enters through clause (b) of the one-parameter bounds of
Lemma~\ref{4.36:lem-one-link-curvature}, which bounds
the change of $F^{\mathrm{far}}$ when the free link variable is moved from $u^0$ to
$u^0e^{it\hat{T}}$, for $\hat{T}\in\mathfrak{g}$ with $|\hat{T}|=1$ and real $|t|\le2g_k$, by
$\eta_{\mathrm{far}}\le1$.
\item[(f)] \emph{The plaquette statements.} These are the bound on the complex plaquette variable
$\mathfrak{U}(\partial q)$, with its constant $\mathfrak{A}$; the further bounds of the same group on
the members of the conditional exponent; and the first clause of the last statement of that
group.
\item[(g)] \emph{The flow inequalities.} There are
$\gamma^{\mathrm{fl}}=\gamma^{\mathrm{fl}}(N,L,M,\mathfrak{p})>0$ and
$s^{\mathrm{fl}}_*=s^{\mathrm{fl}}_*(N,L,M)$ such that for $0<g\le\gamma^{\mathrm{fl}}$, every
admissible $K$, every $s$
with $s^{\mathrm{fl}}_*\le s\le K-1$ and $k=K-s$ the flow inequalities hold with $\bar{B}$ fixed and
anchored at $g_K\in[g_{-}(g),g[$ as in
\eqref{4.36:eq-flow-thresholds-a}--\eqref{4.36:eq-flow-thresholds-b}.
They include
\[
x_k\le g^{-2}+B^{\sharp}(s+1),\qquad x_k\ge x_{\mathrm{th}},\qquad g_k<1 .
\]
\item[(h)] \emph{The translation statements.} They provide a non-empty set $\mathfrak{P}$ of
translations $\tau$ of $\mathbb{T}^{(k)}$; $\operatorname{Av}_p$ is the average of
$\nu^{B}\circ\tau$ over $\tau\in\mathfrak{P}$.
\end{enumerate}

Then there are $\gamma_N=\gamma_N(N,L,M,\mathfrak{p})>0$, $s_*=s_*(N,L,M)$ and constants
$c_1,a_0,C_4,C_5$ with the following properties. The constants depend on $(N,L,M)$ and on
$\mathfrak{p}$ only, and not on $K$, $m$, $s$, $p$ or $g$. The conclusions below hold for
$0<g\le\gamma_N$, every admissible $K$, every $s$ with $s_*\le s\le K-1$, $k=K-s$, and every
plaquette $p$ of $\mathbb{T}^{(k)}$:
\begin{equation}\label{4.36:eq-plaquette-lower-bound-a}
\begin{aligned}
\text{(iii)}\quad
&\operatorname{Av}_p\bigl[\mathbf{1}\{O<a/x_k\}\bigr]
\le\vartheta^{B}(s)+C_4(N)C_5(N)\,a^{(N^2-1)/2}\\
&\qquad\le\tfrac18+C_4(N)C_5(N)\,a^{(N^2-1)/2}\qquad\text{for every }a>0;\\
&\text{with }a_0(N):=\Bigl(\frac{3}{8C_4(N)C_5(N)}\Bigr)^{2/(N^2-1)}\text{ one has}\\
&\operatorname{Av}_p\bigl[\mathbf{1}\{O\ge a_0g_k^2\}\bigr]\ge\tfrac12,\\
&\operatorname{Av}_p\Bigl[\mathbf{1}\Bigl\{O\ge\frac{a_0}{g^{-2}+B^{\sharp}(s+1)}\Bigr\}\Bigr]
\ge\tfrac12;\\
\text{(i)}\quad
&c_1(N)\,g_k^2\le\operatorname{Av}_p[O],\qquad c_1:=a_0/2;\\
\text{(iv)}\quad
&C_4(N):=c^{\flat}_3(N)^{-1}e^{(C'_W(N)+3)/2+1}
=(N-1)^{D/2}c_3(N)^{-1}e^{C'^{\mathrm{op}}_F(N)/2+1},\\
&C_5(N)=\sigma_N\,\omega_D\,\kappa_N^{-D/2}.
\end{aligned}
\end{equation}
The sharper constant $c^{\mathrm{op}}_3(N)^{-1}e^{C'^{\mathrm{op}}_F/2+1}$ is also admissible in
place of $C_4(N)$. The parameters $K,m,s,p,g$ enter no constant, and $k$ enters only through
$g_k$ and $R_k$.
\end{theorem}

\begin{proof}
Throughout, $du$ is Haar probability measure on $G=\mathrm{SU}(N)$.

Since the trace is normalised and the eigenvalues of an element of $\mathrm{SU}(N)$ have modulus
one, $\operatorname{Re}\operatorname{tr}A\le1$ for $A\in G$. Hence $O\ge0$.

\emph{Step~1: non-clean histories and the complement of the good set.} By hypothesis~(b) we
discard two contributions: the terms of the non-clean histories, and the factors
$1-\mathbf{1}_{G(h)}$ of the clean ones. After the $\mathfrak{P}$-average their mass is at most
$\vartheta^{B}(s)$, and $\vartheta^{B}(s)\le\frac18$ for $s\ge s_*$.

The group enters~(b) at one place only: where a large plaquette is estimated
against the action. There the left-hand inequality of the pair in
\eqref{4.36:eq-eigenangle-inequalities-a} comparing $1-\operatorname{Re}\operatorname{tr}A$ with
$|A-\mathbf{1}|^2$ gives, for $A\in G$ and $x\ge0$,
\[
e^{-x(1-\operatorname{Re}\operatorname{tr}A)}\le e^{-x|A-\mathbf{1}|^2/(2N)}.
\]
This is the replacement of $\gamma_0$ by $\gamma_0/N$, and of $c^{B}=\gamma_0(A_1/A_0)^2$ by
$c^{B}/N$, in~(b). It is the origin of the thresholds $\gamma^{B}$, $s^{B}_*$ of~(b) and of
$x_{\mathrm{th}}=x_{\mathrm{th}}(N,L,M)$ in~(g); each is a one-sided condition, and the form of
the bound is unchanged. Accordingly we choose $\gamma_N\le\min\{\gamma^{B},\gamma^{\mathrm{fl}}\}$
and $s_*\ge\max\{s^{B}_*,s^{\mathrm{fl}}_*\}$, so that hypotheses~(b) and~(g) are available
throughout the range of the
conclusions. It remains to treat, for each clean $h$, the contribution of
$\mathbf{1}_{G(h)}\,\mathrm{term}_h$.

\emph{Step~2: one clean history.} Fix a clean $h$ and $V_{\mathrm{off}}$. Traverse $\partial p$
starting from $b_1$, and let $M_0=M_0(V_{\mathrm{off}})\in G$ be the ordered product of the
remaining three bond variables.
\begin{itemize}
\item If $b_1$ is traversed forward, then $V_k(\partial p)=uM_0$.
\item If $b_1$ is traversed backward, then $V_k(\partial p)=u^{-1}M_0$. For $A\in G$ one has
$\operatorname{tr}A^{-1}=\overline{\operatorname{tr}A}$, and
$(u^{-1}M_0)^{-1}=M_0^{-1}u$. Therefore
$\operatorname{Re}\operatorname{tr}(u^{-1}M_0)=\operatorname{Re}\operatorname{tr}(uM_0^{-1})$.
\end{itemize}
In both cases
\[
O=1-\operatorname{Re}\operatorname{tr}(uM)\quad\text{with } M=M(V_{\mathrm{off}})\in G
\]
independent of $u$.

\emph{Normalisation of the one-link density.} We use the factorisation of hypothesis~(a) and
write $W=W(V_{\mathrm{off}})$. All integrands are non-negative and Borel, so Tonelli's theorem
applies throughout. Put
\[
\mathfrak{z}_h(V_{\mathrm{off}}):=\int_W Z_h\,du .
\]
Then
\[
\int dV_{\mathrm{off}}\,\Psi_h\mathfrak{z}_h=\int\mathrm{term}_h\,dV_k\le\int\rho_k<\infty .
\]
Hence for almost every $V_{\mathrm{off}}$ one of two cases occurs. Either
$\Psi_h\mathfrak{z}_h=0$, in which case $\mathbf{1}_{G(h)}\mathrm{term}_h$ vanishes for almost
every $u$ and there is nothing to prove. Or $0<\mathfrak{z}_h<\infty$. In the second case
$f:=\mathbf{1}_WZ_h/\mathfrak{z}_h$ is a Haar probability density, $Z_h<\infty$ almost everywhere
on $W$, and $Z_h>0$.

\emph{The supremum $S^{\sharp}$.} Put
\[
E:=\bigl\{u^0\in W(V_{\mathrm{off}}):\ \mathbf{1}_{G(h)}((u^0,V_{\mathrm{off}}))>0,\
Z_h(u^0)<\infty\bigr\},\qquad S^{\sharp}:=\sup_E f\in[0,\infty],
\]
with $\sup\emptyset:=0$. Then $\mathbf{1}_{G(h)}f\le S^{\sharp}$ almost everywhere. Indeed, $f$
vanishes off $W$, and $\mathbf{1}_{G(h)}$ vanishes off the good set. The remaining points lie in
$\{Z_h=\infty\}\cap W$, which is a null set.

\emph{The integral $I(V_{\mathrm{off}})$.} Fix $a>0$. By hypothesis~(a),
$\mathbf{1}_{G(h)}\mathrm{term}_h=\Psi_h\mathfrak{z}_h\,\mathbf{1}_{G(h)}f$. By Lemma~\ref{4.36:lem-haar-small-sets}~(c),
$u\mapsto uM$ preserves Haar measure, so
\[
\operatorname{Haar}\{u:1-\operatorname{Re}\operatorname{tr}(uM)<a/x_k\}
=\operatorname{Haar}\{g':1-\operatorname{Re}\operatorname{tr}g'<a/x_k\},
\]
and part~(b) of that lemma bounds the right-hand side by $C_5(N)(a/x_k)^{D/2}$. Together
with $\mathbf{1}_{G(h)}f\le S^{\sharp}$ this gives
\begin{equation}\label{4.36:eq-plaquette-lower-bound-1}
\begin{aligned}
I(V_{\mathrm{off}})&:=\int du\,\mathbf{1}\{O<a/x_k\}\,\mathbf{1}_{G(h)}\,\mathrm{term}_h\\
&=\Psi_h\mathfrak{z}_h\int du\,\mathbf{1}\{1-\operatorname{Re}\operatorname{tr}(uM)<a/x_k\}\,
\mathbf{1}_{G(h)}f\\
&\le\Psi_h\mathfrak{z}_h\,S^{\sharp}\,
\operatorname{Haar}\{g':1-\operatorname{Re}\operatorname{tr}g'<a/x_k\}\\
&\le\Psi_h\mathfrak{z}_h\,S^{\sharp}\,C_5(N)\,(a/x_k)^{D/2}.
\end{aligned}
\end{equation}

\emph{Bound on $f(u^0)$ for $u^0\in E$: the data of the mode bound.} Let $u^0\in E$.
\begin{itemize}
\item By hypothesis~(c), $(u^0,V_{\mathrm{off}})\in\mathcal{D}$ and the localised base-point
estimate of Lemma~\ref{4.36:lem-broken-paths} holds at $(u^0,V_{\mathrm{off}})$.
\item By clauses (b) and (c) of Lemma~\ref{4.36:lem-broken-paths}, whose inputs are supplied by
hypothesis~(c),
$B^{\mathrm{op}}_{g_k}(u^0)\subset B^{\mathrm{op}}_{\varrho}(u^0)\subset W$, where $\varrho$ is
the radius of that lemma. The first inclusion holds because $\varrho\ge2g_k$ under the threshold
condition (t1) of \eqref{4.36:eq-flow-thresholds-b}, which holds since $x_k\ge x_{\mathrm{th}}$
(hypothesis~(g)). The segments of clause (c) of that lemma issuing from points of this ball are
windowed in the sense of that clause, whose threshold conditions (t1) and (t2) are met in the
same way.
\item Clause (a) of the one-parameter bounds of Lemma~\ref{4.36:lem-one-link-curvature}, whose
inputs are supplied by hypotheses (c)--(f), holds at every
$u\in B^{\mathrm{op}}_{g_k}(u^0)$ and for every $\hat{T}\in\mathfrak{g}$ with $|\hat{T}|=1$.
Along $\hat{T}$ the map $t\mapsto F^{\mathrm{near}}$ is real-analytic near $0$, hence $C^2$
there, with second derivative at $t=0$ bounded by $C'^{\mathrm{op}}_F(N)\,x_k$. This is
hypothesis~(m1) of \eqref{4.36:eq-mode-bound-a} with
\[
F_1:=F^{\mathrm{near}},\qquad \Lambda:=C'^{\mathrm{op}}_F x_k .
\]
\item Clause (b) of the one-parameter bounds of the same lemma holds at the base point $u^0$. It
gives hypothesis~(m2) with
\[
F_2:=F^{\mathrm{far}},\qquad \eta_2:=\eta_{\mathrm{far}}\le1 .
\]
Indeed, every $u\in B^{\mathrm{op}}_{g_k}(u^0)\setminus\{u^0\}$ has the form $u^0e^{iv}$ with
$v\in\mathfrak{g}$ and $0<|v|<g_k$. So $u=u^0e^{it\hat{T}}$ with
\[
\hat{T}:=v/|v|,\qquad t:=|v|\in(0,g_k)\subset[-2g_k,2g_k] .
\]
At $u=u^0$ the bound is trivial. This is the only point at which hypothesis~(e), the bound on
$F^{\mathrm{far}}$ along one link, enters. It is
applied for $t=|v|<g_k$ measured in the operator norm, and its bound holds whatever the
normalisation of the direction.
\end{itemize}

\emph{Application of the mode bound in the operator norm.} We apply
Lemma~\ref{4.36:lem-mode-bound} in the operator-norm case with the following data: the open set
$W(V_{\mathrm{off}})$; the $\sigma$-finite measure space $(\mathcal{W},m^{V_{\mathrm{off}}}_h)$;
the function $F=F_h=F^{\mathrm{near}}+F^{\mathrm{far}}$ and $Z=Z_h$. Then, with $\Phi=\log Z_h$ as
in that lemma,
\[
\int_We^{\Phi}\,du=\int_W Z_h\,du=\mathfrak{z}_h\in(0,\infty),
\]
the density of that lemma is $f$,
and $Z_h(u^0)<\infty$ because $u^0\in E$.

Take $r:=g_k$. By hypothesis~(g), $g_k<1<\pi/2$, so $r$ is admissible, and
\[
\Lambda r^2=C'^{\mathrm{op}}_F\,x_k\,g_k^2=C'^{\mathrm{op}}_F .
\]
Since $\eta_2\le1$ and $g_k^{-D}=x_k^{D/2}$, the lemma gives
\begin{equation}\label{4.36:eq-plaquette-lower-bound-2}
f(u^0)\le c^{\flat}_3(N)^{-1}e^{C'^{\mathrm{op}}_F/2+1}\,g_k^{-D}=C_4(N)\,x_k^{D/2}.
\end{equation}
Here $C'^{\mathrm{op}}_F=C'_W+3$ (Lemma~\ref{4.36:lem-one-link-curvature}), which gives the first
expression for $C_4(N)$ in \eqref{4.36:eq-plaquette-lower-bound-a}.
Lemma~\ref{4.36:lem-mode-bound} also gives the bound with $c^{\mathrm{op}}_3(N)^{-1}$ in place of
$c^{\flat}_3(N)^{-1}$ (Lemma~\ref{4.36:lem-haar-small-sets}~(a$^{\mathrm{op}}$)); this is the
sharper admissible constant.

\emph{The same bound in the norm $|\cdot|_{\mathfrak{g}}$, and where the factor $\mathsf{n}^2$
goes.} Now apply Lemma~\ref{4.36:lem-mode-bound} in the $\mathfrak{g}$-case, on the ball
$B^{\mathfrak{g}}_{g_k/\mathsf{n}}(u^0)$.
\begin{itemize}
\item \emph{Inclusion.} By the inequality $|X|\le\mathsf{n}|X|_{\mathfrak{g}}$ of
\eqref{4.36:eq-eigenangle-inequalities-a} (equivalently,
Lemma~\ref{4.36:lem-haar-small-sets}~(d)), $|v|\le\mathsf{n}|v|_{\mathfrak{g}}$. Hence
\[
B^{\mathfrak{g}}_{g_k/\mathsf{n}}(u^0)\subset B^{\mathrm{op}}_{g_k}(u^0)\subset W .
\]
\item \emph{Radius.} The radius $g_k/\mathsf{n}$ is admissible: $g_k/\mathsf{n}\le\pi/\sqrt{2N}$
follows from $g_k\le\pi/2$.
\item \emph{Hypothesis (m1).} Along directions $T$ with $|T|_{\mathfrak{g}}=1$,
clause~$(\mathfrak{g})$(a) of Lemma~\ref{4.36:lem-one-link-curvature} gives (m1) with
\[
\Lambda=C'^{\mathfrak{g}}_F x_k=(N-1)(C'_W+3)\,x_k .
\]
\item \emph{Hypothesis (m2).} This holds as before, since the ball lies in
$B^{\mathrm{op}}_{g_k}(u^0)$.
\end{itemize}
With $r=g_k/\mathsf{n}$,
\[
\Lambda r^2=(N-1)(C'_W+3)\,x_k\cdot\frac{g_k^2}{N-1}=C'_W+3 .
\]
Thus the factor $\mathsf{n}^2$ of $C'^{\mathfrak{g}}_F$ is cancelled by the radius. By
Lemma~\ref{4.36:lem-haar-small-sets}~(a),
\[
\operatorname{Haar}\bigl(B^{\mathfrak{g}}_{g_k/\mathsf{n}}\bigr)
\ge c_3(N)\,(g_k/\mathsf{n})^{D}=(N-1)^{-D/2}c_3(N)\,x_k^{-D/2}.
\]
Therefore
\[
f(u^0)\le(N-1)^{D/2}c_3(N)^{-1}e^{(C'_W+3)/2+1}\,x_k^{D/2}.
\]
The power of $\mathsf{n}$ enters $C_4$ as the volume factor $\mathsf{n}^D=(N-1)^{D/2}$. The result
is the second expression for $C_4(N)$ in \eqref{4.36:eq-plaquette-lower-bound-a}, identical to
the constant in \eqref{4.36:eq-plaquette-lower-bound-2}.

\emph{Step~3: anti-concentration.} By \eqref{4.36:eq-plaquette-lower-bound-2},
$S^{\sharp}\le C_4x_k^{D/2}$. Inserting this in \eqref{4.36:eq-plaquette-lower-bound-1} gives
the computation of Lemma~\ref{4.36:lem-anti-concentration}, in which the powers of $x_k$ cancel:
\begin{align*}
I(V_{\mathrm{off}})&\le\Psi_h\mathfrak{z}_h\,C_4x_k^{D/2}\,C_5(a/x_k)^{D/2}\\
&=C_4C_5(N)\,a^{D/2}\int du\,\mathrm{term}_h(u,V_{\mathrm{off}}).
\end{align*}
The equality uses $\Psi_h\mathfrak{z}_h=\int du\,\mathrm{term}_h(u,V_{\mathrm{off}})$, which
follows from hypothesis~(a). When $\Psi_h\mathfrak{z}_h=0$, both sides vanish.

We now integrate over $V_{\mathrm{off}}$, sum over clean $h$, divide by $\mathcal{Z}$, and take
the $\mathfrak{P}$-average. The bound is uniform in the plaquette, and the averaged measure has
total mass at most one, so the bound survives these operations. Adding the contribution
discarded in Step~1 gives the first line of~(iii). The second line follows from
$\vartheta^{B}(s)\le\frac18$ for $s\ge s_*$.

\emph{Step~4: the threshold $a_0$.} By the definition of $a_0$,
\[
C_4C_5\,a_0^{D/2}=C_4C_5\cdot\frac{3}{8C_4C_5}=\frac38 ,
\qquad\text{and}\qquad \frac18+\frac38=\frac12 .
\]
Next, $a_0<1\le x_k$. To see $a_0<1$, write
\[
C_4C_5=(N-1)^{D/2}\kappa_N^{-D/2}\cdot\sigma_N\,\omega_D\,c_3(N)^{-1}
\cdot e^{(C'_W+5)/2}.
\]
We bound the three factors in turn.
\begin{itemize}
\item Since $\kappa_N\le1/(2(N-1)^2)$,
\[
(N-1)^{D/2}\kappa_N^{-D/2}\ge(2(N-1)^3)^{D/2}\ge2^{D/2}.
\]
\item With $c_3(N)$ as fixed in Lemma~\ref{4.36:lem-haar-small-sets},
\[
\sigma_N\,\omega_D\,c_3(N)^{-1}=(\pi/2)^{N^2-N}\ge1 .
\]
\item $e^{(C'_W+5)/2}\ge e$.
\end{itemize}
Hence $C_4C_5\ge2^{D/2}e>3/8$, so $a_0<1$.

Each measure $\nu^{B}\circ\tau$ is a probability, so $\operatorname{Av}_p$ is an average of
probabilities. Then (iii) with $a=a_0$ gives
\[
\operatorname{Av}_p\bigl[\mathbf{1}\{O\ge a_0/x_k\}\bigr]
=1-\operatorname{Av}_p\bigl[\mathbf{1}\{O<a_0/x_k\}\bigr]\ge\frac12,
\]
and $a_0/x_k=a_0g_k^2$.

\emph{Step~5: the lower bound on the mean.} Since $O\ge0$, Markov's inequality and Step~4 give
\[
\operatorname{Av}_p[O]\ge\frac{a_0}{x_k}\,\operatorname{Av}_p\bigl[\mathbf{1}\{O\ge a_0/x_k\}\bigr]
\ge\frac12\,a_0\,g_k^2=c_1\,g_k^2 .
\]
This is~(i).

\emph{Step~6: the bound in terms of $g$ and $s$.} By hypothesis~(g) at the anchor
$g_K\in[g_{-}(g),g[$,
\[
x_k\le g^{-2}+B^{\sharp}(s+1).
\]
Hence $a_0/(g^{-2}+B^{\sharp}(s+1))\le a_0/x_k$, and the event $\{O\ge a_0/x_k\}$ is contained in
$\{O\ge a_0/(g^{-2}+B^{\sharp}(s+1))\}$. Step~4 then gives the last inequality of~(iii).

\emph{Dependence of the constants.} Of the flow, the argument uses only hypothesis~(g), in
particular $x_k\ge x_{\mathrm{th}}$ and $x_k\le g^{-2}+B^{\sharp}(s+1)$. The parameters $K,m,s,p$
enter no constant: every sum
in Lemma~\ref{4.36:lem-one-link-curvature} is weighted as in \eqref{4.36:eq-flow-thresholds-c},
and Lemmas~\ref{4.36:lem-haar-small-sets}, \ref{4.36:lem-mode-bound}
and~\ref{4.36:lem-anti-concentration} concern a single link.
\end{proof}

\begin{remark}[The mode bound on the $\mathfrak{g}$-ball: a weaker constant]
\label{4.36:rem-mode-bound-lie-ball}
Write $\mathsf{n}:=(N-1)^{1/2}$ and $D:=N^2-1$. Work in the setting of the bound on the density
$f$ in the proof of Theorem~\ref{4.36:thm-plaquette-lower-bound}: a clean history $h$, the
configuration $V_{\mathrm{off}}$ off the free link fixed, $f=\mathbf{1}_{W}Z_h/\mathfrak{z}_h$
with $W=W(V_{\mathrm{off}})$, and $u^0$ in the set $E$ of points of $W(V_{\mathrm{off}})$ at
which $\mathbf{1}_{G(h)}((u^0,V_{\mathrm{off}}))>0$ and $Z_h(u^0)<\infty$. In that proof
Lemma~\ref{4.36:lem-mode-bound} is applied in the operator-norm case on
$B^{\mathrm{op}}_{g_k}(u^0)$.
It may instead be applied in the $\mathfrak{g}$-norm case on $B^{\mathfrak{g}}_{g_k}(u^0)$, with
$\Lambda=C'^{\mathfrak{g}}_{F}x_k$, $C'^{\mathfrak{g}}_{F}=(N-1)(C'_W(N)+3)$, and $r=g_k$,
provided the following hold:
\begin{itemize}
\item[(a)] $g_k\le \pi/\sqrt{2N}$;
\item[(b)] $p_0(g_k)\ge 8C_{\mathrm{loc}}\mathsf{n}$ for the degree function $p_0(\cdot)$, where
$C_{\mathrm{loc}}$ is the constant of the window radius $\varrho=\varepsilon_k/(4C_{\mathrm{loc}})$,
so that the inclusion of balls in the window gives
$B^{\mathrm{op}}_{\mathsf{n}g_k}(u^0)\subset W(V_{\mathrm{off}})$;
\item[(c)] the radius and far-size conditions under which Lemma~\ref{4.36:lem-one-link-curvature}
is stated hold at reach $\mathsf{n}g_k$ in place of $g_k$, that is, $\mathsf{n}g_k$ is smaller
than the radius bound in those conditions and $\mathsf{n}\eta_{\mathrm{far}}\le 1$;
\item[(d)] the bound on the far members of the conditional exponent that enters
Lemma~\ref{4.36:lem-one-link-curvature} as an antecedent, whose proof is outstanding and which is
stated for $|t|\le 2g_k$, holds at operator reach $|t|<\mathsf{n}g_k$.
\end{itemize}
Then
\begin{equation}\label{4.36:eq-mode-bound-lie-ball-1}
f(u^0)\le C^{\mathfrak{g}}_4(N)\,x_k^{D/2},\qquad
C^{\mathfrak{g}}_4(N):=c_3(N)^{-1}e^{(N-1)(C'_W(N)+3)/2+1},
\end{equation}
which is a bound of the form used for $f(u^0)$ in the proof of
Theorem~\ref{4.36:thm-plaquette-lower-bound}, with $C^{\mathfrak{g}}_4(N)$ in place of $C_4(N)$.
Hypothesis (d) lies inside the range $|t|\le 2g_k$ of the antecedent as stated if and only if
$\mathsf{n}\le 2$, that is, if and only if $N\le 5$; for $N\ge 6$ it asks for the antecedent on a
larger range than the one on which it is stated, and in either case the antecedent's proof is
outstanding. Moreover
\begin{equation}\label{4.36:eq-mode-bound-lie-ball-2}
\frac{C^{\mathfrak{g}}_4(N)}{C_4(N)}=\frac{e^{(N-2)(C'_W(N)+3)/2}}{(N-1)^{D/2}},
\end{equation}
so that, for $N\ge 3$, $C^{\mathfrak{g}}_4(N)>C_4(N)$ if and only if
\begin{equation}\label{4.36:eq-mode-bound-lie-ball-3}
(N-2)\bigl(C'_W(N)+3\bigr)>(N^2-1)\log(N-1),
\end{equation}
equivalently $C'_W(N)+3>(N+1)(N-1)\log(N-1)/(N-2)$, while at $N=2$ the two constants are equal.
For the constant $C'_W(N)$ of this chapter, whose size is set by $L$ and $M$, while the right side
of \eqref{4.36:eq-mode-bound-lie-ball-3} depends on $N$ alone, the inequality
\eqref{4.36:eq-mode-bound-lie-ball-3} holds at every moderate $N$, and there
$C^{\mathfrak{g}}_4(N)$ is the larger constant.
\end{remark}

\begin{proof}
\emph{Radius.} By (a), $r=g_k\le\pi/\sqrt{2N}$, which is the admissible range of the radius in the
$\mathfrak{g}$-norm case of Lemma~\ref{4.36:lem-mode-bound}.

\emph{Inclusion of balls.} Let $v\in\mathfrak{g}$ have eigenvalues $\lambda_1,\dots,\lambda_N$, so
$\sum_i\lambda_i=0$ and $|v|_{\mathfrak{g}}^2=\operatorname{tr}v^2=N^{-1}\sum_i\lambda_i^2$. By the
Cauchy--Schwarz inequality,
\begin{equation*}
\lambda_1^2=\Bigl(\sum_{i\ge 2}\lambda_i\Bigr)^2\le (N-1)\sum_{i\ge 2}\lambda_i^2,
\end{equation*}
hence $N\lambda_1^2\le (N-1)\sum_i\lambda_i^2$, i.e.\ $\lambda_1^2\le (N-1)|v|_{\mathfrak{g}}^2$; the
same holds for every $\lambda_i$, so $|v|\le\mathsf{n}|v|_{\mathfrak{g}}$. Therefore
$B^{\mathfrak{g}}_{g_k}(u^0)\subset B^{\mathrm{op}}_{\mathsf{n}g_k}(u^0)$, and by (b) this ball lies
in $W(V_{\mathrm{off}})$.

\emph{Hypothesis (m1).} Every $u\in B^{\mathfrak{g}}_{g_k}(u^0)$ lies in
$B^{\mathrm{op}}_{\mathsf{n}g_k}(u^0)$. By (c), clause (a) of
Lemma~\ref{4.36:lem-one-link-curvature}, read along directions of unit $|\cdot|_{\mathfrak{g}}$-norm,
applies at every such $u$ with constant $C'^{\mathfrak{g}}_{F}=(N-1)(C'_W(N)+3)$; this is (m1) of
\eqref{4.36:eq-mode-bound-a} with $F_1:=F^{\mathrm{near}}$ and $\Lambda=C'^{\mathfrak{g}}_{F}x_k$.

\emph{Hypothesis (m2).} Every $u\in B^{\mathfrak{g}}_{g_k}(u^0)\setminus\{u^0\}$ is
$u^0e^{it\hat{T}}$ with $\hat{T}:=v/|v|$ of unit operator norm and $t:=|v|$, and
$t\le\mathsf{n}|v|_{\mathfrak{g}}<\mathsf{n}g_k$. By (d), together with (c), the far bound applies
at this reach at the base $u^0$, which gives (m2) of \eqref{4.36:eq-mode-bound-a} with
$F_2:=F^{\mathrm{far}}$ and $\eta_2:=\mathsf{n}\eta_{\mathrm{far}}$, which is at most $1$ by (c).

\emph{The bound.} Since $x_kg_k^2=1$,
\begin{equation*}
\Lambda r^2=(N-1)\bigl(C'_W(N)+3\bigr)x_kg_k^2=(N-1)\bigl(C'_W(N)+3\bigr).
\end{equation*}
Lemma~\ref{4.36:lem-mode-bound} in the $\mathfrak{g}$-norm case, with the Haar small-ball bound
$\operatorname{Haar}(B^{\mathfrak{g}}_{g_k}(u^0))\ge c_3(N)g_k^{D}$ for $\mathfrak{g}$-balls, valid
for radii at most $\pi/\sqrt{2N}$ and hence at $r=g_k$ by (a), gives, exactly as in the proof
of Theorem~\ref{4.36:thm-plaquette-lower-bound} where the lemma is read on
$B^{\mathfrak{g}}_{g_k/\mathsf{n}}(u^0)$,
\begin{equation*}
f(u^0)\le c_3(N)^{-1}e^{\Lambda r^2/2+1}g_k^{-D}=C^{\mathfrak{g}}_4(N)\,x_k^{D/2},
\end{equation*}
using $g_k^{-D}=x_k^{D/2}$. This is \eqref{4.36:eq-mode-bound-lie-ball-1}.

\emph{Range of (d).} Since $\mathsf{n}=(N-1)^{1/2}$, $\mathsf{n}\le 2$ if and only if
$N-1\le 4$, i.e.\ $N\le 5$; exactly then $|t|<\mathsf{n}g_k$ implies $|t|\le 2g_k$.

\emph{Comparison.} By clause (iv) of Theorem~\ref{4.36:thm-plaquette-lower-bound}
(see \eqref{4.36:eq-plaquette-lower-bound-a}) and $C'^{\mathrm{op}}_{F}(N)=C'_W(N)+3$ from
Lemma~\ref{4.36:lem-one-link-curvature},
$C_4(N)=(N-1)^{D/2}c_3(N)^{-1}e^{(C'_W(N)+3)/2+1}$. Dividing,
\begin{equation*}
\frac{C^{\mathfrak{g}}_4(N)}{C_4(N)}
=\frac{e^{(N-1)(C'_W(N)+3)/2-(C'_W(N)+3)/2}}{(N-1)^{D/2}}
=\frac{e^{(N-2)(C'_W(N)+3)/2}}{(N-1)^{D/2}},
\end{equation*}
which is \eqref{4.36:eq-mode-bound-lie-ball-2}. Its logarithm is
$\tfrac12\bigl[(N-2)(C'_W(N)+3)-(N^2-1)\log(N-1)\bigr]$, which is positive if and only if
\eqref{4.36:eq-mode-bound-lie-ball-3} holds; for $N\ge 3$ dividing by $N-2>0$ and using
$N^2-1=(N+1)(N-1)$ gives the equivalent form. At $N=2$ both $N-2$ and $\log(N-1)$ vanish, so
the ratio is $1$.
\end{proof}

\begin{corollary}[The upper plaquette constant and the values of the lower constants]
\label{4.36:cor-plaquette-constants}
Work at the gauge group $SU(N)$, $N\ge 2$, in the setting of the lower and anti-concentration
bounds for the averaged plaquette (Theorem~\ref{4.36:thm-plaquette-lower-bound}), with its
coupling ceiling $\gamma_N$, its depth $s_*$, its constants $C_4(N)$, $C_5(N)$, the dimension
$D=N^2-1$ of $\mathfrak{g}$, the plaquette variable $O=1-\operatorname{Re}\operatorname{tr}V_k(\partial p)$,
the translation average $\operatorname{Av}_p$, and $x_k=g_k^{-2}$. Let $\varepsilon_k$ be the
small-field radius at level $k$, $p_0(\cdot)$ the degree function, $\gamma_0$ and $c^{B}$ the
constants of the large-field estimates entering $\gamma_N$, and $B_3$ the constant $B_3$ of the
last-level lower bound in (d). Assume the following.
\begin{itemize}
\item[(a)] \emph{The upper bound for the averaged plaquette.} There is a depth term
$\delta^{B}(s)$ with the following property. If the function $u\mapsto 1-\operatorname{Re}\operatorname{tr}u$
on $SU(N)$ is less than $2\varepsilon_k^2$ on the operator-norm window
$|u-\mathbf{1}|<2\varepsilon_k$ and at most $2$ off that window, then, for $0<g\le\gamma_N$ and
for the range of depths determined by the depth function $s_{\mathrm{up}}$ of (b),
\begin{equation*}
\operatorname{Av}_p[O]\le C(N)\,g_k^2\,p_0(g_k)^2+2\delta^{B}(s)
\end{equation*}
holds with $C(N)=2$; the estimate of the complement of the window enters through the factor
$\gamma_0/N$ in the coupling ceiling $\gamma_N$.
\item[(b)] \emph{The depth of the upper bound.} The depth function of the upper bound in (a) is
$s_{\mathrm{up}}(g)=\tfrac14\log_L g^{-2}+C_{\mathrm{up}}$, where
$C_{\mathrm{up}}=C_{\mathrm{up}}(N,L,M)$ is proportional to $N/c^{B}$; and the set $\Xi_K$ of
level-$K$ terms of that depth bookkeeping contains the term $2(N^2-1)\log x_K$.
\item[(c)] \emph{Two constants at $SU(N)$.} The constant $a^{\mu}_0(N)$ has the value
$a^{\mu}_0(N)=\tfrac12\,(4C_4(N)C_5(N))^{-2/D}$, and the constant $a_0(N)$ of
\eqref{4.36:eq-plaquette-constants-1} below satisfies $a_0(N)<1$.
\item[(d)] \emph{The bookkeeping of the last-level lower bound.} With
$r_K:=g_K/(5B_3)$, the last-level lower bound holds with its constant $C'_{\mathrm{low}}$ equal
to $E_{-}(N,L,\mathfrak{p})+4\mathfrak{V}+3$, for every number $\mathfrak{V}$ such that
$\operatorname{Haar}\{u\in SU(N):|u-\mathbf{1}|<r_K\}\ge e^{-\mathfrak{V}}$.
\end{itemize}
Then:
\begin{enumerate}
\item the bound in (a) holds at every $N\ge 2$ with $C(N)=2$, independent of $N$;
\item $C_{\mathrm{up}}$ is a constant fixed after $(N,L,M)$ and before $\gamma$, and $\Xi_K$
contains $2(N^2-1)\log x_K$;
\item $c_1(N)=a_0(N)/2$, where
\begin{equation}\label{4.36:eq-plaquette-constants-1}
a_0(N)=\Bigl(\frac{3}{8C_4(N)C_5(N)}\Bigr)^{2/D},
\qquad
C_4(N)C_5(N)\,a_0(N)^{D/2}=\frac38 ,
\end{equation}
and $a_0(N)<1$;
\item $a^{\mu}_0(N)=\tfrac12\,(4C_4(N)C_5(N))^{-2/D}$;
\item under the coupling ceiling $g_K\le 5\pi B_3/2$, the last-level lower bound holds
with the constant
\begin{equation}\label{4.36:eq-plaquette-constants-2}
C'_{\mathrm{low}}(N):=E_{-}(N,L,\mathfrak{p})+2(N^2-1)\log x_K+c_H(N),
\end{equation}
where
\begin{equation}\label{4.36:eq-plaquette-constants-3}
c_H(N):=4\log c_3(N)^{-1}+4(N^2-1)\log(5B_3\mathsf{n})+3 .
\end{equation}
\end{enumerate}
\end{corollary}

\begin{proof}
\emph{Item (1).} By hypothesis (a) it suffices to see that the two pointwise bounds named there
hold at every $N$. Let $u\in SU(N)$ and $O(u)=1-\operatorname{Re}\operatorname{tr}u$.
If $|u-\mathbf{1}|<2\varepsilon_k$, the right-hand inequality of item (F7) in
\eqref{4.36:eq-eigenangle-inequalities-a} (Lemma~\ref{4.36:lem-eigenangle-inequalities}) gives
\begin{equation*}
O(u)\le\tfrac12\,|u-\mathbf{1}|^2<\tfrac12\,(2\varepsilon_k)^2=2\varepsilon_k^2 ;
\end{equation*}
only this half of (F7) is used. For every $u$, the eigenvalues of $u$ are $e^{i\theta_q}$,
$q=1,\dots,N$, and since $\operatorname{tr}$ is normalised,
$\operatorname{Re}\operatorname{tr}u=N^{-1}\sum_q\cos\theta_q\ge-1$, so $O(u)\le 2$. Neither bound
involves $N$. The estimate of the complement of the window is accounted for by the factor
$\gamma_0/N$ entering the coupling ceiling $\gamma_N$ of
Theorem~\ref{4.36:thm-plaquette-lower-bound}. Hence hypothesis (a) gives the bound with
$C(N)=2$ at every $N$. In the statement the constant is kept as the
symbol $C(N)$; the value $2$ is admissible for it.

\emph{Item (2).} By hypothesis (b), $C_{\mathrm{up}}$ depends on $(N,L,M)$ only, through a
quantity proportional to $N/c^{B}$; it is therefore determined once $(N,L,M)$ are fixed and is
fixed before $\gamma$ in the order of choice. By the same hypothesis, $\Xi_K$ contains the term
$2(N^2-1)\log x_K$.

\emph{Item (3).} Theorem~\ref{4.36:thm-plaquette-lower-bound}(iii) gives, for every $a>0$,
\begin{equation*}
\operatorname{Av}_p\bigl[\mathbf{1}\{O<a/x_k\}\bigr]
\le\tfrac18+C_4(N)C_5(N)\,a^{(N^2-1)/2},
\end{equation*}
and the exponent is $(N^2-1)/2=D/2$. With $a=a_0(N)$ as in \eqref{4.36:eq-plaquette-constants-1},
$a_0(N)^{D/2}=3/(8C_4(N)C_5(N))$, hence $C_4(N)C_5(N)\,a_0(N)^{D/2}=3/8$, which is the identity in
\eqref{4.36:eq-plaquette-constants-1}; thus
$\operatorname{Av}_p[\mathbf{1}\{O<a_0/x_k\}]\le\tfrac18+\tfrac38=\tfrac12$. This cancellation
holds without any condition beyond those of Theorem~\ref{4.36:thm-plaquette-lower-bound}. The
inequality $a_0(N)<1$, which is equivalent to $8C_4(N)C_5(N)>3$, is part of hypothesis (c).
Finally, Theorem~\ref{4.36:thm-plaquette-lower-bound}(i) holds with $c_1:=a_0/2$, so
$c_1(N)=a_0(N)/2$.

\emph{Item (4).} This restates hypothesis (c) for reference.

\emph{Item (5).} Put $r_K=g_K/(5B_3)$. The ceiling $g_K\le 5\pi B_3/2$ is equivalent to
$r_K\le\pi/2$. Lemma~\ref{4.36:lem-haar-small-sets}(d) gives the inclusion of
$\{v:|v|<r_K\}$ in $\{v:|e^{iv}-\mathbf{1}|<r_K\}$, hence
$\operatorname{Haar}\{u:|u-\mathbf{1}|<r_K\}\ge\operatorname{Haar}(B^{\mathrm{op}}_{r_K})$. Since
$r_K\le\pi/2$, clause (a$^{\mathrm{op}}$) of Lemma~\ref{4.36:lem-haar-small-sets} applies at radius
$r_K$, and in the form derived there from clause (a) at radius $r_K/\mathsf{n}$ it gives
\begin{equation*}
\operatorname{Haar}(B^{\mathrm{op}}_{r_K})\ge c^{\flat}_3(N)\,r_K^D=c_3(N)\,(r_K/\mathsf{n})^D .
\end{equation*}
Taking logarithms and using $r_K/\mathsf{n}=g_K/(5B_3\mathsf{n})$ and $\log g_K^{-1}=\tfrac12\log x_K$,
\begin{equation*}
\log\bigl(\operatorname{Haar}\{u:|u-\mathbf{1}|<r_K\}\bigr)^{-1}
\le\log c_3(N)^{-1}+D\log(5B_3\mathsf{n})+\tfrac{D}{2}\log x_K .
\end{equation*}
Hence hypothesis (d) applies with $\mathfrak{V}$ equal to the right-hand side, and
\begin{align*}
E_{-}(N,L,\mathfrak{p})+4\mathfrak{V}+3
&=E_{-}(N,L,\mathfrak{p})+2D\log x_K+4\log c_3(N)^{-1}+4D\log(5B_3\mathsf{n})+3\\
&=E_{-}(N,L,\mathfrak{p})+2(N^2-1)\log x_K+c_H(N),
\end{align*}
which is \eqref{4.36:eq-plaquette-constants-2} with $c_H(N)$ as in
\eqref{4.36:eq-plaquette-constants-3}.
\end{proof}

\begin{remark}[Table of the $SU(N)$ constants beside their $N=2$ values]
\label{4.36:rem-constants-table}
Throughout, $D=N^2-1=\dim\mathfrak{g}$, $\mathsf{n}=(N-1)^{1/2}$ and $x_k=g_k^{-2}$. By
Lemma~\ref{4.36:lem-haar-small-sets}\,(d), for every $u\in SU(N)$,
\begin{equation*}
B^{\mathrm{op}}_{g_k}(u)\supseteq B^{\mathfrak{g}}_{g_k/\mathsf{n}}(u).
\end{equation*}
In the proof of Theorem~\ref{4.36:thm-plaquette-lower-bound} the mode ball is the operator ball
$B^{\mathrm{op}}_{g_k}(u^0)$ about a point $u^0$ of the window $W(V_{\mathrm{off}})$ of the free
link variable. We call this the \emph{operator-ball route}. The variant displayed in
Remark~\ref{4.36:rem-mode-bound-lie-ball} takes the mode ball $B^{\mathfrak{g}}_{g_k}(u^0)$, and we
call it the \emph{$\mathfrak{g}$-ball variant}. Each item below gives a quantity first along the
operator-ball route, then along the $\mathfrak{g}$-ball variant where the two differ, and then its
value at $N=2$. At $N=2$ the operator-ball route and the $\mathfrak{g}$-ball variant give the same
values.
\begin{itemize}
\item \emph{The plaquette observable.} For a plaquette $p$, with $\|\cdot\|$ the normalised trace
norm,
\begin{equation*}
\begin{split}
O:={}&1-\operatorname{Re}\operatorname{tr}V_k(\partial p)
=\tfrac12\|V_k(\partial p)-\mathbf{1}\|^2\\
={}&(2N)^{-1}\operatorname{Tr}\bigl[(V_k(\partial p)-\mathbf{1})^*(V_k(\partial p)-\mathbf{1})\bigr]
\in[0,2].
\end{split}
\end{equation*}
The window condition $O<2\varepsilon_k^2$, with $\varepsilon_k$ the small-field threshold of
level $k$, keeps the form it has at $N=2$, with no $N$-dependent
constant. The $\mathfrak{g}$-ball variant has the same
entries. At $N=2$ we have $O=1-\frac12\operatorname{Tr}V_k(\partial p)$.
\item \emph{The complexified plaquette action.} The action is
\begin{equation*}
\mathfrak{w}(W)=1-\tfrac12(\operatorname{tr}W+\operatorname{tr}W^{-1})
=-\tfrac12\operatorname{tr}\bigl[W^{-1}(W-\mathbf{1})^2\bigr].
\end{equation*}
The smallness of $W-\mathbf{1}$ is measured in $|\cdot|$, and the constant $C_{\mathfrak{w}}$ in the
smallness bound for $\mathfrak{w}$ of the lemma on the complexified plaquette action equals $1$.
These entries are the same in the variant and at $N=2$.
\item \emph{Haar measure of small balls.} On $B^{\mathfrak{g}}_r$ with $0<r\le\pi/\sqrt{2N}$ the
constant has the same form for every $N$:
\begin{equation*}
c_3(N)=\sigma_N\omega_D(2/\pi)^{N^2-N}.
\end{equation*}
The mode-ball constant is $c^{\flat}_3(N)=(N-1)^{-D/2}c_3(N)$, for radii $r\le\pi/2$. The sharper
operator-ball constant is
\begin{equation*}
c^{\mathrm{op}}_3(N)=\sigma_N\omega^{\mathrm{op}}_N(2/\pi)^{N^2-N}.
\end{equation*}
The variant uses $c_3(N)$ with $r\le\pi/\sqrt{2N}$. At $N=2$, $c_3(2)=8/(3\pi^3)$.
\item \emph{Sublevel sets of the action.} The constant in the bound
\begin{equation*}
\operatorname{Haar}\{u\in SU(N):1-\operatorname{Re}\operatorname{tr}u<\tau\}\le C_5(N)\tau^{D/2}
\end{equation*}
of Lemma~\ref{4.36:lem-haar-small-sets}\,(b), valid for every $\tau>0$, has the same form for
every $N$ and is the same in the variant:
\begin{equation*}
C_5(N)=\sigma_N\omega_D\kappa_N^{-D/2},\qquad
\kappa_N=\frac{N^2\sin^2(\pi/N)}{2\pi^2(N-1)^2}.
\end{equation*}
At $N=2$, $C_5(2)=\pi^2/(3\sqrt2)$ and $\kappa_2=2/\pi^2$.
\item \emph{The curvature constant.} In Lemma~\ref{4.36:lem-one-link-curvature} the constant is
$C'^{\mathrm{op}}_F(N)=C'_W(N)+3$ along an operator-unit direction $\hat T$. Along a
$\mathfrak{g}$-unit direction $T$ it is $C'^{\mathfrak{g}}_F(N)=(N-1)(C'_W(N)+3)$, which carries the
factor $\mathsf{n}^2$. In the bound on the normalised one-link density used in the proof of
Theorem~\ref{4.36:thm-plaquette-lower-bound}, $\Lambda$ is the bound on the second derivative of the
near part $F^{\mathrm{near}}$ of the conditional exponent along unit directions in the mode ball,
$r$ is the radius of that ball, and the two enter the exponential factor $e^{\Lambda r^2/2+1}$ of
the bound only through $\Lambda r^2$. With $\Lambda=C'^{\mathrm{op}}_F(N)x_k$ and $r=g_k$,
respectively $\Lambda=C'^{\mathfrak{g}}_F(N)x_k$ and $r=g_k/\mathsf{n}$, either reading gives
\begin{equation*}
\Lambda r^2=C'_W(N)+3.
\end{equation*}
The constant $C'_W(N)$ depends on which of the two readings of the background-field response in
\eqref{4.36:eq-curvature-readings-a} is taken: under the less favourable reading it carries a
factor at most $N\mathsf{n}$, and under the sharper reading a factor at most $\sqrt N$.
The variant has $(N-1)(C'_W(N)+3)$ at $r=g_k$. At $N=2$ the constant is $C'_W(2)+3$; we put
$C'_F:=C'_W(2)+3$ and use this letter only in the values at $N=2$ below.
\item \emph{The mode constant.} Along the operator-ball route,
\begin{equation*}
C_4(N)=(N-1)^{D/2}c_3(N)^{-1}e^{(C'_W(N)+3)/2+1}.
\end{equation*}
The sharper admissible value is
\begin{equation*}
c^{\mathrm{op}}_3(N)^{-1}e^{(C'_W(N)+3)/2+1}.
\end{equation*}
The variant has $c_3(N)^{-1}e^{(N-1)(C'_W(N)+3)/2+1}$. At $N=2$, $C_4(2)=(3\pi^3/8)e^{C'_F/2+1}$.
\item \emph{The product.} Along the operator-ball route, where $\sigma_N\omega_D$ cancels,
\begin{equation*}
C_4C_5(N)=(N-1)^{D/2}\kappa_N^{-D/2}(\pi/2)^{N^2-N}e^{(C'_W(N)+5)/2}.
\end{equation*}
In the variant the product becomes
\begin{equation*}
\kappa_N^{-D/2}(\pi/2)^{N^2-N}e^{(N-1)(C'_W(N)+3)/2+1}.
\end{equation*}
At $N=2$ the product is $\pi^5e^{C'_F/2+1}/(8\sqrt2)$.
\item \emph{Clause (iii) of Theorem~\ref{4.36:thm-plaquette-lower-bound},
\eqref{4.36:eq-plaquette-lower-bound-a}.} The bound is $\vartheta^B(s)+C_4C_5(N)a^{D/2}$ for every
$a>0$, and the factor $x_k^{D/2}$ cancels. This is the same in the variant. At $N=2$ the power is
$a^{3/2}$.
\item \emph{The lower constant.} Along the operator-ball route,
\begin{equation*}
a_0(N)=\Bigl(\frac{3}{8C_4C_5(N)}\Bigr)^{2/D}
=(N-1)^{-1}(3/8)^{2/D}(2/\pi)^{2N/(N+1)}\kappa_Ne^{-(C'_W(N)+5)/D}.
\end{equation*}
In the variant the lower constant becomes
\begin{equation*}
(3/8)^{2/D}(2/\pi)^{2N/(N+1)}\kappa_N\,e^{-((N-1)(C'_W(N)+3)+2)/D}.
\end{equation*}
At $N=2$, $a_0(2)=(8C_4C_5/3)^{-2/3}$.
\item \emph{Derived lower constants.} We have $c_1(N)=a_0(N)/2$ and
$a^{\mu}_0(N)=\frac12(4C_4C_5(N))^{-2/D}$. These are the same in the variant.
\item \emph{Haar measure in the low-action bound.} Let $\mathfrak{B}$ be the low-action set, $\mathfrak{V}$
the number entering the exponent, and $B_3$ the constant fixing the radius, of the low-action bound
in Corollary~\ref{4.36:cor-plaquette-constants}; the Haar measure of $\mathfrak{B}$ enters the
constant $C'_{\mathrm{low}}(N)$ of that corollary. It is bounded below
through $B^{\mathrm{op}}_{g_K/(5B_3)}$ by
\begin{equation*}
\bigl(c_3(N)(g_K/(5B_3\mathsf{n}))^D\bigr)^{4\mathfrak{V}},
\end{equation*}
with
\begin{equation*}
c_H(N)=4\log c_3(N)^{-1}+4D\log(5B_3\mathsf{n})+3.
\end{equation*}
These are the same in the variant. At $N=2$,
\begin{equation*}
c_H(2)=4\log(3\pi^3/8)+12\log(5B_3)+3,
\end{equation*}
and since $c_3(2)^{-1}=3\pi^3/8<12$, $c_H(2)$ is smaller than the value
\begin{equation*}
4\log12+12\log(5B_3)+3
\end{equation*}
obtained by replacing $c_3(2)^{-1}$ with its upper bound $12$.
\item \emph{Second-order constants along $\mathfrak{g}$-unit directions.} The constants $4$, $16$,
$176$ of the lemma giving the second-order bounds keep their values along $\mathfrak{g}$-unit
directions.
\item \emph{Constants of the upper bound and of the large-plaquette cost.} The constants $\gamma_0$
and $c^B=\gamma_0(A_1/A_0)^2$ enter as $\gamma_0/N$ and $c^B/N$. The constant
$C_{\mathrm{up}}(N,L,M)$ is proportional to $N/c^B$, and $C(N)=2$. These are the same in the variant.
\item \emph{Radius of the mode ball.} Along the operator-ball route the radius is $g_k$ in the
operator norm, that is, $\mathsf{n}$ times the radius $g_k/\mathsf{n}$ of the $\mathfrak{g}$-ball
contained in it, and the radius condition is vacuous because
$g_k\le\gamma<1<\pi/2$. The variant needs $g_k\le\pi/\sqrt{2N}$, which is imposed through the
ceiling $\gamma\le\pi/\sqrt{2N}$. At $N=2$
the radius bound is $\pi/2$.
\item \emph{Window, window condition and range of the far-part bound.} Along the operator-ball
route the window is the operator ball $B^{\mathrm{op}}_{\varrho}$ of the window radius $\varrho$,
with window condition $p_0(g_k)\ge8C_{\mathrm{loc}}$, where $C_{\mathrm{loc}}$ is the localisation
constant of the window; this condition is $N$-free in form. The far-part bound, clause (b) of the
one-parameter bounds of Lemma~\ref{4.36:lem-one-link-curvature}, is used only for $|t|<g_k$. The
variant needs $p_0(g_k)\ge8C_{\mathrm{loc}}\mathsf{n}$ and that bound for $|t|$ up to
$\mathsf{n}g_k$. The range $|t|\le\mathsf{n}g_k$ lies within the range $|t|\le2g_k$ of clause (b)
of the one-parameter bounds of Lemma~\ref{4.36:lem-one-link-curvature} if and only if $N\le5$; see
Remark~\ref{4.36:rem-mode-bound-lie-ball}.
\item \emph{Logarithmic coefficient of $E_+$ and $C_{\log}$.} The coefficient of the logarithmic
term of the density constant $E_+$ in the format of $\rho_k$, and the constant $C_{\log}$, both
as they enter Theorem~\ref{4.36:thm-never-large-exponent}, are
\begin{equation*}
\tfrac32(N^2-1)\quad\text{and}\quad C_{\log}=\tfrac92(N^2-1)L^4.
\end{equation*}
At $N=2$ they are $9/2$ and $27L^4/2$.
\item \emph{Constants of Lemma~\ref{4.36:lem-one-link-partition}.} We have $\dim\mathfrak{g}=D$,
$\sigma(0)=\sigma_N$ and $\nu_{\mathfrak{g}}=(D/2)\log(2\pi)$. Further,
\begin{equation*}
r_G=\sqrt6/N,\qquad \varepsilon_G=2\sin(\sqrt6/(2N)),\qquad b_G=\mathsf{n},
\end{equation*}
together with $c_G(N)$ and $\gamma_z(N,A_0,p_0)$.
\item \emph{Non-degeneracy.} The image of $\operatorname{Haar}$ under
$u\mapsto1-\operatorname{Re}\operatorname{tr}u$ has no atoms. This is the non-degeneracy
statement for this image used at $N=2$, which holds unchanged at $SU(N)$.
\end{itemize}
\end{remark}

\begin{proof}
\emph{The two halves of (F7).} Among the inequalities \eqref{4.36:eq-eigenangle-inequalities-a},
(F7) consists of a lower bound of $1-\operatorname{Re}\operatorname{tr}W$ by
$(2N)^{-1}|W-\mathbf{1}|^2$, its left half, and an upper bound of
$1-\operatorname{Re}\operatorname{tr}W$ by $\frac12|W-\mathbf{1}|^2$, its right half, which is used
here in the window $|W-\mathbf{1}|<2\varepsilon_k$.

\emph{The plaquette observable and the complexified action.} Write $V=V_k(\partial p)$. Then
$\operatorname{tr}V^*V=\operatorname{tr}\mathbf{1}=1$. Expanding,
\begin{equation*}
\|V-\mathbf{1}\|^2=\operatorname{tr}(V^*V-V^*-V+\mathbf{1})=2-2\operatorname{Re}\operatorname{tr}V.
\end{equation*}
Since $\operatorname{tr}=N^{-1}\operatorname{Tr}$, this gives the three expressions for $O$. The
range $[0,2]$ holds because $|\operatorname{tr}V|\le1$. The window condition reads only the right
half of (F7), and that half has no $N$-dependent constant. At $N=2$, $\operatorname{Tr}V$ is real, so
$O=1-\frac12\operatorname{Tr}V$.

For $\mathfrak{w}$, expand $W^{-1}(W-\mathbf{1})^2=W-2\cdot\mathbf{1}+W^{-1}$ and take
$-\frac12\operatorname{tr}$. Since $\operatorname{tr}\mathbf{1}=1$, this gives
$1-\frac12(\operatorname{tr}W+\operatorname{tr}W^{-1})$.

\emph{Haar constants.} By Lemma~\ref{4.36:lem-haar-small-sets}\,(a), (d),
\begin{equation*}
\operatorname{Haar}(B^{\mathrm{op}}_r)\ge\operatorname{Haar}(B^{\mathfrak{g}}_{r/\mathsf{n}})
\ge c_3(N)(r/\mathsf{n})^D=(N-1)^{-D/2}c_3(N)\,r^D .
\end{equation*}
This holds whenever $r/\mathsf{n}\le\pi/\sqrt{2N}$. That condition follows from $r\le\pi/2$ because
$(N-1)/(2N)\ge\frac14$ for $N\ge2$. Hence the mode-ball constant is
$c^{\flat}_3(N)=(N-1)^{-D/2}c_3(N)$ for $r\le\pi/2$. By
Lemma~\ref{4.36:lem-haar-small-sets}\,(a$^{\mathrm{op}}$), $c^{\mathrm{op}}_3(N)\ge c^{\flat}_3(N)$,
which is why $c^{\mathrm{op}}_3(N)$ is the sharper value. The constants $c_3(N)$, $C_5(N)$ and
$\kappa_N$ are those of Lemma~\ref{4.36:lem-haar-small-sets}, and the variant uses $c_3(N)$ at its
own radius.

\emph{Curvature constant.} For the operator-ball route take $\Lambda=C'^{\mathrm{op}}_Fx_k$ and
$r=g_k$. Then $\Lambda r^2=(C'_W+3)x_kg_k^2=C'_W+3$. Reading the same bound in $\mathfrak{g}$-unit
directions on $B^{\mathfrak{g}}_{g_k/\mathsf{n}}(u^0)$ gives
\begin{equation*}
\Lambda r^2=(N-1)(C'_W+3)x_k\cdot g_k^2/(N-1)=C'_W+3 .
\end{equation*}
Thus the factor $\mathsf{n}^2$ of $C'^{\mathfrak{g}}_F$ is cancelled by the radius. In the variant
$r=g_k$, so nothing cancels. The factors $N\mathsf{n}$ and $\sqrt N$ carried by $C'_W(N)$ under the
two readings of the background-field response are those of \eqref{4.36:eq-curvature-readings-a}.

\emph{The mode constant and the product.} Theorem~\ref{4.36:thm-plaquette-lower-bound}\,(iv)
defines
\begin{equation*}
C_4(N)=c^{\flat}_3(N)^{-1}e^{(C'_W(N)+3)/2+1}
=(N-1)^{D/2}c_3(N)^{-1}e^{C'^{\mathrm{op}}_F(N)/2+1},
\end{equation*}
and the same statement admits $c^{\mathrm{op}}_3(N)^{-1}e^{C'^{\mathrm{op}}_F/2+1}$. The variant
value is that of Remark~\ref{4.36:rem-mode-bound-lie-ball}. Insert
$c_3(N)^{-1}=(\sigma_N\omega_D)^{-1}(\pi/2)^{N^2-N}$ and
$C_5(N)=\sigma_N\omega_D\kappa_N^{-D/2}$, and use $(C'_W+3)/2+1=(C'_W+5)/2$. The product takes the
displayed form, and $\sigma_N\omega_D$ cancels. The variant product follows in the same way.

\emph{Clause (iii).} By the proof of Theorem~\ref{4.36:thm-plaquette-lower-bound}, the normalised
one-link density of that proof is at most $C_4(N)x_k^{D/2}$ on the good set of each clean term. The
sublevel set has Haar measure at most
$C_5(N)(a/x_k)^{D/2}$. The product is $C_4C_5(N)a^{D/2}$, so $x_k^{D/2}$ cancels. The remaining mass
is at most $\vartheta^B(s)$.

\emph{The lower constant.} Raise the product to the power $-2/D$. The route constant gives
\begin{equation*}
(C_4C_5(N))^{-2/D}=(N-1)^{-1}\kappa_N(\pi/2)^{-2(N^2-N)/D}e^{-(C'_W+5)/D}.
\end{equation*}
Since $2(N^2-N)/(N^2-1)=2N/(N+1)$, the power of $\pi/2$ becomes $(2/\pi)^{2N/(N+1)}$. In the variant
the exponent of $e$ is $-\bigl((N-1)(C'_W+3)/2+1\bigr)\cdot2/D$, which equals
$-((N-1)(C'_W+3)+2)/D$. The identity $c_1(N)=a_0(N)/2$ is Theorem~\ref{4.36:thm-plaquette-lower-bound}\,(i).
The value of $a^{\mu}_0(N)$ is stated in Corollary~\ref{4.36:cor-plaquette-constants}.

\emph{Haar measure in the low-action bound.} The radius $g_K/(5B_3)$ does not exceed $\pi/2$ (the
condition $g_K\le5\pi B_3/2$ that this requires is vacuous under the ceiling on the coupling), so
clause (a$^{\mathrm{op}}$) applies. By
Lemma~\ref{4.36:lem-haar-small-sets}\,(d) and (a$^{\mathrm{op}}$),
\begin{equation*}
\operatorname{Haar}\{|u-\mathbf{1}|<g_K/(5B_3)\}\ge c^{\flat}_3(N)\Bigl(\frac{g_K}{5B_3}\Bigr)^{D}
=c_3(N)\Bigl(\frac{g_K}{5B_3\mathsf{n}}\Bigr)^{D}.
\end{equation*}
The exponent $4\mathfrak{V}$ and the constant $c_H(N)$ are as in
Corollary~\ref{4.36:cor-plaquette-constants}. At $N=2$, $c_3(2)^{-1}=3\pi^3/8<12$ because $\pi^3<32$.

\emph{Second-order constants.} The constants $4$, $16$, $176$ of the lemma giving the second-order
bounds keep their values along $\mathfrak{g}$-unit directions by (F8) in
\eqref{4.36:eq-eigenangle-inequalities-a}.

\emph{Upper-bound and large-plaquette constants.} A large plaquette is paid for by the action
through the left half of (F7): for every $x>0$ and $W\in SU(N)$,
\begin{equation*}
e^{-x(1-\operatorname{Re}\operatorname{tr}W)}\le e^{-x|W-\mathbf{1}|^2/(2N)}.
\end{equation*}
This replaces $\gamma_0$ by $\gamma_0/N$ and $c^B$ by $c^B/N$. The statements
$C(N)=2$ and $C_{\mathrm{up}}(N,L,M)\propto N/c^B$ are those of
Corollary~\ref{4.36:cor-plaquette-constants}.

\emph{Radius, window and range of the far-part bound.} Every
$u\in B^{\mathrm{op}}_{g_k}(u^0)\setminus\{u^0\}$ is $u^0e^{iv}$ with $v\in\mathfrak{g}$ and
$0<|v|<g_k$; put $\hat T:=v/|v|$ and $t:=|v|$, so that $u=u^0e^{it\hat T}$ with
$t\in(0,g_k)\subset[-2g_k,2g_k]$. Here the far-part bound, clause (b) of the one-parameter bounds
of Lemma~\ref{4.36:lem-one-link-curvature}, holds. In the variant that bound is needed for
$|t|$ up to $\mathsf{n}g_k$.
The inequality $\mathsf{n}g_k\le2g_k$ is equivalent to $\mathsf{n}\le2$, that is, to $N\le5$.

\emph{Logarithmic constants.} The values at general $N$ of the coefficient of the logarithmic term
of $E_+$ and of $C_{\log}$ are those with which they enter
Theorem~\ref{4.36:thm-never-large-exponent}. At $N=2$, where $N^2-1=3$, they are
$\frac32\cdot3=\frac92$ and $\frac92\cdot3L^4=\frac{27}{2}L^4$.

\emph{Non-degeneracy.} The non-degeneracy item is the statement, used at $N=2$, that the image of
$\operatorname{Haar}$ under $u\mapsto1-\operatorname{Re}\operatorname{tr}u$ has no atoms, read at
$SU(N)$; its statement involves $N$ only through the group.

\emph{Values at $N=2$.} Here $D=3$, $N^2-N=2$, $\mathsf{n}=1$ and $(N-1)^{D/2}=1$. Hence
$c^{\flat}_3(2)=c_3(2)$, $C'^{\mathfrak{g}}_F(2)=C'^{\mathrm{op}}_F(2)$ and
$\pi/\sqrt{2N}=\pi/2$, so the operator-ball route and the $\mathfrak{g}$-ball variant give the same
values. From $\sin(\pi/2)=1$ we get
$\kappa_2=4/(2\pi^2)=2/\pi^2$, and so $\kappa_2^{-3/2}=\pi^3/(2\sqrt2)$. Let $P_1,P_2,P_3$ denote
the Pauli matrices, which span $\mathfrak{g}$ for $N=2$. For a unit vector $y\in\mathbb{R}^3$ and
$\theta\ge0$ the element $X=\theta\sum_jy_jP_j$ of $\mathfrak{g}$ satisfies $X^2=\theta^2\mathbf{1}$,
so $|X|_{\mathfrak{g}}=\theta$, and
\begin{equation*}
e^{iX}=\cos\theta\,\mathbf{1}+i\sin\theta\,\textstyle\sum_jy_jP_j .
\end{equation*}
Identifying $SU(2)$ with the unit sphere of $\mathbb{R}^4$ through
$q_0\mathbf{1}+i\sum_jq_jP_j\mapsto(q_0,q_1,q_2,q_3)$, with $q\in\mathbb{R}^4$ and $|q|=1$, the map
$X\mapsto e^{iX}$ is therefore the
exponential map at $\mathbf{1}$ of the round unit three-sphere. With the metric of
$|\cdot|_{\mathfrak{g}}$, $SU(2)$ is thus isometric to the unit three-sphere, of volume $2\pi^2$, and
$v\mapsto e^{iv}$ is an isometry to first order at $0$; hence $\sigma_2=1/(2\pi^2)$. With
$\omega_3=4\pi/3$ and $(2/\pi)^{2}=4/\pi^2$,
\begin{equation*}
c_3(2)=\frac{1}{2\pi^2}\cdot\frac{4\pi}{3}\cdot\frac{4}{\pi^2}=\frac{8}{3\pi^3},
\qquad \sigma_2\omega_3=\frac{2}{3\pi}.
\end{equation*}
With this value,
\begin{equation*}
C_5(2)=\frac{2}{3\pi}\cdot\frac{\pi^3}{2\sqrt2}=\frac{\pi^2}{3\sqrt2}.
\end{equation*}
The prefactor of $e^{(C'_W(2)+5)/2}=e^{C'_F/2+1}$ in the route formula for the product is, at
$N=2$,
\begin{equation*}
(2/\pi^2)^{-3/2}(\pi/2)^2=\frac{\pi^5}{8\sqrt2}
=\frac{3\pi^3}{8}\cdot\frac{\pi^2}{3\sqrt2},
\end{equation*}
in agreement with $C_4(2)C_5(2)=\pi^5e^{C'_F/2+1}/(8\sqrt2)$. The power $a^{D/2}$ is $a^{3/2}$, and
$a_0(2)=(8C_4C_5/3)^{-2/3}$.
The remaining values are $\frac32\cdot3=\frac92$ and $\frac92\cdot3L^4=\frac{27}{2}L^4$. In
Lemma~\ref{4.36:lem-one-link-partition} they are $r_G=\sqrt6/2$, $\varepsilon_G=2\sin(\sqrt6/4)$,
$b_G=1$ and $\nu_{\mathfrak{g}}=\frac32\log(2\pi)$.
\end{proof}

\begin{remark}[Order of choice of the $SU(N)$ constants]\label{4.36:rem-order-of-choice}
The constants of this chapter are chosen in the order $N$, $L$, $\mathfrak{p}$, $\gamma$. First $N$ is fixed. Then the odd blocking factor $L$ is fixed; it is never given a value. Then Ba{\l}aban's auxiliary parameters $\mathfrak{p}$ are fixed in Ba{\l}aban's order. Last, the coupling ceiling $\gamma$ is fixed.

A constant is said to be \emph{fixed after $N$ and before $L$} when its dependence on $N$ is through the values listed under (a) below only, so that once $N$ is fixed it is a definite function of $(L,\mathfrak{p})$ of the same form as at $N=2$. The number $N$ enters this order of choice only in the three ways (a)--(c) below.
\begin{itemize}
\item[(a)] \emph{Through values.} $N$ enters through the values of constants that depend on $(N,L,\mathfrak{p})$ in the same way as the corresponding constants at $N=2$ depend on $(L,\mathfrak{p})$. There are two groups of these.
\begin{itemize}
\item[(a1)] Pure numbers depending on $N$ alone, fixed as soon as $N$ is fixed. These are:
\begin{itemize}
\item $D=N^2-1$ and $\mathsf{n}=(N-1)^{1/2}$;
\item the normaliser $\sigma_N$ of the Haar density in exponential coordinates;
\item the volumes $\omega_D$ and $\omega^{\mathrm{op}}_N$ of the unit balls in the two norms;
\item the Haar small-ball and sublevel constants $\kappa_N$, $c_3$, $c^{\flat}_3$, $c^{\mathrm{op}}_3$, $C_5$;
\item the number $(2/\pi)^{N^2-N}$;
\item the constants $c_G(N)$, $\nu_{\mathfrak{g}}$, $r_G$, $\varepsilon_G$ of Lemma~\ref{4.36:lem-one-link-partition};
\item the constant $C_{\mathfrak{g}}=2N^2$;
\item the factor $1/N$ carried by one class of terms of the construction.
\end{itemize}
\item[(a2)] The values at $G=SU(N)$ of the constants of Ba{\l}aban's papers, taken at the generality $G\subset U(N)$ at which they are stated there.
\end{itemize}
These two groups of values fix $c_0(N)$ and $\bar{B}(N,L,\mathfrak{p})$, which here and below denotes $\max\{C_\beta+1,c_5\}$, and hence the following constants:
\begin{itemize}
\item $C_q(N)$, $C'_W(N)$ and the constants $C'^{\mathrm{op}}_F$ and $C'^{\mathfrak{g}}_F$ of Lemma~\ref{4.36:lem-one-link-curvature}, together with the level-$m$ summation constant $\Theta_m$ and the constant $\mathfrak{A}$ of the complex plaquette variable $\mathfrak{U}(\partial q)$;
\item $C_4$, $C_4C_5$, $a_0$, $c_1$, $a^{\mu}_0$, $c_H$ and $C'_{\mathrm{low}}$ of Theorem~\ref{4.36:thm-plaquette-lower-bound} and Corollary~\ref{4.36:cor-plaquette-constants};
\item the ratios $\gamma_0/N$ and $c^{B}/N$, and the constant $C_{\mathrm{up}}(N,L,M)$, which is proportional to $N/c^{B}$, all of Corollary~\ref{4.36:cor-plaquette-constants};
\item $E_0(N,L,\mathfrak{p})$, $E'_1$, $R'_1$, $B_0$ and the further constants of Ba{\l}aban's bounds;
\item the localisation constant $C^{\mathrm{nl}}_{\mathrm{loc}}(N)$, the constant $C_{\mathfrak{r}}$ and the constant $C_{\mathrm{nl}}$ of Theorem~\ref{4.36:thm-never-large-exponent}, the index $\mathrm{nl}$ referring to the never-large term of that theorem;
\item the density constants $E_{\pm}(N,L,\mathfrak{p})$, with the coefficient $\tfrac32(N^2-1)$ in $E_{+}$;
\item the constant $C(N)=2$ of Corollary~\ref{4.36:cor-plaquette-constants}.
\end{itemize}
None of the constants in (a1), (a2) or the list above is compared with any of the following:
\begin{itemize}
\item a power of $L$ alone;
\item $K$, $m$, $k$ or $N_0$;
\item the number of $\mathbf{R}$-operations or of arrests;
\item any $g_k$ from below.
\end{itemize}
\item[(b)] \emph{Through coupling ceilings.} $N$ enters through one-sided upper bounds on $\gamma$, which are fixed last.
\item[(c)] \emph{Through floors on $\mathfrak{p}$.} $N$ enters through lower bounds on the auxiliary parameters $\mathfrak{p}$, each of which leaves the admissible set open upward.
\end{itemize}

No uniformity in $N$ is asserted. Indeed,
\begin{equation}\label{4.36:eq-order-of-choice-1}
C_4C_5(N)\ge 2^{D/2}(N-1)^{3D/2}e,
\qquad
a_0(N)\le \frac{\kappa_N}{N-1}\le \frac{1}{2(N-1)^3},
\end{equation}
and the last bound tends to zero as $N\to\infty$.
\end{remark}

\begin{proof}
We go through the conditions of this chapter in the order of choice. In each case we note which parameter the condition bounds, and in which direction.

\emph{Conditions on $L$.} The integer $L$ is odd and is kept symbolic. The conditions imposed on it are the following.
\begin{itemize}
\item The inequality $6ML^{-2}<ML^{-1}$, which is a lower bound on $L$ by an absolute number.
\item The inequality $M(L-6)>3L^3$. Given $M\ge L^4$, it follows from $L(L-6)>3$, again a lower bound on $L$ by an absolute number.
\item The condition $L^2\ge 20(1+\beta_0)$ among the thresholds \eqref{4.36:eq-flow-thresholds-b}, where $\beta_0$ is the constant appearing in \cite[(3.8)]{B-CMP119} and in the transport bounds \eqref{4.36:eq-transport-every-l-a}. Whenever $\beta_0\le 1/7$, its right side is bounded by an absolute number, so it too is a lower bound on $L$ by an absolute number.
\item The requirement that the integer $\ell_{\circ}=\lfloor\sqrt{L/3}\rfloor$ of \cite{B-CMP122a} satisfy $\ell_{\circ}\ge 1$.
\end{itemize}
The first two together amount to $L\ge 7$, which lies inside the lower bound $L>11$ of Ba{\l}aban's construction; for $\beta_0\le 1/7$ the third amounts to $L\ge 5$; and the fourth holds as soon as $L\ge 3$. None of them involves $N$.

\emph{The condition $L^2\ge 1+\bar B$ is not imposed, at any $N$.} At $SU(N)$ this condition would read
\[
L^2\ge 1+\bar{B}(N,L,\mathfrak{p}),
\qquad\text{with}\qquad
\bar{B}\ge \frac{11N^2}{48\pi^2}\log L,
\]
by Notation~\ref{4.36:not-flow-thresholds}. At fixed $L$ it fails for all sufficiently large $N$, so it would be a lower bound on $L$ that depends on $N$. Already at $N=2$ it compares $L$ with a number fixed after $L$ and $\mathfrak{p}$. That number is $\bar{B}=\max\{C_\beta+1,c_5\}$, in which $C_\beta+1$ is fixed after $\mathfrak{p}$, so that in the order of choice it could not be imposed as a condition on $L$, and no choice of the constants satisfying it is exhibited.

Its role is taken by the transport bounds \eqref{4.36:eq-transport-every-l-a} of Proposition~\ref{4.36:prop-transport-every-l}, whose proof is incomplete at one step. These bounds hold for every $L\ge 2$ under the flow inequalities of Notation~\ref{4.36:not-flow-thresholds}, the anchor $x_K\ge g^{-2}$ and the coupling ceiling $\gamma\le\gamma_{\bar B}$ described below. They replace the condition $L^2\ge1+\bar B$ at every place where it entered that Proposition~\ref{4.36:prop-transport-every-l} examines. Should the step of that proof marked $(\ast)$ reveal a direct comparison of $\bar B$ with $L^2$, the condition would re-enter at that place as a one-sided lower bound on $L$ depending on $N$; it would be neither an upper bound on $L$ nor a lower bound on the coupling.

\emph{Conditions on $\mathfrak{p}$.} These are taken in Ba{\l}aban's order, and each admissible set is open upward unless stated otherwise.
\begin{itemize}
\item The parameter $M=L^{m'}$ is bounded below by the floor on $M$ in \eqref{4.36:eq-flow-thresholds-b}. That floor is the maximum of $L^4$, $RM_1$, $14$ and a number $M_*(d;\delta_0,c_0)$ built from the absolute constants $\delta_0$, $c_0$ of \cite{B-CMP96} (this $c_0$ is not the one-loop constant $c_0(N)$). The floor also carries the two lower bounds on $\delta_0M$ displayed in \eqref{4.36:eq-flow-thresholds-b}.
\item The parameters $\kappa$ and $\kappa_0$ satisfy
\[
\kappa\ge\max\{10,\,8c_d+8\},\qquad \kappa\ge\kappa(d),\qquad \kappa_0\ge 7.
\]
Here $c_d$ and $\kappa(d)$ depend on $d$ only. These lie inside the floors $\kappa\ge 10\log L$ and $\kappa_0\ge 17$ imposed with $\mathfrak{p}$; all are lower bounds on $\kappa$ and $\kappa_0$, open upward.
\item The exponents $q$, $q_0$, $q_1$ satisfy the lower bounds relative to $p_0$ and $r$ imposed with the radii $\alpha_{0,j}$, $\alpha_{1,j}$.
\item The ratio $A_0/A_1$ satisfies
\[
\frac{A_0}{A_1}\ge c(d)(1+\beta_0)B_3L^{d+2},
\]
with $c(d)$ depending on $d$ only and $B_3$ the constant of \cite{B-CMP102b}.
\end{itemize}
Two of the conditions on $\mathfrak{p}$ are upper bounds:
\begin{itemize}
\item the upper bound in \eqref{4.36:eq-flow-thresholds-b} on the auxiliary radius that bounds $g_k$ there;
\item the upper bound $\varepsilon^I_1\le c_0(N)/(2C_{\mathrm{rem}})$, with $C_{\mathrm{rem}}$ the constant introduced in Lemma~\ref{4.2:lem-unmarked}. This ceiling is imposed with the $SU(N)$ values of the constants of Ba{\l}aban's step, not by the conditions examined here; it is recorded because it is an upper bound on a parameter of $\mathfrak{p}$.
\end{itemize}
The remaining conditions of \eqref{4.36:eq-flow-thresholds-b} are upper bounds on $\gamma$: the lower bound $p_0(g_k)\ge 8C_{\mathrm{loc}}$, with $C_{\mathrm{loc}}$ the localisation constant in the window radius $\varepsilon_k/(4C_{\mathrm{loc}})$, the bound of $g_k$ by the auxiliary radius and the further conditions there on the couplings. They are listed among the coupling ceilings below.

The flow inequalities are used with the frozen values $\underline{\lambda}=c_0(N)/4$ in place of $\lambda^{\sharp}$ and $C_\beta+1$ in place of $B^{\sharp}$, with $c_5$ kept. The depth satisfies the one-sided floor $s\ge s_*(N,L,M)$. The torus exponent is bounded below by $m_{\mathbb{T}}(g)$, and only from below.

\emph{Readings that depend on the $SU(N)$ values of Ba{\l}aban's constants.} Two cases must be considered.
\begin{itemize}
\item $\delta_0$ and the absolute constant $c_0$ of \cite{B-CMP96} are used here as stated there. The order of choice does not use that they are independent of $N$: read as functions of $(d,N)$, they give the floor $M\ge M_*(d,N)$, with the other entries of the maximum in \eqref{4.36:eq-flow-thresholds-b} unchanged. This floor is fixed after $(N,L)$ within the choice of $\mathfrak{p}$, and $M=L^{m'}$ remains free upward. It is a one-sided condition on $M$ and not a condition on $L$.
\item $B_3$ is used here as stated in \cite{B-CMP102b}, where, at the generality $G\subset U(N)$, it and two further constants depend on $d$ and $L$ only. The order of choice does not use this: even with $B_3$ read as a function of $(d,L,N)$, the lower bound on $A_0/A_1$ is fixed after $(N,L)$ and is open upward.
\end{itemize}
None of these conditions is two-sided. None bounds $g$ from below. None bounds $L$, $K$, $m$ or $N_0$ from above.

\emph{Coupling ceilings, chosen last.} Each of the following is an upper bound on $\gamma$.
\begin{itemize}
\item $\gamma<1$. Since $0<g_k\le\gamma$ (Notation~\ref{4.36:not-flow-thresholds}), this gives $g_k<1<\pi/2$ and $x_k>1$.
\item $\gamma\le\gamma_z(N,A_0,p_0)$. By Lemma~\ref{4.36:lem-one-link-partition}, every $0<g\le\gamma_z$ satisfies
\[
g\le\min\{e^{-p_0},\,r_G/2\},\qquad g\,p_0(g)\le\varepsilon_G,\qquad p_0(g)\ge 2\mathsf{n}.
\]
At $SU(N)$, moreover, $\gamma_z\le\min\{e^{-p_0},\sqrt6/(2N)\}$, and every such $g$ satisfies $g\,p_0(g)\le 2\sin(\sqrt6/(2N))$. No lower bound on $A_0$ is required.
\item The conditions of \eqref{4.36:eq-flow-thresholds-b} on the couplings. These comprise $p_0(g_k)\ge 8C_{\mathrm{loc}}$, whose form does not involve $N$; the three further coupling conditions listed there; the bound of $g_k$ by the auxiliary radius; and $x_k\ge x_{\mathrm{th}}(N,L,M,\mathfrak{p})$.
\item $\gamma\le\gamma_N(N,L,M,\mathfrak{p})$. Here $\gamma_N$ is the minimum of the following, all taken at the $SU(N)$ values of the constants:
\begin{itemize}
\item $(x_{\mathrm{th}}+c_5)^{-1/2}$, which enters through \eqref{4.36:eq-flow-thresholds-a};
\item $\gamma_z$;
\item the ceiling $g^{\natural}$ with its provisos;
\item the remaining coupling ceilings of the construction.
\end{itemize}
\end{itemize}

The replacement of the condition $L^2\ge 1+\bar{B}$ adds one entry: $x_{\mathrm{th}}$ is replaced by $\max\{x_{\mathrm{th}},\bar B,x_\delta,x_\beta\}$, where $x_\delta$ and $x_\beta$ are the floors entering Proposition~\ref{4.36:prop-transport-every-l} (whose proof is incomplete at one step). In terms of the coupling this reads
\begin{equation}\label{4.36:eq-order-of-choice-2}
\gamma_N\le\gamma_{\bar B}=\bigl(\max\{\bar B,x_\delta,x_\beta\}+c_5\bigr)^{-1/2}.
\end{equation}
This is a number fixed after $(N,L,\mathfrak{p})$. It is of the same kind as the ceiling on $\gamma_H$ and as the function $g_-$:
\[
\gamma_H\le(\gamma^{-2}+c_5+1)^{-1/2},
\qquad
g_-(g)=(g^{-2}+B^{\sharp})^{-1/2}.
\]
A second entry comes from collecting, at $SU(N)$, the conditions on the couplings of the same kind as the third coupling condition listed in \eqref{4.36:eq-flow-thresholds-b}: $x_{\mathrm{th}}$ is further replaced by $\max\{x_{\mathrm{th}},x'_{\mathrm{th}}(N,L,M,\mathfrak{p})\}$, where $x'_{\mathrm{th}}(N,L,M,\mathfrak{p})$ is a number fixed after $(N,L,M,\mathfrak{p})$. Through \eqref{4.36:eq-flow-thresholds-a} this too is an upper bound on $\gamma$, chosen last.

Three further points complete the list.
\begin{itemize}
\item The ratios $\gamma_0/N$ and $c^{B}/N$ enter every large-field threshold, the depth floor $s_*(N,L,M)$ and the depth bound $s_{\mathrm{up}}(g;N)$. The resulting factors $1/N$ are absorbed by the strict exponent gap $2p_0>5r+1$, with no condition on $A_0$.
\item The term
\[
\bigl(C_{\log}(N)\log x_{j-1}+C^{\mathrm{nl}}_{\mathrm{loc}}(N)\bigr)|Z^{+}_j|_j
\]
is absorbed as in \cite{B-CMP122b}.
\item The Haar bound behind $c_H$ uses the ceiling $g_K\le 5\pi B_3/2$, which is again an upper bound on the coupling, chosen last.
\end{itemize}
A bookkeeping of the $SU(N)$ constants that carries the factor $\mathsf{n}=(N-1)^{1/2}$ explicitly in the one-link conditions imposes in addition $g_k\le\pi/\sqrt{2N}$, $p_0(g_k)\ge 8C_{\mathrm{loc}}\mathsf{n}$, $\mathsf{n}g_k$ below the auxiliary radius, $\mathsf{n}\eta_{\mathrm{far}}\le1$ with $\eta_{\mathrm{far}}$ as in Lemma~\ref{4.36:lem-one-link-curvature}, and a further condition at radius $\mathsf{n}g_k$; each of these is again an upper bound on the coupling, and the order of choice described here does not use them.

\emph{Conclusion.} We have shown four things:
\begin{itemize}
\item every condition on $L$ is a lower bound not involving $N$;
\item every condition on $\mathfrak{p}$ is one-sided, and the lower bounds leave the admissible set open upward;
\item every condition involving $\gamma$ is an upper bound chosen last;
\item every constant listed in (a) enters only as a value inside such conditions.
\end{itemize}
None of these constants is compared with a power of $L$ alone, with $K$, $m$, $k$, $N_0$, with the number of $\mathbf{R}$-operations or of arrests, or from below with any $g_k$. This is (a)--(c). The bounds \eqref{4.36:eq-order-of-choice-1} show that the constants $C_4C_5(N)$ and $a_0(N)$ of Theorem~\ref{4.36:thm-plaquette-lower-bound} degenerate as $N$ grows, which is why no uniformity in $N$ is asserted.
\end{proof}

\begin{definition}[Conventions at $SU(N)$: group, norms and coupling flow]
\label{4.36:not-sun-conventions}
\emph{Order of choices.} The integer $N\ge2$ is fixed first.
Next the blocking factor $L$ is chosen; it is odd and is kept
symbolic. Next Ba{\l}aban's auxiliary parameters $\mathfrak{p}$ are chosen, in their fixed
order. At the same stage the auxiliary radius $r_0$ of the thresholds below is fixed, after $N$
and $L$, together with the constants $C_H$, $C_P$, $C_1^{(102)}$, $a_1'$ of the thresholds (t2)
and (t8) below, and before $\gamma$. Last, the
coupling ceiling $\gamma$ is chosen. Every condition imposed in what follows is one-sided.

\emph{Absolute facts.} Here $e$ is Euler's number, and in $e^{x}\ge x^{p}/p!$ the exponent $p$ is
a non-negative integer. The following inequalities are used:
\begin{equation}\label{4.36:eq-sun-conventions-a}
\begin{aligned}
(\mathrm{e}1)\quad & 2<e<3,\qquad \log2<1,\qquad e-1<2,\qquad \frac{1}{e-1}<1,\qquad
\frac{e}{e-1}<2;\\
(\mathrm{e}2)\quad & \log y\le y-1\ \ (y>0),\qquad e^{x}\ge1+x\ \ (x\in\mathbb{R}),\qquad
e^{x}\ge\frac{x^{p}}{p!}\ \ (x\ge0),\\
& L^{u}-1\ge u(L-1)\ \ (u\in\mathbb{Z}_{\ge0}\ \text{or real}\ u\ge1;\ \text{Bernoulli}),\\
& (a+b)^{p}\le2^{p-1}(a^{p}+b^{p})\ \ (a,b\ge0,\ p\ge1),\\
& \sum_{\ell\ge0}e^{-a\ell}=(1-e^{-a})^{-1}\le2\ \ (a\ge\log2),\qquad
(1-L^{-4})^{-1}\le2\ \ (L^{4}\ge2);\\
(\mathrm{e}3)\quad & e^{s}-1\le(e-1)s\qquad(0\le s\le1);\\
(\mathrm{e}4)\quad & (2+\ell)^{1/2}e^{-\ell}\le\sqrt2\qquad(\ell\ge0).
\end{aligned}
\end{equation}

\emph{Group and norms.} We put $G=SU(N)$. Let $\mathfrak{g}$ be the real vector space of traceless
Hermitian $N\times N$ matrices, and let $\mathfrak{g}^{c}=\mathfrak{sl}(N,\mathbb{C})$ be its
complexification. We write $\operatorname{tr}:=N^{-1}\operatorname{Tr}$, so that
$\operatorname{tr}\mathbb{1}=1$ as in Notation~\ref{0.2:not-glossary}, and $|\cdot|$ for the
operator norm on $M_N(\mathbb{C})$ \cite[(19)--(20)]{B-CMP98}. We also put
\begin{equation*}
\|Y\|^{2}:=\operatorname{tr}Y^{*}Y,\qquad |X|_{\mathfrak{g}}:=(\operatorname{tr}X^{2})^{1/2},\qquad
D:=N^{2}-1,\qquad \mathsf{n}:=(N-1)^{1/2}.
\end{equation*}
We shall use the following two elementary matrix inequalities:
\begin{equation}\label{4.36:eq-sun-conventions-2}
|e^{Z}-\mathbb{1}|\le e^{|Z|}-1\quad(Z\in M_N(\mathbb{C})),\qquad
|e^{iX}-\mathbb{1}|\le|X|\quad(X\in\mathfrak{g}).
\end{equation}
For a unitary matrix $v$ we put $\operatorname{Ad}(v)X:=vXv^{-1}$. This map is an isometry of
$(\mathfrak{g}^{c},|\cdot|)$ for every $N$. For $X\in\mathfrak{g}$ and $Y\in\mathfrak{g}^{c}$ we
have
\begin{equation}\label{4.36:eq-sun-conventions-1}
\bigl|\operatorname{Ad}(e^{iX})Y-Y\bigr|=\bigl|e^{\operatorname{ad}_{iX}}Y-Y\bigr|
\le\sum_{n\ge1}\frac{(2|X|)^{n}|Y|}{n!}=(e^{2|X|}-1)|Y|,
\end{equation}
with constants independent of $N$ (compare \cite[(20)]{B-CMP98}).

\emph{Lattice, couplings and scales.} We put $d=4$. At level $k$ we put $\eta:=L^{-k}$, and for
$n\le k$ we put $\ell_n:=L^{-(k-n)}$; these are lengths in the units of level $k$, so
$\ell_k=1$. Throughout, as in the hypotheses of \cite[Theorem~1]{B-CMP122b}, the running
couplings are required to satisfy $0<g_j\le\gamma<1$ for $0\le j\le K$. With $x_j=g_j^{-2}$ as in
Definition~\ref{1.2:def-couplings}, we have $x_j>1$. With the degree function $p_0(\cdot)$ of
Notation~\ref{0.2:not-glossary}, $p_0(g')=A_0(\log g'^{-2})^{p_0}$, we put
\begin{equation*}
\varepsilon_j:=g_j\,p_0(g_j),\qquad
\delta_j:=\frac{A_1}{A_0}\,\varepsilon_j=A_1g_j(\log x_j)^{p_0},
\end{equation*}
\begin{equation*}
\alpha_{0,j}:=C_0g_j(\log x_j)^{q_0},\qquad \alpha_{1,j}:=C_1g_j(\log x_j)^{q_1}.
\end{equation*}
Here the exponents $q_0$, $q_1$ are among the auxiliary parameters $\mathfrak{p}$.
We let $R_j$ be the least non-negative integer power of $L$ that is $\ge(\log x_j)^{r}$. Thus
$R_j\ge1$ and $\log_LR_j\in\mathbb{Z}_{\ge0}$.

\emph{The coupling flow.} The running inverse couplings satisfy, for some constants $\lambda$ and
$B$ with $\lambda\ge\lambda^{\sharp}\ge c_0(N)/4$ and $B\le\bar B$, the two-sided flow bound at
$SU(N)$ of Theorem~\ref{4.1:thm-clause-zero}, with the flow constants of
Definition~\ref{1.2:def-flow-constants}:
\begin{equation}\label{4.36:eq-sun-conventions-3}
\lambda(j-i)-c_5\le x_i-x_j\le B(j-i)\qquad(0\le i\le j\le K).
\end{equation}
Here $c_5=c_5(L,N)\ge0$, and
\begin{equation*}
c_0(N)=\frac{11N^{2}}{12\pi^{2}}\log L,\qquad
\bar B:=\max\{C_\beta+1,\,c_5\}\ge\frac{11N^{2}}{48\pi^{2}}\log L .
\end{equation*}
From here on $\bar B$ thus denotes $\max\{C_\beta+1,c_5\}$, which is not smaller than the
value $C_\beta+1$ given to $\bar B$ in Notation~\ref{0.2:not-glossary}.
These are values independent of $\gamma$ and of $K$, fixed after $(N,L,\mathfrak{p})$, and they
are never compared with a power of $L$. In particular, the condition $L^{2}\ge1+\bar B$ is not
imposed, at any $N$.

We use the following consequences. Indices range over $[0,K]$. In the lines $(\mathrm{F}3)^{\sharp}$,
$(\mathrm{TS}1)$ and $(\mathrm{TS}2)$ we take $j\le n\le k\le K$. Here $\beta_0$ is the constant in
the transport inequalities $(\mathrm{TS}1)$, $(\mathrm{TS}2)$, and $\hat\varepsilon_j$ the scale
transported with $\varepsilon_j$ in them, both fixed in the proposition, cited below, that
replaces the condition $L^{2}\ge1+\bar B$ by a ceiling on the coupling.
The lemma comparing the scales $\delta_n$ along the flow uses the floors
\begin{equation*}
x_\delta:=\max\{e,\,2c_5,\,4p_0c_5\},\qquad x_\beta:=\max\{e,\,2p_0c_5/\beta_0\}.
\end{equation*}
\begin{equation}\label{4.36:eq-sun-conventions-b}
\begin{aligned}
(\mathrm{Fl}1)\quad & x_n\ge x_K+\lambda(K-n)-c_5\ge\gamma^{-2}-c_5\qquad(n\le K);\\
(\mathrm{Fl}2)\quad & x_n\le x_k+B(k-n)\ \ (n\le k),\qquad x_n\le x_j+c_5\ \ (j\le n);\\
(\mathrm{Fl}3)\quad & \delta_n\ge\tfrac12\,\delta_k\,(1+B(k-n))^{-1/2}\qquad
(n\le k,\ x_k\ge x_\delta);\\
(\mathrm{F}3)^{\sharp}\quad & x_j\le(1+(n-j))x_n,\quad g_n\le(1+(n-j))^{1/2}g_j,\quad
\delta_n\ge\tfrac12(1+(k-n))^{-1/2}\delta_k;\\
(\mathrm{TS}1)\quad & \delta_nL^{-(n-j)}\le(1+\beta_0)\delta_j,\qquad
\varepsilon_nL^{-(n-j)}\le(1+\beta_0)\varepsilon_j,\\
& \hat\varepsilon_nL^{-(n-j)}\le(1+\beta_0)\hat\varepsilon_j;\\
(\mathrm{TS}2)\quad & \varepsilon_nL^{-2(n-j)}\le(1+\beta_0)\varepsilon_j,\qquad
\hat\varepsilon_nL^{-2(n-j)}\le(1+\beta_0)\hat\varepsilon_j .
\end{aligned}
\end{equation}
Inequality $(\mathrm{Fl}1)$ uses only the lower member of \eqref{4.36:eq-sun-conventions-3} and
$g_K\le\gamma$; the bound $x_K\ge g^{-2}$ is not used. Consequently every floor of the form
``$x_n\ge x_\bullet$ for all $n\le K$'', with $x_\bullet>0$, follows from
$\gamma\le(x_\bullet+c_5)^{-1/2}$. In $(\mathrm{Fl}3)$, which holds under $x_k\ge x_\delta$, the
factor $\tfrac12$ is the lower bound on $(\log x_n/\log x_k)^{p_0}$ given by the lemma comparing
the scales $\delta_n$ along the flow, and the factor $(1+B(k-n))^{-1/2}$ bounds $g_n/g_k$ from
below by $(\mathrm{Fl}2)$ and $x_k\ge1$.
The inequalities $(\mathrm{F}3)^{\sharp}$, $(\mathrm{TS}1)$ and $(\mathrm{TS}2)$ hold under the
coupling ceiling below. They are taken from the proposition that replaces the
condition $L^{2}\ge1+\bar B$ by a ceiling on the coupling.

\emph{The coupling ceiling.} The primary choice is
\begin{equation*}
\gamma\le\gamma_{\bar B}:=\bigl(\max\{\bar B,\,x_\delta,\,x_\beta\}+c_5\bigr)^{-1/2}.
\end{equation*}
Under this ceiling, $x_n\ge\max\{\bar B,x_\delta,x_\beta\}$ for every $n\le K$, and
$(\mathrm{F}3)^{\sharp}$ holds. In the flow factors below, the
letter $\mathfrak{f}\in\{1,\bar B\}$ takes the value $\mathfrak{f}=1$. The value
$\mathfrak{f}=\bar B(N,L,\mathfrak{p})$ is equally admissible in every flow factor of the later
estimates that is obtained through $(\mathrm{Fl}2)$ and $(\mathrm{Fl}3)$, among them the factor
$\Theta_m$, defined where it is used. The thresholds of the local chart on the cube $\square_0$
(defined where it is used) and the floor on the inverse coupling that accompanies the membership
bound, the bound in the membership norm $\|\cdot\|_{\star}$, however, use $(\mathrm{F}3)^{\sharp}$,
and hence this ceiling, unless, throughout, the inequalities used in place of
$(\mathrm{F}3)^{\sharp}$ are obtained by the alternative argument, which does not use this
ceiling.

\emph{Parameter conditions and thresholds.} In the conditions below, $c_d$ is the constant of the
lattice-animal bound; $\kappa_0$ is the power of $g_j$ in the bound \cite[(2.31)]{B-CMP119} on the
renormalisation-operation pieces $\mathbf{R}^{(j)}$ of the effective action; $C_{\mathrm{loc}}$ is
the constant of the localisation bound on the operators of the step, taken at its value for
$SU(N)$; $C_H$ is the constant of the bound on the response of the minimiser to a one-bond
variation, and $C_P$ the constant of the companion bound used with it;
$C_\chi'=C_\chi'(d,L)$ is the constant, depending on $d$ and $L$ only, of the bound on the
cut-off functions $\chi_{\square}$; $c_\star$ is the radius constant of the membership bound;
$c_B$ is the constant in the radius $c_B\delta_m$ of the data ball $\mathcal{A}'$ of the local
chart on $\square_0$; and $\delta_0$
is the decay constant of the weights $e_y$. The constant $C_1^{(102)}$ is that of
\cite[(172)]{B-CMP102b}, and $a_1'$ is the further constant entering (t8), fixed at the same
stage as $C_1^{(102)}$. The auxiliary parameters satisfy
\begin{equation*}
\begin{aligned}
(\mathrm{Q})\quad & q:=\min\{q_0,q_1\}\ge p_0+3r+2,\quad p_0\ge5r,\quad r\ge2,\\
& \kappa\ge\max\{10,\,8c_d+8\},\quad \kappa_0\ge7 .
\end{aligned}
\end{equation*}
The thresholds are labelled (t1)--(t6), (t8) and (t9), together with (t8)$^{+}$. Of these, (t3), a ceiling on
the coupling formed from
inverse-coupling floors, among them the floors $x_E$, $x_R$, $x_B$ described below, and (t5), a
further condition independent of $N$ imposed together with $(\mathrm{Q})$, are not displayed
here.
\begin{itemize}
\item[(t1)] $p_0(g_k)\ge8C_{\mathrm{loc}}$, with the degree function $p_0(\cdot)$.
\item[(t2)] The radii
\begin{equation}\label{4.36:eq-sun-conventions-4}
R:=\min\Bigl\{\frac{c_\star\delta_k}{4C_\chi'C_H},\ \frac{\varepsilon_k}{16C_{\mathrm{loc}}},\
\frac{r_0}{2}\Bigr\},\qquad r_m:=\min\Bigl\{r_0,\ \frac{c_B\delta_m}{8C_H}\Bigr\}.
\end{equation}
\item[(t4)] The lower bound on $M=L^{m'}$ imposed with the auxiliary parameters
$\mathfrak{p}$; it includes $M\ge L^{4}$ and $\delta_0M\ge16$.
\item[(t6)] $L^{2}\ge20(1+\beta_0)$.
\item[(t8)] The radius bound
\begin{equation*}
r_0\le\min\bigl\{\tfrac12,\ C_1^{(102)}a_1'/4,\ (4C_H)^{-1},\ (2C_P)^{-1}\bigr\}.
\end{equation*}
\item[(t8)$^{+}$] $g_k<r_0$.
\item[(t9)] $R_j\ge L^{2}$. This contains the bound $R_j\ge7$ of the lemma on working cubes,
since the blocking factor $L$ is an odd integer greater than $1$, so that $L\ge3$ and
$L^{2}\ge9$; with (t9) this gives $R_j\ge9$, which exceeds $7$.
\end{itemize}

\emph{The radius $r_0$.} The auxiliary radius $r_0$ is one fixed positive number not exceeding
the minimum in (t8). It is fixed after $(N,L)$, together with $C_H$, $C_P$, $C_1^{(102)}$ and
$a_1'$, and before $\gamma$. The inverse-coupling floors $x_E$, $x_R$, $x_B$ of the bounds on the
regular, renormalisation-operation and boundary pieces $\mathbf{E}^{(j)}$, $\mathbf{R}^{(j)}$,
$\mathbf{B}^{(i)}$ of the effective action, and the constant $C_q(N)$ entering the common constant
of these bounds, depend on $r_0$ only through its numerical value. Any $r_0$ satisfying (t8) is
admissible, and, as the definitions of these floors show where they are given, a smaller $r_0$
only enlarges them.

\emph{Letters with two meanings.}
\begin{itemize}
\item The bare letter $B$ is the flow slope of \eqref{4.36:eq-sun-conventions-3}. The letter
with an argument, $B(\zeta)$, is a logarithmic variable introduced where it is used.
\item $\bar B$ is $\max\{C_\beta+1,c_5\}$, as above, and not the value $C_\beta+1$ of
Notation~\ref{0.2:not-glossary}.
\item $\mathfrak{f}\in\{1,\bar B\}$ is the flow letter above. The symbol
$\mathfrak{b}(\mathsf{N})$ denotes the blocks meeting a bond set $\mathsf{N}$.
\item $r_0$ is the auxiliary radius bounded in (t8); it is not the constant depth offset of the
extraction depth in clause (iv) of Theorem~0.
\item $R$ is the radius of (t2); it is not the ball radius of the observables.
\item $u$ is the free link (also $u_1$) and, in sums, a level gap; which meaning is in force is
said where it is used.
\item $s$ is a scale index and, where this is said, $s:=\log x_k$; neither meaning is the depth
$s$ below the cutoff of Notation~\ref{0.2:not-glossary}.
\item $c$ is a cube in the lattice-animal bound. The symbol $c(q)$ is the Wilson coefficient,
and in an expression $x(c)$ the argument $c$ is a positive number.
\item $e$ is Euler's number. The same letter is also the order of a derivative in one estimate
and the base letter of the weights $e_y$.
\item $C_1^{(102)}$ is the constant of \cite[(172)]{B-CMP102b} entering (t8), while $C_1$ alone
is the constant in $\alpha_{1,j}$.
\item $m$ is the working scale $m(j,n)$ of the working cubes and, in the bounds on boundary
terms, the level to which a boundary term is assigned. No working cube occurs where the second
meaning is in force. Neither meaning is the torus exponent $m$.
\end{itemize}

\emph{Further abbreviations.} We put $\mathfrak{c}:=c_B\delta_m$, the radius of the data ball
$\mathcal{A}'$; $\tau_{n'}:=6+\log_LR_{n'}$; and $c_{\mathrm{exp}}:=\tfrac52L^{2}$.
\end{definition}

\begin{proof}
The flow bound \eqref{4.36:eq-sun-conventions-3} with its constants, including the lower bound
on $\bar B$ displayed above, is taken from
Theorem~\ref{4.1:thm-clause-zero} and Definition~\ref{1.2:def-flow-constants}; $(\mathrm{Fl}3)$ is
taken from the lemma comparing the scales $\delta_n$ along the flow; and $(\mathrm{F}3)^{\sharp}$,
$(\mathrm{TS}1)$, $(\mathrm{TS}2)$ are taken, under the ceiling $\gamma\le\gamma_{\bar B}$, from
the proposition that replaces the condition $L^{2}\ge1+\bar B$ by a ceiling on the coupling. We
prove the remaining assertions.

(e1). For $n\ge1$ we have $n!=\prod_{i=2}^{n}i\ge2^{n-1}$, so $1/n!\le2^{-(n-1)}$, with strict
inequality for $n=3$. Hence
\begin{equation*}
e=\sum_{n\ge0}\frac1{n!}<1+\sum_{n\ge1}2^{-(n-1)}=1+1+\sum_{n\ge1}2^{-n}=3 .
\end{equation*}
Moreover $e\ge1+1+\tfrac12>2$. From $2<e$ we get $\log2<1$. From $e<3$ we get $e-1<2$. From
$e-1>1$ we get $1/(e-1)<1$. Finally, $e/(e-1)<2$ is equivalent to $e<2e-2$, that is, to $e>2$.

(e2). The function $x\mapsto e^{x}$ is convex and $1+x$ is its tangent at $0$, so
$e^{x}\ge1+x$. Putting $x=\log y$ for $y>0$ gives $y\ge1+\log y$. For $x\ge0$ every term of
$\sum_{n\ge0}x^{n}/n!$ is non-negative, and keeping the term $n=p$ gives $e^{x}\ge x^{p}/p!$.
The inequality $L^{u}-1\ge u(L-1)$ is Bernoulli's inequality $(1+h)^{u}\ge1+uh$ for $h=L-1\ge0$.
For integer $u\ge0$ it follows from the binomial expansion of $(1+h)^{u}$, all of whose
terms are non-negative. For real $u\ge1$, the function $\varphi(h)=(1+h)^{u}-1-uh$ on
$[0,\infty)$ satisfies $\varphi(0)=0$ and $\varphi'(h)=u\bigl((1+h)^{u-1}-1\bigr)\ge0$, because
$1+h\ge1$ and $u-1\ge0$; hence $\varphi\ge0$.
For $p\ge1$ the function $t\mapsto t^{p}$ is convex on $[0,\infty)$, so
$((a+b)/2)^{p}\le(a^{p}+b^{p})/2$; multiplying by $2^{p}$ gives the stated bound. For $a>0$ we
have $0<e^{-a}<1$, and the geometric series gives
$\sum_{\ell\ge0}e^{-a\ell}=(1-e^{-a})^{-1}$. If moreover $a\ge\log2$, then $e^{-a}\le\tfrac12$,
so the sum is at most $2$. If $L^{4}\ge2$, then $L^{-4}\le\tfrac12$ and
$(1-L^{-4})^{-1}\le2$.

(e3). The function $s\mapsto e^{s}-1$ is convex on $[0,1]$. It vanishes at $s=0$ and equals
$e-1$ at $s=1$. A convex function lies below its chord, so $e^{s}-1\le(e-1)s$ on $[0,1]$.

(e4). Both sides are positive, so the claim is equivalent to $(2+\ell)e^{-2\ell}\le2$, that is,
to $1+\ell/2\le e^{2\ell}$. For $\ell\ge0$, (e2) gives $e^{2\ell}\ge1+2\ell\ge1+\ell/2$.

The matrix exponential. For $Z\in M_N(\mathbb{C})$ we have
$e^{Z}-\mathbb{1}=\sum_{n\ge1}Z^{n}/n!$, and $|Z^{n}|\le|Z|^{n}$ because the operator norm is
submultiplicative. Hence
\begin{equation*}
|e^{Z}-\mathbb{1}|\le\sum_{n\ge1}\frac{|Z|^{n}}{n!}=e^{|Z|}-1 .
\end{equation*}
Now let $X\in\mathfrak{g}$. Since $X$ is Hermitian, there are a unitary matrix $v$ and real
numbers $\theta_1,\dots,\theta_N$ with $X=v\operatorname{diag}(\theta_1,\dots,\theta_N)v^{-1}$,
and $|X|=\max_i|\theta_i|$. Then
\begin{equation*}
e^{iX}-\mathbb{1}=v\operatorname{diag}\bigl(e^{i\theta_1}-1,\dots,e^{i\theta_N}-1\bigr)v^{-1}.
\end{equation*}
The operator norm is invariant under conjugation by a unitary matrix, and the operator norm of a
diagonal matrix is the largest modulus of its diagonal entries, so
$|e^{iX}-\mathbb{1}|=\max_i|e^{i\theta_i}-1|$. For real $\theta$ we have
$e^{i\theta}-1=\int_0^{\theta}ie^{it}\,dt$, hence $|e^{i\theta}-1|\le|\theta|$. Therefore
$|e^{iX}-\mathbb{1}|\le\max_i|\theta_i|=|X|$. This proves \eqref{4.36:eq-sun-conventions-2}.

The adjoint action. Let $v$ be unitary. Conjugation by $v$ preserves the trace, so
$\operatorname{Ad}(v)$ maps $\mathfrak{g}^{c}$ into itself. For $Y\in\mathfrak{g}^{c}$ we have
$|vYv^{-1}|\le|v|\,|Y|\,|v^{-1}|=|Y|$, because $|v|=|v^{-1}|=1$. Applying the same bound to
$Y=v^{-1}(vYv^{-1})v$ gives the reverse inequality, so $\operatorname{Ad}(v)$ is an isometry of
$(\mathfrak{g}^{c},|\cdot|)$ for every $N$.

Now let $X\in\mathfrak{g}$ and $A=iX$. The matrix function $F(t)=e^{tA}Ye^{-tA}$ and the function
$t\mapsto e^{t\operatorname{ad}_A}Y$ both solve $F'(t)=AF(t)-F(t)A$ with $F(0)=Y$. Hence they
coincide, and at $t=1$ this gives the first equality in \eqref{4.36:eq-sun-conventions-1}.
Since $|\operatorname{ad}_{iX}Y|=|XY-YX|\le2|X|\,|Y|$, induction gives
$|(\operatorname{ad}_{iX})^{n}Y|\le(2|X|)^{n}|Y|$. The term $n=0$ of
$e^{\operatorname{ad}_{iX}}Y=\sum_{n\ge0}(\operatorname{ad}_{iX})^{n}Y/n!$ is $Y$. Estimating the
remaining terms one by one gives the inequality in \eqref{4.36:eq-sun-conventions-1}. The sum
equals $e^{2|X|}-1$ times $|Y|$, and no constant depending on $N$ enters.

Couplings and scales. Since $0<g_j\le\gamma<1$, we have $x_j=g_j^{-2}>1$. The identity for
$\delta_j$ is the definition of $p_0(\cdot)$ inserted in $\varepsilon_j$:
\begin{equation*}
\frac{A_1}{A_0}\,g_jA_0(\log x_j)^{p_0}=A_1g_j(\log x_j)^{p_0}.
\end{equation*}
By definition $R_j=L^{t}$ for some $t\in\mathbb{Z}_{\ge0}$, so $R_j\ge1$ and
$\log_LR_j=t\in\mathbb{Z}_{\ge0}$. Finally, the value of $c_0(N)$ gives
\begin{equation*}
\frac{c_0(N)}{4}=\frac{11N^{2}}{48\pi^{2}}\log L ,
\end{equation*}
so that the lower bound on $\bar B$ displayed above reads $\bar B\ge c_0(N)/4$.

$(\mathrm{Fl}1)$. We apply the lower member of \eqref{4.36:eq-sun-conventions-3} with $(i,j)=(n,K)$,
$0\le n\le K$. It reads $\lambda(K-n)-c_5\le x_n-x_K$, which is the first inequality of
$(\mathrm{Fl}1)$. Since $\log L>0$, we have $c_0(N)\ge0$; as $\lambda^{\sharp}\ge c_0(N)/4$, this
gives $\lambda^{\sharp}\ge0$, hence $\lambda\ge0$. Since also $K-n\ge0$, we have
$\lambda(K-n)\ge0$. Since $0<g_K\le\gamma$, we have $x_K=g_K^{-2}\ge\gamma^{-2}$. Together these
give the second inequality. Only the lower member of the flow bound and $g_K\le\gamma$ were used.

Now let $x_\bullet>0$ and $\gamma\le(x_\bullet+c_5)^{-1/2}$. Then $\gamma^{-2}\ge x_\bullet+c_5$,
and $(\mathrm{Fl}1)$ gives $x_n\ge\gamma^{-2}-c_5\ge x_\bullet$ for every $n\le K$.

$(\mathrm{Fl}2)$. For $n\le k$, the upper member of \eqref{4.36:eq-sun-conventions-3} with
$(i,j)=(n,k)$ gives $x_n-x_k\le B(k-n)$. For $j\le n$, the lower member with $(i,j)$ replaced by
$(j,n)$ gives $x_j-x_n\ge\lambda(n-j)-c_5\ge-c_5$, because $\lambda(n-j)\ge0$. Hence
$x_n\le x_j+c_5$.

The coupling ceiling. We have $\max\{\bar B,x_\delta,x_\beta\}\ge e>0$, because
$x_\delta\ge e$. We apply the floor statement just proved with
$x_\bullet=\max\{\bar B,x_\delta,x_\beta\}$. It shows that $\gamma\le\gamma_{\bar B}$ implies
$x_n\ge\max\{\bar B,x_\delta,x_\beta\}$ for every $n\le K$.
\end{proof}

\begin{definition}[Regions, polymers and cubes of a history in four dimensions]
\label{4.36:not-history-geometry}
Throughout, $d=4$ and the level $k$ is fixed. The inverse couplings $x_j=g_j^{-2}$, the quantities
$\varepsilon_j$, the radii $R_j$ with their exponent $r$, the lengths $\ell_n$ and the threshold
conditions are those of Notation~\ref{4.36:not-sun-conventions}. Lengths are measured in units of
scale $k$; $\ell_n$ ($n\le k$) is the side of a block of scale $n$, and $\ell_k=1$. The constant
$\delta_0>0$ is the decay rate of the weights $e_y$ defined below; it enters the threshold
conditions. None of the objects
below depends on the group $SU(N)$. We write $\operatorname{dist}_\infty$, $\operatorname{diam}_\infty$
and $B_\infty(z,\cdot)$ for distance, diameter and balls centred at $z$ in the maximum norm,
$\operatorname{dist}_k$ for the distance between sets of the lattice of scale $k$, measured in
scale-$k$ units, $A^{c}$ for the complement of a set $A$,
and $(x)_+:=\max\{x,0\}$.

\emph{The marked plaquette.} Let $p$ be a plaquette of $\mathbb{T}^{(k)}$ and $b_1\subset\partial p$
a marked bond. The \emph{free link} is $u_1:=V_k(b_1)$, and we write $V=(u_1,V_{\mathrm{off}})$,
where $V_{\mathrm{off}}$ collects the values of $V_k$ on the bonds other than $b_1$. We put
\begin{equation*}
R_*:=20L^2M\max\{R_k,R_{k+1}\}.
\end{equation*}
A history $h$ is \emph{clean} if $\operatorname{dist}_k(p,\Lambda_k(h)^{c})\ge R_*$ (the region
$\Lambda_k(h)$ is introduced below). Attached to the marked plaquette are the set $\mathcal{N}$,
the functions $\chi_{\square}$ indexed by cubes $\square$, the $u_1$-window $W(V_{\mathrm{off}})$,
an open set of values of $u_1$ depending on $V_{\mathrm{off}}$, the deep set $\mathcal{D}$, the
indicator $\mathbb{1}_{G(h)}$ of the event $G(h)$ and the localisation constant $C_{\mathrm{loc}}$,
which enters the threshold conditions; they belong to the localisation analysis of the marked
plaquette, their definitions are not needed below, and $\varrho:=\varepsilon_k/(4C_{\mathrm{loc}})$.
The window always carries its argument $V_{\mathrm{off}}$; the bare letter $W$ denotes an arrested
region below.

\emph{Regions of a history.} A history $h$ carries regions $\Omega_n=\Omega_n(h)$, $n\le k$, and
$\Lambda_j=\Lambda_j(h)$, $j\le k$. We put
\begin{equation*}
\Gamma_n:=\Omega_n\setminus\Omega_{n+1}\quad(n<k),\qquad \Gamma_k:=\Omega_k,\qquad
Z_j:=\Lambda_j^{c}.
\end{equation*}
Each $\Lambda_j$ is a union of $MR_j$-cubes, that is, of cubes of side $MR_j\ell_j$
\cite{B-CMP119}; $\Lambda_j^{0}$ is $\Lambda_j$ with one
layer of $MR_j$-cubes removed, and $\Lambda_j^{\sim -1}$ is the region of
\cite[(2.30)--(2.31)]{B-CMP119}.

\emph{Arrested regions.} An arrested region $W$ is arrested at step $k_W$ after the integrations
$j=h_W,\dots,j_W-1$. Its data also comprise an integer $N_0^{W}$ and a region $Z''$, both defined
with the construction of arrested regions; their definitions are not needed below. We put
$k_0:=k_W-N_0^{W}$, and $\mathfrak{E}(W):=(k_0,j_W)$ is the interval of levels $i$ with
$k_0<i<j_W$. The number of stages of the operation $\mathbf{R}$
attempted on $W$ is $\mathfrak{N}_W:=j_W-h_W$.

\emph{Polymers.} The classes $\mathbf{D}_j$ and $\mathbf{D}_j^{R}$ are the classes of polymers of
scale $j$ of \cite{B-CMP119}. Every $X\in\mathbf{D}_j^{R}$ belongs to $\mathbf{D}_j$ and satisfies
$X\subset\Lambda_j^{\sim -1}$. For a polymer $X$ we put
$n_X:=\min\{n:\ X\cap\Gamma_n\neq\emptyset\}$; for $X\in\mathbf{D}_j^{R}$ every point $z_X$ of
$X\cap\Gamma_{n_X}$ lies in $\Lambda_j^{\sim -1}$. The function $d_j(X)$ is Ba{\l}aban's
\cite{B-CMP109}; it satisfies
\begin{equation*}
d_j(X)\ \ge\ \frac{\operatorname{diam}_\infty X}{M\ell_j}-2 ,
\end{equation*}
and sums over $\ell$ below run over the set of values of $d_j$.

\emph{Blocks and cubes.} For $y\in\Gamma_n$, $\Delta(y)$ is the block of scale $n$ labelled by $y$;
it has side $\ell_n$. The partition $\pi_n$ is the partition into $M$-cubes of scale $n$, of side
$M\ell_n$. For a
cube $\square\in\pi_n$ and an integer $r\ge0$ (here $r$ counts cube-widths and is not the exponent
of the radii), $\widetilde{\square}^{r}$ is $\square$ enlarged by $r$
cube-widths on every side, a cube of side $(2r+1)M\ell_n$. For $j\le n$ we put
$m=m(j,n):=\max\{j,n-2\}$, and for $\square\in\pi_m$ we put $\square_0:=\widetilde{\square}^{5}$. The
tent functions $h_z$ of scale $j$ are those of \cite{B-CMP119} (not to be confused with a history
$h$); they satisfy $\operatorname{supp}h_z\subset B_\infty(z,\ell_j)$.

\emph{Weights.} The function $d(\cdot,\cdot)$ is the distance of \cite[(2.46)]{B-CMP96}, and
$e_y:=e^{-\delta_0 d(y,b_1)/8}\le 1$. For a set $\mathsf{N}$ of bonds or of points,
$\mathfrak{b}(\mathsf{N})$ is the set of those $y$ for which $\mathsf{N}$ meets the block
$\Delta(y)$ (a set of blocks is identified with its union), and
\begin{equation*}
\hat{e}_{\mathsf{N}}:=\sum_{y\in\mathfrak{b}(\mathsf{N})}e_y^{1/2},\qquad
\check{e}_{\mathsf{N}}:=\max_{y\in\mathfrak{b}(\mathsf{N})}e_y ;
\end{equation*}
the set $\mathsf{N}^{+}$ consists of the blocks at adjacency distance at most $2$ from the blocks
$\Delta(y)$, $y\in\mathfrak{b}(\mathsf{N})$, the adjacency distance of two blocks being the least
number of passages between adjacent blocks along a chain joining them. In particular $X^{+}$ and
$(\square_0)^{+}$ are defined.

\emph{Animal sums.} For $a>0$ we put
\begin{equation*}
\mathfrak{a}(a):=\sup_{c}\sum_{X\in\mathbf{D}_j,\ X\supset c}e^{-a\,d_j(X)},\qquad
\mathfrak{s}(a):=\sum_{\ell}e^{-(a-c_d)\ell},
\end{equation*}
the supremum being over cubes $c$ of scale $j$, where $c_d$ is the animal constant of
item~(vii) below; $\mathfrak{a}(a)$ depends on $j$, which is suppressed from the notation.

\emph{Inputs.} The following properties of these objects are used as stated and are not proved
here. Item~(vi) is a bound of \cite{B-CMP122b}; item~(vii) is taken by statement; the other
items are properties of the
regions, blocks and polymers of a history established with their construction.
\begin{enumerate}
\item[(i)] (Decay of the weights on a layer.) For $\alpha\in[1/288,1/8]$ and $n\le k$,
\begin{equation*}
\sigma_n^{(\alpha)}:=\sum_{y\in\Gamma_n}e^{-\alpha\delta_0 d(b_1,y)}\ \le\ C_\sigma\,e^{-(k-n)} ,
\end{equation*}
where the constant $C_\sigma$ and the condition under which the bound holds are those of the
bound on the decay of the weights on a layer in the localisation analysis of the marked plaquette,
where this bound is established.
\item[(ii)] (Blocks, polymers across regions, clean histories.) (a) Adjacent blocks differ in
scale by at most one. (b) If $X$ is a connected union of
scale-$j$ $M$-cubes meeting $\Omega_a$ and $\Omega_b^{c}$, with $b<a$, then
$d_j(X)\ge L^{a-1-j}-2$. (c) If $h$ is clean, then for $y\in\Gamma_n$ with $n<k$
\begin{equation*}
d(b_1,y)\ \ge\ R_*+M(k-1-n),
\end{equation*}
and for $y\in\Gamma_k$ one has $d(b_1,y)\ge\min\{|y-b_1|,R_*\}$.
\item[(iii)] (The enlarged cube around a point.) Let $j\le n\le k$ and
$x\in\Lambda_j^{\sim -1}\cap\Gamma_n$. Let $m=\max\{j,n-2\}$, let $\square\in\pi_m$ be the cube
containing $x$, and $\square_0=\widetilde{\square}^{5}$. Then
\begin{equation*}
\square_0\subset\Lambda_j\cap\Omega_m\cap\Omega_{n+2}^{c},\qquad
\square_0\subset B_\infty(x,6M\ell_m),
\end{equation*}
and the blocks of $(\square_0)^{+}$ have scale in $[m,n+1]$ and side at most $L^3\ell_m$. This holds
under the conditions $L\ge7$, $M\ge L^4$, $R_j\ge7$, and nothing in it depends on $N$. The proof of
this statement uses of $x$ only that $x\in\Gamma_n$ and
$\operatorname{dist}_\infty(x,\Lambda_j^{c})\ge MR_j\ell_j$, so its conclusion holds for every $x$
with these two properties.
\item[(iv)] (Properties of the regions.) The regions of a history satisfy:
$\Lambda_j^{\sim -1}\subset\Lambda_j\subset\Omega_j$,
and $Z_j$ is increasing in $j$;
\begin{equation*}
\Lambda_i\ \supseteq\ \Omega_{i+1}\setminus\bigcup\{W:\ i\in\mathfrak{E}(W)\};
\end{equation*}
for every arrested region $W$,
\begin{equation*}
W\subset\Omega_{k_W+1}^{c},\qquad
\operatorname{dist}_\infty(W,\Omega_{k_W+1})\ \ge\ 3MR_{k_W+1}\ell_{k_W+1};
\end{equation*}
for the native region $\Omega^{\natural}_{i+1}$ of the construction,
\begin{equation*}
\operatorname{dist}_\infty(\Omega^{\natural}_{i+1},Z_i)\ \ge\ 8MR_{i+1}\ell_{i+1}.
\end{equation*}
\item[(v)] (The radii.) The radii satisfy: (a) $R_{j+1}\le R_j$; (b) $R_j/R_{j+1}\in\{1,L\}$, and
the value $L$ does not occur for two consecutive $j$; (c) $\log x_j\le R_j<L(\log x_j)^{r}$; (d)
$N_0\le 2+\log_L R_k$, where $N_0$ is the offset integer of the construction of arrested regions,
read at level $k$.
\item[(vi)] (The stage count.) $\mathfrak{N}_W\le R_{k_W}$. This is the bound $N\le R_k$ stated in
\cite{B-CMP122b} for the number $N$ of stages of the operation $\mathbf{R}$; there $N$ counts
stages and is not the $N$ of $SU(N)$.
\item[(vii)] (The animal count.) The values of $d_j$ are non-negative integers, and for every
scale-$j$ cube $c$ and every $\ell$,
\begin{equation*}
\#\{X\in\mathbf{D}_j:\ X\supset c,\ d_j(X)=\ell\}\ \le\ e^{c_d\ell},
\end{equation*}
with $c_d$ the animal constant fixed with the threshold conditions; this is a combinatorial fact
about connected sets of cubes in $d=4$, and $e^{c_d\ell}$ is the same factor that enters item~(ix).
\item[(viii)] (The cube count.) A polymer $X\in\mathbf{D}_j$ consists of at most
$2^{d}(3d_j(X)+1)$ cubes of $\pi_j$.
\item[(ix)] (The counting bound.) Let $y\in\Gamma_{n'}$. If $X\in\mathbf{D}_j$ and
$y\in\mathfrak{b}(X^{+})$, then $X$
meets a block of scale in $[n'-2,n'+2]$, and for every $\ell$
\begin{equation*}
\#\{X\in\mathbf{D}_j:\ y\in\mathfrak{b}(X^{+}),\ d_j(X)=\ell\}\ \le\
7^4L^{4(n'+2-j)_+}e^{c_d\ell}.
\end{equation*}
If moreover $X\subset\Omega_j$, as for the polymers of the pieces $\mathbf{E}^{(j)}$ and
$\mathbf{R}^{(j)}$ of the effective action \cite[(2.23)]{B-CMP119}, then $j\le n'+2$ and
$n_X\le n'+2$.
\end{enumerate}

\emph{Consequences.} For every $a\ge\kappa/8$,
\begin{equation}\label{4.36:eq-history-geometry-1}
\mathfrak{a}(a)\ \le\ \mathfrak{s}(a)\ \le\ 2 ,
\end{equation}
and every polymer $X\in\mathbf{D}_j$ contains at most
\begin{equation}\label{4.36:eq-history-geometry-2}
c_\#\,(1+d_j(X))\,M^4
\end{equation}
sites of the lattice of scale $j$, with a constant $c_\#\le 48$.
\end{definition}

\begin{proof}
We prove \eqref{4.36:eq-history-geometry-1}. Fix $a\ge\kappa/8$ and a scale-$j$ cube $c$. All terms
being non-negative, the sum over polymers $X\in\mathbf{D}_j$ containing $c$ may be ordered by the
value $\ell=d_j(X)$. By item~(vii), for each $\ell$ at most $e^{c_d\ell}$ polymers contribute, each
with the weight $e^{-a\ell}$. Hence
\begin{equation*}
\sum_{X\in\mathbf{D}_j,\ X\supset c}e^{-a\,d_j(X)}
=\sum_{\ell}e^{-a\ell}\,\#\{X\in\mathbf{D}_j:\ X\supset c,\ d_j(X)=\ell\}
\le\sum_{\ell}e^{-(a-c_d)\ell}=\mathfrak{s}(a).
\end{equation*}
The right-hand side does not depend on $c$, so taking the supremum over $c$ gives
$\mathfrak{a}(a)\le\mathfrak{s}(a)$.

The threshold conditions of Notation~\ref{4.36:not-sun-conventions}, which fix $c_d$, together
with \eqref{4.36:eq-sun-conventions-a}, give $\kappa\ge 8c_d+8$, that is, $\kappa/8-c_d\ge1$.
Since $a\ge\kappa/8$,
\begin{equation*}
a-c_d\ \ge\ \frac{\kappa}{8}-c_d\ \ge\ 1\ >\ \log 2 .
\end{equation*}
By item~(vii) the values of $d_j$ are non-negative integers, so the sum defining $\mathfrak{s}(a)$ is
bounded by the full geometric series, and since $e^{-(a-c_d)}<1/2$,
\begin{equation*}
\mathfrak{s}(a)\ \le\ \sum_{\ell=0}^{\infty}e^{-(a-c_d)\ell}\ =\ \frac{1}{1-e^{-(a-c_d)}}
\ <\ \frac{1}{1-1/2}\ =\ 2 .
\end{equation*}
This proves \eqref{4.36:eq-history-geometry-1}.

We prove \eqref{4.36:eq-history-geometry-2}. By item~(viii) with $d=4$, a polymer
$X\in\mathbf{D}_j$ consists of at most
\begin{equation*}
2^4\bigl(3d_j(X)+1\bigr)\ =\ 48\,d_j(X)+16\ \le\ 48\bigl(1+d_j(X)\bigr)
\end{equation*}
cubes of $\pi_j$. Each cube of $\pi_j$ has side $M\ell_j$, while the lattice of scale $j$ has
spacing $\ell_j$; in dimension $4$ such a cube therefore contains $M^4$ sites of that lattice.
Multiplying the two counts, $X$ contains at most $48(1+d_j(X))M^4$ sites of scale $j$, which is
\eqref{4.36:eq-history-geometry-2} with $c_\#=48$.
\end{proof}

\begin{definition}[The probing direction along a unit generator]\label{4.36:not-probing-direction}
We use $G=SU(N)$, its Lie algebra $\mathfrak{g}$ of traceless Hermitian $N\times N$ matrices, the
complexification $\mathfrak{g}^{c}=\mathfrak{sl}(N,\mathbb{C})$, the normalised trace
$\operatorname{tr}=N^{-1}\operatorname{Tr}$ and the operator norm $|\cdot|$, all as in
Notation~\ref{4.36:not-sun-conventions}. We also use the marked bond $b_1$ and the splitting
$V=(u_1,V_{\mathrm{off}})$ of a configuration into its value $u_1$ at $b_1$ and its values
$V_{\mathrm{off}}$ on the other bonds, as in Notation~\ref{4.36:not-history-geometry}.

Let $u\in G$, let $\hat{T}\in\mathfrak{g}$ with $|\hat{T}|=1$, and let $\zeta\in\mathbb{C}$. Put
\begin{equation*}
T':=u\hat{T}u^{-1}.
\end{equation*}
The configuration \emph{moved along $\hat{T}$ at the free link} is
\begin{equation*}
V^{\zeta,\hat{T}}:=\bigl(u e^{i\zeta\hat{T}},V_{\mathrm{off}}\bigr).
\end{equation*}
The \emph{one-bond factor} $V'$ is the configuration equal to $e^{i\zeta T'}$ at $b_1$ and to
$\mathbb{1}$ on every other bond. We write $|V'-\mathbb{1}|:=\sup_b|V'(b)-\mathbb{1}|$, the supremum
over bonds of the operator norm; since $V'$ equals $\mathbb{1}$ off $b_1$, this is
$|V'(b_1)-\mathbb{1}|$.

For $|\operatorname{Re}\zeta|<\pi$, let $B(\zeta)$ be the one-bond field given by
\begin{equation*}
B(\zeta)(b_1):=\tfrac{1}{i}\log V'(b_1),
\end{equation*}
with the principal logarithm, and by $B(\zeta)(b)=0$ on every bond $b\neq b_1$.

Let $C_1^{(102)}$ and $a_1'$ denote the two positive constants of the smallness condition
\cite[(172)]{B-CMP102b}, whose product bounds $|V'-\mathbb{1}|$ in that condition. Here $r_0>0$
denotes the auxiliary radius. It is one fixed positive number, and
among the conditions imposed on it are the two ceilings
\begin{equation}\label{4.36:eq-probing-direction-1}
r_0\le\tfrac12,\qquad r_0\le\tfrac14\,C_1^{(102)}a_1'.
\end{equation}

Let $|\zeta|<2r_0$. By the \emph{one-bond response statement} we mean the following statement on the
response of the minimiser to a one-bond variation, which rests on the analysis of the minimiser in
\cite[(190)]{B-CMP102b}. To the variation along $\zeta T'$ at $b_1$ it assigns a response field,
which we denote $\mathcal{H}(\zeta T')$, such that
$U_k(V^{\zeta,\hat{T}})=e^{i\eta\mathcal{H}(\zeta T')}U_k(V)$ modulo gauge, with $\eta=L^{-k}$,
where $V=(u,V_{\mathrm{off}})$; the response field is $\mathfrak{g}^{c}$-valued for a
$\mathfrak{g}^{c}$-valued perturbation and $\mathfrak{g}$-valued for a $\mathfrak{g}$-valued one;
and the $\zeta$-derivative of the response field obeys, on every block $y$, a bound whose velocity
is $C_H e_y$ times the norm of $\partial_\zeta B(\zeta)(b_1)$, with $C_H$ the constant of that
bound. We put
\begin{equation*}
\mathcal{H}_{\zeta}:=\mathcal{H}(\zeta T').
\end{equation*}
With these definitions the following hold.
\begin{enumerate}
\item[(a)] $u e^{i\zeta\hat{T}}=e^{i\zeta T'}u$. Moreover $T'\in\mathfrak{g}$ and $|T'|=1$, and
\begin{equation*}
V^{\zeta,\hat{T}}=V'\cdot(u,V_{\mathrm{off}}),
\end{equation*}
where the product is taken bond by bond.
\item[(b)] $|V'-\mathbb{1}|\le e^{|\zeta|}-1$.
\item[(c)] For $|\operatorname{Re}\zeta|<\pi$ one has $B(\zeta)(b_1)=\zeta T'$. The field $B(\zeta)$
is traceless, and $\partial_\zeta B(\zeta)(b_1)=T'$. For $|\zeta|<2r_0$ the field
$\mathcal{H}_\zeta$ is $\mathfrak{g}^{c}$-valued, and it is $\mathfrak{g}$-valued at real $\zeta$.
\item[(d)] For $|\zeta|<2r_0$,
\begin{equation}\label{4.36:eq-probing-direction-a}
\begin{aligned}
|V'-\mathbb{1}|
&\le e^{|\zeta|}-1\le (e-1)|\zeta|<2(e-1)r_0\\
&\le \tfrac12(e-1)\,C_1^{(102)}a_1'<C_1^{(102)}a_1'.
\end{aligned}
\end{equation}
In particular $V'$ satisfies the smallness condition of \cite[(172)]{B-CMP102b} on the whole disc
$|\zeta|<2r_0$.
\item[(e)] As they are stated, the one-bond response statement and the estimates derived from it
depend on the normalisation of the direction only through three quantities: the bound on
$|V'-\mathbb{1}|$, the derivative $\partial_\zeta B=T'$, and the tracelessness of $B$. Each of the
three is independent of the choice of the unit direction $\hat{T}$ and of $N$: the bound
$e^{|\zeta|}-1$ involves $\hat{T}$ only through $|T'|=1$, and neither $\partial_\zeta B=T'$ nor
$\operatorname{tr}B=0$ introduces a constant; in particular no power of
$\mathsf{n}=(N-1)^{1/2}$ arises. The one place where a norm of
the direction could enter is the velocity in the bound on $\partial_\zeta\mathcal{H}_\zeta$ of the
one-bond response statement. On a block $y$ that velocity is
\begin{equation}\label{4.36:eq-probing-direction-2}
v(y)=C_H e_y\,|T'|=C_H e_y,
\end{equation}
where $C_H$ is the constant of that bound and $e_y=e^{-\delta_0 d(y,b_1)/8}$ is the decay weight
of the block $y$ relative to the marked bond $b_1$.
\end{enumerate}
\end{definition}

\begin{proof}
\emph{Clause (a).} Let $X$ be any $N\times N$ matrix. Since $(uXu^{-1})^n=uX^nu^{-1}$ for every
$n\ge0$, summing the exponential series term by term gives
\begin{equation*}
u e^{X}u^{-1}=\sum_{n\ge0}\frac{(uXu^{-1})^n}{n!}=e^{uXu^{-1}}.
\end{equation*}
With $X=i\zeta\hat{T}$ we have $uXu^{-1}=i\zeta T'$. Hence $u e^{i\zeta\hat{T}}u^{-1}=e^{i\zeta T'}$,
and multiplying on the right by $u$ gives $u e^{i\zeta\hat{T}}=e^{i\zeta T'}u$.

Since $u$ is unitary, $u^{-1}=u^{*}$. Hence
\begin{equation*}
(T')^{*}=(u^{-1})^{*}\hat{T}^{*}u^{*}=u\hat{T}u^{-1}=T',
\end{equation*}
so $T'$ is Hermitian. By cyclicity of the trace, $\operatorname{tr}T'=\operatorname{tr}\hat{T}=0$.
Thus $T'\in\mathfrak{g}$.

For unitary $u$, the map $X\mapsto uXu^{-1}$ is an isometry of $(\mathfrak{g}^{c},|\cdot|)$, as
stated in Notation~\ref{4.36:not-sun-conventions}. Therefore $|T'|=|\hat{T}|=1$.

Finally, compare the two configurations bond by bond. At $b_1$ the product $V'\cdot(u,V_{\mathrm{off}})$
equals $e^{i\zeta T'}u$, which is $u e^{i\zeta\hat{T}}$ by the identity just proved. On every other
bond the product equals $\mathbb{1}$ times the value of $V_{\mathrm{off}}$. These are exactly the
values of $V^{\zeta,\hat{T}}$.

\emph{Clause (b).} On every bond other than $b_1$ we have $V'-\mathbb{1}=0$. At $b_1$ we apply the
exponential bound $|e^{Z}-\mathbb{1}|\le e^{|Z|}-1$, valid for $Z\in M_N(\mathbb{C})$ and stated
in Notation~\ref{4.36:not-sun-conventions}, with $Z=i\zeta T'$. Since $|T'|=1$ by (a), we have
$|Z|=|\zeta|$, and hence $|V'(b_1)-\mathbb{1}|\le e^{|\zeta|}-1$.

\emph{Clause (c).} Since $T'$ is Hermitian, it can be diagonalised by a unitary matrix:
\begin{equation*}
T'=W\operatorname{diag}(t_1,\dots,t_N)W^{*},
\end{equation*}
with $W$ unitary and $t_a$ real. The eigenvalues satisfy $|t_a|\le|T'|=1$ and, by (a),
$\sum_a t_a=\operatorname{Tr}T'=0$. Then
\begin{equation*}
V'(b_1)=W\operatorname{diag}\bigl(e^{w_1},\dots,e^{w_N}\bigr)W^{*},\qquad w_a:=i\zeta t_a.
\end{equation*}
The imaginary parts satisfy
\begin{equation*}
|\operatorname{Im}w_a|=|t_a|\,|\operatorname{Re}\zeta|\le|\operatorname{Re}\zeta|<\pi.
\end{equation*}
So each $e^{w_a}$ has argument in $(-\pi,\pi)$ and does not lie on $(-\infty,0]$, and the principal
logarithm satisfies $\operatorname{Log}e^{w_a}=w_a$. The principal logarithm of the diagonalisable
matrix $V'(b_1)$ is obtained by applying $\operatorname{Log}$ to its eigenvalues in the same
eigenbasis. This gives
\begin{equation*}
\log V'(b_1)=W\operatorname{diag}(w_1,\dots,w_N)W^{*}=i\zeta T',
\end{equation*}
that is, $B(\zeta)(b_1)=\zeta T'$.

Consequently $\operatorname{tr}B(\zeta)(b_1)=\zeta\operatorname{tr}T'=0$, and $B(\zeta)$ vanishes on
the other bonds, so $B(\zeta)$ is traceless. Differentiating $B(\zeta)(b_1)=\zeta T'$ in $\zeta$
gives $\partial_\zeta B(\zeta)(b_1)=T'$.

As a traceless complex matrix, $\zeta T'$ lies in $\mathfrak{sl}(N,\mathbb{C})=\mathfrak{g}^{c}$.
For real $\zeta$ it is moreover Hermitian, so it lies in $\mathfrak{g}$. Let $|\zeta|<2r_0$. Then
$\mathcal{H}_\zeta$ is the response field $\mathcal{H}(\zeta T')$ of this one-bond perturbation, and
the one-bond response statement assigns to a $\mathfrak{g}^{c}$-valued perturbation a
$\mathfrak{g}^{c}$-valued response field and to a $\mathfrak{g}$-valued perturbation a
$\mathfrak{g}$-valued one. Hence $\mathcal{H}_\zeta$ is $\mathfrak{g}^{c}$-valued, and
$\mathfrak{g}$-valued at real $\zeta$.

\emph{Clause (d).} Let $|\zeta|<2r_0$. We verify the links of the chain
\eqref{4.36:eq-probing-direction-a} from left to right.
\begin{itemize}
\item The first inequality is clause (b).
\item By the first ceiling in \eqref{4.36:eq-probing-direction-1}, $|\zeta|<2r_0\le1$, so
$s:=|\zeta|$ lies in $[0,1]$. On $[0,1]$ the function $s\mapsto e^{s}$ is convex, so it lies below
the chord joining $(0,1)$ and $(1,e)$ (see the elementary inequalities
\eqref{4.36:eq-sun-conventions-a}). This gives $e^{s}-1\le(e-1)s$, which is the second inequality.
\item The third inequality is strict: multiply $|\zeta|<2r_0$ by $e-1>0$.
\item The fourth inequality is the second ceiling in \eqref{4.36:eq-probing-direction-1},
multiplied by $2(e-1)$: indeed $2r_0\le\tfrac12 C_1^{(102)}a_1'$.
\item The last inequality is strict. Since $e-1<2$ (again among the elementary inequalities
\eqref{4.36:eq-sun-conventions-a}), we have $\tfrac12(e-1)<1$; as $C_1^{(102)}a_1'>0$, this gives
$\tfrac12(e-1)C_1^{(102)}a_1'<C_1^{(102)}a_1'$.
\end{itemize}
The chain ends in $|V'-\mathbb{1}|<C_1^{(102)}a_1'$ for every $\zeta$ with $|\zeta|<2r_0$. This is
the smallness condition of \cite[(172)]{B-CMP102b}.

\emph{Clause (e).} That the direction enters only through the three quantities named in (e) is the
form in which the one-bond response statement and the estimates derived from it are stated; it
remains to check that none of the three depends on $\hat{T}$ or on $N$. By (b) and (d), the bound on
$|V'-\mathbb{1}|$ is $e^{|\zeta|}-1$, which
contains neither a norm of $\hat{T}$ nor a power of $\mathsf{n}$. By (c), $\partial_\zeta B=T'$
with $|T'|=1$. Also by (c), tracelessness holds without any factor.

In the bound on $\partial_\zeta\mathcal{H}_\zeta$ furnished by the one-bond response statement, the
velocity on a block $y$ is $C_H e_y$ multiplied by the norm of $\partial_\zeta B(b_1)=T'$. Since
$|T'|=1$ by (a), this is \eqref{4.36:eq-probing-direction-2}.
\end{proof}

\begin{lemma}[The auxiliary thresholds as floors on the inverse coupling]\label{4.36:lem-threshold-floors}
We use the notation of Notation~\ref{4.36:not-sun-conventions} and
Notation~\ref{4.36:not-history-geometry}. In particular, for every scale $s\le K$ we have
$x_s=g_s^{-2}>1$, and
\begin{align*}
\varepsilon_s&=g_s\,p_0(g_s)=A_0g_s(\log x_s)^{p_0}, &
\delta_s&=A_1g_s(\log x_s)^{p_0},\\
\alpha_{0,s}&=C_0g_s(\log x_s)^{q_0}, &
\alpha_{1,s}&=C_1g_s(\log x_s)^{q_1}.
\end{align*}
Here $R_s$ is the least power of $L$ not smaller than $(\log x_s)^r$. We put
\begin{align*}
c_{\mathrm{reg}}&:=2B_3L^2(1+\beta_0), & \hat\varepsilon_n&:=c_{\mathrm{reg}}LR_n\varepsilon_n, &
a_m&:=(1+\beta_0)\hat\varepsilon_m,\\
c^{\mathrm{re}}_m&:=C_{\mathrm{re}}(L,M)\,a_m, &
\alpha_{3,j}&:=\frac{\beta_s\min\{\alpha_{0,j},\alpha_{1,j}\}}{2L^2B_3}, &
a^{\flat}&:=\min\Bigl\{\frac{a_1}{2L^2B_3},\frac{a_0}{2L^2B_3^2}\Bigr\}.
\end{align*}
The following are positive constants. $C_{\mathrm{re}}(L,M)$, $c_B$, $C_Q$, $C(M)$ and the number
$\mathsf{o}'$ are the constants of the chart on the window cube $\square_0$ and of its thresholds:
$c_B\delta_m$ is the radius of the ball of chart data, $c^{\mathrm{re}}_m$ bounds the chart datum,
and $\mathsf{o}'$ is the fixed number written $O(1)$ in those thresholds. $B_3$, $\beta_s$, $a_0$ and
$a_1$ are the constants of the spaces of small-field configurations that enter those thresholds
and the ceiling $a^{\flat}$; $C_{\mathrm{loc}}$ is the constant of
Notation~\ref{4.36:not-history-geometry}, $r_0$ is the radius of
Notation~\ref{4.36:not-probing-direction}, and $\delta_0$ is the decay rate of the weights $e_y$
of hypothesis (3) below. Moreover $0<\beta_0\le1$ is the constant of the transport shapes (TS1),
(TS2) of \eqref{4.36:eq-sun-conventions-b}. For $c>0$ we put
\begin{equation*}
\psi(x):=x^{-1/2}(\log x)^{p_0+r},\qquad
x(c):=\max\bigl\{e^{2(p_0+r)},\ \bigl(c\,(4(p_0+r))^{p_0+r}\bigr)^4\bigr\}.
\end{equation*}
Assume the following.
\begin{enumerate}
\item[(1)] The exponents satisfy $q:=\min\{q_0,q_1\}\ge p_0+3r+2$, $p_0\ge5r$ and $r\ge2$.
\item[(2)] The coupling ceiling
\[
\gamma\le\gamma_{\bar B}:=\bigl(\max\{\bar B,x_\delta,x_\beta\}+c_5\bigr)^{-1/2}
\]
holds. Under it,
$x_n\ge\max\{\bar B,x_\delta,x_\beta\}$ for every $n\le K$. Moreover $(\mathrm{F}3)^{\sharp}$
and the transport shapes (TS1), (TS2) of \eqref{4.36:eq-sun-conventions-b} hold for
$j\le n\le k\le K$ and every $L\ge2$.
\item[(3)] $\delta_0M\ge16$, and the weights $e_y=e^{-\delta_0\operatorname{dist}_k(y,b_1)/8}$
satisfy the decay estimate valid under this inequality: for $n<k$, $y\in\Gamma_n$ and
$\ell:=k-1-n\ge0$,
\begin{equation*}
e_y^{1/2}\le e^{-\delta_0R_*/16}e^{-\ell}.
\end{equation*}
\item[(4)] $\mathfrak{c}_T=\mathfrak{c}_T(d,L)>0$ is the constant written $O(dL^2)$ in the
regularity statement for the tower of small-field spaces: that statement bounds the size
$\varepsilon^g_1$ entering it by $\varepsilon^g_1\le\mathfrak{c}_T\sup_{n\le K}\hat\varepsilon_n$,
and it requires $\varepsilon^g_1\le a^{\flat}$.
\end{enumerate}
Then the following hold.

\smallskip\noindent(a) \emph{One-scale sizes.} The function $\psi$ is non-increasing on
$[e^{2(p_0+r)},\infty)$ and tends to $0$ at infinity. For every $s\le K$ with $x_s\ge e$,
every $c>0$ and every $x\ge x(c)$,
\begin{equation}\label{4.36:eq-threshold-floors-a}
\begin{aligned}
&(\mathrm{a}_1)\quad \hat\varepsilon_s<c_{\mathrm{reg}}L^2A_0\,g_s(\log x_s)^{p_0+r},\qquad
(\mathrm{a}_2)\quad \delta_s\le A_1g_s(\log x_s)^{p_0+r},\\
&(\mathrm{a}_3)\quad \min\{\alpha_{0,s},\alpha_{1,s}\}\ge\min\{C_0,C_1\}\,g_s(\log x_s)^{q},\qquad
(\mathrm{a}_4)\quad c\,\psi(x)\le c\,\psi(x(c))\le1.
\end{aligned}
\end{equation}
In particular, $c\,g_s(\log x_s)^{p_0+r}\le1$ whenever $x_s\ge x(c)$.

\smallskip\noindent(b) \emph{The chart thresholds.} Put
\begin{align*}
\mathfrak{k}_1&:=\frac{8(1+\beta_0)^3B_3L^4A_0}{C_0}, &
\mathfrak{k}_2&:=\frac{2B_3^2\,\mathsf{o}'\,M(1+\beta_0)^2c_{\mathrm{reg}}L^2A_0}{C_1},\\
\mathfrak{k}_3&:=\frac{(4B_3\,\mathsf{o}'\,M)^2(1+\beta_0)c_{\mathrm{reg}}L^2A_0}{\beta_s}, &
\mathfrak{k}_5&:=C(M)C_{\mathrm{re}}(1+\beta_0)c_{\mathrm{reg}}L^2A_0,\\
\mathfrak{k}_{\mathrm{im}}&:=\frac{8L^2B_3C_Q(1+\beta_0)\bigl[C_{\mathrm{re}}(1+\beta_0)
c_{\mathrm{reg}}L^2A_0+c_BA_1\bigr]}{\beta_s\min\{C_0,C_1\}}.
\end{align*}
Define the floors
\begin{align*}
x^{(1)}&:=\max\{e,\exp(\mathfrak{k}_1^{1/(q_0-p_0-r)})\}, &
x^{(2)}&:=\max\{e,\exp(\mathfrak{k}_2^{1/(q_1-p_0-r)})\},\\
x^{(4)}&:=\max\{e,\exp(\mathfrak{k}_{\mathrm{im}}^{1/(q-p_0-r)})\}, &
x^{(3)}&:=x(\mathfrak{k}_3),\qquad x^{(5)}:=x(\mathfrak{k}_5).
\end{align*}
Then for $0\le j\le m\le K$ and every integer $n'$ with $m\le n'\le K$,
\begin{equation}\label{4.36:eq-threshold-floors-b}
\begin{aligned}
&(\mathrm{b}_0)\quad \hat\varepsilon_{n'}L^{-2(n'-m)}\le a_m,\\
&(\mathrm{b}_1)\quad x_j\ge x^{(1)}\ \Longrightarrow\ a_mL^{-2(m-j)}\le(1+\beta_0)^2
\hat\varepsilon_j\le\tfrac14\alpha_{0,j},\\
&(\mathrm{b}_2)\quad a_mL^{-(m-j)}\le(1+\beta_0)^2\hat\varepsilon_j,\quad\text{and}\\
&\phantom{(\mathrm{b}_2)\quad}x_j\ge x^{(2)}\ \Longrightarrow\
B_3^2\,\mathsf{o}'\,M(1+\beta_0)^2\hat\varepsilon_j<\tfrac12\alpha_{1,j},\\
&(\mathrm{b}_3)\quad x_m\ge x^{(3)}\ \Longrightarrow\ (4B_3\,\mathsf{o}'\,M)^2a_m\le\beta_s,\\
&(\mathrm{b}_4)\quad x_j\ge x^{(4)}\ \Longrightarrow\
2C_Q\bigl(c^{\mathrm{re}}_m+c_B\delta_m\bigr)L^{j-m}<\tfrac12\alpha_{3,j},\\
&(\mathrm{b}_5)\quad x_m\ge x^{(5)}\ \Longrightarrow\ C(M)\,c^{\mathrm{re}}_m\le1.
\end{aligned}
\end{equation}

\smallskip\noindent(c) \emph{Weights against the radii $\delta$ across levels.}
Put $\mathfrak{k}_6:=4\log(2\sqrt2)/(5\delta_0L^2M)$ and
$x^{(6)}:=\max\{e,\exp(\mathfrak{k}_6^{1/r})\}$. Then
\begin{equation}\label{4.36:eq-threshold-floors-c}
x_k\ge x^{(6)}\ \Longrightarrow\ \frac{e_y^{1/2}}{\delta_n}\le\frac1{\delta_k}
\qquad\text{for all } n<k \text{ and } y\in\Gamma_n.
\end{equation}

\smallskip\noindent(d) \emph{Transport of $\delta$ and comparison with $\alpha_1$.} Put
$x^{(7)}:=\max\{e,\exp((8A_1/C_1)^{1/(q_1-p_0)})\}$. Then for $j\le n'\le K$,
\begin{equation}\label{4.36:eq-threshold-floors-d}
\delta_{n'}L^{-(n'-j)}\le(1+\beta_0)\delta_j\le2\delta_j,\qquad
x_j\ge x^{(7)}\ \Longrightarrow\ 2\delta_j\le\tfrac14\alpha_{1,j}.
\end{equation}

\smallskip\noindent(e) \emph{Further floors.} Put $\mathfrak{k}_{\mathrm{tow}}:=
\mathfrak{c}_Tc_{\mathrm{reg}}L^2A_0/a^{\flat}$ and define
\begin{equation*}
x^{(9)}:=\max\{e,\exp((8C_{\mathrm{loc}}/A_0)^{1/p_0})\},\quad x^{(10)}:=4r_0^{-2},\quad
x^{(11)}:=\exp(L^{2/r}),\quad x^{(\mathrm{tow})}:=x(\mathfrak{k}_{\mathrm{tow}}).
\end{equation*}
Then
\begin{equation}\label{4.36:eq-threshold-floors-e}
\begin{aligned}
&(\mathrm{e}_1)\quad x_k\ge x^{(9)}\ \Longrightarrow\ p_0(g_k)\ge8C_{\mathrm{loc}},\qquad
(\mathrm{e}_2)\quad x_k\ge x^{(10)}\ \Longrightarrow\ g_k\le r_0/2<r_0,\\
&(\mathrm{e}_3)\quad x_j\ge x^{(11)}\ \Longrightarrow\ (\log x_j)^r\ge L^2\
\Longrightarrow\ R_j\ge L^2,\\
&(\mathrm{e}_4)\quad x_n\ge x^{(\mathrm{tow})}\ \text{for all } n\le K\ \Longrightarrow\
\mathfrak{c}_T\sup_{n\le K}\hat\varepsilon_n\le a^{\flat}.
\end{aligned}
\end{equation}
In $(\mathrm{e}_3)$, also $R_j\ge L^2\ge L$. In $(\mathrm{e}_4)$ the numbers $a^{\flat}$ and
$\mathfrak{k}_{\mathrm{tow}}$ depend only on $a_0$, $a_1$, $B_3$, $\beta_0$, $L$, $A_0$ and
$\mathfrak{c}_T(d,L)$, in particular not on $\gamma$ or $K$, and the condition is one-sided.

\smallskip\noindent(f) \emph{Conversion.} For every number $x_\bullet$ and every scale $s\le K$,
\begin{equation*}
\gamma\le(x_\bullet+c_5)^{-1/2}\ \Longrightarrow\ x_s\ge x_\bullet .
\end{equation*}
\end{lemma}

\begin{proof}
We write $a:=p_0+r$. By hypothesis (1), $a\ge6r>0$. We first record the gaps in the exponents.
By (1), $q-p_0-r\ge2r+2>0$. Since $q_0\ge q$ and $q_1\ge q$, both $q_0-p_0-r$ and
$q_1-p_0-r$ are at least $2r+2>0$, and $q_1-p_0\ge3r+2>0$. So every root taken in the
definitions of the floors is a root of positive order.

\emph{Clause (a).} Let $s\le K$ with $x_s\ge e$, so that $\log x_s\ge1$ and $(\log x_s)^r\ge1$.

We prove $(\mathrm{a}_1)$. Write $R_s=L^t$, where $t$ is the least integer not smaller than
$\log_L(\log x_s)^r$; this integer is $\ge0$. Then $t<\log_L(\log x_s)^r+1$, that is,
$R_s<L(\log x_s)^r$. Hence
\begin{equation*}
\hat\varepsilon_s=c_{\mathrm{reg}}LR_s\varepsilon_s<c_{\mathrm{reg}}L^2(\log x_s)^r\cdot
A_0g_s(\log x_s)^{p_0}=c_{\mathrm{reg}}L^2A_0\,g_s(\log x_s)^{p_0+r}.
\end{equation*}

$(\mathrm{a}_2)$ follows from $\delta_s=A_1g_s(\log x_s)^{p_0}$ and $(\log x_s)^r\ge1$.

For $(\mathrm{a}_3)$, note that $q\le q_0$, $q\le q_1$ and $\log x_s\ge1$ give
$(\log x_s)^{q_i}\ge(\log x_s)^q$ for $i=0,1$. Hence
$\alpha_{0,s}\ge C_0g_s(\log x_s)^q$ and $\alpha_{1,s}\ge C_1g_s(\log x_s)^q$, and both right
sides are at least $\min\{C_0,C_1\}g_s(\log x_s)^q$.

For $(\mathrm{a}_4)$, compute for $x>1$
\begin{equation*}
\frac{d}{dx}\bigl[(\log x)^ax^{-1/2}\bigr]=x^{-3/2}(\log x)^{a-1}\bigl(a-\tfrac12\log x\bigr).
\end{equation*}
This is $\le0$ for $\log x\ge2a$. Hence $\psi$ is non-increasing on $[e^{2a},\infty)$.

Next let $x\ge1$ and $y:=x^{1/(4a)}\ge1$. The inequality $\log y\le y$ gives
$(\log x)/(4a)\le x^{1/(4a)}$. Both sides are non-negative, so raising to the power $a$ gives
$(\log x)^a\le(4a)^ax^{1/4}$. Hence $0\le\psi(x)\le(4a)^ax^{-1/4}$ for $x\ge1$, and
$\psi(x)\to0$ as $x\to\infty$.

Now let $c>0$ and $x\ge x(c)$. Since $x(c)\ge e^{2a}$, monotonicity gives
$c\psi(x)\le c\psi(x(c))$. The bound just proved gives
$c\psi(x(c))\le c(4a)^ax(c)^{-1/4}$. This is $\le1$, because
$x(c)\ge(c(4a)^a)^4$, that is, $x(c)^{1/4}\ge c(4a)^a$.

Finally, $g_s(\log g_s^{-2})^{a}=x_s^{-1/2}(\log x_s)^a=\psi(x_s)$. Thus
$c\,g_s(\log x_s)^{a}\le1$ whenever $x_s\ge x(c)$.

\emph{Clause (b).} Every floor $x^{(i)}$ with $i\in\{1,\dots,5\}$ is at least $e$: directly for
$i=1,2,4$, and because $x(c)\ge e^{2a}\ge e$ for $i=3,5$. Hence
$(\mathrm{a}_1)$--$(\mathrm{a}_3)$ apply at every scale on which one of these floors is
assumed.

For $(\mathrm{b}_0)$, apply (TS2) of hypothesis (2) to $\hat\varepsilon$ with the pair
$(j,n)$ replaced by $(m,n')$. This gives
$\hat\varepsilon_{n'}L^{-2(n'-m)}\le(1+\beta_0)\hat\varepsilon_m=a_m$; no floor enters.

For $(\mathrm{b}_1)$, apply (TS2) to $\hat\varepsilon$ with the pair $(j,m)$:
\begin{equation*}
a_mL^{-2(m-j)}=(1+\beta_0)\hat\varepsilon_mL^{-2(m-j)}\le(1+\beta_0)^2\hat\varepsilon_j .
\end{equation*}
By $(\mathrm{a}_1)$ and $c_{\mathrm{reg}}=2B_3L^2(1+\beta_0)$,
\begin{equation*}
(1+\beta_0)^2\hat\varepsilon_j<2(1+\beta_0)^3B_3L^4A_0\,g_j(\log x_j)^{p_0+r}.
\end{equation*}
The right side is at most $\frac14\alpha_{0,j}=\frac14C_0g_j(\log x_j)^{q_0}$ if and only if
$(\log x_j)^{q_0-p_0-r}\ge\mathfrak{k}_1$. This holds when $x_j\ge x^{(1)}$, since then
$\log x_j\ge\mathfrak{k}_1^{1/(q_0-p_0-r)}$.

For $(\mathrm{b}_2)$, the first inequality is (TS1) for $\hat\varepsilon$ with the pair $(j,m)$,
multiplied by $1+\beta_0$. For the second, $(\mathrm{a}_1)$ gives
\begin{equation*}
B_3^2\,\mathsf{o}'\,M(1+\beta_0)^2\hat\varepsilon_j<B_3^2\,\mathsf{o}'\,M(1+\beta_0)^2
c_{\mathrm{reg}}L^2A_0\,g_j(\log x_j)^{p_0+r}.
\end{equation*}
The right side is at most $\frac12C_1g_j(\log x_j)^{q_1}=\frac12\alpha_{1,j}$ if and only if
$(\log x_j)^{q_1-p_0-r}\ge\mathfrak{k}_2$. This holds for $x_j\ge x^{(2)}$.

For $(\mathrm{b}_3)$, no gap enters. By $(\mathrm{a}_1)$ at the scale $m$,
\begin{equation*}
(4B_3\,\mathsf{o}'\,M)^2a_m=(4B_3\,\mathsf{o}'\,M)^2(1+\beta_0)\hat\varepsilon_m
<\beta_s\,\mathfrak{k}_3\,g_m(\log x_m)^{p_0+r}.
\end{equation*}
By the last assertion of clause (a) with $c=\mathfrak{k}_3$, we have
$\mathfrak{k}_3g_m(\log x_m)^{p_0+r}\le1$ when $x_m\ge x(\mathfrak{k}_3)=x^{(3)}$.

For $(\mathrm{b}_4)$, recall $c^{\mathrm{re}}_m=C_{\mathrm{re}}(1+\beta_0)\hat\varepsilon_m$.
Apply (TS1) to $\hat\varepsilon$ and to $\delta$ with the pair $(j,m)$:
\begin{equation*}
2C_Q\bigl(c^{\mathrm{re}}_m+c_B\delta_m\bigr)L^{j-m}\le2C_Q(1+\beta_0)\bigl[C_{\mathrm{re}}
(1+\beta_0)\hat\varepsilon_j+c_B\delta_j\bigr].
\end{equation*}
By $(\mathrm{a}_1)$ and $(\mathrm{a}_2)$ the right side is
\begin{equation*}
<2C_Q(1+\beta_0)\bigl[C_{\mathrm{re}}(1+\beta_0)c_{\mathrm{reg}}L^2A_0+c_BA_1\bigr]\,
g_j(\log x_j)^{p_0+r}.
\end{equation*}
By $(\mathrm{a}_3)$,
$\frac12\alpha_{3,j}\ge\beta_s\min\{C_0,C_1\}\,g_j(\log x_j)^q/(4L^2B_3)$. Dividing by
$g_j(\log x_j)^{p_0+r}>0$, the last display is at most $\frac12\alpha_{3,j}$ provided
$(\log x_j)^{q-p_0-r}\ge\mathfrak{k}_{\mathrm{im}}$. This holds for $x_j\ge x^{(4)}$.

For $(\mathrm{b}_5)$, $(\mathrm{a}_1)$ at the scale $m$ gives
\begin{equation*}
C(M)c^{\mathrm{re}}_m=C(M)C_{\mathrm{re}}(1+\beta_0)\hat\varepsilon_m<
\mathfrak{k}_5\,g_m(\log x_m)^{p_0+r}.
\end{equation*}
This is $\le1$ for $x_m\ge x(\mathfrak{k}_5)=x^{(5)}$, by clause (a). In particular the factor
$1+C(M)c^{\mathrm{re}}_m$ is at most $2$.

\emph{Clause (c).} Let $n<k$, $y\in\Gamma_n$ and $\ell=k-1-n\ge0$. Hypothesis (3) gives
$e_y^{1/2}\le e^{-\delta_0R_*/16}e^{-\ell}$. By hypothesis (2), $x_k\ge x_\delta$ and
$(\mathrm{F}3)^{\sharp}$ holds; it gives $\delta_n\ge\frac12(1+(k-n))^{-1/2}\delta_k$. Since
$1+(k-n)=2+\ell$, this reads $1/\delta_n\le2(2+\ell)^{1/2}/\delta_k$.

Next, $(2+\ell)^{1/2}e^{-\ell}\le\sqrt2$ by \eqref{4.36:eq-sun-conventions-a}. Hence
\begin{equation*}
\frac{e_y^{1/2}}{\delta_n}\le2\sqrt2\,e^{-\delta_0R_*/16}\,\frac{1}{\delta_k},
\end{equation*}
and this is $\le1/\delta_k$ as soon as $e^{\delta_0R_*/16}\ge2\sqrt2$, that is, as soon as
$\delta_0R_*/16\ge\log(2\sqrt2)$.

Now $R_*=20L^2M\max\{R_k,R_{k+1}\}\ge20L^2MR_k$, and $R_k\ge(\log x_k)^r$ by the definition
of $R_k$. Therefore $\delta_0R_*/16\ge\frac54\delta_0L^2M(\log x_k)^r$. This is
$\ge\log(2\sqrt2)$ if $(\log x_k)^r\ge\mathfrak{k}_6$, which holds for $x_k\ge x^{(6)}$.

\emph{Clause (d).} Let $j\le n'\le K$. The first inequality is (TS1) for $\delta$ with the pair
$(j,n')$, and the second follows from $\beta_0\le1$. Further,
\begin{equation*}
2\delta_j=2A_1g_j(\log x_j)^{p_0},\qquad \tfrac14\alpha_{1,j}=\tfrac14C_1g_j(\log x_j)^{q_1}.
\end{equation*}
Since $g_j>0$ and $\log x_j>0$, we get $2\delta_j\le\frac14\alpha_{1,j}$ if and only if
$(\log x_j)^{q_1-p_0}\ge8A_1/C_1$. This holds for $x_j\ge x^{(7)}$.

\emph{Clause (e).} For $(\mathrm{e}_1)$: if $x_k\ge x^{(9)}$, then
$\log x_k\ge(8C_{\mathrm{loc}}/A_0)^{1/p_0}$, so
$p_0(g_k)=A_0(\log x_k)^{p_0}\ge8C_{\mathrm{loc}}$.

For $(\mathrm{e}_2)$: if $x_k\ge4r_0^{-2}$, then $g_k=x_k^{-1/2}\le r_0/2<r_0$.

For $(\mathrm{e}_3)$: if $x_j\ge\exp(L^{2/r})$, then $\log x_j\ge L^{2/r}$, so
$(\log x_j)^r\ge L^2$. Hence $R_j\ge(\log x_j)^r\ge L^2\ge L$, by the definition of $R_j$.

For $(\mathrm{e}_4)$: note $x^{(\mathrm{tow})}\ge e$. By $(\mathrm{a}_1)$,
$\hat\varepsilon_n<c_{\mathrm{reg}}L^2A_0\psi(x_n)$ for each $n\le K$. Hence
\begin{equation*}
\mathfrak{c}_T\hat\varepsilon_n<\mathfrak{c}_Tc_{\mathrm{reg}}L^2A_0\psi(x_n)=
a^{\flat}\,\mathfrak{k}_{\mathrm{tow}}\psi(x_n)\le a^{\flat},
\end{equation*}
where the last step is $(\mathrm{a}_4)$ with $c=\mathfrak{k}_{\mathrm{tow}}$, that is, the
monotonicity of $\psi$ on $[e^{2a},\infty)$. Taking the supremum over the finitely many
$n\le K$ gives $(\mathrm{e}_4)$. The constants $a^{\flat}$ and $\mathfrak{k}_{\mathrm{tow}}$
are built from $a_0$, $a_1$, $B_3$, $\mathfrak{c}_T(d,L)$, $c_{\mathrm{reg}}$, $L$ and $A_0$
alone.

\emph{Clause (f).} By (Fl1) of \eqref{4.36:eq-sun-conventions-b}, $x_s\ge\gamma^{-2}-c_5$ for
every $s\le K$. If $\gamma\le(x_\bullet+c_5)^{-1/2}$, then $\gamma^{-2}\ge x_\bullet+c_5$, and
hence $x_s\ge x_\bullet$.
\end{proof}

\begin{definition}[The near part of the effective action and the standing hypotheses]
\label{4.36:def-near-part}
The plaquette $p$, the marked bond $b_1\subset\partial p$, the free link $u_1=V_k(b_1)$, the
cleanness of a history, the regions $\Lambda_j$, $\Lambda_j^{0}$, $Z_j=\Lambda_j^{c}$, the polymer
classes $\mathbf{D}_j$, $\mathbf{D}_j^{R}$, the tents $h_z$ and the arrests $W$ of a history are
those of Notation~\ref{4.36:not-history-geometry}, as are the $u_1$-window
$W(V_{\mathrm{off}})$, the deep set $\mathcal{D}$ and the indicator $\mathbb{1}_{G(h)}$;
$V=(u_1,V_{\mathrm{off}})$ denotes the level-$k$ configuration split at the marked bond into the
free link $u_1$ and the off-link configuration $V_{\mathrm{off}}$.
The group, the norm $|\cdot|$ and the algebra
$\mathfrak{g}$ are those of Notation~\ref{4.36:not-sun-conventions}. The probing configuration
$V^{t,\hat{T}}$ is that of Notation~\ref{4.36:not-probing-direction}. For a history $h$, $m_h$
denotes its fluctuation measure, and $m_h^{V_{\mathrm{off}}}$ the fluctuation measure whose
$u_1$-free characteristic functions are evaluated at the off-link configuration $V_{\mathrm{off}}$.
The operator-norm window domain $D^{\mathrm{op}}$ is that of the window construction used in this
section. For $r>0$, $B^{\mathrm{op}}_{r}(u^0)$ is the ball of
radius $r$ about $u^0$ in the operator norm $|\cdot|$.

\smallskip
\emph{The near and far parts.} Let $h$ be a clean history, let $F_h$ be the exponent of the term
of the effective density at level $k$ indexed by $h$, and let $A_k(h)$ be the effective action of
$h$ at level $k$. The function $F_h$ is split as
\begin{equation}\label{4.36:eq-near-part-1}
F_h=F^{\mathrm{near}}+F^{\mathrm{far}},\qquad
F^{\mathrm{near}}:=A_k(h)=F^{W}+F^{E}+F^{R}+F^{B}+\mathrm{const},
\end{equation}
where the constant does not depend on $u_1$. The normalisers $\log z_j$ and $z^{(k)}$, which carry
the dimension $D=N^2-1$, enter this constant and the fluctuation measure $m_h$. Since they do not
depend on $u_1$, they are annihilated by the derivative $\partial_t$ along the probing configuration
$V^{t,\hat{T}}$ of Notation~\ref{4.36:not-probing-direction}. The four members of
$F^{\mathrm{near}}$ are the following.
\begin{enumerate}
\item[(W)] The plaquette part is
\begin{equation}\label{4.36:eq-near-part-2}
F^{W}:=-\sum_{q}c(q)\,w_q,\qquad w_q:=1-\operatorname{Re}\operatorname{tr}U_k(\partial q).
\end{equation}
Here $c(q)$ is the coefficient of the plaquette $q$ in the effective action $A_k(h)$ at level $k$.
At complex configurations $w_q$ is continued by the function $\mathfrak{w}$ of the continuation
lemma for the plaquette action.
\item[(E)] The regular part is
\begin{equation}\label{4.36:eq-near-part-3}
F^{E}:=\sum_{j\le k}\ \sum_{z\in\Lambda_j^{0}}\mathfrak{e}_{j,z}\bigl(U_k(V)\bigr),
\end{equation}
with
\begin{equation}\label{4.36:eq-near-part-4}
\mathfrak{e}_{j,z}(U):=\sum_{\substack{X\in\mathbf{D}_j\\ z\in X\subset\Lambda_j}}
\Bigl[\mathbf{E}^{(j)}(X,U,z)-\mathbf{E}^{(j)}(X,\mathbb{1},z)\Bigr]-\beta_j\,A(h_z,U).
\end{equation}
Here $A(h_z,\cdot)$ is the action weighted by the tent $h_z$. The term $A(h_z,\mathbb{1})$ is not
subtracted in \eqref{4.36:eq-near-part-4} because it vanishes: $1-\operatorname{Re}\operatorname{tr}\mathbb{1}=0$,
since $\operatorname{tr}\mathbb{1}=1$.
\item[(R)] The renormalisation part is
\begin{equation}\label{4.36:eq-near-part-5}
F^{R}:=\sum_{j\le k}\ \sum_{X\in\mathbf{D}_j^{R}}
\Bigl[\mathbf{R}^{(j)}\bigl(X,U_k(V)\bigr)-\mathbf{R}^{(j)}(X,\mathbb{1})\Bigr].
\end{equation}
\item[(B)] The boundary part $F^{B}$ is the sum of all terms of boundary type. These are of two
kinds.
\begin{itemize}
\item The first kind consists of the boundary terms $\mathbf{B}^{(i)}(X)$ of the effective action
of \cite[(2.23)]{B-CMP119} with $X\in\mathbf{D}_i$ and
$X\cap Z_i\neq\emptyset$. There is one function for each pair $(i,X)$. Each obeys the bound of
\cite[(2.42)]{B-CMP119}:
\begin{equation}\label{4.36:eq-near-part-6}
\bigl|\mathbf{B}^{(i)}(X)\bigr|<B_0\,e^{-\kappa d_i(X)}
\quad\text{on }\widetilde{U}_i(X),
\end{equation}
with the domain $\widetilde{U}_i(X)$ and the constants $B_0$, $\kappa$ of
\cite[(2.42)]{B-CMP119}, and this bound holds uniformly in the fluctuation variables
$w\in\operatorname{supp}m_h$. The uniformity in $w$ is stated in \cite{B-CMP119} beside that
display; here it is part of the hypothesis on the boundary terms at $SU(N)$ under which the step
is taken.
\item The second kind consists of the terms $\mathbf{C}_{W,j}(X)$ frozen on the arrests $W$
\cite[(1.63), (1.70)]{B-CMP122a}. The constant attached to them there is written here
$C_0^{\mathrm{fr}}$.
\end{itemize}
\end{enumerate}
The far part is
\begin{equation*}
F^{\mathrm{far}}:=Q_h+\sum_{Y}\Psi_Y ,
\end{equation*}
where $Q_h$ and the activities $\Psi_Y$ are those of the far-part bound, in which the far part
is defined and estimated; it plays no role in this section. The members $F^{E}$ and $F^{R}$ are
functions of $U_k(V)$ alone; in particular they do not depend on the fluctuation variables $w$.

\smallskip
\emph{Constants of the second-derivative bound.} The constants
\begin{equation}\label{4.36:eq-near-part-7}
\begin{gathered}
C_q(N):=2C_P\Bigl(\frac{B_3a_1'}{r_0}+C_P\Bigr),\qquad
C_{W'}(N):=36\,C_q(N)\,C_\sigma\,S_{\mathfrak{n}}
\end{gathered}
\end{equation}
are those of the second-derivative bound on the plaquette part, which states, at every $N$,
\begin{equation*}
\bigl|\partial_t^{2}F^{W}\bigr|_0\le C_{W'}(N)\,x_k .
\end{equation*}
Here $r_0$ and $a_1'$ are as in Notation~\ref{4.36:not-probing-direction}, $B_3$ is as in
Lemma~\ref{4.36:lem-threshold-floors}, $C_\sigma$ is the constant of the geometric sum over the
blocks of $\Gamma_n$ in Notation~\ref{4.36:not-history-geometry}, and $C_P$, $S_{\mathfrak{n}}$
are the constants of the second-derivative bound on the plaquette part. The constant
\begin{equation*}
C^{\mathrm{op}}_{F'}(N):=C_{W'}(N)+3
\end{equation*}
is that of the second-derivative bound on the near part along the probing direction, which states
\begin{equation}\label{4.36:eq-near-part-8}
\bigl|\partial_t^{2}F^{\mathrm{near}}(V^{t,\hat{T}},w)\bigr|_0\le C^{\mathrm{op}}_{F'}(N)\,x_k ,
\end{equation}
with $|\cdot|_0$ as in that bound.

\smallskip
\emph{The standing hypotheses.} In what follows in this section the following are fixed:
\begin{itemize}
\item the number of steps $K$ and the level $k$;
\item a clean history $h$ and the off-link configuration $V_{\mathrm{off}}$;
\item fluctuation variables $w\in\operatorname{supp}m_h^{V_{\mathrm{off}}}$;
\item a link $u^0\in W(V_{\mathrm{off}})$ with $(u^0,V_{\mathrm{off}})\in\mathcal{D}$
satisfying the localised condition that accompanies the window domain $D^{\mathrm{op}}$;
\item a link $u\in B^{\mathrm{op}}_{g_k}(u^0)$;
\item a direction $\hat{T}\in\mathfrak{g}$ with $|\hat{T}|=1$.
\end{itemize}
This localised condition holds automatically at the configurations where
$\mathbb{1}_{G(h)}>0$, by a defining condition of the window domain $D^{\mathrm{op}}$. These are
the configurations at which the second-derivative bound on the near part,
\eqref{4.36:eq-near-part-8}, is used in the estimate at $SU(N)$ of the term of the effective
density indexed by $h$.

We take the free link to be $u_1=u$, so that $V=(u,V_{\mathrm{off}})$. This configuration is real
and windowed, that is
$u\in W(V_{\mathrm{off}})$, by the clause of the window domain $D^{\mathrm{op}}$ asserting
that for $u$ in the ball $B^{\mathrm{op}}_{g_k}(u^0)$ about a link $u^0$ as above the configuration
$(u,V_{\mathrm{off}})$ is real and windowed.

The inverse couplings are assumed to satisfy
\begin{equation}\label{4.36:eq-near-part-9}
x_n\ge x_{\mathrm{th}}^{(3)}\qquad\text{for all }n\le K,
\end{equation}
where $x_{\mathrm{th}}^{(3)}$ is the floor defined below through the coupling ceiling
\eqref{4.36:eq-coupling-ceiling-a}. Condition \eqref{4.36:eq-near-part-9} contains the following:
\begin{itemize}
\item the conditions (t1), (t8)$^{+}$ and (t9) of the list of conditions on the auxiliary
parameters and the radii that accompanies the window construction; (t1) and (t8)$^{+}$ are among
the conditions under which the clause of the window domain $D^{\mathrm{op}}$ quoted above holds;
\item the floors \eqref{4.36:eq-threshold-floors-e} of Lemma~\ref{4.36:lem-threshold-floors};
\item the coupling ceilings, equivalently floors on the inverse coupling, that are required by the
charted representation of the pieces, by the membership criterion, by the localisation bound for
the response field, and by the tower.
\end{itemize}
Finally, the conditions (t2), (t4)$'$, (t5) and (t8) of the same list are assumed. The third member
of (t2) is, together with (t1) and (t8)$^{+}$, a condition of the clause of the window domain
$D^{\mathrm{op}}$ quoted above; (t5) is the condition in which the constant $c_d$ of the polymer
count of Notation~\ref{4.36:not-history-geometry} enters; and two members of (t8), the bound
$2r_0\le 1$ and an upper bound on $r_0$ by a fixed multiple of $a_1'$, are used in
Notation~\ref{4.36:not-probing-direction}. The geometric floor $L\ge 7$ of
Notation~\ref{4.36:not-history-geometry} is assumed as well.
\end{definition}

\begin{remark}\label{4.36:thm-regular-bound}
The statement of this theorem is not written out here; parts of its proof are printed below.
\end{remark}

\begin{lemma}[The regular-term bound absorbed into a fraction of the inverse coupling]
\label{4.36:prf-regular-absorption}
This lemma carries out the estimate of the tail terms and the collection step in the proof of
Theorem~\ref{4.36:thm-regular-bound}, together with its clause~(c). The notation is that of
Notation~\ref{4.36:not-sun-conventions}, Notation~\ref{4.36:not-history-geometry} and
Notation~\ref{4.36:not-probing-direction}, and $|\cdot|_0$ is as in Theorem~\ref{4.36:thm-regular-bound}.
Let $h$ be a clean history, $\hat T\in\mathfrak{g}$ with
$|\hat T|=1$, and let $k\le K$ be the level. Assume the standing hypotheses of
Definition~\ref{4.36:def-near-part}, and among them, explicitly, the following.
\begin{itemize}
\item The transfer of Ba{\l}aban's regular-term bound to $SU(N)$: \cite[Theorem~1]{B-CMP109} is stated for
an arbitrary semisimple compact $G\subset U(N)$, and so applies at $G=SU(N)$.
\item Ba{\l}aban's bound on the regular terms \cite[(1.18)]{B-CMP109}: for $X\in\mathbf{D}_j$, the class of
polymers at level $j$, and a site $z$,
$|\mathbf{E}^{(j)}(X,\cdot,z)|\le E_0e^{-\kappa d_j(X)}$ on the domain
$\mathcal{U}^{c}_{j}(X,\alpha_{0,j},\alpha_{1,j})$, which is defined by smallness conditions in the norm
$|\cdot|$ in \cite[(1.11)--(1.16)]{B-CMP109}.
\item Each regular term $\mathbf{E}^{(j)}(X,\cdot,z)$ is gauge invariant and depends on the configuration
only through the bonds of $X$.
\item The membership statement at $SU(N)$ for the ball of the norm $\|\cdot\|_{\star}$ with radius
constant $c_{\star}$, in the form, contained in Ba{\l}aban's hypotheses at $SU(N)$, in which no factor of
$|\hat T|$ enters: if $\mathsf{N}$ is the set of bonds of $X$, $\chi_{\mathsf{N}}$ is the cutoff function
attached to $\mathsf{N}$ in the bound of the next item, and $\|\chi_{\mathsf{N}}\mathcal{H}_\zeta\|_{\star}
\le c_{\star}/2$, then the configurations $e^{i\eta\chi_{\mathsf{N}}\mathcal{H}_\zeta}U_k(V)$ of that
statement, with the parameter $\eta$ of that statement, restricted to $X$, lie in
$\mathcal{U}^{c}_{j}(X,\alpha_{0,j},\alpha_{1,j})$.
\item The bound $\|\chi_{\mathsf{N}}\mathcal{H}_\zeta\|_{\star}\le 2RC_{\chi}'C_H/\delta_k$ for
$|\zeta|\le 2R$, with its constants $C_{\chi}'$ and $C_H$.
\item The second-derivative estimate on the ball: if $\Psi$ is gauge invariant, depends only on the bonds of
$\mathsf{N}$, and has oscillation at most $S$ on the ball for $|\zeta|\le 2R$, then
$|\partial_t^2\Psi(U_k(V^{t,\hat T}))|_0\le S\,\hat e_{\mathsf{N}^{+}}/R^2$. In this estimate the real variable
$t$, $|t|\le 2R$, is identified with the real points $\zeta=t$ of the ball by a clause of the standing
hypotheses of Definition~\ref{4.36:def-near-part}, as in \cite{B-CMP102b}.
\item The geometric statements on polymers: a connected $X\in\mathbf{D}_j$ with $\operatorname{diam}_\infty X\ge
2M\ell_m$ has $d_j(X)\ge 2(L^{m-j}-1)$. For every $n'$, a polymer $X\in\mathbf{D}_j$ that meets
$\Omega_{n'-2}$ and $\Omega^{c}_{n'-3}$ has $L^{n'-3-j}\le d_j(X)+2$. Here $\ell_m$ is the lattice spacing of
scale $m$.
\item The polymer count. Each $X\in\mathbf{D}_j$ is reached from at most $c_{\#}(1+d_j(X))M^4$ sites $z$. For a
block $y\in\Gamma_{n'}$, every $X\in\mathbf{D}_j$ with $X\subset\Lambda_j$ and $y\in\mathfrak{b}(X^{+})$ has
$j\le n'+2$ and $n_X\le n'+2$, and
\begin{equation*}
\#\{X\in\mathbf{D}_j:\ y\in\mathfrak{b}(X^{+}),\ d_j(X)=\ell\}\le 7^4L^{4(n'+2-j)}e^{c_d\ell}.
\end{equation*}
Here $n_X$ is the least $n$ for which $X$ meets $\Gamma_n$, and $c_d$ is the growth constant of the count.
\item The decay of the block-weight sums: $\sigma_{n'}^{(1/16)}\le C_\sigma e^{-(k-n')}$.
\item The threshold conditions on the probing radius $R>0$. $R$ is fixed as a minimum; its first member
gives $2RC_{\chi}'C_H/\delta_k\le c_{\star}/2$, and for each member of the minimum that defines $R$
\begin{equation*}
R^{-2}\le (4C_{\chi}'C_H/c_{\star})^2\delta_k^{-2}+(16C_{\mathrm{loc}})^2\varepsilon_k^{-2}+4r_0^{-2}.
\end{equation*}
\item The threshold conditions on the decay rate: $\kappa\ge 10$ and $\kappa/2-c_d\ge 3c_d+4$.
\item The corresponding bound on the local parts, established in the part of the proof of
Theorem~\ref{4.36:thm-regular-bound} that treats them, and Lemma~\ref{4.36:lem-threshold-floors}.
\end{itemize}
Let $\mathfrak{E}^{(2)}$ and $C_E^m$ be as in Theorem~\ref{4.36:thm-regular-bound}, with $C_E^t$ of \eqref{4.36:eq-regular-absorption-b} below. Then
$|\partial_t^2F^E(V^{t,\hat T})|_0\le\mathfrak{E}^{(2)}$. Moreover
\begin{equation}\label{4.36:eq-regular-absorption-c}
\begin{gathered}
\mathfrak{E}^{(2)}\le\mathcal{C}_1x_k(\log x_k)^{-2p_0}+\mathcal{C}_2,\qquad
\mathcal{C}_2:=(C_E^m+4C_E^t)r_0^{-2},\\
\mathcal{C}_1:=\frac{C_E^m}{A_1^2}+C_E^t\Big[\frac{(4C_{\chi}'C_H/c_{\star})^2}{A_1^2}
+\frac{(16C_{\mathrm{loc}})^2}{A_0^2}\Big],\\
x_E:=\max\big\{\exp\big[(16\mathcal{C}_1)^{1/(2p_0)}\big],\,16\mathcal{C}_2\big\},\\
x_k\ge x_E\ \Longrightarrow\ \mathfrak{E}^{(2)}\le\frac{x_k}{16}+\frac{x_k}{16}=\frac{x_k}{8}.
\end{gathered}
\end{equation}
Here $x_E$ depends only on $(N,L,M,\mathfrak{p})$. The floor $x_k\ge x_E$ is imposed for every
$k\le K$ by the condition $\gamma\le(x_E+c_5)^{-1/2}$.
\end{lemma}

\begin{proof}
For given $(j,n)$, $m=m(j,n)$ is the
scale attached to the pair in the estimate of the local parts in that theorem; recall that $u:=m-j\ge0$
and $n-j\le u+2$. For a site $z\in\Gamma_n$, $\square$ is the cube of $\pi_m$ that contains $z$, and
$\widetilde{\square}^{2}(z)$ is its twofold enlargement.

\smallskip\noindent\textit{The tail terms, one at a time.}
Let $X\in\mathbf{D}_j$ and $z\in X\subset\Lambda_j$, let $n$ be such that $z\in\Gamma_n$, so that
$m=m(j,n)$ is fixed, and assume $X\not\subset\widetilde{\square}^{2}(z)$. The polymer
$X$ is connected. It contains $z\in\square$ and a point outside $\widetilde{\square}^{2}(z)$, so
$\operatorname{diam}_\infty X\ge 2M\ell_m$. The diameter bound of the geometric statements then gives
$d_j(X)\ge 2(L^{m-j}-1)$.

Let $\mathsf{N}$ be the set of bonds of $X$ and let $|\zeta|\le 2R$. By the bound on
$\chi_{\mathsf{N}}\mathcal{H}_\zeta$ and the first member of the minimum defining $R$,
\begin{equation*}
\|\chi_{\mathsf{N}}\mathcal{H}_\zeta\|_{\star}\le 2RC_{\chi}'C_H/\delta_k\le c_{\star}/2 .
\end{equation*}
The membership statement, in the form without a factor of $|\hat T|$, then places the configurations
$e^{i\eta\chi_{\mathsf{N}}\mathcal{H}_\zeta}U_k(V)$, restricted to $X$, in
$\mathcal{U}^{c}_{j}(X,\alpha_{0,j},\alpha_{1,j})$. On that domain \cite[(1.18)]{B-CMP109} holds.
With $u=m-j$ and $n-j\le u+2$, this gives, on the ball of the membership statement for $|\zeta|\le 2R$,
\begin{equation}\label{4.36:eq-regular-absorption-a}
\begin{split}
|\mathbf{E}^{(j)}(X,\cdot,z)|
&\le E_0e^{-\kappa(L^{m-j}-1)}e^{-\kappa d_j(X)/2}\\
&\le e^2L^8E_0L^{-4(n-j)}e^{-(n-j)}e^{-\kappa d_j(X)/2},\\
\bigl|\partial_t^2\bigl[\mathbf{E}^{(j)}(X,U_k(V^{t,\hat T}),z)\bigr]\bigr|_0
&\le \frac{S_X\,\hat e_{X^{+}}}{R^2},\\
S_X&:=2e^2L^8E_0L^{-4(n-j)}e^{-(n-j)}e^{-\kappa d_j(X)/2}.
\end{split}
\end{equation}

The first inequality splits $e^{-\kappa d_j}=e^{-\kappa d_j/2}e^{-\kappa d_j/2}$ and uses
$d_j(X)/2\ge L^{m-j}-1$.

For the second inequality, $n-j\le u+2$ gives
\begin{equation*}
L^{-4(n-j)}e^{-(n-j)}\ge L^{-8}e^{-2}L^{-4u}e^{-u}.
\end{equation*}
Hence it suffices that $e^{-\kappa(L^u-1)}\le L^{-4u}e^{-u}$, that is,
$\kappa(L^u-1)\ge u(4\log L+1)$. At $u=0$ both sides vanish. Let $u\ge1$. Then
\begin{equation*}
L^u-1=(L-1)(L^{u-1}+\dots+1)\ge u(L-1).
\end{equation*}
By the conventions \eqref{4.36:eq-sun-conventions-a}, $\log L\le L-1$ and $1\le L-1$, so
\begin{equation*}
4\log L+1\le 4(L-1)+(L-1)=5(L-1).
\end{equation*}
Therefore any $\kappa\ge5$ suffices, and a fortiori the threshold condition $\kappa\ge10$ does.

For the third line, apply the second-derivative estimate on the ball to
$\Psi:=\mathbf{E}^{(j)}(X,\cdot,z)$. By the hypothesis on the regular terms, this function is gauge invariant
and depends only on the bonds of
$\mathsf{N}$. By the bound just proved, its oscillation on the ball for $|\zeta|\le2R$ is at most
$2\sup|\Psi|\le S_X$. The real variable $t$, $|t|\le 2R$, is identified with the real points of $\zeta$ by the
clause of the standing hypotheses of Definition~\ref{4.36:def-near-part} that identifies real $t$,
$|t|\le 2R$, with the real points of the ball, as in \cite{B-CMP102b}. The estimate then gives the third
line, with $\mathsf{N}^{+}=X^{+}$.

\smallskip\noindent\textit{The tail sum.}
Put
\begin{equation*}
T^E:=\sum_{j}\sum_{z\in\Lambda_j^{0}}
\sum_{\substack{X\ni z,\ X\subset\Lambda_j\\ X\not\subset\widetilde{\square}^{2}(z)}}
\bigl|\partial_t^2\bigl[\mathbf{E}^{(j)}(X,U_k(V^{t,\hat T}),z)\bigr]\bigr|_0 .
\end{equation*}
Insert the third line of \eqref{4.36:eq-regular-absorption-a} and exchange the sums over $z$ and $X$;
the condition $X\not\subset\widetilde{\square}^{2}(z)$ is dropped, all terms being non-negative. By the
polymer count, each $X$ is reached from at most $c_{\#}(1+d_j)M^4$ sites $z$. If $z\in X\cap\Gamma_n$, then
$n\ge n_X$ by the definition of $n_X$, and so
\begin{equation*}
L^{-4(n-j)}e^{-(n-j)}\le L^{-4(n_X-j)}e^{-(n_X-j)}.
\end{equation*}
Write $\hat e_{X^{+}}=\sum_{y\in\mathfrak{b}(X^{+})}e_y^{1/2}$ and exchange the sums over $X$ and $y$. This
leaves, for each block $y$, the sum over the $X$ with $y\in\mathfrak{b}(X^{+})$:
\begin{equation}\label{4.36:eq-regular-absorption-1}
\begin{split}
T^E\le\frac{2e^2L^8E_0c_{\#}M^4}{R^2}&\sum_{n'\le k}\sum_{y\in\Gamma_{n'}}e_y^{1/2}\\
&\times\sum_j\sum_{\substack{X\in\mathbf{D}_j:\ X\subset\Lambda_j\\ y\in\mathfrak{b}(X^{+})}}
(1+d_j)L^{-4(n_X-j)}e^{-(n_X-j)}e^{-\kappa d_j/2}.
\end{split}
\end{equation}

Fix $y\in\Gamma_{n'}$. Since $X\subset\Lambda_j\subset\Omega_j$, the polymer count gives $j\le n'+2$ and
$n_X\le n'+2$ together with the count displayed in the hypotheses. Also $X\subset\Omega_j$, and the
$\Omega_n$ decrease, so $X$ meets no $\Gamma_n$ with $n<j$; hence $n_X\ge j$. Fix the value $d_j(X)=\ell$
and distinguish two cases.

\emph{Case $n_X\ge n'-3$.} Then $L^{4(n'-n_X)}\le L^{12}$ and $e^{n'-n_X}\le e^3$, so
\begin{equation*}
L^{-4(n_X-j)}e^{-(n_X-j)}\le L^{12}e^3L^{-4(n'-j)}e^{-(n'-j)} .
\end{equation*}
Multiply by the count $7^4L^{4(n'+2-j)}e^{c_d\ell}$ and by $(1+\ell)e^{-\kappa\ell/2}$. Since $n'-j\ge-2$,
$e^{-(n'-j)}\le e^2e^{-(n'-j)_+}$ whatever the sign of $n'-j$. These $X$ therefore contribute at most
\begin{equation}\label{4.36:eq-regular-absorption-2}
7^4L^{20}e^3e^{-(n'-j)}(1+\ell)e^{-(\kappa/2-c_d)\ell}
\le 7^4L^{20}e^5(1+\ell)e^{-(\kappa/2-c_d)\ell}e^{-(n'-j)_+}.
\end{equation}

\emph{Case $n_X\le n'-4$.} Since $y\in\mathfrak{b}(X^{+})$ has scale $n'$, $X$ meets a block of scale at
least $n'-2$, and hence meets $\Omega_{n'-2}$. It also meets $\Gamma_{n_X}\subset\Omega^{c}_{n_X+1}$, and
$\Omega^{c}_{n_X+1}\subset\Omega^{c}_{n'-3}$ because $n_X+1\le n'-3$. The geometric statement for
polymers meeting $\Omega_{n'-2}$ and $\Omega^{c}_{n'-3}$ gives $L^{n'-3-j}\le d_j+2$. Hence
\begin{equation*}
L^{4(n'+2-j)}=L^{20}L^{4(n'-3-j)}\le L^{20}(d_j+2)^4 .
\end{equation*}
Next, $e^{(n'-j)_+}\le e^3(d_j+2)$. If $n'-3-j\ge0$, this follows from $e^{n'-3-j}\le L^{n'-3-j}$, since
$e<3\le L$. Otherwise $e^{(n'-j)_+}\le e^3$ and $1\le d_j+2$. Finally $n_X\ge j$ gives
$L^{-4(n_X-j)}e^{-(n_X-j)}\le1$. These $X$ therefore contribute at most
\begin{equation}\label{4.36:eq-regular-absorption-3}
7^4L^{20}e^3(1+\ell)(\ell+2)^5e^{-(\kappa/2-c_d)\ell}e^{-(n'-j)_+}.
\end{equation}

Adding \eqref{4.36:eq-regular-absorption-2} and \eqref{4.36:eq-regular-absorption-3} bounds the
contribution of the $X$ with $d_j=\ell$ by
\begin{equation*}
7^4L^{20}e^5(1+\ell)\big[1+(\ell+2)^5\big]e^{-(\kappa/2-c_d)\ell}e^{-(n'-j)_+}.
\end{equation*}
Here $e^3\le e^5$ was used. Next use $1+(\ell+2)^5\le2(\ell+2)^5$, and then $2L^{20}\le L^{24}$, which holds
because $L^4\ge2$. Summing over the value set of $d_j$ gives
\begin{equation*}
\sum_{\substack{X\in\mathbf{D}_j:\ X\subset\Lambda_j\\ y\in\mathfrak{b}(X^{+})}}
(1+d_j)L^{-4(n_X-j)}e^{-(n_X-j)}e^{-\kappa d_j/2}\le C_{\mathrm{ct}}e^{-(n'-j)_+},
\end{equation*}
with $C_{\mathrm{ct}}$ as in \eqref{4.36:eq-regular-absorption-b} below. The sum over $\ell$ in
$C_{\mathrm{ct}}$ is finite because $\kappa/2-c_d\ge3c_d+4>0$ by the threshold condition on $\kappa$.

Next, split the sum over $j\le n'+2$ into $j\le n'$ and $j\in\{n'+1,n'+2\}$:
\begin{equation*}
\sum_{j\le n'+2}e^{-(n'-j)_+}\le\sum_{u\ge0}e^{-u}+2=\frac{e}{e-1}+2 .
\end{equation*}
The block-weight sum over $\Gamma_{n'}$ that enters \eqref{4.36:eq-regular-absorption-1} is
$\sigma_{n'}^{(1/16)}$, and $\sigma_{n'}^{(1/16)}\le C_\sigma e^{-(k-n')}$ by the decay of the block-weight
sums. Finally $\sum_{n'\le k}e^{-(k-n')}\le e/(e-1)$. Inserting these three bounds in
\eqref{4.36:eq-regular-absorption-1} gives
\begin{equation}\label{4.36:eq-regular-absorption-b}
\begin{split}
C_{\mathrm{ct}}&:=7^4L^{24}e^5\sum_{\ell}(1+\ell)(2+\ell)^5e^{-(\kappa/2-c_d)\ell}<\infty,\\
T^E&\le\frac{C_E^t}{R^2},\qquad
C_E^t:=2e^2L^8E_0c_{\#}M^4C_{\mathrm{ct}}\Big(2+\frac{e}{e-1}\Big)C_\sigma\frac{e}{e-1}.
\end{split}
\end{equation}

\smallskip\noindent\textit{Collection.}
The function $F^E(V^{t,\hat T})$ is a finite sum of real-analytic germs: the local parts
$\mathfrak{e}^{\mathrm{loc}}_{j,z}$ treated in the proof of Theorem~\ref{4.36:thm-regular-bound}, and the
tail terms estimated above. On the finite torus these germs have the common radius
$\min\{\min_{m\le k}r_m,\,2R\}>0$, where the $r_m$ are the radii of the local parts. By the triangle
inequality, the corresponding display and \eqref{4.36:eq-regular-absorption-b},
\begin{equation*}
|\partial_t^2F^E(V^{t,\hat T})|_0\le\mathrm{Main}^E+T^E
\le C_E^m[\delta_k^{-2}+r_0^{-2}]+C_E^tR^{-2}=\mathfrak{E}^{(2)},
\end{equation*}
which is the corresponding display.

\smallskip\noindent\textit{Absorption.}
For each member of the minimum that defines $R$, the threshold condition bounds $R^{-2}$ as displayed in the
hypotheses. Moreover
\begin{equation*}
\delta_k^{-2}=\frac{x_k}{A_1^2(\log x_k)^{2p_0}},\qquad
\varepsilon_k^{-2}=\frac{x_k}{A_0^2(\log x_k)^{2p_0}}.
\end{equation*}
Inserting these into $\mathfrak{E}^{(2)}$ gives
\begin{align*}
\mathfrak{E}^{(2)}
&\le C_E^m\delta_k^{-2}+C_E^mr_0^{-2}
+C_E^t\big[(4C_{\chi}'C_H/c_{\star})^2\delta_k^{-2}+(16C_{\mathrm{loc}})^2\varepsilon_k^{-2}
+4r_0^{-2}\big]\\
&=\mathcal{C}_1x_k(\log x_k)^{-2p_0}+\mathcal{C}_2,
\end{align*}
with $\mathcal{C}_1,\mathcal{C}_2$ as in \eqref{4.36:eq-regular-absorption-c}. This is the first line there.

Let $x_k\ge x_E$. Then $\log x_k\ge(16\mathcal{C}_1)^{1/(2p_0)}>0$, so $(\log x_k)^{2p_0}\ge16\mathcal{C}_1$
and $\mathcal{C}_1x_k(\log x_k)^{-2p_0}\le x_k/16$. Also $\mathcal{C}_2\le x_E/16\le x_k/16$. Adding the two
bounds gives $\mathfrak{E}^{(2)}\le x_k/8$.

The constants $C_E^m$, $C_E^t$, $A_0$, $A_1$, $C_{\chi}'$, $C_H$, $c_{\star}$, $C_{\mathrm{loc}}$ and
$r_0$ entering $x_E$ are $(N,L,M,\mathfrak{p})$-valued, so $x_E$ is an $(N,L,M,\mathfrak{p})$-valued floor on
the inverse coupling. By Lemma~\ref{4.36:lem-threshold-floors}, this floor is imposed for every $k\le K$ by
the condition $\gamma\le(x_E+c_5)^{-1/2}$.
\end{proof}

\begin{remark}[Group-independence of the regular-term bound]\label{4.36:rem-regular-group-free}
The derivation of the bound in Theorem~\ref{4.36:thm-regular-bound} uses none of the following
quantities attached to $SU(N)$. Neither does the absorption of the regular-term bound into a
fraction of the inverse coupling (part~(c) of the same theorem). The excluded
quantities are:
\begin{itemize}
\item the dimension $N^2-1$ of $\mathfrak{g}$;
\item a Casimir;
\item the normalisation $\operatorname{tr}\mathbb{1}=2$;
\item the isomorphism $\operatorname{Ad}(SU(2))\cong SO(3)$;
\item a constant bounding commutators;
\item a normalising constant of a Gaussian or of Haar measure.
\end{itemize}
The integer $N$ enters the bound and its absorption only through the numerical values of
constants each of which is fixed after $N$; among them are the constants $E_1'$, $c_B$, $E_0$,
$C_H$ and $c_{\star}$ entering the proof of Theorem~\ref{4.36:thm-regular-bound}. The constants
$C_\sigma$ and $C_{\mathrm{loc}}$ are used as stated in \cite[Lemma~2.1]{B-CMP96} and in the
proof of Theorem~\ref{4.36:thm-regular-bound} respectively; in these forms neither depends
on $N$.
\end{remark}

\begin{proof}
We go through the ingredients of the two derivations in turn.
\begin{itemize}
\item Derivatives in the deformation parameter are bounded by a Cauchy estimate in one complex
variable. This estimate involves only the radius of analyticity and the size of the function.
\item The chart norm $\|\cdot\|'$ on $\square_0$ is a supremum of operator norms.
\item The membership bound is formulated for the norm $\|\cdot\|_{\star}$, with radius constant
$c_{\star}$. The domains on which it is formulated are sets defined by bounds on the operator
norm $|\cdot|$ of $N\times N$ matrices.
\item The response bound through which the response constant $C_H$ enters is a bound on
operator norms. It is stated per unit of the operator norm $|T'|$ of the conjugated direction
$T'=u\hat{T}u^{-1}$, with $u\in SU(N)$.
\item The adjoint action $X\mapsto uXu^{-1}$ is used only through the identity
$|uXu^{-1}|=|X|$, that is, only as an isometry.
\item Every sum that occurs runs over scales, sites, blocks, cubes and polymers of the lattice of
dimension $d=4$. Each such sum is weighted by the exponential decay factors of the proof of
Theorem~\ref{4.36:thm-regular-bound}, with decay constants $C_\sigma$ and $\delta_0$. The
gauge group indexes none of these sums.
\end{itemize}
Each ingredient is therefore formulated and estimated in terms of operator norms and of the
geometry of the four-dimensional lattice alone. None of the six quantities listed in the
statement above appears in any of them. The constants $E_1'$, $c_B$, $E_0$, $C_H$ and
$c_{\star}$ enter these ingredients only through their numerical values.
\end{proof}

\begin{lemma}[Summability of coupling powers along the flow]\label{4.36:lem-coupling-powers}
Let $N\ge 2$ and $K\ge 0$. Let $g_0,\dots,g_K$ be couplings with $0<g_j<1$, and write $x_j:=g_j^{-2}$, so that $x_j>1$ for $0\le j\le K$. Assume the following two hypotheses.
\begin{itemize}
\item[(F)] The lower member of the two-sided flow inequality of Notation~\ref{4.36:not-sun-conventions} holds: with $c_5\ge 0$ and $\lambda\ge\lambda^{\sharp}$ the constants of that inequality,
\begin{equation}\label{4.36:eq-coupling-powers-1}
\lambda(j-i)-c_5\le x_i-x_j\qquad\text{for }0\le i\le j\le K .
\end{equation}
Here $\lambda^{\sharp}$ is the constant of Notation~\ref{4.36:not-sun-conventions}, which satisfies
\begin{equation*}
\lambda^{\sharp}\ge \frac{c_0(N)}{4}=\frac{11N^2}{48\pi^2}\log L .
\end{equation*}
\item[(Q)] The exponent $\kappa_0$ fixed among the auxiliary parameters $\mathfrak{p}$ in Notation~\ref{4.36:not-sun-conventions} satisfies $\kappa_0\ge 7$.
\end{itemize}
For $0\le j,m\le K$ put
\begin{equation*}
\varpi_{jm}:=\max\Bigl\{1,\Bigl(\frac{x_j}{x_m}\Bigr)^{2}\Bigr\},
\qquad
C_{\varpi}:=2^{\kappa_0/2+1}\Bigl(1+\frac{2}{\lambda^{\sharp}}\Bigr).
\end{equation*}
\begin{itemize}
\item[(a)] If $0\le m\le K$ and $x_m\ge 2c_5$, then
\begin{equation}\label{4.36:eq-coupling-powers-a}
\sum_{0\le j\le m}\varpi_{jm}\,g_j^{\kappa_0}\le C_{\varpi}\,x_m^{1-\kappa_0/2}.
\end{equation}
\item[(b)] Let $0\le n\le k\le K$ and put $m(j,n):=\max\{j,n-2\}$. If $x_k\ge 2c_5$ and, in case $n\ge 2$, also $x_{n-2}\ge 2c_5$, then
\begin{equation}\label{4.36:eq-coupling-powers-b}
\begin{aligned}
\sum_{0\le j\le n}\varpi_{j,m(j,n)}\,g_j^{\kappa_0}
&\le C_{\varpi}\,x_{n-2}^{1-\kappa_0/2}+x_{n-1}^{1-\kappa_0/2}+x_n^{1-\kappa_0/2}\\
&\le (C_{\varpi}+2)\,2^{\kappa_0/2-1}\,x_k^{1-\kappa_0/2},
\end{aligned}
\end{equation}
where terms carrying a negative scale index are absent.
\end{itemize}
\end{lemma}

\begin{proof}
\emph{Part (a).} Fix $m$ with $x_m\ge 2c_5$. For $0\le j\le m$ put $s:=m-j\ge 0$. Hypothesis (F), applied with $(i,j)$ replaced by $(j,m)$, gives
\begin{equation*}
x_j\ge x_m+\lambda s-c_5 .
\end{equation*}
Since $\lambda\ge\lambda^{\sharp}$ and $s\ge 0$, we have $\lambda s\ge\lambda^{\sharp}s$. Moreover $x_m-c_5\ge x_m/2$, because this inequality is equivalent to $x_m\ge 2c_5$. Hence
\begin{equation}\label{4.36:eq-coupling-powers-2}
x_j\ge y_s:=\frac{x_m}{2}+\lambda^{\sharp}s>0 .
\end{equation}
Since $g_j^{\kappa_0}=x_j^{-\kappa_0/2}$, the definition of $\varpi_{jm}$ gives
\begin{equation*}
\varpi_{jm}\,g_j^{\kappa_0}=\max\bigl\{x_j^{-\kappa_0/2},\;x_j^{2-\kappa_0/2}x_m^{-2}\bigr\}.
\end{equation*}
Both entries of the maximum are decreasing functions of $x_j>0$: the exponents $-\kappa_0/2$ and $2-\kappa_0/2$ are negative because $\kappa_0>4$ by (Q). By \eqref{4.36:eq-coupling-powers-2}, and because a maximum of two non-negative numbers is at most their sum,
\begin{equation}\label{4.36:eq-coupling-powers-3}
\varpi_{jm}\,g_j^{\kappa_0}\le y_s^{-\kappa_0/2}+y_s^{2-\kappa_0/2}\,x_m^{-2}.
\end{equation}

We now bound sums of powers of $y_s$. Let $a>1$. Since $y_s$ is increasing in $s$, we have $y_s^{-a}\le\int_{s-1}^{s}y_\sigma^{-a}\,d\sigma$ for $s\ge 1$, where $y_\sigma:=x_m/2+\lambda^{\sharp}\sigma$. Therefore
\begin{equation}\label{4.36:eq-coupling-powers-4}
\begin{aligned}
\sum_{s\ge 0}y_s^{-a}
&\le y_0^{-a}+\int_0^{\infty}\Bigl(\frac{x_m}{2}+\lambda^{\sharp}\sigma\Bigr)^{-a}d\sigma\\
&=\Bigl(\frac{2}{x_m}\Bigr)^{a}+\frac{1}{\lambda^{\sharp}(a-1)}\Bigl(\frac{2}{x_m}\Bigr)^{a-1}.
\end{aligned}
\end{equation}
The sum over $0\le j\le m$ is the sum over $0\le s\le m$, which is at most the sum over all $s\ge 0$. We apply \eqref{4.36:eq-coupling-powers-4} to the two terms of \eqref{4.36:eq-coupling-powers-3} in turn. Throughout we use $x_m\ge 1$, which gives $x_m^{-\kappa_0/2}\le x_m^{1-\kappa_0/2}$.

\emph{First term.} Take $a=\kappa_0/2$. Then $a-1=\kappa_0/2-1\ge 5/2$ by (Q), so $1/(a-1)\le 2/5$. Hence
\begin{align*}
\sum_{s\ge 0}y_s^{-\kappa_0/2}
&\le 2^{\kappa_0/2}x_m^{-\kappa_0/2}
+2^{\kappa_0/2-1}x_m^{1-\kappa_0/2}\cdot\frac{2}{5}\,(\lambda^{\sharp})^{-1}\\
&\le 2^{\kappa_0/2}\Bigl[1+\frac{1}{5\lambda^{\sharp}}\Bigr]x_m^{1-\kappa_0/2}.
\end{align*}

\emph{Second term.} Take $a=\kappa_0/2-2$. Then $a>1$ because $\kappa_0\ge 7$, and $a-1=\kappa_0/2-3\ge 1/2$, so $1/(a-1)\le 2$. Applying \eqref{4.36:eq-coupling-powers-4} and multiplying by $x_m^{-2}$ gives
\begin{align*}
x_m^{-2}\sum_{s\ge 0}y_s^{2-\kappa_0/2}
&\le 2^{\kappa_0/2-2}x_m^{-\kappa_0/2}
+2^{\kappa_0/2-3}\cdot 2(\lambda^{\sharp})^{-1}x_m^{1-\kappa_0/2}\\
&\le 2^{\kappa_0/2-2}\Bigl[1+\frac{1}{\lambda^{\sharp}}\Bigr]x_m^{1-\kappa_0/2}.
\end{align*}

\emph{Conclusion of (a).} Using $1/(5\lambda^{\sharp})\le 1/\lambda^{\sharp}$ in the first bound and adding the two bounds, \eqref{4.36:eq-coupling-powers-3} yields
\begin{align*}
\sum_{0\le j\le m}\varpi_{jm}\,g_j^{\kappa_0}
&\le\bigl[2^{\kappa_0/2}+2^{\kappa_0/2-2}\bigr]\Bigl(1+\frac{1}{\lambda^{\sharp}}\Bigr)x_m^{1-\kappa_0/2}\\
&=\frac{5}{4}\,2^{\kappa_0/2}\Bigl(1+\frac{1}{\lambda^{\sharp}}\Bigr)x_m^{1-\kappa_0/2}\\
&\le C_{\varpi}\,x_m^{1-\kappa_0/2}.
\end{align*}
The last step holds because $5/4\le 2$ and $1+1/\lambda^{\sharp}\le 1+2/\lambda^{\sharp}$. This proves \eqref{4.36:eq-coupling-powers-a}.

\emph{Part (b).} Let $0\le n'\le k$. Hypothesis (F), applied with $(i,j)$ replaced by $(n',k)$, gives
\begin{equation*}
x_{n'}\ge x_k+\lambda(k-n')-c_5\ge x_k-c_5 .
\end{equation*}
Since $x_k\ge 2c_5$, this gives
\begin{equation}\label{4.36:eq-coupling-powers-5}
x_{n'}\ge \frac{x_k}{2}\qquad\text{for }0\le n'\le k .
\end{equation}

We split the sum in \eqref{4.36:eq-coupling-powers-b} into the indices $j\le n-2$, for which $m(j,n)=n-2$, and the indices $j\in\{n-1,n\}$, for which $m(j,n)=j$ and $\varpi_{jj}=1$. Terms carrying a negative scale index are absent; in particular, for $n\le 1$ the first sum is empty.

For $n\ge 2$, the hypothesis $x_{n-2}\ge 2c_5$ allows part (a) to be applied with $m=n-2$. It bounds the sum over $0\le j\le n-2$ by $C_{\varpi}x_{n-2}^{1-\kappa_0/2}$.

For $j\in\{n-1,n\}$ with $j\ge 0$, we have $\varpi_{jj}g_j^{\kappa_0}=x_j^{-\kappa_0/2}\le x_j^{1-\kappa_0/2}$, since $x_j\ge 1$. This proves the first inequality in \eqref{4.36:eq-coupling-powers-b}.

For the second inequality, note that each of the indices $n'\in\{n-2,n-1,n\}$ that occurs satisfies $n'\le n\le k$. By \eqref{4.36:eq-coupling-powers-5}, and because the exponent $1-\kappa_0/2$ is negative,
\begin{equation*}
x_{n'}^{1-\kappa_0/2}\le\Bigl(\frac{x_k}{2}\Bigr)^{1-\kappa_0/2}=2^{\kappa_0/2-1}x_k^{1-\kappa_0/2}.
\end{equation*}
Summing the coefficients $C_{\varpi}$, $1$ and $1$ gives the bound $(C_{\varpi}+2)\,2^{\kappa_0/2-1}x_k^{1-\kappa_0/2}$, which is the second inequality in \eqref{4.36:eq-coupling-powers-b}.
\end{proof}

\begin{remark}\label{4.36:thm-operation-bound}
The statement of this theorem is not written out here; parts of its proof are printed below.
\end{remark}

\begin{proof}[Proof of Theorem~\ref{4.36:thm-operation-bound}, continued]
We keep the notation of the earlier steps of this proof. For $j\le k$ and $X\in D_j^{R}$ with
$X\subset\Lambda_j$, $n_X$ is the scale index and $z_X\in X\cap\Gamma_{n_X}$ the point chosen
for $X$. We put $m:=\max\{j,n_X-2\}$, let $\square\in\pi_m$ be the cube containing $z_X$, and
put $\square_0:=\widetilde{\square}^{5}$. The polymer $X$ is \emph{small} if
$X\subset\widetilde{\square}^{2}$ and a \emph{tail} polymer otherwise. We use the per-polymer
bound with $S_X$, the bound on $\Theta_m$ and the first, second and third counts collected in
the corresponding display.

\emph{The small part, summed.} Each small $X$ is counted once, namely at the triple
$(j,n_X,\square)$ with $\square\ni z_X$. Fix $(j,n)$. The level $m=\max\{j,n-2\}$ depends only
on $(j,n)$, and so do $\varpi_{jm}$ and $\Theta_m$.

First, insert the per-polymer bound. Then sum over the small $X$ inside one cube, using the
first count. Then sum over the cubes $\square$, using the second count. This gives
\begin{align*}
  &\sum_{X\in D_j^{R}\ \text{small},\ n_X=n}
  \bigl|\partial_t^2\bigl[\mathbf{R}^{(j)}\bigl(X,U_k(V^{t,\hat{T}})\bigr)\bigr]\bigr|_0\\
  &\qquad\le 2\cdot 7^4\cdot 13^4\cdot 4e^2\,C_{\sigma}R_1'\,\varpi_{jm}\,g_j^{\kappa_0}\,
  e^{-(k-n)}\,\Theta_m .
\end{align*}
By the bound on $\Theta_m$ and the third count,
\begin{align*}
  \Theta_m&\le\bigl(1+\mathfrak{f}(k-n+2)\bigr)\bigl[C_{\Theta}\delta_k^{-2}+C_{\Theta}'r_0^{-2}\bigr]\\
  &\le\bigl(1+\mathfrak{f}(k-n+2)\bigr)\max\{C_{\Theta},C_{\Theta}'\}
  \bigl[\delta_k^{-2}+r_0^{-2}\bigr].
\end{align*}

Next sum over $j\le n$; the restriction $j\le n\le k$ was established in the earlier steps of the
proof. By \eqref{4.36:eq-coupling-powers-b} of Lemma~\ref{4.36:lem-coupling-powers}, whose
reading at $m=n-2$ uses $x_{n-2}\ge 2c_5$ from the standing hypotheses,
\begin{equation*}
  \sum_{0\le j\le n}\varpi_{j,\max\{j,n-2\}}\,g_j^{\kappa_0}
  \le (C_{\varpi}+2)\,2^{\kappa_0/2-1}x_k^{1-\kappa_0/2}.
\end{equation*}
Finally sum over $n\le k$. With $u:=k-n$,
\begin{equation*}
  \sum_{n\le k}e^{-(k-n)}\bigl(1+\mathfrak{f}(k-n+2)\bigr)
  \le\sum_{u\ge0}e^{-u}\bigl(1+\mathfrak{f}(u+2)\bigr)=:S_{\mathfrak{f}}<\infty .
\end{equation*}
Since $2\cdot 4=8$, multiplying the bounds just obtained and summing gives the bound on
$\mathfrak{r}_{\mathrm{small}}$ in the first line of \eqref{4.36:eq-operation-absorption-a}
below, with the constants $C_R$ and $S_{\mathfrak{f}}$ defined there.

\emph{The tail part, per polymer.} Let $X$ be a tail polymer. Then $X$ meets
$\square$ (it contains $z_X$) and the complement of $\widetilde{\square}^{2}$, so
$\operatorname{diam}_\infty X\ge 2M\ell_m$. The diameter bound for $d_j$ then gives
$d_j(X)\ge 2(L^{m-j}-1)$.

Put $\Psi:=\mathbf{R}^{(j)}(X,\cdot)$, a function of the configuration on the bonds of $X$,
and let $\mathsf{N}$ be the set of bonds of $X$. For $|\zeta|\le 2R$ the perturbed background
$e^{i\eta\chi_{\mathsf{N}}\mathcal{H}_{\zeta}}U_k(V)$, read on the bonds of $X$, with the
parameter $\eta$ and the cut-off $\chi_{\mathsf{N}}$ as in the argument leading to
\eqref{4.36:eq-regular-absorption-a}, lies, at half radius, in the
domain on which \cite[(2.31)]{B-CMP119} states
\begin{equation*}
  |\mathbf{R}^{(j)}(X,\cdot)|\le g_j^{\kappa_0}e^{-\kappa d_j(X)} .
\end{equation*}
Membership in that domain follows from the membership criterion in the norm
$\|\cdot\|_{\star}$ with radius constant $c_{\star}$. The criterion is reached through two facts.
First, the bound $\|\chi_{\mathsf{N}}\mathcal{H}_{\zeta}\|_{\star}\le 2RC_{\chi}'C_H/\delta_k$ on
the response field, supplied by the operator bounds at the free link. Second, the first member of
the minimum defining $R$, which makes this at most $c_{\star}/2$. That $U_k(V)$, restricted to the
bonds of $X\subset\Lambda_j$, lies in that domain with the margin needed to add the perturbation is
one of the statements from which Theorem~\ref{4.36:thm-operation-bound} is proved.

Write $u:=m-j\ge0$. Since $m\ge n_X-2$, we have $n_X-j\le u+2$; moreover
$d_j(X)/2\ge L^{u}-1$. The arithmetic that led to \eqref{4.36:eq-regular-absorption-a}, which
splits $e^{-\kappa d_j(X)}$ into two halves, needs only $\kappa\ge5$, and so holds under the
condition $\kappa\ge10$ on the auxiliary parameters; it therefore applies unchanged
and gives, on the ball of the membership criterion,
\begin{equation*}
  |\mathbf{R}^{(j)}(X,\cdot)|
  \le e^2L^8g_j^{\kappa_0}L^{-4(n_X-j)}e^{-(n_X-j)}e^{-\kappa d_j(X)/2}.
\end{equation*}
The second-derivative bound through the oscillation, from the operator bounds at the free link,
is applied to $\Psi$ with oscillation at most $2\sup|\Psi|$ on $|t|\le 2R$. It gives
\begin{equation}\label{4.36:eq-operation-absorption-2}
  \bigl|\partial_t^2\bigl[\mathbf{R}^{(j)}\bigl(X,U_k(V^{t,\hat{T}})\bigr)\bigr]\bigr|_0
  \le 2e^2L^8g_j^{\kappa_0}L^{-4(n_X-j)}e^{-(n_X-j)}e^{-\kappa d_j(X)/2}\,
  \frac{\hat{e}_{X^{+}}}{R^2},
\end{equation}
where $\hat{e}_{X^{+}}=\sum_{y\in\mathfrak{b}(X^{+})}e_y^{1/2}$.

\emph{The tail sum.} We write
$\omega_j(X):=L^{-4(n_X-j)}e^{-(n_X-j)}e^{-\kappa d_j(X)/2}$. Summing
\eqref{4.36:eq-operation-absorption-2} over the tail polymers, inserting the definition of
$\hat{e}_{X^{+}}$, exchanging the sums, and enlarging the sum over $X\in D_j^{R}$ to the sum over
all $X\in D_j$ with $X\subset\Lambda_j$, we get
\begin{equation*}
  \mathfrak{r}_{\mathrm{tail}}\le\frac{2e^2L^8}{R^2}\sum_{n'\le k}\ \sum_{y\in\Gamma_{n'}}
  e_y^{1/2}\sum_{j}g_j^{\kappa_0}\sum_{X\in D_j:\,X\subset\Lambda_j,\,y\in\mathfrak{b}(X^{+})}
  \omega_j(X).
\end{equation*}

Fix $y\in\Gamma_{n'}$. By the counting statement for blocks below a polymer, as in the argument
leading to \eqref{4.36:eq-regular-absorption-b}, every $X$ occurring satisfies
$j\le n_X\le n'+2$. Hence $j\le k':=\min\{k,n'+2\}$.

The inner sum over $X$ is bounded through \eqref{4.36:eq-regular-absorption-b}. The argument
establishing that display, with its two cases $n_X\ge n'-3$ and $n_X\le n'-4$, bounds the same
sum with the additional factor $1+d_j(X)\ge1$. So, a fortiori,
\begin{equation*}
  \sum_{X\in D_j:\,X\subset\Lambda_j,\,y\in\mathfrak{b}(X^{+})}\omega_j(X)
  \le C_{\mathrm{ct}}e^{-(n'-j)_{+}}\le C_{\mathrm{ct}} .
\end{equation*}

Next, apply \eqref{4.36:eq-coupling-powers-a} of Lemma~\ref{4.36:lem-coupling-powers} at level
$k'$. This is allowed because $x_{k'}\ge 2c_5$ by the standing hypotheses, and $\varpi_{jk'}\ge1$.
The lower member of the flow bound, as in the proof of that lemma, gives
$x_{k'}\ge x_k-c_5\ge x_k/2$ once $x_k\ge 2c_5$. Since $1-\kappa_0/2<0$,
\begin{equation*}
  \sum_{j\le k'}g_j^{\kappa_0}\le C_{\varpi}x_{k'}^{1-\kappa_0/2}
  \le C_{\varpi}2^{\kappa_0/2-1}x_k^{1-\kappa_0/2},
\end{equation*}
uniformly in $n'$. The remaining sum over $y\in\Gamma_{n'}$ of $e_y^{1/2}$ is bounded by the block
sum $\sigma_{n'}^{(1/16)}$, which satisfies $\sigma_{n'}^{(1/16)}\le C_{\sigma}e^{-(k-n')}$. Hence
\begin{equation*}
  \sum_{n'\le k}\sigma_{n'}^{(1/16)}\le C_{\sigma}(1-e^{-1})^{-1}.
\end{equation*}
Altogether, multiplying the factor $2e^2L^8/R^2$ by the three bounds just obtained gives the bound
on $\mathfrak{r}_{\mathrm{tail}}$ in \eqref{4.36:eq-operation-absorption-a} below, with the
constant $C''$ defined there.

\emph{Conclusion.} The bounds obtained in the small and the tail parts read, with their constants
named,
\begin{equation}\label{4.36:eq-operation-absorption-a}
  \begin{split}
    \mathfrak{r}_{\mathrm{small}}&\le C_R\bigl[\delta_k^{-2}+r_0^{-2}\bigr]x_k^{1-\kappa_0/2},\\
    C_R&:=8e^2\cdot 7^4\cdot 13^4\cdot C_{\sigma}R_1'S_{\mathfrak{f}}(C_{\varpi}+2)
    \,2^{\kappa_0/2-1}\max\{C_{\Theta},C_{\Theta}'\},\\
    S_{\mathfrak{f}}&:=\sum_{u\ge0}e^{-u}\bigl(1+\mathfrak{f}(u+2)\bigr)<\infty,\\
    \mathfrak{r}_{\mathrm{tail}}&\le C''R^{-2}x_k^{1-\kappa_0/2},\\
    C''&:=2e^2L^8C_{\mathrm{ct}}C_{\varpi}2^{\kappa_0/2-1}C_{\sigma}(1-e^{-1})^{-1}.
  \end{split}
\end{equation}

The constants $\mathbf{R}^{(j)}(X,1)$ entering $F^R$ do not depend on $t$, so $\partial_t$
annihilates them. Adding the two bounds therefore gives the corresponding display.
Real-analyticity of $t\mapsto F^R(V^{t,\hat{T}})$ near $0$ holds because each term
$t\mapsto\mathbf{R}^{(j)}(X,U_k(V^{t,\hat{T}}))$ is real-analytic near $0$ by the identifications
used above, and there are finitely many $X$ on the finite torus. The independence of $F^R$ from
$w$ holds because $\mathbf{R}$ reads $U_k(V)$ only.
\end{proof}

\begin{lemma}[The operation-term bound absorbed into a unit]\label{4.36:prf-operation-absorption}
Assume the following.
\begin{itemize}
  \item The standing hypotheses of Definition~\ref{4.36:def-near-part} and the hypotheses of
    Theorem~\ref{4.36:thm-operation-bound}, so that, by that theorem,
    the corresponding display holds with the constants $C_R$, $C''$ of
    \eqref{4.36:eq-operation-absorption-a}. Among these hypotheses the following are used again
    in the proof above:
    \begin{itemize}
      \item the bound \cite[(2.31)]{B-CMP119};
      \item the membership criterion in the norm $\|\cdot\|_{\star}$;
      \item the statement that $U_k(V)$, restricted to the bonds of polymers $X\subset\Lambda_j$,
        lies in the domain of \cite[(2.31)]{B-CMP119} with half the radius to spare;
      \item the response-field bound and the second-derivative bound through the oscillation,
        both from the operator bounds at the free link;
      \item the counting statements for polymers and blocks, including the block-sum estimate
        with constant $C_{\sigma}$;
      \item the distance lower bound for a polymer meeting two separated strata, as used in the
        argument leading to \eqref{4.36:eq-regular-absorption-b}, and the animal bound at rate
        $\kappa$, as used in the first count of the corresponding display;
      \item the lower member of the flow bound through Lemma~\ref{4.36:lem-coupling-powers}.
    \end{itemize}
  \item $\kappa_0\ge7$ and $\kappa\ge10$.
  \item The radius $R$ is the minimum of its three defining members; whichever member attains
    the minimum, $R^{-2}$ is at most the sum of the three inverse squares:
    \begin{equation*}
      R^{-2}\le(4C_{\chi}'C_H/c_{\star})^2\delta_k^{-2}+(16C_{\mathrm{loc}})^2\varepsilon_k^{-2}
      +4r_0^{-2}.
    \end{equation*}
\end{itemize}
Put
\begin{equation}\label{4.36:eq-operation-absorption-b}
  \begin{split}
    x_R:=\max\Bigl\{&e,\ 2c_5,\ x_{\delta},\ x_{\beta},\
    \Bigl(\frac{4C_R}{A_1^2}\Bigr)^{2/(\kappa_0-4)},\
    \Bigl(\frac{4C_R}{r_0^2}\Bigr)^{2/(\kappa_0-2)},\\
    &\Bigl(4C''\Bigl[\Bigl(\frac{4C_{\chi}'C_H}{c_{\star}A_1}\Bigr)^2
    +\Bigl(\frac{16C_{\mathrm{loc}}}{A_0}\Bigr)^2\Bigr]\Bigr)^{2/(\kappa_0-4)},\
    \Bigl(\frac{16C''}{r_0^2}\Bigr)^{2/(\kappa_0-2)}\Bigr\}.
  \end{split}
\end{equation}
Here $x_{\delta}$ and $x_{\beta}$ are the floors on the inverse coupling of
Lemma~\ref{4.36:lem-threshold-floors} which, together with $\bar{B}$, express the ceiling
$\gamma\le\gamma_{\bar{B}}$ among the standing hypotheses of Definition~\ref{4.36:def-near-part}.
If $x_n\ge x_R$ for every $n\le K$, then $\mathfrak{r}_{\mathrm{small}}\le\tfrac12$,
$\mathfrak{r}_{\mathrm{tail}}\le\tfrac12$ and $|\partial_t^2F^R(V^{t,\hat{T}})|_0\le1$.
The floor $x_n\ge x_R$ for all $n\le K$ is imposed through the one-sided ceiling
$\gamma\le(x_R+c_5)^{-1/2}$, a member of $\gamma^{(3)}$.
\end{lemma}

\begin{proof}
For $x_k\ge e$ we have $\log x_k\ge1$, hence
\begin{equation*}
  \delta_k^{-2}=\frac{x_k}{A_1^2(\log x_k)^{2p_0}}\le\frac{x_k}{A_1^2},\qquad
  \varepsilon_k^{-2}=\frac{x_k}{A_0^2(\log x_k)^{2p_0}}\le\frac{x_k}{A_0^2}.
\end{equation*}

\emph{The small sum.} Insert the first of these bounds into
$\mathfrak{r}_{\mathrm{small}}\le C_R[\delta_k^{-2}+r_0^{-2}]x_k^{1-\kappa_0/2}$. This gives
\begin{equation*}
  \mathfrak{r}_{\mathrm{small}}\le\frac{C_R}{A_1^2}x_k^{2-\kappa_0/2}
  +C_Rr_0^{-2}x_k^{1-\kappa_0/2}.
\end{equation*}

\emph{The tail sum.} Insert both bounds and the hypothesis on $R^{-2}$ into
$\mathfrak{r}_{\mathrm{tail}}\le C''R^{-2}x_k^{1-\kappa_0/2}$. This gives
\begin{equation*}
  \mathfrak{r}_{\mathrm{tail}}\le C''\Bigl[\Bigl(\frac{4C_{\chi}'C_H}{c_{\star}A_1}\Bigr)^2
  +\Bigl(\frac{16C_{\mathrm{loc}}}{A_0}\Bigr)^2\Bigr]x_k^{2-\kappa_0/2}
  +4C''r_0^{-2}x_k^{1-\kappa_0/2}.
\end{equation*}
No further power of $\log x_k$ enters through $C_R$ or $C''$, since $c_B$, $R_1'$, $C_H$,
$C_{\sigma}$, $C_{\mathrm{ct}}$ and $C_{\varpi}$ do not depend on the coupling.

\emph{Each of the four summands is at most $\tfrac14$.} Since $\kappa_0\ge7$, both exponents
$2-\kappa_0/2$ and $1-\kappa_0/2$ are negative. For a positive constant $a$ and an exponent
$-b<0$,
\begin{equation*}
  a\,x^{-b}\le\tfrac14\iff x\ge(4a)^{1/b}.
\end{equation*}
With $b=\kappa_0/2-2$, so that $1/b=2/(\kappa_0-4)$, this applies to the first and third
summands. With $b=\kappa_0/2-1$, so that $1/b=2/(\kappa_0-2)$, it applies to the second and
fourth summands, with $a=C_Rr_0^{-2}$ and $a=4C''r_0^{-2}$ respectively. For instance,
\begin{equation*}
  \frac{C_R}{A_1^2}x^{2-\kappa_0/2}\le\frac14\quad\Longleftarrow\quad
  x\ge\Bigl(\frac{4C_R}{A_1^2}\Bigr)^{2/(\kappa_0-4)}.
\end{equation*}
Hence each summand is at most $\tfrac14$ as soon as $x_k\ge x_R$, with $x_R$ as in
\eqref{4.36:eq-operation-absorption-b}. A base below $1$ gives a member below $1$, which is
covered by the member $e$. Once $r_0$ is fixed at a value admitted by its defining conditions,
$x_R$ is a number.

Consequently $\mathfrak{r}_{\mathrm{small}}\le\tfrac14+\tfrac14=\tfrac12$ and
$\mathfrak{r}_{\mathrm{tail}}\le\tfrac12$. By the corresponding display,
$|\partial_t^2F^R(V^{t,\hat{T}})|_0\le1$.

The floor is imposed for every $n\le K$, and not only at $n=k$. The reason is that the proof of
Theorem~\ref{4.36:thm-operation-bound} used $x_n\ge 2c_5$, $x_n\ge x_{\delta}$ and $x_n\ge e$ at
the levels $n$ at which its estimates are read. In terms of the reference coupling, the floor is
the one-sided ceiling $\gamma\le(x_R+c_5)^{-1/2}$, one of the members of $\gamma^{(3)}$.
\end{proof}

\begin{remark}[Group-independence of the operation-term bound]\label{4.36:rem-operation-group-free}
The proof of Theorem~\ref{4.36:thm-operation-bound} uses no property of the gauge group that is
special to $N=2$: $N$ enters it only through the values of constants. The proof does not use the
condition $L^2\ge 1+\bar{B}$, provided that the two inputs through which this condition could enter
are supplied for $\mathrm{SU}(N)$ under $\gamma\le\gamma_{\bar{B}}$, a condition implied by
$\gamma\le\gamma^{(3)}$; this supply rests on the proposition on the auxiliary thresholds, subject
to the check, made separately in this chapter, that its clauses apply to these two inputs.
\end{remark}

\begin{proof}
The proof of Theorem~\ref{4.36:thm-operation-bound} draws on the following inputs:
\begin{itemize}
\item clause (d) of the corollary on the local chart $\mathbf{A}\mapsto\mathsf{U}(\mathbf{A})$ on
  $\square_0$, which is a scalar bound on the ball
  $\mathcal{A}'=\{\|\mathbf{A}\|'<c_B\delta_m\}$ of the chart norm, a norm built from
  $|\cdot|_{\mathfrak{g}}$;
\item the bound \cite[(2.31)]{B-CMP119}, applied on the ball in $|\cdot|_{\mathfrak{g}}$ on which
  the membership bound is stated, that is, the bound asserting membership in a ball of the norm
  $\|\cdot\|_{\star}$ whose radius is fixed by the constant $c_{\star}$;
\item the polymer counting bound and the bound on the sums $\sigma_n^{(\alpha)}$;
\item the bound $\mathfrak{a}(a)\le\mathfrak{s}(a)$ on the lattice-animal sums, where
  $\mathfrak{a}(a)$ is the supremum over cubes $c$ of the sum of $e^{-a d_j(X)}$ over the
  polymers $X\supset c$;
\item the lemma on arrests;
\item the bounds on the coupling flow.
\end{itemize}
The bound \cite[(2.31)]{B-CMP119} and the bounds on the coupling flow carry no power of
$\mathsf{n}=(N-1)^{1/2}$, and every other estimate entering the proof of
Theorem~\ref{4.36:thm-operation-bound} involves $N$ only through the constants described next.
The integer $N$ enters only through the values of constants determined by
$(N,L,M,\mathfrak{p})$: among them $R_1'$, $c_B$ and $C_H$, the flow constants
$\lambda^{\sharp}$ (through $C_{\varpi}$) and $c_5$, and $\bar{B}$ where it occurs as the flow
factor below. In the
one-bond factor $V'=e^{itT'}$, $t$ being the parameter of the one-bond variation, the
direction appears only through $|T'|_{\mathfrak{g}}=1$. No
step of the proof therefore specialises to $N=2$.

The factor of the bound on $\Theta_m=4C_H/(c_B\delta_m r_m)$ that is supplied by the coupling
flow is a constant $\mathfrak{f}\in\{1,\bar{B}\}$, according to which of two bounds on the
running inverse couplings of Definition~\ref{1.2:def-couplings} is used. One of these two bounds
is derived from a further bound on the running inverse couplings together with one clause of a
lemma on the coupling flow. Neither of the two bounds uses the condition $L^2\ge 1+\bar{B}$.

Two inputs are the only places where a condition of the type $L^2\ge 1+\bar{B}$ could enter. The
first is clause (d) of the corollary on the local chart, through the two threshold conditions on
the inverse coupling, for the real and for the imaginary part of the chart datum, on which that
clause rests, and through the bound on the coupling flow that feeds these two conditions. The
second is the condition under which the factor $\tfrac12$ in the radius of the membership bound is
absorbed into $1$, namely the floors \eqref{4.36:eq-threshold-floors-d} on the inverse coupling.
In the proof of
Theorem~\ref{4.36:thm-operation-bound}, for $\mathrm{SU}(N)$, both are supplied by clauses (b)
and (c) of the proposition on the auxiliary thresholds, under $\gamma\le\gamma_{\bar{B}}$, a
condition implied by $\gamma\le\gamma^{(3)}$, subject to the check, made separately in this
chapter, that these clauses apply to the two inputs.
\end{proof}

\begin{lemma}[One arrest per cube per level]\label{4.36:lem-arrest-per-cube}
Consider the arrests $W$ of a history and, for each arrest $W$, its three levels $h_W$, $j_W$,
$k_W$, as in Notation~\ref{4.36:not-history-geometry}. Here $W$ is arrested at step $k_W$,
after the integrations $j=h_W,\dots,j_W-1$ of its attempted $\mathbf{R}$-operation. Assume
the following three properties of the arrests, which belong to the arrest construction of
\cite{B-CMP122a}:
\begin{enumerate}
\item[(a)] the family of arrests is laminar: any two arrests are either disjoint or nested;
\item[(b)] if $W$ and $W'$ are arrests with $W\subsetneq W'$, then $h_{W'}\ge k_W$;
\item[(c)] every arrest $W$ satisfies $h_W<j_W\le k_W$.
\end{enumerate}
Fix a level $m$ and a cube $C\in\pi_m$. Then among the arrests $W$ with $W\supseteq C$ at most
one satisfies $m\in(h_W,k_W]$, that is, $h_W<m\le k_W$.
\end{lemma}

\begin{proof}
Let $W$ and $W'$ be two arrests that both contain $C$. The cube $C$ is not empty, so
$W\cap W'\supseteq C$ is not empty, and $W$ and $W'$ are not disjoint. By hypothesis (a) they
are therefore nested. Thus the arrests containing $C$ are pairwise nested, and hence totally
ordered by inclusion. We list them accordingly as
\begin{equation*}
W_1\subsetneq W_2\subsetneq\dots\subsetneq W_s .
\end{equation*}

We show that for $1\le a<b\le s$ the right end of the interval attached to $W_a$ does not
exceed the left end of the interval attached to $W_b$, that is,
\begin{equation}\label{4.36:eq-arrest-per-cube-1}
k_{W_a}\le h_{W_b}.
\end{equation}
For each index $i$ with $a\le i<b$ we have $W_i\subsetneq W_{i+1}$. Hypothesis (b), applied to
this pair, gives $h_{W_{i+1}}\ge k_{W_i}$. For each index $i$ with $a<i<b$, hypothesis (c),
applied to the arrest $W_i$, gives $h_{W_i}<j_{W_i}\le k_{W_i}$, and in particular
$k_{W_i}\ge h_{W_i}$. Applying these two inequalities alternately, from $i=b-1$ down to $i=a$,
yields the chain
\begin{equation*}
\begin{split}
h_{W_b}\ge k_{W_{b-1}}\ge h_{W_{b-1}}\ge k_{W_{b-2}}\ge\dots\ge h_{W_{a+1}}\ge k_{W_a}.
\end{split}
\end{equation*}
In this chain each inequality of the form $h_{W_{i+1}}\ge k_{W_i}$ comes from hypothesis (b),
and each inequality of the form $k_{W_i}\ge h_{W_i}$ comes from hypothesis (c). When $b=a+1$ the
chain reduces to the single inequality $h_{W_{a+1}}\ge k_{W_a}$ given by hypothesis (b). This
proves \eqref{4.36:eq-arrest-per-cube-1}.

Now let $a<b$. Every integer $n\in(h_{W_a},k_{W_a}]$ satisfies $n\le k_{W_a}$, and every
integer $n'\in(h_{W_b},k_{W_b}]$ satisfies $n'>h_{W_b}$. By \eqref{4.36:eq-arrest-per-cube-1}
these two bounds combine to
\begin{equation*}
n\le k_{W_a}\le h_{W_b}<n'.
\end{equation*}
Hence the intervals $(h_{W_a},k_{W_a}]$ and $(h_{W_b},k_{W_b}]$ are disjoint. So the intervals
$(h_{W_i},k_{W_i}]$, $i=1,\dots,s$, are pairwise disjoint, and the fixed level $m$ lies in at
most one of them.

Every arrest $W\supseteq C$ is one of $W_1,\dots,W_s$. It follows that at most one arrest
$W\supseteq C$ satisfies $m\in(h_W,k_W]$.
\end{proof}

\begin{proposition}[The frozen boundary terms at a fixed polymer]\label{4.36:prop-frozen-sum}
Work with the group, norms and coupling flow of Notation~\ref{4.36:not-sun-conventions} and with
the regions, polymers, cubes and arrests of Notation~\ref{4.36:not-history-geometry}. For an
arrest $W$, arrested at step $k_W$ after the integrations of levels $h_W,\dots,j_W-1$, and for a
stage $j\in[h_W,j_W-1]$ of its attempted large-field operation, let $\mathbf{C}_{W,j}(X)$ denote
a boundary term frozen on $W$ by the stage that integrates level $j$, with support $X$. Assume the
arrest structure of Ba{\l}aban's large-field construction, taken here by statement from
\cite{B-CMP122a,B-CMP122b}, namely \textup{(B.0)}--\textup{(B.3)} below, together with
\textup{(a)}--\textup{(e)}.
\begin{itemize}
\item[(B.0)] Every frozen term $\mathbf{C}_{W,j}(X)$ has $X\cap W\neq\emptyset$.
\item[(B.1)] Every arrest $W$ is a union of cubes of $\pi_{k_W}$, hence a union of cubes of
$\pi_m$ for every $m\le k_W$, the partitions $\pi_n$ being nested.
\item[(B.2)] The family of arrests is laminar, all clauses being read componentwise: two arrests
are disjoint or nested; if $W\subsetneq W'$ then $h_{W'}\ge k_W$; and $h_W<j_W\le k_W$.
\item[(B.3)] A frozen term $\mathbf{C}_{W,j}(X)$ has support $X\in\mathbf{D}_m$ with
$h_W\le j\le j_W-1$ and $j+1\le m\le k_W$; it satisfies
$X\cap\Omega_m^{\natural}\neq\emptyset$ and
$X\cap(Z_{j+1}\setminus Z''_{j+1})\neq\emptyset$, where $Z''_{j+1}\subset W$ is the subset with
$\Omega_{j+1}\cap W=W\setminus Z''_{j+1}$; on its domain it obeys
$|\mathbf{C}_{W,j}(X)|\le C_0^{\mathrm{fr}}e^{-\kappa d_m(X)}$, which is the bound of
\cite[(1.63), (1.70)]{B-CMP122a} in the form carried over to $G=SU(N)$ in this chapter; and there
is at most one such term for each quadruple $(W,j,m,X)$.
\item[(a)] For every arrest $W$, the number of its stages, the stage count of its attempted
operation $\mathbf{R}$, satisfies $\mathfrak{N}_W:=j_W-h_W\le R_{k_W}$ \cite{B-CMP122b}.
\item[(b)] $R_{j+1}\le R_j$ and $R_j\ge1$ for every $j$.
\item[(c)] The sets $Z_j=\Lambda_j^{c}$ increase with $j$, and for every $i$
\begin{equation*}
\operatorname{dist}_\infty(\Omega^{\natural}_{i+1},Z_i)\ge 8MR_{i+1}\ell_{i+1}.
\end{equation*}
\item[(d)] Every $X\in\mathbf{D}_j$ is a union of at most $2^4(3d_j(X)+1)\le 48(1+d_j(X))$ cubes
of $\pi_j$, and
\begin{equation*}
d_j(X)\ge \frac{\operatorname{diam}_\infty X}{M\ell_j}-2 .
\end{equation*}
\item[(e)] $\kappa\ge4$.
\end{itemize}
Hypotheses \textup{(b)}--\textup{(d)} are among the conventions of
Notation~\ref{4.36:not-history-geometry}, and \textup{(e)} is obtained from
\eqref{4.36:eq-sun-conventions-a} together with the condition on $\kappa$ fixed with the
parameters $\mathfrak{p}$.
Then for every $m\le k$ and every $X\in\mathbf{D}_m$, on the analyticity domain of the boundary
terms (the domain on which the bound in \textup{(B.3)} holds),
\begin{equation}\label{4.36:eq-frozen-sum-a}
\sum_{W,j}|\mathbf{C}_{W,j}(X)|\le 2c_{\mathrm{fr}}C_0^{\mathrm{fr}}e^{-\kappa d_m(X)/4},
\qquad c_{\mathrm{fr}}:=3^4\cdot48,
\end{equation}
the sum running over all arrests $W$ and all stages $j$ for which a frozen term with support $X$
at level $m$ occurs.
\end{proposition}

\begin{proof}
Fix $m\le k$ and $X\in\mathbf{D}_m$, and write $d_m=d_m(X)$. By the last clause of (B.3), each pair
$(W,j)$ contributes at most one term to the left side of \eqref{4.36:eq-frozen-sum-a}. Call an
arrest $W$ \emph{contributing} if some stage $j$ of $W$ carries a frozen term
$\mathbf{C}_{W,j}(X)$ with support $X$ at level $m$. The proof has three steps: a count of the
contributing arrests, a bound for the sum over the stages of one arrest, and their product.

\emph{Step~$(\alpha)$: the contributing arrests.} For a cube $C\in\pi_m$, its
$3^4$-neighbourhood is the family of the $3^4$ cubes of $\pi_m$ consisting of $C$ and the cubes of
$\pi_m$ whose closures meet the closure of $C$. Put
\begin{equation*}
\mathcal{C}(X):=\{C\in\pi_m:\ C\text{ lies in the }3^4\text{-neighbourhood of a cube of }X\}.
\end{equation*}
By (d), $X$ is a union of at most $48(1+d_m)$ cubes of $\pi_m$, and each of them has $3^4$ cubes in
its $3^4$-neighbourhood, so
\begin{equation}\label{4.36:eq-frozen-sum-1}
|\mathcal{C}(X)|\le 3^4\cdot48\,(1+d_m).
\end{equation}
Let $W$ be contributing, through a stage $j$. By (B.3), $h_W\le j\le m-1$ and $m\le k_W$, so
\begin{equation}\label{4.36:eq-frozen-sum-2}
m\in(h_W,k_W].
\end{equation}
By (B.0), $W$ meets $X$; pick a point $x\in W\cap X$. Since $m\le k_W$, (B.1) shows that $W$ is a
union of cubes of $\pi_m$, so $x$ lies in a cube $C\subset W$ of $\pi_m$; by (d), $x$ also lies in
a cube $C'\subset X$ of $\pi_m$. Two cubes of $\pi_m$ that meet lie in one another's
$3^4$-neighbourhood: as site sets they coincide, and as closed cubes they are either equal or
neighbours. Hence $C\in\mathcal{C}(X)$. For every contributing $W$ fix one such cube and call it
$C_W$; thus $C_W\subset W$ and $C_W\in\mathcal{C}(X)$.

The map $W\mapsto C_W$ is injective on contributing arrests. Indeed, if $W\neq W'$ are contributing
and $C_W=C_{W'}=C$, then $W$ and $W'$ are two arrests containing the cube $C\in\pi_m$, and both
satisfy \eqref{4.36:eq-frozen-sum-2}; this contradicts Lemma~\ref{4.36:lem-arrest-per-cube} (one
arrest per cube per level, proved from (B.2)), which says that among the arrests containing $C$
at most one has $m\in(h_W,k_W]$. With \eqref{4.36:eq-frozen-sum-1} this gives
\begin{equation}\label{4.36:eq-frozen-sum-3}
\#\{W\ \text{contributing at}\ (m,X)\}\le 3^4\cdot48\,(1+d_m).
\end{equation}

\emph{Step~$(\beta)$: the stages of one arrest.} Fix a contributing $W$. Its stages are the
$j\in[h_W,j_W-1]$, and by (a) there are $\mathfrak{N}_W\le R_{k_W}$ of them. By (b) the sequence
$(R_j)$ is non-increasing and $m\le k_W$, so $R_{k_W}\le R_m$; hence $W$ has at most $R_m$ stages.
By (B.3) only stages with $j\le m-1$ contribute at $(m,X)$; we treat $j=m-1$ and $j<m-1$
separately.

The stage $j=m-1$, if present, contributes by (B.3), and since $d_m\ge0$,
\begin{equation*}
|\mathbf{C}_{W,m-1}(X)|\le C_0^{\mathrm{fr}}e^{-\kappa d_m}\le C_0^{\mathrm{fr}}e^{-\kappa d_m/2}.
\end{equation*}

Let now $j+1<m$. By (B.3), $X\cap\Omega_m^{\natural}\neq\emptyset$ and
$X\cap Z_{j+1}\neq\emptyset$, the latter because $X$ meets $Z_{j+1}\setminus Z''_{j+1}$. Since
$j+1\le m-1$ and the sets $Z_i$ increase, (c) gives $Z_{j+1}\subset Z_{m-1}$. The set $X$ therefore
contains a point of $\Omega_m^{\natural}$ and a point of $Z_{m-1}$, and by (c) at $i=m-1$,
\begin{equation*}
\operatorname{diam}_\infty X\ \ge\ \operatorname{dist}_\infty(\Omega_m^{\natural},Z_{m-1})
\ \ge\ 8MR_m\ell_m .
\end{equation*}
The second inequality of (d) at $j=m$ then yields
$d_m\ge 8MR_m\ell_m/(M\ell_m)-2=8R_m-2$, that is $d_m/2\ge 4R_m-1$. Hence, by (B.3),
\begin{equation*}
|\mathbf{C}_{W,j}(X)|\le C_0^{\mathrm{fr}}e^{-\kappa d_m/2}e^{-\kappa d_m/2}
\le C_0^{\mathrm{fr}}e^{-\kappa(4R_m-1)}e^{-\kappa d_m/2}.
\end{equation*}
There are at most $R_m$ such stages, so their total is at most
$R_me^{-\kappa(4R_m-1)}C_0^{\mathrm{fr}}e^{-\kappa d_m/2}$. Now $R_m\ge1$ by (b), and
$\log t\le t-1$ for $t>0$; since $\kappa\ge1$ by (e) and $4R_m-1\ge0$,
\begin{equation*}
\log R_m\le R_m-1\le 4R_m-1\le\kappa(4R_m-1),
\end{equation*}
so $R_me^{-\kappa(4R_m-1)}\le1$ and the stages with $j+1<m$ contribute together at most
$C_0^{\mathrm{fr}}e^{-\kappa d_m/2}$. Adding the stage $j=m-1$,
\begin{equation}\label{4.36:eq-frozen-sum-4}
\sum_j|\mathbf{C}_{W,j}(X)|\le 2C_0^{\mathrm{fr}}e^{-\kappa d_m/2}
\qquad\text{for every contributing }W.
\end{equation}

\emph{Step~$(\gamma)$: the product.} Multiplying \eqref{4.36:eq-frozen-sum-3} by
\eqref{4.36:eq-frozen-sum-4} and using $c_{\mathrm{fr}}=3^4\cdot48$,
\begin{align*}
\sum_{W,j}|\mathbf{C}_{W,j}(X)|
&\le 3^4\cdot48\,(1+d_m)\cdot2C_0^{\mathrm{fr}}e^{-\kappa d_m/2}\\
&=2c_{\mathrm{fr}}C_0^{\mathrm{fr}}\bigl[(1+d_m)e^{-\kappa d_m/4}\bigr]e^{-\kappa d_m/4}.
\end{align*}
Since $\kappa\ge4$ by (e) and $d_m\ge0$,
\begin{equation*}
e^{\kappa d_m/4}\ge 1+\frac{\kappa d_m}{4}\ge 1+d_m ,
\end{equation*}
so the bracket is at most $1$, and \eqref{4.36:eq-frozen-sum-a} follows.
\end{proof}

\begin{remark}
Suppose that, in place of \textup{(B.0)}--\textup{(B.3)}, one takes as input the bound
\begin{equation*}
\sum_{W,j}|\mathbf{C}_{W,j}(X)|\le 2C_0^{\mathrm{fr}}e^{-\kappa d_m(X)/2}
\qquad\text{for every }m\le k\text{ and }X\in\mathbf{D}_m,
\end{equation*}
which exhibits the frozen terms at $(m,X)$ as one further term of the format of the boundary
pieces $\mathbf{B}^{(m)}$ of the effective action \cite[(2.23)]{B-CMP119}. This is the conclusion
of the proposition with $c_{\mathrm{fr}}$
replaced by $1$ and $\kappa/4$ by $\kappa/2$; the change affects values of constants only. Since
$c_{\mathrm{fr}}\ge1$ and $e^{-\kappa d_m(X)/2}\le e^{-\kappa d_m(X)/4}$, this bound implies
\eqref{4.36:eq-frozen-sum-a}, so every estimate that uses \eqref{4.36:eq-frozen-sum-a} with the
exponent $\kappa/4$ holds a fortiori under it.
\end{remark}

\begin{lemma}[A boundary polymer near a far block is long or recent]\label{4.36:lem-long-or-recent}
Use the regions $\Omega_n$, $\Gamma_n=\Omega_n\setminus\Omega_{n+1}$, $\Lambda_j$, $Z_j=\Lambda_j^{c}$ of a label, the blocks $\Delta(y)$, the polymer classes $\mathbf{D}_j$ with the distance function $d_j$, the block sets $\mathfrak{b}(\cdot)$ and $X^{+}$, and the radii $R_j$, all as in Notation~\ref{4.36:not-history-geometry}.

The label also carries its arrested regions: the regions $W$ on which the attempted large-field operation is stopped. For an arrested region $W$ we write:
\begin{itemize}
\item $k_W$ for the step at which $W$ is arrested;
\item $j_W$ for the stage that follows the last integration performed on $W$ before its arrest;
\item $N_0^W$ for the integer parameter of the arrest of $W$ at step $k_W$ that counts how many levels below $k_W$ the interval below begins, so that its left end is $k_0=k_W-N_0^W$;
\item $\mathfrak{E}(W)=(k_0,j_W)$ for the open integer interval attached to $W$.
\end{itemize}

Assume the following statements hold for the label.
\begin{itemize}
\item[(a)] Adjacent blocks differ by at most one scale. That is, if $\Delta(y_1)$ and $\Delta(y_2)$ are adjacent, with $y_1\in\Gamma_{s_1}$ and $y_2\in\Gamma_{s_2}$, then $|s_1-s_2|\le 1$.
\item[(b)] Let $j$ be a scale and let $b<a$. If $X$ is a connected union of scale-$j$ $M$-cubes that meets both $\Omega_a$ and $\Omega_b^{c}$, then $d_j(X)\ge L^{a-1-j}-2$.
\item[(c)] For every $i$,
\begin{equation*}
\Lambda_i\supseteq\Omega_{i+1}\setminus\bigcup\{W:\ i\in\mathfrak{E}(W)\}.
\end{equation*}
\item[(d)] Every arrested region $W$ satisfies $W\subset\Omega^{c}_{k_W+1}$.
\item[(e)] For every $j$, $R_{j+1}\le R_j$.
\item[(f)] For every $j$, $R_j/R_{j+1}\in\{1,L\}$, and the value $L$ is never taken for two consecutive $j$.
\item[(g)] For every arrested region $W$, $N_0^W\le 2+\log_L R_{k_W}$.
\end{itemize}
All of (a)--(g) are among the inputs listed with the geometry of Notation~\ref{4.36:not-history-geometry}: (a) and (b) bound the scales of adjacent blocks and the lengths of polymers, (c) and (d) are the inclusion properties of the label's regions and of its arrested regions, and (e)--(g) are displayed properties of the radii $R_j$ and of the parameter $N_0^W$.

Let $y\in\Gamma_{n'}$, let $i\le n'-4$, and let $X\in \mathbf{D}_i$ with $y\in\mathfrak{b}(X^{+})$ and $X\cap Z_i\neq\emptyset$. Then either $d_i(X)\ge L^{n'-3-i}-2$, or
\begin{equation*}
n'-i\le\tau_{n'}:=6+\log_L R_{n'}.
\end{equation*}
\end{lemma}

\begin{proof}
\emph{The polymer reaches $\Omega_{n'-2}$.} Since $y\in\mathfrak{b}(X^{+})$, the block $\Delta(y)$ lies within adjacency at most $2$ of a block $\Delta(y'')$ that meets $X$. Let $s$ be the scale of $y''$, so $y''\in\Gamma_s$. Apply (a) along each of the at most two adjacency steps from $\Delta(y)$ to $\Delta(y'')$. Since $y\in\Gamma_{n'}$, this gives $s\in[n'-2,n'+2]$.

The block $\Delta(y'')$ of a point $y''\in\Gamma_s$ satisfies $\Delta(y'')\subset\Omega_s$. The regions $\Omega_n$ of the label are nested decreasingly in $n$ (Notation~\ref{4.36:not-history-geometry}), and $s\ge n'-2$, so $\Omega_s\subset\Omega_{n'-2}$. Since $\Delta(y'')$ meets $X$, it follows that $X$ meets $\Omega_{n'-2}$.

\emph{Splitting $Z_i$.} Take complements in (c):
\begin{equation*}
Z_i=\Lambda_i^{c}\subset\Omega^{c}_{i+1}\cup\bigcup\{W:\ i\in\mathfrak{E}(W)\}.
\end{equation*}
Since $X$ meets $Z_i$, it meets either $\Omega^{c}_{i+1}$ or one of the regions $W$ with $i\in\mathfrak{E}(W)$.

\emph{Case 1: $X$ meets $\Omega^{c}_{i+1}$.} From $i\le n'-4$ we get $i+1\le n'-3$. By the nesting of the $\Omega_n$, this gives $\Omega^{c}_{i+1}\subset\Omega^{c}_{n'-3}$. So $X$ meets both $\Omega_{n'-2}$ and $\Omega^{c}_{n'-3}$. Now $X\in \mathbf{D}_i$ is a connected union of scale-$i$ $M$-cubes. Apply (b) with $j=i$, $a=n'-2$ and $b=n'-3$, which satisfies $b<a$. It gives
\begin{equation*}
d_i(X)\ge L^{(n'-2)-1-i}-2=L^{n'-3-i}-2 .
\end{equation*}

\emph{Case 2: $X$ meets a region $W$ with $i\in\mathfrak{E}(W)=(k_0,j_W)$.} Since the interval is open and $k_0=k_W-N_0^W$, we have
\begin{equation*}
i\ge k_0+1=k_W-N_0^W+1 .
\end{equation*}
By (d), $W\subset\Omega^{c}_{k_W+1}$, so $X$ meets $\Omega^{c}_{k_W+1}$.

Suppose first that $k_W+1\le n'-3$. Then $\Omega^{c}_{k_W+1}\subset\Omega^{c}_{n'-3}$, so $X$ meets $\Omega_{n'-2}$ and $\Omega^{c}_{n'-3}$. The argument of Case 1, that is (b) with $(a,b)=(n'-2,n'-3)$, gives $d_i(X)\ge L^{n'-3-i}-2$.

Otherwise $k_W\ge n'-3$, that is, $n'\le k_W+3$. Combining this with the lower bound on $i$ gives
\begin{equation*}
n'-i\le n'-k_W+N_0^W-1\le N_0^W+2 .
\end{equation*}
By (g) at step $k_W$, $N_0^W+2\le 4+\log_L R_{k_W}$.

It remains to show $R_{k_W}\le L^2R_{n'}$. If $n'\le k_W$, then (e), applied repeatedly, gives $R_{k_W}\le R_{n'}\le L^2R_{n'}$. Otherwise $k_W<n'\le k_W+3$, and
\begin{equation*}
\frac{R_{k_W}}{R_{n'}}=\prod_{t=k_W}^{n'-1}\frac{R_t}{R_{t+1}}
\end{equation*}
is a product of $n'-k_W\le 3$ consecutive ratios. By (f) each ratio is $1$ or $L$, and no two consecutive ratios both equal $L$. Hence at most two of these ratios equal $L$, and $R_{k_W}\le L^2R_{n'}$.

Therefore $\log_L R_{k_W}\le 2+\log_L R_{n'}$, and
\begin{equation*}
n'-i\le 4+\log_L R_{k_W}\le 6+\log_L R_{n'}=\tau_{n'} .
\end{equation*}

In every case one of the two alternatives of the lemma holds. The argument uses no information about the number of arrested regions of the label or about how they are nested.
\end{proof}

\begin{theorem}[Second-derivative bound for the boundary terms at $SU(N)$]\label{4.36:thm-boundary-bound}
Let $G=SU(N)$ with the conventions of Notation~\ref{4.36:not-sun-conventions}. Use the geometry of a
history of Notation~\ref{4.36:not-history-geometry} and the probing direction of
Notation~\ref{4.36:not-probing-direction}. Assume the standing hypotheses of
Definition~\ref{4.36:def-near-part}. Among the data fixed there are the following:
\begin{itemize}
\item a level $k\le K$;
\item a clean history $h$, so that the distinguished plaquette lies at distance at least $R_*$ from
the complement of $\Lambda_k(h)$;
\item the off-link configuration $V_{\mathrm{off}}$ and the real, windowed configuration
$V=(u,V_{\mathrm{off}})$;
\item a direction $\hat{T}\in\mathfrak{g}$ with $|\hat{T}|=1$.
\end{itemize}
Moreover $x_n\ge x^{(3)}_{\mathrm{th}}$ for all $n\le K$, as in
\eqref{4.36:eq-coupling-ceiling-a}.

We use the following notation.
\begin{itemize}
\item $\Delta(y)$ is the block labelled by $y$, and $\ell_n$ is the side of a block of scale $n$, so
that $\ell_{n+1}=L\,\ell_n$; the cubes of $\pi_n$ have side $M\ell_n$.
\item $e_y:=e^{-\delta_0 d(y,b_1)/8}$, where $b_1$ is the marked bond.
\item For a bond set $\mathsf{N}$, $\mathfrak{b}(\mathsf{N})$ is the set of blocks meeting
$\mathsf{N}$, and $\mathsf{N}^{+}$ is the $2$-neighbourhood of $\mathsf{N}$; for a polymer $X$ with
set of bonds $\mathsf{N}$ we write $X^{+}:=\mathsf{N}^{+}$.
\item $\hat{e}_{\mathsf{N}}:=\sum_{y\in\mathfrak{b}(\mathsf{N})}e_y^{1/2}$.
\item For a function $\phi$ of $t$, $|\phi|_0$ denotes the modulus of $\phi$ at $t=0$.
\end{itemize}
$F^B$ is the sum of all boundary terms of Definition~\ref{4.36:def-near-part}. These are the
native terms $\mathbf{B}^{(i)}(X)$ and the terms $\mathbf{C}_{W,j}(X)$ frozen on arrests $W$.

Assume in addition the following statements.
\begin{enumerate}
\item[(a)] \emph{Native terms.} For every $i\le k$ and every $X\in\mathbf{D}_i$ with
$X\cap Z_i\neq\emptyset$, the native term $\mathbf{B}^{(i)}(X)$ is gauge invariant and depends on
the configuration on $X$ only. It is analytic on the domain $\widetilde{U}_i(X)$ of
\cite[(2.42)]{B-CMP119} and satisfies there
\[
|\mathbf{B}^{(i)}(X)|<B_0e^{-\kappa d_i(X)},
\]
uniformly in $w\in\operatorname{supp}m_h$.
\item[(b)] \emph{Frozen terms and arrests.} A term $\mathbf{C}_{W,j}(X)$ frozen on the arrest $W$
by the stage integrating level $j$ has support $X\in\mathbf{D}_m$ with $j+1\le m$ and
$X\cap Z_{j+1}\ne\emptyset$. It is gauge invariant and depends on the configuration on $X$ only.
It is analytic on its domain and satisfies there
\[
|\mathbf{C}_{W,j}(X)|\le C_0^{\mathrm{fr}}e^{-\kappa d_m(X)},
\]
as in \cite[(1.63), (1.70)]{B-CMP122a}. The hypotheses of
Proposition~\ref{4.36:prop-frozen-sum} and of Lemma~\ref{4.36:lem-long-or-recent} hold; in
particular every $R_n$ is a power of $L$.
\item[(c)] \emph{Domains.} Let $(i,X)$ be as in (a) or (b), and let $\mathsf{N}$ be the set of
bonds of $X$. Write $\chi_{\mathsf{N}}$ for the cut-off function attached to the bond set
$\mathsf{N}$ in the localisation bound \eqref{4.36:eq-boundary-bound-1} below, and $\eta$ for the
spacing of the lattice carrying $U_k(V)$, so that $e^{i\eta\mathfrak{h}}U$ is the configuration $U$
shifted bondwise by a $\mathfrak{g}^{c}$-valued potential $\mathfrak{h}$; $\mathcal{H}_\zeta$ is the
response field of Notation~\ref{4.36:not-probing-direction}. Here and below $R>0$ is the probing
radius: the radius in the probing parameter $\zeta$ for which the following statements are
assumed, and the same radius with which the floor $x_B$ of
\eqref{4.36:eq-boundary-absorption-b} is computed; it is not the ball radius $R$ of the glossary.
We call the \emph{domain of $(i,X)$} the set of configurations on $X$ that lie in
$\widetilde{U}_i(X)$ and in the domain of \cite[(1.70)]{B-CMP122a} of each frozen term with
support $X$. The following three statements hold.
\begin{itemize}
\item For $|\zeta|\le 2R$ the complexified configuration
$e^{i\eta\chi_{\mathsf{N}}\mathcal{H}_\zeta}U_k(V)$, restricted to $X$, lies in
$\widetilde{U}_i(X)$ at half radius, and in the domain of \cite[(1.70)]{B-CMP122a} of each frozen
term with support $X$ (off the arrests by the membership bound in the norm $\|\cdot\|_{\star}$
with radius constant $c_{\star}$, which places perturbations small in $\|\cdot\|_{\star}$ in the
analyticity domains and whose threshold is the floor \eqref{4.36:eq-threshold-floors-d} of
Lemma~\ref{4.36:lem-threshold-floors}; on the arrests by the domain statement for arrested
regions).
\item The localisation bound of the response field holds:
\begin{equation}\label{4.36:eq-boundary-bound-1}
\|\chi_{\mathsf{N}}\mathcal{H}_\zeta\|_{\star}\le
\frac{C_\chi' C_H\,|\zeta|\,\min\{1,\hat{e}_{\mathsf{N}^{+}}\}}{\delta_k}.
\end{equation}
Here $C_\chi'$ is the constant of this localisation bound; $\delta_k>0$ is the level-$k$
small-field parameter of Ba{\l}aban's analysis, which by Lemma~\ref{4.36:lem-threshold-floors}
satisfies $\delta_k\le A_1g_k(\log x_k)^{p_0+r}$ when $x_k\ge e$; and $C_H$ is the constant of
the response-field bound, which bounds $\partial_\zeta\mathcal{H}_\zeta$ on the block $\Delta(y)$ by
$C_He_y|T'|$.
\item (Real identification.) For every function $\Phi(\cdot,w)$ of the configuration on $X$,
depending on $w\in\operatorname{supp}m_h$ as a parameter, the function
$\zeta\mapsto\Phi(e^{i\eta\chi_{\mathsf{N}}\mathcal{H}_\zeta}U_k(V),w)$ coincides at real
$\zeta=t$, $|t|\le 2R$, with $t\mapsto\Phi(U_k(V^{t,\hat{T}}),w)$.
\end{itemize}
\item[(d)] \emph{Second-derivative estimate from the oscillation.} Let $(i,X)$ be as in (c). Let
$\Phi(\cdot,w)$ be a gauge-invariant function of the configuration on $X$, depending on
$w\in\operatorname{supp}m_h$ as a parameter, that is analytic on the domain of $(i,X)$ and has
oscillation at most $S$ there. Then, for the configurations of (c),
\[
|\partial_t^2\Phi(U_k(V^{t,\hat{T}}),w)|_0\le S\,\hat{e}_{X^{+}}/R^2 .
\]
\item[(e)] \emph{Scales of adjacent blocks.} Let $y\in\Gamma_{n'}$. If a block $\Delta(y'')$ is
within adjacency at most $2$ of $\Delta(y)$, then $y''\in\Gamma_s$ for some
$s\in[n'-2,n'+2]$.
\item[(f)] \emph{Counting.} Every $X\in\mathbf{D}_i$ is a union of cubes of $\pi_i$. For $a>c_d$,
where $c_d$ is the lattice-animal constant of the count below, the animal sum satisfies
\[
\mathfrak{a}(a):=\sup_{c\in\pi_i}\sum_{X\in\mathbf{D}_i,\,X\supset c}e^{-a d_i(X)}
\le\mathfrak{s}(a):=\sum_{l\ge0}e^{-(a-c_d)l}.
\]
For $i,n'\le k$, $y\in\Gamma_{n'}$ and $l\ge0$,
\[
\#\{X\in\mathbf{D}_i:\ y\in\mathfrak{b}(X^{+}),\ d_i(X)=l\}\le 7^4L^{4(n'+2-i)_+}e^{c_dl}.
\]
\item[(g)] \emph{Numerical conditions.} The conditions on $L$, $\kappa$, $c_d$ and $\delta_0M$
among the standing hypotheses of Definition~\ref{4.36:def-near-part} hold; of them the proof,
including the summation over blocks carried out in the proof of the boundary-term bound in two
distance classes, uses only $L\ge3$, $\kappa\ge8$, $\kappa/8-c_d\ge1$ and $\delta_0M\ge16$.
\item[(h)] \emph{Summation over blocks.} The following hold:
\begin{itemize}
\item the sharpened bounds on the coupling flow of Notation~\ref{4.36:not-sun-conventions};
\item the block-summation bound on the sums $\sigma_n^{(\alpha)}$ at
$\alpha\in\{1/32,1/16\}$;
\item the auxiliary parameter $r$ of $\mathfrak{p}$ satisfies $r\ge2$.
\end{itemize}
\end{enumerate}
Then, for every $w\in\operatorname{supp}m_h$, the function $t\mapsto F^B(V^{t,\hat{T}},w)$ is
real-analytic near $t=0$, and
\begin{equation}\label{4.36:eq-boundary-bound-a}
\begin{split}
\bigl|\partial_t^2F^B(V^{t,\hat{T}},w)\bigr|_0
&\le\frac{C_{\mathbf{B}}}{R^2}\,(\log x_k)^{4r}
\Bigl[e^{-\delta_0R_*/128}+e^{-\kappa R_*/(16M)}\Bigr]\\
&\le1 .
\end{split}
\end{equation}
Here $C_{\mathbf{B}}$ is the collection constant of the boundary-term bound in two distance classes,
the statement to which the display \eqref{4.36:eq-boundary-absorption-b} belongs; the bound
contains no factor of $\mathsf{n}=(N-1)^{1/2}$. The last inequality in
\eqref{4.36:eq-boundary-bound-a} holds whenever $x_k\ge x_B$,
with $x_B$ the floor of \eqref{4.36:eq-boundary-absorption-b}. This floor holds for every $k\le K$
when $\gamma\le(x_B+c_5)^{-1/2}$.
\end{theorem}

\begin{proof}[Proof of Theorem~\ref{4.36:thm-boundary-bound}, first part]
Fix $w\in\operatorname{supp}m_h$.

\emph{Step 1: filing of the frozen terms and the oscillation per polymer.}
Let $\mathbf{C}_{W,j}(X)$ be a frozen term with $X\in\mathbf{D}_m$.

By hypothesis (b), $X$ meets $Z_{j+1}$ and $j+1\le m$. The regions $Z_j=\Lambda_j^{c}$ increase
with $j$ (Notation~\ref{4.36:not-history-geometry}), so $Z_{j+1}\subset Z_m$ and
$X\cap Z_m\neq\emptyset$. We file the frozen term at its level $m$, that is, at the pair $(m,X)$.

Consequently the index set
\[
\{(i,X):\ i\le k,\ X\in\mathbf{D}_i,\ X\cap Z_i\neq\emptyset\}
\]
carries native and frozen terms alike. For $(i,X)$ in it we put
\[
\Psi_{i,X}:=\mathbf{B}^{(i)}(X)+\sum_{W,j}\mathbf{C}_{W,j}(X).
\]
The sum runs over the frozen terms filed at $(i,X)$, and either part may be absent. By the
definition of $F^B$ in Definition~\ref{4.36:def-near-part}, $F^B$ is the sum of the
$\Psi_{i,X}$ over this index set.

On the domain of $(i,X)$ defined in hypothesis (c), two bounds hold.
\begin{itemize}
\item Hypothesis (a) gives $|\mathbf{B}^{(i)}(X)|<B_0e^{-\kappa d_i(X)}$.
\item Proposition~\ref{4.36:prop-frozen-sum}, applied at $m=i$, gives
\[
\sum_{W,j}|\mathbf{C}_{W,j}(X)|\le 2c_{\mathrm{fr}}C_0^{\mathrm{fr}}e^{-\kappa d_i(X)/4},
\]
with $c_{\mathrm{fr}}=3^4\cdot48$.
\end{itemize}
For a function $\Phi$ on a set we write $\operatorname{osc}\Phi:=\sup|\Phi(a)-\Phi(b)|$, the
supremum over pairs of points $a,b$ of the set, for its oscillation. The oscillation is at most
twice the supremum of $|\Phi|$, and $e^{-\kappa d_i(X)}\le e^{-\kappa d_i(X)/4}$. Hence the
oscillation of $\Psi_{i,X}$ on the domain of $(i,X)$ is bounded as follows:
\begin{equation}\label{4.36:eq-boundary-bound-b}
\begin{aligned}
\operatorname{osc}\Psi_{i,X}
&\le 2\bigl(B_0e^{-\kappa d_i(X)}+2c_{\mathrm{fr}}C_0^{\mathrm{fr}}e^{-\kappa d_i(X)/4}\bigr)
\le 2B_0^{\sharp}e^{-\kappa d_i(X)/4},\\
B_0^{\sharp}&:=B_0+2c_{\mathrm{fr}}C_0^{\mathrm{fr}}=B_0+7776\,C_0^{\mathrm{fr}} .
\end{aligned}
\end{equation}
The bound is additive. It therefore holds also when a native term and a collection of frozen terms
share one pair $(i,X)$.

\emph{Step 2: analyticity and the second derivative per polymer.}
Fix $(i,X)$ in the index set and let $\mathsf{N}$ be the set of bonds of $X$.

By hypotheses (a) and (b), $\Psi_{i,X}$ is gauge invariant and depends on the configuration on $X$
only. Its native part is analytic on $\widetilde{U}_i(X)$ by \cite[(2.42)]{B-CMP119}. Its frozen
part is analytic on the domain of \cite[(1.70)]{B-CMP122a}.

By the first statement of hypothesis (c), for $|\zeta|\le 2R$ the complexified configuration
$e^{i\eta\chi_{\mathsf{N}}\mathcal{H}_\zeta}U_k(V)$, restricted to $X$, lies in
$\widetilde{U}_i(X)$ at half radius and in the domain of \cite[(1.70)]{B-CMP122a} of each frozen
term filed at $(i,X)$. Off the arrests this membership comes from the membership bound in
$\|\cdot\|_{\star}$; on the arrests it comes from the domain statement for arrested regions. In
both cases the size of the perturbation is controlled by \eqref{4.36:eq-boundary-bound-1}. The
response-field bound enters \eqref{4.36:eq-boundary-bound-1} through the bound
$C_He_y|T'|=C_He_y$ on $\partial_\zeta\mathcal{H}_\zeta$ on the block $\Delta(y)$, because
$|T'|=1$ (Notation~\ref{4.36:not-probing-direction}), so no factor of the norm of a generator
enters.

Hence $\Psi_{i,X}$ is analytic on the domain of $(i,X)$, its oscillation there obeys
\eqref{4.36:eq-boundary-bound-b}, and the configurations of hypothesis (c) lie in that domain for
$|\zeta|\le 2R$, so that the function of $\zeta$ obtained by inserting them is analytic for
$|\zeta|\le 2R$. By the real identification in hypothesis (c), at real $\zeta=t$
with $|t|\le2R$ it coincides with $t\mapsto\Psi_{i,X}(U_k(V^{t,\hat{T}}),w)$.

Apply hypothesis (d) to $\Phi=\Psi_{i,X}$, which is analytic on the domain of $(i,X)$ with
oscillation there bounded by \eqref{4.36:eq-boundary-bound-b}, with
$S:=2B_0^{\sharp}e^{-\kappa d_i(X)/4}$. This gives
\begin{equation}\label{4.36:eq-boundary-bound-2}
\bigl|\partial_t^2\Psi_{i,X}(U_k(V^{t,\hat{T}}),w)\bigr|_0
\le\frac{2B_0^{\sharp}e^{-\kappa d_i(X)/4}\,\hat{e}_{X^{+}}}{R^2},
\end{equation}
uniformly in $w$. The uniformity holds because the bound of \cite[(2.42)]{B-CMP119} in hypothesis
(a) is uniform on $\operatorname{supp}m_h$.

There are finitely many pairs $(i,X)$. Hence $t\mapsto F^B(V^{t,\hat{T}},w)$ is a finite sum of
functions that are analytic in $t$ for $|t|<2R$. It is therefore real-analytic near $t=0$, which is
the first assertion of the theorem.

\emph{Step 3: summation over polymers.}
Sum \eqref{4.36:eq-boundary-bound-2} over the finitely many pairs $(i,X)$ and insert
$\hat{e}_{X^{+}}=\sum_{y\in\mathfrak{b}(X^{+})}e_y^{1/2}$. Then exchange the order of the
non-negative sums, filing each block $y$ of $\mathfrak{b}(X^{+})$ under the set $\Gamma_{n'}$,
$n'\le k$, that contains it. This gives the first inequality of \eqref{4.36:eq-boundary-bound-c}
below, in which, for $n'\le k$ and $y\in\Gamma_{n'}$, the count $\mathfrak{S}(y)$ is
\[
\mathfrak{S}(y):=\sum_{i\le k}\ \sum_{\substack{X\in\mathbf{D}_i:\ y\in\mathfrak{b}(X^{+}),\\
X\cap Z_i\neq\emptyset}}e^{-\kappa d_i(X)/4}.
\]

\emph{Step 4: the count $\mathfrak{S}(y)$.}
Fix $n'\le k$ and $y\in\Gamma_{n'}$. We split the levels $i\le k$ into four ranges. The level
$i=n'-4$ belongs to both the second range and the last two; counting it twice only enlarges the
bound. We first record a consequence of hypothesis (g) and one of hypothesis (e).

First, $\kappa/4-c_d\ge\kappa/8-c_d\ge1$. Hence
\begin{equation}\label{4.36:eq-boundary-bound-4}
\mathfrak{s}(\kappa/4)\le\sum_{l\ge0}e^{-l}=\frac{e}{e-1}\le2 .
\end{equation}

Second, by the definition of $X^{+}$, for $X\in\mathbf{D}_i$ with $y\in\mathfrak{b}(X^{+})$ the
block $\Delta(y)$ lies within adjacency at most $2$ of a block $\Delta(y'')$ meeting $X$. By
hypothesis (e), $y''$ and the blocks of this adjacency chain have scales in $[n'-2,n'+2]$.

Summing the count of hypothesis (f) against $e^{-\kappa l/4}$ over $l\ge0$ gives, for every level
$i$,
\begin{equation}\label{4.36:eq-boundary-bound-5}
\sum_{X\in\mathbf{D}_i:\ y\in\mathfrak{b}(X^{+})}e^{-\kappa d_i(X)/4}
\le 7^4L^{4(n'+2-i)_+}\,\mathfrak{s}(\kappa/4).
\end{equation}

\emph{Range 1: $i\ge n'+5$.}
The blocks of the adjacency chain from $\Delta(y)$ to the block of $X$ that it reaches have scales
in $[n'-2,n'+2]$. They therefore lie in a region whose extent per axis is at most
\[
\ell_{n'}+\ell_{n'+1}+\ell_{n'+2}\le3\ell_{n'+2}\le3L^{-3}\ell_i<\ell_i\le M\ell_i ,
\]
where $3L^{-3}<1$ because $L^3\ge8>3$. This region meets at most $2$ cubes of $\pi_i$ per axis,
hence at most $2^4$ cubes of $\pi_i$ in all. By hypothesis (f), $X$ is a union of cubes of
$\pi_i$, so it contains one of these cubes. Using $\mathfrak{a}(\kappa/4)\le\mathfrak{s}(\kappa/4)$
from hypothesis (f) and \eqref{4.36:eq-boundary-bound-4}, we get per level
\[
\sum_X e^{-\kappa d_i(X)/4}\le2^4\,\mathfrak{a}(\kappa/4)\le2^5 .
\]
There are at most $k-n'$ levels in this range, so Range 1 contributes at most $2^5(k-n')$. This is
an upper bound; whether polymers of boundary type occur in this range is immaterial.

\emph{Range 2: $|i-n'|\le4$.}
This range has nine levels. For each, $(n'+2-i)_+\le6$. By \eqref{4.36:eq-boundary-bound-5} and
\eqref{4.36:eq-boundary-bound-4}, each level contributes at most $2\cdot7^4L^{24}$. Range 2
therefore contributes at most
\[
C_{\mathrm{mid}}:=18\cdot7^4L^{24}.
\]

\emph{Range 3: $i\le n'-4$, long polymers.}
Put $v:=n'-3-i\ge1$. The polymers counted here are those with $d_i(X)\ge L^v-2$.

For such $X$ with $d_i(X)=l$ we have $L^v\le l+2$. Since $4(n'+2-i)=4v+20$, the count of hypothesis
(f) gives
\[
\#\{X:\ d_i(X)=l\}\le7^4L^{20}(L^v)^4e^{c_dl}\le7^4L^{20}(l+2)^4e^{c_dl}.
\]
Moreover, since $l\ge L^v-2$,
\[
e^{-\kappa l/4}=e^{-\kappa l/8}e^{-\kappa l/8}\le e^{-\kappa l/8}e^{-\kappa(L^v-2)/8}.
\]
Hence level $i$ contributes at most $C_{\mathrm{long}}\,e^{-\kappa(L^v-2)/8}$, where
\[
C_{\mathrm{long}}:=7^4L^{20}\sum_{l\ge0}(l+2)^4e^{-(\kappa/8-c_d)l}.
\]

This constant is finite. Indeed $\kappa/8-c_d\ge1$ by hypothesis (g). By the elementary
inequalities \eqref{4.36:eq-sun-conventions-a}, in the form $e^{x}\ge x^4/4!$ for $x\ge0$ taken at
$x=(l+2)/2$, we have
\[
(l+2)^4\le4!\cdot2^4\,e^{(l+2)/2}=384\,e^{(l+2)/2}.
\]
Hence
\begin{equation}\label{4.36:eq-boundary-bound-6}
C_{\mathrm{long}}\le7^4L^{20}\cdot384\,e\sum_{l\ge0}e^{-l/2}\le7^4L^{20}\cdot1152\,e ,
\end{equation}
because $\sum_{l\ge0}e^{-l/2}=(1-e^{-1/2})^{-1}\le3$.

The levels of this range correspond one-to-one to the integers $v\ge1$. By the elementary
inequalities \eqref{4.36:eq-sun-conventions-a}, Bernoulli's inequality
$L^v\ge1+v(L-1)$ gives, for $L\ge3$ and $v\ge1$,
\[
L^v-2\ge1+v(L-1)-2\ge2v-1\ge v .
\]
Together with $\kappa\ge8$ this yields
\[
\sum_{v\ge1}e^{-\kappa(L^v-2)/8}\le\sum_{v\ge1}e^{-v}=\frac{1}{e-1}\le1 .
\]
Range 3 therefore contributes at most $C_{\mathrm{long}}$.

\emph{Range 4: $i\le n'-4$, short polymers.}
Here $d_i(X)<L^{n'-3-i}-2$. The polymer $X$ satisfies $y\in\mathfrak{b}(X^{+})$ and
$X\cap Z_i\neq\emptyset$. By Lemma~\ref{4.36:lem-long-or-recent} it is therefore recent, and
\[
4\le n'-i\le\tau_{n'}:=6+\log_LR_{n'} .
\]
Here $\tau_{n'}$ is an integer, because $R_{n'}$ is a power of $L$ by hypothesis (b).

By \eqref{4.36:eq-boundary-bound-5} and \eqref{4.36:eq-boundary-bound-4}, each such level
contributes at most
\[
7^4L^{4(n'+2-i)}\mathfrak{s}(\kappa/4)\le2\cdot7^4L^{8+4(n'-i)} .
\]
Summing the geometric series over $4\le n'-i\le\tau_{n'}$, we use $1/(1-L^{-4})\le2$ (because
$L^4\ge2$) and $L^{4\tau_{n'}}=L^{24}R_{n'}^4$. Range 4 then contributes at most
\[
\frac{2\cdot7^4L^8L^{4\tau_{n'}}}{1-L^{-4}}\le4\cdot7^4L^{32}R_{n'}^4 .
\]

\emph{Conclusion of the count.}
Adding the four ranges gives, for every $n'\le k$ and $y\in\Gamma_{n'}$,
\[
\mathfrak{S}(y)\le2^5(k-n')+C_{\mathrm{mid}}+C_{\mathrm{long}}+4\cdot7^4L^{32}R_{n'}^4 .
\]
Together with the summation of Step~3 this gives
\begin{equation}\label{4.36:eq-boundary-bound-c}
\begin{gathered}
\bigl|\partial_t^2F^B\bigr|_0\le\frac{2B_0^{\sharp}}{R^2}\sum_{n'\le k}\ \sum_{y\in\Gamma_{n'}}
e_y^{1/2}\,\mathfrak{S}(y),\\
\mathfrak{S}(y)\le C_{\mathfrak{B}}\bigl[1+(k-n')+R_{n'}^4\bigr]\qquad
(n'\le k,\ y\in\Gamma_{n'}),\\
C_{\mathfrak{B}}:=\max\bigl\{2^5,\ C_{\mathrm{mid}}+C_{\mathrm{long}},\ 4\cdot7^4L^{32}\bigr\}.
\end{gathered}
\end{equation}
The constant $C_{\mathfrak{B}}$ depends only on $L$, $\kappa$ and $c_d$, and not on $N$.

It remains to sum \eqref{4.36:eq-boundary-bound-c} over the blocks $y$, using hypotheses (g) and
(h). This leads to the collection constant $C_{\mathbf{B}}$, to the floor
\eqref{4.36:eq-boundary-absorption-b} defining $x_B$, and to the absorption giving
\eqref{4.36:eq-boundary-bound-a}. These steps are carried out in the proof of the boundary-term
bound in two distance classes, which completes the proof.
\end{proof}

\begin{lemma}[The boundary-term bound in two distance classes]\label{4.36:prf-boundary-absorption}
Assume the standing hypotheses of Definition~\ref{4.36:def-near-part}, in the notation of
Notation~\ref{4.36:not-sun-conventions}, Notation~\ref{4.36:not-history-geometry} and
Notation~\ref{4.36:not-probing-direction}. Thus $h$ is a clean label, that is,
$\operatorname{dist}_k(p,\Lambda_k(h)^{c})\ge R_*$ with $R_*=20L^2M\max\{R_k,R_{k+1}\}$; $b_1\subset\partial p$
is the marked bond; and for a block $y$ we write $e_y=e^{-\delta_0 d(y,b_1)/8}$.
Fix $w\in\operatorname{supp}m_h$ and assume the following.
\begin{enumerate}
\item[(a)] The two bounds \eqref{4.36:eq-boundary-bound-c} of the proof of
Theorem~\ref{4.36:thm-boundary-bound}: with
\[
\mathfrak{S}(y):=\sum_{i\le k}\ \sum_{X\in\mathbf{D}_i:\ y\in\mathfrak{b}(X^{+}),\ X\cap Z_i\neq\emptyset}
e^{-\kappa d_i(X)/4},
\]
one has
\begin{gather*}
|\partial_t^2F^B(V^{t,\hat T},w)|_0\le\frac{2B_0^{\sharp}}{R^2}\sum_{n'\le k}\ \sum_{y\in\Gamma_{n'}}
e_y^{1/2}\mathfrak{S}(y),\\
\mathfrak{S}(y)\le C_{\mathfrak{B}}[1+(k-n')+R_{n'}^4]\qquad(y\in\Gamma_{n'}).
\end{gather*}
\item[(b)] The block-sum estimate at $\alpha\in\{1/32,1/16\}$ (both in the range $[1/288,1/8]$ which it
covers), with its sums $\sigma_n^{(\alpha)}$ and its constant $C_\sigma$: for every $n\le k$,
\[
\sum_{y\in\Gamma_n}e^{-\alpha\delta_0d(y,b_1)}\le\sigma_n^{(\alpha)}\le C_\sigma e^{-(k-n)} .
\]
\item[(c)] The location of near blocks: if $y\in\Gamma_{n'}$ and $d(b_1,y)<R_*/4$, then $y\in\Gamma_k$ and
$|y-b_1|<R_*/4$ in scale-$k$ units.
\item[(d)] The lower bound for $d_i$ by the sup-diameter, in the form used below: if $i=k-u\le k$
and $X\in\mathbf{D}_i$ has $\operatorname{diam}_\infty X\ge 2R_*/3$ in scale-$k$ units, then
$d_i(X)\ge\tfrac23R_*L^u/M-2$.
\item[(e)] The polymer count, built on the animal bound $\mathfrak{a}(\theta)\le\mathfrak{s}(\theta)$ with
$\mathfrak{s}(\theta)=\sum_{\ell\ge0}e^{-(\theta-c_d)\ell}$: for $y\in\Gamma_{n'}$, $i\le k$ and $\theta>c_d$,
\[
\sum_{X\in\mathbf{D}_i:\ y\in\mathfrak{b}(X^{+}),\ X\cap Z_i\neq\emptyset}e^{-\theta d_i(X)}
\le 7^4L^{4(n'+2-i)_+}\,\mathfrak{s}(\theta).
\]
\item[(f)] The radii and the flow: property (c) of the list $\check{R}$, namely $R_j<L(\log x_j)^r$
whenever $x_j\ge e$; the lower bound $R_k\ge(\log x_k)^r$; $x_j\ge e$ for all $j\le k$; and, for $n'\le k$,
$x_{n'}\le(1+\mathfrak{f}(k-n'))x_k$, by the flow bounds \eqref{4.36:eq-sun-conventions-b}, with
$\mathfrak{f}=1$ or, provided $x_k\ge1$, with $\mathfrak{f}=\bar B$.
\item[(g)] The parameter conditions $L\ge2$, $M\ge1$, $\delta_0M\ge16$, $\kappa\ge10$,
$\kappa/8-c_d\ge1$ and $r\ge2$.
\end{enumerate}
Put
\begin{align*}
C_A&:=\sum_{u\ge0}e^{-u}\bigl[1+u+L^4\,2^{4r-1}\bigl(1+(\log(1+\mathfrak{f}u))^{4r}\bigr)\bigr],\\
C_{\mathbf{B}}&:=2B_0^{\sharp}C_\sigma\max\{C_{\mathfrak{B}}C_A,\,4\cdot7^4L^8\},\\
C_{R^{-2}}&:=\Bigl(\frac{4C_\chi'C_H}{c_\star A_1}\Bigr)^2+\Bigl(\frac{16C_{\mathrm{loc}}}{A_0}\Bigr)^2
+4r_0^{-2},\\
c_{\mathrm{exp}}&:=\tfrac52L^2,\qquad \mathfrak{A}^{(4)}:=2C_{\mathbf{B}}C_{R^{-2}}.
\end{align*}
Here $C_\chi'$, $C_H$ and $c_\star$ are the constants through which the radius $R$ of
Theorem~\ref{4.36:thm-boundary-bound} is defined.
Then $C_A\le8+L^4\,2^{4r-1}[2+3(4r)!(2\mathfrak{f})^{4r}]$, the map $t\mapsto F^B(V^{t,\hat T},w)$ is
real-analytic near $0$, and
\begin{equation}\label{4.36:eq-boundary-absorption-a}
\begin{aligned}
|\partial_t^2F^B(V^{t,\hat T},w)|_0
&\le\frac{C_{\mathbf{B}}}{R^2}(\log x_k)^{4r}\bigl[e^{-\delta_0R_*/128}+e^{-\kappa R_*/(16M)}\bigr]\\
&\le\mathfrak{A}^{(4)}x_k(\log x_k)^{4r}e^{-c_{\mathrm{exp}}(\log x_k)^r},
\end{aligned}
\end{equation}
and consequently
\begin{equation}\label{4.36:eq-boundary-absorption-b}
\begin{split}
x_k\ge x_B&:=\exp\Bigl(\max\Bigl\{1,\ \frac{2(1+4r)}{c_{\mathrm{exp}}},\
\Bigl(\frac{2\log^{+}(2C_{\mathbf{B}}C_{R^{-2}})}{c_{\mathrm{exp}}}\Bigr)^{1/2}\Bigr\}\Bigr)\\
&\ \Longrightarrow\ |\partial_t^2F^B(V^{t,\hat T},w)|_0\le1 .
\end{split}
\end{equation}
\end{lemma}

\begin{proof}
The blocks $y$ entering the double sum of hypothesis~(a) are split into two classes according to
their distance from the marked bond: class~(A), $d(b_1,y)\ge R_*/4$, and class~(B),
$d(b_1,y)<R_*/4$.

\emph{Class \textup{(A)}.} For $y$ in class~(A),
\[
e_y^{1/2}=e^{-\delta_0d/32}e^{-\delta_0d/32}\le e^{-\delta_0R_*/128}e_y^{1/4},\qquad d=d(b_1,y),
\]
because $\delta_0d/32\ge\delta_0R_*/128$ and $e^{-\delta_0d/32}=e_y^{1/4}$. Since
$e_y^{1/4}=e^{-\delta_0d(y,b_1)/32}$, the sum of $e_y^{1/4}$ over the blocks of class~(A) in
$\Gamma_{n'}$ is at most $\sigma^{(1/32)}_{n'}$, and by hypothesis~(b) with $\alpha=1/32$,
$\sigma^{(1/32)}_{n'}\le C_\sigma e^{-(k-n')}$.

Next the flow. By hypothesis~(f), $x_{n'}\ge e$, so property (c) of $\check{R}$ gives
$R_{n'}<L(\log x_{n'})^r$, and $x_{n'}\le(1+\mathfrak{f}(k-n'))x_k$, with $\mathfrak{f}=1$, or with
$\mathfrak{f}=\bar B$ when $x_k\ge1$. Hence $\log x_{n'}\le\log x_k+\log(1+\mathfrak{f}(k-n'))$, and the
elementary inequalities \eqref{4.36:eq-sun-conventions-a}, in the form
$(s_1+s_2)^{4r}\le2^{4r-1}(s_1^{4r}+s_2^{4r})$ for $s_1,s_2\ge0$, give
\[
R_{n'}^4\le L^4\,2^{4r-1}\bigl[(\log x_k)^{4r}+(\log(1+\mathfrak{f}(k-n')))^{4r}\bigr].
\]
Put $u:=k-n'$. Inserting into the part of the double sum of hypothesis~(a) that runs over class~(A)
the bound on $e_y^{1/2}$, the count $\mathfrak{S}(y)\le C_{\mathfrak{B}}[1+u+R_{n'}^4]$ of
hypothesis~(a), the bound on $R_{n'}^4$ and the bound on $\sigma^{(1/32)}_{n'}$, and then using
$(\log x_k)^{4r}\ge1$ (as $x_k\ge e$) to factor out $(\log x_k)^{4r}$ from every term of the bracket, we
obtain a bound by
\begin{multline*}
e^{-\delta_0R_*/128}C_\sigma C_{\mathfrak{B}}(\log x_k)^{4r}\\
\times\sum_{u=0}^{k}e^{-u}\bigl[1+u+L^4\,2^{4r-1}\bigl(1+(\log(1+\mathfrak{f}u))^{4r}\bigr)\bigr].
\end{multline*}
Extending this sum to all $u\ge0$, so that $K$ enters nowhere, gives
\begin{equation}\label{4.36:eq-boundary-absorption-1}
\sum_{n'\le k}\ \sum_{y\in\Gamma_{n'},\ d(b_1,y)\ge R_*/4}e_y^{1/2}\mathfrak{S}(y)
\le e^{-\delta_0R_*/128}C_\sigma C_{\mathfrak{B}}C_A(\log x_k)^{4r}.
\end{equation}
The constant $C_A$ is finite, with the stated bound. Indeed
$\sum_{u\ge0}e^{-u}=e/(e-1)\le2$; since $ue^{-u/2}\le2/e<2$ for $u\ge0$,
\[
\sum_{u\ge0}ue^{-u}\le2\sum_{u\ge0}e^{-u/2}=\frac{2}{1-e^{-1/2}}\le6 ;
\]
and, by $\log(1+v)\le v$ and by \eqref{4.36:eq-sun-conventions-a} in the form
$v^{4r}\le(4r)!\,e^{v}$ for $v\ge0$, applied to $v=u/2$,
\[
(\log(1+\mathfrak{f}u))^{4r}\le(\mathfrak{f}u)^{4r}=(2\mathfrak{f})^{4r}(u/2)^{4r}
\le(4r)!(2\mathfrak{f})^{4r}e^{u/2},
\]
while $\sum_{u\ge0}e^{-u}e^{u/2}\le3$. Adding,
$C_A\le2+6+L^4\,2^{4r-1}[2+3(4r)!(2\mathfrak{f})^{4r}]$. At $\mathfrak{f}=1$ this bound does not depend
on $N$.

\emph{Class \textup{(B)}.} Let $d(b_1,y)<R_*/4$. By hypothesis~(c), $y\in\Gamma_k$ and $|y-b_1|<R_*/4$ in
scale-$k$ units. The blocks of the adjacency chain from $\Delta(y)$ to a block meeting $X$ have scale at
most $k$, hence sides at most $1$ in scale-$k$ units, so $X$ contains a site $q$ with
$\operatorname{dist}_k(q,p)\le R_*/4+3$. By cleanness, all sites within distance $R_*/4+4$ of $p$ lie in
$\Lambda_k\subset\Omega_k$; here $R_*/4+4\le R_*/3$, which holds because $R_*\ge48$, and this follows from
$R_*\ge20L^2M\ge48$, true for $L\ge2$, $M\ge1$ and $\max\{R_k,R_{k+1}\}\ge1$; the last holds because
$R_k\ge(\log x_k)^r$ and $x_k\ge e$ (hypothesis~(f)). On the other hand, the
polymers of $\mathfrak{S}(y)$ satisfy $X\cap Z_i\neq\emptyset$, and $Z_i\subset Z_k=\Lambda_k^{c}$ for
$i\le k$, so $X$ meets $\Lambda_k^{c}$ at some site $q'$, and $\operatorname{dist}_k(q',p)\ge R_*$ by
cleanness. Therefore, in scale-$k$ units,
\[
\operatorname{diam}_\infty X\ge R_*-\frac{R_*}{4}-3\ge\frac{2R_*}{3},
\]
the last inequality being equivalent to $R_*\ge36$, which holds. Let $u:=k-i\ge0$ and put
$a:=R_*L^u/M$. Hypothesis~(d) gives $d_i(X)\ge\tfrac23a-2$. Since
$R_*/M\ge20L^2$, we have $a\ge20L^2\ge12$, and $\tfrac23a-2\ge a/2$ is equivalent to $a\ge12$. Hence
\[
d_i(X)\ge\frac{R_*L^u}{2M}\ge10L^2L^u .
\]
We now split $e^{-\kappa\ell/4}=e^{-\kappa\ell/8}e^{-\kappa\ell/8}$ with $\ell=d_i(X)$, bound the second
factor by $e^{-\kappa R_*L^u/(16M)}$ using the last display, and apply the count of hypothesis~(e) at
$n'=k$ with $\theta=\kappa/8$, for which $(k+2-i)_+=u+2$. The sum over $\ell$ defining $\mathfrak{s}(\kappa/8)$
has ratio $e^{-(\kappa/8-c_d)}\le e^{-1}$ by hypothesis~(g), so $\mathfrak{s}(\kappa/8)\le e/(e-1)<2$.
Writing $e^{-\kappa R_*L^u/(16M)}=e^{-\kappa R_*/(16M)}e^{-\kappa R_*(L^u-1)/(16M)}$, we get for the
contribution of level $i=k-u$ to $\mathfrak{S}(y)$:
\begin{align*}
\sum_{X}e^{-\kappa d_i(X)/4}
&\le7^4L^{4(u+2)}\mathfrak{s}(\kappa/8)e^{-\kappa R_*L^u/(16M)}\\
&\le2\cdot7^4L^8e^{-\kappa R_*/(16M)}\bigl[L^{4u}e^{-\kappa R_*(L^u-1)/(16M)}\bigr]\\
&\le2\cdot7^4L^8e^{-\kappa R_*/(16M)}e^{-u}.
\end{align*}
The last inequality is seen as follows. Since $\max\{R_k,R_{k+1}\}\ge1$ (hypothesis~(f):
$R_k\ge(\log x_k)^r\ge1$), $\kappa R_*/(16M)\ge\tfrac54\kappa L^2$,
and $L^u-1\ge u(L-1)$. Hence
\[
\frac{\kappa R_*(L^u-1)}{16M}\ge\frac54\kappa L^2(L-1)u\ge5(L-1)u\ge[4(L-1)+1]u\ge(4\log L+1)u,
\]
where the second inequality uses $\kappa L^2\ge4$, which follows from $\kappa\ge1$ and $L\ge2$; the third
uses $L\ge2$; and the fourth uses $\log L\le L-1$. Thus
$L^{4u}e^{-\kappa R_*(L^u-1)/(16M)}\le L^{4u}e^{-(4\log L+1)u}=e^{-u}$. Summing over $u\ge0$ with
$\sum_{u\ge0}e^{-u}=e/(e-1)<2$ gives, for $y$ in class~(B),
\[
\mathfrak{S}(y)\le4\cdot7^4L^8e^{-\kappa R_*/(16M)} .
\]
All blocks of class~(B) lie in $\Gamma_k$, and $e_y^{1/2}=e^{-\delta_0d(y,b_1)/16}$, so by hypothesis~(b)
with $\alpha=1/16$ and $n=k$, the sum of $e_y^{1/2}$ over the blocks of class~(B) is at most
$\sigma_k^{(1/16)}\le C_\sigma$. With $(\log x_k)^{4r}\ge1$,
\begin{equation}\label{4.36:eq-boundary-absorption-2}
\sum_{y\in\Gamma_k,\ d(b_1,y)<R_*/4}e_y^{1/2}\mathfrak{S}(y)
\le4\cdot7^4L^8C_\sigma(\log x_k)^{4r}e^{-\kappa R_*/(16M)} .
\end{equation}

\emph{Collection.} By the first bound of hypothesis~(a), split into the two classes and estimated by
\eqref{4.36:eq-boundary-absorption-1} and \eqref{4.36:eq-boundary-absorption-2},
\[
|\partial_t^2F^B(V^{t,\hat T},w)|_0\le\frac{2B_0^{\sharp}}{R^2}(\log x_k)^{4r}
\bigl[C_\sigma C_{\mathfrak{B}}C_Ae^{-\delta_0R_*/128}+4\cdot7^4L^8C_\sigma e^{-\kappa R_*/(16M)}\bigr].
\]
Bounding both coefficients by $C_\sigma\max\{C_{\mathfrak{B}}C_A,4\cdot7^4L^8\}$ gives the first
inequality of \eqref{4.36:eq-boundary-absorption-a}, with
$C_{\mathbf{B}}=2B_0^{\sharp}C_\sigma\max\{C_{\mathfrak{B}}C_A,4\cdot7^4L^8\}$.

\emph{Real-analyticity.} The label $h$ is fixed and the torus is finite, so $F^B$ is a finite sum of terms
$\Psi_{i,X}$; these are the boundary terms, indexed by $i$ and $X$ as in the sum $\mathfrak{S}(y)$, into
which $F^B$ is decomposed in the proof of Theorem~\ref{4.36:thm-boundary-bound} and whose second
derivatives are estimated there to obtain \eqref{4.36:eq-boundary-bound-c}. In that proof the
configuration moved along $\hat T$ is continued to complex $\zeta$ as
$e^{i\eta\chi_{\mathsf{N}}\mathcal{H}_\zeta}U_k(V)$, with the number $\eta$ and the function
$\chi_{\mathsf{N}}$ attached to the bond set $\mathsf{N}$ fixed there, and $\lceil X$ denotes restriction
to $X$. Each map $t\mapsto\Psi_{i,X}(U_k(V^{t,\hat T}),w)$ is the restriction to real $t$ of
$\zeta\mapsto\Psi_{i,X}(e^{i\eta\chi_{\mathsf{N}}\mathcal{H}_\zeta}U_k(V)\lceil X)$. This map is analytic on
$|\zeta|<2R$, as the composition of the analytic response field $\zeta\mapsto\mathcal{H}_\zeta$ of
Notation~\ref{4.36:not-probing-direction} with $\Psi_{i,X}$, which is analytic on its domain by
\cite[(2.42)]{B-CMP119} and \cite[(1.63)]{B-CMP122a}. On $|t|\le2R$ it equals the real object, by the
identification of the continued and the real configurations on $|t|\le2R$ made earlier in the proof of
Theorem~\ref{4.36:thm-boundary-bound}. Hence
$t\mapsto F^B(V^{t,\hat T},w)$ is real-analytic near $0$.

\emph{Absorption.} We have $R_*\ge20L^2MR_k\ge20L^2M(\log x_k)^r$, by the lower bound on $R_k$ in
hypothesis~(f). The condition $\delta_0M\ge16$ gives
\[
\frac{\delta_0R_*}{128}\ge\frac{20\cdot16}{128}L^2(\log x_k)^r=\frac52L^2(\log x_k)^r,
\]
and the condition $\kappa\ge10$ gives
\[
\frac{\kappa R_*}{16M}\ge\frac{20\kappa}{16}L^2(\log x_k)^r\ge\frac{25}{2}L^2(\log x_k)^r
\ge\frac52L^2(\log x_k)^r .
\]
With $c_{\mathrm{exp}}=\tfrac52L^2$, both exponentials in \eqref{4.36:eq-boundary-absorption-a} are therefore at
most $e^{-c_{\mathrm{exp}}(\log x_k)^r}$.

Recall that $x_k\ge e$ by hypothesis~(f). We claim $R^{-2}\le C_{R^{-2}}x_k$. The radius $R$ is the minimum of three positive
radii $R^{(1)},R^{(2)},R^{(3)}$, whose inverse squares satisfy
\[
(R^{(1)})^{-2}\le\Bigl(\frac{4C_\chi'C_H}{c_\star}\Bigr)^2\delta_k^{-2},\qquad
(R^{(2)})^{-2}\le(16C_{\mathrm{loc}})^2\varepsilon_k^{-2},\qquad
(R^{(3)})^{-2}\le4r_0^{-2},
\]
and $1/\min\{R^{(1)},R^{(2)},R^{(3)}\}^2\le(R^{(1)})^{-2}+(R^{(2)})^{-2}+(R^{(3)})^{-2}$. Using
$\delta_k^{-2}\le A_1^{-2}x_k$ and
$\varepsilon_k^{-2}\le A_0^{-2}x_k$, both valid because $(\log x_k)^{p_0}\ge1$, and
$4r_0^{-2}\le4r_0^{-2}x_k$, the three inverse squares are bounded by the three terms of $C_{R^{-2}}$ times
$x_k$. Inserting $R^{-2}\le C_{R^{-2}}x_k$ and the bound on the two exponentials into
the first line of \eqref{4.36:eq-boundary-absorption-a} gives its second line, with
$\mathfrak{A}^{(4)}=2C_{\mathbf{B}}C_{R^{-2}}$.

\emph{Inversion.} Put $s:=\log x_k\ge1$. Since $r\ge2$, we have $s^r\ge s^2$. Using in addition
$\log s\le s$,
\[
\log\bigl[\mathfrak{A}^{(4)}e^{s}s^{4r}e^{-c_{\mathrm{exp}}s^r}\bigr]
\le\log^{+}\mathfrak{A}^{(4)}+(1+4r)s-c_{\mathrm{exp}}s^2 .
\]
The right side is at most $0$ as soon as $c_{\mathrm{exp}}s^2/2\ge(1+4r)s$ and
$c_{\mathrm{exp}}s^2/2\ge\log^{+}\mathfrak{A}^{(4)}$, that is, as soon as $s\ge2(1+4r)/c_{\mathrm{exp}}$ and
$s\ge(2\log^{+}\mathfrak{A}^{(4)}/c_{\mathrm{exp}})^{1/2}$. Together with $s\ge1$, these three conditions
say exactly that $x_k\ge x_B$. Since $\mathfrak{A}^{(4)}x_k(\log x_k)^{4r}e^{-c_{\mathrm{exp}}(\log x_k)^r}$ is
the expression in brackets, \eqref{4.36:eq-boundary-absorption-a} yields
\eqref{4.36:eq-boundary-absorption-b}.

The number $x_B$ depends on $(N,L,M,\mathfrak{p})$ only. It is a floor on the inverse coupling, and it is
imposed for every $k\le K$ by the ceiling $\gamma\le(x_B+c_5)^{-1/2}$, through the first of the flow bounds
\eqref{4.36:eq-sun-conventions-b}. Since $x_B\ge e$, the same floor supplies the condition $x_{n'}\ge e$ used
for class~(A). Display \eqref{4.36:eq-boundary-absorption-b} is the statement that the boundary contribution
to the second derivative is at most $1$, at the constants of $\mathrm{SU}(N)$.
\end{proof}

\begin{remark}[The boundary-term bound without group identities or arrest caps]\label{4.36:rem-boundary-group-free}
No step of the proof of the second-derivative bound for the boundary terms at $SU(N)$
(Theorem~\ref{4.36:thm-boundary-bound}) uses $N=2$. That proof rests on
Lemma~\ref{4.36:lem-arrest-per-cube}, Lemma~\ref{4.36:lem-long-or-recent} and
Proposition~\ref{4.36:prop-frozen-sum}, together with the membership statements and scalar bounds
recalled below.
The proof imposes no bound on the number of arrests, nor on the depth to which arrests are
nested.
\end{remark}

\begin{proof}
We treat the group, then the arrests.

\emph{Group content.} The group enters at two places. The first is membership of the moved
configuration $V^{t,\hat{T}}$ in domains defined through the operator norm $|\cdot|$ of
$N\times N$ matrices.
This membership is supplied by two membership statements used in the proof of
Theorem~\ref{4.36:thm-boundary-bound}, each of which places the moved configuration in such a
domain: the membership statement with radius constant $c_{\star}$, whose balls are written in the
norm $\|\cdot\|_{\star}$, and the membership statement used for the arrests.
This membership is read along the velocity $T'=u\hat{T}u^{-1}$, with $u\in SU(N)$. Since a unitary
matrix preserves the length of every vector, conjugation by a unitary matrix preserves the operator
norm, and therefore $|T'|=|\hat{T}|=1$. The second place is the
pair of scalar bounds \cite[(2.42)]{B-CMP119} and \cite[(1.70)]{B-CMP122a}. Both the membership
statements and the two scalar bounds hold for every $N$, and $N$ enters only through the values of
the constant $B_0$ of \cite[(2.42)]{B-CMP119}, the constant $C_0^{\mathrm{fr}}$ of
\cite[(1.70)]{B-CMP122a} and the radius constant $c_{\star}$ of the membership statement.

\emph{Letters and numerals.} The letter bounded in the stage-count inequality
$\mathfrak{N}_W\le R_{k_W}$ is the number of stages $\mathfrak{N}_W=j_W-h_W$ of the arrest $W$,
a count belonging to the construction; it is unrelated to the $N$ of $SU(N)$. The numerals $3$
that occur in the proof count layers and levels.

\emph{Flow factors.} Exactly two flow factors enter Theorem~\ref{4.36:thm-boundary-bound}, and no
others: in distance class (A) of the boundary-term bound in two distance classes, used in the proof
of Theorem~\ref{4.36:thm-boundary-bound}, the sharp flow bound invoked there (or, in its place, the
flow bound under which the flow letter $\mathfrak{f}=\bar{B}$ is admissible); and property (c) among
the displayed properties of the radii $R_j$. Neither of them is one of the flow statements of the
uniqueness analysis; in particular, no flow statement of the uniqueness analysis enters the theorem.

\emph{Threshold.} Through the membership statement with radius constant $c_{\star}$, the theorem uses
the same floor on the inverse coupling, \eqref{4.36:eq-threshold-floors-d}, as the bounds on the
pieces coming from $\mathbf{E}^{(j)}$ and $\mathbf{R}^{(j)}$ earlier in this section.

\emph{Arrests.} The multiplicity bound used for the frozen terms is
Lemma~\ref{4.36:lem-arrest-per-cube}. It is stated per cube $C\in\pi_m$. Its proof uses only the
laminarity of the arrests containing $C$, which are pairwise nested and hence totally ordered, and
the ordering of their stage intervals; both hold whatever the number of arrests and the depth of
nesting. It follows that at most one of them has $m\in(h_W,k_W]$, however many arrests there are
and however deeply they are nested.

Lemma~\ref{4.36:lem-long-or-recent} is applied to one arrest $W$ at a time. In its second case it
refers to a single $W$ with $i\in\mathfrak{E}(W)$.

Hence neither statement, and neither Proposition~\ref{4.36:prop-frozen-sum} nor
Theorem~\ref{4.36:thm-boundary-bound}, which rest on them, needs a cap on the number or nesting of
arrests.
\end{proof}

\begin{corollary}[The curvature constant of the near part at $SU(N)$]\label{4.36:cor-near-curvature}
Work under the standing hypotheses of Definition~\ref{4.36:def-near-part}. These fix $K$, $k$, a clean
history $h$, the configuration $V$ with its free link, a fluctuation variable $w$ in the support of
$m_h$, and a direction $\hat{T}\in\mathfrak{g}$ with $|\hat{T}|=1$. They also include the floor
$x_n\ge x_{\mathrm{th}}^{(3)}$ for all $n\le K$. For every $T\in\mathfrak{g}$ write $V^{t,T}$ for the
configuration moved along $T$ at the free link. Assume the following statements.
\begin{itemize}
\item[(a)] The second-derivative bound for the term $F^W$ of
Definition~\ref{4.36:def-near-part}, at every $N$:
\begin{equation*}
\bigl|\partial_t^2 F^W(V^{t,\hat{T}})\bigr|_0\le C_W'(N)\,x_k,
\end{equation*}
with the constant given by
\begin{equation*}
C_W'(N)=36\,C_q(N)\,C_\sigma S_{\mathfrak{n}},\qquad
C_q(N)=2C_P\bigl(B_3a_1'/r_0+C_P\bigr).
\end{equation*}
Here $C_W'(N)$ and $C_q(N)$ are as defined in Definition~\ref{4.36:def-near-part}, and
$C_P$, $B_3$, $a_1'$, $r_0$, $C_\sigma$ and $S_{\mathfrak{n}}$ are the constants fixed in the
statement of this bound for $F^W$.
\item[(b)] The second-derivative bound for the regular terms at $SU(N)$
(Theorem~\ref{4.36:thm-regular-bound}), in the form
$|\partial_t^2F^E(V^{t,\hat{T}})|_0\le x_k/8$ of its absorbed part.
\item[(c)] The second-derivative bound for the renormalisation-operation terms at $SU(N)$
(Theorem~\ref{4.36:thm-operation-bound}), in the absorbed form given by
Lemma~\ref{4.36:prf-operation-absorption}: the two sums
$\mathfrak{r}_{\mathrm{small}}$ and $\mathfrak{r}_{\mathrm{tail}}$ on the right of
the corresponding display are each at most $\tfrac12$.
\item[(d)] The second-derivative bound for the boundary terms at $SU(N)$
(Theorem~\ref{4.36:thm-boundary-bound}), in the form
$|\partial_t^2F^B(V^{t,\hat{T}},w)|_0\le 1$ of the last inequality in
\eqref{4.36:eq-boundary-bound-a}.
\item[(e)] The rescaling identity along $|\cdot|_{\mathfrak{g}}$-unit directions. Let
$T\in\mathfrak{g}$ with $|T|_{\mathfrak{g}}=1$, and let $|T|$ denote its operator norm. Then
$|T|^2\le N-1$, and
\begin{equation*}
\bigl|\partial_t^2F^{\mathrm{near}}(V^{t,T},w)\bigr|_0
=|T|^2\,\bigl|\partial_t^2F^{\mathrm{near}}(V^{t,T/|T|},w)\bigr|_0 .
\end{equation*}
\end{itemize}
Then the following hold.
\begin{itemize}
\item[(i)] We have
\begin{equation}\label{4.36:eq-near-curvature-a}
\begin{gathered}
\bigl|\partial_t^2F^{\mathrm{near}}(V^{t,\hat{T}},w)\bigr|_0
\le C_W'(N)\,x_k+\frac{x_k}{8}+1+1\le C_F^{\prime\,\mathrm{op}}(N)\,x_k,\\
C_F^{\prime\,\mathrm{op}}(N)=C_W'(N)+3 .
\end{gathered}
\end{equation}
\item[(ii)] For every $T\in\mathfrak{g}$ with $|T|_{\mathfrak{g}}=1$ and $T/|T|=\hat{T}$,
\begin{equation*}
\bigl|\partial_t^2F^{\mathrm{near}}(V^{t,T},w)\bigr|_0\le C_F^{\prime\,\mathfrak{g}}(N)\,x_k,
\qquad C_F^{\prime\,\mathfrak{g}}(N)=(N-1)\bigl(C_W'(N)+3\bigr).
\end{equation*}
\end{itemize}
\end{corollary}

\begin{proof}
By Definition~\ref{4.36:def-near-part} the near part decomposes as
\begin{equation*}
F^{\mathrm{near}}=F^W+F^E+F^R+F^B+\mathrm{const}.
\end{equation*}
The additive constant does not depend on the free-link variable. It is therefore the same for
every $V^{t,\hat{T}}$, and its $t$-derivatives vanish. The terms $F^E$ and $F^R$ do not depend
on $w$ by Definition~\ref{4.36:def-near-part}, and $F^W$ is written, as in hypothesis~(a), as a
function of the configuration alone. Only $F^B$ carries the argument $w$.

The second $t$-derivative is linear. By the triangle inequality we therefore have
\begin{equation*}
\begin{split}
\bigl|\partial_t^2F^{\mathrm{near}}(V^{t,\hat{T}},w)\bigr|_0
&\le\bigl|\partial_t^2F^W(V^{t,\hat{T}})\bigr|_0+\bigl|\partial_t^2F^E(V^{t,\hat{T}})\bigr|_0\\
&\quad+\bigl|\partial_t^2F^R(V^{t,\hat{T}})\bigr|_0+\bigl|\partial_t^2F^B(V^{t,\hat{T}},w)\bigr|_0 .
\end{split}
\end{equation*}
We bound the four terms in turn.
\begin{itemize}
\item The first term is at most $C_W'(N)x_k$ by hypothesis~(a).
\item The second term is at most $x_k/8$ by hypothesis~(b).
\item The third term is at most $\mathfrak{r}_{\mathrm{small}}+\mathfrak{r}_{\mathrm{tail}}$ by
the corresponding display. By hypothesis~(c) this is at most
$\tfrac12+\tfrac12=1$.
\item The fourth term is at most $1$ by hypothesis~(d).
\end{itemize}
All four hypotheses are stated under the standing hypotheses of
Definition~\ref{4.36:def-near-part}, including $x_n\ge x_{\mathrm{th}}^{(3)}$ for all $n\le K$.
Adding the four bounds gives the first inequality of \eqref{4.36:eq-near-curvature-a}.

For the second inequality it suffices to show that $x_k/8+2\le 3x_k$. Multiplying by $8$, this
is equivalent to
\begin{equation*}
x_k+16\le 24x_k,\qquad\text{that is,}\qquad 16\le 23x_k .
\end{equation*}
Since $\gamma<1$, we have $x_k=g_k^{-2}>1$. Hence $23x_k>23>16$, and the inequality holds. It
follows that
\begin{equation*}
C_W'(N)\,x_k+\frac{x_k}{8}+2\le\bigl(C_W'(N)+3\bigr)x_k=C_F^{\prime\,\mathrm{op}}(N)\,x_k .
\end{equation*}
This proves~(i).

For~(ii), let $T\in\mathfrak{g}$ with $|T|_{\mathfrak{g}}=1$ and $T/|T|=\hat{T}$, with $|T|$
the operator norm; thus $|\hat{T}|=1$, and part~(i) applies to this direction.
Hypothesis~(e), which is stated for $F^{\mathrm{near}}$ as a whole, gives
\begin{equation*}
\bigl|\partial_t^2F^{\mathrm{near}}(V^{t,T},w)\bigr|_0
=|T|^2\,\bigl|\partial_t^2F^{\mathrm{near}}(V^{t,\hat{T}},w)\bigr|_0 .
\end{equation*}
Since $|T|^2\le N-1$, part~(i) bounds the right-hand side by
$(N-1)\bigl(C_W'(N)+3\bigr)x_k$. The factor $N-1$ is thus taken once, over the sum of the four
members, and not member by member.
\end{proof}

The summand $3$ in $C_F^{\prime\,\mathrm{op}}(N)$ does not depend on $N$. It rounds up
$\tfrac18+1/x_k+1/x_k$. The three members rounded here, the regular, renormalisation-operation and
boundary terms, are those written $E_k$, $R_k$, $B_k$ in the decomposition
$-(1/g_k^2)A+E_k+R_k+B_k$ of the effective action in \cite[(2.23)]{B-CMP119}.

Every dependence of $C_F^{\prime\,\mathrm{op}}(N)$ on $N$ enters through $C_W'(N)$. That constant
is $36\,C_q(N)C_\sigma S_{\mathfrak{n}}$, in which $C_q(N)$ and the constant
$\bar{B}(N,L,\mathfrak{p})$ entering $S_{\mathfrak{n}}$ depend on $N$; it is a value depending on
$(N,L,\mathfrak{p})$. No independence of $N$ is asserted for it. The constant that enters the mode
lemma is $C_W'(N)+3$, whether directions are normalised by $|\hat{T}|=1$ or by
$|T|_{\mathfrak{g}}=1$.

In what follows $N$, $L$, $M$ and the auxiliary parameters $\mathfrak{p}$ are fixed, and $x_j=g_j^{-2}$. We first collect the three absorption floors. For $x>1$ put
\begin{equation}\label{4.36:eq-coupling-ceiling-1}
\mathfrak{r}(x):=\frac{(4C_{\chi}'C_H/c_{\star})^2\,x}{A_1^2(\log x)^{2p_0}}
+\frac{(16C_{\mathrm{loc}})^2\,x}{A_0^2(\log x)^{2p_0}}+4r_0^{-2}.
\end{equation}
The radius $R$ of Notation~\ref{4.36:not-sun-conventions} is the minimum of the three numbers $c_{\star}\delta_k/(4C_{\chi}'C_H)$, $\varepsilon_k/(16C_{\mathrm{loc}})$ and $r_0/2$. By the same notation,
\begin{equation*}
\delta_k^{-2}=\frac{x_k}{A_1^2(\log x_k)^{2p_0}},\qquad
\varepsilon_k^{-2}=\frac{x_k}{A_0^2(\log x_k)^{2p_0}}.
\end{equation*}
Hence $\mathfrak{r}(x_k)$ is the sum of the inverse squares of these three numbers. The inverse square of each of them, and in particular $R^{-2}$, is therefore at most $\mathfrak{r}(x_k)$. The implications below are read together with the lower bound $R_*\ge 20L^2M(\log x_k)^r$ on the radius $R_*$ that enters the threshold \eqref{4.36:eq-threshold-floors-c}.

With these bounds each of the three absorptions is implied by a floor on $x_k$.
\begin{itemize}
\item The absorption $\mathfrak{E}^{(2)}\le x_k/8$ of the regular term holds as soon as $x_k\ge x_E$; see \eqref{4.36:eq-regular-absorption-c}.
\item The bound $|\partial_t^2F^R|_0\le1$ on the operation term holds as soon as $x_k\ge x_R$; see \eqref{4.36:eq-operation-absorption-b}.
\item The bound $|\partial_t^2F^B(V^{t,\hat{T}},w)|_0\le1$ on the boundary term holds as soon as $x_k\ge x_B$; see \eqref{4.36:eq-boundary-absorption-b}.
\end{itemize}
The analogous floor $x_{\mathrm{far}}$ for the far part, above which the size $\eta_{\mathrm{far}}$ of the far part is at most $1$, belongs to the far-part statement and is only named here: it is not a member of the maxima below. We put
\begin{equation*}
x_{\mathrm{abs}}:=\max\{x_E,x_R,x_B\}.
\end{equation*}

\begin{proposition}[A coupling ceiling implying every floor]\label{4.36:prop-coupling-ceiling}
Assume the following.
\begin{enumerate}
\item[(a)] The lower member of the coupling flow in \eqref{4.36:eq-sun-conventions-b} holds at $k=n$. That is, for every $n\le K$,
\begin{equation*}
x_n\ge x_K+\lambda(K-n)-c_5,
\end{equation*}
with $\lambda\ge\lambda^{\sharp}>0$ and $c_5\ge0$. In addition $g_K\le\gamma$.
\item[(b)] The auxiliary parameters satisfy $q:=\min\{q_0,q_1\}\ge p_0+3r+2$, $p_0\ge5r$, $r\ge2$, $\kappa\ge\max\{10,8c_d+8\}$ and $\kappa_0\ge7$.
\item[(c)] The flow comparisons of \eqref{4.36:eq-sun-conventions-b} are taken under the ceiling
\begin{equation*}
\gamma\le\gamma_{\bar B}:=\bigl(\max\{\bar B,x_\delta,x_\beta\}+c_5\bigr)^{-1/2}.
\end{equation*}
Here $x_\delta:=\max\{e,2c_5,4p_0c_5\}$ and $x_\beta:=\max\{e,2p_0c_5/\beta_0\}$ are the floors of Notation~\ref{4.36:not-sun-conventions}.
\item[(d)] $\gamma_N$ is a positive number: the ceiling on $\gamma$ at $\mathrm{SU}(N)$ imposed in the setting of this chapter, distinct from the ceiling $\gamma_{\bar B}$ of (c) and from the ceiling $\gamma^{(3)}$ defined below.
\end{enumerate}
Define
\begin{equation}\label{4.36:eq-coupling-ceiling-a}
\begin{split}
x_{\mathrm{th}}^{(3)}(N,L,M,\mathfrak{p})&:=\max\{x_{\mathrm{abs}},\,x_{\mathrm{comp}}\},\\
x_{\mathrm{comp}}&:=\max\bigl\{e,\ e^{2(p_0+r)},\ 2c_5,\ x_\delta,\ x_\beta,\ \bar B;\
x^{(1)},\dots,x^{(5)};\\
&\qquad\qquad x^{(6)},\ x^{(7)};\ x^{(t1)},\ x^{(t8^{+})},\ x^{(t9)};\ x^{(\mathrm{tow})}\bigr\},
\end{split}
\end{equation}
and
\begin{equation}\label{4.36:eq-coupling-ceiling-b}
\gamma^{(3)}:=\bigl(x_{\mathrm{th}}^{(3)}+c_5\bigr)^{-1/2}.
\end{equation}
The maximum \eqref{4.36:eq-coupling-ceiling-a} has twenty members, grouped as follows.
\begin{itemize}
\item $x_E$, $x_R$, $x_B$ are those of \eqref{4.36:eq-regular-absorption-c}, \eqref{4.36:eq-operation-absorption-b} and \eqref{4.36:eq-boundary-absorption-b}.
\item $\bar B$, $x_\delta$, $x_\beta$ are the floors of the ceiling $\gamma_{\bar B}$ of (c); $e$ and $2c_5$ are the floors entering $x_\delta$ in Notation~\ref{4.36:not-sun-conventions}; $e^{2(p_0+r)}$ is the floor above which the monotonicity of \eqref{4.36:eq-threshold-floors-a} holds.
\item $x^{(1)},\dots,x^{(5)}$ are the chart thresholds of \eqref{4.36:eq-threshold-floors-b}.
\item $x^{(6)}$ and $x^{(7)}$ are those of \eqref{4.36:eq-threshold-floors-c} and \eqref{4.36:eq-threshold-floors-d}.
\item $x^{(t1)}$, $x^{(t8^{+})}$, $x^{(t9)}$ and the tower floor $x^{(\mathrm{tow})}$ are those of \eqref{4.36:eq-threshold-floors-e}.
\end{itemize}
Not members of \eqref{4.36:eq-coupling-ceiling-a}: the floor $x_{\mathrm{far}}$; the radii $R$ and $r_m$, which are definitions, every absorption holding for every member of the minima defining them; and the further conditions of the ceiling on $\gamma$ to which the tower condition of \eqref{4.36:eq-threshold-floors-e} belongs, in particular the condition $\bar BP_1(x_k)\le x_k$ on the function $P_1$ of the bound on the piece $F^W$ of the near part used in Corollary~\ref{4.36:cor-near-curvature}, which belongs to that bound.

Then the following hold.
\begin{enumerate}
\item[(i)] $x_{\mathrm{th}}^{(3)}$ is finite and $\gamma^{(3)}$ is a finite positive number once $(N,L,M,\mathfrak{p})$ are fixed. No member of \eqref{4.36:eq-coupling-ceiling-a} depends on $\gamma$, $K$, $k$ or $x_{\mathrm{th}}^{(3)}$.
\item[(ii)] If $\gamma\le\gamma^{(3)}$, then $x_n\ge x_{\mathrm{th}}^{(3)}$ for every $n\le K$. Hence every floor of \eqref{4.36:eq-coupling-ceiling-a} holds at every scale. In particular, the absorptions \eqref{4.36:eq-regular-absorption-c}, \eqref{4.36:eq-operation-absorption-b} and \eqref{4.36:eq-boundary-absorption-b} hold for every member of the minima defining $R$ and $r_m$. The following also hold at every scale:
\begin{equation*}
p_0(g_k)\ge8C_{\mathrm{loc}},\qquad g_k\le r_0/2<r_0,\qquad R_j\ge L^2,\qquad
\mathfrak{c}_T\sup_{n\le K}\hat\varepsilon_n\le a^{\flat}.
\end{equation*}
\item[(iii)] $\gamma^{(3)}\le\gamma_{\bar B}$. Hence the ceiling of \textup{(c)} is implied by $\gamma\le\gamma^{(3)}$.
\item[(iv)] For every $L$ the window $(0,\min\{\gamma_N,\gamma^{(3)}\}]$ of admissible values of $\gamma$ is non-empty.
\item[(v)] None of $K$, the torus $\mathbb{T}_m$, $s$, the large-field quantity $N_0$, or the number and nesting of the arrests enters $\gamma^{(3)}$ or any member of \eqref{4.36:eq-coupling-ceiling-a}. The condition $\gamma\le\gamma^{(3)}$ is an upper bound only and imposes no floor on $g$ or on the couplings $g_n$. It is not a cap in the sense of Definition~\ref{2.3:def-cap}.
\end{enumerate}
\end{proposition}

\begin{proof}
\emph{Transport to every scale.} Let $\gamma\le\gamma^{(3)}$ and $n\le K$. By hypothesis (a),
\begin{equation}\label{4.36:eq-coupling-ceiling-2}
x_n\ge x_K+\lambda(K-n)-c_5\ge\gamma^{-2}-c_5 .
\end{equation}
The second inequality holds because $x_K=g_K^{-2}\ge\gamma^{-2}$ by $g_K\le\gamma$, and because $\lambda(K-n)\ge0$ by $\lambda>0$ and $n\le K$. From $\gamma\le\gamma^{(3)}$ and \eqref{4.36:eq-coupling-ceiling-b} we get $\gamma^{-2}\ge(\gamma^{(3)})^{-2}=x_{\mathrm{th}}^{(3)}+c_5$. Together with \eqref{4.36:eq-coupling-ceiling-2} this gives $x_n\ge x_{\mathrm{th}}^{(3)}$.

Every member of \eqref{4.36:eq-coupling-ceiling-a} is at most $x_{\mathrm{th}}^{(3)}$, so every floor of \eqref{4.36:eq-coupling-ceiling-a} holds at every scale $n\le K$. We now read off the consequences listed in (ii).
\begin{itemize}
\item Since $x_k\ge x_E,x_R,x_B$, the absorptions \eqref{4.36:eq-regular-absorption-c}, \eqref{4.36:eq-operation-absorption-b} and \eqref{4.36:eq-boundary-absorption-b} hold. For $R$, they hold for every member of its minimum by the bound established after \eqref{4.36:eq-coupling-ceiling-1}: the inverse square of each member of that minimum is at most $\mathfrak{r}(x_k)$. For $r_m$, this is part of \eqref{4.36:eq-operation-absorption-a}, where both members of the minimum defining $r_m$ are treated.
\item Since $x_k\ge x^{(t1)},x^{(t8^{+})}$ and $x_j\ge x^{(t9)}$, the implications displayed in \eqref{4.36:eq-threshold-floors-e} give three thresholds: $p_0(g_k)\ge8C_{\mathrm{loc}}$, $g_k\le r_0/2<r_0$, and $R_j\ge L^2$.
\item Since $x_n\ge x^{(\mathrm{tow})}$ for all $n\le K$, the tower ceiling $\mathfrak{c}_T\sup_{n\le K}\hat\varepsilon_n\le a^{\flat}$ of \eqref{4.36:eq-threshold-floors-e} holds.
\end{itemize}
This proves (ii).

\emph{Finiteness and positivity.} We go through the members of \eqref{4.36:eq-coupling-ceiling-a}. Each is a number once $(N,L,M,\mathfrak{p})$ are fixed, provided the exponents and constants entering it are admissible. We check this in turn.
\begin{itemize}
\item The exponents $1/(q_0-p_0-r)$, $1/(q_1-p_0-r)$, $1/(q-p_0-r)$, $1/(q_1-p_0)$ and $1/p_0$ are positive by hypothesis (b). Indeed $q_0,q_1\ge q\ge p_0+3r+2$. Hence each of $q_0-p_0-r$, $q_1-p_0-r$ and $q-p_0-r$ is at least $2r+2>0$, and $q_1-p_0\ge3r+2>0$. Finally $p_0\ge5r\ge10>0$.
\item The exponents $2/(\kappa_0-4)$ and $2/(\kappa_0-2)$ of $x_R$ have positive denominators, since $\kappa_0\ge7$ gives $\kappa_0-4\ge3$ and $\kappa_0-2\ge5$.
\item The constants $C_R$ and $C''$ of \eqref{4.36:eq-operation-absorption-b} are finite. The series $S_{\mathfrak{f}}$ and $C_{\mathrm{ct}}$ of \eqref{4.36:eq-operation-absorption-a}, from which they are formed, converge because $\kappa/2-c_d>0$. This positivity follows from $\kappa\ge8c_d+8$, which gives $\kappa/2-c_d\ge3c_d+4>0$ for the non-negative constant $c_d$.
\item By \eqref{4.36:eq-threshold-floors-a}, the number $x(c)$ defined there is finite for every $c>0$. This covers the members $x^{(3)}$, $x^{(5)}$ and $x^{(\mathrm{tow})}$, which are of the form $x(c)$.
\item The root $\mathfrak{k}_6^{1/r}$ in $x^{(6)}$ has $r\ge2$.
\item $x^{(t8^{+})}=4r_0^{-2}$ is finite because $r_0>0$.
\item The remaining members are numbers after $(N,L,M,\mathfrak{p})$.
\end{itemize}
Hence $x_{\mathrm{th}}^{(3)}$ is a finite maximum of twenty numbers and $x_{\mathrm{th}}^{(3)}\ge e>0$. Since $c_5\ge0$, we have $x_{\mathrm{th}}^{(3)}+c_5\ge e>0$, and \eqref{4.36:eq-coupling-ceiling-b} defines a finite positive number $\gamma^{(3)}$.

\emph{No circularity.} Every member of \eqref{4.36:eq-coupling-ceiling-a} is an expression in $(N,L,M,\mathfrak{p})$ alone. None reads $\gamma$, $K$, $k$ or $x_{\mathrm{th}}^{(3)}$. Thus $\gamma^{(3)}$ is fixed before $\gamma$: the ceiling $\gamma\le\gamma^{(3)}$ is imposed first, and its consequences in (ii) are read after. This proves (i).

\emph{Comparison with the ceiling of \textup{(c)}.} The numbers $\bar B$, $x_\delta$ and $x_\beta$ are members of \eqref{4.36:eq-coupling-ceiling-a}, so $x_{\mathrm{th}}^{(3)}\ge\max\{\bar B,x_\delta,x_\beta\}$. The map $y\mapsto(y+c_5)^{-1/2}$ is decreasing on $(-c_5,\infty)$, which gives $\gamma^{(3)}\le\gamma_{\bar B}$. Consequently $\gamma\le\gamma^{(3)}$ implies $\gamma\le\gamma_{\bar B}$, so the ceiling under which the flow comparisons of hypothesis (c) are taken is implied by $\gamma\le\gamma^{(3)}$. This proves (iii).

\emph{The window.} By hypothesis (d), $\gamma_N>0$, and by (i), $\gamma^{(3)}>0$. Hence $\min\{\gamma_N,\gamma^{(3)}\}>0$, and the interval $(0,\min\{\gamma_N,\gamma^{(3)}\}]$ of admissible values of $\gamma$ is non-empty. No restriction on $L$ was used, so this holds for every $L$. This proves (iv).

\emph{Dependence and one-sidedness.} By (i), $\gamma^{(3)}$ and every member of \eqref{4.36:eq-coupling-ceiling-a} depend on $(N,L,M,\mathfrak{p})$ only. Hence none of $K$, the torus $\mathbb{T}_m$, $s$, the large-field quantity $N_0$, or the number and nesting of the arrests enters $\gamma^{(3)}$ or any of these members. The condition $\gamma\le\gamma^{(3)}$ is an upper bound on $\gamma$. It imposes no lower bound on $g$ or on the couplings $g_n$, and it bounds none of $L$, $M$ or the auxiliary parameters. It is therefore not a cap in the sense of Definition~\ref{2.3:def-cap}. This proves (v).
\end{proof}

\begin{remark}[Absorbing the flow factor into spare logarithms]\label{4.36:rem-spare-logarithms}
Notation is that of Lemma~\ref{4.36:lem-threshold-floors}: $x_j=g_j^{-2}$, $\bar B$ is the flow constant,
$\mathfrak{k}_1,\mathfrak{k}_2,\mathfrak{k}_3,\mathfrak{k}_{\mathrm{im}},\mathfrak{k}_5,\mathfrak{k}_6$ are the
constants defined there, and $a_m$ is the chart's datum of \eqref{4.36:eq-threshold-floors-b}. The coupling ceiling
$\gamma\le\gamma_{\bar B}$, on which the comparison $(\mathrm{F}3)^{\sharp}$ in \eqref{4.36:eq-sun-conventions-b}
rests, is not assumed. Instead, assume $x_j\ge\max\{x_\delta,x_\beta\}$ for every $j$ (both floors are members of
the threshold of Proposition~\ref{4.36:prop-coupling-ceiling}), and assume the following.
\begin{enumerate}
\item[(a)] The second and third flow bounds of \eqref{4.36:eq-sun-conventions-b}.
\item[(b)] Clauses (iii) and (iv) of the lemma bounding two-scale ratios of the running coupling. Applied to (a),
they give each transport inequality used in Lemma~\ref{4.36:lem-threshold-floors}, with its right-hand side
multiplied by $(1+\bar B)^{1/2}$.
\item[(c)] The condition on the exponents of the auxiliary parameters, which supplies the spare exponents
\[
q-p_0\ge 2r+2,\qquad q-p_0-r\ge 2r+2,\qquad q_1-p_0\ge 2r+2,\qquad q_0-p_0-r\ge 2r+2 .
\]
\end{enumerate}
Then the conclusion of Lemma~\ref{4.36:lem-threshold-floors} holds with its floors modified as follows. The
conditions defining the floors $x^{(1)}$, $x^{(2)}$ and $x^{(4)}$ of \eqref{4.36:eq-threshold-floors-b}, and the
floor $x^{(7)}$ of \eqref{4.36:eq-threshold-floors-d}, are replaced by
\begin{equation}\label{4.36:eq-spare-logarithms-1}
\begin{aligned}
(\log x_j)^{q_0-p_0-r}&\ge(1+\bar B)\,\mathfrak{k}_1, &\qquad
(\log x_j)^{q_1-p_0-r}&\ge(1+\bar B)\,\mathfrak{k}_2,\\
(\log x_j)^{q-p_0-r}&\ge(1+\bar B)\,\mathfrak{k}_{\mathrm{im}}, &\qquad
(\log x_j)^{q_1-p_0}&\ge\frac{8(1+\bar B)^{1/2}A_1}{C_1};
\end{aligned}
\end{equation}
for instance, the floor $x^{(1)}$ becomes
\[
\max\bigl\{e,\exp\bigl(((1+\bar B)\,\mathfrak{k}_1)^{1/(q_0-p_0-r)}\bigr)\bigr\}.
\]
In the floors $x_m\ge x^{(3)}$ and $x_m\ge x^{(5)}$ of \eqref{4.36:eq-threshold-floors-b}, the factor
$(1+\bar B)^{1/2}$ enters through the transport of the datum $a_m$ in \eqref{4.36:eq-threshold-floors-b}, not
through the conditions defining these two floors themselves; in those conditions the constants are replaced by
\begin{equation}\label{4.36:eq-spare-logarithms-2}
\mathfrak{k}_3\mapsto(1+\bar B)^{1/2}\,\mathfrak{k}_3,\qquad
\mathfrak{k}_5\mapsto(1+\bar B)^{1/2}\,\mathfrak{k}_5 .
\end{equation}
In the floor $x_k\ge x^{(6)}$ of \eqref{4.36:eq-threshold-floors-c}, the constant is replaced by
\begin{equation}\label{4.36:eq-spare-logarithms-3}
\mathfrak{k}_6\mapsto\frac{4\log\bigl(2(1+\bar B)^{1/2}\bigr)}{5\delta_0L^2M}.
\end{equation}
Finally, in the gap sum $S_n$ over scales, the sum $S_{\mathfrak{f}}$, the constant $C_A$ and the constant
$\Theta_m$, in each of which the flow constant enters through a parameter $\mathfrak{f}\in\{1,\bar B\}$, the
parameter $\mathfrak{f}$ is set equal to $\bar B$.

The argument of Lemma~\ref{4.36:lem-threshold-floors} is the one adopted in what follows, because it keeps each of
its constants as defined there. In neither argument is the inequality $L^2\ge1+\bar B$ used at any step of the
derivations of the floors collected in Proposition~\ref{4.36:prop-coupling-ceiling}. That inequality occurs only as
one item in the list of hypotheses of the statement to which the floors $x^{(1)},\dots,x^{(5)}$ of
\eqref{4.36:eq-threshold-floors-b} are supplied.
\end{remark}

\begin{proof}
By hypotheses (a) and (b), each transport inequality entering the derivation of the floors of
Lemma~\ref{4.36:lem-threshold-floors} holds with the additional factor $(1+\bar B)^{1/2}$ on its right-hand side.

Consider first the floors defining $x^{(1)}$, $x^{(2)}$, $x^{(4)}$ and $x^{(7)}$. In the first three, the inequality
to be verified involves the datum $a_m$ (in the third through $c^{\mathrm{re}}_m=C_{\mathrm{re}}a_m$) transported
from scale $m$ to scale $j$. The transport of the datum $a_m$ in \eqref{4.36:eq-threshold-floors-b} and the
transport from scale $m$ to scale $j$ each carry the factor $(1+\bar B)^{1/2}$, so the constants $\mathfrak{k}_1$,
$\mathfrak{k}_2$ and $\mathfrak{k}_{\mathrm{im}}$ acquire the factor $1+\bar B$; a term transported only once, such
as $c_B\delta_m$ in the third, is covered because $(1+\bar B)^{1/2}\le1+\bar B$. In the fourth, $\delta$ is
transported once, and the constant $8A_1/C_1$ acquires the factor $(1+\bar B)^{1/2}$. In each of the four cases the
inequality to be verified carries a spare power $(\log x_j)^{\eta}$ (here $\eta$, and below $\theta$, are local
letters), where $\eta$ is, respectively,
\[
q_0-p_0-r,\qquad q_1-p_0-r,\qquad q-p_0-r,\qquad q_1-p_0 ,
\]
and the factor is moved into the constant on the right-hand side of the inequality for $(\log x_j)^{\eta}$. This
gives the four inequalities \eqref{4.36:eq-spare-logarithms-1}. By hypothesis (c), $\eta\ge2r+2$ in the first,
third and fourth cases, and $q_1-p_0-r\ge r+2$ in the second; in each case $\eta>0$.

For $x_j\ge e$, an exponent $\eta>0$ and a number $\theta>0$, the
inequality $(\log x_j)^{\eta}\ge\theta$ is equivalent to $x_j\ge\exp(\theta^{1/\eta})$. Each inequality in
\eqref{4.36:eq-spare-logarithms-1} is therefore the one-sided floor
\[
x_j\ge\max\{e,\exp(\theta^{1/\eta})\},
\]
with $\theta$ its modified constant.

The floors defining $x^{(3)}$ and $x^{(5)}$ involve no change of scale. In them the factor $(1+\bar B)^{1/2}$ does
not enter through the inequality itself. It enters through the bound on the datum $a_m$ in
\eqref{4.36:eq-threshold-floors-b}, and it is absorbed by the replacement \eqref{4.36:eq-spare-logarithms-2}. The
floors become
\[
x_m\ge x\bigl((1+\bar B)^{1/2}\mathfrak{k}_3\bigr),\qquad x_m\ge x\bigl((1+\bar B)^{1/2}\mathfrak{k}_5\bigr),
\]
with $x(\cdot)$ the function of Lemma~\ref{4.36:lem-threshold-floors}.

In the floor defining $x^{(6)}$, the constant $\mathfrak{k}_6$ is replaced by \eqref{4.36:eq-spare-logarithms-3}. The
floor keeps the form $x_k\ge\max\{e,\exp(\mathfrak{k}_6^{1/r})\}$.

In the gap sum $S_n$, the sum $S_{\mathfrak{f}}$ and the constants $C_A$, $\Theta_m$, the parameter $\mathfrak{f}$
is set equal to $\bar B$.

Hence every threshold is again a one-sided floor on the inverse coupling. This is the conclusion of
Lemma~\ref{4.36:lem-threshold-floors}, now with constants depending on $\bar B$ as displayed.
\end{proof}

Throughout, $0<g_j\le\gamma<1$ for $0\le j\le K$, so that $x_j=g_j^{-2}>1$.

The flow constants $\lambda$, $c_5$, $B$ and $\bar B$, with $B\le\bar B$, are those of Notation~\ref{4.36:not-sun-conventions}. So is the two-sided flow bound
\begin{equation}\label{4.36:eq-monotone-substitution-1}
\lambda(j-i)-c_5\;\le\;x_i-x_j\;\le\;B(j-i),\qquad 0\le i\le j\le K,
\end{equation}
together with its consequences displayed in \eqref{4.36:eq-sun-conventions-b}.

The sizes $\varepsilon_j,\delta_j,\hat\varepsilon_j,a_j,c^{\mathrm{re}}_j,c^{\mathrm{im}}_j$ and the radii $\alpha_{\cdot,j},\alpha_{3,j}$ are those of Lemma~\ref{4.36:lem-threshold-floors}. The function $p_0(\cdot)$ is the degree function, and $\beta_0$ is the constant of the transport shapes in \eqref{4.36:eq-sun-conventions-b}. The scale $R_j$ is the least power of $L$ not smaller than $(\log x_j)^r$.

We write $x_\delta:=\max\{e,2c_5,4p_0c_5\}$ and $x_\beta:=\max\{e,2p_0c_5/\beta_0\}$. These are the floors on the inverse coupling under which the two-scale comparison of the logarithmic factors $(\log x)^{p_0}$ holds. We also put
\begin{equation*}
\gamma_{\bar B}:=\bigl(\max\{\bar B,x_\delta,x_\beta\}+c_5\bigr)^{-1/2}
\end{equation*}
for the corresponding coupling ceiling, and $\varpi_{jn}:=\max\{1,(x_j/x_n)^2\}$. Finally, $\vartheta_{jn}$ denotes the two-scale ratio of chart sizes between scales $j$ and $n$ that enters the construction of the localised chart, and $C(L,M)$ the constant accompanying it there. The ratio $\vartheta_{jn}$ is not to be confused with the parameter $\vartheta$ of the separation window $\mathfrak{T}(\mathsf{d};\vartheta)$.

For a real parameter $b\ge0$ let $(\mathrm{R}_b)$ denote the following family of two-scale coupling-ratio bounds. They are required for all $j\le n\le k\le K$, with $m:=n-j$:
\begin{equation}\label{4.36:eq-monotone-substitution-2}
\begin{gathered}
\frac{x_j}{x_n}\le 1+bm,\qquad \frac{g_n}{g_j}\le(1+bm)^{1/2},\qquad
\log x_j-\log x_n\le\log(1+bm),\\
\frac{Y_n}{Y_j}\le(1+\beta_0)(1+bm)^{1/2}\quad
\text{for }Y\in\{\varepsilon,\delta,\hat\varepsilon,a,c^{\mathrm{re}},c^{\mathrm{im}}\},\\
\frac{\delta_k}{\delta_n}\le 2\bigl(1+b(k-n)\bigr)^{1/2},\qquad
\frac{p_0(g_k)}{p_0(g_n)}\le 2,\\
\vartheta_{jn}\le C(L,M)(1+bm)^{1/2},\qquad \varpi_{jn}\le(1+bm)^2 .
\end{gathered}
\end{equation}

A \emph{substitution step} in a derivation is a step
\begin{equation*}
\Phi(\bar B)\;\le\;\Phi(L^2-1),
\end{equation*}
justified by the inequality $1+\bar B\le L^2$. Here $\Phi$ is an expression that is non-decreasing in the slots in which $\bar B$ is replaced. A slot is of \emph{type~1} if it stands inside a member of $(\mathrm{R}_{\bar B})$, that is, if it is produced by a two-scale comparison normalised by $x_n$. It is of \emph{type~2} otherwise: then $\bar B$ enters as a bare datum, such as in $|\beta_j|\le\bar B$, as the bare constant $B$, through an un-normalised difference $x_j-x_n$, or as the number inside a constant.

A factor is said to \emph{vanish with the coupling} if it tends to zero as the running coupling $g_j\to0$, equivalently as $x_j\to\infty$. Examples are a positive power of $g_j$, of $\varepsilon$, $\delta$, $\hat\varepsilon$, $a$, $c^{\mathrm{re}}$, $c^{\mathrm{im}}$ or $\alpha$, a negative power $(\log x_j)^{-s}$, and factors exponentially small in the scales $R_j$.

\begin{lemma}[Monotone substitution for two-scale coupling ratios]
\label{4.36:lem-monotone-substitution}
Assume the two-sided flow bound \eqref{4.36:eq-monotone-substitution-1}. With the notation above the following hold.
\begin{enumerate}
\item[(a)] If $0\le b\le b'$, then $(\mathrm{R}_b)$ implies $(\mathrm{R}_{b'})$.
\item[(b)] If $x_j\ge\max\{x_\delta,x_\beta\}$ for all $j\le K$, then $(\mathrm{R}_{\bar B})$ holds.
\item[(c)] If $\gamma\le\gamma_{\bar B}$, then $(\mathrm{R}_1)$ holds, at every $N$ and for every $\bar B\ge0$.
\item[(d)] Let $\gamma\le\gamma_{\bar B}$. Every substitution step all of whose slots are of type~1 yields its conclusion, with its constant as stated, for every blocking factor $L$ (every odd integer $L>1$), every $N$ and every $\bar B\ge 0$. The following four shapes are instances.
\begin{itemize}
\item[(U-a)] The closing of a transport, for $Y$ in the list of \eqref{4.36:eq-monotone-substitution-2} and an exponent $\sigma\in\{1,2\}$:
\begin{equation}\label{4.36:eq-monotone-substitution-3}
Y_nL^{-\sigma m}\le(1+\beta_0)(1+m)^{1/2}L^{-\sigma m}Y_j\le(1+\beta_0)Y_j .
\end{equation}
\item[(U-b)] A closing of shape (U-a) that feeds an estimate carrying a spare factor $(\log x_j)^s$.
\item[(U-c)] A gap supremum $\sup_u w_u(1+u)^cL^{-\sigma u}$ or a gap sum $\sum_u w_u(1+u)^cL^{-\sigma u}$, with $c>0$ and $\sigma>0$, weights $0\le w_u\le1$, and $u$ ranging over the non-negative integers. Such a quantity is bounded by a finite number depending only on $(c,\sigma,L)$, because, for $\theta:=\min\{1,\sigma\log L/(2c)\}$ and all $u\ge0$,
\begin{equation}\label{4.36:eq-monotone-substitution-4}
(1+u)^cL^{-\sigma u}\le\theta^{-c}e^{-(\sigma/2)u\log L}.
\end{equation}
\item[(U-d)] A pure ratio bounded by a value that does not contain $B$.
\end{itemize}
\item[(e)] Suppose the ceiling $\gamma\le\gamma_{\bar B}$ is not imposed, but $x_j\ge x^{(3)}_{\mathrm{th}}$ for all $j\le K$, so that in particular $x_j\ge\max\{x_\delta,x_\beta\}$, the floors $x_\delta,x_\beta$ being members of the threshold $x^{(3)}_{\mathrm{th}}$ of \eqref{4.36:eq-coupling-ceiling-a}. Then every conclusion of shape (U-a) holds with its constant multiplied by $(1+\bar B)^{1/2}$. This factor is absorbed by a one-sided floor on $x$ wherever the estimate consuming the conclusion carries a spare power of $\log x_j$. The conditions $q:=\min\{q_0,q_1\}\ge p_0+3r+2$, $p_0\ge5r$ and $r\ge2$ on the exponents $q_0,q_1$ of the radii $\alpha_{0,j},\alpha_{1,j}$, the exponent $p_0$ of the degree function and the exponent $r$ of the scales $R_j$ supply such a power for every $\alpha_{\cdot,j}$ and $\alpha_{3,j}$ set against $\varepsilon_j$, $\delta_j$, $\hat\varepsilon_j$. The spare exponents are $q-p_0$, $q-p_0-r$, $q_1-p_0$ and $q_0-p_0-r$, each at least $2r+2$.
\item[(f)] The following two kinds of type-2 occurrence are discharged without the inequality $1+\bar B\le L^2$.
\begin{itemize}
\item An occurrence with a factor vanishing with the coupling on its side, set against a target free of $g$, is a floor on $x$ whose constant depends on $(N,L,\mathfrak{p})$ through $\bar B$.
\item An occurrence inside a constant, or inside a convergent gap sum consumed only as a value, is a value and is never compared with a power of $L$.
\end{itemize}
\item[(g)] Consequently, a substitution step that is not covered by (d)--(f) is of the following form.
\begin{itemize}
\item[(U-e)] A type-2 slot (bare $\bar B$, $|\beta_j|$, or an un-normalised difference) is set against a structural target free of $g$ (a power of $L$ or a numeral). No factor vanishing with the coupling stands on its side, either at the line or where the line is consumed. The step does not merely value a constant or assert a finiteness. Schematically, for a generic non-negative quantity, denoted here by $\mathcal{Q}$:
\begin{equation*}
|\beta_j|\le\bar B\le L^2-1,\quad\text{hence}\quad \bar BL^{-2}\mathcal{Q}\le\mathcal{Q},
\end{equation*}
with this inequality consumed as such.
\end{itemize}
\end{enumerate}
\end{lemma}

\begin{proof}
\emph{(a).} We compare the right-hand sides of \eqref{4.36:eq-monotone-substitution-2} one by one.
\begin{itemize}
\item The members $1+bm$, $(1+bm)^{1/2}$, $\log(1+bm)$, $(1+\beta_0)(1+bm)^{1/2}$, $2(1+b(k-n))^{1/2}$, $C(L,M)(1+bm)^{1/2}$ and $(1+bm)^2$ are non-decreasing in $b$, since $m\ge0$ and $k-n\ge0$.
\item The member $2$ for $p_0(g_k)/p_0(g_n)$ does not depend on $b$.
\end{itemize}
Hence each member of $(\mathrm{R}_b)$ implies the corresponding member of $(\mathrm{R}_{b'})$ for $b\le b'$.

\emph{(b).} Let $j\le n\le k\le K$ and $m=n-j$. By the upper half of \eqref{4.36:eq-monotone-substitution-1} and $B\le\bar B$, we have $x_j-x_n\le\bar Bm$. Since $x_n\ge1$,
\begin{equation*}
x_j\le x_n+\bar Bm\le x_n+\bar Bm\,x_n=(1+\bar Bm)x_n .
\end{equation*}
This is the first member of $(\mathrm{R}_{\bar B})$. The second and third members are equivalent to it, because $g_n/g_j=(x_j/x_n)^{1/2}$ and $\log$ is increasing.

The members for $Y\in\{\varepsilon,\delta,\hat\varepsilon,a,c^{\mathrm{re}},c^{\mathrm{im}}\}$ follow from part~(ii) of the two-scale comparison of the logarithmic factors. That comparison applies because $x_j,x_n\ge\max\{x_\delta,x_\beta\}$, and it is used together with the inequality $x_n\le x_j+c_5$ displayed in \eqref{4.36:eq-sun-conventions-b}. For $Y=\hat\varepsilon$ one uses in addition property~(a) of the scales $R_j$.

The members further down the list, those for $\delta_k/\delta_n$, $p_0(g_k)/p_0(g_n)$ and $\vartheta_{jn}$, follow from part~(i) of the same comparison; for $\vartheta_{jn}$ the member is the form in which that ratio is bounded where the localised chart is constructed. For $\delta_k/\delta_n$ this is the lower bound on $\delta_n$ in terms of $\delta_k$ and $(1+B(k-n))^{-1/2}$ displayed in \eqref{4.36:eq-sun-conventions-b}. That bound holds under $x_k\ge x_\delta$ and is combined with $B\le\bar B$.

Finally, $1\le(1+\bar Bm)^2$ and $(x_j/x_n)^2\le(1+\bar Bm)^2$ give $\varpi_{jn}\le(1+\bar Bm)^2$.

\emph{(c).} Let $\gamma\le\gamma_{\bar B}$. By the lower half of \eqref{4.36:eq-monotone-substitution-1} and $g_K\le\gamma$, we have, for every $n\le K$,
\begin{equation*}
x_n\ge x_K+\lambda(K-n)-c_5\ge\gamma^{-2}-c_5\ge\max\{\bar B,x_\delta,x_\beta\},
\end{equation*}
as displayed in \eqref{4.36:eq-sun-conventions-b}. In particular, $x_n\ge\bar B$ gives
\begin{equation*}
x_j-x_n\le\bar Bm\le x_nm,\qquad\text{that is,}\qquad x_j\le(1+m)x_n .
\end{equation*}
Since also $x_n\ge\max\{x_\delta,x_\beta\}$ for every $n$, the members of $(\mathrm{R}_1)$ other than the one for $\delta_k/\delta_n$ follow from this inequality word for word as in the proof of (b). The member for $\delta_k/\delta_n$ is the bound $\delta_n\ge\frac12(1+(k-n))^{-1/2}\delta_k$ displayed in \eqref{4.36:eq-sun-conventions-b}, which holds under the present ceiling. This is the form of the coupling ratios and transport shapes displayed in \eqref{4.36:eq-sun-conventions-b}. No relation between $\bar B$ and $N$ or $L$ was used.

\emph{(d).} Let $U\le\Phi(L^2-1)$ be the conclusion of a substitution step all of whose slots are of type~1, and let $\gamma\le\gamma_{\bar B}$.

The lines of the derivation before the step use $\bar B$ only through members of $(\mathrm{R}_{\bar B})$. Read with $(\mathrm{R}_b)$ in place of $(\mathrm{R}_{\bar B})$, the same lines prove $U\le\Phi(b)$ for every $b\ge1$ for which $(\mathrm{R}_b)$ holds. By (c), $(\mathrm{R}_1)$ holds, and so $U\le\Phi(1)$.

Since $L$ is an integer exceeding $1$, we have $1\le L^2-1$. As $\Phi$ is non-decreasing in the slots concerned, $\Phi(1)\le\Phi(L^2-1)$. The lines after the step do not involve $b$. Hence the conclusion holds with its constant as stated, for every $L$, $N$ and $\bar B\ge0$.

For (U-a) the expression is $\Phi(b)=(1+\beta_0)(1+bm)^{1/2}L^{-\sigma m}Y_j$. The first inequality of \eqref{4.36:eq-monotone-substitution-3} is the $Y$-member of $(\mathrm{R}_1)$, multiplied by $L^{-\sigma m}Y_j$. The second is part~(v) of the ratio lemma for the sizes $Y_n/Y_j$; it also follows directly, because $1+m\le2^m\le L^{2\sigma m}$ for $m\ge0$, $\sigma\ge1$ and $L\ge2$.

(U-b) is (U-a) followed by lines that do not involve $b$. (U-d) is the case in which $\Phi$ is itself a ratio bound.

For (U-c), note first that $0<\theta\le1$, so that $\theta(1+u)\le1+\theta u$. Hence
\begin{equation*}
(1+u)^c\le\theta^{-c}(1+\theta u)^c\le\theta^{-c}e^{c\theta u}\le\theta^{-c}e^{(\sigma/2)u\log L},
\end{equation*}
the last step because $c\theta\le \sigma\log L/2$. Multiplying by $L^{-\sigma u}=e^{-\sigma u\log L}$ gives \eqref{4.36:eq-monotone-substitution-4}. The right-hand side of \eqref{4.36:eq-monotone-substitution-4} is bounded in $u$ and summable over the non-negative integers $u$, because $\sigma>0$ and $L>1$. Since $0\le w_u\le1$, the gap supremum and the gap sum are bounded by the supremum and the sum of that right-hand side. So they are bounded by finite functions of $(c,\sigma,L)$.

\emph{(e).} Without the ceiling, (b) gives $(\mathrm{R}_{\bar B})$ under the floors $x_\delta,x_\beta$ included in the threshold. Parts~(iii) and~(iv) of the ratio lemma for the sizes then give each closing of shape (U-a) with $(1+\bar Bm)^{1/2}$ in place of $(1+m)^{1/2}$. Since $1+\bar Bm\le(1+\bar B)(1+m)$ for $m\ge0$, each such closing holds with its constant multiplied by $(1+\bar B)^{1/2}$.

Where the consuming estimate carries a spare power $(\log x_j)^s$, the factor is absorbed by requiring $(\log x_j)^s$ to exceed $(1+\bar B)^{1/2}$ or $(1+\bar B)$ times the constant concerned, as the floors displayed in Remark~\ref{4.36:rem-spare-logarithms} specify for each estimate. This is a one-sided floor on $x_j$. The conditions on the exponents give the spare exponents listed in part~(e) above, each at least $2r+2$, for every $\alpha_{\cdot,j}$ and $\alpha_{3,j}$ set against $\varepsilon_j$, $\delta_j$, $\hat\varepsilon_j$.

The resulting floors are displayed in Remark~\ref{4.36:rem-spare-logarithms}.

\emph{(f).} Consider first a type-2 occurrence with a factor vanishing with the coupling on its side, set against a target free of $g$. The inequality holds as soon as that factor is small enough compared with the target divided by the remaining, $g$-free, part of its side. That part depends on $\bar B$, hence on $(N,L,\mathfrak{p})$. Since the factor tends to zero as $x\to\infty$, there is a value of $x$ beyond which it stays below that bound. The condition is therefore a floor on $x$ with a constant depending on $\bar B(N,L,\mathfrak{p})$, and the inequality $1+\bar B\le L^2$ is not used.

Consider next a type-2 occurrence inside a constant, or inside a convergent gap sum consumed only as a value. It enters the derivation only through the finiteness and size of that value. It is never compared with a power of $L$, so again no substitution step is made.

\emph{(g).} A step all of whose slots are of type~1 is covered by (d) under the ceiling $\gamma\le\gamma_{\bar B}$ and by (e) without it. So a substitution step not covered by (d)--(f) has a slot of type~2. If a factor vanishing with the coupling stands on the side of that slot, either at the line or where the line is consumed, or if the slot sits inside a constant or a convergent gap sum consumed as a value, then (f) applies.

Otherwise the slot is a bare $\bar B$, a $|\beta_j|$, or an un-normalised difference. It is set against a target free of $g$, which is a power of $L$ or a numeral, and it does more than value a constant or assert a finiteness. This is the form (U-e).
\end{proof}

\noindent\textbf{Inputs used by statement.}
The proofs of this chapter establish clauses (0) and (i) of Theorem~0, together with the constants
for $SU(N)$. The inputs of the statements concerning $SU(N)$ are entered under (0), and under (i)
where the proof of clause (i) uses them. For each clause, in the order of Theorem~0, two lists
follow. The first names the statements of Ba{\l}aban's papers that these proofs use by statement,
with exact references. The second names the statements of this book used in these proofs whose own
proof is incomplete or outstanding, each marked ``stated; proof incomplete'' or ``proof
outstanding''. Ba{\l}aban's definitions and displayed formulas are cited where they are used and are
not repeated here.

\begin{description}
\item[(0)] \textbf{Published statements cited by statement.}
\begin{itemize}
\item \cite[(1.118)]{B-CMP95}; \cite[p.~229, (2.36)]{B-CMP96}, the bump profiles of the covers and
their partition-of-unity identity.
\item \cite[Lemma~2.1]{B-CMP96}; \cite[Lemma~2.4]{B-CMP96}; \cite[Propositions~2.6--2.7]{B-CMP96}.
\item \cite[Propositions~1--7]{B-CMP98}.
\item \cite[Lemma~1]{B-CMP99a}; \cite[Theorems~2 and~4]{B-CMP99a};
\cite[Proposition~6, pp.~98--99]{B-CMP99a}.
\item \cite[Theorems~3.1--3.4, 3.7, 3.9--3.15]{B-CMP99b};
\cite[Corollaries~3.5, 3.6, 3.8]{B-CMP99b}.
\item \cite[Theorem~1, (7)--(9), p.~279]{B-CMP102b}; \cite[Propositions~3 and~9]{B-CMP102b}.
\item \cite[Theorems~1 and~3]{B-CMP109}; \cite[Lemma~4]{B-CMP109};
\cite[Lemmas~5 and~6]{B-CMP109}.
\item \cite[Lemmas~1, 2 and~3]{B-CMP116}.
\item \cite[Proposition~1]{B-CMP122a}.
\item \cite[Theorem~1, pp.~355--356]{B-CMP122b}; \cite[Proposition~1]{B-CMP122b}.
\item \cite[pp.~178--179]{B-CMP122a}, admissibility of the prepared large-field data ---
printed without proof; cited as printed.
\item \cite[p.~276]{B-CMP119}, the analysis of \cite[\S3]{B-CMP109}, with
\cite[Lemmas~5 and~6]{B-CMP109}, repeated on the cover of \cite[\S3]{B-CMP119} --- printed
without proof; cited as printed.
\item \cite[p.~396, (3.35), (3.37)]{B-CMP99b}, the regularity conditions at the two lowest tower
levels --- printed without proof; cited as printed.
\item \cite[p.~379]{B-CMP122b}, locality of the interaction terms --- printed without proof;
cited as printed.
\item \cite[(1.17)]{B-CMP109}, the inclusion sentence, read at complex pairs --- printed without
proof; cited as printed.
\item \cite[p.~277]{B-CMP109}, the sentence extending the bounds
\cite[(3.26)--(3.28)]{B-CMP109} on block-averaged curvature to near data --- printed without
proof; cited as printed.
\item \cite[\S3, the sentence preceding (3.132)]{B-CMP99b}, the bounds and walk expansion of the
inverse of the constrained covariance --- printed without proof; cited as printed.
\item \cite[\S3, after Theorem~3.9]{B-CMP99b}, the walk expansions of the further operators of
that section --- printed without proof; cited as printed.
\item \cite[\S E, p.~43]{B-CMP98}, analyticity of the derivative bounds for the averaging maps
--- printed without proof; cited as printed.
\item \cite[(3.49), (3.69)]{B-CMP99b}, the second inequality of the current bound for the far
family --- printed without proof; cited as printed.
\item \cite[(1.90)]{B-CMP122b}, the polymer-gas representation of the bracket --- printed
without proof; cited as printed.
\item \cite[p.~386]{B-CMP122b}, the bounds of the large-field terms under the remaining
modifications --- printed without proof; cited as printed.
\item \cite[(1.64)--(1.69), pp.~190--191]{B-CMP122a}, the layerwise gauge clauses of the
large-field spaces --- printed without proof; cited as printed.
\item \cite[p.~192]{B-CMP122a}, the layerwise ceiling as a condition of small coupling ---
printed without proof; cited as printed.
\item \cite[p.~375]{B-CMP122b}, the bounds for old terms beside a new large-field region ---
printed without proof; cited as printed.
\item \cite[p.~192, (1.70)]{B-CMP122a} and \cite[p.~377, (1.68)--(1.69)]{B-CMP122b}, analytic
extension and bounds of the stage terms and of the localized interaction terms on the joined
determining sets --- printed without proof; cited as printed.
\end{itemize}
\textbf{Statements of this book whose proofs are incomplete or outstanding.}
\begin{itemize}
\item Lemma~\ref{4.2:lem-unmarked} --- stated; proof incomplete.
\item The re-blocking step inside admissible boxes --- stated; proof incomplete.
\item Real-factor slot bounds inside admissible boxes --- stated; proof incomplete.
\item The re-blocking step outside admissible boxes --- proof outstanding.
\item Uniqueness and locality of the constrained minimiser --- proof outstanding.
\item The cube gauge condition for the small pairs and the real-background pairs --- proof
outstanding.
\item Positivity dichotomy of the large-field operations --- proof outstanding.
\end{itemize}
\item[(i)] \textbf{Published statements cited by statement.}
\begin{itemize}
\item \cite[Lemma~2.1]{B-CMP96}.
\item \cite[Propositions~2, 6 and~7]{B-CMP98}.
\item \cite[Theorem~2]{B-CMP99a}; \cite[(1.19), p.~79]{B-CMP99a}.
\item \cite[Theorems~3.1--3.3, 3.7, 3.9--3.11 and~3.13--3.15]{B-CMP99b};
\cite[Corollary~3.8]{B-CMP99b}.
\item \cite[Theorem~1]{B-CMP102b}; \cite[Proposition~9]{B-CMP102b}.
\item \cite[Lemma~4]{B-CMP109}; \cite[(3.37)--(3.50)]{B-CMP109}.
\item \cite[Lemmas~1 and~2]{B-CMP116}.
\item \cite[Theorems~1 and~2]{B-CMP119}; \cite[Corollary~3]{B-CMP119}.
\item \cite[Theorem~1]{B-CMP122b}; \cite[Proposition~1]{B-CMP122b}.
\item \cite[p.~276]{B-CMP119}, the analysis of \cite[\S3]{B-CMP109}, with
\cite[Lemmas~5 and~6]{B-CMP109}, repeated on the cover of \cite[\S3]{B-CMP119} --- printed
without proof; cited as printed.
\item \cite[(1.11), (1.23)--(1.48), p.~369]{B-CMP122b}, the absolute bound and gauge invariance
of the boundary terms --- printed without proof; cited as printed.
\item \cite[p.~367, (1.40)--(1.41)]{B-CMP122b}, the analytic solvability of the equation
displayed there --- printed without proof; cited as printed.
\item \cite[p.~375]{B-CMP122b}, the bounds for old terms beside a new large-field region ---
printed without proof; cited as printed.
\item \cite[pp.~178--179]{B-CMP122a}, admissibility of the prepared large-field data ---
printed without proof; cited as printed.
\end{itemize}
\textbf{Statements of this book whose proofs are incomplete or outstanding.}
\begin{itemize}
\item The volume counts displayed as a hypothesis of the relative small-field machine --- proof
outstanding.
\item The two bounds on the data read by the lemmas in the proof of the correlation theorem,
displayed as hypotheses of the relative small-field machine --- proof outstanding.
\item The kernel-difference bound displayed as a hypothesis of the relative small-field machine
--- proof outstanding.
\item Geometry and counting of live large-field regions (regions still carrying their
large-field factor at the level read) --- proof outstanding.
\item A kept term as a functional of the data: its label, norm and law sites --- proof
outstanding.
\item Kept terms add one summand to the probe functional --- proof outstanding.
\item Contribution of the kept terms to the stage totals --- stated; proof incomplete.
\item Kept terms at a large-field operation on a foreign component --- proof outstanding.
\item Combes--Thomas decay of the block covariance --- proof outstanding.
\item Conversion of per-rate bounds to the stage norm --- proof outstanding.
\item Field dependence of the densities in cube gauges --- stated; proof incomplete.
\item Holonomy bound for walks across several cubes --- proof outstanding.
\item Oscillation of a Gaussian piece under small moves --- stated; proof incomplete.
\item Dissolution of the arrested region --- stated; proof incomplete.
\item The $\mathbf{T}$-step layer in a background field --- stated; proof incomplete.
\item Weights, the dated size parameter and heredity --- stated; proof incomplete.
\item Cube families under a stage move --- stated; proof incomplete.
\item Current thresholds across an increase of density --- stated; proof incomplete.
\item The one-point current probe --- stated; proof incomplete.
\item No index-zero margin on source domains --- stated; proof incomplete.
\item Inclusion of the ambient domain at size $\gamma_0^{(k)}$, the dated radius of the arrest
theorem --- stated; proof incomplete.
\item The arrest theorem --- stated; proof incomplete.
\item Newborn terms, space inclusions and continuation of stored domains --- stated; proof
incomplete.
\item Bounds of later steps on live arrested pieces --- stated; proof incomplete.
\item The current radius read at raw points of live pieces --- stated; proof incomplete.
\item Index-zero members under maps, at the real point and at completion --- stated; proof
incomplete.
\item The core-freezing step --- stated; proof incomplete.
\item No lattice point on both boundaries of a stratum --- stated; proof incomplete.
\item The cell-counting bound on the stratified frames of the induction --- stated; proof
incomplete.
\item The stability bound on the near part of the effective action at $N\ge 3$, displayed among
the standing hypotheses of the section on the constants for $SU(N)$ --- stated; proof
incomplete.
\end{itemize}
\item[(iii)] None. The proofs of this chapter concern clauses (0) and (i) only; under the
remaining clauses there is nothing to list.
\item[(ii)] None.
\item[(iv)] None.
\item[(v)] None.
\item[(vi)] None.
\end{description}

The proofs of clauses (0) and (i) form one induction on the level $k$, $0\le k\le K$. At the step
from level $k$ to level $k+1$, each statement of the second lists is applied only to couplings,
effective densities and large-field data of levels at most $k$. Within one step, the sites at
which these statements are applied are taken in a well-founded order. The auxiliary parameters
$\mathfrak{p}$ are chosen in one fixed order after $L$; the constants of
\cite[Theorem~3.1, p.~397]{B-CMP99b} depend on $d$ and $L$ only, and the conditions of
\cite[(162)--(163), p.~304]{B-CMP102b} are one-sided conditions on parameters chosen after $L$.
Hence no statement listed above is used, directly or through another, in its own proof.

%% ---- end inlined file: chapters/c04 ----
